% % % % % % % % % % % % % % 
%!TEX encoding = UTF-8 Unicode
%	U.S. Chartbook
%	Brian W. Dew (brianwdew@gmail.com)
%	Updated: December 15, 2019
%	GitHub repo contains to do list (issues)
%   https://github.com/bdecon/US-chartbook
%
% % % % % % % % % % % % % %
\PassOptionsToPackage{table}{xcolor}
\documentclass{report}

%
% % % % % % Packages % % % % % % % % % 
%
	\usepackage[letterpaper, margin=1.18in]{geometry}
	\usepackage{microtype}
	\usepackage[default]{lato}
	\usepackage{pgfplots, pgfplotstable}
	\usepackage[eulergreek]{sansmath}
	\usepackage{xcolor}
		\definecolor{usaablue}{RGB}{18,57,91}

	\usepackage{array}
	\usepackage{fontawesome5}
	\usepackage{titlesec}
	\usepackage{fancyhdr}
	\usepackage[colorlinks=TRUE, allcolors=usaablue,
		pdfencoding=auto]{hyperref}
	\usetikzlibrary{pgfplots.dateplot, pgfplots.fillbetween, patterns,
	                pgfplots.groupplots, shapes.geometric}

%
% % % % % Document Settings % % % % % % % 
%

	% Paragraph spacing
	\usepackage{parskip}
	\setlength\parindent{0pt}
	\setlength{\parskip}{8pt}
	\makeatletter
		\newcommand{\@minipagerestore}{\setlength{\parskip}{8pt}}
	\makeatother
	
	% Section and Subsection Headings
	\titleformat{\section}
  		{\color{usaablue} \LARGE \seriffont \bfseries}
  		{\thesection}{1em}{}
	\titleformat{\subsection}
  		{\color{usaablue} \seriffont \bfseries \large}
  		{\thesection}{1em}{}
	\titleformat{\subsubsection}
  		{\color{usaablue} \seriffont \bfseries \normalsize}
  		{\thesection}{1em}{}		

% % % % % Graph Settings % % % % % % % 
%
	
	% Header and footer
	
	\pagestyle{fancy}
	\fancyhf{}
	\renewcommand{\headrulewidth}{0pt}
	\rfoot{\color{usaablue}\thepage}
	\cfoot{\color{usaablue}CONFIDENTIAL}

	
	
	% Color square
	\newcommand{\cbox}[1]{
		\begin{tikzpicture} \draw [#1, line width=6](0,0) -- (.2,0);  
		\end{tikzpicture}}
	\newcommand{\colorline}[2]{
		\begin{tikzpicture} \draw [#1, line width=1.8](0,0.2) -- +(0.6,0) node[right, black!80] {#2}; 
		\end{tikzpicture}}
		
	% Table link
	\newcommand{\tbllink}[1]{\href{https://raw.githubusercontent.com/bdecon/US-chartbook/master/chartbook/data/#1}{\faTable}}
	
	% Last two digits of year
	\makeatletter
	\newcommand*\short[1]{\expandafter\@gobbletwo\number\numexpr#1\relax}
	\makeatother	
	
	% Column width and alignment
	\newcolumntype{R}[1]{>{\raggedleft\let\newline\\\arraybackslash\hspace{0pt}}m{#1}}	
	\newcolumntype{C}[1]{>{\centering\let\newline\\\arraybackslash\hspace{0pt}}m{#1}}
	
	% Style for date plots
	\pgfplotsset{compat=newest, 
		scaled y ticks=false,
		axis line style={black!20}, 
		xtick style={black!20}, ytick style={draw=none},
		every tick label/.style={black!50, font=\scriptsize,
			/pgf/number format/assume math mode=true},
		width=12.8cm, height=4.8cm, 
		xticklabel style={align=left}, 
		yticklabel style={text width=0.9em, align=right},       
		axis x line*=bottom, x axis line style={black!50},
	    axis y line=left, y axis line style={opacity=0},
	    ymajorgrids, grid style={very thin, black!10},	        
	    every node near coord/.style={/pgf/number format/fixed,
	    	font=\scriptsize, style={black!70}},
	    legend style={legend columns=-1, draw=none, fill=none,
	    	/tikz/every even column/.append style={column sep=0.3cm}}}
	    	
	
	% stacked diverging bar
	\newcommand{\sbar}[4]{
		\addplot[ybar stacked, bar width=2.45pt, draw opacity=0, fill=#1] 
			table [x=#2, y=#3, col sep=comma]{#4};}

	% stacked diverging bar area legend
	\newcommand{\sbaral}[4]{
		\addplot[ybar stacked, bar width=2.45pt, draw opacity=0, fill=#1, area legend] 
			table [x=#2, y=#3, col sep=comma]{#4};}
			
	% thin stacked diverging bar
	\newcommand{\tsbar}[4]{
		\addplot[ybar stacked, bar width=2.2pt, draw opacity=0, fill=#1] 
			table [x=#2, y=#3, col sep=comma]{#4};}
			
	% custom width stacked diverging bar
	\newcommand{\ctsbar}[5]{
		\addplot[ybar stacked, bar width=#5, draw opacity=0, fill=#1] 
			table [x=#2, y=#3, col sep=comma]{#4};}
			
	% area plot segment
	\newcommand{\abar}[4]{
		\addplot[stack plots=y, area style, draw=none, fill=#1] 
			table [x=#2, y=#3, col sep=comma]{#4}\closedcycle;}
					
	% text node
	\newcommand{\stdnode}[3]{\node[below, align=left, shift=({#1,#2})]{#3};}	
	
	% text node located by data	 
	\newcommand{\absnode}[3]{\node[below right, align=left] at (axis cs: #1,#2) {#3};}   
	
	% multiline text node located by data	 
	\newcommand{\absnodeml}[4]{\node[below right, align=left, text width=#4cm] 
		at (axis cs: #1,#2) {#3};}       
		        
	% Date (X) Axis Tick Marks, one tick per year, every even year labeled
	\newcommand{\dateaxisticks}{
		date coordinates in=x, axis line style={draw=none},
		xmax={2022-03-31},
		max space between ticks=40,	    
		xtick={{1990-01-01}, {1992-01-01}, {1994-01-01}, 
			{1996-01-01}, {1998-01-01}, {2000-01-01}, 
			{2002-01-01}, {2004-01-01}, {2006-01-01},
			{2008-01-01}, {2010-01-01}, {2012-01-01}, {2014-01-01},
		    {2016-01-01}, {2018-01-01}, {2020-01-01}, {2022-01-01}, 
		    {2024-01-01}, {2026-01-01}},
		minor xtick={{1989-01-01}, {1991-01-01}, {1993-01-01},
			{1995-01-01}, {1997-01-01}, {1999-01-01}, 
			{2001-01-01}, {2003-01-01}, {2005-01-01}, {2007-01-01},
		    {2009-01-01}, {2011-01-01}, {2013-01-01}, {2015-01-01},
		    {2017-01-01}, {2019-01-01}, {2021-01-01}, {2023-01-01}, 
		    {2025-01-01}, {2027-01-01}},
		enlarge y limits={0.06}, enlarge x limits={0.01},
		}
		
	% Date (X) Axis Tick Marks, one tick per year, every even year labeled
	\newcommand{\shdateaxisticks}{
		date coordinates in=x, axis line style={draw=none},
		xmax={2022-03-31},
		max space between ticks=40,	    
		xtick={{1990-01-01}, {1995-01-01}, {2000-01-01}, 
			{2005-01-01}, {2010-01-01}, {2015-01-01}, {2020-01-01}},
		minor xtick={},
		enlarge y limits={0.06}, enlarge x limits={0.01},
		}
		
	% Date (X) Axis Tick Marks, one tick per year, every even year labeled
	\newcommand{\ltdateaxisticks}{
		date coordinates in=x, axis line style={draw=none},
		xmax={2022-03-31},
		max space between ticks=40,	    
		xtick={{2013-01-01}, {2014-01-01}, {2015-01-01}, {2016-01-01}, {2017-01-01}, {2018-01-01}, 
		    {2019-01-01}, {2020-01-01}, {2021-01-01}, {2022-01-01}},
		enlarge y limits={0.06}, enlarge x limits={0.01},
		}
		
	% Date (X) Axis Tick Marks, one tick per year, every even year labeled
	\newcommand{\lfdateaxisticks}{
		date coordinates in=x, axis line style={draw=none},
		xmin={2018-01-01}, xmax={2022-03-31},   
		xtick={{2018-01-01}, {2019-01-01}, {2020-01-01}, {2021-01-01}, {2022-01-01}},
		enlarge y limits={value=0.12, upper}, enlarge x limits={0.02}, ymin=0,
		yticklabel style={text width=1.0em},
		height=3.8cm, width=6.4cm,
		}
		
	% Date (X) Axis Tick Marks, one tick per year, every even year labeled
	\newcommand{\tydateaxisticks}{
		date coordinates in=x, axis line style={draw=none},
		xmax={2022-03-31}, max space between ticks=40,	    
		xtick={{2011-01-01}, {2012-01-01}, {2013-01-01}, {2014-01-01}, {2015-01-01}, {2016-01-01}, 
			{2017-01-01}, {2018-01-01}, {2019-01-01}, {2020-01-01}, {2021-01-01}, {2022-01-01}},
		enlarge y limits={0.06}, enlarge x limits={0.01},
		}
		
	% Settings for y label text in horizontal bar charts
	\newcommand{\barylab}[2]{yticklabel style={text width=#1, align=right, 
		style={black!70}, text height=#2},}
	
	% Solid bars at significant  x or y values
	\newcommand{\bbar}[2]{extra #1 ticks = {{#2}}, extra #1 tick labels = ,
		extra #1 tick style = {grid=major, grid style={thick, black!25}},}
		
			% Solid bars at significant  x or y values
	\newcommand{\dbar}[2]{extra #1 ticks = {{#2}}, extra #1 tick labels = ,
		extra #1 tick style = {grid=major, grid style={dashed, thick, black!50}},}
		
	% Standard line
	\newcommand{\stdline}[4]{\addplot[very thick, no markers, color=#1] 
		table [x=#2, y=#3, col sep=comma] {#4};	}
		
	% Thin line
	\newcommand{\thinline}[4]{\addplot[no markers, color=#1] 
		table [x=#2, y=#3, col sep=comma] {#4};	}
		
	% Dashed line
	\newcommand{\dashline}[4]{\addplot[very thick, dashed, no markers, color=#1] 
		table [x=#2, y=#3, col sep=comma] {#4};	}
		
	% Thicker line
	\newcommand{\thickline}[4]{\addplot[ultra thick, no markers, color=#1] 
		table [x=#2, y=#3, col sep=comma] {#4};	}
		
	% Style for bar plots legend symbol		
		\pgfplotsset{/pgfplots/area legend/.style={/pgfplots/legend image code/.code={
            \fill[##1] (0cm, -0.1cm) rectangle (0.6cm, 0.1cm);}},}		
		
	% Additional bar plot settings
	\newcommand{\barplotnogrid}{xbar=0pt, axis line style={draw=none},
	    yticklabel style={align=left, anchor=east},
      		xmajorticks=false, ymajorgrids=false,   
	    ytick=data, tickwidth=0pt, area legend, reverse legend,
	    nodes near coords align={horizontal},}  
		
	% Recession bars		
	\newcommand{\rbars}{
		\fill[color=black!10] (axis cs:{1990-07-01},\pgfkeysvalueof{/pgfplots/ymin}) rectangle 
			(axis cs:{1991-03-01}, \pgfkeysvalueof{/pgfplots/ymax});
		\fill[color=black!10] (axis cs:{2007-12-01},\pgfkeysvalueof{/pgfplots/ymin}) rectangle 
			(axis cs:{2009-07-01}, \pgfkeysvalueof{/pgfplots/ymax});
		\fill[color=black!10] (axis cs:{2001-03-01},\pgfkeysvalueof{/pgfplots/ymin}) rectangle 
			(axis cs:{2001-11-01}, \pgfkeysvalueof{/pgfplots/ymax});
		\fill[color=black!10] (axis cs:{2020-02-01},\pgfkeysvalueof{/pgfplots/ymin}) rectangle 
			(axis cs:{2020-05-01}, \pgfkeysvalueof{/pgfplots/ymax});}
			
	\newcommand{\rebars}{
		\fill[color=black!10] (axis cs:{2007-12-01},\pgfkeysvalueof{/pgfplots/ymin}) rectangle 
			(axis cs:{2009-07-01}, \pgfkeysvalueof{/pgfplots/ymax});
		\fill[color=black!10] (axis cs:{2001-03-01},\pgfkeysvalueof{/pgfplots/ymin}) rectangle 
			(axis cs:{2001-11-01}, \pgfkeysvalueof{/pgfplots/ymax});
		\fill[color=black!10] (axis cs:{2020-02-01},\pgfkeysvalueof{/pgfplots/ymin}) rectangle 
			(axis cs:{2020-05-01}, \pgfkeysvalueof{/pgfplots/ymax});}
			
	\newcommand{\recbars}{
		\fill[color=black!10] (axis cs:{2007-12-01},\pgfkeysvalueof{/pgfplots/ymin}) rectangle 
			(axis cs:{2009-07-01}, \pgfkeysvalueof{/pgfplots/ymax});
		\fill[color=black!10] (axis cs:{2020-02-01},\pgfkeysvalueof{/pgfplots/ymin}) rectangle 
			(axis cs:{2020-05-01}, \pgfkeysvalueof{/pgfplots/ymax});}
			
	\newcommand{\rbar}{
		\fill[color=black!10] (axis cs:{2020-02-01},\pgfkeysvalueof{/pgfplots/ymin}) rectangle 
			(axis cs:{2020-05-01}, \pgfkeysvalueof{/pgfplots/ymax});}
	
	\newfontfamily\seriffont{RobotoSlab}	
	
	\pgfplotstableread[header=true, col sep=comma]{data/cpi_comp.csv}\cpi
	\pgfplotstableread[header=true, col sep=semicolon]{data/ip_comp.csv}\ip
	\pgfplotstableread[header=true, col sep=comma]{data/rs_comp.csv}\rs
	\pgfplotstableread[header=true, col sep=comma]{data/ahe_ind.csv}\ahe
	\pgfplotstableread[header=true, col sep=comma]{data/poor.csv}\poor
	\pgfplotstableread[header=true, col sep=comma]{data/poor2.csv}\pvrt
	\pgfplotstableread[header=true, col sep=comma]{data/spmtbl20.csv}\spm
	\pgfplotstableread[header=true, col sep=semicolon]{data/occs.csv}\occ
	\pgfplotstableread[header=true, col sep=comma]{data/empgroups.csv}\emp
	\pgfplotstableread[header=true, col sep=comma]{data/empgroups2.csv}\empt
	\pgfplotstableread[header=true, col sep=comma]{data/unempgroups.csv}\unemp
	\pgfplotstableread[header=true, col sep=comma]{data/unempgroups2.csv}\unempt
	\pgfplotstableread[header=true, col sep=comma]{data/unempgroups3.csv}\unemptt
	\pgfplotstableread[header=true, col sep=semicolon]{data/cps_educ.csv}\edsh
	\pgfplotstableread[header=true, col sep=comma]{data/cps_educ_tot.csv}\edtot
	\pgfplotstableread[header=true, col sep=comma]{data/cps_age.csv}\agesh
	\pgfplotstableread[header=true, col sep=semicolon]{data/union_ind.csv}\unmem
	\pgfplotstableread[header=true, col sep=semicolon]{data/quits_ind.csv}\quits
	\pgfplotstableread[header=true, col sep=semicolon]{data/state_pa_epop.csv}\paepop
	\pgfplotstableread[header=true, col sep=semicolon]{data/state_pa_epop2.csv}\paepopt
	\pgfplotstableread[header=true, col sep=semicolon]{data/state_pa_epop3.csv}\paepoptt
	\pgfplotstableread[header=true, col sep=semicolon]{data/openings_ind.csv}\opens
	\pgfplotstableread[header=true, col sep=comma]{data/nilf_comp.csv}\nilf
	\pgfplotstableread[header=true, col sep=comma]{data/pinc.csv}\pinc
	\pgfplotstableread[header=true, col sep=comma]{data/unemp_grp.csv}\ungrp
	\pgfplotstableread[header=true, col sep=comma]{data/unemp_grpsh.csv}\ungrpsh
	\pgfplotstableread[header=true, col sep=comma]{data/ce_age.csv}\ceage
	\pgfplotstableread[header=true, col sep=comma]{data/ce_inc.csv}\ceinc
	\pgfplotstableread[header=true, col sep=comma]{data/cpi_monthly.csv}\cpimo
	\pgfplotstableread[header=true, col sep=comma]{data/ccdebtbar.csv}\ccbar
	\pgfplotstableread[header=true, col sep=comma]{data/gdp_rec.csv}\gdprec
	\pgfplotstableread[header=true, col sep=comma]{data/educ_wage_bar.csv}\web
	\pgfplotstableread[header=true, col sep=comma]{data/unemp_reason_mon.csv}\unrsn
	\pgfplotstableread[header=true, col sep=comma]{data/inf_exp_ch.csv}\rych
	\pgfplotstableread[header=true, col sep=comma]{data/inf_exp_ch2.csv}\rytch
	
	% Required for bar plots with individual bar colors for categories
	\pgfplotsset{discard if not/.style 2 args={
        x filter/.code={
            \edef\tempa{\thisrow{#1}}
            \edef\tempb{#2}
            \ifx\tempa\tempb
            \else
                \def\pgfmathresult{inf}
            \fi}}}	
	
% % % % % % % %
%
%  Begin Document
%
% % % % % % % %		
\begin{document}

\chapter*{
		\includegraphics{blue_eagle.png}\\ \color{usaablue} Macroeconomic Chartbook}
\vspace*{-16mm}

\footnotesize \hspace{11mm} \color{usaablue}v0.1, \today \normalsize 

\vspace{11mm}

\normalsize \color{usaablue}Corporate Treasury view on Macroeconomy

\vspace{16mm}

\thispagestyle{empty}

\begin{minipage}{0.36\textwidth}
\subsection*{ {\color{red} \faExclamationTriangle} \seriffont Notes}

{\color{red} \textbf{Very early stage draft:}} \\
Content is draft form. \\ Not to be circulated 
\vfill

\end{minipage}
\begin{minipage}{0.36\textwidth}
\subsection*{{\color{gray} \faUser} Contact}

\textbf{Enterprise Modeling} \  \\
{\color{gray} \faEnvelope} \ edward.blackburne@usaa.com \ \\

\end{minipage}
\newpage
\section*{\hyperlink{toc}{\faList} \ {Contents}}
\hypertarget{toc}{}
\vspace{4mm}

\begin{minipage}{0.4\textwidth}
\begin{description}
\item {\hyperlink{oea}{Overall Economic Activity}}
\begin{description}
\item {\hyperlink{oety}{Types of Activity}}
\item {\hyperlink{oegr}{Economic Growth}}
\item {\hyperlink{oegc}{Components of Growth}}
\end{description}
\item {\hyperlink{ofa}{Overall Financial Activity}}
\begin{description}
\item {\hyperlink{ofsb}{Sectoral Balances}}
\item Borrowing and Liabilities
\item Assets
\item {\hyperlink{ofw}{Wealth}}
\end{description}
\item {\hyperlink{hh}{Households}}
\begin{description}
\item {\hyperlink{hhdem}{Demographics}}
\item {\hyperlink{hhinc}{Income}}
\item {\hyperlink{hhss}{Spending and Saving}}
\item {\hyperlink{hhbs}{Balance Sheets}}
\item {\hyperlink{hhpov}{Poverty}}
\end{description}
\item {\hyperlink{bus}{Businesses}}
\begin{description}
\item {\hyperlink{busin}{Investment}}
\item {\hyperlink{buspr}{Profit}}
\item {\hyperlink{busip}{Industrial Production}}
\item {\hyperlink{busrs}{Retail Sales}}
\end{description}
\item {\hyperlink{gov}{Government}}
\begin{description}
\item Spending and Investment
\item Revenue
\item {\hyperlink{govbs}{Balance Sheets}}
\end{description}
\end{description}
\end{minipage} \hspace{10mm} 
\begin{minipage}{0.4\textwidth}
\begin{description}
\item {\hyperlink{ext}{External Sector}}
\begin{description}
\item Balance of Payments
\item International Investment Position
\item {\hyperlink{extt}{Trade}}
\item {\hyperlink{extfx}{Exchange Rates}}
\end{description}
\item {\hyperlink{lab}{Labor Markets}}
\begin{description}
\item {\hyperlink{labe}{Employment}}
\item {\hyperlink{labu}{Unemployment}}
\item {\hyperlink{labp}{Participation}}
\item {\hyperlink{labf}{Labor Force Flows}}
\item {\hyperlink{labh}{Hours}}
\item {\hyperlink{labw}{Wages}}
\item {\hyperlink{labprod}{Productivity}}
\item {\hyperlink{labun}{Union Membership}}
\end{description}
\item {\hyperlink{cap}{Capital Markets}}
\begin{description}
\item {\hyperlink{capeq}{Equity Markets}}
\item {\hyperlink{capint}{Interest Rates}}
\item {\hyperlink{capmm}{Money and Monetary Policy}}
\end{description}
\item {\hyperlink{pr}{Prices}}
\begin{description}
\item {\hyperlink{prin}{Price Indices}}
\item {\hyperlink{prie}{Inflation Expectations}}
\item {\hyperlink{prco}{Commodities}}
\end{description}
\end{description}
\vspace{1.2cm}

\end{minipage}
\vspace{6mm}


\newpage
\section*{Overall Economic Activity}
\hypertarget{oea}{}
\begin{minipage}{0.61\textwidth}
\small This analysis of the United States economy begins with the most popular measure of economic activity, \href{https://www.bea.gov/data/gdp/gross-domestic-product}{Gross Domestic Product} (GDP). GDP estimates the production of goods and services in a given location during a given period of time. According to the Bureau of Economic Analysis, the seasonally-adjusted annualized value of goods and services produced in the US was \input{text/gdp1.txt}
\end{minipage}\hfill
\begin{minipage}{0.34\textwidth}
\noindent \normalsize \textbf{GDP per capita}\\
\footnotesize{\textit{in thousands of \input{text/gdp_ltdt.txt}dollars}}\\ 
\noindent \begin{tikzpicture}
    \begin{axis}[\shdateaxisticks clip=false, height=5.0cm, width=5.9cm,
        xticklabel={`\short{\year}}, enlarge y limits={lower, 0.2}, 
        enlarge x limits={0.02}]
    \rbars
    \stdline{red!95!black}{date}{value}{data/gdppc.csv}
	\input{text/gdp_pc_node.txt}
    \end{axis}
\end{tikzpicture}\\
\footnotesize{Source: Bureau of Economic Analysis}
\end{minipage}

\hypertarget{oety}{}
\begin{minipage}{0.76\textwidth}
\subsection*{Types of Economic Activity}
\small GDP calculated using the \href{https://www.bea.gov/help/glossary/expenditures-approach}{expenditures approach} is the sum of major types of domestic spending on finished goods and services: consumer spending, private investment, and government spending and investment. To capture only domestic production, foreign spending on US produced goods and services is added, while imports (spending on non-US-produced goods and services) are subtracted.
\end{minipage}

\begin{minipage}{0.37\textwidth}
\normalsize \textbf{Expenditure Types}\\
\footnotesize{\textit{per capita, thousands of \input{text/gdp_ltdt.txt}dollars}}\\
\hspace*{-2mm} \begin{tikzpicture}
	\begin{axis}[\bbar{y}{0}, \shdateaxisticks ytick={0, 10, 20, 30, 40, 50}, 
		yticklabel style={text width=1.0em}, 
		xticklabel={`\short{\year}}, clip=false, height=5.6cm, width=7.0cm]
	\rbars
	\stdnode{1.1cm}{3.35cm}{\footnotesize \color{yellow!45!orange}Consumer\\ \color{yellow!55!orange}Spending}
	\stdline{yellow!45!orange}{date}{DPCERX}{data/gdp_pc_comp.csv}
	\stdnode{1.6cm}{0.68cm}{\footnotesize \color{blue!70!black}Investment}
	\stdline{blue!70!black}{date}{A006RX}{data/gdp_pc_comp.csv}
	\stdnode{1.2cm}{1.4cm}{\footnotesize \color{cyan!60!white}Government}
	\stdline{cyan!60!white}{date}{A822RX}{data/gdp_pc_comp.csv}
	\stdnode{4.4cm}{0.5cm}{\footnotesize \color{green!60!black}Net Exports}
	\stdline{green!60!black}{date}{A019RX}{data/gdp_pc_comp.csv}
	\end{axis}
\end{tikzpicture}\\
\footnotesize{Source: Bureau of Economic Analysis} \hspace{8mm} \tbllink{gdp_pc_comp.csv}
\end{minipage}\hspace{8mm}
\begin{minipage}{0.33\textwidth}
\small \input{text/gdp_pc_comp.txt}
\end{minipage}

\vspace{3mm}

\begin{minipage}{0.77\textwidth}
\normalsize \textbf{Expenditure Types}\\
\footnotesize{\textit{per capita, annualized, \input{text/gdp_ltdt.txt}dollars}}\\
\hspace{-2mm} \rowcolors{1}{}{black!5} \setlength{\tabcolsep}{3.7pt} \color{black!90}
		{\renewcommand{\arraystretch}{1.55}
		 \begin{tabular}{p{53mm} R{13mm} R{13mm} R{13mm} R{13mm}}
			\input data/gdppc_levels.tex \hline
		\end{tabular}} \vspace{-2mm}
		
\footnotesize{Source: Bureau of Economic Analysis}
\end{minipage}
\newpage
\begin{minipage}{0.76\textwidth}
\small The less-discussed \textbf{\href{https://www.bea.gov/help/glossary/income-approach}{income approach}} calculates the level of overall economic activity from the sum of income payments from production and certain expenses incurred in production. Income from production is the payout to labor and capital. Income for labor is referred to as ``compensation of employees, paid'' in the national accounts and includes wages and salaries as well supplements to wages and salaries such as employer-paid health insurance premiums and retirement account contributions. Income received as profit is referred to as the ``net operating surplus'' in national accounts and includes interest payments, rental profits, business proprietor profits, and corporate profits.

Additionally some income does not get paid out to these groups. Taxes on imports and production, such as tariffs, sales tax, property tax, and licensing fees are tracked separately. If the government provides subsidies, which are income payments for production that did not occur, they must be subtracted from income measures of production. Lastly, a growing portion of potential income goes toward replacing and maintaining buildings and equipment used in production. This expense is recorded as ``consumption of fixed capital'' in the national accounts and referred to as depreciation below.

The Bureau of Economic Analysis \href{https://www.bea.gov/data/gdp/gross-domestic-product}{report} seasonally-adjusted and annualized \href{https://www.bea.gov/help/glossary/gross-domestic-income-gdi}{Gross Domestic Income} (GDI) \input{text/gdi_levels.txt}
\end{minipage}
\vspace{1mm}

\begin{minipage}{0.37\textwidth}
\noindent \normalsize \textbf{Income Types}\\
\footnotesize{\textit{per capita, thousands of \input{text/gdi_ltdt.txt}dollars}}\\
\noindent \hspace*{-2mm} \begin{tikzpicture}
	\begin{axis}[\bbar{y}{0}, \shdateaxisticks ytick={0, 10, 20, 30, 40}, 
		yticklabel style={text width=1.0em}, 
		xticklabel={`\short{\year}}, clip=false, height=5.6cm, width=7.0cm]
	\rbars
	\stdnode{1.1cm}{3.0cm}{\footnotesize \color{magenta!90!blue}Labor}
	\stdline{magenta!90!blue}{date}{A4002C}{data/gdi_pc_comp.csv}
	\stdnode{4.3cm}{0.86cm}{\footnotesize \color{teal!60!white}Depreciation}
	\stdline{yellow!60!orange}{date}{W271RC}{data/gdi_pc_comp.csv}
	\stdnode{1.0cm}{1.4cm}{\footnotesize \color{yellow!60!orange}Profit}
	\stdline{teal!60!white}{date}{A262RC}{data/gdi_pc_comp.csv}
	\stdnode{4.1cm}{0.3cm}{\footnotesize \color{violet}Indirect Taxes}
	\stdline{violet}{date}{indirect}{data/gdi_pc_comp.csv}
	\end{axis}
\end{tikzpicture}\\
\footnotesize{Source: Bureau of Economic Analysis} \hspace{8mm} \tbllink{gdi_pc_comp.csv}\\
\end{minipage}\hspace{8mm}
\begin{minipage}{0.33\textwidth}
\small \input{text/gdi_levels_pc.txt}
\end{minipage}

\vspace{3mm}

\begin{minipage}{0.76\textwidth}

\normalsize \textbf{Income Types}\\
\footnotesize{\textit{per capita, annualized, \input{text/gdi_ltdt.txt}dollars}}\\
\hspace{-2mm} \rowcolors{1}{}{black!5} \setlength{\tabcolsep}{3.7pt} \color{black!90}
		{\renewcommand{\arraystretch}{1.6}
		 \begin{tabular}{p{54mm} R{13mm} R{13mm} R{13mm} R{13mm}}
			\input data/gdipc_levels.tex \hline
		\end{tabular}}
\vspace{-2mm}

\footnotesize{Source: Bureau of Economic Analysis}
\end{minipage}
\newpage
\begin{minipage}{0.76\textwidth}
\small The \textbf{production approach} to GDP identifies how individual industries contribute to domestic production by calculating the \href{https://www.bea.gov/help/glossary/value-added}{value added} by each industry during the production process. The value added by an industry or sector group is its sales or gross output minus any \href{https://www.bea.gov/help/glossary/intermediate-inputs}{intermediate inputs} used in production. The Bureau of Economic Analysis \href{https://www.bea.gov/data/gdp/gdp-industry}{report} GDP by industry, which is summarized briefly in this subsection by grouping the various private industries into broad categories. 

\input{text/gdpva_sh.txt}
\end{minipage}
\vspace{1mm}

\begin{minipage}{0.35\textwidth}
\noindent \normalsize \textbf{Production Types}\\
\footnotesize{\textit{per capita, thousands of \input{text/gdpva_ltdt.txt}dollars}}\\
\noindent \hspace*{-2mm} \begin{tikzpicture}
	\begin{axis}[\bbar{y}{0}, \shdateaxisticks ytick={0, 5, 10, 15, 20, 25}, 
		yticklabel style={text width=1.0em}, ymin=1.2, xmin={1998-03-01},
		xticklabel={`\short{\year}}, clip=false, height=6.1cm, width=6.4cm]
	\rebars
	\stdnode{3.75cm}{1.25cm}{\footnotesize \color{orange!80!white}Government}
	\stdline{orange!80!white}{date}{Government}{data/gdpva_pc.csv}
	\stdnode{0.95cm}{3.0cm}{\footnotesize \color{blue!90!white}Other Services}
	\stdline{blue!90!white}{date}{Oth_Serv}{data/gdpva_pc.csv}
	\stdnode{1.45cm}{2.0cm}{\footnotesize \color{red!90!white}FIRE+}
	\stdline{red!90!white}{date}{FIRE}{data/gdpva_pc.csv}
	\stdnode{3.2cm}{4.25cm}{\footnotesize \color{purple!70!blue}Goods + TTU}
	\stdline{purple!70!blue}{date}{GoodsTTU}{data/gdpva_pc.csv}
	\stdnode{2.1cm}{0.3cm}{\footnotesize See text for category descriptions}
	\end{axis}
\end{tikzpicture}\\
\footnotesize{Source: Bureau of Economic Analysis} \hfill \tbllink{gdpva_pc.csv}\\
\end{minipage}\hspace{6mm}
\begin{minipage}{0.365\textwidth}
\small & 2023  Q3 & 2023  Q2 & 2019  Q4 & 2005  Q1 & 1997  (A) \\  \hspace{0.5mm}  {\color{purple!70!blue}\textbf{---}}  Goods  and  TTU   & \$28,296 &27,832&28,257&25,339&21,669\\  \hspace{6mm}  Manufacturing   &8,504&8,330&8,303&7,802&6,430\\  \hspace{6mm}  Construction   &3,635&3,510&3,817&3,860&3,625\\  \hspace{6mm}  Retail  Trade   &5,244&4,983&5,108&4,330&3,409\\  \hspace{0.5mm}  {\color{red!90!white}\textbf{---}}  FIRE+   &21,483&21,168&19,199&13,426&10,449\\  \hspace{6mm}  Finance  \&  Insurance   &5,973&5,893&6,092&nan&3,265\\  \hspace{6mm}  Information   &4,460&4,317&3,362&1,858&1,170\\  \hspace{0.5mm}  {\color{blue!90!white}\textbf{---}}  Other  Services   &23,178&23,130&21,108&15,182&12,850\\  \hspace{6mm}  Education  \&  Healthcare   &7,060&7,015&6,499&4,574&3,931\\  \hspace{6mm}  Professional  \&  Business &10,641&10,580&8,872&6,037&4,900\\  \hspace{0.5mm}  {\color{orange!80!white}\textbf{---}}  Government   &9,337&9,304&9,322&8,661&8,463\\  
\end{minipage}
\vspace{1mm}

\begin{minipage}{0.76\textwidth}

\normalsize \textbf{Production Types}\\
\footnotesize{\textit{per capita, annualized, \input{text/gdpva_ltdt.txt}dollars}}\\
\hspace{-2mm} \rowcolors{1}{}{black!5} \setlength{\tabcolsep}{3.7pt} \color{black!90}
		{\renewcommand{\arraystretch}{1.6}
		 \begin{tabular}{p{38mm} R{12.8mm} R{12.8mm} R{12.8mm} R{12.8mm} R{12.8mm}}
			\input data/gdpva_pc.tex \hline
		\end{tabular}}
\vspace{-2mm}

\footnotesize{Source: Bureau of Economic Analysis}
\end{minipage}
\newpage
\begin{minipage}{0.76\textwidth}
\subsubsection*{Household Inputs to Production}
\small It's useful to consider household inputs when analyzing economic output. For example, is the population growing? Are more people working? Are people working more hours? Is the economy more productive in its use of labor? These questions all add important context to discussion of aggregate output.

The US population is increasing (see {\color{lime!90!green}\textbf{---}}), though at a slower rate than in the past. Over the past 30 years, the US has added 80 million people. This means the amount of goods and services produced needed to increase by 30 percent over the period to maintain a constant standard of living. Employment of any amount (see {\color{green!30!teal!90!black}\textbf{---}}) comes from around half of the population during the peak of an economic expansion as a tighter labor market and higher wages pull people into employment, and falls during a recession as people lose jobs. Hours worked per worker per week (see {\color{blue}\textbf{---}}) are also influenced by economic conditions. During an economic expansion, more hours of work are available to those who want them. During a recession and when demand for goods and services is relatively low, hours of work are cut.
\end{minipage}
\vspace{2mm}

\begin{minipage}{0.32\textwidth}
\normalsize \textbf{Population}\\
\footnotesize{\textit{millions of people}}\\
\hspace*{-2mm} \begin{tikzpicture}
	\begin{axis}[\shdateaxisticks ytick={200, 220, 240, 260, 280, 300, 320, 340}, 
		yticklabel style={text width=1.3em}, 
		enlarge y limits={0.1}, 
		xticklabel={`\short{\year}}, height=4.2cm, width=5.8cm, clip=false]
	\rbars
	\thickline{lime!90!green}{date}{pop}{data/gdpjobslvl.csv}
	\end{axis}
\end{tikzpicture}\\
\footnotesize{Source: BEA}
\end{minipage} \hspace{4mm}
\begin{minipage}{0.32\textwidth}
\normalsize \textbf{Employment Rate}\\
\footnotesize{\textit{share of population}}\\
\hspace*{-2mm} \begin{tikzpicture}
	\begin{axis}[\shdateaxisticks ytick={42, 44, 46, 48, 50, 52}, 
		yticklabel style={text width=1.0em}, 
		enlarge y limits={0.1}, 
		xticklabel={`\short{\year}}, height=4.2cm, width=5.8cm, clip=false]
	\rbars
	\stdline{green!30!teal!90!black}{date}{epop}{data/gdpjobslvl.csv}
	\end{axis}
\end{tikzpicture}\\
\footnotesize{Source: CPS microdata}
\end{minipage}\hspace{4mm}
\begin{minipage}{0.32\textwidth}
\normalsize \textbf{Hours Worked}\\
\footnotesize{\textit{average per worker per week}}\\
\hspace*{-2mm} \begin{tikzpicture}
	\begin{axis}[\shdateaxisticks ytick={36, 37, 38}, 
		yticklabel style={text width=1.0em}, 
		enlarge y limits={0.1}, 
		xticklabel={`\short{\year}}, height=4.2cm, width=5.8cm, clip=false]
	\rbars
	\stdline{blue}{date}{hours}{data/gdpjobslvl.csv}
	\end{axis}
\end{tikzpicture}\\
\footnotesize{Source: CPS microdata}
\end{minipage}

\vspace{4mm}

\begin{minipage}{0.76\textwidth}
\small By multiplying the population, employment rate, and average hours worked, we approximate the aggregate hours worked in the US. We can then divide real GDP by aggregate hours worked to estimate GDP per hour of work, or productivity (see {\color{cyan!60!white}\textbf{---}}). During a recession, economic theory suggests that productivity increases, in the short-run, as less-productive workers disproportionately lose jobs: less-productive firms close and existing firms lay off less-productive workers.
\end{minipage}

\begin{minipage}{0.3\textwidth}
\small \input{text/gdp_per_hour.txt}
\end{minipage}\hspace{8mm}
\begin{minipage}{0.405\textwidth}
\normalsize \textbf{GDP per hour of work}\\
\footnotesize{\textit{\input{text/gdpjobs_ltdt.txt}dollars}}\\
\hspace*{-2mm} \begin{tikzpicture}
	\begin{axis}[\shdateaxisticks ytick={40, 50, 60, 70, 80}, 
		yticklabel style={text width=1.0em}, 
		enlarge y limits={0.1}, 
		xticklabel={`\short{\year}}, height=5.6cm, width=7.4cm, clip=false]
	\rbars
	\thickline{cyan!60!white}{date}{gdpinp}{data/gdpjobslvl.csv}
	\end{axis}
\end{tikzpicture}\\
\footnotesize{Source: Author's calculations} \hfill \tbllink{gdpjobslvl.csv} 
\end{minipage}
\newpage
\hypertarget{oegr}{}
\begin{minipage}{0.76\textwidth}
\subsection*{Economic Growth}
\small Economists are concerned with changes in economic activity. This subsection discusses economic growth, recessions, and contributions to economic growth.

\subsubsection*{Real GDP Growth}
\vspace{-1mm}

\small \input{text/gdp_gr.txt} 
\vspace{1mm}

\normalsize \textbf{Real Gross Domestic Product Growth}\\
\footnotesize{\textit{quarterly growth at seasonally adjusted annual rate, percent}}\\
\hspace*{-2mm} \begin{tikzpicture}
	\begin{axis}[\bbar{y}{0}, \dateaxisticks ytick={-30, -20, -10, 0, 10, 20, 30}, 
		yticklabel style={text width=1.4em}, 
		xticklabel={`\short{\year}}, clip=false, height=5.4cm]
	\rbars
	\tsbar{red!95!black}{date}{A191RL}{data/gdp.csv}
	\absnode{{1991-04-01}}{37}{\scriptsize \color{black!40} Early `90s\\ \scriptsize \color{black!40}  recession}
	\absnode{{2001-12-01}}{37}{\scriptsize \color{black!40} Early `00s\\ \scriptsize \color{black!40}  recession}
	\absnode{{2009-07-01}}{37}{\scriptsize \color{black!40} Great\\ \scriptsize \color{black!40}  recession}
	\node[below left, align=right] at (axis cs: {2020-01-01},37) {\scriptsize \color{black!40} COVID-19\\ \scriptsize \color{black!40}  recession};
	\end{axis}
\end{tikzpicture}
\vspace{-4.5mm}

\footnotesize{\textit{annual growth, percent \hspace{64mm} recent quarterly}}\\
\hspace*{-2mm} \begin{tikzpicture}
    \begin{groupplot}[group style={group size=2 by 1, 
    y descriptions at=edge left, horizontal sep=12pt,},]
    \nextgroupplot[\bbar{y}{0}, \shdateaxisticks 
		yticklabel style={text width=1.0em}, ymin=-4.0, ymax=8.6,
		ytick={-4, -2, 0, 2, 4, 6, 8},
		xticklabel={\year}, clip=false, height=4.8cm, width=10.0cm]
	\rbars
	\ctsbar{red!95!black}{date}{VALUE}{data/gdp_a.csv}{6.8pt}
    \nextgroupplot[ybar=24pt, \bbar{y}{0}, date coordinates in=x, 
    	xtick=data,	xticklabels from table={\gdprec}{Quarter},
    	every node near coord/.append style={yshift=5pt},
		clip=false, height=4.8cm, width=4.0cm, xtick style={draw=none},
		xticklabel style={black!80, yshift=10pt},
		enlarge y limits={0.06}, enlarge x limits={0.2},
		yticklabel style={text width=1.2em}, ymin=-4.0, ymax=8.6,
		ytick={-4, -2, 0, 2, 4, 6, 8},
		axis line style={draw=none}, nodes near coords,
		nodes near coords style={/pgf/number format/.cd, fixed zerofill,
			precision=1, assume math mode}]
	\ctsbar{red!95!black}{date}{A191RL}{\gdprec}{5.4pt}
	\end{groupplot}
\end{tikzpicture}\\
\footnotesize{Source: Bureau of Economic Analysis} \hfill \tbllink{gdp.csv}
\vspace{1mm}

\small \input{text/gdp_gr2.txt} 
\end{minipage}
\newpage
\begin{minipage}{0.76\textwidth}
\subsubsection*{Recessions}
\vspace{-1mm}

\small The long-term pattern in real GDP growth is often described as the business cycle. Periods of economic growth are interrupted every 7--12 years by an \textbf{economic recession}, a period where the amount of economic activity falls. The National Bureau of Economic Research (NBER) \href{https://www.nber.org/research/business-cycle-dating}{identify} four recessions since 1989.

\input{text/recessions.txt}
\vspace{0.5mm}

\normalsize \textbf{US Recessions Since 1989}
\vspace{-4mm}

\footnotesize \rowcolors{1}{}{black!5} \setlength{\tabcolsep}{1.0pt} \color{black!90}
	{\renewcommand{\arraystretch}{1.62}
\hspace*{-2mm} \begin{tabular}{p{29mm} C{16mm} C{16mm} C{13.2mm} C{13.2mm} C{13.2mm} C{13.2mm}}
		\input data/recession.tex \hline
	\end{tabular}}
	\vspace{-1mm}
	
\footnotesize{Sources: NBER, BEA, BLS}\\
\footnotesize{*Percentage point change from average unemployment rate during three years prior to recession to peak unemployment rate. **Months from recession start until unemployment rate returns to pre-recession three-year average.}
\vspace{5mm}

\small The most-reliable indication that the US has entered a recession (see\cbox{black!30}) was identified by \href{https://sites.google.com/site/claudiasahm/}{Claudia Sahm}, and is called the \textbf{Sahm rule}. The Sahm rule indicates the start of a recession (see \begin{tikzpicture} \node[line width=0.3mm, diamond, draw=red, scale=0.4, aspect=0.8] (d) at (0,0) {}; \end{tikzpicture}) when the three-month moving average unemployment rate rises by half a percentage point or more above its low during the previous twelve months (see {\color{blue!60!black}\textbf{---}}). In effect, the Sahm rule identifies increases in unemployment that are significant enough to cause or indicate a recession.  
\vspace{1mm}

\normalsize \textbf{The Sahm Rule}\\
\footnotesize{\textit{three-month moving average of unemployment rate, }}\\
\footnotesize{\textit{percentage points above one-year low}}
\vspace{36mm}

\hspace{3mm} \begin{tikzpicture}[overlay]
	\begin{axis}[\bbar{y}{0}, \dateaxisticks 
		enlarge y limits={0.1}, ymax=6.5, 
		xticklabel={`\short{\year}}, width=13.0cm, height=5.7cm,
		clip=false]
	\fill[color=black!30] (axis cs:{1990-07-01},\pgfkeysvalueof{/pgfplots/ymin}) rectangle 
		(axis cs:{1991-03-01}, \pgfkeysvalueof{/pgfplots/ymax});
	\fill[color=black!30] (axis cs:{2007-12-01},\pgfkeysvalueof{/pgfplots/ymin}) rectangle 
		(axis cs:{2009-07-01}, \pgfkeysvalueof{/pgfplots/ymax});
	\fill[color=black!30] (axis cs:{2001-03-01},\pgfkeysvalueof{/pgfplots/ymin}) rectangle 
		(axis cs:{2001-11-01}, \pgfkeysvalueof{/pgfplots/ymax});
	\fill[color=black!30] (axis cs:{2020-02-01},\pgfkeysvalueof{/pgfplots/ymin}) rectangle 
		(axis cs:{2020-05-01}, \pgfkeysvalueof{/pgfplots/ymax});
	\absnode{{1991-04-01}}{5.8}{\scriptsize \color{black!60} Early `90s\\ \scriptsize \color{black!60}  recession}
	\absnode{{2001-12-01}}{5.8}{\scriptsize \color{black!60} Early `00s\\ \scriptsize \color{black!60}  recession}
	\absnode{{2009-07-01}}{5.8}{\scriptsize \color{black!60} Great\\ \scriptsize \color{black!60}  recession}
	\node[below left, align=right] at (axis cs: {2020-01-01},5.8) {\scriptsize \color{black!60} COVID-19\\ \scriptsize \color{black!60}  recession};
	\stdline{gray}{date}{Bar}{data/sahm_bar.csv}
	\thickline{blue!60!black}{date}{Total}{data/sahm.csv}
	\addplot[only marks, mark=diamond*, mark options={fill=white, 
		draw=red, scale=1.6}] table [x=date, y=Intersect, col sep=comma] {data/sahm_marks.csv};	
	\input{text/sahm_bar_node.txt}
	\end{axis}
\end{tikzpicture}

\footnotesize{Source: Claudia Sahm} \hfill \tbllink{sahm.csv}
\end{minipage}
\newpage
\label{sec:oegc}
\begin{minipage}{0.76\textwidth}
\subsection*{Components of Growth}
\small The \textbf{expenditure approach} also gives insight into the sources of changes in overall economic activity. \input{text/gdp_exp_basic.txt}
\vspace{1mm}

\normalsize \textbf{Real GDP Growth by Expenditure Type}\\
\footnotesize{\textit{percentage point contribution to GDP growth}}
\vspace{4.6cm}

\hspace{4mm} \begin{tikzpicture}[overlay]
	\begin{axis}[\bbar{y}{0}, \dateaxisticks ytick={-30, -20, -10, 0, 10, 20, 30},
		yticklabel style={text width=1.4em}, height=6.6cm, ymax=26.6,
		xticklabel={`\short{\year}}, clip=false, legend style={at={(0.86, 0.99)}, fill=white}]
	\rbars
	\sbar{yellow!80!orange}{date}{DPCERY}{data/comp.csv}
	\sbar{blue!70!black}{date}{A006RY}{data/comp.csv}
	\sbar{cyan!50!white}{date}{A822RY}{data/comp.csv}
	\sbar{green!60!black}{date}{A019RY}{data/comp.csv}
	\stdnode{2.3cm}{0.7cm}{\footnotesize $^*$ Includes change in private inventories}
	\legend{Consumer Spending, Investment$^*$, Government, Net Exports};
	\end{axis}
\end{tikzpicture}

\footnotesize{Source: Bureau of Economic Analysis} \hfill \tbllink{comp.csv}
\vspace{3mm}

\small The \textbf{income approach} enables decomposing annualized production growth into gross labor income (see\cbox{magenta!90!blue}), profit (see\cbox{yellow!60!orange}), indirect taxes less subsidies (see\cbox{violet}), and depreciation (see\cbox{teal!60!white}). 

\input{text/gdi_growth_comp.txt}
\vspace{1mm}

\normalsize \textbf{Real Gross Domestic Income Growth}\\
\footnotesize{\textit{percentage point contribution to GDI growth}}
\vspace{4.4cm}

\hspace{4mm} \begin{tikzpicture}[overlay]
	\begin{axis}[\bbar{y}{0}, \dateaxisticks ytick={-30, -20, -10, 0, 10, 20}, height=6.4cm,
		xticklabel={`\short{\year}}, clip=false, yticklabel style={text width=1.4em}, ymax=17.2,
		legend style={at={(0.66, 0.99)}, fill=white}]
	\rbars
	\sbar{magenta!90!blue}{date}{A4002C}{data/gdi.csv}
	\sbar{yellow!60!orange}{date}{W271RC}{data/gdi.csv}
	\sbar{teal!60!white}{date}{A262RC}{data/gdi.csv}
	\sbar{violet}{date}{indirect}{data/gdi.csv}	
	\legend{Labor, Profit, Depreciation, Indirect Taxes};
	\end{axis}
\end{tikzpicture}

\footnotesize{Source: Bureau of Economic Analysis} \hfill \tbllink{gdi.csv}
\end{minipage}
\newpage
\begin{minipage}{0.76\textwidth}
\small The \textbf{production approach} calculates GDP as the sum of value added--gross output minus intermediate inputs--in each sector. The broad groupings discussed above are used to identify contributions from: goods-producing sectors combined with trade, transportation, and utilities (see\cbox{purple!70!blue}), finance, insurance, and real estate plus information (see\cbox{red!90!white}), other service-providing sectors (see\cbox{blue!90!white}), and government (see\cbox{orange!80!white}).

\input{text/gdpva_growth.txt}
\vspace{1mm}

\normalsize \textbf{Real GDP Growth by Industry Group}\\
\footnotesize{\textit{percentage point contribution to GDP growth}}\\
\hspace*{-2mm} \begin{tikzpicture}
	\begin{axis}[\bbar{y}{0}, \dateaxisticks ytick={-30, -20, -10, 0, 10, 20, 30, 40},
		height=9.6cm,  yticklabel style={text width=1.5em},
		xticklabel={`\short{\year}}, clip=false, legend style={at={(0.74, 1.0)}, fill=white}]
	\rbars
	\draw [dashed] (axis cs:{2004-10-01},\pgfkeysvalueof{/pgfplots/ymin}) -- 
		(axis cs:{2004-10-01},
		\pgfkeysvalueof{/pgfplots/ymax});
	\draw [dashed] (axis cs:{1998-01-01},\pgfkeysvalueof{/pgfplots/ymin}) -- 
		(axis cs:{1998-01-01},
		\pgfkeysvalueof{/pgfplots/ymax});
	\sbar{orange!80!white}{date}{Government}{data/gdpva.csv}
	\sbar{blue!90!white}{date}{Oth_Serv}{data/gdpva.csv}
	\sbar{red!90!white}{date}{FIRE}{data/gdpva.csv}
	\sbar{purple!70!blue}{date}{GoodsTTU}{data/gdpva.csv}
	\stdnode{0.68cm}{0.45cm}{\scriptsize historical data}
	\stdnode{4.1cm}{0.45cm}{\scriptsize annual data}
	\legend{Government, Other Services, FIRE+, Goods and TTU};
	\end{axis}
\end{tikzpicture}\\
\footnotesize{Source: Bureau of Economic Analysis} \hfill \tbllink{gdpva.csv}
\end{minipage}
\newpage
\begin{minipage}{0.76\textwidth}
\small Changes to GDP can also be assigned to changes in \textbf{household inputs}: population (see\cbox{lime!90!green}), employment rates (see\cbox{green!30!teal!90!black}), average hours worked (see\cbox{blue}), and total economy productivity (see\cbox{cyan!60!white}). 

\input{text/gdpjobsch.txt}
\vspace{1mm}

\normalsize \textbf{Real GDP Growth by Household Inputs}\\
\footnotesize{\textit{percentage point contribution to GDP growth}}\\
\hspace*{-2mm} \begin{tikzpicture}
	\begin{axis}[\bbar{y}{0}, \dateaxisticks ytick={-50, -40, -30, -20, -10, 0, 10, 20, 30},
		height=13.6cm,
		xticklabel={`\short{\year}}, clip=false, yticklabel style={text width=1.5em},
		legend style={at={(0.84, 1.0)}, fill=white}]
	\rbars
	\sbar{lime!90!green}{date}{pop_contr}{data/gdpjobs.csv}
	\sbar{green!30!teal!90!black}{date}{epop_contr}{data/gdpjobs.csv}
	\sbar{cyan!60!white}{date}{prod}{data/gdpjobs.csv}
	\sbar{blue}{date}{hours_contr}{data/gdpjobs.csv}	
	\legend{Population, Employment Rate, Productivity, Average Hours};
	\end{axis}
\end{tikzpicture}\\
\footnotesize{Source: Author's Calculations} \hfill \tbllink{gdpjobs.csv}
\end{minipage}
\newpage
\normalsize \textbf{Components of Economic Growth}\\
\footnotesize{\textit{annualized percentage point contribution to real GDP/GDI growth \hspace{20mm} moving averages}}\\
\hspace*{-2mm} \noindent \rowcolors{1}{}{black!5} \setlength{\tabcolsep}{1.5pt} \color{black!90}
{\renewcommand{\arraystretch}{1.6}
	\begin{tabular}{p{4.0mm} p{38mm} R{9.0mm} R{9.0mm} R{9.0mm} R{9.0mm} R{9.0mm} 
	p{0mm} R{9.0mm} R{9.0mm} R{9.0mm} }
    &      & 2023  Q3 & 2023  Q2 & 2023  Q1 & 2022  Q4 & 2022  Q3 & \  \  \  3-    year & \  \  \  10-    year & \  \  \  30-    year \\  \cbox{red!95!black} & \textbf{Gross  Domestic  Product} &4.9&2.1&2.2&2.6&2.7&3.1&2.6&2.6\\  \cbox{yellow!80!orange} & \hspace{2mm}  Consumer  Spending &2.11&0.55&2.54&0.79&1.05&2.58&1.90&1.85\\ & \hspace{4mm}  Durable  Goods &0.53&-0.03&1.07&-0.08&0.08&0.33&0.50&0.48\\ & \hspace{4mm}  Non-Durable  Goods   &0.56&0.14&0.07&0.07&-0.26&0.37&0.46&0.37\\ & \hspace{4mm}  Services   &1.02&0.44&1.40&0.80&1.23&1.88&0.95&1.00\\  \cbox{blue!70!black} & \hspace{2mm}  Gross  Investment &1.74&0.90&-1.69&0.62&-1.45&0.63&0.67&0.70\\ & \hspace{4mm}  Residential   &0.26&-0.09&-0.22&-1.23&-1.41&-0.19&0.07&0.03\\ & \hspace{4mm}  Non-Residential   &0.21&0.98&0.76&0.24&0.62&0.74&0.54&0.60\\ & \hspace{4mm}  Change  in  Inventories   &1.27&0.00&-2.22&1.61&-0.66&0.08&0.06&0.08\\  \cbox{green!60!black} & \hspace{2mm}  Net  Exports   &0.03&0.04&0.58&0.26&2.58&-0.27&-0.26&-0.19\\ & \hspace{4mm}  Exports   &0.59&-1.09&0.76&-0.41&1.80&0.62&0.21&0.45\\ & \hspace{4mm}  Imports   &-0.56&1.13&-0.18&0.66&0.77&-0.90&-0.46&-0.64\\  \cbox{cyan!50!white} & \hspace{2mm}  Government   &0.99&0.57&0.82&0.90&0.49&0.21&0.30&0.25\\ & \hspace{4mm}  Federal   &0.45&0.07&0.33&0.59&0.07&0.07&0.11&0.10\\ & \hspace{4mm}  State  and  Local   &0.53&0.50&0.49&0.31&0.41&0.14&0.19&0.15\\ &&&&&&&&&\\  \cbox{purple!70!blue} & \hspace{2mm}  Goods  and  TTU   &2.53&1.80&0.30&2.14&0.59&0.35&0.63&0.84\\ & \hspace{4mm}  Manufacturing   &0.92&0.64&-0.98&0.22&-0.01&0.16&0.16&0.35\\ & \hspace{4mm}  Construction   &0.65&0.23&0.14&-0.03&-0.67&-0.05&0.05&0.00\\ & \hspace{4mm}  Retail  Trade   &1.36&-0.18&0.69&0.80&0.59&0.06&0.15&0.18\\  \cbox{red!90!white} & \hspace{2mm}  FIRE+   &1.58&0.10&0.29&-0.32&0.22&0.96&0.79&0.78\\ & \hspace{4mm}  Information   &0.75&0.12&0.23&0.40&0.44&0.48&0.41&0.30\\  \cbox{blue!90!white} & \hspace{2mm}  Other  Services   &0.52&0.05&1.50&0.60&1.81&1.65&0.99&0.76\\ & \hspace{4mm}  Education  \&  Healthcare   &0.27&0.12&0.88&0.19&0.50&0.31&0.26&0.22\\ & \hspace{4mm}  Professional  \&  Business &0.38&0.27&0.15&0.54&0.86&0.89&0.60&0.44\\  \cbox{orange!80!white} & \hspace{2mm}  Government   &0.24&0.11&0.16&0.16&0.05&0.18&0.09&0.09\\ &&&&&&&&&\\  \cbox{lime!90!green} & \hspace{2mm}  Population   &0.59&0.47&0.41&0.53&0.53&0.35&0.55&0.84\\  \cbox{green!30!teal!80!black} & \hspace{2mm}  Employment  Rate   &-0.43&0.03&3.41&0.42&1.15&2.83&0.93&0.29\\  \cbox{blue} & \hspace{2mm}  Average  Hours &-0.20&-0.05&-0.12&-0.32&-0.39&0.18&-0.05&-0.03\\  \cbox{cyan!55!white} & \hspace{2mm}  Productivity   &4.90&1.61&-1.42&1.92&1.35&-0.15&1.43&1.62\\ &&&&&&&&&\\ & \textbf{Gross  Domestic  Income}   &1.5&0.5&0.5&-3.0&2.7&3.0&2.4&2.6\\  \cbox{magenta!90!blue} & \hspace{2mm}  Labor   &1.11&1.98&2.37&-0.94&2.62&1.21&1.31&1.28\\  \cbox{yellow!60!orange} & \hspace{2mm}  Profit   &0.28&-1.77&-2.42&-2.21&-0.97&-0.42&0.41&0.72\\  \cbox{teal!60!white} & \hspace{2mm}  Depreciation   &0.26&0.46&0.52&0.31&0.95&0.48&0.46&0.45\\  \cbox{violet} & \hspace{2mm}  Indirect  Taxes   &-0.12&-0.23&0.03&-0.16&0.08&1.68&0.19&0.18\\ 
	\end{tabular}}
	\vspace{-0.5mm}
				
\footnotesize{Source: Bureau of Economic Analysis and Author's Calculations}
\newpage

\noindent \normalsize \textbf{Real GDP Growth by State}\\
\footnotesize{\textit{percentage point change in real GDP}}\\
\vspace{-2mm}
\hspace{-8mm} %% Creator: Matplotlib, PGF backend
%%
%% To include the figure in your LaTeX document, write
%%   \input{<filename>.pgf}
%%
%% Make sure the required packages are loaded in your preamble
%%   \usepackage{pgf}
%%
%% Also ensure that all the required font packages are loaded; for instance,
%% the lmodern package is sometimes necessary when using math font.
%%   \usepackage{lmodern}
%%
%% Figures using additional raster images can only be included by \input if
%% they are in the same directory as the main LaTeX file. For loading figures
%% from other directories you can use the `import` package
%%   \usepackage{import}
%%
%% and then include the figures with
%%   \import{<path to file>}{<filename>.pgf}
%%
%% Matplotlib used the following preamble
%%   
%%   \usepackage{fontspec}
%%   \setmainfont{DejaVuSerif.ttf}[Path=\detokenize{/home/brian/miniconda3/lib/python3.8/site-packages/matplotlib/mpl-data/fonts/ttf/}]
%%   \setsansfont{DejaVuSans.ttf}[Path=\detokenize{/home/brian/miniconda3/lib/python3.8/site-packages/matplotlib/mpl-data/fonts/ttf/}]
%%   \setmonofont{DejaVuSansMono.ttf}[Path=\detokenize{/home/brian/miniconda3/lib/python3.8/site-packages/matplotlib/mpl-data/fonts/ttf/}]
%%   \makeatletter\@ifpackageloaded{underscore}{}{\usepackage[strings]{underscore}}\makeatother
%%
\begingroup%
\makeatletter%
\begin{pgfpicture}%
\pgfpathrectangle{\pgfpointorigin}{\pgfqpoint{6.714028in}{2.186394in}}%
\pgfusepath{use as bounding box, clip}%
\begin{pgfscope}%
\pgfsetbuttcap%
\pgfsetmiterjoin%
\definecolor{currentfill}{rgb}{1.000000,1.000000,1.000000}%
\pgfsetfillcolor{currentfill}%
\pgfsetlinewidth{0.000000pt}%
\definecolor{currentstroke}{rgb}{1.000000,1.000000,1.000000}%
\pgfsetstrokecolor{currentstroke}%
\pgfsetdash{}{0pt}%
\pgfpathmoveto{\pgfqpoint{0.000000in}{0.000000in}}%
\pgfpathlineto{\pgfqpoint{6.714028in}{0.000000in}}%
\pgfpathlineto{\pgfqpoint{6.714028in}{2.186394in}}%
\pgfpathlineto{\pgfqpoint{0.000000in}{2.186394in}}%
\pgfpathlineto{\pgfqpoint{0.000000in}{0.000000in}}%
\pgfpathclose%
\pgfusepath{fill}%
\end{pgfscope}%
\begin{pgfscope}%
\pgfpathrectangle{\pgfqpoint{0.100000in}{0.100000in}}{\pgfqpoint{2.989028in}{1.913466in}}%
\pgfusepath{clip}%
\pgfsetbuttcap%
\pgfsetmiterjoin%
\definecolor{currentfill}{rgb}{0.580392,0.831373,0.644444}%
\pgfsetfillcolor{currentfill}%
\pgfsetlinewidth{0.000000pt}%
\definecolor{currentstroke}{rgb}{0.000000,0.000000,0.000000}%
\pgfsetstrokecolor{currentstroke}%
\pgfsetstrokeopacity{0.000000}%
\pgfsetdash{}{0pt}%
\pgfpathmoveto{\pgfqpoint{1.097679in}{0.548675in}}%
\pgfpathlineto{\pgfqpoint{1.089089in}{0.548815in}}%
\pgfpathlineto{\pgfqpoint{1.083242in}{0.555287in}}%
\pgfpathlineto{\pgfqpoint{1.079266in}{0.572473in}}%
\pgfpathlineto{\pgfqpoint{1.088506in}{0.574627in}}%
\pgfpathlineto{\pgfqpoint{1.097893in}{0.572311in}}%
\pgfpathlineto{\pgfqpoint{1.107446in}{0.565540in}}%
\pgfpathlineto{\pgfqpoint{1.107441in}{0.556177in}}%
\pgfpathlineto{\pgfqpoint{1.097679in}{0.548675in}}%
\pgfpathclose%
\pgfusepath{fill}%
\end{pgfscope}%
\begin{pgfscope}%
\pgfpathrectangle{\pgfqpoint{0.100000in}{0.100000in}}{\pgfqpoint{2.989028in}{1.913466in}}%
\pgfusepath{clip}%
\pgfsetbuttcap%
\pgfsetmiterjoin%
\definecolor{currentfill}{rgb}{0.580392,0.831373,0.644444}%
\pgfsetfillcolor{currentfill}%
\pgfsetlinewidth{0.000000pt}%
\definecolor{currentstroke}{rgb}{0.000000,0.000000,0.000000}%
\pgfsetstrokecolor{currentstroke}%
\pgfsetstrokeopacity{0.000000}%
\pgfsetdash{}{0pt}%
\pgfpathmoveto{\pgfqpoint{1.149970in}{0.446640in}}%
\pgfpathlineto{\pgfqpoint{1.140359in}{0.450300in}}%
\pgfpathlineto{\pgfqpoint{1.140315in}{0.457479in}}%
\pgfpathlineto{\pgfqpoint{1.128705in}{0.463549in}}%
\pgfpathlineto{\pgfqpoint{1.130119in}{0.478835in}}%
\pgfpathlineto{\pgfqpoint{1.133481in}{0.486048in}}%
\pgfpathlineto{\pgfqpoint{1.140581in}{0.479409in}}%
\pgfpathlineto{\pgfqpoint{1.152127in}{0.480518in}}%
\pgfpathlineto{\pgfqpoint{1.149211in}{0.459391in}}%
\pgfpathlineto{\pgfqpoint{1.149970in}{0.446640in}}%
\pgfpathclose%
\pgfusepath{fill}%
\end{pgfscope}%
\begin{pgfscope}%
\pgfpathrectangle{\pgfqpoint{0.100000in}{0.100000in}}{\pgfqpoint{2.989028in}{1.913466in}}%
\pgfusepath{clip}%
\pgfsetbuttcap%
\pgfsetmiterjoin%
\definecolor{currentfill}{rgb}{0.580392,0.831373,0.644444}%
\pgfsetfillcolor{currentfill}%
\pgfsetlinewidth{0.000000pt}%
\definecolor{currentstroke}{rgb}{0.000000,0.000000,0.000000}%
\pgfsetstrokecolor{currentstroke}%
\pgfsetstrokeopacity{0.000000}%
\pgfsetdash{}{0pt}%
\pgfpathmoveto{\pgfqpoint{1.190310in}{0.399517in}}%
\pgfpathlineto{\pgfqpoint{1.178114in}{0.401594in}}%
\pgfpathlineto{\pgfqpoint{1.171369in}{0.412035in}}%
\pgfpathlineto{\pgfqpoint{1.170417in}{0.420620in}}%
\pgfpathlineto{\pgfqpoint{1.190310in}{0.399517in}}%
\pgfpathclose%
\pgfusepath{fill}%
\end{pgfscope}%
\begin{pgfscope}%
\pgfpathrectangle{\pgfqpoint{0.100000in}{0.100000in}}{\pgfqpoint{2.989028in}{1.913466in}}%
\pgfusepath{clip}%
\pgfsetbuttcap%
\pgfsetmiterjoin%
\definecolor{currentfill}{rgb}{0.580392,0.831373,0.644444}%
\pgfsetfillcolor{currentfill}%
\pgfsetlinewidth{0.000000pt}%
\definecolor{currentstroke}{rgb}{0.000000,0.000000,0.000000}%
\pgfsetstrokecolor{currentstroke}%
\pgfsetstrokeopacity{0.000000}%
\pgfsetdash{}{0pt}%
\pgfpathmoveto{\pgfqpoint{1.165317in}{0.401338in}}%
\pgfpathlineto{\pgfqpoint{1.171711in}{0.396147in}}%
\pgfpathlineto{\pgfqpoint{1.173099in}{0.387222in}}%
\pgfpathlineto{\pgfqpoint{1.160227in}{0.387715in}}%
\pgfpathlineto{\pgfqpoint{1.165317in}{0.401338in}}%
\pgfpathclose%
\pgfusepath{fill}%
\end{pgfscope}%
\begin{pgfscope}%
\pgfpathrectangle{\pgfqpoint{0.100000in}{0.100000in}}{\pgfqpoint{2.989028in}{1.913466in}}%
\pgfusepath{clip}%
\pgfsetbuttcap%
\pgfsetmiterjoin%
\definecolor{currentfill}{rgb}{0.580392,0.831373,0.644444}%
\pgfsetfillcolor{currentfill}%
\pgfsetlinewidth{0.000000pt}%
\definecolor{currentstroke}{rgb}{0.000000,0.000000,0.000000}%
\pgfsetstrokecolor{currentstroke}%
\pgfsetstrokeopacity{0.000000}%
\pgfsetdash{}{0pt}%
\pgfpathmoveto{\pgfqpoint{1.194683in}{0.351485in}}%
\pgfpathlineto{\pgfqpoint{1.181837in}{0.356103in}}%
\pgfpathlineto{\pgfqpoint{1.185518in}{0.367298in}}%
\pgfpathlineto{\pgfqpoint{1.180134in}{0.378705in}}%
\pgfpathlineto{\pgfqpoint{1.181825in}{0.386546in}}%
\pgfpathlineto{\pgfqpoint{1.192162in}{0.389535in}}%
\pgfpathlineto{\pgfqpoint{1.191963in}{0.376972in}}%
\pgfpathlineto{\pgfqpoint{1.204190in}{0.370003in}}%
\pgfpathlineto{\pgfqpoint{1.211183in}{0.353155in}}%
\pgfpathlineto{\pgfqpoint{1.203389in}{0.346667in}}%
\pgfpathlineto{\pgfqpoint{1.194683in}{0.351485in}}%
\pgfpathclose%
\pgfusepath{fill}%
\end{pgfscope}%
\begin{pgfscope}%
\pgfpathrectangle{\pgfqpoint{0.100000in}{0.100000in}}{\pgfqpoint{2.989028in}{1.913466in}}%
\pgfusepath{clip}%
\pgfsetbuttcap%
\pgfsetmiterjoin%
\definecolor{currentfill}{rgb}{0.580392,0.831373,0.644444}%
\pgfsetfillcolor{currentfill}%
\pgfsetlinewidth{0.000000pt}%
\definecolor{currentstroke}{rgb}{0.000000,0.000000,0.000000}%
\pgfsetstrokecolor{currentstroke}%
\pgfsetstrokeopacity{0.000000}%
\pgfsetdash{}{0pt}%
\pgfpathmoveto{\pgfqpoint{1.148289in}{0.237569in}}%
\pgfpathlineto{\pgfqpoint{1.142169in}{0.250211in}}%
\pgfpathlineto{\pgfqpoint{1.144879in}{0.258715in}}%
\pgfpathlineto{\pgfqpoint{1.156883in}{0.270485in}}%
\pgfpathlineto{\pgfqpoint{1.158019in}{0.279508in}}%
\pgfpathlineto{\pgfqpoint{1.165875in}{0.299484in}}%
\pgfpathlineto{\pgfqpoint{1.186536in}{0.302897in}}%
\pgfpathlineto{\pgfqpoint{1.190491in}{0.314590in}}%
\pgfpathlineto{\pgfqpoint{1.199600in}{0.310205in}}%
\pgfpathlineto{\pgfqpoint{1.218110in}{0.277388in}}%
\pgfpathlineto{\pgfqpoint{1.218344in}{0.266674in}}%
\pgfpathlineto{\pgfqpoint{1.213112in}{0.257400in}}%
\pgfpathlineto{\pgfqpoint{1.219750in}{0.236108in}}%
\pgfpathlineto{\pgfqpoint{1.215111in}{0.232676in}}%
\pgfpathlineto{\pgfqpoint{1.199311in}{0.233674in}}%
\pgfpathlineto{\pgfqpoint{1.180311in}{0.240024in}}%
\pgfpathlineto{\pgfqpoint{1.164342in}{0.242646in}}%
\pgfpathlineto{\pgfqpoint{1.160076in}{0.238921in}}%
\pgfpathlineto{\pgfqpoint{1.148289in}{0.237569in}}%
\pgfpathclose%
\pgfusepath{fill}%
\end{pgfscope}%
\begin{pgfscope}%
\pgfpathrectangle{\pgfqpoint{0.100000in}{0.100000in}}{\pgfqpoint{2.989028in}{1.913466in}}%
\pgfusepath{clip}%
\pgfsetbuttcap%
\pgfsetmiterjoin%
\definecolor{currentfill}{rgb}{0.352018,0.707266,0.669204}%
\pgfsetfillcolor{currentfill}%
\pgfsetlinewidth{0.000000pt}%
\definecolor{currentstroke}{rgb}{0.000000,0.000000,0.000000}%
\pgfsetstrokecolor{currentstroke}%
\pgfsetstrokeopacity{0.000000}%
\pgfsetdash{}{0pt}%
\pgfpathmoveto{\pgfqpoint{0.810222in}{1.900931in}}%
\pgfpathlineto{\pgfqpoint{0.794722in}{1.839512in}}%
\pgfpathlineto{\pgfqpoint{0.783700in}{1.795504in}}%
\pgfpathlineto{\pgfqpoint{0.770595in}{1.742564in}}%
\pgfpathlineto{\pgfqpoint{0.768022in}{1.729516in}}%
\pgfpathlineto{\pgfqpoint{0.769703in}{1.717547in}}%
\pgfpathlineto{\pgfqpoint{0.767475in}{1.706523in}}%
\pgfpathlineto{\pgfqpoint{0.675998in}{1.730394in}}%
\pgfpathlineto{\pgfqpoint{0.667734in}{1.727622in}}%
\pgfpathlineto{\pgfqpoint{0.660663in}{1.730028in}}%
\pgfpathlineto{\pgfqpoint{0.627921in}{1.730814in}}%
\pgfpathlineto{\pgfqpoint{0.616848in}{1.727692in}}%
\pgfpathlineto{\pgfqpoint{0.605851in}{1.728747in}}%
\pgfpathlineto{\pgfqpoint{0.601190in}{1.733634in}}%
\pgfpathlineto{\pgfqpoint{0.571877in}{1.732568in}}%
\pgfpathlineto{\pgfqpoint{0.567248in}{1.740554in}}%
\pgfpathlineto{\pgfqpoint{0.558460in}{1.744762in}}%
\pgfpathlineto{\pgfqpoint{0.545673in}{1.747325in}}%
\pgfpathlineto{\pgfqpoint{0.523603in}{1.743571in}}%
\pgfpathlineto{\pgfqpoint{0.515456in}{1.747249in}}%
\pgfpathlineto{\pgfqpoint{0.502895in}{1.757031in}}%
\pgfpathlineto{\pgfqpoint{0.506670in}{1.776202in}}%
\pgfpathlineto{\pgfqpoint{0.505347in}{1.785992in}}%
\pgfpathlineto{\pgfqpoint{0.495391in}{1.796386in}}%
\pgfpathlineto{\pgfqpoint{0.489015in}{1.795741in}}%
\pgfpathlineto{\pgfqpoint{0.484440in}{1.806255in}}%
\pgfpathlineto{\pgfqpoint{0.473604in}{1.810504in}}%
\pgfpathlineto{\pgfqpoint{0.465727in}{1.809963in}}%
\pgfpathlineto{\pgfqpoint{0.463179in}{1.820760in}}%
\pgfpathlineto{\pgfqpoint{0.471026in}{1.819611in}}%
\pgfpathlineto{\pgfqpoint{0.470402in}{1.834619in}}%
\pgfpathlineto{\pgfqpoint{0.466982in}{1.843530in}}%
\pgfpathlineto{\pgfqpoint{0.468490in}{1.854935in}}%
\pgfpathlineto{\pgfqpoint{0.472595in}{1.863514in}}%
\pgfpathlineto{\pgfqpoint{0.472364in}{1.879792in}}%
\pgfpathlineto{\pgfqpoint{0.470172in}{1.885699in}}%
\pgfpathlineto{\pgfqpoint{0.473957in}{1.904648in}}%
\pgfpathlineto{\pgfqpoint{0.472865in}{1.916860in}}%
\pgfpathlineto{\pgfqpoint{0.469081in}{1.922716in}}%
\pgfpathlineto{\pgfqpoint{0.469658in}{1.941861in}}%
\pgfpathlineto{\pgfqpoint{0.475162in}{1.955941in}}%
\pgfpathlineto{\pgfqpoint{0.481149in}{1.952495in}}%
\pgfpathlineto{\pgfqpoint{0.500972in}{1.932018in}}%
\pgfpathlineto{\pgfqpoint{0.524993in}{1.920812in}}%
\pgfpathlineto{\pgfqpoint{0.537292in}{1.919486in}}%
\pgfpathlineto{\pgfqpoint{0.548880in}{1.914715in}}%
\pgfpathlineto{\pgfqpoint{0.552168in}{1.898661in}}%
\pgfpathlineto{\pgfqpoint{0.542014in}{1.895620in}}%
\pgfpathlineto{\pgfqpoint{0.535285in}{1.882967in}}%
\pgfpathlineto{\pgfqpoint{0.543185in}{1.883670in}}%
\pgfpathlineto{\pgfqpoint{0.546372in}{1.889349in}}%
\pgfpathlineto{\pgfqpoint{0.557514in}{1.896451in}}%
\pgfpathlineto{\pgfqpoint{0.556939in}{1.885905in}}%
\pgfpathlineto{\pgfqpoint{0.549499in}{1.884204in}}%
\pgfpathlineto{\pgfqpoint{0.550734in}{1.870648in}}%
\pgfpathlineto{\pgfqpoint{0.541674in}{1.857069in}}%
\pgfpathlineto{\pgfqpoint{0.535876in}{1.863526in}}%
\pgfpathlineto{\pgfqpoint{0.518645in}{1.860015in}}%
\pgfpathlineto{\pgfqpoint{0.517738in}{1.852116in}}%
\pgfpathlineto{\pgfqpoint{0.523678in}{1.847303in}}%
\pgfpathlineto{\pgfqpoint{0.532741in}{1.846830in}}%
\pgfpathlineto{\pgfqpoint{0.545271in}{1.857114in}}%
\pgfpathlineto{\pgfqpoint{0.555189in}{1.857968in}}%
\pgfpathlineto{\pgfqpoint{0.555569in}{1.869352in}}%
\pgfpathlineto{\pgfqpoint{0.560680in}{1.886072in}}%
\pgfpathlineto{\pgfqpoint{0.572793in}{1.898506in}}%
\pgfpathlineto{\pgfqpoint{0.568754in}{1.908138in}}%
\pgfpathlineto{\pgfqpoint{0.571513in}{1.918507in}}%
\pgfpathlineto{\pgfqpoint{0.566139in}{1.928842in}}%
\pgfpathlineto{\pgfqpoint{0.573798in}{1.942057in}}%
\pgfpathlineto{\pgfqpoint{0.574945in}{1.950008in}}%
\pgfpathlineto{\pgfqpoint{0.568296in}{1.955234in}}%
\pgfpathlineto{\pgfqpoint{0.569403in}{1.968600in}}%
\pgfpathlineto{\pgfqpoint{0.649072in}{1.944395in}}%
\pgfpathlineto{\pgfqpoint{0.733673in}{1.920671in}}%
\pgfpathlineto{\pgfqpoint{0.810222in}{1.900931in}}%
\pgfpathclose%
\pgfusepath{fill}%
\end{pgfscope}%
\begin{pgfscope}%
\pgfpathrectangle{\pgfqpoint{0.100000in}{0.100000in}}{\pgfqpoint{2.989028in}{1.913466in}}%
\pgfusepath{clip}%
\pgfsetbuttcap%
\pgfsetmiterjoin%
\definecolor{currentfill}{rgb}{0.352018,0.707266,0.669204}%
\pgfsetfillcolor{currentfill}%
\pgfsetlinewidth{0.000000pt}%
\definecolor{currentstroke}{rgb}{0.000000,0.000000,0.000000}%
\pgfsetstrokecolor{currentstroke}%
\pgfsetstrokeopacity{0.000000}%
\pgfsetdash{}{0pt}%
\pgfpathmoveto{\pgfqpoint{0.556297in}{1.922359in}}%
\pgfpathlineto{\pgfqpoint{0.560220in}{1.907235in}}%
\pgfpathlineto{\pgfqpoint{0.566372in}{1.902211in}}%
\pgfpathlineto{\pgfqpoint{0.561463in}{1.895816in}}%
\pgfpathlineto{\pgfqpoint{0.556679in}{1.905186in}}%
\pgfpathlineto{\pgfqpoint{0.556297in}{1.922359in}}%
\pgfpathclose%
\pgfusepath{fill}%
\end{pgfscope}%
\begin{pgfscope}%
\pgfpathrectangle{\pgfqpoint{0.100000in}{0.100000in}}{\pgfqpoint{2.989028in}{1.913466in}}%
\pgfusepath{clip}%
\pgfsetbuttcap%
\pgfsetmiterjoin%
\definecolor{currentfill}{rgb}{0.495502,0.798155,0.645675}%
\pgfsetfillcolor{currentfill}%
\pgfsetlinewidth{0.000000pt}%
\definecolor{currentstroke}{rgb}{0.000000,0.000000,0.000000}%
\pgfsetstrokecolor{currentstroke}%
\pgfsetstrokeopacity{0.000000}%
\pgfsetdash{}{0pt}%
\pgfpathmoveto{\pgfqpoint{0.851305in}{1.891004in}}%
\pgfpathlineto{\pgfqpoint{0.936467in}{1.871893in}}%
\pgfpathlineto{\pgfqpoint{1.016670in}{1.855681in}}%
\pgfpathlineto{\pgfqpoint{1.078378in}{1.844375in}}%
\pgfpathlineto{\pgfqpoint{1.132175in}{1.835337in}}%
\pgfpathlineto{\pgfqpoint{1.186092in}{1.827040in}}%
\pgfpathlineto{\pgfqpoint{1.232007in}{1.820572in}}%
\pgfpathlineto{\pgfqpoint{1.277995in}{1.814641in}}%
\pgfpathlineto{\pgfqpoint{1.324048in}{1.809248in}}%
\pgfpathlineto{\pgfqpoint{1.367448in}{1.804664in}}%
\pgfpathlineto{\pgfqpoint{1.361479in}{1.738522in}}%
\pgfpathlineto{\pgfqpoint{1.352595in}{1.649140in}}%
\pgfpathlineto{\pgfqpoint{1.347936in}{1.603169in}}%
\pgfpathlineto{\pgfqpoint{1.341234in}{1.541215in}}%
\pgfpathlineto{\pgfqpoint{1.293794in}{1.546395in}}%
\pgfpathlineto{\pgfqpoint{1.239479in}{1.552561in}}%
\pgfpathlineto{\pgfqpoint{1.164064in}{1.562882in}}%
\pgfpathlineto{\pgfqpoint{1.130381in}{1.567687in}}%
\pgfpathlineto{\pgfqpoint{1.047396in}{1.580713in}}%
\pgfpathlineto{\pgfqpoint{1.018829in}{1.585939in}}%
\pgfpathlineto{\pgfqpoint{1.012875in}{1.552087in}}%
\pgfpathlineto{\pgfqpoint{1.009620in}{1.554502in}}%
\pgfpathlineto{\pgfqpoint{1.003502in}{1.570788in}}%
\pgfpathlineto{\pgfqpoint{0.995639in}{1.569785in}}%
\pgfpathlineto{\pgfqpoint{0.990011in}{1.561552in}}%
\pgfpathlineto{\pgfqpoint{0.978331in}{1.560451in}}%
\pgfpathlineto{\pgfqpoint{0.975984in}{1.564390in}}%
\pgfpathlineto{\pgfqpoint{0.965043in}{1.563913in}}%
\pgfpathlineto{\pgfqpoint{0.959543in}{1.567727in}}%
\pgfpathlineto{\pgfqpoint{0.951852in}{1.561826in}}%
\pgfpathlineto{\pgfqpoint{0.933172in}{1.567136in}}%
\pgfpathlineto{\pgfqpoint{0.925020in}{1.564285in}}%
\pgfpathlineto{\pgfqpoint{0.920634in}{1.576938in}}%
\pgfpathlineto{\pgfqpoint{0.920429in}{1.588978in}}%
\pgfpathlineto{\pgfqpoint{0.907552in}{1.597637in}}%
\pgfpathlineto{\pgfqpoint{0.909652in}{1.610116in}}%
\pgfpathlineto{\pgfqpoint{0.900483in}{1.630729in}}%
\pgfpathlineto{\pgfqpoint{0.901101in}{1.648251in}}%
\pgfpathlineto{\pgfqpoint{0.893220in}{1.656827in}}%
\pgfpathlineto{\pgfqpoint{0.885827in}{1.649240in}}%
\pgfpathlineto{\pgfqpoint{0.875763in}{1.645136in}}%
\pgfpathlineto{\pgfqpoint{0.868102in}{1.653594in}}%
\pgfpathlineto{\pgfqpoint{0.869348in}{1.667101in}}%
\pgfpathlineto{\pgfqpoint{0.878535in}{1.672422in}}%
\pgfpathlineto{\pgfqpoint{0.876093in}{1.680980in}}%
\pgfpathlineto{\pgfqpoint{0.892063in}{1.722399in}}%
\pgfpathlineto{\pgfqpoint{0.879470in}{1.723437in}}%
\pgfpathlineto{\pgfqpoint{0.878176in}{1.730870in}}%
\pgfpathlineto{\pgfqpoint{0.869003in}{1.737279in}}%
\pgfpathlineto{\pgfqpoint{0.869589in}{1.744445in}}%
\pgfpathlineto{\pgfqpoint{0.864750in}{1.750011in}}%
\pgfpathlineto{\pgfqpoint{0.856873in}{1.770222in}}%
\pgfpathlineto{\pgfqpoint{0.849675in}{1.774343in}}%
\pgfpathlineto{\pgfqpoint{0.844521in}{1.791786in}}%
\pgfpathlineto{\pgfqpoint{0.845737in}{1.803380in}}%
\pgfpathlineto{\pgfqpoint{0.836121in}{1.824733in}}%
\pgfpathlineto{\pgfqpoint{0.851305in}{1.891004in}}%
\pgfpathclose%
\pgfusepath{fill}%
\end{pgfscope}%
\begin{pgfscope}%
\pgfpathrectangle{\pgfqpoint{0.100000in}{0.100000in}}{\pgfqpoint{2.989028in}{1.913466in}}%
\pgfusepath{clip}%
\pgfsetbuttcap%
\pgfsetmiterjoin%
\definecolor{currentfill}{rgb}{0.384006,0.742945,0.654441}%
\pgfsetfillcolor{currentfill}%
\pgfsetlinewidth{0.000000pt}%
\definecolor{currentstroke}{rgb}{0.000000,0.000000,0.000000}%
\pgfsetstrokecolor{currentstroke}%
\pgfsetstrokeopacity{0.000000}%
\pgfsetdash{}{0pt}%
\pgfpathmoveto{\pgfqpoint{2.907175in}{1.550163in}}%
\pgfpathlineto{\pgfqpoint{2.905471in}{1.557411in}}%
\pgfpathlineto{\pgfqpoint{2.895998in}{1.563760in}}%
\pgfpathlineto{\pgfqpoint{2.870964in}{1.645645in}}%
\pgfpathlineto{\pgfqpoint{2.857471in}{1.685161in}}%
\pgfpathlineto{\pgfqpoint{2.868319in}{1.696546in}}%
\pgfpathlineto{\pgfqpoint{2.878864in}{1.724407in}}%
\pgfpathlineto{\pgfqpoint{2.884123in}{1.733311in}}%
\pgfpathlineto{\pgfqpoint{2.880397in}{1.737050in}}%
\pgfpathlineto{\pgfqpoint{2.879149in}{1.761782in}}%
\pgfpathlineto{\pgfqpoint{2.883854in}{1.769364in}}%
\pgfpathlineto{\pgfqpoint{2.881792in}{1.786986in}}%
\pgfpathlineto{\pgfqpoint{2.900570in}{1.844804in}}%
\pgfpathlineto{\pgfqpoint{2.909031in}{1.845120in}}%
\pgfpathlineto{\pgfqpoint{2.912543in}{1.834792in}}%
\pgfpathlineto{\pgfqpoint{2.920062in}{1.831826in}}%
\pgfpathlineto{\pgfqpoint{2.934266in}{1.843974in}}%
\pgfpathlineto{\pgfqpoint{2.945403in}{1.851149in}}%
\pgfpathlineto{\pgfqpoint{2.969943in}{1.838514in}}%
\pgfpathlineto{\pgfqpoint{2.992113in}{1.768360in}}%
\pgfpathlineto{\pgfqpoint{2.996328in}{1.751101in}}%
\pgfpathlineto{\pgfqpoint{3.006039in}{1.749076in}}%
\pgfpathlineto{\pgfqpoint{3.019209in}{1.737354in}}%
\pgfpathlineto{\pgfqpoint{3.018463in}{1.730524in}}%
\pgfpathlineto{\pgfqpoint{3.027435in}{1.722428in}}%
\pgfpathlineto{\pgfqpoint{3.034723in}{1.727073in}}%
\pgfpathlineto{\pgfqpoint{3.049974in}{1.709287in}}%
\pgfpathlineto{\pgfqpoint{3.043158in}{1.695053in}}%
\pgfpathlineto{\pgfqpoint{3.034018in}{1.694771in}}%
\pgfpathlineto{\pgfqpoint{3.026685in}{1.682065in}}%
\pgfpathlineto{\pgfqpoint{3.017802in}{1.680263in}}%
\pgfpathlineto{\pgfqpoint{3.011280in}{1.673497in}}%
\pgfpathlineto{\pgfqpoint{2.991803in}{1.666241in}}%
\pgfpathlineto{\pgfqpoint{2.980014in}{1.654548in}}%
\pgfpathlineto{\pgfqpoint{2.974025in}{1.662941in}}%
\pgfpathlineto{\pgfqpoint{2.968584in}{1.656926in}}%
\pgfpathlineto{\pgfqpoint{2.970074in}{1.632607in}}%
\pgfpathlineto{\pgfqpoint{2.965765in}{1.622974in}}%
\pgfpathlineto{\pgfqpoint{2.956367in}{1.625605in}}%
\pgfpathlineto{\pgfqpoint{2.954850in}{1.615730in}}%
\pgfpathlineto{\pgfqpoint{2.950789in}{1.611669in}}%
\pgfpathlineto{\pgfqpoint{2.942662in}{1.612659in}}%
\pgfpathlineto{\pgfqpoint{2.943174in}{1.603425in}}%
\pgfpathlineto{\pgfqpoint{2.930855in}{1.606048in}}%
\pgfpathlineto{\pgfqpoint{2.924128in}{1.593255in}}%
\pgfpathlineto{\pgfqpoint{2.926668in}{1.586561in}}%
\pgfpathlineto{\pgfqpoint{2.922666in}{1.575429in}}%
\pgfpathlineto{\pgfqpoint{2.916367in}{1.567240in}}%
\pgfpathlineto{\pgfqpoint{2.914781in}{1.550146in}}%
\pgfpathlineto{\pgfqpoint{2.907175in}{1.550163in}}%
\pgfpathclose%
\pgfusepath{fill}%
\end{pgfscope}%
\begin{pgfscope}%
\pgfpathrectangle{\pgfqpoint{0.100000in}{0.100000in}}{\pgfqpoint{2.989028in}{1.913466in}}%
\pgfusepath{clip}%
\pgfsetbuttcap%
\pgfsetmiterjoin%
\definecolor{currentfill}{rgb}{0.384006,0.742945,0.654441}%
\pgfsetfillcolor{currentfill}%
\pgfsetlinewidth{0.000000pt}%
\definecolor{currentstroke}{rgb}{0.000000,0.000000,0.000000}%
\pgfsetstrokecolor{currentstroke}%
\pgfsetstrokeopacity{0.000000}%
\pgfsetdash{}{0pt}%
\pgfpathmoveto{\pgfqpoint{2.995301in}{1.661247in}}%
\pgfpathlineto{\pgfqpoint{3.000863in}{1.667073in}}%
\pgfpathlineto{\pgfqpoint{3.006158in}{1.661597in}}%
\pgfpathlineto{\pgfqpoint{2.996647in}{1.654342in}}%
\pgfpathlineto{\pgfqpoint{2.995301in}{1.661247in}}%
\pgfpathclose%
\pgfusepath{fill}%
\end{pgfscope}%
\begin{pgfscope}%
\pgfpathrectangle{\pgfqpoint{0.100000in}{0.100000in}}{\pgfqpoint{2.989028in}{1.913466in}}%
\pgfusepath{clip}%
\pgfsetbuttcap%
\pgfsetmiterjoin%
\definecolor{currentfill}{rgb}{0.711419,0.883276,0.634833}%
\pgfsetfillcolor{currentfill}%
\pgfsetlinewidth{0.000000pt}%
\definecolor{currentstroke}{rgb}{0.000000,0.000000,0.000000}%
\pgfsetstrokecolor{currentstroke}%
\pgfsetstrokeopacity{0.000000}%
\pgfsetdash{}{0pt}%
\pgfpathmoveto{\pgfqpoint{1.347936in}{1.603169in}}%
\pgfpathlineto{\pgfqpoint{1.352595in}{1.649140in}}%
\pgfpathlineto{\pgfqpoint{1.361479in}{1.738522in}}%
\pgfpathlineto{\pgfqpoint{1.367448in}{1.804664in}}%
\pgfpathlineto{\pgfqpoint{1.416329in}{1.800078in}}%
\pgfpathlineto{\pgfqpoint{1.478862in}{1.795100in}}%
\pgfpathlineto{\pgfqpoint{1.536021in}{1.791418in}}%
\pgfpathlineto{\pgfqpoint{1.587777in}{1.788799in}}%
\pgfpathlineto{\pgfqpoint{1.665014in}{1.786150in}}%
\pgfpathlineto{\pgfqpoint{1.670285in}{1.764381in}}%
\pgfpathlineto{\pgfqpoint{1.668398in}{1.753978in}}%
\pgfpathlineto{\pgfqpoint{1.667777in}{1.732596in}}%
\pgfpathlineto{\pgfqpoint{1.671363in}{1.716582in}}%
\pgfpathlineto{\pgfqpoint{1.679579in}{1.692956in}}%
\pgfpathlineto{\pgfqpoint{1.679560in}{1.663212in}}%
\pgfpathlineto{\pgfqpoint{1.681080in}{1.628658in}}%
\pgfpathlineto{\pgfqpoint{1.683167in}{1.619346in}}%
\pgfpathlineto{\pgfqpoint{1.689274in}{1.609122in}}%
\pgfpathlineto{\pgfqpoint{1.691288in}{1.593192in}}%
\pgfpathlineto{\pgfqpoint{1.690419in}{1.582557in}}%
\pgfpathlineto{\pgfqpoint{1.625668in}{1.583959in}}%
\pgfpathlineto{\pgfqpoint{1.578564in}{1.586347in}}%
\pgfpathlineto{\pgfqpoint{1.509511in}{1.590023in}}%
\pgfpathlineto{\pgfqpoint{1.441415in}{1.594860in}}%
\pgfpathlineto{\pgfqpoint{1.396059in}{1.598519in}}%
\pgfpathlineto{\pgfqpoint{1.347936in}{1.603169in}}%
\pgfpathclose%
\pgfusepath{fill}%
\end{pgfscope}%
\begin{pgfscope}%
\pgfpathrectangle{\pgfqpoint{0.100000in}{0.100000in}}{\pgfqpoint{2.989028in}{1.913466in}}%
\pgfusepath{clip}%
\pgfsetbuttcap%
\pgfsetmiterjoin%
\definecolor{currentfill}{rgb}{0.344022,0.698347,0.672895}%
\pgfsetfillcolor{currentfill}%
\pgfsetlinewidth{0.000000pt}%
\definecolor{currentstroke}{rgb}{0.000000,0.000000,0.000000}%
\pgfsetstrokecolor{currentstroke}%
\pgfsetstrokeopacity{0.000000}%
\pgfsetdash{}{0pt}%
\pgfpathmoveto{\pgfqpoint{1.328363in}{1.410780in}}%
\pgfpathlineto{\pgfqpoint{1.335839in}{1.487840in}}%
\pgfpathlineto{\pgfqpoint{1.341234in}{1.541215in}}%
\pgfpathlineto{\pgfqpoint{1.347936in}{1.603169in}}%
\pgfpathlineto{\pgfqpoint{1.396059in}{1.598519in}}%
\pgfpathlineto{\pgfqpoint{1.441415in}{1.594860in}}%
\pgfpathlineto{\pgfqpoint{1.509511in}{1.590023in}}%
\pgfpathlineto{\pgfqpoint{1.578564in}{1.586347in}}%
\pgfpathlineto{\pgfqpoint{1.625668in}{1.583959in}}%
\pgfpathlineto{\pgfqpoint{1.690419in}{1.582557in}}%
\pgfpathlineto{\pgfqpoint{1.686034in}{1.569764in}}%
\pgfpathlineto{\pgfqpoint{1.677284in}{1.559721in}}%
\pgfpathlineto{\pgfqpoint{1.683985in}{1.548158in}}%
\pgfpathlineto{\pgfqpoint{1.691374in}{1.545688in}}%
\pgfpathlineto{\pgfqpoint{1.694880in}{1.539046in}}%
\pgfpathlineto{\pgfqpoint{1.694075in}{1.490552in}}%
\pgfpathlineto{\pgfqpoint{1.692728in}{1.422304in}}%
\pgfpathlineto{\pgfqpoint{1.687744in}{1.404379in}}%
\pgfpathlineto{\pgfqpoint{1.692201in}{1.394464in}}%
\pgfpathlineto{\pgfqpoint{1.683252in}{1.373157in}}%
\pgfpathlineto{\pgfqpoint{1.692680in}{1.355981in}}%
\pgfpathlineto{\pgfqpoint{1.684670in}{1.357295in}}%
\pgfpathlineto{\pgfqpoint{1.679205in}{1.367982in}}%
\pgfpathlineto{\pgfqpoint{1.659700in}{1.375307in}}%
\pgfpathlineto{\pgfqpoint{1.647400in}{1.381750in}}%
\pgfpathlineto{\pgfqpoint{1.626759in}{1.382286in}}%
\pgfpathlineto{\pgfqpoint{1.619594in}{1.376416in}}%
\pgfpathlineto{\pgfqpoint{1.596231in}{1.387975in}}%
\pgfpathlineto{\pgfqpoint{1.594439in}{1.391630in}}%
\pgfpathlineto{\pgfqpoint{1.512902in}{1.395489in}}%
\pgfpathlineto{\pgfqpoint{1.463365in}{1.398329in}}%
\pgfpathlineto{\pgfqpoint{1.422447in}{1.401533in}}%
\pgfpathlineto{\pgfqpoint{1.354836in}{1.407921in}}%
\pgfpathlineto{\pgfqpoint{1.328363in}{1.410780in}}%
\pgfpathclose%
\pgfusepath{fill}%
\end{pgfscope}%
\begin{pgfscope}%
\pgfpathrectangle{\pgfqpoint{0.100000in}{0.100000in}}{\pgfqpoint{2.989028in}{1.913466in}}%
\pgfusepath{clip}%
\pgfsetbuttcap%
\pgfsetmiterjoin%
\definecolor{currentfill}{rgb}{0.280046,0.626990,0.702422}%
\pgfsetfillcolor{currentfill}%
\pgfsetlinewidth{0.000000pt}%
\definecolor{currentstroke}{rgb}{0.000000,0.000000,0.000000}%
\pgfsetstrokecolor{currentstroke}%
\pgfsetstrokeopacity{0.000000}%
\pgfsetdash{}{0pt}%
\pgfpathmoveto{\pgfqpoint{1.315540in}{1.280295in}}%
\pgfpathlineto{\pgfqpoint{1.272072in}{1.284265in}}%
\pgfpathlineto{\pgfqpoint{1.177321in}{1.295959in}}%
\pgfpathlineto{\pgfqpoint{1.125778in}{1.303357in}}%
\pgfpathlineto{\pgfqpoint{1.070495in}{1.311293in}}%
\pgfpathlineto{\pgfqpoint{1.023928in}{1.318744in}}%
\pgfpathlineto{\pgfqpoint{0.972815in}{1.327496in}}%
\pgfpathlineto{\pgfqpoint{0.984436in}{1.391963in}}%
\pgfpathlineto{\pgfqpoint{0.996208in}{1.458068in}}%
\pgfpathlineto{\pgfqpoint{1.012875in}{1.552087in}}%
\pgfpathlineto{\pgfqpoint{1.018829in}{1.585939in}}%
\pgfpathlineto{\pgfqpoint{1.047396in}{1.580713in}}%
\pgfpathlineto{\pgfqpoint{1.130381in}{1.567687in}}%
\pgfpathlineto{\pgfqpoint{1.164064in}{1.562882in}}%
\pgfpathlineto{\pgfqpoint{1.239479in}{1.552561in}}%
\pgfpathlineto{\pgfqpoint{1.293794in}{1.546395in}}%
\pgfpathlineto{\pgfqpoint{1.341234in}{1.541215in}}%
\pgfpathlineto{\pgfqpoint{1.335839in}{1.487840in}}%
\pgfpathlineto{\pgfqpoint{1.328363in}{1.410780in}}%
\pgfpathlineto{\pgfqpoint{1.321947in}{1.345289in}}%
\pgfpathlineto{\pgfqpoint{1.315540in}{1.280295in}}%
\pgfpathclose%
\pgfusepath{fill}%
\end{pgfscope}%
\begin{pgfscope}%
\pgfpathrectangle{\pgfqpoint{0.100000in}{0.100000in}}{\pgfqpoint{2.989028in}{1.913466in}}%
\pgfusepath{clip}%
\pgfsetbuttcap%
\pgfsetmiterjoin%
\definecolor{currentfill}{rgb}{0.368012,0.725106,0.661822}%
\pgfsetfillcolor{currentfill}%
\pgfsetlinewidth{0.000000pt}%
\definecolor{currentstroke}{rgb}{0.000000,0.000000,0.000000}%
\pgfsetstrokecolor{currentstroke}%
\pgfsetstrokeopacity{0.000000}%
\pgfsetdash{}{0pt}%
\pgfpathmoveto{\pgfqpoint{2.110291in}{1.369221in}}%
\pgfpathlineto{\pgfqpoint{2.055206in}{1.365308in}}%
\pgfpathlineto{\pgfqpoint{1.973100in}{1.361809in}}%
\pgfpathlineto{\pgfqpoint{1.969980in}{1.370103in}}%
\pgfpathlineto{\pgfqpoint{1.951772in}{1.376303in}}%
\pgfpathlineto{\pgfqpoint{1.947756in}{1.388016in}}%
\pgfpathlineto{\pgfqpoint{1.946075in}{1.402507in}}%
\pgfpathlineto{\pgfqpoint{1.950178in}{1.409931in}}%
\pgfpathlineto{\pgfqpoint{1.943697in}{1.417056in}}%
\pgfpathlineto{\pgfqpoint{1.942134in}{1.425554in}}%
\pgfpathlineto{\pgfqpoint{1.940045in}{1.444352in}}%
\pgfpathlineto{\pgfqpoint{1.933835in}{1.454564in}}%
\pgfpathlineto{\pgfqpoint{1.922805in}{1.460295in}}%
\pgfpathlineto{\pgfqpoint{1.910803in}{1.469752in}}%
\pgfpathlineto{\pgfqpoint{1.904599in}{1.480942in}}%
\pgfpathlineto{\pgfqpoint{1.893466in}{1.485448in}}%
\pgfpathlineto{\pgfqpoint{1.886922in}{1.492782in}}%
\pgfpathlineto{\pgfqpoint{1.878994in}{1.494027in}}%
\pgfpathlineto{\pgfqpoint{1.864840in}{1.504916in}}%
\pgfpathlineto{\pgfqpoint{1.867139in}{1.517447in}}%
\pgfpathlineto{\pgfqpoint{1.866696in}{1.541277in}}%
\pgfpathlineto{\pgfqpoint{1.871061in}{1.547836in}}%
\pgfpathlineto{\pgfqpoint{1.867139in}{1.557741in}}%
\pgfpathlineto{\pgfqpoint{1.860230in}{1.559657in}}%
\pgfpathlineto{\pgfqpoint{1.860801in}{1.568355in}}%
\pgfpathlineto{\pgfqpoint{1.869373in}{1.582098in}}%
\pgfpathlineto{\pgfqpoint{1.886348in}{1.592965in}}%
\pgfpathlineto{\pgfqpoint{1.885290in}{1.631611in}}%
\pgfpathlineto{\pgfqpoint{1.893784in}{1.637416in}}%
\pgfpathlineto{\pgfqpoint{1.901822in}{1.633549in}}%
\pgfpathlineto{\pgfqpoint{1.918197in}{1.639211in}}%
\pgfpathlineto{\pgfqpoint{1.949004in}{1.653439in}}%
\pgfpathlineto{\pgfqpoint{1.953017in}{1.649033in}}%
\pgfpathlineto{\pgfqpoint{1.947182in}{1.629079in}}%
\pgfpathlineto{\pgfqpoint{1.955863in}{1.633458in}}%
\pgfpathlineto{\pgfqpoint{1.970731in}{1.629085in}}%
\pgfpathlineto{\pgfqpoint{1.979867in}{1.625436in}}%
\pgfpathlineto{\pgfqpoint{1.984989in}{1.614740in}}%
\pgfpathlineto{\pgfqpoint{2.031860in}{1.604646in}}%
\pgfpathlineto{\pgfqpoint{2.045871in}{1.597722in}}%
\pgfpathlineto{\pgfqpoint{2.060116in}{1.597803in}}%
\pgfpathlineto{\pgfqpoint{2.074752in}{1.594947in}}%
\pgfpathlineto{\pgfqpoint{2.084251in}{1.585189in}}%
\pgfpathlineto{\pgfqpoint{2.093328in}{1.580368in}}%
\pgfpathlineto{\pgfqpoint{2.094996in}{1.566462in}}%
\pgfpathlineto{\pgfqpoint{2.092314in}{1.557730in}}%
\pgfpathlineto{\pgfqpoint{2.102430in}{1.557084in}}%
\pgfpathlineto{\pgfqpoint{2.099022in}{1.546957in}}%
\pgfpathlineto{\pgfqpoint{2.102264in}{1.543360in}}%
\pgfpathlineto{\pgfqpoint{2.105488in}{1.533800in}}%
\pgfpathlineto{\pgfqpoint{2.095629in}{1.528739in}}%
\pgfpathlineto{\pgfqpoint{2.089905in}{1.514637in}}%
\pgfpathlineto{\pgfqpoint{2.088109in}{1.504665in}}%
\pgfpathlineto{\pgfqpoint{2.093600in}{1.502958in}}%
\pgfpathlineto{\pgfqpoint{2.100632in}{1.510436in}}%
\pgfpathlineto{\pgfqpoint{2.106610in}{1.523394in}}%
\pgfpathlineto{\pgfqpoint{2.114698in}{1.527895in}}%
\pgfpathlineto{\pgfqpoint{2.120776in}{1.522041in}}%
\pgfpathlineto{\pgfqpoint{2.114793in}{1.504322in}}%
\pgfpathlineto{\pgfqpoint{2.112916in}{1.490567in}}%
\pgfpathlineto{\pgfqpoint{2.114684in}{1.480679in}}%
\pgfpathlineto{\pgfqpoint{2.109128in}{1.475085in}}%
\pgfpathlineto{\pgfqpoint{2.106353in}{1.461416in}}%
\pgfpathlineto{\pgfqpoint{2.108620in}{1.447050in}}%
\pgfpathlineto{\pgfqpoint{2.102064in}{1.425804in}}%
\pgfpathlineto{\pgfqpoint{2.102202in}{1.415188in}}%
\pgfpathlineto{\pgfqpoint{2.107383in}{1.392182in}}%
\pgfpathlineto{\pgfqpoint{2.110741in}{1.388236in}}%
\pgfpathlineto{\pgfqpoint{2.110291in}{1.369221in}}%
\pgfpathclose%
\pgfusepath{fill}%
\end{pgfscope}%
\begin{pgfscope}%
\pgfpathrectangle{\pgfqpoint{0.100000in}{0.100000in}}{\pgfqpoint{2.989028in}{1.913466in}}%
\pgfusepath{clip}%
\pgfsetbuttcap%
\pgfsetmiterjoin%
\definecolor{currentfill}{rgb}{0.368012,0.725106,0.661822}%
\pgfsetfillcolor{currentfill}%
\pgfsetlinewidth{0.000000pt}%
\definecolor{currentstroke}{rgb}{0.000000,0.000000,0.000000}%
\pgfsetstrokecolor{currentstroke}%
\pgfsetstrokeopacity{0.000000}%
\pgfsetdash{}{0pt}%
\pgfpathmoveto{\pgfqpoint{2.130944in}{1.555675in}}%
\pgfpathlineto{\pgfqpoint{2.132173in}{1.546498in}}%
\pgfpathlineto{\pgfqpoint{2.120909in}{1.522320in}}%
\pgfpathlineto{\pgfqpoint{2.115904in}{1.529323in}}%
\pgfpathlineto{\pgfqpoint{2.130944in}{1.555675in}}%
\pgfpathclose%
\pgfusepath{fill}%
\end{pgfscope}%
\begin{pgfscope}%
\pgfpathrectangle{\pgfqpoint{0.100000in}{0.100000in}}{\pgfqpoint{2.989028in}{1.913466in}}%
\pgfusepath{clip}%
\pgfsetbuttcap%
\pgfsetmiterjoin%
\definecolor{currentfill}{rgb}{0.267128,0.441292,0.697578}%
\pgfsetfillcolor{currentfill}%
\pgfsetlinewidth{0.000000pt}%
\definecolor{currentstroke}{rgb}{0.000000,0.000000,0.000000}%
\pgfsetstrokecolor{currentstroke}%
\pgfsetstrokeopacity{0.000000}%
\pgfsetdash{}{0pt}%
\pgfpathmoveto{\pgfqpoint{0.767475in}{1.706523in}}%
\pgfpathlineto{\pgfqpoint{0.769703in}{1.717547in}}%
\pgfpathlineto{\pgfqpoint{0.768022in}{1.729516in}}%
\pgfpathlineto{\pgfqpoint{0.770595in}{1.742564in}}%
\pgfpathlineto{\pgfqpoint{0.783700in}{1.795504in}}%
\pgfpathlineto{\pgfqpoint{0.794722in}{1.839512in}}%
\pgfpathlineto{\pgfqpoint{0.810222in}{1.900931in}}%
\pgfpathlineto{\pgfqpoint{0.851305in}{1.891004in}}%
\pgfpathlineto{\pgfqpoint{0.836121in}{1.824733in}}%
\pgfpathlineto{\pgfqpoint{0.845737in}{1.803380in}}%
\pgfpathlineto{\pgfqpoint{0.844521in}{1.791786in}}%
\pgfpathlineto{\pgfqpoint{0.849675in}{1.774343in}}%
\pgfpathlineto{\pgfqpoint{0.856873in}{1.770222in}}%
\pgfpathlineto{\pgfqpoint{0.864750in}{1.750011in}}%
\pgfpathlineto{\pgfqpoint{0.869589in}{1.744445in}}%
\pgfpathlineto{\pgfqpoint{0.869003in}{1.737279in}}%
\pgfpathlineto{\pgfqpoint{0.878176in}{1.730870in}}%
\pgfpathlineto{\pgfqpoint{0.879470in}{1.723437in}}%
\pgfpathlineto{\pgfqpoint{0.892063in}{1.722399in}}%
\pgfpathlineto{\pgfqpoint{0.876093in}{1.680980in}}%
\pgfpathlineto{\pgfqpoint{0.878535in}{1.672422in}}%
\pgfpathlineto{\pgfqpoint{0.869348in}{1.667101in}}%
\pgfpathlineto{\pgfqpoint{0.868102in}{1.653594in}}%
\pgfpathlineto{\pgfqpoint{0.875763in}{1.645136in}}%
\pgfpathlineto{\pgfqpoint{0.885827in}{1.649240in}}%
\pgfpathlineto{\pgfqpoint{0.893220in}{1.656827in}}%
\pgfpathlineto{\pgfqpoint{0.901101in}{1.648251in}}%
\pgfpathlineto{\pgfqpoint{0.900483in}{1.630729in}}%
\pgfpathlineto{\pgfqpoint{0.909652in}{1.610116in}}%
\pgfpathlineto{\pgfqpoint{0.907552in}{1.597637in}}%
\pgfpathlineto{\pgfqpoint{0.920429in}{1.588978in}}%
\pgfpathlineto{\pgfqpoint{0.920634in}{1.576938in}}%
\pgfpathlineto{\pgfqpoint{0.925020in}{1.564285in}}%
\pgfpathlineto{\pgfqpoint{0.933172in}{1.567136in}}%
\pgfpathlineto{\pgfqpoint{0.951852in}{1.561826in}}%
\pgfpathlineto{\pgfqpoint{0.959543in}{1.567727in}}%
\pgfpathlineto{\pgfqpoint{0.965043in}{1.563913in}}%
\pgfpathlineto{\pgfqpoint{0.975984in}{1.564390in}}%
\pgfpathlineto{\pgfqpoint{0.978331in}{1.560451in}}%
\pgfpathlineto{\pgfqpoint{0.990011in}{1.561552in}}%
\pgfpathlineto{\pgfqpoint{0.995639in}{1.569785in}}%
\pgfpathlineto{\pgfqpoint{1.003502in}{1.570788in}}%
\pgfpathlineto{\pgfqpoint{1.009620in}{1.554502in}}%
\pgfpathlineto{\pgfqpoint{1.012875in}{1.552087in}}%
\pgfpathlineto{\pgfqpoint{0.996208in}{1.458068in}}%
\pgfpathlineto{\pgfqpoint{0.984436in}{1.391963in}}%
\pgfpathlineto{\pgfqpoint{0.891592in}{1.409886in}}%
\pgfpathlineto{\pgfqpoint{0.841442in}{1.419862in}}%
\pgfpathlineto{\pgfqpoint{0.794534in}{1.430177in}}%
\pgfpathlineto{\pgfqpoint{0.750689in}{1.440103in}}%
\pgfpathlineto{\pgfqpoint{0.699946in}{1.452341in}}%
\pgfpathlineto{\pgfqpoint{0.726119in}{1.559712in}}%
\pgfpathlineto{\pgfqpoint{0.727531in}{1.567559in}}%
\pgfpathlineto{\pgfqpoint{0.739170in}{1.588120in}}%
\pgfpathlineto{\pgfqpoint{0.727118in}{1.600478in}}%
\pgfpathlineto{\pgfqpoint{0.729620in}{1.612613in}}%
\pgfpathlineto{\pgfqpoint{0.734208in}{1.615650in}}%
\pgfpathlineto{\pgfqpoint{0.742402in}{1.628155in}}%
\pgfpathlineto{\pgfqpoint{0.754345in}{1.636806in}}%
\pgfpathlineto{\pgfqpoint{0.754997in}{1.643207in}}%
\pgfpathlineto{\pgfqpoint{0.762206in}{1.649603in}}%
\pgfpathlineto{\pgfqpoint{0.768191in}{1.661575in}}%
\pgfpathlineto{\pgfqpoint{0.780469in}{1.674240in}}%
\pgfpathlineto{\pgfqpoint{0.779597in}{1.686749in}}%
\pgfpathlineto{\pgfqpoint{0.771203in}{1.693698in}}%
\pgfpathlineto{\pgfqpoint{0.767475in}{1.706523in}}%
\pgfpathclose%
\pgfusepath{fill}%
\end{pgfscope}%
\begin{pgfscope}%
\pgfpathrectangle{\pgfqpoint{0.100000in}{0.100000in}}{\pgfqpoint{2.989028in}{1.913466in}}%
\pgfusepath{clip}%
\pgfsetbuttcap%
\pgfsetmiterjoin%
\definecolor{currentfill}{rgb}{0.548558,0.818916,0.644906}%
\pgfsetfillcolor{currentfill}%
\pgfsetlinewidth{0.000000pt}%
\definecolor{currentstroke}{rgb}{0.000000,0.000000,0.000000}%
\pgfsetstrokecolor{currentstroke}%
\pgfsetstrokeopacity{0.000000}%
\pgfsetdash{}{0pt}%
\pgfpathmoveto{\pgfqpoint{2.799865in}{1.496989in}}%
\pgfpathlineto{\pgfqpoint{2.796007in}{1.509167in}}%
\pgfpathlineto{\pgfqpoint{2.788848in}{1.546108in}}%
\pgfpathlineto{\pgfqpoint{2.779593in}{1.557473in}}%
\pgfpathlineto{\pgfqpoint{2.771474in}{1.577911in}}%
\pgfpathlineto{\pgfqpoint{2.774140in}{1.593064in}}%
\pgfpathlineto{\pgfqpoint{2.772002in}{1.604233in}}%
\pgfpathlineto{\pgfqpoint{2.765065in}{1.615202in}}%
\pgfpathlineto{\pgfqpoint{2.764780in}{1.627285in}}%
\pgfpathlineto{\pgfqpoint{2.760752in}{1.640318in}}%
\pgfpathlineto{\pgfqpoint{2.796797in}{1.649189in}}%
\pgfpathlineto{\pgfqpoint{2.843616in}{1.661826in}}%
\pgfpathlineto{\pgfqpoint{2.845485in}{1.654580in}}%
\pgfpathlineto{\pgfqpoint{2.842504in}{1.643043in}}%
\pgfpathlineto{\pgfqpoint{2.849502in}{1.633812in}}%
\pgfpathlineto{\pgfqpoint{2.845785in}{1.622115in}}%
\pgfpathlineto{\pgfqpoint{2.831015in}{1.607411in}}%
\pgfpathlineto{\pgfqpoint{2.835075in}{1.596366in}}%
\pgfpathlineto{\pgfqpoint{2.832066in}{1.573021in}}%
\pgfpathlineto{\pgfqpoint{2.827854in}{1.559858in}}%
\pgfpathlineto{\pgfqpoint{2.833228in}{1.522426in}}%
\pgfpathlineto{\pgfqpoint{2.832504in}{1.509256in}}%
\pgfpathlineto{\pgfqpoint{2.837703in}{1.505030in}}%
\pgfpathlineto{\pgfqpoint{2.799865in}{1.496989in}}%
\pgfpathclose%
\pgfusepath{fill}%
\end{pgfscope}%
\begin{pgfscope}%
\pgfpathrectangle{\pgfqpoint{0.100000in}{0.100000in}}{\pgfqpoint{2.989028in}{1.913466in}}%
\pgfusepath{clip}%
\pgfsetbuttcap%
\pgfsetmiterjoin%
\definecolor{currentfill}{rgb}{0.516724,0.806459,0.645367}%
\pgfsetfillcolor{currentfill}%
\pgfsetlinewidth{0.000000pt}%
\definecolor{currentstroke}{rgb}{0.000000,0.000000,0.000000}%
\pgfsetstrokecolor{currentstroke}%
\pgfsetstrokeopacity{0.000000}%
\pgfsetdash{}{0pt}%
\pgfpathmoveto{\pgfqpoint{1.692728in}{1.422304in}}%
\pgfpathlineto{\pgfqpoint{1.694075in}{1.490552in}}%
\pgfpathlineto{\pgfqpoint{1.694880in}{1.539046in}}%
\pgfpathlineto{\pgfqpoint{1.691374in}{1.545688in}}%
\pgfpathlineto{\pgfqpoint{1.683985in}{1.548158in}}%
\pgfpathlineto{\pgfqpoint{1.677284in}{1.559721in}}%
\pgfpathlineto{\pgfqpoint{1.686034in}{1.569764in}}%
\pgfpathlineto{\pgfqpoint{1.690419in}{1.582557in}}%
\pgfpathlineto{\pgfqpoint{1.691288in}{1.593192in}}%
\pgfpathlineto{\pgfqpoint{1.689274in}{1.609122in}}%
\pgfpathlineto{\pgfqpoint{1.683167in}{1.619346in}}%
\pgfpathlineto{\pgfqpoint{1.681080in}{1.628658in}}%
\pgfpathlineto{\pgfqpoint{1.679560in}{1.663212in}}%
\pgfpathlineto{\pgfqpoint{1.679579in}{1.692956in}}%
\pgfpathlineto{\pgfqpoint{1.671363in}{1.716582in}}%
\pgfpathlineto{\pgfqpoint{1.667777in}{1.732596in}}%
\pgfpathlineto{\pgfqpoint{1.668398in}{1.753978in}}%
\pgfpathlineto{\pgfqpoint{1.670285in}{1.764381in}}%
\pgfpathlineto{\pgfqpoint{1.665014in}{1.786150in}}%
\pgfpathlineto{\pgfqpoint{1.700900in}{1.785431in}}%
\pgfpathlineto{\pgfqpoint{1.755411in}{1.784963in}}%
\pgfpathlineto{\pgfqpoint{1.755709in}{1.809705in}}%
\pgfpathlineto{\pgfqpoint{1.769584in}{1.806981in}}%
\pgfpathlineto{\pgfqpoint{1.776233in}{1.776812in}}%
\pgfpathlineto{\pgfqpoint{1.781136in}{1.765965in}}%
\pgfpathlineto{\pgfqpoint{1.793328in}{1.765645in}}%
\pgfpathlineto{\pgfqpoint{1.796056in}{1.761964in}}%
\pgfpathlineto{\pgfqpoint{1.813062in}{1.760334in}}%
\pgfpathlineto{\pgfqpoint{1.815920in}{1.752855in}}%
\pgfpathlineto{\pgfqpoint{1.827640in}{1.754536in}}%
\pgfpathlineto{\pgfqpoint{1.836742in}{1.761531in}}%
\pgfpathlineto{\pgfqpoint{1.852440in}{1.761270in}}%
\pgfpathlineto{\pgfqpoint{1.862153in}{1.755644in}}%
\pgfpathlineto{\pgfqpoint{1.863270in}{1.750368in}}%
\pgfpathlineto{\pgfqpoint{1.872512in}{1.749264in}}%
\pgfpathlineto{\pgfqpoint{1.878543in}{1.734873in}}%
\pgfpathlineto{\pgfqpoint{1.882434in}{1.743721in}}%
\pgfpathlineto{\pgfqpoint{1.893051in}{1.744265in}}%
\pgfpathlineto{\pgfqpoint{1.895736in}{1.737374in}}%
\pgfpathlineto{\pgfqpoint{1.907682in}{1.734228in}}%
\pgfpathlineto{\pgfqpoint{1.914270in}{1.724302in}}%
\pgfpathlineto{\pgfqpoint{1.928880in}{1.727384in}}%
\pgfpathlineto{\pgfqpoint{1.944957in}{1.739565in}}%
\pgfpathlineto{\pgfqpoint{1.950818in}{1.728827in}}%
\pgfpathlineto{\pgfqpoint{1.977198in}{1.731766in}}%
\pgfpathlineto{\pgfqpoint{1.988474in}{1.724397in}}%
\pgfpathlineto{\pgfqpoint{1.995037in}{1.727028in}}%
\pgfpathlineto{\pgfqpoint{2.000300in}{1.722879in}}%
\pgfpathlineto{\pgfqpoint{1.984687in}{1.713035in}}%
\pgfpathlineto{\pgfqpoint{1.962401in}{1.704265in}}%
\pgfpathlineto{\pgfqpoint{1.940329in}{1.686734in}}%
\pgfpathlineto{\pgfqpoint{1.921193in}{1.663680in}}%
\pgfpathlineto{\pgfqpoint{1.906707in}{1.650050in}}%
\pgfpathlineto{\pgfqpoint{1.894018in}{1.640682in}}%
\pgfpathlineto{\pgfqpoint{1.885290in}{1.631611in}}%
\pgfpathlineto{\pgfqpoint{1.886348in}{1.592965in}}%
\pgfpathlineto{\pgfqpoint{1.869373in}{1.582098in}}%
\pgfpathlineto{\pgfqpoint{1.860801in}{1.568355in}}%
\pgfpathlineto{\pgfqpoint{1.860230in}{1.559657in}}%
\pgfpathlineto{\pgfqpoint{1.867139in}{1.557741in}}%
\pgfpathlineto{\pgfqpoint{1.871061in}{1.547836in}}%
\pgfpathlineto{\pgfqpoint{1.866696in}{1.541277in}}%
\pgfpathlineto{\pgfqpoint{1.867139in}{1.517447in}}%
\pgfpathlineto{\pgfqpoint{1.864840in}{1.504916in}}%
\pgfpathlineto{\pgfqpoint{1.878994in}{1.494027in}}%
\pgfpathlineto{\pgfqpoint{1.886922in}{1.492782in}}%
\pgfpathlineto{\pgfqpoint{1.893466in}{1.485448in}}%
\pgfpathlineto{\pgfqpoint{1.904599in}{1.480942in}}%
\pgfpathlineto{\pgfqpoint{1.910803in}{1.469752in}}%
\pgfpathlineto{\pgfqpoint{1.922805in}{1.460295in}}%
\pgfpathlineto{\pgfqpoint{1.933835in}{1.454564in}}%
\pgfpathlineto{\pgfqpoint{1.940045in}{1.444352in}}%
\pgfpathlineto{\pgfqpoint{1.942134in}{1.425554in}}%
\pgfpathlineto{\pgfqpoint{1.883600in}{1.423428in}}%
\pgfpathlineto{\pgfqpoint{1.833705in}{1.422385in}}%
\pgfpathlineto{\pgfqpoint{1.788246in}{1.421717in}}%
\pgfpathlineto{\pgfqpoint{1.740155in}{1.421790in}}%
\pgfpathlineto{\pgfqpoint{1.692728in}{1.422304in}}%
\pgfpathclose%
\pgfusepath{fill}%
\end{pgfscope}%
\begin{pgfscope}%
\pgfpathrectangle{\pgfqpoint{0.100000in}{0.100000in}}{\pgfqpoint{2.989028in}{1.913466in}}%
\pgfusepath{clip}%
\pgfsetbuttcap%
\pgfsetmiterjoin%
\definecolor{currentfill}{rgb}{0.400000,0.760784,0.647059}%
\pgfsetfillcolor{currentfill}%
\pgfsetlinewidth{0.000000pt}%
\definecolor{currentstroke}{rgb}{0.000000,0.000000,0.000000}%
\pgfsetstrokecolor{currentstroke}%
\pgfsetstrokeopacity{0.000000}%
\pgfsetdash{}{0pt}%
\pgfpathmoveto{\pgfqpoint{0.363713in}{1.549891in}}%
\pgfpathlineto{\pgfqpoint{0.359077in}{1.558421in}}%
\pgfpathlineto{\pgfqpoint{0.362037in}{1.580316in}}%
\pgfpathlineto{\pgfqpoint{0.367750in}{1.591758in}}%
\pgfpathlineto{\pgfqpoint{0.364894in}{1.607237in}}%
\pgfpathlineto{\pgfqpoint{0.370854in}{1.613747in}}%
\pgfpathlineto{\pgfqpoint{0.381727in}{1.631292in}}%
\pgfpathlineto{\pgfqpoint{0.390946in}{1.641883in}}%
\pgfpathlineto{\pgfqpoint{0.404422in}{1.664951in}}%
\pgfpathlineto{\pgfqpoint{0.414758in}{1.690079in}}%
\pgfpathlineto{\pgfqpoint{0.425730in}{1.713646in}}%
\pgfpathlineto{\pgfqpoint{0.427948in}{1.723489in}}%
\pgfpathlineto{\pgfqpoint{0.443085in}{1.751647in}}%
\pgfpathlineto{\pgfqpoint{0.445999in}{1.764014in}}%
\pgfpathlineto{\pgfqpoint{0.452444in}{1.777001in}}%
\pgfpathlineto{\pgfqpoint{0.454747in}{1.792843in}}%
\pgfpathlineto{\pgfqpoint{0.467051in}{1.800568in}}%
\pgfpathlineto{\pgfqpoint{0.481639in}{1.804461in}}%
\pgfpathlineto{\pgfqpoint{0.489015in}{1.795741in}}%
\pgfpathlineto{\pgfqpoint{0.495391in}{1.796386in}}%
\pgfpathlineto{\pgfqpoint{0.505347in}{1.785992in}}%
\pgfpathlineto{\pgfqpoint{0.506670in}{1.776202in}}%
\pgfpathlineto{\pgfqpoint{0.502895in}{1.757031in}}%
\pgfpathlineto{\pgfqpoint{0.515456in}{1.747249in}}%
\pgfpathlineto{\pgfqpoint{0.523603in}{1.743571in}}%
\pgfpathlineto{\pgfqpoint{0.545673in}{1.747325in}}%
\pgfpathlineto{\pgfqpoint{0.558460in}{1.744762in}}%
\pgfpathlineto{\pgfqpoint{0.567248in}{1.740554in}}%
\pgfpathlineto{\pgfqpoint{0.571877in}{1.732568in}}%
\pgfpathlineto{\pgfqpoint{0.601190in}{1.733634in}}%
\pgfpathlineto{\pgfqpoint{0.605851in}{1.728747in}}%
\pgfpathlineto{\pgfqpoint{0.616848in}{1.727692in}}%
\pgfpathlineto{\pgfqpoint{0.627921in}{1.730814in}}%
\pgfpathlineto{\pgfqpoint{0.660663in}{1.730028in}}%
\pgfpathlineto{\pgfqpoint{0.667734in}{1.727622in}}%
\pgfpathlineto{\pgfqpoint{0.675998in}{1.730394in}}%
\pgfpathlineto{\pgfqpoint{0.767475in}{1.706523in}}%
\pgfpathlineto{\pgfqpoint{0.771203in}{1.693698in}}%
\pgfpathlineto{\pgfqpoint{0.779597in}{1.686749in}}%
\pgfpathlineto{\pgfqpoint{0.780469in}{1.674240in}}%
\pgfpathlineto{\pgfqpoint{0.768191in}{1.661575in}}%
\pgfpathlineto{\pgfqpoint{0.762206in}{1.649603in}}%
\pgfpathlineto{\pgfqpoint{0.754997in}{1.643207in}}%
\pgfpathlineto{\pgfqpoint{0.754345in}{1.636806in}}%
\pgfpathlineto{\pgfqpoint{0.742402in}{1.628155in}}%
\pgfpathlineto{\pgfqpoint{0.734208in}{1.615650in}}%
\pgfpathlineto{\pgfqpoint{0.729620in}{1.612613in}}%
\pgfpathlineto{\pgfqpoint{0.727118in}{1.600478in}}%
\pgfpathlineto{\pgfqpoint{0.739170in}{1.588120in}}%
\pgfpathlineto{\pgfqpoint{0.727531in}{1.567559in}}%
\pgfpathlineto{\pgfqpoint{0.726119in}{1.559712in}}%
\pgfpathlineto{\pgfqpoint{0.699946in}{1.452341in}}%
\pgfpathlineto{\pgfqpoint{0.644893in}{1.466447in}}%
\pgfpathlineto{\pgfqpoint{0.591784in}{1.480141in}}%
\pgfpathlineto{\pgfqpoint{0.559744in}{1.489057in}}%
\pgfpathlineto{\pgfqpoint{0.518578in}{1.500784in}}%
\pgfpathlineto{\pgfqpoint{0.452913in}{1.521452in}}%
\pgfpathlineto{\pgfqpoint{0.381532in}{1.543676in}}%
\pgfpathlineto{\pgfqpoint{0.363713in}{1.549891in}}%
\pgfpathclose%
\pgfusepath{fill}%
\end{pgfscope}%
\begin{pgfscope}%
\pgfpathrectangle{\pgfqpoint{0.100000in}{0.100000in}}{\pgfqpoint{2.989028in}{1.913466in}}%
\pgfusepath{clip}%
\pgfsetbuttcap%
\pgfsetmiterjoin%
\definecolor{currentfill}{rgb}{0.453057,0.781546,0.646290}%
\pgfsetfillcolor{currentfill}%
\pgfsetlinewidth{0.000000pt}%
\definecolor{currentstroke}{rgb}{0.000000,0.000000,0.000000}%
\pgfsetstrokecolor{currentstroke}%
\pgfsetstrokeopacity{0.000000}%
\pgfsetdash{}{0pt}%
\pgfpathmoveto{\pgfqpoint{2.837703in}{1.505030in}}%
\pgfpathlineto{\pgfqpoint{2.832504in}{1.509256in}}%
\pgfpathlineto{\pgfqpoint{2.833228in}{1.522426in}}%
\pgfpathlineto{\pgfqpoint{2.827854in}{1.559858in}}%
\pgfpathlineto{\pgfqpoint{2.832066in}{1.573021in}}%
\pgfpathlineto{\pgfqpoint{2.835075in}{1.596366in}}%
\pgfpathlineto{\pgfqpoint{2.831015in}{1.607411in}}%
\pgfpathlineto{\pgfqpoint{2.845785in}{1.622115in}}%
\pgfpathlineto{\pgfqpoint{2.849502in}{1.633812in}}%
\pgfpathlineto{\pgfqpoint{2.842504in}{1.643043in}}%
\pgfpathlineto{\pgfqpoint{2.845485in}{1.654580in}}%
\pgfpathlineto{\pgfqpoint{2.843616in}{1.661826in}}%
\pgfpathlineto{\pgfqpoint{2.845221in}{1.677346in}}%
\pgfpathlineto{\pgfqpoint{2.848221in}{1.682137in}}%
\pgfpathlineto{\pgfqpoint{2.857471in}{1.685161in}}%
\pgfpathlineto{\pgfqpoint{2.870964in}{1.645645in}}%
\pgfpathlineto{\pgfqpoint{2.895998in}{1.563760in}}%
\pgfpathlineto{\pgfqpoint{2.905471in}{1.557411in}}%
\pgfpathlineto{\pgfqpoint{2.907175in}{1.550163in}}%
\pgfpathlineto{\pgfqpoint{2.912175in}{1.547235in}}%
\pgfpathlineto{\pgfqpoint{2.911795in}{1.534090in}}%
\pgfpathlineto{\pgfqpoint{2.906497in}{1.533879in}}%
\pgfpathlineto{\pgfqpoint{2.895764in}{1.525694in}}%
\pgfpathlineto{\pgfqpoint{2.892669in}{1.517482in}}%
\pgfpathlineto{\pgfqpoint{2.863942in}{1.510374in}}%
\pgfpathlineto{\pgfqpoint{2.837703in}{1.505030in}}%
\pgfpathclose%
\pgfusepath{fill}%
\end{pgfscope}%
\begin{pgfscope}%
\pgfpathrectangle{\pgfqpoint{0.100000in}{0.100000in}}{\pgfqpoint{2.989028in}{1.913466in}}%
\pgfusepath{clip}%
\pgfsetbuttcap%
\pgfsetmiterjoin%
\definecolor{currentfill}{rgb}{0.516724,0.806459,0.645367}%
\pgfsetfillcolor{currentfill}%
\pgfsetlinewidth{0.000000pt}%
\definecolor{currentstroke}{rgb}{0.000000,0.000000,0.000000}%
\pgfsetstrokecolor{currentstroke}%
\pgfsetstrokeopacity{0.000000}%
\pgfsetdash{}{0pt}%
\pgfpathmoveto{\pgfqpoint{1.939431in}{1.220143in}}%
\pgfpathlineto{\pgfqpoint{1.924266in}{1.235161in}}%
\pgfpathlineto{\pgfqpoint{1.875789in}{1.232365in}}%
\pgfpathlineto{\pgfqpoint{1.800184in}{1.229874in}}%
\pgfpathlineto{\pgfqpoint{1.724124in}{1.231060in}}%
\pgfpathlineto{\pgfqpoint{1.718785in}{1.240370in}}%
\pgfpathlineto{\pgfqpoint{1.720965in}{1.249511in}}%
\pgfpathlineto{\pgfqpoint{1.719930in}{1.265170in}}%
\pgfpathlineto{\pgfqpoint{1.715910in}{1.280338in}}%
\pgfpathlineto{\pgfqpoint{1.716378in}{1.288348in}}%
\pgfpathlineto{\pgfqpoint{1.707833in}{1.297398in}}%
\pgfpathlineto{\pgfqpoint{1.709692in}{1.309990in}}%
\pgfpathlineto{\pgfqpoint{1.696585in}{1.334864in}}%
\pgfpathlineto{\pgfqpoint{1.692680in}{1.355981in}}%
\pgfpathlineto{\pgfqpoint{1.683252in}{1.373157in}}%
\pgfpathlineto{\pgfqpoint{1.692201in}{1.394464in}}%
\pgfpathlineto{\pgfqpoint{1.687744in}{1.404379in}}%
\pgfpathlineto{\pgfqpoint{1.692728in}{1.422304in}}%
\pgfpathlineto{\pgfqpoint{1.740155in}{1.421790in}}%
\pgfpathlineto{\pgfqpoint{1.788246in}{1.421717in}}%
\pgfpathlineto{\pgfqpoint{1.833705in}{1.422385in}}%
\pgfpathlineto{\pgfqpoint{1.883600in}{1.423428in}}%
\pgfpathlineto{\pgfqpoint{1.942134in}{1.425554in}}%
\pgfpathlineto{\pgfqpoint{1.943697in}{1.417056in}}%
\pgfpathlineto{\pgfqpoint{1.950178in}{1.409931in}}%
\pgfpathlineto{\pgfqpoint{1.946075in}{1.402507in}}%
\pgfpathlineto{\pgfqpoint{1.947756in}{1.388016in}}%
\pgfpathlineto{\pgfqpoint{1.951772in}{1.376303in}}%
\pgfpathlineto{\pgfqpoint{1.969980in}{1.370103in}}%
\pgfpathlineto{\pgfqpoint{1.973100in}{1.361809in}}%
\pgfpathlineto{\pgfqpoint{1.983091in}{1.352497in}}%
\pgfpathlineto{\pgfqpoint{1.987167in}{1.342864in}}%
\pgfpathlineto{\pgfqpoint{1.997291in}{1.336402in}}%
\pgfpathlineto{\pgfqpoint{1.998875in}{1.328623in}}%
\pgfpathlineto{\pgfqpoint{1.996904in}{1.316845in}}%
\pgfpathlineto{\pgfqpoint{1.991751in}{1.313315in}}%
\pgfpathlineto{\pgfqpoint{1.990197in}{1.302104in}}%
\pgfpathlineto{\pgfqpoint{1.975389in}{1.293200in}}%
\pgfpathlineto{\pgfqpoint{1.956087in}{1.288322in}}%
\pgfpathlineto{\pgfqpoint{1.954318in}{1.277106in}}%
\pgfpathlineto{\pgfqpoint{1.961764in}{1.269093in}}%
\pgfpathlineto{\pgfqpoint{1.962068in}{1.259018in}}%
\pgfpathlineto{\pgfqpoint{1.956060in}{1.251098in}}%
\pgfpathlineto{\pgfqpoint{1.952907in}{1.239332in}}%
\pgfpathlineto{\pgfqpoint{1.942465in}{1.235438in}}%
\pgfpathlineto{\pgfqpoint{1.943130in}{1.222326in}}%
\pgfpathlineto{\pgfqpoint{1.939431in}{1.220143in}}%
\pgfpathclose%
\pgfusepath{fill}%
\end{pgfscope}%
\begin{pgfscope}%
\pgfpathrectangle{\pgfqpoint{0.100000in}{0.100000in}}{\pgfqpoint{2.989028in}{1.913466in}}%
\pgfusepath{clip}%
\pgfsetbuttcap%
\pgfsetmiterjoin%
\definecolor{currentfill}{rgb}{0.400000,0.760784,0.647059}%
\pgfsetfillcolor{currentfill}%
\pgfsetlinewidth{0.000000pt}%
\definecolor{currentstroke}{rgb}{0.000000,0.000000,0.000000}%
\pgfsetstrokecolor{currentstroke}%
\pgfsetstrokeopacity{0.000000}%
\pgfsetdash{}{0pt}%
\pgfpathmoveto{\pgfqpoint{2.880364in}{1.467139in}}%
\pgfpathlineto{\pgfqpoint{2.866000in}{1.464872in}}%
\pgfpathlineto{\pgfqpoint{2.800022in}{1.449880in}}%
\pgfpathlineto{\pgfqpoint{2.798869in}{1.451627in}}%
\pgfpathlineto{\pgfqpoint{2.799865in}{1.496989in}}%
\pgfpathlineto{\pgfqpoint{2.837703in}{1.505030in}}%
\pgfpathlineto{\pgfqpoint{2.863942in}{1.510374in}}%
\pgfpathlineto{\pgfqpoint{2.892669in}{1.517482in}}%
\pgfpathlineto{\pgfqpoint{2.895764in}{1.525694in}}%
\pgfpathlineto{\pgfqpoint{2.906497in}{1.533879in}}%
\pgfpathlineto{\pgfqpoint{2.911795in}{1.534090in}}%
\pgfpathlineto{\pgfqpoint{2.918759in}{1.522149in}}%
\pgfpathlineto{\pgfqpoint{2.912440in}{1.504717in}}%
\pgfpathlineto{\pgfqpoint{2.911512in}{1.494495in}}%
\pgfpathlineto{\pgfqpoint{2.924304in}{1.495455in}}%
\pgfpathlineto{\pgfqpoint{2.930105in}{1.490552in}}%
\pgfpathlineto{\pgfqpoint{2.943114in}{1.470509in}}%
\pgfpathlineto{\pgfqpoint{2.949579in}{1.468072in}}%
\pgfpathlineto{\pgfqpoint{2.960410in}{1.468976in}}%
\pgfpathlineto{\pgfqpoint{2.967946in}{1.475787in}}%
\pgfpathlineto{\pgfqpoint{2.972956in}{1.469708in}}%
\pgfpathlineto{\pgfqpoint{2.941451in}{1.453083in}}%
\pgfpathlineto{\pgfqpoint{2.940467in}{1.465013in}}%
\pgfpathlineto{\pgfqpoint{2.926155in}{1.446471in}}%
\pgfpathlineto{\pgfqpoint{2.921123in}{1.443277in}}%
\pgfpathlineto{\pgfqpoint{2.914105in}{1.453976in}}%
\pgfpathlineto{\pgfqpoint{2.912183in}{1.455442in}}%
\pgfpathlineto{\pgfqpoint{2.905650in}{1.458903in}}%
\pgfpathlineto{\pgfqpoint{2.899952in}{1.472940in}}%
\pgfpathlineto{\pgfqpoint{2.880364in}{1.467139in}}%
\pgfpathclose%
\pgfusepath{fill}%
\end{pgfscope}%
\begin{pgfscope}%
\pgfpathrectangle{\pgfqpoint{0.100000in}{0.100000in}}{\pgfqpoint{2.989028in}{1.913466in}}%
\pgfusepath{clip}%
\pgfsetbuttcap%
\pgfsetmiterjoin%
\definecolor{currentfill}{rgb}{0.321261,0.371165,0.664360}%
\pgfsetfillcolor{currentfill}%
\pgfsetlinewidth{0.000000pt}%
\definecolor{currentstroke}{rgb}{0.000000,0.000000,0.000000}%
\pgfsetstrokecolor{currentstroke}%
\pgfsetstrokeopacity{0.000000}%
\pgfsetdash{}{0pt}%
\pgfpathmoveto{\pgfqpoint{1.408935in}{1.205970in}}%
\pgfpathlineto{\pgfqpoint{1.414195in}{1.271196in}}%
\pgfpathlineto{\pgfqpoint{1.384404in}{1.273614in}}%
\pgfpathlineto{\pgfqpoint{1.315540in}{1.280295in}}%
\pgfpathlineto{\pgfqpoint{1.321947in}{1.345289in}}%
\pgfpathlineto{\pgfqpoint{1.328363in}{1.410780in}}%
\pgfpathlineto{\pgfqpoint{1.354836in}{1.407921in}}%
\pgfpathlineto{\pgfqpoint{1.422447in}{1.401533in}}%
\pgfpathlineto{\pgfqpoint{1.463365in}{1.398329in}}%
\pgfpathlineto{\pgfqpoint{1.512902in}{1.395489in}}%
\pgfpathlineto{\pgfqpoint{1.594439in}{1.391630in}}%
\pgfpathlineto{\pgfqpoint{1.596231in}{1.387975in}}%
\pgfpathlineto{\pgfqpoint{1.619594in}{1.376416in}}%
\pgfpathlineto{\pgfqpoint{1.626759in}{1.382286in}}%
\pgfpathlineto{\pgfqpoint{1.647400in}{1.381750in}}%
\pgfpathlineto{\pgfqpoint{1.659700in}{1.375307in}}%
\pgfpathlineto{\pgfqpoint{1.679205in}{1.367982in}}%
\pgfpathlineto{\pgfqpoint{1.684670in}{1.357295in}}%
\pgfpathlineto{\pgfqpoint{1.692680in}{1.355981in}}%
\pgfpathlineto{\pgfqpoint{1.696585in}{1.334864in}}%
\pgfpathlineto{\pgfqpoint{1.709692in}{1.309990in}}%
\pgfpathlineto{\pgfqpoint{1.707833in}{1.297398in}}%
\pgfpathlineto{\pgfqpoint{1.716378in}{1.288348in}}%
\pgfpathlineto{\pgfqpoint{1.715910in}{1.280338in}}%
\pgfpathlineto{\pgfqpoint{1.719930in}{1.265170in}}%
\pgfpathlineto{\pgfqpoint{1.720965in}{1.249511in}}%
\pgfpathlineto{\pgfqpoint{1.718785in}{1.240370in}}%
\pgfpathlineto{\pgfqpoint{1.724124in}{1.231060in}}%
\pgfpathlineto{\pgfqpoint{1.731447in}{1.214132in}}%
\pgfpathlineto{\pgfqpoint{1.738455in}{1.207242in}}%
\pgfpathlineto{\pgfqpoint{1.746810in}{1.192312in}}%
\pgfpathlineto{\pgfqpoint{1.723131in}{1.192066in}}%
\pgfpathlineto{\pgfqpoint{1.671933in}{1.192859in}}%
\pgfpathlineto{\pgfqpoint{1.615367in}{1.194592in}}%
\pgfpathlineto{\pgfqpoint{1.558468in}{1.196775in}}%
\pgfpathlineto{\pgfqpoint{1.474809in}{1.201345in}}%
\pgfpathlineto{\pgfqpoint{1.408935in}{1.205970in}}%
\pgfpathclose%
\pgfusepath{fill}%
\end{pgfscope}%
\begin{pgfscope}%
\pgfpathrectangle{\pgfqpoint{0.100000in}{0.100000in}}{\pgfqpoint{2.989028in}{1.913466in}}%
\pgfusepath{clip}%
\pgfsetbuttcap%
\pgfsetmiterjoin%
\definecolor{currentfill}{rgb}{0.622837,0.847982,0.643829}%
\pgfsetfillcolor{currentfill}%
\pgfsetlinewidth{0.000000pt}%
\definecolor{currentstroke}{rgb}{0.000000,0.000000,0.000000}%
\pgfsetstrokecolor{currentstroke}%
\pgfsetstrokeopacity{0.000000}%
\pgfsetdash{}{0pt}%
\pgfpathmoveto{\pgfqpoint{2.498213in}{1.402714in}}%
\pgfpathlineto{\pgfqpoint{2.524409in}{1.427692in}}%
\pgfpathlineto{\pgfqpoint{2.527665in}{1.436570in}}%
\pgfpathlineto{\pgfqpoint{2.535341in}{1.444171in}}%
\pgfpathlineto{\pgfqpoint{2.529581in}{1.455238in}}%
\pgfpathlineto{\pgfqpoint{2.522360in}{1.461711in}}%
\pgfpathlineto{\pgfqpoint{2.520273in}{1.473182in}}%
\pgfpathlineto{\pgfqpoint{2.547151in}{1.484959in}}%
\pgfpathlineto{\pgfqpoint{2.569405in}{1.488685in}}%
\pgfpathlineto{\pgfqpoint{2.581363in}{1.488934in}}%
\pgfpathlineto{\pgfqpoint{2.590473in}{1.484433in}}%
\pgfpathlineto{\pgfqpoint{2.599352in}{1.488453in}}%
\pgfpathlineto{\pgfqpoint{2.621010in}{1.492956in}}%
\pgfpathlineto{\pgfqpoint{2.628493in}{1.498769in}}%
\pgfpathlineto{\pgfqpoint{2.639589in}{1.511638in}}%
\pgfpathlineto{\pgfqpoint{2.649683in}{1.517325in}}%
\pgfpathlineto{\pgfqpoint{2.650395in}{1.522779in}}%
\pgfpathlineto{\pgfqpoint{2.645097in}{1.535222in}}%
\pgfpathlineto{\pgfqpoint{2.648927in}{1.542544in}}%
\pgfpathlineto{\pgfqpoint{2.643770in}{1.550418in}}%
\pgfpathlineto{\pgfqpoint{2.635873in}{1.550968in}}%
\pgfpathlineto{\pgfqpoint{2.655638in}{1.574751in}}%
\pgfpathlineto{\pgfqpoint{2.657985in}{1.583815in}}%
\pgfpathlineto{\pgfqpoint{2.673541in}{1.606931in}}%
\pgfpathlineto{\pgfqpoint{2.687981in}{1.619440in}}%
\pgfpathlineto{\pgfqpoint{2.697897in}{1.624664in}}%
\pgfpathlineto{\pgfqpoint{2.730342in}{1.632016in}}%
\pgfpathlineto{\pgfqpoint{2.760752in}{1.640318in}}%
\pgfpathlineto{\pgfqpoint{2.764780in}{1.627285in}}%
\pgfpathlineto{\pgfqpoint{2.765065in}{1.615202in}}%
\pgfpathlineto{\pgfqpoint{2.772002in}{1.604233in}}%
\pgfpathlineto{\pgfqpoint{2.774140in}{1.593064in}}%
\pgfpathlineto{\pgfqpoint{2.771474in}{1.577911in}}%
\pgfpathlineto{\pgfqpoint{2.779593in}{1.557473in}}%
\pgfpathlineto{\pgfqpoint{2.788848in}{1.546108in}}%
\pgfpathlineto{\pgfqpoint{2.796007in}{1.509167in}}%
\pgfpathlineto{\pgfqpoint{2.799865in}{1.496989in}}%
\pgfpathlineto{\pgfqpoint{2.798869in}{1.451627in}}%
\pgfpathlineto{\pgfqpoint{2.800022in}{1.449880in}}%
\pgfpathlineto{\pgfqpoint{2.808450in}{1.401092in}}%
\pgfpathlineto{\pgfqpoint{2.813176in}{1.396647in}}%
\pgfpathlineto{\pgfqpoint{2.803010in}{1.386769in}}%
\pgfpathlineto{\pgfqpoint{2.808026in}{1.381100in}}%
\pgfpathlineto{\pgfqpoint{2.803617in}{1.372526in}}%
\pgfpathlineto{\pgfqpoint{2.803654in}{1.368875in}}%
\pgfpathlineto{\pgfqpoint{2.798142in}{1.365581in}}%
\pgfpathlineto{\pgfqpoint{2.795458in}{1.358300in}}%
\pgfpathlineto{\pgfqpoint{2.796310in}{1.378322in}}%
\pgfpathlineto{\pgfqpoint{2.779225in}{1.382730in}}%
\pgfpathlineto{\pgfqpoint{2.752507in}{1.391844in}}%
\pgfpathlineto{\pgfqpoint{2.748728in}{1.396333in}}%
\pgfpathlineto{\pgfqpoint{2.737564in}{1.397408in}}%
\pgfpathlineto{\pgfqpoint{2.731142in}{1.404090in}}%
\pgfpathlineto{\pgfqpoint{2.727676in}{1.417397in}}%
\pgfpathlineto{\pgfqpoint{2.718572in}{1.419052in}}%
\pgfpathlineto{\pgfqpoint{2.712484in}{1.426033in}}%
\pgfpathlineto{\pgfqpoint{2.634972in}{1.410406in}}%
\pgfpathlineto{\pgfqpoint{2.597928in}{1.402745in}}%
\pgfpathlineto{\pgfqpoint{2.541683in}{1.392479in}}%
\pgfpathlineto{\pgfqpoint{2.501180in}{1.385646in}}%
\pgfpathlineto{\pgfqpoint{2.498213in}{1.402714in}}%
\pgfpathclose%
\pgfusepath{fill}%
\end{pgfscope}%
\begin{pgfscope}%
\pgfpathrectangle{\pgfqpoint{0.100000in}{0.100000in}}{\pgfqpoint{2.989028in}{1.913466in}}%
\pgfusepath{clip}%
\pgfsetbuttcap%
\pgfsetmiterjoin%
\definecolor{currentfill}{rgb}{0.622837,0.847982,0.643829}%
\pgfsetfillcolor{currentfill}%
\pgfsetlinewidth{0.000000pt}%
\definecolor{currentstroke}{rgb}{0.000000,0.000000,0.000000}%
\pgfsetstrokecolor{currentstroke}%
\pgfsetstrokeopacity{0.000000}%
\pgfsetdash{}{0pt}%
\pgfpathmoveto{\pgfqpoint{2.809390in}{1.354233in}}%
\pgfpathlineto{\pgfqpoint{2.797403in}{1.350498in}}%
\pgfpathlineto{\pgfqpoint{2.795386in}{1.353932in}}%
\pgfpathlineto{\pgfqpoint{2.799234in}{1.365456in}}%
\pgfpathlineto{\pgfqpoint{2.806360in}{1.366591in}}%
\pgfpathlineto{\pgfqpoint{2.805722in}{1.370219in}}%
\pgfpathlineto{\pgfqpoint{2.812120in}{1.375665in}}%
\pgfpathlineto{\pgfqpoint{2.830618in}{1.380004in}}%
\pgfpathlineto{\pgfqpoint{2.838853in}{1.386571in}}%
\pgfpathlineto{\pgfqpoint{2.857353in}{1.392045in}}%
\pgfpathlineto{\pgfqpoint{2.865790in}{1.390042in}}%
\pgfpathlineto{\pgfqpoint{2.872888in}{1.398866in}}%
\pgfpathlineto{\pgfqpoint{2.883615in}{1.400032in}}%
\pgfpathlineto{\pgfqpoint{2.865314in}{1.382820in}}%
\pgfpathlineto{\pgfqpoint{2.809390in}{1.354233in}}%
\pgfpathclose%
\pgfusepath{fill}%
\end{pgfscope}%
\begin{pgfscope}%
\pgfpathrectangle{\pgfqpoint{0.100000in}{0.100000in}}{\pgfqpoint{2.989028in}{1.913466in}}%
\pgfusepath{clip}%
\pgfsetbuttcap%
\pgfsetmiterjoin%
\definecolor{currentfill}{rgb}{0.296040,0.644829,0.695040}%
\pgfsetfillcolor{currentfill}%
\pgfsetlinewidth{0.000000pt}%
\definecolor{currentstroke}{rgb}{0.000000,0.000000,0.000000}%
\pgfsetstrokecolor{currentstroke}%
\pgfsetstrokeopacity{0.000000}%
\pgfsetdash{}{0pt}%
\pgfpathmoveto{\pgfqpoint{2.540072in}{1.240581in}}%
\pgfpathlineto{\pgfqpoint{2.488234in}{1.232010in}}%
\pgfpathlineto{\pgfqpoint{2.478831in}{1.291252in}}%
\pgfpathlineto{\pgfqpoint{2.464869in}{1.378577in}}%
\pgfpathlineto{\pgfqpoint{2.498213in}{1.402714in}}%
\pgfpathlineto{\pgfqpoint{2.501180in}{1.385646in}}%
\pgfpathlineto{\pgfqpoint{2.541683in}{1.392479in}}%
\pgfpathlineto{\pgfqpoint{2.597928in}{1.402745in}}%
\pgfpathlineto{\pgfqpoint{2.634972in}{1.410406in}}%
\pgfpathlineto{\pgfqpoint{2.712484in}{1.426033in}}%
\pgfpathlineto{\pgfqpoint{2.718572in}{1.419052in}}%
\pgfpathlineto{\pgfqpoint{2.727676in}{1.417397in}}%
\pgfpathlineto{\pgfqpoint{2.731142in}{1.404090in}}%
\pgfpathlineto{\pgfqpoint{2.737564in}{1.397408in}}%
\pgfpathlineto{\pgfqpoint{2.748728in}{1.396333in}}%
\pgfpathlineto{\pgfqpoint{2.752507in}{1.391844in}}%
\pgfpathlineto{\pgfqpoint{2.748671in}{1.388380in}}%
\pgfpathlineto{\pgfqpoint{2.745215in}{1.376126in}}%
\pgfpathlineto{\pgfqpoint{2.737116in}{1.362357in}}%
\pgfpathlineto{\pgfqpoint{2.742552in}{1.356369in}}%
\pgfpathlineto{\pgfqpoint{2.738104in}{1.349953in}}%
\pgfpathlineto{\pgfqpoint{2.740613in}{1.335884in}}%
\pgfpathlineto{\pgfqpoint{2.748625in}{1.328499in}}%
\pgfpathlineto{\pgfqpoint{2.767634in}{1.316500in}}%
\pgfpathlineto{\pgfqpoint{2.752326in}{1.299570in}}%
\pgfpathlineto{\pgfqpoint{2.752111in}{1.293145in}}%
\pgfpathlineto{\pgfqpoint{2.739648in}{1.284858in}}%
\pgfpathlineto{\pgfqpoint{2.725867in}{1.283302in}}%
\pgfpathlineto{\pgfqpoint{2.722458in}{1.276100in}}%
\pgfpathlineto{\pgfqpoint{2.684127in}{1.267806in}}%
\pgfpathlineto{\pgfqpoint{2.639399in}{1.258797in}}%
\pgfpathlineto{\pgfqpoint{2.608655in}{1.253292in}}%
\pgfpathlineto{\pgfqpoint{2.540072in}{1.240581in}}%
\pgfpathclose%
\pgfusepath{fill}%
\end{pgfscope}%
\begin{pgfscope}%
\pgfpathrectangle{\pgfqpoint{0.100000in}{0.100000in}}{\pgfqpoint{2.989028in}{1.913466in}}%
\pgfusepath{clip}%
\pgfsetbuttcap%
\pgfsetmiterjoin%
\definecolor{currentfill}{rgb}{0.410611,0.764937,0.646905}%
\pgfsetfillcolor{currentfill}%
\pgfsetlinewidth{0.000000pt}%
\definecolor{currentstroke}{rgb}{0.000000,0.000000,0.000000}%
\pgfsetstrokecolor{currentstroke}%
\pgfsetstrokeopacity{0.000000}%
\pgfsetdash{}{0pt}%
\pgfpathmoveto{\pgfqpoint{2.800022in}{1.449880in}}%
\pgfpathlineto{\pgfqpoint{2.866000in}{1.464872in}}%
\pgfpathlineto{\pgfqpoint{2.880364in}{1.467139in}}%
\pgfpathlineto{\pgfqpoint{2.889871in}{1.429760in}}%
\pgfpathlineto{\pgfqpoint{2.888366in}{1.423067in}}%
\pgfpathlineto{\pgfqpoint{2.857827in}{1.411244in}}%
\pgfpathlineto{\pgfqpoint{2.839596in}{1.407111in}}%
\pgfpathlineto{\pgfqpoint{2.831849in}{1.397842in}}%
\pgfpathlineto{\pgfqpoint{2.808026in}{1.381100in}}%
\pgfpathlineto{\pgfqpoint{2.803010in}{1.386769in}}%
\pgfpathlineto{\pgfqpoint{2.813176in}{1.396647in}}%
\pgfpathlineto{\pgfqpoint{2.808450in}{1.401092in}}%
\pgfpathlineto{\pgfqpoint{2.800022in}{1.449880in}}%
\pgfpathclose%
\pgfusepath{fill}%
\end{pgfscope}%
\begin{pgfscope}%
\pgfpathrectangle{\pgfqpoint{0.100000in}{0.100000in}}{\pgfqpoint{2.989028in}{1.913466in}}%
\pgfusepath{clip}%
\pgfsetbuttcap%
\pgfsetmiterjoin%
\definecolor{currentfill}{rgb}{0.548558,0.818916,0.644906}%
\pgfsetfillcolor{currentfill}%
\pgfsetlinewidth{0.000000pt}%
\definecolor{currentstroke}{rgb}{0.000000,0.000000,0.000000}%
\pgfsetstrokecolor{currentstroke}%
\pgfsetstrokeopacity{0.000000}%
\pgfsetdash{}{0pt}%
\pgfpathmoveto{\pgfqpoint{2.880364in}{1.467139in}}%
\pgfpathlineto{\pgfqpoint{2.899952in}{1.472940in}}%
\pgfpathlineto{\pgfqpoint{2.905650in}{1.458903in}}%
\pgfpathlineto{\pgfqpoint{2.912183in}{1.455442in}}%
\pgfpathlineto{\pgfqpoint{2.904121in}{1.449529in}}%
\pgfpathlineto{\pgfqpoint{2.905138in}{1.432164in}}%
\pgfpathlineto{\pgfqpoint{2.888366in}{1.423067in}}%
\pgfpathlineto{\pgfqpoint{2.889871in}{1.429760in}}%
\pgfpathlineto{\pgfqpoint{2.880364in}{1.467139in}}%
\pgfpathclose%
\pgfusepath{fill}%
\end{pgfscope}%
\begin{pgfscope}%
\pgfpathrectangle{\pgfqpoint{0.100000in}{0.100000in}}{\pgfqpoint{2.989028in}{1.913466in}}%
\pgfusepath{clip}%
\pgfsetbuttcap%
\pgfsetmiterjoin%
\definecolor{currentfill}{rgb}{0.400000,0.760784,0.647059}%
\pgfsetfillcolor{currentfill}%
\pgfsetlinewidth{0.000000pt}%
\definecolor{currentstroke}{rgb}{0.000000,0.000000,0.000000}%
\pgfsetstrokecolor{currentstroke}%
\pgfsetstrokeopacity{0.000000}%
\pgfsetdash{}{0pt}%
\pgfpathmoveto{\pgfqpoint{2.737440in}{1.278737in}}%
\pgfpathlineto{\pgfqpoint{2.739648in}{1.284858in}}%
\pgfpathlineto{\pgfqpoint{2.752111in}{1.293145in}}%
\pgfpathlineto{\pgfqpoint{2.752326in}{1.299570in}}%
\pgfpathlineto{\pgfqpoint{2.767634in}{1.316500in}}%
\pgfpathlineto{\pgfqpoint{2.748625in}{1.328499in}}%
\pgfpathlineto{\pgfqpoint{2.740613in}{1.335884in}}%
\pgfpathlineto{\pgfqpoint{2.738104in}{1.349953in}}%
\pgfpathlineto{\pgfqpoint{2.742552in}{1.356369in}}%
\pgfpathlineto{\pgfqpoint{2.737116in}{1.362357in}}%
\pgfpathlineto{\pgfqpoint{2.745215in}{1.376126in}}%
\pgfpathlineto{\pgfqpoint{2.748671in}{1.388380in}}%
\pgfpathlineto{\pgfqpoint{2.752507in}{1.391844in}}%
\pgfpathlineto{\pgfqpoint{2.779225in}{1.382730in}}%
\pgfpathlineto{\pgfqpoint{2.796310in}{1.378322in}}%
\pgfpathlineto{\pgfqpoint{2.795458in}{1.358300in}}%
\pgfpathlineto{\pgfqpoint{2.785080in}{1.343116in}}%
\pgfpathlineto{\pgfqpoint{2.793632in}{1.340882in}}%
\pgfpathlineto{\pgfqpoint{2.802511in}{1.334366in}}%
\pgfpathlineto{\pgfqpoint{2.803034in}{1.316545in}}%
\pgfpathlineto{\pgfqpoint{2.800347in}{1.303928in}}%
\pgfpathlineto{\pgfqpoint{2.802135in}{1.293566in}}%
\pgfpathlineto{\pgfqpoint{2.786701in}{1.258597in}}%
\pgfpathlineto{\pgfqpoint{2.781154in}{1.242268in}}%
\pgfpathlineto{\pgfqpoint{2.773409in}{1.250205in}}%
\pgfpathlineto{\pgfqpoint{2.763139in}{1.248853in}}%
\pgfpathlineto{\pgfqpoint{2.737435in}{1.263703in}}%
\pgfpathlineto{\pgfqpoint{2.734805in}{1.271656in}}%
\pgfpathlineto{\pgfqpoint{2.737440in}{1.278737in}}%
\pgfpathclose%
\pgfusepath{fill}%
\end{pgfscope}%
\begin{pgfscope}%
\pgfpathrectangle{\pgfqpoint{0.100000in}{0.100000in}}{\pgfqpoint{2.989028in}{1.913466in}}%
\pgfusepath{clip}%
\pgfsetbuttcap%
\pgfsetmiterjoin%
\definecolor{currentfill}{rgb}{0.400000,0.760784,0.647059}%
\pgfsetfillcolor{currentfill}%
\pgfsetlinewidth{0.000000pt}%
\definecolor{currentstroke}{rgb}{0.000000,0.000000,0.000000}%
\pgfsetstrokecolor{currentstroke}%
\pgfsetstrokeopacity{0.000000}%
\pgfsetdash{}{0pt}%
\pgfpathmoveto{\pgfqpoint{2.122290in}{1.062454in}}%
\pgfpathlineto{\pgfqpoint{2.121088in}{1.078676in}}%
\pgfpathlineto{\pgfqpoint{2.125759in}{1.086724in}}%
\pgfpathlineto{\pgfqpoint{2.122681in}{1.090668in}}%
\pgfpathlineto{\pgfqpoint{2.133907in}{1.103580in}}%
\pgfpathlineto{\pgfqpoint{2.133263in}{1.106692in}}%
\pgfpathlineto{\pgfqpoint{2.144170in}{1.129219in}}%
\pgfpathlineto{\pgfqpoint{2.141966in}{1.140083in}}%
\pgfpathlineto{\pgfqpoint{2.134057in}{1.151457in}}%
\pgfpathlineto{\pgfqpoint{2.139547in}{1.165295in}}%
\pgfpathlineto{\pgfqpoint{2.132431in}{1.256358in}}%
\pgfpathlineto{\pgfqpoint{2.127306in}{1.320243in}}%
\pgfpathlineto{\pgfqpoint{2.134385in}{1.314949in}}%
\pgfpathlineto{\pgfqpoint{2.142283in}{1.315093in}}%
\pgfpathlineto{\pgfqpoint{2.160936in}{1.325898in}}%
\pgfpathlineto{\pgfqpoint{2.218150in}{1.331236in}}%
\pgfpathlineto{\pgfqpoint{2.260498in}{1.335699in}}%
\pgfpathlineto{\pgfqpoint{2.260872in}{1.331556in}}%
\pgfpathlineto{\pgfqpoint{2.270539in}{1.244011in}}%
\pgfpathlineto{\pgfqpoint{2.278860in}{1.162554in}}%
\pgfpathlineto{\pgfqpoint{2.275273in}{1.158730in}}%
\pgfpathlineto{\pgfqpoint{2.280783in}{1.149238in}}%
\pgfpathlineto{\pgfqpoint{2.280755in}{1.142397in}}%
\pgfpathlineto{\pgfqpoint{2.272893in}{1.140679in}}%
\pgfpathlineto{\pgfqpoint{2.264096in}{1.134084in}}%
\pgfpathlineto{\pgfqpoint{2.258139in}{1.136680in}}%
\pgfpathlineto{\pgfqpoint{2.249225in}{1.132459in}}%
\pgfpathlineto{\pgfqpoint{2.251983in}{1.123981in}}%
\pgfpathlineto{\pgfqpoint{2.242820in}{1.115463in}}%
\pgfpathlineto{\pgfqpoint{2.240290in}{1.105610in}}%
\pgfpathlineto{\pgfqpoint{2.233989in}{1.103993in}}%
\pgfpathlineto{\pgfqpoint{2.229287in}{1.096531in}}%
\pgfpathlineto{\pgfqpoint{2.229903in}{1.089016in}}%
\pgfpathlineto{\pgfqpoint{2.224371in}{1.083732in}}%
\pgfpathlineto{\pgfqpoint{2.216048in}{1.084548in}}%
\pgfpathlineto{\pgfqpoint{2.209721in}{1.092656in}}%
\pgfpathlineto{\pgfqpoint{2.199000in}{1.084848in}}%
\pgfpathlineto{\pgfqpoint{2.199753in}{1.078025in}}%
\pgfpathlineto{\pgfqpoint{2.187829in}{1.074039in}}%
\pgfpathlineto{\pgfqpoint{2.184834in}{1.079056in}}%
\pgfpathlineto{\pgfqpoint{2.173226in}{1.073367in}}%
\pgfpathlineto{\pgfqpoint{2.168985in}{1.065222in}}%
\pgfpathlineto{\pgfqpoint{2.155006in}{1.073574in}}%
\pgfpathlineto{\pgfqpoint{2.132850in}{1.068041in}}%
\pgfpathlineto{\pgfqpoint{2.122290in}{1.062454in}}%
\pgfpathclose%
\pgfusepath{fill}%
\end{pgfscope}%
\begin{pgfscope}%
\pgfpathrectangle{\pgfqpoint{0.100000in}{0.100000in}}{\pgfqpoint{2.989028in}{1.913466in}}%
\pgfusepath{clip}%
\pgfsetbuttcap%
\pgfsetmiterjoin%
\definecolor{currentfill}{rgb}{0.200077,0.537793,0.739331}%
\pgfsetfillcolor{currentfill}%
\pgfsetlinewidth{0.000000pt}%
\definecolor{currentstroke}{rgb}{0.000000,0.000000,0.000000}%
\pgfsetstrokecolor{currentstroke}%
\pgfsetstrokeopacity{0.000000}%
\pgfsetdash{}{0pt}%
\pgfpathmoveto{\pgfqpoint{0.559744in}{1.489057in}}%
\pgfpathlineto{\pgfqpoint{0.591784in}{1.480141in}}%
\pgfpathlineto{\pgfqpoint{0.644893in}{1.466447in}}%
\pgfpathlineto{\pgfqpoint{0.699946in}{1.452341in}}%
\pgfpathlineto{\pgfqpoint{0.750689in}{1.440103in}}%
\pgfpathlineto{\pgfqpoint{0.794534in}{1.430177in}}%
\pgfpathlineto{\pgfqpoint{0.841442in}{1.419862in}}%
\pgfpathlineto{\pgfqpoint{0.827889in}{1.355901in}}%
\pgfpathlineto{\pgfqpoint{0.815815in}{1.299101in}}%
\pgfpathlineto{\pgfqpoint{0.796008in}{1.207454in}}%
\pgfpathlineto{\pgfqpoint{0.781137in}{1.138280in}}%
\pgfpathlineto{\pgfqpoint{0.773103in}{1.099678in}}%
\pgfpathlineto{\pgfqpoint{0.762801in}{1.049576in}}%
\pgfpathlineto{\pgfqpoint{0.755674in}{1.039412in}}%
\pgfpathlineto{\pgfqpoint{0.749938in}{1.039072in}}%
\pgfpathlineto{\pgfqpoint{0.745846in}{1.047945in}}%
\pgfpathlineto{\pgfqpoint{0.736451in}{1.051167in}}%
\pgfpathlineto{\pgfqpoint{0.726328in}{1.050037in}}%
\pgfpathlineto{\pgfqpoint{0.723454in}{1.042780in}}%
\pgfpathlineto{\pgfqpoint{0.723155in}{1.017578in}}%
\pgfpathlineto{\pgfqpoint{0.720123in}{1.011825in}}%
\pgfpathlineto{\pgfqpoint{0.721856in}{0.991592in}}%
\pgfpathlineto{\pgfqpoint{0.715525in}{0.978113in}}%
\pgfpathlineto{\pgfqpoint{0.664188in}{1.057164in}}%
\pgfpathlineto{\pgfqpoint{0.613630in}{1.133980in}}%
\pgfpathlineto{\pgfqpoint{0.587312in}{1.174260in}}%
\pgfpathlineto{\pgfqpoint{0.565358in}{1.209028in}}%
\pgfpathlineto{\pgfqpoint{0.537695in}{1.252036in}}%
\pgfpathlineto{\pgfqpoint{0.506381in}{1.300279in}}%
\pgfpathlineto{\pgfqpoint{0.519255in}{1.346069in}}%
\pgfpathlineto{\pgfqpoint{0.545163in}{1.437904in}}%
\pgfpathlineto{\pgfqpoint{0.559744in}{1.489057in}}%
\pgfpathclose%
\pgfusepath{fill}%
\end{pgfscope}%
\begin{pgfscope}%
\pgfpathrectangle{\pgfqpoint{0.100000in}{0.100000in}}{\pgfqpoint{2.989028in}{1.913466in}}%
\pgfusepath{clip}%
\pgfsetbuttcap%
\pgfsetmiterjoin%
\definecolor{currentfill}{rgb}{0.232065,0.573472,0.724567}%
\pgfsetfillcolor{currentfill}%
\pgfsetlinewidth{0.000000pt}%
\definecolor{currentstroke}{rgb}{0.000000,0.000000,0.000000}%
\pgfsetstrokecolor{currentstroke}%
\pgfsetstrokeopacity{0.000000}%
\pgfsetdash{}{0pt}%
\pgfpathmoveto{\pgfqpoint{0.841442in}{1.419862in}}%
\pgfpathlineto{\pgfqpoint{0.891592in}{1.409886in}}%
\pgfpathlineto{\pgfqpoint{0.984436in}{1.391963in}}%
\pgfpathlineto{\pgfqpoint{0.972815in}{1.327496in}}%
\pgfpathlineto{\pgfqpoint{1.023928in}{1.318744in}}%
\pgfpathlineto{\pgfqpoint{1.070495in}{1.311293in}}%
\pgfpathlineto{\pgfqpoint{1.062404in}{1.260326in}}%
\pgfpathlineto{\pgfqpoint{1.053829in}{1.205365in}}%
\pgfpathlineto{\pgfqpoint{1.042350in}{1.133200in}}%
\pgfpathlineto{\pgfqpoint{1.042053in}{1.127151in}}%
\pgfpathlineto{\pgfqpoint{1.030137in}{1.052362in}}%
\pgfpathlineto{\pgfqpoint{0.981086in}{1.059964in}}%
\pgfpathlineto{\pgfqpoint{0.956091in}{1.064984in}}%
\pgfpathlineto{\pgfqpoint{0.865746in}{1.080865in}}%
\pgfpathlineto{\pgfqpoint{0.831724in}{1.087575in}}%
\pgfpathlineto{\pgfqpoint{0.773103in}{1.099678in}}%
\pgfpathlineto{\pgfqpoint{0.781137in}{1.138280in}}%
\pgfpathlineto{\pgfqpoint{0.796008in}{1.207454in}}%
\pgfpathlineto{\pgfqpoint{0.815815in}{1.299101in}}%
\pgfpathlineto{\pgfqpoint{0.827889in}{1.355901in}}%
\pgfpathlineto{\pgfqpoint{0.841442in}{1.419862in}}%
\pgfpathclose%
\pgfusepath{fill}%
\end{pgfscope}%
\begin{pgfscope}%
\pgfpathrectangle{\pgfqpoint{0.100000in}{0.100000in}}{\pgfqpoint{2.989028in}{1.913466in}}%
\pgfusepath{clip}%
\pgfsetbuttcap%
\pgfsetmiterjoin%
\definecolor{currentfill}{rgb}{0.400000,0.760784,0.647059}%
\pgfsetfillcolor{currentfill}%
\pgfsetlinewidth{0.000000pt}%
\definecolor{currentstroke}{rgb}{0.000000,0.000000,0.000000}%
\pgfsetstrokecolor{currentstroke}%
\pgfsetstrokeopacity{0.000000}%
\pgfsetdash{}{0pt}%
\pgfpathmoveto{\pgfqpoint{0.363713in}{1.549891in}}%
\pgfpathlineto{\pgfqpoint{0.381532in}{1.543676in}}%
\pgfpathlineto{\pgfqpoint{0.452913in}{1.521452in}}%
\pgfpathlineto{\pgfqpoint{0.518578in}{1.500784in}}%
\pgfpathlineto{\pgfqpoint{0.559744in}{1.489057in}}%
\pgfpathlineto{\pgfqpoint{0.545163in}{1.437904in}}%
\pgfpathlineto{\pgfqpoint{0.519255in}{1.346069in}}%
\pgfpathlineto{\pgfqpoint{0.506381in}{1.300279in}}%
\pgfpathlineto{\pgfqpoint{0.537695in}{1.252036in}}%
\pgfpathlineto{\pgfqpoint{0.565358in}{1.209028in}}%
\pgfpathlineto{\pgfqpoint{0.587312in}{1.174260in}}%
\pgfpathlineto{\pgfqpoint{0.613630in}{1.133980in}}%
\pgfpathlineto{\pgfqpoint{0.664188in}{1.057164in}}%
\pgfpathlineto{\pgfqpoint{0.715525in}{0.978113in}}%
\pgfpathlineto{\pgfqpoint{0.713464in}{0.970274in}}%
\pgfpathlineto{\pgfqpoint{0.719656in}{0.957789in}}%
\pgfpathlineto{\pgfqpoint{0.720881in}{0.940711in}}%
\pgfpathlineto{\pgfqpoint{0.729874in}{0.932426in}}%
\pgfpathlineto{\pgfqpoint{0.730230in}{0.925750in}}%
\pgfpathlineto{\pgfqpoint{0.714127in}{0.918174in}}%
\pgfpathlineto{\pgfqpoint{0.706455in}{0.910589in}}%
\pgfpathlineto{\pgfqpoint{0.704063in}{0.893839in}}%
\pgfpathlineto{\pgfqpoint{0.700180in}{0.884704in}}%
\pgfpathlineto{\pgfqpoint{0.691998in}{0.877004in}}%
\pgfpathlineto{\pgfqpoint{0.683861in}{0.856982in}}%
\pgfpathlineto{\pgfqpoint{0.695303in}{0.846552in}}%
\pgfpathlineto{\pgfqpoint{0.693826in}{0.837954in}}%
\pgfpathlineto{\pgfqpoint{0.684444in}{0.831972in}}%
\pgfpathlineto{\pgfqpoint{0.677975in}{0.833027in}}%
\pgfpathlineto{\pgfqpoint{0.601622in}{0.843623in}}%
\pgfpathlineto{\pgfqpoint{0.545229in}{0.851623in}}%
\pgfpathlineto{\pgfqpoint{0.547695in}{0.860728in}}%
\pgfpathlineto{\pgfqpoint{0.544024in}{0.875844in}}%
\pgfpathlineto{\pgfqpoint{0.543653in}{0.891092in}}%
\pgfpathlineto{\pgfqpoint{0.541254in}{0.900014in}}%
\pgfpathlineto{\pgfqpoint{0.533864in}{0.912772in}}%
\pgfpathlineto{\pgfqpoint{0.512605in}{0.942205in}}%
\pgfpathlineto{\pgfqpoint{0.505643in}{0.945811in}}%
\pgfpathlineto{\pgfqpoint{0.496672in}{0.945755in}}%
\pgfpathlineto{\pgfqpoint{0.498721in}{0.955040in}}%
\pgfpathlineto{\pgfqpoint{0.494473in}{0.966651in}}%
\pgfpathlineto{\pgfqpoint{0.473608in}{0.972422in}}%
\pgfpathlineto{\pgfqpoint{0.460889in}{0.983108in}}%
\pgfpathlineto{\pgfqpoint{0.459834in}{0.989647in}}%
\pgfpathlineto{\pgfqpoint{0.445189in}{1.005835in}}%
\pgfpathlineto{\pgfqpoint{0.431249in}{1.008932in}}%
\pgfpathlineto{\pgfqpoint{0.418332in}{1.017159in}}%
\pgfpathlineto{\pgfqpoint{0.401344in}{1.020005in}}%
\pgfpathlineto{\pgfqpoint{0.394086in}{1.030977in}}%
\pgfpathlineto{\pgfqpoint{0.400979in}{1.048345in}}%
\pgfpathlineto{\pgfqpoint{0.398871in}{1.052250in}}%
\pgfpathlineto{\pgfqpoint{0.404611in}{1.066728in}}%
\pgfpathlineto{\pgfqpoint{0.394400in}{1.074432in}}%
\pgfpathlineto{\pgfqpoint{0.397713in}{1.088401in}}%
\pgfpathlineto{\pgfqpoint{0.392259in}{1.091971in}}%
\pgfpathlineto{\pgfqpoint{0.387550in}{1.105174in}}%
\pgfpathlineto{\pgfqpoint{0.381867in}{1.109200in}}%
\pgfpathlineto{\pgfqpoint{0.381440in}{1.118736in}}%
\pgfpathlineto{\pgfqpoint{0.376988in}{1.125461in}}%
\pgfpathlineto{\pgfqpoint{0.370279in}{1.148109in}}%
\pgfpathlineto{\pgfqpoint{0.362938in}{1.158956in}}%
\pgfpathlineto{\pgfqpoint{0.364533in}{1.177370in}}%
\pgfpathlineto{\pgfqpoint{0.373166in}{1.179223in}}%
\pgfpathlineto{\pgfqpoint{0.378787in}{1.189211in}}%
\pgfpathlineto{\pgfqpoint{0.375390in}{1.200043in}}%
\pgfpathlineto{\pgfqpoint{0.366209in}{1.201851in}}%
\pgfpathlineto{\pgfqpoint{0.361629in}{1.206918in}}%
\pgfpathlineto{\pgfqpoint{0.354211in}{1.225555in}}%
\pgfpathlineto{\pgfqpoint{0.357678in}{1.232248in}}%
\pgfpathlineto{\pgfqpoint{0.355219in}{1.244711in}}%
\pgfpathlineto{\pgfqpoint{0.360664in}{1.260852in}}%
\pgfpathlineto{\pgfqpoint{0.367043in}{1.254907in}}%
\pgfpathlineto{\pgfqpoint{0.364148in}{1.247890in}}%
\pgfpathlineto{\pgfqpoint{0.375202in}{1.236766in}}%
\pgfpathlineto{\pgfqpoint{0.374497in}{1.253285in}}%
\pgfpathlineto{\pgfqpoint{0.369760in}{1.257715in}}%
\pgfpathlineto{\pgfqpoint{0.375174in}{1.272240in}}%
\pgfpathlineto{\pgfqpoint{0.395652in}{1.269927in}}%
\pgfpathlineto{\pgfqpoint{0.392896in}{1.275328in}}%
\pgfpathlineto{\pgfqpoint{0.379381in}{1.274797in}}%
\pgfpathlineto{\pgfqpoint{0.372951in}{1.282977in}}%
\pgfpathlineto{\pgfqpoint{0.364858in}{1.275713in}}%
\pgfpathlineto{\pgfqpoint{0.360564in}{1.263570in}}%
\pgfpathlineto{\pgfqpoint{0.349102in}{1.279909in}}%
\pgfpathlineto{\pgfqpoint{0.344687in}{1.282884in}}%
\pgfpathlineto{\pgfqpoint{0.346369in}{1.300667in}}%
\pgfpathlineto{\pgfqpoint{0.342875in}{1.311153in}}%
\pgfpathlineto{\pgfqpoint{0.336542in}{1.321006in}}%
\pgfpathlineto{\pgfqpoint{0.323573in}{1.351193in}}%
\pgfpathlineto{\pgfqpoint{0.327813in}{1.357869in}}%
\pgfpathlineto{\pgfqpoint{0.327764in}{1.378965in}}%
\pgfpathlineto{\pgfqpoint{0.334763in}{1.390732in}}%
\pgfpathlineto{\pgfqpoint{0.336368in}{1.409121in}}%
\pgfpathlineto{\pgfqpoint{0.329768in}{1.430191in}}%
\pgfpathlineto{\pgfqpoint{0.321012in}{1.443602in}}%
\pgfpathlineto{\pgfqpoint{0.322574in}{1.455704in}}%
\pgfpathlineto{\pgfqpoint{0.347163in}{1.484994in}}%
\pgfpathlineto{\pgfqpoint{0.348384in}{1.494985in}}%
\pgfpathlineto{\pgfqpoint{0.359464in}{1.514044in}}%
\pgfpathlineto{\pgfqpoint{0.361001in}{1.532109in}}%
\pgfpathlineto{\pgfqpoint{0.357439in}{1.536719in}}%
\pgfpathlineto{\pgfqpoint{0.363713in}{1.549891in}}%
\pgfpathclose%
\pgfusepath{fill}%
\end{pgfscope}%
\begin{pgfscope}%
\pgfpathrectangle{\pgfqpoint{0.100000in}{0.100000in}}{\pgfqpoint{2.989028in}{1.913466in}}%
\pgfusepath{clip}%
\pgfsetbuttcap%
\pgfsetmiterjoin%
\definecolor{currentfill}{rgb}{0.431834,0.773241,0.646597}%
\pgfsetfillcolor{currentfill}%
\pgfsetlinewidth{0.000000pt}%
\definecolor{currentstroke}{rgb}{0.000000,0.000000,0.000000}%
\pgfsetstrokecolor{currentstroke}%
\pgfsetstrokeopacity{0.000000}%
\pgfsetdash{}{0pt}%
\pgfpathmoveto{\pgfqpoint{2.278860in}{1.162554in}}%
\pgfpathlineto{\pgfqpoint{2.270539in}{1.244011in}}%
\pgfpathlineto{\pgfqpoint{2.260872in}{1.331556in}}%
\pgfpathlineto{\pgfqpoint{2.324180in}{1.340984in}}%
\pgfpathlineto{\pgfqpoint{2.340982in}{1.336611in}}%
\pgfpathlineto{\pgfqpoint{2.349041in}{1.331864in}}%
\pgfpathlineto{\pgfqpoint{2.359139in}{1.333180in}}%
\pgfpathlineto{\pgfqpoint{2.372453in}{1.325321in}}%
\pgfpathlineto{\pgfqpoint{2.397252in}{1.337018in}}%
\pgfpathlineto{\pgfqpoint{2.410938in}{1.337409in}}%
\pgfpathlineto{\pgfqpoint{2.426920in}{1.355239in}}%
\pgfpathlineto{\pgfqpoint{2.443180in}{1.366084in}}%
\pgfpathlineto{\pgfqpoint{2.464869in}{1.378577in}}%
\pgfpathlineto{\pgfqpoint{2.478831in}{1.291252in}}%
\pgfpathlineto{\pgfqpoint{2.472344in}{1.285643in}}%
\pgfpathlineto{\pgfqpoint{2.476532in}{1.280492in}}%
\pgfpathlineto{\pgfqpoint{2.478201in}{1.269199in}}%
\pgfpathlineto{\pgfqpoint{2.474951in}{1.258593in}}%
\pgfpathlineto{\pgfqpoint{2.474521in}{1.243089in}}%
\pgfpathlineto{\pgfqpoint{2.471469in}{1.222911in}}%
\pgfpathlineto{\pgfqpoint{2.456028in}{1.204915in}}%
\pgfpathlineto{\pgfqpoint{2.449528in}{1.201094in}}%
\pgfpathlineto{\pgfqpoint{2.444473in}{1.204383in}}%
\pgfpathlineto{\pgfqpoint{2.431963in}{1.187230in}}%
\pgfpathlineto{\pgfqpoint{2.433135in}{1.170714in}}%
\pgfpathlineto{\pgfqpoint{2.419203in}{1.174686in}}%
\pgfpathlineto{\pgfqpoint{2.412606in}{1.158191in}}%
\pgfpathlineto{\pgfqpoint{2.415941in}{1.146502in}}%
\pgfpathlineto{\pgfqpoint{2.410720in}{1.144776in}}%
\pgfpathlineto{\pgfqpoint{2.409990in}{1.135540in}}%
\pgfpathlineto{\pgfqpoint{2.397176in}{1.131815in}}%
\pgfpathlineto{\pgfqpoint{2.390518in}{1.139272in}}%
\pgfpathlineto{\pgfqpoint{2.381950in}{1.142163in}}%
\pgfpathlineto{\pgfqpoint{2.379937in}{1.149724in}}%
\pgfpathlineto{\pgfqpoint{2.372191in}{1.148394in}}%
\pgfpathlineto{\pgfqpoint{2.367116in}{1.141440in}}%
\pgfpathlineto{\pgfqpoint{2.359855in}{1.138994in}}%
\pgfpathlineto{\pgfqpoint{2.347031in}{1.143916in}}%
\pgfpathlineto{\pgfqpoint{2.339938in}{1.138060in}}%
\pgfpathlineto{\pgfqpoint{2.329857in}{1.144988in}}%
\pgfpathlineto{\pgfqpoint{2.310517in}{1.147114in}}%
\pgfpathlineto{\pgfqpoint{2.304684in}{1.159700in}}%
\pgfpathlineto{\pgfqpoint{2.297293in}{1.165276in}}%
\pgfpathlineto{\pgfqpoint{2.290128in}{1.161707in}}%
\pgfpathlineto{\pgfqpoint{2.278860in}{1.162554in}}%
\pgfpathclose%
\pgfusepath{fill}%
\end{pgfscope}%
\begin{pgfscope}%
\pgfpathrectangle{\pgfqpoint{0.100000in}{0.100000in}}{\pgfqpoint{2.989028in}{1.913466in}}%
\pgfusepath{clip}%
\pgfsetbuttcap%
\pgfsetmiterjoin%
\definecolor{currentfill}{rgb}{0.506113,0.802307,0.645521}%
\pgfsetfillcolor{currentfill}%
\pgfsetlinewidth{0.000000pt}%
\definecolor{currentstroke}{rgb}{0.000000,0.000000,0.000000}%
\pgfsetstrokecolor{currentstroke}%
\pgfsetstrokeopacity{0.000000}%
\pgfsetdash{}{0pt}%
\pgfpathmoveto{\pgfqpoint{2.069266in}{1.005020in}}%
\pgfpathlineto{\pgfqpoint{2.061849in}{1.011045in}}%
\pgfpathlineto{\pgfqpoint{2.055801in}{1.008178in}}%
\pgfpathlineto{\pgfqpoint{2.048069in}{1.022616in}}%
\pgfpathlineto{\pgfqpoint{2.052019in}{1.031690in}}%
\pgfpathlineto{\pgfqpoint{2.046326in}{1.041904in}}%
\pgfpathlineto{\pgfqpoint{2.046020in}{1.049946in}}%
\pgfpathlineto{\pgfqpoint{2.038328in}{1.052811in}}%
\pgfpathlineto{\pgfqpoint{2.034760in}{1.058870in}}%
\pgfpathlineto{\pgfqpoint{2.019719in}{1.066437in}}%
\pgfpathlineto{\pgfqpoint{2.006657in}{1.075760in}}%
\pgfpathlineto{\pgfqpoint{2.000587in}{1.082799in}}%
\pgfpathlineto{\pgfqpoint{2.000462in}{1.091383in}}%
\pgfpathlineto{\pgfqpoint{2.008584in}{1.107865in}}%
\pgfpathlineto{\pgfqpoint{2.011082in}{1.120472in}}%
\pgfpathlineto{\pgfqpoint{2.004465in}{1.127595in}}%
\pgfpathlineto{\pgfqpoint{1.995702in}{1.130279in}}%
\pgfpathlineto{\pgfqpoint{1.985070in}{1.124410in}}%
\pgfpathlineto{\pgfqpoint{1.980570in}{1.134486in}}%
\pgfpathlineto{\pgfqpoint{1.978299in}{1.148144in}}%
\pgfpathlineto{\pgfqpoint{1.962630in}{1.160319in}}%
\pgfpathlineto{\pgfqpoint{1.959499in}{1.165722in}}%
\pgfpathlineto{\pgfqpoint{1.945188in}{1.177954in}}%
\pgfpathlineto{\pgfqpoint{1.940703in}{1.186846in}}%
\pgfpathlineto{\pgfqpoint{1.936659in}{1.204480in}}%
\pgfpathlineto{\pgfqpoint{1.939431in}{1.220143in}}%
\pgfpathlineto{\pgfqpoint{1.943130in}{1.222326in}}%
\pgfpathlineto{\pgfqpoint{1.942465in}{1.235438in}}%
\pgfpathlineto{\pgfqpoint{1.952907in}{1.239332in}}%
\pgfpathlineto{\pgfqpoint{1.956060in}{1.251098in}}%
\pgfpathlineto{\pgfqpoint{1.962068in}{1.259018in}}%
\pgfpathlineto{\pgfqpoint{1.961764in}{1.269093in}}%
\pgfpathlineto{\pgfqpoint{1.954318in}{1.277106in}}%
\pgfpathlineto{\pgfqpoint{1.956087in}{1.288322in}}%
\pgfpathlineto{\pgfqpoint{1.975389in}{1.293200in}}%
\pgfpathlineto{\pgfqpoint{1.990197in}{1.302104in}}%
\pgfpathlineto{\pgfqpoint{1.991751in}{1.313315in}}%
\pgfpathlineto{\pgfqpoint{1.996904in}{1.316845in}}%
\pgfpathlineto{\pgfqpoint{1.998875in}{1.328623in}}%
\pgfpathlineto{\pgfqpoint{1.997291in}{1.336402in}}%
\pgfpathlineto{\pgfqpoint{1.987167in}{1.342864in}}%
\pgfpathlineto{\pgfqpoint{1.983091in}{1.352497in}}%
\pgfpathlineto{\pgfqpoint{1.973100in}{1.361809in}}%
\pgfpathlineto{\pgfqpoint{2.055206in}{1.365308in}}%
\pgfpathlineto{\pgfqpoint{2.110291in}{1.369221in}}%
\pgfpathlineto{\pgfqpoint{2.109283in}{1.357634in}}%
\pgfpathlineto{\pgfqpoint{2.118671in}{1.341651in}}%
\pgfpathlineto{\pgfqpoint{2.122612in}{1.327996in}}%
\pgfpathlineto{\pgfqpoint{2.127306in}{1.320243in}}%
\pgfpathlineto{\pgfqpoint{2.132431in}{1.256358in}}%
\pgfpathlineto{\pgfqpoint{2.139547in}{1.165295in}}%
\pgfpathlineto{\pgfqpoint{2.134057in}{1.151457in}}%
\pgfpathlineto{\pgfqpoint{2.141966in}{1.140083in}}%
\pgfpathlineto{\pgfqpoint{2.144170in}{1.129219in}}%
\pgfpathlineto{\pgfqpoint{2.133263in}{1.106692in}}%
\pgfpathlineto{\pgfqpoint{2.133907in}{1.103580in}}%
\pgfpathlineto{\pgfqpoint{2.122681in}{1.090668in}}%
\pgfpathlineto{\pgfqpoint{2.125759in}{1.086724in}}%
\pgfpathlineto{\pgfqpoint{2.121088in}{1.078676in}}%
\pgfpathlineto{\pgfqpoint{2.122290in}{1.062454in}}%
\pgfpathlineto{\pgfqpoint{2.116617in}{1.052499in}}%
\pgfpathlineto{\pgfqpoint{2.121231in}{1.040736in}}%
\pgfpathlineto{\pgfqpoint{2.101900in}{1.034344in}}%
\pgfpathlineto{\pgfqpoint{2.100119in}{1.027393in}}%
\pgfpathlineto{\pgfqpoint{2.105397in}{1.018582in}}%
\pgfpathlineto{\pgfqpoint{2.102969in}{1.012840in}}%
\pgfpathlineto{\pgfqpoint{2.082240in}{1.019933in}}%
\pgfpathlineto{\pgfqpoint{2.075401in}{1.020647in}}%
\pgfpathlineto{\pgfqpoint{2.066874in}{1.009863in}}%
\pgfpathlineto{\pgfqpoint{2.069266in}{1.005020in}}%
\pgfpathclose%
\pgfusepath{fill}%
\end{pgfscope}%
\begin{pgfscope}%
\pgfpathrectangle{\pgfqpoint{0.100000in}{0.100000in}}{\pgfqpoint{2.989028in}{1.913466in}}%
\pgfusepath{clip}%
\pgfsetbuttcap%
\pgfsetmiterjoin%
\definecolor{currentfill}{rgb}{0.774933,0.909112,0.621915}%
\pgfsetfillcolor{currentfill}%
\pgfsetlinewidth{0.000000pt}%
\definecolor{currentstroke}{rgb}{0.000000,0.000000,0.000000}%
\pgfsetstrokecolor{currentstroke}%
\pgfsetstrokeopacity{0.000000}%
\pgfsetdash{}{0pt}%
\pgfpathmoveto{\pgfqpoint{2.672645in}{1.203210in}}%
\pgfpathlineto{\pgfqpoint{2.666941in}{1.211681in}}%
\pgfpathlineto{\pgfqpoint{2.670160in}{1.216421in}}%
\pgfpathlineto{\pgfqpoint{2.678043in}{1.211098in}}%
\pgfpathlineto{\pgfqpoint{2.672645in}{1.203210in}}%
\pgfpathclose%
\pgfusepath{fill}%
\end{pgfscope}%
\begin{pgfscope}%
\pgfpathrectangle{\pgfqpoint{0.100000in}{0.100000in}}{\pgfqpoint{2.989028in}{1.913466in}}%
\pgfusepath{clip}%
\pgfsetbuttcap%
\pgfsetmiterjoin%
\definecolor{currentfill}{rgb}{0.654671,0.860438,0.643368}%
\pgfsetfillcolor{currentfill}%
\pgfsetlinewidth{0.000000pt}%
\definecolor{currentstroke}{rgb}{0.000000,0.000000,0.000000}%
\pgfsetstrokecolor{currentstroke}%
\pgfsetstrokeopacity{0.000000}%
\pgfsetdash{}{0pt}%
\pgfpathmoveto{\pgfqpoint{2.722458in}{1.276100in}}%
\pgfpathlineto{\pgfqpoint{2.725867in}{1.283302in}}%
\pgfpathlineto{\pgfqpoint{2.739648in}{1.284858in}}%
\pgfpathlineto{\pgfqpoint{2.737440in}{1.278737in}}%
\pgfpathlineto{\pgfqpoint{2.732893in}{1.270918in}}%
\pgfpathlineto{\pgfqpoint{2.735974in}{1.261602in}}%
\pgfpathlineto{\pgfqpoint{2.748133in}{1.250439in}}%
\pgfpathlineto{\pgfqpoint{2.750955in}{1.238685in}}%
\pgfpathlineto{\pgfqpoint{2.764963in}{1.224047in}}%
\pgfpathlineto{\pgfqpoint{2.770458in}{1.224674in}}%
\pgfpathlineto{\pgfqpoint{2.777304in}{1.202700in}}%
\pgfpathlineto{\pgfqpoint{2.776181in}{1.202481in}}%
\pgfpathlineto{\pgfqpoint{2.774933in}{1.202236in}}%
\pgfpathlineto{\pgfqpoint{2.744401in}{1.196391in}}%
\pgfpathlineto{\pgfqpoint{2.728066in}{1.254485in}}%
\pgfpathlineto{\pgfqpoint{2.722458in}{1.276100in}}%
\pgfpathclose%
\pgfusepath{fill}%
\end{pgfscope}%
\begin{pgfscope}%
\pgfpathrectangle{\pgfqpoint{0.100000in}{0.100000in}}{\pgfqpoint{2.989028in}{1.913466in}}%
\pgfusepath{clip}%
\pgfsetbuttcap%
\pgfsetmiterjoin%
\definecolor{currentfill}{rgb}{0.580392,0.831373,0.644444}%
\pgfsetfillcolor{currentfill}%
\pgfsetlinewidth{0.000000pt}%
\definecolor{currentstroke}{rgb}{0.000000,0.000000,0.000000}%
\pgfsetstrokecolor{currentstroke}%
\pgfsetstrokeopacity{0.000000}%
\pgfsetdash{}{0pt}%
\pgfpathmoveto{\pgfqpoint{2.437284in}{1.079184in}}%
\pgfpathlineto{\pgfqpoint{2.427636in}{1.079539in}}%
\pgfpathlineto{\pgfqpoint{2.418746in}{1.085662in}}%
\pgfpathlineto{\pgfqpoint{2.406736in}{1.104274in}}%
\pgfpathlineto{\pgfqpoint{2.396513in}{1.114131in}}%
\pgfpathlineto{\pgfqpoint{2.399185in}{1.121744in}}%
\pgfpathlineto{\pgfqpoint{2.397176in}{1.131815in}}%
\pgfpathlineto{\pgfqpoint{2.409990in}{1.135540in}}%
\pgfpathlineto{\pgfqpoint{2.410720in}{1.144776in}}%
\pgfpathlineto{\pgfqpoint{2.415941in}{1.146502in}}%
\pgfpathlineto{\pgfqpoint{2.412606in}{1.158191in}}%
\pgfpathlineto{\pgfqpoint{2.419203in}{1.174686in}}%
\pgfpathlineto{\pgfqpoint{2.433135in}{1.170714in}}%
\pgfpathlineto{\pgfqpoint{2.431963in}{1.187230in}}%
\pgfpathlineto{\pgfqpoint{2.444473in}{1.204383in}}%
\pgfpathlineto{\pgfqpoint{2.449528in}{1.201094in}}%
\pgfpathlineto{\pgfqpoint{2.456028in}{1.204915in}}%
\pgfpathlineto{\pgfqpoint{2.471469in}{1.222911in}}%
\pgfpathlineto{\pgfqpoint{2.474521in}{1.243089in}}%
\pgfpathlineto{\pgfqpoint{2.474951in}{1.258593in}}%
\pgfpathlineto{\pgfqpoint{2.478201in}{1.269199in}}%
\pgfpathlineto{\pgfqpoint{2.476532in}{1.280492in}}%
\pgfpathlineto{\pgfqpoint{2.472344in}{1.285643in}}%
\pgfpathlineto{\pgfqpoint{2.478831in}{1.291252in}}%
\pgfpathlineto{\pgfqpoint{2.488234in}{1.232010in}}%
\pgfpathlineto{\pgfqpoint{2.540072in}{1.240581in}}%
\pgfpathlineto{\pgfqpoint{2.545445in}{1.206768in}}%
\pgfpathlineto{\pgfqpoint{2.564213in}{1.229086in}}%
\pgfpathlineto{\pgfqpoint{2.568635in}{1.226861in}}%
\pgfpathlineto{\pgfqpoint{2.575140in}{1.239793in}}%
\pgfpathlineto{\pgfqpoint{2.584152in}{1.235738in}}%
\pgfpathlineto{\pgfqpoint{2.592010in}{1.236509in}}%
\pgfpathlineto{\pgfqpoint{2.597194in}{1.245499in}}%
\pgfpathlineto{\pgfqpoint{2.604719in}{1.250500in}}%
\pgfpathlineto{\pgfqpoint{2.615180in}{1.246096in}}%
\pgfpathlineto{\pgfqpoint{2.622064in}{1.247618in}}%
\pgfpathlineto{\pgfqpoint{2.631926in}{1.230550in}}%
\pgfpathlineto{\pgfqpoint{2.629074in}{1.217629in}}%
\pgfpathlineto{\pgfqpoint{2.603317in}{1.232177in}}%
\pgfpathlineto{\pgfqpoint{2.598492in}{1.220234in}}%
\pgfpathlineto{\pgfqpoint{2.600084in}{1.214739in}}%
\pgfpathlineto{\pgfqpoint{2.589121in}{1.196164in}}%
\pgfpathlineto{\pgfqpoint{2.583947in}{1.192470in}}%
\pgfpathlineto{\pgfqpoint{2.581611in}{1.184279in}}%
\pgfpathlineto{\pgfqpoint{2.572791in}{1.185136in}}%
\pgfpathlineto{\pgfqpoint{2.566456in}{1.162779in}}%
\pgfpathlineto{\pgfqpoint{2.562911in}{1.157648in}}%
\pgfpathlineto{\pgfqpoint{2.553799in}{1.159359in}}%
\pgfpathlineto{\pgfqpoint{2.544480in}{1.166411in}}%
\pgfpathlineto{\pgfqpoint{2.544157in}{1.155588in}}%
\pgfpathlineto{\pgfqpoint{2.535135in}{1.137399in}}%
\pgfpathlineto{\pgfqpoint{2.532891in}{1.124451in}}%
\pgfpathlineto{\pgfqpoint{2.525969in}{1.115827in}}%
\pgfpathlineto{\pgfqpoint{2.520721in}{1.102055in}}%
\pgfpathlineto{\pgfqpoint{2.520463in}{1.089306in}}%
\pgfpathlineto{\pgfqpoint{2.512358in}{1.086925in}}%
\pgfpathlineto{\pgfqpoint{2.503164in}{1.079708in}}%
\pgfpathlineto{\pgfqpoint{2.494311in}{1.078334in}}%
\pgfpathlineto{\pgfqpoint{2.492271in}{1.072225in}}%
\pgfpathlineto{\pgfqpoint{2.478007in}{1.065970in}}%
\pgfpathlineto{\pgfqpoint{2.470035in}{1.071300in}}%
\pgfpathlineto{\pgfqpoint{2.461126in}{1.061179in}}%
\pgfpathlineto{\pgfqpoint{2.445810in}{1.064162in}}%
\pgfpathlineto{\pgfqpoint{2.436416in}{1.074785in}}%
\pgfpathlineto{\pgfqpoint{2.437284in}{1.079184in}}%
\pgfpathclose%
\pgfusepath{fill}%
\end{pgfscope}%
\begin{pgfscope}%
\pgfpathrectangle{\pgfqpoint{0.100000in}{0.100000in}}{\pgfqpoint{2.989028in}{1.913466in}}%
\pgfusepath{clip}%
\pgfsetbuttcap%
\pgfsetmiterjoin%
\definecolor{currentfill}{rgb}{0.738639,0.894348,0.629296}%
\pgfsetfillcolor{currentfill}%
\pgfsetlinewidth{0.000000pt}%
\definecolor{currentstroke}{rgb}{0.000000,0.000000,0.000000}%
\pgfsetstrokecolor{currentstroke}%
\pgfsetstrokeopacity{0.000000}%
\pgfsetdash{}{0pt}%
\pgfpathmoveto{\pgfqpoint{2.774933in}{1.202236in}}%
\pgfpathlineto{\pgfqpoint{2.774542in}{1.190298in}}%
\pgfpathlineto{\pgfqpoint{2.769947in}{1.184429in}}%
\pgfpathlineto{\pgfqpoint{2.767015in}{1.171401in}}%
\pgfpathlineto{\pgfqpoint{2.753775in}{1.165395in}}%
\pgfpathlineto{\pgfqpoint{2.742667in}{1.163645in}}%
\pgfpathlineto{\pgfqpoint{2.745904in}{1.172214in}}%
\pgfpathlineto{\pgfqpoint{2.728800in}{1.179393in}}%
\pgfpathlineto{\pgfqpoint{2.714904in}{1.188383in}}%
\pgfpathlineto{\pgfqpoint{2.723507in}{1.201821in}}%
\pgfpathlineto{\pgfqpoint{2.716613in}{1.212266in}}%
\pgfpathlineto{\pgfqpoint{2.717424in}{1.221276in}}%
\pgfpathlineto{\pgfqpoint{2.712439in}{1.223749in}}%
\pgfpathlineto{\pgfqpoint{2.708394in}{1.238386in}}%
\pgfpathlineto{\pgfqpoint{2.717960in}{1.258229in}}%
\pgfpathlineto{\pgfqpoint{2.710773in}{1.261469in}}%
\pgfpathlineto{\pgfqpoint{2.708910in}{1.251680in}}%
\pgfpathlineto{\pgfqpoint{2.698651in}{1.248954in}}%
\pgfpathlineto{\pgfqpoint{2.699730in}{1.216746in}}%
\pgfpathlineto{\pgfqpoint{2.697854in}{1.206381in}}%
\pgfpathlineto{\pgfqpoint{2.703000in}{1.191607in}}%
\pgfpathlineto{\pgfqpoint{2.710916in}{1.184499in}}%
\pgfpathlineto{\pgfqpoint{2.707326in}{1.180036in}}%
\pgfpathlineto{\pgfqpoint{2.715407in}{1.173514in}}%
\pgfpathlineto{\pgfqpoint{2.718328in}{1.162929in}}%
\pgfpathlineto{\pgfqpoint{2.703517in}{1.171694in}}%
\pgfpathlineto{\pgfqpoint{2.694143in}{1.170556in}}%
\pgfpathlineto{\pgfqpoint{2.686865in}{1.179573in}}%
\pgfpathlineto{\pgfqpoint{2.679454in}{1.180451in}}%
\pgfpathlineto{\pgfqpoint{2.668921in}{1.175935in}}%
\pgfpathlineto{\pgfqpoint{2.664837in}{1.181563in}}%
\pgfpathlineto{\pgfqpoint{2.670205in}{1.193376in}}%
\pgfpathlineto{\pgfqpoint{2.672645in}{1.203210in}}%
\pgfpathlineto{\pgfqpoint{2.678043in}{1.211098in}}%
\pgfpathlineto{\pgfqpoint{2.670160in}{1.216421in}}%
\pgfpathlineto{\pgfqpoint{2.666941in}{1.211681in}}%
\pgfpathlineto{\pgfqpoint{2.659066in}{1.216492in}}%
\pgfpathlineto{\pgfqpoint{2.645120in}{1.219698in}}%
\pgfpathlineto{\pgfqpoint{2.646385in}{1.226747in}}%
\pgfpathlineto{\pgfqpoint{2.640064in}{1.230836in}}%
\pgfpathlineto{\pgfqpoint{2.631926in}{1.230550in}}%
\pgfpathlineto{\pgfqpoint{2.622064in}{1.247618in}}%
\pgfpathlineto{\pgfqpoint{2.615180in}{1.246096in}}%
\pgfpathlineto{\pgfqpoint{2.604719in}{1.250500in}}%
\pgfpathlineto{\pgfqpoint{2.597194in}{1.245499in}}%
\pgfpathlineto{\pgfqpoint{2.592010in}{1.236509in}}%
\pgfpathlineto{\pgfqpoint{2.584152in}{1.235738in}}%
\pgfpathlineto{\pgfqpoint{2.575140in}{1.239793in}}%
\pgfpathlineto{\pgfqpoint{2.568635in}{1.226861in}}%
\pgfpathlineto{\pgfqpoint{2.564213in}{1.229086in}}%
\pgfpathlineto{\pgfqpoint{2.545445in}{1.206768in}}%
\pgfpathlineto{\pgfqpoint{2.540072in}{1.240581in}}%
\pgfpathlineto{\pgfqpoint{2.608655in}{1.253292in}}%
\pgfpathlineto{\pgfqpoint{2.639399in}{1.258797in}}%
\pgfpathlineto{\pgfqpoint{2.684127in}{1.267806in}}%
\pgfpathlineto{\pgfqpoint{2.722458in}{1.276100in}}%
\pgfpathlineto{\pgfqpoint{2.728066in}{1.254485in}}%
\pgfpathlineto{\pgfqpoint{2.744401in}{1.196391in}}%
\pgfpathlineto{\pgfqpoint{2.774933in}{1.202236in}}%
\pgfpathclose%
\pgfusepath{fill}%
\end{pgfscope}%
\begin{pgfscope}%
\pgfpathrectangle{\pgfqpoint{0.100000in}{0.100000in}}{\pgfqpoint{2.989028in}{1.913466in}}%
\pgfusepath{clip}%
\pgfsetbuttcap%
\pgfsetmiterjoin%
\definecolor{currentfill}{rgb}{0.272049,0.618070,0.706113}%
\pgfsetfillcolor{currentfill}%
\pgfsetlinewidth{0.000000pt}%
\definecolor{currentstroke}{rgb}{0.000000,0.000000,0.000000}%
\pgfsetstrokecolor{currentstroke}%
\pgfsetstrokeopacity{0.000000}%
\pgfsetdash{}{0pt}%
\pgfpathmoveto{\pgfqpoint{1.394404in}{1.009408in}}%
\pgfpathlineto{\pgfqpoint{1.344368in}{1.014183in}}%
\pgfpathlineto{\pgfqpoint{1.292460in}{1.018783in}}%
\pgfpathlineto{\pgfqpoint{1.232477in}{1.025040in}}%
\pgfpathlineto{\pgfqpoint{1.143359in}{1.035639in}}%
\pgfpathlineto{\pgfqpoint{1.111744in}{1.040510in}}%
\pgfpathlineto{\pgfqpoint{1.030137in}{1.052362in}}%
\pgfpathlineto{\pgfqpoint{1.042053in}{1.127151in}}%
\pgfpathlineto{\pgfqpoint{1.042350in}{1.133200in}}%
\pgfpathlineto{\pgfqpoint{1.053829in}{1.205365in}}%
\pgfpathlineto{\pgfqpoint{1.062404in}{1.260326in}}%
\pgfpathlineto{\pgfqpoint{1.070495in}{1.311293in}}%
\pgfpathlineto{\pgfqpoint{1.125778in}{1.303357in}}%
\pgfpathlineto{\pgfqpoint{1.177321in}{1.295959in}}%
\pgfpathlineto{\pgfqpoint{1.272072in}{1.284265in}}%
\pgfpathlineto{\pgfqpoint{1.315540in}{1.280295in}}%
\pgfpathlineto{\pgfqpoint{1.384404in}{1.273614in}}%
\pgfpathlineto{\pgfqpoint{1.414195in}{1.271196in}}%
\pgfpathlineto{\pgfqpoint{1.408935in}{1.205970in}}%
\pgfpathlineto{\pgfqpoint{1.404183in}{1.143169in}}%
\pgfpathlineto{\pgfqpoint{1.400359in}{1.092033in}}%
\pgfpathlineto{\pgfqpoint{1.394404in}{1.009408in}}%
\pgfpathclose%
\pgfusepath{fill}%
\end{pgfscope}%
\begin{pgfscope}%
\pgfpathrectangle{\pgfqpoint{0.100000in}{0.100000in}}{\pgfqpoint{2.989028in}{1.913466in}}%
\pgfusepath{clip}%
\pgfsetbuttcap%
\pgfsetmiterjoin%
\definecolor{currentfill}{rgb}{0.384006,0.742945,0.654441}%
\pgfsetfillcolor{currentfill}%
\pgfsetlinewidth{0.000000pt}%
\definecolor{currentstroke}{rgb}{0.000000,0.000000,0.000000}%
\pgfsetstrokecolor{currentstroke}%
\pgfsetstrokeopacity{0.000000}%
\pgfsetdash{}{0pt}%
\pgfpathmoveto{\pgfqpoint{2.056343in}{0.972329in}}%
\pgfpathlineto{\pgfqpoint{2.058030in}{0.979890in}}%
\pgfpathlineto{\pgfqpoint{2.066773in}{0.978188in}}%
\pgfpathlineto{\pgfqpoint{2.069266in}{1.005020in}}%
\pgfpathlineto{\pgfqpoint{2.066874in}{1.009863in}}%
\pgfpathlineto{\pgfqpoint{2.075401in}{1.020647in}}%
\pgfpathlineto{\pgfqpoint{2.082240in}{1.019933in}}%
\pgfpathlineto{\pgfqpoint{2.102969in}{1.012840in}}%
\pgfpathlineto{\pgfqpoint{2.105397in}{1.018582in}}%
\pgfpathlineto{\pgfqpoint{2.100119in}{1.027393in}}%
\pgfpathlineto{\pgfqpoint{2.101900in}{1.034344in}}%
\pgfpathlineto{\pgfqpoint{2.121231in}{1.040736in}}%
\pgfpathlineto{\pgfqpoint{2.116617in}{1.052499in}}%
\pgfpathlineto{\pgfqpoint{2.122290in}{1.062454in}}%
\pgfpathlineto{\pgfqpoint{2.132850in}{1.068041in}}%
\pgfpathlineto{\pgfqpoint{2.155006in}{1.073574in}}%
\pgfpathlineto{\pgfqpoint{2.168985in}{1.065222in}}%
\pgfpathlineto{\pgfqpoint{2.173226in}{1.073367in}}%
\pgfpathlineto{\pgfqpoint{2.184834in}{1.079056in}}%
\pgfpathlineto{\pgfqpoint{2.187829in}{1.074039in}}%
\pgfpathlineto{\pgfqpoint{2.199753in}{1.078025in}}%
\pgfpathlineto{\pgfqpoint{2.199000in}{1.084848in}}%
\pgfpathlineto{\pgfqpoint{2.209721in}{1.092656in}}%
\pgfpathlineto{\pgfqpoint{2.216048in}{1.084548in}}%
\pgfpathlineto{\pgfqpoint{2.224371in}{1.083732in}}%
\pgfpathlineto{\pgfqpoint{2.229903in}{1.089016in}}%
\pgfpathlineto{\pgfqpoint{2.229287in}{1.096531in}}%
\pgfpathlineto{\pgfqpoint{2.233989in}{1.103993in}}%
\pgfpathlineto{\pgfqpoint{2.240290in}{1.105610in}}%
\pgfpathlineto{\pgfqpoint{2.242820in}{1.115463in}}%
\pgfpathlineto{\pgfqpoint{2.251983in}{1.123981in}}%
\pgfpathlineto{\pgfqpoint{2.249225in}{1.132459in}}%
\pgfpathlineto{\pgfqpoint{2.258139in}{1.136680in}}%
\pgfpathlineto{\pgfqpoint{2.264096in}{1.134084in}}%
\pgfpathlineto{\pgfqpoint{2.272893in}{1.140679in}}%
\pgfpathlineto{\pgfqpoint{2.280755in}{1.142397in}}%
\pgfpathlineto{\pgfqpoint{2.280783in}{1.149238in}}%
\pgfpathlineto{\pgfqpoint{2.275273in}{1.158730in}}%
\pgfpathlineto{\pgfqpoint{2.278860in}{1.162554in}}%
\pgfpathlineto{\pgfqpoint{2.290128in}{1.161707in}}%
\pgfpathlineto{\pgfqpoint{2.297293in}{1.165276in}}%
\pgfpathlineto{\pgfqpoint{2.304684in}{1.159700in}}%
\pgfpathlineto{\pgfqpoint{2.310517in}{1.147114in}}%
\pgfpathlineto{\pgfqpoint{2.329857in}{1.144988in}}%
\pgfpathlineto{\pgfqpoint{2.339938in}{1.138060in}}%
\pgfpathlineto{\pgfqpoint{2.347031in}{1.143916in}}%
\pgfpathlineto{\pgfqpoint{2.359855in}{1.138994in}}%
\pgfpathlineto{\pgfqpoint{2.367116in}{1.141440in}}%
\pgfpathlineto{\pgfqpoint{2.372191in}{1.148394in}}%
\pgfpathlineto{\pgfqpoint{2.379937in}{1.149724in}}%
\pgfpathlineto{\pgfqpoint{2.381950in}{1.142163in}}%
\pgfpathlineto{\pgfqpoint{2.390518in}{1.139272in}}%
\pgfpathlineto{\pgfqpoint{2.397176in}{1.131815in}}%
\pgfpathlineto{\pgfqpoint{2.399185in}{1.121744in}}%
\pgfpathlineto{\pgfqpoint{2.396513in}{1.114131in}}%
\pgfpathlineto{\pgfqpoint{2.406736in}{1.104274in}}%
\pgfpathlineto{\pgfqpoint{2.418746in}{1.085662in}}%
\pgfpathlineto{\pgfqpoint{2.427636in}{1.079539in}}%
\pgfpathlineto{\pgfqpoint{2.437284in}{1.079184in}}%
\pgfpathlineto{\pgfqpoint{2.419469in}{1.058743in}}%
\pgfpathlineto{\pgfqpoint{2.401934in}{1.046371in}}%
\pgfpathlineto{\pgfqpoint{2.395481in}{1.036545in}}%
\pgfpathlineto{\pgfqpoint{2.395590in}{1.031213in}}%
\pgfpathlineto{\pgfqpoint{2.386092in}{1.027126in}}%
\pgfpathlineto{\pgfqpoint{2.383370in}{1.019436in}}%
\pgfpathlineto{\pgfqpoint{2.363584in}{1.011690in}}%
\pgfpathlineto{\pgfqpoint{2.356571in}{1.006660in}}%
\pgfpathlineto{\pgfqpoint{2.355622in}{1.005586in}}%
\pgfpathlineto{\pgfqpoint{2.298708in}{1.000190in}}%
\pgfpathlineto{\pgfqpoint{2.264337in}{0.997299in}}%
\pgfpathlineto{\pgfqpoint{2.207938in}{0.994107in}}%
\pgfpathlineto{\pgfqpoint{2.137674in}{0.987156in}}%
\pgfpathlineto{\pgfqpoint{2.126065in}{0.988764in}}%
\pgfpathlineto{\pgfqpoint{2.128480in}{0.976914in}}%
\pgfpathlineto{\pgfqpoint{2.056343in}{0.972329in}}%
\pgfpathclose%
\pgfusepath{fill}%
\end{pgfscope}%
\begin{pgfscope}%
\pgfpathrectangle{\pgfqpoint{0.100000in}{0.100000in}}{\pgfqpoint{2.989028in}{1.913466in}}%
\pgfusepath{clip}%
\pgfsetbuttcap%
\pgfsetmiterjoin%
\definecolor{currentfill}{rgb}{0.368627,0.309804,0.635294}%
\pgfsetfillcolor{currentfill}%
\pgfsetlinewidth{0.000000pt}%
\definecolor{currentstroke}{rgb}{0.000000,0.000000,0.000000}%
\pgfsetstrokecolor{currentstroke}%
\pgfsetstrokeopacity{0.000000}%
\pgfsetdash{}{0pt}%
\pgfpathmoveto{\pgfqpoint{1.394404in}{1.009408in}}%
\pgfpathlineto{\pgfqpoint{1.400359in}{1.092033in}}%
\pgfpathlineto{\pgfqpoint{1.404183in}{1.143169in}}%
\pgfpathlineto{\pgfqpoint{1.408935in}{1.205970in}}%
\pgfpathlineto{\pgfqpoint{1.474809in}{1.201345in}}%
\pgfpathlineto{\pgfqpoint{1.558468in}{1.196775in}}%
\pgfpathlineto{\pgfqpoint{1.615367in}{1.194592in}}%
\pgfpathlineto{\pgfqpoint{1.671933in}{1.192859in}}%
\pgfpathlineto{\pgfqpoint{1.723131in}{1.192066in}}%
\pgfpathlineto{\pgfqpoint{1.746810in}{1.192312in}}%
\pgfpathlineto{\pgfqpoint{1.757233in}{1.183806in}}%
\pgfpathlineto{\pgfqpoint{1.765397in}{1.185521in}}%
\pgfpathlineto{\pgfqpoint{1.765653in}{1.178102in}}%
\pgfpathlineto{\pgfqpoint{1.759590in}{1.165269in}}%
\pgfpathlineto{\pgfqpoint{1.760247in}{1.157161in}}%
\pgfpathlineto{\pgfqpoint{1.766052in}{1.151814in}}%
\pgfpathlineto{\pgfqpoint{1.771384in}{1.140674in}}%
\pgfpathlineto{\pgfqpoint{1.782201in}{1.134279in}}%
\pgfpathlineto{\pgfqpoint{1.781840in}{1.092287in}}%
\pgfpathlineto{\pgfqpoint{1.782158in}{0.995725in}}%
\pgfpathlineto{\pgfqpoint{1.709656in}{0.996060in}}%
\pgfpathlineto{\pgfqpoint{1.633309in}{0.997396in}}%
\pgfpathlineto{\pgfqpoint{1.577105in}{0.999379in}}%
\pgfpathlineto{\pgfqpoint{1.530191in}{1.001189in}}%
\pgfpathlineto{\pgfqpoint{1.451155in}{1.005849in}}%
\pgfpathlineto{\pgfqpoint{1.394404in}{1.009408in}}%
\pgfpathclose%
\pgfusepath{fill}%
\end{pgfscope}%
\begin{pgfscope}%
\pgfpathrectangle{\pgfqpoint{0.100000in}{0.100000in}}{\pgfqpoint{2.989028in}{1.913466in}}%
\pgfusepath{clip}%
\pgfsetbuttcap%
\pgfsetmiterjoin%
\definecolor{currentfill}{rgb}{0.495502,0.798155,0.645675}%
\pgfsetfillcolor{currentfill}%
\pgfsetlinewidth{0.000000pt}%
\definecolor{currentstroke}{rgb}{0.000000,0.000000,0.000000}%
\pgfsetstrokecolor{currentstroke}%
\pgfsetstrokeopacity{0.000000}%
\pgfsetdash{}{0pt}%
\pgfpathmoveto{\pgfqpoint{2.461189in}{1.020348in}}%
\pgfpathlineto{\pgfqpoint{2.380888in}{1.009068in}}%
\pgfpathlineto{\pgfqpoint{2.356571in}{1.006660in}}%
\pgfpathlineto{\pgfqpoint{2.363584in}{1.011690in}}%
\pgfpathlineto{\pgfqpoint{2.383370in}{1.019436in}}%
\pgfpathlineto{\pgfqpoint{2.386092in}{1.027126in}}%
\pgfpathlineto{\pgfqpoint{2.395590in}{1.031213in}}%
\pgfpathlineto{\pgfqpoint{2.395481in}{1.036545in}}%
\pgfpathlineto{\pgfqpoint{2.401934in}{1.046371in}}%
\pgfpathlineto{\pgfqpoint{2.419469in}{1.058743in}}%
\pgfpathlineto{\pgfqpoint{2.437284in}{1.079184in}}%
\pgfpathlineto{\pgfqpoint{2.436416in}{1.074785in}}%
\pgfpathlineto{\pgfqpoint{2.445810in}{1.064162in}}%
\pgfpathlineto{\pgfqpoint{2.461126in}{1.061179in}}%
\pgfpathlineto{\pgfqpoint{2.470035in}{1.071300in}}%
\pgfpathlineto{\pgfqpoint{2.478007in}{1.065970in}}%
\pgfpathlineto{\pgfqpoint{2.492271in}{1.072225in}}%
\pgfpathlineto{\pgfqpoint{2.494311in}{1.078334in}}%
\pgfpathlineto{\pgfqpoint{2.503164in}{1.079708in}}%
\pgfpathlineto{\pgfqpoint{2.512358in}{1.086925in}}%
\pgfpathlineto{\pgfqpoint{2.520463in}{1.089306in}}%
\pgfpathlineto{\pgfqpoint{2.520721in}{1.102055in}}%
\pgfpathlineto{\pgfqpoint{2.525969in}{1.115827in}}%
\pgfpathlineto{\pgfqpoint{2.532891in}{1.124451in}}%
\pgfpathlineto{\pgfqpoint{2.535135in}{1.137399in}}%
\pgfpathlineto{\pgfqpoint{2.544157in}{1.155588in}}%
\pgfpathlineto{\pgfqpoint{2.544480in}{1.166411in}}%
\pgfpathlineto{\pgfqpoint{2.553799in}{1.159359in}}%
\pgfpathlineto{\pgfqpoint{2.562911in}{1.157648in}}%
\pgfpathlineto{\pgfqpoint{2.566456in}{1.162779in}}%
\pgfpathlineto{\pgfqpoint{2.572791in}{1.185136in}}%
\pgfpathlineto{\pgfqpoint{2.581611in}{1.184279in}}%
\pgfpathlineto{\pgfqpoint{2.583947in}{1.192470in}}%
\pgfpathlineto{\pgfqpoint{2.589121in}{1.196164in}}%
\pgfpathlineto{\pgfqpoint{2.600084in}{1.214739in}}%
\pgfpathlineto{\pgfqpoint{2.598492in}{1.220234in}}%
\pgfpathlineto{\pgfqpoint{2.603317in}{1.232177in}}%
\pgfpathlineto{\pgfqpoint{2.629074in}{1.217629in}}%
\pgfpathlineto{\pgfqpoint{2.631926in}{1.230550in}}%
\pgfpathlineto{\pgfqpoint{2.640064in}{1.230836in}}%
\pgfpathlineto{\pgfqpoint{2.646385in}{1.226747in}}%
\pgfpathlineto{\pgfqpoint{2.645120in}{1.219698in}}%
\pgfpathlineto{\pgfqpoint{2.659066in}{1.216492in}}%
\pgfpathlineto{\pgfqpoint{2.666941in}{1.211681in}}%
\pgfpathlineto{\pgfqpoint{2.672645in}{1.203210in}}%
\pgfpathlineto{\pgfqpoint{2.670205in}{1.193376in}}%
\pgfpathlineto{\pgfqpoint{2.665277in}{1.192570in}}%
\pgfpathlineto{\pgfqpoint{2.662419in}{1.177731in}}%
\pgfpathlineto{\pgfqpoint{2.668680in}{1.171929in}}%
\pgfpathlineto{\pgfqpoint{2.677493in}{1.176619in}}%
\pgfpathlineto{\pgfqpoint{2.685671in}{1.166716in}}%
\pgfpathlineto{\pgfqpoint{2.703942in}{1.164936in}}%
\pgfpathlineto{\pgfqpoint{2.707089in}{1.159239in}}%
\pgfpathlineto{\pgfqpoint{2.723998in}{1.153699in}}%
\pgfpathlineto{\pgfqpoint{2.721912in}{1.147161in}}%
\pgfpathlineto{\pgfqpoint{2.724111in}{1.128683in}}%
\pgfpathlineto{\pgfqpoint{2.730936in}{1.121706in}}%
\pgfpathlineto{\pgfqpoint{2.720866in}{1.118761in}}%
\pgfpathlineto{\pgfqpoint{2.722201in}{1.110873in}}%
\pgfpathlineto{\pgfqpoint{2.732944in}{1.104185in}}%
\pgfpathlineto{\pgfqpoint{2.733912in}{1.097588in}}%
\pgfpathlineto{\pgfqpoint{2.727847in}{1.092630in}}%
\pgfpathlineto{\pgfqpoint{2.715609in}{1.104375in}}%
\pgfpathlineto{\pgfqpoint{2.714399in}{1.095806in}}%
\pgfpathlineto{\pgfqpoint{2.724644in}{1.091731in}}%
\pgfpathlineto{\pgfqpoint{2.725669in}{1.087517in}}%
\pgfpathlineto{\pgfqpoint{2.739721in}{1.093086in}}%
\pgfpathlineto{\pgfqpoint{2.750489in}{1.094559in}}%
\pgfpathlineto{\pgfqpoint{2.761520in}{1.072304in}}%
\pgfpathlineto{\pgfqpoint{2.760289in}{1.072063in}}%
\pgfpathlineto{\pgfqpoint{2.755305in}{1.071033in}}%
\pgfpathlineto{\pgfqpoint{2.753836in}{1.070726in}}%
\pgfpathlineto{\pgfqpoint{2.752865in}{1.070536in}}%
\pgfpathlineto{\pgfqpoint{2.687189in}{1.056913in}}%
\pgfpathlineto{\pgfqpoint{2.628447in}{1.044890in}}%
\pgfpathlineto{\pgfqpoint{2.547229in}{1.030898in}}%
\pgfpathlineto{\pgfqpoint{2.478251in}{1.021788in}}%
\pgfpathlineto{\pgfqpoint{2.461189in}{1.020348in}}%
\pgfpathclose%
\pgfusepath{fill}%
\end{pgfscope}%
\begin{pgfscope}%
\pgfpathrectangle{\pgfqpoint{0.100000in}{0.100000in}}{\pgfqpoint{2.989028in}{1.913466in}}%
\pgfusepath{clip}%
\pgfsetbuttcap%
\pgfsetmiterjoin%
\definecolor{currentfill}{rgb}{0.495502,0.798155,0.645675}%
\pgfsetfillcolor{currentfill}%
\pgfsetlinewidth{0.000000pt}%
\definecolor{currentstroke}{rgb}{0.000000,0.000000,0.000000}%
\pgfsetstrokecolor{currentstroke}%
\pgfsetstrokeopacity{0.000000}%
\pgfsetdash{}{0pt}%
\pgfpathmoveto{\pgfqpoint{2.753775in}{1.165395in}}%
\pgfpathlineto{\pgfqpoint{2.767015in}{1.171401in}}%
\pgfpathlineto{\pgfqpoint{2.756442in}{1.140471in}}%
\pgfpathlineto{\pgfqpoint{2.749450in}{1.124092in}}%
\pgfpathlineto{\pgfqpoint{2.744413in}{1.133361in}}%
\pgfpathlineto{\pgfqpoint{2.753370in}{1.155552in}}%
\pgfpathlineto{\pgfqpoint{2.753775in}{1.165395in}}%
\pgfpathclose%
\pgfusepath{fill}%
\end{pgfscope}%
\begin{pgfscope}%
\pgfpathrectangle{\pgfqpoint{0.100000in}{0.100000in}}{\pgfqpoint{2.989028in}{1.913466in}}%
\pgfusepath{clip}%
\pgfsetbuttcap%
\pgfsetmiterjoin%
\definecolor{currentfill}{rgb}{0.537947,0.814764,0.645060}%
\pgfsetfillcolor{currentfill}%
\pgfsetlinewidth{0.000000pt}%
\definecolor{currentstroke}{rgb}{0.000000,0.000000,0.000000}%
\pgfsetstrokecolor{currentstroke}%
\pgfsetstrokeopacity{0.000000}%
\pgfsetdash{}{0pt}%
\pgfpathmoveto{\pgfqpoint{2.056343in}{0.972329in}}%
\pgfpathlineto{\pgfqpoint{2.053146in}{0.971866in}}%
\pgfpathlineto{\pgfqpoint{2.050131in}{0.971655in}}%
\pgfpathlineto{\pgfqpoint{2.051421in}{0.962397in}}%
\pgfpathlineto{\pgfqpoint{2.043642in}{0.953066in}}%
\pgfpathlineto{\pgfqpoint{2.042109in}{0.938464in}}%
\pgfpathlineto{\pgfqpoint{2.007334in}{0.935892in}}%
\pgfpathlineto{\pgfqpoint{2.010365in}{0.942748in}}%
\pgfpathlineto{\pgfqpoint{2.022904in}{0.955269in}}%
\pgfpathlineto{\pgfqpoint{2.023251in}{0.962527in}}%
\pgfpathlineto{\pgfqpoint{2.017700in}{0.969399in}}%
\pgfpathlineto{\pgfqpoint{1.965941in}{0.966662in}}%
\pgfpathlineto{\pgfqpoint{1.875448in}{0.963789in}}%
\pgfpathlineto{\pgfqpoint{1.822491in}{0.962824in}}%
\pgfpathlineto{\pgfqpoint{1.782462in}{0.962463in}}%
\pgfpathlineto{\pgfqpoint{1.782158in}{0.995725in}}%
\pgfpathlineto{\pgfqpoint{1.781840in}{1.092287in}}%
\pgfpathlineto{\pgfqpoint{1.782201in}{1.134279in}}%
\pgfpathlineto{\pgfqpoint{1.771384in}{1.140674in}}%
\pgfpathlineto{\pgfqpoint{1.766052in}{1.151814in}}%
\pgfpathlineto{\pgfqpoint{1.760247in}{1.157161in}}%
\pgfpathlineto{\pgfqpoint{1.759590in}{1.165269in}}%
\pgfpathlineto{\pgfqpoint{1.765653in}{1.178102in}}%
\pgfpathlineto{\pgfqpoint{1.765397in}{1.185521in}}%
\pgfpathlineto{\pgfqpoint{1.757233in}{1.183806in}}%
\pgfpathlineto{\pgfqpoint{1.746810in}{1.192312in}}%
\pgfpathlineto{\pgfqpoint{1.738455in}{1.207242in}}%
\pgfpathlineto{\pgfqpoint{1.731447in}{1.214132in}}%
\pgfpathlineto{\pgfqpoint{1.724124in}{1.231060in}}%
\pgfpathlineto{\pgfqpoint{1.800184in}{1.229874in}}%
\pgfpathlineto{\pgfqpoint{1.875789in}{1.232365in}}%
\pgfpathlineto{\pgfqpoint{1.924266in}{1.235161in}}%
\pgfpathlineto{\pgfqpoint{1.939431in}{1.220143in}}%
\pgfpathlineto{\pgfqpoint{1.936659in}{1.204480in}}%
\pgfpathlineto{\pgfqpoint{1.940703in}{1.186846in}}%
\pgfpathlineto{\pgfqpoint{1.945188in}{1.177954in}}%
\pgfpathlineto{\pgfqpoint{1.959499in}{1.165722in}}%
\pgfpathlineto{\pgfqpoint{1.962630in}{1.160319in}}%
\pgfpathlineto{\pgfqpoint{1.978299in}{1.148144in}}%
\pgfpathlineto{\pgfqpoint{1.980570in}{1.134486in}}%
\pgfpathlineto{\pgfqpoint{1.985070in}{1.124410in}}%
\pgfpathlineto{\pgfqpoint{1.995702in}{1.130279in}}%
\pgfpathlineto{\pgfqpoint{2.004465in}{1.127595in}}%
\pgfpathlineto{\pgfqpoint{2.011082in}{1.120472in}}%
\pgfpathlineto{\pgfqpoint{2.008584in}{1.107865in}}%
\pgfpathlineto{\pgfqpoint{2.000462in}{1.091383in}}%
\pgfpathlineto{\pgfqpoint{2.000587in}{1.082799in}}%
\pgfpathlineto{\pgfqpoint{2.006657in}{1.075760in}}%
\pgfpathlineto{\pgfqpoint{2.019719in}{1.066437in}}%
\pgfpathlineto{\pgfqpoint{2.034760in}{1.058870in}}%
\pgfpathlineto{\pgfqpoint{2.038328in}{1.052811in}}%
\pgfpathlineto{\pgfqpoint{2.046020in}{1.049946in}}%
\pgfpathlineto{\pgfqpoint{2.046326in}{1.041904in}}%
\pgfpathlineto{\pgfqpoint{2.052019in}{1.031690in}}%
\pgfpathlineto{\pgfqpoint{2.048069in}{1.022616in}}%
\pgfpathlineto{\pgfqpoint{2.055801in}{1.008178in}}%
\pgfpathlineto{\pgfqpoint{2.061849in}{1.011045in}}%
\pgfpathlineto{\pgfqpoint{2.069266in}{1.005020in}}%
\pgfpathlineto{\pgfqpoint{2.066773in}{0.978188in}}%
\pgfpathlineto{\pgfqpoint{2.058030in}{0.979890in}}%
\pgfpathlineto{\pgfqpoint{2.056343in}{0.972329in}}%
\pgfpathclose%
\pgfusepath{fill}%
\end{pgfscope}%
\begin{pgfscope}%
\pgfpathrectangle{\pgfqpoint{0.100000in}{0.100000in}}{\pgfqpoint{2.989028in}{1.913466in}}%
\pgfusepath{clip}%
\pgfsetbuttcap%
\pgfsetmiterjoin%
\definecolor{currentfill}{rgb}{0.384006,0.742945,0.654441}%
\pgfsetfillcolor{currentfill}%
\pgfsetlinewidth{0.000000pt}%
\definecolor{currentstroke}{rgb}{0.000000,0.000000,0.000000}%
\pgfsetstrokecolor{currentstroke}%
\pgfsetstrokeopacity{0.000000}%
\pgfsetdash{}{0pt}%
\pgfpathmoveto{\pgfqpoint{0.715525in}{0.978113in}}%
\pgfpathlineto{\pgfqpoint{0.721856in}{0.991592in}}%
\pgfpathlineto{\pgfqpoint{0.720123in}{1.011825in}}%
\pgfpathlineto{\pgfqpoint{0.723155in}{1.017578in}}%
\pgfpathlineto{\pgfqpoint{0.723454in}{1.042780in}}%
\pgfpathlineto{\pgfqpoint{0.726328in}{1.050037in}}%
\pgfpathlineto{\pgfqpoint{0.736451in}{1.051167in}}%
\pgfpathlineto{\pgfqpoint{0.745846in}{1.047945in}}%
\pgfpathlineto{\pgfqpoint{0.749938in}{1.039072in}}%
\pgfpathlineto{\pgfqpoint{0.755674in}{1.039412in}}%
\pgfpathlineto{\pgfqpoint{0.762801in}{1.049576in}}%
\pgfpathlineto{\pgfqpoint{0.773103in}{1.099678in}}%
\pgfpathlineto{\pgfqpoint{0.831724in}{1.087575in}}%
\pgfpathlineto{\pgfqpoint{0.865746in}{1.080865in}}%
\pgfpathlineto{\pgfqpoint{0.956091in}{1.064984in}}%
\pgfpathlineto{\pgfqpoint{0.981086in}{1.059964in}}%
\pgfpathlineto{\pgfqpoint{1.030137in}{1.052362in}}%
\pgfpathlineto{\pgfqpoint{1.020079in}{0.987601in}}%
\pgfpathlineto{\pgfqpoint{1.009619in}{0.920049in}}%
\pgfpathlineto{\pgfqpoint{0.997573in}{0.844052in}}%
\pgfpathlineto{\pgfqpoint{0.984025in}{0.756815in}}%
\pgfpathlineto{\pgfqpoint{0.973080in}{0.685166in}}%
\pgfpathlineto{\pgfqpoint{0.894687in}{0.697614in}}%
\pgfpathlineto{\pgfqpoint{0.860244in}{0.703513in}}%
\pgfpathlineto{\pgfqpoint{0.844857in}{0.712725in}}%
\pgfpathlineto{\pgfqpoint{0.744535in}{0.773254in}}%
\pgfpathlineto{\pgfqpoint{0.669248in}{0.819189in}}%
\pgfpathlineto{\pgfqpoint{0.671756in}{0.827331in}}%
\pgfpathlineto{\pgfqpoint{0.677975in}{0.833027in}}%
\pgfpathlineto{\pgfqpoint{0.684444in}{0.831972in}}%
\pgfpathlineto{\pgfqpoint{0.693826in}{0.837954in}}%
\pgfpathlineto{\pgfqpoint{0.695303in}{0.846552in}}%
\pgfpathlineto{\pgfqpoint{0.683861in}{0.856982in}}%
\pgfpathlineto{\pgfqpoint{0.691998in}{0.877004in}}%
\pgfpathlineto{\pgfqpoint{0.700180in}{0.884704in}}%
\pgfpathlineto{\pgfqpoint{0.704063in}{0.893839in}}%
\pgfpathlineto{\pgfqpoint{0.706455in}{0.910589in}}%
\pgfpathlineto{\pgfqpoint{0.714127in}{0.918174in}}%
\pgfpathlineto{\pgfqpoint{0.730230in}{0.925750in}}%
\pgfpathlineto{\pgfqpoint{0.729874in}{0.932426in}}%
\pgfpathlineto{\pgfqpoint{0.720881in}{0.940711in}}%
\pgfpathlineto{\pgfqpoint{0.719656in}{0.957789in}}%
\pgfpathlineto{\pgfqpoint{0.713464in}{0.970274in}}%
\pgfpathlineto{\pgfqpoint{0.715525in}{0.978113in}}%
\pgfpathclose%
\pgfusepath{fill}%
\end{pgfscope}%
\begin{pgfscope}%
\pgfpathrectangle{\pgfqpoint{0.100000in}{0.100000in}}{\pgfqpoint{2.989028in}{1.913466in}}%
\pgfusepath{clip}%
\pgfsetbuttcap%
\pgfsetmiterjoin%
\definecolor{currentfill}{rgb}{0.248058,0.591311,0.717186}%
\pgfsetfillcolor{currentfill}%
\pgfsetlinewidth{0.000000pt}%
\definecolor{currentstroke}{rgb}{0.000000,0.000000,0.000000}%
\pgfsetstrokecolor{currentstroke}%
\pgfsetstrokeopacity{0.000000}%
\pgfsetdash{}{0pt}%
\pgfpathmoveto{\pgfqpoint{1.790948in}{0.774957in}}%
\pgfpathlineto{\pgfqpoint{1.776416in}{0.779451in}}%
\pgfpathlineto{\pgfqpoint{1.757795in}{0.793711in}}%
\pgfpathlineto{\pgfqpoint{1.749524in}{0.796776in}}%
\pgfpathlineto{\pgfqpoint{1.744268in}{0.790617in}}%
\pgfpathlineto{\pgfqpoint{1.737631in}{0.790306in}}%
\pgfpathlineto{\pgfqpoint{1.729234in}{0.795535in}}%
\pgfpathlineto{\pgfqpoint{1.716057in}{0.788835in}}%
\pgfpathlineto{\pgfqpoint{1.711367in}{0.792135in}}%
\pgfpathlineto{\pgfqpoint{1.698401in}{0.784318in}}%
\pgfpathlineto{\pgfqpoint{1.684722in}{0.785761in}}%
\pgfpathlineto{\pgfqpoint{1.674172in}{0.790840in}}%
\pgfpathlineto{\pgfqpoint{1.666783in}{0.788928in}}%
\pgfpathlineto{\pgfqpoint{1.654956in}{0.796112in}}%
\pgfpathlineto{\pgfqpoint{1.648377in}{0.783442in}}%
\pgfpathlineto{\pgfqpoint{1.641649in}{0.794338in}}%
\pgfpathlineto{\pgfqpoint{1.628360in}{0.790111in}}%
\pgfpathlineto{\pgfqpoint{1.616736in}{0.800460in}}%
\pgfpathlineto{\pgfqpoint{1.606579in}{0.792118in}}%
\pgfpathlineto{\pgfqpoint{1.601493in}{0.799784in}}%
\pgfpathlineto{\pgfqpoint{1.594161in}{0.802273in}}%
\pgfpathlineto{\pgfqpoint{1.589696in}{0.809664in}}%
\pgfpathlineto{\pgfqpoint{1.580106in}{0.811767in}}%
\pgfpathlineto{\pgfqpoint{1.574573in}{0.806204in}}%
\pgfpathlineto{\pgfqpoint{1.565171in}{0.813414in}}%
\pgfpathlineto{\pgfqpoint{1.560789in}{0.811767in}}%
\pgfpathlineto{\pgfqpoint{1.545207in}{0.817609in}}%
\pgfpathlineto{\pgfqpoint{1.535448in}{0.818261in}}%
\pgfpathlineto{\pgfqpoint{1.531080in}{0.830667in}}%
\pgfpathlineto{\pgfqpoint{1.523256in}{0.829122in}}%
\pgfpathlineto{\pgfqpoint{1.514296in}{0.832183in}}%
\pgfpathlineto{\pgfqpoint{1.508395in}{0.830414in}}%
\pgfpathlineto{\pgfqpoint{1.495741in}{0.844347in}}%
\pgfpathlineto{\pgfqpoint{1.492215in}{0.843436in}}%
\pgfpathlineto{\pgfqpoint{1.495417in}{0.899940in}}%
\pgfpathlineto{\pgfqpoint{1.498906in}{0.969876in}}%
\pgfpathlineto{\pgfqpoint{1.441639in}{0.973080in}}%
\pgfpathlineto{\pgfqpoint{1.385128in}{0.977352in}}%
\pgfpathlineto{\pgfqpoint{1.341469in}{0.981143in}}%
\pgfpathlineto{\pgfqpoint{1.344368in}{1.014183in}}%
\pgfpathlineto{\pgfqpoint{1.394404in}{1.009408in}}%
\pgfpathlineto{\pgfqpoint{1.451155in}{1.005849in}}%
\pgfpathlineto{\pgfqpoint{1.530191in}{1.001189in}}%
\pgfpathlineto{\pgfqpoint{1.577105in}{0.999379in}}%
\pgfpathlineto{\pgfqpoint{1.633309in}{0.997396in}}%
\pgfpathlineto{\pgfqpoint{1.709656in}{0.996060in}}%
\pgfpathlineto{\pgfqpoint{1.782158in}{0.995725in}}%
\pgfpathlineto{\pgfqpoint{1.782462in}{0.962463in}}%
\pgfpathlineto{\pgfqpoint{1.786531in}{0.937404in}}%
\pgfpathlineto{\pgfqpoint{1.792853in}{0.891120in}}%
\pgfpathlineto{\pgfqpoint{1.791549in}{0.812079in}}%
\pgfpathlineto{\pgfqpoint{1.790948in}{0.774957in}}%
\pgfpathclose%
\pgfusepath{fill}%
\end{pgfscope}%
\begin{pgfscope}%
\pgfpathrectangle{\pgfqpoint{0.100000in}{0.100000in}}{\pgfqpoint{2.989028in}{1.913466in}}%
\pgfusepath{clip}%
\pgfsetbuttcap%
\pgfsetmiterjoin%
\definecolor{currentfill}{rgb}{0.431834,0.773241,0.646597}%
\pgfsetfillcolor{currentfill}%
\pgfsetlinewidth{0.000000pt}%
\definecolor{currentstroke}{rgb}{0.000000,0.000000,0.000000}%
\pgfsetstrokecolor{currentstroke}%
\pgfsetstrokeopacity{0.000000}%
\pgfsetdash{}{0pt}%
\pgfpathmoveto{\pgfqpoint{2.335069in}{0.897708in}}%
\pgfpathlineto{\pgfqpoint{2.335112in}{0.912353in}}%
\pgfpathlineto{\pgfqpoint{2.347839in}{0.917977in}}%
\pgfpathlineto{\pgfqpoint{2.348372in}{0.926967in}}%
\pgfpathlineto{\pgfqpoint{2.359767in}{0.937949in}}%
\pgfpathlineto{\pgfqpoint{2.374011in}{0.940198in}}%
\pgfpathlineto{\pgfqpoint{2.386533in}{0.952301in}}%
\pgfpathlineto{\pgfqpoint{2.399610in}{0.957711in}}%
\pgfpathlineto{\pgfqpoint{2.409425in}{0.973951in}}%
\pgfpathlineto{\pgfqpoint{2.418364in}{0.972772in}}%
\pgfpathlineto{\pgfqpoint{2.437256in}{0.987570in}}%
\pgfpathlineto{\pgfqpoint{2.447232in}{0.987849in}}%
\pgfpathlineto{\pgfqpoint{2.451445in}{0.999123in}}%
\pgfpathlineto{\pgfqpoint{2.459375in}{1.006973in}}%
\pgfpathlineto{\pgfqpoint{2.461189in}{1.020348in}}%
\pgfpathlineto{\pgfqpoint{2.478251in}{1.021788in}}%
\pgfpathlineto{\pgfqpoint{2.547229in}{1.030898in}}%
\pgfpathlineto{\pgfqpoint{2.628447in}{1.044890in}}%
\pgfpathlineto{\pgfqpoint{2.687189in}{1.056913in}}%
\pgfpathlineto{\pgfqpoint{2.752865in}{1.070536in}}%
\pgfpathlineto{\pgfqpoint{2.771631in}{1.044770in}}%
\pgfpathlineto{\pgfqpoint{2.761466in}{1.046415in}}%
\pgfpathlineto{\pgfqpoint{2.748417in}{1.043236in}}%
\pgfpathlineto{\pgfqpoint{2.735559in}{1.030105in}}%
\pgfpathlineto{\pgfqpoint{2.726318in}{1.031050in}}%
\pgfpathlineto{\pgfqpoint{2.725149in}{1.023251in}}%
\pgfpathlineto{\pgfqpoint{2.741852in}{1.029418in}}%
\pgfpathlineto{\pgfqpoint{2.758660in}{1.031961in}}%
\pgfpathlineto{\pgfqpoint{2.764901in}{1.028577in}}%
\pgfpathlineto{\pgfqpoint{2.773300in}{1.032435in}}%
\pgfpathlineto{\pgfqpoint{2.777634in}{1.029732in}}%
\pgfpathlineto{\pgfqpoint{2.781470in}{1.016862in}}%
\pgfpathlineto{\pgfqpoint{2.773487in}{1.012876in}}%
\pgfpathlineto{\pgfqpoint{2.768008in}{0.997182in}}%
\pgfpathlineto{\pgfqpoint{2.762263in}{0.991072in}}%
\pgfpathlineto{\pgfqpoint{2.744653in}{0.992409in}}%
\pgfpathlineto{\pgfqpoint{2.745600in}{1.001624in}}%
\pgfpathlineto{\pgfqpoint{2.735982in}{0.997599in}}%
\pgfpathlineto{\pgfqpoint{2.733890in}{0.989913in}}%
\pgfpathlineto{\pgfqpoint{2.740483in}{0.981952in}}%
\pgfpathlineto{\pgfqpoint{2.742722in}{0.968440in}}%
\pgfpathlineto{\pgfqpoint{2.731935in}{0.960741in}}%
\pgfpathlineto{\pgfqpoint{2.743601in}{0.958032in}}%
\pgfpathlineto{\pgfqpoint{2.748881in}{0.963694in}}%
\pgfpathlineto{\pgfqpoint{2.760553in}{0.964383in}}%
\pgfpathlineto{\pgfqpoint{2.754555in}{0.952158in}}%
\pgfpathlineto{\pgfqpoint{2.747208in}{0.946268in}}%
\pgfpathlineto{\pgfqpoint{2.724984in}{0.940370in}}%
\pgfpathlineto{\pgfqpoint{2.702166in}{0.919675in}}%
\pgfpathlineto{\pgfqpoint{2.688096in}{0.899284in}}%
\pgfpathlineto{\pgfqpoint{2.688013in}{0.891003in}}%
\pgfpathlineto{\pgfqpoint{2.682393in}{0.879575in}}%
\pgfpathlineto{\pgfqpoint{2.653548in}{0.872057in}}%
\pgfpathlineto{\pgfqpoint{2.584952in}{0.921230in}}%
\pgfpathlineto{\pgfqpoint{2.524543in}{0.912354in}}%
\pgfpathlineto{\pgfqpoint{2.524037in}{0.920550in}}%
\pgfpathlineto{\pgfqpoint{2.514855in}{0.929779in}}%
\pgfpathlineto{\pgfqpoint{2.507904in}{0.932038in}}%
\pgfpathlineto{\pgfqpoint{2.442269in}{0.925218in}}%
\pgfpathlineto{\pgfqpoint{2.427183in}{0.920102in}}%
\pgfpathlineto{\pgfqpoint{2.399943in}{0.906551in}}%
\pgfpathlineto{\pgfqpoint{2.376402in}{0.902801in}}%
\pgfpathlineto{\pgfqpoint{2.335069in}{0.897708in}}%
\pgfpathclose%
\pgfusepath{fill}%
\end{pgfscope}%
\begin{pgfscope}%
\pgfpathrectangle{\pgfqpoint{0.100000in}{0.100000in}}{\pgfqpoint{2.989028in}{1.913466in}}%
\pgfusepath{clip}%
\pgfsetbuttcap%
\pgfsetmiterjoin%
\definecolor{currentfill}{rgb}{0.352018,0.707266,0.669204}%
\pgfsetfillcolor{currentfill}%
\pgfsetlinewidth{0.000000pt}%
\definecolor{currentstroke}{rgb}{0.000000,0.000000,0.000000}%
\pgfsetstrokecolor{currentstroke}%
\pgfsetstrokeopacity{0.000000}%
\pgfsetdash{}{0pt}%
\pgfpathmoveto{\pgfqpoint{2.335069in}{0.897708in}}%
\pgfpathlineto{\pgfqpoint{2.266375in}{0.890163in}}%
\pgfpathlineto{\pgfqpoint{2.203533in}{0.884512in}}%
\pgfpathlineto{\pgfqpoint{2.138945in}{0.880335in}}%
\pgfpathlineto{\pgfqpoint{2.127858in}{0.878728in}}%
\pgfpathlineto{\pgfqpoint{2.084319in}{0.875379in}}%
\pgfpathlineto{\pgfqpoint{2.014584in}{0.871310in}}%
\pgfpathlineto{\pgfqpoint{2.026990in}{0.881608in}}%
\pgfpathlineto{\pgfqpoint{2.024159in}{0.892449in}}%
\pgfpathlineto{\pgfqpoint{2.026902in}{0.905590in}}%
\pgfpathlineto{\pgfqpoint{2.031092in}{0.911777in}}%
\pgfpathlineto{\pgfqpoint{2.030934in}{0.920384in}}%
\pgfpathlineto{\pgfqpoint{2.042090in}{0.925797in}}%
\pgfpathlineto{\pgfqpoint{2.042109in}{0.938464in}}%
\pgfpathlineto{\pgfqpoint{2.043642in}{0.953066in}}%
\pgfpathlineto{\pgfqpoint{2.051421in}{0.962397in}}%
\pgfpathlineto{\pgfqpoint{2.050131in}{0.971655in}}%
\pgfpathlineto{\pgfqpoint{2.053146in}{0.971866in}}%
\pgfpathlineto{\pgfqpoint{2.056343in}{0.972329in}}%
\pgfpathlineto{\pgfqpoint{2.128480in}{0.976914in}}%
\pgfpathlineto{\pgfqpoint{2.126065in}{0.988764in}}%
\pgfpathlineto{\pgfqpoint{2.137674in}{0.987156in}}%
\pgfpathlineto{\pgfqpoint{2.207938in}{0.994107in}}%
\pgfpathlineto{\pgfqpoint{2.264337in}{0.997299in}}%
\pgfpathlineto{\pgfqpoint{2.298708in}{1.000190in}}%
\pgfpathlineto{\pgfqpoint{2.355622in}{1.005586in}}%
\pgfpathlineto{\pgfqpoint{2.356571in}{1.006660in}}%
\pgfpathlineto{\pgfqpoint{2.380888in}{1.009068in}}%
\pgfpathlineto{\pgfqpoint{2.461189in}{1.020348in}}%
\pgfpathlineto{\pgfqpoint{2.459375in}{1.006973in}}%
\pgfpathlineto{\pgfqpoint{2.451445in}{0.999123in}}%
\pgfpathlineto{\pgfqpoint{2.447232in}{0.987849in}}%
\pgfpathlineto{\pgfqpoint{2.437256in}{0.987570in}}%
\pgfpathlineto{\pgfqpoint{2.418364in}{0.972772in}}%
\pgfpathlineto{\pgfqpoint{2.409425in}{0.973951in}}%
\pgfpathlineto{\pgfqpoint{2.399610in}{0.957711in}}%
\pgfpathlineto{\pgfqpoint{2.386533in}{0.952301in}}%
\pgfpathlineto{\pgfqpoint{2.374011in}{0.940198in}}%
\pgfpathlineto{\pgfqpoint{2.359767in}{0.937949in}}%
\pgfpathlineto{\pgfqpoint{2.348372in}{0.926967in}}%
\pgfpathlineto{\pgfqpoint{2.347839in}{0.917977in}}%
\pgfpathlineto{\pgfqpoint{2.335112in}{0.912353in}}%
\pgfpathlineto{\pgfqpoint{2.335069in}{0.897708in}}%
\pgfpathclose%
\pgfusepath{fill}%
\end{pgfscope}%
\begin{pgfscope}%
\pgfpathrectangle{\pgfqpoint{0.100000in}{0.100000in}}{\pgfqpoint{2.989028in}{1.913466in}}%
\pgfusepath{clip}%
\pgfsetbuttcap%
\pgfsetmiterjoin%
\definecolor{currentfill}{rgb}{0.341561,0.344867,0.651903}%
\pgfsetfillcolor{currentfill}%
\pgfsetlinewidth{0.000000pt}%
\definecolor{currentstroke}{rgb}{0.000000,0.000000,0.000000}%
\pgfsetstrokecolor{currentstroke}%
\pgfsetstrokeopacity{0.000000}%
\pgfsetdash{}{0pt}%
\pgfpathmoveto{\pgfqpoint{1.116758in}{0.694333in}}%
\pgfpathlineto{\pgfqpoint{1.111859in}{0.702205in}}%
\pgfpathlineto{\pgfqpoint{1.113886in}{0.708987in}}%
\pgfpathlineto{\pgfqpoint{1.148324in}{0.704715in}}%
\pgfpathlineto{\pgfqpoint{1.212359in}{0.697442in}}%
\pgfpathlineto{\pgfqpoint{1.258660in}{0.692892in}}%
\pgfpathlineto{\pgfqpoint{1.312135in}{0.687496in}}%
\pgfpathlineto{\pgfqpoint{1.315044in}{0.721225in}}%
\pgfpathlineto{\pgfqpoint{1.324140in}{0.807015in}}%
\pgfpathlineto{\pgfqpoint{1.330041in}{0.867115in}}%
\pgfpathlineto{\pgfqpoint{1.335142in}{0.924525in}}%
\pgfpathlineto{\pgfqpoint{1.339896in}{0.981231in}}%
\pgfpathlineto{\pgfqpoint{1.341469in}{0.981143in}}%
\pgfpathlineto{\pgfqpoint{1.385128in}{0.977352in}}%
\pgfpathlineto{\pgfqpoint{1.441639in}{0.973080in}}%
\pgfpathlineto{\pgfqpoint{1.498906in}{0.969876in}}%
\pgfpathlineto{\pgfqpoint{1.495417in}{0.899940in}}%
\pgfpathlineto{\pgfqpoint{1.492215in}{0.843436in}}%
\pgfpathlineto{\pgfqpoint{1.495741in}{0.844347in}}%
\pgfpathlineto{\pgfqpoint{1.508395in}{0.830414in}}%
\pgfpathlineto{\pgfqpoint{1.514296in}{0.832183in}}%
\pgfpathlineto{\pgfqpoint{1.523256in}{0.829122in}}%
\pgfpathlineto{\pgfqpoint{1.531080in}{0.830667in}}%
\pgfpathlineto{\pgfqpoint{1.535448in}{0.818261in}}%
\pgfpathlineto{\pgfqpoint{1.545207in}{0.817609in}}%
\pgfpathlineto{\pgfqpoint{1.560789in}{0.811767in}}%
\pgfpathlineto{\pgfqpoint{1.565171in}{0.813414in}}%
\pgfpathlineto{\pgfqpoint{1.574573in}{0.806204in}}%
\pgfpathlineto{\pgfqpoint{1.580106in}{0.811767in}}%
\pgfpathlineto{\pgfqpoint{1.589696in}{0.809664in}}%
\pgfpathlineto{\pgfqpoint{1.594161in}{0.802273in}}%
\pgfpathlineto{\pgfqpoint{1.601493in}{0.799784in}}%
\pgfpathlineto{\pgfqpoint{1.606579in}{0.792118in}}%
\pgfpathlineto{\pgfqpoint{1.616736in}{0.800460in}}%
\pgfpathlineto{\pgfqpoint{1.628360in}{0.790111in}}%
\pgfpathlineto{\pgfqpoint{1.641649in}{0.794338in}}%
\pgfpathlineto{\pgfqpoint{1.648377in}{0.783442in}}%
\pgfpathlineto{\pgfqpoint{1.654956in}{0.796112in}}%
\pgfpathlineto{\pgfqpoint{1.666783in}{0.788928in}}%
\pgfpathlineto{\pgfqpoint{1.674172in}{0.790840in}}%
\pgfpathlineto{\pgfqpoint{1.684722in}{0.785761in}}%
\pgfpathlineto{\pgfqpoint{1.698401in}{0.784318in}}%
\pgfpathlineto{\pgfqpoint{1.711367in}{0.792135in}}%
\pgfpathlineto{\pgfqpoint{1.716057in}{0.788835in}}%
\pgfpathlineto{\pgfqpoint{1.729234in}{0.795535in}}%
\pgfpathlineto{\pgfqpoint{1.737631in}{0.790306in}}%
\pgfpathlineto{\pgfqpoint{1.744268in}{0.790617in}}%
\pgfpathlineto{\pgfqpoint{1.749524in}{0.796776in}}%
\pgfpathlineto{\pgfqpoint{1.757795in}{0.793711in}}%
\pgfpathlineto{\pgfqpoint{1.776416in}{0.779451in}}%
\pgfpathlineto{\pgfqpoint{1.790948in}{0.774957in}}%
\pgfpathlineto{\pgfqpoint{1.796775in}{0.769451in}}%
\pgfpathlineto{\pgfqpoint{1.804070in}{0.772449in}}%
\pgfpathlineto{\pgfqpoint{1.815123in}{0.770153in}}%
\pgfpathlineto{\pgfqpoint{1.815340in}{0.735108in}}%
\pgfpathlineto{\pgfqpoint{1.816262in}{0.667326in}}%
\pgfpathlineto{\pgfqpoint{1.823939in}{0.660813in}}%
\pgfpathlineto{\pgfqpoint{1.830092in}{0.648811in}}%
\pgfpathlineto{\pgfqpoint{1.827929in}{0.640781in}}%
\pgfpathlineto{\pgfqpoint{1.832761in}{0.637381in}}%
\pgfpathlineto{\pgfqpoint{1.836688in}{0.622574in}}%
\pgfpathlineto{\pgfqpoint{1.844209in}{0.614644in}}%
\pgfpathlineto{\pgfqpoint{1.845894in}{0.597808in}}%
\pgfpathlineto{\pgfqpoint{1.844586in}{0.590680in}}%
\pgfpathlineto{\pgfqpoint{1.834363in}{0.571815in}}%
\pgfpathlineto{\pgfqpoint{1.833187in}{0.558251in}}%
\pgfpathlineto{\pgfqpoint{1.836655in}{0.555420in}}%
\pgfpathlineto{\pgfqpoint{1.836956in}{0.543537in}}%
\pgfpathlineto{\pgfqpoint{1.833432in}{0.536068in}}%
\pgfpathlineto{\pgfqpoint{1.827950in}{0.533255in}}%
\pgfpathlineto{\pgfqpoint{1.822582in}{0.523496in}}%
\pgfpathlineto{\pgfqpoint{1.829424in}{0.514043in}}%
\pgfpathlineto{\pgfqpoint{1.816144in}{0.513857in}}%
\pgfpathlineto{\pgfqpoint{1.780636in}{0.497611in}}%
\pgfpathlineto{\pgfqpoint{1.787415in}{0.507333in}}%
\pgfpathlineto{\pgfqpoint{1.779206in}{0.512575in}}%
\pgfpathlineto{\pgfqpoint{1.777490in}{0.521488in}}%
\pgfpathlineto{\pgfqpoint{1.761468in}{0.505964in}}%
\pgfpathlineto{\pgfqpoint{1.768579in}{0.495354in}}%
\pgfpathlineto{\pgfqpoint{1.758435in}{0.481844in}}%
\pgfpathlineto{\pgfqpoint{1.752976in}{0.482126in}}%
\pgfpathlineto{\pgfqpoint{1.747841in}{0.467413in}}%
\pgfpathlineto{\pgfqpoint{1.731600in}{0.455840in}}%
\pgfpathlineto{\pgfqpoint{1.716427in}{0.451669in}}%
\pgfpathlineto{\pgfqpoint{1.689964in}{0.440819in}}%
\pgfpathlineto{\pgfqpoint{1.687222in}{0.446871in}}%
\pgfpathlineto{\pgfqpoint{1.670737in}{0.441300in}}%
\pgfpathlineto{\pgfqpoint{1.680876in}{0.431814in}}%
\pgfpathlineto{\pgfqpoint{1.664883in}{0.423573in}}%
\pgfpathlineto{\pgfqpoint{1.647633in}{0.411131in}}%
\pgfpathlineto{\pgfqpoint{1.643490in}{0.416915in}}%
\pgfpathlineto{\pgfqpoint{1.629373in}{0.408246in}}%
\pgfpathlineto{\pgfqpoint{1.643210in}{0.404954in}}%
\pgfpathlineto{\pgfqpoint{1.632777in}{0.391324in}}%
\pgfpathlineto{\pgfqpoint{1.627685in}{0.395380in}}%
\pgfpathlineto{\pgfqpoint{1.620845in}{0.388865in}}%
\pgfpathlineto{\pgfqpoint{1.629369in}{0.383174in}}%
\pgfpathlineto{\pgfqpoint{1.624329in}{0.374856in}}%
\pgfpathlineto{\pgfqpoint{1.618135in}{0.355164in}}%
\pgfpathlineto{\pgfqpoint{1.609200in}{0.336199in}}%
\pgfpathlineto{\pgfqpoint{1.613211in}{0.323767in}}%
\pgfpathlineto{\pgfqpoint{1.620029in}{0.294458in}}%
\pgfpathlineto{\pgfqpoint{1.620643in}{0.282596in}}%
\pgfpathlineto{\pgfqpoint{1.631178in}{0.267038in}}%
\pgfpathlineto{\pgfqpoint{1.623054in}{0.267947in}}%
\pgfpathlineto{\pgfqpoint{1.615164in}{0.260098in}}%
\pgfpathlineto{\pgfqpoint{1.605923in}{0.266267in}}%
\pgfpathlineto{\pgfqpoint{1.602605in}{0.272421in}}%
\pgfpathlineto{\pgfqpoint{1.589453in}{0.275284in}}%
\pgfpathlineto{\pgfqpoint{1.569363in}{0.275635in}}%
\pgfpathlineto{\pgfqpoint{1.554561in}{0.287289in}}%
\pgfpathlineto{\pgfqpoint{1.541120in}{0.289236in}}%
\pgfpathlineto{\pgfqpoint{1.532966in}{0.298497in}}%
\pgfpathlineto{\pgfqpoint{1.515889in}{0.302225in}}%
\pgfpathlineto{\pgfqpoint{1.506552in}{0.332004in}}%
\pgfpathlineto{\pgfqpoint{1.497004in}{0.343934in}}%
\pgfpathlineto{\pgfqpoint{1.498626in}{0.355270in}}%
\pgfpathlineto{\pgfqpoint{1.492706in}{0.363554in}}%
\pgfpathlineto{\pgfqpoint{1.496421in}{0.374877in}}%
\pgfpathlineto{\pgfqpoint{1.493356in}{0.383175in}}%
\pgfpathlineto{\pgfqpoint{1.483762in}{0.386934in}}%
\pgfpathlineto{\pgfqpoint{1.474791in}{0.396502in}}%
\pgfpathlineto{\pgfqpoint{1.471522in}{0.409312in}}%
\pgfpathlineto{\pgfqpoint{1.463033in}{0.420954in}}%
\pgfpathlineto{\pgfqpoint{1.451735in}{0.430012in}}%
\pgfpathlineto{\pgfqpoint{1.441558in}{0.455996in}}%
\pgfpathlineto{\pgfqpoint{1.433340in}{0.484413in}}%
\pgfpathlineto{\pgfqpoint{1.426597in}{0.495650in}}%
\pgfpathlineto{\pgfqpoint{1.414908in}{0.505109in}}%
\pgfpathlineto{\pgfqpoint{1.411987in}{0.511975in}}%
\pgfpathlineto{\pgfqpoint{1.401042in}{0.516238in}}%
\pgfpathlineto{\pgfqpoint{1.390239in}{0.534431in}}%
\pgfpathlineto{\pgfqpoint{1.380374in}{0.533035in}}%
\pgfpathlineto{\pgfqpoint{1.370356in}{0.537579in}}%
\pgfpathlineto{\pgfqpoint{1.356156in}{0.536700in}}%
\pgfpathlineto{\pgfqpoint{1.341731in}{0.544199in}}%
\pgfpathlineto{\pgfqpoint{1.337663in}{0.537067in}}%
\pgfpathlineto{\pgfqpoint{1.320804in}{0.536889in}}%
\pgfpathlineto{\pgfqpoint{1.312222in}{0.523370in}}%
\pgfpathlineto{\pgfqpoint{1.304757in}{0.506631in}}%
\pgfpathlineto{\pgfqpoint{1.288883in}{0.488663in}}%
\pgfpathlineto{\pgfqpoint{1.270899in}{0.496548in}}%
\pgfpathlineto{\pgfqpoint{1.268347in}{0.501760in}}%
\pgfpathlineto{\pgfqpoint{1.257418in}{0.505734in}}%
\pgfpathlineto{\pgfqpoint{1.254061in}{0.511171in}}%
\pgfpathlineto{\pgfqpoint{1.239552in}{0.516660in}}%
\pgfpathlineto{\pgfqpoint{1.231436in}{0.527879in}}%
\pgfpathlineto{\pgfqpoint{1.221957in}{0.533292in}}%
\pgfpathlineto{\pgfqpoint{1.213802in}{0.542738in}}%
\pgfpathlineto{\pgfqpoint{1.207448in}{0.558735in}}%
\pgfpathlineto{\pgfqpoint{1.208159in}{0.580595in}}%
\pgfpathlineto{\pgfqpoint{1.200706in}{0.591641in}}%
\pgfpathlineto{\pgfqpoint{1.199848in}{0.603605in}}%
\pgfpathlineto{\pgfqpoint{1.183312in}{0.621518in}}%
\pgfpathlineto{\pgfqpoint{1.173690in}{0.625342in}}%
\pgfpathlineto{\pgfqpoint{1.163479in}{0.642057in}}%
\pgfpathlineto{\pgfqpoint{1.154778in}{0.648698in}}%
\pgfpathlineto{\pgfqpoint{1.143723in}{0.664869in}}%
\pgfpathlineto{\pgfqpoint{1.132411in}{0.671867in}}%
\pgfpathlineto{\pgfqpoint{1.125003in}{0.689800in}}%
\pgfpathlineto{\pgfqpoint{1.116758in}{0.694333in}}%
\pgfpathclose%
\pgfusepath{fill}%
\end{pgfscope}%
\begin{pgfscope}%
\pgfpathrectangle{\pgfqpoint{0.100000in}{0.100000in}}{\pgfqpoint{2.989028in}{1.913466in}}%
\pgfusepath{clip}%
\pgfsetbuttcap%
\pgfsetmiterjoin%
\definecolor{currentfill}{rgb}{0.280046,0.626990,0.702422}%
\pgfsetfillcolor{currentfill}%
\pgfsetlinewidth{0.000000pt}%
\definecolor{currentstroke}{rgb}{0.000000,0.000000,0.000000}%
\pgfsetstrokecolor{currentstroke}%
\pgfsetstrokeopacity{0.000000}%
\pgfsetdash{}{0pt}%
\pgfpathmoveto{\pgfqpoint{1.030137in}{1.052362in}}%
\pgfpathlineto{\pgfqpoint{1.111744in}{1.040510in}}%
\pgfpathlineto{\pgfqpoint{1.143359in}{1.035639in}}%
\pgfpathlineto{\pgfqpoint{1.232477in}{1.025040in}}%
\pgfpathlineto{\pgfqpoint{1.292460in}{1.018783in}}%
\pgfpathlineto{\pgfqpoint{1.344368in}{1.014183in}}%
\pgfpathlineto{\pgfqpoint{1.341469in}{0.981143in}}%
\pgfpathlineto{\pgfqpoint{1.339896in}{0.981231in}}%
\pgfpathlineto{\pgfqpoint{1.335142in}{0.924525in}}%
\pgfpathlineto{\pgfqpoint{1.330041in}{0.867115in}}%
\pgfpathlineto{\pgfqpoint{1.324140in}{0.807015in}}%
\pgfpathlineto{\pgfqpoint{1.315044in}{0.721225in}}%
\pgfpathlineto{\pgfqpoint{1.312135in}{0.687496in}}%
\pgfpathlineto{\pgfqpoint{1.258660in}{0.692892in}}%
\pgfpathlineto{\pgfqpoint{1.212359in}{0.697442in}}%
\pgfpathlineto{\pgfqpoint{1.148324in}{0.704715in}}%
\pgfpathlineto{\pgfqpoint{1.113886in}{0.708987in}}%
\pgfpathlineto{\pgfqpoint{1.111859in}{0.702205in}}%
\pgfpathlineto{\pgfqpoint{1.116758in}{0.694333in}}%
\pgfpathlineto{\pgfqpoint{1.075370in}{0.699708in}}%
\pgfpathlineto{\pgfqpoint{1.024306in}{0.707037in}}%
\pgfpathlineto{\pgfqpoint{1.019659in}{0.678155in}}%
\pgfpathlineto{\pgfqpoint{0.973080in}{0.685166in}}%
\pgfpathlineto{\pgfqpoint{0.984025in}{0.756815in}}%
\pgfpathlineto{\pgfqpoint{0.997573in}{0.844052in}}%
\pgfpathlineto{\pgfqpoint{1.009619in}{0.920049in}}%
\pgfpathlineto{\pgfqpoint{1.020079in}{0.987601in}}%
\pgfpathlineto{\pgfqpoint{1.030137in}{1.052362in}}%
\pgfpathclose%
\pgfusepath{fill}%
\end{pgfscope}%
\begin{pgfscope}%
\pgfpathrectangle{\pgfqpoint{0.100000in}{0.100000in}}{\pgfqpoint{2.989028in}{1.913466in}}%
\pgfusepath{clip}%
\pgfsetbuttcap%
\pgfsetmiterjoin%
\definecolor{currentfill}{rgb}{0.506113,0.802307,0.645521}%
\pgfsetfillcolor{currentfill}%
\pgfsetlinewidth{0.000000pt}%
\definecolor{currentstroke}{rgb}{0.000000,0.000000,0.000000}%
\pgfsetstrokecolor{currentstroke}%
\pgfsetstrokeopacity{0.000000}%
\pgfsetdash{}{0pt}%
\pgfpathmoveto{\pgfqpoint{2.327687in}{0.632577in}}%
\pgfpathlineto{\pgfqpoint{2.249834in}{0.623930in}}%
\pgfpathlineto{\pgfqpoint{2.181207in}{0.618602in}}%
\pgfpathlineto{\pgfqpoint{2.180349in}{0.610192in}}%
\pgfpathlineto{\pgfqpoint{2.186652in}{0.602188in}}%
\pgfpathlineto{\pgfqpoint{2.194346in}{0.597485in}}%
\pgfpathlineto{\pgfqpoint{2.194230in}{0.585103in}}%
\pgfpathlineto{\pgfqpoint{2.192191in}{0.576831in}}%
\pgfpathlineto{\pgfqpoint{2.185400in}{0.570869in}}%
\pgfpathlineto{\pgfqpoint{2.175936in}{0.571500in}}%
\pgfpathlineto{\pgfqpoint{2.166987in}{0.578887in}}%
\pgfpathlineto{\pgfqpoint{2.165392in}{0.592032in}}%
\pgfpathlineto{\pgfqpoint{2.158727in}{0.599677in}}%
\pgfpathlineto{\pgfqpoint{2.154190in}{0.572315in}}%
\pgfpathlineto{\pgfqpoint{2.138775in}{0.574915in}}%
\pgfpathlineto{\pgfqpoint{2.133415in}{0.622761in}}%
\pgfpathlineto{\pgfqpoint{2.127610in}{0.673185in}}%
\pgfpathlineto{\pgfqpoint{2.129766in}{0.745935in}}%
\pgfpathlineto{\pgfqpoint{2.131104in}{0.795857in}}%
\pgfpathlineto{\pgfqpoint{2.133950in}{0.872025in}}%
\pgfpathlineto{\pgfqpoint{2.127858in}{0.878728in}}%
\pgfpathlineto{\pgfqpoint{2.138945in}{0.880335in}}%
\pgfpathlineto{\pgfqpoint{2.203533in}{0.884512in}}%
\pgfpathlineto{\pgfqpoint{2.266375in}{0.890163in}}%
\pgfpathlineto{\pgfqpoint{2.282908in}{0.832251in}}%
\pgfpathlineto{\pgfqpoint{2.294184in}{0.789743in}}%
\pgfpathlineto{\pgfqpoint{2.304291in}{0.754156in}}%
\pgfpathlineto{\pgfqpoint{2.310136in}{0.739826in}}%
\pgfpathlineto{\pgfqpoint{2.319340in}{0.726491in}}%
\pgfpathlineto{\pgfqpoint{2.317853in}{0.719705in}}%
\pgfpathlineto{\pgfqpoint{2.324428in}{0.716398in}}%
\pgfpathlineto{\pgfqpoint{2.316659in}{0.706111in}}%
\pgfpathlineto{\pgfqpoint{2.314034in}{0.687792in}}%
\pgfpathlineto{\pgfqpoint{2.321628in}{0.666365in}}%
\pgfpathlineto{\pgfqpoint{2.320573in}{0.644805in}}%
\pgfpathlineto{\pgfqpoint{2.327687in}{0.632577in}}%
\pgfpathclose%
\pgfusepath{fill}%
\end{pgfscope}%
\begin{pgfscope}%
\pgfpathrectangle{\pgfqpoint{0.100000in}{0.100000in}}{\pgfqpoint{2.989028in}{1.913466in}}%
\pgfusepath{clip}%
\pgfsetbuttcap%
\pgfsetmiterjoin%
\definecolor{currentfill}{rgb}{0.951942,0.980777,0.674048}%
\pgfsetfillcolor{currentfill}%
\pgfsetlinewidth{0.000000pt}%
\definecolor{currentstroke}{rgb}{0.000000,0.000000,0.000000}%
\pgfsetstrokecolor{currentstroke}%
\pgfsetstrokeopacity{0.000000}%
\pgfsetdash{}{0pt}%
\pgfpathmoveto{\pgfqpoint{2.138775in}{0.574915in}}%
\pgfpathlineto{\pgfqpoint{2.122957in}{0.570389in}}%
\pgfpathlineto{\pgfqpoint{2.111166in}{0.573295in}}%
\pgfpathlineto{\pgfqpoint{2.089268in}{0.566330in}}%
\pgfpathlineto{\pgfqpoint{2.072753in}{0.557360in}}%
\pgfpathlineto{\pgfqpoint{2.065977in}{0.573569in}}%
\pgfpathlineto{\pgfqpoint{2.058992in}{0.580363in}}%
\pgfpathlineto{\pgfqpoint{2.055622in}{0.588493in}}%
\pgfpathlineto{\pgfqpoint{2.060562in}{0.610498in}}%
\pgfpathlineto{\pgfqpoint{2.013761in}{0.607617in}}%
\pgfpathlineto{\pgfqpoint{1.953073in}{0.604986in}}%
\pgfpathlineto{\pgfqpoint{1.956684in}{0.610463in}}%
\pgfpathlineto{\pgfqpoint{1.952221in}{0.620808in}}%
\pgfpathlineto{\pgfqpoint{1.956668in}{0.622916in}}%
\pgfpathlineto{\pgfqpoint{1.955681in}{0.632856in}}%
\pgfpathlineto{\pgfqpoint{1.963432in}{0.642495in}}%
\pgfpathlineto{\pgfqpoint{1.967673in}{0.661207in}}%
\pgfpathlineto{\pgfqpoint{1.981859in}{0.673535in}}%
\pgfpathlineto{\pgfqpoint{1.976601in}{0.685508in}}%
\pgfpathlineto{\pgfqpoint{1.986660in}{0.692400in}}%
\pgfpathlineto{\pgfqpoint{1.977983in}{0.704877in}}%
\pgfpathlineto{\pgfqpoint{1.971703in}{0.733439in}}%
\pgfpathlineto{\pgfqpoint{1.974085in}{0.738625in}}%
\pgfpathlineto{\pgfqpoint{1.977845in}{0.748573in}}%
\pgfpathlineto{\pgfqpoint{1.974345in}{0.759012in}}%
\pgfpathlineto{\pgfqpoint{1.976057in}{0.763762in}}%
\pgfpathlineto{\pgfqpoint{1.969899in}{0.781706in}}%
\pgfpathlineto{\pgfqpoint{1.988077in}{0.813926in}}%
\pgfpathlineto{\pgfqpoint{1.988922in}{0.822501in}}%
\pgfpathlineto{\pgfqpoint{1.997678in}{0.828737in}}%
\pgfpathlineto{\pgfqpoint{2.007021in}{0.849323in}}%
\pgfpathlineto{\pgfqpoint{2.006577in}{0.857690in}}%
\pgfpathlineto{\pgfqpoint{2.018214in}{0.866243in}}%
\pgfpathlineto{\pgfqpoint{2.014584in}{0.871310in}}%
\pgfpathlineto{\pgfqpoint{2.084319in}{0.875379in}}%
\pgfpathlineto{\pgfqpoint{2.127858in}{0.878728in}}%
\pgfpathlineto{\pgfqpoint{2.133950in}{0.872025in}}%
\pgfpathlineto{\pgfqpoint{2.131104in}{0.795857in}}%
\pgfpathlineto{\pgfqpoint{2.129766in}{0.745935in}}%
\pgfpathlineto{\pgfqpoint{2.127610in}{0.673185in}}%
\pgfpathlineto{\pgfqpoint{2.133415in}{0.622761in}}%
\pgfpathlineto{\pgfqpoint{2.138775in}{0.574915in}}%
\pgfpathclose%
\pgfusepath{fill}%
\end{pgfscope}%
\begin{pgfscope}%
\pgfpathrectangle{\pgfqpoint{0.100000in}{0.100000in}}{\pgfqpoint{2.989028in}{1.913466in}}%
\pgfusepath{clip}%
\pgfsetbuttcap%
\pgfsetmiterjoin%
\definecolor{currentfill}{rgb}{0.453057,0.781546,0.646290}%
\pgfsetfillcolor{currentfill}%
\pgfsetlinewidth{0.000000pt}%
\definecolor{currentstroke}{rgb}{0.000000,0.000000,0.000000}%
\pgfsetstrokecolor{currentstroke}%
\pgfsetstrokeopacity{0.000000}%
\pgfsetdash{}{0pt}%
\pgfpathmoveto{\pgfqpoint{2.266375in}{0.890163in}}%
\pgfpathlineto{\pgfqpoint{2.335069in}{0.897708in}}%
\pgfpathlineto{\pgfqpoint{2.376402in}{0.902801in}}%
\pgfpathlineto{\pgfqpoint{2.399943in}{0.906551in}}%
\pgfpathlineto{\pgfqpoint{2.390192in}{0.891171in}}%
\pgfpathlineto{\pgfqpoint{2.390198in}{0.883919in}}%
\pgfpathlineto{\pgfqpoint{2.400235in}{0.880009in}}%
\pgfpathlineto{\pgfqpoint{2.407061in}{0.873695in}}%
\pgfpathlineto{\pgfqpoint{2.417377in}{0.872915in}}%
\pgfpathlineto{\pgfqpoint{2.427019in}{0.855143in}}%
\pgfpathlineto{\pgfqpoint{2.437478in}{0.842609in}}%
\pgfpathlineto{\pgfqpoint{2.456053in}{0.832075in}}%
\pgfpathlineto{\pgfqpoint{2.461332in}{0.823674in}}%
\pgfpathlineto{\pgfqpoint{2.477896in}{0.814587in}}%
\pgfpathlineto{\pgfqpoint{2.479142in}{0.807998in}}%
\pgfpathlineto{\pgfqpoint{2.492579in}{0.794896in}}%
\pgfpathlineto{\pgfqpoint{2.502872in}{0.791165in}}%
\pgfpathlineto{\pgfqpoint{2.510161in}{0.778793in}}%
\pgfpathlineto{\pgfqpoint{2.513387in}{0.764886in}}%
\pgfpathlineto{\pgfqpoint{2.524078in}{0.759511in}}%
\pgfpathlineto{\pgfqpoint{2.532685in}{0.744563in}}%
\pgfpathlineto{\pgfqpoint{2.535460in}{0.733570in}}%
\pgfpathlineto{\pgfqpoint{2.548073in}{0.728890in}}%
\pgfpathlineto{\pgfqpoint{2.537490in}{0.708626in}}%
\pgfpathlineto{\pgfqpoint{2.533507in}{0.696651in}}%
\pgfpathlineto{\pgfqpoint{2.536105in}{0.691037in}}%
\pgfpathlineto{\pgfqpoint{2.531568in}{0.681708in}}%
\pgfpathlineto{\pgfqpoint{2.529639in}{0.668810in}}%
\pgfpathlineto{\pgfqpoint{2.521570in}{0.666408in}}%
\pgfpathlineto{\pgfqpoint{2.525831in}{0.654599in}}%
\pgfpathlineto{\pgfqpoint{2.525559in}{0.639617in}}%
\pgfpathlineto{\pgfqpoint{2.511217in}{0.640178in}}%
\pgfpathlineto{\pgfqpoint{2.501225in}{0.642912in}}%
\pgfpathlineto{\pgfqpoint{2.496902in}{0.637830in}}%
\pgfpathlineto{\pgfqpoint{2.500128in}{0.625826in}}%
\pgfpathlineto{\pgfqpoint{2.499427in}{0.611870in}}%
\pgfpathlineto{\pgfqpoint{2.493123in}{0.610802in}}%
\pgfpathlineto{\pgfqpoint{2.488024in}{0.623836in}}%
\pgfpathlineto{\pgfqpoint{2.436103in}{0.620376in}}%
\pgfpathlineto{\pgfqpoint{2.370635in}{0.616804in}}%
\pgfpathlineto{\pgfqpoint{2.337617in}{0.614480in}}%
\pgfpathlineto{\pgfqpoint{2.327687in}{0.632577in}}%
\pgfpathlineto{\pgfqpoint{2.320573in}{0.644805in}}%
\pgfpathlineto{\pgfqpoint{2.321628in}{0.666365in}}%
\pgfpathlineto{\pgfqpoint{2.314034in}{0.687792in}}%
\pgfpathlineto{\pgfqpoint{2.316659in}{0.706111in}}%
\pgfpathlineto{\pgfqpoint{2.324428in}{0.716398in}}%
\pgfpathlineto{\pgfqpoint{2.317853in}{0.719705in}}%
\pgfpathlineto{\pgfqpoint{2.319340in}{0.726491in}}%
\pgfpathlineto{\pgfqpoint{2.310136in}{0.739826in}}%
\pgfpathlineto{\pgfqpoint{2.304291in}{0.754156in}}%
\pgfpathlineto{\pgfqpoint{2.294184in}{0.789743in}}%
\pgfpathlineto{\pgfqpoint{2.282908in}{0.832251in}}%
\pgfpathlineto{\pgfqpoint{2.266375in}{0.890163in}}%
\pgfpathclose%
\pgfusepath{fill}%
\end{pgfscope}%
\begin{pgfscope}%
\pgfpathrectangle{\pgfqpoint{0.100000in}{0.100000in}}{\pgfqpoint{2.989028in}{1.913466in}}%
\pgfusepath{clip}%
\pgfsetbuttcap%
\pgfsetmiterjoin%
\definecolor{currentfill}{rgb}{0.280046,0.626990,0.702422}%
\pgfsetfillcolor{currentfill}%
\pgfsetlinewidth{0.000000pt}%
\definecolor{currentstroke}{rgb}{0.000000,0.000000,0.000000}%
\pgfsetstrokecolor{currentstroke}%
\pgfsetstrokeopacity{0.000000}%
\pgfsetdash{}{0pt}%
\pgfpathmoveto{\pgfqpoint{2.399943in}{0.906551in}}%
\pgfpathlineto{\pgfqpoint{2.427183in}{0.920102in}}%
\pgfpathlineto{\pgfqpoint{2.442269in}{0.925218in}}%
\pgfpathlineto{\pgfqpoint{2.507904in}{0.932038in}}%
\pgfpathlineto{\pgfqpoint{2.514855in}{0.929779in}}%
\pgfpathlineto{\pgfqpoint{2.524037in}{0.920550in}}%
\pgfpathlineto{\pgfqpoint{2.524543in}{0.912354in}}%
\pgfpathlineto{\pgfqpoint{2.584952in}{0.921230in}}%
\pgfpathlineto{\pgfqpoint{2.653548in}{0.872057in}}%
\pgfpathlineto{\pgfqpoint{2.640684in}{0.858649in}}%
\pgfpathlineto{\pgfqpoint{2.623033in}{0.827511in}}%
\pgfpathlineto{\pgfqpoint{2.628057in}{0.820819in}}%
\pgfpathlineto{\pgfqpoint{2.618641in}{0.807830in}}%
\pgfpathlineto{\pgfqpoint{2.609275in}{0.806359in}}%
\pgfpathlineto{\pgfqpoint{2.609240in}{0.798554in}}%
\pgfpathlineto{\pgfqpoint{2.602464in}{0.790386in}}%
\pgfpathlineto{\pgfqpoint{2.594032in}{0.788730in}}%
\pgfpathlineto{\pgfqpoint{2.595872in}{0.781473in}}%
\pgfpathlineto{\pgfqpoint{2.591167in}{0.775902in}}%
\pgfpathlineto{\pgfqpoint{2.579911in}{0.771089in}}%
\pgfpathlineto{\pgfqpoint{2.566192in}{0.760141in}}%
\pgfpathlineto{\pgfqpoint{2.568782in}{0.753058in}}%
\pgfpathlineto{\pgfqpoint{2.560148in}{0.748611in}}%
\pgfpathlineto{\pgfqpoint{2.552866in}{0.753287in}}%
\pgfpathlineto{\pgfqpoint{2.547540in}{0.732958in}}%
\pgfpathlineto{\pgfqpoint{2.535460in}{0.733570in}}%
\pgfpathlineto{\pgfqpoint{2.532685in}{0.744563in}}%
\pgfpathlineto{\pgfqpoint{2.524078in}{0.759511in}}%
\pgfpathlineto{\pgfqpoint{2.513387in}{0.764886in}}%
\pgfpathlineto{\pgfqpoint{2.510161in}{0.778793in}}%
\pgfpathlineto{\pgfqpoint{2.502872in}{0.791165in}}%
\pgfpathlineto{\pgfqpoint{2.492579in}{0.794896in}}%
\pgfpathlineto{\pgfqpoint{2.479142in}{0.807998in}}%
\pgfpathlineto{\pgfqpoint{2.477896in}{0.814587in}}%
\pgfpathlineto{\pgfqpoint{2.461332in}{0.823674in}}%
\pgfpathlineto{\pgfqpoint{2.456053in}{0.832075in}}%
\pgfpathlineto{\pgfqpoint{2.437478in}{0.842609in}}%
\pgfpathlineto{\pgfqpoint{2.427019in}{0.855143in}}%
\pgfpathlineto{\pgfqpoint{2.417377in}{0.872915in}}%
\pgfpathlineto{\pgfqpoint{2.407061in}{0.873695in}}%
\pgfpathlineto{\pgfqpoint{2.400235in}{0.880009in}}%
\pgfpathlineto{\pgfqpoint{2.390198in}{0.883919in}}%
\pgfpathlineto{\pgfqpoint{2.390192in}{0.891171in}}%
\pgfpathlineto{\pgfqpoint{2.399943in}{0.906551in}}%
\pgfpathclose%
\pgfusepath{fill}%
\end{pgfscope}%
\begin{pgfscope}%
\pgfpathrectangle{\pgfqpoint{0.100000in}{0.100000in}}{\pgfqpoint{2.989028in}{1.913466in}}%
\pgfusepath{clip}%
\pgfsetbuttcap%
\pgfsetmiterjoin%
\definecolor{currentfill}{rgb}{0.959631,0.983852,0.686044}%
\pgfsetfillcolor{currentfill}%
\pgfsetlinewidth{0.000000pt}%
\definecolor{currentstroke}{rgb}{0.000000,0.000000,0.000000}%
\pgfsetstrokecolor{currentstroke}%
\pgfsetstrokeopacity{0.000000}%
\pgfsetdash{}{0pt}%
\pgfpathmoveto{\pgfqpoint{1.782462in}{0.962463in}}%
\pgfpathlineto{\pgfqpoint{1.822491in}{0.962824in}}%
\pgfpathlineto{\pgfqpoint{1.875448in}{0.963789in}}%
\pgfpathlineto{\pgfqpoint{1.965941in}{0.966662in}}%
\pgfpathlineto{\pgfqpoint{2.017700in}{0.969399in}}%
\pgfpathlineto{\pgfqpoint{2.023251in}{0.962527in}}%
\pgfpathlineto{\pgfqpoint{2.022904in}{0.955269in}}%
\pgfpathlineto{\pgfqpoint{2.010365in}{0.942748in}}%
\pgfpathlineto{\pgfqpoint{2.007334in}{0.935892in}}%
\pgfpathlineto{\pgfqpoint{2.042109in}{0.938464in}}%
\pgfpathlineto{\pgfqpoint{2.042090in}{0.925797in}}%
\pgfpathlineto{\pgfqpoint{2.030934in}{0.920384in}}%
\pgfpathlineto{\pgfqpoint{2.031092in}{0.911777in}}%
\pgfpathlineto{\pgfqpoint{2.026902in}{0.905590in}}%
\pgfpathlineto{\pgfqpoint{2.024159in}{0.892449in}}%
\pgfpathlineto{\pgfqpoint{2.026990in}{0.881608in}}%
\pgfpathlineto{\pgfqpoint{2.014584in}{0.871310in}}%
\pgfpathlineto{\pgfqpoint{2.018214in}{0.866243in}}%
\pgfpathlineto{\pgfqpoint{2.006577in}{0.857690in}}%
\pgfpathlineto{\pgfqpoint{2.007021in}{0.849323in}}%
\pgfpathlineto{\pgfqpoint{1.997678in}{0.828737in}}%
\pgfpathlineto{\pgfqpoint{1.988922in}{0.822501in}}%
\pgfpathlineto{\pgfqpoint{1.988077in}{0.813926in}}%
\pgfpathlineto{\pgfqpoint{1.969899in}{0.781706in}}%
\pgfpathlineto{\pgfqpoint{1.976057in}{0.763762in}}%
\pgfpathlineto{\pgfqpoint{1.974345in}{0.759012in}}%
\pgfpathlineto{\pgfqpoint{1.977845in}{0.748573in}}%
\pgfpathlineto{\pgfqpoint{1.974085in}{0.738625in}}%
\pgfpathlineto{\pgfqpoint{1.924378in}{0.736573in}}%
\pgfpathlineto{\pgfqpoint{1.859848in}{0.735510in}}%
\pgfpathlineto{\pgfqpoint{1.815340in}{0.735108in}}%
\pgfpathlineto{\pgfqpoint{1.815123in}{0.770153in}}%
\pgfpathlineto{\pgfqpoint{1.804070in}{0.772449in}}%
\pgfpathlineto{\pgfqpoint{1.796775in}{0.769451in}}%
\pgfpathlineto{\pgfqpoint{1.790948in}{0.774957in}}%
\pgfpathlineto{\pgfqpoint{1.791549in}{0.812079in}}%
\pgfpathlineto{\pgfqpoint{1.792853in}{0.891120in}}%
\pgfpathlineto{\pgfqpoint{1.786531in}{0.937404in}}%
\pgfpathlineto{\pgfqpoint{1.782462in}{0.962463in}}%
\pgfpathclose%
\pgfusepath{fill}%
\end{pgfscope}%
\begin{pgfscope}%
\pgfpathrectangle{\pgfqpoint{0.100000in}{0.100000in}}{\pgfqpoint{2.989028in}{1.913466in}}%
\pgfusepath{clip}%
\pgfsetbuttcap%
\pgfsetmiterjoin%
\definecolor{currentfill}{rgb}{0.212995,0.511419,0.730796}%
\pgfsetfillcolor{currentfill}%
\pgfsetlinewidth{0.000000pt}%
\definecolor{currentstroke}{rgb}{0.000000,0.000000,0.000000}%
\pgfsetstrokecolor{currentstroke}%
\pgfsetstrokeopacity{0.000000}%
\pgfsetdash{}{0pt}%
\pgfpathmoveto{\pgfqpoint{1.815340in}{0.735108in}}%
\pgfpathlineto{\pgfqpoint{1.859848in}{0.735510in}}%
\pgfpathlineto{\pgfqpoint{1.924378in}{0.736573in}}%
\pgfpathlineto{\pgfqpoint{1.974085in}{0.738625in}}%
\pgfpathlineto{\pgfqpoint{1.971703in}{0.733439in}}%
\pgfpathlineto{\pgfqpoint{1.977983in}{0.704877in}}%
\pgfpathlineto{\pgfqpoint{1.986660in}{0.692400in}}%
\pgfpathlineto{\pgfqpoint{1.976601in}{0.685508in}}%
\pgfpathlineto{\pgfqpoint{1.981859in}{0.673535in}}%
\pgfpathlineto{\pgfqpoint{1.967673in}{0.661207in}}%
\pgfpathlineto{\pgfqpoint{1.963432in}{0.642495in}}%
\pgfpathlineto{\pgfqpoint{1.955681in}{0.632856in}}%
\pgfpathlineto{\pgfqpoint{1.956668in}{0.622916in}}%
\pgfpathlineto{\pgfqpoint{1.952221in}{0.620808in}}%
\pgfpathlineto{\pgfqpoint{1.956684in}{0.610463in}}%
\pgfpathlineto{\pgfqpoint{1.953073in}{0.604986in}}%
\pgfpathlineto{\pgfqpoint{2.013761in}{0.607617in}}%
\pgfpathlineto{\pgfqpoint{2.060562in}{0.610498in}}%
\pgfpathlineto{\pgfqpoint{2.055622in}{0.588493in}}%
\pgfpathlineto{\pgfqpoint{2.058992in}{0.580363in}}%
\pgfpathlineto{\pgfqpoint{2.065977in}{0.573569in}}%
\pgfpathlineto{\pgfqpoint{2.072753in}{0.557360in}}%
\pgfpathlineto{\pgfqpoint{2.051350in}{0.561094in}}%
\pgfpathlineto{\pgfqpoint{2.043455in}{0.567234in}}%
\pgfpathlineto{\pgfqpoint{2.034048in}{0.567526in}}%
\pgfpathlineto{\pgfqpoint{2.024158in}{0.554082in}}%
\pgfpathlineto{\pgfqpoint{2.026129in}{0.547960in}}%
\pgfpathlineto{\pgfqpoint{2.057724in}{0.544257in}}%
\pgfpathlineto{\pgfqpoint{2.065995in}{0.537220in}}%
\pgfpathlineto{\pgfqpoint{2.079454in}{0.533600in}}%
\pgfpathlineto{\pgfqpoint{2.073576in}{0.525268in}}%
\pgfpathlineto{\pgfqpoint{2.063672in}{0.518003in}}%
\pgfpathlineto{\pgfqpoint{2.089813in}{0.501664in}}%
\pgfpathlineto{\pgfqpoint{2.101990in}{0.499111in}}%
\pgfpathlineto{\pgfqpoint{2.107054in}{0.485755in}}%
\pgfpathlineto{\pgfqpoint{2.095589in}{0.483298in}}%
\pgfpathlineto{\pgfqpoint{2.074041in}{0.500096in}}%
\pgfpathlineto{\pgfqpoint{2.065616in}{0.502430in}}%
\pgfpathlineto{\pgfqpoint{2.061539in}{0.509066in}}%
\pgfpathlineto{\pgfqpoint{2.049043in}{0.506343in}}%
\pgfpathlineto{\pgfqpoint{2.045130in}{0.497792in}}%
\pgfpathlineto{\pgfqpoint{2.047608in}{0.488264in}}%
\pgfpathlineto{\pgfqpoint{2.039195in}{0.482642in}}%
\pgfpathlineto{\pgfqpoint{2.031517in}{0.496493in}}%
\pgfpathlineto{\pgfqpoint{2.018069in}{0.492363in}}%
\pgfpathlineto{\pgfqpoint{2.013030in}{0.484098in}}%
\pgfpathlineto{\pgfqpoint{2.003479in}{0.486449in}}%
\pgfpathlineto{\pgfqpoint{1.997373in}{0.496886in}}%
\pgfpathlineto{\pgfqpoint{1.981173in}{0.500333in}}%
\pgfpathlineto{\pgfqpoint{1.978115in}{0.505764in}}%
\pgfpathlineto{\pgfqpoint{1.961236in}{0.515225in}}%
\pgfpathlineto{\pgfqpoint{1.957047in}{0.523496in}}%
\pgfpathlineto{\pgfqpoint{1.942910in}{0.520116in}}%
\pgfpathlineto{\pgfqpoint{1.944705in}{0.527689in}}%
\pgfpathlineto{\pgfqpoint{1.927134in}{0.519920in}}%
\pgfpathlineto{\pgfqpoint{1.931854in}{0.511860in}}%
\pgfpathlineto{\pgfqpoint{1.918284in}{0.507095in}}%
\pgfpathlineto{\pgfqpoint{1.900311in}{0.509714in}}%
\pgfpathlineto{\pgfqpoint{1.863934in}{0.522174in}}%
\pgfpathlineto{\pgfqpoint{1.835864in}{0.519701in}}%
\pgfpathlineto{\pgfqpoint{1.833432in}{0.536068in}}%
\pgfpathlineto{\pgfqpoint{1.836956in}{0.543537in}}%
\pgfpathlineto{\pgfqpoint{1.836655in}{0.555420in}}%
\pgfpathlineto{\pgfqpoint{1.833187in}{0.558251in}}%
\pgfpathlineto{\pgfqpoint{1.834363in}{0.571815in}}%
\pgfpathlineto{\pgfqpoint{1.844586in}{0.590680in}}%
\pgfpathlineto{\pgfqpoint{1.845894in}{0.597808in}}%
\pgfpathlineto{\pgfqpoint{1.844209in}{0.614644in}}%
\pgfpathlineto{\pgfqpoint{1.836688in}{0.622574in}}%
\pgfpathlineto{\pgfqpoint{1.832761in}{0.637381in}}%
\pgfpathlineto{\pgfqpoint{1.827929in}{0.640781in}}%
\pgfpathlineto{\pgfqpoint{1.830092in}{0.648811in}}%
\pgfpathlineto{\pgfqpoint{1.823939in}{0.660813in}}%
\pgfpathlineto{\pgfqpoint{1.816262in}{0.667326in}}%
\pgfpathlineto{\pgfqpoint{1.815340in}{0.735108in}}%
\pgfpathclose%
\pgfusepath{fill}%
\end{pgfscope}%
\begin{pgfscope}%
\pgfpathrectangle{\pgfqpoint{0.100000in}{0.100000in}}{\pgfqpoint{2.989028in}{1.913466in}}%
\pgfusepath{clip}%
\pgfsetbuttcap%
\pgfsetmiterjoin%
\definecolor{currentfill}{rgb}{0.212995,0.511419,0.730796}%
\pgfsetfillcolor{currentfill}%
\pgfsetlinewidth{0.000000pt}%
\definecolor{currentstroke}{rgb}{0.000000,0.000000,0.000000}%
\pgfsetstrokecolor{currentstroke}%
\pgfsetstrokeopacity{0.000000}%
\pgfsetdash{}{0pt}%
\pgfpathmoveto{\pgfqpoint{1.934378in}{0.511243in}}%
\pgfpathlineto{\pgfqpoint{1.940817in}{0.515073in}}%
\pgfpathlineto{\pgfqpoint{1.948633in}{0.510545in}}%
\pgfpathlineto{\pgfqpoint{1.944274in}{0.504314in}}%
\pgfpathlineto{\pgfqpoint{1.934378in}{0.511243in}}%
\pgfpathclose%
\pgfusepath{fill}%
\end{pgfscope}%
\begin{pgfscope}%
\pgfpathrectangle{\pgfqpoint{0.100000in}{0.100000in}}{\pgfqpoint{2.989028in}{1.913466in}}%
\pgfusepath{clip}%
\pgfsetbuttcap%
\pgfsetmiterjoin%
\definecolor{currentfill}{rgb}{0.240062,0.582391,0.720877}%
\pgfsetfillcolor{currentfill}%
\pgfsetlinewidth{0.000000pt}%
\definecolor{currentstroke}{rgb}{0.000000,0.000000,0.000000}%
\pgfsetstrokecolor{currentstroke}%
\pgfsetstrokeopacity{0.000000}%
\pgfsetdash{}{0pt}%
\pgfpathmoveto{\pgfqpoint{2.194230in}{0.585103in}}%
\pgfpathlineto{\pgfqpoint{2.194346in}{0.597485in}}%
\pgfpathlineto{\pgfqpoint{2.186652in}{0.602188in}}%
\pgfpathlineto{\pgfqpoint{2.180349in}{0.610192in}}%
\pgfpathlineto{\pgfqpoint{2.181207in}{0.618602in}}%
\pgfpathlineto{\pgfqpoint{2.249834in}{0.623930in}}%
\pgfpathlineto{\pgfqpoint{2.327687in}{0.632577in}}%
\pgfpathlineto{\pgfqpoint{2.337617in}{0.614480in}}%
\pgfpathlineto{\pgfqpoint{2.370635in}{0.616804in}}%
\pgfpathlineto{\pgfqpoint{2.436103in}{0.620376in}}%
\pgfpathlineto{\pgfqpoint{2.488024in}{0.623836in}}%
\pgfpathlineto{\pgfqpoint{2.493123in}{0.610802in}}%
\pgfpathlineto{\pgfqpoint{2.499427in}{0.611870in}}%
\pgfpathlineto{\pgfqpoint{2.500128in}{0.625826in}}%
\pgfpathlineto{\pgfqpoint{2.496902in}{0.637830in}}%
\pgfpathlineto{\pgfqpoint{2.501225in}{0.642912in}}%
\pgfpathlineto{\pgfqpoint{2.511217in}{0.640178in}}%
\pgfpathlineto{\pgfqpoint{2.525559in}{0.639617in}}%
\pgfpathlineto{\pgfqpoint{2.527764in}{0.628892in}}%
\pgfpathlineto{\pgfqpoint{2.532176in}{0.622764in}}%
\pgfpathlineto{\pgfqpoint{2.535621in}{0.609356in}}%
\pgfpathlineto{\pgfqpoint{2.546296in}{0.588607in}}%
\pgfpathlineto{\pgfqpoint{2.546349in}{0.582995in}}%
\pgfpathlineto{\pgfqpoint{2.562126in}{0.558631in}}%
\pgfpathlineto{\pgfqpoint{2.563652in}{0.553598in}}%
\pgfpathlineto{\pgfqpoint{2.587320in}{0.519188in}}%
\pgfpathlineto{\pgfqpoint{2.583570in}{0.518632in}}%
\pgfpathlineto{\pgfqpoint{2.593548in}{0.494158in}}%
\pgfpathlineto{\pgfqpoint{2.620959in}{0.451609in}}%
\pgfpathlineto{\pgfqpoint{2.636873in}{0.420325in}}%
\pgfpathlineto{\pgfqpoint{2.642192in}{0.414848in}}%
\pgfpathlineto{\pgfqpoint{2.651306in}{0.395211in}}%
\pgfpathlineto{\pgfqpoint{2.654468in}{0.363769in}}%
\pgfpathlineto{\pgfqpoint{2.655727in}{0.340258in}}%
\pgfpathlineto{\pgfqpoint{2.654207in}{0.325205in}}%
\pgfpathlineto{\pgfqpoint{2.649320in}{0.314444in}}%
\pgfpathlineto{\pgfqpoint{2.651590in}{0.300332in}}%
\pgfpathlineto{\pgfqpoint{2.646329in}{0.289163in}}%
\pgfpathlineto{\pgfqpoint{2.630707in}{0.280003in}}%
\pgfpathlineto{\pgfqpoint{2.620547in}{0.280672in}}%
\pgfpathlineto{\pgfqpoint{2.613949in}{0.275884in}}%
\pgfpathlineto{\pgfqpoint{2.612059in}{0.288668in}}%
\pgfpathlineto{\pgfqpoint{2.601137in}{0.292042in}}%
\pgfpathlineto{\pgfqpoint{2.591346in}{0.309880in}}%
\pgfpathlineto{\pgfqpoint{2.590241in}{0.318003in}}%
\pgfpathlineto{\pgfqpoint{2.572719in}{0.322936in}}%
\pgfpathlineto{\pgfqpoint{2.561428in}{0.321868in}}%
\pgfpathlineto{\pgfqpoint{2.555040in}{0.333654in}}%
\pgfpathlineto{\pgfqpoint{2.547725in}{0.354839in}}%
\pgfpathlineto{\pgfqpoint{2.537575in}{0.359133in}}%
\pgfpathlineto{\pgfqpoint{2.532072in}{0.371273in}}%
\pgfpathlineto{\pgfqpoint{2.518614in}{0.378428in}}%
\pgfpathlineto{\pgfqpoint{2.510835in}{0.387377in}}%
\pgfpathlineto{\pgfqpoint{2.497416in}{0.411708in}}%
\pgfpathlineto{\pgfqpoint{2.495838in}{0.428833in}}%
\pgfpathlineto{\pgfqpoint{2.503150in}{0.439789in}}%
\pgfpathlineto{\pgfqpoint{2.502419in}{0.447403in}}%
\pgfpathlineto{\pgfqpoint{2.493969in}{0.448217in}}%
\pgfpathlineto{\pgfqpoint{2.486909in}{0.453513in}}%
\pgfpathlineto{\pgfqpoint{2.483042in}{0.447046in}}%
\pgfpathlineto{\pgfqpoint{2.489818in}{0.441777in}}%
\pgfpathlineto{\pgfqpoint{2.484887in}{0.432302in}}%
\pgfpathlineto{\pgfqpoint{2.476908in}{0.440166in}}%
\pgfpathlineto{\pgfqpoint{2.477834in}{0.461974in}}%
\pgfpathlineto{\pgfqpoint{2.481682in}{0.479667in}}%
\pgfpathlineto{\pgfqpoint{2.479704in}{0.510026in}}%
\pgfpathlineto{\pgfqpoint{2.471751in}{0.517257in}}%
\pgfpathlineto{\pgfqpoint{2.467752in}{0.526536in}}%
\pgfpathlineto{\pgfqpoint{2.454052in}{0.526329in}}%
\pgfpathlineto{\pgfqpoint{2.440510in}{0.541597in}}%
\pgfpathlineto{\pgfqpoint{2.431430in}{0.546173in}}%
\pgfpathlineto{\pgfqpoint{2.428737in}{0.555842in}}%
\pgfpathlineto{\pgfqpoint{2.419856in}{0.559238in}}%
\pgfpathlineto{\pgfqpoint{2.412496in}{0.569913in}}%
\pgfpathlineto{\pgfqpoint{2.393052in}{0.578628in}}%
\pgfpathlineto{\pgfqpoint{2.377937in}{0.578850in}}%
\pgfpathlineto{\pgfqpoint{2.371356in}{0.575503in}}%
\pgfpathlineto{\pgfqpoint{2.372991in}{0.565056in}}%
\pgfpathlineto{\pgfqpoint{2.366130in}{0.565551in}}%
\pgfpathlineto{\pgfqpoint{2.345022in}{0.550943in}}%
\pgfpathlineto{\pgfqpoint{2.319668in}{0.545153in}}%
\pgfpathlineto{\pgfqpoint{2.319252in}{0.552317in}}%
\pgfpathlineto{\pgfqpoint{2.313633in}{0.559318in}}%
\pgfpathlineto{\pgfqpoint{2.298558in}{0.568984in}}%
\pgfpathlineto{\pgfqpoint{2.276918in}{0.578894in}}%
\pgfpathlineto{\pgfqpoint{2.253414in}{0.584141in}}%
\pgfpathlineto{\pgfqpoint{2.249002in}{0.591275in}}%
\pgfpathlineto{\pgfqpoint{2.240528in}{0.585325in}}%
\pgfpathlineto{\pgfqpoint{2.230323in}{0.584012in}}%
\pgfpathlineto{\pgfqpoint{2.194801in}{0.574653in}}%
\pgfpathlineto{\pgfqpoint{2.194230in}{0.585103in}}%
\pgfpathclose%
\pgfusepath{fill}%
\end{pgfscope}%
\begin{pgfscope}%
\pgfpathrectangle{\pgfqpoint{0.100000in}{0.100000in}}{\pgfqpoint{2.989028in}{1.913466in}}%
\pgfusepath{clip}%
\pgfsetbuttcap%
\pgfsetmiterjoin%
\definecolor{currentfill}{rgb}{0.240062,0.582391,0.720877}%
\pgfsetfillcolor{currentfill}%
\pgfsetlinewidth{0.000000pt}%
\definecolor{currentstroke}{rgb}{0.000000,0.000000,0.000000}%
\pgfsetstrokecolor{currentstroke}%
\pgfsetstrokeopacity{0.000000}%
\pgfsetdash{}{0pt}%
\pgfpathmoveto{\pgfqpoint{2.590769in}{0.519678in}}%
\pgfpathlineto{\pgfqpoint{2.599959in}{0.508987in}}%
\pgfpathlineto{\pgfqpoint{2.589344in}{0.508314in}}%
\pgfpathlineto{\pgfqpoint{2.590769in}{0.519678in}}%
\pgfpathclose%
\pgfusepath{fill}%
\end{pgfscope}%
\begin{pgfscope}%
\pgfpathrectangle{\pgfqpoint{0.100000in}{0.100000in}}{\pgfqpoint{2.989028in}{1.913466in}}%
\pgfusepath{clip}%
\pgfsetbuttcap%
\pgfsetmiterjoin%
\definecolor{currentfill}{rgb}{0.711419,0.883276,0.634833}%
\pgfsetfillcolor{currentfill}%
\pgfsetlinewidth{0.000000pt}%
\definecolor{currentstroke}{rgb}{0.000000,0.000000,0.000000}%
\pgfsetstrokecolor{currentstroke}%
\pgfsetstrokeopacity{0.000000}%
\pgfsetdash{}{0pt}%
\pgfpathmoveto{\pgfqpoint{2.049740in}{1.740361in}}%
\pgfpathlineto{\pgfqpoint{2.044753in}{1.730656in}}%
\pgfpathlineto{\pgfqpoint{2.032857in}{1.724992in}}%
\pgfpathlineto{\pgfqpoint{2.027705in}{1.717368in}}%
\pgfpathlineto{\pgfqpoint{2.021669in}{1.722859in}}%
\pgfpathlineto{\pgfqpoint{2.049740in}{1.740361in}}%
\pgfpathclose%
\pgfusepath{fill}%
\end{pgfscope}%
\begin{pgfscope}%
\pgfpathrectangle{\pgfqpoint{0.100000in}{0.100000in}}{\pgfqpoint{2.989028in}{1.913466in}}%
\pgfusepath{clip}%
\pgfsetbuttcap%
\pgfsetmiterjoin%
\definecolor{currentfill}{rgb}{0.711419,0.883276,0.634833}%
\pgfsetfillcolor{currentfill}%
\pgfsetlinewidth{0.000000pt}%
\definecolor{currentstroke}{rgb}{0.000000,0.000000,0.000000}%
\pgfsetstrokecolor{currentstroke}%
\pgfsetstrokeopacity{0.000000}%
\pgfsetdash{}{0pt}%
\pgfpathmoveto{\pgfqpoint{2.053784in}{1.681953in}}%
\pgfpathlineto{\pgfqpoint{2.065956in}{1.693308in}}%
\pgfpathlineto{\pgfqpoint{2.084730in}{1.696281in}}%
\pgfpathlineto{\pgfqpoint{2.079559in}{1.688384in}}%
\pgfpathlineto{\pgfqpoint{2.066688in}{1.676964in}}%
\pgfpathlineto{\pgfqpoint{2.059164in}{1.662306in}}%
\pgfpathlineto{\pgfqpoint{2.055977in}{1.670256in}}%
\pgfpathlineto{\pgfqpoint{2.050318in}{1.671398in}}%
\pgfpathlineto{\pgfqpoint{2.049750in}{1.678576in}}%
\pgfpathlineto{\pgfqpoint{2.053784in}{1.681953in}}%
\pgfpathclose%
\pgfusepath{fill}%
\end{pgfscope}%
\begin{pgfscope}%
\pgfpathrectangle{\pgfqpoint{0.100000in}{0.100000in}}{\pgfqpoint{2.989028in}{1.913466in}}%
\pgfusepath{clip}%
\pgfsetbuttcap%
\pgfsetmiterjoin%
\definecolor{currentfill}{rgb}{0.711419,0.883276,0.634833}%
\pgfsetfillcolor{currentfill}%
\pgfsetlinewidth{0.000000pt}%
\definecolor{currentstroke}{rgb}{0.000000,0.000000,0.000000}%
\pgfsetstrokecolor{currentstroke}%
\pgfsetstrokeopacity{0.000000}%
\pgfsetdash{}{0pt}%
\pgfpathmoveto{\pgfqpoint{2.102264in}{1.543360in}}%
\pgfpathlineto{\pgfqpoint{2.099022in}{1.546957in}}%
\pgfpathlineto{\pgfqpoint{2.102430in}{1.557084in}}%
\pgfpathlineto{\pgfqpoint{2.092314in}{1.557730in}}%
\pgfpathlineto{\pgfqpoint{2.094996in}{1.566462in}}%
\pgfpathlineto{\pgfqpoint{2.093328in}{1.580368in}}%
\pgfpathlineto{\pgfqpoint{2.084251in}{1.585189in}}%
\pgfpathlineto{\pgfqpoint{2.074752in}{1.594947in}}%
\pgfpathlineto{\pgfqpoint{2.060116in}{1.597803in}}%
\pgfpathlineto{\pgfqpoint{2.045871in}{1.597722in}}%
\pgfpathlineto{\pgfqpoint{2.031860in}{1.604646in}}%
\pgfpathlineto{\pgfqpoint{1.984989in}{1.614740in}}%
\pgfpathlineto{\pgfqpoint{1.979867in}{1.625436in}}%
\pgfpathlineto{\pgfqpoint{1.970731in}{1.629085in}}%
\pgfpathlineto{\pgfqpoint{1.987989in}{1.637274in}}%
\pgfpathlineto{\pgfqpoint{1.997738in}{1.647490in}}%
\pgfpathlineto{\pgfqpoint{2.015897in}{1.650260in}}%
\pgfpathlineto{\pgfqpoint{2.027071in}{1.660666in}}%
\pgfpathlineto{\pgfqpoint{2.032934in}{1.661083in}}%
\pgfpathlineto{\pgfqpoint{2.037408in}{1.668504in}}%
\pgfpathlineto{\pgfqpoint{2.048315in}{1.677295in}}%
\pgfpathlineto{\pgfqpoint{2.049265in}{1.671081in}}%
\pgfpathlineto{\pgfqpoint{2.054179in}{1.669798in}}%
\pgfpathlineto{\pgfqpoint{2.057877in}{1.662388in}}%
\pgfpathlineto{\pgfqpoint{2.058543in}{1.649753in}}%
\pgfpathlineto{\pgfqpoint{2.069656in}{1.657304in}}%
\pgfpathlineto{\pgfqpoint{2.082606in}{1.658880in}}%
\pgfpathlineto{\pgfqpoint{2.093660in}{1.654925in}}%
\pgfpathlineto{\pgfqpoint{2.108651in}{1.634360in}}%
\pgfpathlineto{\pgfqpoint{2.125023in}{1.637647in}}%
\pgfpathlineto{\pgfqpoint{2.131676in}{1.632143in}}%
\pgfpathlineto{\pgfqpoint{2.142376in}{1.631651in}}%
\pgfpathlineto{\pgfqpoint{2.149485in}{1.641522in}}%
\pgfpathlineto{\pgfqpoint{2.162974in}{1.650248in}}%
\pgfpathlineto{\pgfqpoint{2.175939in}{1.652980in}}%
\pgfpathlineto{\pgfqpoint{2.192026in}{1.653265in}}%
\pgfpathlineto{\pgfqpoint{2.203782in}{1.660005in}}%
\pgfpathlineto{\pgfqpoint{2.213376in}{1.656901in}}%
\pgfpathlineto{\pgfqpoint{2.215419in}{1.642645in}}%
\pgfpathlineto{\pgfqpoint{2.236002in}{1.640413in}}%
\pgfpathlineto{\pgfqpoint{2.247182in}{1.647086in}}%
\pgfpathlineto{\pgfqpoint{2.254932in}{1.632053in}}%
\pgfpathlineto{\pgfqpoint{2.269719in}{1.614670in}}%
\pgfpathlineto{\pgfqpoint{2.249022in}{1.614771in}}%
\pgfpathlineto{\pgfqpoint{2.242482in}{1.612622in}}%
\pgfpathlineto{\pgfqpoint{2.233519in}{1.615414in}}%
\pgfpathlineto{\pgfqpoint{2.232916in}{1.603380in}}%
\pgfpathlineto{\pgfqpoint{2.216640in}{1.612785in}}%
\pgfpathlineto{\pgfqpoint{2.195708in}{1.615659in}}%
\pgfpathlineto{\pgfqpoint{2.189946in}{1.606500in}}%
\pgfpathlineto{\pgfqpoint{2.162559in}{1.602046in}}%
\pgfpathlineto{\pgfqpoint{2.159409in}{1.594288in}}%
\pgfpathlineto{\pgfqpoint{2.140351in}{1.591960in}}%
\pgfpathlineto{\pgfqpoint{2.134590in}{1.584073in}}%
\pgfpathlineto{\pgfqpoint{2.124509in}{1.581958in}}%
\pgfpathlineto{\pgfqpoint{2.116458in}{1.563256in}}%
\pgfpathlineto{\pgfqpoint{2.106264in}{1.545142in}}%
\pgfpathlineto{\pgfqpoint{2.102264in}{1.543360in}}%
\pgfpathclose%
\pgfusepath{fill}%
\end{pgfscope}%
\begin{pgfscope}%
\pgfpathrectangle{\pgfqpoint{0.100000in}{0.100000in}}{\pgfqpoint{2.989028in}{1.913466in}}%
\pgfusepath{clip}%
\pgfsetbuttcap%
\pgfsetmiterjoin%
\definecolor{currentfill}{rgb}{0.711419,0.883276,0.634833}%
\pgfsetfillcolor{currentfill}%
\pgfsetlinewidth{0.000000pt}%
\definecolor{currentstroke}{rgb}{0.000000,0.000000,0.000000}%
\pgfsetstrokecolor{currentstroke}%
\pgfsetstrokeopacity{0.000000}%
\pgfsetdash{}{0pt}%
\pgfpathmoveto{\pgfqpoint{2.160936in}{1.325898in}}%
\pgfpathlineto{\pgfqpoint{2.170655in}{1.336128in}}%
\pgfpathlineto{\pgfqpoint{2.175083in}{1.350960in}}%
\pgfpathlineto{\pgfqpoint{2.180349in}{1.359557in}}%
\pgfpathlineto{\pgfqpoint{2.183578in}{1.371256in}}%
\pgfpathlineto{\pgfqpoint{2.184581in}{1.394584in}}%
\pgfpathlineto{\pgfqpoint{2.179720in}{1.416940in}}%
\pgfpathlineto{\pgfqpoint{2.163654in}{1.451174in}}%
\pgfpathlineto{\pgfqpoint{2.168005in}{1.461938in}}%
\pgfpathlineto{\pgfqpoint{2.162337in}{1.476816in}}%
\pgfpathlineto{\pgfqpoint{2.172059in}{1.497392in}}%
\pgfpathlineto{\pgfqpoint{2.170475in}{1.520335in}}%
\pgfpathlineto{\pgfqpoint{2.177247in}{1.523223in}}%
\pgfpathlineto{\pgfqpoint{2.178098in}{1.534139in}}%
\pgfpathlineto{\pgfqpoint{2.190119in}{1.541146in}}%
\pgfpathlineto{\pgfqpoint{2.196918in}{1.540015in}}%
\pgfpathlineto{\pgfqpoint{2.198829in}{1.528307in}}%
\pgfpathlineto{\pgfqpoint{2.204135in}{1.527848in}}%
\pgfpathlineto{\pgfqpoint{2.208977in}{1.544735in}}%
\pgfpathlineto{\pgfqpoint{2.207331in}{1.557869in}}%
\pgfpathlineto{\pgfqpoint{2.210479in}{1.565411in}}%
\pgfpathlineto{\pgfqpoint{2.227498in}{1.573202in}}%
\pgfpathlineto{\pgfqpoint{2.219743in}{1.575975in}}%
\pgfpathlineto{\pgfqpoint{2.217216in}{1.582668in}}%
\pgfpathlineto{\pgfqpoint{2.222803in}{1.594421in}}%
\pgfpathlineto{\pgfqpoint{2.233821in}{1.598481in}}%
\pgfpathlineto{\pgfqpoint{2.246615in}{1.591546in}}%
\pgfpathlineto{\pgfqpoint{2.258677in}{1.591460in}}%
\pgfpathlineto{\pgfqpoint{2.264279in}{1.583359in}}%
\pgfpathlineto{\pgfqpoint{2.272727in}{1.583922in}}%
\pgfpathlineto{\pgfqpoint{2.298804in}{1.572661in}}%
\pgfpathlineto{\pgfqpoint{2.304014in}{1.561760in}}%
\pgfpathlineto{\pgfqpoint{2.298324in}{1.557985in}}%
\pgfpathlineto{\pgfqpoint{2.300077in}{1.549815in}}%
\pgfpathlineto{\pgfqpoint{2.305699in}{1.546182in}}%
\pgfpathlineto{\pgfqpoint{2.308831in}{1.536139in}}%
\pgfpathlineto{\pgfqpoint{2.308393in}{1.511672in}}%
\pgfpathlineto{\pgfqpoint{2.300971in}{1.505827in}}%
\pgfpathlineto{\pgfqpoint{2.299305in}{1.492975in}}%
\pgfpathlineto{\pgfqpoint{2.285538in}{1.481098in}}%
\pgfpathlineto{\pgfqpoint{2.286355in}{1.466722in}}%
\pgfpathlineto{\pgfqpoint{2.298421in}{1.461685in}}%
\pgfpathlineto{\pgfqpoint{2.312034in}{1.479623in}}%
\pgfpathlineto{\pgfqpoint{2.313159in}{1.486129in}}%
\pgfpathlineto{\pgfqpoint{2.330136in}{1.496942in}}%
\pgfpathlineto{\pgfqpoint{2.340932in}{1.491933in}}%
\pgfpathlineto{\pgfqpoint{2.347707in}{1.480613in}}%
\pgfpathlineto{\pgfqpoint{2.358636in}{1.441310in}}%
\pgfpathlineto{\pgfqpoint{2.364428in}{1.428876in}}%
\pgfpathlineto{\pgfqpoint{2.362615in}{1.423735in}}%
\pgfpathlineto{\pgfqpoint{2.362791in}{1.406232in}}%
\pgfpathlineto{\pgfqpoint{2.352255in}{1.407912in}}%
\pgfpathlineto{\pgfqpoint{2.346305in}{1.394832in}}%
\pgfpathlineto{\pgfqpoint{2.345480in}{1.385940in}}%
\pgfpathlineto{\pgfqpoint{2.337522in}{1.380232in}}%
\pgfpathlineto{\pgfqpoint{2.335753in}{1.362891in}}%
\pgfpathlineto{\pgfqpoint{2.324180in}{1.340984in}}%
\pgfpathlineto{\pgfqpoint{2.260872in}{1.331556in}}%
\pgfpathlineto{\pgfqpoint{2.260498in}{1.335699in}}%
\pgfpathlineto{\pgfqpoint{2.218150in}{1.331236in}}%
\pgfpathlineto{\pgfqpoint{2.160936in}{1.325898in}}%
\pgfpathclose%
\pgfusepath{fill}%
\end{pgfscope}%
\begin{pgfscope}%
\pgfpathrectangle{\pgfqpoint{0.100000in}{0.100000in}}{\pgfqpoint{2.989028in}{1.913466in}}%
\pgfusepath{clip}%
\pgfsetbuttcap%
\pgfsetmiterjoin%
\definecolor{currentfill}{rgb}{0.601615,0.839677,0.644137}%
\pgfsetfillcolor{currentfill}%
\pgfsetlinewidth{0.000000pt}%
\definecolor{currentstroke}{rgb}{0.000000,0.000000,0.000000}%
\pgfsetstrokecolor{currentstroke}%
\pgfsetstrokeopacity{0.000000}%
\pgfsetdash{}{0pt}%
\pgfpathmoveto{\pgfqpoint{0.562031in}{0.446426in}}%
\pgfpathlineto{\pgfqpoint{0.560767in}{0.445314in}}%
\pgfpathlineto{\pgfqpoint{0.554967in}{0.446797in}}%
\pgfpathlineto{\pgfqpoint{0.555400in}{0.449716in}}%
\pgfpathlineto{\pgfqpoint{0.558683in}{0.451053in}}%
\pgfpathlineto{\pgfqpoint{0.563364in}{0.457782in}}%
\pgfpathlineto{\pgfqpoint{0.564142in}{0.462330in}}%
\pgfpathlineto{\pgfqpoint{0.566773in}{0.463522in}}%
\pgfpathlineto{\pgfqpoint{0.571054in}{0.463526in}}%
\pgfpathlineto{\pgfqpoint{0.575530in}{0.478285in}}%
\pgfpathlineto{\pgfqpoint{0.577283in}{0.479499in}}%
\pgfpathlineto{\pgfqpoint{0.575711in}{0.482384in}}%
\pgfpathlineto{\pgfqpoint{0.572030in}{0.479466in}}%
\pgfpathlineto{\pgfqpoint{0.561215in}{0.483400in}}%
\pgfpathlineto{\pgfqpoint{0.554760in}{0.490029in}}%
\pgfpathlineto{\pgfqpoint{0.557575in}{0.493053in}}%
\pgfpathlineto{\pgfqpoint{0.556295in}{0.497319in}}%
\pgfpathlineto{\pgfqpoint{0.557083in}{0.502359in}}%
\pgfpathlineto{\pgfqpoint{0.555695in}{0.508250in}}%
\pgfpathlineto{\pgfqpoint{0.555473in}{0.514148in}}%
\pgfpathlineto{\pgfqpoint{0.561686in}{0.513161in}}%
\pgfpathlineto{\pgfqpoint{0.562916in}{0.514899in}}%
\pgfpathlineto{\pgfqpoint{0.566006in}{0.514731in}}%
\pgfpathlineto{\pgfqpoint{0.569565in}{0.508441in}}%
\pgfpathlineto{\pgfqpoint{0.572996in}{0.503835in}}%
\pgfpathlineto{\pgfqpoint{0.570248in}{0.500837in}}%
\pgfpathlineto{\pgfqpoint{0.575000in}{0.499687in}}%
\pgfpathlineto{\pgfqpoint{0.576083in}{0.501158in}}%
\pgfpathlineto{\pgfqpoint{0.573542in}{0.507123in}}%
\pgfpathlineto{\pgfqpoint{0.568190in}{0.512068in}}%
\pgfpathlineto{\pgfqpoint{0.568845in}{0.514140in}}%
\pgfpathlineto{\pgfqpoint{0.564574in}{0.519376in}}%
\pgfpathlineto{\pgfqpoint{0.568935in}{0.520439in}}%
\pgfpathlineto{\pgfqpoint{0.565157in}{0.531650in}}%
\pgfpathlineto{\pgfqpoint{0.568229in}{0.533204in}}%
\pgfpathlineto{\pgfqpoint{0.568978in}{0.535892in}}%
\pgfpathlineto{\pgfqpoint{0.566980in}{0.539085in}}%
\pgfpathlineto{\pgfqpoint{0.571104in}{0.539884in}}%
\pgfpathlineto{\pgfqpoint{0.571845in}{0.542299in}}%
\pgfpathlineto{\pgfqpoint{0.576765in}{0.538674in}}%
\pgfpathlineto{\pgfqpoint{0.577857in}{0.542605in}}%
\pgfpathlineto{\pgfqpoint{0.580443in}{0.544334in}}%
\pgfpathlineto{\pgfqpoint{0.593023in}{0.546739in}}%
\pgfpathlineto{\pgfqpoint{0.596885in}{0.545889in}}%
\pgfpathlineto{\pgfqpoint{0.600120in}{0.550663in}}%
\pgfpathlineto{\pgfqpoint{0.607797in}{0.553794in}}%
\pgfpathlineto{\pgfqpoint{0.613710in}{0.552606in}}%
\pgfpathlineto{\pgfqpoint{0.616407in}{0.548449in}}%
\pgfpathlineto{\pgfqpoint{0.616266in}{0.541844in}}%
\pgfpathlineto{\pgfqpoint{0.619292in}{0.540347in}}%
\pgfpathlineto{\pgfqpoint{0.625464in}{0.540556in}}%
\pgfpathlineto{\pgfqpoint{0.633370in}{0.542274in}}%
\pgfpathlineto{\pgfqpoint{0.633272in}{0.540162in}}%
\pgfpathlineto{\pgfqpoint{0.643266in}{0.534139in}}%
\pgfpathlineto{\pgfqpoint{0.650255in}{0.535754in}}%
\pgfpathlineto{\pgfqpoint{0.652313in}{0.537844in}}%
\pgfpathlineto{\pgfqpoint{0.654516in}{0.543197in}}%
\pgfpathlineto{\pgfqpoint{0.657248in}{0.546625in}}%
\pgfpathlineto{\pgfqpoint{0.657492in}{0.551762in}}%
\pgfpathlineto{\pgfqpoint{0.654990in}{0.553782in}}%
\pgfpathlineto{\pgfqpoint{0.658695in}{0.555608in}}%
\pgfpathlineto{\pgfqpoint{0.664920in}{0.552865in}}%
\pgfpathlineto{\pgfqpoint{0.666948in}{0.554559in}}%
\pgfpathlineto{\pgfqpoint{0.667305in}{0.561826in}}%
\pgfpathlineto{\pgfqpoint{0.663281in}{0.560307in}}%
\pgfpathlineto{\pgfqpoint{0.660490in}{0.562825in}}%
\pgfpathlineto{\pgfqpoint{0.655960in}{0.564395in}}%
\pgfpathlineto{\pgfqpoint{0.648770in}{0.564133in}}%
\pgfpathlineto{\pgfqpoint{0.641880in}{0.566521in}}%
\pgfpathlineto{\pgfqpoint{0.641535in}{0.572469in}}%
\pgfpathlineto{\pgfqpoint{0.635373in}{0.578787in}}%
\pgfpathlineto{\pgfqpoint{0.627559in}{0.581439in}}%
\pgfpathlineto{\pgfqpoint{0.620739in}{0.593671in}}%
\pgfpathlineto{\pgfqpoint{0.621298in}{0.598197in}}%
\pgfpathlineto{\pgfqpoint{0.624971in}{0.600261in}}%
\pgfpathlineto{\pgfqpoint{0.624430in}{0.604545in}}%
\pgfpathlineto{\pgfqpoint{0.625282in}{0.608808in}}%
\pgfpathlineto{\pgfqpoint{0.628188in}{0.605890in}}%
\pgfpathlineto{\pgfqpoint{0.632025in}{0.607107in}}%
\pgfpathlineto{\pgfqpoint{0.631833in}{0.610810in}}%
\pgfpathlineto{\pgfqpoint{0.626514in}{0.618243in}}%
\pgfpathlineto{\pgfqpoint{0.624956in}{0.625911in}}%
\pgfpathlineto{\pgfqpoint{0.629409in}{0.626496in}}%
\pgfpathlineto{\pgfqpoint{0.636074in}{0.626025in}}%
\pgfpathlineto{\pgfqpoint{0.638172in}{0.623460in}}%
\pgfpathlineto{\pgfqpoint{0.644147in}{0.624378in}}%
\pgfpathlineto{\pgfqpoint{0.649700in}{0.623173in}}%
\pgfpathlineto{\pgfqpoint{0.653454in}{0.621055in}}%
\pgfpathlineto{\pgfqpoint{0.656052in}{0.622023in}}%
\pgfpathlineto{\pgfqpoint{0.662951in}{0.619906in}}%
\pgfpathlineto{\pgfqpoint{0.663014in}{0.618032in}}%
\pgfpathlineto{\pgfqpoint{0.668268in}{0.618665in}}%
\pgfpathlineto{\pgfqpoint{0.675313in}{0.614092in}}%
\pgfpathlineto{\pgfqpoint{0.675039in}{0.611057in}}%
\pgfpathlineto{\pgfqpoint{0.671687in}{0.609656in}}%
\pgfpathlineto{\pgfqpoint{0.668124in}{0.606323in}}%
\pgfpathlineto{\pgfqpoint{0.667694in}{0.601746in}}%
\pgfpathlineto{\pgfqpoint{0.675369in}{0.595909in}}%
\pgfpathlineto{\pgfqpoint{0.674141in}{0.593112in}}%
\pgfpathlineto{\pgfqpoint{0.679602in}{0.590964in}}%
\pgfpathlineto{\pgfqpoint{0.680943in}{0.587455in}}%
\pgfpathlineto{\pgfqpoint{0.687873in}{0.590243in}}%
\pgfpathlineto{\pgfqpoint{0.690925in}{0.586589in}}%
\pgfpathlineto{\pgfqpoint{0.692334in}{0.588801in}}%
\pgfpathlineto{\pgfqpoint{0.688186in}{0.595999in}}%
\pgfpathlineto{\pgfqpoint{0.689660in}{0.598155in}}%
\pgfpathlineto{\pgfqpoint{0.690313in}{0.603841in}}%
\pgfpathlineto{\pgfqpoint{0.688661in}{0.606342in}}%
\pgfpathlineto{\pgfqpoint{0.689901in}{0.609704in}}%
\pgfpathlineto{\pgfqpoint{0.693961in}{0.609388in}}%
\pgfpathlineto{\pgfqpoint{0.693121in}{0.603668in}}%
\pgfpathlineto{\pgfqpoint{0.690529in}{0.601760in}}%
\pgfpathlineto{\pgfqpoint{0.690712in}{0.594065in}}%
\pgfpathlineto{\pgfqpoint{0.695364in}{0.593146in}}%
\pgfpathlineto{\pgfqpoint{0.695421in}{0.587424in}}%
\pgfpathlineto{\pgfqpoint{0.700501in}{0.583713in}}%
\pgfpathlineto{\pgfqpoint{0.702836in}{0.586053in}}%
\pgfpathlineto{\pgfqpoint{0.700284in}{0.593176in}}%
\pgfpathlineto{\pgfqpoint{0.697644in}{0.594979in}}%
\pgfpathlineto{\pgfqpoint{0.695047in}{0.593793in}}%
\pgfpathlineto{\pgfqpoint{0.692763in}{0.595306in}}%
\pgfpathlineto{\pgfqpoint{0.693161in}{0.601889in}}%
\pgfpathlineto{\pgfqpoint{0.696974in}{0.604581in}}%
\pgfpathlineto{\pgfqpoint{0.700726in}{0.604333in}}%
\pgfpathlineto{\pgfqpoint{0.699172in}{0.608058in}}%
\pgfpathlineto{\pgfqpoint{0.693547in}{0.610911in}}%
\pgfpathlineto{\pgfqpoint{0.686098in}{0.622277in}}%
\pgfpathlineto{\pgfqpoint{0.689816in}{0.627728in}}%
\pgfpathlineto{\pgfqpoint{0.691885in}{0.632863in}}%
\pgfpathlineto{\pgfqpoint{0.692048in}{0.643174in}}%
\pgfpathlineto{\pgfqpoint{0.691251in}{0.651872in}}%
\pgfpathlineto{\pgfqpoint{0.688978in}{0.657584in}}%
\pgfpathlineto{\pgfqpoint{0.689431in}{0.663294in}}%
\pgfpathlineto{\pgfqpoint{0.695244in}{0.668247in}}%
\pgfpathlineto{\pgfqpoint{0.701026in}{0.674133in}}%
\pgfpathlineto{\pgfqpoint{0.715519in}{0.662108in}}%
\pgfpathlineto{\pgfqpoint{0.724995in}{0.661221in}}%
\pgfpathlineto{\pgfqpoint{0.732817in}{0.664216in}}%
\pgfpathlineto{\pgfqpoint{0.739683in}{0.668159in}}%
\pgfpathlineto{\pgfqpoint{0.741141in}{0.670038in}}%
\pgfpathlineto{\pgfqpoint{0.758090in}{0.674322in}}%
\pgfpathlineto{\pgfqpoint{0.759562in}{0.671184in}}%
\pgfpathlineto{\pgfqpoint{0.765321in}{0.668615in}}%
\pgfpathlineto{\pgfqpoint{0.776313in}{0.669403in}}%
\pgfpathlineto{\pgfqpoint{0.778320in}{0.668702in}}%
\pgfpathlineto{\pgfqpoint{0.783260in}{0.670576in}}%
\pgfpathlineto{\pgfqpoint{0.785246in}{0.667335in}}%
\pgfpathlineto{\pgfqpoint{0.789990in}{0.664374in}}%
\pgfpathlineto{\pgfqpoint{0.792363in}{0.664709in}}%
\pgfpathlineto{\pgfqpoint{0.796956in}{0.661151in}}%
\pgfpathlineto{\pgfqpoint{0.804717in}{0.661805in}}%
\pgfpathlineto{\pgfqpoint{0.812469in}{0.664685in}}%
\pgfpathlineto{\pgfqpoint{0.818800in}{0.655475in}}%
\pgfpathlineto{\pgfqpoint{0.818309in}{0.653408in}}%
\pgfpathlineto{\pgfqpoint{0.811323in}{0.651032in}}%
\pgfpathlineto{\pgfqpoint{0.813190in}{0.647440in}}%
\pgfpathlineto{\pgfqpoint{0.819621in}{0.651221in}}%
\pgfpathlineto{\pgfqpoint{0.822246in}{0.650305in}}%
\pgfpathlineto{\pgfqpoint{0.823363in}{0.645783in}}%
\pgfpathlineto{\pgfqpoint{0.819524in}{0.643870in}}%
\pgfpathlineto{\pgfqpoint{0.822266in}{0.638154in}}%
\pgfpathlineto{\pgfqpoint{0.826248in}{0.639120in}}%
\pgfpathlineto{\pgfqpoint{0.832026in}{0.636170in}}%
\pgfpathlineto{\pgfqpoint{0.837551in}{0.627982in}}%
\pgfpathlineto{\pgfqpoint{0.832689in}{0.625708in}}%
\pgfpathlineto{\pgfqpoint{0.836985in}{0.619675in}}%
\pgfpathlineto{\pgfqpoint{0.833366in}{0.618264in}}%
\pgfpathlineto{\pgfqpoint{0.837868in}{0.612534in}}%
\pgfpathlineto{\pgfqpoint{0.843358in}{0.613145in}}%
\pgfpathlineto{\pgfqpoint{0.846249in}{0.610567in}}%
\pgfpathlineto{\pgfqpoint{0.853562in}{0.606765in}}%
\pgfpathlineto{\pgfqpoint{0.862995in}{0.593271in}}%
\pgfpathlineto{\pgfqpoint{0.862782in}{0.590043in}}%
\pgfpathlineto{\pgfqpoint{0.876928in}{0.579706in}}%
\pgfpathlineto{\pgfqpoint{0.877233in}{0.576049in}}%
\pgfpathlineto{\pgfqpoint{0.881003in}{0.570578in}}%
\pgfpathlineto{\pgfqpoint{0.892585in}{0.567843in}}%
\pgfpathlineto{\pgfqpoint{0.896881in}{0.565899in}}%
\pgfpathlineto{\pgfqpoint{0.900922in}{0.559868in}}%
\pgfpathlineto{\pgfqpoint{0.901442in}{0.554585in}}%
\pgfpathlineto{\pgfqpoint{0.904915in}{0.550279in}}%
\pgfpathlineto{\pgfqpoint{0.905958in}{0.545423in}}%
\pgfpathlineto{\pgfqpoint{0.909213in}{0.543649in}}%
\pgfpathlineto{\pgfqpoint{0.893089in}{0.514548in}}%
\pgfpathlineto{\pgfqpoint{0.859622in}{0.454133in}}%
\pgfpathlineto{\pgfqpoint{0.810518in}{0.365475in}}%
\pgfpathlineto{\pgfqpoint{0.791294in}{0.330756in}}%
\pgfpathlineto{\pgfqpoint{0.795388in}{0.326093in}}%
\pgfpathlineto{\pgfqpoint{0.797192in}{0.327581in}}%
\pgfpathlineto{\pgfqpoint{0.800850in}{0.322152in}}%
\pgfpathlineto{\pgfqpoint{0.805927in}{0.323727in}}%
\pgfpathlineto{\pgfqpoint{0.812667in}{0.320475in}}%
\pgfpathlineto{\pgfqpoint{0.808301in}{0.315393in}}%
\pgfpathlineto{\pgfqpoint{0.811672in}{0.308535in}}%
\pgfpathlineto{\pgfqpoint{0.810976in}{0.305196in}}%
\pgfpathlineto{\pgfqpoint{0.816353in}{0.287587in}}%
\pgfpathlineto{\pgfqpoint{0.814077in}{0.279580in}}%
\pgfpathlineto{\pgfqpoint{0.824350in}{0.281562in}}%
\pgfpathlineto{\pgfqpoint{0.826872in}{0.280263in}}%
\pgfpathlineto{\pgfqpoint{0.829610in}{0.282376in}}%
\pgfpathlineto{\pgfqpoint{0.831558in}{0.286368in}}%
\pgfpathlineto{\pgfqpoint{0.834347in}{0.288657in}}%
\pgfpathlineto{\pgfqpoint{0.846252in}{0.288367in}}%
\pgfpathlineto{\pgfqpoint{0.848943in}{0.280683in}}%
\pgfpathlineto{\pgfqpoint{0.846634in}{0.273938in}}%
\pgfpathlineto{\pgfqpoint{0.849257in}{0.271712in}}%
\pgfpathlineto{\pgfqpoint{0.850636in}{0.267905in}}%
\pgfpathlineto{\pgfqpoint{0.850336in}{0.260746in}}%
\pgfpathlineto{\pgfqpoint{0.853836in}{0.255613in}}%
\pgfpathlineto{\pgfqpoint{0.855755in}{0.247568in}}%
\pgfpathlineto{\pgfqpoint{0.854770in}{0.243223in}}%
\pgfpathlineto{\pgfqpoint{0.855958in}{0.238819in}}%
\pgfpathlineto{\pgfqpoint{0.857671in}{0.213318in}}%
\pgfpathlineto{\pgfqpoint{0.855216in}{0.211309in}}%
\pgfpathlineto{\pgfqpoint{0.858455in}{0.208533in}}%
\pgfpathlineto{\pgfqpoint{0.855910in}{0.205157in}}%
\pgfpathlineto{\pgfqpoint{0.858206in}{0.202405in}}%
\pgfpathlineto{\pgfqpoint{0.856727in}{0.197707in}}%
\pgfpathlineto{\pgfqpoint{0.860034in}{0.196433in}}%
\pgfpathlineto{\pgfqpoint{0.864287in}{0.189287in}}%
\pgfpathlineto{\pgfqpoint{0.867447in}{0.187134in}}%
\pgfpathlineto{\pgfqpoint{0.868274in}{0.183987in}}%
\pgfpathlineto{\pgfqpoint{0.870105in}{0.182623in}}%
\pgfpathlineto{\pgfqpoint{0.869723in}{0.180010in}}%
\pgfpathlineto{\pgfqpoint{0.873670in}{0.178116in}}%
\pgfpathlineto{\pgfqpoint{0.872696in}{0.173088in}}%
\pgfpathlineto{\pgfqpoint{0.869298in}{0.170454in}}%
\pgfpathlineto{\pgfqpoint{0.867755in}{0.165560in}}%
\pgfpathlineto{\pgfqpoint{0.867345in}{0.159006in}}%
\pgfpathlineto{\pgfqpoint{0.865458in}{0.158379in}}%
\pgfpathlineto{\pgfqpoint{0.859289in}{0.152634in}}%
\pgfpathlineto{\pgfqpoint{0.853806in}{0.151118in}}%
\pgfpathlineto{\pgfqpoint{0.851190in}{0.153435in}}%
\pgfpathlineto{\pgfqpoint{0.853718in}{0.159583in}}%
\pgfpathlineto{\pgfqpoint{0.852946in}{0.162601in}}%
\pgfpathlineto{\pgfqpoint{0.856568in}{0.164131in}}%
\pgfpathlineto{\pgfqpoint{0.859956in}{0.173027in}}%
\pgfpathlineto{\pgfqpoint{0.858525in}{0.180625in}}%
\pgfpathlineto{\pgfqpoint{0.856235in}{0.182419in}}%
\pgfpathlineto{\pgfqpoint{0.848578in}{0.181838in}}%
\pgfpathlineto{\pgfqpoint{0.850121in}{0.179853in}}%
\pgfpathlineto{\pgfqpoint{0.844579in}{0.174155in}}%
\pgfpathlineto{\pgfqpoint{0.842793in}{0.177323in}}%
\pgfpathlineto{\pgfqpoint{0.843258in}{0.181462in}}%
\pgfpathlineto{\pgfqpoint{0.846439in}{0.181386in}}%
\pgfpathlineto{\pgfqpoint{0.849272in}{0.184426in}}%
\pgfpathlineto{\pgfqpoint{0.850948in}{0.188866in}}%
\pgfpathlineto{\pgfqpoint{0.857264in}{0.187613in}}%
\pgfpathlineto{\pgfqpoint{0.852208in}{0.189984in}}%
\pgfpathlineto{\pgfqpoint{0.852653in}{0.193161in}}%
\pgfpathlineto{\pgfqpoint{0.850549in}{0.194596in}}%
\pgfpathlineto{\pgfqpoint{0.849687in}{0.198923in}}%
\pgfpathlineto{\pgfqpoint{0.851181in}{0.201177in}}%
\pgfpathlineto{\pgfqpoint{0.847619in}{0.208179in}}%
\pgfpathlineto{\pgfqpoint{0.847919in}{0.212633in}}%
\pgfpathlineto{\pgfqpoint{0.843746in}{0.216599in}}%
\pgfpathlineto{\pgfqpoint{0.841488in}{0.219808in}}%
\pgfpathlineto{\pgfqpoint{0.844606in}{0.222866in}}%
\pgfpathlineto{\pgfqpoint{0.845761in}{0.227437in}}%
\pgfpathlineto{\pgfqpoint{0.844730in}{0.230438in}}%
\pgfpathlineto{\pgfqpoint{0.845936in}{0.234017in}}%
\pgfpathlineto{\pgfqpoint{0.844399in}{0.245686in}}%
\pgfpathlineto{\pgfqpoint{0.842100in}{0.250980in}}%
\pgfpathlineto{\pgfqpoint{0.839471in}{0.252970in}}%
\pgfpathlineto{\pgfqpoint{0.840590in}{0.259857in}}%
\pgfpathlineto{\pgfqpoint{0.839801in}{0.265235in}}%
\pgfpathlineto{\pgfqpoint{0.841230in}{0.268685in}}%
\pgfpathlineto{\pgfqpoint{0.838272in}{0.269169in}}%
\pgfpathlineto{\pgfqpoint{0.837479in}{0.260789in}}%
\pgfpathlineto{\pgfqpoint{0.834154in}{0.251441in}}%
\pgfpathlineto{\pgfqpoint{0.831484in}{0.253570in}}%
\pgfpathlineto{\pgfqpoint{0.831127in}{0.256642in}}%
\pgfpathlineto{\pgfqpoint{0.827861in}{0.260839in}}%
\pgfpathlineto{\pgfqpoint{0.829521in}{0.263587in}}%
\pgfpathlineto{\pgfqpoint{0.829005in}{0.269772in}}%
\pgfpathlineto{\pgfqpoint{0.824778in}{0.267597in}}%
\pgfpathlineto{\pgfqpoint{0.825607in}{0.264242in}}%
\pgfpathlineto{\pgfqpoint{0.824811in}{0.259983in}}%
\pgfpathlineto{\pgfqpoint{0.820066in}{0.259990in}}%
\pgfpathlineto{\pgfqpoint{0.817939in}{0.262630in}}%
\pgfpathlineto{\pgfqpoint{0.814729in}{0.263134in}}%
\pgfpathlineto{\pgfqpoint{0.812569in}{0.265977in}}%
\pgfpathlineto{\pgfqpoint{0.808852in}{0.273782in}}%
\pgfpathlineto{\pgfqpoint{0.807648in}{0.279842in}}%
\pgfpathlineto{\pgfqpoint{0.808418in}{0.282418in}}%
\pgfpathlineto{\pgfqpoint{0.806620in}{0.288091in}}%
\pgfpathlineto{\pgfqpoint{0.803639in}{0.291375in}}%
\pgfpathlineto{\pgfqpoint{0.798483in}{0.299768in}}%
\pgfpathlineto{\pgfqpoint{0.794744in}{0.307057in}}%
\pgfpathlineto{\pgfqpoint{0.798559in}{0.307167in}}%
\pgfpathlineto{\pgfqpoint{0.800804in}{0.308698in}}%
\pgfpathlineto{\pgfqpoint{0.801246in}{0.313435in}}%
\pgfpathlineto{\pgfqpoint{0.790450in}{0.314007in}}%
\pgfpathlineto{\pgfqpoint{0.781267in}{0.324264in}}%
\pgfpathlineto{\pgfqpoint{0.783682in}{0.326962in}}%
\pgfpathlineto{\pgfqpoint{0.779241in}{0.327423in}}%
\pgfpathlineto{\pgfqpoint{0.770423in}{0.336959in}}%
\pgfpathlineto{\pgfqpoint{0.756451in}{0.341741in}}%
\pgfpathlineto{\pgfqpoint{0.754570in}{0.347345in}}%
\pgfpathlineto{\pgfqpoint{0.751671in}{0.350264in}}%
\pgfpathlineto{\pgfqpoint{0.751539in}{0.353005in}}%
\pgfpathlineto{\pgfqpoint{0.749358in}{0.355041in}}%
\pgfpathlineto{\pgfqpoint{0.752220in}{0.358442in}}%
\pgfpathlineto{\pgfqpoint{0.745693in}{0.358770in}}%
\pgfpathlineto{\pgfqpoint{0.744034in}{0.363173in}}%
\pgfpathlineto{\pgfqpoint{0.741063in}{0.365047in}}%
\pgfpathlineto{\pgfqpoint{0.747012in}{0.367309in}}%
\pgfpathlineto{\pgfqpoint{0.743918in}{0.371345in}}%
\pgfpathlineto{\pgfqpoint{0.738224in}{0.373038in}}%
\pgfpathlineto{\pgfqpoint{0.739110in}{0.380491in}}%
\pgfpathlineto{\pgfqpoint{0.737478in}{0.383130in}}%
\pgfpathlineto{\pgfqpoint{0.734538in}{0.384988in}}%
\pgfpathlineto{\pgfqpoint{0.732024in}{0.383124in}}%
\pgfpathlineto{\pgfqpoint{0.730538in}{0.384877in}}%
\pgfpathlineto{\pgfqpoint{0.726206in}{0.385379in}}%
\pgfpathlineto{\pgfqpoint{0.726555in}{0.392143in}}%
\pgfpathlineto{\pgfqpoint{0.720275in}{0.387574in}}%
\pgfpathlineto{\pgfqpoint{0.720407in}{0.378160in}}%
\pgfpathlineto{\pgfqpoint{0.713666in}{0.376355in}}%
\pgfpathlineto{\pgfqpoint{0.709263in}{0.372596in}}%
\pgfpathlineto{\pgfqpoint{0.705725in}{0.371804in}}%
\pgfpathlineto{\pgfqpoint{0.701895in}{0.375593in}}%
\pgfpathlineto{\pgfqpoint{0.698485in}{0.375090in}}%
\pgfpathlineto{\pgfqpoint{0.699118in}{0.378551in}}%
\pgfpathlineto{\pgfqpoint{0.694270in}{0.376688in}}%
\pgfpathlineto{\pgfqpoint{0.689779in}{0.376478in}}%
\pgfpathlineto{\pgfqpoint{0.684523in}{0.374589in}}%
\pgfpathlineto{\pgfqpoint{0.673855in}{0.375652in}}%
\pgfpathlineto{\pgfqpoint{0.670430in}{0.374953in}}%
\pgfpathlineto{\pgfqpoint{0.667959in}{0.378065in}}%
\pgfpathlineto{\pgfqpoint{0.663397in}{0.378376in}}%
\pgfpathlineto{\pgfqpoint{0.662449in}{0.382409in}}%
\pgfpathlineto{\pgfqpoint{0.665167in}{0.384606in}}%
\pgfpathlineto{\pgfqpoint{0.668357in}{0.384399in}}%
\pgfpathlineto{\pgfqpoint{0.671019in}{0.381403in}}%
\pgfpathlineto{\pgfqpoint{0.672772in}{0.386005in}}%
\pgfpathlineto{\pgfqpoint{0.670326in}{0.391102in}}%
\pgfpathlineto{\pgfqpoint{0.675863in}{0.395589in}}%
\pgfpathlineto{\pgfqpoint{0.681411in}{0.397243in}}%
\pgfpathlineto{\pgfqpoint{0.685273in}{0.399940in}}%
\pgfpathlineto{\pgfqpoint{0.687775in}{0.402846in}}%
\pgfpathlineto{\pgfqpoint{0.689172in}{0.407609in}}%
\pgfpathlineto{\pgfqpoint{0.693583in}{0.406312in}}%
\pgfpathlineto{\pgfqpoint{0.703148in}{0.406979in}}%
\pgfpathlineto{\pgfqpoint{0.705201in}{0.401756in}}%
\pgfpathlineto{\pgfqpoint{0.708556in}{0.401529in}}%
\pgfpathlineto{\pgfqpoint{0.714662in}{0.393534in}}%
\pgfpathlineto{\pgfqpoint{0.714673in}{0.396448in}}%
\pgfpathlineto{\pgfqpoint{0.712540in}{0.397396in}}%
\pgfpathlineto{\pgfqpoint{0.708636in}{0.406940in}}%
\pgfpathlineto{\pgfqpoint{0.710536in}{0.408236in}}%
\pgfpathlineto{\pgfqpoint{0.705322in}{0.413215in}}%
\pgfpathlineto{\pgfqpoint{0.700663in}{0.414184in}}%
\pgfpathlineto{\pgfqpoint{0.696223in}{0.412463in}}%
\pgfpathlineto{\pgfqpoint{0.688619in}{0.413804in}}%
\pgfpathlineto{\pgfqpoint{0.685013in}{0.411392in}}%
\pgfpathlineto{\pgfqpoint{0.676801in}{0.409529in}}%
\pgfpathlineto{\pgfqpoint{0.676050in}{0.406847in}}%
\pgfpathlineto{\pgfqpoint{0.669741in}{0.406115in}}%
\pgfpathlineto{\pgfqpoint{0.668122in}{0.403049in}}%
\pgfpathlineto{\pgfqpoint{0.664513in}{0.400803in}}%
\pgfpathlineto{\pgfqpoint{0.660154in}{0.401189in}}%
\pgfpathlineto{\pgfqpoint{0.657977in}{0.399018in}}%
\pgfpathlineto{\pgfqpoint{0.655231in}{0.399119in}}%
\pgfpathlineto{\pgfqpoint{0.652607in}{0.401074in}}%
\pgfpathlineto{\pgfqpoint{0.648019in}{0.400725in}}%
\pgfpathlineto{\pgfqpoint{0.646938in}{0.398924in}}%
\pgfpathlineto{\pgfqpoint{0.642119in}{0.400773in}}%
\pgfpathlineto{\pgfqpoint{0.638027in}{0.396000in}}%
\pgfpathlineto{\pgfqpoint{0.641950in}{0.391382in}}%
\pgfpathlineto{\pgfqpoint{0.643460in}{0.387584in}}%
\pgfpathlineto{\pgfqpoint{0.642068in}{0.384367in}}%
\pgfpathlineto{\pgfqpoint{0.636493in}{0.382148in}}%
\pgfpathlineto{\pgfqpoint{0.633263in}{0.383954in}}%
\pgfpathlineto{\pgfqpoint{0.629198in}{0.382976in}}%
\pgfpathlineto{\pgfqpoint{0.628708in}{0.380157in}}%
\pgfpathlineto{\pgfqpoint{0.623759in}{0.381154in}}%
\pgfpathlineto{\pgfqpoint{0.624855in}{0.378337in}}%
\pgfpathlineto{\pgfqpoint{0.617555in}{0.377743in}}%
\pgfpathlineto{\pgfqpoint{0.612830in}{0.381319in}}%
\pgfpathlineto{\pgfqpoint{0.610306in}{0.379045in}}%
\pgfpathlineto{\pgfqpoint{0.603590in}{0.381466in}}%
\pgfpathlineto{\pgfqpoint{0.601918in}{0.379014in}}%
\pgfpathlineto{\pgfqpoint{0.598588in}{0.377923in}}%
\pgfpathlineto{\pgfqpoint{0.595917in}{0.380797in}}%
\pgfpathlineto{\pgfqpoint{0.586784in}{0.380083in}}%
\pgfpathlineto{\pgfqpoint{0.586794in}{0.375843in}}%
\pgfpathlineto{\pgfqpoint{0.583125in}{0.374721in}}%
\pgfpathlineto{\pgfqpoint{0.580270in}{0.376900in}}%
\pgfpathlineto{\pgfqpoint{0.576497in}{0.375689in}}%
\pgfpathlineto{\pgfqpoint{0.571614in}{0.375401in}}%
\pgfpathlineto{\pgfqpoint{0.564340in}{0.378466in}}%
\pgfpathlineto{\pgfqpoint{0.560174in}{0.375305in}}%
\pgfpathlineto{\pgfqpoint{0.554216in}{0.379530in}}%
\pgfpathlineto{\pgfqpoint{0.552268in}{0.377929in}}%
\pgfpathlineto{\pgfqpoint{0.550900in}{0.373791in}}%
\pgfpathlineto{\pgfqpoint{0.547071in}{0.371365in}}%
\pgfpathlineto{\pgfqpoint{0.540825in}{0.374142in}}%
\pgfpathlineto{\pgfqpoint{0.530479in}{0.381003in}}%
\pgfpathlineto{\pgfqpoint{0.527463in}{0.381613in}}%
\pgfpathlineto{\pgfqpoint{0.525762in}{0.380183in}}%
\pgfpathlineto{\pgfqpoint{0.521211in}{0.381155in}}%
\pgfpathlineto{\pgfqpoint{0.518991in}{0.383290in}}%
\pgfpathlineto{\pgfqpoint{0.514445in}{0.384441in}}%
\pgfpathlineto{\pgfqpoint{0.508238in}{0.384507in}}%
\pgfpathlineto{\pgfqpoint{0.505773in}{0.387054in}}%
\pgfpathlineto{\pgfqpoint{0.510497in}{0.389690in}}%
\pgfpathlineto{\pgfqpoint{0.509395in}{0.392291in}}%
\pgfpathlineto{\pgfqpoint{0.506149in}{0.391707in}}%
\pgfpathlineto{\pgfqpoint{0.504528in}{0.389662in}}%
\pgfpathlineto{\pgfqpoint{0.496930in}{0.386730in}}%
\pgfpathlineto{\pgfqpoint{0.493088in}{0.387727in}}%
\pgfpathlineto{\pgfqpoint{0.491471in}{0.393681in}}%
\pgfpathlineto{\pgfqpoint{0.486679in}{0.391121in}}%
\pgfpathlineto{\pgfqpoint{0.485286in}{0.398667in}}%
\pgfpathlineto{\pgfqpoint{0.483015in}{0.396722in}}%
\pgfpathlineto{\pgfqpoint{0.483395in}{0.393819in}}%
\pgfpathlineto{\pgfqpoint{0.478173in}{0.394586in}}%
\pgfpathlineto{\pgfqpoint{0.483605in}{0.399646in}}%
\pgfpathlineto{\pgfqpoint{0.489110in}{0.396639in}}%
\pgfpathlineto{\pgfqpoint{0.490291in}{0.398148in}}%
\pgfpathlineto{\pgfqpoint{0.496082in}{0.396339in}}%
\pgfpathlineto{\pgfqpoint{0.495720in}{0.398399in}}%
\pgfpathlineto{\pgfqpoint{0.511521in}{0.399081in}}%
\pgfpathlineto{\pgfqpoint{0.519230in}{0.395868in}}%
\pgfpathlineto{\pgfqpoint{0.522212in}{0.392698in}}%
\pgfpathlineto{\pgfqpoint{0.521327in}{0.389931in}}%
\pgfpathlineto{\pgfqpoint{0.524323in}{0.388883in}}%
\pgfpathlineto{\pgfqpoint{0.531765in}{0.393251in}}%
\pgfpathlineto{\pgfqpoint{0.541392in}{0.393531in}}%
\pgfpathlineto{\pgfqpoint{0.550136in}{0.390898in}}%
\pgfpathlineto{\pgfqpoint{0.555089in}{0.391340in}}%
\pgfpathlineto{\pgfqpoint{0.556933in}{0.387255in}}%
\pgfpathlineto{\pgfqpoint{0.559703in}{0.391758in}}%
\pgfpathlineto{\pgfqpoint{0.561407in}{0.392733in}}%
\pgfpathlineto{\pgfqpoint{0.569288in}{0.394691in}}%
\pgfpathlineto{\pgfqpoint{0.572217in}{0.393494in}}%
\pgfpathlineto{\pgfqpoint{0.575385in}{0.394915in}}%
\pgfpathlineto{\pgfqpoint{0.579163in}{0.394172in}}%
\pgfpathlineto{\pgfqpoint{0.580055in}{0.395827in}}%
\pgfpathlineto{\pgfqpoint{0.587939in}{0.403420in}}%
\pgfpathlineto{\pgfqpoint{0.591278in}{0.404545in}}%
\pgfpathlineto{\pgfqpoint{0.593251in}{0.408644in}}%
\pgfpathlineto{\pgfqpoint{0.603722in}{0.411817in}}%
\pgfpathlineto{\pgfqpoint{0.605033in}{0.414238in}}%
\pgfpathlineto{\pgfqpoint{0.590192in}{0.417974in}}%
\pgfpathlineto{\pgfqpoint{0.590776in}{0.421993in}}%
\pgfpathlineto{\pgfqpoint{0.589250in}{0.427301in}}%
\pgfpathlineto{\pgfqpoint{0.585794in}{0.426860in}}%
\pgfpathlineto{\pgfqpoint{0.583443in}{0.419972in}}%
\pgfpathlineto{\pgfqpoint{0.580714in}{0.419404in}}%
\pgfpathlineto{\pgfqpoint{0.578938in}{0.421634in}}%
\pgfpathlineto{\pgfqpoint{0.581010in}{0.431388in}}%
\pgfpathlineto{\pgfqpoint{0.578947in}{0.435475in}}%
\pgfpathlineto{\pgfqpoint{0.574950in}{0.440293in}}%
\pgfpathlineto{\pgfqpoint{0.576871in}{0.443969in}}%
\pgfpathlineto{\pgfqpoint{0.568436in}{0.444102in}}%
\pgfpathlineto{\pgfqpoint{0.568628in}{0.445528in}}%
\pgfpathlineto{\pgfqpoint{0.563595in}{0.447526in}}%
\pgfpathlineto{\pgfqpoint{0.562031in}{0.446426in}}%
\pgfpathclose%
\pgfusepath{fill}%
\end{pgfscope}%
\begin{pgfscope}%
\pgfpathrectangle{\pgfqpoint{0.100000in}{0.100000in}}{\pgfqpoint{2.989028in}{1.913466in}}%
\pgfusepath{clip}%
\pgfsetbuttcap%
\pgfsetmiterjoin%
\definecolor{currentfill}{rgb}{0.601615,0.839677,0.644137}%
\pgfsetfillcolor{currentfill}%
\pgfsetlinewidth{0.000000pt}%
\definecolor{currentstroke}{rgb}{0.000000,0.000000,0.000000}%
\pgfsetstrokecolor{currentstroke}%
\pgfsetstrokeopacity{0.000000}%
\pgfsetdash{}{0pt}%
\pgfpathmoveto{\pgfqpoint{0.546979in}{0.517466in}}%
\pgfpathlineto{\pgfqpoint{0.545936in}{0.515633in}}%
\pgfpathlineto{\pgfqpoint{0.548696in}{0.511730in}}%
\pgfpathlineto{\pgfqpoint{0.544334in}{0.508079in}}%
\pgfpathlineto{\pgfqpoint{0.544455in}{0.504939in}}%
\pgfpathlineto{\pgfqpoint{0.542483in}{0.504198in}}%
\pgfpathlineto{\pgfqpoint{0.535704in}{0.508881in}}%
\pgfpathlineto{\pgfqpoint{0.530872in}{0.519132in}}%
\pgfpathlineto{\pgfqpoint{0.530549in}{0.522094in}}%
\pgfpathlineto{\pgfqpoint{0.531791in}{0.525640in}}%
\pgfpathlineto{\pgfqpoint{0.537027in}{0.520216in}}%
\pgfpathlineto{\pgfqpoint{0.538628in}{0.521262in}}%
\pgfpathlineto{\pgfqpoint{0.543070in}{0.520267in}}%
\pgfpathlineto{\pgfqpoint{0.546979in}{0.517466in}}%
\pgfpathclose%
\pgfusepath{fill}%
\end{pgfscope}%
\begin{pgfscope}%
\pgfpathrectangle{\pgfqpoint{0.100000in}{0.100000in}}{\pgfqpoint{2.989028in}{1.913466in}}%
\pgfusepath{clip}%
\pgfsetbuttcap%
\pgfsetmiterjoin%
\definecolor{currentfill}{rgb}{0.601615,0.839677,0.644137}%
\pgfsetfillcolor{currentfill}%
\pgfsetlinewidth{0.000000pt}%
\definecolor{currentstroke}{rgb}{0.000000,0.000000,0.000000}%
\pgfsetstrokecolor{currentstroke}%
\pgfsetstrokeopacity{0.000000}%
\pgfsetdash{}{0pt}%
\pgfpathmoveto{\pgfqpoint{0.466584in}{0.398705in}}%
\pgfpathlineto{\pgfqpoint{0.461820in}{0.397932in}}%
\pgfpathlineto{\pgfqpoint{0.458372in}{0.399525in}}%
\pgfpathlineto{\pgfqpoint{0.456865in}{0.402120in}}%
\pgfpathlineto{\pgfqpoint{0.458573in}{0.405842in}}%
\pgfpathlineto{\pgfqpoint{0.462377in}{0.404858in}}%
\pgfpathlineto{\pgfqpoint{0.468630in}{0.407066in}}%
\pgfpathlineto{\pgfqpoint{0.470820in}{0.404016in}}%
\pgfpathlineto{\pgfqpoint{0.472944in}{0.404579in}}%
\pgfpathlineto{\pgfqpoint{0.478084in}{0.402681in}}%
\pgfpathlineto{\pgfqpoint{0.480305in}{0.400269in}}%
\pgfpathlineto{\pgfqpoint{0.477595in}{0.393963in}}%
\pgfpathlineto{\pgfqpoint{0.474899in}{0.392626in}}%
\pgfpathlineto{\pgfqpoint{0.472476in}{0.393346in}}%
\pgfpathlineto{\pgfqpoint{0.466584in}{0.398705in}}%
\pgfpathclose%
\pgfusepath{fill}%
\end{pgfscope}%
\begin{pgfscope}%
\pgfpathrectangle{\pgfqpoint{0.100000in}{0.100000in}}{\pgfqpoint{2.989028in}{1.913466in}}%
\pgfusepath{clip}%
\pgfsetbuttcap%
\pgfsetmiterjoin%
\definecolor{currentfill}{rgb}{0.601615,0.839677,0.644137}%
\pgfsetfillcolor{currentfill}%
\pgfsetlinewidth{0.000000pt}%
\definecolor{currentstroke}{rgb}{0.000000,0.000000,0.000000}%
\pgfsetstrokecolor{currentstroke}%
\pgfsetstrokeopacity{0.000000}%
\pgfsetdash{}{0pt}%
\pgfpathmoveto{\pgfqpoint{0.423985in}{0.404436in}}%
\pgfpathlineto{\pgfqpoint{0.417234in}{0.407656in}}%
\pgfpathlineto{\pgfqpoint{0.410987in}{0.408614in}}%
\pgfpathlineto{\pgfqpoint{0.409018in}{0.409985in}}%
\pgfpathlineto{\pgfqpoint{0.411260in}{0.412685in}}%
\pgfpathlineto{\pgfqpoint{0.415239in}{0.409312in}}%
\pgfpathlineto{\pgfqpoint{0.418576in}{0.409777in}}%
\pgfpathlineto{\pgfqpoint{0.424004in}{0.412113in}}%
\pgfpathlineto{\pgfqpoint{0.424442in}{0.415120in}}%
\pgfpathlineto{\pgfqpoint{0.431109in}{0.413632in}}%
\pgfpathlineto{\pgfqpoint{0.429548in}{0.409103in}}%
\pgfpathlineto{\pgfqpoint{0.433794in}{0.408772in}}%
\pgfpathlineto{\pgfqpoint{0.430231in}{0.403385in}}%
\pgfpathlineto{\pgfqpoint{0.426781in}{0.405131in}}%
\pgfpathlineto{\pgfqpoint{0.423985in}{0.404436in}}%
\pgfpathclose%
\pgfusepath{fill}%
\end{pgfscope}%
\begin{pgfscope}%
\pgfpathrectangle{\pgfqpoint{0.100000in}{0.100000in}}{\pgfqpoint{2.989028in}{1.913466in}}%
\pgfusepath{clip}%
\pgfsetbuttcap%
\pgfsetmiterjoin%
\definecolor{currentfill}{rgb}{0.601615,0.839677,0.644137}%
\pgfsetfillcolor{currentfill}%
\pgfsetlinewidth{0.000000pt}%
\definecolor{currentstroke}{rgb}{0.000000,0.000000,0.000000}%
\pgfsetstrokecolor{currentstroke}%
\pgfsetstrokeopacity{0.000000}%
\pgfsetdash{}{0pt}%
\pgfpathmoveto{\pgfqpoint{0.412124in}{0.415568in}}%
\pgfpathlineto{\pgfqpoint{0.409643in}{0.414219in}}%
\pgfpathlineto{\pgfqpoint{0.402882in}{0.416293in}}%
\pgfpathlineto{\pgfqpoint{0.397830in}{0.415092in}}%
\pgfpathlineto{\pgfqpoint{0.395264in}{0.418504in}}%
\pgfpathlineto{\pgfqpoint{0.396383in}{0.420066in}}%
\pgfpathlineto{\pgfqpoint{0.400212in}{0.420327in}}%
\pgfpathlineto{\pgfqpoint{0.403140in}{0.419346in}}%
\pgfpathlineto{\pgfqpoint{0.406331in}{0.420898in}}%
\pgfpathlineto{\pgfqpoint{0.411223in}{0.418823in}}%
\pgfpathlineto{\pgfqpoint{0.412124in}{0.415568in}}%
\pgfpathclose%
\pgfusepath{fill}%
\end{pgfscope}%
\begin{pgfscope}%
\pgfpathrectangle{\pgfqpoint{0.100000in}{0.100000in}}{\pgfqpoint{2.989028in}{1.913466in}}%
\pgfusepath{clip}%
\pgfsetbuttcap%
\pgfsetmiterjoin%
\definecolor{currentfill}{rgb}{0.601615,0.839677,0.644137}%
\pgfsetfillcolor{currentfill}%
\pgfsetlinewidth{0.000000pt}%
\definecolor{currentstroke}{rgb}{0.000000,0.000000,0.000000}%
\pgfsetstrokecolor{currentstroke}%
\pgfsetstrokeopacity{0.000000}%
\pgfsetdash{}{0pt}%
\pgfpathmoveto{\pgfqpoint{0.332429in}{0.465782in}}%
\pgfpathlineto{\pgfqpoint{0.332250in}{0.462436in}}%
\pgfpathlineto{\pgfqpoint{0.327314in}{0.460334in}}%
\pgfpathlineto{\pgfqpoint{0.322807in}{0.465597in}}%
\pgfpathlineto{\pgfqpoint{0.327862in}{0.467350in}}%
\pgfpathlineto{\pgfqpoint{0.332429in}{0.465782in}}%
\pgfpathclose%
\pgfusepath{fill}%
\end{pgfscope}%
\begin{pgfscope}%
\pgfpathrectangle{\pgfqpoint{0.100000in}{0.100000in}}{\pgfqpoint{2.989028in}{1.913466in}}%
\pgfusepath{clip}%
\pgfsetbuttcap%
\pgfsetmiterjoin%
\definecolor{currentfill}{rgb}{0.601615,0.839677,0.644137}%
\pgfsetfillcolor{currentfill}%
\pgfsetlinewidth{0.000000pt}%
\definecolor{currentstroke}{rgb}{0.000000,0.000000,0.000000}%
\pgfsetstrokecolor{currentstroke}%
\pgfsetstrokeopacity{0.000000}%
\pgfsetdash{}{0pt}%
\pgfpathmoveto{\pgfqpoint{0.291722in}{0.484589in}}%
\pgfpathlineto{\pgfqpoint{0.292771in}{0.486043in}}%
\pgfpathlineto{\pgfqpoint{0.298825in}{0.485792in}}%
\pgfpathlineto{\pgfqpoint{0.299887in}{0.487956in}}%
\pgfpathlineto{\pgfqpoint{0.302543in}{0.486068in}}%
\pgfpathlineto{\pgfqpoint{0.301009in}{0.482435in}}%
\pgfpathlineto{\pgfqpoint{0.298753in}{0.482127in}}%
\pgfpathlineto{\pgfqpoint{0.293361in}{0.483267in}}%
\pgfpathlineto{\pgfqpoint{0.291722in}{0.484589in}}%
\pgfpathclose%
\pgfusepath{fill}%
\end{pgfscope}%
\begin{pgfscope}%
\pgfpathrectangle{\pgfqpoint{0.100000in}{0.100000in}}{\pgfqpoint{2.989028in}{1.913466in}}%
\pgfusepath{clip}%
\pgfsetbuttcap%
\pgfsetmiterjoin%
\definecolor{currentfill}{rgb}{0.601615,0.839677,0.644137}%
\pgfsetfillcolor{currentfill}%
\pgfsetlinewidth{0.000000pt}%
\definecolor{currentstroke}{rgb}{0.000000,0.000000,0.000000}%
\pgfsetstrokecolor{currentstroke}%
\pgfsetstrokeopacity{0.000000}%
\pgfsetdash{}{0pt}%
\pgfpathmoveto{\pgfqpoint{0.283137in}{0.495848in}}%
\pgfpathlineto{\pgfqpoint{0.282770in}{0.499125in}}%
\pgfpathlineto{\pgfqpoint{0.285929in}{0.498882in}}%
\pgfpathlineto{\pgfqpoint{0.285709in}{0.503484in}}%
\pgfpathlineto{\pgfqpoint{0.288886in}{0.501000in}}%
\pgfpathlineto{\pgfqpoint{0.287787in}{0.497411in}}%
\pgfpathlineto{\pgfqpoint{0.283137in}{0.495848in}}%
\pgfpathclose%
\pgfusepath{fill}%
\end{pgfscope}%
\begin{pgfscope}%
\pgfpathrectangle{\pgfqpoint{0.100000in}{0.100000in}}{\pgfqpoint{2.989028in}{1.913466in}}%
\pgfusepath{clip}%
\pgfsetbuttcap%
\pgfsetmiterjoin%
\definecolor{currentfill}{rgb}{0.601615,0.839677,0.644137}%
\pgfsetfillcolor{currentfill}%
\pgfsetlinewidth{0.000000pt}%
\definecolor{currentstroke}{rgb}{0.000000,0.000000,0.000000}%
\pgfsetstrokecolor{currentstroke}%
\pgfsetstrokeopacity{0.000000}%
\pgfsetdash{}{0pt}%
\pgfpathmoveto{\pgfqpoint{0.561438in}{0.612633in}}%
\pgfpathlineto{\pgfqpoint{0.556400in}{0.613721in}}%
\pgfpathlineto{\pgfqpoint{0.555318in}{0.617088in}}%
\pgfpathlineto{\pgfqpoint{0.557049in}{0.620453in}}%
\pgfpathlineto{\pgfqpoint{0.560927in}{0.622211in}}%
\pgfpathlineto{\pgfqpoint{0.565367in}{0.613572in}}%
\pgfpathlineto{\pgfqpoint{0.569395in}{0.613189in}}%
\pgfpathlineto{\pgfqpoint{0.572399in}{0.610523in}}%
\pgfpathlineto{\pgfqpoint{0.572340in}{0.606855in}}%
\pgfpathlineto{\pgfqpoint{0.570143in}{0.605097in}}%
\pgfpathlineto{\pgfqpoint{0.573716in}{0.594249in}}%
\pgfpathlineto{\pgfqpoint{0.577625in}{0.590269in}}%
\pgfpathlineto{\pgfqpoint{0.573628in}{0.589006in}}%
\pgfpathlineto{\pgfqpoint{0.570393in}{0.593472in}}%
\pgfpathlineto{\pgfqpoint{0.563363in}{0.592089in}}%
\pgfpathlineto{\pgfqpoint{0.565549in}{0.596509in}}%
\pgfpathlineto{\pgfqpoint{0.564531in}{0.598269in}}%
\pgfpathlineto{\pgfqpoint{0.565128in}{0.602876in}}%
\pgfpathlineto{\pgfqpoint{0.563735in}{0.611204in}}%
\pgfpathlineto{\pgfqpoint{0.561438in}{0.612633in}}%
\pgfpathclose%
\pgfusepath{fill}%
\end{pgfscope}%
\begin{pgfscope}%
\pgfpathrectangle{\pgfqpoint{0.100000in}{0.100000in}}{\pgfqpoint{2.989028in}{1.913466in}}%
\pgfusepath{clip}%
\pgfsetbuttcap%
\pgfsetmiterjoin%
\definecolor{currentfill}{rgb}{0.601615,0.839677,0.644137}%
\pgfsetfillcolor{currentfill}%
\pgfsetlinewidth{0.000000pt}%
\definecolor{currentstroke}{rgb}{0.000000,0.000000,0.000000}%
\pgfsetstrokecolor{currentstroke}%
\pgfsetstrokeopacity{0.000000}%
\pgfsetdash{}{0pt}%
\pgfpathmoveto{\pgfqpoint{0.738441in}{0.364487in}}%
\pgfpathlineto{\pgfqpoint{0.731138in}{0.364569in}}%
\pgfpathlineto{\pgfqpoint{0.731700in}{0.368235in}}%
\pgfpathlineto{\pgfqpoint{0.734623in}{0.369546in}}%
\pgfpathlineto{\pgfqpoint{0.738441in}{0.364487in}}%
\pgfpathclose%
\pgfusepath{fill}%
\end{pgfscope}%
\begin{pgfscope}%
\pgfpathrectangle{\pgfqpoint{0.100000in}{0.100000in}}{\pgfqpoint{2.989028in}{1.913466in}}%
\pgfusepath{clip}%
\pgfsetbuttcap%
\pgfsetmiterjoin%
\definecolor{currentfill}{rgb}{0.601615,0.839677,0.644137}%
\pgfsetfillcolor{currentfill}%
\pgfsetlinewidth{0.000000pt}%
\definecolor{currentstroke}{rgb}{0.000000,0.000000,0.000000}%
\pgfsetstrokecolor{currentstroke}%
\pgfsetstrokeopacity{0.000000}%
\pgfsetdash{}{0pt}%
\pgfpathmoveto{\pgfqpoint{0.728306in}{0.367807in}}%
\pgfpathlineto{\pgfqpoint{0.723074in}{0.366541in}}%
\pgfpathlineto{\pgfqpoint{0.718102in}{0.364373in}}%
\pgfpathlineto{\pgfqpoint{0.716567in}{0.366927in}}%
\pgfpathlineto{\pgfqpoint{0.725117in}{0.368657in}}%
\pgfpathlineto{\pgfqpoint{0.728306in}{0.367807in}}%
\pgfpathclose%
\pgfusepath{fill}%
\end{pgfscope}%
\begin{pgfscope}%
\pgfpathrectangle{\pgfqpoint{0.100000in}{0.100000in}}{\pgfqpoint{2.989028in}{1.913466in}}%
\pgfusepath{clip}%
\pgfsetbuttcap%
\pgfsetmiterjoin%
\definecolor{currentfill}{rgb}{0.601615,0.839677,0.644137}%
\pgfsetfillcolor{currentfill}%
\pgfsetlinewidth{0.000000pt}%
\definecolor{currentstroke}{rgb}{0.000000,0.000000,0.000000}%
\pgfsetstrokecolor{currentstroke}%
\pgfsetstrokeopacity{0.000000}%
\pgfsetdash{}{0pt}%
\pgfpathmoveto{\pgfqpoint{0.648103in}{0.365212in}}%
\pgfpathlineto{\pgfqpoint{0.648303in}{0.362676in}}%
\pgfpathlineto{\pgfqpoint{0.645288in}{0.361100in}}%
\pgfpathlineto{\pgfqpoint{0.636875in}{0.364714in}}%
\pgfpathlineto{\pgfqpoint{0.635570in}{0.363261in}}%
\pgfpathlineto{\pgfqpoint{0.633535in}{0.370138in}}%
\pgfpathlineto{\pgfqpoint{0.637750in}{0.371095in}}%
\pgfpathlineto{\pgfqpoint{0.639698in}{0.368269in}}%
\pgfpathlineto{\pgfqpoint{0.643929in}{0.371807in}}%
\pgfpathlineto{\pgfqpoint{0.648103in}{0.365212in}}%
\pgfpathclose%
\pgfusepath{fill}%
\end{pgfscope}%
\begin{pgfscope}%
\pgfpathrectangle{\pgfqpoint{0.100000in}{0.100000in}}{\pgfqpoint{2.989028in}{1.913466in}}%
\pgfusepath{clip}%
\pgfsetbuttcap%
\pgfsetmiterjoin%
\definecolor{currentfill}{rgb}{0.601615,0.839677,0.644137}%
\pgfsetfillcolor{currentfill}%
\pgfsetlinewidth{0.000000pt}%
\definecolor{currentstroke}{rgb}{0.000000,0.000000,0.000000}%
\pgfsetstrokecolor{currentstroke}%
\pgfsetstrokeopacity{0.000000}%
\pgfsetdash{}{0pt}%
\pgfpathmoveto{\pgfqpoint{0.837890in}{0.224669in}}%
\pgfpathlineto{\pgfqpoint{0.830835in}{0.222042in}}%
\pgfpathlineto{\pgfqpoint{0.827098in}{0.221885in}}%
\pgfpathlineto{\pgfqpoint{0.830139in}{0.227457in}}%
\pgfpathlineto{\pgfqpoint{0.832738in}{0.229535in}}%
\pgfpathlineto{\pgfqpoint{0.832654in}{0.235306in}}%
\pgfpathlineto{\pgfqpoint{0.835479in}{0.245723in}}%
\pgfpathlineto{\pgfqpoint{0.838081in}{0.247702in}}%
\pgfpathlineto{\pgfqpoint{0.844087in}{0.244837in}}%
\pgfpathlineto{\pgfqpoint{0.843210in}{0.243196in}}%
\pgfpathlineto{\pgfqpoint{0.843665in}{0.235310in}}%
\pgfpathlineto{\pgfqpoint{0.841718in}{0.233182in}}%
\pgfpathlineto{\pgfqpoint{0.840881in}{0.227581in}}%
\pgfpathlineto{\pgfqpoint{0.837890in}{0.224669in}}%
\pgfpathclose%
\pgfusepath{fill}%
\end{pgfscope}%
\begin{pgfscope}%
\pgfpathrectangle{\pgfqpoint{0.100000in}{0.100000in}}{\pgfqpoint{2.989028in}{1.913466in}}%
\pgfusepath{clip}%
\pgfsetbuttcap%
\pgfsetmiterjoin%
\definecolor{currentfill}{rgb}{0.601615,0.839677,0.644137}%
\pgfsetfillcolor{currentfill}%
\pgfsetlinewidth{0.000000pt}%
\definecolor{currentstroke}{rgb}{0.000000,0.000000,0.000000}%
\pgfsetstrokecolor{currentstroke}%
\pgfsetstrokeopacity{0.000000}%
\pgfsetdash{}{0pt}%
\pgfpathmoveto{\pgfqpoint{0.822266in}{0.250492in}}%
\pgfpathlineto{\pgfqpoint{0.826399in}{0.243355in}}%
\pgfpathlineto{\pgfqpoint{0.825633in}{0.241653in}}%
\pgfpathlineto{\pgfqpoint{0.831161in}{0.239922in}}%
\pgfpathlineto{\pgfqpoint{0.829574in}{0.233362in}}%
\pgfpathlineto{\pgfqpoint{0.827167in}{0.233665in}}%
\pgfpathlineto{\pgfqpoint{0.820842in}{0.245101in}}%
\pgfpathlineto{\pgfqpoint{0.821742in}{0.239954in}}%
\pgfpathlineto{\pgfqpoint{0.820664in}{0.236990in}}%
\pgfpathlineto{\pgfqpoint{0.817986in}{0.235671in}}%
\pgfpathlineto{\pgfqpoint{0.816463in}{0.237092in}}%
\pgfpathlineto{\pgfqpoint{0.815613in}{0.243682in}}%
\pgfpathlineto{\pgfqpoint{0.816115in}{0.247956in}}%
\pgfpathlineto{\pgfqpoint{0.814331in}{0.249822in}}%
\pgfpathlineto{\pgfqpoint{0.819043in}{0.255220in}}%
\pgfpathlineto{\pgfqpoint{0.822135in}{0.256822in}}%
\pgfpathlineto{\pgfqpoint{0.823165in}{0.254717in}}%
\pgfpathlineto{\pgfqpoint{0.826391in}{0.256346in}}%
\pgfpathlineto{\pgfqpoint{0.828728in}{0.251922in}}%
\pgfpathlineto{\pgfqpoint{0.833783in}{0.246295in}}%
\pgfpathlineto{\pgfqpoint{0.833093in}{0.243746in}}%
\pgfpathlineto{\pgfqpoint{0.830286in}{0.240969in}}%
\pgfpathlineto{\pgfqpoint{0.827691in}{0.242277in}}%
\pgfpathlineto{\pgfqpoint{0.822266in}{0.250492in}}%
\pgfpathclose%
\pgfusepath{fill}%
\end{pgfscope}%
\begin{pgfscope}%
\pgfpathrectangle{\pgfqpoint{0.100000in}{0.100000in}}{\pgfqpoint{2.989028in}{1.913466in}}%
\pgfusepath{clip}%
\pgfsetbuttcap%
\pgfsetmiterjoin%
\definecolor{currentfill}{rgb}{0.601615,0.839677,0.644137}%
\pgfsetfillcolor{currentfill}%
\pgfsetlinewidth{0.000000pt}%
\definecolor{currentstroke}{rgb}{0.000000,0.000000,0.000000}%
\pgfsetstrokecolor{currentstroke}%
\pgfsetstrokeopacity{0.000000}%
\pgfsetdash{}{0pt}%
\pgfpathmoveto{\pgfqpoint{0.617161in}{0.351684in}}%
\pgfpathlineto{\pgfqpoint{0.612041in}{0.351018in}}%
\pgfpathlineto{\pgfqpoint{0.604830in}{0.348246in}}%
\pgfpathlineto{\pgfqpoint{0.603148in}{0.349752in}}%
\pgfpathlineto{\pgfqpoint{0.610680in}{0.351454in}}%
\pgfpathlineto{\pgfqpoint{0.608661in}{0.352939in}}%
\pgfpathlineto{\pgfqpoint{0.603991in}{0.354027in}}%
\pgfpathlineto{\pgfqpoint{0.602744in}{0.357507in}}%
\pgfpathlineto{\pgfqpoint{0.605094in}{0.360506in}}%
\pgfpathlineto{\pgfqpoint{0.607073in}{0.367387in}}%
\pgfpathlineto{\pgfqpoint{0.614301in}{0.368669in}}%
\pgfpathlineto{\pgfqpoint{0.617801in}{0.365696in}}%
\pgfpathlineto{\pgfqpoint{0.621088in}{0.368573in}}%
\pgfpathlineto{\pgfqpoint{0.625211in}{0.367867in}}%
\pgfpathlineto{\pgfqpoint{0.628174in}{0.362553in}}%
\pgfpathlineto{\pgfqpoint{0.630683in}{0.367277in}}%
\pgfpathlineto{\pgfqpoint{0.637804in}{0.356228in}}%
\pgfpathlineto{\pgfqpoint{0.635232in}{0.355345in}}%
\pgfpathlineto{\pgfqpoint{0.633653in}{0.352950in}}%
\pgfpathlineto{\pgfqpoint{0.636636in}{0.350764in}}%
\pgfpathlineto{\pgfqpoint{0.632011in}{0.348715in}}%
\pgfpathlineto{\pgfqpoint{0.628656in}{0.349843in}}%
\pgfpathlineto{\pgfqpoint{0.622321in}{0.353583in}}%
\pgfpathlineto{\pgfqpoint{0.617161in}{0.351684in}}%
\pgfpathclose%
\pgfusepath{fill}%
\end{pgfscope}%
\begin{pgfscope}%
\pgfpathrectangle{\pgfqpoint{0.100000in}{0.100000in}}{\pgfqpoint{2.989028in}{1.913466in}}%
\pgfusepath{clip}%
\pgfsetbuttcap%
\pgfsetmiterjoin%
\definecolor{currentfill}{rgb}{0.601615,0.839677,0.644137}%
\pgfsetfillcolor{currentfill}%
\pgfsetlinewidth{0.000000pt}%
\definecolor{currentstroke}{rgb}{0.000000,0.000000,0.000000}%
\pgfsetstrokecolor{currentstroke}%
\pgfsetstrokeopacity{0.000000}%
\pgfsetdash{}{0pt}%
\pgfpathmoveto{\pgfqpoint{0.817967in}{0.203589in}}%
\pgfpathlineto{\pgfqpoint{0.816952in}{0.213793in}}%
\pgfpathlineto{\pgfqpoint{0.819688in}{0.217539in}}%
\pgfpathlineto{\pgfqpoint{0.816762in}{0.217800in}}%
\pgfpathlineto{\pgfqpoint{0.817832in}{0.225086in}}%
\pgfpathlineto{\pgfqpoint{0.819093in}{0.229412in}}%
\pgfpathlineto{\pgfqpoint{0.817997in}{0.235224in}}%
\pgfpathlineto{\pgfqpoint{0.820677in}{0.236349in}}%
\pgfpathlineto{\pgfqpoint{0.821551in}{0.237852in}}%
\pgfpathlineto{\pgfqpoint{0.824191in}{0.237282in}}%
\pgfpathlineto{\pgfqpoint{0.826476in}{0.232626in}}%
\pgfpathlineto{\pgfqpoint{0.826901in}{0.228360in}}%
\pgfpathlineto{\pgfqpoint{0.824105in}{0.215583in}}%
\pgfpathlineto{\pgfqpoint{0.817967in}{0.203589in}}%
\pgfpathclose%
\pgfusepath{fill}%
\end{pgfscope}%
\begin{pgfscope}%
\pgfpathrectangle{\pgfqpoint{0.100000in}{0.100000in}}{\pgfqpoint{2.989028in}{1.913466in}}%
\pgfusepath{clip}%
\pgfsetbuttcap%
\pgfsetmiterjoin%
\definecolor{currentfill}{rgb}{0.601615,0.839677,0.644137}%
\pgfsetfillcolor{currentfill}%
\pgfsetlinewidth{0.000000pt}%
\definecolor{currentstroke}{rgb}{0.000000,0.000000,0.000000}%
\pgfsetstrokecolor{currentstroke}%
\pgfsetstrokeopacity{0.000000}%
\pgfsetdash{}{0pt}%
\pgfpathmoveto{\pgfqpoint{0.817360in}{0.234737in}}%
\pgfpathlineto{\pgfqpoint{0.817219in}{0.229762in}}%
\pgfpathlineto{\pgfqpoint{0.814925in}{0.227503in}}%
\pgfpathlineto{\pgfqpoint{0.812273in}{0.228338in}}%
\pgfpathlineto{\pgfqpoint{0.815625in}{0.235565in}}%
\pgfpathlineto{\pgfqpoint{0.817360in}{0.234737in}}%
\pgfpathclose%
\pgfusepath{fill}%
\end{pgfscope}%
\begin{pgfscope}%
\pgfpathrectangle{\pgfqpoint{0.100000in}{0.100000in}}{\pgfqpoint{2.989028in}{1.913466in}}%
\pgfusepath{clip}%
\pgfsetbuttcap%
\pgfsetmiterjoin%
\definecolor{currentfill}{rgb}{0.601615,0.839677,0.644137}%
\pgfsetfillcolor{currentfill}%
\pgfsetlinewidth{0.000000pt}%
\definecolor{currentstroke}{rgb}{0.000000,0.000000,0.000000}%
\pgfsetstrokecolor{currentstroke}%
\pgfsetstrokeopacity{0.000000}%
\pgfsetdash{}{0pt}%
\pgfpathmoveto{\pgfqpoint{0.845513in}{0.212648in}}%
\pgfpathlineto{\pgfqpoint{0.843557in}{0.204322in}}%
\pgfpathlineto{\pgfqpoint{0.838998in}{0.201690in}}%
\pgfpathlineto{\pgfqpoint{0.833017in}{0.204115in}}%
\pgfpathlineto{\pgfqpoint{0.834665in}{0.207350in}}%
\pgfpathlineto{\pgfqpoint{0.834632in}{0.209591in}}%
\pgfpathlineto{\pgfqpoint{0.836855in}{0.213398in}}%
\pgfpathlineto{\pgfqpoint{0.835292in}{0.212632in}}%
\pgfpathlineto{\pgfqpoint{0.836432in}{0.214491in}}%
\pgfpathlineto{\pgfqpoint{0.834267in}{0.218715in}}%
\pgfpathlineto{\pgfqpoint{0.836723in}{0.219493in}}%
\pgfpathlineto{\pgfqpoint{0.842410in}{0.214551in}}%
\pgfpathlineto{\pgfqpoint{0.845513in}{0.212648in}}%
\pgfpathclose%
\pgfusepath{fill}%
\end{pgfscope}%
\begin{pgfscope}%
\pgfpathrectangle{\pgfqpoint{0.100000in}{0.100000in}}{\pgfqpoint{2.989028in}{1.913466in}}%
\pgfusepath{clip}%
\pgfsetbuttcap%
\pgfsetmiterjoin%
\definecolor{currentfill}{rgb}{0.601615,0.839677,0.644137}%
\pgfsetfillcolor{currentfill}%
\pgfsetlinewidth{0.000000pt}%
\definecolor{currentstroke}{rgb}{0.000000,0.000000,0.000000}%
\pgfsetstrokecolor{currentstroke}%
\pgfsetstrokeopacity{0.000000}%
\pgfsetdash{}{0pt}%
\pgfpathmoveto{\pgfqpoint{0.825557in}{0.197926in}}%
\pgfpathlineto{\pgfqpoint{0.822979in}{0.196923in}}%
\pgfpathlineto{\pgfqpoint{0.824451in}{0.205709in}}%
\pgfpathlineto{\pgfqpoint{0.827546in}{0.207760in}}%
\pgfpathlineto{\pgfqpoint{0.826791in}{0.214343in}}%
\pgfpathlineto{\pgfqpoint{0.828049in}{0.217262in}}%
\pgfpathlineto{\pgfqpoint{0.830536in}{0.218356in}}%
\pgfpathlineto{\pgfqpoint{0.833148in}{0.215493in}}%
\pgfpathlineto{\pgfqpoint{0.835389in}{0.211607in}}%
\pgfpathlineto{\pgfqpoint{0.828184in}{0.203995in}}%
\pgfpathlineto{\pgfqpoint{0.825557in}{0.197926in}}%
\pgfpathclose%
\pgfusepath{fill}%
\end{pgfscope}%
\begin{pgfscope}%
\pgfpathrectangle{\pgfqpoint{0.100000in}{0.100000in}}{\pgfqpoint{2.989028in}{1.913466in}}%
\pgfusepath{clip}%
\pgfsetbuttcap%
\pgfsetmiterjoin%
\definecolor{currentfill}{rgb}{0.601615,0.839677,0.644137}%
\pgfsetfillcolor{currentfill}%
\pgfsetlinewidth{0.000000pt}%
\definecolor{currentstroke}{rgb}{0.000000,0.000000,0.000000}%
\pgfsetstrokecolor{currentstroke}%
\pgfsetstrokeopacity{0.000000}%
\pgfsetdash{}{0pt}%
\pgfpathmoveto{\pgfqpoint{0.846657in}{0.207083in}}%
\pgfpathlineto{\pgfqpoint{0.847921in}{0.201020in}}%
\pgfpathlineto{\pgfqpoint{0.842275in}{0.201714in}}%
\pgfpathlineto{\pgfqpoint{0.846657in}{0.207083in}}%
\pgfpathclose%
\pgfusepath{fill}%
\end{pgfscope}%
\begin{pgfscope}%
\pgfpathrectangle{\pgfqpoint{0.100000in}{0.100000in}}{\pgfqpoint{2.989028in}{1.913466in}}%
\pgfusepath{clip}%
\pgfsetbuttcap%
\pgfsetmiterjoin%
\definecolor{currentfill}{rgb}{0.601615,0.839677,0.644137}%
\pgfsetfillcolor{currentfill}%
\pgfsetlinewidth{0.000000pt}%
\definecolor{currentstroke}{rgb}{0.000000,0.000000,0.000000}%
\pgfsetstrokecolor{currentstroke}%
\pgfsetstrokeopacity{0.000000}%
\pgfsetdash{}{0pt}%
\pgfpathmoveto{\pgfqpoint{0.849351in}{0.197819in}}%
\pgfpathlineto{\pgfqpoint{0.850276in}{0.194079in}}%
\pgfpathlineto{\pgfqpoint{0.852047in}{0.192882in}}%
\pgfpathlineto{\pgfqpoint{0.851725in}{0.189362in}}%
\pgfpathlineto{\pgfqpoint{0.849251in}{0.188177in}}%
\pgfpathlineto{\pgfqpoint{0.847242in}{0.193433in}}%
\pgfpathlineto{\pgfqpoint{0.849351in}{0.197819in}}%
\pgfpathclose%
\pgfusepath{fill}%
\end{pgfscope}%
\begin{pgfscope}%
\pgfpathrectangle{\pgfqpoint{0.100000in}{0.100000in}}{\pgfqpoint{2.989028in}{1.913466in}}%
\pgfusepath{clip}%
\pgfsetbuttcap%
\pgfsetmiterjoin%
\definecolor{currentfill}{rgb}{0.601615,0.839677,0.644137}%
\pgfsetfillcolor{currentfill}%
\pgfsetlinewidth{0.000000pt}%
\definecolor{currentstroke}{rgb}{0.000000,0.000000,0.000000}%
\pgfsetstrokecolor{currentstroke}%
\pgfsetstrokeopacity{0.000000}%
\pgfsetdash{}{0pt}%
\pgfpathmoveto{\pgfqpoint{0.844640in}{0.199254in}}%
\pgfpathlineto{\pgfqpoint{0.843104in}{0.194909in}}%
\pgfpathlineto{\pgfqpoint{0.840719in}{0.195069in}}%
\pgfpathlineto{\pgfqpoint{0.839215in}{0.198652in}}%
\pgfpathlineto{\pgfqpoint{0.841665in}{0.200349in}}%
\pgfpathlineto{\pgfqpoint{0.844640in}{0.199254in}}%
\pgfpathclose%
\pgfusepath{fill}%
\end{pgfscope}%
\begin{pgfscope}%
\pgfpathrectangle{\pgfqpoint{0.100000in}{0.100000in}}{\pgfqpoint{2.989028in}{1.913466in}}%
\pgfusepath{clip}%
\pgfsetbuttcap%
\pgfsetmiterjoin%
\definecolor{currentfill}{rgb}{0.601615,0.839677,0.644137}%
\pgfsetfillcolor{currentfill}%
\pgfsetlinewidth{0.000000pt}%
\definecolor{currentstroke}{rgb}{0.000000,0.000000,0.000000}%
\pgfsetstrokecolor{currentstroke}%
\pgfsetstrokeopacity{0.000000}%
\pgfsetdash{}{0pt}%
\pgfpathmoveto{\pgfqpoint{0.838205in}{0.176854in}}%
\pgfpathlineto{\pgfqpoint{0.840333in}{0.174527in}}%
\pgfpathlineto{\pgfqpoint{0.840802in}{0.169622in}}%
\pgfpathlineto{\pgfqpoint{0.841922in}{0.168197in}}%
\pgfpathlineto{\pgfqpoint{0.840092in}{0.163758in}}%
\pgfpathlineto{\pgfqpoint{0.837686in}{0.161260in}}%
\pgfpathlineto{\pgfqpoint{0.836405in}{0.156623in}}%
\pgfpathlineto{\pgfqpoint{0.833774in}{0.163848in}}%
\pgfpathlineto{\pgfqpoint{0.833338in}{0.170408in}}%
\pgfpathlineto{\pgfqpoint{0.831865in}{0.173673in}}%
\pgfpathlineto{\pgfqpoint{0.826795in}{0.176321in}}%
\pgfpathlineto{\pgfqpoint{0.831551in}{0.181841in}}%
\pgfpathlineto{\pgfqpoint{0.828380in}{0.184201in}}%
\pgfpathlineto{\pgfqpoint{0.831174in}{0.186552in}}%
\pgfpathlineto{\pgfqpoint{0.834882in}{0.195579in}}%
\pgfpathlineto{\pgfqpoint{0.831038in}{0.198762in}}%
\pgfpathlineto{\pgfqpoint{0.832519in}{0.201818in}}%
\pgfpathlineto{\pgfqpoint{0.837487in}{0.199107in}}%
\pgfpathlineto{\pgfqpoint{0.837990in}{0.194829in}}%
\pgfpathlineto{\pgfqpoint{0.836111in}{0.192394in}}%
\pgfpathlineto{\pgfqpoint{0.839836in}{0.189098in}}%
\pgfpathlineto{\pgfqpoint{0.841009in}{0.183965in}}%
\pgfpathlineto{\pgfqpoint{0.840850in}{0.176765in}}%
\pgfpathlineto{\pgfqpoint{0.838205in}{0.176854in}}%
\pgfpathclose%
\pgfusepath{fill}%
\end{pgfscope}%
\begin{pgfscope}%
\pgfpathrectangle{\pgfqpoint{0.100000in}{0.100000in}}{\pgfqpoint{2.989028in}{1.913466in}}%
\pgfusepath{clip}%
\pgfsetbuttcap%
\pgfsetmiterjoin%
\definecolor{currentfill}{rgb}{0.601615,0.839677,0.644137}%
\pgfsetfillcolor{currentfill}%
\pgfsetlinewidth{0.000000pt}%
\definecolor{currentstroke}{rgb}{0.000000,0.000000,0.000000}%
\pgfsetstrokecolor{currentstroke}%
\pgfsetstrokeopacity{0.000000}%
\pgfsetdash{}{0pt}%
\pgfpathmoveto{\pgfqpoint{0.847302in}{0.194926in}}%
\pgfpathlineto{\pgfqpoint{0.846143in}{0.192429in}}%
\pgfpathlineto{\pgfqpoint{0.847231in}{0.188374in}}%
\pgfpathlineto{\pgfqpoint{0.844445in}{0.188033in}}%
\pgfpathlineto{\pgfqpoint{0.841524in}{0.189911in}}%
\pgfpathlineto{\pgfqpoint{0.842472in}{0.194046in}}%
\pgfpathlineto{\pgfqpoint{0.847302in}{0.194926in}}%
\pgfpathclose%
\pgfusepath{fill}%
\end{pgfscope}%
\begin{pgfscope}%
\pgfpathrectangle{\pgfqpoint{0.100000in}{0.100000in}}{\pgfqpoint{2.989028in}{1.913466in}}%
\pgfusepath{clip}%
\pgfsetbuttcap%
\pgfsetmiterjoin%
\definecolor{currentfill}{rgb}{0.601615,0.839677,0.644137}%
\pgfsetfillcolor{currentfill}%
\pgfsetlinewidth{0.000000pt}%
\definecolor{currentstroke}{rgb}{0.000000,0.000000,0.000000}%
\pgfsetstrokecolor{currentstroke}%
\pgfsetstrokeopacity{0.000000}%
\pgfsetdash{}{0pt}%
\pgfpathmoveto{\pgfqpoint{0.834316in}{0.195209in}}%
\pgfpathlineto{\pgfqpoint{0.828167in}{0.192451in}}%
\pgfpathlineto{\pgfqpoint{0.826076in}{0.193444in}}%
\pgfpathlineto{\pgfqpoint{0.830270in}{0.197082in}}%
\pgfpathlineto{\pgfqpoint{0.834316in}{0.195209in}}%
\pgfpathclose%
\pgfusepath{fill}%
\end{pgfscope}%
\begin{pgfscope}%
\pgfpathrectangle{\pgfqpoint{0.100000in}{0.100000in}}{\pgfqpoint{2.989028in}{1.913466in}}%
\pgfusepath{clip}%
\pgfsetbuttcap%
\pgfsetmiterjoin%
\definecolor{currentfill}{rgb}{0.601615,0.839677,0.644137}%
\pgfsetfillcolor{currentfill}%
\pgfsetlinewidth{0.000000pt}%
\definecolor{currentstroke}{rgb}{0.000000,0.000000,0.000000}%
\pgfsetstrokecolor{currentstroke}%
\pgfsetstrokeopacity{0.000000}%
\pgfsetdash{}{0pt}%
\pgfpathmoveto{\pgfqpoint{0.847369in}{0.168302in}}%
\pgfpathlineto{\pgfqpoint{0.845903in}{0.172455in}}%
\pgfpathlineto{\pgfqpoint{0.849004in}{0.173477in}}%
\pgfpathlineto{\pgfqpoint{0.849943in}{0.177955in}}%
\pgfpathlineto{\pgfqpoint{0.853303in}{0.177960in}}%
\pgfpathlineto{\pgfqpoint{0.853289in}{0.181079in}}%
\pgfpathlineto{\pgfqpoint{0.857959in}{0.180341in}}%
\pgfpathlineto{\pgfqpoint{0.858755in}{0.171887in}}%
\pgfpathlineto{\pgfqpoint{0.856046in}{0.166498in}}%
\pgfpathlineto{\pgfqpoint{0.852027in}{0.163131in}}%
\pgfpathlineto{\pgfqpoint{0.847369in}{0.168302in}}%
\pgfpathclose%
\pgfusepath{fill}%
\end{pgfscope}%
\begin{pgfscope}%
\pgfpathrectangle{\pgfqpoint{0.100000in}{0.100000in}}{\pgfqpoint{2.989028in}{1.913466in}}%
\pgfusepath{clip}%
\pgfsetbuttcap%
\pgfsetmiterjoin%
\definecolor{currentfill}{rgb}{0.601615,0.839677,0.644137}%
\pgfsetfillcolor{currentfill}%
\pgfsetlinewidth{0.000000pt}%
\definecolor{currentstroke}{rgb}{0.000000,0.000000,0.000000}%
\pgfsetstrokecolor{currentstroke}%
\pgfsetstrokeopacity{0.000000}%
\pgfsetdash{}{0pt}%
\pgfpathmoveto{\pgfqpoint{0.845553in}{0.171496in}}%
\pgfpathlineto{\pgfqpoint{0.846878in}{0.167564in}}%
\pgfpathlineto{\pgfqpoint{0.843612in}{0.167086in}}%
\pgfpathlineto{\pgfqpoint{0.845553in}{0.171496in}}%
\pgfpathclose%
\pgfusepath{fill}%
\end{pgfscope}%
\begin{pgfscope}%
\pgfpathrectangle{\pgfqpoint{0.100000in}{0.100000in}}{\pgfqpoint{2.989028in}{1.913466in}}%
\pgfusepath{clip}%
\pgfsetbuttcap%
\pgfsetmiterjoin%
\definecolor{currentfill}{rgb}{0.601615,0.839677,0.644137}%
\pgfsetfillcolor{currentfill}%
\pgfsetlinewidth{0.000000pt}%
\definecolor{currentstroke}{rgb}{0.000000,0.000000,0.000000}%
\pgfsetstrokecolor{currentstroke}%
\pgfsetstrokeopacity{0.000000}%
\pgfsetdash{}{0pt}%
\pgfpathmoveto{\pgfqpoint{0.848306in}{0.165713in}}%
\pgfpathlineto{\pgfqpoint{0.847851in}{0.160700in}}%
\pgfpathlineto{\pgfqpoint{0.844406in}{0.161116in}}%
\pgfpathlineto{\pgfqpoint{0.845520in}{0.163537in}}%
\pgfpathlineto{\pgfqpoint{0.848306in}{0.165713in}}%
\pgfpathclose%
\pgfusepath{fill}%
\end{pgfscope}%
\begin{pgfscope}%
\pgfsetbuttcap%
\pgfsetmiterjoin%
\definecolor{currentfill}{rgb}{1.000000,1.000000,1.000000}%
\pgfsetfillcolor{currentfill}%
\pgfsetlinewidth{1.003750pt}%
\definecolor{currentstroke}{rgb}{0.827451,0.827451,0.827451}%
\pgfsetstrokecolor{currentstroke}%
\pgfsetdash{}{0pt}%
\pgfpathmoveto{\pgfqpoint{1.015239in}{1.931644in}}%
\pgfpathlineto{\pgfqpoint{2.626989in}{1.931644in}}%
\pgfpathlineto{\pgfqpoint{2.626989in}{2.127644in}}%
\pgfpathlineto{\pgfqpoint{1.015239in}{2.127644in}}%
\pgfpathlineto{\pgfqpoint{1.015239in}{1.931644in}}%
\pgfpathclose%
\pgfusepath{stroke,fill}%
\end{pgfscope}%
\begin{pgfscope}%
\definecolor{textcolor}{rgb}{0.000000,0.000000,0.000000}%
\pgfsetstrokecolor{textcolor}%
\pgfsetfillcolor{textcolor}%
\pgftext[x=1.056489in,y=1.994331in,left,base]{\color{textcolor}\setmainfont{Lato}\rmfamily\fontsize{9.000000}{10.800000}\selectfont Quarterly Growth, 2023 Q3}%
\end{pgfscope}%
\begin{pgfscope}%
\pgfpathrectangle{\pgfqpoint{3.625000in}{0.100000in}}{\pgfqpoint{2.989028in}{1.913466in}}%
\pgfusepath{clip}%
\pgfsetbuttcap%
\pgfsetmiterjoin%
\definecolor{currentfill}{rgb}{0.829373,0.931257,0.610842}%
\pgfsetfillcolor{currentfill}%
\pgfsetlinewidth{0.000000pt}%
\definecolor{currentstroke}{rgb}{0.000000,0.000000,0.000000}%
\pgfsetstrokecolor{currentstroke}%
\pgfsetstrokeopacity{0.000000}%
\pgfsetdash{}{0pt}%
\pgfpathmoveto{\pgfqpoint{4.622679in}{0.548675in}}%
\pgfpathlineto{\pgfqpoint{4.614089in}{0.548815in}}%
\pgfpathlineto{\pgfqpoint{4.608242in}{0.555287in}}%
\pgfpathlineto{\pgfqpoint{4.604266in}{0.572473in}}%
\pgfpathlineto{\pgfqpoint{4.613506in}{0.574627in}}%
\pgfpathlineto{\pgfqpoint{4.622893in}{0.572311in}}%
\pgfpathlineto{\pgfqpoint{4.632446in}{0.565540in}}%
\pgfpathlineto{\pgfqpoint{4.632441in}{0.556177in}}%
\pgfpathlineto{\pgfqpoint{4.622679in}{0.548675in}}%
\pgfpathclose%
\pgfusepath{fill}%
\end{pgfscope}%
\begin{pgfscope}%
\pgfpathrectangle{\pgfqpoint{3.625000in}{0.100000in}}{\pgfqpoint{2.989028in}{1.913466in}}%
\pgfusepath{clip}%
\pgfsetbuttcap%
\pgfsetmiterjoin%
\definecolor{currentfill}{rgb}{0.829373,0.931257,0.610842}%
\pgfsetfillcolor{currentfill}%
\pgfsetlinewidth{0.000000pt}%
\definecolor{currentstroke}{rgb}{0.000000,0.000000,0.000000}%
\pgfsetstrokecolor{currentstroke}%
\pgfsetstrokeopacity{0.000000}%
\pgfsetdash{}{0pt}%
\pgfpathmoveto{\pgfqpoint{4.674970in}{0.446640in}}%
\pgfpathlineto{\pgfqpoint{4.665359in}{0.450300in}}%
\pgfpathlineto{\pgfqpoint{4.665315in}{0.457479in}}%
\pgfpathlineto{\pgfqpoint{4.653705in}{0.463549in}}%
\pgfpathlineto{\pgfqpoint{4.655119in}{0.478835in}}%
\pgfpathlineto{\pgfqpoint{4.658481in}{0.486048in}}%
\pgfpathlineto{\pgfqpoint{4.665581in}{0.479409in}}%
\pgfpathlineto{\pgfqpoint{4.677127in}{0.480518in}}%
\pgfpathlineto{\pgfqpoint{4.674211in}{0.459391in}}%
\pgfpathlineto{\pgfqpoint{4.674970in}{0.446640in}}%
\pgfpathclose%
\pgfusepath{fill}%
\end{pgfscope}%
\begin{pgfscope}%
\pgfpathrectangle{\pgfqpoint{3.625000in}{0.100000in}}{\pgfqpoint{2.989028in}{1.913466in}}%
\pgfusepath{clip}%
\pgfsetbuttcap%
\pgfsetmiterjoin%
\definecolor{currentfill}{rgb}{0.829373,0.931257,0.610842}%
\pgfsetfillcolor{currentfill}%
\pgfsetlinewidth{0.000000pt}%
\definecolor{currentstroke}{rgb}{0.000000,0.000000,0.000000}%
\pgfsetstrokecolor{currentstroke}%
\pgfsetstrokeopacity{0.000000}%
\pgfsetdash{}{0pt}%
\pgfpathmoveto{\pgfqpoint{4.715310in}{0.399517in}}%
\pgfpathlineto{\pgfqpoint{4.703114in}{0.401594in}}%
\pgfpathlineto{\pgfqpoint{4.696369in}{0.412035in}}%
\pgfpathlineto{\pgfqpoint{4.695417in}{0.420620in}}%
\pgfpathlineto{\pgfqpoint{4.715310in}{0.399517in}}%
\pgfpathclose%
\pgfusepath{fill}%
\end{pgfscope}%
\begin{pgfscope}%
\pgfpathrectangle{\pgfqpoint{3.625000in}{0.100000in}}{\pgfqpoint{2.989028in}{1.913466in}}%
\pgfusepath{clip}%
\pgfsetbuttcap%
\pgfsetmiterjoin%
\definecolor{currentfill}{rgb}{0.829373,0.931257,0.610842}%
\pgfsetfillcolor{currentfill}%
\pgfsetlinewidth{0.000000pt}%
\definecolor{currentstroke}{rgb}{0.000000,0.000000,0.000000}%
\pgfsetstrokecolor{currentstroke}%
\pgfsetstrokeopacity{0.000000}%
\pgfsetdash{}{0pt}%
\pgfpathmoveto{\pgfqpoint{4.690317in}{0.401338in}}%
\pgfpathlineto{\pgfqpoint{4.696711in}{0.396147in}}%
\pgfpathlineto{\pgfqpoint{4.698099in}{0.387222in}}%
\pgfpathlineto{\pgfqpoint{4.685227in}{0.387715in}}%
\pgfpathlineto{\pgfqpoint{4.690317in}{0.401338in}}%
\pgfpathclose%
\pgfusepath{fill}%
\end{pgfscope}%
\begin{pgfscope}%
\pgfpathrectangle{\pgfqpoint{3.625000in}{0.100000in}}{\pgfqpoint{2.989028in}{1.913466in}}%
\pgfusepath{clip}%
\pgfsetbuttcap%
\pgfsetmiterjoin%
\definecolor{currentfill}{rgb}{0.829373,0.931257,0.610842}%
\pgfsetfillcolor{currentfill}%
\pgfsetlinewidth{0.000000pt}%
\definecolor{currentstroke}{rgb}{0.000000,0.000000,0.000000}%
\pgfsetstrokecolor{currentstroke}%
\pgfsetstrokeopacity{0.000000}%
\pgfsetdash{}{0pt}%
\pgfpathmoveto{\pgfqpoint{4.719683in}{0.351485in}}%
\pgfpathlineto{\pgfqpoint{4.706837in}{0.356103in}}%
\pgfpathlineto{\pgfqpoint{4.710518in}{0.367298in}}%
\pgfpathlineto{\pgfqpoint{4.705134in}{0.378705in}}%
\pgfpathlineto{\pgfqpoint{4.706825in}{0.386546in}}%
\pgfpathlineto{\pgfqpoint{4.717162in}{0.389535in}}%
\pgfpathlineto{\pgfqpoint{4.716963in}{0.376972in}}%
\pgfpathlineto{\pgfqpoint{4.729190in}{0.370003in}}%
\pgfpathlineto{\pgfqpoint{4.736183in}{0.353155in}}%
\pgfpathlineto{\pgfqpoint{4.728389in}{0.346667in}}%
\pgfpathlineto{\pgfqpoint{4.719683in}{0.351485in}}%
\pgfpathclose%
\pgfusepath{fill}%
\end{pgfscope}%
\begin{pgfscope}%
\pgfpathrectangle{\pgfqpoint{3.625000in}{0.100000in}}{\pgfqpoint{2.989028in}{1.913466in}}%
\pgfusepath{clip}%
\pgfsetbuttcap%
\pgfsetmiterjoin%
\definecolor{currentfill}{rgb}{0.829373,0.931257,0.610842}%
\pgfsetfillcolor{currentfill}%
\pgfsetlinewidth{0.000000pt}%
\definecolor{currentstroke}{rgb}{0.000000,0.000000,0.000000}%
\pgfsetstrokecolor{currentstroke}%
\pgfsetstrokeopacity{0.000000}%
\pgfsetdash{}{0pt}%
\pgfpathmoveto{\pgfqpoint{4.673289in}{0.237569in}}%
\pgfpathlineto{\pgfqpoint{4.667169in}{0.250211in}}%
\pgfpathlineto{\pgfqpoint{4.669879in}{0.258715in}}%
\pgfpathlineto{\pgfqpoint{4.681883in}{0.270485in}}%
\pgfpathlineto{\pgfqpoint{4.683019in}{0.279508in}}%
\pgfpathlineto{\pgfqpoint{4.690875in}{0.299484in}}%
\pgfpathlineto{\pgfqpoint{4.711536in}{0.302897in}}%
\pgfpathlineto{\pgfqpoint{4.715491in}{0.314590in}}%
\pgfpathlineto{\pgfqpoint{4.724600in}{0.310205in}}%
\pgfpathlineto{\pgfqpoint{4.743110in}{0.277388in}}%
\pgfpathlineto{\pgfqpoint{4.743344in}{0.266674in}}%
\pgfpathlineto{\pgfqpoint{4.738112in}{0.257400in}}%
\pgfpathlineto{\pgfqpoint{4.744750in}{0.236108in}}%
\pgfpathlineto{\pgfqpoint{4.740111in}{0.232676in}}%
\pgfpathlineto{\pgfqpoint{4.724311in}{0.233674in}}%
\pgfpathlineto{\pgfqpoint{4.705311in}{0.240024in}}%
\pgfpathlineto{\pgfqpoint{4.689342in}{0.242646in}}%
\pgfpathlineto{\pgfqpoint{4.685076in}{0.238921in}}%
\pgfpathlineto{\pgfqpoint{4.673289in}{0.237569in}}%
\pgfpathclose%
\pgfusepath{fill}%
\end{pgfscope}%
\begin{pgfscope}%
\pgfpathrectangle{\pgfqpoint{3.625000in}{0.100000in}}{\pgfqpoint{2.989028in}{1.913466in}}%
\pgfusepath{clip}%
\pgfsetbuttcap%
\pgfsetmiterjoin%
\definecolor{currentfill}{rgb}{0.421223,0.769089,0.646751}%
\pgfsetfillcolor{currentfill}%
\pgfsetlinewidth{0.000000pt}%
\definecolor{currentstroke}{rgb}{0.000000,0.000000,0.000000}%
\pgfsetstrokecolor{currentstroke}%
\pgfsetstrokeopacity{0.000000}%
\pgfsetdash{}{0pt}%
\pgfpathmoveto{\pgfqpoint{4.335222in}{1.900931in}}%
\pgfpathlineto{\pgfqpoint{4.319722in}{1.839512in}}%
\pgfpathlineto{\pgfqpoint{4.308700in}{1.795504in}}%
\pgfpathlineto{\pgfqpoint{4.295595in}{1.742564in}}%
\pgfpathlineto{\pgfqpoint{4.293022in}{1.729516in}}%
\pgfpathlineto{\pgfqpoint{4.294703in}{1.717547in}}%
\pgfpathlineto{\pgfqpoint{4.292475in}{1.706523in}}%
\pgfpathlineto{\pgfqpoint{4.200998in}{1.730394in}}%
\pgfpathlineto{\pgfqpoint{4.192734in}{1.727622in}}%
\pgfpathlineto{\pgfqpoint{4.185663in}{1.730028in}}%
\pgfpathlineto{\pgfqpoint{4.152921in}{1.730814in}}%
\pgfpathlineto{\pgfqpoint{4.141848in}{1.727692in}}%
\pgfpathlineto{\pgfqpoint{4.130851in}{1.728747in}}%
\pgfpathlineto{\pgfqpoint{4.126190in}{1.733634in}}%
\pgfpathlineto{\pgfqpoint{4.096877in}{1.732568in}}%
\pgfpathlineto{\pgfqpoint{4.092248in}{1.740554in}}%
\pgfpathlineto{\pgfqpoint{4.083460in}{1.744762in}}%
\pgfpathlineto{\pgfqpoint{4.070673in}{1.747325in}}%
\pgfpathlineto{\pgfqpoint{4.048603in}{1.743571in}}%
\pgfpathlineto{\pgfqpoint{4.040456in}{1.747249in}}%
\pgfpathlineto{\pgfqpoint{4.027895in}{1.757031in}}%
\pgfpathlineto{\pgfqpoint{4.031670in}{1.776202in}}%
\pgfpathlineto{\pgfqpoint{4.030347in}{1.785992in}}%
\pgfpathlineto{\pgfqpoint{4.020391in}{1.796386in}}%
\pgfpathlineto{\pgfqpoint{4.014015in}{1.795741in}}%
\pgfpathlineto{\pgfqpoint{4.009440in}{1.806255in}}%
\pgfpathlineto{\pgfqpoint{3.998604in}{1.810504in}}%
\pgfpathlineto{\pgfqpoint{3.990727in}{1.809963in}}%
\pgfpathlineto{\pgfqpoint{3.988179in}{1.820760in}}%
\pgfpathlineto{\pgfqpoint{3.996026in}{1.819611in}}%
\pgfpathlineto{\pgfqpoint{3.995402in}{1.834619in}}%
\pgfpathlineto{\pgfqpoint{3.991982in}{1.843530in}}%
\pgfpathlineto{\pgfqpoint{3.993490in}{1.854935in}}%
\pgfpathlineto{\pgfqpoint{3.997595in}{1.863514in}}%
\pgfpathlineto{\pgfqpoint{3.997364in}{1.879792in}}%
\pgfpathlineto{\pgfqpoint{3.995172in}{1.885699in}}%
\pgfpathlineto{\pgfqpoint{3.998957in}{1.904648in}}%
\pgfpathlineto{\pgfqpoint{3.997865in}{1.916860in}}%
\pgfpathlineto{\pgfqpoint{3.994081in}{1.922716in}}%
\pgfpathlineto{\pgfqpoint{3.994658in}{1.941861in}}%
\pgfpathlineto{\pgfqpoint{4.000162in}{1.955941in}}%
\pgfpathlineto{\pgfqpoint{4.006149in}{1.952495in}}%
\pgfpathlineto{\pgfqpoint{4.025972in}{1.932018in}}%
\pgfpathlineto{\pgfqpoint{4.049993in}{1.920812in}}%
\pgfpathlineto{\pgfqpoint{4.062292in}{1.919486in}}%
\pgfpathlineto{\pgfqpoint{4.073880in}{1.914715in}}%
\pgfpathlineto{\pgfqpoint{4.077168in}{1.898661in}}%
\pgfpathlineto{\pgfqpoint{4.067014in}{1.895620in}}%
\pgfpathlineto{\pgfqpoint{4.060285in}{1.882967in}}%
\pgfpathlineto{\pgfqpoint{4.068185in}{1.883670in}}%
\pgfpathlineto{\pgfqpoint{4.071372in}{1.889349in}}%
\pgfpathlineto{\pgfqpoint{4.082514in}{1.896451in}}%
\pgfpathlineto{\pgfqpoint{4.081939in}{1.885905in}}%
\pgfpathlineto{\pgfqpoint{4.074499in}{1.884204in}}%
\pgfpathlineto{\pgfqpoint{4.075734in}{1.870648in}}%
\pgfpathlineto{\pgfqpoint{4.066674in}{1.857069in}}%
\pgfpathlineto{\pgfqpoint{4.060876in}{1.863526in}}%
\pgfpathlineto{\pgfqpoint{4.043645in}{1.860015in}}%
\pgfpathlineto{\pgfqpoint{4.042738in}{1.852116in}}%
\pgfpathlineto{\pgfqpoint{4.048678in}{1.847303in}}%
\pgfpathlineto{\pgfqpoint{4.057741in}{1.846830in}}%
\pgfpathlineto{\pgfqpoint{4.070271in}{1.857114in}}%
\pgfpathlineto{\pgfqpoint{4.080189in}{1.857968in}}%
\pgfpathlineto{\pgfqpoint{4.080569in}{1.869352in}}%
\pgfpathlineto{\pgfqpoint{4.085680in}{1.886072in}}%
\pgfpathlineto{\pgfqpoint{4.097793in}{1.898506in}}%
\pgfpathlineto{\pgfqpoint{4.093754in}{1.908138in}}%
\pgfpathlineto{\pgfqpoint{4.096513in}{1.918507in}}%
\pgfpathlineto{\pgfqpoint{4.091139in}{1.928842in}}%
\pgfpathlineto{\pgfqpoint{4.098798in}{1.942057in}}%
\pgfpathlineto{\pgfqpoint{4.099945in}{1.950008in}}%
\pgfpathlineto{\pgfqpoint{4.093296in}{1.955234in}}%
\pgfpathlineto{\pgfqpoint{4.094403in}{1.968600in}}%
\pgfpathlineto{\pgfqpoint{4.174072in}{1.944395in}}%
\pgfpathlineto{\pgfqpoint{4.258673in}{1.920671in}}%
\pgfpathlineto{\pgfqpoint{4.335222in}{1.900931in}}%
\pgfpathclose%
\pgfusepath{fill}%
\end{pgfscope}%
\begin{pgfscope}%
\pgfpathrectangle{\pgfqpoint{3.625000in}{0.100000in}}{\pgfqpoint{2.989028in}{1.913466in}}%
\pgfusepath{clip}%
\pgfsetbuttcap%
\pgfsetmiterjoin%
\definecolor{currentfill}{rgb}{0.421223,0.769089,0.646751}%
\pgfsetfillcolor{currentfill}%
\pgfsetlinewidth{0.000000pt}%
\definecolor{currentstroke}{rgb}{0.000000,0.000000,0.000000}%
\pgfsetstrokecolor{currentstroke}%
\pgfsetstrokeopacity{0.000000}%
\pgfsetdash{}{0pt}%
\pgfpathmoveto{\pgfqpoint{4.081297in}{1.922359in}}%
\pgfpathlineto{\pgfqpoint{4.085220in}{1.907235in}}%
\pgfpathlineto{\pgfqpoint{4.091372in}{1.902211in}}%
\pgfpathlineto{\pgfqpoint{4.086463in}{1.895816in}}%
\pgfpathlineto{\pgfqpoint{4.081679in}{1.905186in}}%
\pgfpathlineto{\pgfqpoint{4.081297in}{1.922359in}}%
\pgfpathclose%
\pgfusepath{fill}%
\end{pgfscope}%
\begin{pgfscope}%
\pgfpathrectangle{\pgfqpoint{3.625000in}{0.100000in}}{\pgfqpoint{2.989028in}{1.913466in}}%
\pgfusepath{clip}%
\pgfsetbuttcap%
\pgfsetmiterjoin%
\definecolor{currentfill}{rgb}{0.484890,0.794002,0.645829}%
\pgfsetfillcolor{currentfill}%
\pgfsetlinewidth{0.000000pt}%
\definecolor{currentstroke}{rgb}{0.000000,0.000000,0.000000}%
\pgfsetstrokecolor{currentstroke}%
\pgfsetstrokeopacity{0.000000}%
\pgfsetdash{}{0pt}%
\pgfpathmoveto{\pgfqpoint{4.376305in}{1.891004in}}%
\pgfpathlineto{\pgfqpoint{4.461467in}{1.871893in}}%
\pgfpathlineto{\pgfqpoint{4.541670in}{1.855681in}}%
\pgfpathlineto{\pgfqpoint{4.603378in}{1.844375in}}%
\pgfpathlineto{\pgfqpoint{4.657175in}{1.835337in}}%
\pgfpathlineto{\pgfqpoint{4.711092in}{1.827040in}}%
\pgfpathlineto{\pgfqpoint{4.757007in}{1.820572in}}%
\pgfpathlineto{\pgfqpoint{4.802995in}{1.814641in}}%
\pgfpathlineto{\pgfqpoint{4.849048in}{1.809248in}}%
\pgfpathlineto{\pgfqpoint{4.892448in}{1.804664in}}%
\pgfpathlineto{\pgfqpoint{4.886479in}{1.738522in}}%
\pgfpathlineto{\pgfqpoint{4.877595in}{1.649140in}}%
\pgfpathlineto{\pgfqpoint{4.872936in}{1.603169in}}%
\pgfpathlineto{\pgfqpoint{4.866234in}{1.541215in}}%
\pgfpathlineto{\pgfqpoint{4.818794in}{1.546395in}}%
\pgfpathlineto{\pgfqpoint{4.764479in}{1.552561in}}%
\pgfpathlineto{\pgfqpoint{4.689064in}{1.562882in}}%
\pgfpathlineto{\pgfqpoint{4.655381in}{1.567687in}}%
\pgfpathlineto{\pgfqpoint{4.572396in}{1.580713in}}%
\pgfpathlineto{\pgfqpoint{4.543829in}{1.585939in}}%
\pgfpathlineto{\pgfqpoint{4.537875in}{1.552087in}}%
\pgfpathlineto{\pgfqpoint{4.534620in}{1.554502in}}%
\pgfpathlineto{\pgfqpoint{4.528502in}{1.570788in}}%
\pgfpathlineto{\pgfqpoint{4.520639in}{1.569785in}}%
\pgfpathlineto{\pgfqpoint{4.515011in}{1.561552in}}%
\pgfpathlineto{\pgfqpoint{4.503331in}{1.560451in}}%
\pgfpathlineto{\pgfqpoint{4.500984in}{1.564390in}}%
\pgfpathlineto{\pgfqpoint{4.490043in}{1.563913in}}%
\pgfpathlineto{\pgfqpoint{4.484543in}{1.567727in}}%
\pgfpathlineto{\pgfqpoint{4.476852in}{1.561826in}}%
\pgfpathlineto{\pgfqpoint{4.458172in}{1.567136in}}%
\pgfpathlineto{\pgfqpoint{4.450020in}{1.564285in}}%
\pgfpathlineto{\pgfqpoint{4.445634in}{1.576938in}}%
\pgfpathlineto{\pgfqpoint{4.445429in}{1.588978in}}%
\pgfpathlineto{\pgfqpoint{4.432552in}{1.597637in}}%
\pgfpathlineto{\pgfqpoint{4.434652in}{1.610116in}}%
\pgfpathlineto{\pgfqpoint{4.425483in}{1.630729in}}%
\pgfpathlineto{\pgfqpoint{4.426101in}{1.648251in}}%
\pgfpathlineto{\pgfqpoint{4.418220in}{1.656827in}}%
\pgfpathlineto{\pgfqpoint{4.410827in}{1.649240in}}%
\pgfpathlineto{\pgfqpoint{4.400763in}{1.645136in}}%
\pgfpathlineto{\pgfqpoint{4.393102in}{1.653594in}}%
\pgfpathlineto{\pgfqpoint{4.394348in}{1.667101in}}%
\pgfpathlineto{\pgfqpoint{4.403535in}{1.672422in}}%
\pgfpathlineto{\pgfqpoint{4.401093in}{1.680980in}}%
\pgfpathlineto{\pgfqpoint{4.417063in}{1.722399in}}%
\pgfpathlineto{\pgfqpoint{4.404470in}{1.723437in}}%
\pgfpathlineto{\pgfqpoint{4.403176in}{1.730870in}}%
\pgfpathlineto{\pgfqpoint{4.394003in}{1.737279in}}%
\pgfpathlineto{\pgfqpoint{4.394589in}{1.744445in}}%
\pgfpathlineto{\pgfqpoint{4.389750in}{1.750011in}}%
\pgfpathlineto{\pgfqpoint{4.381873in}{1.770222in}}%
\pgfpathlineto{\pgfqpoint{4.374675in}{1.774343in}}%
\pgfpathlineto{\pgfqpoint{4.369521in}{1.791786in}}%
\pgfpathlineto{\pgfqpoint{4.370737in}{1.803380in}}%
\pgfpathlineto{\pgfqpoint{4.361121in}{1.824733in}}%
\pgfpathlineto{\pgfqpoint{4.376305in}{1.891004in}}%
\pgfpathclose%
\pgfusepath{fill}%
\end{pgfscope}%
\begin{pgfscope}%
\pgfpathrectangle{\pgfqpoint{3.625000in}{0.100000in}}{\pgfqpoint{2.989028in}{1.913466in}}%
\pgfusepath{clip}%
\pgfsetbuttcap%
\pgfsetmiterjoin%
\definecolor{currentfill}{rgb}{0.684198,0.872203,0.640369}%
\pgfsetfillcolor{currentfill}%
\pgfsetlinewidth{0.000000pt}%
\definecolor{currentstroke}{rgb}{0.000000,0.000000,0.000000}%
\pgfsetstrokecolor{currentstroke}%
\pgfsetstrokeopacity{0.000000}%
\pgfsetdash{}{0pt}%
\pgfpathmoveto{\pgfqpoint{6.432175in}{1.550163in}}%
\pgfpathlineto{\pgfqpoint{6.430471in}{1.557411in}}%
\pgfpathlineto{\pgfqpoint{6.420998in}{1.563760in}}%
\pgfpathlineto{\pgfqpoint{6.395964in}{1.645645in}}%
\pgfpathlineto{\pgfqpoint{6.382471in}{1.685161in}}%
\pgfpathlineto{\pgfqpoint{6.393319in}{1.696546in}}%
\pgfpathlineto{\pgfqpoint{6.403864in}{1.724407in}}%
\pgfpathlineto{\pgfqpoint{6.409123in}{1.733311in}}%
\pgfpathlineto{\pgfqpoint{6.405397in}{1.737050in}}%
\pgfpathlineto{\pgfqpoint{6.404149in}{1.761782in}}%
\pgfpathlineto{\pgfqpoint{6.408854in}{1.769364in}}%
\pgfpathlineto{\pgfqpoint{6.406792in}{1.786986in}}%
\pgfpathlineto{\pgfqpoint{6.425570in}{1.844804in}}%
\pgfpathlineto{\pgfqpoint{6.434031in}{1.845120in}}%
\pgfpathlineto{\pgfqpoint{6.437543in}{1.834792in}}%
\pgfpathlineto{\pgfqpoint{6.445062in}{1.831826in}}%
\pgfpathlineto{\pgfqpoint{6.459266in}{1.843974in}}%
\pgfpathlineto{\pgfqpoint{6.470403in}{1.851149in}}%
\pgfpathlineto{\pgfqpoint{6.494943in}{1.838514in}}%
\pgfpathlineto{\pgfqpoint{6.517113in}{1.768360in}}%
\pgfpathlineto{\pgfqpoint{6.521328in}{1.751101in}}%
\pgfpathlineto{\pgfqpoint{6.531039in}{1.749076in}}%
\pgfpathlineto{\pgfqpoint{6.544209in}{1.737354in}}%
\pgfpathlineto{\pgfqpoint{6.543463in}{1.730524in}}%
\pgfpathlineto{\pgfqpoint{6.552435in}{1.722428in}}%
\pgfpathlineto{\pgfqpoint{6.559723in}{1.727073in}}%
\pgfpathlineto{\pgfqpoint{6.574974in}{1.709287in}}%
\pgfpathlineto{\pgfqpoint{6.568158in}{1.695053in}}%
\pgfpathlineto{\pgfqpoint{6.559018in}{1.694771in}}%
\pgfpathlineto{\pgfqpoint{6.551685in}{1.682065in}}%
\pgfpathlineto{\pgfqpoint{6.542802in}{1.680263in}}%
\pgfpathlineto{\pgfqpoint{6.536280in}{1.673497in}}%
\pgfpathlineto{\pgfqpoint{6.516803in}{1.666241in}}%
\pgfpathlineto{\pgfqpoint{6.505014in}{1.654548in}}%
\pgfpathlineto{\pgfqpoint{6.499025in}{1.662941in}}%
\pgfpathlineto{\pgfqpoint{6.493584in}{1.656926in}}%
\pgfpathlineto{\pgfqpoint{6.495074in}{1.632607in}}%
\pgfpathlineto{\pgfqpoint{6.490765in}{1.622974in}}%
\pgfpathlineto{\pgfqpoint{6.481367in}{1.625605in}}%
\pgfpathlineto{\pgfqpoint{6.479850in}{1.615730in}}%
\pgfpathlineto{\pgfqpoint{6.475789in}{1.611669in}}%
\pgfpathlineto{\pgfqpoint{6.467662in}{1.612659in}}%
\pgfpathlineto{\pgfqpoint{6.468174in}{1.603425in}}%
\pgfpathlineto{\pgfqpoint{6.455855in}{1.606048in}}%
\pgfpathlineto{\pgfqpoint{6.449128in}{1.593255in}}%
\pgfpathlineto{\pgfqpoint{6.451668in}{1.586561in}}%
\pgfpathlineto{\pgfqpoint{6.447666in}{1.575429in}}%
\pgfpathlineto{\pgfqpoint{6.441367in}{1.567240in}}%
\pgfpathlineto{\pgfqpoint{6.439781in}{1.550146in}}%
\pgfpathlineto{\pgfqpoint{6.432175in}{1.550163in}}%
\pgfpathclose%
\pgfusepath{fill}%
\end{pgfscope}%
\begin{pgfscope}%
\pgfpathrectangle{\pgfqpoint{3.625000in}{0.100000in}}{\pgfqpoint{2.989028in}{1.913466in}}%
\pgfusepath{clip}%
\pgfsetbuttcap%
\pgfsetmiterjoin%
\definecolor{currentfill}{rgb}{0.684198,0.872203,0.640369}%
\pgfsetfillcolor{currentfill}%
\pgfsetlinewidth{0.000000pt}%
\definecolor{currentstroke}{rgb}{0.000000,0.000000,0.000000}%
\pgfsetstrokecolor{currentstroke}%
\pgfsetstrokeopacity{0.000000}%
\pgfsetdash{}{0pt}%
\pgfpathmoveto{\pgfqpoint{6.520301in}{1.661247in}}%
\pgfpathlineto{\pgfqpoint{6.525863in}{1.667073in}}%
\pgfpathlineto{\pgfqpoint{6.531158in}{1.661597in}}%
\pgfpathlineto{\pgfqpoint{6.521647in}{1.654342in}}%
\pgfpathlineto{\pgfqpoint{6.520301in}{1.661247in}}%
\pgfpathclose%
\pgfusepath{fill}%
\end{pgfscope}%
\begin{pgfscope}%
\pgfpathrectangle{\pgfqpoint{3.625000in}{0.100000in}}{\pgfqpoint{2.989028in}{1.913466in}}%
\pgfusepath{clip}%
\pgfsetbuttcap%
\pgfsetmiterjoin%
\definecolor{currentfill}{rgb}{0.240062,0.476355,0.714187}%
\pgfsetfillcolor{currentfill}%
\pgfsetlinewidth{0.000000pt}%
\definecolor{currentstroke}{rgb}{0.000000,0.000000,0.000000}%
\pgfsetstrokecolor{currentstroke}%
\pgfsetstrokeopacity{0.000000}%
\pgfsetdash{}{0pt}%
\pgfpathmoveto{\pgfqpoint{4.872936in}{1.603169in}}%
\pgfpathlineto{\pgfqpoint{4.877595in}{1.649140in}}%
\pgfpathlineto{\pgfqpoint{4.886479in}{1.738522in}}%
\pgfpathlineto{\pgfqpoint{4.892448in}{1.804664in}}%
\pgfpathlineto{\pgfqpoint{4.941329in}{1.800078in}}%
\pgfpathlineto{\pgfqpoint{5.003862in}{1.795100in}}%
\pgfpathlineto{\pgfqpoint{5.061021in}{1.791418in}}%
\pgfpathlineto{\pgfqpoint{5.112777in}{1.788799in}}%
\pgfpathlineto{\pgfqpoint{5.190014in}{1.786150in}}%
\pgfpathlineto{\pgfqpoint{5.195285in}{1.764381in}}%
\pgfpathlineto{\pgfqpoint{5.193398in}{1.753978in}}%
\pgfpathlineto{\pgfqpoint{5.192777in}{1.732596in}}%
\pgfpathlineto{\pgfqpoint{5.196363in}{1.716582in}}%
\pgfpathlineto{\pgfqpoint{5.204579in}{1.692956in}}%
\pgfpathlineto{\pgfqpoint{5.204560in}{1.663212in}}%
\pgfpathlineto{\pgfqpoint{5.206080in}{1.628658in}}%
\pgfpathlineto{\pgfqpoint{5.208167in}{1.619346in}}%
\pgfpathlineto{\pgfqpoint{5.214274in}{1.609122in}}%
\pgfpathlineto{\pgfqpoint{5.216288in}{1.593192in}}%
\pgfpathlineto{\pgfqpoint{5.215419in}{1.582557in}}%
\pgfpathlineto{\pgfqpoint{5.150668in}{1.583959in}}%
\pgfpathlineto{\pgfqpoint{5.103564in}{1.586347in}}%
\pgfpathlineto{\pgfqpoint{5.034511in}{1.590023in}}%
\pgfpathlineto{\pgfqpoint{4.966415in}{1.594860in}}%
\pgfpathlineto{\pgfqpoint{4.921059in}{1.598519in}}%
\pgfpathlineto{\pgfqpoint{4.872936in}{1.603169in}}%
\pgfpathclose%
\pgfusepath{fill}%
\end{pgfscope}%
\begin{pgfscope}%
\pgfpathrectangle{\pgfqpoint{3.625000in}{0.100000in}}{\pgfqpoint{2.989028in}{1.913466in}}%
\pgfusepath{clip}%
\pgfsetbuttcap%
\pgfsetmiterjoin%
\definecolor{currentfill}{rgb}{0.665283,0.864591,0.643214}%
\pgfsetfillcolor{currentfill}%
\pgfsetlinewidth{0.000000pt}%
\definecolor{currentstroke}{rgb}{0.000000,0.000000,0.000000}%
\pgfsetstrokecolor{currentstroke}%
\pgfsetstrokeopacity{0.000000}%
\pgfsetdash{}{0pt}%
\pgfpathmoveto{\pgfqpoint{4.853363in}{1.410780in}}%
\pgfpathlineto{\pgfqpoint{4.860839in}{1.487840in}}%
\pgfpathlineto{\pgfqpoint{4.866234in}{1.541215in}}%
\pgfpathlineto{\pgfqpoint{4.872936in}{1.603169in}}%
\pgfpathlineto{\pgfqpoint{4.921059in}{1.598519in}}%
\pgfpathlineto{\pgfqpoint{4.966415in}{1.594860in}}%
\pgfpathlineto{\pgfqpoint{5.034511in}{1.590023in}}%
\pgfpathlineto{\pgfqpoint{5.103564in}{1.586347in}}%
\pgfpathlineto{\pgfqpoint{5.150668in}{1.583959in}}%
\pgfpathlineto{\pgfqpoint{5.215419in}{1.582557in}}%
\pgfpathlineto{\pgfqpoint{5.211034in}{1.569764in}}%
\pgfpathlineto{\pgfqpoint{5.202284in}{1.559721in}}%
\pgfpathlineto{\pgfqpoint{5.208985in}{1.548158in}}%
\pgfpathlineto{\pgfqpoint{5.216374in}{1.545688in}}%
\pgfpathlineto{\pgfqpoint{5.219880in}{1.539046in}}%
\pgfpathlineto{\pgfqpoint{5.219075in}{1.490552in}}%
\pgfpathlineto{\pgfqpoint{5.217728in}{1.422304in}}%
\pgfpathlineto{\pgfqpoint{5.212744in}{1.404379in}}%
\pgfpathlineto{\pgfqpoint{5.217201in}{1.394464in}}%
\pgfpathlineto{\pgfqpoint{5.208252in}{1.373157in}}%
\pgfpathlineto{\pgfqpoint{5.217680in}{1.355981in}}%
\pgfpathlineto{\pgfqpoint{5.209670in}{1.357295in}}%
\pgfpathlineto{\pgfqpoint{5.204205in}{1.367982in}}%
\pgfpathlineto{\pgfqpoint{5.184700in}{1.375307in}}%
\pgfpathlineto{\pgfqpoint{5.172400in}{1.381750in}}%
\pgfpathlineto{\pgfqpoint{5.151759in}{1.382286in}}%
\pgfpathlineto{\pgfqpoint{5.144594in}{1.376416in}}%
\pgfpathlineto{\pgfqpoint{5.121231in}{1.387975in}}%
\pgfpathlineto{\pgfqpoint{5.119439in}{1.391630in}}%
\pgfpathlineto{\pgfqpoint{5.037902in}{1.395489in}}%
\pgfpathlineto{\pgfqpoint{4.988365in}{1.398329in}}%
\pgfpathlineto{\pgfqpoint{4.947447in}{1.401533in}}%
\pgfpathlineto{\pgfqpoint{4.879836in}{1.407921in}}%
\pgfpathlineto{\pgfqpoint{4.853363in}{1.410780in}}%
\pgfpathclose%
\pgfusepath{fill}%
\end{pgfscope}%
\begin{pgfscope}%
\pgfpathrectangle{\pgfqpoint{3.625000in}{0.100000in}}{\pgfqpoint{2.989028in}{1.913466in}}%
\pgfusepath{clip}%
\pgfsetbuttcap%
\pgfsetmiterjoin%
\definecolor{currentfill}{rgb}{0.294195,0.406228,0.680969}%
\pgfsetfillcolor{currentfill}%
\pgfsetlinewidth{0.000000pt}%
\definecolor{currentstroke}{rgb}{0.000000,0.000000,0.000000}%
\pgfsetstrokecolor{currentstroke}%
\pgfsetstrokeopacity{0.000000}%
\pgfsetdash{}{0pt}%
\pgfpathmoveto{\pgfqpoint{4.840540in}{1.280295in}}%
\pgfpathlineto{\pgfqpoint{4.797072in}{1.284265in}}%
\pgfpathlineto{\pgfqpoint{4.702321in}{1.295959in}}%
\pgfpathlineto{\pgfqpoint{4.650778in}{1.303357in}}%
\pgfpathlineto{\pgfqpoint{4.595495in}{1.311293in}}%
\pgfpathlineto{\pgfqpoint{4.548928in}{1.318744in}}%
\pgfpathlineto{\pgfqpoint{4.497815in}{1.327496in}}%
\pgfpathlineto{\pgfqpoint{4.509436in}{1.391963in}}%
\pgfpathlineto{\pgfqpoint{4.521208in}{1.458068in}}%
\pgfpathlineto{\pgfqpoint{4.537875in}{1.552087in}}%
\pgfpathlineto{\pgfqpoint{4.543829in}{1.585939in}}%
\pgfpathlineto{\pgfqpoint{4.572396in}{1.580713in}}%
\pgfpathlineto{\pgfqpoint{4.655381in}{1.567687in}}%
\pgfpathlineto{\pgfqpoint{4.689064in}{1.562882in}}%
\pgfpathlineto{\pgfqpoint{4.764479in}{1.552561in}}%
\pgfpathlineto{\pgfqpoint{4.818794in}{1.546395in}}%
\pgfpathlineto{\pgfqpoint{4.866234in}{1.541215in}}%
\pgfpathlineto{\pgfqpoint{4.860839in}{1.487840in}}%
\pgfpathlineto{\pgfqpoint{4.853363in}{1.410780in}}%
\pgfpathlineto{\pgfqpoint{4.846947in}{1.345289in}}%
\pgfpathlineto{\pgfqpoint{4.840540in}{1.280295in}}%
\pgfpathclose%
\pgfusepath{fill}%
\end{pgfscope}%
\begin{pgfscope}%
\pgfpathrectangle{\pgfqpoint{3.625000in}{0.100000in}}{\pgfqpoint{2.989028in}{1.913466in}}%
\pgfusepath{clip}%
\pgfsetbuttcap%
\pgfsetmiterjoin%
\definecolor{currentfill}{rgb}{0.978854,0.991542,0.716032}%
\pgfsetfillcolor{currentfill}%
\pgfsetlinewidth{0.000000pt}%
\definecolor{currentstroke}{rgb}{0.000000,0.000000,0.000000}%
\pgfsetstrokecolor{currentstroke}%
\pgfsetstrokeopacity{0.000000}%
\pgfsetdash{}{0pt}%
\pgfpathmoveto{\pgfqpoint{5.635291in}{1.369221in}}%
\pgfpathlineto{\pgfqpoint{5.580206in}{1.365308in}}%
\pgfpathlineto{\pgfqpoint{5.498100in}{1.361809in}}%
\pgfpathlineto{\pgfqpoint{5.494980in}{1.370103in}}%
\pgfpathlineto{\pgfqpoint{5.476772in}{1.376303in}}%
\pgfpathlineto{\pgfqpoint{5.472756in}{1.388016in}}%
\pgfpathlineto{\pgfqpoint{5.471075in}{1.402507in}}%
\pgfpathlineto{\pgfqpoint{5.475178in}{1.409931in}}%
\pgfpathlineto{\pgfqpoint{5.468697in}{1.417056in}}%
\pgfpathlineto{\pgfqpoint{5.467134in}{1.425554in}}%
\pgfpathlineto{\pgfqpoint{5.465045in}{1.444352in}}%
\pgfpathlineto{\pgfqpoint{5.458835in}{1.454564in}}%
\pgfpathlineto{\pgfqpoint{5.447805in}{1.460295in}}%
\pgfpathlineto{\pgfqpoint{5.435803in}{1.469752in}}%
\pgfpathlineto{\pgfqpoint{5.429599in}{1.480942in}}%
\pgfpathlineto{\pgfqpoint{5.418466in}{1.485448in}}%
\pgfpathlineto{\pgfqpoint{5.411922in}{1.492782in}}%
\pgfpathlineto{\pgfqpoint{5.403994in}{1.494027in}}%
\pgfpathlineto{\pgfqpoint{5.389840in}{1.504916in}}%
\pgfpathlineto{\pgfqpoint{5.392139in}{1.517447in}}%
\pgfpathlineto{\pgfqpoint{5.391696in}{1.541277in}}%
\pgfpathlineto{\pgfqpoint{5.396061in}{1.547836in}}%
\pgfpathlineto{\pgfqpoint{5.392139in}{1.557741in}}%
\pgfpathlineto{\pgfqpoint{5.385230in}{1.559657in}}%
\pgfpathlineto{\pgfqpoint{5.385801in}{1.568355in}}%
\pgfpathlineto{\pgfqpoint{5.394373in}{1.582098in}}%
\pgfpathlineto{\pgfqpoint{5.411348in}{1.592965in}}%
\pgfpathlineto{\pgfqpoint{5.410290in}{1.631611in}}%
\pgfpathlineto{\pgfqpoint{5.418784in}{1.637416in}}%
\pgfpathlineto{\pgfqpoint{5.426822in}{1.633549in}}%
\pgfpathlineto{\pgfqpoint{5.443197in}{1.639211in}}%
\pgfpathlineto{\pgfqpoint{5.474004in}{1.653439in}}%
\pgfpathlineto{\pgfqpoint{5.478017in}{1.649033in}}%
\pgfpathlineto{\pgfqpoint{5.472182in}{1.629079in}}%
\pgfpathlineto{\pgfqpoint{5.480863in}{1.633458in}}%
\pgfpathlineto{\pgfqpoint{5.495731in}{1.629085in}}%
\pgfpathlineto{\pgfqpoint{5.504867in}{1.625436in}}%
\pgfpathlineto{\pgfqpoint{5.509989in}{1.614740in}}%
\pgfpathlineto{\pgfqpoint{5.556860in}{1.604646in}}%
\pgfpathlineto{\pgfqpoint{5.570871in}{1.597722in}}%
\pgfpathlineto{\pgfqpoint{5.585116in}{1.597803in}}%
\pgfpathlineto{\pgfqpoint{5.599752in}{1.594947in}}%
\pgfpathlineto{\pgfqpoint{5.609251in}{1.585189in}}%
\pgfpathlineto{\pgfqpoint{5.618328in}{1.580368in}}%
\pgfpathlineto{\pgfqpoint{5.619996in}{1.566462in}}%
\pgfpathlineto{\pgfqpoint{5.617314in}{1.557730in}}%
\pgfpathlineto{\pgfqpoint{5.627430in}{1.557084in}}%
\pgfpathlineto{\pgfqpoint{5.624022in}{1.546957in}}%
\pgfpathlineto{\pgfqpoint{5.627264in}{1.543360in}}%
\pgfpathlineto{\pgfqpoint{5.630488in}{1.533800in}}%
\pgfpathlineto{\pgfqpoint{5.620629in}{1.528739in}}%
\pgfpathlineto{\pgfqpoint{5.614905in}{1.514637in}}%
\pgfpathlineto{\pgfqpoint{5.613109in}{1.504665in}}%
\pgfpathlineto{\pgfqpoint{5.618600in}{1.502958in}}%
\pgfpathlineto{\pgfqpoint{5.625632in}{1.510436in}}%
\pgfpathlineto{\pgfqpoint{5.631610in}{1.523394in}}%
\pgfpathlineto{\pgfqpoint{5.639698in}{1.527895in}}%
\pgfpathlineto{\pgfqpoint{5.645776in}{1.522041in}}%
\pgfpathlineto{\pgfqpoint{5.639793in}{1.504322in}}%
\pgfpathlineto{\pgfqpoint{5.637916in}{1.490567in}}%
\pgfpathlineto{\pgfqpoint{5.639684in}{1.480679in}}%
\pgfpathlineto{\pgfqpoint{5.634128in}{1.475085in}}%
\pgfpathlineto{\pgfqpoint{5.631353in}{1.461416in}}%
\pgfpathlineto{\pgfqpoint{5.633620in}{1.447050in}}%
\pgfpathlineto{\pgfqpoint{5.627064in}{1.425804in}}%
\pgfpathlineto{\pgfqpoint{5.627202in}{1.415188in}}%
\pgfpathlineto{\pgfqpoint{5.632383in}{1.392182in}}%
\pgfpathlineto{\pgfqpoint{5.635741in}{1.388236in}}%
\pgfpathlineto{\pgfqpoint{5.635291in}{1.369221in}}%
\pgfpathclose%
\pgfusepath{fill}%
\end{pgfscope}%
\begin{pgfscope}%
\pgfpathrectangle{\pgfqpoint{3.625000in}{0.100000in}}{\pgfqpoint{2.989028in}{1.913466in}}%
\pgfusepath{clip}%
\pgfsetbuttcap%
\pgfsetmiterjoin%
\definecolor{currentfill}{rgb}{0.978854,0.991542,0.716032}%
\pgfsetfillcolor{currentfill}%
\pgfsetlinewidth{0.000000pt}%
\definecolor{currentstroke}{rgb}{0.000000,0.000000,0.000000}%
\pgfsetstrokecolor{currentstroke}%
\pgfsetstrokeopacity{0.000000}%
\pgfsetdash{}{0pt}%
\pgfpathmoveto{\pgfqpoint{5.655944in}{1.555675in}}%
\pgfpathlineto{\pgfqpoint{5.657173in}{1.546498in}}%
\pgfpathlineto{\pgfqpoint{5.645909in}{1.522320in}}%
\pgfpathlineto{\pgfqpoint{5.640904in}{1.529323in}}%
\pgfpathlineto{\pgfqpoint{5.655944in}{1.555675in}}%
\pgfpathclose%
\pgfusepath{fill}%
\end{pgfscope}%
\begin{pgfscope}%
\pgfpathrectangle{\pgfqpoint{3.625000in}{0.100000in}}{\pgfqpoint{2.989028in}{1.913466in}}%
\pgfusepath{clip}%
\pgfsetbuttcap%
\pgfsetmiterjoin%
\definecolor{currentfill}{rgb}{0.537947,0.814764,0.645060}%
\pgfsetfillcolor{currentfill}%
\pgfsetlinewidth{0.000000pt}%
\definecolor{currentstroke}{rgb}{0.000000,0.000000,0.000000}%
\pgfsetstrokecolor{currentstroke}%
\pgfsetstrokeopacity{0.000000}%
\pgfsetdash{}{0pt}%
\pgfpathmoveto{\pgfqpoint{4.292475in}{1.706523in}}%
\pgfpathlineto{\pgfqpoint{4.294703in}{1.717547in}}%
\pgfpathlineto{\pgfqpoint{4.293022in}{1.729516in}}%
\pgfpathlineto{\pgfqpoint{4.295595in}{1.742564in}}%
\pgfpathlineto{\pgfqpoint{4.308700in}{1.795504in}}%
\pgfpathlineto{\pgfqpoint{4.319722in}{1.839512in}}%
\pgfpathlineto{\pgfqpoint{4.335222in}{1.900931in}}%
\pgfpathlineto{\pgfqpoint{4.376305in}{1.891004in}}%
\pgfpathlineto{\pgfqpoint{4.361121in}{1.824733in}}%
\pgfpathlineto{\pgfqpoint{4.370737in}{1.803380in}}%
\pgfpathlineto{\pgfqpoint{4.369521in}{1.791786in}}%
\pgfpathlineto{\pgfqpoint{4.374675in}{1.774343in}}%
\pgfpathlineto{\pgfqpoint{4.381873in}{1.770222in}}%
\pgfpathlineto{\pgfqpoint{4.389750in}{1.750011in}}%
\pgfpathlineto{\pgfqpoint{4.394589in}{1.744445in}}%
\pgfpathlineto{\pgfqpoint{4.394003in}{1.737279in}}%
\pgfpathlineto{\pgfqpoint{4.403176in}{1.730870in}}%
\pgfpathlineto{\pgfqpoint{4.404470in}{1.723437in}}%
\pgfpathlineto{\pgfqpoint{4.417063in}{1.722399in}}%
\pgfpathlineto{\pgfqpoint{4.401093in}{1.680980in}}%
\pgfpathlineto{\pgfqpoint{4.403535in}{1.672422in}}%
\pgfpathlineto{\pgfqpoint{4.394348in}{1.667101in}}%
\pgfpathlineto{\pgfqpoint{4.393102in}{1.653594in}}%
\pgfpathlineto{\pgfqpoint{4.400763in}{1.645136in}}%
\pgfpathlineto{\pgfqpoint{4.410827in}{1.649240in}}%
\pgfpathlineto{\pgfqpoint{4.418220in}{1.656827in}}%
\pgfpathlineto{\pgfqpoint{4.426101in}{1.648251in}}%
\pgfpathlineto{\pgfqpoint{4.425483in}{1.630729in}}%
\pgfpathlineto{\pgfqpoint{4.434652in}{1.610116in}}%
\pgfpathlineto{\pgfqpoint{4.432552in}{1.597637in}}%
\pgfpathlineto{\pgfqpoint{4.445429in}{1.588978in}}%
\pgfpathlineto{\pgfqpoint{4.445634in}{1.576938in}}%
\pgfpathlineto{\pgfqpoint{4.450020in}{1.564285in}}%
\pgfpathlineto{\pgfqpoint{4.458172in}{1.567136in}}%
\pgfpathlineto{\pgfqpoint{4.476852in}{1.561826in}}%
\pgfpathlineto{\pgfqpoint{4.484543in}{1.567727in}}%
\pgfpathlineto{\pgfqpoint{4.490043in}{1.563913in}}%
\pgfpathlineto{\pgfqpoint{4.500984in}{1.564390in}}%
\pgfpathlineto{\pgfqpoint{4.503331in}{1.560451in}}%
\pgfpathlineto{\pgfqpoint{4.515011in}{1.561552in}}%
\pgfpathlineto{\pgfqpoint{4.520639in}{1.569785in}}%
\pgfpathlineto{\pgfqpoint{4.528502in}{1.570788in}}%
\pgfpathlineto{\pgfqpoint{4.534620in}{1.554502in}}%
\pgfpathlineto{\pgfqpoint{4.537875in}{1.552087in}}%
\pgfpathlineto{\pgfqpoint{4.521208in}{1.458068in}}%
\pgfpathlineto{\pgfqpoint{4.509436in}{1.391963in}}%
\pgfpathlineto{\pgfqpoint{4.416592in}{1.409886in}}%
\pgfpathlineto{\pgfqpoint{4.366442in}{1.419862in}}%
\pgfpathlineto{\pgfqpoint{4.319534in}{1.430177in}}%
\pgfpathlineto{\pgfqpoint{4.275689in}{1.440103in}}%
\pgfpathlineto{\pgfqpoint{4.224946in}{1.452341in}}%
\pgfpathlineto{\pgfqpoint{4.251119in}{1.559712in}}%
\pgfpathlineto{\pgfqpoint{4.252531in}{1.567559in}}%
\pgfpathlineto{\pgfqpoint{4.264170in}{1.588120in}}%
\pgfpathlineto{\pgfqpoint{4.252118in}{1.600478in}}%
\pgfpathlineto{\pgfqpoint{4.254620in}{1.612613in}}%
\pgfpathlineto{\pgfqpoint{4.259208in}{1.615650in}}%
\pgfpathlineto{\pgfqpoint{4.267402in}{1.628155in}}%
\pgfpathlineto{\pgfqpoint{4.279345in}{1.636806in}}%
\pgfpathlineto{\pgfqpoint{4.279997in}{1.643207in}}%
\pgfpathlineto{\pgfqpoint{4.287206in}{1.649603in}}%
\pgfpathlineto{\pgfqpoint{4.293191in}{1.661575in}}%
\pgfpathlineto{\pgfqpoint{4.305469in}{1.674240in}}%
\pgfpathlineto{\pgfqpoint{4.304597in}{1.686749in}}%
\pgfpathlineto{\pgfqpoint{4.296203in}{1.693698in}}%
\pgfpathlineto{\pgfqpoint{4.292475in}{1.706523in}}%
\pgfpathclose%
\pgfusepath{fill}%
\end{pgfscope}%
\begin{pgfscope}%
\pgfpathrectangle{\pgfqpoint{3.625000in}{0.100000in}}{\pgfqpoint{2.989028in}{1.913466in}}%
\pgfusepath{clip}%
\pgfsetbuttcap%
\pgfsetmiterjoin%
\definecolor{currentfill}{rgb}{0.802153,0.920185,0.616378}%
\pgfsetfillcolor{currentfill}%
\pgfsetlinewidth{0.000000pt}%
\definecolor{currentstroke}{rgb}{0.000000,0.000000,0.000000}%
\pgfsetstrokecolor{currentstroke}%
\pgfsetstrokeopacity{0.000000}%
\pgfsetdash{}{0pt}%
\pgfpathmoveto{\pgfqpoint{6.324865in}{1.496989in}}%
\pgfpathlineto{\pgfqpoint{6.321007in}{1.509167in}}%
\pgfpathlineto{\pgfqpoint{6.313848in}{1.546108in}}%
\pgfpathlineto{\pgfqpoint{6.304593in}{1.557473in}}%
\pgfpathlineto{\pgfqpoint{6.296474in}{1.577911in}}%
\pgfpathlineto{\pgfqpoint{6.299140in}{1.593064in}}%
\pgfpathlineto{\pgfqpoint{6.297002in}{1.604233in}}%
\pgfpathlineto{\pgfqpoint{6.290065in}{1.615202in}}%
\pgfpathlineto{\pgfqpoint{6.289780in}{1.627285in}}%
\pgfpathlineto{\pgfqpoint{6.285752in}{1.640318in}}%
\pgfpathlineto{\pgfqpoint{6.321797in}{1.649189in}}%
\pgfpathlineto{\pgfqpoint{6.368616in}{1.661826in}}%
\pgfpathlineto{\pgfqpoint{6.370485in}{1.654580in}}%
\pgfpathlineto{\pgfqpoint{6.367504in}{1.643043in}}%
\pgfpathlineto{\pgfqpoint{6.374502in}{1.633812in}}%
\pgfpathlineto{\pgfqpoint{6.370785in}{1.622115in}}%
\pgfpathlineto{\pgfqpoint{6.356015in}{1.607411in}}%
\pgfpathlineto{\pgfqpoint{6.360075in}{1.596366in}}%
\pgfpathlineto{\pgfqpoint{6.357066in}{1.573021in}}%
\pgfpathlineto{\pgfqpoint{6.352854in}{1.559858in}}%
\pgfpathlineto{\pgfqpoint{6.358228in}{1.522426in}}%
\pgfpathlineto{\pgfqpoint{6.357504in}{1.509256in}}%
\pgfpathlineto{\pgfqpoint{6.362703in}{1.505030in}}%
\pgfpathlineto{\pgfqpoint{6.324865in}{1.496989in}}%
\pgfpathclose%
\pgfusepath{fill}%
\end{pgfscope}%
\begin{pgfscope}%
\pgfpathrectangle{\pgfqpoint{3.625000in}{0.100000in}}{\pgfqpoint{2.989028in}{1.913466in}}%
\pgfusepath{clip}%
\pgfsetbuttcap%
\pgfsetmiterjoin%
\definecolor{currentfill}{rgb}{0.948097,0.979239,0.668051}%
\pgfsetfillcolor{currentfill}%
\pgfsetlinewidth{0.000000pt}%
\definecolor{currentstroke}{rgb}{0.000000,0.000000,0.000000}%
\pgfsetstrokecolor{currentstroke}%
\pgfsetstrokeopacity{0.000000}%
\pgfsetdash{}{0pt}%
\pgfpathmoveto{\pgfqpoint{5.217728in}{1.422304in}}%
\pgfpathlineto{\pgfqpoint{5.219075in}{1.490552in}}%
\pgfpathlineto{\pgfqpoint{5.219880in}{1.539046in}}%
\pgfpathlineto{\pgfqpoint{5.216374in}{1.545688in}}%
\pgfpathlineto{\pgfqpoint{5.208985in}{1.548158in}}%
\pgfpathlineto{\pgfqpoint{5.202284in}{1.559721in}}%
\pgfpathlineto{\pgfqpoint{5.211034in}{1.569764in}}%
\pgfpathlineto{\pgfqpoint{5.215419in}{1.582557in}}%
\pgfpathlineto{\pgfqpoint{5.216288in}{1.593192in}}%
\pgfpathlineto{\pgfqpoint{5.214274in}{1.609122in}}%
\pgfpathlineto{\pgfqpoint{5.208167in}{1.619346in}}%
\pgfpathlineto{\pgfqpoint{5.206080in}{1.628658in}}%
\pgfpathlineto{\pgfqpoint{5.204560in}{1.663212in}}%
\pgfpathlineto{\pgfqpoint{5.204579in}{1.692956in}}%
\pgfpathlineto{\pgfqpoint{5.196363in}{1.716582in}}%
\pgfpathlineto{\pgfqpoint{5.192777in}{1.732596in}}%
\pgfpathlineto{\pgfqpoint{5.193398in}{1.753978in}}%
\pgfpathlineto{\pgfqpoint{5.195285in}{1.764381in}}%
\pgfpathlineto{\pgfqpoint{5.190014in}{1.786150in}}%
\pgfpathlineto{\pgfqpoint{5.225900in}{1.785431in}}%
\pgfpathlineto{\pgfqpoint{5.280411in}{1.784963in}}%
\pgfpathlineto{\pgfqpoint{5.280709in}{1.809705in}}%
\pgfpathlineto{\pgfqpoint{5.294584in}{1.806981in}}%
\pgfpathlineto{\pgfqpoint{5.301233in}{1.776812in}}%
\pgfpathlineto{\pgfqpoint{5.306136in}{1.765965in}}%
\pgfpathlineto{\pgfqpoint{5.318328in}{1.765645in}}%
\pgfpathlineto{\pgfqpoint{5.321056in}{1.761964in}}%
\pgfpathlineto{\pgfqpoint{5.338062in}{1.760334in}}%
\pgfpathlineto{\pgfqpoint{5.340920in}{1.752855in}}%
\pgfpathlineto{\pgfqpoint{5.352640in}{1.754536in}}%
\pgfpathlineto{\pgfqpoint{5.361742in}{1.761531in}}%
\pgfpathlineto{\pgfqpoint{5.377440in}{1.761270in}}%
\pgfpathlineto{\pgfqpoint{5.387153in}{1.755644in}}%
\pgfpathlineto{\pgfqpoint{5.388270in}{1.750368in}}%
\pgfpathlineto{\pgfqpoint{5.397512in}{1.749264in}}%
\pgfpathlineto{\pgfqpoint{5.403543in}{1.734873in}}%
\pgfpathlineto{\pgfqpoint{5.407434in}{1.743721in}}%
\pgfpathlineto{\pgfqpoint{5.418051in}{1.744265in}}%
\pgfpathlineto{\pgfqpoint{5.420736in}{1.737374in}}%
\pgfpathlineto{\pgfqpoint{5.432682in}{1.734228in}}%
\pgfpathlineto{\pgfqpoint{5.439270in}{1.724302in}}%
\pgfpathlineto{\pgfqpoint{5.453880in}{1.727384in}}%
\pgfpathlineto{\pgfqpoint{5.469957in}{1.739565in}}%
\pgfpathlineto{\pgfqpoint{5.475818in}{1.728827in}}%
\pgfpathlineto{\pgfqpoint{5.502198in}{1.731766in}}%
\pgfpathlineto{\pgfqpoint{5.513474in}{1.724397in}}%
\pgfpathlineto{\pgfqpoint{5.520037in}{1.727028in}}%
\pgfpathlineto{\pgfqpoint{5.525300in}{1.722879in}}%
\pgfpathlineto{\pgfqpoint{5.509687in}{1.713035in}}%
\pgfpathlineto{\pgfqpoint{5.487401in}{1.704265in}}%
\pgfpathlineto{\pgfqpoint{5.465329in}{1.686734in}}%
\pgfpathlineto{\pgfqpoint{5.446193in}{1.663680in}}%
\pgfpathlineto{\pgfqpoint{5.431707in}{1.650050in}}%
\pgfpathlineto{\pgfqpoint{5.419018in}{1.640682in}}%
\pgfpathlineto{\pgfqpoint{5.410290in}{1.631611in}}%
\pgfpathlineto{\pgfqpoint{5.411348in}{1.592965in}}%
\pgfpathlineto{\pgfqpoint{5.394373in}{1.582098in}}%
\pgfpathlineto{\pgfqpoint{5.385801in}{1.568355in}}%
\pgfpathlineto{\pgfqpoint{5.385230in}{1.559657in}}%
\pgfpathlineto{\pgfqpoint{5.392139in}{1.557741in}}%
\pgfpathlineto{\pgfqpoint{5.396061in}{1.547836in}}%
\pgfpathlineto{\pgfqpoint{5.391696in}{1.541277in}}%
\pgfpathlineto{\pgfqpoint{5.392139in}{1.517447in}}%
\pgfpathlineto{\pgfqpoint{5.389840in}{1.504916in}}%
\pgfpathlineto{\pgfqpoint{5.403994in}{1.494027in}}%
\pgfpathlineto{\pgfqpoint{5.411922in}{1.492782in}}%
\pgfpathlineto{\pgfqpoint{5.418466in}{1.485448in}}%
\pgfpathlineto{\pgfqpoint{5.429599in}{1.480942in}}%
\pgfpathlineto{\pgfqpoint{5.435803in}{1.469752in}}%
\pgfpathlineto{\pgfqpoint{5.447805in}{1.460295in}}%
\pgfpathlineto{\pgfqpoint{5.458835in}{1.454564in}}%
\pgfpathlineto{\pgfqpoint{5.465045in}{1.444352in}}%
\pgfpathlineto{\pgfqpoint{5.467134in}{1.425554in}}%
\pgfpathlineto{\pgfqpoint{5.408600in}{1.423428in}}%
\pgfpathlineto{\pgfqpoint{5.358705in}{1.422385in}}%
\pgfpathlineto{\pgfqpoint{5.313246in}{1.421717in}}%
\pgfpathlineto{\pgfqpoint{5.265155in}{1.421790in}}%
\pgfpathlineto{\pgfqpoint{5.217728in}{1.422304in}}%
\pgfpathclose%
\pgfusepath{fill}%
\end{pgfscope}%
\begin{pgfscope}%
\pgfpathrectangle{\pgfqpoint{3.625000in}{0.100000in}}{\pgfqpoint{2.989028in}{1.913466in}}%
\pgfusepath{clip}%
\pgfsetbuttcap%
\pgfsetmiterjoin%
\definecolor{currentfill}{rgb}{0.711419,0.883276,0.634833}%
\pgfsetfillcolor{currentfill}%
\pgfsetlinewidth{0.000000pt}%
\definecolor{currentstroke}{rgb}{0.000000,0.000000,0.000000}%
\pgfsetstrokecolor{currentstroke}%
\pgfsetstrokeopacity{0.000000}%
\pgfsetdash{}{0pt}%
\pgfpathmoveto{\pgfqpoint{3.888713in}{1.549891in}}%
\pgfpathlineto{\pgfqpoint{3.884077in}{1.558421in}}%
\pgfpathlineto{\pgfqpoint{3.887037in}{1.580316in}}%
\pgfpathlineto{\pgfqpoint{3.892750in}{1.591758in}}%
\pgfpathlineto{\pgfqpoint{3.889894in}{1.607237in}}%
\pgfpathlineto{\pgfqpoint{3.895854in}{1.613747in}}%
\pgfpathlineto{\pgfqpoint{3.906727in}{1.631292in}}%
\pgfpathlineto{\pgfqpoint{3.915946in}{1.641883in}}%
\pgfpathlineto{\pgfqpoint{3.929422in}{1.664951in}}%
\pgfpathlineto{\pgfqpoint{3.939758in}{1.690079in}}%
\pgfpathlineto{\pgfqpoint{3.950730in}{1.713646in}}%
\pgfpathlineto{\pgfqpoint{3.952948in}{1.723489in}}%
\pgfpathlineto{\pgfqpoint{3.968085in}{1.751647in}}%
\pgfpathlineto{\pgfqpoint{3.970999in}{1.764014in}}%
\pgfpathlineto{\pgfqpoint{3.977444in}{1.777001in}}%
\pgfpathlineto{\pgfqpoint{3.979747in}{1.792843in}}%
\pgfpathlineto{\pgfqpoint{3.992051in}{1.800568in}}%
\pgfpathlineto{\pgfqpoint{4.006639in}{1.804461in}}%
\pgfpathlineto{\pgfqpoint{4.014015in}{1.795741in}}%
\pgfpathlineto{\pgfqpoint{4.020391in}{1.796386in}}%
\pgfpathlineto{\pgfqpoint{4.030347in}{1.785992in}}%
\pgfpathlineto{\pgfqpoint{4.031670in}{1.776202in}}%
\pgfpathlineto{\pgfqpoint{4.027895in}{1.757031in}}%
\pgfpathlineto{\pgfqpoint{4.040456in}{1.747249in}}%
\pgfpathlineto{\pgfqpoint{4.048603in}{1.743571in}}%
\pgfpathlineto{\pgfqpoint{4.070673in}{1.747325in}}%
\pgfpathlineto{\pgfqpoint{4.083460in}{1.744762in}}%
\pgfpathlineto{\pgfqpoint{4.092248in}{1.740554in}}%
\pgfpathlineto{\pgfqpoint{4.096877in}{1.732568in}}%
\pgfpathlineto{\pgfqpoint{4.126190in}{1.733634in}}%
\pgfpathlineto{\pgfqpoint{4.130851in}{1.728747in}}%
\pgfpathlineto{\pgfqpoint{4.141848in}{1.727692in}}%
\pgfpathlineto{\pgfqpoint{4.152921in}{1.730814in}}%
\pgfpathlineto{\pgfqpoint{4.185663in}{1.730028in}}%
\pgfpathlineto{\pgfqpoint{4.192734in}{1.727622in}}%
\pgfpathlineto{\pgfqpoint{4.200998in}{1.730394in}}%
\pgfpathlineto{\pgfqpoint{4.292475in}{1.706523in}}%
\pgfpathlineto{\pgfqpoint{4.296203in}{1.693698in}}%
\pgfpathlineto{\pgfqpoint{4.304597in}{1.686749in}}%
\pgfpathlineto{\pgfqpoint{4.305469in}{1.674240in}}%
\pgfpathlineto{\pgfqpoint{4.293191in}{1.661575in}}%
\pgfpathlineto{\pgfqpoint{4.287206in}{1.649603in}}%
\pgfpathlineto{\pgfqpoint{4.279997in}{1.643207in}}%
\pgfpathlineto{\pgfqpoint{4.279345in}{1.636806in}}%
\pgfpathlineto{\pgfqpoint{4.267402in}{1.628155in}}%
\pgfpathlineto{\pgfqpoint{4.259208in}{1.615650in}}%
\pgfpathlineto{\pgfqpoint{4.254620in}{1.612613in}}%
\pgfpathlineto{\pgfqpoint{4.252118in}{1.600478in}}%
\pgfpathlineto{\pgfqpoint{4.264170in}{1.588120in}}%
\pgfpathlineto{\pgfqpoint{4.252531in}{1.567559in}}%
\pgfpathlineto{\pgfqpoint{4.251119in}{1.559712in}}%
\pgfpathlineto{\pgfqpoint{4.224946in}{1.452341in}}%
\pgfpathlineto{\pgfqpoint{4.169893in}{1.466447in}}%
\pgfpathlineto{\pgfqpoint{4.116784in}{1.480141in}}%
\pgfpathlineto{\pgfqpoint{4.084744in}{1.489057in}}%
\pgfpathlineto{\pgfqpoint{4.043578in}{1.500784in}}%
\pgfpathlineto{\pgfqpoint{3.977913in}{1.521452in}}%
\pgfpathlineto{\pgfqpoint{3.906532in}{1.543676in}}%
\pgfpathlineto{\pgfqpoint{3.888713in}{1.549891in}}%
\pgfpathclose%
\pgfusepath{fill}%
\end{pgfscope}%
\begin{pgfscope}%
\pgfpathrectangle{\pgfqpoint{3.625000in}{0.100000in}}{\pgfqpoint{2.989028in}{1.913466in}}%
\pgfusepath{clip}%
\pgfsetbuttcap%
\pgfsetmiterjoin%
\definecolor{currentfill}{rgb}{0.925029,0.970012,0.632065}%
\pgfsetfillcolor{currentfill}%
\pgfsetlinewidth{0.000000pt}%
\definecolor{currentstroke}{rgb}{0.000000,0.000000,0.000000}%
\pgfsetstrokecolor{currentstroke}%
\pgfsetstrokeopacity{0.000000}%
\pgfsetdash{}{0pt}%
\pgfpathmoveto{\pgfqpoint{6.362703in}{1.505030in}}%
\pgfpathlineto{\pgfqpoint{6.357504in}{1.509256in}}%
\pgfpathlineto{\pgfqpoint{6.358228in}{1.522426in}}%
\pgfpathlineto{\pgfqpoint{6.352854in}{1.559858in}}%
\pgfpathlineto{\pgfqpoint{6.357066in}{1.573021in}}%
\pgfpathlineto{\pgfqpoint{6.360075in}{1.596366in}}%
\pgfpathlineto{\pgfqpoint{6.356015in}{1.607411in}}%
\pgfpathlineto{\pgfqpoint{6.370785in}{1.622115in}}%
\pgfpathlineto{\pgfqpoint{6.374502in}{1.633812in}}%
\pgfpathlineto{\pgfqpoint{6.367504in}{1.643043in}}%
\pgfpathlineto{\pgfqpoint{6.370485in}{1.654580in}}%
\pgfpathlineto{\pgfqpoint{6.368616in}{1.661826in}}%
\pgfpathlineto{\pgfqpoint{6.370221in}{1.677346in}}%
\pgfpathlineto{\pgfqpoint{6.373221in}{1.682137in}}%
\pgfpathlineto{\pgfqpoint{6.382471in}{1.685161in}}%
\pgfpathlineto{\pgfqpoint{6.395964in}{1.645645in}}%
\pgfpathlineto{\pgfqpoint{6.420998in}{1.563760in}}%
\pgfpathlineto{\pgfqpoint{6.430471in}{1.557411in}}%
\pgfpathlineto{\pgfqpoint{6.432175in}{1.550163in}}%
\pgfpathlineto{\pgfqpoint{6.437175in}{1.547235in}}%
\pgfpathlineto{\pgfqpoint{6.436795in}{1.534090in}}%
\pgfpathlineto{\pgfqpoint{6.431497in}{1.533879in}}%
\pgfpathlineto{\pgfqpoint{6.420764in}{1.525694in}}%
\pgfpathlineto{\pgfqpoint{6.417669in}{1.517482in}}%
\pgfpathlineto{\pgfqpoint{6.388942in}{1.510374in}}%
\pgfpathlineto{\pgfqpoint{6.362703in}{1.505030in}}%
\pgfpathclose%
\pgfusepath{fill}%
\end{pgfscope}%
\begin{pgfscope}%
\pgfpathrectangle{\pgfqpoint{3.625000in}{0.100000in}}{\pgfqpoint{2.989028in}{1.913466in}}%
\pgfusepath{clip}%
\pgfsetbuttcap%
\pgfsetmiterjoin%
\definecolor{currentfill}{rgb}{0.901961,0.960784,0.596078}%
\pgfsetfillcolor{currentfill}%
\pgfsetlinewidth{0.000000pt}%
\definecolor{currentstroke}{rgb}{0.000000,0.000000,0.000000}%
\pgfsetstrokecolor{currentstroke}%
\pgfsetstrokeopacity{0.000000}%
\pgfsetdash{}{0pt}%
\pgfpathmoveto{\pgfqpoint{5.464431in}{1.220143in}}%
\pgfpathlineto{\pgfqpoint{5.449266in}{1.235161in}}%
\pgfpathlineto{\pgfqpoint{5.400789in}{1.232365in}}%
\pgfpathlineto{\pgfqpoint{5.325184in}{1.229874in}}%
\pgfpathlineto{\pgfqpoint{5.249124in}{1.231060in}}%
\pgfpathlineto{\pgfqpoint{5.243785in}{1.240370in}}%
\pgfpathlineto{\pgfqpoint{5.245965in}{1.249511in}}%
\pgfpathlineto{\pgfqpoint{5.244930in}{1.265170in}}%
\pgfpathlineto{\pgfqpoint{5.240910in}{1.280338in}}%
\pgfpathlineto{\pgfqpoint{5.241378in}{1.288348in}}%
\pgfpathlineto{\pgfqpoint{5.232833in}{1.297398in}}%
\pgfpathlineto{\pgfqpoint{5.234692in}{1.309990in}}%
\pgfpathlineto{\pgfqpoint{5.221585in}{1.334864in}}%
\pgfpathlineto{\pgfqpoint{5.217680in}{1.355981in}}%
\pgfpathlineto{\pgfqpoint{5.208252in}{1.373157in}}%
\pgfpathlineto{\pgfqpoint{5.217201in}{1.394464in}}%
\pgfpathlineto{\pgfqpoint{5.212744in}{1.404379in}}%
\pgfpathlineto{\pgfqpoint{5.217728in}{1.422304in}}%
\pgfpathlineto{\pgfqpoint{5.265155in}{1.421790in}}%
\pgfpathlineto{\pgfqpoint{5.313246in}{1.421717in}}%
\pgfpathlineto{\pgfqpoint{5.358705in}{1.422385in}}%
\pgfpathlineto{\pgfqpoint{5.408600in}{1.423428in}}%
\pgfpathlineto{\pgfqpoint{5.467134in}{1.425554in}}%
\pgfpathlineto{\pgfqpoint{5.468697in}{1.417056in}}%
\pgfpathlineto{\pgfqpoint{5.475178in}{1.409931in}}%
\pgfpathlineto{\pgfqpoint{5.471075in}{1.402507in}}%
\pgfpathlineto{\pgfqpoint{5.472756in}{1.388016in}}%
\pgfpathlineto{\pgfqpoint{5.476772in}{1.376303in}}%
\pgfpathlineto{\pgfqpoint{5.494980in}{1.370103in}}%
\pgfpathlineto{\pgfqpoint{5.498100in}{1.361809in}}%
\pgfpathlineto{\pgfqpoint{5.508091in}{1.352497in}}%
\pgfpathlineto{\pgfqpoint{5.512167in}{1.342864in}}%
\pgfpathlineto{\pgfqpoint{5.522291in}{1.336402in}}%
\pgfpathlineto{\pgfqpoint{5.523875in}{1.328623in}}%
\pgfpathlineto{\pgfqpoint{5.521904in}{1.316845in}}%
\pgfpathlineto{\pgfqpoint{5.516751in}{1.313315in}}%
\pgfpathlineto{\pgfqpoint{5.515197in}{1.302104in}}%
\pgfpathlineto{\pgfqpoint{5.500389in}{1.293200in}}%
\pgfpathlineto{\pgfqpoint{5.481087in}{1.288322in}}%
\pgfpathlineto{\pgfqpoint{5.479318in}{1.277106in}}%
\pgfpathlineto{\pgfqpoint{5.486764in}{1.269093in}}%
\pgfpathlineto{\pgfqpoint{5.487068in}{1.259018in}}%
\pgfpathlineto{\pgfqpoint{5.481060in}{1.251098in}}%
\pgfpathlineto{\pgfqpoint{5.477907in}{1.239332in}}%
\pgfpathlineto{\pgfqpoint{5.467465in}{1.235438in}}%
\pgfpathlineto{\pgfqpoint{5.468130in}{1.222326in}}%
\pgfpathlineto{\pgfqpoint{5.464431in}{1.220143in}}%
\pgfpathclose%
\pgfusepath{fill}%
\end{pgfscope}%
\begin{pgfscope}%
\pgfpathrectangle{\pgfqpoint{3.625000in}{0.100000in}}{\pgfqpoint{2.989028in}{1.913466in}}%
\pgfusepath{clip}%
\pgfsetbuttcap%
\pgfsetmiterjoin%
\definecolor{currentfill}{rgb}{0.684198,0.872203,0.640369}%
\pgfsetfillcolor{currentfill}%
\pgfsetlinewidth{0.000000pt}%
\definecolor{currentstroke}{rgb}{0.000000,0.000000,0.000000}%
\pgfsetstrokecolor{currentstroke}%
\pgfsetstrokeopacity{0.000000}%
\pgfsetdash{}{0pt}%
\pgfpathmoveto{\pgfqpoint{6.405364in}{1.467139in}}%
\pgfpathlineto{\pgfqpoint{6.391000in}{1.464872in}}%
\pgfpathlineto{\pgfqpoint{6.325022in}{1.449880in}}%
\pgfpathlineto{\pgfqpoint{6.323869in}{1.451627in}}%
\pgfpathlineto{\pgfqpoint{6.324865in}{1.496989in}}%
\pgfpathlineto{\pgfqpoint{6.362703in}{1.505030in}}%
\pgfpathlineto{\pgfqpoint{6.388942in}{1.510374in}}%
\pgfpathlineto{\pgfqpoint{6.417669in}{1.517482in}}%
\pgfpathlineto{\pgfqpoint{6.420764in}{1.525694in}}%
\pgfpathlineto{\pgfqpoint{6.431497in}{1.533879in}}%
\pgfpathlineto{\pgfqpoint{6.436795in}{1.534090in}}%
\pgfpathlineto{\pgfqpoint{6.443759in}{1.522149in}}%
\pgfpathlineto{\pgfqpoint{6.437440in}{1.504717in}}%
\pgfpathlineto{\pgfqpoint{6.436512in}{1.494495in}}%
\pgfpathlineto{\pgfqpoint{6.449304in}{1.495455in}}%
\pgfpathlineto{\pgfqpoint{6.455105in}{1.490552in}}%
\pgfpathlineto{\pgfqpoint{6.468114in}{1.470509in}}%
\pgfpathlineto{\pgfqpoint{6.474579in}{1.468072in}}%
\pgfpathlineto{\pgfqpoint{6.485410in}{1.468976in}}%
\pgfpathlineto{\pgfqpoint{6.492946in}{1.475787in}}%
\pgfpathlineto{\pgfqpoint{6.497956in}{1.469708in}}%
\pgfpathlineto{\pgfqpoint{6.466451in}{1.453083in}}%
\pgfpathlineto{\pgfqpoint{6.465467in}{1.465013in}}%
\pgfpathlineto{\pgfqpoint{6.451155in}{1.446471in}}%
\pgfpathlineto{\pgfqpoint{6.446123in}{1.443277in}}%
\pgfpathlineto{\pgfqpoint{6.439105in}{1.453976in}}%
\pgfpathlineto{\pgfqpoint{6.437183in}{1.455442in}}%
\pgfpathlineto{\pgfqpoint{6.430650in}{1.458903in}}%
\pgfpathlineto{\pgfqpoint{6.424952in}{1.472940in}}%
\pgfpathlineto{\pgfqpoint{6.405364in}{1.467139in}}%
\pgfpathclose%
\pgfusepath{fill}%
\end{pgfscope}%
\begin{pgfscope}%
\pgfpathrectangle{\pgfqpoint{3.625000in}{0.100000in}}{\pgfqpoint{2.989028in}{1.913466in}}%
\pgfusepath{clip}%
\pgfsetbuttcap%
\pgfsetmiterjoin%
\definecolor{currentfill}{rgb}{0.304037,0.653749,0.691349}%
\pgfsetfillcolor{currentfill}%
\pgfsetlinewidth{0.000000pt}%
\definecolor{currentstroke}{rgb}{0.000000,0.000000,0.000000}%
\pgfsetstrokecolor{currentstroke}%
\pgfsetstrokeopacity{0.000000}%
\pgfsetdash{}{0pt}%
\pgfpathmoveto{\pgfqpoint{4.933935in}{1.205970in}}%
\pgfpathlineto{\pgfqpoint{4.939195in}{1.271196in}}%
\pgfpathlineto{\pgfqpoint{4.909404in}{1.273614in}}%
\pgfpathlineto{\pgfqpoint{4.840540in}{1.280295in}}%
\pgfpathlineto{\pgfqpoint{4.846947in}{1.345289in}}%
\pgfpathlineto{\pgfqpoint{4.853363in}{1.410780in}}%
\pgfpathlineto{\pgfqpoint{4.879836in}{1.407921in}}%
\pgfpathlineto{\pgfqpoint{4.947447in}{1.401533in}}%
\pgfpathlineto{\pgfqpoint{4.988365in}{1.398329in}}%
\pgfpathlineto{\pgfqpoint{5.037902in}{1.395489in}}%
\pgfpathlineto{\pgfqpoint{5.119439in}{1.391630in}}%
\pgfpathlineto{\pgfqpoint{5.121231in}{1.387975in}}%
\pgfpathlineto{\pgfqpoint{5.144594in}{1.376416in}}%
\pgfpathlineto{\pgfqpoint{5.151759in}{1.382286in}}%
\pgfpathlineto{\pgfqpoint{5.172400in}{1.381750in}}%
\pgfpathlineto{\pgfqpoint{5.184700in}{1.375307in}}%
\pgfpathlineto{\pgfqpoint{5.204205in}{1.367982in}}%
\pgfpathlineto{\pgfqpoint{5.209670in}{1.357295in}}%
\pgfpathlineto{\pgfqpoint{5.217680in}{1.355981in}}%
\pgfpathlineto{\pgfqpoint{5.221585in}{1.334864in}}%
\pgfpathlineto{\pgfqpoint{5.234692in}{1.309990in}}%
\pgfpathlineto{\pgfqpoint{5.232833in}{1.297398in}}%
\pgfpathlineto{\pgfqpoint{5.241378in}{1.288348in}}%
\pgfpathlineto{\pgfqpoint{5.240910in}{1.280338in}}%
\pgfpathlineto{\pgfqpoint{5.244930in}{1.265170in}}%
\pgfpathlineto{\pgfqpoint{5.245965in}{1.249511in}}%
\pgfpathlineto{\pgfqpoint{5.243785in}{1.240370in}}%
\pgfpathlineto{\pgfqpoint{5.249124in}{1.231060in}}%
\pgfpathlineto{\pgfqpoint{5.256447in}{1.214132in}}%
\pgfpathlineto{\pgfqpoint{5.263455in}{1.207242in}}%
\pgfpathlineto{\pgfqpoint{5.271810in}{1.192312in}}%
\pgfpathlineto{\pgfqpoint{5.248131in}{1.192066in}}%
\pgfpathlineto{\pgfqpoint{5.196933in}{1.192859in}}%
\pgfpathlineto{\pgfqpoint{5.140367in}{1.194592in}}%
\pgfpathlineto{\pgfqpoint{5.083468in}{1.196775in}}%
\pgfpathlineto{\pgfqpoint{4.999809in}{1.201345in}}%
\pgfpathlineto{\pgfqpoint{4.933935in}{1.205970in}}%
\pgfpathclose%
\pgfusepath{fill}%
\end{pgfscope}%
\begin{pgfscope}%
\pgfpathrectangle{\pgfqpoint{3.625000in}{0.100000in}}{\pgfqpoint{2.989028in}{1.913466in}}%
\pgfusepath{clip}%
\pgfsetbuttcap%
\pgfsetmiterjoin%
\definecolor{currentfill}{rgb}{0.932718,0.973087,0.644060}%
\pgfsetfillcolor{currentfill}%
\pgfsetlinewidth{0.000000pt}%
\definecolor{currentstroke}{rgb}{0.000000,0.000000,0.000000}%
\pgfsetstrokecolor{currentstroke}%
\pgfsetstrokeopacity{0.000000}%
\pgfsetdash{}{0pt}%
\pgfpathmoveto{\pgfqpoint{6.023213in}{1.402714in}}%
\pgfpathlineto{\pgfqpoint{6.049409in}{1.427692in}}%
\pgfpathlineto{\pgfqpoint{6.052665in}{1.436570in}}%
\pgfpathlineto{\pgfqpoint{6.060341in}{1.444171in}}%
\pgfpathlineto{\pgfqpoint{6.054581in}{1.455238in}}%
\pgfpathlineto{\pgfqpoint{6.047360in}{1.461711in}}%
\pgfpathlineto{\pgfqpoint{6.045273in}{1.473182in}}%
\pgfpathlineto{\pgfqpoint{6.072151in}{1.484959in}}%
\pgfpathlineto{\pgfqpoint{6.094405in}{1.488685in}}%
\pgfpathlineto{\pgfqpoint{6.106363in}{1.488934in}}%
\pgfpathlineto{\pgfqpoint{6.115473in}{1.484433in}}%
\pgfpathlineto{\pgfqpoint{6.124352in}{1.488453in}}%
\pgfpathlineto{\pgfqpoint{6.146010in}{1.492956in}}%
\pgfpathlineto{\pgfqpoint{6.153493in}{1.498769in}}%
\pgfpathlineto{\pgfqpoint{6.164589in}{1.511638in}}%
\pgfpathlineto{\pgfqpoint{6.174683in}{1.517325in}}%
\pgfpathlineto{\pgfqpoint{6.175395in}{1.522779in}}%
\pgfpathlineto{\pgfqpoint{6.170097in}{1.535222in}}%
\pgfpathlineto{\pgfqpoint{6.173927in}{1.542544in}}%
\pgfpathlineto{\pgfqpoint{6.168770in}{1.550418in}}%
\pgfpathlineto{\pgfqpoint{6.160873in}{1.550968in}}%
\pgfpathlineto{\pgfqpoint{6.180638in}{1.574751in}}%
\pgfpathlineto{\pgfqpoint{6.182985in}{1.583815in}}%
\pgfpathlineto{\pgfqpoint{6.198541in}{1.606931in}}%
\pgfpathlineto{\pgfqpoint{6.212981in}{1.619440in}}%
\pgfpathlineto{\pgfqpoint{6.222897in}{1.624664in}}%
\pgfpathlineto{\pgfqpoint{6.255342in}{1.632016in}}%
\pgfpathlineto{\pgfqpoint{6.285752in}{1.640318in}}%
\pgfpathlineto{\pgfqpoint{6.289780in}{1.627285in}}%
\pgfpathlineto{\pgfqpoint{6.290065in}{1.615202in}}%
\pgfpathlineto{\pgfqpoint{6.297002in}{1.604233in}}%
\pgfpathlineto{\pgfqpoint{6.299140in}{1.593064in}}%
\pgfpathlineto{\pgfqpoint{6.296474in}{1.577911in}}%
\pgfpathlineto{\pgfqpoint{6.304593in}{1.557473in}}%
\pgfpathlineto{\pgfqpoint{6.313848in}{1.546108in}}%
\pgfpathlineto{\pgfqpoint{6.321007in}{1.509167in}}%
\pgfpathlineto{\pgfqpoint{6.324865in}{1.496989in}}%
\pgfpathlineto{\pgfqpoint{6.323869in}{1.451627in}}%
\pgfpathlineto{\pgfqpoint{6.325022in}{1.449880in}}%
\pgfpathlineto{\pgfqpoint{6.333450in}{1.401092in}}%
\pgfpathlineto{\pgfqpoint{6.338176in}{1.396647in}}%
\pgfpathlineto{\pgfqpoint{6.328010in}{1.386769in}}%
\pgfpathlineto{\pgfqpoint{6.333026in}{1.381100in}}%
\pgfpathlineto{\pgfqpoint{6.328617in}{1.372526in}}%
\pgfpathlineto{\pgfqpoint{6.328654in}{1.368875in}}%
\pgfpathlineto{\pgfqpoint{6.323142in}{1.365581in}}%
\pgfpathlineto{\pgfqpoint{6.320458in}{1.358300in}}%
\pgfpathlineto{\pgfqpoint{6.321310in}{1.378322in}}%
\pgfpathlineto{\pgfqpoint{6.304225in}{1.382730in}}%
\pgfpathlineto{\pgfqpoint{6.277507in}{1.391844in}}%
\pgfpathlineto{\pgfqpoint{6.273728in}{1.396333in}}%
\pgfpathlineto{\pgfqpoint{6.262564in}{1.397408in}}%
\pgfpathlineto{\pgfqpoint{6.256142in}{1.404090in}}%
\pgfpathlineto{\pgfqpoint{6.252676in}{1.417397in}}%
\pgfpathlineto{\pgfqpoint{6.243572in}{1.419052in}}%
\pgfpathlineto{\pgfqpoint{6.237484in}{1.426033in}}%
\pgfpathlineto{\pgfqpoint{6.159972in}{1.410406in}}%
\pgfpathlineto{\pgfqpoint{6.122928in}{1.402745in}}%
\pgfpathlineto{\pgfqpoint{6.066683in}{1.392479in}}%
\pgfpathlineto{\pgfqpoint{6.026180in}{1.385646in}}%
\pgfpathlineto{\pgfqpoint{6.023213in}{1.402714in}}%
\pgfpathclose%
\pgfusepath{fill}%
\end{pgfscope}%
\begin{pgfscope}%
\pgfpathrectangle{\pgfqpoint{3.625000in}{0.100000in}}{\pgfqpoint{2.989028in}{1.913466in}}%
\pgfusepath{clip}%
\pgfsetbuttcap%
\pgfsetmiterjoin%
\definecolor{currentfill}{rgb}{0.932718,0.973087,0.644060}%
\pgfsetfillcolor{currentfill}%
\pgfsetlinewidth{0.000000pt}%
\definecolor{currentstroke}{rgb}{0.000000,0.000000,0.000000}%
\pgfsetstrokecolor{currentstroke}%
\pgfsetstrokeopacity{0.000000}%
\pgfsetdash{}{0pt}%
\pgfpathmoveto{\pgfqpoint{6.334390in}{1.354233in}}%
\pgfpathlineto{\pgfqpoint{6.322403in}{1.350498in}}%
\pgfpathlineto{\pgfqpoint{6.320386in}{1.353932in}}%
\pgfpathlineto{\pgfqpoint{6.324234in}{1.365456in}}%
\pgfpathlineto{\pgfqpoint{6.331360in}{1.366591in}}%
\pgfpathlineto{\pgfqpoint{6.330722in}{1.370219in}}%
\pgfpathlineto{\pgfqpoint{6.337120in}{1.375665in}}%
\pgfpathlineto{\pgfqpoint{6.355618in}{1.380004in}}%
\pgfpathlineto{\pgfqpoint{6.363853in}{1.386571in}}%
\pgfpathlineto{\pgfqpoint{6.382353in}{1.392045in}}%
\pgfpathlineto{\pgfqpoint{6.390790in}{1.390042in}}%
\pgfpathlineto{\pgfqpoint{6.397888in}{1.398866in}}%
\pgfpathlineto{\pgfqpoint{6.408615in}{1.400032in}}%
\pgfpathlineto{\pgfqpoint{6.390314in}{1.382820in}}%
\pgfpathlineto{\pgfqpoint{6.334390in}{1.354233in}}%
\pgfpathclose%
\pgfusepath{fill}%
\end{pgfscope}%
\begin{pgfscope}%
\pgfpathrectangle{\pgfqpoint{3.625000in}{0.100000in}}{\pgfqpoint{2.989028in}{1.913466in}}%
\pgfusepath{clip}%
\pgfsetbuttcap%
\pgfsetmiterjoin%
\definecolor{currentfill}{rgb}{0.774933,0.909112,0.621915}%
\pgfsetfillcolor{currentfill}%
\pgfsetlinewidth{0.000000pt}%
\definecolor{currentstroke}{rgb}{0.000000,0.000000,0.000000}%
\pgfsetstrokecolor{currentstroke}%
\pgfsetstrokeopacity{0.000000}%
\pgfsetdash{}{0pt}%
\pgfpathmoveto{\pgfqpoint{6.065072in}{1.240581in}}%
\pgfpathlineto{\pgfqpoint{6.013234in}{1.232010in}}%
\pgfpathlineto{\pgfqpoint{6.003831in}{1.291252in}}%
\pgfpathlineto{\pgfqpoint{5.989869in}{1.378577in}}%
\pgfpathlineto{\pgfqpoint{6.023213in}{1.402714in}}%
\pgfpathlineto{\pgfqpoint{6.026180in}{1.385646in}}%
\pgfpathlineto{\pgfqpoint{6.066683in}{1.392479in}}%
\pgfpathlineto{\pgfqpoint{6.122928in}{1.402745in}}%
\pgfpathlineto{\pgfqpoint{6.159972in}{1.410406in}}%
\pgfpathlineto{\pgfqpoint{6.237484in}{1.426033in}}%
\pgfpathlineto{\pgfqpoint{6.243572in}{1.419052in}}%
\pgfpathlineto{\pgfqpoint{6.252676in}{1.417397in}}%
\pgfpathlineto{\pgfqpoint{6.256142in}{1.404090in}}%
\pgfpathlineto{\pgfqpoint{6.262564in}{1.397408in}}%
\pgfpathlineto{\pgfqpoint{6.273728in}{1.396333in}}%
\pgfpathlineto{\pgfqpoint{6.277507in}{1.391844in}}%
\pgfpathlineto{\pgfqpoint{6.273671in}{1.388380in}}%
\pgfpathlineto{\pgfqpoint{6.270215in}{1.376126in}}%
\pgfpathlineto{\pgfqpoint{6.262116in}{1.362357in}}%
\pgfpathlineto{\pgfqpoint{6.267552in}{1.356369in}}%
\pgfpathlineto{\pgfqpoint{6.263104in}{1.349953in}}%
\pgfpathlineto{\pgfqpoint{6.265613in}{1.335884in}}%
\pgfpathlineto{\pgfqpoint{6.273625in}{1.328499in}}%
\pgfpathlineto{\pgfqpoint{6.292634in}{1.316500in}}%
\pgfpathlineto{\pgfqpoint{6.277326in}{1.299570in}}%
\pgfpathlineto{\pgfqpoint{6.277111in}{1.293145in}}%
\pgfpathlineto{\pgfqpoint{6.264648in}{1.284858in}}%
\pgfpathlineto{\pgfqpoint{6.250867in}{1.283302in}}%
\pgfpathlineto{\pgfqpoint{6.247458in}{1.276100in}}%
\pgfpathlineto{\pgfqpoint{6.209127in}{1.267806in}}%
\pgfpathlineto{\pgfqpoint{6.164399in}{1.258797in}}%
\pgfpathlineto{\pgfqpoint{6.133655in}{1.253292in}}%
\pgfpathlineto{\pgfqpoint{6.065072in}{1.240581in}}%
\pgfpathclose%
\pgfusepath{fill}%
\end{pgfscope}%
\begin{pgfscope}%
\pgfpathrectangle{\pgfqpoint{3.625000in}{0.100000in}}{\pgfqpoint{2.989028in}{1.913466in}}%
\pgfusepath{clip}%
\pgfsetbuttcap%
\pgfsetmiterjoin%
\definecolor{currentfill}{rgb}{0.747712,0.898039,0.627451}%
\pgfsetfillcolor{currentfill}%
\pgfsetlinewidth{0.000000pt}%
\definecolor{currentstroke}{rgb}{0.000000,0.000000,0.000000}%
\pgfsetstrokecolor{currentstroke}%
\pgfsetstrokeopacity{0.000000}%
\pgfsetdash{}{0pt}%
\pgfpathmoveto{\pgfqpoint{6.325022in}{1.449880in}}%
\pgfpathlineto{\pgfqpoint{6.391000in}{1.464872in}}%
\pgfpathlineto{\pgfqpoint{6.405364in}{1.467139in}}%
\pgfpathlineto{\pgfqpoint{6.414871in}{1.429760in}}%
\pgfpathlineto{\pgfqpoint{6.413366in}{1.423067in}}%
\pgfpathlineto{\pgfqpoint{6.382827in}{1.411244in}}%
\pgfpathlineto{\pgfqpoint{6.364596in}{1.407111in}}%
\pgfpathlineto{\pgfqpoint{6.356849in}{1.397842in}}%
\pgfpathlineto{\pgfqpoint{6.333026in}{1.381100in}}%
\pgfpathlineto{\pgfqpoint{6.328010in}{1.386769in}}%
\pgfpathlineto{\pgfqpoint{6.338176in}{1.396647in}}%
\pgfpathlineto{\pgfqpoint{6.333450in}{1.401092in}}%
\pgfpathlineto{\pgfqpoint{6.325022in}{1.449880in}}%
\pgfpathclose%
\pgfusepath{fill}%
\end{pgfscope}%
\begin{pgfscope}%
\pgfpathrectangle{\pgfqpoint{3.625000in}{0.100000in}}{\pgfqpoint{2.989028in}{1.913466in}}%
\pgfusepath{clip}%
\pgfsetbuttcap%
\pgfsetmiterjoin%
\definecolor{currentfill}{rgb}{0.936563,0.974625,0.650058}%
\pgfsetfillcolor{currentfill}%
\pgfsetlinewidth{0.000000pt}%
\definecolor{currentstroke}{rgb}{0.000000,0.000000,0.000000}%
\pgfsetstrokecolor{currentstroke}%
\pgfsetstrokeopacity{0.000000}%
\pgfsetdash{}{0pt}%
\pgfpathmoveto{\pgfqpoint{6.405364in}{1.467139in}}%
\pgfpathlineto{\pgfqpoint{6.424952in}{1.472940in}}%
\pgfpathlineto{\pgfqpoint{6.430650in}{1.458903in}}%
\pgfpathlineto{\pgfqpoint{6.437183in}{1.455442in}}%
\pgfpathlineto{\pgfqpoint{6.429121in}{1.449529in}}%
\pgfpathlineto{\pgfqpoint{6.430138in}{1.432164in}}%
\pgfpathlineto{\pgfqpoint{6.413366in}{1.423067in}}%
\pgfpathlineto{\pgfqpoint{6.414871in}{1.429760in}}%
\pgfpathlineto{\pgfqpoint{6.405364in}{1.467139in}}%
\pgfpathclose%
\pgfusepath{fill}%
\end{pgfscope}%
\begin{pgfscope}%
\pgfpathrectangle{\pgfqpoint{3.625000in}{0.100000in}}{\pgfqpoint{2.989028in}{1.913466in}}%
\pgfusepath{clip}%
\pgfsetbuttcap%
\pgfsetmiterjoin%
\definecolor{currentfill}{rgb}{0.856594,0.942330,0.605306}%
\pgfsetfillcolor{currentfill}%
\pgfsetlinewidth{0.000000pt}%
\definecolor{currentstroke}{rgb}{0.000000,0.000000,0.000000}%
\pgfsetstrokecolor{currentstroke}%
\pgfsetstrokeopacity{0.000000}%
\pgfsetdash{}{0pt}%
\pgfpathmoveto{\pgfqpoint{6.262440in}{1.278737in}}%
\pgfpathlineto{\pgfqpoint{6.264648in}{1.284858in}}%
\pgfpathlineto{\pgfqpoint{6.277111in}{1.293145in}}%
\pgfpathlineto{\pgfqpoint{6.277326in}{1.299570in}}%
\pgfpathlineto{\pgfqpoint{6.292634in}{1.316500in}}%
\pgfpathlineto{\pgfqpoint{6.273625in}{1.328499in}}%
\pgfpathlineto{\pgfqpoint{6.265613in}{1.335884in}}%
\pgfpathlineto{\pgfqpoint{6.263104in}{1.349953in}}%
\pgfpathlineto{\pgfqpoint{6.267552in}{1.356369in}}%
\pgfpathlineto{\pgfqpoint{6.262116in}{1.362357in}}%
\pgfpathlineto{\pgfqpoint{6.270215in}{1.376126in}}%
\pgfpathlineto{\pgfqpoint{6.273671in}{1.388380in}}%
\pgfpathlineto{\pgfqpoint{6.277507in}{1.391844in}}%
\pgfpathlineto{\pgfqpoint{6.304225in}{1.382730in}}%
\pgfpathlineto{\pgfqpoint{6.321310in}{1.378322in}}%
\pgfpathlineto{\pgfqpoint{6.320458in}{1.358300in}}%
\pgfpathlineto{\pgfqpoint{6.310080in}{1.343116in}}%
\pgfpathlineto{\pgfqpoint{6.318632in}{1.340882in}}%
\pgfpathlineto{\pgfqpoint{6.327511in}{1.334366in}}%
\pgfpathlineto{\pgfqpoint{6.328034in}{1.316545in}}%
\pgfpathlineto{\pgfqpoint{6.325347in}{1.303928in}}%
\pgfpathlineto{\pgfqpoint{6.327135in}{1.293566in}}%
\pgfpathlineto{\pgfqpoint{6.311701in}{1.258597in}}%
\pgfpathlineto{\pgfqpoint{6.306154in}{1.242268in}}%
\pgfpathlineto{\pgfqpoint{6.298409in}{1.250205in}}%
\pgfpathlineto{\pgfqpoint{6.288139in}{1.248853in}}%
\pgfpathlineto{\pgfqpoint{6.262435in}{1.263703in}}%
\pgfpathlineto{\pgfqpoint{6.259805in}{1.271656in}}%
\pgfpathlineto{\pgfqpoint{6.262440in}{1.278737in}}%
\pgfpathclose%
\pgfusepath{fill}%
\end{pgfscope}%
\begin{pgfscope}%
\pgfpathrectangle{\pgfqpoint{3.625000in}{0.100000in}}{\pgfqpoint{2.989028in}{1.913466in}}%
\pgfusepath{clip}%
\pgfsetbuttcap%
\pgfsetmiterjoin%
\definecolor{currentfill}{rgb}{0.905805,0.962322,0.602076}%
\pgfsetfillcolor{currentfill}%
\pgfsetlinewidth{0.000000pt}%
\definecolor{currentstroke}{rgb}{0.000000,0.000000,0.000000}%
\pgfsetstrokecolor{currentstroke}%
\pgfsetstrokeopacity{0.000000}%
\pgfsetdash{}{0pt}%
\pgfpathmoveto{\pgfqpoint{5.647290in}{1.062454in}}%
\pgfpathlineto{\pgfqpoint{5.646088in}{1.078676in}}%
\pgfpathlineto{\pgfqpoint{5.650759in}{1.086724in}}%
\pgfpathlineto{\pgfqpoint{5.647681in}{1.090668in}}%
\pgfpathlineto{\pgfqpoint{5.658907in}{1.103580in}}%
\pgfpathlineto{\pgfqpoint{5.658263in}{1.106692in}}%
\pgfpathlineto{\pgfqpoint{5.669170in}{1.129219in}}%
\pgfpathlineto{\pgfqpoint{5.666966in}{1.140083in}}%
\pgfpathlineto{\pgfqpoint{5.659057in}{1.151457in}}%
\pgfpathlineto{\pgfqpoint{5.664547in}{1.165295in}}%
\pgfpathlineto{\pgfqpoint{5.657431in}{1.256358in}}%
\pgfpathlineto{\pgfqpoint{5.652306in}{1.320243in}}%
\pgfpathlineto{\pgfqpoint{5.659385in}{1.314949in}}%
\pgfpathlineto{\pgfqpoint{5.667283in}{1.315093in}}%
\pgfpathlineto{\pgfqpoint{5.685936in}{1.325898in}}%
\pgfpathlineto{\pgfqpoint{5.743150in}{1.331236in}}%
\pgfpathlineto{\pgfqpoint{5.785498in}{1.335699in}}%
\pgfpathlineto{\pgfqpoint{5.785872in}{1.331556in}}%
\pgfpathlineto{\pgfqpoint{5.795539in}{1.244011in}}%
\pgfpathlineto{\pgfqpoint{5.803860in}{1.162554in}}%
\pgfpathlineto{\pgfqpoint{5.800273in}{1.158730in}}%
\pgfpathlineto{\pgfqpoint{5.805783in}{1.149238in}}%
\pgfpathlineto{\pgfqpoint{5.805755in}{1.142397in}}%
\pgfpathlineto{\pgfqpoint{5.797893in}{1.140679in}}%
\pgfpathlineto{\pgfqpoint{5.789096in}{1.134084in}}%
\pgfpathlineto{\pgfqpoint{5.783139in}{1.136680in}}%
\pgfpathlineto{\pgfqpoint{5.774225in}{1.132459in}}%
\pgfpathlineto{\pgfqpoint{5.776983in}{1.123981in}}%
\pgfpathlineto{\pgfqpoint{5.767820in}{1.115463in}}%
\pgfpathlineto{\pgfqpoint{5.765290in}{1.105610in}}%
\pgfpathlineto{\pgfqpoint{5.758989in}{1.103993in}}%
\pgfpathlineto{\pgfqpoint{5.754287in}{1.096531in}}%
\pgfpathlineto{\pgfqpoint{5.754903in}{1.089016in}}%
\pgfpathlineto{\pgfqpoint{5.749371in}{1.083732in}}%
\pgfpathlineto{\pgfqpoint{5.741048in}{1.084548in}}%
\pgfpathlineto{\pgfqpoint{5.734721in}{1.092656in}}%
\pgfpathlineto{\pgfqpoint{5.724000in}{1.084848in}}%
\pgfpathlineto{\pgfqpoint{5.724753in}{1.078025in}}%
\pgfpathlineto{\pgfqpoint{5.712829in}{1.074039in}}%
\pgfpathlineto{\pgfqpoint{5.709834in}{1.079056in}}%
\pgfpathlineto{\pgfqpoint{5.698226in}{1.073367in}}%
\pgfpathlineto{\pgfqpoint{5.693985in}{1.065222in}}%
\pgfpathlineto{\pgfqpoint{5.680006in}{1.073574in}}%
\pgfpathlineto{\pgfqpoint{5.657850in}{1.068041in}}%
\pgfpathlineto{\pgfqpoint{5.647290in}{1.062454in}}%
\pgfpathclose%
\pgfusepath{fill}%
\end{pgfscope}%
\begin{pgfscope}%
\pgfpathrectangle{\pgfqpoint{3.625000in}{0.100000in}}{\pgfqpoint{2.989028in}{1.913466in}}%
\pgfusepath{clip}%
\pgfsetbuttcap%
\pgfsetmiterjoin%
\definecolor{currentfill}{rgb}{0.720492,0.886967,0.632987}%
\pgfsetfillcolor{currentfill}%
\pgfsetlinewidth{0.000000pt}%
\definecolor{currentstroke}{rgb}{0.000000,0.000000,0.000000}%
\pgfsetstrokecolor{currentstroke}%
\pgfsetstrokeopacity{0.000000}%
\pgfsetdash{}{0pt}%
\pgfpathmoveto{\pgfqpoint{4.084744in}{1.489057in}}%
\pgfpathlineto{\pgfqpoint{4.116784in}{1.480141in}}%
\pgfpathlineto{\pgfqpoint{4.169893in}{1.466447in}}%
\pgfpathlineto{\pgfqpoint{4.224946in}{1.452341in}}%
\pgfpathlineto{\pgfqpoint{4.275689in}{1.440103in}}%
\pgfpathlineto{\pgfqpoint{4.319534in}{1.430177in}}%
\pgfpathlineto{\pgfqpoint{4.366442in}{1.419862in}}%
\pgfpathlineto{\pgfqpoint{4.352889in}{1.355901in}}%
\pgfpathlineto{\pgfqpoint{4.340815in}{1.299101in}}%
\pgfpathlineto{\pgfqpoint{4.321008in}{1.207454in}}%
\pgfpathlineto{\pgfqpoint{4.306137in}{1.138280in}}%
\pgfpathlineto{\pgfqpoint{4.298103in}{1.099678in}}%
\pgfpathlineto{\pgfqpoint{4.287801in}{1.049576in}}%
\pgfpathlineto{\pgfqpoint{4.280674in}{1.039412in}}%
\pgfpathlineto{\pgfqpoint{4.274938in}{1.039072in}}%
\pgfpathlineto{\pgfqpoint{4.270846in}{1.047945in}}%
\pgfpathlineto{\pgfqpoint{4.261451in}{1.051167in}}%
\pgfpathlineto{\pgfqpoint{4.251328in}{1.050037in}}%
\pgfpathlineto{\pgfqpoint{4.248454in}{1.042780in}}%
\pgfpathlineto{\pgfqpoint{4.248155in}{1.017578in}}%
\pgfpathlineto{\pgfqpoint{4.245123in}{1.011825in}}%
\pgfpathlineto{\pgfqpoint{4.246856in}{0.991592in}}%
\pgfpathlineto{\pgfqpoint{4.240525in}{0.978113in}}%
\pgfpathlineto{\pgfqpoint{4.189188in}{1.057164in}}%
\pgfpathlineto{\pgfqpoint{4.138630in}{1.133980in}}%
\pgfpathlineto{\pgfqpoint{4.112312in}{1.174260in}}%
\pgfpathlineto{\pgfqpoint{4.090358in}{1.209028in}}%
\pgfpathlineto{\pgfqpoint{4.062695in}{1.252036in}}%
\pgfpathlineto{\pgfqpoint{4.031381in}{1.300279in}}%
\pgfpathlineto{\pgfqpoint{4.044255in}{1.346069in}}%
\pgfpathlineto{\pgfqpoint{4.070163in}{1.437904in}}%
\pgfpathlineto{\pgfqpoint{4.084744in}{1.489057in}}%
\pgfpathclose%
\pgfusepath{fill}%
\end{pgfscope}%
\begin{pgfscope}%
\pgfpathrectangle{\pgfqpoint{3.625000in}{0.100000in}}{\pgfqpoint{2.989028in}{1.913466in}}%
\pgfusepath{clip}%
\pgfsetbuttcap%
\pgfsetmiterjoin%
\definecolor{currentfill}{rgb}{0.675125,0.868512,0.642215}%
\pgfsetfillcolor{currentfill}%
\pgfsetlinewidth{0.000000pt}%
\definecolor{currentstroke}{rgb}{0.000000,0.000000,0.000000}%
\pgfsetstrokecolor{currentstroke}%
\pgfsetstrokeopacity{0.000000}%
\pgfsetdash{}{0pt}%
\pgfpathmoveto{\pgfqpoint{4.366442in}{1.419862in}}%
\pgfpathlineto{\pgfqpoint{4.416592in}{1.409886in}}%
\pgfpathlineto{\pgfqpoint{4.509436in}{1.391963in}}%
\pgfpathlineto{\pgfqpoint{4.497815in}{1.327496in}}%
\pgfpathlineto{\pgfqpoint{4.548928in}{1.318744in}}%
\pgfpathlineto{\pgfqpoint{4.595495in}{1.311293in}}%
\pgfpathlineto{\pgfqpoint{4.587404in}{1.260326in}}%
\pgfpathlineto{\pgfqpoint{4.578829in}{1.205365in}}%
\pgfpathlineto{\pgfqpoint{4.567350in}{1.133200in}}%
\pgfpathlineto{\pgfqpoint{4.567053in}{1.127151in}}%
\pgfpathlineto{\pgfqpoint{4.555137in}{1.052362in}}%
\pgfpathlineto{\pgfqpoint{4.506086in}{1.059964in}}%
\pgfpathlineto{\pgfqpoint{4.481091in}{1.064984in}}%
\pgfpathlineto{\pgfqpoint{4.390746in}{1.080865in}}%
\pgfpathlineto{\pgfqpoint{4.356724in}{1.087575in}}%
\pgfpathlineto{\pgfqpoint{4.298103in}{1.099678in}}%
\pgfpathlineto{\pgfqpoint{4.306137in}{1.138280in}}%
\pgfpathlineto{\pgfqpoint{4.321008in}{1.207454in}}%
\pgfpathlineto{\pgfqpoint{4.340815in}{1.299101in}}%
\pgfpathlineto{\pgfqpoint{4.352889in}{1.355901in}}%
\pgfpathlineto{\pgfqpoint{4.366442in}{1.419862in}}%
\pgfpathclose%
\pgfusepath{fill}%
\end{pgfscope}%
\begin{pgfscope}%
\pgfpathrectangle{\pgfqpoint{3.625000in}{0.100000in}}{\pgfqpoint{2.989028in}{1.913466in}}%
\pgfusepath{clip}%
\pgfsetbuttcap%
\pgfsetmiterjoin%
\definecolor{currentfill}{rgb}{0.774933,0.909112,0.621915}%
\pgfsetfillcolor{currentfill}%
\pgfsetlinewidth{0.000000pt}%
\definecolor{currentstroke}{rgb}{0.000000,0.000000,0.000000}%
\pgfsetstrokecolor{currentstroke}%
\pgfsetstrokeopacity{0.000000}%
\pgfsetdash{}{0pt}%
\pgfpathmoveto{\pgfqpoint{3.888713in}{1.549891in}}%
\pgfpathlineto{\pgfqpoint{3.906532in}{1.543676in}}%
\pgfpathlineto{\pgfqpoint{3.977913in}{1.521452in}}%
\pgfpathlineto{\pgfqpoint{4.043578in}{1.500784in}}%
\pgfpathlineto{\pgfqpoint{4.084744in}{1.489057in}}%
\pgfpathlineto{\pgfqpoint{4.070163in}{1.437904in}}%
\pgfpathlineto{\pgfqpoint{4.044255in}{1.346069in}}%
\pgfpathlineto{\pgfqpoint{4.031381in}{1.300279in}}%
\pgfpathlineto{\pgfqpoint{4.062695in}{1.252036in}}%
\pgfpathlineto{\pgfqpoint{4.090358in}{1.209028in}}%
\pgfpathlineto{\pgfqpoint{4.112312in}{1.174260in}}%
\pgfpathlineto{\pgfqpoint{4.138630in}{1.133980in}}%
\pgfpathlineto{\pgfqpoint{4.189188in}{1.057164in}}%
\pgfpathlineto{\pgfqpoint{4.240525in}{0.978113in}}%
\pgfpathlineto{\pgfqpoint{4.238464in}{0.970274in}}%
\pgfpathlineto{\pgfqpoint{4.244656in}{0.957789in}}%
\pgfpathlineto{\pgfqpoint{4.245881in}{0.940711in}}%
\pgfpathlineto{\pgfqpoint{4.254874in}{0.932426in}}%
\pgfpathlineto{\pgfqpoint{4.255230in}{0.925750in}}%
\pgfpathlineto{\pgfqpoint{4.239127in}{0.918174in}}%
\pgfpathlineto{\pgfqpoint{4.231455in}{0.910589in}}%
\pgfpathlineto{\pgfqpoint{4.229063in}{0.893839in}}%
\pgfpathlineto{\pgfqpoint{4.225180in}{0.884704in}}%
\pgfpathlineto{\pgfqpoint{4.216998in}{0.877004in}}%
\pgfpathlineto{\pgfqpoint{4.208861in}{0.856982in}}%
\pgfpathlineto{\pgfqpoint{4.220303in}{0.846552in}}%
\pgfpathlineto{\pgfqpoint{4.218826in}{0.837954in}}%
\pgfpathlineto{\pgfqpoint{4.209444in}{0.831972in}}%
\pgfpathlineto{\pgfqpoint{4.202975in}{0.833027in}}%
\pgfpathlineto{\pgfqpoint{4.126622in}{0.843623in}}%
\pgfpathlineto{\pgfqpoint{4.070229in}{0.851623in}}%
\pgfpathlineto{\pgfqpoint{4.072695in}{0.860728in}}%
\pgfpathlineto{\pgfqpoint{4.069024in}{0.875844in}}%
\pgfpathlineto{\pgfqpoint{4.068653in}{0.891092in}}%
\pgfpathlineto{\pgfqpoint{4.066254in}{0.900014in}}%
\pgfpathlineto{\pgfqpoint{4.058864in}{0.912772in}}%
\pgfpathlineto{\pgfqpoint{4.037605in}{0.942205in}}%
\pgfpathlineto{\pgfqpoint{4.030643in}{0.945811in}}%
\pgfpathlineto{\pgfqpoint{4.021672in}{0.945755in}}%
\pgfpathlineto{\pgfqpoint{4.023721in}{0.955040in}}%
\pgfpathlineto{\pgfqpoint{4.019473in}{0.966651in}}%
\pgfpathlineto{\pgfqpoint{3.998608in}{0.972422in}}%
\pgfpathlineto{\pgfqpoint{3.985889in}{0.983108in}}%
\pgfpathlineto{\pgfqpoint{3.984834in}{0.989647in}}%
\pgfpathlineto{\pgfqpoint{3.970189in}{1.005835in}}%
\pgfpathlineto{\pgfqpoint{3.956249in}{1.008932in}}%
\pgfpathlineto{\pgfqpoint{3.943332in}{1.017159in}}%
\pgfpathlineto{\pgfqpoint{3.926344in}{1.020005in}}%
\pgfpathlineto{\pgfqpoint{3.919086in}{1.030977in}}%
\pgfpathlineto{\pgfqpoint{3.925979in}{1.048345in}}%
\pgfpathlineto{\pgfqpoint{3.923871in}{1.052250in}}%
\pgfpathlineto{\pgfqpoint{3.929611in}{1.066728in}}%
\pgfpathlineto{\pgfqpoint{3.919400in}{1.074432in}}%
\pgfpathlineto{\pgfqpoint{3.922713in}{1.088401in}}%
\pgfpathlineto{\pgfqpoint{3.917259in}{1.091971in}}%
\pgfpathlineto{\pgfqpoint{3.912550in}{1.105174in}}%
\pgfpathlineto{\pgfqpoint{3.906867in}{1.109200in}}%
\pgfpathlineto{\pgfqpoint{3.906440in}{1.118736in}}%
\pgfpathlineto{\pgfqpoint{3.901988in}{1.125461in}}%
\pgfpathlineto{\pgfqpoint{3.895279in}{1.148109in}}%
\pgfpathlineto{\pgfqpoint{3.887938in}{1.158956in}}%
\pgfpathlineto{\pgfqpoint{3.889533in}{1.177370in}}%
\pgfpathlineto{\pgfqpoint{3.898166in}{1.179223in}}%
\pgfpathlineto{\pgfqpoint{3.903787in}{1.189211in}}%
\pgfpathlineto{\pgfqpoint{3.900390in}{1.200043in}}%
\pgfpathlineto{\pgfqpoint{3.891209in}{1.201851in}}%
\pgfpathlineto{\pgfqpoint{3.886629in}{1.206918in}}%
\pgfpathlineto{\pgfqpoint{3.879211in}{1.225555in}}%
\pgfpathlineto{\pgfqpoint{3.882678in}{1.232248in}}%
\pgfpathlineto{\pgfqpoint{3.880219in}{1.244711in}}%
\pgfpathlineto{\pgfqpoint{3.885664in}{1.260852in}}%
\pgfpathlineto{\pgfqpoint{3.892043in}{1.254907in}}%
\pgfpathlineto{\pgfqpoint{3.889148in}{1.247890in}}%
\pgfpathlineto{\pgfqpoint{3.900202in}{1.236766in}}%
\pgfpathlineto{\pgfqpoint{3.899497in}{1.253285in}}%
\pgfpathlineto{\pgfqpoint{3.894760in}{1.257715in}}%
\pgfpathlineto{\pgfqpoint{3.900174in}{1.272240in}}%
\pgfpathlineto{\pgfqpoint{3.920652in}{1.269927in}}%
\pgfpathlineto{\pgfqpoint{3.917896in}{1.275328in}}%
\pgfpathlineto{\pgfqpoint{3.904381in}{1.274797in}}%
\pgfpathlineto{\pgfqpoint{3.897951in}{1.282977in}}%
\pgfpathlineto{\pgfqpoint{3.889858in}{1.275713in}}%
\pgfpathlineto{\pgfqpoint{3.885564in}{1.263570in}}%
\pgfpathlineto{\pgfqpoint{3.874102in}{1.279909in}}%
\pgfpathlineto{\pgfqpoint{3.869687in}{1.282884in}}%
\pgfpathlineto{\pgfqpoint{3.871369in}{1.300667in}}%
\pgfpathlineto{\pgfqpoint{3.867875in}{1.311153in}}%
\pgfpathlineto{\pgfqpoint{3.861542in}{1.321006in}}%
\pgfpathlineto{\pgfqpoint{3.848573in}{1.351193in}}%
\pgfpathlineto{\pgfqpoint{3.852813in}{1.357869in}}%
\pgfpathlineto{\pgfqpoint{3.852764in}{1.378965in}}%
\pgfpathlineto{\pgfqpoint{3.859763in}{1.390732in}}%
\pgfpathlineto{\pgfqpoint{3.861368in}{1.409121in}}%
\pgfpathlineto{\pgfqpoint{3.854768in}{1.430191in}}%
\pgfpathlineto{\pgfqpoint{3.846012in}{1.443602in}}%
\pgfpathlineto{\pgfqpoint{3.847574in}{1.455704in}}%
\pgfpathlineto{\pgfqpoint{3.872163in}{1.484994in}}%
\pgfpathlineto{\pgfqpoint{3.873384in}{1.494985in}}%
\pgfpathlineto{\pgfqpoint{3.884464in}{1.514044in}}%
\pgfpathlineto{\pgfqpoint{3.886001in}{1.532109in}}%
\pgfpathlineto{\pgfqpoint{3.882439in}{1.536719in}}%
\pgfpathlineto{\pgfqpoint{3.888713in}{1.549891in}}%
\pgfpathclose%
\pgfusepath{fill}%
\end{pgfscope}%
\begin{pgfscope}%
\pgfpathrectangle{\pgfqpoint{3.625000in}{0.100000in}}{\pgfqpoint{2.989028in}{1.913466in}}%
\pgfusepath{clip}%
\pgfsetbuttcap%
\pgfsetmiterjoin%
\definecolor{currentfill}{rgb}{0.847520,0.938639,0.607151}%
\pgfsetfillcolor{currentfill}%
\pgfsetlinewidth{0.000000pt}%
\definecolor{currentstroke}{rgb}{0.000000,0.000000,0.000000}%
\pgfsetstrokecolor{currentstroke}%
\pgfsetstrokeopacity{0.000000}%
\pgfsetdash{}{0pt}%
\pgfpathmoveto{\pgfqpoint{5.803860in}{1.162554in}}%
\pgfpathlineto{\pgfqpoint{5.795539in}{1.244011in}}%
\pgfpathlineto{\pgfqpoint{5.785872in}{1.331556in}}%
\pgfpathlineto{\pgfqpoint{5.849180in}{1.340984in}}%
\pgfpathlineto{\pgfqpoint{5.865982in}{1.336611in}}%
\pgfpathlineto{\pgfqpoint{5.874041in}{1.331864in}}%
\pgfpathlineto{\pgfqpoint{5.884139in}{1.333180in}}%
\pgfpathlineto{\pgfqpoint{5.897453in}{1.325321in}}%
\pgfpathlineto{\pgfqpoint{5.922252in}{1.337018in}}%
\pgfpathlineto{\pgfqpoint{5.935938in}{1.337409in}}%
\pgfpathlineto{\pgfqpoint{5.951920in}{1.355239in}}%
\pgfpathlineto{\pgfqpoint{5.968180in}{1.366084in}}%
\pgfpathlineto{\pgfqpoint{5.989869in}{1.378577in}}%
\pgfpathlineto{\pgfqpoint{6.003831in}{1.291252in}}%
\pgfpathlineto{\pgfqpoint{5.997344in}{1.285643in}}%
\pgfpathlineto{\pgfqpoint{6.001532in}{1.280492in}}%
\pgfpathlineto{\pgfqpoint{6.003201in}{1.269199in}}%
\pgfpathlineto{\pgfqpoint{5.999951in}{1.258593in}}%
\pgfpathlineto{\pgfqpoint{5.999521in}{1.243089in}}%
\pgfpathlineto{\pgfqpoint{5.996469in}{1.222911in}}%
\pgfpathlineto{\pgfqpoint{5.981028in}{1.204915in}}%
\pgfpathlineto{\pgfqpoint{5.974528in}{1.201094in}}%
\pgfpathlineto{\pgfqpoint{5.969473in}{1.204383in}}%
\pgfpathlineto{\pgfqpoint{5.956963in}{1.187230in}}%
\pgfpathlineto{\pgfqpoint{5.958135in}{1.170714in}}%
\pgfpathlineto{\pgfqpoint{5.944203in}{1.174686in}}%
\pgfpathlineto{\pgfqpoint{5.937606in}{1.158191in}}%
\pgfpathlineto{\pgfqpoint{5.940941in}{1.146502in}}%
\pgfpathlineto{\pgfqpoint{5.935720in}{1.144776in}}%
\pgfpathlineto{\pgfqpoint{5.934990in}{1.135540in}}%
\pgfpathlineto{\pgfqpoint{5.922176in}{1.131815in}}%
\pgfpathlineto{\pgfqpoint{5.915518in}{1.139272in}}%
\pgfpathlineto{\pgfqpoint{5.906950in}{1.142163in}}%
\pgfpathlineto{\pgfqpoint{5.904937in}{1.149724in}}%
\pgfpathlineto{\pgfqpoint{5.897191in}{1.148394in}}%
\pgfpathlineto{\pgfqpoint{5.892116in}{1.141440in}}%
\pgfpathlineto{\pgfqpoint{5.884855in}{1.138994in}}%
\pgfpathlineto{\pgfqpoint{5.872031in}{1.143916in}}%
\pgfpathlineto{\pgfqpoint{5.864938in}{1.138060in}}%
\pgfpathlineto{\pgfqpoint{5.854857in}{1.144988in}}%
\pgfpathlineto{\pgfqpoint{5.835517in}{1.147114in}}%
\pgfpathlineto{\pgfqpoint{5.829684in}{1.159700in}}%
\pgfpathlineto{\pgfqpoint{5.822293in}{1.165276in}}%
\pgfpathlineto{\pgfqpoint{5.815128in}{1.161707in}}%
\pgfpathlineto{\pgfqpoint{5.803860in}{1.162554in}}%
\pgfpathclose%
\pgfusepath{fill}%
\end{pgfscope}%
\begin{pgfscope}%
\pgfpathrectangle{\pgfqpoint{3.625000in}{0.100000in}}{\pgfqpoint{2.989028in}{1.913466in}}%
\pgfusepath{clip}%
\pgfsetbuttcap%
\pgfsetmiterjoin%
\definecolor{currentfill}{rgb}{0.829373,0.931257,0.610842}%
\pgfsetfillcolor{currentfill}%
\pgfsetlinewidth{0.000000pt}%
\definecolor{currentstroke}{rgb}{0.000000,0.000000,0.000000}%
\pgfsetstrokecolor{currentstroke}%
\pgfsetstrokeopacity{0.000000}%
\pgfsetdash{}{0pt}%
\pgfpathmoveto{\pgfqpoint{5.594266in}{1.005020in}}%
\pgfpathlineto{\pgfqpoint{5.586849in}{1.011045in}}%
\pgfpathlineto{\pgfqpoint{5.580801in}{1.008178in}}%
\pgfpathlineto{\pgfqpoint{5.573069in}{1.022616in}}%
\pgfpathlineto{\pgfqpoint{5.577019in}{1.031690in}}%
\pgfpathlineto{\pgfqpoint{5.571326in}{1.041904in}}%
\pgfpathlineto{\pgfqpoint{5.571020in}{1.049946in}}%
\pgfpathlineto{\pgfqpoint{5.563328in}{1.052811in}}%
\pgfpathlineto{\pgfqpoint{5.559760in}{1.058870in}}%
\pgfpathlineto{\pgfqpoint{5.544719in}{1.066437in}}%
\pgfpathlineto{\pgfqpoint{5.531657in}{1.075760in}}%
\pgfpathlineto{\pgfqpoint{5.525587in}{1.082799in}}%
\pgfpathlineto{\pgfqpoint{5.525462in}{1.091383in}}%
\pgfpathlineto{\pgfqpoint{5.533584in}{1.107865in}}%
\pgfpathlineto{\pgfqpoint{5.536082in}{1.120472in}}%
\pgfpathlineto{\pgfqpoint{5.529465in}{1.127595in}}%
\pgfpathlineto{\pgfqpoint{5.520702in}{1.130279in}}%
\pgfpathlineto{\pgfqpoint{5.510070in}{1.124410in}}%
\pgfpathlineto{\pgfqpoint{5.505570in}{1.134486in}}%
\pgfpathlineto{\pgfqpoint{5.503299in}{1.148144in}}%
\pgfpathlineto{\pgfqpoint{5.487630in}{1.160319in}}%
\pgfpathlineto{\pgfqpoint{5.484499in}{1.165722in}}%
\pgfpathlineto{\pgfqpoint{5.470188in}{1.177954in}}%
\pgfpathlineto{\pgfqpoint{5.465703in}{1.186846in}}%
\pgfpathlineto{\pgfqpoint{5.461659in}{1.204480in}}%
\pgfpathlineto{\pgfqpoint{5.464431in}{1.220143in}}%
\pgfpathlineto{\pgfqpoint{5.468130in}{1.222326in}}%
\pgfpathlineto{\pgfqpoint{5.467465in}{1.235438in}}%
\pgfpathlineto{\pgfqpoint{5.477907in}{1.239332in}}%
\pgfpathlineto{\pgfqpoint{5.481060in}{1.251098in}}%
\pgfpathlineto{\pgfqpoint{5.487068in}{1.259018in}}%
\pgfpathlineto{\pgfqpoint{5.486764in}{1.269093in}}%
\pgfpathlineto{\pgfqpoint{5.479318in}{1.277106in}}%
\pgfpathlineto{\pgfqpoint{5.481087in}{1.288322in}}%
\pgfpathlineto{\pgfqpoint{5.500389in}{1.293200in}}%
\pgfpathlineto{\pgfqpoint{5.515197in}{1.302104in}}%
\pgfpathlineto{\pgfqpoint{5.516751in}{1.313315in}}%
\pgfpathlineto{\pgfqpoint{5.521904in}{1.316845in}}%
\pgfpathlineto{\pgfqpoint{5.523875in}{1.328623in}}%
\pgfpathlineto{\pgfqpoint{5.522291in}{1.336402in}}%
\pgfpathlineto{\pgfqpoint{5.512167in}{1.342864in}}%
\pgfpathlineto{\pgfqpoint{5.508091in}{1.352497in}}%
\pgfpathlineto{\pgfqpoint{5.498100in}{1.361809in}}%
\pgfpathlineto{\pgfqpoint{5.580206in}{1.365308in}}%
\pgfpathlineto{\pgfqpoint{5.635291in}{1.369221in}}%
\pgfpathlineto{\pgfqpoint{5.634283in}{1.357634in}}%
\pgfpathlineto{\pgfqpoint{5.643671in}{1.341651in}}%
\pgfpathlineto{\pgfqpoint{5.647612in}{1.327996in}}%
\pgfpathlineto{\pgfqpoint{5.652306in}{1.320243in}}%
\pgfpathlineto{\pgfqpoint{5.657431in}{1.256358in}}%
\pgfpathlineto{\pgfqpoint{5.664547in}{1.165295in}}%
\pgfpathlineto{\pgfqpoint{5.659057in}{1.151457in}}%
\pgfpathlineto{\pgfqpoint{5.666966in}{1.140083in}}%
\pgfpathlineto{\pgfqpoint{5.669170in}{1.129219in}}%
\pgfpathlineto{\pgfqpoint{5.658263in}{1.106692in}}%
\pgfpathlineto{\pgfqpoint{5.658907in}{1.103580in}}%
\pgfpathlineto{\pgfqpoint{5.647681in}{1.090668in}}%
\pgfpathlineto{\pgfqpoint{5.650759in}{1.086724in}}%
\pgfpathlineto{\pgfqpoint{5.646088in}{1.078676in}}%
\pgfpathlineto{\pgfqpoint{5.647290in}{1.062454in}}%
\pgfpathlineto{\pgfqpoint{5.641617in}{1.052499in}}%
\pgfpathlineto{\pgfqpoint{5.646231in}{1.040736in}}%
\pgfpathlineto{\pgfqpoint{5.626900in}{1.034344in}}%
\pgfpathlineto{\pgfqpoint{5.625119in}{1.027393in}}%
\pgfpathlineto{\pgfqpoint{5.630397in}{1.018582in}}%
\pgfpathlineto{\pgfqpoint{5.627969in}{1.012840in}}%
\pgfpathlineto{\pgfqpoint{5.607240in}{1.019933in}}%
\pgfpathlineto{\pgfqpoint{5.600401in}{1.020647in}}%
\pgfpathlineto{\pgfqpoint{5.591874in}{1.009863in}}%
\pgfpathlineto{\pgfqpoint{5.594266in}{1.005020in}}%
\pgfpathclose%
\pgfusepath{fill}%
\end{pgfscope}%
\begin{pgfscope}%
\pgfpathrectangle{\pgfqpoint{3.625000in}{0.100000in}}{\pgfqpoint{2.989028in}{1.913466in}}%
\pgfusepath{clip}%
\pgfsetbuttcap%
\pgfsetmiterjoin%
\definecolor{currentfill}{rgb}{0.901961,0.960784,0.596078}%
\pgfsetfillcolor{currentfill}%
\pgfsetlinewidth{0.000000pt}%
\definecolor{currentstroke}{rgb}{0.000000,0.000000,0.000000}%
\pgfsetstrokecolor{currentstroke}%
\pgfsetstrokeopacity{0.000000}%
\pgfsetdash{}{0pt}%
\pgfpathmoveto{\pgfqpoint{6.197645in}{1.203210in}}%
\pgfpathlineto{\pgfqpoint{6.191941in}{1.211681in}}%
\pgfpathlineto{\pgfqpoint{6.195160in}{1.216421in}}%
\pgfpathlineto{\pgfqpoint{6.203043in}{1.211098in}}%
\pgfpathlineto{\pgfqpoint{6.197645in}{1.203210in}}%
\pgfpathclose%
\pgfusepath{fill}%
\end{pgfscope}%
\begin{pgfscope}%
\pgfpathrectangle{\pgfqpoint{3.625000in}{0.100000in}}{\pgfqpoint{2.989028in}{1.913466in}}%
\pgfusepath{clip}%
\pgfsetbuttcap%
\pgfsetmiterjoin%
\definecolor{currentfill}{rgb}{0.997616,0.926105,0.625067}%
\pgfsetfillcolor{currentfill}%
\pgfsetlinewidth{0.000000pt}%
\definecolor{currentstroke}{rgb}{0.000000,0.000000,0.000000}%
\pgfsetstrokecolor{currentstroke}%
\pgfsetstrokeopacity{0.000000}%
\pgfsetdash{}{0pt}%
\pgfpathmoveto{\pgfqpoint{6.247458in}{1.276100in}}%
\pgfpathlineto{\pgfqpoint{6.250867in}{1.283302in}}%
\pgfpathlineto{\pgfqpoint{6.264648in}{1.284858in}}%
\pgfpathlineto{\pgfqpoint{6.262440in}{1.278737in}}%
\pgfpathlineto{\pgfqpoint{6.257893in}{1.270918in}}%
\pgfpathlineto{\pgfqpoint{6.260974in}{1.261602in}}%
\pgfpathlineto{\pgfqpoint{6.273133in}{1.250439in}}%
\pgfpathlineto{\pgfqpoint{6.275955in}{1.238685in}}%
\pgfpathlineto{\pgfqpoint{6.289963in}{1.224047in}}%
\pgfpathlineto{\pgfqpoint{6.295458in}{1.224674in}}%
\pgfpathlineto{\pgfqpoint{6.302304in}{1.202700in}}%
\pgfpathlineto{\pgfqpoint{6.301181in}{1.202481in}}%
\pgfpathlineto{\pgfqpoint{6.299933in}{1.202236in}}%
\pgfpathlineto{\pgfqpoint{6.269401in}{1.196391in}}%
\pgfpathlineto{\pgfqpoint{6.253066in}{1.254485in}}%
\pgfpathlineto{\pgfqpoint{6.247458in}{1.276100in}}%
\pgfpathclose%
\pgfusepath{fill}%
\end{pgfscope}%
\begin{pgfscope}%
\pgfpathrectangle{\pgfqpoint{3.625000in}{0.100000in}}{\pgfqpoint{2.989028in}{1.913466in}}%
\pgfusepath{clip}%
\pgfsetbuttcap%
\pgfsetmiterjoin%
\definecolor{currentfill}{rgb}{0.368012,0.725106,0.661822}%
\pgfsetfillcolor{currentfill}%
\pgfsetlinewidth{0.000000pt}%
\definecolor{currentstroke}{rgb}{0.000000,0.000000,0.000000}%
\pgfsetstrokecolor{currentstroke}%
\pgfsetstrokeopacity{0.000000}%
\pgfsetdash{}{0pt}%
\pgfpathmoveto{\pgfqpoint{5.962284in}{1.079184in}}%
\pgfpathlineto{\pgfqpoint{5.952636in}{1.079539in}}%
\pgfpathlineto{\pgfqpoint{5.943746in}{1.085662in}}%
\pgfpathlineto{\pgfqpoint{5.931736in}{1.104274in}}%
\pgfpathlineto{\pgfqpoint{5.921513in}{1.114131in}}%
\pgfpathlineto{\pgfqpoint{5.924185in}{1.121744in}}%
\pgfpathlineto{\pgfqpoint{5.922176in}{1.131815in}}%
\pgfpathlineto{\pgfqpoint{5.934990in}{1.135540in}}%
\pgfpathlineto{\pgfqpoint{5.935720in}{1.144776in}}%
\pgfpathlineto{\pgfqpoint{5.940941in}{1.146502in}}%
\pgfpathlineto{\pgfqpoint{5.937606in}{1.158191in}}%
\pgfpathlineto{\pgfqpoint{5.944203in}{1.174686in}}%
\pgfpathlineto{\pgfqpoint{5.958135in}{1.170714in}}%
\pgfpathlineto{\pgfqpoint{5.956963in}{1.187230in}}%
\pgfpathlineto{\pgfqpoint{5.969473in}{1.204383in}}%
\pgfpathlineto{\pgfqpoint{5.974528in}{1.201094in}}%
\pgfpathlineto{\pgfqpoint{5.981028in}{1.204915in}}%
\pgfpathlineto{\pgfqpoint{5.996469in}{1.222911in}}%
\pgfpathlineto{\pgfqpoint{5.999521in}{1.243089in}}%
\pgfpathlineto{\pgfqpoint{5.999951in}{1.258593in}}%
\pgfpathlineto{\pgfqpoint{6.003201in}{1.269199in}}%
\pgfpathlineto{\pgfqpoint{6.001532in}{1.280492in}}%
\pgfpathlineto{\pgfqpoint{5.997344in}{1.285643in}}%
\pgfpathlineto{\pgfqpoint{6.003831in}{1.291252in}}%
\pgfpathlineto{\pgfqpoint{6.013234in}{1.232010in}}%
\pgfpathlineto{\pgfqpoint{6.065072in}{1.240581in}}%
\pgfpathlineto{\pgfqpoint{6.070445in}{1.206768in}}%
\pgfpathlineto{\pgfqpoint{6.089213in}{1.229086in}}%
\pgfpathlineto{\pgfqpoint{6.093635in}{1.226861in}}%
\pgfpathlineto{\pgfqpoint{6.100140in}{1.239793in}}%
\pgfpathlineto{\pgfqpoint{6.109152in}{1.235738in}}%
\pgfpathlineto{\pgfqpoint{6.117010in}{1.236509in}}%
\pgfpathlineto{\pgfqpoint{6.122194in}{1.245499in}}%
\pgfpathlineto{\pgfqpoint{6.129719in}{1.250500in}}%
\pgfpathlineto{\pgfqpoint{6.140180in}{1.246096in}}%
\pgfpathlineto{\pgfqpoint{6.147064in}{1.247618in}}%
\pgfpathlineto{\pgfqpoint{6.156926in}{1.230550in}}%
\pgfpathlineto{\pgfqpoint{6.154074in}{1.217629in}}%
\pgfpathlineto{\pgfqpoint{6.128317in}{1.232177in}}%
\pgfpathlineto{\pgfqpoint{6.123492in}{1.220234in}}%
\pgfpathlineto{\pgfqpoint{6.125084in}{1.214739in}}%
\pgfpathlineto{\pgfqpoint{6.114121in}{1.196164in}}%
\pgfpathlineto{\pgfqpoint{6.108947in}{1.192470in}}%
\pgfpathlineto{\pgfqpoint{6.106611in}{1.184279in}}%
\pgfpathlineto{\pgfqpoint{6.097791in}{1.185136in}}%
\pgfpathlineto{\pgfqpoint{6.091456in}{1.162779in}}%
\pgfpathlineto{\pgfqpoint{6.087911in}{1.157648in}}%
\pgfpathlineto{\pgfqpoint{6.078799in}{1.159359in}}%
\pgfpathlineto{\pgfqpoint{6.069480in}{1.166411in}}%
\pgfpathlineto{\pgfqpoint{6.069157in}{1.155588in}}%
\pgfpathlineto{\pgfqpoint{6.060135in}{1.137399in}}%
\pgfpathlineto{\pgfqpoint{6.057891in}{1.124451in}}%
\pgfpathlineto{\pgfqpoint{6.050969in}{1.115827in}}%
\pgfpathlineto{\pgfqpoint{6.045721in}{1.102055in}}%
\pgfpathlineto{\pgfqpoint{6.045463in}{1.089306in}}%
\pgfpathlineto{\pgfqpoint{6.037358in}{1.086925in}}%
\pgfpathlineto{\pgfqpoint{6.028164in}{1.079708in}}%
\pgfpathlineto{\pgfqpoint{6.019311in}{1.078334in}}%
\pgfpathlineto{\pgfqpoint{6.017271in}{1.072225in}}%
\pgfpathlineto{\pgfqpoint{6.003007in}{1.065970in}}%
\pgfpathlineto{\pgfqpoint{5.995035in}{1.071300in}}%
\pgfpathlineto{\pgfqpoint{5.986126in}{1.061179in}}%
\pgfpathlineto{\pgfqpoint{5.970810in}{1.064162in}}%
\pgfpathlineto{\pgfqpoint{5.961416in}{1.074785in}}%
\pgfpathlineto{\pgfqpoint{5.962284in}{1.079184in}}%
\pgfpathclose%
\pgfusepath{fill}%
\end{pgfscope}%
\begin{pgfscope}%
\pgfpathrectangle{\pgfqpoint{3.625000in}{0.100000in}}{\pgfqpoint{2.989028in}{1.913466in}}%
\pgfusepath{clip}%
\pgfsetbuttcap%
\pgfsetmiterjoin%
\definecolor{currentfill}{rgb}{0.829373,0.931257,0.610842}%
\pgfsetfillcolor{currentfill}%
\pgfsetlinewidth{0.000000pt}%
\definecolor{currentstroke}{rgb}{0.000000,0.000000,0.000000}%
\pgfsetstrokecolor{currentstroke}%
\pgfsetstrokeopacity{0.000000}%
\pgfsetdash{}{0pt}%
\pgfpathmoveto{\pgfqpoint{6.299933in}{1.202236in}}%
\pgfpathlineto{\pgfqpoint{6.299542in}{1.190298in}}%
\pgfpathlineto{\pgfqpoint{6.294947in}{1.184429in}}%
\pgfpathlineto{\pgfqpoint{6.292015in}{1.171401in}}%
\pgfpathlineto{\pgfqpoint{6.278775in}{1.165395in}}%
\pgfpathlineto{\pgfqpoint{6.267667in}{1.163645in}}%
\pgfpathlineto{\pgfqpoint{6.270904in}{1.172214in}}%
\pgfpathlineto{\pgfqpoint{6.253800in}{1.179393in}}%
\pgfpathlineto{\pgfqpoint{6.239904in}{1.188383in}}%
\pgfpathlineto{\pgfqpoint{6.248507in}{1.201821in}}%
\pgfpathlineto{\pgfqpoint{6.241613in}{1.212266in}}%
\pgfpathlineto{\pgfqpoint{6.242424in}{1.221276in}}%
\pgfpathlineto{\pgfqpoint{6.237439in}{1.223749in}}%
\pgfpathlineto{\pgfqpoint{6.233394in}{1.238386in}}%
\pgfpathlineto{\pgfqpoint{6.242960in}{1.258229in}}%
\pgfpathlineto{\pgfqpoint{6.235773in}{1.261469in}}%
\pgfpathlineto{\pgfqpoint{6.233910in}{1.251680in}}%
\pgfpathlineto{\pgfqpoint{6.223651in}{1.248954in}}%
\pgfpathlineto{\pgfqpoint{6.224730in}{1.216746in}}%
\pgfpathlineto{\pgfqpoint{6.222854in}{1.206381in}}%
\pgfpathlineto{\pgfqpoint{6.228000in}{1.191607in}}%
\pgfpathlineto{\pgfqpoint{6.235916in}{1.184499in}}%
\pgfpathlineto{\pgfqpoint{6.232326in}{1.180036in}}%
\pgfpathlineto{\pgfqpoint{6.240407in}{1.173514in}}%
\pgfpathlineto{\pgfqpoint{6.243328in}{1.162929in}}%
\pgfpathlineto{\pgfqpoint{6.228517in}{1.171694in}}%
\pgfpathlineto{\pgfqpoint{6.219143in}{1.170556in}}%
\pgfpathlineto{\pgfqpoint{6.211865in}{1.179573in}}%
\pgfpathlineto{\pgfqpoint{6.204454in}{1.180451in}}%
\pgfpathlineto{\pgfqpoint{6.193921in}{1.175935in}}%
\pgfpathlineto{\pgfqpoint{6.189837in}{1.181563in}}%
\pgfpathlineto{\pgfqpoint{6.195205in}{1.193376in}}%
\pgfpathlineto{\pgfqpoint{6.197645in}{1.203210in}}%
\pgfpathlineto{\pgfqpoint{6.203043in}{1.211098in}}%
\pgfpathlineto{\pgfqpoint{6.195160in}{1.216421in}}%
\pgfpathlineto{\pgfqpoint{6.191941in}{1.211681in}}%
\pgfpathlineto{\pgfqpoint{6.184066in}{1.216492in}}%
\pgfpathlineto{\pgfqpoint{6.170120in}{1.219698in}}%
\pgfpathlineto{\pgfqpoint{6.171385in}{1.226747in}}%
\pgfpathlineto{\pgfqpoint{6.165064in}{1.230836in}}%
\pgfpathlineto{\pgfqpoint{6.156926in}{1.230550in}}%
\pgfpathlineto{\pgfqpoint{6.147064in}{1.247618in}}%
\pgfpathlineto{\pgfqpoint{6.140180in}{1.246096in}}%
\pgfpathlineto{\pgfqpoint{6.129719in}{1.250500in}}%
\pgfpathlineto{\pgfqpoint{6.122194in}{1.245499in}}%
\pgfpathlineto{\pgfqpoint{6.117010in}{1.236509in}}%
\pgfpathlineto{\pgfqpoint{6.109152in}{1.235738in}}%
\pgfpathlineto{\pgfqpoint{6.100140in}{1.239793in}}%
\pgfpathlineto{\pgfqpoint{6.093635in}{1.226861in}}%
\pgfpathlineto{\pgfqpoint{6.089213in}{1.229086in}}%
\pgfpathlineto{\pgfqpoint{6.070445in}{1.206768in}}%
\pgfpathlineto{\pgfqpoint{6.065072in}{1.240581in}}%
\pgfpathlineto{\pgfqpoint{6.133655in}{1.253292in}}%
\pgfpathlineto{\pgfqpoint{6.164399in}{1.258797in}}%
\pgfpathlineto{\pgfqpoint{6.209127in}{1.267806in}}%
\pgfpathlineto{\pgfqpoint{6.247458in}{1.276100in}}%
\pgfpathlineto{\pgfqpoint{6.253066in}{1.254485in}}%
\pgfpathlineto{\pgfqpoint{6.269401in}{1.196391in}}%
\pgfpathlineto{\pgfqpoint{6.299933in}{1.202236in}}%
\pgfpathclose%
\pgfusepath{fill}%
\end{pgfscope}%
\begin{pgfscope}%
\pgfpathrectangle{\pgfqpoint{3.625000in}{0.100000in}}{\pgfqpoint{2.989028in}{1.913466in}}%
\pgfusepath{clip}%
\pgfsetbuttcap%
\pgfsetmiterjoin%
\definecolor{currentfill}{rgb}{0.633449,0.852134,0.643676}%
\pgfsetfillcolor{currentfill}%
\pgfsetlinewidth{0.000000pt}%
\definecolor{currentstroke}{rgb}{0.000000,0.000000,0.000000}%
\pgfsetstrokecolor{currentstroke}%
\pgfsetstrokeopacity{0.000000}%
\pgfsetdash{}{0pt}%
\pgfpathmoveto{\pgfqpoint{4.919404in}{1.009408in}}%
\pgfpathlineto{\pgfqpoint{4.869368in}{1.014183in}}%
\pgfpathlineto{\pgfqpoint{4.817460in}{1.018783in}}%
\pgfpathlineto{\pgfqpoint{4.757477in}{1.025040in}}%
\pgfpathlineto{\pgfqpoint{4.668359in}{1.035639in}}%
\pgfpathlineto{\pgfqpoint{4.636744in}{1.040510in}}%
\pgfpathlineto{\pgfqpoint{4.555137in}{1.052362in}}%
\pgfpathlineto{\pgfqpoint{4.567053in}{1.127151in}}%
\pgfpathlineto{\pgfqpoint{4.567350in}{1.133200in}}%
\pgfpathlineto{\pgfqpoint{4.578829in}{1.205365in}}%
\pgfpathlineto{\pgfqpoint{4.587404in}{1.260326in}}%
\pgfpathlineto{\pgfqpoint{4.595495in}{1.311293in}}%
\pgfpathlineto{\pgfqpoint{4.650778in}{1.303357in}}%
\pgfpathlineto{\pgfqpoint{4.702321in}{1.295959in}}%
\pgfpathlineto{\pgfqpoint{4.797072in}{1.284265in}}%
\pgfpathlineto{\pgfqpoint{4.840540in}{1.280295in}}%
\pgfpathlineto{\pgfqpoint{4.909404in}{1.273614in}}%
\pgfpathlineto{\pgfqpoint{4.939195in}{1.271196in}}%
\pgfpathlineto{\pgfqpoint{4.933935in}{1.205970in}}%
\pgfpathlineto{\pgfqpoint{4.929183in}{1.143169in}}%
\pgfpathlineto{\pgfqpoint{4.925359in}{1.092033in}}%
\pgfpathlineto{\pgfqpoint{4.919404in}{1.009408in}}%
\pgfpathclose%
\pgfusepath{fill}%
\end{pgfscope}%
\begin{pgfscope}%
\pgfpathrectangle{\pgfqpoint{3.625000in}{0.100000in}}{\pgfqpoint{2.989028in}{1.913466in}}%
\pgfusepath{clip}%
\pgfsetbuttcap%
\pgfsetmiterjoin%
\definecolor{currentfill}{rgb}{0.569781,0.827220,0.644598}%
\pgfsetfillcolor{currentfill}%
\pgfsetlinewidth{0.000000pt}%
\definecolor{currentstroke}{rgb}{0.000000,0.000000,0.000000}%
\pgfsetstrokecolor{currentstroke}%
\pgfsetstrokeopacity{0.000000}%
\pgfsetdash{}{0pt}%
\pgfpathmoveto{\pgfqpoint{5.581343in}{0.972329in}}%
\pgfpathlineto{\pgfqpoint{5.583030in}{0.979890in}}%
\pgfpathlineto{\pgfqpoint{5.591773in}{0.978188in}}%
\pgfpathlineto{\pgfqpoint{5.594266in}{1.005020in}}%
\pgfpathlineto{\pgfqpoint{5.591874in}{1.009863in}}%
\pgfpathlineto{\pgfqpoint{5.600401in}{1.020647in}}%
\pgfpathlineto{\pgfqpoint{5.607240in}{1.019933in}}%
\pgfpathlineto{\pgfqpoint{5.627969in}{1.012840in}}%
\pgfpathlineto{\pgfqpoint{5.630397in}{1.018582in}}%
\pgfpathlineto{\pgfqpoint{5.625119in}{1.027393in}}%
\pgfpathlineto{\pgfqpoint{5.626900in}{1.034344in}}%
\pgfpathlineto{\pgfqpoint{5.646231in}{1.040736in}}%
\pgfpathlineto{\pgfqpoint{5.641617in}{1.052499in}}%
\pgfpathlineto{\pgfqpoint{5.647290in}{1.062454in}}%
\pgfpathlineto{\pgfqpoint{5.657850in}{1.068041in}}%
\pgfpathlineto{\pgfqpoint{5.680006in}{1.073574in}}%
\pgfpathlineto{\pgfqpoint{5.693985in}{1.065222in}}%
\pgfpathlineto{\pgfqpoint{5.698226in}{1.073367in}}%
\pgfpathlineto{\pgfqpoint{5.709834in}{1.079056in}}%
\pgfpathlineto{\pgfqpoint{5.712829in}{1.074039in}}%
\pgfpathlineto{\pgfqpoint{5.724753in}{1.078025in}}%
\pgfpathlineto{\pgfqpoint{5.724000in}{1.084848in}}%
\pgfpathlineto{\pgfqpoint{5.734721in}{1.092656in}}%
\pgfpathlineto{\pgfqpoint{5.741048in}{1.084548in}}%
\pgfpathlineto{\pgfqpoint{5.749371in}{1.083732in}}%
\pgfpathlineto{\pgfqpoint{5.754903in}{1.089016in}}%
\pgfpathlineto{\pgfqpoint{5.754287in}{1.096531in}}%
\pgfpathlineto{\pgfqpoint{5.758989in}{1.103993in}}%
\pgfpathlineto{\pgfqpoint{5.765290in}{1.105610in}}%
\pgfpathlineto{\pgfqpoint{5.767820in}{1.115463in}}%
\pgfpathlineto{\pgfqpoint{5.776983in}{1.123981in}}%
\pgfpathlineto{\pgfqpoint{5.774225in}{1.132459in}}%
\pgfpathlineto{\pgfqpoint{5.783139in}{1.136680in}}%
\pgfpathlineto{\pgfqpoint{5.789096in}{1.134084in}}%
\pgfpathlineto{\pgfqpoint{5.797893in}{1.140679in}}%
\pgfpathlineto{\pgfqpoint{5.805755in}{1.142397in}}%
\pgfpathlineto{\pgfqpoint{5.805783in}{1.149238in}}%
\pgfpathlineto{\pgfqpoint{5.800273in}{1.158730in}}%
\pgfpathlineto{\pgfqpoint{5.803860in}{1.162554in}}%
\pgfpathlineto{\pgfqpoint{5.815128in}{1.161707in}}%
\pgfpathlineto{\pgfqpoint{5.822293in}{1.165276in}}%
\pgfpathlineto{\pgfqpoint{5.829684in}{1.159700in}}%
\pgfpathlineto{\pgfqpoint{5.835517in}{1.147114in}}%
\pgfpathlineto{\pgfqpoint{5.854857in}{1.144988in}}%
\pgfpathlineto{\pgfqpoint{5.864938in}{1.138060in}}%
\pgfpathlineto{\pgfqpoint{5.872031in}{1.143916in}}%
\pgfpathlineto{\pgfqpoint{5.884855in}{1.138994in}}%
\pgfpathlineto{\pgfqpoint{5.892116in}{1.141440in}}%
\pgfpathlineto{\pgfqpoint{5.897191in}{1.148394in}}%
\pgfpathlineto{\pgfqpoint{5.904937in}{1.149724in}}%
\pgfpathlineto{\pgfqpoint{5.906950in}{1.142163in}}%
\pgfpathlineto{\pgfqpoint{5.915518in}{1.139272in}}%
\pgfpathlineto{\pgfqpoint{5.922176in}{1.131815in}}%
\pgfpathlineto{\pgfqpoint{5.924185in}{1.121744in}}%
\pgfpathlineto{\pgfqpoint{5.921513in}{1.114131in}}%
\pgfpathlineto{\pgfqpoint{5.931736in}{1.104274in}}%
\pgfpathlineto{\pgfqpoint{5.943746in}{1.085662in}}%
\pgfpathlineto{\pgfqpoint{5.952636in}{1.079539in}}%
\pgfpathlineto{\pgfqpoint{5.962284in}{1.079184in}}%
\pgfpathlineto{\pgfqpoint{5.944469in}{1.058743in}}%
\pgfpathlineto{\pgfqpoint{5.926934in}{1.046371in}}%
\pgfpathlineto{\pgfqpoint{5.920481in}{1.036545in}}%
\pgfpathlineto{\pgfqpoint{5.920590in}{1.031213in}}%
\pgfpathlineto{\pgfqpoint{5.911092in}{1.027126in}}%
\pgfpathlineto{\pgfqpoint{5.908370in}{1.019436in}}%
\pgfpathlineto{\pgfqpoint{5.888584in}{1.011690in}}%
\pgfpathlineto{\pgfqpoint{5.881571in}{1.006660in}}%
\pgfpathlineto{\pgfqpoint{5.880622in}{1.005586in}}%
\pgfpathlineto{\pgfqpoint{5.823708in}{1.000190in}}%
\pgfpathlineto{\pgfqpoint{5.789337in}{0.997299in}}%
\pgfpathlineto{\pgfqpoint{5.732938in}{0.994107in}}%
\pgfpathlineto{\pgfqpoint{5.662674in}{0.987156in}}%
\pgfpathlineto{\pgfqpoint{5.651065in}{0.988764in}}%
\pgfpathlineto{\pgfqpoint{5.653480in}{0.976914in}}%
\pgfpathlineto{\pgfqpoint{5.581343in}{0.972329in}}%
\pgfpathclose%
\pgfusepath{fill}%
\end{pgfscope}%
\begin{pgfscope}%
\pgfpathrectangle{\pgfqpoint{3.625000in}{0.100000in}}{\pgfqpoint{2.989028in}{1.913466in}}%
\pgfusepath{clip}%
\pgfsetbuttcap%
\pgfsetmiterjoin%
\definecolor{currentfill}{rgb}{0.280046,0.626990,0.702422}%
\pgfsetfillcolor{currentfill}%
\pgfsetlinewidth{0.000000pt}%
\definecolor{currentstroke}{rgb}{0.000000,0.000000,0.000000}%
\pgfsetstrokecolor{currentstroke}%
\pgfsetstrokeopacity{0.000000}%
\pgfsetdash{}{0pt}%
\pgfpathmoveto{\pgfqpoint{4.919404in}{1.009408in}}%
\pgfpathlineto{\pgfqpoint{4.925359in}{1.092033in}}%
\pgfpathlineto{\pgfqpoint{4.929183in}{1.143169in}}%
\pgfpathlineto{\pgfqpoint{4.933935in}{1.205970in}}%
\pgfpathlineto{\pgfqpoint{4.999809in}{1.201345in}}%
\pgfpathlineto{\pgfqpoint{5.083468in}{1.196775in}}%
\pgfpathlineto{\pgfqpoint{5.140367in}{1.194592in}}%
\pgfpathlineto{\pgfqpoint{5.196933in}{1.192859in}}%
\pgfpathlineto{\pgfqpoint{5.248131in}{1.192066in}}%
\pgfpathlineto{\pgfqpoint{5.271810in}{1.192312in}}%
\pgfpathlineto{\pgfqpoint{5.282233in}{1.183806in}}%
\pgfpathlineto{\pgfqpoint{5.290397in}{1.185521in}}%
\pgfpathlineto{\pgfqpoint{5.290653in}{1.178102in}}%
\pgfpathlineto{\pgfqpoint{5.284590in}{1.165269in}}%
\pgfpathlineto{\pgfqpoint{5.285247in}{1.157161in}}%
\pgfpathlineto{\pgfqpoint{5.291052in}{1.151814in}}%
\pgfpathlineto{\pgfqpoint{5.296384in}{1.140674in}}%
\pgfpathlineto{\pgfqpoint{5.307201in}{1.134279in}}%
\pgfpathlineto{\pgfqpoint{5.306840in}{1.092287in}}%
\pgfpathlineto{\pgfqpoint{5.307158in}{0.995725in}}%
\pgfpathlineto{\pgfqpoint{5.234656in}{0.996060in}}%
\pgfpathlineto{\pgfqpoint{5.158309in}{0.997396in}}%
\pgfpathlineto{\pgfqpoint{5.102105in}{0.999379in}}%
\pgfpathlineto{\pgfqpoint{5.055191in}{1.001189in}}%
\pgfpathlineto{\pgfqpoint{4.976155in}{1.005849in}}%
\pgfpathlineto{\pgfqpoint{4.919404in}{1.009408in}}%
\pgfpathclose%
\pgfusepath{fill}%
\end{pgfscope}%
\begin{pgfscope}%
\pgfpathrectangle{\pgfqpoint{3.625000in}{0.100000in}}{\pgfqpoint{2.989028in}{1.913466in}}%
\pgfusepath{clip}%
\pgfsetbuttcap%
\pgfsetmiterjoin%
\definecolor{currentfill}{rgb}{0.802153,0.920185,0.616378}%
\pgfsetfillcolor{currentfill}%
\pgfsetlinewidth{0.000000pt}%
\definecolor{currentstroke}{rgb}{0.000000,0.000000,0.000000}%
\pgfsetstrokecolor{currentstroke}%
\pgfsetstrokeopacity{0.000000}%
\pgfsetdash{}{0pt}%
\pgfpathmoveto{\pgfqpoint{5.986189in}{1.020348in}}%
\pgfpathlineto{\pgfqpoint{5.905888in}{1.009068in}}%
\pgfpathlineto{\pgfqpoint{5.881571in}{1.006660in}}%
\pgfpathlineto{\pgfqpoint{5.888584in}{1.011690in}}%
\pgfpathlineto{\pgfqpoint{5.908370in}{1.019436in}}%
\pgfpathlineto{\pgfqpoint{5.911092in}{1.027126in}}%
\pgfpathlineto{\pgfqpoint{5.920590in}{1.031213in}}%
\pgfpathlineto{\pgfqpoint{5.920481in}{1.036545in}}%
\pgfpathlineto{\pgfqpoint{5.926934in}{1.046371in}}%
\pgfpathlineto{\pgfqpoint{5.944469in}{1.058743in}}%
\pgfpathlineto{\pgfqpoint{5.962284in}{1.079184in}}%
\pgfpathlineto{\pgfqpoint{5.961416in}{1.074785in}}%
\pgfpathlineto{\pgfqpoint{5.970810in}{1.064162in}}%
\pgfpathlineto{\pgfqpoint{5.986126in}{1.061179in}}%
\pgfpathlineto{\pgfqpoint{5.995035in}{1.071300in}}%
\pgfpathlineto{\pgfqpoint{6.003007in}{1.065970in}}%
\pgfpathlineto{\pgfqpoint{6.017271in}{1.072225in}}%
\pgfpathlineto{\pgfqpoint{6.019311in}{1.078334in}}%
\pgfpathlineto{\pgfqpoint{6.028164in}{1.079708in}}%
\pgfpathlineto{\pgfqpoint{6.037358in}{1.086925in}}%
\pgfpathlineto{\pgfqpoint{6.045463in}{1.089306in}}%
\pgfpathlineto{\pgfqpoint{6.045721in}{1.102055in}}%
\pgfpathlineto{\pgfqpoint{6.050969in}{1.115827in}}%
\pgfpathlineto{\pgfqpoint{6.057891in}{1.124451in}}%
\pgfpathlineto{\pgfqpoint{6.060135in}{1.137399in}}%
\pgfpathlineto{\pgfqpoint{6.069157in}{1.155588in}}%
\pgfpathlineto{\pgfqpoint{6.069480in}{1.166411in}}%
\pgfpathlineto{\pgfqpoint{6.078799in}{1.159359in}}%
\pgfpathlineto{\pgfqpoint{6.087911in}{1.157648in}}%
\pgfpathlineto{\pgfqpoint{6.091456in}{1.162779in}}%
\pgfpathlineto{\pgfqpoint{6.097791in}{1.185136in}}%
\pgfpathlineto{\pgfqpoint{6.106611in}{1.184279in}}%
\pgfpathlineto{\pgfqpoint{6.108947in}{1.192470in}}%
\pgfpathlineto{\pgfqpoint{6.114121in}{1.196164in}}%
\pgfpathlineto{\pgfqpoint{6.125084in}{1.214739in}}%
\pgfpathlineto{\pgfqpoint{6.123492in}{1.220234in}}%
\pgfpathlineto{\pgfqpoint{6.128317in}{1.232177in}}%
\pgfpathlineto{\pgfqpoint{6.154074in}{1.217629in}}%
\pgfpathlineto{\pgfqpoint{6.156926in}{1.230550in}}%
\pgfpathlineto{\pgfqpoint{6.165064in}{1.230836in}}%
\pgfpathlineto{\pgfqpoint{6.171385in}{1.226747in}}%
\pgfpathlineto{\pgfqpoint{6.170120in}{1.219698in}}%
\pgfpathlineto{\pgfqpoint{6.184066in}{1.216492in}}%
\pgfpathlineto{\pgfqpoint{6.191941in}{1.211681in}}%
\pgfpathlineto{\pgfqpoint{6.197645in}{1.203210in}}%
\pgfpathlineto{\pgfqpoint{6.195205in}{1.193376in}}%
\pgfpathlineto{\pgfqpoint{6.190277in}{1.192570in}}%
\pgfpathlineto{\pgfqpoint{6.187419in}{1.177731in}}%
\pgfpathlineto{\pgfqpoint{6.193680in}{1.171929in}}%
\pgfpathlineto{\pgfqpoint{6.202493in}{1.176619in}}%
\pgfpathlineto{\pgfqpoint{6.210671in}{1.166716in}}%
\pgfpathlineto{\pgfqpoint{6.228942in}{1.164936in}}%
\pgfpathlineto{\pgfqpoint{6.232089in}{1.159239in}}%
\pgfpathlineto{\pgfqpoint{6.248998in}{1.153699in}}%
\pgfpathlineto{\pgfqpoint{6.246912in}{1.147161in}}%
\pgfpathlineto{\pgfqpoint{6.249111in}{1.128683in}}%
\pgfpathlineto{\pgfqpoint{6.255936in}{1.121706in}}%
\pgfpathlineto{\pgfqpoint{6.245866in}{1.118761in}}%
\pgfpathlineto{\pgfqpoint{6.247201in}{1.110873in}}%
\pgfpathlineto{\pgfqpoint{6.257944in}{1.104185in}}%
\pgfpathlineto{\pgfqpoint{6.258912in}{1.097588in}}%
\pgfpathlineto{\pgfqpoint{6.252847in}{1.092630in}}%
\pgfpathlineto{\pgfqpoint{6.240609in}{1.104375in}}%
\pgfpathlineto{\pgfqpoint{6.239399in}{1.095806in}}%
\pgfpathlineto{\pgfqpoint{6.249644in}{1.091731in}}%
\pgfpathlineto{\pgfqpoint{6.250669in}{1.087517in}}%
\pgfpathlineto{\pgfqpoint{6.264721in}{1.093086in}}%
\pgfpathlineto{\pgfqpoint{6.275489in}{1.094559in}}%
\pgfpathlineto{\pgfqpoint{6.286520in}{1.072304in}}%
\pgfpathlineto{\pgfqpoint{6.285289in}{1.072063in}}%
\pgfpathlineto{\pgfqpoint{6.280305in}{1.071033in}}%
\pgfpathlineto{\pgfqpoint{6.278836in}{1.070726in}}%
\pgfpathlineto{\pgfqpoint{6.277865in}{1.070536in}}%
\pgfpathlineto{\pgfqpoint{6.212189in}{1.056913in}}%
\pgfpathlineto{\pgfqpoint{6.153447in}{1.044890in}}%
\pgfpathlineto{\pgfqpoint{6.072229in}{1.030898in}}%
\pgfpathlineto{\pgfqpoint{6.003251in}{1.021788in}}%
\pgfpathlineto{\pgfqpoint{5.986189in}{1.020348in}}%
\pgfpathclose%
\pgfusepath{fill}%
\end{pgfscope}%
\begin{pgfscope}%
\pgfpathrectangle{\pgfqpoint{3.625000in}{0.100000in}}{\pgfqpoint{2.989028in}{1.913466in}}%
\pgfusepath{clip}%
\pgfsetbuttcap%
\pgfsetmiterjoin%
\definecolor{currentfill}{rgb}{0.802153,0.920185,0.616378}%
\pgfsetfillcolor{currentfill}%
\pgfsetlinewidth{0.000000pt}%
\definecolor{currentstroke}{rgb}{0.000000,0.000000,0.000000}%
\pgfsetstrokecolor{currentstroke}%
\pgfsetstrokeopacity{0.000000}%
\pgfsetdash{}{0pt}%
\pgfpathmoveto{\pgfqpoint{6.278775in}{1.165395in}}%
\pgfpathlineto{\pgfqpoint{6.292015in}{1.171401in}}%
\pgfpathlineto{\pgfqpoint{6.281442in}{1.140471in}}%
\pgfpathlineto{\pgfqpoint{6.274450in}{1.124092in}}%
\pgfpathlineto{\pgfqpoint{6.269413in}{1.133361in}}%
\pgfpathlineto{\pgfqpoint{6.278370in}{1.155552in}}%
\pgfpathlineto{\pgfqpoint{6.278775in}{1.165395in}}%
\pgfpathclose%
\pgfusepath{fill}%
\end{pgfscope}%
\begin{pgfscope}%
\pgfpathrectangle{\pgfqpoint{3.625000in}{0.100000in}}{\pgfqpoint{2.989028in}{1.913466in}}%
\pgfusepath{clip}%
\pgfsetbuttcap%
\pgfsetmiterjoin%
\definecolor{currentfill}{rgb}{0.829373,0.931257,0.610842}%
\pgfsetfillcolor{currentfill}%
\pgfsetlinewidth{0.000000pt}%
\definecolor{currentstroke}{rgb}{0.000000,0.000000,0.000000}%
\pgfsetstrokecolor{currentstroke}%
\pgfsetstrokeopacity{0.000000}%
\pgfsetdash{}{0pt}%
\pgfpathmoveto{\pgfqpoint{5.581343in}{0.972329in}}%
\pgfpathlineto{\pgfqpoint{5.578146in}{0.971866in}}%
\pgfpathlineto{\pgfqpoint{5.575131in}{0.971655in}}%
\pgfpathlineto{\pgfqpoint{5.576421in}{0.962397in}}%
\pgfpathlineto{\pgfqpoint{5.568642in}{0.953066in}}%
\pgfpathlineto{\pgfqpoint{5.567109in}{0.938464in}}%
\pgfpathlineto{\pgfqpoint{5.532334in}{0.935892in}}%
\pgfpathlineto{\pgfqpoint{5.535365in}{0.942748in}}%
\pgfpathlineto{\pgfqpoint{5.547904in}{0.955269in}}%
\pgfpathlineto{\pgfqpoint{5.548251in}{0.962527in}}%
\pgfpathlineto{\pgfqpoint{5.542700in}{0.969399in}}%
\pgfpathlineto{\pgfqpoint{5.490941in}{0.966662in}}%
\pgfpathlineto{\pgfqpoint{5.400448in}{0.963789in}}%
\pgfpathlineto{\pgfqpoint{5.347491in}{0.962824in}}%
\pgfpathlineto{\pgfqpoint{5.307462in}{0.962463in}}%
\pgfpathlineto{\pgfqpoint{5.307158in}{0.995725in}}%
\pgfpathlineto{\pgfqpoint{5.306840in}{1.092287in}}%
\pgfpathlineto{\pgfqpoint{5.307201in}{1.134279in}}%
\pgfpathlineto{\pgfqpoint{5.296384in}{1.140674in}}%
\pgfpathlineto{\pgfqpoint{5.291052in}{1.151814in}}%
\pgfpathlineto{\pgfqpoint{5.285247in}{1.157161in}}%
\pgfpathlineto{\pgfqpoint{5.284590in}{1.165269in}}%
\pgfpathlineto{\pgfqpoint{5.290653in}{1.178102in}}%
\pgfpathlineto{\pgfqpoint{5.290397in}{1.185521in}}%
\pgfpathlineto{\pgfqpoint{5.282233in}{1.183806in}}%
\pgfpathlineto{\pgfqpoint{5.271810in}{1.192312in}}%
\pgfpathlineto{\pgfqpoint{5.263455in}{1.207242in}}%
\pgfpathlineto{\pgfqpoint{5.256447in}{1.214132in}}%
\pgfpathlineto{\pgfqpoint{5.249124in}{1.231060in}}%
\pgfpathlineto{\pgfqpoint{5.325184in}{1.229874in}}%
\pgfpathlineto{\pgfqpoint{5.400789in}{1.232365in}}%
\pgfpathlineto{\pgfqpoint{5.449266in}{1.235161in}}%
\pgfpathlineto{\pgfqpoint{5.464431in}{1.220143in}}%
\pgfpathlineto{\pgfqpoint{5.461659in}{1.204480in}}%
\pgfpathlineto{\pgfqpoint{5.465703in}{1.186846in}}%
\pgfpathlineto{\pgfqpoint{5.470188in}{1.177954in}}%
\pgfpathlineto{\pgfqpoint{5.484499in}{1.165722in}}%
\pgfpathlineto{\pgfqpoint{5.487630in}{1.160319in}}%
\pgfpathlineto{\pgfqpoint{5.503299in}{1.148144in}}%
\pgfpathlineto{\pgfqpoint{5.505570in}{1.134486in}}%
\pgfpathlineto{\pgfqpoint{5.510070in}{1.124410in}}%
\pgfpathlineto{\pgfqpoint{5.520702in}{1.130279in}}%
\pgfpathlineto{\pgfqpoint{5.529465in}{1.127595in}}%
\pgfpathlineto{\pgfqpoint{5.536082in}{1.120472in}}%
\pgfpathlineto{\pgfqpoint{5.533584in}{1.107865in}}%
\pgfpathlineto{\pgfqpoint{5.525462in}{1.091383in}}%
\pgfpathlineto{\pgfqpoint{5.525587in}{1.082799in}}%
\pgfpathlineto{\pgfqpoint{5.531657in}{1.075760in}}%
\pgfpathlineto{\pgfqpoint{5.544719in}{1.066437in}}%
\pgfpathlineto{\pgfqpoint{5.559760in}{1.058870in}}%
\pgfpathlineto{\pgfqpoint{5.563328in}{1.052811in}}%
\pgfpathlineto{\pgfqpoint{5.571020in}{1.049946in}}%
\pgfpathlineto{\pgfqpoint{5.571326in}{1.041904in}}%
\pgfpathlineto{\pgfqpoint{5.577019in}{1.031690in}}%
\pgfpathlineto{\pgfqpoint{5.573069in}{1.022616in}}%
\pgfpathlineto{\pgfqpoint{5.580801in}{1.008178in}}%
\pgfpathlineto{\pgfqpoint{5.586849in}{1.011045in}}%
\pgfpathlineto{\pgfqpoint{5.594266in}{1.005020in}}%
\pgfpathlineto{\pgfqpoint{5.591773in}{0.978188in}}%
\pgfpathlineto{\pgfqpoint{5.583030in}{0.979890in}}%
\pgfpathlineto{\pgfqpoint{5.581343in}{0.972329in}}%
\pgfpathclose%
\pgfusepath{fill}%
\end{pgfscope}%
\begin{pgfscope}%
\pgfpathrectangle{\pgfqpoint{3.625000in}{0.100000in}}{\pgfqpoint{2.989028in}{1.913466in}}%
\pgfusepath{clip}%
\pgfsetbuttcap%
\pgfsetmiterjoin%
\definecolor{currentfill}{rgb}{0.702345,0.879585,0.636678}%
\pgfsetfillcolor{currentfill}%
\pgfsetlinewidth{0.000000pt}%
\definecolor{currentstroke}{rgb}{0.000000,0.000000,0.000000}%
\pgfsetstrokecolor{currentstroke}%
\pgfsetstrokeopacity{0.000000}%
\pgfsetdash{}{0pt}%
\pgfpathmoveto{\pgfqpoint{4.240525in}{0.978113in}}%
\pgfpathlineto{\pgfqpoint{4.246856in}{0.991592in}}%
\pgfpathlineto{\pgfqpoint{4.245123in}{1.011825in}}%
\pgfpathlineto{\pgfqpoint{4.248155in}{1.017578in}}%
\pgfpathlineto{\pgfqpoint{4.248454in}{1.042780in}}%
\pgfpathlineto{\pgfqpoint{4.251328in}{1.050037in}}%
\pgfpathlineto{\pgfqpoint{4.261451in}{1.051167in}}%
\pgfpathlineto{\pgfqpoint{4.270846in}{1.047945in}}%
\pgfpathlineto{\pgfqpoint{4.274938in}{1.039072in}}%
\pgfpathlineto{\pgfqpoint{4.280674in}{1.039412in}}%
\pgfpathlineto{\pgfqpoint{4.287801in}{1.049576in}}%
\pgfpathlineto{\pgfqpoint{4.298103in}{1.099678in}}%
\pgfpathlineto{\pgfqpoint{4.356724in}{1.087575in}}%
\pgfpathlineto{\pgfqpoint{4.390746in}{1.080865in}}%
\pgfpathlineto{\pgfqpoint{4.481091in}{1.064984in}}%
\pgfpathlineto{\pgfqpoint{4.506086in}{1.059964in}}%
\pgfpathlineto{\pgfqpoint{4.555137in}{1.052362in}}%
\pgfpathlineto{\pgfqpoint{4.545079in}{0.987601in}}%
\pgfpathlineto{\pgfqpoint{4.534619in}{0.920049in}}%
\pgfpathlineto{\pgfqpoint{4.522573in}{0.844052in}}%
\pgfpathlineto{\pgfqpoint{4.509025in}{0.756815in}}%
\pgfpathlineto{\pgfqpoint{4.498080in}{0.685166in}}%
\pgfpathlineto{\pgfqpoint{4.419687in}{0.697614in}}%
\pgfpathlineto{\pgfqpoint{4.385244in}{0.703513in}}%
\pgfpathlineto{\pgfqpoint{4.369857in}{0.712725in}}%
\pgfpathlineto{\pgfqpoint{4.269535in}{0.773254in}}%
\pgfpathlineto{\pgfqpoint{4.194248in}{0.819189in}}%
\pgfpathlineto{\pgfqpoint{4.196756in}{0.827331in}}%
\pgfpathlineto{\pgfqpoint{4.202975in}{0.833027in}}%
\pgfpathlineto{\pgfqpoint{4.209444in}{0.831972in}}%
\pgfpathlineto{\pgfqpoint{4.218826in}{0.837954in}}%
\pgfpathlineto{\pgfqpoint{4.220303in}{0.846552in}}%
\pgfpathlineto{\pgfqpoint{4.208861in}{0.856982in}}%
\pgfpathlineto{\pgfqpoint{4.216998in}{0.877004in}}%
\pgfpathlineto{\pgfqpoint{4.225180in}{0.884704in}}%
\pgfpathlineto{\pgfqpoint{4.229063in}{0.893839in}}%
\pgfpathlineto{\pgfqpoint{4.231455in}{0.910589in}}%
\pgfpathlineto{\pgfqpoint{4.239127in}{0.918174in}}%
\pgfpathlineto{\pgfqpoint{4.255230in}{0.925750in}}%
\pgfpathlineto{\pgfqpoint{4.254874in}{0.932426in}}%
\pgfpathlineto{\pgfqpoint{4.245881in}{0.940711in}}%
\pgfpathlineto{\pgfqpoint{4.244656in}{0.957789in}}%
\pgfpathlineto{\pgfqpoint{4.238464in}{0.970274in}}%
\pgfpathlineto{\pgfqpoint{4.240525in}{0.978113in}}%
\pgfpathclose%
\pgfusepath{fill}%
\end{pgfscope}%
\begin{pgfscope}%
\pgfpathrectangle{\pgfqpoint{3.625000in}{0.100000in}}{\pgfqpoint{2.989028in}{1.913466in}}%
\pgfusepath{clip}%
\pgfsetbuttcap%
\pgfsetmiterjoin%
\definecolor{currentfill}{rgb}{0.264052,0.609150,0.709804}%
\pgfsetfillcolor{currentfill}%
\pgfsetlinewidth{0.000000pt}%
\definecolor{currentstroke}{rgb}{0.000000,0.000000,0.000000}%
\pgfsetstrokecolor{currentstroke}%
\pgfsetstrokeopacity{0.000000}%
\pgfsetdash{}{0pt}%
\pgfpathmoveto{\pgfqpoint{5.315948in}{0.774957in}}%
\pgfpathlineto{\pgfqpoint{5.301416in}{0.779451in}}%
\pgfpathlineto{\pgfqpoint{5.282795in}{0.793711in}}%
\pgfpathlineto{\pgfqpoint{5.274524in}{0.796776in}}%
\pgfpathlineto{\pgfqpoint{5.269268in}{0.790617in}}%
\pgfpathlineto{\pgfqpoint{5.262631in}{0.790306in}}%
\pgfpathlineto{\pgfqpoint{5.254234in}{0.795535in}}%
\pgfpathlineto{\pgfqpoint{5.241057in}{0.788835in}}%
\pgfpathlineto{\pgfqpoint{5.236367in}{0.792135in}}%
\pgfpathlineto{\pgfqpoint{5.223401in}{0.784318in}}%
\pgfpathlineto{\pgfqpoint{5.209722in}{0.785761in}}%
\pgfpathlineto{\pgfqpoint{5.199172in}{0.790840in}}%
\pgfpathlineto{\pgfqpoint{5.191783in}{0.788928in}}%
\pgfpathlineto{\pgfqpoint{5.179956in}{0.796112in}}%
\pgfpathlineto{\pgfqpoint{5.173377in}{0.783442in}}%
\pgfpathlineto{\pgfqpoint{5.166649in}{0.794338in}}%
\pgfpathlineto{\pgfqpoint{5.153360in}{0.790111in}}%
\pgfpathlineto{\pgfqpoint{5.141736in}{0.800460in}}%
\pgfpathlineto{\pgfqpoint{5.131579in}{0.792118in}}%
\pgfpathlineto{\pgfqpoint{5.126493in}{0.799784in}}%
\pgfpathlineto{\pgfqpoint{5.119161in}{0.802273in}}%
\pgfpathlineto{\pgfqpoint{5.114696in}{0.809664in}}%
\pgfpathlineto{\pgfqpoint{5.105106in}{0.811767in}}%
\pgfpathlineto{\pgfqpoint{5.099573in}{0.806204in}}%
\pgfpathlineto{\pgfqpoint{5.090171in}{0.813414in}}%
\pgfpathlineto{\pgfqpoint{5.085789in}{0.811767in}}%
\pgfpathlineto{\pgfqpoint{5.070207in}{0.817609in}}%
\pgfpathlineto{\pgfqpoint{5.060448in}{0.818261in}}%
\pgfpathlineto{\pgfqpoint{5.056080in}{0.830667in}}%
\pgfpathlineto{\pgfqpoint{5.048256in}{0.829122in}}%
\pgfpathlineto{\pgfqpoint{5.039296in}{0.832183in}}%
\pgfpathlineto{\pgfqpoint{5.033395in}{0.830414in}}%
\pgfpathlineto{\pgfqpoint{5.020741in}{0.844347in}}%
\pgfpathlineto{\pgfqpoint{5.017215in}{0.843436in}}%
\pgfpathlineto{\pgfqpoint{5.020417in}{0.899940in}}%
\pgfpathlineto{\pgfqpoint{5.023906in}{0.969876in}}%
\pgfpathlineto{\pgfqpoint{4.966639in}{0.973080in}}%
\pgfpathlineto{\pgfqpoint{4.910128in}{0.977352in}}%
\pgfpathlineto{\pgfqpoint{4.866469in}{0.981143in}}%
\pgfpathlineto{\pgfqpoint{4.869368in}{1.014183in}}%
\pgfpathlineto{\pgfqpoint{4.919404in}{1.009408in}}%
\pgfpathlineto{\pgfqpoint{4.976155in}{1.005849in}}%
\pgfpathlineto{\pgfqpoint{5.055191in}{1.001189in}}%
\pgfpathlineto{\pgfqpoint{5.102105in}{0.999379in}}%
\pgfpathlineto{\pgfqpoint{5.158309in}{0.997396in}}%
\pgfpathlineto{\pgfqpoint{5.234656in}{0.996060in}}%
\pgfpathlineto{\pgfqpoint{5.307158in}{0.995725in}}%
\pgfpathlineto{\pgfqpoint{5.307462in}{0.962463in}}%
\pgfpathlineto{\pgfqpoint{5.311531in}{0.937404in}}%
\pgfpathlineto{\pgfqpoint{5.317853in}{0.891120in}}%
\pgfpathlineto{\pgfqpoint{5.316549in}{0.812079in}}%
\pgfpathlineto{\pgfqpoint{5.315948in}{0.774957in}}%
\pgfpathclose%
\pgfusepath{fill}%
\end{pgfscope}%
\begin{pgfscope}%
\pgfpathrectangle{\pgfqpoint{3.625000in}{0.100000in}}{\pgfqpoint{2.989028in}{1.913466in}}%
\pgfusepath{clip}%
\pgfsetbuttcap%
\pgfsetmiterjoin%
\definecolor{currentfill}{rgb}{0.601615,0.839677,0.644137}%
\pgfsetfillcolor{currentfill}%
\pgfsetlinewidth{0.000000pt}%
\definecolor{currentstroke}{rgb}{0.000000,0.000000,0.000000}%
\pgfsetstrokecolor{currentstroke}%
\pgfsetstrokeopacity{0.000000}%
\pgfsetdash{}{0pt}%
\pgfpathmoveto{\pgfqpoint{5.860069in}{0.897708in}}%
\pgfpathlineto{\pgfqpoint{5.860112in}{0.912353in}}%
\pgfpathlineto{\pgfqpoint{5.872839in}{0.917977in}}%
\pgfpathlineto{\pgfqpoint{5.873372in}{0.926967in}}%
\pgfpathlineto{\pgfqpoint{5.884767in}{0.937949in}}%
\pgfpathlineto{\pgfqpoint{5.899011in}{0.940198in}}%
\pgfpathlineto{\pgfqpoint{5.911533in}{0.952301in}}%
\pgfpathlineto{\pgfqpoint{5.924610in}{0.957711in}}%
\pgfpathlineto{\pgfqpoint{5.934425in}{0.973951in}}%
\pgfpathlineto{\pgfqpoint{5.943364in}{0.972772in}}%
\pgfpathlineto{\pgfqpoint{5.962256in}{0.987570in}}%
\pgfpathlineto{\pgfqpoint{5.972232in}{0.987849in}}%
\pgfpathlineto{\pgfqpoint{5.976445in}{0.999123in}}%
\pgfpathlineto{\pgfqpoint{5.984375in}{1.006973in}}%
\pgfpathlineto{\pgfqpoint{5.986189in}{1.020348in}}%
\pgfpathlineto{\pgfqpoint{6.003251in}{1.021788in}}%
\pgfpathlineto{\pgfqpoint{6.072229in}{1.030898in}}%
\pgfpathlineto{\pgfqpoint{6.153447in}{1.044890in}}%
\pgfpathlineto{\pgfqpoint{6.212189in}{1.056913in}}%
\pgfpathlineto{\pgfqpoint{6.277865in}{1.070536in}}%
\pgfpathlineto{\pgfqpoint{6.296631in}{1.044770in}}%
\pgfpathlineto{\pgfqpoint{6.286466in}{1.046415in}}%
\pgfpathlineto{\pgfqpoint{6.273417in}{1.043236in}}%
\pgfpathlineto{\pgfqpoint{6.260559in}{1.030105in}}%
\pgfpathlineto{\pgfqpoint{6.251318in}{1.031050in}}%
\pgfpathlineto{\pgfqpoint{6.250149in}{1.023251in}}%
\pgfpathlineto{\pgfqpoint{6.266852in}{1.029418in}}%
\pgfpathlineto{\pgfqpoint{6.283660in}{1.031961in}}%
\pgfpathlineto{\pgfqpoint{6.289901in}{1.028577in}}%
\pgfpathlineto{\pgfqpoint{6.298300in}{1.032435in}}%
\pgfpathlineto{\pgfqpoint{6.302634in}{1.029732in}}%
\pgfpathlineto{\pgfqpoint{6.306470in}{1.016862in}}%
\pgfpathlineto{\pgfqpoint{6.298487in}{1.012876in}}%
\pgfpathlineto{\pgfqpoint{6.293008in}{0.997182in}}%
\pgfpathlineto{\pgfqpoint{6.287263in}{0.991072in}}%
\pgfpathlineto{\pgfqpoint{6.269653in}{0.992409in}}%
\pgfpathlineto{\pgfqpoint{6.270600in}{1.001624in}}%
\pgfpathlineto{\pgfqpoint{6.260982in}{0.997599in}}%
\pgfpathlineto{\pgfqpoint{6.258890in}{0.989913in}}%
\pgfpathlineto{\pgfqpoint{6.265483in}{0.981952in}}%
\pgfpathlineto{\pgfqpoint{6.267722in}{0.968440in}}%
\pgfpathlineto{\pgfqpoint{6.256935in}{0.960741in}}%
\pgfpathlineto{\pgfqpoint{6.268601in}{0.958032in}}%
\pgfpathlineto{\pgfqpoint{6.273881in}{0.963694in}}%
\pgfpathlineto{\pgfqpoint{6.285553in}{0.964383in}}%
\pgfpathlineto{\pgfqpoint{6.279555in}{0.952158in}}%
\pgfpathlineto{\pgfqpoint{6.272208in}{0.946268in}}%
\pgfpathlineto{\pgfqpoint{6.249984in}{0.940370in}}%
\pgfpathlineto{\pgfqpoint{6.227166in}{0.919675in}}%
\pgfpathlineto{\pgfqpoint{6.213096in}{0.899284in}}%
\pgfpathlineto{\pgfqpoint{6.213013in}{0.891003in}}%
\pgfpathlineto{\pgfqpoint{6.207393in}{0.879575in}}%
\pgfpathlineto{\pgfqpoint{6.178548in}{0.872057in}}%
\pgfpathlineto{\pgfqpoint{6.109952in}{0.921230in}}%
\pgfpathlineto{\pgfqpoint{6.049543in}{0.912354in}}%
\pgfpathlineto{\pgfqpoint{6.049037in}{0.920550in}}%
\pgfpathlineto{\pgfqpoint{6.039855in}{0.929779in}}%
\pgfpathlineto{\pgfqpoint{6.032904in}{0.932038in}}%
\pgfpathlineto{\pgfqpoint{5.967269in}{0.925218in}}%
\pgfpathlineto{\pgfqpoint{5.952183in}{0.920102in}}%
\pgfpathlineto{\pgfqpoint{5.924943in}{0.906551in}}%
\pgfpathlineto{\pgfqpoint{5.901402in}{0.902801in}}%
\pgfpathlineto{\pgfqpoint{5.860069in}{0.897708in}}%
\pgfpathclose%
\pgfusepath{fill}%
\end{pgfscope}%
\begin{pgfscope}%
\pgfpathrectangle{\pgfqpoint{3.625000in}{0.100000in}}{\pgfqpoint{2.989028in}{1.913466in}}%
\pgfusepath{clip}%
\pgfsetbuttcap%
\pgfsetmiterjoin%
\definecolor{currentfill}{rgb}{0.601615,0.839677,0.644137}%
\pgfsetfillcolor{currentfill}%
\pgfsetlinewidth{0.000000pt}%
\definecolor{currentstroke}{rgb}{0.000000,0.000000,0.000000}%
\pgfsetstrokecolor{currentstroke}%
\pgfsetstrokeopacity{0.000000}%
\pgfsetdash{}{0pt}%
\pgfpathmoveto{\pgfqpoint{5.860069in}{0.897708in}}%
\pgfpathlineto{\pgfqpoint{5.791375in}{0.890163in}}%
\pgfpathlineto{\pgfqpoint{5.728533in}{0.884512in}}%
\pgfpathlineto{\pgfqpoint{5.663945in}{0.880335in}}%
\pgfpathlineto{\pgfqpoint{5.652858in}{0.878728in}}%
\pgfpathlineto{\pgfqpoint{5.609319in}{0.875379in}}%
\pgfpathlineto{\pgfqpoint{5.539584in}{0.871310in}}%
\pgfpathlineto{\pgfqpoint{5.551990in}{0.881608in}}%
\pgfpathlineto{\pgfqpoint{5.549159in}{0.892449in}}%
\pgfpathlineto{\pgfqpoint{5.551902in}{0.905590in}}%
\pgfpathlineto{\pgfqpoint{5.556092in}{0.911777in}}%
\pgfpathlineto{\pgfqpoint{5.555934in}{0.920384in}}%
\pgfpathlineto{\pgfqpoint{5.567090in}{0.925797in}}%
\pgfpathlineto{\pgfqpoint{5.567109in}{0.938464in}}%
\pgfpathlineto{\pgfqpoint{5.568642in}{0.953066in}}%
\pgfpathlineto{\pgfqpoint{5.576421in}{0.962397in}}%
\pgfpathlineto{\pgfqpoint{5.575131in}{0.971655in}}%
\pgfpathlineto{\pgfqpoint{5.578146in}{0.971866in}}%
\pgfpathlineto{\pgfqpoint{5.581343in}{0.972329in}}%
\pgfpathlineto{\pgfqpoint{5.653480in}{0.976914in}}%
\pgfpathlineto{\pgfqpoint{5.651065in}{0.988764in}}%
\pgfpathlineto{\pgfqpoint{5.662674in}{0.987156in}}%
\pgfpathlineto{\pgfqpoint{5.732938in}{0.994107in}}%
\pgfpathlineto{\pgfqpoint{5.789337in}{0.997299in}}%
\pgfpathlineto{\pgfqpoint{5.823708in}{1.000190in}}%
\pgfpathlineto{\pgfqpoint{5.880622in}{1.005586in}}%
\pgfpathlineto{\pgfqpoint{5.881571in}{1.006660in}}%
\pgfpathlineto{\pgfqpoint{5.905888in}{1.009068in}}%
\pgfpathlineto{\pgfqpoint{5.986189in}{1.020348in}}%
\pgfpathlineto{\pgfqpoint{5.984375in}{1.006973in}}%
\pgfpathlineto{\pgfqpoint{5.976445in}{0.999123in}}%
\pgfpathlineto{\pgfqpoint{5.972232in}{0.987849in}}%
\pgfpathlineto{\pgfqpoint{5.962256in}{0.987570in}}%
\pgfpathlineto{\pgfqpoint{5.943364in}{0.972772in}}%
\pgfpathlineto{\pgfqpoint{5.934425in}{0.973951in}}%
\pgfpathlineto{\pgfqpoint{5.924610in}{0.957711in}}%
\pgfpathlineto{\pgfqpoint{5.911533in}{0.952301in}}%
\pgfpathlineto{\pgfqpoint{5.899011in}{0.940198in}}%
\pgfpathlineto{\pgfqpoint{5.884767in}{0.937949in}}%
\pgfpathlineto{\pgfqpoint{5.873372in}{0.926967in}}%
\pgfpathlineto{\pgfqpoint{5.872839in}{0.917977in}}%
\pgfpathlineto{\pgfqpoint{5.860112in}{0.912353in}}%
\pgfpathlineto{\pgfqpoint{5.860069in}{0.897708in}}%
\pgfpathclose%
\pgfusepath{fill}%
\end{pgfscope}%
\begin{pgfscope}%
\pgfpathrectangle{\pgfqpoint{3.625000in}{0.100000in}}{\pgfqpoint{2.989028in}{1.913466in}}%
\pgfusepath{clip}%
\pgfsetbuttcap%
\pgfsetmiterjoin%
\definecolor{currentfill}{rgb}{0.212995,0.511419,0.730796}%
\pgfsetfillcolor{currentfill}%
\pgfsetlinewidth{0.000000pt}%
\definecolor{currentstroke}{rgb}{0.000000,0.000000,0.000000}%
\pgfsetstrokecolor{currentstroke}%
\pgfsetstrokeopacity{0.000000}%
\pgfsetdash{}{0pt}%
\pgfpathmoveto{\pgfqpoint{4.641758in}{0.694333in}}%
\pgfpathlineto{\pgfqpoint{4.636859in}{0.702205in}}%
\pgfpathlineto{\pgfqpoint{4.638886in}{0.708987in}}%
\pgfpathlineto{\pgfqpoint{4.673324in}{0.704715in}}%
\pgfpathlineto{\pgfqpoint{4.737359in}{0.697442in}}%
\pgfpathlineto{\pgfqpoint{4.783660in}{0.692892in}}%
\pgfpathlineto{\pgfqpoint{4.837135in}{0.687496in}}%
\pgfpathlineto{\pgfqpoint{4.840044in}{0.721225in}}%
\pgfpathlineto{\pgfqpoint{4.849140in}{0.807015in}}%
\pgfpathlineto{\pgfqpoint{4.855041in}{0.867115in}}%
\pgfpathlineto{\pgfqpoint{4.860142in}{0.924525in}}%
\pgfpathlineto{\pgfqpoint{4.864896in}{0.981231in}}%
\pgfpathlineto{\pgfqpoint{4.866469in}{0.981143in}}%
\pgfpathlineto{\pgfqpoint{4.910128in}{0.977352in}}%
\pgfpathlineto{\pgfqpoint{4.966639in}{0.973080in}}%
\pgfpathlineto{\pgfqpoint{5.023906in}{0.969876in}}%
\pgfpathlineto{\pgfqpoint{5.020417in}{0.899940in}}%
\pgfpathlineto{\pgfqpoint{5.017215in}{0.843436in}}%
\pgfpathlineto{\pgfqpoint{5.020741in}{0.844347in}}%
\pgfpathlineto{\pgfqpoint{5.033395in}{0.830414in}}%
\pgfpathlineto{\pgfqpoint{5.039296in}{0.832183in}}%
\pgfpathlineto{\pgfqpoint{5.048256in}{0.829122in}}%
\pgfpathlineto{\pgfqpoint{5.056080in}{0.830667in}}%
\pgfpathlineto{\pgfqpoint{5.060448in}{0.818261in}}%
\pgfpathlineto{\pgfqpoint{5.070207in}{0.817609in}}%
\pgfpathlineto{\pgfqpoint{5.085789in}{0.811767in}}%
\pgfpathlineto{\pgfqpoint{5.090171in}{0.813414in}}%
\pgfpathlineto{\pgfqpoint{5.099573in}{0.806204in}}%
\pgfpathlineto{\pgfqpoint{5.105106in}{0.811767in}}%
\pgfpathlineto{\pgfqpoint{5.114696in}{0.809664in}}%
\pgfpathlineto{\pgfqpoint{5.119161in}{0.802273in}}%
\pgfpathlineto{\pgfqpoint{5.126493in}{0.799784in}}%
\pgfpathlineto{\pgfqpoint{5.131579in}{0.792118in}}%
\pgfpathlineto{\pgfqpoint{5.141736in}{0.800460in}}%
\pgfpathlineto{\pgfqpoint{5.153360in}{0.790111in}}%
\pgfpathlineto{\pgfqpoint{5.166649in}{0.794338in}}%
\pgfpathlineto{\pgfqpoint{5.173377in}{0.783442in}}%
\pgfpathlineto{\pgfqpoint{5.179956in}{0.796112in}}%
\pgfpathlineto{\pgfqpoint{5.191783in}{0.788928in}}%
\pgfpathlineto{\pgfqpoint{5.199172in}{0.790840in}}%
\pgfpathlineto{\pgfqpoint{5.209722in}{0.785761in}}%
\pgfpathlineto{\pgfqpoint{5.223401in}{0.784318in}}%
\pgfpathlineto{\pgfqpoint{5.236367in}{0.792135in}}%
\pgfpathlineto{\pgfqpoint{5.241057in}{0.788835in}}%
\pgfpathlineto{\pgfqpoint{5.254234in}{0.795535in}}%
\pgfpathlineto{\pgfqpoint{5.262631in}{0.790306in}}%
\pgfpathlineto{\pgfqpoint{5.269268in}{0.790617in}}%
\pgfpathlineto{\pgfqpoint{5.274524in}{0.796776in}}%
\pgfpathlineto{\pgfqpoint{5.282795in}{0.793711in}}%
\pgfpathlineto{\pgfqpoint{5.301416in}{0.779451in}}%
\pgfpathlineto{\pgfqpoint{5.315948in}{0.774957in}}%
\pgfpathlineto{\pgfqpoint{5.321775in}{0.769451in}}%
\pgfpathlineto{\pgfqpoint{5.329070in}{0.772449in}}%
\pgfpathlineto{\pgfqpoint{5.340123in}{0.770153in}}%
\pgfpathlineto{\pgfqpoint{5.340340in}{0.735108in}}%
\pgfpathlineto{\pgfqpoint{5.341262in}{0.667326in}}%
\pgfpathlineto{\pgfqpoint{5.348939in}{0.660813in}}%
\pgfpathlineto{\pgfqpoint{5.355092in}{0.648811in}}%
\pgfpathlineto{\pgfqpoint{5.352929in}{0.640781in}}%
\pgfpathlineto{\pgfqpoint{5.357761in}{0.637381in}}%
\pgfpathlineto{\pgfqpoint{5.361688in}{0.622574in}}%
\pgfpathlineto{\pgfqpoint{5.369209in}{0.614644in}}%
\pgfpathlineto{\pgfqpoint{5.370894in}{0.597808in}}%
\pgfpathlineto{\pgfqpoint{5.369586in}{0.590680in}}%
\pgfpathlineto{\pgfqpoint{5.359363in}{0.571815in}}%
\pgfpathlineto{\pgfqpoint{5.358187in}{0.558251in}}%
\pgfpathlineto{\pgfqpoint{5.361655in}{0.555420in}}%
\pgfpathlineto{\pgfqpoint{5.361956in}{0.543537in}}%
\pgfpathlineto{\pgfqpoint{5.358432in}{0.536068in}}%
\pgfpathlineto{\pgfqpoint{5.352950in}{0.533255in}}%
\pgfpathlineto{\pgfqpoint{5.347582in}{0.523496in}}%
\pgfpathlineto{\pgfqpoint{5.354424in}{0.514043in}}%
\pgfpathlineto{\pgfqpoint{5.341144in}{0.513857in}}%
\pgfpathlineto{\pgfqpoint{5.305636in}{0.497611in}}%
\pgfpathlineto{\pgfqpoint{5.312415in}{0.507333in}}%
\pgfpathlineto{\pgfqpoint{5.304206in}{0.512575in}}%
\pgfpathlineto{\pgfqpoint{5.302490in}{0.521488in}}%
\pgfpathlineto{\pgfqpoint{5.286468in}{0.505964in}}%
\pgfpathlineto{\pgfqpoint{5.293579in}{0.495354in}}%
\pgfpathlineto{\pgfqpoint{5.283435in}{0.481844in}}%
\pgfpathlineto{\pgfqpoint{5.277976in}{0.482126in}}%
\pgfpathlineto{\pgfqpoint{5.272841in}{0.467413in}}%
\pgfpathlineto{\pgfqpoint{5.256600in}{0.455840in}}%
\pgfpathlineto{\pgfqpoint{5.241427in}{0.451669in}}%
\pgfpathlineto{\pgfqpoint{5.214964in}{0.440819in}}%
\pgfpathlineto{\pgfqpoint{5.212222in}{0.446871in}}%
\pgfpathlineto{\pgfqpoint{5.195737in}{0.441300in}}%
\pgfpathlineto{\pgfqpoint{5.205876in}{0.431814in}}%
\pgfpathlineto{\pgfqpoint{5.189883in}{0.423573in}}%
\pgfpathlineto{\pgfqpoint{5.172633in}{0.411131in}}%
\pgfpathlineto{\pgfqpoint{5.168490in}{0.416915in}}%
\pgfpathlineto{\pgfqpoint{5.154373in}{0.408246in}}%
\pgfpathlineto{\pgfqpoint{5.168210in}{0.404954in}}%
\pgfpathlineto{\pgfqpoint{5.157777in}{0.391324in}}%
\pgfpathlineto{\pgfqpoint{5.152685in}{0.395380in}}%
\pgfpathlineto{\pgfqpoint{5.145845in}{0.388865in}}%
\pgfpathlineto{\pgfqpoint{5.154369in}{0.383174in}}%
\pgfpathlineto{\pgfqpoint{5.149329in}{0.374856in}}%
\pgfpathlineto{\pgfqpoint{5.143135in}{0.355164in}}%
\pgfpathlineto{\pgfqpoint{5.134200in}{0.336199in}}%
\pgfpathlineto{\pgfqpoint{5.138211in}{0.323767in}}%
\pgfpathlineto{\pgfqpoint{5.145029in}{0.294458in}}%
\pgfpathlineto{\pgfqpoint{5.145643in}{0.282596in}}%
\pgfpathlineto{\pgfqpoint{5.156178in}{0.267038in}}%
\pgfpathlineto{\pgfqpoint{5.148054in}{0.267947in}}%
\pgfpathlineto{\pgfqpoint{5.140164in}{0.260098in}}%
\pgfpathlineto{\pgfqpoint{5.130923in}{0.266267in}}%
\pgfpathlineto{\pgfqpoint{5.127605in}{0.272421in}}%
\pgfpathlineto{\pgfqpoint{5.114453in}{0.275284in}}%
\pgfpathlineto{\pgfqpoint{5.094363in}{0.275635in}}%
\pgfpathlineto{\pgfqpoint{5.079561in}{0.287289in}}%
\pgfpathlineto{\pgfqpoint{5.066120in}{0.289236in}}%
\pgfpathlineto{\pgfqpoint{5.057966in}{0.298497in}}%
\pgfpathlineto{\pgfqpoint{5.040889in}{0.302225in}}%
\pgfpathlineto{\pgfqpoint{5.031552in}{0.332004in}}%
\pgfpathlineto{\pgfqpoint{5.022004in}{0.343934in}}%
\pgfpathlineto{\pgfqpoint{5.023626in}{0.355270in}}%
\pgfpathlineto{\pgfqpoint{5.017706in}{0.363554in}}%
\pgfpathlineto{\pgfqpoint{5.021421in}{0.374877in}}%
\pgfpathlineto{\pgfqpoint{5.018356in}{0.383175in}}%
\pgfpathlineto{\pgfqpoint{5.008762in}{0.386934in}}%
\pgfpathlineto{\pgfqpoint{4.999791in}{0.396502in}}%
\pgfpathlineto{\pgfqpoint{4.996522in}{0.409312in}}%
\pgfpathlineto{\pgfqpoint{4.988033in}{0.420954in}}%
\pgfpathlineto{\pgfqpoint{4.976735in}{0.430012in}}%
\pgfpathlineto{\pgfqpoint{4.966558in}{0.455996in}}%
\pgfpathlineto{\pgfqpoint{4.958340in}{0.484413in}}%
\pgfpathlineto{\pgfqpoint{4.951597in}{0.495650in}}%
\pgfpathlineto{\pgfqpoint{4.939908in}{0.505109in}}%
\pgfpathlineto{\pgfqpoint{4.936987in}{0.511975in}}%
\pgfpathlineto{\pgfqpoint{4.926042in}{0.516238in}}%
\pgfpathlineto{\pgfqpoint{4.915239in}{0.534431in}}%
\pgfpathlineto{\pgfqpoint{4.905374in}{0.533035in}}%
\pgfpathlineto{\pgfqpoint{4.895356in}{0.537579in}}%
\pgfpathlineto{\pgfqpoint{4.881156in}{0.536700in}}%
\pgfpathlineto{\pgfqpoint{4.866731in}{0.544199in}}%
\pgfpathlineto{\pgfqpoint{4.862663in}{0.537067in}}%
\pgfpathlineto{\pgfqpoint{4.845804in}{0.536889in}}%
\pgfpathlineto{\pgfqpoint{4.837222in}{0.523370in}}%
\pgfpathlineto{\pgfqpoint{4.829757in}{0.506631in}}%
\pgfpathlineto{\pgfqpoint{4.813883in}{0.488663in}}%
\pgfpathlineto{\pgfqpoint{4.795899in}{0.496548in}}%
\pgfpathlineto{\pgfqpoint{4.793347in}{0.501760in}}%
\pgfpathlineto{\pgfqpoint{4.782418in}{0.505734in}}%
\pgfpathlineto{\pgfqpoint{4.779061in}{0.511171in}}%
\pgfpathlineto{\pgfqpoint{4.764552in}{0.516660in}}%
\pgfpathlineto{\pgfqpoint{4.756436in}{0.527879in}}%
\pgfpathlineto{\pgfqpoint{4.746957in}{0.533292in}}%
\pgfpathlineto{\pgfqpoint{4.738802in}{0.542738in}}%
\pgfpathlineto{\pgfqpoint{4.732448in}{0.558735in}}%
\pgfpathlineto{\pgfqpoint{4.733159in}{0.580595in}}%
\pgfpathlineto{\pgfqpoint{4.725706in}{0.591641in}}%
\pgfpathlineto{\pgfqpoint{4.724848in}{0.603605in}}%
\pgfpathlineto{\pgfqpoint{4.708312in}{0.621518in}}%
\pgfpathlineto{\pgfqpoint{4.698690in}{0.625342in}}%
\pgfpathlineto{\pgfqpoint{4.688479in}{0.642057in}}%
\pgfpathlineto{\pgfqpoint{4.679778in}{0.648698in}}%
\pgfpathlineto{\pgfqpoint{4.668723in}{0.664869in}}%
\pgfpathlineto{\pgfqpoint{4.657411in}{0.671867in}}%
\pgfpathlineto{\pgfqpoint{4.650003in}{0.689800in}}%
\pgfpathlineto{\pgfqpoint{4.641758in}{0.694333in}}%
\pgfpathclose%
\pgfusepath{fill}%
\end{pgfscope}%
\begin{pgfscope}%
\pgfpathrectangle{\pgfqpoint{3.625000in}{0.100000in}}{\pgfqpoint{2.989028in}{1.913466in}}%
\pgfusepath{clip}%
\pgfsetbuttcap%
\pgfsetmiterjoin%
\definecolor{currentfill}{rgb}{0.474279,0.789850,0.645982}%
\pgfsetfillcolor{currentfill}%
\pgfsetlinewidth{0.000000pt}%
\definecolor{currentstroke}{rgb}{0.000000,0.000000,0.000000}%
\pgfsetstrokecolor{currentstroke}%
\pgfsetstrokeopacity{0.000000}%
\pgfsetdash{}{0pt}%
\pgfpathmoveto{\pgfqpoint{4.555137in}{1.052362in}}%
\pgfpathlineto{\pgfqpoint{4.636744in}{1.040510in}}%
\pgfpathlineto{\pgfqpoint{4.668359in}{1.035639in}}%
\pgfpathlineto{\pgfqpoint{4.757477in}{1.025040in}}%
\pgfpathlineto{\pgfqpoint{4.817460in}{1.018783in}}%
\pgfpathlineto{\pgfqpoint{4.869368in}{1.014183in}}%
\pgfpathlineto{\pgfqpoint{4.866469in}{0.981143in}}%
\pgfpathlineto{\pgfqpoint{4.864896in}{0.981231in}}%
\pgfpathlineto{\pgfqpoint{4.860142in}{0.924525in}}%
\pgfpathlineto{\pgfqpoint{4.855041in}{0.867115in}}%
\pgfpathlineto{\pgfqpoint{4.849140in}{0.807015in}}%
\pgfpathlineto{\pgfqpoint{4.840044in}{0.721225in}}%
\pgfpathlineto{\pgfqpoint{4.837135in}{0.687496in}}%
\pgfpathlineto{\pgfqpoint{4.783660in}{0.692892in}}%
\pgfpathlineto{\pgfqpoint{4.737359in}{0.697442in}}%
\pgfpathlineto{\pgfqpoint{4.673324in}{0.704715in}}%
\pgfpathlineto{\pgfqpoint{4.638886in}{0.708987in}}%
\pgfpathlineto{\pgfqpoint{4.636859in}{0.702205in}}%
\pgfpathlineto{\pgfqpoint{4.641758in}{0.694333in}}%
\pgfpathlineto{\pgfqpoint{4.600370in}{0.699708in}}%
\pgfpathlineto{\pgfqpoint{4.549306in}{0.707037in}}%
\pgfpathlineto{\pgfqpoint{4.544659in}{0.678155in}}%
\pgfpathlineto{\pgfqpoint{4.498080in}{0.685166in}}%
\pgfpathlineto{\pgfqpoint{4.509025in}{0.756815in}}%
\pgfpathlineto{\pgfqpoint{4.522573in}{0.844052in}}%
\pgfpathlineto{\pgfqpoint{4.534619in}{0.920049in}}%
\pgfpathlineto{\pgfqpoint{4.545079in}{0.987601in}}%
\pgfpathlineto{\pgfqpoint{4.555137in}{1.052362in}}%
\pgfpathclose%
\pgfusepath{fill}%
\end{pgfscope}%
\begin{pgfscope}%
\pgfpathrectangle{\pgfqpoint{3.625000in}{0.100000in}}{\pgfqpoint{2.989028in}{1.913466in}}%
\pgfusepath{clip}%
\pgfsetbuttcap%
\pgfsetmiterjoin%
\definecolor{currentfill}{rgb}{0.756786,0.901730,0.625606}%
\pgfsetfillcolor{currentfill}%
\pgfsetlinewidth{0.000000pt}%
\definecolor{currentstroke}{rgb}{0.000000,0.000000,0.000000}%
\pgfsetstrokecolor{currentstroke}%
\pgfsetstrokeopacity{0.000000}%
\pgfsetdash{}{0pt}%
\pgfpathmoveto{\pgfqpoint{5.852687in}{0.632577in}}%
\pgfpathlineto{\pgfqpoint{5.774834in}{0.623930in}}%
\pgfpathlineto{\pgfqpoint{5.706207in}{0.618602in}}%
\pgfpathlineto{\pgfqpoint{5.705349in}{0.610192in}}%
\pgfpathlineto{\pgfqpoint{5.711652in}{0.602188in}}%
\pgfpathlineto{\pgfqpoint{5.719346in}{0.597485in}}%
\pgfpathlineto{\pgfqpoint{5.719230in}{0.585103in}}%
\pgfpathlineto{\pgfqpoint{5.717191in}{0.576831in}}%
\pgfpathlineto{\pgfqpoint{5.710400in}{0.570869in}}%
\pgfpathlineto{\pgfqpoint{5.700936in}{0.571500in}}%
\pgfpathlineto{\pgfqpoint{5.691987in}{0.578887in}}%
\pgfpathlineto{\pgfqpoint{5.690392in}{0.592032in}}%
\pgfpathlineto{\pgfqpoint{5.683727in}{0.599677in}}%
\pgfpathlineto{\pgfqpoint{5.679190in}{0.572315in}}%
\pgfpathlineto{\pgfqpoint{5.663775in}{0.574915in}}%
\pgfpathlineto{\pgfqpoint{5.658415in}{0.622761in}}%
\pgfpathlineto{\pgfqpoint{5.652610in}{0.673185in}}%
\pgfpathlineto{\pgfqpoint{5.654766in}{0.745935in}}%
\pgfpathlineto{\pgfqpoint{5.656104in}{0.795857in}}%
\pgfpathlineto{\pgfqpoint{5.658950in}{0.872025in}}%
\pgfpathlineto{\pgfqpoint{5.652858in}{0.878728in}}%
\pgfpathlineto{\pgfqpoint{5.663945in}{0.880335in}}%
\pgfpathlineto{\pgfqpoint{5.728533in}{0.884512in}}%
\pgfpathlineto{\pgfqpoint{5.791375in}{0.890163in}}%
\pgfpathlineto{\pgfqpoint{5.807908in}{0.832251in}}%
\pgfpathlineto{\pgfqpoint{5.819184in}{0.789743in}}%
\pgfpathlineto{\pgfqpoint{5.829291in}{0.754156in}}%
\pgfpathlineto{\pgfqpoint{5.835136in}{0.739826in}}%
\pgfpathlineto{\pgfqpoint{5.844340in}{0.726491in}}%
\pgfpathlineto{\pgfqpoint{5.842853in}{0.719705in}}%
\pgfpathlineto{\pgfqpoint{5.849428in}{0.716398in}}%
\pgfpathlineto{\pgfqpoint{5.841659in}{0.706111in}}%
\pgfpathlineto{\pgfqpoint{5.839034in}{0.687792in}}%
\pgfpathlineto{\pgfqpoint{5.846628in}{0.666365in}}%
\pgfpathlineto{\pgfqpoint{5.845573in}{0.644805in}}%
\pgfpathlineto{\pgfqpoint{5.852687in}{0.632577in}}%
\pgfpathclose%
\pgfusepath{fill}%
\end{pgfscope}%
\begin{pgfscope}%
\pgfpathrectangle{\pgfqpoint{3.625000in}{0.100000in}}{\pgfqpoint{2.989028in}{1.913466in}}%
\pgfusepath{clip}%
\pgfsetbuttcap%
\pgfsetmiterjoin%
\definecolor{currentfill}{rgb}{0.963476,0.985390,0.692042}%
\pgfsetfillcolor{currentfill}%
\pgfsetlinewidth{0.000000pt}%
\definecolor{currentstroke}{rgb}{0.000000,0.000000,0.000000}%
\pgfsetstrokecolor{currentstroke}%
\pgfsetstrokeopacity{0.000000}%
\pgfsetdash{}{0pt}%
\pgfpathmoveto{\pgfqpoint{5.663775in}{0.574915in}}%
\pgfpathlineto{\pgfqpoint{5.647957in}{0.570389in}}%
\pgfpathlineto{\pgfqpoint{5.636166in}{0.573295in}}%
\pgfpathlineto{\pgfqpoint{5.614268in}{0.566330in}}%
\pgfpathlineto{\pgfqpoint{5.597753in}{0.557360in}}%
\pgfpathlineto{\pgfqpoint{5.590977in}{0.573569in}}%
\pgfpathlineto{\pgfqpoint{5.583992in}{0.580363in}}%
\pgfpathlineto{\pgfqpoint{5.580622in}{0.588493in}}%
\pgfpathlineto{\pgfqpoint{5.585562in}{0.610498in}}%
\pgfpathlineto{\pgfqpoint{5.538761in}{0.607617in}}%
\pgfpathlineto{\pgfqpoint{5.478073in}{0.604986in}}%
\pgfpathlineto{\pgfqpoint{5.481684in}{0.610463in}}%
\pgfpathlineto{\pgfqpoint{5.477221in}{0.620808in}}%
\pgfpathlineto{\pgfqpoint{5.481668in}{0.622916in}}%
\pgfpathlineto{\pgfqpoint{5.480681in}{0.632856in}}%
\pgfpathlineto{\pgfqpoint{5.488432in}{0.642495in}}%
\pgfpathlineto{\pgfqpoint{5.492673in}{0.661207in}}%
\pgfpathlineto{\pgfqpoint{5.506859in}{0.673535in}}%
\pgfpathlineto{\pgfqpoint{5.501601in}{0.685508in}}%
\pgfpathlineto{\pgfqpoint{5.511660in}{0.692400in}}%
\pgfpathlineto{\pgfqpoint{5.502983in}{0.704877in}}%
\pgfpathlineto{\pgfqpoint{5.496703in}{0.733439in}}%
\pgfpathlineto{\pgfqpoint{5.499085in}{0.738625in}}%
\pgfpathlineto{\pgfqpoint{5.502845in}{0.748573in}}%
\pgfpathlineto{\pgfqpoint{5.499345in}{0.759012in}}%
\pgfpathlineto{\pgfqpoint{5.501057in}{0.763762in}}%
\pgfpathlineto{\pgfqpoint{5.494899in}{0.781706in}}%
\pgfpathlineto{\pgfqpoint{5.513077in}{0.813926in}}%
\pgfpathlineto{\pgfqpoint{5.513922in}{0.822501in}}%
\pgfpathlineto{\pgfqpoint{5.522678in}{0.828737in}}%
\pgfpathlineto{\pgfqpoint{5.532021in}{0.849323in}}%
\pgfpathlineto{\pgfqpoint{5.531577in}{0.857690in}}%
\pgfpathlineto{\pgfqpoint{5.543214in}{0.866243in}}%
\pgfpathlineto{\pgfqpoint{5.539584in}{0.871310in}}%
\pgfpathlineto{\pgfqpoint{5.609319in}{0.875379in}}%
\pgfpathlineto{\pgfqpoint{5.652858in}{0.878728in}}%
\pgfpathlineto{\pgfqpoint{5.658950in}{0.872025in}}%
\pgfpathlineto{\pgfqpoint{5.656104in}{0.795857in}}%
\pgfpathlineto{\pgfqpoint{5.654766in}{0.745935in}}%
\pgfpathlineto{\pgfqpoint{5.652610in}{0.673185in}}%
\pgfpathlineto{\pgfqpoint{5.658415in}{0.622761in}}%
\pgfpathlineto{\pgfqpoint{5.663775in}{0.574915in}}%
\pgfpathclose%
\pgfusepath{fill}%
\end{pgfscope}%
\begin{pgfscope}%
\pgfpathrectangle{\pgfqpoint{3.625000in}{0.100000in}}{\pgfqpoint{2.989028in}{1.913466in}}%
\pgfusepath{clip}%
\pgfsetbuttcap%
\pgfsetmiterjoin%
\definecolor{currentfill}{rgb}{0.936563,0.974625,0.650058}%
\pgfsetfillcolor{currentfill}%
\pgfsetlinewidth{0.000000pt}%
\definecolor{currentstroke}{rgb}{0.000000,0.000000,0.000000}%
\pgfsetstrokecolor{currentstroke}%
\pgfsetstrokeopacity{0.000000}%
\pgfsetdash{}{0pt}%
\pgfpathmoveto{\pgfqpoint{5.791375in}{0.890163in}}%
\pgfpathlineto{\pgfqpoint{5.860069in}{0.897708in}}%
\pgfpathlineto{\pgfqpoint{5.901402in}{0.902801in}}%
\pgfpathlineto{\pgfqpoint{5.924943in}{0.906551in}}%
\pgfpathlineto{\pgfqpoint{5.915192in}{0.891171in}}%
\pgfpathlineto{\pgfqpoint{5.915198in}{0.883919in}}%
\pgfpathlineto{\pgfqpoint{5.925235in}{0.880009in}}%
\pgfpathlineto{\pgfqpoint{5.932061in}{0.873695in}}%
\pgfpathlineto{\pgfqpoint{5.942377in}{0.872915in}}%
\pgfpathlineto{\pgfqpoint{5.952019in}{0.855143in}}%
\pgfpathlineto{\pgfqpoint{5.962478in}{0.842609in}}%
\pgfpathlineto{\pgfqpoint{5.981053in}{0.832075in}}%
\pgfpathlineto{\pgfqpoint{5.986332in}{0.823674in}}%
\pgfpathlineto{\pgfqpoint{6.002896in}{0.814587in}}%
\pgfpathlineto{\pgfqpoint{6.004142in}{0.807998in}}%
\pgfpathlineto{\pgfqpoint{6.017579in}{0.794896in}}%
\pgfpathlineto{\pgfqpoint{6.027872in}{0.791165in}}%
\pgfpathlineto{\pgfqpoint{6.035161in}{0.778793in}}%
\pgfpathlineto{\pgfqpoint{6.038387in}{0.764886in}}%
\pgfpathlineto{\pgfqpoint{6.049078in}{0.759511in}}%
\pgfpathlineto{\pgfqpoint{6.057685in}{0.744563in}}%
\pgfpathlineto{\pgfqpoint{6.060460in}{0.733570in}}%
\pgfpathlineto{\pgfqpoint{6.073073in}{0.728890in}}%
\pgfpathlineto{\pgfqpoint{6.062490in}{0.708626in}}%
\pgfpathlineto{\pgfqpoint{6.058507in}{0.696651in}}%
\pgfpathlineto{\pgfqpoint{6.061105in}{0.691037in}}%
\pgfpathlineto{\pgfqpoint{6.056568in}{0.681708in}}%
\pgfpathlineto{\pgfqpoint{6.054639in}{0.668810in}}%
\pgfpathlineto{\pgfqpoint{6.046570in}{0.666408in}}%
\pgfpathlineto{\pgfqpoint{6.050831in}{0.654599in}}%
\pgfpathlineto{\pgfqpoint{6.050559in}{0.639617in}}%
\pgfpathlineto{\pgfqpoint{6.036217in}{0.640178in}}%
\pgfpathlineto{\pgfqpoint{6.026225in}{0.642912in}}%
\pgfpathlineto{\pgfqpoint{6.021902in}{0.637830in}}%
\pgfpathlineto{\pgfqpoint{6.025128in}{0.625826in}}%
\pgfpathlineto{\pgfqpoint{6.024427in}{0.611870in}}%
\pgfpathlineto{\pgfqpoint{6.018123in}{0.610802in}}%
\pgfpathlineto{\pgfqpoint{6.013024in}{0.623836in}}%
\pgfpathlineto{\pgfqpoint{5.961103in}{0.620376in}}%
\pgfpathlineto{\pgfqpoint{5.895635in}{0.616804in}}%
\pgfpathlineto{\pgfqpoint{5.862617in}{0.614480in}}%
\pgfpathlineto{\pgfqpoint{5.852687in}{0.632577in}}%
\pgfpathlineto{\pgfqpoint{5.845573in}{0.644805in}}%
\pgfpathlineto{\pgfqpoint{5.846628in}{0.666365in}}%
\pgfpathlineto{\pgfqpoint{5.839034in}{0.687792in}}%
\pgfpathlineto{\pgfqpoint{5.841659in}{0.706111in}}%
\pgfpathlineto{\pgfqpoint{5.849428in}{0.716398in}}%
\pgfpathlineto{\pgfqpoint{5.842853in}{0.719705in}}%
\pgfpathlineto{\pgfqpoint{5.844340in}{0.726491in}}%
\pgfpathlineto{\pgfqpoint{5.835136in}{0.739826in}}%
\pgfpathlineto{\pgfqpoint{5.829291in}{0.754156in}}%
\pgfpathlineto{\pgfqpoint{5.819184in}{0.789743in}}%
\pgfpathlineto{\pgfqpoint{5.807908in}{0.832251in}}%
\pgfpathlineto{\pgfqpoint{5.791375in}{0.890163in}}%
\pgfpathclose%
\pgfusepath{fill}%
\end{pgfscope}%
\begin{pgfscope}%
\pgfpathrectangle{\pgfqpoint{3.625000in}{0.100000in}}{\pgfqpoint{2.989028in}{1.913466in}}%
\pgfusepath{clip}%
\pgfsetbuttcap%
\pgfsetmiterjoin%
\definecolor{currentfill}{rgb}{0.559170,0.823068,0.644752}%
\pgfsetfillcolor{currentfill}%
\pgfsetlinewidth{0.000000pt}%
\definecolor{currentstroke}{rgb}{0.000000,0.000000,0.000000}%
\pgfsetstrokecolor{currentstroke}%
\pgfsetstrokeopacity{0.000000}%
\pgfsetdash{}{0pt}%
\pgfpathmoveto{\pgfqpoint{5.924943in}{0.906551in}}%
\pgfpathlineto{\pgfqpoint{5.952183in}{0.920102in}}%
\pgfpathlineto{\pgfqpoint{5.967269in}{0.925218in}}%
\pgfpathlineto{\pgfqpoint{6.032904in}{0.932038in}}%
\pgfpathlineto{\pgfqpoint{6.039855in}{0.929779in}}%
\pgfpathlineto{\pgfqpoint{6.049037in}{0.920550in}}%
\pgfpathlineto{\pgfqpoint{6.049543in}{0.912354in}}%
\pgfpathlineto{\pgfqpoint{6.109952in}{0.921230in}}%
\pgfpathlineto{\pgfqpoint{6.178548in}{0.872057in}}%
\pgfpathlineto{\pgfqpoint{6.165684in}{0.858649in}}%
\pgfpathlineto{\pgfqpoint{6.148033in}{0.827511in}}%
\pgfpathlineto{\pgfqpoint{6.153057in}{0.820819in}}%
\pgfpathlineto{\pgfqpoint{6.143641in}{0.807830in}}%
\pgfpathlineto{\pgfqpoint{6.134275in}{0.806359in}}%
\pgfpathlineto{\pgfqpoint{6.134240in}{0.798554in}}%
\pgfpathlineto{\pgfqpoint{6.127464in}{0.790386in}}%
\pgfpathlineto{\pgfqpoint{6.119032in}{0.788730in}}%
\pgfpathlineto{\pgfqpoint{6.120872in}{0.781473in}}%
\pgfpathlineto{\pgfqpoint{6.116167in}{0.775902in}}%
\pgfpathlineto{\pgfqpoint{6.104911in}{0.771089in}}%
\pgfpathlineto{\pgfqpoint{6.091192in}{0.760141in}}%
\pgfpathlineto{\pgfqpoint{6.093782in}{0.753058in}}%
\pgfpathlineto{\pgfqpoint{6.085148in}{0.748611in}}%
\pgfpathlineto{\pgfqpoint{6.077866in}{0.753287in}}%
\pgfpathlineto{\pgfqpoint{6.072540in}{0.732958in}}%
\pgfpathlineto{\pgfqpoint{6.060460in}{0.733570in}}%
\pgfpathlineto{\pgfqpoint{6.057685in}{0.744563in}}%
\pgfpathlineto{\pgfqpoint{6.049078in}{0.759511in}}%
\pgfpathlineto{\pgfqpoint{6.038387in}{0.764886in}}%
\pgfpathlineto{\pgfqpoint{6.035161in}{0.778793in}}%
\pgfpathlineto{\pgfqpoint{6.027872in}{0.791165in}}%
\pgfpathlineto{\pgfqpoint{6.017579in}{0.794896in}}%
\pgfpathlineto{\pgfqpoint{6.004142in}{0.807998in}}%
\pgfpathlineto{\pgfqpoint{6.002896in}{0.814587in}}%
\pgfpathlineto{\pgfqpoint{5.986332in}{0.823674in}}%
\pgfpathlineto{\pgfqpoint{5.981053in}{0.832075in}}%
\pgfpathlineto{\pgfqpoint{5.962478in}{0.842609in}}%
\pgfpathlineto{\pgfqpoint{5.952019in}{0.855143in}}%
\pgfpathlineto{\pgfqpoint{5.942377in}{0.872915in}}%
\pgfpathlineto{\pgfqpoint{5.932061in}{0.873695in}}%
\pgfpathlineto{\pgfqpoint{5.925235in}{0.880009in}}%
\pgfpathlineto{\pgfqpoint{5.915198in}{0.883919in}}%
\pgfpathlineto{\pgfqpoint{5.915192in}{0.891171in}}%
\pgfpathlineto{\pgfqpoint{5.924943in}{0.906551in}}%
\pgfpathclose%
\pgfusepath{fill}%
\end{pgfscope}%
\begin{pgfscope}%
\pgfpathrectangle{\pgfqpoint{3.625000in}{0.100000in}}{\pgfqpoint{2.989028in}{1.913466in}}%
\pgfusepath{clip}%
\pgfsetbuttcap%
\pgfsetmiterjoin%
\definecolor{currentfill}{rgb}{0.847520,0.938639,0.607151}%
\pgfsetfillcolor{currentfill}%
\pgfsetlinewidth{0.000000pt}%
\definecolor{currentstroke}{rgb}{0.000000,0.000000,0.000000}%
\pgfsetstrokecolor{currentstroke}%
\pgfsetstrokeopacity{0.000000}%
\pgfsetdash{}{0pt}%
\pgfpathmoveto{\pgfqpoint{5.307462in}{0.962463in}}%
\pgfpathlineto{\pgfqpoint{5.347491in}{0.962824in}}%
\pgfpathlineto{\pgfqpoint{5.400448in}{0.963789in}}%
\pgfpathlineto{\pgfqpoint{5.490941in}{0.966662in}}%
\pgfpathlineto{\pgfqpoint{5.542700in}{0.969399in}}%
\pgfpathlineto{\pgfqpoint{5.548251in}{0.962527in}}%
\pgfpathlineto{\pgfqpoint{5.547904in}{0.955269in}}%
\pgfpathlineto{\pgfqpoint{5.535365in}{0.942748in}}%
\pgfpathlineto{\pgfqpoint{5.532334in}{0.935892in}}%
\pgfpathlineto{\pgfqpoint{5.567109in}{0.938464in}}%
\pgfpathlineto{\pgfqpoint{5.567090in}{0.925797in}}%
\pgfpathlineto{\pgfqpoint{5.555934in}{0.920384in}}%
\pgfpathlineto{\pgfqpoint{5.556092in}{0.911777in}}%
\pgfpathlineto{\pgfqpoint{5.551902in}{0.905590in}}%
\pgfpathlineto{\pgfqpoint{5.549159in}{0.892449in}}%
\pgfpathlineto{\pgfqpoint{5.551990in}{0.881608in}}%
\pgfpathlineto{\pgfqpoint{5.539584in}{0.871310in}}%
\pgfpathlineto{\pgfqpoint{5.543214in}{0.866243in}}%
\pgfpathlineto{\pgfqpoint{5.531577in}{0.857690in}}%
\pgfpathlineto{\pgfqpoint{5.532021in}{0.849323in}}%
\pgfpathlineto{\pgfqpoint{5.522678in}{0.828737in}}%
\pgfpathlineto{\pgfqpoint{5.513922in}{0.822501in}}%
\pgfpathlineto{\pgfqpoint{5.513077in}{0.813926in}}%
\pgfpathlineto{\pgfqpoint{5.494899in}{0.781706in}}%
\pgfpathlineto{\pgfqpoint{5.501057in}{0.763762in}}%
\pgfpathlineto{\pgfqpoint{5.499345in}{0.759012in}}%
\pgfpathlineto{\pgfqpoint{5.502845in}{0.748573in}}%
\pgfpathlineto{\pgfqpoint{5.499085in}{0.738625in}}%
\pgfpathlineto{\pgfqpoint{5.449378in}{0.736573in}}%
\pgfpathlineto{\pgfqpoint{5.384848in}{0.735510in}}%
\pgfpathlineto{\pgfqpoint{5.340340in}{0.735108in}}%
\pgfpathlineto{\pgfqpoint{5.340123in}{0.770153in}}%
\pgfpathlineto{\pgfqpoint{5.329070in}{0.772449in}}%
\pgfpathlineto{\pgfqpoint{5.321775in}{0.769451in}}%
\pgfpathlineto{\pgfqpoint{5.315948in}{0.774957in}}%
\pgfpathlineto{\pgfqpoint{5.316549in}{0.812079in}}%
\pgfpathlineto{\pgfqpoint{5.317853in}{0.891120in}}%
\pgfpathlineto{\pgfqpoint{5.311531in}{0.937404in}}%
\pgfpathlineto{\pgfqpoint{5.307462in}{0.962463in}}%
\pgfpathclose%
\pgfusepath{fill}%
\end{pgfscope}%
\begin{pgfscope}%
\pgfpathrectangle{\pgfqpoint{3.625000in}{0.100000in}}{\pgfqpoint{2.989028in}{1.913466in}}%
\pgfusepath{clip}%
\pgfsetbuttcap%
\pgfsetmiterjoin%
\definecolor{currentfill}{rgb}{0.548558,0.818916,0.644906}%
\pgfsetfillcolor{currentfill}%
\pgfsetlinewidth{0.000000pt}%
\definecolor{currentstroke}{rgb}{0.000000,0.000000,0.000000}%
\pgfsetstrokecolor{currentstroke}%
\pgfsetstrokeopacity{0.000000}%
\pgfsetdash{}{0pt}%
\pgfpathmoveto{\pgfqpoint{5.340340in}{0.735108in}}%
\pgfpathlineto{\pgfqpoint{5.384848in}{0.735510in}}%
\pgfpathlineto{\pgfqpoint{5.449378in}{0.736573in}}%
\pgfpathlineto{\pgfqpoint{5.499085in}{0.738625in}}%
\pgfpathlineto{\pgfqpoint{5.496703in}{0.733439in}}%
\pgfpathlineto{\pgfqpoint{5.502983in}{0.704877in}}%
\pgfpathlineto{\pgfqpoint{5.511660in}{0.692400in}}%
\pgfpathlineto{\pgfqpoint{5.501601in}{0.685508in}}%
\pgfpathlineto{\pgfqpoint{5.506859in}{0.673535in}}%
\pgfpathlineto{\pgfqpoint{5.492673in}{0.661207in}}%
\pgfpathlineto{\pgfqpoint{5.488432in}{0.642495in}}%
\pgfpathlineto{\pgfqpoint{5.480681in}{0.632856in}}%
\pgfpathlineto{\pgfqpoint{5.481668in}{0.622916in}}%
\pgfpathlineto{\pgfqpoint{5.477221in}{0.620808in}}%
\pgfpathlineto{\pgfqpoint{5.481684in}{0.610463in}}%
\pgfpathlineto{\pgfqpoint{5.478073in}{0.604986in}}%
\pgfpathlineto{\pgfqpoint{5.538761in}{0.607617in}}%
\pgfpathlineto{\pgfqpoint{5.585562in}{0.610498in}}%
\pgfpathlineto{\pgfqpoint{5.580622in}{0.588493in}}%
\pgfpathlineto{\pgfqpoint{5.583992in}{0.580363in}}%
\pgfpathlineto{\pgfqpoint{5.590977in}{0.573569in}}%
\pgfpathlineto{\pgfqpoint{5.597753in}{0.557360in}}%
\pgfpathlineto{\pgfqpoint{5.576350in}{0.561094in}}%
\pgfpathlineto{\pgfqpoint{5.568455in}{0.567234in}}%
\pgfpathlineto{\pgfqpoint{5.559048in}{0.567526in}}%
\pgfpathlineto{\pgfqpoint{5.549158in}{0.554082in}}%
\pgfpathlineto{\pgfqpoint{5.551129in}{0.547960in}}%
\pgfpathlineto{\pgfqpoint{5.582724in}{0.544257in}}%
\pgfpathlineto{\pgfqpoint{5.590995in}{0.537220in}}%
\pgfpathlineto{\pgfqpoint{5.604454in}{0.533600in}}%
\pgfpathlineto{\pgfqpoint{5.598576in}{0.525268in}}%
\pgfpathlineto{\pgfqpoint{5.588672in}{0.518003in}}%
\pgfpathlineto{\pgfqpoint{5.614813in}{0.501664in}}%
\pgfpathlineto{\pgfqpoint{5.626990in}{0.499111in}}%
\pgfpathlineto{\pgfqpoint{5.632054in}{0.485755in}}%
\pgfpathlineto{\pgfqpoint{5.620589in}{0.483298in}}%
\pgfpathlineto{\pgfqpoint{5.599041in}{0.500096in}}%
\pgfpathlineto{\pgfqpoint{5.590616in}{0.502430in}}%
\pgfpathlineto{\pgfqpoint{5.586539in}{0.509066in}}%
\pgfpathlineto{\pgfqpoint{5.574043in}{0.506343in}}%
\pgfpathlineto{\pgfqpoint{5.570130in}{0.497792in}}%
\pgfpathlineto{\pgfqpoint{5.572608in}{0.488264in}}%
\pgfpathlineto{\pgfqpoint{5.564195in}{0.482642in}}%
\pgfpathlineto{\pgfqpoint{5.556517in}{0.496493in}}%
\pgfpathlineto{\pgfqpoint{5.543069in}{0.492363in}}%
\pgfpathlineto{\pgfqpoint{5.538030in}{0.484098in}}%
\pgfpathlineto{\pgfqpoint{5.528479in}{0.486449in}}%
\pgfpathlineto{\pgfqpoint{5.522373in}{0.496886in}}%
\pgfpathlineto{\pgfqpoint{5.506173in}{0.500333in}}%
\pgfpathlineto{\pgfqpoint{5.503115in}{0.505764in}}%
\pgfpathlineto{\pgfqpoint{5.486236in}{0.515225in}}%
\pgfpathlineto{\pgfqpoint{5.482047in}{0.523496in}}%
\pgfpathlineto{\pgfqpoint{5.467910in}{0.520116in}}%
\pgfpathlineto{\pgfqpoint{5.469705in}{0.527689in}}%
\pgfpathlineto{\pgfqpoint{5.452134in}{0.519920in}}%
\pgfpathlineto{\pgfqpoint{5.456854in}{0.511860in}}%
\pgfpathlineto{\pgfqpoint{5.443284in}{0.507095in}}%
\pgfpathlineto{\pgfqpoint{5.425311in}{0.509714in}}%
\pgfpathlineto{\pgfqpoint{5.388934in}{0.522174in}}%
\pgfpathlineto{\pgfqpoint{5.360864in}{0.519701in}}%
\pgfpathlineto{\pgfqpoint{5.358432in}{0.536068in}}%
\pgfpathlineto{\pgfqpoint{5.361956in}{0.543537in}}%
\pgfpathlineto{\pgfqpoint{5.361655in}{0.555420in}}%
\pgfpathlineto{\pgfqpoint{5.358187in}{0.558251in}}%
\pgfpathlineto{\pgfqpoint{5.359363in}{0.571815in}}%
\pgfpathlineto{\pgfqpoint{5.369586in}{0.590680in}}%
\pgfpathlineto{\pgfqpoint{5.370894in}{0.597808in}}%
\pgfpathlineto{\pgfqpoint{5.369209in}{0.614644in}}%
\pgfpathlineto{\pgfqpoint{5.361688in}{0.622574in}}%
\pgfpathlineto{\pgfqpoint{5.357761in}{0.637381in}}%
\pgfpathlineto{\pgfqpoint{5.352929in}{0.640781in}}%
\pgfpathlineto{\pgfqpoint{5.355092in}{0.648811in}}%
\pgfpathlineto{\pgfqpoint{5.348939in}{0.660813in}}%
\pgfpathlineto{\pgfqpoint{5.341262in}{0.667326in}}%
\pgfpathlineto{\pgfqpoint{5.340340in}{0.735108in}}%
\pgfpathclose%
\pgfusepath{fill}%
\end{pgfscope}%
\begin{pgfscope}%
\pgfpathrectangle{\pgfqpoint{3.625000in}{0.100000in}}{\pgfqpoint{2.989028in}{1.913466in}}%
\pgfusepath{clip}%
\pgfsetbuttcap%
\pgfsetmiterjoin%
\definecolor{currentfill}{rgb}{0.548558,0.818916,0.644906}%
\pgfsetfillcolor{currentfill}%
\pgfsetlinewidth{0.000000pt}%
\definecolor{currentstroke}{rgb}{0.000000,0.000000,0.000000}%
\pgfsetstrokecolor{currentstroke}%
\pgfsetstrokeopacity{0.000000}%
\pgfsetdash{}{0pt}%
\pgfpathmoveto{\pgfqpoint{5.459378in}{0.511243in}}%
\pgfpathlineto{\pgfqpoint{5.465817in}{0.515073in}}%
\pgfpathlineto{\pgfqpoint{5.473633in}{0.510545in}}%
\pgfpathlineto{\pgfqpoint{5.469274in}{0.504314in}}%
\pgfpathlineto{\pgfqpoint{5.459378in}{0.511243in}}%
\pgfpathclose%
\pgfusepath{fill}%
\end{pgfscope}%
\begin{pgfscope}%
\pgfpathrectangle{\pgfqpoint{3.625000in}{0.100000in}}{\pgfqpoint{2.989028in}{1.913466in}}%
\pgfusepath{clip}%
\pgfsetbuttcap%
\pgfsetmiterjoin%
\definecolor{currentfill}{rgb}{0.384006,0.742945,0.654441}%
\pgfsetfillcolor{currentfill}%
\pgfsetlinewidth{0.000000pt}%
\definecolor{currentstroke}{rgb}{0.000000,0.000000,0.000000}%
\pgfsetstrokecolor{currentstroke}%
\pgfsetstrokeopacity{0.000000}%
\pgfsetdash{}{0pt}%
\pgfpathmoveto{\pgfqpoint{5.719230in}{0.585103in}}%
\pgfpathlineto{\pgfqpoint{5.719346in}{0.597485in}}%
\pgfpathlineto{\pgfqpoint{5.711652in}{0.602188in}}%
\pgfpathlineto{\pgfqpoint{5.705349in}{0.610192in}}%
\pgfpathlineto{\pgfqpoint{5.706207in}{0.618602in}}%
\pgfpathlineto{\pgfqpoint{5.774834in}{0.623930in}}%
\pgfpathlineto{\pgfqpoint{5.852687in}{0.632577in}}%
\pgfpathlineto{\pgfqpoint{5.862617in}{0.614480in}}%
\pgfpathlineto{\pgfqpoint{5.895635in}{0.616804in}}%
\pgfpathlineto{\pgfqpoint{5.961103in}{0.620376in}}%
\pgfpathlineto{\pgfqpoint{6.013024in}{0.623836in}}%
\pgfpathlineto{\pgfqpoint{6.018123in}{0.610802in}}%
\pgfpathlineto{\pgfqpoint{6.024427in}{0.611870in}}%
\pgfpathlineto{\pgfqpoint{6.025128in}{0.625826in}}%
\pgfpathlineto{\pgfqpoint{6.021902in}{0.637830in}}%
\pgfpathlineto{\pgfqpoint{6.026225in}{0.642912in}}%
\pgfpathlineto{\pgfqpoint{6.036217in}{0.640178in}}%
\pgfpathlineto{\pgfqpoint{6.050559in}{0.639617in}}%
\pgfpathlineto{\pgfqpoint{6.052764in}{0.628892in}}%
\pgfpathlineto{\pgfqpoint{6.057176in}{0.622764in}}%
\pgfpathlineto{\pgfqpoint{6.060621in}{0.609356in}}%
\pgfpathlineto{\pgfqpoint{6.071296in}{0.588607in}}%
\pgfpathlineto{\pgfqpoint{6.071349in}{0.582995in}}%
\pgfpathlineto{\pgfqpoint{6.087126in}{0.558631in}}%
\pgfpathlineto{\pgfqpoint{6.088652in}{0.553598in}}%
\pgfpathlineto{\pgfqpoint{6.112320in}{0.519188in}}%
\pgfpathlineto{\pgfqpoint{6.108570in}{0.518632in}}%
\pgfpathlineto{\pgfqpoint{6.118548in}{0.494158in}}%
\pgfpathlineto{\pgfqpoint{6.145959in}{0.451609in}}%
\pgfpathlineto{\pgfqpoint{6.161873in}{0.420325in}}%
\pgfpathlineto{\pgfqpoint{6.167192in}{0.414848in}}%
\pgfpathlineto{\pgfqpoint{6.176306in}{0.395211in}}%
\pgfpathlineto{\pgfqpoint{6.179468in}{0.363769in}}%
\pgfpathlineto{\pgfqpoint{6.180727in}{0.340258in}}%
\pgfpathlineto{\pgfqpoint{6.179207in}{0.325205in}}%
\pgfpathlineto{\pgfqpoint{6.174320in}{0.314444in}}%
\pgfpathlineto{\pgfqpoint{6.176590in}{0.300332in}}%
\pgfpathlineto{\pgfqpoint{6.171329in}{0.289163in}}%
\pgfpathlineto{\pgfqpoint{6.155707in}{0.280003in}}%
\pgfpathlineto{\pgfqpoint{6.145547in}{0.280672in}}%
\pgfpathlineto{\pgfqpoint{6.138949in}{0.275884in}}%
\pgfpathlineto{\pgfqpoint{6.137059in}{0.288668in}}%
\pgfpathlineto{\pgfqpoint{6.126137in}{0.292042in}}%
\pgfpathlineto{\pgfqpoint{6.116346in}{0.309880in}}%
\pgfpathlineto{\pgfqpoint{6.115241in}{0.318003in}}%
\pgfpathlineto{\pgfqpoint{6.097719in}{0.322936in}}%
\pgfpathlineto{\pgfqpoint{6.086428in}{0.321868in}}%
\pgfpathlineto{\pgfqpoint{6.080040in}{0.333654in}}%
\pgfpathlineto{\pgfqpoint{6.072725in}{0.354839in}}%
\pgfpathlineto{\pgfqpoint{6.062575in}{0.359133in}}%
\pgfpathlineto{\pgfqpoint{6.057072in}{0.371273in}}%
\pgfpathlineto{\pgfqpoint{6.043614in}{0.378428in}}%
\pgfpathlineto{\pgfqpoint{6.035835in}{0.387377in}}%
\pgfpathlineto{\pgfqpoint{6.022416in}{0.411708in}}%
\pgfpathlineto{\pgfqpoint{6.020838in}{0.428833in}}%
\pgfpathlineto{\pgfqpoint{6.028150in}{0.439789in}}%
\pgfpathlineto{\pgfqpoint{6.027419in}{0.447403in}}%
\pgfpathlineto{\pgfqpoint{6.018969in}{0.448217in}}%
\pgfpathlineto{\pgfqpoint{6.011909in}{0.453513in}}%
\pgfpathlineto{\pgfqpoint{6.008042in}{0.447046in}}%
\pgfpathlineto{\pgfqpoint{6.014818in}{0.441777in}}%
\pgfpathlineto{\pgfqpoint{6.009887in}{0.432302in}}%
\pgfpathlineto{\pgfqpoint{6.001908in}{0.440166in}}%
\pgfpathlineto{\pgfqpoint{6.002834in}{0.461974in}}%
\pgfpathlineto{\pgfqpoint{6.006682in}{0.479667in}}%
\pgfpathlineto{\pgfqpoint{6.004704in}{0.510026in}}%
\pgfpathlineto{\pgfqpoint{5.996751in}{0.517257in}}%
\pgfpathlineto{\pgfqpoint{5.992752in}{0.526536in}}%
\pgfpathlineto{\pgfqpoint{5.979052in}{0.526329in}}%
\pgfpathlineto{\pgfqpoint{5.965510in}{0.541597in}}%
\pgfpathlineto{\pgfqpoint{5.956430in}{0.546173in}}%
\pgfpathlineto{\pgfqpoint{5.953737in}{0.555842in}}%
\pgfpathlineto{\pgfqpoint{5.944856in}{0.559238in}}%
\pgfpathlineto{\pgfqpoint{5.937496in}{0.569913in}}%
\pgfpathlineto{\pgfqpoint{5.918052in}{0.578628in}}%
\pgfpathlineto{\pgfqpoint{5.902937in}{0.578850in}}%
\pgfpathlineto{\pgfqpoint{5.896356in}{0.575503in}}%
\pgfpathlineto{\pgfqpoint{5.897991in}{0.565056in}}%
\pgfpathlineto{\pgfqpoint{5.891130in}{0.565551in}}%
\pgfpathlineto{\pgfqpoint{5.870022in}{0.550943in}}%
\pgfpathlineto{\pgfqpoint{5.844668in}{0.545153in}}%
\pgfpathlineto{\pgfqpoint{5.844252in}{0.552317in}}%
\pgfpathlineto{\pgfqpoint{5.838633in}{0.559318in}}%
\pgfpathlineto{\pgfqpoint{5.823558in}{0.568984in}}%
\pgfpathlineto{\pgfqpoint{5.801918in}{0.578894in}}%
\pgfpathlineto{\pgfqpoint{5.778414in}{0.584141in}}%
\pgfpathlineto{\pgfqpoint{5.774002in}{0.591275in}}%
\pgfpathlineto{\pgfqpoint{5.765528in}{0.585325in}}%
\pgfpathlineto{\pgfqpoint{5.755323in}{0.584012in}}%
\pgfpathlineto{\pgfqpoint{5.719801in}{0.574653in}}%
\pgfpathlineto{\pgfqpoint{5.719230in}{0.585103in}}%
\pgfpathclose%
\pgfusepath{fill}%
\end{pgfscope}%
\begin{pgfscope}%
\pgfpathrectangle{\pgfqpoint{3.625000in}{0.100000in}}{\pgfqpoint{2.989028in}{1.913466in}}%
\pgfusepath{clip}%
\pgfsetbuttcap%
\pgfsetmiterjoin%
\definecolor{currentfill}{rgb}{0.384006,0.742945,0.654441}%
\pgfsetfillcolor{currentfill}%
\pgfsetlinewidth{0.000000pt}%
\definecolor{currentstroke}{rgb}{0.000000,0.000000,0.000000}%
\pgfsetstrokecolor{currentstroke}%
\pgfsetstrokeopacity{0.000000}%
\pgfsetdash{}{0pt}%
\pgfpathmoveto{\pgfqpoint{6.115769in}{0.519678in}}%
\pgfpathlineto{\pgfqpoint{6.124959in}{0.508987in}}%
\pgfpathlineto{\pgfqpoint{6.114344in}{0.508314in}}%
\pgfpathlineto{\pgfqpoint{6.115769in}{0.519678in}}%
\pgfpathclose%
\pgfusepath{fill}%
\end{pgfscope}%
\begin{pgfscope}%
\pgfpathrectangle{\pgfqpoint{3.625000in}{0.100000in}}{\pgfqpoint{2.989028in}{1.913466in}}%
\pgfusepath{clip}%
\pgfsetbuttcap%
\pgfsetmiterjoin%
\definecolor{currentfill}{rgb}{0.874740,0.949712,0.601615}%
\pgfsetfillcolor{currentfill}%
\pgfsetlinewidth{0.000000pt}%
\definecolor{currentstroke}{rgb}{0.000000,0.000000,0.000000}%
\pgfsetstrokecolor{currentstroke}%
\pgfsetstrokeopacity{0.000000}%
\pgfsetdash{}{0pt}%
\pgfpathmoveto{\pgfqpoint{5.574740in}{1.740361in}}%
\pgfpathlineto{\pgfqpoint{5.569753in}{1.730656in}}%
\pgfpathlineto{\pgfqpoint{5.557857in}{1.724992in}}%
\pgfpathlineto{\pgfqpoint{5.552705in}{1.717368in}}%
\pgfpathlineto{\pgfqpoint{5.546669in}{1.722859in}}%
\pgfpathlineto{\pgfqpoint{5.574740in}{1.740361in}}%
\pgfpathclose%
\pgfusepath{fill}%
\end{pgfscope}%
\begin{pgfscope}%
\pgfpathrectangle{\pgfqpoint{3.625000in}{0.100000in}}{\pgfqpoint{2.989028in}{1.913466in}}%
\pgfusepath{clip}%
\pgfsetbuttcap%
\pgfsetmiterjoin%
\definecolor{currentfill}{rgb}{0.874740,0.949712,0.601615}%
\pgfsetfillcolor{currentfill}%
\pgfsetlinewidth{0.000000pt}%
\definecolor{currentstroke}{rgb}{0.000000,0.000000,0.000000}%
\pgfsetstrokecolor{currentstroke}%
\pgfsetstrokeopacity{0.000000}%
\pgfsetdash{}{0pt}%
\pgfpathmoveto{\pgfqpoint{5.578784in}{1.681953in}}%
\pgfpathlineto{\pgfqpoint{5.590956in}{1.693308in}}%
\pgfpathlineto{\pgfqpoint{5.609730in}{1.696281in}}%
\pgfpathlineto{\pgfqpoint{5.604559in}{1.688384in}}%
\pgfpathlineto{\pgfqpoint{5.591688in}{1.676964in}}%
\pgfpathlineto{\pgfqpoint{5.584164in}{1.662306in}}%
\pgfpathlineto{\pgfqpoint{5.580977in}{1.670256in}}%
\pgfpathlineto{\pgfqpoint{5.575318in}{1.671398in}}%
\pgfpathlineto{\pgfqpoint{5.574750in}{1.678576in}}%
\pgfpathlineto{\pgfqpoint{5.578784in}{1.681953in}}%
\pgfpathclose%
\pgfusepath{fill}%
\end{pgfscope}%
\begin{pgfscope}%
\pgfpathrectangle{\pgfqpoint{3.625000in}{0.100000in}}{\pgfqpoint{2.989028in}{1.913466in}}%
\pgfusepath{clip}%
\pgfsetbuttcap%
\pgfsetmiterjoin%
\definecolor{currentfill}{rgb}{0.874740,0.949712,0.601615}%
\pgfsetfillcolor{currentfill}%
\pgfsetlinewidth{0.000000pt}%
\definecolor{currentstroke}{rgb}{0.000000,0.000000,0.000000}%
\pgfsetstrokecolor{currentstroke}%
\pgfsetstrokeopacity{0.000000}%
\pgfsetdash{}{0pt}%
\pgfpathmoveto{\pgfqpoint{5.627264in}{1.543360in}}%
\pgfpathlineto{\pgfqpoint{5.624022in}{1.546957in}}%
\pgfpathlineto{\pgfqpoint{5.627430in}{1.557084in}}%
\pgfpathlineto{\pgfqpoint{5.617314in}{1.557730in}}%
\pgfpathlineto{\pgfqpoint{5.619996in}{1.566462in}}%
\pgfpathlineto{\pgfqpoint{5.618328in}{1.580368in}}%
\pgfpathlineto{\pgfqpoint{5.609251in}{1.585189in}}%
\pgfpathlineto{\pgfqpoint{5.599752in}{1.594947in}}%
\pgfpathlineto{\pgfqpoint{5.585116in}{1.597803in}}%
\pgfpathlineto{\pgfqpoint{5.570871in}{1.597722in}}%
\pgfpathlineto{\pgfqpoint{5.556860in}{1.604646in}}%
\pgfpathlineto{\pgfqpoint{5.509989in}{1.614740in}}%
\pgfpathlineto{\pgfqpoint{5.504867in}{1.625436in}}%
\pgfpathlineto{\pgfqpoint{5.495731in}{1.629085in}}%
\pgfpathlineto{\pgfqpoint{5.512989in}{1.637274in}}%
\pgfpathlineto{\pgfqpoint{5.522738in}{1.647490in}}%
\pgfpathlineto{\pgfqpoint{5.540897in}{1.650260in}}%
\pgfpathlineto{\pgfqpoint{5.552071in}{1.660666in}}%
\pgfpathlineto{\pgfqpoint{5.557934in}{1.661083in}}%
\pgfpathlineto{\pgfqpoint{5.562408in}{1.668504in}}%
\pgfpathlineto{\pgfqpoint{5.573315in}{1.677295in}}%
\pgfpathlineto{\pgfqpoint{5.574265in}{1.671081in}}%
\pgfpathlineto{\pgfqpoint{5.579179in}{1.669798in}}%
\pgfpathlineto{\pgfqpoint{5.582877in}{1.662388in}}%
\pgfpathlineto{\pgfqpoint{5.583543in}{1.649753in}}%
\pgfpathlineto{\pgfqpoint{5.594656in}{1.657304in}}%
\pgfpathlineto{\pgfqpoint{5.607606in}{1.658880in}}%
\pgfpathlineto{\pgfqpoint{5.618660in}{1.654925in}}%
\pgfpathlineto{\pgfqpoint{5.633651in}{1.634360in}}%
\pgfpathlineto{\pgfqpoint{5.650023in}{1.637647in}}%
\pgfpathlineto{\pgfqpoint{5.656676in}{1.632143in}}%
\pgfpathlineto{\pgfqpoint{5.667376in}{1.631651in}}%
\pgfpathlineto{\pgfqpoint{5.674485in}{1.641522in}}%
\pgfpathlineto{\pgfqpoint{5.687974in}{1.650248in}}%
\pgfpathlineto{\pgfqpoint{5.700939in}{1.652980in}}%
\pgfpathlineto{\pgfqpoint{5.717026in}{1.653265in}}%
\pgfpathlineto{\pgfqpoint{5.728782in}{1.660005in}}%
\pgfpathlineto{\pgfqpoint{5.738376in}{1.656901in}}%
\pgfpathlineto{\pgfqpoint{5.740419in}{1.642645in}}%
\pgfpathlineto{\pgfqpoint{5.761002in}{1.640413in}}%
\pgfpathlineto{\pgfqpoint{5.772182in}{1.647086in}}%
\pgfpathlineto{\pgfqpoint{5.779932in}{1.632053in}}%
\pgfpathlineto{\pgfqpoint{5.794719in}{1.614670in}}%
\pgfpathlineto{\pgfqpoint{5.774022in}{1.614771in}}%
\pgfpathlineto{\pgfqpoint{5.767482in}{1.612622in}}%
\pgfpathlineto{\pgfqpoint{5.758519in}{1.615414in}}%
\pgfpathlineto{\pgfqpoint{5.757916in}{1.603380in}}%
\pgfpathlineto{\pgfqpoint{5.741640in}{1.612785in}}%
\pgfpathlineto{\pgfqpoint{5.720708in}{1.615659in}}%
\pgfpathlineto{\pgfqpoint{5.714946in}{1.606500in}}%
\pgfpathlineto{\pgfqpoint{5.687559in}{1.602046in}}%
\pgfpathlineto{\pgfqpoint{5.684409in}{1.594288in}}%
\pgfpathlineto{\pgfqpoint{5.665351in}{1.591960in}}%
\pgfpathlineto{\pgfqpoint{5.659590in}{1.584073in}}%
\pgfpathlineto{\pgfqpoint{5.649509in}{1.581958in}}%
\pgfpathlineto{\pgfqpoint{5.641458in}{1.563256in}}%
\pgfpathlineto{\pgfqpoint{5.631264in}{1.545142in}}%
\pgfpathlineto{\pgfqpoint{5.627264in}{1.543360in}}%
\pgfpathclose%
\pgfusepath{fill}%
\end{pgfscope}%
\begin{pgfscope}%
\pgfpathrectangle{\pgfqpoint{3.625000in}{0.100000in}}{\pgfqpoint{2.989028in}{1.913466in}}%
\pgfusepath{clip}%
\pgfsetbuttcap%
\pgfsetmiterjoin%
\definecolor{currentfill}{rgb}{0.874740,0.949712,0.601615}%
\pgfsetfillcolor{currentfill}%
\pgfsetlinewidth{0.000000pt}%
\definecolor{currentstroke}{rgb}{0.000000,0.000000,0.000000}%
\pgfsetstrokecolor{currentstroke}%
\pgfsetstrokeopacity{0.000000}%
\pgfsetdash{}{0pt}%
\pgfpathmoveto{\pgfqpoint{5.685936in}{1.325898in}}%
\pgfpathlineto{\pgfqpoint{5.695655in}{1.336128in}}%
\pgfpathlineto{\pgfqpoint{5.700083in}{1.350960in}}%
\pgfpathlineto{\pgfqpoint{5.705349in}{1.359557in}}%
\pgfpathlineto{\pgfqpoint{5.708578in}{1.371256in}}%
\pgfpathlineto{\pgfqpoint{5.709581in}{1.394584in}}%
\pgfpathlineto{\pgfqpoint{5.704720in}{1.416940in}}%
\pgfpathlineto{\pgfqpoint{5.688654in}{1.451174in}}%
\pgfpathlineto{\pgfqpoint{5.693005in}{1.461938in}}%
\pgfpathlineto{\pgfqpoint{5.687337in}{1.476816in}}%
\pgfpathlineto{\pgfqpoint{5.697059in}{1.497392in}}%
\pgfpathlineto{\pgfqpoint{5.695475in}{1.520335in}}%
\pgfpathlineto{\pgfqpoint{5.702247in}{1.523223in}}%
\pgfpathlineto{\pgfqpoint{5.703098in}{1.534139in}}%
\pgfpathlineto{\pgfqpoint{5.715119in}{1.541146in}}%
\pgfpathlineto{\pgfqpoint{5.721918in}{1.540015in}}%
\pgfpathlineto{\pgfqpoint{5.723829in}{1.528307in}}%
\pgfpathlineto{\pgfqpoint{5.729135in}{1.527848in}}%
\pgfpathlineto{\pgfqpoint{5.733977in}{1.544735in}}%
\pgfpathlineto{\pgfqpoint{5.732331in}{1.557869in}}%
\pgfpathlineto{\pgfqpoint{5.735479in}{1.565411in}}%
\pgfpathlineto{\pgfqpoint{5.752498in}{1.573202in}}%
\pgfpathlineto{\pgfqpoint{5.744743in}{1.575975in}}%
\pgfpathlineto{\pgfqpoint{5.742216in}{1.582668in}}%
\pgfpathlineto{\pgfqpoint{5.747803in}{1.594421in}}%
\pgfpathlineto{\pgfqpoint{5.758821in}{1.598481in}}%
\pgfpathlineto{\pgfqpoint{5.771615in}{1.591546in}}%
\pgfpathlineto{\pgfqpoint{5.783677in}{1.591460in}}%
\pgfpathlineto{\pgfqpoint{5.789279in}{1.583359in}}%
\pgfpathlineto{\pgfqpoint{5.797727in}{1.583922in}}%
\pgfpathlineto{\pgfqpoint{5.823804in}{1.572661in}}%
\pgfpathlineto{\pgfqpoint{5.829014in}{1.561760in}}%
\pgfpathlineto{\pgfqpoint{5.823324in}{1.557985in}}%
\pgfpathlineto{\pgfqpoint{5.825077in}{1.549815in}}%
\pgfpathlineto{\pgfqpoint{5.830699in}{1.546182in}}%
\pgfpathlineto{\pgfqpoint{5.833831in}{1.536139in}}%
\pgfpathlineto{\pgfqpoint{5.833393in}{1.511672in}}%
\pgfpathlineto{\pgfqpoint{5.825971in}{1.505827in}}%
\pgfpathlineto{\pgfqpoint{5.824305in}{1.492975in}}%
\pgfpathlineto{\pgfqpoint{5.810538in}{1.481098in}}%
\pgfpathlineto{\pgfqpoint{5.811355in}{1.466722in}}%
\pgfpathlineto{\pgfqpoint{5.823421in}{1.461685in}}%
\pgfpathlineto{\pgfqpoint{5.837034in}{1.479623in}}%
\pgfpathlineto{\pgfqpoint{5.838159in}{1.486129in}}%
\pgfpathlineto{\pgfqpoint{5.855136in}{1.496942in}}%
\pgfpathlineto{\pgfqpoint{5.865932in}{1.491933in}}%
\pgfpathlineto{\pgfqpoint{5.872707in}{1.480613in}}%
\pgfpathlineto{\pgfqpoint{5.883636in}{1.441310in}}%
\pgfpathlineto{\pgfqpoint{5.889428in}{1.428876in}}%
\pgfpathlineto{\pgfqpoint{5.887615in}{1.423735in}}%
\pgfpathlineto{\pgfqpoint{5.887791in}{1.406232in}}%
\pgfpathlineto{\pgfqpoint{5.877255in}{1.407912in}}%
\pgfpathlineto{\pgfqpoint{5.871305in}{1.394832in}}%
\pgfpathlineto{\pgfqpoint{5.870480in}{1.385940in}}%
\pgfpathlineto{\pgfqpoint{5.862522in}{1.380232in}}%
\pgfpathlineto{\pgfqpoint{5.860753in}{1.362891in}}%
\pgfpathlineto{\pgfqpoint{5.849180in}{1.340984in}}%
\pgfpathlineto{\pgfqpoint{5.785872in}{1.331556in}}%
\pgfpathlineto{\pgfqpoint{5.785498in}{1.335699in}}%
\pgfpathlineto{\pgfqpoint{5.743150in}{1.331236in}}%
\pgfpathlineto{\pgfqpoint{5.685936in}{1.325898in}}%
\pgfpathclose%
\pgfusepath{fill}%
\end{pgfscope}%
\begin{pgfscope}%
\pgfpathrectangle{\pgfqpoint{3.625000in}{0.100000in}}{\pgfqpoint{2.989028in}{1.913466in}}%
\pgfusepath{clip}%
\pgfsetbuttcap%
\pgfsetmiterjoin%
\definecolor{currentfill}{rgb}{0.288043,0.635909,0.698731}%
\pgfsetfillcolor{currentfill}%
\pgfsetlinewidth{0.000000pt}%
\definecolor{currentstroke}{rgb}{0.000000,0.000000,0.000000}%
\pgfsetstrokecolor{currentstroke}%
\pgfsetstrokeopacity{0.000000}%
\pgfsetdash{}{0pt}%
\pgfpathmoveto{\pgfqpoint{4.087031in}{0.446426in}}%
\pgfpathlineto{\pgfqpoint{4.085767in}{0.445314in}}%
\pgfpathlineto{\pgfqpoint{4.079967in}{0.446797in}}%
\pgfpathlineto{\pgfqpoint{4.080400in}{0.449716in}}%
\pgfpathlineto{\pgfqpoint{4.083683in}{0.451053in}}%
\pgfpathlineto{\pgfqpoint{4.088364in}{0.457782in}}%
\pgfpathlineto{\pgfqpoint{4.089142in}{0.462330in}}%
\pgfpathlineto{\pgfqpoint{4.091773in}{0.463522in}}%
\pgfpathlineto{\pgfqpoint{4.096054in}{0.463526in}}%
\pgfpathlineto{\pgfqpoint{4.100530in}{0.478285in}}%
\pgfpathlineto{\pgfqpoint{4.102283in}{0.479499in}}%
\pgfpathlineto{\pgfqpoint{4.100711in}{0.482384in}}%
\pgfpathlineto{\pgfqpoint{4.097030in}{0.479466in}}%
\pgfpathlineto{\pgfqpoint{4.086215in}{0.483400in}}%
\pgfpathlineto{\pgfqpoint{4.079760in}{0.490029in}}%
\pgfpathlineto{\pgfqpoint{4.082575in}{0.493053in}}%
\pgfpathlineto{\pgfqpoint{4.081295in}{0.497319in}}%
\pgfpathlineto{\pgfqpoint{4.082083in}{0.502359in}}%
\pgfpathlineto{\pgfqpoint{4.080695in}{0.508250in}}%
\pgfpathlineto{\pgfqpoint{4.080473in}{0.514148in}}%
\pgfpathlineto{\pgfqpoint{4.086686in}{0.513161in}}%
\pgfpathlineto{\pgfqpoint{4.087916in}{0.514899in}}%
\pgfpathlineto{\pgfqpoint{4.091006in}{0.514731in}}%
\pgfpathlineto{\pgfqpoint{4.094565in}{0.508441in}}%
\pgfpathlineto{\pgfqpoint{4.097996in}{0.503835in}}%
\pgfpathlineto{\pgfqpoint{4.095248in}{0.500837in}}%
\pgfpathlineto{\pgfqpoint{4.100000in}{0.499687in}}%
\pgfpathlineto{\pgfqpoint{4.101083in}{0.501158in}}%
\pgfpathlineto{\pgfqpoint{4.098542in}{0.507123in}}%
\pgfpathlineto{\pgfqpoint{4.093190in}{0.512068in}}%
\pgfpathlineto{\pgfqpoint{4.093845in}{0.514140in}}%
\pgfpathlineto{\pgfqpoint{4.089574in}{0.519376in}}%
\pgfpathlineto{\pgfqpoint{4.093935in}{0.520439in}}%
\pgfpathlineto{\pgfqpoint{4.090157in}{0.531650in}}%
\pgfpathlineto{\pgfqpoint{4.093229in}{0.533204in}}%
\pgfpathlineto{\pgfqpoint{4.093978in}{0.535892in}}%
\pgfpathlineto{\pgfqpoint{4.091980in}{0.539085in}}%
\pgfpathlineto{\pgfqpoint{4.096104in}{0.539884in}}%
\pgfpathlineto{\pgfqpoint{4.096845in}{0.542299in}}%
\pgfpathlineto{\pgfqpoint{4.101765in}{0.538674in}}%
\pgfpathlineto{\pgfqpoint{4.102857in}{0.542605in}}%
\pgfpathlineto{\pgfqpoint{4.105443in}{0.544334in}}%
\pgfpathlineto{\pgfqpoint{4.118023in}{0.546739in}}%
\pgfpathlineto{\pgfqpoint{4.121885in}{0.545889in}}%
\pgfpathlineto{\pgfqpoint{4.125120in}{0.550663in}}%
\pgfpathlineto{\pgfqpoint{4.132797in}{0.553794in}}%
\pgfpathlineto{\pgfqpoint{4.138710in}{0.552606in}}%
\pgfpathlineto{\pgfqpoint{4.141407in}{0.548449in}}%
\pgfpathlineto{\pgfqpoint{4.141266in}{0.541844in}}%
\pgfpathlineto{\pgfqpoint{4.144292in}{0.540347in}}%
\pgfpathlineto{\pgfqpoint{4.150464in}{0.540556in}}%
\pgfpathlineto{\pgfqpoint{4.158370in}{0.542274in}}%
\pgfpathlineto{\pgfqpoint{4.158272in}{0.540162in}}%
\pgfpathlineto{\pgfqpoint{4.168266in}{0.534139in}}%
\pgfpathlineto{\pgfqpoint{4.175255in}{0.535754in}}%
\pgfpathlineto{\pgfqpoint{4.177313in}{0.537844in}}%
\pgfpathlineto{\pgfqpoint{4.179516in}{0.543197in}}%
\pgfpathlineto{\pgfqpoint{4.182248in}{0.546625in}}%
\pgfpathlineto{\pgfqpoint{4.182492in}{0.551762in}}%
\pgfpathlineto{\pgfqpoint{4.179990in}{0.553782in}}%
\pgfpathlineto{\pgfqpoint{4.183695in}{0.555608in}}%
\pgfpathlineto{\pgfqpoint{4.189920in}{0.552865in}}%
\pgfpathlineto{\pgfqpoint{4.191948in}{0.554559in}}%
\pgfpathlineto{\pgfqpoint{4.192305in}{0.561826in}}%
\pgfpathlineto{\pgfqpoint{4.188281in}{0.560307in}}%
\pgfpathlineto{\pgfqpoint{4.185490in}{0.562825in}}%
\pgfpathlineto{\pgfqpoint{4.180960in}{0.564395in}}%
\pgfpathlineto{\pgfqpoint{4.173770in}{0.564133in}}%
\pgfpathlineto{\pgfqpoint{4.166880in}{0.566521in}}%
\pgfpathlineto{\pgfqpoint{4.166535in}{0.572469in}}%
\pgfpathlineto{\pgfqpoint{4.160373in}{0.578787in}}%
\pgfpathlineto{\pgfqpoint{4.152559in}{0.581439in}}%
\pgfpathlineto{\pgfqpoint{4.145739in}{0.593671in}}%
\pgfpathlineto{\pgfqpoint{4.146298in}{0.598197in}}%
\pgfpathlineto{\pgfqpoint{4.149971in}{0.600261in}}%
\pgfpathlineto{\pgfqpoint{4.149430in}{0.604545in}}%
\pgfpathlineto{\pgfqpoint{4.150282in}{0.608808in}}%
\pgfpathlineto{\pgfqpoint{4.153188in}{0.605890in}}%
\pgfpathlineto{\pgfqpoint{4.157025in}{0.607107in}}%
\pgfpathlineto{\pgfqpoint{4.156833in}{0.610810in}}%
\pgfpathlineto{\pgfqpoint{4.151514in}{0.618243in}}%
\pgfpathlineto{\pgfqpoint{4.149956in}{0.625911in}}%
\pgfpathlineto{\pgfqpoint{4.154409in}{0.626496in}}%
\pgfpathlineto{\pgfqpoint{4.161074in}{0.626025in}}%
\pgfpathlineto{\pgfqpoint{4.163172in}{0.623460in}}%
\pgfpathlineto{\pgfqpoint{4.169147in}{0.624378in}}%
\pgfpathlineto{\pgfqpoint{4.174700in}{0.623173in}}%
\pgfpathlineto{\pgfqpoint{4.178454in}{0.621055in}}%
\pgfpathlineto{\pgfqpoint{4.181052in}{0.622023in}}%
\pgfpathlineto{\pgfqpoint{4.187951in}{0.619906in}}%
\pgfpathlineto{\pgfqpoint{4.188014in}{0.618032in}}%
\pgfpathlineto{\pgfqpoint{4.193268in}{0.618665in}}%
\pgfpathlineto{\pgfqpoint{4.200313in}{0.614092in}}%
\pgfpathlineto{\pgfqpoint{4.200039in}{0.611057in}}%
\pgfpathlineto{\pgfqpoint{4.196687in}{0.609656in}}%
\pgfpathlineto{\pgfqpoint{4.193124in}{0.606323in}}%
\pgfpathlineto{\pgfqpoint{4.192694in}{0.601746in}}%
\pgfpathlineto{\pgfqpoint{4.200369in}{0.595909in}}%
\pgfpathlineto{\pgfqpoint{4.199141in}{0.593112in}}%
\pgfpathlineto{\pgfqpoint{4.204602in}{0.590964in}}%
\pgfpathlineto{\pgfqpoint{4.205943in}{0.587455in}}%
\pgfpathlineto{\pgfqpoint{4.212873in}{0.590243in}}%
\pgfpathlineto{\pgfqpoint{4.215925in}{0.586589in}}%
\pgfpathlineto{\pgfqpoint{4.217334in}{0.588801in}}%
\pgfpathlineto{\pgfqpoint{4.213186in}{0.595999in}}%
\pgfpathlineto{\pgfqpoint{4.214660in}{0.598155in}}%
\pgfpathlineto{\pgfqpoint{4.215313in}{0.603841in}}%
\pgfpathlineto{\pgfqpoint{4.213661in}{0.606342in}}%
\pgfpathlineto{\pgfqpoint{4.214901in}{0.609704in}}%
\pgfpathlineto{\pgfqpoint{4.218961in}{0.609388in}}%
\pgfpathlineto{\pgfqpoint{4.218121in}{0.603668in}}%
\pgfpathlineto{\pgfqpoint{4.215529in}{0.601760in}}%
\pgfpathlineto{\pgfqpoint{4.215712in}{0.594065in}}%
\pgfpathlineto{\pgfqpoint{4.220364in}{0.593146in}}%
\pgfpathlineto{\pgfqpoint{4.220421in}{0.587424in}}%
\pgfpathlineto{\pgfqpoint{4.225501in}{0.583713in}}%
\pgfpathlineto{\pgfqpoint{4.227836in}{0.586053in}}%
\pgfpathlineto{\pgfqpoint{4.225284in}{0.593176in}}%
\pgfpathlineto{\pgfqpoint{4.222644in}{0.594979in}}%
\pgfpathlineto{\pgfqpoint{4.220047in}{0.593793in}}%
\pgfpathlineto{\pgfqpoint{4.217763in}{0.595306in}}%
\pgfpathlineto{\pgfqpoint{4.218161in}{0.601889in}}%
\pgfpathlineto{\pgfqpoint{4.221974in}{0.604581in}}%
\pgfpathlineto{\pgfqpoint{4.225726in}{0.604333in}}%
\pgfpathlineto{\pgfqpoint{4.224172in}{0.608058in}}%
\pgfpathlineto{\pgfqpoint{4.218547in}{0.610911in}}%
\pgfpathlineto{\pgfqpoint{4.211098in}{0.622277in}}%
\pgfpathlineto{\pgfqpoint{4.214816in}{0.627728in}}%
\pgfpathlineto{\pgfqpoint{4.216885in}{0.632863in}}%
\pgfpathlineto{\pgfqpoint{4.217048in}{0.643174in}}%
\pgfpathlineto{\pgfqpoint{4.216251in}{0.651872in}}%
\pgfpathlineto{\pgfqpoint{4.213978in}{0.657584in}}%
\pgfpathlineto{\pgfqpoint{4.214431in}{0.663294in}}%
\pgfpathlineto{\pgfqpoint{4.220244in}{0.668247in}}%
\pgfpathlineto{\pgfqpoint{4.226026in}{0.674133in}}%
\pgfpathlineto{\pgfqpoint{4.240519in}{0.662108in}}%
\pgfpathlineto{\pgfqpoint{4.249995in}{0.661221in}}%
\pgfpathlineto{\pgfqpoint{4.257817in}{0.664216in}}%
\pgfpathlineto{\pgfqpoint{4.264683in}{0.668159in}}%
\pgfpathlineto{\pgfqpoint{4.266141in}{0.670038in}}%
\pgfpathlineto{\pgfqpoint{4.283090in}{0.674322in}}%
\pgfpathlineto{\pgfqpoint{4.284562in}{0.671184in}}%
\pgfpathlineto{\pgfqpoint{4.290321in}{0.668615in}}%
\pgfpathlineto{\pgfqpoint{4.301313in}{0.669403in}}%
\pgfpathlineto{\pgfqpoint{4.303320in}{0.668702in}}%
\pgfpathlineto{\pgfqpoint{4.308260in}{0.670576in}}%
\pgfpathlineto{\pgfqpoint{4.310246in}{0.667335in}}%
\pgfpathlineto{\pgfqpoint{4.314990in}{0.664374in}}%
\pgfpathlineto{\pgfqpoint{4.317363in}{0.664709in}}%
\pgfpathlineto{\pgfqpoint{4.321956in}{0.661151in}}%
\pgfpathlineto{\pgfqpoint{4.329717in}{0.661805in}}%
\pgfpathlineto{\pgfqpoint{4.337469in}{0.664685in}}%
\pgfpathlineto{\pgfqpoint{4.343800in}{0.655475in}}%
\pgfpathlineto{\pgfqpoint{4.343309in}{0.653408in}}%
\pgfpathlineto{\pgfqpoint{4.336323in}{0.651032in}}%
\pgfpathlineto{\pgfqpoint{4.338190in}{0.647440in}}%
\pgfpathlineto{\pgfqpoint{4.344621in}{0.651221in}}%
\pgfpathlineto{\pgfqpoint{4.347246in}{0.650305in}}%
\pgfpathlineto{\pgfqpoint{4.348363in}{0.645783in}}%
\pgfpathlineto{\pgfqpoint{4.344524in}{0.643870in}}%
\pgfpathlineto{\pgfqpoint{4.347266in}{0.638154in}}%
\pgfpathlineto{\pgfqpoint{4.351248in}{0.639120in}}%
\pgfpathlineto{\pgfqpoint{4.357026in}{0.636170in}}%
\pgfpathlineto{\pgfqpoint{4.362551in}{0.627982in}}%
\pgfpathlineto{\pgfqpoint{4.357689in}{0.625708in}}%
\pgfpathlineto{\pgfqpoint{4.361985in}{0.619675in}}%
\pgfpathlineto{\pgfqpoint{4.358366in}{0.618264in}}%
\pgfpathlineto{\pgfqpoint{4.362868in}{0.612534in}}%
\pgfpathlineto{\pgfqpoint{4.368358in}{0.613145in}}%
\pgfpathlineto{\pgfqpoint{4.371249in}{0.610567in}}%
\pgfpathlineto{\pgfqpoint{4.378562in}{0.606765in}}%
\pgfpathlineto{\pgfqpoint{4.387995in}{0.593271in}}%
\pgfpathlineto{\pgfqpoint{4.387782in}{0.590043in}}%
\pgfpathlineto{\pgfqpoint{4.401928in}{0.579706in}}%
\pgfpathlineto{\pgfqpoint{4.402233in}{0.576049in}}%
\pgfpathlineto{\pgfqpoint{4.406003in}{0.570578in}}%
\pgfpathlineto{\pgfqpoint{4.417585in}{0.567843in}}%
\pgfpathlineto{\pgfqpoint{4.421881in}{0.565899in}}%
\pgfpathlineto{\pgfqpoint{4.425922in}{0.559868in}}%
\pgfpathlineto{\pgfqpoint{4.426442in}{0.554585in}}%
\pgfpathlineto{\pgfqpoint{4.429915in}{0.550279in}}%
\pgfpathlineto{\pgfqpoint{4.430958in}{0.545423in}}%
\pgfpathlineto{\pgfqpoint{4.434213in}{0.543649in}}%
\pgfpathlineto{\pgfqpoint{4.418089in}{0.514548in}}%
\pgfpathlineto{\pgfqpoint{4.384622in}{0.454133in}}%
\pgfpathlineto{\pgfqpoint{4.335518in}{0.365475in}}%
\pgfpathlineto{\pgfqpoint{4.316294in}{0.330756in}}%
\pgfpathlineto{\pgfqpoint{4.320388in}{0.326093in}}%
\pgfpathlineto{\pgfqpoint{4.322192in}{0.327581in}}%
\pgfpathlineto{\pgfqpoint{4.325850in}{0.322152in}}%
\pgfpathlineto{\pgfqpoint{4.330927in}{0.323727in}}%
\pgfpathlineto{\pgfqpoint{4.337667in}{0.320475in}}%
\pgfpathlineto{\pgfqpoint{4.333301in}{0.315393in}}%
\pgfpathlineto{\pgfqpoint{4.336672in}{0.308535in}}%
\pgfpathlineto{\pgfqpoint{4.335976in}{0.305196in}}%
\pgfpathlineto{\pgfqpoint{4.341353in}{0.287587in}}%
\pgfpathlineto{\pgfqpoint{4.339077in}{0.279580in}}%
\pgfpathlineto{\pgfqpoint{4.349350in}{0.281562in}}%
\pgfpathlineto{\pgfqpoint{4.351872in}{0.280263in}}%
\pgfpathlineto{\pgfqpoint{4.354610in}{0.282376in}}%
\pgfpathlineto{\pgfqpoint{4.356558in}{0.286368in}}%
\pgfpathlineto{\pgfqpoint{4.359347in}{0.288657in}}%
\pgfpathlineto{\pgfqpoint{4.371252in}{0.288367in}}%
\pgfpathlineto{\pgfqpoint{4.373943in}{0.280683in}}%
\pgfpathlineto{\pgfqpoint{4.371634in}{0.273938in}}%
\pgfpathlineto{\pgfqpoint{4.374257in}{0.271712in}}%
\pgfpathlineto{\pgfqpoint{4.375636in}{0.267905in}}%
\pgfpathlineto{\pgfqpoint{4.375336in}{0.260746in}}%
\pgfpathlineto{\pgfqpoint{4.378836in}{0.255613in}}%
\pgfpathlineto{\pgfqpoint{4.380755in}{0.247568in}}%
\pgfpathlineto{\pgfqpoint{4.379770in}{0.243223in}}%
\pgfpathlineto{\pgfqpoint{4.380958in}{0.238819in}}%
\pgfpathlineto{\pgfqpoint{4.382671in}{0.213318in}}%
\pgfpathlineto{\pgfqpoint{4.380216in}{0.211309in}}%
\pgfpathlineto{\pgfqpoint{4.383455in}{0.208533in}}%
\pgfpathlineto{\pgfqpoint{4.380910in}{0.205157in}}%
\pgfpathlineto{\pgfqpoint{4.383206in}{0.202405in}}%
\pgfpathlineto{\pgfqpoint{4.381727in}{0.197707in}}%
\pgfpathlineto{\pgfqpoint{4.385034in}{0.196433in}}%
\pgfpathlineto{\pgfqpoint{4.389287in}{0.189287in}}%
\pgfpathlineto{\pgfqpoint{4.392447in}{0.187134in}}%
\pgfpathlineto{\pgfqpoint{4.393274in}{0.183987in}}%
\pgfpathlineto{\pgfqpoint{4.395105in}{0.182623in}}%
\pgfpathlineto{\pgfqpoint{4.394723in}{0.180010in}}%
\pgfpathlineto{\pgfqpoint{4.398670in}{0.178116in}}%
\pgfpathlineto{\pgfqpoint{4.397696in}{0.173088in}}%
\pgfpathlineto{\pgfqpoint{4.394298in}{0.170454in}}%
\pgfpathlineto{\pgfqpoint{4.392755in}{0.165560in}}%
\pgfpathlineto{\pgfqpoint{4.392345in}{0.159006in}}%
\pgfpathlineto{\pgfqpoint{4.390458in}{0.158379in}}%
\pgfpathlineto{\pgfqpoint{4.384289in}{0.152634in}}%
\pgfpathlineto{\pgfqpoint{4.378806in}{0.151118in}}%
\pgfpathlineto{\pgfqpoint{4.376190in}{0.153435in}}%
\pgfpathlineto{\pgfqpoint{4.378718in}{0.159583in}}%
\pgfpathlineto{\pgfqpoint{4.377946in}{0.162601in}}%
\pgfpathlineto{\pgfqpoint{4.381568in}{0.164131in}}%
\pgfpathlineto{\pgfqpoint{4.384956in}{0.173027in}}%
\pgfpathlineto{\pgfqpoint{4.383525in}{0.180625in}}%
\pgfpathlineto{\pgfqpoint{4.381235in}{0.182419in}}%
\pgfpathlineto{\pgfqpoint{4.373578in}{0.181838in}}%
\pgfpathlineto{\pgfqpoint{4.375121in}{0.179853in}}%
\pgfpathlineto{\pgfqpoint{4.369579in}{0.174155in}}%
\pgfpathlineto{\pgfqpoint{4.367793in}{0.177323in}}%
\pgfpathlineto{\pgfqpoint{4.368258in}{0.181462in}}%
\pgfpathlineto{\pgfqpoint{4.371439in}{0.181386in}}%
\pgfpathlineto{\pgfqpoint{4.374272in}{0.184426in}}%
\pgfpathlineto{\pgfqpoint{4.375948in}{0.188866in}}%
\pgfpathlineto{\pgfqpoint{4.382264in}{0.187613in}}%
\pgfpathlineto{\pgfqpoint{4.377208in}{0.189984in}}%
\pgfpathlineto{\pgfqpoint{4.377653in}{0.193161in}}%
\pgfpathlineto{\pgfqpoint{4.375549in}{0.194596in}}%
\pgfpathlineto{\pgfqpoint{4.374687in}{0.198923in}}%
\pgfpathlineto{\pgfqpoint{4.376181in}{0.201177in}}%
\pgfpathlineto{\pgfqpoint{4.372619in}{0.208179in}}%
\pgfpathlineto{\pgfqpoint{4.372919in}{0.212633in}}%
\pgfpathlineto{\pgfqpoint{4.368746in}{0.216599in}}%
\pgfpathlineto{\pgfqpoint{4.366488in}{0.219808in}}%
\pgfpathlineto{\pgfqpoint{4.369606in}{0.222866in}}%
\pgfpathlineto{\pgfqpoint{4.370761in}{0.227437in}}%
\pgfpathlineto{\pgfqpoint{4.369730in}{0.230438in}}%
\pgfpathlineto{\pgfqpoint{4.370936in}{0.234017in}}%
\pgfpathlineto{\pgfqpoint{4.369399in}{0.245686in}}%
\pgfpathlineto{\pgfqpoint{4.367100in}{0.250980in}}%
\pgfpathlineto{\pgfqpoint{4.364471in}{0.252970in}}%
\pgfpathlineto{\pgfqpoint{4.365590in}{0.259857in}}%
\pgfpathlineto{\pgfqpoint{4.364801in}{0.265235in}}%
\pgfpathlineto{\pgfqpoint{4.366230in}{0.268685in}}%
\pgfpathlineto{\pgfqpoint{4.363272in}{0.269169in}}%
\pgfpathlineto{\pgfqpoint{4.362479in}{0.260789in}}%
\pgfpathlineto{\pgfqpoint{4.359154in}{0.251441in}}%
\pgfpathlineto{\pgfqpoint{4.356484in}{0.253570in}}%
\pgfpathlineto{\pgfqpoint{4.356127in}{0.256642in}}%
\pgfpathlineto{\pgfqpoint{4.352861in}{0.260839in}}%
\pgfpathlineto{\pgfqpoint{4.354521in}{0.263587in}}%
\pgfpathlineto{\pgfqpoint{4.354005in}{0.269772in}}%
\pgfpathlineto{\pgfqpoint{4.349778in}{0.267597in}}%
\pgfpathlineto{\pgfqpoint{4.350607in}{0.264242in}}%
\pgfpathlineto{\pgfqpoint{4.349811in}{0.259983in}}%
\pgfpathlineto{\pgfqpoint{4.345066in}{0.259990in}}%
\pgfpathlineto{\pgfqpoint{4.342939in}{0.262630in}}%
\pgfpathlineto{\pgfqpoint{4.339729in}{0.263134in}}%
\pgfpathlineto{\pgfqpoint{4.337569in}{0.265977in}}%
\pgfpathlineto{\pgfqpoint{4.333852in}{0.273782in}}%
\pgfpathlineto{\pgfqpoint{4.332648in}{0.279842in}}%
\pgfpathlineto{\pgfqpoint{4.333418in}{0.282418in}}%
\pgfpathlineto{\pgfqpoint{4.331620in}{0.288091in}}%
\pgfpathlineto{\pgfqpoint{4.328639in}{0.291375in}}%
\pgfpathlineto{\pgfqpoint{4.323483in}{0.299768in}}%
\pgfpathlineto{\pgfqpoint{4.319744in}{0.307057in}}%
\pgfpathlineto{\pgfqpoint{4.323559in}{0.307167in}}%
\pgfpathlineto{\pgfqpoint{4.325804in}{0.308698in}}%
\pgfpathlineto{\pgfqpoint{4.326246in}{0.313435in}}%
\pgfpathlineto{\pgfqpoint{4.315450in}{0.314007in}}%
\pgfpathlineto{\pgfqpoint{4.306267in}{0.324264in}}%
\pgfpathlineto{\pgfqpoint{4.308682in}{0.326962in}}%
\pgfpathlineto{\pgfqpoint{4.304241in}{0.327423in}}%
\pgfpathlineto{\pgfqpoint{4.295423in}{0.336959in}}%
\pgfpathlineto{\pgfqpoint{4.281451in}{0.341741in}}%
\pgfpathlineto{\pgfqpoint{4.279570in}{0.347345in}}%
\pgfpathlineto{\pgfqpoint{4.276671in}{0.350264in}}%
\pgfpathlineto{\pgfqpoint{4.276539in}{0.353005in}}%
\pgfpathlineto{\pgfqpoint{4.274358in}{0.355041in}}%
\pgfpathlineto{\pgfqpoint{4.277220in}{0.358442in}}%
\pgfpathlineto{\pgfqpoint{4.270693in}{0.358770in}}%
\pgfpathlineto{\pgfqpoint{4.269034in}{0.363173in}}%
\pgfpathlineto{\pgfqpoint{4.266063in}{0.365047in}}%
\pgfpathlineto{\pgfqpoint{4.272012in}{0.367309in}}%
\pgfpathlineto{\pgfqpoint{4.268918in}{0.371345in}}%
\pgfpathlineto{\pgfqpoint{4.263224in}{0.373038in}}%
\pgfpathlineto{\pgfqpoint{4.264110in}{0.380491in}}%
\pgfpathlineto{\pgfqpoint{4.262478in}{0.383130in}}%
\pgfpathlineto{\pgfqpoint{4.259538in}{0.384988in}}%
\pgfpathlineto{\pgfqpoint{4.257024in}{0.383124in}}%
\pgfpathlineto{\pgfqpoint{4.255538in}{0.384877in}}%
\pgfpathlineto{\pgfqpoint{4.251206in}{0.385379in}}%
\pgfpathlineto{\pgfqpoint{4.251555in}{0.392143in}}%
\pgfpathlineto{\pgfqpoint{4.245275in}{0.387574in}}%
\pgfpathlineto{\pgfqpoint{4.245407in}{0.378160in}}%
\pgfpathlineto{\pgfqpoint{4.238666in}{0.376355in}}%
\pgfpathlineto{\pgfqpoint{4.234263in}{0.372596in}}%
\pgfpathlineto{\pgfqpoint{4.230725in}{0.371804in}}%
\pgfpathlineto{\pgfqpoint{4.226895in}{0.375593in}}%
\pgfpathlineto{\pgfqpoint{4.223485in}{0.375090in}}%
\pgfpathlineto{\pgfqpoint{4.224118in}{0.378551in}}%
\pgfpathlineto{\pgfqpoint{4.219270in}{0.376688in}}%
\pgfpathlineto{\pgfqpoint{4.214779in}{0.376478in}}%
\pgfpathlineto{\pgfqpoint{4.209523in}{0.374589in}}%
\pgfpathlineto{\pgfqpoint{4.198855in}{0.375652in}}%
\pgfpathlineto{\pgfqpoint{4.195430in}{0.374953in}}%
\pgfpathlineto{\pgfqpoint{4.192959in}{0.378065in}}%
\pgfpathlineto{\pgfqpoint{4.188397in}{0.378376in}}%
\pgfpathlineto{\pgfqpoint{4.187449in}{0.382409in}}%
\pgfpathlineto{\pgfqpoint{4.190167in}{0.384606in}}%
\pgfpathlineto{\pgfqpoint{4.193357in}{0.384399in}}%
\pgfpathlineto{\pgfqpoint{4.196019in}{0.381403in}}%
\pgfpathlineto{\pgfqpoint{4.197772in}{0.386005in}}%
\pgfpathlineto{\pgfqpoint{4.195326in}{0.391102in}}%
\pgfpathlineto{\pgfqpoint{4.200863in}{0.395589in}}%
\pgfpathlineto{\pgfqpoint{4.206411in}{0.397243in}}%
\pgfpathlineto{\pgfqpoint{4.210273in}{0.399940in}}%
\pgfpathlineto{\pgfqpoint{4.212775in}{0.402846in}}%
\pgfpathlineto{\pgfqpoint{4.214172in}{0.407609in}}%
\pgfpathlineto{\pgfqpoint{4.218583in}{0.406312in}}%
\pgfpathlineto{\pgfqpoint{4.228148in}{0.406979in}}%
\pgfpathlineto{\pgfqpoint{4.230201in}{0.401756in}}%
\pgfpathlineto{\pgfqpoint{4.233556in}{0.401529in}}%
\pgfpathlineto{\pgfqpoint{4.239662in}{0.393534in}}%
\pgfpathlineto{\pgfqpoint{4.239673in}{0.396448in}}%
\pgfpathlineto{\pgfqpoint{4.237540in}{0.397396in}}%
\pgfpathlineto{\pgfqpoint{4.233636in}{0.406940in}}%
\pgfpathlineto{\pgfqpoint{4.235536in}{0.408236in}}%
\pgfpathlineto{\pgfqpoint{4.230322in}{0.413215in}}%
\pgfpathlineto{\pgfqpoint{4.225663in}{0.414184in}}%
\pgfpathlineto{\pgfqpoint{4.221223in}{0.412463in}}%
\pgfpathlineto{\pgfqpoint{4.213619in}{0.413804in}}%
\pgfpathlineto{\pgfqpoint{4.210013in}{0.411392in}}%
\pgfpathlineto{\pgfqpoint{4.201801in}{0.409529in}}%
\pgfpathlineto{\pgfqpoint{4.201050in}{0.406847in}}%
\pgfpathlineto{\pgfqpoint{4.194741in}{0.406115in}}%
\pgfpathlineto{\pgfqpoint{4.193122in}{0.403049in}}%
\pgfpathlineto{\pgfqpoint{4.189513in}{0.400803in}}%
\pgfpathlineto{\pgfqpoint{4.185154in}{0.401189in}}%
\pgfpathlineto{\pgfqpoint{4.182977in}{0.399018in}}%
\pgfpathlineto{\pgfqpoint{4.180231in}{0.399119in}}%
\pgfpathlineto{\pgfqpoint{4.177607in}{0.401074in}}%
\pgfpathlineto{\pgfqpoint{4.173019in}{0.400725in}}%
\pgfpathlineto{\pgfqpoint{4.171938in}{0.398924in}}%
\pgfpathlineto{\pgfqpoint{4.167119in}{0.400773in}}%
\pgfpathlineto{\pgfqpoint{4.163027in}{0.396000in}}%
\pgfpathlineto{\pgfqpoint{4.166950in}{0.391382in}}%
\pgfpathlineto{\pgfqpoint{4.168460in}{0.387584in}}%
\pgfpathlineto{\pgfqpoint{4.167068in}{0.384367in}}%
\pgfpathlineto{\pgfqpoint{4.161493in}{0.382148in}}%
\pgfpathlineto{\pgfqpoint{4.158263in}{0.383954in}}%
\pgfpathlineto{\pgfqpoint{4.154198in}{0.382976in}}%
\pgfpathlineto{\pgfqpoint{4.153708in}{0.380157in}}%
\pgfpathlineto{\pgfqpoint{4.148759in}{0.381154in}}%
\pgfpathlineto{\pgfqpoint{4.149855in}{0.378337in}}%
\pgfpathlineto{\pgfqpoint{4.142555in}{0.377743in}}%
\pgfpathlineto{\pgfqpoint{4.137830in}{0.381319in}}%
\pgfpathlineto{\pgfqpoint{4.135306in}{0.379045in}}%
\pgfpathlineto{\pgfqpoint{4.128590in}{0.381466in}}%
\pgfpathlineto{\pgfqpoint{4.126918in}{0.379014in}}%
\pgfpathlineto{\pgfqpoint{4.123588in}{0.377923in}}%
\pgfpathlineto{\pgfqpoint{4.120917in}{0.380797in}}%
\pgfpathlineto{\pgfqpoint{4.111784in}{0.380083in}}%
\pgfpathlineto{\pgfqpoint{4.111794in}{0.375843in}}%
\pgfpathlineto{\pgfqpoint{4.108125in}{0.374721in}}%
\pgfpathlineto{\pgfqpoint{4.105270in}{0.376900in}}%
\pgfpathlineto{\pgfqpoint{4.101497in}{0.375689in}}%
\pgfpathlineto{\pgfqpoint{4.096614in}{0.375401in}}%
\pgfpathlineto{\pgfqpoint{4.089340in}{0.378466in}}%
\pgfpathlineto{\pgfqpoint{4.085174in}{0.375305in}}%
\pgfpathlineto{\pgfqpoint{4.079216in}{0.379530in}}%
\pgfpathlineto{\pgfqpoint{4.077268in}{0.377929in}}%
\pgfpathlineto{\pgfqpoint{4.075900in}{0.373791in}}%
\pgfpathlineto{\pgfqpoint{4.072071in}{0.371365in}}%
\pgfpathlineto{\pgfqpoint{4.065825in}{0.374142in}}%
\pgfpathlineto{\pgfqpoint{4.055479in}{0.381003in}}%
\pgfpathlineto{\pgfqpoint{4.052463in}{0.381613in}}%
\pgfpathlineto{\pgfqpoint{4.050762in}{0.380183in}}%
\pgfpathlineto{\pgfqpoint{4.046211in}{0.381155in}}%
\pgfpathlineto{\pgfqpoint{4.043991in}{0.383290in}}%
\pgfpathlineto{\pgfqpoint{4.039445in}{0.384441in}}%
\pgfpathlineto{\pgfqpoint{4.033238in}{0.384507in}}%
\pgfpathlineto{\pgfqpoint{4.030773in}{0.387054in}}%
\pgfpathlineto{\pgfqpoint{4.035497in}{0.389690in}}%
\pgfpathlineto{\pgfqpoint{4.034395in}{0.392291in}}%
\pgfpathlineto{\pgfqpoint{4.031149in}{0.391707in}}%
\pgfpathlineto{\pgfqpoint{4.029528in}{0.389662in}}%
\pgfpathlineto{\pgfqpoint{4.021930in}{0.386730in}}%
\pgfpathlineto{\pgfqpoint{4.018088in}{0.387727in}}%
\pgfpathlineto{\pgfqpoint{4.016471in}{0.393681in}}%
\pgfpathlineto{\pgfqpoint{4.011679in}{0.391121in}}%
\pgfpathlineto{\pgfqpoint{4.010286in}{0.398667in}}%
\pgfpathlineto{\pgfqpoint{4.008015in}{0.396722in}}%
\pgfpathlineto{\pgfqpoint{4.008395in}{0.393819in}}%
\pgfpathlineto{\pgfqpoint{4.003173in}{0.394586in}}%
\pgfpathlineto{\pgfqpoint{4.008605in}{0.399646in}}%
\pgfpathlineto{\pgfqpoint{4.014110in}{0.396639in}}%
\pgfpathlineto{\pgfqpoint{4.015291in}{0.398148in}}%
\pgfpathlineto{\pgfqpoint{4.021082in}{0.396339in}}%
\pgfpathlineto{\pgfqpoint{4.020720in}{0.398399in}}%
\pgfpathlineto{\pgfqpoint{4.036521in}{0.399081in}}%
\pgfpathlineto{\pgfqpoint{4.044230in}{0.395868in}}%
\pgfpathlineto{\pgfqpoint{4.047212in}{0.392698in}}%
\pgfpathlineto{\pgfqpoint{4.046327in}{0.389931in}}%
\pgfpathlineto{\pgfqpoint{4.049323in}{0.388883in}}%
\pgfpathlineto{\pgfqpoint{4.056765in}{0.393251in}}%
\pgfpathlineto{\pgfqpoint{4.066392in}{0.393531in}}%
\pgfpathlineto{\pgfqpoint{4.075136in}{0.390898in}}%
\pgfpathlineto{\pgfqpoint{4.080089in}{0.391340in}}%
\pgfpathlineto{\pgfqpoint{4.081933in}{0.387255in}}%
\pgfpathlineto{\pgfqpoint{4.084703in}{0.391758in}}%
\pgfpathlineto{\pgfqpoint{4.086407in}{0.392733in}}%
\pgfpathlineto{\pgfqpoint{4.094288in}{0.394691in}}%
\pgfpathlineto{\pgfqpoint{4.097217in}{0.393494in}}%
\pgfpathlineto{\pgfqpoint{4.100385in}{0.394915in}}%
\pgfpathlineto{\pgfqpoint{4.104163in}{0.394172in}}%
\pgfpathlineto{\pgfqpoint{4.105055in}{0.395827in}}%
\pgfpathlineto{\pgfqpoint{4.112939in}{0.403420in}}%
\pgfpathlineto{\pgfqpoint{4.116278in}{0.404545in}}%
\pgfpathlineto{\pgfqpoint{4.118251in}{0.408644in}}%
\pgfpathlineto{\pgfqpoint{4.128722in}{0.411817in}}%
\pgfpathlineto{\pgfqpoint{4.130033in}{0.414238in}}%
\pgfpathlineto{\pgfqpoint{4.115192in}{0.417974in}}%
\pgfpathlineto{\pgfqpoint{4.115776in}{0.421993in}}%
\pgfpathlineto{\pgfqpoint{4.114250in}{0.427301in}}%
\pgfpathlineto{\pgfqpoint{4.110794in}{0.426860in}}%
\pgfpathlineto{\pgfqpoint{4.108443in}{0.419972in}}%
\pgfpathlineto{\pgfqpoint{4.105714in}{0.419404in}}%
\pgfpathlineto{\pgfqpoint{4.103938in}{0.421634in}}%
\pgfpathlineto{\pgfqpoint{4.106010in}{0.431388in}}%
\pgfpathlineto{\pgfqpoint{4.103947in}{0.435475in}}%
\pgfpathlineto{\pgfqpoint{4.099950in}{0.440293in}}%
\pgfpathlineto{\pgfqpoint{4.101871in}{0.443969in}}%
\pgfpathlineto{\pgfqpoint{4.093436in}{0.444102in}}%
\pgfpathlineto{\pgfqpoint{4.093628in}{0.445528in}}%
\pgfpathlineto{\pgfqpoint{4.088595in}{0.447526in}}%
\pgfpathlineto{\pgfqpoint{4.087031in}{0.446426in}}%
\pgfpathclose%
\pgfusepath{fill}%
\end{pgfscope}%
\begin{pgfscope}%
\pgfpathrectangle{\pgfqpoint{3.625000in}{0.100000in}}{\pgfqpoint{2.989028in}{1.913466in}}%
\pgfusepath{clip}%
\pgfsetbuttcap%
\pgfsetmiterjoin%
\definecolor{currentfill}{rgb}{0.288043,0.635909,0.698731}%
\pgfsetfillcolor{currentfill}%
\pgfsetlinewidth{0.000000pt}%
\definecolor{currentstroke}{rgb}{0.000000,0.000000,0.000000}%
\pgfsetstrokecolor{currentstroke}%
\pgfsetstrokeopacity{0.000000}%
\pgfsetdash{}{0pt}%
\pgfpathmoveto{\pgfqpoint{4.071979in}{0.517466in}}%
\pgfpathlineto{\pgfqpoint{4.070936in}{0.515633in}}%
\pgfpathlineto{\pgfqpoint{4.073696in}{0.511730in}}%
\pgfpathlineto{\pgfqpoint{4.069334in}{0.508079in}}%
\pgfpathlineto{\pgfqpoint{4.069455in}{0.504939in}}%
\pgfpathlineto{\pgfqpoint{4.067483in}{0.504198in}}%
\pgfpathlineto{\pgfqpoint{4.060704in}{0.508881in}}%
\pgfpathlineto{\pgfqpoint{4.055872in}{0.519132in}}%
\pgfpathlineto{\pgfqpoint{4.055549in}{0.522094in}}%
\pgfpathlineto{\pgfqpoint{4.056791in}{0.525640in}}%
\pgfpathlineto{\pgfqpoint{4.062027in}{0.520216in}}%
\pgfpathlineto{\pgfqpoint{4.063628in}{0.521262in}}%
\pgfpathlineto{\pgfqpoint{4.068070in}{0.520267in}}%
\pgfpathlineto{\pgfqpoint{4.071979in}{0.517466in}}%
\pgfpathclose%
\pgfusepath{fill}%
\end{pgfscope}%
\begin{pgfscope}%
\pgfpathrectangle{\pgfqpoint{3.625000in}{0.100000in}}{\pgfqpoint{2.989028in}{1.913466in}}%
\pgfusepath{clip}%
\pgfsetbuttcap%
\pgfsetmiterjoin%
\definecolor{currentfill}{rgb}{0.288043,0.635909,0.698731}%
\pgfsetfillcolor{currentfill}%
\pgfsetlinewidth{0.000000pt}%
\definecolor{currentstroke}{rgb}{0.000000,0.000000,0.000000}%
\pgfsetstrokecolor{currentstroke}%
\pgfsetstrokeopacity{0.000000}%
\pgfsetdash{}{0pt}%
\pgfpathmoveto{\pgfqpoint{3.991584in}{0.398705in}}%
\pgfpathlineto{\pgfqpoint{3.986820in}{0.397932in}}%
\pgfpathlineto{\pgfqpoint{3.983372in}{0.399525in}}%
\pgfpathlineto{\pgfqpoint{3.981865in}{0.402120in}}%
\pgfpathlineto{\pgfqpoint{3.983573in}{0.405842in}}%
\pgfpathlineto{\pgfqpoint{3.987377in}{0.404858in}}%
\pgfpathlineto{\pgfqpoint{3.993630in}{0.407066in}}%
\pgfpathlineto{\pgfqpoint{3.995820in}{0.404016in}}%
\pgfpathlineto{\pgfqpoint{3.997944in}{0.404579in}}%
\pgfpathlineto{\pgfqpoint{4.003084in}{0.402681in}}%
\pgfpathlineto{\pgfqpoint{4.005305in}{0.400269in}}%
\pgfpathlineto{\pgfqpoint{4.002595in}{0.393963in}}%
\pgfpathlineto{\pgfqpoint{3.999899in}{0.392626in}}%
\pgfpathlineto{\pgfqpoint{3.997476in}{0.393346in}}%
\pgfpathlineto{\pgfqpoint{3.991584in}{0.398705in}}%
\pgfpathclose%
\pgfusepath{fill}%
\end{pgfscope}%
\begin{pgfscope}%
\pgfpathrectangle{\pgfqpoint{3.625000in}{0.100000in}}{\pgfqpoint{2.989028in}{1.913466in}}%
\pgfusepath{clip}%
\pgfsetbuttcap%
\pgfsetmiterjoin%
\definecolor{currentfill}{rgb}{0.288043,0.635909,0.698731}%
\pgfsetfillcolor{currentfill}%
\pgfsetlinewidth{0.000000pt}%
\definecolor{currentstroke}{rgb}{0.000000,0.000000,0.000000}%
\pgfsetstrokecolor{currentstroke}%
\pgfsetstrokeopacity{0.000000}%
\pgfsetdash{}{0pt}%
\pgfpathmoveto{\pgfqpoint{3.948985in}{0.404436in}}%
\pgfpathlineto{\pgfqpoint{3.942234in}{0.407656in}}%
\pgfpathlineto{\pgfqpoint{3.935987in}{0.408614in}}%
\pgfpathlineto{\pgfqpoint{3.934018in}{0.409985in}}%
\pgfpathlineto{\pgfqpoint{3.936260in}{0.412685in}}%
\pgfpathlineto{\pgfqpoint{3.940239in}{0.409312in}}%
\pgfpathlineto{\pgfqpoint{3.943576in}{0.409777in}}%
\pgfpathlineto{\pgfqpoint{3.949004in}{0.412113in}}%
\pgfpathlineto{\pgfqpoint{3.949442in}{0.415120in}}%
\pgfpathlineto{\pgfqpoint{3.956109in}{0.413632in}}%
\pgfpathlineto{\pgfqpoint{3.954548in}{0.409103in}}%
\pgfpathlineto{\pgfqpoint{3.958794in}{0.408772in}}%
\pgfpathlineto{\pgfqpoint{3.955231in}{0.403385in}}%
\pgfpathlineto{\pgfqpoint{3.951781in}{0.405131in}}%
\pgfpathlineto{\pgfqpoint{3.948985in}{0.404436in}}%
\pgfpathclose%
\pgfusepath{fill}%
\end{pgfscope}%
\begin{pgfscope}%
\pgfpathrectangle{\pgfqpoint{3.625000in}{0.100000in}}{\pgfqpoint{2.989028in}{1.913466in}}%
\pgfusepath{clip}%
\pgfsetbuttcap%
\pgfsetmiterjoin%
\definecolor{currentfill}{rgb}{0.288043,0.635909,0.698731}%
\pgfsetfillcolor{currentfill}%
\pgfsetlinewidth{0.000000pt}%
\definecolor{currentstroke}{rgb}{0.000000,0.000000,0.000000}%
\pgfsetstrokecolor{currentstroke}%
\pgfsetstrokeopacity{0.000000}%
\pgfsetdash{}{0pt}%
\pgfpathmoveto{\pgfqpoint{3.937124in}{0.415568in}}%
\pgfpathlineto{\pgfqpoint{3.934643in}{0.414219in}}%
\pgfpathlineto{\pgfqpoint{3.927882in}{0.416293in}}%
\pgfpathlineto{\pgfqpoint{3.922830in}{0.415092in}}%
\pgfpathlineto{\pgfqpoint{3.920264in}{0.418504in}}%
\pgfpathlineto{\pgfqpoint{3.921383in}{0.420066in}}%
\pgfpathlineto{\pgfqpoint{3.925212in}{0.420327in}}%
\pgfpathlineto{\pgfqpoint{3.928140in}{0.419346in}}%
\pgfpathlineto{\pgfqpoint{3.931331in}{0.420898in}}%
\pgfpathlineto{\pgfqpoint{3.936223in}{0.418823in}}%
\pgfpathlineto{\pgfqpoint{3.937124in}{0.415568in}}%
\pgfpathclose%
\pgfusepath{fill}%
\end{pgfscope}%
\begin{pgfscope}%
\pgfpathrectangle{\pgfqpoint{3.625000in}{0.100000in}}{\pgfqpoint{2.989028in}{1.913466in}}%
\pgfusepath{clip}%
\pgfsetbuttcap%
\pgfsetmiterjoin%
\definecolor{currentfill}{rgb}{0.288043,0.635909,0.698731}%
\pgfsetfillcolor{currentfill}%
\pgfsetlinewidth{0.000000pt}%
\definecolor{currentstroke}{rgb}{0.000000,0.000000,0.000000}%
\pgfsetstrokecolor{currentstroke}%
\pgfsetstrokeopacity{0.000000}%
\pgfsetdash{}{0pt}%
\pgfpathmoveto{\pgfqpoint{3.857429in}{0.465782in}}%
\pgfpathlineto{\pgfqpoint{3.857250in}{0.462436in}}%
\pgfpathlineto{\pgfqpoint{3.852314in}{0.460334in}}%
\pgfpathlineto{\pgfqpoint{3.847807in}{0.465597in}}%
\pgfpathlineto{\pgfqpoint{3.852862in}{0.467350in}}%
\pgfpathlineto{\pgfqpoint{3.857429in}{0.465782in}}%
\pgfpathclose%
\pgfusepath{fill}%
\end{pgfscope}%
\begin{pgfscope}%
\pgfpathrectangle{\pgfqpoint{3.625000in}{0.100000in}}{\pgfqpoint{2.989028in}{1.913466in}}%
\pgfusepath{clip}%
\pgfsetbuttcap%
\pgfsetmiterjoin%
\definecolor{currentfill}{rgb}{0.288043,0.635909,0.698731}%
\pgfsetfillcolor{currentfill}%
\pgfsetlinewidth{0.000000pt}%
\definecolor{currentstroke}{rgb}{0.000000,0.000000,0.000000}%
\pgfsetstrokecolor{currentstroke}%
\pgfsetstrokeopacity{0.000000}%
\pgfsetdash{}{0pt}%
\pgfpathmoveto{\pgfqpoint{3.816722in}{0.484589in}}%
\pgfpathlineto{\pgfqpoint{3.817771in}{0.486043in}}%
\pgfpathlineto{\pgfqpoint{3.823825in}{0.485792in}}%
\pgfpathlineto{\pgfqpoint{3.824887in}{0.487956in}}%
\pgfpathlineto{\pgfqpoint{3.827543in}{0.486068in}}%
\pgfpathlineto{\pgfqpoint{3.826009in}{0.482435in}}%
\pgfpathlineto{\pgfqpoint{3.823753in}{0.482127in}}%
\pgfpathlineto{\pgfqpoint{3.818361in}{0.483267in}}%
\pgfpathlineto{\pgfqpoint{3.816722in}{0.484589in}}%
\pgfpathclose%
\pgfusepath{fill}%
\end{pgfscope}%
\begin{pgfscope}%
\pgfpathrectangle{\pgfqpoint{3.625000in}{0.100000in}}{\pgfqpoint{2.989028in}{1.913466in}}%
\pgfusepath{clip}%
\pgfsetbuttcap%
\pgfsetmiterjoin%
\definecolor{currentfill}{rgb}{0.288043,0.635909,0.698731}%
\pgfsetfillcolor{currentfill}%
\pgfsetlinewidth{0.000000pt}%
\definecolor{currentstroke}{rgb}{0.000000,0.000000,0.000000}%
\pgfsetstrokecolor{currentstroke}%
\pgfsetstrokeopacity{0.000000}%
\pgfsetdash{}{0pt}%
\pgfpathmoveto{\pgfqpoint{3.808137in}{0.495848in}}%
\pgfpathlineto{\pgfqpoint{3.807770in}{0.499125in}}%
\pgfpathlineto{\pgfqpoint{3.810929in}{0.498882in}}%
\pgfpathlineto{\pgfqpoint{3.810709in}{0.503484in}}%
\pgfpathlineto{\pgfqpoint{3.813886in}{0.501000in}}%
\pgfpathlineto{\pgfqpoint{3.812787in}{0.497411in}}%
\pgfpathlineto{\pgfqpoint{3.808137in}{0.495848in}}%
\pgfpathclose%
\pgfusepath{fill}%
\end{pgfscope}%
\begin{pgfscope}%
\pgfpathrectangle{\pgfqpoint{3.625000in}{0.100000in}}{\pgfqpoint{2.989028in}{1.913466in}}%
\pgfusepath{clip}%
\pgfsetbuttcap%
\pgfsetmiterjoin%
\definecolor{currentfill}{rgb}{0.288043,0.635909,0.698731}%
\pgfsetfillcolor{currentfill}%
\pgfsetlinewidth{0.000000pt}%
\definecolor{currentstroke}{rgb}{0.000000,0.000000,0.000000}%
\pgfsetstrokecolor{currentstroke}%
\pgfsetstrokeopacity{0.000000}%
\pgfsetdash{}{0pt}%
\pgfpathmoveto{\pgfqpoint{4.086438in}{0.612633in}}%
\pgfpathlineto{\pgfqpoint{4.081400in}{0.613721in}}%
\pgfpathlineto{\pgfqpoint{4.080318in}{0.617088in}}%
\pgfpathlineto{\pgfqpoint{4.082049in}{0.620453in}}%
\pgfpathlineto{\pgfqpoint{4.085927in}{0.622211in}}%
\pgfpathlineto{\pgfqpoint{4.090367in}{0.613572in}}%
\pgfpathlineto{\pgfqpoint{4.094395in}{0.613189in}}%
\pgfpathlineto{\pgfqpoint{4.097399in}{0.610523in}}%
\pgfpathlineto{\pgfqpoint{4.097340in}{0.606855in}}%
\pgfpathlineto{\pgfqpoint{4.095143in}{0.605097in}}%
\pgfpathlineto{\pgfqpoint{4.098716in}{0.594249in}}%
\pgfpathlineto{\pgfqpoint{4.102625in}{0.590269in}}%
\pgfpathlineto{\pgfqpoint{4.098628in}{0.589006in}}%
\pgfpathlineto{\pgfqpoint{4.095393in}{0.593472in}}%
\pgfpathlineto{\pgfqpoint{4.088363in}{0.592089in}}%
\pgfpathlineto{\pgfqpoint{4.090549in}{0.596509in}}%
\pgfpathlineto{\pgfqpoint{4.089531in}{0.598269in}}%
\pgfpathlineto{\pgfqpoint{4.090128in}{0.602876in}}%
\pgfpathlineto{\pgfqpoint{4.088735in}{0.611204in}}%
\pgfpathlineto{\pgfqpoint{4.086438in}{0.612633in}}%
\pgfpathclose%
\pgfusepath{fill}%
\end{pgfscope}%
\begin{pgfscope}%
\pgfpathrectangle{\pgfqpoint{3.625000in}{0.100000in}}{\pgfqpoint{2.989028in}{1.913466in}}%
\pgfusepath{clip}%
\pgfsetbuttcap%
\pgfsetmiterjoin%
\definecolor{currentfill}{rgb}{0.288043,0.635909,0.698731}%
\pgfsetfillcolor{currentfill}%
\pgfsetlinewidth{0.000000pt}%
\definecolor{currentstroke}{rgb}{0.000000,0.000000,0.000000}%
\pgfsetstrokecolor{currentstroke}%
\pgfsetstrokeopacity{0.000000}%
\pgfsetdash{}{0pt}%
\pgfpathmoveto{\pgfqpoint{4.263441in}{0.364487in}}%
\pgfpathlineto{\pgfqpoint{4.256138in}{0.364569in}}%
\pgfpathlineto{\pgfqpoint{4.256700in}{0.368235in}}%
\pgfpathlineto{\pgfqpoint{4.259623in}{0.369546in}}%
\pgfpathlineto{\pgfqpoint{4.263441in}{0.364487in}}%
\pgfpathclose%
\pgfusepath{fill}%
\end{pgfscope}%
\begin{pgfscope}%
\pgfpathrectangle{\pgfqpoint{3.625000in}{0.100000in}}{\pgfqpoint{2.989028in}{1.913466in}}%
\pgfusepath{clip}%
\pgfsetbuttcap%
\pgfsetmiterjoin%
\definecolor{currentfill}{rgb}{0.288043,0.635909,0.698731}%
\pgfsetfillcolor{currentfill}%
\pgfsetlinewidth{0.000000pt}%
\definecolor{currentstroke}{rgb}{0.000000,0.000000,0.000000}%
\pgfsetstrokecolor{currentstroke}%
\pgfsetstrokeopacity{0.000000}%
\pgfsetdash{}{0pt}%
\pgfpathmoveto{\pgfqpoint{4.253306in}{0.367807in}}%
\pgfpathlineto{\pgfqpoint{4.248074in}{0.366541in}}%
\pgfpathlineto{\pgfqpoint{4.243102in}{0.364373in}}%
\pgfpathlineto{\pgfqpoint{4.241567in}{0.366927in}}%
\pgfpathlineto{\pgfqpoint{4.250117in}{0.368657in}}%
\pgfpathlineto{\pgfqpoint{4.253306in}{0.367807in}}%
\pgfpathclose%
\pgfusepath{fill}%
\end{pgfscope}%
\begin{pgfscope}%
\pgfpathrectangle{\pgfqpoint{3.625000in}{0.100000in}}{\pgfqpoint{2.989028in}{1.913466in}}%
\pgfusepath{clip}%
\pgfsetbuttcap%
\pgfsetmiterjoin%
\definecolor{currentfill}{rgb}{0.288043,0.635909,0.698731}%
\pgfsetfillcolor{currentfill}%
\pgfsetlinewidth{0.000000pt}%
\definecolor{currentstroke}{rgb}{0.000000,0.000000,0.000000}%
\pgfsetstrokecolor{currentstroke}%
\pgfsetstrokeopacity{0.000000}%
\pgfsetdash{}{0pt}%
\pgfpathmoveto{\pgfqpoint{4.173103in}{0.365212in}}%
\pgfpathlineto{\pgfqpoint{4.173303in}{0.362676in}}%
\pgfpathlineto{\pgfqpoint{4.170288in}{0.361100in}}%
\pgfpathlineto{\pgfqpoint{4.161875in}{0.364714in}}%
\pgfpathlineto{\pgfqpoint{4.160570in}{0.363261in}}%
\pgfpathlineto{\pgfqpoint{4.158535in}{0.370138in}}%
\pgfpathlineto{\pgfqpoint{4.162750in}{0.371095in}}%
\pgfpathlineto{\pgfqpoint{4.164698in}{0.368269in}}%
\pgfpathlineto{\pgfqpoint{4.168929in}{0.371807in}}%
\pgfpathlineto{\pgfqpoint{4.173103in}{0.365212in}}%
\pgfpathclose%
\pgfusepath{fill}%
\end{pgfscope}%
\begin{pgfscope}%
\pgfpathrectangle{\pgfqpoint{3.625000in}{0.100000in}}{\pgfqpoint{2.989028in}{1.913466in}}%
\pgfusepath{clip}%
\pgfsetbuttcap%
\pgfsetmiterjoin%
\definecolor{currentfill}{rgb}{0.288043,0.635909,0.698731}%
\pgfsetfillcolor{currentfill}%
\pgfsetlinewidth{0.000000pt}%
\definecolor{currentstroke}{rgb}{0.000000,0.000000,0.000000}%
\pgfsetstrokecolor{currentstroke}%
\pgfsetstrokeopacity{0.000000}%
\pgfsetdash{}{0pt}%
\pgfpathmoveto{\pgfqpoint{4.362890in}{0.224669in}}%
\pgfpathlineto{\pgfqpoint{4.355835in}{0.222042in}}%
\pgfpathlineto{\pgfqpoint{4.352098in}{0.221885in}}%
\pgfpathlineto{\pgfqpoint{4.355139in}{0.227457in}}%
\pgfpathlineto{\pgfqpoint{4.357738in}{0.229535in}}%
\pgfpathlineto{\pgfqpoint{4.357654in}{0.235306in}}%
\pgfpathlineto{\pgfqpoint{4.360479in}{0.245723in}}%
\pgfpathlineto{\pgfqpoint{4.363081in}{0.247702in}}%
\pgfpathlineto{\pgfqpoint{4.369087in}{0.244837in}}%
\pgfpathlineto{\pgfqpoint{4.368210in}{0.243196in}}%
\pgfpathlineto{\pgfqpoint{4.368665in}{0.235310in}}%
\pgfpathlineto{\pgfqpoint{4.366718in}{0.233182in}}%
\pgfpathlineto{\pgfqpoint{4.365881in}{0.227581in}}%
\pgfpathlineto{\pgfqpoint{4.362890in}{0.224669in}}%
\pgfpathclose%
\pgfusepath{fill}%
\end{pgfscope}%
\begin{pgfscope}%
\pgfpathrectangle{\pgfqpoint{3.625000in}{0.100000in}}{\pgfqpoint{2.989028in}{1.913466in}}%
\pgfusepath{clip}%
\pgfsetbuttcap%
\pgfsetmiterjoin%
\definecolor{currentfill}{rgb}{0.288043,0.635909,0.698731}%
\pgfsetfillcolor{currentfill}%
\pgfsetlinewidth{0.000000pt}%
\definecolor{currentstroke}{rgb}{0.000000,0.000000,0.000000}%
\pgfsetstrokecolor{currentstroke}%
\pgfsetstrokeopacity{0.000000}%
\pgfsetdash{}{0pt}%
\pgfpathmoveto{\pgfqpoint{4.347266in}{0.250492in}}%
\pgfpathlineto{\pgfqpoint{4.351399in}{0.243355in}}%
\pgfpathlineto{\pgfqpoint{4.350633in}{0.241653in}}%
\pgfpathlineto{\pgfqpoint{4.356161in}{0.239922in}}%
\pgfpathlineto{\pgfqpoint{4.354574in}{0.233362in}}%
\pgfpathlineto{\pgfqpoint{4.352167in}{0.233665in}}%
\pgfpathlineto{\pgfqpoint{4.345842in}{0.245101in}}%
\pgfpathlineto{\pgfqpoint{4.346742in}{0.239954in}}%
\pgfpathlineto{\pgfqpoint{4.345664in}{0.236990in}}%
\pgfpathlineto{\pgfqpoint{4.342986in}{0.235671in}}%
\pgfpathlineto{\pgfqpoint{4.341463in}{0.237092in}}%
\pgfpathlineto{\pgfqpoint{4.340613in}{0.243682in}}%
\pgfpathlineto{\pgfqpoint{4.341115in}{0.247956in}}%
\pgfpathlineto{\pgfqpoint{4.339331in}{0.249822in}}%
\pgfpathlineto{\pgfqpoint{4.344043in}{0.255220in}}%
\pgfpathlineto{\pgfqpoint{4.347135in}{0.256822in}}%
\pgfpathlineto{\pgfqpoint{4.348165in}{0.254717in}}%
\pgfpathlineto{\pgfqpoint{4.351391in}{0.256346in}}%
\pgfpathlineto{\pgfqpoint{4.353728in}{0.251922in}}%
\pgfpathlineto{\pgfqpoint{4.358783in}{0.246295in}}%
\pgfpathlineto{\pgfqpoint{4.358093in}{0.243746in}}%
\pgfpathlineto{\pgfqpoint{4.355286in}{0.240969in}}%
\pgfpathlineto{\pgfqpoint{4.352691in}{0.242277in}}%
\pgfpathlineto{\pgfqpoint{4.347266in}{0.250492in}}%
\pgfpathclose%
\pgfusepath{fill}%
\end{pgfscope}%
\begin{pgfscope}%
\pgfpathrectangle{\pgfqpoint{3.625000in}{0.100000in}}{\pgfqpoint{2.989028in}{1.913466in}}%
\pgfusepath{clip}%
\pgfsetbuttcap%
\pgfsetmiterjoin%
\definecolor{currentfill}{rgb}{0.288043,0.635909,0.698731}%
\pgfsetfillcolor{currentfill}%
\pgfsetlinewidth{0.000000pt}%
\definecolor{currentstroke}{rgb}{0.000000,0.000000,0.000000}%
\pgfsetstrokecolor{currentstroke}%
\pgfsetstrokeopacity{0.000000}%
\pgfsetdash{}{0pt}%
\pgfpathmoveto{\pgfqpoint{4.142161in}{0.351684in}}%
\pgfpathlineto{\pgfqpoint{4.137041in}{0.351018in}}%
\pgfpathlineto{\pgfqpoint{4.129830in}{0.348246in}}%
\pgfpathlineto{\pgfqpoint{4.128148in}{0.349752in}}%
\pgfpathlineto{\pgfqpoint{4.135680in}{0.351454in}}%
\pgfpathlineto{\pgfqpoint{4.133661in}{0.352939in}}%
\pgfpathlineto{\pgfqpoint{4.128991in}{0.354027in}}%
\pgfpathlineto{\pgfqpoint{4.127744in}{0.357507in}}%
\pgfpathlineto{\pgfqpoint{4.130094in}{0.360506in}}%
\pgfpathlineto{\pgfqpoint{4.132073in}{0.367387in}}%
\pgfpathlineto{\pgfqpoint{4.139301in}{0.368669in}}%
\pgfpathlineto{\pgfqpoint{4.142801in}{0.365696in}}%
\pgfpathlineto{\pgfqpoint{4.146088in}{0.368573in}}%
\pgfpathlineto{\pgfqpoint{4.150211in}{0.367867in}}%
\pgfpathlineto{\pgfqpoint{4.153174in}{0.362553in}}%
\pgfpathlineto{\pgfqpoint{4.155683in}{0.367277in}}%
\pgfpathlineto{\pgfqpoint{4.162804in}{0.356228in}}%
\pgfpathlineto{\pgfqpoint{4.160232in}{0.355345in}}%
\pgfpathlineto{\pgfqpoint{4.158653in}{0.352950in}}%
\pgfpathlineto{\pgfqpoint{4.161636in}{0.350764in}}%
\pgfpathlineto{\pgfqpoint{4.157011in}{0.348715in}}%
\pgfpathlineto{\pgfqpoint{4.153656in}{0.349843in}}%
\pgfpathlineto{\pgfqpoint{4.147321in}{0.353583in}}%
\pgfpathlineto{\pgfqpoint{4.142161in}{0.351684in}}%
\pgfpathclose%
\pgfusepath{fill}%
\end{pgfscope}%
\begin{pgfscope}%
\pgfpathrectangle{\pgfqpoint{3.625000in}{0.100000in}}{\pgfqpoint{2.989028in}{1.913466in}}%
\pgfusepath{clip}%
\pgfsetbuttcap%
\pgfsetmiterjoin%
\definecolor{currentfill}{rgb}{0.288043,0.635909,0.698731}%
\pgfsetfillcolor{currentfill}%
\pgfsetlinewidth{0.000000pt}%
\definecolor{currentstroke}{rgb}{0.000000,0.000000,0.000000}%
\pgfsetstrokecolor{currentstroke}%
\pgfsetstrokeopacity{0.000000}%
\pgfsetdash{}{0pt}%
\pgfpathmoveto{\pgfqpoint{4.342967in}{0.203589in}}%
\pgfpathlineto{\pgfqpoint{4.341952in}{0.213793in}}%
\pgfpathlineto{\pgfqpoint{4.344688in}{0.217539in}}%
\pgfpathlineto{\pgfqpoint{4.341762in}{0.217800in}}%
\pgfpathlineto{\pgfqpoint{4.342832in}{0.225086in}}%
\pgfpathlineto{\pgfqpoint{4.344093in}{0.229412in}}%
\pgfpathlineto{\pgfqpoint{4.342997in}{0.235224in}}%
\pgfpathlineto{\pgfqpoint{4.345677in}{0.236349in}}%
\pgfpathlineto{\pgfqpoint{4.346551in}{0.237852in}}%
\pgfpathlineto{\pgfqpoint{4.349191in}{0.237282in}}%
\pgfpathlineto{\pgfqpoint{4.351476in}{0.232626in}}%
\pgfpathlineto{\pgfqpoint{4.351901in}{0.228360in}}%
\pgfpathlineto{\pgfqpoint{4.349105in}{0.215583in}}%
\pgfpathlineto{\pgfqpoint{4.342967in}{0.203589in}}%
\pgfpathclose%
\pgfusepath{fill}%
\end{pgfscope}%
\begin{pgfscope}%
\pgfpathrectangle{\pgfqpoint{3.625000in}{0.100000in}}{\pgfqpoint{2.989028in}{1.913466in}}%
\pgfusepath{clip}%
\pgfsetbuttcap%
\pgfsetmiterjoin%
\definecolor{currentfill}{rgb}{0.288043,0.635909,0.698731}%
\pgfsetfillcolor{currentfill}%
\pgfsetlinewidth{0.000000pt}%
\definecolor{currentstroke}{rgb}{0.000000,0.000000,0.000000}%
\pgfsetstrokecolor{currentstroke}%
\pgfsetstrokeopacity{0.000000}%
\pgfsetdash{}{0pt}%
\pgfpathmoveto{\pgfqpoint{4.342360in}{0.234737in}}%
\pgfpathlineto{\pgfqpoint{4.342219in}{0.229762in}}%
\pgfpathlineto{\pgfqpoint{4.339925in}{0.227503in}}%
\pgfpathlineto{\pgfqpoint{4.337273in}{0.228338in}}%
\pgfpathlineto{\pgfqpoint{4.340625in}{0.235565in}}%
\pgfpathlineto{\pgfqpoint{4.342360in}{0.234737in}}%
\pgfpathclose%
\pgfusepath{fill}%
\end{pgfscope}%
\begin{pgfscope}%
\pgfpathrectangle{\pgfqpoint{3.625000in}{0.100000in}}{\pgfqpoint{2.989028in}{1.913466in}}%
\pgfusepath{clip}%
\pgfsetbuttcap%
\pgfsetmiterjoin%
\definecolor{currentfill}{rgb}{0.288043,0.635909,0.698731}%
\pgfsetfillcolor{currentfill}%
\pgfsetlinewidth{0.000000pt}%
\definecolor{currentstroke}{rgb}{0.000000,0.000000,0.000000}%
\pgfsetstrokecolor{currentstroke}%
\pgfsetstrokeopacity{0.000000}%
\pgfsetdash{}{0pt}%
\pgfpathmoveto{\pgfqpoint{4.370513in}{0.212648in}}%
\pgfpathlineto{\pgfqpoint{4.368557in}{0.204322in}}%
\pgfpathlineto{\pgfqpoint{4.363998in}{0.201690in}}%
\pgfpathlineto{\pgfqpoint{4.358017in}{0.204115in}}%
\pgfpathlineto{\pgfqpoint{4.359665in}{0.207350in}}%
\pgfpathlineto{\pgfqpoint{4.359632in}{0.209591in}}%
\pgfpathlineto{\pgfqpoint{4.361855in}{0.213398in}}%
\pgfpathlineto{\pgfqpoint{4.360292in}{0.212632in}}%
\pgfpathlineto{\pgfqpoint{4.361432in}{0.214491in}}%
\pgfpathlineto{\pgfqpoint{4.359267in}{0.218715in}}%
\pgfpathlineto{\pgfqpoint{4.361723in}{0.219493in}}%
\pgfpathlineto{\pgfqpoint{4.367410in}{0.214551in}}%
\pgfpathlineto{\pgfqpoint{4.370513in}{0.212648in}}%
\pgfpathclose%
\pgfusepath{fill}%
\end{pgfscope}%
\begin{pgfscope}%
\pgfpathrectangle{\pgfqpoint{3.625000in}{0.100000in}}{\pgfqpoint{2.989028in}{1.913466in}}%
\pgfusepath{clip}%
\pgfsetbuttcap%
\pgfsetmiterjoin%
\definecolor{currentfill}{rgb}{0.288043,0.635909,0.698731}%
\pgfsetfillcolor{currentfill}%
\pgfsetlinewidth{0.000000pt}%
\definecolor{currentstroke}{rgb}{0.000000,0.000000,0.000000}%
\pgfsetstrokecolor{currentstroke}%
\pgfsetstrokeopacity{0.000000}%
\pgfsetdash{}{0pt}%
\pgfpathmoveto{\pgfqpoint{4.350557in}{0.197926in}}%
\pgfpathlineto{\pgfqpoint{4.347979in}{0.196923in}}%
\pgfpathlineto{\pgfqpoint{4.349451in}{0.205709in}}%
\pgfpathlineto{\pgfqpoint{4.352546in}{0.207760in}}%
\pgfpathlineto{\pgfqpoint{4.351791in}{0.214343in}}%
\pgfpathlineto{\pgfqpoint{4.353049in}{0.217262in}}%
\pgfpathlineto{\pgfqpoint{4.355536in}{0.218356in}}%
\pgfpathlineto{\pgfqpoint{4.358148in}{0.215493in}}%
\pgfpathlineto{\pgfqpoint{4.360389in}{0.211607in}}%
\pgfpathlineto{\pgfqpoint{4.353184in}{0.203995in}}%
\pgfpathlineto{\pgfqpoint{4.350557in}{0.197926in}}%
\pgfpathclose%
\pgfusepath{fill}%
\end{pgfscope}%
\begin{pgfscope}%
\pgfpathrectangle{\pgfqpoint{3.625000in}{0.100000in}}{\pgfqpoint{2.989028in}{1.913466in}}%
\pgfusepath{clip}%
\pgfsetbuttcap%
\pgfsetmiterjoin%
\definecolor{currentfill}{rgb}{0.288043,0.635909,0.698731}%
\pgfsetfillcolor{currentfill}%
\pgfsetlinewidth{0.000000pt}%
\definecolor{currentstroke}{rgb}{0.000000,0.000000,0.000000}%
\pgfsetstrokecolor{currentstroke}%
\pgfsetstrokeopacity{0.000000}%
\pgfsetdash{}{0pt}%
\pgfpathmoveto{\pgfqpoint{4.371657in}{0.207083in}}%
\pgfpathlineto{\pgfqpoint{4.372921in}{0.201020in}}%
\pgfpathlineto{\pgfqpoint{4.367275in}{0.201714in}}%
\pgfpathlineto{\pgfqpoint{4.371657in}{0.207083in}}%
\pgfpathclose%
\pgfusepath{fill}%
\end{pgfscope}%
\begin{pgfscope}%
\pgfpathrectangle{\pgfqpoint{3.625000in}{0.100000in}}{\pgfqpoint{2.989028in}{1.913466in}}%
\pgfusepath{clip}%
\pgfsetbuttcap%
\pgfsetmiterjoin%
\definecolor{currentfill}{rgb}{0.288043,0.635909,0.698731}%
\pgfsetfillcolor{currentfill}%
\pgfsetlinewidth{0.000000pt}%
\definecolor{currentstroke}{rgb}{0.000000,0.000000,0.000000}%
\pgfsetstrokecolor{currentstroke}%
\pgfsetstrokeopacity{0.000000}%
\pgfsetdash{}{0pt}%
\pgfpathmoveto{\pgfqpoint{4.374351in}{0.197819in}}%
\pgfpathlineto{\pgfqpoint{4.375276in}{0.194079in}}%
\pgfpathlineto{\pgfqpoint{4.377047in}{0.192882in}}%
\pgfpathlineto{\pgfqpoint{4.376725in}{0.189362in}}%
\pgfpathlineto{\pgfqpoint{4.374251in}{0.188177in}}%
\pgfpathlineto{\pgfqpoint{4.372242in}{0.193433in}}%
\pgfpathlineto{\pgfqpoint{4.374351in}{0.197819in}}%
\pgfpathclose%
\pgfusepath{fill}%
\end{pgfscope}%
\begin{pgfscope}%
\pgfpathrectangle{\pgfqpoint{3.625000in}{0.100000in}}{\pgfqpoint{2.989028in}{1.913466in}}%
\pgfusepath{clip}%
\pgfsetbuttcap%
\pgfsetmiterjoin%
\definecolor{currentfill}{rgb}{0.288043,0.635909,0.698731}%
\pgfsetfillcolor{currentfill}%
\pgfsetlinewidth{0.000000pt}%
\definecolor{currentstroke}{rgb}{0.000000,0.000000,0.000000}%
\pgfsetstrokecolor{currentstroke}%
\pgfsetstrokeopacity{0.000000}%
\pgfsetdash{}{0pt}%
\pgfpathmoveto{\pgfqpoint{4.369640in}{0.199254in}}%
\pgfpathlineto{\pgfqpoint{4.368104in}{0.194909in}}%
\pgfpathlineto{\pgfqpoint{4.365719in}{0.195069in}}%
\pgfpathlineto{\pgfqpoint{4.364215in}{0.198652in}}%
\pgfpathlineto{\pgfqpoint{4.366665in}{0.200349in}}%
\pgfpathlineto{\pgfqpoint{4.369640in}{0.199254in}}%
\pgfpathclose%
\pgfusepath{fill}%
\end{pgfscope}%
\begin{pgfscope}%
\pgfpathrectangle{\pgfqpoint{3.625000in}{0.100000in}}{\pgfqpoint{2.989028in}{1.913466in}}%
\pgfusepath{clip}%
\pgfsetbuttcap%
\pgfsetmiterjoin%
\definecolor{currentfill}{rgb}{0.288043,0.635909,0.698731}%
\pgfsetfillcolor{currentfill}%
\pgfsetlinewidth{0.000000pt}%
\definecolor{currentstroke}{rgb}{0.000000,0.000000,0.000000}%
\pgfsetstrokecolor{currentstroke}%
\pgfsetstrokeopacity{0.000000}%
\pgfsetdash{}{0pt}%
\pgfpathmoveto{\pgfqpoint{4.363205in}{0.176854in}}%
\pgfpathlineto{\pgfqpoint{4.365333in}{0.174527in}}%
\pgfpathlineto{\pgfqpoint{4.365802in}{0.169622in}}%
\pgfpathlineto{\pgfqpoint{4.366922in}{0.168197in}}%
\pgfpathlineto{\pgfqpoint{4.365092in}{0.163758in}}%
\pgfpathlineto{\pgfqpoint{4.362686in}{0.161260in}}%
\pgfpathlineto{\pgfqpoint{4.361405in}{0.156623in}}%
\pgfpathlineto{\pgfqpoint{4.358774in}{0.163848in}}%
\pgfpathlineto{\pgfqpoint{4.358338in}{0.170408in}}%
\pgfpathlineto{\pgfqpoint{4.356865in}{0.173673in}}%
\pgfpathlineto{\pgfqpoint{4.351795in}{0.176321in}}%
\pgfpathlineto{\pgfqpoint{4.356551in}{0.181841in}}%
\pgfpathlineto{\pgfqpoint{4.353380in}{0.184201in}}%
\pgfpathlineto{\pgfqpoint{4.356174in}{0.186552in}}%
\pgfpathlineto{\pgfqpoint{4.359882in}{0.195579in}}%
\pgfpathlineto{\pgfqpoint{4.356038in}{0.198762in}}%
\pgfpathlineto{\pgfqpoint{4.357519in}{0.201818in}}%
\pgfpathlineto{\pgfqpoint{4.362487in}{0.199107in}}%
\pgfpathlineto{\pgfqpoint{4.362990in}{0.194829in}}%
\pgfpathlineto{\pgfqpoint{4.361111in}{0.192394in}}%
\pgfpathlineto{\pgfqpoint{4.364836in}{0.189098in}}%
\pgfpathlineto{\pgfqpoint{4.366009in}{0.183965in}}%
\pgfpathlineto{\pgfqpoint{4.365850in}{0.176765in}}%
\pgfpathlineto{\pgfqpoint{4.363205in}{0.176854in}}%
\pgfpathclose%
\pgfusepath{fill}%
\end{pgfscope}%
\begin{pgfscope}%
\pgfpathrectangle{\pgfqpoint{3.625000in}{0.100000in}}{\pgfqpoint{2.989028in}{1.913466in}}%
\pgfusepath{clip}%
\pgfsetbuttcap%
\pgfsetmiterjoin%
\definecolor{currentfill}{rgb}{0.288043,0.635909,0.698731}%
\pgfsetfillcolor{currentfill}%
\pgfsetlinewidth{0.000000pt}%
\definecolor{currentstroke}{rgb}{0.000000,0.000000,0.000000}%
\pgfsetstrokecolor{currentstroke}%
\pgfsetstrokeopacity{0.000000}%
\pgfsetdash{}{0pt}%
\pgfpathmoveto{\pgfqpoint{4.372302in}{0.194926in}}%
\pgfpathlineto{\pgfqpoint{4.371143in}{0.192429in}}%
\pgfpathlineto{\pgfqpoint{4.372231in}{0.188374in}}%
\pgfpathlineto{\pgfqpoint{4.369445in}{0.188033in}}%
\pgfpathlineto{\pgfqpoint{4.366524in}{0.189911in}}%
\pgfpathlineto{\pgfqpoint{4.367472in}{0.194046in}}%
\pgfpathlineto{\pgfqpoint{4.372302in}{0.194926in}}%
\pgfpathclose%
\pgfusepath{fill}%
\end{pgfscope}%
\begin{pgfscope}%
\pgfpathrectangle{\pgfqpoint{3.625000in}{0.100000in}}{\pgfqpoint{2.989028in}{1.913466in}}%
\pgfusepath{clip}%
\pgfsetbuttcap%
\pgfsetmiterjoin%
\definecolor{currentfill}{rgb}{0.288043,0.635909,0.698731}%
\pgfsetfillcolor{currentfill}%
\pgfsetlinewidth{0.000000pt}%
\definecolor{currentstroke}{rgb}{0.000000,0.000000,0.000000}%
\pgfsetstrokecolor{currentstroke}%
\pgfsetstrokeopacity{0.000000}%
\pgfsetdash{}{0pt}%
\pgfpathmoveto{\pgfqpoint{4.359316in}{0.195209in}}%
\pgfpathlineto{\pgfqpoint{4.353167in}{0.192451in}}%
\pgfpathlineto{\pgfqpoint{4.351076in}{0.193444in}}%
\pgfpathlineto{\pgfqpoint{4.355270in}{0.197082in}}%
\pgfpathlineto{\pgfqpoint{4.359316in}{0.195209in}}%
\pgfpathclose%
\pgfusepath{fill}%
\end{pgfscope}%
\begin{pgfscope}%
\pgfpathrectangle{\pgfqpoint{3.625000in}{0.100000in}}{\pgfqpoint{2.989028in}{1.913466in}}%
\pgfusepath{clip}%
\pgfsetbuttcap%
\pgfsetmiterjoin%
\definecolor{currentfill}{rgb}{0.288043,0.635909,0.698731}%
\pgfsetfillcolor{currentfill}%
\pgfsetlinewidth{0.000000pt}%
\definecolor{currentstroke}{rgb}{0.000000,0.000000,0.000000}%
\pgfsetstrokecolor{currentstroke}%
\pgfsetstrokeopacity{0.000000}%
\pgfsetdash{}{0pt}%
\pgfpathmoveto{\pgfqpoint{4.372369in}{0.168302in}}%
\pgfpathlineto{\pgfqpoint{4.370903in}{0.172455in}}%
\pgfpathlineto{\pgfqpoint{4.374004in}{0.173477in}}%
\pgfpathlineto{\pgfqpoint{4.374943in}{0.177955in}}%
\pgfpathlineto{\pgfqpoint{4.378303in}{0.177960in}}%
\pgfpathlineto{\pgfqpoint{4.378289in}{0.181079in}}%
\pgfpathlineto{\pgfqpoint{4.382959in}{0.180341in}}%
\pgfpathlineto{\pgfqpoint{4.383755in}{0.171887in}}%
\pgfpathlineto{\pgfqpoint{4.381046in}{0.166498in}}%
\pgfpathlineto{\pgfqpoint{4.377027in}{0.163131in}}%
\pgfpathlineto{\pgfqpoint{4.372369in}{0.168302in}}%
\pgfpathclose%
\pgfusepath{fill}%
\end{pgfscope}%
\begin{pgfscope}%
\pgfpathrectangle{\pgfqpoint{3.625000in}{0.100000in}}{\pgfqpoint{2.989028in}{1.913466in}}%
\pgfusepath{clip}%
\pgfsetbuttcap%
\pgfsetmiterjoin%
\definecolor{currentfill}{rgb}{0.288043,0.635909,0.698731}%
\pgfsetfillcolor{currentfill}%
\pgfsetlinewidth{0.000000pt}%
\definecolor{currentstroke}{rgb}{0.000000,0.000000,0.000000}%
\pgfsetstrokecolor{currentstroke}%
\pgfsetstrokeopacity{0.000000}%
\pgfsetdash{}{0pt}%
\pgfpathmoveto{\pgfqpoint{4.370553in}{0.171496in}}%
\pgfpathlineto{\pgfqpoint{4.371878in}{0.167564in}}%
\pgfpathlineto{\pgfqpoint{4.368612in}{0.167086in}}%
\pgfpathlineto{\pgfqpoint{4.370553in}{0.171496in}}%
\pgfpathclose%
\pgfusepath{fill}%
\end{pgfscope}%
\begin{pgfscope}%
\pgfpathrectangle{\pgfqpoint{3.625000in}{0.100000in}}{\pgfqpoint{2.989028in}{1.913466in}}%
\pgfusepath{clip}%
\pgfsetbuttcap%
\pgfsetmiterjoin%
\definecolor{currentfill}{rgb}{0.288043,0.635909,0.698731}%
\pgfsetfillcolor{currentfill}%
\pgfsetlinewidth{0.000000pt}%
\definecolor{currentstroke}{rgb}{0.000000,0.000000,0.000000}%
\pgfsetstrokecolor{currentstroke}%
\pgfsetstrokeopacity{0.000000}%
\pgfsetdash{}{0pt}%
\pgfpathmoveto{\pgfqpoint{4.373306in}{0.165713in}}%
\pgfpathlineto{\pgfqpoint{4.372851in}{0.160700in}}%
\pgfpathlineto{\pgfqpoint{4.369406in}{0.161116in}}%
\pgfpathlineto{\pgfqpoint{4.370520in}{0.163537in}}%
\pgfpathlineto{\pgfqpoint{4.373306in}{0.165713in}}%
\pgfpathclose%
\pgfusepath{fill}%
\end{pgfscope}%
\begin{pgfscope}%
\pgfsetbuttcap%
\pgfsetmiterjoin%
\definecolor{currentfill}{rgb}{1.000000,1.000000,1.000000}%
\pgfsetfillcolor{currentfill}%
\pgfsetlinewidth{1.003750pt}%
\definecolor{currentstroke}{rgb}{0.827451,0.827451,0.827451}%
\pgfsetstrokecolor{currentstroke}%
\pgfsetdash{}{0pt}%
\pgfpathmoveto{\pgfqpoint{4.540239in}{1.931644in}}%
\pgfpathlineto{\pgfqpoint{6.189614in}{1.931644in}}%
\pgfpathlineto{\pgfqpoint{6.189614in}{2.127644in}}%
\pgfpathlineto{\pgfqpoint{4.540239in}{2.127644in}}%
\pgfpathlineto{\pgfqpoint{4.540239in}{1.931644in}}%
\pgfpathclose%
\pgfusepath{stroke,fill}%
\end{pgfscope}%
\begin{pgfscope}%
\definecolor{textcolor}{rgb}{0.000000,0.000000,0.000000}%
\pgfsetstrokecolor{textcolor}%
\pgfsetfillcolor{textcolor}%
\pgftext[x=4.581489in,y=1.994331in,left,base]{\color{textcolor}\setmainfont{Lato}\rmfamily\fontsize{9.000000}{10.800000}\selectfont One-Year Growth, 2023 Q3*}%
\end{pgfscope}%
\begin{pgfscope}%
\pgfpathrectangle{\pgfqpoint{3.019583in}{0.169444in}}{\pgfqpoint{0.896708in}{1.339426in}}%
\pgfusepath{clip}%
\pgfsetbuttcap%
\pgfsetmiterjoin%
\definecolor{currentfill}{rgb}{0.619608,0.003922,0.258824}%
\pgfsetfillcolor{currentfill}%
\pgfsetlinewidth{0.000000pt}%
\definecolor{currentstroke}{rgb}{0.000000,0.000000,0.000000}%
\pgfsetstrokecolor{currentstroke}%
\pgfsetstrokeopacity{0.000000}%
\pgfsetdash{}{0pt}%
\pgfpathmoveto{\pgfqpoint{3.279629in}{0.169444in}}%
\pgfpathlineto{\pgfqpoint{3.243760in}{0.169444in}}%
\pgfpathlineto{\pgfqpoint{3.243760in}{0.190373in}}%
\pgfpathlineto{\pgfqpoint{3.279629in}{0.190373in}}%
\pgfpathlineto{\pgfqpoint{3.279629in}{0.169444in}}%
\pgfpathclose%
\pgfusepath{fill}%
\end{pgfscope}%
\begin{pgfscope}%
\pgfpathrectangle{\pgfqpoint{3.019583in}{0.169444in}}{\pgfqpoint{0.896708in}{1.339426in}}%
\pgfusepath{clip}%
\pgfsetbuttcap%
\pgfsetmiterjoin%
\definecolor{currentfill}{rgb}{0.653441,0.041446,0.266820}%
\pgfsetfillcolor{currentfill}%
\pgfsetlinewidth{0.000000pt}%
\definecolor{currentstroke}{rgb}{0.000000,0.000000,0.000000}%
\pgfsetstrokecolor{currentstroke}%
\pgfsetstrokeopacity{0.000000}%
\pgfsetdash{}{0pt}%
\pgfpathmoveto{\pgfqpoint{3.279629in}{0.190373in}}%
\pgfpathlineto{\pgfqpoint{3.243760in}{0.190373in}}%
\pgfpathlineto{\pgfqpoint{3.243760in}{0.211302in}}%
\pgfpathlineto{\pgfqpoint{3.279629in}{0.211302in}}%
\pgfpathlineto{\pgfqpoint{3.279629in}{0.190373in}}%
\pgfpathclose%
\pgfusepath{fill}%
\end{pgfscope}%
\begin{pgfscope}%
\pgfpathrectangle{\pgfqpoint{3.019583in}{0.169444in}}{\pgfqpoint{0.896708in}{1.339426in}}%
\pgfusepath{clip}%
\pgfsetbuttcap%
\pgfsetmiterjoin%
\definecolor{currentfill}{rgb}{0.687274,0.078970,0.274817}%
\pgfsetfillcolor{currentfill}%
\pgfsetlinewidth{0.000000pt}%
\definecolor{currentstroke}{rgb}{0.000000,0.000000,0.000000}%
\pgfsetstrokecolor{currentstroke}%
\pgfsetstrokeopacity{0.000000}%
\pgfsetdash{}{0pt}%
\pgfpathmoveto{\pgfqpoint{3.279629in}{0.211302in}}%
\pgfpathlineto{\pgfqpoint{3.243760in}{0.211302in}}%
\pgfpathlineto{\pgfqpoint{3.243760in}{0.232230in}}%
\pgfpathlineto{\pgfqpoint{3.279629in}{0.232230in}}%
\pgfpathlineto{\pgfqpoint{3.279629in}{0.211302in}}%
\pgfpathclose%
\pgfusepath{fill}%
\end{pgfscope}%
\begin{pgfscope}%
\pgfpathrectangle{\pgfqpoint{3.019583in}{0.169444in}}{\pgfqpoint{0.896708in}{1.339426in}}%
\pgfusepath{clip}%
\pgfsetbuttcap%
\pgfsetmiterjoin%
\definecolor{currentfill}{rgb}{0.721107,0.116494,0.282814}%
\pgfsetfillcolor{currentfill}%
\pgfsetlinewidth{0.000000pt}%
\definecolor{currentstroke}{rgb}{0.000000,0.000000,0.000000}%
\pgfsetstrokecolor{currentstroke}%
\pgfsetstrokeopacity{0.000000}%
\pgfsetdash{}{0pt}%
\pgfpathmoveto{\pgfqpoint{3.279629in}{0.232230in}}%
\pgfpathlineto{\pgfqpoint{3.243760in}{0.232230in}}%
\pgfpathlineto{\pgfqpoint{3.243760in}{0.253159in}}%
\pgfpathlineto{\pgfqpoint{3.279629in}{0.253159in}}%
\pgfpathlineto{\pgfqpoint{3.279629in}{0.232230in}}%
\pgfpathclose%
\pgfusepath{fill}%
\end{pgfscope}%
\begin{pgfscope}%
\pgfpathrectangle{\pgfqpoint{3.019583in}{0.169444in}}{\pgfqpoint{0.896708in}{1.339426in}}%
\pgfusepath{clip}%
\pgfsetbuttcap%
\pgfsetmiterjoin%
\definecolor{currentfill}{rgb}{0.754940,0.154018,0.290811}%
\pgfsetfillcolor{currentfill}%
\pgfsetlinewidth{0.000000pt}%
\definecolor{currentstroke}{rgb}{0.000000,0.000000,0.000000}%
\pgfsetstrokecolor{currentstroke}%
\pgfsetstrokeopacity{0.000000}%
\pgfsetdash{}{0pt}%
\pgfpathmoveto{\pgfqpoint{3.279629in}{0.253159in}}%
\pgfpathlineto{\pgfqpoint{3.243760in}{0.253159in}}%
\pgfpathlineto{\pgfqpoint{3.243760in}{0.274087in}}%
\pgfpathlineto{\pgfqpoint{3.279629in}{0.274087in}}%
\pgfpathlineto{\pgfqpoint{3.279629in}{0.253159in}}%
\pgfpathclose%
\pgfusepath{fill}%
\end{pgfscope}%
\begin{pgfscope}%
\pgfpathrectangle{\pgfqpoint{3.019583in}{0.169444in}}{\pgfqpoint{0.896708in}{1.339426in}}%
\pgfusepath{clip}%
\pgfsetbuttcap%
\pgfsetmiterjoin%
\definecolor{currentfill}{rgb}{0.788774,0.191542,0.298808}%
\pgfsetfillcolor{currentfill}%
\pgfsetlinewidth{0.000000pt}%
\definecolor{currentstroke}{rgb}{0.000000,0.000000,0.000000}%
\pgfsetstrokecolor{currentstroke}%
\pgfsetstrokeopacity{0.000000}%
\pgfsetdash{}{0pt}%
\pgfpathmoveto{\pgfqpoint{3.279629in}{0.274087in}}%
\pgfpathlineto{\pgfqpoint{3.243760in}{0.274087in}}%
\pgfpathlineto{\pgfqpoint{3.243760in}{0.295016in}}%
\pgfpathlineto{\pgfqpoint{3.279629in}{0.295016in}}%
\pgfpathlineto{\pgfqpoint{3.279629in}{0.274087in}}%
\pgfpathclose%
\pgfusepath{fill}%
\end{pgfscope}%
\begin{pgfscope}%
\pgfpathrectangle{\pgfqpoint{3.019583in}{0.169444in}}{\pgfqpoint{0.896708in}{1.339426in}}%
\pgfusepath{clip}%
\pgfsetbuttcap%
\pgfsetmiterjoin%
\definecolor{currentfill}{rgb}{0.822607,0.229066,0.306805}%
\pgfsetfillcolor{currentfill}%
\pgfsetlinewidth{0.000000pt}%
\definecolor{currentstroke}{rgb}{0.000000,0.000000,0.000000}%
\pgfsetstrokecolor{currentstroke}%
\pgfsetstrokeopacity{0.000000}%
\pgfsetdash{}{0pt}%
\pgfpathmoveto{\pgfqpoint{3.279629in}{0.295016in}}%
\pgfpathlineto{\pgfqpoint{3.243760in}{0.295016in}}%
\pgfpathlineto{\pgfqpoint{3.243760in}{0.315944in}}%
\pgfpathlineto{\pgfqpoint{3.279629in}{0.315944in}}%
\pgfpathlineto{\pgfqpoint{3.279629in}{0.295016in}}%
\pgfpathclose%
\pgfusepath{fill}%
\end{pgfscope}%
\begin{pgfscope}%
\pgfpathrectangle{\pgfqpoint{3.019583in}{0.169444in}}{\pgfqpoint{0.896708in}{1.339426in}}%
\pgfusepath{clip}%
\pgfsetbuttcap%
\pgfsetmiterjoin%
\definecolor{currentfill}{rgb}{0.847213,0.261207,0.305190}%
\pgfsetfillcolor{currentfill}%
\pgfsetlinewidth{0.000000pt}%
\definecolor{currentstroke}{rgb}{0.000000,0.000000,0.000000}%
\pgfsetstrokecolor{currentstroke}%
\pgfsetstrokeopacity{0.000000}%
\pgfsetdash{}{0pt}%
\pgfpathmoveto{\pgfqpoint{3.279629in}{0.315944in}}%
\pgfpathlineto{\pgfqpoint{3.243760in}{0.315944in}}%
\pgfpathlineto{\pgfqpoint{3.243760in}{0.336873in}}%
\pgfpathlineto{\pgfqpoint{3.279629in}{0.336873in}}%
\pgfpathlineto{\pgfqpoint{3.279629in}{0.315944in}}%
\pgfpathclose%
\pgfusepath{fill}%
\end{pgfscope}%
\begin{pgfscope}%
\pgfpathrectangle{\pgfqpoint{3.019583in}{0.169444in}}{\pgfqpoint{0.896708in}{1.339426in}}%
\pgfusepath{clip}%
\pgfsetbuttcap%
\pgfsetmiterjoin%
\definecolor{currentfill}{rgb}{0.866282,0.290119,0.297809}%
\pgfsetfillcolor{currentfill}%
\pgfsetlinewidth{0.000000pt}%
\definecolor{currentstroke}{rgb}{0.000000,0.000000,0.000000}%
\pgfsetstrokecolor{currentstroke}%
\pgfsetstrokeopacity{0.000000}%
\pgfsetdash{}{0pt}%
\pgfpathmoveto{\pgfqpoint{3.279629in}{0.336873in}}%
\pgfpathlineto{\pgfqpoint{3.243760in}{0.336873in}}%
\pgfpathlineto{\pgfqpoint{3.243760in}{0.357801in}}%
\pgfpathlineto{\pgfqpoint{3.279629in}{0.357801in}}%
\pgfpathlineto{\pgfqpoint{3.279629in}{0.336873in}}%
\pgfpathclose%
\pgfusepath{fill}%
\end{pgfscope}%
\begin{pgfscope}%
\pgfpathrectangle{\pgfqpoint{3.019583in}{0.169444in}}{\pgfqpoint{0.896708in}{1.339426in}}%
\pgfusepath{clip}%
\pgfsetbuttcap%
\pgfsetmiterjoin%
\definecolor{currentfill}{rgb}{0.885352,0.319031,0.290427}%
\pgfsetfillcolor{currentfill}%
\pgfsetlinewidth{0.000000pt}%
\definecolor{currentstroke}{rgb}{0.000000,0.000000,0.000000}%
\pgfsetstrokecolor{currentstroke}%
\pgfsetstrokeopacity{0.000000}%
\pgfsetdash{}{0pt}%
\pgfpathmoveto{\pgfqpoint{3.279629in}{0.357801in}}%
\pgfpathlineto{\pgfqpoint{3.243760in}{0.357801in}}%
\pgfpathlineto{\pgfqpoint{3.243760in}{0.378730in}}%
\pgfpathlineto{\pgfqpoint{3.279629in}{0.378730in}}%
\pgfpathlineto{\pgfqpoint{3.279629in}{0.357801in}}%
\pgfpathclose%
\pgfusepath{fill}%
\end{pgfscope}%
\begin{pgfscope}%
\pgfpathrectangle{\pgfqpoint{3.019583in}{0.169444in}}{\pgfqpoint{0.896708in}{1.339426in}}%
\pgfusepath{clip}%
\pgfsetbuttcap%
\pgfsetmiterjoin%
\definecolor{currentfill}{rgb}{0.904421,0.347943,0.283045}%
\pgfsetfillcolor{currentfill}%
\pgfsetlinewidth{0.000000pt}%
\definecolor{currentstroke}{rgb}{0.000000,0.000000,0.000000}%
\pgfsetstrokecolor{currentstroke}%
\pgfsetstrokeopacity{0.000000}%
\pgfsetdash{}{0pt}%
\pgfpathmoveto{\pgfqpoint{3.279629in}{0.378730in}}%
\pgfpathlineto{\pgfqpoint{3.243760in}{0.378730in}}%
\pgfpathlineto{\pgfqpoint{3.243760in}{0.399658in}}%
\pgfpathlineto{\pgfqpoint{3.279629in}{0.399658in}}%
\pgfpathlineto{\pgfqpoint{3.279629in}{0.378730in}}%
\pgfpathclose%
\pgfusepath{fill}%
\end{pgfscope}%
\begin{pgfscope}%
\pgfpathrectangle{\pgfqpoint{3.019583in}{0.169444in}}{\pgfqpoint{0.896708in}{1.339426in}}%
\pgfusepath{clip}%
\pgfsetbuttcap%
\pgfsetmiterjoin%
\definecolor{currentfill}{rgb}{0.923491,0.376855,0.275663}%
\pgfsetfillcolor{currentfill}%
\pgfsetlinewidth{0.000000pt}%
\definecolor{currentstroke}{rgb}{0.000000,0.000000,0.000000}%
\pgfsetstrokecolor{currentstroke}%
\pgfsetstrokeopacity{0.000000}%
\pgfsetdash{}{0pt}%
\pgfpathmoveto{\pgfqpoint{3.279629in}{0.399658in}}%
\pgfpathlineto{\pgfqpoint{3.243760in}{0.399658in}}%
\pgfpathlineto{\pgfqpoint{3.243760in}{0.420587in}}%
\pgfpathlineto{\pgfqpoint{3.279629in}{0.420587in}}%
\pgfpathlineto{\pgfqpoint{3.279629in}{0.399658in}}%
\pgfpathclose%
\pgfusepath{fill}%
\end{pgfscope}%
\begin{pgfscope}%
\pgfpathrectangle{\pgfqpoint{3.019583in}{0.169444in}}{\pgfqpoint{0.896708in}{1.339426in}}%
\pgfusepath{clip}%
\pgfsetbuttcap%
\pgfsetmiterjoin%
\definecolor{currentfill}{rgb}{0.942561,0.405767,0.268281}%
\pgfsetfillcolor{currentfill}%
\pgfsetlinewidth{0.000000pt}%
\definecolor{currentstroke}{rgb}{0.000000,0.000000,0.000000}%
\pgfsetstrokecolor{currentstroke}%
\pgfsetstrokeopacity{0.000000}%
\pgfsetdash{}{0pt}%
\pgfpathmoveto{\pgfqpoint{3.279629in}{0.420587in}}%
\pgfpathlineto{\pgfqpoint{3.243760in}{0.420587in}}%
\pgfpathlineto{\pgfqpoint{3.243760in}{0.441515in}}%
\pgfpathlineto{\pgfqpoint{3.279629in}{0.441515in}}%
\pgfpathlineto{\pgfqpoint{3.279629in}{0.420587in}}%
\pgfpathclose%
\pgfusepath{fill}%
\end{pgfscope}%
\begin{pgfscope}%
\pgfpathrectangle{\pgfqpoint{3.019583in}{0.169444in}}{\pgfqpoint{0.896708in}{1.339426in}}%
\pgfusepath{clip}%
\pgfsetbuttcap%
\pgfsetmiterjoin%
\definecolor{currentfill}{rgb}{0.958247,0.437447,0.267359}%
\pgfsetfillcolor{currentfill}%
\pgfsetlinewidth{0.000000pt}%
\definecolor{currentstroke}{rgb}{0.000000,0.000000,0.000000}%
\pgfsetstrokecolor{currentstroke}%
\pgfsetstrokeopacity{0.000000}%
\pgfsetdash{}{0pt}%
\pgfpathmoveto{\pgfqpoint{3.279629in}{0.441515in}}%
\pgfpathlineto{\pgfqpoint{3.243760in}{0.441515in}}%
\pgfpathlineto{\pgfqpoint{3.243760in}{0.462444in}}%
\pgfpathlineto{\pgfqpoint{3.279629in}{0.462444in}}%
\pgfpathlineto{\pgfqpoint{3.279629in}{0.441515in}}%
\pgfpathclose%
\pgfusepath{fill}%
\end{pgfscope}%
\begin{pgfscope}%
\pgfpathrectangle{\pgfqpoint{3.019583in}{0.169444in}}{\pgfqpoint{0.896708in}{1.339426in}}%
\pgfusepath{clip}%
\pgfsetbuttcap%
\pgfsetmiterjoin%
\definecolor{currentfill}{rgb}{0.963783,0.477432,0.285813}%
\pgfsetfillcolor{currentfill}%
\pgfsetlinewidth{0.000000pt}%
\definecolor{currentstroke}{rgb}{0.000000,0.000000,0.000000}%
\pgfsetstrokecolor{currentstroke}%
\pgfsetstrokeopacity{0.000000}%
\pgfsetdash{}{0pt}%
\pgfpathmoveto{\pgfqpoint{3.279629in}{0.462444in}}%
\pgfpathlineto{\pgfqpoint{3.243760in}{0.462444in}}%
\pgfpathlineto{\pgfqpoint{3.243760in}{0.483372in}}%
\pgfpathlineto{\pgfqpoint{3.279629in}{0.483372in}}%
\pgfpathlineto{\pgfqpoint{3.279629in}{0.462444in}}%
\pgfpathclose%
\pgfusepath{fill}%
\end{pgfscope}%
\begin{pgfscope}%
\pgfpathrectangle{\pgfqpoint{3.019583in}{0.169444in}}{\pgfqpoint{0.896708in}{1.339426in}}%
\pgfusepath{clip}%
\pgfsetbuttcap%
\pgfsetmiterjoin%
\definecolor{currentfill}{rgb}{0.969319,0.517416,0.304268}%
\pgfsetfillcolor{currentfill}%
\pgfsetlinewidth{0.000000pt}%
\definecolor{currentstroke}{rgb}{0.000000,0.000000,0.000000}%
\pgfsetstrokecolor{currentstroke}%
\pgfsetstrokeopacity{0.000000}%
\pgfsetdash{}{0pt}%
\pgfpathmoveto{\pgfqpoint{3.279629in}{0.483372in}}%
\pgfpathlineto{\pgfqpoint{3.243760in}{0.483372in}}%
\pgfpathlineto{\pgfqpoint{3.243760in}{0.504301in}}%
\pgfpathlineto{\pgfqpoint{3.279629in}{0.504301in}}%
\pgfpathlineto{\pgfqpoint{3.279629in}{0.483372in}}%
\pgfpathclose%
\pgfusepath{fill}%
\end{pgfscope}%
\begin{pgfscope}%
\pgfpathrectangle{\pgfqpoint{3.019583in}{0.169444in}}{\pgfqpoint{0.896708in}{1.339426in}}%
\pgfusepath{clip}%
\pgfsetbuttcap%
\pgfsetmiterjoin%
\definecolor{currentfill}{rgb}{0.974856,0.557401,0.322722}%
\pgfsetfillcolor{currentfill}%
\pgfsetlinewidth{0.000000pt}%
\definecolor{currentstroke}{rgb}{0.000000,0.000000,0.000000}%
\pgfsetstrokecolor{currentstroke}%
\pgfsetstrokeopacity{0.000000}%
\pgfsetdash{}{0pt}%
\pgfpathmoveto{\pgfqpoint{3.279629in}{0.504301in}}%
\pgfpathlineto{\pgfqpoint{3.243760in}{0.504301in}}%
\pgfpathlineto{\pgfqpoint{3.243760in}{0.525230in}}%
\pgfpathlineto{\pgfqpoint{3.279629in}{0.525230in}}%
\pgfpathlineto{\pgfqpoint{3.279629in}{0.504301in}}%
\pgfpathclose%
\pgfusepath{fill}%
\end{pgfscope}%
\begin{pgfscope}%
\pgfpathrectangle{\pgfqpoint{3.019583in}{0.169444in}}{\pgfqpoint{0.896708in}{1.339426in}}%
\pgfusepath{clip}%
\pgfsetbuttcap%
\pgfsetmiterjoin%
\definecolor{currentfill}{rgb}{0.980392,0.597386,0.341176}%
\pgfsetfillcolor{currentfill}%
\pgfsetlinewidth{0.000000pt}%
\definecolor{currentstroke}{rgb}{0.000000,0.000000,0.000000}%
\pgfsetstrokecolor{currentstroke}%
\pgfsetstrokeopacity{0.000000}%
\pgfsetdash{}{0pt}%
\pgfpathmoveto{\pgfqpoint{3.279629in}{0.525230in}}%
\pgfpathlineto{\pgfqpoint{3.243760in}{0.525230in}}%
\pgfpathlineto{\pgfqpoint{3.243760in}{0.546158in}}%
\pgfpathlineto{\pgfqpoint{3.279629in}{0.546158in}}%
\pgfpathlineto{\pgfqpoint{3.279629in}{0.525230in}}%
\pgfpathclose%
\pgfusepath{fill}%
\end{pgfscope}%
\begin{pgfscope}%
\pgfpathrectangle{\pgfqpoint{3.019583in}{0.169444in}}{\pgfqpoint{0.896708in}{1.339426in}}%
\pgfusepath{clip}%
\pgfsetbuttcap%
\pgfsetmiterjoin%
\definecolor{currentfill}{rgb}{0.985928,0.637370,0.359631}%
\pgfsetfillcolor{currentfill}%
\pgfsetlinewidth{0.000000pt}%
\definecolor{currentstroke}{rgb}{0.000000,0.000000,0.000000}%
\pgfsetstrokecolor{currentstroke}%
\pgfsetstrokeopacity{0.000000}%
\pgfsetdash{}{0pt}%
\pgfpathmoveto{\pgfqpoint{3.279629in}{0.546158in}}%
\pgfpathlineto{\pgfqpoint{3.243760in}{0.546158in}}%
\pgfpathlineto{\pgfqpoint{3.243760in}{0.567087in}}%
\pgfpathlineto{\pgfqpoint{3.279629in}{0.567087in}}%
\pgfpathlineto{\pgfqpoint{3.279629in}{0.546158in}}%
\pgfpathclose%
\pgfusepath{fill}%
\end{pgfscope}%
\begin{pgfscope}%
\pgfpathrectangle{\pgfqpoint{3.019583in}{0.169444in}}{\pgfqpoint{0.896708in}{1.339426in}}%
\pgfusepath{clip}%
\pgfsetbuttcap%
\pgfsetmiterjoin%
\definecolor{currentfill}{rgb}{0.991465,0.677355,0.378085}%
\pgfsetfillcolor{currentfill}%
\pgfsetlinewidth{0.000000pt}%
\definecolor{currentstroke}{rgb}{0.000000,0.000000,0.000000}%
\pgfsetstrokecolor{currentstroke}%
\pgfsetstrokeopacity{0.000000}%
\pgfsetdash{}{0pt}%
\pgfpathmoveto{\pgfqpoint{3.279629in}{0.567087in}}%
\pgfpathlineto{\pgfqpoint{3.243760in}{0.567087in}}%
\pgfpathlineto{\pgfqpoint{3.243760in}{0.588015in}}%
\pgfpathlineto{\pgfqpoint{3.279629in}{0.588015in}}%
\pgfpathlineto{\pgfqpoint{3.279629in}{0.567087in}}%
\pgfpathclose%
\pgfusepath{fill}%
\end{pgfscope}%
\begin{pgfscope}%
\pgfpathrectangle{\pgfqpoint{3.019583in}{0.169444in}}{\pgfqpoint{0.896708in}{1.339426in}}%
\pgfusepath{clip}%
\pgfsetbuttcap%
\pgfsetmiterjoin%
\definecolor{currentfill}{rgb}{0.992695,0.709266,0.402999}%
\pgfsetfillcolor{currentfill}%
\pgfsetlinewidth{0.000000pt}%
\definecolor{currentstroke}{rgb}{0.000000,0.000000,0.000000}%
\pgfsetstrokecolor{currentstroke}%
\pgfsetstrokeopacity{0.000000}%
\pgfsetdash{}{0pt}%
\pgfpathmoveto{\pgfqpoint{3.279629in}{0.588015in}}%
\pgfpathlineto{\pgfqpoint{3.243760in}{0.588015in}}%
\pgfpathlineto{\pgfqpoint{3.243760in}{0.608944in}}%
\pgfpathlineto{\pgfqpoint{3.279629in}{0.608944in}}%
\pgfpathlineto{\pgfqpoint{3.279629in}{0.588015in}}%
\pgfpathclose%
\pgfusepath{fill}%
\end{pgfscope}%
\begin{pgfscope}%
\pgfpathrectangle{\pgfqpoint{3.019583in}{0.169444in}}{\pgfqpoint{0.896708in}{1.339426in}}%
\pgfusepath{clip}%
\pgfsetbuttcap%
\pgfsetmiterjoin%
\definecolor{currentfill}{rgb}{0.993310,0.740023,0.428835}%
\pgfsetfillcolor{currentfill}%
\pgfsetlinewidth{0.000000pt}%
\definecolor{currentstroke}{rgb}{0.000000,0.000000,0.000000}%
\pgfsetstrokecolor{currentstroke}%
\pgfsetstrokeopacity{0.000000}%
\pgfsetdash{}{0pt}%
\pgfpathmoveto{\pgfqpoint{3.279629in}{0.608944in}}%
\pgfpathlineto{\pgfqpoint{3.243760in}{0.608944in}}%
\pgfpathlineto{\pgfqpoint{3.243760in}{0.629872in}}%
\pgfpathlineto{\pgfqpoint{3.279629in}{0.629872in}}%
\pgfpathlineto{\pgfqpoint{3.279629in}{0.608944in}}%
\pgfpathclose%
\pgfusepath{fill}%
\end{pgfscope}%
\begin{pgfscope}%
\pgfpathrectangle{\pgfqpoint{3.019583in}{0.169444in}}{\pgfqpoint{0.896708in}{1.339426in}}%
\pgfusepath{clip}%
\pgfsetbuttcap%
\pgfsetmiterjoin%
\definecolor{currentfill}{rgb}{0.993925,0.770780,0.454671}%
\pgfsetfillcolor{currentfill}%
\pgfsetlinewidth{0.000000pt}%
\definecolor{currentstroke}{rgb}{0.000000,0.000000,0.000000}%
\pgfsetstrokecolor{currentstroke}%
\pgfsetstrokeopacity{0.000000}%
\pgfsetdash{}{0pt}%
\pgfpathmoveto{\pgfqpoint{3.279629in}{0.629872in}}%
\pgfpathlineto{\pgfqpoint{3.243760in}{0.629872in}}%
\pgfpathlineto{\pgfqpoint{3.243760in}{0.650801in}}%
\pgfpathlineto{\pgfqpoint{3.279629in}{0.650801in}}%
\pgfpathlineto{\pgfqpoint{3.279629in}{0.629872in}}%
\pgfpathclose%
\pgfusepath{fill}%
\end{pgfscope}%
\begin{pgfscope}%
\pgfpathrectangle{\pgfqpoint{3.019583in}{0.169444in}}{\pgfqpoint{0.896708in}{1.339426in}}%
\pgfusepath{clip}%
\pgfsetbuttcap%
\pgfsetmiterjoin%
\definecolor{currentfill}{rgb}{0.994541,0.801538,0.480507}%
\pgfsetfillcolor{currentfill}%
\pgfsetlinewidth{0.000000pt}%
\definecolor{currentstroke}{rgb}{0.000000,0.000000,0.000000}%
\pgfsetstrokecolor{currentstroke}%
\pgfsetstrokeopacity{0.000000}%
\pgfsetdash{}{0pt}%
\pgfpathmoveto{\pgfqpoint{3.279629in}{0.650801in}}%
\pgfpathlineto{\pgfqpoint{3.243760in}{0.650801in}}%
\pgfpathlineto{\pgfqpoint{3.243760in}{0.671729in}}%
\pgfpathlineto{\pgfqpoint{3.279629in}{0.671729in}}%
\pgfpathlineto{\pgfqpoint{3.279629in}{0.650801in}}%
\pgfpathclose%
\pgfusepath{fill}%
\end{pgfscope}%
\begin{pgfscope}%
\pgfpathrectangle{\pgfqpoint{3.019583in}{0.169444in}}{\pgfqpoint{0.896708in}{1.339426in}}%
\pgfusepath{clip}%
\pgfsetbuttcap%
\pgfsetmiterjoin%
\definecolor{currentfill}{rgb}{0.995156,0.832295,0.506344}%
\pgfsetfillcolor{currentfill}%
\pgfsetlinewidth{0.000000pt}%
\definecolor{currentstroke}{rgb}{0.000000,0.000000,0.000000}%
\pgfsetstrokecolor{currentstroke}%
\pgfsetstrokeopacity{0.000000}%
\pgfsetdash{}{0pt}%
\pgfpathmoveto{\pgfqpoint{3.279629in}{0.671729in}}%
\pgfpathlineto{\pgfqpoint{3.243760in}{0.671729in}}%
\pgfpathlineto{\pgfqpoint{3.243760in}{0.692658in}}%
\pgfpathlineto{\pgfqpoint{3.279629in}{0.692658in}}%
\pgfpathlineto{\pgfqpoint{3.279629in}{0.671729in}}%
\pgfpathclose%
\pgfusepath{fill}%
\end{pgfscope}%
\begin{pgfscope}%
\pgfpathrectangle{\pgfqpoint{3.019583in}{0.169444in}}{\pgfqpoint{0.896708in}{1.339426in}}%
\pgfusepath{clip}%
\pgfsetbuttcap%
\pgfsetmiterjoin%
\definecolor{currentfill}{rgb}{0.995771,0.863053,0.532180}%
\pgfsetfillcolor{currentfill}%
\pgfsetlinewidth{0.000000pt}%
\definecolor{currentstroke}{rgb}{0.000000,0.000000,0.000000}%
\pgfsetstrokecolor{currentstroke}%
\pgfsetstrokeopacity{0.000000}%
\pgfsetdash{}{0pt}%
\pgfpathmoveto{\pgfqpoint{3.279629in}{0.692658in}}%
\pgfpathlineto{\pgfqpoint{3.243760in}{0.692658in}}%
\pgfpathlineto{\pgfqpoint{3.243760in}{0.713586in}}%
\pgfpathlineto{\pgfqpoint{3.279629in}{0.713586in}}%
\pgfpathlineto{\pgfqpoint{3.279629in}{0.692658in}}%
\pgfpathclose%
\pgfusepath{fill}%
\end{pgfscope}%
\begin{pgfscope}%
\pgfpathrectangle{\pgfqpoint{3.019583in}{0.169444in}}{\pgfqpoint{0.896708in}{1.339426in}}%
\pgfusepath{clip}%
\pgfsetbuttcap%
\pgfsetmiterjoin%
\definecolor{currentfill}{rgb}{0.996386,0.887966,0.561092}%
\pgfsetfillcolor{currentfill}%
\pgfsetlinewidth{0.000000pt}%
\definecolor{currentstroke}{rgb}{0.000000,0.000000,0.000000}%
\pgfsetstrokecolor{currentstroke}%
\pgfsetstrokeopacity{0.000000}%
\pgfsetdash{}{0pt}%
\pgfpathmoveto{\pgfqpoint{3.279629in}{0.713586in}}%
\pgfpathlineto{\pgfqpoint{3.243760in}{0.713586in}}%
\pgfpathlineto{\pgfqpoint{3.243760in}{0.734515in}}%
\pgfpathlineto{\pgfqpoint{3.279629in}{0.734515in}}%
\pgfpathlineto{\pgfqpoint{3.279629in}{0.713586in}}%
\pgfpathclose%
\pgfusepath{fill}%
\end{pgfscope}%
\begin{pgfscope}%
\pgfpathrectangle{\pgfqpoint{3.019583in}{0.169444in}}{\pgfqpoint{0.896708in}{1.339426in}}%
\pgfusepath{clip}%
\pgfsetbuttcap%
\pgfsetmiterjoin%
\definecolor{currentfill}{rgb}{0.997001,0.907036,0.593080}%
\pgfsetfillcolor{currentfill}%
\pgfsetlinewidth{0.000000pt}%
\definecolor{currentstroke}{rgb}{0.000000,0.000000,0.000000}%
\pgfsetstrokecolor{currentstroke}%
\pgfsetstrokeopacity{0.000000}%
\pgfsetdash{}{0pt}%
\pgfpathmoveto{\pgfqpoint{3.279629in}{0.734515in}}%
\pgfpathlineto{\pgfqpoint{3.243760in}{0.734515in}}%
\pgfpathlineto{\pgfqpoint{3.243760in}{0.755443in}}%
\pgfpathlineto{\pgfqpoint{3.279629in}{0.755443in}}%
\pgfpathlineto{\pgfqpoint{3.279629in}{0.734515in}}%
\pgfpathclose%
\pgfusepath{fill}%
\end{pgfscope}%
\begin{pgfscope}%
\pgfpathrectangle{\pgfqpoint{3.019583in}{0.169444in}}{\pgfqpoint{0.896708in}{1.339426in}}%
\pgfusepath{clip}%
\pgfsetbuttcap%
\pgfsetmiterjoin%
\definecolor{currentfill}{rgb}{0.997616,0.926105,0.625067}%
\pgfsetfillcolor{currentfill}%
\pgfsetlinewidth{0.000000pt}%
\definecolor{currentstroke}{rgb}{0.000000,0.000000,0.000000}%
\pgfsetstrokecolor{currentstroke}%
\pgfsetstrokeopacity{0.000000}%
\pgfsetdash{}{0pt}%
\pgfpathmoveto{\pgfqpoint{3.279629in}{0.755443in}}%
\pgfpathlineto{\pgfqpoint{3.243760in}{0.755443in}}%
\pgfpathlineto{\pgfqpoint{3.243760in}{0.776372in}}%
\pgfpathlineto{\pgfqpoint{3.279629in}{0.776372in}}%
\pgfpathlineto{\pgfqpoint{3.279629in}{0.755443in}}%
\pgfpathclose%
\pgfusepath{fill}%
\end{pgfscope}%
\begin{pgfscope}%
\pgfpathrectangle{\pgfqpoint{3.019583in}{0.169444in}}{\pgfqpoint{0.896708in}{1.339426in}}%
\pgfusepath{clip}%
\pgfsetbuttcap%
\pgfsetmiterjoin%
\definecolor{currentfill}{rgb}{0.998231,0.945175,0.657055}%
\pgfsetfillcolor{currentfill}%
\pgfsetlinewidth{0.000000pt}%
\definecolor{currentstroke}{rgb}{0.000000,0.000000,0.000000}%
\pgfsetstrokecolor{currentstroke}%
\pgfsetstrokeopacity{0.000000}%
\pgfsetdash{}{0pt}%
\pgfpathmoveto{\pgfqpoint{3.279629in}{0.776372in}}%
\pgfpathlineto{\pgfqpoint{3.243760in}{0.776372in}}%
\pgfpathlineto{\pgfqpoint{3.243760in}{0.797300in}}%
\pgfpathlineto{\pgfqpoint{3.279629in}{0.797300in}}%
\pgfpathlineto{\pgfqpoint{3.279629in}{0.776372in}}%
\pgfpathclose%
\pgfusepath{fill}%
\end{pgfscope}%
\begin{pgfscope}%
\pgfpathrectangle{\pgfqpoint{3.019583in}{0.169444in}}{\pgfqpoint{0.896708in}{1.339426in}}%
\pgfusepath{clip}%
\pgfsetbuttcap%
\pgfsetmiterjoin%
\definecolor{currentfill}{rgb}{0.998847,0.964245,0.689043}%
\pgfsetfillcolor{currentfill}%
\pgfsetlinewidth{0.000000pt}%
\definecolor{currentstroke}{rgb}{0.000000,0.000000,0.000000}%
\pgfsetstrokecolor{currentstroke}%
\pgfsetstrokeopacity{0.000000}%
\pgfsetdash{}{0pt}%
\pgfpathmoveto{\pgfqpoint{3.279629in}{0.797300in}}%
\pgfpathlineto{\pgfqpoint{3.243760in}{0.797300in}}%
\pgfpathlineto{\pgfqpoint{3.243760in}{0.818229in}}%
\pgfpathlineto{\pgfqpoint{3.279629in}{0.818229in}}%
\pgfpathlineto{\pgfqpoint{3.279629in}{0.797300in}}%
\pgfpathclose%
\pgfusepath{fill}%
\end{pgfscope}%
\begin{pgfscope}%
\pgfpathrectangle{\pgfqpoint{3.019583in}{0.169444in}}{\pgfqpoint{0.896708in}{1.339426in}}%
\pgfusepath{clip}%
\pgfsetbuttcap%
\pgfsetmiterjoin%
\definecolor{currentfill}{rgb}{0.999462,0.983314,0.721030}%
\pgfsetfillcolor{currentfill}%
\pgfsetlinewidth{0.000000pt}%
\definecolor{currentstroke}{rgb}{0.000000,0.000000,0.000000}%
\pgfsetstrokecolor{currentstroke}%
\pgfsetstrokeopacity{0.000000}%
\pgfsetdash{}{0pt}%
\pgfpathmoveto{\pgfqpoint{3.279629in}{0.818229in}}%
\pgfpathlineto{\pgfqpoint{3.243760in}{0.818229in}}%
\pgfpathlineto{\pgfqpoint{3.243760in}{0.839157in}}%
\pgfpathlineto{\pgfqpoint{3.279629in}{0.839157in}}%
\pgfpathlineto{\pgfqpoint{3.279629in}{0.818229in}}%
\pgfpathclose%
\pgfusepath{fill}%
\end{pgfscope}%
\begin{pgfscope}%
\pgfpathrectangle{\pgfqpoint{3.019583in}{0.169444in}}{\pgfqpoint{0.896708in}{1.339426in}}%
\pgfusepath{clip}%
\pgfsetbuttcap%
\pgfsetmiterjoin%
\definecolor{currentfill}{rgb}{0.998078,0.999231,0.746021}%
\pgfsetfillcolor{currentfill}%
\pgfsetlinewidth{0.000000pt}%
\definecolor{currentstroke}{rgb}{0.000000,0.000000,0.000000}%
\pgfsetstrokecolor{currentstroke}%
\pgfsetstrokeopacity{0.000000}%
\pgfsetdash{}{0pt}%
\pgfpathmoveto{\pgfqpoint{3.279629in}{0.839157in}}%
\pgfpathlineto{\pgfqpoint{3.243760in}{0.839157in}}%
\pgfpathlineto{\pgfqpoint{3.243760in}{0.860086in}}%
\pgfpathlineto{\pgfqpoint{3.279629in}{0.860086in}}%
\pgfpathlineto{\pgfqpoint{3.279629in}{0.839157in}}%
\pgfpathclose%
\pgfusepath{fill}%
\end{pgfscope}%
\begin{pgfscope}%
\pgfpathrectangle{\pgfqpoint{3.019583in}{0.169444in}}{\pgfqpoint{0.896708in}{1.339426in}}%
\pgfusepath{clip}%
\pgfsetbuttcap%
\pgfsetmiterjoin%
\definecolor{currentfill}{rgb}{0.982699,0.993080,0.722030}%
\pgfsetfillcolor{currentfill}%
\pgfsetlinewidth{0.000000pt}%
\definecolor{currentstroke}{rgb}{0.000000,0.000000,0.000000}%
\pgfsetstrokecolor{currentstroke}%
\pgfsetstrokeopacity{0.000000}%
\pgfsetdash{}{0pt}%
\pgfpathmoveto{\pgfqpoint{3.279629in}{0.860086in}}%
\pgfpathlineto{\pgfqpoint{3.243760in}{0.860086in}}%
\pgfpathlineto{\pgfqpoint{3.243760in}{0.881015in}}%
\pgfpathlineto{\pgfqpoint{3.279629in}{0.881015in}}%
\pgfpathlineto{\pgfqpoint{3.279629in}{0.860086in}}%
\pgfpathclose%
\pgfusepath{fill}%
\end{pgfscope}%
\begin{pgfscope}%
\pgfpathrectangle{\pgfqpoint{3.019583in}{0.169444in}}{\pgfqpoint{0.896708in}{1.339426in}}%
\pgfusepath{clip}%
\pgfsetbuttcap%
\pgfsetmiterjoin%
\definecolor{currentfill}{rgb}{0.967320,0.986928,0.698039}%
\pgfsetfillcolor{currentfill}%
\pgfsetlinewidth{0.000000pt}%
\definecolor{currentstroke}{rgb}{0.000000,0.000000,0.000000}%
\pgfsetstrokecolor{currentstroke}%
\pgfsetstrokeopacity{0.000000}%
\pgfsetdash{}{0pt}%
\pgfpathmoveto{\pgfqpoint{3.279629in}{0.881015in}}%
\pgfpathlineto{\pgfqpoint{3.243760in}{0.881015in}}%
\pgfpathlineto{\pgfqpoint{3.243760in}{0.901943in}}%
\pgfpathlineto{\pgfqpoint{3.279629in}{0.901943in}}%
\pgfpathlineto{\pgfqpoint{3.279629in}{0.881015in}}%
\pgfpathclose%
\pgfusepath{fill}%
\end{pgfscope}%
\begin{pgfscope}%
\pgfpathrectangle{\pgfqpoint{3.019583in}{0.169444in}}{\pgfqpoint{0.896708in}{1.339426in}}%
\pgfusepath{clip}%
\pgfsetbuttcap%
\pgfsetmiterjoin%
\definecolor{currentfill}{rgb}{0.951942,0.980777,0.674048}%
\pgfsetfillcolor{currentfill}%
\pgfsetlinewidth{0.000000pt}%
\definecolor{currentstroke}{rgb}{0.000000,0.000000,0.000000}%
\pgfsetstrokecolor{currentstroke}%
\pgfsetstrokeopacity{0.000000}%
\pgfsetdash{}{0pt}%
\pgfpathmoveto{\pgfqpoint{3.279629in}{0.901943in}}%
\pgfpathlineto{\pgfqpoint{3.243760in}{0.901943in}}%
\pgfpathlineto{\pgfqpoint{3.243760in}{0.922872in}}%
\pgfpathlineto{\pgfqpoint{3.279629in}{0.922872in}}%
\pgfpathlineto{\pgfqpoint{3.279629in}{0.901943in}}%
\pgfpathclose%
\pgfusepath{fill}%
\end{pgfscope}%
\begin{pgfscope}%
\pgfpathrectangle{\pgfqpoint{3.019583in}{0.169444in}}{\pgfqpoint{0.896708in}{1.339426in}}%
\pgfusepath{clip}%
\pgfsetbuttcap%
\pgfsetmiterjoin%
\definecolor{currentfill}{rgb}{0.936563,0.974625,0.650058}%
\pgfsetfillcolor{currentfill}%
\pgfsetlinewidth{0.000000pt}%
\definecolor{currentstroke}{rgb}{0.000000,0.000000,0.000000}%
\pgfsetstrokecolor{currentstroke}%
\pgfsetstrokeopacity{0.000000}%
\pgfsetdash{}{0pt}%
\pgfpathmoveto{\pgfqpoint{3.279629in}{0.922872in}}%
\pgfpathlineto{\pgfqpoint{3.243760in}{0.922872in}}%
\pgfpathlineto{\pgfqpoint{3.243760in}{0.943800in}}%
\pgfpathlineto{\pgfqpoint{3.279629in}{0.943800in}}%
\pgfpathlineto{\pgfqpoint{3.279629in}{0.922872in}}%
\pgfpathclose%
\pgfusepath{fill}%
\end{pgfscope}%
\begin{pgfscope}%
\pgfpathrectangle{\pgfqpoint{3.019583in}{0.169444in}}{\pgfqpoint{0.896708in}{1.339426in}}%
\pgfusepath{clip}%
\pgfsetbuttcap%
\pgfsetmiterjoin%
\definecolor{currentfill}{rgb}{0.921184,0.968474,0.626067}%
\pgfsetfillcolor{currentfill}%
\pgfsetlinewidth{0.000000pt}%
\definecolor{currentstroke}{rgb}{0.000000,0.000000,0.000000}%
\pgfsetstrokecolor{currentstroke}%
\pgfsetstrokeopacity{0.000000}%
\pgfsetdash{}{0pt}%
\pgfpathmoveto{\pgfqpoint{3.279629in}{0.943800in}}%
\pgfpathlineto{\pgfqpoint{3.243760in}{0.943800in}}%
\pgfpathlineto{\pgfqpoint{3.243760in}{0.964729in}}%
\pgfpathlineto{\pgfqpoint{3.279629in}{0.964729in}}%
\pgfpathlineto{\pgfqpoint{3.279629in}{0.943800in}}%
\pgfpathclose%
\pgfusepath{fill}%
\end{pgfscope}%
\begin{pgfscope}%
\pgfpathrectangle{\pgfqpoint{3.019583in}{0.169444in}}{\pgfqpoint{0.896708in}{1.339426in}}%
\pgfusepath{clip}%
\pgfsetbuttcap%
\pgfsetmiterjoin%
\definecolor{currentfill}{rgb}{0.905805,0.962322,0.602076}%
\pgfsetfillcolor{currentfill}%
\pgfsetlinewidth{0.000000pt}%
\definecolor{currentstroke}{rgb}{0.000000,0.000000,0.000000}%
\pgfsetstrokecolor{currentstroke}%
\pgfsetstrokeopacity{0.000000}%
\pgfsetdash{}{0pt}%
\pgfpathmoveto{\pgfqpoint{3.279629in}{0.964729in}}%
\pgfpathlineto{\pgfqpoint{3.243760in}{0.964729in}}%
\pgfpathlineto{\pgfqpoint{3.243760in}{0.985657in}}%
\pgfpathlineto{\pgfqpoint{3.279629in}{0.985657in}}%
\pgfpathlineto{\pgfqpoint{3.279629in}{0.964729in}}%
\pgfpathclose%
\pgfusepath{fill}%
\end{pgfscope}%
\begin{pgfscope}%
\pgfpathrectangle{\pgfqpoint{3.019583in}{0.169444in}}{\pgfqpoint{0.896708in}{1.339426in}}%
\pgfusepath{clip}%
\pgfsetbuttcap%
\pgfsetmiterjoin%
\definecolor{currentfill}{rgb}{0.874740,0.949712,0.601615}%
\pgfsetfillcolor{currentfill}%
\pgfsetlinewidth{0.000000pt}%
\definecolor{currentstroke}{rgb}{0.000000,0.000000,0.000000}%
\pgfsetstrokecolor{currentstroke}%
\pgfsetstrokeopacity{0.000000}%
\pgfsetdash{}{0pt}%
\pgfpathmoveto{\pgfqpoint{3.279629in}{0.985657in}}%
\pgfpathlineto{\pgfqpoint{3.243760in}{0.985657in}}%
\pgfpathlineto{\pgfqpoint{3.243760in}{1.006586in}}%
\pgfpathlineto{\pgfqpoint{3.279629in}{1.006586in}}%
\pgfpathlineto{\pgfqpoint{3.279629in}{0.985657in}}%
\pgfpathclose%
\pgfusepath{fill}%
\end{pgfscope}%
\begin{pgfscope}%
\pgfpathrectangle{\pgfqpoint{3.019583in}{0.169444in}}{\pgfqpoint{0.896708in}{1.339426in}}%
\pgfusepath{clip}%
\pgfsetbuttcap%
\pgfsetmiterjoin%
\definecolor{currentfill}{rgb}{0.838447,0.934948,0.608997}%
\pgfsetfillcolor{currentfill}%
\pgfsetlinewidth{0.000000pt}%
\definecolor{currentstroke}{rgb}{0.000000,0.000000,0.000000}%
\pgfsetstrokecolor{currentstroke}%
\pgfsetstrokeopacity{0.000000}%
\pgfsetdash{}{0pt}%
\pgfpathmoveto{\pgfqpoint{3.279629in}{1.006586in}}%
\pgfpathlineto{\pgfqpoint{3.243760in}{1.006586in}}%
\pgfpathlineto{\pgfqpoint{3.243760in}{1.027514in}}%
\pgfpathlineto{\pgfqpoint{3.279629in}{1.027514in}}%
\pgfpathlineto{\pgfqpoint{3.279629in}{1.006586in}}%
\pgfpathclose%
\pgfusepath{fill}%
\end{pgfscope}%
\begin{pgfscope}%
\pgfpathrectangle{\pgfqpoint{3.019583in}{0.169444in}}{\pgfqpoint{0.896708in}{1.339426in}}%
\pgfusepath{clip}%
\pgfsetbuttcap%
\pgfsetmiterjoin%
\definecolor{currentfill}{rgb}{0.802153,0.920185,0.616378}%
\pgfsetfillcolor{currentfill}%
\pgfsetlinewidth{0.000000pt}%
\definecolor{currentstroke}{rgb}{0.000000,0.000000,0.000000}%
\pgfsetstrokecolor{currentstroke}%
\pgfsetstrokeopacity{0.000000}%
\pgfsetdash{}{0pt}%
\pgfpathmoveto{\pgfqpoint{3.279629in}{1.027514in}}%
\pgfpathlineto{\pgfqpoint{3.243760in}{1.027514in}}%
\pgfpathlineto{\pgfqpoint{3.243760in}{1.048443in}}%
\pgfpathlineto{\pgfqpoint{3.279629in}{1.048443in}}%
\pgfpathlineto{\pgfqpoint{3.279629in}{1.027514in}}%
\pgfpathclose%
\pgfusepath{fill}%
\end{pgfscope}%
\begin{pgfscope}%
\pgfpathrectangle{\pgfqpoint{3.019583in}{0.169444in}}{\pgfqpoint{0.896708in}{1.339426in}}%
\pgfusepath{clip}%
\pgfsetbuttcap%
\pgfsetmiterjoin%
\definecolor{currentfill}{rgb}{0.765859,0.905421,0.623760}%
\pgfsetfillcolor{currentfill}%
\pgfsetlinewidth{0.000000pt}%
\definecolor{currentstroke}{rgb}{0.000000,0.000000,0.000000}%
\pgfsetstrokecolor{currentstroke}%
\pgfsetstrokeopacity{0.000000}%
\pgfsetdash{}{0pt}%
\pgfpathmoveto{\pgfqpoint{3.279629in}{1.048443in}}%
\pgfpathlineto{\pgfqpoint{3.243760in}{1.048443in}}%
\pgfpathlineto{\pgfqpoint{3.243760in}{1.069371in}}%
\pgfpathlineto{\pgfqpoint{3.279629in}{1.069371in}}%
\pgfpathlineto{\pgfqpoint{3.279629in}{1.048443in}}%
\pgfpathclose%
\pgfusepath{fill}%
\end{pgfscope}%
\begin{pgfscope}%
\pgfpathrectangle{\pgfqpoint{3.019583in}{0.169444in}}{\pgfqpoint{0.896708in}{1.339426in}}%
\pgfusepath{clip}%
\pgfsetbuttcap%
\pgfsetmiterjoin%
\definecolor{currentfill}{rgb}{0.729566,0.890657,0.631142}%
\pgfsetfillcolor{currentfill}%
\pgfsetlinewidth{0.000000pt}%
\definecolor{currentstroke}{rgb}{0.000000,0.000000,0.000000}%
\pgfsetstrokecolor{currentstroke}%
\pgfsetstrokeopacity{0.000000}%
\pgfsetdash{}{0pt}%
\pgfpathmoveto{\pgfqpoint{3.279629in}{1.069371in}}%
\pgfpathlineto{\pgfqpoint{3.243760in}{1.069371in}}%
\pgfpathlineto{\pgfqpoint{3.243760in}{1.090300in}}%
\pgfpathlineto{\pgfqpoint{3.279629in}{1.090300in}}%
\pgfpathlineto{\pgfqpoint{3.279629in}{1.069371in}}%
\pgfpathclose%
\pgfusepath{fill}%
\end{pgfscope}%
\begin{pgfscope}%
\pgfpathrectangle{\pgfqpoint{3.019583in}{0.169444in}}{\pgfqpoint{0.896708in}{1.339426in}}%
\pgfusepath{clip}%
\pgfsetbuttcap%
\pgfsetmiterjoin%
\definecolor{currentfill}{rgb}{0.693272,0.875894,0.638524}%
\pgfsetfillcolor{currentfill}%
\pgfsetlinewidth{0.000000pt}%
\definecolor{currentstroke}{rgb}{0.000000,0.000000,0.000000}%
\pgfsetstrokecolor{currentstroke}%
\pgfsetstrokeopacity{0.000000}%
\pgfsetdash{}{0pt}%
\pgfpathmoveto{\pgfqpoint{3.279629in}{1.090300in}}%
\pgfpathlineto{\pgfqpoint{3.243760in}{1.090300in}}%
\pgfpathlineto{\pgfqpoint{3.243760in}{1.111228in}}%
\pgfpathlineto{\pgfqpoint{3.279629in}{1.111228in}}%
\pgfpathlineto{\pgfqpoint{3.279629in}{1.090300in}}%
\pgfpathclose%
\pgfusepath{fill}%
\end{pgfscope}%
\begin{pgfscope}%
\pgfpathrectangle{\pgfqpoint{3.019583in}{0.169444in}}{\pgfqpoint{0.896708in}{1.339426in}}%
\pgfusepath{clip}%
\pgfsetbuttcap%
\pgfsetmiterjoin%
\definecolor{currentfill}{rgb}{0.654671,0.860438,0.643368}%
\pgfsetfillcolor{currentfill}%
\pgfsetlinewidth{0.000000pt}%
\definecolor{currentstroke}{rgb}{0.000000,0.000000,0.000000}%
\pgfsetstrokecolor{currentstroke}%
\pgfsetstrokeopacity{0.000000}%
\pgfsetdash{}{0pt}%
\pgfpathmoveto{\pgfqpoint{3.279629in}{1.111228in}}%
\pgfpathlineto{\pgfqpoint{3.243760in}{1.111228in}}%
\pgfpathlineto{\pgfqpoint{3.243760in}{1.132157in}}%
\pgfpathlineto{\pgfqpoint{3.279629in}{1.132157in}}%
\pgfpathlineto{\pgfqpoint{3.279629in}{1.111228in}}%
\pgfpathclose%
\pgfusepath{fill}%
\end{pgfscope}%
\begin{pgfscope}%
\pgfpathrectangle{\pgfqpoint{3.019583in}{0.169444in}}{\pgfqpoint{0.896708in}{1.339426in}}%
\pgfusepath{clip}%
\pgfsetbuttcap%
\pgfsetmiterjoin%
\definecolor{currentfill}{rgb}{0.612226,0.843829,0.643983}%
\pgfsetfillcolor{currentfill}%
\pgfsetlinewidth{0.000000pt}%
\definecolor{currentstroke}{rgb}{0.000000,0.000000,0.000000}%
\pgfsetstrokecolor{currentstroke}%
\pgfsetstrokeopacity{0.000000}%
\pgfsetdash{}{0pt}%
\pgfpathmoveto{\pgfqpoint{3.279629in}{1.132157in}}%
\pgfpathlineto{\pgfqpoint{3.243760in}{1.132157in}}%
\pgfpathlineto{\pgfqpoint{3.243760in}{1.153085in}}%
\pgfpathlineto{\pgfqpoint{3.279629in}{1.153085in}}%
\pgfpathlineto{\pgfqpoint{3.279629in}{1.132157in}}%
\pgfpathclose%
\pgfusepath{fill}%
\end{pgfscope}%
\begin{pgfscope}%
\pgfpathrectangle{\pgfqpoint{3.019583in}{0.169444in}}{\pgfqpoint{0.896708in}{1.339426in}}%
\pgfusepath{clip}%
\pgfsetbuttcap%
\pgfsetmiterjoin%
\definecolor{currentfill}{rgb}{0.569781,0.827220,0.644598}%
\pgfsetfillcolor{currentfill}%
\pgfsetlinewidth{0.000000pt}%
\definecolor{currentstroke}{rgb}{0.000000,0.000000,0.000000}%
\pgfsetstrokecolor{currentstroke}%
\pgfsetstrokeopacity{0.000000}%
\pgfsetdash{}{0pt}%
\pgfpathmoveto{\pgfqpoint{3.279629in}{1.153085in}}%
\pgfpathlineto{\pgfqpoint{3.243760in}{1.153085in}}%
\pgfpathlineto{\pgfqpoint{3.243760in}{1.174014in}}%
\pgfpathlineto{\pgfqpoint{3.279629in}{1.174014in}}%
\pgfpathlineto{\pgfqpoint{3.279629in}{1.153085in}}%
\pgfpathclose%
\pgfusepath{fill}%
\end{pgfscope}%
\begin{pgfscope}%
\pgfpathrectangle{\pgfqpoint{3.019583in}{0.169444in}}{\pgfqpoint{0.896708in}{1.339426in}}%
\pgfusepath{clip}%
\pgfsetbuttcap%
\pgfsetmiterjoin%
\definecolor{currentfill}{rgb}{0.527336,0.810611,0.645213}%
\pgfsetfillcolor{currentfill}%
\pgfsetlinewidth{0.000000pt}%
\definecolor{currentstroke}{rgb}{0.000000,0.000000,0.000000}%
\pgfsetstrokecolor{currentstroke}%
\pgfsetstrokeopacity{0.000000}%
\pgfsetdash{}{0pt}%
\pgfpathmoveto{\pgfqpoint{3.279629in}{1.174014in}}%
\pgfpathlineto{\pgfqpoint{3.243760in}{1.174014in}}%
\pgfpathlineto{\pgfqpoint{3.243760in}{1.194943in}}%
\pgfpathlineto{\pgfqpoint{3.279629in}{1.194943in}}%
\pgfpathlineto{\pgfqpoint{3.279629in}{1.174014in}}%
\pgfpathclose%
\pgfusepath{fill}%
\end{pgfscope}%
\begin{pgfscope}%
\pgfpathrectangle{\pgfqpoint{3.019583in}{0.169444in}}{\pgfqpoint{0.896708in}{1.339426in}}%
\pgfusepath{clip}%
\pgfsetbuttcap%
\pgfsetmiterjoin%
\definecolor{currentfill}{rgb}{0.484890,0.794002,0.645829}%
\pgfsetfillcolor{currentfill}%
\pgfsetlinewidth{0.000000pt}%
\definecolor{currentstroke}{rgb}{0.000000,0.000000,0.000000}%
\pgfsetstrokecolor{currentstroke}%
\pgfsetstrokeopacity{0.000000}%
\pgfsetdash{}{0pt}%
\pgfpathmoveto{\pgfqpoint{3.279629in}{1.194943in}}%
\pgfpathlineto{\pgfqpoint{3.243760in}{1.194943in}}%
\pgfpathlineto{\pgfqpoint{3.243760in}{1.215871in}}%
\pgfpathlineto{\pgfqpoint{3.279629in}{1.215871in}}%
\pgfpathlineto{\pgfqpoint{3.279629in}{1.194943in}}%
\pgfpathclose%
\pgfusepath{fill}%
\end{pgfscope}%
\begin{pgfscope}%
\pgfpathrectangle{\pgfqpoint{3.019583in}{0.169444in}}{\pgfqpoint{0.896708in}{1.339426in}}%
\pgfusepath{clip}%
\pgfsetbuttcap%
\pgfsetmiterjoin%
\definecolor{currentfill}{rgb}{0.442445,0.777393,0.646444}%
\pgfsetfillcolor{currentfill}%
\pgfsetlinewidth{0.000000pt}%
\definecolor{currentstroke}{rgb}{0.000000,0.000000,0.000000}%
\pgfsetstrokecolor{currentstroke}%
\pgfsetstrokeopacity{0.000000}%
\pgfsetdash{}{0pt}%
\pgfpathmoveto{\pgfqpoint{3.279629in}{1.215871in}}%
\pgfpathlineto{\pgfqpoint{3.243760in}{1.215871in}}%
\pgfpathlineto{\pgfqpoint{3.243760in}{1.236800in}}%
\pgfpathlineto{\pgfqpoint{3.279629in}{1.236800in}}%
\pgfpathlineto{\pgfqpoint{3.279629in}{1.215871in}}%
\pgfpathclose%
\pgfusepath{fill}%
\end{pgfscope}%
\begin{pgfscope}%
\pgfpathrectangle{\pgfqpoint{3.019583in}{0.169444in}}{\pgfqpoint{0.896708in}{1.339426in}}%
\pgfusepath{clip}%
\pgfsetbuttcap%
\pgfsetmiterjoin%
\definecolor{currentfill}{rgb}{0.400000,0.760784,0.647059}%
\pgfsetfillcolor{currentfill}%
\pgfsetlinewidth{0.000000pt}%
\definecolor{currentstroke}{rgb}{0.000000,0.000000,0.000000}%
\pgfsetstrokecolor{currentstroke}%
\pgfsetstrokeopacity{0.000000}%
\pgfsetdash{}{0pt}%
\pgfpathmoveto{\pgfqpoint{3.279629in}{1.236800in}}%
\pgfpathlineto{\pgfqpoint{3.243760in}{1.236800in}}%
\pgfpathlineto{\pgfqpoint{3.243760in}{1.257728in}}%
\pgfpathlineto{\pgfqpoint{3.279629in}{1.257728in}}%
\pgfpathlineto{\pgfqpoint{3.279629in}{1.236800in}}%
\pgfpathclose%
\pgfusepath{fill}%
\end{pgfscope}%
\begin{pgfscope}%
\pgfpathrectangle{\pgfqpoint{3.019583in}{0.169444in}}{\pgfqpoint{0.896708in}{1.339426in}}%
\pgfusepath{clip}%
\pgfsetbuttcap%
\pgfsetmiterjoin%
\definecolor{currentfill}{rgb}{0.368012,0.725106,0.661822}%
\pgfsetfillcolor{currentfill}%
\pgfsetlinewidth{0.000000pt}%
\definecolor{currentstroke}{rgb}{0.000000,0.000000,0.000000}%
\pgfsetstrokecolor{currentstroke}%
\pgfsetstrokeopacity{0.000000}%
\pgfsetdash{}{0pt}%
\pgfpathmoveto{\pgfqpoint{3.279629in}{1.257728in}}%
\pgfpathlineto{\pgfqpoint{3.243760in}{1.257728in}}%
\pgfpathlineto{\pgfqpoint{3.243760in}{1.278657in}}%
\pgfpathlineto{\pgfqpoint{3.279629in}{1.278657in}}%
\pgfpathlineto{\pgfqpoint{3.279629in}{1.257728in}}%
\pgfpathclose%
\pgfusepath{fill}%
\end{pgfscope}%
\begin{pgfscope}%
\pgfpathrectangle{\pgfqpoint{3.019583in}{0.169444in}}{\pgfqpoint{0.896708in}{1.339426in}}%
\pgfusepath{clip}%
\pgfsetbuttcap%
\pgfsetmiterjoin%
\definecolor{currentfill}{rgb}{0.336025,0.689427,0.676586}%
\pgfsetfillcolor{currentfill}%
\pgfsetlinewidth{0.000000pt}%
\definecolor{currentstroke}{rgb}{0.000000,0.000000,0.000000}%
\pgfsetstrokecolor{currentstroke}%
\pgfsetstrokeopacity{0.000000}%
\pgfsetdash{}{0pt}%
\pgfpathmoveto{\pgfqpoint{3.279629in}{1.278657in}}%
\pgfpathlineto{\pgfqpoint{3.243760in}{1.278657in}}%
\pgfpathlineto{\pgfqpoint{3.243760in}{1.299585in}}%
\pgfpathlineto{\pgfqpoint{3.279629in}{1.299585in}}%
\pgfpathlineto{\pgfqpoint{3.279629in}{1.278657in}}%
\pgfpathclose%
\pgfusepath{fill}%
\end{pgfscope}%
\begin{pgfscope}%
\pgfpathrectangle{\pgfqpoint{3.019583in}{0.169444in}}{\pgfqpoint{0.896708in}{1.339426in}}%
\pgfusepath{clip}%
\pgfsetbuttcap%
\pgfsetmiterjoin%
\definecolor{currentfill}{rgb}{0.304037,0.653749,0.691349}%
\pgfsetfillcolor{currentfill}%
\pgfsetlinewidth{0.000000pt}%
\definecolor{currentstroke}{rgb}{0.000000,0.000000,0.000000}%
\pgfsetstrokecolor{currentstroke}%
\pgfsetstrokeopacity{0.000000}%
\pgfsetdash{}{0pt}%
\pgfpathmoveto{\pgfqpoint{3.279629in}{1.299585in}}%
\pgfpathlineto{\pgfqpoint{3.243760in}{1.299585in}}%
\pgfpathlineto{\pgfqpoint{3.243760in}{1.320514in}}%
\pgfpathlineto{\pgfqpoint{3.279629in}{1.320514in}}%
\pgfpathlineto{\pgfqpoint{3.279629in}{1.299585in}}%
\pgfpathclose%
\pgfusepath{fill}%
\end{pgfscope}%
\begin{pgfscope}%
\pgfpathrectangle{\pgfqpoint{3.019583in}{0.169444in}}{\pgfqpoint{0.896708in}{1.339426in}}%
\pgfusepath{clip}%
\pgfsetbuttcap%
\pgfsetmiterjoin%
\definecolor{currentfill}{rgb}{0.272049,0.618070,0.706113}%
\pgfsetfillcolor{currentfill}%
\pgfsetlinewidth{0.000000pt}%
\definecolor{currentstroke}{rgb}{0.000000,0.000000,0.000000}%
\pgfsetstrokecolor{currentstroke}%
\pgfsetstrokeopacity{0.000000}%
\pgfsetdash{}{0pt}%
\pgfpathmoveto{\pgfqpoint{3.279629in}{1.320514in}}%
\pgfpathlineto{\pgfqpoint{3.243760in}{1.320514in}}%
\pgfpathlineto{\pgfqpoint{3.243760in}{1.341442in}}%
\pgfpathlineto{\pgfqpoint{3.279629in}{1.341442in}}%
\pgfpathlineto{\pgfqpoint{3.279629in}{1.320514in}}%
\pgfpathclose%
\pgfusepath{fill}%
\end{pgfscope}%
\begin{pgfscope}%
\pgfpathrectangle{\pgfqpoint{3.019583in}{0.169444in}}{\pgfqpoint{0.896708in}{1.339426in}}%
\pgfusepath{clip}%
\pgfsetbuttcap%
\pgfsetmiterjoin%
\definecolor{currentfill}{rgb}{0.240062,0.582391,0.720877}%
\pgfsetfillcolor{currentfill}%
\pgfsetlinewidth{0.000000pt}%
\definecolor{currentstroke}{rgb}{0.000000,0.000000,0.000000}%
\pgfsetstrokecolor{currentstroke}%
\pgfsetstrokeopacity{0.000000}%
\pgfsetdash{}{0pt}%
\pgfpathmoveto{\pgfqpoint{3.279629in}{1.341442in}}%
\pgfpathlineto{\pgfqpoint{3.243760in}{1.341442in}}%
\pgfpathlineto{\pgfqpoint{3.243760in}{1.362371in}}%
\pgfpathlineto{\pgfqpoint{3.279629in}{1.362371in}}%
\pgfpathlineto{\pgfqpoint{3.279629in}{1.341442in}}%
\pgfpathclose%
\pgfusepath{fill}%
\end{pgfscope}%
\begin{pgfscope}%
\pgfpathrectangle{\pgfqpoint{3.019583in}{0.169444in}}{\pgfqpoint{0.896708in}{1.339426in}}%
\pgfusepath{clip}%
\pgfsetbuttcap%
\pgfsetmiterjoin%
\definecolor{currentfill}{rgb}{0.208074,0.546713,0.735640}%
\pgfsetfillcolor{currentfill}%
\pgfsetlinewidth{0.000000pt}%
\definecolor{currentstroke}{rgb}{0.000000,0.000000,0.000000}%
\pgfsetstrokecolor{currentstroke}%
\pgfsetstrokeopacity{0.000000}%
\pgfsetdash{}{0pt}%
\pgfpathmoveto{\pgfqpoint{3.279629in}{1.362371in}}%
\pgfpathlineto{\pgfqpoint{3.243760in}{1.362371in}}%
\pgfpathlineto{\pgfqpoint{3.243760in}{1.383299in}}%
\pgfpathlineto{\pgfqpoint{3.279629in}{1.383299in}}%
\pgfpathlineto{\pgfqpoint{3.279629in}{1.362371in}}%
\pgfpathclose%
\pgfusepath{fill}%
\end{pgfscope}%
\begin{pgfscope}%
\pgfpathrectangle{\pgfqpoint{3.019583in}{0.169444in}}{\pgfqpoint{0.896708in}{1.339426in}}%
\pgfusepath{clip}%
\pgfsetbuttcap%
\pgfsetmiterjoin%
\definecolor{currentfill}{rgb}{0.212995,0.511419,0.730796}%
\pgfsetfillcolor{currentfill}%
\pgfsetlinewidth{0.000000pt}%
\definecolor{currentstroke}{rgb}{0.000000,0.000000,0.000000}%
\pgfsetstrokecolor{currentstroke}%
\pgfsetstrokeopacity{0.000000}%
\pgfsetdash{}{0pt}%
\pgfpathmoveto{\pgfqpoint{3.279629in}{1.383299in}}%
\pgfpathlineto{\pgfqpoint{3.243760in}{1.383299in}}%
\pgfpathlineto{\pgfqpoint{3.243760in}{1.404228in}}%
\pgfpathlineto{\pgfqpoint{3.279629in}{1.404228in}}%
\pgfpathlineto{\pgfqpoint{3.279629in}{1.383299in}}%
\pgfpathclose%
\pgfusepath{fill}%
\end{pgfscope}%
\begin{pgfscope}%
\pgfpathrectangle{\pgfqpoint{3.019583in}{0.169444in}}{\pgfqpoint{0.896708in}{1.339426in}}%
\pgfusepath{clip}%
\pgfsetbuttcap%
\pgfsetmiterjoin%
\definecolor{currentfill}{rgb}{0.240062,0.476355,0.714187}%
\pgfsetfillcolor{currentfill}%
\pgfsetlinewidth{0.000000pt}%
\definecolor{currentstroke}{rgb}{0.000000,0.000000,0.000000}%
\pgfsetstrokecolor{currentstroke}%
\pgfsetstrokeopacity{0.000000}%
\pgfsetdash{}{0pt}%
\pgfpathmoveto{\pgfqpoint{3.279629in}{1.404228in}}%
\pgfpathlineto{\pgfqpoint{3.243760in}{1.404228in}}%
\pgfpathlineto{\pgfqpoint{3.243760in}{1.425156in}}%
\pgfpathlineto{\pgfqpoint{3.279629in}{1.425156in}}%
\pgfpathlineto{\pgfqpoint{3.279629in}{1.404228in}}%
\pgfpathclose%
\pgfusepath{fill}%
\end{pgfscope}%
\begin{pgfscope}%
\pgfpathrectangle{\pgfqpoint{3.019583in}{0.169444in}}{\pgfqpoint{0.896708in}{1.339426in}}%
\pgfusepath{clip}%
\pgfsetbuttcap%
\pgfsetmiterjoin%
\definecolor{currentfill}{rgb}{0.267128,0.441292,0.697578}%
\pgfsetfillcolor{currentfill}%
\pgfsetlinewidth{0.000000pt}%
\definecolor{currentstroke}{rgb}{0.000000,0.000000,0.000000}%
\pgfsetstrokecolor{currentstroke}%
\pgfsetstrokeopacity{0.000000}%
\pgfsetdash{}{0pt}%
\pgfpathmoveto{\pgfqpoint{3.279629in}{1.425156in}}%
\pgfpathlineto{\pgfqpoint{3.243760in}{1.425156in}}%
\pgfpathlineto{\pgfqpoint{3.243760in}{1.446085in}}%
\pgfpathlineto{\pgfqpoint{3.279629in}{1.446085in}}%
\pgfpathlineto{\pgfqpoint{3.279629in}{1.425156in}}%
\pgfpathclose%
\pgfusepath{fill}%
\end{pgfscope}%
\begin{pgfscope}%
\pgfpathrectangle{\pgfqpoint{3.019583in}{0.169444in}}{\pgfqpoint{0.896708in}{1.339426in}}%
\pgfusepath{clip}%
\pgfsetbuttcap%
\pgfsetmiterjoin%
\definecolor{currentfill}{rgb}{0.294195,0.406228,0.680969}%
\pgfsetfillcolor{currentfill}%
\pgfsetlinewidth{0.000000pt}%
\definecolor{currentstroke}{rgb}{0.000000,0.000000,0.000000}%
\pgfsetstrokecolor{currentstroke}%
\pgfsetstrokeopacity{0.000000}%
\pgfsetdash{}{0pt}%
\pgfpathmoveto{\pgfqpoint{3.279629in}{1.446085in}}%
\pgfpathlineto{\pgfqpoint{3.243760in}{1.446085in}}%
\pgfpathlineto{\pgfqpoint{3.243760in}{1.467013in}}%
\pgfpathlineto{\pgfqpoint{3.279629in}{1.467013in}}%
\pgfpathlineto{\pgfqpoint{3.279629in}{1.446085in}}%
\pgfpathclose%
\pgfusepath{fill}%
\end{pgfscope}%
\begin{pgfscope}%
\pgfpathrectangle{\pgfqpoint{3.019583in}{0.169444in}}{\pgfqpoint{0.896708in}{1.339426in}}%
\pgfusepath{clip}%
\pgfsetbuttcap%
\pgfsetmiterjoin%
\definecolor{currentfill}{rgb}{0.321261,0.371165,0.664360}%
\pgfsetfillcolor{currentfill}%
\pgfsetlinewidth{0.000000pt}%
\definecolor{currentstroke}{rgb}{0.000000,0.000000,0.000000}%
\pgfsetstrokecolor{currentstroke}%
\pgfsetstrokeopacity{0.000000}%
\pgfsetdash{}{0pt}%
\pgfpathmoveto{\pgfqpoint{3.279629in}{1.467013in}}%
\pgfpathlineto{\pgfqpoint{3.243760in}{1.467013in}}%
\pgfpathlineto{\pgfqpoint{3.243760in}{1.487942in}}%
\pgfpathlineto{\pgfqpoint{3.279629in}{1.487942in}}%
\pgfpathlineto{\pgfqpoint{3.279629in}{1.467013in}}%
\pgfpathclose%
\pgfusepath{fill}%
\end{pgfscope}%
\begin{pgfscope}%
\pgfpathrectangle{\pgfqpoint{3.019583in}{0.169444in}}{\pgfqpoint{0.896708in}{1.339426in}}%
\pgfusepath{clip}%
\pgfsetbuttcap%
\pgfsetmiterjoin%
\definecolor{currentfill}{rgb}{0.348328,0.336101,0.647751}%
\pgfsetfillcolor{currentfill}%
\pgfsetlinewidth{0.000000pt}%
\definecolor{currentstroke}{rgb}{0.000000,0.000000,0.000000}%
\pgfsetstrokecolor{currentstroke}%
\pgfsetstrokeopacity{0.000000}%
\pgfsetdash{}{0pt}%
\pgfpathmoveto{\pgfqpoint{3.279629in}{1.487942in}}%
\pgfpathlineto{\pgfqpoint{3.243760in}{1.487942in}}%
\pgfpathlineto{\pgfqpoint{3.243760in}{1.508871in}}%
\pgfpathlineto{\pgfqpoint{3.279629in}{1.508871in}}%
\pgfpathlineto{\pgfqpoint{3.279629in}{1.487942in}}%
\pgfpathclose%
\pgfusepath{fill}%
\end{pgfscope}%
\begin{pgfscope}%
\pgfpathrectangle{\pgfqpoint{3.019583in}{0.169444in}}{\pgfqpoint{0.896708in}{1.339426in}}%
\pgfusepath{clip}%
\pgfsetbuttcap%
\pgfsetmiterjoin%
\definecolor{currentfill}{rgb}{0.368627,0.309804,0.635294}%
\pgfsetfillcolor{currentfill}%
\pgfsetlinewidth{0.000000pt}%
\definecolor{currentstroke}{rgb}{0.000000,0.000000,0.000000}%
\pgfsetstrokecolor{currentstroke}%
\pgfsetstrokeopacity{0.000000}%
\pgfsetdash{}{0pt}%
\pgfpathmoveto{\pgfqpoint{3.279629in}{1.508871in}}%
\pgfpathlineto{\pgfqpoint{3.243760in}{1.508871in}}%
\pgfpathlineto{\pgfqpoint{3.243760in}{1.529799in}}%
\pgfpathlineto{\pgfqpoint{3.279629in}{1.529799in}}%
\pgfpathlineto{\pgfqpoint{3.279629in}{1.508871in}}%
\pgfpathclose%
\pgfusepath{fill}%
\end{pgfscope}%
\begin{pgfscope}%
\pgfpathrectangle{\pgfqpoint{3.019583in}{0.169444in}}{\pgfqpoint{0.896708in}{1.339426in}}%
\pgfusepath{clip}%
\pgfsetbuttcap%
\pgfsetmiterjoin%
\definecolor{currentfill}{rgb}{0.967320,0.986928,0.698039}%
\pgfsetfillcolor{currentfill}%
\pgfsetlinewidth{0.000000pt}%
\definecolor{currentstroke}{rgb}{0.000000,0.000000,0.000000}%
\pgfsetstrokecolor{currentstroke}%
\pgfsetstrokeopacity{0.000000}%
\pgfsetdash{}{0pt}%
\pgfpathmoveto{\pgfqpoint{3.163057in}{0.881015in}}%
\pgfpathlineto{\pgfqpoint{3.198925in}{0.881015in}}%
\pgfpathlineto{\pgfqpoint{3.198925in}{0.922872in}}%
\pgfpathlineto{\pgfqpoint{3.163057in}{0.922872in}}%
\pgfpathlineto{\pgfqpoint{3.163057in}{0.881015in}}%
\pgfpathclose%
\pgfusepath{fill}%
\end{pgfscope}%
\begin{pgfscope}%
\pgfpathrectangle{\pgfqpoint{3.019583in}{0.169444in}}{\pgfqpoint{0.896708in}{1.339426in}}%
\pgfusepath{clip}%
\pgfsetbuttcap%
\pgfsetmiterjoin%
\definecolor{currentfill}{rgb}{0.838447,0.934948,0.608997}%
\pgfsetfillcolor{currentfill}%
\pgfsetlinewidth{0.000000pt}%
\definecolor{currentstroke}{rgb}{0.000000,0.000000,0.000000}%
\pgfsetstrokecolor{currentstroke}%
\pgfsetstrokeopacity{0.000000}%
\pgfsetdash{}{0pt}%
\pgfpathmoveto{\pgfqpoint{3.180991in}{1.006586in}}%
\pgfpathlineto{\pgfqpoint{3.198925in}{1.006586in}}%
\pgfpathlineto{\pgfqpoint{3.198925in}{1.048443in}}%
\pgfpathlineto{\pgfqpoint{3.180991in}{1.048443in}}%
\pgfpathlineto{\pgfqpoint{3.180991in}{1.006586in}}%
\pgfpathclose%
\pgfusepath{fill}%
\end{pgfscope}%
\begin{pgfscope}%
\pgfpathrectangle{\pgfqpoint{3.019583in}{0.169444in}}{\pgfqpoint{0.896708in}{1.339426in}}%
\pgfusepath{clip}%
\pgfsetbuttcap%
\pgfsetmiterjoin%
\definecolor{currentfill}{rgb}{0.765859,0.905421,0.623760}%
\pgfsetfillcolor{currentfill}%
\pgfsetlinewidth{0.000000pt}%
\definecolor{currentstroke}{rgb}{0.000000,0.000000,0.000000}%
\pgfsetstrokecolor{currentstroke}%
\pgfsetstrokeopacity{0.000000}%
\pgfsetdash{}{0pt}%
\pgfpathmoveto{\pgfqpoint{3.145122in}{1.048443in}}%
\pgfpathlineto{\pgfqpoint{3.198925in}{1.048443in}}%
\pgfpathlineto{\pgfqpoint{3.198925in}{1.090300in}}%
\pgfpathlineto{\pgfqpoint{3.145122in}{1.090300in}}%
\pgfpathlineto{\pgfqpoint{3.145122in}{1.048443in}}%
\pgfpathclose%
\pgfusepath{fill}%
\end{pgfscope}%
\begin{pgfscope}%
\pgfpathrectangle{\pgfqpoint{3.019583in}{0.169444in}}{\pgfqpoint{0.896708in}{1.339426in}}%
\pgfusepath{clip}%
\pgfsetbuttcap%
\pgfsetmiterjoin%
\definecolor{currentfill}{rgb}{0.693272,0.875894,0.638524}%
\pgfsetfillcolor{currentfill}%
\pgfsetlinewidth{0.000000pt}%
\definecolor{currentstroke}{rgb}{0.000000,0.000000,0.000000}%
\pgfsetstrokecolor{currentstroke}%
\pgfsetstrokeopacity{0.000000}%
\pgfsetdash{}{0pt}%
\pgfpathmoveto{\pgfqpoint{3.163057in}{1.090300in}}%
\pgfpathlineto{\pgfqpoint{3.198925in}{1.090300in}}%
\pgfpathlineto{\pgfqpoint{3.198925in}{1.132157in}}%
\pgfpathlineto{\pgfqpoint{3.163057in}{1.132157in}}%
\pgfpathlineto{\pgfqpoint{3.163057in}{1.090300in}}%
\pgfpathclose%
\pgfusepath{fill}%
\end{pgfscope}%
\begin{pgfscope}%
\pgfpathrectangle{\pgfqpoint{3.019583in}{0.169444in}}{\pgfqpoint{0.896708in}{1.339426in}}%
\pgfusepath{clip}%
\pgfsetbuttcap%
\pgfsetmiterjoin%
\definecolor{currentfill}{rgb}{0.612226,0.843829,0.643983}%
\pgfsetfillcolor{currentfill}%
\pgfsetlinewidth{0.000000pt}%
\definecolor{currentstroke}{rgb}{0.000000,0.000000,0.000000}%
\pgfsetstrokecolor{currentstroke}%
\pgfsetstrokeopacity{0.000000}%
\pgfsetdash{}{0pt}%
\pgfpathmoveto{\pgfqpoint{3.091320in}{1.132157in}}%
\pgfpathlineto{\pgfqpoint{3.198925in}{1.132157in}}%
\pgfpathlineto{\pgfqpoint{3.198925in}{1.174014in}}%
\pgfpathlineto{\pgfqpoint{3.091320in}{1.174014in}}%
\pgfpathlineto{\pgfqpoint{3.091320in}{1.132157in}}%
\pgfpathclose%
\pgfusepath{fill}%
\end{pgfscope}%
\begin{pgfscope}%
\pgfpathrectangle{\pgfqpoint{3.019583in}{0.169444in}}{\pgfqpoint{0.896708in}{1.339426in}}%
\pgfusepath{clip}%
\pgfsetbuttcap%
\pgfsetmiterjoin%
\definecolor{currentfill}{rgb}{0.527336,0.810611,0.645213}%
\pgfsetfillcolor{currentfill}%
\pgfsetlinewidth{0.000000pt}%
\definecolor{currentstroke}{rgb}{0.000000,0.000000,0.000000}%
\pgfsetstrokecolor{currentstroke}%
\pgfsetstrokeopacity{0.000000}%
\pgfsetdash{}{0pt}%
\pgfpathmoveto{\pgfqpoint{3.055452in}{1.174014in}}%
\pgfpathlineto{\pgfqpoint{3.198925in}{1.174014in}}%
\pgfpathlineto{\pgfqpoint{3.198925in}{1.215871in}}%
\pgfpathlineto{\pgfqpoint{3.055452in}{1.215871in}}%
\pgfpathlineto{\pgfqpoint{3.055452in}{1.174014in}}%
\pgfpathclose%
\pgfusepath{fill}%
\end{pgfscope}%
\begin{pgfscope}%
\pgfpathrectangle{\pgfqpoint{3.019583in}{0.169444in}}{\pgfqpoint{0.896708in}{1.339426in}}%
\pgfusepath{clip}%
\pgfsetbuttcap%
\pgfsetmiterjoin%
\definecolor{currentfill}{rgb}{0.442445,0.777393,0.646444}%
\pgfsetfillcolor{currentfill}%
\pgfsetlinewidth{0.000000pt}%
\definecolor{currentstroke}{rgb}{0.000000,0.000000,0.000000}%
\pgfsetstrokecolor{currentstroke}%
\pgfsetstrokeopacity{0.000000}%
\pgfsetdash{}{0pt}%
\pgfpathmoveto{\pgfqpoint{3.001649in}{1.215871in}}%
\pgfpathlineto{\pgfqpoint{3.198925in}{1.215871in}}%
\pgfpathlineto{\pgfqpoint{3.198925in}{1.257728in}}%
\pgfpathlineto{\pgfqpoint{3.001649in}{1.257728in}}%
\pgfpathlineto{\pgfqpoint{3.001649in}{1.215871in}}%
\pgfpathclose%
\pgfusepath{fill}%
\end{pgfscope}%
\begin{pgfscope}%
\pgfpathrectangle{\pgfqpoint{3.019583in}{0.169444in}}{\pgfqpoint{0.896708in}{1.339426in}}%
\pgfusepath{clip}%
\pgfsetbuttcap%
\pgfsetmiterjoin%
\definecolor{currentfill}{rgb}{0.368012,0.725106,0.661822}%
\pgfsetfillcolor{currentfill}%
\pgfsetlinewidth{0.000000pt}%
\definecolor{currentstroke}{rgb}{0.000000,0.000000,0.000000}%
\pgfsetstrokecolor{currentstroke}%
\pgfsetstrokeopacity{0.000000}%
\pgfsetdash{}{0pt}%
\pgfpathmoveto{\pgfqpoint{3.127188in}{1.257728in}}%
\pgfpathlineto{\pgfqpoint{3.198925in}{1.257728in}}%
\pgfpathlineto{\pgfqpoint{3.198925in}{1.299585in}}%
\pgfpathlineto{\pgfqpoint{3.127188in}{1.299585in}}%
\pgfpathlineto{\pgfqpoint{3.127188in}{1.257728in}}%
\pgfpathclose%
\pgfusepath{fill}%
\end{pgfscope}%
\begin{pgfscope}%
\pgfpathrectangle{\pgfqpoint{3.019583in}{0.169444in}}{\pgfqpoint{0.896708in}{1.339426in}}%
\pgfusepath{clip}%
\pgfsetbuttcap%
\pgfsetmiterjoin%
\definecolor{currentfill}{rgb}{0.304037,0.653749,0.691349}%
\pgfsetfillcolor{currentfill}%
\pgfsetlinewidth{0.000000pt}%
\definecolor{currentstroke}{rgb}{0.000000,0.000000,0.000000}%
\pgfsetstrokecolor{currentstroke}%
\pgfsetstrokeopacity{0.000000}%
\pgfsetdash{}{0pt}%
\pgfpathmoveto{\pgfqpoint{3.091320in}{1.299585in}}%
\pgfpathlineto{\pgfqpoint{3.198925in}{1.299585in}}%
\pgfpathlineto{\pgfqpoint{3.198925in}{1.341442in}}%
\pgfpathlineto{\pgfqpoint{3.091320in}{1.341442in}}%
\pgfpathlineto{\pgfqpoint{3.091320in}{1.299585in}}%
\pgfpathclose%
\pgfusepath{fill}%
\end{pgfscope}%
\begin{pgfscope}%
\pgfpathrectangle{\pgfqpoint{3.019583in}{0.169444in}}{\pgfqpoint{0.896708in}{1.339426in}}%
\pgfusepath{clip}%
\pgfsetbuttcap%
\pgfsetmiterjoin%
\definecolor{currentfill}{rgb}{0.240062,0.582391,0.720877}%
\pgfsetfillcolor{currentfill}%
\pgfsetlinewidth{0.000000pt}%
\definecolor{currentstroke}{rgb}{0.000000,0.000000,0.000000}%
\pgfsetstrokecolor{currentstroke}%
\pgfsetstrokeopacity{0.000000}%
\pgfsetdash{}{0pt}%
\pgfpathmoveto{\pgfqpoint{3.145122in}{1.341442in}}%
\pgfpathlineto{\pgfqpoint{3.198925in}{1.341442in}}%
\pgfpathlineto{\pgfqpoint{3.198925in}{1.383299in}}%
\pgfpathlineto{\pgfqpoint{3.145122in}{1.383299in}}%
\pgfpathlineto{\pgfqpoint{3.145122in}{1.341442in}}%
\pgfpathclose%
\pgfusepath{fill}%
\end{pgfscope}%
\begin{pgfscope}%
\pgfpathrectangle{\pgfqpoint{3.019583in}{0.169444in}}{\pgfqpoint{0.896708in}{1.339426in}}%
\pgfusepath{clip}%
\pgfsetbuttcap%
\pgfsetmiterjoin%
\definecolor{currentfill}{rgb}{0.212995,0.511419,0.730796}%
\pgfsetfillcolor{currentfill}%
\pgfsetlinewidth{0.000000pt}%
\definecolor{currentstroke}{rgb}{0.000000,0.000000,0.000000}%
\pgfsetstrokecolor{currentstroke}%
\pgfsetstrokeopacity{0.000000}%
\pgfsetdash{}{0pt}%
\pgfpathmoveto{\pgfqpoint{3.180991in}{1.383299in}}%
\pgfpathlineto{\pgfqpoint{3.198925in}{1.383299in}}%
\pgfpathlineto{\pgfqpoint{3.198925in}{1.425156in}}%
\pgfpathlineto{\pgfqpoint{3.180991in}{1.425156in}}%
\pgfpathlineto{\pgfqpoint{3.180991in}{1.383299in}}%
\pgfpathclose%
\pgfusepath{fill}%
\end{pgfscope}%
\begin{pgfscope}%
\pgfpathrectangle{\pgfqpoint{3.019583in}{0.169444in}}{\pgfqpoint{0.896708in}{1.339426in}}%
\pgfusepath{clip}%
\pgfsetbuttcap%
\pgfsetmiterjoin%
\definecolor{currentfill}{rgb}{0.267128,0.441292,0.697578}%
\pgfsetfillcolor{currentfill}%
\pgfsetlinewidth{0.000000pt}%
\definecolor{currentstroke}{rgb}{0.000000,0.000000,0.000000}%
\pgfsetstrokecolor{currentstroke}%
\pgfsetstrokeopacity{0.000000}%
\pgfsetdash{}{0pt}%
\pgfpathmoveto{\pgfqpoint{3.180991in}{1.425156in}}%
\pgfpathlineto{\pgfqpoint{3.198925in}{1.425156in}}%
\pgfpathlineto{\pgfqpoint{3.198925in}{1.467013in}}%
\pgfpathlineto{\pgfqpoint{3.180991in}{1.467013in}}%
\pgfpathlineto{\pgfqpoint{3.180991in}{1.425156in}}%
\pgfpathclose%
\pgfusepath{fill}%
\end{pgfscope}%
\begin{pgfscope}%
\pgfpathrectangle{\pgfqpoint{3.019583in}{0.169444in}}{\pgfqpoint{0.896708in}{1.339426in}}%
\pgfusepath{clip}%
\pgfsetbuttcap%
\pgfsetmiterjoin%
\definecolor{currentfill}{rgb}{0.321261,0.371165,0.664360}%
\pgfsetfillcolor{currentfill}%
\pgfsetlinewidth{0.000000pt}%
\definecolor{currentstroke}{rgb}{0.000000,0.000000,0.000000}%
\pgfsetstrokecolor{currentstroke}%
\pgfsetstrokeopacity{0.000000}%
\pgfsetdash{}{0pt}%
\pgfpathmoveto{\pgfqpoint{3.145122in}{1.467013in}}%
\pgfpathlineto{\pgfqpoint{3.198925in}{1.467013in}}%
\pgfpathlineto{\pgfqpoint{3.198925in}{1.508871in}}%
\pgfpathlineto{\pgfqpoint{3.145122in}{1.508871in}}%
\pgfpathlineto{\pgfqpoint{3.145122in}{1.467013in}}%
\pgfpathclose%
\pgfusepath{fill}%
\end{pgfscope}%
\begin{pgfscope}%
\pgfpathrectangle{\pgfqpoint{3.019583in}{0.169444in}}{\pgfqpoint{0.896708in}{1.339426in}}%
\pgfusepath{clip}%
\pgfsetbuttcap%
\pgfsetmiterjoin%
\definecolor{currentfill}{rgb}{0.619608,0.003922,0.258824}%
\pgfsetfillcolor{currentfill}%
\pgfsetlinewidth{0.000000pt}%
\definecolor{currentstroke}{rgb}{0.000000,0.000000,0.000000}%
\pgfsetstrokecolor{currentstroke}%
\pgfsetstrokeopacity{0.000000}%
\pgfsetdash{}{0pt}%
\pgfpathmoveto{\pgfqpoint{3.521740in}{0.169444in}}%
\pgfpathlineto{\pgfqpoint{3.485872in}{0.169444in}}%
\pgfpathlineto{\pgfqpoint{3.485872in}{0.190373in}}%
\pgfpathlineto{\pgfqpoint{3.521740in}{0.190373in}}%
\pgfpathlineto{\pgfqpoint{3.521740in}{0.169444in}}%
\pgfpathclose%
\pgfusepath{fill}%
\end{pgfscope}%
\begin{pgfscope}%
\pgfpathrectangle{\pgfqpoint{3.019583in}{0.169444in}}{\pgfqpoint{0.896708in}{1.339426in}}%
\pgfusepath{clip}%
\pgfsetbuttcap%
\pgfsetmiterjoin%
\definecolor{currentfill}{rgb}{0.653441,0.041446,0.266820}%
\pgfsetfillcolor{currentfill}%
\pgfsetlinewidth{0.000000pt}%
\definecolor{currentstroke}{rgb}{0.000000,0.000000,0.000000}%
\pgfsetstrokecolor{currentstroke}%
\pgfsetstrokeopacity{0.000000}%
\pgfsetdash{}{0pt}%
\pgfpathmoveto{\pgfqpoint{3.521740in}{0.190373in}}%
\pgfpathlineto{\pgfqpoint{3.485872in}{0.190373in}}%
\pgfpathlineto{\pgfqpoint{3.485872in}{0.211302in}}%
\pgfpathlineto{\pgfqpoint{3.521740in}{0.211302in}}%
\pgfpathlineto{\pgfqpoint{3.521740in}{0.190373in}}%
\pgfpathclose%
\pgfusepath{fill}%
\end{pgfscope}%
\begin{pgfscope}%
\pgfpathrectangle{\pgfqpoint{3.019583in}{0.169444in}}{\pgfqpoint{0.896708in}{1.339426in}}%
\pgfusepath{clip}%
\pgfsetbuttcap%
\pgfsetmiterjoin%
\definecolor{currentfill}{rgb}{0.687274,0.078970,0.274817}%
\pgfsetfillcolor{currentfill}%
\pgfsetlinewidth{0.000000pt}%
\definecolor{currentstroke}{rgb}{0.000000,0.000000,0.000000}%
\pgfsetstrokecolor{currentstroke}%
\pgfsetstrokeopacity{0.000000}%
\pgfsetdash{}{0pt}%
\pgfpathmoveto{\pgfqpoint{3.521740in}{0.211302in}}%
\pgfpathlineto{\pgfqpoint{3.485872in}{0.211302in}}%
\pgfpathlineto{\pgfqpoint{3.485872in}{0.232230in}}%
\pgfpathlineto{\pgfqpoint{3.521740in}{0.232230in}}%
\pgfpathlineto{\pgfqpoint{3.521740in}{0.211302in}}%
\pgfpathclose%
\pgfusepath{fill}%
\end{pgfscope}%
\begin{pgfscope}%
\pgfpathrectangle{\pgfqpoint{3.019583in}{0.169444in}}{\pgfqpoint{0.896708in}{1.339426in}}%
\pgfusepath{clip}%
\pgfsetbuttcap%
\pgfsetmiterjoin%
\definecolor{currentfill}{rgb}{0.721107,0.116494,0.282814}%
\pgfsetfillcolor{currentfill}%
\pgfsetlinewidth{0.000000pt}%
\definecolor{currentstroke}{rgb}{0.000000,0.000000,0.000000}%
\pgfsetstrokecolor{currentstroke}%
\pgfsetstrokeopacity{0.000000}%
\pgfsetdash{}{0pt}%
\pgfpathmoveto{\pgfqpoint{3.521740in}{0.232230in}}%
\pgfpathlineto{\pgfqpoint{3.485872in}{0.232230in}}%
\pgfpathlineto{\pgfqpoint{3.485872in}{0.253159in}}%
\pgfpathlineto{\pgfqpoint{3.521740in}{0.253159in}}%
\pgfpathlineto{\pgfqpoint{3.521740in}{0.232230in}}%
\pgfpathclose%
\pgfusepath{fill}%
\end{pgfscope}%
\begin{pgfscope}%
\pgfpathrectangle{\pgfqpoint{3.019583in}{0.169444in}}{\pgfqpoint{0.896708in}{1.339426in}}%
\pgfusepath{clip}%
\pgfsetbuttcap%
\pgfsetmiterjoin%
\definecolor{currentfill}{rgb}{0.754940,0.154018,0.290811}%
\pgfsetfillcolor{currentfill}%
\pgfsetlinewidth{0.000000pt}%
\definecolor{currentstroke}{rgb}{0.000000,0.000000,0.000000}%
\pgfsetstrokecolor{currentstroke}%
\pgfsetstrokeopacity{0.000000}%
\pgfsetdash{}{0pt}%
\pgfpathmoveto{\pgfqpoint{3.521740in}{0.253159in}}%
\pgfpathlineto{\pgfqpoint{3.485872in}{0.253159in}}%
\pgfpathlineto{\pgfqpoint{3.485872in}{0.274087in}}%
\pgfpathlineto{\pgfqpoint{3.521740in}{0.274087in}}%
\pgfpathlineto{\pgfqpoint{3.521740in}{0.253159in}}%
\pgfpathclose%
\pgfusepath{fill}%
\end{pgfscope}%
\begin{pgfscope}%
\pgfpathrectangle{\pgfqpoint{3.019583in}{0.169444in}}{\pgfqpoint{0.896708in}{1.339426in}}%
\pgfusepath{clip}%
\pgfsetbuttcap%
\pgfsetmiterjoin%
\definecolor{currentfill}{rgb}{0.788774,0.191542,0.298808}%
\pgfsetfillcolor{currentfill}%
\pgfsetlinewidth{0.000000pt}%
\definecolor{currentstroke}{rgb}{0.000000,0.000000,0.000000}%
\pgfsetstrokecolor{currentstroke}%
\pgfsetstrokeopacity{0.000000}%
\pgfsetdash{}{0pt}%
\pgfpathmoveto{\pgfqpoint{3.521740in}{0.274087in}}%
\pgfpathlineto{\pgfqpoint{3.485872in}{0.274087in}}%
\pgfpathlineto{\pgfqpoint{3.485872in}{0.295016in}}%
\pgfpathlineto{\pgfqpoint{3.521740in}{0.295016in}}%
\pgfpathlineto{\pgfqpoint{3.521740in}{0.274087in}}%
\pgfpathclose%
\pgfusepath{fill}%
\end{pgfscope}%
\begin{pgfscope}%
\pgfpathrectangle{\pgfqpoint{3.019583in}{0.169444in}}{\pgfqpoint{0.896708in}{1.339426in}}%
\pgfusepath{clip}%
\pgfsetbuttcap%
\pgfsetmiterjoin%
\definecolor{currentfill}{rgb}{0.822607,0.229066,0.306805}%
\pgfsetfillcolor{currentfill}%
\pgfsetlinewidth{0.000000pt}%
\definecolor{currentstroke}{rgb}{0.000000,0.000000,0.000000}%
\pgfsetstrokecolor{currentstroke}%
\pgfsetstrokeopacity{0.000000}%
\pgfsetdash{}{0pt}%
\pgfpathmoveto{\pgfqpoint{3.521740in}{0.295016in}}%
\pgfpathlineto{\pgfqpoint{3.485872in}{0.295016in}}%
\pgfpathlineto{\pgfqpoint{3.485872in}{0.315944in}}%
\pgfpathlineto{\pgfqpoint{3.521740in}{0.315944in}}%
\pgfpathlineto{\pgfqpoint{3.521740in}{0.295016in}}%
\pgfpathclose%
\pgfusepath{fill}%
\end{pgfscope}%
\begin{pgfscope}%
\pgfpathrectangle{\pgfqpoint{3.019583in}{0.169444in}}{\pgfqpoint{0.896708in}{1.339426in}}%
\pgfusepath{clip}%
\pgfsetbuttcap%
\pgfsetmiterjoin%
\definecolor{currentfill}{rgb}{0.847213,0.261207,0.305190}%
\pgfsetfillcolor{currentfill}%
\pgfsetlinewidth{0.000000pt}%
\definecolor{currentstroke}{rgb}{0.000000,0.000000,0.000000}%
\pgfsetstrokecolor{currentstroke}%
\pgfsetstrokeopacity{0.000000}%
\pgfsetdash{}{0pt}%
\pgfpathmoveto{\pgfqpoint{3.521740in}{0.315944in}}%
\pgfpathlineto{\pgfqpoint{3.485872in}{0.315944in}}%
\pgfpathlineto{\pgfqpoint{3.485872in}{0.336873in}}%
\pgfpathlineto{\pgfqpoint{3.521740in}{0.336873in}}%
\pgfpathlineto{\pgfqpoint{3.521740in}{0.315944in}}%
\pgfpathclose%
\pgfusepath{fill}%
\end{pgfscope}%
\begin{pgfscope}%
\pgfpathrectangle{\pgfqpoint{3.019583in}{0.169444in}}{\pgfqpoint{0.896708in}{1.339426in}}%
\pgfusepath{clip}%
\pgfsetbuttcap%
\pgfsetmiterjoin%
\definecolor{currentfill}{rgb}{0.866282,0.290119,0.297809}%
\pgfsetfillcolor{currentfill}%
\pgfsetlinewidth{0.000000pt}%
\definecolor{currentstroke}{rgb}{0.000000,0.000000,0.000000}%
\pgfsetstrokecolor{currentstroke}%
\pgfsetstrokeopacity{0.000000}%
\pgfsetdash{}{0pt}%
\pgfpathmoveto{\pgfqpoint{3.521740in}{0.336873in}}%
\pgfpathlineto{\pgfqpoint{3.485872in}{0.336873in}}%
\pgfpathlineto{\pgfqpoint{3.485872in}{0.357801in}}%
\pgfpathlineto{\pgfqpoint{3.521740in}{0.357801in}}%
\pgfpathlineto{\pgfqpoint{3.521740in}{0.336873in}}%
\pgfpathclose%
\pgfusepath{fill}%
\end{pgfscope}%
\begin{pgfscope}%
\pgfpathrectangle{\pgfqpoint{3.019583in}{0.169444in}}{\pgfqpoint{0.896708in}{1.339426in}}%
\pgfusepath{clip}%
\pgfsetbuttcap%
\pgfsetmiterjoin%
\definecolor{currentfill}{rgb}{0.885352,0.319031,0.290427}%
\pgfsetfillcolor{currentfill}%
\pgfsetlinewidth{0.000000pt}%
\definecolor{currentstroke}{rgb}{0.000000,0.000000,0.000000}%
\pgfsetstrokecolor{currentstroke}%
\pgfsetstrokeopacity{0.000000}%
\pgfsetdash{}{0pt}%
\pgfpathmoveto{\pgfqpoint{3.521740in}{0.357801in}}%
\pgfpathlineto{\pgfqpoint{3.485872in}{0.357801in}}%
\pgfpathlineto{\pgfqpoint{3.485872in}{0.378730in}}%
\pgfpathlineto{\pgfqpoint{3.521740in}{0.378730in}}%
\pgfpathlineto{\pgfqpoint{3.521740in}{0.357801in}}%
\pgfpathclose%
\pgfusepath{fill}%
\end{pgfscope}%
\begin{pgfscope}%
\pgfpathrectangle{\pgfqpoint{3.019583in}{0.169444in}}{\pgfqpoint{0.896708in}{1.339426in}}%
\pgfusepath{clip}%
\pgfsetbuttcap%
\pgfsetmiterjoin%
\definecolor{currentfill}{rgb}{0.904421,0.347943,0.283045}%
\pgfsetfillcolor{currentfill}%
\pgfsetlinewidth{0.000000pt}%
\definecolor{currentstroke}{rgb}{0.000000,0.000000,0.000000}%
\pgfsetstrokecolor{currentstroke}%
\pgfsetstrokeopacity{0.000000}%
\pgfsetdash{}{0pt}%
\pgfpathmoveto{\pgfqpoint{3.521740in}{0.378730in}}%
\pgfpathlineto{\pgfqpoint{3.485872in}{0.378730in}}%
\pgfpathlineto{\pgfqpoint{3.485872in}{0.399658in}}%
\pgfpathlineto{\pgfqpoint{3.521740in}{0.399658in}}%
\pgfpathlineto{\pgfqpoint{3.521740in}{0.378730in}}%
\pgfpathclose%
\pgfusepath{fill}%
\end{pgfscope}%
\begin{pgfscope}%
\pgfpathrectangle{\pgfqpoint{3.019583in}{0.169444in}}{\pgfqpoint{0.896708in}{1.339426in}}%
\pgfusepath{clip}%
\pgfsetbuttcap%
\pgfsetmiterjoin%
\definecolor{currentfill}{rgb}{0.923491,0.376855,0.275663}%
\pgfsetfillcolor{currentfill}%
\pgfsetlinewidth{0.000000pt}%
\definecolor{currentstroke}{rgb}{0.000000,0.000000,0.000000}%
\pgfsetstrokecolor{currentstroke}%
\pgfsetstrokeopacity{0.000000}%
\pgfsetdash{}{0pt}%
\pgfpathmoveto{\pgfqpoint{3.521740in}{0.399658in}}%
\pgfpathlineto{\pgfqpoint{3.485872in}{0.399658in}}%
\pgfpathlineto{\pgfqpoint{3.485872in}{0.420587in}}%
\pgfpathlineto{\pgfqpoint{3.521740in}{0.420587in}}%
\pgfpathlineto{\pgfqpoint{3.521740in}{0.399658in}}%
\pgfpathclose%
\pgfusepath{fill}%
\end{pgfscope}%
\begin{pgfscope}%
\pgfpathrectangle{\pgfqpoint{3.019583in}{0.169444in}}{\pgfqpoint{0.896708in}{1.339426in}}%
\pgfusepath{clip}%
\pgfsetbuttcap%
\pgfsetmiterjoin%
\definecolor{currentfill}{rgb}{0.942561,0.405767,0.268281}%
\pgfsetfillcolor{currentfill}%
\pgfsetlinewidth{0.000000pt}%
\definecolor{currentstroke}{rgb}{0.000000,0.000000,0.000000}%
\pgfsetstrokecolor{currentstroke}%
\pgfsetstrokeopacity{0.000000}%
\pgfsetdash{}{0pt}%
\pgfpathmoveto{\pgfqpoint{3.521740in}{0.420587in}}%
\pgfpathlineto{\pgfqpoint{3.485872in}{0.420587in}}%
\pgfpathlineto{\pgfqpoint{3.485872in}{0.441515in}}%
\pgfpathlineto{\pgfqpoint{3.521740in}{0.441515in}}%
\pgfpathlineto{\pgfqpoint{3.521740in}{0.420587in}}%
\pgfpathclose%
\pgfusepath{fill}%
\end{pgfscope}%
\begin{pgfscope}%
\pgfpathrectangle{\pgfqpoint{3.019583in}{0.169444in}}{\pgfqpoint{0.896708in}{1.339426in}}%
\pgfusepath{clip}%
\pgfsetbuttcap%
\pgfsetmiterjoin%
\definecolor{currentfill}{rgb}{0.958247,0.437447,0.267359}%
\pgfsetfillcolor{currentfill}%
\pgfsetlinewidth{0.000000pt}%
\definecolor{currentstroke}{rgb}{0.000000,0.000000,0.000000}%
\pgfsetstrokecolor{currentstroke}%
\pgfsetstrokeopacity{0.000000}%
\pgfsetdash{}{0pt}%
\pgfpathmoveto{\pgfqpoint{3.521740in}{0.441515in}}%
\pgfpathlineto{\pgfqpoint{3.485872in}{0.441515in}}%
\pgfpathlineto{\pgfqpoint{3.485872in}{0.462444in}}%
\pgfpathlineto{\pgfqpoint{3.521740in}{0.462444in}}%
\pgfpathlineto{\pgfqpoint{3.521740in}{0.441515in}}%
\pgfpathclose%
\pgfusepath{fill}%
\end{pgfscope}%
\begin{pgfscope}%
\pgfpathrectangle{\pgfqpoint{3.019583in}{0.169444in}}{\pgfqpoint{0.896708in}{1.339426in}}%
\pgfusepath{clip}%
\pgfsetbuttcap%
\pgfsetmiterjoin%
\definecolor{currentfill}{rgb}{0.963783,0.477432,0.285813}%
\pgfsetfillcolor{currentfill}%
\pgfsetlinewidth{0.000000pt}%
\definecolor{currentstroke}{rgb}{0.000000,0.000000,0.000000}%
\pgfsetstrokecolor{currentstroke}%
\pgfsetstrokeopacity{0.000000}%
\pgfsetdash{}{0pt}%
\pgfpathmoveto{\pgfqpoint{3.521740in}{0.462444in}}%
\pgfpathlineto{\pgfqpoint{3.485872in}{0.462444in}}%
\pgfpathlineto{\pgfqpoint{3.485872in}{0.483372in}}%
\pgfpathlineto{\pgfqpoint{3.521740in}{0.483372in}}%
\pgfpathlineto{\pgfqpoint{3.521740in}{0.462444in}}%
\pgfpathclose%
\pgfusepath{fill}%
\end{pgfscope}%
\begin{pgfscope}%
\pgfpathrectangle{\pgfqpoint{3.019583in}{0.169444in}}{\pgfqpoint{0.896708in}{1.339426in}}%
\pgfusepath{clip}%
\pgfsetbuttcap%
\pgfsetmiterjoin%
\definecolor{currentfill}{rgb}{0.969319,0.517416,0.304268}%
\pgfsetfillcolor{currentfill}%
\pgfsetlinewidth{0.000000pt}%
\definecolor{currentstroke}{rgb}{0.000000,0.000000,0.000000}%
\pgfsetstrokecolor{currentstroke}%
\pgfsetstrokeopacity{0.000000}%
\pgfsetdash{}{0pt}%
\pgfpathmoveto{\pgfqpoint{3.521740in}{0.483372in}}%
\pgfpathlineto{\pgfqpoint{3.485872in}{0.483372in}}%
\pgfpathlineto{\pgfqpoint{3.485872in}{0.504301in}}%
\pgfpathlineto{\pgfqpoint{3.521740in}{0.504301in}}%
\pgfpathlineto{\pgfqpoint{3.521740in}{0.483372in}}%
\pgfpathclose%
\pgfusepath{fill}%
\end{pgfscope}%
\begin{pgfscope}%
\pgfpathrectangle{\pgfqpoint{3.019583in}{0.169444in}}{\pgfqpoint{0.896708in}{1.339426in}}%
\pgfusepath{clip}%
\pgfsetbuttcap%
\pgfsetmiterjoin%
\definecolor{currentfill}{rgb}{0.974856,0.557401,0.322722}%
\pgfsetfillcolor{currentfill}%
\pgfsetlinewidth{0.000000pt}%
\definecolor{currentstroke}{rgb}{0.000000,0.000000,0.000000}%
\pgfsetstrokecolor{currentstroke}%
\pgfsetstrokeopacity{0.000000}%
\pgfsetdash{}{0pt}%
\pgfpathmoveto{\pgfqpoint{3.521740in}{0.504301in}}%
\pgfpathlineto{\pgfqpoint{3.485872in}{0.504301in}}%
\pgfpathlineto{\pgfqpoint{3.485872in}{0.525230in}}%
\pgfpathlineto{\pgfqpoint{3.521740in}{0.525230in}}%
\pgfpathlineto{\pgfqpoint{3.521740in}{0.504301in}}%
\pgfpathclose%
\pgfusepath{fill}%
\end{pgfscope}%
\begin{pgfscope}%
\pgfpathrectangle{\pgfqpoint{3.019583in}{0.169444in}}{\pgfqpoint{0.896708in}{1.339426in}}%
\pgfusepath{clip}%
\pgfsetbuttcap%
\pgfsetmiterjoin%
\definecolor{currentfill}{rgb}{0.980392,0.597386,0.341176}%
\pgfsetfillcolor{currentfill}%
\pgfsetlinewidth{0.000000pt}%
\definecolor{currentstroke}{rgb}{0.000000,0.000000,0.000000}%
\pgfsetstrokecolor{currentstroke}%
\pgfsetstrokeopacity{0.000000}%
\pgfsetdash{}{0pt}%
\pgfpathmoveto{\pgfqpoint{3.521740in}{0.525230in}}%
\pgfpathlineto{\pgfqpoint{3.485872in}{0.525230in}}%
\pgfpathlineto{\pgfqpoint{3.485872in}{0.546158in}}%
\pgfpathlineto{\pgfqpoint{3.521740in}{0.546158in}}%
\pgfpathlineto{\pgfqpoint{3.521740in}{0.525230in}}%
\pgfpathclose%
\pgfusepath{fill}%
\end{pgfscope}%
\begin{pgfscope}%
\pgfpathrectangle{\pgfqpoint{3.019583in}{0.169444in}}{\pgfqpoint{0.896708in}{1.339426in}}%
\pgfusepath{clip}%
\pgfsetbuttcap%
\pgfsetmiterjoin%
\definecolor{currentfill}{rgb}{0.985928,0.637370,0.359631}%
\pgfsetfillcolor{currentfill}%
\pgfsetlinewidth{0.000000pt}%
\definecolor{currentstroke}{rgb}{0.000000,0.000000,0.000000}%
\pgfsetstrokecolor{currentstroke}%
\pgfsetstrokeopacity{0.000000}%
\pgfsetdash{}{0pt}%
\pgfpathmoveto{\pgfqpoint{3.521740in}{0.546158in}}%
\pgfpathlineto{\pgfqpoint{3.485872in}{0.546158in}}%
\pgfpathlineto{\pgfqpoint{3.485872in}{0.567087in}}%
\pgfpathlineto{\pgfqpoint{3.521740in}{0.567087in}}%
\pgfpathlineto{\pgfqpoint{3.521740in}{0.546158in}}%
\pgfpathclose%
\pgfusepath{fill}%
\end{pgfscope}%
\begin{pgfscope}%
\pgfpathrectangle{\pgfqpoint{3.019583in}{0.169444in}}{\pgfqpoint{0.896708in}{1.339426in}}%
\pgfusepath{clip}%
\pgfsetbuttcap%
\pgfsetmiterjoin%
\definecolor{currentfill}{rgb}{0.991465,0.677355,0.378085}%
\pgfsetfillcolor{currentfill}%
\pgfsetlinewidth{0.000000pt}%
\definecolor{currentstroke}{rgb}{0.000000,0.000000,0.000000}%
\pgfsetstrokecolor{currentstroke}%
\pgfsetstrokeopacity{0.000000}%
\pgfsetdash{}{0pt}%
\pgfpathmoveto{\pgfqpoint{3.521740in}{0.567087in}}%
\pgfpathlineto{\pgfqpoint{3.485872in}{0.567087in}}%
\pgfpathlineto{\pgfqpoint{3.485872in}{0.588015in}}%
\pgfpathlineto{\pgfqpoint{3.521740in}{0.588015in}}%
\pgfpathlineto{\pgfqpoint{3.521740in}{0.567087in}}%
\pgfpathclose%
\pgfusepath{fill}%
\end{pgfscope}%
\begin{pgfscope}%
\pgfpathrectangle{\pgfqpoint{3.019583in}{0.169444in}}{\pgfqpoint{0.896708in}{1.339426in}}%
\pgfusepath{clip}%
\pgfsetbuttcap%
\pgfsetmiterjoin%
\definecolor{currentfill}{rgb}{0.992695,0.709266,0.402999}%
\pgfsetfillcolor{currentfill}%
\pgfsetlinewidth{0.000000pt}%
\definecolor{currentstroke}{rgb}{0.000000,0.000000,0.000000}%
\pgfsetstrokecolor{currentstroke}%
\pgfsetstrokeopacity{0.000000}%
\pgfsetdash{}{0pt}%
\pgfpathmoveto{\pgfqpoint{3.521740in}{0.588015in}}%
\pgfpathlineto{\pgfqpoint{3.485872in}{0.588015in}}%
\pgfpathlineto{\pgfqpoint{3.485872in}{0.608944in}}%
\pgfpathlineto{\pgfqpoint{3.521740in}{0.608944in}}%
\pgfpathlineto{\pgfqpoint{3.521740in}{0.588015in}}%
\pgfpathclose%
\pgfusepath{fill}%
\end{pgfscope}%
\begin{pgfscope}%
\pgfpathrectangle{\pgfqpoint{3.019583in}{0.169444in}}{\pgfqpoint{0.896708in}{1.339426in}}%
\pgfusepath{clip}%
\pgfsetbuttcap%
\pgfsetmiterjoin%
\definecolor{currentfill}{rgb}{0.993310,0.740023,0.428835}%
\pgfsetfillcolor{currentfill}%
\pgfsetlinewidth{0.000000pt}%
\definecolor{currentstroke}{rgb}{0.000000,0.000000,0.000000}%
\pgfsetstrokecolor{currentstroke}%
\pgfsetstrokeopacity{0.000000}%
\pgfsetdash{}{0pt}%
\pgfpathmoveto{\pgfqpoint{3.521740in}{0.608944in}}%
\pgfpathlineto{\pgfqpoint{3.485872in}{0.608944in}}%
\pgfpathlineto{\pgfqpoint{3.485872in}{0.629872in}}%
\pgfpathlineto{\pgfqpoint{3.521740in}{0.629872in}}%
\pgfpathlineto{\pgfqpoint{3.521740in}{0.608944in}}%
\pgfpathclose%
\pgfusepath{fill}%
\end{pgfscope}%
\begin{pgfscope}%
\pgfpathrectangle{\pgfqpoint{3.019583in}{0.169444in}}{\pgfqpoint{0.896708in}{1.339426in}}%
\pgfusepath{clip}%
\pgfsetbuttcap%
\pgfsetmiterjoin%
\definecolor{currentfill}{rgb}{0.993925,0.770780,0.454671}%
\pgfsetfillcolor{currentfill}%
\pgfsetlinewidth{0.000000pt}%
\definecolor{currentstroke}{rgb}{0.000000,0.000000,0.000000}%
\pgfsetstrokecolor{currentstroke}%
\pgfsetstrokeopacity{0.000000}%
\pgfsetdash{}{0pt}%
\pgfpathmoveto{\pgfqpoint{3.521740in}{0.629872in}}%
\pgfpathlineto{\pgfqpoint{3.485872in}{0.629872in}}%
\pgfpathlineto{\pgfqpoint{3.485872in}{0.650801in}}%
\pgfpathlineto{\pgfqpoint{3.521740in}{0.650801in}}%
\pgfpathlineto{\pgfqpoint{3.521740in}{0.629872in}}%
\pgfpathclose%
\pgfusepath{fill}%
\end{pgfscope}%
\begin{pgfscope}%
\pgfpathrectangle{\pgfqpoint{3.019583in}{0.169444in}}{\pgfqpoint{0.896708in}{1.339426in}}%
\pgfusepath{clip}%
\pgfsetbuttcap%
\pgfsetmiterjoin%
\definecolor{currentfill}{rgb}{0.994541,0.801538,0.480507}%
\pgfsetfillcolor{currentfill}%
\pgfsetlinewidth{0.000000pt}%
\definecolor{currentstroke}{rgb}{0.000000,0.000000,0.000000}%
\pgfsetstrokecolor{currentstroke}%
\pgfsetstrokeopacity{0.000000}%
\pgfsetdash{}{0pt}%
\pgfpathmoveto{\pgfqpoint{3.521740in}{0.650801in}}%
\pgfpathlineto{\pgfqpoint{3.485872in}{0.650801in}}%
\pgfpathlineto{\pgfqpoint{3.485872in}{0.671729in}}%
\pgfpathlineto{\pgfqpoint{3.521740in}{0.671729in}}%
\pgfpathlineto{\pgfqpoint{3.521740in}{0.650801in}}%
\pgfpathclose%
\pgfusepath{fill}%
\end{pgfscope}%
\begin{pgfscope}%
\pgfpathrectangle{\pgfqpoint{3.019583in}{0.169444in}}{\pgfqpoint{0.896708in}{1.339426in}}%
\pgfusepath{clip}%
\pgfsetbuttcap%
\pgfsetmiterjoin%
\definecolor{currentfill}{rgb}{0.995156,0.832295,0.506344}%
\pgfsetfillcolor{currentfill}%
\pgfsetlinewidth{0.000000pt}%
\definecolor{currentstroke}{rgb}{0.000000,0.000000,0.000000}%
\pgfsetstrokecolor{currentstroke}%
\pgfsetstrokeopacity{0.000000}%
\pgfsetdash{}{0pt}%
\pgfpathmoveto{\pgfqpoint{3.521740in}{0.671729in}}%
\pgfpathlineto{\pgfqpoint{3.485872in}{0.671729in}}%
\pgfpathlineto{\pgfqpoint{3.485872in}{0.692658in}}%
\pgfpathlineto{\pgfqpoint{3.521740in}{0.692658in}}%
\pgfpathlineto{\pgfqpoint{3.521740in}{0.671729in}}%
\pgfpathclose%
\pgfusepath{fill}%
\end{pgfscope}%
\begin{pgfscope}%
\pgfpathrectangle{\pgfqpoint{3.019583in}{0.169444in}}{\pgfqpoint{0.896708in}{1.339426in}}%
\pgfusepath{clip}%
\pgfsetbuttcap%
\pgfsetmiterjoin%
\definecolor{currentfill}{rgb}{0.995771,0.863053,0.532180}%
\pgfsetfillcolor{currentfill}%
\pgfsetlinewidth{0.000000pt}%
\definecolor{currentstroke}{rgb}{0.000000,0.000000,0.000000}%
\pgfsetstrokecolor{currentstroke}%
\pgfsetstrokeopacity{0.000000}%
\pgfsetdash{}{0pt}%
\pgfpathmoveto{\pgfqpoint{3.521740in}{0.692658in}}%
\pgfpathlineto{\pgfqpoint{3.485872in}{0.692658in}}%
\pgfpathlineto{\pgfqpoint{3.485872in}{0.713586in}}%
\pgfpathlineto{\pgfqpoint{3.521740in}{0.713586in}}%
\pgfpathlineto{\pgfqpoint{3.521740in}{0.692658in}}%
\pgfpathclose%
\pgfusepath{fill}%
\end{pgfscope}%
\begin{pgfscope}%
\pgfpathrectangle{\pgfqpoint{3.019583in}{0.169444in}}{\pgfqpoint{0.896708in}{1.339426in}}%
\pgfusepath{clip}%
\pgfsetbuttcap%
\pgfsetmiterjoin%
\definecolor{currentfill}{rgb}{0.996386,0.887966,0.561092}%
\pgfsetfillcolor{currentfill}%
\pgfsetlinewidth{0.000000pt}%
\definecolor{currentstroke}{rgb}{0.000000,0.000000,0.000000}%
\pgfsetstrokecolor{currentstroke}%
\pgfsetstrokeopacity{0.000000}%
\pgfsetdash{}{0pt}%
\pgfpathmoveto{\pgfqpoint{3.521740in}{0.713586in}}%
\pgfpathlineto{\pgfqpoint{3.485872in}{0.713586in}}%
\pgfpathlineto{\pgfqpoint{3.485872in}{0.734515in}}%
\pgfpathlineto{\pgfqpoint{3.521740in}{0.734515in}}%
\pgfpathlineto{\pgfqpoint{3.521740in}{0.713586in}}%
\pgfpathclose%
\pgfusepath{fill}%
\end{pgfscope}%
\begin{pgfscope}%
\pgfpathrectangle{\pgfqpoint{3.019583in}{0.169444in}}{\pgfqpoint{0.896708in}{1.339426in}}%
\pgfusepath{clip}%
\pgfsetbuttcap%
\pgfsetmiterjoin%
\definecolor{currentfill}{rgb}{0.997001,0.907036,0.593080}%
\pgfsetfillcolor{currentfill}%
\pgfsetlinewidth{0.000000pt}%
\definecolor{currentstroke}{rgb}{0.000000,0.000000,0.000000}%
\pgfsetstrokecolor{currentstroke}%
\pgfsetstrokeopacity{0.000000}%
\pgfsetdash{}{0pt}%
\pgfpathmoveto{\pgfqpoint{3.521740in}{0.734515in}}%
\pgfpathlineto{\pgfqpoint{3.485872in}{0.734515in}}%
\pgfpathlineto{\pgfqpoint{3.485872in}{0.755443in}}%
\pgfpathlineto{\pgfqpoint{3.521740in}{0.755443in}}%
\pgfpathlineto{\pgfqpoint{3.521740in}{0.734515in}}%
\pgfpathclose%
\pgfusepath{fill}%
\end{pgfscope}%
\begin{pgfscope}%
\pgfpathrectangle{\pgfqpoint{3.019583in}{0.169444in}}{\pgfqpoint{0.896708in}{1.339426in}}%
\pgfusepath{clip}%
\pgfsetbuttcap%
\pgfsetmiterjoin%
\definecolor{currentfill}{rgb}{0.997616,0.926105,0.625067}%
\pgfsetfillcolor{currentfill}%
\pgfsetlinewidth{0.000000pt}%
\definecolor{currentstroke}{rgb}{0.000000,0.000000,0.000000}%
\pgfsetstrokecolor{currentstroke}%
\pgfsetstrokeopacity{0.000000}%
\pgfsetdash{}{0pt}%
\pgfpathmoveto{\pgfqpoint{3.521740in}{0.755443in}}%
\pgfpathlineto{\pgfqpoint{3.485872in}{0.755443in}}%
\pgfpathlineto{\pgfqpoint{3.485872in}{0.776372in}}%
\pgfpathlineto{\pgfqpoint{3.521740in}{0.776372in}}%
\pgfpathlineto{\pgfqpoint{3.521740in}{0.755443in}}%
\pgfpathclose%
\pgfusepath{fill}%
\end{pgfscope}%
\begin{pgfscope}%
\pgfpathrectangle{\pgfqpoint{3.019583in}{0.169444in}}{\pgfqpoint{0.896708in}{1.339426in}}%
\pgfusepath{clip}%
\pgfsetbuttcap%
\pgfsetmiterjoin%
\definecolor{currentfill}{rgb}{0.998231,0.945175,0.657055}%
\pgfsetfillcolor{currentfill}%
\pgfsetlinewidth{0.000000pt}%
\definecolor{currentstroke}{rgb}{0.000000,0.000000,0.000000}%
\pgfsetstrokecolor{currentstroke}%
\pgfsetstrokeopacity{0.000000}%
\pgfsetdash{}{0pt}%
\pgfpathmoveto{\pgfqpoint{3.521740in}{0.776372in}}%
\pgfpathlineto{\pgfqpoint{3.485872in}{0.776372in}}%
\pgfpathlineto{\pgfqpoint{3.485872in}{0.797300in}}%
\pgfpathlineto{\pgfqpoint{3.521740in}{0.797300in}}%
\pgfpathlineto{\pgfqpoint{3.521740in}{0.776372in}}%
\pgfpathclose%
\pgfusepath{fill}%
\end{pgfscope}%
\begin{pgfscope}%
\pgfpathrectangle{\pgfqpoint{3.019583in}{0.169444in}}{\pgfqpoint{0.896708in}{1.339426in}}%
\pgfusepath{clip}%
\pgfsetbuttcap%
\pgfsetmiterjoin%
\definecolor{currentfill}{rgb}{0.998847,0.964245,0.689043}%
\pgfsetfillcolor{currentfill}%
\pgfsetlinewidth{0.000000pt}%
\definecolor{currentstroke}{rgb}{0.000000,0.000000,0.000000}%
\pgfsetstrokecolor{currentstroke}%
\pgfsetstrokeopacity{0.000000}%
\pgfsetdash{}{0pt}%
\pgfpathmoveto{\pgfqpoint{3.521740in}{0.797300in}}%
\pgfpathlineto{\pgfqpoint{3.485872in}{0.797300in}}%
\pgfpathlineto{\pgfqpoint{3.485872in}{0.818229in}}%
\pgfpathlineto{\pgfqpoint{3.521740in}{0.818229in}}%
\pgfpathlineto{\pgfqpoint{3.521740in}{0.797300in}}%
\pgfpathclose%
\pgfusepath{fill}%
\end{pgfscope}%
\begin{pgfscope}%
\pgfpathrectangle{\pgfqpoint{3.019583in}{0.169444in}}{\pgfqpoint{0.896708in}{1.339426in}}%
\pgfusepath{clip}%
\pgfsetbuttcap%
\pgfsetmiterjoin%
\definecolor{currentfill}{rgb}{0.999462,0.983314,0.721030}%
\pgfsetfillcolor{currentfill}%
\pgfsetlinewidth{0.000000pt}%
\definecolor{currentstroke}{rgb}{0.000000,0.000000,0.000000}%
\pgfsetstrokecolor{currentstroke}%
\pgfsetstrokeopacity{0.000000}%
\pgfsetdash{}{0pt}%
\pgfpathmoveto{\pgfqpoint{3.521740in}{0.818229in}}%
\pgfpathlineto{\pgfqpoint{3.485872in}{0.818229in}}%
\pgfpathlineto{\pgfqpoint{3.485872in}{0.839157in}}%
\pgfpathlineto{\pgfqpoint{3.521740in}{0.839157in}}%
\pgfpathlineto{\pgfqpoint{3.521740in}{0.818229in}}%
\pgfpathclose%
\pgfusepath{fill}%
\end{pgfscope}%
\begin{pgfscope}%
\pgfpathrectangle{\pgfqpoint{3.019583in}{0.169444in}}{\pgfqpoint{0.896708in}{1.339426in}}%
\pgfusepath{clip}%
\pgfsetbuttcap%
\pgfsetmiterjoin%
\definecolor{currentfill}{rgb}{0.998078,0.999231,0.746021}%
\pgfsetfillcolor{currentfill}%
\pgfsetlinewidth{0.000000pt}%
\definecolor{currentstroke}{rgb}{0.000000,0.000000,0.000000}%
\pgfsetstrokecolor{currentstroke}%
\pgfsetstrokeopacity{0.000000}%
\pgfsetdash{}{0pt}%
\pgfpathmoveto{\pgfqpoint{3.521740in}{0.839157in}}%
\pgfpathlineto{\pgfqpoint{3.485872in}{0.839157in}}%
\pgfpathlineto{\pgfqpoint{3.485872in}{0.860086in}}%
\pgfpathlineto{\pgfqpoint{3.521740in}{0.860086in}}%
\pgfpathlineto{\pgfqpoint{3.521740in}{0.839157in}}%
\pgfpathclose%
\pgfusepath{fill}%
\end{pgfscope}%
\begin{pgfscope}%
\pgfpathrectangle{\pgfqpoint{3.019583in}{0.169444in}}{\pgfqpoint{0.896708in}{1.339426in}}%
\pgfusepath{clip}%
\pgfsetbuttcap%
\pgfsetmiterjoin%
\definecolor{currentfill}{rgb}{0.982699,0.993080,0.722030}%
\pgfsetfillcolor{currentfill}%
\pgfsetlinewidth{0.000000pt}%
\definecolor{currentstroke}{rgb}{0.000000,0.000000,0.000000}%
\pgfsetstrokecolor{currentstroke}%
\pgfsetstrokeopacity{0.000000}%
\pgfsetdash{}{0pt}%
\pgfpathmoveto{\pgfqpoint{3.521740in}{0.860086in}}%
\pgfpathlineto{\pgfqpoint{3.485872in}{0.860086in}}%
\pgfpathlineto{\pgfqpoint{3.485872in}{0.881015in}}%
\pgfpathlineto{\pgfqpoint{3.521740in}{0.881015in}}%
\pgfpathlineto{\pgfqpoint{3.521740in}{0.860086in}}%
\pgfpathclose%
\pgfusepath{fill}%
\end{pgfscope}%
\begin{pgfscope}%
\pgfpathrectangle{\pgfqpoint{3.019583in}{0.169444in}}{\pgfqpoint{0.896708in}{1.339426in}}%
\pgfusepath{clip}%
\pgfsetbuttcap%
\pgfsetmiterjoin%
\definecolor{currentfill}{rgb}{0.967320,0.986928,0.698039}%
\pgfsetfillcolor{currentfill}%
\pgfsetlinewidth{0.000000pt}%
\definecolor{currentstroke}{rgb}{0.000000,0.000000,0.000000}%
\pgfsetstrokecolor{currentstroke}%
\pgfsetstrokeopacity{0.000000}%
\pgfsetdash{}{0pt}%
\pgfpathmoveto{\pgfqpoint{3.521740in}{0.881015in}}%
\pgfpathlineto{\pgfqpoint{3.485872in}{0.881015in}}%
\pgfpathlineto{\pgfqpoint{3.485872in}{0.901943in}}%
\pgfpathlineto{\pgfqpoint{3.521740in}{0.901943in}}%
\pgfpathlineto{\pgfqpoint{3.521740in}{0.881015in}}%
\pgfpathclose%
\pgfusepath{fill}%
\end{pgfscope}%
\begin{pgfscope}%
\pgfpathrectangle{\pgfqpoint{3.019583in}{0.169444in}}{\pgfqpoint{0.896708in}{1.339426in}}%
\pgfusepath{clip}%
\pgfsetbuttcap%
\pgfsetmiterjoin%
\definecolor{currentfill}{rgb}{0.951942,0.980777,0.674048}%
\pgfsetfillcolor{currentfill}%
\pgfsetlinewidth{0.000000pt}%
\definecolor{currentstroke}{rgb}{0.000000,0.000000,0.000000}%
\pgfsetstrokecolor{currentstroke}%
\pgfsetstrokeopacity{0.000000}%
\pgfsetdash{}{0pt}%
\pgfpathmoveto{\pgfqpoint{3.521740in}{0.901943in}}%
\pgfpathlineto{\pgfqpoint{3.485872in}{0.901943in}}%
\pgfpathlineto{\pgfqpoint{3.485872in}{0.922872in}}%
\pgfpathlineto{\pgfqpoint{3.521740in}{0.922872in}}%
\pgfpathlineto{\pgfqpoint{3.521740in}{0.901943in}}%
\pgfpathclose%
\pgfusepath{fill}%
\end{pgfscope}%
\begin{pgfscope}%
\pgfpathrectangle{\pgfqpoint{3.019583in}{0.169444in}}{\pgfqpoint{0.896708in}{1.339426in}}%
\pgfusepath{clip}%
\pgfsetbuttcap%
\pgfsetmiterjoin%
\definecolor{currentfill}{rgb}{0.936563,0.974625,0.650058}%
\pgfsetfillcolor{currentfill}%
\pgfsetlinewidth{0.000000pt}%
\definecolor{currentstroke}{rgb}{0.000000,0.000000,0.000000}%
\pgfsetstrokecolor{currentstroke}%
\pgfsetstrokeopacity{0.000000}%
\pgfsetdash{}{0pt}%
\pgfpathmoveto{\pgfqpoint{3.521740in}{0.922872in}}%
\pgfpathlineto{\pgfqpoint{3.485872in}{0.922872in}}%
\pgfpathlineto{\pgfqpoint{3.485872in}{0.943800in}}%
\pgfpathlineto{\pgfqpoint{3.521740in}{0.943800in}}%
\pgfpathlineto{\pgfqpoint{3.521740in}{0.922872in}}%
\pgfpathclose%
\pgfusepath{fill}%
\end{pgfscope}%
\begin{pgfscope}%
\pgfpathrectangle{\pgfqpoint{3.019583in}{0.169444in}}{\pgfqpoint{0.896708in}{1.339426in}}%
\pgfusepath{clip}%
\pgfsetbuttcap%
\pgfsetmiterjoin%
\definecolor{currentfill}{rgb}{0.921184,0.968474,0.626067}%
\pgfsetfillcolor{currentfill}%
\pgfsetlinewidth{0.000000pt}%
\definecolor{currentstroke}{rgb}{0.000000,0.000000,0.000000}%
\pgfsetstrokecolor{currentstroke}%
\pgfsetstrokeopacity{0.000000}%
\pgfsetdash{}{0pt}%
\pgfpathmoveto{\pgfqpoint{3.521740in}{0.943800in}}%
\pgfpathlineto{\pgfqpoint{3.485872in}{0.943800in}}%
\pgfpathlineto{\pgfqpoint{3.485872in}{0.964729in}}%
\pgfpathlineto{\pgfqpoint{3.521740in}{0.964729in}}%
\pgfpathlineto{\pgfqpoint{3.521740in}{0.943800in}}%
\pgfpathclose%
\pgfusepath{fill}%
\end{pgfscope}%
\begin{pgfscope}%
\pgfpathrectangle{\pgfqpoint{3.019583in}{0.169444in}}{\pgfqpoint{0.896708in}{1.339426in}}%
\pgfusepath{clip}%
\pgfsetbuttcap%
\pgfsetmiterjoin%
\definecolor{currentfill}{rgb}{0.905805,0.962322,0.602076}%
\pgfsetfillcolor{currentfill}%
\pgfsetlinewidth{0.000000pt}%
\definecolor{currentstroke}{rgb}{0.000000,0.000000,0.000000}%
\pgfsetstrokecolor{currentstroke}%
\pgfsetstrokeopacity{0.000000}%
\pgfsetdash{}{0pt}%
\pgfpathmoveto{\pgfqpoint{3.521740in}{0.964729in}}%
\pgfpathlineto{\pgfqpoint{3.485872in}{0.964729in}}%
\pgfpathlineto{\pgfqpoint{3.485872in}{0.985657in}}%
\pgfpathlineto{\pgfqpoint{3.521740in}{0.985657in}}%
\pgfpathlineto{\pgfqpoint{3.521740in}{0.964729in}}%
\pgfpathclose%
\pgfusepath{fill}%
\end{pgfscope}%
\begin{pgfscope}%
\pgfpathrectangle{\pgfqpoint{3.019583in}{0.169444in}}{\pgfqpoint{0.896708in}{1.339426in}}%
\pgfusepath{clip}%
\pgfsetbuttcap%
\pgfsetmiterjoin%
\definecolor{currentfill}{rgb}{0.874740,0.949712,0.601615}%
\pgfsetfillcolor{currentfill}%
\pgfsetlinewidth{0.000000pt}%
\definecolor{currentstroke}{rgb}{0.000000,0.000000,0.000000}%
\pgfsetstrokecolor{currentstroke}%
\pgfsetstrokeopacity{0.000000}%
\pgfsetdash{}{0pt}%
\pgfpathmoveto{\pgfqpoint{3.521740in}{0.985657in}}%
\pgfpathlineto{\pgfqpoint{3.485872in}{0.985657in}}%
\pgfpathlineto{\pgfqpoint{3.485872in}{1.006586in}}%
\pgfpathlineto{\pgfqpoint{3.521740in}{1.006586in}}%
\pgfpathlineto{\pgfqpoint{3.521740in}{0.985657in}}%
\pgfpathclose%
\pgfusepath{fill}%
\end{pgfscope}%
\begin{pgfscope}%
\pgfpathrectangle{\pgfqpoint{3.019583in}{0.169444in}}{\pgfqpoint{0.896708in}{1.339426in}}%
\pgfusepath{clip}%
\pgfsetbuttcap%
\pgfsetmiterjoin%
\definecolor{currentfill}{rgb}{0.838447,0.934948,0.608997}%
\pgfsetfillcolor{currentfill}%
\pgfsetlinewidth{0.000000pt}%
\definecolor{currentstroke}{rgb}{0.000000,0.000000,0.000000}%
\pgfsetstrokecolor{currentstroke}%
\pgfsetstrokeopacity{0.000000}%
\pgfsetdash{}{0pt}%
\pgfpathmoveto{\pgfqpoint{3.521740in}{1.006586in}}%
\pgfpathlineto{\pgfqpoint{3.485872in}{1.006586in}}%
\pgfpathlineto{\pgfqpoint{3.485872in}{1.027514in}}%
\pgfpathlineto{\pgfqpoint{3.521740in}{1.027514in}}%
\pgfpathlineto{\pgfqpoint{3.521740in}{1.006586in}}%
\pgfpathclose%
\pgfusepath{fill}%
\end{pgfscope}%
\begin{pgfscope}%
\pgfpathrectangle{\pgfqpoint{3.019583in}{0.169444in}}{\pgfqpoint{0.896708in}{1.339426in}}%
\pgfusepath{clip}%
\pgfsetbuttcap%
\pgfsetmiterjoin%
\definecolor{currentfill}{rgb}{0.802153,0.920185,0.616378}%
\pgfsetfillcolor{currentfill}%
\pgfsetlinewidth{0.000000pt}%
\definecolor{currentstroke}{rgb}{0.000000,0.000000,0.000000}%
\pgfsetstrokecolor{currentstroke}%
\pgfsetstrokeopacity{0.000000}%
\pgfsetdash{}{0pt}%
\pgfpathmoveto{\pgfqpoint{3.521740in}{1.027514in}}%
\pgfpathlineto{\pgfqpoint{3.485872in}{1.027514in}}%
\pgfpathlineto{\pgfqpoint{3.485872in}{1.048443in}}%
\pgfpathlineto{\pgfqpoint{3.521740in}{1.048443in}}%
\pgfpathlineto{\pgfqpoint{3.521740in}{1.027514in}}%
\pgfpathclose%
\pgfusepath{fill}%
\end{pgfscope}%
\begin{pgfscope}%
\pgfpathrectangle{\pgfqpoint{3.019583in}{0.169444in}}{\pgfqpoint{0.896708in}{1.339426in}}%
\pgfusepath{clip}%
\pgfsetbuttcap%
\pgfsetmiterjoin%
\definecolor{currentfill}{rgb}{0.765859,0.905421,0.623760}%
\pgfsetfillcolor{currentfill}%
\pgfsetlinewidth{0.000000pt}%
\definecolor{currentstroke}{rgb}{0.000000,0.000000,0.000000}%
\pgfsetstrokecolor{currentstroke}%
\pgfsetstrokeopacity{0.000000}%
\pgfsetdash{}{0pt}%
\pgfpathmoveto{\pgfqpoint{3.521740in}{1.048443in}}%
\pgfpathlineto{\pgfqpoint{3.485872in}{1.048443in}}%
\pgfpathlineto{\pgfqpoint{3.485872in}{1.069371in}}%
\pgfpathlineto{\pgfqpoint{3.521740in}{1.069371in}}%
\pgfpathlineto{\pgfqpoint{3.521740in}{1.048443in}}%
\pgfpathclose%
\pgfusepath{fill}%
\end{pgfscope}%
\begin{pgfscope}%
\pgfpathrectangle{\pgfqpoint{3.019583in}{0.169444in}}{\pgfqpoint{0.896708in}{1.339426in}}%
\pgfusepath{clip}%
\pgfsetbuttcap%
\pgfsetmiterjoin%
\definecolor{currentfill}{rgb}{0.729566,0.890657,0.631142}%
\pgfsetfillcolor{currentfill}%
\pgfsetlinewidth{0.000000pt}%
\definecolor{currentstroke}{rgb}{0.000000,0.000000,0.000000}%
\pgfsetstrokecolor{currentstroke}%
\pgfsetstrokeopacity{0.000000}%
\pgfsetdash{}{0pt}%
\pgfpathmoveto{\pgfqpoint{3.521740in}{1.069371in}}%
\pgfpathlineto{\pgfqpoint{3.485872in}{1.069371in}}%
\pgfpathlineto{\pgfqpoint{3.485872in}{1.090300in}}%
\pgfpathlineto{\pgfqpoint{3.521740in}{1.090300in}}%
\pgfpathlineto{\pgfqpoint{3.521740in}{1.069371in}}%
\pgfpathclose%
\pgfusepath{fill}%
\end{pgfscope}%
\begin{pgfscope}%
\pgfpathrectangle{\pgfqpoint{3.019583in}{0.169444in}}{\pgfqpoint{0.896708in}{1.339426in}}%
\pgfusepath{clip}%
\pgfsetbuttcap%
\pgfsetmiterjoin%
\definecolor{currentfill}{rgb}{0.693272,0.875894,0.638524}%
\pgfsetfillcolor{currentfill}%
\pgfsetlinewidth{0.000000pt}%
\definecolor{currentstroke}{rgb}{0.000000,0.000000,0.000000}%
\pgfsetstrokecolor{currentstroke}%
\pgfsetstrokeopacity{0.000000}%
\pgfsetdash{}{0pt}%
\pgfpathmoveto{\pgfqpoint{3.521740in}{1.090300in}}%
\pgfpathlineto{\pgfqpoint{3.485872in}{1.090300in}}%
\pgfpathlineto{\pgfqpoint{3.485872in}{1.111228in}}%
\pgfpathlineto{\pgfqpoint{3.521740in}{1.111228in}}%
\pgfpathlineto{\pgfqpoint{3.521740in}{1.090300in}}%
\pgfpathclose%
\pgfusepath{fill}%
\end{pgfscope}%
\begin{pgfscope}%
\pgfpathrectangle{\pgfqpoint{3.019583in}{0.169444in}}{\pgfqpoint{0.896708in}{1.339426in}}%
\pgfusepath{clip}%
\pgfsetbuttcap%
\pgfsetmiterjoin%
\definecolor{currentfill}{rgb}{0.654671,0.860438,0.643368}%
\pgfsetfillcolor{currentfill}%
\pgfsetlinewidth{0.000000pt}%
\definecolor{currentstroke}{rgb}{0.000000,0.000000,0.000000}%
\pgfsetstrokecolor{currentstroke}%
\pgfsetstrokeopacity{0.000000}%
\pgfsetdash{}{0pt}%
\pgfpathmoveto{\pgfqpoint{3.521740in}{1.111228in}}%
\pgfpathlineto{\pgfqpoint{3.485872in}{1.111228in}}%
\pgfpathlineto{\pgfqpoint{3.485872in}{1.132157in}}%
\pgfpathlineto{\pgfqpoint{3.521740in}{1.132157in}}%
\pgfpathlineto{\pgfqpoint{3.521740in}{1.111228in}}%
\pgfpathclose%
\pgfusepath{fill}%
\end{pgfscope}%
\begin{pgfscope}%
\pgfpathrectangle{\pgfqpoint{3.019583in}{0.169444in}}{\pgfqpoint{0.896708in}{1.339426in}}%
\pgfusepath{clip}%
\pgfsetbuttcap%
\pgfsetmiterjoin%
\definecolor{currentfill}{rgb}{0.612226,0.843829,0.643983}%
\pgfsetfillcolor{currentfill}%
\pgfsetlinewidth{0.000000pt}%
\definecolor{currentstroke}{rgb}{0.000000,0.000000,0.000000}%
\pgfsetstrokecolor{currentstroke}%
\pgfsetstrokeopacity{0.000000}%
\pgfsetdash{}{0pt}%
\pgfpathmoveto{\pgfqpoint{3.521740in}{1.132157in}}%
\pgfpathlineto{\pgfqpoint{3.485872in}{1.132157in}}%
\pgfpathlineto{\pgfqpoint{3.485872in}{1.153085in}}%
\pgfpathlineto{\pgfqpoint{3.521740in}{1.153085in}}%
\pgfpathlineto{\pgfqpoint{3.521740in}{1.132157in}}%
\pgfpathclose%
\pgfusepath{fill}%
\end{pgfscope}%
\begin{pgfscope}%
\pgfpathrectangle{\pgfqpoint{3.019583in}{0.169444in}}{\pgfqpoint{0.896708in}{1.339426in}}%
\pgfusepath{clip}%
\pgfsetbuttcap%
\pgfsetmiterjoin%
\definecolor{currentfill}{rgb}{0.569781,0.827220,0.644598}%
\pgfsetfillcolor{currentfill}%
\pgfsetlinewidth{0.000000pt}%
\definecolor{currentstroke}{rgb}{0.000000,0.000000,0.000000}%
\pgfsetstrokecolor{currentstroke}%
\pgfsetstrokeopacity{0.000000}%
\pgfsetdash{}{0pt}%
\pgfpathmoveto{\pgfqpoint{3.521740in}{1.153085in}}%
\pgfpathlineto{\pgfqpoint{3.485872in}{1.153085in}}%
\pgfpathlineto{\pgfqpoint{3.485872in}{1.174014in}}%
\pgfpathlineto{\pgfqpoint{3.521740in}{1.174014in}}%
\pgfpathlineto{\pgfqpoint{3.521740in}{1.153085in}}%
\pgfpathclose%
\pgfusepath{fill}%
\end{pgfscope}%
\begin{pgfscope}%
\pgfpathrectangle{\pgfqpoint{3.019583in}{0.169444in}}{\pgfqpoint{0.896708in}{1.339426in}}%
\pgfusepath{clip}%
\pgfsetbuttcap%
\pgfsetmiterjoin%
\definecolor{currentfill}{rgb}{0.527336,0.810611,0.645213}%
\pgfsetfillcolor{currentfill}%
\pgfsetlinewidth{0.000000pt}%
\definecolor{currentstroke}{rgb}{0.000000,0.000000,0.000000}%
\pgfsetstrokecolor{currentstroke}%
\pgfsetstrokeopacity{0.000000}%
\pgfsetdash{}{0pt}%
\pgfpathmoveto{\pgfqpoint{3.521740in}{1.174014in}}%
\pgfpathlineto{\pgfqpoint{3.485872in}{1.174014in}}%
\pgfpathlineto{\pgfqpoint{3.485872in}{1.194943in}}%
\pgfpathlineto{\pgfqpoint{3.521740in}{1.194943in}}%
\pgfpathlineto{\pgfqpoint{3.521740in}{1.174014in}}%
\pgfpathclose%
\pgfusepath{fill}%
\end{pgfscope}%
\begin{pgfscope}%
\pgfpathrectangle{\pgfqpoint{3.019583in}{0.169444in}}{\pgfqpoint{0.896708in}{1.339426in}}%
\pgfusepath{clip}%
\pgfsetbuttcap%
\pgfsetmiterjoin%
\definecolor{currentfill}{rgb}{0.484890,0.794002,0.645829}%
\pgfsetfillcolor{currentfill}%
\pgfsetlinewidth{0.000000pt}%
\definecolor{currentstroke}{rgb}{0.000000,0.000000,0.000000}%
\pgfsetstrokecolor{currentstroke}%
\pgfsetstrokeopacity{0.000000}%
\pgfsetdash{}{0pt}%
\pgfpathmoveto{\pgfqpoint{3.521740in}{1.194943in}}%
\pgfpathlineto{\pgfqpoint{3.485872in}{1.194943in}}%
\pgfpathlineto{\pgfqpoint{3.485872in}{1.215871in}}%
\pgfpathlineto{\pgfqpoint{3.521740in}{1.215871in}}%
\pgfpathlineto{\pgfqpoint{3.521740in}{1.194943in}}%
\pgfpathclose%
\pgfusepath{fill}%
\end{pgfscope}%
\begin{pgfscope}%
\pgfpathrectangle{\pgfqpoint{3.019583in}{0.169444in}}{\pgfqpoint{0.896708in}{1.339426in}}%
\pgfusepath{clip}%
\pgfsetbuttcap%
\pgfsetmiterjoin%
\definecolor{currentfill}{rgb}{0.442445,0.777393,0.646444}%
\pgfsetfillcolor{currentfill}%
\pgfsetlinewidth{0.000000pt}%
\definecolor{currentstroke}{rgb}{0.000000,0.000000,0.000000}%
\pgfsetstrokecolor{currentstroke}%
\pgfsetstrokeopacity{0.000000}%
\pgfsetdash{}{0pt}%
\pgfpathmoveto{\pgfqpoint{3.521740in}{1.215871in}}%
\pgfpathlineto{\pgfqpoint{3.485872in}{1.215871in}}%
\pgfpathlineto{\pgfqpoint{3.485872in}{1.236800in}}%
\pgfpathlineto{\pgfqpoint{3.521740in}{1.236800in}}%
\pgfpathlineto{\pgfqpoint{3.521740in}{1.215871in}}%
\pgfpathclose%
\pgfusepath{fill}%
\end{pgfscope}%
\begin{pgfscope}%
\pgfpathrectangle{\pgfqpoint{3.019583in}{0.169444in}}{\pgfqpoint{0.896708in}{1.339426in}}%
\pgfusepath{clip}%
\pgfsetbuttcap%
\pgfsetmiterjoin%
\definecolor{currentfill}{rgb}{0.400000,0.760784,0.647059}%
\pgfsetfillcolor{currentfill}%
\pgfsetlinewidth{0.000000pt}%
\definecolor{currentstroke}{rgb}{0.000000,0.000000,0.000000}%
\pgfsetstrokecolor{currentstroke}%
\pgfsetstrokeopacity{0.000000}%
\pgfsetdash{}{0pt}%
\pgfpathmoveto{\pgfqpoint{3.521740in}{1.236800in}}%
\pgfpathlineto{\pgfqpoint{3.485872in}{1.236800in}}%
\pgfpathlineto{\pgfqpoint{3.485872in}{1.257728in}}%
\pgfpathlineto{\pgfqpoint{3.521740in}{1.257728in}}%
\pgfpathlineto{\pgfqpoint{3.521740in}{1.236800in}}%
\pgfpathclose%
\pgfusepath{fill}%
\end{pgfscope}%
\begin{pgfscope}%
\pgfpathrectangle{\pgfqpoint{3.019583in}{0.169444in}}{\pgfqpoint{0.896708in}{1.339426in}}%
\pgfusepath{clip}%
\pgfsetbuttcap%
\pgfsetmiterjoin%
\definecolor{currentfill}{rgb}{0.368012,0.725106,0.661822}%
\pgfsetfillcolor{currentfill}%
\pgfsetlinewidth{0.000000pt}%
\definecolor{currentstroke}{rgb}{0.000000,0.000000,0.000000}%
\pgfsetstrokecolor{currentstroke}%
\pgfsetstrokeopacity{0.000000}%
\pgfsetdash{}{0pt}%
\pgfpathmoveto{\pgfqpoint{3.521740in}{1.257728in}}%
\pgfpathlineto{\pgfqpoint{3.485872in}{1.257728in}}%
\pgfpathlineto{\pgfqpoint{3.485872in}{1.278657in}}%
\pgfpathlineto{\pgfqpoint{3.521740in}{1.278657in}}%
\pgfpathlineto{\pgfqpoint{3.521740in}{1.257728in}}%
\pgfpathclose%
\pgfusepath{fill}%
\end{pgfscope}%
\begin{pgfscope}%
\pgfpathrectangle{\pgfqpoint{3.019583in}{0.169444in}}{\pgfqpoint{0.896708in}{1.339426in}}%
\pgfusepath{clip}%
\pgfsetbuttcap%
\pgfsetmiterjoin%
\definecolor{currentfill}{rgb}{0.336025,0.689427,0.676586}%
\pgfsetfillcolor{currentfill}%
\pgfsetlinewidth{0.000000pt}%
\definecolor{currentstroke}{rgb}{0.000000,0.000000,0.000000}%
\pgfsetstrokecolor{currentstroke}%
\pgfsetstrokeopacity{0.000000}%
\pgfsetdash{}{0pt}%
\pgfpathmoveto{\pgfqpoint{3.521740in}{1.278657in}}%
\pgfpathlineto{\pgfqpoint{3.485872in}{1.278657in}}%
\pgfpathlineto{\pgfqpoint{3.485872in}{1.299585in}}%
\pgfpathlineto{\pgfqpoint{3.521740in}{1.299585in}}%
\pgfpathlineto{\pgfqpoint{3.521740in}{1.278657in}}%
\pgfpathclose%
\pgfusepath{fill}%
\end{pgfscope}%
\begin{pgfscope}%
\pgfpathrectangle{\pgfqpoint{3.019583in}{0.169444in}}{\pgfqpoint{0.896708in}{1.339426in}}%
\pgfusepath{clip}%
\pgfsetbuttcap%
\pgfsetmiterjoin%
\definecolor{currentfill}{rgb}{0.304037,0.653749,0.691349}%
\pgfsetfillcolor{currentfill}%
\pgfsetlinewidth{0.000000pt}%
\definecolor{currentstroke}{rgb}{0.000000,0.000000,0.000000}%
\pgfsetstrokecolor{currentstroke}%
\pgfsetstrokeopacity{0.000000}%
\pgfsetdash{}{0pt}%
\pgfpathmoveto{\pgfqpoint{3.521740in}{1.299585in}}%
\pgfpathlineto{\pgfqpoint{3.485872in}{1.299585in}}%
\pgfpathlineto{\pgfqpoint{3.485872in}{1.320514in}}%
\pgfpathlineto{\pgfqpoint{3.521740in}{1.320514in}}%
\pgfpathlineto{\pgfqpoint{3.521740in}{1.299585in}}%
\pgfpathclose%
\pgfusepath{fill}%
\end{pgfscope}%
\begin{pgfscope}%
\pgfpathrectangle{\pgfqpoint{3.019583in}{0.169444in}}{\pgfqpoint{0.896708in}{1.339426in}}%
\pgfusepath{clip}%
\pgfsetbuttcap%
\pgfsetmiterjoin%
\definecolor{currentfill}{rgb}{0.272049,0.618070,0.706113}%
\pgfsetfillcolor{currentfill}%
\pgfsetlinewidth{0.000000pt}%
\definecolor{currentstroke}{rgb}{0.000000,0.000000,0.000000}%
\pgfsetstrokecolor{currentstroke}%
\pgfsetstrokeopacity{0.000000}%
\pgfsetdash{}{0pt}%
\pgfpathmoveto{\pgfqpoint{3.521740in}{1.320514in}}%
\pgfpathlineto{\pgfqpoint{3.485872in}{1.320514in}}%
\pgfpathlineto{\pgfqpoint{3.485872in}{1.341442in}}%
\pgfpathlineto{\pgfqpoint{3.521740in}{1.341442in}}%
\pgfpathlineto{\pgfqpoint{3.521740in}{1.320514in}}%
\pgfpathclose%
\pgfusepath{fill}%
\end{pgfscope}%
\begin{pgfscope}%
\pgfpathrectangle{\pgfqpoint{3.019583in}{0.169444in}}{\pgfqpoint{0.896708in}{1.339426in}}%
\pgfusepath{clip}%
\pgfsetbuttcap%
\pgfsetmiterjoin%
\definecolor{currentfill}{rgb}{0.240062,0.582391,0.720877}%
\pgfsetfillcolor{currentfill}%
\pgfsetlinewidth{0.000000pt}%
\definecolor{currentstroke}{rgb}{0.000000,0.000000,0.000000}%
\pgfsetstrokecolor{currentstroke}%
\pgfsetstrokeopacity{0.000000}%
\pgfsetdash{}{0pt}%
\pgfpathmoveto{\pgfqpoint{3.521740in}{1.341442in}}%
\pgfpathlineto{\pgfqpoint{3.485872in}{1.341442in}}%
\pgfpathlineto{\pgfqpoint{3.485872in}{1.362371in}}%
\pgfpathlineto{\pgfqpoint{3.521740in}{1.362371in}}%
\pgfpathlineto{\pgfqpoint{3.521740in}{1.341442in}}%
\pgfpathclose%
\pgfusepath{fill}%
\end{pgfscope}%
\begin{pgfscope}%
\pgfpathrectangle{\pgfqpoint{3.019583in}{0.169444in}}{\pgfqpoint{0.896708in}{1.339426in}}%
\pgfusepath{clip}%
\pgfsetbuttcap%
\pgfsetmiterjoin%
\definecolor{currentfill}{rgb}{0.208074,0.546713,0.735640}%
\pgfsetfillcolor{currentfill}%
\pgfsetlinewidth{0.000000pt}%
\definecolor{currentstroke}{rgb}{0.000000,0.000000,0.000000}%
\pgfsetstrokecolor{currentstroke}%
\pgfsetstrokeopacity{0.000000}%
\pgfsetdash{}{0pt}%
\pgfpathmoveto{\pgfqpoint{3.521740in}{1.362371in}}%
\pgfpathlineto{\pgfqpoint{3.485872in}{1.362371in}}%
\pgfpathlineto{\pgfqpoint{3.485872in}{1.383299in}}%
\pgfpathlineto{\pgfqpoint{3.521740in}{1.383299in}}%
\pgfpathlineto{\pgfqpoint{3.521740in}{1.362371in}}%
\pgfpathclose%
\pgfusepath{fill}%
\end{pgfscope}%
\begin{pgfscope}%
\pgfpathrectangle{\pgfqpoint{3.019583in}{0.169444in}}{\pgfqpoint{0.896708in}{1.339426in}}%
\pgfusepath{clip}%
\pgfsetbuttcap%
\pgfsetmiterjoin%
\definecolor{currentfill}{rgb}{0.212995,0.511419,0.730796}%
\pgfsetfillcolor{currentfill}%
\pgfsetlinewidth{0.000000pt}%
\definecolor{currentstroke}{rgb}{0.000000,0.000000,0.000000}%
\pgfsetstrokecolor{currentstroke}%
\pgfsetstrokeopacity{0.000000}%
\pgfsetdash{}{0pt}%
\pgfpathmoveto{\pgfqpoint{3.521740in}{1.383299in}}%
\pgfpathlineto{\pgfqpoint{3.485872in}{1.383299in}}%
\pgfpathlineto{\pgfqpoint{3.485872in}{1.404228in}}%
\pgfpathlineto{\pgfqpoint{3.521740in}{1.404228in}}%
\pgfpathlineto{\pgfqpoint{3.521740in}{1.383299in}}%
\pgfpathclose%
\pgfusepath{fill}%
\end{pgfscope}%
\begin{pgfscope}%
\pgfpathrectangle{\pgfqpoint{3.019583in}{0.169444in}}{\pgfqpoint{0.896708in}{1.339426in}}%
\pgfusepath{clip}%
\pgfsetbuttcap%
\pgfsetmiterjoin%
\definecolor{currentfill}{rgb}{0.240062,0.476355,0.714187}%
\pgfsetfillcolor{currentfill}%
\pgfsetlinewidth{0.000000pt}%
\definecolor{currentstroke}{rgb}{0.000000,0.000000,0.000000}%
\pgfsetstrokecolor{currentstroke}%
\pgfsetstrokeopacity{0.000000}%
\pgfsetdash{}{0pt}%
\pgfpathmoveto{\pgfqpoint{3.521740in}{1.404228in}}%
\pgfpathlineto{\pgfqpoint{3.485872in}{1.404228in}}%
\pgfpathlineto{\pgfqpoint{3.485872in}{1.425156in}}%
\pgfpathlineto{\pgfqpoint{3.521740in}{1.425156in}}%
\pgfpathlineto{\pgfqpoint{3.521740in}{1.404228in}}%
\pgfpathclose%
\pgfusepath{fill}%
\end{pgfscope}%
\begin{pgfscope}%
\pgfpathrectangle{\pgfqpoint{3.019583in}{0.169444in}}{\pgfqpoint{0.896708in}{1.339426in}}%
\pgfusepath{clip}%
\pgfsetbuttcap%
\pgfsetmiterjoin%
\definecolor{currentfill}{rgb}{0.267128,0.441292,0.697578}%
\pgfsetfillcolor{currentfill}%
\pgfsetlinewidth{0.000000pt}%
\definecolor{currentstroke}{rgb}{0.000000,0.000000,0.000000}%
\pgfsetstrokecolor{currentstroke}%
\pgfsetstrokeopacity{0.000000}%
\pgfsetdash{}{0pt}%
\pgfpathmoveto{\pgfqpoint{3.521740in}{1.425156in}}%
\pgfpathlineto{\pgfqpoint{3.485872in}{1.425156in}}%
\pgfpathlineto{\pgfqpoint{3.485872in}{1.446085in}}%
\pgfpathlineto{\pgfqpoint{3.521740in}{1.446085in}}%
\pgfpathlineto{\pgfqpoint{3.521740in}{1.425156in}}%
\pgfpathclose%
\pgfusepath{fill}%
\end{pgfscope}%
\begin{pgfscope}%
\pgfpathrectangle{\pgfqpoint{3.019583in}{0.169444in}}{\pgfqpoint{0.896708in}{1.339426in}}%
\pgfusepath{clip}%
\pgfsetbuttcap%
\pgfsetmiterjoin%
\definecolor{currentfill}{rgb}{0.294195,0.406228,0.680969}%
\pgfsetfillcolor{currentfill}%
\pgfsetlinewidth{0.000000pt}%
\definecolor{currentstroke}{rgb}{0.000000,0.000000,0.000000}%
\pgfsetstrokecolor{currentstroke}%
\pgfsetstrokeopacity{0.000000}%
\pgfsetdash{}{0pt}%
\pgfpathmoveto{\pgfqpoint{3.521740in}{1.446085in}}%
\pgfpathlineto{\pgfqpoint{3.485872in}{1.446085in}}%
\pgfpathlineto{\pgfqpoint{3.485872in}{1.467013in}}%
\pgfpathlineto{\pgfqpoint{3.521740in}{1.467013in}}%
\pgfpathlineto{\pgfqpoint{3.521740in}{1.446085in}}%
\pgfpathclose%
\pgfusepath{fill}%
\end{pgfscope}%
\begin{pgfscope}%
\pgfpathrectangle{\pgfqpoint{3.019583in}{0.169444in}}{\pgfqpoint{0.896708in}{1.339426in}}%
\pgfusepath{clip}%
\pgfsetbuttcap%
\pgfsetmiterjoin%
\definecolor{currentfill}{rgb}{0.321261,0.371165,0.664360}%
\pgfsetfillcolor{currentfill}%
\pgfsetlinewidth{0.000000pt}%
\definecolor{currentstroke}{rgb}{0.000000,0.000000,0.000000}%
\pgfsetstrokecolor{currentstroke}%
\pgfsetstrokeopacity{0.000000}%
\pgfsetdash{}{0pt}%
\pgfpathmoveto{\pgfqpoint{3.521740in}{1.467013in}}%
\pgfpathlineto{\pgfqpoint{3.485872in}{1.467013in}}%
\pgfpathlineto{\pgfqpoint{3.485872in}{1.487942in}}%
\pgfpathlineto{\pgfqpoint{3.521740in}{1.487942in}}%
\pgfpathlineto{\pgfqpoint{3.521740in}{1.467013in}}%
\pgfpathclose%
\pgfusepath{fill}%
\end{pgfscope}%
\begin{pgfscope}%
\pgfpathrectangle{\pgfqpoint{3.019583in}{0.169444in}}{\pgfqpoint{0.896708in}{1.339426in}}%
\pgfusepath{clip}%
\pgfsetbuttcap%
\pgfsetmiterjoin%
\definecolor{currentfill}{rgb}{0.348328,0.336101,0.647751}%
\pgfsetfillcolor{currentfill}%
\pgfsetlinewidth{0.000000pt}%
\definecolor{currentstroke}{rgb}{0.000000,0.000000,0.000000}%
\pgfsetstrokecolor{currentstroke}%
\pgfsetstrokeopacity{0.000000}%
\pgfsetdash{}{0pt}%
\pgfpathmoveto{\pgfqpoint{3.521740in}{1.487942in}}%
\pgfpathlineto{\pgfqpoint{3.485872in}{1.487942in}}%
\pgfpathlineto{\pgfqpoint{3.485872in}{1.508871in}}%
\pgfpathlineto{\pgfqpoint{3.521740in}{1.508871in}}%
\pgfpathlineto{\pgfqpoint{3.521740in}{1.487942in}}%
\pgfpathclose%
\pgfusepath{fill}%
\end{pgfscope}%
\begin{pgfscope}%
\pgfpathrectangle{\pgfqpoint{3.019583in}{0.169444in}}{\pgfqpoint{0.896708in}{1.339426in}}%
\pgfusepath{clip}%
\pgfsetbuttcap%
\pgfsetmiterjoin%
\definecolor{currentfill}{rgb}{0.368627,0.309804,0.635294}%
\pgfsetfillcolor{currentfill}%
\pgfsetlinewidth{0.000000pt}%
\definecolor{currentstroke}{rgb}{0.000000,0.000000,0.000000}%
\pgfsetstrokecolor{currentstroke}%
\pgfsetstrokeopacity{0.000000}%
\pgfsetdash{}{0pt}%
\pgfpathmoveto{\pgfqpoint{3.521740in}{1.508871in}}%
\pgfpathlineto{\pgfqpoint{3.485872in}{1.508871in}}%
\pgfpathlineto{\pgfqpoint{3.485872in}{1.529799in}}%
\pgfpathlineto{\pgfqpoint{3.521740in}{1.529799in}}%
\pgfpathlineto{\pgfqpoint{3.521740in}{1.508871in}}%
\pgfpathclose%
\pgfusepath{fill}%
\end{pgfscope}%
\begin{pgfscope}%
\pgfpathrectangle{\pgfqpoint{3.019583in}{0.169444in}}{\pgfqpoint{0.896708in}{1.339426in}}%
\pgfusepath{clip}%
\pgfsetbuttcap%
\pgfsetmiterjoin%
\definecolor{currentfill}{rgb}{0.997616,0.926105,0.625067}%
\pgfsetfillcolor{currentfill}%
\pgfsetlinewidth{0.000000pt}%
\definecolor{currentstroke}{rgb}{0.000000,0.000000,0.000000}%
\pgfsetstrokecolor{currentstroke}%
\pgfsetstrokeopacity{0.000000}%
\pgfsetdash{}{0pt}%
\pgfpathmoveto{\pgfqpoint{3.566575in}{0.755443in}}%
\pgfpathlineto{\pgfqpoint{3.584510in}{0.755443in}}%
\pgfpathlineto{\pgfqpoint{3.584510in}{0.797300in}}%
\pgfpathlineto{\pgfqpoint{3.566575in}{0.797300in}}%
\pgfpathlineto{\pgfqpoint{3.566575in}{0.755443in}}%
\pgfpathclose%
\pgfusepath{fill}%
\end{pgfscope}%
\begin{pgfscope}%
\pgfpathrectangle{\pgfqpoint{3.019583in}{0.169444in}}{\pgfqpoint{0.896708in}{1.339426in}}%
\pgfusepath{clip}%
\pgfsetbuttcap%
\pgfsetmiterjoin%
\definecolor{currentfill}{rgb}{0.998078,0.999231,0.746021}%
\pgfsetfillcolor{currentfill}%
\pgfsetlinewidth{0.000000pt}%
\definecolor{currentstroke}{rgb}{0.000000,0.000000,0.000000}%
\pgfsetstrokecolor{currentstroke}%
\pgfsetstrokeopacity{0.000000}%
\pgfsetdash{}{0pt}%
\pgfpathmoveto{\pgfqpoint{3.566575in}{0.839157in}}%
\pgfpathlineto{\pgfqpoint{3.584510in}{0.839157in}}%
\pgfpathlineto{\pgfqpoint{3.584510in}{0.881015in}}%
\pgfpathlineto{\pgfqpoint{3.566575in}{0.881015in}}%
\pgfpathlineto{\pgfqpoint{3.566575in}{0.839157in}}%
\pgfpathclose%
\pgfusepath{fill}%
\end{pgfscope}%
\begin{pgfscope}%
\pgfpathrectangle{\pgfqpoint{3.019583in}{0.169444in}}{\pgfqpoint{0.896708in}{1.339426in}}%
\pgfusepath{clip}%
\pgfsetbuttcap%
\pgfsetmiterjoin%
\definecolor{currentfill}{rgb}{0.967320,0.986928,0.698039}%
\pgfsetfillcolor{currentfill}%
\pgfsetlinewidth{0.000000pt}%
\definecolor{currentstroke}{rgb}{0.000000,0.000000,0.000000}%
\pgfsetstrokecolor{currentstroke}%
\pgfsetstrokeopacity{0.000000}%
\pgfsetdash{}{0pt}%
\pgfpathmoveto{\pgfqpoint{3.566575in}{0.881015in}}%
\pgfpathlineto{\pgfqpoint{3.602444in}{0.881015in}}%
\pgfpathlineto{\pgfqpoint{3.602444in}{0.922872in}}%
\pgfpathlineto{\pgfqpoint{3.566575in}{0.922872in}}%
\pgfpathlineto{\pgfqpoint{3.566575in}{0.881015in}}%
\pgfpathclose%
\pgfusepath{fill}%
\end{pgfscope}%
\begin{pgfscope}%
\pgfpathrectangle{\pgfqpoint{3.019583in}{0.169444in}}{\pgfqpoint{0.896708in}{1.339426in}}%
\pgfusepath{clip}%
\pgfsetbuttcap%
\pgfsetmiterjoin%
\definecolor{currentfill}{rgb}{0.936563,0.974625,0.650058}%
\pgfsetfillcolor{currentfill}%
\pgfsetlinewidth{0.000000pt}%
\definecolor{currentstroke}{rgb}{0.000000,0.000000,0.000000}%
\pgfsetstrokecolor{currentstroke}%
\pgfsetstrokeopacity{0.000000}%
\pgfsetdash{}{0pt}%
\pgfpathmoveto{\pgfqpoint{3.566575in}{0.922872in}}%
\pgfpathlineto{\pgfqpoint{3.638312in}{0.922872in}}%
\pgfpathlineto{\pgfqpoint{3.638312in}{0.964729in}}%
\pgfpathlineto{\pgfqpoint{3.566575in}{0.964729in}}%
\pgfpathlineto{\pgfqpoint{3.566575in}{0.922872in}}%
\pgfpathclose%
\pgfusepath{fill}%
\end{pgfscope}%
\begin{pgfscope}%
\pgfpathrectangle{\pgfqpoint{3.019583in}{0.169444in}}{\pgfqpoint{0.896708in}{1.339426in}}%
\pgfusepath{clip}%
\pgfsetbuttcap%
\pgfsetmiterjoin%
\definecolor{currentfill}{rgb}{0.905805,0.962322,0.602076}%
\pgfsetfillcolor{currentfill}%
\pgfsetlinewidth{0.000000pt}%
\definecolor{currentstroke}{rgb}{0.000000,0.000000,0.000000}%
\pgfsetstrokecolor{currentstroke}%
\pgfsetstrokeopacity{0.000000}%
\pgfsetdash{}{0pt}%
\pgfpathmoveto{\pgfqpoint{3.566575in}{0.964729in}}%
\pgfpathlineto{\pgfqpoint{3.692115in}{0.964729in}}%
\pgfpathlineto{\pgfqpoint{3.692115in}{1.006586in}}%
\pgfpathlineto{\pgfqpoint{3.566575in}{1.006586in}}%
\pgfpathlineto{\pgfqpoint{3.566575in}{0.964729in}}%
\pgfpathclose%
\pgfusepath{fill}%
\end{pgfscope}%
\begin{pgfscope}%
\pgfpathrectangle{\pgfqpoint{3.019583in}{0.169444in}}{\pgfqpoint{0.896708in}{1.339426in}}%
\pgfusepath{clip}%
\pgfsetbuttcap%
\pgfsetmiterjoin%
\definecolor{currentfill}{rgb}{0.838447,0.934948,0.608997}%
\pgfsetfillcolor{currentfill}%
\pgfsetlinewidth{0.000000pt}%
\definecolor{currentstroke}{rgb}{0.000000,0.000000,0.000000}%
\pgfsetstrokecolor{currentstroke}%
\pgfsetstrokeopacity{0.000000}%
\pgfsetdash{}{0pt}%
\pgfpathmoveto{\pgfqpoint{3.566575in}{1.006586in}}%
\pgfpathlineto{\pgfqpoint{3.710049in}{1.006586in}}%
\pgfpathlineto{\pgfqpoint{3.710049in}{1.048443in}}%
\pgfpathlineto{\pgfqpoint{3.566575in}{1.048443in}}%
\pgfpathlineto{\pgfqpoint{3.566575in}{1.006586in}}%
\pgfpathclose%
\pgfusepath{fill}%
\end{pgfscope}%
\begin{pgfscope}%
\pgfpathrectangle{\pgfqpoint{3.019583in}{0.169444in}}{\pgfqpoint{0.896708in}{1.339426in}}%
\pgfusepath{clip}%
\pgfsetbuttcap%
\pgfsetmiterjoin%
\definecolor{currentfill}{rgb}{0.765859,0.905421,0.623760}%
\pgfsetfillcolor{currentfill}%
\pgfsetlinewidth{0.000000pt}%
\definecolor{currentstroke}{rgb}{0.000000,0.000000,0.000000}%
\pgfsetstrokecolor{currentstroke}%
\pgfsetstrokeopacity{0.000000}%
\pgfsetdash{}{0pt}%
\pgfpathmoveto{\pgfqpoint{3.566575in}{1.048443in}}%
\pgfpathlineto{\pgfqpoint{3.656246in}{1.048443in}}%
\pgfpathlineto{\pgfqpoint{3.656246in}{1.090300in}}%
\pgfpathlineto{\pgfqpoint{3.566575in}{1.090300in}}%
\pgfpathlineto{\pgfqpoint{3.566575in}{1.048443in}}%
\pgfpathclose%
\pgfusepath{fill}%
\end{pgfscope}%
\begin{pgfscope}%
\pgfpathrectangle{\pgfqpoint{3.019583in}{0.169444in}}{\pgfqpoint{0.896708in}{1.339426in}}%
\pgfusepath{clip}%
\pgfsetbuttcap%
\pgfsetmiterjoin%
\definecolor{currentfill}{rgb}{0.693272,0.875894,0.638524}%
\pgfsetfillcolor{currentfill}%
\pgfsetlinewidth{0.000000pt}%
\definecolor{currentstroke}{rgb}{0.000000,0.000000,0.000000}%
\pgfsetstrokecolor{currentstroke}%
\pgfsetstrokeopacity{0.000000}%
\pgfsetdash{}{0pt}%
\pgfpathmoveto{\pgfqpoint{3.566575in}{1.090300in}}%
\pgfpathlineto{\pgfqpoint{3.656246in}{1.090300in}}%
\pgfpathlineto{\pgfqpoint{3.656246in}{1.132157in}}%
\pgfpathlineto{\pgfqpoint{3.566575in}{1.132157in}}%
\pgfpathlineto{\pgfqpoint{3.566575in}{1.090300in}}%
\pgfpathclose%
\pgfusepath{fill}%
\end{pgfscope}%
\begin{pgfscope}%
\pgfpathrectangle{\pgfqpoint{3.019583in}{0.169444in}}{\pgfqpoint{0.896708in}{1.339426in}}%
\pgfusepath{clip}%
\pgfsetbuttcap%
\pgfsetmiterjoin%
\definecolor{currentfill}{rgb}{0.612226,0.843829,0.643983}%
\pgfsetfillcolor{currentfill}%
\pgfsetlinewidth{0.000000pt}%
\definecolor{currentstroke}{rgb}{0.000000,0.000000,0.000000}%
\pgfsetstrokecolor{currentstroke}%
\pgfsetstrokeopacity{0.000000}%
\pgfsetdash{}{0pt}%
\pgfpathmoveto{\pgfqpoint{3.566575in}{1.132157in}}%
\pgfpathlineto{\pgfqpoint{3.674180in}{1.132157in}}%
\pgfpathlineto{\pgfqpoint{3.674180in}{1.174014in}}%
\pgfpathlineto{\pgfqpoint{3.566575in}{1.174014in}}%
\pgfpathlineto{\pgfqpoint{3.566575in}{1.132157in}}%
\pgfpathclose%
\pgfusepath{fill}%
\end{pgfscope}%
\begin{pgfscope}%
\pgfpathrectangle{\pgfqpoint{3.019583in}{0.169444in}}{\pgfqpoint{0.896708in}{1.339426in}}%
\pgfusepath{clip}%
\pgfsetbuttcap%
\pgfsetmiterjoin%
\definecolor{currentfill}{rgb}{0.527336,0.810611,0.645213}%
\pgfsetfillcolor{currentfill}%
\pgfsetlinewidth{0.000000pt}%
\definecolor{currentstroke}{rgb}{0.000000,0.000000,0.000000}%
\pgfsetstrokecolor{currentstroke}%
\pgfsetstrokeopacity{0.000000}%
\pgfsetdash{}{0pt}%
\pgfpathmoveto{\pgfqpoint{3.566575in}{1.174014in}}%
\pgfpathlineto{\pgfqpoint{3.602444in}{1.174014in}}%
\pgfpathlineto{\pgfqpoint{3.602444in}{1.215871in}}%
\pgfpathlineto{\pgfqpoint{3.566575in}{1.215871in}}%
\pgfpathlineto{\pgfqpoint{3.566575in}{1.174014in}}%
\pgfpathclose%
\pgfusepath{fill}%
\end{pgfscope}%
\begin{pgfscope}%
\pgfpathrectangle{\pgfqpoint{3.019583in}{0.169444in}}{\pgfqpoint{0.896708in}{1.339426in}}%
\pgfusepath{clip}%
\pgfsetbuttcap%
\pgfsetmiterjoin%
\definecolor{currentfill}{rgb}{0.442445,0.777393,0.646444}%
\pgfsetfillcolor{currentfill}%
\pgfsetlinewidth{0.000000pt}%
\definecolor{currentstroke}{rgb}{0.000000,0.000000,0.000000}%
\pgfsetstrokecolor{currentstroke}%
\pgfsetstrokeopacity{0.000000}%
\pgfsetdash{}{0pt}%
\pgfpathmoveto{\pgfqpoint{3.566575in}{1.215871in}}%
\pgfpathlineto{\pgfqpoint{3.602444in}{1.215871in}}%
\pgfpathlineto{\pgfqpoint{3.602444in}{1.257728in}}%
\pgfpathlineto{\pgfqpoint{3.566575in}{1.257728in}}%
\pgfpathlineto{\pgfqpoint{3.566575in}{1.215871in}}%
\pgfpathclose%
\pgfusepath{fill}%
\end{pgfscope}%
\begin{pgfscope}%
\pgfpathrectangle{\pgfqpoint{3.019583in}{0.169444in}}{\pgfqpoint{0.896708in}{1.339426in}}%
\pgfusepath{clip}%
\pgfsetbuttcap%
\pgfsetmiterjoin%
\definecolor{currentfill}{rgb}{0.368012,0.725106,0.661822}%
\pgfsetfillcolor{currentfill}%
\pgfsetlinewidth{0.000000pt}%
\definecolor{currentstroke}{rgb}{0.000000,0.000000,0.000000}%
\pgfsetstrokecolor{currentstroke}%
\pgfsetstrokeopacity{0.000000}%
\pgfsetdash{}{0pt}%
\pgfpathmoveto{\pgfqpoint{3.566575in}{1.257728in}}%
\pgfpathlineto{\pgfqpoint{3.584510in}{1.257728in}}%
\pgfpathlineto{\pgfqpoint{3.584510in}{1.299585in}}%
\pgfpathlineto{\pgfqpoint{3.566575in}{1.299585in}}%
\pgfpathlineto{\pgfqpoint{3.566575in}{1.257728in}}%
\pgfpathclose%
\pgfusepath{fill}%
\end{pgfscope}%
\begin{pgfscope}%
\pgfpathrectangle{\pgfqpoint{3.019583in}{0.169444in}}{\pgfqpoint{0.896708in}{1.339426in}}%
\pgfusepath{clip}%
\pgfsetbuttcap%
\pgfsetmiterjoin%
\definecolor{currentfill}{rgb}{0.304037,0.653749,0.691349}%
\pgfsetfillcolor{currentfill}%
\pgfsetlinewidth{0.000000pt}%
\definecolor{currentstroke}{rgb}{0.000000,0.000000,0.000000}%
\pgfsetstrokecolor{currentstroke}%
\pgfsetstrokeopacity{0.000000}%
\pgfsetdash{}{0pt}%
\pgfpathmoveto{\pgfqpoint{3.566575in}{1.299585in}}%
\pgfpathlineto{\pgfqpoint{3.638312in}{1.299585in}}%
\pgfpathlineto{\pgfqpoint{3.638312in}{1.341442in}}%
\pgfpathlineto{\pgfqpoint{3.566575in}{1.341442in}}%
\pgfpathlineto{\pgfqpoint{3.566575in}{1.299585in}}%
\pgfpathclose%
\pgfusepath{fill}%
\end{pgfscope}%
\begin{pgfscope}%
\pgfpathrectangle{\pgfqpoint{3.019583in}{0.169444in}}{\pgfqpoint{0.896708in}{1.339426in}}%
\pgfusepath{clip}%
\pgfsetbuttcap%
\pgfsetmiterjoin%
\definecolor{currentfill}{rgb}{0.212995,0.511419,0.730796}%
\pgfsetfillcolor{currentfill}%
\pgfsetlinewidth{0.000000pt}%
\definecolor{currentstroke}{rgb}{0.000000,0.000000,0.000000}%
\pgfsetstrokecolor{currentstroke}%
\pgfsetstrokeopacity{0.000000}%
\pgfsetdash{}{0pt}%
\pgfpathmoveto{\pgfqpoint{3.566575in}{1.383299in}}%
\pgfpathlineto{\pgfqpoint{3.602444in}{1.383299in}}%
\pgfpathlineto{\pgfqpoint{3.602444in}{1.425156in}}%
\pgfpathlineto{\pgfqpoint{3.566575in}{1.425156in}}%
\pgfpathlineto{\pgfqpoint{3.566575in}{1.383299in}}%
\pgfpathclose%
\pgfusepath{fill}%
\end{pgfscope}%
\begin{pgfscope}%
\pgfpathrectangle{\pgfqpoint{3.019583in}{0.169444in}}{\pgfqpoint{0.896708in}{1.339426in}}%
\pgfusepath{clip}%
\pgfsetbuttcap%
\pgfsetmiterjoin%
\definecolor{currentfill}{rgb}{0.267128,0.441292,0.697578}%
\pgfsetfillcolor{currentfill}%
\pgfsetlinewidth{0.000000pt}%
\definecolor{currentstroke}{rgb}{0.000000,0.000000,0.000000}%
\pgfsetstrokecolor{currentstroke}%
\pgfsetstrokeopacity{0.000000}%
\pgfsetdash{}{0pt}%
\pgfpathmoveto{\pgfqpoint{3.566575in}{1.425156in}}%
\pgfpathlineto{\pgfqpoint{3.584510in}{1.425156in}}%
\pgfpathlineto{\pgfqpoint{3.584510in}{1.467013in}}%
\pgfpathlineto{\pgfqpoint{3.566575in}{1.467013in}}%
\pgfpathlineto{\pgfqpoint{3.566575in}{1.425156in}}%
\pgfpathclose%
\pgfusepath{fill}%
\end{pgfscope}%
\begin{pgfscope}%
\definecolor{textcolor}{rgb}{0.000000,0.000000,0.000000}%
\pgfsetstrokecolor{textcolor}%
\pgfsetfillcolor{textcolor}%
\pgftext[x=3.378267in,y=0.169444in,,]{\color{textcolor}\setmainfont{Lato}\rmfamily\fontsize{8.000000}{9.600000}\selectfont -8}%
\end{pgfscope}%
\begin{pgfscope}%
\definecolor{textcolor}{rgb}{0.000000,0.000000,0.000000}%
\pgfsetstrokecolor{textcolor}%
\pgfsetfillcolor{textcolor}%
\pgftext[x=3.378267in,y=0.504301in,,]{\color{textcolor}\setmainfont{Lato}\rmfamily\fontsize{8.000000}{9.600000}\selectfont -4}%
\end{pgfscope}%
\begin{pgfscope}%
\definecolor{textcolor}{rgb}{0.000000,0.000000,0.000000}%
\pgfsetstrokecolor{textcolor}%
\pgfsetfillcolor{textcolor}%
\pgftext[x=3.378267in,y=0.839157in,,]{\color{textcolor}\setmainfont{Lato}\rmfamily\fontsize{8.000000}{9.600000}\selectfont 0}%
\end{pgfscope}%
\begin{pgfscope}%
\definecolor{textcolor}{rgb}{0.000000,0.000000,0.000000}%
\pgfsetstrokecolor{textcolor}%
\pgfsetfillcolor{textcolor}%
\pgftext[x=3.378267in,y=1.174014in,,]{\color{textcolor}\setmainfont{Lato}\rmfamily\fontsize{8.000000}{9.600000}\selectfont 4}%
\end{pgfscope}%
\begin{pgfscope}%
\definecolor{textcolor}{rgb}{0.000000,0.000000,0.000000}%
\pgfsetstrokecolor{textcolor}%
\pgfsetfillcolor{textcolor}%
\pgftext[x=3.378267in,y=1.508871in,,]{\color{textcolor}\setmainfont{Lato}\rmfamily\fontsize{8.000000}{9.600000}\selectfont 8+}%
\end{pgfscope}%
\end{pgfpicture}%
\makeatother%
\endgroup%
 \\
\footnotesize{Source: Bureau of Economic Analysis}\\

\vspace{2mm}

\begin{minipage}{0.76\textwidth}

\small \input{text/gdp_state.txt} 
\end{minipage}

\vspace{2mm}

\normalsize \textbf{Real GDP Growth by State}\\
\footnotesize{\textit{quarterly growth at seasonally adjusted annualized rate \hspace{20mm} total growth, \input{text/gdp_state_date.txt}}\\ 

\vspace{-4.5mm}
\hspace{-2mm} \noindent \rowcolors{1}{}{black!5} \setlength{\tabcolsep}{3.7pt} \color{black!90}
		{\renewcommand{\arraystretch}{1.44}
		 \begin{tabular}{p{30mm} R{7mm} R{7mm} R{7mm} R{7mm} R{7mm} p{0mm} R{9mm} R{9mm} R{10mm} }
		\input data/gdp_state_table.tex \hline
		\end{tabular}}	\\

\vspace{-3mm}	
\footnotesize{Source: Bureau of Economic Analysis}

\newpage
\hypertarget{ofa}{}
\begin{minipage}{0.76\textwidth}
\section*{Financial Accounts}
\small The Federal Reserve \href{https://www.federalreserve.gov/releases/z1/}{report} transactions and levels for US assets and liabilities in the \textbf{financial accounts}. This section provides a high-level overview of the financial accounts, including liabilities, sectoral balances, wealth, and investment. 
\subsection*{Liabilities}r
\small \textbf{Total liabilities} is the outstanding debt at a point in time. While the vast majority of US borrowing comes from US lending, that is the net result on balance sheets is zero, it is still useful to consider how total liabilities vary by major sector and over time. 

\input{text/liabsh.txt}
\vspace{1mm}

\normalsize \textbf{Total Liabilities by Sector}\\
\footnotesize{\textit{share of GDP, percent}}\\
\hspace*{-2mm} \begin{tikzpicture}
	\begin{axis}[\bbar{y}{0}, \dateaxisticks ytick={0, 25, 50, 75, 100, 125, 150, 175}, 
		yticklabel style={text width=1.4em}, ymin=9,
		xticklabel={`\short{\year}}, clip=false, height=6.4cm, width=12.2cm]
	\rbars
	\thickline{blue!70!white}{date}{FL314190005.Q}{data/liabsh.csv}
	\absnode{{2014-05-01}}{113}{\footnotesize \color{blue}Federal Gov.}
	\thickline{cyan!60!white}{date}{FL214190005.Q}{data/liabsh.csv}
	\absnode{{2011-05-01}}{39}{\footnotesize \color{cyan}State \& Local Gov.}
	\thickline{orange!90!red}{date}{FL154190005.Q}{data/liabsh.csv}
	\absnode{{2014-10-01}}{75}{\footnotesize \color{orange!90!red}Household}
	\thickline{green!72!black}{date}{Business}{data/liabsh.csv}
	\absnode{{1991-10-01}}{116}{\footnotesize \color{green!72!black}Private Business}
	\input{text/liabsh_node.txt}
	\end{axis}
\end{tikzpicture}\\
\footnotesize{Source: Federal Reserve, Bureau of Economic Analysis}  \hfill \tbllink{liabsh.csv}
\end{minipage}
\newpage
\begin{minipage}{0.76\textwidth}
\small Calculating the contribution of different sectors to the \textbf{total change in borrowing} can highlight risks to economic activity. For example, the tech bubble of the late 1990s and early 2000s shows up as a large increase in corporate borrowing. The housing bubble from the 1990s to 2007 shows up as an increase in household borrowing. Government borrowing increased following the collapse of the housing bubble, in an effort to compensate for the massive fall in wage income. Keep in mind, however, that the vast majority of liabilities in the domestic economy are to other domestic parties.

&   2023  Q3 & `23  Q2 & `23  Q1 & `22  Q4 & `22  Q3 &4-year&10-year&30-year\\  Total  Real  Debt  Growth &1.40&0.66&-0.48&-0.16&0.59&3.50&3.31&3.79\\  \hspace{-2mm}\cbox{orange!90!red}Household  \&  Nonprofit &-0.14&-0.01&-0.03&-0.01&0.10&0.43&0.41&0.93\\  \hspace{3mm}  Home  Mortgages &-0.03&0.07&0.11&0.09&0.13&0.32&0.18&0.60\\  \hspace{3mm}  Consumer  Credit &0.02&0.09&0.11&0.11&0.05&0.08&0.20&0.27\\  \hspace{-2mm}\cbox{green!72!black}Business &-0.10&-0.10&-0.04&0.12&0.21&0.62&1.07&1.09\\  \hspace{3mm}Corporate  Business &-0.09&-0.12&-0.09&0.08&0.15&0.44&0.72&0.68\\  \hspace{5mm}  Debt  Securities &-0.08&-0.15&-0.36&-0.55&-0.66&0.05&0.34&0.42\\  \hspace{5mm}  Loans &-0.01&0.03&0.27&0.64&0.81&0.38&0.38&0.26\\  \hspace{3mm}Noncorporate  Business &-0.00&0.02&0.04&0.04&0.06&0.18&0.36&0.41\\  \hspace{5mm}  Commercial  Mortgages &-0.01&-0.02&-0.03&-0.05&-0.06&-0.01&0.05&0.10\\  \hspace{-2mm}\cbox{cyan!60!white}State  \&  Local  Government &-0.17&-0.20&-0.26&-0.35&-0.34&-0.13&-0.13&0.09\\  \hspace{-2mm}\cbox{blue!70!white}Federal  Government &1.80&0.97&-0.14&0.08&0.63&2.58&1.95&1.69\\  
\vspace{2mm}

\normalsize \textbf{Real Debt Growth}\\
\footnotesize{\textit{contribution to one-year percent change in liabilities, adjusted by PCE price deflator}}\\
\hspace*{-3mm} \begin{tikzpicture}
	\begin{axis}[\bbar{y}{0}, \dateaxisticks ytick={-2, 0, 2, 4, 6, 8, 10}, 
		xmin = 1989-01-01, height=6.5cm, width=13.0cm, enlarge y limits={0.05}, 
		xticklabel={`\short{\year}}, yticklabel style={text width=1.0em}, clip=false, 
		legend style={at={(0.95, 1.12)}}]
	\rbars
	\sbar{blue!70!white}{date}{FL314190005.Q}{data/liabgr.csv}
	\sbar{cyan!52!white}{date}{FL214190005.Q}{data/liabgr.csv}
	\sbar{orange!70!white}{date}{FL154190005.Q}{data/liabgr.csv}
	\sbar{lime!70!white}{date}{FL104190005.Q}{data/liabgr.csv}
	\sbar{green!72!black}{date}{FL114190005.Q}{data/liabgr.csv}
	\legend{Fed. Gov.,State/Loc. Gov., Households, Corp bus., Non-corp. bus.};
	\end{axis}
\end{tikzpicture}\\
\footnotesize{Source: Federal Reserve, Bureau of Economic Analysis}  \hfill \tbllink{liabgr.csv}

\vspace{4mm}

\normalsize \textbf{Real Debt Growth}\\
\footnotesize{\textit{contribution to one-year real growth \hspace{48mm} moving averages}}\\ 
\rowcolors{1}{}{black!5} \setlength{\tabcolsep}{3.1pt} \color{black!90}
		{\renewcommand{\arraystretch}{1.55}
		 \begin{tabular}{p{38mm} R{6.7mm} R{6.7mm} R{6.7mm} R{6.7mm} R{6.7mm} 
		   R{7.8mm} R{6.7mm} R{6.7mm} }
			 \input data/liabgr.tex \hline
		\end{tabular}}	\vspace{-2mm}
		
\footnotesize{Source: Federal Reserve, Bureau of Economic Analysis}
\end{minipage}
\newpage
\hypertarget{ofsb}{}
\begin{minipage}{0.76\textwidth}
\subsection*{Sectoral Balances}
\small A high-level overview of US financial activities can be provided by the \textbf{sectoral financial balances}. This measure divides the world into three sectors: the US private sector (see\cbox{green!70!black}), the US government (see\cbox{yellow!70!orange}), and the rest of the world (see\cbox{blue!90!black}), then examines the net lending and borrowing between the groups. In any given period, one sector's borrowing is another sector's lending, and as such the total will always sum to zero. 

Note that a sector runs a surplus in a given accounting period when its aggregate income is higher than its outlays. When this occurs the sector becomes a lender of the savings created by having more income than expenses. Likewise, if a sector spends more than it takes in, it borrows the difference. Therefore, when the government sector runs a deficit and is a net borrower, it is creating a surplus for other sectors because it is taking in less income through taxes than it is creating through spending. 
\vspace{1mm}

\normalsize \textbf{Sectoral Financial Balance}\\
\footnotesize{\textit{net lending (+) or borrowing (-), NIPA basis, by sector, as share of GDP}}
\vspace{3.8cm}

\hspace{4mm} \begin{tikzpicture}[overlay]
	\begin{axis}[\bbar{y}{0}, \dateaxisticks ytick={-30, -20, -10, 0, 10, 20},
		height=5.8cm, ymax=17.4, clip=false,
		xticklabel={`\short{\year}}, yticklabel style={text width=1.4em}, 
		legend cell align={left}, reverse legend,
		legend style={legend columns=1, at={(0.36, 0.99)}}]
	\rbars
	\sbar{yellow!70!orange}{date}{GOV}{data/sectbal.csv}
	\sbar{green!70!black}{date}{PRIV}{data/sectbal.csv}
	\sbar{blue!90!black}{date}{ROW}{data/sectbal.csv}
    \node[align=left] (source) at (axis cs:2017-06-15,10.5) {%
      	\scriptsize TCJA \\ \scriptsize repatriation};
    \node (destination) at (axis cs:2017-11-01,3.1) {};
    \draw[-] (source)--(destination);
	\legend{ \ Government, \ Private, \ Rest of World};
	\end{axis}
\end{tikzpicture}

\footnotesize{Source: Bureau of Economic Analysis} \hfill \tbllink{sectbal.csv} 
\vspace{3mm}

\small \input{text/sectbal.txt} 

\input{text/sectbal2.txt}
\vspace{1mm}

\normalsize \textbf{Domestic Private Sector Financial Balance}\\
\footnotesize{\textit{net lending (+) or borrowing (-), NIPA basis, by sector, share of GDP, percent}}
\vspace{2.6cm}

\hspace{4mm} \begin{tikzpicture}[overlay]
	\begin{axis}[\bbar{y}{0}, \dateaxisticks ytick={-5, 0, 5, 10, 15, 20},
		height=4.6cm, ymax=18.6, clip=false, 
		xticklabel={`\short{\year}}, yticklabel style={text width=1.4em}, 
		legend cell align={left}, 
		legend style={legend columns=1, at={(0.35, 0.99)}}]
	\rbars
	\sbar{purple!50!red}{date}{W995RC}{data/sectbal2.csv}
	\sbar{orange!90!yellow}{date}{W996RC}{data/sectbal2.csv}
	\legend{ \ Private Businesses, \ Households};
	\end{axis}
\end{tikzpicture}

\footnotesize{Source: Bureau of Economic Analysis} \hfill \tbllink{sectbal2.csv}
\end{minipage}
\newpage
\begin{minipage}{0.76\textwidth}
\subsection*{Wealth}
\hypertarget{ofw}{}
\small \textbf{Net wealth} is the sum of domestic tangible assets, such as land (excluding public land), structures, and equipment, minus foreign financial claims on these assets, plus domestic claims on foreign assets. &   2023  Q3 & `23  Q2 & `23  Q1 & `22  Q3 &2019& 2005  --`07 &1989\\  US  Net  Wealth &513.9&527.1&509.0&514.2&472.6&478.7&383.1\\  \hspace{2mm}  Households  \&  Nonprofits &213.1&215.6&207.6&217.5&182.5&217.8&169.2\\  \hspace{2mm}  Noncorporate  Businesses &70.7&71.9&71.6&76.3&68.1&73.2&71.1\\  \hspace{2mm}  Domestic  Corporations &217.8&228.3&214.3&206.4&199.5&130.4&74.4\\  \hspace{2mm}  Federal  Government &16.5&16.7&16.7&16.8&16.8&18.3&25.7\\  \hspace{2mm}  State  \&  Local  Government &57.2&58.1&58.4&58.7&55.3&48.7&41.9\\  \hspace{2mm}  Net  Claims  on  ROW &-61.5&-63.5&-59.6&-61.4&-49.6&-9.7&0.8\\  \hspace{5mm}  US  Claims  on  ROW &99.7&104.1&102.9&92.1&106.8&86.6&28.6\\  \hspace{5mm}  Less:  ROW  Claims  on  US &161.3&167.6&162.5&153.5&156.4&96.3&27.8\\   

\input{text/wealthgdp.txt}
\vspace{1mm}

\normalsize \textbf{Net Wealth to GDP Ratio}\\
\footnotesize{\textit{total US wealth divided by GDP}}\\
\hspace*{-3mm} \begin{tikzpicture}
	\begin{axis}[\bbar{y}{0}, \dateaxisticks ytick={0, 1, 2, 3, 4, 5, 6}, ymin=0.15,
		height=5.4cm, width=13.0cm, enlarge y limits={0.05}, 
		xticklabel={`\short{\year}}, clip=false, 
		legend style={at={(0.72, 1.01)}, fill=white}]
	\rbars
	\sbar{cyan!35!white}{date}{Other}{data/wealthgdp.csv}
	\sbar{green!80!blue}{date}{Real Estate}{data/wealthgdp.csv}
	\sbar{magenta!50!violet}{date}{Corporate Equities}{data/wealthgdp.csv}
	\legend{Other (Net), Residential Real Estate, Corporate Equities};
	\end{axis}
\end{tikzpicture}\\
\footnotesize{Source: Federal Reserve} \hfill \tbllink{wealthgdp.csv}
\vspace{2mm}

\small The tangible assets of each major domestic sector are show below, along with the summary of financial claims between the rest of the world and the US.
\vspace{1mm}

\normalsize \textbf{Derivation of US Net Wealth}\\
\footnotesize{\textit{share of GDP, percentage points}}\\ 
\rowcolors{1}{}{black!5} \setlength{\tabcolsep}{3.5pt} \color{black!90}
		{\renewcommand{\arraystretch}{1.65}
		 \begin{tabular}{p{40mm} R{7.8mm} R{7.8mm} R{7.8mm} R{7.8mm} 
		   R{8.2mm} R{8.2mm} R{8.2mm} }
			 \input data/uswealth.tex \hline
		\end{tabular}}	\vspace{-2mm}
		
\footnotesize{Source: Federal Reserve, Bureau of Economic Analysis}
\end{minipage}
\newpage
\begin{minipage}{0.76\textwidth}
\subsection*{Investment}

\small Investment is the exchange of a liquid asset, such as cash, for a less-liquid asset, such as a building. 

Economic data use varying definitions of asset, for purposes of measuring investment. \textbf{\href{https://www.bea.gov/help/glossary/gross-private-fixed-investment}{Gross private fixed investment}}, as measured in the national accounts, includes residential fixed investment, such as the construction and improvement of houses, apartment buildings, and other residential property (see\cbox{blue!90!black}), but does not include spending on durable goods, such as automobiles, appliances, or furniture. 

Non-residential private fixed investment includes the construction and improvement of offices, warehouses, factories, and other commercial and industrial property (see\cbox{yellow!50!orange}), as well as purchases of equipment, software, and intellectual property products. 

Additionally, the change in private inventories (see\cbox{red}) at the end of the accounting period, whether intentional or unintentional, affects GDP growth in the period. Inventory investment is grouped in the national accounts with gross private investment, but is not fixed investment.

\input{text/inv_text.txt} 
\vspace{2mm} 

\normalsize \textbf{Private Gross Investment}\\
\footnotesize{\textit{percentage point contribution to real GDP growth}}
\vspace{5.2cm}

\hspace{4mm} \begin{tikzpicture}[overlay]
	\begin{axis}[\bbar{y}{0}, \dateaxisticks ytick={-10, -5, 0, 5, 10},
		yticklabel style={text width=1.4em}, height=7.2cm, 
		reverse legend, legend cell align=left,
		xticklabel={`\short{\year}}, clip=false, ymax=8.9,
		legend style={at={(0.46, 0.99)}, legend columns=1}]
	\rbars
	\sbar{yellow!50!orange}{date}{A008RY}{data/inv.csv}
	\sbar{blue!90!black}{date}{A011RY}{data/inv.csv}
	\sbar{red}{date}{A014RY}{data/inv.csv}
	\legend{ \ Non-Residential, \ Residential, \ Change in Private Inventories};
	\end{axis}
\end{tikzpicture}

\footnotesize{Source: Bureau of Economic Analysis} \hfill \tbllink{inv.csv} 
\end{minipage}
\newpage
\begin{minipage}{0.76\textwidth}
\subsubsection*{Net Fixed Investment}
\vspace{-1.5mm}
\small Gross investment in fixed assets covers both depreciation, the wearing down of existing assets, and new investment. Fixed investment less depreciation is referred to as \textbf{net fixed investment}, and positive net investment represents new or expanded investment. 

\begin{minipage}{0.445\textwidth}
\small \input{text/z1_nfi_main.txt}
\end{minipage} \hfill
\begin{minipage}{0.49\textwidth}
\normalsize \textbf{Net Fixed Investment}\\
\footnotesize{\textit{share of GDP, percent}}\\
\hspace*{-2mm} \begin{tikzpicture}
	\begin{axis}[\bbar{y}{0}, \shdateaxisticks ytick={0, 2, 4, 6, 8, 10}, ymin=0.2, 
		yticklabel style={text width=1.0em},
		xticklabel={`\short{\year}}, clip=false, height=5.5cm, width=6.35cm]
	\rbars
	\absnode{{2009-06-01}}{6.8}{\footnotesize \color{purple}Fixed\\\footnotesize \color{purple}Investment}
	\stdline{purple}{date}{TOT}{data/z1_nfi.csv}
	\absnode{{1991-01-01}}{4.1}{\footnotesize \color{blue!80!green!95!white}Consumer\\\footnotesize \color{blue!80!green!95!white}Durables}
	\stdline{blue!80!green!95!white}{date}{DG}{data/z1_nfi.csv}
	\input{text/z1_nfi_nodes.txt}
	\end{axis}
\end{tikzpicture}\\
\footnotesize{Source: Federal Reserve, BEA} \hfill \tbllink{z1_nfi.csv}
\end{minipage}

\small \input{text/z1_nfi_sect.txt}
\vspace{0.5mm}

\normalsize \textbf{Net Fixed Investment by Sector or Type}\\
\footnotesize{\textit{share GDP, percent}}\\
\hspace*{-2mm} \begin{tikzpicture}
	\begin{axis}[\bbar{y}{0}, \dateaxisticks ytick={0, 2, 4, 6, 8, 10}, 
		reverse legend, ymin=0.15,
		height=4.9cm, width=12.9cm, yticklabel style={text width=1.0em},
		xticklabel={`\short{\year}}, clip=false, legend cell align={left}, 
		legend style={legend columns=1, at={(0.89, 1.0)}}]
	\rbars
	\sbar{violet!85!black}{date}{GOV}{data/z1_nfi.csv}
	\sbar{magenta!90!white}{date}{BUS}{data/z1_nfi.csv}
	\sbar{cyan!40!white!95!white}{date}{HH}{data/z1_nfi.csv}
	\sbar{blue!80!green!95!white}{date}{DG}{data/z1_nfi.csv}
	\legend{Government, Business, Household, Consumer Durables};
	\end{axis}
\end{tikzpicture}\\
\footnotesize{Source: Federal Reserve, Bureau of Economic Analysis} \hfill \tbllink{z1_nfi.csv}

\rowcolors{1}{}{black!5} \setlength{\tabcolsep}{3.5pt} \color{black!90}
		{\renewcommand{\arraystretch}{1.5}
		 \begin{tabular}{p{42mm} R{7.8mm} R{7.8mm} R{7.8mm} R{7.8mm} 
		   R{8.2mm} R{8.2mm} R{8.2mm} }
			 \input data/z1_nfi.tex \hline
		\end{tabular}}	\vspace{-2mm}
		
\footnotesize{Source: Federal Reserve, Bureau of Economic Analysis}
\end{minipage}
\newpage
\hypertarget{hh}{}
\begin{minipage}{0.76\textwidth}
\section*{Households}
\small This section covers the household sector of the economy. Households are the source of labor for production and savings for investment. Households are also the primary consumers in the economy. The core topics in the households section include demographics, personal and household income and outlays, consumer sentiment, residential investment, household balance sheets, home ownership, housing, and poverty.

\hypertarget{hhdem}{}
\subsection*{Demographics}
\vspace{-1.0mm}

\small Demographics provide a foundation for examining the household sector. Demographics provide insight on the structure and characteristics of the population. The demographics subsection covers population, population growth, household formation and headship, age, life expectancy, and education.

\subsubsection*{Population}
\small The Census Bureau provides \href{https://www.census.gov/programs-surveys/popest.html}{estimates} and \href{https://www.census.gov/programs-surveys/popproj.html}{projections} of the \textbf{US population}. Population levels and growth rates affect the economy and are critical pieces of information in determining and evaluating economic policies and outcomes. Population projections are based on assumptions, for example about the future level of net migration to the US, but are useful for thinking about future US demographics.

\input{text/pop_est.txt}
\vspace{1mm}

\normalsize \textbf{Population Estimates and Projections}\\
\footnotesize{\textit{resident population, in millions of people}}\\
\hspace*{-2mm} \begin{tikzpicture}
	\begin{axis}[\bbar{y}{0}, \dateaxisticks xmin=1990-01-01, xmax=2027-07-01,
		xticklabel={`\short{\year}}, clip=false, clip mode=individual,
		height=6.2cm, width=12.8cm, yticklabel style={text width=1.6em}]
	\rbars
	\input{text/pop_est_projbar.txt}
	\addplot[thick, mark=*, color=magenta!80!white] 
		table [x=date, y=PROJ_POP, col sep=comma] {data/pop_est_proj.csv};	
	\addplot[thick, mark=triangle, color=cyan!80!white] 
		table [x=date, y=PROJ_POP, col sep=comma] {data/pop_u65_proj.csv};	
	\thickline{green!70!blue}{date}{TOT_POP}{data/pop_u65.csv}
	\thickline{violet}{date}{TOT_POP}{data/pop_est.csv}
	\absnode{{2010-01-01}}{266}{\small \color{green!70!blue}\textbf{Under Age 65}}
	\absnode{{1993-01-01}}{282}{\small \color{violet}\textbf{Total}}
	\absnode{{2024-12-01}}{{340}}{\footnotesize \color{magenta!80!white}\textbf{2017}\\ 
		\small \color{magenta!80!white}\textbf{Proj.}}
	\absnode{{2022-10-01}}{{275}}{\footnotesize \color{cyan!80!white}\textbf{2017}\\ 
		\small \color{cyan!80!white}\textbf{Proj.}}
	\end{axis}
\end{tikzpicture}\\
\footnotesize{Source: Census Bureau} \hfill \tbllink{pop_est.csv}
\vspace{2mm}

\normalsize \textbf{Population Estimates and Projections}\\
\footnotesize{\textit{resident population, in millions of people \hspace{58mm} Projected}}\\
\noindent \hspace*{-2mm} \rowcolors{1}{}{black!5} \setlength{\tabcolsep}{6.1pt} \color{black!90}
		{\renewcommand{\arraystretch}{1.6}
		 \begin{tabular}{p{31mm} R{7.5mm} R{7.5mm} R{7.5mm} R{7.5mm} R{7.5mm} R{7.5mm} | R{8.8mm}}
			 \input data/pop_table.tex \hline
		\end{tabular}}\vspace{-2mm}
		
\footnotesize{Source: Census Bureau}
\end{minipage}
\newpage
\begin{minipage}{0.76\textwidth}
\subsubsection*{Population Growth}
\small \input{text/pop_comp.txt}
\vspace{1mm}

\normalsize \textbf{Components of US Population Growth}\\
\footnotesize{\textit{in millions of people}}\\
\hspace*{-2mm} \begin{tikzpicture}
	\begin{axis}[\bbar{y}{0}, \shdateaxisticks max space between ticks=20,
		xticklabel={`\short{\year}}, clip=false, clip mode=individual,
		height=5.0cm, width=6.7cm, yticklabel style={text width=1.0em},		
		fill between/on layer={main}]
	\rbars
	\addplot[name path=F, ultra thick, no markers, color=lime!60!green] 
		table [x=date, y=BIRTHS, col sep=comma] {data/popcomp.csv};
	\addplot[name path=G, ultra thick, no markers, color=orange!80!red] 
		table [x=date, y=DEATHS, col sep=comma] {data/popcomp.csv};	
	\addplot[style={pattern=crosshatch, pattern color=cyan!80!blue}]
		fill between[of=F and G];
	\abar{blue!75!black}{date}{NETMIG}{data/popcomp.csv}
	\absnode{{2013-01-01}}{4.55}{\footnotesize \color{green!90!black}\textbf{Births}}
	\absnode{{2013-01-01}}{2.6}{\footnotesize \color{orange!80!red}\textbf{Deaths}}
	\absnode{{1995-01-01}}{3.5}{\footnotesize \color{blue!80!cyan}\textbf{Natural Increase}}
	\absnode{{1997-01-01}}{0.7}{\footnotesize \color{white}\textbf{Net Migration}}
	\end{axis}
\end{tikzpicture} \hfill
\begin{tikzpicture}
	\begin{axis}[\bbar{y}{0}, \shdateaxisticks max space between ticks=30,
		xmin={1988-01-01},
		xticklabel={`\short{\year}}, clip=false, clip mode=individual, ymin=0.1,
		height=5.0cm, width=6.7cm, yticklabel style={text width=1.0em}]
	\rbars
	\addplot[stack plots=y, area style, draw=none, 
		style={pattern=crosshatch, pattern color=cyan!80!blue}] 
			table [x=date, y=NATURALINC, col sep=comma]{data/popcomp.csv}\closedcycle;
	\abar{blue!75!black}{date}{NETMIG}{data/popcomp.csv}
	\absnode{{1995-01-01}}{1.1}{\footnotesize \color{blue!80!cyan}\textbf{Natural Increase}}
	\absnode{{1994-01-01}}{2.3}{\footnotesize \color{white}\textbf{Net Migration}}
	\end{axis}
\end{tikzpicture}\\
\footnotesize{Source: Census Bureau} \hfill \tbllink{popcomp.csv}
\vspace{1mm}

\subsubsection*{Related Measures}
\small There are multiple measures of population, based on different definitions. \input{text/pop_universe.txt}The Bureau of Economic Analysis (BEA) use midyear resident population estimates from the Census Bureau for per capita measures. Chartbook measures that use both population and BEA data use the resident population, while chartbook measures based on the Current Population Survey use the civilian noninstitionalized population. 

The Census Bureau further divides the population into those living in households and those living in group quarters. \input{text/pop_hh.txt} 

Lastly, an important related concept, \textbf{households}, is \href{https://www.census.gov/housing/hvs/index.html}{measured} as occupied housing units. The number of households varies over time, separately from the population, as people make changes in their living arrangements. \input{text/hh_tot.txt} 
\end{minipage}
\newpage
\begin{minipage}{0.76\textwidth}
\subsubsection*{Household Formation}
\small Households are measured as \textbf{occupied housing units}, whether occupied by the owner or rented. \input{text/hhform_lvl.txt}Since 1989, the US has experienced the boom and bust of a major housing bubble. By 2016, the end result of the bubble bursting was a shortage of housing, as housing units per capita fell from 1995 to 2016. 
\vspace{1mm}
 
\normalsize \textbf{Housing by Type}\\
\footnotesize{\textit{one-year moving averages}}\\
\footnotesize{\textit{millions of units \hspace{3.2cm} change since 1989, units per capita, percentage points}}\\
\hspace*{-2mm} \begin{tikzpicture}
    \begin{groupplot}[group style={group size=2 by 1, horizontal sep=44pt,}]
    \nextgroupplot[\shdateaxisticks \bbar{y}{0}, 
		enlarge y limits={0.02}, ymin=1, yticklabel style={text width=1.0em},
		xticklabel={`\short{\year}}, height=5.0cm, width=5.0cm, clip=false]
    \rbars 
	\thickline{teal!68!green!95!black}{date}{Owner}{data/hhform_lvl.csv}
	\absnode{{2010-01-01}}{74.5}{\footnotesize \color{teal!68!green!95!black}Owned}
	\thickline{cyan!45!white}{date}{Renter}{data/hhform_lvl.csv}
	\absnode{{2010-01-01}}{55}{\footnotesize \color{cyan!65!white}Rented}
	\thickline{orange!80!yellow}{date}{Vacant}{data/hhform_lvl.csv}
	\absnode{{2010-01-01}}{16.6}{\footnotesize \color{orange!80!yellow}Vacant}
	\input{text/hhform_lvl_nodes.txt}
    \nextgroupplot[\bbar{y}{0}, \shdateaxisticks 
		enlarge y limits={0.12}, ytick={-2, -1, 0, 1, 2, 3, 4},
		xticklabel={`\short{\year}}, height=5.0cm, width=7.5cm, clip=false]
	\rbars
	\stdline{teal!68!green!95!black}{date}{Owner}{data/hhform_idx.csv}
	\absnode{{1993-09-01}}{1.65}{\footnotesize \color{teal!68!green!95!black}Owned}
	\stdline{cyan!45!white}{date}{Renter}{data/hhform_idx.csv}
	\absnode{{1993-09-01}}{-1.25}{\footnotesize \color{cyan!65!white}Rented}
	\stdline{orange!80!yellow}{date}{Vacant}{data/hhform_idx.csv}
	\absnode{{2011-09-01}}{1.7}{\footnotesize \color{orange!80!yellow}Vacant}
	\input{text/hhform_idx_nodes.txt}
	\end{groupplot}
	\end{tikzpicture}\\
\footnotesize{Source: Census Bureau} \hfill \tbllink{hhform_lvl.csv} \ \tbllink{hhform_idx.csv} 
\vspace{3mm}

\small \textbf{Household formation} measures the change in occupied housing units. During the housing bubble, housing construction exceeded population growth and the homeownership rate increased. Following the collapse of the housing bubble, household formation was often below population growth and homeownership decreased as foreclosures converted homeowners into renters.

\input{text/hhform.txt}
\vspace{1mm}

\normalsize \textbf{Contributions to Household Formation}\\
\footnotesize{\textit{one-year moving average of annual growth rates, percent}}
\vspace{3.8cm}

\hspace{3mm} \begin{tikzpicture}[overlay]
	\begin{axis}[\bbar{y}{0}, \dateaxisticks ytick={-2, -1, 0, 1, 2, 3, 4},
		xticklabel={`\short{\year}}, clip=false, height=5.8cm, 
		legend cell align=left, reverse legend, ymax=2.8,
		legend style={at={(0.38, 0.99)}, legend columns=1}]
	\rbars
	\draw [dashed] (axis cs:{2002-01-01},\pgfkeysvalueof{/pgfplots/ymin}) -- (axis cs:{2002-01-01},
		\pgfkeysvalueof{/pgfplots/ymax});
	\absnode{{1997-07-01}}{-0.95}{\scriptsize historical}
	\absnode{{2002-06-01}}{-0.95}{\scriptsize revised}
	\sbaral{cyan!45!white}{date}{Renter}{data/hhform.csv}
	\sbaral{teal!68!green!95!black}{date}{Owner}{data/hhform.csv}
	\stdline{black}{date}{pop}{data/hhform.csv}
	\legend{Rented, Owned, Population Growth};
	\end{axis}
\end{tikzpicture}

\footnotesize{Source: Census Bureau, Housing Vacancies and Homeownership}  \hfill \tbllink{hhform.csv}
\end{minipage}
\newpage
\begin{minipage}{0.76\textwidth}
\subsubsection*{Headship Rate}
\small Individual decisions about starting a household or living with family are influenced by economic conditions. The ratio of households to people age 16 or older is referred to as the \href{https://www.federalreserve.gov/econres/notes/feds-notes/an-early-evaluation-of-the-effects-of-the-pandemic-on-living-arrangements-and-household-formation-20200807.htm}{aggregate headship} rate. 
\end{minipage}

\begin{minipage}{0.25\textwidth}
\small The headship rate fell following the collapse of the housing bubble and during the COVID-19 pandemic, as more people moved in with family. \input{text/headship.txt}
\end{minipage}\hspace{6mm}
\begin{minipage}{0.47\textwidth}
\normalsize \textbf{Aggregate Headship Rate}\\
\footnotesize{\textit{people age 16 or older, percent, seasonally adjusted}}\\
\hspace*{-2mm} \begin{tikzpicture}
	\begin{axis}[\shdateaxisticks ytick={49, 50, 51}, 
		enlarge y limits={0.1}, yticklabel style={text width=1.0em},
		xticklabel={`\short{\year}}, height=4.0cm, width=7.6cm, clip=false]
	\rebars
	\stdline{purple!70!violet}{date}{value}{data/headship.csv}
	\input{text/headship_node.txt}
	\end{axis}
\end{tikzpicture}\\
\footnotesize{Source: Author's Calculations from CPS} \hfill \tbllink{headship.csv} 
\end{minipage}
\newpage
\begin{minipage}{0.76\textwidth}
\subsubsection*{Age}
\small In discussions on demographics, \textbf{aging} is often described as a serious headwind to economic growth in major advanced economies. The increased share of many countries' population that is of retirement age means a smaller share are working and borrowing and a larger share are receiving pension benefits and lending to the financial system. These trends can be overcome by a workforce that is more efficiently able to provide goods and services. In part due to a shorter life-expectancy in the US, this problem is more pronounced in Japan and western Europe, but is still an important issue for the US.

\begin{minipage}{0.41\textwidth}
\small \input{text/med_age.txt}
\end{minipage}\hfill
\begin{minipage}{0.53\textwidth}
\normalsize \textbf{Median Age}\\
\footnotesize{\textit{by group, in years}}\\
\noindent \hspace*{-2mm} \begin{tikzpicture}
	\begin{axis}[\bbar{y}{0}, \shdateaxisticks  
		xticklabel={`\short{\year}}, clip=false, 
		yticklabel style={text width=1.2em}, xmin={1986-10-01},
		height=5.8cm, width=6.6cm, enlarge y limits={0.15}]
	\rbars
	\stdline{blue!60!cyan}{date}{ALL}{data/median_age.csv}
	\absnode{{1992-01-01}}{32}{\footnotesize \color{blue!60!cyan}All Ages}
	\stdline{violet}{date}{AGE16PLUS}{data/median_age.csv}
	\absnode{{1992-01-01}}{43.5}{\footnotesize \color{violet}Age\\\footnotesize \color{violet}16+}
	\stdline{red}{date}{EMP}{data/median_age.csv}
	\absnode{{2010-07-01}}{43.5}{\footnotesize \color{red}Employed}
	\input{text/med_age_node.txt}
	\end{axis}
\end{tikzpicture}\\
\footnotesize{Source: Author's Calculations from CPS} \hfill \tbllink{median_age.csv}
\end{minipage}
\vspace{3mm}

\small Economic indicators are sometimes based on specific \textbf{age groups}. As examples, labor statistics often exclude those under age 16, the retirement-age population is above a certain age, such as 64, and a popular measure of labor market slack is the age 25--54 employment rate. It is therefore useful to know what share of the overall population is in each major age group, and how the current age distribution compares with the past. 
\end{minipage}
\vspace{0.5mm}

\begin{minipage}{0.35\textwidth}
\normalsize \textbf{Age Groups}\\
\footnotesize{\textit{share of population, percent}}\\
\hspace*{-4mm} \begin{tikzpicture}
  	\begin{axis}[\barplotnogrid axis y line=left, \barylab{1.1cm}{1.5ex}
    	width=5.6cm, bar width=1.8ex, height=8.2cm,
    	enlarge y limits={abs=0.35cm}, enlarge x limits=0.05, \bbar{x}{0},
		yticklabels from table={\agesh}{name}, xmin=0, clip=false,
		yticklabel style={font=\footnotesize}, nodes near coords, 
		nodes near coords style={/pgf/number format/.cd, fixed zerofill,
			precision=1, assume math mode},
		legend style={text=black!70, at={(0.29, 1.08)}, anchor=north, legend columns=-1, 
				fill=none, draw=none,
		        /tikz/every even column/.append style={column sep=0.4cm}}]
  	\addplot[fill=cyan!90!blue!55!white, draw=none] 
  		table [y expr=-\coordindex, x index=2] {\agesh};
  	\addplot[fill=blue!75!violet!90!white, draw=none] 
  		table [y expr=-\coordindex, x index=1] {\agesh};
 	\draw [black!35!white, dashed] (-12, -6.5) -- 
 		(\pgfkeysvalueof{/pgfplots/xmax}, -6.5);
  	\legend{1989, \input{text/cps_age_dt.txt}}
  	\end{axis}
\end{tikzpicture}\\
\footnotesize{Source: Author's Calculations from CPS}
\end{minipage} \hspace{4.5mm}
\begin{minipage}{0.37\textwidth}
\small \input{text/cps_age.txt}
\end{minipage}
\newpage
\begin{minipage}{0.76\textwidth}
\small Mapping American Community Survey data to commuter zones gives insight on the \textbf{age of local areas}. \input{text/cps_u18.txt} 

\input{text/cps_o64.txt}

\end{minipage}
\vspace{1mm}

\noindent \normalsize \textbf{Age Group Share of Commuter Zone Population, \input{text/acs_ftfy_year.txt}}

\hspace{2mm} Age 0 to 17 \hspace{62mm} Age 65+

\vspace{-3mm}
\hspace{-9mm} %% Creator: Matplotlib, PGF backend
%%
%% To include the figure in your LaTeX document, write
%%   \input{<filename>.pgf}
%%
%% Make sure the required packages are loaded in your preamble
%%   \usepackage{pgf}
%%
%% Also ensure that all the required font packages are loaded; for instance,
%% the lmodern package is sometimes necessary when using math font.
%%   \usepackage{lmodern}
%%
%% Figures using additional raster images can only be included by \input if
%% they are in the same directory as the main LaTeX file. For loading figures
%% from other directories you can use the `import` package
%%   \usepackage{import}
%%
%% and then include the figures with
%%   \import{<path to file>}{<filename>.pgf}
%%
%% Matplotlib used the following preamble
%%   
%%   \usepackage{fontspec}
%%   \setmainfont{DejaVuSerif.ttf}[Path=\detokenize{/home/brian/miniconda3/lib/python3.8/site-packages/matplotlib/mpl-data/fonts/ttf/}]
%%   \setsansfont{DejaVuSans.ttf}[Path=\detokenize{/home/brian/miniconda3/lib/python3.8/site-packages/matplotlib/mpl-data/fonts/ttf/}]
%%   \setmonofont{DejaVuSansMono.ttf}[Path=\detokenize{/home/brian/miniconda3/lib/python3.8/site-packages/matplotlib/mpl-data/fonts/ttf/}]
%%   \makeatletter\@ifpackageloaded{underscore}{}{\usepackage[strings]{underscore}}\makeatother
%%
\begingroup%
\makeatletter%
\begin{pgfpicture}%
\pgfpathrectangle{\pgfpointorigin}{\pgfqpoint{3.347362in}{2.125000in}}%
\pgfusepath{use as bounding box, clip}%
\begin{pgfscope}%
\pgfsetbuttcap%
\pgfsetmiterjoin%
\pgfsetlinewidth{0.000000pt}%
\definecolor{currentstroke}{rgb}{0.000000,0.000000,0.000000}%
\pgfsetstrokecolor{currentstroke}%
\pgfsetstrokeopacity{0.000000}%
\pgfsetdash{}{0pt}%
\pgfpathmoveto{\pgfqpoint{0.000000in}{0.000000in}}%
\pgfpathlineto{\pgfqpoint{3.347362in}{0.000000in}}%
\pgfpathlineto{\pgfqpoint{3.347362in}{2.125000in}}%
\pgfpathlineto{\pgfqpoint{0.000000in}{2.125000in}}%
\pgfpathlineto{\pgfqpoint{0.000000in}{0.000000in}}%
\pgfpathclose%
\pgfusepath{}%
\end{pgfscope}%
\begin{pgfscope}%
\pgfpathrectangle{\pgfqpoint{0.100000in}{0.100000in}}{\pgfqpoint{3.007045in}{1.925000in}}%
\pgfusepath{clip}%
\pgfsetbuttcap%
\pgfsetmiterjoin%
\pgfsetlinewidth{0.000000pt}%
\definecolor{currentstroke}{rgb}{0.000000,0.000000,0.000000}%
\pgfsetstrokecolor{currentstroke}%
\pgfsetstrokeopacity{0.000000}%
\pgfsetdash{}{0pt}%
\pgfpathmoveto{\pgfqpoint{0.100000in}{0.100000in}}%
\pgfpathlineto{\pgfqpoint{3.107045in}{0.100000in}}%
\pgfpathlineto{\pgfqpoint{3.107045in}{2.025000in}}%
\pgfpathlineto{\pgfqpoint{0.100000in}{2.025000in}}%
\pgfpathlineto{\pgfqpoint{0.100000in}{0.100000in}}%
\pgfpathclose%
\pgfusepath{}%
\end{pgfscope}%
\begin{pgfscope}%
\pgfpathrectangle{\pgfqpoint{0.100000in}{0.100000in}}{\pgfqpoint{3.007045in}{1.925000in}}%
\pgfusepath{clip}%
\pgfsetbuttcap%
\pgfsetmiterjoin%
\definecolor{currentfill}{rgb}{0.366644,0.646182,0.818547}%
\pgfsetfillcolor{currentfill}%
\pgfsetlinewidth{0.000000pt}%
\definecolor{currentstroke}{rgb}{0.000000,0.000000,0.000000}%
\pgfsetstrokecolor{currentstroke}%
\pgfsetstrokeopacity{0.000000}%
\pgfsetdash{}{0pt}%
\pgfpathmoveto{\pgfqpoint{1.787185in}{1.275241in}}%
\pgfpathlineto{\pgfqpoint{1.764595in}{1.275306in}}%
\pgfpathlineto{\pgfqpoint{1.764683in}{1.298101in}}%
\pgfpathlineto{\pgfqpoint{1.747728in}{1.298216in}}%
\pgfpathlineto{\pgfqpoint{1.744815in}{1.304579in}}%
\pgfpathlineto{\pgfqpoint{1.744960in}{1.321802in}}%
\pgfpathlineto{\pgfqpoint{1.784860in}{1.321668in}}%
\pgfpathlineto{\pgfqpoint{1.784942in}{1.304507in}}%
\pgfpathlineto{\pgfqpoint{1.787084in}{1.304515in}}%
\pgfpathlineto{\pgfqpoint{1.787185in}{1.275241in}}%
\pgfpathclose%
\pgfusepath{fill}%
\end{pgfscope}%
\begin{pgfscope}%
\pgfpathrectangle{\pgfqpoint{0.100000in}{0.100000in}}{\pgfqpoint{3.007045in}{1.925000in}}%
\pgfusepath{clip}%
\pgfsetbuttcap%
\pgfsetmiterjoin%
\definecolor{currentfill}{rgb}{0.401953,0.670296,0.832326}%
\pgfsetfillcolor{currentfill}%
\pgfsetlinewidth{0.000000pt}%
\definecolor{currentstroke}{rgb}{0.000000,0.000000,0.000000}%
\pgfsetstrokecolor{currentstroke}%
\pgfsetstrokeopacity{0.000000}%
\pgfsetdash{}{0pt}%
\pgfpathmoveto{\pgfqpoint{2.287383in}{1.019436in}}%
\pgfpathlineto{\pgfqpoint{2.292665in}{1.018214in}}%
\pgfpathlineto{\pgfqpoint{2.306270in}{1.030998in}}%
\pgfpathlineto{\pgfqpoint{2.302206in}{1.033847in}}%
\pgfpathlineto{\pgfqpoint{2.301947in}{1.038398in}}%
\pgfpathlineto{\pgfqpoint{2.310463in}{1.047437in}}%
\pgfpathlineto{\pgfqpoint{2.318137in}{1.056117in}}%
\pgfpathlineto{\pgfqpoint{2.321129in}{1.054650in}}%
\pgfpathlineto{\pgfqpoint{2.333634in}{1.044388in}}%
\pgfpathlineto{\pgfqpoint{2.332415in}{1.038665in}}%
\pgfpathlineto{\pgfqpoint{2.334697in}{1.028538in}}%
\pgfpathlineto{\pgfqpoint{2.332746in}{1.021710in}}%
\pgfpathlineto{\pgfqpoint{2.336554in}{1.012826in}}%
\pgfpathlineto{\pgfqpoint{2.341210in}{1.008147in}}%
\pgfpathlineto{\pgfqpoint{2.340163in}{1.003243in}}%
\pgfpathlineto{\pgfqpoint{2.336026in}{1.001374in}}%
\pgfpathlineto{\pgfqpoint{2.338194in}{0.987373in}}%
\pgfpathlineto{\pgfqpoint{2.334580in}{0.979199in}}%
\pgfpathlineto{\pgfqpoint{2.325950in}{0.982785in}}%
\pgfpathlineto{\pgfqpoint{2.317969in}{0.990753in}}%
\pgfpathlineto{\pgfqpoint{2.319984in}{0.992600in}}%
\pgfpathlineto{\pgfqpoint{2.315188in}{1.000641in}}%
\pgfpathlineto{\pgfqpoint{2.311789in}{1.005478in}}%
\pgfpathlineto{\pgfqpoint{2.285799in}{1.004151in}}%
\pgfpathlineto{\pgfqpoint{2.287383in}{1.019436in}}%
\pgfpathclose%
\pgfusepath{fill}%
\end{pgfscope}%
\begin{pgfscope}%
\pgfpathrectangle{\pgfqpoint{0.100000in}{0.100000in}}{\pgfqpoint{3.007045in}{1.925000in}}%
\pgfusepath{clip}%
\pgfsetbuttcap%
\pgfsetmiterjoin%
\definecolor{currentfill}{rgb}{0.166967,0.480692,0.729150}%
\pgfsetfillcolor{currentfill}%
\pgfsetlinewidth{0.000000pt}%
\definecolor{currentstroke}{rgb}{0.000000,0.000000,0.000000}%
\pgfsetstrokecolor{currentstroke}%
\pgfsetstrokeopacity{0.000000}%
\pgfsetdash{}{0pt}%
\pgfpathmoveto{\pgfqpoint{1.471341in}{0.833174in}}%
\pgfpathlineto{\pgfqpoint{1.469347in}{0.801604in}}%
\pgfpathlineto{\pgfqpoint{1.440696in}{0.803294in}}%
\pgfpathlineto{\pgfqpoint{1.442781in}{0.834883in}}%
\pgfpathlineto{\pgfqpoint{1.471341in}{0.833174in}}%
\pgfpathclose%
\pgfusepath{fill}%
\end{pgfscope}%
\begin{pgfscope}%
\pgfpathrectangle{\pgfqpoint{0.100000in}{0.100000in}}{\pgfqpoint{3.007045in}{1.925000in}}%
\pgfusepath{clip}%
\pgfsetbuttcap%
\pgfsetmiterjoin%
\definecolor{currentfill}{rgb}{0.417086,0.680631,0.838231}%
\pgfsetfillcolor{currentfill}%
\pgfsetlinewidth{0.000000pt}%
\definecolor{currentstroke}{rgb}{0.000000,0.000000,0.000000}%
\pgfsetstrokecolor{currentstroke}%
\pgfsetstrokeopacity{0.000000}%
\pgfsetdash{}{0pt}%
\pgfpathmoveto{\pgfqpoint{1.408374in}{1.098989in}}%
\pgfpathlineto{\pgfqpoint{1.408351in}{1.098572in}}%
\pgfpathlineto{\pgfqpoint{1.405795in}{1.064151in}}%
\pgfpathlineto{\pgfqpoint{1.405343in}{1.057957in}}%
\pgfpathlineto{\pgfqpoint{1.351564in}{1.062383in}}%
\pgfpathlineto{\pgfqpoint{1.300552in}{1.067248in}}%
\pgfpathlineto{\pgfqpoint{1.301086in}{1.073162in}}%
\pgfpathlineto{\pgfqpoint{1.271819in}{1.087203in}}%
\pgfpathlineto{\pgfqpoint{1.252696in}{1.088003in}}%
\pgfpathlineto{\pgfqpoint{1.251074in}{1.090450in}}%
\pgfpathlineto{\pgfqpoint{1.253600in}{1.112871in}}%
\pgfpathlineto{\pgfqpoint{1.259199in}{1.112273in}}%
\pgfpathlineto{\pgfqpoint{1.261017in}{1.129389in}}%
\pgfpathlineto{\pgfqpoint{1.306529in}{1.124802in}}%
\pgfpathlineto{\pgfqpoint{1.334664in}{1.121680in}}%
\pgfpathlineto{\pgfqpoint{1.333267in}{1.105178in}}%
\pgfpathlineto{\pgfqpoint{1.372398in}{1.101802in}}%
\pgfpathlineto{\pgfqpoint{1.408374in}{1.098989in}}%
\pgfpathclose%
\pgfusepath{fill}%
\end{pgfscope}%
\begin{pgfscope}%
\pgfpathrectangle{\pgfqpoint{0.100000in}{0.100000in}}{\pgfqpoint{3.007045in}{1.925000in}}%
\pgfusepath{clip}%
\pgfsetbuttcap%
\pgfsetmiterjoin%
\definecolor{currentfill}{rgb}{0.381776,0.656517,0.824452}%
\pgfsetfillcolor{currentfill}%
\pgfsetlinewidth{0.000000pt}%
\definecolor{currentstroke}{rgb}{0.000000,0.000000,0.000000}%
\pgfsetstrokecolor{currentstroke}%
\pgfsetstrokeopacity{0.000000}%
\pgfsetdash{}{0pt}%
\pgfpathmoveto{\pgfqpoint{1.769446in}{1.131864in}}%
\pgfpathlineto{\pgfqpoint{1.764032in}{1.131806in}}%
\pgfpathlineto{\pgfqpoint{1.763130in}{1.141541in}}%
\pgfpathlineto{\pgfqpoint{1.763181in}{1.160646in}}%
\pgfpathlineto{\pgfqpoint{1.743298in}{1.160738in}}%
\pgfpathlineto{\pgfqpoint{1.743716in}{1.176139in}}%
\pgfpathlineto{\pgfqpoint{1.755095in}{1.176082in}}%
\pgfpathlineto{\pgfqpoint{1.766486in}{1.176026in}}%
\pgfpathlineto{\pgfqpoint{1.769538in}{1.174111in}}%
\pgfpathlineto{\pgfqpoint{1.767101in}{1.168173in}}%
\pgfpathlineto{\pgfqpoint{1.772963in}{1.162326in}}%
\pgfpathlineto{\pgfqpoint{1.780164in}{1.150731in}}%
\pgfpathlineto{\pgfqpoint{1.777411in}{1.146354in}}%
\pgfpathlineto{\pgfqpoint{1.776978in}{1.132450in}}%
\pgfpathlineto{\pgfqpoint{1.769446in}{1.131864in}}%
\pgfpathclose%
\pgfusepath{fill}%
\end{pgfscope}%
\begin{pgfscope}%
\pgfpathrectangle{\pgfqpoint{0.100000in}{0.100000in}}{\pgfqpoint{3.007045in}{1.925000in}}%
\pgfusepath{clip}%
\pgfsetbuttcap%
\pgfsetmiterjoin%
\definecolor{currentfill}{rgb}{0.346467,0.632403,0.810673}%
\pgfsetfillcolor{currentfill}%
\pgfsetlinewidth{0.000000pt}%
\definecolor{currentstroke}{rgb}{0.000000,0.000000,0.000000}%
\pgfsetstrokecolor{currentstroke}%
\pgfsetstrokeopacity{0.000000}%
\pgfsetdash{}{0pt}%
\pgfpathmoveto{\pgfqpoint{2.317320in}{0.826818in}}%
\pgfpathlineto{\pgfqpoint{2.298990in}{0.824653in}}%
\pgfpathlineto{\pgfqpoint{2.297886in}{0.828687in}}%
\pgfpathlineto{\pgfqpoint{2.290820in}{0.827401in}}%
\pgfpathlineto{\pgfqpoint{2.279304in}{0.827039in}}%
\pgfpathlineto{\pgfqpoint{2.271751in}{0.844592in}}%
\pgfpathlineto{\pgfqpoint{2.275437in}{0.845959in}}%
\pgfpathlineto{\pgfqpoint{2.282179in}{0.853795in}}%
\pgfpathlineto{\pgfqpoint{2.284621in}{0.861875in}}%
\pgfpathlineto{\pgfqpoint{2.287764in}{0.864843in}}%
\pgfpathlineto{\pgfqpoint{2.286587in}{0.869005in}}%
\pgfpathlineto{\pgfqpoint{2.311838in}{0.871613in}}%
\pgfpathlineto{\pgfqpoint{2.311402in}{0.880390in}}%
\pgfpathlineto{\pgfqpoint{2.305031in}{0.879714in}}%
\pgfpathlineto{\pgfqpoint{2.306322in}{0.886459in}}%
\pgfpathlineto{\pgfqpoint{2.310101in}{0.886859in}}%
\pgfpathlineto{\pgfqpoint{2.313442in}{0.898402in}}%
\pgfpathlineto{\pgfqpoint{2.332547in}{0.900599in}}%
\pgfpathlineto{\pgfqpoint{2.333466in}{0.891984in}}%
\pgfpathlineto{\pgfqpoint{2.335741in}{0.890096in}}%
\pgfpathlineto{\pgfqpoint{2.332579in}{0.883358in}}%
\pgfpathlineto{\pgfqpoint{2.334077in}{0.871563in}}%
\pgfpathlineto{\pgfqpoint{2.337348in}{0.840608in}}%
\pgfpathlineto{\pgfqpoint{2.322977in}{0.839287in}}%
\pgfpathlineto{\pgfqpoint{2.320886in}{0.830332in}}%
\pgfpathlineto{\pgfqpoint{2.317320in}{0.826818in}}%
\pgfpathclose%
\pgfusepath{fill}%
\end{pgfscope}%
\begin{pgfscope}%
\pgfpathrectangle{\pgfqpoint{0.100000in}{0.100000in}}{\pgfqpoint{3.007045in}{1.925000in}}%
\pgfusepath{clip}%
\pgfsetbuttcap%
\pgfsetmiterjoin%
\definecolor{currentfill}{rgb}{0.244106,0.557832,0.768889}%
\pgfsetfillcolor{currentfill}%
\pgfsetlinewidth{0.000000pt}%
\definecolor{currentstroke}{rgb}{0.000000,0.000000,0.000000}%
\pgfsetstrokecolor{currentstroke}%
\pgfsetstrokeopacity{0.000000}%
\pgfsetdash{}{0pt}%
\pgfpathmoveto{\pgfqpoint{1.447246in}{1.572779in}}%
\pgfpathlineto{\pgfqpoint{1.445034in}{1.544102in}}%
\pgfpathlineto{\pgfqpoint{1.400379in}{1.547840in}}%
\pgfpathlineto{\pgfqpoint{1.401374in}{1.559284in}}%
\pgfpathlineto{\pgfqpoint{1.406341in}{1.607594in}}%
\pgfpathlineto{\pgfqpoint{1.403892in}{1.607820in}}%
\pgfpathlineto{\pgfqpoint{1.405844in}{1.630008in}}%
\pgfpathlineto{\pgfqpoint{1.409091in}{1.629777in}}%
\pgfpathlineto{\pgfqpoint{1.410950in}{1.652859in}}%
\pgfpathlineto{\pgfqpoint{1.448770in}{1.649792in}}%
\pgfpathlineto{\pgfqpoint{1.450846in}{1.649617in}}%
\pgfpathlineto{\pgfqpoint{1.449255in}{1.626605in}}%
\pgfpathlineto{\pgfqpoint{1.451489in}{1.626431in}}%
\pgfpathlineto{\pgfqpoint{1.450280in}{1.611196in}}%
\pgfpathlineto{\pgfqpoint{1.447246in}{1.572779in}}%
\pgfpathclose%
\pgfusepath{fill}%
\end{pgfscope}%
\begin{pgfscope}%
\pgfpathrectangle{\pgfqpoint{0.100000in}{0.100000in}}{\pgfqpoint{3.007045in}{1.925000in}}%
\pgfusepath{clip}%
\pgfsetbuttcap%
\pgfsetmiterjoin%
\definecolor{currentfill}{rgb}{0.730950,0.839477,0.921323}%
\pgfsetfillcolor{currentfill}%
\pgfsetlinewidth{0.000000pt}%
\definecolor{currentstroke}{rgb}{0.000000,0.000000,0.000000}%
\pgfsetstrokecolor{currentstroke}%
\pgfsetstrokeopacity{0.000000}%
\pgfsetdash{}{0pt}%
\pgfpathmoveto{\pgfqpoint{2.765291in}{1.169198in}}%
\pgfpathlineto{\pgfqpoint{2.770616in}{1.173606in}}%
\pgfpathlineto{\pgfqpoint{2.789514in}{1.179938in}}%
\pgfpathlineto{\pgfqpoint{2.785420in}{1.167710in}}%
\pgfpathlineto{\pgfqpoint{2.782226in}{1.168112in}}%
\pgfpathlineto{\pgfqpoint{2.778730in}{1.161873in}}%
\pgfpathlineto{\pgfqpoint{2.776737in}{1.148749in}}%
\pgfpathlineto{\pgfqpoint{2.778073in}{1.146484in}}%
\pgfpathlineto{\pgfqpoint{2.774225in}{1.130181in}}%
\pgfpathlineto{\pgfqpoint{2.772699in}{1.120440in}}%
\pgfpathlineto{\pgfqpoint{2.768773in}{1.114106in}}%
\pgfpathlineto{\pgfqpoint{2.765038in}{1.113269in}}%
\pgfpathlineto{\pgfqpoint{2.760275in}{1.121718in}}%
\pgfpathlineto{\pgfqpoint{2.760413in}{1.142327in}}%
\pgfpathlineto{\pgfqpoint{2.763321in}{1.150894in}}%
\pgfpathlineto{\pgfqpoint{2.763523in}{1.158416in}}%
\pgfpathlineto{\pgfqpoint{2.768180in}{1.159798in}}%
\pgfpathlineto{\pgfqpoint{2.769212in}{1.165623in}}%
\pgfpathlineto{\pgfqpoint{2.765291in}{1.169198in}}%
\pgfpathclose%
\pgfusepath{fill}%
\end{pgfscope}%
\begin{pgfscope}%
\pgfpathrectangle{\pgfqpoint{0.100000in}{0.100000in}}{\pgfqpoint{3.007045in}{1.925000in}}%
\pgfusepath{clip}%
\pgfsetbuttcap%
\pgfsetmiterjoin%
\definecolor{currentfill}{rgb}{0.627605,0.795556,0.885152}%
\pgfsetfillcolor{currentfill}%
\pgfsetlinewidth{0.000000pt}%
\definecolor{currentstroke}{rgb}{0.000000,0.000000,0.000000}%
\pgfsetstrokecolor{currentstroke}%
\pgfsetstrokeopacity{0.000000}%
\pgfsetdash{}{0pt}%
\pgfpathmoveto{\pgfqpoint{2.090405in}{1.513286in}}%
\pgfpathlineto{\pgfqpoint{2.073413in}{1.512435in}}%
\pgfpathlineto{\pgfqpoint{2.073048in}{1.518236in}}%
\pgfpathlineto{\pgfqpoint{2.043986in}{1.516491in}}%
\pgfpathlineto{\pgfqpoint{2.042126in}{1.545379in}}%
\pgfpathlineto{\pgfqpoint{2.032705in}{1.544832in}}%
\pgfpathlineto{\pgfqpoint{2.031351in}{1.567858in}}%
\pgfpathlineto{\pgfqpoint{2.054639in}{1.569212in}}%
\pgfpathlineto{\pgfqpoint{2.054884in}{1.563482in}}%
\pgfpathlineto{\pgfqpoint{2.066412in}{1.564272in}}%
\pgfpathlineto{\pgfqpoint{2.067108in}{1.552783in}}%
\pgfpathlineto{\pgfqpoint{2.069326in}{1.547117in}}%
\pgfpathlineto{\pgfqpoint{2.076668in}{1.547607in}}%
\pgfpathlineto{\pgfqpoint{2.077656in}{1.530311in}}%
\pgfpathlineto{\pgfqpoint{2.088935in}{1.530951in}}%
\pgfpathlineto{\pgfqpoint{2.090405in}{1.513286in}}%
\pgfpathclose%
\pgfusepath{fill}%
\end{pgfscope}%
\begin{pgfscope}%
\pgfpathrectangle{\pgfqpoint{0.100000in}{0.100000in}}{\pgfqpoint{3.007045in}{1.925000in}}%
\pgfusepath{clip}%
\pgfsetbuttcap%
\pgfsetmiterjoin%
\definecolor{currentfill}{rgb}{0.825928,0.891795,0.953741}%
\pgfsetfillcolor{currentfill}%
\pgfsetlinewidth{0.000000pt}%
\definecolor{currentstroke}{rgb}{0.000000,0.000000,0.000000}%
\pgfsetstrokecolor{currentstroke}%
\pgfsetstrokeopacity{0.000000}%
\pgfsetdash{}{0pt}%
\pgfpathmoveto{\pgfqpoint{2.836904in}{1.462520in}}%
\pgfpathlineto{\pgfqpoint{2.816097in}{1.458112in}}%
\pgfpathlineto{\pgfqpoint{2.815641in}{1.458005in}}%
\pgfpathlineto{\pgfqpoint{2.814539in}{1.460221in}}%
\pgfpathlineto{\pgfqpoint{2.815796in}{1.505313in}}%
\pgfpathlineto{\pgfqpoint{2.813736in}{1.511035in}}%
\pgfpathlineto{\pgfqpoint{2.807383in}{1.541913in}}%
\pgfpathlineto{\pgfqpoint{2.825651in}{1.545463in}}%
\pgfpathlineto{\pgfqpoint{2.828637in}{1.543127in}}%
\pgfpathlineto{\pgfqpoint{2.828599in}{1.533431in}}%
\pgfpathlineto{\pgfqpoint{2.822507in}{1.532326in}}%
\pgfpathlineto{\pgfqpoint{2.824976in}{1.520081in}}%
\pgfpathlineto{\pgfqpoint{2.829097in}{1.520876in}}%
\pgfpathlineto{\pgfqpoint{2.831656in}{1.508759in}}%
\pgfpathlineto{\pgfqpoint{2.827848in}{1.505323in}}%
\pgfpathlineto{\pgfqpoint{2.832224in}{1.502258in}}%
\pgfpathlineto{\pgfqpoint{2.833508in}{1.484777in}}%
\pgfpathlineto{\pgfqpoint{2.836904in}{1.462520in}}%
\pgfpathclose%
\pgfusepath{fill}%
\end{pgfscope}%
\begin{pgfscope}%
\pgfpathrectangle{\pgfqpoint{0.100000in}{0.100000in}}{\pgfqpoint{3.007045in}{1.925000in}}%
\pgfusepath{clip}%
\pgfsetbuttcap%
\pgfsetmiterjoin%
\definecolor{currentfill}{rgb}{0.516863,0.735748,0.860192}%
\pgfsetfillcolor{currentfill}%
\pgfsetlinewidth{0.000000pt}%
\definecolor{currentstroke}{rgb}{0.000000,0.000000,0.000000}%
\pgfsetstrokecolor{currentstroke}%
\pgfsetstrokeopacity{0.000000}%
\pgfsetdash{}{0pt}%
\pgfpathmoveto{\pgfqpoint{2.429643in}{1.109432in}}%
\pgfpathlineto{\pgfqpoint{2.425369in}{1.103463in}}%
\pgfpathlineto{\pgfqpoint{2.420050in}{1.106990in}}%
\pgfpathlineto{\pgfqpoint{2.421148in}{1.110912in}}%
\pgfpathlineto{\pgfqpoint{2.410869in}{1.119825in}}%
\pgfpathlineto{\pgfqpoint{2.411679in}{1.127319in}}%
\pgfpathlineto{\pgfqpoint{2.401483in}{1.125255in}}%
\pgfpathlineto{\pgfqpoint{2.396061in}{1.120463in}}%
\pgfpathlineto{\pgfqpoint{2.398193in}{1.116290in}}%
\pgfpathlineto{\pgfqpoint{2.392249in}{1.108544in}}%
\pgfpathlineto{\pgfqpoint{2.383396in}{1.107611in}}%
\pgfpathlineto{\pgfqpoint{2.379152in}{1.110286in}}%
\pgfpathlineto{\pgfqpoint{2.370833in}{1.107548in}}%
\pgfpathlineto{\pgfqpoint{2.359204in}{1.116262in}}%
\pgfpathlineto{\pgfqpoint{2.365037in}{1.122682in}}%
\pgfpathlineto{\pgfqpoint{2.364836in}{1.127691in}}%
\pgfpathlineto{\pgfqpoint{2.368680in}{1.131320in}}%
\pgfpathlineto{\pgfqpoint{2.372393in}{1.126861in}}%
\pgfpathlineto{\pgfqpoint{2.378692in}{1.128851in}}%
\pgfpathlineto{\pgfqpoint{2.377693in}{1.133822in}}%
\pgfpathlineto{\pgfqpoint{2.379735in}{1.139390in}}%
\pgfpathlineto{\pgfqpoint{2.384672in}{1.146612in}}%
\pgfpathlineto{\pgfqpoint{2.385587in}{1.155405in}}%
\pgfpathlineto{\pgfqpoint{2.392277in}{1.158355in}}%
\pgfpathlineto{\pgfqpoint{2.395723in}{1.148881in}}%
\pgfpathlineto{\pgfqpoint{2.401519in}{1.149007in}}%
\pgfpathlineto{\pgfqpoint{2.404127in}{1.159387in}}%
\pgfpathlineto{\pgfqpoint{2.403876in}{1.166074in}}%
\pgfpathlineto{\pgfqpoint{2.407704in}{1.166293in}}%
\pgfpathlineto{\pgfqpoint{2.407930in}{1.161976in}}%
\pgfpathlineto{\pgfqpoint{2.413032in}{1.162228in}}%
\pgfpathlineto{\pgfqpoint{2.414411in}{1.156435in}}%
\pgfpathlineto{\pgfqpoint{2.420478in}{1.156867in}}%
\pgfpathlineto{\pgfqpoint{2.420941in}{1.150887in}}%
\pgfpathlineto{\pgfqpoint{2.428157in}{1.152335in}}%
\pgfpathlineto{\pgfqpoint{2.437602in}{1.145887in}}%
\pgfpathlineto{\pgfqpoint{2.438948in}{1.139413in}}%
\pgfpathlineto{\pgfqpoint{2.431774in}{1.134957in}}%
\pgfpathlineto{\pgfqpoint{2.428704in}{1.126732in}}%
\pgfpathlineto{\pgfqpoint{2.428858in}{1.122661in}}%
\pgfpathlineto{\pgfqpoint{2.433994in}{1.115719in}}%
\pgfpathlineto{\pgfqpoint{2.429643in}{1.109432in}}%
\pgfpathclose%
\pgfusepath{fill}%
\end{pgfscope}%
\begin{pgfscope}%
\pgfpathrectangle{\pgfqpoint{0.100000in}{0.100000in}}{\pgfqpoint{3.007045in}{1.925000in}}%
\pgfusepath{clip}%
\pgfsetbuttcap%
\pgfsetmiterjoin%
\definecolor{currentfill}{rgb}{0.765398,0.854118,0.933379}%
\pgfsetfillcolor{currentfill}%
\pgfsetlinewidth{0.000000pt}%
\definecolor{currentstroke}{rgb}{0.000000,0.000000,0.000000}%
\pgfsetstrokecolor{currentstroke}%
\pgfsetstrokeopacity{0.000000}%
\pgfsetdash{}{0pt}%
\pgfpathmoveto{\pgfqpoint{1.194470in}{0.906193in}}%
\pgfpathlineto{\pgfqpoint{1.166043in}{0.909544in}}%
\pgfpathlineto{\pgfqpoint{1.171643in}{0.956432in}}%
\pgfpathlineto{\pgfqpoint{1.164202in}{0.962160in}}%
\pgfpathlineto{\pgfqpoint{1.165289in}{0.970829in}}%
\pgfpathlineto{\pgfqpoint{1.173700in}{0.972532in}}%
\pgfpathlineto{\pgfqpoint{1.139908in}{0.976686in}}%
\pgfpathlineto{\pgfqpoint{1.141795in}{0.990914in}}%
\pgfpathlineto{\pgfqpoint{1.102849in}{0.996413in}}%
\pgfpathlineto{\pgfqpoint{1.108498in}{1.034565in}}%
\pgfpathlineto{\pgfqpoint{1.118519in}{1.040536in}}%
\pgfpathlineto{\pgfqpoint{1.120573in}{1.045885in}}%
\pgfpathlineto{\pgfqpoint{1.169817in}{1.038957in}}%
\pgfpathlineto{\pgfqpoint{1.235453in}{1.031177in}}%
\pgfpathlineto{\pgfqpoint{1.235393in}{1.026018in}}%
\pgfpathlineto{\pgfqpoint{1.231839in}{1.013157in}}%
\pgfpathlineto{\pgfqpoint{1.225227in}{1.010762in}}%
\pgfpathlineto{\pgfqpoint{1.225457in}{0.993376in}}%
\pgfpathlineto{\pgfqpoint{1.223675in}{0.983885in}}%
\pgfpathlineto{\pgfqpoint{1.221151in}{0.983550in}}%
\pgfpathlineto{\pgfqpoint{1.218512in}{0.975812in}}%
\pgfpathlineto{\pgfqpoint{1.209230in}{0.966315in}}%
\pgfpathlineto{\pgfqpoint{1.201592in}{0.967633in}}%
\pgfpathlineto{\pgfqpoint{1.194470in}{0.906193in}}%
\pgfpathclose%
\pgfusepath{fill}%
\end{pgfscope}%
\begin{pgfscope}%
\pgfpathrectangle{\pgfqpoint{0.100000in}{0.100000in}}{\pgfqpoint{3.007045in}{1.925000in}}%
\pgfusepath{clip}%
\pgfsetbuttcap%
\pgfsetmiterjoin%
\definecolor{currentfill}{rgb}{0.381776,0.656517,0.824452}%
\pgfsetfillcolor{currentfill}%
\pgfsetlinewidth{0.000000pt}%
\definecolor{currentstroke}{rgb}{0.000000,0.000000,0.000000}%
\pgfsetstrokecolor{currentstroke}%
\pgfsetstrokeopacity{0.000000}%
\pgfsetdash{}{0pt}%
\pgfpathmoveto{\pgfqpoint{2.022844in}{1.028625in}}%
\pgfpathlineto{\pgfqpoint{2.005455in}{1.027915in}}%
\pgfpathlineto{\pgfqpoint{2.005554in}{1.024952in}}%
\pgfpathlineto{\pgfqpoint{1.995876in}{1.024384in}}%
\pgfpathlineto{\pgfqpoint{1.993598in}{1.031047in}}%
\pgfpathlineto{\pgfqpoint{1.993117in}{1.046153in}}%
\pgfpathlineto{\pgfqpoint{1.981728in}{1.045692in}}%
\pgfpathlineto{\pgfqpoint{1.973010in}{1.044275in}}%
\pgfpathlineto{\pgfqpoint{1.972796in}{1.051471in}}%
\pgfpathlineto{\pgfqpoint{1.975453in}{1.054398in}}%
\pgfpathlineto{\pgfqpoint{1.974361in}{1.084952in}}%
\pgfpathlineto{\pgfqpoint{1.990680in}{1.085683in}}%
\pgfpathlineto{\pgfqpoint{2.001501in}{1.072967in}}%
\pgfpathlineto{\pgfqpoint{2.010046in}{1.075951in}}%
\pgfpathlineto{\pgfqpoint{2.014367in}{1.079611in}}%
\pgfpathlineto{\pgfqpoint{2.020734in}{1.079515in}}%
\pgfpathlineto{\pgfqpoint{2.024754in}{1.078038in}}%
\pgfpathlineto{\pgfqpoint{2.031875in}{1.071974in}}%
\pgfpathlineto{\pgfqpoint{2.034604in}{1.066814in}}%
\pgfpathlineto{\pgfqpoint{2.040262in}{1.068526in}}%
\pgfpathlineto{\pgfqpoint{2.045818in}{1.065000in}}%
\pgfpathlineto{\pgfqpoint{2.054473in}{1.056687in}}%
\pgfpathlineto{\pgfqpoint{2.058159in}{1.055368in}}%
\pgfpathlineto{\pgfqpoint{2.060389in}{1.049493in}}%
\pgfpathlineto{\pgfqpoint{2.058254in}{1.047208in}}%
\pgfpathlineto{\pgfqpoint{2.053627in}{1.049022in}}%
\pgfpathlineto{\pgfqpoint{2.025561in}{1.047408in}}%
\pgfpathlineto{\pgfqpoint{2.026620in}{1.028636in}}%
\pgfpathlineto{\pgfqpoint{2.022844in}{1.028625in}}%
\pgfpathclose%
\pgfusepath{fill}%
\end{pgfscope}%
\begin{pgfscope}%
\pgfpathrectangle{\pgfqpoint{0.100000in}{0.100000in}}{\pgfqpoint{3.007045in}{1.925000in}}%
\pgfusepath{clip}%
\pgfsetbuttcap%
\pgfsetmiterjoin%
\definecolor{currentfill}{rgb}{0.541961,0.749527,0.865606}%
\pgfsetfillcolor{currentfill}%
\pgfsetlinewidth{0.000000pt}%
\definecolor{currentstroke}{rgb}{0.000000,0.000000,0.000000}%
\pgfsetstrokecolor{currentstroke}%
\pgfsetstrokeopacity{0.000000}%
\pgfsetdash{}{0pt}%
\pgfpathmoveto{\pgfqpoint{2.309942in}{0.708924in}}%
\pgfpathlineto{\pgfqpoint{2.309841in}{0.700126in}}%
\pgfpathlineto{\pgfqpoint{2.301073in}{0.698033in}}%
\pgfpathlineto{\pgfqpoint{2.297933in}{0.689896in}}%
\pgfpathlineto{\pgfqpoint{2.290454in}{0.689136in}}%
\pgfpathlineto{\pgfqpoint{2.289927in}{0.694879in}}%
\pgfpathlineto{\pgfqpoint{2.284645in}{0.694345in}}%
\pgfpathlineto{\pgfqpoint{2.283592in}{0.699581in}}%
\pgfpathlineto{\pgfqpoint{2.277794in}{0.699248in}}%
\pgfpathlineto{\pgfqpoint{2.276404in}{0.718014in}}%
\pgfpathlineto{\pgfqpoint{2.274136in}{0.722813in}}%
\pgfpathlineto{\pgfqpoint{2.281053in}{0.729035in}}%
\pgfpathlineto{\pgfqpoint{2.285996in}{0.729613in}}%
\pgfpathlineto{\pgfqpoint{2.285327in}{0.735328in}}%
\pgfpathlineto{\pgfqpoint{2.290941in}{0.736895in}}%
\pgfpathlineto{\pgfqpoint{2.290204in}{0.743581in}}%
\pgfpathlineto{\pgfqpoint{2.295738in}{0.746225in}}%
\pgfpathlineto{\pgfqpoint{2.312802in}{0.748173in}}%
\pgfpathlineto{\pgfqpoint{2.321001in}{0.750113in}}%
\pgfpathlineto{\pgfqpoint{2.323388in}{0.743508in}}%
\pgfpathlineto{\pgfqpoint{2.330141in}{0.735377in}}%
\pgfpathlineto{\pgfqpoint{2.327177in}{0.732545in}}%
\pgfpathlineto{\pgfqpoint{2.311941in}{0.730610in}}%
\pgfpathlineto{\pgfqpoint{2.312554in}{0.726822in}}%
\pgfpathlineto{\pgfqpoint{2.306778in}{0.726142in}}%
\pgfpathlineto{\pgfqpoint{2.308027in}{0.714591in}}%
\pgfpathlineto{\pgfqpoint{2.309942in}{0.708924in}}%
\pgfpathclose%
\pgfusepath{fill}%
\end{pgfscope}%
\begin{pgfscope}%
\pgfpathrectangle{\pgfqpoint{0.100000in}{0.100000in}}{\pgfqpoint{3.007045in}{1.925000in}}%
\pgfusepath{clip}%
\pgfsetbuttcap%
\pgfsetmiterjoin%
\definecolor{currentfill}{rgb}{0.296025,0.597955,0.790988}%
\pgfsetfillcolor{currentfill}%
\pgfsetlinewidth{0.000000pt}%
\definecolor{currentstroke}{rgb}{0.000000,0.000000,0.000000}%
\pgfsetstrokecolor{currentstroke}%
\pgfsetstrokeopacity{0.000000}%
\pgfsetdash{}{0pt}%
\pgfpathmoveto{\pgfqpoint{0.829274in}{1.562729in}}%
\pgfpathlineto{\pgfqpoint{0.828248in}{1.558524in}}%
\pgfpathlineto{\pgfqpoint{0.823175in}{1.550976in}}%
\pgfpathlineto{\pgfqpoint{0.824049in}{1.547593in}}%
\pgfpathlineto{\pgfqpoint{0.819243in}{1.542361in}}%
\pgfpathlineto{\pgfqpoint{0.809080in}{1.497950in}}%
\pgfpathlineto{\pgfqpoint{0.811385in}{1.497272in}}%
\pgfpathlineto{\pgfqpoint{0.798116in}{1.438303in}}%
\pgfpathlineto{\pgfqpoint{0.734607in}{1.453046in}}%
\pgfpathlineto{\pgfqpoint{0.703710in}{1.460774in}}%
\pgfpathlineto{\pgfqpoint{0.703346in}{1.460924in}}%
\pgfpathlineto{\pgfqpoint{0.729959in}{1.568628in}}%
\pgfpathlineto{\pgfqpoint{0.733534in}{1.580107in}}%
\pgfpathlineto{\pgfqpoint{0.736446in}{1.580574in}}%
\pgfpathlineto{\pgfqpoint{0.739760in}{1.573820in}}%
\pgfpathlineto{\pgfqpoint{0.746496in}{1.573138in}}%
\pgfpathlineto{\pgfqpoint{0.749229in}{1.584343in}}%
\pgfpathlineto{\pgfqpoint{0.755702in}{1.582709in}}%
\pgfpathlineto{\pgfqpoint{0.758886in}{1.587771in}}%
\pgfpathlineto{\pgfqpoint{0.762552in}{1.586835in}}%
\pgfpathlineto{\pgfqpoint{0.763908in}{1.592366in}}%
\pgfpathlineto{\pgfqpoint{0.768216in}{1.591309in}}%
\pgfpathlineto{\pgfqpoint{0.771340in}{1.602215in}}%
\pgfpathlineto{\pgfqpoint{0.775430in}{1.609514in}}%
\pgfpathlineto{\pgfqpoint{0.783037in}{1.611433in}}%
\pgfpathlineto{\pgfqpoint{0.780379in}{1.600195in}}%
\pgfpathlineto{\pgfqpoint{0.777625in}{1.600841in}}%
\pgfpathlineto{\pgfqpoint{0.774977in}{1.589715in}}%
\pgfpathlineto{\pgfqpoint{0.779546in}{1.588394in}}%
\pgfpathlineto{\pgfqpoint{0.781320in}{1.594007in}}%
\pgfpathlineto{\pgfqpoint{0.808189in}{1.587646in}}%
\pgfpathlineto{\pgfqpoint{0.811192in}{1.587476in}}%
\pgfpathlineto{\pgfqpoint{0.820142in}{1.591987in}}%
\pgfpathlineto{\pgfqpoint{0.824845in}{1.588003in}}%
\pgfpathlineto{\pgfqpoint{0.824028in}{1.581837in}}%
\pgfpathlineto{\pgfqpoint{0.829905in}{1.577513in}}%
\pgfpathlineto{\pgfqpoint{0.827488in}{1.571131in}}%
\pgfpathlineto{\pgfqpoint{0.829274in}{1.562729in}}%
\pgfpathclose%
\pgfusepath{fill}%
\end{pgfscope}%
\begin{pgfscope}%
\pgfpathrectangle{\pgfqpoint{0.100000in}{0.100000in}}{\pgfqpoint{3.007045in}{1.925000in}}%
\pgfusepath{clip}%
\pgfsetbuttcap%
\pgfsetmiterjoin%
\definecolor{currentfill}{rgb}{0.231926,0.545652,0.762614}%
\pgfsetfillcolor{currentfill}%
\pgfsetlinewidth{0.000000pt}%
\definecolor{currentstroke}{rgb}{0.000000,0.000000,0.000000}%
\pgfsetstrokecolor{currentstroke}%
\pgfsetstrokeopacity{0.000000}%
\pgfsetdash{}{0pt}%
\pgfpathmoveto{\pgfqpoint{1.910648in}{0.511564in}}%
\pgfpathlineto{\pgfqpoint{1.903450in}{0.513741in}}%
\pgfpathlineto{\pgfqpoint{1.888369in}{0.520432in}}%
\pgfpathlineto{\pgfqpoint{1.876777in}{0.523740in}}%
\pgfpathlineto{\pgfqpoint{1.868009in}{0.522526in}}%
\pgfpathlineto{\pgfqpoint{1.859158in}{0.523013in}}%
\pgfpathlineto{\pgfqpoint{1.843789in}{0.520220in}}%
\pgfpathlineto{\pgfqpoint{1.839662in}{0.517608in}}%
\pgfpathlineto{\pgfqpoint{1.834496in}{0.525796in}}%
\pgfpathlineto{\pgfqpoint{1.840585in}{0.532547in}}%
\pgfpathlineto{\pgfqpoint{1.842229in}{0.537436in}}%
\pgfpathlineto{\pgfqpoint{1.846781in}{0.541896in}}%
\pgfpathlineto{\pgfqpoint{1.846835in}{0.557616in}}%
\pgfpathlineto{\pgfqpoint{1.843280in}{0.560440in}}%
\pgfpathlineto{\pgfqpoint{1.847088in}{0.567676in}}%
\pgfpathlineto{\pgfqpoint{1.844939in}{0.573686in}}%
\pgfpathlineto{\pgfqpoint{1.847738in}{0.581014in}}%
\pgfpathlineto{\pgfqpoint{1.850658in}{0.583534in}}%
\pgfpathlineto{\pgfqpoint{1.853934in}{0.596710in}}%
\pgfpathlineto{\pgfqpoint{1.856414in}{0.600398in}}%
\pgfpathlineto{\pgfqpoint{1.854154in}{0.604333in}}%
\pgfpathlineto{\pgfqpoint{1.854631in}{0.617364in}}%
\pgfpathlineto{\pgfqpoint{1.855514in}{0.623416in}}%
\pgfpathlineto{\pgfqpoint{1.860787in}{0.623525in}}%
\pgfpathlineto{\pgfqpoint{1.863641in}{0.629386in}}%
\pgfpathlineto{\pgfqpoint{1.872399in}{0.629573in}}%
\pgfpathlineto{\pgfqpoint{1.883960in}{0.629796in}}%
\pgfpathlineto{\pgfqpoint{1.886870in}{0.628687in}}%
\pgfpathlineto{\pgfqpoint{1.892685in}{0.626874in}}%
\pgfpathlineto{\pgfqpoint{1.895632in}{0.622616in}}%
\pgfpathlineto{\pgfqpoint{1.896550in}{0.598540in}}%
\pgfpathlineto{\pgfqpoint{1.909431in}{0.599271in}}%
\pgfpathlineto{\pgfqpoint{1.910006in}{0.578352in}}%
\pgfpathlineto{\pgfqpoint{1.908086in}{0.571754in}}%
\pgfpathlineto{\pgfqpoint{1.909910in}{0.542133in}}%
\pgfpathlineto{\pgfqpoint{1.910648in}{0.511564in}}%
\pgfpathclose%
\pgfusepath{fill}%
\end{pgfscope}%
\begin{pgfscope}%
\pgfpathrectangle{\pgfqpoint{0.100000in}{0.100000in}}{\pgfqpoint{3.007045in}{1.925000in}}%
\pgfusepath{clip}%
\pgfsetbuttcap%
\pgfsetmiterjoin%
\definecolor{currentfill}{rgb}{0.784591,0.864237,0.939962}%
\pgfsetfillcolor{currentfill}%
\pgfsetlinewidth{0.000000pt}%
\definecolor{currentstroke}{rgb}{0.000000,0.000000,0.000000}%
\pgfsetstrokecolor{currentstroke}%
\pgfsetstrokeopacity{0.000000}%
\pgfsetdash{}{0pt}%
\pgfpathmoveto{\pgfqpoint{0.677176in}{1.047687in}}%
\pgfpathlineto{\pgfqpoint{0.672324in}{1.047770in}}%
\pgfpathlineto{\pgfqpoint{0.654898in}{1.051840in}}%
\pgfpathlineto{\pgfqpoint{0.593532in}{1.066855in}}%
\pgfpathlineto{\pgfqpoint{0.573744in}{1.071996in}}%
\pgfpathlineto{\pgfqpoint{0.554172in}{1.076466in}}%
\pgfpathlineto{\pgfqpoint{0.557260in}{1.082859in}}%
\pgfpathlineto{\pgfqpoint{0.556142in}{1.096690in}}%
\pgfpathlineto{\pgfqpoint{0.557824in}{1.101480in}}%
\pgfpathlineto{\pgfqpoint{0.556496in}{1.110688in}}%
\pgfpathlineto{\pgfqpoint{0.558630in}{1.113262in}}%
\pgfpathlineto{\pgfqpoint{0.553369in}{1.124234in}}%
\pgfpathlineto{\pgfqpoint{0.554104in}{1.127797in}}%
\pgfpathlineto{\pgfqpoint{0.551392in}{1.137382in}}%
\pgfpathlineto{\pgfqpoint{0.551954in}{1.148285in}}%
\pgfpathlineto{\pgfqpoint{0.554321in}{1.151293in}}%
\pgfpathlineto{\pgfqpoint{0.552868in}{1.161985in}}%
\pgfpathlineto{\pgfqpoint{0.546742in}{1.170366in}}%
\pgfpathlineto{\pgfqpoint{0.543660in}{1.171437in}}%
\pgfpathlineto{\pgfqpoint{0.543712in}{1.184135in}}%
\pgfpathlineto{\pgfqpoint{0.540581in}{1.185427in}}%
\pgfpathlineto{\pgfqpoint{0.543469in}{1.191912in}}%
\pgfpathlineto{\pgfqpoint{0.539193in}{1.195591in}}%
\pgfpathlineto{\pgfqpoint{0.537345in}{1.200997in}}%
\pgfpathlineto{\pgfqpoint{0.532644in}{1.204682in}}%
\pgfpathlineto{\pgfqpoint{0.532191in}{1.211551in}}%
\pgfpathlineto{\pgfqpoint{0.530069in}{1.215499in}}%
\pgfpathlineto{\pgfqpoint{0.522889in}{1.217818in}}%
\pgfpathlineto{\pgfqpoint{0.527157in}{1.220456in}}%
\pgfpathlineto{\pgfqpoint{0.527593in}{1.228218in}}%
\pgfpathlineto{\pgfqpoint{0.524524in}{1.231551in}}%
\pgfpathlineto{\pgfqpoint{0.525021in}{1.240725in}}%
\pgfpathlineto{\pgfqpoint{0.514807in}{1.248637in}}%
\pgfpathlineto{\pgfqpoint{0.513336in}{1.257851in}}%
\pgfpathlineto{\pgfqpoint{0.514663in}{1.260246in}}%
\pgfpathlineto{\pgfqpoint{0.519485in}{1.264180in}}%
\pgfpathlineto{\pgfqpoint{0.522497in}{1.269886in}}%
\pgfpathlineto{\pgfqpoint{0.520542in}{1.277526in}}%
\pgfpathlineto{\pgfqpoint{0.524302in}{1.283464in}}%
\pgfpathlineto{\pgfqpoint{0.533766in}{1.269179in}}%
\pgfpathlineto{\pgfqpoint{0.540209in}{1.259221in}}%
\pgfpathlineto{\pgfqpoint{0.567833in}{1.216455in}}%
\pgfpathlineto{\pgfqpoint{0.590725in}{1.181138in}}%
\pgfpathlineto{\pgfqpoint{0.616671in}{1.141013in}}%
\pgfpathlineto{\pgfqpoint{0.667024in}{1.063268in}}%
\pgfpathlineto{\pgfqpoint{0.677176in}{1.047687in}}%
\pgfpathclose%
\pgfusepath{fill}%
\end{pgfscope}%
\begin{pgfscope}%
\pgfpathrectangle{\pgfqpoint{0.100000in}{0.100000in}}{\pgfqpoint{3.007045in}{1.925000in}}%
\pgfusepath{clip}%
\pgfsetbuttcap%
\pgfsetmiterjoin%
\definecolor{currentfill}{rgb}{0.260715,0.573841,0.777209}%
\pgfsetfillcolor{currentfill}%
\pgfsetlinewidth{0.000000pt}%
\definecolor{currentstroke}{rgb}{0.000000,0.000000,0.000000}%
\pgfsetstrokecolor{currentstroke}%
\pgfsetstrokeopacity{0.000000}%
\pgfsetdash{}{0pt}%
\pgfpathmoveto{\pgfqpoint{2.442052in}{0.679496in}}%
\pgfpathlineto{\pgfqpoint{2.428996in}{0.678445in}}%
\pgfpathlineto{\pgfqpoint{2.425943in}{0.682862in}}%
\pgfpathlineto{\pgfqpoint{2.421279in}{0.683375in}}%
\pgfpathlineto{\pgfqpoint{2.420802in}{0.691745in}}%
\pgfpathlineto{\pgfqpoint{2.422816in}{0.692785in}}%
\pgfpathlineto{\pgfqpoint{2.420123in}{0.701779in}}%
\pgfpathlineto{\pgfqpoint{2.437033in}{0.704008in}}%
\pgfpathlineto{\pgfqpoint{2.447789in}{0.700910in}}%
\pgfpathlineto{\pgfqpoint{2.448668in}{0.692985in}}%
\pgfpathlineto{\pgfqpoint{2.445730in}{0.689304in}}%
\pgfpathlineto{\pgfqpoint{2.443946in}{0.680145in}}%
\pgfpathlineto{\pgfqpoint{2.442052in}{0.679496in}}%
\pgfpathclose%
\pgfusepath{fill}%
\end{pgfscope}%
\begin{pgfscope}%
\pgfpathrectangle{\pgfqpoint{0.100000in}{0.100000in}}{\pgfqpoint{3.007045in}{1.925000in}}%
\pgfusepath{clip}%
\pgfsetbuttcap%
\pgfsetmiterjoin%
\definecolor{currentfill}{rgb}{0.622684,0.793464,0.883429}%
\pgfsetfillcolor{currentfill}%
\pgfsetlinewidth{0.000000pt}%
\definecolor{currentstroke}{rgb}{0.000000,0.000000,0.000000}%
\pgfsetstrokecolor{currentstroke}%
\pgfsetstrokeopacity{0.000000}%
\pgfsetdash{}{0pt}%
\pgfpathmoveto{\pgfqpoint{0.404299in}{1.600395in}}%
\pgfpathlineto{\pgfqpoint{0.415190in}{1.591253in}}%
\pgfpathlineto{\pgfqpoint{0.419385in}{1.591902in}}%
\pgfpathlineto{\pgfqpoint{0.426249in}{1.588654in}}%
\pgfpathlineto{\pgfqpoint{0.429770in}{1.590236in}}%
\pgfpathlineto{\pgfqpoint{0.436255in}{1.587889in}}%
\pgfpathlineto{\pgfqpoint{0.446705in}{1.586464in}}%
\pgfpathlineto{\pgfqpoint{0.460693in}{1.593122in}}%
\pgfpathlineto{\pgfqpoint{0.465330in}{1.591720in}}%
\pgfpathlineto{\pgfqpoint{0.469189in}{1.594432in}}%
\pgfpathlineto{\pgfqpoint{0.474662in}{1.592743in}}%
\pgfpathlineto{\pgfqpoint{0.455007in}{1.530523in}}%
\pgfpathlineto{\pgfqpoint{0.397486in}{1.548285in}}%
\pgfpathlineto{\pgfqpoint{0.383232in}{1.552564in}}%
\pgfpathlineto{\pgfqpoint{0.384587in}{1.562231in}}%
\pgfpathlineto{\pgfqpoint{0.389591in}{1.567203in}}%
\pgfpathlineto{\pgfqpoint{0.387396in}{1.575981in}}%
\pgfpathlineto{\pgfqpoint{0.381248in}{1.578474in}}%
\pgfpathlineto{\pgfqpoint{0.384929in}{1.586076in}}%
\pgfpathlineto{\pgfqpoint{0.390435in}{1.585257in}}%
\pgfpathlineto{\pgfqpoint{0.395865in}{1.592654in}}%
\pgfpathlineto{\pgfqpoint{0.404299in}{1.600395in}}%
\pgfpathclose%
\pgfusepath{fill}%
\end{pgfscope}%
\begin{pgfscope}%
\pgfpathrectangle{\pgfqpoint{0.100000in}{0.100000in}}{\pgfqpoint{3.007045in}{1.925000in}}%
\pgfusepath{clip}%
\pgfsetbuttcap%
\pgfsetmiterjoin%
\definecolor{currentfill}{rgb}{0.560784,0.759862,0.869666}%
\pgfsetfillcolor{currentfill}%
\pgfsetlinewidth{0.000000pt}%
\definecolor{currentstroke}{rgb}{0.000000,0.000000,0.000000}%
\pgfsetstrokecolor{currentstroke}%
\pgfsetstrokeopacity{0.000000}%
\pgfsetdash{}{0pt}%
\pgfpathmoveto{\pgfqpoint{2.295738in}{0.746225in}}%
\pgfpathlineto{\pgfqpoint{2.290204in}{0.743581in}}%
\pgfpathlineto{\pgfqpoint{2.290941in}{0.736895in}}%
\pgfpathlineto{\pgfqpoint{2.285327in}{0.735328in}}%
\pgfpathlineto{\pgfqpoint{2.285996in}{0.729613in}}%
\pgfpathlineto{\pgfqpoint{2.281053in}{0.729035in}}%
\pgfpathlineto{\pgfqpoint{2.279667in}{0.746313in}}%
\pgfpathlineto{\pgfqpoint{2.252087in}{0.743598in}}%
\pgfpathlineto{\pgfqpoint{2.247289in}{0.747109in}}%
\pgfpathlineto{\pgfqpoint{2.243973in}{0.753410in}}%
\pgfpathlineto{\pgfqpoint{2.242492in}{0.770671in}}%
\pgfpathlineto{\pgfqpoint{2.244576in}{0.775293in}}%
\pgfpathlineto{\pgfqpoint{2.249772in}{0.779385in}}%
\pgfpathlineto{\pgfqpoint{2.248013in}{0.785562in}}%
\pgfpathlineto{\pgfqpoint{2.257385in}{0.796619in}}%
\pgfpathlineto{\pgfqpoint{2.255427in}{0.799367in}}%
\pgfpathlineto{\pgfqpoint{2.258979in}{0.805773in}}%
\pgfpathlineto{\pgfqpoint{2.264190in}{0.805986in}}%
\pgfpathlineto{\pgfqpoint{2.267838in}{0.800489in}}%
\pgfpathlineto{\pgfqpoint{2.278832in}{0.801530in}}%
\pgfpathlineto{\pgfqpoint{2.281013in}{0.795971in}}%
\pgfpathlineto{\pgfqpoint{2.287725in}{0.796480in}}%
\pgfpathlineto{\pgfqpoint{2.289813in}{0.770791in}}%
\pgfpathlineto{\pgfqpoint{2.293144in}{0.771183in}}%
\pgfpathlineto{\pgfqpoint{2.295738in}{0.746225in}}%
\pgfpathclose%
\pgfusepath{fill}%
\end{pgfscope}%
\begin{pgfscope}%
\pgfpathrectangle{\pgfqpoint{0.100000in}{0.100000in}}{\pgfqpoint{3.007045in}{1.925000in}}%
\pgfusepath{clip}%
\pgfsetbuttcap%
\pgfsetmiterjoin%
\definecolor{currentfill}{rgb}{0.498039,0.725413,0.856132}%
\pgfsetfillcolor{currentfill}%
\pgfsetlinewidth{0.000000pt}%
\definecolor{currentstroke}{rgb}{0.000000,0.000000,0.000000}%
\pgfsetstrokecolor{currentstroke}%
\pgfsetstrokeopacity{0.000000}%
\pgfsetdash{}{0pt}%
\pgfpathmoveto{\pgfqpoint{2.067386in}{1.211625in}}%
\pgfpathlineto{\pgfqpoint{2.066190in}{1.226950in}}%
\pgfpathlineto{\pgfqpoint{2.060239in}{1.229418in}}%
\pgfpathlineto{\pgfqpoint{2.058847in}{1.247095in}}%
\pgfpathlineto{\pgfqpoint{2.065605in}{1.247690in}}%
\pgfpathlineto{\pgfqpoint{2.072749in}{1.252662in}}%
\pgfpathlineto{\pgfqpoint{2.075117in}{1.258653in}}%
\pgfpathlineto{\pgfqpoint{2.073487in}{1.281826in}}%
\pgfpathlineto{\pgfqpoint{2.090562in}{1.283138in}}%
\pgfpathlineto{\pgfqpoint{2.107101in}{1.284740in}}%
\pgfpathlineto{\pgfqpoint{2.107905in}{1.276923in}}%
\pgfpathlineto{\pgfqpoint{2.110362in}{1.252216in}}%
\pgfpathlineto{\pgfqpoint{2.099165in}{1.251343in}}%
\pgfpathlineto{\pgfqpoint{2.100744in}{1.229280in}}%
\pgfpathlineto{\pgfqpoint{2.094960in}{1.228863in}}%
\pgfpathlineto{\pgfqpoint{2.087435in}{1.213369in}}%
\pgfpathlineto{\pgfqpoint{2.067386in}{1.211625in}}%
\pgfpathclose%
\pgfusepath{fill}%
\end{pgfscope}%
\begin{pgfscope}%
\pgfpathrectangle{\pgfqpoint{0.100000in}{0.100000in}}{\pgfqpoint{3.007045in}{1.925000in}}%
\pgfusepath{clip}%
\pgfsetbuttcap%
\pgfsetmiterjoin%
\definecolor{currentfill}{rgb}{0.256286,0.570012,0.775163}%
\pgfsetfillcolor{currentfill}%
\pgfsetlinewidth{0.000000pt}%
\definecolor{currentstroke}{rgb}{0.000000,0.000000,0.000000}%
\pgfsetstrokecolor{currentstroke}%
\pgfsetstrokeopacity{0.000000}%
\pgfsetdash{}{0pt}%
\pgfpathmoveto{\pgfqpoint{1.952131in}{0.625210in}}%
\pgfpathlineto{\pgfqpoint{1.961233in}{0.620328in}}%
\pgfpathlineto{\pgfqpoint{1.958991in}{0.616692in}}%
\pgfpathlineto{\pgfqpoint{1.958935in}{0.610678in}}%
\pgfpathlineto{\pgfqpoint{1.954440in}{0.605869in}}%
\pgfpathlineto{\pgfqpoint{1.954185in}{0.597454in}}%
\pgfpathlineto{\pgfqpoint{1.930070in}{0.596671in}}%
\pgfpathlineto{\pgfqpoint{1.927440in}{0.602431in}}%
\pgfpathlineto{\pgfqpoint{1.921122in}{0.606604in}}%
\pgfpathlineto{\pgfqpoint{1.915657in}{0.600227in}}%
\pgfpathlineto{\pgfqpoint{1.909431in}{0.599271in}}%
\pgfpathlineto{\pgfqpoint{1.896550in}{0.598540in}}%
\pgfpathlineto{\pgfqpoint{1.895632in}{0.622616in}}%
\pgfpathlineto{\pgfqpoint{1.892685in}{0.626874in}}%
\pgfpathlineto{\pgfqpoint{1.886870in}{0.628687in}}%
\pgfpathlineto{\pgfqpoint{1.890720in}{0.630209in}}%
\pgfpathlineto{\pgfqpoint{1.901476in}{0.640393in}}%
\pgfpathlineto{\pgfqpoint{1.896682in}{0.645552in}}%
\pgfpathlineto{\pgfqpoint{1.890709in}{0.647179in}}%
\pgfpathlineto{\pgfqpoint{1.886783in}{0.652756in}}%
\pgfpathlineto{\pgfqpoint{1.906746in}{0.653283in}}%
\pgfpathlineto{\pgfqpoint{1.906637in}{0.659096in}}%
\pgfpathlineto{\pgfqpoint{1.921110in}{0.659462in}}%
\pgfpathlineto{\pgfqpoint{1.922010in}{0.647995in}}%
\pgfpathlineto{\pgfqpoint{1.926639in}{0.641691in}}%
\pgfpathlineto{\pgfqpoint{1.933198in}{0.637846in}}%
\pgfpathlineto{\pgfqpoint{1.933963in}{0.633867in}}%
\pgfpathlineto{\pgfqpoint{1.942513in}{0.628295in}}%
\pgfpathlineto{\pgfqpoint{1.952131in}{0.625210in}}%
\pgfpathclose%
\pgfusepath{fill}%
\end{pgfscope}%
\begin{pgfscope}%
\pgfpathrectangle{\pgfqpoint{0.100000in}{0.100000in}}{\pgfqpoint{3.007045in}{1.925000in}}%
\pgfusepath{clip}%
\pgfsetbuttcap%
\pgfsetmiterjoin%
\definecolor{currentfill}{rgb}{0.429020,0.687520,0.841246}%
\pgfsetfillcolor{currentfill}%
\pgfsetlinewidth{0.000000pt}%
\definecolor{currentstroke}{rgb}{0.000000,0.000000,0.000000}%
\pgfsetstrokecolor{currentstroke}%
\pgfsetstrokeopacity{0.000000}%
\pgfsetdash{}{0pt}%
\pgfpathmoveto{\pgfqpoint{2.131121in}{1.073626in}}%
\pgfpathlineto{\pgfqpoint{2.128732in}{1.074544in}}%
\pgfpathlineto{\pgfqpoint{2.116398in}{1.073683in}}%
\pgfpathlineto{\pgfqpoint{2.114940in}{1.096476in}}%
\pgfpathlineto{\pgfqpoint{2.097811in}{1.095356in}}%
\pgfpathlineto{\pgfqpoint{2.096418in}{1.118342in}}%
\pgfpathlineto{\pgfqpoint{2.119277in}{1.119542in}}%
\pgfpathlineto{\pgfqpoint{2.118010in}{1.123238in}}%
\pgfpathlineto{\pgfqpoint{2.117823in}{1.135815in}}%
\pgfpathlineto{\pgfqpoint{2.135746in}{1.137366in}}%
\pgfpathlineto{\pgfqpoint{2.136969in}{1.118978in}}%
\pgfpathlineto{\pgfqpoint{2.135045in}{1.113976in}}%
\pgfpathlineto{\pgfqpoint{2.134603in}{1.106062in}}%
\pgfpathlineto{\pgfqpoint{2.136503in}{1.101602in}}%
\pgfpathlineto{\pgfqpoint{2.134514in}{1.098251in}}%
\pgfpathlineto{\pgfqpoint{2.137521in}{1.089690in}}%
\pgfpathlineto{\pgfqpoint{2.131121in}{1.073626in}}%
\pgfpathclose%
\pgfusepath{fill}%
\end{pgfscope}%
\begin{pgfscope}%
\pgfpathrectangle{\pgfqpoint{0.100000in}{0.100000in}}{\pgfqpoint{3.007045in}{1.925000in}}%
\pgfusepath{clip}%
\pgfsetbuttcap%
\pgfsetmiterjoin%
\definecolor{currentfill}{rgb}{0.642368,0.801830,0.890319}%
\pgfsetfillcolor{currentfill}%
\pgfsetlinewidth{0.000000pt}%
\definecolor{currentstroke}{rgb}{0.000000,0.000000,0.000000}%
\pgfsetstrokecolor{currentstroke}%
\pgfsetstrokeopacity{0.000000}%
\pgfsetdash{}{0pt}%
\pgfpathmoveto{\pgfqpoint{2.450037in}{1.161421in}}%
\pgfpathlineto{\pgfqpoint{2.440950in}{1.158535in}}%
\pgfpathlineto{\pgfqpoint{2.435873in}{1.154519in}}%
\pgfpathlineto{\pgfqpoint{2.437602in}{1.145887in}}%
\pgfpathlineto{\pgfqpoint{2.428157in}{1.152335in}}%
\pgfpathlineto{\pgfqpoint{2.420941in}{1.150887in}}%
\pgfpathlineto{\pgfqpoint{2.420478in}{1.156867in}}%
\pgfpathlineto{\pgfqpoint{2.414411in}{1.156435in}}%
\pgfpathlineto{\pgfqpoint{2.413032in}{1.162228in}}%
\pgfpathlineto{\pgfqpoint{2.407930in}{1.161976in}}%
\pgfpathlineto{\pgfqpoint{2.407704in}{1.166293in}}%
\pgfpathlineto{\pgfqpoint{2.413821in}{1.167145in}}%
\pgfpathlineto{\pgfqpoint{2.413102in}{1.179720in}}%
\pgfpathlineto{\pgfqpoint{2.418866in}{1.180027in}}%
\pgfpathlineto{\pgfqpoint{2.418057in}{1.191742in}}%
\pgfpathlineto{\pgfqpoint{2.419602in}{1.197771in}}%
\pgfpathlineto{\pgfqpoint{2.417290in}{1.203518in}}%
\pgfpathlineto{\pgfqpoint{2.416942in}{1.209349in}}%
\pgfpathlineto{\pgfqpoint{2.422680in}{1.209735in}}%
\pgfpathlineto{\pgfqpoint{2.422257in}{1.215672in}}%
\pgfpathlineto{\pgfqpoint{2.427935in}{1.216104in}}%
\pgfpathlineto{\pgfqpoint{2.428341in}{1.210131in}}%
\pgfpathlineto{\pgfqpoint{2.439065in}{1.211016in}}%
\pgfpathlineto{\pgfqpoint{2.439887in}{1.202342in}}%
\pgfpathlineto{\pgfqpoint{2.446995in}{1.200159in}}%
\pgfpathlineto{\pgfqpoint{2.446120in}{1.194088in}}%
\pgfpathlineto{\pgfqpoint{2.447277in}{1.183280in}}%
\pgfpathlineto{\pgfqpoint{2.441696in}{1.172351in}}%
\pgfpathlineto{\pgfqpoint{2.450037in}{1.161421in}}%
\pgfpathclose%
\pgfusepath{fill}%
\end{pgfscope}%
\begin{pgfscope}%
\pgfpathrectangle{\pgfqpoint{0.100000in}{0.100000in}}{\pgfqpoint{3.007045in}{1.925000in}}%
\pgfusepath{clip}%
\pgfsetbuttcap%
\pgfsetmiterjoin%
\definecolor{currentfill}{rgb}{0.361599,0.642737,0.816578}%
\pgfsetfillcolor{currentfill}%
\pgfsetlinewidth{0.000000pt}%
\definecolor{currentstroke}{rgb}{0.000000,0.000000,0.000000}%
\pgfsetstrokecolor{currentstroke}%
\pgfsetstrokeopacity{0.000000}%
\pgfsetdash{}{0pt}%
\pgfpathmoveto{\pgfqpoint{1.594392in}{1.173857in}}%
\pgfpathlineto{\pgfqpoint{1.623649in}{1.172778in}}%
\pgfpathlineto{\pgfqpoint{1.623055in}{1.155630in}}%
\pgfpathlineto{\pgfqpoint{1.622833in}{1.149886in}}%
\pgfpathlineto{\pgfqpoint{1.594289in}{1.150925in}}%
\pgfpathlineto{\pgfqpoint{1.594074in}{1.145204in}}%
\pgfpathlineto{\pgfqpoint{1.565600in}{1.146405in}}%
\pgfpathlineto{\pgfqpoint{1.567058in}{1.175010in}}%
\pgfpathlineto{\pgfqpoint{1.594392in}{1.173857in}}%
\pgfpathclose%
\pgfusepath{fill}%
\end{pgfscope}%
\begin{pgfscope}%
\pgfpathrectangle{\pgfqpoint{0.100000in}{0.100000in}}{\pgfqpoint{3.007045in}{1.925000in}}%
\pgfusepath{clip}%
\pgfsetbuttcap%
\pgfsetmiterjoin%
\definecolor{currentfill}{rgb}{0.491765,0.721968,0.854779}%
\pgfsetfillcolor{currentfill}%
\pgfsetlinewidth{0.000000pt}%
\definecolor{currentstroke}{rgb}{0.000000,0.000000,0.000000}%
\pgfsetstrokecolor{currentstroke}%
\pgfsetstrokeopacity{0.000000}%
\pgfsetdash{}{0pt}%
\pgfpathmoveto{\pgfqpoint{2.695089in}{1.071478in}}%
\pgfpathlineto{\pgfqpoint{2.691815in}{1.071736in}}%
\pgfpathlineto{\pgfqpoint{2.691645in}{1.081236in}}%
\pgfpathlineto{\pgfqpoint{2.685692in}{1.079312in}}%
\pgfpathlineto{\pgfqpoint{2.680930in}{1.081273in}}%
\pgfpathlineto{\pgfqpoint{2.674370in}{1.085389in}}%
\pgfpathlineto{\pgfqpoint{2.667768in}{1.085077in}}%
\pgfpathlineto{\pgfqpoint{2.665311in}{1.095015in}}%
\pgfpathlineto{\pgfqpoint{2.670092in}{1.099194in}}%
\pgfpathlineto{\pgfqpoint{2.646388in}{1.101662in}}%
\pgfpathlineto{\pgfqpoint{2.645322in}{1.106286in}}%
\pgfpathlineto{\pgfqpoint{2.649572in}{1.112889in}}%
\pgfpathlineto{\pgfqpoint{2.649818in}{1.126222in}}%
\pgfpathlineto{\pgfqpoint{2.644603in}{1.131634in}}%
\pgfpathlineto{\pgfqpoint{2.647609in}{1.142718in}}%
\pgfpathlineto{\pgfqpoint{2.654207in}{1.139998in}}%
\pgfpathlineto{\pgfqpoint{2.658319in}{1.135059in}}%
\pgfpathlineto{\pgfqpoint{2.663369in}{1.133914in}}%
\pgfpathlineto{\pgfqpoint{2.665522in}{1.153025in}}%
\pgfpathlineto{\pgfqpoint{2.679314in}{1.148103in}}%
\pgfpathlineto{\pgfqpoint{2.681074in}{1.143862in}}%
\pgfpathlineto{\pgfqpoint{2.685472in}{1.142255in}}%
\pgfpathlineto{\pgfqpoint{2.692040in}{1.154397in}}%
\pgfpathlineto{\pgfqpoint{2.697515in}{1.156406in}}%
\pgfpathlineto{\pgfqpoint{2.701679in}{1.149182in}}%
\pgfpathlineto{\pgfqpoint{2.706818in}{1.146620in}}%
\pgfpathlineto{\pgfqpoint{2.713398in}{1.148167in}}%
\pgfpathlineto{\pgfqpoint{2.716860in}{1.144435in}}%
\pgfpathlineto{\pgfqpoint{2.723470in}{1.137350in}}%
\pgfpathlineto{\pgfqpoint{2.724698in}{1.127666in}}%
\pgfpathlineto{\pgfqpoint{2.713382in}{1.123521in}}%
\pgfpathlineto{\pgfqpoint{2.716713in}{1.110416in}}%
\pgfpathlineto{\pgfqpoint{2.711789in}{1.111134in}}%
\pgfpathlineto{\pgfqpoint{2.704080in}{1.100551in}}%
\pgfpathlineto{\pgfqpoint{2.713742in}{1.098160in}}%
\pgfpathlineto{\pgfqpoint{2.716714in}{1.091684in}}%
\pgfpathlineto{\pgfqpoint{2.695089in}{1.071478in}}%
\pgfpathclose%
\pgfusepath{fill}%
\end{pgfscope}%
\begin{pgfscope}%
\pgfpathrectangle{\pgfqpoint{0.100000in}{0.100000in}}{\pgfqpoint{3.007045in}{1.925000in}}%
\pgfusepath{clip}%
\pgfsetbuttcap%
\pgfsetmiterjoin%
\definecolor{currentfill}{rgb}{0.290980,0.594510,0.789020}%
\pgfsetfillcolor{currentfill}%
\pgfsetlinewidth{0.000000pt}%
\definecolor{currentstroke}{rgb}{0.000000,0.000000,0.000000}%
\pgfsetstrokecolor{currentstroke}%
\pgfsetstrokeopacity{0.000000}%
\pgfsetdash{}{0pt}%
\pgfpathmoveto{\pgfqpoint{1.681563in}{0.619854in}}%
\pgfpathlineto{\pgfqpoint{1.668209in}{0.612458in}}%
\pgfpathlineto{\pgfqpoint{1.654308in}{0.604797in}}%
\pgfpathlineto{\pgfqpoint{1.642936in}{0.624588in}}%
\pgfpathlineto{\pgfqpoint{1.639177in}{0.622324in}}%
\pgfpathlineto{\pgfqpoint{1.635041in}{0.627484in}}%
\pgfpathlineto{\pgfqpoint{1.620418in}{0.653867in}}%
\pgfpathlineto{\pgfqpoint{1.616016in}{0.651438in}}%
\pgfpathlineto{\pgfqpoint{1.603281in}{0.674775in}}%
\pgfpathlineto{\pgfqpoint{1.611326in}{0.679134in}}%
\pgfpathlineto{\pgfqpoint{1.624705in}{0.686669in}}%
\pgfpathlineto{\pgfqpoint{1.628904in}{0.682713in}}%
\pgfpathlineto{\pgfqpoint{1.633343in}{0.684748in}}%
\pgfpathlineto{\pgfqpoint{1.655346in}{0.689754in}}%
\pgfpathlineto{\pgfqpoint{1.663204in}{0.675422in}}%
\pgfpathlineto{\pgfqpoint{1.665733in}{0.676821in}}%
\pgfpathlineto{\pgfqpoint{1.675361in}{0.659431in}}%
\pgfpathlineto{\pgfqpoint{1.663193in}{0.652637in}}%
\pgfpathlineto{\pgfqpoint{1.681563in}{0.619854in}}%
\pgfpathclose%
\pgfusepath{fill}%
\end{pgfscope}%
\begin{pgfscope}%
\pgfpathrectangle{\pgfqpoint{0.100000in}{0.100000in}}{\pgfqpoint{3.007045in}{1.925000in}}%
\pgfusepath{clip}%
\pgfsetbuttcap%
\pgfsetmiterjoin%
\definecolor{currentfill}{rgb}{0.676817,0.816471,0.902376}%
\pgfsetfillcolor{currentfill}%
\pgfsetlinewidth{0.000000pt}%
\definecolor{currentstroke}{rgb}{0.000000,0.000000,0.000000}%
\pgfsetstrokecolor{currentstroke}%
\pgfsetstrokeopacity{0.000000}%
\pgfsetdash{}{0pt}%
\pgfpathmoveto{\pgfqpoint{2.498857in}{1.196503in}}%
\pgfpathlineto{\pgfqpoint{2.494442in}{1.196721in}}%
\pgfpathlineto{\pgfqpoint{2.487686in}{1.206566in}}%
\pgfpathlineto{\pgfqpoint{2.488733in}{1.208153in}}%
\pgfpathlineto{\pgfqpoint{2.500713in}{1.220454in}}%
\pgfpathlineto{\pgfqpoint{2.505896in}{1.220258in}}%
\pgfpathlineto{\pgfqpoint{2.507751in}{1.225110in}}%
\pgfpathlineto{\pgfqpoint{2.505584in}{1.228741in}}%
\pgfpathlineto{\pgfqpoint{2.510163in}{1.234931in}}%
\pgfpathlineto{\pgfqpoint{2.507767in}{1.239655in}}%
\pgfpathlineto{\pgfqpoint{2.554976in}{1.247539in}}%
\pgfpathlineto{\pgfqpoint{2.560339in}{1.213292in}}%
\pgfpathlineto{\pgfqpoint{2.549426in}{1.216536in}}%
\pgfpathlineto{\pgfqpoint{2.543719in}{1.212838in}}%
\pgfpathlineto{\pgfqpoint{2.538153in}{1.216728in}}%
\pgfpathlineto{\pgfqpoint{2.532780in}{1.212166in}}%
\pgfpathlineto{\pgfqpoint{2.525644in}{1.210611in}}%
\pgfpathlineto{\pgfqpoint{2.523965in}{1.201691in}}%
\pgfpathlineto{\pgfqpoint{2.520381in}{1.200511in}}%
\pgfpathlineto{\pgfqpoint{2.505727in}{1.202402in}}%
\pgfpathlineto{\pgfqpoint{2.498857in}{1.196503in}}%
\pgfpathclose%
\pgfusepath{fill}%
\end{pgfscope}%
\begin{pgfscope}%
\pgfpathrectangle{\pgfqpoint{0.100000in}{0.100000in}}{\pgfqpoint{3.007045in}{1.925000in}}%
\pgfusepath{clip}%
\pgfsetbuttcap%
\pgfsetmiterjoin%
\definecolor{currentfill}{rgb}{0.598431,0.780531,0.877785}%
\pgfsetfillcolor{currentfill}%
\pgfsetlinewidth{0.000000pt}%
\definecolor{currentstroke}{rgb}{0.000000,0.000000,0.000000}%
\pgfsetstrokecolor{currentstroke}%
\pgfsetstrokeopacity{0.000000}%
\pgfsetdash{}{0pt}%
\pgfpathmoveto{\pgfqpoint{2.677192in}{1.366767in}}%
\pgfpathlineto{\pgfqpoint{2.679466in}{1.361047in}}%
\pgfpathlineto{\pgfqpoint{2.686221in}{1.354519in}}%
\pgfpathlineto{\pgfqpoint{2.688600in}{1.346784in}}%
\pgfpathlineto{\pgfqpoint{2.692918in}{1.348382in}}%
\pgfpathlineto{\pgfqpoint{2.694956in}{1.344848in}}%
\pgfpathlineto{\pgfqpoint{2.680799in}{1.334029in}}%
\pgfpathlineto{\pgfqpoint{2.670437in}{1.328641in}}%
\pgfpathlineto{\pgfqpoint{2.664036in}{1.332084in}}%
\pgfpathlineto{\pgfqpoint{2.648248in}{1.330436in}}%
\pgfpathlineto{\pgfqpoint{2.645950in}{1.339764in}}%
\pgfpathlineto{\pgfqpoint{2.654171in}{1.354581in}}%
\pgfpathlineto{\pgfqpoint{2.661650in}{1.359044in}}%
\pgfpathlineto{\pgfqpoint{2.661755in}{1.363168in}}%
\pgfpathlineto{\pgfqpoint{2.669738in}{1.366595in}}%
\pgfpathlineto{\pgfqpoint{2.677192in}{1.366767in}}%
\pgfpathclose%
\pgfusepath{fill}%
\end{pgfscope}%
\begin{pgfscope}%
\pgfpathrectangle{\pgfqpoint{0.100000in}{0.100000in}}{\pgfqpoint{3.007045in}{1.925000in}}%
\pgfusepath{clip}%
\pgfsetbuttcap%
\pgfsetmiterjoin%
\definecolor{currentfill}{rgb}{0.567059,0.763306,0.871019}%
\pgfsetfillcolor{currentfill}%
\pgfsetlinewidth{0.000000pt}%
\definecolor{currentstroke}{rgb}{0.000000,0.000000,0.000000}%
\pgfsetstrokecolor{currentstroke}%
\pgfsetstrokeopacity{0.000000}%
\pgfsetdash{}{0pt}%
\pgfpathmoveto{\pgfqpoint{2.141224in}{0.664018in}}%
\pgfpathlineto{\pgfqpoint{2.139795in}{0.676959in}}%
\pgfpathlineto{\pgfqpoint{2.141241in}{0.722568in}}%
\pgfpathlineto{\pgfqpoint{2.141820in}{0.745927in}}%
\pgfpathlineto{\pgfqpoint{2.141932in}{0.750061in}}%
\pgfpathlineto{\pgfqpoint{2.151248in}{0.751063in}}%
\pgfpathlineto{\pgfqpoint{2.154101in}{0.754354in}}%
\pgfpathlineto{\pgfqpoint{2.158791in}{0.754738in}}%
\pgfpathlineto{\pgfqpoint{2.163129in}{0.759429in}}%
\pgfpathlineto{\pgfqpoint{2.168941in}{0.762907in}}%
\pgfpathlineto{\pgfqpoint{2.169944in}{0.753927in}}%
\pgfpathlineto{\pgfqpoint{2.176652in}{0.754493in}}%
\pgfpathlineto{\pgfqpoint{2.175967in}{0.749888in}}%
\pgfpathlineto{\pgfqpoint{2.172078in}{0.746133in}}%
\pgfpathlineto{\pgfqpoint{2.171759in}{0.737367in}}%
\pgfpathlineto{\pgfqpoint{2.174689in}{0.726922in}}%
\pgfpathlineto{\pgfqpoint{2.173613in}{0.721470in}}%
\pgfpathlineto{\pgfqpoint{2.178513in}{0.718958in}}%
\pgfpathlineto{\pgfqpoint{2.189892in}{0.719957in}}%
\pgfpathlineto{\pgfqpoint{2.190855in}{0.708411in}}%
\pgfpathlineto{\pgfqpoint{2.193997in}{0.705851in}}%
\pgfpathlineto{\pgfqpoint{2.191190in}{0.704948in}}%
\pgfpathlineto{\pgfqpoint{2.191937in}{0.696888in}}%
\pgfpathlineto{\pgfqpoint{2.186226in}{0.696389in}}%
\pgfpathlineto{\pgfqpoint{2.186978in}{0.688005in}}%
\pgfpathlineto{\pgfqpoint{2.184529in}{0.686866in}}%
\pgfpathlineto{\pgfqpoint{2.161771in}{0.684999in}}%
\pgfpathlineto{\pgfqpoint{2.156603in}{0.673169in}}%
\pgfpathlineto{\pgfqpoint{2.161161in}{0.669683in}}%
\pgfpathlineto{\pgfqpoint{2.162408in}{0.665683in}}%
\pgfpathlineto{\pgfqpoint{2.141224in}{0.664018in}}%
\pgfpathclose%
\pgfusepath{fill}%
\end{pgfscope}%
\begin{pgfscope}%
\pgfpathrectangle{\pgfqpoint{0.100000in}{0.100000in}}{\pgfqpoint{3.007045in}{1.925000in}}%
\pgfusepath{clip}%
\pgfsetbuttcap%
\pgfsetmiterjoin%
\definecolor{currentfill}{rgb}{0.412042,0.677186,0.836263}%
\pgfsetfillcolor{currentfill}%
\pgfsetlinewidth{0.000000pt}%
\definecolor{currentstroke}{rgb}{0.000000,0.000000,0.000000}%
\pgfsetstrokecolor{currentstroke}%
\pgfsetstrokeopacity{0.000000}%
\pgfsetdash{}{0pt}%
\pgfpathmoveto{\pgfqpoint{1.984723in}{0.815682in}}%
\pgfpathlineto{\pgfqpoint{1.980611in}{0.811051in}}%
\pgfpathlineto{\pgfqpoint{1.984377in}{0.804898in}}%
\pgfpathlineto{\pgfqpoint{1.974688in}{0.804739in}}%
\pgfpathlineto{\pgfqpoint{1.967563in}{0.808897in}}%
\pgfpathlineto{\pgfqpoint{1.966388in}{0.815925in}}%
\pgfpathlineto{\pgfqpoint{1.962924in}{0.816378in}}%
\pgfpathlineto{\pgfqpoint{1.961337in}{0.821725in}}%
\pgfpathlineto{\pgfqpoint{1.951971in}{0.821985in}}%
\pgfpathlineto{\pgfqpoint{1.951736in}{0.838426in}}%
\pgfpathlineto{\pgfqpoint{1.955232in}{0.842505in}}%
\pgfpathlineto{\pgfqpoint{1.955013in}{0.847952in}}%
\pgfpathlineto{\pgfqpoint{1.951519in}{0.850943in}}%
\pgfpathlineto{\pgfqpoint{1.951312in}{0.867300in}}%
\pgfpathlineto{\pgfqpoint{1.945345in}{0.867212in}}%
\pgfpathlineto{\pgfqpoint{1.945167in}{0.874395in}}%
\pgfpathlineto{\pgfqpoint{1.950043in}{0.876695in}}%
\pgfpathlineto{\pgfqpoint{1.956926in}{0.874466in}}%
\pgfpathlineto{\pgfqpoint{1.956817in}{0.878848in}}%
\pgfpathlineto{\pgfqpoint{1.963244in}{0.878965in}}%
\pgfpathlineto{\pgfqpoint{1.968811in}{0.876703in}}%
\pgfpathlineto{\pgfqpoint{1.972810in}{0.872450in}}%
\pgfpathlineto{\pgfqpoint{1.980505in}{0.873918in}}%
\pgfpathlineto{\pgfqpoint{1.980621in}{0.867935in}}%
\pgfpathlineto{\pgfqpoint{1.983593in}{0.865113in}}%
\pgfpathlineto{\pgfqpoint{1.983824in}{0.856375in}}%
\pgfpathlineto{\pgfqpoint{1.986708in}{0.856447in}}%
\pgfpathlineto{\pgfqpoint{1.987166in}{0.839058in}}%
\pgfpathlineto{\pgfqpoint{1.990126in}{0.839097in}}%
\pgfpathlineto{\pgfqpoint{1.990401in}{0.830295in}}%
\pgfpathlineto{\pgfqpoint{1.987469in}{0.830190in}}%
\pgfpathlineto{\pgfqpoint{1.986497in}{0.818543in}}%
\pgfpathlineto{\pgfqpoint{1.984723in}{0.815682in}}%
\pgfpathclose%
\pgfusepath{fill}%
\end{pgfscope}%
\begin{pgfscope}%
\pgfpathrectangle{\pgfqpoint{0.100000in}{0.100000in}}{\pgfqpoint{3.007045in}{1.925000in}}%
\pgfusepath{clip}%
\pgfsetbuttcap%
\pgfsetmiterjoin%
\definecolor{currentfill}{rgb}{0.401953,0.670296,0.832326}%
\pgfsetfillcolor{currentfill}%
\pgfsetlinewidth{0.000000pt}%
\definecolor{currentstroke}{rgb}{0.000000,0.000000,0.000000}%
\pgfsetstrokecolor{currentstroke}%
\pgfsetstrokeopacity{0.000000}%
\pgfsetdash{}{0pt}%
\pgfpathmoveto{\pgfqpoint{1.408374in}{1.098989in}}%
\pgfpathlineto{\pgfqpoint{1.372398in}{1.101802in}}%
\pgfpathlineto{\pgfqpoint{1.333267in}{1.105178in}}%
\pgfpathlineto{\pgfqpoint{1.334664in}{1.121680in}}%
\pgfpathlineto{\pgfqpoint{1.351745in}{1.120688in}}%
\pgfpathlineto{\pgfqpoint{1.353166in}{1.131470in}}%
\pgfpathlineto{\pgfqpoint{1.355807in}{1.160158in}}%
\pgfpathlineto{\pgfqpoint{1.358836in}{1.188865in}}%
\pgfpathlineto{\pgfqpoint{1.414767in}{1.184726in}}%
\pgfpathlineto{\pgfqpoint{1.414739in}{1.184340in}}%
\pgfpathlineto{\pgfqpoint{1.410528in}{1.127149in}}%
\pgfpathlineto{\pgfqpoint{1.408374in}{1.098989in}}%
\pgfpathclose%
\pgfusepath{fill}%
\end{pgfscope}%
\begin{pgfscope}%
\pgfpathrectangle{\pgfqpoint{0.100000in}{0.100000in}}{\pgfqpoint{3.007045in}{1.925000in}}%
\pgfusepath{clip}%
\pgfsetbuttcap%
\pgfsetmiterjoin%
\definecolor{currentfill}{rgb}{0.248166,0.561892,0.770980}%
\pgfsetfillcolor{currentfill}%
\pgfsetlinewidth{0.000000pt}%
\definecolor{currentstroke}{rgb}{0.000000,0.000000,0.000000}%
\pgfsetstrokecolor{currentstroke}%
\pgfsetstrokeopacity{0.000000}%
\pgfsetdash{}{0pt}%
\pgfpathmoveto{\pgfqpoint{2.013407in}{0.843839in}}%
\pgfpathlineto{\pgfqpoint{2.013158in}{0.836265in}}%
\pgfpathlineto{\pgfqpoint{2.003733in}{0.832951in}}%
\pgfpathlineto{\pgfqpoint{2.003512in}{0.827202in}}%
\pgfpathlineto{\pgfqpoint{1.993574in}{0.816135in}}%
\pgfpathlineto{\pgfqpoint{1.993621in}{0.816052in}}%
\pgfpathlineto{\pgfqpoint{1.984723in}{0.815682in}}%
\pgfpathlineto{\pgfqpoint{1.986497in}{0.818543in}}%
\pgfpathlineto{\pgfqpoint{1.987469in}{0.830190in}}%
\pgfpathlineto{\pgfqpoint{1.990401in}{0.830295in}}%
\pgfpathlineto{\pgfqpoint{1.990126in}{0.839097in}}%
\pgfpathlineto{\pgfqpoint{1.987166in}{0.839058in}}%
\pgfpathlineto{\pgfqpoint{1.986708in}{0.856447in}}%
\pgfpathlineto{\pgfqpoint{1.983824in}{0.856375in}}%
\pgfpathlineto{\pgfqpoint{1.983593in}{0.865113in}}%
\pgfpathlineto{\pgfqpoint{1.980621in}{0.867935in}}%
\pgfpathlineto{\pgfqpoint{1.980505in}{0.873918in}}%
\pgfpathlineto{\pgfqpoint{1.986161in}{0.874058in}}%
\pgfpathlineto{\pgfqpoint{1.985207in}{0.903228in}}%
\pgfpathlineto{\pgfqpoint{2.008210in}{0.904668in}}%
\pgfpathlineto{\pgfqpoint{2.025623in}{0.904867in}}%
\pgfpathlineto{\pgfqpoint{2.033514in}{0.905154in}}%
\pgfpathlineto{\pgfqpoint{2.037206in}{0.901885in}}%
\pgfpathlineto{\pgfqpoint{2.034157in}{0.893004in}}%
\pgfpathlineto{\pgfqpoint{2.037839in}{0.890956in}}%
\pgfpathlineto{\pgfqpoint{2.036782in}{0.884180in}}%
\pgfpathlineto{\pgfqpoint{2.032840in}{0.884060in}}%
\pgfpathlineto{\pgfqpoint{2.030962in}{0.877859in}}%
\pgfpathlineto{\pgfqpoint{2.026701in}{0.874367in}}%
\pgfpathlineto{\pgfqpoint{2.029442in}{0.869945in}}%
\pgfpathlineto{\pgfqpoint{2.020371in}{0.864609in}}%
\pgfpathlineto{\pgfqpoint{2.018227in}{0.853487in}}%
\pgfpathlineto{\pgfqpoint{2.012143in}{0.849788in}}%
\pgfpathlineto{\pgfqpoint{2.013407in}{0.843839in}}%
\pgfpathclose%
\pgfusepath{fill}%
\end{pgfscope}%
\begin{pgfscope}%
\pgfpathrectangle{\pgfqpoint{0.100000in}{0.100000in}}{\pgfqpoint{3.007045in}{1.925000in}}%
\pgfusepath{clip}%
\pgfsetbuttcap%
\pgfsetmiterjoin%
\definecolor{currentfill}{rgb}{0.376732,0.653072,0.822484}%
\pgfsetfillcolor{currentfill}%
\pgfsetlinewidth{0.000000pt}%
\definecolor{currentstroke}{rgb}{0.000000,0.000000,0.000000}%
\pgfsetstrokecolor{currentstroke}%
\pgfsetstrokeopacity{0.000000}%
\pgfsetdash{}{0pt}%
\pgfpathmoveto{\pgfqpoint{2.283647in}{1.145438in}}%
\pgfpathlineto{\pgfqpoint{2.275093in}{1.139834in}}%
\pgfpathlineto{\pgfqpoint{2.273431in}{1.154396in}}%
\pgfpathlineto{\pgfqpoint{2.261118in}{1.153025in}}%
\pgfpathlineto{\pgfqpoint{2.259410in}{1.171566in}}%
\pgfpathlineto{\pgfqpoint{2.253429in}{1.166775in}}%
\pgfpathlineto{\pgfqpoint{2.247318in}{1.166002in}}%
\pgfpathlineto{\pgfqpoint{2.245949in}{1.180376in}}%
\pgfpathlineto{\pgfqpoint{2.248039in}{1.187423in}}%
\pgfpathlineto{\pgfqpoint{2.264834in}{1.189201in}}%
\pgfpathlineto{\pgfqpoint{2.264298in}{1.193927in}}%
\pgfpathlineto{\pgfqpoint{2.288706in}{1.196382in}}%
\pgfpathlineto{\pgfqpoint{2.291530in}{1.169095in}}%
\pgfpathlineto{\pgfqpoint{2.288006in}{1.165183in}}%
\pgfpathlineto{\pgfqpoint{2.291496in}{1.161362in}}%
\pgfpathlineto{\pgfqpoint{2.289964in}{1.156159in}}%
\pgfpathlineto{\pgfqpoint{2.294975in}{1.154537in}}%
\pgfpathlineto{\pgfqpoint{2.290373in}{1.148367in}}%
\pgfpathlineto{\pgfqpoint{2.283647in}{1.145438in}}%
\pgfpathclose%
\pgfusepath{fill}%
\end{pgfscope}%
\begin{pgfscope}%
\pgfpathrectangle{\pgfqpoint{0.100000in}{0.100000in}}{\pgfqpoint{3.007045in}{1.925000in}}%
\pgfusepath{clip}%
\pgfsetbuttcap%
\pgfsetmiterjoin%
\definecolor{currentfill}{rgb}{0.435294,0.690965,0.842599}%
\pgfsetfillcolor{currentfill}%
\pgfsetlinewidth{0.000000pt}%
\definecolor{currentstroke}{rgb}{0.000000,0.000000,0.000000}%
\pgfsetstrokecolor{currentstroke}%
\pgfsetstrokeopacity{0.000000}%
\pgfsetdash{}{0pt}%
\pgfpathmoveto{\pgfqpoint{1.440696in}{0.803294in}}%
\pgfpathlineto{\pgfqpoint{1.469347in}{0.801604in}}%
\pgfpathlineto{\pgfqpoint{1.495136in}{0.800124in}}%
\pgfpathlineto{\pgfqpoint{1.514372in}{0.798412in}}%
\pgfpathlineto{\pgfqpoint{1.525327in}{0.794045in}}%
\pgfpathlineto{\pgfqpoint{1.526290in}{0.791703in}}%
\pgfpathlineto{\pgfqpoint{1.525314in}{0.769593in}}%
\pgfpathlineto{\pgfqpoint{1.523941in}{0.740360in}}%
\pgfpathlineto{\pgfqpoint{1.486605in}{0.742533in}}%
\pgfpathlineto{\pgfqpoint{1.465814in}{0.743950in}}%
\pgfpathlineto{\pgfqpoint{1.467658in}{0.772689in}}%
\pgfpathlineto{\pgfqpoint{1.438908in}{0.774476in}}%
\pgfpathlineto{\pgfqpoint{1.440696in}{0.803294in}}%
\pgfpathclose%
\pgfusepath{fill}%
\end{pgfscope}%
\begin{pgfscope}%
\pgfpathrectangle{\pgfqpoint{0.100000in}{0.100000in}}{\pgfqpoint{3.007045in}{1.925000in}}%
\pgfusepath{clip}%
\pgfsetbuttcap%
\pgfsetmiterjoin%
\definecolor{currentfill}{rgb}{0.341423,0.628958,0.808704}%
\pgfsetfillcolor{currentfill}%
\pgfsetlinewidth{0.000000pt}%
\definecolor{currentstroke}{rgb}{0.000000,0.000000,0.000000}%
\pgfsetstrokecolor{currentstroke}%
\pgfsetstrokeopacity{0.000000}%
\pgfsetdash{}{0pt}%
\pgfpathmoveto{\pgfqpoint{1.963991in}{0.607879in}}%
\pgfpathlineto{\pgfqpoint{1.958935in}{0.610678in}}%
\pgfpathlineto{\pgfqpoint{1.958991in}{0.616692in}}%
\pgfpathlineto{\pgfqpoint{1.961233in}{0.620328in}}%
\pgfpathlineto{\pgfqpoint{1.952131in}{0.625210in}}%
\pgfpathlineto{\pgfqpoint{1.949979in}{0.638422in}}%
\pgfpathlineto{\pgfqpoint{1.954469in}{0.643065in}}%
\pgfpathlineto{\pgfqpoint{1.959044in}{0.656761in}}%
\pgfpathlineto{\pgfqpoint{1.962071in}{0.656269in}}%
\pgfpathlineto{\pgfqpoint{1.965330in}{0.661044in}}%
\pgfpathlineto{\pgfqpoint{1.965268in}{0.666642in}}%
\pgfpathlineto{\pgfqpoint{1.969490in}{0.679240in}}%
\pgfpathlineto{\pgfqpoint{1.969121in}{0.688091in}}%
\pgfpathlineto{\pgfqpoint{1.978768in}{0.688413in}}%
\pgfpathlineto{\pgfqpoint{1.989120in}{0.691955in}}%
\pgfpathlineto{\pgfqpoint{1.987667in}{0.684255in}}%
\pgfpathlineto{\pgfqpoint{1.995075in}{0.682326in}}%
\pgfpathlineto{\pgfqpoint{1.993029in}{0.676683in}}%
\pgfpathlineto{\pgfqpoint{1.986876in}{0.672737in}}%
\pgfpathlineto{\pgfqpoint{1.986197in}{0.669162in}}%
\pgfpathlineto{\pgfqpoint{1.978327in}{0.664514in}}%
\pgfpathlineto{\pgfqpoint{1.981608in}{0.653144in}}%
\pgfpathlineto{\pgfqpoint{1.989789in}{0.649517in}}%
\pgfpathlineto{\pgfqpoint{2.013321in}{0.650638in}}%
\pgfpathlineto{\pgfqpoint{2.019201in}{0.650937in}}%
\pgfpathlineto{\pgfqpoint{2.020031in}{0.633536in}}%
\pgfpathlineto{\pgfqpoint{1.998230in}{0.632422in}}%
\pgfpathlineto{\pgfqpoint{1.993008in}{0.630363in}}%
\pgfpathlineto{\pgfqpoint{1.985698in}{0.633433in}}%
\pgfpathlineto{\pgfqpoint{1.972763in}{0.629511in}}%
\pgfpathlineto{\pgfqpoint{1.971855in}{0.623291in}}%
\pgfpathlineto{\pgfqpoint{1.967749in}{0.623734in}}%
\pgfpathlineto{\pgfqpoint{1.964384in}{0.616831in}}%
\pgfpathlineto{\pgfqpoint{1.967968in}{0.612514in}}%
\pgfpathlineto{\pgfqpoint{1.963991in}{0.607879in}}%
\pgfpathclose%
\pgfusepath{fill}%
\end{pgfscope}%
\begin{pgfscope}%
\pgfpathrectangle{\pgfqpoint{0.100000in}{0.100000in}}{\pgfqpoint{3.007045in}{1.925000in}}%
\pgfusepath{clip}%
\pgfsetbuttcap%
\pgfsetmiterjoin%
\definecolor{currentfill}{rgb}{0.401953,0.670296,0.832326}%
\pgfsetfillcolor{currentfill}%
\pgfsetlinewidth{0.000000pt}%
\definecolor{currentstroke}{rgb}{0.000000,0.000000,0.000000}%
\pgfsetstrokecolor{currentstroke}%
\pgfsetstrokeopacity{0.000000}%
\pgfsetdash{}{0pt}%
\pgfpathmoveto{\pgfqpoint{1.990535in}{1.152936in}}%
\pgfpathlineto{\pgfqpoint{1.967023in}{1.152094in}}%
\pgfpathlineto{\pgfqpoint{1.967135in}{1.146275in}}%
\pgfpathlineto{\pgfqpoint{1.963320in}{1.146179in}}%
\pgfpathlineto{\pgfqpoint{1.955717in}{1.146050in}}%
\pgfpathlineto{\pgfqpoint{1.953772in}{1.157587in}}%
\pgfpathlineto{\pgfqpoint{1.939480in}{1.158394in}}%
\pgfpathlineto{\pgfqpoint{1.938875in}{1.179366in}}%
\pgfpathlineto{\pgfqpoint{1.932193in}{1.179188in}}%
\pgfpathlineto{\pgfqpoint{1.931854in}{1.198270in}}%
\pgfpathlineto{\pgfqpoint{1.926181in}{1.198155in}}%
\pgfpathlineto{\pgfqpoint{1.925598in}{1.218425in}}%
\pgfpathlineto{\pgfqpoint{1.948406in}{1.218688in}}%
\pgfpathlineto{\pgfqpoint{1.948122in}{1.215494in}}%
\pgfpathlineto{\pgfqpoint{1.977971in}{1.216245in}}%
\pgfpathlineto{\pgfqpoint{1.978770in}{1.193340in}}%
\pgfpathlineto{\pgfqpoint{1.996210in}{1.193778in}}%
\pgfpathlineto{\pgfqpoint{1.992899in}{1.184317in}}%
\pgfpathlineto{\pgfqpoint{1.997146in}{1.174029in}}%
\pgfpathlineto{\pgfqpoint{1.995501in}{1.164493in}}%
\pgfpathlineto{\pgfqpoint{1.979066in}{1.163974in}}%
\pgfpathlineto{\pgfqpoint{1.986624in}{1.158384in}}%
\pgfpathlineto{\pgfqpoint{1.990535in}{1.152936in}}%
\pgfpathclose%
\pgfusepath{fill}%
\end{pgfscope}%
\begin{pgfscope}%
\pgfpathrectangle{\pgfqpoint{0.100000in}{0.100000in}}{\pgfqpoint{3.007045in}{1.925000in}}%
\pgfusepath{clip}%
\pgfsetbuttcap%
\pgfsetmiterjoin%
\definecolor{currentfill}{rgb}{0.622684,0.793464,0.883429}%
\pgfsetfillcolor{currentfill}%
\pgfsetlinewidth{0.000000pt}%
\definecolor{currentstroke}{rgb}{0.000000,0.000000,0.000000}%
\pgfsetstrokecolor{currentstroke}%
\pgfsetstrokeopacity{0.000000}%
\pgfsetdash{}{0pt}%
\pgfpathmoveto{\pgfqpoint{0.856506in}{1.901711in}}%
\pgfpathlineto{\pgfqpoint{0.848850in}{1.869247in}}%
\pgfpathlineto{\pgfqpoint{0.817161in}{1.876758in}}%
\pgfpathlineto{\pgfqpoint{0.822653in}{1.899281in}}%
\pgfpathlineto{\pgfqpoint{0.812070in}{1.901824in}}%
\pgfpathlineto{\pgfqpoint{0.814525in}{1.911735in}}%
\pgfpathlineto{\pgfqpoint{0.856506in}{1.901711in}}%
\pgfpathclose%
\pgfusepath{fill}%
\end{pgfscope}%
\begin{pgfscope}%
\pgfpathrectangle{\pgfqpoint{0.100000in}{0.100000in}}{\pgfqpoint{3.007045in}{1.925000in}}%
\pgfusepath{clip}%
\pgfsetbuttcap%
\pgfsetmiterjoin%
\definecolor{currentfill}{rgb}{0.460392,0.704744,0.848012}%
\pgfsetfillcolor{currentfill}%
\pgfsetlinewidth{0.000000pt}%
\definecolor{currentstroke}{rgb}{0.000000,0.000000,0.000000}%
\pgfsetstrokecolor{currentstroke}%
\pgfsetstrokeopacity{0.000000}%
\pgfsetdash{}{0pt}%
\pgfpathmoveto{\pgfqpoint{2.349537in}{1.370183in}}%
\pgfpathlineto{\pgfqpoint{2.343612in}{1.373779in}}%
\pgfpathlineto{\pgfqpoint{2.320460in}{1.370056in}}%
\pgfpathlineto{\pgfqpoint{2.303185in}{1.367262in}}%
\pgfpathlineto{\pgfqpoint{2.299946in}{1.390294in}}%
\pgfpathlineto{\pgfqpoint{2.296340in}{1.413257in}}%
\pgfpathlineto{\pgfqpoint{2.286455in}{1.412044in}}%
\pgfpathlineto{\pgfqpoint{2.283625in}{1.434436in}}%
\pgfpathlineto{\pgfqpoint{2.283546in}{1.435118in}}%
\pgfpathlineto{\pgfqpoint{2.304504in}{1.437926in}}%
\pgfpathlineto{\pgfqpoint{2.303651in}{1.443639in}}%
\pgfpathlineto{\pgfqpoint{2.326140in}{1.446645in}}%
\pgfpathlineto{\pgfqpoint{2.331232in}{1.448101in}}%
\pgfpathlineto{\pgfqpoint{2.330370in}{1.453796in}}%
\pgfpathlineto{\pgfqpoint{2.342275in}{1.455683in}}%
\pgfpathlineto{\pgfqpoint{2.338611in}{1.478397in}}%
\pgfpathlineto{\pgfqpoint{2.362681in}{1.482619in}}%
\pgfpathlineto{\pgfqpoint{2.364694in}{1.474272in}}%
\pgfpathlineto{\pgfqpoint{2.368291in}{1.465466in}}%
\pgfpathlineto{\pgfqpoint{2.371281in}{1.451948in}}%
\pgfpathlineto{\pgfqpoint{2.373711in}{1.444991in}}%
\pgfpathlineto{\pgfqpoint{2.377887in}{1.438966in}}%
\pgfpathlineto{\pgfqpoint{2.376582in}{1.425757in}}%
\pgfpathlineto{\pgfqpoint{2.377704in}{1.422799in}}%
\pgfpathlineto{\pgfqpoint{2.376418in}{1.412332in}}%
\pgfpathlineto{\pgfqpoint{2.367395in}{1.416634in}}%
\pgfpathlineto{\pgfqpoint{2.363048in}{1.413679in}}%
\pgfpathlineto{\pgfqpoint{2.360076in}{1.403017in}}%
\pgfpathlineto{\pgfqpoint{2.361096in}{1.399839in}}%
\pgfpathlineto{\pgfqpoint{2.359012in}{1.392852in}}%
\pgfpathlineto{\pgfqpoint{2.352893in}{1.389860in}}%
\pgfpathlineto{\pgfqpoint{2.350480in}{1.384293in}}%
\pgfpathlineto{\pgfqpoint{2.351493in}{1.374397in}}%
\pgfpathlineto{\pgfqpoint{2.349537in}{1.370183in}}%
\pgfpathclose%
\pgfusepath{fill}%
\end{pgfscope}%
\begin{pgfscope}%
\pgfpathrectangle{\pgfqpoint{0.100000in}{0.100000in}}{\pgfqpoint{3.007045in}{1.925000in}}%
\pgfusepath{clip}%
\pgfsetbuttcap%
\pgfsetmiterjoin%
\definecolor{currentfill}{rgb}{0.681738,0.818562,0.904098}%
\pgfsetfillcolor{currentfill}%
\pgfsetlinewidth{0.000000pt}%
\definecolor{currentstroke}{rgb}{0.000000,0.000000,0.000000}%
\pgfsetstrokecolor{currentstroke}%
\pgfsetstrokeopacity{0.000000}%
\pgfsetdash{}{0pt}%
\pgfpathmoveto{\pgfqpoint{0.652693in}{1.955469in}}%
\pgfpathlineto{\pgfqpoint{0.654305in}{1.933747in}}%
\pgfpathlineto{\pgfqpoint{0.650462in}{1.932235in}}%
\pgfpathlineto{\pgfqpoint{0.643432in}{1.933133in}}%
\pgfpathlineto{\pgfqpoint{0.623985in}{1.938935in}}%
\pgfpathlineto{\pgfqpoint{0.576681in}{1.953612in}}%
\pgfpathlineto{\pgfqpoint{0.577191in}{1.962609in}}%
\pgfpathlineto{\pgfqpoint{0.572821in}{1.961633in}}%
\pgfpathlineto{\pgfqpoint{0.569082in}{1.972907in}}%
\pgfpathlineto{\pgfqpoint{0.572639in}{1.979931in}}%
\pgfpathlineto{\pgfqpoint{0.614766in}{1.966556in}}%
\pgfpathlineto{\pgfqpoint{0.652693in}{1.955469in}}%
\pgfpathclose%
\pgfusepath{fill}%
\end{pgfscope}%
\begin{pgfscope}%
\pgfpathrectangle{\pgfqpoint{0.100000in}{0.100000in}}{\pgfqpoint{3.007045in}{1.925000in}}%
\pgfusepath{clip}%
\pgfsetbuttcap%
\pgfsetmiterjoin%
\definecolor{currentfill}{rgb}{0.199446,0.513172,0.745882}%
\pgfsetfillcolor{currentfill}%
\pgfsetlinewidth{0.000000pt}%
\definecolor{currentstroke}{rgb}{0.000000,0.000000,0.000000}%
\pgfsetstrokecolor{currentstroke}%
\pgfsetstrokeopacity{0.000000}%
\pgfsetdash{}{0pt}%
\pgfpathmoveto{\pgfqpoint{1.035972in}{1.058230in}}%
\pgfpathlineto{\pgfqpoint{1.025816in}{0.993300in}}%
\pgfpathlineto{\pgfqpoint{1.100725in}{0.982093in}}%
\pgfpathlineto{\pgfqpoint{1.117536in}{0.979664in}}%
\pgfpathlineto{\pgfqpoint{1.111260in}{0.934336in}}%
\pgfpathlineto{\pgfqpoint{1.049496in}{0.943259in}}%
\pgfpathlineto{\pgfqpoint{1.046069in}{0.920555in}}%
\pgfpathlineto{\pgfqpoint{1.015230in}{0.925309in}}%
\pgfpathlineto{\pgfqpoint{1.011345in}{0.900484in}}%
\pgfpathlineto{\pgfqpoint{1.003150in}{0.848168in}}%
\pgfpathlineto{\pgfqpoint{0.986809in}{0.850694in}}%
\pgfpathlineto{\pgfqpoint{0.978685in}{0.845988in}}%
\pgfpathlineto{\pgfqpoint{0.975090in}{0.841991in}}%
\pgfpathlineto{\pgfqpoint{0.971108in}{0.842140in}}%
\pgfpathlineto{\pgfqpoint{0.960834in}{0.834656in}}%
\pgfpathlineto{\pgfqpoint{0.954625in}{0.842378in}}%
\pgfpathlineto{\pgfqpoint{0.949199in}{0.843379in}}%
\pgfpathlineto{\pgfqpoint{0.953797in}{0.870982in}}%
\pgfpathlineto{\pgfqpoint{0.913272in}{0.877904in}}%
\pgfpathlineto{\pgfqpoint{0.916273in}{0.895048in}}%
\pgfpathlineto{\pgfqpoint{0.947453in}{1.073147in}}%
\pgfpathlineto{\pgfqpoint{0.986309in}{1.066160in}}%
\pgfpathlineto{\pgfqpoint{1.035972in}{1.058230in}}%
\pgfpathclose%
\pgfusepath{fill}%
\end{pgfscope}%
\begin{pgfscope}%
\pgfpathrectangle{\pgfqpoint{0.100000in}{0.100000in}}{\pgfqpoint{3.007045in}{1.925000in}}%
\pgfusepath{clip}%
\pgfsetbuttcap%
\pgfsetmiterjoin%
\definecolor{currentfill}{rgb}{0.548235,0.752972,0.866959}%
\pgfsetfillcolor{currentfill}%
\pgfsetlinewidth{0.000000pt}%
\definecolor{currentstroke}{rgb}{0.000000,0.000000,0.000000}%
\pgfsetstrokecolor{currentstroke}%
\pgfsetstrokeopacity{0.000000}%
\pgfsetdash{}{0pt}%
\pgfpathmoveto{\pgfqpoint{1.907975in}{1.607918in}}%
\pgfpathlineto{\pgfqpoint{1.930881in}{1.608680in}}%
\pgfpathlineto{\pgfqpoint{1.959683in}{1.609718in}}%
\pgfpathlineto{\pgfqpoint{1.960172in}{1.598201in}}%
\pgfpathlineto{\pgfqpoint{1.971604in}{1.598763in}}%
\pgfpathlineto{\pgfqpoint{1.973425in}{1.558760in}}%
\pgfpathlineto{\pgfqpoint{1.961960in}{1.558341in}}%
\pgfpathlineto{\pgfqpoint{1.962214in}{1.552559in}}%
\pgfpathlineto{\pgfqpoint{1.933521in}{1.551363in}}%
\pgfpathlineto{\pgfqpoint{1.933694in}{1.545701in}}%
\pgfpathlineto{\pgfqpoint{1.905032in}{1.544903in}}%
\pgfpathlineto{\pgfqpoint{1.904215in}{1.573411in}}%
\pgfpathlineto{\pgfqpoint{1.909946in}{1.573607in}}%
\pgfpathlineto{\pgfqpoint{1.909099in}{1.596409in}}%
\pgfpathlineto{\pgfqpoint{1.907975in}{1.607918in}}%
\pgfpathclose%
\pgfusepath{fill}%
\end{pgfscope}%
\begin{pgfscope}%
\pgfpathrectangle{\pgfqpoint{0.100000in}{0.100000in}}{\pgfqpoint{3.007045in}{1.925000in}}%
\pgfusepath{clip}%
\pgfsetbuttcap%
\pgfsetmiterjoin%
\definecolor{currentfill}{rgb}{0.504314,0.728858,0.857486}%
\pgfsetfillcolor{currentfill}%
\pgfsetlinewidth{0.000000pt}%
\definecolor{currentstroke}{rgb}{0.000000,0.000000,0.000000}%
\pgfsetstrokecolor{currentstroke}%
\pgfsetstrokeopacity{0.000000}%
\pgfsetdash{}{0pt}%
\pgfpathmoveto{\pgfqpoint{1.454083in}{0.686391in}}%
\pgfpathlineto{\pgfqpoint{1.425000in}{0.688410in}}%
\pgfpathlineto{\pgfqpoint{1.427496in}{0.717282in}}%
\pgfpathlineto{\pgfqpoint{1.429539in}{0.746331in}}%
\pgfpathlineto{\pgfqpoint{1.437038in}{0.746267in}}%
\pgfpathlineto{\pgfqpoint{1.438908in}{0.774476in}}%
\pgfpathlineto{\pgfqpoint{1.467658in}{0.772689in}}%
\pgfpathlineto{\pgfqpoint{1.465814in}{0.743950in}}%
\pgfpathlineto{\pgfqpoint{1.458235in}{0.744453in}}%
\pgfpathlineto{\pgfqpoint{1.454083in}{0.686391in}}%
\pgfpathclose%
\pgfusepath{fill}%
\end{pgfscope}%
\begin{pgfscope}%
\pgfpathrectangle{\pgfqpoint{0.100000in}{0.100000in}}{\pgfqpoint{3.007045in}{1.925000in}}%
\pgfusepath{clip}%
\pgfsetbuttcap%
\pgfsetmiterjoin%
\definecolor{currentfill}{rgb}{0.356555,0.639293,0.814610}%
\pgfsetfillcolor{currentfill}%
\pgfsetlinewidth{0.000000pt}%
\definecolor{currentstroke}{rgb}{0.000000,0.000000,0.000000}%
\pgfsetstrokecolor{currentstroke}%
\pgfsetstrokeopacity{0.000000}%
\pgfsetdash{}{0pt}%
\pgfpathmoveto{\pgfqpoint{2.079034in}{0.765390in}}%
\pgfpathlineto{\pgfqpoint{2.086624in}{0.765838in}}%
\pgfpathlineto{\pgfqpoint{2.087936in}{0.742436in}}%
\pgfpathlineto{\pgfqpoint{2.065136in}{0.740946in}}%
\pgfpathlineto{\pgfqpoint{2.065317in}{0.737960in}}%
\pgfpathlineto{\pgfqpoint{2.052237in}{0.736873in}}%
\pgfpathlineto{\pgfqpoint{2.031102in}{0.744404in}}%
\pgfpathlineto{\pgfqpoint{2.025625in}{0.749814in}}%
\pgfpathlineto{\pgfqpoint{2.031123in}{0.752696in}}%
\pgfpathlineto{\pgfqpoint{2.030481in}{0.763778in}}%
\pgfpathlineto{\pgfqpoint{2.033226in}{0.766351in}}%
\pgfpathlineto{\pgfqpoint{2.039047in}{0.766251in}}%
\pgfpathlineto{\pgfqpoint{2.062975in}{0.759697in}}%
\pgfpathlineto{\pgfqpoint{2.065349in}{0.764544in}}%
\pgfpathlineto{\pgfqpoint{2.079034in}{0.765390in}}%
\pgfpathclose%
\pgfusepath{fill}%
\end{pgfscope}%
\begin{pgfscope}%
\pgfpathrectangle{\pgfqpoint{0.100000in}{0.100000in}}{\pgfqpoint{3.007045in}{1.925000in}}%
\pgfusepath{clip}%
\pgfsetbuttcap%
\pgfsetmiterjoin%
\definecolor{currentfill}{rgb}{0.554510,0.756417,0.868312}%
\pgfsetfillcolor{currentfill}%
\pgfsetlinewidth{0.000000pt}%
\definecolor{currentstroke}{rgb}{0.000000,0.000000,0.000000}%
\pgfsetstrokecolor{currentstroke}%
\pgfsetstrokeopacity{0.000000}%
\pgfsetdash{}{0pt}%
\pgfpathmoveto{\pgfqpoint{1.973511in}{1.274071in}}%
\pgfpathlineto{\pgfqpoint{1.952591in}{1.273282in}}%
\pgfpathlineto{\pgfqpoint{1.952380in}{1.279020in}}%
\pgfpathlineto{\pgfqpoint{1.935244in}{1.278471in}}%
\pgfpathlineto{\pgfqpoint{1.923843in}{1.278117in}}%
\pgfpathlineto{\pgfqpoint{1.923061in}{1.300995in}}%
\pgfpathlineto{\pgfqpoint{1.928815in}{1.301237in}}%
\pgfpathlineto{\pgfqpoint{1.927902in}{1.324319in}}%
\pgfpathlineto{\pgfqpoint{1.950752in}{1.325057in}}%
\pgfpathlineto{\pgfqpoint{1.950534in}{1.330806in}}%
\pgfpathlineto{\pgfqpoint{1.973436in}{1.331705in}}%
\pgfpathlineto{\pgfqpoint{1.973965in}{1.320197in}}%
\pgfpathlineto{\pgfqpoint{1.974408in}{1.308715in}}%
\pgfpathlineto{\pgfqpoint{1.980128in}{1.308936in}}%
\pgfpathlineto{\pgfqpoint{1.980445in}{1.299472in}}%
\pgfpathlineto{\pgfqpoint{1.967742in}{1.296340in}}%
\pgfpathlineto{\pgfqpoint{1.964916in}{1.284730in}}%
\pgfpathlineto{\pgfqpoint{1.971270in}{1.279448in}}%
\pgfpathlineto{\pgfqpoint{1.973511in}{1.274071in}}%
\pgfpathclose%
\pgfusepath{fill}%
\end{pgfscope}%
\begin{pgfscope}%
\pgfpathrectangle{\pgfqpoint{0.100000in}{0.100000in}}{\pgfqpoint{3.007045in}{1.925000in}}%
\pgfusepath{clip}%
\pgfsetbuttcap%
\pgfsetmiterjoin%
\definecolor{currentfill}{rgb}{0.396909,0.666851,0.830358}%
\pgfsetfillcolor{currentfill}%
\pgfsetlinewidth{0.000000pt}%
\definecolor{currentstroke}{rgb}{0.000000,0.000000,0.000000}%
\pgfsetstrokecolor{currentstroke}%
\pgfsetstrokeopacity{0.000000}%
\pgfsetdash{}{0pt}%
\pgfpathmoveto{\pgfqpoint{1.765864in}{1.580217in}}%
\pgfpathlineto{\pgfqpoint{1.766199in}{1.556278in}}%
\pgfpathlineto{\pgfqpoint{1.720360in}{1.556564in}}%
\pgfpathlineto{\pgfqpoint{1.714573in}{1.556685in}}%
\pgfpathlineto{\pgfqpoint{1.714182in}{1.568146in}}%
\pgfpathlineto{\pgfqpoint{1.714476in}{1.591318in}}%
\pgfpathlineto{\pgfqpoint{1.714002in}{1.597084in}}%
\pgfpathlineto{\pgfqpoint{1.713584in}{1.614400in}}%
\pgfpathlineto{\pgfqpoint{1.713868in}{1.637553in}}%
\pgfpathlineto{\pgfqpoint{1.718828in}{1.637478in}}%
\pgfpathlineto{\pgfqpoint{1.718799in}{1.643280in}}%
\pgfpathlineto{\pgfqpoint{1.764876in}{1.642920in}}%
\pgfpathlineto{\pgfqpoint{1.765646in}{1.613993in}}%
\pgfpathlineto{\pgfqpoint{1.765864in}{1.580217in}}%
\pgfpathclose%
\pgfusepath{fill}%
\end{pgfscope}%
\begin{pgfscope}%
\pgfpathrectangle{\pgfqpoint{0.100000in}{0.100000in}}{\pgfqpoint{3.007045in}{1.925000in}}%
\pgfusepath{clip}%
\pgfsetbuttcap%
\pgfsetmiterjoin%
\definecolor{currentfill}{rgb}{0.346467,0.632403,0.810673}%
\pgfsetfillcolor{currentfill}%
\pgfsetlinewidth{0.000000pt}%
\definecolor{currentstroke}{rgb}{0.000000,0.000000,0.000000}%
\pgfsetstrokecolor{currentstroke}%
\pgfsetstrokeopacity{0.000000}%
\pgfsetdash{}{0pt}%
\pgfpathmoveto{\pgfqpoint{2.364743in}{1.174580in}}%
\pgfpathlineto{\pgfqpoint{2.353578in}{1.170749in}}%
\pgfpathlineto{\pgfqpoint{2.340324in}{1.169258in}}%
\pgfpathlineto{\pgfqpoint{2.338836in}{1.184076in}}%
\pgfpathlineto{\pgfqpoint{2.331627in}{1.183748in}}%
\pgfpathlineto{\pgfqpoint{2.330635in}{1.204428in}}%
\pgfpathlineto{\pgfqpoint{2.346250in}{1.205103in}}%
\pgfpathlineto{\pgfqpoint{2.345734in}{1.216099in}}%
\pgfpathlineto{\pgfqpoint{2.366064in}{1.217273in}}%
\pgfpathlineto{\pgfqpoint{2.366865in}{1.205475in}}%
\pgfpathlineto{\pgfqpoint{2.362689in}{1.195710in}}%
\pgfpathlineto{\pgfqpoint{2.363339in}{1.187311in}}%
\pgfpathlineto{\pgfqpoint{2.365382in}{1.186478in}}%
\pgfpathlineto{\pgfqpoint{2.364743in}{1.174580in}}%
\pgfpathclose%
\pgfusepath{fill}%
\end{pgfscope}%
\begin{pgfscope}%
\pgfpathrectangle{\pgfqpoint{0.100000in}{0.100000in}}{\pgfqpoint{3.007045in}{1.925000in}}%
\pgfusepath{clip}%
\pgfsetbuttcap%
\pgfsetmiterjoin%
\definecolor{currentfill}{rgb}{0.805260,0.878016,0.946851}%
\pgfsetfillcolor{currentfill}%
\pgfsetlinewidth{0.000000pt}%
\definecolor{currentstroke}{rgb}{0.000000,0.000000,0.000000}%
\pgfsetstrokecolor{currentstroke}%
\pgfsetstrokeopacity{0.000000}%
\pgfsetdash{}{0pt}%
\pgfpathmoveto{\pgfqpoint{3.013596in}{1.761178in}}%
\pgfpathlineto{\pgfqpoint{3.033217in}{1.758207in}}%
\pgfpathlineto{\pgfqpoint{3.030957in}{1.753391in}}%
\pgfpathlineto{\pgfqpoint{3.036928in}{1.745059in}}%
\pgfpathlineto{\pgfqpoint{3.036155in}{1.740067in}}%
\pgfpathlineto{\pgfqpoint{3.045228in}{1.732473in}}%
\pgfpathlineto{\pgfqpoint{3.046987in}{1.737585in}}%
\pgfpathlineto{\pgfqpoint{3.052581in}{1.737319in}}%
\pgfpathlineto{\pgfqpoint{3.064704in}{1.724060in}}%
\pgfpathlineto{\pgfqpoint{3.067422in}{1.717798in}}%
\pgfpathlineto{\pgfqpoint{3.063897in}{1.711796in}}%
\pgfpathlineto{\pgfqpoint{3.061530in}{1.704211in}}%
\pgfpathlineto{\pgfqpoint{3.052328in}{1.703378in}}%
\pgfpathlineto{\pgfqpoint{3.052800in}{1.697836in}}%
\pgfpathlineto{\pgfqpoint{3.045823in}{1.697647in}}%
\pgfpathlineto{\pgfqpoint{3.044375in}{1.690587in}}%
\pgfpathlineto{\pgfqpoint{3.036932in}{1.689276in}}%
\pgfpathlineto{\pgfqpoint{3.034646in}{1.679524in}}%
\pgfpathlineto{\pgfqpoint{3.031696in}{1.678779in}}%
\pgfpathlineto{\pgfqpoint{3.025401in}{1.689136in}}%
\pgfpathlineto{\pgfqpoint{3.014639in}{1.710091in}}%
\pgfpathlineto{\pgfqpoint{3.019840in}{1.712826in}}%
\pgfpathlineto{\pgfqpoint{3.010773in}{1.730803in}}%
\pgfpathlineto{\pgfqpoint{3.015744in}{1.733166in}}%
\pgfpathlineto{\pgfqpoint{3.003662in}{1.755298in}}%
\pgfpathlineto{\pgfqpoint{3.013596in}{1.761178in}}%
\pgfpathclose%
\pgfusepath{fill}%
\end{pgfscope}%
\begin{pgfscope}%
\pgfpathrectangle{\pgfqpoint{0.100000in}{0.100000in}}{\pgfqpoint{3.007045in}{1.925000in}}%
\pgfusepath{clip}%
\pgfsetbuttcap%
\pgfsetmiterjoin%
\definecolor{currentfill}{rgb}{0.417086,0.680631,0.838231}%
\pgfsetfillcolor{currentfill}%
\pgfsetlinewidth{0.000000pt}%
\definecolor{currentstroke}{rgb}{0.000000,0.000000,0.000000}%
\pgfsetstrokecolor{currentstroke}%
\pgfsetstrokeopacity{0.000000}%
\pgfsetdash{}{0pt}%
\pgfpathmoveto{\pgfqpoint{2.378309in}{1.338718in}}%
\pgfpathlineto{\pgfqpoint{2.365754in}{1.334376in}}%
\pgfpathlineto{\pgfqpoint{2.346907in}{1.331869in}}%
\pgfpathlineto{\pgfqpoint{2.342797in}{1.334180in}}%
\pgfpathlineto{\pgfqpoint{2.345273in}{1.312326in}}%
\pgfpathlineto{\pgfqpoint{2.322615in}{1.309553in}}%
\pgfpathlineto{\pgfqpoint{2.305531in}{1.307364in}}%
\pgfpathlineto{\pgfqpoint{2.303481in}{1.324508in}}%
\pgfpathlineto{\pgfqpoint{2.297909in}{1.323832in}}%
\pgfpathlineto{\pgfqpoint{2.295342in}{1.329269in}}%
\pgfpathlineto{\pgfqpoint{2.292953in}{1.341716in}}%
\pgfpathlineto{\pgfqpoint{2.323909in}{1.346635in}}%
\pgfpathlineto{\pgfqpoint{2.320460in}{1.370056in}}%
\pgfpathlineto{\pgfqpoint{2.343612in}{1.373779in}}%
\pgfpathlineto{\pgfqpoint{2.349537in}{1.370183in}}%
\pgfpathlineto{\pgfqpoint{2.343956in}{1.363335in}}%
\pgfpathlineto{\pgfqpoint{2.338987in}{1.354529in}}%
\pgfpathlineto{\pgfqpoint{2.338422in}{1.346684in}}%
\pgfpathlineto{\pgfqpoint{2.344381in}{1.346977in}}%
\pgfpathlineto{\pgfqpoint{2.358766in}{1.342811in}}%
\pgfpathlineto{\pgfqpoint{2.364318in}{1.338530in}}%
\pgfpathlineto{\pgfqpoint{2.375956in}{1.341494in}}%
\pgfpathlineto{\pgfqpoint{2.378309in}{1.338718in}}%
\pgfpathclose%
\pgfusepath{fill}%
\end{pgfscope}%
\begin{pgfscope}%
\pgfpathrectangle{\pgfqpoint{0.100000in}{0.100000in}}{\pgfqpoint{3.007045in}{1.925000in}}%
\pgfusepath{clip}%
\pgfsetbuttcap%
\pgfsetmiterjoin%
\definecolor{currentfill}{rgb}{0.696501,0.824837,0.909266}%
\pgfsetfillcolor{currentfill}%
\pgfsetlinewidth{0.000000pt}%
\definecolor{currentstroke}{rgb}{0.000000,0.000000,0.000000}%
\pgfsetstrokecolor{currentstroke}%
\pgfsetstrokeopacity{0.000000}%
\pgfsetdash{}{0pt}%
\pgfpathmoveto{\pgfqpoint{2.000756in}{1.600279in}}%
\pgfpathlineto{\pgfqpoint{2.002328in}{1.571903in}}%
\pgfpathlineto{\pgfqpoint{2.003003in}{1.560543in}}%
\pgfpathlineto{\pgfqpoint{1.973425in}{1.558760in}}%
\pgfpathlineto{\pgfqpoint{1.971604in}{1.598763in}}%
\pgfpathlineto{\pgfqpoint{1.960172in}{1.598201in}}%
\pgfpathlineto{\pgfqpoint{1.959683in}{1.609718in}}%
\pgfpathlineto{\pgfqpoint{1.930881in}{1.608680in}}%
\pgfpathlineto{\pgfqpoint{1.929386in}{1.648471in}}%
\pgfpathlineto{\pgfqpoint{1.937520in}{1.651178in}}%
\pgfpathlineto{\pgfqpoint{1.944689in}{1.656422in}}%
\pgfpathlineto{\pgfqpoint{1.949266in}{1.655981in}}%
\pgfpathlineto{\pgfqpoint{1.955129in}{1.661825in}}%
\pgfpathlineto{\pgfqpoint{1.960258in}{1.663453in}}%
\pgfpathlineto{\pgfqpoint{1.965406in}{1.659054in}}%
\pgfpathlineto{\pgfqpoint{1.959719in}{1.649522in}}%
\pgfpathlineto{\pgfqpoint{1.959203in}{1.638150in}}%
\pgfpathlineto{\pgfqpoint{1.968638in}{1.644005in}}%
\pgfpathlineto{\pgfqpoint{1.975507in}{1.638992in}}%
\pgfpathlineto{\pgfqpoint{1.976511in}{1.616232in}}%
\pgfpathlineto{\pgfqpoint{1.982301in}{1.616494in}}%
\pgfpathlineto{\pgfqpoint{1.982576in}{1.610763in}}%
\pgfpathlineto{\pgfqpoint{1.988231in}{1.611105in}}%
\pgfpathlineto{\pgfqpoint{1.988873in}{1.599607in}}%
\pgfpathlineto{\pgfqpoint{2.000756in}{1.600279in}}%
\pgfpathclose%
\pgfusepath{fill}%
\end{pgfscope}%
\begin{pgfscope}%
\pgfpathrectangle{\pgfqpoint{0.100000in}{0.100000in}}{\pgfqpoint{3.007045in}{1.925000in}}%
\pgfusepath{clip}%
\pgfsetbuttcap%
\pgfsetmiterjoin%
\definecolor{currentfill}{rgb}{0.447843,0.697855,0.845306}%
\pgfsetfillcolor{currentfill}%
\pgfsetlinewidth{0.000000pt}%
\definecolor{currentstroke}{rgb}{0.000000,0.000000,0.000000}%
\pgfsetstrokecolor{currentstroke}%
\pgfsetstrokeopacity{0.000000}%
\pgfsetdash{}{0pt}%
\pgfpathmoveto{\pgfqpoint{2.212181in}{1.234336in}}%
\pgfpathlineto{\pgfqpoint{2.212424in}{1.232038in}}%
\pgfpathlineto{\pgfqpoint{2.189712in}{1.229760in}}%
\pgfpathlineto{\pgfqpoint{2.189499in}{1.232103in}}%
\pgfpathlineto{\pgfqpoint{2.169632in}{1.230355in}}%
\pgfpathlineto{\pgfqpoint{2.171191in}{1.213179in}}%
\pgfpathlineto{\pgfqpoint{2.154631in}{1.211728in}}%
\pgfpathlineto{\pgfqpoint{2.154351in}{1.223276in}}%
\pgfpathlineto{\pgfqpoint{2.147976in}{1.224103in}}%
\pgfpathlineto{\pgfqpoint{2.149319in}{1.206705in}}%
\pgfpathlineto{\pgfqpoint{2.128961in}{1.204856in}}%
\pgfpathlineto{\pgfqpoint{2.126956in}{1.227545in}}%
\pgfpathlineto{\pgfqpoint{2.125964in}{1.244692in}}%
\pgfpathlineto{\pgfqpoint{2.146371in}{1.246681in}}%
\pgfpathlineto{\pgfqpoint{2.145056in}{1.262796in}}%
\pgfpathlineto{\pgfqpoint{2.166328in}{1.264572in}}%
\pgfpathlineto{\pgfqpoint{2.165690in}{1.271203in}}%
\pgfpathlineto{\pgfqpoint{2.171318in}{1.271716in}}%
\pgfpathlineto{\pgfqpoint{2.173650in}{1.276857in}}%
\pgfpathlineto{\pgfqpoint{2.190977in}{1.278330in}}%
\pgfpathlineto{\pgfqpoint{2.192034in}{1.266835in}}%
\pgfpathlineto{\pgfqpoint{2.195116in}{1.264251in}}%
\pgfpathlineto{\pgfqpoint{2.202654in}{1.264977in}}%
\pgfpathlineto{\pgfqpoint{2.204170in}{1.247864in}}%
\pgfpathlineto{\pgfqpoint{2.211166in}{1.244697in}}%
\pgfpathlineto{\pgfqpoint{2.212181in}{1.234336in}}%
\pgfpathclose%
\pgfusepath{fill}%
\end{pgfscope}%
\begin{pgfscope}%
\pgfpathrectangle{\pgfqpoint{0.100000in}{0.100000in}}{\pgfqpoint{3.007045in}{1.925000in}}%
\pgfusepath{clip}%
\pgfsetbuttcap%
\pgfsetmiterjoin%
\definecolor{currentfill}{rgb}{0.296025,0.597955,0.790988}%
\pgfsetfillcolor{currentfill}%
\pgfsetlinewidth{0.000000pt}%
\definecolor{currentstroke}{rgb}{0.000000,0.000000,0.000000}%
\pgfsetstrokecolor{currentstroke}%
\pgfsetstrokeopacity{0.000000}%
\pgfsetdash{}{0pt}%
\pgfpathmoveto{\pgfqpoint{1.851931in}{0.616696in}}%
\pgfpathlineto{\pgfqpoint{1.834269in}{0.615251in}}%
\pgfpathlineto{\pgfqpoint{1.821877in}{0.611209in}}%
\pgfpathlineto{\pgfqpoint{1.810103in}{0.619771in}}%
\pgfpathlineto{\pgfqpoint{1.804231in}{0.625251in}}%
\pgfpathlineto{\pgfqpoint{1.800480in}{0.631756in}}%
\pgfpathlineto{\pgfqpoint{1.787706in}{0.634802in}}%
\pgfpathlineto{\pgfqpoint{1.779925in}{0.639399in}}%
\pgfpathlineto{\pgfqpoint{1.776581in}{0.643123in}}%
\pgfpathlineto{\pgfqpoint{1.773764in}{0.656330in}}%
\pgfpathlineto{\pgfqpoint{1.775875in}{0.660550in}}%
\pgfpathlineto{\pgfqpoint{1.803211in}{0.660558in}}%
\pgfpathlineto{\pgfqpoint{1.799830in}{0.669148in}}%
\pgfpathlineto{\pgfqpoint{1.827733in}{0.669764in}}%
\pgfpathlineto{\pgfqpoint{1.835531in}{0.660856in}}%
\pgfpathlineto{\pgfqpoint{1.838775in}{0.656201in}}%
\pgfpathlineto{\pgfqpoint{1.839877in}{0.650274in}}%
\pgfpathlineto{\pgfqpoint{1.837976in}{0.644603in}}%
\pgfpathlineto{\pgfqpoint{1.843180in}{0.636011in}}%
\pgfpathlineto{\pgfqpoint{1.847519in}{0.631354in}}%
\pgfpathlineto{\pgfqpoint{1.846842in}{0.625571in}}%
\pgfpathlineto{\pgfqpoint{1.851931in}{0.616696in}}%
\pgfpathclose%
\pgfusepath{fill}%
\end{pgfscope}%
\begin{pgfscope}%
\pgfpathrectangle{\pgfqpoint{0.100000in}{0.100000in}}{\pgfqpoint{3.007045in}{1.925000in}}%
\pgfusepath{clip}%
\pgfsetbuttcap%
\pgfsetmiterjoin%
\definecolor{currentfill}{rgb}{0.637447,0.799739,0.888597}%
\pgfsetfillcolor{currentfill}%
\pgfsetlinewidth{0.000000pt}%
\definecolor{currentstroke}{rgb}{0.000000,0.000000,0.000000}%
\pgfsetstrokecolor{currentstroke}%
\pgfsetstrokeopacity{0.000000}%
\pgfsetdash{}{0pt}%
\pgfpathmoveto{\pgfqpoint{2.639205in}{0.819156in}}%
\pgfpathlineto{\pgfqpoint{2.628686in}{0.826071in}}%
\pgfpathlineto{\pgfqpoint{2.612678in}{0.829037in}}%
\pgfpathlineto{\pgfqpoint{2.597720in}{0.838651in}}%
\pgfpathlineto{\pgfqpoint{2.590755in}{0.837224in}}%
\pgfpathlineto{\pgfqpoint{2.595239in}{0.853039in}}%
\pgfpathlineto{\pgfqpoint{2.593582in}{0.855868in}}%
\pgfpathlineto{\pgfqpoint{2.597384in}{0.859907in}}%
\pgfpathlineto{\pgfqpoint{2.592778in}{0.867711in}}%
\pgfpathlineto{\pgfqpoint{2.596705in}{0.871224in}}%
\pgfpathlineto{\pgfqpoint{2.580510in}{0.880092in}}%
\pgfpathlineto{\pgfqpoint{2.582316in}{0.884878in}}%
\pgfpathlineto{\pgfqpoint{2.577854in}{0.890686in}}%
\pgfpathlineto{\pgfqpoint{2.571314in}{0.892287in}}%
\pgfpathlineto{\pgfqpoint{2.582994in}{0.902658in}}%
\pgfpathlineto{\pgfqpoint{2.595973in}{0.904310in}}%
\pgfpathlineto{\pgfqpoint{2.607323in}{0.893730in}}%
\pgfpathlineto{\pgfqpoint{2.613555in}{0.916458in}}%
\pgfpathlineto{\pgfqpoint{2.637551in}{0.899069in}}%
\pgfpathlineto{\pgfqpoint{2.671279in}{0.875019in}}%
\pgfpathlineto{\pgfqpoint{2.663098in}{0.870424in}}%
\pgfpathlineto{\pgfqpoint{2.657985in}{0.865431in}}%
\pgfpathlineto{\pgfqpoint{2.651920in}{0.856646in}}%
\pgfpathlineto{\pgfqpoint{2.646223in}{0.845821in}}%
\pgfpathlineto{\pgfqpoint{2.644013in}{0.838187in}}%
\pgfpathlineto{\pgfqpoint{2.644164in}{0.824590in}}%
\pgfpathlineto{\pgfqpoint{2.639205in}{0.819156in}}%
\pgfpathclose%
\pgfusepath{fill}%
\end{pgfscope}%
\begin{pgfscope}%
\pgfpathrectangle{\pgfqpoint{0.100000in}{0.100000in}}{\pgfqpoint{3.007045in}{1.925000in}}%
\pgfusepath{clip}%
\pgfsetbuttcap%
\pgfsetmiterjoin%
\definecolor{currentfill}{rgb}{0.770319,0.856209,0.935102}%
\pgfsetfillcolor{currentfill}%
\pgfsetlinewidth{0.000000pt}%
\definecolor{currentstroke}{rgb}{0.000000,0.000000,0.000000}%
\pgfsetstrokecolor{currentstroke}%
\pgfsetstrokeopacity{0.000000}%
\pgfsetdash{}{0pt}%
\pgfpathmoveto{\pgfqpoint{2.192935in}{1.614287in}}%
\pgfpathlineto{\pgfqpoint{2.188229in}{1.662007in}}%
\pgfpathlineto{\pgfqpoint{2.204739in}{1.662788in}}%
\pgfpathlineto{\pgfqpoint{2.213687in}{1.668241in}}%
\pgfpathlineto{\pgfqpoint{2.228170in}{1.671708in}}%
\pgfpathlineto{\pgfqpoint{2.226016in}{1.665581in}}%
\pgfpathlineto{\pgfqpoint{2.227291in}{1.652489in}}%
\pgfpathlineto{\pgfqpoint{2.236977in}{1.650715in}}%
\pgfpathlineto{\pgfqpoint{2.243546in}{1.654372in}}%
\pgfpathlineto{\pgfqpoint{2.247528in}{1.650275in}}%
\pgfpathlineto{\pgfqpoint{2.255094in}{1.656649in}}%
\pgfpathlineto{\pgfqpoint{2.269105in}{1.658515in}}%
\pgfpathlineto{\pgfqpoint{2.268196in}{1.652678in}}%
\pgfpathlineto{\pgfqpoint{2.270778in}{1.646529in}}%
\pgfpathlineto{\pgfqpoint{2.269357in}{1.637766in}}%
\pgfpathlineto{\pgfqpoint{2.275952in}{1.634432in}}%
\pgfpathlineto{\pgfqpoint{2.278409in}{1.628086in}}%
\pgfpathlineto{\pgfqpoint{2.287170in}{1.626728in}}%
\pgfpathlineto{\pgfqpoint{2.296641in}{1.634077in}}%
\pgfpathlineto{\pgfqpoint{2.302457in}{1.627925in}}%
\pgfpathlineto{\pgfqpoint{2.301033in}{1.623548in}}%
\pgfpathlineto{\pgfqpoint{2.294668in}{1.624039in}}%
\pgfpathlineto{\pgfqpoint{2.288149in}{1.622604in}}%
\pgfpathlineto{\pgfqpoint{2.283389in}{1.624709in}}%
\pgfpathlineto{\pgfqpoint{2.272380in}{1.623457in}}%
\pgfpathlineto{\pgfqpoint{2.261434in}{1.619276in}}%
\pgfpathlineto{\pgfqpoint{2.254060in}{1.620885in}}%
\pgfpathlineto{\pgfqpoint{2.251911in}{1.624956in}}%
\pgfpathlineto{\pgfqpoint{2.246315in}{1.624264in}}%
\pgfpathlineto{\pgfqpoint{2.244624in}{1.611374in}}%
\pgfpathlineto{\pgfqpoint{2.236991in}{1.615863in}}%
\pgfpathlineto{\pgfqpoint{2.232484in}{1.620710in}}%
\pgfpathlineto{\pgfqpoint{2.225249in}{1.622997in}}%
\pgfpathlineto{\pgfqpoint{2.213158in}{1.624960in}}%
\pgfpathlineto{\pgfqpoint{2.207852in}{1.624039in}}%
\pgfpathlineto{\pgfqpoint{2.201875in}{1.614996in}}%
\pgfpathlineto{\pgfqpoint{2.192935in}{1.614287in}}%
\pgfpathclose%
\pgfusepath{fill}%
\end{pgfscope}%
\begin{pgfscope}%
\pgfpathrectangle{\pgfqpoint{0.100000in}{0.100000in}}{\pgfqpoint{3.007045in}{1.925000in}}%
\pgfusepath{clip}%
\pgfsetbuttcap%
\pgfsetmiterjoin%
\definecolor{currentfill}{rgb}{0.391865,0.663406,0.828389}%
\pgfsetfillcolor{currentfill}%
\pgfsetlinewidth{0.000000pt}%
\definecolor{currentstroke}{rgb}{0.000000,0.000000,0.000000}%
\pgfsetstrokecolor{currentstroke}%
\pgfsetstrokeopacity{0.000000}%
\pgfsetdash{}{0pt}%
\pgfpathmoveto{\pgfqpoint{2.085014in}{0.905571in}}%
\pgfpathlineto{\pgfqpoint{2.069407in}{0.904987in}}%
\pgfpathlineto{\pgfqpoint{2.068214in}{0.916226in}}%
\pgfpathlineto{\pgfqpoint{2.061506in}{0.919581in}}%
\pgfpathlineto{\pgfqpoint{2.051337in}{0.919577in}}%
\pgfpathlineto{\pgfqpoint{2.045514in}{0.912506in}}%
\pgfpathlineto{\pgfqpoint{2.043371in}{0.916596in}}%
\pgfpathlineto{\pgfqpoint{2.042608in}{0.925600in}}%
\pgfpathlineto{\pgfqpoint{2.049440in}{0.927189in}}%
\pgfpathlineto{\pgfqpoint{2.052220in}{0.930837in}}%
\pgfpathlineto{\pgfqpoint{2.051897in}{0.936223in}}%
\pgfpathlineto{\pgfqpoint{2.054350in}{0.941105in}}%
\pgfpathlineto{\pgfqpoint{2.055729in}{0.949079in}}%
\pgfpathlineto{\pgfqpoint{2.060132in}{0.952292in}}%
\pgfpathlineto{\pgfqpoint{2.054510in}{0.958420in}}%
\pgfpathlineto{\pgfqpoint{2.062788in}{0.960749in}}%
\pgfpathlineto{\pgfqpoint{2.061897in}{0.966247in}}%
\pgfpathlineto{\pgfqpoint{2.061582in}{0.976889in}}%
\pgfpathlineto{\pgfqpoint{2.064440in}{0.977028in}}%
\pgfpathlineto{\pgfqpoint{2.068022in}{0.977336in}}%
\pgfpathlineto{\pgfqpoint{2.071414in}{0.971616in}}%
\pgfpathlineto{\pgfqpoint{2.068381in}{0.961543in}}%
\pgfpathlineto{\pgfqpoint{2.065703in}{0.958280in}}%
\pgfpathlineto{\pgfqpoint{2.083172in}{0.958841in}}%
\pgfpathlineto{\pgfqpoint{2.083454in}{0.947053in}}%
\pgfpathlineto{\pgfqpoint{2.074636in}{0.931335in}}%
\pgfpathlineto{\pgfqpoint{2.089884in}{0.925436in}}%
\pgfpathlineto{\pgfqpoint{2.090507in}{0.908195in}}%
\pgfpathlineto{\pgfqpoint{2.085014in}{0.905571in}}%
\pgfpathclose%
\pgfusepath{fill}%
\end{pgfscope}%
\begin{pgfscope}%
\pgfpathrectangle{\pgfqpoint{0.100000in}{0.100000in}}{\pgfqpoint{3.007045in}{1.925000in}}%
\pgfusepath{clip}%
\pgfsetbuttcap%
\pgfsetmiterjoin%
\definecolor{currentfill}{rgb}{0.454118,0.701300,0.846659}%
\pgfsetfillcolor{currentfill}%
\pgfsetlinewidth{0.000000pt}%
\definecolor{currentstroke}{rgb}{0.000000,0.000000,0.000000}%
\pgfsetstrokecolor{currentstroke}%
\pgfsetstrokeopacity{0.000000}%
\pgfsetdash{}{0pt}%
\pgfpathmoveto{\pgfqpoint{2.492448in}{0.960338in}}%
\pgfpathlineto{\pgfqpoint{2.498902in}{0.973883in}}%
\pgfpathlineto{\pgfqpoint{2.494999in}{0.975158in}}%
\pgfpathlineto{\pgfqpoint{2.489274in}{0.973776in}}%
\pgfpathlineto{\pgfqpoint{2.480170in}{0.978949in}}%
\pgfpathlineto{\pgfqpoint{2.474213in}{0.983871in}}%
\pgfpathlineto{\pgfqpoint{2.479657in}{0.995365in}}%
\pgfpathlineto{\pgfqpoint{2.486521in}{0.996045in}}%
\pgfpathlineto{\pgfqpoint{2.494094in}{0.994158in}}%
\pgfpathlineto{\pgfqpoint{2.499152in}{0.989884in}}%
\pgfpathlineto{\pgfqpoint{2.506001in}{0.993623in}}%
\pgfpathlineto{\pgfqpoint{2.510978in}{0.993811in}}%
\pgfpathlineto{\pgfqpoint{2.522081in}{0.997367in}}%
\pgfpathlineto{\pgfqpoint{2.532109in}{0.998638in}}%
\pgfpathlineto{\pgfqpoint{2.533363in}{0.985590in}}%
\pgfpathlineto{\pgfqpoint{2.531826in}{0.973469in}}%
\pgfpathlineto{\pgfqpoint{2.535380in}{0.962820in}}%
\pgfpathlineto{\pgfqpoint{2.523665in}{0.961730in}}%
\pgfpathlineto{\pgfqpoint{2.523136in}{0.963616in}}%
\pgfpathlineto{\pgfqpoint{2.492448in}{0.960338in}}%
\pgfpathclose%
\pgfusepath{fill}%
\end{pgfscope}%
\begin{pgfscope}%
\pgfpathrectangle{\pgfqpoint{0.100000in}{0.100000in}}{\pgfqpoint{3.007045in}{1.925000in}}%
\pgfusepath{clip}%
\pgfsetbuttcap%
\pgfsetmiterjoin%
\definecolor{currentfill}{rgb}{0.460392,0.704744,0.848012}%
\pgfsetfillcolor{currentfill}%
\pgfsetlinewidth{0.000000pt}%
\definecolor{currentstroke}{rgb}{0.000000,0.000000,0.000000}%
\pgfsetstrokecolor{currentstroke}%
\pgfsetstrokeopacity{0.000000}%
\pgfsetdash{}{0pt}%
\pgfpathmoveto{\pgfqpoint{2.350943in}{1.269110in}}%
\pgfpathlineto{\pgfqpoint{2.345898in}{1.268387in}}%
\pgfpathlineto{\pgfqpoint{2.327732in}{1.268420in}}%
\pgfpathlineto{\pgfqpoint{2.322615in}{1.309553in}}%
\pgfpathlineto{\pgfqpoint{2.345273in}{1.312326in}}%
\pgfpathlineto{\pgfqpoint{2.342797in}{1.334180in}}%
\pgfpathlineto{\pgfqpoint{2.346907in}{1.331869in}}%
\pgfpathlineto{\pgfqpoint{2.365754in}{1.334376in}}%
\pgfpathlineto{\pgfqpoint{2.371168in}{1.333253in}}%
\pgfpathlineto{\pgfqpoint{2.375866in}{1.305055in}}%
\pgfpathlineto{\pgfqpoint{2.361935in}{1.302986in}}%
\pgfpathlineto{\pgfqpoint{2.364496in}{1.284008in}}%
\pgfpathlineto{\pgfqpoint{2.349350in}{1.280970in}}%
\pgfpathlineto{\pgfqpoint{2.350943in}{1.269110in}}%
\pgfpathclose%
\pgfusepath{fill}%
\end{pgfscope}%
\begin{pgfscope}%
\pgfpathrectangle{\pgfqpoint{0.100000in}{0.100000in}}{\pgfqpoint{3.007045in}{1.925000in}}%
\pgfusepath{clip}%
\pgfsetbuttcap%
\pgfsetmiterjoin%
\definecolor{currentfill}{rgb}{0.666974,0.812288,0.898931}%
\pgfsetfillcolor{currentfill}%
\pgfsetlinewidth{0.000000pt}%
\definecolor{currentstroke}{rgb}{0.000000,0.000000,0.000000}%
\pgfsetstrokecolor{currentstroke}%
\pgfsetstrokeopacity{0.000000}%
\pgfsetdash{}{0pt}%
\pgfpathmoveto{\pgfqpoint{2.880647in}{1.472875in}}%
\pgfpathlineto{\pgfqpoint{2.836904in}{1.462520in}}%
\pgfpathlineto{\pgfqpoint{2.833508in}{1.484777in}}%
\pgfpathlineto{\pgfqpoint{2.832224in}{1.502258in}}%
\pgfpathlineto{\pgfqpoint{2.827848in}{1.505323in}}%
\pgfpathlineto{\pgfqpoint{2.831656in}{1.508759in}}%
\pgfpathlineto{\pgfqpoint{2.862374in}{1.515277in}}%
\pgfpathlineto{\pgfqpoint{2.865846in}{1.512168in}}%
\pgfpathlineto{\pgfqpoint{2.865008in}{1.506125in}}%
\pgfpathlineto{\pgfqpoint{2.867051in}{1.490726in}}%
\pgfpathlineto{\pgfqpoint{2.873175in}{1.485821in}}%
\pgfpathlineto{\pgfqpoint{2.872055in}{1.481134in}}%
\pgfpathlineto{\pgfqpoint{2.878530in}{1.481285in}}%
\pgfpathlineto{\pgfqpoint{2.880647in}{1.472875in}}%
\pgfpathclose%
\pgfusepath{fill}%
\end{pgfscope}%
\begin{pgfscope}%
\pgfpathrectangle{\pgfqpoint{0.100000in}{0.100000in}}{\pgfqpoint{3.007045in}{1.925000in}}%
\pgfusepath{clip}%
\pgfsetbuttcap%
\pgfsetmiterjoin%
\definecolor{currentfill}{rgb}{0.316201,0.611734,0.798862}%
\pgfsetfillcolor{currentfill}%
\pgfsetlinewidth{0.000000pt}%
\definecolor{currentstroke}{rgb}{0.000000,0.000000,0.000000}%
\pgfsetstrokecolor{currentstroke}%
\pgfsetstrokeopacity{0.000000}%
\pgfsetdash{}{0pt}%
\pgfpathmoveto{\pgfqpoint{1.631910in}{1.063521in}}%
\pgfpathlineto{\pgfqpoint{1.632419in}{1.080689in}}%
\pgfpathlineto{\pgfqpoint{1.621015in}{1.081045in}}%
\pgfpathlineto{\pgfqpoint{1.621825in}{1.109765in}}%
\pgfpathlineto{\pgfqpoint{1.650255in}{1.108902in}}%
\pgfpathlineto{\pgfqpoint{1.672984in}{1.108364in}}%
\pgfpathlineto{\pgfqpoint{1.672876in}{1.102642in}}%
\pgfpathlineto{\pgfqpoint{1.678559in}{1.102524in}}%
\pgfpathlineto{\pgfqpoint{1.677388in}{1.091095in}}%
\pgfpathlineto{\pgfqpoint{1.676880in}{1.073769in}}%
\pgfpathlineto{\pgfqpoint{1.660705in}{1.074262in}}%
\pgfpathlineto{\pgfqpoint{1.660467in}{1.062737in}}%
\pgfpathlineto{\pgfqpoint{1.631910in}{1.063521in}}%
\pgfpathclose%
\pgfusepath{fill}%
\end{pgfscope}%
\begin{pgfscope}%
\pgfpathrectangle{\pgfqpoint{0.100000in}{0.100000in}}{\pgfqpoint{3.007045in}{1.925000in}}%
\pgfusepath{clip}%
\pgfsetbuttcap%
\pgfsetmiterjoin%
\definecolor{currentfill}{rgb}{0.114802,0.424437,0.695194}%
\pgfsetfillcolor{currentfill}%
\pgfsetlinewidth{0.000000pt}%
\definecolor{currentstroke}{rgb}{0.000000,0.000000,0.000000}%
\pgfsetstrokecolor{currentstroke}%
\pgfsetstrokeopacity{0.000000}%
\pgfsetdash{}{0pt}%
\pgfpathmoveto{\pgfqpoint{1.449237in}{1.655533in}}%
\pgfpathlineto{\pgfqpoint{1.448770in}{1.649792in}}%
\pgfpathlineto{\pgfqpoint{1.410950in}{1.652859in}}%
\pgfpathlineto{\pgfqpoint{1.379987in}{1.655666in}}%
\pgfpathlineto{\pgfqpoint{1.382174in}{1.678842in}}%
\pgfpathlineto{\pgfqpoint{1.379578in}{1.679082in}}%
\pgfpathlineto{\pgfqpoint{1.381806in}{1.702240in}}%
\pgfpathlineto{\pgfqpoint{1.407305in}{1.699840in}}%
\pgfpathlineto{\pgfqpoint{1.409324in}{1.722778in}}%
\pgfpathlineto{\pgfqpoint{1.429757in}{1.720971in}}%
\pgfpathlineto{\pgfqpoint{1.433959in}{1.727722in}}%
\pgfpathlineto{\pgfqpoint{1.441706in}{1.725605in}}%
\pgfpathlineto{\pgfqpoint{1.447405in}{1.726710in}}%
\pgfpathlineto{\pgfqpoint{1.446373in}{1.716181in}}%
\pgfpathlineto{\pgfqpoint{1.448796in}{1.712766in}}%
\pgfpathlineto{\pgfqpoint{1.447485in}{1.696400in}}%
\pgfpathlineto{\pgfqpoint{1.450238in}{1.696191in}}%
\pgfpathlineto{\pgfqpoint{1.448415in}{1.673164in}}%
\pgfpathlineto{\pgfqpoint{1.450582in}{1.673001in}}%
\pgfpathlineto{\pgfqpoint{1.449237in}{1.655533in}}%
\pgfpathclose%
\pgfusepath{fill}%
\end{pgfscope}%
\begin{pgfscope}%
\pgfpathrectangle{\pgfqpoint{0.100000in}{0.100000in}}{\pgfqpoint{3.007045in}{1.925000in}}%
\pgfusepath{clip}%
\pgfsetbuttcap%
\pgfsetmiterjoin%
\definecolor{currentfill}{rgb}{0.412042,0.677186,0.836263}%
\pgfsetfillcolor{currentfill}%
\pgfsetlinewidth{0.000000pt}%
\definecolor{currentstroke}{rgb}{0.000000,0.000000,0.000000}%
\pgfsetstrokecolor{currentstroke}%
\pgfsetstrokeopacity{0.000000}%
\pgfsetdash{}{0pt}%
\pgfpathmoveto{\pgfqpoint{0.842950in}{0.338119in}}%
\pgfpathlineto{\pgfqpoint{0.844097in}{0.340724in}}%
\pgfpathlineto{\pgfqpoint{0.845519in}{0.341453in}}%
\pgfpathlineto{\pgfqpoint{0.848054in}{0.340022in}}%
\pgfpathlineto{\pgfqpoint{0.847662in}{0.338845in}}%
\pgfpathlineto{\pgfqpoint{0.850122in}{0.336629in}}%
\pgfpathlineto{\pgfqpoint{0.844580in}{0.336908in}}%
\pgfpathlineto{\pgfqpoint{0.842950in}{0.338119in}}%
\pgfpathclose%
\pgfusepath{fill}%
\end{pgfscope}%
\begin{pgfscope}%
\pgfpathrectangle{\pgfqpoint{0.100000in}{0.100000in}}{\pgfqpoint{3.007045in}{1.925000in}}%
\pgfusepath{clip}%
\pgfsetbuttcap%
\pgfsetmiterjoin%
\definecolor{currentfill}{rgb}{0.412042,0.677186,0.836263}%
\pgfsetfillcolor{currentfill}%
\pgfsetlinewidth{0.000000pt}%
\definecolor{currentstroke}{rgb}{0.000000,0.000000,0.000000}%
\pgfsetstrokecolor{currentstroke}%
\pgfsetstrokeopacity{0.000000}%
\pgfsetdash{}{0pt}%
\pgfpathmoveto{\pgfqpoint{0.824088in}{0.336340in}}%
\pgfpathlineto{\pgfqpoint{0.826449in}{0.337990in}}%
\pgfpathlineto{\pgfqpoint{0.828677in}{0.338112in}}%
\pgfpathlineto{\pgfqpoint{0.837136in}{0.340612in}}%
\pgfpathlineto{\pgfqpoint{0.837987in}{0.342496in}}%
\pgfpathlineto{\pgfqpoint{0.839837in}{0.341361in}}%
\pgfpathlineto{\pgfqpoint{0.839821in}{0.339455in}}%
\pgfpathlineto{\pgfqpoint{0.834944in}{0.338315in}}%
\pgfpathlineto{\pgfqpoint{0.831383in}{0.335875in}}%
\pgfpathlineto{\pgfqpoint{0.829597in}{0.336005in}}%
\pgfpathlineto{\pgfqpoint{0.828843in}{0.334684in}}%
\pgfpathlineto{\pgfqpoint{0.825115in}{0.335470in}}%
\pgfpathlineto{\pgfqpoint{0.824088in}{0.336340in}}%
\pgfpathclose%
\pgfusepath{fill}%
\end{pgfscope}%
\begin{pgfscope}%
\pgfpathrectangle{\pgfqpoint{0.100000in}{0.100000in}}{\pgfqpoint{3.007045in}{1.925000in}}%
\pgfusepath{clip}%
\pgfsetbuttcap%
\pgfsetmiterjoin%
\definecolor{currentfill}{rgb}{0.412042,0.677186,0.836263}%
\pgfsetfillcolor{currentfill}%
\pgfsetlinewidth{0.000000pt}%
\definecolor{currentstroke}{rgb}{0.000000,0.000000,0.000000}%
\pgfsetstrokecolor{currentstroke}%
\pgfsetstrokeopacity{0.000000}%
\pgfsetdash{}{0pt}%
\pgfpathmoveto{\pgfqpoint{0.829895in}{0.342488in}}%
\pgfpathlineto{\pgfqpoint{0.832685in}{0.348119in}}%
\pgfpathlineto{\pgfqpoint{0.833319in}{0.346252in}}%
\pgfpathlineto{\pgfqpoint{0.829895in}{0.342488in}}%
\pgfpathclose%
\pgfusepath{fill}%
\end{pgfscope}%
\begin{pgfscope}%
\pgfpathrectangle{\pgfqpoint{0.100000in}{0.100000in}}{\pgfqpoint{3.007045in}{1.925000in}}%
\pgfusepath{clip}%
\pgfsetbuttcap%
\pgfsetmiterjoin%
\definecolor{currentfill}{rgb}{0.412042,0.677186,0.836263}%
\pgfsetfillcolor{currentfill}%
\pgfsetlinewidth{0.000000pt}%
\definecolor{currentstroke}{rgb}{0.000000,0.000000,0.000000}%
\pgfsetstrokecolor{currentstroke}%
\pgfsetstrokeopacity{0.000000}%
\pgfsetdash{}{0pt}%
\pgfpathmoveto{\pgfqpoint{0.880633in}{0.393554in}}%
\pgfpathlineto{\pgfqpoint{0.884245in}{0.391268in}}%
\pgfpathlineto{\pgfqpoint{0.886512in}{0.394965in}}%
\pgfpathlineto{\pgfqpoint{0.895631in}{0.389258in}}%
\pgfpathlineto{\pgfqpoint{0.894498in}{0.387400in}}%
\pgfpathlineto{\pgfqpoint{0.899657in}{0.384205in}}%
\pgfpathlineto{\pgfqpoint{0.914264in}{0.375762in}}%
\pgfpathlineto{\pgfqpoint{0.912562in}{0.370845in}}%
\pgfpathlineto{\pgfqpoint{0.909515in}{0.369615in}}%
\pgfpathlineto{\pgfqpoint{0.910203in}{0.367541in}}%
\pgfpathlineto{\pgfqpoint{0.907829in}{0.364423in}}%
\pgfpathlineto{\pgfqpoint{0.910185in}{0.362649in}}%
\pgfpathlineto{\pgfqpoint{0.912886in}{0.364295in}}%
\pgfpathlineto{\pgfqpoint{0.913444in}{0.363209in}}%
\pgfpathlineto{\pgfqpoint{0.916564in}{0.361368in}}%
\pgfpathlineto{\pgfqpoint{0.915470in}{0.359574in}}%
\pgfpathlineto{\pgfqpoint{0.917368in}{0.358439in}}%
\pgfpathlineto{\pgfqpoint{0.916310in}{0.356673in}}%
\pgfpathlineto{\pgfqpoint{0.917983in}{0.355693in}}%
\pgfpathlineto{\pgfqpoint{0.913691in}{0.348399in}}%
\pgfpathlineto{\pgfqpoint{0.914753in}{0.346694in}}%
\pgfpathlineto{\pgfqpoint{0.912069in}{0.341923in}}%
\pgfpathlineto{\pgfqpoint{0.920301in}{0.337334in}}%
\pgfpathlineto{\pgfqpoint{0.911029in}{0.320607in}}%
\pgfpathlineto{\pgfqpoint{0.901489in}{0.303399in}}%
\pgfpathlineto{\pgfqpoint{0.895195in}{0.292046in}}%
\pgfpathlineto{\pgfqpoint{0.891908in}{0.297356in}}%
\pgfpathlineto{\pgfqpoint{0.895042in}{0.297561in}}%
\pgfpathlineto{\pgfqpoint{0.895865in}{0.299933in}}%
\pgfpathlineto{\pgfqpoint{0.894005in}{0.301147in}}%
\pgfpathlineto{\pgfqpoint{0.893616in}{0.299471in}}%
\pgfpathlineto{\pgfqpoint{0.889618in}{0.300158in}}%
\pgfpathlineto{\pgfqpoint{0.888432in}{0.302470in}}%
\pgfpathlineto{\pgfqpoint{0.880446in}{0.309836in}}%
\pgfpathlineto{\pgfqpoint{0.876880in}{0.310671in}}%
\pgfpathlineto{\pgfqpoint{0.872039in}{0.312853in}}%
\pgfpathlineto{\pgfqpoint{0.867319in}{0.314295in}}%
\pgfpathlineto{\pgfqpoint{0.866304in}{0.315660in}}%
\pgfpathlineto{\pgfqpoint{0.866224in}{0.318437in}}%
\pgfpathlineto{\pgfqpoint{0.865182in}{0.320626in}}%
\pgfpathlineto{\pgfqpoint{0.862537in}{0.322620in}}%
\pgfpathlineto{\pgfqpoint{0.862785in}{0.325009in}}%
\pgfpathlineto{\pgfqpoint{0.861746in}{0.326681in}}%
\pgfpathlineto{\pgfqpoint{0.864218in}{0.329401in}}%
\pgfpathlineto{\pgfqpoint{0.864631in}{0.330890in}}%
\pgfpathlineto{\pgfqpoint{0.863184in}{0.332010in}}%
\pgfpathlineto{\pgfqpoint{0.861073in}{0.330560in}}%
\pgfpathlineto{\pgfqpoint{0.858822in}{0.331645in}}%
\pgfpathlineto{\pgfqpoint{0.857771in}{0.330884in}}%
\pgfpathlineto{\pgfqpoint{0.855435in}{0.335097in}}%
\pgfpathlineto{\pgfqpoint{0.853673in}{0.335957in}}%
\pgfpathlineto{\pgfqpoint{0.855410in}{0.337624in}}%
\pgfpathlineto{\pgfqpoint{0.855863in}{0.341580in}}%
\pgfpathlineto{\pgfqpoint{0.855648in}{0.344279in}}%
\pgfpathlineto{\pgfqpoint{0.852144in}{0.345090in}}%
\pgfpathlineto{\pgfqpoint{0.853807in}{0.346777in}}%
\pgfpathlineto{\pgfqpoint{0.850106in}{0.348448in}}%
\pgfpathlineto{\pgfqpoint{0.850605in}{0.352065in}}%
\pgfpathlineto{\pgfqpoint{0.848839in}{0.353027in}}%
\pgfpathlineto{\pgfqpoint{0.848400in}{0.354633in}}%
\pgfpathlineto{\pgfqpoint{0.846238in}{0.354335in}}%
\pgfpathlineto{\pgfqpoint{0.845449in}{0.352516in}}%
\pgfpathlineto{\pgfqpoint{0.842638in}{0.355946in}}%
\pgfpathlineto{\pgfqpoint{0.843920in}{0.358581in}}%
\pgfpathlineto{\pgfqpoint{0.839850in}{0.357610in}}%
\pgfpathlineto{\pgfqpoint{0.835041in}{0.357359in}}%
\pgfpathlineto{\pgfqpoint{0.836568in}{0.359811in}}%
\pgfpathlineto{\pgfqpoint{0.838051in}{0.358886in}}%
\pgfpathlineto{\pgfqpoint{0.838232in}{0.362073in}}%
\pgfpathlineto{\pgfqpoint{0.836781in}{0.362125in}}%
\pgfpathlineto{\pgfqpoint{0.832732in}{0.359491in}}%
\pgfpathlineto{\pgfqpoint{0.829761in}{0.356685in}}%
\pgfpathlineto{\pgfqpoint{0.830722in}{0.356227in}}%
\pgfpathlineto{\pgfqpoint{0.832645in}{0.357725in}}%
\pgfpathlineto{\pgfqpoint{0.834243in}{0.356277in}}%
\pgfpathlineto{\pgfqpoint{0.832157in}{0.354595in}}%
\pgfpathlineto{\pgfqpoint{0.830411in}{0.354246in}}%
\pgfpathlineto{\pgfqpoint{0.827722in}{0.355289in}}%
\pgfpathlineto{\pgfqpoint{0.828994in}{0.351862in}}%
\pgfpathlineto{\pgfqpoint{0.830872in}{0.350681in}}%
\pgfpathlineto{\pgfqpoint{0.825792in}{0.348726in}}%
\pgfpathlineto{\pgfqpoint{0.825045in}{0.347737in}}%
\pgfpathlineto{\pgfqpoint{0.828029in}{0.344167in}}%
\pgfpathlineto{\pgfqpoint{0.827087in}{0.340105in}}%
\pgfpathlineto{\pgfqpoint{0.825414in}{0.340100in}}%
\pgfpathlineto{\pgfqpoint{0.823966in}{0.342216in}}%
\pgfpathlineto{\pgfqpoint{0.824700in}{0.344679in}}%
\pgfpathlineto{\pgfqpoint{0.823443in}{0.345830in}}%
\pgfpathlineto{\pgfqpoint{0.821225in}{0.343325in}}%
\pgfpathlineto{\pgfqpoint{0.819433in}{0.344425in}}%
\pgfpathlineto{\pgfqpoint{0.821535in}{0.347132in}}%
\pgfpathlineto{\pgfqpoint{0.825526in}{0.353057in}}%
\pgfpathlineto{\pgfqpoint{0.824494in}{0.353761in}}%
\pgfpathlineto{\pgfqpoint{0.828821in}{0.359977in}}%
\pgfpathlineto{\pgfqpoint{0.827922in}{0.360561in}}%
\pgfpathlineto{\pgfqpoint{0.829442in}{0.362834in}}%
\pgfpathlineto{\pgfqpoint{0.832191in}{0.361126in}}%
\pgfpathlineto{\pgfqpoint{0.836391in}{0.367548in}}%
\pgfpathlineto{\pgfqpoint{0.840140in}{0.372751in}}%
\pgfpathlineto{\pgfqpoint{0.852394in}{0.364648in}}%
\pgfpathlineto{\pgfqpoint{0.853061in}{0.365671in}}%
\pgfpathlineto{\pgfqpoint{0.855608in}{0.363995in}}%
\pgfpathlineto{\pgfqpoint{0.864498in}{0.378218in}}%
\pgfpathlineto{\pgfqpoint{0.865448in}{0.380135in}}%
\pgfpathlineto{\pgfqpoint{0.870796in}{0.376637in}}%
\pgfpathlineto{\pgfqpoint{0.878779in}{0.389457in}}%
\pgfpathlineto{\pgfqpoint{0.880633in}{0.393554in}}%
\pgfpathclose%
\pgfusepath{fill}%
\end{pgfscope}%
\begin{pgfscope}%
\pgfpathrectangle{\pgfqpoint{0.100000in}{0.100000in}}{\pgfqpoint{3.007045in}{1.925000in}}%
\pgfusepath{clip}%
\pgfsetbuttcap%
\pgfsetmiterjoin%
\definecolor{currentfill}{rgb}{0.523137,0.739193,0.861546}%
\pgfsetfillcolor{currentfill}%
\pgfsetlinewidth{0.000000pt}%
\definecolor{currentstroke}{rgb}{0.000000,0.000000,0.000000}%
\pgfsetstrokecolor{currentstroke}%
\pgfsetstrokeopacity{0.000000}%
\pgfsetdash{}{0pt}%
\pgfpathmoveto{\pgfqpoint{2.453185in}{1.084073in}}%
\pgfpathlineto{\pgfqpoint{2.456556in}{1.087012in}}%
\pgfpathlineto{\pgfqpoint{2.458232in}{1.094688in}}%
\pgfpathlineto{\pgfqpoint{2.463061in}{1.098120in}}%
\pgfpathlineto{\pgfqpoint{2.459342in}{1.100923in}}%
\pgfpathlineto{\pgfqpoint{2.471855in}{1.105407in}}%
\pgfpathlineto{\pgfqpoint{2.476705in}{1.103396in}}%
\pgfpathlineto{\pgfqpoint{2.484474in}{1.093009in}}%
\pgfpathlineto{\pgfqpoint{2.489686in}{1.089390in}}%
\pgfpathlineto{\pgfqpoint{2.487921in}{1.088350in}}%
\pgfpathlineto{\pgfqpoint{2.483702in}{1.075698in}}%
\pgfpathlineto{\pgfqpoint{2.475445in}{1.067081in}}%
\pgfpathlineto{\pgfqpoint{2.469086in}{1.065710in}}%
\pgfpathlineto{\pgfqpoint{2.462901in}{1.070185in}}%
\pgfpathlineto{\pgfqpoint{2.455254in}{1.073522in}}%
\pgfpathlineto{\pgfqpoint{2.452938in}{1.079250in}}%
\pgfpathlineto{\pgfqpoint{2.453185in}{1.084073in}}%
\pgfpathclose%
\pgfusepath{fill}%
\end{pgfscope}%
\begin{pgfscope}%
\pgfpathrectangle{\pgfqpoint{0.100000in}{0.100000in}}{\pgfqpoint{3.007045in}{1.925000in}}%
\pgfusepath{clip}%
\pgfsetbuttcap%
\pgfsetmiterjoin%
\definecolor{currentfill}{rgb}{0.535686,0.746082,0.864252}%
\pgfsetfillcolor{currentfill}%
\pgfsetlinewidth{0.000000pt}%
\definecolor{currentstroke}{rgb}{0.000000,0.000000,0.000000}%
\pgfsetstrokecolor{currentstroke}%
\pgfsetstrokeopacity{0.000000}%
\pgfsetdash{}{0pt}%
\pgfpathmoveto{\pgfqpoint{2.533154in}{0.479245in}}%
\pgfpathlineto{\pgfqpoint{2.531241in}{0.493155in}}%
\pgfpathlineto{\pgfqpoint{2.521853in}{0.495228in}}%
\pgfpathlineto{\pgfqpoint{2.517736in}{0.501110in}}%
\pgfpathlineto{\pgfqpoint{2.522009in}{0.510090in}}%
\pgfpathlineto{\pgfqpoint{2.512188in}{0.520102in}}%
\pgfpathlineto{\pgfqpoint{2.549990in}{0.525560in}}%
\pgfpathlineto{\pgfqpoint{2.550293in}{0.531506in}}%
\pgfpathlineto{\pgfqpoint{2.547864in}{0.546632in}}%
\pgfpathlineto{\pgfqpoint{2.550313in}{0.542005in}}%
\pgfpathlineto{\pgfqpoint{2.555178in}{0.540975in}}%
\pgfpathlineto{\pgfqpoint{2.557842in}{0.535907in}}%
\pgfpathlineto{\pgfqpoint{2.567244in}{0.529939in}}%
\pgfpathlineto{\pgfqpoint{2.570842in}{0.520187in}}%
\pgfpathlineto{\pgfqpoint{2.576314in}{0.520862in}}%
\pgfpathlineto{\pgfqpoint{2.581927in}{0.518998in}}%
\pgfpathlineto{\pgfqpoint{2.584158in}{0.521537in}}%
\pgfpathlineto{\pgfqpoint{2.592396in}{0.508481in}}%
\pgfpathlineto{\pgfqpoint{2.599396in}{0.502567in}}%
\pgfpathlineto{\pgfqpoint{2.599254in}{0.497364in}}%
\pgfpathlineto{\pgfqpoint{2.602369in}{0.492044in}}%
\pgfpathlineto{\pgfqpoint{2.609094in}{0.445229in}}%
\pgfpathlineto{\pgfqpoint{2.593275in}{0.442838in}}%
\pgfpathlineto{\pgfqpoint{2.592017in}{0.447697in}}%
\pgfpathlineto{\pgfqpoint{2.587084in}{0.454555in}}%
\pgfpathlineto{\pgfqpoint{2.581515in}{0.455951in}}%
\pgfpathlineto{\pgfqpoint{2.580449in}{0.459760in}}%
\pgfpathlineto{\pgfqpoint{2.575198in}{0.465448in}}%
\pgfpathlineto{\pgfqpoint{2.565859in}{0.472579in}}%
\pgfpathlineto{\pgfqpoint{2.562765in}{0.480015in}}%
\pgfpathlineto{\pgfqpoint{2.556944in}{0.479141in}}%
\pgfpathlineto{\pgfqpoint{2.556078in}{0.484967in}}%
\pgfpathlineto{\pgfqpoint{2.538576in}{0.482276in}}%
\pgfpathlineto{\pgfqpoint{2.533154in}{0.479245in}}%
\pgfpathclose%
\pgfusepath{fill}%
\end{pgfscope}%
\begin{pgfscope}%
\pgfpathrectangle{\pgfqpoint{0.100000in}{0.100000in}}{\pgfqpoint{3.007045in}{1.925000in}}%
\pgfusepath{clip}%
\pgfsetbuttcap%
\pgfsetmiterjoin%
\definecolor{currentfill}{rgb}{0.154787,0.468512,0.722876}%
\pgfsetfillcolor{currentfill}%
\pgfsetlinewidth{0.000000pt}%
\definecolor{currentstroke}{rgb}{0.000000,0.000000,0.000000}%
\pgfsetstrokecolor{currentstroke}%
\pgfsetstrokeopacity{0.000000}%
\pgfsetdash{}{0pt}%
\pgfpathmoveto{\pgfqpoint{1.453227in}{1.011793in}}%
\pgfpathlineto{\pgfqpoint{1.459593in}{1.011387in}}%
\pgfpathlineto{\pgfqpoint{1.456990in}{0.978623in}}%
\pgfpathlineto{\pgfqpoint{1.421577in}{0.981072in}}%
\pgfpathlineto{\pgfqpoint{1.419469in}{0.951860in}}%
\pgfpathlineto{\pgfqpoint{1.390781in}{0.954030in}}%
\pgfpathlineto{\pgfqpoint{1.393120in}{0.983287in}}%
\pgfpathlineto{\pgfqpoint{1.348759in}{0.987002in}}%
\pgfpathlineto{\pgfqpoint{1.351668in}{1.019814in}}%
\pgfpathlineto{\pgfqpoint{1.367585in}{1.018108in}}%
\pgfpathlineto{\pgfqpoint{1.401987in}{1.015175in}}%
\pgfpathlineto{\pgfqpoint{1.404039in}{1.041199in}}%
\pgfpathlineto{\pgfqpoint{1.429358in}{1.039255in}}%
\pgfpathlineto{\pgfqpoint{1.427563in}{1.013412in}}%
\pgfpathlineto{\pgfqpoint{1.453227in}{1.011793in}}%
\pgfpathclose%
\pgfusepath{fill}%
\end{pgfscope}%
\begin{pgfscope}%
\pgfpathrectangle{\pgfqpoint{0.100000in}{0.100000in}}{\pgfqpoint{3.007045in}{1.925000in}}%
\pgfusepath{clip}%
\pgfsetbuttcap%
\pgfsetmiterjoin%
\definecolor{currentfill}{rgb}{0.290980,0.594510,0.789020}%
\pgfsetfillcolor{currentfill}%
\pgfsetlinewidth{0.000000pt}%
\definecolor{currentstroke}{rgb}{0.000000,0.000000,0.000000}%
\pgfsetstrokecolor{currentstroke}%
\pgfsetstrokeopacity{0.000000}%
\pgfsetdash{}{0pt}%
\pgfpathmoveto{\pgfqpoint{2.218727in}{1.043400in}}%
\pgfpathlineto{\pgfqpoint{2.223300in}{1.033829in}}%
\pgfpathlineto{\pgfqpoint{2.216520in}{1.033633in}}%
\pgfpathlineto{\pgfqpoint{2.212201in}{1.031659in}}%
\pgfpathlineto{\pgfqpoint{2.209941in}{1.021313in}}%
\pgfpathlineto{\pgfqpoint{2.200660in}{1.024227in}}%
\pgfpathlineto{\pgfqpoint{2.189732in}{1.023554in}}%
\pgfpathlineto{\pgfqpoint{2.179536in}{1.021312in}}%
\pgfpathlineto{\pgfqpoint{2.174509in}{1.028592in}}%
\pgfpathlineto{\pgfqpoint{2.171148in}{1.034806in}}%
\pgfpathlineto{\pgfqpoint{2.174985in}{1.043028in}}%
\pgfpathlineto{\pgfqpoint{2.172018in}{1.046217in}}%
\pgfpathlineto{\pgfqpoint{2.170940in}{1.054187in}}%
\pgfpathlineto{\pgfqpoint{2.164342in}{1.057050in}}%
\pgfpathlineto{\pgfqpoint{2.171424in}{1.072069in}}%
\pgfpathlineto{\pgfqpoint{2.172010in}{1.077414in}}%
\pgfpathlineto{\pgfqpoint{2.184734in}{1.072536in}}%
\pgfpathlineto{\pgfqpoint{2.185322in}{1.077996in}}%
\pgfpathlineto{\pgfqpoint{2.197741in}{1.085673in}}%
\pgfpathlineto{\pgfqpoint{2.200736in}{1.080556in}}%
\pgfpathlineto{\pgfqpoint{2.205878in}{1.081086in}}%
\pgfpathlineto{\pgfqpoint{2.205852in}{1.076709in}}%
\pgfpathlineto{\pgfqpoint{2.204976in}{1.073343in}}%
\pgfpathlineto{\pgfqpoint{2.208703in}{1.064571in}}%
\pgfpathlineto{\pgfqpoint{2.207346in}{1.056620in}}%
\pgfpathlineto{\pgfqpoint{2.210809in}{1.047577in}}%
\pgfpathlineto{\pgfqpoint{2.217524in}{1.046026in}}%
\pgfpathlineto{\pgfqpoint{2.218727in}{1.043400in}}%
\pgfpathclose%
\pgfusepath{fill}%
\end{pgfscope}%
\begin{pgfscope}%
\pgfpathrectangle{\pgfqpoint{0.100000in}{0.100000in}}{\pgfqpoint{3.007045in}{1.925000in}}%
\pgfusepath{clip}%
\pgfsetbuttcap%
\pgfsetmiterjoin%
\definecolor{currentfill}{rgb}{0.381776,0.656517,0.824452}%
\pgfsetfillcolor{currentfill}%
\pgfsetlinewidth{0.000000pt}%
\definecolor{currentstroke}{rgb}{0.000000,0.000000,0.000000}%
\pgfsetstrokecolor{currentstroke}%
\pgfsetstrokeopacity{0.000000}%
\pgfsetdash{}{0pt}%
\pgfpathmoveto{\pgfqpoint{1.865473in}{1.453759in}}%
\pgfpathlineto{\pgfqpoint{1.882427in}{1.454160in}}%
\pgfpathlineto{\pgfqpoint{1.893937in}{1.453516in}}%
\pgfpathlineto{\pgfqpoint{1.894560in}{1.431480in}}%
\pgfpathlineto{\pgfqpoint{1.866993in}{1.430735in}}%
\pgfpathlineto{\pgfqpoint{1.837085in}{1.430186in}}%
\pgfpathlineto{\pgfqpoint{1.836747in}{1.453219in}}%
\pgfpathlineto{\pgfqpoint{1.865473in}{1.453759in}}%
\pgfpathclose%
\pgfusepath{fill}%
\end{pgfscope}%
\begin{pgfscope}%
\pgfpathrectangle{\pgfqpoint{0.100000in}{0.100000in}}{\pgfqpoint{3.007045in}{1.925000in}}%
\pgfusepath{clip}%
\pgfsetbuttcap%
\pgfsetmiterjoin%
\definecolor{currentfill}{rgb}{0.223806,0.537532,0.758431}%
\pgfsetfillcolor{currentfill}%
\pgfsetlinewidth{0.000000pt}%
\definecolor{currentstroke}{rgb}{0.000000,0.000000,0.000000}%
\pgfsetstrokecolor{currentstroke}%
\pgfsetstrokeopacity{0.000000}%
\pgfsetdash{}{0pt}%
\pgfpathmoveto{\pgfqpoint{1.674656in}{1.382036in}}%
\pgfpathlineto{\pgfqpoint{1.667711in}{1.384730in}}%
\pgfpathlineto{\pgfqpoint{1.668294in}{1.409172in}}%
\pgfpathlineto{\pgfqpoint{1.656789in}{1.409421in}}%
\pgfpathlineto{\pgfqpoint{1.657257in}{1.431226in}}%
\pgfpathlineto{\pgfqpoint{1.647409in}{1.431503in}}%
\pgfpathlineto{\pgfqpoint{1.648008in}{1.454576in}}%
\pgfpathlineto{\pgfqpoint{1.682306in}{1.453680in}}%
\pgfpathlineto{\pgfqpoint{1.703058in}{1.453389in}}%
\pgfpathlineto{\pgfqpoint{1.702699in}{1.430317in}}%
\pgfpathlineto{\pgfqpoint{1.695709in}{1.430437in}}%
\pgfpathlineto{\pgfqpoint{1.695898in}{1.425999in}}%
\pgfpathlineto{\pgfqpoint{1.699167in}{1.423436in}}%
\pgfpathlineto{\pgfqpoint{1.698770in}{1.417142in}}%
\pgfpathlineto{\pgfqpoint{1.701304in}{1.411915in}}%
\pgfpathlineto{\pgfqpoint{1.702210in}{1.402774in}}%
\pgfpathlineto{\pgfqpoint{1.685257in}{1.403080in}}%
\pgfpathlineto{\pgfqpoint{1.684732in}{1.378050in}}%
\pgfpathlineto{\pgfqpoint{1.674656in}{1.382036in}}%
\pgfpathclose%
\pgfusepath{fill}%
\end{pgfscope}%
\begin{pgfscope}%
\pgfpathrectangle{\pgfqpoint{0.100000in}{0.100000in}}{\pgfqpoint{3.007045in}{1.925000in}}%
\pgfusepath{clip}%
\pgfsetbuttcap%
\pgfsetmiterjoin%
\definecolor{currentfill}{rgb}{0.371688,0.649627,0.820515}%
\pgfsetfillcolor{currentfill}%
\pgfsetlinewidth{0.000000pt}%
\definecolor{currentstroke}{rgb}{0.000000,0.000000,0.000000}%
\pgfsetstrokecolor{currentstroke}%
\pgfsetstrokeopacity{0.000000}%
\pgfsetdash{}{0pt}%
\pgfpathmoveto{\pgfqpoint{1.698475in}{1.199567in}}%
\pgfpathlineto{\pgfqpoint{1.709799in}{1.199381in}}%
\pgfpathlineto{\pgfqpoint{1.709420in}{1.170767in}}%
\pgfpathlineto{\pgfqpoint{1.652076in}{1.171950in}}%
\pgfpathlineto{\pgfqpoint{1.652796in}{1.200602in}}%
\pgfpathlineto{\pgfqpoint{1.698475in}{1.199567in}}%
\pgfpathclose%
\pgfusepath{fill}%
\end{pgfscope}%
\begin{pgfscope}%
\pgfpathrectangle{\pgfqpoint{0.100000in}{0.100000in}}{\pgfqpoint{3.007045in}{1.925000in}}%
\pgfusepath{clip}%
\pgfsetbuttcap%
\pgfsetmiterjoin%
\definecolor{currentfill}{rgb}{0.523137,0.739193,0.861546}%
\pgfsetfillcolor{currentfill}%
\pgfsetlinewidth{0.000000pt}%
\definecolor{currentstroke}{rgb}{0.000000,0.000000,0.000000}%
\pgfsetstrokecolor{currentstroke}%
\pgfsetstrokeopacity{0.000000}%
\pgfsetdash{}{0pt}%
\pgfpathmoveto{\pgfqpoint{2.013813in}{1.217404in}}%
\pgfpathlineto{\pgfqpoint{2.001133in}{1.217077in}}%
\pgfpathlineto{\pgfqpoint{2.000904in}{1.222841in}}%
\pgfpathlineto{\pgfqpoint{1.977807in}{1.222266in}}%
\pgfpathlineto{\pgfqpoint{1.977034in}{1.245635in}}%
\pgfpathlineto{\pgfqpoint{1.982776in}{1.245666in}}%
\pgfpathlineto{\pgfqpoint{1.981668in}{1.274180in}}%
\pgfpathlineto{\pgfqpoint{1.998849in}{1.274698in}}%
\pgfpathlineto{\pgfqpoint{1.998647in}{1.280445in}}%
\pgfpathlineto{\pgfqpoint{2.021108in}{1.281507in}}%
\pgfpathlineto{\pgfqpoint{2.021691in}{1.270010in}}%
\pgfpathlineto{\pgfqpoint{2.022785in}{1.247040in}}%
\pgfpathlineto{\pgfqpoint{2.028561in}{1.247285in}}%
\pgfpathlineto{\pgfqpoint{2.029027in}{1.239956in}}%
\pgfpathlineto{\pgfqpoint{2.027473in}{1.235391in}}%
\pgfpathlineto{\pgfqpoint{2.021248in}{1.230411in}}%
\pgfpathlineto{\pgfqpoint{2.017601in}{1.220842in}}%
\pgfpathlineto{\pgfqpoint{2.013813in}{1.217404in}}%
\pgfpathclose%
\pgfusepath{fill}%
\end{pgfscope}%
\begin{pgfscope}%
\pgfpathrectangle{\pgfqpoint{0.100000in}{0.100000in}}{\pgfqpoint{3.007045in}{1.925000in}}%
\pgfusepath{clip}%
\pgfsetbuttcap%
\pgfsetmiterjoin%
\definecolor{currentfill}{rgb}{0.326290,0.618624,0.802799}%
\pgfsetfillcolor{currentfill}%
\pgfsetlinewidth{0.000000pt}%
\definecolor{currentstroke}{rgb}{0.000000,0.000000,0.000000}%
\pgfsetstrokecolor{currentstroke}%
\pgfsetstrokeopacity{0.000000}%
\pgfsetdash{}{0pt}%
\pgfpathmoveto{\pgfqpoint{2.291623in}{1.065687in}}%
\pgfpathlineto{\pgfqpoint{2.292317in}{1.056941in}}%
\pgfpathlineto{\pgfqpoint{2.285344in}{1.059690in}}%
\pgfpathlineto{\pgfqpoint{2.281792in}{1.056352in}}%
\pgfpathlineto{\pgfqpoint{2.276076in}{1.059145in}}%
\pgfpathlineto{\pgfqpoint{2.270021in}{1.058081in}}%
\pgfpathlineto{\pgfqpoint{2.267818in}{1.063028in}}%
\pgfpathlineto{\pgfqpoint{2.262113in}{1.062729in}}%
\pgfpathlineto{\pgfqpoint{2.261909in}{1.072347in}}%
\pgfpathlineto{\pgfqpoint{2.259515in}{1.072291in}}%
\pgfpathlineto{\pgfqpoint{2.253561in}{1.079316in}}%
\pgfpathlineto{\pgfqpoint{2.258362in}{1.081298in}}%
\pgfpathlineto{\pgfqpoint{2.265212in}{1.092380in}}%
\pgfpathlineto{\pgfqpoint{2.270164in}{1.089523in}}%
\pgfpathlineto{\pgfqpoint{2.281936in}{1.092901in}}%
\pgfpathlineto{\pgfqpoint{2.283241in}{1.088142in}}%
\pgfpathlineto{\pgfqpoint{2.289598in}{1.088435in}}%
\pgfpathlineto{\pgfqpoint{2.291416in}{1.086166in}}%
\pgfpathlineto{\pgfqpoint{2.291623in}{1.065687in}}%
\pgfpathclose%
\pgfusepath{fill}%
\end{pgfscope}%
\begin{pgfscope}%
\pgfpathrectangle{\pgfqpoint{0.100000in}{0.100000in}}{\pgfqpoint{3.007045in}{1.925000in}}%
\pgfusepath{clip}%
\pgfsetbuttcap%
\pgfsetmiterjoin%
\definecolor{currentfill}{rgb}{0.535686,0.746082,0.864252}%
\pgfsetfillcolor{currentfill}%
\pgfsetlinewidth{0.000000pt}%
\definecolor{currentstroke}{rgb}{0.000000,0.000000,0.000000}%
\pgfsetstrokecolor{currentstroke}%
\pgfsetstrokeopacity{0.000000}%
\pgfsetdash{}{0pt}%
\pgfpathmoveto{\pgfqpoint{2.406992in}{0.858796in}}%
\pgfpathlineto{\pgfqpoint{2.409663in}{0.861811in}}%
\pgfpathlineto{\pgfqpoint{2.422010in}{0.860947in}}%
\pgfpathlineto{\pgfqpoint{2.421120in}{0.857223in}}%
\pgfpathlineto{\pgfqpoint{2.431219in}{0.848011in}}%
\pgfpathlineto{\pgfqpoint{2.440442in}{0.846333in}}%
\pgfpathlineto{\pgfqpoint{2.435742in}{0.842463in}}%
\pgfpathlineto{\pgfqpoint{2.430794in}{0.832351in}}%
\pgfpathlineto{\pgfqpoint{2.430059in}{0.826283in}}%
\pgfpathlineto{\pgfqpoint{2.424340in}{0.826019in}}%
\pgfpathlineto{\pgfqpoint{2.414961in}{0.829005in}}%
\pgfpathlineto{\pgfqpoint{2.408663in}{0.823968in}}%
\pgfpathlineto{\pgfqpoint{2.403550in}{0.830476in}}%
\pgfpathlineto{\pgfqpoint{2.392091in}{0.836488in}}%
\pgfpathlineto{\pgfqpoint{2.387448in}{0.834640in}}%
\pgfpathlineto{\pgfqpoint{2.380905in}{0.840878in}}%
\pgfpathlineto{\pgfqpoint{2.382703in}{0.849302in}}%
\pgfpathlineto{\pgfqpoint{2.386057in}{0.845572in}}%
\pgfpathlineto{\pgfqpoint{2.397316in}{0.844754in}}%
\pgfpathlineto{\pgfqpoint{2.399267in}{0.840607in}}%
\pgfpathlineto{\pgfqpoint{2.406734in}{0.847180in}}%
\pgfpathlineto{\pgfqpoint{2.407803in}{0.851715in}}%
\pgfpathlineto{\pgfqpoint{2.404649in}{0.856728in}}%
\pgfpathlineto{\pgfqpoint{2.406992in}{0.858796in}}%
\pgfpathclose%
\pgfusepath{fill}%
\end{pgfscope}%
\begin{pgfscope}%
\pgfpathrectangle{\pgfqpoint{0.100000in}{0.100000in}}{\pgfqpoint{3.007045in}{1.925000in}}%
\pgfusepath{clip}%
\pgfsetbuttcap%
\pgfsetmiterjoin%
\definecolor{currentfill}{rgb}{0.376732,0.653072,0.822484}%
\pgfsetfillcolor{currentfill}%
\pgfsetlinewidth{0.000000pt}%
\definecolor{currentstroke}{rgb}{0.000000,0.000000,0.000000}%
\pgfsetstrokecolor{currentstroke}%
\pgfsetstrokeopacity{0.000000}%
\pgfsetdash{}{0pt}%
\pgfpathmoveto{\pgfqpoint{1.565129in}{1.341345in}}%
\pgfpathlineto{\pgfqpoint{1.543010in}{1.342313in}}%
\pgfpathlineto{\pgfqpoint{1.519516in}{1.343574in}}%
\pgfpathlineto{\pgfqpoint{1.520062in}{1.366510in}}%
\pgfpathlineto{\pgfqpoint{1.521365in}{1.389410in}}%
\pgfpathlineto{\pgfqpoint{1.521452in}{1.403789in}}%
\pgfpathlineto{\pgfqpoint{1.566978in}{1.401393in}}%
\pgfpathlineto{\pgfqpoint{1.565129in}{1.341345in}}%
\pgfpathclose%
\pgfusepath{fill}%
\end{pgfscope}%
\begin{pgfscope}%
\pgfpathrectangle{\pgfqpoint{0.100000in}{0.100000in}}{\pgfqpoint{3.007045in}{1.925000in}}%
\pgfusepath{clip}%
\pgfsetbuttcap%
\pgfsetmiterjoin%
\definecolor{currentfill}{rgb}{0.162907,0.476632,0.727059}%
\pgfsetfillcolor{currentfill}%
\pgfsetlinewidth{0.000000pt}%
\definecolor{currentstroke}{rgb}{0.000000,0.000000,0.000000}%
\pgfsetstrokecolor{currentstroke}%
\pgfsetstrokeopacity{0.000000}%
\pgfsetdash{}{0pt}%
\pgfpathmoveto{\pgfqpoint{1.716390in}{1.322247in}}%
\pgfpathlineto{\pgfqpoint{1.719990in}{1.317507in}}%
\pgfpathlineto{\pgfqpoint{1.717094in}{1.310143in}}%
\pgfpathlineto{\pgfqpoint{1.701191in}{1.310424in}}%
\pgfpathlineto{\pgfqpoint{1.701276in}{1.314251in}}%
\pgfpathlineto{\pgfqpoint{1.678558in}{1.314735in}}%
\pgfpathlineto{\pgfqpoint{1.672961in}{1.314877in}}%
\pgfpathlineto{\pgfqpoint{1.673497in}{1.337786in}}%
\pgfpathlineto{\pgfqpoint{1.683057in}{1.337584in}}%
\pgfpathlineto{\pgfqpoint{1.683358in}{1.349049in}}%
\pgfpathlineto{\pgfqpoint{1.688055in}{1.349861in}}%
\pgfpathlineto{\pgfqpoint{1.706125in}{1.349454in}}%
\pgfpathlineto{\pgfqpoint{1.706612in}{1.341900in}}%
\pgfpathlineto{\pgfqpoint{1.713888in}{1.330082in}}%
\pgfpathlineto{\pgfqpoint{1.716907in}{1.329278in}}%
\pgfpathlineto{\pgfqpoint{1.716390in}{1.322247in}}%
\pgfpathclose%
\pgfusepath{fill}%
\end{pgfscope}%
\begin{pgfscope}%
\pgfpathrectangle{\pgfqpoint{0.100000in}{0.100000in}}{\pgfqpoint{3.007045in}{1.925000in}}%
\pgfusepath{clip}%
\pgfsetbuttcap%
\pgfsetmiterjoin%
\definecolor{currentfill}{rgb}{0.351511,0.635848,0.812641}%
\pgfsetfillcolor{currentfill}%
\pgfsetlinewidth{0.000000pt}%
\definecolor{currentstroke}{rgb}{0.000000,0.000000,0.000000}%
\pgfsetstrokecolor{currentstroke}%
\pgfsetstrokeopacity{0.000000}%
\pgfsetdash{}{0pt}%
\pgfpathmoveto{\pgfqpoint{1.898889in}{1.500432in}}%
\pgfpathlineto{\pgfqpoint{1.889296in}{1.502422in}}%
\pgfpathlineto{\pgfqpoint{1.884271in}{1.505208in}}%
\pgfpathlineto{\pgfqpoint{1.876275in}{1.505700in}}%
\pgfpathlineto{\pgfqpoint{1.876437in}{1.499997in}}%
\pgfpathlineto{\pgfqpoint{1.864831in}{1.497891in}}%
\pgfpathlineto{\pgfqpoint{1.864898in}{1.494975in}}%
\pgfpathlineto{\pgfqpoint{1.853444in}{1.494748in}}%
\pgfpathlineto{\pgfqpoint{1.853375in}{1.499519in}}%
\pgfpathlineto{\pgfqpoint{1.841921in}{1.499291in}}%
\pgfpathlineto{\pgfqpoint{1.823694in}{1.499014in}}%
\pgfpathlineto{\pgfqpoint{1.824600in}{1.510542in}}%
\pgfpathlineto{\pgfqpoint{1.818879in}{1.510489in}}%
\pgfpathlineto{\pgfqpoint{1.818599in}{1.527791in}}%
\pgfpathlineto{\pgfqpoint{1.807185in}{1.527732in}}%
\pgfpathlineto{\pgfqpoint{1.806790in}{1.547902in}}%
\pgfpathlineto{\pgfqpoint{1.813546in}{1.550942in}}%
\pgfpathlineto{\pgfqpoint{1.816671in}{1.557157in}}%
\pgfpathlineto{\pgfqpoint{1.830410in}{1.548866in}}%
\pgfpathlineto{\pgfqpoint{1.838457in}{1.548803in}}%
\pgfpathlineto{\pgfqpoint{1.841735in}{1.545791in}}%
\pgfpathlineto{\pgfqpoint{1.841127in}{1.578186in}}%
\pgfpathlineto{\pgfqpoint{1.872272in}{1.578521in}}%
\pgfpathlineto{\pgfqpoint{1.870534in}{1.568228in}}%
\pgfpathlineto{\pgfqpoint{1.875816in}{1.567788in}}%
\pgfpathlineto{\pgfqpoint{1.881731in}{1.559757in}}%
\pgfpathlineto{\pgfqpoint{1.881663in}{1.556766in}}%
\pgfpathlineto{\pgfqpoint{1.876801in}{1.549513in}}%
\pgfpathlineto{\pgfqpoint{1.876967in}{1.544122in}}%
\pgfpathlineto{\pgfqpoint{1.905032in}{1.544903in}}%
\pgfpathlineto{\pgfqpoint{1.906167in}{1.539095in}}%
\pgfpathlineto{\pgfqpoint{1.907419in}{1.500598in}}%
\pgfpathlineto{\pgfqpoint{1.898889in}{1.500432in}}%
\pgfpathclose%
\pgfusepath{fill}%
\end{pgfscope}%
\begin{pgfscope}%
\pgfpathrectangle{\pgfqpoint{0.100000in}{0.100000in}}{\pgfqpoint{3.007045in}{1.925000in}}%
\pgfusepath{clip}%
\pgfsetbuttcap%
\pgfsetmiterjoin%
\definecolor{currentfill}{rgb}{0.326290,0.618624,0.802799}%
\pgfsetfillcolor{currentfill}%
\pgfsetlinewidth{0.000000pt}%
\definecolor{currentstroke}{rgb}{0.000000,0.000000,0.000000}%
\pgfsetstrokecolor{currentstroke}%
\pgfsetstrokeopacity{0.000000}%
\pgfsetdash{}{0pt}%
\pgfpathmoveto{\pgfqpoint{2.037435in}{0.611183in}}%
\pgfpathlineto{\pgfqpoint{2.043709in}{0.592924in}}%
\pgfpathlineto{\pgfqpoint{2.046065in}{0.559905in}}%
\pgfpathlineto{\pgfqpoint{2.042685in}{0.564352in}}%
\pgfpathlineto{\pgfqpoint{2.032903in}{0.562233in}}%
\pgfpathlineto{\pgfqpoint{2.029617in}{0.558203in}}%
\pgfpathlineto{\pgfqpoint{2.023323in}{0.554899in}}%
\pgfpathlineto{\pgfqpoint{2.007733in}{0.550626in}}%
\pgfpathlineto{\pgfqpoint{2.005096in}{0.547382in}}%
\pgfpathlineto{\pgfqpoint{2.002634in}{0.548120in}}%
\pgfpathlineto{\pgfqpoint{1.990185in}{0.543999in}}%
\pgfpathlineto{\pgfqpoint{1.976175in}{0.548583in}}%
\pgfpathlineto{\pgfqpoint{1.975186in}{0.557014in}}%
\pgfpathlineto{\pgfqpoint{1.970123in}{0.557651in}}%
\pgfpathlineto{\pgfqpoint{1.966602in}{0.562051in}}%
\pgfpathlineto{\pgfqpoint{1.965180in}{0.570849in}}%
\pgfpathlineto{\pgfqpoint{1.959526in}{0.575009in}}%
\pgfpathlineto{\pgfqpoint{1.958167in}{0.578691in}}%
\pgfpathlineto{\pgfqpoint{1.959192in}{0.587885in}}%
\pgfpathlineto{\pgfqpoint{1.954185in}{0.597454in}}%
\pgfpathlineto{\pgfqpoint{1.954440in}{0.605869in}}%
\pgfpathlineto{\pgfqpoint{1.958935in}{0.610678in}}%
\pgfpathlineto{\pgfqpoint{1.963991in}{0.607879in}}%
\pgfpathlineto{\pgfqpoint{2.010177in}{0.609791in}}%
\pgfpathlineto{\pgfqpoint{2.037435in}{0.611183in}}%
\pgfpathclose%
\pgfusepath{fill}%
\end{pgfscope}%
\begin{pgfscope}%
\pgfpathrectangle{\pgfqpoint{0.100000in}{0.100000in}}{\pgfqpoint{3.007045in}{1.925000in}}%
\pgfusepath{clip}%
\pgfsetbuttcap%
\pgfsetmiterjoin%
\definecolor{currentfill}{rgb}{0.301069,0.601399,0.792957}%
\pgfsetfillcolor{currentfill}%
\pgfsetlinewidth{0.000000pt}%
\definecolor{currentstroke}{rgb}{0.000000,0.000000,0.000000}%
\pgfsetstrokecolor{currentstroke}%
\pgfsetstrokeopacity{0.000000}%
\pgfsetdash{}{0pt}%
\pgfpathmoveto{\pgfqpoint{1.596180in}{1.004385in}}%
\pgfpathlineto{\pgfqpoint{1.561807in}{1.005883in}}%
\pgfpathlineto{\pgfqpoint{1.562899in}{1.031213in}}%
\pgfpathlineto{\pgfqpoint{1.534548in}{1.032337in}}%
\pgfpathlineto{\pgfqpoint{1.534859in}{1.055660in}}%
\pgfpathlineto{\pgfqpoint{1.563272in}{1.054193in}}%
\pgfpathlineto{\pgfqpoint{1.563552in}{1.060256in}}%
\pgfpathlineto{\pgfqpoint{1.591676in}{1.059068in}}%
\pgfpathlineto{\pgfqpoint{1.591838in}{1.053004in}}%
\pgfpathlineto{\pgfqpoint{1.590957in}{1.030045in}}%
\pgfpathlineto{\pgfqpoint{1.596980in}{1.029821in}}%
\pgfpathlineto{\pgfqpoint{1.596180in}{1.004385in}}%
\pgfpathclose%
\pgfusepath{fill}%
\end{pgfscope}%
\begin{pgfscope}%
\pgfpathrectangle{\pgfqpoint{0.100000in}{0.100000in}}{\pgfqpoint{3.007045in}{1.925000in}}%
\pgfusepath{clip}%
\pgfsetbuttcap%
\pgfsetmiterjoin%
\definecolor{currentfill}{rgb}{0.485490,0.718524,0.853426}%
\pgfsetfillcolor{currentfill}%
\pgfsetlinewidth{0.000000pt}%
\definecolor{currentstroke}{rgb}{0.000000,0.000000,0.000000}%
\pgfsetstrokecolor{currentstroke}%
\pgfsetstrokeopacity{0.000000}%
\pgfsetdash{}{0pt}%
\pgfpathmoveto{\pgfqpoint{2.099536in}{1.444895in}}%
\pgfpathlineto{\pgfqpoint{2.088107in}{1.444073in}}%
\pgfpathlineto{\pgfqpoint{2.087655in}{1.449833in}}%
\pgfpathlineto{\pgfqpoint{2.058743in}{1.447959in}}%
\pgfpathlineto{\pgfqpoint{2.047314in}{1.447880in}}%
\pgfpathlineto{\pgfqpoint{2.046797in}{1.456021in}}%
\pgfpathlineto{\pgfqpoint{2.050014in}{1.463525in}}%
\pgfpathlineto{\pgfqpoint{2.049555in}{1.470541in}}%
\pgfpathlineto{\pgfqpoint{2.029155in}{1.469240in}}%
\pgfpathlineto{\pgfqpoint{2.028120in}{1.486633in}}%
\pgfpathlineto{\pgfqpoint{2.045773in}{1.487575in}}%
\pgfpathlineto{\pgfqpoint{2.043986in}{1.516491in}}%
\pgfpathlineto{\pgfqpoint{2.073048in}{1.518236in}}%
\pgfpathlineto{\pgfqpoint{2.073413in}{1.512435in}}%
\pgfpathlineto{\pgfqpoint{2.090405in}{1.513286in}}%
\pgfpathlineto{\pgfqpoint{2.092961in}{1.513610in}}%
\pgfpathlineto{\pgfqpoint{2.094574in}{1.490855in}}%
\pgfpathlineto{\pgfqpoint{2.101634in}{1.491321in}}%
\pgfpathlineto{\pgfqpoint{2.103472in}{1.468301in}}%
\pgfpathlineto{\pgfqpoint{2.097758in}{1.467853in}}%
\pgfpathlineto{\pgfqpoint{2.099536in}{1.444895in}}%
\pgfpathclose%
\pgfusepath{fill}%
\end{pgfscope}%
\begin{pgfscope}%
\pgfpathrectangle{\pgfqpoint{0.100000in}{0.100000in}}{\pgfqpoint{3.007045in}{1.925000in}}%
\pgfusepath{clip}%
\pgfsetbuttcap%
\pgfsetmiterjoin%
\definecolor{currentfill}{rgb}{0.466667,0.708189,0.849366}%
\pgfsetfillcolor{currentfill}%
\pgfsetlinewidth{0.000000pt}%
\definecolor{currentstroke}{rgb}{0.000000,0.000000,0.000000}%
\pgfsetstrokecolor{currentstroke}%
\pgfsetstrokeopacity{0.000000}%
\pgfsetdash{}{0pt}%
\pgfpathmoveto{\pgfqpoint{2.639205in}{0.819156in}}%
\pgfpathlineto{\pgfqpoint{2.636694in}{0.812102in}}%
\pgfpathlineto{\pgfqpoint{2.628381in}{0.812335in}}%
\pgfpathlineto{\pgfqpoint{2.623221in}{0.806040in}}%
\pgfpathlineto{\pgfqpoint{2.626021in}{0.803536in}}%
\pgfpathlineto{\pgfqpoint{2.619999in}{0.798689in}}%
\pgfpathlineto{\pgfqpoint{2.618917in}{0.795608in}}%
\pgfpathlineto{\pgfqpoint{2.611802in}{0.789793in}}%
\pgfpathlineto{\pgfqpoint{2.611478in}{0.786111in}}%
\pgfpathlineto{\pgfqpoint{2.605521in}{0.780009in}}%
\pgfpathlineto{\pgfqpoint{2.599364in}{0.777779in}}%
\pgfpathlineto{\pgfqpoint{2.588994in}{0.768511in}}%
\pgfpathlineto{\pgfqpoint{2.581007in}{0.767715in}}%
\pgfpathlineto{\pgfqpoint{2.576021in}{0.772345in}}%
\pgfpathlineto{\pgfqpoint{2.572033in}{0.772743in}}%
\pgfpathlineto{\pgfqpoint{2.567570in}{0.778188in}}%
\pgfpathlineto{\pgfqpoint{2.559434in}{0.779009in}}%
\pgfpathlineto{\pgfqpoint{2.554010in}{0.786819in}}%
\pgfpathlineto{\pgfqpoint{2.548566in}{0.790531in}}%
\pgfpathlineto{\pgfqpoint{2.542055in}{0.797900in}}%
\pgfpathlineto{\pgfqpoint{2.556125in}{0.810180in}}%
\pgfpathlineto{\pgfqpoint{2.570619in}{0.823049in}}%
\pgfpathlineto{\pgfqpoint{2.575563in}{0.818701in}}%
\pgfpathlineto{\pgfqpoint{2.584582in}{0.822960in}}%
\pgfpathlineto{\pgfqpoint{2.584761in}{0.832670in}}%
\pgfpathlineto{\pgfqpoint{2.590755in}{0.837224in}}%
\pgfpathlineto{\pgfqpoint{2.597720in}{0.838651in}}%
\pgfpathlineto{\pgfqpoint{2.612678in}{0.829037in}}%
\pgfpathlineto{\pgfqpoint{2.628686in}{0.826071in}}%
\pgfpathlineto{\pgfqpoint{2.639205in}{0.819156in}}%
\pgfpathclose%
\pgfusepath{fill}%
\end{pgfscope}%
\begin{pgfscope}%
\pgfpathrectangle{\pgfqpoint{0.100000in}{0.100000in}}{\pgfqpoint{3.007045in}{1.925000in}}%
\pgfusepath{clip}%
\pgfsetbuttcap%
\pgfsetmiterjoin%
\definecolor{currentfill}{rgb}{0.396909,0.666851,0.830358}%
\pgfsetfillcolor{currentfill}%
\pgfsetlinewidth{0.000000pt}%
\definecolor{currentstroke}{rgb}{0.000000,0.000000,0.000000}%
\pgfsetstrokecolor{currentstroke}%
\pgfsetstrokeopacity{0.000000}%
\pgfsetdash{}{0pt}%
\pgfpathmoveto{\pgfqpoint{2.675874in}{0.992417in}}%
\pgfpathlineto{\pgfqpoint{2.673153in}{0.980638in}}%
\pgfpathlineto{\pgfqpoint{2.674050in}{0.976293in}}%
\pgfpathlineto{\pgfqpoint{2.667375in}{0.972671in}}%
\pgfpathlineto{\pgfqpoint{2.661577in}{0.967658in}}%
\pgfpathlineto{\pgfqpoint{2.657195in}{0.967322in}}%
\pgfpathlineto{\pgfqpoint{2.653930in}{0.970205in}}%
\pgfpathlineto{\pgfqpoint{2.650737in}{0.964955in}}%
\pgfpathlineto{\pgfqpoint{2.641093in}{0.964571in}}%
\pgfpathlineto{\pgfqpoint{2.625940in}{0.956912in}}%
\pgfpathlineto{\pgfqpoint{2.618329in}{0.960663in}}%
\pgfpathlineto{\pgfqpoint{2.608908in}{0.972810in}}%
\pgfpathlineto{\pgfqpoint{2.608759in}{0.975733in}}%
\pgfpathlineto{\pgfqpoint{2.597898in}{0.973667in}}%
\pgfpathlineto{\pgfqpoint{2.594940in}{0.995133in}}%
\pgfpathlineto{\pgfqpoint{2.611008in}{0.997970in}}%
\pgfpathlineto{\pgfqpoint{2.608855in}{1.001837in}}%
\pgfpathlineto{\pgfqpoint{2.605433in}{1.023762in}}%
\pgfpathlineto{\pgfqpoint{2.610860in}{1.024597in}}%
\pgfpathlineto{\pgfqpoint{2.608257in}{1.044199in}}%
\pgfpathlineto{\pgfqpoint{2.643747in}{1.050644in}}%
\pgfpathlineto{\pgfqpoint{2.641967in}{1.044405in}}%
\pgfpathlineto{\pgfqpoint{2.645335in}{1.026823in}}%
\pgfpathlineto{\pgfqpoint{2.650549in}{1.026900in}}%
\pgfpathlineto{\pgfqpoint{2.654897in}{1.034173in}}%
\pgfpathlineto{\pgfqpoint{2.664233in}{1.034590in}}%
\pgfpathlineto{\pgfqpoint{2.671365in}{1.033005in}}%
\pgfpathlineto{\pgfqpoint{2.663029in}{1.005596in}}%
\pgfpathlineto{\pgfqpoint{2.668059in}{1.000440in}}%
\pgfpathlineto{\pgfqpoint{2.672444in}{0.992955in}}%
\pgfpathlineto{\pgfqpoint{2.675874in}{0.992417in}}%
\pgfpathclose%
\pgfusepath{fill}%
\end{pgfscope}%
\begin{pgfscope}%
\pgfpathrectangle{\pgfqpoint{0.100000in}{0.100000in}}{\pgfqpoint{3.007045in}{1.925000in}}%
\pgfusepath{clip}%
\pgfsetbuttcap%
\pgfsetmiterjoin%
\definecolor{currentfill}{rgb}{0.401953,0.670296,0.832326}%
\pgfsetfillcolor{currentfill}%
\pgfsetlinewidth{0.000000pt}%
\definecolor{currentstroke}{rgb}{0.000000,0.000000,0.000000}%
\pgfsetstrokecolor{currentstroke}%
\pgfsetstrokeopacity{0.000000}%
\pgfsetdash{}{0pt}%
\pgfpathmoveto{\pgfqpoint{1.584352in}{1.202908in}}%
\pgfpathlineto{\pgfqpoint{1.561531in}{1.203886in}}%
\pgfpathlineto{\pgfqpoint{1.563506in}{1.247230in}}%
\pgfpathlineto{\pgfqpoint{1.551658in}{1.248467in}}%
\pgfpathlineto{\pgfqpoint{1.552431in}{1.273213in}}%
\pgfpathlineto{\pgfqpoint{1.563271in}{1.272731in}}%
\pgfpathlineto{\pgfqpoint{1.564548in}{1.295579in}}%
\pgfpathlineto{\pgfqpoint{1.587220in}{1.294590in}}%
\pgfpathlineto{\pgfqpoint{1.586068in}{1.271711in}}%
\pgfpathlineto{\pgfqpoint{1.587386in}{1.271651in}}%
\pgfpathlineto{\pgfqpoint{1.586476in}{1.248770in}}%
\pgfpathlineto{\pgfqpoint{1.584352in}{1.202908in}}%
\pgfpathclose%
\pgfusepath{fill}%
\end{pgfscope}%
\begin{pgfscope}%
\pgfpathrectangle{\pgfqpoint{0.100000in}{0.100000in}}{\pgfqpoint{3.007045in}{1.925000in}}%
\pgfusepath{clip}%
\pgfsetbuttcap%
\pgfsetmiterjoin%
\definecolor{currentfill}{rgb}{0.361599,0.642737,0.816578}%
\pgfsetfillcolor{currentfill}%
\pgfsetlinewidth{0.000000pt}%
\definecolor{currentstroke}{rgb}{0.000000,0.000000,0.000000}%
\pgfsetstrokecolor{currentstroke}%
\pgfsetstrokeopacity{0.000000}%
\pgfsetdash{}{0pt}%
\pgfpathmoveto{\pgfqpoint{0.713573in}{1.866288in}}%
\pgfpathlineto{\pgfqpoint{0.708324in}{1.867692in}}%
\pgfpathlineto{\pgfqpoint{0.702364in}{1.864190in}}%
\pgfpathlineto{\pgfqpoint{0.700822in}{1.858600in}}%
\pgfpathlineto{\pgfqpoint{0.695822in}{1.854965in}}%
\pgfpathlineto{\pgfqpoint{0.691919in}{1.848024in}}%
\pgfpathlineto{\pgfqpoint{0.689171in}{1.848788in}}%
\pgfpathlineto{\pgfqpoint{0.682133in}{1.844866in}}%
\pgfpathlineto{\pgfqpoint{0.680569in}{1.839333in}}%
\pgfpathlineto{\pgfqpoint{0.665649in}{1.843218in}}%
\pgfpathlineto{\pgfqpoint{0.662469in}{1.839071in}}%
\pgfpathlineto{\pgfqpoint{0.657936in}{1.837275in}}%
\pgfpathlineto{\pgfqpoint{0.656028in}{1.830896in}}%
\pgfpathlineto{\pgfqpoint{0.652989in}{1.834668in}}%
\pgfpathlineto{\pgfqpoint{0.640168in}{1.838234in}}%
\pgfpathlineto{\pgfqpoint{0.635771in}{1.843623in}}%
\pgfpathlineto{\pgfqpoint{0.629717in}{1.846964in}}%
\pgfpathlineto{\pgfqpoint{0.627702in}{1.851025in}}%
\pgfpathlineto{\pgfqpoint{0.620178in}{1.855997in}}%
\pgfpathlineto{\pgfqpoint{0.618659in}{1.862354in}}%
\pgfpathlineto{\pgfqpoint{0.615087in}{1.868865in}}%
\pgfpathlineto{\pgfqpoint{0.619147in}{1.875248in}}%
\pgfpathlineto{\pgfqpoint{0.617908in}{1.888997in}}%
\pgfpathlineto{\pgfqpoint{0.622765in}{1.897664in}}%
\pgfpathlineto{\pgfqpoint{0.631834in}{1.900038in}}%
\pgfpathlineto{\pgfqpoint{0.634619in}{1.904375in}}%
\pgfpathlineto{\pgfqpoint{0.630805in}{1.915337in}}%
\pgfpathlineto{\pgfqpoint{0.633546in}{1.923434in}}%
\pgfpathlineto{\pgfqpoint{0.643966in}{1.926557in}}%
\pgfpathlineto{\pgfqpoint{0.650253in}{1.923532in}}%
\pgfpathlineto{\pgfqpoint{0.650154in}{1.917880in}}%
\pgfpathlineto{\pgfqpoint{0.654234in}{1.907098in}}%
\pgfpathlineto{\pgfqpoint{0.659239in}{1.897832in}}%
\pgfpathlineto{\pgfqpoint{0.657890in}{1.895448in}}%
\pgfpathlineto{\pgfqpoint{0.666800in}{1.883212in}}%
\pgfpathlineto{\pgfqpoint{0.667442in}{1.879016in}}%
\pgfpathlineto{\pgfqpoint{0.675310in}{1.876711in}}%
\pgfpathlineto{\pgfqpoint{0.676269in}{1.883139in}}%
\pgfpathlineto{\pgfqpoint{0.683348in}{1.884220in}}%
\pgfpathlineto{\pgfqpoint{0.685273in}{1.876771in}}%
\pgfpathlineto{\pgfqpoint{0.695862in}{1.879153in}}%
\pgfpathlineto{\pgfqpoint{0.699664in}{1.876125in}}%
\pgfpathlineto{\pgfqpoint{0.710511in}{1.880181in}}%
\pgfpathlineto{\pgfqpoint{0.713541in}{1.878149in}}%
\pgfpathlineto{\pgfqpoint{0.713573in}{1.866288in}}%
\pgfpathclose%
\pgfusepath{fill}%
\end{pgfscope}%
\begin{pgfscope}%
\pgfpathrectangle{\pgfqpoint{0.100000in}{0.100000in}}{\pgfqpoint{3.007045in}{1.925000in}}%
\pgfusepath{clip}%
\pgfsetbuttcap%
\pgfsetmiterjoin%
\definecolor{currentfill}{rgb}{0.371688,0.649627,0.820515}%
\pgfsetfillcolor{currentfill}%
\pgfsetlinewidth{0.000000pt}%
\definecolor{currentstroke}{rgb}{0.000000,0.000000,0.000000}%
\pgfsetstrokecolor{currentstroke}%
\pgfsetstrokeopacity{0.000000}%
\pgfsetdash{}{0pt}%
\pgfpathmoveto{\pgfqpoint{2.087534in}{0.559447in}}%
\pgfpathlineto{\pgfqpoint{2.082124in}{0.561840in}}%
\pgfpathlineto{\pgfqpoint{2.077352in}{0.576813in}}%
\pgfpathlineto{\pgfqpoint{2.072286in}{0.580627in}}%
\pgfpathlineto{\pgfqpoint{2.067276in}{0.590315in}}%
\pgfpathlineto{\pgfqpoint{2.067754in}{0.598346in}}%
\pgfpathlineto{\pgfqpoint{2.070931in}{0.604028in}}%
\pgfpathlineto{\pgfqpoint{2.072666in}{0.613237in}}%
\pgfpathlineto{\pgfqpoint{2.076881in}{0.613497in}}%
\pgfpathlineto{\pgfqpoint{2.085520in}{0.614773in}}%
\pgfpathlineto{\pgfqpoint{2.094715in}{0.614481in}}%
\pgfpathlineto{\pgfqpoint{2.095084in}{0.608430in}}%
\pgfpathlineto{\pgfqpoint{2.123918in}{0.610337in}}%
\pgfpathlineto{\pgfqpoint{2.121151in}{0.645024in}}%
\pgfpathlineto{\pgfqpoint{2.143322in}{0.646719in}}%
\pgfpathlineto{\pgfqpoint{2.146781in}{0.617810in}}%
\pgfpathlineto{\pgfqpoint{2.151467in}{0.577078in}}%
\pgfpathlineto{\pgfqpoint{2.147381in}{0.572588in}}%
\pgfpathlineto{\pgfqpoint{2.139247in}{0.575561in}}%
\pgfpathlineto{\pgfqpoint{2.134440in}{0.573344in}}%
\pgfpathlineto{\pgfqpoint{2.121286in}{0.575705in}}%
\pgfpathlineto{\pgfqpoint{2.113119in}{0.573784in}}%
\pgfpathlineto{\pgfqpoint{2.098215in}{0.568019in}}%
\pgfpathlineto{\pgfqpoint{2.087534in}{0.559447in}}%
\pgfpathclose%
\pgfusepath{fill}%
\end{pgfscope}%
\begin{pgfscope}%
\pgfpathrectangle{\pgfqpoint{0.100000in}{0.100000in}}{\pgfqpoint{3.007045in}{1.925000in}}%
\pgfusepath{clip}%
\pgfsetbuttcap%
\pgfsetmiterjoin%
\definecolor{currentfill}{rgb}{0.968627,0.984314,1.000000}%
\pgfsetfillcolor{currentfill}%
\pgfsetlinewidth{0.000000pt}%
\definecolor{currentstroke}{rgb}{0.000000,0.000000,0.000000}%
\pgfsetstrokecolor{currentstroke}%
\pgfsetstrokeopacity{0.000000}%
\pgfsetdash{}{0pt}%
\pgfpathmoveto{\pgfqpoint{2.990796in}{1.458726in}}%
\pgfpathlineto{\pgfqpoint{2.997029in}{1.452912in}}%
\pgfpathlineto{\pgfqpoint{2.995205in}{1.449674in}}%
\pgfpathlineto{\pgfqpoint{2.988017in}{1.448533in}}%
\pgfpathlineto{\pgfqpoint{2.992375in}{1.455283in}}%
\pgfpathlineto{\pgfqpoint{2.990796in}{1.458726in}}%
\pgfpathclose%
\pgfusepath{fill}%
\end{pgfscope}%
\begin{pgfscope}%
\pgfpathrectangle{\pgfqpoint{0.100000in}{0.100000in}}{\pgfqpoint{3.007045in}{1.925000in}}%
\pgfusepath{clip}%
\pgfsetbuttcap%
\pgfsetmiterjoin%
\definecolor{currentfill}{rgb}{0.691580,0.822745,0.907543}%
\pgfsetfillcolor{currentfill}%
\pgfsetlinewidth{0.000000pt}%
\definecolor{currentstroke}{rgb}{0.000000,0.000000,0.000000}%
\pgfsetstrokecolor{currentstroke}%
\pgfsetstrokeopacity{0.000000}%
\pgfsetdash{}{0pt}%
\pgfpathmoveto{\pgfqpoint{2.537236in}{1.043373in}}%
\pgfpathlineto{\pgfqpoint{2.526060in}{1.056473in}}%
\pgfpathlineto{\pgfqpoint{2.521063in}{1.052157in}}%
\pgfpathlineto{\pgfqpoint{2.511594in}{1.060569in}}%
\pgfpathlineto{\pgfqpoint{2.508760in}{1.065976in}}%
\pgfpathlineto{\pgfqpoint{2.512374in}{1.068702in}}%
\pgfpathlineto{\pgfqpoint{2.502786in}{1.075737in}}%
\pgfpathlineto{\pgfqpoint{2.510605in}{1.081728in}}%
\pgfpathlineto{\pgfqpoint{2.509239in}{1.087027in}}%
\pgfpathlineto{\pgfqpoint{2.516736in}{1.108073in}}%
\pgfpathlineto{\pgfqpoint{2.527642in}{1.106923in}}%
\pgfpathlineto{\pgfqpoint{2.535565in}{1.108737in}}%
\pgfpathlineto{\pgfqpoint{2.544389in}{1.103816in}}%
\pgfpathlineto{\pgfqpoint{2.554928in}{1.112611in}}%
\pgfpathlineto{\pgfqpoint{2.559027in}{1.119923in}}%
\pgfpathlineto{\pgfqpoint{2.566770in}{1.118700in}}%
\pgfpathlineto{\pgfqpoint{2.567398in}{1.112031in}}%
\pgfpathlineto{\pgfqpoint{2.578351in}{1.105297in}}%
\pgfpathlineto{\pgfqpoint{2.573073in}{1.099388in}}%
\pgfpathlineto{\pgfqpoint{2.570652in}{1.100774in}}%
\pgfpathlineto{\pgfqpoint{2.564677in}{1.094671in}}%
\pgfpathlineto{\pgfqpoint{2.562917in}{1.088318in}}%
\pgfpathlineto{\pgfqpoint{2.564592in}{1.081776in}}%
\pgfpathlineto{\pgfqpoint{2.574988in}{1.077569in}}%
\pgfpathlineto{\pgfqpoint{2.578887in}{1.072582in}}%
\pgfpathlineto{\pgfqpoint{2.578534in}{1.060078in}}%
\pgfpathlineto{\pgfqpoint{2.572984in}{1.054602in}}%
\pgfpathlineto{\pgfqpoint{2.547424in}{1.055944in}}%
\pgfpathlineto{\pgfqpoint{2.543457in}{1.053394in}}%
\pgfpathlineto{\pgfqpoint{2.537236in}{1.043373in}}%
\pgfpathclose%
\pgfusepath{fill}%
\end{pgfscope}%
\begin{pgfscope}%
\pgfpathrectangle{\pgfqpoint{0.100000in}{0.100000in}}{\pgfqpoint{3.007045in}{1.925000in}}%
\pgfusepath{clip}%
\pgfsetbuttcap%
\pgfsetmiterjoin%
\definecolor{currentfill}{rgb}{0.260715,0.573841,0.777209}%
\pgfsetfillcolor{currentfill}%
\pgfsetlinewidth{0.000000pt}%
\definecolor{currentstroke}{rgb}{0.000000,0.000000,0.000000}%
\pgfsetstrokecolor{currentstroke}%
\pgfsetstrokeopacity{0.000000}%
\pgfsetdash{}{0pt}%
\pgfpathmoveto{\pgfqpoint{1.072883in}{1.296139in}}%
\pgfpathlineto{\pgfqpoint{1.067620in}{1.298977in}}%
\pgfpathlineto{\pgfqpoint{1.068540in}{1.305641in}}%
\pgfpathlineto{\pgfqpoint{1.067332in}{1.310685in}}%
\pgfpathlineto{\pgfqpoint{1.057913in}{1.312141in}}%
\pgfpathlineto{\pgfqpoint{1.057220in}{1.307768in}}%
\pgfpathlineto{\pgfqpoint{1.052269in}{1.305131in}}%
\pgfpathlineto{\pgfqpoint{1.050961in}{1.311223in}}%
\pgfpathlineto{\pgfqpoint{1.033610in}{1.309798in}}%
\pgfpathlineto{\pgfqpoint{1.027613in}{1.314125in}}%
\pgfpathlineto{\pgfqpoint{1.029601in}{1.326060in}}%
\pgfpathlineto{\pgfqpoint{1.005052in}{1.330045in}}%
\pgfpathlineto{\pgfqpoint{0.978339in}{1.334944in}}%
\pgfpathlineto{\pgfqpoint{0.981285in}{1.351407in}}%
\pgfpathlineto{\pgfqpoint{0.972493in}{1.347031in}}%
\pgfpathlineto{\pgfqpoint{0.969963in}{1.352958in}}%
\pgfpathlineto{\pgfqpoint{0.970085in}{1.358692in}}%
\pgfpathlineto{\pgfqpoint{0.963794in}{1.361956in}}%
\pgfpathlineto{\pgfqpoint{0.960688in}{1.366618in}}%
\pgfpathlineto{\pgfqpoint{0.965270in}{1.373029in}}%
\pgfpathlineto{\pgfqpoint{0.968734in}{1.380626in}}%
\pgfpathlineto{\pgfqpoint{0.967550in}{1.392277in}}%
\pgfpathlineto{\pgfqpoint{0.965948in}{1.394443in}}%
\pgfpathlineto{\pgfqpoint{0.968744in}{1.400234in}}%
\pgfpathlineto{\pgfqpoint{0.967728in}{1.403968in}}%
\pgfpathlineto{\pgfqpoint{0.989963in}{1.400079in}}%
\pgfpathlineto{\pgfqpoint{1.001890in}{1.466220in}}%
\pgfpathlineto{\pgfqpoint{1.005297in}{1.485446in}}%
\pgfpathlineto{\pgfqpoint{1.016022in}{1.482462in}}%
\pgfpathlineto{\pgfqpoint{1.015288in}{1.478274in}}%
\pgfpathlineto{\pgfqpoint{1.026443in}{1.476331in}}%
\pgfpathlineto{\pgfqpoint{1.027092in}{1.480000in}}%
\pgfpathlineto{\pgfqpoint{1.038193in}{1.478061in}}%
\pgfpathlineto{\pgfqpoint{1.039144in}{1.483711in}}%
\pgfpathlineto{\pgfqpoint{1.052724in}{1.481408in}}%
\pgfpathlineto{\pgfqpoint{1.053944in}{1.487032in}}%
\pgfpathlineto{\pgfqpoint{1.067969in}{1.484605in}}%
\pgfpathlineto{\pgfqpoint{1.069796in}{1.477757in}}%
\pgfpathlineto{\pgfqpoint{1.069165in}{1.468405in}}%
\pgfpathlineto{\pgfqpoint{1.070708in}{1.461628in}}%
\pgfpathlineto{\pgfqpoint{1.075159in}{1.453148in}}%
\pgfpathlineto{\pgfqpoint{1.079954in}{1.447589in}}%
\pgfpathlineto{\pgfqpoint{1.085388in}{1.435891in}}%
\pgfpathlineto{\pgfqpoint{1.092414in}{1.429227in}}%
\pgfpathlineto{\pgfqpoint{1.089929in}{1.412247in}}%
\pgfpathlineto{\pgfqpoint{1.089437in}{1.400986in}}%
\pgfpathlineto{\pgfqpoint{1.139789in}{1.393330in}}%
\pgfpathlineto{\pgfqpoint{1.162889in}{1.390045in}}%
\pgfpathlineto{\pgfqpoint{1.162230in}{1.378562in}}%
\pgfpathlineto{\pgfqpoint{1.158048in}{1.350480in}}%
\pgfpathlineto{\pgfqpoint{1.137581in}{1.353507in}}%
\pgfpathlineto{\pgfqpoint{1.135759in}{1.336378in}}%
\pgfpathlineto{\pgfqpoint{1.131992in}{1.310519in}}%
\pgfpathlineto{\pgfqpoint{1.102069in}{1.314741in}}%
\pgfpathlineto{\pgfqpoint{1.076315in}{1.318750in}}%
\pgfpathlineto{\pgfqpoint{1.072883in}{1.296139in}}%
\pgfpathclose%
\pgfusepath{fill}%
\end{pgfscope}%
\begin{pgfscope}%
\pgfpathrectangle{\pgfqpoint{0.100000in}{0.100000in}}{\pgfqpoint{3.007045in}{1.925000in}}%
\pgfusepath{clip}%
\pgfsetbuttcap%
\pgfsetmiterjoin%
\definecolor{currentfill}{rgb}{0.560784,0.759862,0.869666}%
\pgfsetfillcolor{currentfill}%
\pgfsetlinewidth{0.000000pt}%
\definecolor{currentstroke}{rgb}{0.000000,0.000000,0.000000}%
\pgfsetstrokecolor{currentstroke}%
\pgfsetstrokeopacity{0.000000}%
\pgfsetdash{}{0pt}%
\pgfpathmoveto{\pgfqpoint{2.117083in}{1.446231in}}%
\pgfpathlineto{\pgfqpoint{2.099536in}{1.444895in}}%
\pgfpathlineto{\pgfqpoint{2.097758in}{1.467853in}}%
\pgfpathlineto{\pgfqpoint{2.103472in}{1.468301in}}%
\pgfpathlineto{\pgfqpoint{2.101634in}{1.491321in}}%
\pgfpathlineto{\pgfqpoint{2.108968in}{1.491853in}}%
\pgfpathlineto{\pgfqpoint{2.108536in}{1.497608in}}%
\pgfpathlineto{\pgfqpoint{2.124853in}{1.498938in}}%
\pgfpathlineto{\pgfqpoint{2.126881in}{1.489860in}}%
\pgfpathlineto{\pgfqpoint{2.121101in}{1.483818in}}%
\pgfpathlineto{\pgfqpoint{2.118059in}{1.468764in}}%
\pgfpathlineto{\pgfqpoint{2.120658in}{1.456120in}}%
\pgfpathlineto{\pgfqpoint{2.117083in}{1.446231in}}%
\pgfpathclose%
\pgfusepath{fill}%
\end{pgfscope}%
\begin{pgfscope}%
\pgfpathrectangle{\pgfqpoint{0.100000in}{0.100000in}}{\pgfqpoint{3.007045in}{1.925000in}}%
\pgfusepath{clip}%
\pgfsetbuttcap%
\pgfsetmiterjoin%
\definecolor{currentfill}{rgb}{0.150727,0.464452,0.720784}%
\pgfsetfillcolor{currentfill}%
\pgfsetlinewidth{0.000000pt}%
\definecolor{currentstroke}{rgb}{0.000000,0.000000,0.000000}%
\pgfsetstrokecolor{currentstroke}%
\pgfsetstrokeopacity{0.000000}%
\pgfsetdash{}{0pt}%
\pgfpathmoveto{\pgfqpoint{1.330767in}{0.811238in}}%
\pgfpathlineto{\pgfqpoint{1.333803in}{0.843365in}}%
\pgfpathlineto{\pgfqpoint{1.361996in}{0.840962in}}%
\pgfpathlineto{\pgfqpoint{1.364391in}{0.869531in}}%
\pgfpathlineto{\pgfqpoint{1.392891in}{0.867332in}}%
\pgfpathlineto{\pgfqpoint{1.390704in}{0.838664in}}%
\pgfpathlineto{\pgfqpoint{1.385676in}{0.839061in}}%
\pgfpathlineto{\pgfqpoint{1.383417in}{0.806844in}}%
\pgfpathlineto{\pgfqpoint{1.330767in}{0.811238in}}%
\pgfpathclose%
\pgfusepath{fill}%
\end{pgfscope}%
\begin{pgfscope}%
\pgfpathrectangle{\pgfqpoint{0.100000in}{0.100000in}}{\pgfqpoint{3.007045in}{1.925000in}}%
\pgfusepath{clip}%
\pgfsetbuttcap%
\pgfsetmiterjoin%
\definecolor{currentfill}{rgb}{0.541961,0.749527,0.865606}%
\pgfsetfillcolor{currentfill}%
\pgfsetlinewidth{0.000000pt}%
\definecolor{currentstroke}{rgb}{0.000000,0.000000,0.000000}%
\pgfsetstrokecolor{currentstroke}%
\pgfsetstrokeopacity{0.000000}%
\pgfsetdash{}{0pt}%
\pgfpathmoveto{\pgfqpoint{0.481516in}{1.783693in}}%
\pgfpathlineto{\pgfqpoint{0.465616in}{1.788403in}}%
\pgfpathlineto{\pgfqpoint{0.454831in}{1.792649in}}%
\pgfpathlineto{\pgfqpoint{0.462406in}{1.810296in}}%
\pgfpathlineto{\pgfqpoint{0.462308in}{1.820138in}}%
\pgfpathlineto{\pgfqpoint{0.467265in}{1.814731in}}%
\pgfpathlineto{\pgfqpoint{0.473024in}{1.816689in}}%
\pgfpathlineto{\pgfqpoint{0.479649in}{1.816372in}}%
\pgfpathlineto{\pgfqpoint{0.476951in}{1.821800in}}%
\pgfpathlineto{\pgfqpoint{0.468427in}{1.820044in}}%
\pgfpathlineto{\pgfqpoint{0.462073in}{1.828534in}}%
\pgfpathlineto{\pgfqpoint{0.465398in}{1.837277in}}%
\pgfpathlineto{\pgfqpoint{0.473406in}{1.837959in}}%
\pgfpathlineto{\pgfqpoint{0.472606in}{1.850993in}}%
\pgfpathlineto{\pgfqpoint{0.469231in}{1.854603in}}%
\pgfpathlineto{\pgfqpoint{0.471018in}{1.865773in}}%
\pgfpathlineto{\pgfqpoint{0.474450in}{1.863703in}}%
\pgfpathlineto{\pgfqpoint{0.482840in}{1.866326in}}%
\pgfpathlineto{\pgfqpoint{0.472793in}{1.873445in}}%
\pgfpathlineto{\pgfqpoint{0.474678in}{1.891400in}}%
\pgfpathlineto{\pgfqpoint{0.472885in}{1.900270in}}%
\pgfpathlineto{\pgfqpoint{0.474818in}{1.908630in}}%
\pgfpathlineto{\pgfqpoint{0.487655in}{1.903257in}}%
\pgfpathlineto{\pgfqpoint{0.510848in}{1.895622in}}%
\pgfpathlineto{\pgfqpoint{0.505442in}{1.879195in}}%
\pgfpathlineto{\pgfqpoint{0.502572in}{1.867867in}}%
\pgfpathlineto{\pgfqpoint{0.514994in}{1.864013in}}%
\pgfpathlineto{\pgfqpoint{0.515078in}{1.857792in}}%
\pgfpathlineto{\pgfqpoint{0.510872in}{1.844976in}}%
\pgfpathlineto{\pgfqpoint{0.501743in}{1.847839in}}%
\pgfpathlineto{\pgfqpoint{0.493956in}{1.821847in}}%
\pgfpathlineto{\pgfqpoint{0.542761in}{1.806508in}}%
\pgfpathlineto{\pgfqpoint{0.536092in}{1.785540in}}%
\pgfpathlineto{\pgfqpoint{0.529311in}{1.780896in}}%
\pgfpathlineto{\pgfqpoint{0.524572in}{1.784544in}}%
\pgfpathlineto{\pgfqpoint{0.510888in}{1.782623in}}%
\pgfpathlineto{\pgfqpoint{0.507744in}{1.796488in}}%
\pgfpathlineto{\pgfqpoint{0.498370in}{1.806808in}}%
\pgfpathlineto{\pgfqpoint{0.488870in}{1.806888in}}%
\pgfpathlineto{\pgfqpoint{0.481516in}{1.783693in}}%
\pgfpathclose%
\pgfusepath{fill}%
\end{pgfscope}%
\begin{pgfscope}%
\pgfpathrectangle{\pgfqpoint{0.100000in}{0.100000in}}{\pgfqpoint{3.007045in}{1.925000in}}%
\pgfusepath{clip}%
\pgfsetbuttcap%
\pgfsetmiterjoin%
\definecolor{currentfill}{rgb}{0.560784,0.759862,0.869666}%
\pgfsetfillcolor{currentfill}%
\pgfsetlinewidth{0.000000pt}%
\definecolor{currentstroke}{rgb}{0.000000,0.000000,0.000000}%
\pgfsetstrokecolor{currentstroke}%
\pgfsetstrokeopacity{0.000000}%
\pgfsetdash{}{0pt}%
\pgfpathmoveto{\pgfqpoint{0.812473in}{1.770492in}}%
\pgfpathlineto{\pgfqpoint{0.807745in}{1.750662in}}%
\pgfpathlineto{\pgfqpoint{0.794807in}{1.753888in}}%
\pgfpathlineto{\pgfqpoint{0.789165in}{1.749329in}}%
\pgfpathlineto{\pgfqpoint{0.774865in}{1.752771in}}%
\pgfpathlineto{\pgfqpoint{0.773003in}{1.745268in}}%
\pgfpathlineto{\pgfqpoint{0.765843in}{1.746677in}}%
\pgfpathlineto{\pgfqpoint{0.765426in}{1.755324in}}%
\pgfpathlineto{\pgfqpoint{0.760765in}{1.761473in}}%
\pgfpathlineto{\pgfqpoint{0.758923in}{1.767006in}}%
\pgfpathlineto{\pgfqpoint{0.752113in}{1.767443in}}%
\pgfpathlineto{\pgfqpoint{0.748808in}{1.770376in}}%
\pgfpathlineto{\pgfqpoint{0.740246in}{1.767188in}}%
\pgfpathlineto{\pgfqpoint{0.724939in}{1.766494in}}%
\pgfpathlineto{\pgfqpoint{0.724517in}{1.778735in}}%
\pgfpathlineto{\pgfqpoint{0.726344in}{1.778628in}}%
\pgfpathlineto{\pgfqpoint{0.735917in}{1.779747in}}%
\pgfpathlineto{\pgfqpoint{0.740338in}{1.787082in}}%
\pgfpathlineto{\pgfqpoint{0.746029in}{1.809361in}}%
\pgfpathlineto{\pgfqpoint{0.786311in}{1.799097in}}%
\pgfpathlineto{\pgfqpoint{0.784213in}{1.790565in}}%
\pgfpathlineto{\pgfqpoint{0.794601in}{1.781859in}}%
\pgfpathlineto{\pgfqpoint{0.808364in}{1.778465in}}%
\pgfpathlineto{\pgfqpoint{0.813833in}{1.776182in}}%
\pgfpathlineto{\pgfqpoint{0.812473in}{1.770492in}}%
\pgfpathclose%
\pgfusepath{fill}%
\end{pgfscope}%
\begin{pgfscope}%
\pgfpathrectangle{\pgfqpoint{0.100000in}{0.100000in}}{\pgfqpoint{3.007045in}{1.925000in}}%
\pgfusepath{clip}%
\pgfsetbuttcap%
\pgfsetmiterjoin%
\definecolor{currentfill}{rgb}{0.579608,0.770196,0.873725}%
\pgfsetfillcolor{currentfill}%
\pgfsetlinewidth{0.000000pt}%
\definecolor{currentstroke}{rgb}{0.000000,0.000000,0.000000}%
\pgfsetstrokecolor{currentstroke}%
\pgfsetstrokeopacity{0.000000}%
\pgfsetdash{}{0pt}%
\pgfpathmoveto{\pgfqpoint{0.748143in}{0.777058in}}%
\pgfpathlineto{\pgfqpoint{0.754327in}{0.807329in}}%
\pgfpathlineto{\pgfqpoint{0.795773in}{0.799049in}}%
\pgfpathlineto{\pgfqpoint{0.851184in}{0.788621in}}%
\pgfpathlineto{\pgfqpoint{0.912796in}{0.778126in}}%
\pgfpathlineto{\pgfqpoint{0.911810in}{0.772480in}}%
\pgfpathlineto{\pgfqpoint{0.927312in}{0.769778in}}%
\pgfpathlineto{\pgfqpoint{0.989331in}{0.759721in}}%
\pgfpathlineto{\pgfqpoint{0.978016in}{0.688001in}}%
\pgfpathlineto{\pgfqpoint{0.924756in}{0.696630in}}%
\pgfpathlineto{\pgfqpoint{0.864661in}{0.706945in}}%
\pgfpathlineto{\pgfqpoint{0.849482in}{0.716019in}}%
\pgfpathlineto{\pgfqpoint{0.766721in}{0.765622in}}%
\pgfpathlineto{\pgfqpoint{0.748143in}{0.777058in}}%
\pgfpathclose%
\pgfusepath{fill}%
\end{pgfscope}%
\begin{pgfscope}%
\pgfpathrectangle{\pgfqpoint{0.100000in}{0.100000in}}{\pgfqpoint{3.007045in}{1.925000in}}%
\pgfusepath{clip}%
\pgfsetbuttcap%
\pgfsetmiterjoin%
\definecolor{currentfill}{rgb}{0.341423,0.628958,0.808704}%
\pgfsetfillcolor{currentfill}%
\pgfsetlinewidth{0.000000pt}%
\definecolor{currentstroke}{rgb}{0.000000,0.000000,0.000000}%
\pgfsetstrokecolor{currentstroke}%
\pgfsetstrokeopacity{0.000000}%
\pgfsetdash{}{0pt}%
\pgfpathmoveto{\pgfqpoint{2.644431in}{1.224390in}}%
\pgfpathlineto{\pgfqpoint{2.643347in}{1.220746in}}%
\pgfpathlineto{\pgfqpoint{2.637469in}{1.212817in}}%
\pgfpathlineto{\pgfqpoint{2.634232in}{1.214856in}}%
\pgfpathlineto{\pgfqpoint{2.621536in}{1.211734in}}%
\pgfpathlineto{\pgfqpoint{2.616014in}{1.216485in}}%
\pgfpathlineto{\pgfqpoint{2.614276in}{1.211702in}}%
\pgfpathlineto{\pgfqpoint{2.609375in}{1.212759in}}%
\pgfpathlineto{\pgfqpoint{2.604647in}{1.216037in}}%
\pgfpathlineto{\pgfqpoint{2.589864in}{1.218745in}}%
\pgfpathlineto{\pgfqpoint{2.585351in}{1.220532in}}%
\pgfpathlineto{\pgfqpoint{2.590743in}{1.227736in}}%
\pgfpathlineto{\pgfqpoint{2.592721in}{1.234363in}}%
\pgfpathlineto{\pgfqpoint{2.598024in}{1.242727in}}%
\pgfpathlineto{\pgfqpoint{2.607688in}{1.243264in}}%
\pgfpathlineto{\pgfqpoint{2.614371in}{1.241096in}}%
\pgfpathlineto{\pgfqpoint{2.644431in}{1.224390in}}%
\pgfpathclose%
\pgfusepath{fill}%
\end{pgfscope}%
\begin{pgfscope}%
\pgfpathrectangle{\pgfqpoint{0.100000in}{0.100000in}}{\pgfqpoint{3.007045in}{1.925000in}}%
\pgfusepath{clip}%
\pgfsetbuttcap%
\pgfsetmiterjoin%
\definecolor{currentfill}{rgb}{0.158847,0.472572,0.724967}%
\pgfsetfillcolor{currentfill}%
\pgfsetlinewidth{0.000000pt}%
\definecolor{currentstroke}{rgb}{0.000000,0.000000,0.000000}%
\pgfsetstrokecolor{currentstroke}%
\pgfsetstrokeopacity{0.000000}%
\pgfsetdash{}{0pt}%
\pgfpathmoveto{\pgfqpoint{1.654308in}{0.604797in}}%
\pgfpathlineto{\pgfqpoint{1.643382in}{0.598629in}}%
\pgfpathlineto{\pgfqpoint{1.639929in}{0.589583in}}%
\pgfpathlineto{\pgfqpoint{1.622469in}{0.597909in}}%
\pgfpathlineto{\pgfqpoint{1.610932in}{0.600948in}}%
\pgfpathlineto{\pgfqpoint{1.606454in}{0.609350in}}%
\pgfpathlineto{\pgfqpoint{1.576374in}{0.610043in}}%
\pgfpathlineto{\pgfqpoint{1.571686in}{0.614167in}}%
\pgfpathlineto{\pgfqpoint{1.569891in}{0.623670in}}%
\pgfpathlineto{\pgfqpoint{1.586901in}{0.635387in}}%
\pgfpathlineto{\pgfqpoint{1.616016in}{0.651438in}}%
\pgfpathlineto{\pgfqpoint{1.620418in}{0.653867in}}%
\pgfpathlineto{\pgfqpoint{1.635041in}{0.627484in}}%
\pgfpathlineto{\pgfqpoint{1.639177in}{0.622324in}}%
\pgfpathlineto{\pgfqpoint{1.642936in}{0.624588in}}%
\pgfpathlineto{\pgfqpoint{1.654308in}{0.604797in}}%
\pgfpathclose%
\pgfusepath{fill}%
\end{pgfscope}%
\begin{pgfscope}%
\pgfpathrectangle{\pgfqpoint{0.100000in}{0.100000in}}{\pgfqpoint{3.007045in}{1.925000in}}%
\pgfusepath{clip}%
\pgfsetbuttcap%
\pgfsetmiterjoin%
\definecolor{currentfill}{rgb}{0.031373,0.241015,0.497978}%
\pgfsetfillcolor{currentfill}%
\pgfsetlinewidth{0.000000pt}%
\definecolor{currentstroke}{rgb}{0.000000,0.000000,0.000000}%
\pgfsetstrokecolor{currentstroke}%
\pgfsetstrokeopacity{0.000000}%
\pgfsetdash{}{0pt}%
\pgfpathmoveto{\pgfqpoint{1.521148in}{0.313814in}}%
\pgfpathlineto{\pgfqpoint{1.530987in}{0.320341in}}%
\pgfpathlineto{\pgfqpoint{1.534731in}{0.327636in}}%
\pgfpathlineto{\pgfqpoint{1.566652in}{0.326194in}}%
\pgfpathlineto{\pgfqpoint{1.592904in}{0.325092in}}%
\pgfpathlineto{\pgfqpoint{1.592550in}{0.313976in}}%
\pgfpathlineto{\pgfqpoint{1.599953in}{0.312537in}}%
\pgfpathlineto{\pgfqpoint{1.625397in}{0.311924in}}%
\pgfpathlineto{\pgfqpoint{1.627117in}{0.305150in}}%
\pgfpathlineto{\pgfqpoint{1.625131in}{0.302682in}}%
\pgfpathlineto{\pgfqpoint{1.630801in}{0.295513in}}%
\pgfpathlineto{\pgfqpoint{1.630560in}{0.289662in}}%
\pgfpathlineto{\pgfqpoint{1.633409in}{0.284357in}}%
\pgfpathlineto{\pgfqpoint{1.633735in}{0.277987in}}%
\pgfpathlineto{\pgfqpoint{1.641756in}{0.274374in}}%
\pgfpathlineto{\pgfqpoint{1.642081in}{0.269076in}}%
\pgfpathlineto{\pgfqpoint{1.628506in}{0.265666in}}%
\pgfpathlineto{\pgfqpoint{1.624844in}{0.260718in}}%
\pgfpathlineto{\pgfqpoint{1.619737in}{0.263995in}}%
\pgfpathlineto{\pgfqpoint{1.609852in}{0.273649in}}%
\pgfpathlineto{\pgfqpoint{1.597319in}{0.276746in}}%
\pgfpathlineto{\pgfqpoint{1.586571in}{0.275707in}}%
\pgfpathlineto{\pgfqpoint{1.575367in}{0.277979in}}%
\pgfpathlineto{\pgfqpoint{1.562903in}{0.288739in}}%
\pgfpathlineto{\pgfqpoint{1.550740in}{0.289846in}}%
\pgfpathlineto{\pgfqpoint{1.542376in}{0.299233in}}%
\pgfpathlineto{\pgfqpoint{1.540217in}{0.298637in}}%
\pgfpathlineto{\pgfqpoint{1.524231in}{0.303862in}}%
\pgfpathlineto{\pgfqpoint{1.525543in}{0.307224in}}%
\pgfpathlineto{\pgfqpoint{1.521148in}{0.313814in}}%
\pgfpathclose%
\pgfusepath{fill}%
\end{pgfscope}%
\begin{pgfscope}%
\pgfpathrectangle{\pgfqpoint{0.100000in}{0.100000in}}{\pgfqpoint{3.007045in}{1.925000in}}%
\pgfusepath{clip}%
\pgfsetbuttcap%
\pgfsetmiterjoin%
\definecolor{currentfill}{rgb}{0.406997,0.673741,0.834295}%
\pgfsetfillcolor{currentfill}%
\pgfsetlinewidth{0.000000pt}%
\definecolor{currentstroke}{rgb}{0.000000,0.000000,0.000000}%
\pgfsetstrokecolor{currentstroke}%
\pgfsetstrokeopacity{0.000000}%
\pgfsetdash{}{0pt}%
\pgfpathmoveto{\pgfqpoint{2.096418in}{1.118342in}}%
\pgfpathlineto{\pgfqpoint{2.095572in}{1.132813in}}%
\pgfpathlineto{\pgfqpoint{2.072910in}{1.131159in}}%
\pgfpathlineto{\pgfqpoint{2.073272in}{1.125381in}}%
\pgfpathlineto{\pgfqpoint{2.067309in}{1.125374in}}%
\pgfpathlineto{\pgfqpoint{2.065520in}{1.156638in}}%
\pgfpathlineto{\pgfqpoint{2.071143in}{1.157025in}}%
\pgfpathlineto{\pgfqpoint{2.105221in}{1.159153in}}%
\pgfpathlineto{\pgfqpoint{2.105428in}{1.156288in}}%
\pgfpathlineto{\pgfqpoint{2.142303in}{1.159619in}}%
\pgfpathlineto{\pgfqpoint{2.148381in}{1.158667in}}%
\pgfpathlineto{\pgfqpoint{2.151777in}{1.152326in}}%
\pgfpathlineto{\pgfqpoint{2.152006in}{1.147645in}}%
\pgfpathlineto{\pgfqpoint{2.155147in}{1.142366in}}%
\pgfpathlineto{\pgfqpoint{2.154997in}{1.139057in}}%
\pgfpathlineto{\pgfqpoint{2.135746in}{1.137366in}}%
\pgfpathlineto{\pgfqpoint{2.117823in}{1.135815in}}%
\pgfpathlineto{\pgfqpoint{2.118010in}{1.123238in}}%
\pgfpathlineto{\pgfqpoint{2.119277in}{1.119542in}}%
\pgfpathlineto{\pgfqpoint{2.096418in}{1.118342in}}%
\pgfpathclose%
\pgfusepath{fill}%
\end{pgfscope}%
\begin{pgfscope}%
\pgfpathrectangle{\pgfqpoint{0.100000in}{0.100000in}}{\pgfqpoint{3.007045in}{1.925000in}}%
\pgfusepath{clip}%
\pgfsetbuttcap%
\pgfsetmiterjoin%
\definecolor{currentfill}{rgb}{0.441569,0.694410,0.843952}%
\pgfsetfillcolor{currentfill}%
\pgfsetlinewidth{0.000000pt}%
\definecolor{currentstroke}{rgb}{0.000000,0.000000,0.000000}%
\pgfsetstrokecolor{currentstroke}%
\pgfsetstrokeopacity{0.000000}%
\pgfsetdash{}{0pt}%
\pgfpathmoveto{\pgfqpoint{2.547659in}{0.980532in}}%
\pgfpathlineto{\pgfqpoint{2.537459in}{0.985766in}}%
\pgfpathlineto{\pgfqpoint{2.533363in}{0.985590in}}%
\pgfpathlineto{\pgfqpoint{2.532109in}{0.998638in}}%
\pgfpathlineto{\pgfqpoint{2.522081in}{0.997367in}}%
\pgfpathlineto{\pgfqpoint{2.520017in}{1.015115in}}%
\pgfpathlineto{\pgfqpoint{2.514084in}{1.019211in}}%
\pgfpathlineto{\pgfqpoint{2.513186in}{1.021977in}}%
\pgfpathlineto{\pgfqpoint{2.515896in}{1.030173in}}%
\pgfpathlineto{\pgfqpoint{2.531083in}{1.032312in}}%
\pgfpathlineto{\pgfqpoint{2.547798in}{1.034084in}}%
\pgfpathlineto{\pgfqpoint{2.561819in}{1.036295in}}%
\pgfpathlineto{\pgfqpoint{2.566736in}{1.001545in}}%
\pgfpathlineto{\pgfqpoint{2.557574in}{1.001121in}}%
\pgfpathlineto{\pgfqpoint{2.549576in}{0.995557in}}%
\pgfpathlineto{\pgfqpoint{2.549802in}{0.987470in}}%
\pgfpathlineto{\pgfqpoint{2.545915in}{0.985956in}}%
\pgfpathlineto{\pgfqpoint{2.547659in}{0.980532in}}%
\pgfpathclose%
\pgfusepath{fill}%
\end{pgfscope}%
\begin{pgfscope}%
\pgfpathrectangle{\pgfqpoint{0.100000in}{0.100000in}}{\pgfqpoint{3.007045in}{1.925000in}}%
\pgfusepath{clip}%
\pgfsetbuttcap%
\pgfsetmiterjoin%
\definecolor{currentfill}{rgb}{0.366644,0.646182,0.818547}%
\pgfsetfillcolor{currentfill}%
\pgfsetlinewidth{0.000000pt}%
\definecolor{currentstroke}{rgb}{0.000000,0.000000,0.000000}%
\pgfsetstrokecolor{currentstroke}%
\pgfsetstrokeopacity{0.000000}%
\pgfsetdash{}{0pt}%
\pgfpathmoveto{\pgfqpoint{2.815570in}{1.283815in}}%
\pgfpathlineto{\pgfqpoint{2.811323in}{1.284246in}}%
\pgfpathlineto{\pgfqpoint{2.806809in}{1.301783in}}%
\pgfpathlineto{\pgfqpoint{2.791685in}{1.322963in}}%
\pgfpathlineto{\pgfqpoint{2.789776in}{1.325534in}}%
\pgfpathlineto{\pgfqpoint{2.799417in}{1.334517in}}%
\pgfpathlineto{\pgfqpoint{2.804907in}{1.347247in}}%
\pgfpathlineto{\pgfqpoint{2.809088in}{1.348480in}}%
\pgfpathlineto{\pgfqpoint{2.816490in}{1.347154in}}%
\pgfpathlineto{\pgfqpoint{2.819196in}{1.340247in}}%
\pgfpathlineto{\pgfqpoint{2.819265in}{1.328442in}}%
\pgfpathlineto{\pgfqpoint{2.821303in}{1.305330in}}%
\pgfpathlineto{\pgfqpoint{2.820085in}{1.298464in}}%
\pgfpathlineto{\pgfqpoint{2.815570in}{1.283815in}}%
\pgfpathclose%
\pgfusepath{fill}%
\end{pgfscope}%
\begin{pgfscope}%
\pgfpathrectangle{\pgfqpoint{0.100000in}{0.100000in}}{\pgfqpoint{3.007045in}{1.925000in}}%
\pgfusepath{clip}%
\pgfsetbuttcap%
\pgfsetmiterjoin%
\definecolor{currentfill}{rgb}{0.270804,0.580730,0.781146}%
\pgfsetfillcolor{currentfill}%
\pgfsetlinewidth{0.000000pt}%
\definecolor{currentstroke}{rgb}{0.000000,0.000000,0.000000}%
\pgfsetstrokecolor{currentstroke}%
\pgfsetstrokeopacity{0.000000}%
\pgfsetdash{}{0pt}%
\pgfpathmoveto{\pgfqpoint{1.633052in}{1.541386in}}%
\pgfpathlineto{\pgfqpoint{1.632109in}{1.506861in}}%
\pgfpathlineto{\pgfqpoint{1.597846in}{1.508282in}}%
\pgfpathlineto{\pgfqpoint{1.598571in}{1.525674in}}%
\pgfpathlineto{\pgfqpoint{1.557925in}{1.527548in}}%
\pgfpathlineto{\pgfqpoint{1.559045in}{1.550570in}}%
\pgfpathlineto{\pgfqpoint{1.552729in}{1.550914in}}%
\pgfpathlineto{\pgfqpoint{1.554716in}{1.596972in}}%
\pgfpathlineto{\pgfqpoint{1.587550in}{1.595331in}}%
\pgfpathlineto{\pgfqpoint{1.634891in}{1.593241in}}%
\pgfpathlineto{\pgfqpoint{1.633052in}{1.541386in}}%
\pgfpathclose%
\pgfusepath{fill}%
\end{pgfscope}%
\begin{pgfscope}%
\pgfpathrectangle{\pgfqpoint{0.100000in}{0.100000in}}{\pgfqpoint{3.007045in}{1.925000in}}%
\pgfusepath{clip}%
\pgfsetbuttcap%
\pgfsetmiterjoin%
\definecolor{currentfill}{rgb}{0.270804,0.580730,0.781146}%
\pgfsetfillcolor{currentfill}%
\pgfsetlinewidth{0.000000pt}%
\definecolor{currentstroke}{rgb}{0.000000,0.000000,0.000000}%
\pgfsetstrokecolor{currentstroke}%
\pgfsetstrokeopacity{0.000000}%
\pgfsetdash{}{0pt}%
\pgfpathmoveto{\pgfqpoint{1.827733in}{0.669764in}}%
\pgfpathlineto{\pgfqpoint{1.799830in}{0.669148in}}%
\pgfpathlineto{\pgfqpoint{1.803211in}{0.660558in}}%
\pgfpathlineto{\pgfqpoint{1.775875in}{0.660550in}}%
\pgfpathlineto{\pgfqpoint{1.773171in}{0.660592in}}%
\pgfpathlineto{\pgfqpoint{1.773020in}{0.706433in}}%
\pgfpathlineto{\pgfqpoint{1.763758in}{0.708620in}}%
\pgfpathlineto{\pgfqpoint{1.763868in}{0.730641in}}%
\pgfpathlineto{\pgfqpoint{1.787896in}{0.730793in}}%
\pgfpathlineto{\pgfqpoint{1.788689in}{0.729118in}}%
\pgfpathlineto{\pgfqpoint{1.825570in}{0.729490in}}%
\pgfpathlineto{\pgfqpoint{1.826203in}{0.670688in}}%
\pgfpathlineto{\pgfqpoint{1.827733in}{0.669764in}}%
\pgfpathclose%
\pgfusepath{fill}%
\end{pgfscope}%
\begin{pgfscope}%
\pgfpathrectangle{\pgfqpoint{0.100000in}{0.100000in}}{\pgfqpoint{3.007045in}{1.925000in}}%
\pgfusepath{clip}%
\pgfsetbuttcap%
\pgfsetmiterjoin%
\definecolor{currentfill}{rgb}{0.093272,0.396878,0.673664}%
\pgfsetfillcolor{currentfill}%
\pgfsetlinewidth{0.000000pt}%
\definecolor{currentstroke}{rgb}{0.000000,0.000000,0.000000}%
\pgfsetstrokecolor{currentstroke}%
\pgfsetstrokeopacity{0.000000}%
\pgfsetdash{}{0pt}%
\pgfpathmoveto{\pgfqpoint{0.877896in}{0.551221in}}%
\pgfpathlineto{\pgfqpoint{0.873853in}{0.546928in}}%
\pgfpathlineto{\pgfqpoint{0.875594in}{0.545533in}}%
\pgfpathlineto{\pgfqpoint{0.874170in}{0.543737in}}%
\pgfpathlineto{\pgfqpoint{0.876010in}{0.542324in}}%
\pgfpathlineto{\pgfqpoint{0.874603in}{0.540599in}}%
\pgfpathlineto{\pgfqpoint{0.876407in}{0.539190in}}%
\pgfpathlineto{\pgfqpoint{0.874516in}{0.537747in}}%
\pgfpathlineto{\pgfqpoint{0.870380in}{0.532422in}}%
\pgfpathlineto{\pgfqpoint{0.873947in}{0.529693in}}%
\pgfpathlineto{\pgfqpoint{0.872574in}{0.527836in}}%
\pgfpathlineto{\pgfqpoint{0.873845in}{0.526915in}}%
\pgfpathlineto{\pgfqpoint{0.867902in}{0.519376in}}%
\pgfpathlineto{\pgfqpoint{0.865978in}{0.517880in}}%
\pgfpathlineto{\pgfqpoint{0.860718in}{0.522144in}}%
\pgfpathlineto{\pgfqpoint{0.857919in}{0.518605in}}%
\pgfpathlineto{\pgfqpoint{0.852921in}{0.522581in}}%
\pgfpathlineto{\pgfqpoint{0.848387in}{0.517598in}}%
\pgfpathlineto{\pgfqpoint{0.845093in}{0.520174in}}%
\pgfpathlineto{\pgfqpoint{0.846505in}{0.521977in}}%
\pgfpathlineto{\pgfqpoint{0.844770in}{0.523398in}}%
\pgfpathlineto{\pgfqpoint{0.846127in}{0.525116in}}%
\pgfpathlineto{\pgfqpoint{0.842673in}{0.527910in}}%
\pgfpathlineto{\pgfqpoint{0.839868in}{0.524417in}}%
\pgfpathlineto{\pgfqpoint{0.836381in}{0.527215in}}%
\pgfpathlineto{\pgfqpoint{0.839219in}{0.530767in}}%
\pgfpathlineto{\pgfqpoint{0.832267in}{0.536448in}}%
\pgfpathlineto{\pgfqpoint{0.826641in}{0.529420in}}%
\pgfpathlineto{\pgfqpoint{0.818104in}{0.536633in}}%
\pgfpathlineto{\pgfqpoint{0.815410in}{0.533363in}}%
\pgfpathlineto{\pgfqpoint{0.810396in}{0.537663in}}%
\pgfpathlineto{\pgfqpoint{0.803129in}{0.529197in}}%
\pgfpathlineto{\pgfqpoint{0.801466in}{0.530632in}}%
\pgfpathlineto{\pgfqpoint{0.799988in}{0.528978in}}%
\pgfpathlineto{\pgfqpoint{0.790016in}{0.537796in}}%
\pgfpathlineto{\pgfqpoint{0.775661in}{0.550783in}}%
\pgfpathlineto{\pgfqpoint{0.768863in}{0.557118in}}%
\pgfpathlineto{\pgfqpoint{0.771871in}{0.560318in}}%
\pgfpathlineto{\pgfqpoint{0.770248in}{0.561865in}}%
\pgfpathlineto{\pgfqpoint{0.773326in}{0.565108in}}%
\pgfpathlineto{\pgfqpoint{0.771216in}{0.567098in}}%
\pgfpathlineto{\pgfqpoint{0.777383in}{0.573526in}}%
\pgfpathlineto{\pgfqpoint{0.776856in}{0.574080in}}%
\pgfpathlineto{\pgfqpoint{0.782998in}{0.580535in}}%
\pgfpathlineto{\pgfqpoint{0.784268in}{0.583082in}}%
\pgfpathlineto{\pgfqpoint{0.786955in}{0.580900in}}%
\pgfpathlineto{\pgfqpoint{0.788384in}{0.577765in}}%
\pgfpathlineto{\pgfqpoint{0.786641in}{0.576148in}}%
\pgfpathlineto{\pgfqpoint{0.784316in}{0.575077in}}%
\pgfpathlineto{\pgfqpoint{0.782913in}{0.573171in}}%
\pgfpathlineto{\pgfqpoint{0.780593in}{0.573219in}}%
\pgfpathlineto{\pgfqpoint{0.779806in}{0.568879in}}%
\pgfpathlineto{\pgfqpoint{0.781934in}{0.567331in}}%
\pgfpathlineto{\pgfqpoint{0.785364in}{0.563844in}}%
\pgfpathlineto{\pgfqpoint{0.788254in}{0.562327in}}%
\pgfpathlineto{\pgfqpoint{0.787949in}{0.560680in}}%
\pgfpathlineto{\pgfqpoint{0.789850in}{0.558307in}}%
\pgfpathlineto{\pgfqpoint{0.792509in}{0.557623in}}%
\pgfpathlineto{\pgfqpoint{0.793630in}{0.555844in}}%
\pgfpathlineto{\pgfqpoint{0.792534in}{0.553561in}}%
\pgfpathlineto{\pgfqpoint{0.797356in}{0.555466in}}%
\pgfpathlineto{\pgfqpoint{0.800166in}{0.557163in}}%
\pgfpathlineto{\pgfqpoint{0.801439in}{0.554540in}}%
\pgfpathlineto{\pgfqpoint{0.804446in}{0.553565in}}%
\pgfpathlineto{\pgfqpoint{0.804535in}{0.556258in}}%
\pgfpathlineto{\pgfqpoint{0.802504in}{0.559954in}}%
\pgfpathlineto{\pgfqpoint{0.799473in}{0.561628in}}%
\pgfpathlineto{\pgfqpoint{0.801889in}{0.563583in}}%
\pgfpathlineto{\pgfqpoint{0.801986in}{0.567136in}}%
\pgfpathlineto{\pgfqpoint{0.802817in}{0.569907in}}%
\pgfpathlineto{\pgfqpoint{0.801188in}{0.572131in}}%
\pgfpathlineto{\pgfqpoint{0.802671in}{0.576011in}}%
\pgfpathlineto{\pgfqpoint{0.804683in}{0.576334in}}%
\pgfpathlineto{\pgfqpoint{0.806090in}{0.575510in}}%
\pgfpathlineto{\pgfqpoint{0.805398in}{0.570181in}}%
\pgfpathlineto{\pgfqpoint{0.803112in}{0.567234in}}%
\pgfpathlineto{\pgfqpoint{0.803364in}{0.561395in}}%
\pgfpathlineto{\pgfqpoint{0.807010in}{0.559720in}}%
\pgfpathlineto{\pgfqpoint{0.807723in}{0.558079in}}%
\pgfpathlineto{\pgfqpoint{0.807544in}{0.554762in}}%
\pgfpathlineto{\pgfqpoint{0.812317in}{0.550291in}}%
\pgfpathlineto{\pgfqpoint{0.814525in}{0.552070in}}%
\pgfpathlineto{\pgfqpoint{0.815449in}{0.555639in}}%
\pgfpathlineto{\pgfqpoint{0.812956in}{0.557139in}}%
\pgfpathlineto{\pgfqpoint{0.812165in}{0.560243in}}%
\pgfpathlineto{\pgfqpoint{0.809060in}{0.561120in}}%
\pgfpathlineto{\pgfqpoint{0.807259in}{0.560638in}}%
\pgfpathlineto{\pgfqpoint{0.805657in}{0.562677in}}%
\pgfpathlineto{\pgfqpoint{0.805494in}{0.568104in}}%
\pgfpathlineto{\pgfqpoint{0.807400in}{0.569549in}}%
\pgfpathlineto{\pgfqpoint{0.811973in}{0.571987in}}%
\pgfpathlineto{\pgfqpoint{0.811391in}{0.574218in}}%
\pgfpathlineto{\pgfqpoint{0.806259in}{0.576839in}}%
\pgfpathlineto{\pgfqpoint{0.807373in}{0.578266in}}%
\pgfpathlineto{\pgfqpoint{0.803872in}{0.581047in}}%
\pgfpathlineto{\pgfqpoint{0.801672in}{0.583533in}}%
\pgfpathlineto{\pgfqpoint{0.799005in}{0.588521in}}%
\pgfpathlineto{\pgfqpoint{0.802634in}{0.593263in}}%
\pgfpathlineto{\pgfqpoint{0.804432in}{0.598840in}}%
\pgfpathlineto{\pgfqpoint{0.804689in}{0.601267in}}%
\pgfpathlineto{\pgfqpoint{0.803744in}{0.605764in}}%
\pgfpathlineto{\pgfqpoint{0.803822in}{0.608129in}}%
\pgfpathlineto{\pgfqpoint{0.803267in}{0.616806in}}%
\pgfpathlineto{\pgfqpoint{0.809438in}{0.610810in}}%
\pgfpathlineto{\pgfqpoint{0.813031in}{0.614554in}}%
\pgfpathlineto{\pgfqpoint{0.825701in}{0.602628in}}%
\pgfpathlineto{\pgfqpoint{0.826925in}{0.603954in}}%
\pgfpathlineto{\pgfqpoint{0.832109in}{0.599208in}}%
\pgfpathlineto{\pgfqpoint{0.830895in}{0.597873in}}%
\pgfpathlineto{\pgfqpoint{0.834401in}{0.594707in}}%
\pgfpathlineto{\pgfqpoint{0.835607in}{0.596048in}}%
\pgfpathlineto{\pgfqpoint{0.838673in}{0.593306in}}%
\pgfpathlineto{\pgfqpoint{0.837473in}{0.591959in}}%
\pgfpathlineto{\pgfqpoint{0.844619in}{0.585660in}}%
\pgfpathlineto{\pgfqpoint{0.845795in}{0.587009in}}%
\pgfpathlineto{\pgfqpoint{0.848948in}{0.584275in}}%
\pgfpathlineto{\pgfqpoint{0.847780in}{0.582922in}}%
\pgfpathlineto{\pgfqpoint{0.859713in}{0.572782in}}%
\pgfpathlineto{\pgfqpoint{0.867050in}{0.566744in}}%
\pgfpathlineto{\pgfqpoint{0.865552in}{0.564907in}}%
\pgfpathlineto{\pgfqpoint{0.872509in}{0.559210in}}%
\pgfpathlineto{\pgfqpoint{0.871135in}{0.557489in}}%
\pgfpathlineto{\pgfqpoint{0.878340in}{0.551804in}}%
\pgfpathlineto{\pgfqpoint{0.877896in}{0.551221in}}%
\pgfpathclose%
\pgfusepath{fill}%
\end{pgfscope}%
\begin{pgfscope}%
\pgfpathrectangle{\pgfqpoint{0.100000in}{0.100000in}}{\pgfqpoint{3.007045in}{1.925000in}}%
\pgfusepath{clip}%
\pgfsetbuttcap%
\pgfsetmiterjoin%
\definecolor{currentfill}{rgb}{0.031373,0.285675,0.564291}%
\pgfsetfillcolor{currentfill}%
\pgfsetlinewidth{0.000000pt}%
\definecolor{currentstroke}{rgb}{0.000000,0.000000,0.000000}%
\pgfsetstrokecolor{currentstroke}%
\pgfsetstrokeopacity{0.000000}%
\pgfsetdash{}{0pt}%
\pgfpathmoveto{\pgfqpoint{1.517880in}{1.481436in}}%
\pgfpathlineto{\pgfqpoint{1.516674in}{1.460211in}}%
\pgfpathlineto{\pgfqpoint{1.517464in}{1.451554in}}%
\pgfpathlineto{\pgfqpoint{1.511045in}{1.452183in}}%
\pgfpathlineto{\pgfqpoint{1.506479in}{1.454584in}}%
\pgfpathlineto{\pgfqpoint{1.496174in}{1.453745in}}%
\pgfpathlineto{\pgfqpoint{1.490267in}{1.460196in}}%
\pgfpathlineto{\pgfqpoint{1.483456in}{1.461953in}}%
\pgfpathlineto{\pgfqpoint{1.485701in}{1.483457in}}%
\pgfpathlineto{\pgfqpoint{1.517880in}{1.481436in}}%
\pgfpathclose%
\pgfusepath{fill}%
\end{pgfscope}%
\begin{pgfscope}%
\pgfpathrectangle{\pgfqpoint{0.100000in}{0.100000in}}{\pgfqpoint{3.007045in}{1.925000in}}%
\pgfusepath{clip}%
\pgfsetbuttcap%
\pgfsetmiterjoin%
\definecolor{currentfill}{rgb}{0.412042,0.677186,0.836263}%
\pgfsetfillcolor{currentfill}%
\pgfsetlinewidth{0.000000pt}%
\definecolor{currentstroke}{rgb}{0.000000,0.000000,0.000000}%
\pgfsetstrokecolor{currentstroke}%
\pgfsetstrokeopacity{0.000000}%
\pgfsetdash{}{0pt}%
\pgfpathmoveto{\pgfqpoint{2.006146in}{1.322530in}}%
\pgfpathlineto{\pgfqpoint{2.000310in}{1.318595in}}%
\pgfpathlineto{\pgfqpoint{1.996691in}{1.320570in}}%
\pgfpathlineto{\pgfqpoint{1.973965in}{1.320197in}}%
\pgfpathlineto{\pgfqpoint{1.973436in}{1.331705in}}%
\pgfpathlineto{\pgfqpoint{1.972422in}{1.354713in}}%
\pgfpathlineto{\pgfqpoint{1.983756in}{1.355205in}}%
\pgfpathlineto{\pgfqpoint{1.983506in}{1.360972in}}%
\pgfpathlineto{\pgfqpoint{1.992817in}{1.361377in}}%
\pgfpathlineto{\pgfqpoint{1.995792in}{1.357804in}}%
\pgfpathlineto{\pgfqpoint{1.997258in}{1.351410in}}%
\pgfpathlineto{\pgfqpoint{2.001078in}{1.349385in}}%
\pgfpathlineto{\pgfqpoint{2.008917in}{1.344598in}}%
\pgfpathlineto{\pgfqpoint{2.010393in}{1.337303in}}%
\pgfpathlineto{\pgfqpoint{2.009028in}{1.324274in}}%
\pgfpathlineto{\pgfqpoint{2.006146in}{1.322530in}}%
\pgfpathclose%
\pgfusepath{fill}%
\end{pgfscope}%
\begin{pgfscope}%
\pgfpathrectangle{\pgfqpoint{0.100000in}{0.100000in}}{\pgfqpoint{3.007045in}{1.925000in}}%
\pgfusepath{clip}%
\pgfsetbuttcap%
\pgfsetmiterjoin%
\definecolor{currentfill}{rgb}{0.351511,0.635848,0.812641}%
\pgfsetfillcolor{currentfill}%
\pgfsetlinewidth{0.000000pt}%
\definecolor{currentstroke}{rgb}{0.000000,0.000000,0.000000}%
\pgfsetstrokecolor{currentstroke}%
\pgfsetstrokeopacity{0.000000}%
\pgfsetdash{}{0pt}%
\pgfpathmoveto{\pgfqpoint{1.561531in}{1.203886in}}%
\pgfpathlineto{\pgfqpoint{1.584352in}{1.202908in}}%
\pgfpathlineto{\pgfqpoint{1.595536in}{1.202468in}}%
\pgfpathlineto{\pgfqpoint{1.594392in}{1.173857in}}%
\pgfpathlineto{\pgfqpoint{1.567058in}{1.175010in}}%
\pgfpathlineto{\pgfqpoint{1.537445in}{1.176396in}}%
\pgfpathlineto{\pgfqpoint{1.539001in}{1.204970in}}%
\pgfpathlineto{\pgfqpoint{1.561531in}{1.203886in}}%
\pgfpathclose%
\pgfusepath{fill}%
\end{pgfscope}%
\begin{pgfscope}%
\pgfpathrectangle{\pgfqpoint{0.100000in}{0.100000in}}{\pgfqpoint{3.007045in}{1.925000in}}%
\pgfusepath{clip}%
\pgfsetbuttcap%
\pgfsetmiterjoin%
\definecolor{currentfill}{rgb}{0.516863,0.735748,0.860192}%
\pgfsetfillcolor{currentfill}%
\pgfsetlinewidth{0.000000pt}%
\definecolor{currentstroke}{rgb}{0.000000,0.000000,0.000000}%
\pgfsetstrokecolor{currentstroke}%
\pgfsetstrokeopacity{0.000000}%
\pgfsetdash{}{0pt}%
\pgfpathmoveto{\pgfqpoint{1.995501in}{1.164493in}}%
\pgfpathlineto{\pgfqpoint{1.997146in}{1.174029in}}%
\pgfpathlineto{\pgfqpoint{1.992899in}{1.184317in}}%
\pgfpathlineto{\pgfqpoint{1.996210in}{1.193778in}}%
\pgfpathlineto{\pgfqpoint{1.978770in}{1.193340in}}%
\pgfpathlineto{\pgfqpoint{1.977971in}{1.216245in}}%
\pgfpathlineto{\pgfqpoint{1.977807in}{1.222266in}}%
\pgfpathlineto{\pgfqpoint{2.000904in}{1.222841in}}%
\pgfpathlineto{\pgfqpoint{2.001133in}{1.217077in}}%
\pgfpathlineto{\pgfqpoint{2.013813in}{1.217404in}}%
\pgfpathlineto{\pgfqpoint{2.008276in}{1.212623in}}%
\pgfpathlineto{\pgfqpoint{2.009474in}{1.209217in}}%
\pgfpathlineto{\pgfqpoint{2.014487in}{1.208775in}}%
\pgfpathlineto{\pgfqpoint{2.024381in}{1.212989in}}%
\pgfpathlineto{\pgfqpoint{2.025601in}{1.191743in}}%
\pgfpathlineto{\pgfqpoint{2.029984in}{1.174625in}}%
\pgfpathlineto{\pgfqpoint{2.018445in}{1.173868in}}%
\pgfpathlineto{\pgfqpoint{2.019642in}{1.156899in}}%
\pgfpathlineto{\pgfqpoint{2.011513in}{1.150677in}}%
\pgfpathlineto{\pgfqpoint{2.002639in}{1.150292in}}%
\pgfpathlineto{\pgfqpoint{1.997266in}{1.151743in}}%
\pgfpathlineto{\pgfqpoint{1.995501in}{1.164493in}}%
\pgfpathclose%
\pgfusepath{fill}%
\end{pgfscope}%
\begin{pgfscope}%
\pgfpathrectangle{\pgfqpoint{0.100000in}{0.100000in}}{\pgfqpoint{3.007045in}{1.925000in}}%
\pgfusepath{clip}%
\pgfsetbuttcap%
\pgfsetmiterjoin%
\definecolor{currentfill}{rgb}{0.429020,0.687520,0.841246}%
\pgfsetfillcolor{currentfill}%
\pgfsetlinewidth{0.000000pt}%
\definecolor{currentstroke}{rgb}{0.000000,0.000000,0.000000}%
\pgfsetstrokecolor{currentstroke}%
\pgfsetstrokeopacity{0.000000}%
\pgfsetdash{}{0pt}%
\pgfpathmoveto{\pgfqpoint{2.664388in}{1.534234in}}%
\pgfpathlineto{\pgfqpoint{2.661128in}{1.541996in}}%
\pgfpathlineto{\pgfqpoint{2.664794in}{1.548806in}}%
\pgfpathlineto{\pgfqpoint{2.662584in}{1.552711in}}%
\pgfpathlineto{\pgfqpoint{2.655181in}{1.554501in}}%
\pgfpathlineto{\pgfqpoint{2.651982in}{1.562017in}}%
\pgfpathlineto{\pgfqpoint{2.659262in}{1.569668in}}%
\pgfpathlineto{\pgfqpoint{2.658006in}{1.573468in}}%
\pgfpathlineto{\pgfqpoint{2.668332in}{1.581558in}}%
\pgfpathlineto{\pgfqpoint{2.671705in}{1.586619in}}%
\pgfpathlineto{\pgfqpoint{2.673065in}{1.592568in}}%
\pgfpathlineto{\pgfqpoint{2.682395in}{1.607388in}}%
\pgfpathlineto{\pgfqpoint{2.696457in}{1.624732in}}%
\pgfpathlineto{\pgfqpoint{2.702880in}{1.630799in}}%
\pgfpathlineto{\pgfqpoint{2.709073in}{1.634394in}}%
\pgfpathlineto{\pgfqpoint{2.713963in}{1.634075in}}%
\pgfpathlineto{\pgfqpoint{2.718437in}{1.632224in}}%
\pgfpathlineto{\pgfqpoint{2.735305in}{1.582994in}}%
\pgfpathlineto{\pgfqpoint{2.735921in}{1.578293in}}%
\pgfpathlineto{\pgfqpoint{2.721573in}{1.573055in}}%
\pgfpathlineto{\pgfqpoint{2.712187in}{1.569616in}}%
\pgfpathlineto{\pgfqpoint{2.706539in}{1.571514in}}%
\pgfpathlineto{\pgfqpoint{2.716233in}{1.541048in}}%
\pgfpathlineto{\pgfqpoint{2.699199in}{1.524078in}}%
\pgfpathlineto{\pgfqpoint{2.688054in}{1.525089in}}%
\pgfpathlineto{\pgfqpoint{2.684161in}{1.539019in}}%
\pgfpathlineto{\pgfqpoint{2.672204in}{1.537694in}}%
\pgfpathlineto{\pgfqpoint{2.664388in}{1.534234in}}%
\pgfpathclose%
\pgfusepath{fill}%
\end{pgfscope}%
\begin{pgfscope}%
\pgfpathrectangle{\pgfqpoint{0.100000in}{0.100000in}}{\pgfqpoint{3.007045in}{1.925000in}}%
\pgfusepath{clip}%
\pgfsetbuttcap%
\pgfsetmiterjoin%
\definecolor{currentfill}{rgb}{0.361599,0.642737,0.816578}%
\pgfsetfillcolor{currentfill}%
\pgfsetlinewidth{0.000000pt}%
\definecolor{currentstroke}{rgb}{0.000000,0.000000,0.000000}%
\pgfsetstrokecolor{currentstroke}%
\pgfsetstrokeopacity{0.000000}%
\pgfsetdash{}{0pt}%
\pgfpathmoveto{\pgfqpoint{1.907064in}{1.059811in}}%
\pgfpathlineto{\pgfqpoint{1.900128in}{1.055674in}}%
\pgfpathlineto{\pgfqpoint{1.896036in}{1.061694in}}%
\pgfpathlineto{\pgfqpoint{1.872518in}{1.061815in}}%
\pgfpathlineto{\pgfqpoint{1.872670in}{1.072352in}}%
\pgfpathlineto{\pgfqpoint{1.871618in}{1.080502in}}%
\pgfpathlineto{\pgfqpoint{1.872048in}{1.113887in}}%
\pgfpathlineto{\pgfqpoint{1.883385in}{1.113543in}}%
\pgfpathlineto{\pgfqpoint{1.883415in}{1.107800in}}%
\pgfpathlineto{\pgfqpoint{1.894852in}{1.097146in}}%
\pgfpathlineto{\pgfqpoint{1.906297in}{1.097100in}}%
\pgfpathlineto{\pgfqpoint{1.906201in}{1.091563in}}%
\pgfpathlineto{\pgfqpoint{1.911673in}{1.091410in}}%
\pgfpathlineto{\pgfqpoint{1.917287in}{1.088655in}}%
\pgfpathlineto{\pgfqpoint{1.918382in}{1.070506in}}%
\pgfpathlineto{\pgfqpoint{1.906913in}{1.070484in}}%
\pgfpathlineto{\pgfqpoint{1.907064in}{1.059811in}}%
\pgfpathclose%
\pgfusepath{fill}%
\end{pgfscope}%
\begin{pgfscope}%
\pgfpathrectangle{\pgfqpoint{0.100000in}{0.100000in}}{\pgfqpoint{3.007045in}{1.925000in}}%
\pgfusepath{clip}%
\pgfsetbuttcap%
\pgfsetmiterjoin%
\definecolor{currentfill}{rgb}{0.296025,0.597955,0.790988}%
\pgfsetfillcolor{currentfill}%
\pgfsetlinewidth{0.000000pt}%
\definecolor{currentstroke}{rgb}{0.000000,0.000000,0.000000}%
\pgfsetstrokecolor{currentstroke}%
\pgfsetstrokeopacity{0.000000}%
\pgfsetdash{}{0pt}%
\pgfpathmoveto{\pgfqpoint{2.045514in}{0.912506in}}%
\pgfpathlineto{\pgfqpoint{2.038540in}{0.912838in}}%
\pgfpathlineto{\pgfqpoint{2.039932in}{0.907365in}}%
\pgfpathlineto{\pgfqpoint{2.033514in}{0.905154in}}%
\pgfpathlineto{\pgfqpoint{2.025623in}{0.904867in}}%
\pgfpathlineto{\pgfqpoint{2.023624in}{0.941619in}}%
\pgfpathlineto{\pgfqpoint{2.021688in}{0.947698in}}%
\pgfpathlineto{\pgfqpoint{2.026615in}{0.954215in}}%
\pgfpathlineto{\pgfqpoint{2.033532in}{0.960342in}}%
\pgfpathlineto{\pgfqpoint{2.034251in}{0.967797in}}%
\pgfpathlineto{\pgfqpoint{2.025516in}{0.974832in}}%
\pgfpathlineto{\pgfqpoint{2.028875in}{0.983778in}}%
\pgfpathlineto{\pgfqpoint{2.038888in}{0.984228in}}%
\pgfpathlineto{\pgfqpoint{2.039044in}{0.968405in}}%
\pgfpathlineto{\pgfqpoint{2.050739in}{0.968865in}}%
\pgfpathlineto{\pgfqpoint{2.059156in}{0.970494in}}%
\pgfpathlineto{\pgfqpoint{2.061897in}{0.966247in}}%
\pgfpathlineto{\pgfqpoint{2.062788in}{0.960749in}}%
\pgfpathlineto{\pgfqpoint{2.054510in}{0.958420in}}%
\pgfpathlineto{\pgfqpoint{2.060132in}{0.952292in}}%
\pgfpathlineto{\pgfqpoint{2.055729in}{0.949079in}}%
\pgfpathlineto{\pgfqpoint{2.054350in}{0.941105in}}%
\pgfpathlineto{\pgfqpoint{2.051897in}{0.936223in}}%
\pgfpathlineto{\pgfqpoint{2.052220in}{0.930837in}}%
\pgfpathlineto{\pgfqpoint{2.049440in}{0.927189in}}%
\pgfpathlineto{\pgfqpoint{2.042608in}{0.925600in}}%
\pgfpathlineto{\pgfqpoint{2.043371in}{0.916596in}}%
\pgfpathlineto{\pgfqpoint{2.045514in}{0.912506in}}%
\pgfpathclose%
\pgfusepath{fill}%
\end{pgfscope}%
\begin{pgfscope}%
\pgfpathrectangle{\pgfqpoint{0.100000in}{0.100000in}}{\pgfqpoint{3.007045in}{1.925000in}}%
\pgfusepath{clip}%
\pgfsetbuttcap%
\pgfsetmiterjoin%
\definecolor{currentfill}{rgb}{0.376732,0.653072,0.822484}%
\pgfsetfillcolor{currentfill}%
\pgfsetlinewidth{0.000000pt}%
\definecolor{currentstroke}{rgb}{0.000000,0.000000,0.000000}%
\pgfsetstrokecolor{currentstroke}%
\pgfsetstrokeopacity{0.000000}%
\pgfsetdash{}{0pt}%
\pgfpathmoveto{\pgfqpoint{2.512057in}{0.734800in}}%
\pgfpathlineto{\pgfqpoint{2.509344in}{0.740112in}}%
\pgfpathlineto{\pgfqpoint{2.485433in}{0.742115in}}%
\pgfpathlineto{\pgfqpoint{2.478955in}{0.739426in}}%
\pgfpathlineto{\pgfqpoint{2.475118in}{0.743801in}}%
\pgfpathlineto{\pgfqpoint{2.476049in}{0.748944in}}%
\pgfpathlineto{\pgfqpoint{2.469526in}{0.751792in}}%
\pgfpathlineto{\pgfqpoint{2.460691in}{0.751658in}}%
\pgfpathlineto{\pgfqpoint{2.463033in}{0.756433in}}%
\pgfpathlineto{\pgfqpoint{2.469397in}{0.761860in}}%
\pgfpathlineto{\pgfqpoint{2.467299in}{0.764175in}}%
\pgfpathlineto{\pgfqpoint{2.470222in}{0.769671in}}%
\pgfpathlineto{\pgfqpoint{2.475814in}{0.775674in}}%
\pgfpathlineto{\pgfqpoint{2.482589in}{0.774200in}}%
\pgfpathlineto{\pgfqpoint{2.485475in}{0.775665in}}%
\pgfpathlineto{\pgfqpoint{2.488225in}{0.782556in}}%
\pgfpathlineto{\pgfqpoint{2.500186in}{0.786765in}}%
\pgfpathlineto{\pgfqpoint{2.505569in}{0.784566in}}%
\pgfpathlineto{\pgfqpoint{2.516660in}{0.795332in}}%
\pgfpathlineto{\pgfqpoint{2.519777in}{0.793317in}}%
\pgfpathlineto{\pgfqpoint{2.519904in}{0.788430in}}%
\pgfpathlineto{\pgfqpoint{2.525012in}{0.782861in}}%
\pgfpathlineto{\pgfqpoint{2.526324in}{0.773753in}}%
\pgfpathlineto{\pgfqpoint{2.529482in}{0.767093in}}%
\pgfpathlineto{\pgfqpoint{2.522083in}{0.760192in}}%
\pgfpathlineto{\pgfqpoint{2.523508in}{0.752566in}}%
\pgfpathlineto{\pgfqpoint{2.528796in}{0.747623in}}%
\pgfpathlineto{\pgfqpoint{2.530406in}{0.743490in}}%
\pgfpathlineto{\pgfqpoint{2.512057in}{0.734800in}}%
\pgfpathclose%
\pgfusepath{fill}%
\end{pgfscope}%
\begin{pgfscope}%
\pgfpathrectangle{\pgfqpoint{0.100000in}{0.100000in}}{\pgfqpoint{3.007045in}{1.925000in}}%
\pgfusepath{clip}%
\pgfsetbuttcap%
\pgfsetmiterjoin%
\definecolor{currentfill}{rgb}{0.187266,0.500992,0.739608}%
\pgfsetfillcolor{currentfill}%
\pgfsetlinewidth{0.000000pt}%
\definecolor{currentstroke}{rgb}{0.000000,0.000000,0.000000}%
\pgfsetstrokecolor{currentstroke}%
\pgfsetstrokeopacity{0.000000}%
\pgfsetdash{}{0pt}%
\pgfpathmoveto{\pgfqpoint{1.215312in}{0.701261in}}%
\pgfpathlineto{\pgfqpoint{1.212900in}{0.676457in}}%
\pgfpathlineto{\pgfqpoint{1.207983in}{0.628535in}}%
\pgfpathlineto{\pgfqpoint{1.205698in}{0.612940in}}%
\pgfpathlineto{\pgfqpoint{1.201869in}{0.611048in}}%
\pgfpathlineto{\pgfqpoint{1.194762in}{0.619908in}}%
\pgfpathlineto{\pgfqpoint{1.189798in}{0.624237in}}%
\pgfpathlineto{\pgfqpoint{1.180001in}{0.628883in}}%
\pgfpathlineto{\pgfqpoint{1.179940in}{0.630968in}}%
\pgfpathlineto{\pgfqpoint{1.171824in}{0.638589in}}%
\pgfpathlineto{\pgfqpoint{1.169910in}{0.645050in}}%
\pgfpathlineto{\pgfqpoint{1.160977in}{0.651771in}}%
\pgfpathlineto{\pgfqpoint{1.152306in}{0.666060in}}%
\pgfpathlineto{\pgfqpoint{1.145301in}{0.669338in}}%
\pgfpathlineto{\pgfqpoint{1.138853in}{0.674660in}}%
\pgfpathlineto{\pgfqpoint{1.131339in}{0.693210in}}%
\pgfpathlineto{\pgfqpoint{1.123511in}{0.697668in}}%
\pgfpathlineto{\pgfqpoint{1.080529in}{0.703340in}}%
\pgfpathlineto{\pgfqpoint{1.087735in}{0.757358in}}%
\pgfpathlineto{\pgfqpoint{1.089302in}{0.768831in}}%
\pgfpathlineto{\pgfqpoint{1.112021in}{0.765725in}}%
\pgfpathlineto{\pgfqpoint{1.112390in}{0.768586in}}%
\pgfpathlineto{\pgfqpoint{1.144512in}{0.779815in}}%
\pgfpathlineto{\pgfqpoint{1.141737in}{0.773832in}}%
\pgfpathlineto{\pgfqpoint{1.133803in}{0.710900in}}%
\pgfpathlineto{\pgfqpoint{1.186737in}{0.704479in}}%
\pgfpathlineto{\pgfqpoint{1.215312in}{0.701261in}}%
\pgfpathclose%
\pgfusepath{fill}%
\end{pgfscope}%
\begin{pgfscope}%
\pgfpathrectangle{\pgfqpoint{0.100000in}{0.100000in}}{\pgfqpoint{3.007045in}{1.925000in}}%
\pgfusepath{clip}%
\pgfsetbuttcap%
\pgfsetmiterjoin%
\definecolor{currentfill}{rgb}{0.256286,0.570012,0.775163}%
\pgfsetfillcolor{currentfill}%
\pgfsetlinewidth{0.000000pt}%
\definecolor{currentstroke}{rgb}{0.000000,0.000000,0.000000}%
\pgfsetstrokecolor{currentstroke}%
\pgfsetstrokeopacity{0.000000}%
\pgfsetdash{}{0pt}%
\pgfpathmoveto{\pgfqpoint{2.321001in}{0.750113in}}%
\pgfpathlineto{\pgfqpoint{2.317396in}{0.757985in}}%
\pgfpathlineto{\pgfqpoint{2.335233in}{0.759815in}}%
\pgfpathlineto{\pgfqpoint{2.333138in}{0.783244in}}%
\pgfpathlineto{\pgfqpoint{2.352426in}{0.785269in}}%
\pgfpathlineto{\pgfqpoint{2.352032in}{0.771233in}}%
\pgfpathlineto{\pgfqpoint{2.354579in}{0.763090in}}%
\pgfpathlineto{\pgfqpoint{2.361889in}{0.757434in}}%
\pgfpathlineto{\pgfqpoint{2.367977in}{0.755519in}}%
\pgfpathlineto{\pgfqpoint{2.360665in}{0.742287in}}%
\pgfpathlineto{\pgfqpoint{2.350162in}{0.739256in}}%
\pgfpathlineto{\pgfqpoint{2.349352in}{0.736500in}}%
\pgfpathlineto{\pgfqpoint{2.351237in}{0.719223in}}%
\pgfpathlineto{\pgfqpoint{2.336528in}{0.717430in}}%
\pgfpathlineto{\pgfqpoint{2.329526in}{0.710196in}}%
\pgfpathlineto{\pgfqpoint{2.330327in}{0.705545in}}%
\pgfpathlineto{\pgfqpoint{2.323124in}{0.704669in}}%
\pgfpathlineto{\pgfqpoint{2.318447in}{0.709939in}}%
\pgfpathlineto{\pgfqpoint{2.309942in}{0.708924in}}%
\pgfpathlineto{\pgfqpoint{2.308027in}{0.714591in}}%
\pgfpathlineto{\pgfqpoint{2.306778in}{0.726142in}}%
\pgfpathlineto{\pgfqpoint{2.312554in}{0.726822in}}%
\pgfpathlineto{\pgfqpoint{2.311941in}{0.730610in}}%
\pgfpathlineto{\pgfqpoint{2.327177in}{0.732545in}}%
\pgfpathlineto{\pgfqpoint{2.330141in}{0.735377in}}%
\pgfpathlineto{\pgfqpoint{2.323388in}{0.743508in}}%
\pgfpathlineto{\pgfqpoint{2.321001in}{0.750113in}}%
\pgfpathclose%
\pgfusepath{fill}%
\end{pgfscope}%
\begin{pgfscope}%
\pgfpathrectangle{\pgfqpoint{0.100000in}{0.100000in}}{\pgfqpoint{3.007045in}{1.925000in}}%
\pgfusepath{clip}%
\pgfsetbuttcap%
\pgfsetmiterjoin%
\definecolor{currentfill}{rgb}{0.361599,0.642737,0.816578}%
\pgfsetfillcolor{currentfill}%
\pgfsetlinewidth{0.000000pt}%
\definecolor{currentstroke}{rgb}{0.000000,0.000000,0.000000}%
\pgfsetstrokecolor{currentstroke}%
\pgfsetstrokeopacity{0.000000}%
\pgfsetdash{}{0pt}%
\pgfpathmoveto{\pgfqpoint{2.252087in}{0.743598in}}%
\pgfpathlineto{\pgfqpoint{2.250388in}{0.740313in}}%
\pgfpathlineto{\pgfqpoint{2.233682in}{0.738645in}}%
\pgfpathlineto{\pgfqpoint{2.233994in}{0.735748in}}%
\pgfpathlineto{\pgfqpoint{2.222669in}{0.734886in}}%
\pgfpathlineto{\pgfqpoint{2.217078in}{0.734332in}}%
\pgfpathlineto{\pgfqpoint{2.215985in}{0.745794in}}%
\pgfpathlineto{\pgfqpoint{2.223920in}{0.746448in}}%
\pgfpathlineto{\pgfqpoint{2.222615in}{0.760540in}}%
\pgfpathlineto{\pgfqpoint{2.217222in}{0.763282in}}%
\pgfpathlineto{\pgfqpoint{2.213720in}{0.767453in}}%
\pgfpathlineto{\pgfqpoint{2.213183in}{0.772767in}}%
\pgfpathlineto{\pgfqpoint{2.204331in}{0.776956in}}%
\pgfpathlineto{\pgfqpoint{2.197889in}{0.784340in}}%
\pgfpathlineto{\pgfqpoint{2.198442in}{0.789193in}}%
\pgfpathlineto{\pgfqpoint{2.202677in}{0.796254in}}%
\pgfpathlineto{\pgfqpoint{2.207356in}{0.797594in}}%
\pgfpathlineto{\pgfqpoint{2.209929in}{0.801633in}}%
\pgfpathlineto{\pgfqpoint{2.210327in}{0.807475in}}%
\pgfpathlineto{\pgfqpoint{2.213109in}{0.813406in}}%
\pgfpathlineto{\pgfqpoint{2.206348in}{0.816022in}}%
\pgfpathlineto{\pgfqpoint{2.204539in}{0.819210in}}%
\pgfpathlineto{\pgfqpoint{2.202564in}{0.841744in}}%
\pgfpathlineto{\pgfqpoint{2.202469in}{0.842698in}}%
\pgfpathlineto{\pgfqpoint{2.231240in}{0.844683in}}%
\pgfpathlineto{\pgfqpoint{2.236928in}{0.845086in}}%
\pgfpathlineto{\pgfqpoint{2.247672in}{0.832584in}}%
\pgfpathlineto{\pgfqpoint{2.246429in}{0.828970in}}%
\pgfpathlineto{\pgfqpoint{2.249130in}{0.825060in}}%
\pgfpathlineto{\pgfqpoint{2.254073in}{0.825883in}}%
\pgfpathlineto{\pgfqpoint{2.257720in}{0.820289in}}%
\pgfpathlineto{\pgfqpoint{2.262319in}{0.816941in}}%
\pgfpathlineto{\pgfqpoint{2.263640in}{0.812852in}}%
\pgfpathlineto{\pgfqpoint{2.258979in}{0.805773in}}%
\pgfpathlineto{\pgfqpoint{2.255427in}{0.799367in}}%
\pgfpathlineto{\pgfqpoint{2.257385in}{0.796619in}}%
\pgfpathlineto{\pgfqpoint{2.248013in}{0.785562in}}%
\pgfpathlineto{\pgfqpoint{2.249772in}{0.779385in}}%
\pgfpathlineto{\pgfqpoint{2.244576in}{0.775293in}}%
\pgfpathlineto{\pgfqpoint{2.242492in}{0.770671in}}%
\pgfpathlineto{\pgfqpoint{2.243973in}{0.753410in}}%
\pgfpathlineto{\pgfqpoint{2.247289in}{0.747109in}}%
\pgfpathlineto{\pgfqpoint{2.252087in}{0.743598in}}%
\pgfpathclose%
\pgfusepath{fill}%
\end{pgfscope}%
\begin{pgfscope}%
\pgfpathrectangle{\pgfqpoint{0.100000in}{0.100000in}}{\pgfqpoint{3.007045in}{1.925000in}}%
\pgfusepath{clip}%
\pgfsetbuttcap%
\pgfsetmiterjoin%
\definecolor{currentfill}{rgb}{0.472941,0.711634,0.850719}%
\pgfsetfillcolor{currentfill}%
\pgfsetlinewidth{0.000000pt}%
\definecolor{currentstroke}{rgb}{0.000000,0.000000,0.000000}%
\pgfsetstrokecolor{currentstroke}%
\pgfsetstrokeopacity{0.000000}%
\pgfsetdash{}{0pt}%
\pgfpathmoveto{\pgfqpoint{1.710630in}{1.795591in}}%
\pgfpathlineto{\pgfqpoint{1.710357in}{1.776679in}}%
\pgfpathlineto{\pgfqpoint{1.710937in}{1.765107in}}%
\pgfpathlineto{\pgfqpoint{1.676648in}{1.765752in}}%
\pgfpathlineto{\pgfqpoint{1.677958in}{1.760345in}}%
\pgfpathlineto{\pgfqpoint{1.676638in}{1.755236in}}%
\pgfpathlineto{\pgfqpoint{1.678294in}{1.748130in}}%
\pgfpathlineto{\pgfqpoint{1.677065in}{1.742344in}}%
\pgfpathlineto{\pgfqpoint{1.643166in}{1.743370in}}%
\pgfpathlineto{\pgfqpoint{1.625869in}{1.743967in}}%
\pgfpathlineto{\pgfqpoint{1.626307in}{1.755616in}}%
\pgfpathlineto{\pgfqpoint{1.625579in}{1.767286in}}%
\pgfpathlineto{\pgfqpoint{1.596746in}{1.768470in}}%
\pgfpathlineto{\pgfqpoint{1.595955in}{1.780137in}}%
\pgfpathlineto{\pgfqpoint{1.596789in}{1.798977in}}%
\pgfpathlineto{\pgfqpoint{1.638516in}{1.797389in}}%
\pgfpathlineto{\pgfqpoint{1.674481in}{1.796361in}}%
\pgfpathlineto{\pgfqpoint{1.710630in}{1.795591in}}%
\pgfpathclose%
\pgfusepath{fill}%
\end{pgfscope}%
\begin{pgfscope}%
\pgfpathrectangle{\pgfqpoint{0.100000in}{0.100000in}}{\pgfqpoint{3.007045in}{1.925000in}}%
\pgfusepath{clip}%
\pgfsetbuttcap%
\pgfsetmiterjoin%
\definecolor{currentfill}{rgb}{0.031373,0.196355,0.431665}%
\pgfsetfillcolor{currentfill}%
\pgfsetlinewidth{0.000000pt}%
\definecolor{currentstroke}{rgb}{0.000000,0.000000,0.000000}%
\pgfsetstrokecolor{currentstroke}%
\pgfsetstrokeopacity{0.000000}%
\pgfsetdash{}{0pt}%
\pgfpathmoveto{\pgfqpoint{1.521148in}{0.313814in}}%
\pgfpathlineto{\pgfqpoint{1.519248in}{0.324225in}}%
\pgfpathlineto{\pgfqpoint{1.516036in}{0.332359in}}%
\pgfpathlineto{\pgfqpoint{1.509870in}{0.338738in}}%
\pgfpathlineto{\pgfqpoint{1.505834in}{0.347871in}}%
\pgfpathlineto{\pgfqpoint{1.507664in}{0.355223in}}%
\pgfpathlineto{\pgfqpoint{1.503531in}{0.366333in}}%
\pgfpathlineto{\pgfqpoint{1.505470in}{0.376130in}}%
\pgfpathlineto{\pgfqpoint{1.502792in}{0.376996in}}%
\pgfpathlineto{\pgfqpoint{1.502216in}{0.384359in}}%
\pgfpathlineto{\pgfqpoint{1.491739in}{0.388953in}}%
\pgfpathlineto{\pgfqpoint{1.486685in}{0.396463in}}%
\pgfpathlineto{\pgfqpoint{1.482967in}{0.398437in}}%
\pgfpathlineto{\pgfqpoint{1.482372in}{0.405110in}}%
\pgfpathlineto{\pgfqpoint{1.476970in}{0.411788in}}%
\pgfpathlineto{\pgfqpoint{1.472047in}{0.422119in}}%
\pgfpathlineto{\pgfqpoint{1.464650in}{0.426083in}}%
\pgfpathlineto{\pgfqpoint{1.470461in}{0.425819in}}%
\pgfpathlineto{\pgfqpoint{1.512921in}{0.424052in}}%
\pgfpathlineto{\pgfqpoint{1.512666in}{0.412373in}}%
\pgfpathlineto{\pgfqpoint{1.547367in}{0.412624in}}%
\pgfpathlineto{\pgfqpoint{1.545690in}{0.365433in}}%
\pgfpathlineto{\pgfqpoint{1.560356in}{0.365181in}}%
\pgfpathlineto{\pgfqpoint{1.563594in}{0.356706in}}%
\pgfpathlineto{\pgfqpoint{1.564709in}{0.344532in}}%
\pgfpathlineto{\pgfqpoint{1.567623in}{0.344460in}}%
\pgfpathlineto{\pgfqpoint{1.566652in}{0.326194in}}%
\pgfpathlineto{\pgfqpoint{1.534731in}{0.327636in}}%
\pgfpathlineto{\pgfqpoint{1.530987in}{0.320341in}}%
\pgfpathlineto{\pgfqpoint{1.521148in}{0.313814in}}%
\pgfpathclose%
\pgfusepath{fill}%
\end{pgfscope}%
\begin{pgfscope}%
\pgfpathrectangle{\pgfqpoint{0.100000in}{0.100000in}}{\pgfqpoint{3.007045in}{1.925000in}}%
\pgfusepath{clip}%
\pgfsetbuttcap%
\pgfsetmiterjoin%
\definecolor{currentfill}{rgb}{0.429020,0.687520,0.841246}%
\pgfsetfillcolor{currentfill}%
\pgfsetlinewidth{0.000000pt}%
\definecolor{currentstroke}{rgb}{0.000000,0.000000,0.000000}%
\pgfsetstrokecolor{currentstroke}%
\pgfsetstrokeopacity{0.000000}%
\pgfsetdash{}{0pt}%
\pgfpathmoveto{\pgfqpoint{2.503300in}{0.656930in}}%
\pgfpathlineto{\pgfqpoint{2.498304in}{0.661191in}}%
\pgfpathlineto{\pgfqpoint{2.497961in}{0.669319in}}%
\pgfpathlineto{\pgfqpoint{2.503124in}{0.674785in}}%
\pgfpathlineto{\pgfqpoint{2.508359in}{0.677599in}}%
\pgfpathlineto{\pgfqpoint{2.505019in}{0.681519in}}%
\pgfpathlineto{\pgfqpoint{2.511887in}{0.680666in}}%
\pgfpathlineto{\pgfqpoint{2.518569in}{0.682840in}}%
\pgfpathlineto{\pgfqpoint{2.523188in}{0.681053in}}%
\pgfpathlineto{\pgfqpoint{2.528105in}{0.691406in}}%
\pgfpathlineto{\pgfqpoint{2.524353in}{0.695341in}}%
\pgfpathlineto{\pgfqpoint{2.529648in}{0.698799in}}%
\pgfpathlineto{\pgfqpoint{2.532623in}{0.707455in}}%
\pgfpathlineto{\pgfqpoint{2.536542in}{0.704088in}}%
\pgfpathlineto{\pgfqpoint{2.545964in}{0.706161in}}%
\pgfpathlineto{\pgfqpoint{2.550608in}{0.701257in}}%
\pgfpathlineto{\pgfqpoint{2.551853in}{0.698652in}}%
\pgfpathlineto{\pgfqpoint{2.548410in}{0.688938in}}%
\pgfpathlineto{\pgfqpoint{2.550045in}{0.680517in}}%
\pgfpathlineto{\pgfqpoint{2.543234in}{0.669808in}}%
\pgfpathlineto{\pgfqpoint{2.542300in}{0.662555in}}%
\pgfpathlineto{\pgfqpoint{2.533990in}{0.667284in}}%
\pgfpathlineto{\pgfqpoint{2.521337in}{0.669031in}}%
\pgfpathlineto{\pgfqpoint{2.514003in}{0.660961in}}%
\pgfpathlineto{\pgfqpoint{2.506695in}{0.661778in}}%
\pgfpathlineto{\pgfqpoint{2.503300in}{0.656930in}}%
\pgfpathclose%
\pgfusepath{fill}%
\end{pgfscope}%
\begin{pgfscope}%
\pgfpathrectangle{\pgfqpoint{0.100000in}{0.100000in}}{\pgfqpoint{3.007045in}{1.925000in}}%
\pgfusepath{clip}%
\pgfsetbuttcap%
\pgfsetmiterjoin%
\definecolor{currentfill}{rgb}{0.376732,0.653072,0.822484}%
\pgfsetfillcolor{currentfill}%
\pgfsetlinewidth{0.000000pt}%
\definecolor{currentstroke}{rgb}{0.000000,0.000000,0.000000}%
\pgfsetstrokecolor{currentstroke}%
\pgfsetstrokeopacity{0.000000}%
\pgfsetdash{}{0pt}%
\pgfpathmoveto{\pgfqpoint{1.326529in}{1.324213in}}%
\pgfpathlineto{\pgfqpoint{1.333387in}{1.392952in}}%
\pgfpathlineto{\pgfqpoint{1.335923in}{1.418503in}}%
\pgfpathlineto{\pgfqpoint{1.386507in}{1.413753in}}%
\pgfpathlineto{\pgfqpoint{1.396545in}{1.412828in}}%
\pgfpathlineto{\pgfqpoint{1.396236in}{1.398628in}}%
\pgfpathlineto{\pgfqpoint{1.394290in}{1.375922in}}%
\pgfpathlineto{\pgfqpoint{1.395452in}{1.375827in}}%
\pgfpathlineto{\pgfqpoint{1.393848in}{1.352944in}}%
\pgfpathlineto{\pgfqpoint{1.396518in}{1.346955in}}%
\pgfpathlineto{\pgfqpoint{1.395435in}{1.334971in}}%
\pgfpathlineto{\pgfqpoint{1.396817in}{1.327688in}}%
\pgfpathlineto{\pgfqpoint{1.395381in}{1.309582in}}%
\pgfpathlineto{\pgfqpoint{1.359322in}{1.312674in}}%
\pgfpathlineto{\pgfqpoint{1.358996in}{1.309848in}}%
\pgfpathlineto{\pgfqpoint{1.325420in}{1.312996in}}%
\pgfpathlineto{\pgfqpoint{1.326529in}{1.324213in}}%
\pgfpathclose%
\pgfusepath{fill}%
\end{pgfscope}%
\begin{pgfscope}%
\pgfpathrectangle{\pgfqpoint{0.100000in}{0.100000in}}{\pgfqpoint{3.007045in}{1.925000in}}%
\pgfusepath{clip}%
\pgfsetbuttcap%
\pgfsetmiterjoin%
\definecolor{currentfill}{rgb}{0.270804,0.580730,0.781146}%
\pgfsetfillcolor{currentfill}%
\pgfsetlinewidth{0.000000pt}%
\definecolor{currentstroke}{rgb}{0.000000,0.000000,0.000000}%
\pgfsetstrokecolor{currentstroke}%
\pgfsetstrokeopacity{0.000000}%
\pgfsetdash{}{0pt}%
\pgfpathmoveto{\pgfqpoint{2.011147in}{0.717467in}}%
\pgfpathlineto{\pgfqpoint{2.010976in}{0.714602in}}%
\pgfpathlineto{\pgfqpoint{2.003446in}{0.714212in}}%
\pgfpathlineto{\pgfqpoint{2.003803in}{0.705398in}}%
\pgfpathlineto{\pgfqpoint{1.998667in}{0.705092in}}%
\pgfpathlineto{\pgfqpoint{1.997671in}{0.713980in}}%
\pgfpathlineto{\pgfqpoint{1.993233in}{0.713829in}}%
\pgfpathlineto{\pgfqpoint{1.988496in}{0.714303in}}%
\pgfpathlineto{\pgfqpoint{1.986961in}{0.717838in}}%
\pgfpathlineto{\pgfqpoint{1.991838in}{0.723800in}}%
\pgfpathlineto{\pgfqpoint{1.985938in}{0.725835in}}%
\pgfpathlineto{\pgfqpoint{1.986138in}{0.729844in}}%
\pgfpathlineto{\pgfqpoint{1.990829in}{0.734444in}}%
\pgfpathlineto{\pgfqpoint{1.989679in}{0.740203in}}%
\pgfpathlineto{\pgfqpoint{1.985200in}{0.741858in}}%
\pgfpathlineto{\pgfqpoint{1.988976in}{0.750911in}}%
\pgfpathlineto{\pgfqpoint{1.988710in}{0.756370in}}%
\pgfpathlineto{\pgfqpoint{1.985739in}{0.761415in}}%
\pgfpathlineto{\pgfqpoint{1.985442in}{0.766767in}}%
\pgfpathlineto{\pgfqpoint{1.980431in}{0.778978in}}%
\pgfpathlineto{\pgfqpoint{1.985316in}{0.780979in}}%
\pgfpathlineto{\pgfqpoint{1.980623in}{0.785127in}}%
\pgfpathlineto{\pgfqpoint{1.984613in}{0.793228in}}%
\pgfpathlineto{\pgfqpoint{1.992182in}{0.792843in}}%
\pgfpathlineto{\pgfqpoint{1.988015in}{0.799119in}}%
\pgfpathlineto{\pgfqpoint{1.991139in}{0.803150in}}%
\pgfpathlineto{\pgfqpoint{1.986772in}{0.806255in}}%
\pgfpathlineto{\pgfqpoint{1.996950in}{0.809911in}}%
\pgfpathlineto{\pgfqpoint{1.998299in}{0.813452in}}%
\pgfpathlineto{\pgfqpoint{1.993621in}{0.816052in}}%
\pgfpathlineto{\pgfqpoint{1.993574in}{0.816135in}}%
\pgfpathlineto{\pgfqpoint{2.009788in}{0.816800in}}%
\pgfpathlineto{\pgfqpoint{2.010153in}{0.808128in}}%
\pgfpathlineto{\pgfqpoint{2.021640in}{0.808564in}}%
\pgfpathlineto{\pgfqpoint{2.022165in}{0.797050in}}%
\pgfpathlineto{\pgfqpoint{2.023674in}{0.765200in}}%
\pgfpathlineto{\pgfqpoint{2.018021in}{0.764899in}}%
\pgfpathlineto{\pgfqpoint{2.018223in}{0.761023in}}%
\pgfpathlineto{\pgfqpoint{2.009327in}{0.760562in}}%
\pgfpathlineto{\pgfqpoint{2.013096in}{0.749239in}}%
\pgfpathlineto{\pgfqpoint{2.013725in}{0.737732in}}%
\pgfpathlineto{\pgfqpoint{2.008261in}{0.731716in}}%
\pgfpathlineto{\pgfqpoint{2.011096in}{0.729886in}}%
\pgfpathlineto{\pgfqpoint{2.011147in}{0.717467in}}%
\pgfpathclose%
\pgfusepath{fill}%
\end{pgfscope}%
\begin{pgfscope}%
\pgfpathrectangle{\pgfqpoint{0.100000in}{0.100000in}}{\pgfqpoint{3.007045in}{1.925000in}}%
\pgfusepath{clip}%
\pgfsetbuttcap%
\pgfsetmiterjoin%
\definecolor{currentfill}{rgb}{0.662053,0.810196,0.897209}%
\pgfsetfillcolor{currentfill}%
\pgfsetlinewidth{0.000000pt}%
\definecolor{currentstroke}{rgb}{0.000000,0.000000,0.000000}%
\pgfsetstrokecolor{currentstroke}%
\pgfsetstrokeopacity{0.000000}%
\pgfsetdash{}{0pt}%
\pgfpathmoveto{\pgfqpoint{0.416964in}{1.700001in}}%
\pgfpathlineto{\pgfqpoint{0.427465in}{1.723088in}}%
\pgfpathlineto{\pgfqpoint{0.429108in}{1.730577in}}%
\pgfpathlineto{\pgfqpoint{0.441718in}{1.753288in}}%
\pgfpathlineto{\pgfqpoint{0.445758in}{1.765737in}}%
\pgfpathlineto{\pgfqpoint{0.453664in}{1.784658in}}%
\pgfpathlineto{\pgfqpoint{0.454831in}{1.792649in}}%
\pgfpathlineto{\pgfqpoint{0.465616in}{1.788403in}}%
\pgfpathlineto{\pgfqpoint{0.481516in}{1.783693in}}%
\pgfpathlineto{\pgfqpoint{0.480074in}{1.779212in}}%
\pgfpathlineto{\pgfqpoint{0.475846in}{1.775523in}}%
\pgfpathlineto{\pgfqpoint{0.475014in}{1.769793in}}%
\pgfpathlineto{\pgfqpoint{0.469874in}{1.763357in}}%
\pgfpathlineto{\pgfqpoint{0.465504in}{1.749616in}}%
\pgfpathlineto{\pgfqpoint{0.451249in}{1.754244in}}%
\pgfpathlineto{\pgfqpoint{0.448349in}{1.745454in}}%
\pgfpathlineto{\pgfqpoint{0.451037in}{1.744559in}}%
\pgfpathlineto{\pgfqpoint{0.444068in}{1.723357in}}%
\pgfpathlineto{\pgfqpoint{0.449171in}{1.720447in}}%
\pgfpathlineto{\pgfqpoint{0.443505in}{1.702226in}}%
\pgfpathlineto{\pgfqpoint{0.437975in}{1.704042in}}%
\pgfpathlineto{\pgfqpoint{0.436719in}{1.699407in}}%
\pgfpathlineto{\pgfqpoint{0.432427in}{1.695452in}}%
\pgfpathlineto{\pgfqpoint{0.416964in}{1.700001in}}%
\pgfpathclose%
\pgfusepath{fill}%
\end{pgfscope}%
\begin{pgfscope}%
\pgfpathrectangle{\pgfqpoint{0.100000in}{0.100000in}}{\pgfqpoint{3.007045in}{1.925000in}}%
\pgfusepath{clip}%
\pgfsetbuttcap%
\pgfsetmiterjoin%
\definecolor{currentfill}{rgb}{0.093272,0.396878,0.673664}%
\pgfsetfillcolor{currentfill}%
\pgfsetlinewidth{0.000000pt}%
\definecolor{currentstroke}{rgb}{0.000000,0.000000,0.000000}%
\pgfsetstrokecolor{currentstroke}%
\pgfsetstrokeopacity{0.000000}%
\pgfsetdash{}{0pt}%
\pgfpathmoveto{\pgfqpoint{0.377147in}{0.493245in}}%
\pgfpathlineto{\pgfqpoint{0.376520in}{0.497403in}}%
\pgfpathlineto{\pgfqpoint{0.377477in}{0.497571in}}%
\pgfpathlineto{\pgfqpoint{0.377147in}{0.493245in}}%
\pgfpathclose%
\pgfusepath{fill}%
\end{pgfscope}%
\begin{pgfscope}%
\pgfpathrectangle{\pgfqpoint{0.100000in}{0.100000in}}{\pgfqpoint{3.007045in}{1.925000in}}%
\pgfusepath{clip}%
\pgfsetbuttcap%
\pgfsetmiterjoin%
\definecolor{currentfill}{rgb}{0.093272,0.396878,0.673664}%
\pgfsetfillcolor{currentfill}%
\pgfsetlinewidth{0.000000pt}%
\definecolor{currentstroke}{rgb}{0.000000,0.000000,0.000000}%
\pgfsetstrokecolor{currentstroke}%
\pgfsetstrokeopacity{0.000000}%
\pgfsetdash{}{0pt}%
\pgfpathmoveto{\pgfqpoint{0.405385in}{0.465085in}}%
\pgfpathlineto{\pgfqpoint{0.405017in}{0.468510in}}%
\pgfpathlineto{\pgfqpoint{0.406749in}{0.467215in}}%
\pgfpathlineto{\pgfqpoint{0.407796in}{0.467652in}}%
\pgfpathlineto{\pgfqpoint{0.407185in}{0.472036in}}%
\pgfpathlineto{\pgfqpoint{0.408305in}{0.472399in}}%
\pgfpathlineto{\pgfqpoint{0.410270in}{0.469591in}}%
\pgfpathlineto{\pgfqpoint{0.409867in}{0.468623in}}%
\pgfpathlineto{\pgfqpoint{0.410756in}{0.466052in}}%
\pgfpathlineto{\pgfqpoint{0.406364in}{0.465975in}}%
\pgfpathlineto{\pgfqpoint{0.405385in}{0.465085in}}%
\pgfpathclose%
\pgfusepath{fill}%
\end{pgfscope}%
\begin{pgfscope}%
\pgfpathrectangle{\pgfqpoint{0.100000in}{0.100000in}}{\pgfqpoint{3.007045in}{1.925000in}}%
\pgfusepath{clip}%
\pgfsetbuttcap%
\pgfsetmiterjoin%
\definecolor{currentfill}{rgb}{0.093272,0.396878,0.673664}%
\pgfsetfillcolor{currentfill}%
\pgfsetlinewidth{0.000000pt}%
\definecolor{currentstroke}{rgb}{0.000000,0.000000,0.000000}%
\pgfsetstrokecolor{currentstroke}%
\pgfsetstrokeopacity{0.000000}%
\pgfsetdash{}{0pt}%
\pgfpathmoveto{\pgfqpoint{0.415072in}{0.452560in}}%
\pgfpathlineto{\pgfqpoint{0.414645in}{0.455401in}}%
\pgfpathlineto{\pgfqpoint{0.419140in}{0.455415in}}%
\pgfpathlineto{\pgfqpoint{0.420979in}{0.457297in}}%
\pgfpathlineto{\pgfqpoint{0.422131in}{0.456956in}}%
\pgfpathlineto{\pgfqpoint{0.423763in}{0.455028in}}%
\pgfpathlineto{\pgfqpoint{0.422007in}{0.454672in}}%
\pgfpathlineto{\pgfqpoint{0.421045in}{0.453407in}}%
\pgfpathlineto{\pgfqpoint{0.422643in}{0.451021in}}%
\pgfpathlineto{\pgfqpoint{0.423045in}{0.449284in}}%
\pgfpathlineto{\pgfqpoint{0.421244in}{0.449179in}}%
\pgfpathlineto{\pgfqpoint{0.420180in}{0.450490in}}%
\pgfpathlineto{\pgfqpoint{0.417429in}{0.452432in}}%
\pgfpathlineto{\pgfqpoint{0.415072in}{0.452560in}}%
\pgfpathclose%
\pgfusepath{fill}%
\end{pgfscope}%
\begin{pgfscope}%
\pgfpathrectangle{\pgfqpoint{0.100000in}{0.100000in}}{\pgfqpoint{3.007045in}{1.925000in}}%
\pgfusepath{clip}%
\pgfsetbuttcap%
\pgfsetmiterjoin%
\definecolor{currentfill}{rgb}{0.093272,0.396878,0.673664}%
\pgfsetfillcolor{currentfill}%
\pgfsetlinewidth{0.000000pt}%
\definecolor{currentstroke}{rgb}{0.000000,0.000000,0.000000}%
\pgfsetstrokecolor{currentstroke}%
\pgfsetstrokeopacity{0.000000}%
\pgfsetdash{}{0pt}%
\pgfpathmoveto{\pgfqpoint{0.412280in}{0.458704in}}%
\pgfpathlineto{\pgfqpoint{0.412274in}{0.460723in}}%
\pgfpathlineto{\pgfqpoint{0.408425in}{0.462631in}}%
\pgfpathlineto{\pgfqpoint{0.410054in}{0.463833in}}%
\pgfpathlineto{\pgfqpoint{0.412335in}{0.461427in}}%
\pgfpathlineto{\pgfqpoint{0.415525in}{0.460193in}}%
\pgfpathlineto{\pgfqpoint{0.417642in}{0.461549in}}%
\pgfpathlineto{\pgfqpoint{0.418112in}{0.459372in}}%
\pgfpathlineto{\pgfqpoint{0.415856in}{0.459267in}}%
\pgfpathlineto{\pgfqpoint{0.414135in}{0.457815in}}%
\pgfpathlineto{\pgfqpoint{0.412280in}{0.458704in}}%
\pgfpathclose%
\pgfusepath{fill}%
\end{pgfscope}%
\begin{pgfscope}%
\pgfpathrectangle{\pgfqpoint{0.100000in}{0.100000in}}{\pgfqpoint{3.007045in}{1.925000in}}%
\pgfusepath{clip}%
\pgfsetbuttcap%
\pgfsetmiterjoin%
\definecolor{currentfill}{rgb}{0.093272,0.396878,0.673664}%
\pgfsetfillcolor{currentfill}%
\pgfsetlinewidth{0.000000pt}%
\definecolor{currentstroke}{rgb}{0.000000,0.000000,0.000000}%
\pgfsetstrokecolor{currentstroke}%
\pgfsetstrokeopacity{0.000000}%
\pgfsetdash{}{0pt}%
\pgfpathmoveto{\pgfqpoint{0.389271in}{0.496402in}}%
\pgfpathlineto{\pgfqpoint{0.388843in}{0.498609in}}%
\pgfpathlineto{\pgfqpoint{0.391517in}{0.499040in}}%
\pgfpathlineto{\pgfqpoint{0.392195in}{0.498040in}}%
\pgfpathlineto{\pgfqpoint{0.391561in}{0.496110in}}%
\pgfpathlineto{\pgfqpoint{0.389271in}{0.496402in}}%
\pgfpathclose%
\pgfusepath{fill}%
\end{pgfscope}%
\begin{pgfscope}%
\pgfpathrectangle{\pgfqpoint{0.100000in}{0.100000in}}{\pgfqpoint{3.007045in}{1.925000in}}%
\pgfusepath{clip}%
\pgfsetbuttcap%
\pgfsetmiterjoin%
\definecolor{currentfill}{rgb}{0.093272,0.396878,0.673664}%
\pgfsetfillcolor{currentfill}%
\pgfsetlinewidth{0.000000pt}%
\definecolor{currentstroke}{rgb}{0.000000,0.000000,0.000000}%
\pgfsetstrokecolor{currentstroke}%
\pgfsetstrokeopacity{0.000000}%
\pgfsetdash{}{0pt}%
\pgfpathmoveto{\pgfqpoint{0.424629in}{0.447465in}}%
\pgfpathlineto{\pgfqpoint{0.424048in}{0.449531in}}%
\pgfpathlineto{\pgfqpoint{0.427830in}{0.447738in}}%
\pgfpathlineto{\pgfqpoint{0.426724in}{0.446647in}}%
\pgfpathlineto{\pgfqpoint{0.424629in}{0.447465in}}%
\pgfpathclose%
\pgfusepath{fill}%
\end{pgfscope}%
\begin{pgfscope}%
\pgfpathrectangle{\pgfqpoint{0.100000in}{0.100000in}}{\pgfqpoint{3.007045in}{1.925000in}}%
\pgfusepath{clip}%
\pgfsetbuttcap%
\pgfsetmiterjoin%
\definecolor{currentfill}{rgb}{0.093272,0.396878,0.673664}%
\pgfsetfillcolor{currentfill}%
\pgfsetlinewidth{0.000000pt}%
\definecolor{currentstroke}{rgb}{0.000000,0.000000,0.000000}%
\pgfsetstrokecolor{currentstroke}%
\pgfsetstrokeopacity{0.000000}%
\pgfsetdash{}{0pt}%
\pgfpathmoveto{\pgfqpoint{0.380099in}{0.509690in}}%
\pgfpathlineto{\pgfqpoint{0.380270in}{0.511273in}}%
\pgfpathlineto{\pgfqpoint{0.382192in}{0.510622in}}%
\pgfpathlineto{\pgfqpoint{0.381722in}{0.509141in}}%
\pgfpathlineto{\pgfqpoint{0.380099in}{0.509690in}}%
\pgfpathclose%
\pgfusepath{fill}%
\end{pgfscope}%
\begin{pgfscope}%
\pgfpathrectangle{\pgfqpoint{0.100000in}{0.100000in}}{\pgfqpoint{3.007045in}{1.925000in}}%
\pgfusepath{clip}%
\pgfsetbuttcap%
\pgfsetmiterjoin%
\definecolor{currentfill}{rgb}{0.093272,0.396878,0.673664}%
\pgfsetfillcolor{currentfill}%
\pgfsetlinewidth{0.000000pt}%
\definecolor{currentstroke}{rgb}{0.000000,0.000000,0.000000}%
\pgfsetstrokecolor{currentstroke}%
\pgfsetstrokeopacity{0.000000}%
\pgfsetdash{}{0pt}%
\pgfpathmoveto{\pgfqpoint{0.428107in}{0.449417in}}%
\pgfpathlineto{\pgfqpoint{0.427960in}{0.451570in}}%
\pgfpathlineto{\pgfqpoint{0.429777in}{0.452594in}}%
\pgfpathlineto{\pgfqpoint{0.430464in}{0.450596in}}%
\pgfpathlineto{\pgfqpoint{0.428107in}{0.449417in}}%
\pgfpathclose%
\pgfusepath{fill}%
\end{pgfscope}%
\begin{pgfscope}%
\pgfpathrectangle{\pgfqpoint{0.100000in}{0.100000in}}{\pgfqpoint{3.007045in}{1.925000in}}%
\pgfusepath{clip}%
\pgfsetbuttcap%
\pgfsetmiterjoin%
\definecolor{currentfill}{rgb}{0.093272,0.396878,0.673664}%
\pgfsetfillcolor{currentfill}%
\pgfsetlinewidth{0.000000pt}%
\definecolor{currentstroke}{rgb}{0.000000,0.000000,0.000000}%
\pgfsetstrokecolor{currentstroke}%
\pgfsetstrokeopacity{0.000000}%
\pgfsetdash{}{0pt}%
\pgfpathmoveto{\pgfqpoint{0.369259in}{0.522705in}}%
\pgfpathlineto{\pgfqpoint{0.369319in}{0.524685in}}%
\pgfpathlineto{\pgfqpoint{0.372098in}{0.524246in}}%
\pgfpathlineto{\pgfqpoint{0.373941in}{0.523135in}}%
\pgfpathlineto{\pgfqpoint{0.376137in}{0.523717in}}%
\pgfpathlineto{\pgfqpoint{0.377311in}{0.522297in}}%
\pgfpathlineto{\pgfqpoint{0.374074in}{0.522045in}}%
\pgfpathlineto{\pgfqpoint{0.373276in}{0.520714in}}%
\pgfpathlineto{\pgfqpoint{0.371566in}{0.523261in}}%
\pgfpathlineto{\pgfqpoint{0.369259in}{0.522705in}}%
\pgfpathclose%
\pgfusepath{fill}%
\end{pgfscope}%
\begin{pgfscope}%
\pgfpathrectangle{\pgfqpoint{0.100000in}{0.100000in}}{\pgfqpoint{3.007045in}{1.925000in}}%
\pgfusepath{clip}%
\pgfsetbuttcap%
\pgfsetmiterjoin%
\definecolor{currentfill}{rgb}{0.093272,0.396878,0.673664}%
\pgfsetfillcolor{currentfill}%
\pgfsetlinewidth{0.000000pt}%
\definecolor{currentstroke}{rgb}{0.000000,0.000000,0.000000}%
\pgfsetstrokecolor{currentstroke}%
\pgfsetstrokeopacity{0.000000}%
\pgfsetdash{}{0pt}%
\pgfpathmoveto{\pgfqpoint{0.465052in}{0.415402in}}%
\pgfpathlineto{\pgfqpoint{0.464568in}{0.416834in}}%
\pgfpathlineto{\pgfqpoint{0.468037in}{0.416734in}}%
\pgfpathlineto{\pgfqpoint{0.469080in}{0.414904in}}%
\pgfpathlineto{\pgfqpoint{0.468333in}{0.414126in}}%
\pgfpathlineto{\pgfqpoint{0.465052in}{0.415402in}}%
\pgfpathclose%
\pgfusepath{fill}%
\end{pgfscope}%
\begin{pgfscope}%
\pgfpathrectangle{\pgfqpoint{0.100000in}{0.100000in}}{\pgfqpoint{3.007045in}{1.925000in}}%
\pgfusepath{clip}%
\pgfsetbuttcap%
\pgfsetmiterjoin%
\definecolor{currentfill}{rgb}{0.093272,0.396878,0.673664}%
\pgfsetfillcolor{currentfill}%
\pgfsetlinewidth{0.000000pt}%
\definecolor{currentstroke}{rgb}{0.000000,0.000000,0.000000}%
\pgfsetstrokecolor{currentstroke}%
\pgfsetstrokeopacity{0.000000}%
\pgfsetdash{}{0pt}%
\pgfpathmoveto{\pgfqpoint{0.351227in}{0.571842in}}%
\pgfpathlineto{\pgfqpoint{0.350829in}{0.574370in}}%
\pgfpathlineto{\pgfqpoint{0.351503in}{0.575341in}}%
\pgfpathlineto{\pgfqpoint{0.353352in}{0.574845in}}%
\pgfpathlineto{\pgfqpoint{0.353069in}{0.573567in}}%
\pgfpathlineto{\pgfqpoint{0.351227in}{0.571842in}}%
\pgfpathclose%
\pgfusepath{fill}%
\end{pgfscope}%
\begin{pgfscope}%
\pgfpathrectangle{\pgfqpoint{0.100000in}{0.100000in}}{\pgfqpoint{3.007045in}{1.925000in}}%
\pgfusepath{clip}%
\pgfsetbuttcap%
\pgfsetmiterjoin%
\definecolor{currentfill}{rgb}{0.093272,0.396878,0.673664}%
\pgfsetfillcolor{currentfill}%
\pgfsetlinewidth{0.000000pt}%
\definecolor{currentstroke}{rgb}{0.000000,0.000000,0.000000}%
\pgfsetstrokecolor{currentstroke}%
\pgfsetstrokeopacity{0.000000}%
\pgfsetdash{}{0pt}%
\pgfpathmoveto{\pgfqpoint{0.486654in}{0.401495in}}%
\pgfpathlineto{\pgfqpoint{0.487380in}{0.403172in}}%
\pgfpathlineto{\pgfqpoint{0.490544in}{0.402448in}}%
\pgfpathlineto{\pgfqpoint{0.489751in}{0.400770in}}%
\pgfpathlineto{\pgfqpoint{0.486654in}{0.401495in}}%
\pgfpathclose%
\pgfusepath{fill}%
\end{pgfscope}%
\begin{pgfscope}%
\pgfpathrectangle{\pgfqpoint{0.100000in}{0.100000in}}{\pgfqpoint{3.007045in}{1.925000in}}%
\pgfusepath{clip}%
\pgfsetbuttcap%
\pgfsetmiterjoin%
\definecolor{currentfill}{rgb}{0.093272,0.396878,0.673664}%
\pgfsetfillcolor{currentfill}%
\pgfsetlinewidth{0.000000pt}%
\definecolor{currentstroke}{rgb}{0.000000,0.000000,0.000000}%
\pgfsetstrokecolor{currentstroke}%
\pgfsetstrokeopacity{0.000000}%
\pgfsetdash{}{0pt}%
\pgfpathmoveto{\pgfqpoint{0.355532in}{0.583863in}}%
\pgfpathlineto{\pgfqpoint{0.355346in}{0.585775in}}%
\pgfpathlineto{\pgfqpoint{0.352345in}{0.588784in}}%
\pgfpathlineto{\pgfqpoint{0.354611in}{0.589330in}}%
\pgfpathlineto{\pgfqpoint{0.353122in}{0.591113in}}%
\pgfpathlineto{\pgfqpoint{0.354013in}{0.592136in}}%
\pgfpathlineto{\pgfqpoint{0.356001in}{0.592468in}}%
\pgfpathlineto{\pgfqpoint{0.357361in}{0.589332in}}%
\pgfpathlineto{\pgfqpoint{0.359106in}{0.586517in}}%
\pgfpathlineto{\pgfqpoint{0.359311in}{0.584053in}}%
\pgfpathlineto{\pgfqpoint{0.357819in}{0.583219in}}%
\pgfpathlineto{\pgfqpoint{0.355532in}{0.583863in}}%
\pgfpathclose%
\pgfusepath{fill}%
\end{pgfscope}%
\begin{pgfscope}%
\pgfpathrectangle{\pgfqpoint{0.100000in}{0.100000in}}{\pgfqpoint{3.007045in}{1.925000in}}%
\pgfusepath{clip}%
\pgfsetbuttcap%
\pgfsetmiterjoin%
\definecolor{currentfill}{rgb}{0.093272,0.396878,0.673664}%
\pgfsetfillcolor{currentfill}%
\pgfsetlinewidth{0.000000pt}%
\definecolor{currentstroke}{rgb}{0.000000,0.000000,0.000000}%
\pgfsetstrokecolor{currentstroke}%
\pgfsetstrokeopacity{0.000000}%
\pgfsetdash{}{0pt}%
\pgfpathmoveto{\pgfqpoint{0.528105in}{0.380782in}}%
\pgfpathlineto{\pgfqpoint{0.527637in}{0.383157in}}%
\pgfpathlineto{\pgfqpoint{0.528964in}{0.383518in}}%
\pgfpathlineto{\pgfqpoint{0.532084in}{0.382526in}}%
\pgfpathlineto{\pgfqpoint{0.533585in}{0.381267in}}%
\pgfpathlineto{\pgfqpoint{0.535312in}{0.380961in}}%
\pgfpathlineto{\pgfqpoint{0.536723in}{0.379546in}}%
\pgfpathlineto{\pgfqpoint{0.538639in}{0.381047in}}%
\pgfpathlineto{\pgfqpoint{0.541713in}{0.382292in}}%
\pgfpathlineto{\pgfqpoint{0.544259in}{0.379787in}}%
\pgfpathlineto{\pgfqpoint{0.542790in}{0.385194in}}%
\pgfpathlineto{\pgfqpoint{0.544767in}{0.386107in}}%
\pgfpathlineto{\pgfqpoint{0.546961in}{0.385637in}}%
\pgfpathlineto{\pgfqpoint{0.550504in}{0.383713in}}%
\pgfpathlineto{\pgfqpoint{0.549689in}{0.382189in}}%
\pgfpathlineto{\pgfqpoint{0.551450in}{0.380178in}}%
\pgfpathlineto{\pgfqpoint{0.553278in}{0.380675in}}%
\pgfpathlineto{\pgfqpoint{0.552656in}{0.378168in}}%
\pgfpathlineto{\pgfqpoint{0.549658in}{0.377895in}}%
\pgfpathlineto{\pgfqpoint{0.548439in}{0.374718in}}%
\pgfpathlineto{\pgfqpoint{0.543340in}{0.376546in}}%
\pgfpathlineto{\pgfqpoint{0.541469in}{0.376728in}}%
\pgfpathlineto{\pgfqpoint{0.540243in}{0.375113in}}%
\pgfpathlineto{\pgfqpoint{0.538796in}{0.377500in}}%
\pgfpathlineto{\pgfqpoint{0.537382in}{0.378201in}}%
\pgfpathlineto{\pgfqpoint{0.533254in}{0.378923in}}%
\pgfpathlineto{\pgfqpoint{0.531395in}{0.379938in}}%
\pgfpathlineto{\pgfqpoint{0.528105in}{0.380782in}}%
\pgfpathclose%
\pgfusepath{fill}%
\end{pgfscope}%
\begin{pgfscope}%
\pgfpathrectangle{\pgfqpoint{0.100000in}{0.100000in}}{\pgfqpoint{3.007045in}{1.925000in}}%
\pgfusepath{clip}%
\pgfsetbuttcap%
\pgfsetmiterjoin%
\definecolor{currentfill}{rgb}{0.093272,0.396878,0.673664}%
\pgfsetfillcolor{currentfill}%
\pgfsetlinewidth{0.000000pt}%
\definecolor{currentstroke}{rgb}{0.000000,0.000000,0.000000}%
\pgfsetstrokecolor{currentstroke}%
\pgfsetstrokeopacity{0.000000}%
\pgfsetdash{}{0pt}%
\pgfpathmoveto{\pgfqpoint{0.500539in}{0.394151in}}%
\pgfpathlineto{\pgfqpoint{0.500162in}{0.396419in}}%
\pgfpathlineto{\pgfqpoint{0.502003in}{0.396647in}}%
\pgfpathlineto{\pgfqpoint{0.502049in}{0.394805in}}%
\pgfpathlineto{\pgfqpoint{0.500539in}{0.394151in}}%
\pgfpathclose%
\pgfusepath{fill}%
\end{pgfscope}%
\begin{pgfscope}%
\pgfpathrectangle{\pgfqpoint{0.100000in}{0.100000in}}{\pgfqpoint{3.007045in}{1.925000in}}%
\pgfusepath{clip}%
\pgfsetbuttcap%
\pgfsetmiterjoin%
\definecolor{currentfill}{rgb}{0.093272,0.396878,0.673664}%
\pgfsetfillcolor{currentfill}%
\pgfsetlinewidth{0.000000pt}%
\definecolor{currentstroke}{rgb}{0.000000,0.000000,0.000000}%
\pgfsetstrokecolor{currentstroke}%
\pgfsetstrokeopacity{0.000000}%
\pgfsetdash{}{0pt}%
\pgfpathmoveto{\pgfqpoint{0.507877in}{0.388268in}}%
\pgfpathlineto{\pgfqpoint{0.508331in}{0.389076in}}%
\pgfpathlineto{\pgfqpoint{0.513005in}{0.389208in}}%
\pgfpathlineto{\pgfqpoint{0.515777in}{0.391018in}}%
\pgfpathlineto{\pgfqpoint{0.519331in}{0.390800in}}%
\pgfpathlineto{\pgfqpoint{0.521540in}{0.387996in}}%
\pgfpathlineto{\pgfqpoint{0.523561in}{0.390645in}}%
\pgfpathlineto{\pgfqpoint{0.525758in}{0.391376in}}%
\pgfpathlineto{\pgfqpoint{0.528197in}{0.390858in}}%
\pgfpathlineto{\pgfqpoint{0.530050in}{0.389563in}}%
\pgfpathlineto{\pgfqpoint{0.530693in}{0.387435in}}%
\pgfpathlineto{\pgfqpoint{0.530317in}{0.386066in}}%
\pgfpathlineto{\pgfqpoint{0.528322in}{0.384961in}}%
\pgfpathlineto{\pgfqpoint{0.521632in}{0.386913in}}%
\pgfpathlineto{\pgfqpoint{0.517279in}{0.385660in}}%
\pgfpathlineto{\pgfqpoint{0.515618in}{0.386416in}}%
\pgfpathlineto{\pgfqpoint{0.507877in}{0.388268in}}%
\pgfpathclose%
\pgfusepath{fill}%
\end{pgfscope}%
\begin{pgfscope}%
\pgfpathrectangle{\pgfqpoint{0.100000in}{0.100000in}}{\pgfqpoint{3.007045in}{1.925000in}}%
\pgfusepath{clip}%
\pgfsetbuttcap%
\pgfsetmiterjoin%
\definecolor{currentfill}{rgb}{0.093272,0.396878,0.673664}%
\pgfsetfillcolor{currentfill}%
\pgfsetlinewidth{0.000000pt}%
\definecolor{currentstroke}{rgb}{0.000000,0.000000,0.000000}%
\pgfsetstrokecolor{currentstroke}%
\pgfsetstrokeopacity{0.000000}%
\pgfsetdash{}{0pt}%
\pgfpathmoveto{\pgfqpoint{0.453614in}{0.421300in}}%
\pgfpathlineto{\pgfqpoint{0.451899in}{0.423133in}}%
\pgfpathlineto{\pgfqpoint{0.450880in}{0.425283in}}%
\pgfpathlineto{\pgfqpoint{0.449241in}{0.426901in}}%
\pgfpathlineto{\pgfqpoint{0.449633in}{0.429133in}}%
\pgfpathlineto{\pgfqpoint{0.453091in}{0.425783in}}%
\pgfpathlineto{\pgfqpoint{0.455129in}{0.421602in}}%
\pgfpathlineto{\pgfqpoint{0.453614in}{0.421300in}}%
\pgfpathclose%
\pgfusepath{fill}%
\end{pgfscope}%
\begin{pgfscope}%
\pgfpathrectangle{\pgfqpoint{0.100000in}{0.100000in}}{\pgfqpoint{3.007045in}{1.925000in}}%
\pgfusepath{clip}%
\pgfsetbuttcap%
\pgfsetmiterjoin%
\definecolor{currentfill}{rgb}{0.093272,0.396878,0.673664}%
\pgfsetfillcolor{currentfill}%
\pgfsetlinewidth{0.000000pt}%
\definecolor{currentstroke}{rgb}{0.000000,0.000000,0.000000}%
\pgfsetstrokecolor{currentstroke}%
\pgfsetstrokeopacity{0.000000}%
\pgfsetdash{}{0pt}%
\pgfpathmoveto{\pgfqpoint{0.444349in}{0.432536in}}%
\pgfpathlineto{\pgfqpoint{0.443210in}{0.434576in}}%
\pgfpathlineto{\pgfqpoint{0.440911in}{0.435309in}}%
\pgfpathlineto{\pgfqpoint{0.440187in}{0.438580in}}%
\pgfpathlineto{\pgfqpoint{0.442735in}{0.437137in}}%
\pgfpathlineto{\pgfqpoint{0.443790in}{0.435233in}}%
\pgfpathlineto{\pgfqpoint{0.445554in}{0.435638in}}%
\pgfpathlineto{\pgfqpoint{0.448969in}{0.435304in}}%
\pgfpathlineto{\pgfqpoint{0.450814in}{0.436419in}}%
\pgfpathlineto{\pgfqpoint{0.452553in}{0.435695in}}%
\pgfpathlineto{\pgfqpoint{0.453462in}{0.433672in}}%
\pgfpathlineto{\pgfqpoint{0.453143in}{0.432430in}}%
\pgfpathlineto{\pgfqpoint{0.451239in}{0.431442in}}%
\pgfpathlineto{\pgfqpoint{0.449605in}{0.432844in}}%
\pgfpathlineto{\pgfqpoint{0.448738in}{0.430447in}}%
\pgfpathlineto{\pgfqpoint{0.447526in}{0.430880in}}%
\pgfpathlineto{\pgfqpoint{0.445671in}{0.433246in}}%
\pgfpathlineto{\pgfqpoint{0.444349in}{0.432536in}}%
\pgfpathclose%
\pgfusepath{fill}%
\end{pgfscope}%
\begin{pgfscope}%
\pgfpathrectangle{\pgfqpoint{0.100000in}{0.100000in}}{\pgfqpoint{3.007045in}{1.925000in}}%
\pgfusepath{clip}%
\pgfsetbuttcap%
\pgfsetmiterjoin%
\definecolor{currentfill}{rgb}{0.579608,0.770196,0.873725}%
\pgfsetfillcolor{currentfill}%
\pgfsetlinewidth{0.000000pt}%
\definecolor{currentstroke}{rgb}{0.000000,0.000000,0.000000}%
\pgfsetstrokecolor{currentstroke}%
\pgfsetstrokeopacity{0.000000}%
\pgfsetdash{}{0pt}%
\pgfpathmoveto{\pgfqpoint{1.876321in}{0.768157in}}%
\pgfpathlineto{\pgfqpoint{1.870545in}{0.767215in}}%
\pgfpathlineto{\pgfqpoint{1.856034in}{0.767401in}}%
\pgfpathlineto{\pgfqpoint{1.856051in}{0.769311in}}%
\pgfpathlineto{\pgfqpoint{1.842812in}{0.769497in}}%
\pgfpathlineto{\pgfqpoint{1.836570in}{0.777705in}}%
\pgfpathlineto{\pgfqpoint{1.829258in}{0.781767in}}%
\pgfpathlineto{\pgfqpoint{1.829586in}{0.787129in}}%
\pgfpathlineto{\pgfqpoint{1.828440in}{0.791464in}}%
\pgfpathlineto{\pgfqpoint{1.824314in}{0.796016in}}%
\pgfpathlineto{\pgfqpoint{1.822337in}{0.805756in}}%
\pgfpathlineto{\pgfqpoint{1.821492in}{0.816167in}}%
\pgfpathlineto{\pgfqpoint{1.813675in}{0.816167in}}%
\pgfpathlineto{\pgfqpoint{1.812989in}{0.826934in}}%
\pgfpathlineto{\pgfqpoint{1.830412in}{0.826752in}}%
\pgfpathlineto{\pgfqpoint{1.830433in}{0.847601in}}%
\pgfpathlineto{\pgfqpoint{1.833670in}{0.850751in}}%
\pgfpathlineto{\pgfqpoint{1.842321in}{0.853020in}}%
\pgfpathlineto{\pgfqpoint{1.859472in}{0.853168in}}%
\pgfpathlineto{\pgfqpoint{1.865214in}{0.855241in}}%
\pgfpathlineto{\pgfqpoint{1.882494in}{0.855536in}}%
\pgfpathlineto{\pgfqpoint{1.882114in}{0.846605in}}%
\pgfpathlineto{\pgfqpoint{1.887792in}{0.846596in}}%
\pgfpathlineto{\pgfqpoint{1.891528in}{0.843729in}}%
\pgfpathlineto{\pgfqpoint{1.892466in}{0.838016in}}%
\pgfpathlineto{\pgfqpoint{1.898412in}{0.836142in}}%
\pgfpathlineto{\pgfqpoint{1.899400in}{0.832281in}}%
\pgfpathlineto{\pgfqpoint{1.899320in}{0.814784in}}%
\pgfpathlineto{\pgfqpoint{1.887940in}{0.814920in}}%
\pgfpathlineto{\pgfqpoint{1.887657in}{0.795131in}}%
\pgfpathlineto{\pgfqpoint{1.884971in}{0.791103in}}%
\pgfpathlineto{\pgfqpoint{1.876541in}{0.789626in}}%
\pgfpathlineto{\pgfqpoint{1.876321in}{0.768157in}}%
\pgfpathclose%
\pgfusepath{fill}%
\end{pgfscope}%
\begin{pgfscope}%
\pgfpathrectangle{\pgfqpoint{0.100000in}{0.100000in}}{\pgfqpoint{3.007045in}{1.925000in}}%
\pgfusepath{clip}%
\pgfsetbuttcap%
\pgfsetmiterjoin%
\definecolor{currentfill}{rgb}{0.346467,0.632403,0.810673}%
\pgfsetfillcolor{currentfill}%
\pgfsetlinewidth{0.000000pt}%
\definecolor{currentstroke}{rgb}{0.000000,0.000000,0.000000}%
\pgfsetstrokecolor{currentstroke}%
\pgfsetstrokeopacity{0.000000}%
\pgfsetdash{}{0pt}%
\pgfpathmoveto{\pgfqpoint{0.596303in}{1.682378in}}%
\pgfpathlineto{\pgfqpoint{0.563962in}{1.691905in}}%
\pgfpathlineto{\pgfqpoint{0.533873in}{1.701067in}}%
\pgfpathlineto{\pgfqpoint{0.536189in}{1.704696in}}%
\pgfpathlineto{\pgfqpoint{0.535329in}{1.713270in}}%
\pgfpathlineto{\pgfqpoint{0.541696in}{1.715411in}}%
\pgfpathlineto{\pgfqpoint{0.541409in}{1.723741in}}%
\pgfpathlineto{\pgfqpoint{0.545552in}{1.726930in}}%
\pgfpathlineto{\pgfqpoint{0.547023in}{1.735645in}}%
\pgfpathlineto{\pgfqpoint{0.540688in}{1.747112in}}%
\pgfpathlineto{\pgfqpoint{0.542490in}{1.755480in}}%
\pgfpathlineto{\pgfqpoint{0.545786in}{1.757536in}}%
\pgfpathlineto{\pgfqpoint{0.551653in}{1.755806in}}%
\pgfpathlineto{\pgfqpoint{0.561232in}{1.755167in}}%
\pgfpathlineto{\pgfqpoint{0.558944in}{1.759738in}}%
\pgfpathlineto{\pgfqpoint{0.563789in}{1.776383in}}%
\pgfpathlineto{\pgfqpoint{0.567701in}{1.775221in}}%
\pgfpathlineto{\pgfqpoint{0.616044in}{1.760740in}}%
\pgfpathlineto{\pgfqpoint{0.640972in}{1.753699in}}%
\pgfpathlineto{\pgfqpoint{0.637219in}{1.740644in}}%
\pgfpathlineto{\pgfqpoint{0.630951in}{1.740790in}}%
\pgfpathlineto{\pgfqpoint{0.619982in}{1.737934in}}%
\pgfpathlineto{\pgfqpoint{0.610892in}{1.738691in}}%
\pgfpathlineto{\pgfqpoint{0.606293in}{1.736988in}}%
\pgfpathlineto{\pgfqpoint{0.608690in}{1.732151in}}%
\pgfpathlineto{\pgfqpoint{0.608895in}{1.725281in}}%
\pgfpathlineto{\pgfqpoint{0.602494in}{1.724955in}}%
\pgfpathlineto{\pgfqpoint{0.595993in}{1.712179in}}%
\pgfpathlineto{\pgfqpoint{0.598195in}{1.705394in}}%
\pgfpathlineto{\pgfqpoint{0.593379in}{1.692150in}}%
\pgfpathlineto{\pgfqpoint{0.593360in}{1.686222in}}%
\pgfpathlineto{\pgfqpoint{0.596303in}{1.682378in}}%
\pgfpathclose%
\pgfusepath{fill}%
\end{pgfscope}%
\begin{pgfscope}%
\pgfpathrectangle{\pgfqpoint{0.100000in}{0.100000in}}{\pgfqpoint{3.007045in}{1.925000in}}%
\pgfusepath{clip}%
\pgfsetbuttcap%
\pgfsetmiterjoin%
\definecolor{currentfill}{rgb}{0.108651,0.416563,0.689043}%
\pgfsetfillcolor{currentfill}%
\pgfsetlinewidth{0.000000pt}%
\definecolor{currentstroke}{rgb}{0.000000,0.000000,0.000000}%
\pgfsetstrokecolor{currentstroke}%
\pgfsetstrokeopacity{0.000000}%
\pgfsetdash{}{0pt}%
\pgfpathmoveto{\pgfqpoint{1.473659in}{0.890543in}}%
\pgfpathlineto{\pgfqpoint{1.444186in}{0.892364in}}%
\pgfpathlineto{\pgfqpoint{1.415306in}{0.894451in}}%
\pgfpathlineto{\pgfqpoint{1.417413in}{0.923222in}}%
\pgfpathlineto{\pgfqpoint{1.419469in}{0.951860in}}%
\pgfpathlineto{\pgfqpoint{1.477103in}{0.948069in}}%
\pgfpathlineto{\pgfqpoint{1.505836in}{0.946347in}}%
\pgfpathlineto{\pgfqpoint{1.505202in}{0.934838in}}%
\pgfpathlineto{\pgfqpoint{1.504252in}{0.917589in}}%
\pgfpathlineto{\pgfqpoint{1.475346in}{0.919280in}}%
\pgfpathlineto{\pgfqpoint{1.473659in}{0.890543in}}%
\pgfpathclose%
\pgfusepath{fill}%
\end{pgfscope}%
\begin{pgfscope}%
\pgfpathrectangle{\pgfqpoint{0.100000in}{0.100000in}}{\pgfqpoint{3.007045in}{1.925000in}}%
\pgfusepath{clip}%
\pgfsetbuttcap%
\pgfsetmiterjoin%
\definecolor{currentfill}{rgb}{0.504314,0.728858,0.857486}%
\pgfsetfillcolor{currentfill}%
\pgfsetlinewidth{0.000000pt}%
\definecolor{currentstroke}{rgb}{0.000000,0.000000,0.000000}%
\pgfsetstrokecolor{currentstroke}%
\pgfsetstrokeopacity{0.000000}%
\pgfsetdash{}{0pt}%
\pgfpathmoveto{\pgfqpoint{1.586901in}{0.635387in}}%
\pgfpathlineto{\pgfqpoint{1.569891in}{0.623670in}}%
\pgfpathlineto{\pgfqpoint{1.571686in}{0.614167in}}%
\pgfpathlineto{\pgfqpoint{1.576374in}{0.610043in}}%
\pgfpathlineto{\pgfqpoint{1.575754in}{0.602872in}}%
\pgfpathlineto{\pgfqpoint{1.546174in}{0.604074in}}%
\pgfpathlineto{\pgfqpoint{1.538952in}{0.605702in}}%
\pgfpathlineto{\pgfqpoint{1.540604in}{0.640218in}}%
\pgfpathlineto{\pgfqpoint{1.532433in}{0.642045in}}%
\pgfpathlineto{\pgfqpoint{1.527773in}{0.637409in}}%
\pgfpathlineto{\pgfqpoint{1.520806in}{0.642095in}}%
\pgfpathlineto{\pgfqpoint{1.510361in}{0.642115in}}%
\pgfpathlineto{\pgfqpoint{1.505217in}{0.648714in}}%
\pgfpathlineto{\pgfqpoint{1.507461in}{0.683119in}}%
\pgfpathlineto{\pgfqpoint{1.551784in}{0.680740in}}%
\pgfpathlineto{\pgfqpoint{1.565148in}{0.655103in}}%
\pgfpathlineto{\pgfqpoint{1.575122in}{0.656245in}}%
\pgfpathlineto{\pgfqpoint{1.586901in}{0.635387in}}%
\pgfpathclose%
\pgfusepath{fill}%
\end{pgfscope}%
\begin{pgfscope}%
\pgfpathrectangle{\pgfqpoint{0.100000in}{0.100000in}}{\pgfqpoint{3.007045in}{1.925000in}}%
\pgfusepath{clip}%
\pgfsetbuttcap%
\pgfsetmiterjoin%
\definecolor{currentfill}{rgb}{0.548235,0.752972,0.866959}%
\pgfsetfillcolor{currentfill}%
\pgfsetlinewidth{0.000000pt}%
\definecolor{currentstroke}{rgb}{0.000000,0.000000,0.000000}%
\pgfsetstrokecolor{currentstroke}%
\pgfsetstrokeopacity{0.000000}%
\pgfsetdash{}{0pt}%
\pgfpathmoveto{\pgfqpoint{2.442205in}{0.558231in}}%
\pgfpathlineto{\pgfqpoint{2.433773in}{0.561843in}}%
\pgfpathlineto{\pgfqpoint{2.427536in}{0.571164in}}%
\pgfpathlineto{\pgfqpoint{2.419054in}{0.576060in}}%
\pgfpathlineto{\pgfqpoint{2.409749in}{0.579100in}}%
\pgfpathlineto{\pgfqpoint{2.405641in}{0.582074in}}%
\pgfpathlineto{\pgfqpoint{2.408711in}{0.592974in}}%
\pgfpathlineto{\pgfqpoint{2.415159in}{0.598741in}}%
\pgfpathlineto{\pgfqpoint{2.413904in}{0.601638in}}%
\pgfpathlineto{\pgfqpoint{2.418230in}{0.606839in}}%
\pgfpathlineto{\pgfqpoint{2.416778in}{0.611640in}}%
\pgfpathlineto{\pgfqpoint{2.423706in}{0.617445in}}%
\pgfpathlineto{\pgfqpoint{2.422704in}{0.621971in}}%
\pgfpathlineto{\pgfqpoint{2.481583in}{0.625736in}}%
\pgfpathlineto{\pgfqpoint{2.488716in}{0.626214in}}%
\pgfpathlineto{\pgfqpoint{2.492884in}{0.596677in}}%
\pgfpathlineto{\pgfqpoint{2.488025in}{0.585155in}}%
\pgfpathlineto{\pgfqpoint{2.487697in}{0.577489in}}%
\pgfpathlineto{\pgfqpoint{2.485231in}{0.574838in}}%
\pgfpathlineto{\pgfqpoint{2.479736in}{0.576043in}}%
\pgfpathlineto{\pgfqpoint{2.476107in}{0.580682in}}%
\pgfpathlineto{\pgfqpoint{2.471783in}{0.578648in}}%
\pgfpathlineto{\pgfqpoint{2.470809in}{0.572873in}}%
\pgfpathlineto{\pgfqpoint{2.446412in}{0.569336in}}%
\pgfpathlineto{\pgfqpoint{2.444618in}{0.558566in}}%
\pgfpathlineto{\pgfqpoint{2.442205in}{0.558231in}}%
\pgfpathclose%
\pgfusepath{fill}%
\end{pgfscope}%
\begin{pgfscope}%
\pgfpathrectangle{\pgfqpoint{0.100000in}{0.100000in}}{\pgfqpoint{3.007045in}{1.925000in}}%
\pgfusepath{clip}%
\pgfsetbuttcap%
\pgfsetmiterjoin%
\definecolor{currentfill}{rgb}{0.454118,0.701300,0.846659}%
\pgfsetfillcolor{currentfill}%
\pgfsetlinewidth{0.000000pt}%
\definecolor{currentstroke}{rgb}{0.000000,0.000000,0.000000}%
\pgfsetstrokecolor{currentstroke}%
\pgfsetstrokeopacity{0.000000}%
\pgfsetdash{}{0pt}%
\pgfpathmoveto{\pgfqpoint{2.135867in}{0.998378in}}%
\pgfpathlineto{\pgfqpoint{2.141030in}{0.985137in}}%
\pgfpathlineto{\pgfqpoint{2.140155in}{0.982191in}}%
\pgfpathlineto{\pgfqpoint{2.117091in}{0.980740in}}%
\pgfpathlineto{\pgfqpoint{2.116101in}{0.997143in}}%
\pgfpathlineto{\pgfqpoint{2.135867in}{0.998378in}}%
\pgfpathclose%
\pgfusepath{fill}%
\end{pgfscope}%
\begin{pgfscope}%
\pgfpathrectangle{\pgfqpoint{0.100000in}{0.100000in}}{\pgfqpoint{3.007045in}{1.925000in}}%
\pgfusepath{clip}%
\pgfsetbuttcap%
\pgfsetmiterjoin%
\definecolor{currentfill}{rgb}{0.472941,0.711634,0.850719}%
\pgfsetfillcolor{currentfill}%
\pgfsetlinewidth{0.000000pt}%
\definecolor{currentstroke}{rgb}{0.000000,0.000000,0.000000}%
\pgfsetstrokecolor{currentstroke}%
\pgfsetstrokeopacity{0.000000}%
\pgfsetdash{}{0pt}%
\pgfpathmoveto{\pgfqpoint{1.861397in}{1.624038in}}%
\pgfpathlineto{\pgfqpoint{1.861813in}{1.606700in}}%
\pgfpathlineto{\pgfqpoint{1.844441in}{1.606089in}}%
\pgfpathlineto{\pgfqpoint{1.844395in}{1.612207in}}%
\pgfpathlineto{\pgfqpoint{1.826933in}{1.611742in}}%
\pgfpathlineto{\pgfqpoint{1.827033in}{1.605997in}}%
\pgfpathlineto{\pgfqpoint{1.801088in}{1.605771in}}%
\pgfpathlineto{\pgfqpoint{1.802129in}{1.614963in}}%
\pgfpathlineto{\pgfqpoint{1.798627in}{1.617392in}}%
\pgfpathlineto{\pgfqpoint{1.793675in}{1.615402in}}%
\pgfpathlineto{\pgfqpoint{1.782448in}{1.621431in}}%
\pgfpathlineto{\pgfqpoint{1.782023in}{1.648669in}}%
\pgfpathlineto{\pgfqpoint{1.787806in}{1.648648in}}%
\pgfpathlineto{\pgfqpoint{1.787190in}{1.689017in}}%
\pgfpathlineto{\pgfqpoint{1.798628in}{1.689186in}}%
\pgfpathlineto{\pgfqpoint{1.798346in}{1.718105in}}%
\pgfpathlineto{\pgfqpoint{1.827086in}{1.718431in}}%
\pgfpathlineto{\pgfqpoint{1.827059in}{1.721935in}}%
\pgfpathlineto{\pgfqpoint{1.858121in}{1.722001in}}%
\pgfpathlineto{\pgfqpoint{1.859255in}{1.710576in}}%
\pgfpathlineto{\pgfqpoint{1.859346in}{1.688805in}}%
\pgfpathlineto{\pgfqpoint{1.860483in}{1.664391in}}%
\pgfpathlineto{\pgfqpoint{1.860671in}{1.638052in}}%
\pgfpathlineto{\pgfqpoint{1.861397in}{1.624038in}}%
\pgfpathclose%
\pgfusepath{fill}%
\end{pgfscope}%
\begin{pgfscope}%
\pgfpathrectangle{\pgfqpoint{0.100000in}{0.100000in}}{\pgfqpoint{3.007045in}{1.925000in}}%
\pgfusepath{clip}%
\pgfsetbuttcap%
\pgfsetmiterjoin%
\definecolor{currentfill}{rgb}{0.321246,0.615179,0.800830}%
\pgfsetfillcolor{currentfill}%
\pgfsetlinewidth{0.000000pt}%
\definecolor{currentstroke}{rgb}{0.000000,0.000000,0.000000}%
\pgfsetstrokecolor{currentstroke}%
\pgfsetstrokeopacity{0.000000}%
\pgfsetdash{}{0pt}%
\pgfpathmoveto{\pgfqpoint{1.348840in}{1.549898in}}%
\pgfpathlineto{\pgfqpoint{1.303187in}{1.554906in}}%
\pgfpathlineto{\pgfqpoint{1.305764in}{1.578064in}}%
\pgfpathlineto{\pgfqpoint{1.308128in}{1.577801in}}%
\pgfpathlineto{\pgfqpoint{1.311526in}{1.606441in}}%
\pgfpathlineto{\pgfqpoint{1.313493in}{1.606222in}}%
\pgfpathlineto{\pgfqpoint{1.316018in}{1.629285in}}%
\pgfpathlineto{\pgfqpoint{1.318468in}{1.629024in}}%
\pgfpathlineto{\pgfqpoint{1.335389in}{1.627188in}}%
\pgfpathlineto{\pgfqpoint{1.334802in}{1.621794in}}%
\pgfpathlineto{\pgfqpoint{1.340576in}{1.621215in}}%
\pgfpathlineto{\pgfqpoint{1.339970in}{1.615629in}}%
\pgfpathlineto{\pgfqpoint{1.351714in}{1.614337in}}%
\pgfpathlineto{\pgfqpoint{1.355655in}{1.612365in}}%
\pgfpathlineto{\pgfqpoint{1.403892in}{1.607820in}}%
\pgfpathlineto{\pgfqpoint{1.406341in}{1.607594in}}%
\pgfpathlineto{\pgfqpoint{1.401374in}{1.559284in}}%
\pgfpathlineto{\pgfqpoint{1.351084in}{1.564019in}}%
\pgfpathlineto{\pgfqpoint{1.348840in}{1.549898in}}%
\pgfpathclose%
\pgfusepath{fill}%
\end{pgfscope}%
\begin{pgfscope}%
\pgfpathrectangle{\pgfqpoint{0.100000in}{0.100000in}}{\pgfqpoint{3.007045in}{1.925000in}}%
\pgfusepath{clip}%
\pgfsetbuttcap%
\pgfsetmiterjoin%
\definecolor{currentfill}{rgb}{0.642368,0.801830,0.890319}%
\pgfsetfillcolor{currentfill}%
\pgfsetlinewidth{0.000000pt}%
\definecolor{currentstroke}{rgb}{0.000000,0.000000,0.000000}%
\pgfsetstrokecolor{currentstroke}%
\pgfsetstrokeopacity{0.000000}%
\pgfsetdash{}{0pt}%
\pgfpathmoveto{\pgfqpoint{2.199018in}{1.471645in}}%
\pgfpathlineto{\pgfqpoint{2.180375in}{1.470113in}}%
\pgfpathlineto{\pgfqpoint{2.177932in}{1.479844in}}%
\pgfpathlineto{\pgfqpoint{2.175073in}{1.484492in}}%
\pgfpathlineto{\pgfqpoint{2.178913in}{1.490543in}}%
\pgfpathlineto{\pgfqpoint{2.184775in}{1.505418in}}%
\pgfpathlineto{\pgfqpoint{2.185259in}{1.517003in}}%
\pgfpathlineto{\pgfqpoint{2.204865in}{1.518579in}}%
\pgfpathlineto{\pgfqpoint{2.207015in}{1.495613in}}%
\pgfpathlineto{\pgfqpoint{2.196560in}{1.494814in}}%
\pgfpathlineto{\pgfqpoint{2.199018in}{1.471645in}}%
\pgfpathclose%
\pgfusepath{fill}%
\end{pgfscope}%
\begin{pgfscope}%
\pgfpathrectangle{\pgfqpoint{0.100000in}{0.100000in}}{\pgfqpoint{3.007045in}{1.925000in}}%
\pgfusepath{clip}%
\pgfsetbuttcap%
\pgfsetmiterjoin%
\definecolor{currentfill}{rgb}{0.084045,0.385067,0.664437}%
\pgfsetfillcolor{currentfill}%
\pgfsetlinewidth{0.000000pt}%
\definecolor{currentstroke}{rgb}{0.000000,0.000000,0.000000}%
\pgfsetstrokecolor{currentstroke}%
\pgfsetstrokeopacity{0.000000}%
\pgfsetdash{}{0pt}%
\pgfpathmoveto{\pgfqpoint{1.011802in}{1.226459in}}%
\pgfpathlineto{\pgfqpoint{1.010713in}{1.219317in}}%
\pgfpathlineto{\pgfqpoint{1.007174in}{1.215093in}}%
\pgfpathlineto{\pgfqpoint{1.007785in}{1.211607in}}%
\pgfpathlineto{\pgfqpoint{1.001745in}{1.203071in}}%
\pgfpathlineto{\pgfqpoint{0.998322in}{1.188795in}}%
\pgfpathlineto{\pgfqpoint{1.000947in}{1.182123in}}%
\pgfpathlineto{\pgfqpoint{0.999346in}{1.179674in}}%
\pgfpathlineto{\pgfqpoint{1.002802in}{1.166666in}}%
\pgfpathlineto{\pgfqpoint{1.001312in}{1.163987in}}%
\pgfpathlineto{\pgfqpoint{0.968213in}{1.169632in}}%
\pgfpathlineto{\pgfqpoint{0.936286in}{1.175309in}}%
\pgfpathlineto{\pgfqpoint{0.936407in}{1.175969in}}%
\pgfpathlineto{\pgfqpoint{0.942870in}{1.209747in}}%
\pgfpathlineto{\pgfqpoint{0.914996in}{1.214937in}}%
\pgfpathlineto{\pgfqpoint{0.907070in}{1.217261in}}%
\pgfpathlineto{\pgfqpoint{0.910572in}{1.236135in}}%
\pgfpathlineto{\pgfqpoint{0.915780in}{1.238145in}}%
\pgfpathlineto{\pgfqpoint{0.924573in}{1.236589in}}%
\pgfpathlineto{\pgfqpoint{0.927198in}{1.241278in}}%
\pgfpathlineto{\pgfqpoint{0.929754in}{1.256515in}}%
\pgfpathlineto{\pgfqpoint{0.937487in}{1.261054in}}%
\pgfpathlineto{\pgfqpoint{0.934985in}{1.263476in}}%
\pgfpathlineto{\pgfqpoint{0.974041in}{1.256437in}}%
\pgfpathlineto{\pgfqpoint{0.990893in}{1.253051in}}%
\pgfpathlineto{\pgfqpoint{1.022492in}{1.247657in}}%
\pgfpathlineto{\pgfqpoint{1.016300in}{1.244056in}}%
\pgfpathlineto{\pgfqpoint{1.015989in}{1.239279in}}%
\pgfpathlineto{\pgfqpoint{1.012248in}{1.233395in}}%
\pgfpathlineto{\pgfqpoint{1.011802in}{1.226459in}}%
\pgfpathclose%
\pgfusepath{fill}%
\end{pgfscope}%
\begin{pgfscope}%
\pgfpathrectangle{\pgfqpoint{0.100000in}{0.100000in}}{\pgfqpoint{3.007045in}{1.925000in}}%
\pgfusepath{clip}%
\pgfsetbuttcap%
\pgfsetmiterjoin%
\definecolor{currentfill}{rgb}{0.706344,0.829020,0.912710}%
\pgfsetfillcolor{currentfill}%
\pgfsetlinewidth{0.000000pt}%
\definecolor{currentstroke}{rgb}{0.000000,0.000000,0.000000}%
\pgfsetstrokecolor{currentstroke}%
\pgfsetstrokeopacity{0.000000}%
\pgfsetdash{}{0pt}%
\pgfpathmoveto{\pgfqpoint{0.397486in}{1.548285in}}%
\pgfpathlineto{\pgfqpoint{0.393613in}{1.542963in}}%
\pgfpathlineto{\pgfqpoint{0.385318in}{1.540362in}}%
\pgfpathlineto{\pgfqpoint{0.384988in}{1.532462in}}%
\pgfpathlineto{\pgfqpoint{0.380582in}{1.527032in}}%
\pgfpathlineto{\pgfqpoint{0.381897in}{1.521054in}}%
\pgfpathlineto{\pgfqpoint{0.381028in}{1.514238in}}%
\pgfpathlineto{\pgfqpoint{0.378145in}{1.511673in}}%
\pgfpathlineto{\pgfqpoint{0.372999in}{1.513265in}}%
\pgfpathlineto{\pgfqpoint{0.374697in}{1.518495in}}%
\pgfpathlineto{\pgfqpoint{0.360884in}{1.523094in}}%
\pgfpathlineto{\pgfqpoint{0.362297in}{1.539826in}}%
\pgfpathlineto{\pgfqpoint{0.358579in}{1.545963in}}%
\pgfpathlineto{\pgfqpoint{0.364222in}{1.554694in}}%
\pgfpathlineto{\pgfqpoint{0.365111in}{1.559248in}}%
\pgfpathlineto{\pgfqpoint{0.360771in}{1.567349in}}%
\pgfpathlineto{\pgfqpoint{0.362017in}{1.572568in}}%
\pgfpathlineto{\pgfqpoint{0.361686in}{1.582534in}}%
\pgfpathlineto{\pgfqpoint{0.365712in}{1.592867in}}%
\pgfpathlineto{\pgfqpoint{0.368947in}{1.597569in}}%
\pgfpathlineto{\pgfqpoint{0.369577in}{1.603351in}}%
\pgfpathlineto{\pgfqpoint{0.366697in}{1.609297in}}%
\pgfpathlineto{\pgfqpoint{0.367089in}{1.616757in}}%
\pgfpathlineto{\pgfqpoint{0.372778in}{1.622728in}}%
\pgfpathlineto{\pgfqpoint{0.382988in}{1.619317in}}%
\pgfpathlineto{\pgfqpoint{0.386071in}{1.611691in}}%
\pgfpathlineto{\pgfqpoint{0.382499in}{1.599727in}}%
\pgfpathlineto{\pgfqpoint{0.389344in}{1.598489in}}%
\pgfpathlineto{\pgfqpoint{0.397064in}{1.602137in}}%
\pgfpathlineto{\pgfqpoint{0.404299in}{1.600395in}}%
\pgfpathlineto{\pgfqpoint{0.395865in}{1.592654in}}%
\pgfpathlineto{\pgfqpoint{0.390435in}{1.585257in}}%
\pgfpathlineto{\pgfqpoint{0.384929in}{1.586076in}}%
\pgfpathlineto{\pgfqpoint{0.381248in}{1.578474in}}%
\pgfpathlineto{\pgfqpoint{0.387396in}{1.575981in}}%
\pgfpathlineto{\pgfqpoint{0.389591in}{1.567203in}}%
\pgfpathlineto{\pgfqpoint{0.384587in}{1.562231in}}%
\pgfpathlineto{\pgfqpoint{0.383232in}{1.552564in}}%
\pgfpathlineto{\pgfqpoint{0.397486in}{1.548285in}}%
\pgfpathclose%
\pgfusepath{fill}%
\end{pgfscope}%
\begin{pgfscope}%
\pgfpathrectangle{\pgfqpoint{0.100000in}{0.100000in}}{\pgfqpoint{3.007045in}{1.925000in}}%
\pgfusepath{clip}%
\pgfsetbuttcap%
\pgfsetmiterjoin%
\definecolor{currentfill}{rgb}{0.541961,0.749527,0.865606}%
\pgfsetfillcolor{currentfill}%
\pgfsetlinewidth{0.000000pt}%
\definecolor{currentstroke}{rgb}{0.000000,0.000000,0.000000}%
\pgfsetstrokecolor{currentstroke}%
\pgfsetstrokeopacity{0.000000}%
\pgfsetdash{}{0pt}%
\pgfpathmoveto{\pgfqpoint{2.283471in}{1.248041in}}%
\pgfpathlineto{\pgfqpoint{2.285375in}{1.228022in}}%
\pgfpathlineto{\pgfqpoint{2.265786in}{1.225833in}}%
\pgfpathlineto{\pgfqpoint{2.266331in}{1.211577in}}%
\pgfpathlineto{\pgfqpoint{2.262289in}{1.211080in}}%
\pgfpathlineto{\pgfqpoint{2.247429in}{1.209448in}}%
\pgfpathlineto{\pgfqpoint{2.247375in}{1.219993in}}%
\pgfpathlineto{\pgfqpoint{2.233016in}{1.218375in}}%
\pgfpathlineto{\pgfqpoint{2.231222in}{1.236425in}}%
\pgfpathlineto{\pgfqpoint{2.228217in}{1.264785in}}%
\pgfpathlineto{\pgfqpoint{2.239459in}{1.266008in}}%
\pgfpathlineto{\pgfqpoint{2.248890in}{1.267034in}}%
\pgfpathlineto{\pgfqpoint{2.249581in}{1.261349in}}%
\pgfpathlineto{\pgfqpoint{2.261811in}{1.262682in}}%
\pgfpathlineto{\pgfqpoint{2.262780in}{1.245532in}}%
\pgfpathlineto{\pgfqpoint{2.283471in}{1.248041in}}%
\pgfpathclose%
\pgfusepath{fill}%
\end{pgfscope}%
\begin{pgfscope}%
\pgfpathrectangle{\pgfqpoint{0.100000in}{0.100000in}}{\pgfqpoint{3.007045in}{1.925000in}}%
\pgfusepath{clip}%
\pgfsetbuttcap%
\pgfsetmiterjoin%
\definecolor{currentfill}{rgb}{0.554510,0.756417,0.868312}%
\pgfsetfillcolor{currentfill}%
\pgfsetlinewidth{0.000000pt}%
\definecolor{currentstroke}{rgb}{0.000000,0.000000,0.000000}%
\pgfsetstrokecolor{currentstroke}%
\pgfsetstrokeopacity{0.000000}%
\pgfsetdash{}{0pt}%
\pgfpathmoveto{\pgfqpoint{2.358324in}{1.060460in}}%
\pgfpathlineto{\pgfqpoint{2.350874in}{1.053026in}}%
\pgfpathlineto{\pgfqpoint{2.347601in}{1.052815in}}%
\pgfpathlineto{\pgfqpoint{2.340592in}{1.057330in}}%
\pgfpathlineto{\pgfqpoint{2.334876in}{1.050418in}}%
\pgfpathlineto{\pgfqpoint{2.333634in}{1.044388in}}%
\pgfpathlineto{\pgfqpoint{2.321129in}{1.054650in}}%
\pgfpathlineto{\pgfqpoint{2.322848in}{1.065287in}}%
\pgfpathlineto{\pgfqpoint{2.327579in}{1.069371in}}%
\pgfpathlineto{\pgfqpoint{2.316608in}{1.083386in}}%
\pgfpathlineto{\pgfqpoint{2.328598in}{1.094602in}}%
\pgfpathlineto{\pgfqpoint{2.342975in}{1.090758in}}%
\pgfpathlineto{\pgfqpoint{2.350883in}{1.086320in}}%
\pgfpathlineto{\pgfqpoint{2.358589in}{1.084899in}}%
\pgfpathlineto{\pgfqpoint{2.353799in}{1.080602in}}%
\pgfpathlineto{\pgfqpoint{2.350488in}{1.072384in}}%
\pgfpathlineto{\pgfqpoint{2.352145in}{1.068229in}}%
\pgfpathlineto{\pgfqpoint{2.358324in}{1.060460in}}%
\pgfpathclose%
\pgfusepath{fill}%
\end{pgfscope}%
\begin{pgfscope}%
\pgfpathrectangle{\pgfqpoint{0.100000in}{0.100000in}}{\pgfqpoint{3.007045in}{1.925000in}}%
\pgfusepath{clip}%
\pgfsetbuttcap%
\pgfsetmiterjoin%
\definecolor{currentfill}{rgb}{0.401953,0.670296,0.832326}%
\pgfsetfillcolor{currentfill}%
\pgfsetlinewidth{0.000000pt}%
\definecolor{currentstroke}{rgb}{0.000000,0.000000,0.000000}%
\pgfsetstrokecolor{currentstroke}%
\pgfsetstrokeopacity{0.000000}%
\pgfsetdash{}{0pt}%
\pgfpathmoveto{\pgfqpoint{1.835531in}{0.660856in}}%
\pgfpathlineto{\pgfqpoint{1.860218in}{0.661283in}}%
\pgfpathlineto{\pgfqpoint{1.860335in}{0.652602in}}%
\pgfpathlineto{\pgfqpoint{1.866139in}{0.652691in}}%
\pgfpathlineto{\pgfqpoint{1.866438in}{0.640996in}}%
\pgfpathlineto{\pgfqpoint{1.872143in}{0.641119in}}%
\pgfpathlineto{\pgfqpoint{1.872399in}{0.629573in}}%
\pgfpathlineto{\pgfqpoint{1.863641in}{0.629386in}}%
\pgfpathlineto{\pgfqpoint{1.860787in}{0.623525in}}%
\pgfpathlineto{\pgfqpoint{1.855514in}{0.623416in}}%
\pgfpathlineto{\pgfqpoint{1.854631in}{0.617364in}}%
\pgfpathlineto{\pgfqpoint{1.851931in}{0.616696in}}%
\pgfpathlineto{\pgfqpoint{1.846842in}{0.625571in}}%
\pgfpathlineto{\pgfqpoint{1.847519in}{0.631354in}}%
\pgfpathlineto{\pgfqpoint{1.843180in}{0.636011in}}%
\pgfpathlineto{\pgfqpoint{1.837976in}{0.644603in}}%
\pgfpathlineto{\pgfqpoint{1.839877in}{0.650274in}}%
\pgfpathlineto{\pgfqpoint{1.838775in}{0.656201in}}%
\pgfpathlineto{\pgfqpoint{1.835531in}{0.660856in}}%
\pgfpathclose%
\pgfusepath{fill}%
\end{pgfscope}%
\begin{pgfscope}%
\pgfpathrectangle{\pgfqpoint{0.100000in}{0.100000in}}{\pgfqpoint{3.007045in}{1.925000in}}%
\pgfusepath{clip}%
\pgfsetbuttcap%
\pgfsetmiterjoin%
\definecolor{currentfill}{rgb}{0.321246,0.615179,0.800830}%
\pgfsetfillcolor{currentfill}%
\pgfsetlinewidth{0.000000pt}%
\definecolor{currentstroke}{rgb}{0.000000,0.000000,0.000000}%
\pgfsetstrokecolor{currentstroke}%
\pgfsetstrokeopacity{0.000000}%
\pgfsetdash{}{0pt}%
\pgfpathmoveto{\pgfqpoint{1.460453in}{1.255605in}}%
\pgfpathlineto{\pgfqpoint{1.420426in}{1.258524in}}%
\pgfpathlineto{\pgfqpoint{1.420694in}{1.261942in}}%
\pgfpathlineto{\pgfqpoint{1.390844in}{1.264339in}}%
\pgfpathlineto{\pgfqpoint{1.392134in}{1.280968in}}%
\pgfpathlineto{\pgfqpoint{1.394532in}{1.283801in}}%
\pgfpathlineto{\pgfqpoint{1.395479in}{1.295197in}}%
\pgfpathlineto{\pgfqpoint{1.422895in}{1.292969in}}%
\pgfpathlineto{\pgfqpoint{1.423796in}{1.304378in}}%
\pgfpathlineto{\pgfqpoint{1.427276in}{1.304079in}}%
\pgfpathlineto{\pgfqpoint{1.455754in}{1.302023in}}%
\pgfpathlineto{\pgfqpoint{1.462537in}{1.301517in}}%
\pgfpathlineto{\pgfqpoint{1.460924in}{1.278668in}}%
\pgfpathlineto{\pgfqpoint{1.461934in}{1.278597in}}%
\pgfpathlineto{\pgfqpoint{1.460453in}{1.255605in}}%
\pgfpathclose%
\pgfusepath{fill}%
\end{pgfscope}%
\begin{pgfscope}%
\pgfpathrectangle{\pgfqpoint{0.100000in}{0.100000in}}{\pgfqpoint{3.007045in}{1.925000in}}%
\pgfusepath{clip}%
\pgfsetbuttcap%
\pgfsetmiterjoin%
\definecolor{currentfill}{rgb}{0.306113,0.604844,0.794925}%
\pgfsetfillcolor{currentfill}%
\pgfsetlinewidth{0.000000pt}%
\definecolor{currentstroke}{rgb}{0.000000,0.000000,0.000000}%
\pgfsetstrokecolor{currentstroke}%
\pgfsetstrokeopacity{0.000000}%
\pgfsetdash{}{0pt}%
\pgfpathmoveto{\pgfqpoint{2.191739in}{0.996087in}}%
\pgfpathlineto{\pgfqpoint{2.189732in}{1.023554in}}%
\pgfpathlineto{\pgfqpoint{2.200660in}{1.024227in}}%
\pgfpathlineto{\pgfqpoint{2.209941in}{1.021313in}}%
\pgfpathlineto{\pgfqpoint{2.212201in}{1.031659in}}%
\pgfpathlineto{\pgfqpoint{2.216520in}{1.033633in}}%
\pgfpathlineto{\pgfqpoint{2.223300in}{1.033829in}}%
\pgfpathlineto{\pgfqpoint{2.218727in}{1.043400in}}%
\pgfpathlineto{\pgfqpoint{2.234875in}{1.045920in}}%
\pgfpathlineto{\pgfqpoint{2.241309in}{1.038598in}}%
\pgfpathlineto{\pgfqpoint{2.241044in}{1.027825in}}%
\pgfpathlineto{\pgfqpoint{2.236965in}{1.019520in}}%
\pgfpathlineto{\pgfqpoint{2.247870in}{1.007837in}}%
\pgfpathlineto{\pgfqpoint{2.248936in}{1.000483in}}%
\pgfpathlineto{\pgfqpoint{2.230102in}{0.999971in}}%
\pgfpathlineto{\pgfqpoint{2.216487in}{0.998935in}}%
\pgfpathlineto{\pgfqpoint{2.191739in}{0.996087in}}%
\pgfpathclose%
\pgfusepath{fill}%
\end{pgfscope}%
\begin{pgfscope}%
\pgfpathrectangle{\pgfqpoint{0.100000in}{0.100000in}}{\pgfqpoint{3.007045in}{1.925000in}}%
\pgfusepath{clip}%
\pgfsetbuttcap%
\pgfsetmiterjoin%
\definecolor{currentfill}{rgb}{0.142607,0.456332,0.716601}%
\pgfsetfillcolor{currentfill}%
\pgfsetlinewidth{0.000000pt}%
\definecolor{currentstroke}{rgb}{0.000000,0.000000,0.000000}%
\pgfsetstrokecolor{currentstroke}%
\pgfsetstrokeopacity{0.000000}%
\pgfsetdash{}{0pt}%
\pgfpathmoveto{\pgfqpoint{0.829274in}{1.562729in}}%
\pgfpathlineto{\pgfqpoint{0.834605in}{1.565249in}}%
\pgfpathlineto{\pgfqpoint{0.836443in}{1.560436in}}%
\pgfpathlineto{\pgfqpoint{0.842609in}{1.559070in}}%
\pgfpathlineto{\pgfqpoint{0.849873in}{1.555251in}}%
\pgfpathlineto{\pgfqpoint{0.860657in}{1.551817in}}%
\pgfpathlineto{\pgfqpoint{0.861425in}{1.547079in}}%
\pgfpathlineto{\pgfqpoint{0.872520in}{1.537698in}}%
\pgfpathlineto{\pgfqpoint{0.878844in}{1.527019in}}%
\pgfpathlineto{\pgfqpoint{0.882414in}{1.527693in}}%
\pgfpathlineto{\pgfqpoint{0.884911in}{1.517908in}}%
\pgfpathlineto{\pgfqpoint{0.883755in}{1.512410in}}%
\pgfpathlineto{\pgfqpoint{0.896614in}{1.509734in}}%
\pgfpathlineto{\pgfqpoint{0.895520in}{1.504390in}}%
\pgfpathlineto{\pgfqpoint{0.912230in}{1.501011in}}%
\pgfpathlineto{\pgfqpoint{0.909954in}{1.489763in}}%
\pgfpathlineto{\pgfqpoint{0.898811in}{1.492021in}}%
\pgfpathlineto{\pgfqpoint{0.894533in}{1.469396in}}%
\pgfpathlineto{\pgfqpoint{0.896143in}{1.463222in}}%
\pgfpathlineto{\pgfqpoint{0.892728in}{1.460530in}}%
\pgfpathlineto{\pgfqpoint{0.887062in}{1.465893in}}%
\pgfpathlineto{\pgfqpoint{0.882125in}{1.465664in}}%
\pgfpathlineto{\pgfqpoint{0.884523in}{1.477298in}}%
\pgfpathlineto{\pgfqpoint{0.887326in}{1.476723in}}%
\pgfpathlineto{\pgfqpoint{0.891952in}{1.499400in}}%
\pgfpathlineto{\pgfqpoint{0.877767in}{1.502355in}}%
\pgfpathlineto{\pgfqpoint{0.873074in}{1.479709in}}%
\pgfpathlineto{\pgfqpoint{0.869518in}{1.474687in}}%
\pgfpathlineto{\pgfqpoint{0.861547in}{1.476417in}}%
\pgfpathlineto{\pgfqpoint{0.858463in}{1.461570in}}%
\pgfpathlineto{\pgfqpoint{0.851744in}{1.461596in}}%
\pgfpathlineto{\pgfqpoint{0.850305in}{1.455334in}}%
\pgfpathlineto{\pgfqpoint{0.840009in}{1.457550in}}%
\pgfpathlineto{\pgfqpoint{0.834301in}{1.430229in}}%
\pgfpathlineto{\pgfqpoint{0.808116in}{1.436493in}}%
\pgfpathlineto{\pgfqpoint{0.798116in}{1.438303in}}%
\pgfpathlineto{\pgfqpoint{0.811385in}{1.497272in}}%
\pgfpathlineto{\pgfqpoint{0.809080in}{1.497950in}}%
\pgfpathlineto{\pgfqpoint{0.819243in}{1.542361in}}%
\pgfpathlineto{\pgfqpoint{0.824049in}{1.547593in}}%
\pgfpathlineto{\pgfqpoint{0.823175in}{1.550976in}}%
\pgfpathlineto{\pgfqpoint{0.828248in}{1.558524in}}%
\pgfpathlineto{\pgfqpoint{0.829274in}{1.562729in}}%
\pgfpathclose%
\pgfusepath{fill}%
\end{pgfscope}%
\begin{pgfscope}%
\pgfpathrectangle{\pgfqpoint{0.100000in}{0.100000in}}{\pgfqpoint{3.007045in}{1.925000in}}%
\pgfusepath{clip}%
\pgfsetbuttcap%
\pgfsetmiterjoin%
\definecolor{currentfill}{rgb}{0.696501,0.824837,0.909266}%
\pgfsetfillcolor{currentfill}%
\pgfsetlinewidth{0.000000pt}%
\definecolor{currentstroke}{rgb}{0.000000,0.000000,0.000000}%
\pgfsetstrokecolor{currentstroke}%
\pgfsetstrokeopacity{0.000000}%
\pgfsetdash{}{0pt}%
\pgfpathmoveto{\pgfqpoint{0.429023in}{1.481634in}}%
\pgfpathlineto{\pgfqpoint{0.424324in}{1.471482in}}%
\pgfpathlineto{\pgfqpoint{0.420836in}{1.470570in}}%
\pgfpathlineto{\pgfqpoint{0.412778in}{1.458416in}}%
\pgfpathlineto{\pgfqpoint{0.407285in}{1.452726in}}%
\pgfpathlineto{\pgfqpoint{0.407911in}{1.446662in}}%
\pgfpathlineto{\pgfqpoint{0.399108in}{1.444237in}}%
\pgfpathlineto{\pgfqpoint{0.395246in}{1.441320in}}%
\pgfpathlineto{\pgfqpoint{0.390077in}{1.441086in}}%
\pgfpathlineto{\pgfqpoint{0.385422in}{1.437519in}}%
\pgfpathlineto{\pgfqpoint{0.384276in}{1.433228in}}%
\pgfpathlineto{\pgfqpoint{0.387316in}{1.430255in}}%
\pgfpathlineto{\pgfqpoint{0.385058in}{1.423808in}}%
\pgfpathlineto{\pgfqpoint{0.384397in}{1.412475in}}%
\pgfpathlineto{\pgfqpoint{0.355032in}{1.421828in}}%
\pgfpathlineto{\pgfqpoint{0.355543in}{1.423387in}}%
\pgfpathlineto{\pgfqpoint{0.332528in}{1.430977in}}%
\pgfpathlineto{\pgfqpoint{0.330534in}{1.438279in}}%
\pgfpathlineto{\pgfqpoint{0.327431in}{1.441471in}}%
\pgfpathlineto{\pgfqpoint{0.321640in}{1.452224in}}%
\pgfpathlineto{\pgfqpoint{0.323491in}{1.455583in}}%
\pgfpathlineto{\pgfqpoint{0.323252in}{1.464314in}}%
\pgfpathlineto{\pgfqpoint{0.332049in}{1.475470in}}%
\pgfpathlineto{\pgfqpoint{0.345804in}{1.489825in}}%
\pgfpathlineto{\pgfqpoint{0.349534in}{1.496541in}}%
\pgfpathlineto{\pgfqpoint{0.348214in}{1.499832in}}%
\pgfpathlineto{\pgfqpoint{0.355538in}{1.511740in}}%
\pgfpathlineto{\pgfqpoint{0.360884in}{1.523094in}}%
\pgfpathlineto{\pgfqpoint{0.374697in}{1.518495in}}%
\pgfpathlineto{\pgfqpoint{0.372999in}{1.513265in}}%
\pgfpathlineto{\pgfqpoint{0.378145in}{1.511673in}}%
\pgfpathlineto{\pgfqpoint{0.385689in}{1.509253in}}%
\pgfpathlineto{\pgfqpoint{0.386307in}{1.493716in}}%
\pgfpathlineto{\pgfqpoint{0.391169in}{1.489496in}}%
\pgfpathlineto{\pgfqpoint{0.391596in}{1.486250in}}%
\pgfpathlineto{\pgfqpoint{0.397986in}{1.484025in}}%
\pgfpathlineto{\pgfqpoint{0.399966in}{1.478464in}}%
\pgfpathlineto{\pgfqpoint{0.407088in}{1.477930in}}%
\pgfpathlineto{\pgfqpoint{0.405051in}{1.484139in}}%
\pgfpathlineto{\pgfqpoint{0.408114in}{1.487902in}}%
\pgfpathlineto{\pgfqpoint{0.414605in}{1.487528in}}%
\pgfpathlineto{\pgfqpoint{0.428385in}{1.494011in}}%
\pgfpathlineto{\pgfqpoint{0.431912in}{1.491802in}}%
\pgfpathlineto{\pgfqpoint{0.429023in}{1.481634in}}%
\pgfpathclose%
\pgfusepath{fill}%
\end{pgfscope}%
\begin{pgfscope}%
\pgfpathrectangle{\pgfqpoint{0.100000in}{0.100000in}}{\pgfqpoint{3.007045in}{1.925000in}}%
\pgfusepath{clip}%
\pgfsetbuttcap%
\pgfsetmiterjoin%
\definecolor{currentfill}{rgb}{0.401953,0.670296,0.832326}%
\pgfsetfillcolor{currentfill}%
\pgfsetlinewidth{0.000000pt}%
\definecolor{currentstroke}{rgb}{0.000000,0.000000,0.000000}%
\pgfsetstrokecolor{currentstroke}%
\pgfsetstrokeopacity{0.000000}%
\pgfsetdash{}{0pt}%
\pgfpathmoveto{\pgfqpoint{1.889132in}{1.155695in}}%
\pgfpathlineto{\pgfqpoint{1.884869in}{1.157037in}}%
\pgfpathlineto{\pgfqpoint{1.882047in}{1.149212in}}%
\pgfpathlineto{\pgfqpoint{1.877817in}{1.149309in}}%
\pgfpathlineto{\pgfqpoint{1.875772in}{1.154964in}}%
\pgfpathlineto{\pgfqpoint{1.870359in}{1.156300in}}%
\pgfpathlineto{\pgfqpoint{1.867949in}{1.161555in}}%
\pgfpathlineto{\pgfqpoint{1.862920in}{1.162275in}}%
\pgfpathlineto{\pgfqpoint{1.859930in}{1.165889in}}%
\pgfpathlineto{\pgfqpoint{1.860786in}{1.170457in}}%
\pgfpathlineto{\pgfqpoint{1.859241in}{1.178993in}}%
\pgfpathlineto{\pgfqpoint{1.855189in}{1.180142in}}%
\pgfpathlineto{\pgfqpoint{1.855015in}{1.197513in}}%
\pgfpathlineto{\pgfqpoint{1.854945in}{1.201835in}}%
\pgfpathlineto{\pgfqpoint{1.880439in}{1.202637in}}%
\pgfpathlineto{\pgfqpoint{1.880879in}{1.202651in}}%
\pgfpathlineto{\pgfqpoint{1.880876in}{1.180425in}}%
\pgfpathlineto{\pgfqpoint{1.889421in}{1.180509in}}%
\pgfpathlineto{\pgfqpoint{1.889132in}{1.155695in}}%
\pgfpathclose%
\pgfusepath{fill}%
\end{pgfscope}%
\begin{pgfscope}%
\pgfpathrectangle{\pgfqpoint{0.100000in}{0.100000in}}{\pgfqpoint{3.007045in}{1.925000in}}%
\pgfusepath{clip}%
\pgfsetbuttcap%
\pgfsetmiterjoin%
\definecolor{currentfill}{rgb}{0.248166,0.561892,0.770980}%
\pgfsetfillcolor{currentfill}%
\pgfsetlinewidth{0.000000pt}%
\definecolor{currentstroke}{rgb}{0.000000,0.000000,0.000000}%
\pgfsetstrokecolor{currentstroke}%
\pgfsetstrokeopacity{0.000000}%
\pgfsetdash{}{0pt}%
\pgfpathmoveto{\pgfqpoint{1.564310in}{0.846924in}}%
\pgfpathlineto{\pgfqpoint{1.564075in}{0.841182in}}%
\pgfpathlineto{\pgfqpoint{1.573014in}{0.840820in}}%
\pgfpathlineto{\pgfqpoint{1.572553in}{0.829339in}}%
\pgfpathlineto{\pgfqpoint{1.575401in}{0.829239in}}%
\pgfpathlineto{\pgfqpoint{1.574901in}{0.817578in}}%
\pgfpathlineto{\pgfqpoint{1.570469in}{0.816172in}}%
\pgfpathlineto{\pgfqpoint{1.560430in}{0.818079in}}%
\pgfpathlineto{\pgfqpoint{1.556351in}{0.822025in}}%
\pgfpathlineto{\pgfqpoint{1.543414in}{0.822716in}}%
\pgfpathlineto{\pgfqpoint{1.542619in}{0.830859in}}%
\pgfpathlineto{\pgfqpoint{1.542518in}{0.841248in}}%
\pgfpathlineto{\pgfqpoint{1.545976in}{0.847321in}}%
\pgfpathlineto{\pgfqpoint{1.564310in}{0.846924in}}%
\pgfpathclose%
\pgfusepath{fill}%
\end{pgfscope}%
\begin{pgfscope}%
\pgfpathrectangle{\pgfqpoint{0.100000in}{0.100000in}}{\pgfqpoint{3.007045in}{1.925000in}}%
\pgfusepath{clip}%
\pgfsetbuttcap%
\pgfsetmiterjoin%
\definecolor{currentfill}{rgb}{0.260715,0.573841,0.777209}%
\pgfsetfillcolor{currentfill}%
\pgfsetlinewidth{0.000000pt}%
\definecolor{currentstroke}{rgb}{0.000000,0.000000,0.000000}%
\pgfsetstrokecolor{currentstroke}%
\pgfsetstrokeopacity{0.000000}%
\pgfsetdash{}{0pt}%
\pgfpathmoveto{\pgfqpoint{1.721870in}{1.430045in}}%
\pgfpathlineto{\pgfqpoint{1.702699in}{1.430317in}}%
\pgfpathlineto{\pgfqpoint{1.703058in}{1.453389in}}%
\pgfpathlineto{\pgfqpoint{1.703851in}{1.499277in}}%
\pgfpathlineto{\pgfqpoint{1.703956in}{1.505078in}}%
\pgfpathlineto{\pgfqpoint{1.743883in}{1.504541in}}%
\pgfpathlineto{\pgfqpoint{1.744270in}{1.498734in}}%
\pgfpathlineto{\pgfqpoint{1.744188in}{1.475811in}}%
\pgfpathlineto{\pgfqpoint{1.721890in}{1.476069in}}%
\pgfpathlineto{\pgfqpoint{1.722176in}{1.453096in}}%
\pgfpathlineto{\pgfqpoint{1.721870in}{1.430045in}}%
\pgfpathclose%
\pgfusepath{fill}%
\end{pgfscope}%
\begin{pgfscope}%
\pgfpathrectangle{\pgfqpoint{0.100000in}{0.100000in}}{\pgfqpoint{3.007045in}{1.925000in}}%
\pgfusepath{clip}%
\pgfsetbuttcap%
\pgfsetmiterjoin%
\definecolor{currentfill}{rgb}{0.460392,0.704744,0.848012}%
\pgfsetfillcolor{currentfill}%
\pgfsetlinewidth{0.000000pt}%
\definecolor{currentstroke}{rgb}{0.000000,0.000000,0.000000}%
\pgfsetstrokecolor{currentstroke}%
\pgfsetstrokeopacity{0.000000}%
\pgfsetdash{}{0pt}%
\pgfpathmoveto{\pgfqpoint{1.234424in}{1.392114in}}%
\pgfpathlineto{\pgfqpoint{1.231571in}{1.369429in}}%
\pgfpathlineto{\pgfqpoint{1.226347in}{1.324123in}}%
\pgfpathlineto{\pgfqpoint{1.213860in}{1.325677in}}%
\pgfpathlineto{\pgfqpoint{1.210701in}{1.299766in}}%
\pgfpathlineto{\pgfqpoint{1.184259in}{1.303392in}}%
\pgfpathlineto{\pgfqpoint{1.131992in}{1.310519in}}%
\pgfpathlineto{\pgfqpoint{1.135759in}{1.336378in}}%
\pgfpathlineto{\pgfqpoint{1.137581in}{1.353507in}}%
\pgfpathlineto{\pgfqpoint{1.158048in}{1.350480in}}%
\pgfpathlineto{\pgfqpoint{1.162230in}{1.378562in}}%
\pgfpathlineto{\pgfqpoint{1.162889in}{1.390045in}}%
\pgfpathlineto{\pgfqpoint{1.164443in}{1.401325in}}%
\pgfpathlineto{\pgfqpoint{1.207752in}{1.395280in}}%
\pgfpathlineto{\pgfqpoint{1.234424in}{1.392114in}}%
\pgfpathclose%
\pgfusepath{fill}%
\end{pgfscope}%
\begin{pgfscope}%
\pgfpathrectangle{\pgfqpoint{0.100000in}{0.100000in}}{\pgfqpoint{3.007045in}{1.925000in}}%
\pgfusepath{clip}%
\pgfsetbuttcap%
\pgfsetmiterjoin%
\definecolor{currentfill}{rgb}{0.396909,0.666851,0.830358}%
\pgfsetfillcolor{currentfill}%
\pgfsetlinewidth{0.000000pt}%
\definecolor{currentstroke}{rgb}{0.000000,0.000000,0.000000}%
\pgfsetstrokecolor{currentstroke}%
\pgfsetstrokeopacity{0.000000}%
\pgfsetdash{}{0pt}%
\pgfpathmoveto{\pgfqpoint{1.652796in}{1.200602in}}%
\pgfpathlineto{\pgfqpoint{1.653431in}{1.223538in}}%
\pgfpathlineto{\pgfqpoint{1.654636in}{1.269418in}}%
\pgfpathlineto{\pgfqpoint{1.677475in}{1.268823in}}%
\pgfpathlineto{\pgfqpoint{1.699583in}{1.268350in}}%
\pgfpathlineto{\pgfqpoint{1.699296in}{1.251149in}}%
\pgfpathlineto{\pgfqpoint{1.698475in}{1.199567in}}%
\pgfpathlineto{\pgfqpoint{1.652796in}{1.200602in}}%
\pgfpathclose%
\pgfusepath{fill}%
\end{pgfscope}%
\begin{pgfscope}%
\pgfpathrectangle{\pgfqpoint{0.100000in}{0.100000in}}{\pgfqpoint{3.007045in}{1.925000in}}%
\pgfusepath{clip}%
\pgfsetbuttcap%
\pgfsetmiterjoin%
\definecolor{currentfill}{rgb}{0.441569,0.694410,0.843952}%
\pgfsetfillcolor{currentfill}%
\pgfsetlinewidth{0.000000pt}%
\definecolor{currentstroke}{rgb}{0.000000,0.000000,0.000000}%
\pgfsetstrokecolor{currentstroke}%
\pgfsetstrokeopacity{0.000000}%
\pgfsetdash{}{0pt}%
\pgfpathmoveto{\pgfqpoint{1.650048in}{1.679481in}}%
\pgfpathlineto{\pgfqpoint{1.638529in}{1.679874in}}%
\pgfpathlineto{\pgfqpoint{1.637567in}{1.685608in}}%
\pgfpathlineto{\pgfqpoint{1.638404in}{1.708645in}}%
\pgfpathlineto{\pgfqpoint{1.642987in}{1.708503in}}%
\pgfpathlineto{\pgfqpoint{1.643745in}{1.731622in}}%
\pgfpathlineto{\pgfqpoint{1.643166in}{1.743370in}}%
\pgfpathlineto{\pgfqpoint{1.677065in}{1.742344in}}%
\pgfpathlineto{\pgfqpoint{1.676789in}{1.740990in}}%
\pgfpathlineto{\pgfqpoint{1.705532in}{1.740473in}}%
\pgfpathlineto{\pgfqpoint{1.706015in}{1.718669in}}%
\pgfpathlineto{\pgfqpoint{1.711793in}{1.718564in}}%
\pgfpathlineto{\pgfqpoint{1.711752in}{1.712779in}}%
\pgfpathlineto{\pgfqpoint{1.734818in}{1.712519in}}%
\pgfpathlineto{\pgfqpoint{1.734882in}{1.718354in}}%
\pgfpathlineto{\pgfqpoint{1.740597in}{1.718328in}}%
\pgfpathlineto{\pgfqpoint{1.740624in}{1.724147in}}%
\pgfpathlineto{\pgfqpoint{1.746314in}{1.723984in}}%
\pgfpathlineto{\pgfqpoint{1.746352in}{1.706621in}}%
\pgfpathlineto{\pgfqpoint{1.747387in}{1.695072in}}%
\pgfpathlineto{\pgfqpoint{1.724288in}{1.695202in}}%
\pgfpathlineto{\pgfqpoint{1.689129in}{1.695752in}}%
\pgfpathlineto{\pgfqpoint{1.689494in}{1.678334in}}%
\pgfpathlineto{\pgfqpoint{1.650048in}{1.679481in}}%
\pgfpathclose%
\pgfusepath{fill}%
\end{pgfscope}%
\begin{pgfscope}%
\pgfpathrectangle{\pgfqpoint{0.100000in}{0.100000in}}{\pgfqpoint{3.007045in}{1.925000in}}%
\pgfusepath{clip}%
\pgfsetbuttcap%
\pgfsetmiterjoin%
\definecolor{currentfill}{rgb}{0.346467,0.632403,0.810673}%
\pgfsetfillcolor{currentfill}%
\pgfsetlinewidth{0.000000pt}%
\definecolor{currentstroke}{rgb}{0.000000,0.000000,0.000000}%
\pgfsetstrokecolor{currentstroke}%
\pgfsetstrokeopacity{0.000000}%
\pgfsetdash{}{0pt}%
\pgfpathmoveto{\pgfqpoint{1.767169in}{1.475679in}}%
\pgfpathlineto{\pgfqpoint{1.744188in}{1.475811in}}%
\pgfpathlineto{\pgfqpoint{1.744270in}{1.498734in}}%
\pgfpathlineto{\pgfqpoint{1.755330in}{1.498635in}}%
\pgfpathlineto{\pgfqpoint{1.755381in}{1.509019in}}%
\pgfpathlineto{\pgfqpoint{1.749628in}{1.512662in}}%
\pgfpathlineto{\pgfqpoint{1.749650in}{1.521778in}}%
\pgfpathlineto{\pgfqpoint{1.783670in}{1.521817in}}%
\pgfpathlineto{\pgfqpoint{1.795889in}{1.521880in}}%
\pgfpathlineto{\pgfqpoint{1.795936in}{1.510293in}}%
\pgfpathlineto{\pgfqpoint{1.789742in}{1.510236in}}%
\pgfpathlineto{\pgfqpoint{1.790073in}{1.492987in}}%
\pgfpathlineto{\pgfqpoint{1.782199in}{1.492987in}}%
\pgfpathlineto{\pgfqpoint{1.778660in}{1.495730in}}%
\pgfpathlineto{\pgfqpoint{1.778661in}{1.481448in}}%
\pgfpathlineto{\pgfqpoint{1.767240in}{1.481415in}}%
\pgfpathlineto{\pgfqpoint{1.767169in}{1.475679in}}%
\pgfpathclose%
\pgfusepath{fill}%
\end{pgfscope}%
\begin{pgfscope}%
\pgfpathrectangle{\pgfqpoint{0.100000in}{0.100000in}}{\pgfqpoint{3.007045in}{1.925000in}}%
\pgfusepath{clip}%
\pgfsetbuttcap%
\pgfsetmiterjoin%
\definecolor{currentfill}{rgb}{0.721107,0.835294,0.917878}%
\pgfsetfillcolor{currentfill}%
\pgfsetlinewidth{0.000000pt}%
\definecolor{currentstroke}{rgb}{0.000000,0.000000,0.000000}%
\pgfsetstrokecolor{currentstroke}%
\pgfsetstrokeopacity{0.000000}%
\pgfsetdash{}{0pt}%
\pgfpathmoveto{\pgfqpoint{1.948847in}{1.744739in}}%
\pgfpathlineto{\pgfqpoint{1.956874in}{1.748840in}}%
\pgfpathlineto{\pgfqpoint{1.961605in}{1.738804in}}%
\pgfpathlineto{\pgfqpoint{1.978390in}{1.739623in}}%
\pgfpathlineto{\pgfqpoint{1.988931in}{1.741585in}}%
\pgfpathlineto{\pgfqpoint{1.994032in}{1.740034in}}%
\pgfpathlineto{\pgfqpoint{1.995567in}{1.736342in}}%
\pgfpathlineto{\pgfqpoint{2.001096in}{1.733792in}}%
\pgfpathlineto{\pgfqpoint{2.005613in}{1.736504in}}%
\pgfpathlineto{\pgfqpoint{2.012380in}{1.736151in}}%
\pgfpathlineto{\pgfqpoint{2.011685in}{1.732318in}}%
\pgfpathlineto{\pgfqpoint{1.997324in}{1.723369in}}%
\pgfpathlineto{\pgfqpoint{1.987813in}{1.719063in}}%
\pgfpathlineto{\pgfqpoint{1.972465in}{1.713418in}}%
\pgfpathlineto{\pgfqpoint{1.963898in}{1.707804in}}%
\pgfpathlineto{\pgfqpoint{1.951342in}{1.696561in}}%
\pgfpathlineto{\pgfqpoint{1.950215in}{1.724734in}}%
\pgfpathlineto{\pgfqpoint{1.948847in}{1.744739in}}%
\pgfpathclose%
\pgfusepath{fill}%
\end{pgfscope}%
\begin{pgfscope}%
\pgfpathrectangle{\pgfqpoint{0.100000in}{0.100000in}}{\pgfqpoint{3.007045in}{1.925000in}}%
\pgfusepath{clip}%
\pgfsetbuttcap%
\pgfsetmiterjoin%
\definecolor{currentfill}{rgb}{0.406997,0.673741,0.834295}%
\pgfsetfillcolor{currentfill}%
\pgfsetlinewidth{0.000000pt}%
\definecolor{currentstroke}{rgb}{0.000000,0.000000,0.000000}%
\pgfsetstrokecolor{currentstroke}%
\pgfsetstrokeopacity{0.000000}%
\pgfsetdash{}{0pt}%
\pgfpathmoveto{\pgfqpoint{2.385587in}{1.155405in}}%
\pgfpathlineto{\pgfqpoint{2.380064in}{1.147699in}}%
\pgfpathlineto{\pgfqpoint{2.374431in}{1.146809in}}%
\pgfpathlineto{\pgfqpoint{2.370900in}{1.172765in}}%
\pgfpathlineto{\pgfqpoint{2.364743in}{1.174580in}}%
\pgfpathlineto{\pgfqpoint{2.365382in}{1.186478in}}%
\pgfpathlineto{\pgfqpoint{2.363339in}{1.187311in}}%
\pgfpathlineto{\pgfqpoint{2.362689in}{1.195710in}}%
\pgfpathlineto{\pgfqpoint{2.366865in}{1.205475in}}%
\pgfpathlineto{\pgfqpoint{2.381162in}{1.206723in}}%
\pgfpathlineto{\pgfqpoint{2.393794in}{1.205858in}}%
\pgfpathlineto{\pgfqpoint{2.394275in}{1.199273in}}%
\pgfpathlineto{\pgfqpoint{2.405917in}{1.200954in}}%
\pgfpathlineto{\pgfqpoint{2.405836in}{1.202691in}}%
\pgfpathlineto{\pgfqpoint{2.417290in}{1.203518in}}%
\pgfpathlineto{\pgfqpoint{2.419602in}{1.197771in}}%
\pgfpathlineto{\pgfqpoint{2.418057in}{1.191742in}}%
\pgfpathlineto{\pgfqpoint{2.418866in}{1.180027in}}%
\pgfpathlineto{\pgfqpoint{2.413102in}{1.179720in}}%
\pgfpathlineto{\pgfqpoint{2.413821in}{1.167145in}}%
\pgfpathlineto{\pgfqpoint{2.407704in}{1.166293in}}%
\pgfpathlineto{\pgfqpoint{2.403876in}{1.166074in}}%
\pgfpathlineto{\pgfqpoint{2.404127in}{1.159387in}}%
\pgfpathlineto{\pgfqpoint{2.401519in}{1.149007in}}%
\pgfpathlineto{\pgfqpoint{2.395723in}{1.148881in}}%
\pgfpathlineto{\pgfqpoint{2.392277in}{1.158355in}}%
\pgfpathlineto{\pgfqpoint{2.385587in}{1.155405in}}%
\pgfpathclose%
\pgfusepath{fill}%
\end{pgfscope}%
\begin{pgfscope}%
\pgfpathrectangle{\pgfqpoint{0.100000in}{0.100000in}}{\pgfqpoint{3.007045in}{1.925000in}}%
\pgfusepath{clip}%
\pgfsetbuttcap%
\pgfsetmiterjoin%
\definecolor{currentfill}{rgb}{0.406997,0.673741,0.834295}%
\pgfsetfillcolor{currentfill}%
\pgfsetlinewidth{0.000000pt}%
\definecolor{currentstroke}{rgb}{0.000000,0.000000,0.000000}%
\pgfsetstrokecolor{currentstroke}%
\pgfsetstrokeopacity{0.000000}%
\pgfsetdash{}{0pt}%
\pgfpathmoveto{\pgfqpoint{2.018995in}{1.370976in}}%
\pgfpathlineto{\pgfqpoint{2.020468in}{1.350608in}}%
\pgfpathlineto{\pgfqpoint{2.001078in}{1.349385in}}%
\pgfpathlineto{\pgfqpoint{1.997258in}{1.351410in}}%
\pgfpathlineto{\pgfqpoint{1.995792in}{1.357804in}}%
\pgfpathlineto{\pgfqpoint{1.992817in}{1.361377in}}%
\pgfpathlineto{\pgfqpoint{1.983506in}{1.360972in}}%
\pgfpathlineto{\pgfqpoint{1.983756in}{1.355205in}}%
\pgfpathlineto{\pgfqpoint{1.972422in}{1.354713in}}%
\pgfpathlineto{\pgfqpoint{1.938345in}{1.353373in}}%
\pgfpathlineto{\pgfqpoint{1.937355in}{1.364751in}}%
\pgfpathlineto{\pgfqpoint{1.936000in}{1.405159in}}%
\pgfpathlineto{\pgfqpoint{1.934699in}{1.432796in}}%
\pgfpathlineto{\pgfqpoint{1.953526in}{1.433541in}}%
\pgfpathlineto{\pgfqpoint{1.954322in}{1.428443in}}%
\pgfpathlineto{\pgfqpoint{1.954703in}{1.423577in}}%
\pgfpathlineto{\pgfqpoint{1.961879in}{1.417685in}}%
\pgfpathlineto{\pgfqpoint{1.956575in}{1.409480in}}%
\pgfpathlineto{\pgfqpoint{1.958754in}{1.394373in}}%
\pgfpathlineto{\pgfqpoint{1.960884in}{1.393037in}}%
\pgfpathlineto{\pgfqpoint{1.963562in}{1.383583in}}%
\pgfpathlineto{\pgfqpoint{1.969155in}{1.380236in}}%
\pgfpathlineto{\pgfqpoint{1.980017in}{1.377888in}}%
\pgfpathlineto{\pgfqpoint{1.984280in}{1.369339in}}%
\pgfpathlineto{\pgfqpoint{2.018995in}{1.370976in}}%
\pgfpathclose%
\pgfusepath{fill}%
\end{pgfscope}%
\begin{pgfscope}%
\pgfpathrectangle{\pgfqpoint{0.100000in}{0.100000in}}{\pgfqpoint{3.007045in}{1.925000in}}%
\pgfusepath{clip}%
\pgfsetbuttcap%
\pgfsetmiterjoin%
\definecolor{currentfill}{rgb}{0.632526,0.797647,0.886874}%
\pgfsetfillcolor{currentfill}%
\pgfsetlinewidth{0.000000pt}%
\definecolor{currentstroke}{rgb}{0.000000,0.000000,0.000000}%
\pgfsetstrokecolor{currentstroke}%
\pgfsetstrokeopacity{0.000000}%
\pgfsetdash{}{0pt}%
\pgfpathmoveto{\pgfqpoint{2.604673in}{1.078358in}}%
\pgfpathlineto{\pgfqpoint{2.599528in}{1.076992in}}%
\pgfpathlineto{\pgfqpoint{2.596091in}{1.081050in}}%
\pgfpathlineto{\pgfqpoint{2.591032in}{1.080451in}}%
\pgfpathlineto{\pgfqpoint{2.584592in}{1.075050in}}%
\pgfpathlineto{\pgfqpoint{2.578887in}{1.072582in}}%
\pgfpathlineto{\pgfqpoint{2.574988in}{1.077569in}}%
\pgfpathlineto{\pgfqpoint{2.564592in}{1.081776in}}%
\pgfpathlineto{\pgfqpoint{2.562917in}{1.088318in}}%
\pgfpathlineto{\pgfqpoint{2.564677in}{1.094671in}}%
\pgfpathlineto{\pgfqpoint{2.570652in}{1.100774in}}%
\pgfpathlineto{\pgfqpoint{2.573073in}{1.099388in}}%
\pgfpathlineto{\pgfqpoint{2.578351in}{1.105297in}}%
\pgfpathlineto{\pgfqpoint{2.579736in}{1.109975in}}%
\pgfpathlineto{\pgfqpoint{2.584797in}{1.114957in}}%
\pgfpathlineto{\pgfqpoint{2.587434in}{1.124379in}}%
\pgfpathlineto{\pgfqpoint{2.593692in}{1.126779in}}%
\pgfpathlineto{\pgfqpoint{2.598994in}{1.120903in}}%
\pgfpathlineto{\pgfqpoint{2.602817in}{1.120583in}}%
\pgfpathlineto{\pgfqpoint{2.609095in}{1.111445in}}%
\pgfpathlineto{\pgfqpoint{2.612770in}{1.112672in}}%
\pgfpathlineto{\pgfqpoint{2.621502in}{1.105990in}}%
\pgfpathlineto{\pgfqpoint{2.626460in}{1.104802in}}%
\pgfpathlineto{\pgfqpoint{2.623729in}{1.094309in}}%
\pgfpathlineto{\pgfqpoint{2.616873in}{1.090165in}}%
\pgfpathlineto{\pgfqpoint{2.614846in}{1.077543in}}%
\pgfpathlineto{\pgfqpoint{2.604673in}{1.078358in}}%
\pgfpathclose%
\pgfusepath{fill}%
\end{pgfscope}%
\begin{pgfscope}%
\pgfpathrectangle{\pgfqpoint{0.100000in}{0.100000in}}{\pgfqpoint{3.007045in}{1.925000in}}%
\pgfusepath{clip}%
\pgfsetbuttcap%
\pgfsetmiterjoin%
\definecolor{currentfill}{rgb}{0.396909,0.666851,0.830358}%
\pgfsetfillcolor{currentfill}%
\pgfsetlinewidth{0.000000pt}%
\definecolor{currentstroke}{rgb}{0.000000,0.000000,0.000000}%
\pgfsetstrokecolor{currentstroke}%
\pgfsetstrokeopacity{0.000000}%
\pgfsetdash{}{0pt}%
\pgfpathmoveto{\pgfqpoint{0.530203in}{1.618670in}}%
\pgfpathlineto{\pgfqpoint{0.573938in}{1.605496in}}%
\pgfpathlineto{\pgfqpoint{0.596099in}{1.599150in}}%
\pgfpathlineto{\pgfqpoint{0.588372in}{1.571730in}}%
\pgfpathlineto{\pgfqpoint{0.586719in}{1.572193in}}%
\pgfpathlineto{\pgfqpoint{0.578445in}{1.544847in}}%
\pgfpathlineto{\pgfqpoint{0.605512in}{1.537505in}}%
\pgfpathlineto{\pgfqpoint{0.592563in}{1.489427in}}%
\pgfpathlineto{\pgfqpoint{0.562393in}{1.497838in}}%
\pgfpathlineto{\pgfqpoint{0.520911in}{1.509711in}}%
\pgfpathlineto{\pgfqpoint{0.534635in}{1.557261in}}%
\pgfpathlineto{\pgfqpoint{0.513001in}{1.563845in}}%
\pgfpathlineto{\pgfqpoint{0.524578in}{1.602426in}}%
\pgfpathlineto{\pgfqpoint{0.530203in}{1.618670in}}%
\pgfpathclose%
\pgfusepath{fill}%
\end{pgfscope}%
\begin{pgfscope}%
\pgfpathrectangle{\pgfqpoint{0.100000in}{0.100000in}}{\pgfqpoint{3.007045in}{1.925000in}}%
\pgfusepath{clip}%
\pgfsetbuttcap%
\pgfsetmiterjoin%
\definecolor{currentfill}{rgb}{0.270804,0.580730,0.781146}%
\pgfsetfillcolor{currentfill}%
\pgfsetlinewidth{0.000000pt}%
\definecolor{currentstroke}{rgb}{0.000000,0.000000,0.000000}%
\pgfsetstrokecolor{currentstroke}%
\pgfsetstrokeopacity{0.000000}%
\pgfsetdash{}{0pt}%
\pgfpathmoveto{\pgfqpoint{1.799463in}{0.919527in}}%
\pgfpathlineto{\pgfqpoint{1.795605in}{0.946063in}}%
\pgfpathlineto{\pgfqpoint{1.771725in}{0.946010in}}%
\pgfpathlineto{\pgfqpoint{1.772350in}{0.979565in}}%
\pgfpathlineto{\pgfqpoint{1.792522in}{0.979422in}}%
\pgfpathlineto{\pgfqpoint{1.792483in}{0.985928in}}%
\pgfpathlineto{\pgfqpoint{1.821509in}{0.984906in}}%
\pgfpathlineto{\pgfqpoint{1.821240in}{0.968478in}}%
\pgfpathlineto{\pgfqpoint{1.832387in}{0.968604in}}%
\pgfpathlineto{\pgfqpoint{1.832421in}{0.955990in}}%
\pgfpathlineto{\pgfqpoint{1.843659in}{0.955855in}}%
\pgfpathlineto{\pgfqpoint{1.845759in}{0.949170in}}%
\pgfpathlineto{\pgfqpoint{1.849514in}{0.944275in}}%
\pgfpathlineto{\pgfqpoint{1.855176in}{0.941374in}}%
\pgfpathlineto{\pgfqpoint{1.855032in}{0.933900in}}%
\pgfpathlineto{\pgfqpoint{1.851232in}{0.928170in}}%
\pgfpathlineto{\pgfqpoint{1.851596in}{0.920385in}}%
\pgfpathlineto{\pgfqpoint{1.830635in}{0.919840in}}%
\pgfpathlineto{\pgfqpoint{1.819183in}{0.919101in}}%
\pgfpathlineto{\pgfqpoint{1.799463in}{0.919527in}}%
\pgfpathclose%
\pgfusepath{fill}%
\end{pgfscope}%
\begin{pgfscope}%
\pgfpathrectangle{\pgfqpoint{0.100000in}{0.100000in}}{\pgfqpoint{3.007045in}{1.925000in}}%
\pgfusepath{clip}%
\pgfsetbuttcap%
\pgfsetmiterjoin%
\definecolor{currentfill}{rgb}{0.435294,0.690965,0.842599}%
\pgfsetfillcolor{currentfill}%
\pgfsetlinewidth{0.000000pt}%
\definecolor{currentstroke}{rgb}{0.000000,0.000000,0.000000}%
\pgfsetstrokecolor{currentstroke}%
\pgfsetstrokeopacity{0.000000}%
\pgfsetdash{}{0pt}%
\pgfpathmoveto{\pgfqpoint{2.584761in}{0.832670in}}%
\pgfpathlineto{\pgfqpoint{2.584582in}{0.822960in}}%
\pgfpathlineto{\pgfqpoint{2.575563in}{0.818701in}}%
\pgfpathlineto{\pgfqpoint{2.570619in}{0.823049in}}%
\pgfpathlineto{\pgfqpoint{2.556125in}{0.810180in}}%
\pgfpathlineto{\pgfqpoint{2.545370in}{0.817990in}}%
\pgfpathlineto{\pgfqpoint{2.532363in}{0.822168in}}%
\pgfpathlineto{\pgfqpoint{2.525782in}{0.825781in}}%
\pgfpathlineto{\pgfqpoint{2.521533in}{0.825837in}}%
\pgfpathlineto{\pgfqpoint{2.530038in}{0.837997in}}%
\pgfpathlineto{\pgfqpoint{2.516115in}{0.841918in}}%
\pgfpathlineto{\pgfqpoint{2.506929in}{0.849528in}}%
\pgfpathlineto{\pgfqpoint{2.503173in}{0.844767in}}%
\pgfpathlineto{\pgfqpoint{2.492621in}{0.846691in}}%
\pgfpathlineto{\pgfqpoint{2.488445in}{0.853368in}}%
\pgfpathlineto{\pgfqpoint{2.482424in}{0.851583in}}%
\pgfpathlineto{\pgfqpoint{2.486659in}{0.864555in}}%
\pgfpathlineto{\pgfqpoint{2.484386in}{0.869099in}}%
\pgfpathlineto{\pgfqpoint{2.485317in}{0.874629in}}%
\pgfpathlineto{\pgfqpoint{2.490797in}{0.879764in}}%
\pgfpathlineto{\pgfqpoint{2.496666in}{0.891862in}}%
\pgfpathlineto{\pgfqpoint{2.502098in}{0.888878in}}%
\pgfpathlineto{\pgfqpoint{2.508933in}{0.891663in}}%
\pgfpathlineto{\pgfqpoint{2.508172in}{0.896146in}}%
\pgfpathlineto{\pgfqpoint{2.537783in}{0.898762in}}%
\pgfpathlineto{\pgfqpoint{2.538562in}{0.893175in}}%
\pgfpathlineto{\pgfqpoint{2.545276in}{0.888835in}}%
\pgfpathlineto{\pgfqpoint{2.541892in}{0.875274in}}%
\pgfpathlineto{\pgfqpoint{2.551190in}{0.869149in}}%
\pgfpathlineto{\pgfqpoint{2.555420in}{0.871742in}}%
\pgfpathlineto{\pgfqpoint{2.557919in}{0.868208in}}%
\pgfpathlineto{\pgfqpoint{2.558099in}{0.855702in}}%
\pgfpathlineto{\pgfqpoint{2.560497in}{0.852920in}}%
\pgfpathlineto{\pgfqpoint{2.559661in}{0.847879in}}%
\pgfpathlineto{\pgfqpoint{2.566054in}{0.842571in}}%
\pgfpathlineto{\pgfqpoint{2.571176in}{0.834427in}}%
\pgfpathlineto{\pgfqpoint{2.577826in}{0.830633in}}%
\pgfpathlineto{\pgfqpoint{2.584761in}{0.832670in}}%
\pgfpathclose%
\pgfusepath{fill}%
\end{pgfscope}%
\begin{pgfscope}%
\pgfpathrectangle{\pgfqpoint{0.100000in}{0.100000in}}{\pgfqpoint{3.007045in}{1.925000in}}%
\pgfusepath{clip}%
\pgfsetbuttcap%
\pgfsetmiterjoin%
\definecolor{currentfill}{rgb}{0.793449,0.870142,0.942914}%
\pgfsetfillcolor{currentfill}%
\pgfsetlinewidth{0.000000pt}%
\definecolor{currentstroke}{rgb}{0.000000,0.000000,0.000000}%
\pgfsetstrokecolor{currentstroke}%
\pgfsetstrokeopacity{0.000000}%
\pgfsetdash{}{0pt}%
\pgfpathmoveto{\pgfqpoint{2.413670in}{0.911416in}}%
\pgfpathlineto{\pgfqpoint{2.413500in}{0.907511in}}%
\pgfpathlineto{\pgfqpoint{2.408265in}{0.902647in}}%
\pgfpathlineto{\pgfqpoint{2.403903in}{0.896059in}}%
\pgfpathlineto{\pgfqpoint{2.402990in}{0.891059in}}%
\pgfpathlineto{\pgfqpoint{2.391995in}{0.891918in}}%
\pgfpathlineto{\pgfqpoint{2.389462in}{0.896298in}}%
\pgfpathlineto{\pgfqpoint{2.385967in}{0.896018in}}%
\pgfpathlineto{\pgfqpoint{2.385392in}{0.899513in}}%
\pgfpathlineto{\pgfqpoint{2.391763in}{0.908652in}}%
\pgfpathlineto{\pgfqpoint{2.383505in}{0.917672in}}%
\pgfpathlineto{\pgfqpoint{2.378628in}{0.917255in}}%
\pgfpathlineto{\pgfqpoint{2.379773in}{0.923563in}}%
\pgfpathlineto{\pgfqpoint{2.381105in}{0.925986in}}%
\pgfpathlineto{\pgfqpoint{2.392521in}{0.928792in}}%
\pgfpathlineto{\pgfqpoint{2.398481in}{0.931422in}}%
\pgfpathlineto{\pgfqpoint{2.406157in}{0.920864in}}%
\pgfpathlineto{\pgfqpoint{2.410454in}{0.917486in}}%
\pgfpathlineto{\pgfqpoint{2.413670in}{0.911416in}}%
\pgfpathclose%
\pgfusepath{fill}%
\end{pgfscope}%
\begin{pgfscope}%
\pgfpathrectangle{\pgfqpoint{0.100000in}{0.100000in}}{\pgfqpoint{3.007045in}{1.925000in}}%
\pgfusepath{clip}%
\pgfsetbuttcap%
\pgfsetmiterjoin%
\definecolor{currentfill}{rgb}{0.504314,0.728858,0.857486}%
\pgfsetfillcolor{currentfill}%
\pgfsetlinewidth{0.000000pt}%
\definecolor{currentstroke}{rgb}{0.000000,0.000000,0.000000}%
\pgfsetstrokecolor{currentstroke}%
\pgfsetstrokeopacity{0.000000}%
\pgfsetdash{}{0pt}%
\pgfpathmoveto{\pgfqpoint{1.923472in}{0.954935in}}%
\pgfpathlineto{\pgfqpoint{1.921677in}{0.946416in}}%
\pgfpathlineto{\pgfqpoint{1.921721in}{0.941554in}}%
\pgfpathlineto{\pgfqpoint{1.915999in}{0.941499in}}%
\pgfpathlineto{\pgfqpoint{1.916049in}{0.935756in}}%
\pgfpathlineto{\pgfqpoint{1.910330in}{0.935711in}}%
\pgfpathlineto{\pgfqpoint{1.910317in}{0.941444in}}%
\pgfpathlineto{\pgfqpoint{1.887586in}{0.941336in}}%
\pgfpathlineto{\pgfqpoint{1.884698in}{0.944172in}}%
\pgfpathlineto{\pgfqpoint{1.884426in}{0.969017in}}%
\pgfpathlineto{\pgfqpoint{1.886046in}{0.969526in}}%
\pgfpathlineto{\pgfqpoint{1.890382in}{0.969613in}}%
\pgfpathlineto{\pgfqpoint{1.890280in}{0.989939in}}%
\pgfpathlineto{\pgfqpoint{1.924728in}{0.990132in}}%
\pgfpathlineto{\pgfqpoint{1.924925in}{0.970656in}}%
\pgfpathlineto{\pgfqpoint{1.923336in}{0.970590in}}%
\pgfpathlineto{\pgfqpoint{1.923472in}{0.954935in}}%
\pgfpathclose%
\pgfusepath{fill}%
\end{pgfscope}%
\begin{pgfscope}%
\pgfpathrectangle{\pgfqpoint{0.100000in}{0.100000in}}{\pgfqpoint{3.007045in}{1.925000in}}%
\pgfusepath{clip}%
\pgfsetbuttcap%
\pgfsetmiterjoin%
\definecolor{currentfill}{rgb}{0.652211,0.806013,0.893764}%
\pgfsetfillcolor{currentfill}%
\pgfsetlinewidth{0.000000pt}%
\definecolor{currentstroke}{rgb}{0.000000,0.000000,0.000000}%
\pgfsetstrokecolor{currentstroke}%
\pgfsetstrokeopacity{0.000000}%
\pgfsetdash{}{0pt}%
\pgfpathmoveto{\pgfqpoint{2.680930in}{1.081273in}}%
\pgfpathlineto{\pgfqpoint{2.683306in}{1.069582in}}%
\pgfpathlineto{\pgfqpoint{2.679621in}{1.057585in}}%
\pgfpathlineto{\pgfqpoint{2.672729in}{1.056217in}}%
\pgfpathlineto{\pgfqpoint{2.643747in}{1.050644in}}%
\pgfpathlineto{\pgfqpoint{2.608257in}{1.044199in}}%
\pgfpathlineto{\pgfqpoint{2.604048in}{1.043452in}}%
\pgfpathlineto{\pgfqpoint{2.604673in}{1.078358in}}%
\pgfpathlineto{\pgfqpoint{2.614846in}{1.077543in}}%
\pgfpathlineto{\pgfqpoint{2.616873in}{1.090165in}}%
\pgfpathlineto{\pgfqpoint{2.623729in}{1.094309in}}%
\pgfpathlineto{\pgfqpoint{2.626460in}{1.104802in}}%
\pgfpathlineto{\pgfqpoint{2.633816in}{1.102244in}}%
\pgfpathlineto{\pgfqpoint{2.641373in}{1.123470in}}%
\pgfpathlineto{\pgfqpoint{2.641185in}{1.127114in}}%
\pgfpathlineto{\pgfqpoint{2.644603in}{1.131634in}}%
\pgfpathlineto{\pgfqpoint{2.649818in}{1.126222in}}%
\pgfpathlineto{\pgfqpoint{2.649572in}{1.112889in}}%
\pgfpathlineto{\pgfqpoint{2.645322in}{1.106286in}}%
\pgfpathlineto{\pgfqpoint{2.646388in}{1.101662in}}%
\pgfpathlineto{\pgfqpoint{2.670092in}{1.099194in}}%
\pgfpathlineto{\pgfqpoint{2.665311in}{1.095015in}}%
\pgfpathlineto{\pgfqpoint{2.667768in}{1.085077in}}%
\pgfpathlineto{\pgfqpoint{2.674370in}{1.085389in}}%
\pgfpathlineto{\pgfqpoint{2.680930in}{1.081273in}}%
\pgfpathclose%
\pgfusepath{fill}%
\end{pgfscope}%
\begin{pgfscope}%
\pgfpathrectangle{\pgfqpoint{0.100000in}{0.100000in}}{\pgfqpoint{3.007045in}{1.925000in}}%
\pgfusepath{clip}%
\pgfsetbuttcap%
\pgfsetmiterjoin%
\definecolor{currentfill}{rgb}{0.491765,0.721968,0.854779}%
\pgfsetfillcolor{currentfill}%
\pgfsetlinewidth{0.000000pt}%
\definecolor{currentstroke}{rgb}{0.000000,0.000000,0.000000}%
\pgfsetstrokecolor{currentstroke}%
\pgfsetstrokeopacity{0.000000}%
\pgfsetdash{}{0pt}%
\pgfpathmoveto{\pgfqpoint{0.878844in}{1.527019in}}%
\pgfpathlineto{\pgfqpoint{0.872520in}{1.537698in}}%
\pgfpathlineto{\pgfqpoint{0.861425in}{1.547079in}}%
\pgfpathlineto{\pgfqpoint{0.860657in}{1.551817in}}%
\pgfpathlineto{\pgfqpoint{0.849873in}{1.555251in}}%
\pgfpathlineto{\pgfqpoint{0.842609in}{1.559070in}}%
\pgfpathlineto{\pgfqpoint{0.836443in}{1.560436in}}%
\pgfpathlineto{\pgfqpoint{0.834605in}{1.565249in}}%
\pgfpathlineto{\pgfqpoint{0.829274in}{1.562729in}}%
\pgfpathlineto{\pgfqpoint{0.827488in}{1.571131in}}%
\pgfpathlineto{\pgfqpoint{0.829905in}{1.577513in}}%
\pgfpathlineto{\pgfqpoint{0.824028in}{1.581837in}}%
\pgfpathlineto{\pgfqpoint{0.824845in}{1.588003in}}%
\pgfpathlineto{\pgfqpoint{0.820142in}{1.591987in}}%
\pgfpathlineto{\pgfqpoint{0.824171in}{1.596209in}}%
\pgfpathlineto{\pgfqpoint{0.825412in}{1.600994in}}%
\pgfpathlineto{\pgfqpoint{0.822960in}{1.604421in}}%
\pgfpathlineto{\pgfqpoint{0.824810in}{1.609423in}}%
\pgfpathlineto{\pgfqpoint{0.830622in}{1.610723in}}%
\pgfpathlineto{\pgfqpoint{0.834363in}{1.617027in}}%
\pgfpathlineto{\pgfqpoint{0.841346in}{1.613085in}}%
\pgfpathlineto{\pgfqpoint{0.853704in}{1.621066in}}%
\pgfpathlineto{\pgfqpoint{0.857696in}{1.638685in}}%
\pgfpathlineto{\pgfqpoint{0.860069in}{1.641241in}}%
\pgfpathlineto{\pgfqpoint{0.865574in}{1.646714in}}%
\pgfpathlineto{\pgfqpoint{0.859318in}{1.660383in}}%
\pgfpathlineto{\pgfqpoint{0.865364in}{1.658676in}}%
\pgfpathlineto{\pgfqpoint{0.871070in}{1.663297in}}%
\pgfpathlineto{\pgfqpoint{0.873352in}{1.663504in}}%
\pgfpathlineto{\pgfqpoint{0.879064in}{1.657147in}}%
\pgfpathlineto{\pgfqpoint{0.890559in}{1.659065in}}%
\pgfpathlineto{\pgfqpoint{0.894600in}{1.662965in}}%
\pgfpathlineto{\pgfqpoint{0.901414in}{1.666132in}}%
\pgfpathlineto{\pgfqpoint{0.906088in}{1.658804in}}%
\pgfpathlineto{\pgfqpoint{0.903714in}{1.653787in}}%
\pgfpathlineto{\pgfqpoint{0.906555in}{1.644348in}}%
\pgfpathlineto{\pgfqpoint{0.905525in}{1.640271in}}%
\pgfpathlineto{\pgfqpoint{0.910956in}{1.624666in}}%
\pgfpathlineto{\pgfqpoint{0.914945in}{1.620213in}}%
\pgfpathlineto{\pgfqpoint{0.912228in}{1.607641in}}%
\pgfpathlineto{\pgfqpoint{0.916159in}{1.603777in}}%
\pgfpathlineto{\pgfqpoint{0.923418in}{1.601071in}}%
\pgfpathlineto{\pgfqpoint{0.926655in}{1.595511in}}%
\pgfpathlineto{\pgfqpoint{0.928393in}{1.581446in}}%
\pgfpathlineto{\pgfqpoint{0.927577in}{1.576545in}}%
\pgfpathlineto{\pgfqpoint{0.933731in}{1.569068in}}%
\pgfpathlineto{\pgfqpoint{0.935411in}{1.569991in}}%
\pgfpathlineto{\pgfqpoint{0.933521in}{1.560492in}}%
\pgfpathlineto{\pgfqpoint{0.905554in}{1.566329in}}%
\pgfpathlineto{\pgfqpoint{0.903257in}{1.555054in}}%
\pgfpathlineto{\pgfqpoint{0.900875in}{1.550400in}}%
\pgfpathlineto{\pgfqpoint{0.897650in}{1.536577in}}%
\pgfpathlineto{\pgfqpoint{0.891068in}{1.533017in}}%
\pgfpathlineto{\pgfqpoint{0.881395in}{1.530161in}}%
\pgfpathlineto{\pgfqpoint{0.878844in}{1.527019in}}%
\pgfpathclose%
\pgfusepath{fill}%
\end{pgfscope}%
\begin{pgfscope}%
\pgfpathrectangle{\pgfqpoint{0.100000in}{0.100000in}}{\pgfqpoint{3.007045in}{1.925000in}}%
\pgfusepath{clip}%
\pgfsetbuttcap%
\pgfsetmiterjoin%
\definecolor{currentfill}{rgb}{0.142607,0.456332,0.716601}%
\pgfsetfillcolor{currentfill}%
\pgfsetlinewidth{0.000000pt}%
\definecolor{currentstroke}{rgb}{0.000000,0.000000,0.000000}%
\pgfsetstrokecolor{currentstroke}%
\pgfsetstrokeopacity{0.000000}%
\pgfsetdash{}{0pt}%
\pgfpathmoveto{\pgfqpoint{1.002522in}{1.140541in}}%
\pgfpathlineto{\pgfqpoint{0.997732in}{1.138462in}}%
\pgfpathlineto{\pgfqpoint{0.992394in}{1.131729in}}%
\pgfpathlineto{\pgfqpoint{0.990253in}{1.132805in}}%
\pgfpathlineto{\pgfqpoint{0.982614in}{1.127004in}}%
\pgfpathlineto{\pgfqpoint{0.975147in}{1.126984in}}%
\pgfpathlineto{\pgfqpoint{0.967830in}{1.113017in}}%
\pgfpathlineto{\pgfqpoint{0.961951in}{1.110649in}}%
\pgfpathlineto{\pgfqpoint{0.958933in}{1.107125in}}%
\pgfpathlineto{\pgfqpoint{0.937370in}{1.110856in}}%
\pgfpathlineto{\pgfqpoint{0.870763in}{1.123009in}}%
\pgfpathlineto{\pgfqpoint{0.854337in}{1.126666in}}%
\pgfpathlineto{\pgfqpoint{0.857375in}{1.143704in}}%
\pgfpathlineto{\pgfqpoint{0.863012in}{1.142510in}}%
\pgfpathlineto{\pgfqpoint{0.864002in}{1.147526in}}%
\pgfpathlineto{\pgfqpoint{0.869732in}{1.146966in}}%
\pgfpathlineto{\pgfqpoint{0.872431in}{1.163681in}}%
\pgfpathlineto{\pgfqpoint{0.877415in}{1.166653in}}%
\pgfpathlineto{\pgfqpoint{0.882037in}{1.174273in}}%
\pgfpathlineto{\pgfqpoint{0.878976in}{1.180671in}}%
\pgfpathlineto{\pgfqpoint{0.875518in}{1.183498in}}%
\pgfpathlineto{\pgfqpoint{0.875936in}{1.191519in}}%
\pgfpathlineto{\pgfqpoint{0.880713in}{1.197892in}}%
\pgfpathlineto{\pgfqpoint{0.887613in}{1.196226in}}%
\pgfpathlineto{\pgfqpoint{0.892855in}{1.198668in}}%
\pgfpathlineto{\pgfqpoint{0.894249in}{1.207621in}}%
\pgfpathlineto{\pgfqpoint{0.900045in}{1.212653in}}%
\pgfpathlineto{\pgfqpoint{0.903422in}{1.212075in}}%
\pgfpathlineto{\pgfqpoint{0.907070in}{1.217261in}}%
\pgfpathlineto{\pgfqpoint{0.914996in}{1.214937in}}%
\pgfpathlineto{\pgfqpoint{0.942870in}{1.209747in}}%
\pgfpathlineto{\pgfqpoint{0.936407in}{1.175969in}}%
\pgfpathlineto{\pgfqpoint{0.913573in}{1.180156in}}%
\pgfpathlineto{\pgfqpoint{0.907847in}{1.175524in}}%
\pgfpathlineto{\pgfqpoint{0.904834in}{1.157752in}}%
\pgfpathlineto{\pgfqpoint{0.963729in}{1.147104in}}%
\pgfpathlineto{\pgfqpoint{1.002522in}{1.140541in}}%
\pgfpathclose%
\pgfusepath{fill}%
\end{pgfscope}%
\begin{pgfscope}%
\pgfpathrectangle{\pgfqpoint{0.100000in}{0.100000in}}{\pgfqpoint{3.007045in}{1.925000in}}%
\pgfusepath{clip}%
\pgfsetbuttcap%
\pgfsetmiterjoin%
\definecolor{currentfill}{rgb}{0.361599,0.642737,0.816578}%
\pgfsetfillcolor{currentfill}%
\pgfsetlinewidth{0.000000pt}%
\definecolor{currentstroke}{rgb}{0.000000,0.000000,0.000000}%
\pgfsetstrokecolor{currentstroke}%
\pgfsetstrokeopacity{0.000000}%
\pgfsetdash{}{0pt}%
\pgfpathmoveto{\pgfqpoint{1.768458in}{1.001242in}}%
\pgfpathlineto{\pgfqpoint{1.750868in}{1.001271in}}%
\pgfpathlineto{\pgfqpoint{1.749585in}{0.997555in}}%
\pgfpathlineto{\pgfqpoint{1.749489in}{0.974797in}}%
\pgfpathlineto{\pgfqpoint{1.729327in}{0.974933in}}%
\pgfpathlineto{\pgfqpoint{1.729236in}{0.963455in}}%
\pgfpathlineto{\pgfqpoint{1.719212in}{0.963573in}}%
\pgfpathlineto{\pgfqpoint{1.719649in}{1.001488in}}%
\pgfpathlineto{\pgfqpoint{1.721571in}{1.001483in}}%
\pgfpathlineto{\pgfqpoint{1.722263in}{1.049928in}}%
\pgfpathlineto{\pgfqpoint{1.767740in}{1.049552in}}%
\pgfpathlineto{\pgfqpoint{1.768402in}{1.026569in}}%
\pgfpathlineto{\pgfqpoint{1.768458in}{1.001242in}}%
\pgfpathclose%
\pgfusepath{fill}%
\end{pgfscope}%
\begin{pgfscope}%
\pgfpathrectangle{\pgfqpoint{0.100000in}{0.100000in}}{\pgfqpoint{3.007045in}{1.925000in}}%
\pgfusepath{clip}%
\pgfsetbuttcap%
\pgfsetmiterjoin%
\definecolor{currentfill}{rgb}{0.460392,0.704744,0.848012}%
\pgfsetfillcolor{currentfill}%
\pgfsetlinewidth{0.000000pt}%
\definecolor{currentstroke}{rgb}{0.000000,0.000000,0.000000}%
\pgfsetstrokecolor{currentstroke}%
\pgfsetstrokeopacity{0.000000}%
\pgfsetdash{}{0pt}%
\pgfpathmoveto{\pgfqpoint{1.192875in}{1.732905in}}%
\pgfpathlineto{\pgfqpoint{1.188345in}{1.720102in}}%
\pgfpathlineto{\pgfqpoint{1.188167in}{1.711464in}}%
\pgfpathlineto{\pgfqpoint{1.189559in}{1.703789in}}%
\pgfpathlineto{\pgfqpoint{1.186086in}{1.695738in}}%
\pgfpathlineto{\pgfqpoint{1.187986in}{1.693328in}}%
\pgfpathlineto{\pgfqpoint{1.189998in}{1.686616in}}%
\pgfpathlineto{\pgfqpoint{1.153729in}{1.691500in}}%
\pgfpathlineto{\pgfqpoint{1.119627in}{1.696907in}}%
\pgfpathlineto{\pgfqpoint{1.119016in}{1.693144in}}%
\pgfpathlineto{\pgfqpoint{1.102971in}{1.695747in}}%
\pgfpathlineto{\pgfqpoint{1.102514in}{1.707496in}}%
\pgfpathlineto{\pgfqpoint{1.104710in}{1.709101in}}%
\pgfpathlineto{\pgfqpoint{1.107755in}{1.728110in}}%
\pgfpathlineto{\pgfqpoint{1.103070in}{1.731808in}}%
\pgfpathlineto{\pgfqpoint{1.097369in}{1.732746in}}%
\pgfpathlineto{\pgfqpoint{1.092154in}{1.736531in}}%
\pgfpathlineto{\pgfqpoint{1.094250in}{1.746968in}}%
\pgfpathlineto{\pgfqpoint{1.107591in}{1.752201in}}%
\pgfpathlineto{\pgfqpoint{1.110393in}{1.763495in}}%
\pgfpathlineto{\pgfqpoint{1.123379in}{1.762908in}}%
\pgfpathlineto{\pgfqpoint{1.126744in}{1.761331in}}%
\pgfpathlineto{\pgfqpoint{1.135562in}{1.764609in}}%
\pgfpathlineto{\pgfqpoint{1.141833in}{1.762846in}}%
\pgfpathlineto{\pgfqpoint{1.148828in}{1.763186in}}%
\pgfpathlineto{\pgfqpoint{1.152002in}{1.758091in}}%
\pgfpathlineto{\pgfqpoint{1.156612in}{1.751283in}}%
\pgfpathlineto{\pgfqpoint{1.164883in}{1.747708in}}%
\pgfpathlineto{\pgfqpoint{1.170837in}{1.747018in}}%
\pgfpathlineto{\pgfqpoint{1.174368in}{1.744062in}}%
\pgfpathlineto{\pgfqpoint{1.187367in}{1.743646in}}%
\pgfpathlineto{\pgfqpoint{1.193102in}{1.740744in}}%
\pgfpathlineto{\pgfqpoint{1.192875in}{1.732905in}}%
\pgfpathclose%
\pgfusepath{fill}%
\end{pgfscope}%
\begin{pgfscope}%
\pgfpathrectangle{\pgfqpoint{0.100000in}{0.100000in}}{\pgfqpoint{3.007045in}{1.925000in}}%
\pgfusepath{clip}%
\pgfsetbuttcap%
\pgfsetmiterjoin%
\definecolor{currentfill}{rgb}{0.256286,0.570012,0.775163}%
\pgfsetfillcolor{currentfill}%
\pgfsetlinewidth{0.000000pt}%
\definecolor{currentstroke}{rgb}{0.000000,0.000000,0.000000}%
\pgfsetstrokecolor{currentstroke}%
\pgfsetstrokeopacity{0.000000}%
\pgfsetdash{}{0pt}%
\pgfpathmoveto{\pgfqpoint{1.685427in}{0.912299in}}%
\pgfpathlineto{\pgfqpoint{1.685760in}{0.932254in}}%
\pgfpathlineto{\pgfqpoint{1.685998in}{0.946658in}}%
\pgfpathlineto{\pgfqpoint{1.675476in}{0.946791in}}%
\pgfpathlineto{\pgfqpoint{1.675597in}{0.952567in}}%
\pgfpathlineto{\pgfqpoint{1.669953in}{0.952684in}}%
\pgfpathlineto{\pgfqpoint{1.670084in}{0.958402in}}%
\pgfpathlineto{\pgfqpoint{1.664393in}{0.958534in}}%
\pgfpathlineto{\pgfqpoint{1.664644in}{0.969981in}}%
\pgfpathlineto{\pgfqpoint{1.671642in}{0.972779in}}%
\pgfpathlineto{\pgfqpoint{1.670693in}{0.976025in}}%
\pgfpathlineto{\pgfqpoint{1.642047in}{0.976241in}}%
\pgfpathlineto{\pgfqpoint{1.642737in}{1.002945in}}%
\pgfpathlineto{\pgfqpoint{1.659290in}{1.002553in}}%
\pgfpathlineto{\pgfqpoint{1.719649in}{1.001488in}}%
\pgfpathlineto{\pgfqpoint{1.719212in}{0.963573in}}%
\pgfpathlineto{\pgfqpoint{1.729236in}{0.963455in}}%
\pgfpathlineto{\pgfqpoint{1.729327in}{0.974933in}}%
\pgfpathlineto{\pgfqpoint{1.749489in}{0.974797in}}%
\pgfpathlineto{\pgfqpoint{1.749585in}{0.997555in}}%
\pgfpathlineto{\pgfqpoint{1.750868in}{1.001271in}}%
\pgfpathlineto{\pgfqpoint{1.768458in}{1.001242in}}%
\pgfpathlineto{\pgfqpoint{1.772364in}{0.997534in}}%
\pgfpathlineto{\pgfqpoint{1.772350in}{0.979565in}}%
\pgfpathlineto{\pgfqpoint{1.771725in}{0.946010in}}%
\pgfpathlineto{\pgfqpoint{1.766020in}{0.946003in}}%
\pgfpathlineto{\pgfqpoint{1.766027in}{0.940252in}}%
\pgfpathlineto{\pgfqpoint{1.761242in}{0.940278in}}%
\pgfpathlineto{\pgfqpoint{1.757448in}{0.932993in}}%
\pgfpathlineto{\pgfqpoint{1.758114in}{0.923040in}}%
\pgfpathlineto{\pgfqpoint{1.746777in}{0.922489in}}%
\pgfpathlineto{\pgfqpoint{1.740818in}{0.919900in}}%
\pgfpathlineto{\pgfqpoint{1.737553in}{0.926013in}}%
\pgfpathlineto{\pgfqpoint{1.731364in}{0.926038in}}%
\pgfpathlineto{\pgfqpoint{1.708573in}{0.926325in}}%
\pgfpathlineto{\pgfqpoint{1.708380in}{0.911948in}}%
\pgfpathlineto{\pgfqpoint{1.685427in}{0.912299in}}%
\pgfpathclose%
\pgfusepath{fill}%
\end{pgfscope}%
\begin{pgfscope}%
\pgfpathrectangle{\pgfqpoint{0.100000in}{0.100000in}}{\pgfqpoint{3.007045in}{1.925000in}}%
\pgfusepath{clip}%
\pgfsetbuttcap%
\pgfsetmiterjoin%
\definecolor{currentfill}{rgb}{0.031373,0.285675,0.564291}%
\pgfsetfillcolor{currentfill}%
\pgfsetlinewidth{0.000000pt}%
\definecolor{currentstroke}{rgb}{0.000000,0.000000,0.000000}%
\pgfsetstrokecolor{currentstroke}%
\pgfsetstrokeopacity{0.000000}%
\pgfsetdash{}{0pt}%
\pgfpathmoveto{\pgfqpoint{1.551371in}{1.482830in}}%
\pgfpathlineto{\pgfqpoint{1.552579in}{1.474853in}}%
\pgfpathlineto{\pgfqpoint{1.546378in}{1.477641in}}%
\pgfpathlineto{\pgfqpoint{1.542308in}{1.477127in}}%
\pgfpathlineto{\pgfqpoint{1.538292in}{1.482046in}}%
\pgfpathlineto{\pgfqpoint{1.517980in}{1.483159in}}%
\pgfpathlineto{\pgfqpoint{1.517880in}{1.481436in}}%
\pgfpathlineto{\pgfqpoint{1.485701in}{1.483457in}}%
\pgfpathlineto{\pgfqpoint{1.480026in}{1.483812in}}%
\pgfpathlineto{\pgfqpoint{1.481358in}{1.495327in}}%
\pgfpathlineto{\pgfqpoint{1.482943in}{1.518221in}}%
\pgfpathlineto{\pgfqpoint{1.484012in}{1.521937in}}%
\pgfpathlineto{\pgfqpoint{1.489275in}{1.520056in}}%
\pgfpathlineto{\pgfqpoint{1.498577in}{1.523374in}}%
\pgfpathlineto{\pgfqpoint{1.501256in}{1.521545in}}%
\pgfpathlineto{\pgfqpoint{1.506367in}{1.525118in}}%
\pgfpathlineto{\pgfqpoint{1.512562in}{1.521360in}}%
\pgfpathlineto{\pgfqpoint{1.519350in}{1.527670in}}%
\pgfpathlineto{\pgfqpoint{1.518970in}{1.529708in}}%
\pgfpathlineto{\pgfqpoint{1.553171in}{1.527805in}}%
\pgfpathlineto{\pgfqpoint{1.551887in}{1.504699in}}%
\pgfpathlineto{\pgfqpoint{1.551371in}{1.482830in}}%
\pgfpathclose%
\pgfusepath{fill}%
\end{pgfscope}%
\begin{pgfscope}%
\pgfpathrectangle{\pgfqpoint{0.100000in}{0.100000in}}{\pgfqpoint{3.007045in}{1.925000in}}%
\pgfusepath{clip}%
\pgfsetbuttcap%
\pgfsetmiterjoin%
\definecolor{currentfill}{rgb}{0.326290,0.618624,0.802799}%
\pgfsetfillcolor{currentfill}%
\pgfsetlinewidth{0.000000pt}%
\definecolor{currentstroke}{rgb}{0.000000,0.000000,0.000000}%
\pgfsetstrokecolor{currentstroke}%
\pgfsetstrokeopacity{0.000000}%
\pgfsetdash{}{0pt}%
\pgfpathmoveto{\pgfqpoint{2.020031in}{0.633536in}}%
\pgfpathlineto{\pgfqpoint{2.019201in}{0.650937in}}%
\pgfpathlineto{\pgfqpoint{2.013321in}{0.650638in}}%
\pgfpathlineto{\pgfqpoint{2.013078in}{0.656435in}}%
\pgfpathlineto{\pgfqpoint{2.038852in}{0.657801in}}%
\pgfpathlineto{\pgfqpoint{2.047327in}{0.661747in}}%
\pgfpathlineto{\pgfqpoint{2.055851in}{0.662822in}}%
\pgfpathlineto{\pgfqpoint{2.056646in}{0.646958in}}%
\pgfpathlineto{\pgfqpoint{2.058100in}{0.638257in}}%
\pgfpathlineto{\pgfqpoint{2.053466in}{0.638280in}}%
\pgfpathlineto{\pgfqpoint{2.049764in}{0.635164in}}%
\pgfpathlineto{\pgfqpoint{2.020031in}{0.633536in}}%
\pgfpathclose%
\pgfusepath{fill}%
\end{pgfscope}%
\begin{pgfscope}%
\pgfpathrectangle{\pgfqpoint{0.100000in}{0.100000in}}{\pgfqpoint{3.007045in}{1.925000in}}%
\pgfusepath{clip}%
\pgfsetbuttcap%
\pgfsetmiterjoin%
\definecolor{currentfill}{rgb}{0.031373,0.297855,0.582376}%
\pgfsetfillcolor{currentfill}%
\pgfsetlinewidth{0.000000pt}%
\definecolor{currentstroke}{rgb}{0.000000,0.000000,0.000000}%
\pgfsetstrokecolor{currentstroke}%
\pgfsetstrokeopacity{0.000000}%
\pgfsetdash{}{0pt}%
\pgfpathmoveto{\pgfqpoint{1.386813in}{0.624056in}}%
\pgfpathlineto{\pgfqpoint{1.351088in}{0.627381in}}%
\pgfpathlineto{\pgfqpoint{1.349725in}{0.635825in}}%
\pgfpathlineto{\pgfqpoint{1.337100in}{0.644492in}}%
\pgfpathlineto{\pgfqpoint{1.331766in}{0.642740in}}%
\pgfpathlineto{\pgfqpoint{1.333802in}{0.666433in}}%
\pgfpathlineto{\pgfqpoint{1.332017in}{0.666591in}}%
\pgfpathlineto{\pgfqpoint{1.334479in}{0.695294in}}%
\pgfpathlineto{\pgfqpoint{1.319633in}{0.696676in}}%
\pgfpathlineto{\pgfqpoint{1.322173in}{0.725422in}}%
\pgfpathlineto{\pgfqpoint{1.370186in}{0.721452in}}%
\pgfpathlineto{\pgfqpoint{1.398836in}{0.719387in}}%
\pgfpathlineto{\pgfqpoint{1.396350in}{0.690465in}}%
\pgfpathlineto{\pgfqpoint{1.391806in}{0.690752in}}%
\pgfpathlineto{\pgfqpoint{1.386813in}{0.624056in}}%
\pgfpathclose%
\pgfusepath{fill}%
\end{pgfscope}%
\begin{pgfscope}%
\pgfpathrectangle{\pgfqpoint{0.100000in}{0.100000in}}{\pgfqpoint{3.007045in}{1.925000in}}%
\pgfusepath{clip}%
\pgfsetbuttcap%
\pgfsetmiterjoin%
\definecolor{currentfill}{rgb}{0.491765,0.721968,0.854779}%
\pgfsetfillcolor{currentfill}%
\pgfsetlinewidth{0.000000pt}%
\definecolor{currentstroke}{rgb}{0.000000,0.000000,0.000000}%
\pgfsetstrokecolor{currentstroke}%
\pgfsetstrokeopacity{0.000000}%
\pgfsetdash{}{0pt}%
\pgfpathmoveto{\pgfqpoint{2.044424in}{1.118250in}}%
\pgfpathlineto{\pgfqpoint{2.020049in}{1.117235in}}%
\pgfpathlineto{\pgfqpoint{2.018572in}{1.121893in}}%
\pgfpathlineto{\pgfqpoint{2.022871in}{1.129850in}}%
\pgfpathlineto{\pgfqpoint{2.015419in}{1.134149in}}%
\pgfpathlineto{\pgfqpoint{2.005193in}{1.136744in}}%
\pgfpathlineto{\pgfqpoint{2.000724in}{1.130408in}}%
\pgfpathlineto{\pgfqpoint{1.994783in}{1.133102in}}%
\pgfpathlineto{\pgfqpoint{1.991482in}{1.141762in}}%
\pgfpathlineto{\pgfqpoint{1.992952in}{1.144098in}}%
\pgfpathlineto{\pgfqpoint{1.990535in}{1.152936in}}%
\pgfpathlineto{\pgfqpoint{1.986624in}{1.158384in}}%
\pgfpathlineto{\pgfqpoint{1.979066in}{1.163974in}}%
\pgfpathlineto{\pgfqpoint{1.995501in}{1.164493in}}%
\pgfpathlineto{\pgfqpoint{1.997266in}{1.151743in}}%
\pgfpathlineto{\pgfqpoint{2.002639in}{1.150292in}}%
\pgfpathlineto{\pgfqpoint{2.011513in}{1.150677in}}%
\pgfpathlineto{\pgfqpoint{2.019642in}{1.156899in}}%
\pgfpathlineto{\pgfqpoint{2.018445in}{1.173868in}}%
\pgfpathlineto{\pgfqpoint{2.029984in}{1.174625in}}%
\pgfpathlineto{\pgfqpoint{2.049876in}{1.175940in}}%
\pgfpathlineto{\pgfqpoint{2.050719in}{1.164389in}}%
\pgfpathlineto{\pgfqpoint{2.070604in}{1.165647in}}%
\pgfpathlineto{\pgfqpoint{2.071143in}{1.157025in}}%
\pgfpathlineto{\pgfqpoint{2.065520in}{1.156638in}}%
\pgfpathlineto{\pgfqpoint{2.067309in}{1.125374in}}%
\pgfpathlineto{\pgfqpoint{2.049708in}{1.124375in}}%
\pgfpathlineto{\pgfqpoint{2.050157in}{1.118638in}}%
\pgfpathlineto{\pgfqpoint{2.044424in}{1.118250in}}%
\pgfpathclose%
\pgfusepath{fill}%
\end{pgfscope}%
\begin{pgfscope}%
\pgfpathrectangle{\pgfqpoint{0.100000in}{0.100000in}}{\pgfqpoint{3.007045in}{1.925000in}}%
\pgfusepath{clip}%
\pgfsetbuttcap%
\pgfsetmiterjoin%
\definecolor{currentfill}{rgb}{0.396909,0.666851,0.830358}%
\pgfsetfillcolor{currentfill}%
\pgfsetlinewidth{0.000000pt}%
\definecolor{currentstroke}{rgb}{0.000000,0.000000,0.000000}%
\pgfsetstrokecolor{currentstroke}%
\pgfsetstrokeopacity{0.000000}%
\pgfsetdash{}{0pt}%
\pgfpathmoveto{\pgfqpoint{1.968862in}{0.741315in}}%
\pgfpathlineto{\pgfqpoint{1.935072in}{0.740260in}}%
\pgfpathlineto{\pgfqpoint{1.932837in}{0.741766in}}%
\pgfpathlineto{\pgfqpoint{1.930951in}{0.750384in}}%
\pgfpathlineto{\pgfqpoint{1.931371in}{0.753155in}}%
\pgfpathlineto{\pgfqpoint{1.939283in}{0.759647in}}%
\pgfpathlineto{\pgfqpoint{1.938599in}{0.766214in}}%
\pgfpathlineto{\pgfqpoint{1.968042in}{0.766663in}}%
\pgfpathlineto{\pgfqpoint{1.968862in}{0.741315in}}%
\pgfpathclose%
\pgfusepath{fill}%
\end{pgfscope}%
\begin{pgfscope}%
\pgfpathrectangle{\pgfqpoint{0.100000in}{0.100000in}}{\pgfqpoint{3.007045in}{1.925000in}}%
\pgfusepath{clip}%
\pgfsetbuttcap%
\pgfsetmiterjoin%
\definecolor{currentfill}{rgb}{0.356555,0.639293,0.814610}%
\pgfsetfillcolor{currentfill}%
\pgfsetlinewidth{0.000000pt}%
\definecolor{currentstroke}{rgb}{0.000000,0.000000,0.000000}%
\pgfsetstrokecolor{currentstroke}%
\pgfsetstrokeopacity{0.000000}%
\pgfsetdash{}{0pt}%
\pgfpathmoveto{\pgfqpoint{1.455754in}{1.302023in}}%
\pgfpathlineto{\pgfqpoint{1.427276in}{1.304079in}}%
\pgfpathlineto{\pgfqpoint{1.428992in}{1.326930in}}%
\pgfpathlineto{\pgfqpoint{1.424949in}{1.327230in}}%
\pgfpathlineto{\pgfqpoint{1.426380in}{1.344786in}}%
\pgfpathlineto{\pgfqpoint{1.428115in}{1.350443in}}%
\pgfpathlineto{\pgfqpoint{1.429987in}{1.373229in}}%
\pgfpathlineto{\pgfqpoint{1.428602in}{1.373346in}}%
\pgfpathlineto{\pgfqpoint{1.430345in}{1.395850in}}%
\pgfpathlineto{\pgfqpoint{1.430688in}{1.409985in}}%
\pgfpathlineto{\pgfqpoint{1.471839in}{1.406868in}}%
\pgfpathlineto{\pgfqpoint{1.521452in}{1.403789in}}%
\pgfpathlineto{\pgfqpoint{1.521365in}{1.389410in}}%
\pgfpathlineto{\pgfqpoint{1.520062in}{1.366510in}}%
\pgfpathlineto{\pgfqpoint{1.519516in}{1.343574in}}%
\pgfpathlineto{\pgfqpoint{1.486374in}{1.345738in}}%
\pgfpathlineto{\pgfqpoint{1.485080in}{1.322839in}}%
\pgfpathlineto{\pgfqpoint{1.457408in}{1.324892in}}%
\pgfpathlineto{\pgfqpoint{1.455754in}{1.302023in}}%
\pgfpathclose%
\pgfusepath{fill}%
\end{pgfscope}%
\begin{pgfscope}%
\pgfpathrectangle{\pgfqpoint{0.100000in}{0.100000in}}{\pgfqpoint{3.007045in}{1.925000in}}%
\pgfusepath{clip}%
\pgfsetbuttcap%
\pgfsetmiterjoin%
\definecolor{currentfill}{rgb}{0.510588,0.732303,0.858839}%
\pgfsetfillcolor{currentfill}%
\pgfsetlinewidth{0.000000pt}%
\definecolor{currentstroke}{rgb}{0.000000,0.000000,0.000000}%
\pgfsetstrokecolor{currentstroke}%
\pgfsetstrokeopacity{0.000000}%
\pgfsetdash{}{0pt}%
\pgfpathmoveto{\pgfqpoint{2.292953in}{1.341716in}}%
\pgfpathlineto{\pgfqpoint{2.273138in}{1.338797in}}%
\pgfpathlineto{\pgfqpoint{2.271732in}{1.342879in}}%
\pgfpathlineto{\pgfqpoint{2.269373in}{1.363358in}}%
\pgfpathlineto{\pgfqpoint{2.275068in}{1.363856in}}%
\pgfpathlineto{\pgfqpoint{2.272022in}{1.386856in}}%
\pgfpathlineto{\pgfqpoint{2.299946in}{1.390294in}}%
\pgfpathlineto{\pgfqpoint{2.303185in}{1.367262in}}%
\pgfpathlineto{\pgfqpoint{2.320460in}{1.370056in}}%
\pgfpathlineto{\pgfqpoint{2.323909in}{1.346635in}}%
\pgfpathlineto{\pgfqpoint{2.292953in}{1.341716in}}%
\pgfpathclose%
\pgfusepath{fill}%
\end{pgfscope}%
\begin{pgfscope}%
\pgfpathrectangle{\pgfqpoint{0.100000in}{0.100000in}}{\pgfqpoint{3.007045in}{1.925000in}}%
\pgfusepath{clip}%
\pgfsetbuttcap%
\pgfsetmiterjoin%
\definecolor{currentfill}{rgb}{0.516863,0.735748,0.860192}%
\pgfsetfillcolor{currentfill}%
\pgfsetlinewidth{0.000000pt}%
\definecolor{currentstroke}{rgb}{0.000000,0.000000,0.000000}%
\pgfsetstrokecolor{currentstroke}%
\pgfsetstrokeopacity{0.000000}%
\pgfsetdash{}{0pt}%
\pgfpathmoveto{\pgfqpoint{2.741068in}{1.332531in}}%
\pgfpathlineto{\pgfqpoint{2.720167in}{1.343542in}}%
\pgfpathlineto{\pgfqpoint{2.725895in}{1.348689in}}%
\pgfpathlineto{\pgfqpoint{2.711768in}{1.357635in}}%
\pgfpathlineto{\pgfqpoint{2.723365in}{1.366527in}}%
\pgfpathlineto{\pgfqpoint{2.720526in}{1.370514in}}%
\pgfpathlineto{\pgfqpoint{2.723613in}{1.374911in}}%
\pgfpathlineto{\pgfqpoint{2.727132in}{1.374300in}}%
\pgfpathlineto{\pgfqpoint{2.735363in}{1.367506in}}%
\pgfpathlineto{\pgfqpoint{2.732190in}{1.364248in}}%
\pgfpathlineto{\pgfqpoint{2.738618in}{1.356791in}}%
\pgfpathlineto{\pgfqpoint{2.746998in}{1.361954in}}%
\pgfpathlineto{\pgfqpoint{2.753651in}{1.370464in}}%
\pgfpathlineto{\pgfqpoint{2.758253in}{1.365099in}}%
\pgfpathlineto{\pgfqpoint{2.756767in}{1.359194in}}%
\pgfpathlineto{\pgfqpoint{2.753885in}{1.357653in}}%
\pgfpathlineto{\pgfqpoint{2.755066in}{1.346497in}}%
\pgfpathlineto{\pgfqpoint{2.746311in}{1.336376in}}%
\pgfpathlineto{\pgfqpoint{2.741068in}{1.332531in}}%
\pgfpathclose%
\pgfusepath{fill}%
\end{pgfscope}%
\begin{pgfscope}%
\pgfpathrectangle{\pgfqpoint{0.100000in}{0.100000in}}{\pgfqpoint{3.007045in}{1.925000in}}%
\pgfusepath{clip}%
\pgfsetbuttcap%
\pgfsetmiterjoin%
\definecolor{currentfill}{rgb}{0.447843,0.697855,0.845306}%
\pgfsetfillcolor{currentfill}%
\pgfsetlinewidth{0.000000pt}%
\definecolor{currentstroke}{rgb}{0.000000,0.000000,0.000000}%
\pgfsetstrokecolor{currentstroke}%
\pgfsetstrokeopacity{0.000000}%
\pgfsetdash{}{0pt}%
\pgfpathmoveto{\pgfqpoint{1.538952in}{0.605702in}}%
\pgfpathlineto{\pgfqpoint{1.546174in}{0.604074in}}%
\pgfpathlineto{\pgfqpoint{1.544968in}{0.575956in}}%
\pgfpathlineto{\pgfqpoint{1.525500in}{0.576932in}}%
\pgfpathlineto{\pgfqpoint{1.515156in}{0.577420in}}%
\pgfpathlineto{\pgfqpoint{1.516589in}{0.606731in}}%
\pgfpathlineto{\pgfqpoint{1.538952in}{0.605702in}}%
\pgfpathclose%
\pgfusepath{fill}%
\end{pgfscope}%
\begin{pgfscope}%
\pgfpathrectangle{\pgfqpoint{0.100000in}{0.100000in}}{\pgfqpoint{3.007045in}{1.925000in}}%
\pgfusepath{clip}%
\pgfsetbuttcap%
\pgfsetmiterjoin%
\definecolor{currentfill}{rgb}{0.447843,0.697855,0.845306}%
\pgfsetfillcolor{currentfill}%
\pgfsetlinewidth{0.000000pt}%
\definecolor{currentstroke}{rgb}{0.000000,0.000000,0.000000}%
\pgfsetstrokecolor{currentstroke}%
\pgfsetstrokeopacity{0.000000}%
\pgfsetdash{}{0pt}%
\pgfpathmoveto{\pgfqpoint{1.880439in}{1.202637in}}%
\pgfpathlineto{\pgfqpoint{1.854945in}{1.201835in}}%
\pgfpathlineto{\pgfqpoint{1.853670in}{1.237862in}}%
\pgfpathlineto{\pgfqpoint{1.886676in}{1.239179in}}%
\pgfpathlineto{\pgfqpoint{1.889157in}{1.227613in}}%
\pgfpathlineto{\pgfqpoint{1.888599in}{1.223016in}}%
\pgfpathlineto{\pgfqpoint{1.879993in}{1.222750in}}%
\pgfpathlineto{\pgfqpoint{1.880439in}{1.202637in}}%
\pgfpathclose%
\pgfusepath{fill}%
\end{pgfscope}%
\begin{pgfscope}%
\pgfpathrectangle{\pgfqpoint{0.100000in}{0.100000in}}{\pgfqpoint{3.007045in}{1.925000in}}%
\pgfusepath{clip}%
\pgfsetbuttcap%
\pgfsetmiterjoin%
\definecolor{currentfill}{rgb}{0.657132,0.808105,0.895486}%
\pgfsetfillcolor{currentfill}%
\pgfsetlinewidth{0.000000pt}%
\definecolor{currentstroke}{rgb}{0.000000,0.000000,0.000000}%
\pgfsetstrokecolor{currentstroke}%
\pgfsetstrokeopacity{0.000000}%
\pgfsetdash{}{0pt}%
\pgfpathmoveto{\pgfqpoint{2.612770in}{1.112672in}}%
\pgfpathlineto{\pgfqpoint{2.609095in}{1.111445in}}%
\pgfpathlineto{\pgfqpoint{2.602817in}{1.120583in}}%
\pgfpathlineto{\pgfqpoint{2.598994in}{1.120903in}}%
\pgfpathlineto{\pgfqpoint{2.593692in}{1.126779in}}%
\pgfpathlineto{\pgfqpoint{2.587434in}{1.124379in}}%
\pgfpathlineto{\pgfqpoint{2.584797in}{1.114957in}}%
\pgfpathlineto{\pgfqpoint{2.579736in}{1.109975in}}%
\pgfpathlineto{\pgfqpoint{2.578351in}{1.105297in}}%
\pgfpathlineto{\pgfqpoint{2.567398in}{1.112031in}}%
\pgfpathlineto{\pgfqpoint{2.566770in}{1.118700in}}%
\pgfpathlineto{\pgfqpoint{2.559027in}{1.119923in}}%
\pgfpathlineto{\pgfqpoint{2.554928in}{1.112611in}}%
\pgfpathlineto{\pgfqpoint{2.544389in}{1.103816in}}%
\pgfpathlineto{\pgfqpoint{2.535565in}{1.108737in}}%
\pgfpathlineto{\pgfqpoint{2.538295in}{1.117062in}}%
\pgfpathlineto{\pgfqpoint{2.547549in}{1.130953in}}%
\pgfpathlineto{\pgfqpoint{2.548827in}{1.135761in}}%
\pgfpathlineto{\pgfqpoint{2.549531in}{1.143698in}}%
\pgfpathlineto{\pgfqpoint{2.555307in}{1.150523in}}%
\pgfpathlineto{\pgfqpoint{2.553841in}{1.152723in}}%
\pgfpathlineto{\pgfqpoint{2.558583in}{1.161971in}}%
\pgfpathlineto{\pgfqpoint{2.558875in}{1.172717in}}%
\pgfpathlineto{\pgfqpoint{2.565022in}{1.171056in}}%
\pgfpathlineto{\pgfqpoint{2.569128in}{1.165501in}}%
\pgfpathlineto{\pgfqpoint{2.577995in}{1.164006in}}%
\pgfpathlineto{\pgfqpoint{2.581715in}{1.169068in}}%
\pgfpathlineto{\pgfqpoint{2.609006in}{1.155780in}}%
\pgfpathlineto{\pgfqpoint{2.608923in}{1.148127in}}%
\pgfpathlineto{\pgfqpoint{2.606282in}{1.144642in}}%
\pgfpathlineto{\pgfqpoint{2.616773in}{1.129743in}}%
\pgfpathlineto{\pgfqpoint{2.617784in}{1.122931in}}%
\pgfpathlineto{\pgfqpoint{2.613630in}{1.118818in}}%
\pgfpathlineto{\pgfqpoint{2.612770in}{1.112672in}}%
\pgfpathclose%
\pgfusepath{fill}%
\end{pgfscope}%
\begin{pgfscope}%
\pgfpathrectangle{\pgfqpoint{0.100000in}{0.100000in}}{\pgfqpoint{3.007045in}{1.925000in}}%
\pgfusepath{clip}%
\pgfsetbuttcap%
\pgfsetmiterjoin%
\definecolor{currentfill}{rgb}{0.447843,0.697855,0.845306}%
\pgfsetfillcolor{currentfill}%
\pgfsetlinewidth{0.000000pt}%
\definecolor{currentstroke}{rgb}{0.000000,0.000000,0.000000}%
\pgfsetstrokecolor{currentstroke}%
\pgfsetstrokeopacity{0.000000}%
\pgfsetdash{}{0pt}%
\pgfpathmoveto{\pgfqpoint{1.525500in}{0.576932in}}%
\pgfpathlineto{\pgfqpoint{1.524943in}{0.562745in}}%
\pgfpathlineto{\pgfqpoint{1.498977in}{0.564304in}}%
\pgfpathlineto{\pgfqpoint{1.478167in}{0.565413in}}%
\pgfpathlineto{\pgfqpoint{1.479753in}{0.593339in}}%
\pgfpathlineto{\pgfqpoint{1.481225in}{0.618416in}}%
\pgfpathlineto{\pgfqpoint{1.510333in}{0.616812in}}%
\pgfpathlineto{\pgfqpoint{1.509816in}{0.607078in}}%
\pgfpathlineto{\pgfqpoint{1.516589in}{0.606731in}}%
\pgfpathlineto{\pgfqpoint{1.515156in}{0.577420in}}%
\pgfpathlineto{\pgfqpoint{1.525500in}{0.576932in}}%
\pgfpathclose%
\pgfusepath{fill}%
\end{pgfscope}%
\begin{pgfscope}%
\pgfpathrectangle{\pgfqpoint{0.100000in}{0.100000in}}{\pgfqpoint{3.007045in}{1.925000in}}%
\pgfusepath{clip}%
\pgfsetbuttcap%
\pgfsetmiterjoin%
\definecolor{currentfill}{rgb}{0.721107,0.835294,0.917878}%
\pgfsetfillcolor{currentfill}%
\pgfsetlinewidth{0.000000pt}%
\definecolor{currentstroke}{rgb}{0.000000,0.000000,0.000000}%
\pgfsetstrokecolor{currentstroke}%
\pgfsetstrokeopacity{0.000000}%
\pgfsetdash{}{0pt}%
\pgfpathmoveto{\pgfqpoint{2.671279in}{0.875019in}}%
\pgfpathlineto{\pgfqpoint{2.637551in}{0.899069in}}%
\pgfpathlineto{\pgfqpoint{2.641884in}{0.909457in}}%
\pgfpathlineto{\pgfqpoint{2.647369in}{0.914791in}}%
\pgfpathlineto{\pgfqpoint{2.647207in}{0.927059in}}%
\pgfpathlineto{\pgfqpoint{2.640458in}{0.935526in}}%
\pgfpathlineto{\pgfqpoint{2.645656in}{0.937482in}}%
\pgfpathlineto{\pgfqpoint{2.661874in}{0.940850in}}%
\pgfpathlineto{\pgfqpoint{2.668672in}{0.934954in}}%
\pgfpathlineto{\pgfqpoint{2.678361in}{0.924174in}}%
\pgfpathlineto{\pgfqpoint{2.683843in}{0.935913in}}%
\pgfpathlineto{\pgfqpoint{2.706932in}{0.940219in}}%
\pgfpathlineto{\pgfqpoint{2.714846in}{0.926915in}}%
\pgfpathlineto{\pgfqpoint{2.719169in}{0.923665in}}%
\pgfpathlineto{\pgfqpoint{2.709222in}{0.910276in}}%
\pgfpathlineto{\pgfqpoint{2.705650in}{0.901507in}}%
\pgfpathlineto{\pgfqpoint{2.702584in}{0.882702in}}%
\pgfpathlineto{\pgfqpoint{2.691067in}{0.882886in}}%
\pgfpathlineto{\pgfqpoint{2.680626in}{0.880256in}}%
\pgfpathlineto{\pgfqpoint{2.671279in}{0.875019in}}%
\pgfpathclose%
\pgfusepath{fill}%
\end{pgfscope}%
\begin{pgfscope}%
\pgfpathrectangle{\pgfqpoint{0.100000in}{0.100000in}}{\pgfqpoint{3.007045in}{1.925000in}}%
\pgfusepath{clip}%
\pgfsetbuttcap%
\pgfsetmiterjoin%
\definecolor{currentfill}{rgb}{0.541961,0.749527,0.865606}%
\pgfsetfillcolor{currentfill}%
\pgfsetlinewidth{0.000000pt}%
\definecolor{currentstroke}{rgb}{0.000000,0.000000,0.000000}%
\pgfsetstrokecolor{currentstroke}%
\pgfsetstrokeopacity{0.000000}%
\pgfsetdash{}{0pt}%
\pgfpathmoveto{\pgfqpoint{0.540209in}{1.259221in}}%
\pgfpathlineto{\pgfqpoint{0.533766in}{1.269179in}}%
\pgfpathlineto{\pgfqpoint{0.536090in}{1.281820in}}%
\pgfpathlineto{\pgfqpoint{0.533368in}{1.282853in}}%
\pgfpathlineto{\pgfqpoint{0.536518in}{1.297107in}}%
\pgfpathlineto{\pgfqpoint{0.542583in}{1.297900in}}%
\pgfpathlineto{\pgfqpoint{0.544151in}{1.303815in}}%
\pgfpathlineto{\pgfqpoint{0.522474in}{1.309901in}}%
\pgfpathlineto{\pgfqpoint{0.522542in}{1.311873in}}%
\pgfpathlineto{\pgfqpoint{0.510601in}{1.315136in}}%
\pgfpathlineto{\pgfqpoint{0.516565in}{1.336204in}}%
\pgfpathlineto{\pgfqpoint{0.521615in}{1.353758in}}%
\pgfpathlineto{\pgfqpoint{0.535789in}{1.402977in}}%
\pgfpathlineto{\pgfqpoint{0.547844in}{1.446415in}}%
\pgfpathlineto{\pgfqpoint{0.562393in}{1.497838in}}%
\pgfpathlineto{\pgfqpoint{0.592563in}{1.489427in}}%
\pgfpathlineto{\pgfqpoint{0.594261in}{1.488948in}}%
\pgfpathlineto{\pgfqpoint{0.584193in}{1.452083in}}%
\pgfpathlineto{\pgfqpoint{0.575970in}{1.423296in}}%
\pgfpathlineto{\pgfqpoint{0.576936in}{1.423028in}}%
\pgfpathlineto{\pgfqpoint{0.569070in}{1.394980in}}%
\pgfpathlineto{\pgfqpoint{0.567831in}{1.395284in}}%
\pgfpathlineto{\pgfqpoint{0.558741in}{1.362309in}}%
\pgfpathlineto{\pgfqpoint{0.594471in}{1.352638in}}%
\pgfpathlineto{\pgfqpoint{0.646486in}{1.339327in}}%
\pgfpathlineto{\pgfqpoint{0.648824in}{1.333373in}}%
\pgfpathlineto{\pgfqpoint{0.644679in}{1.322619in}}%
\pgfpathlineto{\pgfqpoint{0.644313in}{1.314171in}}%
\pgfpathlineto{\pgfqpoint{0.641368in}{1.309018in}}%
\pgfpathlineto{\pgfqpoint{0.632081in}{1.306373in}}%
\pgfpathlineto{\pgfqpoint{0.625081in}{1.300716in}}%
\pgfpathlineto{\pgfqpoint{0.625302in}{1.292738in}}%
\pgfpathlineto{\pgfqpoint{0.621234in}{1.289668in}}%
\pgfpathlineto{\pgfqpoint{0.620230in}{1.284320in}}%
\pgfpathlineto{\pgfqpoint{0.615468in}{1.284198in}}%
\pgfpathlineto{\pgfqpoint{0.562826in}{1.298061in}}%
\pgfpathlineto{\pgfqpoint{0.556383in}{1.290907in}}%
\pgfpathlineto{\pgfqpoint{0.554814in}{1.285052in}}%
\pgfpathlineto{\pgfqpoint{0.560395in}{1.283604in}}%
\pgfpathlineto{\pgfqpoint{0.552667in}{1.255842in}}%
\pgfpathlineto{\pgfqpoint{0.540209in}{1.259221in}}%
\pgfpathclose%
\pgfusepath{fill}%
\end{pgfscope}%
\begin{pgfscope}%
\pgfpathrectangle{\pgfqpoint{0.100000in}{0.100000in}}{\pgfqpoint{3.007045in}{1.925000in}}%
\pgfusepath{clip}%
\pgfsetbuttcap%
\pgfsetmiterjoin%
\definecolor{currentfill}{rgb}{0.290980,0.594510,0.789020}%
\pgfsetfillcolor{currentfill}%
\pgfsetlinewidth{0.000000pt}%
\definecolor{currentstroke}{rgb}{0.000000,0.000000,0.000000}%
\pgfsetstrokecolor{currentstroke}%
\pgfsetstrokeopacity{0.000000}%
\pgfsetdash{}{0pt}%
\pgfpathmoveto{\pgfqpoint{1.591111in}{1.462290in}}%
\pgfpathlineto{\pgfqpoint{1.585428in}{1.462527in}}%
\pgfpathlineto{\pgfqpoint{1.586215in}{1.479799in}}%
\pgfpathlineto{\pgfqpoint{1.596903in}{1.479359in}}%
\pgfpathlineto{\pgfqpoint{1.597846in}{1.508282in}}%
\pgfpathlineto{\pgfqpoint{1.632109in}{1.506861in}}%
\pgfpathlineto{\pgfqpoint{1.637825in}{1.506641in}}%
\pgfpathlineto{\pgfqpoint{1.637786in}{1.500890in}}%
\pgfpathlineto{\pgfqpoint{1.637041in}{1.477807in}}%
\pgfpathlineto{\pgfqpoint{1.614383in}{1.478679in}}%
\pgfpathlineto{\pgfqpoint{1.613769in}{1.461576in}}%
\pgfpathlineto{\pgfqpoint{1.591111in}{1.462290in}}%
\pgfpathclose%
\pgfusepath{fill}%
\end{pgfscope}%
\begin{pgfscope}%
\pgfpathrectangle{\pgfqpoint{0.100000in}{0.100000in}}{\pgfqpoint{3.007045in}{1.925000in}}%
\pgfusepath{clip}%
\pgfsetbuttcap%
\pgfsetmiterjoin%
\definecolor{currentfill}{rgb}{0.730950,0.839477,0.921323}%
\pgfsetfillcolor{currentfill}%
\pgfsetlinewidth{0.000000pt}%
\definecolor{currentstroke}{rgb}{0.000000,0.000000,0.000000}%
\pgfsetstrokecolor{currentstroke}%
\pgfsetstrokeopacity{0.000000}%
\pgfsetdash{}{0pt}%
\pgfpathmoveto{\pgfqpoint{2.735476in}{1.140922in}}%
\pgfpathlineto{\pgfqpoint{2.729649in}{1.140417in}}%
\pgfpathlineto{\pgfqpoint{2.725113in}{1.146693in}}%
\pgfpathlineto{\pgfqpoint{2.720436in}{1.148513in}}%
\pgfpathlineto{\pgfqpoint{2.716860in}{1.144435in}}%
\pgfpathlineto{\pgfqpoint{2.713398in}{1.148167in}}%
\pgfpathlineto{\pgfqpoint{2.706818in}{1.146620in}}%
\pgfpathlineto{\pgfqpoint{2.701679in}{1.149182in}}%
\pgfpathlineto{\pgfqpoint{2.697515in}{1.156406in}}%
\pgfpathlineto{\pgfqpoint{2.690990in}{1.162746in}}%
\pgfpathlineto{\pgfqpoint{2.695805in}{1.170685in}}%
\pgfpathlineto{\pgfqpoint{2.697143in}{1.177380in}}%
\pgfpathlineto{\pgfqpoint{2.704729in}{1.171795in}}%
\pgfpathlineto{\pgfqpoint{2.718443in}{1.173001in}}%
\pgfpathlineto{\pgfqpoint{2.728305in}{1.164174in}}%
\pgfpathlineto{\pgfqpoint{2.740946in}{1.160176in}}%
\pgfpathlineto{\pgfqpoint{2.739016in}{1.152520in}}%
\pgfpathlineto{\pgfqpoint{2.740280in}{1.146999in}}%
\pgfpathlineto{\pgfqpoint{2.738357in}{1.140814in}}%
\pgfpathlineto{\pgfqpoint{2.735476in}{1.140922in}}%
\pgfpathclose%
\pgfusepath{fill}%
\end{pgfscope}%
\begin{pgfscope}%
\pgfpathrectangle{\pgfqpoint{0.100000in}{0.100000in}}{\pgfqpoint{3.007045in}{1.925000in}}%
\pgfusepath{clip}%
\pgfsetbuttcap%
\pgfsetmiterjoin%
\definecolor{currentfill}{rgb}{0.435294,0.690965,0.842599}%
\pgfsetfillcolor{currentfill}%
\pgfsetlinewidth{0.000000pt}%
\definecolor{currentstroke}{rgb}{0.000000,0.000000,0.000000}%
\pgfsetstrokecolor{currentstroke}%
\pgfsetstrokeopacity{0.000000}%
\pgfsetdash{}{0pt}%
\pgfpathmoveto{\pgfqpoint{1.631910in}{1.063521in}}%
\pgfpathlineto{\pgfqpoint{1.631737in}{1.051813in}}%
\pgfpathlineto{\pgfqpoint{1.591838in}{1.053004in}}%
\pgfpathlineto{\pgfqpoint{1.591676in}{1.059068in}}%
\pgfpathlineto{\pgfqpoint{1.593165in}{1.116517in}}%
\pgfpathlineto{\pgfqpoint{1.593668in}{1.127954in}}%
\pgfpathlineto{\pgfqpoint{1.622153in}{1.126974in}}%
\pgfpathlineto{\pgfqpoint{1.621825in}{1.109765in}}%
\pgfpathlineto{\pgfqpoint{1.621015in}{1.081045in}}%
\pgfpathlineto{\pgfqpoint{1.632419in}{1.080689in}}%
\pgfpathlineto{\pgfqpoint{1.631910in}{1.063521in}}%
\pgfpathclose%
\pgfusepath{fill}%
\end{pgfscope}%
\begin{pgfscope}%
\pgfpathrectangle{\pgfqpoint{0.100000in}{0.100000in}}{\pgfqpoint{3.007045in}{1.925000in}}%
\pgfusepath{clip}%
\pgfsetbuttcap%
\pgfsetmiterjoin%
\definecolor{currentfill}{rgb}{0.280892,0.587620,0.785083}%
\pgfsetfillcolor{currentfill}%
\pgfsetlinewidth{0.000000pt}%
\definecolor{currentstroke}{rgb}{0.000000,0.000000,0.000000}%
\pgfsetstrokecolor{currentstroke}%
\pgfsetstrokeopacity{0.000000}%
\pgfsetdash{}{0pt}%
\pgfpathmoveto{\pgfqpoint{2.397002in}{0.647345in}}%
\pgfpathlineto{\pgfqpoint{2.375927in}{0.644931in}}%
\pgfpathlineto{\pgfqpoint{2.366514in}{0.643850in}}%
\pgfpathlineto{\pgfqpoint{2.365478in}{0.655570in}}%
\pgfpathlineto{\pgfqpoint{2.359552in}{0.655091in}}%
\pgfpathlineto{\pgfqpoint{2.358460in}{0.666668in}}%
\pgfpathlineto{\pgfqpoint{2.366771in}{0.667222in}}%
\pgfpathlineto{\pgfqpoint{2.370955in}{0.669671in}}%
\pgfpathlineto{\pgfqpoint{2.367869in}{0.675418in}}%
\pgfpathlineto{\pgfqpoint{2.367559in}{0.680287in}}%
\pgfpathlineto{\pgfqpoint{2.362322in}{0.683671in}}%
\pgfpathlineto{\pgfqpoint{2.357840in}{0.689197in}}%
\pgfpathlineto{\pgfqpoint{2.356903in}{0.698990in}}%
\pgfpathlineto{\pgfqpoint{2.394856in}{0.702917in}}%
\pgfpathlineto{\pgfqpoint{2.395249in}{0.698742in}}%
\pgfpathlineto{\pgfqpoint{2.402438in}{0.696626in}}%
\pgfpathlineto{\pgfqpoint{2.404240in}{0.684768in}}%
\pgfpathlineto{\pgfqpoint{2.408654in}{0.685257in}}%
\pgfpathlineto{\pgfqpoint{2.412962in}{0.682194in}}%
\pgfpathlineto{\pgfqpoint{2.414666in}{0.666517in}}%
\pgfpathlineto{\pgfqpoint{2.395188in}{0.664380in}}%
\pgfpathlineto{\pgfqpoint{2.397002in}{0.647345in}}%
\pgfpathclose%
\pgfusepath{fill}%
\end{pgfscope}%
\begin{pgfscope}%
\pgfpathrectangle{\pgfqpoint{0.100000in}{0.100000in}}{\pgfqpoint{3.007045in}{1.925000in}}%
\pgfusepath{clip}%
\pgfsetbuttcap%
\pgfsetmiterjoin%
\definecolor{currentfill}{rgb}{0.686659,0.820654,0.905821}%
\pgfsetfillcolor{currentfill}%
\pgfsetlinewidth{0.000000pt}%
\definecolor{currentstroke}{rgb}{0.000000,0.000000,0.000000}%
\pgfsetstrokecolor{currentstroke}%
\pgfsetstrokeopacity{0.000000}%
\pgfsetdash{}{0pt}%
\pgfpathmoveto{\pgfqpoint{2.209634in}{1.589574in}}%
\pgfpathlineto{\pgfqpoint{2.206269in}{1.590978in}}%
\pgfpathlineto{\pgfqpoint{2.208079in}{1.601957in}}%
\pgfpathlineto{\pgfqpoint{2.211439in}{1.600059in}}%
\pgfpathlineto{\pgfqpoint{2.212615in}{1.592362in}}%
\pgfpathlineto{\pgfqpoint{2.209634in}{1.589574in}}%
\pgfpathclose%
\pgfusepath{fill}%
\end{pgfscope}%
\begin{pgfscope}%
\pgfpathrectangle{\pgfqpoint{0.100000in}{0.100000in}}{\pgfqpoint{3.007045in}{1.925000in}}%
\pgfusepath{clip}%
\pgfsetbuttcap%
\pgfsetmiterjoin%
\definecolor{currentfill}{rgb}{0.686659,0.820654,0.905821}%
\pgfsetfillcolor{currentfill}%
\pgfsetlinewidth{0.000000pt}%
\definecolor{currentstroke}{rgb}{0.000000,0.000000,0.000000}%
\pgfsetstrokecolor{currentstroke}%
\pgfsetstrokeopacity{0.000000}%
\pgfsetdash{}{0pt}%
\pgfpathmoveto{\pgfqpoint{2.271744in}{1.600141in}}%
\pgfpathlineto{\pgfqpoint{2.269589in}{1.599880in}}%
\pgfpathlineto{\pgfqpoint{2.273175in}{1.571785in}}%
\pgfpathlineto{\pgfqpoint{2.250552in}{1.569368in}}%
\pgfpathlineto{\pgfqpoint{2.251176in}{1.563619in}}%
\pgfpathlineto{\pgfqpoint{2.245496in}{1.563087in}}%
\pgfpathlineto{\pgfqpoint{2.228428in}{1.561303in}}%
\pgfpathlineto{\pgfqpoint{2.227849in}{1.567100in}}%
\pgfpathlineto{\pgfqpoint{2.220169in}{1.566322in}}%
\pgfpathlineto{\pgfqpoint{2.223475in}{1.574310in}}%
\pgfpathlineto{\pgfqpoint{2.230581in}{1.578556in}}%
\pgfpathlineto{\pgfqpoint{2.234669in}{1.583276in}}%
\pgfpathlineto{\pgfqpoint{2.231732in}{1.586161in}}%
\pgfpathlineto{\pgfqpoint{2.230119in}{1.592502in}}%
\pgfpathlineto{\pgfqpoint{2.237120in}{1.602372in}}%
\pgfpathlineto{\pgfqpoint{2.243327in}{1.604920in}}%
\pgfpathlineto{\pgfqpoint{2.244332in}{1.607881in}}%
\pgfpathlineto{\pgfqpoint{2.259869in}{1.600439in}}%
\pgfpathlineto{\pgfqpoint{2.265685in}{1.601992in}}%
\pgfpathlineto{\pgfqpoint{2.271744in}{1.600141in}}%
\pgfpathclose%
\pgfusepath{fill}%
\end{pgfscope}%
\begin{pgfscope}%
\pgfpathrectangle{\pgfqpoint{0.100000in}{0.100000in}}{\pgfqpoint{3.007045in}{1.925000in}}%
\pgfusepath{clip}%
\pgfsetbuttcap%
\pgfsetmiterjoin%
\definecolor{currentfill}{rgb}{0.652211,0.806013,0.893764}%
\pgfsetfillcolor{currentfill}%
\pgfsetlinewidth{0.000000pt}%
\definecolor{currentstroke}{rgb}{0.000000,0.000000,0.000000}%
\pgfsetstrokecolor{currentstroke}%
\pgfsetstrokeopacity{0.000000}%
\pgfsetdash{}{0pt}%
\pgfpathmoveto{\pgfqpoint{2.415134in}{0.621502in}}%
\pgfpathlineto{\pgfqpoint{2.422704in}{0.621971in}}%
\pgfpathlineto{\pgfqpoint{2.423706in}{0.617445in}}%
\pgfpathlineto{\pgfqpoint{2.416778in}{0.611640in}}%
\pgfpathlineto{\pgfqpoint{2.418230in}{0.606839in}}%
\pgfpathlineto{\pgfqpoint{2.413904in}{0.601638in}}%
\pgfpathlineto{\pgfqpoint{2.415159in}{0.598741in}}%
\pgfpathlineto{\pgfqpoint{2.408711in}{0.592974in}}%
\pgfpathlineto{\pgfqpoint{2.405641in}{0.582074in}}%
\pgfpathlineto{\pgfqpoint{2.395790in}{0.579657in}}%
\pgfpathlineto{\pgfqpoint{2.389487in}{0.580324in}}%
\pgfpathlineto{\pgfqpoint{2.384731in}{0.575233in}}%
\pgfpathlineto{\pgfqpoint{2.386694in}{0.571228in}}%
\pgfpathlineto{\pgfqpoint{2.374564in}{0.573201in}}%
\pgfpathlineto{\pgfqpoint{2.347979in}{0.570392in}}%
\pgfpathlineto{\pgfqpoint{2.340428in}{0.571556in}}%
\pgfpathlineto{\pgfqpoint{2.339480in}{0.580648in}}%
\pgfpathlineto{\pgfqpoint{2.324742in}{0.580445in}}%
\pgfpathlineto{\pgfqpoint{2.322486in}{0.604682in}}%
\pgfpathlineto{\pgfqpoint{2.316671in}{0.604111in}}%
\pgfpathlineto{\pgfqpoint{2.317482in}{0.595459in}}%
\pgfpathlineto{\pgfqpoint{2.296561in}{0.593554in}}%
\pgfpathlineto{\pgfqpoint{2.288959in}{0.589362in}}%
\pgfpathlineto{\pgfqpoint{2.295880in}{0.596792in}}%
\pgfpathlineto{\pgfqpoint{2.291737in}{0.606164in}}%
\pgfpathlineto{\pgfqpoint{2.294588in}{0.610608in}}%
\pgfpathlineto{\pgfqpoint{2.284322in}{0.609983in}}%
\pgfpathlineto{\pgfqpoint{2.282576in}{0.629115in}}%
\pgfpathlineto{\pgfqpoint{2.341073in}{0.635732in}}%
\pgfpathlineto{\pgfqpoint{2.345288in}{0.628691in}}%
\pgfpathlineto{\pgfqpoint{2.346074in}{0.624253in}}%
\pgfpathlineto{\pgfqpoint{2.351014in}{0.617483in}}%
\pgfpathlineto{\pgfqpoint{2.378692in}{0.619214in}}%
\pgfpathlineto{\pgfqpoint{2.415134in}{0.621502in}}%
\pgfpathclose%
\pgfusepath{fill}%
\end{pgfscope}%
\begin{pgfscope}%
\pgfpathrectangle{\pgfqpoint{0.100000in}{0.100000in}}{\pgfqpoint{3.007045in}{1.925000in}}%
\pgfusepath{clip}%
\pgfsetbuttcap%
\pgfsetmiterjoin%
\definecolor{currentfill}{rgb}{0.114802,0.424437,0.695194}%
\pgfsetfillcolor{currentfill}%
\pgfsetlinewidth{0.000000pt}%
\definecolor{currentstroke}{rgb}{0.000000,0.000000,0.000000}%
\pgfsetstrokecolor{currentstroke}%
\pgfsetstrokeopacity{0.000000}%
\pgfsetdash{}{0pt}%
\pgfpathmoveto{\pgfqpoint{1.448770in}{1.649792in}}%
\pgfpathlineto{\pgfqpoint{1.449237in}{1.655533in}}%
\pgfpathlineto{\pgfqpoint{1.466420in}{1.654245in}}%
\pgfpathlineto{\pgfqpoint{1.465987in}{1.648477in}}%
\pgfpathlineto{\pgfqpoint{1.485073in}{1.647065in}}%
\pgfpathlineto{\pgfqpoint{1.483880in}{1.629904in}}%
\pgfpathlineto{\pgfqpoint{1.495332in}{1.629121in}}%
\pgfpathlineto{\pgfqpoint{1.495994in}{1.623275in}}%
\pgfpathlineto{\pgfqpoint{1.482806in}{1.614036in}}%
\pgfpathlineto{\pgfqpoint{1.480364in}{1.610710in}}%
\pgfpathlineto{\pgfqpoint{1.470594in}{1.606845in}}%
\pgfpathlineto{\pgfqpoint{1.459052in}{1.608999in}}%
\pgfpathlineto{\pgfqpoint{1.455103in}{1.611872in}}%
\pgfpathlineto{\pgfqpoint{1.450280in}{1.611196in}}%
\pgfpathlineto{\pgfqpoint{1.451489in}{1.626431in}}%
\pgfpathlineto{\pgfqpoint{1.449255in}{1.626605in}}%
\pgfpathlineto{\pgfqpoint{1.450846in}{1.649617in}}%
\pgfpathlineto{\pgfqpoint{1.448770in}{1.649792in}}%
\pgfpathclose%
\pgfusepath{fill}%
\end{pgfscope}%
\begin{pgfscope}%
\pgfpathrectangle{\pgfqpoint{0.100000in}{0.100000in}}{\pgfqpoint{3.007045in}{1.925000in}}%
\pgfusepath{clip}%
\pgfsetbuttcap%
\pgfsetmiterjoin%
\definecolor{currentfill}{rgb}{0.351511,0.635848,0.812641}%
\pgfsetfillcolor{currentfill}%
\pgfsetlinewidth{0.000000pt}%
\definecolor{currentstroke}{rgb}{0.000000,0.000000,0.000000}%
\pgfsetstrokecolor{currentstroke}%
\pgfsetstrokeopacity{0.000000}%
\pgfsetdash{}{0pt}%
\pgfpathmoveto{\pgfqpoint{1.408351in}{1.098572in}}%
\pgfpathlineto{\pgfqpoint{1.408374in}{1.098989in}}%
\pgfpathlineto{\pgfqpoint{1.410528in}{1.127149in}}%
\pgfpathlineto{\pgfqpoint{1.439278in}{1.125169in}}%
\pgfpathlineto{\pgfqpoint{1.480135in}{1.122427in}}%
\pgfpathlineto{\pgfqpoint{1.478511in}{1.093767in}}%
\pgfpathlineto{\pgfqpoint{1.432949in}{1.096770in}}%
\pgfpathlineto{\pgfqpoint{1.408351in}{1.098572in}}%
\pgfpathclose%
\pgfusepath{fill}%
\end{pgfscope}%
\begin{pgfscope}%
\pgfpathrectangle{\pgfqpoint{0.100000in}{0.100000in}}{\pgfqpoint{3.007045in}{1.925000in}}%
\pgfusepath{clip}%
\pgfsetbuttcap%
\pgfsetmiterjoin%
\definecolor{currentfill}{rgb}{0.296025,0.597955,0.790988}%
\pgfsetfillcolor{currentfill}%
\pgfsetlinewidth{0.000000pt}%
\definecolor{currentstroke}{rgb}{0.000000,0.000000,0.000000}%
\pgfsetstrokecolor{currentstroke}%
\pgfsetstrokeopacity{0.000000}%
\pgfsetdash{}{0pt}%
\pgfpathmoveto{\pgfqpoint{1.753660in}{1.413597in}}%
\pgfpathlineto{\pgfqpoint{1.753560in}{1.390798in}}%
\pgfpathlineto{\pgfqpoint{1.730762in}{1.390911in}}%
\pgfpathlineto{\pgfqpoint{1.697951in}{1.391327in}}%
\pgfpathlineto{\pgfqpoint{1.699908in}{1.394517in}}%
\pgfpathlineto{\pgfqpoint{1.699433in}{1.400596in}}%
\pgfpathlineto{\pgfqpoint{1.702210in}{1.402774in}}%
\pgfpathlineto{\pgfqpoint{1.701304in}{1.411915in}}%
\pgfpathlineto{\pgfqpoint{1.698770in}{1.417142in}}%
\pgfpathlineto{\pgfqpoint{1.699167in}{1.423436in}}%
\pgfpathlineto{\pgfqpoint{1.695898in}{1.425999in}}%
\pgfpathlineto{\pgfqpoint{1.695709in}{1.430437in}}%
\pgfpathlineto{\pgfqpoint{1.702699in}{1.430317in}}%
\pgfpathlineto{\pgfqpoint{1.721870in}{1.430045in}}%
\pgfpathlineto{\pgfqpoint{1.731061in}{1.429926in}}%
\pgfpathlineto{\pgfqpoint{1.730861in}{1.413920in}}%
\pgfpathlineto{\pgfqpoint{1.753660in}{1.413597in}}%
\pgfpathclose%
\pgfusepath{fill}%
\end{pgfscope}%
\begin{pgfscope}%
\pgfpathrectangle{\pgfqpoint{0.100000in}{0.100000in}}{\pgfqpoint{3.007045in}{1.925000in}}%
\pgfusepath{clip}%
\pgfsetbuttcap%
\pgfsetmiterjoin%
\definecolor{currentfill}{rgb}{0.701423,0.826928,0.910988}%
\pgfsetfillcolor{currentfill}%
\pgfsetlinewidth{0.000000pt}%
\definecolor{currentstroke}{rgb}{0.000000,0.000000,0.000000}%
\pgfsetstrokecolor{currentstroke}%
\pgfsetstrokeopacity{0.000000}%
\pgfsetdash{}{0pt}%
\pgfpathmoveto{\pgfqpoint{2.867151in}{1.544262in}}%
\pgfpathlineto{\pgfqpoint{2.869055in}{1.532209in}}%
\pgfpathlineto{\pgfqpoint{2.871705in}{1.532829in}}%
\pgfpathlineto{\pgfqpoint{2.875745in}{1.523338in}}%
\pgfpathlineto{\pgfqpoint{2.879160in}{1.518963in}}%
\pgfpathlineto{\pgfqpoint{2.862374in}{1.515277in}}%
\pgfpathlineto{\pgfqpoint{2.831656in}{1.508759in}}%
\pgfpathlineto{\pgfqpoint{2.829097in}{1.520876in}}%
\pgfpathlineto{\pgfqpoint{2.824976in}{1.520081in}}%
\pgfpathlineto{\pgfqpoint{2.822507in}{1.532326in}}%
\pgfpathlineto{\pgfqpoint{2.828599in}{1.533431in}}%
\pgfpathlineto{\pgfqpoint{2.828637in}{1.543127in}}%
\pgfpathlineto{\pgfqpoint{2.835404in}{1.542568in}}%
\pgfpathlineto{\pgfqpoint{2.846989in}{1.546225in}}%
\pgfpathlineto{\pgfqpoint{2.850891in}{1.542069in}}%
\pgfpathlineto{\pgfqpoint{2.860753in}{1.546240in}}%
\pgfpathlineto{\pgfqpoint{2.867151in}{1.544262in}}%
\pgfpathclose%
\pgfusepath{fill}%
\end{pgfscope}%
\begin{pgfscope}%
\pgfpathrectangle{\pgfqpoint{0.100000in}{0.100000in}}{\pgfqpoint{3.007045in}{1.925000in}}%
\pgfusepath{clip}%
\pgfsetbuttcap%
\pgfsetmiterjoin%
\definecolor{currentfill}{rgb}{0.529412,0.742637,0.862899}%
\pgfsetfillcolor{currentfill}%
\pgfsetlinewidth{0.000000pt}%
\definecolor{currentstroke}{rgb}{0.000000,0.000000,0.000000}%
\pgfsetstrokecolor{currentstroke}%
\pgfsetstrokeopacity{0.000000}%
\pgfsetdash{}{0pt}%
\pgfpathmoveto{\pgfqpoint{1.147775in}{0.452101in}}%
\pgfpathlineto{\pgfqpoint{1.141912in}{0.462482in}}%
\pgfpathlineto{\pgfqpoint{1.134776in}{0.466628in}}%
\pgfpathlineto{\pgfqpoint{1.136483in}{0.471729in}}%
\pgfpathlineto{\pgfqpoint{1.136838in}{0.481546in}}%
\pgfpathlineto{\pgfqpoint{1.140040in}{0.487692in}}%
\pgfpathlineto{\pgfqpoint{1.147213in}{0.482417in}}%
\pgfpathlineto{\pgfqpoint{1.157196in}{0.484131in}}%
\pgfpathlineto{\pgfqpoint{1.159892in}{0.482747in}}%
\pgfpathlineto{\pgfqpoint{1.158835in}{0.470012in}}%
\pgfpathlineto{\pgfqpoint{1.155158in}{0.464696in}}%
\pgfpathlineto{\pgfqpoint{1.158256in}{0.454687in}}%
\pgfpathlineto{\pgfqpoint{1.157522in}{0.448171in}}%
\pgfpathlineto{\pgfqpoint{1.147775in}{0.452101in}}%
\pgfpathclose%
\pgfusepath{fill}%
\end{pgfscope}%
\begin{pgfscope}%
\pgfpathrectangle{\pgfqpoint{0.100000in}{0.100000in}}{\pgfqpoint{3.007045in}{1.925000in}}%
\pgfusepath{clip}%
\pgfsetbuttcap%
\pgfsetmiterjoin%
\definecolor{currentfill}{rgb}{0.429020,0.687520,0.841246}%
\pgfsetfillcolor{currentfill}%
\pgfsetlinewidth{0.000000pt}%
\definecolor{currentstroke}{rgb}{0.000000,0.000000,0.000000}%
\pgfsetstrokecolor{currentstroke}%
\pgfsetstrokeopacity{0.000000}%
\pgfsetdash{}{0pt}%
\pgfpathmoveto{\pgfqpoint{0.495869in}{0.921526in}}%
\pgfpathlineto{\pgfqpoint{0.488883in}{0.925335in}}%
\pgfpathlineto{\pgfqpoint{0.489033in}{0.933786in}}%
\pgfpathlineto{\pgfqpoint{0.495364in}{0.929362in}}%
\pgfpathlineto{\pgfqpoint{0.495869in}{0.921526in}}%
\pgfpathclose%
\pgfusepath{fill}%
\end{pgfscope}%
\begin{pgfscope}%
\pgfpathrectangle{\pgfqpoint{0.100000in}{0.100000in}}{\pgfqpoint{3.007045in}{1.925000in}}%
\pgfusepath{clip}%
\pgfsetbuttcap%
\pgfsetmiterjoin%
\definecolor{currentfill}{rgb}{0.429020,0.687520,0.841246}%
\pgfsetfillcolor{currentfill}%
\pgfsetlinewidth{0.000000pt}%
\definecolor{currentstroke}{rgb}{0.000000,0.000000,0.000000}%
\pgfsetstrokecolor{currentstroke}%
\pgfsetstrokeopacity{0.000000}%
\pgfsetdash{}{0pt}%
\pgfpathmoveto{\pgfqpoint{0.536458in}{0.917204in}}%
\pgfpathlineto{\pgfqpoint{0.528897in}{0.929641in}}%
\pgfpathlineto{\pgfqpoint{0.522315in}{0.935716in}}%
\pgfpathlineto{\pgfqpoint{0.515735in}{0.945831in}}%
\pgfpathlineto{\pgfqpoint{0.511433in}{0.949293in}}%
\pgfpathlineto{\pgfqpoint{0.504379in}{0.947346in}}%
\pgfpathlineto{\pgfqpoint{0.498832in}{0.951061in}}%
\pgfpathlineto{\pgfqpoint{0.501532in}{0.957039in}}%
\pgfpathlineto{\pgfqpoint{0.499515in}{0.967768in}}%
\pgfpathlineto{\pgfqpoint{0.496988in}{0.971830in}}%
\pgfpathlineto{\pgfqpoint{0.486075in}{0.974392in}}%
\pgfpathlineto{\pgfqpoint{0.481940in}{0.973863in}}%
\pgfpathlineto{\pgfqpoint{0.463162in}{0.988331in}}%
\pgfpathlineto{\pgfqpoint{0.461865in}{0.997033in}}%
\pgfpathlineto{\pgfqpoint{0.453370in}{1.006642in}}%
\pgfpathlineto{\pgfqpoint{0.456601in}{1.011141in}}%
\pgfpathlineto{\pgfqpoint{0.464444in}{1.039777in}}%
\pgfpathlineto{\pgfqpoint{0.472706in}{1.036029in}}%
\pgfpathlineto{\pgfqpoint{0.473306in}{1.031380in}}%
\pgfpathlineto{\pgfqpoint{0.515049in}{1.020412in}}%
\pgfpathlineto{\pgfqpoint{0.557933in}{1.009610in}}%
\pgfpathlineto{\pgfqpoint{0.573744in}{1.071996in}}%
\pgfpathlineto{\pgfqpoint{0.593532in}{1.066855in}}%
\pgfpathlineto{\pgfqpoint{0.654898in}{1.051840in}}%
\pgfpathlineto{\pgfqpoint{0.672324in}{1.047770in}}%
\pgfpathlineto{\pgfqpoint{0.677176in}{1.047687in}}%
\pgfpathlineto{\pgfqpoint{0.718691in}{0.983740in}}%
\pgfpathlineto{\pgfqpoint{0.716895in}{0.976496in}}%
\pgfpathlineto{\pgfqpoint{0.719619in}{0.967635in}}%
\pgfpathlineto{\pgfqpoint{0.723182in}{0.963163in}}%
\pgfpathlineto{\pgfqpoint{0.723368in}{0.955271in}}%
\pgfpathlineto{\pgfqpoint{0.726621in}{0.944753in}}%
\pgfpathlineto{\pgfqpoint{0.735192in}{0.932928in}}%
\pgfpathlineto{\pgfqpoint{0.740331in}{0.928918in}}%
\pgfpathlineto{\pgfqpoint{0.762651in}{0.922152in}}%
\pgfpathlineto{\pgfqpoint{0.767196in}{0.926377in}}%
\pgfpathlineto{\pgfqpoint{0.778392in}{0.924882in}}%
\pgfpathlineto{\pgfqpoint{0.774157in}{0.904247in}}%
\pgfpathlineto{\pgfqpoint{0.765905in}{0.863949in}}%
\pgfpathlineto{\pgfqpoint{0.732137in}{0.870987in}}%
\pgfpathlineto{\pgfqpoint{0.733331in}{0.876625in}}%
\pgfpathlineto{\pgfqpoint{0.716592in}{0.880125in}}%
\pgfpathlineto{\pgfqpoint{0.710499in}{0.852063in}}%
\pgfpathlineto{\pgfqpoint{0.689447in}{0.856914in}}%
\pgfpathlineto{\pgfqpoint{0.687845in}{0.862115in}}%
\pgfpathlineto{\pgfqpoint{0.691629in}{0.872688in}}%
\pgfpathlineto{\pgfqpoint{0.691055in}{0.881658in}}%
\pgfpathlineto{\pgfqpoint{0.673193in}{0.887262in}}%
\pgfpathlineto{\pgfqpoint{0.618157in}{0.899792in}}%
\pgfpathlineto{\pgfqpoint{0.567402in}{0.912130in}}%
\pgfpathlineto{\pgfqpoint{0.556193in}{0.915257in}}%
\pgfpathlineto{\pgfqpoint{0.550214in}{0.920766in}}%
\pgfpathlineto{\pgfqpoint{0.543010in}{0.923557in}}%
\pgfpathlineto{\pgfqpoint{0.538502in}{0.921211in}}%
\pgfpathlineto{\pgfqpoint{0.536458in}{0.917204in}}%
\pgfpathclose%
\pgfusepath{fill}%
\end{pgfscope}%
\begin{pgfscope}%
\pgfpathrectangle{\pgfqpoint{0.100000in}{0.100000in}}{\pgfqpoint{3.007045in}{1.925000in}}%
\pgfusepath{clip}%
\pgfsetbuttcap%
\pgfsetmiterjoin%
\definecolor{currentfill}{rgb}{0.510588,0.732303,0.858839}%
\pgfsetfillcolor{currentfill}%
\pgfsetlinewidth{0.000000pt}%
\definecolor{currentstroke}{rgb}{0.000000,0.000000,0.000000}%
\pgfsetstrokecolor{currentstroke}%
\pgfsetstrokeopacity{0.000000}%
\pgfsetdash{}{0pt}%
\pgfpathmoveto{\pgfqpoint{2.241536in}{0.660624in}}%
\pgfpathlineto{\pgfqpoint{2.252934in}{0.661782in}}%
\pgfpathlineto{\pgfqpoint{2.252110in}{0.670372in}}%
\pgfpathlineto{\pgfqpoint{2.254972in}{0.670665in}}%
\pgfpathlineto{\pgfqpoint{2.253069in}{0.691091in}}%
\pgfpathlineto{\pgfqpoint{2.254880in}{0.697001in}}%
\pgfpathlineto{\pgfqpoint{2.260621in}{0.697611in}}%
\pgfpathlineto{\pgfqpoint{2.261248in}{0.691899in}}%
\pgfpathlineto{\pgfqpoint{2.278425in}{0.693743in}}%
\pgfpathlineto{\pgfqpoint{2.277794in}{0.699248in}}%
\pgfpathlineto{\pgfqpoint{2.283592in}{0.699581in}}%
\pgfpathlineto{\pgfqpoint{2.284645in}{0.694345in}}%
\pgfpathlineto{\pgfqpoint{2.289927in}{0.694879in}}%
\pgfpathlineto{\pgfqpoint{2.290454in}{0.689136in}}%
\pgfpathlineto{\pgfqpoint{2.297933in}{0.689896in}}%
\pgfpathlineto{\pgfqpoint{2.298312in}{0.683587in}}%
\pgfpathlineto{\pgfqpoint{2.294624in}{0.672055in}}%
\pgfpathlineto{\pgfqpoint{2.272250in}{0.669798in}}%
\pgfpathlineto{\pgfqpoint{2.270062in}{0.663726in}}%
\pgfpathlineto{\pgfqpoint{2.273936in}{0.628314in}}%
\pgfpathlineto{\pgfqpoint{2.245494in}{0.625658in}}%
\pgfpathlineto{\pgfqpoint{2.244783in}{0.625988in}}%
\pgfpathlineto{\pgfqpoint{2.241536in}{0.660624in}}%
\pgfpathclose%
\pgfusepath{fill}%
\end{pgfscope}%
\begin{pgfscope}%
\pgfpathrectangle{\pgfqpoint{0.100000in}{0.100000in}}{\pgfqpoint{3.007045in}{1.925000in}}%
\pgfusepath{clip}%
\pgfsetbuttcap%
\pgfsetmiterjoin%
\definecolor{currentfill}{rgb}{0.235986,0.549712,0.764706}%
\pgfsetfillcolor{currentfill}%
\pgfsetlinewidth{0.000000pt}%
\definecolor{currentstroke}{rgb}{0.000000,0.000000,0.000000}%
\pgfsetstrokecolor{currentstroke}%
\pgfsetstrokeopacity{0.000000}%
\pgfsetdash{}{0pt}%
\pgfpathmoveto{\pgfqpoint{2.033514in}{0.905154in}}%
\pgfpathlineto{\pgfqpoint{2.039932in}{0.907365in}}%
\pgfpathlineto{\pgfqpoint{2.038540in}{0.912838in}}%
\pgfpathlineto{\pgfqpoint{2.045514in}{0.912506in}}%
\pgfpathlineto{\pgfqpoint{2.051337in}{0.919577in}}%
\pgfpathlineto{\pgfqpoint{2.061506in}{0.919581in}}%
\pgfpathlineto{\pgfqpoint{2.068214in}{0.916226in}}%
\pgfpathlineto{\pgfqpoint{2.069407in}{0.904987in}}%
\pgfpathlineto{\pgfqpoint{2.085014in}{0.905571in}}%
\pgfpathlineto{\pgfqpoint{2.085939in}{0.878995in}}%
\pgfpathlineto{\pgfqpoint{2.077608in}{0.878428in}}%
\pgfpathlineto{\pgfqpoint{2.078219in}{0.868579in}}%
\pgfpathlineto{\pgfqpoint{2.081112in}{0.868758in}}%
\pgfpathlineto{\pgfqpoint{2.082243in}{0.851523in}}%
\pgfpathlineto{\pgfqpoint{2.085442in}{0.845945in}}%
\pgfpathlineto{\pgfqpoint{2.072679in}{0.849028in}}%
\pgfpathlineto{\pgfqpoint{2.033456in}{0.846814in}}%
\pgfpathlineto{\pgfqpoint{2.027867in}{0.843557in}}%
\pgfpathlineto{\pgfqpoint{2.028291in}{0.837950in}}%
\pgfpathlineto{\pgfqpoint{2.022978in}{0.837629in}}%
\pgfpathlineto{\pgfqpoint{2.013407in}{0.843839in}}%
\pgfpathlineto{\pgfqpoint{2.012143in}{0.849788in}}%
\pgfpathlineto{\pgfqpoint{2.018227in}{0.853487in}}%
\pgfpathlineto{\pgfqpoint{2.020371in}{0.864609in}}%
\pgfpathlineto{\pgfqpoint{2.029442in}{0.869945in}}%
\pgfpathlineto{\pgfqpoint{2.026701in}{0.874367in}}%
\pgfpathlineto{\pgfqpoint{2.030962in}{0.877859in}}%
\pgfpathlineto{\pgfqpoint{2.032840in}{0.884060in}}%
\pgfpathlineto{\pgfqpoint{2.036782in}{0.884180in}}%
\pgfpathlineto{\pgfqpoint{2.037839in}{0.890956in}}%
\pgfpathlineto{\pgfqpoint{2.034157in}{0.893004in}}%
\pgfpathlineto{\pgfqpoint{2.037206in}{0.901885in}}%
\pgfpathlineto{\pgfqpoint{2.033514in}{0.905154in}}%
\pgfpathclose%
\pgfusepath{fill}%
\end{pgfscope}%
\begin{pgfscope}%
\pgfpathrectangle{\pgfqpoint{0.100000in}{0.100000in}}{\pgfqpoint{3.007045in}{1.925000in}}%
\pgfusepath{clip}%
\pgfsetbuttcap%
\pgfsetmiterjoin%
\definecolor{currentfill}{rgb}{0.391865,0.663406,0.828389}%
\pgfsetfillcolor{currentfill}%
\pgfsetlinewidth{0.000000pt}%
\definecolor{currentstroke}{rgb}{0.000000,0.000000,0.000000}%
\pgfsetstrokecolor{currentstroke}%
\pgfsetstrokeopacity{0.000000}%
\pgfsetdash{}{0pt}%
\pgfpathmoveto{\pgfqpoint{2.372484in}{1.056245in}}%
\pgfpathlineto{\pgfqpoint{2.363870in}{1.060375in}}%
\pgfpathlineto{\pgfqpoint{2.358324in}{1.060460in}}%
\pgfpathlineto{\pgfqpoint{2.352145in}{1.068229in}}%
\pgfpathlineto{\pgfqpoint{2.350488in}{1.072384in}}%
\pgfpathlineto{\pgfqpoint{2.353799in}{1.080602in}}%
\pgfpathlineto{\pgfqpoint{2.358589in}{1.084899in}}%
\pgfpathlineto{\pgfqpoint{2.362396in}{1.092621in}}%
\pgfpathlineto{\pgfqpoint{2.367300in}{1.093259in}}%
\pgfpathlineto{\pgfqpoint{2.371064in}{1.101491in}}%
\pgfpathlineto{\pgfqpoint{2.370833in}{1.107548in}}%
\pgfpathlineto{\pgfqpoint{2.379152in}{1.110286in}}%
\pgfpathlineto{\pgfqpoint{2.383396in}{1.107611in}}%
\pgfpathlineto{\pgfqpoint{2.392249in}{1.108544in}}%
\pgfpathlineto{\pgfqpoint{2.396141in}{1.102157in}}%
\pgfpathlineto{\pgfqpoint{2.389841in}{1.100458in}}%
\pgfpathlineto{\pgfqpoint{2.381551in}{1.091127in}}%
\pgfpathlineto{\pgfqpoint{2.382872in}{1.085587in}}%
\pgfpathlineto{\pgfqpoint{2.391489in}{1.082969in}}%
\pgfpathlineto{\pgfqpoint{2.400437in}{1.076137in}}%
\pgfpathlineto{\pgfqpoint{2.392967in}{1.076431in}}%
\pgfpathlineto{\pgfqpoint{2.392156in}{1.068534in}}%
\pgfpathlineto{\pgfqpoint{2.384647in}{1.067115in}}%
\pgfpathlineto{\pgfqpoint{2.378157in}{1.062296in}}%
\pgfpathlineto{\pgfqpoint{2.375741in}{1.064118in}}%
\pgfpathlineto{\pgfqpoint{2.370868in}{1.060399in}}%
\pgfpathlineto{\pgfqpoint{2.372484in}{1.056245in}}%
\pgfpathclose%
\pgfusepath{fill}%
\end{pgfscope}%
\begin{pgfscope}%
\pgfpathrectangle{\pgfqpoint{0.100000in}{0.100000in}}{\pgfqpoint{3.007045in}{1.925000in}}%
\pgfusepath{clip}%
\pgfsetbuttcap%
\pgfsetmiterjoin%
\definecolor{currentfill}{rgb}{0.793449,0.870142,0.942914}%
\pgfsetfillcolor{currentfill}%
\pgfsetlinewidth{0.000000pt}%
\definecolor{currentstroke}{rgb}{0.000000,0.000000,0.000000}%
\pgfsetstrokecolor{currentstroke}%
\pgfsetstrokeopacity{0.000000}%
\pgfsetdash{}{0pt}%
\pgfpathmoveto{\pgfqpoint{2.623105in}{0.353228in}}%
\pgfpathlineto{\pgfqpoint{2.628250in}{0.323017in}}%
\pgfpathlineto{\pgfqpoint{2.599950in}{0.318474in}}%
\pgfpathlineto{\pgfqpoint{2.594969in}{0.321852in}}%
\pgfpathlineto{\pgfqpoint{2.580427in}{0.321862in}}%
\pgfpathlineto{\pgfqpoint{2.573972in}{0.324804in}}%
\pgfpathlineto{\pgfqpoint{2.569745in}{0.333328in}}%
\pgfpathlineto{\pgfqpoint{2.565657in}{0.347157in}}%
\pgfpathlineto{\pgfqpoint{2.562738in}{0.352174in}}%
\pgfpathlineto{\pgfqpoint{2.555247in}{0.358065in}}%
\pgfpathlineto{\pgfqpoint{2.547708in}{0.358122in}}%
\pgfpathlineto{\pgfqpoint{2.543685in}{0.367836in}}%
\pgfpathlineto{\pgfqpoint{2.546700in}{0.369561in}}%
\pgfpathlineto{\pgfqpoint{2.547508in}{0.376728in}}%
\pgfpathlineto{\pgfqpoint{2.577000in}{0.380965in}}%
\pgfpathlineto{\pgfqpoint{2.579681in}{0.363910in}}%
\pgfpathlineto{\pgfqpoint{2.597057in}{0.366782in}}%
\pgfpathlineto{\pgfqpoint{2.599929in}{0.349210in}}%
\pgfpathlineto{\pgfqpoint{2.623105in}{0.353228in}}%
\pgfpathclose%
\pgfusepath{fill}%
\end{pgfscope}%
\begin{pgfscope}%
\pgfpathrectangle{\pgfqpoint{0.100000in}{0.100000in}}{\pgfqpoint{3.007045in}{1.925000in}}%
\pgfusepath{clip}%
\pgfsetbuttcap%
\pgfsetmiterjoin%
\definecolor{currentfill}{rgb}{0.485490,0.718524,0.853426}%
\pgfsetfillcolor{currentfill}%
\pgfsetlinewidth{0.000000pt}%
\definecolor{currentstroke}{rgb}{0.000000,0.000000,0.000000}%
\pgfsetstrokecolor{currentstroke}%
\pgfsetstrokeopacity{0.000000}%
\pgfsetdash{}{0pt}%
\pgfpathmoveto{\pgfqpoint{0.440258in}{1.107424in}}%
\pgfpathlineto{\pgfqpoint{0.406868in}{1.117231in}}%
\pgfpathlineto{\pgfqpoint{0.382061in}{1.124859in}}%
\pgfpathlineto{\pgfqpoint{0.377822in}{1.132382in}}%
\pgfpathlineto{\pgfqpoint{0.377753in}{1.140606in}}%
\pgfpathlineto{\pgfqpoint{0.375122in}{1.142276in}}%
\pgfpathlineto{\pgfqpoint{0.370888in}{1.155493in}}%
\pgfpathlineto{\pgfqpoint{0.366684in}{1.159674in}}%
\pgfpathlineto{\pgfqpoint{0.363995in}{1.165891in}}%
\pgfpathlineto{\pgfqpoint{0.365227in}{1.184145in}}%
\pgfpathlineto{\pgfqpoint{0.371166in}{1.184004in}}%
\pgfpathlineto{\pgfqpoint{0.374775in}{1.187299in}}%
\pgfpathlineto{\pgfqpoint{0.378870in}{1.195364in}}%
\pgfpathlineto{\pgfqpoint{0.377503in}{1.204628in}}%
\pgfpathlineto{\pgfqpoint{0.365785in}{1.209702in}}%
\pgfpathlineto{\pgfqpoint{0.361242in}{1.215875in}}%
\pgfpathlineto{\pgfqpoint{0.359301in}{1.222111in}}%
\pgfpathlineto{\pgfqpoint{0.359639in}{1.227533in}}%
\pgfpathlineto{\pgfqpoint{0.368463in}{1.226772in}}%
\pgfpathlineto{\pgfqpoint{0.368846in}{1.236860in}}%
\pgfpathlineto{\pgfqpoint{0.370754in}{1.240933in}}%
\pgfpathlineto{\pgfqpoint{0.377116in}{1.242201in}}%
\pgfpathlineto{\pgfqpoint{0.384479in}{1.238341in}}%
\pgfpathlineto{\pgfqpoint{0.388058in}{1.239323in}}%
\pgfpathlineto{\pgfqpoint{0.407661in}{1.233230in}}%
\pgfpathlineto{\pgfqpoint{0.408920in}{1.226474in}}%
\pgfpathlineto{\pgfqpoint{0.404848in}{1.220444in}}%
\pgfpathlineto{\pgfqpoint{0.404869in}{1.211825in}}%
\pgfpathlineto{\pgfqpoint{0.413366in}{1.209180in}}%
\pgfpathlineto{\pgfqpoint{0.413081in}{1.203149in}}%
\pgfpathlineto{\pgfqpoint{0.409170in}{1.194540in}}%
\pgfpathlineto{\pgfqpoint{0.412185in}{1.187494in}}%
\pgfpathlineto{\pgfqpoint{0.421673in}{1.178083in}}%
\pgfpathlineto{\pgfqpoint{0.433028in}{1.157462in}}%
\pgfpathlineto{\pgfqpoint{0.430464in}{1.147237in}}%
\pgfpathlineto{\pgfqpoint{0.426809in}{1.147328in}}%
\pgfpathlineto{\pgfqpoint{0.423734in}{1.137658in}}%
\pgfpathlineto{\pgfqpoint{0.428309in}{1.126977in}}%
\pgfpathlineto{\pgfqpoint{0.432260in}{1.122040in}}%
\pgfpathlineto{\pgfqpoint{0.435906in}{1.120684in}}%
\pgfpathlineto{\pgfqpoint{0.440258in}{1.107424in}}%
\pgfpathclose%
\pgfusepath{fill}%
\end{pgfscope}%
\begin{pgfscope}%
\pgfpathrectangle{\pgfqpoint{0.100000in}{0.100000in}}{\pgfqpoint{3.007045in}{1.925000in}}%
\pgfusepath{clip}%
\pgfsetbuttcap%
\pgfsetmiterjoin%
\definecolor{currentfill}{rgb}{0.412042,0.677186,0.836263}%
\pgfsetfillcolor{currentfill}%
\pgfsetlinewidth{0.000000pt}%
\definecolor{currentstroke}{rgb}{0.000000,0.000000,0.000000}%
\pgfsetstrokecolor{currentstroke}%
\pgfsetstrokeopacity{0.000000}%
\pgfsetdash{}{0pt}%
\pgfpathmoveto{\pgfqpoint{0.966052in}{0.192399in}}%
\pgfpathlineto{\pgfqpoint{0.964721in}{0.193135in}}%
\pgfpathlineto{\pgfqpoint{0.963493in}{0.195565in}}%
\pgfpathlineto{\pgfqpoint{0.958655in}{0.195219in}}%
\pgfpathlineto{\pgfqpoint{0.957247in}{0.193413in}}%
\pgfpathlineto{\pgfqpoint{0.955692in}{0.192812in}}%
\pgfpathlineto{\pgfqpoint{0.954245in}{0.194809in}}%
\pgfpathlineto{\pgfqpoint{0.951030in}{0.194631in}}%
\pgfpathlineto{\pgfqpoint{0.952673in}{0.196998in}}%
\pgfpathlineto{\pgfqpoint{0.954685in}{0.197089in}}%
\pgfpathlineto{\pgfqpoint{0.954041in}{0.199019in}}%
\pgfpathlineto{\pgfqpoint{0.954053in}{0.201532in}}%
\pgfpathlineto{\pgfqpoint{0.954782in}{0.203274in}}%
\pgfpathlineto{\pgfqpoint{0.953594in}{0.205042in}}%
\pgfpathlineto{\pgfqpoint{0.954325in}{0.206175in}}%
\pgfpathlineto{\pgfqpoint{0.956732in}{0.206941in}}%
\pgfpathlineto{\pgfqpoint{0.957233in}{0.209806in}}%
\pgfpathlineto{\pgfqpoint{0.954751in}{0.209244in}}%
\pgfpathlineto{\pgfqpoint{0.954872in}{0.213244in}}%
\pgfpathlineto{\pgfqpoint{0.953569in}{0.215767in}}%
\pgfpathlineto{\pgfqpoint{0.954140in}{0.218651in}}%
\pgfpathlineto{\pgfqpoint{0.953788in}{0.220128in}}%
\pgfpathlineto{\pgfqpoint{0.952256in}{0.220622in}}%
\pgfpathlineto{\pgfqpoint{0.949173in}{0.223004in}}%
\pgfpathlineto{\pgfqpoint{0.948744in}{0.225034in}}%
\pgfpathlineto{\pgfqpoint{0.949609in}{0.226010in}}%
\pgfpathlineto{\pgfqpoint{0.948660in}{0.227534in}}%
\pgfpathlineto{\pgfqpoint{0.949321in}{0.229307in}}%
\pgfpathlineto{\pgfqpoint{0.949075in}{0.234015in}}%
\pgfpathlineto{\pgfqpoint{0.949908in}{0.234459in}}%
\pgfpathlineto{\pgfqpoint{0.950927in}{0.237714in}}%
\pgfpathlineto{\pgfqpoint{0.949190in}{0.237528in}}%
\pgfpathlineto{\pgfqpoint{0.949382in}{0.239911in}}%
\pgfpathlineto{\pgfqpoint{0.951032in}{0.247085in}}%
\pgfpathlineto{\pgfqpoint{0.951672in}{0.251760in}}%
\pgfpathlineto{\pgfqpoint{0.949767in}{0.251172in}}%
\pgfpathlineto{\pgfqpoint{0.950024in}{0.249097in}}%
\pgfpathlineto{\pgfqpoint{0.948795in}{0.248793in}}%
\pgfpathlineto{\pgfqpoint{0.948847in}{0.246567in}}%
\pgfpathlineto{\pgfqpoint{0.947286in}{0.241526in}}%
\pgfpathlineto{\pgfqpoint{0.946946in}{0.236340in}}%
\pgfpathlineto{\pgfqpoint{0.945783in}{0.232618in}}%
\pgfpathlineto{\pgfqpoint{0.944897in}{0.227958in}}%
\pgfpathlineto{\pgfqpoint{0.942843in}{0.225828in}}%
\pgfpathlineto{\pgfqpoint{0.941004in}{0.228956in}}%
\pgfpathlineto{\pgfqpoint{0.941264in}{0.231409in}}%
\pgfpathlineto{\pgfqpoint{0.939615in}{0.232446in}}%
\pgfpathlineto{\pgfqpoint{0.935598in}{0.233694in}}%
\pgfpathlineto{\pgfqpoint{0.937893in}{0.236489in}}%
\pgfpathlineto{\pgfqpoint{0.938389in}{0.240874in}}%
\pgfpathlineto{\pgfqpoint{0.938282in}{0.243236in}}%
\pgfpathlineto{\pgfqpoint{0.935953in}{0.243285in}}%
\pgfpathlineto{\pgfqpoint{0.934974in}{0.245815in}}%
\pgfpathlineto{\pgfqpoint{0.933942in}{0.246488in}}%
\pgfpathlineto{\pgfqpoint{0.934379in}{0.249543in}}%
\pgfpathlineto{\pgfqpoint{0.931764in}{0.250440in}}%
\pgfpathlineto{\pgfqpoint{0.930132in}{0.252504in}}%
\pgfpathlineto{\pgfqpoint{0.929218in}{0.251104in}}%
\pgfpathlineto{\pgfqpoint{0.931463in}{0.248759in}}%
\pgfpathlineto{\pgfqpoint{0.931593in}{0.245629in}}%
\pgfpathlineto{\pgfqpoint{0.933934in}{0.242455in}}%
\pgfpathlineto{\pgfqpoint{0.931815in}{0.241739in}}%
\pgfpathlineto{\pgfqpoint{0.932820in}{0.240335in}}%
\pgfpathlineto{\pgfqpoint{0.934512in}{0.240731in}}%
\pgfpathlineto{\pgfqpoint{0.934225in}{0.234374in}}%
\pgfpathlineto{\pgfqpoint{0.930700in}{0.234187in}}%
\pgfpathlineto{\pgfqpoint{0.931316in}{0.236252in}}%
\pgfpathlineto{\pgfqpoint{0.928812in}{0.235626in}}%
\pgfpathlineto{\pgfqpoint{0.925421in}{0.233892in}}%
\pgfpathlineto{\pgfqpoint{0.925428in}{0.237038in}}%
\pgfpathlineto{\pgfqpoint{0.924460in}{0.239424in}}%
\pgfpathlineto{\pgfqpoint{0.922320in}{0.240051in}}%
\pgfpathlineto{\pgfqpoint{0.919294in}{0.247052in}}%
\pgfpathlineto{\pgfqpoint{0.919630in}{0.248516in}}%
\pgfpathlineto{\pgfqpoint{0.917767in}{0.252248in}}%
\pgfpathlineto{\pgfqpoint{0.917999in}{0.256789in}}%
\pgfpathlineto{\pgfqpoint{0.916349in}{0.261630in}}%
\pgfpathlineto{\pgfqpoint{0.913473in}{0.265334in}}%
\pgfpathlineto{\pgfqpoint{0.912193in}{0.268097in}}%
\pgfpathlineto{\pgfqpoint{0.908339in}{0.273326in}}%
\pgfpathlineto{\pgfqpoint{0.904581in}{0.280446in}}%
\pgfpathlineto{\pgfqpoint{0.907013in}{0.280057in}}%
\pgfpathlineto{\pgfqpoint{0.910459in}{0.281718in}}%
\pgfpathlineto{\pgfqpoint{0.910823in}{0.286283in}}%
\pgfpathlineto{\pgfqpoint{0.911674in}{0.287962in}}%
\pgfpathlineto{\pgfqpoint{0.907295in}{0.286400in}}%
\pgfpathlineto{\pgfqpoint{0.905838in}{0.288835in}}%
\pgfpathlineto{\pgfqpoint{0.902783in}{0.288080in}}%
\pgfpathlineto{\pgfqpoint{0.901946in}{0.286374in}}%
\pgfpathlineto{\pgfqpoint{0.897692in}{0.289155in}}%
\pgfpathlineto{\pgfqpoint{0.895195in}{0.292046in}}%
\pgfpathlineto{\pgfqpoint{0.901489in}{0.303399in}}%
\pgfpathlineto{\pgfqpoint{0.905283in}{0.298997in}}%
\pgfpathlineto{\pgfqpoint{0.907695in}{0.299493in}}%
\pgfpathlineto{\pgfqpoint{0.910640in}{0.295087in}}%
\pgfpathlineto{\pgfqpoint{0.915744in}{0.296619in}}%
\pgfpathlineto{\pgfqpoint{0.920918in}{0.294241in}}%
\pgfpathlineto{\pgfqpoint{0.921920in}{0.292840in}}%
\pgfpathlineto{\pgfqpoint{0.917937in}{0.288586in}}%
\pgfpathlineto{\pgfqpoint{0.918519in}{0.285751in}}%
\pgfpathlineto{\pgfqpoint{0.921178in}{0.282019in}}%
\pgfpathlineto{\pgfqpoint{0.920601in}{0.278454in}}%
\pgfpathlineto{\pgfqpoint{0.925898in}{0.261651in}}%
\pgfpathlineto{\pgfqpoint{0.924680in}{0.256502in}}%
\pgfpathlineto{\pgfqpoint{0.923189in}{0.253738in}}%
\pgfpathlineto{\pgfqpoint{0.924096in}{0.253355in}}%
\pgfpathlineto{\pgfqpoint{0.927167in}{0.254345in}}%
\pgfpathlineto{\pgfqpoint{0.932410in}{0.254987in}}%
\pgfpathlineto{\pgfqpoint{0.936447in}{0.254101in}}%
\pgfpathlineto{\pgfqpoint{0.938989in}{0.256067in}}%
\pgfpathlineto{\pgfqpoint{0.940309in}{0.258917in}}%
\pgfpathlineto{\pgfqpoint{0.942720in}{0.259083in}}%
\pgfpathlineto{\pgfqpoint{0.946116in}{0.262219in}}%
\pgfpathlineto{\pgfqpoint{0.948733in}{0.261499in}}%
\pgfpathlineto{\pgfqpoint{0.954529in}{0.261913in}}%
\pgfpathlineto{\pgfqpoint{0.955655in}{0.261214in}}%
\pgfpathlineto{\pgfqpoint{0.957563in}{0.255964in}}%
\pgfpathlineto{\pgfqpoint{0.957437in}{0.253580in}}%
\pgfpathlineto{\pgfqpoint{0.955678in}{0.251817in}}%
\pgfpathlineto{\pgfqpoint{0.955644in}{0.247945in}}%
\pgfpathlineto{\pgfqpoint{0.958149in}{0.245787in}}%
\pgfpathlineto{\pgfqpoint{0.958180in}{0.243117in}}%
\pgfpathlineto{\pgfqpoint{0.959385in}{0.241184in}}%
\pgfpathlineto{\pgfqpoint{0.958932in}{0.238807in}}%
\pgfpathlineto{\pgfqpoint{0.959232in}{0.235056in}}%
\pgfpathlineto{\pgfqpoint{0.962530in}{0.230044in}}%
\pgfpathlineto{\pgfqpoint{0.962897in}{0.225216in}}%
\pgfpathlineto{\pgfqpoint{0.964318in}{0.222393in}}%
\pgfpathlineto{\pgfqpoint{0.963412in}{0.220631in}}%
\pgfpathlineto{\pgfqpoint{0.964300in}{0.213106in}}%
\pgfpathlineto{\pgfqpoint{0.964057in}{0.210474in}}%
\pgfpathlineto{\pgfqpoint{0.964942in}{0.205541in}}%
\pgfpathlineto{\pgfqpoint{0.965301in}{0.196768in}}%
\pgfpathlineto{\pgfqpoint{0.966052in}{0.192399in}}%
\pgfpathclose%
\pgfusepath{fill}%
\end{pgfscope}%
\begin{pgfscope}%
\pgfpathrectangle{\pgfqpoint{0.100000in}{0.100000in}}{\pgfqpoint{3.007045in}{1.925000in}}%
\pgfusepath{clip}%
\pgfsetbuttcap%
\pgfsetmiterjoin%
\definecolor{currentfill}{rgb}{0.412042,0.677186,0.836263}%
\pgfsetfillcolor{currentfill}%
\pgfsetlinewidth{0.000000pt}%
\definecolor{currentstroke}{rgb}{0.000000,0.000000,0.000000}%
\pgfsetstrokecolor{currentstroke}%
\pgfsetstrokeopacity{0.000000}%
\pgfsetdash{}{0pt}%
\pgfpathmoveto{\pgfqpoint{0.936626in}{0.197068in}}%
\pgfpathlineto{\pgfqpoint{0.938877in}{0.201656in}}%
\pgfpathlineto{\pgfqpoint{0.940558in}{0.202970in}}%
\pgfpathlineto{\pgfqpoint{0.942011in}{0.205039in}}%
\pgfpathlineto{\pgfqpoint{0.942374in}{0.211894in}}%
\pgfpathlineto{\pgfqpoint{0.944008in}{0.216194in}}%
\pgfpathlineto{\pgfqpoint{0.944696in}{0.219988in}}%
\pgfpathlineto{\pgfqpoint{0.945984in}{0.221447in}}%
\pgfpathlineto{\pgfqpoint{0.945425in}{0.223366in}}%
\pgfpathlineto{\pgfqpoint{0.945878in}{0.227745in}}%
\pgfpathlineto{\pgfqpoint{0.947540in}{0.225877in}}%
\pgfpathlineto{\pgfqpoint{0.947168in}{0.222342in}}%
\pgfpathlineto{\pgfqpoint{0.948688in}{0.221616in}}%
\pgfpathlineto{\pgfqpoint{0.952311in}{0.218681in}}%
\pgfpathlineto{\pgfqpoint{0.952653in}{0.214058in}}%
\pgfpathlineto{\pgfqpoint{0.952346in}{0.208455in}}%
\pgfpathlineto{\pgfqpoint{0.950987in}{0.209391in}}%
\pgfpathlineto{\pgfqpoint{0.951266in}{0.216399in}}%
\pgfpathlineto{\pgfqpoint{0.950054in}{0.216109in}}%
\pgfpathlineto{\pgfqpoint{0.948419in}{0.213095in}}%
\pgfpathlineto{\pgfqpoint{0.949973in}{0.212546in}}%
\pgfpathlineto{\pgfqpoint{0.949906in}{0.210661in}}%
\pgfpathlineto{\pgfqpoint{0.948880in}{0.207820in}}%
\pgfpathlineto{\pgfqpoint{0.950125in}{0.207070in}}%
\pgfpathlineto{\pgfqpoint{0.949683in}{0.203304in}}%
\pgfpathlineto{\pgfqpoint{0.947074in}{0.204157in}}%
\pgfpathlineto{\pgfqpoint{0.948879in}{0.200841in}}%
\pgfpathlineto{\pgfqpoint{0.947064in}{0.199809in}}%
\pgfpathlineto{\pgfqpoint{0.946264in}{0.200919in}}%
\pgfpathlineto{\pgfqpoint{0.944697in}{0.199414in}}%
\pgfpathlineto{\pgfqpoint{0.942282in}{0.199077in}}%
\pgfpathlineto{\pgfqpoint{0.939613in}{0.197250in}}%
\pgfpathlineto{\pgfqpoint{0.936626in}{0.197068in}}%
\pgfpathclose%
\pgfusepath{fill}%
\end{pgfscope}%
\begin{pgfscope}%
\pgfpathrectangle{\pgfqpoint{0.100000in}{0.100000in}}{\pgfqpoint{3.007045in}{1.925000in}}%
\pgfusepath{clip}%
\pgfsetbuttcap%
\pgfsetmiterjoin%
\definecolor{currentfill}{rgb}{0.412042,0.677186,0.836263}%
\pgfsetfillcolor{currentfill}%
\pgfsetlinewidth{0.000000pt}%
\definecolor{currentstroke}{rgb}{0.000000,0.000000,0.000000}%
\pgfsetstrokecolor{currentstroke}%
\pgfsetstrokeopacity{0.000000}%
\pgfsetdash{}{0pt}%
\pgfpathmoveto{\pgfqpoint{0.924230in}{0.224566in}}%
\pgfpathlineto{\pgfqpoint{0.924247in}{0.226972in}}%
\pgfpathlineto{\pgfqpoint{0.925304in}{0.230453in}}%
\pgfpathlineto{\pgfqpoint{0.928832in}{0.230375in}}%
\pgfpathlineto{\pgfqpoint{0.930194in}{0.229696in}}%
\pgfpathlineto{\pgfqpoint{0.931247in}{0.231217in}}%
\pgfpathlineto{\pgfqpoint{0.932446in}{0.229785in}}%
\pgfpathlineto{\pgfqpoint{0.935561in}{0.230884in}}%
\pgfpathlineto{\pgfqpoint{0.937650in}{0.227795in}}%
\pgfpathlineto{\pgfqpoint{0.938131in}{0.224744in}}%
\pgfpathlineto{\pgfqpoint{0.941809in}{0.222536in}}%
\pgfpathlineto{\pgfqpoint{0.943001in}{0.220916in}}%
\pgfpathlineto{\pgfqpoint{0.942398in}{0.218706in}}%
\pgfpathlineto{\pgfqpoint{0.941010in}{0.217751in}}%
\pgfpathlineto{\pgfqpoint{0.939636in}{0.215477in}}%
\pgfpathlineto{\pgfqpoint{0.936886in}{0.216547in}}%
\pgfpathlineto{\pgfqpoint{0.934778in}{0.220682in}}%
\pgfpathlineto{\pgfqpoint{0.933152in}{0.222753in}}%
\pgfpathlineto{\pgfqpoint{0.931898in}{0.225478in}}%
\pgfpathlineto{\pgfqpoint{0.928775in}{0.222468in}}%
\pgfpathlineto{\pgfqpoint{0.924230in}{0.224566in}}%
\pgfpathclose%
\pgfusepath{fill}%
\end{pgfscope}%
\begin{pgfscope}%
\pgfpathrectangle{\pgfqpoint{0.100000in}{0.100000in}}{\pgfqpoint{3.007045in}{1.925000in}}%
\pgfusepath{clip}%
\pgfsetbuttcap%
\pgfsetmiterjoin%
\definecolor{currentfill}{rgb}{0.535686,0.746082,0.864252}%
\pgfsetfillcolor{currentfill}%
\pgfsetlinewidth{0.000000pt}%
\definecolor{currentstroke}{rgb}{0.000000,0.000000,0.000000}%
\pgfsetstrokecolor{currentstroke}%
\pgfsetstrokeopacity{0.000000}%
\pgfsetdash{}{0pt}%
\pgfpathmoveto{\pgfqpoint{2.803517in}{1.405975in}}%
\pgfpathlineto{\pgfqpoint{2.803333in}{1.417011in}}%
\pgfpathlineto{\pgfqpoint{2.822125in}{1.424128in}}%
\pgfpathlineto{\pgfqpoint{2.824695in}{1.409031in}}%
\pgfpathlineto{\pgfqpoint{2.829259in}{1.404526in}}%
\pgfpathlineto{\pgfqpoint{2.819156in}{1.394508in}}%
\pgfpathlineto{\pgfqpoint{2.824324in}{1.387968in}}%
\pgfpathlineto{\pgfqpoint{2.826877in}{1.382803in}}%
\pgfpathlineto{\pgfqpoint{2.833070in}{1.385813in}}%
\pgfpathlineto{\pgfqpoint{2.846367in}{1.387920in}}%
\pgfpathlineto{\pgfqpoint{2.853875in}{1.393696in}}%
\pgfpathlineto{\pgfqpoint{2.867459in}{1.397196in}}%
\pgfpathlineto{\pgfqpoint{2.873958in}{1.399810in}}%
\pgfpathlineto{\pgfqpoint{2.883054in}{1.409613in}}%
\pgfpathlineto{\pgfqpoint{2.884754in}{1.413273in}}%
\pgfpathlineto{\pgfqpoint{2.890760in}{1.408191in}}%
\pgfpathlineto{\pgfqpoint{2.895504in}{1.410186in}}%
\pgfpathlineto{\pgfqpoint{2.899044in}{1.407001in}}%
\pgfpathlineto{\pgfqpoint{2.904883in}{1.413872in}}%
\pgfpathlineto{\pgfqpoint{2.909280in}{1.413649in}}%
\pgfpathlineto{\pgfqpoint{2.882440in}{1.391754in}}%
\pgfpathlineto{\pgfqpoint{2.865813in}{1.380465in}}%
\pgfpathlineto{\pgfqpoint{2.859459in}{1.375218in}}%
\pgfpathlineto{\pgfqpoint{2.849223in}{1.369251in}}%
\pgfpathlineto{\pgfqpoint{2.845081in}{1.369000in}}%
\pgfpathlineto{\pgfqpoint{2.835181in}{1.363297in}}%
\pgfpathlineto{\pgfqpoint{2.824782in}{1.361319in}}%
\pgfpathlineto{\pgfqpoint{2.818078in}{1.357032in}}%
\pgfpathlineto{\pgfqpoint{2.810932in}{1.358841in}}%
\pgfpathlineto{\pgfqpoint{2.808914in}{1.354582in}}%
\pgfpathlineto{\pgfqpoint{2.802639in}{1.349500in}}%
\pgfpathlineto{\pgfqpoint{2.803807in}{1.360108in}}%
\pgfpathlineto{\pgfqpoint{2.810062in}{1.361931in}}%
\pgfpathlineto{\pgfqpoint{2.812809in}{1.384940in}}%
\pgfpathlineto{\pgfqpoint{2.810869in}{1.395086in}}%
\pgfpathlineto{\pgfqpoint{2.805188in}{1.401167in}}%
\pgfpathlineto{\pgfqpoint{2.803517in}{1.405975in}}%
\pgfpathclose%
\pgfusepath{fill}%
\end{pgfscope}%
\begin{pgfscope}%
\pgfpathrectangle{\pgfqpoint{0.100000in}{0.100000in}}{\pgfqpoint{3.007045in}{1.925000in}}%
\pgfusepath{clip}%
\pgfsetbuttcap%
\pgfsetmiterjoin%
\definecolor{currentfill}{rgb}{0.429020,0.687520,0.841246}%
\pgfsetfillcolor{currentfill}%
\pgfsetlinewidth{0.000000pt}%
\definecolor{currentstroke}{rgb}{0.000000,0.000000,0.000000}%
\pgfsetstrokecolor{currentstroke}%
\pgfsetstrokeopacity{0.000000}%
\pgfsetdash{}{0pt}%
\pgfpathmoveto{\pgfqpoint{0.981446in}{1.653571in}}%
\pgfpathlineto{\pgfqpoint{0.972309in}{1.661350in}}%
\pgfpathlineto{\pgfqpoint{0.973053in}{1.678906in}}%
\pgfpathlineto{\pgfqpoint{0.969357in}{1.685011in}}%
\pgfpathlineto{\pgfqpoint{0.971079in}{1.692732in}}%
\pgfpathlineto{\pgfqpoint{0.973346in}{1.694790in}}%
\pgfpathlineto{\pgfqpoint{0.982127in}{1.696489in}}%
\pgfpathlineto{\pgfqpoint{0.985045in}{1.700348in}}%
\pgfpathlineto{\pgfqpoint{0.987484in}{1.710382in}}%
\pgfpathlineto{\pgfqpoint{0.983438in}{1.717518in}}%
\pgfpathlineto{\pgfqpoint{0.977851in}{1.718526in}}%
\pgfpathlineto{\pgfqpoint{0.979570in}{1.727393in}}%
\pgfpathlineto{\pgfqpoint{0.967944in}{1.729505in}}%
\pgfpathlineto{\pgfqpoint{0.972359in}{1.752059in}}%
\pgfpathlineto{\pgfqpoint{0.960752in}{1.754525in}}%
\pgfpathlineto{\pgfqpoint{0.964785in}{1.774681in}}%
\pgfpathlineto{\pgfqpoint{0.962444in}{1.782513in}}%
\pgfpathlineto{\pgfqpoint{0.963166in}{1.790823in}}%
\pgfpathlineto{\pgfqpoint{0.967284in}{1.791973in}}%
\pgfpathlineto{\pgfqpoint{0.969863in}{1.800673in}}%
\pgfpathlineto{\pgfqpoint{0.974188in}{1.804466in}}%
\pgfpathlineto{\pgfqpoint{0.975711in}{1.795488in}}%
\pgfpathlineto{\pgfqpoint{0.973893in}{1.786912in}}%
\pgfpathlineto{\pgfqpoint{0.976971in}{1.780376in}}%
\pgfpathlineto{\pgfqpoint{0.987897in}{1.779953in}}%
\pgfpathlineto{\pgfqpoint{0.994400in}{1.776622in}}%
\pgfpathlineto{\pgfqpoint{0.994442in}{1.773321in}}%
\pgfpathlineto{\pgfqpoint{0.999503in}{1.769057in}}%
\pgfpathlineto{\pgfqpoint{1.009841in}{1.767764in}}%
\pgfpathlineto{\pgfqpoint{1.005933in}{1.746597in}}%
\pgfpathlineto{\pgfqpoint{1.010013in}{1.742368in}}%
\pgfpathlineto{\pgfqpoint{1.016513in}{1.740285in}}%
\pgfpathlineto{\pgfqpoint{1.013851in}{1.726227in}}%
\pgfpathlineto{\pgfqpoint{1.019738in}{1.725093in}}%
\pgfpathlineto{\pgfqpoint{1.019449in}{1.720433in}}%
\pgfpathlineto{\pgfqpoint{1.025028in}{1.717016in}}%
\pgfpathlineto{\pgfqpoint{1.027516in}{1.705647in}}%
\pgfpathlineto{\pgfqpoint{1.030764in}{1.699834in}}%
\pgfpathlineto{\pgfqpoint{1.033247in}{1.689867in}}%
\pgfpathlineto{\pgfqpoint{1.037077in}{1.690684in}}%
\pgfpathlineto{\pgfqpoint{1.040593in}{1.686492in}}%
\pgfpathlineto{\pgfqpoint{1.036887in}{1.681674in}}%
\pgfpathlineto{\pgfqpoint{1.037941in}{1.673119in}}%
\pgfpathlineto{\pgfqpoint{1.025044in}{1.675119in}}%
\pgfpathlineto{\pgfqpoint{1.024131in}{1.670887in}}%
\pgfpathlineto{\pgfqpoint{1.019996in}{1.667445in}}%
\pgfpathlineto{\pgfqpoint{1.016570in}{1.660701in}}%
\pgfpathlineto{\pgfqpoint{1.010507in}{1.658210in}}%
\pgfpathlineto{\pgfqpoint{1.004464in}{1.652753in}}%
\pgfpathlineto{\pgfqpoint{0.997011in}{1.655494in}}%
\pgfpathlineto{\pgfqpoint{0.994682in}{1.658385in}}%
\pgfpathlineto{\pgfqpoint{0.988214in}{1.659319in}}%
\pgfpathlineto{\pgfqpoint{0.981446in}{1.653571in}}%
\pgfpathclose%
\pgfusepath{fill}%
\end{pgfscope}%
\begin{pgfscope}%
\pgfpathrectangle{\pgfqpoint{0.100000in}{0.100000in}}{\pgfqpoint{3.007045in}{1.925000in}}%
\pgfusepath{clip}%
\pgfsetbuttcap%
\pgfsetmiterjoin%
\definecolor{currentfill}{rgb}{0.371688,0.649627,0.820515}%
\pgfsetfillcolor{currentfill}%
\pgfsetlinewidth{0.000000pt}%
\definecolor{currentstroke}{rgb}{0.000000,0.000000,0.000000}%
\pgfsetstrokecolor{currentstroke}%
\pgfsetstrokeopacity{0.000000}%
\pgfsetdash{}{0pt}%
\pgfpathmoveto{\pgfqpoint{1.653431in}{1.223538in}}%
\pgfpathlineto{\pgfqpoint{1.630558in}{1.224197in}}%
\pgfpathlineto{\pgfqpoint{1.631190in}{1.247136in}}%
\pgfpathlineto{\pgfqpoint{1.632031in}{1.278421in}}%
\pgfpathlineto{\pgfqpoint{1.635249in}{1.282362in}}%
\pgfpathlineto{\pgfqpoint{1.643733in}{1.288581in}}%
\pgfpathlineto{\pgfqpoint{1.655235in}{1.292354in}}%
\pgfpathlineto{\pgfqpoint{1.654636in}{1.269418in}}%
\pgfpathlineto{\pgfqpoint{1.653431in}{1.223538in}}%
\pgfpathclose%
\pgfusepath{fill}%
\end{pgfscope}%
\begin{pgfscope}%
\pgfpathrectangle{\pgfqpoint{0.100000in}{0.100000in}}{\pgfqpoint{3.007045in}{1.925000in}}%
\pgfusepath{clip}%
\pgfsetbuttcap%
\pgfsetmiterjoin%
\definecolor{currentfill}{rgb}{0.260715,0.573841,0.777209}%
\pgfsetfillcolor{currentfill}%
\pgfsetlinewidth{0.000000pt}%
\definecolor{currentstroke}{rgb}{0.000000,0.000000,0.000000}%
\pgfsetstrokecolor{currentstroke}%
\pgfsetstrokeopacity{0.000000}%
\pgfsetdash{}{0pt}%
\pgfpathmoveto{\pgfqpoint{0.630991in}{1.684466in}}%
\pgfpathlineto{\pgfqpoint{0.628849in}{1.673291in}}%
\pgfpathlineto{\pgfqpoint{0.619536in}{1.640423in}}%
\pgfpathlineto{\pgfqpoint{0.608444in}{1.643504in}}%
\pgfpathlineto{\pgfqpoint{0.609918in}{1.648785in}}%
\pgfpathlineto{\pgfqpoint{0.605211in}{1.653481in}}%
\pgfpathlineto{\pgfqpoint{0.588656in}{1.658125in}}%
\pgfpathlineto{\pgfqpoint{0.594595in}{1.678719in}}%
\pgfpathlineto{\pgfqpoint{0.596303in}{1.682378in}}%
\pgfpathlineto{\pgfqpoint{0.593360in}{1.686222in}}%
\pgfpathlineto{\pgfqpoint{0.593379in}{1.692150in}}%
\pgfpathlineto{\pgfqpoint{0.598195in}{1.705394in}}%
\pgfpathlineto{\pgfqpoint{0.595993in}{1.712179in}}%
\pgfpathlineto{\pgfqpoint{0.602494in}{1.724955in}}%
\pgfpathlineto{\pgfqpoint{0.608895in}{1.725281in}}%
\pgfpathlineto{\pgfqpoint{0.608690in}{1.732151in}}%
\pgfpathlineto{\pgfqpoint{0.606293in}{1.736988in}}%
\pgfpathlineto{\pgfqpoint{0.610892in}{1.738691in}}%
\pgfpathlineto{\pgfqpoint{0.619982in}{1.737934in}}%
\pgfpathlineto{\pgfqpoint{0.630951in}{1.740790in}}%
\pgfpathlineto{\pgfqpoint{0.620630in}{1.705467in}}%
\pgfpathlineto{\pgfqpoint{0.626202in}{1.703795in}}%
\pgfpathlineto{\pgfqpoint{0.624611in}{1.698303in}}%
\pgfpathlineto{\pgfqpoint{0.630102in}{1.696622in}}%
\pgfpathlineto{\pgfqpoint{0.625624in}{1.685979in}}%
\pgfpathlineto{\pgfqpoint{0.630991in}{1.684466in}}%
\pgfpathclose%
\pgfusepath{fill}%
\end{pgfscope}%
\begin{pgfscope}%
\pgfpathrectangle{\pgfqpoint{0.100000in}{0.100000in}}{\pgfqpoint{3.007045in}{1.925000in}}%
\pgfusepath{clip}%
\pgfsetbuttcap%
\pgfsetmiterjoin%
\definecolor{currentfill}{rgb}{0.290980,0.594510,0.789020}%
\pgfsetfillcolor{currentfill}%
\pgfsetlinewidth{0.000000pt}%
\definecolor{currentstroke}{rgb}{0.000000,0.000000,0.000000}%
\pgfsetstrokecolor{currentstroke}%
\pgfsetstrokeopacity{0.000000}%
\pgfsetdash{}{0pt}%
\pgfpathmoveto{\pgfqpoint{1.510401in}{1.206498in}}%
\pgfpathlineto{\pgfqpoint{1.453477in}{1.210184in}}%
\pgfpathlineto{\pgfqpoint{1.455104in}{1.233038in}}%
\pgfpathlineto{\pgfqpoint{1.418635in}{1.235643in}}%
\pgfpathlineto{\pgfqpoint{1.420426in}{1.258524in}}%
\pgfpathlineto{\pgfqpoint{1.460453in}{1.255605in}}%
\pgfpathlineto{\pgfqpoint{1.483897in}{1.254208in}}%
\pgfpathlineto{\pgfqpoint{1.482455in}{1.231145in}}%
\pgfpathlineto{\pgfqpoint{1.511500in}{1.229428in}}%
\pgfpathlineto{\pgfqpoint{1.510401in}{1.206498in}}%
\pgfpathclose%
\pgfusepath{fill}%
\end{pgfscope}%
\begin{pgfscope}%
\pgfpathrectangle{\pgfqpoint{0.100000in}{0.100000in}}{\pgfqpoint{3.007045in}{1.925000in}}%
\pgfusepath{clip}%
\pgfsetbuttcap%
\pgfsetmiterjoin%
\definecolor{currentfill}{rgb}{0.479216,0.715079,0.852072}%
\pgfsetfillcolor{currentfill}%
\pgfsetlinewidth{0.000000pt}%
\definecolor{currentstroke}{rgb}{0.000000,0.000000,0.000000}%
\pgfsetstrokecolor{currentstroke}%
\pgfsetstrokeopacity{0.000000}%
\pgfsetdash{}{0pt}%
\pgfpathmoveto{\pgfqpoint{1.035972in}{1.058230in}}%
\pgfpathlineto{\pgfqpoint{0.986309in}{1.066160in}}%
\pgfpathlineto{\pgfqpoint{0.947453in}{1.073147in}}%
\pgfpathlineto{\pgfqpoint{0.915481in}{1.078785in}}%
\pgfpathlineto{\pgfqpoint{0.923032in}{1.079857in}}%
\pgfpathlineto{\pgfqpoint{0.926450in}{1.083530in}}%
\pgfpathlineto{\pgfqpoint{0.932657in}{1.082202in}}%
\pgfpathlineto{\pgfqpoint{0.941624in}{1.085677in}}%
\pgfpathlineto{\pgfqpoint{0.944689in}{1.092843in}}%
\pgfpathlineto{\pgfqpoint{0.951769in}{1.095219in}}%
\pgfpathlineto{\pgfqpoint{0.958933in}{1.107125in}}%
\pgfpathlineto{\pgfqpoint{0.961951in}{1.110649in}}%
\pgfpathlineto{\pgfqpoint{0.967830in}{1.113017in}}%
\pgfpathlineto{\pgfqpoint{0.975147in}{1.126984in}}%
\pgfpathlineto{\pgfqpoint{0.982614in}{1.127004in}}%
\pgfpathlineto{\pgfqpoint{0.990253in}{1.132805in}}%
\pgfpathlineto{\pgfqpoint{0.992394in}{1.131729in}}%
\pgfpathlineto{\pgfqpoint{0.997732in}{1.138462in}}%
\pgfpathlineto{\pgfqpoint{1.002522in}{1.140541in}}%
\pgfpathlineto{\pgfqpoint{1.005017in}{1.143547in}}%
\pgfpathlineto{\pgfqpoint{1.002930in}{1.149720in}}%
\pgfpathlineto{\pgfqpoint{1.001312in}{1.163987in}}%
\pgfpathlineto{\pgfqpoint{1.050444in}{1.156057in}}%
\pgfpathlineto{\pgfqpoint{1.048160in}{1.141447in}}%
\pgfpathlineto{\pgfqpoint{1.047972in}{1.134089in}}%
\pgfpathlineto{\pgfqpoint{1.045110in}{1.115632in}}%
\pgfpathlineto{\pgfqpoint{1.056798in}{1.114784in}}%
\pgfpathlineto{\pgfqpoint{1.085598in}{1.110401in}}%
\pgfpathlineto{\pgfqpoint{1.087280in}{1.105261in}}%
\pgfpathlineto{\pgfqpoint{1.098980in}{1.105934in}}%
\pgfpathlineto{\pgfqpoint{1.104848in}{1.099688in}}%
\pgfpathlineto{\pgfqpoint{1.103247in}{1.096115in}}%
\pgfpathlineto{\pgfqpoint{1.094867in}{1.088434in}}%
\pgfpathlineto{\pgfqpoint{1.093042in}{1.080708in}}%
\pgfpathlineto{\pgfqpoint{1.088472in}{1.075836in}}%
\pgfpathlineto{\pgfqpoint{1.083485in}{1.074787in}}%
\pgfpathlineto{\pgfqpoint{1.077407in}{1.066910in}}%
\pgfpathlineto{\pgfqpoint{1.073504in}{1.056252in}}%
\pgfpathlineto{\pgfqpoint{1.070635in}{1.052989in}}%
\pgfpathlineto{\pgfqpoint{1.035972in}{1.058230in}}%
\pgfpathclose%
\pgfusepath{fill}%
\end{pgfscope}%
\begin{pgfscope}%
\pgfpathrectangle{\pgfqpoint{0.100000in}{0.100000in}}{\pgfqpoint{3.007045in}{1.925000in}}%
\pgfusepath{clip}%
\pgfsetbuttcap%
\pgfsetmiterjoin%
\definecolor{currentfill}{rgb}{0.417086,0.680631,0.838231}%
\pgfsetfillcolor{currentfill}%
\pgfsetlinewidth{0.000000pt}%
\definecolor{currentstroke}{rgb}{0.000000,0.000000,0.000000}%
\pgfsetstrokecolor{currentstroke}%
\pgfsetstrokeopacity{0.000000}%
\pgfsetdash{}{0pt}%
\pgfpathmoveto{\pgfqpoint{2.241623in}{0.891209in}}%
\pgfpathlineto{\pgfqpoint{2.240204in}{0.883283in}}%
\pgfpathlineto{\pgfqpoint{2.241813in}{0.880528in}}%
\pgfpathlineto{\pgfqpoint{2.241299in}{0.867393in}}%
\pgfpathlineto{\pgfqpoint{2.244466in}{0.861498in}}%
\pgfpathlineto{\pgfqpoint{2.238722in}{0.856984in}}%
\pgfpathlineto{\pgfqpoint{2.232658in}{0.857963in}}%
\pgfpathlineto{\pgfqpoint{2.231240in}{0.844683in}}%
\pgfpathlineto{\pgfqpoint{2.202469in}{0.842698in}}%
\pgfpathlineto{\pgfqpoint{2.202564in}{0.841744in}}%
\pgfpathlineto{\pgfqpoint{2.179696in}{0.840159in}}%
\pgfpathlineto{\pgfqpoint{2.178278in}{0.857446in}}%
\pgfpathlineto{\pgfqpoint{2.183086in}{0.867245in}}%
\pgfpathlineto{\pgfqpoint{2.182587in}{0.873244in}}%
\pgfpathlineto{\pgfqpoint{2.189352in}{0.871261in}}%
\pgfpathlineto{\pgfqpoint{2.193931in}{0.875265in}}%
\pgfpathlineto{\pgfqpoint{2.193121in}{0.887325in}}%
\pgfpathlineto{\pgfqpoint{2.213336in}{0.888623in}}%
\pgfpathlineto{\pgfqpoint{2.212159in}{0.906542in}}%
\pgfpathlineto{\pgfqpoint{2.216592in}{0.906228in}}%
\pgfpathlineto{\pgfqpoint{2.223886in}{0.911812in}}%
\pgfpathlineto{\pgfqpoint{2.223823in}{0.914321in}}%
\pgfpathlineto{\pgfqpoint{2.235740in}{0.906849in}}%
\pgfpathlineto{\pgfqpoint{2.237551in}{0.900340in}}%
\pgfpathlineto{\pgfqpoint{2.240394in}{0.900117in}}%
\pgfpathlineto{\pgfqpoint{2.241623in}{0.891209in}}%
\pgfpathclose%
\pgfusepath{fill}%
\end{pgfscope}%
\begin{pgfscope}%
\pgfpathrectangle{\pgfqpoint{0.100000in}{0.100000in}}{\pgfqpoint{3.007045in}{1.925000in}}%
\pgfusepath{clip}%
\pgfsetbuttcap%
\pgfsetmiterjoin%
\definecolor{currentfill}{rgb}{0.361599,0.642737,0.816578}%
\pgfsetfillcolor{currentfill}%
\pgfsetlinewidth{0.000000pt}%
\definecolor{currentstroke}{rgb}{0.000000,0.000000,0.000000}%
\pgfsetstrokecolor{currentstroke}%
\pgfsetstrokeopacity{0.000000}%
\pgfsetdash{}{0pt}%
\pgfpathmoveto{\pgfqpoint{1.062663in}{1.231515in}}%
\pgfpathlineto{\pgfqpoint{1.072883in}{1.296139in}}%
\pgfpathlineto{\pgfqpoint{1.076315in}{1.318750in}}%
\pgfpathlineto{\pgfqpoint{1.102069in}{1.314741in}}%
\pgfpathlineto{\pgfqpoint{1.131992in}{1.310519in}}%
\pgfpathlineto{\pgfqpoint{1.184259in}{1.303392in}}%
\pgfpathlineto{\pgfqpoint{1.183883in}{1.298271in}}%
\pgfpathlineto{\pgfqpoint{1.192988in}{1.292052in}}%
\pgfpathlineto{\pgfqpoint{1.193426in}{1.288196in}}%
\pgfpathlineto{\pgfqpoint{1.188332in}{1.274642in}}%
\pgfpathlineto{\pgfqpoint{1.189706in}{1.265642in}}%
\pgfpathlineto{\pgfqpoint{1.189821in}{1.258740in}}%
\pgfpathlineto{\pgfqpoint{1.186728in}{1.236591in}}%
\pgfpathlineto{\pgfqpoint{1.186601in}{1.231472in}}%
\pgfpathlineto{\pgfqpoint{1.166255in}{1.233746in}}%
\pgfpathlineto{\pgfqpoint{1.167476in}{1.245055in}}%
\pgfpathlineto{\pgfqpoint{1.153503in}{1.246887in}}%
\pgfpathlineto{\pgfqpoint{1.151850in}{1.235425in}}%
\pgfpathlineto{\pgfqpoint{1.146223in}{1.236379in}}%
\pgfpathlineto{\pgfqpoint{1.145415in}{1.230549in}}%
\pgfpathlineto{\pgfqpoint{1.120000in}{1.233895in}}%
\pgfpathlineto{\pgfqpoint{1.118778in}{1.225384in}}%
\pgfpathlineto{\pgfqpoint{1.087872in}{1.229867in}}%
\pgfpathlineto{\pgfqpoint{1.087451in}{1.227034in}}%
\pgfpathlineto{\pgfqpoint{1.062663in}{1.231515in}}%
\pgfpathclose%
\pgfusepath{fill}%
\end{pgfscope}%
\begin{pgfscope}%
\pgfpathrectangle{\pgfqpoint{0.100000in}{0.100000in}}{\pgfqpoint{3.007045in}{1.925000in}}%
\pgfusepath{clip}%
\pgfsetbuttcap%
\pgfsetmiterjoin%
\definecolor{currentfill}{rgb}{0.050211,0.341761,0.630604}%
\pgfsetfillcolor{currentfill}%
\pgfsetlinewidth{0.000000pt}%
\definecolor{currentstroke}{rgb}{0.000000,0.000000,0.000000}%
\pgfsetstrokecolor{currentstroke}%
\pgfsetstrokeopacity{0.000000}%
\pgfsetdash{}{0pt}%
\pgfpathmoveto{\pgfqpoint{1.441506in}{0.486973in}}%
\pgfpathlineto{\pgfqpoint{1.473940in}{0.485183in}}%
\pgfpathlineto{\pgfqpoint{1.470461in}{0.425819in}}%
\pgfpathlineto{\pgfqpoint{1.464650in}{0.426083in}}%
\pgfpathlineto{\pgfqpoint{1.460228in}{0.431872in}}%
\pgfpathlineto{\pgfqpoint{1.458257in}{0.442114in}}%
\pgfpathlineto{\pgfqpoint{1.454336in}{0.450453in}}%
\pgfpathlineto{\pgfqpoint{1.455112in}{0.452787in}}%
\pgfpathlineto{\pgfqpoint{1.449609in}{0.458112in}}%
\pgfpathlineto{\pgfqpoint{1.447489in}{0.469258in}}%
\pgfpathlineto{\pgfqpoint{1.441856in}{0.477516in}}%
\pgfpathlineto{\pgfqpoint{1.441506in}{0.486973in}}%
\pgfpathclose%
\pgfusepath{fill}%
\end{pgfscope}%
\begin{pgfscope}%
\pgfpathrectangle{\pgfqpoint{0.100000in}{0.100000in}}{\pgfqpoint{3.007045in}{1.925000in}}%
\pgfusepath{clip}%
\pgfsetbuttcap%
\pgfsetmiterjoin%
\definecolor{currentfill}{rgb}{0.371688,0.649627,0.820515}%
\pgfsetfillcolor{currentfill}%
\pgfsetlinewidth{0.000000pt}%
\definecolor{currentstroke}{rgb}{0.000000,0.000000,0.000000}%
\pgfsetstrokecolor{currentstroke}%
\pgfsetstrokeopacity{0.000000}%
\pgfsetdash{}{0pt}%
\pgfpathmoveto{\pgfqpoint{1.775875in}{0.660550in}}%
\pgfpathlineto{\pgfqpoint{1.773764in}{0.656330in}}%
\pgfpathlineto{\pgfqpoint{1.776581in}{0.643123in}}%
\pgfpathlineto{\pgfqpoint{1.779925in}{0.639399in}}%
\pgfpathlineto{\pgfqpoint{1.772144in}{0.632693in}}%
\pgfpathlineto{\pgfqpoint{1.766321in}{0.636474in}}%
\pgfpathlineto{\pgfqpoint{1.763851in}{0.642226in}}%
\pgfpathlineto{\pgfqpoint{1.756860in}{0.643794in}}%
\pgfpathlineto{\pgfqpoint{1.735433in}{0.640512in}}%
\pgfpathlineto{\pgfqpoint{1.730423in}{0.638033in}}%
\pgfpathlineto{\pgfqpoint{1.732103in}{0.644712in}}%
\pgfpathlineto{\pgfqpoint{1.723459in}{0.650585in}}%
\pgfpathlineto{\pgfqpoint{1.722872in}{0.654828in}}%
\pgfpathlineto{\pgfqpoint{1.718592in}{0.656205in}}%
\pgfpathlineto{\pgfqpoint{1.712375in}{0.673683in}}%
\pgfpathlineto{\pgfqpoint{1.703043in}{0.687922in}}%
\pgfpathlineto{\pgfqpoint{1.694814in}{0.693189in}}%
\pgfpathlineto{\pgfqpoint{1.712066in}{0.694873in}}%
\pgfpathlineto{\pgfqpoint{1.712393in}{0.726751in}}%
\pgfpathlineto{\pgfqpoint{1.720316in}{0.726579in}}%
\pgfpathlineto{\pgfqpoint{1.725321in}{0.722021in}}%
\pgfpathlineto{\pgfqpoint{1.728058in}{0.722452in}}%
\pgfpathlineto{\pgfqpoint{1.736915in}{0.718690in}}%
\pgfpathlineto{\pgfqpoint{1.735350in}{0.734610in}}%
\pgfpathlineto{\pgfqpoint{1.755158in}{0.734645in}}%
\pgfpathlineto{\pgfqpoint{1.762705in}{0.734578in}}%
\pgfpathlineto{\pgfqpoint{1.763868in}{0.730641in}}%
\pgfpathlineto{\pgfqpoint{1.763758in}{0.708620in}}%
\pgfpathlineto{\pgfqpoint{1.773020in}{0.706433in}}%
\pgfpathlineto{\pgfqpoint{1.773171in}{0.660592in}}%
\pgfpathlineto{\pgfqpoint{1.775875in}{0.660550in}}%
\pgfpathclose%
\pgfusepath{fill}%
\end{pgfscope}%
\begin{pgfscope}%
\pgfpathrectangle{\pgfqpoint{0.100000in}{0.100000in}}{\pgfqpoint{3.007045in}{1.925000in}}%
\pgfusepath{clip}%
\pgfsetbuttcap%
\pgfsetmiterjoin%
\definecolor{currentfill}{rgb}{0.361599,0.642737,0.816578}%
\pgfsetfillcolor{currentfill}%
\pgfsetlinewidth{0.000000pt}%
\definecolor{currentstroke}{rgb}{0.000000,0.000000,0.000000}%
\pgfsetstrokecolor{currentstroke}%
\pgfsetstrokeopacity{0.000000}%
\pgfsetdash{}{0pt}%
\pgfpathmoveto{\pgfqpoint{0.407690in}{1.233195in}}%
\pgfpathlineto{\pgfqpoint{0.404604in}{1.238315in}}%
\pgfpathlineto{\pgfqpoint{0.409870in}{1.255826in}}%
\pgfpathlineto{\pgfqpoint{0.409743in}{1.259767in}}%
\pgfpathlineto{\pgfqpoint{0.413844in}{1.273461in}}%
\pgfpathlineto{\pgfqpoint{0.409876in}{1.274771in}}%
\pgfpathlineto{\pgfqpoint{0.409917in}{1.279165in}}%
\pgfpathlineto{\pgfqpoint{0.414282in}{1.280552in}}%
\pgfpathlineto{\pgfqpoint{0.417431in}{1.287454in}}%
\pgfpathlineto{\pgfqpoint{0.412425in}{1.289011in}}%
\pgfpathlineto{\pgfqpoint{0.416471in}{1.302466in}}%
\pgfpathlineto{\pgfqpoint{0.404388in}{1.306594in}}%
\pgfpathlineto{\pgfqpoint{0.400092in}{1.304827in}}%
\pgfpathlineto{\pgfqpoint{0.396070in}{1.307741in}}%
\pgfpathlineto{\pgfqpoint{0.395623in}{1.317666in}}%
\pgfpathlineto{\pgfqpoint{0.393792in}{1.321038in}}%
\pgfpathlineto{\pgfqpoint{0.393337in}{1.331101in}}%
\pgfpathlineto{\pgfqpoint{0.388567in}{1.334258in}}%
\pgfpathlineto{\pgfqpoint{0.392426in}{1.337198in}}%
\pgfpathlineto{\pgfqpoint{0.417135in}{1.329622in}}%
\pgfpathlineto{\pgfqpoint{0.421047in}{1.323928in}}%
\pgfpathlineto{\pgfqpoint{0.420345in}{1.320339in}}%
\pgfpathlineto{\pgfqpoint{0.424973in}{1.314267in}}%
\pgfpathlineto{\pgfqpoint{0.431828in}{1.313269in}}%
\pgfpathlineto{\pgfqpoint{0.435140in}{1.324290in}}%
\pgfpathlineto{\pgfqpoint{0.439260in}{1.328328in}}%
\pgfpathlineto{\pgfqpoint{0.445384in}{1.329930in}}%
\pgfpathlineto{\pgfqpoint{0.450179in}{1.340731in}}%
\pgfpathlineto{\pgfqpoint{0.460387in}{1.347970in}}%
\pgfpathlineto{\pgfqpoint{0.465761in}{1.347164in}}%
\pgfpathlineto{\pgfqpoint{0.480726in}{1.347606in}}%
\pgfpathlineto{\pgfqpoint{0.486181in}{1.350484in}}%
\pgfpathlineto{\pgfqpoint{0.489895in}{1.349175in}}%
\pgfpathlineto{\pgfqpoint{0.492015in}{1.343273in}}%
\pgfpathlineto{\pgfqpoint{0.516565in}{1.336204in}}%
\pgfpathlineto{\pgfqpoint{0.510601in}{1.315136in}}%
\pgfpathlineto{\pgfqpoint{0.522542in}{1.311873in}}%
\pgfpathlineto{\pgfqpoint{0.522474in}{1.309901in}}%
\pgfpathlineto{\pgfqpoint{0.544151in}{1.303815in}}%
\pgfpathlineto{\pgfqpoint{0.542583in}{1.297900in}}%
\pgfpathlineto{\pgfqpoint{0.536518in}{1.297107in}}%
\pgfpathlineto{\pgfqpoint{0.533368in}{1.282853in}}%
\pgfpathlineto{\pgfqpoint{0.536090in}{1.281820in}}%
\pgfpathlineto{\pgfqpoint{0.533766in}{1.269179in}}%
\pgfpathlineto{\pgfqpoint{0.524302in}{1.283464in}}%
\pgfpathlineto{\pgfqpoint{0.520542in}{1.277526in}}%
\pgfpathlineto{\pgfqpoint{0.522497in}{1.269886in}}%
\pgfpathlineto{\pgfqpoint{0.519485in}{1.264180in}}%
\pgfpathlineto{\pgfqpoint{0.514663in}{1.260246in}}%
\pgfpathlineto{\pgfqpoint{0.510697in}{1.267522in}}%
\pgfpathlineto{\pgfqpoint{0.502964in}{1.265519in}}%
\pgfpathlineto{\pgfqpoint{0.497623in}{1.272306in}}%
\pgfpathlineto{\pgfqpoint{0.488781in}{1.270732in}}%
\pgfpathlineto{\pgfqpoint{0.475080in}{1.259356in}}%
\pgfpathlineto{\pgfqpoint{0.464303in}{1.248933in}}%
\pgfpathlineto{\pgfqpoint{0.455151in}{1.243275in}}%
\pgfpathlineto{\pgfqpoint{0.446079in}{1.262743in}}%
\pgfpathlineto{\pgfqpoint{0.440085in}{1.241248in}}%
\pgfpathlineto{\pgfqpoint{0.436885in}{1.243705in}}%
\pgfpathlineto{\pgfqpoint{0.424991in}{1.242670in}}%
\pgfpathlineto{\pgfqpoint{0.407690in}{1.233195in}}%
\pgfpathclose%
\pgfusepath{fill}%
\end{pgfscope}%
\begin{pgfscope}%
\pgfpathrectangle{\pgfqpoint{0.100000in}{0.100000in}}{\pgfqpoint{3.007045in}{1.925000in}}%
\pgfusepath{clip}%
\pgfsetbuttcap%
\pgfsetmiterjoin%
\definecolor{currentfill}{rgb}{0.296025,0.597955,0.790988}%
\pgfsetfillcolor{currentfill}%
\pgfsetlinewidth{0.000000pt}%
\definecolor{currentstroke}{rgb}{0.000000,0.000000,0.000000}%
\pgfsetstrokecolor{currentstroke}%
\pgfsetstrokeopacity{0.000000}%
\pgfsetdash{}{0pt}%
\pgfpathmoveto{\pgfqpoint{1.588665in}{1.618076in}}%
\pgfpathlineto{\pgfqpoint{1.587099in}{1.618157in}}%
\pgfpathlineto{\pgfqpoint{1.588163in}{1.641202in}}%
\pgfpathlineto{\pgfqpoint{1.569331in}{1.642166in}}%
\pgfpathlineto{\pgfqpoint{1.570614in}{1.665351in}}%
\pgfpathlineto{\pgfqpoint{1.569042in}{1.665415in}}%
\pgfpathlineto{\pgfqpoint{1.570193in}{1.688446in}}%
\pgfpathlineto{\pgfqpoint{1.579862in}{1.687980in}}%
\pgfpathlineto{\pgfqpoint{1.614445in}{1.686458in}}%
\pgfpathlineto{\pgfqpoint{1.615660in}{1.680679in}}%
\pgfpathlineto{\pgfqpoint{1.638529in}{1.679874in}}%
\pgfpathlineto{\pgfqpoint{1.650048in}{1.679481in}}%
\pgfpathlineto{\pgfqpoint{1.650635in}{1.662065in}}%
\pgfpathlineto{\pgfqpoint{1.649938in}{1.638890in}}%
\pgfpathlineto{\pgfqpoint{1.633892in}{1.639443in}}%
\pgfpathlineto{\pgfqpoint{1.633089in}{1.616328in}}%
\pgfpathlineto{\pgfqpoint{1.588665in}{1.618076in}}%
\pgfpathclose%
\pgfusepath{fill}%
\end{pgfscope}%
\begin{pgfscope}%
\pgfpathrectangle{\pgfqpoint{0.100000in}{0.100000in}}{\pgfqpoint{3.007045in}{1.925000in}}%
\pgfusepath{clip}%
\pgfsetbuttcap%
\pgfsetmiterjoin%
\definecolor{currentfill}{rgb}{0.790496,0.868174,0.941930}%
\pgfsetfillcolor{currentfill}%
\pgfsetlinewidth{0.000000pt}%
\definecolor{currentstroke}{rgb}{0.000000,0.000000,0.000000}%
\pgfsetstrokecolor{currentstroke}%
\pgfsetstrokeopacity{0.000000}%
\pgfsetdash{}{0pt}%
\pgfpathmoveto{\pgfqpoint{2.867151in}{1.544262in}}%
\pgfpathlineto{\pgfqpoint{2.860753in}{1.546240in}}%
\pgfpathlineto{\pgfqpoint{2.850891in}{1.542069in}}%
\pgfpathlineto{\pgfqpoint{2.846989in}{1.546225in}}%
\pgfpathlineto{\pgfqpoint{2.835404in}{1.542568in}}%
\pgfpathlineto{\pgfqpoint{2.828637in}{1.543127in}}%
\pgfpathlineto{\pgfqpoint{2.825651in}{1.545463in}}%
\pgfpathlineto{\pgfqpoint{2.829267in}{1.550294in}}%
\pgfpathlineto{\pgfqpoint{2.829522in}{1.557864in}}%
\pgfpathlineto{\pgfqpoint{2.826130in}{1.558806in}}%
\pgfpathlineto{\pgfqpoint{2.827305in}{1.571893in}}%
\pgfpathlineto{\pgfqpoint{2.821240in}{1.572694in}}%
\pgfpathlineto{\pgfqpoint{2.821564in}{1.578723in}}%
\pgfpathlineto{\pgfqpoint{2.811802in}{1.580671in}}%
\pgfpathlineto{\pgfqpoint{2.815683in}{1.583923in}}%
\pgfpathlineto{\pgfqpoint{2.819866in}{1.593654in}}%
\pgfpathlineto{\pgfqpoint{2.811115in}{1.594372in}}%
\pgfpathlineto{\pgfqpoint{2.808838in}{1.600051in}}%
\pgfpathlineto{\pgfqpoint{2.808017in}{1.607896in}}%
\pgfpathlineto{\pgfqpoint{2.810189in}{1.620019in}}%
\pgfpathlineto{\pgfqpoint{2.817950in}{1.618314in}}%
\pgfpathlineto{\pgfqpoint{2.819281in}{1.624179in}}%
\pgfpathlineto{\pgfqpoint{2.825014in}{1.621968in}}%
\pgfpathlineto{\pgfqpoint{2.826283in}{1.627777in}}%
\pgfpathlineto{\pgfqpoint{2.837240in}{1.625008in}}%
\pgfpathlineto{\pgfqpoint{2.834224in}{1.609178in}}%
\pgfpathlineto{\pgfqpoint{2.837477in}{1.608579in}}%
\pgfpathlineto{\pgfqpoint{2.848338in}{1.611817in}}%
\pgfpathlineto{\pgfqpoint{2.846880in}{1.617175in}}%
\pgfpathlineto{\pgfqpoint{2.849656in}{1.622164in}}%
\pgfpathlineto{\pgfqpoint{2.854814in}{1.623529in}}%
\pgfpathlineto{\pgfqpoint{2.858546in}{1.629287in}}%
\pgfpathlineto{\pgfqpoint{2.862905in}{1.625347in}}%
\pgfpathlineto{\pgfqpoint{2.868383in}{1.624903in}}%
\pgfpathlineto{\pgfqpoint{2.869747in}{1.621836in}}%
\pgfpathlineto{\pgfqpoint{2.875904in}{1.623095in}}%
\pgfpathlineto{\pgfqpoint{2.877796in}{1.621050in}}%
\pgfpathlineto{\pgfqpoint{2.882477in}{1.612503in}}%
\pgfpathlineto{\pgfqpoint{2.885308in}{1.602759in}}%
\pgfpathlineto{\pgfqpoint{2.876355in}{1.598852in}}%
\pgfpathlineto{\pgfqpoint{2.881617in}{1.586847in}}%
\pgfpathlineto{\pgfqpoint{2.875977in}{1.584608in}}%
\pgfpathlineto{\pgfqpoint{2.876946in}{1.580426in}}%
\pgfpathlineto{\pgfqpoint{2.870513in}{1.573950in}}%
\pgfpathlineto{\pgfqpoint{2.867917in}{1.577405in}}%
\pgfpathlineto{\pgfqpoint{2.863308in}{1.564293in}}%
\pgfpathlineto{\pgfqpoint{2.861539in}{1.556030in}}%
\pgfpathlineto{\pgfqpoint{2.867535in}{1.548390in}}%
\pgfpathlineto{\pgfqpoint{2.867151in}{1.544262in}}%
\pgfpathclose%
\pgfusepath{fill}%
\end{pgfscope}%
\begin{pgfscope}%
\pgfpathrectangle{\pgfqpoint{0.100000in}{0.100000in}}{\pgfqpoint{3.007045in}{1.925000in}}%
\pgfusepath{clip}%
\pgfsetbuttcap%
\pgfsetmiterjoin%
\definecolor{currentfill}{rgb}{0.554510,0.756417,0.868312}%
\pgfsetfillcolor{currentfill}%
\pgfsetlinewidth{0.000000pt}%
\definecolor{currentstroke}{rgb}{0.000000,0.000000,0.000000}%
\pgfsetstrokecolor{currentstroke}%
\pgfsetstrokeopacity{0.000000}%
\pgfsetdash{}{0pt}%
\pgfpathmoveto{\pgfqpoint{2.552895in}{1.336088in}}%
\pgfpathlineto{\pgfqpoint{2.550029in}{1.332847in}}%
\pgfpathlineto{\pgfqpoint{2.540085in}{1.329055in}}%
\pgfpathlineto{\pgfqpoint{2.533624in}{1.328350in}}%
\pgfpathlineto{\pgfqpoint{2.528039in}{1.339978in}}%
\pgfpathlineto{\pgfqpoint{2.513140in}{1.337398in}}%
\pgfpathlineto{\pgfqpoint{2.509694in}{1.358087in}}%
\pgfpathlineto{\pgfqpoint{2.503675in}{1.356445in}}%
\pgfpathlineto{\pgfqpoint{2.484276in}{1.354588in}}%
\pgfpathlineto{\pgfqpoint{2.479266in}{1.385678in}}%
\pgfpathlineto{\pgfqpoint{2.494271in}{1.396072in}}%
\pgfpathlineto{\pgfqpoint{2.507781in}{1.407380in}}%
\pgfpathlineto{\pgfqpoint{2.516225in}{1.413233in}}%
\pgfpathlineto{\pgfqpoint{2.526196in}{1.422773in}}%
\pgfpathlineto{\pgfqpoint{2.534317in}{1.431974in}}%
\pgfpathlineto{\pgfqpoint{2.539326in}{1.435418in}}%
\pgfpathlineto{\pgfqpoint{2.543334in}{1.433990in}}%
\pgfpathlineto{\pgfqpoint{2.549510in}{1.399000in}}%
\pgfpathlineto{\pgfqpoint{2.556395in}{1.400149in}}%
\pgfpathlineto{\pgfqpoint{2.558369in}{1.389969in}}%
\pgfpathlineto{\pgfqpoint{2.556898in}{1.389007in}}%
\pgfpathlineto{\pgfqpoint{2.558934in}{1.375533in}}%
\pgfpathlineto{\pgfqpoint{2.559925in}{1.364184in}}%
\pgfpathlineto{\pgfqpoint{2.554563in}{1.361326in}}%
\pgfpathlineto{\pgfqpoint{2.555437in}{1.355961in}}%
\pgfpathlineto{\pgfqpoint{2.549986in}{1.354398in}}%
\pgfpathlineto{\pgfqpoint{2.552895in}{1.336088in}}%
\pgfpathclose%
\pgfusepath{fill}%
\end{pgfscope}%
\begin{pgfscope}%
\pgfpathrectangle{\pgfqpoint{0.100000in}{0.100000in}}{\pgfqpoint{3.007045in}{1.925000in}}%
\pgfusepath{clip}%
\pgfsetbuttcap%
\pgfsetmiterjoin%
\definecolor{currentfill}{rgb}{0.187266,0.500992,0.739608}%
\pgfsetfillcolor{currentfill}%
\pgfsetlinewidth{0.000000pt}%
\definecolor{currentstroke}{rgb}{0.000000,0.000000,0.000000}%
\pgfsetstrokecolor{currentstroke}%
\pgfsetstrokeopacity{0.000000}%
\pgfsetdash{}{0pt}%
\pgfpathmoveto{\pgfqpoint{1.677475in}{1.268823in}}%
\pgfpathlineto{\pgfqpoint{1.654636in}{1.269418in}}%
\pgfpathlineto{\pgfqpoint{1.655235in}{1.292354in}}%
\pgfpathlineto{\pgfqpoint{1.643733in}{1.288581in}}%
\pgfpathlineto{\pgfqpoint{1.643851in}{1.292672in}}%
\pgfpathlineto{\pgfqpoint{1.638646in}{1.292816in}}%
\pgfpathlineto{\pgfqpoint{1.638392in}{1.287167in}}%
\pgfpathlineto{\pgfqpoint{1.621155in}{1.286648in}}%
\pgfpathlineto{\pgfqpoint{1.618684in}{1.284817in}}%
\pgfpathlineto{\pgfqpoint{1.610083in}{1.285141in}}%
\pgfpathlineto{\pgfqpoint{1.609594in}{1.293697in}}%
\pgfpathlineto{\pgfqpoint{1.610215in}{1.316599in}}%
\pgfpathlineto{\pgfqpoint{1.610638in}{1.328045in}}%
\pgfpathlineto{\pgfqpoint{1.633318in}{1.327373in}}%
\pgfpathlineto{\pgfqpoint{1.633082in}{1.315887in}}%
\pgfpathlineto{\pgfqpoint{1.672961in}{1.314877in}}%
\pgfpathlineto{\pgfqpoint{1.678558in}{1.314735in}}%
\pgfpathlineto{\pgfqpoint{1.677475in}{1.268823in}}%
\pgfpathclose%
\pgfusepath{fill}%
\end{pgfscope}%
\begin{pgfscope}%
\pgfpathrectangle{\pgfqpoint{0.100000in}{0.100000in}}{\pgfqpoint{3.007045in}{1.925000in}}%
\pgfusepath{clip}%
\pgfsetbuttcap%
\pgfsetmiterjoin%
\definecolor{currentfill}{rgb}{0.491765,0.721968,0.854779}%
\pgfsetfillcolor{currentfill}%
\pgfsetlinewidth{0.000000pt}%
\definecolor{currentstroke}{rgb}{0.000000,0.000000,0.000000}%
\pgfsetstrokecolor{currentstroke}%
\pgfsetstrokeopacity{0.000000}%
\pgfsetdash{}{0pt}%
\pgfpathmoveto{\pgfqpoint{1.779266in}{1.452736in}}%
\pgfpathlineto{\pgfqpoint{1.778980in}{1.469921in}}%
\pgfpathlineto{\pgfqpoint{1.767569in}{1.469929in}}%
\pgfpathlineto{\pgfqpoint{1.767169in}{1.475679in}}%
\pgfpathlineto{\pgfqpoint{1.767240in}{1.481415in}}%
\pgfpathlineto{\pgfqpoint{1.778661in}{1.481448in}}%
\pgfpathlineto{\pgfqpoint{1.778660in}{1.495730in}}%
\pgfpathlineto{\pgfqpoint{1.782199in}{1.492987in}}%
\pgfpathlineto{\pgfqpoint{1.790073in}{1.492987in}}%
\pgfpathlineto{\pgfqpoint{1.822887in}{1.493279in}}%
\pgfpathlineto{\pgfqpoint{1.823694in}{1.499014in}}%
\pgfpathlineto{\pgfqpoint{1.841921in}{1.499291in}}%
\pgfpathlineto{\pgfqpoint{1.842251in}{1.476305in}}%
\pgfpathlineto{\pgfqpoint{1.830743in}{1.476124in}}%
\pgfpathlineto{\pgfqpoint{1.831065in}{1.453126in}}%
\pgfpathlineto{\pgfqpoint{1.779266in}{1.452736in}}%
\pgfpathclose%
\pgfusepath{fill}%
\end{pgfscope}%
\begin{pgfscope}%
\pgfpathrectangle{\pgfqpoint{0.100000in}{0.100000in}}{\pgfqpoint{3.007045in}{1.925000in}}%
\pgfusepath{clip}%
\pgfsetbuttcap%
\pgfsetmiterjoin%
\definecolor{currentfill}{rgb}{0.356555,0.639293,0.814610}%
\pgfsetfillcolor{currentfill}%
\pgfsetlinewidth{0.000000pt}%
\definecolor{currentstroke}{rgb}{0.000000,0.000000,0.000000}%
\pgfsetstrokecolor{currentstroke}%
\pgfsetstrokeopacity{0.000000}%
\pgfsetdash{}{0pt}%
\pgfpathmoveto{\pgfqpoint{1.890498in}{1.239370in}}%
\pgfpathlineto{\pgfqpoint{1.889907in}{1.259660in}}%
\pgfpathlineto{\pgfqpoint{1.867080in}{1.259096in}}%
\pgfpathlineto{\pgfqpoint{1.866638in}{1.276378in}}%
\pgfpathlineto{\pgfqpoint{1.855298in}{1.276147in}}%
\pgfpathlineto{\pgfqpoint{1.854836in}{1.299029in}}%
\pgfpathlineto{\pgfqpoint{1.905724in}{1.300373in}}%
\pgfpathlineto{\pgfqpoint{1.923061in}{1.300995in}}%
\pgfpathlineto{\pgfqpoint{1.923843in}{1.278117in}}%
\pgfpathlineto{\pgfqpoint{1.935244in}{1.278471in}}%
\pgfpathlineto{\pgfqpoint{1.935891in}{1.242279in}}%
\pgfpathlineto{\pgfqpoint{1.925211in}{1.241392in}}%
\pgfpathlineto{\pgfqpoint{1.890498in}{1.239370in}}%
\pgfpathclose%
\pgfusepath{fill}%
\end{pgfscope}%
\begin{pgfscope}%
\pgfpathrectangle{\pgfqpoint{0.100000in}{0.100000in}}{\pgfqpoint{3.007045in}{1.925000in}}%
\pgfusepath{clip}%
\pgfsetbuttcap%
\pgfsetmiterjoin%
\definecolor{currentfill}{rgb}{0.037908,0.326013,0.618301}%
\pgfsetfillcolor{currentfill}%
\pgfsetlinewidth{0.000000pt}%
\definecolor{currentstroke}{rgb}{0.000000,0.000000,0.000000}%
\pgfsetstrokecolor{currentstroke}%
\pgfsetstrokeopacity{0.000000}%
\pgfsetdash{}{0pt}%
\pgfpathmoveto{\pgfqpoint{0.967728in}{1.403968in}}%
\pgfpathlineto{\pgfqpoint{0.968744in}{1.400234in}}%
\pgfpathlineto{\pgfqpoint{0.965948in}{1.394443in}}%
\pgfpathlineto{\pgfqpoint{0.967550in}{1.392277in}}%
\pgfpathlineto{\pgfqpoint{0.968734in}{1.380626in}}%
\pgfpathlineto{\pgfqpoint{0.965270in}{1.373029in}}%
\pgfpathlineto{\pgfqpoint{0.960688in}{1.366618in}}%
\pgfpathlineto{\pgfqpoint{0.951603in}{1.367889in}}%
\pgfpathlineto{\pgfqpoint{0.947822in}{1.365378in}}%
\pgfpathlineto{\pgfqpoint{0.942058in}{1.369750in}}%
\pgfpathlineto{\pgfqpoint{0.934476in}{1.365698in}}%
\pgfpathlineto{\pgfqpoint{0.924294in}{1.367661in}}%
\pgfpathlineto{\pgfqpoint{0.908601in}{1.353242in}}%
\pgfpathlineto{\pgfqpoint{0.892707in}{1.351230in}}%
\pgfpathlineto{\pgfqpoint{0.832155in}{1.363654in}}%
\pgfpathlineto{\pgfqpoint{0.845804in}{1.427734in}}%
\pgfpathlineto{\pgfqpoint{0.853680in}{1.425681in}}%
\pgfpathlineto{\pgfqpoint{0.895826in}{1.417699in}}%
\pgfpathlineto{\pgfqpoint{0.900100in}{1.438993in}}%
\pgfpathlineto{\pgfqpoint{0.916709in}{1.435786in}}%
\pgfpathlineto{\pgfqpoint{0.917819in}{1.441425in}}%
\pgfpathlineto{\pgfqpoint{0.925285in}{1.439963in}}%
\pgfpathlineto{\pgfqpoint{0.926367in}{1.445605in}}%
\pgfpathlineto{\pgfqpoint{0.931789in}{1.444541in}}%
\pgfpathlineto{\pgfqpoint{0.934955in}{1.436848in}}%
\pgfpathlineto{\pgfqpoint{0.941439in}{1.428069in}}%
\pgfpathlineto{\pgfqpoint{0.948249in}{1.426793in}}%
\pgfpathlineto{\pgfqpoint{0.952072in}{1.424129in}}%
\pgfpathlineto{\pgfqpoint{0.955057in}{1.434288in}}%
\pgfpathlineto{\pgfqpoint{0.968314in}{1.431784in}}%
\pgfpathlineto{\pgfqpoint{0.969777in}{1.427234in}}%
\pgfpathlineto{\pgfqpoint{0.968279in}{1.422237in}}%
\pgfpathlineto{\pgfqpoint{0.964506in}{1.418586in}}%
\pgfpathlineto{\pgfqpoint{0.964652in}{1.410841in}}%
\pgfpathlineto{\pgfqpoint{0.967940in}{1.408953in}}%
\pgfpathlineto{\pgfqpoint{0.967728in}{1.403968in}}%
\pgfpathclose%
\pgfusepath{fill}%
\end{pgfscope}%
\begin{pgfscope}%
\pgfpathrectangle{\pgfqpoint{0.100000in}{0.100000in}}{\pgfqpoint{3.007045in}{1.925000in}}%
\pgfusepath{clip}%
\pgfsetbuttcap%
\pgfsetmiterjoin%
\definecolor{currentfill}{rgb}{0.381776,0.656517,0.824452}%
\pgfsetfillcolor{currentfill}%
\pgfsetlinewidth{0.000000pt}%
\definecolor{currentstroke}{rgb}{0.000000,0.000000,0.000000}%
\pgfsetstrokecolor{currentstroke}%
\pgfsetstrokeopacity{0.000000}%
\pgfsetdash{}{0pt}%
\pgfpathmoveto{\pgfqpoint{2.359204in}{1.116262in}}%
\pgfpathlineto{\pgfqpoint{2.353384in}{1.121802in}}%
\pgfpathlineto{\pgfqpoint{2.347373in}{1.122044in}}%
\pgfpathlineto{\pgfqpoint{2.339247in}{1.131363in}}%
\pgfpathlineto{\pgfqpoint{2.341607in}{1.134315in}}%
\pgfpathlineto{\pgfqpoint{2.338933in}{1.139975in}}%
\pgfpathlineto{\pgfqpoint{2.332965in}{1.140303in}}%
\pgfpathlineto{\pgfqpoint{2.326824in}{1.138484in}}%
\pgfpathlineto{\pgfqpoint{2.323384in}{1.154002in}}%
\pgfpathlineto{\pgfqpoint{2.331133in}{1.151744in}}%
\pgfpathlineto{\pgfqpoint{2.336776in}{1.153283in}}%
\pgfpathlineto{\pgfqpoint{2.342358in}{1.152328in}}%
\pgfpathlineto{\pgfqpoint{2.348520in}{1.145975in}}%
\pgfpathlineto{\pgfqpoint{2.354011in}{1.144847in}}%
\pgfpathlineto{\pgfqpoint{2.356171in}{1.148869in}}%
\pgfpathlineto{\pgfqpoint{2.360800in}{1.150692in}}%
\pgfpathlineto{\pgfqpoint{2.374431in}{1.146809in}}%
\pgfpathlineto{\pgfqpoint{2.380064in}{1.147699in}}%
\pgfpathlineto{\pgfqpoint{2.385587in}{1.155405in}}%
\pgfpathlineto{\pgfqpoint{2.384672in}{1.146612in}}%
\pgfpathlineto{\pgfqpoint{2.379735in}{1.139390in}}%
\pgfpathlineto{\pgfqpoint{2.377693in}{1.133822in}}%
\pgfpathlineto{\pgfqpoint{2.378692in}{1.128851in}}%
\pgfpathlineto{\pgfqpoint{2.372393in}{1.126861in}}%
\pgfpathlineto{\pgfqpoint{2.368680in}{1.131320in}}%
\pgfpathlineto{\pgfqpoint{2.364836in}{1.127691in}}%
\pgfpathlineto{\pgfqpoint{2.365037in}{1.122682in}}%
\pgfpathlineto{\pgfqpoint{2.359204in}{1.116262in}}%
\pgfpathclose%
\pgfusepath{fill}%
\end{pgfscope}%
\begin{pgfscope}%
\pgfpathrectangle{\pgfqpoint{0.100000in}{0.100000in}}{\pgfqpoint{3.007045in}{1.925000in}}%
\pgfusepath{clip}%
\pgfsetbuttcap%
\pgfsetmiterjoin%
\definecolor{currentfill}{rgb}{0.331334,0.622068,0.804767}%
\pgfsetfillcolor{currentfill}%
\pgfsetlinewidth{0.000000pt}%
\definecolor{currentstroke}{rgb}{0.000000,0.000000,0.000000}%
\pgfsetstrokecolor{currentstroke}%
\pgfsetstrokeopacity{0.000000}%
\pgfsetdash{}{0pt}%
\pgfpathmoveto{\pgfqpoint{2.382677in}{0.742228in}}%
\pgfpathlineto{\pgfqpoint{2.378977in}{0.740714in}}%
\pgfpathlineto{\pgfqpoint{2.375605in}{0.733287in}}%
\pgfpathlineto{\pgfqpoint{2.368826in}{0.730586in}}%
\pgfpathlineto{\pgfqpoint{2.363559in}{0.734763in}}%
\pgfpathlineto{\pgfqpoint{2.360665in}{0.742287in}}%
\pgfpathlineto{\pgfqpoint{2.367977in}{0.755519in}}%
\pgfpathlineto{\pgfqpoint{2.373071in}{0.752288in}}%
\pgfpathlineto{\pgfqpoint{2.375145in}{0.759957in}}%
\pgfpathlineto{\pgfqpoint{2.375546in}{0.768744in}}%
\pgfpathlineto{\pgfqpoint{2.379793in}{0.770306in}}%
\pgfpathlineto{\pgfqpoint{2.377909in}{0.787069in}}%
\pgfpathlineto{\pgfqpoint{2.387448in}{0.788188in}}%
\pgfpathlineto{\pgfqpoint{2.390153in}{0.787078in}}%
\pgfpathlineto{\pgfqpoint{2.390888in}{0.783936in}}%
\pgfpathlineto{\pgfqpoint{2.405437in}{0.788426in}}%
\pgfpathlineto{\pgfqpoint{2.411719in}{0.790134in}}%
\pgfpathlineto{\pgfqpoint{2.417835in}{0.773607in}}%
\pgfpathlineto{\pgfqpoint{2.415340in}{0.771423in}}%
\pgfpathlineto{\pgfqpoint{2.427990in}{0.752079in}}%
\pgfpathlineto{\pgfqpoint{2.434230in}{0.742118in}}%
\pgfpathlineto{\pgfqpoint{2.432076in}{0.743806in}}%
\pgfpathlineto{\pgfqpoint{2.425722in}{0.735586in}}%
\pgfpathlineto{\pgfqpoint{2.422656in}{0.729020in}}%
\pgfpathlineto{\pgfqpoint{2.427863in}{0.724959in}}%
\pgfpathlineto{\pgfqpoint{2.425917in}{0.719822in}}%
\pgfpathlineto{\pgfqpoint{2.410702in}{0.718714in}}%
\pgfpathlineto{\pgfqpoint{2.409404in}{0.729924in}}%
\pgfpathlineto{\pgfqpoint{2.395922in}{0.728401in}}%
\pgfpathlineto{\pgfqpoint{2.395370in}{0.738234in}}%
\pgfpathlineto{\pgfqpoint{2.392692in}{0.741268in}}%
\pgfpathlineto{\pgfqpoint{2.382677in}{0.742228in}}%
\pgfpathclose%
\pgfusepath{fill}%
\end{pgfscope}%
\begin{pgfscope}%
\pgfpathrectangle{\pgfqpoint{0.100000in}{0.100000in}}{\pgfqpoint{3.007045in}{1.925000in}}%
\pgfusepath{clip}%
\pgfsetbuttcap%
\pgfsetmiterjoin%
\definecolor{currentfill}{rgb}{0.270804,0.580730,0.781146}%
\pgfsetfillcolor{currentfill}%
\pgfsetlinewidth{0.000000pt}%
\definecolor{currentstroke}{rgb}{0.000000,0.000000,0.000000}%
\pgfsetstrokecolor{currentstroke}%
\pgfsetstrokeopacity{0.000000}%
\pgfsetdash{}{0pt}%
\pgfpathmoveto{\pgfqpoint{1.668319in}{1.540277in}}%
\pgfpathlineto{\pgfqpoint{1.633052in}{1.541386in}}%
\pgfpathlineto{\pgfqpoint{1.634891in}{1.593241in}}%
\pgfpathlineto{\pgfqpoint{1.669507in}{1.592201in}}%
\pgfpathlineto{\pgfqpoint{1.700191in}{1.591568in}}%
\pgfpathlineto{\pgfqpoint{1.699144in}{1.583940in}}%
\pgfpathlineto{\pgfqpoint{1.695859in}{1.579141in}}%
\pgfpathlineto{\pgfqpoint{1.687099in}{1.572723in}}%
\pgfpathlineto{\pgfqpoint{1.686157in}{1.569986in}}%
\pgfpathlineto{\pgfqpoint{1.694140in}{1.556870in}}%
\pgfpathlineto{\pgfqpoint{1.701484in}{1.554430in}}%
\pgfpathlineto{\pgfqpoint{1.703819in}{1.551137in}}%
\pgfpathlineto{\pgfqpoint{1.679482in}{1.551623in}}%
\pgfpathlineto{\pgfqpoint{1.678758in}{1.549649in}}%
\pgfpathlineto{\pgfqpoint{1.668554in}{1.549944in}}%
\pgfpathlineto{\pgfqpoint{1.668319in}{1.540277in}}%
\pgfpathclose%
\pgfusepath{fill}%
\end{pgfscope}%
\begin{pgfscope}%
\pgfpathrectangle{\pgfqpoint{0.100000in}{0.100000in}}{\pgfqpoint{3.007045in}{1.925000in}}%
\pgfusepath{clip}%
\pgfsetbuttcap%
\pgfsetmiterjoin%
\definecolor{currentfill}{rgb}{0.585882,0.773641,0.875079}%
\pgfsetfillcolor{currentfill}%
\pgfsetlinewidth{0.000000pt}%
\definecolor{currentstroke}{rgb}{0.000000,0.000000,0.000000}%
\pgfsetstrokecolor{currentstroke}%
\pgfsetstrokeopacity{0.000000}%
\pgfsetdash{}{0pt}%
\pgfpathmoveto{\pgfqpoint{2.493292in}{1.298554in}}%
\pgfpathlineto{\pgfqpoint{2.486521in}{1.293715in}}%
\pgfpathlineto{\pgfqpoint{2.476757in}{1.293301in}}%
\pgfpathlineto{\pgfqpoint{2.473535in}{1.295735in}}%
\pgfpathlineto{\pgfqpoint{2.472733in}{1.301176in}}%
\pgfpathlineto{\pgfqpoint{2.464330in}{1.299934in}}%
\pgfpathlineto{\pgfqpoint{2.461738in}{1.316890in}}%
\pgfpathlineto{\pgfqpoint{2.466018in}{1.317526in}}%
\pgfpathlineto{\pgfqpoint{2.460651in}{1.350974in}}%
\pgfpathlineto{\pgfqpoint{2.484276in}{1.354588in}}%
\pgfpathlineto{\pgfqpoint{2.503675in}{1.356445in}}%
\pgfpathlineto{\pgfqpoint{2.509694in}{1.358087in}}%
\pgfpathlineto{\pgfqpoint{2.513140in}{1.337398in}}%
\pgfpathlineto{\pgfqpoint{2.506874in}{1.324911in}}%
\pgfpathlineto{\pgfqpoint{2.508763in}{1.315516in}}%
\pgfpathlineto{\pgfqpoint{2.491069in}{1.312382in}}%
\pgfpathlineto{\pgfqpoint{2.493292in}{1.298554in}}%
\pgfpathclose%
\pgfusepath{fill}%
\end{pgfscope}%
\begin{pgfscope}%
\pgfpathrectangle{\pgfqpoint{0.100000in}{0.100000in}}{\pgfqpoint{3.007045in}{1.925000in}}%
\pgfusepath{clip}%
\pgfsetbuttcap%
\pgfsetmiterjoin%
\definecolor{currentfill}{rgb}{0.479216,0.715079,0.852072}%
\pgfsetfillcolor{currentfill}%
\pgfsetlinewidth{0.000000pt}%
\definecolor{currentstroke}{rgb}{0.000000,0.000000,0.000000}%
\pgfsetstrokecolor{currentstroke}%
\pgfsetstrokeopacity{0.000000}%
\pgfsetdash{}{0pt}%
\pgfpathmoveto{\pgfqpoint{1.386507in}{1.413753in}}%
\pgfpathlineto{\pgfqpoint{1.335923in}{1.418503in}}%
\pgfpathlineto{\pgfqpoint{1.339144in}{1.451533in}}%
\pgfpathlineto{\pgfqpoint{1.343611in}{1.496079in}}%
\pgfpathlineto{\pgfqpoint{1.348840in}{1.549898in}}%
\pgfpathlineto{\pgfqpoint{1.351084in}{1.564019in}}%
\pgfpathlineto{\pgfqpoint{1.401374in}{1.559284in}}%
\pgfpathlineto{\pgfqpoint{1.400379in}{1.547840in}}%
\pgfpathlineto{\pgfqpoint{1.445034in}{1.544102in}}%
\pgfpathlineto{\pgfqpoint{1.442332in}{1.509325in}}%
\pgfpathlineto{\pgfqpoint{1.439374in}{1.475257in}}%
\pgfpathlineto{\pgfqpoint{1.437357in}{1.456409in}}%
\pgfpathlineto{\pgfqpoint{1.431599in}{1.456389in}}%
\pgfpathlineto{\pgfqpoint{1.429680in}{1.455658in}}%
\pgfpathlineto{\pgfqpoint{1.399572in}{1.458189in}}%
\pgfpathlineto{\pgfqpoint{1.395146in}{1.457695in}}%
\pgfpathlineto{\pgfqpoint{1.390013in}{1.453572in}}%
\pgfpathlineto{\pgfqpoint{1.386507in}{1.413753in}}%
\pgfpathclose%
\pgfusepath{fill}%
\end{pgfscope}%
\begin{pgfscope}%
\pgfpathrectangle{\pgfqpoint{0.100000in}{0.100000in}}{\pgfqpoint{3.007045in}{1.925000in}}%
\pgfusepath{clip}%
\pgfsetbuttcap%
\pgfsetmiterjoin%
\definecolor{currentfill}{rgb}{0.435294,0.690965,0.842599}%
\pgfsetfillcolor{currentfill}%
\pgfsetlinewidth{0.000000pt}%
\definecolor{currentstroke}{rgb}{0.000000,0.000000,0.000000}%
\pgfsetstrokecolor{currentstroke}%
\pgfsetstrokeopacity{0.000000}%
\pgfsetdash{}{0pt}%
\pgfpathmoveto{\pgfqpoint{2.406434in}{1.016470in}}%
\pgfpathlineto{\pgfqpoint{2.399781in}{1.010130in}}%
\pgfpathlineto{\pgfqpoint{2.396925in}{1.003949in}}%
\pgfpathlineto{\pgfqpoint{2.392615in}{1.001420in}}%
\pgfpathlineto{\pgfqpoint{2.388640in}{1.002575in}}%
\pgfpathlineto{\pgfqpoint{2.380497in}{0.997923in}}%
\pgfpathlineto{\pgfqpoint{2.372662in}{0.995532in}}%
\pgfpathlineto{\pgfqpoint{2.361719in}{1.000172in}}%
\pgfpathlineto{\pgfqpoint{2.358542in}{0.999323in}}%
\pgfpathlineto{\pgfqpoint{2.353818in}{1.009495in}}%
\pgfpathlineto{\pgfqpoint{2.355024in}{1.013333in}}%
\pgfpathlineto{\pgfqpoint{2.360603in}{1.018528in}}%
\pgfpathlineto{\pgfqpoint{2.362675in}{1.024583in}}%
\pgfpathlineto{\pgfqpoint{2.376071in}{1.035433in}}%
\pgfpathlineto{\pgfqpoint{2.378235in}{1.029769in}}%
\pgfpathlineto{\pgfqpoint{2.378204in}{1.022084in}}%
\pgfpathlineto{\pgfqpoint{2.380873in}{1.017870in}}%
\pgfpathlineto{\pgfqpoint{2.373385in}{1.015073in}}%
\pgfpathlineto{\pgfqpoint{2.372965in}{1.012422in}}%
\pgfpathlineto{\pgfqpoint{2.406434in}{1.016470in}}%
\pgfpathclose%
\pgfusepath{fill}%
\end{pgfscope}%
\begin{pgfscope}%
\pgfpathrectangle{\pgfqpoint{0.100000in}{0.100000in}}{\pgfqpoint{3.007045in}{1.925000in}}%
\pgfusepath{clip}%
\pgfsetbuttcap%
\pgfsetmiterjoin%
\definecolor{currentfill}{rgb}{0.784591,0.864237,0.939962}%
\pgfsetfillcolor{currentfill}%
\pgfsetlinewidth{0.000000pt}%
\definecolor{currentstroke}{rgb}{0.000000,0.000000,0.000000}%
\pgfsetstrokecolor{currentstroke}%
\pgfsetstrokeopacity{0.000000}%
\pgfsetdash{}{0pt}%
\pgfpathmoveto{\pgfqpoint{2.385967in}{0.896018in}}%
\pgfpathlineto{\pgfqpoint{2.377707in}{0.892532in}}%
\pgfpathlineto{\pgfqpoint{2.375790in}{0.888043in}}%
\pgfpathlineto{\pgfqpoint{2.371189in}{0.888722in}}%
\pgfpathlineto{\pgfqpoint{2.366727in}{0.881594in}}%
\pgfpathlineto{\pgfqpoint{2.360142in}{0.881259in}}%
\pgfpathlineto{\pgfqpoint{2.362404in}{0.891696in}}%
\pgfpathlineto{\pgfqpoint{2.356919in}{0.901553in}}%
\pgfpathlineto{\pgfqpoint{2.348646in}{0.902478in}}%
\pgfpathlineto{\pgfqpoint{2.348504in}{0.918229in}}%
\pgfpathlineto{\pgfqpoint{2.351726in}{0.921503in}}%
\pgfpathlineto{\pgfqpoint{2.356471in}{0.920502in}}%
\pgfpathlineto{\pgfqpoint{2.361956in}{0.924269in}}%
\pgfpathlineto{\pgfqpoint{2.366318in}{0.919795in}}%
\pgfpathlineto{\pgfqpoint{2.372864in}{0.923192in}}%
\pgfpathlineto{\pgfqpoint{2.379773in}{0.923563in}}%
\pgfpathlineto{\pgfqpoint{2.378628in}{0.917255in}}%
\pgfpathlineto{\pgfqpoint{2.383505in}{0.917672in}}%
\pgfpathlineto{\pgfqpoint{2.391763in}{0.908652in}}%
\pgfpathlineto{\pgfqpoint{2.385392in}{0.899513in}}%
\pgfpathlineto{\pgfqpoint{2.385967in}{0.896018in}}%
\pgfpathclose%
\pgfusepath{fill}%
\end{pgfscope}%
\begin{pgfscope}%
\pgfpathrectangle{\pgfqpoint{0.100000in}{0.100000in}}{\pgfqpoint{3.007045in}{1.925000in}}%
\pgfusepath{clip}%
\pgfsetbuttcap%
\pgfsetmiterjoin%
\definecolor{currentfill}{rgb}{0.548235,0.752972,0.866959}%
\pgfsetfillcolor{currentfill}%
\pgfsetlinewidth{0.000000pt}%
\definecolor{currentstroke}{rgb}{0.000000,0.000000,0.000000}%
\pgfsetstrokecolor{currentstroke}%
\pgfsetstrokeopacity{0.000000}%
\pgfsetdash{}{0pt}%
\pgfpathmoveto{\pgfqpoint{2.410399in}{1.309530in}}%
\pgfpathlineto{\pgfqpoint{2.407666in}{1.313899in}}%
\pgfpathlineto{\pgfqpoint{2.407064in}{1.318707in}}%
\pgfpathlineto{\pgfqpoint{2.411263in}{1.323487in}}%
\pgfpathlineto{\pgfqpoint{2.416075in}{1.324178in}}%
\pgfpathlineto{\pgfqpoint{2.415460in}{1.329120in}}%
\pgfpathlineto{\pgfqpoint{2.420098in}{1.329773in}}%
\pgfpathlineto{\pgfqpoint{2.419430in}{1.334702in}}%
\pgfpathlineto{\pgfqpoint{2.414804in}{1.334085in}}%
\pgfpathlineto{\pgfqpoint{2.413516in}{1.344212in}}%
\pgfpathlineto{\pgfqpoint{2.424873in}{1.344692in}}%
\pgfpathlineto{\pgfqpoint{2.430966in}{1.350397in}}%
\pgfpathlineto{\pgfqpoint{2.437543in}{1.358914in}}%
\pgfpathlineto{\pgfqpoint{2.456731in}{1.373827in}}%
\pgfpathlineto{\pgfqpoint{2.461753in}{1.375576in}}%
\pgfpathlineto{\pgfqpoint{2.479266in}{1.385678in}}%
\pgfpathlineto{\pgfqpoint{2.484276in}{1.354588in}}%
\pgfpathlineto{\pgfqpoint{2.460651in}{1.350974in}}%
\pgfpathlineto{\pgfqpoint{2.466018in}{1.317526in}}%
\pgfpathlineto{\pgfqpoint{2.461738in}{1.316890in}}%
\pgfpathlineto{\pgfqpoint{2.445530in}{1.314439in}}%
\pgfpathlineto{\pgfqpoint{2.446146in}{1.309080in}}%
\pgfpathlineto{\pgfqpoint{2.434855in}{1.307908in}}%
\pgfpathlineto{\pgfqpoint{2.434169in}{1.312710in}}%
\pgfpathlineto{\pgfqpoint{2.410399in}{1.309530in}}%
\pgfpathclose%
\pgfusepath{fill}%
\end{pgfscope}%
\begin{pgfscope}%
\pgfpathrectangle{\pgfqpoint{0.100000in}{0.100000in}}{\pgfqpoint{3.007045in}{1.925000in}}%
\pgfusepath{clip}%
\pgfsetbuttcap%
\pgfsetmiterjoin%
\definecolor{currentfill}{rgb}{0.195386,0.509112,0.743791}%
\pgfsetfillcolor{currentfill}%
\pgfsetlinewidth{0.000000pt}%
\definecolor{currentstroke}{rgb}{0.000000,0.000000,0.000000}%
\pgfsetstrokecolor{currentstroke}%
\pgfsetstrokeopacity{0.000000}%
\pgfsetdash{}{0pt}%
\pgfpathmoveto{\pgfqpoint{0.630991in}{1.684466in}}%
\pgfpathlineto{\pgfqpoint{0.625624in}{1.685979in}}%
\pgfpathlineto{\pgfqpoint{0.630102in}{1.696622in}}%
\pgfpathlineto{\pgfqpoint{0.624611in}{1.698303in}}%
\pgfpathlineto{\pgfqpoint{0.626202in}{1.703795in}}%
\pgfpathlineto{\pgfqpoint{0.620630in}{1.705467in}}%
\pgfpathlineto{\pgfqpoint{0.630951in}{1.740790in}}%
\pgfpathlineto{\pgfqpoint{0.637219in}{1.740644in}}%
\pgfpathlineto{\pgfqpoint{0.640972in}{1.753699in}}%
\pgfpathlineto{\pgfqpoint{0.651182in}{1.791391in}}%
\pgfpathlineto{\pgfqpoint{0.662371in}{1.789397in}}%
\pgfpathlineto{\pgfqpoint{0.668689in}{1.793331in}}%
\pgfpathlineto{\pgfqpoint{0.670581in}{1.789487in}}%
\pgfpathlineto{\pgfqpoint{0.675279in}{1.792196in}}%
\pgfpathlineto{\pgfqpoint{0.702590in}{1.784664in}}%
\pgfpathlineto{\pgfqpoint{0.726344in}{1.778628in}}%
\pgfpathlineto{\pgfqpoint{0.724517in}{1.778735in}}%
\pgfpathlineto{\pgfqpoint{0.724939in}{1.766494in}}%
\pgfpathlineto{\pgfqpoint{0.740246in}{1.767188in}}%
\pgfpathlineto{\pgfqpoint{0.737197in}{1.757418in}}%
\pgfpathlineto{\pgfqpoint{0.742787in}{1.756014in}}%
\pgfpathlineto{\pgfqpoint{0.743350in}{1.746702in}}%
\pgfpathlineto{\pgfqpoint{0.746177in}{1.746042in}}%
\pgfpathlineto{\pgfqpoint{0.741011in}{1.724091in}}%
\pgfpathlineto{\pgfqpoint{0.724313in}{1.728496in}}%
\pgfpathlineto{\pgfqpoint{0.722204in}{1.719436in}}%
\pgfpathlineto{\pgfqpoint{0.718256in}{1.717477in}}%
\pgfpathlineto{\pgfqpoint{0.715135in}{1.709465in}}%
\pgfpathlineto{\pgfqpoint{0.712881in}{1.710060in}}%
\pgfpathlineto{\pgfqpoint{0.709253in}{1.696086in}}%
\pgfpathlineto{\pgfqpoint{0.704883in}{1.694237in}}%
\pgfpathlineto{\pgfqpoint{0.694542in}{1.697073in}}%
\pgfpathlineto{\pgfqpoint{0.693304in}{1.692341in}}%
\pgfpathlineto{\pgfqpoint{0.681006in}{1.694917in}}%
\pgfpathlineto{\pgfqpoint{0.680362in}{1.684919in}}%
\pgfpathlineto{\pgfqpoint{0.685286in}{1.683617in}}%
\pgfpathlineto{\pgfqpoint{0.683125in}{1.670416in}}%
\pgfpathlineto{\pgfqpoint{0.630991in}{1.684466in}}%
\pgfpathclose%
\pgfusepath{fill}%
\end{pgfscope}%
\begin{pgfscope}%
\pgfpathrectangle{\pgfqpoint{0.100000in}{0.100000in}}{\pgfqpoint{3.007045in}{1.925000in}}%
\pgfusepath{clip}%
\pgfsetbuttcap%
\pgfsetmiterjoin%
\definecolor{currentfill}{rgb}{0.627605,0.795556,0.885152}%
\pgfsetfillcolor{currentfill}%
\pgfsetlinewidth{0.000000pt}%
\definecolor{currentstroke}{rgb}{0.000000,0.000000,0.000000}%
\pgfsetstrokecolor{currentstroke}%
\pgfsetstrokeopacity{0.000000}%
\pgfsetdash{}{0pt}%
\pgfpathmoveto{\pgfqpoint{1.351564in}{1.062383in}}%
\pgfpathlineto{\pgfqpoint{1.347270in}{1.020215in}}%
\pgfpathlineto{\pgfqpoint{1.298942in}{1.024507in}}%
\pgfpathlineto{\pgfqpoint{1.277959in}{1.026457in}}%
\pgfpathlineto{\pgfqpoint{1.238883in}{1.030797in}}%
\pgfpathlineto{\pgfqpoint{1.241118in}{1.050301in}}%
\pgfpathlineto{\pgfqpoint{1.242720in}{1.057254in}}%
\pgfpathlineto{\pgfqpoint{1.241493in}{1.067690in}}%
\pgfpathlineto{\pgfqpoint{1.236457in}{1.074727in}}%
\pgfpathlineto{\pgfqpoint{1.225618in}{1.073383in}}%
\pgfpathlineto{\pgfqpoint{1.231022in}{1.088493in}}%
\pgfpathlineto{\pgfqpoint{1.229068in}{1.091653in}}%
\pgfpathlineto{\pgfqpoint{1.232103in}{1.090942in}}%
\pgfpathlineto{\pgfqpoint{1.241682in}{1.094643in}}%
\pgfpathlineto{\pgfqpoint{1.245702in}{1.097910in}}%
\pgfpathlineto{\pgfqpoint{1.251074in}{1.090450in}}%
\pgfpathlineto{\pgfqpoint{1.252696in}{1.088003in}}%
\pgfpathlineto{\pgfqpoint{1.271819in}{1.087203in}}%
\pgfpathlineto{\pgfqpoint{1.301086in}{1.073162in}}%
\pgfpathlineto{\pgfqpoint{1.300552in}{1.067248in}}%
\pgfpathlineto{\pgfqpoint{1.351564in}{1.062383in}}%
\pgfpathclose%
\pgfusepath{fill}%
\end{pgfscope}%
\begin{pgfscope}%
\pgfpathrectangle{\pgfqpoint{0.100000in}{0.100000in}}{\pgfqpoint{3.007045in}{1.925000in}}%
\pgfusepath{clip}%
\pgfsetbuttcap%
\pgfsetmiterjoin%
\definecolor{currentfill}{rgb}{0.429020,0.687520,0.841246}%
\pgfsetfillcolor{currentfill}%
\pgfsetlinewidth{0.000000pt}%
\definecolor{currentstroke}{rgb}{0.000000,0.000000,0.000000}%
\pgfsetstrokecolor{currentstroke}%
\pgfsetstrokeopacity{0.000000}%
\pgfsetdash{}{0pt}%
\pgfpathmoveto{\pgfqpoint{1.972796in}{1.051471in}}%
\pgfpathlineto{\pgfqpoint{1.964595in}{1.051309in}}%
\pgfpathlineto{\pgfqpoint{1.964484in}{1.057050in}}%
\pgfpathlineto{\pgfqpoint{1.953009in}{1.056717in}}%
\pgfpathlineto{\pgfqpoint{1.952290in}{1.080650in}}%
\pgfpathlineto{\pgfqpoint{1.946082in}{1.080499in}}%
\pgfpathlineto{\pgfqpoint{1.945763in}{1.089375in}}%
\pgfpathlineto{\pgfqpoint{1.944978in}{1.116719in}}%
\pgfpathlineto{\pgfqpoint{1.949655in}{1.115080in}}%
\pgfpathlineto{\pgfqpoint{1.956351in}{1.117548in}}%
\pgfpathlineto{\pgfqpoint{1.956143in}{1.126474in}}%
\pgfpathlineto{\pgfqpoint{1.963877in}{1.126688in}}%
\pgfpathlineto{\pgfqpoint{1.963320in}{1.146179in}}%
\pgfpathlineto{\pgfqpoint{1.967135in}{1.146275in}}%
\pgfpathlineto{\pgfqpoint{1.967023in}{1.152094in}}%
\pgfpathlineto{\pgfqpoint{1.990535in}{1.152936in}}%
\pgfpathlineto{\pgfqpoint{1.992952in}{1.144098in}}%
\pgfpathlineto{\pgfqpoint{1.991482in}{1.141762in}}%
\pgfpathlineto{\pgfqpoint{1.994783in}{1.133102in}}%
\pgfpathlineto{\pgfqpoint{2.000724in}{1.130408in}}%
\pgfpathlineto{\pgfqpoint{2.005193in}{1.136744in}}%
\pgfpathlineto{\pgfqpoint{2.015419in}{1.134149in}}%
\pgfpathlineto{\pgfqpoint{2.022871in}{1.129850in}}%
\pgfpathlineto{\pgfqpoint{2.018572in}{1.121893in}}%
\pgfpathlineto{\pgfqpoint{2.020049in}{1.117235in}}%
\pgfpathlineto{\pgfqpoint{2.044424in}{1.118250in}}%
\pgfpathlineto{\pgfqpoint{2.046290in}{1.089601in}}%
\pgfpathlineto{\pgfqpoint{2.029064in}{1.088879in}}%
\pgfpathlineto{\pgfqpoint{2.029399in}{1.083124in}}%
\pgfpathlineto{\pgfqpoint{2.020734in}{1.079515in}}%
\pgfpathlineto{\pgfqpoint{2.014367in}{1.079611in}}%
\pgfpathlineto{\pgfqpoint{2.010046in}{1.075951in}}%
\pgfpathlineto{\pgfqpoint{2.001501in}{1.072967in}}%
\pgfpathlineto{\pgfqpoint{1.990680in}{1.085683in}}%
\pgfpathlineto{\pgfqpoint{1.974361in}{1.084952in}}%
\pgfpathlineto{\pgfqpoint{1.975453in}{1.054398in}}%
\pgfpathlineto{\pgfqpoint{1.972796in}{1.051471in}}%
\pgfpathclose%
\pgfusepath{fill}%
\end{pgfscope}%
\begin{pgfscope}%
\pgfpathrectangle{\pgfqpoint{0.100000in}{0.100000in}}{\pgfqpoint{3.007045in}{1.925000in}}%
\pgfusepath{clip}%
\pgfsetbuttcap%
\pgfsetmiterjoin%
\definecolor{currentfill}{rgb}{0.311157,0.608289,0.796894}%
\pgfsetfillcolor{currentfill}%
\pgfsetlinewidth{0.000000pt}%
\definecolor{currentstroke}{rgb}{0.000000,0.000000,0.000000}%
\pgfsetstrokecolor{currentstroke}%
\pgfsetstrokeopacity{0.000000}%
\pgfsetdash{}{0pt}%
\pgfpathmoveto{\pgfqpoint{1.792356in}{1.069644in}}%
\pgfpathlineto{\pgfqpoint{1.768296in}{1.069654in}}%
\pgfpathlineto{\pgfqpoint{1.768960in}{1.084214in}}%
\pgfpathlineto{\pgfqpoint{1.768947in}{1.112944in}}%
\pgfpathlineto{\pgfqpoint{1.769426in}{1.115817in}}%
\pgfpathlineto{\pgfqpoint{1.769446in}{1.131864in}}%
\pgfpathlineto{\pgfqpoint{1.776978in}{1.132450in}}%
\pgfpathlineto{\pgfqpoint{1.777411in}{1.146354in}}%
\pgfpathlineto{\pgfqpoint{1.780164in}{1.150731in}}%
\pgfpathlineto{\pgfqpoint{1.772963in}{1.162326in}}%
\pgfpathlineto{\pgfqpoint{1.767101in}{1.168173in}}%
\pgfpathlineto{\pgfqpoint{1.792583in}{1.168136in}}%
\pgfpathlineto{\pgfqpoint{1.792627in}{1.163111in}}%
\pgfpathlineto{\pgfqpoint{1.812471in}{1.163163in}}%
\pgfpathlineto{\pgfqpoint{1.812537in}{1.167939in}}%
\pgfpathlineto{\pgfqpoint{1.835366in}{1.168017in}}%
\pgfpathlineto{\pgfqpoint{1.835689in}{1.147108in}}%
\pgfpathlineto{\pgfqpoint{1.840824in}{1.149749in}}%
\pgfpathlineto{\pgfqpoint{1.848856in}{1.148030in}}%
\pgfpathlineto{\pgfqpoint{1.849349in}{1.129915in}}%
\pgfpathlineto{\pgfqpoint{1.849007in}{1.104433in}}%
\pgfpathlineto{\pgfqpoint{1.820514in}{1.104792in}}%
\pgfpathlineto{\pgfqpoint{1.820190in}{1.081544in}}%
\pgfpathlineto{\pgfqpoint{1.821187in}{1.069825in}}%
\pgfpathlineto{\pgfqpoint{1.811752in}{1.070677in}}%
\pgfpathlineto{\pgfqpoint{1.792345in}{1.071169in}}%
\pgfpathlineto{\pgfqpoint{1.792356in}{1.069644in}}%
\pgfpathclose%
\pgfusepath{fill}%
\end{pgfscope}%
\begin{pgfscope}%
\pgfpathrectangle{\pgfqpoint{0.100000in}{0.100000in}}{\pgfqpoint{3.007045in}{1.925000in}}%
\pgfusepath{clip}%
\pgfsetbuttcap%
\pgfsetmiterjoin%
\definecolor{currentfill}{rgb}{0.183206,0.496932,0.737516}%
\pgfsetfillcolor{currentfill}%
\pgfsetlinewidth{0.000000pt}%
\definecolor{currentstroke}{rgb}{0.000000,0.000000,0.000000}%
\pgfsetstrokecolor{currentstroke}%
\pgfsetstrokeopacity{0.000000}%
\pgfsetdash{}{0pt}%
\pgfpathmoveto{\pgfqpoint{1.285700in}{1.824835in}}%
\pgfpathlineto{\pgfqpoint{1.331685in}{1.819464in}}%
\pgfpathlineto{\pgfqpoint{1.331869in}{1.813477in}}%
\pgfpathlineto{\pgfqpoint{1.329938in}{1.796177in}}%
\pgfpathlineto{\pgfqpoint{1.332164in}{1.790080in}}%
\pgfpathlineto{\pgfqpoint{1.295653in}{1.794279in}}%
\pgfpathlineto{\pgfqpoint{1.286131in}{1.795421in}}%
\pgfpathlineto{\pgfqpoint{1.288237in}{1.812721in}}%
\pgfpathlineto{\pgfqpoint{1.282530in}{1.813419in}}%
\pgfpathlineto{\pgfqpoint{1.285700in}{1.824835in}}%
\pgfpathclose%
\pgfusepath{fill}%
\end{pgfscope}%
\begin{pgfscope}%
\pgfpathrectangle{\pgfqpoint{0.100000in}{0.100000in}}{\pgfqpoint{3.007045in}{1.925000in}}%
\pgfusepath{clip}%
\pgfsetbuttcap%
\pgfsetmiterjoin%
\definecolor{currentfill}{rgb}{0.401953,0.670296,0.832326}%
\pgfsetfillcolor{currentfill}%
\pgfsetlinewidth{0.000000pt}%
\definecolor{currentstroke}{rgb}{0.000000,0.000000,0.000000}%
\pgfsetstrokecolor{currentstroke}%
\pgfsetstrokeopacity{0.000000}%
\pgfsetdash{}{0pt}%
\pgfpathmoveto{\pgfqpoint{2.568608in}{0.576327in}}%
\pgfpathlineto{\pgfqpoint{2.562689in}{0.572385in}}%
\pgfpathlineto{\pgfqpoint{2.551269in}{0.570455in}}%
\pgfpathlineto{\pgfqpoint{2.549830in}{0.579481in}}%
\pgfpathlineto{\pgfqpoint{2.545862in}{0.584302in}}%
\pgfpathlineto{\pgfqpoint{2.532632in}{0.582139in}}%
\pgfpathlineto{\pgfqpoint{2.526180in}{0.575201in}}%
\pgfpathlineto{\pgfqpoint{2.520191in}{0.572396in}}%
\pgfpathlineto{\pgfqpoint{2.516161in}{0.600413in}}%
\pgfpathlineto{\pgfqpoint{2.492884in}{0.596677in}}%
\pgfpathlineto{\pgfqpoint{2.488716in}{0.626214in}}%
\pgfpathlineto{\pgfqpoint{2.502750in}{0.627127in}}%
\pgfpathlineto{\pgfqpoint{2.504409in}{0.641856in}}%
\pgfpathlineto{\pgfqpoint{2.488999in}{0.640420in}}%
\pgfpathlineto{\pgfqpoint{2.487249in}{0.654855in}}%
\pgfpathlineto{\pgfqpoint{2.503300in}{0.656930in}}%
\pgfpathlineto{\pgfqpoint{2.506695in}{0.661778in}}%
\pgfpathlineto{\pgfqpoint{2.514003in}{0.660961in}}%
\pgfpathlineto{\pgfqpoint{2.521337in}{0.669031in}}%
\pgfpathlineto{\pgfqpoint{2.533990in}{0.667284in}}%
\pgfpathlineto{\pgfqpoint{2.542300in}{0.662555in}}%
\pgfpathlineto{\pgfqpoint{2.545263in}{0.656150in}}%
\pgfpathlineto{\pgfqpoint{2.543774in}{0.645180in}}%
\pgfpathlineto{\pgfqpoint{2.546129in}{0.642161in}}%
\pgfpathlineto{\pgfqpoint{2.546297in}{0.634841in}}%
\pgfpathlineto{\pgfqpoint{2.551061in}{0.622627in}}%
\pgfpathlineto{\pgfqpoint{2.552074in}{0.616494in}}%
\pgfpathlineto{\pgfqpoint{2.560429in}{0.594452in}}%
\pgfpathlineto{\pgfqpoint{2.568608in}{0.576327in}}%
\pgfpathclose%
\pgfusepath{fill}%
\end{pgfscope}%
\begin{pgfscope}%
\pgfpathrectangle{\pgfqpoint{0.100000in}{0.100000in}}{\pgfqpoint{3.007045in}{1.925000in}}%
\pgfusepath{clip}%
\pgfsetbuttcap%
\pgfsetmiterjoin%
\definecolor{currentfill}{rgb}{0.311157,0.608289,0.796894}%
\pgfsetfillcolor{currentfill}%
\pgfsetlinewidth{0.000000pt}%
\definecolor{currentstroke}{rgb}{0.000000,0.000000,0.000000}%
\pgfsetstrokecolor{currentstroke}%
\pgfsetstrokeopacity{0.000000}%
\pgfsetdash{}{0pt}%
\pgfpathmoveto{\pgfqpoint{1.342266in}{0.937180in}}%
\pgfpathlineto{\pgfqpoint{1.341578in}{0.929505in}}%
\pgfpathlineto{\pgfqpoint{1.337618in}{0.885562in}}%
\pgfpathlineto{\pgfqpoint{1.324670in}{0.886744in}}%
\pgfpathlineto{\pgfqpoint{1.324074in}{0.880866in}}%
\pgfpathlineto{\pgfqpoint{1.318390in}{0.881414in}}%
\pgfpathlineto{\pgfqpoint{1.317814in}{0.875699in}}%
\pgfpathlineto{\pgfqpoint{1.312137in}{0.876255in}}%
\pgfpathlineto{\pgfqpoint{1.311575in}{0.870525in}}%
\pgfpathlineto{\pgfqpoint{1.300241in}{0.871662in}}%
\pgfpathlineto{\pgfqpoint{1.299641in}{0.865888in}}%
\pgfpathlineto{\pgfqpoint{1.286524in}{0.867195in}}%
\pgfpathlineto{\pgfqpoint{1.276783in}{0.868180in}}%
\pgfpathlineto{\pgfqpoint{1.280556in}{0.903442in}}%
\pgfpathlineto{\pgfqpoint{1.295559in}{0.908588in}}%
\pgfpathlineto{\pgfqpoint{1.307382in}{0.907359in}}%
\pgfpathlineto{\pgfqpoint{1.308509in}{0.919262in}}%
\pgfpathlineto{\pgfqpoint{1.303832in}{0.919724in}}%
\pgfpathlineto{\pgfqpoint{1.292797in}{0.945962in}}%
\pgfpathlineto{\pgfqpoint{1.290299in}{0.944870in}}%
\pgfpathlineto{\pgfqpoint{1.271875in}{0.946736in}}%
\pgfpathlineto{\pgfqpoint{1.273425in}{0.954257in}}%
\pgfpathlineto{\pgfqpoint{1.273752in}{0.963960in}}%
\pgfpathlineto{\pgfqpoint{1.271190in}{0.975810in}}%
\pgfpathlineto{\pgfqpoint{1.293808in}{0.973457in}}%
\pgfpathlineto{\pgfqpoint{1.293521in}{0.970602in}}%
\pgfpathlineto{\pgfqpoint{1.304909in}{0.969450in}}%
\pgfpathlineto{\pgfqpoint{1.304327in}{0.963661in}}%
\pgfpathlineto{\pgfqpoint{1.327114in}{0.961542in}}%
\pgfpathlineto{\pgfqpoint{1.324480in}{0.938840in}}%
\pgfpathlineto{\pgfqpoint{1.342266in}{0.937180in}}%
\pgfpathclose%
\pgfusepath{fill}%
\end{pgfscope}%
\begin{pgfscope}%
\pgfpathrectangle{\pgfqpoint{0.100000in}{0.100000in}}{\pgfqpoint{3.007045in}{1.925000in}}%
\pgfusepath{clip}%
\pgfsetbuttcap%
\pgfsetmiterjoin%
\definecolor{currentfill}{rgb}{0.711265,0.831111,0.914433}%
\pgfsetfillcolor{currentfill}%
\pgfsetlinewidth{0.000000pt}%
\definecolor{currentstroke}{rgb}{0.000000,0.000000,0.000000}%
\pgfsetstrokecolor{currentstroke}%
\pgfsetstrokeopacity{0.000000}%
\pgfsetdash{}{0pt}%
\pgfpathmoveto{\pgfqpoint{2.625680in}{1.382482in}}%
\pgfpathlineto{\pgfqpoint{2.619453in}{1.411977in}}%
\pgfpathlineto{\pgfqpoint{2.612717in}{1.410639in}}%
\pgfpathlineto{\pgfqpoint{2.608126in}{1.441492in}}%
\pgfpathlineto{\pgfqpoint{2.609773in}{1.449142in}}%
\pgfpathlineto{\pgfqpoint{2.634427in}{1.453744in}}%
\pgfpathlineto{\pgfqpoint{2.635144in}{1.448720in}}%
\pgfpathlineto{\pgfqpoint{2.647379in}{1.448629in}}%
\pgfpathlineto{\pgfqpoint{2.646021in}{1.453686in}}%
\pgfpathlineto{\pgfqpoint{2.661352in}{1.457374in}}%
\pgfpathlineto{\pgfqpoint{2.656543in}{1.461343in}}%
\pgfpathlineto{\pgfqpoint{2.675608in}{1.465295in}}%
\pgfpathlineto{\pgfqpoint{2.679109in}{1.451464in}}%
\pgfpathlineto{\pgfqpoint{2.680763in}{1.444336in}}%
\pgfpathlineto{\pgfqpoint{2.672580in}{1.444093in}}%
\pgfpathlineto{\pgfqpoint{2.667190in}{1.440512in}}%
\pgfpathlineto{\pgfqpoint{2.669988in}{1.422124in}}%
\pgfpathlineto{\pgfqpoint{2.652249in}{1.418613in}}%
\pgfpathlineto{\pgfqpoint{2.660061in}{1.392938in}}%
\pgfpathlineto{\pgfqpoint{2.656505in}{1.389182in}}%
\pgfpathlineto{\pgfqpoint{2.625680in}{1.382482in}}%
\pgfpathclose%
\pgfusepath{fill}%
\end{pgfscope}%
\begin{pgfscope}%
\pgfpathrectangle{\pgfqpoint{0.100000in}{0.100000in}}{\pgfqpoint{3.007045in}{1.925000in}}%
\pgfusepath{clip}%
\pgfsetbuttcap%
\pgfsetmiterjoin%
\definecolor{currentfill}{rgb}{0.346467,0.632403,0.810673}%
\pgfsetfillcolor{currentfill}%
\pgfsetlinewidth{0.000000pt}%
\definecolor{currentstroke}{rgb}{0.000000,0.000000,0.000000}%
\pgfsetstrokecolor{currentstroke}%
\pgfsetstrokeopacity{0.000000}%
\pgfsetdash{}{0pt}%
\pgfpathmoveto{\pgfqpoint{0.683738in}{1.619619in}}%
\pgfpathlineto{\pgfqpoint{0.687625in}{1.619603in}}%
\pgfpathlineto{\pgfqpoint{0.695577in}{1.627414in}}%
\pgfpathlineto{\pgfqpoint{0.698781in}{1.628508in}}%
\pgfpathlineto{\pgfqpoint{0.716216in}{1.624121in}}%
\pgfpathlineto{\pgfqpoint{0.720108in}{1.619262in}}%
\pgfpathlineto{\pgfqpoint{0.718698in}{1.613732in}}%
\pgfpathlineto{\pgfqpoint{0.730914in}{1.610662in}}%
\pgfpathlineto{\pgfqpoint{0.733470in}{1.622014in}}%
\pgfpathlineto{\pgfqpoint{0.738581in}{1.626973in}}%
\pgfpathlineto{\pgfqpoint{0.746366in}{1.637288in}}%
\pgfpathlineto{\pgfqpoint{0.754151in}{1.641671in}}%
\pgfpathlineto{\pgfqpoint{0.766659in}{1.638492in}}%
\pgfpathlineto{\pgfqpoint{0.763942in}{1.627183in}}%
\pgfpathlineto{\pgfqpoint{0.767582in}{1.623021in}}%
\pgfpathlineto{\pgfqpoint{0.765709in}{1.615018in}}%
\pgfpathlineto{\pgfqpoint{0.770478in}{1.610681in}}%
\pgfpathlineto{\pgfqpoint{0.775430in}{1.609514in}}%
\pgfpathlineto{\pgfqpoint{0.771340in}{1.602215in}}%
\pgfpathlineto{\pgfqpoint{0.768216in}{1.591309in}}%
\pgfpathlineto{\pgfqpoint{0.763908in}{1.592366in}}%
\pgfpathlineto{\pgfqpoint{0.762552in}{1.586835in}}%
\pgfpathlineto{\pgfqpoint{0.758886in}{1.587771in}}%
\pgfpathlineto{\pgfqpoint{0.755702in}{1.582709in}}%
\pgfpathlineto{\pgfqpoint{0.749229in}{1.584343in}}%
\pgfpathlineto{\pgfqpoint{0.746496in}{1.573138in}}%
\pgfpathlineto{\pgfqpoint{0.739760in}{1.573820in}}%
\pgfpathlineto{\pgfqpoint{0.736446in}{1.580574in}}%
\pgfpathlineto{\pgfqpoint{0.733534in}{1.580107in}}%
\pgfpathlineto{\pgfqpoint{0.729959in}{1.568628in}}%
\pgfpathlineto{\pgfqpoint{0.703346in}{1.460924in}}%
\pgfpathlineto{\pgfqpoint{0.647669in}{1.474848in}}%
\pgfpathlineto{\pgfqpoint{0.651488in}{1.492874in}}%
\pgfpathlineto{\pgfqpoint{0.662044in}{1.533668in}}%
\pgfpathlineto{\pgfqpoint{0.661467in}{1.533821in}}%
\pgfpathlineto{\pgfqpoint{0.675770in}{1.589097in}}%
\pgfpathlineto{\pgfqpoint{0.680340in}{1.605736in}}%
\pgfpathlineto{\pgfqpoint{0.683738in}{1.619619in}}%
\pgfpathclose%
\pgfusepath{fill}%
\end{pgfscope}%
\begin{pgfscope}%
\pgfpathrectangle{\pgfqpoint{0.100000in}{0.100000in}}{\pgfqpoint{3.007045in}{1.925000in}}%
\pgfusepath{clip}%
\pgfsetbuttcap%
\pgfsetmiterjoin%
\definecolor{currentfill}{rgb}{0.417086,0.680631,0.838231}%
\pgfsetfillcolor{currentfill}%
\pgfsetlinewidth{0.000000pt}%
\definecolor{currentstroke}{rgb}{0.000000,0.000000,0.000000}%
\pgfsetstrokecolor{currentstroke}%
\pgfsetstrokeopacity{0.000000}%
\pgfsetdash{}{0pt}%
\pgfpathmoveto{\pgfqpoint{2.402438in}{0.696626in}}%
\pgfpathlineto{\pgfqpoint{2.395249in}{0.698742in}}%
\pgfpathlineto{\pgfqpoint{2.394856in}{0.702917in}}%
\pgfpathlineto{\pgfqpoint{2.389600in}{0.714570in}}%
\pgfpathlineto{\pgfqpoint{2.386745in}{0.718055in}}%
\pgfpathlineto{\pgfqpoint{2.389930in}{0.724977in}}%
\pgfpathlineto{\pgfqpoint{2.395922in}{0.728401in}}%
\pgfpathlineto{\pgfqpoint{2.409404in}{0.729924in}}%
\pgfpathlineto{\pgfqpoint{2.410702in}{0.718714in}}%
\pgfpathlineto{\pgfqpoint{2.425917in}{0.719822in}}%
\pgfpathlineto{\pgfqpoint{2.428514in}{0.718344in}}%
\pgfpathlineto{\pgfqpoint{2.429949in}{0.711648in}}%
\pgfpathlineto{\pgfqpoint{2.433519in}{0.709598in}}%
\pgfpathlineto{\pgfqpoint{2.437033in}{0.704008in}}%
\pgfpathlineto{\pgfqpoint{2.420123in}{0.701779in}}%
\pgfpathlineto{\pgfqpoint{2.412679in}{0.701290in}}%
\pgfpathlineto{\pgfqpoint{2.413072in}{0.697996in}}%
\pgfpathlineto{\pgfqpoint{2.402438in}{0.696626in}}%
\pgfpathclose%
\pgfusepath{fill}%
\end{pgfscope}%
\begin{pgfscope}%
\pgfpathrectangle{\pgfqpoint{0.100000in}{0.100000in}}{\pgfqpoint{3.007045in}{1.925000in}}%
\pgfusepath{clip}%
\pgfsetbuttcap%
\pgfsetmiterjoin%
\definecolor{currentfill}{rgb}{0.326290,0.618624,0.802799}%
\pgfsetfillcolor{currentfill}%
\pgfsetlinewidth{0.000000pt}%
\definecolor{currentstroke}{rgb}{0.000000,0.000000,0.000000}%
\pgfsetstrokecolor{currentstroke}%
\pgfsetstrokeopacity{0.000000}%
\pgfsetdash{}{0pt}%
\pgfpathmoveto{\pgfqpoint{1.511195in}{1.206454in}}%
\pgfpathlineto{\pgfqpoint{1.510401in}{1.206498in}}%
\pgfpathlineto{\pgfqpoint{1.511500in}{1.229428in}}%
\pgfpathlineto{\pgfqpoint{1.482455in}{1.231145in}}%
\pgfpathlineto{\pgfqpoint{1.483897in}{1.254208in}}%
\pgfpathlineto{\pgfqpoint{1.511547in}{1.252543in}}%
\pgfpathlineto{\pgfqpoint{1.512812in}{1.275327in}}%
\pgfpathlineto{\pgfqpoint{1.552431in}{1.273213in}}%
\pgfpathlineto{\pgfqpoint{1.551658in}{1.248467in}}%
\pgfpathlineto{\pgfqpoint{1.563506in}{1.247230in}}%
\pgfpathlineto{\pgfqpoint{1.561531in}{1.203886in}}%
\pgfpathlineto{\pgfqpoint{1.539001in}{1.204970in}}%
\pgfpathlineto{\pgfqpoint{1.511195in}{1.206454in}}%
\pgfpathclose%
\pgfusepath{fill}%
\end{pgfscope}%
\begin{pgfscope}%
\pgfpathrectangle{\pgfqpoint{0.100000in}{0.100000in}}{\pgfqpoint{3.007045in}{1.925000in}}%
\pgfusepath{clip}%
\pgfsetbuttcap%
\pgfsetmiterjoin%
\definecolor{currentfill}{rgb}{0.341423,0.628958,0.808704}%
\pgfsetfillcolor{currentfill}%
\pgfsetlinewidth{0.000000pt}%
\definecolor{currentstroke}{rgb}{0.000000,0.000000,0.000000}%
\pgfsetstrokecolor{currentstroke}%
\pgfsetstrokeopacity{0.000000}%
\pgfsetdash{}{0pt}%
\pgfpathmoveto{\pgfqpoint{1.298841in}{1.455523in}}%
\pgfpathlineto{\pgfqpoint{1.339144in}{1.451533in}}%
\pgfpathlineto{\pgfqpoint{1.335923in}{1.418503in}}%
\pgfpathlineto{\pgfqpoint{1.333387in}{1.392952in}}%
\pgfpathlineto{\pgfqpoint{1.292810in}{1.397021in}}%
\pgfpathlineto{\pgfqpoint{1.293878in}{1.409890in}}%
\pgfpathlineto{\pgfqpoint{1.298841in}{1.455523in}}%
\pgfpathclose%
\pgfusepath{fill}%
\end{pgfscope}%
\begin{pgfscope}%
\pgfpathrectangle{\pgfqpoint{0.100000in}{0.100000in}}{\pgfqpoint{3.007045in}{1.925000in}}%
\pgfusepath{clip}%
\pgfsetbuttcap%
\pgfsetmiterjoin%
\definecolor{currentfill}{rgb}{0.632526,0.797647,0.886874}%
\pgfsetfillcolor{currentfill}%
\pgfsetlinewidth{0.000000pt}%
\definecolor{currentstroke}{rgb}{0.000000,0.000000,0.000000}%
\pgfsetstrokecolor{currentstroke}%
\pgfsetstrokeopacity{0.000000}%
\pgfsetdash{}{0pt}%
\pgfpathmoveto{\pgfqpoint{0.374195in}{1.289386in}}%
\pgfpathlineto{\pgfqpoint{0.377156in}{1.290848in}}%
\pgfpathlineto{\pgfqpoint{0.377641in}{1.299051in}}%
\pgfpathlineto{\pgfqpoint{0.376145in}{1.309067in}}%
\pgfpathlineto{\pgfqpoint{0.371022in}{1.321407in}}%
\pgfpathlineto{\pgfqpoint{0.373315in}{1.325445in}}%
\pgfpathlineto{\pgfqpoint{0.382079in}{1.325306in}}%
\pgfpathlineto{\pgfqpoint{0.388700in}{1.330958in}}%
\pgfpathlineto{\pgfqpoint{0.388567in}{1.334258in}}%
\pgfpathlineto{\pgfqpoint{0.393337in}{1.331101in}}%
\pgfpathlineto{\pgfqpoint{0.393792in}{1.321038in}}%
\pgfpathlineto{\pgfqpoint{0.395623in}{1.317666in}}%
\pgfpathlineto{\pgfqpoint{0.396070in}{1.307741in}}%
\pgfpathlineto{\pgfqpoint{0.400092in}{1.304827in}}%
\pgfpathlineto{\pgfqpoint{0.404388in}{1.306594in}}%
\pgfpathlineto{\pgfqpoint{0.416471in}{1.302466in}}%
\pgfpathlineto{\pgfqpoint{0.412425in}{1.289011in}}%
\pgfpathlineto{\pgfqpoint{0.417431in}{1.287454in}}%
\pgfpathlineto{\pgfqpoint{0.414282in}{1.280552in}}%
\pgfpathlineto{\pgfqpoint{0.409917in}{1.279165in}}%
\pgfpathlineto{\pgfqpoint{0.409876in}{1.274771in}}%
\pgfpathlineto{\pgfqpoint{0.413844in}{1.273461in}}%
\pgfpathlineto{\pgfqpoint{0.409743in}{1.259767in}}%
\pgfpathlineto{\pgfqpoint{0.409870in}{1.255826in}}%
\pgfpathlineto{\pgfqpoint{0.404604in}{1.238315in}}%
\pgfpathlineto{\pgfqpoint{0.407690in}{1.233195in}}%
\pgfpathlineto{\pgfqpoint{0.407661in}{1.233230in}}%
\pgfpathlineto{\pgfqpoint{0.388058in}{1.239323in}}%
\pgfpathlineto{\pgfqpoint{0.384479in}{1.238341in}}%
\pgfpathlineto{\pgfqpoint{0.377116in}{1.242201in}}%
\pgfpathlineto{\pgfqpoint{0.370754in}{1.240933in}}%
\pgfpathlineto{\pgfqpoint{0.368846in}{1.236860in}}%
\pgfpathlineto{\pgfqpoint{0.368463in}{1.226772in}}%
\pgfpathlineto{\pgfqpoint{0.359639in}{1.227533in}}%
\pgfpathlineto{\pgfqpoint{0.359301in}{1.222111in}}%
\pgfpathlineto{\pgfqpoint{0.355662in}{1.228408in}}%
\pgfpathlineto{\pgfqpoint{0.355631in}{1.232502in}}%
\pgfpathlineto{\pgfqpoint{0.358662in}{1.239132in}}%
\pgfpathlineto{\pgfqpoint{0.358656in}{1.248068in}}%
\pgfpathlineto{\pgfqpoint{0.356187in}{1.251923in}}%
\pgfpathlineto{\pgfqpoint{0.359786in}{1.258991in}}%
\pgfpathlineto{\pgfqpoint{0.361404in}{1.267963in}}%
\pgfpathlineto{\pgfqpoint{0.367391in}{1.268171in}}%
\pgfpathlineto{\pgfqpoint{0.365394in}{1.253654in}}%
\pgfpathlineto{\pgfqpoint{0.369838in}{1.250965in}}%
\pgfpathlineto{\pgfqpoint{0.373837in}{1.244770in}}%
\pgfpathlineto{\pgfqpoint{0.376610in}{1.245915in}}%
\pgfpathlineto{\pgfqpoint{0.376740in}{1.254309in}}%
\pgfpathlineto{\pgfqpoint{0.370828in}{1.266174in}}%
\pgfpathlineto{\pgfqpoint{0.373731in}{1.272170in}}%
\pgfpathlineto{\pgfqpoint{0.370644in}{1.277375in}}%
\pgfpathlineto{\pgfqpoint{0.379239in}{1.281168in}}%
\pgfpathlineto{\pgfqpoint{0.378012in}{1.285740in}}%
\pgfpathlineto{\pgfqpoint{0.374195in}{1.289386in}}%
\pgfpathclose%
\pgfusepath{fill}%
\end{pgfscope}%
\begin{pgfscope}%
\pgfpathrectangle{\pgfqpoint{0.100000in}{0.100000in}}{\pgfqpoint{3.007045in}{1.925000in}}%
\pgfusepath{clip}%
\pgfsetbuttcap%
\pgfsetmiterjoin%
\definecolor{currentfill}{rgb}{0.632526,0.797647,0.886874}%
\pgfsetfillcolor{currentfill}%
\pgfsetlinewidth{0.000000pt}%
\definecolor{currentstroke}{rgb}{0.000000,0.000000,0.000000}%
\pgfsetstrokecolor{currentstroke}%
\pgfsetstrokeopacity{0.000000}%
\pgfsetdash{}{0pt}%
\pgfpathmoveto{\pgfqpoint{0.369150in}{1.288189in}}%
\pgfpathlineto{\pgfqpoint{0.366091in}{1.274818in}}%
\pgfpathlineto{\pgfqpoint{0.361852in}{1.270865in}}%
\pgfpathlineto{\pgfqpoint{0.357053in}{1.277694in}}%
\pgfpathlineto{\pgfqpoint{0.354298in}{1.277874in}}%
\pgfpathlineto{\pgfqpoint{0.349286in}{1.287909in}}%
\pgfpathlineto{\pgfqpoint{0.344216in}{1.290431in}}%
\pgfpathlineto{\pgfqpoint{0.347199in}{1.297896in}}%
\pgfpathlineto{\pgfqpoint{0.347482in}{1.307924in}}%
\pgfpathlineto{\pgfqpoint{0.352980in}{1.307630in}}%
\pgfpathlineto{\pgfqpoint{0.358679in}{1.298249in}}%
\pgfpathlineto{\pgfqpoint{0.366648in}{1.294372in}}%
\pgfpathlineto{\pgfqpoint{0.369150in}{1.288189in}}%
\pgfpathclose%
\pgfusepath{fill}%
\end{pgfscope}%
\begin{pgfscope}%
\pgfpathrectangle{\pgfqpoint{0.100000in}{0.100000in}}{\pgfqpoint{3.007045in}{1.925000in}}%
\pgfusepath{clip}%
\pgfsetbuttcap%
\pgfsetmiterjoin%
\definecolor{currentfill}{rgb}{0.396909,0.666851,0.830358}%
\pgfsetfillcolor{currentfill}%
\pgfsetlinewidth{0.000000pt}%
\definecolor{currentstroke}{rgb}{0.000000,0.000000,0.000000}%
\pgfsetstrokecolor{currentstroke}%
\pgfsetstrokeopacity{0.000000}%
\pgfsetdash{}{0pt}%
\pgfpathmoveto{\pgfqpoint{2.025542in}{1.310434in}}%
\pgfpathlineto{\pgfqpoint{2.009626in}{1.309620in}}%
\pgfpathlineto{\pgfqpoint{2.006146in}{1.322530in}}%
\pgfpathlineto{\pgfqpoint{2.009028in}{1.324274in}}%
\pgfpathlineto{\pgfqpoint{2.010393in}{1.337303in}}%
\pgfpathlineto{\pgfqpoint{2.008917in}{1.344598in}}%
\pgfpathlineto{\pgfqpoint{2.001078in}{1.349385in}}%
\pgfpathlineto{\pgfqpoint{2.020468in}{1.350608in}}%
\pgfpathlineto{\pgfqpoint{2.018995in}{1.370976in}}%
\pgfpathlineto{\pgfqpoint{2.046233in}{1.372218in}}%
\pgfpathlineto{\pgfqpoint{2.078237in}{1.373879in}}%
\pgfpathlineto{\pgfqpoint{2.079874in}{1.351505in}}%
\pgfpathlineto{\pgfqpoint{2.085588in}{1.351908in}}%
\pgfpathlineto{\pgfqpoint{2.085344in}{1.346129in}}%
\pgfpathlineto{\pgfqpoint{2.087331in}{1.317506in}}%
\pgfpathlineto{\pgfqpoint{2.059573in}{1.315459in}}%
\pgfpathlineto{\pgfqpoint{2.059792in}{1.312587in}}%
\pgfpathlineto{\pgfqpoint{2.025542in}{1.310434in}}%
\pgfpathclose%
\pgfusepath{fill}%
\end{pgfscope}%
\begin{pgfscope}%
\pgfpathrectangle{\pgfqpoint{0.100000in}{0.100000in}}{\pgfqpoint{3.007045in}{1.925000in}}%
\pgfusepath{clip}%
\pgfsetbuttcap%
\pgfsetmiterjoin%
\definecolor{currentfill}{rgb}{0.447843,0.697855,0.845306}%
\pgfsetfillcolor{currentfill}%
\pgfsetlinewidth{0.000000pt}%
\definecolor{currentstroke}{rgb}{0.000000,0.000000,0.000000}%
\pgfsetstrokecolor{currentstroke}%
\pgfsetstrokeopacity{0.000000}%
\pgfsetdash{}{0pt}%
\pgfpathmoveto{\pgfqpoint{2.169528in}{1.143936in}}%
\pgfpathlineto{\pgfqpoint{2.175194in}{1.144388in}}%
\pgfpathlineto{\pgfqpoint{2.175011in}{1.140154in}}%
\pgfpathlineto{\pgfqpoint{2.170118in}{1.138155in}}%
\pgfpathlineto{\pgfqpoint{2.168758in}{1.129637in}}%
\pgfpathlineto{\pgfqpoint{2.171500in}{1.120112in}}%
\pgfpathlineto{\pgfqpoint{2.160227in}{1.118407in}}%
\pgfpathlineto{\pgfqpoint{2.154476in}{1.114634in}}%
\pgfpathlineto{\pgfqpoint{2.150356in}{1.119886in}}%
\pgfpathlineto{\pgfqpoint{2.136969in}{1.118978in}}%
\pgfpathlineto{\pgfqpoint{2.135746in}{1.137366in}}%
\pgfpathlineto{\pgfqpoint{2.154997in}{1.139057in}}%
\pgfpathlineto{\pgfqpoint{2.155147in}{1.142366in}}%
\pgfpathlineto{\pgfqpoint{2.169528in}{1.143936in}}%
\pgfpathclose%
\pgfusepath{fill}%
\end{pgfscope}%
\begin{pgfscope}%
\pgfpathrectangle{\pgfqpoint{0.100000in}{0.100000in}}{\pgfqpoint{3.007045in}{1.925000in}}%
\pgfusepath{clip}%
\pgfsetbuttcap%
\pgfsetmiterjoin%
\definecolor{currentfill}{rgb}{0.296025,0.597955,0.790988}%
\pgfsetfillcolor{currentfill}%
\pgfsetlinewidth{0.000000pt}%
\definecolor{currentstroke}{rgb}{0.000000,0.000000,0.000000}%
\pgfsetstrokecolor{currentstroke}%
\pgfsetstrokeopacity{0.000000}%
\pgfsetdash{}{0pt}%
\pgfpathmoveto{\pgfqpoint{1.735784in}{0.842654in}}%
\pgfpathlineto{\pgfqpoint{1.735756in}{0.836902in}}%
\pgfpathlineto{\pgfqpoint{1.729891in}{0.836933in}}%
\pgfpathlineto{\pgfqpoint{1.729695in}{0.813812in}}%
\pgfpathlineto{\pgfqpoint{1.718089in}{0.813933in}}%
\pgfpathlineto{\pgfqpoint{1.695332in}{0.814267in}}%
\pgfpathlineto{\pgfqpoint{1.695616in}{0.831530in}}%
\pgfpathlineto{\pgfqpoint{1.689815in}{0.831695in}}%
\pgfpathlineto{\pgfqpoint{1.689906in}{0.837361in}}%
\pgfpathlineto{\pgfqpoint{1.672844in}{0.837714in}}%
\pgfpathlineto{\pgfqpoint{1.672960in}{0.843473in}}%
\pgfpathlineto{\pgfqpoint{1.667244in}{0.843587in}}%
\pgfpathlineto{\pgfqpoint{1.667600in}{0.860848in}}%
\pgfpathlineto{\pgfqpoint{1.670026in}{0.865689in}}%
\pgfpathlineto{\pgfqpoint{1.674047in}{0.866744in}}%
\pgfpathlineto{\pgfqpoint{1.678967in}{0.860772in}}%
\pgfpathlineto{\pgfqpoint{1.683133in}{0.863693in}}%
\pgfpathlineto{\pgfqpoint{1.691602in}{0.864102in}}%
\pgfpathlineto{\pgfqpoint{1.691843in}{0.877659in}}%
\pgfpathlineto{\pgfqpoint{1.694493in}{0.877607in}}%
\pgfpathlineto{\pgfqpoint{1.694674in}{0.889113in}}%
\pgfpathlineto{\pgfqpoint{1.719456in}{0.888768in}}%
\pgfpathlineto{\pgfqpoint{1.719231in}{0.879638in}}%
\pgfpathlineto{\pgfqpoint{1.713476in}{0.872989in}}%
\pgfpathlineto{\pgfqpoint{1.713003in}{0.848605in}}%
\pgfpathlineto{\pgfqpoint{1.724413in}{0.848474in}}%
\pgfpathlineto{\pgfqpoint{1.724353in}{0.842730in}}%
\pgfpathlineto{\pgfqpoint{1.735784in}{0.842654in}}%
\pgfpathclose%
\pgfusepath{fill}%
\end{pgfscope}%
\begin{pgfscope}%
\pgfpathrectangle{\pgfqpoint{0.100000in}{0.100000in}}{\pgfqpoint{3.007045in}{1.925000in}}%
\pgfusepath{clip}%
\pgfsetbuttcap%
\pgfsetmiterjoin%
\definecolor{currentfill}{rgb}{0.466667,0.708189,0.849366}%
\pgfsetfillcolor{currentfill}%
\pgfsetlinewidth{0.000000pt}%
\definecolor{currentstroke}{rgb}{0.000000,0.000000,0.000000}%
\pgfsetstrokecolor{currentstroke}%
\pgfsetstrokeopacity{0.000000}%
\pgfsetdash{}{0pt}%
\pgfpathmoveto{\pgfqpoint{2.643747in}{1.050644in}}%
\pgfpathlineto{\pgfqpoint{2.672729in}{1.056217in}}%
\pgfpathlineto{\pgfqpoint{2.674093in}{1.045577in}}%
\pgfpathlineto{\pgfqpoint{2.671365in}{1.033005in}}%
\pgfpathlineto{\pgfqpoint{2.664233in}{1.034590in}}%
\pgfpathlineto{\pgfqpoint{2.654897in}{1.034173in}}%
\pgfpathlineto{\pgfqpoint{2.650549in}{1.026900in}}%
\pgfpathlineto{\pgfqpoint{2.645335in}{1.026823in}}%
\pgfpathlineto{\pgfqpoint{2.641967in}{1.044405in}}%
\pgfpathlineto{\pgfqpoint{2.643747in}{1.050644in}}%
\pgfpathclose%
\pgfusepath{fill}%
\end{pgfscope}%
\begin{pgfscope}%
\pgfpathrectangle{\pgfqpoint{0.100000in}{0.100000in}}{\pgfqpoint{3.007045in}{1.925000in}}%
\pgfusepath{clip}%
\pgfsetbuttcap%
\pgfsetmiterjoin%
\definecolor{currentfill}{rgb}{0.311157,0.608289,0.796894}%
\pgfsetfillcolor{currentfill}%
\pgfsetlinewidth{0.000000pt}%
\definecolor{currentstroke}{rgb}{0.000000,0.000000,0.000000}%
\pgfsetstrokecolor{currentstroke}%
\pgfsetstrokeopacity{0.000000}%
\pgfsetdash{}{0pt}%
\pgfpathmoveto{\pgfqpoint{1.003150in}{0.848168in}}%
\pgfpathlineto{\pgfqpoint{0.997334in}{0.810949in}}%
\pgfpathlineto{\pgfqpoint{0.989331in}{0.759721in}}%
\pgfpathlineto{\pgfqpoint{0.927312in}{0.769778in}}%
\pgfpathlineto{\pgfqpoint{0.911810in}{0.772480in}}%
\pgfpathlineto{\pgfqpoint{0.912796in}{0.778126in}}%
\pgfpathlineto{\pgfqpoint{0.920548in}{0.822552in}}%
\pgfpathlineto{\pgfqpoint{0.923055in}{0.834676in}}%
\pgfpathlineto{\pgfqpoint{0.930194in}{0.836921in}}%
\pgfpathlineto{\pgfqpoint{0.938796in}{0.837657in}}%
\pgfpathlineto{\pgfqpoint{0.947992in}{0.836130in}}%
\pgfpathlineto{\pgfqpoint{0.949199in}{0.843379in}}%
\pgfpathlineto{\pgfqpoint{0.954625in}{0.842378in}}%
\pgfpathlineto{\pgfqpoint{0.960834in}{0.834656in}}%
\pgfpathlineto{\pgfqpoint{0.971108in}{0.842140in}}%
\pgfpathlineto{\pgfqpoint{0.975090in}{0.841991in}}%
\pgfpathlineto{\pgfqpoint{0.978685in}{0.845988in}}%
\pgfpathlineto{\pgfqpoint{0.986809in}{0.850694in}}%
\pgfpathlineto{\pgfqpoint{1.003150in}{0.848168in}}%
\pgfpathclose%
\pgfusepath{fill}%
\end{pgfscope}%
\begin{pgfscope}%
\pgfpathrectangle{\pgfqpoint{0.100000in}{0.100000in}}{\pgfqpoint{3.007045in}{1.925000in}}%
\pgfusepath{clip}%
\pgfsetbuttcap%
\pgfsetmiterjoin%
\definecolor{currentfill}{rgb}{0.331334,0.622068,0.804767}%
\pgfsetfillcolor{currentfill}%
\pgfsetlinewidth{0.000000pt}%
\definecolor{currentstroke}{rgb}{0.000000,0.000000,0.000000}%
\pgfsetstrokecolor{currentstroke}%
\pgfsetstrokeopacity{0.000000}%
\pgfsetdash{}{0pt}%
\pgfpathmoveto{\pgfqpoint{1.730423in}{0.638033in}}%
\pgfpathlineto{\pgfqpoint{1.734214in}{0.624182in}}%
\pgfpathlineto{\pgfqpoint{1.731113in}{0.622526in}}%
\pgfpathlineto{\pgfqpoint{1.728210in}{0.613654in}}%
\pgfpathlineto{\pgfqpoint{1.728806in}{0.610818in}}%
\pgfpathlineto{\pgfqpoint{1.716755in}{0.610772in}}%
\pgfpathlineto{\pgfqpoint{1.701562in}{0.603092in}}%
\pgfpathlineto{\pgfqpoint{1.696915in}{0.623980in}}%
\pgfpathlineto{\pgfqpoint{1.697477in}{0.628654in}}%
\pgfpathlineto{\pgfqpoint{1.681563in}{0.619854in}}%
\pgfpathlineto{\pgfqpoint{1.663193in}{0.652637in}}%
\pgfpathlineto{\pgfqpoint{1.675361in}{0.659431in}}%
\pgfpathlineto{\pgfqpoint{1.665733in}{0.676821in}}%
\pgfpathlineto{\pgfqpoint{1.694814in}{0.693189in}}%
\pgfpathlineto{\pgfqpoint{1.703043in}{0.687922in}}%
\pgfpathlineto{\pgfqpoint{1.712375in}{0.673683in}}%
\pgfpathlineto{\pgfqpoint{1.718592in}{0.656205in}}%
\pgfpathlineto{\pgfqpoint{1.722872in}{0.654828in}}%
\pgfpathlineto{\pgfqpoint{1.723459in}{0.650585in}}%
\pgfpathlineto{\pgfqpoint{1.732103in}{0.644712in}}%
\pgfpathlineto{\pgfqpoint{1.730423in}{0.638033in}}%
\pgfpathclose%
\pgfusepath{fill}%
\end{pgfscope}%
\begin{pgfscope}%
\pgfpathrectangle{\pgfqpoint{0.100000in}{0.100000in}}{\pgfqpoint{3.007045in}{1.925000in}}%
\pgfusepath{clip}%
\pgfsetbuttcap%
\pgfsetmiterjoin%
\definecolor{currentfill}{rgb}{0.441569,0.694410,0.843952}%
\pgfsetfillcolor{currentfill}%
\pgfsetlinewidth{0.000000pt}%
\definecolor{currentstroke}{rgb}{0.000000,0.000000,0.000000}%
\pgfsetstrokecolor{currentstroke}%
\pgfsetstrokeopacity{0.000000}%
\pgfsetdash{}{0pt}%
\pgfpathmoveto{\pgfqpoint{1.643166in}{1.743370in}}%
\pgfpathlineto{\pgfqpoint{1.643745in}{1.731622in}}%
\pgfpathlineto{\pgfqpoint{1.642987in}{1.708503in}}%
\pgfpathlineto{\pgfqpoint{1.638404in}{1.708645in}}%
\pgfpathlineto{\pgfqpoint{1.615296in}{1.709477in}}%
\pgfpathlineto{\pgfqpoint{1.614607in}{1.721113in}}%
\pgfpathlineto{\pgfqpoint{1.580095in}{1.722634in}}%
\pgfpathlineto{\pgfqpoint{1.557103in}{1.723835in}}%
\pgfpathlineto{\pgfqpoint{1.556105in}{1.735501in}}%
\pgfpathlineto{\pgfqpoint{1.557416in}{1.758816in}}%
\pgfpathlineto{\pgfqpoint{1.573086in}{1.757987in}}%
\pgfpathlineto{\pgfqpoint{1.596215in}{1.756877in}}%
\pgfpathlineto{\pgfqpoint{1.596746in}{1.768470in}}%
\pgfpathlineto{\pgfqpoint{1.625579in}{1.767286in}}%
\pgfpathlineto{\pgfqpoint{1.626307in}{1.755616in}}%
\pgfpathlineto{\pgfqpoint{1.625869in}{1.743967in}}%
\pgfpathlineto{\pgfqpoint{1.643166in}{1.743370in}}%
\pgfpathclose%
\pgfusepath{fill}%
\end{pgfscope}%
\begin{pgfscope}%
\pgfpathrectangle{\pgfqpoint{0.100000in}{0.100000in}}{\pgfqpoint{3.007045in}{1.925000in}}%
\pgfusepath{clip}%
\pgfsetbuttcap%
\pgfsetmiterjoin%
\definecolor{currentfill}{rgb}{0.376732,0.653072,0.822484}%
\pgfsetfillcolor{currentfill}%
\pgfsetlinewidth{0.000000pt}%
\definecolor{currentstroke}{rgb}{0.000000,0.000000,0.000000}%
\pgfsetstrokecolor{currentstroke}%
\pgfsetstrokeopacity{0.000000}%
\pgfsetdash{}{0pt}%
\pgfpathmoveto{\pgfqpoint{2.060711in}{0.508116in}}%
\pgfpathlineto{\pgfqpoint{2.060425in}{0.500127in}}%
\pgfpathlineto{\pgfqpoint{2.056411in}{0.498585in}}%
\pgfpathlineto{\pgfqpoint{2.060074in}{0.492771in}}%
\pgfpathlineto{\pgfqpoint{2.058740in}{0.489589in}}%
\pgfpathlineto{\pgfqpoint{2.051172in}{0.484014in}}%
\pgfpathlineto{\pgfqpoint{2.048181in}{0.488038in}}%
\pgfpathlineto{\pgfqpoint{2.046171in}{0.495142in}}%
\pgfpathlineto{\pgfqpoint{2.042744in}{0.497419in}}%
\pgfpathlineto{\pgfqpoint{2.040294in}{0.492004in}}%
\pgfpathlineto{\pgfqpoint{2.034473in}{0.496696in}}%
\pgfpathlineto{\pgfqpoint{2.026621in}{0.490352in}}%
\pgfpathlineto{\pgfqpoint{2.026344in}{0.485680in}}%
\pgfpathlineto{\pgfqpoint{2.021444in}{0.485001in}}%
\pgfpathlineto{\pgfqpoint{2.014806in}{0.481223in}}%
\pgfpathlineto{\pgfqpoint{2.007886in}{0.488255in}}%
\pgfpathlineto{\pgfqpoint{2.000041in}{0.488340in}}%
\pgfpathlineto{\pgfqpoint{1.989313in}{0.491975in}}%
\pgfpathlineto{\pgfqpoint{1.985959in}{0.494974in}}%
\pgfpathlineto{\pgfqpoint{1.991281in}{0.500512in}}%
\pgfpathlineto{\pgfqpoint{1.989147in}{0.507119in}}%
\pgfpathlineto{\pgfqpoint{1.973270in}{0.512709in}}%
\pgfpathlineto{\pgfqpoint{1.968012in}{0.516450in}}%
\pgfpathlineto{\pgfqpoint{1.968031in}{0.523710in}}%
\pgfpathlineto{\pgfqpoint{1.958964in}{0.523855in}}%
\pgfpathlineto{\pgfqpoint{1.953040in}{0.522060in}}%
\pgfpathlineto{\pgfqpoint{1.962543in}{0.530961in}}%
\pgfpathlineto{\pgfqpoint{1.967294in}{0.537756in}}%
\pgfpathlineto{\pgfqpoint{1.976086in}{0.538805in}}%
\pgfpathlineto{\pgfqpoint{1.982084in}{0.530926in}}%
\pgfpathlineto{\pgfqpoint{1.981476in}{0.527516in}}%
\pgfpathlineto{\pgfqpoint{1.993498in}{0.524333in}}%
\pgfpathlineto{\pgfqpoint{1.998738in}{0.526353in}}%
\pgfpathlineto{\pgfqpoint{1.997286in}{0.531031in}}%
\pgfpathlineto{\pgfqpoint{1.992985in}{0.532196in}}%
\pgfpathlineto{\pgfqpoint{1.988235in}{0.542241in}}%
\pgfpathlineto{\pgfqpoint{1.990185in}{0.543999in}}%
\pgfpathlineto{\pgfqpoint{2.002634in}{0.548120in}}%
\pgfpathlineto{\pgfqpoint{2.005096in}{0.547382in}}%
\pgfpathlineto{\pgfqpoint{2.005579in}{0.541500in}}%
\pgfpathlineto{\pgfqpoint{2.010082in}{0.536853in}}%
\pgfpathlineto{\pgfqpoint{2.030059in}{0.536911in}}%
\pgfpathlineto{\pgfqpoint{2.034075in}{0.531818in}}%
\pgfpathlineto{\pgfqpoint{2.040150in}{0.528704in}}%
\pgfpathlineto{\pgfqpoint{2.042354in}{0.524234in}}%
\pgfpathlineto{\pgfqpoint{2.050269in}{0.524553in}}%
\pgfpathlineto{\pgfqpoint{2.053131in}{0.522603in}}%
\pgfpathlineto{\pgfqpoint{2.051671in}{0.516053in}}%
\pgfpathlineto{\pgfqpoint{2.054795in}{0.510464in}}%
\pgfpathlineto{\pgfqpoint{2.060711in}{0.508116in}}%
\pgfpathclose%
\pgfusepath{fill}%
\end{pgfscope}%
\begin{pgfscope}%
\pgfpathrectangle{\pgfqpoint{0.100000in}{0.100000in}}{\pgfqpoint{3.007045in}{1.925000in}}%
\pgfusepath{clip}%
\pgfsetbuttcap%
\pgfsetmiterjoin%
\definecolor{currentfill}{rgb}{0.244106,0.557832,0.768889}%
\pgfsetfillcolor{currentfill}%
\pgfsetlinewidth{0.000000pt}%
\definecolor{currentstroke}{rgb}{0.000000,0.000000,0.000000}%
\pgfsetstrokecolor{currentstroke}%
\pgfsetstrokeopacity{0.000000}%
\pgfsetdash{}{0pt}%
\pgfpathmoveto{\pgfqpoint{1.502115in}{0.878788in}}%
\pgfpathlineto{\pgfqpoint{1.504252in}{0.917589in}}%
\pgfpathlineto{\pgfqpoint{1.505202in}{0.934838in}}%
\pgfpathlineto{\pgfqpoint{1.510237in}{0.939913in}}%
\pgfpathlineto{\pgfqpoint{1.515519in}{0.937448in}}%
\pgfpathlineto{\pgfqpoint{1.520249in}{0.932634in}}%
\pgfpathlineto{\pgfqpoint{1.528739in}{0.932611in}}%
\pgfpathlineto{\pgfqpoint{1.532268in}{0.939536in}}%
\pgfpathlineto{\pgfqpoint{1.538682in}{0.941826in}}%
\pgfpathlineto{\pgfqpoint{1.538310in}{0.928586in}}%
\pgfpathlineto{\pgfqpoint{1.578070in}{0.926846in}}%
\pgfpathlineto{\pgfqpoint{1.577863in}{0.909634in}}%
\pgfpathlineto{\pgfqpoint{1.576781in}{0.879670in}}%
\pgfpathlineto{\pgfqpoint{1.569929in}{0.881210in}}%
\pgfpathlineto{\pgfqpoint{1.514160in}{0.883851in}}%
\pgfpathlineto{\pgfqpoint{1.513852in}{0.878138in}}%
\pgfpathlineto{\pgfqpoint{1.502115in}{0.878788in}}%
\pgfpathclose%
\pgfusepath{fill}%
\end{pgfscope}%
\begin{pgfscope}%
\pgfpathrectangle{\pgfqpoint{0.100000in}{0.100000in}}{\pgfqpoint{3.007045in}{1.925000in}}%
\pgfusepath{clip}%
\pgfsetbuttcap%
\pgfsetmiterjoin%
\definecolor{currentfill}{rgb}{0.491765,0.721968,0.854779}%
\pgfsetfillcolor{currentfill}%
\pgfsetlinewidth{0.000000pt}%
\definecolor{currentstroke}{rgb}{0.000000,0.000000,0.000000}%
\pgfsetstrokecolor{currentstroke}%
\pgfsetstrokeopacity{0.000000}%
\pgfsetdash{}{0pt}%
\pgfpathmoveto{\pgfqpoint{2.523965in}{1.201691in}}%
\pgfpathlineto{\pgfqpoint{2.534247in}{1.194214in}}%
\pgfpathlineto{\pgfqpoint{2.532223in}{1.188646in}}%
\pgfpathlineto{\pgfqpoint{2.532999in}{1.178683in}}%
\pgfpathlineto{\pgfqpoint{2.527409in}{1.177624in}}%
\pgfpathlineto{\pgfqpoint{2.525737in}{1.174087in}}%
\pgfpathlineto{\pgfqpoint{2.517598in}{1.175060in}}%
\pgfpathlineto{\pgfqpoint{2.514455in}{1.181489in}}%
\pgfpathlineto{\pgfqpoint{2.507565in}{1.184564in}}%
\pgfpathlineto{\pgfqpoint{2.500337in}{1.190919in}}%
\pgfpathlineto{\pgfqpoint{2.498857in}{1.196503in}}%
\pgfpathlineto{\pgfqpoint{2.505727in}{1.202402in}}%
\pgfpathlineto{\pgfqpoint{2.520381in}{1.200511in}}%
\pgfpathlineto{\pgfqpoint{2.523965in}{1.201691in}}%
\pgfpathclose%
\pgfusepath{fill}%
\end{pgfscope}%
\begin{pgfscope}%
\pgfpathrectangle{\pgfqpoint{0.100000in}{0.100000in}}{\pgfqpoint{3.007045in}{1.925000in}}%
\pgfusepath{clip}%
\pgfsetbuttcap%
\pgfsetmiterjoin%
\definecolor{currentfill}{rgb}{0.441569,0.694410,0.843952}%
\pgfsetfillcolor{currentfill}%
\pgfsetlinewidth{0.000000pt}%
\definecolor{currentstroke}{rgb}{0.000000,0.000000,0.000000}%
\pgfsetstrokecolor{currentstroke}%
\pgfsetstrokeopacity{0.000000}%
\pgfsetdash{}{0pt}%
\pgfpathmoveto{\pgfqpoint{2.362396in}{1.092621in}}%
\pgfpathlineto{\pgfqpoint{2.358589in}{1.084899in}}%
\pgfpathlineto{\pgfqpoint{2.350883in}{1.086320in}}%
\pgfpathlineto{\pgfqpoint{2.342975in}{1.090758in}}%
\pgfpathlineto{\pgfqpoint{2.328598in}{1.094602in}}%
\pgfpathlineto{\pgfqpoint{2.316608in}{1.083386in}}%
\pgfpathlineto{\pgfqpoint{2.314503in}{1.080433in}}%
\pgfpathlineto{\pgfqpoint{2.300122in}{1.093676in}}%
\pgfpathlineto{\pgfqpoint{2.294182in}{1.093073in}}%
\pgfpathlineto{\pgfqpoint{2.289598in}{1.088435in}}%
\pgfpathlineto{\pgfqpoint{2.283241in}{1.088142in}}%
\pgfpathlineto{\pgfqpoint{2.281936in}{1.092901in}}%
\pgfpathlineto{\pgfqpoint{2.287143in}{1.100325in}}%
\pgfpathlineto{\pgfqpoint{2.288667in}{1.106429in}}%
\pgfpathlineto{\pgfqpoint{2.288467in}{1.113713in}}%
\pgfpathlineto{\pgfqpoint{2.292494in}{1.120777in}}%
\pgfpathlineto{\pgfqpoint{2.287382in}{1.127701in}}%
\pgfpathlineto{\pgfqpoint{2.285748in}{1.132441in}}%
\pgfpathlineto{\pgfqpoint{2.282277in}{1.134362in}}%
\pgfpathlineto{\pgfqpoint{2.288955in}{1.139422in}}%
\pgfpathlineto{\pgfqpoint{2.296120in}{1.144067in}}%
\pgfpathlineto{\pgfqpoint{2.297453in}{1.137486in}}%
\pgfpathlineto{\pgfqpoint{2.308461in}{1.129091in}}%
\pgfpathlineto{\pgfqpoint{2.317016in}{1.135906in}}%
\pgfpathlineto{\pgfqpoint{2.326824in}{1.138484in}}%
\pgfpathlineto{\pgfqpoint{2.332965in}{1.140303in}}%
\pgfpathlineto{\pgfqpoint{2.338933in}{1.139975in}}%
\pgfpathlineto{\pgfqpoint{2.341607in}{1.134315in}}%
\pgfpathlineto{\pgfqpoint{2.339247in}{1.131363in}}%
\pgfpathlineto{\pgfqpoint{2.347373in}{1.122044in}}%
\pgfpathlineto{\pgfqpoint{2.336897in}{1.108708in}}%
\pgfpathlineto{\pgfqpoint{2.344556in}{1.097617in}}%
\pgfpathlineto{\pgfqpoint{2.355335in}{1.096669in}}%
\pgfpathlineto{\pgfqpoint{2.362396in}{1.092621in}}%
\pgfpathclose%
\pgfusepath{fill}%
\end{pgfscope}%
\begin{pgfscope}%
\pgfpathrectangle{\pgfqpoint{0.100000in}{0.100000in}}{\pgfqpoint{3.007045in}{1.925000in}}%
\pgfusepath{clip}%
\pgfsetbuttcap%
\pgfsetmiterjoin%
\definecolor{currentfill}{rgb}{0.376732,0.653072,0.822484}%
\pgfsetfillcolor{currentfill}%
\pgfsetlinewidth{0.000000pt}%
\definecolor{currentstroke}{rgb}{0.000000,0.000000,0.000000}%
\pgfsetstrokecolor{currentstroke}%
\pgfsetstrokeopacity{0.000000}%
\pgfsetdash{}{0pt}%
\pgfpathmoveto{\pgfqpoint{2.231222in}{1.236425in}}%
\pgfpathlineto{\pgfqpoint{2.212181in}{1.234336in}}%
\pgfpathlineto{\pgfqpoint{2.211166in}{1.244697in}}%
\pgfpathlineto{\pgfqpoint{2.204170in}{1.247864in}}%
\pgfpathlineto{\pgfqpoint{2.202654in}{1.264977in}}%
\pgfpathlineto{\pgfqpoint{2.195116in}{1.264251in}}%
\pgfpathlineto{\pgfqpoint{2.192034in}{1.266835in}}%
\pgfpathlineto{\pgfqpoint{2.190977in}{1.278330in}}%
\pgfpathlineto{\pgfqpoint{2.173650in}{1.276857in}}%
\pgfpathlineto{\pgfqpoint{2.172150in}{1.293918in}}%
\pgfpathlineto{\pgfqpoint{2.195043in}{1.295965in}}%
\pgfpathlineto{\pgfqpoint{2.196588in}{1.278782in}}%
\pgfpathlineto{\pgfqpoint{2.211433in}{1.280191in}}%
\pgfpathlineto{\pgfqpoint{2.210917in}{1.285878in}}%
\pgfpathlineto{\pgfqpoint{2.221882in}{1.287153in}}%
\pgfpathlineto{\pgfqpoint{2.224496in}{1.264365in}}%
\pgfpathlineto{\pgfqpoint{2.228217in}{1.264785in}}%
\pgfpathlineto{\pgfqpoint{2.231222in}{1.236425in}}%
\pgfpathclose%
\pgfusepath{fill}%
\end{pgfscope}%
\begin{pgfscope}%
\pgfpathrectangle{\pgfqpoint{0.100000in}{0.100000in}}{\pgfqpoint{3.007045in}{1.925000in}}%
\pgfusepath{clip}%
\pgfsetbuttcap%
\pgfsetmiterjoin%
\definecolor{currentfill}{rgb}{0.422745,0.684075,0.839892}%
\pgfsetfillcolor{currentfill}%
\pgfsetlinewidth{0.000000pt}%
\definecolor{currentstroke}{rgb}{0.000000,0.000000,0.000000}%
\pgfsetstrokecolor{currentstroke}%
\pgfsetstrokeopacity{0.000000}%
\pgfsetdash{}{0pt}%
\pgfpathmoveto{\pgfqpoint{1.980505in}{0.873918in}}%
\pgfpathlineto{\pgfqpoint{1.972810in}{0.872450in}}%
\pgfpathlineto{\pgfqpoint{1.968811in}{0.876703in}}%
\pgfpathlineto{\pgfqpoint{1.963244in}{0.878965in}}%
\pgfpathlineto{\pgfqpoint{1.956817in}{0.878848in}}%
\pgfpathlineto{\pgfqpoint{1.956926in}{0.874466in}}%
\pgfpathlineto{\pgfqpoint{1.950043in}{0.876695in}}%
\pgfpathlineto{\pgfqpoint{1.945167in}{0.874395in}}%
\pgfpathlineto{\pgfqpoint{1.941240in}{0.873795in}}%
\pgfpathlineto{\pgfqpoint{1.933967in}{0.876483in}}%
\pgfpathlineto{\pgfqpoint{1.927950in}{0.876198in}}%
\pgfpathlineto{\pgfqpoint{1.927775in}{0.895762in}}%
\pgfpathlineto{\pgfqpoint{1.941745in}{0.895933in}}%
\pgfpathlineto{\pgfqpoint{1.944537in}{0.902457in}}%
\pgfpathlineto{\pgfqpoint{1.944456in}{0.907639in}}%
\pgfpathlineto{\pgfqpoint{1.967250in}{0.907985in}}%
\pgfpathlineto{\pgfqpoint{1.966993in}{0.915480in}}%
\pgfpathlineto{\pgfqpoint{1.969720in}{0.923885in}}%
\pgfpathlineto{\pgfqpoint{1.975423in}{0.932334in}}%
\pgfpathlineto{\pgfqpoint{1.984303in}{0.932387in}}%
\pgfpathlineto{\pgfqpoint{1.985207in}{0.903228in}}%
\pgfpathlineto{\pgfqpoint{1.986161in}{0.874058in}}%
\pgfpathlineto{\pgfqpoint{1.980505in}{0.873918in}}%
\pgfpathclose%
\pgfusepath{fill}%
\end{pgfscope}%
\begin{pgfscope}%
\pgfpathrectangle{\pgfqpoint{0.100000in}{0.100000in}}{\pgfqpoint{3.007045in}{1.925000in}}%
\pgfusepath{clip}%
\pgfsetbuttcap%
\pgfsetmiterjoin%
\definecolor{currentfill}{rgb}{0.346467,0.632403,0.810673}%
\pgfsetfillcolor{currentfill}%
\pgfsetlinewidth{0.000000pt}%
\definecolor{currentstroke}{rgb}{0.000000,0.000000,0.000000}%
\pgfsetstrokecolor{currentstroke}%
\pgfsetstrokeopacity{0.000000}%
\pgfsetdash{}{0pt}%
\pgfpathmoveto{\pgfqpoint{1.553928in}{0.768166in}}%
\pgfpathlineto{\pgfqpoint{1.555230in}{0.797022in}}%
\pgfpathlineto{\pgfqpoint{1.556351in}{0.822025in}}%
\pgfpathlineto{\pgfqpoint{1.560430in}{0.818079in}}%
\pgfpathlineto{\pgfqpoint{1.570469in}{0.816172in}}%
\pgfpathlineto{\pgfqpoint{1.574901in}{0.817578in}}%
\pgfpathlineto{\pgfqpoint{1.581430in}{0.811070in}}%
\pgfpathlineto{\pgfqpoint{1.586260in}{0.813360in}}%
\pgfpathlineto{\pgfqpoint{1.588268in}{0.817057in}}%
\pgfpathlineto{\pgfqpoint{1.598818in}{0.813845in}}%
\pgfpathlineto{\pgfqpoint{1.605264in}{0.805706in}}%
\pgfpathlineto{\pgfqpoint{1.610850in}{0.805245in}}%
\pgfpathlineto{\pgfqpoint{1.608887in}{0.798669in}}%
\pgfpathlineto{\pgfqpoint{1.607913in}{0.770715in}}%
\pgfpathlineto{\pgfqpoint{1.583500in}{0.771553in}}%
\pgfpathlineto{\pgfqpoint{1.583365in}{0.766882in}}%
\pgfpathlineto{\pgfqpoint{1.553928in}{0.768166in}}%
\pgfpathclose%
\pgfusepath{fill}%
\end{pgfscope}%
\begin{pgfscope}%
\pgfpathrectangle{\pgfqpoint{0.100000in}{0.100000in}}{\pgfqpoint{3.007045in}{1.925000in}}%
\pgfusepath{clip}%
\pgfsetbuttcap%
\pgfsetmiterjoin%
\definecolor{currentfill}{rgb}{0.326290,0.618624,0.802799}%
\pgfsetfillcolor{currentfill}%
\pgfsetlinewidth{0.000000pt}%
\definecolor{currentstroke}{rgb}{0.000000,0.000000,0.000000}%
\pgfsetstrokecolor{currentstroke}%
\pgfsetstrokeopacity{0.000000}%
\pgfsetdash{}{0pt}%
\pgfpathmoveto{\pgfqpoint{1.872670in}{1.072352in}}%
\pgfpathlineto{\pgfqpoint{1.849904in}{1.072618in}}%
\pgfpathlineto{\pgfqpoint{1.849656in}{1.062106in}}%
\pgfpathlineto{\pgfqpoint{1.846715in}{1.062147in}}%
\pgfpathlineto{\pgfqpoint{1.843696in}{1.056469in}}%
\pgfpathlineto{\pgfqpoint{1.834159in}{1.056680in}}%
\pgfpathlineto{\pgfqpoint{1.834274in}{1.060508in}}%
\pgfpathlineto{\pgfqpoint{1.820896in}{1.060894in}}%
\pgfpathlineto{\pgfqpoint{1.821187in}{1.069825in}}%
\pgfpathlineto{\pgfqpoint{1.820190in}{1.081544in}}%
\pgfpathlineto{\pgfqpoint{1.820514in}{1.104792in}}%
\pgfpathlineto{\pgfqpoint{1.849007in}{1.104433in}}%
\pgfpathlineto{\pgfqpoint{1.849349in}{1.129915in}}%
\pgfpathlineto{\pgfqpoint{1.872212in}{1.129384in}}%
\pgfpathlineto{\pgfqpoint{1.872048in}{1.113887in}}%
\pgfpathlineto{\pgfqpoint{1.871618in}{1.080502in}}%
\pgfpathlineto{\pgfqpoint{1.872670in}{1.072352in}}%
\pgfpathclose%
\pgfusepath{fill}%
\end{pgfscope}%
\begin{pgfscope}%
\pgfpathrectangle{\pgfqpoint{0.100000in}{0.100000in}}{\pgfqpoint{3.007045in}{1.925000in}}%
\pgfusepath{clip}%
\pgfsetbuttcap%
\pgfsetmiterjoin%
\definecolor{currentfill}{rgb}{0.567059,0.763306,0.871019}%
\pgfsetfillcolor{currentfill}%
\pgfsetlinewidth{0.000000pt}%
\definecolor{currentstroke}{rgb}{0.000000,0.000000,0.000000}%
\pgfsetstrokecolor{currentstroke}%
\pgfsetstrokeopacity{0.000000}%
\pgfsetdash{}{0pt}%
\pgfpathmoveto{\pgfqpoint{0.433096in}{0.983854in}}%
\pgfpathlineto{\pgfqpoint{0.425437in}{0.987474in}}%
\pgfpathlineto{\pgfqpoint{0.427947in}{0.992997in}}%
\pgfpathlineto{\pgfqpoint{0.436545in}{0.987083in}}%
\pgfpathlineto{\pgfqpoint{0.433096in}{0.983854in}}%
\pgfpathclose%
\pgfusepath{fill}%
\end{pgfscope}%
\begin{pgfscope}%
\pgfpathrectangle{\pgfqpoint{0.100000in}{0.100000in}}{\pgfqpoint{3.007045in}{1.925000in}}%
\pgfusepath{clip}%
\pgfsetbuttcap%
\pgfsetmiterjoin%
\definecolor{currentfill}{rgb}{0.567059,0.763306,0.871019}%
\pgfsetfillcolor{currentfill}%
\pgfsetlinewidth{0.000000pt}%
\definecolor{currentstroke}{rgb}{0.000000,0.000000,0.000000}%
\pgfsetstrokecolor{currentstroke}%
\pgfsetstrokeopacity{0.000000}%
\pgfsetdash{}{0pt}%
\pgfpathmoveto{\pgfqpoint{0.411502in}{0.985534in}}%
\pgfpathlineto{\pgfqpoint{0.407141in}{0.993168in}}%
\pgfpathlineto{\pgfqpoint{0.412033in}{0.994414in}}%
\pgfpathlineto{\pgfqpoint{0.419874in}{0.989359in}}%
\pgfpathlineto{\pgfqpoint{0.418366in}{0.986829in}}%
\pgfpathlineto{\pgfqpoint{0.411502in}{0.985534in}}%
\pgfpathclose%
\pgfusepath{fill}%
\end{pgfscope}%
\begin{pgfscope}%
\pgfpathrectangle{\pgfqpoint{0.100000in}{0.100000in}}{\pgfqpoint{3.007045in}{1.925000in}}%
\pgfusepath{clip}%
\pgfsetbuttcap%
\pgfsetmiterjoin%
\definecolor{currentfill}{rgb}{0.567059,0.763306,0.871019}%
\pgfsetfillcolor{currentfill}%
\pgfsetlinewidth{0.000000pt}%
\definecolor{currentstroke}{rgb}{0.000000,0.000000,0.000000}%
\pgfsetstrokecolor{currentstroke}%
\pgfsetstrokeopacity{0.000000}%
\pgfsetdash{}{0pt}%
\pgfpathmoveto{\pgfqpoint{0.464444in}{1.039777in}}%
\pgfpathlineto{\pgfqpoint{0.456601in}{1.011141in}}%
\pgfpathlineto{\pgfqpoint{0.453370in}{1.006642in}}%
\pgfpathlineto{\pgfqpoint{0.447005in}{1.011573in}}%
\pgfpathlineto{\pgfqpoint{0.440321in}{1.011780in}}%
\pgfpathlineto{\pgfqpoint{0.433003in}{1.014607in}}%
\pgfpathlineto{\pgfqpoint{0.427568in}{1.019544in}}%
\pgfpathlineto{\pgfqpoint{0.416032in}{1.023886in}}%
\pgfpathlineto{\pgfqpoint{0.403356in}{1.025371in}}%
\pgfpathlineto{\pgfqpoint{0.401817in}{1.031469in}}%
\pgfpathlineto{\pgfqpoint{0.395689in}{1.036947in}}%
\pgfpathlineto{\pgfqpoint{0.400361in}{1.044762in}}%
\pgfpathlineto{\pgfqpoint{0.399420in}{1.048155in}}%
\pgfpathlineto{\pgfqpoint{0.402613in}{1.053373in}}%
\pgfpathlineto{\pgfqpoint{0.400409in}{1.057912in}}%
\pgfpathlineto{\pgfqpoint{0.403882in}{1.063816in}}%
\pgfpathlineto{\pgfqpoint{0.405983in}{1.073145in}}%
\pgfpathlineto{\pgfqpoint{0.400206in}{1.075730in}}%
\pgfpathlineto{\pgfqpoint{0.395182in}{1.083867in}}%
\pgfpathlineto{\pgfqpoint{0.399136in}{1.090944in}}%
\pgfpathlineto{\pgfqpoint{0.398305in}{1.096127in}}%
\pgfpathlineto{\pgfqpoint{0.393442in}{1.098296in}}%
\pgfpathlineto{\pgfqpoint{0.388518in}{1.111365in}}%
\pgfpathlineto{\pgfqpoint{0.382876in}{1.115594in}}%
\pgfpathlineto{\pgfqpoint{0.382061in}{1.124859in}}%
\pgfpathlineto{\pgfqpoint{0.406868in}{1.117231in}}%
\pgfpathlineto{\pgfqpoint{0.440258in}{1.107424in}}%
\pgfpathlineto{\pgfqpoint{0.441277in}{1.107129in}}%
\pgfpathlineto{\pgfqpoint{0.438131in}{1.096030in}}%
\pgfpathlineto{\pgfqpoint{0.443707in}{1.094459in}}%
\pgfpathlineto{\pgfqpoint{0.442087in}{1.088884in}}%
\pgfpathlineto{\pgfqpoint{0.445664in}{1.083905in}}%
\pgfpathlineto{\pgfqpoint{0.451207in}{1.080317in}}%
\pgfpathlineto{\pgfqpoint{0.449631in}{1.074733in}}%
\pgfpathlineto{\pgfqpoint{0.453296in}{1.073681in}}%
\pgfpathlineto{\pgfqpoint{0.451737in}{1.068135in}}%
\pgfpathlineto{\pgfqpoint{0.459131in}{1.065999in}}%
\pgfpathlineto{\pgfqpoint{0.457558in}{1.060425in}}%
\pgfpathlineto{\pgfqpoint{0.463546in}{1.059117in}}%
\pgfpathlineto{\pgfqpoint{0.461538in}{1.053350in}}%
\pgfpathlineto{\pgfqpoint{0.465944in}{1.051393in}}%
\pgfpathlineto{\pgfqpoint{0.464444in}{1.039777in}}%
\pgfpathclose%
\pgfusepath{fill}%
\end{pgfscope}%
\begin{pgfscope}%
\pgfpathrectangle{\pgfqpoint{0.100000in}{0.100000in}}{\pgfqpoint{3.007045in}{1.925000in}}%
\pgfusepath{clip}%
\pgfsetbuttcap%
\pgfsetmiterjoin%
\definecolor{currentfill}{rgb}{0.735871,0.841569,0.923045}%
\pgfsetfillcolor{currentfill}%
\pgfsetlinewidth{0.000000pt}%
\definecolor{currentstroke}{rgb}{0.000000,0.000000,0.000000}%
\pgfsetstrokecolor{currentstroke}%
\pgfsetstrokeopacity{0.000000}%
\pgfsetdash{}{0pt}%
\pgfpathmoveto{\pgfqpoint{2.825651in}{1.545463in}}%
\pgfpathlineto{\pgfqpoint{2.807383in}{1.541913in}}%
\pgfpathlineto{\pgfqpoint{2.803917in}{1.557332in}}%
\pgfpathlineto{\pgfqpoint{2.794553in}{1.564410in}}%
\pgfpathlineto{\pgfqpoint{2.793875in}{1.572201in}}%
\pgfpathlineto{\pgfqpoint{2.787488in}{1.586749in}}%
\pgfpathlineto{\pgfqpoint{2.787394in}{1.596652in}}%
\pgfpathlineto{\pgfqpoint{2.789757in}{1.602899in}}%
\pgfpathlineto{\pgfqpoint{2.787188in}{1.608949in}}%
\pgfpathlineto{\pgfqpoint{2.786419in}{1.617481in}}%
\pgfpathlineto{\pgfqpoint{2.780749in}{1.624141in}}%
\pgfpathlineto{\pgfqpoint{2.780372in}{1.636703in}}%
\pgfpathlineto{\pgfqpoint{2.777816in}{1.638392in}}%
\pgfpathlineto{\pgfqpoint{2.778407in}{1.644079in}}%
\pgfpathlineto{\pgfqpoint{2.776771in}{1.649915in}}%
\pgfpathlineto{\pgfqpoint{2.809280in}{1.658194in}}%
\pgfpathlineto{\pgfqpoint{2.812663in}{1.658624in}}%
\pgfpathlineto{\pgfqpoint{2.816481in}{1.647652in}}%
\pgfpathlineto{\pgfqpoint{2.815249in}{1.643795in}}%
\pgfpathlineto{\pgfqpoint{2.823531in}{1.641882in}}%
\pgfpathlineto{\pgfqpoint{2.822013in}{1.635422in}}%
\pgfpathlineto{\pgfqpoint{2.827912in}{1.634106in}}%
\pgfpathlineto{\pgfqpoint{2.826283in}{1.627777in}}%
\pgfpathlineto{\pgfqpoint{2.825014in}{1.621968in}}%
\pgfpathlineto{\pgfqpoint{2.819281in}{1.624179in}}%
\pgfpathlineto{\pgfqpoint{2.817950in}{1.618314in}}%
\pgfpathlineto{\pgfqpoint{2.810189in}{1.620019in}}%
\pgfpathlineto{\pgfqpoint{2.808017in}{1.607896in}}%
\pgfpathlineto{\pgfqpoint{2.808838in}{1.600051in}}%
\pgfpathlineto{\pgfqpoint{2.811115in}{1.594372in}}%
\pgfpathlineto{\pgfqpoint{2.819866in}{1.593654in}}%
\pgfpathlineto{\pgfqpoint{2.815683in}{1.583923in}}%
\pgfpathlineto{\pgfqpoint{2.811802in}{1.580671in}}%
\pgfpathlineto{\pgfqpoint{2.821564in}{1.578723in}}%
\pgfpathlineto{\pgfqpoint{2.821240in}{1.572694in}}%
\pgfpathlineto{\pgfqpoint{2.827305in}{1.571893in}}%
\pgfpathlineto{\pgfqpoint{2.826130in}{1.558806in}}%
\pgfpathlineto{\pgfqpoint{2.829522in}{1.557864in}}%
\pgfpathlineto{\pgfqpoint{2.829267in}{1.550294in}}%
\pgfpathlineto{\pgfqpoint{2.825651in}{1.545463in}}%
\pgfpathclose%
\pgfusepath{fill}%
\end{pgfscope}%
\begin{pgfscope}%
\pgfpathrectangle{\pgfqpoint{0.100000in}{0.100000in}}{\pgfqpoint{3.007045in}{1.925000in}}%
\pgfusepath{clip}%
\pgfsetbuttcap%
\pgfsetmiterjoin%
\definecolor{currentfill}{rgb}{0.401953,0.670296,0.832326}%
\pgfsetfillcolor{currentfill}%
\pgfsetlinewidth{0.000000pt}%
\definecolor{currentstroke}{rgb}{0.000000,0.000000,0.000000}%
\pgfsetstrokecolor{currentstroke}%
\pgfsetstrokeopacity{0.000000}%
\pgfsetdash{}{0pt}%
\pgfpathmoveto{\pgfqpoint{1.969121in}{0.688091in}}%
\pgfpathlineto{\pgfqpoint{1.966832in}{0.693993in}}%
\pgfpathlineto{\pgfqpoint{1.970011in}{0.695703in}}%
\pgfpathlineto{\pgfqpoint{1.968808in}{0.707431in}}%
\pgfpathlineto{\pgfqpoint{1.970224in}{0.710328in}}%
\pgfpathlineto{\pgfqpoint{1.988311in}{0.711075in}}%
\pgfpathlineto{\pgfqpoint{1.993233in}{0.713829in}}%
\pgfpathlineto{\pgfqpoint{1.997671in}{0.713980in}}%
\pgfpathlineto{\pgfqpoint{1.998667in}{0.705092in}}%
\pgfpathlineto{\pgfqpoint{2.003803in}{0.705398in}}%
\pgfpathlineto{\pgfqpoint{2.003446in}{0.714212in}}%
\pgfpathlineto{\pgfqpoint{2.010976in}{0.714602in}}%
\pgfpathlineto{\pgfqpoint{2.011147in}{0.717467in}}%
\pgfpathlineto{\pgfqpoint{2.020779in}{0.710545in}}%
\pgfpathlineto{\pgfqpoint{2.018058in}{0.708669in}}%
\pgfpathlineto{\pgfqpoint{2.016613in}{0.701553in}}%
\pgfpathlineto{\pgfqpoint{2.012854in}{0.699525in}}%
\pgfpathlineto{\pgfqpoint{2.011898in}{0.691368in}}%
\pgfpathlineto{\pgfqpoint{2.013470in}{0.675004in}}%
\pgfpathlineto{\pgfqpoint{2.013078in}{0.656435in}}%
\pgfpathlineto{\pgfqpoint{2.013321in}{0.650638in}}%
\pgfpathlineto{\pgfqpoint{1.989789in}{0.649517in}}%
\pgfpathlineto{\pgfqpoint{1.981608in}{0.653144in}}%
\pgfpathlineto{\pgfqpoint{1.978327in}{0.664514in}}%
\pgfpathlineto{\pgfqpoint{1.986197in}{0.669162in}}%
\pgfpathlineto{\pgfqpoint{1.986876in}{0.672737in}}%
\pgfpathlineto{\pgfqpoint{1.993029in}{0.676683in}}%
\pgfpathlineto{\pgfqpoint{1.995075in}{0.682326in}}%
\pgfpathlineto{\pgfqpoint{1.987667in}{0.684255in}}%
\pgfpathlineto{\pgfqpoint{1.989120in}{0.691955in}}%
\pgfpathlineto{\pgfqpoint{1.978768in}{0.688413in}}%
\pgfpathlineto{\pgfqpoint{1.969121in}{0.688091in}}%
\pgfpathclose%
\pgfusepath{fill}%
\end{pgfscope}%
\begin{pgfscope}%
\pgfpathrectangle{\pgfqpoint{0.100000in}{0.100000in}}{\pgfqpoint{3.007045in}{1.925000in}}%
\pgfusepath{clip}%
\pgfsetbuttcap%
\pgfsetmiterjoin%
\definecolor{currentfill}{rgb}{0.059439,0.353572,0.639831}%
\pgfsetfillcolor{currentfill}%
\pgfsetlinewidth{0.000000pt}%
\definecolor{currentstroke}{rgb}{0.000000,0.000000,0.000000}%
\pgfsetstrokecolor{currentstroke}%
\pgfsetstrokeopacity{0.000000}%
\pgfsetdash{}{0pt}%
\pgfpathmoveto{\pgfqpoint{0.441277in}{1.107129in}}%
\pgfpathlineto{\pgfqpoint{0.440258in}{1.107424in}}%
\pgfpathlineto{\pgfqpoint{0.435906in}{1.120684in}}%
\pgfpathlineto{\pgfqpoint{0.432260in}{1.122040in}}%
\pgfpathlineto{\pgfqpoint{0.428309in}{1.126977in}}%
\pgfpathlineto{\pgfqpoint{0.423734in}{1.137658in}}%
\pgfpathlineto{\pgfqpoint{0.426809in}{1.147328in}}%
\pgfpathlineto{\pgfqpoint{0.430464in}{1.147237in}}%
\pgfpathlineto{\pgfqpoint{0.433028in}{1.157462in}}%
\pgfpathlineto{\pgfqpoint{0.421673in}{1.178083in}}%
\pgfpathlineto{\pgfqpoint{0.450589in}{1.194118in}}%
\pgfpathlineto{\pgfqpoint{0.461047in}{1.194974in}}%
\pgfpathlineto{\pgfqpoint{0.473609in}{1.193492in}}%
\pgfpathlineto{\pgfqpoint{0.492430in}{1.204219in}}%
\pgfpathlineto{\pgfqpoint{0.498018in}{1.202710in}}%
\pgfpathlineto{\pgfqpoint{0.502749in}{1.206653in}}%
\pgfpathlineto{\pgfqpoint{0.503936in}{1.210812in}}%
\pgfpathlineto{\pgfqpoint{0.521553in}{1.220832in}}%
\pgfpathlineto{\pgfqpoint{0.522889in}{1.217818in}}%
\pgfpathlineto{\pgfqpoint{0.530069in}{1.215499in}}%
\pgfpathlineto{\pgfqpoint{0.532191in}{1.211551in}}%
\pgfpathlineto{\pgfqpoint{0.532644in}{1.204682in}}%
\pgfpathlineto{\pgfqpoint{0.537345in}{1.200997in}}%
\pgfpathlineto{\pgfqpoint{0.539193in}{1.195591in}}%
\pgfpathlineto{\pgfqpoint{0.543469in}{1.191912in}}%
\pgfpathlineto{\pgfqpoint{0.540581in}{1.185427in}}%
\pgfpathlineto{\pgfqpoint{0.543712in}{1.184135in}}%
\pgfpathlineto{\pgfqpoint{0.543660in}{1.171437in}}%
\pgfpathlineto{\pgfqpoint{0.546742in}{1.170366in}}%
\pgfpathlineto{\pgfqpoint{0.552868in}{1.161985in}}%
\pgfpathlineto{\pgfqpoint{0.554321in}{1.151293in}}%
\pgfpathlineto{\pgfqpoint{0.551954in}{1.148285in}}%
\pgfpathlineto{\pgfqpoint{0.551392in}{1.137382in}}%
\pgfpathlineto{\pgfqpoint{0.554104in}{1.127797in}}%
\pgfpathlineto{\pgfqpoint{0.553369in}{1.124234in}}%
\pgfpathlineto{\pgfqpoint{0.558630in}{1.113262in}}%
\pgfpathlineto{\pgfqpoint{0.556496in}{1.110688in}}%
\pgfpathlineto{\pgfqpoint{0.557824in}{1.101480in}}%
\pgfpathlineto{\pgfqpoint{0.556142in}{1.096690in}}%
\pgfpathlineto{\pgfqpoint{0.557260in}{1.082859in}}%
\pgfpathlineto{\pgfqpoint{0.554172in}{1.076466in}}%
\pgfpathlineto{\pgfqpoint{0.522151in}{1.084892in}}%
\pgfpathlineto{\pgfqpoint{0.475084in}{1.097667in}}%
\pgfpathlineto{\pgfqpoint{0.441277in}{1.107129in}}%
\pgfpathclose%
\pgfusepath{fill}%
\end{pgfscope}%
\begin{pgfscope}%
\pgfpathrectangle{\pgfqpoint{0.100000in}{0.100000in}}{\pgfqpoint{3.007045in}{1.925000in}}%
\pgfusepath{clip}%
\pgfsetbuttcap%
\pgfsetmiterjoin%
\definecolor{currentfill}{rgb}{0.417086,0.680631,0.838231}%
\pgfsetfillcolor{currentfill}%
\pgfsetlinewidth{0.000000pt}%
\definecolor{currentstroke}{rgb}{0.000000,0.000000,0.000000}%
\pgfsetstrokecolor{currentstroke}%
\pgfsetstrokeopacity{0.000000}%
\pgfsetdash{}{0pt}%
\pgfpathmoveto{\pgfqpoint{0.856506in}{1.901711in}}%
\pgfpathlineto{\pgfqpoint{0.892466in}{1.893289in}}%
\pgfpathlineto{\pgfqpoint{0.913025in}{1.888772in}}%
\pgfpathlineto{\pgfqpoint{0.912328in}{1.878097in}}%
\pgfpathlineto{\pgfqpoint{0.913309in}{1.873552in}}%
\pgfpathlineto{\pgfqpoint{0.910046in}{1.867809in}}%
\pgfpathlineto{\pgfqpoint{0.901194in}{1.867977in}}%
\pgfpathlineto{\pgfqpoint{0.901644in}{1.862062in}}%
\pgfpathlineto{\pgfqpoint{0.896692in}{1.839394in}}%
\pgfpathlineto{\pgfqpoint{0.889217in}{1.841046in}}%
\pgfpathlineto{\pgfqpoint{0.886583in}{1.827387in}}%
\pgfpathlineto{\pgfqpoint{0.880330in}{1.828801in}}%
\pgfpathlineto{\pgfqpoint{0.877958in}{1.823506in}}%
\pgfpathlineto{\pgfqpoint{0.873246in}{1.821840in}}%
\pgfpathlineto{\pgfqpoint{0.862566in}{1.825406in}}%
\pgfpathlineto{\pgfqpoint{0.863113in}{1.830232in}}%
\pgfpathlineto{\pgfqpoint{0.860078in}{1.837613in}}%
\pgfpathlineto{\pgfqpoint{0.859696in}{1.849285in}}%
\pgfpathlineto{\pgfqpoint{0.845385in}{1.846522in}}%
\pgfpathlineto{\pgfqpoint{0.844425in}{1.850601in}}%
\pgfpathlineto{\pgfqpoint{0.848850in}{1.869247in}}%
\pgfpathlineto{\pgfqpoint{0.856506in}{1.901711in}}%
\pgfpathclose%
\pgfusepath{fill}%
\end{pgfscope}%
\begin{pgfscope}%
\pgfpathrectangle{\pgfqpoint{0.100000in}{0.100000in}}{\pgfqpoint{3.007045in}{1.925000in}}%
\pgfusepath{clip}%
\pgfsetbuttcap%
\pgfsetmiterjoin%
\definecolor{currentfill}{rgb}{0.386820,0.659962,0.826421}%
\pgfsetfillcolor{currentfill}%
\pgfsetlinewidth{0.000000pt}%
\definecolor{currentstroke}{rgb}{0.000000,0.000000,0.000000}%
\pgfsetstrokecolor{currentstroke}%
\pgfsetstrokeopacity{0.000000}%
\pgfsetdash{}{0pt}%
\pgfpathmoveto{\pgfqpoint{2.113075in}{0.848655in}}%
\pgfpathlineto{\pgfqpoint{2.112919in}{0.865058in}}%
\pgfpathlineto{\pgfqpoint{2.110591in}{0.871650in}}%
\pgfpathlineto{\pgfqpoint{2.106779in}{0.871379in}}%
\pgfpathlineto{\pgfqpoint{2.106178in}{0.880380in}}%
\pgfpathlineto{\pgfqpoint{2.108156in}{0.880516in}}%
\pgfpathlineto{\pgfqpoint{2.107278in}{0.897155in}}%
\pgfpathlineto{\pgfqpoint{2.110613in}{0.897382in}}%
\pgfpathlineto{\pgfqpoint{2.114334in}{0.905115in}}%
\pgfpathlineto{\pgfqpoint{2.117945in}{0.906798in}}%
\pgfpathlineto{\pgfqpoint{2.129176in}{0.907553in}}%
\pgfpathlineto{\pgfqpoint{2.129095in}{0.910042in}}%
\pgfpathlineto{\pgfqpoint{2.142348in}{0.908782in}}%
\pgfpathlineto{\pgfqpoint{2.147248in}{0.909640in}}%
\pgfpathlineto{\pgfqpoint{2.150055in}{0.903459in}}%
\pgfpathlineto{\pgfqpoint{2.151351in}{0.884373in}}%
\pgfpathlineto{\pgfqpoint{2.139583in}{0.883616in}}%
\pgfpathlineto{\pgfqpoint{2.145830in}{0.876417in}}%
\pgfpathlineto{\pgfqpoint{2.144807in}{0.847920in}}%
\pgfpathlineto{\pgfqpoint{2.123839in}{0.846519in}}%
\pgfpathlineto{\pgfqpoint{2.123649in}{0.849414in}}%
\pgfpathlineto{\pgfqpoint{2.113075in}{0.848655in}}%
\pgfpathclose%
\pgfusepath{fill}%
\end{pgfscope}%
\begin{pgfscope}%
\pgfpathrectangle{\pgfqpoint{0.100000in}{0.100000in}}{\pgfqpoint{3.007045in}{1.925000in}}%
\pgfusepath{clip}%
\pgfsetbuttcap%
\pgfsetmiterjoin%
\definecolor{currentfill}{rgb}{0.888658,0.933133,0.974410}%
\pgfsetfillcolor{currentfill}%
\pgfsetlinewidth{0.000000pt}%
\definecolor{currentstroke}{rgb}{0.000000,0.000000,0.000000}%
\pgfsetstrokecolor{currentstroke}%
\pgfsetstrokeopacity{0.000000}%
\pgfsetdash{}{0pt}%
\pgfpathmoveto{\pgfqpoint{0.556363in}{1.910254in}}%
\pgfpathlineto{\pgfqpoint{0.556850in}{1.906039in}}%
\pgfpathlineto{\pgfqpoint{0.549783in}{1.902072in}}%
\pgfpathlineto{\pgfqpoint{0.546025in}{1.896043in}}%
\pgfpathlineto{\pgfqpoint{0.536363in}{1.893635in}}%
\pgfpathlineto{\pgfqpoint{0.512685in}{1.901288in}}%
\pgfpathlineto{\pgfqpoint{0.510848in}{1.895622in}}%
\pgfpathlineto{\pgfqpoint{0.487655in}{1.903257in}}%
\pgfpathlineto{\pgfqpoint{0.474818in}{1.908630in}}%
\pgfpathlineto{\pgfqpoint{0.476195in}{1.922659in}}%
\pgfpathlineto{\pgfqpoint{0.474995in}{1.928796in}}%
\pgfpathlineto{\pgfqpoint{0.470655in}{1.935298in}}%
\pgfpathlineto{\pgfqpoint{0.470537in}{1.940585in}}%
\pgfpathlineto{\pgfqpoint{0.473848in}{1.954693in}}%
\pgfpathlineto{\pgfqpoint{0.481028in}{1.966614in}}%
\pgfpathlineto{\pgfqpoint{0.489160in}{1.956885in}}%
\pgfpathlineto{\pgfqpoint{0.494976in}{1.953083in}}%
\pgfpathlineto{\pgfqpoint{0.502419in}{1.944351in}}%
\pgfpathlineto{\pgfqpoint{0.509754in}{1.940761in}}%
\pgfpathlineto{\pgfqpoint{0.515790in}{1.939825in}}%
\pgfpathlineto{\pgfqpoint{0.520125in}{1.935808in}}%
\pgfpathlineto{\pgfqpoint{0.533931in}{1.930174in}}%
\pgfpathlineto{\pgfqpoint{0.539348in}{1.932466in}}%
\pgfpathlineto{\pgfqpoint{0.542272in}{1.925230in}}%
\pgfpathlineto{\pgfqpoint{0.555512in}{1.925149in}}%
\pgfpathlineto{\pgfqpoint{0.557364in}{1.921639in}}%
\pgfpathlineto{\pgfqpoint{0.555607in}{1.915548in}}%
\pgfpathlineto{\pgfqpoint{0.556363in}{1.910254in}}%
\pgfpathclose%
\pgfusepath{fill}%
\end{pgfscope}%
\begin{pgfscope}%
\pgfpathrectangle{\pgfqpoint{0.100000in}{0.100000in}}{\pgfqpoint{3.007045in}{1.925000in}}%
\pgfusepath{clip}%
\pgfsetbuttcap%
\pgfsetmiterjoin%
\definecolor{currentfill}{rgb}{0.396909,0.666851,0.830358}%
\pgfsetfillcolor{currentfill}%
\pgfsetlinewidth{0.000000pt}%
\definecolor{currentstroke}{rgb}{0.000000,0.000000,0.000000}%
\pgfsetstrokecolor{currentstroke}%
\pgfsetstrokeopacity{0.000000}%
\pgfsetdash{}{0pt}%
\pgfpathmoveto{\pgfqpoint{2.288393in}{1.199350in}}%
\pgfpathlineto{\pgfqpoint{2.288706in}{1.196382in}}%
\pgfpathlineto{\pgfqpoint{2.264298in}{1.193927in}}%
\pgfpathlineto{\pgfqpoint{2.262289in}{1.211080in}}%
\pgfpathlineto{\pgfqpoint{2.266331in}{1.211577in}}%
\pgfpathlineto{\pgfqpoint{2.265786in}{1.225833in}}%
\pgfpathlineto{\pgfqpoint{2.285375in}{1.228022in}}%
\pgfpathlineto{\pgfqpoint{2.288393in}{1.199350in}}%
\pgfpathclose%
\pgfusepath{fill}%
\end{pgfscope}%
\begin{pgfscope}%
\pgfpathrectangle{\pgfqpoint{0.100000in}{0.100000in}}{\pgfqpoint{3.007045in}{1.925000in}}%
\pgfusepath{clip}%
\pgfsetbuttcap%
\pgfsetmiterjoin%
\definecolor{currentfill}{rgb}{0.466667,0.708189,0.849366}%
\pgfsetfillcolor{currentfill}%
\pgfsetlinewidth{0.000000pt}%
\definecolor{currentstroke}{rgb}{0.000000,0.000000,0.000000}%
\pgfsetstrokecolor{currentstroke}%
\pgfsetstrokeopacity{0.000000}%
\pgfsetdash{}{0pt}%
\pgfpathmoveto{\pgfqpoint{1.905032in}{1.544903in}}%
\pgfpathlineto{\pgfqpoint{1.876967in}{1.544122in}}%
\pgfpathlineto{\pgfqpoint{1.876801in}{1.549513in}}%
\pgfpathlineto{\pgfqpoint{1.881663in}{1.556766in}}%
\pgfpathlineto{\pgfqpoint{1.881731in}{1.559757in}}%
\pgfpathlineto{\pgfqpoint{1.875816in}{1.567788in}}%
\pgfpathlineto{\pgfqpoint{1.870534in}{1.568228in}}%
\pgfpathlineto{\pgfqpoint{1.872272in}{1.578521in}}%
\pgfpathlineto{\pgfqpoint{1.874776in}{1.580802in}}%
\pgfpathlineto{\pgfqpoint{1.877917in}{1.589345in}}%
\pgfpathlineto{\pgfqpoint{1.881222in}{1.592179in}}%
\pgfpathlineto{\pgfqpoint{1.894425in}{1.598113in}}%
\pgfpathlineto{\pgfqpoint{1.896908in}{1.602038in}}%
\pgfpathlineto{\pgfqpoint{1.896759in}{1.607543in}}%
\pgfpathlineto{\pgfqpoint{1.907975in}{1.607918in}}%
\pgfpathlineto{\pgfqpoint{1.909099in}{1.596409in}}%
\pgfpathlineto{\pgfqpoint{1.909946in}{1.573607in}}%
\pgfpathlineto{\pgfqpoint{1.904215in}{1.573411in}}%
\pgfpathlineto{\pgfqpoint{1.905032in}{1.544903in}}%
\pgfpathclose%
\pgfusepath{fill}%
\end{pgfscope}%
\begin{pgfscope}%
\pgfpathrectangle{\pgfqpoint{0.100000in}{0.100000in}}{\pgfqpoint{3.007045in}{1.925000in}}%
\pgfusepath{clip}%
\pgfsetbuttcap%
\pgfsetmiterjoin%
\definecolor{currentfill}{rgb}{0.270804,0.580730,0.781146}%
\pgfsetfillcolor{currentfill}%
\pgfsetlinewidth{0.000000pt}%
\definecolor{currentstroke}{rgb}{0.000000,0.000000,0.000000}%
\pgfsetstrokecolor{currentstroke}%
\pgfsetstrokeopacity{0.000000}%
\pgfsetdash{}{0pt}%
\pgfpathmoveto{\pgfqpoint{2.052237in}{0.736873in}}%
\pgfpathlineto{\pgfqpoint{2.048959in}{0.728755in}}%
\pgfpathlineto{\pgfqpoint{2.040123in}{0.721921in}}%
\pgfpathlineto{\pgfqpoint{2.027190in}{0.716066in}}%
\pgfpathlineto{\pgfqpoint{2.020779in}{0.710545in}}%
\pgfpathlineto{\pgfqpoint{2.011147in}{0.717467in}}%
\pgfpathlineto{\pgfqpoint{2.011096in}{0.729886in}}%
\pgfpathlineto{\pgfqpoint{2.008261in}{0.731716in}}%
\pgfpathlineto{\pgfqpoint{2.013725in}{0.737732in}}%
\pgfpathlineto{\pgfqpoint{2.013096in}{0.749239in}}%
\pgfpathlineto{\pgfqpoint{2.009327in}{0.760562in}}%
\pgfpathlineto{\pgfqpoint{2.018223in}{0.761023in}}%
\pgfpathlineto{\pgfqpoint{2.018021in}{0.764899in}}%
\pgfpathlineto{\pgfqpoint{2.023674in}{0.765200in}}%
\pgfpathlineto{\pgfqpoint{2.023847in}{0.761087in}}%
\pgfpathlineto{\pgfqpoint{2.030481in}{0.763778in}}%
\pgfpathlineto{\pgfqpoint{2.031123in}{0.752696in}}%
\pgfpathlineto{\pgfqpoint{2.025625in}{0.749814in}}%
\pgfpathlineto{\pgfqpoint{2.031102in}{0.744404in}}%
\pgfpathlineto{\pgfqpoint{2.052237in}{0.736873in}}%
\pgfpathclose%
\pgfusepath{fill}%
\end{pgfscope}%
\begin{pgfscope}%
\pgfpathrectangle{\pgfqpoint{0.100000in}{0.100000in}}{\pgfqpoint{3.007045in}{1.925000in}}%
\pgfusepath{clip}%
\pgfsetbuttcap%
\pgfsetmiterjoin%
\definecolor{currentfill}{rgb}{0.381776,0.656517,0.824452}%
\pgfsetfillcolor{currentfill}%
\pgfsetlinewidth{0.000000pt}%
\definecolor{currentstroke}{rgb}{0.000000,0.000000,0.000000}%
\pgfsetstrokecolor{currentstroke}%
\pgfsetstrokeopacity{0.000000}%
\pgfsetdash{}{0pt}%
\pgfpathmoveto{\pgfqpoint{1.827033in}{1.605997in}}%
\pgfpathlineto{\pgfqpoint{1.827186in}{1.594478in}}%
\pgfpathlineto{\pgfqpoint{1.829097in}{1.594472in}}%
\pgfpathlineto{\pgfqpoint{1.829469in}{1.583832in}}%
\pgfpathlineto{\pgfqpoint{1.802241in}{1.583541in}}%
\pgfpathlineto{\pgfqpoint{1.806019in}{1.580465in}}%
\pgfpathlineto{\pgfqpoint{1.765864in}{1.580217in}}%
\pgfpathlineto{\pgfqpoint{1.765646in}{1.613993in}}%
\pgfpathlineto{\pgfqpoint{1.764876in}{1.642920in}}%
\pgfpathlineto{\pgfqpoint{1.764918in}{1.648716in}}%
\pgfpathlineto{\pgfqpoint{1.782023in}{1.648669in}}%
\pgfpathlineto{\pgfqpoint{1.782448in}{1.621431in}}%
\pgfpathlineto{\pgfqpoint{1.793675in}{1.615402in}}%
\pgfpathlineto{\pgfqpoint{1.798627in}{1.617392in}}%
\pgfpathlineto{\pgfqpoint{1.802129in}{1.614963in}}%
\pgfpathlineto{\pgfqpoint{1.801088in}{1.605771in}}%
\pgfpathlineto{\pgfqpoint{1.827033in}{1.605997in}}%
\pgfpathclose%
\pgfusepath{fill}%
\end{pgfscope}%
\begin{pgfscope}%
\pgfpathrectangle{\pgfqpoint{0.100000in}{0.100000in}}{\pgfqpoint{3.007045in}{1.925000in}}%
\pgfusepath{clip}%
\pgfsetbuttcap%
\pgfsetmiterjoin%
\definecolor{currentfill}{rgb}{0.321246,0.615179,0.800830}%
\pgfsetfillcolor{currentfill}%
\pgfsetlinewidth{0.000000pt}%
\definecolor{currentstroke}{rgb}{0.000000,0.000000,0.000000}%
\pgfsetstrokecolor{currentstroke}%
\pgfsetstrokeopacity{0.000000}%
\pgfsetdash{}{0pt}%
\pgfpathmoveto{\pgfqpoint{1.876321in}{0.768157in}}%
\pgfpathlineto{\pgfqpoint{1.876274in}{0.763349in}}%
\pgfpathlineto{\pgfqpoint{1.884045in}{0.763316in}}%
\pgfpathlineto{\pgfqpoint{1.884006in}{0.739484in}}%
\pgfpathlineto{\pgfqpoint{1.838686in}{0.738812in}}%
\pgfpathlineto{\pgfqpoint{1.835563in}{0.743537in}}%
\pgfpathlineto{\pgfqpoint{1.836073in}{0.747122in}}%
\pgfpathlineto{\pgfqpoint{1.841164in}{0.748497in}}%
\pgfpathlineto{\pgfqpoint{1.846121in}{0.759629in}}%
\pgfpathlineto{\pgfqpoint{1.842406in}{0.765728in}}%
\pgfpathlineto{\pgfqpoint{1.842812in}{0.769497in}}%
\pgfpathlineto{\pgfqpoint{1.856051in}{0.769311in}}%
\pgfpathlineto{\pgfqpoint{1.856034in}{0.767401in}}%
\pgfpathlineto{\pgfqpoint{1.870545in}{0.767215in}}%
\pgfpathlineto{\pgfqpoint{1.876321in}{0.768157in}}%
\pgfpathclose%
\pgfusepath{fill}%
\end{pgfscope}%
\begin{pgfscope}%
\pgfpathrectangle{\pgfqpoint{0.100000in}{0.100000in}}{\pgfqpoint{3.007045in}{1.925000in}}%
\pgfusepath{clip}%
\pgfsetbuttcap%
\pgfsetmiterjoin%
\definecolor{currentfill}{rgb}{0.585882,0.773641,0.875079}%
\pgfsetfillcolor{currentfill}%
\pgfsetlinewidth{0.000000pt}%
\definecolor{currentstroke}{rgb}{0.000000,0.000000,0.000000}%
\pgfsetstrokecolor{currentstroke}%
\pgfsetstrokeopacity{0.000000}%
\pgfsetdash{}{0pt}%
\pgfpathmoveto{\pgfqpoint{2.061841in}{1.396171in}}%
\pgfpathlineto{\pgfqpoint{2.044679in}{1.394919in}}%
\pgfpathlineto{\pgfqpoint{2.021967in}{1.394387in}}%
\pgfpathlineto{\pgfqpoint{2.022119in}{1.391513in}}%
\pgfpathlineto{\pgfqpoint{1.993663in}{1.389926in}}%
\pgfpathlineto{\pgfqpoint{1.992238in}{1.415509in}}%
\pgfpathlineto{\pgfqpoint{1.986647in}{1.415927in}}%
\pgfpathlineto{\pgfqpoint{1.980936in}{1.413027in}}%
\pgfpathlineto{\pgfqpoint{1.980060in}{1.429597in}}%
\pgfpathlineto{\pgfqpoint{1.979503in}{1.438168in}}%
\pgfpathlineto{\pgfqpoint{1.996750in}{1.439088in}}%
\pgfpathlineto{\pgfqpoint{1.996384in}{1.444829in}}%
\pgfpathlineto{\pgfqpoint{2.021528in}{1.446204in}}%
\pgfpathlineto{\pgfqpoint{2.030394in}{1.446824in}}%
\pgfpathlineto{\pgfqpoint{2.029155in}{1.469240in}}%
\pgfpathlineto{\pgfqpoint{2.049555in}{1.470541in}}%
\pgfpathlineto{\pgfqpoint{2.050014in}{1.463525in}}%
\pgfpathlineto{\pgfqpoint{2.046797in}{1.456021in}}%
\pgfpathlineto{\pgfqpoint{2.047314in}{1.447880in}}%
\pgfpathlineto{\pgfqpoint{2.058743in}{1.447959in}}%
\pgfpathlineto{\pgfqpoint{2.061841in}{1.396171in}}%
\pgfpathclose%
\pgfusepath{fill}%
\end{pgfscope}%
\begin{pgfscope}%
\pgfpathrectangle{\pgfqpoint{0.100000in}{0.100000in}}{\pgfqpoint{3.007045in}{1.925000in}}%
\pgfusepath{clip}%
\pgfsetbuttcap%
\pgfsetmiterjoin%
\definecolor{currentfill}{rgb}{0.485490,0.718524,0.853426}%
\pgfsetfillcolor{currentfill}%
\pgfsetlinewidth{0.000000pt}%
\definecolor{currentstroke}{rgb}{0.000000,0.000000,0.000000}%
\pgfsetstrokecolor{currentstroke}%
\pgfsetstrokeopacity{0.000000}%
\pgfsetdash{}{0pt}%
\pgfpathmoveto{\pgfqpoint{1.082392in}{1.585392in}}%
\pgfpathlineto{\pgfqpoint{1.114575in}{1.580480in}}%
\pgfpathlineto{\pgfqpoint{1.115456in}{1.579920in}}%
\pgfpathlineto{\pgfqpoint{1.154054in}{1.574061in}}%
\pgfpathlineto{\pgfqpoint{1.169718in}{1.571914in}}%
\pgfpathlineto{\pgfqpoint{1.170691in}{1.563235in}}%
\pgfpathlineto{\pgfqpoint{1.174953in}{1.552958in}}%
\pgfpathlineto{\pgfqpoint{1.180687in}{1.550302in}}%
\pgfpathlineto{\pgfqpoint{1.188802in}{1.543168in}}%
\pgfpathlineto{\pgfqpoint{1.193646in}{1.534813in}}%
\pgfpathlineto{\pgfqpoint{1.197692in}{1.530943in}}%
\pgfpathlineto{\pgfqpoint{1.199309in}{1.522550in}}%
\pgfpathlineto{\pgfqpoint{1.197829in}{1.512182in}}%
\pgfpathlineto{\pgfqpoint{1.131931in}{1.521748in}}%
\pgfpathlineto{\pgfqpoint{1.131047in}{1.515926in}}%
\pgfpathlineto{\pgfqpoint{1.119707in}{1.517659in}}%
\pgfpathlineto{\pgfqpoint{1.118845in}{1.511867in}}%
\pgfpathlineto{\pgfqpoint{1.113003in}{1.512751in}}%
\pgfpathlineto{\pgfqpoint{1.112286in}{1.507103in}}%
\pgfpathlineto{\pgfqpoint{1.103807in}{1.508400in}}%
\pgfpathlineto{\pgfqpoint{1.102940in}{1.502630in}}%
\pgfpathlineto{\pgfqpoint{1.092509in}{1.504079in}}%
\pgfpathlineto{\pgfqpoint{1.087640in}{1.512697in}}%
\pgfpathlineto{\pgfqpoint{1.083266in}{1.515471in}}%
\pgfpathlineto{\pgfqpoint{1.078184in}{1.513711in}}%
\pgfpathlineto{\pgfqpoint{1.076162in}{1.510378in}}%
\pgfpathlineto{\pgfqpoint{1.070042in}{1.506895in}}%
\pgfpathlineto{\pgfqpoint{1.067064in}{1.517262in}}%
\pgfpathlineto{\pgfqpoint{1.062108in}{1.518116in}}%
\pgfpathlineto{\pgfqpoint{1.059788in}{1.522455in}}%
\pgfpathlineto{\pgfqpoint{1.061138in}{1.530583in}}%
\pgfpathlineto{\pgfqpoint{1.058870in}{1.534775in}}%
\pgfpathlineto{\pgfqpoint{1.057765in}{1.542345in}}%
\pgfpathlineto{\pgfqpoint{1.053058in}{1.551890in}}%
\pgfpathlineto{\pgfqpoint{1.054900in}{1.557314in}}%
\pgfpathlineto{\pgfqpoint{1.051087in}{1.562499in}}%
\pgfpathlineto{\pgfqpoint{1.037496in}{1.564829in}}%
\pgfpathlineto{\pgfqpoint{1.038445in}{1.570273in}}%
\pgfpathlineto{\pgfqpoint{1.020483in}{1.573439in}}%
\pgfpathlineto{\pgfqpoint{1.024388in}{1.595276in}}%
\pgfpathlineto{\pgfqpoint{1.040422in}{1.591840in}}%
\pgfpathlineto{\pgfqpoint{1.082392in}{1.585392in}}%
\pgfpathclose%
\pgfusepath{fill}%
\end{pgfscope}%
\begin{pgfscope}%
\pgfpathrectangle{\pgfqpoint{0.100000in}{0.100000in}}{\pgfqpoint{3.007045in}{1.925000in}}%
\pgfusepath{clip}%
\pgfsetbuttcap%
\pgfsetmiterjoin%
\definecolor{currentfill}{rgb}{0.381776,0.656517,0.824452}%
\pgfsetfillcolor{currentfill}%
\pgfsetlinewidth{0.000000pt}%
\definecolor{currentstroke}{rgb}{0.000000,0.000000,0.000000}%
\pgfsetstrokecolor{currentstroke}%
\pgfsetstrokeopacity{0.000000}%
\pgfsetdash{}{0pt}%
\pgfpathmoveto{\pgfqpoint{2.143299in}{0.800190in}}%
\pgfpathlineto{\pgfqpoint{2.139121in}{0.799917in}}%
\pgfpathlineto{\pgfqpoint{2.136914in}{0.795797in}}%
\pgfpathlineto{\pgfqpoint{2.129211in}{0.792870in}}%
\pgfpathlineto{\pgfqpoint{2.117772in}{0.794063in}}%
\pgfpathlineto{\pgfqpoint{2.117152in}{0.802691in}}%
\pgfpathlineto{\pgfqpoint{2.105570in}{0.801993in}}%
\pgfpathlineto{\pgfqpoint{2.105804in}{0.798159in}}%
\pgfpathlineto{\pgfqpoint{2.091654in}{0.795102in}}%
\pgfpathlineto{\pgfqpoint{2.074344in}{0.794001in}}%
\pgfpathlineto{\pgfqpoint{2.073753in}{0.803617in}}%
\pgfpathlineto{\pgfqpoint{2.072446in}{0.823053in}}%
\pgfpathlineto{\pgfqpoint{2.086876in}{0.823888in}}%
\pgfpathlineto{\pgfqpoint{2.085917in}{0.838289in}}%
\pgfpathlineto{\pgfqpoint{2.103193in}{0.839402in}}%
\pgfpathlineto{\pgfqpoint{2.110699in}{0.840860in}}%
\pgfpathlineto{\pgfqpoint{2.110261in}{0.847530in}}%
\pgfpathlineto{\pgfqpoint{2.113075in}{0.848655in}}%
\pgfpathlineto{\pgfqpoint{2.123649in}{0.849414in}}%
\pgfpathlineto{\pgfqpoint{2.123839in}{0.846519in}}%
\pgfpathlineto{\pgfqpoint{2.144807in}{0.847920in}}%
\pgfpathlineto{\pgfqpoint{2.144570in}{0.838481in}}%
\pgfpathlineto{\pgfqpoint{2.143299in}{0.800190in}}%
\pgfpathclose%
\pgfusepath{fill}%
\end{pgfscope}%
\begin{pgfscope}%
\pgfpathrectangle{\pgfqpoint{0.100000in}{0.100000in}}{\pgfqpoint{3.007045in}{1.925000in}}%
\pgfusepath{clip}%
\pgfsetbuttcap%
\pgfsetmiterjoin%
\definecolor{currentfill}{rgb}{0.805260,0.878016,0.946851}%
\pgfsetfillcolor{currentfill}%
\pgfsetlinewidth{0.000000pt}%
\definecolor{currentstroke}{rgb}{0.000000,0.000000,0.000000}%
\pgfsetstrokecolor{currentstroke}%
\pgfsetstrokeopacity{0.000000}%
\pgfsetdash{}{0pt}%
\pgfpathmoveto{\pgfqpoint{2.898990in}{1.789167in}}%
\pgfpathlineto{\pgfqpoint{2.897996in}{1.797410in}}%
\pgfpathlineto{\pgfqpoint{2.917566in}{1.855793in}}%
\pgfpathlineto{\pgfqpoint{2.926613in}{1.855322in}}%
\pgfpathlineto{\pgfqpoint{2.928868in}{1.845060in}}%
\pgfpathlineto{\pgfqpoint{2.936844in}{1.841940in}}%
\pgfpathlineto{\pgfqpoint{2.948719in}{1.853007in}}%
\pgfpathlineto{\pgfqpoint{2.957139in}{1.855503in}}%
\pgfpathlineto{\pgfqpoint{2.956602in}{1.859933in}}%
\pgfpathlineto{\pgfqpoint{2.962545in}{1.861779in}}%
\pgfpathlineto{\pgfqpoint{2.977427in}{1.855560in}}%
\pgfpathlineto{\pgfqpoint{2.982323in}{1.850596in}}%
\pgfpathlineto{\pgfqpoint{2.987134in}{1.849352in}}%
\pgfpathlineto{\pgfqpoint{2.996362in}{1.819764in}}%
\pgfpathlineto{\pgfqpoint{3.009567in}{1.778140in}}%
\pgfpathlineto{\pgfqpoint{3.009737in}{1.774049in}}%
\pgfpathlineto{\pgfqpoint{3.013596in}{1.761178in}}%
\pgfpathlineto{\pgfqpoint{3.003662in}{1.755298in}}%
\pgfpathlineto{\pgfqpoint{2.988875in}{1.746178in}}%
\pgfpathlineto{\pgfqpoint{2.987728in}{1.746399in}}%
\pgfpathlineto{\pgfqpoint{2.981003in}{1.769545in}}%
\pgfpathlineto{\pgfqpoint{2.972249in}{1.797222in}}%
\pgfpathlineto{\pgfqpoint{2.955175in}{1.793081in}}%
\pgfpathlineto{\pgfqpoint{2.951724in}{1.804308in}}%
\pgfpathlineto{\pgfqpoint{2.912260in}{1.792925in}}%
\pgfpathlineto{\pgfqpoint{2.898990in}{1.789167in}}%
\pgfpathclose%
\pgfusepath{fill}%
\end{pgfscope}%
\begin{pgfscope}%
\pgfpathrectangle{\pgfqpoint{0.100000in}{0.100000in}}{\pgfqpoint{3.007045in}{1.925000in}}%
\pgfusepath{clip}%
\pgfsetbuttcap%
\pgfsetmiterjoin%
\definecolor{currentfill}{rgb}{0.351511,0.635848,0.812641}%
\pgfsetfillcolor{currentfill}%
\pgfsetlinewidth{0.000000pt}%
\definecolor{currentstroke}{rgb}{0.000000,0.000000,0.000000}%
\pgfsetstrokecolor{currentstroke}%
\pgfsetstrokeopacity{0.000000}%
\pgfsetdash{}{0pt}%
\pgfpathmoveto{\pgfqpoint{2.278491in}{1.292509in}}%
\pgfpathlineto{\pgfqpoint{2.273138in}{1.338797in}}%
\pgfpathlineto{\pgfqpoint{2.292953in}{1.341716in}}%
\pgfpathlineto{\pgfqpoint{2.295342in}{1.329269in}}%
\pgfpathlineto{\pgfqpoint{2.297909in}{1.323832in}}%
\pgfpathlineto{\pgfqpoint{2.303481in}{1.324508in}}%
\pgfpathlineto{\pgfqpoint{2.305531in}{1.307364in}}%
\pgfpathlineto{\pgfqpoint{2.299961in}{1.306677in}}%
\pgfpathlineto{\pgfqpoint{2.301283in}{1.295237in}}%
\pgfpathlineto{\pgfqpoint{2.278491in}{1.292509in}}%
\pgfpathclose%
\pgfusepath{fill}%
\end{pgfscope}%
\begin{pgfscope}%
\pgfpathrectangle{\pgfqpoint{0.100000in}{0.100000in}}{\pgfqpoint{3.007045in}{1.925000in}}%
\pgfusepath{clip}%
\pgfsetbuttcap%
\pgfsetmiterjoin%
\definecolor{currentfill}{rgb}{0.361599,0.642737,0.816578}%
\pgfsetfillcolor{currentfill}%
\pgfsetlinewidth{0.000000pt}%
\definecolor{currentstroke}{rgb}{0.000000,0.000000,0.000000}%
\pgfsetstrokecolor{currentstroke}%
\pgfsetstrokeopacity{0.000000}%
\pgfsetdash{}{0pt}%
\pgfpathmoveto{\pgfqpoint{1.102849in}{0.996413in}}%
\pgfpathlineto{\pgfqpoint{1.100725in}{0.982093in}}%
\pgfpathlineto{\pgfqpoint{1.025816in}{0.993300in}}%
\pgfpathlineto{\pgfqpoint{1.035972in}{1.058230in}}%
\pgfpathlineto{\pgfqpoint{1.070635in}{1.052989in}}%
\pgfpathlineto{\pgfqpoint{1.073504in}{1.056252in}}%
\pgfpathlineto{\pgfqpoint{1.077407in}{1.066910in}}%
\pgfpathlineto{\pgfqpoint{1.083485in}{1.074787in}}%
\pgfpathlineto{\pgfqpoint{1.088472in}{1.075836in}}%
\pgfpathlineto{\pgfqpoint{1.093042in}{1.080708in}}%
\pgfpathlineto{\pgfqpoint{1.094867in}{1.088434in}}%
\pgfpathlineto{\pgfqpoint{1.103247in}{1.096115in}}%
\pgfpathlineto{\pgfqpoint{1.104848in}{1.099688in}}%
\pgfpathlineto{\pgfqpoint{1.112363in}{1.107267in}}%
\pgfpathlineto{\pgfqpoint{1.121625in}{1.109932in}}%
\pgfpathlineto{\pgfqpoint{1.123300in}{1.108295in}}%
\pgfpathlineto{\pgfqpoint{1.124305in}{1.096464in}}%
\pgfpathlineto{\pgfqpoint{1.121195in}{1.073917in}}%
\pgfpathlineto{\pgfqpoint{1.139523in}{1.071438in}}%
\pgfpathlineto{\pgfqpoint{1.139253in}{1.069444in}}%
\pgfpathlineto{\pgfqpoint{1.162663in}{1.066636in}}%
\pgfpathlineto{\pgfqpoint{1.161246in}{1.055657in}}%
\pgfpathlineto{\pgfqpoint{1.169817in}{1.038957in}}%
\pgfpathlineto{\pgfqpoint{1.120573in}{1.045885in}}%
\pgfpathlineto{\pgfqpoint{1.118519in}{1.040536in}}%
\pgfpathlineto{\pgfqpoint{1.108498in}{1.034565in}}%
\pgfpathlineto{\pgfqpoint{1.102849in}{0.996413in}}%
\pgfpathclose%
\pgfusepath{fill}%
\end{pgfscope}%
\begin{pgfscope}%
\pgfpathrectangle{\pgfqpoint{0.100000in}{0.100000in}}{\pgfqpoint{3.007045in}{1.925000in}}%
\pgfusepath{clip}%
\pgfsetbuttcap%
\pgfsetmiterjoin%
\definecolor{currentfill}{rgb}{0.760477,0.852026,0.931657}%
\pgfsetfillcolor{currentfill}%
\pgfsetlinewidth{0.000000pt}%
\definecolor{currentstroke}{rgb}{0.000000,0.000000,0.000000}%
\pgfsetstrokecolor{currentstroke}%
\pgfsetstrokeopacity{0.000000}%
\pgfsetdash{}{0pt}%
\pgfpathmoveto{\pgfqpoint{2.033429in}{1.724806in}}%
\pgfpathlineto{\pgfqpoint{2.028764in}{1.727403in}}%
\pgfpathlineto{\pgfqpoint{2.031261in}{1.731523in}}%
\pgfpathlineto{\pgfqpoint{2.049089in}{1.741926in}}%
\pgfpathlineto{\pgfqpoint{2.055441in}{1.748010in}}%
\pgfpathlineto{\pgfqpoint{2.063105in}{1.751149in}}%
\pgfpathlineto{\pgfqpoint{2.058606in}{1.742131in}}%
\pgfpathlineto{\pgfqpoint{2.053691in}{1.738050in}}%
\pgfpathlineto{\pgfqpoint{2.042387in}{1.732402in}}%
\pgfpathlineto{\pgfqpoint{2.044192in}{1.729855in}}%
\pgfpathlineto{\pgfqpoint{2.033429in}{1.724806in}}%
\pgfpathclose%
\pgfusepath{fill}%
\end{pgfscope}%
\begin{pgfscope}%
\pgfpathrectangle{\pgfqpoint{0.100000in}{0.100000in}}{\pgfqpoint{3.007045in}{1.925000in}}%
\pgfusepath{clip}%
\pgfsetbuttcap%
\pgfsetmiterjoin%
\definecolor{currentfill}{rgb}{0.760477,0.852026,0.931657}%
\pgfsetfillcolor{currentfill}%
\pgfsetlinewidth{0.000000pt}%
\definecolor{currentstroke}{rgb}{0.000000,0.000000,0.000000}%
\pgfsetstrokecolor{currentstroke}%
\pgfsetstrokeopacity{0.000000}%
\pgfsetdash{}{0pt}%
\pgfpathmoveto{\pgfqpoint{2.088047in}{1.667928in}}%
\pgfpathlineto{\pgfqpoint{2.088768in}{1.658202in}}%
\pgfpathlineto{\pgfqpoint{2.091627in}{1.652681in}}%
\pgfpathlineto{\pgfqpoint{2.085981in}{1.652184in}}%
\pgfpathlineto{\pgfqpoint{2.087227in}{1.635040in}}%
\pgfpathlineto{\pgfqpoint{2.047301in}{1.632207in}}%
\pgfpathlineto{\pgfqpoint{2.046234in}{1.649481in}}%
\pgfpathlineto{\pgfqpoint{2.051916in}{1.649810in}}%
\pgfpathlineto{\pgfqpoint{2.048453in}{1.655274in}}%
\pgfpathlineto{\pgfqpoint{2.047218in}{1.673072in}}%
\pgfpathlineto{\pgfqpoint{2.048913in}{1.677724in}}%
\pgfpathlineto{\pgfqpoint{2.058859in}{1.686684in}}%
\pgfpathlineto{\pgfqpoint{2.062270in}{1.688002in}}%
\pgfpathlineto{\pgfqpoint{2.068764in}{1.697183in}}%
\pgfpathlineto{\pgfqpoint{2.077462in}{1.703060in}}%
\pgfpathlineto{\pgfqpoint{2.089417in}{1.705852in}}%
\pgfpathlineto{\pgfqpoint{2.096001in}{1.705831in}}%
\pgfpathlineto{\pgfqpoint{2.100493in}{1.701354in}}%
\pgfpathlineto{\pgfqpoint{2.091665in}{1.700318in}}%
\pgfpathlineto{\pgfqpoint{2.090480in}{1.696264in}}%
\pgfpathlineto{\pgfqpoint{2.085473in}{1.693327in}}%
\pgfpathlineto{\pgfqpoint{2.078290in}{1.686386in}}%
\pgfpathlineto{\pgfqpoint{2.078268in}{1.682642in}}%
\pgfpathlineto{\pgfqpoint{2.069709in}{1.670334in}}%
\pgfpathlineto{\pgfqpoint{2.068597in}{1.661571in}}%
\pgfpathlineto{\pgfqpoint{2.071052in}{1.658913in}}%
\pgfpathlineto{\pgfqpoint{2.078864in}{1.668386in}}%
\pgfpathlineto{\pgfqpoint{2.082700in}{1.671038in}}%
\pgfpathlineto{\pgfqpoint{2.088047in}{1.667928in}}%
\pgfpathclose%
\pgfusepath{fill}%
\end{pgfscope}%
\begin{pgfscope}%
\pgfpathrectangle{\pgfqpoint{0.100000in}{0.100000in}}{\pgfqpoint{3.007045in}{1.925000in}}%
\pgfusepath{clip}%
\pgfsetbuttcap%
\pgfsetmiterjoin%
\definecolor{currentfill}{rgb}{0.391865,0.663406,0.828389}%
\pgfsetfillcolor{currentfill}%
\pgfsetlinewidth{0.000000pt}%
\definecolor{currentstroke}{rgb}{0.000000,0.000000,0.000000}%
\pgfsetstrokecolor{currentstroke}%
\pgfsetstrokeopacity{0.000000}%
\pgfsetdash{}{0pt}%
\pgfpathmoveto{\pgfqpoint{1.251074in}{1.090450in}}%
\pgfpathlineto{\pgfqpoint{1.245702in}{1.097910in}}%
\pgfpathlineto{\pgfqpoint{1.241682in}{1.094643in}}%
\pgfpathlineto{\pgfqpoint{1.232103in}{1.090942in}}%
\pgfpathlineto{\pgfqpoint{1.229068in}{1.091653in}}%
\pgfpathlineto{\pgfqpoint{1.223896in}{1.096979in}}%
\pgfpathlineto{\pgfqpoint{1.220511in}{1.109227in}}%
\pgfpathlineto{\pgfqpoint{1.205700in}{1.130953in}}%
\pgfpathlineto{\pgfqpoint{1.213054in}{1.137702in}}%
\pgfpathlineto{\pgfqpoint{1.213461in}{1.141563in}}%
\pgfpathlineto{\pgfqpoint{1.209760in}{1.146848in}}%
\pgfpathlineto{\pgfqpoint{1.242494in}{1.143237in}}%
\pgfpathlineto{\pgfqpoint{1.245658in}{1.171547in}}%
\pgfpathlineto{\pgfqpoint{1.310620in}{1.164502in}}%
\pgfpathlineto{\pgfqpoint{1.308723in}{1.147516in}}%
\pgfpathlineto{\pgfqpoint{1.306529in}{1.124802in}}%
\pgfpathlineto{\pgfqpoint{1.261017in}{1.129389in}}%
\pgfpathlineto{\pgfqpoint{1.259199in}{1.112273in}}%
\pgfpathlineto{\pgfqpoint{1.253600in}{1.112871in}}%
\pgfpathlineto{\pgfqpoint{1.251074in}{1.090450in}}%
\pgfpathclose%
\pgfusepath{fill}%
\end{pgfscope}%
\begin{pgfscope}%
\pgfpathrectangle{\pgfqpoint{0.100000in}{0.100000in}}{\pgfqpoint{3.007045in}{1.925000in}}%
\pgfusepath{clip}%
\pgfsetbuttcap%
\pgfsetmiterjoin%
\definecolor{currentfill}{rgb}{0.647290,0.803922,0.892042}%
\pgfsetfillcolor{currentfill}%
\pgfsetlinewidth{0.000000pt}%
\definecolor{currentstroke}{rgb}{0.000000,0.000000,0.000000}%
\pgfsetstrokecolor{currentstroke}%
\pgfsetstrokeopacity{0.000000}%
\pgfsetdash{}{0pt}%
\pgfpathmoveto{\pgfqpoint{2.498296in}{1.267386in}}%
\pgfpathlineto{\pgfqpoint{2.480209in}{1.264505in}}%
\pgfpathlineto{\pgfqpoint{2.478385in}{1.281786in}}%
\pgfpathlineto{\pgfqpoint{2.474602in}{1.281312in}}%
\pgfpathlineto{\pgfqpoint{2.474283in}{1.290040in}}%
\pgfpathlineto{\pgfqpoint{2.476757in}{1.293301in}}%
\pgfpathlineto{\pgfqpoint{2.486521in}{1.293715in}}%
\pgfpathlineto{\pgfqpoint{2.493292in}{1.298554in}}%
\pgfpathlineto{\pgfqpoint{2.498296in}{1.267386in}}%
\pgfpathclose%
\pgfusepath{fill}%
\end{pgfscope}%
\begin{pgfscope}%
\pgfpathrectangle{\pgfqpoint{0.100000in}{0.100000in}}{\pgfqpoint{3.007045in}{1.925000in}}%
\pgfusepath{clip}%
\pgfsetbuttcap%
\pgfsetmiterjoin%
\definecolor{currentfill}{rgb}{0.454118,0.701300,0.846659}%
\pgfsetfillcolor{currentfill}%
\pgfsetlinewidth{0.000000pt}%
\definecolor{currentstroke}{rgb}{0.000000,0.000000,0.000000}%
\pgfsetstrokecolor{currentstroke}%
\pgfsetstrokeopacity{0.000000}%
\pgfsetdash{}{0pt}%
\pgfpathmoveto{\pgfqpoint{2.161409in}{0.940303in}}%
\pgfpathlineto{\pgfqpoint{2.160890in}{0.928692in}}%
\pgfpathlineto{\pgfqpoint{2.166375in}{0.924298in}}%
\pgfpathlineto{\pgfqpoint{2.163243in}{0.916924in}}%
\pgfpathlineto{\pgfqpoint{2.149635in}{0.914318in}}%
\pgfpathlineto{\pgfqpoint{2.147248in}{0.909640in}}%
\pgfpathlineto{\pgfqpoint{2.142348in}{0.908782in}}%
\pgfpathlineto{\pgfqpoint{2.129095in}{0.910042in}}%
\pgfpathlineto{\pgfqpoint{2.125902in}{0.914003in}}%
\pgfpathlineto{\pgfqpoint{2.120897in}{0.915606in}}%
\pgfpathlineto{\pgfqpoint{2.114744in}{0.920210in}}%
\pgfpathlineto{\pgfqpoint{2.114171in}{0.933474in}}%
\pgfpathlineto{\pgfqpoint{2.116000in}{0.935950in}}%
\pgfpathlineto{\pgfqpoint{2.136863in}{0.936941in}}%
\pgfpathlineto{\pgfqpoint{2.143491in}{0.939640in}}%
\pgfpathlineto{\pgfqpoint{2.144945in}{0.937174in}}%
\pgfpathlineto{\pgfqpoint{2.151093in}{0.939458in}}%
\pgfpathlineto{\pgfqpoint{2.161409in}{0.940303in}}%
\pgfpathclose%
\pgfusepath{fill}%
\end{pgfscope}%
\begin{pgfscope}%
\pgfpathrectangle{\pgfqpoint{0.100000in}{0.100000in}}{\pgfqpoint{3.007045in}{1.925000in}}%
\pgfusepath{clip}%
\pgfsetbuttcap%
\pgfsetmiterjoin%
\definecolor{currentfill}{rgb}{0.191326,0.505052,0.741699}%
\pgfsetfillcolor{currentfill}%
\pgfsetlinewidth{0.000000pt}%
\definecolor{currentstroke}{rgb}{0.000000,0.000000,0.000000}%
\pgfsetstrokecolor{currentstroke}%
\pgfsetstrokeopacity{0.000000}%
\pgfsetdash{}{0pt}%
\pgfpathmoveto{\pgfqpoint{1.282225in}{0.694354in}}%
\pgfpathlineto{\pgfqpoint{1.265329in}{0.695996in}}%
\pgfpathlineto{\pgfqpoint{1.219256in}{0.700832in}}%
\pgfpathlineto{\pgfqpoint{1.222778in}{0.735139in}}%
\pgfpathlineto{\pgfqpoint{1.194858in}{0.738140in}}%
\pgfpathlineto{\pgfqpoint{1.198272in}{0.767309in}}%
\pgfpathlineto{\pgfqpoint{1.200282in}{0.767082in}}%
\pgfpathlineto{\pgfqpoint{1.201518in}{0.778286in}}%
\pgfpathlineto{\pgfqpoint{1.224219in}{0.776176in}}%
\pgfpathlineto{\pgfqpoint{1.226820in}{0.793144in}}%
\pgfpathlineto{\pgfqpoint{1.231742in}{0.838561in}}%
\pgfpathlineto{\pgfqpoint{1.237556in}{0.837894in}}%
\pgfpathlineto{\pgfqpoint{1.236948in}{0.832126in}}%
\pgfpathlineto{\pgfqpoint{1.271248in}{0.828230in}}%
\pgfpathlineto{\pgfqpoint{1.271855in}{0.834010in}}%
\pgfpathlineto{\pgfqpoint{1.283287in}{0.832832in}}%
\pgfpathlineto{\pgfqpoint{1.288966in}{0.832202in}}%
\pgfpathlineto{\pgfqpoint{1.287264in}{0.814931in}}%
\pgfpathlineto{\pgfqpoint{1.294107in}{0.814272in}}%
\pgfpathlineto{\pgfqpoint{1.292993in}{0.803535in}}%
\pgfpathlineto{\pgfqpoint{1.304431in}{0.802517in}}%
\pgfpathlineto{\pgfqpoint{1.303863in}{0.796794in}}%
\pgfpathlineto{\pgfqpoint{1.292309in}{0.797855in}}%
\pgfpathlineto{\pgfqpoint{1.291172in}{0.786411in}}%
\pgfpathlineto{\pgfqpoint{1.288691in}{0.786623in}}%
\pgfpathlineto{\pgfqpoint{1.285953in}{0.758267in}}%
\pgfpathlineto{\pgfqpoint{1.283283in}{0.758525in}}%
\pgfpathlineto{\pgfqpoint{1.280477in}{0.729323in}}%
\pgfpathlineto{\pgfqpoint{1.285541in}{0.728841in}}%
\pgfpathlineto{\pgfqpoint{1.282225in}{0.694354in}}%
\pgfpathclose%
\pgfusepath{fill}%
\end{pgfscope}%
\begin{pgfscope}%
\pgfpathrectangle{\pgfqpoint{0.100000in}{0.100000in}}{\pgfqpoint{3.007045in}{1.925000in}}%
\pgfusepath{clip}%
\pgfsetbuttcap%
\pgfsetmiterjoin%
\definecolor{currentfill}{rgb}{0.460392,0.704744,0.848012}%
\pgfsetfillcolor{currentfill}%
\pgfsetlinewidth{0.000000pt}%
\definecolor{currentstroke}{rgb}{0.000000,0.000000,0.000000}%
\pgfsetstrokecolor{currentstroke}%
\pgfsetstrokeopacity{0.000000}%
\pgfsetdash{}{0pt}%
\pgfpathmoveto{\pgfqpoint{1.193280in}{1.837329in}}%
\pgfpathlineto{\pgfqpoint{1.189784in}{1.820439in}}%
\pgfpathlineto{\pgfqpoint{1.187177in}{1.820830in}}%
\pgfpathlineto{\pgfqpoint{1.185494in}{1.809331in}}%
\pgfpathlineto{\pgfqpoint{1.182880in}{1.799848in}}%
\pgfpathlineto{\pgfqpoint{1.180237in}{1.801764in}}%
\pgfpathlineto{\pgfqpoint{1.177030in}{1.787023in}}%
\pgfpathlineto{\pgfqpoint{1.174638in}{1.770993in}}%
\pgfpathlineto{\pgfqpoint{1.167096in}{1.773103in}}%
\pgfpathlineto{\pgfqpoint{1.165616in}{1.768587in}}%
\pgfpathlineto{\pgfqpoint{1.154039in}{1.770483in}}%
\pgfpathlineto{\pgfqpoint{1.152002in}{1.758091in}}%
\pgfpathlineto{\pgfqpoint{1.148828in}{1.763186in}}%
\pgfpathlineto{\pgfqpoint{1.141833in}{1.762846in}}%
\pgfpathlineto{\pgfqpoint{1.135562in}{1.764609in}}%
\pgfpathlineto{\pgfqpoint{1.126744in}{1.761331in}}%
\pgfpathlineto{\pgfqpoint{1.123379in}{1.762908in}}%
\pgfpathlineto{\pgfqpoint{1.127771in}{1.788966in}}%
\pgfpathlineto{\pgfqpoint{1.119270in}{1.790180in}}%
\pgfpathlineto{\pgfqpoint{1.120232in}{1.796028in}}%
\pgfpathlineto{\pgfqpoint{1.115258in}{1.796858in}}%
\pgfpathlineto{\pgfqpoint{1.115936in}{1.802454in}}%
\pgfpathlineto{\pgfqpoint{1.081806in}{1.808264in}}%
\pgfpathlineto{\pgfqpoint{1.080804in}{1.802553in}}%
\pgfpathlineto{\pgfqpoint{1.071976in}{1.804099in}}%
\pgfpathlineto{\pgfqpoint{1.070974in}{1.798414in}}%
\pgfpathlineto{\pgfqpoint{1.045450in}{1.802919in}}%
\pgfpathlineto{\pgfqpoint{1.046502in}{1.808652in}}%
\pgfpathlineto{\pgfqpoint{1.052373in}{1.807583in}}%
\pgfpathlineto{\pgfqpoint{1.061807in}{1.858712in}}%
\pgfpathlineto{\pgfqpoint{1.107359in}{1.850859in}}%
\pgfpathlineto{\pgfqpoint{1.160277in}{1.842278in}}%
\pgfpathlineto{\pgfqpoint{1.193280in}{1.837329in}}%
\pgfpathclose%
\pgfusepath{fill}%
\end{pgfscope}%
\begin{pgfscope}%
\pgfpathrectangle{\pgfqpoint{0.100000in}{0.100000in}}{\pgfqpoint{3.007045in}{1.925000in}}%
\pgfusepath{clip}%
\pgfsetbuttcap%
\pgfsetmiterjoin%
\definecolor{currentfill}{rgb}{0.726028,0.837386,0.919600}%
\pgfsetfillcolor{currentfill}%
\pgfsetlinewidth{0.000000pt}%
\definecolor{currentstroke}{rgb}{0.000000,0.000000,0.000000}%
\pgfsetstrokecolor{currentstroke}%
\pgfsetstrokeopacity{0.000000}%
\pgfsetdash{}{0pt}%
\pgfpathmoveto{\pgfqpoint{2.755737in}{0.991248in}}%
\pgfpathlineto{\pgfqpoint{2.755915in}{0.984607in}}%
\pgfpathlineto{\pgfqpoint{2.744917in}{0.981075in}}%
\pgfpathlineto{\pgfqpoint{2.737785in}{0.985361in}}%
\pgfpathlineto{\pgfqpoint{2.724069in}{0.990555in}}%
\pgfpathlineto{\pgfqpoint{2.723536in}{0.997192in}}%
\pgfpathlineto{\pgfqpoint{2.725670in}{0.999941in}}%
\pgfpathlineto{\pgfqpoint{2.722156in}{1.004752in}}%
\pgfpathlineto{\pgfqpoint{2.720462in}{1.011491in}}%
\pgfpathlineto{\pgfqpoint{2.710520in}{1.014800in}}%
\pgfpathlineto{\pgfqpoint{2.709850in}{1.020460in}}%
\pgfpathlineto{\pgfqpoint{2.705428in}{1.026285in}}%
\pgfpathlineto{\pgfqpoint{2.709028in}{1.030469in}}%
\pgfpathlineto{\pgfqpoint{2.715659in}{1.028449in}}%
\pgfpathlineto{\pgfqpoint{2.716657in}{1.024638in}}%
\pgfpathlineto{\pgfqpoint{2.722671in}{1.024298in}}%
\pgfpathlineto{\pgfqpoint{2.728149in}{1.020939in}}%
\pgfpathlineto{\pgfqpoint{2.731963in}{1.024371in}}%
\pgfpathlineto{\pgfqpoint{2.735129in}{1.019199in}}%
\pgfpathlineto{\pgfqpoint{2.738228in}{1.024682in}}%
\pgfpathlineto{\pgfqpoint{2.741385in}{1.024328in}}%
\pgfpathlineto{\pgfqpoint{2.743345in}{1.029859in}}%
\pgfpathlineto{\pgfqpoint{2.751811in}{1.031450in}}%
\pgfpathlineto{\pgfqpoint{2.758023in}{1.035375in}}%
\pgfpathlineto{\pgfqpoint{2.761386in}{1.033527in}}%
\pgfpathlineto{\pgfqpoint{2.767606in}{1.037752in}}%
\pgfpathlineto{\pgfqpoint{2.775648in}{1.039823in}}%
\pgfpathlineto{\pgfqpoint{2.782087in}{1.034841in}}%
\pgfpathlineto{\pgfqpoint{2.784638in}{1.040494in}}%
\pgfpathlineto{\pgfqpoint{2.789579in}{1.040321in}}%
\pgfpathlineto{\pgfqpoint{2.793501in}{1.036335in}}%
\pgfpathlineto{\pgfqpoint{2.797982in}{1.024799in}}%
\pgfpathlineto{\pgfqpoint{2.796768in}{1.016487in}}%
\pgfpathlineto{\pgfqpoint{2.790807in}{1.014831in}}%
\pgfpathlineto{\pgfqpoint{2.786449in}{1.006708in}}%
\pgfpathlineto{\pgfqpoint{2.785854in}{1.002606in}}%
\pgfpathlineto{\pgfqpoint{2.779981in}{0.995686in}}%
\pgfpathlineto{\pgfqpoint{2.772764in}{0.995498in}}%
\pgfpathlineto{\pgfqpoint{2.764411in}{0.997827in}}%
\pgfpathlineto{\pgfqpoint{2.762533in}{0.994980in}}%
\pgfpathlineto{\pgfqpoint{2.754523in}{0.995316in}}%
\pgfpathlineto{\pgfqpoint{2.755737in}{0.991248in}}%
\pgfpathclose%
\pgfusepath{fill}%
\end{pgfscope}%
\begin{pgfscope}%
\pgfpathrectangle{\pgfqpoint{0.100000in}{0.100000in}}{\pgfqpoint{3.007045in}{1.925000in}}%
\pgfusepath{clip}%
\pgfsetbuttcap%
\pgfsetmiterjoin%
\definecolor{currentfill}{rgb}{0.726028,0.837386,0.919600}%
\pgfsetfillcolor{currentfill}%
\pgfsetlinewidth{0.000000pt}%
\definecolor{currentstroke}{rgb}{0.000000,0.000000,0.000000}%
\pgfsetstrokecolor{currentstroke}%
\pgfsetstrokeopacity{0.000000}%
\pgfsetdash{}{0pt}%
\pgfpathmoveto{\pgfqpoint{2.787444in}{1.058403in}}%
\pgfpathlineto{\pgfqpoint{2.798434in}{1.043326in}}%
\pgfpathlineto{\pgfqpoint{2.793014in}{1.044103in}}%
\pgfpathlineto{\pgfqpoint{2.786756in}{1.058061in}}%
\pgfpathlineto{\pgfqpoint{2.787444in}{1.058403in}}%
\pgfpathclose%
\pgfusepath{fill}%
\end{pgfscope}%
\begin{pgfscope}%
\pgfpathrectangle{\pgfqpoint{0.100000in}{0.100000in}}{\pgfqpoint{3.007045in}{1.925000in}}%
\pgfusepath{clip}%
\pgfsetbuttcap%
\pgfsetmiterjoin%
\definecolor{currentfill}{rgb}{0.579608,0.770196,0.873725}%
\pgfsetfillcolor{currentfill}%
\pgfsetlinewidth{0.000000pt}%
\definecolor{currentstroke}{rgb}{0.000000,0.000000,0.000000}%
\pgfsetstrokecolor{currentstroke}%
\pgfsetstrokeopacity{0.000000}%
\pgfsetdash{}{0pt}%
\pgfpathmoveto{\pgfqpoint{1.619645in}{0.516644in}}%
\pgfpathlineto{\pgfqpoint{1.621659in}{0.515303in}}%
\pgfpathlineto{\pgfqpoint{1.638148in}{0.525132in}}%
\pgfpathlineto{\pgfqpoint{1.648009in}{0.514439in}}%
\pgfpathlineto{\pgfqpoint{1.644405in}{0.511356in}}%
\pgfpathlineto{\pgfqpoint{1.641992in}{0.498317in}}%
\pgfpathlineto{\pgfqpoint{1.619767in}{0.480550in}}%
\pgfpathlineto{\pgfqpoint{1.613261in}{0.488228in}}%
\pgfpathlineto{\pgfqpoint{1.605940in}{0.497218in}}%
\pgfpathlineto{\pgfqpoint{1.619645in}{0.516644in}}%
\pgfpathclose%
\pgfusepath{fill}%
\end{pgfscope}%
\begin{pgfscope}%
\pgfpathrectangle{\pgfqpoint{0.100000in}{0.100000in}}{\pgfqpoint{3.007045in}{1.925000in}}%
\pgfusepath{clip}%
\pgfsetbuttcap%
\pgfsetmiterjoin%
\definecolor{currentfill}{rgb}{0.134487,0.448212,0.712418}%
\pgfsetfillcolor{currentfill}%
\pgfsetlinewidth{0.000000pt}%
\definecolor{currentstroke}{rgb}{0.000000,0.000000,0.000000}%
\pgfsetstrokecolor{currentstroke}%
\pgfsetstrokeopacity{0.000000}%
\pgfsetdash{}{0pt}%
\pgfpathmoveto{\pgfqpoint{1.590530in}{1.433489in}}%
\pgfpathlineto{\pgfqpoint{1.647409in}{1.431503in}}%
\pgfpathlineto{\pgfqpoint{1.657257in}{1.431226in}}%
\pgfpathlineto{\pgfqpoint{1.656789in}{1.409421in}}%
\pgfpathlineto{\pgfqpoint{1.624011in}{1.410336in}}%
\pgfpathlineto{\pgfqpoint{1.622340in}{1.404124in}}%
\pgfpathlineto{\pgfqpoint{1.622963in}{1.397529in}}%
\pgfpathlineto{\pgfqpoint{1.621472in}{1.391045in}}%
\pgfpathlineto{\pgfqpoint{1.618866in}{1.388633in}}%
\pgfpathlineto{\pgfqpoint{1.612271in}{1.391873in}}%
\pgfpathlineto{\pgfqpoint{1.604815in}{1.396442in}}%
\pgfpathlineto{\pgfqpoint{1.603438in}{1.399865in}}%
\pgfpathlineto{\pgfqpoint{1.602167in}{1.402735in}}%
\pgfpathlineto{\pgfqpoint{1.594388in}{1.405377in}}%
\pgfpathlineto{\pgfqpoint{1.590580in}{1.410117in}}%
\pgfpathlineto{\pgfqpoint{1.586747in}{1.410853in}}%
\pgfpathlineto{\pgfqpoint{1.584652in}{1.417090in}}%
\pgfpathlineto{\pgfqpoint{1.576855in}{1.422263in}}%
\pgfpathlineto{\pgfqpoint{1.572614in}{1.429955in}}%
\pgfpathlineto{\pgfqpoint{1.568472in}{1.430306in}}%
\pgfpathlineto{\pgfqpoint{1.566436in}{1.434578in}}%
\pgfpathlineto{\pgfqpoint{1.590530in}{1.433489in}}%
\pgfpathclose%
\pgfusepath{fill}%
\end{pgfscope}%
\begin{pgfscope}%
\pgfpathrectangle{\pgfqpoint{0.100000in}{0.100000in}}{\pgfqpoint{3.007045in}{1.925000in}}%
\pgfusepath{clip}%
\pgfsetbuttcap%
\pgfsetmiterjoin%
\definecolor{currentfill}{rgb}{0.301069,0.601399,0.792957}%
\pgfsetfillcolor{currentfill}%
\pgfsetlinewidth{0.000000pt}%
\definecolor{currentstroke}{rgb}{0.000000,0.000000,0.000000}%
\pgfsetstrokecolor{currentstroke}%
\pgfsetstrokeopacity{0.000000}%
\pgfsetdash{}{0pt}%
\pgfpathmoveto{\pgfqpoint{2.256834in}{1.132262in}}%
\pgfpathlineto{\pgfqpoint{2.245396in}{1.130991in}}%
\pgfpathlineto{\pgfqpoint{2.240875in}{1.128586in}}%
\pgfpathlineto{\pgfqpoint{2.239682in}{1.138991in}}%
\pgfpathlineto{\pgfqpoint{2.235786in}{1.140573in}}%
\pgfpathlineto{\pgfqpoint{2.219718in}{1.138972in}}%
\pgfpathlineto{\pgfqpoint{2.218022in}{1.153845in}}%
\pgfpathlineto{\pgfqpoint{2.215737in}{1.157449in}}%
\pgfpathlineto{\pgfqpoint{2.227838in}{1.158721in}}%
\pgfpathlineto{\pgfqpoint{2.225671in}{1.177961in}}%
\pgfpathlineto{\pgfqpoint{2.245949in}{1.180376in}}%
\pgfpathlineto{\pgfqpoint{2.247318in}{1.166002in}}%
\pgfpathlineto{\pgfqpoint{2.253429in}{1.166775in}}%
\pgfpathlineto{\pgfqpoint{2.259410in}{1.171566in}}%
\pgfpathlineto{\pgfqpoint{2.261118in}{1.153025in}}%
\pgfpathlineto{\pgfqpoint{2.253275in}{1.146329in}}%
\pgfpathlineto{\pgfqpoint{2.249568in}{1.145924in}}%
\pgfpathlineto{\pgfqpoint{2.250151in}{1.140216in}}%
\pgfpathlineto{\pgfqpoint{2.256292in}{1.137004in}}%
\pgfpathlineto{\pgfqpoint{2.256834in}{1.132262in}}%
\pgfpathclose%
\pgfusepath{fill}%
\end{pgfscope}%
\begin{pgfscope}%
\pgfpathrectangle{\pgfqpoint{0.100000in}{0.100000in}}{\pgfqpoint{3.007045in}{1.925000in}}%
\pgfusepath{clip}%
\pgfsetbuttcap%
\pgfsetmiterjoin%
\definecolor{currentfill}{rgb}{0.326290,0.618624,0.802799}%
\pgfsetfillcolor{currentfill}%
\pgfsetlinewidth{0.000000pt}%
\definecolor{currentstroke}{rgb}{0.000000,0.000000,0.000000}%
\pgfsetstrokecolor{currentstroke}%
\pgfsetstrokeopacity{0.000000}%
\pgfsetdash{}{0pt}%
\pgfpathmoveto{\pgfqpoint{1.719649in}{1.001488in}}%
\pgfpathlineto{\pgfqpoint{1.659290in}{1.002553in}}%
\pgfpathlineto{\pgfqpoint{1.659741in}{1.033948in}}%
\pgfpathlineto{\pgfqpoint{1.692552in}{1.033341in}}%
\pgfpathlineto{\pgfqpoint{1.693377in}{1.073520in}}%
\pgfpathlineto{\pgfqpoint{1.701914in}{1.073348in}}%
\pgfpathlineto{\pgfqpoint{1.702042in}{1.079070in}}%
\pgfpathlineto{\pgfqpoint{1.722658in}{1.078703in}}%
\pgfpathlineto{\pgfqpoint{1.722556in}{1.070090in}}%
\pgfpathlineto{\pgfqpoint{1.722263in}{1.049928in}}%
\pgfpathlineto{\pgfqpoint{1.721571in}{1.001483in}}%
\pgfpathlineto{\pgfqpoint{1.719649in}{1.001488in}}%
\pgfpathclose%
\pgfusepath{fill}%
\end{pgfscope}%
\begin{pgfscope}%
\pgfpathrectangle{\pgfqpoint{0.100000in}{0.100000in}}{\pgfqpoint{3.007045in}{1.925000in}}%
\pgfusepath{clip}%
\pgfsetbuttcap%
\pgfsetmiterjoin%
\definecolor{currentfill}{rgb}{0.671895,0.814379,0.900654}%
\pgfsetfillcolor{currentfill}%
\pgfsetlinewidth{0.000000pt}%
\definecolor{currentstroke}{rgb}{0.000000,0.000000,0.000000}%
\pgfsetstrokecolor{currentstroke}%
\pgfsetstrokeopacity{0.000000}%
\pgfsetdash{}{0pt}%
\pgfpathmoveto{\pgfqpoint{2.406434in}{1.016470in}}%
\pgfpathlineto{\pgfqpoint{2.372965in}{1.012422in}}%
\pgfpathlineto{\pgfqpoint{2.373385in}{1.015073in}}%
\pgfpathlineto{\pgfqpoint{2.380873in}{1.017870in}}%
\pgfpathlineto{\pgfqpoint{2.397238in}{1.025171in}}%
\pgfpathlineto{\pgfqpoint{2.399471in}{1.032629in}}%
\pgfpathlineto{\pgfqpoint{2.409387in}{1.036508in}}%
\pgfpathlineto{\pgfqpoint{2.409125in}{1.041720in}}%
\pgfpathlineto{\pgfqpoint{2.415182in}{1.047525in}}%
\pgfpathlineto{\pgfqpoint{2.415558in}{1.052709in}}%
\pgfpathlineto{\pgfqpoint{2.423037in}{1.058784in}}%
\pgfpathlineto{\pgfqpoint{2.435154in}{1.067105in}}%
\pgfpathlineto{\pgfqpoint{2.440019in}{1.062167in}}%
\pgfpathlineto{\pgfqpoint{2.446669in}{1.052394in}}%
\pgfpathlineto{\pgfqpoint{2.437440in}{1.045801in}}%
\pgfpathlineto{\pgfqpoint{2.439575in}{1.041742in}}%
\pgfpathlineto{\pgfqpoint{2.434200in}{1.039055in}}%
\pgfpathlineto{\pgfqpoint{2.422280in}{1.037551in}}%
\pgfpathlineto{\pgfqpoint{2.416735in}{1.033753in}}%
\pgfpathlineto{\pgfqpoint{2.413702in}{1.026726in}}%
\pgfpathlineto{\pgfqpoint{2.407690in}{1.022099in}}%
\pgfpathlineto{\pgfqpoint{2.406434in}{1.016470in}}%
\pgfpathclose%
\pgfusepath{fill}%
\end{pgfscope}%
\begin{pgfscope}%
\pgfpathrectangle{\pgfqpoint{0.100000in}{0.100000in}}{\pgfqpoint{3.007045in}{1.925000in}}%
\pgfusepath{clip}%
\pgfsetbuttcap%
\pgfsetmiterjoin%
\definecolor{currentfill}{rgb}{0.108651,0.416563,0.689043}%
\pgfsetfillcolor{currentfill}%
\pgfsetlinewidth{0.000000pt}%
\definecolor{currentstroke}{rgb}{0.000000,0.000000,0.000000}%
\pgfsetstrokecolor{currentstroke}%
\pgfsetstrokeopacity{0.000000}%
\pgfsetdash{}{0pt}%
\pgfpathmoveto{\pgfqpoint{1.419469in}{0.951860in}}%
\pgfpathlineto{\pgfqpoint{1.417413in}{0.923222in}}%
\pgfpathlineto{\pgfqpoint{1.388570in}{0.925898in}}%
\pgfpathlineto{\pgfqpoint{1.341578in}{0.929505in}}%
\pgfpathlineto{\pgfqpoint{1.342266in}{0.937180in}}%
\pgfpathlineto{\pgfqpoint{1.346674in}{0.987189in}}%
\pgfpathlineto{\pgfqpoint{1.348759in}{0.987002in}}%
\pgfpathlineto{\pgfqpoint{1.393120in}{0.983287in}}%
\pgfpathlineto{\pgfqpoint{1.390781in}{0.954030in}}%
\pgfpathlineto{\pgfqpoint{1.419469in}{0.951860in}}%
\pgfpathclose%
\pgfusepath{fill}%
\end{pgfscope}%
\begin{pgfscope}%
\pgfpathrectangle{\pgfqpoint{0.100000in}{0.100000in}}{\pgfqpoint{3.007045in}{1.925000in}}%
\pgfusepath{clip}%
\pgfsetbuttcap%
\pgfsetmiterjoin%
\definecolor{currentfill}{rgb}{0.256286,0.570012,0.775163}%
\pgfsetfillcolor{currentfill}%
\pgfsetlinewidth{0.000000pt}%
\definecolor{currentstroke}{rgb}{0.000000,0.000000,0.000000}%
\pgfsetstrokecolor{currentstroke}%
\pgfsetstrokeopacity{0.000000}%
\pgfsetdash{}{0pt}%
\pgfpathmoveto{\pgfqpoint{1.605884in}{0.418493in}}%
\pgfpathlineto{\pgfqpoint{1.600656in}{0.414463in}}%
\pgfpathlineto{\pgfqpoint{1.601772in}{0.410582in}}%
\pgfpathlineto{\pgfqpoint{1.547367in}{0.412624in}}%
\pgfpathlineto{\pgfqpoint{1.549167in}{0.452106in}}%
\pgfpathlineto{\pgfqpoint{1.550168in}{0.481716in}}%
\pgfpathlineto{\pgfqpoint{1.514641in}{0.483384in}}%
\pgfpathlineto{\pgfqpoint{1.516499in}{0.519146in}}%
\pgfpathlineto{\pgfqpoint{1.505414in}{0.519674in}}%
\pgfpathlineto{\pgfqpoint{1.506378in}{0.538351in}}%
\pgfpathlineto{\pgfqpoint{1.530916in}{0.536323in}}%
\pgfpathlineto{\pgfqpoint{1.545579in}{0.528102in}}%
\pgfpathlineto{\pgfqpoint{1.546459in}{0.551872in}}%
\pgfpathlineto{\pgfqpoint{1.565594in}{0.551126in}}%
\pgfpathlineto{\pgfqpoint{1.575104in}{0.537324in}}%
\pgfpathlineto{\pgfqpoint{1.582069in}{0.543767in}}%
\pgfpathlineto{\pgfqpoint{1.597051in}{0.530605in}}%
\pgfpathlineto{\pgfqpoint{1.598646in}{0.524138in}}%
\pgfpathlineto{\pgfqpoint{1.606065in}{0.530962in}}%
\pgfpathlineto{\pgfqpoint{1.615672in}{0.518019in}}%
\pgfpathlineto{\pgfqpoint{1.619645in}{0.516644in}}%
\pgfpathlineto{\pgfqpoint{1.605940in}{0.497218in}}%
\pgfpathlineto{\pgfqpoint{1.613261in}{0.488228in}}%
\pgfpathlineto{\pgfqpoint{1.585508in}{0.466431in}}%
\pgfpathlineto{\pgfqpoint{1.595954in}{0.453216in}}%
\pgfpathlineto{\pgfqpoint{1.590934in}{0.451560in}}%
\pgfpathlineto{\pgfqpoint{1.606432in}{0.418922in}}%
\pgfpathlineto{\pgfqpoint{1.605884in}{0.418493in}}%
\pgfpathclose%
\pgfusepath{fill}%
\end{pgfscope}%
\begin{pgfscope}%
\pgfpathrectangle{\pgfqpoint{0.100000in}{0.100000in}}{\pgfqpoint{3.007045in}{1.925000in}}%
\pgfusepath{clip}%
\pgfsetbuttcap%
\pgfsetmiterjoin%
\definecolor{currentfill}{rgb}{0.248166,0.561892,0.770980}%
\pgfsetfillcolor{currentfill}%
\pgfsetlinewidth{0.000000pt}%
\definecolor{currentstroke}{rgb}{0.000000,0.000000,0.000000}%
\pgfsetstrokecolor{currentstroke}%
\pgfsetstrokeopacity{0.000000}%
\pgfsetdash{}{0pt}%
\pgfpathmoveto{\pgfqpoint{1.542619in}{0.830859in}}%
\pgfpathlineto{\pgfqpoint{1.533732in}{0.839225in}}%
\pgfpathlineto{\pgfqpoint{1.531947in}{0.833744in}}%
\pgfpathlineto{\pgfqpoint{1.526280in}{0.836772in}}%
\pgfpathlineto{\pgfqpoint{1.521539in}{0.834387in}}%
\pgfpathlineto{\pgfqpoint{1.516379in}{0.834892in}}%
\pgfpathlineto{\pgfqpoint{1.504435in}{0.848612in}}%
\pgfpathlineto{\pgfqpoint{1.500565in}{0.847796in}}%
\pgfpathlineto{\pgfqpoint{1.501084in}{0.860066in}}%
\pgfpathlineto{\pgfqpoint{1.502115in}{0.878788in}}%
\pgfpathlineto{\pgfqpoint{1.513852in}{0.878138in}}%
\pgfpathlineto{\pgfqpoint{1.514160in}{0.883851in}}%
\pgfpathlineto{\pgfqpoint{1.569929in}{0.881210in}}%
\pgfpathlineto{\pgfqpoint{1.576781in}{0.879670in}}%
\pgfpathlineto{\pgfqpoint{1.576210in}{0.863673in}}%
\pgfpathlineto{\pgfqpoint{1.565031in}{0.864162in}}%
\pgfpathlineto{\pgfqpoint{1.564310in}{0.846924in}}%
\pgfpathlineto{\pgfqpoint{1.545976in}{0.847321in}}%
\pgfpathlineto{\pgfqpoint{1.542518in}{0.841248in}}%
\pgfpathlineto{\pgfqpoint{1.542619in}{0.830859in}}%
\pgfpathclose%
\pgfusepath{fill}%
\end{pgfscope}%
\begin{pgfscope}%
\pgfpathrectangle{\pgfqpoint{0.100000in}{0.100000in}}{\pgfqpoint{3.007045in}{1.925000in}}%
\pgfusepath{clip}%
\pgfsetbuttcap%
\pgfsetmiterjoin%
\definecolor{currentfill}{rgb}{0.376732,0.653072,0.822484}%
\pgfsetfillcolor{currentfill}%
\pgfsetlinewidth{0.000000pt}%
\definecolor{currentstroke}{rgb}{0.000000,0.000000,0.000000}%
\pgfsetstrokecolor{currentstroke}%
\pgfsetstrokeopacity{0.000000}%
\pgfsetdash{}{0pt}%
\pgfpathmoveto{\pgfqpoint{1.702839in}{1.116383in}}%
\pgfpathlineto{\pgfqpoint{1.702564in}{1.102052in}}%
\pgfpathlineto{\pgfqpoint{1.678559in}{1.102524in}}%
\pgfpathlineto{\pgfqpoint{1.672876in}{1.102642in}}%
\pgfpathlineto{\pgfqpoint{1.672984in}{1.108364in}}%
\pgfpathlineto{\pgfqpoint{1.650255in}{1.108902in}}%
\pgfpathlineto{\pgfqpoint{1.651571in}{1.154782in}}%
\pgfpathlineto{\pgfqpoint{1.652076in}{1.171950in}}%
\pgfpathlineto{\pgfqpoint{1.709420in}{1.170767in}}%
\pgfpathlineto{\pgfqpoint{1.719727in}{1.170628in}}%
\pgfpathlineto{\pgfqpoint{1.719484in}{1.147603in}}%
\pgfpathlineto{\pgfqpoint{1.712256in}{1.145853in}}%
\pgfpathlineto{\pgfqpoint{1.708830in}{1.147175in}}%
\pgfpathlineto{\pgfqpoint{1.701412in}{1.144959in}}%
\pgfpathlineto{\pgfqpoint{1.701186in}{1.136436in}}%
\pgfpathlineto{\pgfqpoint{1.695539in}{1.136550in}}%
\pgfpathlineto{\pgfqpoint{1.695279in}{1.122232in}}%
\pgfpathlineto{\pgfqpoint{1.700977in}{1.122113in}}%
\pgfpathlineto{\pgfqpoint{1.702839in}{1.116383in}}%
\pgfpathclose%
\pgfusepath{fill}%
\end{pgfscope}%
\begin{pgfscope}%
\pgfpathrectangle{\pgfqpoint{0.100000in}{0.100000in}}{\pgfqpoint{3.007045in}{1.925000in}}%
\pgfusepath{clip}%
\pgfsetbuttcap%
\pgfsetmiterjoin%
\definecolor{currentfill}{rgb}{0.460392,0.704744,0.848012}%
\pgfsetfillcolor{currentfill}%
\pgfsetlinewidth{0.000000pt}%
\definecolor{currentstroke}{rgb}{0.000000,0.000000,0.000000}%
\pgfsetstrokecolor{currentstroke}%
\pgfsetstrokeopacity{0.000000}%
\pgfsetdash{}{0pt}%
\pgfpathmoveto{\pgfqpoint{2.505332in}{0.935315in}}%
\pgfpathlineto{\pgfqpoint{2.472948in}{0.932025in}}%
\pgfpathlineto{\pgfqpoint{2.472311in}{0.936614in}}%
\pgfpathlineto{\pgfqpoint{2.460723in}{0.945065in}}%
\pgfpathlineto{\pgfqpoint{2.455447in}{0.943242in}}%
\pgfpathlineto{\pgfqpoint{2.453980in}{0.946316in}}%
\pgfpathlineto{\pgfqpoint{2.455786in}{0.951722in}}%
\pgfpathlineto{\pgfqpoint{2.459120in}{0.952751in}}%
\pgfpathlineto{\pgfqpoint{2.469542in}{0.954025in}}%
\pgfpathlineto{\pgfqpoint{2.476332in}{0.956183in}}%
\pgfpathlineto{\pgfqpoint{2.480332in}{0.961447in}}%
\pgfpathlineto{\pgfqpoint{2.486422in}{0.958962in}}%
\pgfpathlineto{\pgfqpoint{2.492448in}{0.960338in}}%
\pgfpathlineto{\pgfqpoint{2.523136in}{0.963616in}}%
\pgfpathlineto{\pgfqpoint{2.523665in}{0.961730in}}%
\pgfpathlineto{\pgfqpoint{2.524892in}{0.952370in}}%
\pgfpathlineto{\pgfqpoint{2.523035in}{0.944991in}}%
\pgfpathlineto{\pgfqpoint{2.523780in}{0.937038in}}%
\pgfpathlineto{\pgfqpoint{2.505332in}{0.935315in}}%
\pgfpathclose%
\pgfusepath{fill}%
\end{pgfscope}%
\begin{pgfscope}%
\pgfpathrectangle{\pgfqpoint{0.100000in}{0.100000in}}{\pgfqpoint{3.007045in}{1.925000in}}%
\pgfusepath{clip}%
\pgfsetbuttcap%
\pgfsetmiterjoin%
\definecolor{currentfill}{rgb}{0.479216,0.715079,0.852072}%
\pgfsetfillcolor{currentfill}%
\pgfsetlinewidth{0.000000pt}%
\definecolor{currentstroke}{rgb}{0.000000,0.000000,0.000000}%
\pgfsetstrokecolor{currentstroke}%
\pgfsetstrokeopacity{0.000000}%
\pgfsetdash{}{0pt}%
\pgfpathmoveto{\pgfqpoint{2.452246in}{0.846939in}}%
\pgfpathlineto{\pgfqpoint{2.447000in}{0.848563in}}%
\pgfpathlineto{\pgfqpoint{2.440442in}{0.846333in}}%
\pgfpathlineto{\pgfqpoint{2.431219in}{0.848011in}}%
\pgfpathlineto{\pgfqpoint{2.421120in}{0.857223in}}%
\pgfpathlineto{\pgfqpoint{2.422010in}{0.860947in}}%
\pgfpathlineto{\pgfqpoint{2.409663in}{0.861811in}}%
\pgfpathlineto{\pgfqpoint{2.406992in}{0.858796in}}%
\pgfpathlineto{\pgfqpoint{2.404051in}{0.865135in}}%
\pgfpathlineto{\pgfqpoint{2.403763in}{0.872065in}}%
\pgfpathlineto{\pgfqpoint{2.399169in}{0.874931in}}%
\pgfpathlineto{\pgfqpoint{2.402644in}{0.888304in}}%
\pgfpathlineto{\pgfqpoint{2.403995in}{0.889257in}}%
\pgfpathlineto{\pgfqpoint{2.410336in}{0.885056in}}%
\pgfpathlineto{\pgfqpoint{2.413604in}{0.885230in}}%
\pgfpathlineto{\pgfqpoint{2.420077in}{0.879039in}}%
\pgfpathlineto{\pgfqpoint{2.423854in}{0.877561in}}%
\pgfpathlineto{\pgfqpoint{2.429221in}{0.879173in}}%
\pgfpathlineto{\pgfqpoint{2.439934in}{0.864954in}}%
\pgfpathlineto{\pgfqpoint{2.442202in}{0.858557in}}%
\pgfpathlineto{\pgfqpoint{2.449961in}{0.851491in}}%
\pgfpathlineto{\pgfqpoint{2.452246in}{0.846939in}}%
\pgfpathclose%
\pgfusepath{fill}%
\end{pgfscope}%
\begin{pgfscope}%
\pgfpathrectangle{\pgfqpoint{0.100000in}{0.100000in}}{\pgfqpoint{3.007045in}{1.925000in}}%
\pgfusepath{clip}%
\pgfsetbuttcap%
\pgfsetmiterjoin%
\definecolor{currentfill}{rgb}{0.429020,0.687520,0.841246}%
\pgfsetfillcolor{currentfill}%
\pgfsetlinewidth{0.000000pt}%
\definecolor{currentstroke}{rgb}{0.000000,0.000000,0.000000}%
\pgfsetstrokecolor{currentstroke}%
\pgfsetstrokeopacity{0.000000}%
\pgfsetdash{}{0pt}%
\pgfpathmoveto{\pgfqpoint{0.816917in}{1.293210in}}%
\pgfpathlineto{\pgfqpoint{0.811909in}{1.269809in}}%
\pgfpathlineto{\pgfqpoint{0.799942in}{1.214078in}}%
\pgfpathlineto{\pgfqpoint{0.752166in}{1.224502in}}%
\pgfpathlineto{\pgfqpoint{0.714021in}{1.265975in}}%
\pgfpathlineto{\pgfqpoint{0.679545in}{1.274162in}}%
\pgfpathlineto{\pgfqpoint{0.643335in}{1.283223in}}%
\pgfpathlineto{\pgfqpoint{0.620230in}{1.284320in}}%
\pgfpathlineto{\pgfqpoint{0.621234in}{1.289668in}}%
\pgfpathlineto{\pgfqpoint{0.625302in}{1.292738in}}%
\pgfpathlineto{\pgfqpoint{0.625081in}{1.300716in}}%
\pgfpathlineto{\pgfqpoint{0.632081in}{1.306373in}}%
\pgfpathlineto{\pgfqpoint{0.641368in}{1.309018in}}%
\pgfpathlineto{\pgfqpoint{0.644313in}{1.314171in}}%
\pgfpathlineto{\pgfqpoint{0.644679in}{1.322619in}}%
\pgfpathlineto{\pgfqpoint{0.648824in}{1.333373in}}%
\pgfpathlineto{\pgfqpoint{0.646486in}{1.339327in}}%
\pgfpathlineto{\pgfqpoint{0.671162in}{1.376694in}}%
\pgfpathlineto{\pgfqpoint{0.682240in}{1.373981in}}%
\pgfpathlineto{\pgfqpoint{0.687900in}{1.396831in}}%
\pgfpathlineto{\pgfqpoint{0.703710in}{1.460774in}}%
\pgfpathlineto{\pgfqpoint{0.734607in}{1.453046in}}%
\pgfpathlineto{\pgfqpoint{0.798116in}{1.438303in}}%
\pgfpathlineto{\pgfqpoint{0.808116in}{1.436493in}}%
\pgfpathlineto{\pgfqpoint{0.834301in}{1.430229in}}%
\pgfpathlineto{\pgfqpoint{0.845804in}{1.427734in}}%
\pgfpathlineto{\pgfqpoint{0.832155in}{1.363654in}}%
\pgfpathlineto{\pgfqpoint{0.816917in}{1.293210in}}%
\pgfpathclose%
\pgfusepath{fill}%
\end{pgfscope}%
\begin{pgfscope}%
\pgfpathrectangle{\pgfqpoint{0.100000in}{0.100000in}}{\pgfqpoint{3.007045in}{1.925000in}}%
\pgfusepath{clip}%
\pgfsetbuttcap%
\pgfsetmiterjoin%
\definecolor{currentfill}{rgb}{0.031373,0.285675,0.564291}%
\pgfsetfillcolor{currentfill}%
\pgfsetlinewidth{0.000000pt}%
\definecolor{currentstroke}{rgb}{0.000000,0.000000,0.000000}%
\pgfsetstrokecolor{currentstroke}%
\pgfsetstrokeopacity{0.000000}%
\pgfsetdash{}{0pt}%
\pgfpathmoveto{\pgfqpoint{1.566436in}{1.434578in}}%
\pgfpathlineto{\pgfqpoint{1.568472in}{1.430306in}}%
\pgfpathlineto{\pgfqpoint{1.572614in}{1.429955in}}%
\pgfpathlineto{\pgfqpoint{1.576855in}{1.422263in}}%
\pgfpathlineto{\pgfqpoint{1.584652in}{1.417090in}}%
\pgfpathlineto{\pgfqpoint{1.586747in}{1.410853in}}%
\pgfpathlineto{\pgfqpoint{1.590580in}{1.410117in}}%
\pgfpathlineto{\pgfqpoint{1.594388in}{1.405377in}}%
\pgfpathlineto{\pgfqpoint{1.602167in}{1.402735in}}%
\pgfpathlineto{\pgfqpoint{1.603438in}{1.399865in}}%
\pgfpathlineto{\pgfqpoint{1.566978in}{1.401393in}}%
\pgfpathlineto{\pgfqpoint{1.521452in}{1.403789in}}%
\pgfpathlineto{\pgfqpoint{1.521204in}{1.412401in}}%
\pgfpathlineto{\pgfqpoint{1.522496in}{1.435436in}}%
\pgfpathlineto{\pgfqpoint{1.522608in}{1.451050in}}%
\pgfpathlineto{\pgfqpoint{1.531144in}{1.449134in}}%
\pgfpathlineto{\pgfqpoint{1.541165in}{1.448356in}}%
\pgfpathlineto{\pgfqpoint{1.551671in}{1.451429in}}%
\pgfpathlineto{\pgfqpoint{1.550851in}{1.435352in}}%
\pgfpathlineto{\pgfqpoint{1.566436in}{1.434578in}}%
\pgfpathclose%
\pgfusepath{fill}%
\end{pgfscope}%
\begin{pgfscope}%
\pgfpathrectangle{\pgfqpoint{0.100000in}{0.100000in}}{\pgfqpoint{3.007045in}{1.925000in}}%
\pgfusepath{clip}%
\pgfsetbuttcap%
\pgfsetmiterjoin%
\definecolor{currentfill}{rgb}{0.447843,0.697855,0.845306}%
\pgfsetfillcolor{currentfill}%
\pgfsetlinewidth{0.000000pt}%
\definecolor{currentstroke}{rgb}{0.000000,0.000000,0.000000}%
\pgfsetstrokecolor{currentstroke}%
\pgfsetstrokeopacity{0.000000}%
\pgfsetdash{}{0pt}%
\pgfpathmoveto{\pgfqpoint{2.359552in}{0.655091in}}%
\pgfpathlineto{\pgfqpoint{2.343997in}{0.653135in}}%
\pgfpathlineto{\pgfqpoint{2.345074in}{0.640987in}}%
\pgfpathlineto{\pgfqpoint{2.339044in}{0.640515in}}%
\pgfpathlineto{\pgfqpoint{2.334338in}{0.646016in}}%
\pgfpathlineto{\pgfqpoint{2.332928in}{0.653635in}}%
\pgfpathlineto{\pgfqpoint{2.334926in}{0.671284in}}%
\pgfpathlineto{\pgfqpoint{2.333467in}{0.676210in}}%
\pgfpathlineto{\pgfqpoint{2.329091in}{0.680831in}}%
\pgfpathlineto{\pgfqpoint{2.312691in}{0.679850in}}%
\pgfpathlineto{\pgfqpoint{2.313278in}{0.674089in}}%
\pgfpathlineto{\pgfqpoint{2.294624in}{0.672055in}}%
\pgfpathlineto{\pgfqpoint{2.298312in}{0.683587in}}%
\pgfpathlineto{\pgfqpoint{2.297933in}{0.689896in}}%
\pgfpathlineto{\pgfqpoint{2.301073in}{0.698033in}}%
\pgfpathlineto{\pgfqpoint{2.309841in}{0.700126in}}%
\pgfpathlineto{\pgfqpoint{2.309942in}{0.708924in}}%
\pgfpathlineto{\pgfqpoint{2.318447in}{0.709939in}}%
\pgfpathlineto{\pgfqpoint{2.323124in}{0.704669in}}%
\pgfpathlineto{\pgfqpoint{2.330327in}{0.705545in}}%
\pgfpathlineto{\pgfqpoint{2.331743in}{0.699866in}}%
\pgfpathlineto{\pgfqpoint{2.339099in}{0.697275in}}%
\pgfpathlineto{\pgfqpoint{2.356903in}{0.698990in}}%
\pgfpathlineto{\pgfqpoint{2.357840in}{0.689197in}}%
\pgfpathlineto{\pgfqpoint{2.362322in}{0.683671in}}%
\pgfpathlineto{\pgfqpoint{2.367559in}{0.680287in}}%
\pgfpathlineto{\pgfqpoint{2.367869in}{0.675418in}}%
\pgfpathlineto{\pgfqpoint{2.370955in}{0.669671in}}%
\pgfpathlineto{\pgfqpoint{2.366771in}{0.667222in}}%
\pgfpathlineto{\pgfqpoint{2.358460in}{0.666668in}}%
\pgfpathlineto{\pgfqpoint{2.359552in}{0.655091in}}%
\pgfpathclose%
\pgfusepath{fill}%
\end{pgfscope}%
\begin{pgfscope}%
\pgfpathrectangle{\pgfqpoint{0.100000in}{0.100000in}}{\pgfqpoint{3.007045in}{1.925000in}}%
\pgfusepath{clip}%
\pgfsetbuttcap%
\pgfsetmiterjoin%
\definecolor{currentfill}{rgb}{0.681738,0.818562,0.904098}%
\pgfsetfillcolor{currentfill}%
\pgfsetlinewidth{0.000000pt}%
\definecolor{currentstroke}{rgb}{0.000000,0.000000,0.000000}%
\pgfsetstrokecolor{currentstroke}%
\pgfsetstrokeopacity{0.000000}%
\pgfsetdash{}{0pt}%
\pgfpathmoveto{\pgfqpoint{0.923792in}{1.677916in}}%
\pgfpathlineto{\pgfqpoint{0.921890in}{1.674121in}}%
\pgfpathlineto{\pgfqpoint{0.914135in}{1.673457in}}%
\pgfpathlineto{\pgfqpoint{0.908010in}{1.668930in}}%
\pgfpathlineto{\pgfqpoint{0.901414in}{1.666132in}}%
\pgfpathlineto{\pgfqpoint{0.894600in}{1.662965in}}%
\pgfpathlineto{\pgfqpoint{0.890559in}{1.659065in}}%
\pgfpathlineto{\pgfqpoint{0.879064in}{1.657147in}}%
\pgfpathlineto{\pgfqpoint{0.873352in}{1.663504in}}%
\pgfpathlineto{\pgfqpoint{0.871070in}{1.663297in}}%
\pgfpathlineto{\pgfqpoint{0.875599in}{1.669863in}}%
\pgfpathlineto{\pgfqpoint{0.874072in}{1.677313in}}%
\pgfpathlineto{\pgfqpoint{0.877655in}{1.681396in}}%
\pgfpathlineto{\pgfqpoint{0.883784in}{1.683111in}}%
\pgfpathlineto{\pgfqpoint{0.881292in}{1.694653in}}%
\pgfpathlineto{\pgfqpoint{0.885040in}{1.701682in}}%
\pgfpathlineto{\pgfqpoint{0.885687in}{1.708694in}}%
\pgfpathlineto{\pgfqpoint{0.891175in}{1.716853in}}%
\pgfpathlineto{\pgfqpoint{0.897388in}{1.731726in}}%
\pgfpathlineto{\pgfqpoint{0.895843in}{1.733369in}}%
\pgfpathlineto{\pgfqpoint{0.884196in}{1.733912in}}%
\pgfpathlineto{\pgfqpoint{0.884881in}{1.738405in}}%
\pgfpathlineto{\pgfqpoint{0.873285in}{1.749159in}}%
\pgfpathlineto{\pgfqpoint{0.873084in}{1.756371in}}%
\pgfpathlineto{\pgfqpoint{0.869191in}{1.761494in}}%
\pgfpathlineto{\pgfqpoint{0.861882in}{1.780983in}}%
\pgfpathlineto{\pgfqpoint{0.853297in}{1.785851in}}%
\pgfpathlineto{\pgfqpoint{0.853094in}{1.789354in}}%
\pgfpathlineto{\pgfqpoint{0.845810in}{1.797316in}}%
\pgfpathlineto{\pgfqpoint{0.851627in}{1.798562in}}%
\pgfpathlineto{\pgfqpoint{0.859543in}{1.795823in}}%
\pgfpathlineto{\pgfqpoint{0.866931in}{1.795554in}}%
\pgfpathlineto{\pgfqpoint{0.878913in}{1.786617in}}%
\pgfpathlineto{\pgfqpoint{0.878165in}{1.781128in}}%
\pgfpathlineto{\pgfqpoint{0.883118in}{1.777032in}}%
\pgfpathlineto{\pgfqpoint{0.890065in}{1.775890in}}%
\pgfpathlineto{\pgfqpoint{0.900754in}{1.769170in}}%
\pgfpathlineto{\pgfqpoint{0.906982in}{1.761729in}}%
\pgfpathlineto{\pgfqpoint{0.910196in}{1.762234in}}%
\pgfpathlineto{\pgfqpoint{0.921428in}{1.759848in}}%
\pgfpathlineto{\pgfqpoint{0.926970in}{1.761628in}}%
\pgfpathlineto{\pgfqpoint{0.926220in}{1.769071in}}%
\pgfpathlineto{\pgfqpoint{0.927695in}{1.775797in}}%
\pgfpathlineto{\pgfqpoint{0.925310in}{1.782011in}}%
\pgfpathlineto{\pgfqpoint{0.927953in}{1.789677in}}%
\pgfpathlineto{\pgfqpoint{0.948336in}{1.785640in}}%
\pgfpathlineto{\pgfqpoint{0.942694in}{1.758174in}}%
\pgfpathlineto{\pgfqpoint{0.949935in}{1.756722in}}%
\pgfpathlineto{\pgfqpoint{0.945308in}{1.734086in}}%
\pgfpathlineto{\pgfqpoint{0.940260in}{1.735106in}}%
\pgfpathlineto{\pgfqpoint{0.936258in}{1.730024in}}%
\pgfpathlineto{\pgfqpoint{0.927484in}{1.729880in}}%
\pgfpathlineto{\pgfqpoint{0.926725in}{1.726127in}}%
\pgfpathlineto{\pgfqpoint{0.921093in}{1.727298in}}%
\pgfpathlineto{\pgfqpoint{0.915399in}{1.718366in}}%
\pgfpathlineto{\pgfqpoint{0.914970in}{1.712349in}}%
\pgfpathlineto{\pgfqpoint{0.917545in}{1.708081in}}%
\pgfpathlineto{\pgfqpoint{0.917747in}{1.701623in}}%
\pgfpathlineto{\pgfqpoint{0.914111in}{1.694407in}}%
\pgfpathlineto{\pgfqpoint{0.913359in}{1.689243in}}%
\pgfpathlineto{\pgfqpoint{0.917372in}{1.685810in}}%
\pgfpathlineto{\pgfqpoint{0.918191in}{1.680868in}}%
\pgfpathlineto{\pgfqpoint{0.923792in}{1.677916in}}%
\pgfpathclose%
\pgfusepath{fill}%
\end{pgfscope}%
\begin{pgfscope}%
\pgfpathrectangle{\pgfqpoint{0.100000in}{0.100000in}}{\pgfqpoint{3.007045in}{1.925000in}}%
\pgfusepath{clip}%
\pgfsetbuttcap%
\pgfsetmiterjoin%
\definecolor{currentfill}{rgb}{0.541961,0.749527,0.865606}%
\pgfsetfillcolor{currentfill}%
\pgfsetlinewidth{0.000000pt}%
\definecolor{currentstroke}{rgb}{0.000000,0.000000,0.000000}%
\pgfsetstrokecolor{currentstroke}%
\pgfsetstrokeopacity{0.000000}%
\pgfsetdash{}{0pt}%
\pgfpathmoveto{\pgfqpoint{1.844441in}{1.606089in}}%
\pgfpathlineto{\pgfqpoint{1.861813in}{1.606700in}}%
\pgfpathlineto{\pgfqpoint{1.861397in}{1.624038in}}%
\pgfpathlineto{\pgfqpoint{1.896284in}{1.624795in}}%
\pgfpathlineto{\pgfqpoint{1.896759in}{1.607543in}}%
\pgfpathlineto{\pgfqpoint{1.896908in}{1.602038in}}%
\pgfpathlineto{\pgfqpoint{1.894425in}{1.598113in}}%
\pgfpathlineto{\pgfqpoint{1.881222in}{1.592179in}}%
\pgfpathlineto{\pgfqpoint{1.877917in}{1.589345in}}%
\pgfpathlineto{\pgfqpoint{1.874776in}{1.580802in}}%
\pgfpathlineto{\pgfqpoint{1.872272in}{1.578521in}}%
\pgfpathlineto{\pgfqpoint{1.841127in}{1.578186in}}%
\pgfpathlineto{\pgfqpoint{1.840662in}{1.594594in}}%
\pgfpathlineto{\pgfqpoint{1.844608in}{1.594634in}}%
\pgfpathlineto{\pgfqpoint{1.844441in}{1.606089in}}%
\pgfpathclose%
\pgfusepath{fill}%
\end{pgfscope}%
\begin{pgfscope}%
\pgfpathrectangle{\pgfqpoint{0.100000in}{0.100000in}}{\pgfqpoint{3.007045in}{1.925000in}}%
\pgfusepath{clip}%
\pgfsetbuttcap%
\pgfsetmiterjoin%
\definecolor{currentfill}{rgb}{0.301069,0.601399,0.792957}%
\pgfsetfillcolor{currentfill}%
\pgfsetlinewidth{0.000000pt}%
\definecolor{currentstroke}{rgb}{0.000000,0.000000,0.000000}%
\pgfsetstrokecolor{currentstroke}%
\pgfsetstrokeopacity{0.000000}%
\pgfsetdash{}{0pt}%
\pgfpathmoveto{\pgfqpoint{1.561807in}{1.005883in}}%
\pgfpathlineto{\pgfqpoint{1.537864in}{1.006990in}}%
\pgfpathlineto{\pgfqpoint{1.533400in}{1.007220in}}%
\pgfpathlineto{\pgfqpoint{1.534548in}{1.032337in}}%
\pgfpathlineto{\pgfqpoint{1.562899in}{1.031213in}}%
\pgfpathlineto{\pgfqpoint{1.561807in}{1.005883in}}%
\pgfpathclose%
\pgfusepath{fill}%
\end{pgfscope}%
\begin{pgfscope}%
\pgfpathrectangle{\pgfqpoint{0.100000in}{0.100000in}}{\pgfqpoint{3.007045in}{1.925000in}}%
\pgfusepath{clip}%
\pgfsetbuttcap%
\pgfsetmiterjoin%
\definecolor{currentfill}{rgb}{0.371688,0.649627,0.820515}%
\pgfsetfillcolor{currentfill}%
\pgfsetlinewidth{0.000000pt}%
\definecolor{currentstroke}{rgb}{0.000000,0.000000,0.000000}%
\pgfsetstrokecolor{currentstroke}%
\pgfsetstrokeopacity{0.000000}%
\pgfsetdash{}{0pt}%
\pgfpathmoveto{\pgfqpoint{2.174509in}{1.028592in}}%
\pgfpathlineto{\pgfqpoint{2.164956in}{1.024377in}}%
\pgfpathlineto{\pgfqpoint{2.156379in}{1.026602in}}%
\pgfpathlineto{\pgfqpoint{2.155476in}{1.018069in}}%
\pgfpathlineto{\pgfqpoint{2.152882in}{1.014762in}}%
\pgfpathlineto{\pgfqpoint{2.146362in}{1.013286in}}%
\pgfpathlineto{\pgfqpoint{2.132817in}{1.006098in}}%
\pgfpathlineto{\pgfqpoint{2.127814in}{1.020206in}}%
\pgfpathlineto{\pgfqpoint{2.129826in}{1.024933in}}%
\pgfpathlineto{\pgfqpoint{2.127608in}{1.032537in}}%
\pgfpathlineto{\pgfqpoint{2.119654in}{1.040653in}}%
\pgfpathlineto{\pgfqpoint{2.123429in}{1.044085in}}%
\pgfpathlineto{\pgfqpoint{2.133456in}{1.046059in}}%
\pgfpathlineto{\pgfqpoint{2.128798in}{1.058421in}}%
\pgfpathlineto{\pgfqpoint{2.134908in}{1.067855in}}%
\pgfpathlineto{\pgfqpoint{2.141157in}{1.068895in}}%
\pgfpathlineto{\pgfqpoint{2.139094in}{1.074231in}}%
\pgfpathlineto{\pgfqpoint{2.144755in}{1.073731in}}%
\pgfpathlineto{\pgfqpoint{2.152373in}{1.076112in}}%
\pgfpathlineto{\pgfqpoint{2.156206in}{1.071747in}}%
\pgfpathlineto{\pgfqpoint{2.158994in}{1.078092in}}%
\pgfpathlineto{\pgfqpoint{2.165193in}{1.079905in}}%
\pgfpathlineto{\pgfqpoint{2.172010in}{1.077414in}}%
\pgfpathlineto{\pgfqpoint{2.171424in}{1.072069in}}%
\pgfpathlineto{\pgfqpoint{2.164342in}{1.057050in}}%
\pgfpathlineto{\pgfqpoint{2.170940in}{1.054187in}}%
\pgfpathlineto{\pgfqpoint{2.172018in}{1.046217in}}%
\pgfpathlineto{\pgfqpoint{2.174985in}{1.043028in}}%
\pgfpathlineto{\pgfqpoint{2.171148in}{1.034806in}}%
\pgfpathlineto{\pgfqpoint{2.174509in}{1.028592in}}%
\pgfpathclose%
\pgfusepath{fill}%
\end{pgfscope}%
\begin{pgfscope}%
\pgfpathrectangle{\pgfqpoint{0.100000in}{0.100000in}}{\pgfqpoint{3.007045in}{1.925000in}}%
\pgfusepath{clip}%
\pgfsetbuttcap%
\pgfsetmiterjoin%
\definecolor{currentfill}{rgb}{0.171027,0.484752,0.731242}%
\pgfsetfillcolor{currentfill}%
\pgfsetlinewidth{0.000000pt}%
\definecolor{currentstroke}{rgb}{0.000000,0.000000,0.000000}%
\pgfsetstrokecolor{currentstroke}%
\pgfsetstrokeopacity{0.000000}%
\pgfsetdash{}{0pt}%
\pgfpathmoveto{\pgfqpoint{1.381806in}{1.702240in}}%
\pgfpathlineto{\pgfqpoint{1.379578in}{1.679082in}}%
\pgfpathlineto{\pgfqpoint{1.382174in}{1.678842in}}%
\pgfpathlineto{\pgfqpoint{1.379987in}{1.655666in}}%
\pgfpathlineto{\pgfqpoint{1.371277in}{1.656487in}}%
\pgfpathlineto{\pgfqpoint{1.370707in}{1.650612in}}%
\pgfpathlineto{\pgfqpoint{1.359600in}{1.651719in}}%
\pgfpathlineto{\pgfqpoint{1.360257in}{1.658366in}}%
\pgfpathlineto{\pgfqpoint{1.346246in}{1.659789in}}%
\pgfpathlineto{\pgfqpoint{1.343644in}{1.662962in}}%
\pgfpathlineto{\pgfqpoint{1.335063in}{1.663866in}}%
\pgfpathlineto{\pgfqpoint{1.336444in}{1.675484in}}%
\pgfpathlineto{\pgfqpoint{1.317293in}{1.677601in}}%
\pgfpathlineto{\pgfqpoint{1.315821in}{1.681622in}}%
\pgfpathlineto{\pgfqpoint{1.310066in}{1.682269in}}%
\pgfpathlineto{\pgfqpoint{1.310508in}{1.686101in}}%
\pgfpathlineto{\pgfqpoint{1.304728in}{1.686756in}}%
\pgfpathlineto{\pgfqpoint{1.306273in}{1.700261in}}%
\pgfpathlineto{\pgfqpoint{1.302614in}{1.700688in}}%
\pgfpathlineto{\pgfqpoint{1.305247in}{1.723687in}}%
\pgfpathlineto{\pgfqpoint{1.307173in}{1.723478in}}%
\pgfpathlineto{\pgfqpoint{1.309172in}{1.740759in}}%
\pgfpathlineto{\pgfqpoint{1.314888in}{1.740131in}}%
\pgfpathlineto{\pgfqpoint{1.314224in}{1.734361in}}%
\pgfpathlineto{\pgfqpoint{1.325673in}{1.733068in}}%
\pgfpathlineto{\pgfqpoint{1.325013in}{1.727299in}}%
\pgfpathlineto{\pgfqpoint{1.342199in}{1.725419in}}%
\pgfpathlineto{\pgfqpoint{1.341580in}{1.719656in}}%
\pgfpathlineto{\pgfqpoint{1.345202in}{1.719293in}}%
\pgfpathlineto{\pgfqpoint{1.343969in}{1.707734in}}%
\pgfpathlineto{\pgfqpoint{1.361096in}{1.705933in}}%
\pgfpathlineto{\pgfqpoint{1.364804in}{1.703953in}}%
\pgfpathlineto{\pgfqpoint{1.381806in}{1.702240in}}%
\pgfpathclose%
\pgfusepath{fill}%
\end{pgfscope}%
\begin{pgfscope}%
\pgfpathrectangle{\pgfqpoint{0.100000in}{0.100000in}}{\pgfqpoint{3.007045in}{1.925000in}}%
\pgfusepath{clip}%
\pgfsetbuttcap%
\pgfsetmiterjoin%
\definecolor{currentfill}{rgb}{0.598431,0.780531,0.877785}%
\pgfsetfillcolor{currentfill}%
\pgfsetlinewidth{0.000000pt}%
\definecolor{currentstroke}{rgb}{0.000000,0.000000,0.000000}%
\pgfsetstrokecolor{currentstroke}%
\pgfsetstrokeopacity{0.000000}%
\pgfsetdash{}{0pt}%
\pgfpathmoveto{\pgfqpoint{2.128961in}{1.204856in}}%
\pgfpathlineto{\pgfqpoint{2.128488in}{1.192008in}}%
\pgfpathlineto{\pgfqpoint{2.123761in}{1.189439in}}%
\pgfpathlineto{\pgfqpoint{2.103104in}{1.187835in}}%
\pgfpathlineto{\pgfqpoint{2.105221in}{1.159153in}}%
\pgfpathlineto{\pgfqpoint{2.071143in}{1.157025in}}%
\pgfpathlineto{\pgfqpoint{2.070604in}{1.165647in}}%
\pgfpathlineto{\pgfqpoint{2.076427in}{1.165836in}}%
\pgfpathlineto{\pgfqpoint{2.075079in}{1.186079in}}%
\pgfpathlineto{\pgfqpoint{2.069337in}{1.185778in}}%
\pgfpathlineto{\pgfqpoint{2.068581in}{1.195334in}}%
\pgfpathlineto{\pgfqpoint{2.064705in}{1.195908in}}%
\pgfpathlineto{\pgfqpoint{2.064257in}{1.202733in}}%
\pgfpathlineto{\pgfqpoint{2.068009in}{1.203031in}}%
\pgfpathlineto{\pgfqpoint{2.067386in}{1.211625in}}%
\pgfpathlineto{\pgfqpoint{2.087435in}{1.213369in}}%
\pgfpathlineto{\pgfqpoint{2.094960in}{1.228863in}}%
\pgfpathlineto{\pgfqpoint{2.100744in}{1.229280in}}%
\pgfpathlineto{\pgfqpoint{2.099165in}{1.251343in}}%
\pgfpathlineto{\pgfqpoint{2.110362in}{1.252216in}}%
\pgfpathlineto{\pgfqpoint{2.107905in}{1.276923in}}%
\pgfpathlineto{\pgfqpoint{2.113613in}{1.277542in}}%
\pgfpathlineto{\pgfqpoint{2.116838in}{1.244113in}}%
\pgfpathlineto{\pgfqpoint{2.125964in}{1.244692in}}%
\pgfpathlineto{\pgfqpoint{2.126956in}{1.227545in}}%
\pgfpathlineto{\pgfqpoint{2.128961in}{1.204856in}}%
\pgfpathclose%
\pgfusepath{fill}%
\end{pgfscope}%
\begin{pgfscope}%
\pgfpathrectangle{\pgfqpoint{0.100000in}{0.100000in}}{\pgfqpoint{3.007045in}{1.925000in}}%
\pgfusepath{clip}%
\pgfsetbuttcap%
\pgfsetmiterjoin%
\definecolor{currentfill}{rgb}{0.171027,0.484752,0.731242}%
\pgfsetfillcolor{currentfill}%
\pgfsetlinewidth{0.000000pt}%
\definecolor{currentstroke}{rgb}{0.000000,0.000000,0.000000}%
\pgfsetstrokecolor{currentstroke}%
\pgfsetstrokeopacity{0.000000}%
\pgfsetdash{}{0pt}%
\pgfpathmoveto{\pgfqpoint{1.381806in}{1.702240in}}%
\pgfpathlineto{\pgfqpoint{1.364804in}{1.703953in}}%
\pgfpathlineto{\pgfqpoint{1.361096in}{1.705933in}}%
\pgfpathlineto{\pgfqpoint{1.343969in}{1.707734in}}%
\pgfpathlineto{\pgfqpoint{1.345202in}{1.719293in}}%
\pgfpathlineto{\pgfqpoint{1.341580in}{1.719656in}}%
\pgfpathlineto{\pgfqpoint{1.342199in}{1.725419in}}%
\pgfpathlineto{\pgfqpoint{1.325013in}{1.727299in}}%
\pgfpathlineto{\pgfqpoint{1.325673in}{1.733068in}}%
\pgfpathlineto{\pgfqpoint{1.314224in}{1.734361in}}%
\pgfpathlineto{\pgfqpoint{1.314888in}{1.740131in}}%
\pgfpathlineto{\pgfqpoint{1.317232in}{1.745061in}}%
\pgfpathlineto{\pgfqpoint{1.318715in}{1.758150in}}%
\pgfpathlineto{\pgfqpoint{1.324957in}{1.756945in}}%
\pgfpathlineto{\pgfqpoint{1.335725in}{1.759869in}}%
\pgfpathlineto{\pgfqpoint{1.341818in}{1.758120in}}%
\pgfpathlineto{\pgfqpoint{1.351556in}{1.758299in}}%
\pgfpathlineto{\pgfqpoint{1.355902in}{1.752996in}}%
\pgfpathlineto{\pgfqpoint{1.366019in}{1.751221in}}%
\pgfpathlineto{\pgfqpoint{1.369270in}{1.748112in}}%
\pgfpathlineto{\pgfqpoint{1.376795in}{1.745675in}}%
\pgfpathlineto{\pgfqpoint{1.386114in}{1.754885in}}%
\pgfpathlineto{\pgfqpoint{1.389403in}{1.754089in}}%
\pgfpathlineto{\pgfqpoint{1.390611in}{1.748693in}}%
\pgfpathlineto{\pgfqpoint{1.394175in}{1.746414in}}%
\pgfpathlineto{\pgfqpoint{1.403722in}{1.748968in}}%
\pgfpathlineto{\pgfqpoint{1.407533in}{1.753794in}}%
\pgfpathlineto{\pgfqpoint{1.422411in}{1.752597in}}%
\pgfpathlineto{\pgfqpoint{1.425483in}{1.768138in}}%
\pgfpathlineto{\pgfqpoint{1.422948in}{1.768339in}}%
\pgfpathlineto{\pgfqpoint{1.423931in}{1.779899in}}%
\pgfpathlineto{\pgfqpoint{1.452658in}{1.777477in}}%
\pgfpathlineto{\pgfqpoint{1.464168in}{1.776657in}}%
\pgfpathlineto{\pgfqpoint{1.463268in}{1.765000in}}%
\pgfpathlineto{\pgfqpoint{1.465533in}{1.764842in}}%
\pgfpathlineto{\pgfqpoint{1.463751in}{1.741609in}}%
\pgfpathlineto{\pgfqpoint{1.466009in}{1.741430in}}%
\pgfpathlineto{\pgfqpoint{1.465127in}{1.729765in}}%
\pgfpathlineto{\pgfqpoint{1.442195in}{1.731563in}}%
\pgfpathlineto{\pgfqpoint{1.441706in}{1.725605in}}%
\pgfpathlineto{\pgfqpoint{1.433959in}{1.727722in}}%
\pgfpathlineto{\pgfqpoint{1.429757in}{1.720971in}}%
\pgfpathlineto{\pgfqpoint{1.409324in}{1.722778in}}%
\pgfpathlineto{\pgfqpoint{1.407305in}{1.699840in}}%
\pgfpathlineto{\pgfqpoint{1.381806in}{1.702240in}}%
\pgfpathclose%
\pgfusepath{fill}%
\end{pgfscope}%
\begin{pgfscope}%
\pgfpathrectangle{\pgfqpoint{0.100000in}{0.100000in}}{\pgfqpoint{3.007045in}{1.925000in}}%
\pgfusepath{clip}%
\pgfsetbuttcap%
\pgfsetmiterjoin%
\definecolor{currentfill}{rgb}{0.968627,0.984314,1.000000}%
\pgfsetfillcolor{currentfill}%
\pgfsetlinewidth{0.000000pt}%
\definecolor{currentstroke}{rgb}{0.000000,0.000000,0.000000}%
\pgfsetstrokecolor{currentstroke}%
\pgfsetstrokeopacity{0.000000}%
\pgfsetdash{}{0pt}%
\pgfpathmoveto{\pgfqpoint{2.962613in}{1.457141in}}%
\pgfpathlineto{\pgfqpoint{2.972406in}{1.450863in}}%
\pgfpathlineto{\pgfqpoint{2.960151in}{1.447007in}}%
\pgfpathlineto{\pgfqpoint{2.958968in}{1.452964in}}%
\pgfpathlineto{\pgfqpoint{2.962613in}{1.457141in}}%
\pgfpathclose%
\pgfusepath{fill}%
\end{pgfscope}%
\begin{pgfscope}%
\pgfpathrectangle{\pgfqpoint{0.100000in}{0.100000in}}{\pgfqpoint{3.007045in}{1.925000in}}%
\pgfusepath{clip}%
\pgfsetbuttcap%
\pgfsetmiterjoin%
\definecolor{currentfill}{rgb}{0.617255,0.790865,0.881845}%
\pgfsetfillcolor{currentfill}%
\pgfsetlinewidth{0.000000pt}%
\definecolor{currentstroke}{rgb}{0.000000,0.000000,0.000000}%
\pgfsetstrokecolor{currentstroke}%
\pgfsetstrokeopacity{0.000000}%
\pgfsetdash{}{0pt}%
\pgfpathmoveto{\pgfqpoint{2.644603in}{1.131634in}}%
\pgfpathlineto{\pgfqpoint{2.641185in}{1.127114in}}%
\pgfpathlineto{\pgfqpoint{2.641373in}{1.123470in}}%
\pgfpathlineto{\pgfqpoint{2.633816in}{1.102244in}}%
\pgfpathlineto{\pgfqpoint{2.626460in}{1.104802in}}%
\pgfpathlineto{\pgfqpoint{2.621502in}{1.105990in}}%
\pgfpathlineto{\pgfqpoint{2.612770in}{1.112672in}}%
\pgfpathlineto{\pgfqpoint{2.613630in}{1.118818in}}%
\pgfpathlineto{\pgfqpoint{2.617784in}{1.122931in}}%
\pgfpathlineto{\pgfqpoint{2.616773in}{1.129743in}}%
\pgfpathlineto{\pgfqpoint{2.606282in}{1.144642in}}%
\pgfpathlineto{\pgfqpoint{2.608923in}{1.148127in}}%
\pgfpathlineto{\pgfqpoint{2.609006in}{1.155780in}}%
\pgfpathlineto{\pgfqpoint{2.613166in}{1.160149in}}%
\pgfpathlineto{\pgfqpoint{2.619686in}{1.171111in}}%
\pgfpathlineto{\pgfqpoint{2.621002in}{1.178691in}}%
\pgfpathlineto{\pgfqpoint{2.624882in}{1.186857in}}%
\pgfpathlineto{\pgfqpoint{2.631429in}{1.181706in}}%
\pgfpathlineto{\pgfqpoint{2.634662in}{1.181683in}}%
\pgfpathlineto{\pgfqpoint{2.644402in}{1.195106in}}%
\pgfpathlineto{\pgfqpoint{2.648785in}{1.188481in}}%
\pgfpathlineto{\pgfqpoint{2.658001in}{1.178784in}}%
\pgfpathlineto{\pgfqpoint{2.663143in}{1.179525in}}%
\pgfpathlineto{\pgfqpoint{2.650904in}{1.157743in}}%
\pgfpathlineto{\pgfqpoint{2.656779in}{1.157595in}}%
\pgfpathlineto{\pgfqpoint{2.665522in}{1.153025in}}%
\pgfpathlineto{\pgfqpoint{2.663369in}{1.133914in}}%
\pgfpathlineto{\pgfqpoint{2.658319in}{1.135059in}}%
\pgfpathlineto{\pgfqpoint{2.654207in}{1.139998in}}%
\pgfpathlineto{\pgfqpoint{2.647609in}{1.142718in}}%
\pgfpathlineto{\pgfqpoint{2.644603in}{1.131634in}}%
\pgfpathclose%
\pgfusepath{fill}%
\end{pgfscope}%
\begin{pgfscope}%
\pgfpathrectangle{\pgfqpoint{0.100000in}{0.100000in}}{\pgfqpoint{3.007045in}{1.925000in}}%
\pgfusepath{clip}%
\pgfsetbuttcap%
\pgfsetmiterjoin%
\definecolor{currentfill}{rgb}{0.199446,0.513172,0.745882}%
\pgfsetfillcolor{currentfill}%
\pgfsetlinewidth{0.000000pt}%
\definecolor{currentstroke}{rgb}{0.000000,0.000000,0.000000}%
\pgfsetstrokecolor{currentstroke}%
\pgfsetstrokeopacity{0.000000}%
\pgfsetdash{}{0pt}%
\pgfpathmoveto{\pgfqpoint{1.444186in}{0.892364in}}%
\pgfpathlineto{\pgfqpoint{1.442027in}{0.863792in}}%
\pgfpathlineto{\pgfqpoint{1.421405in}{0.865168in}}%
\pgfpathlineto{\pgfqpoint{1.392891in}{0.867332in}}%
\pgfpathlineto{\pgfqpoint{1.364391in}{0.869531in}}%
\pgfpathlineto{\pgfqpoint{1.336407in}{0.871962in}}%
\pgfpathlineto{\pgfqpoint{1.337618in}{0.885562in}}%
\pgfpathlineto{\pgfqpoint{1.341578in}{0.929505in}}%
\pgfpathlineto{\pgfqpoint{1.388570in}{0.925898in}}%
\pgfpathlineto{\pgfqpoint{1.417413in}{0.923222in}}%
\pgfpathlineto{\pgfqpoint{1.415306in}{0.894451in}}%
\pgfpathlineto{\pgfqpoint{1.444186in}{0.892364in}}%
\pgfpathclose%
\pgfusepath{fill}%
\end{pgfscope}%
\begin{pgfscope}%
\pgfpathrectangle{\pgfqpoint{0.100000in}{0.100000in}}{\pgfqpoint{3.007045in}{1.925000in}}%
\pgfusepath{clip}%
\pgfsetbuttcap%
\pgfsetmiterjoin%
\definecolor{currentfill}{rgb}{0.326290,0.618624,0.802799}%
\pgfsetfillcolor{currentfill}%
\pgfsetlinewidth{0.000000pt}%
\definecolor{currentstroke}{rgb}{0.000000,0.000000,0.000000}%
\pgfsetstrokecolor{currentstroke}%
\pgfsetstrokeopacity{0.000000}%
\pgfsetdash{}{0pt}%
\pgfpathmoveto{\pgfqpoint{1.918185in}{0.809351in}}%
\pgfpathlineto{\pgfqpoint{1.909437in}{0.814697in}}%
\pgfpathlineto{\pgfqpoint{1.899320in}{0.814784in}}%
\pgfpathlineto{\pgfqpoint{1.899400in}{0.832281in}}%
\pgfpathlineto{\pgfqpoint{1.898412in}{0.836142in}}%
\pgfpathlineto{\pgfqpoint{1.892466in}{0.838016in}}%
\pgfpathlineto{\pgfqpoint{1.891528in}{0.843729in}}%
\pgfpathlineto{\pgfqpoint{1.887792in}{0.846596in}}%
\pgfpathlineto{\pgfqpoint{1.882114in}{0.846605in}}%
\pgfpathlineto{\pgfqpoint{1.882494in}{0.855536in}}%
\pgfpathlineto{\pgfqpoint{1.865214in}{0.855241in}}%
\pgfpathlineto{\pgfqpoint{1.864117in}{0.861985in}}%
\pgfpathlineto{\pgfqpoint{1.871864in}{0.868767in}}%
\pgfpathlineto{\pgfqpoint{1.884196in}{0.881637in}}%
\pgfpathlineto{\pgfqpoint{1.888066in}{0.882118in}}%
\pgfpathlineto{\pgfqpoint{1.887849in}{0.901272in}}%
\pgfpathlineto{\pgfqpoint{1.907800in}{0.901293in}}%
\pgfpathlineto{\pgfqpoint{1.907859in}{0.895543in}}%
\pgfpathlineto{\pgfqpoint{1.927775in}{0.895762in}}%
\pgfpathlineto{\pgfqpoint{1.927950in}{0.876198in}}%
\pgfpathlineto{\pgfqpoint{1.933967in}{0.876483in}}%
\pgfpathlineto{\pgfqpoint{1.941240in}{0.873795in}}%
\pgfpathlineto{\pgfqpoint{1.945167in}{0.874395in}}%
\pgfpathlineto{\pgfqpoint{1.945345in}{0.867212in}}%
\pgfpathlineto{\pgfqpoint{1.951312in}{0.867300in}}%
\pgfpathlineto{\pgfqpoint{1.951519in}{0.850943in}}%
\pgfpathlineto{\pgfqpoint{1.955013in}{0.847952in}}%
\pgfpathlineto{\pgfqpoint{1.955232in}{0.842505in}}%
\pgfpathlineto{\pgfqpoint{1.951736in}{0.838426in}}%
\pgfpathlineto{\pgfqpoint{1.924362in}{0.838176in}}%
\pgfpathlineto{\pgfqpoint{1.924611in}{0.820832in}}%
\pgfpathlineto{\pgfqpoint{1.923831in}{0.809696in}}%
\pgfpathlineto{\pgfqpoint{1.918185in}{0.809351in}}%
\pgfpathclose%
\pgfusepath{fill}%
\end{pgfscope}%
\begin{pgfscope}%
\pgfpathrectangle{\pgfqpoint{0.100000in}{0.100000in}}{\pgfqpoint{3.007045in}{1.925000in}}%
\pgfusepath{clip}%
\pgfsetbuttcap%
\pgfsetmiterjoin%
\definecolor{currentfill}{rgb}{0.183206,0.496932,0.737516}%
\pgfsetfillcolor{currentfill}%
\pgfsetlinewidth{0.000000pt}%
\definecolor{currentstroke}{rgb}{0.000000,0.000000,0.000000}%
\pgfsetstrokecolor{currentstroke}%
\pgfsetstrokeopacity{0.000000}%
\pgfsetdash{}{0pt}%
\pgfpathmoveto{\pgfqpoint{1.311526in}{1.606441in}}%
\pgfpathlineto{\pgfqpoint{1.308128in}{1.577801in}}%
\pgfpathlineto{\pgfqpoint{1.305764in}{1.578064in}}%
\pgfpathlineto{\pgfqpoint{1.303187in}{1.554906in}}%
\pgfpathlineto{\pgfqpoint{1.262489in}{1.559658in}}%
\pgfpathlineto{\pgfqpoint{1.257251in}{1.559802in}}%
\pgfpathlineto{\pgfqpoint{1.246154in}{1.561183in}}%
\pgfpathlineto{\pgfqpoint{1.247672in}{1.573410in}}%
\pgfpathlineto{\pgfqpoint{1.224309in}{1.576380in}}%
\pgfpathlineto{\pgfqpoint{1.227555in}{1.587451in}}%
\pgfpathlineto{\pgfqpoint{1.230291in}{1.608576in}}%
\pgfpathlineto{\pgfqpoint{1.221965in}{1.610407in}}%
\pgfpathlineto{\pgfqpoint{1.222383in}{1.622797in}}%
\pgfpathlineto{\pgfqpoint{1.212290in}{1.629948in}}%
\pgfpathlineto{\pgfqpoint{1.200926in}{1.631495in}}%
\pgfpathlineto{\pgfqpoint{1.201711in}{1.637190in}}%
\pgfpathlineto{\pgfqpoint{1.197757in}{1.637726in}}%
\pgfpathlineto{\pgfqpoint{1.201035in}{1.646271in}}%
\pgfpathlineto{\pgfqpoint{1.195568in}{1.655776in}}%
\pgfpathlineto{\pgfqpoint{1.192292in}{1.666923in}}%
\pgfpathlineto{\pgfqpoint{1.188964in}{1.669306in}}%
\pgfpathlineto{\pgfqpoint{1.190856in}{1.681412in}}%
\pgfpathlineto{\pgfqpoint{1.189998in}{1.686616in}}%
\pgfpathlineto{\pgfqpoint{1.187986in}{1.693328in}}%
\pgfpathlineto{\pgfqpoint{1.240870in}{1.686663in}}%
\pgfpathlineto{\pgfqpoint{1.240737in}{1.685718in}}%
\pgfpathlineto{\pgfqpoint{1.269340in}{1.682158in}}%
\pgfpathlineto{\pgfqpoint{1.267668in}{1.681381in}}%
\pgfpathlineto{\pgfqpoint{1.264869in}{1.658439in}}%
\pgfpathlineto{\pgfqpoint{1.263286in}{1.658626in}}%
\pgfpathlineto{\pgfqpoint{1.260561in}{1.635790in}}%
\pgfpathlineto{\pgfqpoint{1.258817in}{1.635994in}}%
\pgfpathlineto{\pgfqpoint{1.255957in}{1.613050in}}%
\pgfpathlineto{\pgfqpoint{1.311526in}{1.606441in}}%
\pgfpathclose%
\pgfusepath{fill}%
\end{pgfscope}%
\begin{pgfscope}%
\pgfpathrectangle{\pgfqpoint{0.100000in}{0.100000in}}{\pgfqpoint{3.007045in}{1.925000in}}%
\pgfusepath{clip}%
\pgfsetbuttcap%
\pgfsetmiterjoin%
\definecolor{currentfill}{rgb}{0.290980,0.594510,0.789020}%
\pgfsetfillcolor{currentfill}%
\pgfsetlinewidth{0.000000pt}%
\definecolor{currentstroke}{rgb}{0.000000,0.000000,0.000000}%
\pgfsetstrokecolor{currentstroke}%
\pgfsetstrokeopacity{0.000000}%
\pgfsetdash{}{0pt}%
\pgfpathmoveto{\pgfqpoint{2.498935in}{0.687218in}}%
\pgfpathlineto{\pgfqpoint{2.492063in}{0.691931in}}%
\pgfpathlineto{\pgfqpoint{2.477917in}{0.698299in}}%
\pgfpathlineto{\pgfqpoint{2.457652in}{0.694920in}}%
\pgfpathlineto{\pgfqpoint{2.455203in}{0.702082in}}%
\pgfpathlineto{\pgfqpoint{2.447789in}{0.700910in}}%
\pgfpathlineto{\pgfqpoint{2.437033in}{0.704008in}}%
\pgfpathlineto{\pgfqpoint{2.433519in}{0.709598in}}%
\pgfpathlineto{\pgfqpoint{2.429949in}{0.711648in}}%
\pgfpathlineto{\pgfqpoint{2.428514in}{0.718344in}}%
\pgfpathlineto{\pgfqpoint{2.425917in}{0.719822in}}%
\pgfpathlineto{\pgfqpoint{2.427863in}{0.724959in}}%
\pgfpathlineto{\pgfqpoint{2.422656in}{0.729020in}}%
\pgfpathlineto{\pgfqpoint{2.425722in}{0.735586in}}%
\pgfpathlineto{\pgfqpoint{2.432076in}{0.743806in}}%
\pgfpathlineto{\pgfqpoint{2.434230in}{0.742118in}}%
\pgfpathlineto{\pgfqpoint{2.444853in}{0.725072in}}%
\pgfpathlineto{\pgfqpoint{2.451822in}{0.729156in}}%
\pgfpathlineto{\pgfqpoint{2.458380in}{0.737732in}}%
\pgfpathlineto{\pgfqpoint{2.457386in}{0.739684in}}%
\pgfpathlineto{\pgfqpoint{2.460691in}{0.751658in}}%
\pgfpathlineto{\pgfqpoint{2.469526in}{0.751792in}}%
\pgfpathlineto{\pgfqpoint{2.476049in}{0.748944in}}%
\pgfpathlineto{\pgfqpoint{2.475118in}{0.743801in}}%
\pgfpathlineto{\pgfqpoint{2.478955in}{0.739426in}}%
\pgfpathlineto{\pgfqpoint{2.485433in}{0.742115in}}%
\pgfpathlineto{\pgfqpoint{2.489704in}{0.732476in}}%
\pgfpathlineto{\pgfqpoint{2.489650in}{0.715531in}}%
\pgfpathlineto{\pgfqpoint{2.496462in}{0.714530in}}%
\pgfpathlineto{\pgfqpoint{2.499221in}{0.711875in}}%
\pgfpathlineto{\pgfqpoint{2.496078in}{0.707079in}}%
\pgfpathlineto{\pgfqpoint{2.498935in}{0.687218in}}%
\pgfpathclose%
\pgfusepath{fill}%
\end{pgfscope}%
\begin{pgfscope}%
\pgfpathrectangle{\pgfqpoint{0.100000in}{0.100000in}}{\pgfqpoint{3.007045in}{1.925000in}}%
\pgfusepath{clip}%
\pgfsetbuttcap%
\pgfsetmiterjoin%
\definecolor{currentfill}{rgb}{0.429020,0.687520,0.841246}%
\pgfsetfillcolor{currentfill}%
\pgfsetlinewidth{0.000000pt}%
\definecolor{currentstroke}{rgb}{0.000000,0.000000,0.000000}%
\pgfsetstrokecolor{currentstroke}%
\pgfsetstrokeopacity{0.000000}%
\pgfsetdash{}{0pt}%
\pgfpathmoveto{\pgfqpoint{2.741244in}{1.095623in}}%
\pgfpathlineto{\pgfqpoint{2.738779in}{1.094235in}}%
\pgfpathlineto{\pgfqpoint{2.723519in}{1.072852in}}%
\pgfpathlineto{\pgfqpoint{2.723871in}{1.066185in}}%
\pgfpathlineto{\pgfqpoint{2.701345in}{1.062144in}}%
\pgfpathlineto{\pgfqpoint{2.695089in}{1.071478in}}%
\pgfpathlineto{\pgfqpoint{2.716714in}{1.091684in}}%
\pgfpathlineto{\pgfqpoint{2.713742in}{1.098160in}}%
\pgfpathlineto{\pgfqpoint{2.704080in}{1.100551in}}%
\pgfpathlineto{\pgfqpoint{2.711789in}{1.111134in}}%
\pgfpathlineto{\pgfqpoint{2.716713in}{1.110416in}}%
\pgfpathlineto{\pgfqpoint{2.713382in}{1.123521in}}%
\pgfpathlineto{\pgfqpoint{2.724698in}{1.127666in}}%
\pgfpathlineto{\pgfqpoint{2.723470in}{1.137350in}}%
\pgfpathlineto{\pgfqpoint{2.716860in}{1.144435in}}%
\pgfpathlineto{\pgfqpoint{2.720436in}{1.148513in}}%
\pgfpathlineto{\pgfqpoint{2.725113in}{1.146693in}}%
\pgfpathlineto{\pgfqpoint{2.729649in}{1.140417in}}%
\pgfpathlineto{\pgfqpoint{2.735476in}{1.140922in}}%
\pgfpathlineto{\pgfqpoint{2.739950in}{1.134391in}}%
\pgfpathlineto{\pgfqpoint{2.744726in}{1.133904in}}%
\pgfpathlineto{\pgfqpoint{2.747294in}{1.127106in}}%
\pgfpathlineto{\pgfqpoint{2.744962in}{1.123172in}}%
\pgfpathlineto{\pgfqpoint{2.740631in}{1.125917in}}%
\pgfpathlineto{\pgfqpoint{2.742495in}{1.113720in}}%
\pgfpathlineto{\pgfqpoint{2.750051in}{1.107970in}}%
\pgfpathlineto{\pgfqpoint{2.749079in}{1.102041in}}%
\pgfpathlineto{\pgfqpoint{2.745560in}{1.100962in}}%
\pgfpathlineto{\pgfqpoint{2.741244in}{1.095623in}}%
\pgfpathclose%
\pgfusepath{fill}%
\end{pgfscope}%
\begin{pgfscope}%
\pgfpathrectangle{\pgfqpoint{0.100000in}{0.100000in}}{\pgfqpoint{3.007045in}{1.925000in}}%
\pgfusepath{clip}%
\pgfsetbuttcap%
\pgfsetmiterjoin%
\definecolor{currentfill}{rgb}{0.516863,0.735748,0.860192}%
\pgfsetfillcolor{currentfill}%
\pgfsetlinewidth{0.000000pt}%
\definecolor{currentstroke}{rgb}{0.000000,0.000000,0.000000}%
\pgfsetstrokecolor{currentstroke}%
\pgfsetstrokeopacity{0.000000}%
\pgfsetdash{}{0pt}%
\pgfpathmoveto{\pgfqpoint{2.386694in}{0.571228in}}%
\pgfpathlineto{\pgfqpoint{2.386777in}{0.566973in}}%
\pgfpathlineto{\pgfqpoint{2.379860in}{0.568127in}}%
\pgfpathlineto{\pgfqpoint{2.372613in}{0.564198in}}%
\pgfpathlineto{\pgfqpoint{2.358211in}{0.553384in}}%
\pgfpathlineto{\pgfqpoint{2.349321in}{0.550333in}}%
\pgfpathlineto{\pgfqpoint{2.344406in}{0.550144in}}%
\pgfpathlineto{\pgfqpoint{2.336004in}{0.546671in}}%
\pgfpathlineto{\pgfqpoint{2.332928in}{0.547891in}}%
\pgfpathlineto{\pgfqpoint{2.332519in}{0.554903in}}%
\pgfpathlineto{\pgfqpoint{2.324536in}{0.563497in}}%
\pgfpathlineto{\pgfqpoint{2.320470in}{0.564006in}}%
\pgfpathlineto{\pgfqpoint{2.314549in}{0.569015in}}%
\pgfpathlineto{\pgfqpoint{2.302272in}{0.575567in}}%
\pgfpathlineto{\pgfqpoint{2.289539in}{0.581381in}}%
\pgfpathlineto{\pgfqpoint{2.288959in}{0.589362in}}%
\pgfpathlineto{\pgfqpoint{2.296561in}{0.593554in}}%
\pgfpathlineto{\pgfqpoint{2.317482in}{0.595459in}}%
\pgfpathlineto{\pgfqpoint{2.316671in}{0.604111in}}%
\pgfpathlineto{\pgfqpoint{2.322486in}{0.604682in}}%
\pgfpathlineto{\pgfqpoint{2.324742in}{0.580445in}}%
\pgfpathlineto{\pgfqpoint{2.339480in}{0.580648in}}%
\pgfpathlineto{\pgfqpoint{2.340428in}{0.571556in}}%
\pgfpathlineto{\pgfqpoint{2.347979in}{0.570392in}}%
\pgfpathlineto{\pgfqpoint{2.374564in}{0.573201in}}%
\pgfpathlineto{\pgfqpoint{2.386694in}{0.571228in}}%
\pgfpathclose%
\pgfusepath{fill}%
\end{pgfscope}%
\begin{pgfscope}%
\pgfpathrectangle{\pgfqpoint{0.100000in}{0.100000in}}{\pgfqpoint{3.007045in}{1.925000in}}%
\pgfusepath{clip}%
\pgfsetbuttcap%
\pgfsetmiterjoin%
\definecolor{currentfill}{rgb}{0.652211,0.806013,0.893764}%
\pgfsetfillcolor{currentfill}%
\pgfsetlinewidth{0.000000pt}%
\definecolor{currentstroke}{rgb}{0.000000,0.000000,0.000000}%
\pgfsetstrokecolor{currentstroke}%
\pgfsetstrokeopacity{0.000000}%
\pgfsetdash{}{0pt}%
\pgfpathmoveto{\pgfqpoint{2.578534in}{1.060078in}}%
\pgfpathlineto{\pgfqpoint{2.578144in}{1.039001in}}%
\pgfpathlineto{\pgfqpoint{2.561819in}{1.036295in}}%
\pgfpathlineto{\pgfqpoint{2.547798in}{1.034084in}}%
\pgfpathlineto{\pgfqpoint{2.531083in}{1.032312in}}%
\pgfpathlineto{\pgfqpoint{2.529927in}{1.036274in}}%
\pgfpathlineto{\pgfqpoint{2.537514in}{1.041038in}}%
\pgfpathlineto{\pgfqpoint{2.537236in}{1.043373in}}%
\pgfpathlineto{\pgfqpoint{2.543457in}{1.053394in}}%
\pgfpathlineto{\pgfqpoint{2.547424in}{1.055944in}}%
\pgfpathlineto{\pgfqpoint{2.572984in}{1.054602in}}%
\pgfpathlineto{\pgfqpoint{2.578534in}{1.060078in}}%
\pgfpathclose%
\pgfusepath{fill}%
\end{pgfscope}%
\begin{pgfscope}%
\pgfpathrectangle{\pgfqpoint{0.100000in}{0.100000in}}{\pgfqpoint{3.007045in}{1.925000in}}%
\pgfusepath{clip}%
\pgfsetbuttcap%
\pgfsetmiterjoin%
\definecolor{currentfill}{rgb}{0.311157,0.608289,0.796894}%
\pgfsetfillcolor{currentfill}%
\pgfsetlinewidth{0.000000pt}%
\definecolor{currentstroke}{rgb}{0.000000,0.000000,0.000000}%
\pgfsetstrokecolor{currentstroke}%
\pgfsetstrokeopacity{0.000000}%
\pgfsetdash{}{0pt}%
\pgfpathmoveto{\pgfqpoint{1.286524in}{0.867195in}}%
\pgfpathlineto{\pgfqpoint{1.283287in}{0.832832in}}%
\pgfpathlineto{\pgfqpoint{1.271855in}{0.834010in}}%
\pgfpathlineto{\pgfqpoint{1.271248in}{0.828230in}}%
\pgfpathlineto{\pgfqpoint{1.236948in}{0.832126in}}%
\pgfpathlineto{\pgfqpoint{1.237556in}{0.837894in}}%
\pgfpathlineto{\pgfqpoint{1.231742in}{0.838561in}}%
\pgfpathlineto{\pgfqpoint{1.233670in}{0.855551in}}%
\pgfpathlineto{\pgfqpoint{1.210832in}{0.858083in}}%
\pgfpathlineto{\pgfqpoint{1.212799in}{0.875172in}}%
\pgfpathlineto{\pgfqpoint{1.214018in}{0.875058in}}%
\pgfpathlineto{\pgfqpoint{1.217220in}{0.903588in}}%
\pgfpathlineto{\pgfqpoint{1.218495in}{0.915021in}}%
\pgfpathlineto{\pgfqpoint{1.281006in}{0.908273in}}%
\pgfpathlineto{\pgfqpoint{1.280556in}{0.903442in}}%
\pgfpathlineto{\pgfqpoint{1.276783in}{0.868180in}}%
\pgfpathlineto{\pgfqpoint{1.286524in}{0.867195in}}%
\pgfpathclose%
\pgfusepath{fill}%
\end{pgfscope}%
\begin{pgfscope}%
\pgfpathrectangle{\pgfqpoint{0.100000in}{0.100000in}}{\pgfqpoint{3.007045in}{1.925000in}}%
\pgfusepath{clip}%
\pgfsetbuttcap%
\pgfsetmiterjoin%
\definecolor{currentfill}{rgb}{0.632526,0.797647,0.886874}%
\pgfsetfillcolor{currentfill}%
\pgfsetlinewidth{0.000000pt}%
\definecolor{currentstroke}{rgb}{0.000000,0.000000,0.000000}%
\pgfsetstrokecolor{currentstroke}%
\pgfsetstrokeopacity{0.000000}%
\pgfsetdash{}{0pt}%
\pgfpathmoveto{\pgfqpoint{2.959007in}{1.458513in}}%
\pgfpathlineto{\pgfqpoint{2.957672in}{1.458775in}}%
\pgfpathlineto{\pgfqpoint{2.959346in}{1.462818in}}%
\pgfpathlineto{\pgfqpoint{2.955719in}{1.471405in}}%
\pgfpathlineto{\pgfqpoint{2.948844in}{1.463490in}}%
\pgfpathlineto{\pgfqpoint{2.944479in}{1.471321in}}%
\pgfpathlineto{\pgfqpoint{2.933228in}{1.482168in}}%
\pgfpathlineto{\pgfqpoint{2.929528in}{1.490206in}}%
\pgfpathlineto{\pgfqpoint{2.917977in}{1.479575in}}%
\pgfpathlineto{\pgfqpoint{2.916591in}{1.481492in}}%
\pgfpathlineto{\pgfqpoint{2.896973in}{1.475592in}}%
\pgfpathlineto{\pgfqpoint{2.880647in}{1.472875in}}%
\pgfpathlineto{\pgfqpoint{2.878530in}{1.481285in}}%
\pgfpathlineto{\pgfqpoint{2.872055in}{1.481134in}}%
\pgfpathlineto{\pgfqpoint{2.873175in}{1.485821in}}%
\pgfpathlineto{\pgfqpoint{2.867051in}{1.490726in}}%
\pgfpathlineto{\pgfqpoint{2.865008in}{1.506125in}}%
\pgfpathlineto{\pgfqpoint{2.865846in}{1.512168in}}%
\pgfpathlineto{\pgfqpoint{2.862374in}{1.515277in}}%
\pgfpathlineto{\pgfqpoint{2.879160in}{1.518963in}}%
\pgfpathlineto{\pgfqpoint{2.909180in}{1.525824in}}%
\pgfpathlineto{\pgfqpoint{2.914625in}{1.535783in}}%
\pgfpathlineto{\pgfqpoint{2.922960in}{1.542355in}}%
\pgfpathlineto{\pgfqpoint{2.928476in}{1.542999in}}%
\pgfpathlineto{\pgfqpoint{2.934550in}{1.531865in}}%
\pgfpathlineto{\pgfqpoint{2.942818in}{1.530677in}}%
\pgfpathlineto{\pgfqpoint{2.939653in}{1.526607in}}%
\pgfpathlineto{\pgfqpoint{2.932124in}{1.521534in}}%
\pgfpathlineto{\pgfqpoint{2.933679in}{1.519559in}}%
\pgfpathlineto{\pgfqpoint{2.928376in}{1.511191in}}%
\pgfpathlineto{\pgfqpoint{2.930098in}{1.505248in}}%
\pgfpathlineto{\pgfqpoint{2.935189in}{1.506115in}}%
\pgfpathlineto{\pgfqpoint{2.944880in}{1.501863in}}%
\pgfpathlineto{\pgfqpoint{2.950388in}{1.495521in}}%
\pgfpathlineto{\pgfqpoint{2.949426in}{1.488196in}}%
\pgfpathlineto{\pgfqpoint{2.958190in}{1.486046in}}%
\pgfpathlineto{\pgfqpoint{2.959875in}{1.479333in}}%
\pgfpathlineto{\pgfqpoint{2.966716in}{1.476431in}}%
\pgfpathlineto{\pgfqpoint{2.973137in}{1.477451in}}%
\pgfpathlineto{\pgfqpoint{2.988983in}{1.486082in}}%
\pgfpathlineto{\pgfqpoint{2.991233in}{1.478471in}}%
\pgfpathlineto{\pgfqpoint{2.968516in}{1.467337in}}%
\pgfpathlineto{\pgfqpoint{2.967427in}{1.463763in}}%
\pgfpathlineto{\pgfqpoint{2.959007in}{1.458513in}}%
\pgfpathclose%
\pgfusepath{fill}%
\end{pgfscope}%
\begin{pgfscope}%
\pgfpathrectangle{\pgfqpoint{0.100000in}{0.100000in}}{\pgfqpoint{3.007045in}{1.925000in}}%
\pgfusepath{clip}%
\pgfsetbuttcap%
\pgfsetmiterjoin%
\definecolor{currentfill}{rgb}{0.093272,0.396878,0.673664}%
\pgfsetfillcolor{currentfill}%
\pgfsetlinewidth{0.000000pt}%
\definecolor{currentstroke}{rgb}{0.000000,0.000000,0.000000}%
\pgfsetstrokecolor{currentstroke}%
\pgfsetstrokeopacity{0.000000}%
\pgfsetdash{}{0pt}%
\pgfpathmoveto{\pgfqpoint{0.678803in}{0.558557in}}%
\pgfpathlineto{\pgfqpoint{0.679997in}{0.561480in}}%
\pgfpathlineto{\pgfqpoint{0.680052in}{0.564426in}}%
\pgfpathlineto{\pgfqpoint{0.679031in}{0.568728in}}%
\pgfpathlineto{\pgfqpoint{0.679664in}{0.571354in}}%
\pgfpathlineto{\pgfqpoint{0.678267in}{0.576636in}}%
\pgfpathlineto{\pgfqpoint{0.676653in}{0.578838in}}%
\pgfpathlineto{\pgfqpoint{0.671826in}{0.579804in}}%
\pgfpathlineto{\pgfqpoint{0.671241in}{0.581751in}}%
\pgfpathlineto{\pgfqpoint{0.671704in}{0.585285in}}%
\pgfpathlineto{\pgfqpoint{0.673540in}{0.587189in}}%
\pgfpathlineto{\pgfqpoint{0.675922in}{0.587938in}}%
\pgfpathlineto{\pgfqpoint{0.678005in}{0.589551in}}%
\pgfpathlineto{\pgfqpoint{0.678705in}{0.588737in}}%
\pgfpathlineto{\pgfqpoint{0.677833in}{0.586473in}}%
\pgfpathlineto{\pgfqpoint{0.680118in}{0.580300in}}%
\pgfpathlineto{\pgfqpoint{0.681622in}{0.579560in}}%
\pgfpathlineto{\pgfqpoint{0.686508in}{0.578300in}}%
\pgfpathlineto{\pgfqpoint{0.687618in}{0.575930in}}%
\pgfpathlineto{\pgfqpoint{0.687530in}{0.573360in}}%
\pgfpathlineto{\pgfqpoint{0.685839in}{0.571449in}}%
\pgfpathlineto{\pgfqpoint{0.687044in}{0.567758in}}%
\pgfpathlineto{\pgfqpoint{0.687213in}{0.565076in}}%
\pgfpathlineto{\pgfqpoint{0.688961in}{0.561585in}}%
\pgfpathlineto{\pgfqpoint{0.691199in}{0.560100in}}%
\pgfpathlineto{\pgfqpoint{0.692475in}{0.557535in}}%
\pgfpathlineto{\pgfqpoint{0.691676in}{0.556369in}}%
\pgfpathlineto{\pgfqpoint{0.688983in}{0.555610in}}%
\pgfpathlineto{\pgfqpoint{0.687659in}{0.557811in}}%
\pgfpathlineto{\pgfqpoint{0.685611in}{0.559618in}}%
\pgfpathlineto{\pgfqpoint{0.682975in}{0.560083in}}%
\pgfpathlineto{\pgfqpoint{0.678803in}{0.558557in}}%
\pgfpathclose%
\pgfusepath{fill}%
\end{pgfscope}%
\begin{pgfscope}%
\pgfpathrectangle{\pgfqpoint{0.100000in}{0.100000in}}{\pgfqpoint{3.007045in}{1.925000in}}%
\pgfusepath{clip}%
\pgfsetbuttcap%
\pgfsetmiterjoin%
\definecolor{currentfill}{rgb}{0.093272,0.396878,0.673664}%
\pgfsetfillcolor{currentfill}%
\pgfsetlinewidth{0.000000pt}%
\definecolor{currentstroke}{rgb}{0.000000,0.000000,0.000000}%
\pgfsetstrokecolor{currentstroke}%
\pgfsetstrokeopacity{0.000000}%
\pgfsetdash{}{0pt}%
\pgfpathmoveto{\pgfqpoint{0.784268in}{0.583082in}}%
\pgfpathlineto{\pgfqpoint{0.782998in}{0.580535in}}%
\pgfpathlineto{\pgfqpoint{0.776856in}{0.574080in}}%
\pgfpathlineto{\pgfqpoint{0.777383in}{0.573526in}}%
\pgfpathlineto{\pgfqpoint{0.771216in}{0.567098in}}%
\pgfpathlineto{\pgfqpoint{0.773326in}{0.565108in}}%
\pgfpathlineto{\pgfqpoint{0.770248in}{0.561865in}}%
\pgfpathlineto{\pgfqpoint{0.771871in}{0.560318in}}%
\pgfpathlineto{\pgfqpoint{0.768863in}{0.557118in}}%
\pgfpathlineto{\pgfqpoint{0.775661in}{0.550783in}}%
\pgfpathlineto{\pgfqpoint{0.790016in}{0.537796in}}%
\pgfpathlineto{\pgfqpoint{0.799988in}{0.528978in}}%
\pgfpathlineto{\pgfqpoint{0.801466in}{0.530632in}}%
\pgfpathlineto{\pgfqpoint{0.803129in}{0.529197in}}%
\pgfpathlineto{\pgfqpoint{0.804798in}{0.527778in}}%
\pgfpathlineto{\pgfqpoint{0.800473in}{0.522699in}}%
\pgfpathlineto{\pgfqpoint{0.798922in}{0.524065in}}%
\pgfpathlineto{\pgfqpoint{0.793312in}{0.517519in}}%
\pgfpathlineto{\pgfqpoint{0.791670in}{0.518914in}}%
\pgfpathlineto{\pgfqpoint{0.788832in}{0.515575in}}%
\pgfpathlineto{\pgfqpoint{0.787134in}{0.516986in}}%
\pgfpathlineto{\pgfqpoint{0.775774in}{0.503541in}}%
\pgfpathlineto{\pgfqpoint{0.777669in}{0.501945in}}%
\pgfpathlineto{\pgfqpoint{0.773225in}{0.497185in}}%
\pgfpathlineto{\pgfqpoint{0.771756in}{0.498478in}}%
\pgfpathlineto{\pgfqpoint{0.770286in}{0.496757in}}%
\pgfpathlineto{\pgfqpoint{0.768805in}{0.498004in}}%
\pgfpathlineto{\pgfqpoint{0.767409in}{0.496302in}}%
\pgfpathlineto{\pgfqpoint{0.764086in}{0.499185in}}%
\pgfpathlineto{\pgfqpoint{0.762810in}{0.497757in}}%
\pgfpathlineto{\pgfqpoint{0.761153in}{0.499198in}}%
\pgfpathlineto{\pgfqpoint{0.756937in}{0.494067in}}%
\pgfpathlineto{\pgfqpoint{0.755431in}{0.495345in}}%
\pgfpathlineto{\pgfqpoint{0.751165in}{0.490590in}}%
\pgfpathlineto{\pgfqpoint{0.741587in}{0.499176in}}%
\pgfpathlineto{\pgfqpoint{0.742883in}{0.500799in}}%
\pgfpathlineto{\pgfqpoint{0.741434in}{0.502134in}}%
\pgfpathlineto{\pgfqpoint{0.742884in}{0.503844in}}%
\pgfpathlineto{\pgfqpoint{0.741312in}{0.505279in}}%
\pgfpathlineto{\pgfqpoint{0.742738in}{0.506896in}}%
\pgfpathlineto{\pgfqpoint{0.741050in}{0.508543in}}%
\pgfpathlineto{\pgfqpoint{0.746586in}{0.510216in}}%
\pgfpathlineto{\pgfqpoint{0.746712in}{0.512326in}}%
\pgfpathlineto{\pgfqpoint{0.745027in}{0.513729in}}%
\pgfpathlineto{\pgfqpoint{0.746433in}{0.515178in}}%
\pgfpathlineto{\pgfqpoint{0.748958in}{0.513680in}}%
\pgfpathlineto{\pgfqpoint{0.748481in}{0.510045in}}%
\pgfpathlineto{\pgfqpoint{0.750210in}{0.505718in}}%
\pgfpathlineto{\pgfqpoint{0.755351in}{0.502374in}}%
\pgfpathlineto{\pgfqpoint{0.757588in}{0.502208in}}%
\pgfpathlineto{\pgfqpoint{0.761221in}{0.502993in}}%
\pgfpathlineto{\pgfqpoint{0.764189in}{0.504114in}}%
\pgfpathlineto{\pgfqpoint{0.765417in}{0.505253in}}%
\pgfpathlineto{\pgfqpoint{0.767984in}{0.511139in}}%
\pgfpathlineto{\pgfqpoint{0.770424in}{0.514447in}}%
\pgfpathlineto{\pgfqpoint{0.771090in}{0.519688in}}%
\pgfpathlineto{\pgfqpoint{0.769265in}{0.521681in}}%
\pgfpathlineto{\pgfqpoint{0.771550in}{0.523004in}}%
\pgfpathlineto{\pgfqpoint{0.773952in}{0.519957in}}%
\pgfpathlineto{\pgfqpoint{0.778165in}{0.520246in}}%
\pgfpathlineto{\pgfqpoint{0.780075in}{0.522157in}}%
\pgfpathlineto{\pgfqpoint{0.781283in}{0.524482in}}%
\pgfpathlineto{\pgfqpoint{0.780153in}{0.528966in}}%
\pgfpathlineto{\pgfqpoint{0.776623in}{0.527750in}}%
\pgfpathlineto{\pgfqpoint{0.774378in}{0.527915in}}%
\pgfpathlineto{\pgfqpoint{0.773486in}{0.530158in}}%
\pgfpathlineto{\pgfqpoint{0.770430in}{0.530368in}}%
\pgfpathlineto{\pgfqpoint{0.767348in}{0.531211in}}%
\pgfpathlineto{\pgfqpoint{0.764853in}{0.531128in}}%
\pgfpathlineto{\pgfqpoint{0.757318in}{0.529341in}}%
\pgfpathlineto{\pgfqpoint{0.759189in}{0.532678in}}%
\pgfpathlineto{\pgfqpoint{0.759643in}{0.537554in}}%
\pgfpathlineto{\pgfqpoint{0.757195in}{0.537867in}}%
\pgfpathlineto{\pgfqpoint{0.757549in}{0.535521in}}%
\pgfpathlineto{\pgfqpoint{0.755714in}{0.534840in}}%
\pgfpathlineto{\pgfqpoint{0.755754in}{0.537318in}}%
\pgfpathlineto{\pgfqpoint{0.754836in}{0.539913in}}%
\pgfpathlineto{\pgfqpoint{0.749644in}{0.544946in}}%
\pgfpathlineto{\pgfqpoint{0.743760in}{0.546795in}}%
\pgfpathlineto{\pgfqpoint{0.742220in}{0.547611in}}%
\pgfpathlineto{\pgfqpoint{0.735804in}{0.558236in}}%
\pgfpathlineto{\pgfqpoint{0.734630in}{0.560668in}}%
\pgfpathlineto{\pgfqpoint{0.734596in}{0.562774in}}%
\pgfpathlineto{\pgfqpoint{0.736774in}{0.565731in}}%
\pgfpathlineto{\pgfqpoint{0.738729in}{0.567328in}}%
\pgfpathlineto{\pgfqpoint{0.739066in}{0.574040in}}%
\pgfpathlineto{\pgfqpoint{0.742356in}{0.572549in}}%
\pgfpathlineto{\pgfqpoint{0.746458in}{0.574359in}}%
\pgfpathlineto{\pgfqpoint{0.742466in}{0.581233in}}%
\pgfpathlineto{\pgfqpoint{0.740777in}{0.583420in}}%
\pgfpathlineto{\pgfqpoint{0.739815in}{0.588701in}}%
\pgfpathlineto{\pgfqpoint{0.738785in}{0.592111in}}%
\pgfpathlineto{\pgfqpoint{0.740864in}{0.594187in}}%
\pgfpathlineto{\pgfqpoint{0.750011in}{0.591581in}}%
\pgfpathlineto{\pgfqpoint{0.750589in}{0.589917in}}%
\pgfpathlineto{\pgfqpoint{0.753942in}{0.589343in}}%
\pgfpathlineto{\pgfqpoint{0.753820in}{0.591119in}}%
\pgfpathlineto{\pgfqpoint{0.762748in}{0.589868in}}%
\pgfpathlineto{\pgfqpoint{0.764937in}{0.585166in}}%
\pgfpathlineto{\pgfqpoint{0.766490in}{0.583615in}}%
\pgfpathlineto{\pgfqpoint{0.768297in}{0.583840in}}%
\pgfpathlineto{\pgfqpoint{0.767184in}{0.586211in}}%
\pgfpathlineto{\pgfqpoint{0.770001in}{0.588417in}}%
\pgfpathlineto{\pgfqpoint{0.777196in}{0.586185in}}%
\pgfpathlineto{\pgfqpoint{0.781761in}{0.584481in}}%
\pgfpathlineto{\pgfqpoint{0.784268in}{0.583082in}}%
\pgfpathclose%
\pgfusepath{fill}%
\end{pgfscope}%
\begin{pgfscope}%
\pgfpathrectangle{\pgfqpoint{0.100000in}{0.100000in}}{\pgfqpoint{3.007045in}{1.925000in}}%
\pgfusepath{clip}%
\pgfsetbuttcap%
\pgfsetmiterjoin%
\definecolor{currentfill}{rgb}{0.031373,0.285675,0.564291}%
\pgfsetfillcolor{currentfill}%
\pgfsetlinewidth{0.000000pt}%
\definecolor{currentstroke}{rgb}{0.000000,0.000000,0.000000}%
\pgfsetstrokecolor{currentstroke}%
\pgfsetstrokeopacity{0.000000}%
\pgfsetdash{}{0pt}%
\pgfpathmoveto{\pgfqpoint{1.521452in}{1.403789in}}%
\pgfpathlineto{\pgfqpoint{1.471839in}{1.406868in}}%
\pgfpathlineto{\pgfqpoint{1.473591in}{1.432656in}}%
\pgfpathlineto{\pgfqpoint{1.475452in}{1.459167in}}%
\pgfpathlineto{\pgfqpoint{1.483456in}{1.461953in}}%
\pgfpathlineto{\pgfqpoint{1.490267in}{1.460196in}}%
\pgfpathlineto{\pgfqpoint{1.496174in}{1.453745in}}%
\pgfpathlineto{\pgfqpoint{1.506479in}{1.454584in}}%
\pgfpathlineto{\pgfqpoint{1.511045in}{1.452183in}}%
\pgfpathlineto{\pgfqpoint{1.517464in}{1.451554in}}%
\pgfpathlineto{\pgfqpoint{1.522608in}{1.451050in}}%
\pgfpathlineto{\pgfqpoint{1.522496in}{1.435436in}}%
\pgfpathlineto{\pgfqpoint{1.521204in}{1.412401in}}%
\pgfpathlineto{\pgfqpoint{1.521452in}{1.403789in}}%
\pgfpathclose%
\pgfusepath{fill}%
\end{pgfscope}%
\begin{pgfscope}%
\pgfpathrectangle{\pgfqpoint{0.100000in}{0.100000in}}{\pgfqpoint{3.007045in}{1.925000in}}%
\pgfusepath{clip}%
\pgfsetbuttcap%
\pgfsetmiterjoin%
\definecolor{currentfill}{rgb}{0.485490,0.718524,0.853426}%
\pgfsetfillcolor{currentfill}%
\pgfsetlinewidth{0.000000pt}%
\definecolor{currentstroke}{rgb}{0.000000,0.000000,0.000000}%
\pgfsetstrokecolor{currentstroke}%
\pgfsetstrokeopacity{0.000000}%
\pgfsetdash{}{0pt}%
\pgfpathmoveto{\pgfqpoint{0.763180in}{1.655715in}}%
\pgfpathlineto{\pgfqpoint{0.759381in}{1.652797in}}%
\pgfpathlineto{\pgfqpoint{0.758646in}{1.647388in}}%
\pgfpathlineto{\pgfqpoint{0.754151in}{1.641671in}}%
\pgfpathlineto{\pgfqpoint{0.746366in}{1.637288in}}%
\pgfpathlineto{\pgfqpoint{0.738581in}{1.626973in}}%
\pgfpathlineto{\pgfqpoint{0.733470in}{1.622014in}}%
\pgfpathlineto{\pgfqpoint{0.730914in}{1.610662in}}%
\pgfpathlineto{\pgfqpoint{0.718698in}{1.613732in}}%
\pgfpathlineto{\pgfqpoint{0.720108in}{1.619262in}}%
\pgfpathlineto{\pgfqpoint{0.716216in}{1.624121in}}%
\pgfpathlineto{\pgfqpoint{0.698781in}{1.628508in}}%
\pgfpathlineto{\pgfqpoint{0.695577in}{1.627414in}}%
\pgfpathlineto{\pgfqpoint{0.687625in}{1.619603in}}%
\pgfpathlineto{\pgfqpoint{0.683738in}{1.619619in}}%
\pgfpathlineto{\pgfqpoint{0.671575in}{1.622748in}}%
\pgfpathlineto{\pgfqpoint{0.677402in}{1.629304in}}%
\pgfpathlineto{\pgfqpoint{0.678278in}{1.634204in}}%
\pgfpathlineto{\pgfqpoint{0.685898in}{1.641698in}}%
\pgfpathlineto{\pgfqpoint{0.684947in}{1.644999in}}%
\pgfpathlineto{\pgfqpoint{0.678193in}{1.649100in}}%
\pgfpathlineto{\pgfqpoint{0.678258in}{1.651609in}}%
\pgfpathlineto{\pgfqpoint{0.689599in}{1.651922in}}%
\pgfpathlineto{\pgfqpoint{0.689503in}{1.657799in}}%
\pgfpathlineto{\pgfqpoint{0.694447in}{1.659187in}}%
\pgfpathlineto{\pgfqpoint{0.695518in}{1.663103in}}%
\pgfpathlineto{\pgfqpoint{0.690132in}{1.668259in}}%
\pgfpathlineto{\pgfqpoint{0.683125in}{1.670416in}}%
\pgfpathlineto{\pgfqpoint{0.685286in}{1.683617in}}%
\pgfpathlineto{\pgfqpoint{0.680362in}{1.684919in}}%
\pgfpathlineto{\pgfqpoint{0.681006in}{1.694917in}}%
\pgfpathlineto{\pgfqpoint{0.693304in}{1.692341in}}%
\pgfpathlineto{\pgfqpoint{0.694542in}{1.697073in}}%
\pgfpathlineto{\pgfqpoint{0.704883in}{1.694237in}}%
\pgfpathlineto{\pgfqpoint{0.709253in}{1.696086in}}%
\pgfpathlineto{\pgfqpoint{0.712881in}{1.710060in}}%
\pgfpathlineto{\pgfqpoint{0.715135in}{1.709465in}}%
\pgfpathlineto{\pgfqpoint{0.718256in}{1.717477in}}%
\pgfpathlineto{\pgfqpoint{0.722204in}{1.719436in}}%
\pgfpathlineto{\pgfqpoint{0.732306in}{1.716868in}}%
\pgfpathlineto{\pgfqpoint{0.727679in}{1.706314in}}%
\pgfpathlineto{\pgfqpoint{0.729208in}{1.700974in}}%
\pgfpathlineto{\pgfqpoint{0.727519in}{1.694386in}}%
\pgfpathlineto{\pgfqpoint{0.728377in}{1.682025in}}%
\pgfpathlineto{\pgfqpoint{0.732751in}{1.676168in}}%
\pgfpathlineto{\pgfqpoint{0.733120in}{1.669073in}}%
\pgfpathlineto{\pgfqpoint{0.742633in}{1.666720in}}%
\pgfpathlineto{\pgfqpoint{0.741293in}{1.661247in}}%
\pgfpathlineto{\pgfqpoint{0.763180in}{1.655715in}}%
\pgfpathclose%
\pgfusepath{fill}%
\end{pgfscope}%
\begin{pgfscope}%
\pgfpathrectangle{\pgfqpoint{0.100000in}{0.100000in}}{\pgfqpoint{3.007045in}{1.925000in}}%
\pgfusepath{clip}%
\pgfsetbuttcap%
\pgfsetmiterjoin%
\definecolor{currentfill}{rgb}{0.351511,0.635848,0.812641}%
\pgfsetfillcolor{currentfill}%
\pgfsetlinewidth{0.000000pt}%
\definecolor{currentstroke}{rgb}{0.000000,0.000000,0.000000}%
\pgfsetstrokecolor{currentstroke}%
\pgfsetstrokeopacity{0.000000}%
\pgfsetdash{}{0pt}%
\pgfpathmoveto{\pgfqpoint{1.509458in}{1.149192in}}%
\pgfpathlineto{\pgfqpoint{1.507527in}{1.120611in}}%
\pgfpathlineto{\pgfqpoint{1.480135in}{1.122427in}}%
\pgfpathlineto{\pgfqpoint{1.439278in}{1.125169in}}%
\pgfpathlineto{\pgfqpoint{1.441646in}{1.153706in}}%
\pgfpathlineto{\pgfqpoint{1.446054in}{1.153418in}}%
\pgfpathlineto{\pgfqpoint{1.448195in}{1.181868in}}%
\pgfpathlineto{\pgfqpoint{1.509436in}{1.177873in}}%
\pgfpathlineto{\pgfqpoint{1.510385in}{1.177845in}}%
\pgfpathlineto{\pgfqpoint{1.508618in}{1.149222in}}%
\pgfpathlineto{\pgfqpoint{1.509458in}{1.149192in}}%
\pgfpathclose%
\pgfusepath{fill}%
\end{pgfscope}%
\begin{pgfscope}%
\pgfpathrectangle{\pgfqpoint{0.100000in}{0.100000in}}{\pgfqpoint{3.007045in}{1.925000in}}%
\pgfusepath{clip}%
\pgfsetbuttcap%
\pgfsetmiterjoin%
\definecolor{currentfill}{rgb}{0.435294,0.690965,0.842599}%
\pgfsetfillcolor{currentfill}%
\pgfsetlinewidth{0.000000pt}%
\definecolor{currentstroke}{rgb}{0.000000,0.000000,0.000000}%
\pgfsetstrokecolor{currentstroke}%
\pgfsetstrokeopacity{0.000000}%
\pgfsetdash{}{0pt}%
\pgfpathmoveto{\pgfqpoint{2.070604in}{1.165647in}}%
\pgfpathlineto{\pgfqpoint{2.050719in}{1.164389in}}%
\pgfpathlineto{\pgfqpoint{2.049876in}{1.175940in}}%
\pgfpathlineto{\pgfqpoint{2.029984in}{1.174625in}}%
\pgfpathlineto{\pgfqpoint{2.025601in}{1.191743in}}%
\pgfpathlineto{\pgfqpoint{2.024381in}{1.212989in}}%
\pgfpathlineto{\pgfqpoint{2.025473in}{1.215274in}}%
\pgfpathlineto{\pgfqpoint{2.032407in}{1.214497in}}%
\pgfpathlineto{\pgfqpoint{2.039739in}{1.216849in}}%
\pgfpathlineto{\pgfqpoint{2.044100in}{1.215046in}}%
\pgfpathlineto{\pgfqpoint{2.043262in}{1.228036in}}%
\pgfpathlineto{\pgfqpoint{2.060239in}{1.229418in}}%
\pgfpathlineto{\pgfqpoint{2.066190in}{1.226950in}}%
\pgfpathlineto{\pgfqpoint{2.067386in}{1.211625in}}%
\pgfpathlineto{\pgfqpoint{2.068009in}{1.203031in}}%
\pgfpathlineto{\pgfqpoint{2.064257in}{1.202733in}}%
\pgfpathlineto{\pgfqpoint{2.064705in}{1.195908in}}%
\pgfpathlineto{\pgfqpoint{2.068581in}{1.195334in}}%
\pgfpathlineto{\pgfqpoint{2.069337in}{1.185778in}}%
\pgfpathlineto{\pgfqpoint{2.075079in}{1.186079in}}%
\pgfpathlineto{\pgfqpoint{2.076427in}{1.165836in}}%
\pgfpathlineto{\pgfqpoint{2.070604in}{1.165647in}}%
\pgfpathclose%
\pgfusepath{fill}%
\end{pgfscope}%
\begin{pgfscope}%
\pgfpathrectangle{\pgfqpoint{0.100000in}{0.100000in}}{\pgfqpoint{3.007045in}{1.925000in}}%
\pgfusepath{clip}%
\pgfsetbuttcap%
\pgfsetmiterjoin%
\definecolor{currentfill}{rgb}{0.491765,0.721968,0.854779}%
\pgfsetfillcolor{currentfill}%
\pgfsetlinewidth{0.000000pt}%
\definecolor{currentstroke}{rgb}{0.000000,0.000000,0.000000}%
\pgfsetstrokecolor{currentstroke}%
\pgfsetstrokeopacity{0.000000}%
\pgfsetdash{}{0pt}%
\pgfpathmoveto{\pgfqpoint{2.803517in}{1.405975in}}%
\pgfpathlineto{\pgfqpoint{2.794033in}{1.391447in}}%
\pgfpathlineto{\pgfqpoint{2.768627in}{1.400117in}}%
\pgfpathlineto{\pgfqpoint{2.765754in}{1.403877in}}%
\pgfpathlineto{\pgfqpoint{2.760698in}{1.403338in}}%
\pgfpathlineto{\pgfqpoint{2.752320in}{1.406665in}}%
\pgfpathlineto{\pgfqpoint{2.747986in}{1.413389in}}%
\pgfpathlineto{\pgfqpoint{2.745672in}{1.421591in}}%
\pgfpathlineto{\pgfqpoint{2.739778in}{1.426983in}}%
\pgfpathlineto{\pgfqpoint{2.768356in}{1.454967in}}%
\pgfpathlineto{\pgfqpoint{2.776040in}{1.453025in}}%
\pgfpathlineto{\pgfqpoint{2.787399in}{1.454470in}}%
\pgfpathlineto{\pgfqpoint{2.787830in}{1.459571in}}%
\pgfpathlineto{\pgfqpoint{2.794361in}{1.458299in}}%
\pgfpathlineto{\pgfqpoint{2.795846in}{1.455210in}}%
\pgfpathlineto{\pgfqpoint{2.806124in}{1.452725in}}%
\pgfpathlineto{\pgfqpoint{2.815310in}{1.453061in}}%
\pgfpathlineto{\pgfqpoint{2.815641in}{1.458005in}}%
\pgfpathlineto{\pgfqpoint{2.816097in}{1.458112in}}%
\pgfpathlineto{\pgfqpoint{2.822125in}{1.424128in}}%
\pgfpathlineto{\pgfqpoint{2.803333in}{1.417011in}}%
\pgfpathlineto{\pgfqpoint{2.803517in}{1.405975in}}%
\pgfpathclose%
\pgfusepath{fill}%
\end{pgfscope}%
\begin{pgfscope}%
\pgfpathrectangle{\pgfqpoint{0.100000in}{0.100000in}}{\pgfqpoint{3.007045in}{1.925000in}}%
\pgfusepath{clip}%
\pgfsetbuttcap%
\pgfsetmiterjoin%
\definecolor{currentfill}{rgb}{0.485490,0.718524,0.853426}%
\pgfsetfillcolor{currentfill}%
\pgfsetlinewidth{0.000000pt}%
\definecolor{currentstroke}{rgb}{0.000000,0.000000,0.000000}%
\pgfsetstrokecolor{currentstroke}%
\pgfsetstrokeopacity{0.000000}%
\pgfsetdash{}{0pt}%
\pgfpathmoveto{\pgfqpoint{1.219256in}{0.700832in}}%
\pgfpathlineto{\pgfqpoint{1.215312in}{0.701261in}}%
\pgfpathlineto{\pgfqpoint{1.186737in}{0.704479in}}%
\pgfpathlineto{\pgfqpoint{1.133803in}{0.710900in}}%
\pgfpathlineto{\pgfqpoint{1.141737in}{0.773832in}}%
\pgfpathlineto{\pgfqpoint{1.144512in}{0.779815in}}%
\pgfpathlineto{\pgfqpoint{1.147002in}{0.802056in}}%
\pgfpathlineto{\pgfqpoint{1.146247in}{0.808107in}}%
\pgfpathlineto{\pgfqpoint{1.147690in}{0.819949in}}%
\pgfpathlineto{\pgfqpoint{1.165157in}{0.817109in}}%
\pgfpathlineto{\pgfqpoint{1.166565in}{0.828636in}}%
\pgfpathlineto{\pgfqpoint{1.173550in}{0.827766in}}%
\pgfpathlineto{\pgfqpoint{1.177035in}{0.856282in}}%
\pgfpathlineto{\pgfqpoint{1.210171in}{0.852394in}}%
\pgfpathlineto{\pgfqpoint{1.210832in}{0.858083in}}%
\pgfpathlineto{\pgfqpoint{1.233670in}{0.855551in}}%
\pgfpathlineto{\pgfqpoint{1.231742in}{0.838561in}}%
\pgfpathlineto{\pgfqpoint{1.226820in}{0.793144in}}%
\pgfpathlineto{\pgfqpoint{1.224219in}{0.776176in}}%
\pgfpathlineto{\pgfqpoint{1.201518in}{0.778286in}}%
\pgfpathlineto{\pgfqpoint{1.200282in}{0.767082in}}%
\pgfpathlineto{\pgfqpoint{1.198272in}{0.767309in}}%
\pgfpathlineto{\pgfqpoint{1.194858in}{0.738140in}}%
\pgfpathlineto{\pgfqpoint{1.222778in}{0.735139in}}%
\pgfpathlineto{\pgfqpoint{1.219256in}{0.700832in}}%
\pgfpathclose%
\pgfusepath{fill}%
\end{pgfscope}%
\begin{pgfscope}%
\pgfpathrectangle{\pgfqpoint{0.100000in}{0.100000in}}{\pgfqpoint{3.007045in}{1.925000in}}%
\pgfusepath{clip}%
\pgfsetbuttcap%
\pgfsetmiterjoin%
\definecolor{currentfill}{rgb}{0.422745,0.684075,0.839892}%
\pgfsetfillcolor{currentfill}%
\pgfsetlinewidth{0.000000pt}%
\definecolor{currentstroke}{rgb}{0.000000,0.000000,0.000000}%
\pgfsetstrokecolor{currentstroke}%
\pgfsetstrokeopacity{0.000000}%
\pgfsetdash{}{0pt}%
\pgfpathmoveto{\pgfqpoint{2.394856in}{0.702917in}}%
\pgfpathlineto{\pgfqpoint{2.356903in}{0.698990in}}%
\pgfpathlineto{\pgfqpoint{2.339099in}{0.697275in}}%
\pgfpathlineto{\pgfqpoint{2.331743in}{0.699866in}}%
\pgfpathlineto{\pgfqpoint{2.330327in}{0.705545in}}%
\pgfpathlineto{\pgfqpoint{2.329526in}{0.710196in}}%
\pgfpathlineto{\pgfqpoint{2.336528in}{0.717430in}}%
\pgfpathlineto{\pgfqpoint{2.351237in}{0.719223in}}%
\pgfpathlineto{\pgfqpoint{2.349352in}{0.736500in}}%
\pgfpathlineto{\pgfqpoint{2.350162in}{0.739256in}}%
\pgfpathlineto{\pgfqpoint{2.360665in}{0.742287in}}%
\pgfpathlineto{\pgfqpoint{2.363559in}{0.734763in}}%
\pgfpathlineto{\pgfqpoint{2.368826in}{0.730586in}}%
\pgfpathlineto{\pgfqpoint{2.375605in}{0.733287in}}%
\pgfpathlineto{\pgfqpoint{2.378977in}{0.740714in}}%
\pgfpathlineto{\pgfqpoint{2.382677in}{0.742228in}}%
\pgfpathlineto{\pgfqpoint{2.392692in}{0.741268in}}%
\pgfpathlineto{\pgfqpoint{2.395370in}{0.738234in}}%
\pgfpathlineto{\pgfqpoint{2.395922in}{0.728401in}}%
\pgfpathlineto{\pgfqpoint{2.389930in}{0.724977in}}%
\pgfpathlineto{\pgfqpoint{2.386745in}{0.718055in}}%
\pgfpathlineto{\pgfqpoint{2.389600in}{0.714570in}}%
\pgfpathlineto{\pgfqpoint{2.394856in}{0.702917in}}%
\pgfpathclose%
\pgfusepath{fill}%
\end{pgfscope}%
\begin{pgfscope}%
\pgfpathrectangle{\pgfqpoint{0.100000in}{0.100000in}}{\pgfqpoint{3.007045in}{1.925000in}}%
\pgfusepath{clip}%
\pgfsetbuttcap%
\pgfsetmiterjoin%
\definecolor{currentfill}{rgb}{0.784591,0.864237,0.939962}%
\pgfsetfillcolor{currentfill}%
\pgfsetlinewidth{0.000000pt}%
\definecolor{currentstroke}{rgb}{0.000000,0.000000,0.000000}%
\pgfsetstrokecolor{currentstroke}%
\pgfsetstrokeopacity{0.000000}%
\pgfsetdash{}{0pt}%
\pgfpathmoveto{\pgfqpoint{2.313368in}{1.507664in}}%
\pgfpathlineto{\pgfqpoint{2.298344in}{1.505689in}}%
\pgfpathlineto{\pgfqpoint{2.291082in}{1.562710in}}%
\pgfpathlineto{\pgfqpoint{2.290081in}{1.574045in}}%
\pgfpathlineto{\pgfqpoint{2.273175in}{1.571785in}}%
\pgfpathlineto{\pgfqpoint{2.269589in}{1.599880in}}%
\pgfpathlineto{\pgfqpoint{2.271744in}{1.600141in}}%
\pgfpathlineto{\pgfqpoint{2.277048in}{1.592666in}}%
\pgfpathlineto{\pgfqpoint{2.285970in}{1.592606in}}%
\pgfpathlineto{\pgfqpoint{2.292882in}{1.588599in}}%
\pgfpathlineto{\pgfqpoint{2.308293in}{1.585034in}}%
\pgfpathlineto{\pgfqpoint{2.311509in}{1.579156in}}%
\pgfpathlineto{\pgfqpoint{2.316771in}{1.573334in}}%
\pgfpathlineto{\pgfqpoint{2.320179in}{1.565731in}}%
\pgfpathlineto{\pgfqpoint{2.311941in}{1.567027in}}%
\pgfpathlineto{\pgfqpoint{2.313141in}{1.558957in}}%
\pgfpathlineto{\pgfqpoint{2.319127in}{1.556401in}}%
\pgfpathlineto{\pgfqpoint{2.322517in}{1.545363in}}%
\pgfpathlineto{\pgfqpoint{2.321374in}{1.538390in}}%
\pgfpathlineto{\pgfqpoint{2.322792in}{1.520480in}}%
\pgfpathlineto{\pgfqpoint{2.318985in}{1.515613in}}%
\pgfpathlineto{\pgfqpoint{2.314903in}{1.514937in}}%
\pgfpathlineto{\pgfqpoint{2.313368in}{1.507664in}}%
\pgfpathclose%
\pgfusepath{fill}%
\end{pgfscope}%
\begin{pgfscope}%
\pgfpathrectangle{\pgfqpoint{0.100000in}{0.100000in}}{\pgfqpoint{3.007045in}{1.925000in}}%
\pgfusepath{clip}%
\pgfsetbuttcap%
\pgfsetmiterjoin%
\definecolor{currentfill}{rgb}{0.326290,0.618624,0.802799}%
\pgfsetfillcolor{currentfill}%
\pgfsetlinewidth{0.000000pt}%
\definecolor{currentstroke}{rgb}{0.000000,0.000000,0.000000}%
\pgfsetstrokecolor{currentstroke}%
\pgfsetstrokeopacity{0.000000}%
\pgfsetdash{}{0pt}%
\pgfpathmoveto{\pgfqpoint{2.327732in}{1.268420in}}%
\pgfpathlineto{\pgfqpoint{2.322114in}{1.267499in}}%
\pgfpathlineto{\pgfqpoint{2.322524in}{1.264076in}}%
\pgfpathlineto{\pgfqpoint{2.305349in}{1.261913in}}%
\pgfpathlineto{\pgfqpoint{2.306160in}{1.255224in}}%
\pgfpathlineto{\pgfqpoint{2.301562in}{1.253044in}}%
\pgfpathlineto{\pgfqpoint{2.283146in}{1.250837in}}%
\pgfpathlineto{\pgfqpoint{2.278491in}{1.292509in}}%
\pgfpathlineto{\pgfqpoint{2.301283in}{1.295237in}}%
\pgfpathlineto{\pgfqpoint{2.299961in}{1.306677in}}%
\pgfpathlineto{\pgfqpoint{2.305531in}{1.307364in}}%
\pgfpathlineto{\pgfqpoint{2.322615in}{1.309553in}}%
\pgfpathlineto{\pgfqpoint{2.327732in}{1.268420in}}%
\pgfpathclose%
\pgfusepath{fill}%
\end{pgfscope}%
\begin{pgfscope}%
\pgfpathrectangle{\pgfqpoint{0.100000in}{0.100000in}}{\pgfqpoint{3.007045in}{1.925000in}}%
\pgfusepath{clip}%
\pgfsetbuttcap%
\pgfsetmiterjoin%
\definecolor{currentfill}{rgb}{0.585882,0.773641,0.875079}%
\pgfsetfillcolor{currentfill}%
\pgfsetlinewidth{0.000000pt}%
\definecolor{currentstroke}{rgb}{0.000000,0.000000,0.000000}%
\pgfsetstrokecolor{currentstroke}%
\pgfsetstrokeopacity{0.000000}%
\pgfsetdash{}{0pt}%
\pgfpathmoveto{\pgfqpoint{2.625680in}{1.382482in}}%
\pgfpathlineto{\pgfqpoint{2.626550in}{1.378374in}}%
\pgfpathlineto{\pgfqpoint{2.607627in}{1.374496in}}%
\pgfpathlineto{\pgfqpoint{2.604647in}{1.373975in}}%
\pgfpathlineto{\pgfqpoint{2.595492in}{1.381828in}}%
\pgfpathlineto{\pgfqpoint{2.581784in}{1.380075in}}%
\pgfpathlineto{\pgfqpoint{2.558934in}{1.375533in}}%
\pgfpathlineto{\pgfqpoint{2.556898in}{1.389007in}}%
\pgfpathlineto{\pgfqpoint{2.558369in}{1.389969in}}%
\pgfpathlineto{\pgfqpoint{2.556395in}{1.400149in}}%
\pgfpathlineto{\pgfqpoint{2.549510in}{1.399000in}}%
\pgfpathlineto{\pgfqpoint{2.543334in}{1.433990in}}%
\pgfpathlineto{\pgfqpoint{2.546681in}{1.434064in}}%
\pgfpathlineto{\pgfqpoint{2.552432in}{1.429022in}}%
\pgfpathlineto{\pgfqpoint{2.557784in}{1.431044in}}%
\pgfpathlineto{\pgfqpoint{2.570755in}{1.438725in}}%
\pgfpathlineto{\pgfqpoint{2.572079in}{1.437927in}}%
\pgfpathlineto{\pgfqpoint{2.601922in}{1.443430in}}%
\pgfpathlineto{\pgfqpoint{2.608126in}{1.441492in}}%
\pgfpathlineto{\pgfqpoint{2.612717in}{1.410639in}}%
\pgfpathlineto{\pgfqpoint{2.619453in}{1.411977in}}%
\pgfpathlineto{\pgfqpoint{2.625680in}{1.382482in}}%
\pgfpathclose%
\pgfusepath{fill}%
\end{pgfscope}%
\begin{pgfscope}%
\pgfpathrectangle{\pgfqpoint{0.100000in}{0.100000in}}{\pgfqpoint{3.007045in}{1.925000in}}%
\pgfusepath{clip}%
\pgfsetbuttcap%
\pgfsetmiterjoin%
\definecolor{currentfill}{rgb}{0.592157,0.777086,0.876432}%
\pgfsetfillcolor{currentfill}%
\pgfsetlinewidth{0.000000pt}%
\definecolor{currentstroke}{rgb}{0.000000,0.000000,0.000000}%
\pgfsetstrokecolor{currentstroke}%
\pgfsetstrokeopacity{0.000000}%
\pgfsetdash{}{0pt}%
\pgfpathmoveto{\pgfqpoint{2.252299in}{1.407809in}}%
\pgfpathlineto{\pgfqpoint{2.263712in}{1.409069in}}%
\pgfpathlineto{\pgfqpoint{2.261163in}{1.431916in}}%
\pgfpathlineto{\pgfqpoint{2.283625in}{1.434436in}}%
\pgfpathlineto{\pgfqpoint{2.286455in}{1.412044in}}%
\pgfpathlineto{\pgfqpoint{2.296340in}{1.413257in}}%
\pgfpathlineto{\pgfqpoint{2.299946in}{1.390294in}}%
\pgfpathlineto{\pgfqpoint{2.272022in}{1.386856in}}%
\pgfpathlineto{\pgfqpoint{2.254932in}{1.384910in}}%
\pgfpathlineto{\pgfqpoint{2.252299in}{1.407809in}}%
\pgfpathclose%
\pgfusepath{fill}%
\end{pgfscope}%
\begin{pgfscope}%
\pgfpathrectangle{\pgfqpoint{0.100000in}{0.100000in}}{\pgfqpoint{3.007045in}{1.925000in}}%
\pgfusepath{clip}%
\pgfsetbuttcap%
\pgfsetmiterjoin%
\definecolor{currentfill}{rgb}{0.290980,0.594510,0.789020}%
\pgfsetfillcolor{currentfill}%
\pgfsetlinewidth{0.000000pt}%
\definecolor{currentstroke}{rgb}{0.000000,0.000000,0.000000}%
\pgfsetstrokecolor{currentstroke}%
\pgfsetstrokeopacity{0.000000}%
\pgfsetdash{}{0pt}%
\pgfpathmoveto{\pgfqpoint{1.677065in}{1.742344in}}%
\pgfpathlineto{\pgfqpoint{1.678294in}{1.748130in}}%
\pgfpathlineto{\pgfqpoint{1.676638in}{1.755236in}}%
\pgfpathlineto{\pgfqpoint{1.677958in}{1.760345in}}%
\pgfpathlineto{\pgfqpoint{1.676648in}{1.765752in}}%
\pgfpathlineto{\pgfqpoint{1.710937in}{1.765107in}}%
\pgfpathlineto{\pgfqpoint{1.745680in}{1.764366in}}%
\pgfpathlineto{\pgfqpoint{1.745890in}{1.729762in}}%
\pgfpathlineto{\pgfqpoint{1.746314in}{1.723984in}}%
\pgfpathlineto{\pgfqpoint{1.740624in}{1.724147in}}%
\pgfpathlineto{\pgfqpoint{1.740597in}{1.718328in}}%
\pgfpathlineto{\pgfqpoint{1.734882in}{1.718354in}}%
\pgfpathlineto{\pgfqpoint{1.734818in}{1.712519in}}%
\pgfpathlineto{\pgfqpoint{1.711752in}{1.712779in}}%
\pgfpathlineto{\pgfqpoint{1.711793in}{1.718564in}}%
\pgfpathlineto{\pgfqpoint{1.706015in}{1.718669in}}%
\pgfpathlineto{\pgfqpoint{1.705532in}{1.740473in}}%
\pgfpathlineto{\pgfqpoint{1.676789in}{1.740990in}}%
\pgfpathlineto{\pgfqpoint{1.677065in}{1.742344in}}%
\pgfpathclose%
\pgfusepath{fill}%
\end{pgfscope}%
\begin{pgfscope}%
\pgfpathrectangle{\pgfqpoint{0.100000in}{0.100000in}}{\pgfqpoint{3.007045in}{1.925000in}}%
\pgfusepath{clip}%
\pgfsetbuttcap%
\pgfsetmiterjoin%
\definecolor{currentfill}{rgb}{0.460392,0.704744,0.848012}%
\pgfsetfillcolor{currentfill}%
\pgfsetlinewidth{0.000000pt}%
\definecolor{currentstroke}{rgb}{0.000000,0.000000,0.000000}%
\pgfsetstrokecolor{currentstroke}%
\pgfsetstrokeopacity{0.000000}%
\pgfsetdash{}{0pt}%
\pgfpathmoveto{\pgfqpoint{2.667563in}{1.268330in}}%
\pgfpathlineto{\pgfqpoint{2.655502in}{1.265959in}}%
\pgfpathlineto{\pgfqpoint{2.652060in}{1.280358in}}%
\pgfpathlineto{\pgfqpoint{2.638827in}{1.299653in}}%
\pgfpathlineto{\pgfqpoint{2.636603in}{1.298784in}}%
\pgfpathlineto{\pgfqpoint{2.632407in}{1.305014in}}%
\pgfpathlineto{\pgfqpoint{2.627758in}{1.303889in}}%
\pgfpathlineto{\pgfqpoint{2.625195in}{1.307053in}}%
\pgfpathlineto{\pgfqpoint{2.628110in}{1.312979in}}%
\pgfpathlineto{\pgfqpoint{2.626236in}{1.316064in}}%
\pgfpathlineto{\pgfqpoint{2.633083in}{1.330389in}}%
\pgfpathlineto{\pgfqpoint{2.645950in}{1.339764in}}%
\pgfpathlineto{\pgfqpoint{2.648248in}{1.330436in}}%
\pgfpathlineto{\pgfqpoint{2.664036in}{1.332084in}}%
\pgfpathlineto{\pgfqpoint{2.670437in}{1.328641in}}%
\pgfpathlineto{\pgfqpoint{2.680799in}{1.334029in}}%
\pgfpathlineto{\pgfqpoint{2.690299in}{1.329057in}}%
\pgfpathlineto{\pgfqpoint{2.684366in}{1.322430in}}%
\pgfpathlineto{\pgfqpoint{2.693385in}{1.305552in}}%
\pgfpathlineto{\pgfqpoint{2.687398in}{1.300784in}}%
\pgfpathlineto{\pgfqpoint{2.690460in}{1.296019in}}%
\pgfpathlineto{\pgfqpoint{2.697015in}{1.296543in}}%
\pgfpathlineto{\pgfqpoint{2.701481in}{1.290448in}}%
\pgfpathlineto{\pgfqpoint{2.706007in}{1.289043in}}%
\pgfpathlineto{\pgfqpoint{2.716037in}{1.278304in}}%
\pgfpathlineto{\pgfqpoint{2.667563in}{1.268330in}}%
\pgfpathclose%
\pgfusepath{fill}%
\end{pgfscope}%
\begin{pgfscope}%
\pgfpathrectangle{\pgfqpoint{0.100000in}{0.100000in}}{\pgfqpoint{3.007045in}{1.925000in}}%
\pgfusepath{clip}%
\pgfsetbuttcap%
\pgfsetmiterjoin%
\definecolor{currentfill}{rgb}{0.381776,0.656517,0.824452}%
\pgfsetfillcolor{currentfill}%
\pgfsetlinewidth{0.000000pt}%
\definecolor{currentstroke}{rgb}{0.000000,0.000000,0.000000}%
\pgfsetstrokecolor{currentstroke}%
\pgfsetstrokeopacity{0.000000}%
\pgfsetdash{}{0pt}%
\pgfpathmoveto{\pgfqpoint{1.395381in}{1.309582in}}%
\pgfpathlineto{\pgfqpoint{1.396484in}{1.306583in}}%
\pgfpathlineto{\pgfqpoint{1.395479in}{1.295197in}}%
\pgfpathlineto{\pgfqpoint{1.394532in}{1.283801in}}%
\pgfpathlineto{\pgfqpoint{1.392134in}{1.280968in}}%
\pgfpathlineto{\pgfqpoint{1.346503in}{1.285013in}}%
\pgfpathlineto{\pgfqpoint{1.294492in}{1.290219in}}%
\pgfpathlineto{\pgfqpoint{1.262254in}{1.293547in}}%
\pgfpathlineto{\pgfqpoint{1.267098in}{1.336682in}}%
\pgfpathlineto{\pgfqpoint{1.297671in}{1.333074in}}%
\pgfpathlineto{\pgfqpoint{1.297072in}{1.327304in}}%
\pgfpathlineto{\pgfqpoint{1.326529in}{1.324213in}}%
\pgfpathlineto{\pgfqpoint{1.325420in}{1.312996in}}%
\pgfpathlineto{\pgfqpoint{1.358996in}{1.309848in}}%
\pgfpathlineto{\pgfqpoint{1.359322in}{1.312674in}}%
\pgfpathlineto{\pgfqpoint{1.395381in}{1.309582in}}%
\pgfpathclose%
\pgfusepath{fill}%
\end{pgfscope}%
\begin{pgfscope}%
\pgfpathrectangle{\pgfqpoint{0.100000in}{0.100000in}}{\pgfqpoint{3.007045in}{1.925000in}}%
\pgfusepath{clip}%
\pgfsetbuttcap%
\pgfsetmiterjoin%
\definecolor{currentfill}{rgb}{0.849550,0.907543,0.961615}%
\pgfsetfillcolor{currentfill}%
\pgfsetlinewidth{0.000000pt}%
\definecolor{currentstroke}{rgb}{0.000000,0.000000,0.000000}%
\pgfsetstrokecolor{currentstroke}%
\pgfsetstrokeopacity{0.000000}%
\pgfsetdash{}{0pt}%
\pgfpathmoveto{\pgfqpoint{1.201403in}{1.171929in}}%
\pgfpathlineto{\pgfqpoint{1.200990in}{1.169461in}}%
\pgfpathlineto{\pgfqpoint{1.208767in}{1.163511in}}%
\pgfpathlineto{\pgfqpoint{1.213448in}{1.158253in}}%
\pgfpathlineto{\pgfqpoint{1.213877in}{1.153875in}}%
\pgfpathlineto{\pgfqpoint{1.209760in}{1.146848in}}%
\pgfpathlineto{\pgfqpoint{1.213461in}{1.141563in}}%
\pgfpathlineto{\pgfqpoint{1.213054in}{1.137702in}}%
\pgfpathlineto{\pgfqpoint{1.205700in}{1.130953in}}%
\pgfpathlineto{\pgfqpoint{1.202180in}{1.129872in}}%
\pgfpathlineto{\pgfqpoint{1.193414in}{1.130886in}}%
\pgfpathlineto{\pgfqpoint{1.189161in}{1.136735in}}%
\pgfpathlineto{\pgfqpoint{1.187185in}{1.151775in}}%
\pgfpathlineto{\pgfqpoint{1.188060in}{1.158199in}}%
\pgfpathlineto{\pgfqpoint{1.193507in}{1.162201in}}%
\pgfpathlineto{\pgfqpoint{1.186288in}{1.164150in}}%
\pgfpathlineto{\pgfqpoint{1.180201in}{1.170746in}}%
\pgfpathlineto{\pgfqpoint{1.181782in}{1.174522in}}%
\pgfpathlineto{\pgfqpoint{1.201403in}{1.171929in}}%
\pgfpathclose%
\pgfusepath{fill}%
\end{pgfscope}%
\begin{pgfscope}%
\pgfpathrectangle{\pgfqpoint{0.100000in}{0.100000in}}{\pgfqpoint{3.007045in}{1.925000in}}%
\pgfusepath{clip}%
\pgfsetbuttcap%
\pgfsetmiterjoin%
\definecolor{currentfill}{rgb}{0.516863,0.735748,0.860192}%
\pgfsetfillcolor{currentfill}%
\pgfsetlinewidth{0.000000pt}%
\definecolor{currentstroke}{rgb}{0.000000,0.000000,0.000000}%
\pgfsetstrokecolor{currentstroke}%
\pgfsetstrokeopacity{0.000000}%
\pgfsetdash{}{0pt}%
\pgfpathmoveto{\pgfqpoint{2.114171in}{0.933474in}}%
\pgfpathlineto{\pgfqpoint{2.108828in}{0.933214in}}%
\pgfpathlineto{\pgfqpoint{2.108338in}{0.951144in}}%
\pgfpathlineto{\pgfqpoint{2.110066in}{0.954874in}}%
\pgfpathlineto{\pgfqpoint{2.116626in}{0.955211in}}%
\pgfpathlineto{\pgfqpoint{2.115651in}{0.980661in}}%
\pgfpathlineto{\pgfqpoint{2.117091in}{0.980740in}}%
\pgfpathlineto{\pgfqpoint{2.140155in}{0.982191in}}%
\pgfpathlineto{\pgfqpoint{2.140247in}{0.978276in}}%
\pgfpathlineto{\pgfqpoint{2.146159in}{0.970611in}}%
\pgfpathlineto{\pgfqpoint{2.146928in}{0.965917in}}%
\pgfpathlineto{\pgfqpoint{2.158552in}{0.966327in}}%
\pgfpathlineto{\pgfqpoint{2.166450in}{0.964466in}}%
\pgfpathlineto{\pgfqpoint{2.169480in}{0.968102in}}%
\pgfpathlineto{\pgfqpoint{2.169549in}{0.973787in}}%
\pgfpathlineto{\pgfqpoint{2.181638in}{0.973954in}}%
\pgfpathlineto{\pgfqpoint{2.188596in}{0.973610in}}%
\pgfpathlineto{\pgfqpoint{2.187491in}{0.967104in}}%
\pgfpathlineto{\pgfqpoint{2.188978in}{0.963509in}}%
\pgfpathlineto{\pgfqpoint{2.188043in}{0.950490in}}%
\pgfpathlineto{\pgfqpoint{2.170276in}{0.951262in}}%
\pgfpathlineto{\pgfqpoint{2.161575in}{0.944740in}}%
\pgfpathlineto{\pgfqpoint{2.161409in}{0.940303in}}%
\pgfpathlineto{\pgfqpoint{2.151093in}{0.939458in}}%
\pgfpathlineto{\pgfqpoint{2.144945in}{0.937174in}}%
\pgfpathlineto{\pgfqpoint{2.143491in}{0.939640in}}%
\pgfpathlineto{\pgfqpoint{2.136863in}{0.936941in}}%
\pgfpathlineto{\pgfqpoint{2.116000in}{0.935950in}}%
\pgfpathlineto{\pgfqpoint{2.114171in}{0.933474in}}%
\pgfpathclose%
\pgfusepath{fill}%
\end{pgfscope}%
\begin{pgfscope}%
\pgfpathrectangle{\pgfqpoint{0.100000in}{0.100000in}}{\pgfqpoint{3.007045in}{1.925000in}}%
\pgfusepath{clip}%
\pgfsetbuttcap%
\pgfsetmiterjoin%
\definecolor{currentfill}{rgb}{0.417086,0.680631,0.838231}%
\pgfsetfillcolor{currentfill}%
\pgfsetlinewidth{0.000000pt}%
\definecolor{currentstroke}{rgb}{0.000000,0.000000,0.000000}%
\pgfsetstrokecolor{currentstroke}%
\pgfsetstrokeopacity{0.000000}%
\pgfsetdash{}{0pt}%
\pgfpathmoveto{\pgfqpoint{0.920301in}{0.337334in}}%
\pgfpathlineto{\pgfqpoint{0.912069in}{0.341923in}}%
\pgfpathlineto{\pgfqpoint{0.914753in}{0.346694in}}%
\pgfpathlineto{\pgfqpoint{0.913691in}{0.348399in}}%
\pgfpathlineto{\pgfqpoint{0.917983in}{0.355693in}}%
\pgfpathlineto{\pgfqpoint{0.916310in}{0.356673in}}%
\pgfpathlineto{\pgfqpoint{0.917368in}{0.358439in}}%
\pgfpathlineto{\pgfqpoint{0.915470in}{0.359574in}}%
\pgfpathlineto{\pgfqpoint{0.916564in}{0.361368in}}%
\pgfpathlineto{\pgfqpoint{0.913444in}{0.363209in}}%
\pgfpathlineto{\pgfqpoint{0.912886in}{0.364295in}}%
\pgfpathlineto{\pgfqpoint{0.910185in}{0.362649in}}%
\pgfpathlineto{\pgfqpoint{0.907829in}{0.364423in}}%
\pgfpathlineto{\pgfqpoint{0.910203in}{0.367541in}}%
\pgfpathlineto{\pgfqpoint{0.909515in}{0.369615in}}%
\pgfpathlineto{\pgfqpoint{0.912562in}{0.370845in}}%
\pgfpathlineto{\pgfqpoint{0.914264in}{0.375762in}}%
\pgfpathlineto{\pgfqpoint{0.899657in}{0.384205in}}%
\pgfpathlineto{\pgfqpoint{0.894498in}{0.387400in}}%
\pgfpathlineto{\pgfqpoint{0.895631in}{0.389258in}}%
\pgfpathlineto{\pgfqpoint{0.886512in}{0.394965in}}%
\pgfpathlineto{\pgfqpoint{0.884245in}{0.391268in}}%
\pgfpathlineto{\pgfqpoint{0.880633in}{0.393554in}}%
\pgfpathlineto{\pgfqpoint{0.886359in}{0.402751in}}%
\pgfpathlineto{\pgfqpoint{0.875033in}{0.409830in}}%
\pgfpathlineto{\pgfqpoint{0.871795in}{0.412013in}}%
\pgfpathlineto{\pgfqpoint{0.869886in}{0.409139in}}%
\pgfpathlineto{\pgfqpoint{0.856121in}{0.418510in}}%
\pgfpathlineto{\pgfqpoint{0.847814in}{0.419834in}}%
\pgfpathlineto{\pgfqpoint{0.827109in}{0.423257in}}%
\pgfpathlineto{\pgfqpoint{0.826027in}{0.421875in}}%
\pgfpathlineto{\pgfqpoint{0.816153in}{0.429462in}}%
\pgfpathlineto{\pgfqpoint{0.809857in}{0.421160in}}%
\pgfpathlineto{\pgfqpoint{0.807974in}{0.419972in}}%
\pgfpathlineto{\pgfqpoint{0.805001in}{0.422296in}}%
\pgfpathlineto{\pgfqpoint{0.803696in}{0.420629in}}%
\pgfpathlineto{\pgfqpoint{0.800378in}{0.423260in}}%
\pgfpathlineto{\pgfqpoint{0.798438in}{0.420763in}}%
\pgfpathlineto{\pgfqpoint{0.796550in}{0.422002in}}%
\pgfpathlineto{\pgfqpoint{0.787335in}{0.429069in}}%
\pgfpathlineto{\pgfqpoint{0.776021in}{0.437944in}}%
\pgfpathlineto{\pgfqpoint{0.768555in}{0.444177in}}%
\pgfpathlineto{\pgfqpoint{0.770089in}{0.446056in}}%
\pgfpathlineto{\pgfqpoint{0.757361in}{0.456738in}}%
\pgfpathlineto{\pgfqpoint{0.755991in}{0.455118in}}%
\pgfpathlineto{\pgfqpoint{0.749271in}{0.460813in}}%
\pgfpathlineto{\pgfqpoint{0.747517in}{0.458771in}}%
\pgfpathlineto{\pgfqpoint{0.741635in}{0.463832in}}%
\pgfpathlineto{\pgfqpoint{0.733232in}{0.471329in}}%
\pgfpathlineto{\pgfqpoint{0.737384in}{0.476047in}}%
\pgfpathlineto{\pgfqpoint{0.736206in}{0.477104in}}%
\pgfpathlineto{\pgfqpoint{0.741905in}{0.483472in}}%
\pgfpathlineto{\pgfqpoint{0.742309in}{0.483156in}}%
\pgfpathlineto{\pgfqpoint{0.748071in}{0.489521in}}%
\pgfpathlineto{\pgfqpoint{0.750170in}{0.491012in}}%
\pgfpathlineto{\pgfqpoint{0.751165in}{0.490590in}}%
\pgfpathlineto{\pgfqpoint{0.755431in}{0.495345in}}%
\pgfpathlineto{\pgfqpoint{0.756937in}{0.494067in}}%
\pgfpathlineto{\pgfqpoint{0.761153in}{0.499198in}}%
\pgfpathlineto{\pgfqpoint{0.762810in}{0.497757in}}%
\pgfpathlineto{\pgfqpoint{0.764086in}{0.499185in}}%
\pgfpathlineto{\pgfqpoint{0.767409in}{0.496302in}}%
\pgfpathlineto{\pgfqpoint{0.768805in}{0.498004in}}%
\pgfpathlineto{\pgfqpoint{0.770286in}{0.496757in}}%
\pgfpathlineto{\pgfqpoint{0.771756in}{0.498478in}}%
\pgfpathlineto{\pgfqpoint{0.773225in}{0.497185in}}%
\pgfpathlineto{\pgfqpoint{0.777669in}{0.501945in}}%
\pgfpathlineto{\pgfqpoint{0.775774in}{0.503541in}}%
\pgfpathlineto{\pgfqpoint{0.787134in}{0.516986in}}%
\pgfpathlineto{\pgfqpoint{0.788832in}{0.515575in}}%
\pgfpathlineto{\pgfqpoint{0.791670in}{0.518914in}}%
\pgfpathlineto{\pgfqpoint{0.793312in}{0.517519in}}%
\pgfpathlineto{\pgfqpoint{0.798922in}{0.524065in}}%
\pgfpathlineto{\pgfqpoint{0.800473in}{0.522699in}}%
\pgfpathlineto{\pgfqpoint{0.804798in}{0.527778in}}%
\pgfpathlineto{\pgfqpoint{0.803129in}{0.529197in}}%
\pgfpathlineto{\pgfqpoint{0.810396in}{0.537663in}}%
\pgfpathlineto{\pgfqpoint{0.815410in}{0.533363in}}%
\pgfpathlineto{\pgfqpoint{0.818104in}{0.536633in}}%
\pgfpathlineto{\pgfqpoint{0.826641in}{0.529420in}}%
\pgfpathlineto{\pgfqpoint{0.832267in}{0.536448in}}%
\pgfpathlineto{\pgfqpoint{0.839219in}{0.530767in}}%
\pgfpathlineto{\pgfqpoint{0.836381in}{0.527215in}}%
\pgfpathlineto{\pgfqpoint{0.839868in}{0.524417in}}%
\pgfpathlineto{\pgfqpoint{0.842673in}{0.527910in}}%
\pgfpathlineto{\pgfqpoint{0.846127in}{0.525116in}}%
\pgfpathlineto{\pgfqpoint{0.844770in}{0.523398in}}%
\pgfpathlineto{\pgfqpoint{0.846505in}{0.521977in}}%
\pgfpathlineto{\pgfqpoint{0.845093in}{0.520174in}}%
\pgfpathlineto{\pgfqpoint{0.848387in}{0.517598in}}%
\pgfpathlineto{\pgfqpoint{0.852921in}{0.522581in}}%
\pgfpathlineto{\pgfqpoint{0.857919in}{0.518605in}}%
\pgfpathlineto{\pgfqpoint{0.860718in}{0.522144in}}%
\pgfpathlineto{\pgfqpoint{0.865978in}{0.517880in}}%
\pgfpathlineto{\pgfqpoint{0.867902in}{0.519376in}}%
\pgfpathlineto{\pgfqpoint{0.873845in}{0.526915in}}%
\pgfpathlineto{\pgfqpoint{0.872574in}{0.527836in}}%
\pgfpathlineto{\pgfqpoint{0.873947in}{0.529693in}}%
\pgfpathlineto{\pgfqpoint{0.870380in}{0.532422in}}%
\pgfpathlineto{\pgfqpoint{0.874516in}{0.537747in}}%
\pgfpathlineto{\pgfqpoint{0.876407in}{0.539190in}}%
\pgfpathlineto{\pgfqpoint{0.874603in}{0.540599in}}%
\pgfpathlineto{\pgfqpoint{0.876010in}{0.542324in}}%
\pgfpathlineto{\pgfqpoint{0.874170in}{0.543737in}}%
\pgfpathlineto{\pgfqpoint{0.875594in}{0.545533in}}%
\pgfpathlineto{\pgfqpoint{0.873853in}{0.546928in}}%
\pgfpathlineto{\pgfqpoint{0.877896in}{0.551221in}}%
\pgfpathlineto{\pgfqpoint{0.895892in}{0.537283in}}%
\pgfpathlineto{\pgfqpoint{0.919823in}{0.519872in}}%
\pgfpathlineto{\pgfqpoint{0.941118in}{0.505304in}}%
\pgfpathlineto{\pgfqpoint{0.944682in}{0.503522in}}%
\pgfpathlineto{\pgfqpoint{0.952820in}{0.498342in}}%
\pgfpathlineto{\pgfqpoint{0.959476in}{0.508253in}}%
\pgfpathlineto{\pgfqpoint{0.966890in}{0.503916in}}%
\pgfpathlineto{\pgfqpoint{0.989639in}{0.489920in}}%
\pgfpathlineto{\pgfqpoint{1.001047in}{0.482991in}}%
\pgfpathlineto{\pgfqpoint{0.992914in}{0.468320in}}%
\pgfpathlineto{\pgfqpoint{0.981844in}{0.448351in}}%
\pgfpathlineto{\pgfqpoint{0.970915in}{0.428636in}}%
\pgfpathlineto{\pgfqpoint{0.960932in}{0.410628in}}%
\pgfpathlineto{\pgfqpoint{0.950240in}{0.391342in}}%
\pgfpathlineto{\pgfqpoint{0.939656in}{0.372249in}}%
\pgfpathlineto{\pgfqpoint{0.929170in}{0.353333in}}%
\pgfpathlineto{\pgfqpoint{0.920301in}{0.337334in}}%
\pgfpathclose%
\pgfusepath{fill}%
\end{pgfscope}%
\begin{pgfscope}%
\pgfpathrectangle{\pgfqpoint{0.100000in}{0.100000in}}{\pgfqpoint{3.007045in}{1.925000in}}%
\pgfusepath{clip}%
\pgfsetbuttcap%
\pgfsetmiterjoin%
\definecolor{currentfill}{rgb}{0.396909,0.666851,0.830358}%
\pgfsetfillcolor{currentfill}%
\pgfsetlinewidth{0.000000pt}%
\definecolor{currentstroke}{rgb}{0.000000,0.000000,0.000000}%
\pgfsetstrokecolor{currentstroke}%
\pgfsetstrokeopacity{0.000000}%
\pgfsetdash{}{0pt}%
\pgfpathmoveto{\pgfqpoint{1.924728in}{0.990132in}}%
\pgfpathlineto{\pgfqpoint{1.890280in}{0.989939in}}%
\pgfpathlineto{\pgfqpoint{1.859810in}{0.990030in}}%
\pgfpathlineto{\pgfqpoint{1.859842in}{1.001580in}}%
\pgfpathlineto{\pgfqpoint{1.845487in}{1.001557in}}%
\pgfpathlineto{\pgfqpoint{1.845462in}{1.020387in}}%
\pgfpathlineto{\pgfqpoint{1.844374in}{1.020395in}}%
\pgfpathlineto{\pgfqpoint{1.844766in}{1.050707in}}%
\pgfpathlineto{\pgfqpoint{1.843696in}{1.056469in}}%
\pgfpathlineto{\pgfqpoint{1.846715in}{1.062147in}}%
\pgfpathlineto{\pgfqpoint{1.849656in}{1.062106in}}%
\pgfpathlineto{\pgfqpoint{1.849904in}{1.072618in}}%
\pgfpathlineto{\pgfqpoint{1.872670in}{1.072352in}}%
\pgfpathlineto{\pgfqpoint{1.872518in}{1.061815in}}%
\pgfpathlineto{\pgfqpoint{1.896036in}{1.061694in}}%
\pgfpathlineto{\pgfqpoint{1.900128in}{1.055674in}}%
\pgfpathlineto{\pgfqpoint{1.907064in}{1.059811in}}%
\pgfpathlineto{\pgfqpoint{1.907105in}{1.050239in}}%
\pgfpathlineto{\pgfqpoint{1.915821in}{1.046178in}}%
\pgfpathlineto{\pgfqpoint{1.915878in}{1.043286in}}%
\pgfpathlineto{\pgfqpoint{1.916864in}{1.007392in}}%
\pgfpathlineto{\pgfqpoint{1.925420in}{1.007434in}}%
\pgfpathlineto{\pgfqpoint{1.924728in}{0.990132in}}%
\pgfpathclose%
\pgfusepath{fill}%
\end{pgfscope}%
\begin{pgfscope}%
\pgfpathrectangle{\pgfqpoint{0.100000in}{0.100000in}}{\pgfqpoint{3.007045in}{1.925000in}}%
\pgfusepath{clip}%
\pgfsetbuttcap%
\pgfsetmiterjoin%
\definecolor{currentfill}{rgb}{0.275848,0.584175,0.783114}%
\pgfsetfillcolor{currentfill}%
\pgfsetlinewidth{0.000000pt}%
\definecolor{currentstroke}{rgb}{0.000000,0.000000,0.000000}%
\pgfsetstrokecolor{currentstroke}%
\pgfsetstrokeopacity{0.000000}%
\pgfsetdash{}{0pt}%
\pgfpathmoveto{\pgfqpoint{1.829586in}{0.787129in}}%
\pgfpathlineto{\pgfqpoint{1.825054in}{0.789133in}}%
\pgfpathlineto{\pgfqpoint{1.814189in}{0.790353in}}%
\pgfpathlineto{\pgfqpoint{1.806094in}{0.799451in}}%
\pgfpathlineto{\pgfqpoint{1.801001in}{0.799445in}}%
\pgfpathlineto{\pgfqpoint{1.800644in}{0.779398in}}%
\pgfpathlineto{\pgfqpoint{1.786299in}{0.783647in}}%
\pgfpathlineto{\pgfqpoint{1.779488in}{0.786480in}}%
\pgfpathlineto{\pgfqpoint{1.774881in}{0.791133in}}%
\pgfpathlineto{\pgfqpoint{1.760297in}{0.800742in}}%
\pgfpathlineto{\pgfqpoint{1.756132in}{0.794865in}}%
\pgfpathlineto{\pgfqpoint{1.749934in}{0.794479in}}%
\pgfpathlineto{\pgfqpoint{1.742020in}{0.795983in}}%
\pgfpathlineto{\pgfqpoint{1.740181in}{0.799040in}}%
\pgfpathlineto{\pgfqpoint{1.730486in}{0.795042in}}%
\pgfpathlineto{\pgfqpoint{1.726173in}{0.795162in}}%
\pgfpathlineto{\pgfqpoint{1.720879in}{0.800114in}}%
\pgfpathlineto{\pgfqpoint{1.720987in}{0.811838in}}%
\pgfpathlineto{\pgfqpoint{1.718089in}{0.813933in}}%
\pgfpathlineto{\pgfqpoint{1.729695in}{0.813812in}}%
\pgfpathlineto{\pgfqpoint{1.729891in}{0.836933in}}%
\pgfpathlineto{\pgfqpoint{1.735756in}{0.836902in}}%
\pgfpathlineto{\pgfqpoint{1.735784in}{0.842654in}}%
\pgfpathlineto{\pgfqpoint{1.744358in}{0.842618in}}%
\pgfpathlineto{\pgfqpoint{1.744379in}{0.848364in}}%
\pgfpathlineto{\pgfqpoint{1.775709in}{0.848268in}}%
\pgfpathlineto{\pgfqpoint{1.775725in}{0.836774in}}%
\pgfpathlineto{\pgfqpoint{1.801665in}{0.836897in}}%
\pgfpathlineto{\pgfqpoint{1.801936in}{0.851532in}}%
\pgfpathlineto{\pgfqpoint{1.809383in}{0.849499in}}%
\pgfpathlineto{\pgfqpoint{1.825651in}{0.849322in}}%
\pgfpathlineto{\pgfqpoint{1.830433in}{0.847601in}}%
\pgfpathlineto{\pgfqpoint{1.830412in}{0.826752in}}%
\pgfpathlineto{\pgfqpoint{1.812989in}{0.826934in}}%
\pgfpathlineto{\pgfqpoint{1.813675in}{0.816167in}}%
\pgfpathlineto{\pgfqpoint{1.821492in}{0.816167in}}%
\pgfpathlineto{\pgfqpoint{1.822337in}{0.805756in}}%
\pgfpathlineto{\pgfqpoint{1.824314in}{0.796016in}}%
\pgfpathlineto{\pgfqpoint{1.828440in}{0.791464in}}%
\pgfpathlineto{\pgfqpoint{1.829586in}{0.787129in}}%
\pgfpathclose%
\pgfusepath{fill}%
\end{pgfscope}%
\begin{pgfscope}%
\pgfpathrectangle{\pgfqpoint{0.100000in}{0.100000in}}{\pgfqpoint{3.007045in}{1.925000in}}%
\pgfusepath{clip}%
\pgfsetbuttcap%
\pgfsetmiterjoin%
\definecolor{currentfill}{rgb}{0.336378,0.625513,0.806736}%
\pgfsetfillcolor{currentfill}%
\pgfsetlinewidth{0.000000pt}%
\definecolor{currentstroke}{rgb}{0.000000,0.000000,0.000000}%
\pgfsetstrokecolor{currentstroke}%
\pgfsetstrokeopacity{0.000000}%
\pgfsetdash{}{0pt}%
\pgfpathmoveto{\pgfqpoint{1.798597in}{1.275256in}}%
\pgfpathlineto{\pgfqpoint{1.787185in}{1.275241in}}%
\pgfpathlineto{\pgfqpoint{1.787084in}{1.304515in}}%
\pgfpathlineto{\pgfqpoint{1.784942in}{1.304507in}}%
\pgfpathlineto{\pgfqpoint{1.784860in}{1.321668in}}%
\pgfpathlineto{\pgfqpoint{1.813376in}{1.321894in}}%
\pgfpathlineto{\pgfqpoint{1.813145in}{1.344757in}}%
\pgfpathlineto{\pgfqpoint{1.847419in}{1.345161in}}%
\pgfpathlineto{\pgfqpoint{1.858723in}{1.345342in}}%
\pgfpathlineto{\pgfqpoint{1.859156in}{1.322514in}}%
\pgfpathlineto{\pgfqpoint{1.853457in}{1.322433in}}%
\pgfpathlineto{\pgfqpoint{1.853741in}{1.305138in}}%
\pgfpathlineto{\pgfqpoint{1.854836in}{1.299029in}}%
\pgfpathlineto{\pgfqpoint{1.855298in}{1.276147in}}%
\pgfpathlineto{\pgfqpoint{1.843908in}{1.275998in}}%
\pgfpathlineto{\pgfqpoint{1.826979in}{1.275432in}}%
\pgfpathlineto{\pgfqpoint{1.798597in}{1.275256in}}%
\pgfpathclose%
\pgfusepath{fill}%
\end{pgfscope}%
\begin{pgfscope}%
\pgfpathrectangle{\pgfqpoint{0.100000in}{0.100000in}}{\pgfqpoint{3.007045in}{1.925000in}}%
\pgfusepath{clip}%
\pgfsetbuttcap%
\pgfsetmiterjoin%
\definecolor{currentfill}{rgb}{0.108651,0.416563,0.689043}%
\pgfsetfillcolor{currentfill}%
\pgfsetlinewidth{0.000000pt}%
\definecolor{currentstroke}{rgb}{0.000000,0.000000,0.000000}%
\pgfsetstrokecolor{currentstroke}%
\pgfsetstrokeopacity{0.000000}%
\pgfsetdash{}{0pt}%
\pgfpathmoveto{\pgfqpoint{1.502115in}{0.878788in}}%
\pgfpathlineto{\pgfqpoint{1.501084in}{0.860066in}}%
\pgfpathlineto{\pgfqpoint{1.478574in}{0.861428in}}%
\pgfpathlineto{\pgfqpoint{1.471813in}{0.861849in}}%
\pgfpathlineto{\pgfqpoint{1.473659in}{0.890543in}}%
\pgfpathlineto{\pgfqpoint{1.475346in}{0.919280in}}%
\pgfpathlineto{\pgfqpoint{1.504252in}{0.917589in}}%
\pgfpathlineto{\pgfqpoint{1.502115in}{0.878788in}}%
\pgfpathclose%
\pgfusepath{fill}%
\end{pgfscope}%
\begin{pgfscope}%
\pgfpathrectangle{\pgfqpoint{0.100000in}{0.100000in}}{\pgfqpoint{3.007045in}{1.925000in}}%
\pgfusepath{clip}%
\pgfsetbuttcap%
\pgfsetmiterjoin%
\definecolor{currentfill}{rgb}{0.191326,0.505052,0.741699}%
\pgfsetfillcolor{currentfill}%
\pgfsetlinewidth{0.000000pt}%
\definecolor{currentstroke}{rgb}{0.000000,0.000000,0.000000}%
\pgfsetstrokecolor{currentstroke}%
\pgfsetstrokeopacity{0.000000}%
\pgfsetdash{}{0pt}%
\pgfpathmoveto{\pgfqpoint{1.594488in}{0.908994in}}%
\pgfpathlineto{\pgfqpoint{1.577863in}{0.909634in}}%
\pgfpathlineto{\pgfqpoint{1.578070in}{0.926846in}}%
\pgfpathlineto{\pgfqpoint{1.538310in}{0.928586in}}%
\pgfpathlineto{\pgfqpoint{1.538682in}{0.941826in}}%
\pgfpathlineto{\pgfqpoint{1.539102in}{0.951802in}}%
\pgfpathlineto{\pgfqpoint{1.561768in}{0.950581in}}%
\pgfpathlineto{\pgfqpoint{1.562805in}{0.979407in}}%
\pgfpathlineto{\pgfqpoint{1.553847in}{0.991951in}}%
\pgfpathlineto{\pgfqpoint{1.547875in}{0.992381in}}%
\pgfpathlineto{\pgfqpoint{1.541765in}{0.998932in}}%
\pgfpathlineto{\pgfqpoint{1.537864in}{1.006990in}}%
\pgfpathlineto{\pgfqpoint{1.561807in}{1.005883in}}%
\pgfpathlineto{\pgfqpoint{1.596180in}{1.004385in}}%
\pgfpathlineto{\pgfqpoint{1.642737in}{1.002945in}}%
\pgfpathlineto{\pgfqpoint{1.642047in}{0.976241in}}%
\pgfpathlineto{\pgfqpoint{1.641329in}{0.947962in}}%
\pgfpathlineto{\pgfqpoint{1.601428in}{0.949236in}}%
\pgfpathlineto{\pgfqpoint{1.600584in}{0.920240in}}%
\pgfpathlineto{\pgfqpoint{1.594900in}{0.920448in}}%
\pgfpathlineto{\pgfqpoint{1.594488in}{0.908994in}}%
\pgfpathclose%
\pgfusepath{fill}%
\end{pgfscope}%
\begin{pgfscope}%
\pgfpathrectangle{\pgfqpoint{0.100000in}{0.100000in}}{\pgfqpoint{3.007045in}{1.925000in}}%
\pgfusepath{clip}%
\pgfsetbuttcap%
\pgfsetmiterjoin%
\definecolor{currentfill}{rgb}{0.031373,0.188235,0.419608}%
\pgfsetfillcolor{currentfill}%
\pgfsetlinewidth{0.000000pt}%
\definecolor{currentstroke}{rgb}{0.000000,0.000000,0.000000}%
\pgfsetstrokecolor{currentstroke}%
\pgfsetstrokeopacity{0.000000}%
\pgfsetdash{}{0pt}%
\pgfpathmoveto{\pgfqpoint{0.907070in}{1.217261in}}%
\pgfpathlineto{\pgfqpoint{0.903422in}{1.212075in}}%
\pgfpathlineto{\pgfqpoint{0.900045in}{1.212653in}}%
\pgfpathlineto{\pgfqpoint{0.894249in}{1.207621in}}%
\pgfpathlineto{\pgfqpoint{0.892855in}{1.198668in}}%
\pgfpathlineto{\pgfqpoint{0.887613in}{1.196226in}}%
\pgfpathlineto{\pgfqpoint{0.880713in}{1.197892in}}%
\pgfpathlineto{\pgfqpoint{0.875936in}{1.191519in}}%
\pgfpathlineto{\pgfqpoint{0.850288in}{1.196720in}}%
\pgfpathlineto{\pgfqpoint{0.798495in}{1.207353in}}%
\pgfpathlineto{\pgfqpoint{0.799942in}{1.214078in}}%
\pgfpathlineto{\pgfqpoint{0.811909in}{1.269809in}}%
\pgfpathlineto{\pgfqpoint{0.816917in}{1.293210in}}%
\pgfpathlineto{\pgfqpoint{0.878712in}{1.280320in}}%
\pgfpathlineto{\pgfqpoint{0.901338in}{1.275958in}}%
\pgfpathlineto{\pgfqpoint{0.910706in}{1.281361in}}%
\pgfpathlineto{\pgfqpoint{0.914434in}{1.291420in}}%
\pgfpathlineto{\pgfqpoint{0.913219in}{1.297816in}}%
\pgfpathlineto{\pgfqpoint{0.914647in}{1.310630in}}%
\pgfpathlineto{\pgfqpoint{0.923286in}{1.309346in}}%
\pgfpathlineto{\pgfqpoint{0.925436in}{1.305547in}}%
\pgfpathlineto{\pgfqpoint{0.937610in}{1.310762in}}%
\pgfpathlineto{\pgfqpoint{0.947451in}{1.310370in}}%
\pgfpathlineto{\pgfqpoint{0.943668in}{1.304520in}}%
\pgfpathlineto{\pgfqpoint{0.950156in}{1.292847in}}%
\pgfpathlineto{\pgfqpoint{0.957665in}{1.289984in}}%
\pgfpathlineto{\pgfqpoint{0.959849in}{1.282582in}}%
\pgfpathlineto{\pgfqpoint{0.957377in}{1.275606in}}%
\pgfpathlineto{\pgfqpoint{0.964301in}{1.266613in}}%
\pgfpathlineto{\pgfqpoint{0.963825in}{1.264056in}}%
\pgfpathlineto{\pgfqpoint{0.973337in}{1.262333in}}%
\pgfpathlineto{\pgfqpoint{0.974041in}{1.256437in}}%
\pgfpathlineto{\pgfqpoint{0.934985in}{1.263476in}}%
\pgfpathlineto{\pgfqpoint{0.937487in}{1.261054in}}%
\pgfpathlineto{\pgfqpoint{0.929754in}{1.256515in}}%
\pgfpathlineto{\pgfqpoint{0.927198in}{1.241278in}}%
\pgfpathlineto{\pgfqpoint{0.924573in}{1.236589in}}%
\pgfpathlineto{\pgfqpoint{0.915780in}{1.238145in}}%
\pgfpathlineto{\pgfqpoint{0.910572in}{1.236135in}}%
\pgfpathlineto{\pgfqpoint{0.907070in}{1.217261in}}%
\pgfpathclose%
\pgfusepath{fill}%
\end{pgfscope}%
\begin{pgfscope}%
\pgfpathrectangle{\pgfqpoint{0.100000in}{0.100000in}}{\pgfqpoint{3.007045in}{1.925000in}}%
\pgfusepath{clip}%
\pgfsetbuttcap%
\pgfsetmiterjoin%
\definecolor{currentfill}{rgb}{0.460392,0.704744,0.848012}%
\pgfsetfillcolor{currentfill}%
\pgfsetlinewidth{0.000000pt}%
\definecolor{currentstroke}{rgb}{0.000000,0.000000,0.000000}%
\pgfsetstrokecolor{currentstroke}%
\pgfsetstrokeopacity{0.000000}%
\pgfsetdash{}{0pt}%
\pgfpathmoveto{\pgfqpoint{2.375866in}{1.305055in}}%
\pgfpathlineto{\pgfqpoint{2.371168in}{1.333253in}}%
\pgfpathlineto{\pgfqpoint{2.365754in}{1.334376in}}%
\pgfpathlineto{\pgfqpoint{2.378309in}{1.338718in}}%
\pgfpathlineto{\pgfqpoint{2.383725in}{1.333912in}}%
\pgfpathlineto{\pgfqpoint{2.389914in}{1.332610in}}%
\pgfpathlineto{\pgfqpoint{2.400156in}{1.337527in}}%
\pgfpathlineto{\pgfqpoint{2.411268in}{1.344518in}}%
\pgfpathlineto{\pgfqpoint{2.413516in}{1.344212in}}%
\pgfpathlineto{\pgfqpoint{2.414804in}{1.334085in}}%
\pgfpathlineto{\pgfqpoint{2.419430in}{1.334702in}}%
\pgfpathlineto{\pgfqpoint{2.420098in}{1.329773in}}%
\pgfpathlineto{\pgfqpoint{2.415460in}{1.329120in}}%
\pgfpathlineto{\pgfqpoint{2.416075in}{1.324178in}}%
\pgfpathlineto{\pgfqpoint{2.411263in}{1.323487in}}%
\pgfpathlineto{\pgfqpoint{2.407064in}{1.318707in}}%
\pgfpathlineto{\pgfqpoint{2.407666in}{1.313899in}}%
\pgfpathlineto{\pgfqpoint{2.394589in}{1.312199in}}%
\pgfpathlineto{\pgfqpoint{2.395446in}{1.307499in}}%
\pgfpathlineto{\pgfqpoint{2.375866in}{1.305055in}}%
\pgfpathclose%
\pgfusepath{fill}%
\end{pgfscope}%
\begin{pgfscope}%
\pgfpathrectangle{\pgfqpoint{0.100000in}{0.100000in}}{\pgfqpoint{3.007045in}{1.925000in}}%
\pgfusepath{clip}%
\pgfsetbuttcap%
\pgfsetmiterjoin%
\definecolor{currentfill}{rgb}{0.183206,0.496932,0.737516}%
\pgfsetfillcolor{currentfill}%
\pgfsetlinewidth{0.000000pt}%
\definecolor{currentstroke}{rgb}{0.000000,0.000000,0.000000}%
\pgfsetstrokecolor{currentstroke}%
\pgfsetstrokeopacity{0.000000}%
\pgfsetdash{}{0pt}%
\pgfpathmoveto{\pgfqpoint{1.359600in}{1.651719in}}%
\pgfpathlineto{\pgfqpoint{1.355655in}{1.612365in}}%
\pgfpathlineto{\pgfqpoint{1.351714in}{1.614337in}}%
\pgfpathlineto{\pgfqpoint{1.339970in}{1.615629in}}%
\pgfpathlineto{\pgfqpoint{1.340576in}{1.621215in}}%
\pgfpathlineto{\pgfqpoint{1.334802in}{1.621794in}}%
\pgfpathlineto{\pgfqpoint{1.335389in}{1.627188in}}%
\pgfpathlineto{\pgfqpoint{1.318468in}{1.629024in}}%
\pgfpathlineto{\pgfqpoint{1.320048in}{1.651933in}}%
\pgfpathlineto{\pgfqpoint{1.327908in}{1.651084in}}%
\pgfpathlineto{\pgfqpoint{1.329154in}{1.662581in}}%
\pgfpathlineto{\pgfqpoint{1.335063in}{1.663866in}}%
\pgfpathlineto{\pgfqpoint{1.343644in}{1.662962in}}%
\pgfpathlineto{\pgfqpoint{1.346246in}{1.659789in}}%
\pgfpathlineto{\pgfqpoint{1.360257in}{1.658366in}}%
\pgfpathlineto{\pgfqpoint{1.359600in}{1.651719in}}%
\pgfpathclose%
\pgfusepath{fill}%
\end{pgfscope}%
\begin{pgfscope}%
\pgfpathrectangle{\pgfqpoint{0.100000in}{0.100000in}}{\pgfqpoint{3.007045in}{1.925000in}}%
\pgfusepath{clip}%
\pgfsetbuttcap%
\pgfsetmiterjoin%
\definecolor{currentfill}{rgb}{0.592157,0.777086,0.876432}%
\pgfsetfillcolor{currentfill}%
\pgfsetlinewidth{0.000000pt}%
\definecolor{currentstroke}{rgb}{0.000000,0.000000,0.000000}%
\pgfsetstrokecolor{currentstroke}%
\pgfsetstrokeopacity{0.000000}%
\pgfsetdash{}{0pt}%
\pgfpathmoveto{\pgfqpoint{2.765291in}{1.169198in}}%
\pgfpathlineto{\pgfqpoint{2.759352in}{1.165916in}}%
\pgfpathlineto{\pgfqpoint{2.756954in}{1.169348in}}%
\pgfpathlineto{\pgfqpoint{2.759204in}{1.174214in}}%
\pgfpathlineto{\pgfqpoint{2.757805in}{1.180412in}}%
\pgfpathlineto{\pgfqpoint{2.751639in}{1.179119in}}%
\pgfpathlineto{\pgfqpoint{2.751955in}{1.186579in}}%
\pgfpathlineto{\pgfqpoint{2.751060in}{1.190437in}}%
\pgfpathlineto{\pgfqpoint{2.753781in}{1.202785in}}%
\pgfpathlineto{\pgfqpoint{2.758832in}{1.209240in}}%
\pgfpathlineto{\pgfqpoint{2.745705in}{1.255995in}}%
\pgfpathlineto{\pgfqpoint{2.751195in}{1.256851in}}%
\pgfpathlineto{\pgfqpoint{2.757050in}{1.263166in}}%
\pgfpathlineto{\pgfqpoint{2.761478in}{1.260531in}}%
\pgfpathlineto{\pgfqpoint{2.765354in}{1.253983in}}%
\pgfpathlineto{\pgfqpoint{2.766799in}{1.245106in}}%
\pgfpathlineto{\pgfqpoint{2.770462in}{1.242822in}}%
\pgfpathlineto{\pgfqpoint{2.773747in}{1.236358in}}%
\pgfpathlineto{\pgfqpoint{2.781281in}{1.230541in}}%
\pgfpathlineto{\pgfqpoint{2.786967in}{1.229663in}}%
\pgfpathlineto{\pgfqpoint{2.793235in}{1.209405in}}%
\pgfpathlineto{\pgfqpoint{2.791000in}{1.184936in}}%
\pgfpathlineto{\pgfqpoint{2.789514in}{1.179938in}}%
\pgfpathlineto{\pgfqpoint{2.770616in}{1.173606in}}%
\pgfpathlineto{\pgfqpoint{2.765291in}{1.169198in}}%
\pgfpathclose%
\pgfusepath{fill}%
\end{pgfscope}%
\begin{pgfscope}%
\pgfpathrectangle{\pgfqpoint{0.100000in}{0.100000in}}{\pgfqpoint{3.007045in}{1.925000in}}%
\pgfusepath{clip}%
\pgfsetbuttcap%
\pgfsetmiterjoin%
\definecolor{currentfill}{rgb}{0.321246,0.615179,0.800830}%
\pgfsetfillcolor{currentfill}%
\pgfsetlinewidth{0.000000pt}%
\definecolor{currentstroke}{rgb}{0.000000,0.000000,0.000000}%
\pgfsetstrokecolor{currentstroke}%
\pgfsetstrokeopacity{0.000000}%
\pgfsetdash{}{0pt}%
\pgfpathmoveto{\pgfqpoint{1.779241in}{1.344556in}}%
\pgfpathlineto{\pgfqpoint{1.739541in}{1.344779in}}%
\pgfpathlineto{\pgfqpoint{1.739732in}{1.362155in}}%
\pgfpathlineto{\pgfqpoint{1.736283in}{1.367899in}}%
\pgfpathlineto{\pgfqpoint{1.730520in}{1.367928in}}%
\pgfpathlineto{\pgfqpoint{1.730762in}{1.390911in}}%
\pgfpathlineto{\pgfqpoint{1.753560in}{1.390798in}}%
\pgfpathlineto{\pgfqpoint{1.776480in}{1.390760in}}%
\pgfpathlineto{\pgfqpoint{1.776477in}{1.362015in}}%
\pgfpathlineto{\pgfqpoint{1.779293in}{1.361999in}}%
\pgfpathlineto{\pgfqpoint{1.779241in}{1.344556in}}%
\pgfpathclose%
\pgfusepath{fill}%
\end{pgfscope}%
\begin{pgfscope}%
\pgfpathrectangle{\pgfqpoint{0.100000in}{0.100000in}}{\pgfqpoint{3.007045in}{1.925000in}}%
\pgfusepath{clip}%
\pgfsetbuttcap%
\pgfsetmiterjoin%
\definecolor{currentfill}{rgb}{0.031757,0.318139,0.612149}%
\pgfsetfillcolor{currentfill}%
\pgfsetlinewidth{0.000000pt}%
\definecolor{currentstroke}{rgb}{0.000000,0.000000,0.000000}%
\pgfsetstrokecolor{currentstroke}%
\pgfsetstrokeopacity{0.000000}%
\pgfsetdash{}{0pt}%
\pgfpathmoveto{\pgfqpoint{1.011802in}{1.226459in}}%
\pgfpathlineto{\pgfqpoint{1.057846in}{1.219034in}}%
\pgfpathlineto{\pgfqpoint{1.061002in}{1.220921in}}%
\pgfpathlineto{\pgfqpoint{1.059704in}{1.212390in}}%
\pgfpathlineto{\pgfqpoint{1.050444in}{1.156057in}}%
\pgfpathlineto{\pgfqpoint{1.001312in}{1.163987in}}%
\pgfpathlineto{\pgfqpoint{1.002802in}{1.166666in}}%
\pgfpathlineto{\pgfqpoint{0.999346in}{1.179674in}}%
\pgfpathlineto{\pgfqpoint{1.000947in}{1.182123in}}%
\pgfpathlineto{\pgfqpoint{0.998322in}{1.188795in}}%
\pgfpathlineto{\pgfqpoint{1.001745in}{1.203071in}}%
\pgfpathlineto{\pgfqpoint{1.007785in}{1.211607in}}%
\pgfpathlineto{\pgfqpoint{1.007174in}{1.215093in}}%
\pgfpathlineto{\pgfqpoint{1.010713in}{1.219317in}}%
\pgfpathlineto{\pgfqpoint{1.011802in}{1.226459in}}%
\pgfpathclose%
\pgfusepath{fill}%
\end{pgfscope}%
\begin{pgfscope}%
\pgfpathrectangle{\pgfqpoint{0.100000in}{0.100000in}}{\pgfqpoint{3.007045in}{1.925000in}}%
\pgfusepath{clip}%
\pgfsetbuttcap%
\pgfsetmiterjoin%
\definecolor{currentfill}{rgb}{0.585882,0.773641,0.875079}%
\pgfsetfillcolor{currentfill}%
\pgfsetlinewidth{0.000000pt}%
\definecolor{currentstroke}{rgb}{0.000000,0.000000,0.000000}%
\pgfsetstrokecolor{currentstroke}%
\pgfsetstrokeopacity{0.000000}%
\pgfsetdash{}{0pt}%
\pgfpathmoveto{\pgfqpoint{2.334077in}{0.871563in}}%
\pgfpathlineto{\pgfqpoint{2.332579in}{0.883358in}}%
\pgfpathlineto{\pgfqpoint{2.335741in}{0.890096in}}%
\pgfpathlineto{\pgfqpoint{2.333466in}{0.891984in}}%
\pgfpathlineto{\pgfqpoint{2.332547in}{0.900599in}}%
\pgfpathlineto{\pgfqpoint{2.348646in}{0.902478in}}%
\pgfpathlineto{\pgfqpoint{2.356919in}{0.901553in}}%
\pgfpathlineto{\pgfqpoint{2.362404in}{0.891696in}}%
\pgfpathlineto{\pgfqpoint{2.360142in}{0.881259in}}%
\pgfpathlineto{\pgfqpoint{2.358828in}{0.878027in}}%
\pgfpathlineto{\pgfqpoint{2.351405in}{0.875886in}}%
\pgfpathlineto{\pgfqpoint{2.349378in}{0.873307in}}%
\pgfpathlineto{\pgfqpoint{2.334077in}{0.871563in}}%
\pgfpathclose%
\pgfusepath{fill}%
\end{pgfscope}%
\begin{pgfscope}%
\pgfpathrectangle{\pgfqpoint{0.100000in}{0.100000in}}{\pgfqpoint{3.007045in}{1.925000in}}%
\pgfusepath{clip}%
\pgfsetbuttcap%
\pgfsetmiterjoin%
\definecolor{currentfill}{rgb}{0.691580,0.822745,0.907543}%
\pgfsetfillcolor{currentfill}%
\pgfsetlinewidth{0.000000pt}%
\definecolor{currentstroke}{rgb}{0.000000,0.000000,0.000000}%
\pgfsetstrokecolor{currentstroke}%
\pgfsetstrokeopacity{0.000000}%
\pgfsetdash{}{0pt}%
\pgfpathmoveto{\pgfqpoint{0.372778in}{1.622728in}}%
\pgfpathlineto{\pgfqpoint{0.383904in}{1.641105in}}%
\pgfpathlineto{\pgfqpoint{0.396803in}{1.656959in}}%
\pgfpathlineto{\pgfqpoint{0.406441in}{1.674762in}}%
\pgfpathlineto{\pgfqpoint{0.416059in}{1.671542in}}%
\pgfpathlineto{\pgfqpoint{0.423095in}{1.674955in}}%
\pgfpathlineto{\pgfqpoint{0.432010in}{1.670306in}}%
\pgfpathlineto{\pgfqpoint{0.434242in}{1.663279in}}%
\pgfpathlineto{\pgfqpoint{0.441475in}{1.657513in}}%
\pgfpathlineto{\pgfqpoint{0.451005in}{1.654418in}}%
\pgfpathlineto{\pgfqpoint{0.447535in}{1.643476in}}%
\pgfpathlineto{\pgfqpoint{0.447557in}{1.638893in}}%
\pgfpathlineto{\pgfqpoint{0.464341in}{1.633883in}}%
\pgfpathlineto{\pgfqpoint{0.462254in}{1.627100in}}%
\pgfpathlineto{\pgfqpoint{0.490196in}{1.618627in}}%
\pgfpathlineto{\pgfqpoint{0.493741in}{1.612121in}}%
\pgfpathlineto{\pgfqpoint{0.493766in}{1.605432in}}%
\pgfpathlineto{\pgfqpoint{0.487657in}{1.599535in}}%
\pgfpathlineto{\pgfqpoint{0.485002in}{1.594976in}}%
\pgfpathlineto{\pgfqpoint{0.476029in}{1.597225in}}%
\pgfpathlineto{\pgfqpoint{0.474662in}{1.592743in}}%
\pgfpathlineto{\pgfqpoint{0.469189in}{1.594432in}}%
\pgfpathlineto{\pgfqpoint{0.465330in}{1.591720in}}%
\pgfpathlineto{\pgfqpoint{0.460693in}{1.593122in}}%
\pgfpathlineto{\pgfqpoint{0.446705in}{1.586464in}}%
\pgfpathlineto{\pgfqpoint{0.436255in}{1.587889in}}%
\pgfpathlineto{\pgfqpoint{0.429770in}{1.590236in}}%
\pgfpathlineto{\pgfqpoint{0.426249in}{1.588654in}}%
\pgfpathlineto{\pgfqpoint{0.419385in}{1.591902in}}%
\pgfpathlineto{\pgfqpoint{0.415190in}{1.591253in}}%
\pgfpathlineto{\pgfqpoint{0.404299in}{1.600395in}}%
\pgfpathlineto{\pgfqpoint{0.397064in}{1.602137in}}%
\pgfpathlineto{\pgfqpoint{0.389344in}{1.598489in}}%
\pgfpathlineto{\pgfqpoint{0.382499in}{1.599727in}}%
\pgfpathlineto{\pgfqpoint{0.386071in}{1.611691in}}%
\pgfpathlineto{\pgfqpoint{0.382988in}{1.619317in}}%
\pgfpathlineto{\pgfqpoint{0.372778in}{1.622728in}}%
\pgfpathclose%
\pgfusepath{fill}%
\end{pgfscope}%
\begin{pgfscope}%
\pgfpathrectangle{\pgfqpoint{0.100000in}{0.100000in}}{\pgfqpoint{3.007045in}{1.925000in}}%
\pgfusepath{clip}%
\pgfsetbuttcap%
\pgfsetmiterjoin%
\definecolor{currentfill}{rgb}{0.760477,0.852026,0.931657}%
\pgfsetfillcolor{currentfill}%
\pgfsetlinewidth{0.000000pt}%
\definecolor{currentstroke}{rgb}{0.000000,0.000000,0.000000}%
\pgfsetstrokecolor{currentstroke}%
\pgfsetstrokeopacity{0.000000}%
\pgfsetdash{}{0pt}%
\pgfpathmoveto{\pgfqpoint{2.605371in}{0.522544in}}%
\pgfpathlineto{\pgfqpoint{2.591696in}{0.520372in}}%
\pgfpathlineto{\pgfqpoint{2.592396in}{0.508481in}}%
\pgfpathlineto{\pgfqpoint{2.584158in}{0.521537in}}%
\pgfpathlineto{\pgfqpoint{2.581927in}{0.518998in}}%
\pgfpathlineto{\pgfqpoint{2.576314in}{0.520862in}}%
\pgfpathlineto{\pgfqpoint{2.570842in}{0.520187in}}%
\pgfpathlineto{\pgfqpoint{2.567244in}{0.529939in}}%
\pgfpathlineto{\pgfqpoint{2.557842in}{0.535907in}}%
\pgfpathlineto{\pgfqpoint{2.555178in}{0.540975in}}%
\pgfpathlineto{\pgfqpoint{2.550313in}{0.542005in}}%
\pgfpathlineto{\pgfqpoint{2.547864in}{0.546632in}}%
\pgfpathlineto{\pgfqpoint{2.541214in}{0.552000in}}%
\pgfpathlineto{\pgfqpoint{2.538085in}{0.559398in}}%
\pgfpathlineto{\pgfqpoint{2.533920in}{0.561053in}}%
\pgfpathlineto{\pgfqpoint{2.522149in}{0.556021in}}%
\pgfpathlineto{\pgfqpoint{2.520191in}{0.572396in}}%
\pgfpathlineto{\pgfqpoint{2.526180in}{0.575201in}}%
\pgfpathlineto{\pgfqpoint{2.532632in}{0.582139in}}%
\pgfpathlineto{\pgfqpoint{2.545862in}{0.584302in}}%
\pgfpathlineto{\pgfqpoint{2.549830in}{0.579481in}}%
\pgfpathlineto{\pgfqpoint{2.551269in}{0.570455in}}%
\pgfpathlineto{\pgfqpoint{2.562689in}{0.572385in}}%
\pgfpathlineto{\pgfqpoint{2.568608in}{0.576327in}}%
\pgfpathlineto{\pgfqpoint{2.581840in}{0.553861in}}%
\pgfpathlineto{\pgfqpoint{2.588561in}{0.543384in}}%
\pgfpathlineto{\pgfqpoint{2.605371in}{0.522544in}}%
\pgfpathclose%
\pgfusepath{fill}%
\end{pgfscope}%
\begin{pgfscope}%
\pgfpathrectangle{\pgfqpoint{0.100000in}{0.100000in}}{\pgfqpoint{3.007045in}{1.925000in}}%
\pgfusepath{clip}%
\pgfsetbuttcap%
\pgfsetmiterjoin%
\definecolor{currentfill}{rgb}{0.422745,0.684075,0.839892}%
\pgfsetfillcolor{currentfill}%
\pgfsetlinewidth{0.000000pt}%
\definecolor{currentstroke}{rgb}{0.000000,0.000000,0.000000}%
\pgfsetstrokecolor{currentstroke}%
\pgfsetstrokeopacity{0.000000}%
\pgfsetdash{}{0pt}%
\pgfpathmoveto{\pgfqpoint{0.778392in}{0.924882in}}%
\pgfpathlineto{\pgfqpoint{0.767196in}{0.926377in}}%
\pgfpathlineto{\pgfqpoint{0.762651in}{0.922152in}}%
\pgfpathlineto{\pgfqpoint{0.740331in}{0.928918in}}%
\pgfpathlineto{\pgfqpoint{0.735192in}{0.932928in}}%
\pgfpathlineto{\pgfqpoint{0.726621in}{0.944753in}}%
\pgfpathlineto{\pgfqpoint{0.723368in}{0.955271in}}%
\pgfpathlineto{\pgfqpoint{0.723182in}{0.963163in}}%
\pgfpathlineto{\pgfqpoint{0.719619in}{0.967635in}}%
\pgfpathlineto{\pgfqpoint{0.716895in}{0.976496in}}%
\pgfpathlineto{\pgfqpoint{0.718691in}{0.983740in}}%
\pgfpathlineto{\pgfqpoint{0.677176in}{1.047687in}}%
\pgfpathlineto{\pgfqpoint{0.667024in}{1.063268in}}%
\pgfpathlineto{\pgfqpoint{0.616671in}{1.141013in}}%
\pgfpathlineto{\pgfqpoint{0.590725in}{1.181138in}}%
\pgfpathlineto{\pgfqpoint{0.567833in}{1.216455in}}%
\pgfpathlineto{\pgfqpoint{0.571645in}{1.215280in}}%
\pgfpathlineto{\pgfqpoint{0.614328in}{1.243704in}}%
\pgfpathlineto{\pgfqpoint{0.596506in}{1.278590in}}%
\pgfpathlineto{\pgfqpoint{0.597849in}{1.283744in}}%
\pgfpathlineto{\pgfqpoint{0.615468in}{1.284198in}}%
\pgfpathlineto{\pgfqpoint{0.620230in}{1.284320in}}%
\pgfpathlineto{\pgfqpoint{0.643335in}{1.283223in}}%
\pgfpathlineto{\pgfqpoint{0.679545in}{1.274162in}}%
\pgfpathlineto{\pgfqpoint{0.714021in}{1.265975in}}%
\pgfpathlineto{\pgfqpoint{0.752166in}{1.224502in}}%
\pgfpathlineto{\pgfqpoint{0.799942in}{1.214078in}}%
\pgfpathlineto{\pgfqpoint{0.798495in}{1.207353in}}%
\pgfpathlineto{\pgfqpoint{0.788118in}{1.158139in}}%
\pgfpathlineto{\pgfqpoint{0.776912in}{1.106019in}}%
\pgfpathlineto{\pgfqpoint{0.836260in}{1.093783in}}%
\pgfpathlineto{\pgfqpoint{0.854802in}{1.090140in}}%
\pgfpathlineto{\pgfqpoint{0.845589in}{1.070038in}}%
\pgfpathlineto{\pgfqpoint{0.842643in}{1.059879in}}%
\pgfpathlineto{\pgfqpoint{0.842909in}{1.052579in}}%
\pgfpathlineto{\pgfqpoint{0.839109in}{1.048947in}}%
\pgfpathlineto{\pgfqpoint{0.832280in}{1.046410in}}%
\pgfpathlineto{\pgfqpoint{0.818051in}{1.044246in}}%
\pgfpathlineto{\pgfqpoint{0.809522in}{1.041947in}}%
\pgfpathlineto{\pgfqpoint{0.807519in}{1.038421in}}%
\pgfpathlineto{\pgfqpoint{0.803014in}{1.039936in}}%
\pgfpathlineto{\pgfqpoint{0.800098in}{1.036517in}}%
\pgfpathlineto{\pgfqpoint{0.800968in}{1.030983in}}%
\pgfpathlineto{\pgfqpoint{0.798014in}{1.020963in}}%
\pgfpathlineto{\pgfqpoint{0.794371in}{1.003188in}}%
\pgfpathlineto{\pgfqpoint{0.778392in}{0.924882in}}%
\pgfpathclose%
\pgfusepath{fill}%
\end{pgfscope}%
\begin{pgfscope}%
\pgfpathrectangle{\pgfqpoint{0.100000in}{0.100000in}}{\pgfqpoint{3.007045in}{1.925000in}}%
\pgfusepath{clip}%
\pgfsetbuttcap%
\pgfsetmiterjoin%
\definecolor{currentfill}{rgb}{0.381776,0.656517,0.824452}%
\pgfsetfillcolor{currentfill}%
\pgfsetlinewidth{0.000000pt}%
\definecolor{currentstroke}{rgb}{0.000000,0.000000,0.000000}%
\pgfsetstrokecolor{currentstroke}%
\pgfsetstrokeopacity{0.000000}%
\pgfsetdash{}{0pt}%
\pgfpathmoveto{\pgfqpoint{2.480209in}{1.264505in}}%
\pgfpathlineto{\pgfqpoint{2.463037in}{1.262554in}}%
\pgfpathlineto{\pgfqpoint{2.457021in}{1.264583in}}%
\pgfpathlineto{\pgfqpoint{2.440547in}{1.262665in}}%
\pgfpathlineto{\pgfqpoint{2.441070in}{1.257982in}}%
\pgfpathlineto{\pgfqpoint{2.415184in}{1.255293in}}%
\pgfpathlineto{\pgfqpoint{2.414093in}{1.259940in}}%
\pgfpathlineto{\pgfqpoint{2.411854in}{1.281854in}}%
\pgfpathlineto{\pgfqpoint{2.409849in}{1.281241in}}%
\pgfpathlineto{\pgfqpoint{2.408904in}{1.287720in}}%
\pgfpathlineto{\pgfqpoint{2.413572in}{1.288426in}}%
\pgfpathlineto{\pgfqpoint{2.410399in}{1.309530in}}%
\pgfpathlineto{\pgfqpoint{2.434169in}{1.312710in}}%
\pgfpathlineto{\pgfqpoint{2.434855in}{1.307908in}}%
\pgfpathlineto{\pgfqpoint{2.446146in}{1.309080in}}%
\pgfpathlineto{\pgfqpoint{2.445530in}{1.314439in}}%
\pgfpathlineto{\pgfqpoint{2.461738in}{1.316890in}}%
\pgfpathlineto{\pgfqpoint{2.464330in}{1.299934in}}%
\pgfpathlineto{\pgfqpoint{2.472733in}{1.301176in}}%
\pgfpathlineto{\pgfqpoint{2.473535in}{1.295735in}}%
\pgfpathlineto{\pgfqpoint{2.476757in}{1.293301in}}%
\pgfpathlineto{\pgfqpoint{2.474283in}{1.290040in}}%
\pgfpathlineto{\pgfqpoint{2.474602in}{1.281312in}}%
\pgfpathlineto{\pgfqpoint{2.478385in}{1.281786in}}%
\pgfpathlineto{\pgfqpoint{2.480209in}{1.264505in}}%
\pgfpathclose%
\pgfusepath{fill}%
\end{pgfscope}%
\begin{pgfscope}%
\pgfpathrectangle{\pgfqpoint{0.100000in}{0.100000in}}{\pgfqpoint{3.007045in}{1.925000in}}%
\pgfusepath{clip}%
\pgfsetbuttcap%
\pgfsetmiterjoin%
\definecolor{currentfill}{rgb}{0.676817,0.816471,0.902376}%
\pgfsetfillcolor{currentfill}%
\pgfsetlinewidth{0.000000pt}%
\definecolor{currentstroke}{rgb}{0.000000,0.000000,0.000000}%
\pgfsetstrokecolor{currentstroke}%
\pgfsetstrokeopacity{0.000000}%
\pgfsetdash{}{0pt}%
\pgfpathmoveto{\pgfqpoint{2.514084in}{1.019211in}}%
\pgfpathlineto{\pgfqpoint{2.520017in}{1.015115in}}%
\pgfpathlineto{\pgfqpoint{2.522081in}{0.997367in}}%
\pgfpathlineto{\pgfqpoint{2.510978in}{0.993811in}}%
\pgfpathlineto{\pgfqpoint{2.506001in}{0.993623in}}%
\pgfpathlineto{\pgfqpoint{2.499152in}{0.989884in}}%
\pgfpathlineto{\pgfqpoint{2.494094in}{0.994158in}}%
\pgfpathlineto{\pgfqpoint{2.486521in}{0.996045in}}%
\pgfpathlineto{\pgfqpoint{2.490105in}{1.002416in}}%
\pgfpathlineto{\pgfqpoint{2.479391in}{1.010461in}}%
\pgfpathlineto{\pgfqpoint{2.473634in}{1.012683in}}%
\pgfpathlineto{\pgfqpoint{2.475252in}{1.017742in}}%
\pgfpathlineto{\pgfqpoint{2.475027in}{1.025720in}}%
\pgfpathlineto{\pgfqpoint{2.492133in}{1.027383in}}%
\pgfpathlineto{\pgfqpoint{2.492311in}{1.025615in}}%
\pgfpathlineto{\pgfqpoint{2.500787in}{1.015480in}}%
\pgfpathlineto{\pgfqpoint{2.506602in}{1.020237in}}%
\pgfpathlineto{\pgfqpoint{2.514084in}{1.019211in}}%
\pgfpathclose%
\pgfusepath{fill}%
\end{pgfscope}%
\begin{pgfscope}%
\pgfpathrectangle{\pgfqpoint{0.100000in}{0.100000in}}{\pgfqpoint{3.007045in}{1.925000in}}%
\pgfusepath{clip}%
\pgfsetbuttcap%
\pgfsetmiterjoin%
\definecolor{currentfill}{rgb}{0.376732,0.653072,0.822484}%
\pgfsetfillcolor{currentfill}%
\pgfsetlinewidth{0.000000pt}%
\definecolor{currentstroke}{rgb}{0.000000,0.000000,0.000000}%
\pgfsetstrokecolor{currentstroke}%
\pgfsetstrokeopacity{0.000000}%
\pgfsetdash{}{0pt}%
\pgfpathmoveto{\pgfqpoint{1.396518in}{1.346955in}}%
\pgfpathlineto{\pgfqpoint{1.426380in}{1.344786in}}%
\pgfpathlineto{\pgfqpoint{1.424949in}{1.327230in}}%
\pgfpathlineto{\pgfqpoint{1.428992in}{1.326930in}}%
\pgfpathlineto{\pgfqpoint{1.427276in}{1.304079in}}%
\pgfpathlineto{\pgfqpoint{1.423796in}{1.304378in}}%
\pgfpathlineto{\pgfqpoint{1.422895in}{1.292969in}}%
\pgfpathlineto{\pgfqpoint{1.395479in}{1.295197in}}%
\pgfpathlineto{\pgfqpoint{1.396484in}{1.306583in}}%
\pgfpathlineto{\pgfqpoint{1.395381in}{1.309582in}}%
\pgfpathlineto{\pgfqpoint{1.396817in}{1.327688in}}%
\pgfpathlineto{\pgfqpoint{1.395435in}{1.334971in}}%
\pgfpathlineto{\pgfqpoint{1.396518in}{1.346955in}}%
\pgfpathclose%
\pgfusepath{fill}%
\end{pgfscope}%
\begin{pgfscope}%
\pgfpathrectangle{\pgfqpoint{0.100000in}{0.100000in}}{\pgfqpoint{3.007045in}{1.925000in}}%
\pgfusepath{clip}%
\pgfsetbuttcap%
\pgfsetmiterjoin%
\definecolor{currentfill}{rgb}{0.296025,0.597955,0.790988}%
\pgfsetfillcolor{currentfill}%
\pgfsetlinewidth{0.000000pt}%
\definecolor{currentstroke}{rgb}{0.000000,0.000000,0.000000}%
\pgfsetstrokecolor{currentstroke}%
\pgfsetstrokeopacity{0.000000}%
\pgfsetdash{}{0pt}%
\pgfpathmoveto{\pgfqpoint{2.285799in}{1.004151in}}%
\pgfpathlineto{\pgfqpoint{2.285270in}{0.997505in}}%
\pgfpathlineto{\pgfqpoint{2.278118in}{0.995018in}}%
\pgfpathlineto{\pgfqpoint{2.275809in}{0.988285in}}%
\pgfpathlineto{\pgfqpoint{2.271083in}{0.992318in}}%
\pgfpathlineto{\pgfqpoint{2.263992in}{0.994959in}}%
\pgfpathlineto{\pgfqpoint{2.258356in}{0.992865in}}%
\pgfpathlineto{\pgfqpoint{2.258514in}{1.000440in}}%
\pgfpathlineto{\pgfqpoint{2.248936in}{1.000483in}}%
\pgfpathlineto{\pgfqpoint{2.247870in}{1.007837in}}%
\pgfpathlineto{\pgfqpoint{2.236965in}{1.019520in}}%
\pgfpathlineto{\pgfqpoint{2.241044in}{1.027825in}}%
\pgfpathlineto{\pgfqpoint{2.241309in}{1.038598in}}%
\pgfpathlineto{\pgfqpoint{2.234875in}{1.045920in}}%
\pgfpathlineto{\pgfqpoint{2.238225in}{1.047271in}}%
\pgfpathlineto{\pgfqpoint{2.239874in}{1.053973in}}%
\pgfpathlineto{\pgfqpoint{2.260372in}{1.054200in}}%
\pgfpathlineto{\pgfqpoint{2.259090in}{1.046090in}}%
\pgfpathlineto{\pgfqpoint{2.260460in}{1.038312in}}%
\pgfpathlineto{\pgfqpoint{2.263654in}{1.034861in}}%
\pgfpathlineto{\pgfqpoint{2.269307in}{1.034406in}}%
\pgfpathlineto{\pgfqpoint{2.271241in}{1.027836in}}%
\pgfpathlineto{\pgfqpoint{2.274407in}{1.023601in}}%
\pgfpathlineto{\pgfqpoint{2.286020in}{1.023971in}}%
\pgfpathlineto{\pgfqpoint{2.287383in}{1.019436in}}%
\pgfpathlineto{\pgfqpoint{2.285799in}{1.004151in}}%
\pgfpathclose%
\pgfusepath{fill}%
\end{pgfscope}%
\begin{pgfscope}%
\pgfpathrectangle{\pgfqpoint{0.100000in}{0.100000in}}{\pgfqpoint{3.007045in}{1.925000in}}%
\pgfusepath{clip}%
\pgfsetbuttcap%
\pgfsetmiterjoin%
\definecolor{currentfill}{rgb}{0.376732,0.653072,0.822484}%
\pgfsetfillcolor{currentfill}%
\pgfsetlinewidth{0.000000pt}%
\definecolor{currentstroke}{rgb}{0.000000,0.000000,0.000000}%
\pgfsetstrokecolor{currentstroke}%
\pgfsetstrokeopacity{0.000000}%
\pgfsetdash{}{0pt}%
\pgfpathmoveto{\pgfqpoint{2.688736in}{0.994006in}}%
\pgfpathlineto{\pgfqpoint{2.675874in}{0.992417in}}%
\pgfpathlineto{\pgfqpoint{2.672444in}{0.992955in}}%
\pgfpathlineto{\pgfqpoint{2.668059in}{1.000440in}}%
\pgfpathlineto{\pgfqpoint{2.663029in}{1.005596in}}%
\pgfpathlineto{\pgfqpoint{2.671365in}{1.033005in}}%
\pgfpathlineto{\pgfqpoint{2.679615in}{1.030971in}}%
\pgfpathlineto{\pgfqpoint{2.689119in}{1.033149in}}%
\pgfpathlineto{\pgfqpoint{2.697499in}{1.030502in}}%
\pgfpathlineto{\pgfqpoint{2.698517in}{1.027123in}}%
\pgfpathlineto{\pgfqpoint{2.705428in}{1.026285in}}%
\pgfpathlineto{\pgfqpoint{2.709850in}{1.020460in}}%
\pgfpathlineto{\pgfqpoint{2.710520in}{1.014800in}}%
\pgfpathlineto{\pgfqpoint{2.708269in}{1.015278in}}%
\pgfpathlineto{\pgfqpoint{2.688736in}{0.994006in}}%
\pgfpathclose%
\pgfusepath{fill}%
\end{pgfscope}%
\begin{pgfscope}%
\pgfpathrectangle{\pgfqpoint{0.100000in}{0.100000in}}{\pgfqpoint{3.007045in}{1.925000in}}%
\pgfusepath{clip}%
\pgfsetbuttcap%
\pgfsetmiterjoin%
\definecolor{currentfill}{rgb}{0.142607,0.456332,0.716601}%
\pgfsetfillcolor{currentfill}%
\pgfsetlinewidth{0.000000pt}%
\definecolor{currentstroke}{rgb}{0.000000,0.000000,0.000000}%
\pgfsetstrokecolor{currentstroke}%
\pgfsetstrokeopacity{0.000000}%
\pgfsetdash{}{0pt}%
\pgfpathmoveto{\pgfqpoint{0.892728in}{1.460530in}}%
\pgfpathlineto{\pgfqpoint{0.895172in}{1.457614in}}%
\pgfpathlineto{\pgfqpoint{0.903388in}{1.455868in}}%
\pgfpathlineto{\pgfqpoint{0.900100in}{1.438993in}}%
\pgfpathlineto{\pgfqpoint{0.895826in}{1.417699in}}%
\pgfpathlineto{\pgfqpoint{0.853680in}{1.425681in}}%
\pgfpathlineto{\pgfqpoint{0.845804in}{1.427734in}}%
\pgfpathlineto{\pgfqpoint{0.834301in}{1.430229in}}%
\pgfpathlineto{\pgfqpoint{0.840009in}{1.457550in}}%
\pgfpathlineto{\pgfqpoint{0.850305in}{1.455334in}}%
\pgfpathlineto{\pgfqpoint{0.851744in}{1.461596in}}%
\pgfpathlineto{\pgfqpoint{0.858463in}{1.461570in}}%
\pgfpathlineto{\pgfqpoint{0.861547in}{1.476417in}}%
\pgfpathlineto{\pgfqpoint{0.869518in}{1.474687in}}%
\pgfpathlineto{\pgfqpoint{0.873074in}{1.479709in}}%
\pgfpathlineto{\pgfqpoint{0.877767in}{1.502355in}}%
\pgfpathlineto{\pgfqpoint{0.891952in}{1.499400in}}%
\pgfpathlineto{\pgfqpoint{0.887326in}{1.476723in}}%
\pgfpathlineto{\pgfqpoint{0.884523in}{1.477298in}}%
\pgfpathlineto{\pgfqpoint{0.882125in}{1.465664in}}%
\pgfpathlineto{\pgfqpoint{0.887062in}{1.465893in}}%
\pgfpathlineto{\pgfqpoint{0.892728in}{1.460530in}}%
\pgfpathclose%
\pgfusepath{fill}%
\end{pgfscope}%
\begin{pgfscope}%
\pgfpathrectangle{\pgfqpoint{0.100000in}{0.100000in}}{\pgfqpoint{3.007045in}{1.925000in}}%
\pgfusepath{clip}%
\pgfsetbuttcap%
\pgfsetmiterjoin%
\definecolor{currentfill}{rgb}{0.498039,0.725413,0.856132}%
\pgfsetfillcolor{currentfill}%
\pgfsetlinewidth{0.000000pt}%
\definecolor{currentstroke}{rgb}{0.000000,0.000000,0.000000}%
\pgfsetstrokecolor{currentstroke}%
\pgfsetstrokeopacity{0.000000}%
\pgfsetdash{}{0pt}%
\pgfpathmoveto{\pgfqpoint{2.286587in}{0.869005in}}%
\pgfpathlineto{\pgfqpoint{2.279674in}{0.894582in}}%
\pgfpathlineto{\pgfqpoint{2.265651in}{0.893385in}}%
\pgfpathlineto{\pgfqpoint{2.264140in}{0.908633in}}%
\pgfpathlineto{\pgfqpoint{2.261069in}{0.912869in}}%
\pgfpathlineto{\pgfqpoint{2.262417in}{0.916982in}}%
\pgfpathlineto{\pgfqpoint{2.261477in}{0.928516in}}%
\pgfpathlineto{\pgfqpoint{2.264733in}{0.928140in}}%
\pgfpathlineto{\pgfqpoint{2.278483in}{0.930848in}}%
\pgfpathlineto{\pgfqpoint{2.285385in}{0.933858in}}%
\pgfpathlineto{\pgfqpoint{2.287005in}{0.929042in}}%
\pgfpathlineto{\pgfqpoint{2.297460in}{0.921002in}}%
\pgfpathlineto{\pgfqpoint{2.299822in}{0.927995in}}%
\pgfpathlineto{\pgfqpoint{2.308274in}{0.925835in}}%
\pgfpathlineto{\pgfqpoint{2.313145in}{0.917931in}}%
\pgfpathlineto{\pgfqpoint{2.309785in}{0.907283in}}%
\pgfpathlineto{\pgfqpoint{2.313442in}{0.898402in}}%
\pgfpathlineto{\pgfqpoint{2.310101in}{0.886859in}}%
\pgfpathlineto{\pgfqpoint{2.306322in}{0.886459in}}%
\pgfpathlineto{\pgfqpoint{2.305031in}{0.879714in}}%
\pgfpathlineto{\pgfqpoint{2.311402in}{0.880390in}}%
\pgfpathlineto{\pgfqpoint{2.311838in}{0.871613in}}%
\pgfpathlineto{\pgfqpoint{2.286587in}{0.869005in}}%
\pgfpathclose%
\pgfusepath{fill}%
\end{pgfscope}%
\begin{pgfscope}%
\pgfpathrectangle{\pgfqpoint{0.100000in}{0.100000in}}{\pgfqpoint{3.007045in}{1.925000in}}%
\pgfusepath{clip}%
\pgfsetbuttcap%
\pgfsetmiterjoin%
\definecolor{currentfill}{rgb}{0.361599,0.642737,0.816578}%
\pgfsetfillcolor{currentfill}%
\pgfsetlinewidth{0.000000pt}%
\definecolor{currentstroke}{rgb}{0.000000,0.000000,0.000000}%
\pgfsetstrokecolor{currentstroke}%
\pgfsetstrokeopacity{0.000000}%
\pgfsetdash{}{0pt}%
\pgfpathmoveto{\pgfqpoint{2.239459in}{1.266008in}}%
\pgfpathlineto{\pgfqpoint{2.228217in}{1.264785in}}%
\pgfpathlineto{\pgfqpoint{2.224496in}{1.264365in}}%
\pgfpathlineto{\pgfqpoint{2.221882in}{1.287153in}}%
\pgfpathlineto{\pgfqpoint{2.210917in}{1.285878in}}%
\pgfpathlineto{\pgfqpoint{2.211433in}{1.280191in}}%
\pgfpathlineto{\pgfqpoint{2.196588in}{1.278782in}}%
\pgfpathlineto{\pgfqpoint{2.195043in}{1.295965in}}%
\pgfpathlineto{\pgfqpoint{2.215428in}{1.298031in}}%
\pgfpathlineto{\pgfqpoint{2.213486in}{1.315239in}}%
\pgfpathlineto{\pgfqpoint{2.233410in}{1.317320in}}%
\pgfpathlineto{\pgfqpoint{2.234453in}{1.308012in}}%
\pgfpathlineto{\pgfqpoint{2.233530in}{1.300204in}}%
\pgfpathlineto{\pgfqpoint{2.234886in}{1.288646in}}%
\pgfpathlineto{\pgfqpoint{2.236831in}{1.288875in}}%
\pgfpathlineto{\pgfqpoint{2.239459in}{1.266008in}}%
\pgfpathclose%
\pgfusepath{fill}%
\end{pgfscope}%
\begin{pgfscope}%
\pgfpathrectangle{\pgfqpoint{0.100000in}{0.100000in}}{\pgfqpoint{3.007045in}{1.925000in}}%
\pgfusepath{clip}%
\pgfsetbuttcap%
\pgfsetmiterjoin%
\definecolor{currentfill}{rgb}{0.637447,0.799739,0.888597}%
\pgfsetfillcolor{currentfill}%
\pgfsetlinewidth{0.000000pt}%
\definecolor{currentstroke}{rgb}{0.000000,0.000000,0.000000}%
\pgfsetstrokecolor{currentstroke}%
\pgfsetstrokeopacity{0.000000}%
\pgfsetdash{}{0pt}%
\pgfpathmoveto{\pgfqpoint{0.374195in}{1.289386in}}%
\pgfpathlineto{\pgfqpoint{0.369150in}{1.288189in}}%
\pgfpathlineto{\pgfqpoint{0.366648in}{1.294372in}}%
\pgfpathlineto{\pgfqpoint{0.358679in}{1.298249in}}%
\pgfpathlineto{\pgfqpoint{0.352980in}{1.307630in}}%
\pgfpathlineto{\pgfqpoint{0.347482in}{1.307924in}}%
\pgfpathlineto{\pgfqpoint{0.344373in}{1.319665in}}%
\pgfpathlineto{\pgfqpoint{0.339643in}{1.325177in}}%
\pgfpathlineto{\pgfqpoint{0.335632in}{1.333256in}}%
\pgfpathlineto{\pgfqpoint{0.333488in}{1.341560in}}%
\pgfpathlineto{\pgfqpoint{0.327012in}{1.352333in}}%
\pgfpathlineto{\pgfqpoint{0.324629in}{1.358271in}}%
\pgfpathlineto{\pgfqpoint{0.328917in}{1.365788in}}%
\pgfpathlineto{\pgfqpoint{0.328154in}{1.376358in}}%
\pgfpathlineto{\pgfqpoint{0.328676in}{1.386959in}}%
\pgfpathlineto{\pgfqpoint{0.331174in}{1.392985in}}%
\pgfpathlineto{\pgfqpoint{0.335747in}{1.398988in}}%
\pgfpathlineto{\pgfqpoint{0.337136in}{1.407814in}}%
\pgfpathlineto{\pgfqpoint{0.337272in}{1.417649in}}%
\pgfpathlineto{\pgfqpoint{0.332528in}{1.430977in}}%
\pgfpathlineto{\pgfqpoint{0.355543in}{1.423387in}}%
\pgfpathlineto{\pgfqpoint{0.355032in}{1.421828in}}%
\pgfpathlineto{\pgfqpoint{0.384397in}{1.412475in}}%
\pgfpathlineto{\pgfqpoint{0.380635in}{1.401231in}}%
\pgfpathlineto{\pgfqpoint{0.381027in}{1.394938in}}%
\pgfpathlineto{\pgfqpoint{0.378819in}{1.386738in}}%
\pgfpathlineto{\pgfqpoint{0.386077in}{1.384484in}}%
\pgfpathlineto{\pgfqpoint{0.383653in}{1.376276in}}%
\pgfpathlineto{\pgfqpoint{0.379169in}{1.366194in}}%
\pgfpathlineto{\pgfqpoint{0.383879in}{1.360469in}}%
\pgfpathlineto{\pgfqpoint{0.389453in}{1.357782in}}%
\pgfpathlineto{\pgfqpoint{0.387533in}{1.347636in}}%
\pgfpathlineto{\pgfqpoint{0.390745in}{1.343728in}}%
\pgfpathlineto{\pgfqpoint{0.392426in}{1.337198in}}%
\pgfpathlineto{\pgfqpoint{0.388567in}{1.334258in}}%
\pgfpathlineto{\pgfqpoint{0.388700in}{1.330958in}}%
\pgfpathlineto{\pgfqpoint{0.382079in}{1.325306in}}%
\pgfpathlineto{\pgfqpoint{0.373315in}{1.325445in}}%
\pgfpathlineto{\pgfqpoint{0.371022in}{1.321407in}}%
\pgfpathlineto{\pgfqpoint{0.376145in}{1.309067in}}%
\pgfpathlineto{\pgfqpoint{0.377641in}{1.299051in}}%
\pgfpathlineto{\pgfqpoint{0.377156in}{1.290848in}}%
\pgfpathlineto{\pgfqpoint{0.374195in}{1.289386in}}%
\pgfpathclose%
\pgfusepath{fill}%
\end{pgfscope}%
\begin{pgfscope}%
\pgfpathrectangle{\pgfqpoint{0.100000in}{0.100000in}}{\pgfqpoint{3.007045in}{1.925000in}}%
\pgfusepath{clip}%
\pgfsetbuttcap%
\pgfsetmiterjoin%
\definecolor{currentfill}{rgb}{0.579608,0.770196,0.873725}%
\pgfsetfillcolor{currentfill}%
\pgfsetlinewidth{0.000000pt}%
\definecolor{currentstroke}{rgb}{0.000000,0.000000,0.000000}%
\pgfsetstrokecolor{currentstroke}%
\pgfsetstrokeopacity{0.000000}%
\pgfsetdash{}{0pt}%
\pgfpathmoveto{\pgfqpoint{2.458232in}{1.094688in}}%
\pgfpathlineto{\pgfqpoint{2.446820in}{1.093542in}}%
\pgfpathlineto{\pgfqpoint{2.441207in}{1.098026in}}%
\pgfpathlineto{\pgfqpoint{2.440031in}{1.105220in}}%
\pgfpathlineto{\pgfqpoint{2.435308in}{1.112304in}}%
\pgfpathlineto{\pgfqpoint{2.429643in}{1.109432in}}%
\pgfpathlineto{\pgfqpoint{2.433994in}{1.115719in}}%
\pgfpathlineto{\pgfqpoint{2.428858in}{1.122661in}}%
\pgfpathlineto{\pgfqpoint{2.428704in}{1.126732in}}%
\pgfpathlineto{\pgfqpoint{2.431774in}{1.134957in}}%
\pgfpathlineto{\pgfqpoint{2.438948in}{1.139413in}}%
\pgfpathlineto{\pgfqpoint{2.437602in}{1.145887in}}%
\pgfpathlineto{\pgfqpoint{2.435873in}{1.154519in}}%
\pgfpathlineto{\pgfqpoint{2.440950in}{1.158535in}}%
\pgfpathlineto{\pgfqpoint{2.450037in}{1.161421in}}%
\pgfpathlineto{\pgfqpoint{2.458391in}{1.154238in}}%
\pgfpathlineto{\pgfqpoint{2.463450in}{1.158871in}}%
\pgfpathlineto{\pgfqpoint{2.466955in}{1.155210in}}%
\pgfpathlineto{\pgfqpoint{2.480993in}{1.155854in}}%
\pgfpathlineto{\pgfqpoint{2.487851in}{1.166255in}}%
\pgfpathlineto{\pgfqpoint{2.494103in}{1.161327in}}%
\pgfpathlineto{\pgfqpoint{2.497293in}{1.156805in}}%
\pgfpathlineto{\pgfqpoint{2.496539in}{1.150737in}}%
\pgfpathlineto{\pgfqpoint{2.481701in}{1.138352in}}%
\pgfpathlineto{\pgfqpoint{2.477515in}{1.130920in}}%
\pgfpathlineto{\pgfqpoint{2.476974in}{1.118028in}}%
\pgfpathlineto{\pgfqpoint{2.470294in}{1.118309in}}%
\pgfpathlineto{\pgfqpoint{2.467580in}{1.113861in}}%
\pgfpathlineto{\pgfqpoint{2.471855in}{1.105407in}}%
\pgfpathlineto{\pgfqpoint{2.459342in}{1.100923in}}%
\pgfpathlineto{\pgfqpoint{2.463061in}{1.098120in}}%
\pgfpathlineto{\pgfqpoint{2.458232in}{1.094688in}}%
\pgfpathclose%
\pgfusepath{fill}%
\end{pgfscope}%
\begin{pgfscope}%
\pgfpathrectangle{\pgfqpoint{0.100000in}{0.100000in}}{\pgfqpoint{3.007045in}{1.925000in}}%
\pgfusepath{clip}%
\pgfsetbuttcap%
\pgfsetmiterjoin%
\definecolor{currentfill}{rgb}{0.460392,0.704744,0.848012}%
\pgfsetfillcolor{currentfill}%
\pgfsetlinewidth{0.000000pt}%
\definecolor{currentstroke}{rgb}{0.000000,0.000000,0.000000}%
\pgfsetstrokecolor{currentstroke}%
\pgfsetstrokeopacity{0.000000}%
\pgfsetdash{}{0pt}%
\pgfpathmoveto{\pgfqpoint{1.298841in}{1.455523in}}%
\pgfpathlineto{\pgfqpoint{1.293878in}{1.409890in}}%
\pgfpathlineto{\pgfqpoint{1.292810in}{1.397021in}}%
\pgfpathlineto{\pgfqpoint{1.273764in}{1.398937in}}%
\pgfpathlineto{\pgfqpoint{1.272590in}{1.387487in}}%
\pgfpathlineto{\pgfqpoint{1.270231in}{1.387736in}}%
\pgfpathlineto{\pgfqpoint{1.267095in}{1.378734in}}%
\pgfpathlineto{\pgfqpoint{1.255653in}{1.381012in}}%
\pgfpathlineto{\pgfqpoint{1.256175in}{1.385798in}}%
\pgfpathlineto{\pgfqpoint{1.259419in}{1.388949in}}%
\pgfpathlineto{\pgfqpoint{1.234424in}{1.392114in}}%
\pgfpathlineto{\pgfqpoint{1.207752in}{1.395280in}}%
\pgfpathlineto{\pgfqpoint{1.164443in}{1.401325in}}%
\pgfpathlineto{\pgfqpoint{1.163522in}{1.401453in}}%
\pgfpathlineto{\pgfqpoint{1.166601in}{1.424158in}}%
\pgfpathlineto{\pgfqpoint{1.168621in}{1.423872in}}%
\pgfpathlineto{\pgfqpoint{1.171716in}{1.446544in}}%
\pgfpathlineto{\pgfqpoint{1.174139in}{1.469186in}}%
\pgfpathlineto{\pgfqpoint{1.173551in}{1.471182in}}%
\pgfpathlineto{\pgfqpoint{1.193681in}{1.468364in}}%
\pgfpathlineto{\pgfqpoint{1.218116in}{1.464571in}}%
\pgfpathlineto{\pgfqpoint{1.245609in}{1.461365in}}%
\pgfpathlineto{\pgfqpoint{1.298841in}{1.455523in}}%
\pgfpathclose%
\pgfusepath{fill}%
\end{pgfscope}%
\begin{pgfscope}%
\pgfpathrectangle{\pgfqpoint{0.100000in}{0.100000in}}{\pgfqpoint{3.007045in}{1.925000in}}%
\pgfusepath{clip}%
\pgfsetbuttcap%
\pgfsetmiterjoin%
\definecolor{currentfill}{rgb}{0.371688,0.649627,0.820515}%
\pgfsetfillcolor{currentfill}%
\pgfsetlinewidth{0.000000pt}%
\definecolor{currentstroke}{rgb}{0.000000,0.000000,0.000000}%
\pgfsetstrokecolor{currentstroke}%
\pgfsetstrokeopacity{0.000000}%
\pgfsetdash{}{0pt}%
\pgfpathmoveto{\pgfqpoint{1.756771in}{1.198935in}}%
\pgfpathlineto{\pgfqpoint{1.755176in}{1.198941in}}%
\pgfpathlineto{\pgfqpoint{1.709799in}{1.199381in}}%
\pgfpathlineto{\pgfqpoint{1.698475in}{1.199567in}}%
\pgfpathlineto{\pgfqpoint{1.699296in}{1.251149in}}%
\pgfpathlineto{\pgfqpoint{1.730711in}{1.250739in}}%
\pgfpathlineto{\pgfqpoint{1.727958in}{1.247635in}}%
\pgfpathlineto{\pgfqpoint{1.732568in}{1.242722in}}%
\pgfpathlineto{\pgfqpoint{1.734017in}{1.237625in}}%
\pgfpathlineto{\pgfqpoint{1.741287in}{1.219682in}}%
\pgfpathlineto{\pgfqpoint{1.747673in}{1.215377in}}%
\pgfpathlineto{\pgfqpoint{1.748211in}{1.211203in}}%
\pgfpathlineto{\pgfqpoint{1.751959in}{1.207454in}}%
\pgfpathlineto{\pgfqpoint{1.751069in}{1.202842in}}%
\pgfpathlineto{\pgfqpoint{1.756771in}{1.198935in}}%
\pgfpathclose%
\pgfusepath{fill}%
\end{pgfscope}%
\begin{pgfscope}%
\pgfpathrectangle{\pgfqpoint{0.100000in}{0.100000in}}{\pgfqpoint{3.007045in}{1.925000in}}%
\pgfusepath{clip}%
\pgfsetbuttcap%
\pgfsetmiterjoin%
\definecolor{currentfill}{rgb}{0.087120,0.389004,0.667512}%
\pgfsetfillcolor{currentfill}%
\pgfsetlinewidth{0.000000pt}%
\definecolor{currentstroke}{rgb}{0.000000,0.000000,0.000000}%
\pgfsetstrokecolor{currentstroke}%
\pgfsetstrokeopacity{0.000000}%
\pgfsetdash{}{0pt}%
\pgfpathmoveto{\pgfqpoint{1.370186in}{0.721452in}}%
\pgfpathlineto{\pgfqpoint{1.322173in}{0.725422in}}%
\pgfpathlineto{\pgfqpoint{1.319633in}{0.696676in}}%
\pgfpathlineto{\pgfqpoint{1.319120in}{0.690958in}}%
\pgfpathlineto{\pgfqpoint{1.282225in}{0.694354in}}%
\pgfpathlineto{\pgfqpoint{1.285541in}{0.728841in}}%
\pgfpathlineto{\pgfqpoint{1.280477in}{0.729323in}}%
\pgfpathlineto{\pgfqpoint{1.283283in}{0.758525in}}%
\pgfpathlineto{\pgfqpoint{1.285953in}{0.758267in}}%
\pgfpathlineto{\pgfqpoint{1.288691in}{0.786623in}}%
\pgfpathlineto{\pgfqpoint{1.291172in}{0.786411in}}%
\pgfpathlineto{\pgfqpoint{1.292309in}{0.797855in}}%
\pgfpathlineto{\pgfqpoint{1.303863in}{0.796794in}}%
\pgfpathlineto{\pgfqpoint{1.328991in}{0.794521in}}%
\pgfpathlineto{\pgfqpoint{1.327701in}{0.782533in}}%
\pgfpathlineto{\pgfqpoint{1.381713in}{0.778104in}}%
\pgfpathlineto{\pgfqpoint{1.379483in}{0.749728in}}%
\pgfpathlineto{\pgfqpoint{1.372150in}{0.750256in}}%
\pgfpathlineto{\pgfqpoint{1.370186in}{0.721452in}}%
\pgfpathclose%
\pgfusepath{fill}%
\end{pgfscope}%
\begin{pgfscope}%
\pgfpathrectangle{\pgfqpoint{0.100000in}{0.100000in}}{\pgfqpoint{3.007045in}{1.925000in}}%
\pgfusepath{clip}%
\pgfsetbuttcap%
\pgfsetmiterjoin%
\definecolor{currentfill}{rgb}{0.422745,0.684075,0.839892}%
\pgfsetfillcolor{currentfill}%
\pgfsetlinewidth{0.000000pt}%
\definecolor{currentstroke}{rgb}{0.000000,0.000000,0.000000}%
\pgfsetstrokecolor{currentstroke}%
\pgfsetstrokeopacity{0.000000}%
\pgfsetdash{}{0pt}%
\pgfpathmoveto{\pgfqpoint{2.139901in}{1.326614in}}%
\pgfpathlineto{\pgfqpoint{2.141986in}{1.299655in}}%
\pgfpathlineto{\pgfqpoint{2.117970in}{1.297350in}}%
\pgfpathlineto{\pgfqpoint{2.118514in}{1.291662in}}%
\pgfpathlineto{\pgfqpoint{2.107035in}{1.290511in}}%
\pgfpathlineto{\pgfqpoint{2.107101in}{1.284740in}}%
\pgfpathlineto{\pgfqpoint{2.090562in}{1.283138in}}%
\pgfpathlineto{\pgfqpoint{2.087331in}{1.317506in}}%
\pgfpathlineto{\pgfqpoint{2.085344in}{1.346129in}}%
\pgfpathlineto{\pgfqpoint{2.085588in}{1.351908in}}%
\pgfpathlineto{\pgfqpoint{2.079874in}{1.351505in}}%
\pgfpathlineto{\pgfqpoint{2.078237in}{1.373879in}}%
\pgfpathlineto{\pgfqpoint{2.122132in}{1.377026in}}%
\pgfpathlineto{\pgfqpoint{2.121650in}{1.363916in}}%
\pgfpathlineto{\pgfqpoint{2.123749in}{1.358371in}}%
\pgfpathlineto{\pgfqpoint{2.130184in}{1.350131in}}%
\pgfpathlineto{\pgfqpoint{2.134976in}{1.335269in}}%
\pgfpathlineto{\pgfqpoint{2.139901in}{1.326614in}}%
\pgfpathclose%
\pgfusepath{fill}%
\end{pgfscope}%
\begin{pgfscope}%
\pgfpathrectangle{\pgfqpoint{0.100000in}{0.100000in}}{\pgfqpoint{3.007045in}{1.925000in}}%
\pgfusepath{clip}%
\pgfsetbuttcap%
\pgfsetmiterjoin%
\definecolor{currentfill}{rgb}{0.371688,0.649627,0.820515}%
\pgfsetfillcolor{currentfill}%
\pgfsetlinewidth{0.000000pt}%
\definecolor{currentstroke}{rgb}{0.000000,0.000000,0.000000}%
\pgfsetstrokecolor{currentstroke}%
\pgfsetstrokeopacity{0.000000}%
\pgfsetdash{}{0pt}%
\pgfpathmoveto{\pgfqpoint{1.652076in}{1.171950in}}%
\pgfpathlineto{\pgfqpoint{1.651571in}{1.154782in}}%
\pgfpathlineto{\pgfqpoint{1.623055in}{1.155630in}}%
\pgfpathlineto{\pgfqpoint{1.623649in}{1.172778in}}%
\pgfpathlineto{\pgfqpoint{1.594392in}{1.173857in}}%
\pgfpathlineto{\pgfqpoint{1.595536in}{1.202468in}}%
\pgfpathlineto{\pgfqpoint{1.607157in}{1.202049in}}%
\pgfpathlineto{\pgfqpoint{1.652796in}{1.200602in}}%
\pgfpathlineto{\pgfqpoint{1.652076in}{1.171950in}}%
\pgfpathclose%
\pgfusepath{fill}%
\end{pgfscope}%
\begin{pgfscope}%
\pgfpathrectangle{\pgfqpoint{0.100000in}{0.100000in}}{\pgfqpoint{3.007045in}{1.925000in}}%
\pgfusepath{clip}%
\pgfsetbuttcap%
\pgfsetmiterjoin%
\definecolor{currentfill}{rgb}{0.074817,0.373256,0.655210}%
\pgfsetfillcolor{currentfill}%
\pgfsetlinewidth{0.000000pt}%
\definecolor{currentstroke}{rgb}{0.000000,0.000000,0.000000}%
\pgfsetstrokecolor{currentstroke}%
\pgfsetstrokeopacity{0.000000}%
\pgfsetdash{}{0pt}%
\pgfpathmoveto{\pgfqpoint{1.426380in}{1.344786in}}%
\pgfpathlineto{\pgfqpoint{1.396518in}{1.346955in}}%
\pgfpathlineto{\pgfqpoint{1.393848in}{1.352944in}}%
\pgfpathlineto{\pgfqpoint{1.395452in}{1.375827in}}%
\pgfpathlineto{\pgfqpoint{1.394290in}{1.375922in}}%
\pgfpathlineto{\pgfqpoint{1.396236in}{1.398628in}}%
\pgfpathlineto{\pgfqpoint{1.396545in}{1.412828in}}%
\pgfpathlineto{\pgfqpoint{1.386507in}{1.413753in}}%
\pgfpathlineto{\pgfqpoint{1.390013in}{1.453572in}}%
\pgfpathlineto{\pgfqpoint{1.395146in}{1.457695in}}%
\pgfpathlineto{\pgfqpoint{1.399572in}{1.458189in}}%
\pgfpathlineto{\pgfqpoint{1.429680in}{1.455658in}}%
\pgfpathlineto{\pgfqpoint{1.431599in}{1.456389in}}%
\pgfpathlineto{\pgfqpoint{1.430442in}{1.441625in}}%
\pgfpathlineto{\pgfqpoint{1.431343in}{1.435804in}}%
\pgfpathlineto{\pgfqpoint{1.473591in}{1.432656in}}%
\pgfpathlineto{\pgfqpoint{1.471839in}{1.406868in}}%
\pgfpathlineto{\pgfqpoint{1.430688in}{1.409985in}}%
\pgfpathlineto{\pgfqpoint{1.430345in}{1.395850in}}%
\pgfpathlineto{\pgfqpoint{1.428602in}{1.373346in}}%
\pgfpathlineto{\pgfqpoint{1.429987in}{1.373229in}}%
\pgfpathlineto{\pgfqpoint{1.428115in}{1.350443in}}%
\pgfpathlineto{\pgfqpoint{1.426380in}{1.344786in}}%
\pgfpathclose%
\pgfusepath{fill}%
\end{pgfscope}%
\begin{pgfscope}%
\pgfpathrectangle{\pgfqpoint{0.100000in}{0.100000in}}{\pgfqpoint{3.007045in}{1.925000in}}%
\pgfusepath{clip}%
\pgfsetbuttcap%
\pgfsetmiterjoin%
\definecolor{currentfill}{rgb}{0.541961,0.749527,0.865606}%
\pgfsetfillcolor{currentfill}%
\pgfsetlinewidth{0.000000pt}%
\definecolor{currentstroke}{rgb}{0.000000,0.000000,0.000000}%
\pgfsetstrokecolor{currentstroke}%
\pgfsetstrokeopacity{0.000000}%
\pgfsetdash{}{0pt}%
\pgfpathmoveto{\pgfqpoint{1.797461in}{1.775190in}}%
\pgfpathlineto{\pgfqpoint{1.804817in}{1.775072in}}%
\pgfpathlineto{\pgfqpoint{1.805689in}{1.772077in}}%
\pgfpathlineto{\pgfqpoint{1.823819in}{1.770365in}}%
\pgfpathlineto{\pgfqpoint{1.825676in}{1.763110in}}%
\pgfpathlineto{\pgfqpoint{1.840111in}{1.765309in}}%
\pgfpathlineto{\pgfqpoint{1.840141in}{1.768138in}}%
\pgfpathlineto{\pgfqpoint{1.851487in}{1.771989in}}%
\pgfpathlineto{\pgfqpoint{1.856761in}{1.771102in}}%
\pgfpathlineto{\pgfqpoint{1.857189in}{1.733722in}}%
\pgfpathlineto{\pgfqpoint{1.858121in}{1.722001in}}%
\pgfpathlineto{\pgfqpoint{1.827059in}{1.721935in}}%
\pgfpathlineto{\pgfqpoint{1.827086in}{1.718431in}}%
\pgfpathlineto{\pgfqpoint{1.798346in}{1.718105in}}%
\pgfpathlineto{\pgfqpoint{1.797671in}{1.752902in}}%
\pgfpathlineto{\pgfqpoint{1.797461in}{1.775190in}}%
\pgfpathclose%
\pgfusepath{fill}%
\end{pgfscope}%
\begin{pgfscope}%
\pgfpathrectangle{\pgfqpoint{0.100000in}{0.100000in}}{\pgfqpoint{3.007045in}{1.925000in}}%
\pgfusepath{clip}%
\pgfsetbuttcap%
\pgfsetmiterjoin%
\definecolor{currentfill}{rgb}{0.326290,0.618624,0.802799}%
\pgfsetfillcolor{currentfill}%
\pgfsetlinewidth{0.000000pt}%
\definecolor{currentstroke}{rgb}{0.000000,0.000000,0.000000}%
\pgfsetstrokecolor{currentstroke}%
\pgfsetstrokeopacity{0.000000}%
\pgfsetdash{}{0pt}%
\pgfpathmoveto{\pgfqpoint{2.145056in}{1.262796in}}%
\pgfpathlineto{\pgfqpoint{2.146371in}{1.246681in}}%
\pgfpathlineto{\pgfqpoint{2.125964in}{1.244692in}}%
\pgfpathlineto{\pgfqpoint{2.116838in}{1.244113in}}%
\pgfpathlineto{\pgfqpoint{2.113613in}{1.277542in}}%
\pgfpathlineto{\pgfqpoint{2.107905in}{1.276923in}}%
\pgfpathlineto{\pgfqpoint{2.107101in}{1.284740in}}%
\pgfpathlineto{\pgfqpoint{2.107035in}{1.290511in}}%
\pgfpathlineto{\pgfqpoint{2.118514in}{1.291662in}}%
\pgfpathlineto{\pgfqpoint{2.117970in}{1.297350in}}%
\pgfpathlineto{\pgfqpoint{2.141986in}{1.299655in}}%
\pgfpathlineto{\pgfqpoint{2.145056in}{1.262796in}}%
\pgfpathclose%
\pgfusepath{fill}%
\end{pgfscope}%
\begin{pgfscope}%
\pgfpathrectangle{\pgfqpoint{0.100000in}{0.100000in}}{\pgfqpoint{3.007045in}{1.925000in}}%
\pgfusepath{clip}%
\pgfsetbuttcap%
\pgfsetmiterjoin%
\definecolor{currentfill}{rgb}{0.406997,0.673741,0.834295}%
\pgfsetfillcolor{currentfill}%
\pgfsetlinewidth{0.000000pt}%
\definecolor{currentstroke}{rgb}{0.000000,0.000000,0.000000}%
\pgfsetstrokecolor{currentstroke}%
\pgfsetstrokeopacity{0.000000}%
\pgfsetdash{}{0pt}%
\pgfpathmoveto{\pgfqpoint{2.060711in}{0.508116in}}%
\pgfpathlineto{\pgfqpoint{2.054795in}{0.510464in}}%
\pgfpathlineto{\pgfqpoint{2.051671in}{0.516053in}}%
\pgfpathlineto{\pgfqpoint{2.053131in}{0.522603in}}%
\pgfpathlineto{\pgfqpoint{2.050269in}{0.524553in}}%
\pgfpathlineto{\pgfqpoint{2.042354in}{0.524234in}}%
\pgfpathlineto{\pgfqpoint{2.040150in}{0.528704in}}%
\pgfpathlineto{\pgfqpoint{2.034075in}{0.531818in}}%
\pgfpathlineto{\pgfqpoint{2.030059in}{0.536911in}}%
\pgfpathlineto{\pgfqpoint{2.010082in}{0.536853in}}%
\pgfpathlineto{\pgfqpoint{2.005579in}{0.541500in}}%
\pgfpathlineto{\pgfqpoint{2.005096in}{0.547382in}}%
\pgfpathlineto{\pgfqpoint{2.007733in}{0.550626in}}%
\pgfpathlineto{\pgfqpoint{2.023323in}{0.554899in}}%
\pgfpathlineto{\pgfqpoint{2.029617in}{0.558203in}}%
\pgfpathlineto{\pgfqpoint{2.032903in}{0.562233in}}%
\pgfpathlineto{\pgfqpoint{2.042685in}{0.564352in}}%
\pgfpathlineto{\pgfqpoint{2.046065in}{0.559905in}}%
\pgfpathlineto{\pgfqpoint{2.043709in}{0.592924in}}%
\pgfpathlineto{\pgfqpoint{2.037435in}{0.611183in}}%
\pgfpathlineto{\pgfqpoint{2.072666in}{0.613237in}}%
\pgfpathlineto{\pgfqpoint{2.070931in}{0.604028in}}%
\pgfpathlineto{\pgfqpoint{2.067754in}{0.598346in}}%
\pgfpathlineto{\pgfqpoint{2.067276in}{0.590315in}}%
\pgfpathlineto{\pgfqpoint{2.072286in}{0.580627in}}%
\pgfpathlineto{\pgfqpoint{2.077352in}{0.576813in}}%
\pgfpathlineto{\pgfqpoint{2.082124in}{0.561840in}}%
\pgfpathlineto{\pgfqpoint{2.087534in}{0.559447in}}%
\pgfpathlineto{\pgfqpoint{2.082198in}{0.557399in}}%
\pgfpathlineto{\pgfqpoint{2.076836in}{0.548759in}}%
\pgfpathlineto{\pgfqpoint{2.071543in}{0.549265in}}%
\pgfpathlineto{\pgfqpoint{2.070214in}{0.542954in}}%
\pgfpathlineto{\pgfqpoint{2.075634in}{0.542956in}}%
\pgfpathlineto{\pgfqpoint{2.077930in}{0.538591in}}%
\pgfpathlineto{\pgfqpoint{2.084493in}{0.539170in}}%
\pgfpathlineto{\pgfqpoint{2.085010in}{0.546770in}}%
\pgfpathlineto{\pgfqpoint{2.089855in}{0.550151in}}%
\pgfpathlineto{\pgfqpoint{2.100634in}{0.542949in}}%
\pgfpathlineto{\pgfqpoint{2.094872in}{0.530452in}}%
\pgfpathlineto{\pgfqpoint{2.081945in}{0.530204in}}%
\pgfpathlineto{\pgfqpoint{2.084051in}{0.524154in}}%
\pgfpathlineto{\pgfqpoint{2.081293in}{0.516251in}}%
\pgfpathlineto{\pgfqpoint{2.087768in}{0.513497in}}%
\pgfpathlineto{\pgfqpoint{2.093764in}{0.507672in}}%
\pgfpathlineto{\pgfqpoint{2.109818in}{0.505073in}}%
\pgfpathlineto{\pgfqpoint{2.117129in}{0.495373in}}%
\pgfpathlineto{\pgfqpoint{2.122084in}{0.494513in}}%
\pgfpathlineto{\pgfqpoint{2.119160in}{0.486899in}}%
\pgfpathlineto{\pgfqpoint{2.113765in}{0.483500in}}%
\pgfpathlineto{\pgfqpoint{2.109844in}{0.485263in}}%
\pgfpathlineto{\pgfqpoint{2.102674in}{0.484267in}}%
\pgfpathlineto{\pgfqpoint{2.093049in}{0.495483in}}%
\pgfpathlineto{\pgfqpoint{2.079626in}{0.500114in}}%
\pgfpathlineto{\pgfqpoint{2.072262in}{0.500507in}}%
\pgfpathlineto{\pgfqpoint{2.073553in}{0.504658in}}%
\pgfpathlineto{\pgfqpoint{2.072286in}{0.510828in}}%
\pgfpathlineto{\pgfqpoint{2.060711in}{0.508116in}}%
\pgfpathclose%
\pgfusepath{fill}%
\end{pgfscope}%
\begin{pgfscope}%
\pgfpathrectangle{\pgfqpoint{0.100000in}{0.100000in}}{\pgfqpoint{3.007045in}{1.925000in}}%
\pgfusepath{clip}%
\pgfsetbuttcap%
\pgfsetmiterjoin%
\definecolor{currentfill}{rgb}{0.406997,0.673741,0.834295}%
\pgfsetfillcolor{currentfill}%
\pgfsetlinewidth{0.000000pt}%
\definecolor{currentstroke}{rgb}{0.000000,0.000000,0.000000}%
\pgfsetstrokecolor{currentstroke}%
\pgfsetstrokeopacity{0.000000}%
\pgfsetdash{}{0pt}%
\pgfpathmoveto{\pgfqpoint{1.938345in}{1.353373in}}%
\pgfpathlineto{\pgfqpoint{1.904114in}{1.352301in}}%
\pgfpathlineto{\pgfqpoint{1.881224in}{1.351626in}}%
\pgfpathlineto{\pgfqpoint{1.881395in}{1.345922in}}%
\pgfpathlineto{\pgfqpoint{1.869963in}{1.345597in}}%
\pgfpathlineto{\pgfqpoint{1.869522in}{1.362845in}}%
\pgfpathlineto{\pgfqpoint{1.868339in}{1.362827in}}%
\pgfpathlineto{\pgfqpoint{1.867818in}{1.391658in}}%
\pgfpathlineto{\pgfqpoint{1.913408in}{1.392837in}}%
\pgfpathlineto{\pgfqpoint{1.913064in}{1.404438in}}%
\pgfpathlineto{\pgfqpoint{1.936000in}{1.405159in}}%
\pgfpathlineto{\pgfqpoint{1.937355in}{1.364751in}}%
\pgfpathlineto{\pgfqpoint{1.938345in}{1.353373in}}%
\pgfpathclose%
\pgfusepath{fill}%
\end{pgfscope}%
\begin{pgfscope}%
\pgfpathrectangle{\pgfqpoint{0.100000in}{0.100000in}}{\pgfqpoint{3.007045in}{1.925000in}}%
\pgfusepath{clip}%
\pgfsetbuttcap%
\pgfsetmiterjoin%
\definecolor{currentfill}{rgb}{0.311157,0.608289,0.796894}%
\pgfsetfillcolor{currentfill}%
\pgfsetlinewidth{0.000000pt}%
\definecolor{currentstroke}{rgb}{0.000000,0.000000,0.000000}%
\pgfsetstrokecolor{currentstroke}%
\pgfsetstrokeopacity{0.000000}%
\pgfsetdash{}{0pt}%
\pgfpathmoveto{\pgfqpoint{2.405437in}{0.788426in}}%
\pgfpathlineto{\pgfqpoint{2.390888in}{0.783936in}}%
\pgfpathlineto{\pgfqpoint{2.390153in}{0.787078in}}%
\pgfpathlineto{\pgfqpoint{2.390080in}{0.791606in}}%
\pgfpathlineto{\pgfqpoint{2.386212in}{0.799739in}}%
\pgfpathlineto{\pgfqpoint{2.382395in}{0.804196in}}%
\pgfpathlineto{\pgfqpoint{2.375369in}{0.793380in}}%
\pgfpathlineto{\pgfqpoint{2.365423in}{0.795571in}}%
\pgfpathlineto{\pgfqpoint{2.357695in}{0.794657in}}%
\pgfpathlineto{\pgfqpoint{2.355976in}{0.788053in}}%
\pgfpathlineto{\pgfqpoint{2.352410in}{0.787703in}}%
\pgfpathlineto{\pgfqpoint{2.346769in}{0.793773in}}%
\pgfpathlineto{\pgfqpoint{2.343332in}{0.803548in}}%
\pgfpathlineto{\pgfqpoint{2.332358in}{0.802455in}}%
\pgfpathlineto{\pgfqpoint{2.327623in}{0.805924in}}%
\pgfpathlineto{\pgfqpoint{2.326336in}{0.819580in}}%
\pgfpathlineto{\pgfqpoint{2.318665in}{0.820809in}}%
\pgfpathlineto{\pgfqpoint{2.317320in}{0.826818in}}%
\pgfpathlineto{\pgfqpoint{2.320886in}{0.830332in}}%
\pgfpathlineto{\pgfqpoint{2.322977in}{0.839287in}}%
\pgfpathlineto{\pgfqpoint{2.337348in}{0.840608in}}%
\pgfpathlineto{\pgfqpoint{2.334077in}{0.871563in}}%
\pgfpathlineto{\pgfqpoint{2.349378in}{0.873307in}}%
\pgfpathlineto{\pgfqpoint{2.351405in}{0.875886in}}%
\pgfpathlineto{\pgfqpoint{2.358828in}{0.878027in}}%
\pgfpathlineto{\pgfqpoint{2.359050in}{0.873951in}}%
\pgfpathlineto{\pgfqpoint{2.365118in}{0.869291in}}%
\pgfpathlineto{\pgfqpoint{2.371582in}{0.867326in}}%
\pgfpathlineto{\pgfqpoint{2.375538in}{0.858536in}}%
\pgfpathlineto{\pgfqpoint{2.372820in}{0.852600in}}%
\pgfpathlineto{\pgfqpoint{2.380206in}{0.847045in}}%
\pgfpathlineto{\pgfqpoint{2.382703in}{0.849302in}}%
\pgfpathlineto{\pgfqpoint{2.380905in}{0.840878in}}%
\pgfpathlineto{\pgfqpoint{2.387448in}{0.834640in}}%
\pgfpathlineto{\pgfqpoint{2.392091in}{0.836488in}}%
\pgfpathlineto{\pgfqpoint{2.403550in}{0.830476in}}%
\pgfpathlineto{\pgfqpoint{2.408663in}{0.823968in}}%
\pgfpathlineto{\pgfqpoint{2.417379in}{0.810753in}}%
\pgfpathlineto{\pgfqpoint{2.403890in}{0.805755in}}%
\pgfpathlineto{\pgfqpoint{2.405437in}{0.788426in}}%
\pgfpathclose%
\pgfusepath{fill}%
\end{pgfscope}%
\begin{pgfscope}%
\pgfpathrectangle{\pgfqpoint{0.100000in}{0.100000in}}{\pgfqpoint{3.007045in}{1.925000in}}%
\pgfusepath{clip}%
\pgfsetbuttcap%
\pgfsetmiterjoin%
\definecolor{currentfill}{rgb}{0.787543,0.866205,0.940946}%
\pgfsetfillcolor{currentfill}%
\pgfsetlinewidth{0.000000pt}%
\definecolor{currentstroke}{rgb}{0.000000,0.000000,0.000000}%
\pgfsetstrokecolor{currentstroke}%
\pgfsetstrokeopacity{0.000000}%
\pgfsetdash{}{0pt}%
\pgfpathmoveto{\pgfqpoint{2.418790in}{0.913875in}}%
\pgfpathlineto{\pgfqpoint{2.413670in}{0.911416in}}%
\pgfpathlineto{\pgfqpoint{2.410454in}{0.917486in}}%
\pgfpathlineto{\pgfqpoint{2.406157in}{0.920864in}}%
\pgfpathlineto{\pgfqpoint{2.398481in}{0.931422in}}%
\pgfpathlineto{\pgfqpoint{2.392521in}{0.928792in}}%
\pgfpathlineto{\pgfqpoint{2.381105in}{0.925986in}}%
\pgfpathlineto{\pgfqpoint{2.379773in}{0.923563in}}%
\pgfpathlineto{\pgfqpoint{2.372864in}{0.923192in}}%
\pgfpathlineto{\pgfqpoint{2.366318in}{0.919795in}}%
\pgfpathlineto{\pgfqpoint{2.361956in}{0.924269in}}%
\pgfpathlineto{\pgfqpoint{2.361650in}{0.931386in}}%
\pgfpathlineto{\pgfqpoint{2.364223in}{0.935952in}}%
\pgfpathlineto{\pgfqpoint{2.373522in}{0.943638in}}%
\pgfpathlineto{\pgfqpoint{2.379259in}{0.944820in}}%
\pgfpathlineto{\pgfqpoint{2.388077in}{0.945524in}}%
\pgfpathlineto{\pgfqpoint{2.399943in}{0.955887in}}%
\pgfpathlineto{\pgfqpoint{2.403998in}{0.954773in}}%
\pgfpathlineto{\pgfqpoint{2.405701in}{0.945045in}}%
\pgfpathlineto{\pgfqpoint{2.420735in}{0.932481in}}%
\pgfpathlineto{\pgfqpoint{2.421150in}{0.927599in}}%
\pgfpathlineto{\pgfqpoint{2.416370in}{0.915795in}}%
\pgfpathlineto{\pgfqpoint{2.418790in}{0.913875in}}%
\pgfpathclose%
\pgfusepath{fill}%
\end{pgfscope}%
\begin{pgfscope}%
\pgfpathrectangle{\pgfqpoint{0.100000in}{0.100000in}}{\pgfqpoint{3.007045in}{1.925000in}}%
\pgfusepath{clip}%
\pgfsetbuttcap%
\pgfsetmiterjoin%
\definecolor{currentfill}{rgb}{0.454118,0.701300,0.846659}%
\pgfsetfillcolor{currentfill}%
\pgfsetlinewidth{0.000000pt}%
\definecolor{currentstroke}{rgb}{0.000000,0.000000,0.000000}%
\pgfsetstrokecolor{currentstroke}%
\pgfsetstrokeopacity{0.000000}%
\pgfsetdash{}{0pt}%
\pgfpathmoveto{\pgfqpoint{2.044679in}{1.394919in}}%
\pgfpathlineto{\pgfqpoint{2.046233in}{1.372218in}}%
\pgfpathlineto{\pgfqpoint{2.018995in}{1.370976in}}%
\pgfpathlineto{\pgfqpoint{1.984280in}{1.369339in}}%
\pgfpathlineto{\pgfqpoint{1.980017in}{1.377888in}}%
\pgfpathlineto{\pgfqpoint{1.969155in}{1.380236in}}%
\pgfpathlineto{\pgfqpoint{1.963562in}{1.383583in}}%
\pgfpathlineto{\pgfqpoint{1.960884in}{1.393037in}}%
\pgfpathlineto{\pgfqpoint{1.958754in}{1.394373in}}%
\pgfpathlineto{\pgfqpoint{1.956575in}{1.409480in}}%
\pgfpathlineto{\pgfqpoint{1.961879in}{1.417685in}}%
\pgfpathlineto{\pgfqpoint{1.954703in}{1.423577in}}%
\pgfpathlineto{\pgfqpoint{1.954322in}{1.428443in}}%
\pgfpathlineto{\pgfqpoint{1.980060in}{1.429597in}}%
\pgfpathlineto{\pgfqpoint{1.980936in}{1.413027in}}%
\pgfpathlineto{\pgfqpoint{1.986647in}{1.415927in}}%
\pgfpathlineto{\pgfqpoint{1.992238in}{1.415509in}}%
\pgfpathlineto{\pgfqpoint{1.993663in}{1.389926in}}%
\pgfpathlineto{\pgfqpoint{2.022119in}{1.391513in}}%
\pgfpathlineto{\pgfqpoint{2.021967in}{1.394387in}}%
\pgfpathlineto{\pgfqpoint{2.044679in}{1.394919in}}%
\pgfpathclose%
\pgfusepath{fill}%
\end{pgfscope}%
\begin{pgfscope}%
\pgfpathrectangle{\pgfqpoint{0.100000in}{0.100000in}}{\pgfqpoint{3.007045in}{1.925000in}}%
\pgfusepath{clip}%
\pgfsetbuttcap%
\pgfsetmiterjoin%
\definecolor{currentfill}{rgb}{0.381776,0.656517,0.824452}%
\pgfsetfillcolor{currentfill}%
\pgfsetlinewidth{0.000000pt}%
\definecolor{currentstroke}{rgb}{0.000000,0.000000,0.000000}%
\pgfsetstrokecolor{currentstroke}%
\pgfsetstrokeopacity{0.000000}%
\pgfsetdash{}{0pt}%
\pgfpathmoveto{\pgfqpoint{1.923472in}{0.954935in}}%
\pgfpathlineto{\pgfqpoint{1.923336in}{0.970590in}}%
\pgfpathlineto{\pgfqpoint{1.924925in}{0.970656in}}%
\pgfpathlineto{\pgfqpoint{1.924728in}{0.990132in}}%
\pgfpathlineto{\pgfqpoint{1.925420in}{1.007434in}}%
\pgfpathlineto{\pgfqpoint{1.916864in}{1.007392in}}%
\pgfpathlineto{\pgfqpoint{1.915878in}{1.043286in}}%
\pgfpathlineto{\pgfqpoint{1.941697in}{1.043763in}}%
\pgfpathlineto{\pgfqpoint{1.942096in}{1.032256in}}%
\pgfpathlineto{\pgfqpoint{1.972430in}{1.032823in}}%
\pgfpathlineto{\pgfqpoint{1.972674in}{1.026092in}}%
\pgfpathlineto{\pgfqpoint{1.975616in}{1.018755in}}%
\pgfpathlineto{\pgfqpoint{1.981501in}{1.015319in}}%
\pgfpathlineto{\pgfqpoint{1.981594in}{1.012166in}}%
\pgfpathlineto{\pgfqpoint{1.971148in}{1.011094in}}%
\pgfpathlineto{\pgfqpoint{1.971458in}{0.993921in}}%
\pgfpathlineto{\pgfqpoint{1.977205in}{0.994031in}}%
\pgfpathlineto{\pgfqpoint{1.977561in}{0.972537in}}%
\pgfpathlineto{\pgfqpoint{1.962703in}{0.971887in}}%
\pgfpathlineto{\pgfqpoint{1.960431in}{0.971833in}}%
\pgfpathlineto{\pgfqpoint{1.960613in}{0.961138in}}%
\pgfpathlineto{\pgfqpoint{1.954694in}{0.961038in}}%
\pgfpathlineto{\pgfqpoint{1.954901in}{0.955275in}}%
\pgfpathlineto{\pgfqpoint{1.923472in}{0.954935in}}%
\pgfpathclose%
\pgfusepath{fill}%
\end{pgfscope}%
\begin{pgfscope}%
\pgfpathrectangle{\pgfqpoint{0.100000in}{0.100000in}}{\pgfqpoint{3.007045in}{1.925000in}}%
\pgfusepath{clip}%
\pgfsetbuttcap%
\pgfsetmiterjoin%
\definecolor{currentfill}{rgb}{0.366644,0.646182,0.818547}%
\pgfsetfillcolor{currentfill}%
\pgfsetlinewidth{0.000000pt}%
\definecolor{currentstroke}{rgb}{0.000000,0.000000,0.000000}%
\pgfsetstrokecolor{currentstroke}%
\pgfsetstrokeopacity{0.000000}%
\pgfsetdash{}{0pt}%
\pgfpathmoveto{\pgfqpoint{1.018543in}{1.560865in}}%
\pgfpathlineto{\pgfqpoint{1.020483in}{1.573439in}}%
\pgfpathlineto{\pgfqpoint{1.038445in}{1.570273in}}%
\pgfpathlineto{\pgfqpoint{1.037496in}{1.564829in}}%
\pgfpathlineto{\pgfqpoint{1.051087in}{1.562499in}}%
\pgfpathlineto{\pgfqpoint{1.054900in}{1.557314in}}%
\pgfpathlineto{\pgfqpoint{1.053058in}{1.551890in}}%
\pgfpathlineto{\pgfqpoint{1.057765in}{1.542345in}}%
\pgfpathlineto{\pgfqpoint{1.058870in}{1.534775in}}%
\pgfpathlineto{\pgfqpoint{1.061138in}{1.530583in}}%
\pgfpathlineto{\pgfqpoint{1.059788in}{1.522455in}}%
\pgfpathlineto{\pgfqpoint{1.053944in}{1.487032in}}%
\pgfpathlineto{\pgfqpoint{1.052724in}{1.481408in}}%
\pgfpathlineto{\pgfqpoint{1.039144in}{1.483711in}}%
\pgfpathlineto{\pgfqpoint{1.038193in}{1.478061in}}%
\pgfpathlineto{\pgfqpoint{1.027092in}{1.480000in}}%
\pgfpathlineto{\pgfqpoint{1.026443in}{1.476331in}}%
\pgfpathlineto{\pgfqpoint{1.015288in}{1.478274in}}%
\pgfpathlineto{\pgfqpoint{1.016022in}{1.482462in}}%
\pgfpathlineto{\pgfqpoint{1.005297in}{1.485446in}}%
\pgfpathlineto{\pgfqpoint{1.007407in}{1.497511in}}%
\pgfpathlineto{\pgfqpoint{1.000614in}{1.501315in}}%
\pgfpathlineto{\pgfqpoint{0.999466in}{1.507030in}}%
\pgfpathlineto{\pgfqpoint{0.992126in}{1.508384in}}%
\pgfpathlineto{\pgfqpoint{0.995747in}{1.527955in}}%
\pgfpathlineto{\pgfqpoint{1.005109in}{1.526977in}}%
\pgfpathlineto{\pgfqpoint{1.012932in}{1.528900in}}%
\pgfpathlineto{\pgfqpoint{1.018543in}{1.560865in}}%
\pgfpathclose%
\pgfusepath{fill}%
\end{pgfscope}%
\begin{pgfscope}%
\pgfpathrectangle{\pgfqpoint{0.100000in}{0.100000in}}{\pgfqpoint{3.007045in}{1.925000in}}%
\pgfusepath{clip}%
\pgfsetbuttcap%
\pgfsetmiterjoin%
\definecolor{currentfill}{rgb}{0.326290,0.618624,0.802799}%
\pgfsetfillcolor{currentfill}%
\pgfsetlinewidth{0.000000pt}%
\definecolor{currentstroke}{rgb}{0.000000,0.000000,0.000000}%
\pgfsetstrokecolor{currentstroke}%
\pgfsetstrokeopacity{0.000000}%
\pgfsetdash{}{0pt}%
\pgfpathmoveto{\pgfqpoint{1.578119in}{0.708542in}}%
\pgfpathlineto{\pgfqpoint{1.577602in}{0.694365in}}%
\pgfpathlineto{\pgfqpoint{1.551784in}{0.680740in}}%
\pgfpathlineto{\pgfqpoint{1.507461in}{0.683119in}}%
\pgfpathlineto{\pgfqpoint{1.482882in}{0.684491in}}%
\pgfpathlineto{\pgfqpoint{1.486605in}{0.742533in}}%
\pgfpathlineto{\pgfqpoint{1.523941in}{0.740360in}}%
\pgfpathlineto{\pgfqpoint{1.573659in}{0.737980in}}%
\pgfpathlineto{\pgfqpoint{1.572479in}{0.708902in}}%
\pgfpathlineto{\pgfqpoint{1.578119in}{0.708542in}}%
\pgfpathclose%
\pgfusepath{fill}%
\end{pgfscope}%
\begin{pgfscope}%
\pgfpathrectangle{\pgfqpoint{0.100000in}{0.100000in}}{\pgfqpoint{3.007045in}{1.925000in}}%
\pgfusepath{clip}%
\pgfsetbuttcap%
\pgfsetmiterjoin%
\definecolor{currentfill}{rgb}{0.142607,0.456332,0.716601}%
\pgfsetfillcolor{currentfill}%
\pgfsetlinewidth{0.000000pt}%
\definecolor{currentstroke}{rgb}{0.000000,0.000000,0.000000}%
\pgfsetstrokecolor{currentstroke}%
\pgfsetstrokeopacity{0.000000}%
\pgfsetdash{}{0pt}%
\pgfpathmoveto{\pgfqpoint{0.936407in}{1.175969in}}%
\pgfpathlineto{\pgfqpoint{0.936286in}{1.175309in}}%
\pgfpathlineto{\pgfqpoint{0.968213in}{1.169632in}}%
\pgfpathlineto{\pgfqpoint{1.001312in}{1.163987in}}%
\pgfpathlineto{\pgfqpoint{1.002930in}{1.149720in}}%
\pgfpathlineto{\pgfqpoint{1.005017in}{1.143547in}}%
\pgfpathlineto{\pgfqpoint{1.002522in}{1.140541in}}%
\pgfpathlineto{\pgfqpoint{0.963729in}{1.147104in}}%
\pgfpathlineto{\pgfqpoint{0.904834in}{1.157752in}}%
\pgfpathlineto{\pgfqpoint{0.907847in}{1.175524in}}%
\pgfpathlineto{\pgfqpoint{0.913573in}{1.180156in}}%
\pgfpathlineto{\pgfqpoint{0.936407in}{1.175969in}}%
\pgfpathclose%
\pgfusepath{fill}%
\end{pgfscope}%
\begin{pgfscope}%
\pgfpathrectangle{\pgfqpoint{0.100000in}{0.100000in}}{\pgfqpoint{3.007045in}{1.925000in}}%
\pgfusepath{clip}%
\pgfsetbuttcap%
\pgfsetmiterjoin%
\definecolor{currentfill}{rgb}{0.301069,0.601399,0.792957}%
\pgfsetfillcolor{currentfill}%
\pgfsetlinewidth{0.000000pt}%
\definecolor{currentstroke}{rgb}{0.000000,0.000000,0.000000}%
\pgfsetstrokecolor{currentstroke}%
\pgfsetstrokeopacity{0.000000}%
\pgfsetdash{}{0pt}%
\pgfpathmoveto{\pgfqpoint{1.786299in}{0.783647in}}%
\pgfpathlineto{\pgfqpoint{1.800644in}{0.779398in}}%
\pgfpathlineto{\pgfqpoint{1.801001in}{0.799445in}}%
\pgfpathlineto{\pgfqpoint{1.806094in}{0.799451in}}%
\pgfpathlineto{\pgfqpoint{1.814189in}{0.790353in}}%
\pgfpathlineto{\pgfqpoint{1.825054in}{0.789133in}}%
\pgfpathlineto{\pgfqpoint{1.829586in}{0.787129in}}%
\pgfpathlineto{\pgfqpoint{1.829258in}{0.781767in}}%
\pgfpathlineto{\pgfqpoint{1.836570in}{0.777705in}}%
\pgfpathlineto{\pgfqpoint{1.842812in}{0.769497in}}%
\pgfpathlineto{\pgfqpoint{1.842406in}{0.765728in}}%
\pgfpathlineto{\pgfqpoint{1.846121in}{0.759629in}}%
\pgfpathlineto{\pgfqpoint{1.841164in}{0.748497in}}%
\pgfpathlineto{\pgfqpoint{1.836073in}{0.747122in}}%
\pgfpathlineto{\pgfqpoint{1.835563in}{0.743537in}}%
\pgfpathlineto{\pgfqpoint{1.838686in}{0.738812in}}%
\pgfpathlineto{\pgfqpoint{1.825469in}{0.738647in}}%
\pgfpathlineto{\pgfqpoint{1.825570in}{0.729490in}}%
\pgfpathlineto{\pgfqpoint{1.788689in}{0.729118in}}%
\pgfpathlineto{\pgfqpoint{1.787896in}{0.730793in}}%
\pgfpathlineto{\pgfqpoint{1.763868in}{0.730641in}}%
\pgfpathlineto{\pgfqpoint{1.762705in}{0.734578in}}%
\pgfpathlineto{\pgfqpoint{1.755158in}{0.734645in}}%
\pgfpathlineto{\pgfqpoint{1.755274in}{0.762070in}}%
\pgfpathlineto{\pgfqpoint{1.767169in}{0.763420in}}%
\pgfpathlineto{\pgfqpoint{1.779169in}{0.763215in}}%
\pgfpathlineto{\pgfqpoint{1.786323in}{0.758894in}}%
\pgfpathlineto{\pgfqpoint{1.786299in}{0.783647in}}%
\pgfpathclose%
\pgfusepath{fill}%
\end{pgfscope}%
\begin{pgfscope}%
\pgfpathrectangle{\pgfqpoint{0.100000in}{0.100000in}}{\pgfqpoint{3.007045in}{1.925000in}}%
\pgfusepath{clip}%
\pgfsetbuttcap%
\pgfsetmiterjoin%
\definecolor{currentfill}{rgb}{0.750634,0.847843,0.928212}%
\pgfsetfillcolor{currentfill}%
\pgfsetlinewidth{0.000000pt}%
\definecolor{currentstroke}{rgb}{0.000000,0.000000,0.000000}%
\pgfsetstrokecolor{currentstroke}%
\pgfsetstrokeopacity{0.000000}%
\pgfsetdash{}{0pt}%
\pgfpathmoveto{\pgfqpoint{2.517736in}{0.501110in}}%
\pgfpathlineto{\pgfqpoint{2.508454in}{0.501648in}}%
\pgfpathlineto{\pgfqpoint{2.495655in}{0.499863in}}%
\pgfpathlineto{\pgfqpoint{2.494027in}{0.512423in}}%
\pgfpathlineto{\pgfqpoint{2.486048in}{0.518807in}}%
\pgfpathlineto{\pgfqpoint{2.488304in}{0.521528in}}%
\pgfpathlineto{\pgfqpoint{2.494299in}{0.520921in}}%
\pgfpathlineto{\pgfqpoint{2.498437in}{0.523892in}}%
\pgfpathlineto{\pgfqpoint{2.496864in}{0.535136in}}%
\pgfpathlineto{\pgfqpoint{2.504526in}{0.536259in}}%
\pgfpathlineto{\pgfqpoint{2.501841in}{0.554053in}}%
\pgfpathlineto{\pgfqpoint{2.513059in}{0.555566in}}%
\pgfpathlineto{\pgfqpoint{2.513609in}{0.551975in}}%
\pgfpathlineto{\pgfqpoint{2.522149in}{0.556021in}}%
\pgfpathlineto{\pgfqpoint{2.533920in}{0.561053in}}%
\pgfpathlineto{\pgfqpoint{2.538085in}{0.559398in}}%
\pgfpathlineto{\pgfqpoint{2.541214in}{0.552000in}}%
\pgfpathlineto{\pgfqpoint{2.547864in}{0.546632in}}%
\pgfpathlineto{\pgfqpoint{2.550293in}{0.531506in}}%
\pgfpathlineto{\pgfqpoint{2.549990in}{0.525560in}}%
\pgfpathlineto{\pgfqpoint{2.512188in}{0.520102in}}%
\pgfpathlineto{\pgfqpoint{2.522009in}{0.510090in}}%
\pgfpathlineto{\pgfqpoint{2.517736in}{0.501110in}}%
\pgfpathclose%
\pgfusepath{fill}%
\end{pgfscope}%
\begin{pgfscope}%
\pgfpathrectangle{\pgfqpoint{0.100000in}{0.100000in}}{\pgfqpoint{3.007045in}{1.925000in}}%
\pgfusepath{clip}%
\pgfsetbuttcap%
\pgfsetmiterjoin%
\definecolor{currentfill}{rgb}{0.341423,0.628958,0.808704}%
\pgfsetfillcolor{currentfill}%
\pgfsetlinewidth{0.000000pt}%
\definecolor{currentstroke}{rgb}{0.000000,0.000000,0.000000}%
\pgfsetstrokecolor{currentstroke}%
\pgfsetstrokeopacity{0.000000}%
\pgfsetdash{}{0pt}%
\pgfpathmoveto{\pgfqpoint{1.685427in}{0.912299in}}%
\pgfpathlineto{\pgfqpoint{1.708380in}{0.911948in}}%
\pgfpathlineto{\pgfqpoint{1.708573in}{0.926325in}}%
\pgfpathlineto{\pgfqpoint{1.731364in}{0.926038in}}%
\pgfpathlineto{\pgfqpoint{1.731290in}{0.917416in}}%
\pgfpathlineto{\pgfqpoint{1.734135in}{0.917404in}}%
\pgfpathlineto{\pgfqpoint{1.734052in}{0.905918in}}%
\pgfpathlineto{\pgfqpoint{1.728181in}{0.905962in}}%
\pgfpathlineto{\pgfqpoint{1.728070in}{0.900172in}}%
\pgfpathlineto{\pgfqpoint{1.725208in}{0.894438in}}%
\pgfpathlineto{\pgfqpoint{1.719521in}{0.894505in}}%
\pgfpathlineto{\pgfqpoint{1.719456in}{0.888768in}}%
\pgfpathlineto{\pgfqpoint{1.694674in}{0.889113in}}%
\pgfpathlineto{\pgfqpoint{1.691799in}{0.898304in}}%
\pgfpathlineto{\pgfqpoint{1.684983in}{0.896595in}}%
\pgfpathlineto{\pgfqpoint{1.685427in}{0.912299in}}%
\pgfpathclose%
\pgfusepath{fill}%
\end{pgfscope}%
\begin{pgfscope}%
\pgfpathrectangle{\pgfqpoint{0.100000in}{0.100000in}}{\pgfqpoint{3.007045in}{1.925000in}}%
\pgfusepath{clip}%
\pgfsetbuttcap%
\pgfsetmiterjoin%
\definecolor{currentfill}{rgb}{0.840692,0.901638,0.958662}%
\pgfsetfillcolor{currentfill}%
\pgfsetlinewidth{0.000000pt}%
\definecolor{currentstroke}{rgb}{0.000000,0.000000,0.000000}%
\pgfsetstrokecolor{currentstroke}%
\pgfsetstrokeopacity{0.000000}%
\pgfsetdash{}{0pt}%
\pgfpathmoveto{\pgfqpoint{2.474213in}{0.983871in}}%
\pgfpathlineto{\pgfqpoint{2.468967in}{0.985644in}}%
\pgfpathlineto{\pgfqpoint{2.465463in}{0.979232in}}%
\pgfpathlineto{\pgfqpoint{2.462402in}{0.980829in}}%
\pgfpathlineto{\pgfqpoint{2.458453in}{0.990986in}}%
\pgfpathlineto{\pgfqpoint{2.460956in}{0.992711in}}%
\pgfpathlineto{\pgfqpoint{2.463756in}{1.002325in}}%
\pgfpathlineto{\pgfqpoint{2.458161in}{1.009811in}}%
\pgfpathlineto{\pgfqpoint{2.461449in}{1.016280in}}%
\pgfpathlineto{\pgfqpoint{2.462097in}{1.022483in}}%
\pgfpathlineto{\pgfqpoint{2.466964in}{1.026275in}}%
\pgfpathlineto{\pgfqpoint{2.476402in}{1.027509in}}%
\pgfpathlineto{\pgfqpoint{2.475027in}{1.025720in}}%
\pgfpathlineto{\pgfqpoint{2.475252in}{1.017742in}}%
\pgfpathlineto{\pgfqpoint{2.473634in}{1.012683in}}%
\pgfpathlineto{\pgfqpoint{2.479391in}{1.010461in}}%
\pgfpathlineto{\pgfqpoint{2.490105in}{1.002416in}}%
\pgfpathlineto{\pgfqpoint{2.486521in}{0.996045in}}%
\pgfpathlineto{\pgfqpoint{2.479657in}{0.995365in}}%
\pgfpathlineto{\pgfqpoint{2.474213in}{0.983871in}}%
\pgfpathclose%
\pgfusepath{fill}%
\end{pgfscope}%
\begin{pgfscope}%
\pgfpathrectangle{\pgfqpoint{0.100000in}{0.100000in}}{\pgfqpoint{3.007045in}{1.925000in}}%
\pgfusepath{clip}%
\pgfsetbuttcap%
\pgfsetmiterjoin%
\definecolor{currentfill}{rgb}{0.270804,0.580730,0.781146}%
\pgfsetfillcolor{currentfill}%
\pgfsetlinewidth{0.000000pt}%
\definecolor{currentstroke}{rgb}{0.000000,0.000000,0.000000}%
\pgfsetstrokecolor{currentstroke}%
\pgfsetstrokeopacity{0.000000}%
\pgfsetdash{}{0pt}%
\pgfpathmoveto{\pgfqpoint{1.527075in}{1.552396in}}%
\pgfpathlineto{\pgfqpoint{1.552729in}{1.550914in}}%
\pgfpathlineto{\pgfqpoint{1.559045in}{1.550570in}}%
\pgfpathlineto{\pgfqpoint{1.557925in}{1.527548in}}%
\pgfpathlineto{\pgfqpoint{1.553171in}{1.527805in}}%
\pgfpathlineto{\pgfqpoint{1.518970in}{1.529708in}}%
\pgfpathlineto{\pgfqpoint{1.518468in}{1.535998in}}%
\pgfpathlineto{\pgfqpoint{1.524818in}{1.537110in}}%
\pgfpathlineto{\pgfqpoint{1.523533in}{1.542968in}}%
\pgfpathlineto{\pgfqpoint{1.527075in}{1.552396in}}%
\pgfpathclose%
\pgfusepath{fill}%
\end{pgfscope}%
\begin{pgfscope}%
\pgfpathrectangle{\pgfqpoint{0.100000in}{0.100000in}}{\pgfqpoint{3.007045in}{1.925000in}}%
\pgfusepath{clip}%
\pgfsetbuttcap%
\pgfsetmiterjoin%
\definecolor{currentfill}{rgb}{0.031373,0.285675,0.564291}%
\pgfsetfillcolor{currentfill}%
\pgfsetlinewidth{0.000000pt}%
\definecolor{currentstroke}{rgb}{0.000000,0.000000,0.000000}%
\pgfsetstrokecolor{currentstroke}%
\pgfsetstrokeopacity{0.000000}%
\pgfsetdash{}{0pt}%
\pgfpathmoveto{\pgfqpoint{1.265329in}{0.695996in}}%
\pgfpathlineto{\pgfqpoint{1.282225in}{0.694354in}}%
\pgfpathlineto{\pgfqpoint{1.319120in}{0.690958in}}%
\pgfpathlineto{\pgfqpoint{1.319633in}{0.696676in}}%
\pgfpathlineto{\pgfqpoint{1.334479in}{0.695294in}}%
\pgfpathlineto{\pgfqpoint{1.332017in}{0.666591in}}%
\pgfpathlineto{\pgfqpoint{1.333802in}{0.666433in}}%
\pgfpathlineto{\pgfqpoint{1.331766in}{0.642740in}}%
\pgfpathlineto{\pgfqpoint{1.328596in}{0.641304in}}%
\pgfpathlineto{\pgfqpoint{1.322677in}{0.647278in}}%
\pgfpathlineto{\pgfqpoint{1.318428in}{0.649073in}}%
\pgfpathlineto{\pgfqpoint{1.282209in}{0.612003in}}%
\pgfpathlineto{\pgfqpoint{1.255027in}{0.637292in}}%
\pgfpathlineto{\pgfqpoint{1.265329in}{0.695996in}}%
\pgfpathclose%
\pgfusepath{fill}%
\end{pgfscope}%
\begin{pgfscope}%
\pgfpathrectangle{\pgfqpoint{0.100000in}{0.100000in}}{\pgfqpoint{3.007045in}{1.925000in}}%
\pgfusepath{clip}%
\pgfsetbuttcap%
\pgfsetmiterjoin%
\definecolor{currentfill}{rgb}{0.406997,0.673741,0.834295}%
\pgfsetfillcolor{currentfill}%
\pgfsetlinewidth{0.000000pt}%
\definecolor{currentstroke}{rgb}{0.000000,0.000000,0.000000}%
\pgfsetstrokecolor{currentstroke}%
\pgfsetstrokeopacity{0.000000}%
\pgfsetdash{}{0pt}%
\pgfpathmoveto{\pgfqpoint{2.742724in}{1.035580in}}%
\pgfpathlineto{\pgfqpoint{2.738199in}{1.043269in}}%
\pgfpathlineto{\pgfqpoint{2.738529in}{1.052500in}}%
\pgfpathlineto{\pgfqpoint{2.733124in}{1.055886in}}%
\pgfpathlineto{\pgfqpoint{2.724911in}{1.056794in}}%
\pgfpathlineto{\pgfqpoint{2.723871in}{1.066185in}}%
\pgfpathlineto{\pgfqpoint{2.723519in}{1.072852in}}%
\pgfpathlineto{\pgfqpoint{2.738779in}{1.094235in}}%
\pgfpathlineto{\pgfqpoint{2.741244in}{1.095623in}}%
\pgfpathlineto{\pgfqpoint{2.746605in}{1.094479in}}%
\pgfpathlineto{\pgfqpoint{2.750900in}{1.100328in}}%
\pgfpathlineto{\pgfqpoint{2.761478in}{1.098464in}}%
\pgfpathlineto{\pgfqpoint{2.766403in}{1.100448in}}%
\pgfpathlineto{\pgfqpoint{2.771017in}{1.090586in}}%
\pgfpathlineto{\pgfqpoint{2.777353in}{1.079986in}}%
\pgfpathlineto{\pgfqpoint{2.787444in}{1.058403in}}%
\pgfpathlineto{\pgfqpoint{2.786756in}{1.058061in}}%
\pgfpathlineto{\pgfqpoint{2.780444in}{1.069950in}}%
\pgfpathlineto{\pgfqpoint{2.777839in}{1.065735in}}%
\pgfpathlineto{\pgfqpoint{2.788743in}{1.048436in}}%
\pgfpathlineto{\pgfqpoint{2.782567in}{1.052153in}}%
\pgfpathlineto{\pgfqpoint{2.776592in}{1.052297in}}%
\pgfpathlineto{\pgfqpoint{2.766420in}{1.044865in}}%
\pgfpathlineto{\pgfqpoint{2.757806in}{1.041934in}}%
\pgfpathlineto{\pgfqpoint{2.752708in}{1.036350in}}%
\pgfpathlineto{\pgfqpoint{2.742724in}{1.035580in}}%
\pgfpathclose%
\pgfusepath{fill}%
\end{pgfscope}%
\begin{pgfscope}%
\pgfpathrectangle{\pgfqpoint{0.100000in}{0.100000in}}{\pgfqpoint{3.007045in}{1.925000in}}%
\pgfusepath{clip}%
\pgfsetbuttcap%
\pgfsetmiterjoin%
\definecolor{currentfill}{rgb}{0.529412,0.742637,0.862899}%
\pgfsetfillcolor{currentfill}%
\pgfsetlinewidth{0.000000pt}%
\definecolor{currentstroke}{rgb}{0.000000,0.000000,0.000000}%
\pgfsetstrokecolor{currentstroke}%
\pgfsetstrokeopacity{0.000000}%
\pgfsetdash{}{0pt}%
\pgfpathmoveto{\pgfqpoint{1.685760in}{0.932254in}}%
\pgfpathlineto{\pgfqpoint{1.658027in}{0.932849in}}%
\pgfpathlineto{\pgfqpoint{1.646710in}{0.935974in}}%
\pgfpathlineto{\pgfqpoint{1.647016in}{0.947800in}}%
\pgfpathlineto{\pgfqpoint{1.641329in}{0.947962in}}%
\pgfpathlineto{\pgfqpoint{1.642047in}{0.976241in}}%
\pgfpathlineto{\pgfqpoint{1.670693in}{0.976025in}}%
\pgfpathlineto{\pgfqpoint{1.671642in}{0.972779in}}%
\pgfpathlineto{\pgfqpoint{1.664644in}{0.969981in}}%
\pgfpathlineto{\pgfqpoint{1.664393in}{0.958534in}}%
\pgfpathlineto{\pgfqpoint{1.670084in}{0.958402in}}%
\pgfpathlineto{\pgfqpoint{1.669953in}{0.952684in}}%
\pgfpathlineto{\pgfqpoint{1.675597in}{0.952567in}}%
\pgfpathlineto{\pgfqpoint{1.675476in}{0.946791in}}%
\pgfpathlineto{\pgfqpoint{1.685998in}{0.946658in}}%
\pgfpathlineto{\pgfqpoint{1.685760in}{0.932254in}}%
\pgfpathclose%
\pgfusepath{fill}%
\end{pgfscope}%
\begin{pgfscope}%
\pgfpathrectangle{\pgfqpoint{0.100000in}{0.100000in}}{\pgfqpoint{3.007045in}{1.925000in}}%
\pgfusepath{clip}%
\pgfsetbuttcap%
\pgfsetmiterjoin%
\definecolor{currentfill}{rgb}{0.187266,0.500992,0.739608}%
\pgfsetfillcolor{currentfill}%
\pgfsetlinewidth{0.000000pt}%
\definecolor{currentstroke}{rgb}{0.000000,0.000000,0.000000}%
\pgfsetstrokecolor{currentstroke}%
\pgfsetstrokeopacity{0.000000}%
\pgfsetdash{}{0pt}%
\pgfpathmoveto{\pgfqpoint{1.564548in}{1.295579in}}%
\pgfpathlineto{\pgfqpoint{1.563271in}{1.272731in}}%
\pgfpathlineto{\pgfqpoint{1.552431in}{1.273213in}}%
\pgfpathlineto{\pgfqpoint{1.512812in}{1.275327in}}%
\pgfpathlineto{\pgfqpoint{1.514112in}{1.321020in}}%
\pgfpathlineto{\pgfqpoint{1.485080in}{1.322839in}}%
\pgfpathlineto{\pgfqpoint{1.486374in}{1.345738in}}%
\pgfpathlineto{\pgfqpoint{1.519516in}{1.343574in}}%
\pgfpathlineto{\pgfqpoint{1.543010in}{1.342313in}}%
\pgfpathlineto{\pgfqpoint{1.541840in}{1.319542in}}%
\pgfpathlineto{\pgfqpoint{1.565174in}{1.318418in}}%
\pgfpathlineto{\pgfqpoint{1.564548in}{1.295579in}}%
\pgfpathclose%
\pgfusepath{fill}%
\end{pgfscope}%
\begin{pgfscope}%
\pgfpathrectangle{\pgfqpoint{0.100000in}{0.100000in}}{\pgfqpoint{3.007045in}{1.925000in}}%
\pgfusepath{clip}%
\pgfsetbuttcap%
\pgfsetmiterjoin%
\definecolor{currentfill}{rgb}{0.265759,0.577286,0.779177}%
\pgfsetfillcolor{currentfill}%
\pgfsetlinewidth{0.000000pt}%
\definecolor{currentstroke}{rgb}{0.000000,0.000000,0.000000}%
\pgfsetstrokecolor{currentstroke}%
\pgfsetstrokeopacity{0.000000}%
\pgfsetdash{}{0pt}%
\pgfpathmoveto{\pgfqpoint{2.025623in}{0.904867in}}%
\pgfpathlineto{\pgfqpoint{2.008210in}{0.904668in}}%
\pgfpathlineto{\pgfqpoint{1.985207in}{0.903228in}}%
\pgfpathlineto{\pgfqpoint{1.984303in}{0.932387in}}%
\pgfpathlineto{\pgfqpoint{1.975423in}{0.932334in}}%
\pgfpathlineto{\pgfqpoint{1.966949in}{0.932026in}}%
\pgfpathlineto{\pgfqpoint{1.966442in}{0.954472in}}%
\pgfpathlineto{\pgfqpoint{1.971212in}{0.955492in}}%
\pgfpathlineto{\pgfqpoint{1.969967in}{0.964093in}}%
\pgfpathlineto{\pgfqpoint{1.962703in}{0.971887in}}%
\pgfpathlineto{\pgfqpoint{1.977561in}{0.972537in}}%
\pgfpathlineto{\pgfqpoint{2.025516in}{0.974832in}}%
\pgfpathlineto{\pgfqpoint{2.034251in}{0.967797in}}%
\pgfpathlineto{\pgfqpoint{2.033532in}{0.960342in}}%
\pgfpathlineto{\pgfqpoint{2.026615in}{0.954215in}}%
\pgfpathlineto{\pgfqpoint{2.021688in}{0.947698in}}%
\pgfpathlineto{\pgfqpoint{2.023624in}{0.941619in}}%
\pgfpathlineto{\pgfqpoint{2.025623in}{0.904867in}}%
\pgfpathclose%
\pgfusepath{fill}%
\end{pgfscope}%
\begin{pgfscope}%
\pgfpathrectangle{\pgfqpoint{0.100000in}{0.100000in}}{\pgfqpoint{3.007045in}{1.925000in}}%
\pgfusepath{clip}%
\pgfsetbuttcap%
\pgfsetmiterjoin%
\definecolor{currentfill}{rgb}{0.479216,0.715079,0.852072}%
\pgfsetfillcolor{currentfill}%
\pgfsetlinewidth{0.000000pt}%
\definecolor{currentstroke}{rgb}{0.000000,0.000000,0.000000}%
\pgfsetstrokecolor{currentstroke}%
\pgfsetstrokeopacity{0.000000}%
\pgfsetdash{}{0pt}%
\pgfpathmoveto{\pgfqpoint{1.169817in}{1.038957in}}%
\pgfpathlineto{\pgfqpoint{1.161246in}{1.055657in}}%
\pgfpathlineto{\pgfqpoint{1.162663in}{1.066636in}}%
\pgfpathlineto{\pgfqpoint{1.139253in}{1.069444in}}%
\pgfpathlineto{\pgfqpoint{1.139523in}{1.071438in}}%
\pgfpathlineto{\pgfqpoint{1.140887in}{1.088045in}}%
\pgfpathlineto{\pgfqpoint{1.143557in}{1.105192in}}%
\pgfpathlineto{\pgfqpoint{1.150792in}{1.105369in}}%
\pgfpathlineto{\pgfqpoint{1.154865in}{1.136033in}}%
\pgfpathlineto{\pgfqpoint{1.193414in}{1.130886in}}%
\pgfpathlineto{\pgfqpoint{1.202180in}{1.129872in}}%
\pgfpathlineto{\pgfqpoint{1.205700in}{1.130953in}}%
\pgfpathlineto{\pgfqpoint{1.220511in}{1.109227in}}%
\pgfpathlineto{\pgfqpoint{1.223896in}{1.096979in}}%
\pgfpathlineto{\pgfqpoint{1.229068in}{1.091653in}}%
\pgfpathlineto{\pgfqpoint{1.231022in}{1.088493in}}%
\pgfpathlineto{\pgfqpoint{1.225618in}{1.073383in}}%
\pgfpathlineto{\pgfqpoint{1.236457in}{1.074727in}}%
\pgfpathlineto{\pgfqpoint{1.241493in}{1.067690in}}%
\pgfpathlineto{\pgfqpoint{1.242720in}{1.057254in}}%
\pgfpathlineto{\pgfqpoint{1.241118in}{1.050301in}}%
\pgfpathlineto{\pgfqpoint{1.238883in}{1.030797in}}%
\pgfpathlineto{\pgfqpoint{1.235453in}{1.031177in}}%
\pgfpathlineto{\pgfqpoint{1.169817in}{1.038957in}}%
\pgfpathclose%
\pgfusepath{fill}%
\end{pgfscope}%
\begin{pgfscope}%
\pgfpathrectangle{\pgfqpoint{0.100000in}{0.100000in}}{\pgfqpoint{3.007045in}{1.925000in}}%
\pgfusepath{clip}%
\pgfsetbuttcap%
\pgfsetmiterjoin%
\definecolor{currentfill}{rgb}{0.535686,0.746082,0.864252}%
\pgfsetfillcolor{currentfill}%
\pgfsetlinewidth{0.000000pt}%
\definecolor{currentstroke}{rgb}{0.000000,0.000000,0.000000}%
\pgfsetstrokecolor{currentstroke}%
\pgfsetstrokeopacity{0.000000}%
\pgfsetdash{}{0pt}%
\pgfpathmoveto{\pgfqpoint{1.853670in}{1.237862in}}%
\pgfpathlineto{\pgfqpoint{1.844539in}{1.237709in}}%
\pgfpathlineto{\pgfqpoint{1.843908in}{1.275998in}}%
\pgfpathlineto{\pgfqpoint{1.855298in}{1.276147in}}%
\pgfpathlineto{\pgfqpoint{1.866638in}{1.276378in}}%
\pgfpathlineto{\pgfqpoint{1.867080in}{1.259096in}}%
\pgfpathlineto{\pgfqpoint{1.889907in}{1.259660in}}%
\pgfpathlineto{\pgfqpoint{1.890498in}{1.239370in}}%
\pgfpathlineto{\pgfqpoint{1.886676in}{1.239179in}}%
\pgfpathlineto{\pgfqpoint{1.853670in}{1.237862in}}%
\pgfpathclose%
\pgfusepath{fill}%
\end{pgfscope}%
\begin{pgfscope}%
\pgfpathrectangle{\pgfqpoint{0.100000in}{0.100000in}}{\pgfqpoint{3.007045in}{1.925000in}}%
\pgfusepath{clip}%
\pgfsetbuttcap%
\pgfsetmiterjoin%
\definecolor{currentfill}{rgb}{0.435294,0.690965,0.842599}%
\pgfsetfillcolor{currentfill}%
\pgfsetlinewidth{0.000000pt}%
\definecolor{currentstroke}{rgb}{0.000000,0.000000,0.000000}%
\pgfsetstrokecolor{currentstroke}%
\pgfsetstrokeopacity{0.000000}%
\pgfsetdash{}{0pt}%
\pgfpathmoveto{\pgfqpoint{1.638148in}{0.525132in}}%
\pgfpathlineto{\pgfqpoint{1.621659in}{0.515303in}}%
\pgfpathlineto{\pgfqpoint{1.619645in}{0.516644in}}%
\pgfpathlineto{\pgfqpoint{1.615672in}{0.518019in}}%
\pgfpathlineto{\pgfqpoint{1.606065in}{0.530962in}}%
\pgfpathlineto{\pgfqpoint{1.598646in}{0.524138in}}%
\pgfpathlineto{\pgfqpoint{1.597051in}{0.530605in}}%
\pgfpathlineto{\pgfqpoint{1.582069in}{0.543767in}}%
\pgfpathlineto{\pgfqpoint{1.575104in}{0.537324in}}%
\pgfpathlineto{\pgfqpoint{1.565594in}{0.551126in}}%
\pgfpathlineto{\pgfqpoint{1.566327in}{0.575160in}}%
\pgfpathlineto{\pgfqpoint{1.576662in}{0.574910in}}%
\pgfpathlineto{\pgfqpoint{1.592881in}{0.569267in}}%
\pgfpathlineto{\pgfqpoint{1.593576in}{0.574175in}}%
\pgfpathlineto{\pgfqpoint{1.602998in}{0.592878in}}%
\pgfpathlineto{\pgfqpoint{1.610932in}{0.600948in}}%
\pgfpathlineto{\pgfqpoint{1.622469in}{0.597909in}}%
\pgfpathlineto{\pgfqpoint{1.639929in}{0.589583in}}%
\pgfpathlineto{\pgfqpoint{1.643382in}{0.598629in}}%
\pgfpathlineto{\pgfqpoint{1.654308in}{0.604797in}}%
\pgfpathlineto{\pgfqpoint{1.668209in}{0.612458in}}%
\pgfpathlineto{\pgfqpoint{1.670600in}{0.609460in}}%
\pgfpathlineto{\pgfqpoint{1.672091in}{0.600546in}}%
\pgfpathlineto{\pgfqpoint{1.678179in}{0.593593in}}%
\pgfpathlineto{\pgfqpoint{1.679583in}{0.587280in}}%
\pgfpathlineto{\pgfqpoint{1.659758in}{0.576115in}}%
\pgfpathlineto{\pgfqpoint{1.662849in}{0.573743in}}%
\pgfpathlineto{\pgfqpoint{1.665066in}{0.567789in}}%
\pgfpathlineto{\pgfqpoint{1.672046in}{0.560070in}}%
\pgfpathlineto{\pgfqpoint{1.677967in}{0.558409in}}%
\pgfpathlineto{\pgfqpoint{1.672597in}{0.554943in}}%
\pgfpathlineto{\pgfqpoint{1.668934in}{0.549504in}}%
\pgfpathlineto{\pgfqpoint{1.665918in}{0.546597in}}%
\pgfpathlineto{\pgfqpoint{1.655528in}{0.542512in}}%
\pgfpathlineto{\pgfqpoint{1.638148in}{0.525132in}}%
\pgfpathclose%
\pgfusepath{fill}%
\end{pgfscope}%
\begin{pgfscope}%
\pgfpathrectangle{\pgfqpoint{0.100000in}{0.100000in}}{\pgfqpoint{3.007045in}{1.925000in}}%
\pgfusepath{clip}%
\pgfsetbuttcap%
\pgfsetmiterjoin%
\definecolor{currentfill}{rgb}{0.311157,0.608289,0.796894}%
\pgfsetfillcolor{currentfill}%
\pgfsetlinewidth{0.000000pt}%
\definecolor{currentstroke}{rgb}{0.000000,0.000000,0.000000}%
\pgfsetstrokecolor{currentstroke}%
\pgfsetstrokeopacity{0.000000}%
\pgfsetdash{}{0pt}%
\pgfpathmoveto{\pgfqpoint{1.923673in}{0.668219in}}%
\pgfpathlineto{\pgfqpoint{1.922946in}{0.694307in}}%
\pgfpathlineto{\pgfqpoint{1.916979in}{0.699909in}}%
\pgfpathlineto{\pgfqpoint{1.916505in}{0.717286in}}%
\pgfpathlineto{\pgfqpoint{1.910744in}{0.717166in}}%
\pgfpathlineto{\pgfqpoint{1.906378in}{0.722687in}}%
\pgfpathlineto{\pgfqpoint{1.899049in}{0.722750in}}%
\pgfpathlineto{\pgfqpoint{1.898651in}{0.739626in}}%
\pgfpathlineto{\pgfqpoint{1.935072in}{0.740260in}}%
\pgfpathlineto{\pgfqpoint{1.968862in}{0.741315in}}%
\pgfpathlineto{\pgfqpoint{1.970228in}{0.741380in}}%
\pgfpathlineto{\pgfqpoint{1.967255in}{0.733621in}}%
\pgfpathlineto{\pgfqpoint{1.964704in}{0.732579in}}%
\pgfpathlineto{\pgfqpoint{1.959700in}{0.720849in}}%
\pgfpathlineto{\pgfqpoint{1.962387in}{0.712972in}}%
\pgfpathlineto{\pgfqpoint{1.970799in}{0.713279in}}%
\pgfpathlineto{\pgfqpoint{1.970224in}{0.710328in}}%
\pgfpathlineto{\pgfqpoint{1.968808in}{0.707431in}}%
\pgfpathlineto{\pgfqpoint{1.970011in}{0.695703in}}%
\pgfpathlineto{\pgfqpoint{1.966832in}{0.693993in}}%
\pgfpathlineto{\pgfqpoint{1.969121in}{0.688091in}}%
\pgfpathlineto{\pgfqpoint{1.969490in}{0.679240in}}%
\pgfpathlineto{\pgfqpoint{1.965268in}{0.666642in}}%
\pgfpathlineto{\pgfqpoint{1.963857in}{0.672381in}}%
\pgfpathlineto{\pgfqpoint{1.953769in}{0.665525in}}%
\pgfpathlineto{\pgfqpoint{1.949500in}{0.671101in}}%
\pgfpathlineto{\pgfqpoint{1.946656in}{0.668963in}}%
\pgfpathlineto{\pgfqpoint{1.923673in}{0.668219in}}%
\pgfpathclose%
\pgfusepath{fill}%
\end{pgfscope}%
\begin{pgfscope}%
\pgfpathrectangle{\pgfqpoint{0.100000in}{0.100000in}}{\pgfqpoint{3.007045in}{1.925000in}}%
\pgfusepath{clip}%
\pgfsetbuttcap%
\pgfsetmiterjoin%
\definecolor{currentfill}{rgb}{0.765398,0.854118,0.933379}%
\pgfsetfillcolor{currentfill}%
\pgfsetlinewidth{0.000000pt}%
\definecolor{currentstroke}{rgb}{0.000000,0.000000,0.000000}%
\pgfsetstrokecolor{currentstroke}%
\pgfsetstrokeopacity{0.000000}%
\pgfsetdash{}{0pt}%
\pgfpathmoveto{\pgfqpoint{2.192935in}{1.614287in}}%
\pgfpathlineto{\pgfqpoint{2.174903in}{1.611101in}}%
\pgfpathlineto{\pgfqpoint{2.172417in}{1.608214in}}%
\pgfpathlineto{\pgfqpoint{2.171847in}{1.600827in}}%
\pgfpathlineto{\pgfqpoint{2.163982in}{1.596790in}}%
\pgfpathlineto{\pgfqpoint{2.160048in}{1.590165in}}%
\pgfpathlineto{\pgfqpoint{2.155554in}{1.591476in}}%
\pgfpathlineto{\pgfqpoint{2.160765in}{1.598791in}}%
\pgfpathlineto{\pgfqpoint{2.163198in}{1.604703in}}%
\pgfpathlineto{\pgfqpoint{2.152294in}{1.599865in}}%
\pgfpathlineto{\pgfqpoint{2.149656in}{1.593559in}}%
\pgfpathlineto{\pgfqpoint{2.144154in}{1.589747in}}%
\pgfpathlineto{\pgfqpoint{2.139521in}{1.593745in}}%
\pgfpathlineto{\pgfqpoint{2.134646in}{1.588277in}}%
\pgfpathlineto{\pgfqpoint{2.131069in}{1.580709in}}%
\pgfpathlineto{\pgfqpoint{2.128131in}{1.580457in}}%
\pgfpathlineto{\pgfqpoint{2.126263in}{1.603472in}}%
\pgfpathlineto{\pgfqpoint{2.123850in}{1.609038in}}%
\pgfpathlineto{\pgfqpoint{2.112389in}{1.608131in}}%
\pgfpathlineto{\pgfqpoint{2.111036in}{1.625398in}}%
\pgfpathlineto{\pgfqpoint{2.088086in}{1.623582in}}%
\pgfpathlineto{\pgfqpoint{2.087227in}{1.635040in}}%
\pgfpathlineto{\pgfqpoint{2.085981in}{1.652184in}}%
\pgfpathlineto{\pgfqpoint{2.091627in}{1.652681in}}%
\pgfpathlineto{\pgfqpoint{2.088768in}{1.658202in}}%
\pgfpathlineto{\pgfqpoint{2.088047in}{1.667928in}}%
\pgfpathlineto{\pgfqpoint{2.094537in}{1.668479in}}%
\pgfpathlineto{\pgfqpoint{2.107303in}{1.662856in}}%
\pgfpathlineto{\pgfqpoint{2.109669in}{1.658235in}}%
\pgfpathlineto{\pgfqpoint{2.122391in}{1.643026in}}%
\pgfpathlineto{\pgfqpoint{2.132464in}{1.643757in}}%
\pgfpathlineto{\pgfqpoint{2.136667in}{1.646574in}}%
\pgfpathlineto{\pgfqpoint{2.143748in}{1.640811in}}%
\pgfpathlineto{\pgfqpoint{2.153377in}{1.649798in}}%
\pgfpathlineto{\pgfqpoint{2.156972in}{1.644308in}}%
\pgfpathlineto{\pgfqpoint{2.161790in}{1.650391in}}%
\pgfpathlineto{\pgfqpoint{2.176831in}{1.658426in}}%
\pgfpathlineto{\pgfqpoint{2.188229in}{1.662007in}}%
\pgfpathlineto{\pgfqpoint{2.192935in}{1.614287in}}%
\pgfpathclose%
\pgfusepath{fill}%
\end{pgfscope}%
\begin{pgfscope}%
\pgfpathrectangle{\pgfqpoint{0.100000in}{0.100000in}}{\pgfqpoint{3.007045in}{1.925000in}}%
\pgfusepath{clip}%
\pgfsetbuttcap%
\pgfsetmiterjoin%
\definecolor{currentfill}{rgb}{0.632526,0.797647,0.886874}%
\pgfsetfillcolor{currentfill}%
\pgfsetlinewidth{0.000000pt}%
\definecolor{currentstroke}{rgb}{0.000000,0.000000,0.000000}%
\pgfsetstrokecolor{currentstroke}%
\pgfsetstrokeopacity{0.000000}%
\pgfsetdash{}{0pt}%
\pgfpathmoveto{\pgfqpoint{2.753651in}{1.370464in}}%
\pgfpathlineto{\pgfqpoint{2.746998in}{1.361954in}}%
\pgfpathlineto{\pgfqpoint{2.738618in}{1.356791in}}%
\pgfpathlineto{\pgfqpoint{2.732190in}{1.364248in}}%
\pgfpathlineto{\pgfqpoint{2.735363in}{1.367506in}}%
\pgfpathlineto{\pgfqpoint{2.727132in}{1.374300in}}%
\pgfpathlineto{\pgfqpoint{2.723613in}{1.374911in}}%
\pgfpathlineto{\pgfqpoint{2.720526in}{1.370514in}}%
\pgfpathlineto{\pgfqpoint{2.723365in}{1.366527in}}%
\pgfpathlineto{\pgfqpoint{2.711768in}{1.357635in}}%
\pgfpathlineto{\pgfqpoint{2.701015in}{1.357842in}}%
\pgfpathlineto{\pgfqpoint{2.698167in}{1.352791in}}%
\pgfpathlineto{\pgfqpoint{2.698167in}{1.347287in}}%
\pgfpathlineto{\pgfqpoint{2.694956in}{1.344848in}}%
\pgfpathlineto{\pgfqpoint{2.692918in}{1.348382in}}%
\pgfpathlineto{\pgfqpoint{2.688600in}{1.346784in}}%
\pgfpathlineto{\pgfqpoint{2.686221in}{1.354519in}}%
\pgfpathlineto{\pgfqpoint{2.679466in}{1.361047in}}%
\pgfpathlineto{\pgfqpoint{2.677192in}{1.366767in}}%
\pgfpathlineto{\pgfqpoint{2.679506in}{1.367379in}}%
\pgfpathlineto{\pgfqpoint{2.686454in}{1.378941in}}%
\pgfpathlineto{\pgfqpoint{2.691171in}{1.379954in}}%
\pgfpathlineto{\pgfqpoint{2.692396in}{1.393752in}}%
\pgfpathlineto{\pgfqpoint{2.691983in}{1.402891in}}%
\pgfpathlineto{\pgfqpoint{2.695962in}{1.404095in}}%
\pgfpathlineto{\pgfqpoint{2.727566in}{1.410162in}}%
\pgfpathlineto{\pgfqpoint{2.721527in}{1.432985in}}%
\pgfpathlineto{\pgfqpoint{2.727454in}{1.434295in}}%
\pgfpathlineto{\pgfqpoint{2.732166in}{1.431244in}}%
\pgfpathlineto{\pgfqpoint{2.733867in}{1.427016in}}%
\pgfpathlineto{\pgfqpoint{2.739778in}{1.426983in}}%
\pgfpathlineto{\pgfqpoint{2.745672in}{1.421591in}}%
\pgfpathlineto{\pgfqpoint{2.747986in}{1.413389in}}%
\pgfpathlineto{\pgfqpoint{2.752320in}{1.406665in}}%
\pgfpathlineto{\pgfqpoint{2.760698in}{1.403338in}}%
\pgfpathlineto{\pgfqpoint{2.765754in}{1.403877in}}%
\pgfpathlineto{\pgfqpoint{2.768627in}{1.400117in}}%
\pgfpathlineto{\pgfqpoint{2.763125in}{1.394125in}}%
\pgfpathlineto{\pgfqpoint{2.762147in}{1.386718in}}%
\pgfpathlineto{\pgfqpoint{2.753651in}{1.370464in}}%
\pgfpathclose%
\pgfusepath{fill}%
\end{pgfscope}%
\begin{pgfscope}%
\pgfpathrectangle{\pgfqpoint{0.100000in}{0.100000in}}{\pgfqpoint{3.007045in}{1.925000in}}%
\pgfusepath{clip}%
\pgfsetbuttcap%
\pgfsetmiterjoin%
\definecolor{currentfill}{rgb}{0.351511,0.635848,0.812641}%
\pgfsetfillcolor{currentfill}%
\pgfsetlinewidth{0.000000pt}%
\definecolor{currentstroke}{rgb}{0.000000,0.000000,0.000000}%
\pgfsetstrokecolor{currentstroke}%
\pgfsetstrokeopacity{0.000000}%
\pgfsetdash{}{0pt}%
\pgfpathmoveto{\pgfqpoint{1.565600in}{1.146405in}}%
\pgfpathlineto{\pgfqpoint{1.594074in}{1.145204in}}%
\pgfpathlineto{\pgfqpoint{1.593668in}{1.127954in}}%
\pgfpathlineto{\pgfqpoint{1.593165in}{1.116517in}}%
\pgfpathlineto{\pgfqpoint{1.536726in}{1.118994in}}%
\pgfpathlineto{\pgfqpoint{1.507527in}{1.120611in}}%
\pgfpathlineto{\pgfqpoint{1.509458in}{1.149192in}}%
\pgfpathlineto{\pgfqpoint{1.565600in}{1.146405in}}%
\pgfpathclose%
\pgfusepath{fill}%
\end{pgfscope}%
\begin{pgfscope}%
\pgfpathrectangle{\pgfqpoint{0.100000in}{0.100000in}}{\pgfqpoint{3.007045in}{1.925000in}}%
\pgfusepath{clip}%
\pgfsetbuttcap%
\pgfsetmiterjoin%
\definecolor{currentfill}{rgb}{0.093272,0.396878,0.673664}%
\pgfsetfillcolor{currentfill}%
\pgfsetlinewidth{0.000000pt}%
\definecolor{currentstroke}{rgb}{0.000000,0.000000,0.000000}%
\pgfsetstrokecolor{currentstroke}%
\pgfsetstrokeopacity{0.000000}%
\pgfsetdash{}{0pt}%
\pgfpathmoveto{\pgfqpoint{0.558440in}{0.376370in}}%
\pgfpathlineto{\pgfqpoint{0.557442in}{0.379760in}}%
\pgfpathlineto{\pgfqpoint{0.559412in}{0.380419in}}%
\pgfpathlineto{\pgfqpoint{0.561194in}{0.379545in}}%
\pgfpathlineto{\pgfqpoint{0.561660in}{0.376234in}}%
\pgfpathlineto{\pgfqpoint{0.560353in}{0.375728in}}%
\pgfpathlineto{\pgfqpoint{0.558440in}{0.376370in}}%
\pgfpathclose%
\pgfusepath{fill}%
\end{pgfscope}%
\begin{pgfscope}%
\pgfpathrectangle{\pgfqpoint{0.100000in}{0.100000in}}{\pgfqpoint{3.007045in}{1.925000in}}%
\pgfusepath{clip}%
\pgfsetbuttcap%
\pgfsetmiterjoin%
\definecolor{currentfill}{rgb}{0.093272,0.396878,0.673664}%
\pgfsetfillcolor{currentfill}%
\pgfsetlinewidth{0.000000pt}%
\definecolor{currentstroke}{rgb}{0.000000,0.000000,0.000000}%
\pgfsetstrokecolor{currentstroke}%
\pgfsetstrokeopacity{0.000000}%
\pgfsetdash{}{0pt}%
\pgfpathmoveto{\pgfqpoint{0.562662in}{0.374546in}}%
\pgfpathlineto{\pgfqpoint{0.563801in}{0.376835in}}%
\pgfpathlineto{\pgfqpoint{0.565086in}{0.376399in}}%
\pgfpathlineto{\pgfqpoint{0.565958in}{0.374725in}}%
\pgfpathlineto{\pgfqpoint{0.562662in}{0.374546in}}%
\pgfpathclose%
\pgfusepath{fill}%
\end{pgfscope}%
\begin{pgfscope}%
\pgfpathrectangle{\pgfqpoint{0.100000in}{0.100000in}}{\pgfqpoint{3.007045in}{1.925000in}}%
\pgfusepath{clip}%
\pgfsetbuttcap%
\pgfsetmiterjoin%
\definecolor{currentfill}{rgb}{0.093272,0.396878,0.673664}%
\pgfsetfillcolor{currentfill}%
\pgfsetlinewidth{0.000000pt}%
\definecolor{currentstroke}{rgb}{0.000000,0.000000,0.000000}%
\pgfsetstrokecolor{currentstroke}%
\pgfsetstrokeopacity{0.000000}%
\pgfsetdash{}{0pt}%
\pgfpathmoveto{\pgfqpoint{0.630180in}{0.347609in}}%
\pgfpathlineto{\pgfqpoint{0.628737in}{0.345721in}}%
\pgfpathlineto{\pgfqpoint{0.627895in}{0.348703in}}%
\pgfpathlineto{\pgfqpoint{0.625581in}{0.348130in}}%
\pgfpathlineto{\pgfqpoint{0.628285in}{0.352021in}}%
\pgfpathlineto{\pgfqpoint{0.631284in}{0.352146in}}%
\pgfpathlineto{\pgfqpoint{0.630506in}{0.350620in}}%
\pgfpathlineto{\pgfqpoint{0.631605in}{0.349217in}}%
\pgfpathlineto{\pgfqpoint{0.634087in}{0.348313in}}%
\pgfpathlineto{\pgfqpoint{0.632370in}{0.345968in}}%
\pgfpathlineto{\pgfqpoint{0.631280in}{0.347967in}}%
\pgfpathlineto{\pgfqpoint{0.630180in}{0.347609in}}%
\pgfpathclose%
\pgfusepath{fill}%
\end{pgfscope}%
\begin{pgfscope}%
\pgfpathrectangle{\pgfqpoint{0.100000in}{0.100000in}}{\pgfqpoint{3.007045in}{1.925000in}}%
\pgfusepath{clip}%
\pgfsetbuttcap%
\pgfsetmiterjoin%
\definecolor{currentfill}{rgb}{0.093272,0.396878,0.673664}%
\pgfsetfillcolor{currentfill}%
\pgfsetlinewidth{0.000000pt}%
\definecolor{currentstroke}{rgb}{0.000000,0.000000,0.000000}%
\pgfsetstrokecolor{currentstroke}%
\pgfsetstrokeopacity{0.000000}%
\pgfsetdash{}{0pt}%
\pgfpathmoveto{\pgfqpoint{0.606094in}{0.356546in}}%
\pgfpathlineto{\pgfqpoint{0.606285in}{0.359081in}}%
\pgfpathlineto{\pgfqpoint{0.608596in}{0.359253in}}%
\pgfpathlineto{\pgfqpoint{0.608292in}{0.357179in}}%
\pgfpathlineto{\pgfqpoint{0.606094in}{0.356546in}}%
\pgfpathclose%
\pgfusepath{fill}%
\end{pgfscope}%
\begin{pgfscope}%
\pgfpathrectangle{\pgfqpoint{0.100000in}{0.100000in}}{\pgfqpoint{3.007045in}{1.925000in}}%
\pgfusepath{clip}%
\pgfsetbuttcap%
\pgfsetmiterjoin%
\definecolor{currentfill}{rgb}{0.093272,0.396878,0.673664}%
\pgfsetfillcolor{currentfill}%
\pgfsetlinewidth{0.000000pt}%
\definecolor{currentstroke}{rgb}{0.000000,0.000000,0.000000}%
\pgfsetstrokecolor{currentstroke}%
\pgfsetstrokeopacity{0.000000}%
\pgfsetdash{}{0pt}%
\pgfpathmoveto{\pgfqpoint{0.628222in}{0.338692in}}%
\pgfpathlineto{\pgfqpoint{0.634226in}{0.342115in}}%
\pgfpathlineto{\pgfqpoint{0.635742in}{0.340232in}}%
\pgfpathlineto{\pgfqpoint{0.633794in}{0.339431in}}%
\pgfpathlineto{\pgfqpoint{0.628222in}{0.338692in}}%
\pgfpathclose%
\pgfusepath{fill}%
\end{pgfscope}%
\begin{pgfscope}%
\pgfpathrectangle{\pgfqpoint{0.100000in}{0.100000in}}{\pgfqpoint{3.007045in}{1.925000in}}%
\pgfusepath{clip}%
\pgfsetbuttcap%
\pgfsetmiterjoin%
\definecolor{currentfill}{rgb}{0.093272,0.396878,0.673664}%
\pgfsetfillcolor{currentfill}%
\pgfsetlinewidth{0.000000pt}%
\definecolor{currentstroke}{rgb}{0.000000,0.000000,0.000000}%
\pgfsetstrokecolor{currentstroke}%
\pgfsetstrokeopacity{0.000000}%
\pgfsetdash{}{0pt}%
\pgfpathmoveto{\pgfqpoint{0.646196in}{0.346313in}}%
\pgfpathlineto{\pgfqpoint{0.645147in}{0.347754in}}%
\pgfpathlineto{\pgfqpoint{0.648533in}{0.349423in}}%
\pgfpathlineto{\pgfqpoint{0.647621in}{0.351703in}}%
\pgfpathlineto{\pgfqpoint{0.645804in}{0.351697in}}%
\pgfpathlineto{\pgfqpoint{0.644687in}{0.352622in}}%
\pgfpathlineto{\pgfqpoint{0.640184in}{0.353124in}}%
\pgfpathlineto{\pgfqpoint{0.636431in}{0.353197in}}%
\pgfpathlineto{\pgfqpoint{0.634892in}{0.354431in}}%
\pgfpathlineto{\pgfqpoint{0.633377in}{0.353497in}}%
\pgfpathlineto{\pgfqpoint{0.631296in}{0.353562in}}%
\pgfpathlineto{\pgfqpoint{0.631571in}{0.355560in}}%
\pgfpathlineto{\pgfqpoint{0.627061in}{0.355667in}}%
\pgfpathlineto{\pgfqpoint{0.624673in}{0.356334in}}%
\pgfpathlineto{\pgfqpoint{0.622641in}{0.358315in}}%
\pgfpathlineto{\pgfqpoint{0.624401in}{0.359859in}}%
\pgfpathlineto{\pgfqpoint{0.627993in}{0.361739in}}%
\pgfpathlineto{\pgfqpoint{0.626212in}{0.363517in}}%
\pgfpathlineto{\pgfqpoint{0.624339in}{0.363616in}}%
\pgfpathlineto{\pgfqpoint{0.620521in}{0.360581in}}%
\pgfpathlineto{\pgfqpoint{0.617999in}{0.360096in}}%
\pgfpathlineto{\pgfqpoint{0.615778in}{0.360562in}}%
\pgfpathlineto{\pgfqpoint{0.614484in}{0.358255in}}%
\pgfpathlineto{\pgfqpoint{0.612634in}{0.358313in}}%
\pgfpathlineto{\pgfqpoint{0.608448in}{0.361014in}}%
\pgfpathlineto{\pgfqpoint{0.608279in}{0.362605in}}%
\pgfpathlineto{\pgfqpoint{0.609968in}{0.362762in}}%
\pgfpathlineto{\pgfqpoint{0.610410in}{0.364673in}}%
\pgfpathlineto{\pgfqpoint{0.610099in}{0.367347in}}%
\pgfpathlineto{\pgfqpoint{0.606743in}{0.363415in}}%
\pgfpathlineto{\pgfqpoint{0.606556in}{0.361958in}}%
\pgfpathlineto{\pgfqpoint{0.603910in}{0.362284in}}%
\pgfpathlineto{\pgfqpoint{0.602151in}{0.363752in}}%
\pgfpathlineto{\pgfqpoint{0.603264in}{0.366602in}}%
\pgfpathlineto{\pgfqpoint{0.602235in}{0.369208in}}%
\pgfpathlineto{\pgfqpoint{0.601172in}{0.368904in}}%
\pgfpathlineto{\pgfqpoint{0.600989in}{0.365044in}}%
\pgfpathlineto{\pgfqpoint{0.595049in}{0.365340in}}%
\pgfpathlineto{\pgfqpoint{0.596365in}{0.363673in}}%
\pgfpathlineto{\pgfqpoint{0.594450in}{0.363174in}}%
\pgfpathlineto{\pgfqpoint{0.593766in}{0.364861in}}%
\pgfpathlineto{\pgfqpoint{0.592021in}{0.364091in}}%
\pgfpathlineto{\pgfqpoint{0.589288in}{0.365861in}}%
\pgfpathlineto{\pgfqpoint{0.585892in}{0.369053in}}%
\pgfpathlineto{\pgfqpoint{0.581946in}{0.370113in}}%
\pgfpathlineto{\pgfqpoint{0.578510in}{0.369321in}}%
\pgfpathlineto{\pgfqpoint{0.576281in}{0.370552in}}%
\pgfpathlineto{\pgfqpoint{0.575149in}{0.374376in}}%
\pgfpathlineto{\pgfqpoint{0.576135in}{0.376449in}}%
\pgfpathlineto{\pgfqpoint{0.580707in}{0.375972in}}%
\pgfpathlineto{\pgfqpoint{0.586276in}{0.378070in}}%
\pgfpathlineto{\pgfqpoint{0.587415in}{0.376136in}}%
\pgfpathlineto{\pgfqpoint{0.589357in}{0.375160in}}%
\pgfpathlineto{\pgfqpoint{0.590661in}{0.375506in}}%
\pgfpathlineto{\pgfqpoint{0.595051in}{0.373915in}}%
\pgfpathlineto{\pgfqpoint{0.597649in}{0.371586in}}%
\pgfpathlineto{\pgfqpoint{0.596738in}{0.368224in}}%
\pgfpathlineto{\pgfqpoint{0.598337in}{0.367723in}}%
\pgfpathlineto{\pgfqpoint{0.600669in}{0.370774in}}%
\pgfpathlineto{\pgfqpoint{0.604276in}{0.369985in}}%
\pgfpathlineto{\pgfqpoint{0.605343in}{0.367909in}}%
\pgfpathlineto{\pgfqpoint{0.607337in}{0.368055in}}%
\pgfpathlineto{\pgfqpoint{0.620695in}{0.370601in}}%
\pgfpathlineto{\pgfqpoint{0.625712in}{0.370166in}}%
\pgfpathlineto{\pgfqpoint{0.628012in}{0.370317in}}%
\pgfpathlineto{\pgfqpoint{0.633395in}{0.367875in}}%
\pgfpathlineto{\pgfqpoint{0.636602in}{0.364776in}}%
\pgfpathlineto{\pgfqpoint{0.635557in}{0.363488in}}%
\pgfpathlineto{\pgfqpoint{0.635523in}{0.359748in}}%
\pgfpathlineto{\pgfqpoint{0.637405in}{0.360249in}}%
\pgfpathlineto{\pgfqpoint{0.640275in}{0.358605in}}%
\pgfpathlineto{\pgfqpoint{0.639530in}{0.356749in}}%
\pgfpathlineto{\pgfqpoint{0.640742in}{0.355326in}}%
\pgfpathlineto{\pgfqpoint{0.642306in}{0.356377in}}%
\pgfpathlineto{\pgfqpoint{0.641680in}{0.358719in}}%
\pgfpathlineto{\pgfqpoint{0.640362in}{0.360039in}}%
\pgfpathlineto{\pgfqpoint{0.641503in}{0.361170in}}%
\pgfpathlineto{\pgfqpoint{0.644772in}{0.362563in}}%
\pgfpathlineto{\pgfqpoint{0.646911in}{0.363995in}}%
\pgfpathlineto{\pgfqpoint{0.651915in}{0.364863in}}%
\pgfpathlineto{\pgfqpoint{0.654420in}{0.364388in}}%
\pgfpathlineto{\pgfqpoint{0.656995in}{0.364705in}}%
\pgfpathlineto{\pgfqpoint{0.659517in}{0.364305in}}%
\pgfpathlineto{\pgfqpoint{0.665903in}{0.362144in}}%
\pgfpathlineto{\pgfqpoint{0.665627in}{0.363168in}}%
\pgfpathlineto{\pgfqpoint{0.669120in}{0.362795in}}%
\pgfpathlineto{\pgfqpoint{0.670704in}{0.361448in}}%
\pgfpathlineto{\pgfqpoint{0.668605in}{0.361721in}}%
\pgfpathlineto{\pgfqpoint{0.667310in}{0.360198in}}%
\pgfpathlineto{\pgfqpoint{0.664908in}{0.362318in}}%
\pgfpathlineto{\pgfqpoint{0.663515in}{0.360688in}}%
\pgfpathlineto{\pgfqpoint{0.660386in}{0.363372in}}%
\pgfpathlineto{\pgfqpoint{0.659049in}{0.361833in}}%
\pgfpathlineto{\pgfqpoint{0.655995in}{0.364392in}}%
\pgfpathlineto{\pgfqpoint{0.653338in}{0.361411in}}%
\pgfpathlineto{\pgfqpoint{0.654252in}{0.360591in}}%
\pgfpathlineto{\pgfqpoint{0.649075in}{0.354518in}}%
\pgfpathlineto{\pgfqpoint{0.647002in}{0.353387in}}%
\pgfpathlineto{\pgfqpoint{0.650028in}{0.350774in}}%
\pgfpathlineto{\pgfqpoint{0.646196in}{0.346313in}}%
\pgfpathclose%
\pgfusepath{fill}%
\end{pgfscope}%
\begin{pgfscope}%
\pgfpathrectangle{\pgfqpoint{0.100000in}{0.100000in}}{\pgfqpoint{3.007045in}{1.925000in}}%
\pgfusepath{clip}%
\pgfsetbuttcap%
\pgfsetmiterjoin%
\definecolor{currentfill}{rgb}{0.296025,0.597955,0.790988}%
\pgfsetfillcolor{currentfill}%
\pgfsetlinewidth{0.000000pt}%
\definecolor{currentstroke}{rgb}{0.000000,0.000000,0.000000}%
\pgfsetstrokecolor{currentstroke}%
\pgfsetstrokeopacity{0.000000}%
\pgfsetdash{}{0pt}%
\pgfpathmoveto{\pgfqpoint{1.588665in}{1.618076in}}%
\pgfpathlineto{\pgfqpoint{1.587550in}{1.595331in}}%
\pgfpathlineto{\pgfqpoint{1.554716in}{1.596972in}}%
\pgfpathlineto{\pgfqpoint{1.518744in}{1.599173in}}%
\pgfpathlineto{\pgfqpoint{1.514131in}{1.607335in}}%
\pgfpathlineto{\pgfqpoint{1.514093in}{1.614101in}}%
\pgfpathlineto{\pgfqpoint{1.517258in}{1.619082in}}%
\pgfpathlineto{\pgfqpoint{1.518046in}{1.627144in}}%
\pgfpathlineto{\pgfqpoint{1.516525in}{1.631523in}}%
\pgfpathlineto{\pgfqpoint{1.519039in}{1.638417in}}%
\pgfpathlineto{\pgfqpoint{1.518209in}{1.642479in}}%
\pgfpathlineto{\pgfqpoint{1.514103in}{1.645425in}}%
\pgfpathlineto{\pgfqpoint{1.540553in}{1.643789in}}%
\pgfpathlineto{\pgfqpoint{1.569331in}{1.642166in}}%
\pgfpathlineto{\pgfqpoint{1.588163in}{1.641202in}}%
\pgfpathlineto{\pgfqpoint{1.587099in}{1.618157in}}%
\pgfpathlineto{\pgfqpoint{1.588665in}{1.618076in}}%
\pgfpathclose%
\pgfusepath{fill}%
\end{pgfscope}%
\begin{pgfscope}%
\pgfpathrectangle{\pgfqpoint{0.100000in}{0.100000in}}{\pgfqpoint{3.007045in}{1.925000in}}%
\pgfusepath{clip}%
\pgfsetbuttcap%
\pgfsetmiterjoin%
\definecolor{currentfill}{rgb}{0.071742,0.369319,0.652134}%
\pgfsetfillcolor{currentfill}%
\pgfsetlinewidth{0.000000pt}%
\definecolor{currentstroke}{rgb}{0.000000,0.000000,0.000000}%
\pgfsetstrokecolor{currentstroke}%
\pgfsetstrokeopacity{0.000000}%
\pgfsetdash{}{0pt}%
\pgfpathmoveto{\pgfqpoint{1.509153in}{1.008633in}}%
\pgfpathlineto{\pgfqpoint{1.507272in}{0.975589in}}%
\pgfpathlineto{\pgfqpoint{1.456990in}{0.978623in}}%
\pgfpathlineto{\pgfqpoint{1.459593in}{1.011387in}}%
\pgfpathlineto{\pgfqpoint{1.453227in}{1.011793in}}%
\pgfpathlineto{\pgfqpoint{1.454880in}{1.037434in}}%
\pgfpathlineto{\pgfqpoint{1.476575in}{1.036016in}}%
\pgfpathlineto{\pgfqpoint{1.476919in}{1.041763in}}%
\pgfpathlineto{\pgfqpoint{1.505407in}{1.040070in}}%
\pgfpathlineto{\pgfqpoint{1.505902in}{1.034236in}}%
\pgfpathlineto{\pgfqpoint{1.504587in}{1.008926in}}%
\pgfpathlineto{\pgfqpoint{1.509153in}{1.008633in}}%
\pgfpathclose%
\pgfusepath{fill}%
\end{pgfscope}%
\begin{pgfscope}%
\pgfpathrectangle{\pgfqpoint{0.100000in}{0.100000in}}{\pgfqpoint{3.007045in}{1.925000in}}%
\pgfusepath{clip}%
\pgfsetbuttcap%
\pgfsetmiterjoin%
\definecolor{currentfill}{rgb}{0.265759,0.577286,0.779177}%
\pgfsetfillcolor{currentfill}%
\pgfsetlinewidth{0.000000pt}%
\definecolor{currentstroke}{rgb}{0.000000,0.000000,0.000000}%
\pgfsetstrokecolor{currentstroke}%
\pgfsetstrokeopacity{0.000000}%
\pgfsetdash{}{0pt}%
\pgfpathmoveto{\pgfqpoint{1.753660in}{1.413597in}}%
\pgfpathlineto{\pgfqpoint{1.730861in}{1.413920in}}%
\pgfpathlineto{\pgfqpoint{1.731061in}{1.429926in}}%
\pgfpathlineto{\pgfqpoint{1.721870in}{1.430045in}}%
\pgfpathlineto{\pgfqpoint{1.722176in}{1.453096in}}%
\pgfpathlineto{\pgfqpoint{1.721890in}{1.476069in}}%
\pgfpathlineto{\pgfqpoint{1.744188in}{1.475811in}}%
\pgfpathlineto{\pgfqpoint{1.767169in}{1.475679in}}%
\pgfpathlineto{\pgfqpoint{1.767569in}{1.469929in}}%
\pgfpathlineto{\pgfqpoint{1.778980in}{1.469921in}}%
\pgfpathlineto{\pgfqpoint{1.779266in}{1.452736in}}%
\pgfpathlineto{\pgfqpoint{1.779303in}{1.429779in}}%
\pgfpathlineto{\pgfqpoint{1.776405in}{1.429765in}}%
\pgfpathlineto{\pgfqpoint{1.753741in}{1.429789in}}%
\pgfpathlineto{\pgfqpoint{1.753660in}{1.413597in}}%
\pgfpathclose%
\pgfusepath{fill}%
\end{pgfscope}%
\begin{pgfscope}%
\pgfpathrectangle{\pgfqpoint{0.100000in}{0.100000in}}{\pgfqpoint{3.007045in}{1.925000in}}%
\pgfusepath{clip}%
\pgfsetbuttcap%
\pgfsetmiterjoin%
\definecolor{currentfill}{rgb}{0.535686,0.746082,0.864252}%
\pgfsetfillcolor{currentfill}%
\pgfsetlinewidth{0.000000pt}%
\definecolor{currentstroke}{rgb}{0.000000,0.000000,0.000000}%
\pgfsetstrokecolor{currentstroke}%
\pgfsetstrokeopacity{0.000000}%
\pgfsetdash{}{0pt}%
\pgfpathmoveto{\pgfqpoint{1.798792in}{1.236593in}}%
\pgfpathlineto{\pgfqpoint{1.790743in}{1.236578in}}%
\pgfpathlineto{\pgfqpoint{1.734017in}{1.237625in}}%
\pgfpathlineto{\pgfqpoint{1.732568in}{1.242722in}}%
\pgfpathlineto{\pgfqpoint{1.727958in}{1.247635in}}%
\pgfpathlineto{\pgfqpoint{1.730711in}{1.250739in}}%
\pgfpathlineto{\pgfqpoint{1.731811in}{1.258499in}}%
\pgfpathlineto{\pgfqpoint{1.753153in}{1.258349in}}%
\pgfpathlineto{\pgfqpoint{1.753309in}{1.275315in}}%
\pgfpathlineto{\pgfqpoint{1.764595in}{1.275306in}}%
\pgfpathlineto{\pgfqpoint{1.787185in}{1.275241in}}%
\pgfpathlineto{\pgfqpoint{1.798597in}{1.275256in}}%
\pgfpathlineto{\pgfqpoint{1.798792in}{1.236593in}}%
\pgfpathclose%
\pgfusepath{fill}%
\end{pgfscope}%
\begin{pgfscope}%
\pgfpathrectangle{\pgfqpoint{0.100000in}{0.100000in}}{\pgfqpoint{3.007045in}{1.925000in}}%
\pgfusepath{clip}%
\pgfsetbuttcap%
\pgfsetmiterjoin%
\definecolor{currentfill}{rgb}{0.260715,0.573841,0.777209}%
\pgfsetfillcolor{currentfill}%
\pgfsetlinewidth{0.000000pt}%
\definecolor{currentstroke}{rgb}{0.000000,0.000000,0.000000}%
\pgfsetstrokecolor{currentstroke}%
\pgfsetstrokeopacity{0.000000}%
\pgfsetdash{}{0pt}%
\pgfpathmoveto{\pgfqpoint{1.634574in}{0.311704in}}%
\pgfpathlineto{\pgfqpoint{1.632618in}{0.311751in}}%
\pgfpathlineto{\pgfqpoint{1.628886in}{0.326991in}}%
\pgfpathlineto{\pgfqpoint{1.629356in}{0.349314in}}%
\pgfpathlineto{\pgfqpoint{1.631677in}{0.363791in}}%
\pgfpathlineto{\pgfqpoint{1.647149in}{0.392795in}}%
\pgfpathlineto{\pgfqpoint{1.651127in}{0.395908in}}%
\pgfpathlineto{\pgfqpoint{1.655482in}{0.405989in}}%
\pgfpathlineto{\pgfqpoint{1.661560in}{0.414117in}}%
\pgfpathlineto{\pgfqpoint{1.664160in}{0.412781in}}%
\pgfpathlineto{\pgfqpoint{1.662731in}{0.409312in}}%
\pgfpathlineto{\pgfqpoint{1.652282in}{0.397408in}}%
\pgfpathlineto{\pgfqpoint{1.643553in}{0.383920in}}%
\pgfpathlineto{\pgfqpoint{1.637879in}{0.372867in}}%
\pgfpathlineto{\pgfqpoint{1.633473in}{0.362042in}}%
\pgfpathlineto{\pgfqpoint{1.630879in}{0.352084in}}%
\pgfpathlineto{\pgfqpoint{1.629805in}{0.336635in}}%
\pgfpathlineto{\pgfqpoint{1.631675in}{0.322573in}}%
\pgfpathlineto{\pgfqpoint{1.634574in}{0.311704in}}%
\pgfpathclose%
\pgfusepath{fill}%
\end{pgfscope}%
\begin{pgfscope}%
\pgfpathrectangle{\pgfqpoint{0.100000in}{0.100000in}}{\pgfqpoint{3.007045in}{1.925000in}}%
\pgfusepath{clip}%
\pgfsetbuttcap%
\pgfsetmiterjoin%
\definecolor{currentfill}{rgb}{0.260715,0.573841,0.777209}%
\pgfsetfillcolor{currentfill}%
\pgfsetlinewidth{0.000000pt}%
\definecolor{currentstroke}{rgb}{0.000000,0.000000,0.000000}%
\pgfsetstrokecolor{currentstroke}%
\pgfsetstrokeopacity{0.000000}%
\pgfsetdash{}{0pt}%
\pgfpathmoveto{\pgfqpoint{1.625397in}{0.311924in}}%
\pgfpathlineto{\pgfqpoint{1.599953in}{0.312537in}}%
\pgfpathlineto{\pgfqpoint{1.592550in}{0.313976in}}%
\pgfpathlineto{\pgfqpoint{1.592904in}{0.325092in}}%
\pgfpathlineto{\pgfqpoint{1.566652in}{0.326194in}}%
\pgfpathlineto{\pgfqpoint{1.567623in}{0.344460in}}%
\pgfpathlineto{\pgfqpoint{1.564709in}{0.344532in}}%
\pgfpathlineto{\pgfqpoint{1.563594in}{0.356706in}}%
\pgfpathlineto{\pgfqpoint{1.560356in}{0.365181in}}%
\pgfpathlineto{\pgfqpoint{1.545690in}{0.365433in}}%
\pgfpathlineto{\pgfqpoint{1.547367in}{0.412624in}}%
\pgfpathlineto{\pgfqpoint{1.601772in}{0.410582in}}%
\pgfpathlineto{\pgfqpoint{1.600656in}{0.414463in}}%
\pgfpathlineto{\pgfqpoint{1.605884in}{0.418493in}}%
\pgfpathlineto{\pgfqpoint{1.620569in}{0.414900in}}%
\pgfpathlineto{\pgfqpoint{1.632363in}{0.431928in}}%
\pgfpathlineto{\pgfqpoint{1.645296in}{0.442654in}}%
\pgfpathlineto{\pgfqpoint{1.651860in}{0.442260in}}%
\pgfpathlineto{\pgfqpoint{1.655495in}{0.438831in}}%
\pgfpathlineto{\pgfqpoint{1.661661in}{0.438978in}}%
\pgfpathlineto{\pgfqpoint{1.667536in}{0.431953in}}%
\pgfpathlineto{\pgfqpoint{1.666217in}{0.423659in}}%
\pgfpathlineto{\pgfqpoint{1.665504in}{0.419888in}}%
\pgfpathlineto{\pgfqpoint{1.658131in}{0.413570in}}%
\pgfpathlineto{\pgfqpoint{1.654267in}{0.413508in}}%
\pgfpathlineto{\pgfqpoint{1.651083in}{0.406992in}}%
\pgfpathlineto{\pgfqpoint{1.642502in}{0.393804in}}%
\pgfpathlineto{\pgfqpoint{1.637586in}{0.397797in}}%
\pgfpathlineto{\pgfqpoint{1.631800in}{0.395290in}}%
\pgfpathlineto{\pgfqpoint{1.632028in}{0.388063in}}%
\pgfpathlineto{\pgfqpoint{1.638379in}{0.385411in}}%
\pgfpathlineto{\pgfqpoint{1.634179in}{0.376996in}}%
\pgfpathlineto{\pgfqpoint{1.628882in}{0.361108in}}%
\pgfpathlineto{\pgfqpoint{1.626077in}{0.346812in}}%
\pgfpathlineto{\pgfqpoint{1.619237in}{0.338192in}}%
\pgfpathlineto{\pgfqpoint{1.618817in}{0.328087in}}%
\pgfpathlineto{\pgfqpoint{1.623729in}{0.322672in}}%
\pgfpathlineto{\pgfqpoint{1.625397in}{0.311924in}}%
\pgfpathclose%
\pgfusepath{fill}%
\end{pgfscope}%
\begin{pgfscope}%
\pgfpathrectangle{\pgfqpoint{0.100000in}{0.100000in}}{\pgfqpoint{3.007045in}{1.925000in}}%
\pgfusepath{clip}%
\pgfsetbuttcap%
\pgfsetmiterjoin%
\definecolor{currentfill}{rgb}{0.290980,0.594510,0.789020}%
\pgfsetfillcolor{currentfill}%
\pgfsetlinewidth{0.000000pt}%
\definecolor{currentstroke}{rgb}{0.000000,0.000000,0.000000}%
\pgfsetstrokecolor{currentstroke}%
\pgfsetstrokeopacity{0.000000}%
\pgfsetdash{}{0pt}%
\pgfpathmoveto{\pgfqpoint{1.703058in}{1.453389in}}%
\pgfpathlineto{\pgfqpoint{1.682306in}{1.453680in}}%
\pgfpathlineto{\pgfqpoint{1.682752in}{1.476661in}}%
\pgfpathlineto{\pgfqpoint{1.637239in}{1.477803in}}%
\pgfpathlineto{\pgfqpoint{1.637041in}{1.477807in}}%
\pgfpathlineto{\pgfqpoint{1.637786in}{1.500890in}}%
\pgfpathlineto{\pgfqpoint{1.671967in}{1.499872in}}%
\pgfpathlineto{\pgfqpoint{1.703851in}{1.499277in}}%
\pgfpathlineto{\pgfqpoint{1.703058in}{1.453389in}}%
\pgfpathclose%
\pgfusepath{fill}%
\end{pgfscope}%
\begin{pgfscope}%
\pgfpathrectangle{\pgfqpoint{0.100000in}{0.100000in}}{\pgfqpoint{3.007045in}{1.925000in}}%
\pgfusepath{clip}%
\pgfsetbuttcap%
\pgfsetmiterjoin%
\definecolor{currentfill}{rgb}{0.585882,0.773641,0.875079}%
\pgfsetfillcolor{currentfill}%
\pgfsetlinewidth{0.000000pt}%
\definecolor{currentstroke}{rgb}{0.000000,0.000000,0.000000}%
\pgfsetstrokecolor{currentstroke}%
\pgfsetstrokeopacity{0.000000}%
\pgfsetdash{}{0pt}%
\pgfpathmoveto{\pgfqpoint{1.944456in}{0.907639in}}%
\pgfpathlineto{\pgfqpoint{1.944537in}{0.902457in}}%
\pgfpathlineto{\pgfqpoint{1.941745in}{0.895933in}}%
\pgfpathlineto{\pgfqpoint{1.927775in}{0.895762in}}%
\pgfpathlineto{\pgfqpoint{1.907859in}{0.895543in}}%
\pgfpathlineto{\pgfqpoint{1.907800in}{0.901293in}}%
\pgfpathlineto{\pgfqpoint{1.887849in}{0.901272in}}%
\pgfpathlineto{\pgfqpoint{1.889712in}{0.906997in}}%
\pgfpathlineto{\pgfqpoint{1.889632in}{0.918440in}}%
\pgfpathlineto{\pgfqpoint{1.899200in}{0.918414in}}%
\pgfpathlineto{\pgfqpoint{1.899149in}{0.923195in}}%
\pgfpathlineto{\pgfqpoint{1.910553in}{0.923358in}}%
\pgfpathlineto{\pgfqpoint{1.910330in}{0.935711in}}%
\pgfpathlineto{\pgfqpoint{1.916049in}{0.935756in}}%
\pgfpathlineto{\pgfqpoint{1.915999in}{0.941499in}}%
\pgfpathlineto{\pgfqpoint{1.921721in}{0.941554in}}%
\pgfpathlineto{\pgfqpoint{1.921677in}{0.946416in}}%
\pgfpathlineto{\pgfqpoint{1.929258in}{0.940159in}}%
\pgfpathlineto{\pgfqpoint{1.927752in}{0.935378in}}%
\pgfpathlineto{\pgfqpoint{1.935018in}{0.933586in}}%
\pgfpathlineto{\pgfqpoint{1.940723in}{0.927923in}}%
\pgfpathlineto{\pgfqpoint{1.939696in}{0.925677in}}%
\pgfpathlineto{\pgfqpoint{1.944264in}{0.919360in}}%
\pgfpathlineto{\pgfqpoint{1.944456in}{0.907639in}}%
\pgfpathclose%
\pgfusepath{fill}%
\end{pgfscope}%
\begin{pgfscope}%
\pgfpathrectangle{\pgfqpoint{0.100000in}{0.100000in}}{\pgfqpoint{3.007045in}{1.925000in}}%
\pgfusepath{clip}%
\pgfsetbuttcap%
\pgfsetmiterjoin%
\definecolor{currentfill}{rgb}{0.441569,0.694410,0.843952}%
\pgfsetfillcolor{currentfill}%
\pgfsetlinewidth{0.000000pt}%
\definecolor{currentstroke}{rgb}{0.000000,0.000000,0.000000}%
\pgfsetstrokecolor{currentstroke}%
\pgfsetstrokeopacity{0.000000}%
\pgfsetdash{}{0pt}%
\pgfpathmoveto{\pgfqpoint{2.482424in}{0.851583in}}%
\pgfpathlineto{\pgfqpoint{2.474703in}{0.848311in}}%
\pgfpathlineto{\pgfqpoint{2.477899in}{0.842105in}}%
\pgfpathlineto{\pgfqpoint{2.475134in}{0.834677in}}%
\pgfpathlineto{\pgfqpoint{2.480732in}{0.829446in}}%
\pgfpathlineto{\pgfqpoint{2.479069in}{0.826299in}}%
\pgfpathlineto{\pgfqpoint{2.473536in}{0.831913in}}%
\pgfpathlineto{\pgfqpoint{2.468037in}{0.838774in}}%
\pgfpathlineto{\pgfqpoint{2.452246in}{0.846939in}}%
\pgfpathlineto{\pgfqpoint{2.449961in}{0.851491in}}%
\pgfpathlineto{\pgfqpoint{2.442202in}{0.858557in}}%
\pgfpathlineto{\pgfqpoint{2.439934in}{0.864954in}}%
\pgfpathlineto{\pgfqpoint{2.429221in}{0.879173in}}%
\pgfpathlineto{\pgfqpoint{2.423854in}{0.877561in}}%
\pgfpathlineto{\pgfqpoint{2.420077in}{0.879039in}}%
\pgfpathlineto{\pgfqpoint{2.413604in}{0.885230in}}%
\pgfpathlineto{\pgfqpoint{2.410336in}{0.885056in}}%
\pgfpathlineto{\pgfqpoint{2.403995in}{0.889257in}}%
\pgfpathlineto{\pgfqpoint{2.402990in}{0.891059in}}%
\pgfpathlineto{\pgfqpoint{2.403903in}{0.896059in}}%
\pgfpathlineto{\pgfqpoint{2.408265in}{0.902647in}}%
\pgfpathlineto{\pgfqpoint{2.413500in}{0.907511in}}%
\pgfpathlineto{\pgfqpoint{2.413670in}{0.911416in}}%
\pgfpathlineto{\pgfqpoint{2.418790in}{0.913875in}}%
\pgfpathlineto{\pgfqpoint{2.435961in}{0.922177in}}%
\pgfpathlineto{\pgfqpoint{2.452389in}{0.929337in}}%
\pgfpathlineto{\pgfqpoint{2.459680in}{0.930712in}}%
\pgfpathlineto{\pgfqpoint{2.462329in}{0.908061in}}%
\pgfpathlineto{\pgfqpoint{2.476040in}{0.897437in}}%
\pgfpathlineto{\pgfqpoint{2.481602in}{0.894314in}}%
\pgfpathlineto{\pgfqpoint{2.496666in}{0.891862in}}%
\pgfpathlineto{\pgfqpoint{2.490797in}{0.879764in}}%
\pgfpathlineto{\pgfqpoint{2.485317in}{0.874629in}}%
\pgfpathlineto{\pgfqpoint{2.484386in}{0.869099in}}%
\pgfpathlineto{\pgfqpoint{2.486659in}{0.864555in}}%
\pgfpathlineto{\pgfqpoint{2.482424in}{0.851583in}}%
\pgfpathclose%
\pgfusepath{fill}%
\end{pgfscope}%
\begin{pgfscope}%
\pgfpathrectangle{\pgfqpoint{0.100000in}{0.100000in}}{\pgfqpoint{3.007045in}{1.925000in}}%
\pgfusepath{clip}%
\pgfsetbuttcap%
\pgfsetmiterjoin%
\definecolor{currentfill}{rgb}{0.406997,0.673741,0.834295}%
\pgfsetfillcolor{currentfill}%
\pgfsetlinewidth{0.000000pt}%
\definecolor{currentstroke}{rgb}{0.000000,0.000000,0.000000}%
\pgfsetstrokecolor{currentstroke}%
\pgfsetstrokeopacity{0.000000}%
\pgfsetdash{}{0pt}%
\pgfpathmoveto{\pgfqpoint{2.253561in}{1.079316in}}%
\pgfpathlineto{\pgfqpoint{2.249127in}{1.081944in}}%
\pgfpathlineto{\pgfqpoint{2.246327in}{1.087565in}}%
\pgfpathlineto{\pgfqpoint{2.241377in}{1.090992in}}%
\pgfpathlineto{\pgfqpoint{2.229804in}{1.090057in}}%
\pgfpathlineto{\pgfqpoint{2.224931in}{1.092417in}}%
\pgfpathlineto{\pgfqpoint{2.223926in}{1.097599in}}%
\pgfpathlineto{\pgfqpoint{2.219675in}{1.101261in}}%
\pgfpathlineto{\pgfqpoint{2.217176in}{1.095188in}}%
\pgfpathlineto{\pgfqpoint{2.214133in}{1.095845in}}%
\pgfpathlineto{\pgfqpoint{2.213695in}{1.101544in}}%
\pgfpathlineto{\pgfqpoint{2.207958in}{1.101054in}}%
\pgfpathlineto{\pgfqpoint{2.207678in}{1.104917in}}%
\pgfpathlineto{\pgfqpoint{2.202065in}{1.104184in}}%
\pgfpathlineto{\pgfqpoint{2.201173in}{1.112813in}}%
\pgfpathlineto{\pgfqpoint{2.220327in}{1.114564in}}%
\pgfpathlineto{\pgfqpoint{2.218133in}{1.136791in}}%
\pgfpathlineto{\pgfqpoint{2.219718in}{1.138972in}}%
\pgfpathlineto{\pgfqpoint{2.235786in}{1.140573in}}%
\pgfpathlineto{\pgfqpoint{2.239682in}{1.138991in}}%
\pgfpathlineto{\pgfqpoint{2.240875in}{1.128586in}}%
\pgfpathlineto{\pgfqpoint{2.245396in}{1.130991in}}%
\pgfpathlineto{\pgfqpoint{2.256834in}{1.132262in}}%
\pgfpathlineto{\pgfqpoint{2.264234in}{1.131721in}}%
\pgfpathlineto{\pgfqpoint{2.264410in}{1.127604in}}%
\pgfpathlineto{\pgfqpoint{2.272501in}{1.126660in}}%
\pgfpathlineto{\pgfqpoint{2.277544in}{1.133106in}}%
\pgfpathlineto{\pgfqpoint{2.282277in}{1.134362in}}%
\pgfpathlineto{\pgfqpoint{2.285748in}{1.132441in}}%
\pgfpathlineto{\pgfqpoint{2.287382in}{1.127701in}}%
\pgfpathlineto{\pgfqpoint{2.292494in}{1.120777in}}%
\pgfpathlineto{\pgfqpoint{2.288467in}{1.113713in}}%
\pgfpathlineto{\pgfqpoint{2.288667in}{1.106429in}}%
\pgfpathlineto{\pgfqpoint{2.287143in}{1.100325in}}%
\pgfpathlineto{\pgfqpoint{2.281936in}{1.092901in}}%
\pgfpathlineto{\pgfqpoint{2.270164in}{1.089523in}}%
\pgfpathlineto{\pgfqpoint{2.265212in}{1.092380in}}%
\pgfpathlineto{\pgfqpoint{2.258362in}{1.081298in}}%
\pgfpathlineto{\pgfqpoint{2.253561in}{1.079316in}}%
\pgfpathclose%
\pgfusepath{fill}%
\end{pgfscope}%
\begin{pgfscope}%
\pgfpathrectangle{\pgfqpoint{0.100000in}{0.100000in}}{\pgfqpoint{3.007045in}{1.925000in}}%
\pgfusepath{clip}%
\pgfsetbuttcap%
\pgfsetmiterjoin%
\definecolor{currentfill}{rgb}{0.366644,0.646182,0.818547}%
\pgfsetfillcolor{currentfill}%
\pgfsetlinewidth{0.000000pt}%
\definecolor{currentstroke}{rgb}{0.000000,0.000000,0.000000}%
\pgfsetstrokecolor{currentstroke}%
\pgfsetstrokeopacity{0.000000}%
\pgfsetdash{}{0pt}%
\pgfpathmoveto{\pgfqpoint{1.610638in}{1.328045in}}%
\pgfpathlineto{\pgfqpoint{1.610215in}{1.316599in}}%
\pgfpathlineto{\pgfqpoint{1.587433in}{1.317426in}}%
\pgfpathlineto{\pgfqpoint{1.588274in}{1.340363in}}%
\pgfpathlineto{\pgfqpoint{1.565129in}{1.341345in}}%
\pgfpathlineto{\pgfqpoint{1.566978in}{1.401393in}}%
\pgfpathlineto{\pgfqpoint{1.603438in}{1.399865in}}%
\pgfpathlineto{\pgfqpoint{1.604815in}{1.396442in}}%
\pgfpathlineto{\pgfqpoint{1.612271in}{1.391873in}}%
\pgfpathlineto{\pgfqpoint{1.611668in}{1.362462in}}%
\pgfpathlineto{\pgfqpoint{1.610638in}{1.328045in}}%
\pgfpathclose%
\pgfusepath{fill}%
\end{pgfscope}%
\begin{pgfscope}%
\pgfpathrectangle{\pgfqpoint{0.100000in}{0.100000in}}{\pgfqpoint{3.007045in}{1.925000in}}%
\pgfusepath{clip}%
\pgfsetbuttcap%
\pgfsetmiterjoin%
\definecolor{currentfill}{rgb}{0.211626,0.525352,0.752157}%
\pgfsetfillcolor{currentfill}%
\pgfsetlinewidth{0.000000pt}%
\definecolor{currentstroke}{rgb}{0.000000,0.000000,0.000000}%
\pgfsetstrokecolor{currentstroke}%
\pgfsetstrokeopacity{0.000000}%
\pgfsetdash{}{0pt}%
\pgfpathmoveto{\pgfqpoint{2.283146in}{1.250837in}}%
\pgfpathlineto{\pgfqpoint{2.283471in}{1.248041in}}%
\pgfpathlineto{\pgfqpoint{2.262780in}{1.245532in}}%
\pgfpathlineto{\pgfqpoint{2.261811in}{1.262682in}}%
\pgfpathlineto{\pgfqpoint{2.249581in}{1.261349in}}%
\pgfpathlineto{\pgfqpoint{2.248890in}{1.267034in}}%
\pgfpathlineto{\pgfqpoint{2.239459in}{1.266008in}}%
\pgfpathlineto{\pgfqpoint{2.236831in}{1.288875in}}%
\pgfpathlineto{\pgfqpoint{2.234886in}{1.288646in}}%
\pgfpathlineto{\pgfqpoint{2.233530in}{1.300204in}}%
\pgfpathlineto{\pgfqpoint{2.234453in}{1.308012in}}%
\pgfpathlineto{\pgfqpoint{2.240301in}{1.306723in}}%
\pgfpathlineto{\pgfqpoint{2.257315in}{1.308409in}}%
\pgfpathlineto{\pgfqpoint{2.253575in}{1.340858in}}%
\pgfpathlineto{\pgfqpoint{2.271732in}{1.342879in}}%
\pgfpathlineto{\pgfqpoint{2.273138in}{1.338797in}}%
\pgfpathlineto{\pgfqpoint{2.278491in}{1.292509in}}%
\pgfpathlineto{\pgfqpoint{2.283146in}{1.250837in}}%
\pgfpathclose%
\pgfusepath{fill}%
\end{pgfscope}%
\begin{pgfscope}%
\pgfpathrectangle{\pgfqpoint{0.100000in}{0.100000in}}{\pgfqpoint{3.007045in}{1.925000in}}%
\pgfusepath{clip}%
\pgfsetbuttcap%
\pgfsetmiterjoin%
\definecolor{currentfill}{rgb}{0.371688,0.649627,0.820515}%
\pgfsetfillcolor{currentfill}%
\pgfsetlinewidth{0.000000pt}%
\definecolor{currentstroke}{rgb}{0.000000,0.000000,0.000000}%
\pgfsetstrokecolor{currentstroke}%
\pgfsetstrokeopacity{0.000000}%
\pgfsetdash{}{0pt}%
\pgfpathmoveto{\pgfqpoint{2.380873in}{1.017870in}}%
\pgfpathlineto{\pgfqpoint{2.378204in}{1.022084in}}%
\pgfpathlineto{\pgfqpoint{2.378235in}{1.029769in}}%
\pgfpathlineto{\pgfqpoint{2.376071in}{1.035433in}}%
\pgfpathlineto{\pgfqpoint{2.371954in}{1.050139in}}%
\pgfpathlineto{\pgfqpoint{2.372484in}{1.056245in}}%
\pgfpathlineto{\pgfqpoint{2.370868in}{1.060399in}}%
\pgfpathlineto{\pgfqpoint{2.375741in}{1.064118in}}%
\pgfpathlineto{\pgfqpoint{2.378157in}{1.062296in}}%
\pgfpathlineto{\pgfqpoint{2.384647in}{1.067115in}}%
\pgfpathlineto{\pgfqpoint{2.392156in}{1.068534in}}%
\pgfpathlineto{\pgfqpoint{2.392967in}{1.076431in}}%
\pgfpathlineto{\pgfqpoint{2.400437in}{1.076137in}}%
\pgfpathlineto{\pgfqpoint{2.410452in}{1.072460in}}%
\pgfpathlineto{\pgfqpoint{2.412270in}{1.064935in}}%
\pgfpathlineto{\pgfqpoint{2.417226in}{1.060057in}}%
\pgfpathlineto{\pgfqpoint{2.423037in}{1.058784in}}%
\pgfpathlineto{\pgfqpoint{2.415558in}{1.052709in}}%
\pgfpathlineto{\pgfqpoint{2.415182in}{1.047525in}}%
\pgfpathlineto{\pgfqpoint{2.409125in}{1.041720in}}%
\pgfpathlineto{\pgfqpoint{2.409387in}{1.036508in}}%
\pgfpathlineto{\pgfqpoint{2.399471in}{1.032629in}}%
\pgfpathlineto{\pgfqpoint{2.397238in}{1.025171in}}%
\pgfpathlineto{\pgfqpoint{2.380873in}{1.017870in}}%
\pgfpathclose%
\pgfusepath{fill}%
\end{pgfscope}%
\begin{pgfscope}%
\pgfpathrectangle{\pgfqpoint{0.100000in}{0.100000in}}{\pgfqpoint{3.007045in}{1.925000in}}%
\pgfusepath{clip}%
\pgfsetbuttcap%
\pgfsetmiterjoin%
\definecolor{currentfill}{rgb}{0.346467,0.632403,0.810673}%
\pgfsetfillcolor{currentfill}%
\pgfsetlinewidth{0.000000pt}%
\definecolor{currentstroke}{rgb}{0.000000,0.000000,0.000000}%
\pgfsetstrokecolor{currentstroke}%
\pgfsetstrokeopacity{0.000000}%
\pgfsetdash{}{0pt}%
\pgfpathmoveto{\pgfqpoint{2.508172in}{0.896146in}}%
\pgfpathlineto{\pgfqpoint{2.508933in}{0.891663in}}%
\pgfpathlineto{\pgfqpoint{2.502098in}{0.888878in}}%
\pgfpathlineto{\pgfqpoint{2.496666in}{0.891862in}}%
\pgfpathlineto{\pgfqpoint{2.481602in}{0.894314in}}%
\pgfpathlineto{\pgfqpoint{2.476040in}{0.897437in}}%
\pgfpathlineto{\pgfqpoint{2.462329in}{0.908061in}}%
\pgfpathlineto{\pgfqpoint{2.459680in}{0.930712in}}%
\pgfpathlineto{\pgfqpoint{2.452389in}{0.929337in}}%
\pgfpathlineto{\pgfqpoint{2.451924in}{0.935651in}}%
\pgfpathlineto{\pgfqpoint{2.455447in}{0.943242in}}%
\pgfpathlineto{\pgfqpoint{2.460723in}{0.945065in}}%
\pgfpathlineto{\pgfqpoint{2.472311in}{0.936614in}}%
\pgfpathlineto{\pgfqpoint{2.472948in}{0.932025in}}%
\pgfpathlineto{\pgfqpoint{2.505332in}{0.935315in}}%
\pgfpathlineto{\pgfqpoint{2.504209in}{0.927233in}}%
\pgfpathlineto{\pgfqpoint{2.500183in}{0.925858in}}%
\pgfpathlineto{\pgfqpoint{2.503825in}{0.913860in}}%
\pgfpathlineto{\pgfqpoint{2.505466in}{0.902631in}}%
\pgfpathlineto{\pgfqpoint{2.508172in}{0.896146in}}%
\pgfpathclose%
\pgfusepath{fill}%
\end{pgfscope}%
\begin{pgfscope}%
\pgfpathrectangle{\pgfqpoint{0.100000in}{0.100000in}}{\pgfqpoint{3.007045in}{1.925000in}}%
\pgfusepath{clip}%
\pgfsetbuttcap%
\pgfsetmiterjoin%
\definecolor{currentfill}{rgb}{0.093272,0.396878,0.673664}%
\pgfsetfillcolor{currentfill}%
\pgfsetlinewidth{0.000000pt}%
\definecolor{currentstroke}{rgb}{0.000000,0.000000,0.000000}%
\pgfsetstrokecolor{currentstroke}%
\pgfsetstrokeopacity{0.000000}%
\pgfsetdash{}{0pt}%
\pgfpathmoveto{\pgfqpoint{0.928770in}{0.169774in}}%
\pgfpathlineto{\pgfqpoint{0.929087in}{0.172237in}}%
\pgfpathlineto{\pgfqpoint{0.930780in}{0.170565in}}%
\pgfpathlineto{\pgfqpoint{0.928770in}{0.169774in}}%
\pgfpathclose%
\pgfusepath{fill}%
\end{pgfscope}%
\begin{pgfscope}%
\pgfpathrectangle{\pgfqpoint{0.100000in}{0.100000in}}{\pgfqpoint{3.007045in}{1.925000in}}%
\pgfusepath{clip}%
\pgfsetbuttcap%
\pgfsetmiterjoin%
\definecolor{currentfill}{rgb}{0.093272,0.396878,0.673664}%
\pgfsetfillcolor{currentfill}%
\pgfsetlinewidth{0.000000pt}%
\definecolor{currentstroke}{rgb}{0.000000,0.000000,0.000000}%
\pgfsetstrokecolor{currentstroke}%
\pgfsetstrokeopacity{0.000000}%
\pgfsetdash{}{0pt}%
\pgfpathmoveto{\pgfqpoint{0.927632in}{0.179344in}}%
\pgfpathlineto{\pgfqpoint{0.926956in}{0.181069in}}%
\pgfpathlineto{\pgfqpoint{0.926534in}{0.189086in}}%
\pgfpathlineto{\pgfqpoint{0.927077in}{0.189993in}}%
\pgfpathlineto{\pgfqpoint{0.926777in}{0.193425in}}%
\pgfpathlineto{\pgfqpoint{0.927516in}{0.194491in}}%
\pgfpathlineto{\pgfqpoint{0.926137in}{0.195868in}}%
\pgfpathlineto{\pgfqpoint{0.925991in}{0.197780in}}%
\pgfpathlineto{\pgfqpoint{0.928814in}{0.198830in}}%
\pgfpathlineto{\pgfqpoint{0.929984in}{0.200875in}}%
\pgfpathlineto{\pgfqpoint{0.928490in}{0.201725in}}%
\pgfpathlineto{\pgfqpoint{0.928312in}{0.203305in}}%
\pgfpathlineto{\pgfqpoint{0.929750in}{0.205372in}}%
\pgfpathlineto{\pgfqpoint{0.928020in}{0.209653in}}%
\pgfpathlineto{\pgfqpoint{0.932348in}{0.213098in}}%
\pgfpathlineto{\pgfqpoint{0.934167in}{0.209712in}}%
\pgfpathlineto{\pgfqpoint{0.937918in}{0.207074in}}%
\pgfpathlineto{\pgfqpoint{0.935283in}{0.206082in}}%
\pgfpathlineto{\pgfqpoint{0.936138in}{0.203890in}}%
\pgfpathlineto{\pgfqpoint{0.935915in}{0.201464in}}%
\pgfpathlineto{\pgfqpoint{0.934812in}{0.198205in}}%
\pgfpathlineto{\pgfqpoint{0.934411in}{0.195163in}}%
\pgfpathlineto{\pgfqpoint{0.933099in}{0.192388in}}%
\pgfpathlineto{\pgfqpoint{0.933499in}{0.191353in}}%
\pgfpathlineto{\pgfqpoint{0.932518in}{0.188498in}}%
\pgfpathlineto{\pgfqpoint{0.930488in}{0.187744in}}%
\pgfpathlineto{\pgfqpoint{0.931046in}{0.186350in}}%
\pgfpathlineto{\pgfqpoint{0.929012in}{0.184009in}}%
\pgfpathlineto{\pgfqpoint{0.927632in}{0.179344in}}%
\pgfpathclose%
\pgfusepath{fill}%
\end{pgfscope}%
\begin{pgfscope}%
\pgfpathrectangle{\pgfqpoint{0.100000in}{0.100000in}}{\pgfqpoint{3.007045in}{1.925000in}}%
\pgfusepath{clip}%
\pgfsetbuttcap%
\pgfsetmiterjoin%
\definecolor{currentfill}{rgb}{0.093272,0.396878,0.673664}%
\pgfsetfillcolor{currentfill}%
\pgfsetlinewidth{0.000000pt}%
\definecolor{currentstroke}{rgb}{0.000000,0.000000,0.000000}%
\pgfsetstrokecolor{currentstroke}%
\pgfsetstrokeopacity{0.000000}%
\pgfsetdash{}{0pt}%
\pgfpathmoveto{\pgfqpoint{0.921995in}{0.203350in}}%
\pgfpathlineto{\pgfqpoint{0.924523in}{0.205630in}}%
\pgfpathlineto{\pgfqpoint{0.924613in}{0.208652in}}%
\pgfpathlineto{\pgfqpoint{0.925424in}{0.210159in}}%
\pgfpathlineto{\pgfqpoint{0.927073in}{0.209363in}}%
\pgfpathlineto{\pgfqpoint{0.926551in}{0.205362in}}%
\pgfpathlineto{\pgfqpoint{0.925876in}{0.203745in}}%
\pgfpathlineto{\pgfqpoint{0.924259in}{0.202862in}}%
\pgfpathlineto{\pgfqpoint{0.921995in}{0.203350in}}%
\pgfpathclose%
\pgfusepath{fill}%
\end{pgfscope}%
\begin{pgfscope}%
\pgfpathrectangle{\pgfqpoint{0.100000in}{0.100000in}}{\pgfqpoint{3.007045in}{1.925000in}}%
\pgfusepath{clip}%
\pgfsetbuttcap%
\pgfsetmiterjoin%
\definecolor{currentfill}{rgb}{0.093272,0.396878,0.673664}%
\pgfsetfillcolor{currentfill}%
\pgfsetlinewidth{0.000000pt}%
\definecolor{currentstroke}{rgb}{0.000000,0.000000,0.000000}%
\pgfsetstrokecolor{currentstroke}%
\pgfsetstrokeopacity{0.000000}%
\pgfsetdash{}{0pt}%
\pgfpathmoveto{\pgfqpoint{0.924230in}{0.224566in}}%
\pgfpathlineto{\pgfqpoint{0.928775in}{0.222468in}}%
\pgfpathlineto{\pgfqpoint{0.931898in}{0.225478in}}%
\pgfpathlineto{\pgfqpoint{0.933152in}{0.222753in}}%
\pgfpathlineto{\pgfqpoint{0.934778in}{0.220682in}}%
\pgfpathlineto{\pgfqpoint{0.936886in}{0.216547in}}%
\pgfpathlineto{\pgfqpoint{0.939636in}{0.215477in}}%
\pgfpathlineto{\pgfqpoint{0.940303in}{0.214612in}}%
\pgfpathlineto{\pgfqpoint{0.938629in}{0.208282in}}%
\pgfpathlineto{\pgfqpoint{0.936276in}{0.208683in}}%
\pgfpathlineto{\pgfqpoint{0.931068in}{0.218130in}}%
\pgfpathlineto{\pgfqpoint{0.930637in}{0.216360in}}%
\pgfpathlineto{\pgfqpoint{0.931418in}{0.214692in}}%
\pgfpathlineto{\pgfqpoint{0.930183in}{0.212277in}}%
\pgfpathlineto{\pgfqpoint{0.927651in}{0.210591in}}%
\pgfpathlineto{\pgfqpoint{0.926290in}{0.211782in}}%
\pgfpathlineto{\pgfqpoint{0.925901in}{0.214793in}}%
\pgfpathlineto{\pgfqpoint{0.926598in}{0.218053in}}%
\pgfpathlineto{\pgfqpoint{0.924230in}{0.224566in}}%
\pgfpathclose%
\pgfusepath{fill}%
\end{pgfscope}%
\begin{pgfscope}%
\pgfpathrectangle{\pgfqpoint{0.100000in}{0.100000in}}{\pgfqpoint{3.007045in}{1.925000in}}%
\pgfusepath{clip}%
\pgfsetbuttcap%
\pgfsetmiterjoin%
\definecolor{currentfill}{rgb}{0.093272,0.396878,0.673664}%
\pgfsetfillcolor{currentfill}%
\pgfsetlinewidth{0.000000pt}%
\definecolor{currentstroke}{rgb}{0.000000,0.000000,0.000000}%
\pgfsetstrokecolor{currentstroke}%
\pgfsetstrokeopacity{0.000000}%
\pgfsetdash{}{0pt}%
\pgfpathmoveto{\pgfqpoint{0.949543in}{0.170923in}}%
\pgfpathlineto{\pgfqpoint{0.948341in}{0.174468in}}%
\pgfpathlineto{\pgfqpoint{0.951022in}{0.176053in}}%
\pgfpathlineto{\pgfqpoint{0.953034in}{0.175105in}}%
\pgfpathlineto{\pgfqpoint{0.953826in}{0.173319in}}%
\pgfpathlineto{\pgfqpoint{0.951936in}{0.170704in}}%
\pgfpathlineto{\pgfqpoint{0.949543in}{0.170923in}}%
\pgfpathclose%
\pgfusepath{fill}%
\end{pgfscope}%
\begin{pgfscope}%
\pgfpathrectangle{\pgfqpoint{0.100000in}{0.100000in}}{\pgfqpoint{3.007045in}{1.925000in}}%
\pgfusepath{clip}%
\pgfsetbuttcap%
\pgfsetmiterjoin%
\definecolor{currentfill}{rgb}{0.093272,0.396878,0.673664}%
\pgfsetfillcolor{currentfill}%
\pgfsetlinewidth{0.000000pt}%
\definecolor{currentstroke}{rgb}{0.000000,0.000000,0.000000}%
\pgfsetstrokecolor{currentstroke}%
\pgfsetstrokeopacity{0.000000}%
\pgfsetdash{}{0pt}%
\pgfpathmoveto{\pgfqpoint{0.941866in}{0.183097in}}%
\pgfpathlineto{\pgfqpoint{0.939091in}{0.182574in}}%
\pgfpathlineto{\pgfqpoint{0.938900in}{0.178993in}}%
\pgfpathlineto{\pgfqpoint{0.936985in}{0.178377in}}%
\pgfpathlineto{\pgfqpoint{0.937239in}{0.175986in}}%
\pgfpathlineto{\pgfqpoint{0.935359in}{0.173870in}}%
\pgfpathlineto{\pgfqpoint{0.934641in}{0.176663in}}%
\pgfpathlineto{\pgfqpoint{0.933054in}{0.175528in}}%
\pgfpathlineto{\pgfqpoint{0.932146in}{0.176239in}}%
\pgfpathlineto{\pgfqpoint{0.933479in}{0.179062in}}%
\pgfpathlineto{\pgfqpoint{0.935593in}{0.180734in}}%
\pgfpathlineto{\pgfqpoint{0.937053in}{0.180678in}}%
\pgfpathlineto{\pgfqpoint{0.938116in}{0.183219in}}%
\pgfpathlineto{\pgfqpoint{0.936870in}{0.183669in}}%
\pgfpathlineto{\pgfqpoint{0.935994in}{0.185420in}}%
\pgfpathlineto{\pgfqpoint{0.936562in}{0.187485in}}%
\pgfpathlineto{\pgfqpoint{0.936865in}{0.191439in}}%
\pgfpathlineto{\pgfqpoint{0.939368in}{0.192607in}}%
\pgfpathlineto{\pgfqpoint{0.942659in}{0.187128in}}%
\pgfpathlineto{\pgfqpoint{0.943420in}{0.184913in}}%
\pgfpathlineto{\pgfqpoint{0.941866in}{0.183097in}}%
\pgfpathclose%
\pgfusepath{fill}%
\end{pgfscope}%
\begin{pgfscope}%
\pgfpathrectangle{\pgfqpoint{0.100000in}{0.100000in}}{\pgfqpoint{3.007045in}{1.925000in}}%
\pgfusepath{clip}%
\pgfsetbuttcap%
\pgfsetmiterjoin%
\definecolor{currentfill}{rgb}{0.093272,0.396878,0.673664}%
\pgfsetfillcolor{currentfill}%
\pgfsetlinewidth{0.000000pt}%
\definecolor{currentstroke}{rgb}{0.000000,0.000000,0.000000}%
\pgfsetstrokecolor{currentstroke}%
\pgfsetstrokeopacity{0.000000}%
\pgfsetdash{}{0pt}%
\pgfpathmoveto{\pgfqpoint{0.951030in}{0.194631in}}%
\pgfpathlineto{\pgfqpoint{0.954245in}{0.194809in}}%
\pgfpathlineto{\pgfqpoint{0.955692in}{0.192812in}}%
\pgfpathlineto{\pgfqpoint{0.957247in}{0.193413in}}%
\pgfpathlineto{\pgfqpoint{0.958655in}{0.195219in}}%
\pgfpathlineto{\pgfqpoint{0.963493in}{0.195565in}}%
\pgfpathlineto{\pgfqpoint{0.964721in}{0.193135in}}%
\pgfpathlineto{\pgfqpoint{0.966052in}{0.192399in}}%
\pgfpathlineto{\pgfqpoint{0.966135in}{0.189796in}}%
\pgfpathlineto{\pgfqpoint{0.963633in}{0.186915in}}%
\pgfpathlineto{\pgfqpoint{0.966486in}{0.184667in}}%
\pgfpathlineto{\pgfqpoint{0.964468in}{0.180678in}}%
\pgfpathlineto{\pgfqpoint{0.966318in}{0.178588in}}%
\pgfpathlineto{\pgfqpoint{0.965355in}{0.173338in}}%
\pgfpathlineto{\pgfqpoint{0.967092in}{0.172775in}}%
\pgfpathlineto{\pgfqpoint{0.969300in}{0.170288in}}%
\pgfpathlineto{\pgfqpoint{0.972690in}{0.165137in}}%
\pgfpathlineto{\pgfqpoint{0.972584in}{0.162951in}}%
\pgfpathlineto{\pgfqpoint{0.970864in}{0.160909in}}%
\pgfpathlineto{\pgfqpoint{0.968598in}{0.161733in}}%
\pgfpathlineto{\pgfqpoint{0.965820in}{0.161252in}}%
\pgfpathlineto{\pgfqpoint{0.966014in}{0.159668in}}%
\pgfpathlineto{\pgfqpoint{0.964181in}{0.158832in}}%
\pgfpathlineto{\pgfqpoint{0.962912in}{0.161570in}}%
\pgfpathlineto{\pgfqpoint{0.960550in}{0.161270in}}%
\pgfpathlineto{\pgfqpoint{0.958681in}{0.160262in}}%
\pgfpathlineto{\pgfqpoint{0.958436in}{0.159148in}}%
\pgfpathlineto{\pgfqpoint{0.956199in}{0.158645in}}%
\pgfpathlineto{\pgfqpoint{0.956678in}{0.156700in}}%
\pgfpathlineto{\pgfqpoint{0.955416in}{0.154694in}}%
\pgfpathlineto{\pgfqpoint{0.953820in}{0.154455in}}%
\pgfpathlineto{\pgfqpoint{0.952217in}{0.157810in}}%
\pgfpathlineto{\pgfqpoint{0.955289in}{0.158119in}}%
\pgfpathlineto{\pgfqpoint{0.956244in}{0.159613in}}%
\pgfpathlineto{\pgfqpoint{0.958348in}{0.160762in}}%
\pgfpathlineto{\pgfqpoint{0.959317in}{0.163914in}}%
\pgfpathlineto{\pgfqpoint{0.960954in}{0.164908in}}%
\pgfpathlineto{\pgfqpoint{0.959218in}{0.166639in}}%
\pgfpathlineto{\pgfqpoint{0.958611in}{0.164391in}}%
\pgfpathlineto{\pgfqpoint{0.955800in}{0.164056in}}%
\pgfpathlineto{\pgfqpoint{0.954820in}{0.161400in}}%
\pgfpathlineto{\pgfqpoint{0.952564in}{0.162138in}}%
\pgfpathlineto{\pgfqpoint{0.952336in}{0.166328in}}%
\pgfpathlineto{\pgfqpoint{0.950396in}{0.166275in}}%
\pgfpathlineto{\pgfqpoint{0.951336in}{0.169331in}}%
\pgfpathlineto{\pgfqpoint{0.955934in}{0.171121in}}%
\pgfpathlineto{\pgfqpoint{0.955518in}{0.168806in}}%
\pgfpathlineto{\pgfqpoint{0.956516in}{0.167155in}}%
\pgfpathlineto{\pgfqpoint{0.957594in}{0.173312in}}%
\pgfpathlineto{\pgfqpoint{0.958602in}{0.175079in}}%
\pgfpathlineto{\pgfqpoint{0.961488in}{0.177541in}}%
\pgfpathlineto{\pgfqpoint{0.957917in}{0.179285in}}%
\pgfpathlineto{\pgfqpoint{0.958031in}{0.180854in}}%
\pgfpathlineto{\pgfqpoint{0.956711in}{0.183035in}}%
\pgfpathlineto{\pgfqpoint{0.956083in}{0.186705in}}%
\pgfpathlineto{\pgfqpoint{0.957478in}{0.186958in}}%
\pgfpathlineto{\pgfqpoint{0.957475in}{0.188952in}}%
\pgfpathlineto{\pgfqpoint{0.955586in}{0.189502in}}%
\pgfpathlineto{\pgfqpoint{0.954193in}{0.191220in}}%
\pgfpathlineto{\pgfqpoint{0.953928in}{0.193026in}}%
\pgfpathlineto{\pgfqpoint{0.952297in}{0.192656in}}%
\pgfpathlineto{\pgfqpoint{0.951030in}{0.194631in}}%
\pgfpathclose%
\pgfusepath{fill}%
\end{pgfscope}%
\begin{pgfscope}%
\pgfpathrectangle{\pgfqpoint{0.100000in}{0.100000in}}{\pgfqpoint{3.007045in}{1.925000in}}%
\pgfusepath{clip}%
\pgfsetbuttcap%
\pgfsetmiterjoin%
\definecolor{currentfill}{rgb}{0.093272,0.396878,0.673664}%
\pgfsetfillcolor{currentfill}%
\pgfsetlinewidth{0.000000pt}%
\definecolor{currentstroke}{rgb}{0.000000,0.000000,0.000000}%
\pgfsetstrokecolor{currentstroke}%
\pgfsetstrokeopacity{0.000000}%
\pgfsetdash{}{0pt}%
\pgfpathmoveto{\pgfqpoint{0.942241in}{0.179713in}}%
\pgfpathlineto{\pgfqpoint{0.944515in}{0.186094in}}%
\pgfpathlineto{\pgfqpoint{0.944156in}{0.190021in}}%
\pgfpathlineto{\pgfqpoint{0.944652in}{0.190685in}}%
\pgfpathlineto{\pgfqpoint{0.943856in}{0.194420in}}%
\pgfpathlineto{\pgfqpoint{0.945632in}{0.194990in}}%
\pgfpathlineto{\pgfqpoint{0.951121in}{0.190206in}}%
\pgfpathlineto{\pgfqpoint{0.954271in}{0.188027in}}%
\pgfpathlineto{\pgfqpoint{0.954701in}{0.185800in}}%
\pgfpathlineto{\pgfqpoint{0.953794in}{0.184014in}}%
\pgfpathlineto{\pgfqpoint{0.955507in}{0.182225in}}%
\pgfpathlineto{\pgfqpoint{0.956828in}{0.176810in}}%
\pgfpathlineto{\pgfqpoint{0.953089in}{0.176262in}}%
\pgfpathlineto{\pgfqpoint{0.951273in}{0.177255in}}%
\pgfpathlineto{\pgfqpoint{0.953068in}{0.180485in}}%
\pgfpathlineto{\pgfqpoint{0.951111in}{0.180223in}}%
\pgfpathlineto{\pgfqpoint{0.949807in}{0.178518in}}%
\pgfpathlineto{\pgfqpoint{0.947328in}{0.177601in}}%
\pgfpathlineto{\pgfqpoint{0.945308in}{0.179613in}}%
\pgfpathlineto{\pgfqpoint{0.944254in}{0.178896in}}%
\pgfpathlineto{\pgfqpoint{0.942241in}{0.179713in}}%
\pgfpathclose%
\pgfusepath{fill}%
\end{pgfscope}%
\begin{pgfscope}%
\pgfpathrectangle{\pgfqpoint{0.100000in}{0.100000in}}{\pgfqpoint{3.007045in}{1.925000in}}%
\pgfusepath{clip}%
\pgfsetbuttcap%
\pgfsetmiterjoin%
\definecolor{currentfill}{rgb}{0.735871,0.841569,0.923045}%
\pgfsetfillcolor{currentfill}%
\pgfsetlinewidth{0.000000pt}%
\definecolor{currentstroke}{rgb}{0.000000,0.000000,0.000000}%
\pgfsetstrokecolor{currentstroke}%
\pgfsetstrokeopacity{0.000000}%
\pgfsetdash{}{0pt}%
\pgfpathmoveto{\pgfqpoint{0.981446in}{1.653571in}}%
\pgfpathlineto{\pgfqpoint{0.966686in}{1.656395in}}%
\pgfpathlineto{\pgfqpoint{0.964271in}{1.652566in}}%
\pgfpathlineto{\pgfqpoint{0.957302in}{1.649942in}}%
\pgfpathlineto{\pgfqpoint{0.955593in}{1.658524in}}%
\pgfpathlineto{\pgfqpoint{0.944841in}{1.665499in}}%
\pgfpathlineto{\pgfqpoint{0.941856in}{1.670351in}}%
\pgfpathlineto{\pgfqpoint{0.936043in}{1.670714in}}%
\pgfpathlineto{\pgfqpoint{0.926926in}{1.665686in}}%
\pgfpathlineto{\pgfqpoint{0.923792in}{1.677916in}}%
\pgfpathlineto{\pgfqpoint{0.918191in}{1.680868in}}%
\pgfpathlineto{\pgfqpoint{0.917372in}{1.685810in}}%
\pgfpathlineto{\pgfqpoint{0.913359in}{1.689243in}}%
\pgfpathlineto{\pgfqpoint{0.914111in}{1.694407in}}%
\pgfpathlineto{\pgfqpoint{0.917747in}{1.701623in}}%
\pgfpathlineto{\pgfqpoint{0.917545in}{1.708081in}}%
\pgfpathlineto{\pgfqpoint{0.914970in}{1.712349in}}%
\pgfpathlineto{\pgfqpoint{0.915399in}{1.718366in}}%
\pgfpathlineto{\pgfqpoint{0.921093in}{1.727298in}}%
\pgfpathlineto{\pgfqpoint{0.926725in}{1.726127in}}%
\pgfpathlineto{\pgfqpoint{0.927484in}{1.729880in}}%
\pgfpathlineto{\pgfqpoint{0.936258in}{1.730024in}}%
\pgfpathlineto{\pgfqpoint{0.940260in}{1.735106in}}%
\pgfpathlineto{\pgfqpoint{0.945308in}{1.734086in}}%
\pgfpathlineto{\pgfqpoint{0.949935in}{1.756722in}}%
\pgfpathlineto{\pgfqpoint{0.942694in}{1.758174in}}%
\pgfpathlineto{\pgfqpoint{0.948336in}{1.785640in}}%
\pgfpathlineto{\pgfqpoint{0.962444in}{1.782513in}}%
\pgfpathlineto{\pgfqpoint{0.964785in}{1.774681in}}%
\pgfpathlineto{\pgfqpoint{0.960752in}{1.754525in}}%
\pgfpathlineto{\pgfqpoint{0.972359in}{1.752059in}}%
\pgfpathlineto{\pgfqpoint{0.967944in}{1.729505in}}%
\pgfpathlineto{\pgfqpoint{0.979570in}{1.727393in}}%
\pgfpathlineto{\pgfqpoint{0.977851in}{1.718526in}}%
\pgfpathlineto{\pgfqpoint{0.983438in}{1.717518in}}%
\pgfpathlineto{\pgfqpoint{0.987484in}{1.710382in}}%
\pgfpathlineto{\pgfqpoint{0.985045in}{1.700348in}}%
\pgfpathlineto{\pgfqpoint{0.982127in}{1.696489in}}%
\pgfpathlineto{\pgfqpoint{0.973346in}{1.694790in}}%
\pgfpathlineto{\pgfqpoint{0.971079in}{1.692732in}}%
\pgfpathlineto{\pgfqpoint{0.969357in}{1.685011in}}%
\pgfpathlineto{\pgfqpoint{0.973053in}{1.678906in}}%
\pgfpathlineto{\pgfqpoint{0.972309in}{1.661350in}}%
\pgfpathlineto{\pgfqpoint{0.981446in}{1.653571in}}%
\pgfpathclose%
\pgfusepath{fill}%
\end{pgfscope}%
\begin{pgfscope}%
\pgfpathrectangle{\pgfqpoint{0.100000in}{0.100000in}}{\pgfqpoint{3.007045in}{1.925000in}}%
\pgfusepath{clip}%
\pgfsetbuttcap%
\pgfsetmiterjoin%
\definecolor{currentfill}{rgb}{0.179146,0.492872,0.735425}%
\pgfsetfillcolor{currentfill}%
\pgfsetlinewidth{0.000000pt}%
\definecolor{currentstroke}{rgb}{0.000000,0.000000,0.000000}%
\pgfsetstrokecolor{currentstroke}%
\pgfsetstrokeopacity{0.000000}%
\pgfsetdash{}{0pt}%
\pgfpathmoveto{\pgfqpoint{1.631190in}{1.247136in}}%
\pgfpathlineto{\pgfqpoint{1.586476in}{1.248770in}}%
\pgfpathlineto{\pgfqpoint{1.587386in}{1.271651in}}%
\pgfpathlineto{\pgfqpoint{1.586068in}{1.271711in}}%
\pgfpathlineto{\pgfqpoint{1.587220in}{1.294590in}}%
\pgfpathlineto{\pgfqpoint{1.609594in}{1.293697in}}%
\pgfpathlineto{\pgfqpoint{1.610083in}{1.285141in}}%
\pgfpathlineto{\pgfqpoint{1.618684in}{1.284817in}}%
\pgfpathlineto{\pgfqpoint{1.621155in}{1.286648in}}%
\pgfpathlineto{\pgfqpoint{1.638392in}{1.287167in}}%
\pgfpathlineto{\pgfqpoint{1.638646in}{1.292816in}}%
\pgfpathlineto{\pgfqpoint{1.643851in}{1.292672in}}%
\pgfpathlineto{\pgfqpoint{1.643733in}{1.288581in}}%
\pgfpathlineto{\pgfqpoint{1.635249in}{1.282362in}}%
\pgfpathlineto{\pgfqpoint{1.632031in}{1.278421in}}%
\pgfpathlineto{\pgfqpoint{1.631190in}{1.247136in}}%
\pgfpathclose%
\pgfusepath{fill}%
\end{pgfscope}%
\begin{pgfscope}%
\pgfpathrectangle{\pgfqpoint{0.100000in}{0.100000in}}{\pgfqpoint{3.007045in}{1.925000in}}%
\pgfusepath{clip}%
\pgfsetbuttcap%
\pgfsetmiterjoin%
\definecolor{currentfill}{rgb}{0.260715,0.573841,0.777209}%
\pgfsetfillcolor{currentfill}%
\pgfsetlinewidth{0.000000pt}%
\definecolor{currentstroke}{rgb}{0.000000,0.000000,0.000000}%
\pgfsetstrokecolor{currentstroke}%
\pgfsetstrokeopacity{0.000000}%
\pgfsetdash{}{0pt}%
\pgfpathmoveto{\pgfqpoint{0.683125in}{1.670416in}}%
\pgfpathlineto{\pgfqpoint{0.690132in}{1.668259in}}%
\pgfpathlineto{\pgfqpoint{0.695518in}{1.663103in}}%
\pgfpathlineto{\pgfqpoint{0.694447in}{1.659187in}}%
\pgfpathlineto{\pgfqpoint{0.689503in}{1.657799in}}%
\pgfpathlineto{\pgfqpoint{0.689599in}{1.651922in}}%
\pgfpathlineto{\pgfqpoint{0.678258in}{1.651609in}}%
\pgfpathlineto{\pgfqpoint{0.678193in}{1.649100in}}%
\pgfpathlineto{\pgfqpoint{0.684947in}{1.644999in}}%
\pgfpathlineto{\pgfqpoint{0.685898in}{1.641698in}}%
\pgfpathlineto{\pgfqpoint{0.678278in}{1.634204in}}%
\pgfpathlineto{\pgfqpoint{0.677402in}{1.629304in}}%
\pgfpathlineto{\pgfqpoint{0.671575in}{1.622748in}}%
\pgfpathlineto{\pgfqpoint{0.683738in}{1.619619in}}%
\pgfpathlineto{\pgfqpoint{0.680340in}{1.605736in}}%
\pgfpathlineto{\pgfqpoint{0.653418in}{1.613391in}}%
\pgfpathlineto{\pgfqpoint{0.651898in}{1.607795in}}%
\pgfpathlineto{\pgfqpoint{0.613279in}{1.618283in}}%
\pgfpathlineto{\pgfqpoint{0.619536in}{1.640423in}}%
\pgfpathlineto{\pgfqpoint{0.628849in}{1.673291in}}%
\pgfpathlineto{\pgfqpoint{0.630991in}{1.684466in}}%
\pgfpathlineto{\pgfqpoint{0.683125in}{1.670416in}}%
\pgfpathclose%
\pgfusepath{fill}%
\end{pgfscope}%
\begin{pgfscope}%
\pgfpathrectangle{\pgfqpoint{0.100000in}{0.100000in}}{\pgfqpoint{3.007045in}{1.925000in}}%
\pgfusepath{clip}%
\pgfsetbuttcap%
\pgfsetmiterjoin%
\definecolor{currentfill}{rgb}{0.031757,0.318139,0.612149}%
\pgfsetfillcolor{currentfill}%
\pgfsetlinewidth{0.000000pt}%
\definecolor{currentstroke}{rgb}{0.000000,0.000000,0.000000}%
\pgfsetstrokecolor{currentstroke}%
\pgfsetstrokeopacity{0.000000}%
\pgfsetdash{}{0pt}%
\pgfpathmoveto{\pgfqpoint{1.072883in}{1.296139in}}%
\pgfpathlineto{\pgfqpoint{1.062663in}{1.231515in}}%
\pgfpathlineto{\pgfqpoint{1.061002in}{1.220921in}}%
\pgfpathlineto{\pgfqpoint{1.057846in}{1.219034in}}%
\pgfpathlineto{\pgfqpoint{1.011802in}{1.226459in}}%
\pgfpathlineto{\pgfqpoint{1.012248in}{1.233395in}}%
\pgfpathlineto{\pgfqpoint{1.015989in}{1.239279in}}%
\pgfpathlineto{\pgfqpoint{1.016300in}{1.244056in}}%
\pgfpathlineto{\pgfqpoint{1.022492in}{1.247657in}}%
\pgfpathlineto{\pgfqpoint{0.990893in}{1.253051in}}%
\pgfpathlineto{\pgfqpoint{0.974041in}{1.256437in}}%
\pgfpathlineto{\pgfqpoint{0.973337in}{1.262333in}}%
\pgfpathlineto{\pgfqpoint{0.981761in}{1.312933in}}%
\pgfpathlineto{\pgfqpoint{0.989953in}{1.316140in}}%
\pgfpathlineto{\pgfqpoint{0.994498in}{1.314874in}}%
\pgfpathlineto{\pgfqpoint{1.008923in}{1.315572in}}%
\pgfpathlineto{\pgfqpoint{1.015074in}{1.317685in}}%
\pgfpathlineto{\pgfqpoint{1.019386in}{1.315423in}}%
\pgfpathlineto{\pgfqpoint{1.027613in}{1.314125in}}%
\pgfpathlineto{\pgfqpoint{1.033610in}{1.309798in}}%
\pgfpathlineto{\pgfqpoint{1.050961in}{1.311223in}}%
\pgfpathlineto{\pgfqpoint{1.052269in}{1.305131in}}%
\pgfpathlineto{\pgfqpoint{1.057220in}{1.307768in}}%
\pgfpathlineto{\pgfqpoint{1.057913in}{1.312141in}}%
\pgfpathlineto{\pgfqpoint{1.067332in}{1.310685in}}%
\pgfpathlineto{\pgfqpoint{1.068540in}{1.305641in}}%
\pgfpathlineto{\pgfqpoint{1.067620in}{1.298977in}}%
\pgfpathlineto{\pgfqpoint{1.072883in}{1.296139in}}%
\pgfpathclose%
\pgfusepath{fill}%
\end{pgfscope}%
\begin{pgfscope}%
\pgfpathrectangle{\pgfqpoint{0.100000in}{0.100000in}}{\pgfqpoint{3.007045in}{1.925000in}}%
\pgfusepath{clip}%
\pgfsetbuttcap%
\pgfsetmiterjoin%
\definecolor{currentfill}{rgb}{0.280892,0.587620,0.785083}%
\pgfsetfillcolor{currentfill}%
\pgfsetlinewidth{0.000000pt}%
\definecolor{currentstroke}{rgb}{0.000000,0.000000,0.000000}%
\pgfsetstrokecolor{currentstroke}%
\pgfsetstrokeopacity{0.000000}%
\pgfsetdash{}{0pt}%
\pgfpathmoveto{\pgfqpoint{1.772350in}{0.979565in}}%
\pgfpathlineto{\pgfqpoint{1.772364in}{0.997534in}}%
\pgfpathlineto{\pgfqpoint{1.768458in}{1.001242in}}%
\pgfpathlineto{\pgfqpoint{1.768402in}{1.026569in}}%
\pgfpathlineto{\pgfqpoint{1.767740in}{1.049552in}}%
\pgfpathlineto{\pgfqpoint{1.768296in}{1.069654in}}%
\pgfpathlineto{\pgfqpoint{1.792356in}{1.069644in}}%
\pgfpathlineto{\pgfqpoint{1.792358in}{1.025328in}}%
\pgfpathlineto{\pgfqpoint{1.820409in}{1.024532in}}%
\pgfpathlineto{\pgfqpoint{1.822007in}{1.020648in}}%
\pgfpathlineto{\pgfqpoint{1.821509in}{0.984906in}}%
\pgfpathlineto{\pgfqpoint{1.792483in}{0.985928in}}%
\pgfpathlineto{\pgfqpoint{1.792522in}{0.979422in}}%
\pgfpathlineto{\pgfqpoint{1.772350in}{0.979565in}}%
\pgfpathclose%
\pgfusepath{fill}%
\end{pgfscope}%
\begin{pgfscope}%
\pgfpathrectangle{\pgfqpoint{0.100000in}{0.100000in}}{\pgfqpoint{3.007045in}{1.925000in}}%
\pgfusepath{clip}%
\pgfsetbuttcap%
\pgfsetmiterjoin%
\definecolor{currentfill}{rgb}{0.567059,0.763306,0.871019}%
\pgfsetfillcolor{currentfill}%
\pgfsetlinewidth{0.000000pt}%
\definecolor{currentstroke}{rgb}{0.000000,0.000000,0.000000}%
\pgfsetstrokecolor{currentstroke}%
\pgfsetstrokeopacity{0.000000}%
\pgfsetdash{}{0pt}%
\pgfpathmoveto{\pgfqpoint{2.584120in}{0.229934in}}%
\pgfpathlineto{\pgfqpoint{2.584920in}{0.232373in}}%
\pgfpathlineto{\pgfqpoint{2.594250in}{0.236772in}}%
\pgfpathlineto{\pgfqpoint{2.597851in}{0.242148in}}%
\pgfpathlineto{\pgfqpoint{2.603807in}{0.246137in}}%
\pgfpathlineto{\pgfqpoint{2.612413in}{0.245063in}}%
\pgfpathlineto{\pgfqpoint{2.616307in}{0.242603in}}%
\pgfpathlineto{\pgfqpoint{2.608919in}{0.238771in}}%
\pgfpathlineto{\pgfqpoint{2.602623in}{0.240537in}}%
\pgfpathlineto{\pgfqpoint{2.601276in}{0.237478in}}%
\pgfpathlineto{\pgfqpoint{2.591656in}{0.231982in}}%
\pgfpathlineto{\pgfqpoint{2.584120in}{0.229934in}}%
\pgfpathclose%
\pgfusepath{fill}%
\end{pgfscope}%
\begin{pgfscope}%
\pgfpathrectangle{\pgfqpoint{0.100000in}{0.100000in}}{\pgfqpoint{3.007045in}{1.925000in}}%
\pgfusepath{clip}%
\pgfsetbuttcap%
\pgfsetmiterjoin%
\definecolor{currentfill}{rgb}{0.567059,0.763306,0.871019}%
\pgfsetfillcolor{currentfill}%
\pgfsetlinewidth{0.000000pt}%
\definecolor{currentstroke}{rgb}{0.000000,0.000000,0.000000}%
\pgfsetstrokecolor{currentstroke}%
\pgfsetstrokeopacity{0.000000}%
\pgfsetdash{}{0pt}%
\pgfpathmoveto{\pgfqpoint{2.599950in}{0.318474in}}%
\pgfpathlineto{\pgfqpoint{2.628250in}{0.323017in}}%
\pgfpathlineto{\pgfqpoint{2.623105in}{0.353228in}}%
\pgfpathlineto{\pgfqpoint{2.622244in}{0.358169in}}%
\pgfpathlineto{\pgfqpoint{2.661292in}{0.364487in}}%
\pgfpathlineto{\pgfqpoint{2.670514in}{0.365060in}}%
\pgfpathlineto{\pgfqpoint{2.671006in}{0.349163in}}%
\pgfpathlineto{\pgfqpoint{2.673290in}{0.330949in}}%
\pgfpathlineto{\pgfqpoint{2.672635in}{0.325576in}}%
\pgfpathlineto{\pgfqpoint{2.666787in}{0.322924in}}%
\pgfpathlineto{\pgfqpoint{2.663898in}{0.305786in}}%
\pgfpathlineto{\pgfqpoint{2.665452in}{0.296790in}}%
\pgfpathlineto{\pgfqpoint{2.656966in}{0.286274in}}%
\pgfpathlineto{\pgfqpoint{2.652560in}{0.288018in}}%
\pgfpathlineto{\pgfqpoint{2.650535in}{0.283393in}}%
\pgfpathlineto{\pgfqpoint{2.645495in}{0.280171in}}%
\pgfpathlineto{\pgfqpoint{2.636303in}{0.281237in}}%
\pgfpathlineto{\pgfqpoint{2.633609in}{0.278212in}}%
\pgfpathlineto{\pgfqpoint{2.622499in}{0.274880in}}%
\pgfpathlineto{\pgfqpoint{2.616272in}{0.281159in}}%
\pgfpathlineto{\pgfqpoint{2.616427in}{0.288354in}}%
\pgfpathlineto{\pgfqpoint{2.614021in}{0.297409in}}%
\pgfpathlineto{\pgfqpoint{2.611003in}{0.299907in}}%
\pgfpathlineto{\pgfqpoint{2.607596in}{0.307893in}}%
\pgfpathlineto{\pgfqpoint{2.601344in}{0.312735in}}%
\pgfpathlineto{\pgfqpoint{2.599950in}{0.318474in}}%
\pgfpathclose%
\pgfusepath{fill}%
\end{pgfscope}%
\begin{pgfscope}%
\pgfpathrectangle{\pgfqpoint{0.100000in}{0.100000in}}{\pgfqpoint{3.007045in}{1.925000in}}%
\pgfusepath{clip}%
\pgfsetbuttcap%
\pgfsetmiterjoin%
\definecolor{currentfill}{rgb}{0.361599,0.642737,0.816578}%
\pgfsetfillcolor{currentfill}%
\pgfsetlinewidth{0.000000pt}%
\definecolor{currentstroke}{rgb}{0.000000,0.000000,0.000000}%
\pgfsetstrokecolor{currentstroke}%
\pgfsetstrokeopacity{0.000000}%
\pgfsetdash{}{0pt}%
\pgfpathmoveto{\pgfqpoint{2.139795in}{0.676959in}}%
\pgfpathlineto{\pgfqpoint{2.130566in}{0.673645in}}%
\pgfpathlineto{\pgfqpoint{2.115516in}{0.670799in}}%
\pgfpathlineto{\pgfqpoint{2.113560in}{0.697106in}}%
\pgfpathlineto{\pgfqpoint{2.090671in}{0.695575in}}%
\pgfpathlineto{\pgfqpoint{2.087936in}{0.742436in}}%
\pgfpathlineto{\pgfqpoint{2.116053in}{0.743899in}}%
\pgfpathlineto{\pgfqpoint{2.141820in}{0.745927in}}%
\pgfpathlineto{\pgfqpoint{2.141241in}{0.722568in}}%
\pgfpathlineto{\pgfqpoint{2.139795in}{0.676959in}}%
\pgfpathclose%
\pgfusepath{fill}%
\end{pgfscope}%
\begin{pgfscope}%
\pgfpathrectangle{\pgfqpoint{0.100000in}{0.100000in}}{\pgfqpoint{3.007045in}{1.925000in}}%
\pgfusepath{clip}%
\pgfsetbuttcap%
\pgfsetmiterjoin%
\definecolor{currentfill}{rgb}{0.485490,0.718524,0.853426}%
\pgfsetfillcolor{currentfill}%
\pgfsetlinewidth{0.000000pt}%
\definecolor{currentstroke}{rgb}{0.000000,0.000000,0.000000}%
\pgfsetstrokecolor{currentstroke}%
\pgfsetstrokeopacity{0.000000}%
\pgfsetdash{}{0pt}%
\pgfpathmoveto{\pgfqpoint{1.054085in}{0.578918in}}%
\pgfpathlineto{\pgfqpoint{1.057990in}{0.585618in}}%
\pgfpathlineto{\pgfqpoint{1.069716in}{0.583261in}}%
\pgfpathlineto{\pgfqpoint{1.065530in}{0.579554in}}%
\pgfpathlineto{\pgfqpoint{1.060064in}{0.581478in}}%
\pgfpathlineto{\pgfqpoint{1.054085in}{0.578918in}}%
\pgfpathclose%
\pgfusepath{fill}%
\end{pgfscope}%
\begin{pgfscope}%
\pgfpathrectangle{\pgfqpoint{0.100000in}{0.100000in}}{\pgfqpoint{3.007045in}{1.925000in}}%
\pgfusepath{clip}%
\pgfsetbuttcap%
\pgfsetmiterjoin%
\definecolor{currentfill}{rgb}{0.485490,0.718524,0.853426}%
\pgfsetfillcolor{currentfill}%
\pgfsetlinewidth{0.000000pt}%
\definecolor{currentstroke}{rgb}{0.000000,0.000000,0.000000}%
\pgfsetstrokecolor{currentstroke}%
\pgfsetstrokeopacity{0.000000}%
\pgfsetdash{}{0pt}%
\pgfpathmoveto{\pgfqpoint{1.094825in}{0.552023in}}%
\pgfpathlineto{\pgfqpoint{1.087738in}{0.562041in}}%
\pgfpathlineto{\pgfqpoint{1.084656in}{0.570735in}}%
\pgfpathlineto{\pgfqpoint{1.085398in}{0.574624in}}%
\pgfpathlineto{\pgfqpoint{1.094010in}{0.578238in}}%
\pgfpathlineto{\pgfqpoint{1.104098in}{0.576218in}}%
\pgfpathlineto{\pgfqpoint{1.107468in}{0.572007in}}%
\pgfpathlineto{\pgfqpoint{1.113289in}{0.568890in}}%
\pgfpathlineto{\pgfqpoint{1.115269in}{0.562992in}}%
\pgfpathlineto{\pgfqpoint{1.112904in}{0.558446in}}%
\pgfpathlineto{\pgfqpoint{1.108126in}{0.556840in}}%
\pgfpathlineto{\pgfqpoint{1.104455in}{0.552095in}}%
\pgfpathlineto{\pgfqpoint{1.094825in}{0.552023in}}%
\pgfpathclose%
\pgfusepath{fill}%
\end{pgfscope}%
\begin{pgfscope}%
\pgfpathrectangle{\pgfqpoint{0.100000in}{0.100000in}}{\pgfqpoint{3.007045in}{1.925000in}}%
\pgfusepath{clip}%
\pgfsetbuttcap%
\pgfsetmiterjoin%
\definecolor{currentfill}{rgb}{0.296025,0.597955,0.790988}%
\pgfsetfillcolor{currentfill}%
\pgfsetlinewidth{0.000000pt}%
\definecolor{currentstroke}{rgb}{0.000000,0.000000,0.000000}%
\pgfsetstrokecolor{currentstroke}%
\pgfsetstrokeopacity{0.000000}%
\pgfsetdash{}{0pt}%
\pgfpathmoveto{\pgfqpoint{1.638404in}{1.708645in}}%
\pgfpathlineto{\pgfqpoint{1.637567in}{1.685608in}}%
\pgfpathlineto{\pgfqpoint{1.638529in}{1.679874in}}%
\pgfpathlineto{\pgfqpoint{1.615660in}{1.680679in}}%
\pgfpathlineto{\pgfqpoint{1.614445in}{1.686458in}}%
\pgfpathlineto{\pgfqpoint{1.615296in}{1.709477in}}%
\pgfpathlineto{\pgfqpoint{1.638404in}{1.708645in}}%
\pgfpathclose%
\pgfusepath{fill}%
\end{pgfscope}%
\begin{pgfscope}%
\pgfpathrectangle{\pgfqpoint{0.100000in}{0.100000in}}{\pgfqpoint{3.007045in}{1.925000in}}%
\pgfusepath{clip}%
\pgfsetbuttcap%
\pgfsetmiterjoin%
\definecolor{currentfill}{rgb}{0.371688,0.649627,0.820515}%
\pgfsetfillcolor{currentfill}%
\pgfsetlinewidth{0.000000pt}%
\definecolor{currentstroke}{rgb}{0.000000,0.000000,0.000000}%
\pgfsetstrokecolor{currentstroke}%
\pgfsetstrokeopacity{0.000000}%
\pgfsetdash{}{0pt}%
\pgfpathmoveto{\pgfqpoint{1.245609in}{1.461365in}}%
\pgfpathlineto{\pgfqpoint{1.218116in}{1.464571in}}%
\pgfpathlineto{\pgfqpoint{1.193681in}{1.468364in}}%
\pgfpathlineto{\pgfqpoint{1.197829in}{1.512182in}}%
\pgfpathlineto{\pgfqpoint{1.199309in}{1.522550in}}%
\pgfpathlineto{\pgfqpoint{1.197692in}{1.530943in}}%
\pgfpathlineto{\pgfqpoint{1.193646in}{1.534813in}}%
\pgfpathlineto{\pgfqpoint{1.188802in}{1.543168in}}%
\pgfpathlineto{\pgfqpoint{1.180687in}{1.550302in}}%
\pgfpathlineto{\pgfqpoint{1.174953in}{1.552958in}}%
\pgfpathlineto{\pgfqpoint{1.170691in}{1.563235in}}%
\pgfpathlineto{\pgfqpoint{1.169718in}{1.571914in}}%
\pgfpathlineto{\pgfqpoint{1.154054in}{1.574061in}}%
\pgfpathlineto{\pgfqpoint{1.160708in}{1.586055in}}%
\pgfpathlineto{\pgfqpoint{1.163161in}{1.587525in}}%
\pgfpathlineto{\pgfqpoint{1.135376in}{1.591485in}}%
\pgfpathlineto{\pgfqpoint{1.138319in}{1.605460in}}%
\pgfpathlineto{\pgfqpoint{1.140335in}{1.607233in}}%
\pgfpathlineto{\pgfqpoint{1.155409in}{1.604740in}}%
\pgfpathlineto{\pgfqpoint{1.167306in}{1.606868in}}%
\pgfpathlineto{\pgfqpoint{1.169858in}{1.624279in}}%
\pgfpathlineto{\pgfqpoint{1.171947in}{1.631734in}}%
\pgfpathlineto{\pgfqpoint{1.179523in}{1.630622in}}%
\pgfpathlineto{\pgfqpoint{1.184148in}{1.635775in}}%
\pgfpathlineto{\pgfqpoint{1.189792in}{1.634979in}}%
\pgfpathlineto{\pgfqpoint{1.190334in}{1.638786in}}%
\pgfpathlineto{\pgfqpoint{1.197757in}{1.637726in}}%
\pgfpathlineto{\pgfqpoint{1.201711in}{1.637190in}}%
\pgfpathlineto{\pgfqpoint{1.200926in}{1.631495in}}%
\pgfpathlineto{\pgfqpoint{1.212290in}{1.629948in}}%
\pgfpathlineto{\pgfqpoint{1.222383in}{1.622797in}}%
\pgfpathlineto{\pgfqpoint{1.221965in}{1.610407in}}%
\pgfpathlineto{\pgfqpoint{1.230291in}{1.608576in}}%
\pgfpathlineto{\pgfqpoint{1.227555in}{1.587451in}}%
\pgfpathlineto{\pgfqpoint{1.224309in}{1.576380in}}%
\pgfpathlineto{\pgfqpoint{1.247672in}{1.573410in}}%
\pgfpathlineto{\pgfqpoint{1.246154in}{1.561183in}}%
\pgfpathlineto{\pgfqpoint{1.257251in}{1.559802in}}%
\pgfpathlineto{\pgfqpoint{1.255504in}{1.539978in}}%
\pgfpathlineto{\pgfqpoint{1.251044in}{1.505633in}}%
\pgfpathlineto{\pgfqpoint{1.245609in}{1.461365in}}%
\pgfpathclose%
\pgfusepath{fill}%
\end{pgfscope}%
\begin{pgfscope}%
\pgfpathrectangle{\pgfqpoint{0.100000in}{0.100000in}}{\pgfqpoint{3.007045in}{1.925000in}}%
\pgfusepath{clip}%
\pgfsetbuttcap%
\pgfsetmiterjoin%
\definecolor{currentfill}{rgb}{0.371688,0.649627,0.820515}%
\pgfsetfillcolor{currentfill}%
\pgfsetlinewidth{0.000000pt}%
\definecolor{currentstroke}{rgb}{0.000000,0.000000,0.000000}%
\pgfsetstrokecolor{currentstroke}%
\pgfsetstrokeopacity{0.000000}%
\pgfsetdash{}{0pt}%
\pgfpathmoveto{\pgfqpoint{1.836747in}{1.453219in}}%
\pgfpathlineto{\pgfqpoint{1.831065in}{1.453126in}}%
\pgfpathlineto{\pgfqpoint{1.830743in}{1.476124in}}%
\pgfpathlineto{\pgfqpoint{1.842251in}{1.476305in}}%
\pgfpathlineto{\pgfqpoint{1.841921in}{1.499291in}}%
\pgfpathlineto{\pgfqpoint{1.853375in}{1.499519in}}%
\pgfpathlineto{\pgfqpoint{1.853444in}{1.494748in}}%
\pgfpathlineto{\pgfqpoint{1.864898in}{1.494975in}}%
\pgfpathlineto{\pgfqpoint{1.865473in}{1.453759in}}%
\pgfpathlineto{\pgfqpoint{1.836747in}{1.453219in}}%
\pgfpathclose%
\pgfusepath{fill}%
\end{pgfscope}%
\begin{pgfscope}%
\pgfpathrectangle{\pgfqpoint{0.100000in}{0.100000in}}{\pgfqpoint{3.007045in}{1.925000in}}%
\pgfusepath{clip}%
\pgfsetbuttcap%
\pgfsetmiterjoin%
\definecolor{currentfill}{rgb}{0.351511,0.635848,0.812641}%
\pgfsetfillcolor{currentfill}%
\pgfsetlinewidth{0.000000pt}%
\definecolor{currentstroke}{rgb}{0.000000,0.000000,0.000000}%
\pgfsetstrokecolor{currentstroke}%
\pgfsetstrokeopacity{0.000000}%
\pgfsetdash{}{0pt}%
\pgfpathmoveto{\pgfqpoint{1.507527in}{1.120611in}}%
\pgfpathlineto{\pgfqpoint{1.536726in}{1.118994in}}%
\pgfpathlineto{\pgfqpoint{1.535329in}{1.090418in}}%
\pgfpathlineto{\pgfqpoint{1.502124in}{1.092214in}}%
\pgfpathlineto{\pgfqpoint{1.478511in}{1.093767in}}%
\pgfpathlineto{\pgfqpoint{1.480135in}{1.122427in}}%
\pgfpathlineto{\pgfqpoint{1.507527in}{1.120611in}}%
\pgfpathclose%
\pgfusepath{fill}%
\end{pgfscope}%
\begin{pgfscope}%
\pgfpathrectangle{\pgfqpoint{0.100000in}{0.100000in}}{\pgfqpoint{3.007045in}{1.925000in}}%
\pgfusepath{clip}%
\pgfsetbuttcap%
\pgfsetmiterjoin%
\definecolor{currentfill}{rgb}{0.059439,0.353572,0.639831}%
\pgfsetfillcolor{currentfill}%
\pgfsetlinewidth{0.000000pt}%
\definecolor{currentstroke}{rgb}{0.000000,0.000000,0.000000}%
\pgfsetstrokecolor{currentstroke}%
\pgfsetstrokeopacity{0.000000}%
\pgfsetdash{}{0pt}%
\pgfpathmoveto{\pgfqpoint{1.450280in}{1.611196in}}%
\pgfpathlineto{\pgfqpoint{1.455103in}{1.611872in}}%
\pgfpathlineto{\pgfqpoint{1.459052in}{1.608999in}}%
\pgfpathlineto{\pgfqpoint{1.470594in}{1.606845in}}%
\pgfpathlineto{\pgfqpoint{1.480364in}{1.610710in}}%
\pgfpathlineto{\pgfqpoint{1.482806in}{1.614036in}}%
\pgfpathlineto{\pgfqpoint{1.495994in}{1.623275in}}%
\pgfpathlineto{\pgfqpoint{1.503158in}{1.631150in}}%
\pgfpathlineto{\pgfqpoint{1.506325in}{1.629275in}}%
\pgfpathlineto{\pgfqpoint{1.516525in}{1.631523in}}%
\pgfpathlineto{\pgfqpoint{1.518046in}{1.627144in}}%
\pgfpathlineto{\pgfqpoint{1.517258in}{1.619082in}}%
\pgfpathlineto{\pgfqpoint{1.514093in}{1.614101in}}%
\pgfpathlineto{\pgfqpoint{1.514131in}{1.607335in}}%
\pgfpathlineto{\pgfqpoint{1.518744in}{1.599173in}}%
\pgfpathlineto{\pgfqpoint{1.523621in}{1.593929in}}%
\pgfpathlineto{\pgfqpoint{1.526579in}{1.583044in}}%
\pgfpathlineto{\pgfqpoint{1.522749in}{1.580062in}}%
\pgfpathlineto{\pgfqpoint{1.518319in}{1.571035in}}%
\pgfpathlineto{\pgfqpoint{1.524166in}{1.567570in}}%
\pgfpathlineto{\pgfqpoint{1.494665in}{1.569370in}}%
\pgfpathlineto{\pgfqpoint{1.447246in}{1.572779in}}%
\pgfpathlineto{\pgfqpoint{1.450280in}{1.611196in}}%
\pgfpathclose%
\pgfusepath{fill}%
\end{pgfscope}%
\begin{pgfscope}%
\pgfpathrectangle{\pgfqpoint{0.100000in}{0.100000in}}{\pgfqpoint{3.007045in}{1.925000in}}%
\pgfusepath{clip}%
\pgfsetbuttcap%
\pgfsetmiterjoin%
\definecolor{currentfill}{rgb}{0.778685,0.860300,0.937993}%
\pgfsetfillcolor{currentfill}%
\pgfsetlinewidth{0.000000pt}%
\definecolor{currentstroke}{rgb}{0.000000,0.000000,0.000000}%
\pgfsetstrokecolor{currentstroke}%
\pgfsetstrokeopacity{0.000000}%
\pgfsetdash{}{0pt}%
\pgfpathmoveto{\pgfqpoint{1.139523in}{1.071438in}}%
\pgfpathlineto{\pgfqpoint{1.121195in}{1.073917in}}%
\pgfpathlineto{\pgfqpoint{1.124305in}{1.096464in}}%
\pgfpathlineto{\pgfqpoint{1.123300in}{1.108295in}}%
\pgfpathlineto{\pgfqpoint{1.121625in}{1.109932in}}%
\pgfpathlineto{\pgfqpoint{1.123423in}{1.117299in}}%
\pgfpathlineto{\pgfqpoint{1.123455in}{1.127357in}}%
\pgfpathlineto{\pgfqpoint{1.121297in}{1.132386in}}%
\pgfpathlineto{\pgfqpoint{1.128210in}{1.131426in}}%
\pgfpathlineto{\pgfqpoint{1.136431in}{1.191196in}}%
\pgfpathlineto{\pgfqpoint{1.142160in}{1.192960in}}%
\pgfpathlineto{\pgfqpoint{1.143492in}{1.187243in}}%
\pgfpathlineto{\pgfqpoint{1.146599in}{1.183208in}}%
\pgfpathlineto{\pgfqpoint{1.157547in}{1.181738in}}%
\pgfpathlineto{\pgfqpoint{1.164471in}{1.172432in}}%
\pgfpathlineto{\pgfqpoint{1.169868in}{1.170865in}}%
\pgfpathlineto{\pgfqpoint{1.172957in}{1.174640in}}%
\pgfpathlineto{\pgfqpoint{1.177735in}{1.174146in}}%
\pgfpathlineto{\pgfqpoint{1.180201in}{1.170746in}}%
\pgfpathlineto{\pgfqpoint{1.186288in}{1.164150in}}%
\pgfpathlineto{\pgfqpoint{1.193507in}{1.162201in}}%
\pgfpathlineto{\pgfqpoint{1.188060in}{1.158199in}}%
\pgfpathlineto{\pgfqpoint{1.187185in}{1.151775in}}%
\pgfpathlineto{\pgfqpoint{1.189161in}{1.136735in}}%
\pgfpathlineto{\pgfqpoint{1.193414in}{1.130886in}}%
\pgfpathlineto{\pgfqpoint{1.154865in}{1.136033in}}%
\pgfpathlineto{\pgfqpoint{1.150792in}{1.105369in}}%
\pgfpathlineto{\pgfqpoint{1.143557in}{1.105192in}}%
\pgfpathlineto{\pgfqpoint{1.140887in}{1.088045in}}%
\pgfpathlineto{\pgfqpoint{1.139523in}{1.071438in}}%
\pgfpathclose%
\pgfusepath{fill}%
\end{pgfscope}%
\begin{pgfscope}%
\pgfpathrectangle{\pgfqpoint{0.100000in}{0.100000in}}{\pgfqpoint{3.007045in}{1.925000in}}%
\pgfusepath{clip}%
\pgfsetbuttcap%
\pgfsetmiterjoin%
\definecolor{currentfill}{rgb}{0.031373,0.285675,0.564291}%
\pgfsetfillcolor{currentfill}%
\pgfsetlinewidth{0.000000pt}%
\definecolor{currentstroke}{rgb}{0.000000,0.000000,0.000000}%
\pgfsetstrokecolor{currentstroke}%
\pgfsetstrokeopacity{0.000000}%
\pgfsetdash{}{0pt}%
\pgfpathmoveto{\pgfqpoint{1.447246in}{1.572779in}}%
\pgfpathlineto{\pgfqpoint{1.494665in}{1.569370in}}%
\pgfpathlineto{\pgfqpoint{1.524166in}{1.567570in}}%
\pgfpathlineto{\pgfqpoint{1.526953in}{1.561321in}}%
\pgfpathlineto{\pgfqpoint{1.524722in}{1.557253in}}%
\pgfpathlineto{\pgfqpoint{1.527075in}{1.552396in}}%
\pgfpathlineto{\pgfqpoint{1.523533in}{1.542968in}}%
\pgfpathlineto{\pgfqpoint{1.524818in}{1.537110in}}%
\pgfpathlineto{\pgfqpoint{1.518468in}{1.535998in}}%
\pgfpathlineto{\pgfqpoint{1.518970in}{1.529708in}}%
\pgfpathlineto{\pgfqpoint{1.519350in}{1.527670in}}%
\pgfpathlineto{\pgfqpoint{1.512562in}{1.521360in}}%
\pgfpathlineto{\pgfqpoint{1.506367in}{1.525118in}}%
\pgfpathlineto{\pgfqpoint{1.501256in}{1.521545in}}%
\pgfpathlineto{\pgfqpoint{1.498577in}{1.523374in}}%
\pgfpathlineto{\pgfqpoint{1.489275in}{1.520056in}}%
\pgfpathlineto{\pgfqpoint{1.484012in}{1.521937in}}%
\pgfpathlineto{\pgfqpoint{1.479161in}{1.518521in}}%
\pgfpathlineto{\pgfqpoint{1.473297in}{1.518824in}}%
\pgfpathlineto{\pgfqpoint{1.464688in}{1.511398in}}%
\pgfpathlineto{\pgfqpoint{1.458166in}{1.512912in}}%
\pgfpathlineto{\pgfqpoint{1.452730in}{1.510166in}}%
\pgfpathlineto{\pgfqpoint{1.442332in}{1.509325in}}%
\pgfpathlineto{\pgfqpoint{1.445034in}{1.544102in}}%
\pgfpathlineto{\pgfqpoint{1.447246in}{1.572779in}}%
\pgfpathclose%
\pgfusepath{fill}%
\end{pgfscope}%
\begin{pgfscope}%
\pgfpathrectangle{\pgfqpoint{0.100000in}{0.100000in}}{\pgfqpoint{3.007045in}{1.925000in}}%
\pgfusepath{clip}%
\pgfsetbuttcap%
\pgfsetmiterjoin%
\definecolor{currentfill}{rgb}{0.535686,0.746082,0.864252}%
\pgfsetfillcolor{currentfill}%
\pgfsetlinewidth{0.000000pt}%
\definecolor{currentstroke}{rgb}{0.000000,0.000000,0.000000}%
\pgfsetstrokecolor{currentstroke}%
\pgfsetstrokeopacity{0.000000}%
\pgfsetdash{}{0pt}%
\pgfpathmoveto{\pgfqpoint{2.379259in}{0.944820in}}%
\pgfpathlineto{\pgfqpoint{2.373522in}{0.943638in}}%
\pgfpathlineto{\pgfqpoint{2.364223in}{0.935952in}}%
\pgfpathlineto{\pgfqpoint{2.354830in}{0.942630in}}%
\pgfpathlineto{\pgfqpoint{2.341668in}{0.946890in}}%
\pgfpathlineto{\pgfqpoint{2.329679in}{0.943109in}}%
\pgfpathlineto{\pgfqpoint{2.327024in}{0.948030in}}%
\pgfpathlineto{\pgfqpoint{2.321351in}{0.950113in}}%
\pgfpathlineto{\pgfqpoint{2.317579in}{0.954327in}}%
\pgfpathlineto{\pgfqpoint{2.322488in}{0.960571in}}%
\pgfpathlineto{\pgfqpoint{2.319716in}{0.965996in}}%
\pgfpathlineto{\pgfqpoint{2.308631in}{0.975465in}}%
\pgfpathlineto{\pgfqpoint{2.309261in}{0.983763in}}%
\pgfpathlineto{\pgfqpoint{2.317969in}{0.990753in}}%
\pgfpathlineto{\pgfqpoint{2.325950in}{0.982785in}}%
\pgfpathlineto{\pgfqpoint{2.334580in}{0.979199in}}%
\pgfpathlineto{\pgfqpoint{2.338194in}{0.987373in}}%
\pgfpathlineto{\pgfqpoint{2.336026in}{1.001374in}}%
\pgfpathlineto{\pgfqpoint{2.340163in}{1.003243in}}%
\pgfpathlineto{\pgfqpoint{2.341210in}{1.008147in}}%
\pgfpathlineto{\pgfqpoint{2.353818in}{1.009495in}}%
\pgfpathlineto{\pgfqpoint{2.358542in}{0.999323in}}%
\pgfpathlineto{\pgfqpoint{2.361719in}{1.000172in}}%
\pgfpathlineto{\pgfqpoint{2.372662in}{0.995532in}}%
\pgfpathlineto{\pgfqpoint{2.370660in}{0.984801in}}%
\pgfpathlineto{\pgfqpoint{2.376395in}{0.971244in}}%
\pgfpathlineto{\pgfqpoint{2.369682in}{0.964768in}}%
\pgfpathlineto{\pgfqpoint{2.374820in}{0.956140in}}%
\pgfpathlineto{\pgfqpoint{2.378270in}{0.952775in}}%
\pgfpathlineto{\pgfqpoint{2.379259in}{0.944820in}}%
\pgfpathclose%
\pgfusepath{fill}%
\end{pgfscope}%
\begin{pgfscope}%
\pgfpathrectangle{\pgfqpoint{0.100000in}{0.100000in}}{\pgfqpoint{3.007045in}{1.925000in}}%
\pgfusepath{clip}%
\pgfsetbuttcap%
\pgfsetmiterjoin%
\definecolor{currentfill}{rgb}{0.351511,0.635848,0.812641}%
\pgfsetfillcolor{currentfill}%
\pgfsetlinewidth{0.000000pt}%
\definecolor{currentstroke}{rgb}{0.000000,0.000000,0.000000}%
\pgfsetstrokecolor{currentstroke}%
\pgfsetstrokeopacity{0.000000}%
\pgfsetdash{}{0pt}%
\pgfpathmoveto{\pgfqpoint{1.511195in}{1.206454in}}%
\pgfpathlineto{\pgfqpoint{1.509436in}{1.177873in}}%
\pgfpathlineto{\pgfqpoint{1.448195in}{1.181868in}}%
\pgfpathlineto{\pgfqpoint{1.446935in}{1.181960in}}%
\pgfpathlineto{\pgfqpoint{1.449096in}{1.210470in}}%
\pgfpathlineto{\pgfqpoint{1.453477in}{1.210184in}}%
\pgfpathlineto{\pgfqpoint{1.510401in}{1.206498in}}%
\pgfpathlineto{\pgfqpoint{1.511195in}{1.206454in}}%
\pgfpathclose%
\pgfusepath{fill}%
\end{pgfscope}%
\begin{pgfscope}%
\pgfpathrectangle{\pgfqpoint{0.100000in}{0.100000in}}{\pgfqpoint{3.007045in}{1.925000in}}%
\pgfusepath{clip}%
\pgfsetbuttcap%
\pgfsetmiterjoin%
\definecolor{currentfill}{rgb}{0.447843,0.697855,0.845306}%
\pgfsetfillcolor{currentfill}%
\pgfsetlinewidth{0.000000pt}%
\definecolor{currentstroke}{rgb}{0.000000,0.000000,0.000000}%
\pgfsetstrokecolor{currentstroke}%
\pgfsetstrokeopacity{0.000000}%
\pgfsetdash{}{0pt}%
\pgfpathmoveto{\pgfqpoint{2.585937in}{0.924191in}}%
\pgfpathlineto{\pgfqpoint{2.599449in}{0.926351in}}%
\pgfpathlineto{\pgfqpoint{2.613555in}{0.916458in}}%
\pgfpathlineto{\pgfqpoint{2.607323in}{0.893730in}}%
\pgfpathlineto{\pgfqpoint{2.595973in}{0.904310in}}%
\pgfpathlineto{\pgfqpoint{2.582994in}{0.902658in}}%
\pgfpathlineto{\pgfqpoint{2.571314in}{0.892287in}}%
\pgfpathlineto{\pgfqpoint{2.567423in}{0.900852in}}%
\pgfpathlineto{\pgfqpoint{2.562238in}{0.907450in}}%
\pgfpathlineto{\pgfqpoint{2.551891in}{0.919345in}}%
\pgfpathlineto{\pgfqpoint{2.585937in}{0.924191in}}%
\pgfpathclose%
\pgfusepath{fill}%
\end{pgfscope}%
\begin{pgfscope}%
\pgfpathrectangle{\pgfqpoint{0.100000in}{0.100000in}}{\pgfqpoint{3.007045in}{1.925000in}}%
\pgfusepath{clip}%
\pgfsetbuttcap%
\pgfsetmiterjoin%
\definecolor{currentfill}{rgb}{0.056363,0.349635,0.636755}%
\pgfsetfillcolor{currentfill}%
\pgfsetlinewidth{0.000000pt}%
\definecolor{currentstroke}{rgb}{0.000000,0.000000,0.000000}%
\pgfsetstrokecolor{currentstroke}%
\pgfsetstrokeopacity{0.000000}%
\pgfsetdash{}{0pt}%
\pgfpathmoveto{\pgfqpoint{1.473940in}{0.485183in}}%
\pgfpathlineto{\pgfqpoint{1.441506in}{0.486973in}}%
\pgfpathlineto{\pgfqpoint{1.435567in}{0.493311in}}%
\pgfpathlineto{\pgfqpoint{1.434480in}{0.498439in}}%
\pgfpathlineto{\pgfqpoint{1.422748in}{0.507253in}}%
\pgfpathlineto{\pgfqpoint{1.418771in}{0.514227in}}%
\pgfpathlineto{\pgfqpoint{1.415118in}{0.514942in}}%
\pgfpathlineto{\pgfqpoint{1.410285in}{0.521202in}}%
\pgfpathlineto{\pgfqpoint{1.406272in}{0.522948in}}%
\pgfpathlineto{\pgfqpoint{1.401030in}{0.533989in}}%
\pgfpathlineto{\pgfqpoint{1.394375in}{0.537885in}}%
\pgfpathlineto{\pgfqpoint{1.390298in}{0.536623in}}%
\pgfpathlineto{\pgfqpoint{1.369726in}{0.539824in}}%
\pgfpathlineto{\pgfqpoint{1.364524in}{0.539321in}}%
\pgfpathlineto{\pgfqpoint{1.349461in}{0.546748in}}%
\pgfpathlineto{\pgfqpoint{1.336243in}{0.559459in}}%
\pgfpathlineto{\pgfqpoint{1.337515in}{0.574775in}}%
\pgfpathlineto{\pgfqpoint{1.350354in}{0.573776in}}%
\pgfpathlineto{\pgfqpoint{1.352050in}{0.594662in}}%
\pgfpathlineto{\pgfqpoint{1.363729in}{0.593646in}}%
\pgfpathlineto{\pgfqpoint{1.364071in}{0.597523in}}%
\pgfpathlineto{\pgfqpoint{1.386971in}{0.595647in}}%
\pgfpathlineto{\pgfqpoint{1.390664in}{0.594316in}}%
\pgfpathlineto{\pgfqpoint{1.391278in}{0.588321in}}%
\pgfpathlineto{\pgfqpoint{1.385588in}{0.580234in}}%
\pgfpathlineto{\pgfqpoint{1.389923in}{0.575693in}}%
\pgfpathlineto{\pgfqpoint{1.383908in}{0.571430in}}%
\pgfpathlineto{\pgfqpoint{1.429685in}{0.568203in}}%
\pgfpathlineto{\pgfqpoint{1.444627in}{0.567276in}}%
\pgfpathlineto{\pgfqpoint{1.441878in}{0.523060in}}%
\pgfpathlineto{\pgfqpoint{1.475915in}{0.520991in}}%
\pgfpathlineto{\pgfqpoint{1.473940in}{0.485183in}}%
\pgfpathclose%
\pgfusepath{fill}%
\end{pgfscope}%
\begin{pgfscope}%
\pgfpathrectangle{\pgfqpoint{0.100000in}{0.100000in}}{\pgfqpoint{3.007045in}{1.925000in}}%
\pgfusepath{clip}%
\pgfsetbuttcap%
\pgfsetmiterjoin%
\definecolor{currentfill}{rgb}{0.647290,0.803922,0.892042}%
\pgfsetfillcolor{currentfill}%
\pgfsetlinewidth{0.000000pt}%
\definecolor{currentstroke}{rgb}{0.000000,0.000000,0.000000}%
\pgfsetstrokecolor{currentstroke}%
\pgfsetstrokeopacity{0.000000}%
\pgfsetdash{}{0pt}%
\pgfpathmoveto{\pgfqpoint{2.559174in}{1.248298in}}%
\pgfpathlineto{\pgfqpoint{2.559394in}{1.254953in}}%
\pgfpathlineto{\pgfqpoint{2.556847in}{1.259865in}}%
\pgfpathlineto{\pgfqpoint{2.561814in}{1.273283in}}%
\pgfpathlineto{\pgfqpoint{2.567920in}{1.283023in}}%
\pgfpathlineto{\pgfqpoint{2.574047in}{1.302973in}}%
\pgfpathlineto{\pgfqpoint{2.576932in}{1.318538in}}%
\pgfpathlineto{\pgfqpoint{2.599096in}{1.322608in}}%
\pgfpathlineto{\pgfqpoint{2.606011in}{1.320523in}}%
\pgfpathlineto{\pgfqpoint{2.609389in}{1.317404in}}%
\pgfpathlineto{\pgfqpoint{2.608640in}{1.313232in}}%
\pgfpathlineto{\pgfqpoint{2.612866in}{1.308899in}}%
\pgfpathlineto{\pgfqpoint{2.608820in}{1.295800in}}%
\pgfpathlineto{\pgfqpoint{2.610752in}{1.291249in}}%
\pgfpathlineto{\pgfqpoint{2.616493in}{1.288353in}}%
\pgfpathlineto{\pgfqpoint{2.613885in}{1.282571in}}%
\pgfpathlineto{\pgfqpoint{2.614370in}{1.276431in}}%
\pgfpathlineto{\pgfqpoint{2.612460in}{1.264832in}}%
\pgfpathlineto{\pgfqpoint{2.609634in}{1.257414in}}%
\pgfpathlineto{\pgfqpoint{2.559174in}{1.248298in}}%
\pgfpathclose%
\pgfusepath{fill}%
\end{pgfscope}%
\begin{pgfscope}%
\pgfpathrectangle{\pgfqpoint{0.100000in}{0.100000in}}{\pgfqpoint{3.007045in}{1.925000in}}%
\pgfusepath{clip}%
\pgfsetbuttcap%
\pgfsetmiterjoin%
\definecolor{currentfill}{rgb}{0.301069,0.601399,0.792957}%
\pgfsetfillcolor{currentfill}%
\pgfsetlinewidth{0.000000pt}%
\definecolor{currentstroke}{rgb}{0.000000,0.000000,0.000000}%
\pgfsetstrokecolor{currentstroke}%
\pgfsetstrokeopacity{0.000000}%
\pgfsetdash{}{0pt}%
\pgfpathmoveto{\pgfqpoint{1.963320in}{1.146179in}}%
\pgfpathlineto{\pgfqpoint{1.963877in}{1.126688in}}%
\pgfpathlineto{\pgfqpoint{1.956143in}{1.126474in}}%
\pgfpathlineto{\pgfqpoint{1.956351in}{1.117548in}}%
\pgfpathlineto{\pgfqpoint{1.949655in}{1.115080in}}%
\pgfpathlineto{\pgfqpoint{1.944978in}{1.116719in}}%
\pgfpathlineto{\pgfqpoint{1.944438in}{1.140114in}}%
\pgfpathlineto{\pgfqpoint{1.920082in}{1.139613in}}%
\pgfpathlineto{\pgfqpoint{1.920007in}{1.151193in}}%
\pgfpathlineto{\pgfqpoint{1.909292in}{1.151308in}}%
\pgfpathlineto{\pgfqpoint{1.909303in}{1.157973in}}%
\pgfpathlineto{\pgfqpoint{1.939480in}{1.158394in}}%
\pgfpathlineto{\pgfqpoint{1.953772in}{1.157587in}}%
\pgfpathlineto{\pgfqpoint{1.955717in}{1.146050in}}%
\pgfpathlineto{\pgfqpoint{1.963320in}{1.146179in}}%
\pgfpathclose%
\pgfusepath{fill}%
\end{pgfscope}%
\begin{pgfscope}%
\pgfpathrectangle{\pgfqpoint{0.100000in}{0.100000in}}{\pgfqpoint{3.007045in}{1.925000in}}%
\pgfusepath{clip}%
\pgfsetbuttcap%
\pgfsetmiterjoin%
\definecolor{currentfill}{rgb}{0.429020,0.687520,0.841246}%
\pgfsetfillcolor{currentfill}%
\pgfsetlinewidth{0.000000pt}%
\definecolor{currentstroke}{rgb}{0.000000,0.000000,0.000000}%
\pgfsetstrokecolor{currentstroke}%
\pgfsetstrokeopacity{0.000000}%
\pgfsetdash{}{0pt}%
\pgfpathmoveto{\pgfqpoint{1.835366in}{1.168017in}}%
\pgfpathlineto{\pgfqpoint{1.812537in}{1.167939in}}%
\pgfpathlineto{\pgfqpoint{1.812471in}{1.163163in}}%
\pgfpathlineto{\pgfqpoint{1.792627in}{1.163111in}}%
\pgfpathlineto{\pgfqpoint{1.792583in}{1.168136in}}%
\pgfpathlineto{\pgfqpoint{1.767101in}{1.168173in}}%
\pgfpathlineto{\pgfqpoint{1.769538in}{1.174111in}}%
\pgfpathlineto{\pgfqpoint{1.766486in}{1.176026in}}%
\pgfpathlineto{\pgfqpoint{1.755095in}{1.176082in}}%
\pgfpathlineto{\pgfqpoint{1.755176in}{1.198941in}}%
\pgfpathlineto{\pgfqpoint{1.756771in}{1.198935in}}%
\pgfpathlineto{\pgfqpoint{1.762123in}{1.192600in}}%
\pgfpathlineto{\pgfqpoint{1.768029in}{1.189813in}}%
\pgfpathlineto{\pgfqpoint{1.772661in}{1.193312in}}%
\pgfpathlineto{\pgfqpoint{1.770162in}{1.207494in}}%
\pgfpathlineto{\pgfqpoint{1.792177in}{1.207263in}}%
\pgfpathlineto{\pgfqpoint{1.792230in}{1.201536in}}%
\pgfpathlineto{\pgfqpoint{1.811728in}{1.201377in}}%
\pgfpathlineto{\pgfqpoint{1.811789in}{1.208062in}}%
\pgfpathlineto{\pgfqpoint{1.834576in}{1.208087in}}%
\pgfpathlineto{\pgfqpoint{1.834849in}{1.196667in}}%
\pgfpathlineto{\pgfqpoint{1.835366in}{1.168017in}}%
\pgfpathclose%
\pgfusepath{fill}%
\end{pgfscope}%
\begin{pgfscope}%
\pgfpathrectangle{\pgfqpoint{0.100000in}{0.100000in}}{\pgfqpoint{3.007045in}{1.925000in}}%
\pgfusepath{clip}%
\pgfsetbuttcap%
\pgfsetmiterjoin%
\definecolor{currentfill}{rgb}{0.406997,0.673741,0.834295}%
\pgfsetfillcolor{currentfill}%
\pgfsetlinewidth{0.000000pt}%
\definecolor{currentstroke}{rgb}{0.000000,0.000000,0.000000}%
\pgfsetstrokecolor{currentstroke}%
\pgfsetstrokeopacity{0.000000}%
\pgfsetdash{}{0pt}%
\pgfpathmoveto{\pgfqpoint{2.068022in}{0.977336in}}%
\pgfpathlineto{\pgfqpoint{2.064440in}{0.977028in}}%
\pgfpathlineto{\pgfqpoint{2.061582in}{0.976889in}}%
\pgfpathlineto{\pgfqpoint{2.061897in}{0.966247in}}%
\pgfpathlineto{\pgfqpoint{2.059156in}{0.970494in}}%
\pgfpathlineto{\pgfqpoint{2.050739in}{0.968865in}}%
\pgfpathlineto{\pgfqpoint{2.039044in}{0.968405in}}%
\pgfpathlineto{\pgfqpoint{2.038888in}{0.984228in}}%
\pgfpathlineto{\pgfqpoint{2.028875in}{0.983778in}}%
\pgfpathlineto{\pgfqpoint{2.028045in}{0.989826in}}%
\pgfpathlineto{\pgfqpoint{2.022047in}{1.002653in}}%
\pgfpathlineto{\pgfqpoint{2.029271in}{1.013848in}}%
\pgfpathlineto{\pgfqpoint{2.023343in}{1.013571in}}%
\pgfpathlineto{\pgfqpoint{2.022844in}{1.028625in}}%
\pgfpathlineto{\pgfqpoint{2.026620in}{1.028636in}}%
\pgfpathlineto{\pgfqpoint{2.025561in}{1.047408in}}%
\pgfpathlineto{\pgfqpoint{2.053627in}{1.049022in}}%
\pgfpathlineto{\pgfqpoint{2.058254in}{1.047208in}}%
\pgfpathlineto{\pgfqpoint{2.064198in}{1.036075in}}%
\pgfpathlineto{\pgfqpoint{2.059538in}{1.028875in}}%
\pgfpathlineto{\pgfqpoint{2.067443in}{1.014097in}}%
\pgfpathlineto{\pgfqpoint{2.072603in}{1.010193in}}%
\pgfpathlineto{\pgfqpoint{2.077300in}{1.012074in}}%
\pgfpathlineto{\pgfqpoint{2.081022in}{1.010070in}}%
\pgfpathlineto{\pgfqpoint{2.082880in}{1.008816in}}%
\pgfpathlineto{\pgfqpoint{2.081317in}{1.001182in}}%
\pgfpathlineto{\pgfqpoint{2.082341in}{0.997150in}}%
\pgfpathlineto{\pgfqpoint{2.078631in}{0.992336in}}%
\pgfpathlineto{\pgfqpoint{2.080971in}{0.989204in}}%
\pgfpathlineto{\pgfqpoint{2.077777in}{0.982593in}}%
\pgfpathlineto{\pgfqpoint{2.070320in}{0.985858in}}%
\pgfpathlineto{\pgfqpoint{2.068022in}{0.977336in}}%
\pgfpathclose%
\pgfusepath{fill}%
\end{pgfscope}%
\begin{pgfscope}%
\pgfpathrectangle{\pgfqpoint{0.100000in}{0.100000in}}{\pgfqpoint{3.007045in}{1.925000in}}%
\pgfusepath{clip}%
\pgfsetbuttcap%
\pgfsetmiterjoin%
\definecolor{currentfill}{rgb}{0.171027,0.484752,0.731242}%
\pgfsetfillcolor{currentfill}%
\pgfsetlinewidth{0.000000pt}%
\definecolor{currentstroke}{rgb}{0.000000,0.000000,0.000000}%
\pgfsetstrokecolor{currentstroke}%
\pgfsetstrokeopacity{0.000000}%
\pgfsetdash{}{0pt}%
\pgfpathmoveto{\pgfqpoint{2.139795in}{0.676959in}}%
\pgfpathlineto{\pgfqpoint{2.141224in}{0.664018in}}%
\pgfpathlineto{\pgfqpoint{2.143322in}{0.646719in}}%
\pgfpathlineto{\pgfqpoint{2.121151in}{0.645024in}}%
\pgfpathlineto{\pgfqpoint{2.089615in}{0.643002in}}%
\pgfpathlineto{\pgfqpoint{2.088016in}{0.667046in}}%
\pgfpathlineto{\pgfqpoint{2.073893in}{0.665116in}}%
\pgfpathlineto{\pgfqpoint{2.072910in}{0.680552in}}%
\pgfpathlineto{\pgfqpoint{2.068533in}{0.682516in}}%
\pgfpathlineto{\pgfqpoint{2.067842in}{0.694087in}}%
\pgfpathlineto{\pgfqpoint{2.090671in}{0.695575in}}%
\pgfpathlineto{\pgfqpoint{2.113560in}{0.697106in}}%
\pgfpathlineto{\pgfqpoint{2.115516in}{0.670799in}}%
\pgfpathlineto{\pgfqpoint{2.130566in}{0.673645in}}%
\pgfpathlineto{\pgfqpoint{2.139795in}{0.676959in}}%
\pgfpathclose%
\pgfusepath{fill}%
\end{pgfscope}%
\begin{pgfscope}%
\pgfpathrectangle{\pgfqpoint{0.100000in}{0.100000in}}{\pgfqpoint{3.007045in}{1.925000in}}%
\pgfusepath{clip}%
\pgfsetbuttcap%
\pgfsetmiterjoin%
\definecolor{currentfill}{rgb}{0.183206,0.496932,0.737516}%
\pgfsetfillcolor{currentfill}%
\pgfsetlinewidth{0.000000pt}%
\definecolor{currentstroke}{rgb}{0.000000,0.000000,0.000000}%
\pgfsetstrokecolor{currentstroke}%
\pgfsetstrokeopacity{0.000000}%
\pgfsetdash{}{0pt}%
\pgfpathmoveto{\pgfqpoint{1.314888in}{1.740131in}}%
\pgfpathlineto{\pgfqpoint{1.309172in}{1.740759in}}%
\pgfpathlineto{\pgfqpoint{1.307173in}{1.723478in}}%
\pgfpathlineto{\pgfqpoint{1.305247in}{1.723687in}}%
\pgfpathlineto{\pgfqpoint{1.302614in}{1.700688in}}%
\pgfpathlineto{\pgfqpoint{1.283442in}{1.702901in}}%
\pgfpathlineto{\pgfqpoint{1.282748in}{1.697095in}}%
\pgfpathlineto{\pgfqpoint{1.277023in}{1.697772in}}%
\pgfpathlineto{\pgfqpoint{1.277730in}{1.703596in}}%
\pgfpathlineto{\pgfqpoint{1.272031in}{1.704281in}}%
\pgfpathlineto{\pgfqpoint{1.268284in}{1.704741in}}%
\pgfpathlineto{\pgfqpoint{1.269690in}{1.716224in}}%
\pgfpathlineto{\pgfqpoint{1.263981in}{1.716912in}}%
\pgfpathlineto{\pgfqpoint{1.267101in}{1.728242in}}%
\pgfpathlineto{\pgfqpoint{1.269873in}{1.750687in}}%
\pgfpathlineto{\pgfqpoint{1.265560in}{1.751233in}}%
\pgfpathlineto{\pgfqpoint{1.263629in}{1.757459in}}%
\pgfpathlineto{\pgfqpoint{1.265605in}{1.764425in}}%
\pgfpathlineto{\pgfqpoint{1.272584in}{1.760976in}}%
\pgfpathlineto{\pgfqpoint{1.289529in}{1.757826in}}%
\pgfpathlineto{\pgfqpoint{1.291202in}{1.771653in}}%
\pgfpathlineto{\pgfqpoint{1.292930in}{1.771456in}}%
\pgfpathlineto{\pgfqpoint{1.295653in}{1.794279in}}%
\pgfpathlineto{\pgfqpoint{1.332164in}{1.790080in}}%
\pgfpathlineto{\pgfqpoint{1.341649in}{1.789040in}}%
\pgfpathlineto{\pgfqpoint{1.341017in}{1.783290in}}%
\pgfpathlineto{\pgfqpoint{1.346744in}{1.782663in}}%
\pgfpathlineto{\pgfqpoint{1.346116in}{1.776879in}}%
\pgfpathlineto{\pgfqpoint{1.371755in}{1.774248in}}%
\pgfpathlineto{\pgfqpoint{1.369270in}{1.748112in}}%
\pgfpathlineto{\pgfqpoint{1.366019in}{1.751221in}}%
\pgfpathlineto{\pgfqpoint{1.355902in}{1.752996in}}%
\pgfpathlineto{\pgfqpoint{1.351556in}{1.758299in}}%
\pgfpathlineto{\pgfqpoint{1.341818in}{1.758120in}}%
\pgfpathlineto{\pgfqpoint{1.335725in}{1.759869in}}%
\pgfpathlineto{\pgfqpoint{1.324957in}{1.756945in}}%
\pgfpathlineto{\pgfqpoint{1.318715in}{1.758150in}}%
\pgfpathlineto{\pgfqpoint{1.317232in}{1.745061in}}%
\pgfpathlineto{\pgfqpoint{1.314888in}{1.740131in}}%
\pgfpathclose%
\pgfusepath{fill}%
\end{pgfscope}%
\begin{pgfscope}%
\pgfpathrectangle{\pgfqpoint{0.100000in}{0.100000in}}{\pgfqpoint{3.007045in}{1.925000in}}%
\pgfusepath{clip}%
\pgfsetbuttcap%
\pgfsetmiterjoin%
\definecolor{currentfill}{rgb}{0.031373,0.261315,0.528120}%
\pgfsetfillcolor{currentfill}%
\pgfsetlinewidth{0.000000pt}%
\definecolor{currentstroke}{rgb}{0.000000,0.000000,0.000000}%
\pgfsetstrokecolor{currentstroke}%
\pgfsetstrokeopacity{0.000000}%
\pgfsetdash{}{0pt}%
\pgfpathmoveto{\pgfqpoint{1.512921in}{0.424052in}}%
\pgfpathlineto{\pgfqpoint{1.470461in}{0.425819in}}%
\pgfpathlineto{\pgfqpoint{1.473940in}{0.485183in}}%
\pgfpathlineto{\pgfqpoint{1.514641in}{0.483384in}}%
\pgfpathlineto{\pgfqpoint{1.513462in}{0.453257in}}%
\pgfpathlineto{\pgfqpoint{1.514242in}{0.453233in}}%
\pgfpathlineto{\pgfqpoint{1.512921in}{0.424052in}}%
\pgfpathclose%
\pgfusepath{fill}%
\end{pgfscope}%
\begin{pgfscope}%
\pgfpathrectangle{\pgfqpoint{0.100000in}{0.100000in}}{\pgfqpoint{3.007045in}{1.925000in}}%
\pgfusepath{clip}%
\pgfsetbuttcap%
\pgfsetmiterjoin%
\definecolor{currentfill}{rgb}{0.031373,0.285675,0.564291}%
\pgfsetfillcolor{currentfill}%
\pgfsetlinewidth{0.000000pt}%
\definecolor{currentstroke}{rgb}{0.000000,0.000000,0.000000}%
\pgfsetstrokecolor{currentstroke}%
\pgfsetstrokeopacity{0.000000}%
\pgfsetdash{}{0pt}%
\pgfpathmoveto{\pgfqpoint{1.219256in}{0.700832in}}%
\pgfpathlineto{\pgfqpoint{1.265329in}{0.695996in}}%
\pgfpathlineto{\pgfqpoint{1.255027in}{0.637292in}}%
\pgfpathlineto{\pgfqpoint{1.205698in}{0.612940in}}%
\pgfpathlineto{\pgfqpoint{1.207983in}{0.628535in}}%
\pgfpathlineto{\pgfqpoint{1.212900in}{0.676457in}}%
\pgfpathlineto{\pgfqpoint{1.215312in}{0.701261in}}%
\pgfpathlineto{\pgfqpoint{1.219256in}{0.700832in}}%
\pgfpathclose%
\pgfusepath{fill}%
\end{pgfscope}%
\begin{pgfscope}%
\pgfpathrectangle{\pgfqpoint{0.100000in}{0.100000in}}{\pgfqpoint{3.007045in}{1.925000in}}%
\pgfusepath{clip}%
\pgfsetbuttcap%
\pgfsetmiterjoin%
\definecolor{currentfill}{rgb}{0.311157,0.608289,0.796894}%
\pgfsetfillcolor{currentfill}%
\pgfsetlinewidth{0.000000pt}%
\definecolor{currentstroke}{rgb}{0.000000,0.000000,0.000000}%
\pgfsetstrokecolor{currentstroke}%
\pgfsetstrokeopacity{0.000000}%
\pgfsetdash{}{0pt}%
\pgfpathmoveto{\pgfqpoint{1.347270in}{1.020215in}}%
\pgfpathlineto{\pgfqpoint{1.351668in}{1.019814in}}%
\pgfpathlineto{\pgfqpoint{1.348759in}{0.987002in}}%
\pgfpathlineto{\pgfqpoint{1.346674in}{0.987189in}}%
\pgfpathlineto{\pgfqpoint{1.342266in}{0.937180in}}%
\pgfpathlineto{\pgfqpoint{1.324480in}{0.938840in}}%
\pgfpathlineto{\pgfqpoint{1.327114in}{0.961542in}}%
\pgfpathlineto{\pgfqpoint{1.304327in}{0.963661in}}%
\pgfpathlineto{\pgfqpoint{1.304909in}{0.969450in}}%
\pgfpathlineto{\pgfqpoint{1.293521in}{0.970602in}}%
\pgfpathlineto{\pgfqpoint{1.293808in}{0.973457in}}%
\pgfpathlineto{\pgfqpoint{1.298942in}{1.024507in}}%
\pgfpathlineto{\pgfqpoint{1.347270in}{1.020215in}}%
\pgfpathclose%
\pgfusepath{fill}%
\end{pgfscope}%
\begin{pgfscope}%
\pgfpathrectangle{\pgfqpoint{0.100000in}{0.100000in}}{\pgfqpoint{3.007045in}{1.925000in}}%
\pgfusepath{clip}%
\pgfsetbuttcap%
\pgfsetmiterjoin%
\definecolor{currentfill}{rgb}{0.093272,0.396878,0.673664}%
\pgfsetfillcolor{currentfill}%
\pgfsetlinewidth{0.000000pt}%
\definecolor{currentstroke}{rgb}{0.000000,0.000000,0.000000}%
\pgfsetstrokecolor{currentstroke}%
\pgfsetstrokeopacity{0.000000}%
\pgfsetdash{}{0pt}%
\pgfpathmoveto{\pgfqpoint{0.610456in}{0.538425in}}%
\pgfpathlineto{\pgfqpoint{0.607882in}{0.541289in}}%
\pgfpathlineto{\pgfqpoint{0.608049in}{0.546486in}}%
\pgfpathlineto{\pgfqpoint{0.607286in}{0.548897in}}%
\pgfpathlineto{\pgfqpoint{0.610143in}{0.549575in}}%
\pgfpathlineto{\pgfqpoint{0.608906in}{0.546215in}}%
\pgfpathlineto{\pgfqpoint{0.609173in}{0.543568in}}%
\pgfpathlineto{\pgfqpoint{0.611044in}{0.540017in}}%
\pgfpathlineto{\pgfqpoint{0.610456in}{0.538425in}}%
\pgfpathclose%
\pgfusepath{fill}%
\end{pgfscope}%
\begin{pgfscope}%
\pgfpathrectangle{\pgfqpoint{0.100000in}{0.100000in}}{\pgfqpoint{3.007045in}{1.925000in}}%
\pgfusepath{clip}%
\pgfsetbuttcap%
\pgfsetmiterjoin%
\definecolor{currentfill}{rgb}{0.093272,0.396878,0.673664}%
\pgfsetfillcolor{currentfill}%
\pgfsetlinewidth{0.000000pt}%
\definecolor{currentstroke}{rgb}{0.000000,0.000000,0.000000}%
\pgfsetstrokecolor{currentstroke}%
\pgfsetstrokeopacity{0.000000}%
\pgfsetdash{}{0pt}%
\pgfpathmoveto{\pgfqpoint{0.652105in}{0.476446in}}%
\pgfpathlineto{\pgfqpoint{0.651830in}{0.477944in}}%
\pgfpathlineto{\pgfqpoint{0.649670in}{0.479666in}}%
\pgfpathlineto{\pgfqpoint{0.648834in}{0.482194in}}%
\pgfpathlineto{\pgfqpoint{0.648697in}{0.484544in}}%
\pgfpathlineto{\pgfqpoint{0.647047in}{0.487579in}}%
\pgfpathlineto{\pgfqpoint{0.647342in}{0.491406in}}%
\pgfpathlineto{\pgfqpoint{0.648359in}{0.493435in}}%
\pgfpathlineto{\pgfqpoint{0.649744in}{0.493109in}}%
\pgfpathlineto{\pgfqpoint{0.652856in}{0.488673in}}%
\pgfpathlineto{\pgfqpoint{0.656744in}{0.489158in}}%
\pgfpathlineto{\pgfqpoint{0.659223in}{0.488519in}}%
\pgfpathlineto{\pgfqpoint{0.659628in}{0.487236in}}%
\pgfpathlineto{\pgfqpoint{0.661768in}{0.486250in}}%
\pgfpathlineto{\pgfqpoint{0.661612in}{0.484269in}}%
\pgfpathlineto{\pgfqpoint{0.663795in}{0.482269in}}%
\pgfpathlineto{\pgfqpoint{0.663710in}{0.479661in}}%
\pgfpathlineto{\pgfqpoint{0.661944in}{0.477630in}}%
\pgfpathlineto{\pgfqpoint{0.660195in}{0.476581in}}%
\pgfpathlineto{\pgfqpoint{0.659577in}{0.472416in}}%
\pgfpathlineto{\pgfqpoint{0.657595in}{0.473029in}}%
\pgfpathlineto{\pgfqpoint{0.655097in}{0.475183in}}%
\pgfpathlineto{\pgfqpoint{0.652105in}{0.476446in}}%
\pgfpathclose%
\pgfusepath{fill}%
\end{pgfscope}%
\begin{pgfscope}%
\pgfpathrectangle{\pgfqpoint{0.100000in}{0.100000in}}{\pgfqpoint{3.007045in}{1.925000in}}%
\pgfusepath{clip}%
\pgfsetbuttcap%
\pgfsetmiterjoin%
\definecolor{currentfill}{rgb}{0.093272,0.396878,0.673664}%
\pgfsetfillcolor{currentfill}%
\pgfsetlinewidth{0.000000pt}%
\definecolor{currentstroke}{rgb}{0.000000,0.000000,0.000000}%
\pgfsetstrokecolor{currentstroke}%
\pgfsetstrokeopacity{0.000000}%
\pgfsetdash{}{0pt}%
\pgfpathmoveto{\pgfqpoint{0.809857in}{0.421160in}}%
\pgfpathlineto{\pgfqpoint{0.796605in}{0.403352in}}%
\pgfpathlineto{\pgfqpoint{0.796295in}{0.405189in}}%
\pgfpathlineto{\pgfqpoint{0.792854in}{0.407923in}}%
\pgfpathlineto{\pgfqpoint{0.789615in}{0.403627in}}%
\pgfpathlineto{\pgfqpoint{0.788663in}{0.401187in}}%
\pgfpathlineto{\pgfqpoint{0.784950in}{0.396276in}}%
\pgfpathlineto{\pgfqpoint{0.775612in}{0.403305in}}%
\pgfpathlineto{\pgfqpoint{0.761232in}{0.414538in}}%
\pgfpathlineto{\pgfqpoint{0.750064in}{0.423440in}}%
\pgfpathlineto{\pgfqpoint{0.751507in}{0.425233in}}%
\pgfpathlineto{\pgfqpoint{0.748296in}{0.427861in}}%
\pgfpathlineto{\pgfqpoint{0.746839in}{0.426097in}}%
\pgfpathlineto{\pgfqpoint{0.745138in}{0.427348in}}%
\pgfpathlineto{\pgfqpoint{0.743807in}{0.425686in}}%
\pgfpathlineto{\pgfqpoint{0.742145in}{0.427074in}}%
\pgfpathlineto{\pgfqpoint{0.743447in}{0.428714in}}%
\pgfpathlineto{\pgfqpoint{0.734378in}{0.436156in}}%
\pgfpathlineto{\pgfqpoint{0.733009in}{0.434561in}}%
\pgfpathlineto{\pgfqpoint{0.731767in}{0.435552in}}%
\pgfpathlineto{\pgfqpoint{0.730325in}{0.433863in}}%
\pgfpathlineto{\pgfqpoint{0.728779in}{0.435116in}}%
\pgfpathlineto{\pgfqpoint{0.726974in}{0.433888in}}%
\pgfpathlineto{\pgfqpoint{0.724234in}{0.430706in}}%
\pgfpathlineto{\pgfqpoint{0.722451in}{0.432254in}}%
\pgfpathlineto{\pgfqpoint{0.721076in}{0.430661in}}%
\pgfpathlineto{\pgfqpoint{0.719488in}{0.432032in}}%
\pgfpathlineto{\pgfqpoint{0.716770in}{0.428913in}}%
\pgfpathlineto{\pgfqpoint{0.715184in}{0.430242in}}%
\pgfpathlineto{\pgfqpoint{0.713230in}{0.429162in}}%
\pgfpathlineto{\pgfqpoint{0.710582in}{0.426191in}}%
\pgfpathlineto{\pgfqpoint{0.709062in}{0.427491in}}%
\pgfpathlineto{\pgfqpoint{0.706317in}{0.424353in}}%
\pgfpathlineto{\pgfqpoint{0.704186in}{0.426137in}}%
\pgfpathlineto{\pgfqpoint{0.701434in}{0.423011in}}%
\pgfpathlineto{\pgfqpoint{0.699846in}{0.424401in}}%
\pgfpathlineto{\pgfqpoint{0.697072in}{0.421245in}}%
\pgfpathlineto{\pgfqpoint{0.695145in}{0.420096in}}%
\pgfpathlineto{\pgfqpoint{0.693623in}{0.421487in}}%
\pgfpathlineto{\pgfqpoint{0.691710in}{0.420379in}}%
\pgfpathlineto{\pgfqpoint{0.688938in}{0.417253in}}%
\pgfpathlineto{\pgfqpoint{0.687384in}{0.418587in}}%
\pgfpathlineto{\pgfqpoint{0.686127in}{0.417159in}}%
\pgfpathlineto{\pgfqpoint{0.684613in}{0.418474in}}%
\pgfpathlineto{\pgfqpoint{0.681584in}{0.415593in}}%
\pgfpathlineto{\pgfqpoint{0.681431in}{0.415610in}}%
\pgfpathlineto{\pgfqpoint{0.678710in}{0.418051in}}%
\pgfpathlineto{\pgfqpoint{0.677618in}{0.416243in}}%
\pgfpathlineto{\pgfqpoint{0.676726in}{0.415156in}}%
\pgfpathlineto{\pgfqpoint{0.670768in}{0.416525in}}%
\pgfpathlineto{\pgfqpoint{0.671401in}{0.419366in}}%
\pgfpathlineto{\pgfqpoint{0.674143in}{0.420400in}}%
\pgfpathlineto{\pgfqpoint{0.677137in}{0.423980in}}%
\pgfpathlineto{\pgfqpoint{0.678234in}{0.426008in}}%
\pgfpathlineto{\pgfqpoint{0.678253in}{0.428853in}}%
\pgfpathlineto{\pgfqpoint{0.679315in}{0.431241in}}%
\pgfpathlineto{\pgfqpoint{0.680432in}{0.430540in}}%
\pgfpathlineto{\pgfqpoint{0.681892in}{0.432397in}}%
\pgfpathlineto{\pgfqpoint{0.685663in}{0.432497in}}%
\pgfpathlineto{\pgfqpoint{0.686784in}{0.434521in}}%
\pgfpathlineto{\pgfqpoint{0.688419in}{0.439211in}}%
\pgfpathlineto{\pgfqpoint{0.688889in}{0.441645in}}%
\pgfpathlineto{\pgfqpoint{0.690056in}{0.443550in}}%
\pgfpathlineto{\pgfqpoint{0.690186in}{0.446179in}}%
\pgfpathlineto{\pgfqpoint{0.692680in}{0.447617in}}%
\pgfpathlineto{\pgfqpoint{0.692600in}{0.450868in}}%
\pgfpathlineto{\pgfqpoint{0.690991in}{0.451746in}}%
\pgfpathlineto{\pgfqpoint{0.688110in}{0.448749in}}%
\pgfpathlineto{\pgfqpoint{0.685848in}{0.448554in}}%
\pgfpathlineto{\pgfqpoint{0.684889in}{0.449901in}}%
\pgfpathlineto{\pgfqpoint{0.680145in}{0.450430in}}%
\pgfpathlineto{\pgfqpoint{0.675467in}{0.452772in}}%
\pgfpathlineto{\pgfqpoint{0.673367in}{0.454382in}}%
\pgfpathlineto{\pgfqpoint{0.670995in}{0.457386in}}%
\pgfpathlineto{\pgfqpoint{0.670409in}{0.459079in}}%
\pgfpathlineto{\pgfqpoint{0.671394in}{0.461167in}}%
\pgfpathlineto{\pgfqpoint{0.672846in}{0.461844in}}%
\pgfpathlineto{\pgfqpoint{0.672427in}{0.464293in}}%
\pgfpathlineto{\pgfqpoint{0.672655in}{0.467035in}}%
\pgfpathlineto{\pgfqpoint{0.673992in}{0.467451in}}%
\pgfpathlineto{\pgfqpoint{0.672827in}{0.471052in}}%
\pgfpathlineto{\pgfqpoint{0.673952in}{0.471769in}}%
\pgfpathlineto{\pgfqpoint{0.671748in}{0.473544in}}%
\pgfpathlineto{\pgfqpoint{0.671417in}{0.476039in}}%
\pgfpathlineto{\pgfqpoint{0.672034in}{0.478136in}}%
\pgfpathlineto{\pgfqpoint{0.673606in}{0.477783in}}%
\pgfpathlineto{\pgfqpoint{0.674154in}{0.479291in}}%
\pgfpathlineto{\pgfqpoint{0.671865in}{0.479944in}}%
\pgfpathlineto{\pgfqpoint{0.671055in}{0.482481in}}%
\pgfpathlineto{\pgfqpoint{0.674291in}{0.481708in}}%
\pgfpathlineto{\pgfqpoint{0.678236in}{0.483360in}}%
\pgfpathlineto{\pgfqpoint{0.680888in}{0.483413in}}%
\pgfpathlineto{\pgfqpoint{0.679590in}{0.486507in}}%
\pgfpathlineto{\pgfqpoint{0.680253in}{0.488047in}}%
\pgfpathlineto{\pgfqpoint{0.682631in}{0.488869in}}%
\pgfpathlineto{\pgfqpoint{0.682850in}{0.492468in}}%
\pgfpathlineto{\pgfqpoint{0.680782in}{0.491150in}}%
\pgfpathlineto{\pgfqpoint{0.679720in}{0.494316in}}%
\pgfpathlineto{\pgfqpoint{0.682387in}{0.496907in}}%
\pgfpathlineto{\pgfqpoint{0.680696in}{0.499993in}}%
\pgfpathlineto{\pgfqpoint{0.681929in}{0.501414in}}%
\pgfpathlineto{\pgfqpoint{0.683631in}{0.501161in}}%
\pgfpathlineto{\pgfqpoint{0.684549in}{0.502655in}}%
\pgfpathlineto{\pgfqpoint{0.684007in}{0.504296in}}%
\pgfpathlineto{\pgfqpoint{0.681235in}{0.504565in}}%
\pgfpathlineto{\pgfqpoint{0.682519in}{0.506916in}}%
\pgfpathlineto{\pgfqpoint{0.686294in}{0.507512in}}%
\pgfpathlineto{\pgfqpoint{0.687166in}{0.510160in}}%
\pgfpathlineto{\pgfqpoint{0.689933in}{0.507661in}}%
\pgfpathlineto{\pgfqpoint{0.692175in}{0.508261in}}%
\pgfpathlineto{\pgfqpoint{0.693238in}{0.510668in}}%
\pgfpathlineto{\pgfqpoint{0.695178in}{0.511813in}}%
\pgfpathlineto{\pgfqpoint{0.704163in}{0.514029in}}%
\pgfpathlineto{\pgfqpoint{0.708310in}{0.514388in}}%
\pgfpathlineto{\pgfqpoint{0.709700in}{0.512813in}}%
\pgfpathlineto{\pgfqpoint{0.712038in}{0.513756in}}%
\pgfpathlineto{\pgfqpoint{0.712689in}{0.516016in}}%
\pgfpathlineto{\pgfqpoint{0.717386in}{0.519817in}}%
\pgfpathlineto{\pgfqpoint{0.721882in}{0.521512in}}%
\pgfpathlineto{\pgfqpoint{0.726704in}{0.521249in}}%
\pgfpathlineto{\pgfqpoint{0.729612in}{0.518760in}}%
\pgfpathlineto{\pgfqpoint{0.730444in}{0.514976in}}%
\pgfpathlineto{\pgfqpoint{0.730775in}{0.511205in}}%
\pgfpathlineto{\pgfqpoint{0.731879in}{0.508992in}}%
\pgfpathlineto{\pgfqpoint{0.734192in}{0.507876in}}%
\pgfpathlineto{\pgfqpoint{0.738215in}{0.508943in}}%
\pgfpathlineto{\pgfqpoint{0.741050in}{0.508543in}}%
\pgfpathlineto{\pgfqpoint{0.742738in}{0.506896in}}%
\pgfpathlineto{\pgfqpoint{0.741312in}{0.505279in}}%
\pgfpathlineto{\pgfqpoint{0.742884in}{0.503844in}}%
\pgfpathlineto{\pgfqpoint{0.741434in}{0.502134in}}%
\pgfpathlineto{\pgfqpoint{0.742883in}{0.500799in}}%
\pgfpathlineto{\pgfqpoint{0.741587in}{0.499176in}}%
\pgfpathlineto{\pgfqpoint{0.751165in}{0.490590in}}%
\pgfpathlineto{\pgfqpoint{0.750170in}{0.491012in}}%
\pgfpathlineto{\pgfqpoint{0.748071in}{0.489521in}}%
\pgfpathlineto{\pgfqpoint{0.742309in}{0.483156in}}%
\pgfpathlineto{\pgfqpoint{0.741905in}{0.483472in}}%
\pgfpathlineto{\pgfqpoint{0.736206in}{0.477104in}}%
\pgfpathlineto{\pgfqpoint{0.737384in}{0.476047in}}%
\pgfpathlineto{\pgfqpoint{0.733232in}{0.471329in}}%
\pgfpathlineto{\pgfqpoint{0.741635in}{0.463832in}}%
\pgfpathlineto{\pgfqpoint{0.747517in}{0.458771in}}%
\pgfpathlineto{\pgfqpoint{0.749271in}{0.460813in}}%
\pgfpathlineto{\pgfqpoint{0.755991in}{0.455118in}}%
\pgfpathlineto{\pgfqpoint{0.757361in}{0.456738in}}%
\pgfpathlineto{\pgfqpoint{0.770089in}{0.446056in}}%
\pgfpathlineto{\pgfqpoint{0.768555in}{0.444177in}}%
\pgfpathlineto{\pgfqpoint{0.776021in}{0.437944in}}%
\pgfpathlineto{\pgfqpoint{0.787335in}{0.429069in}}%
\pgfpathlineto{\pgfqpoint{0.796550in}{0.422002in}}%
\pgfpathlineto{\pgfqpoint{0.798438in}{0.420763in}}%
\pgfpathlineto{\pgfqpoint{0.800378in}{0.423260in}}%
\pgfpathlineto{\pgfqpoint{0.803696in}{0.420629in}}%
\pgfpathlineto{\pgfqpoint{0.805001in}{0.422296in}}%
\pgfpathlineto{\pgfqpoint{0.807974in}{0.419972in}}%
\pgfpathlineto{\pgfqpoint{0.809857in}{0.421160in}}%
\pgfpathclose%
\pgfusepath{fill}%
\end{pgfscope}%
\begin{pgfscope}%
\pgfpathrectangle{\pgfqpoint{0.100000in}{0.100000in}}{\pgfqpoint{3.007045in}{1.925000in}}%
\pgfusepath{clip}%
\pgfsetbuttcap%
\pgfsetmiterjoin%
\definecolor{currentfill}{rgb}{0.491765,0.721968,0.854779}%
\pgfsetfillcolor{currentfill}%
\pgfsetlinewidth{0.000000pt}%
\definecolor{currentstroke}{rgb}{0.000000,0.000000,0.000000}%
\pgfsetstrokecolor{currentstroke}%
\pgfsetstrokeopacity{0.000000}%
\pgfsetdash{}{0pt}%
\pgfpathmoveto{\pgfqpoint{0.860069in}{1.641241in}}%
\pgfpathlineto{\pgfqpoint{0.857696in}{1.638685in}}%
\pgfpathlineto{\pgfqpoint{0.853704in}{1.621066in}}%
\pgfpathlineto{\pgfqpoint{0.841346in}{1.613085in}}%
\pgfpathlineto{\pgfqpoint{0.834363in}{1.617027in}}%
\pgfpathlineto{\pgfqpoint{0.830622in}{1.610723in}}%
\pgfpathlineto{\pgfqpoint{0.824810in}{1.609423in}}%
\pgfpathlineto{\pgfqpoint{0.822960in}{1.604421in}}%
\pgfpathlineto{\pgfqpoint{0.825412in}{1.600994in}}%
\pgfpathlineto{\pgfqpoint{0.824171in}{1.596209in}}%
\pgfpathlineto{\pgfqpoint{0.820142in}{1.591987in}}%
\pgfpathlineto{\pgfqpoint{0.811192in}{1.587476in}}%
\pgfpathlineto{\pgfqpoint{0.808189in}{1.587646in}}%
\pgfpathlineto{\pgfqpoint{0.781320in}{1.594007in}}%
\pgfpathlineto{\pgfqpoint{0.779546in}{1.588394in}}%
\pgfpathlineto{\pgfqpoint{0.774977in}{1.589715in}}%
\pgfpathlineto{\pgfqpoint{0.777625in}{1.600841in}}%
\pgfpathlineto{\pgfqpoint{0.780379in}{1.600195in}}%
\pgfpathlineto{\pgfqpoint{0.783037in}{1.611433in}}%
\pgfpathlineto{\pgfqpoint{0.775430in}{1.609514in}}%
\pgfpathlineto{\pgfqpoint{0.770478in}{1.610681in}}%
\pgfpathlineto{\pgfqpoint{0.765709in}{1.615018in}}%
\pgfpathlineto{\pgfqpoint{0.767582in}{1.623021in}}%
\pgfpathlineto{\pgfqpoint{0.763942in}{1.627183in}}%
\pgfpathlineto{\pgfqpoint{0.766659in}{1.638492in}}%
\pgfpathlineto{\pgfqpoint{0.754151in}{1.641671in}}%
\pgfpathlineto{\pgfqpoint{0.758646in}{1.647388in}}%
\pgfpathlineto{\pgfqpoint{0.759381in}{1.652797in}}%
\pgfpathlineto{\pgfqpoint{0.763180in}{1.655715in}}%
\pgfpathlineto{\pgfqpoint{0.766636in}{1.659171in}}%
\pgfpathlineto{\pgfqpoint{0.770539in}{1.666797in}}%
\pgfpathlineto{\pgfqpoint{0.786144in}{1.663045in}}%
\pgfpathlineto{\pgfqpoint{0.786641in}{1.653470in}}%
\pgfpathlineto{\pgfqpoint{0.792690in}{1.650559in}}%
\pgfpathlineto{\pgfqpoint{0.802276in}{1.654267in}}%
\pgfpathlineto{\pgfqpoint{0.860069in}{1.641241in}}%
\pgfpathclose%
\pgfusepath{fill}%
\end{pgfscope}%
\begin{pgfscope}%
\pgfpathrectangle{\pgfqpoint{0.100000in}{0.100000in}}{\pgfqpoint{3.007045in}{1.925000in}}%
\pgfusepath{clip}%
\pgfsetbuttcap%
\pgfsetmiterjoin%
\definecolor{currentfill}{rgb}{0.548235,0.752972,0.866959}%
\pgfsetfillcolor{currentfill}%
\pgfsetlinewidth{0.000000pt}%
\definecolor{currentstroke}{rgb}{0.000000,0.000000,0.000000}%
\pgfsetstrokecolor{currentstroke}%
\pgfsetstrokeopacity{0.000000}%
\pgfsetdash{}{0pt}%
\pgfpathmoveto{\pgfqpoint{2.081022in}{1.010070in}}%
\pgfpathlineto{\pgfqpoint{2.077300in}{1.012074in}}%
\pgfpathlineto{\pgfqpoint{2.072603in}{1.010193in}}%
\pgfpathlineto{\pgfqpoint{2.067443in}{1.014097in}}%
\pgfpathlineto{\pgfqpoint{2.059538in}{1.028875in}}%
\pgfpathlineto{\pgfqpoint{2.064198in}{1.036075in}}%
\pgfpathlineto{\pgfqpoint{2.058254in}{1.047208in}}%
\pgfpathlineto{\pgfqpoint{2.060389in}{1.049493in}}%
\pgfpathlineto{\pgfqpoint{2.058159in}{1.055368in}}%
\pgfpathlineto{\pgfqpoint{2.054473in}{1.056687in}}%
\pgfpathlineto{\pgfqpoint{2.045818in}{1.065000in}}%
\pgfpathlineto{\pgfqpoint{2.040262in}{1.068526in}}%
\pgfpathlineto{\pgfqpoint{2.034604in}{1.066814in}}%
\pgfpathlineto{\pgfqpoint{2.031875in}{1.071974in}}%
\pgfpathlineto{\pgfqpoint{2.024754in}{1.078038in}}%
\pgfpathlineto{\pgfqpoint{2.020734in}{1.079515in}}%
\pgfpathlineto{\pgfqpoint{2.029399in}{1.083124in}}%
\pgfpathlineto{\pgfqpoint{2.029064in}{1.088879in}}%
\pgfpathlineto{\pgfqpoint{2.046290in}{1.089601in}}%
\pgfpathlineto{\pgfqpoint{2.075029in}{1.090942in}}%
\pgfpathlineto{\pgfqpoint{2.075282in}{1.085153in}}%
\pgfpathlineto{\pgfqpoint{2.098280in}{1.086724in}}%
\pgfpathlineto{\pgfqpoint{2.099179in}{1.072352in}}%
\pgfpathlineto{\pgfqpoint{2.100489in}{1.052130in}}%
\pgfpathlineto{\pgfqpoint{2.101574in}{1.034888in}}%
\pgfpathlineto{\pgfqpoint{2.091233in}{1.034083in}}%
\pgfpathlineto{\pgfqpoint{2.090698in}{1.026843in}}%
\pgfpathlineto{\pgfqpoint{2.085427in}{1.025478in}}%
\pgfpathlineto{\pgfqpoint{2.078167in}{1.013748in}}%
\pgfpathlineto{\pgfqpoint{2.081022in}{1.010070in}}%
\pgfpathclose%
\pgfusepath{fill}%
\end{pgfscope}%
\begin{pgfscope}%
\pgfpathrectangle{\pgfqpoint{0.100000in}{0.100000in}}{\pgfqpoint{3.007045in}{1.925000in}}%
\pgfusepath{clip}%
\pgfsetbuttcap%
\pgfsetmiterjoin%
\definecolor{currentfill}{rgb}{0.401953,0.670296,0.832326}%
\pgfsetfillcolor{currentfill}%
\pgfsetlinewidth{0.000000pt}%
\definecolor{currentstroke}{rgb}{0.000000,0.000000,0.000000}%
\pgfsetstrokecolor{currentstroke}%
\pgfsetstrokeopacity{0.000000}%
\pgfsetdash{}{0pt}%
\pgfpathmoveto{\pgfqpoint{2.721216in}{1.279405in}}%
\pgfpathlineto{\pgfqpoint{2.716371in}{1.278367in}}%
\pgfpathlineto{\pgfqpoint{2.716037in}{1.278304in}}%
\pgfpathlineto{\pgfqpoint{2.706007in}{1.289043in}}%
\pgfpathlineto{\pgfqpoint{2.701481in}{1.290448in}}%
\pgfpathlineto{\pgfqpoint{2.697015in}{1.296543in}}%
\pgfpathlineto{\pgfqpoint{2.690460in}{1.296019in}}%
\pgfpathlineto{\pgfqpoint{2.687398in}{1.300784in}}%
\pgfpathlineto{\pgfqpoint{2.693385in}{1.305552in}}%
\pgfpathlineto{\pgfqpoint{2.684366in}{1.322430in}}%
\pgfpathlineto{\pgfqpoint{2.690299in}{1.329057in}}%
\pgfpathlineto{\pgfqpoint{2.680799in}{1.334029in}}%
\pgfpathlineto{\pgfqpoint{2.694956in}{1.344848in}}%
\pgfpathlineto{\pgfqpoint{2.698167in}{1.347287in}}%
\pgfpathlineto{\pgfqpoint{2.698167in}{1.352791in}}%
\pgfpathlineto{\pgfqpoint{2.701015in}{1.357842in}}%
\pgfpathlineto{\pgfqpoint{2.711768in}{1.357635in}}%
\pgfpathlineto{\pgfqpoint{2.725895in}{1.348689in}}%
\pgfpathlineto{\pgfqpoint{2.720167in}{1.343542in}}%
\pgfpathlineto{\pgfqpoint{2.741068in}{1.332531in}}%
\pgfpathlineto{\pgfqpoint{2.735728in}{1.317531in}}%
\pgfpathlineto{\pgfqpoint{2.725428in}{1.306338in}}%
\pgfpathlineto{\pgfqpoint{2.726798in}{1.302007in}}%
\pgfpathlineto{\pgfqpoint{2.725049in}{1.295828in}}%
\pgfpathlineto{\pgfqpoint{2.725778in}{1.289195in}}%
\pgfpathlineto{\pgfqpoint{2.724100in}{1.282597in}}%
\pgfpathlineto{\pgfqpoint{2.721216in}{1.279405in}}%
\pgfpathclose%
\pgfusepath{fill}%
\end{pgfscope}%
\begin{pgfscope}%
\pgfpathrectangle{\pgfqpoint{0.100000in}{0.100000in}}{\pgfqpoint{3.007045in}{1.925000in}}%
\pgfusepath{clip}%
\pgfsetbuttcap%
\pgfsetmiterjoin%
\definecolor{currentfill}{rgb}{0.248166,0.561892,0.770980}%
\pgfsetfillcolor{currentfill}%
\pgfsetlinewidth{0.000000pt}%
\definecolor{currentstroke}{rgb}{0.000000,0.000000,0.000000}%
\pgfsetstrokecolor{currentstroke}%
\pgfsetstrokeopacity{0.000000}%
\pgfsetdash{}{0pt}%
\pgfpathmoveto{\pgfqpoint{2.422704in}{0.621971in}}%
\pgfpathlineto{\pgfqpoint{2.415134in}{0.621502in}}%
\pgfpathlineto{\pgfqpoint{2.412441in}{0.646597in}}%
\pgfpathlineto{\pgfqpoint{2.397277in}{0.644995in}}%
\pgfpathlineto{\pgfqpoint{2.397002in}{0.647345in}}%
\pgfpathlineto{\pgfqpoint{2.395188in}{0.664380in}}%
\pgfpathlineto{\pgfqpoint{2.414666in}{0.666517in}}%
\pgfpathlineto{\pgfqpoint{2.412962in}{0.682194in}}%
\pgfpathlineto{\pgfqpoint{2.408654in}{0.685257in}}%
\pgfpathlineto{\pgfqpoint{2.404240in}{0.684768in}}%
\pgfpathlineto{\pgfqpoint{2.402438in}{0.696626in}}%
\pgfpathlineto{\pgfqpoint{2.413072in}{0.697996in}}%
\pgfpathlineto{\pgfqpoint{2.412679in}{0.701290in}}%
\pgfpathlineto{\pgfqpoint{2.420123in}{0.701779in}}%
\pgfpathlineto{\pgfqpoint{2.422816in}{0.692785in}}%
\pgfpathlineto{\pgfqpoint{2.420802in}{0.691745in}}%
\pgfpathlineto{\pgfqpoint{2.421279in}{0.683375in}}%
\pgfpathlineto{\pgfqpoint{2.425943in}{0.682862in}}%
\pgfpathlineto{\pgfqpoint{2.428996in}{0.678445in}}%
\pgfpathlineto{\pgfqpoint{2.442052in}{0.679496in}}%
\pgfpathlineto{\pgfqpoint{2.443165in}{0.674743in}}%
\pgfpathlineto{\pgfqpoint{2.450078in}{0.667343in}}%
\pgfpathlineto{\pgfqpoint{2.450133in}{0.661231in}}%
\pgfpathlineto{\pgfqpoint{2.454390in}{0.661837in}}%
\pgfpathlineto{\pgfqpoint{2.457148in}{0.641063in}}%
\pgfpathlineto{\pgfqpoint{2.464523in}{0.639748in}}%
\pgfpathlineto{\pgfqpoint{2.469936in}{0.633604in}}%
\pgfpathlineto{\pgfqpoint{2.479352in}{0.633603in}}%
\pgfpathlineto{\pgfqpoint{2.481583in}{0.625736in}}%
\pgfpathlineto{\pgfqpoint{2.422704in}{0.621971in}}%
\pgfpathclose%
\pgfusepath{fill}%
\end{pgfscope}%
\begin{pgfscope}%
\pgfpathrectangle{\pgfqpoint{0.100000in}{0.100000in}}{\pgfqpoint{3.007045in}{1.925000in}}%
\pgfusepath{clip}%
\pgfsetbuttcap%
\pgfsetmiterjoin%
\definecolor{currentfill}{rgb}{0.356555,0.639293,0.814610}%
\pgfsetfillcolor{currentfill}%
\pgfsetlinewidth{0.000000pt}%
\definecolor{currentstroke}{rgb}{0.000000,0.000000,0.000000}%
\pgfsetstrokecolor{currentstroke}%
\pgfsetstrokeopacity{0.000000}%
\pgfsetdash{}{0pt}%
\pgfpathmoveto{\pgfqpoint{2.030481in}{0.763778in}}%
\pgfpathlineto{\pgfqpoint{2.023847in}{0.761087in}}%
\pgfpathlineto{\pgfqpoint{2.023674in}{0.765200in}}%
\pgfpathlineto{\pgfqpoint{2.022165in}{0.797050in}}%
\pgfpathlineto{\pgfqpoint{2.036938in}{0.797675in}}%
\pgfpathlineto{\pgfqpoint{2.030424in}{0.792915in}}%
\pgfpathlineto{\pgfqpoint{2.039761in}{0.792018in}}%
\pgfpathlineto{\pgfqpoint{2.039467in}{0.797809in}}%
\pgfpathlineto{\pgfqpoint{2.050947in}{0.799400in}}%
\pgfpathlineto{\pgfqpoint{2.050685in}{0.804294in}}%
\pgfpathlineto{\pgfqpoint{2.059401in}{0.802761in}}%
\pgfpathlineto{\pgfqpoint{2.073753in}{0.803617in}}%
\pgfpathlineto{\pgfqpoint{2.074344in}{0.794001in}}%
\pgfpathlineto{\pgfqpoint{2.075429in}{0.776700in}}%
\pgfpathlineto{\pgfqpoint{2.078426in}{0.775122in}}%
\pgfpathlineto{\pgfqpoint{2.079034in}{0.765390in}}%
\pgfpathlineto{\pgfqpoint{2.065349in}{0.764544in}}%
\pgfpathlineto{\pgfqpoint{2.062975in}{0.759697in}}%
\pgfpathlineto{\pgfqpoint{2.039047in}{0.766251in}}%
\pgfpathlineto{\pgfqpoint{2.033226in}{0.766351in}}%
\pgfpathlineto{\pgfqpoint{2.030481in}{0.763778in}}%
\pgfpathclose%
\pgfusepath{fill}%
\end{pgfscope}%
\begin{pgfscope}%
\pgfpathrectangle{\pgfqpoint{0.100000in}{0.100000in}}{\pgfqpoint{3.007045in}{1.925000in}}%
\pgfusepath{clip}%
\pgfsetbuttcap%
\pgfsetmiterjoin%
\definecolor{currentfill}{rgb}{0.207566,0.521292,0.750065}%
\pgfsetfillcolor{currentfill}%
\pgfsetlinewidth{0.000000pt}%
\definecolor{currentstroke}{rgb}{0.000000,0.000000,0.000000}%
\pgfsetstrokecolor{currentstroke}%
\pgfsetstrokeopacity{0.000000}%
\pgfsetdash{}{0pt}%
\pgfpathmoveto{\pgfqpoint{1.611326in}{0.679134in}}%
\pgfpathlineto{\pgfqpoint{1.607101in}{0.688968in}}%
\pgfpathlineto{\pgfqpoint{1.600832in}{0.707630in}}%
\pgfpathlineto{\pgfqpoint{1.578119in}{0.708542in}}%
\pgfpathlineto{\pgfqpoint{1.572479in}{0.708902in}}%
\pgfpathlineto{\pgfqpoint{1.573659in}{0.737980in}}%
\pgfpathlineto{\pgfqpoint{1.581896in}{0.737526in}}%
\pgfpathlineto{\pgfqpoint{1.582071in}{0.741206in}}%
\pgfpathlineto{\pgfqpoint{1.610081in}{0.739788in}}%
\pgfpathlineto{\pgfqpoint{1.611193in}{0.768407in}}%
\pgfpathlineto{\pgfqpoint{1.607913in}{0.770715in}}%
\pgfpathlineto{\pgfqpoint{1.608887in}{0.798669in}}%
\pgfpathlineto{\pgfqpoint{1.614953in}{0.795791in}}%
\pgfpathlineto{\pgfqpoint{1.625906in}{0.804863in}}%
\pgfpathlineto{\pgfqpoint{1.631785in}{0.798613in}}%
\pgfpathlineto{\pgfqpoint{1.636022in}{0.799536in}}%
\pgfpathlineto{\pgfqpoint{1.634992in}{0.767696in}}%
\pgfpathlineto{\pgfqpoint{1.640734in}{0.767480in}}%
\pgfpathlineto{\pgfqpoint{1.639106in}{0.738203in}}%
\pgfpathlineto{\pgfqpoint{1.659429in}{0.737338in}}%
\pgfpathlineto{\pgfqpoint{1.658448in}{0.708465in}}%
\pgfpathlineto{\pgfqpoint{1.655743in}{0.708579in}}%
\pgfpathlineto{\pgfqpoint{1.655346in}{0.689754in}}%
\pgfpathlineto{\pgfqpoint{1.633343in}{0.684748in}}%
\pgfpathlineto{\pgfqpoint{1.628904in}{0.682713in}}%
\pgfpathlineto{\pgfqpoint{1.624705in}{0.686669in}}%
\pgfpathlineto{\pgfqpoint{1.611326in}{0.679134in}}%
\pgfpathclose%
\pgfusepath{fill}%
\end{pgfscope}%
\begin{pgfscope}%
\pgfpathrectangle{\pgfqpoint{0.100000in}{0.100000in}}{\pgfqpoint{3.007045in}{1.925000in}}%
\pgfusepath{clip}%
\pgfsetbuttcap%
\pgfsetmiterjoin%
\definecolor{currentfill}{rgb}{0.366644,0.646182,0.818547}%
\pgfsetfillcolor{currentfill}%
\pgfsetlinewidth{0.000000pt}%
\definecolor{currentstroke}{rgb}{0.000000,0.000000,0.000000}%
\pgfsetstrokecolor{currentstroke}%
\pgfsetstrokeopacity{0.000000}%
\pgfsetdash{}{0pt}%
\pgfpathmoveto{\pgfqpoint{2.265651in}{0.893385in}}%
\pgfpathlineto{\pgfqpoint{2.241623in}{0.891209in}}%
\pgfpathlineto{\pgfqpoint{2.240394in}{0.900117in}}%
\pgfpathlineto{\pgfqpoint{2.237551in}{0.900340in}}%
\pgfpathlineto{\pgfqpoint{2.235740in}{0.906849in}}%
\pgfpathlineto{\pgfqpoint{2.223823in}{0.914321in}}%
\pgfpathlineto{\pgfqpoint{2.220248in}{0.918653in}}%
\pgfpathlineto{\pgfqpoint{2.219736in}{0.935214in}}%
\pgfpathlineto{\pgfqpoint{2.226333in}{0.935993in}}%
\pgfpathlineto{\pgfqpoint{2.228851in}{0.933436in}}%
\pgfpathlineto{\pgfqpoint{2.240739in}{0.937396in}}%
\pgfpathlineto{\pgfqpoint{2.251280in}{0.935449in}}%
\pgfpathlineto{\pgfqpoint{2.254813in}{0.936868in}}%
\pgfpathlineto{\pgfqpoint{2.257324in}{0.931388in}}%
\pgfpathlineto{\pgfqpoint{2.261477in}{0.928516in}}%
\pgfpathlineto{\pgfqpoint{2.262417in}{0.916982in}}%
\pgfpathlineto{\pgfqpoint{2.261069in}{0.912869in}}%
\pgfpathlineto{\pgfqpoint{2.264140in}{0.908633in}}%
\pgfpathlineto{\pgfqpoint{2.265651in}{0.893385in}}%
\pgfpathclose%
\pgfusepath{fill}%
\end{pgfscope}%
\begin{pgfscope}%
\pgfpathrectangle{\pgfqpoint{0.100000in}{0.100000in}}{\pgfqpoint{3.007045in}{1.925000in}}%
\pgfusepath{clip}%
\pgfsetbuttcap%
\pgfsetmiterjoin%
\definecolor{currentfill}{rgb}{0.516863,0.735748,0.860192}%
\pgfsetfillcolor{currentfill}%
\pgfsetlinewidth{0.000000pt}%
\definecolor{currentstroke}{rgb}{0.000000,0.000000,0.000000}%
\pgfsetstrokecolor{currentstroke}%
\pgfsetstrokeopacity{0.000000}%
\pgfsetdash{}{0pt}%
\pgfpathmoveto{\pgfqpoint{2.141820in}{0.745927in}}%
\pgfpathlineto{\pgfqpoint{2.116053in}{0.743899in}}%
\pgfpathlineto{\pgfqpoint{2.114414in}{0.767709in}}%
\pgfpathlineto{\pgfqpoint{2.122360in}{0.768291in}}%
\pgfpathlineto{\pgfqpoint{2.121230in}{0.782758in}}%
\pgfpathlineto{\pgfqpoint{2.128772in}{0.785353in}}%
\pgfpathlineto{\pgfqpoint{2.131570in}{0.789934in}}%
\pgfpathlineto{\pgfqpoint{2.129211in}{0.792870in}}%
\pgfpathlineto{\pgfqpoint{2.136914in}{0.795797in}}%
\pgfpathlineto{\pgfqpoint{2.139121in}{0.799917in}}%
\pgfpathlineto{\pgfqpoint{2.143299in}{0.800190in}}%
\pgfpathlineto{\pgfqpoint{2.142921in}{0.786166in}}%
\pgfpathlineto{\pgfqpoint{2.166832in}{0.787387in}}%
\pgfpathlineto{\pgfqpoint{2.168941in}{0.762907in}}%
\pgfpathlineto{\pgfqpoint{2.163129in}{0.759429in}}%
\pgfpathlineto{\pgfqpoint{2.158791in}{0.754738in}}%
\pgfpathlineto{\pgfqpoint{2.154101in}{0.754354in}}%
\pgfpathlineto{\pgfqpoint{2.151248in}{0.751063in}}%
\pgfpathlineto{\pgfqpoint{2.141932in}{0.750061in}}%
\pgfpathlineto{\pgfqpoint{2.141820in}{0.745927in}}%
\pgfpathclose%
\pgfusepath{fill}%
\end{pgfscope}%
\begin{pgfscope}%
\pgfpathrectangle{\pgfqpoint{0.100000in}{0.100000in}}{\pgfqpoint{3.007045in}{1.925000in}}%
\pgfusepath{clip}%
\pgfsetbuttcap%
\pgfsetmiterjoin%
\definecolor{currentfill}{rgb}{0.657132,0.808105,0.895486}%
\pgfsetfillcolor{currentfill}%
\pgfsetlinewidth{0.000000pt}%
\definecolor{currentstroke}{rgb}{0.000000,0.000000,0.000000}%
\pgfsetstrokecolor{currentstroke}%
\pgfsetstrokeopacity{0.000000}%
\pgfsetdash{}{0pt}%
\pgfpathmoveto{\pgfqpoint{2.695962in}{1.404095in}}%
\pgfpathlineto{\pgfqpoint{2.689741in}{1.426141in}}%
\pgfpathlineto{\pgfqpoint{2.669988in}{1.422124in}}%
\pgfpathlineto{\pgfqpoint{2.667190in}{1.440512in}}%
\pgfpathlineto{\pgfqpoint{2.672580in}{1.444093in}}%
\pgfpathlineto{\pgfqpoint{2.680763in}{1.444336in}}%
\pgfpathlineto{\pgfqpoint{2.679109in}{1.451464in}}%
\pgfpathlineto{\pgfqpoint{2.697553in}{1.455914in}}%
\pgfpathlineto{\pgfqpoint{2.700669in}{1.446056in}}%
\pgfpathlineto{\pgfqpoint{2.710631in}{1.447455in}}%
\pgfpathlineto{\pgfqpoint{2.711428in}{1.444046in}}%
\pgfpathlineto{\pgfqpoint{2.721879in}{1.446280in}}%
\pgfpathlineto{\pgfqpoint{2.723883in}{1.436359in}}%
\pgfpathlineto{\pgfqpoint{2.727454in}{1.434295in}}%
\pgfpathlineto{\pgfqpoint{2.721527in}{1.432985in}}%
\pgfpathlineto{\pgfqpoint{2.727566in}{1.410162in}}%
\pgfpathlineto{\pgfqpoint{2.695962in}{1.404095in}}%
\pgfpathclose%
\pgfusepath{fill}%
\end{pgfscope}%
\begin{pgfscope}%
\pgfpathrectangle{\pgfqpoint{0.100000in}{0.100000in}}{\pgfqpoint{3.007045in}{1.925000in}}%
\pgfusepath{clip}%
\pgfsetbuttcap%
\pgfsetmiterjoin%
\definecolor{currentfill}{rgb}{0.316201,0.611734,0.798862}%
\pgfsetfillcolor{currentfill}%
\pgfsetlinewidth{0.000000pt}%
\definecolor{currentstroke}{rgb}{0.000000,0.000000,0.000000}%
\pgfsetstrokecolor{currentstroke}%
\pgfsetstrokeopacity{0.000000}%
\pgfsetdash{}{0pt}%
\pgfpathmoveto{\pgfqpoint{1.700709in}{0.760907in}}%
\pgfpathlineto{\pgfqpoint{1.695830in}{0.764016in}}%
\pgfpathlineto{\pgfqpoint{1.670614in}{0.764978in}}%
\pgfpathlineto{\pgfqpoint{1.664958in}{0.765812in}}%
\pgfpathlineto{\pgfqpoint{1.665670in}{0.801052in}}%
\pgfpathlineto{\pgfqpoint{1.668937in}{0.795613in}}%
\pgfpathlineto{\pgfqpoint{1.678257in}{0.793022in}}%
\pgfpathlineto{\pgfqpoint{1.683354in}{0.804224in}}%
\pgfpathlineto{\pgfqpoint{1.695332in}{0.814267in}}%
\pgfpathlineto{\pgfqpoint{1.718089in}{0.813933in}}%
\pgfpathlineto{\pgfqpoint{1.720987in}{0.811838in}}%
\pgfpathlineto{\pgfqpoint{1.720879in}{0.800114in}}%
\pgfpathlineto{\pgfqpoint{1.726173in}{0.795162in}}%
\pgfpathlineto{\pgfqpoint{1.730486in}{0.795042in}}%
\pgfpathlineto{\pgfqpoint{1.725337in}{0.792070in}}%
\pgfpathlineto{\pgfqpoint{1.724887in}{0.764449in}}%
\pgfpathlineto{\pgfqpoint{1.700709in}{0.760907in}}%
\pgfpathclose%
\pgfusepath{fill}%
\end{pgfscope}%
\begin{pgfscope}%
\pgfpathrectangle{\pgfqpoint{0.100000in}{0.100000in}}{\pgfqpoint{3.007045in}{1.925000in}}%
\pgfusepath{clip}%
\pgfsetbuttcap%
\pgfsetmiterjoin%
\definecolor{currentfill}{rgb}{0.396909,0.666851,0.830358}%
\pgfsetfillcolor{currentfill}%
\pgfsetlinewidth{0.000000pt}%
\definecolor{currentstroke}{rgb}{0.000000,0.000000,0.000000}%
\pgfsetstrokecolor{currentstroke}%
\pgfsetstrokeopacity{0.000000}%
\pgfsetdash{}{0pt}%
\pgfpathmoveto{\pgfqpoint{2.151467in}{0.577078in}}%
\pgfpathlineto{\pgfqpoint{2.146781in}{0.617810in}}%
\pgfpathlineto{\pgfqpoint{2.143322in}{0.646719in}}%
\pgfpathlineto{\pgfqpoint{2.141224in}{0.664018in}}%
\pgfpathlineto{\pgfqpoint{2.162408in}{0.665683in}}%
\pgfpathlineto{\pgfqpoint{2.161161in}{0.669683in}}%
\pgfpathlineto{\pgfqpoint{2.156603in}{0.673169in}}%
\pgfpathlineto{\pgfqpoint{2.161771in}{0.684999in}}%
\pgfpathlineto{\pgfqpoint{2.184529in}{0.686866in}}%
\pgfpathlineto{\pgfqpoint{2.186978in}{0.688005in}}%
\pgfpathlineto{\pgfqpoint{2.186226in}{0.696389in}}%
\pgfpathlineto{\pgfqpoint{2.191937in}{0.696888in}}%
\pgfpathlineto{\pgfqpoint{2.191190in}{0.704948in}}%
\pgfpathlineto{\pgfqpoint{2.193997in}{0.705851in}}%
\pgfpathlineto{\pgfqpoint{2.202255in}{0.702095in}}%
\pgfpathlineto{\pgfqpoint{2.211673in}{0.692874in}}%
\pgfpathlineto{\pgfqpoint{2.226858in}{0.694238in}}%
\pgfpathlineto{\pgfqpoint{2.228144in}{0.679871in}}%
\pgfpathlineto{\pgfqpoint{2.194774in}{0.676909in}}%
\pgfpathlineto{\pgfqpoint{2.194637in}{0.668142in}}%
\pgfpathlineto{\pgfqpoint{2.191786in}{0.667860in}}%
\pgfpathlineto{\pgfqpoint{2.192871in}{0.650628in}}%
\pgfpathlineto{\pgfqpoint{2.187693in}{0.649271in}}%
\pgfpathlineto{\pgfqpoint{2.184550in}{0.645723in}}%
\pgfpathlineto{\pgfqpoint{2.185584in}{0.641031in}}%
\pgfpathlineto{\pgfqpoint{2.191464in}{0.637662in}}%
\pgfpathlineto{\pgfqpoint{2.193761in}{0.621372in}}%
\pgfpathlineto{\pgfqpoint{2.192424in}{0.612482in}}%
\pgfpathlineto{\pgfqpoint{2.199079in}{0.604683in}}%
\pgfpathlineto{\pgfqpoint{2.206456in}{0.600906in}}%
\pgfpathlineto{\pgfqpoint{2.207470in}{0.596991in}}%
\pgfpathlineto{\pgfqpoint{2.205014in}{0.590892in}}%
\pgfpathlineto{\pgfqpoint{2.210055in}{0.585128in}}%
\pgfpathlineto{\pgfqpoint{2.206102in}{0.581924in}}%
\pgfpathlineto{\pgfqpoint{2.202262in}{0.574181in}}%
\pgfpathlineto{\pgfqpoint{2.194379in}{0.571471in}}%
\pgfpathlineto{\pgfqpoint{2.188678in}{0.573144in}}%
\pgfpathlineto{\pgfqpoint{2.179172in}{0.581268in}}%
\pgfpathlineto{\pgfqpoint{2.177184in}{0.585669in}}%
\pgfpathlineto{\pgfqpoint{2.178736in}{0.589471in}}%
\pgfpathlineto{\pgfqpoint{2.176721in}{0.596910in}}%
\pgfpathlineto{\pgfqpoint{2.173439in}{0.598944in}}%
\pgfpathlineto{\pgfqpoint{2.169431in}{0.595956in}}%
\pgfpathlineto{\pgfqpoint{2.167717in}{0.586842in}}%
\pgfpathlineto{\pgfqpoint{2.168005in}{0.577878in}}%
\pgfpathlineto{\pgfqpoint{2.166722in}{0.574427in}}%
\pgfpathlineto{\pgfqpoint{2.154103in}{0.578631in}}%
\pgfpathlineto{\pgfqpoint{2.151467in}{0.577078in}}%
\pgfpathclose%
\pgfusepath{fill}%
\end{pgfscope}%
\begin{pgfscope}%
\pgfpathrectangle{\pgfqpoint{0.100000in}{0.100000in}}{\pgfqpoint{3.007045in}{1.925000in}}%
\pgfusepath{clip}%
\pgfsetbuttcap%
\pgfsetmiterjoin%
\definecolor{currentfill}{rgb}{0.510588,0.732303,0.858839}%
\pgfsetfillcolor{currentfill}%
\pgfsetlinewidth{0.000000pt}%
\definecolor{currentstroke}{rgb}{0.000000,0.000000,0.000000}%
\pgfsetstrokecolor{currentstroke}%
\pgfsetstrokeopacity{0.000000}%
\pgfsetdash{}{0pt}%
\pgfpathmoveto{\pgfqpoint{1.217220in}{0.903588in}}%
\pgfpathlineto{\pgfqpoint{1.214018in}{0.875058in}}%
\pgfpathlineto{\pgfqpoint{1.212799in}{0.875172in}}%
\pgfpathlineto{\pgfqpoint{1.210832in}{0.858083in}}%
\pgfpathlineto{\pgfqpoint{1.210171in}{0.852394in}}%
\pgfpathlineto{\pgfqpoint{1.177035in}{0.856282in}}%
\pgfpathlineto{\pgfqpoint{1.150364in}{0.859555in}}%
\pgfpathlineto{\pgfqpoint{1.151795in}{0.871370in}}%
\pgfpathlineto{\pgfqpoint{1.133587in}{0.879798in}}%
\pgfpathlineto{\pgfqpoint{1.127835in}{0.883763in}}%
\pgfpathlineto{\pgfqpoint{1.082425in}{0.889855in}}%
\pgfpathlineto{\pgfqpoint{1.047397in}{0.894955in}}%
\pgfpathlineto{\pgfqpoint{1.011345in}{0.900484in}}%
\pgfpathlineto{\pgfqpoint{1.015230in}{0.925309in}}%
\pgfpathlineto{\pgfqpoint{1.046069in}{0.920555in}}%
\pgfpathlineto{\pgfqpoint{1.049496in}{0.943259in}}%
\pgfpathlineto{\pgfqpoint{1.111260in}{0.934336in}}%
\pgfpathlineto{\pgfqpoint{1.117536in}{0.979664in}}%
\pgfpathlineto{\pgfqpoint{1.100725in}{0.982093in}}%
\pgfpathlineto{\pgfqpoint{1.102849in}{0.996413in}}%
\pgfpathlineto{\pgfqpoint{1.141795in}{0.990914in}}%
\pgfpathlineto{\pgfqpoint{1.139908in}{0.976686in}}%
\pgfpathlineto{\pgfqpoint{1.173700in}{0.972532in}}%
\pgfpathlineto{\pgfqpoint{1.165289in}{0.970829in}}%
\pgfpathlineto{\pgfqpoint{1.164202in}{0.962160in}}%
\pgfpathlineto{\pgfqpoint{1.171643in}{0.956432in}}%
\pgfpathlineto{\pgfqpoint{1.166043in}{0.909544in}}%
\pgfpathlineto{\pgfqpoint{1.194470in}{0.906193in}}%
\pgfpathlineto{\pgfqpoint{1.217220in}{0.903588in}}%
\pgfpathclose%
\pgfusepath{fill}%
\end{pgfscope}%
\begin{pgfscope}%
\pgfpathrectangle{\pgfqpoint{0.100000in}{0.100000in}}{\pgfqpoint{3.007045in}{1.925000in}}%
\pgfusepath{clip}%
\pgfsetbuttcap%
\pgfsetmiterjoin%
\definecolor{currentfill}{rgb}{0.535686,0.746082,0.864252}%
\pgfsetfillcolor{currentfill}%
\pgfsetlinewidth{0.000000pt}%
\definecolor{currentstroke}{rgb}{0.000000,0.000000,0.000000}%
\pgfsetstrokecolor{currentstroke}%
\pgfsetstrokeopacity{0.000000}%
\pgfsetdash{}{0pt}%
\pgfpathmoveto{\pgfqpoint{1.677967in}{0.558409in}}%
\pgfpathlineto{\pgfqpoint{1.672046in}{0.560070in}}%
\pgfpathlineto{\pgfqpoint{1.665066in}{0.567789in}}%
\pgfpathlineto{\pgfqpoint{1.662849in}{0.573743in}}%
\pgfpathlineto{\pgfqpoint{1.659758in}{0.576115in}}%
\pgfpathlineto{\pgfqpoint{1.679583in}{0.587280in}}%
\pgfpathlineto{\pgfqpoint{1.678179in}{0.593593in}}%
\pgfpathlineto{\pgfqpoint{1.672091in}{0.600546in}}%
\pgfpathlineto{\pgfqpoint{1.670600in}{0.609460in}}%
\pgfpathlineto{\pgfqpoint{1.668209in}{0.612458in}}%
\pgfpathlineto{\pgfqpoint{1.681563in}{0.619854in}}%
\pgfpathlineto{\pgfqpoint{1.697477in}{0.628654in}}%
\pgfpathlineto{\pgfqpoint{1.696915in}{0.623980in}}%
\pgfpathlineto{\pgfqpoint{1.701562in}{0.603092in}}%
\pgfpathlineto{\pgfqpoint{1.705582in}{0.593025in}}%
\pgfpathlineto{\pgfqpoint{1.723044in}{0.595566in}}%
\pgfpathlineto{\pgfqpoint{1.724776in}{0.579988in}}%
\pgfpathlineto{\pgfqpoint{1.726035in}{0.554362in}}%
\pgfpathlineto{\pgfqpoint{1.709639in}{0.553402in}}%
\pgfpathlineto{\pgfqpoint{1.709428in}{0.556522in}}%
\pgfpathlineto{\pgfqpoint{1.704636in}{0.564588in}}%
\pgfpathlineto{\pgfqpoint{1.696897in}{0.563937in}}%
\pgfpathlineto{\pgfqpoint{1.690544in}{0.561282in}}%
\pgfpathlineto{\pgfqpoint{1.677967in}{0.558409in}}%
\pgfpathclose%
\pgfusepath{fill}%
\end{pgfscope}%
\begin{pgfscope}%
\pgfpathrectangle{\pgfqpoint{0.100000in}{0.100000in}}{\pgfqpoint{3.007045in}{1.925000in}}%
\pgfusepath{clip}%
\pgfsetbuttcap%
\pgfsetmiterjoin%
\definecolor{currentfill}{rgb}{0.516863,0.735748,0.860192}%
\pgfsetfillcolor{currentfill}%
\pgfsetlinewidth{0.000000pt}%
\definecolor{currentstroke}{rgb}{0.000000,0.000000,0.000000}%
\pgfsetstrokecolor{currentstroke}%
\pgfsetstrokeopacity{0.000000}%
\pgfsetdash{}{0pt}%
\pgfpathmoveto{\pgfqpoint{2.509239in}{1.087027in}}%
\pgfpathlineto{\pgfqpoint{2.495503in}{1.095488in}}%
\pgfpathlineto{\pgfqpoint{2.490450in}{1.093853in}}%
\pgfpathlineto{\pgfqpoint{2.489686in}{1.089390in}}%
\pgfpathlineto{\pgfqpoint{2.484474in}{1.093009in}}%
\pgfpathlineto{\pgfqpoint{2.476705in}{1.103396in}}%
\pgfpathlineto{\pgfqpoint{2.471855in}{1.105407in}}%
\pgfpathlineto{\pgfqpoint{2.467580in}{1.113861in}}%
\pgfpathlineto{\pgfqpoint{2.470294in}{1.118309in}}%
\pgfpathlineto{\pgfqpoint{2.476974in}{1.118028in}}%
\pgfpathlineto{\pgfqpoint{2.477515in}{1.130920in}}%
\pgfpathlineto{\pgfqpoint{2.481701in}{1.138352in}}%
\pgfpathlineto{\pgfqpoint{2.490413in}{1.135206in}}%
\pgfpathlineto{\pgfqpoint{2.494710in}{1.137649in}}%
\pgfpathlineto{\pgfqpoint{2.501225in}{1.130708in}}%
\pgfpathlineto{\pgfqpoint{2.510167in}{1.130498in}}%
\pgfpathlineto{\pgfqpoint{2.522292in}{1.144995in}}%
\pgfpathlineto{\pgfqpoint{2.526634in}{1.142930in}}%
\pgfpathlineto{\pgfqpoint{2.527646in}{1.135652in}}%
\pgfpathlineto{\pgfqpoint{2.533411in}{1.132073in}}%
\pgfpathlineto{\pgfqpoint{2.538421in}{1.132193in}}%
\pgfpathlineto{\pgfqpoint{2.548827in}{1.135761in}}%
\pgfpathlineto{\pgfqpoint{2.547549in}{1.130953in}}%
\pgfpathlineto{\pgfqpoint{2.538295in}{1.117062in}}%
\pgfpathlineto{\pgfqpoint{2.535565in}{1.108737in}}%
\pgfpathlineto{\pgfqpoint{2.527642in}{1.106923in}}%
\pgfpathlineto{\pgfqpoint{2.516736in}{1.108073in}}%
\pgfpathlineto{\pgfqpoint{2.509239in}{1.087027in}}%
\pgfpathclose%
\pgfusepath{fill}%
\end{pgfscope}%
\begin{pgfscope}%
\pgfpathrectangle{\pgfqpoint{0.100000in}{0.100000in}}{\pgfqpoint{3.007045in}{1.925000in}}%
\pgfusepath{clip}%
\pgfsetbuttcap%
\pgfsetmiterjoin%
\definecolor{currentfill}{rgb}{0.701423,0.826928,0.910988}%
\pgfsetfillcolor{currentfill}%
\pgfsetlinewidth{0.000000pt}%
\definecolor{currentstroke}{rgb}{0.000000,0.000000,0.000000}%
\pgfsetstrokecolor{currentstroke}%
\pgfsetstrokeopacity{0.000000}%
\pgfsetdash{}{0pt}%
\pgfpathmoveto{\pgfqpoint{2.768356in}{1.454967in}}%
\pgfpathlineto{\pgfqpoint{2.739778in}{1.426983in}}%
\pgfpathlineto{\pgfqpoint{2.733867in}{1.427016in}}%
\pgfpathlineto{\pgfqpoint{2.732166in}{1.431244in}}%
\pgfpathlineto{\pgfqpoint{2.727454in}{1.434295in}}%
\pgfpathlineto{\pgfqpoint{2.723883in}{1.436359in}}%
\pgfpathlineto{\pgfqpoint{2.721879in}{1.446280in}}%
\pgfpathlineto{\pgfqpoint{2.711428in}{1.444046in}}%
\pgfpathlineto{\pgfqpoint{2.710631in}{1.447455in}}%
\pgfpathlineto{\pgfqpoint{2.700669in}{1.446056in}}%
\pgfpathlineto{\pgfqpoint{2.697553in}{1.455914in}}%
\pgfpathlineto{\pgfqpoint{2.692078in}{1.475538in}}%
\pgfpathlineto{\pgfqpoint{2.719947in}{1.482951in}}%
\pgfpathlineto{\pgfqpoint{2.722365in}{1.491084in}}%
\pgfpathlineto{\pgfqpoint{2.726801in}{1.495572in}}%
\pgfpathlineto{\pgfqpoint{2.737175in}{1.492249in}}%
\pgfpathlineto{\pgfqpoint{2.737431in}{1.497269in}}%
\pgfpathlineto{\pgfqpoint{2.743391in}{1.496243in}}%
\pgfpathlineto{\pgfqpoint{2.749311in}{1.495307in}}%
\pgfpathlineto{\pgfqpoint{2.749558in}{1.490026in}}%
\pgfpathlineto{\pgfqpoint{2.753121in}{1.482442in}}%
\pgfpathlineto{\pgfqpoint{2.750879in}{1.474563in}}%
\pgfpathlineto{\pgfqpoint{2.756645in}{1.469554in}}%
\pgfpathlineto{\pgfqpoint{2.766813in}{1.466787in}}%
\pgfpathlineto{\pgfqpoint{2.763804in}{1.456101in}}%
\pgfpathlineto{\pgfqpoint{2.768356in}{1.454967in}}%
\pgfpathclose%
\pgfusepath{fill}%
\end{pgfscope}%
\begin{pgfscope}%
\pgfpathrectangle{\pgfqpoint{0.100000in}{0.100000in}}{\pgfqpoint{3.007045in}{1.925000in}}%
\pgfusepath{clip}%
\pgfsetbuttcap%
\pgfsetmiterjoin%
\definecolor{currentfill}{rgb}{0.296025,0.597955,0.790988}%
\pgfsetfillcolor{currentfill}%
\pgfsetlinewidth{0.000000pt}%
\definecolor{currentstroke}{rgb}{0.000000,0.000000,0.000000}%
\pgfsetstrokecolor{currentstroke}%
\pgfsetstrokeopacity{0.000000}%
\pgfsetdash{}{0pt}%
\pgfpathmoveto{\pgfqpoint{1.614445in}{1.686458in}}%
\pgfpathlineto{\pgfqpoint{1.579862in}{1.687980in}}%
\pgfpathlineto{\pgfqpoint{1.580956in}{1.710973in}}%
\pgfpathlineto{\pgfqpoint{1.580095in}{1.722634in}}%
\pgfpathlineto{\pgfqpoint{1.614607in}{1.721113in}}%
\pgfpathlineto{\pgfqpoint{1.615296in}{1.709477in}}%
\pgfpathlineto{\pgfqpoint{1.614445in}{1.686458in}}%
\pgfpathclose%
\pgfusepath{fill}%
\end{pgfscope}%
\begin{pgfscope}%
\pgfpathrectangle{\pgfqpoint{0.100000in}{0.100000in}}{\pgfqpoint{3.007045in}{1.925000in}}%
\pgfusepath{clip}%
\pgfsetbuttcap%
\pgfsetmiterjoin%
\definecolor{currentfill}{rgb}{0.256286,0.570012,0.775163}%
\pgfsetfillcolor{currentfill}%
\pgfsetlinewidth{0.000000pt}%
\definecolor{currentstroke}{rgb}{0.000000,0.000000,0.000000}%
\pgfsetstrokecolor{currentstroke}%
\pgfsetstrokeopacity{0.000000}%
\pgfsetdash{}{0pt}%
\pgfpathmoveto{\pgfqpoint{1.627761in}{0.861901in}}%
\pgfpathlineto{\pgfqpoint{1.627401in}{0.850427in}}%
\pgfpathlineto{\pgfqpoint{1.604562in}{0.851135in}}%
\pgfpathlineto{\pgfqpoint{1.604926in}{0.862611in}}%
\pgfpathlineto{\pgfqpoint{1.576210in}{0.863673in}}%
\pgfpathlineto{\pgfqpoint{1.576781in}{0.879670in}}%
\pgfpathlineto{\pgfqpoint{1.577863in}{0.909634in}}%
\pgfpathlineto{\pgfqpoint{1.594488in}{0.908994in}}%
\pgfpathlineto{\pgfqpoint{1.594438in}{0.897482in}}%
\pgfpathlineto{\pgfqpoint{1.609800in}{0.896970in}}%
\pgfpathlineto{\pgfqpoint{1.619336in}{0.893761in}}%
\pgfpathlineto{\pgfqpoint{1.628500in}{0.893638in}}%
\pgfpathlineto{\pgfqpoint{1.627761in}{0.861901in}}%
\pgfpathclose%
\pgfusepath{fill}%
\end{pgfscope}%
\begin{pgfscope}%
\pgfpathrectangle{\pgfqpoint{0.100000in}{0.100000in}}{\pgfqpoint{3.007045in}{1.925000in}}%
\pgfusepath{clip}%
\pgfsetbuttcap%
\pgfsetmiterjoin%
\definecolor{currentfill}{rgb}{0.341423,0.628958,0.808704}%
\pgfsetfillcolor{currentfill}%
\pgfsetlinewidth{0.000000pt}%
\definecolor{currentstroke}{rgb}{0.000000,0.000000,0.000000}%
\pgfsetstrokecolor{currentstroke}%
\pgfsetstrokeopacity{0.000000}%
\pgfsetdash{}{0pt}%
\pgfpathmoveto{\pgfqpoint{1.821649in}{1.429987in}}%
\pgfpathlineto{\pgfqpoint{1.822052in}{1.390899in}}%
\pgfpathlineto{\pgfqpoint{1.776480in}{1.390760in}}%
\pgfpathlineto{\pgfqpoint{1.753560in}{1.390798in}}%
\pgfpathlineto{\pgfqpoint{1.753660in}{1.413597in}}%
\pgfpathlineto{\pgfqpoint{1.753741in}{1.429789in}}%
\pgfpathlineto{\pgfqpoint{1.776405in}{1.429765in}}%
\pgfpathlineto{\pgfqpoint{1.776471in}{1.413556in}}%
\pgfpathlineto{\pgfqpoint{1.799116in}{1.413633in}}%
\pgfpathlineto{\pgfqpoint{1.799028in}{1.429855in}}%
\pgfpathlineto{\pgfqpoint{1.821649in}{1.429987in}}%
\pgfpathclose%
\pgfusepath{fill}%
\end{pgfscope}%
\begin{pgfscope}%
\pgfpathrectangle{\pgfqpoint{0.100000in}{0.100000in}}{\pgfqpoint{3.007045in}{1.925000in}}%
\pgfusepath{clip}%
\pgfsetbuttcap%
\pgfsetmiterjoin%
\definecolor{currentfill}{rgb}{0.381776,0.656517,0.824452}%
\pgfsetfillcolor{currentfill}%
\pgfsetlinewidth{0.000000pt}%
\definecolor{currentstroke}{rgb}{0.000000,0.000000,0.000000}%
\pgfsetstrokecolor{currentstroke}%
\pgfsetstrokeopacity{0.000000}%
\pgfsetdash{}{0pt}%
\pgfpathmoveto{\pgfqpoint{2.370833in}{1.107548in}}%
\pgfpathlineto{\pgfqpoint{2.371064in}{1.101491in}}%
\pgfpathlineto{\pgfqpoint{2.367300in}{1.093259in}}%
\pgfpathlineto{\pgfqpoint{2.362396in}{1.092621in}}%
\pgfpathlineto{\pgfqpoint{2.355335in}{1.096669in}}%
\pgfpathlineto{\pgfqpoint{2.344556in}{1.097617in}}%
\pgfpathlineto{\pgfqpoint{2.336897in}{1.108708in}}%
\pgfpathlineto{\pgfqpoint{2.347373in}{1.122044in}}%
\pgfpathlineto{\pgfqpoint{2.353384in}{1.121802in}}%
\pgfpathlineto{\pgfqpoint{2.359204in}{1.116262in}}%
\pgfpathlineto{\pgfqpoint{2.370833in}{1.107548in}}%
\pgfpathclose%
\pgfusepath{fill}%
\end{pgfscope}%
\begin{pgfscope}%
\pgfpathrectangle{\pgfqpoint{0.100000in}{0.100000in}}{\pgfqpoint{3.007045in}{1.925000in}}%
\pgfusepath{clip}%
\pgfsetbuttcap%
\pgfsetmiterjoin%
\definecolor{currentfill}{rgb}{0.356555,0.639293,0.814610}%
\pgfsetfillcolor{currentfill}%
\pgfsetlinewidth{0.000000pt}%
\definecolor{currentstroke}{rgb}{0.000000,0.000000,0.000000}%
\pgfsetstrokecolor{currentstroke}%
\pgfsetstrokeopacity{0.000000}%
\pgfsetdash{}{0pt}%
\pgfpathmoveto{\pgfqpoint{2.172150in}{1.293918in}}%
\pgfpathlineto{\pgfqpoint{2.173650in}{1.276857in}}%
\pgfpathlineto{\pgfqpoint{2.171318in}{1.271716in}}%
\pgfpathlineto{\pgfqpoint{2.165690in}{1.271203in}}%
\pgfpathlineto{\pgfqpoint{2.166328in}{1.264572in}}%
\pgfpathlineto{\pgfqpoint{2.145056in}{1.262796in}}%
\pgfpathlineto{\pgfqpoint{2.141986in}{1.299655in}}%
\pgfpathlineto{\pgfqpoint{2.139901in}{1.326614in}}%
\pgfpathlineto{\pgfqpoint{2.145664in}{1.322773in}}%
\pgfpathlineto{\pgfqpoint{2.155565in}{1.322406in}}%
\pgfpathlineto{\pgfqpoint{2.168897in}{1.329206in}}%
\pgfpathlineto{\pgfqpoint{2.172150in}{1.293918in}}%
\pgfpathclose%
\pgfusepath{fill}%
\end{pgfscope}%
\begin{pgfscope}%
\pgfpathrectangle{\pgfqpoint{0.100000in}{0.100000in}}{\pgfqpoint{3.007045in}{1.925000in}}%
\pgfusepath{clip}%
\pgfsetbuttcap%
\pgfsetmiterjoin%
\definecolor{currentfill}{rgb}{0.316201,0.611734,0.798862}%
\pgfsetfillcolor{currentfill}%
\pgfsetlinewidth{0.000000pt}%
\definecolor{currentstroke}{rgb}{0.000000,0.000000,0.000000}%
\pgfsetstrokecolor{currentstroke}%
\pgfsetstrokeopacity{0.000000}%
\pgfsetdash{}{0pt}%
\pgfpathmoveto{\pgfqpoint{1.553928in}{0.768166in}}%
\pgfpathlineto{\pgfqpoint{1.583365in}{0.766882in}}%
\pgfpathlineto{\pgfqpoint{1.583500in}{0.771553in}}%
\pgfpathlineto{\pgfqpoint{1.607913in}{0.770715in}}%
\pgfpathlineto{\pgfqpoint{1.611193in}{0.768407in}}%
\pgfpathlineto{\pgfqpoint{1.610081in}{0.739788in}}%
\pgfpathlineto{\pgfqpoint{1.582071in}{0.741206in}}%
\pgfpathlineto{\pgfqpoint{1.581896in}{0.737526in}}%
\pgfpathlineto{\pgfqpoint{1.573659in}{0.737980in}}%
\pgfpathlineto{\pgfqpoint{1.523941in}{0.740360in}}%
\pgfpathlineto{\pgfqpoint{1.525314in}{0.769593in}}%
\pgfpathlineto{\pgfqpoint{1.553928in}{0.768166in}}%
\pgfpathclose%
\pgfusepath{fill}%
\end{pgfscope}%
\begin{pgfscope}%
\pgfpathrectangle{\pgfqpoint{0.100000in}{0.100000in}}{\pgfqpoint{3.007045in}{1.925000in}}%
\pgfusepath{clip}%
\pgfsetbuttcap%
\pgfsetmiterjoin%
\definecolor{currentfill}{rgb}{0.361599,0.642737,0.816578}%
\pgfsetfillcolor{currentfill}%
\pgfsetlinewidth{0.000000pt}%
\definecolor{currentstroke}{rgb}{0.000000,0.000000,0.000000}%
\pgfsetstrokecolor{currentstroke}%
\pgfsetstrokeopacity{0.000000}%
\pgfsetdash{}{0pt}%
\pgfpathmoveto{\pgfqpoint{1.939480in}{1.158394in}}%
\pgfpathlineto{\pgfqpoint{1.909303in}{1.157973in}}%
\pgfpathlineto{\pgfqpoint{1.909292in}{1.151308in}}%
\pgfpathlineto{\pgfqpoint{1.903384in}{1.151278in}}%
\pgfpathlineto{\pgfqpoint{1.889132in}{1.155695in}}%
\pgfpathlineto{\pgfqpoint{1.889421in}{1.180509in}}%
\pgfpathlineto{\pgfqpoint{1.880876in}{1.180425in}}%
\pgfpathlineto{\pgfqpoint{1.880879in}{1.202651in}}%
\pgfpathlineto{\pgfqpoint{1.906198in}{1.203318in}}%
\pgfpathlineto{\pgfqpoint{1.906383in}{1.197788in}}%
\pgfpathlineto{\pgfqpoint{1.926181in}{1.198155in}}%
\pgfpathlineto{\pgfqpoint{1.931854in}{1.198270in}}%
\pgfpathlineto{\pgfqpoint{1.932193in}{1.179188in}}%
\pgfpathlineto{\pgfqpoint{1.938875in}{1.179366in}}%
\pgfpathlineto{\pgfqpoint{1.939480in}{1.158394in}}%
\pgfpathclose%
\pgfusepath{fill}%
\end{pgfscope}%
\begin{pgfscope}%
\pgfpathrectangle{\pgfqpoint{0.100000in}{0.100000in}}{\pgfqpoint{3.007045in}{1.925000in}}%
\pgfusepath{clip}%
\pgfsetbuttcap%
\pgfsetmiterjoin%
\definecolor{currentfill}{rgb}{0.306113,0.604844,0.794925}%
\pgfsetfillcolor{currentfill}%
\pgfsetlinewidth{0.000000pt}%
\definecolor{currentstroke}{rgb}{0.000000,0.000000,0.000000}%
\pgfsetstrokecolor{currentstroke}%
\pgfsetstrokeopacity{0.000000}%
\pgfsetdash{}{0pt}%
\pgfpathmoveto{\pgfqpoint{1.719231in}{0.879638in}}%
\pgfpathlineto{\pgfqpoint{1.719456in}{0.888768in}}%
\pgfpathlineto{\pgfqpoint{1.719521in}{0.894505in}}%
\pgfpathlineto{\pgfqpoint{1.725208in}{0.894438in}}%
\pgfpathlineto{\pgfqpoint{1.728070in}{0.900172in}}%
\pgfpathlineto{\pgfqpoint{1.728181in}{0.905962in}}%
\pgfpathlineto{\pgfqpoint{1.734052in}{0.905918in}}%
\pgfpathlineto{\pgfqpoint{1.734135in}{0.917404in}}%
\pgfpathlineto{\pgfqpoint{1.731290in}{0.917416in}}%
\pgfpathlineto{\pgfqpoint{1.731364in}{0.926038in}}%
\pgfpathlineto{\pgfqpoint{1.737553in}{0.926013in}}%
\pgfpathlineto{\pgfqpoint{1.740818in}{0.919900in}}%
\pgfpathlineto{\pgfqpoint{1.746777in}{0.922489in}}%
\pgfpathlineto{\pgfqpoint{1.758114in}{0.923040in}}%
\pgfpathlineto{\pgfqpoint{1.757448in}{0.932993in}}%
\pgfpathlineto{\pgfqpoint{1.761242in}{0.940278in}}%
\pgfpathlineto{\pgfqpoint{1.766027in}{0.940252in}}%
\pgfpathlineto{\pgfqpoint{1.766020in}{0.946003in}}%
\pgfpathlineto{\pgfqpoint{1.771725in}{0.946010in}}%
\pgfpathlineto{\pgfqpoint{1.795605in}{0.946063in}}%
\pgfpathlineto{\pgfqpoint{1.799463in}{0.919527in}}%
\pgfpathlineto{\pgfqpoint{1.800606in}{0.911576in}}%
\pgfpathlineto{\pgfqpoint{1.765492in}{0.911520in}}%
\pgfpathlineto{\pgfqpoint{1.765215in}{0.904084in}}%
\pgfpathlineto{\pgfqpoint{1.768002in}{0.897857in}}%
\pgfpathlineto{\pgfqpoint{1.763675in}{0.893735in}}%
\pgfpathlineto{\pgfqpoint{1.759980in}{0.886994in}}%
\pgfpathlineto{\pgfqpoint{1.751555in}{0.890916in}}%
\pgfpathlineto{\pgfqpoint{1.741580in}{0.885648in}}%
\pgfpathlineto{\pgfqpoint{1.732284in}{0.882217in}}%
\pgfpathlineto{\pgfqpoint{1.725397in}{0.881888in}}%
\pgfpathlineto{\pgfqpoint{1.719231in}{0.879638in}}%
\pgfpathclose%
\pgfusepath{fill}%
\end{pgfscope}%
\begin{pgfscope}%
\pgfpathrectangle{\pgfqpoint{0.100000in}{0.100000in}}{\pgfqpoint{3.007045in}{1.925000in}}%
\pgfusepath{clip}%
\pgfsetbuttcap%
\pgfsetmiterjoin%
\definecolor{currentfill}{rgb}{0.447843,0.697855,0.845306}%
\pgfsetfillcolor{currentfill}%
\pgfsetlinewidth{0.000000pt}%
\definecolor{currentstroke}{rgb}{0.000000,0.000000,0.000000}%
\pgfsetstrokecolor{currentstroke}%
\pgfsetstrokeopacity{0.000000}%
\pgfsetdash{}{0pt}%
\pgfpathmoveto{\pgfqpoint{1.756771in}{1.198935in}}%
\pgfpathlineto{\pgfqpoint{1.751069in}{1.202842in}}%
\pgfpathlineto{\pgfqpoint{1.751959in}{1.207454in}}%
\pgfpathlineto{\pgfqpoint{1.748211in}{1.211203in}}%
\pgfpathlineto{\pgfqpoint{1.747673in}{1.215377in}}%
\pgfpathlineto{\pgfqpoint{1.741287in}{1.219682in}}%
\pgfpathlineto{\pgfqpoint{1.734017in}{1.237625in}}%
\pgfpathlineto{\pgfqpoint{1.790743in}{1.236578in}}%
\pgfpathlineto{\pgfqpoint{1.792476in}{1.224430in}}%
\pgfpathlineto{\pgfqpoint{1.792177in}{1.207263in}}%
\pgfpathlineto{\pgfqpoint{1.770162in}{1.207494in}}%
\pgfpathlineto{\pgfqpoint{1.772661in}{1.193312in}}%
\pgfpathlineto{\pgfqpoint{1.768029in}{1.189813in}}%
\pgfpathlineto{\pgfqpoint{1.762123in}{1.192600in}}%
\pgfpathlineto{\pgfqpoint{1.756771in}{1.198935in}}%
\pgfpathclose%
\pgfusepath{fill}%
\end{pgfscope}%
\begin{pgfscope}%
\pgfpathrectangle{\pgfqpoint{0.100000in}{0.100000in}}{\pgfqpoint{3.007045in}{1.925000in}}%
\pgfusepath{clip}%
\pgfsetbuttcap%
\pgfsetmiterjoin%
\definecolor{currentfill}{rgb}{0.573333,0.766751,0.872372}%
\pgfsetfillcolor{currentfill}%
\pgfsetlinewidth{0.000000pt}%
\definecolor{currentstroke}{rgb}{0.000000,0.000000,0.000000}%
\pgfsetstrokecolor{currentstroke}%
\pgfsetstrokeopacity{0.000000}%
\pgfsetdash{}{0pt}%
\pgfpathmoveto{\pgfqpoint{2.285385in}{0.933858in}}%
\pgfpathlineto{\pgfqpoint{2.278483in}{0.930848in}}%
\pgfpathlineto{\pgfqpoint{2.264733in}{0.928140in}}%
\pgfpathlineto{\pgfqpoint{2.261477in}{0.928516in}}%
\pgfpathlineto{\pgfqpoint{2.257324in}{0.931388in}}%
\pgfpathlineto{\pgfqpoint{2.254813in}{0.936868in}}%
\pgfpathlineto{\pgfqpoint{2.254335in}{0.941757in}}%
\pgfpathlineto{\pgfqpoint{2.257521in}{0.951222in}}%
\pgfpathlineto{\pgfqpoint{2.251263in}{0.956468in}}%
\pgfpathlineto{\pgfqpoint{2.248417in}{0.962697in}}%
\pgfpathlineto{\pgfqpoint{2.253187in}{0.966596in}}%
\pgfpathlineto{\pgfqpoint{2.258371in}{0.965969in}}%
\pgfpathlineto{\pgfqpoint{2.261147in}{0.968764in}}%
\pgfpathlineto{\pgfqpoint{2.263169in}{0.965439in}}%
\pgfpathlineto{\pgfqpoint{2.267876in}{0.965337in}}%
\pgfpathlineto{\pgfqpoint{2.269975in}{0.962101in}}%
\pgfpathlineto{\pgfqpoint{2.277359in}{0.967152in}}%
\pgfpathlineto{\pgfqpoint{2.285629in}{0.965628in}}%
\pgfpathlineto{\pgfqpoint{2.291000in}{0.961775in}}%
\pgfpathlineto{\pgfqpoint{2.293939in}{0.955321in}}%
\pgfpathlineto{\pgfqpoint{2.292762in}{0.944860in}}%
\pgfpathlineto{\pgfqpoint{2.285385in}{0.933858in}}%
\pgfpathclose%
\pgfusepath{fill}%
\end{pgfscope}%
\begin{pgfscope}%
\pgfpathrectangle{\pgfqpoint{0.100000in}{0.100000in}}{\pgfqpoint{3.007045in}{1.925000in}}%
\pgfusepath{clip}%
\pgfsetbuttcap%
\pgfsetmiterjoin%
\definecolor{currentfill}{rgb}{0.265759,0.577286,0.779177}%
\pgfsetfillcolor{currentfill}%
\pgfsetlinewidth{0.000000pt}%
\definecolor{currentstroke}{rgb}{0.000000,0.000000,0.000000}%
\pgfsetstrokecolor{currentstroke}%
\pgfsetstrokeopacity{0.000000}%
\pgfsetdash{}{0pt}%
\pgfpathmoveto{\pgfqpoint{1.193280in}{1.837329in}}%
\pgfpathlineto{\pgfqpoint{1.244920in}{1.830121in}}%
\pgfpathlineto{\pgfqpoint{1.285700in}{1.824835in}}%
\pgfpathlineto{\pgfqpoint{1.282530in}{1.813419in}}%
\pgfpathlineto{\pgfqpoint{1.288237in}{1.812721in}}%
\pgfpathlineto{\pgfqpoint{1.286131in}{1.795421in}}%
\pgfpathlineto{\pgfqpoint{1.295653in}{1.794279in}}%
\pgfpathlineto{\pgfqpoint{1.292930in}{1.771456in}}%
\pgfpathlineto{\pgfqpoint{1.291202in}{1.771653in}}%
\pgfpathlineto{\pgfqpoint{1.289529in}{1.757826in}}%
\pgfpathlineto{\pgfqpoint{1.272584in}{1.760976in}}%
\pgfpathlineto{\pgfqpoint{1.265605in}{1.764425in}}%
\pgfpathlineto{\pgfqpoint{1.263629in}{1.757459in}}%
\pgfpathlineto{\pgfqpoint{1.260272in}{1.755961in}}%
\pgfpathlineto{\pgfqpoint{1.255859in}{1.748834in}}%
\pgfpathlineto{\pgfqpoint{1.246943in}{1.744290in}}%
\pgfpathlineto{\pgfqpoint{1.241819in}{1.745330in}}%
\pgfpathlineto{\pgfqpoint{1.240243in}{1.741597in}}%
\pgfpathlineto{\pgfqpoint{1.224257in}{1.742711in}}%
\pgfpathlineto{\pgfqpoint{1.217097in}{1.745725in}}%
\pgfpathlineto{\pgfqpoint{1.214896in}{1.741346in}}%
\pgfpathlineto{\pgfqpoint{1.206785in}{1.744306in}}%
\pgfpathlineto{\pgfqpoint{1.201819in}{1.739713in}}%
\pgfpathlineto{\pgfqpoint{1.195585in}{1.736909in}}%
\pgfpathlineto{\pgfqpoint{1.192875in}{1.732905in}}%
\pgfpathlineto{\pgfqpoint{1.193102in}{1.740744in}}%
\pgfpathlineto{\pgfqpoint{1.187367in}{1.743646in}}%
\pgfpathlineto{\pgfqpoint{1.174368in}{1.744062in}}%
\pgfpathlineto{\pgfqpoint{1.170837in}{1.747018in}}%
\pgfpathlineto{\pgfqpoint{1.164883in}{1.747708in}}%
\pgfpathlineto{\pgfqpoint{1.156612in}{1.751283in}}%
\pgfpathlineto{\pgfqpoint{1.152002in}{1.758091in}}%
\pgfpathlineto{\pgfqpoint{1.154039in}{1.770483in}}%
\pgfpathlineto{\pgfqpoint{1.165616in}{1.768587in}}%
\pgfpathlineto{\pgfqpoint{1.167096in}{1.773103in}}%
\pgfpathlineto{\pgfqpoint{1.174638in}{1.770993in}}%
\pgfpathlineto{\pgfqpoint{1.177030in}{1.787023in}}%
\pgfpathlineto{\pgfqpoint{1.180237in}{1.801764in}}%
\pgfpathlineto{\pgfqpoint{1.182880in}{1.799848in}}%
\pgfpathlineto{\pgfqpoint{1.185494in}{1.809331in}}%
\pgfpathlineto{\pgfqpoint{1.187177in}{1.820830in}}%
\pgfpathlineto{\pgfqpoint{1.189784in}{1.820439in}}%
\pgfpathlineto{\pgfqpoint{1.193280in}{1.837329in}}%
\pgfpathclose%
\pgfusepath{fill}%
\end{pgfscope}%
\begin{pgfscope}%
\pgfpathrectangle{\pgfqpoint{0.100000in}{0.100000in}}{\pgfqpoint{3.007045in}{1.925000in}}%
\pgfusepath{clip}%
\pgfsetbuttcap%
\pgfsetmiterjoin%
\definecolor{currentfill}{rgb}{0.760477,0.852026,0.931657}%
\pgfsetfillcolor{currentfill}%
\pgfsetlinewidth{0.000000pt}%
\definecolor{currentstroke}{rgb}{0.000000,0.000000,0.000000}%
\pgfsetstrokecolor{currentstroke}%
\pgfsetstrokeopacity{0.000000}%
\pgfsetdash{}{0pt}%
\pgfpathmoveto{\pgfqpoint{2.642481in}{0.444772in}}%
\pgfpathlineto{\pgfqpoint{2.615609in}{0.440573in}}%
\pgfpathlineto{\pgfqpoint{2.614673in}{0.446167in}}%
\pgfpathlineto{\pgfqpoint{2.609094in}{0.445229in}}%
\pgfpathlineto{\pgfqpoint{2.602369in}{0.492044in}}%
\pgfpathlineto{\pgfqpoint{2.599254in}{0.497364in}}%
\pgfpathlineto{\pgfqpoint{2.599396in}{0.502567in}}%
\pgfpathlineto{\pgfqpoint{2.592396in}{0.508481in}}%
\pgfpathlineto{\pgfqpoint{2.591696in}{0.520372in}}%
\pgfpathlineto{\pgfqpoint{2.605371in}{0.522544in}}%
\pgfpathlineto{\pgfqpoint{2.616569in}{0.510498in}}%
\pgfpathlineto{\pgfqpoint{2.620884in}{0.502421in}}%
\pgfpathlineto{\pgfqpoint{2.617307in}{0.495527in}}%
\pgfpathlineto{\pgfqpoint{2.618107in}{0.489509in}}%
\pgfpathlineto{\pgfqpoint{2.622871in}{0.477008in}}%
\pgfpathlineto{\pgfqpoint{2.636790in}{0.456290in}}%
\pgfpathlineto{\pgfqpoint{2.642481in}{0.444772in}}%
\pgfpathclose%
\pgfusepath{fill}%
\end{pgfscope}%
\begin{pgfscope}%
\pgfpathrectangle{\pgfqpoint{0.100000in}{0.100000in}}{\pgfqpoint{3.007045in}{1.925000in}}%
\pgfusepath{clip}%
\pgfsetbuttcap%
\pgfsetmiterjoin%
\definecolor{currentfill}{rgb}{0.321246,0.615179,0.800830}%
\pgfsetfillcolor{currentfill}%
\pgfsetlinewidth{0.000000pt}%
\definecolor{currentstroke}{rgb}{0.000000,0.000000,0.000000}%
\pgfsetstrokecolor{currentstroke}%
\pgfsetstrokeopacity{0.000000}%
\pgfsetdash{}{0pt}%
\pgfpathmoveto{\pgfqpoint{1.865214in}{0.855241in}}%
\pgfpathlineto{\pgfqpoint{1.859472in}{0.853168in}}%
\pgfpathlineto{\pgfqpoint{1.842321in}{0.853020in}}%
\pgfpathlineto{\pgfqpoint{1.842357in}{0.871145in}}%
\pgfpathlineto{\pgfqpoint{1.822375in}{0.871354in}}%
\pgfpathlineto{\pgfqpoint{1.818712in}{0.873240in}}%
\pgfpathlineto{\pgfqpoint{1.818790in}{0.879030in}}%
\pgfpathlineto{\pgfqpoint{1.821794in}{0.883783in}}%
\pgfpathlineto{\pgfqpoint{1.822064in}{0.899000in}}%
\pgfpathlineto{\pgfqpoint{1.824480in}{0.903095in}}%
\pgfpathlineto{\pgfqpoint{1.821906in}{0.907584in}}%
\pgfpathlineto{\pgfqpoint{1.821996in}{0.913307in}}%
\pgfpathlineto{\pgfqpoint{1.826781in}{0.913265in}}%
\pgfpathlineto{\pgfqpoint{1.830635in}{0.919840in}}%
\pgfpathlineto{\pgfqpoint{1.851596in}{0.920385in}}%
\pgfpathlineto{\pgfqpoint{1.851607in}{0.918465in}}%
\pgfpathlineto{\pgfqpoint{1.889632in}{0.918440in}}%
\pgfpathlineto{\pgfqpoint{1.889712in}{0.906997in}}%
\pgfpathlineto{\pgfqpoint{1.887849in}{0.901272in}}%
\pgfpathlineto{\pgfqpoint{1.888066in}{0.882118in}}%
\pgfpathlineto{\pgfqpoint{1.884196in}{0.881637in}}%
\pgfpathlineto{\pgfqpoint{1.871864in}{0.868767in}}%
\pgfpathlineto{\pgfqpoint{1.864117in}{0.861985in}}%
\pgfpathlineto{\pgfqpoint{1.865214in}{0.855241in}}%
\pgfpathclose%
\pgfusepath{fill}%
\end{pgfscope}%
\begin{pgfscope}%
\pgfpathrectangle{\pgfqpoint{0.100000in}{0.100000in}}{\pgfqpoint{3.007045in}{1.925000in}}%
\pgfusepath{clip}%
\pgfsetbuttcap%
\pgfsetmiterjoin%
\definecolor{currentfill}{rgb}{0.371688,0.649627,0.820515}%
\pgfsetfillcolor{currentfill}%
\pgfsetlinewidth{0.000000pt}%
\definecolor{currentstroke}{rgb}{0.000000,0.000000,0.000000}%
\pgfsetstrokecolor{currentstroke}%
\pgfsetstrokeopacity{0.000000}%
\pgfsetdash{}{0pt}%
\pgfpathmoveto{\pgfqpoint{0.912796in}{0.778126in}}%
\pgfpathlineto{\pgfqpoint{0.851184in}{0.788621in}}%
\pgfpathlineto{\pgfqpoint{0.795773in}{0.799049in}}%
\pgfpathlineto{\pgfqpoint{0.754327in}{0.807329in}}%
\pgfpathlineto{\pgfqpoint{0.765905in}{0.863949in}}%
\pgfpathlineto{\pgfqpoint{0.774157in}{0.904247in}}%
\pgfpathlineto{\pgfqpoint{0.805861in}{0.897916in}}%
\pgfpathlineto{\pgfqpoint{0.829566in}{0.885369in}}%
\pgfpathlineto{\pgfqpoint{0.837638in}{0.894986in}}%
\pgfpathlineto{\pgfqpoint{0.860696in}{0.887387in}}%
\pgfpathlineto{\pgfqpoint{0.873121in}{0.885088in}}%
\pgfpathlineto{\pgfqpoint{0.875838in}{0.894759in}}%
\pgfpathlineto{\pgfqpoint{0.863233in}{0.897098in}}%
\pgfpathlineto{\pgfqpoint{0.868902in}{0.911853in}}%
\pgfpathlineto{\pgfqpoint{0.876178in}{0.916207in}}%
\pgfpathlineto{\pgfqpoint{0.883290in}{0.911640in}}%
\pgfpathlineto{\pgfqpoint{0.890889in}{0.911616in}}%
\pgfpathlineto{\pgfqpoint{0.896277in}{0.906751in}}%
\pgfpathlineto{\pgfqpoint{0.903247in}{0.903940in}}%
\pgfpathlineto{\pgfqpoint{0.911088in}{0.896114in}}%
\pgfpathlineto{\pgfqpoint{0.916273in}{0.895048in}}%
\pgfpathlineto{\pgfqpoint{0.913272in}{0.877904in}}%
\pgfpathlineto{\pgfqpoint{0.953797in}{0.870982in}}%
\pgfpathlineto{\pgfqpoint{0.949199in}{0.843379in}}%
\pgfpathlineto{\pgfqpoint{0.947992in}{0.836130in}}%
\pgfpathlineto{\pgfqpoint{0.938796in}{0.837657in}}%
\pgfpathlineto{\pgfqpoint{0.930194in}{0.836921in}}%
\pgfpathlineto{\pgfqpoint{0.923055in}{0.834676in}}%
\pgfpathlineto{\pgfqpoint{0.920548in}{0.822552in}}%
\pgfpathlineto{\pgfqpoint{0.912796in}{0.778126in}}%
\pgfpathclose%
\pgfusepath{fill}%
\end{pgfscope}%
\begin{pgfscope}%
\pgfpathrectangle{\pgfqpoint{0.100000in}{0.100000in}}{\pgfqpoint{3.007045in}{1.925000in}}%
\pgfusepath{clip}%
\pgfsetbuttcap%
\pgfsetmiterjoin%
\definecolor{currentfill}{rgb}{0.290980,0.594510,0.789020}%
\pgfsetfillcolor{currentfill}%
\pgfsetlinewidth{0.000000pt}%
\definecolor{currentstroke}{rgb}{0.000000,0.000000,0.000000}%
\pgfsetstrokecolor{currentstroke}%
\pgfsetstrokeopacity{0.000000}%
\pgfsetdash{}{0pt}%
\pgfpathmoveto{\pgfqpoint{1.710630in}{1.795591in}}%
\pgfpathlineto{\pgfqpoint{1.765577in}{1.795047in}}%
\pgfpathlineto{\pgfqpoint{1.765638in}{1.820857in}}%
\pgfpathlineto{\pgfqpoint{1.769776in}{1.818763in}}%
\pgfpathlineto{\pgfqpoint{1.774181in}{1.819903in}}%
\pgfpathlineto{\pgfqpoint{1.779963in}{1.814824in}}%
\pgfpathlineto{\pgfqpoint{1.786263in}{1.787384in}}%
\pgfpathlineto{\pgfqpoint{1.785945in}{1.780293in}}%
\pgfpathlineto{\pgfqpoint{1.790535in}{1.776268in}}%
\pgfpathlineto{\pgfqpoint{1.797461in}{1.775190in}}%
\pgfpathlineto{\pgfqpoint{1.797671in}{1.752902in}}%
\pgfpathlineto{\pgfqpoint{1.762966in}{1.752742in}}%
\pgfpathlineto{\pgfqpoint{1.762948in}{1.764365in}}%
\pgfpathlineto{\pgfqpoint{1.745680in}{1.764366in}}%
\pgfpathlineto{\pgfqpoint{1.710937in}{1.765107in}}%
\pgfpathlineto{\pgfqpoint{1.710357in}{1.776679in}}%
\pgfpathlineto{\pgfqpoint{1.710630in}{1.795591in}}%
\pgfpathclose%
\pgfusepath{fill}%
\end{pgfscope}%
\begin{pgfscope}%
\pgfpathrectangle{\pgfqpoint{0.100000in}{0.100000in}}{\pgfqpoint{3.007045in}{1.925000in}}%
\pgfusepath{clip}%
\pgfsetbuttcap%
\pgfsetmiterjoin%
\definecolor{currentfill}{rgb}{0.491765,0.721968,0.854779}%
\pgfsetfillcolor{currentfill}%
\pgfsetlinewidth{0.000000pt}%
\definecolor{currentstroke}{rgb}{0.000000,0.000000,0.000000}%
\pgfsetstrokecolor{currentstroke}%
\pgfsetstrokeopacity{0.000000}%
\pgfsetdash{}{0pt}%
\pgfpathmoveto{\pgfqpoint{0.851627in}{1.798562in}}%
\pgfpathlineto{\pgfqpoint{0.847085in}{1.804318in}}%
\pgfpathlineto{\pgfqpoint{0.850775in}{1.806624in}}%
\pgfpathlineto{\pgfqpoint{0.850754in}{1.813504in}}%
\pgfpathlineto{\pgfqpoint{0.847161in}{1.818241in}}%
\pgfpathlineto{\pgfqpoint{0.840821in}{1.835141in}}%
\pgfpathlineto{\pgfqpoint{0.844425in}{1.850601in}}%
\pgfpathlineto{\pgfqpoint{0.845385in}{1.846522in}}%
\pgfpathlineto{\pgfqpoint{0.859696in}{1.849285in}}%
\pgfpathlineto{\pgfqpoint{0.860078in}{1.837613in}}%
\pgfpathlineto{\pgfqpoint{0.863113in}{1.830232in}}%
\pgfpathlineto{\pgfqpoint{0.862566in}{1.825406in}}%
\pgfpathlineto{\pgfqpoint{0.873246in}{1.821840in}}%
\pgfpathlineto{\pgfqpoint{0.877958in}{1.823506in}}%
\pgfpathlineto{\pgfqpoint{0.880330in}{1.828801in}}%
\pgfpathlineto{\pgfqpoint{0.886583in}{1.827387in}}%
\pgfpathlineto{\pgfqpoint{0.889217in}{1.841046in}}%
\pgfpathlineto{\pgfqpoint{0.896692in}{1.839394in}}%
\pgfpathlineto{\pgfqpoint{0.901644in}{1.862062in}}%
\pgfpathlineto{\pgfqpoint{0.901194in}{1.867977in}}%
\pgfpathlineto{\pgfqpoint{0.910046in}{1.867809in}}%
\pgfpathlineto{\pgfqpoint{0.913309in}{1.873552in}}%
\pgfpathlineto{\pgfqpoint{0.912328in}{1.878097in}}%
\pgfpathlineto{\pgfqpoint{0.913025in}{1.888772in}}%
\pgfpathlineto{\pgfqpoint{0.941334in}{1.882540in}}%
\pgfpathlineto{\pgfqpoint{0.941225in}{1.874500in}}%
\pgfpathlineto{\pgfqpoint{0.943633in}{1.869658in}}%
\pgfpathlineto{\pgfqpoint{0.949121in}{1.871112in}}%
\pgfpathlineto{\pgfqpoint{0.953218in}{1.860951in}}%
\pgfpathlineto{\pgfqpoint{0.949618in}{1.854808in}}%
\pgfpathlineto{\pgfqpoint{0.961059in}{1.847574in}}%
\pgfpathlineto{\pgfqpoint{0.959285in}{1.841945in}}%
\pgfpathlineto{\pgfqpoint{0.964152in}{1.834807in}}%
\pgfpathlineto{\pgfqpoint{0.961768in}{1.833418in}}%
\pgfpathlineto{\pgfqpoint{0.967054in}{1.825877in}}%
\pgfpathlineto{\pgfqpoint{0.966264in}{1.820814in}}%
\pgfpathlineto{\pgfqpoint{0.975124in}{1.816423in}}%
\pgfpathlineto{\pgfqpoint{0.978555in}{1.808720in}}%
\pgfpathlineto{\pgfqpoint{0.974188in}{1.804466in}}%
\pgfpathlineto{\pgfqpoint{0.969863in}{1.800673in}}%
\pgfpathlineto{\pgfqpoint{0.967284in}{1.791973in}}%
\pgfpathlineto{\pgfqpoint{0.963166in}{1.790823in}}%
\pgfpathlineto{\pgfqpoint{0.962444in}{1.782513in}}%
\pgfpathlineto{\pgfqpoint{0.948336in}{1.785640in}}%
\pgfpathlineto{\pgfqpoint{0.927953in}{1.789677in}}%
\pgfpathlineto{\pgfqpoint{0.925310in}{1.782011in}}%
\pgfpathlineto{\pgfqpoint{0.927695in}{1.775797in}}%
\pgfpathlineto{\pgfqpoint{0.926220in}{1.769071in}}%
\pgfpathlineto{\pgfqpoint{0.926970in}{1.761628in}}%
\pgfpathlineto{\pgfqpoint{0.921428in}{1.759848in}}%
\pgfpathlineto{\pgfqpoint{0.910196in}{1.762234in}}%
\pgfpathlineto{\pgfqpoint{0.906982in}{1.761729in}}%
\pgfpathlineto{\pgfqpoint{0.900754in}{1.769170in}}%
\pgfpathlineto{\pgfqpoint{0.890065in}{1.775890in}}%
\pgfpathlineto{\pgfqpoint{0.883118in}{1.777032in}}%
\pgfpathlineto{\pgfqpoint{0.878165in}{1.781128in}}%
\pgfpathlineto{\pgfqpoint{0.878913in}{1.786617in}}%
\pgfpathlineto{\pgfqpoint{0.866931in}{1.795554in}}%
\pgfpathlineto{\pgfqpoint{0.859543in}{1.795823in}}%
\pgfpathlineto{\pgfqpoint{0.851627in}{1.798562in}}%
\pgfpathclose%
\pgfusepath{fill}%
\end{pgfscope}%
\begin{pgfscope}%
\pgfpathrectangle{\pgfqpoint{0.100000in}{0.100000in}}{\pgfqpoint{3.007045in}{1.925000in}}%
\pgfusepath{clip}%
\pgfsetbuttcap%
\pgfsetmiterjoin%
\definecolor{currentfill}{rgb}{0.031373,0.285675,0.564291}%
\pgfsetfillcolor{currentfill}%
\pgfsetlinewidth{0.000000pt}%
\definecolor{currentstroke}{rgb}{0.000000,0.000000,0.000000}%
\pgfsetstrokecolor{currentstroke}%
\pgfsetstrokeopacity{0.000000}%
\pgfsetdash{}{0pt}%
\pgfpathmoveto{\pgfqpoint{1.336243in}{0.559459in}}%
\pgfpathlineto{\pgfqpoint{1.349461in}{0.546748in}}%
\pgfpathlineto{\pgfqpoint{1.344620in}{0.539783in}}%
\pgfpathlineto{\pgfqpoint{1.328189in}{0.539358in}}%
\pgfpathlineto{\pgfqpoint{1.323294in}{0.529791in}}%
\pgfpathlineto{\pgfqpoint{1.319302in}{0.525390in}}%
\pgfpathlineto{\pgfqpoint{1.316511in}{0.514705in}}%
\pgfpathlineto{\pgfqpoint{1.305003in}{0.500040in}}%
\pgfpathlineto{\pgfqpoint{1.299683in}{0.496151in}}%
\pgfpathlineto{\pgfqpoint{1.298360in}{0.491273in}}%
\pgfpathlineto{\pgfqpoint{1.291128in}{0.491622in}}%
\pgfpathlineto{\pgfqpoint{1.278806in}{0.498569in}}%
\pgfpathlineto{\pgfqpoint{1.273825in}{0.505053in}}%
\pgfpathlineto{\pgfqpoint{1.264636in}{0.507489in}}%
\pgfpathlineto{\pgfqpoint{1.261257in}{0.513470in}}%
\pgfpathlineto{\pgfqpoint{1.250316in}{0.516616in}}%
\pgfpathlineto{\pgfqpoint{1.241232in}{0.523169in}}%
\pgfpathlineto{\pgfqpoint{1.236587in}{0.531465in}}%
\pgfpathlineto{\pgfqpoint{1.231030in}{0.533550in}}%
\pgfpathlineto{\pgfqpoint{1.221419in}{0.542714in}}%
\pgfpathlineto{\pgfqpoint{1.219216in}{0.552154in}}%
\pgfpathlineto{\pgfqpoint{1.214219in}{0.561569in}}%
\pgfpathlineto{\pgfqpoint{1.213608in}{0.571913in}}%
\pgfpathlineto{\pgfqpoint{1.214655in}{0.583241in}}%
\pgfpathlineto{\pgfqpoint{1.207041in}{0.594453in}}%
\pgfpathlineto{\pgfqpoint{1.207348in}{0.601413in}}%
\pgfpathlineto{\pgfqpoint{1.204717in}{0.609116in}}%
\pgfpathlineto{\pgfqpoint{1.201869in}{0.611048in}}%
\pgfpathlineto{\pgfqpoint{1.205698in}{0.612940in}}%
\pgfpathlineto{\pgfqpoint{1.255027in}{0.637292in}}%
\pgfpathlineto{\pgfqpoint{1.282209in}{0.612003in}}%
\pgfpathlineto{\pgfqpoint{1.336243in}{0.559459in}}%
\pgfpathclose%
\pgfusepath{fill}%
\end{pgfscope}%
\begin{pgfscope}%
\pgfpathrectangle{\pgfqpoint{0.100000in}{0.100000in}}{\pgfqpoint{3.007045in}{1.925000in}}%
\pgfusepath{clip}%
\pgfsetbuttcap%
\pgfsetmiterjoin%
\definecolor{currentfill}{rgb}{0.336378,0.625513,0.806736}%
\pgfsetfillcolor{currentfill}%
\pgfsetlinewidth{0.000000pt}%
\definecolor{currentstroke}{rgb}{0.000000,0.000000,0.000000}%
\pgfsetstrokecolor{currentstroke}%
\pgfsetstrokeopacity{0.000000}%
\pgfsetdash{}{0pt}%
\pgfpathmoveto{\pgfqpoint{2.279304in}{0.827039in}}%
\pgfpathlineto{\pgfqpoint{2.271634in}{0.825054in}}%
\pgfpathlineto{\pgfqpoint{2.262319in}{0.816941in}}%
\pgfpathlineto{\pgfqpoint{2.257720in}{0.820289in}}%
\pgfpathlineto{\pgfqpoint{2.254073in}{0.825883in}}%
\pgfpathlineto{\pgfqpoint{2.249130in}{0.825060in}}%
\pgfpathlineto{\pgfqpoint{2.246429in}{0.828970in}}%
\pgfpathlineto{\pgfqpoint{2.247672in}{0.832584in}}%
\pgfpathlineto{\pgfqpoint{2.236928in}{0.845086in}}%
\pgfpathlineto{\pgfqpoint{2.231240in}{0.844683in}}%
\pgfpathlineto{\pgfqpoint{2.232658in}{0.857963in}}%
\pgfpathlineto{\pgfqpoint{2.238722in}{0.856984in}}%
\pgfpathlineto{\pgfqpoint{2.244466in}{0.861498in}}%
\pgfpathlineto{\pgfqpoint{2.241299in}{0.867393in}}%
\pgfpathlineto{\pgfqpoint{2.241813in}{0.880528in}}%
\pgfpathlineto{\pgfqpoint{2.240204in}{0.883283in}}%
\pgfpathlineto{\pgfqpoint{2.241623in}{0.891209in}}%
\pgfpathlineto{\pgfqpoint{2.265651in}{0.893385in}}%
\pgfpathlineto{\pgfqpoint{2.279674in}{0.894582in}}%
\pgfpathlineto{\pgfqpoint{2.286587in}{0.869005in}}%
\pgfpathlineto{\pgfqpoint{2.287764in}{0.864843in}}%
\pgfpathlineto{\pgfqpoint{2.284621in}{0.861875in}}%
\pgfpathlineto{\pgfqpoint{2.282179in}{0.853795in}}%
\pgfpathlineto{\pgfqpoint{2.275437in}{0.845959in}}%
\pgfpathlineto{\pgfqpoint{2.271751in}{0.844592in}}%
\pgfpathlineto{\pgfqpoint{2.279304in}{0.827039in}}%
\pgfpathclose%
\pgfusepath{fill}%
\end{pgfscope}%
\begin{pgfscope}%
\pgfpathrectangle{\pgfqpoint{0.100000in}{0.100000in}}{\pgfqpoint{3.007045in}{1.925000in}}%
\pgfusepath{clip}%
\pgfsetbuttcap%
\pgfsetmiterjoin%
\definecolor{currentfill}{rgb}{0.541961,0.749527,0.865606}%
\pgfsetfillcolor{currentfill}%
\pgfsetlinewidth{0.000000pt}%
\definecolor{currentstroke}{rgb}{0.000000,0.000000,0.000000}%
\pgfsetstrokecolor{currentstroke}%
\pgfsetstrokeopacity{0.000000}%
\pgfsetdash{}{0pt}%
\pgfpathmoveto{\pgfqpoint{0.521615in}{1.353758in}}%
\pgfpathlineto{\pgfqpoint{0.517066in}{1.358013in}}%
\pgfpathlineto{\pgfqpoint{0.520278in}{1.368998in}}%
\pgfpathlineto{\pgfqpoint{0.517118in}{1.374961in}}%
\pgfpathlineto{\pgfqpoint{0.518020in}{1.379672in}}%
\pgfpathlineto{\pgfqpoint{0.512166in}{1.383346in}}%
\pgfpathlineto{\pgfqpoint{0.503882in}{1.397397in}}%
\pgfpathlineto{\pgfqpoint{0.495392in}{1.401958in}}%
\pgfpathlineto{\pgfqpoint{0.488838in}{1.398768in}}%
\pgfpathlineto{\pgfqpoint{0.482087in}{1.399330in}}%
\pgfpathlineto{\pgfqpoint{0.479961in}{1.402391in}}%
\pgfpathlineto{\pgfqpoint{0.483507in}{1.414404in}}%
\pgfpathlineto{\pgfqpoint{0.470680in}{1.418127in}}%
\pgfpathlineto{\pgfqpoint{0.479723in}{1.447096in}}%
\pgfpathlineto{\pgfqpoint{0.484393in}{1.464821in}}%
\pgfpathlineto{\pgfqpoint{0.429023in}{1.481634in}}%
\pgfpathlineto{\pgfqpoint{0.431912in}{1.491802in}}%
\pgfpathlineto{\pgfqpoint{0.428385in}{1.494011in}}%
\pgfpathlineto{\pgfqpoint{0.414605in}{1.487528in}}%
\pgfpathlineto{\pgfqpoint{0.408114in}{1.487902in}}%
\pgfpathlineto{\pgfqpoint{0.405051in}{1.484139in}}%
\pgfpathlineto{\pgfqpoint{0.407088in}{1.477930in}}%
\pgfpathlineto{\pgfqpoint{0.399966in}{1.478464in}}%
\pgfpathlineto{\pgfqpoint{0.397986in}{1.484025in}}%
\pgfpathlineto{\pgfqpoint{0.391596in}{1.486250in}}%
\pgfpathlineto{\pgfqpoint{0.391169in}{1.489496in}}%
\pgfpathlineto{\pgfqpoint{0.386307in}{1.493716in}}%
\pgfpathlineto{\pgfqpoint{0.385689in}{1.509253in}}%
\pgfpathlineto{\pgfqpoint{0.378145in}{1.511673in}}%
\pgfpathlineto{\pgfqpoint{0.381028in}{1.514238in}}%
\pgfpathlineto{\pgfqpoint{0.381897in}{1.521054in}}%
\pgfpathlineto{\pgfqpoint{0.380582in}{1.527032in}}%
\pgfpathlineto{\pgfqpoint{0.384988in}{1.532462in}}%
\pgfpathlineto{\pgfqpoint{0.385318in}{1.540362in}}%
\pgfpathlineto{\pgfqpoint{0.393613in}{1.542963in}}%
\pgfpathlineto{\pgfqpoint{0.397486in}{1.548285in}}%
\pgfpathlineto{\pgfqpoint{0.455007in}{1.530523in}}%
\pgfpathlineto{\pgfqpoint{0.474662in}{1.592743in}}%
\pgfpathlineto{\pgfqpoint{0.476029in}{1.597225in}}%
\pgfpathlineto{\pgfqpoint{0.485002in}{1.594976in}}%
\pgfpathlineto{\pgfqpoint{0.487657in}{1.599535in}}%
\pgfpathlineto{\pgfqpoint{0.493766in}{1.605432in}}%
\pgfpathlineto{\pgfqpoint{0.493741in}{1.612121in}}%
\pgfpathlineto{\pgfqpoint{0.490196in}{1.618627in}}%
\pgfpathlineto{\pgfqpoint{0.492514in}{1.625981in}}%
\pgfpathlineto{\pgfqpoint{0.499509in}{1.627848in}}%
\pgfpathlineto{\pgfqpoint{0.530203in}{1.618670in}}%
\pgfpathlineto{\pgfqpoint{0.524578in}{1.602426in}}%
\pgfpathlineto{\pgfqpoint{0.513001in}{1.563845in}}%
\pgfpathlineto{\pgfqpoint{0.534635in}{1.557261in}}%
\pgfpathlineto{\pgfqpoint{0.520911in}{1.509711in}}%
\pgfpathlineto{\pgfqpoint{0.562393in}{1.497838in}}%
\pgfpathlineto{\pgfqpoint{0.547844in}{1.446415in}}%
\pgfpathlineto{\pgfqpoint{0.535789in}{1.402977in}}%
\pgfpathlineto{\pgfqpoint{0.521615in}{1.353758in}}%
\pgfpathclose%
\pgfusepath{fill}%
\end{pgfscope}%
\begin{pgfscope}%
\pgfpathrectangle{\pgfqpoint{0.100000in}{0.100000in}}{\pgfqpoint{3.007045in}{1.925000in}}%
\pgfusepath{clip}%
\pgfsetbuttcap%
\pgfsetmiterjoin%
\definecolor{currentfill}{rgb}{0.391865,0.663406,0.828389}%
\pgfsetfillcolor{currentfill}%
\pgfsetlinewidth{0.000000pt}%
\definecolor{currentstroke}{rgb}{0.000000,0.000000,0.000000}%
\pgfsetstrokecolor{currentstroke}%
\pgfsetstrokeopacity{0.000000}%
\pgfsetdash{}{0pt}%
\pgfpathmoveto{\pgfqpoint{2.213486in}{1.315239in}}%
\pgfpathlineto{\pgfqpoint{2.215428in}{1.298031in}}%
\pgfpathlineto{\pgfqpoint{2.195043in}{1.295965in}}%
\pgfpathlineto{\pgfqpoint{2.172150in}{1.293918in}}%
\pgfpathlineto{\pgfqpoint{2.168897in}{1.329206in}}%
\pgfpathlineto{\pgfqpoint{2.180297in}{1.338873in}}%
\pgfpathlineto{\pgfqpoint{2.184660in}{1.346048in}}%
\pgfpathlineto{\pgfqpoint{2.188143in}{1.357873in}}%
\pgfpathlineto{\pgfqpoint{2.193342in}{1.366778in}}%
\pgfpathlineto{\pgfqpoint{2.200252in}{1.367456in}}%
\pgfpathlineto{\pgfqpoint{2.201376in}{1.356174in}}%
\pgfpathlineto{\pgfqpoint{2.223806in}{1.358262in}}%
\pgfpathlineto{\pgfqpoint{2.225748in}{1.340539in}}%
\pgfpathlineto{\pgfqpoint{2.224488in}{1.337759in}}%
\pgfpathlineto{\pgfqpoint{2.211201in}{1.336530in}}%
\pgfpathlineto{\pgfqpoint{2.213486in}{1.315239in}}%
\pgfpathclose%
\pgfusepath{fill}%
\end{pgfscope}%
\begin{pgfscope}%
\pgfpathrectangle{\pgfqpoint{0.100000in}{0.100000in}}{\pgfqpoint{3.007045in}{1.925000in}}%
\pgfusepath{clip}%
\pgfsetbuttcap%
\pgfsetmiterjoin%
\definecolor{currentfill}{rgb}{0.331334,0.622068,0.804767}%
\pgfsetfillcolor{currentfill}%
\pgfsetlinewidth{0.000000pt}%
\definecolor{currentstroke}{rgb}{0.000000,0.000000,0.000000}%
\pgfsetstrokecolor{currentstroke}%
\pgfsetstrokeopacity{0.000000}%
\pgfsetdash{}{0pt}%
\pgfpathmoveto{\pgfqpoint{2.013078in}{0.656435in}}%
\pgfpathlineto{\pgfqpoint{2.013470in}{0.675004in}}%
\pgfpathlineto{\pgfqpoint{2.011898in}{0.691368in}}%
\pgfpathlineto{\pgfqpoint{2.012854in}{0.699525in}}%
\pgfpathlineto{\pgfqpoint{2.016613in}{0.701553in}}%
\pgfpathlineto{\pgfqpoint{2.018058in}{0.708669in}}%
\pgfpathlineto{\pgfqpoint{2.020779in}{0.710545in}}%
\pgfpathlineto{\pgfqpoint{2.027190in}{0.716066in}}%
\pgfpathlineto{\pgfqpoint{2.040123in}{0.721921in}}%
\pgfpathlineto{\pgfqpoint{2.048959in}{0.728755in}}%
\pgfpathlineto{\pgfqpoint{2.052237in}{0.736873in}}%
\pgfpathlineto{\pgfqpoint{2.065317in}{0.737960in}}%
\pgfpathlineto{\pgfqpoint{2.065136in}{0.740946in}}%
\pgfpathlineto{\pgfqpoint{2.087936in}{0.742436in}}%
\pgfpathlineto{\pgfqpoint{2.090671in}{0.695575in}}%
\pgfpathlineto{\pgfqpoint{2.067842in}{0.694087in}}%
\pgfpathlineto{\pgfqpoint{2.068533in}{0.682516in}}%
\pgfpathlineto{\pgfqpoint{2.072910in}{0.680552in}}%
\pgfpathlineto{\pgfqpoint{2.073893in}{0.665116in}}%
\pgfpathlineto{\pgfqpoint{2.055851in}{0.662822in}}%
\pgfpathlineto{\pgfqpoint{2.047327in}{0.661747in}}%
\pgfpathlineto{\pgfqpoint{2.038852in}{0.657801in}}%
\pgfpathlineto{\pgfqpoint{2.013078in}{0.656435in}}%
\pgfpathclose%
\pgfusepath{fill}%
\end{pgfscope}%
\begin{pgfscope}%
\pgfpathrectangle{\pgfqpoint{0.100000in}{0.100000in}}{\pgfqpoint{3.007045in}{1.925000in}}%
\pgfusepath{clip}%
\pgfsetbuttcap%
\pgfsetmiterjoin%
\definecolor{currentfill}{rgb}{0.316201,0.611734,0.798862}%
\pgfsetfillcolor{currentfill}%
\pgfsetlinewidth{0.000000pt}%
\definecolor{currentstroke}{rgb}{0.000000,0.000000,0.000000}%
\pgfsetstrokecolor{currentstroke}%
\pgfsetstrokeopacity{0.000000}%
\pgfsetdash{}{0pt}%
\pgfpathmoveto{\pgfqpoint{1.886783in}{0.652756in}}%
\pgfpathlineto{\pgfqpoint{1.890709in}{0.647179in}}%
\pgfpathlineto{\pgfqpoint{1.896682in}{0.645552in}}%
\pgfpathlineto{\pgfqpoint{1.901476in}{0.640393in}}%
\pgfpathlineto{\pgfqpoint{1.890720in}{0.630209in}}%
\pgfpathlineto{\pgfqpoint{1.886870in}{0.628687in}}%
\pgfpathlineto{\pgfqpoint{1.883960in}{0.629796in}}%
\pgfpathlineto{\pgfqpoint{1.872399in}{0.629573in}}%
\pgfpathlineto{\pgfqpoint{1.872143in}{0.641119in}}%
\pgfpathlineto{\pgfqpoint{1.866438in}{0.640996in}}%
\pgfpathlineto{\pgfqpoint{1.866139in}{0.652691in}}%
\pgfpathlineto{\pgfqpoint{1.860335in}{0.652602in}}%
\pgfpathlineto{\pgfqpoint{1.860218in}{0.661283in}}%
\pgfpathlineto{\pgfqpoint{1.835531in}{0.660856in}}%
\pgfpathlineto{\pgfqpoint{1.827733in}{0.669764in}}%
\pgfpathlineto{\pgfqpoint{1.826203in}{0.670688in}}%
\pgfpathlineto{\pgfqpoint{1.825570in}{0.729490in}}%
\pgfpathlineto{\pgfqpoint{1.825469in}{0.738647in}}%
\pgfpathlineto{\pgfqpoint{1.838686in}{0.738812in}}%
\pgfpathlineto{\pgfqpoint{1.884006in}{0.739484in}}%
\pgfpathlineto{\pgfqpoint{1.898651in}{0.739626in}}%
\pgfpathlineto{\pgfqpoint{1.899049in}{0.722750in}}%
\pgfpathlineto{\pgfqpoint{1.893306in}{0.722562in}}%
\pgfpathlineto{\pgfqpoint{1.890590in}{0.711973in}}%
\pgfpathlineto{\pgfqpoint{1.890952in}{0.702305in}}%
\pgfpathlineto{\pgfqpoint{1.896662in}{0.702398in}}%
\pgfpathlineto{\pgfqpoint{1.897220in}{0.687919in}}%
\pgfpathlineto{\pgfqpoint{1.895056in}{0.682024in}}%
\pgfpathlineto{\pgfqpoint{1.888095in}{0.681930in}}%
\pgfpathlineto{\pgfqpoint{1.891446in}{0.668175in}}%
\pgfpathlineto{\pgfqpoint{1.888359in}{0.662427in}}%
\pgfpathlineto{\pgfqpoint{1.886783in}{0.652756in}}%
\pgfpathclose%
\pgfusepath{fill}%
\end{pgfscope}%
\begin{pgfscope}%
\pgfpathrectangle{\pgfqpoint{0.100000in}{0.100000in}}{\pgfqpoint{3.007045in}{1.925000in}}%
\pgfusepath{clip}%
\pgfsetbuttcap%
\pgfsetmiterjoin%
\definecolor{currentfill}{rgb}{0.548235,0.752972,0.866959}%
\pgfsetfillcolor{currentfill}%
\pgfsetlinewidth{0.000000pt}%
\definecolor{currentstroke}{rgb}{0.000000,0.000000,0.000000}%
\pgfsetstrokecolor{currentstroke}%
\pgfsetstrokeopacity{0.000000}%
\pgfsetdash{}{0pt}%
\pgfpathmoveto{\pgfqpoint{2.399943in}{0.955887in}}%
\pgfpathlineto{\pgfqpoint{2.388077in}{0.945524in}}%
\pgfpathlineto{\pgfqpoint{2.379259in}{0.944820in}}%
\pgfpathlineto{\pgfqpoint{2.378270in}{0.952775in}}%
\pgfpathlineto{\pgfqpoint{2.374820in}{0.956140in}}%
\pgfpathlineto{\pgfqpoint{2.369682in}{0.964768in}}%
\pgfpathlineto{\pgfqpoint{2.376395in}{0.971244in}}%
\pgfpathlineto{\pgfqpoint{2.370660in}{0.984801in}}%
\pgfpathlineto{\pgfqpoint{2.372662in}{0.995532in}}%
\pgfpathlineto{\pgfqpoint{2.380497in}{0.997923in}}%
\pgfpathlineto{\pgfqpoint{2.388640in}{1.002575in}}%
\pgfpathlineto{\pgfqpoint{2.392615in}{1.001420in}}%
\pgfpathlineto{\pgfqpoint{2.394844in}{0.994687in}}%
\pgfpathlineto{\pgfqpoint{2.399497in}{0.998207in}}%
\pgfpathlineto{\pgfqpoint{2.404364in}{0.992821in}}%
\pgfpathlineto{\pgfqpoint{2.400322in}{0.986669in}}%
\pgfpathlineto{\pgfqpoint{2.406942in}{0.983024in}}%
\pgfpathlineto{\pgfqpoint{2.416656in}{0.974681in}}%
\pgfpathlineto{\pgfqpoint{2.414244in}{0.964014in}}%
\pgfpathlineto{\pgfqpoint{2.404400in}{0.961065in}}%
\pgfpathlineto{\pgfqpoint{2.399943in}{0.955887in}}%
\pgfpathclose%
\pgfusepath{fill}%
\end{pgfscope}%
\begin{pgfscope}%
\pgfpathrectangle{\pgfqpoint{0.100000in}{0.100000in}}{\pgfqpoint{3.007045in}{1.925000in}}%
\pgfusepath{clip}%
\pgfsetbuttcap%
\pgfsetmiterjoin%
\definecolor{currentfill}{rgb}{0.296025,0.597955,0.790988}%
\pgfsetfillcolor{currentfill}%
\pgfsetlinewidth{0.000000pt}%
\definecolor{currentstroke}{rgb}{0.000000,0.000000,0.000000}%
\pgfsetstrokecolor{currentstroke}%
\pgfsetstrokeopacity{0.000000}%
\pgfsetdash{}{0pt}%
\pgfpathmoveto{\pgfqpoint{1.092509in}{1.504079in}}%
\pgfpathlineto{\pgfqpoint{1.102940in}{1.502630in}}%
\pgfpathlineto{\pgfqpoint{1.103807in}{1.508400in}}%
\pgfpathlineto{\pgfqpoint{1.112286in}{1.507103in}}%
\pgfpathlineto{\pgfqpoint{1.113003in}{1.512751in}}%
\pgfpathlineto{\pgfqpoint{1.118845in}{1.511867in}}%
\pgfpathlineto{\pgfqpoint{1.119707in}{1.517659in}}%
\pgfpathlineto{\pgfqpoint{1.131047in}{1.515926in}}%
\pgfpathlineto{\pgfqpoint{1.131931in}{1.521748in}}%
\pgfpathlineto{\pgfqpoint{1.197829in}{1.512182in}}%
\pgfpathlineto{\pgfqpoint{1.193681in}{1.468364in}}%
\pgfpathlineto{\pgfqpoint{1.173551in}{1.471182in}}%
\pgfpathlineto{\pgfqpoint{1.170426in}{1.469676in}}%
\pgfpathlineto{\pgfqpoint{1.143778in}{1.473368in}}%
\pgfpathlineto{\pgfqpoint{1.135060in}{1.473922in}}%
\pgfpathlineto{\pgfqpoint{1.118840in}{1.481195in}}%
\pgfpathlineto{\pgfqpoint{1.119290in}{1.484034in}}%
\pgfpathlineto{\pgfqpoint{1.108398in}{1.488604in}}%
\pgfpathlineto{\pgfqpoint{1.097179in}{1.490393in}}%
\pgfpathlineto{\pgfqpoint{1.098004in}{1.495582in}}%
\pgfpathlineto{\pgfqpoint{1.092509in}{1.504079in}}%
\pgfpathclose%
\pgfusepath{fill}%
\end{pgfscope}%
\begin{pgfscope}%
\pgfpathrectangle{\pgfqpoint{0.100000in}{0.100000in}}{\pgfqpoint{3.007045in}{1.925000in}}%
\pgfusepath{clip}%
\pgfsetbuttcap%
\pgfsetmiterjoin%
\definecolor{currentfill}{rgb}{0.396909,0.666851,0.830358}%
\pgfsetfillcolor{currentfill}%
\pgfsetlinewidth{0.000000pt}%
\definecolor{currentstroke}{rgb}{0.000000,0.000000,0.000000}%
\pgfsetstrokecolor{currentstroke}%
\pgfsetstrokeopacity{0.000000}%
\pgfsetdash{}{0pt}%
\pgfpathmoveto{\pgfqpoint{2.119150in}{1.400006in}}%
\pgfpathlineto{\pgfqpoint{2.073275in}{1.396646in}}%
\pgfpathlineto{\pgfqpoint{2.061841in}{1.396171in}}%
\pgfpathlineto{\pgfqpoint{2.058743in}{1.447959in}}%
\pgfpathlineto{\pgfqpoint{2.087655in}{1.449833in}}%
\pgfpathlineto{\pgfqpoint{2.088107in}{1.444073in}}%
\pgfpathlineto{\pgfqpoint{2.099536in}{1.444895in}}%
\pgfpathlineto{\pgfqpoint{2.117083in}{1.446231in}}%
\pgfpathlineto{\pgfqpoint{2.116939in}{1.440991in}}%
\pgfpathlineto{\pgfqpoint{2.114427in}{1.436064in}}%
\pgfpathlineto{\pgfqpoint{2.113005in}{1.425983in}}%
\pgfpathlineto{\pgfqpoint{2.119150in}{1.400006in}}%
\pgfpathclose%
\pgfusepath{fill}%
\end{pgfscope}%
\begin{pgfscope}%
\pgfpathrectangle{\pgfqpoint{0.100000in}{0.100000in}}{\pgfqpoint{3.007045in}{1.925000in}}%
\pgfusepath{clip}%
\pgfsetbuttcap%
\pgfsetmiterjoin%
\definecolor{currentfill}{rgb}{0.504314,0.728858,0.857486}%
\pgfsetfillcolor{currentfill}%
\pgfsetlinewidth{0.000000pt}%
\definecolor{currentstroke}{rgb}{0.000000,0.000000,0.000000}%
\pgfsetstrokecolor{currentstroke}%
\pgfsetstrokeopacity{0.000000}%
\pgfsetdash{}{0pt}%
\pgfpathmoveto{\pgfqpoint{0.618157in}{0.899792in}}%
\pgfpathlineto{\pgfqpoint{0.613032in}{0.877119in}}%
\pgfpathlineto{\pgfqpoint{0.611833in}{0.877391in}}%
\pgfpathlineto{\pgfqpoint{0.604721in}{0.848024in}}%
\pgfpathlineto{\pgfqpoint{0.548111in}{0.856043in}}%
\pgfpathlineto{\pgfqpoint{0.548819in}{0.861290in}}%
\pgfpathlineto{\pgfqpoint{0.543625in}{0.866151in}}%
\pgfpathlineto{\pgfqpoint{0.546518in}{0.879093in}}%
\pgfpathlineto{\pgfqpoint{0.547290in}{0.888052in}}%
\pgfpathlineto{\pgfqpoint{0.545659in}{0.899091in}}%
\pgfpathlineto{\pgfqpoint{0.540654in}{0.912060in}}%
\pgfpathlineto{\pgfqpoint{0.536458in}{0.917204in}}%
\pgfpathlineto{\pgfqpoint{0.538502in}{0.921211in}}%
\pgfpathlineto{\pgfqpoint{0.543010in}{0.923557in}}%
\pgfpathlineto{\pgfqpoint{0.550214in}{0.920766in}}%
\pgfpathlineto{\pgfqpoint{0.556193in}{0.915257in}}%
\pgfpathlineto{\pgfqpoint{0.567402in}{0.912130in}}%
\pgfpathlineto{\pgfqpoint{0.618157in}{0.899792in}}%
\pgfpathclose%
\pgfusepath{fill}%
\end{pgfscope}%
\begin{pgfscope}%
\pgfpathrectangle{\pgfqpoint{0.100000in}{0.100000in}}{\pgfqpoint{3.007045in}{1.925000in}}%
\pgfusepath{clip}%
\pgfsetbuttcap%
\pgfsetmiterjoin%
\definecolor{currentfill}{rgb}{0.745713,0.845752,0.926490}%
\pgfsetfillcolor{currentfill}%
\pgfsetlinewidth{0.000000pt}%
\definecolor{currentstroke}{rgb}{0.000000,0.000000,0.000000}%
\pgfsetstrokecolor{currentstroke}%
\pgfsetstrokeopacity{0.000000}%
\pgfsetdash{}{0pt}%
\pgfpathmoveto{\pgfqpoint{2.000756in}{1.600279in}}%
\pgfpathlineto{\pgfqpoint{1.988873in}{1.599607in}}%
\pgfpathlineto{\pgfqpoint{1.988231in}{1.611105in}}%
\pgfpathlineto{\pgfqpoint{1.982576in}{1.610763in}}%
\pgfpathlineto{\pgfqpoint{1.982301in}{1.616494in}}%
\pgfpathlineto{\pgfqpoint{1.976511in}{1.616232in}}%
\pgfpathlineto{\pgfqpoint{1.975507in}{1.638992in}}%
\pgfpathlineto{\pgfqpoint{1.980609in}{1.637782in}}%
\pgfpathlineto{\pgfqpoint{1.984352in}{1.640321in}}%
\pgfpathlineto{\pgfqpoint{1.998858in}{1.646114in}}%
\pgfpathlineto{\pgfqpoint{2.009289in}{1.656368in}}%
\pgfpathlineto{\pgfqpoint{2.026300in}{1.659095in}}%
\pgfpathlineto{\pgfqpoint{2.033497in}{1.663365in}}%
\pgfpathlineto{\pgfqpoint{2.037806in}{1.669257in}}%
\pgfpathlineto{\pgfqpoint{2.044910in}{1.670501in}}%
\pgfpathlineto{\pgfqpoint{2.047218in}{1.673072in}}%
\pgfpathlineto{\pgfqpoint{2.048453in}{1.655274in}}%
\pgfpathlineto{\pgfqpoint{2.051916in}{1.649810in}}%
\pgfpathlineto{\pgfqpoint{2.046234in}{1.649481in}}%
\pgfpathlineto{\pgfqpoint{2.047301in}{1.632207in}}%
\pgfpathlineto{\pgfqpoint{2.048726in}{1.610846in}}%
\pgfpathlineto{\pgfqpoint{2.043917in}{1.613272in}}%
\pgfpathlineto{\pgfqpoint{2.004866in}{1.621636in}}%
\pgfpathlineto{\pgfqpoint{2.006057in}{1.600572in}}%
\pgfpathlineto{\pgfqpoint{2.000756in}{1.600279in}}%
\pgfpathclose%
\pgfusepath{fill}%
\end{pgfscope}%
\begin{pgfscope}%
\pgfpathrectangle{\pgfqpoint{0.100000in}{0.100000in}}{\pgfqpoint{3.007045in}{1.925000in}}%
\pgfusepath{clip}%
\pgfsetbuttcap%
\pgfsetmiterjoin%
\definecolor{currentfill}{rgb}{0.401953,0.670296,0.832326}%
\pgfsetfillcolor{currentfill}%
\pgfsetlinewidth{0.000000pt}%
\definecolor{currentstroke}{rgb}{0.000000,0.000000,0.000000}%
\pgfsetstrokecolor{currentstroke}%
\pgfsetstrokeopacity{0.000000}%
\pgfsetdash{}{0pt}%
\pgfpathmoveto{\pgfqpoint{1.330895in}{1.191503in}}%
\pgfpathlineto{\pgfqpoint{1.358836in}{1.188865in}}%
\pgfpathlineto{\pgfqpoint{1.355807in}{1.160158in}}%
\pgfpathlineto{\pgfqpoint{1.353166in}{1.131470in}}%
\pgfpathlineto{\pgfqpoint{1.351745in}{1.120688in}}%
\pgfpathlineto{\pgfqpoint{1.334664in}{1.121680in}}%
\pgfpathlineto{\pgfqpoint{1.306529in}{1.124802in}}%
\pgfpathlineto{\pgfqpoint{1.308723in}{1.147516in}}%
\pgfpathlineto{\pgfqpoint{1.325845in}{1.145704in}}%
\pgfpathlineto{\pgfqpoint{1.330895in}{1.191503in}}%
\pgfpathclose%
\pgfusepath{fill}%
\end{pgfscope}%
\begin{pgfscope}%
\pgfpathrectangle{\pgfqpoint{0.100000in}{0.100000in}}{\pgfqpoint{3.007045in}{1.925000in}}%
\pgfusepath{clip}%
\pgfsetbuttcap%
\pgfsetmiterjoin%
\definecolor{currentfill}{rgb}{0.316201,0.611734,0.798862}%
\pgfsetfillcolor{currentfill}%
\pgfsetlinewidth{0.000000pt}%
\definecolor{currentstroke}{rgb}{0.000000,0.000000,0.000000}%
\pgfsetstrokecolor{currentstroke}%
\pgfsetstrokeopacity{0.000000}%
\pgfsetdash{}{0pt}%
\pgfpathmoveto{\pgfqpoint{1.768296in}{1.069654in}}%
\pgfpathlineto{\pgfqpoint{1.767740in}{1.049552in}}%
\pgfpathlineto{\pgfqpoint{1.722263in}{1.049928in}}%
\pgfpathlineto{\pgfqpoint{1.722556in}{1.070090in}}%
\pgfpathlineto{\pgfqpoint{1.745387in}{1.069764in}}%
\pgfpathlineto{\pgfqpoint{1.746105in}{1.095854in}}%
\pgfpathlineto{\pgfqpoint{1.746586in}{1.115888in}}%
\pgfpathlineto{\pgfqpoint{1.769426in}{1.115817in}}%
\pgfpathlineto{\pgfqpoint{1.768947in}{1.112944in}}%
\pgfpathlineto{\pgfqpoint{1.768960in}{1.084214in}}%
\pgfpathlineto{\pgfqpoint{1.768296in}{1.069654in}}%
\pgfpathclose%
\pgfusepath{fill}%
\end{pgfscope}%
\begin{pgfscope}%
\pgfpathrectangle{\pgfqpoint{0.100000in}{0.100000in}}{\pgfqpoint{3.007045in}{1.925000in}}%
\pgfusepath{clip}%
\pgfsetbuttcap%
\pgfsetmiterjoin%
\definecolor{currentfill}{rgb}{0.516863,0.735748,0.860192}%
\pgfsetfillcolor{currentfill}%
\pgfsetlinewidth{0.000000pt}%
\definecolor{currentstroke}{rgb}{0.000000,0.000000,0.000000}%
\pgfsetstrokecolor{currentstroke}%
\pgfsetstrokeopacity{0.000000}%
\pgfsetdash{}{0pt}%
\pgfpathmoveto{\pgfqpoint{1.980060in}{1.429597in}}%
\pgfpathlineto{\pgfqpoint{1.954322in}{1.428443in}}%
\pgfpathlineto{\pgfqpoint{1.953526in}{1.433541in}}%
\pgfpathlineto{\pgfqpoint{1.934699in}{1.432796in}}%
\pgfpathlineto{\pgfqpoint{1.928977in}{1.432597in}}%
\pgfpathlineto{\pgfqpoint{1.928151in}{1.455464in}}%
\pgfpathlineto{\pgfqpoint{1.949409in}{1.456290in}}%
\pgfpathlineto{\pgfqpoint{1.942334in}{1.465087in}}%
\pgfpathlineto{\pgfqpoint{1.947232in}{1.465292in}}%
\pgfpathlineto{\pgfqpoint{1.947712in}{1.470634in}}%
\pgfpathlineto{\pgfqpoint{1.963568in}{1.471696in}}%
\pgfpathlineto{\pgfqpoint{1.966485in}{1.477622in}}%
\pgfpathlineto{\pgfqpoint{1.994624in}{1.478781in}}%
\pgfpathlineto{\pgfqpoint{1.994330in}{1.484956in}}%
\pgfpathlineto{\pgfqpoint{2.013699in}{1.486047in}}%
\pgfpathlineto{\pgfqpoint{2.008307in}{1.475250in}}%
\pgfpathlineto{\pgfqpoint{2.012462in}{1.466036in}}%
\pgfpathlineto{\pgfqpoint{2.012501in}{1.460181in}}%
\pgfpathlineto{\pgfqpoint{2.017195in}{1.455365in}}%
\pgfpathlineto{\pgfqpoint{2.021528in}{1.446204in}}%
\pgfpathlineto{\pgfqpoint{1.996384in}{1.444829in}}%
\pgfpathlineto{\pgfqpoint{1.996750in}{1.439088in}}%
\pgfpathlineto{\pgfqpoint{1.979503in}{1.438168in}}%
\pgfpathlineto{\pgfqpoint{1.980060in}{1.429597in}}%
\pgfpathclose%
\pgfusepath{fill}%
\end{pgfscope}%
\begin{pgfscope}%
\pgfpathrectangle{\pgfqpoint{0.100000in}{0.100000in}}{\pgfqpoint{3.007045in}{1.925000in}}%
\pgfusepath{clip}%
\pgfsetbuttcap%
\pgfsetmiterjoin%
\definecolor{currentfill}{rgb}{0.435294,0.690965,0.842599}%
\pgfsetfillcolor{currentfill}%
\pgfsetlinewidth{0.000000pt}%
\definecolor{currentstroke}{rgb}{0.000000,0.000000,0.000000}%
\pgfsetstrokecolor{currentstroke}%
\pgfsetstrokeopacity{0.000000}%
\pgfsetdash{}{0pt}%
\pgfpathmoveto{\pgfqpoint{1.045450in}{1.802919in}}%
\pgfpathlineto{\pgfqpoint{1.043735in}{1.793392in}}%
\pgfpathlineto{\pgfqpoint{1.018457in}{1.797935in}}%
\pgfpathlineto{\pgfqpoint{1.016942in}{1.805156in}}%
\pgfpathlineto{\pgfqpoint{1.011784in}{1.809131in}}%
\pgfpathlineto{\pgfqpoint{0.975124in}{1.816423in}}%
\pgfpathlineto{\pgfqpoint{0.966264in}{1.820814in}}%
\pgfpathlineto{\pgfqpoint{0.967054in}{1.825877in}}%
\pgfpathlineto{\pgfqpoint{0.961768in}{1.833418in}}%
\pgfpathlineto{\pgfqpoint{0.964152in}{1.834807in}}%
\pgfpathlineto{\pgfqpoint{0.959285in}{1.841945in}}%
\pgfpathlineto{\pgfqpoint{0.961059in}{1.847574in}}%
\pgfpathlineto{\pgfqpoint{0.949618in}{1.854808in}}%
\pgfpathlineto{\pgfqpoint{0.953218in}{1.860951in}}%
\pgfpathlineto{\pgfqpoint{0.949121in}{1.871112in}}%
\pgfpathlineto{\pgfqpoint{0.943633in}{1.869658in}}%
\pgfpathlineto{\pgfqpoint{0.941225in}{1.874500in}}%
\pgfpathlineto{\pgfqpoint{0.941334in}{1.882540in}}%
\pgfpathlineto{\pgfqpoint{0.981051in}{1.874255in}}%
\pgfpathlineto{\pgfqpoint{1.010982in}{1.868345in}}%
\pgfpathlineto{\pgfqpoint{1.061807in}{1.858712in}}%
\pgfpathlineto{\pgfqpoint{1.052373in}{1.807583in}}%
\pgfpathlineto{\pgfqpoint{1.046502in}{1.808652in}}%
\pgfpathlineto{\pgfqpoint{1.045450in}{1.802919in}}%
\pgfpathclose%
\pgfusepath{fill}%
\end{pgfscope}%
\begin{pgfscope}%
\pgfpathrectangle{\pgfqpoint{0.100000in}{0.100000in}}{\pgfqpoint{3.007045in}{1.925000in}}%
\pgfusepath{clip}%
\pgfsetbuttcap%
\pgfsetmiterjoin%
\definecolor{currentfill}{rgb}{0.240046,0.553772,0.766797}%
\pgfsetfillcolor{currentfill}%
\pgfsetlinewidth{0.000000pt}%
\definecolor{currentstroke}{rgb}{0.000000,0.000000,0.000000}%
\pgfsetstrokecolor{currentstroke}%
\pgfsetstrokeopacity{0.000000}%
\pgfsetdash{}{0pt}%
\pgfpathmoveto{\pgfqpoint{1.522557in}{1.725806in}}%
\pgfpathlineto{\pgfqpoint{1.465127in}{1.729765in}}%
\pgfpathlineto{\pgfqpoint{1.466009in}{1.741430in}}%
\pgfpathlineto{\pgfqpoint{1.463751in}{1.741609in}}%
\pgfpathlineto{\pgfqpoint{1.465533in}{1.764842in}}%
\pgfpathlineto{\pgfqpoint{1.463268in}{1.765000in}}%
\pgfpathlineto{\pgfqpoint{1.464168in}{1.776657in}}%
\pgfpathlineto{\pgfqpoint{1.452658in}{1.777477in}}%
\pgfpathlineto{\pgfqpoint{1.453609in}{1.789081in}}%
\pgfpathlineto{\pgfqpoint{1.457171in}{1.788804in}}%
\pgfpathlineto{\pgfqpoint{1.457650in}{1.794593in}}%
\pgfpathlineto{\pgfqpoint{1.463375in}{1.794138in}}%
\pgfpathlineto{\pgfqpoint{1.464391in}{1.806947in}}%
\pgfpathlineto{\pgfqpoint{1.508366in}{1.803801in}}%
\pgfpathlineto{\pgfqpoint{1.544920in}{1.801546in}}%
\pgfpathlineto{\pgfqpoint{1.543838in}{1.782832in}}%
\pgfpathlineto{\pgfqpoint{1.544810in}{1.771147in}}%
\pgfpathlineto{\pgfqpoint{1.539033in}{1.771449in}}%
\pgfpathlineto{\pgfqpoint{1.540108in}{1.759776in}}%
\pgfpathlineto{\pgfqpoint{1.538782in}{1.736495in}}%
\pgfpathlineto{\pgfqpoint{1.539856in}{1.724787in}}%
\pgfpathlineto{\pgfqpoint{1.522557in}{1.725806in}}%
\pgfpathclose%
\pgfusepath{fill}%
\end{pgfscope}%
\begin{pgfscope}%
\pgfpathrectangle{\pgfqpoint{0.100000in}{0.100000in}}{\pgfqpoint{3.007045in}{1.925000in}}%
\pgfusepath{clip}%
\pgfsetbuttcap%
\pgfsetmiterjoin%
\definecolor{currentfill}{rgb}{0.124029,0.436248,0.704421}%
\pgfsetfillcolor{currentfill}%
\pgfsetlinewidth{0.000000pt}%
\definecolor{currentstroke}{rgb}{0.000000,0.000000,0.000000}%
\pgfsetstrokecolor{currentstroke}%
\pgfsetstrokeopacity{0.000000}%
\pgfsetdash{}{0pt}%
\pgfpathmoveto{\pgfqpoint{1.591111in}{1.462290in}}%
\pgfpathlineto{\pgfqpoint{1.590530in}{1.433489in}}%
\pgfpathlineto{\pgfqpoint{1.566436in}{1.434578in}}%
\pgfpathlineto{\pgfqpoint{1.550851in}{1.435352in}}%
\pgfpathlineto{\pgfqpoint{1.551671in}{1.451429in}}%
\pgfpathlineto{\pgfqpoint{1.541165in}{1.448356in}}%
\pgfpathlineto{\pgfqpoint{1.531144in}{1.449134in}}%
\pgfpathlineto{\pgfqpoint{1.522608in}{1.451050in}}%
\pgfpathlineto{\pgfqpoint{1.517464in}{1.451554in}}%
\pgfpathlineto{\pgfqpoint{1.516674in}{1.460211in}}%
\pgfpathlineto{\pgfqpoint{1.517880in}{1.481436in}}%
\pgfpathlineto{\pgfqpoint{1.517980in}{1.483159in}}%
\pgfpathlineto{\pgfqpoint{1.538292in}{1.482046in}}%
\pgfpathlineto{\pgfqpoint{1.542308in}{1.477127in}}%
\pgfpathlineto{\pgfqpoint{1.546378in}{1.477641in}}%
\pgfpathlineto{\pgfqpoint{1.552579in}{1.474853in}}%
\pgfpathlineto{\pgfqpoint{1.551371in}{1.482830in}}%
\pgfpathlineto{\pgfqpoint{1.555401in}{1.480969in}}%
\pgfpathlineto{\pgfqpoint{1.586215in}{1.479799in}}%
\pgfpathlineto{\pgfqpoint{1.585428in}{1.462527in}}%
\pgfpathlineto{\pgfqpoint{1.591111in}{1.462290in}}%
\pgfpathclose%
\pgfusepath{fill}%
\end{pgfscope}%
\begin{pgfscope}%
\pgfpathrectangle{\pgfqpoint{0.100000in}{0.100000in}}{\pgfqpoint{3.007045in}{1.925000in}}%
\pgfusepath{clip}%
\pgfsetbuttcap%
\pgfsetmiterjoin%
\definecolor{currentfill}{rgb}{0.479216,0.715079,0.852072}%
\pgfsetfillcolor{currentfill}%
\pgfsetlinewidth{0.000000pt}%
\definecolor{currentstroke}{rgb}{0.000000,0.000000,0.000000}%
\pgfsetstrokecolor{currentstroke}%
\pgfsetstrokeopacity{0.000000}%
\pgfsetdash{}{0pt}%
\pgfpathmoveto{\pgfqpoint{2.429643in}{1.109432in}}%
\pgfpathlineto{\pgfqpoint{2.435308in}{1.112304in}}%
\pgfpathlineto{\pgfqpoint{2.440031in}{1.105220in}}%
\pgfpathlineto{\pgfqpoint{2.441207in}{1.098026in}}%
\pgfpathlineto{\pgfqpoint{2.446820in}{1.093542in}}%
\pgfpathlineto{\pgfqpoint{2.458232in}{1.094688in}}%
\pgfpathlineto{\pgfqpoint{2.456556in}{1.087012in}}%
\pgfpathlineto{\pgfqpoint{2.453185in}{1.084073in}}%
\pgfpathlineto{\pgfqpoint{2.450852in}{1.085445in}}%
\pgfpathlineto{\pgfqpoint{2.435154in}{1.067105in}}%
\pgfpathlineto{\pgfqpoint{2.423037in}{1.058784in}}%
\pgfpathlineto{\pgfqpoint{2.417226in}{1.060057in}}%
\pgfpathlineto{\pgfqpoint{2.412270in}{1.064935in}}%
\pgfpathlineto{\pgfqpoint{2.410452in}{1.072460in}}%
\pgfpathlineto{\pgfqpoint{2.400437in}{1.076137in}}%
\pgfpathlineto{\pgfqpoint{2.391489in}{1.082969in}}%
\pgfpathlineto{\pgfqpoint{2.382872in}{1.085587in}}%
\pgfpathlineto{\pgfqpoint{2.381551in}{1.091127in}}%
\pgfpathlineto{\pgfqpoint{2.389841in}{1.100458in}}%
\pgfpathlineto{\pgfqpoint{2.396141in}{1.102157in}}%
\pgfpathlineto{\pgfqpoint{2.392249in}{1.108544in}}%
\pgfpathlineto{\pgfqpoint{2.398193in}{1.116290in}}%
\pgfpathlineto{\pgfqpoint{2.396061in}{1.120463in}}%
\pgfpathlineto{\pgfqpoint{2.401483in}{1.125255in}}%
\pgfpathlineto{\pgfqpoint{2.411679in}{1.127319in}}%
\pgfpathlineto{\pgfqpoint{2.410869in}{1.119825in}}%
\pgfpathlineto{\pgfqpoint{2.421148in}{1.110912in}}%
\pgfpathlineto{\pgfqpoint{2.420050in}{1.106990in}}%
\pgfpathlineto{\pgfqpoint{2.425369in}{1.103463in}}%
\pgfpathlineto{\pgfqpoint{2.429643in}{1.109432in}}%
\pgfpathclose%
\pgfusepath{fill}%
\end{pgfscope}%
\begin{pgfscope}%
\pgfpathrectangle{\pgfqpoint{0.100000in}{0.100000in}}{\pgfqpoint{3.007045in}{1.925000in}}%
\pgfusepath{clip}%
\pgfsetbuttcap%
\pgfsetmiterjoin%
\definecolor{currentfill}{rgb}{0.498039,0.725413,0.856132}%
\pgfsetfillcolor{currentfill}%
\pgfsetlinewidth{0.000000pt}%
\definecolor{currentstroke}{rgb}{0.000000,0.000000,0.000000}%
\pgfsetstrokecolor{currentstroke}%
\pgfsetstrokeopacity{0.000000}%
\pgfsetdash{}{0pt}%
\pgfpathmoveto{\pgfqpoint{1.844917in}{1.391266in}}%
\pgfpathlineto{\pgfqpoint{1.822052in}{1.390899in}}%
\pgfpathlineto{\pgfqpoint{1.821649in}{1.429987in}}%
\pgfpathlineto{\pgfqpoint{1.837085in}{1.430186in}}%
\pgfpathlineto{\pgfqpoint{1.866993in}{1.430735in}}%
\pgfpathlineto{\pgfqpoint{1.867818in}{1.391658in}}%
\pgfpathlineto{\pgfqpoint{1.844917in}{1.391266in}}%
\pgfpathclose%
\pgfusepath{fill}%
\end{pgfscope}%
\begin{pgfscope}%
\pgfpathrectangle{\pgfqpoint{0.100000in}{0.100000in}}{\pgfqpoint{3.007045in}{1.925000in}}%
\pgfusepath{clip}%
\pgfsetbuttcap%
\pgfsetmiterjoin%
\definecolor{currentfill}{rgb}{0.627605,0.795556,0.885152}%
\pgfsetfillcolor{currentfill}%
\pgfsetlinewidth{0.000000pt}%
\definecolor{currentstroke}{rgb}{0.000000,0.000000,0.000000}%
\pgfsetstrokecolor{currentstroke}%
\pgfsetstrokeopacity{0.000000}%
\pgfsetdash{}{0pt}%
\pgfpathmoveto{\pgfqpoint{2.463450in}{1.158871in}}%
\pgfpathlineto{\pgfqpoint{2.458391in}{1.154238in}}%
\pgfpathlineto{\pgfqpoint{2.450037in}{1.161421in}}%
\pgfpathlineto{\pgfqpoint{2.441696in}{1.172351in}}%
\pgfpathlineto{\pgfqpoint{2.447277in}{1.183280in}}%
\pgfpathlineto{\pgfqpoint{2.446120in}{1.194088in}}%
\pgfpathlineto{\pgfqpoint{2.446995in}{1.200159in}}%
\pgfpathlineto{\pgfqpoint{2.439887in}{1.202342in}}%
\pgfpathlineto{\pgfqpoint{2.439065in}{1.211016in}}%
\pgfpathlineto{\pgfqpoint{2.439738in}{1.214050in}}%
\pgfpathlineto{\pgfqpoint{2.445630in}{1.214029in}}%
\pgfpathlineto{\pgfqpoint{2.445563in}{1.221325in}}%
\pgfpathlineto{\pgfqpoint{2.456449in}{1.222465in}}%
\pgfpathlineto{\pgfqpoint{2.465782in}{1.225372in}}%
\pgfpathlineto{\pgfqpoint{2.467867in}{1.223907in}}%
\pgfpathlineto{\pgfqpoint{2.478476in}{1.225076in}}%
\pgfpathlineto{\pgfqpoint{2.476467in}{1.219730in}}%
\pgfpathlineto{\pgfqpoint{2.485346in}{1.213775in}}%
\pgfpathlineto{\pgfqpoint{2.488733in}{1.208153in}}%
\pgfpathlineto{\pgfqpoint{2.487686in}{1.206566in}}%
\pgfpathlineto{\pgfqpoint{2.494442in}{1.196721in}}%
\pgfpathlineto{\pgfqpoint{2.489753in}{1.191584in}}%
\pgfpathlineto{\pgfqpoint{2.484306in}{1.188484in}}%
\pgfpathlineto{\pgfqpoint{2.478749in}{1.189995in}}%
\pgfpathlineto{\pgfqpoint{2.469733in}{1.178659in}}%
\pgfpathlineto{\pgfqpoint{2.463999in}{1.180859in}}%
\pgfpathlineto{\pgfqpoint{2.460301in}{1.177194in}}%
\pgfpathlineto{\pgfqpoint{2.461626in}{1.162571in}}%
\pgfpathlineto{\pgfqpoint{2.463450in}{1.158871in}}%
\pgfpathclose%
\pgfusepath{fill}%
\end{pgfscope}%
\begin{pgfscope}%
\pgfpathrectangle{\pgfqpoint{0.100000in}{0.100000in}}{\pgfqpoint{3.007045in}{1.925000in}}%
\pgfusepath{clip}%
\pgfsetbuttcap%
\pgfsetmiterjoin%
\definecolor{currentfill}{rgb}{0.510588,0.732303,0.858839}%
\pgfsetfillcolor{currentfill}%
\pgfsetlinewidth{0.000000pt}%
\definecolor{currentstroke}{rgb}{0.000000,0.000000,0.000000}%
\pgfsetstrokecolor{currentstroke}%
\pgfsetstrokeopacity{0.000000}%
\pgfsetdash{}{0pt}%
\pgfpathmoveto{\pgfqpoint{1.898889in}{1.500432in}}%
\pgfpathlineto{\pgfqpoint{1.907419in}{1.500598in}}%
\pgfpathlineto{\pgfqpoint{1.906167in}{1.539095in}}%
\pgfpathlineto{\pgfqpoint{1.905032in}{1.544903in}}%
\pgfpathlineto{\pgfqpoint{1.933694in}{1.545701in}}%
\pgfpathlineto{\pgfqpoint{1.933521in}{1.551363in}}%
\pgfpathlineto{\pgfqpoint{1.962214in}{1.552559in}}%
\pgfpathlineto{\pgfqpoint{1.964639in}{1.500815in}}%
\pgfpathlineto{\pgfqpoint{1.970355in}{1.501053in}}%
\pgfpathlineto{\pgfqpoint{1.970609in}{1.495307in}}%
\pgfpathlineto{\pgfqpoint{1.993541in}{1.496568in}}%
\pgfpathlineto{\pgfqpoint{1.994330in}{1.484956in}}%
\pgfpathlineto{\pgfqpoint{1.994624in}{1.478781in}}%
\pgfpathlineto{\pgfqpoint{1.966485in}{1.477622in}}%
\pgfpathlineto{\pgfqpoint{1.963568in}{1.471696in}}%
\pgfpathlineto{\pgfqpoint{1.947712in}{1.470634in}}%
\pgfpathlineto{\pgfqpoint{1.947232in}{1.465292in}}%
\pgfpathlineto{\pgfqpoint{1.942334in}{1.465087in}}%
\pgfpathlineto{\pgfqpoint{1.949409in}{1.456290in}}%
\pgfpathlineto{\pgfqpoint{1.928151in}{1.455464in}}%
\pgfpathlineto{\pgfqpoint{1.911530in}{1.454907in}}%
\pgfpathlineto{\pgfqpoint{1.910791in}{1.477769in}}%
\pgfpathlineto{\pgfqpoint{1.921271in}{1.478116in}}%
\pgfpathlineto{\pgfqpoint{1.917765in}{1.487298in}}%
\pgfpathlineto{\pgfqpoint{1.911427in}{1.491667in}}%
\pgfpathlineto{\pgfqpoint{1.903859in}{1.493796in}}%
\pgfpathlineto{\pgfqpoint{1.898889in}{1.500432in}}%
\pgfpathclose%
\pgfusepath{fill}%
\end{pgfscope}%
\begin{pgfscope}%
\pgfpathrectangle{\pgfqpoint{0.100000in}{0.100000in}}{\pgfqpoint{3.007045in}{1.925000in}}%
\pgfusepath{clip}%
\pgfsetbuttcap%
\pgfsetmiterjoin%
\definecolor{currentfill}{rgb}{0.341423,0.628958,0.808704}%
\pgfsetfillcolor{currentfill}%
\pgfsetlinewidth{0.000000pt}%
\definecolor{currentstroke}{rgb}{0.000000,0.000000,0.000000}%
\pgfsetstrokecolor{currentstroke}%
\pgfsetstrokeopacity{0.000000}%
\pgfsetdash{}{0pt}%
\pgfpathmoveto{\pgfqpoint{1.632363in}{0.431928in}}%
\pgfpathlineto{\pgfqpoint{1.620569in}{0.414900in}}%
\pgfpathlineto{\pgfqpoint{1.605884in}{0.418493in}}%
\pgfpathlineto{\pgfqpoint{1.606432in}{0.418922in}}%
\pgfpathlineto{\pgfqpoint{1.590934in}{0.451560in}}%
\pgfpathlineto{\pgfqpoint{1.595954in}{0.453216in}}%
\pgfpathlineto{\pgfqpoint{1.585508in}{0.466431in}}%
\pgfpathlineto{\pgfqpoint{1.613261in}{0.488228in}}%
\pgfpathlineto{\pgfqpoint{1.619767in}{0.480550in}}%
\pgfpathlineto{\pgfqpoint{1.611272in}{0.473936in}}%
\pgfpathlineto{\pgfqpoint{1.621451in}{0.460671in}}%
\pgfpathlineto{\pgfqpoint{1.609200in}{0.451305in}}%
\pgfpathlineto{\pgfqpoint{1.613769in}{0.442713in}}%
\pgfpathlineto{\pgfqpoint{1.621822in}{0.440642in}}%
\pgfpathlineto{\pgfqpoint{1.622101in}{0.436398in}}%
\pgfpathlineto{\pgfqpoint{1.632363in}{0.431928in}}%
\pgfpathclose%
\pgfusepath{fill}%
\end{pgfscope}%
\begin{pgfscope}%
\pgfpathrectangle{\pgfqpoint{0.100000in}{0.100000in}}{\pgfqpoint{3.007045in}{1.925000in}}%
\pgfusepath{clip}%
\pgfsetbuttcap%
\pgfsetmiterjoin%
\definecolor{currentfill}{rgb}{0.366644,0.646182,0.818547}%
\pgfsetfillcolor{currentfill}%
\pgfsetlinewidth{0.000000pt}%
\definecolor{currentstroke}{rgb}{0.000000,0.000000,0.000000}%
\pgfsetstrokecolor{currentstroke}%
\pgfsetstrokeopacity{0.000000}%
\pgfsetdash{}{0pt}%
\pgfpathmoveto{\pgfqpoint{2.516660in}{0.795332in}}%
\pgfpathlineto{\pgfqpoint{2.505569in}{0.784566in}}%
\pgfpathlineto{\pgfqpoint{2.500186in}{0.786765in}}%
\pgfpathlineto{\pgfqpoint{2.488225in}{0.782556in}}%
\pgfpathlineto{\pgfqpoint{2.485475in}{0.775665in}}%
\pgfpathlineto{\pgfqpoint{2.482589in}{0.774200in}}%
\pgfpathlineto{\pgfqpoint{2.475814in}{0.775674in}}%
\pgfpathlineto{\pgfqpoint{2.470222in}{0.769671in}}%
\pgfpathlineto{\pgfqpoint{2.464508in}{0.773648in}}%
\pgfpathlineto{\pgfqpoint{2.464577in}{0.779288in}}%
\pgfpathlineto{\pgfqpoint{2.461072in}{0.785247in}}%
\pgfpathlineto{\pgfqpoint{2.454423in}{0.792007in}}%
\pgfpathlineto{\pgfqpoint{2.449695in}{0.794403in}}%
\pgfpathlineto{\pgfqpoint{2.448640in}{0.798533in}}%
\pgfpathlineto{\pgfqpoint{2.441136in}{0.811210in}}%
\pgfpathlineto{\pgfqpoint{2.442122in}{0.815272in}}%
\pgfpathlineto{\pgfqpoint{2.446193in}{0.816402in}}%
\pgfpathlineto{\pgfqpoint{2.449159in}{0.822767in}}%
\pgfpathlineto{\pgfqpoint{2.452424in}{0.825463in}}%
\pgfpathlineto{\pgfqpoint{2.460350in}{0.826839in}}%
\pgfpathlineto{\pgfqpoint{2.468621in}{0.832178in}}%
\pgfpathlineto{\pgfqpoint{2.473536in}{0.831913in}}%
\pgfpathlineto{\pgfqpoint{2.479069in}{0.826299in}}%
\pgfpathlineto{\pgfqpoint{2.480732in}{0.829446in}}%
\pgfpathlineto{\pgfqpoint{2.475134in}{0.834677in}}%
\pgfpathlineto{\pgfqpoint{2.477899in}{0.842105in}}%
\pgfpathlineto{\pgfqpoint{2.474703in}{0.848311in}}%
\pgfpathlineto{\pgfqpoint{2.482424in}{0.851583in}}%
\pgfpathlineto{\pgfqpoint{2.488445in}{0.853368in}}%
\pgfpathlineto{\pgfqpoint{2.492621in}{0.846691in}}%
\pgfpathlineto{\pgfqpoint{2.503173in}{0.844767in}}%
\pgfpathlineto{\pgfqpoint{2.506929in}{0.849528in}}%
\pgfpathlineto{\pgfqpoint{2.516115in}{0.841918in}}%
\pgfpathlineto{\pgfqpoint{2.530038in}{0.837997in}}%
\pgfpathlineto{\pgfqpoint{2.521533in}{0.825837in}}%
\pgfpathlineto{\pgfqpoint{2.503376in}{0.803575in}}%
\pgfpathlineto{\pgfqpoint{2.507037in}{0.798645in}}%
\pgfpathlineto{\pgfqpoint{2.516660in}{0.795332in}}%
\pgfpathclose%
\pgfusepath{fill}%
\end{pgfscope}%
\begin{pgfscope}%
\pgfpathrectangle{\pgfqpoint{0.100000in}{0.100000in}}{\pgfqpoint{3.007045in}{1.925000in}}%
\pgfusepath{clip}%
\pgfsetbuttcap%
\pgfsetmiterjoin%
\definecolor{currentfill}{rgb}{0.417086,0.680631,0.838231}%
\pgfsetfillcolor{currentfill}%
\pgfsetlinewidth{0.000000pt}%
\definecolor{currentstroke}{rgb}{0.000000,0.000000,0.000000}%
\pgfsetstrokecolor{currentstroke}%
\pgfsetstrokeopacity{0.000000}%
\pgfsetdash{}{0pt}%
\pgfpathmoveto{\pgfqpoint{2.719169in}{0.923665in}}%
\pgfpathlineto{\pgfqpoint{2.714846in}{0.926915in}}%
\pgfpathlineto{\pgfqpoint{2.706932in}{0.940219in}}%
\pgfpathlineto{\pgfqpoint{2.705959in}{0.953701in}}%
\pgfpathlineto{\pgfqpoint{2.699683in}{0.959124in}}%
\pgfpathlineto{\pgfqpoint{2.696911in}{0.966874in}}%
\pgfpathlineto{\pgfqpoint{2.693069in}{0.968298in}}%
\pgfpathlineto{\pgfqpoint{2.688736in}{0.994006in}}%
\pgfpathlineto{\pgfqpoint{2.708269in}{1.015278in}}%
\pgfpathlineto{\pgfqpoint{2.710520in}{1.014800in}}%
\pgfpathlineto{\pgfqpoint{2.720462in}{1.011491in}}%
\pgfpathlineto{\pgfqpoint{2.722156in}{1.004752in}}%
\pgfpathlineto{\pgfqpoint{2.725670in}{0.999941in}}%
\pgfpathlineto{\pgfqpoint{2.723536in}{0.997192in}}%
\pgfpathlineto{\pgfqpoint{2.724069in}{0.990555in}}%
\pgfpathlineto{\pgfqpoint{2.737785in}{0.985361in}}%
\pgfpathlineto{\pgfqpoint{2.744917in}{0.981075in}}%
\pgfpathlineto{\pgfqpoint{2.755915in}{0.984607in}}%
\pgfpathlineto{\pgfqpoint{2.755737in}{0.991248in}}%
\pgfpathlineto{\pgfqpoint{2.762672in}{0.991320in}}%
\pgfpathlineto{\pgfqpoint{2.764291in}{0.988344in}}%
\pgfpathlineto{\pgfqpoint{2.761049in}{0.976325in}}%
\pgfpathlineto{\pgfqpoint{2.757169in}{0.971840in}}%
\pgfpathlineto{\pgfqpoint{2.760232in}{0.968674in}}%
\pgfpathlineto{\pgfqpoint{2.769233in}{0.972061in}}%
\pgfpathlineto{\pgfqpoint{2.773670in}{0.970838in}}%
\pgfpathlineto{\pgfqpoint{2.777250in}{0.973729in}}%
\pgfpathlineto{\pgfqpoint{2.783649in}{0.973537in}}%
\pgfpathlineto{\pgfqpoint{2.771520in}{0.950478in}}%
\pgfpathlineto{\pgfqpoint{2.767312in}{0.947670in}}%
\pgfpathlineto{\pgfqpoint{2.760927in}{0.948913in}}%
\pgfpathlineto{\pgfqpoint{2.754301in}{0.947904in}}%
\pgfpathlineto{\pgfqpoint{2.736174in}{0.939738in}}%
\pgfpathlineto{\pgfqpoint{2.719169in}{0.923665in}}%
\pgfpathclose%
\pgfusepath{fill}%
\end{pgfscope}%
\begin{pgfscope}%
\pgfpathrectangle{\pgfqpoint{0.100000in}{0.100000in}}{\pgfqpoint{3.007045in}{1.925000in}}%
\pgfusepath{clip}%
\pgfsetbuttcap%
\pgfsetmiterjoin%
\definecolor{currentfill}{rgb}{0.280892,0.587620,0.785083}%
\pgfsetfillcolor{currentfill}%
\pgfsetlinewidth{0.000000pt}%
\definecolor{currentstroke}{rgb}{0.000000,0.000000,0.000000}%
\pgfsetstrokecolor{currentstroke}%
\pgfsetstrokeopacity{0.000000}%
\pgfsetdash{}{0pt}%
\pgfpathmoveto{\pgfqpoint{1.481225in}{0.618416in}}%
\pgfpathlineto{\pgfqpoint{1.479753in}{0.593339in}}%
\pgfpathlineto{\pgfqpoint{1.431514in}{0.595996in}}%
\pgfpathlineto{\pgfqpoint{1.433065in}{0.620988in}}%
\pgfpathlineto{\pgfqpoint{1.415290in}{0.621980in}}%
\pgfpathlineto{\pgfqpoint{1.417736in}{0.651743in}}%
\pgfpathlineto{\pgfqpoint{1.420501in}{0.688714in}}%
\pgfpathlineto{\pgfqpoint{1.425000in}{0.688410in}}%
\pgfpathlineto{\pgfqpoint{1.454083in}{0.686391in}}%
\pgfpathlineto{\pgfqpoint{1.478210in}{0.684741in}}%
\pgfpathlineto{\pgfqpoint{1.476325in}{0.651539in}}%
\pgfpathlineto{\pgfqpoint{1.483290in}{0.651066in}}%
\pgfpathlineto{\pgfqpoint{1.481225in}{0.618416in}}%
\pgfpathclose%
\pgfusepath{fill}%
\end{pgfscope}%
\begin{pgfscope}%
\pgfpathrectangle{\pgfqpoint{0.100000in}{0.100000in}}{\pgfqpoint{3.007045in}{1.925000in}}%
\pgfusepath{clip}%
\pgfsetbuttcap%
\pgfsetmiterjoin%
\definecolor{currentfill}{rgb}{0.371688,0.649627,0.820515}%
\pgfsetfillcolor{currentfill}%
\pgfsetlinewidth{0.000000pt}%
\definecolor{currentstroke}{rgb}{0.000000,0.000000,0.000000}%
\pgfsetstrokecolor{currentstroke}%
\pgfsetstrokeopacity{0.000000}%
\pgfsetdash{}{0pt}%
\pgfpathmoveto{\pgfqpoint{2.382703in}{0.849302in}}%
\pgfpathlineto{\pgfqpoint{2.380206in}{0.847045in}}%
\pgfpathlineto{\pgfqpoint{2.372820in}{0.852600in}}%
\pgfpathlineto{\pgfqpoint{2.375538in}{0.858536in}}%
\pgfpathlineto{\pgfqpoint{2.371582in}{0.867326in}}%
\pgfpathlineto{\pgfqpoint{2.365118in}{0.869291in}}%
\pgfpathlineto{\pgfqpoint{2.359050in}{0.873951in}}%
\pgfpathlineto{\pgfqpoint{2.358828in}{0.878027in}}%
\pgfpathlineto{\pgfqpoint{2.360142in}{0.881259in}}%
\pgfpathlineto{\pgfqpoint{2.366727in}{0.881594in}}%
\pgfpathlineto{\pgfqpoint{2.371189in}{0.888722in}}%
\pgfpathlineto{\pgfqpoint{2.375790in}{0.888043in}}%
\pgfpathlineto{\pgfqpoint{2.377707in}{0.892532in}}%
\pgfpathlineto{\pgfqpoint{2.385967in}{0.896018in}}%
\pgfpathlineto{\pgfqpoint{2.389462in}{0.896298in}}%
\pgfpathlineto{\pgfqpoint{2.391995in}{0.891918in}}%
\pgfpathlineto{\pgfqpoint{2.402990in}{0.891059in}}%
\pgfpathlineto{\pgfqpoint{2.403995in}{0.889257in}}%
\pgfpathlineto{\pgfqpoint{2.402644in}{0.888304in}}%
\pgfpathlineto{\pgfqpoint{2.399169in}{0.874931in}}%
\pgfpathlineto{\pgfqpoint{2.403763in}{0.872065in}}%
\pgfpathlineto{\pgfqpoint{2.404051in}{0.865135in}}%
\pgfpathlineto{\pgfqpoint{2.406992in}{0.858796in}}%
\pgfpathlineto{\pgfqpoint{2.404649in}{0.856728in}}%
\pgfpathlineto{\pgfqpoint{2.407803in}{0.851715in}}%
\pgfpathlineto{\pgfqpoint{2.406734in}{0.847180in}}%
\pgfpathlineto{\pgfqpoint{2.399267in}{0.840607in}}%
\pgfpathlineto{\pgfqpoint{2.397316in}{0.844754in}}%
\pgfpathlineto{\pgfqpoint{2.386057in}{0.845572in}}%
\pgfpathlineto{\pgfqpoint{2.382703in}{0.849302in}}%
\pgfpathclose%
\pgfusepath{fill}%
\end{pgfscope}%
\begin{pgfscope}%
\pgfpathrectangle{\pgfqpoint{0.100000in}{0.100000in}}{\pgfqpoint{3.007045in}{1.925000in}}%
\pgfusepath{clip}%
\pgfsetbuttcap%
\pgfsetmiterjoin%
\definecolor{currentfill}{rgb}{0.510588,0.732303,0.858839}%
\pgfsetfillcolor{currentfill}%
\pgfsetlinewidth{0.000000pt}%
\definecolor{currentstroke}{rgb}{0.000000,0.000000,0.000000}%
\pgfsetstrokecolor{currentstroke}%
\pgfsetstrokeopacity{0.000000}%
\pgfsetdash{}{0pt}%
\pgfpathmoveto{\pgfqpoint{1.746105in}{1.095854in}}%
\pgfpathlineto{\pgfqpoint{1.723295in}{1.096039in}}%
\pgfpathlineto{\pgfqpoint{1.723731in}{1.116097in}}%
\pgfpathlineto{\pgfqpoint{1.702839in}{1.116383in}}%
\pgfpathlineto{\pgfqpoint{1.700977in}{1.122113in}}%
\pgfpathlineto{\pgfqpoint{1.695279in}{1.122232in}}%
\pgfpathlineto{\pgfqpoint{1.695539in}{1.136550in}}%
\pgfpathlineto{\pgfqpoint{1.701186in}{1.136436in}}%
\pgfpathlineto{\pgfqpoint{1.701412in}{1.144959in}}%
\pgfpathlineto{\pgfqpoint{1.708830in}{1.147175in}}%
\pgfpathlineto{\pgfqpoint{1.712256in}{1.145853in}}%
\pgfpathlineto{\pgfqpoint{1.719484in}{1.147603in}}%
\pgfpathlineto{\pgfqpoint{1.719727in}{1.170628in}}%
\pgfpathlineto{\pgfqpoint{1.732247in}{1.170499in}}%
\pgfpathlineto{\pgfqpoint{1.732315in}{1.176225in}}%
\pgfpathlineto{\pgfqpoint{1.743716in}{1.176139in}}%
\pgfpathlineto{\pgfqpoint{1.743298in}{1.160738in}}%
\pgfpathlineto{\pgfqpoint{1.763181in}{1.160646in}}%
\pgfpathlineto{\pgfqpoint{1.763130in}{1.141541in}}%
\pgfpathlineto{\pgfqpoint{1.764032in}{1.131806in}}%
\pgfpathlineto{\pgfqpoint{1.769446in}{1.131864in}}%
\pgfpathlineto{\pgfqpoint{1.769426in}{1.115817in}}%
\pgfpathlineto{\pgfqpoint{1.746586in}{1.115888in}}%
\pgfpathlineto{\pgfqpoint{1.746105in}{1.095854in}}%
\pgfpathclose%
\pgfusepath{fill}%
\end{pgfscope}%
\begin{pgfscope}%
\pgfpathrectangle{\pgfqpoint{0.100000in}{0.100000in}}{\pgfqpoint{3.007045in}{1.925000in}}%
\pgfusepath{clip}%
\pgfsetbuttcap%
\pgfsetmiterjoin%
\definecolor{currentfill}{rgb}{0.301069,0.601399,0.792957}%
\pgfsetfillcolor{currentfill}%
\pgfsetlinewidth{0.000000pt}%
\definecolor{currentstroke}{rgb}{0.000000,0.000000,0.000000}%
\pgfsetstrokecolor{currentstroke}%
\pgfsetstrokeopacity{0.000000}%
\pgfsetdash{}{0pt}%
\pgfpathmoveto{\pgfqpoint{1.704139in}{0.513794in}}%
\pgfpathlineto{\pgfqpoint{1.696363in}{0.505939in}}%
\pgfpathlineto{\pgfqpoint{1.693878in}{0.498881in}}%
\pgfpathlineto{\pgfqpoint{1.676720in}{0.488443in}}%
\pgfpathlineto{\pgfqpoint{1.681469in}{0.494214in}}%
\pgfpathlineto{\pgfqpoint{1.668257in}{0.510605in}}%
\pgfpathlineto{\pgfqpoint{1.663618in}{0.514403in}}%
\pgfpathlineto{\pgfqpoint{1.648009in}{0.514439in}}%
\pgfpathlineto{\pgfqpoint{1.638148in}{0.525132in}}%
\pgfpathlineto{\pgfqpoint{1.655528in}{0.542512in}}%
\pgfpathlineto{\pgfqpoint{1.665918in}{0.546597in}}%
\pgfpathlineto{\pgfqpoint{1.668934in}{0.549504in}}%
\pgfpathlineto{\pgfqpoint{1.675326in}{0.548748in}}%
\pgfpathlineto{\pgfqpoint{1.683561in}{0.532352in}}%
\pgfpathlineto{\pgfqpoint{1.690540in}{0.526743in}}%
\pgfpathlineto{\pgfqpoint{1.695319in}{0.526005in}}%
\pgfpathlineto{\pgfqpoint{1.699409in}{0.516089in}}%
\pgfpathlineto{\pgfqpoint{1.704139in}{0.513794in}}%
\pgfpathclose%
\pgfusepath{fill}%
\end{pgfscope}%
\begin{pgfscope}%
\pgfpathrectangle{\pgfqpoint{0.100000in}{0.100000in}}{\pgfqpoint{3.007045in}{1.925000in}}%
\pgfusepath{clip}%
\pgfsetbuttcap%
\pgfsetmiterjoin%
\definecolor{currentfill}{rgb}{0.429020,0.687520,0.841246}%
\pgfsetfillcolor{currentfill}%
\pgfsetlinewidth{0.000000pt}%
\definecolor{currentstroke}{rgb}{0.000000,0.000000,0.000000}%
\pgfsetstrokecolor{currentstroke}%
\pgfsetstrokeopacity{0.000000}%
\pgfsetdash{}{0pt}%
\pgfpathmoveto{\pgfqpoint{0.646486in}{1.339327in}}%
\pgfpathlineto{\pgfqpoint{0.594471in}{1.352638in}}%
\pgfpathlineto{\pgfqpoint{0.558741in}{1.362309in}}%
\pgfpathlineto{\pgfqpoint{0.567831in}{1.395284in}}%
\pgfpathlineto{\pgfqpoint{0.569070in}{1.394980in}}%
\pgfpathlineto{\pgfqpoint{0.576936in}{1.423028in}}%
\pgfpathlineto{\pgfqpoint{0.575970in}{1.423296in}}%
\pgfpathlineto{\pgfqpoint{0.584193in}{1.452083in}}%
\pgfpathlineto{\pgfqpoint{0.594261in}{1.488948in}}%
\pgfpathlineto{\pgfqpoint{0.623923in}{1.480760in}}%
\pgfpathlineto{\pgfqpoint{0.647669in}{1.474848in}}%
\pgfpathlineto{\pgfqpoint{0.703346in}{1.460924in}}%
\pgfpathlineto{\pgfqpoint{0.703710in}{1.460774in}}%
\pgfpathlineto{\pgfqpoint{0.687900in}{1.396831in}}%
\pgfpathlineto{\pgfqpoint{0.682240in}{1.373981in}}%
\pgfpathlineto{\pgfqpoint{0.671162in}{1.376694in}}%
\pgfpathlineto{\pgfqpoint{0.646486in}{1.339327in}}%
\pgfpathclose%
\pgfusepath{fill}%
\end{pgfscope}%
\begin{pgfscope}%
\pgfpathrectangle{\pgfqpoint{0.100000in}{0.100000in}}{\pgfqpoint{3.007045in}{1.925000in}}%
\pgfusepath{clip}%
\pgfsetbuttcap%
\pgfsetmiterjoin%
\definecolor{currentfill}{rgb}{0.187266,0.500992,0.739608}%
\pgfsetfillcolor{currentfill}%
\pgfsetlinewidth{0.000000pt}%
\definecolor{currentstroke}{rgb}{0.000000,0.000000,0.000000}%
\pgfsetstrokecolor{currentstroke}%
\pgfsetstrokeopacity{0.000000}%
\pgfsetdash{}{0pt}%
\pgfpathmoveto{\pgfqpoint{1.674656in}{1.382036in}}%
\pgfpathlineto{\pgfqpoint{1.673964in}{1.354981in}}%
\pgfpathlineto{\pgfqpoint{1.656900in}{1.355436in}}%
\pgfpathlineto{\pgfqpoint{1.657059in}{1.361181in}}%
\pgfpathlineto{\pgfqpoint{1.611668in}{1.362462in}}%
\pgfpathlineto{\pgfqpoint{1.612271in}{1.391873in}}%
\pgfpathlineto{\pgfqpoint{1.618866in}{1.388633in}}%
\pgfpathlineto{\pgfqpoint{1.621472in}{1.391045in}}%
\pgfpathlineto{\pgfqpoint{1.622963in}{1.397529in}}%
\pgfpathlineto{\pgfqpoint{1.622340in}{1.404124in}}%
\pgfpathlineto{\pgfqpoint{1.624011in}{1.410336in}}%
\pgfpathlineto{\pgfqpoint{1.656789in}{1.409421in}}%
\pgfpathlineto{\pgfqpoint{1.668294in}{1.409172in}}%
\pgfpathlineto{\pgfqpoint{1.667711in}{1.384730in}}%
\pgfpathlineto{\pgfqpoint{1.674656in}{1.382036in}}%
\pgfpathclose%
\pgfusepath{fill}%
\end{pgfscope}%
\begin{pgfscope}%
\pgfpathrectangle{\pgfqpoint{0.100000in}{0.100000in}}{\pgfqpoint{3.007045in}{1.925000in}}%
\pgfusepath{clip}%
\pgfsetbuttcap%
\pgfsetmiterjoin%
\definecolor{currentfill}{rgb}{0.265759,0.577286,0.779177}%
\pgfsetfillcolor{currentfill}%
\pgfsetlinewidth{0.000000pt}%
\definecolor{currentstroke}{rgb}{0.000000,0.000000,0.000000}%
\pgfsetstrokecolor{currentstroke}%
\pgfsetstrokeopacity{0.000000}%
\pgfsetdash{}{0pt}%
\pgfpathmoveto{\pgfqpoint{2.675874in}{0.992417in}}%
\pgfpathlineto{\pgfqpoint{2.688736in}{0.994006in}}%
\pgfpathlineto{\pgfqpoint{2.693069in}{0.968298in}}%
\pgfpathlineto{\pgfqpoint{2.696911in}{0.966874in}}%
\pgfpathlineto{\pgfqpoint{2.699683in}{0.959124in}}%
\pgfpathlineto{\pgfqpoint{2.705959in}{0.953701in}}%
\pgfpathlineto{\pgfqpoint{2.706932in}{0.940219in}}%
\pgfpathlineto{\pgfqpoint{2.683843in}{0.935913in}}%
\pgfpathlineto{\pgfqpoint{2.678361in}{0.924174in}}%
\pgfpathlineto{\pgfqpoint{2.668672in}{0.934954in}}%
\pgfpathlineto{\pgfqpoint{2.661874in}{0.940850in}}%
\pgfpathlineto{\pgfqpoint{2.652361in}{0.948046in}}%
\pgfpathlineto{\pgfqpoint{2.650566in}{0.955387in}}%
\pgfpathlineto{\pgfqpoint{2.650737in}{0.964955in}}%
\pgfpathlineto{\pgfqpoint{2.653930in}{0.970205in}}%
\pgfpathlineto{\pgfqpoint{2.657195in}{0.967322in}}%
\pgfpathlineto{\pgfqpoint{2.661577in}{0.967658in}}%
\pgfpathlineto{\pgfqpoint{2.667375in}{0.972671in}}%
\pgfpathlineto{\pgfqpoint{2.674050in}{0.976293in}}%
\pgfpathlineto{\pgfqpoint{2.673153in}{0.980638in}}%
\pgfpathlineto{\pgfqpoint{2.675874in}{0.992417in}}%
\pgfpathclose%
\pgfusepath{fill}%
\end{pgfscope}%
\begin{pgfscope}%
\pgfpathrectangle{\pgfqpoint{0.100000in}{0.100000in}}{\pgfqpoint{3.007045in}{1.925000in}}%
\pgfusepath{clip}%
\pgfsetbuttcap%
\pgfsetmiterjoin%
\definecolor{currentfill}{rgb}{0.351511,0.635848,0.812641}%
\pgfsetfillcolor{currentfill}%
\pgfsetlinewidth{0.000000pt}%
\definecolor{currentstroke}{rgb}{0.000000,0.000000,0.000000}%
\pgfsetstrokecolor{currentstroke}%
\pgfsetstrokeopacity{0.000000}%
\pgfsetdash{}{0pt}%
\pgfpathmoveto{\pgfqpoint{2.339044in}{0.640515in}}%
\pgfpathlineto{\pgfqpoint{2.341073in}{0.635732in}}%
\pgfpathlineto{\pgfqpoint{2.282576in}{0.629115in}}%
\pgfpathlineto{\pgfqpoint{2.273936in}{0.628314in}}%
\pgfpathlineto{\pgfqpoint{2.270062in}{0.663726in}}%
\pgfpathlineto{\pgfqpoint{2.272250in}{0.669798in}}%
\pgfpathlineto{\pgfqpoint{2.294624in}{0.672055in}}%
\pgfpathlineto{\pgfqpoint{2.313278in}{0.674089in}}%
\pgfpathlineto{\pgfqpoint{2.312691in}{0.679850in}}%
\pgfpathlineto{\pgfqpoint{2.329091in}{0.680831in}}%
\pgfpathlineto{\pgfqpoint{2.333467in}{0.676210in}}%
\pgfpathlineto{\pgfqpoint{2.334926in}{0.671284in}}%
\pgfpathlineto{\pgfqpoint{2.332928in}{0.653635in}}%
\pgfpathlineto{\pgfqpoint{2.334338in}{0.646016in}}%
\pgfpathlineto{\pgfqpoint{2.339044in}{0.640515in}}%
\pgfpathclose%
\pgfusepath{fill}%
\end{pgfscope}%
\begin{pgfscope}%
\pgfpathrectangle{\pgfqpoint{0.100000in}{0.100000in}}{\pgfqpoint{3.007045in}{1.925000in}}%
\pgfusepath{clip}%
\pgfsetbuttcap%
\pgfsetmiterjoin%
\definecolor{currentfill}{rgb}{0.454118,0.701300,0.846659}%
\pgfsetfillcolor{currentfill}%
\pgfsetlinewidth{0.000000pt}%
\definecolor{currentstroke}{rgb}{0.000000,0.000000,0.000000}%
\pgfsetstrokecolor{currentstroke}%
\pgfsetstrokeopacity{0.000000}%
\pgfsetdash{}{0pt}%
\pgfpathmoveto{\pgfqpoint{0.746029in}{1.809361in}}%
\pgfpathlineto{\pgfqpoint{0.701521in}{1.821189in}}%
\pgfpathlineto{\pgfqpoint{0.713636in}{1.864996in}}%
\pgfpathlineto{\pgfqpoint{0.719784in}{1.864006in}}%
\pgfpathlineto{\pgfqpoint{0.730840in}{1.859613in}}%
\pgfpathlineto{\pgfqpoint{0.735184in}{1.852084in}}%
\pgfpathlineto{\pgfqpoint{0.744909in}{1.856652in}}%
\pgfpathlineto{\pgfqpoint{0.748542in}{1.851631in}}%
\pgfpathlineto{\pgfqpoint{0.749934in}{1.845470in}}%
\pgfpathlineto{\pgfqpoint{0.759267in}{1.846349in}}%
\pgfpathlineto{\pgfqpoint{0.762060in}{1.843258in}}%
\pgfpathlineto{\pgfqpoint{0.768647in}{1.844176in}}%
\pgfpathlineto{\pgfqpoint{0.773351in}{1.839140in}}%
\pgfpathlineto{\pgfqpoint{0.777516in}{1.855475in}}%
\pgfpathlineto{\pgfqpoint{0.780353in}{1.860745in}}%
\pgfpathlineto{\pgfqpoint{0.789000in}{1.894475in}}%
\pgfpathlineto{\pgfqpoint{0.783449in}{1.895856in}}%
\pgfpathlineto{\pgfqpoint{0.787685in}{1.900755in}}%
\pgfpathlineto{\pgfqpoint{0.789229in}{1.906864in}}%
\pgfpathlineto{\pgfqpoint{0.797625in}{1.916024in}}%
\pgfpathlineto{\pgfqpoint{0.814525in}{1.911735in}}%
\pgfpathlineto{\pgfqpoint{0.812070in}{1.901824in}}%
\pgfpathlineto{\pgfqpoint{0.822653in}{1.899281in}}%
\pgfpathlineto{\pgfqpoint{0.817161in}{1.876758in}}%
\pgfpathlineto{\pgfqpoint{0.848850in}{1.869247in}}%
\pgfpathlineto{\pgfqpoint{0.844425in}{1.850601in}}%
\pgfpathlineto{\pgfqpoint{0.840821in}{1.835141in}}%
\pgfpathlineto{\pgfqpoint{0.847161in}{1.818241in}}%
\pgfpathlineto{\pgfqpoint{0.850754in}{1.813504in}}%
\pgfpathlineto{\pgfqpoint{0.850775in}{1.806624in}}%
\pgfpathlineto{\pgfqpoint{0.847085in}{1.804318in}}%
\pgfpathlineto{\pgfqpoint{0.851627in}{1.798562in}}%
\pgfpathlineto{\pgfqpoint{0.845810in}{1.797316in}}%
\pgfpathlineto{\pgfqpoint{0.853094in}{1.789354in}}%
\pgfpathlineto{\pgfqpoint{0.853297in}{1.785851in}}%
\pgfpathlineto{\pgfqpoint{0.861882in}{1.780983in}}%
\pgfpathlineto{\pgfqpoint{0.869191in}{1.761494in}}%
\pgfpathlineto{\pgfqpoint{0.873084in}{1.756371in}}%
\pgfpathlineto{\pgfqpoint{0.844253in}{1.762934in}}%
\pgfpathlineto{\pgfqpoint{0.812473in}{1.770492in}}%
\pgfpathlineto{\pgfqpoint{0.813833in}{1.776182in}}%
\pgfpathlineto{\pgfqpoint{0.808364in}{1.778465in}}%
\pgfpathlineto{\pgfqpoint{0.794601in}{1.781859in}}%
\pgfpathlineto{\pgfqpoint{0.784213in}{1.790565in}}%
\pgfpathlineto{\pgfqpoint{0.786311in}{1.799097in}}%
\pgfpathlineto{\pgfqpoint{0.746029in}{1.809361in}}%
\pgfpathclose%
\pgfusepath{fill}%
\end{pgfscope}%
\begin{pgfscope}%
\pgfpathrectangle{\pgfqpoint{0.100000in}{0.100000in}}{\pgfqpoint{3.007045in}{1.925000in}}%
\pgfusepath{clip}%
\pgfsetbuttcap%
\pgfsetmiterjoin%
\definecolor{currentfill}{rgb}{0.306113,0.604844,0.794925}%
\pgfsetfillcolor{currentfill}%
\pgfsetlinewidth{0.000000pt}%
\definecolor{currentstroke}{rgb}{0.000000,0.000000,0.000000}%
\pgfsetstrokecolor{currentstroke}%
\pgfsetstrokeopacity{0.000000}%
\pgfsetdash{}{0pt}%
\pgfpathmoveto{\pgfqpoint{1.839662in}{0.517608in}}%
\pgfpathlineto{\pgfqpoint{1.824849in}{0.515348in}}%
\pgfpathlineto{\pgfqpoint{1.809818in}{0.508665in}}%
\pgfpathlineto{\pgfqpoint{1.809447in}{0.530410in}}%
\pgfpathlineto{\pgfqpoint{1.804575in}{0.530516in}}%
\pgfpathlineto{\pgfqpoint{1.804317in}{0.545414in}}%
\pgfpathlineto{\pgfqpoint{1.795599in}{0.545307in}}%
\pgfpathlineto{\pgfqpoint{1.787662in}{0.570458in}}%
\pgfpathlineto{\pgfqpoint{1.774291in}{0.576196in}}%
\pgfpathlineto{\pgfqpoint{1.770975in}{0.581765in}}%
\pgfpathlineto{\pgfqpoint{1.765743in}{0.581820in}}%
\pgfpathlineto{\pgfqpoint{1.764790in}{0.589806in}}%
\pgfpathlineto{\pgfqpoint{1.760890in}{0.592713in}}%
\pgfpathlineto{\pgfqpoint{1.775795in}{0.607167in}}%
\pgfpathlineto{\pgfqpoint{1.781271in}{0.614134in}}%
\pgfpathlineto{\pgfqpoint{1.786479in}{0.611427in}}%
\pgfpathlineto{\pgfqpoint{1.793823in}{0.610228in}}%
\pgfpathlineto{\pgfqpoint{1.803223in}{0.606687in}}%
\pgfpathlineto{\pgfqpoint{1.821877in}{0.611209in}}%
\pgfpathlineto{\pgfqpoint{1.834269in}{0.615251in}}%
\pgfpathlineto{\pgfqpoint{1.851931in}{0.616696in}}%
\pgfpathlineto{\pgfqpoint{1.854631in}{0.617364in}}%
\pgfpathlineto{\pgfqpoint{1.854154in}{0.604333in}}%
\pgfpathlineto{\pgfqpoint{1.856414in}{0.600398in}}%
\pgfpathlineto{\pgfqpoint{1.853934in}{0.596710in}}%
\pgfpathlineto{\pgfqpoint{1.850658in}{0.583534in}}%
\pgfpathlineto{\pgfqpoint{1.847738in}{0.581014in}}%
\pgfpathlineto{\pgfqpoint{1.844939in}{0.573686in}}%
\pgfpathlineto{\pgfqpoint{1.847088in}{0.567676in}}%
\pgfpathlineto{\pgfqpoint{1.843280in}{0.560440in}}%
\pgfpathlineto{\pgfqpoint{1.846835in}{0.557616in}}%
\pgfpathlineto{\pgfqpoint{1.846781in}{0.541896in}}%
\pgfpathlineto{\pgfqpoint{1.842229in}{0.537436in}}%
\pgfpathlineto{\pgfqpoint{1.840585in}{0.532547in}}%
\pgfpathlineto{\pgfqpoint{1.834496in}{0.525796in}}%
\pgfpathlineto{\pgfqpoint{1.839662in}{0.517608in}}%
\pgfpathclose%
\pgfusepath{fill}%
\end{pgfscope}%
\begin{pgfscope}%
\pgfpathrectangle{\pgfqpoint{0.100000in}{0.100000in}}{\pgfqpoint{3.007045in}{1.925000in}}%
\pgfusepath{clip}%
\pgfsetbuttcap%
\pgfsetmiterjoin%
\definecolor{currentfill}{rgb}{0.260715,0.573841,0.777209}%
\pgfsetfillcolor{currentfill}%
\pgfsetlinewidth{0.000000pt}%
\definecolor{currentstroke}{rgb}{0.000000,0.000000,0.000000}%
\pgfsetstrokecolor{currentstroke}%
\pgfsetstrokeopacity{0.000000}%
\pgfsetdash{}{0pt}%
\pgfpathmoveto{\pgfqpoint{2.697515in}{1.156406in}}%
\pgfpathlineto{\pgfqpoint{2.692040in}{1.154397in}}%
\pgfpathlineto{\pgfqpoint{2.685472in}{1.142255in}}%
\pgfpathlineto{\pgfqpoint{2.681074in}{1.143862in}}%
\pgfpathlineto{\pgfqpoint{2.679314in}{1.148103in}}%
\pgfpathlineto{\pgfqpoint{2.665522in}{1.153025in}}%
\pgfpathlineto{\pgfqpoint{2.656779in}{1.157595in}}%
\pgfpathlineto{\pgfqpoint{2.650904in}{1.157743in}}%
\pgfpathlineto{\pgfqpoint{2.663143in}{1.179525in}}%
\pgfpathlineto{\pgfqpoint{2.662595in}{1.183012in}}%
\pgfpathlineto{\pgfqpoint{2.668861in}{1.192884in}}%
\pgfpathlineto{\pgfqpoint{2.678549in}{1.188294in}}%
\pgfpathlineto{\pgfqpoint{2.680316in}{1.180573in}}%
\pgfpathlineto{\pgfqpoint{2.682867in}{1.178238in}}%
\pgfpathlineto{\pgfqpoint{2.691552in}{1.182331in}}%
\pgfpathlineto{\pgfqpoint{2.695566in}{1.180728in}}%
\pgfpathlineto{\pgfqpoint{2.697143in}{1.177380in}}%
\pgfpathlineto{\pgfqpoint{2.695805in}{1.170685in}}%
\pgfpathlineto{\pgfqpoint{2.690990in}{1.162746in}}%
\pgfpathlineto{\pgfqpoint{2.697515in}{1.156406in}}%
\pgfpathclose%
\pgfusepath{fill}%
\end{pgfscope}%
\begin{pgfscope}%
\pgfpathrectangle{\pgfqpoint{0.100000in}{0.100000in}}{\pgfqpoint{3.007045in}{1.925000in}}%
\pgfusepath{clip}%
\pgfsetbuttcap%
\pgfsetmiterjoin%
\definecolor{currentfill}{rgb}{0.056363,0.349635,0.636755}%
\pgfsetfillcolor{currentfill}%
\pgfsetlinewidth{0.000000pt}%
\definecolor{currentstroke}{rgb}{0.000000,0.000000,0.000000}%
\pgfsetstrokecolor{currentstroke}%
\pgfsetstrokeopacity{0.000000}%
\pgfsetdash{}{0pt}%
\pgfpathmoveto{\pgfqpoint{1.534548in}{1.032337in}}%
\pgfpathlineto{\pgfqpoint{1.533400in}{1.007220in}}%
\pgfpathlineto{\pgfqpoint{1.509153in}{1.008633in}}%
\pgfpathlineto{\pgfqpoint{1.504587in}{1.008926in}}%
\pgfpathlineto{\pgfqpoint{1.505902in}{1.034236in}}%
\pgfpathlineto{\pgfqpoint{1.505407in}{1.040070in}}%
\pgfpathlineto{\pgfqpoint{1.476919in}{1.041763in}}%
\pgfpathlineto{\pgfqpoint{1.478439in}{1.065027in}}%
\pgfpathlineto{\pgfqpoint{1.478495in}{1.076500in}}%
\pgfpathlineto{\pgfqpoint{1.501163in}{1.075036in}}%
\pgfpathlineto{\pgfqpoint{1.502124in}{1.092214in}}%
\pgfpathlineto{\pgfqpoint{1.535329in}{1.090418in}}%
\pgfpathlineto{\pgfqpoint{1.536034in}{1.090382in}}%
\pgfpathlineto{\pgfqpoint{1.534600in}{1.061649in}}%
\pgfpathlineto{\pgfqpoint{1.534859in}{1.055660in}}%
\pgfpathlineto{\pgfqpoint{1.534548in}{1.032337in}}%
\pgfpathclose%
\pgfusepath{fill}%
\end{pgfscope}%
\begin{pgfscope}%
\pgfpathrectangle{\pgfqpoint{0.100000in}{0.100000in}}{\pgfqpoint{3.007045in}{1.925000in}}%
\pgfusepath{clip}%
\pgfsetbuttcap%
\pgfsetmiterjoin%
\definecolor{currentfill}{rgb}{0.504314,0.728858,0.857486}%
\pgfsetfillcolor{currentfill}%
\pgfsetlinewidth{0.000000pt}%
\definecolor{currentstroke}{rgb}{0.000000,0.000000,0.000000}%
\pgfsetstrokecolor{currentstroke}%
\pgfsetstrokeopacity{0.000000}%
\pgfsetdash{}{0pt}%
\pgfpathmoveto{\pgfqpoint{1.482882in}{0.684491in}}%
\pgfpathlineto{\pgfqpoint{1.478210in}{0.684741in}}%
\pgfpathlineto{\pgfqpoint{1.454083in}{0.686391in}}%
\pgfpathlineto{\pgfqpoint{1.458235in}{0.744453in}}%
\pgfpathlineto{\pgfqpoint{1.465814in}{0.743950in}}%
\pgfpathlineto{\pgfqpoint{1.486605in}{0.742533in}}%
\pgfpathlineto{\pgfqpoint{1.482882in}{0.684491in}}%
\pgfpathclose%
\pgfusepath{fill}%
\end{pgfscope}%
\begin{pgfscope}%
\pgfpathrectangle{\pgfqpoint{0.100000in}{0.100000in}}{\pgfqpoint{3.007045in}{1.925000in}}%
\pgfusepath{clip}%
\pgfsetbuttcap%
\pgfsetmiterjoin%
\definecolor{currentfill}{rgb}{0.056363,0.349635,0.636755}%
\pgfsetfillcolor{currentfill}%
\pgfsetlinewidth{0.000000pt}%
\definecolor{currentstroke}{rgb}{0.000000,0.000000,0.000000}%
\pgfsetstrokecolor{currentstroke}%
\pgfsetstrokeopacity{0.000000}%
\pgfsetdash{}{0pt}%
\pgfpathmoveto{\pgfqpoint{1.476919in}{1.041763in}}%
\pgfpathlineto{\pgfqpoint{1.476575in}{1.036016in}}%
\pgfpathlineto{\pgfqpoint{1.454880in}{1.037434in}}%
\pgfpathlineto{\pgfqpoint{1.453744in}{1.037504in}}%
\pgfpathlineto{\pgfqpoint{1.455281in}{1.060433in}}%
\pgfpathlineto{\pgfqpoint{1.405795in}{1.064151in}}%
\pgfpathlineto{\pgfqpoint{1.408351in}{1.098572in}}%
\pgfpathlineto{\pgfqpoint{1.432949in}{1.096770in}}%
\pgfpathlineto{\pgfqpoint{1.478511in}{1.093767in}}%
\pgfpathlineto{\pgfqpoint{1.502124in}{1.092214in}}%
\pgfpathlineto{\pgfqpoint{1.501163in}{1.075036in}}%
\pgfpathlineto{\pgfqpoint{1.478495in}{1.076500in}}%
\pgfpathlineto{\pgfqpoint{1.478439in}{1.065027in}}%
\pgfpathlineto{\pgfqpoint{1.476919in}{1.041763in}}%
\pgfpathclose%
\pgfusepath{fill}%
\end{pgfscope}%
\begin{pgfscope}%
\pgfpathrectangle{\pgfqpoint{0.100000in}{0.100000in}}{\pgfqpoint{3.007045in}{1.925000in}}%
\pgfusepath{clip}%
\pgfsetbuttcap%
\pgfsetmiterjoin%
\definecolor{currentfill}{rgb}{0.429020,0.687520,0.841246}%
\pgfsetfillcolor{currentfill}%
\pgfsetlinewidth{0.000000pt}%
\definecolor{currentstroke}{rgb}{0.000000,0.000000,0.000000}%
\pgfsetstrokecolor{currentstroke}%
\pgfsetstrokeopacity{0.000000}%
\pgfsetdash{}{0pt}%
\pgfpathmoveto{\pgfqpoint{1.917287in}{1.088655in}}%
\pgfpathlineto{\pgfqpoint{1.945763in}{1.089375in}}%
\pgfpathlineto{\pgfqpoint{1.946082in}{1.080499in}}%
\pgfpathlineto{\pgfqpoint{1.952290in}{1.080650in}}%
\pgfpathlineto{\pgfqpoint{1.953009in}{1.056717in}}%
\pgfpathlineto{\pgfqpoint{1.964484in}{1.057050in}}%
\pgfpathlineto{\pgfqpoint{1.964595in}{1.051309in}}%
\pgfpathlineto{\pgfqpoint{1.972796in}{1.051471in}}%
\pgfpathlineto{\pgfqpoint{1.973010in}{1.044275in}}%
\pgfpathlineto{\pgfqpoint{1.981728in}{1.045692in}}%
\pgfpathlineto{\pgfqpoint{1.993117in}{1.046153in}}%
\pgfpathlineto{\pgfqpoint{1.993598in}{1.031047in}}%
\pgfpathlineto{\pgfqpoint{1.995876in}{1.024384in}}%
\pgfpathlineto{\pgfqpoint{1.996111in}{1.017492in}}%
\pgfpathlineto{\pgfqpoint{1.994242in}{1.009746in}}%
\pgfpathlineto{\pgfqpoint{1.984516in}{1.009519in}}%
\pgfpathlineto{\pgfqpoint{1.981594in}{1.012166in}}%
\pgfpathlineto{\pgfqpoint{1.981501in}{1.015319in}}%
\pgfpathlineto{\pgfqpoint{1.975616in}{1.018755in}}%
\pgfpathlineto{\pgfqpoint{1.972674in}{1.026092in}}%
\pgfpathlineto{\pgfqpoint{1.972430in}{1.032823in}}%
\pgfpathlineto{\pgfqpoint{1.942096in}{1.032256in}}%
\pgfpathlineto{\pgfqpoint{1.941697in}{1.043763in}}%
\pgfpathlineto{\pgfqpoint{1.915878in}{1.043286in}}%
\pgfpathlineto{\pgfqpoint{1.915821in}{1.046178in}}%
\pgfpathlineto{\pgfqpoint{1.907105in}{1.050239in}}%
\pgfpathlineto{\pgfqpoint{1.907064in}{1.059811in}}%
\pgfpathlineto{\pgfqpoint{1.906913in}{1.070484in}}%
\pgfpathlineto{\pgfqpoint{1.918382in}{1.070506in}}%
\pgfpathlineto{\pgfqpoint{1.917287in}{1.088655in}}%
\pgfpathclose%
\pgfusepath{fill}%
\end{pgfscope}%
\begin{pgfscope}%
\pgfpathrectangle{\pgfqpoint{0.100000in}{0.100000in}}{\pgfqpoint{3.007045in}{1.925000in}}%
\pgfusepath{clip}%
\pgfsetbuttcap%
\pgfsetmiterjoin%
\definecolor{currentfill}{rgb}{0.376732,0.653072,0.822484}%
\pgfsetfillcolor{currentfill}%
\pgfsetlinewidth{0.000000pt}%
\definecolor{currentstroke}{rgb}{0.000000,0.000000,0.000000}%
\pgfsetstrokecolor{currentstroke}%
\pgfsetstrokeopacity{0.000000}%
\pgfsetdash{}{0pt}%
\pgfpathmoveto{\pgfqpoint{2.375927in}{0.644931in}}%
\pgfpathlineto{\pgfqpoint{2.378692in}{0.619214in}}%
\pgfpathlineto{\pgfqpoint{2.351014in}{0.617483in}}%
\pgfpathlineto{\pgfqpoint{2.346074in}{0.624253in}}%
\pgfpathlineto{\pgfqpoint{2.345288in}{0.628691in}}%
\pgfpathlineto{\pgfqpoint{2.341073in}{0.635732in}}%
\pgfpathlineto{\pgfqpoint{2.339044in}{0.640515in}}%
\pgfpathlineto{\pgfqpoint{2.345074in}{0.640987in}}%
\pgfpathlineto{\pgfqpoint{2.343997in}{0.653135in}}%
\pgfpathlineto{\pgfqpoint{2.359552in}{0.655091in}}%
\pgfpathlineto{\pgfqpoint{2.365478in}{0.655570in}}%
\pgfpathlineto{\pgfqpoint{2.366514in}{0.643850in}}%
\pgfpathlineto{\pgfqpoint{2.375927in}{0.644931in}}%
\pgfpathclose%
\pgfusepath{fill}%
\end{pgfscope}%
\begin{pgfscope}%
\pgfpathrectangle{\pgfqpoint{0.100000in}{0.100000in}}{\pgfqpoint{3.007045in}{1.925000in}}%
\pgfusepath{clip}%
\pgfsetbuttcap%
\pgfsetmiterjoin%
\definecolor{currentfill}{rgb}{0.166967,0.480692,0.729150}%
\pgfsetfillcolor{currentfill}%
\pgfsetlinewidth{0.000000pt}%
\definecolor{currentstroke}{rgb}{0.000000,0.000000,0.000000}%
\pgfsetstrokecolor{currentstroke}%
\pgfsetstrokeopacity{0.000000}%
\pgfsetdash{}{0pt}%
\pgfpathmoveto{\pgfqpoint{1.704139in}{0.513794in}}%
\pgfpathlineto{\pgfqpoint{1.699409in}{0.516089in}}%
\pgfpathlineto{\pgfqpoint{1.695319in}{0.526005in}}%
\pgfpathlineto{\pgfqpoint{1.690540in}{0.526743in}}%
\pgfpathlineto{\pgfqpoint{1.683561in}{0.532352in}}%
\pgfpathlineto{\pgfqpoint{1.675326in}{0.548748in}}%
\pgfpathlineto{\pgfqpoint{1.668934in}{0.549504in}}%
\pgfpathlineto{\pgfqpoint{1.672597in}{0.554943in}}%
\pgfpathlineto{\pgfqpoint{1.677967in}{0.558409in}}%
\pgfpathlineto{\pgfqpoint{1.690544in}{0.561282in}}%
\pgfpathlineto{\pgfqpoint{1.696897in}{0.563937in}}%
\pgfpathlineto{\pgfqpoint{1.704636in}{0.564588in}}%
\pgfpathlineto{\pgfqpoint{1.709428in}{0.556522in}}%
\pgfpathlineto{\pgfqpoint{1.709639in}{0.553402in}}%
\pgfpathlineto{\pgfqpoint{1.726035in}{0.554362in}}%
\pgfpathlineto{\pgfqpoint{1.724776in}{0.579988in}}%
\pgfpathlineto{\pgfqpoint{1.737969in}{0.571808in}}%
\pgfpathlineto{\pgfqpoint{1.751572in}{0.571734in}}%
\pgfpathlineto{\pgfqpoint{1.753626in}{0.595060in}}%
\pgfpathlineto{\pgfqpoint{1.757986in}{0.598171in}}%
\pgfpathlineto{\pgfqpoint{1.760890in}{0.592713in}}%
\pgfpathlineto{\pgfqpoint{1.764790in}{0.589806in}}%
\pgfpathlineto{\pgfqpoint{1.765743in}{0.581820in}}%
\pgfpathlineto{\pgfqpoint{1.770975in}{0.581765in}}%
\pgfpathlineto{\pgfqpoint{1.774291in}{0.576196in}}%
\pgfpathlineto{\pgfqpoint{1.787662in}{0.570458in}}%
\pgfpathlineto{\pgfqpoint{1.795599in}{0.545307in}}%
\pgfpathlineto{\pgfqpoint{1.804317in}{0.545414in}}%
\pgfpathlineto{\pgfqpoint{1.804575in}{0.530516in}}%
\pgfpathlineto{\pgfqpoint{1.809447in}{0.530410in}}%
\pgfpathlineto{\pgfqpoint{1.809818in}{0.508665in}}%
\pgfpathlineto{\pgfqpoint{1.794978in}{0.501878in}}%
\pgfpathlineto{\pgfqpoint{1.797904in}{0.509417in}}%
\pgfpathlineto{\pgfqpoint{1.787377in}{0.506124in}}%
\pgfpathlineto{\pgfqpoint{1.790301in}{0.517701in}}%
\pgfpathlineto{\pgfqpoint{1.787535in}{0.523539in}}%
\pgfpathlineto{\pgfqpoint{1.783015in}{0.521586in}}%
\pgfpathlineto{\pgfqpoint{1.779708in}{0.515842in}}%
\pgfpathlineto{\pgfqpoint{1.774796in}{0.517671in}}%
\pgfpathlineto{\pgfqpoint{1.771403in}{0.513633in}}%
\pgfpathlineto{\pgfqpoint{1.771383in}{0.507282in}}%
\pgfpathlineto{\pgfqpoint{1.777596in}{0.504206in}}%
\pgfpathlineto{\pgfqpoint{1.778463in}{0.493573in}}%
\pgfpathlineto{\pgfqpoint{1.788184in}{0.493368in}}%
\pgfpathlineto{\pgfqpoint{1.765360in}{0.476226in}}%
\pgfpathlineto{\pgfqpoint{1.769052in}{0.484340in}}%
\pgfpathlineto{\pgfqpoint{1.758777in}{0.502145in}}%
\pgfpathlineto{\pgfqpoint{1.759753in}{0.508170in}}%
\pgfpathlineto{\pgfqpoint{1.757144in}{0.510663in}}%
\pgfpathlineto{\pgfqpoint{1.747709in}{0.509837in}}%
\pgfpathlineto{\pgfqpoint{1.745433in}{0.500436in}}%
\pgfpathlineto{\pgfqpoint{1.740375in}{0.500409in}}%
\pgfpathlineto{\pgfqpoint{1.738650in}{0.493869in}}%
\pgfpathlineto{\pgfqpoint{1.734233in}{0.490553in}}%
\pgfpathlineto{\pgfqpoint{1.728488in}{0.492886in}}%
\pgfpathlineto{\pgfqpoint{1.722932in}{0.488802in}}%
\pgfpathlineto{\pgfqpoint{1.717041in}{0.493709in}}%
\pgfpathlineto{\pgfqpoint{1.714651in}{0.500322in}}%
\pgfpathlineto{\pgfqpoint{1.710733in}{0.503026in}}%
\pgfpathlineto{\pgfqpoint{1.710725in}{0.510845in}}%
\pgfpathlineto{\pgfqpoint{1.704139in}{0.513794in}}%
\pgfpathclose%
\pgfusepath{fill}%
\end{pgfscope}%
\begin{pgfscope}%
\pgfpathrectangle{\pgfqpoint{0.100000in}{0.100000in}}{\pgfqpoint{3.007045in}{1.925000in}}%
\pgfusepath{clip}%
\pgfsetbuttcap%
\pgfsetmiterjoin%
\definecolor{currentfill}{rgb}{0.275848,0.584175,0.783114}%
\pgfsetfillcolor{currentfill}%
\pgfsetlinewidth{0.000000pt}%
\definecolor{currentstroke}{rgb}{0.000000,0.000000,0.000000}%
\pgfsetstrokecolor{currentstroke}%
\pgfsetstrokeopacity{0.000000}%
\pgfsetdash{}{0pt}%
\pgfpathmoveto{\pgfqpoint{1.631785in}{0.798613in}}%
\pgfpathlineto{\pgfqpoint{1.625906in}{0.804863in}}%
\pgfpathlineto{\pgfqpoint{1.614953in}{0.795791in}}%
\pgfpathlineto{\pgfqpoint{1.608887in}{0.798669in}}%
\pgfpathlineto{\pgfqpoint{1.610850in}{0.805245in}}%
\pgfpathlineto{\pgfqpoint{1.605264in}{0.805706in}}%
\pgfpathlineto{\pgfqpoint{1.598818in}{0.813845in}}%
\pgfpathlineto{\pgfqpoint{1.588268in}{0.817057in}}%
\pgfpathlineto{\pgfqpoint{1.586260in}{0.813360in}}%
\pgfpathlineto{\pgfqpoint{1.581430in}{0.811070in}}%
\pgfpathlineto{\pgfqpoint{1.574901in}{0.817578in}}%
\pgfpathlineto{\pgfqpoint{1.575401in}{0.829239in}}%
\pgfpathlineto{\pgfqpoint{1.572553in}{0.829339in}}%
\pgfpathlineto{\pgfqpoint{1.573014in}{0.840820in}}%
\pgfpathlineto{\pgfqpoint{1.564075in}{0.841182in}}%
\pgfpathlineto{\pgfqpoint{1.564310in}{0.846924in}}%
\pgfpathlineto{\pgfqpoint{1.565031in}{0.864162in}}%
\pgfpathlineto{\pgfqpoint{1.576210in}{0.863673in}}%
\pgfpathlineto{\pgfqpoint{1.604926in}{0.862611in}}%
\pgfpathlineto{\pgfqpoint{1.604562in}{0.851135in}}%
\pgfpathlineto{\pgfqpoint{1.627401in}{0.850427in}}%
\pgfpathlineto{\pgfqpoint{1.633138in}{0.850249in}}%
\pgfpathlineto{\pgfqpoint{1.631785in}{0.798613in}}%
\pgfpathclose%
\pgfusepath{fill}%
\end{pgfscope}%
\begin{pgfscope}%
\pgfpathrectangle{\pgfqpoint{0.100000in}{0.100000in}}{\pgfqpoint{3.007045in}{1.925000in}}%
\pgfusepath{clip}%
\pgfsetbuttcap%
\pgfsetmiterjoin%
\definecolor{currentfill}{rgb}{0.691580,0.822745,0.907543}%
\pgfsetfillcolor{currentfill}%
\pgfsetlinewidth{0.000000pt}%
\definecolor{currentstroke}{rgb}{0.000000,0.000000,0.000000}%
\pgfsetstrokecolor{currentstroke}%
\pgfsetstrokeopacity{0.000000}%
\pgfsetdash{}{0pt}%
\pgfpathmoveto{\pgfqpoint{2.815641in}{1.458005in}}%
\pgfpathlineto{\pgfqpoint{2.815310in}{1.453061in}}%
\pgfpathlineto{\pgfqpoint{2.806124in}{1.452725in}}%
\pgfpathlineto{\pgfqpoint{2.795846in}{1.455210in}}%
\pgfpathlineto{\pgfqpoint{2.794361in}{1.458299in}}%
\pgfpathlineto{\pgfqpoint{2.787830in}{1.459571in}}%
\pgfpathlineto{\pgfqpoint{2.787399in}{1.454470in}}%
\pgfpathlineto{\pgfqpoint{2.776040in}{1.453025in}}%
\pgfpathlineto{\pgfqpoint{2.768356in}{1.454967in}}%
\pgfpathlineto{\pgfqpoint{2.763804in}{1.456101in}}%
\pgfpathlineto{\pgfqpoint{2.766813in}{1.466787in}}%
\pgfpathlineto{\pgfqpoint{2.756645in}{1.469554in}}%
\pgfpathlineto{\pgfqpoint{2.750879in}{1.474563in}}%
\pgfpathlineto{\pgfqpoint{2.753121in}{1.482442in}}%
\pgfpathlineto{\pgfqpoint{2.749558in}{1.490026in}}%
\pgfpathlineto{\pgfqpoint{2.749311in}{1.495307in}}%
\pgfpathlineto{\pgfqpoint{2.759273in}{1.493756in}}%
\pgfpathlineto{\pgfqpoint{2.766923in}{1.496257in}}%
\pgfpathlineto{\pgfqpoint{2.774908in}{1.505824in}}%
\pgfpathlineto{\pgfqpoint{2.756305in}{1.557859in}}%
\pgfpathlineto{\pgfqpoint{2.763396in}{1.560589in}}%
\pgfpathlineto{\pgfqpoint{2.793875in}{1.572201in}}%
\pgfpathlineto{\pgfqpoint{2.794553in}{1.564410in}}%
\pgfpathlineto{\pgfqpoint{2.803917in}{1.557332in}}%
\pgfpathlineto{\pgfqpoint{2.807383in}{1.541913in}}%
\pgfpathlineto{\pgfqpoint{2.813736in}{1.511035in}}%
\pgfpathlineto{\pgfqpoint{2.815796in}{1.505313in}}%
\pgfpathlineto{\pgfqpoint{2.814539in}{1.460221in}}%
\pgfpathlineto{\pgfqpoint{2.815641in}{1.458005in}}%
\pgfpathclose%
\pgfusepath{fill}%
\end{pgfscope}%
\begin{pgfscope}%
\pgfpathrectangle{\pgfqpoint{0.100000in}{0.100000in}}{\pgfqpoint{3.007045in}{1.925000in}}%
\pgfusepath{clip}%
\pgfsetbuttcap%
\pgfsetmiterjoin%
\definecolor{currentfill}{rgb}{0.361599,0.642737,0.816578}%
\pgfsetfillcolor{currentfill}%
\pgfsetlinewidth{0.000000pt}%
\definecolor{currentstroke}{rgb}{0.000000,0.000000,0.000000}%
\pgfsetstrokecolor{currentstroke}%
\pgfsetstrokeopacity{0.000000}%
\pgfsetdash{}{0pt}%
\pgfpathmoveto{\pgfqpoint{1.154054in}{1.574061in}}%
\pgfpathlineto{\pgfqpoint{1.115456in}{1.579920in}}%
\pgfpathlineto{\pgfqpoint{1.114575in}{1.580480in}}%
\pgfpathlineto{\pgfqpoint{1.082392in}{1.585392in}}%
\pgfpathlineto{\pgfqpoint{1.084238in}{1.596483in}}%
\pgfpathlineto{\pgfqpoint{1.072000in}{1.598524in}}%
\pgfpathlineto{\pgfqpoint{1.073953in}{1.610065in}}%
\pgfpathlineto{\pgfqpoint{1.079975in}{1.609049in}}%
\pgfpathlineto{\pgfqpoint{1.081871in}{1.620429in}}%
\pgfpathlineto{\pgfqpoint{1.094101in}{1.624217in}}%
\pgfpathlineto{\pgfqpoint{1.099750in}{1.623315in}}%
\pgfpathlineto{\pgfqpoint{1.101489in}{1.634717in}}%
\pgfpathlineto{\pgfqpoint{1.105971in}{1.645829in}}%
\pgfpathlineto{\pgfqpoint{1.109783in}{1.645289in}}%
\pgfpathlineto{\pgfqpoint{1.110938in}{1.650768in}}%
\pgfpathlineto{\pgfqpoint{1.102412in}{1.652202in}}%
\pgfpathlineto{\pgfqpoint{1.101102in}{1.658254in}}%
\pgfpathlineto{\pgfqpoint{1.102024in}{1.663968in}}%
\pgfpathlineto{\pgfqpoint{1.113428in}{1.662137in}}%
\pgfpathlineto{\pgfqpoint{1.119016in}{1.693144in}}%
\pgfpathlineto{\pgfqpoint{1.119627in}{1.696907in}}%
\pgfpathlineto{\pgfqpoint{1.153729in}{1.691500in}}%
\pgfpathlineto{\pgfqpoint{1.189998in}{1.686616in}}%
\pgfpathlineto{\pgfqpoint{1.190856in}{1.681412in}}%
\pgfpathlineto{\pgfqpoint{1.188964in}{1.669306in}}%
\pgfpathlineto{\pgfqpoint{1.192292in}{1.666923in}}%
\pgfpathlineto{\pgfqpoint{1.195568in}{1.655776in}}%
\pgfpathlineto{\pgfqpoint{1.201035in}{1.646271in}}%
\pgfpathlineto{\pgfqpoint{1.197757in}{1.637726in}}%
\pgfpathlineto{\pgfqpoint{1.190334in}{1.638786in}}%
\pgfpathlineto{\pgfqpoint{1.189792in}{1.634979in}}%
\pgfpathlineto{\pgfqpoint{1.184148in}{1.635775in}}%
\pgfpathlineto{\pgfqpoint{1.179523in}{1.630622in}}%
\pgfpathlineto{\pgfqpoint{1.171947in}{1.631734in}}%
\pgfpathlineto{\pgfqpoint{1.169858in}{1.624279in}}%
\pgfpathlineto{\pgfqpoint{1.167306in}{1.606868in}}%
\pgfpathlineto{\pgfqpoint{1.155409in}{1.604740in}}%
\pgfpathlineto{\pgfqpoint{1.140335in}{1.607233in}}%
\pgfpathlineto{\pgfqpoint{1.138319in}{1.605460in}}%
\pgfpathlineto{\pgfqpoint{1.135376in}{1.591485in}}%
\pgfpathlineto{\pgfqpoint{1.163161in}{1.587525in}}%
\pgfpathlineto{\pgfqpoint{1.160708in}{1.586055in}}%
\pgfpathlineto{\pgfqpoint{1.154054in}{1.574061in}}%
\pgfpathclose%
\pgfusepath{fill}%
\end{pgfscope}%
\begin{pgfscope}%
\pgfpathrectangle{\pgfqpoint{0.100000in}{0.100000in}}{\pgfqpoint{3.007045in}{1.925000in}}%
\pgfusepath{clip}%
\pgfsetbuttcap%
\pgfsetmiterjoin%
\definecolor{currentfill}{rgb}{0.579608,0.770196,0.873725}%
\pgfsetfillcolor{currentfill}%
\pgfsetlinewidth{0.000000pt}%
\definecolor{currentstroke}{rgb}{0.000000,0.000000,0.000000}%
\pgfsetstrokecolor{currentstroke}%
\pgfsetstrokeopacity{0.000000}%
\pgfsetdash{}{0pt}%
\pgfpathmoveto{\pgfqpoint{2.427990in}{0.752079in}}%
\pgfpathlineto{\pgfqpoint{2.415340in}{0.771423in}}%
\pgfpathlineto{\pgfqpoint{2.417835in}{0.773607in}}%
\pgfpathlineto{\pgfqpoint{2.411719in}{0.790134in}}%
\pgfpathlineto{\pgfqpoint{2.405437in}{0.788426in}}%
\pgfpathlineto{\pgfqpoint{2.403890in}{0.805755in}}%
\pgfpathlineto{\pgfqpoint{2.417379in}{0.810753in}}%
\pgfpathlineto{\pgfqpoint{2.424800in}{0.803148in}}%
\pgfpathlineto{\pgfqpoint{2.432103in}{0.811727in}}%
\pgfpathlineto{\pgfqpoint{2.441136in}{0.811210in}}%
\pgfpathlineto{\pgfqpoint{2.448640in}{0.798533in}}%
\pgfpathlineto{\pgfqpoint{2.443204in}{0.795028in}}%
\pgfpathlineto{\pgfqpoint{2.441806in}{0.790953in}}%
\pgfpathlineto{\pgfqpoint{2.437607in}{0.788428in}}%
\pgfpathlineto{\pgfqpoint{2.432968in}{0.782406in}}%
\pgfpathlineto{\pgfqpoint{2.433317in}{0.776989in}}%
\pgfpathlineto{\pgfqpoint{2.436633in}{0.771428in}}%
\pgfpathlineto{\pgfqpoint{2.441164in}{0.768340in}}%
\pgfpathlineto{\pgfqpoint{2.441826in}{0.762195in}}%
\pgfpathlineto{\pgfqpoint{2.427990in}{0.752079in}}%
\pgfpathclose%
\pgfusepath{fill}%
\end{pgfscope}%
\begin{pgfscope}%
\pgfpathrectangle{\pgfqpoint{0.100000in}{0.100000in}}{\pgfqpoint{3.007045in}{1.925000in}}%
\pgfusepath{clip}%
\pgfsetbuttcap%
\pgfsetmiterjoin%
\definecolor{currentfill}{rgb}{0.183206,0.496932,0.737516}%
\pgfsetfillcolor{currentfill}%
\pgfsetlinewidth{0.000000pt}%
\definecolor{currentstroke}{rgb}{0.000000,0.000000,0.000000}%
\pgfsetstrokecolor{currentstroke}%
\pgfsetstrokeopacity{0.000000}%
\pgfsetdash{}{0pt}%
\pgfpathmoveto{\pgfqpoint{1.263629in}{1.757459in}}%
\pgfpathlineto{\pgfqpoint{1.265560in}{1.751233in}}%
\pgfpathlineto{\pgfqpoint{1.269873in}{1.750687in}}%
\pgfpathlineto{\pgfqpoint{1.267101in}{1.728242in}}%
\pgfpathlineto{\pgfqpoint{1.263981in}{1.716912in}}%
\pgfpathlineto{\pgfqpoint{1.269690in}{1.716224in}}%
\pgfpathlineto{\pgfqpoint{1.268284in}{1.704741in}}%
\pgfpathlineto{\pgfqpoint{1.272031in}{1.704281in}}%
\pgfpathlineto{\pgfqpoint{1.269340in}{1.682158in}}%
\pgfpathlineto{\pgfqpoint{1.240737in}{1.685718in}}%
\pgfpathlineto{\pgfqpoint{1.240870in}{1.686663in}}%
\pgfpathlineto{\pgfqpoint{1.187986in}{1.693328in}}%
\pgfpathlineto{\pgfqpoint{1.186086in}{1.695738in}}%
\pgfpathlineto{\pgfqpoint{1.189559in}{1.703789in}}%
\pgfpathlineto{\pgfqpoint{1.188167in}{1.711464in}}%
\pgfpathlineto{\pgfqpoint{1.188345in}{1.720102in}}%
\pgfpathlineto{\pgfqpoint{1.192875in}{1.732905in}}%
\pgfpathlineto{\pgfqpoint{1.195585in}{1.736909in}}%
\pgfpathlineto{\pgfqpoint{1.201819in}{1.739713in}}%
\pgfpathlineto{\pgfqpoint{1.206785in}{1.744306in}}%
\pgfpathlineto{\pgfqpoint{1.214896in}{1.741346in}}%
\pgfpathlineto{\pgfqpoint{1.217097in}{1.745725in}}%
\pgfpathlineto{\pgfqpoint{1.224257in}{1.742711in}}%
\pgfpathlineto{\pgfqpoint{1.240243in}{1.741597in}}%
\pgfpathlineto{\pgfqpoint{1.241819in}{1.745330in}}%
\pgfpathlineto{\pgfqpoint{1.246943in}{1.744290in}}%
\pgfpathlineto{\pgfqpoint{1.255859in}{1.748834in}}%
\pgfpathlineto{\pgfqpoint{1.260272in}{1.755961in}}%
\pgfpathlineto{\pgfqpoint{1.263629in}{1.757459in}}%
\pgfpathclose%
\pgfusepath{fill}%
\end{pgfscope}%
\begin{pgfscope}%
\pgfpathrectangle{\pgfqpoint{0.100000in}{0.100000in}}{\pgfqpoint{3.007045in}{1.925000in}}%
\pgfusepath{clip}%
\pgfsetbuttcap%
\pgfsetmiterjoin%
\definecolor{currentfill}{rgb}{0.466667,0.708189,0.849366}%
\pgfsetfillcolor{currentfill}%
\pgfsetlinewidth{0.000000pt}%
\definecolor{currentstroke}{rgb}{0.000000,0.000000,0.000000}%
\pgfsetstrokecolor{currentstroke}%
\pgfsetstrokeopacity{0.000000}%
\pgfsetdash{}{0pt}%
\pgfpathmoveto{\pgfqpoint{2.644431in}{1.224390in}}%
\pgfpathlineto{\pgfqpoint{2.614371in}{1.241096in}}%
\pgfpathlineto{\pgfqpoint{2.607688in}{1.243264in}}%
\pgfpathlineto{\pgfqpoint{2.608755in}{1.249542in}}%
\pgfpathlineto{\pgfqpoint{2.613340in}{1.256035in}}%
\pgfpathlineto{\pgfqpoint{2.609634in}{1.257414in}}%
\pgfpathlineto{\pgfqpoint{2.612460in}{1.264832in}}%
\pgfpathlineto{\pgfqpoint{2.614370in}{1.276431in}}%
\pgfpathlineto{\pgfqpoint{2.613885in}{1.282571in}}%
\pgfpathlineto{\pgfqpoint{2.616493in}{1.288353in}}%
\pgfpathlineto{\pgfqpoint{2.627614in}{1.286206in}}%
\pgfpathlineto{\pgfqpoint{2.631148in}{1.284175in}}%
\pgfpathlineto{\pgfqpoint{2.636603in}{1.298784in}}%
\pgfpathlineto{\pgfqpoint{2.638827in}{1.299653in}}%
\pgfpathlineto{\pgfqpoint{2.652060in}{1.280358in}}%
\pgfpathlineto{\pgfqpoint{2.655502in}{1.265959in}}%
\pgfpathlineto{\pgfqpoint{2.651264in}{1.258406in}}%
\pgfpathlineto{\pgfqpoint{2.649600in}{1.238259in}}%
\pgfpathlineto{\pgfqpoint{2.646870in}{1.236936in}}%
\pgfpathlineto{\pgfqpoint{2.644431in}{1.224390in}}%
\pgfpathclose%
\pgfusepath{fill}%
\end{pgfscope}%
\begin{pgfscope}%
\pgfpathrectangle{\pgfqpoint{0.100000in}{0.100000in}}{\pgfqpoint{3.007045in}{1.925000in}}%
\pgfusepath{clip}%
\pgfsetbuttcap%
\pgfsetmiterjoin%
\definecolor{currentfill}{rgb}{0.093272,0.396878,0.673664}%
\pgfsetfillcolor{currentfill}%
\pgfsetlinewidth{0.000000pt}%
\definecolor{currentstroke}{rgb}{0.000000,0.000000,0.000000}%
\pgfsetstrokecolor{currentstroke}%
\pgfsetstrokeopacity{0.000000}%
\pgfsetdash{}{0pt}%
\pgfpathmoveto{\pgfqpoint{0.677092in}{0.410736in}}%
\pgfpathlineto{\pgfqpoint{0.679449in}{0.413309in}}%
\pgfpathlineto{\pgfqpoint{0.683920in}{0.412319in}}%
\pgfpathlineto{\pgfqpoint{0.681471in}{0.410739in}}%
\pgfpathlineto{\pgfqpoint{0.677092in}{0.410736in}}%
\pgfpathclose%
\pgfusepath{fill}%
\end{pgfscope}%
\begin{pgfscope}%
\pgfpathrectangle{\pgfqpoint{0.100000in}{0.100000in}}{\pgfqpoint{3.007045in}{1.925000in}}%
\pgfusepath{clip}%
\pgfsetbuttcap%
\pgfsetmiterjoin%
\definecolor{currentfill}{rgb}{0.093272,0.396878,0.673664}%
\pgfsetfillcolor{currentfill}%
\pgfsetlinewidth{0.000000pt}%
\definecolor{currentstroke}{rgb}{0.000000,0.000000,0.000000}%
\pgfsetstrokecolor{currentstroke}%
\pgfsetstrokeopacity{0.000000}%
\pgfsetdash{}{0pt}%
\pgfpathmoveto{\pgfqpoint{0.784950in}{0.396276in}}%
\pgfpathlineto{\pgfqpoint{0.784096in}{0.394320in}}%
\pgfpathlineto{\pgfqpoint{0.778986in}{0.387487in}}%
\pgfpathlineto{\pgfqpoint{0.777700in}{0.388433in}}%
\pgfpathlineto{\pgfqpoint{0.772579in}{0.381674in}}%
\pgfpathlineto{\pgfqpoint{0.771354in}{0.382565in}}%
\pgfpathlineto{\pgfqpoint{0.766473in}{0.376050in}}%
\pgfpathlineto{\pgfqpoint{0.765233in}{0.376990in}}%
\pgfpathlineto{\pgfqpoint{0.763938in}{0.375291in}}%
\pgfpathlineto{\pgfqpoint{0.762258in}{0.376547in}}%
\pgfpathlineto{\pgfqpoint{0.759707in}{0.373154in}}%
\pgfpathlineto{\pgfqpoint{0.758045in}{0.374400in}}%
\pgfpathlineto{\pgfqpoint{0.756787in}{0.372685in}}%
\pgfpathlineto{\pgfqpoint{0.753894in}{0.374816in}}%
\pgfpathlineto{\pgfqpoint{0.751397in}{0.371451in}}%
\pgfpathlineto{\pgfqpoint{0.748076in}{0.373959in}}%
\pgfpathlineto{\pgfqpoint{0.745512in}{0.370591in}}%
\pgfpathlineto{\pgfqpoint{0.746018in}{0.370169in}}%
\pgfpathlineto{\pgfqpoint{0.741011in}{0.363689in}}%
\pgfpathlineto{\pgfqpoint{0.740274in}{0.361603in}}%
\pgfpathlineto{\pgfqpoint{0.743549in}{0.359090in}}%
\pgfpathlineto{\pgfqpoint{0.740865in}{0.355528in}}%
\pgfpathlineto{\pgfqpoint{0.738983in}{0.355722in}}%
\pgfpathlineto{\pgfqpoint{0.737813in}{0.354884in}}%
\pgfpathlineto{\pgfqpoint{0.734357in}{0.355834in}}%
\pgfpathlineto{\pgfqpoint{0.731400in}{0.357708in}}%
\pgfpathlineto{\pgfqpoint{0.730107in}{0.357059in}}%
\pgfpathlineto{\pgfqpoint{0.728487in}{0.358327in}}%
\pgfpathlineto{\pgfqpoint{0.726759in}{0.358371in}}%
\pgfpathlineto{\pgfqpoint{0.722429in}{0.352828in}}%
\pgfpathlineto{\pgfqpoint{0.719341in}{0.353168in}}%
\pgfpathlineto{\pgfqpoint{0.717812in}{0.354377in}}%
\pgfpathlineto{\pgfqpoint{0.715201in}{0.352581in}}%
\pgfpathlineto{\pgfqpoint{0.712047in}{0.351741in}}%
\pgfpathlineto{\pgfqpoint{0.711910in}{0.351797in}}%
\pgfpathlineto{\pgfqpoint{0.712359in}{0.352736in}}%
\pgfpathlineto{\pgfqpoint{0.708693in}{0.353375in}}%
\pgfpathlineto{\pgfqpoint{0.707132in}{0.354637in}}%
\pgfpathlineto{\pgfqpoint{0.704495in}{0.354085in}}%
\pgfpathlineto{\pgfqpoint{0.703476in}{0.352730in}}%
\pgfpathlineto{\pgfqpoint{0.701818in}{0.353160in}}%
\pgfpathlineto{\pgfqpoint{0.698993in}{0.352386in}}%
\pgfpathlineto{\pgfqpoint{0.697922in}{0.351049in}}%
\pgfpathlineto{\pgfqpoint{0.701239in}{0.348319in}}%
\pgfpathlineto{\pgfqpoint{0.697491in}{0.346230in}}%
\pgfpathlineto{\pgfqpoint{0.696501in}{0.348097in}}%
\pgfpathlineto{\pgfqpoint{0.693475in}{0.346785in}}%
\pgfpathlineto{\pgfqpoint{0.692417in}{0.347973in}}%
\pgfpathlineto{\pgfqpoint{0.683563in}{0.349485in}}%
\pgfpathlineto{\pgfqpoint{0.682454in}{0.347098in}}%
\pgfpathlineto{\pgfqpoint{0.680426in}{0.348551in}}%
\pgfpathlineto{\pgfqpoint{0.680359in}{0.350534in}}%
\pgfpathlineto{\pgfqpoint{0.678074in}{0.351002in}}%
\pgfpathlineto{\pgfqpoint{0.677768in}{0.349141in}}%
\pgfpathlineto{\pgfqpoint{0.676340in}{0.347791in}}%
\pgfpathlineto{\pgfqpoint{0.672925in}{0.348867in}}%
\pgfpathlineto{\pgfqpoint{0.670574in}{0.351024in}}%
\pgfpathlineto{\pgfqpoint{0.667546in}{0.349762in}}%
\pgfpathlineto{\pgfqpoint{0.669799in}{0.346596in}}%
\pgfpathlineto{\pgfqpoint{0.664976in}{0.345492in}}%
\pgfpathlineto{\pgfqpoint{0.662946in}{0.343281in}}%
\pgfpathlineto{\pgfqpoint{0.662231in}{0.345505in}}%
\pgfpathlineto{\pgfqpoint{0.659923in}{0.345201in}}%
\pgfpathlineto{\pgfqpoint{0.658186in}{0.346446in}}%
\pgfpathlineto{\pgfqpoint{0.656689in}{0.346479in}}%
\pgfpathlineto{\pgfqpoint{0.652000in}{0.348708in}}%
\pgfpathlineto{\pgfqpoint{0.649972in}{0.347731in}}%
\pgfpathlineto{\pgfqpoint{0.648426in}{0.347772in}}%
\pgfpathlineto{\pgfqpoint{0.646196in}{0.346313in}}%
\pgfpathlineto{\pgfqpoint{0.650028in}{0.350774in}}%
\pgfpathlineto{\pgfqpoint{0.647002in}{0.353387in}}%
\pgfpathlineto{\pgfqpoint{0.649075in}{0.354518in}}%
\pgfpathlineto{\pgfqpoint{0.654252in}{0.360591in}}%
\pgfpathlineto{\pgfqpoint{0.653338in}{0.361411in}}%
\pgfpathlineto{\pgfqpoint{0.655995in}{0.364392in}}%
\pgfpathlineto{\pgfqpoint{0.659049in}{0.361833in}}%
\pgfpathlineto{\pgfqpoint{0.660386in}{0.363372in}}%
\pgfpathlineto{\pgfqpoint{0.663515in}{0.360688in}}%
\pgfpathlineto{\pgfqpoint{0.664908in}{0.362318in}}%
\pgfpathlineto{\pgfqpoint{0.667310in}{0.360198in}}%
\pgfpathlineto{\pgfqpoint{0.668605in}{0.361721in}}%
\pgfpathlineto{\pgfqpoint{0.670704in}{0.361448in}}%
\pgfpathlineto{\pgfqpoint{0.672992in}{0.359286in}}%
\pgfpathlineto{\pgfqpoint{0.675516in}{0.363060in}}%
\pgfpathlineto{\pgfqpoint{0.680318in}{0.364924in}}%
\pgfpathlineto{\pgfqpoint{0.685284in}{0.365748in}}%
\pgfpathlineto{\pgfqpoint{0.687937in}{0.364755in}}%
\pgfpathlineto{\pgfqpoint{0.690707in}{0.366122in}}%
\pgfpathlineto{\pgfqpoint{0.694569in}{0.365827in}}%
\pgfpathlineto{\pgfqpoint{0.695772in}{0.367709in}}%
\pgfpathlineto{\pgfqpoint{0.704688in}{0.375601in}}%
\pgfpathlineto{\pgfqpoint{0.706614in}{0.375694in}}%
\pgfpathlineto{\pgfqpoint{0.706494in}{0.377302in}}%
\pgfpathlineto{\pgfqpoint{0.709397in}{0.380342in}}%
\pgfpathlineto{\pgfqpoint{0.713956in}{0.381377in}}%
\pgfpathlineto{\pgfqpoint{0.723393in}{0.373987in}}%
\pgfpathlineto{\pgfqpoint{0.727658in}{0.379320in}}%
\pgfpathlineto{\pgfqpoint{0.720329in}{0.385044in}}%
\pgfpathlineto{\pgfqpoint{0.707689in}{0.387813in}}%
\pgfpathlineto{\pgfqpoint{0.705303in}{0.388793in}}%
\pgfpathlineto{\pgfqpoint{0.704791in}{0.390712in}}%
\pgfpathlineto{\pgfqpoint{0.705523in}{0.393336in}}%
\pgfpathlineto{\pgfqpoint{0.704478in}{0.394497in}}%
\pgfpathlineto{\pgfqpoint{0.707497in}{0.397391in}}%
\pgfpathlineto{\pgfqpoint{0.706468in}{0.397884in}}%
\pgfpathlineto{\pgfqpoint{0.700476in}{0.396199in}}%
\pgfpathlineto{\pgfqpoint{0.698526in}{0.391552in}}%
\pgfpathlineto{\pgfqpoint{0.697132in}{0.389984in}}%
\pgfpathlineto{\pgfqpoint{0.695255in}{0.390285in}}%
\pgfpathlineto{\pgfqpoint{0.694284in}{0.392577in}}%
\pgfpathlineto{\pgfqpoint{0.695108in}{0.393300in}}%
\pgfpathlineto{\pgfqpoint{0.695145in}{0.396106in}}%
\pgfpathlineto{\pgfqpoint{0.696329in}{0.402560in}}%
\pgfpathlineto{\pgfqpoint{0.695632in}{0.403858in}}%
\pgfpathlineto{\pgfqpoint{0.697277in}{0.406002in}}%
\pgfpathlineto{\pgfqpoint{0.695522in}{0.406874in}}%
\pgfpathlineto{\pgfqpoint{0.694538in}{0.405951in}}%
\pgfpathlineto{\pgfqpoint{0.691532in}{0.405166in}}%
\pgfpathlineto{\pgfqpoint{0.691789in}{0.407114in}}%
\pgfpathlineto{\pgfqpoint{0.690598in}{0.411229in}}%
\pgfpathlineto{\pgfqpoint{0.692288in}{0.413752in}}%
\pgfpathlineto{\pgfqpoint{0.690110in}{0.414243in}}%
\pgfpathlineto{\pgfqpoint{0.685710in}{0.414422in}}%
\pgfpathlineto{\pgfqpoint{0.681584in}{0.415593in}}%
\pgfpathlineto{\pgfqpoint{0.684613in}{0.418474in}}%
\pgfpathlineto{\pgfqpoint{0.686127in}{0.417159in}}%
\pgfpathlineto{\pgfqpoint{0.687384in}{0.418587in}}%
\pgfpathlineto{\pgfqpoint{0.688938in}{0.417253in}}%
\pgfpathlineto{\pgfqpoint{0.691710in}{0.420379in}}%
\pgfpathlineto{\pgfqpoint{0.693623in}{0.421487in}}%
\pgfpathlineto{\pgfqpoint{0.695145in}{0.420096in}}%
\pgfpathlineto{\pgfqpoint{0.697072in}{0.421245in}}%
\pgfpathlineto{\pgfqpoint{0.699846in}{0.424401in}}%
\pgfpathlineto{\pgfqpoint{0.701434in}{0.423011in}}%
\pgfpathlineto{\pgfqpoint{0.704186in}{0.426137in}}%
\pgfpathlineto{\pgfqpoint{0.706317in}{0.424353in}}%
\pgfpathlineto{\pgfqpoint{0.709062in}{0.427491in}}%
\pgfpathlineto{\pgfqpoint{0.710582in}{0.426191in}}%
\pgfpathlineto{\pgfqpoint{0.713230in}{0.429162in}}%
\pgfpathlineto{\pgfqpoint{0.715184in}{0.430242in}}%
\pgfpathlineto{\pgfqpoint{0.716770in}{0.428913in}}%
\pgfpathlineto{\pgfqpoint{0.719488in}{0.432032in}}%
\pgfpathlineto{\pgfqpoint{0.721076in}{0.430661in}}%
\pgfpathlineto{\pgfqpoint{0.722451in}{0.432254in}}%
\pgfpathlineto{\pgfqpoint{0.724234in}{0.430706in}}%
\pgfpathlineto{\pgfqpoint{0.726974in}{0.433888in}}%
\pgfpathlineto{\pgfqpoint{0.728779in}{0.435116in}}%
\pgfpathlineto{\pgfqpoint{0.730325in}{0.433863in}}%
\pgfpathlineto{\pgfqpoint{0.731767in}{0.435552in}}%
\pgfpathlineto{\pgfqpoint{0.733009in}{0.434561in}}%
\pgfpathlineto{\pgfqpoint{0.734378in}{0.436156in}}%
\pgfpathlineto{\pgfqpoint{0.743447in}{0.428714in}}%
\pgfpathlineto{\pgfqpoint{0.742145in}{0.427074in}}%
\pgfpathlineto{\pgfqpoint{0.743807in}{0.425686in}}%
\pgfpathlineto{\pgfqpoint{0.745138in}{0.427348in}}%
\pgfpathlineto{\pgfqpoint{0.746839in}{0.426097in}}%
\pgfpathlineto{\pgfqpoint{0.748296in}{0.427861in}}%
\pgfpathlineto{\pgfqpoint{0.751507in}{0.425233in}}%
\pgfpathlineto{\pgfqpoint{0.750064in}{0.423440in}}%
\pgfpathlineto{\pgfqpoint{0.761232in}{0.414538in}}%
\pgfpathlineto{\pgfqpoint{0.775612in}{0.403305in}}%
\pgfpathlineto{\pgfqpoint{0.784950in}{0.396276in}}%
\pgfpathclose%
\pgfusepath{fill}%
\end{pgfscope}%
\begin{pgfscope}%
\pgfpathrectangle{\pgfqpoint{0.100000in}{0.100000in}}{\pgfqpoint{3.007045in}{1.925000in}}%
\pgfusepath{clip}%
\pgfsetbuttcap%
\pgfsetmiterjoin%
\definecolor{currentfill}{rgb}{0.093272,0.396878,0.673664}%
\pgfsetfillcolor{currentfill}%
\pgfsetlinewidth{0.000000pt}%
\definecolor{currentstroke}{rgb}{0.000000,0.000000,0.000000}%
\pgfsetstrokecolor{currentstroke}%
\pgfsetstrokeopacity{0.000000}%
\pgfsetdash{}{0pt}%
\pgfpathmoveto{\pgfqpoint{0.677618in}{0.416243in}}%
\pgfpathlineto{\pgfqpoint{0.678710in}{0.418051in}}%
\pgfpathlineto{\pgfqpoint{0.681431in}{0.415610in}}%
\pgfpathlineto{\pgfqpoint{0.677618in}{0.416243in}}%
\pgfpathclose%
\pgfusepath{fill}%
\end{pgfscope}%
\begin{pgfscope}%
\pgfpathrectangle{\pgfqpoint{0.100000in}{0.100000in}}{\pgfqpoint{3.007045in}{1.925000in}}%
\pgfusepath{clip}%
\pgfsetbuttcap%
\pgfsetmiterjoin%
\definecolor{currentfill}{rgb}{0.735871,0.841569,0.923045}%
\pgfsetfillcolor{currentfill}%
\pgfsetlinewidth{0.000000pt}%
\definecolor{currentstroke}{rgb}{0.000000,0.000000,0.000000}%
\pgfsetstrokecolor{currentstroke}%
\pgfsetstrokeopacity{0.000000}%
\pgfsetdash{}{0pt}%
\pgfpathmoveto{\pgfqpoint{0.957302in}{1.649942in}}%
\pgfpathlineto{\pgfqpoint{0.955356in}{1.640352in}}%
\pgfpathlineto{\pgfqpoint{0.959812in}{1.636263in}}%
\pgfpathlineto{\pgfqpoint{0.963238in}{1.636470in}}%
\pgfpathlineto{\pgfqpoint{0.964235in}{1.629990in}}%
\pgfpathlineto{\pgfqpoint{0.960968in}{1.612983in}}%
\pgfpathlineto{\pgfqpoint{0.966618in}{1.611886in}}%
\pgfpathlineto{\pgfqpoint{0.965491in}{1.606172in}}%
\pgfpathlineto{\pgfqpoint{0.972074in}{1.604894in}}%
\pgfpathlineto{\pgfqpoint{0.971016in}{1.593385in}}%
\pgfpathlineto{\pgfqpoint{0.975154in}{1.589701in}}%
\pgfpathlineto{\pgfqpoint{0.988238in}{1.587058in}}%
\pgfpathlineto{\pgfqpoint{0.987365in}{1.582354in}}%
\pgfpathlineto{\pgfqpoint{1.001653in}{1.579735in}}%
\pgfpathlineto{\pgfqpoint{0.998004in}{1.571906in}}%
\pgfpathlineto{\pgfqpoint{0.989150in}{1.572014in}}%
\pgfpathlineto{\pgfqpoint{0.983095in}{1.569716in}}%
\pgfpathlineto{\pgfqpoint{0.981565in}{1.573774in}}%
\pgfpathlineto{\pgfqpoint{0.970525in}{1.573027in}}%
\pgfpathlineto{\pgfqpoint{0.962156in}{1.577626in}}%
\pgfpathlineto{\pgfqpoint{0.958428in}{1.575728in}}%
\pgfpathlineto{\pgfqpoint{0.956111in}{1.570703in}}%
\pgfpathlineto{\pgfqpoint{0.952519in}{1.573584in}}%
\pgfpathlineto{\pgfqpoint{0.938309in}{1.576703in}}%
\pgfpathlineto{\pgfqpoint{0.935411in}{1.569991in}}%
\pgfpathlineto{\pgfqpoint{0.933731in}{1.569068in}}%
\pgfpathlineto{\pgfqpoint{0.927577in}{1.576545in}}%
\pgfpathlineto{\pgfqpoint{0.928393in}{1.581446in}}%
\pgfpathlineto{\pgfqpoint{0.926655in}{1.595511in}}%
\pgfpathlineto{\pgfqpoint{0.923418in}{1.601071in}}%
\pgfpathlineto{\pgfqpoint{0.916159in}{1.603777in}}%
\pgfpathlineto{\pgfqpoint{0.912228in}{1.607641in}}%
\pgfpathlineto{\pgfqpoint{0.914945in}{1.620213in}}%
\pgfpathlineto{\pgfqpoint{0.910956in}{1.624666in}}%
\pgfpathlineto{\pgfqpoint{0.905525in}{1.640271in}}%
\pgfpathlineto{\pgfqpoint{0.906555in}{1.644348in}}%
\pgfpathlineto{\pgfqpoint{0.903714in}{1.653787in}}%
\pgfpathlineto{\pgfqpoint{0.906088in}{1.658804in}}%
\pgfpathlineto{\pgfqpoint{0.901414in}{1.666132in}}%
\pgfpathlineto{\pgfqpoint{0.908010in}{1.668930in}}%
\pgfpathlineto{\pgfqpoint{0.914135in}{1.673457in}}%
\pgfpathlineto{\pgfqpoint{0.921890in}{1.674121in}}%
\pgfpathlineto{\pgfqpoint{0.923792in}{1.677916in}}%
\pgfpathlineto{\pgfqpoint{0.926926in}{1.665686in}}%
\pgfpathlineto{\pgfqpoint{0.936043in}{1.670714in}}%
\pgfpathlineto{\pgfqpoint{0.941856in}{1.670351in}}%
\pgfpathlineto{\pgfqpoint{0.944841in}{1.665499in}}%
\pgfpathlineto{\pgfqpoint{0.955593in}{1.658524in}}%
\pgfpathlineto{\pgfqpoint{0.957302in}{1.649942in}}%
\pgfpathclose%
\pgfusepath{fill}%
\end{pgfscope}%
\begin{pgfscope}%
\pgfpathrectangle{\pgfqpoint{0.100000in}{0.100000in}}{\pgfqpoint{3.007045in}{1.925000in}}%
\pgfusepath{clip}%
\pgfsetbuttcap%
\pgfsetmiterjoin%
\definecolor{currentfill}{rgb}{0.485490,0.718524,0.853426}%
\pgfsetfillcolor{currentfill}%
\pgfsetlinewidth{0.000000pt}%
\definecolor{currentstroke}{rgb}{0.000000,0.000000,0.000000}%
\pgfsetstrokecolor{currentstroke}%
\pgfsetstrokeopacity{0.000000}%
\pgfsetdash{}{0pt}%
\pgfpathmoveto{\pgfqpoint{0.741011in}{1.724091in}}%
\pgfpathlineto{\pgfqpoint{0.771689in}{1.716185in}}%
\pgfpathlineto{\pgfqpoint{0.774980in}{1.703762in}}%
\pgfpathlineto{\pgfqpoint{0.778737in}{1.702487in}}%
\pgfpathlineto{\pgfqpoint{0.784656in}{1.695189in}}%
\pgfpathlineto{\pgfqpoint{0.786067in}{1.686768in}}%
\pgfpathlineto{\pgfqpoint{0.780091in}{1.680597in}}%
\pgfpathlineto{\pgfqpoint{0.770539in}{1.666797in}}%
\pgfpathlineto{\pgfqpoint{0.766636in}{1.659171in}}%
\pgfpathlineto{\pgfqpoint{0.763180in}{1.655715in}}%
\pgfpathlineto{\pgfqpoint{0.741293in}{1.661247in}}%
\pgfpathlineto{\pgfqpoint{0.742633in}{1.666720in}}%
\pgfpathlineto{\pgfqpoint{0.733120in}{1.669073in}}%
\pgfpathlineto{\pgfqpoint{0.732751in}{1.676168in}}%
\pgfpathlineto{\pgfqpoint{0.728377in}{1.682025in}}%
\pgfpathlineto{\pgfqpoint{0.727519in}{1.694386in}}%
\pgfpathlineto{\pgfqpoint{0.729208in}{1.700974in}}%
\pgfpathlineto{\pgfqpoint{0.727679in}{1.706314in}}%
\pgfpathlineto{\pgfqpoint{0.732306in}{1.716868in}}%
\pgfpathlineto{\pgfqpoint{0.722204in}{1.719436in}}%
\pgfpathlineto{\pgfqpoint{0.724313in}{1.728496in}}%
\pgfpathlineto{\pgfqpoint{0.741011in}{1.724091in}}%
\pgfpathclose%
\pgfusepath{fill}%
\end{pgfscope}%
\begin{pgfscope}%
\pgfpathrectangle{\pgfqpoint{0.100000in}{0.100000in}}{\pgfqpoint{3.007045in}{1.925000in}}%
\pgfusepath{clip}%
\pgfsetbuttcap%
\pgfsetmiterjoin%
\definecolor{currentfill}{rgb}{0.351511,0.635848,0.812641}%
\pgfsetfillcolor{currentfill}%
\pgfsetlinewidth{0.000000pt}%
\definecolor{currentstroke}{rgb}{0.000000,0.000000,0.000000}%
\pgfsetstrokecolor{currentstroke}%
\pgfsetstrokeopacity{0.000000}%
\pgfsetdash{}{0pt}%
\pgfpathmoveto{\pgfqpoint{1.439278in}{1.125169in}}%
\pgfpathlineto{\pgfqpoint{1.410528in}{1.127149in}}%
\pgfpathlineto{\pgfqpoint{1.414739in}{1.184340in}}%
\pgfpathlineto{\pgfqpoint{1.446935in}{1.181960in}}%
\pgfpathlineto{\pgfqpoint{1.448195in}{1.181868in}}%
\pgfpathlineto{\pgfqpoint{1.446054in}{1.153418in}}%
\pgfpathlineto{\pgfqpoint{1.441646in}{1.153706in}}%
\pgfpathlineto{\pgfqpoint{1.439278in}{1.125169in}}%
\pgfpathclose%
\pgfusepath{fill}%
\end{pgfscope}%
\begin{pgfscope}%
\pgfpathrectangle{\pgfqpoint{0.100000in}{0.100000in}}{\pgfqpoint{3.007045in}{1.925000in}}%
\pgfusepath{clip}%
\pgfsetbuttcap%
\pgfsetmiterjoin%
\definecolor{currentfill}{rgb}{0.223806,0.537532,0.758431}%
\pgfsetfillcolor{currentfill}%
\pgfsetlinewidth{0.000000pt}%
\definecolor{currentstroke}{rgb}{0.000000,0.000000,0.000000}%
\pgfsetstrokecolor{currentstroke}%
\pgfsetstrokeopacity{0.000000}%
\pgfsetdash{}{0pt}%
\pgfpathmoveto{\pgfqpoint{1.712393in}{0.726751in}}%
\pgfpathlineto{\pgfqpoint{1.712066in}{0.694873in}}%
\pgfpathlineto{\pgfqpoint{1.694814in}{0.693189in}}%
\pgfpathlineto{\pgfqpoint{1.665733in}{0.676821in}}%
\pgfpathlineto{\pgfqpoint{1.663204in}{0.675422in}}%
\pgfpathlineto{\pgfqpoint{1.655346in}{0.689754in}}%
\pgfpathlineto{\pgfqpoint{1.655743in}{0.708579in}}%
\pgfpathlineto{\pgfqpoint{1.658448in}{0.708465in}}%
\pgfpathlineto{\pgfqpoint{1.659429in}{0.737338in}}%
\pgfpathlineto{\pgfqpoint{1.639106in}{0.738203in}}%
\pgfpathlineto{\pgfqpoint{1.640734in}{0.767480in}}%
\pgfpathlineto{\pgfqpoint{1.634992in}{0.767696in}}%
\pgfpathlineto{\pgfqpoint{1.636022in}{0.799536in}}%
\pgfpathlineto{\pgfqpoint{1.638052in}{0.793421in}}%
\pgfpathlineto{\pgfqpoint{1.641958in}{0.793032in}}%
\pgfpathlineto{\pgfqpoint{1.649125in}{0.798383in}}%
\pgfpathlineto{\pgfqpoint{1.652538in}{0.797937in}}%
\pgfpathlineto{\pgfqpoint{1.651130in}{0.791612in}}%
\pgfpathlineto{\pgfqpoint{1.653822in}{0.786416in}}%
\pgfpathlineto{\pgfqpoint{1.656967in}{0.786599in}}%
\pgfpathlineto{\pgfqpoint{1.665670in}{0.801052in}}%
\pgfpathlineto{\pgfqpoint{1.664958in}{0.765812in}}%
\pgfpathlineto{\pgfqpoint{1.670614in}{0.764978in}}%
\pgfpathlineto{\pgfqpoint{1.695830in}{0.764016in}}%
\pgfpathlineto{\pgfqpoint{1.700709in}{0.760907in}}%
\pgfpathlineto{\pgfqpoint{1.700123in}{0.727126in}}%
\pgfpathlineto{\pgfqpoint{1.712393in}{0.726751in}}%
\pgfpathclose%
\pgfusepath{fill}%
\end{pgfscope}%
\begin{pgfscope}%
\pgfpathrectangle{\pgfqpoint{0.100000in}{0.100000in}}{\pgfqpoint{3.007045in}{1.925000in}}%
\pgfusepath{clip}%
\pgfsetbuttcap%
\pgfsetmiterjoin%
\definecolor{currentfill}{rgb}{0.280892,0.587620,0.785083}%
\pgfsetfillcolor{currentfill}%
\pgfsetlinewidth{0.000000pt}%
\definecolor{currentstroke}{rgb}{0.000000,0.000000,0.000000}%
\pgfsetstrokecolor{currentstroke}%
\pgfsetstrokeopacity{0.000000}%
\pgfsetdash{}{0pt}%
\pgfpathmoveto{\pgfqpoint{1.665670in}{0.801052in}}%
\pgfpathlineto{\pgfqpoint{1.656967in}{0.786599in}}%
\pgfpathlineto{\pgfqpoint{1.653822in}{0.786416in}}%
\pgfpathlineto{\pgfqpoint{1.651130in}{0.791612in}}%
\pgfpathlineto{\pgfqpoint{1.652538in}{0.797937in}}%
\pgfpathlineto{\pgfqpoint{1.649125in}{0.798383in}}%
\pgfpathlineto{\pgfqpoint{1.641958in}{0.793032in}}%
\pgfpathlineto{\pgfqpoint{1.638052in}{0.793421in}}%
\pgfpathlineto{\pgfqpoint{1.636022in}{0.799536in}}%
\pgfpathlineto{\pgfqpoint{1.631785in}{0.798613in}}%
\pgfpathlineto{\pgfqpoint{1.633138in}{0.850249in}}%
\pgfpathlineto{\pgfqpoint{1.627401in}{0.850427in}}%
\pgfpathlineto{\pgfqpoint{1.627761in}{0.861901in}}%
\pgfpathlineto{\pgfqpoint{1.667600in}{0.860848in}}%
\pgfpathlineto{\pgfqpoint{1.667244in}{0.843587in}}%
\pgfpathlineto{\pgfqpoint{1.672960in}{0.843473in}}%
\pgfpathlineto{\pgfqpoint{1.672844in}{0.837714in}}%
\pgfpathlineto{\pgfqpoint{1.689906in}{0.837361in}}%
\pgfpathlineto{\pgfqpoint{1.689815in}{0.831695in}}%
\pgfpathlineto{\pgfqpoint{1.695616in}{0.831530in}}%
\pgfpathlineto{\pgfqpoint{1.695332in}{0.814267in}}%
\pgfpathlineto{\pgfqpoint{1.683354in}{0.804224in}}%
\pgfpathlineto{\pgfqpoint{1.678257in}{0.793022in}}%
\pgfpathlineto{\pgfqpoint{1.668937in}{0.795613in}}%
\pgfpathlineto{\pgfqpoint{1.665670in}{0.801052in}}%
\pgfpathclose%
\pgfusepath{fill}%
\end{pgfscope}%
\begin{pgfscope}%
\pgfpathrectangle{\pgfqpoint{0.100000in}{0.100000in}}{\pgfqpoint{3.007045in}{1.925000in}}%
\pgfusepath{clip}%
\pgfsetbuttcap%
\pgfsetmiterjoin%
\definecolor{currentfill}{rgb}{0.485490,0.718524,0.853426}%
\pgfsetfillcolor{currentfill}%
\pgfsetlinewidth{0.000000pt}%
\definecolor{currentstroke}{rgb}{0.000000,0.000000,0.000000}%
\pgfsetstrokecolor{currentstroke}%
\pgfsetstrokeopacity{0.000000}%
\pgfsetdash{}{0pt}%
\pgfpathmoveto{\pgfqpoint{1.146247in}{0.808107in}}%
\pgfpathlineto{\pgfqpoint{1.147002in}{0.802056in}}%
\pgfpathlineto{\pgfqpoint{1.144512in}{0.779815in}}%
\pgfpathlineto{\pgfqpoint{1.112390in}{0.768586in}}%
\pgfpathlineto{\pgfqpoint{1.112021in}{0.765725in}}%
\pgfpathlineto{\pgfqpoint{1.089302in}{0.768831in}}%
\pgfpathlineto{\pgfqpoint{1.087735in}{0.757358in}}%
\pgfpathlineto{\pgfqpoint{1.070654in}{0.759703in}}%
\pgfpathlineto{\pgfqpoint{1.063758in}{0.762298in}}%
\pgfpathlineto{\pgfqpoint{1.065981in}{0.771962in}}%
\pgfpathlineto{\pgfqpoint{1.063832in}{0.777370in}}%
\pgfpathlineto{\pgfqpoint{1.064806in}{0.782977in}}%
\pgfpathlineto{\pgfqpoint{1.061362in}{0.785689in}}%
\pgfpathlineto{\pgfqpoint{1.062578in}{0.792196in}}%
\pgfpathlineto{\pgfqpoint{1.061462in}{0.800953in}}%
\pgfpathlineto{\pgfqpoint{1.054598in}{0.801868in}}%
\pgfpathlineto{\pgfqpoint{1.057206in}{0.819980in}}%
\pgfpathlineto{\pgfqpoint{1.060540in}{0.819258in}}%
\pgfpathlineto{\pgfqpoint{1.115027in}{0.811851in}}%
\pgfpathlineto{\pgfqpoint{1.146247in}{0.808107in}}%
\pgfpathclose%
\pgfusepath{fill}%
\end{pgfscope}%
\begin{pgfscope}%
\pgfpathrectangle{\pgfqpoint{0.100000in}{0.100000in}}{\pgfqpoint{3.007045in}{1.925000in}}%
\pgfusepath{clip}%
\pgfsetbuttcap%
\pgfsetmiterjoin%
\definecolor{currentfill}{rgb}{0.396909,0.666851,0.830358}%
\pgfsetfillcolor{currentfill}%
\pgfsetlinewidth{0.000000pt}%
\definecolor{currentstroke}{rgb}{0.000000,0.000000,0.000000}%
\pgfsetstrokecolor{currentstroke}%
\pgfsetstrokeopacity{0.000000}%
\pgfsetdash{}{0pt}%
\pgfpathmoveto{\pgfqpoint{1.718799in}{1.643280in}}%
\pgfpathlineto{\pgfqpoint{1.718828in}{1.637478in}}%
\pgfpathlineto{\pgfqpoint{1.713868in}{1.637553in}}%
\pgfpathlineto{\pgfqpoint{1.713584in}{1.614400in}}%
\pgfpathlineto{\pgfqpoint{1.714002in}{1.597084in}}%
\pgfpathlineto{\pgfqpoint{1.699675in}{1.597325in}}%
\pgfpathlineto{\pgfqpoint{1.700191in}{1.591568in}}%
\pgfpathlineto{\pgfqpoint{1.669507in}{1.592201in}}%
\pgfpathlineto{\pgfqpoint{1.668146in}{1.596346in}}%
\pgfpathlineto{\pgfqpoint{1.667675in}{1.615294in}}%
\pgfpathlineto{\pgfqpoint{1.668280in}{1.638340in}}%
\pgfpathlineto{\pgfqpoint{1.649938in}{1.638890in}}%
\pgfpathlineto{\pgfqpoint{1.650635in}{1.662065in}}%
\pgfpathlineto{\pgfqpoint{1.650048in}{1.679481in}}%
\pgfpathlineto{\pgfqpoint{1.689494in}{1.678334in}}%
\pgfpathlineto{\pgfqpoint{1.689129in}{1.695752in}}%
\pgfpathlineto{\pgfqpoint{1.724288in}{1.695202in}}%
\pgfpathlineto{\pgfqpoint{1.724037in}{1.672065in}}%
\pgfpathlineto{\pgfqpoint{1.718261in}{1.672097in}}%
\pgfpathlineto{\pgfqpoint{1.719048in}{1.660538in}}%
\pgfpathlineto{\pgfqpoint{1.718799in}{1.643280in}}%
\pgfpathclose%
\pgfusepath{fill}%
\end{pgfscope}%
\begin{pgfscope}%
\pgfpathrectangle{\pgfqpoint{0.100000in}{0.100000in}}{\pgfqpoint{3.007045in}{1.925000in}}%
\pgfusepath{clip}%
\pgfsetbuttcap%
\pgfsetmiterjoin%
\definecolor{currentfill}{rgb}{0.784591,0.864237,0.939962}%
\pgfsetfillcolor{currentfill}%
\pgfsetlinewidth{0.000000pt}%
\definecolor{currentstroke}{rgb}{0.000000,0.000000,0.000000}%
\pgfsetstrokecolor{currentstroke}%
\pgfsetstrokeopacity{0.000000}%
\pgfsetdash{}{0pt}%
\pgfpathmoveto{\pgfqpoint{2.298344in}{1.505689in}}%
\pgfpathlineto{\pgfqpoint{2.264263in}{1.501550in}}%
\pgfpathlineto{\pgfqpoint{2.252766in}{1.500298in}}%
\pgfpathlineto{\pgfqpoint{2.250233in}{1.523286in}}%
\pgfpathlineto{\pgfqpoint{2.247824in}{1.546163in}}%
\pgfpathlineto{\pgfqpoint{2.245496in}{1.563087in}}%
\pgfpathlineto{\pgfqpoint{2.251176in}{1.563619in}}%
\pgfpathlineto{\pgfqpoint{2.250552in}{1.569368in}}%
\pgfpathlineto{\pgfqpoint{2.273175in}{1.571785in}}%
\pgfpathlineto{\pgfqpoint{2.290081in}{1.574045in}}%
\pgfpathlineto{\pgfqpoint{2.291082in}{1.562710in}}%
\pgfpathlineto{\pgfqpoint{2.298344in}{1.505689in}}%
\pgfpathclose%
\pgfusepath{fill}%
\end{pgfscope}%
\begin{pgfscope}%
\pgfpathrectangle{\pgfqpoint{0.100000in}{0.100000in}}{\pgfqpoint{3.007045in}{1.925000in}}%
\pgfusepath{clip}%
\pgfsetbuttcap%
\pgfsetmiterjoin%
\definecolor{currentfill}{rgb}{0.867266,0.919354,0.967520}%
\pgfsetfillcolor{currentfill}%
\pgfsetlinewidth{0.000000pt}%
\definecolor{currentstroke}{rgb}{0.000000,0.000000,0.000000}%
\pgfsetstrokecolor{currentstroke}%
\pgfsetstrokeopacity{0.000000}%
\pgfsetdash{}{0pt}%
\pgfpathmoveto{\pgfqpoint{2.574578in}{0.398555in}}%
\pgfpathlineto{\pgfqpoint{2.577000in}{0.380965in}}%
\pgfpathlineto{\pgfqpoint{2.547508in}{0.376728in}}%
\pgfpathlineto{\pgfqpoint{2.547051in}{0.383388in}}%
\pgfpathlineto{\pgfqpoint{2.542240in}{0.387941in}}%
\pgfpathlineto{\pgfqpoint{2.539505in}{0.384621in}}%
\pgfpathlineto{\pgfqpoint{2.542207in}{0.377186in}}%
\pgfpathlineto{\pgfqpoint{2.536927in}{0.376168in}}%
\pgfpathlineto{\pgfqpoint{2.534833in}{0.376193in}}%
\pgfpathlineto{\pgfqpoint{2.521960in}{0.392832in}}%
\pgfpathlineto{\pgfqpoint{2.512567in}{0.407613in}}%
\pgfpathlineto{\pgfqpoint{2.503920in}{0.416388in}}%
\pgfpathlineto{\pgfqpoint{2.505714in}{0.422690in}}%
\pgfpathlineto{\pgfqpoint{2.510105in}{0.430828in}}%
\pgfpathlineto{\pgfqpoint{2.539586in}{0.435076in}}%
\pgfpathlineto{\pgfqpoint{2.542441in}{0.414576in}}%
\pgfpathlineto{\pgfqpoint{2.571480in}{0.418988in}}%
\pgfpathlineto{\pgfqpoint{2.574578in}{0.398555in}}%
\pgfpathclose%
\pgfusepath{fill}%
\end{pgfscope}%
\begin{pgfscope}%
\pgfpathrectangle{\pgfqpoint{0.100000in}{0.100000in}}{\pgfqpoint{3.007045in}{1.925000in}}%
\pgfusepath{clip}%
\pgfsetbuttcap%
\pgfsetmiterjoin%
\definecolor{currentfill}{rgb}{0.686659,0.820654,0.905821}%
\pgfsetfillcolor{currentfill}%
\pgfsetlinewidth{0.000000pt}%
\definecolor{currentstroke}{rgb}{0.000000,0.000000,0.000000}%
\pgfsetstrokecolor{currentstroke}%
\pgfsetstrokeopacity{0.000000}%
\pgfsetdash{}{0pt}%
\pgfpathmoveto{\pgfqpoint{2.555437in}{1.355961in}}%
\pgfpathlineto{\pgfqpoint{2.554563in}{1.361326in}}%
\pgfpathlineto{\pgfqpoint{2.559925in}{1.364184in}}%
\pgfpathlineto{\pgfqpoint{2.558934in}{1.375533in}}%
\pgfpathlineto{\pgfqpoint{2.581784in}{1.380075in}}%
\pgfpathlineto{\pgfqpoint{2.595492in}{1.381828in}}%
\pgfpathlineto{\pgfqpoint{2.604647in}{1.373975in}}%
\pgfpathlineto{\pgfqpoint{2.607627in}{1.374496in}}%
\pgfpathlineto{\pgfqpoint{2.608964in}{1.367546in}}%
\pgfpathlineto{\pgfqpoint{2.605715in}{1.356835in}}%
\pgfpathlineto{\pgfqpoint{2.573880in}{1.354042in}}%
\pgfpathlineto{\pgfqpoint{2.561808in}{1.358520in}}%
\pgfpathlineto{\pgfqpoint{2.555437in}{1.355961in}}%
\pgfpathclose%
\pgfusepath{fill}%
\end{pgfscope}%
\begin{pgfscope}%
\pgfpathrectangle{\pgfqpoint{0.100000in}{0.100000in}}{\pgfqpoint{3.007045in}{1.925000in}}%
\pgfusepath{clip}%
\pgfsetbuttcap%
\pgfsetmiterjoin%
\definecolor{currentfill}{rgb}{0.516863,0.735748,0.860192}%
\pgfsetfillcolor{currentfill}%
\pgfsetlinewidth{0.000000pt}%
\definecolor{currentstroke}{rgb}{0.000000,0.000000,0.000000}%
\pgfsetstrokecolor{currentstroke}%
\pgfsetstrokeopacity{0.000000}%
\pgfsetdash{}{0pt}%
\pgfpathmoveto{\pgfqpoint{2.169528in}{1.143936in}}%
\pgfpathlineto{\pgfqpoint{2.155147in}{1.142366in}}%
\pgfpathlineto{\pgfqpoint{2.152006in}{1.147645in}}%
\pgfpathlineto{\pgfqpoint{2.151777in}{1.152326in}}%
\pgfpathlineto{\pgfqpoint{2.148381in}{1.158667in}}%
\pgfpathlineto{\pgfqpoint{2.150340in}{1.161666in}}%
\pgfpathlineto{\pgfqpoint{2.148720in}{1.168849in}}%
\pgfpathlineto{\pgfqpoint{2.152293in}{1.171574in}}%
\pgfpathlineto{\pgfqpoint{2.149319in}{1.206705in}}%
\pgfpathlineto{\pgfqpoint{2.147976in}{1.224103in}}%
\pgfpathlineto{\pgfqpoint{2.154351in}{1.223276in}}%
\pgfpathlineto{\pgfqpoint{2.154631in}{1.211728in}}%
\pgfpathlineto{\pgfqpoint{2.171191in}{1.213179in}}%
\pgfpathlineto{\pgfqpoint{2.169632in}{1.230355in}}%
\pgfpathlineto{\pgfqpoint{2.189499in}{1.232103in}}%
\pgfpathlineto{\pgfqpoint{2.189712in}{1.229760in}}%
\pgfpathlineto{\pgfqpoint{2.193431in}{1.193828in}}%
\pgfpathlineto{\pgfqpoint{2.195596in}{1.188844in}}%
\pgfpathlineto{\pgfqpoint{2.194423in}{1.183326in}}%
\pgfpathlineto{\pgfqpoint{2.181588in}{1.182377in}}%
\pgfpathlineto{\pgfqpoint{2.182227in}{1.173751in}}%
\pgfpathlineto{\pgfqpoint{2.176581in}{1.173291in}}%
\pgfpathlineto{\pgfqpoint{2.177539in}{1.161835in}}%
\pgfpathlineto{\pgfqpoint{2.168048in}{1.161256in}}%
\pgfpathlineto{\pgfqpoint{2.169528in}{1.143936in}}%
\pgfpathclose%
\pgfusepath{fill}%
\end{pgfscope}%
\begin{pgfscope}%
\pgfpathrectangle{\pgfqpoint{0.100000in}{0.100000in}}{\pgfqpoint{3.007045in}{1.925000in}}%
\pgfusepath{clip}%
\pgfsetbuttcap%
\pgfsetmiterjoin%
\definecolor{currentfill}{rgb}{0.341423,0.628958,0.808704}%
\pgfsetfillcolor{currentfill}%
\pgfsetlinewidth{0.000000pt}%
\definecolor{currentstroke}{rgb}{0.000000,0.000000,0.000000}%
\pgfsetstrokecolor{currentstroke}%
\pgfsetstrokeopacity{0.000000}%
\pgfsetdash{}{0pt}%
\pgfpathmoveto{\pgfqpoint{1.333387in}{1.392952in}}%
\pgfpathlineto{\pgfqpoint{1.326529in}{1.324213in}}%
\pgfpathlineto{\pgfqpoint{1.297072in}{1.327304in}}%
\pgfpathlineto{\pgfqpoint{1.297671in}{1.333074in}}%
\pgfpathlineto{\pgfqpoint{1.267098in}{1.336682in}}%
\pgfpathlineto{\pgfqpoint{1.269866in}{1.358232in}}%
\pgfpathlineto{\pgfqpoint{1.272590in}{1.387487in}}%
\pgfpathlineto{\pgfqpoint{1.273764in}{1.398937in}}%
\pgfpathlineto{\pgfqpoint{1.292810in}{1.397021in}}%
\pgfpathlineto{\pgfqpoint{1.333387in}{1.392952in}}%
\pgfpathclose%
\pgfusepath{fill}%
\end{pgfscope}%
\begin{pgfscope}%
\pgfpathrectangle{\pgfqpoint{0.100000in}{0.100000in}}{\pgfqpoint{3.007045in}{1.925000in}}%
\pgfusepath{clip}%
\pgfsetbuttcap%
\pgfsetmiterjoin%
\definecolor{currentfill}{rgb}{0.154787,0.468512,0.722876}%
\pgfsetfillcolor{currentfill}%
\pgfsetlinewidth{0.000000pt}%
\definecolor{currentstroke}{rgb}{0.000000,0.000000,0.000000}%
\pgfsetstrokecolor{currentstroke}%
\pgfsetstrokeopacity{0.000000}%
\pgfsetdash{}{0pt}%
\pgfpathmoveto{\pgfqpoint{1.392891in}{0.867332in}}%
\pgfpathlineto{\pgfqpoint{1.421405in}{0.865168in}}%
\pgfpathlineto{\pgfqpoint{1.419358in}{0.836488in}}%
\pgfpathlineto{\pgfqpoint{1.442781in}{0.834883in}}%
\pgfpathlineto{\pgfqpoint{1.440696in}{0.803294in}}%
\pgfpathlineto{\pgfqpoint{1.383417in}{0.806844in}}%
\pgfpathlineto{\pgfqpoint{1.385676in}{0.839061in}}%
\pgfpathlineto{\pgfqpoint{1.390704in}{0.838664in}}%
\pgfpathlineto{\pgfqpoint{1.392891in}{0.867332in}}%
\pgfpathclose%
\pgfusepath{fill}%
\end{pgfscope}%
\begin{pgfscope}%
\pgfpathrectangle{\pgfqpoint{0.100000in}{0.100000in}}{\pgfqpoint{3.007045in}{1.925000in}}%
\pgfusepath{clip}%
\pgfsetbuttcap%
\pgfsetmiterjoin%
\definecolor{currentfill}{rgb}{0.637447,0.799739,0.888597}%
\pgfsetfillcolor{currentfill}%
\pgfsetlinewidth{0.000000pt}%
\definecolor{currentstroke}{rgb}{0.000000,0.000000,0.000000}%
\pgfsetstrokecolor{currentstroke}%
\pgfsetstrokeopacity{0.000000}%
\pgfsetdash{}{0pt}%
\pgfpathmoveto{\pgfqpoint{1.180201in}{1.170746in}}%
\pgfpathlineto{\pgfqpoint{1.177735in}{1.174146in}}%
\pgfpathlineto{\pgfqpoint{1.172957in}{1.174640in}}%
\pgfpathlineto{\pgfqpoint{1.169868in}{1.170865in}}%
\pgfpathlineto{\pgfqpoint{1.164471in}{1.172432in}}%
\pgfpathlineto{\pgfqpoint{1.157547in}{1.181738in}}%
\pgfpathlineto{\pgfqpoint{1.146599in}{1.183208in}}%
\pgfpathlineto{\pgfqpoint{1.143492in}{1.187243in}}%
\pgfpathlineto{\pgfqpoint{1.142160in}{1.192960in}}%
\pgfpathlineto{\pgfqpoint{1.139110in}{1.197239in}}%
\pgfpathlineto{\pgfqpoint{1.141305in}{1.200383in}}%
\pgfpathlineto{\pgfqpoint{1.091973in}{1.207431in}}%
\pgfpathlineto{\pgfqpoint{1.059704in}{1.212390in}}%
\pgfpathlineto{\pgfqpoint{1.061002in}{1.220921in}}%
\pgfpathlineto{\pgfqpoint{1.062663in}{1.231515in}}%
\pgfpathlineto{\pgfqpoint{1.087451in}{1.227034in}}%
\pgfpathlineto{\pgfqpoint{1.087872in}{1.229867in}}%
\pgfpathlineto{\pgfqpoint{1.118778in}{1.225384in}}%
\pgfpathlineto{\pgfqpoint{1.120000in}{1.233895in}}%
\pgfpathlineto{\pgfqpoint{1.145415in}{1.230549in}}%
\pgfpathlineto{\pgfqpoint{1.146223in}{1.236379in}}%
\pgfpathlineto{\pgfqpoint{1.151850in}{1.235425in}}%
\pgfpathlineto{\pgfqpoint{1.153503in}{1.246887in}}%
\pgfpathlineto{\pgfqpoint{1.167476in}{1.245055in}}%
\pgfpathlineto{\pgfqpoint{1.166255in}{1.233746in}}%
\pgfpathlineto{\pgfqpoint{1.186601in}{1.231472in}}%
\pgfpathlineto{\pgfqpoint{1.205530in}{1.228275in}}%
\pgfpathlineto{\pgfqpoint{1.212843in}{1.220333in}}%
\pgfpathlineto{\pgfqpoint{1.214794in}{1.212322in}}%
\pgfpathlineto{\pgfqpoint{1.220004in}{1.212317in}}%
\pgfpathlineto{\pgfqpoint{1.226581in}{1.206947in}}%
\pgfpathlineto{\pgfqpoint{1.216380in}{1.193870in}}%
\pgfpathlineto{\pgfqpoint{1.203807in}{1.188442in}}%
\pgfpathlineto{\pgfqpoint{1.202987in}{1.177221in}}%
\pgfpathlineto{\pgfqpoint{1.201403in}{1.171929in}}%
\pgfpathlineto{\pgfqpoint{1.181782in}{1.174522in}}%
\pgfpathlineto{\pgfqpoint{1.180201in}{1.170746in}}%
\pgfpathclose%
\pgfusepath{fill}%
\end{pgfscope}%
\begin{pgfscope}%
\pgfpathrectangle{\pgfqpoint{0.100000in}{0.100000in}}{\pgfqpoint{3.007045in}{1.925000in}}%
\pgfusepath{clip}%
\pgfsetbuttcap%
\pgfsetmiterjoin%
\definecolor{currentfill}{rgb}{0.435294,0.690965,0.842599}%
\pgfsetfillcolor{currentfill}%
\pgfsetlinewidth{0.000000pt}%
\definecolor{currentstroke}{rgb}{0.000000,0.000000,0.000000}%
\pgfsetstrokecolor{currentstroke}%
\pgfsetstrokeopacity{0.000000}%
\pgfsetdash{}{0pt}%
\pgfpathmoveto{\pgfqpoint{2.593275in}{0.442838in}}%
\pgfpathlineto{\pgfqpoint{2.590385in}{0.434304in}}%
\pgfpathlineto{\pgfqpoint{2.597645in}{0.425944in}}%
\pgfpathlineto{\pgfqpoint{2.600616in}{0.426095in}}%
\pgfpathlineto{\pgfqpoint{2.607766in}{0.416135in}}%
\pgfpathlineto{\pgfqpoint{2.596119in}{0.413805in}}%
\pgfpathlineto{\pgfqpoint{2.597054in}{0.407964in}}%
\pgfpathlineto{\pgfqpoint{2.591178in}{0.407023in}}%
\pgfpathlineto{\pgfqpoint{2.592108in}{0.401133in}}%
\pgfpathlineto{\pgfqpoint{2.574578in}{0.398555in}}%
\pgfpathlineto{\pgfqpoint{2.571480in}{0.418988in}}%
\pgfpathlineto{\pgfqpoint{2.542441in}{0.414576in}}%
\pgfpathlineto{\pgfqpoint{2.539586in}{0.435076in}}%
\pgfpathlineto{\pgfqpoint{2.534472in}{0.469934in}}%
\pgfpathlineto{\pgfqpoint{2.530722in}{0.475320in}}%
\pgfpathlineto{\pgfqpoint{2.533154in}{0.479245in}}%
\pgfpathlineto{\pgfqpoint{2.538576in}{0.482276in}}%
\pgfpathlineto{\pgfqpoint{2.556078in}{0.484967in}}%
\pgfpathlineto{\pgfqpoint{2.556944in}{0.479141in}}%
\pgfpathlineto{\pgfqpoint{2.562765in}{0.480015in}}%
\pgfpathlineto{\pgfqpoint{2.565859in}{0.472579in}}%
\pgfpathlineto{\pgfqpoint{2.575198in}{0.465448in}}%
\pgfpathlineto{\pgfqpoint{2.580449in}{0.459760in}}%
\pgfpathlineto{\pgfqpoint{2.581515in}{0.455951in}}%
\pgfpathlineto{\pgfqpoint{2.587084in}{0.454555in}}%
\pgfpathlineto{\pgfqpoint{2.592017in}{0.447697in}}%
\pgfpathlineto{\pgfqpoint{2.593275in}{0.442838in}}%
\pgfpathclose%
\pgfusepath{fill}%
\end{pgfscope}%
\begin{pgfscope}%
\pgfpathrectangle{\pgfqpoint{0.100000in}{0.100000in}}{\pgfqpoint{3.007045in}{1.925000in}}%
\pgfusepath{clip}%
\pgfsetbuttcap%
\pgfsetmiterjoin%
\definecolor{currentfill}{rgb}{0.040984,0.329950,0.621376}%
\pgfsetfillcolor{currentfill}%
\pgfsetlinewidth{0.000000pt}%
\definecolor{currentstroke}{rgb}{0.000000,0.000000,0.000000}%
\pgfsetstrokecolor{currentstroke}%
\pgfsetstrokeopacity{0.000000}%
\pgfsetdash{}{0pt}%
\pgfpathmoveto{\pgfqpoint{0.441277in}{1.107129in}}%
\pgfpathlineto{\pgfqpoint{0.475084in}{1.097667in}}%
\pgfpathlineto{\pgfqpoint{0.522151in}{1.084892in}}%
\pgfpathlineto{\pgfqpoint{0.554172in}{1.076466in}}%
\pgfpathlineto{\pgfqpoint{0.573744in}{1.071996in}}%
\pgfpathlineto{\pgfqpoint{0.557933in}{1.009610in}}%
\pgfpathlineto{\pgfqpoint{0.515049in}{1.020412in}}%
\pgfpathlineto{\pgfqpoint{0.473306in}{1.031380in}}%
\pgfpathlineto{\pgfqpoint{0.472706in}{1.036029in}}%
\pgfpathlineto{\pgfqpoint{0.464444in}{1.039777in}}%
\pgfpathlineto{\pgfqpoint{0.465944in}{1.051393in}}%
\pgfpathlineto{\pgfqpoint{0.461538in}{1.053350in}}%
\pgfpathlineto{\pgfqpoint{0.463546in}{1.059117in}}%
\pgfpathlineto{\pgfqpoint{0.457558in}{1.060425in}}%
\pgfpathlineto{\pgfqpoint{0.459131in}{1.065999in}}%
\pgfpathlineto{\pgfqpoint{0.451737in}{1.068135in}}%
\pgfpathlineto{\pgfqpoint{0.453296in}{1.073681in}}%
\pgfpathlineto{\pgfqpoint{0.449631in}{1.074733in}}%
\pgfpathlineto{\pgfqpoint{0.451207in}{1.080317in}}%
\pgfpathlineto{\pgfqpoint{0.445664in}{1.083905in}}%
\pgfpathlineto{\pgfqpoint{0.442087in}{1.088884in}}%
\pgfpathlineto{\pgfqpoint{0.443707in}{1.094459in}}%
\pgfpathlineto{\pgfqpoint{0.438131in}{1.096030in}}%
\pgfpathlineto{\pgfqpoint{0.441277in}{1.107129in}}%
\pgfpathclose%
\pgfusepath{fill}%
\end{pgfscope}%
\begin{pgfscope}%
\pgfpathrectangle{\pgfqpoint{0.100000in}{0.100000in}}{\pgfqpoint{3.007045in}{1.925000in}}%
\pgfusepath{clip}%
\pgfsetbuttcap%
\pgfsetmiterjoin%
\definecolor{currentfill}{rgb}{0.535686,0.746082,0.864252}%
\pgfsetfillcolor{currentfill}%
\pgfsetlinewidth{0.000000pt}%
\definecolor{currentstroke}{rgb}{0.000000,0.000000,0.000000}%
\pgfsetstrokecolor{currentstroke}%
\pgfsetstrokeopacity{0.000000}%
\pgfsetdash{}{0pt}%
\pgfpathmoveto{\pgfqpoint{1.844539in}{1.237709in}}%
\pgfpathlineto{\pgfqpoint{1.833661in}{1.237372in}}%
\pgfpathlineto{\pgfqpoint{1.798792in}{1.236593in}}%
\pgfpathlineto{\pgfqpoint{1.798597in}{1.275256in}}%
\pgfpathlineto{\pgfqpoint{1.826979in}{1.275432in}}%
\pgfpathlineto{\pgfqpoint{1.843908in}{1.275998in}}%
\pgfpathlineto{\pgfqpoint{1.844539in}{1.237709in}}%
\pgfpathclose%
\pgfusepath{fill}%
\end{pgfscope}%
\begin{pgfscope}%
\pgfpathrectangle{\pgfqpoint{0.100000in}{0.100000in}}{\pgfqpoint{3.007045in}{1.925000in}}%
\pgfusepath{clip}%
\pgfsetbuttcap%
\pgfsetmiterjoin%
\definecolor{currentfill}{rgb}{0.429020,0.687520,0.841246}%
\pgfsetfillcolor{currentfill}%
\pgfsetlinewidth{0.000000pt}%
\definecolor{currentstroke}{rgb}{0.000000,0.000000,0.000000}%
\pgfsetstrokecolor{currentstroke}%
\pgfsetstrokeopacity{0.000000}%
\pgfsetdash{}{0pt}%
\pgfpathmoveto{\pgfqpoint{2.287383in}{1.019436in}}%
\pgfpathlineto{\pgfqpoint{2.286020in}{1.023971in}}%
\pgfpathlineto{\pgfqpoint{2.274407in}{1.023601in}}%
\pgfpathlineto{\pgfqpoint{2.271241in}{1.027836in}}%
\pgfpathlineto{\pgfqpoint{2.269307in}{1.034406in}}%
\pgfpathlineto{\pgfqpoint{2.263654in}{1.034861in}}%
\pgfpathlineto{\pgfqpoint{2.260460in}{1.038312in}}%
\pgfpathlineto{\pgfqpoint{2.259090in}{1.046090in}}%
\pgfpathlineto{\pgfqpoint{2.260372in}{1.054200in}}%
\pgfpathlineto{\pgfqpoint{2.261228in}{1.057564in}}%
\pgfpathlineto{\pgfqpoint{2.270021in}{1.058081in}}%
\pgfpathlineto{\pgfqpoint{2.276076in}{1.059145in}}%
\pgfpathlineto{\pgfqpoint{2.281792in}{1.056352in}}%
\pgfpathlineto{\pgfqpoint{2.285344in}{1.059690in}}%
\pgfpathlineto{\pgfqpoint{2.292317in}{1.056941in}}%
\pgfpathlineto{\pgfqpoint{2.291623in}{1.065687in}}%
\pgfpathlineto{\pgfqpoint{2.301227in}{1.065647in}}%
\pgfpathlineto{\pgfqpoint{2.302296in}{1.059040in}}%
\pgfpathlineto{\pgfqpoint{2.309458in}{1.055417in}}%
\pgfpathlineto{\pgfqpoint{2.310463in}{1.047437in}}%
\pgfpathlineto{\pgfqpoint{2.301947in}{1.038398in}}%
\pgfpathlineto{\pgfqpoint{2.302206in}{1.033847in}}%
\pgfpathlineto{\pgfqpoint{2.306270in}{1.030998in}}%
\pgfpathlineto{\pgfqpoint{2.292665in}{1.018214in}}%
\pgfpathlineto{\pgfqpoint{2.287383in}{1.019436in}}%
\pgfpathclose%
\pgfusepath{fill}%
\end{pgfscope}%
\begin{pgfscope}%
\pgfpathrectangle{\pgfqpoint{0.100000in}{0.100000in}}{\pgfqpoint{3.007045in}{1.925000in}}%
\pgfusepath{clip}%
\pgfsetbuttcap%
\pgfsetmiterjoin%
\definecolor{currentfill}{rgb}{0.108651,0.416563,0.689043}%
\pgfsetfillcolor{currentfill}%
\pgfsetlinewidth{0.000000pt}%
\definecolor{currentstroke}{rgb}{0.000000,0.000000,0.000000}%
\pgfsetstrokecolor{currentstroke}%
\pgfsetstrokeopacity{0.000000}%
\pgfsetdash{}{0pt}%
\pgfpathmoveto{\pgfqpoint{1.442781in}{0.834883in}}%
\pgfpathlineto{\pgfqpoint{1.419358in}{0.836488in}}%
\pgfpathlineto{\pgfqpoint{1.421405in}{0.865168in}}%
\pgfpathlineto{\pgfqpoint{1.442027in}{0.863792in}}%
\pgfpathlineto{\pgfqpoint{1.444186in}{0.892364in}}%
\pgfpathlineto{\pgfqpoint{1.473659in}{0.890543in}}%
\pgfpathlineto{\pgfqpoint{1.471813in}{0.861849in}}%
\pgfpathlineto{\pgfqpoint{1.478574in}{0.861428in}}%
\pgfpathlineto{\pgfqpoint{1.476764in}{0.832809in}}%
\pgfpathlineto{\pgfqpoint{1.471341in}{0.833174in}}%
\pgfpathlineto{\pgfqpoint{1.442781in}{0.834883in}}%
\pgfpathclose%
\pgfusepath{fill}%
\end{pgfscope}%
\begin{pgfscope}%
\pgfpathrectangle{\pgfqpoint{0.100000in}{0.100000in}}{\pgfqpoint{3.007045in}{1.925000in}}%
\pgfusepath{clip}%
\pgfsetbuttcap%
\pgfsetmiterjoin%
\definecolor{currentfill}{rgb}{0.479216,0.715079,0.852072}%
\pgfsetfillcolor{currentfill}%
\pgfsetlinewidth{0.000000pt}%
\definecolor{currentstroke}{rgb}{0.000000,0.000000,0.000000}%
\pgfsetstrokecolor{currentstroke}%
\pgfsetstrokeopacity{0.000000}%
\pgfsetdash{}{0pt}%
\pgfpathmoveto{\pgfqpoint{1.405343in}{1.057957in}}%
\pgfpathlineto{\pgfqpoint{1.404039in}{1.041199in}}%
\pgfpathlineto{\pgfqpoint{1.401987in}{1.015175in}}%
\pgfpathlineto{\pgfqpoint{1.367585in}{1.018108in}}%
\pgfpathlineto{\pgfqpoint{1.351668in}{1.019814in}}%
\pgfpathlineto{\pgfqpoint{1.347270in}{1.020215in}}%
\pgfpathlineto{\pgfqpoint{1.351564in}{1.062383in}}%
\pgfpathlineto{\pgfqpoint{1.405343in}{1.057957in}}%
\pgfpathclose%
\pgfusepath{fill}%
\end{pgfscope}%
\begin{pgfscope}%
\pgfpathrectangle{\pgfqpoint{0.100000in}{0.100000in}}{\pgfqpoint{3.007045in}{1.925000in}}%
\pgfusepath{clip}%
\pgfsetbuttcap%
\pgfsetmiterjoin%
\definecolor{currentfill}{rgb}{0.381776,0.656517,0.824452}%
\pgfsetfillcolor{currentfill}%
\pgfsetlinewidth{0.000000pt}%
\definecolor{currentstroke}{rgb}{0.000000,0.000000,0.000000}%
\pgfsetstrokecolor{currentstroke}%
\pgfsetstrokeopacity{0.000000}%
\pgfsetdash{}{0pt}%
\pgfpathmoveto{\pgfqpoint{2.254813in}{0.936868in}}%
\pgfpathlineto{\pgfqpoint{2.251280in}{0.935449in}}%
\pgfpathlineto{\pgfqpoint{2.240739in}{0.937396in}}%
\pgfpathlineto{\pgfqpoint{2.228851in}{0.933436in}}%
\pgfpathlineto{\pgfqpoint{2.226333in}{0.935993in}}%
\pgfpathlineto{\pgfqpoint{2.219736in}{0.935214in}}%
\pgfpathlineto{\pgfqpoint{2.210327in}{0.936644in}}%
\pgfpathlineto{\pgfqpoint{2.188078in}{0.943307in}}%
\pgfpathlineto{\pgfqpoint{2.188043in}{0.950490in}}%
\pgfpathlineto{\pgfqpoint{2.188978in}{0.963509in}}%
\pgfpathlineto{\pgfqpoint{2.187491in}{0.967104in}}%
\pgfpathlineto{\pgfqpoint{2.188596in}{0.973610in}}%
\pgfpathlineto{\pgfqpoint{2.181638in}{0.973954in}}%
\pgfpathlineto{\pgfqpoint{2.183256in}{0.977798in}}%
\pgfpathlineto{\pgfqpoint{2.189194in}{0.983871in}}%
\pgfpathlineto{\pgfqpoint{2.188886in}{0.995787in}}%
\pgfpathlineto{\pgfqpoint{2.191739in}{0.996087in}}%
\pgfpathlineto{\pgfqpoint{2.216487in}{0.998935in}}%
\pgfpathlineto{\pgfqpoint{2.230102in}{0.999971in}}%
\pgfpathlineto{\pgfqpoint{2.248936in}{1.000483in}}%
\pgfpathlineto{\pgfqpoint{2.258514in}{1.000440in}}%
\pgfpathlineto{\pgfqpoint{2.258356in}{0.992865in}}%
\pgfpathlineto{\pgfqpoint{2.258038in}{0.978829in}}%
\pgfpathlineto{\pgfqpoint{2.261071in}{0.973710in}}%
\pgfpathlineto{\pgfqpoint{2.261147in}{0.968764in}}%
\pgfpathlineto{\pgfqpoint{2.258371in}{0.965969in}}%
\pgfpathlineto{\pgfqpoint{2.253187in}{0.966596in}}%
\pgfpathlineto{\pgfqpoint{2.248417in}{0.962697in}}%
\pgfpathlineto{\pgfqpoint{2.251263in}{0.956468in}}%
\pgfpathlineto{\pgfqpoint{2.257521in}{0.951222in}}%
\pgfpathlineto{\pgfqpoint{2.254335in}{0.941757in}}%
\pgfpathlineto{\pgfqpoint{2.254813in}{0.936868in}}%
\pgfpathclose%
\pgfusepath{fill}%
\end{pgfscope}%
\begin{pgfscope}%
\pgfpathrectangle{\pgfqpoint{0.100000in}{0.100000in}}{\pgfqpoint{3.007045in}{1.925000in}}%
\pgfusepath{clip}%
\pgfsetbuttcap%
\pgfsetmiterjoin%
\definecolor{currentfill}{rgb}{0.130427,0.444152,0.710327}%
\pgfsetfillcolor{currentfill}%
\pgfsetlinewidth{0.000000pt}%
\definecolor{currentstroke}{rgb}{0.000000,0.000000,0.000000}%
\pgfsetstrokecolor{currentstroke}%
\pgfsetstrokeopacity{0.000000}%
\pgfsetdash{}{0pt}%
\pgfpathmoveto{\pgfqpoint{0.421673in}{1.178083in}}%
\pgfpathlineto{\pgfqpoint{0.412185in}{1.187494in}}%
\pgfpathlineto{\pgfqpoint{0.409170in}{1.194540in}}%
\pgfpathlineto{\pgfqpoint{0.413081in}{1.203149in}}%
\pgfpathlineto{\pgfqpoint{0.413366in}{1.209180in}}%
\pgfpathlineto{\pgfqpoint{0.404869in}{1.211825in}}%
\pgfpathlineto{\pgfqpoint{0.404848in}{1.220444in}}%
\pgfpathlineto{\pgfqpoint{0.408920in}{1.226474in}}%
\pgfpathlineto{\pgfqpoint{0.407661in}{1.233230in}}%
\pgfpathlineto{\pgfqpoint{0.407690in}{1.233195in}}%
\pgfpathlineto{\pgfqpoint{0.424991in}{1.242670in}}%
\pgfpathlineto{\pgfqpoint{0.436885in}{1.243705in}}%
\pgfpathlineto{\pgfqpoint{0.440085in}{1.241248in}}%
\pgfpathlineto{\pgfqpoint{0.446079in}{1.262743in}}%
\pgfpathlineto{\pgfqpoint{0.455151in}{1.243275in}}%
\pgfpathlineto{\pgfqpoint{0.464303in}{1.248933in}}%
\pgfpathlineto{\pgfqpoint{0.475080in}{1.259356in}}%
\pgfpathlineto{\pgfqpoint{0.488781in}{1.270732in}}%
\pgfpathlineto{\pgfqpoint{0.497623in}{1.272306in}}%
\pgfpathlineto{\pgfqpoint{0.502964in}{1.265519in}}%
\pgfpathlineto{\pgfqpoint{0.510697in}{1.267522in}}%
\pgfpathlineto{\pgfqpoint{0.514663in}{1.260246in}}%
\pgfpathlineto{\pgfqpoint{0.513336in}{1.257851in}}%
\pgfpathlineto{\pgfqpoint{0.514807in}{1.248637in}}%
\pgfpathlineto{\pgfqpoint{0.525021in}{1.240725in}}%
\pgfpathlineto{\pgfqpoint{0.524524in}{1.231551in}}%
\pgfpathlineto{\pgfqpoint{0.527593in}{1.228218in}}%
\pgfpathlineto{\pgfqpoint{0.527157in}{1.220456in}}%
\pgfpathlineto{\pgfqpoint{0.522889in}{1.217818in}}%
\pgfpathlineto{\pgfqpoint{0.521553in}{1.220832in}}%
\pgfpathlineto{\pgfqpoint{0.503936in}{1.210812in}}%
\pgfpathlineto{\pgfqpoint{0.502749in}{1.206653in}}%
\pgfpathlineto{\pgfqpoint{0.498018in}{1.202710in}}%
\pgfpathlineto{\pgfqpoint{0.492430in}{1.204219in}}%
\pgfpathlineto{\pgfqpoint{0.473609in}{1.193492in}}%
\pgfpathlineto{\pgfqpoint{0.461047in}{1.194974in}}%
\pgfpathlineto{\pgfqpoint{0.450589in}{1.194118in}}%
\pgfpathlineto{\pgfqpoint{0.421673in}{1.178083in}}%
\pgfpathclose%
\pgfusepath{fill}%
\end{pgfscope}%
\begin{pgfscope}%
\pgfpathrectangle{\pgfqpoint{0.100000in}{0.100000in}}{\pgfqpoint{3.007045in}{1.925000in}}%
\pgfusepath{clip}%
\pgfsetbuttcap%
\pgfsetmiterjoin%
\definecolor{currentfill}{rgb}{0.114802,0.424437,0.695194}%
\pgfsetfillcolor{currentfill}%
\pgfsetlinewidth{0.000000pt}%
\definecolor{currentstroke}{rgb}{0.000000,0.000000,0.000000}%
\pgfsetstrokecolor{currentstroke}%
\pgfsetstrokeopacity{0.000000}%
\pgfsetdash{}{0pt}%
\pgfpathmoveto{\pgfqpoint{1.403892in}{1.607820in}}%
\pgfpathlineto{\pgfqpoint{1.355655in}{1.612365in}}%
\pgfpathlineto{\pgfqpoint{1.359600in}{1.651719in}}%
\pgfpathlineto{\pgfqpoint{1.370707in}{1.650612in}}%
\pgfpathlineto{\pgfqpoint{1.371277in}{1.656487in}}%
\pgfpathlineto{\pgfqpoint{1.379987in}{1.655666in}}%
\pgfpathlineto{\pgfqpoint{1.410950in}{1.652859in}}%
\pgfpathlineto{\pgfqpoint{1.409091in}{1.629777in}}%
\pgfpathlineto{\pgfqpoint{1.405844in}{1.630008in}}%
\pgfpathlineto{\pgfqpoint{1.403892in}{1.607820in}}%
\pgfpathclose%
\pgfusepath{fill}%
\end{pgfscope}%
\begin{pgfscope}%
\pgfpathrectangle{\pgfqpoint{0.100000in}{0.100000in}}{\pgfqpoint{3.007045in}{1.925000in}}%
\pgfusepath{clip}%
\pgfsetbuttcap%
\pgfsetmiterjoin%
\definecolor{currentfill}{rgb}{0.529412,0.742637,0.862899}%
\pgfsetfillcolor{currentfill}%
\pgfsetlinewidth{0.000000pt}%
\definecolor{currentstroke}{rgb}{0.000000,0.000000,0.000000}%
\pgfsetstrokecolor{currentstroke}%
\pgfsetstrokeopacity{0.000000}%
\pgfsetdash{}{0pt}%
\pgfpathmoveto{\pgfqpoint{2.751955in}{1.186579in}}%
\pgfpathlineto{\pgfqpoint{2.748588in}{1.187139in}}%
\pgfpathlineto{\pgfqpoint{2.739288in}{1.183920in}}%
\pgfpathlineto{\pgfqpoint{2.728190in}{1.196931in}}%
\pgfpathlineto{\pgfqpoint{2.729962in}{1.199364in}}%
\pgfpathlineto{\pgfqpoint{2.728781in}{1.207289in}}%
\pgfpathlineto{\pgfqpoint{2.734998in}{1.209843in}}%
\pgfpathlineto{\pgfqpoint{2.724066in}{1.216278in}}%
\pgfpathlineto{\pgfqpoint{2.728601in}{1.227398in}}%
\pgfpathlineto{\pgfqpoint{2.725297in}{1.231605in}}%
\pgfpathlineto{\pgfqpoint{2.725039in}{1.237269in}}%
\pgfpathlineto{\pgfqpoint{2.721783in}{1.241309in}}%
\pgfpathlineto{\pgfqpoint{2.724195in}{1.253074in}}%
\pgfpathlineto{\pgfqpoint{2.730272in}{1.258328in}}%
\pgfpathlineto{\pgfqpoint{2.738533in}{1.258792in}}%
\pgfpathlineto{\pgfqpoint{2.744257in}{1.261114in}}%
\pgfpathlineto{\pgfqpoint{2.745705in}{1.255995in}}%
\pgfpathlineto{\pgfqpoint{2.758832in}{1.209240in}}%
\pgfpathlineto{\pgfqpoint{2.753781in}{1.202785in}}%
\pgfpathlineto{\pgfqpoint{2.751060in}{1.190437in}}%
\pgfpathlineto{\pgfqpoint{2.751955in}{1.186579in}}%
\pgfpathclose%
\pgfusepath{fill}%
\end{pgfscope}%
\begin{pgfscope}%
\pgfpathrectangle{\pgfqpoint{0.100000in}{0.100000in}}{\pgfqpoint{3.007045in}{1.925000in}}%
\pgfusepath{clip}%
\pgfsetbuttcap%
\pgfsetmiterjoin%
\definecolor{currentfill}{rgb}{0.441569,0.694410,0.843952}%
\pgfsetfillcolor{currentfill}%
\pgfsetlinewidth{0.000000pt}%
\definecolor{currentstroke}{rgb}{0.000000,0.000000,0.000000}%
\pgfsetstrokecolor{currentstroke}%
\pgfsetstrokeopacity{0.000000}%
\pgfsetdash{}{0pt}%
\pgfpathmoveto{\pgfqpoint{2.124853in}{1.498938in}}%
\pgfpathlineto{\pgfqpoint{2.108536in}{1.497608in}}%
\pgfpathlineto{\pgfqpoint{2.108968in}{1.491853in}}%
\pgfpathlineto{\pgfqpoint{2.101634in}{1.491321in}}%
\pgfpathlineto{\pgfqpoint{2.094574in}{1.490855in}}%
\pgfpathlineto{\pgfqpoint{2.092961in}{1.513610in}}%
\pgfpathlineto{\pgfqpoint{2.090405in}{1.513286in}}%
\pgfpathlineto{\pgfqpoint{2.088935in}{1.530951in}}%
\pgfpathlineto{\pgfqpoint{2.077656in}{1.530311in}}%
\pgfpathlineto{\pgfqpoint{2.076668in}{1.547607in}}%
\pgfpathlineto{\pgfqpoint{2.069326in}{1.547117in}}%
\pgfpathlineto{\pgfqpoint{2.067108in}{1.552783in}}%
\pgfpathlineto{\pgfqpoint{2.066412in}{1.564272in}}%
\pgfpathlineto{\pgfqpoint{2.078017in}{1.564986in}}%
\pgfpathlineto{\pgfqpoint{2.083707in}{1.565240in}}%
\pgfpathlineto{\pgfqpoint{2.084400in}{1.553797in}}%
\pgfpathlineto{\pgfqpoint{2.090120in}{1.553996in}}%
\pgfpathlineto{\pgfqpoint{2.093706in}{1.548428in}}%
\pgfpathlineto{\pgfqpoint{2.094036in}{1.542696in}}%
\pgfpathlineto{\pgfqpoint{2.111077in}{1.540159in}}%
\pgfpathlineto{\pgfqpoint{2.102036in}{1.523219in}}%
\pgfpathlineto{\pgfqpoint{2.100053in}{1.512797in}}%
\pgfpathlineto{\pgfqpoint{2.104771in}{1.510691in}}%
\pgfpathlineto{\pgfqpoint{2.108185in}{1.516689in}}%
\pgfpathlineto{\pgfqpoint{2.112193in}{1.518509in}}%
\pgfpathlineto{\pgfqpoint{2.118815in}{1.532207in}}%
\pgfpathlineto{\pgfqpoint{2.124890in}{1.534323in}}%
\pgfpathlineto{\pgfqpoint{2.127996in}{1.537667in}}%
\pgfpathlineto{\pgfqpoint{2.133801in}{1.550138in}}%
\pgfpathlineto{\pgfqpoint{2.134382in}{1.555481in}}%
\pgfpathlineto{\pgfqpoint{2.143696in}{1.559230in}}%
\pgfpathlineto{\pgfqpoint{2.143649in}{1.552093in}}%
\pgfpathlineto{\pgfqpoint{2.137836in}{1.542472in}}%
\pgfpathlineto{\pgfqpoint{2.137692in}{1.536286in}}%
\pgfpathlineto{\pgfqpoint{2.134764in}{1.533978in}}%
\pgfpathlineto{\pgfqpoint{2.128245in}{1.517619in}}%
\pgfpathlineto{\pgfqpoint{2.124853in}{1.498938in}}%
\pgfpathclose%
\pgfusepath{fill}%
\end{pgfscope}%
\begin{pgfscope}%
\pgfpathrectangle{\pgfqpoint{0.100000in}{0.100000in}}{\pgfqpoint{3.007045in}{1.925000in}}%
\pgfusepath{clip}%
\pgfsetbuttcap%
\pgfsetmiterjoin%
\definecolor{currentfill}{rgb}{0.802307,0.876048,0.945867}%
\pgfsetfillcolor{currentfill}%
\pgfsetlinewidth{0.000000pt}%
\definecolor{currentstroke}{rgb}{0.000000,0.000000,0.000000}%
\pgfsetstrokecolor{currentstroke}%
\pgfsetstrokeopacity{0.000000}%
\pgfsetdash{}{0pt}%
\pgfpathmoveto{\pgfqpoint{2.636603in}{1.298784in}}%
\pgfpathlineto{\pgfqpoint{2.631148in}{1.284175in}}%
\pgfpathlineto{\pgfqpoint{2.627614in}{1.286206in}}%
\pgfpathlineto{\pgfqpoint{2.616493in}{1.288353in}}%
\pgfpathlineto{\pgfqpoint{2.610752in}{1.291249in}}%
\pgfpathlineto{\pgfqpoint{2.608820in}{1.295800in}}%
\pgfpathlineto{\pgfqpoint{2.612866in}{1.308899in}}%
\pgfpathlineto{\pgfqpoint{2.608640in}{1.313232in}}%
\pgfpathlineto{\pgfqpoint{2.609389in}{1.317404in}}%
\pgfpathlineto{\pgfqpoint{2.606011in}{1.320523in}}%
\pgfpathlineto{\pgfqpoint{2.599096in}{1.322608in}}%
\pgfpathlineto{\pgfqpoint{2.576932in}{1.318538in}}%
\pgfpathlineto{\pgfqpoint{2.574642in}{1.330275in}}%
\pgfpathlineto{\pgfqpoint{2.554469in}{1.327047in}}%
\pgfpathlineto{\pgfqpoint{2.552895in}{1.336088in}}%
\pgfpathlineto{\pgfqpoint{2.549986in}{1.354398in}}%
\pgfpathlineto{\pgfqpoint{2.555437in}{1.355961in}}%
\pgfpathlineto{\pgfqpoint{2.561808in}{1.358520in}}%
\pgfpathlineto{\pgfqpoint{2.573880in}{1.354042in}}%
\pgfpathlineto{\pgfqpoint{2.605715in}{1.356835in}}%
\pgfpathlineto{\pgfqpoint{2.609177in}{1.353242in}}%
\pgfpathlineto{\pgfqpoint{2.621003in}{1.357074in}}%
\pgfpathlineto{\pgfqpoint{2.623988in}{1.352894in}}%
\pgfpathlineto{\pgfqpoint{2.627962in}{1.353837in}}%
\pgfpathlineto{\pgfqpoint{2.637385in}{1.345841in}}%
\pgfpathlineto{\pgfqpoint{2.645166in}{1.349038in}}%
\pgfpathlineto{\pgfqpoint{2.654171in}{1.354581in}}%
\pgfpathlineto{\pgfqpoint{2.645950in}{1.339764in}}%
\pgfpathlineto{\pgfqpoint{2.633083in}{1.330389in}}%
\pgfpathlineto{\pgfqpoint{2.626236in}{1.316064in}}%
\pgfpathlineto{\pgfqpoint{2.628110in}{1.312979in}}%
\pgfpathlineto{\pgfqpoint{2.625195in}{1.307053in}}%
\pgfpathlineto{\pgfqpoint{2.627758in}{1.303889in}}%
\pgfpathlineto{\pgfqpoint{2.632407in}{1.305014in}}%
\pgfpathlineto{\pgfqpoint{2.636603in}{1.298784in}}%
\pgfpathclose%
\pgfusepath{fill}%
\end{pgfscope}%
\begin{pgfscope}%
\pgfpathrectangle{\pgfqpoint{0.100000in}{0.100000in}}{\pgfqpoint{3.007045in}{1.925000in}}%
\pgfusepath{clip}%
\pgfsetbuttcap%
\pgfsetmiterjoin%
\definecolor{currentfill}{rgb}{0.412042,0.677186,0.836263}%
\pgfsetfillcolor{currentfill}%
\pgfsetlinewidth{0.000000pt}%
\definecolor{currentstroke}{rgb}{0.000000,0.000000,0.000000}%
\pgfsetstrokecolor{currentstroke}%
\pgfsetstrokeopacity{0.000000}%
\pgfsetdash{}{0pt}%
\pgfpathmoveto{\pgfqpoint{1.847419in}{1.345161in}}%
\pgfpathlineto{\pgfqpoint{1.813145in}{1.344757in}}%
\pgfpathlineto{\pgfqpoint{1.779241in}{1.344556in}}%
\pgfpathlineto{\pgfqpoint{1.779293in}{1.361999in}}%
\pgfpathlineto{\pgfqpoint{1.776477in}{1.362015in}}%
\pgfpathlineto{\pgfqpoint{1.776480in}{1.390760in}}%
\pgfpathlineto{\pgfqpoint{1.822052in}{1.390899in}}%
\pgfpathlineto{\pgfqpoint{1.844917in}{1.391266in}}%
\pgfpathlineto{\pgfqpoint{1.845372in}{1.362394in}}%
\pgfpathlineto{\pgfqpoint{1.847146in}{1.362415in}}%
\pgfpathlineto{\pgfqpoint{1.847419in}{1.345161in}}%
\pgfpathclose%
\pgfusepath{fill}%
\end{pgfscope}%
\begin{pgfscope}%
\pgfpathrectangle{\pgfqpoint{0.100000in}{0.100000in}}{\pgfqpoint{3.007045in}{1.925000in}}%
\pgfusepath{clip}%
\pgfsetbuttcap%
\pgfsetmiterjoin%
\definecolor{currentfill}{rgb}{0.554510,0.756417,0.868312}%
\pgfsetfillcolor{currentfill}%
\pgfsetlinewidth{0.000000pt}%
\definecolor{currentstroke}{rgb}{0.000000,0.000000,0.000000}%
\pgfsetstrokecolor{currentstroke}%
\pgfsetstrokeopacity{0.000000}%
\pgfsetdash{}{0pt}%
\pgfpathmoveto{\pgfqpoint{2.905539in}{1.431014in}}%
\pgfpathlineto{\pgfqpoint{2.901522in}{1.431102in}}%
\pgfpathlineto{\pgfqpoint{2.893781in}{1.426472in}}%
\pgfpathlineto{\pgfqpoint{2.889192in}{1.426210in}}%
\pgfpathlineto{\pgfqpoint{2.881149in}{1.420623in}}%
\pgfpathlineto{\pgfqpoint{2.877803in}{1.421007in}}%
\pgfpathlineto{\pgfqpoint{2.855094in}{1.414428in}}%
\pgfpathlineto{\pgfqpoint{2.824324in}{1.387968in}}%
\pgfpathlineto{\pgfqpoint{2.819156in}{1.394508in}}%
\pgfpathlineto{\pgfqpoint{2.829259in}{1.404526in}}%
\pgfpathlineto{\pgfqpoint{2.824695in}{1.409031in}}%
\pgfpathlineto{\pgfqpoint{2.822125in}{1.424128in}}%
\pgfpathlineto{\pgfqpoint{2.816097in}{1.458112in}}%
\pgfpathlineto{\pgfqpoint{2.836904in}{1.462520in}}%
\pgfpathlineto{\pgfqpoint{2.880647in}{1.472875in}}%
\pgfpathlineto{\pgfqpoint{2.896973in}{1.475592in}}%
\pgfpathlineto{\pgfqpoint{2.903419in}{1.453280in}}%
\pgfpathlineto{\pgfqpoint{2.906901in}{1.437874in}}%
\pgfpathlineto{\pgfqpoint{2.905539in}{1.431014in}}%
\pgfpathclose%
\pgfusepath{fill}%
\end{pgfscope}%
\begin{pgfscope}%
\pgfpathrectangle{\pgfqpoint{0.100000in}{0.100000in}}{\pgfqpoint{3.007045in}{1.925000in}}%
\pgfusepath{clip}%
\pgfsetbuttcap%
\pgfsetmiterjoin%
\definecolor{currentfill}{rgb}{0.498039,0.725413,0.856132}%
\pgfsetfillcolor{currentfill}%
\pgfsetlinewidth{0.000000pt}%
\definecolor{currentstroke}{rgb}{0.000000,0.000000,0.000000}%
\pgfsetstrokecolor{currentstroke}%
\pgfsetstrokeopacity{0.000000}%
\pgfsetdash{}{0pt}%
\pgfpathmoveto{\pgfqpoint{1.977807in}{1.222266in}}%
\pgfpathlineto{\pgfqpoint{1.977971in}{1.216245in}}%
\pgfpathlineto{\pgfqpoint{1.948122in}{1.215494in}}%
\pgfpathlineto{\pgfqpoint{1.948406in}{1.218688in}}%
\pgfpathlineto{\pgfqpoint{1.925598in}{1.218425in}}%
\pgfpathlineto{\pgfqpoint{1.925211in}{1.241392in}}%
\pgfpathlineto{\pgfqpoint{1.935891in}{1.242279in}}%
\pgfpathlineto{\pgfqpoint{1.935244in}{1.278471in}}%
\pgfpathlineto{\pgfqpoint{1.952380in}{1.279020in}}%
\pgfpathlineto{\pgfqpoint{1.952591in}{1.273282in}}%
\pgfpathlineto{\pgfqpoint{1.973511in}{1.274071in}}%
\pgfpathlineto{\pgfqpoint{1.981668in}{1.274180in}}%
\pgfpathlineto{\pgfqpoint{1.982776in}{1.245666in}}%
\pgfpathlineto{\pgfqpoint{1.977034in}{1.245635in}}%
\pgfpathlineto{\pgfqpoint{1.977807in}{1.222266in}}%
\pgfpathclose%
\pgfusepath{fill}%
\end{pgfscope}%
\begin{pgfscope}%
\pgfpathrectangle{\pgfqpoint{0.100000in}{0.100000in}}{\pgfqpoint{3.007045in}{1.925000in}}%
\pgfusepath{clip}%
\pgfsetbuttcap%
\pgfsetmiterjoin%
\definecolor{currentfill}{rgb}{0.510588,0.732303,0.858839}%
\pgfsetfillcolor{currentfill}%
\pgfsetlinewidth{0.000000pt}%
\definecolor{currentstroke}{rgb}{0.000000,0.000000,0.000000}%
\pgfsetstrokecolor{currentstroke}%
\pgfsetstrokeopacity{0.000000}%
\pgfsetdash{}{0pt}%
\pgfpathmoveto{\pgfqpoint{2.348646in}{0.902478in}}%
\pgfpathlineto{\pgfqpoint{2.332547in}{0.900599in}}%
\pgfpathlineto{\pgfqpoint{2.313442in}{0.898402in}}%
\pgfpathlineto{\pgfqpoint{2.309785in}{0.907283in}}%
\pgfpathlineto{\pgfqpoint{2.313145in}{0.917931in}}%
\pgfpathlineto{\pgfqpoint{2.308274in}{0.925835in}}%
\pgfpathlineto{\pgfqpoint{2.313025in}{0.928226in}}%
\pgfpathlineto{\pgfqpoint{2.319657in}{0.939932in}}%
\pgfpathlineto{\pgfqpoint{2.318851in}{0.944884in}}%
\pgfpathlineto{\pgfqpoint{2.321351in}{0.950113in}}%
\pgfpathlineto{\pgfqpoint{2.327024in}{0.948030in}}%
\pgfpathlineto{\pgfqpoint{2.329679in}{0.943109in}}%
\pgfpathlineto{\pgfqpoint{2.341668in}{0.946890in}}%
\pgfpathlineto{\pgfqpoint{2.354830in}{0.942630in}}%
\pgfpathlineto{\pgfqpoint{2.364223in}{0.935952in}}%
\pgfpathlineto{\pgfqpoint{2.361650in}{0.931386in}}%
\pgfpathlineto{\pgfqpoint{2.361956in}{0.924269in}}%
\pgfpathlineto{\pgfqpoint{2.356471in}{0.920502in}}%
\pgfpathlineto{\pgfqpoint{2.351726in}{0.921503in}}%
\pgfpathlineto{\pgfqpoint{2.348504in}{0.918229in}}%
\pgfpathlineto{\pgfqpoint{2.348646in}{0.902478in}}%
\pgfpathclose%
\pgfusepath{fill}%
\end{pgfscope}%
\begin{pgfscope}%
\pgfpathrectangle{\pgfqpoint{0.100000in}{0.100000in}}{\pgfqpoint{3.007045in}{1.925000in}}%
\pgfusepath{clip}%
\pgfsetbuttcap%
\pgfsetmiterjoin%
\definecolor{currentfill}{rgb}{0.541961,0.749527,0.865606}%
\pgfsetfillcolor{currentfill}%
\pgfsetlinewidth{0.000000pt}%
\definecolor{currentstroke}{rgb}{0.000000,0.000000,0.000000}%
\pgfsetstrokecolor{currentstroke}%
\pgfsetstrokeopacity{0.000000}%
\pgfsetdash{}{0pt}%
\pgfpathmoveto{\pgfqpoint{1.280556in}{0.903442in}}%
\pgfpathlineto{\pgfqpoint{1.281006in}{0.908273in}}%
\pgfpathlineto{\pgfqpoint{1.218495in}{0.915021in}}%
\pgfpathlineto{\pgfqpoint{1.217220in}{0.903588in}}%
\pgfpathlineto{\pgfqpoint{1.194470in}{0.906193in}}%
\pgfpathlineto{\pgfqpoint{1.201592in}{0.967633in}}%
\pgfpathlineto{\pgfqpoint{1.209230in}{0.966315in}}%
\pgfpathlineto{\pgfqpoint{1.218512in}{0.975812in}}%
\pgfpathlineto{\pgfqpoint{1.221151in}{0.983550in}}%
\pgfpathlineto{\pgfqpoint{1.223675in}{0.983885in}}%
\pgfpathlineto{\pgfqpoint{1.225457in}{0.993376in}}%
\pgfpathlineto{\pgfqpoint{1.225227in}{1.010762in}}%
\pgfpathlineto{\pgfqpoint{1.231839in}{1.013157in}}%
\pgfpathlineto{\pgfqpoint{1.235393in}{1.026018in}}%
\pgfpathlineto{\pgfqpoint{1.235453in}{1.031177in}}%
\pgfpathlineto{\pgfqpoint{1.238883in}{1.030797in}}%
\pgfpathlineto{\pgfqpoint{1.277959in}{1.026457in}}%
\pgfpathlineto{\pgfqpoint{1.298942in}{1.024507in}}%
\pgfpathlineto{\pgfqpoint{1.293808in}{0.973457in}}%
\pgfpathlineto{\pgfqpoint{1.271190in}{0.975810in}}%
\pgfpathlineto{\pgfqpoint{1.273752in}{0.963960in}}%
\pgfpathlineto{\pgfqpoint{1.273425in}{0.954257in}}%
\pgfpathlineto{\pgfqpoint{1.271875in}{0.946736in}}%
\pgfpathlineto{\pgfqpoint{1.290299in}{0.944870in}}%
\pgfpathlineto{\pgfqpoint{1.292797in}{0.945962in}}%
\pgfpathlineto{\pgfqpoint{1.303832in}{0.919724in}}%
\pgfpathlineto{\pgfqpoint{1.308509in}{0.919262in}}%
\pgfpathlineto{\pgfqpoint{1.307382in}{0.907359in}}%
\pgfpathlineto{\pgfqpoint{1.295559in}{0.908588in}}%
\pgfpathlineto{\pgfqpoint{1.280556in}{0.903442in}}%
\pgfpathclose%
\pgfusepath{fill}%
\end{pgfscope}%
\begin{pgfscope}%
\pgfpathrectangle{\pgfqpoint{0.100000in}{0.100000in}}{\pgfqpoint{3.007045in}{1.925000in}}%
\pgfusepath{clip}%
\pgfsetbuttcap%
\pgfsetmiterjoin%
\definecolor{currentfill}{rgb}{0.326290,0.618624,0.802799}%
\pgfsetfillcolor{currentfill}%
\pgfsetlinewidth{0.000000pt}%
\definecolor{currentstroke}{rgb}{0.000000,0.000000,0.000000}%
\pgfsetstrokecolor{currentstroke}%
\pgfsetstrokeopacity{0.000000}%
\pgfsetdash{}{0pt}%
\pgfpathmoveto{\pgfqpoint{1.905724in}{1.300373in}}%
\pgfpathlineto{\pgfqpoint{1.854836in}{1.299029in}}%
\pgfpathlineto{\pgfqpoint{1.853741in}{1.305138in}}%
\pgfpathlineto{\pgfqpoint{1.853457in}{1.322433in}}%
\pgfpathlineto{\pgfqpoint{1.859156in}{1.322514in}}%
\pgfpathlineto{\pgfqpoint{1.858723in}{1.345342in}}%
\pgfpathlineto{\pgfqpoint{1.869963in}{1.345597in}}%
\pgfpathlineto{\pgfqpoint{1.881395in}{1.345922in}}%
\pgfpathlineto{\pgfqpoint{1.881224in}{1.351626in}}%
\pgfpathlineto{\pgfqpoint{1.904114in}{1.352301in}}%
\pgfpathlineto{\pgfqpoint{1.905724in}{1.300373in}}%
\pgfpathclose%
\pgfusepath{fill}%
\end{pgfscope}%
\begin{pgfscope}%
\pgfpathrectangle{\pgfqpoint{0.100000in}{0.100000in}}{\pgfqpoint{3.007045in}{1.925000in}}%
\pgfusepath{clip}%
\pgfsetbuttcap%
\pgfsetmiterjoin%
\definecolor{currentfill}{rgb}{0.429020,0.687520,0.841246}%
\pgfsetfillcolor{currentfill}%
\pgfsetlinewidth{0.000000pt}%
\definecolor{currentstroke}{rgb}{0.000000,0.000000,0.000000}%
\pgfsetstrokecolor{currentstroke}%
\pgfsetstrokeopacity{0.000000}%
\pgfsetdash{}{0pt}%
\pgfpathmoveto{\pgfqpoint{0.615468in}{1.284198in}}%
\pgfpathlineto{\pgfqpoint{0.597849in}{1.283744in}}%
\pgfpathlineto{\pgfqpoint{0.596506in}{1.278590in}}%
\pgfpathlineto{\pgfqpoint{0.614328in}{1.243704in}}%
\pgfpathlineto{\pgfqpoint{0.571645in}{1.215280in}}%
\pgfpathlineto{\pgfqpoint{0.567833in}{1.216455in}}%
\pgfpathlineto{\pgfqpoint{0.540209in}{1.259221in}}%
\pgfpathlineto{\pgfqpoint{0.552667in}{1.255842in}}%
\pgfpathlineto{\pgfqpoint{0.560395in}{1.283604in}}%
\pgfpathlineto{\pgfqpoint{0.554814in}{1.285052in}}%
\pgfpathlineto{\pgfqpoint{0.556383in}{1.290907in}}%
\pgfpathlineto{\pgfqpoint{0.562826in}{1.298061in}}%
\pgfpathlineto{\pgfqpoint{0.615468in}{1.284198in}}%
\pgfpathclose%
\pgfusepath{fill}%
\end{pgfscope}%
\begin{pgfscope}%
\pgfpathrectangle{\pgfqpoint{0.100000in}{0.100000in}}{\pgfqpoint{3.007045in}{1.925000in}}%
\pgfusepath{clip}%
\pgfsetbuttcap%
\pgfsetmiterjoin%
\definecolor{currentfill}{rgb}{0.676817,0.816471,0.902376}%
\pgfsetfillcolor{currentfill}%
\pgfsetlinewidth{0.000000pt}%
\definecolor{currentstroke}{rgb}{0.000000,0.000000,0.000000}%
\pgfsetstrokecolor{currentstroke}%
\pgfsetstrokeopacity{0.000000}%
\pgfsetdash{}{0pt}%
\pgfpathmoveto{\pgfqpoint{2.514084in}{1.019211in}}%
\pgfpathlineto{\pgfqpoint{2.506602in}{1.020237in}}%
\pgfpathlineto{\pgfqpoint{2.500787in}{1.015480in}}%
\pgfpathlineto{\pgfqpoint{2.492311in}{1.025615in}}%
\pgfpathlineto{\pgfqpoint{2.492133in}{1.027383in}}%
\pgfpathlineto{\pgfqpoint{2.475027in}{1.025720in}}%
\pgfpathlineto{\pgfqpoint{2.476402in}{1.027509in}}%
\pgfpathlineto{\pgfqpoint{2.481176in}{1.031217in}}%
\pgfpathlineto{\pgfqpoint{2.481841in}{1.034709in}}%
\pgfpathlineto{\pgfqpoint{2.495062in}{1.040390in}}%
\pgfpathlineto{\pgfqpoint{2.503517in}{1.042080in}}%
\pgfpathlineto{\pgfqpoint{2.505363in}{1.044619in}}%
\pgfpathlineto{\pgfqpoint{2.521063in}{1.052157in}}%
\pgfpathlineto{\pgfqpoint{2.526060in}{1.056473in}}%
\pgfpathlineto{\pgfqpoint{2.537236in}{1.043373in}}%
\pgfpathlineto{\pgfqpoint{2.537514in}{1.041038in}}%
\pgfpathlineto{\pgfqpoint{2.529927in}{1.036274in}}%
\pgfpathlineto{\pgfqpoint{2.531083in}{1.032312in}}%
\pgfpathlineto{\pgfqpoint{2.515896in}{1.030173in}}%
\pgfpathlineto{\pgfqpoint{2.513186in}{1.021977in}}%
\pgfpathlineto{\pgfqpoint{2.514084in}{1.019211in}}%
\pgfpathclose%
\pgfusepath{fill}%
\end{pgfscope}%
\begin{pgfscope}%
\pgfpathrectangle{\pgfqpoint{0.100000in}{0.100000in}}{\pgfqpoint{3.007045in}{1.925000in}}%
\pgfusepath{clip}%
\pgfsetbuttcap%
\pgfsetmiterjoin%
\definecolor{currentfill}{rgb}{0.376732,0.653072,0.822484}%
\pgfsetfillcolor{currentfill}%
\pgfsetlinewidth{0.000000pt}%
\definecolor{currentstroke}{rgb}{0.000000,0.000000,0.000000}%
\pgfsetstrokecolor{currentstroke}%
\pgfsetstrokeopacity{0.000000}%
\pgfsetdash{}{0pt}%
\pgfpathmoveto{\pgfqpoint{2.407666in}{1.313899in}}%
\pgfpathlineto{\pgfqpoint{2.410399in}{1.309530in}}%
\pgfpathlineto{\pgfqpoint{2.413572in}{1.288426in}}%
\pgfpathlineto{\pgfqpoint{2.408904in}{1.287720in}}%
\pgfpathlineto{\pgfqpoint{2.409849in}{1.281241in}}%
\pgfpathlineto{\pgfqpoint{2.402337in}{1.279030in}}%
\pgfpathlineto{\pgfqpoint{2.389048in}{1.277134in}}%
\pgfpathlineto{\pgfqpoint{2.390704in}{1.263733in}}%
\pgfpathlineto{\pgfqpoint{2.376485in}{1.262685in}}%
\pgfpathlineto{\pgfqpoint{2.376250in}{1.266421in}}%
\pgfpathlineto{\pgfqpoint{2.359900in}{1.266193in}}%
\pgfpathlineto{\pgfqpoint{2.359392in}{1.270321in}}%
\pgfpathlineto{\pgfqpoint{2.350943in}{1.269110in}}%
\pgfpathlineto{\pgfqpoint{2.349350in}{1.280970in}}%
\pgfpathlineto{\pgfqpoint{2.364496in}{1.284008in}}%
\pgfpathlineto{\pgfqpoint{2.361935in}{1.302986in}}%
\pgfpathlineto{\pgfqpoint{2.375866in}{1.305055in}}%
\pgfpathlineto{\pgfqpoint{2.395446in}{1.307499in}}%
\pgfpathlineto{\pgfqpoint{2.394589in}{1.312199in}}%
\pgfpathlineto{\pgfqpoint{2.407666in}{1.313899in}}%
\pgfpathclose%
\pgfusepath{fill}%
\end{pgfscope}%
\begin{pgfscope}%
\pgfpathrectangle{\pgfqpoint{0.100000in}{0.100000in}}{\pgfqpoint{3.007045in}{1.925000in}}%
\pgfusepath{clip}%
\pgfsetbuttcap%
\pgfsetmiterjoin%
\definecolor{currentfill}{rgb}{0.666974,0.812288,0.898931}%
\pgfsetfillcolor{currentfill}%
\pgfsetlinewidth{0.000000pt}%
\definecolor{currentstroke}{rgb}{0.000000,0.000000,0.000000}%
\pgfsetstrokecolor{currentstroke}%
\pgfsetstrokeopacity{0.000000}%
\pgfsetdash{}{0pt}%
\pgfpathmoveto{\pgfqpoint{2.285385in}{0.933858in}}%
\pgfpathlineto{\pgfqpoint{2.292762in}{0.944860in}}%
\pgfpathlineto{\pgfqpoint{2.293939in}{0.955321in}}%
\pgfpathlineto{\pgfqpoint{2.291000in}{0.961775in}}%
\pgfpathlineto{\pgfqpoint{2.291608in}{0.970240in}}%
\pgfpathlineto{\pgfqpoint{2.296056in}{0.970309in}}%
\pgfpathlineto{\pgfqpoint{2.299083in}{0.977051in}}%
\pgfpathlineto{\pgfqpoint{2.296204in}{0.983928in}}%
\pgfpathlineto{\pgfqpoint{2.295566in}{0.990662in}}%
\pgfpathlineto{\pgfqpoint{2.298509in}{1.001042in}}%
\pgfpathlineto{\pgfqpoint{2.303650in}{1.003101in}}%
\pgfpathlineto{\pgfqpoint{2.315188in}{1.000641in}}%
\pgfpathlineto{\pgfqpoint{2.319984in}{0.992600in}}%
\pgfpathlineto{\pgfqpoint{2.317969in}{0.990753in}}%
\pgfpathlineto{\pgfqpoint{2.309261in}{0.983763in}}%
\pgfpathlineto{\pgfqpoint{2.308631in}{0.975465in}}%
\pgfpathlineto{\pgfqpoint{2.319716in}{0.965996in}}%
\pgfpathlineto{\pgfqpoint{2.322488in}{0.960571in}}%
\pgfpathlineto{\pgfqpoint{2.317579in}{0.954327in}}%
\pgfpathlineto{\pgfqpoint{2.321351in}{0.950113in}}%
\pgfpathlineto{\pgfqpoint{2.318851in}{0.944884in}}%
\pgfpathlineto{\pgfqpoint{2.319657in}{0.939932in}}%
\pgfpathlineto{\pgfqpoint{2.313025in}{0.928226in}}%
\pgfpathlineto{\pgfqpoint{2.308274in}{0.925835in}}%
\pgfpathlineto{\pgfqpoint{2.299822in}{0.927995in}}%
\pgfpathlineto{\pgfqpoint{2.297460in}{0.921002in}}%
\pgfpathlineto{\pgfqpoint{2.287005in}{0.929042in}}%
\pgfpathlineto{\pgfqpoint{2.285385in}{0.933858in}}%
\pgfpathclose%
\pgfusepath{fill}%
\end{pgfscope}%
\begin{pgfscope}%
\pgfpathrectangle{\pgfqpoint{0.100000in}{0.100000in}}{\pgfqpoint{3.007045in}{1.925000in}}%
\pgfusepath{clip}%
\pgfsetbuttcap%
\pgfsetmiterjoin%
\definecolor{currentfill}{rgb}{0.657132,0.808105,0.895486}%
\pgfsetfillcolor{currentfill}%
\pgfsetlinewidth{0.000000pt}%
\definecolor{currentstroke}{rgb}{0.000000,0.000000,0.000000}%
\pgfsetstrokecolor{currentstroke}%
\pgfsetstrokeopacity{0.000000}%
\pgfsetdash{}{0pt}%
\pgfpathmoveto{\pgfqpoint{2.742724in}{1.035580in}}%
\pgfpathlineto{\pgfqpoint{2.743345in}{1.029859in}}%
\pgfpathlineto{\pgfqpoint{2.741385in}{1.024328in}}%
\pgfpathlineto{\pgfqpoint{2.738228in}{1.024682in}}%
\pgfpathlineto{\pgfqpoint{2.735129in}{1.019199in}}%
\pgfpathlineto{\pgfqpoint{2.731963in}{1.024371in}}%
\pgfpathlineto{\pgfqpoint{2.728149in}{1.020939in}}%
\pgfpathlineto{\pgfqpoint{2.722671in}{1.024298in}}%
\pgfpathlineto{\pgfqpoint{2.716657in}{1.024638in}}%
\pgfpathlineto{\pgfqpoint{2.715659in}{1.028449in}}%
\pgfpathlineto{\pgfqpoint{2.709028in}{1.030469in}}%
\pgfpathlineto{\pgfqpoint{2.705428in}{1.026285in}}%
\pgfpathlineto{\pgfqpoint{2.698517in}{1.027123in}}%
\pgfpathlineto{\pgfqpoint{2.697499in}{1.030502in}}%
\pgfpathlineto{\pgfqpoint{2.689119in}{1.033149in}}%
\pgfpathlineto{\pgfqpoint{2.679615in}{1.030971in}}%
\pgfpathlineto{\pgfqpoint{2.671365in}{1.033005in}}%
\pgfpathlineto{\pgfqpoint{2.674093in}{1.045577in}}%
\pgfpathlineto{\pgfqpoint{2.672729in}{1.056217in}}%
\pgfpathlineto{\pgfqpoint{2.679621in}{1.057585in}}%
\pgfpathlineto{\pgfqpoint{2.683306in}{1.069582in}}%
\pgfpathlineto{\pgfqpoint{2.680930in}{1.081273in}}%
\pgfpathlineto{\pgfqpoint{2.685692in}{1.079312in}}%
\pgfpathlineto{\pgfqpoint{2.691645in}{1.081236in}}%
\pgfpathlineto{\pgfqpoint{2.691815in}{1.071736in}}%
\pgfpathlineto{\pgfqpoint{2.695089in}{1.071478in}}%
\pgfpathlineto{\pgfqpoint{2.701345in}{1.062144in}}%
\pgfpathlineto{\pgfqpoint{2.723871in}{1.066185in}}%
\pgfpathlineto{\pgfqpoint{2.724911in}{1.056794in}}%
\pgfpathlineto{\pgfqpoint{2.733124in}{1.055886in}}%
\pgfpathlineto{\pgfqpoint{2.738529in}{1.052500in}}%
\pgfpathlineto{\pgfqpoint{2.738199in}{1.043269in}}%
\pgfpathlineto{\pgfqpoint{2.742724in}{1.035580in}}%
\pgfpathclose%
\pgfusepath{fill}%
\end{pgfscope}%
\begin{pgfscope}%
\pgfpathrectangle{\pgfqpoint{0.100000in}{0.100000in}}{\pgfqpoint{3.007045in}{1.925000in}}%
\pgfusepath{clip}%
\pgfsetbuttcap%
\pgfsetmiterjoin%
\definecolor{currentfill}{rgb}{0.175087,0.488812,0.733333}%
\pgfsetfillcolor{currentfill}%
\pgfsetlinewidth{0.000000pt}%
\definecolor{currentstroke}{rgb}{0.000000,0.000000,0.000000}%
\pgfsetstrokecolor{currentstroke}%
\pgfsetstrokeopacity{0.000000}%
\pgfsetdash{}{0pt}%
\pgfpathmoveto{\pgfqpoint{1.538682in}{0.941826in}}%
\pgfpathlineto{\pgfqpoint{1.532268in}{0.939536in}}%
\pgfpathlineto{\pgfqpoint{1.528739in}{0.932611in}}%
\pgfpathlineto{\pgfqpoint{1.520249in}{0.932634in}}%
\pgfpathlineto{\pgfqpoint{1.515519in}{0.937448in}}%
\pgfpathlineto{\pgfqpoint{1.510237in}{0.939913in}}%
\pgfpathlineto{\pgfqpoint{1.505202in}{0.934838in}}%
\pgfpathlineto{\pgfqpoint{1.505836in}{0.946347in}}%
\pgfpathlineto{\pgfqpoint{1.507272in}{0.975589in}}%
\pgfpathlineto{\pgfqpoint{1.509153in}{1.008633in}}%
\pgfpathlineto{\pgfqpoint{1.533400in}{1.007220in}}%
\pgfpathlineto{\pgfqpoint{1.537864in}{1.006990in}}%
\pgfpathlineto{\pgfqpoint{1.541765in}{0.998932in}}%
\pgfpathlineto{\pgfqpoint{1.547875in}{0.992381in}}%
\pgfpathlineto{\pgfqpoint{1.553847in}{0.991951in}}%
\pgfpathlineto{\pgfqpoint{1.562805in}{0.979407in}}%
\pgfpathlineto{\pgfqpoint{1.561768in}{0.950581in}}%
\pgfpathlineto{\pgfqpoint{1.539102in}{0.951802in}}%
\pgfpathlineto{\pgfqpoint{1.538682in}{0.941826in}}%
\pgfpathclose%
\pgfusepath{fill}%
\end{pgfscope}%
\begin{pgfscope}%
\pgfpathrectangle{\pgfqpoint{0.100000in}{0.100000in}}{\pgfqpoint{3.007045in}{1.925000in}}%
\pgfusepath{clip}%
\pgfsetbuttcap%
\pgfsetmiterjoin%
\definecolor{currentfill}{rgb}{0.336378,0.625513,0.806736}%
\pgfsetfillcolor{currentfill}%
\pgfsetlinewidth{0.000000pt}%
\definecolor{currentstroke}{rgb}{0.000000,0.000000,0.000000}%
\pgfsetstrokecolor{currentstroke}%
\pgfsetstrokeopacity{0.000000}%
\pgfsetdash{}{0pt}%
\pgfpathmoveto{\pgfqpoint{2.585937in}{0.924191in}}%
\pgfpathlineto{\pgfqpoint{2.551891in}{0.919345in}}%
\pgfpathlineto{\pgfqpoint{2.562238in}{0.907450in}}%
\pgfpathlineto{\pgfqpoint{2.555178in}{0.902655in}}%
\pgfpathlineto{\pgfqpoint{2.555983in}{0.898307in}}%
\pgfpathlineto{\pgfqpoint{2.550816in}{0.897393in}}%
\pgfpathlineto{\pgfqpoint{2.545027in}{0.899687in}}%
\pgfpathlineto{\pgfqpoint{2.538562in}{0.893175in}}%
\pgfpathlineto{\pgfqpoint{2.537783in}{0.898762in}}%
\pgfpathlineto{\pgfqpoint{2.508172in}{0.896146in}}%
\pgfpathlineto{\pgfqpoint{2.505466in}{0.902631in}}%
\pgfpathlineto{\pgfqpoint{2.503825in}{0.913860in}}%
\pgfpathlineto{\pgfqpoint{2.500183in}{0.925858in}}%
\pgfpathlineto{\pgfqpoint{2.504209in}{0.927233in}}%
\pgfpathlineto{\pgfqpoint{2.505332in}{0.935315in}}%
\pgfpathlineto{\pgfqpoint{2.523780in}{0.937038in}}%
\pgfpathlineto{\pgfqpoint{2.523035in}{0.944991in}}%
\pgfpathlineto{\pgfqpoint{2.524892in}{0.952370in}}%
\pgfpathlineto{\pgfqpoint{2.523665in}{0.961730in}}%
\pgfpathlineto{\pgfqpoint{2.535380in}{0.962820in}}%
\pgfpathlineto{\pgfqpoint{2.531826in}{0.973469in}}%
\pgfpathlineto{\pgfqpoint{2.533363in}{0.985590in}}%
\pgfpathlineto{\pgfqpoint{2.537459in}{0.985766in}}%
\pgfpathlineto{\pgfqpoint{2.547659in}{0.980532in}}%
\pgfpathlineto{\pgfqpoint{2.555149in}{0.977464in}}%
\pgfpathlineto{\pgfqpoint{2.562648in}{0.972115in}}%
\pgfpathlineto{\pgfqpoint{2.564813in}{0.967383in}}%
\pgfpathlineto{\pgfqpoint{2.569136in}{0.965197in}}%
\pgfpathlineto{\pgfqpoint{2.573216in}{0.959889in}}%
\pgfpathlineto{\pgfqpoint{2.571933in}{0.952085in}}%
\pgfpathlineto{\pgfqpoint{2.575191in}{0.947306in}}%
\pgfpathlineto{\pgfqpoint{2.575141in}{0.942834in}}%
\pgfpathlineto{\pgfqpoint{2.581758in}{0.944406in}}%
\pgfpathlineto{\pgfqpoint{2.586120in}{0.939582in}}%
\pgfpathlineto{\pgfqpoint{2.589091in}{0.931283in}}%
\pgfpathlineto{\pgfqpoint{2.585937in}{0.924191in}}%
\pgfpathclose%
\pgfusepath{fill}%
\end{pgfscope}%
\begin{pgfscope}%
\pgfpathrectangle{\pgfqpoint{0.100000in}{0.100000in}}{\pgfqpoint{3.007045in}{1.925000in}}%
\pgfusepath{clip}%
\pgfsetbuttcap%
\pgfsetmiterjoin%
\definecolor{currentfill}{rgb}{0.235986,0.549712,0.764706}%
\pgfsetfillcolor{currentfill}%
\pgfsetlinewidth{0.000000pt}%
\definecolor{currentstroke}{rgb}{0.000000,0.000000,0.000000}%
\pgfsetstrokecolor{currentstroke}%
\pgfsetstrokeopacity{0.000000}%
\pgfsetdash{}{0pt}%
\pgfpathmoveto{\pgfqpoint{1.712393in}{0.726751in}}%
\pgfpathlineto{\pgfqpoint{1.700123in}{0.727126in}}%
\pgfpathlineto{\pgfqpoint{1.700709in}{0.760907in}}%
\pgfpathlineto{\pgfqpoint{1.724887in}{0.764449in}}%
\pgfpathlineto{\pgfqpoint{1.725337in}{0.792070in}}%
\pgfpathlineto{\pgfqpoint{1.730486in}{0.795042in}}%
\pgfpathlineto{\pgfqpoint{1.740181in}{0.799040in}}%
\pgfpathlineto{\pgfqpoint{1.742020in}{0.795983in}}%
\pgfpathlineto{\pgfqpoint{1.749934in}{0.794479in}}%
\pgfpathlineto{\pgfqpoint{1.756132in}{0.794865in}}%
\pgfpathlineto{\pgfqpoint{1.760297in}{0.800742in}}%
\pgfpathlineto{\pgfqpoint{1.774881in}{0.791133in}}%
\pgfpathlineto{\pgfqpoint{1.779488in}{0.786480in}}%
\pgfpathlineto{\pgfqpoint{1.786299in}{0.783647in}}%
\pgfpathlineto{\pgfqpoint{1.786323in}{0.758894in}}%
\pgfpathlineto{\pgfqpoint{1.779169in}{0.763215in}}%
\pgfpathlineto{\pgfqpoint{1.767169in}{0.763420in}}%
\pgfpathlineto{\pgfqpoint{1.755274in}{0.762070in}}%
\pgfpathlineto{\pgfqpoint{1.755158in}{0.734645in}}%
\pgfpathlineto{\pgfqpoint{1.735350in}{0.734610in}}%
\pgfpathlineto{\pgfqpoint{1.736915in}{0.718690in}}%
\pgfpathlineto{\pgfqpoint{1.728058in}{0.722452in}}%
\pgfpathlineto{\pgfqpoint{1.725321in}{0.722021in}}%
\pgfpathlineto{\pgfqpoint{1.720316in}{0.726579in}}%
\pgfpathlineto{\pgfqpoint{1.712393in}{0.726751in}}%
\pgfpathclose%
\pgfusepath{fill}%
\end{pgfscope}%
\begin{pgfscope}%
\pgfpathrectangle{\pgfqpoint{0.100000in}{0.100000in}}{\pgfqpoint{3.007045in}{1.925000in}}%
\pgfusepath{clip}%
\pgfsetbuttcap%
\pgfsetmiterjoin%
\definecolor{currentfill}{rgb}{0.610980,0.787420,0.880492}%
\pgfsetfillcolor{currentfill}%
\pgfsetlinewidth{0.000000pt}%
\definecolor{currentstroke}{rgb}{0.000000,0.000000,0.000000}%
\pgfsetstrokecolor{currentstroke}%
\pgfsetstrokeopacity{0.000000}%
\pgfsetdash{}{0pt}%
\pgfpathmoveto{\pgfqpoint{2.484306in}{1.188484in}}%
\pgfpathlineto{\pgfqpoint{2.487519in}{1.178106in}}%
\pgfpathlineto{\pgfqpoint{2.485316in}{1.174226in}}%
\pgfpathlineto{\pgfqpoint{2.489719in}{1.169881in}}%
\pgfpathlineto{\pgfqpoint{2.487851in}{1.166255in}}%
\pgfpathlineto{\pgfqpoint{2.480993in}{1.155854in}}%
\pgfpathlineto{\pgfqpoint{2.466955in}{1.155210in}}%
\pgfpathlineto{\pgfqpoint{2.463450in}{1.158871in}}%
\pgfpathlineto{\pgfqpoint{2.461626in}{1.162571in}}%
\pgfpathlineto{\pgfqpoint{2.460301in}{1.177194in}}%
\pgfpathlineto{\pgfqpoint{2.463999in}{1.180859in}}%
\pgfpathlineto{\pgfqpoint{2.469733in}{1.178659in}}%
\pgfpathlineto{\pgfqpoint{2.478749in}{1.189995in}}%
\pgfpathlineto{\pgfqpoint{2.484306in}{1.188484in}}%
\pgfpathclose%
\pgfusepath{fill}%
\end{pgfscope}%
\begin{pgfscope}%
\pgfpathrectangle{\pgfqpoint{0.100000in}{0.100000in}}{\pgfqpoint{3.007045in}{1.925000in}}%
\pgfusepath{clip}%
\pgfsetbuttcap%
\pgfsetmiterjoin%
\definecolor{currentfill}{rgb}{0.117878,0.428374,0.698270}%
\pgfsetfillcolor{currentfill}%
\pgfsetlinewidth{0.000000pt}%
\definecolor{currentstroke}{rgb}{0.000000,0.000000,0.000000}%
\pgfsetstrokecolor{currentstroke}%
\pgfsetstrokeopacity{0.000000}%
\pgfsetdash{}{0pt}%
\pgfpathmoveto{\pgfqpoint{1.369270in}{1.748112in}}%
\pgfpathlineto{\pgfqpoint{1.371755in}{1.774248in}}%
\pgfpathlineto{\pgfqpoint{1.375731in}{1.814864in}}%
\pgfpathlineto{\pgfqpoint{1.424243in}{1.810279in}}%
\pgfpathlineto{\pgfqpoint{1.464391in}{1.806947in}}%
\pgfpathlineto{\pgfqpoint{1.463375in}{1.794138in}}%
\pgfpathlineto{\pgfqpoint{1.457650in}{1.794593in}}%
\pgfpathlineto{\pgfqpoint{1.457171in}{1.788804in}}%
\pgfpathlineto{\pgfqpoint{1.453609in}{1.789081in}}%
\pgfpathlineto{\pgfqpoint{1.452658in}{1.777477in}}%
\pgfpathlineto{\pgfqpoint{1.423931in}{1.779899in}}%
\pgfpathlineto{\pgfqpoint{1.422948in}{1.768339in}}%
\pgfpathlineto{\pgfqpoint{1.425483in}{1.768138in}}%
\pgfpathlineto{\pgfqpoint{1.422411in}{1.752597in}}%
\pgfpathlineto{\pgfqpoint{1.407533in}{1.753794in}}%
\pgfpathlineto{\pgfqpoint{1.403722in}{1.748968in}}%
\pgfpathlineto{\pgfqpoint{1.394175in}{1.746414in}}%
\pgfpathlineto{\pgfqpoint{1.390611in}{1.748693in}}%
\pgfpathlineto{\pgfqpoint{1.389403in}{1.754089in}}%
\pgfpathlineto{\pgfqpoint{1.386114in}{1.754885in}}%
\pgfpathlineto{\pgfqpoint{1.376795in}{1.745675in}}%
\pgfpathlineto{\pgfqpoint{1.369270in}{1.748112in}}%
\pgfpathclose%
\pgfusepath{fill}%
\end{pgfscope}%
\begin{pgfscope}%
\pgfpathrectangle{\pgfqpoint{0.100000in}{0.100000in}}{\pgfqpoint{3.007045in}{1.925000in}}%
\pgfusepath{clip}%
\pgfsetbuttcap%
\pgfsetmiterjoin%
\definecolor{currentfill}{rgb}{0.280892,0.587620,0.785083}%
\pgfsetfillcolor{currentfill}%
\pgfsetlinewidth{0.000000pt}%
\definecolor{currentstroke}{rgb}{0.000000,0.000000,0.000000}%
\pgfsetstrokecolor{currentstroke}%
\pgfsetstrokeopacity{0.000000}%
\pgfsetdash{}{0pt}%
\pgfpathmoveto{\pgfqpoint{1.330767in}{0.811238in}}%
\pgfpathlineto{\pgfqpoint{1.383417in}{0.806844in}}%
\pgfpathlineto{\pgfqpoint{1.440696in}{0.803294in}}%
\pgfpathlineto{\pgfqpoint{1.438908in}{0.774476in}}%
\pgfpathlineto{\pgfqpoint{1.437038in}{0.746267in}}%
\pgfpathlineto{\pgfqpoint{1.429539in}{0.746331in}}%
\pgfpathlineto{\pgfqpoint{1.379483in}{0.749728in}}%
\pgfpathlineto{\pgfqpoint{1.381713in}{0.778104in}}%
\pgfpathlineto{\pgfqpoint{1.327701in}{0.782533in}}%
\pgfpathlineto{\pgfqpoint{1.328991in}{0.794521in}}%
\pgfpathlineto{\pgfqpoint{1.330767in}{0.811238in}}%
\pgfpathclose%
\pgfusepath{fill}%
\end{pgfscope}%
\begin{pgfscope}%
\pgfpathrectangle{\pgfqpoint{0.100000in}{0.100000in}}{\pgfqpoint{3.007045in}{1.925000in}}%
\pgfusepath{clip}%
\pgfsetbuttcap%
\pgfsetmiterjoin%
\definecolor{currentfill}{rgb}{0.265759,0.577286,0.779177}%
\pgfsetfillcolor{currentfill}%
\pgfsetlinewidth{0.000000pt}%
\definecolor{currentstroke}{rgb}{0.000000,0.000000,0.000000}%
\pgfsetstrokecolor{currentstroke}%
\pgfsetstrokeopacity{0.000000}%
\pgfsetdash{}{0pt}%
\pgfpathmoveto{\pgfqpoint{2.637551in}{0.899069in}}%
\pgfpathlineto{\pgfqpoint{2.613555in}{0.916458in}}%
\pgfpathlineto{\pgfqpoint{2.599449in}{0.926351in}}%
\pgfpathlineto{\pgfqpoint{2.585937in}{0.924191in}}%
\pgfpathlineto{\pgfqpoint{2.589091in}{0.931283in}}%
\pgfpathlineto{\pgfqpoint{2.586120in}{0.939582in}}%
\pgfpathlineto{\pgfqpoint{2.581758in}{0.944406in}}%
\pgfpathlineto{\pgfqpoint{2.575141in}{0.942834in}}%
\pgfpathlineto{\pgfqpoint{2.575191in}{0.947306in}}%
\pgfpathlineto{\pgfqpoint{2.571933in}{0.952085in}}%
\pgfpathlineto{\pgfqpoint{2.573216in}{0.959889in}}%
\pgfpathlineto{\pgfqpoint{2.569136in}{0.965197in}}%
\pgfpathlineto{\pgfqpoint{2.564813in}{0.967383in}}%
\pgfpathlineto{\pgfqpoint{2.597898in}{0.973667in}}%
\pgfpathlineto{\pgfqpoint{2.608759in}{0.975733in}}%
\pgfpathlineto{\pgfqpoint{2.608908in}{0.972810in}}%
\pgfpathlineto{\pgfqpoint{2.618329in}{0.960663in}}%
\pgfpathlineto{\pgfqpoint{2.625940in}{0.956912in}}%
\pgfpathlineto{\pgfqpoint{2.641093in}{0.964571in}}%
\pgfpathlineto{\pgfqpoint{2.650737in}{0.964955in}}%
\pgfpathlineto{\pgfqpoint{2.650566in}{0.955387in}}%
\pgfpathlineto{\pgfqpoint{2.652361in}{0.948046in}}%
\pgfpathlineto{\pgfqpoint{2.661874in}{0.940850in}}%
\pgfpathlineto{\pgfqpoint{2.645656in}{0.937482in}}%
\pgfpathlineto{\pgfqpoint{2.640458in}{0.935526in}}%
\pgfpathlineto{\pgfqpoint{2.647207in}{0.927059in}}%
\pgfpathlineto{\pgfqpoint{2.647369in}{0.914791in}}%
\pgfpathlineto{\pgfqpoint{2.641884in}{0.909457in}}%
\pgfpathlineto{\pgfqpoint{2.637551in}{0.899069in}}%
\pgfpathclose%
\pgfusepath{fill}%
\end{pgfscope}%
\begin{pgfscope}%
\pgfpathrectangle{\pgfqpoint{0.100000in}{0.100000in}}{\pgfqpoint{3.007045in}{1.925000in}}%
\pgfusepath{clip}%
\pgfsetbuttcap%
\pgfsetmiterjoin%
\definecolor{currentfill}{rgb}{0.391865,0.663406,0.828389}%
\pgfsetfillcolor{currentfill}%
\pgfsetlinewidth{0.000000pt}%
\definecolor{currentstroke}{rgb}{0.000000,0.000000,0.000000}%
\pgfsetstrokecolor{currentstroke}%
\pgfsetstrokeopacity{0.000000}%
\pgfsetdash{}{0pt}%
\pgfpathmoveto{\pgfqpoint{1.653431in}{1.223538in}}%
\pgfpathlineto{\pgfqpoint{1.652796in}{1.200602in}}%
\pgfpathlineto{\pgfqpoint{1.607157in}{1.202049in}}%
\pgfpathlineto{\pgfqpoint{1.608008in}{1.224954in}}%
\pgfpathlineto{\pgfqpoint{1.630558in}{1.224197in}}%
\pgfpathlineto{\pgfqpoint{1.653431in}{1.223538in}}%
\pgfpathclose%
\pgfusepath{fill}%
\end{pgfscope}%
\begin{pgfscope}%
\pgfpathrectangle{\pgfqpoint{0.100000in}{0.100000in}}{\pgfqpoint{3.007045in}{1.925000in}}%
\pgfusepath{clip}%
\pgfsetbuttcap%
\pgfsetmiterjoin%
\definecolor{currentfill}{rgb}{0.454118,0.701300,0.846659}%
\pgfsetfillcolor{currentfill}%
\pgfsetlinewidth{0.000000pt}%
\definecolor{currentstroke}{rgb}{0.000000,0.000000,0.000000}%
\pgfsetstrokecolor{currentstroke}%
\pgfsetstrokeopacity{0.000000}%
\pgfsetdash{}{0pt}%
\pgfpathmoveto{\pgfqpoint{2.108338in}{0.951144in}}%
\pgfpathlineto{\pgfqpoint{2.098903in}{0.955455in}}%
\pgfpathlineto{\pgfqpoint{2.093449in}{0.959312in}}%
\pgfpathlineto{\pgfqpoint{2.083172in}{0.958841in}}%
\pgfpathlineto{\pgfqpoint{2.065703in}{0.958280in}}%
\pgfpathlineto{\pgfqpoint{2.068381in}{0.961543in}}%
\pgfpathlineto{\pgfqpoint{2.071414in}{0.971616in}}%
\pgfpathlineto{\pgfqpoint{2.068022in}{0.977336in}}%
\pgfpathlineto{\pgfqpoint{2.070320in}{0.985858in}}%
\pgfpathlineto{\pgfqpoint{2.077777in}{0.982593in}}%
\pgfpathlineto{\pgfqpoint{2.080971in}{0.989204in}}%
\pgfpathlineto{\pgfqpoint{2.078631in}{0.992336in}}%
\pgfpathlineto{\pgfqpoint{2.082341in}{0.997150in}}%
\pgfpathlineto{\pgfqpoint{2.098746in}{0.997437in}}%
\pgfpathlineto{\pgfqpoint{2.097982in}{1.008842in}}%
\pgfpathlineto{\pgfqpoint{2.115261in}{1.009765in}}%
\pgfpathlineto{\pgfqpoint{2.116101in}{0.997143in}}%
\pgfpathlineto{\pgfqpoint{2.117091in}{0.980740in}}%
\pgfpathlineto{\pgfqpoint{2.115651in}{0.980661in}}%
\pgfpathlineto{\pgfqpoint{2.116626in}{0.955211in}}%
\pgfpathlineto{\pgfqpoint{2.110066in}{0.954874in}}%
\pgfpathlineto{\pgfqpoint{2.108338in}{0.951144in}}%
\pgfpathclose%
\pgfusepath{fill}%
\end{pgfscope}%
\begin{pgfscope}%
\pgfpathrectangle{\pgfqpoint{0.100000in}{0.100000in}}{\pgfqpoint{3.007045in}{1.925000in}}%
\pgfusepath{clip}%
\pgfsetbuttcap%
\pgfsetmiterjoin%
\definecolor{currentfill}{rgb}{0.472941,0.711634,0.850719}%
\pgfsetfillcolor{currentfill}%
\pgfsetlinewidth{0.000000pt}%
\definecolor{currentstroke}{rgb}{0.000000,0.000000,0.000000}%
\pgfsetstrokecolor{currentstroke}%
\pgfsetstrokeopacity{0.000000}%
\pgfsetdash{}{0pt}%
\pgfpathmoveto{\pgfqpoint{1.184259in}{1.303392in}}%
\pgfpathlineto{\pgfqpoint{1.210701in}{1.299766in}}%
\pgfpathlineto{\pgfqpoint{1.213860in}{1.325677in}}%
\pgfpathlineto{\pgfqpoint{1.226347in}{1.324123in}}%
\pgfpathlineto{\pgfqpoint{1.231571in}{1.369429in}}%
\pgfpathlineto{\pgfqpoint{1.234424in}{1.392114in}}%
\pgfpathlineto{\pgfqpoint{1.259419in}{1.388949in}}%
\pgfpathlineto{\pgfqpoint{1.256175in}{1.385798in}}%
\pgfpathlineto{\pgfqpoint{1.255653in}{1.381012in}}%
\pgfpathlineto{\pgfqpoint{1.267095in}{1.378734in}}%
\pgfpathlineto{\pgfqpoint{1.270231in}{1.387736in}}%
\pgfpathlineto{\pgfqpoint{1.272590in}{1.387487in}}%
\pgfpathlineto{\pgfqpoint{1.269866in}{1.358232in}}%
\pgfpathlineto{\pgfqpoint{1.267098in}{1.336682in}}%
\pgfpathlineto{\pgfqpoint{1.262254in}{1.293547in}}%
\pgfpathlineto{\pgfqpoint{1.217135in}{1.298862in}}%
\pgfpathlineto{\pgfqpoint{1.216879in}{1.294670in}}%
\pgfpathlineto{\pgfqpoint{1.222250in}{1.286251in}}%
\pgfpathlineto{\pgfqpoint{1.229691in}{1.263388in}}%
\pgfpathlineto{\pgfqpoint{1.222855in}{1.255180in}}%
\pgfpathlineto{\pgfqpoint{1.220648in}{1.256986in}}%
\pgfpathlineto{\pgfqpoint{1.213800in}{1.255082in}}%
\pgfpathlineto{\pgfqpoint{1.205670in}{1.255974in}}%
\pgfpathlineto{\pgfqpoint{1.192494in}{1.261100in}}%
\pgfpathlineto{\pgfqpoint{1.189706in}{1.265642in}}%
\pgfpathlineto{\pgfqpoint{1.188332in}{1.274642in}}%
\pgfpathlineto{\pgfqpoint{1.193426in}{1.288196in}}%
\pgfpathlineto{\pgfqpoint{1.192988in}{1.292052in}}%
\pgfpathlineto{\pgfqpoint{1.183883in}{1.298271in}}%
\pgfpathlineto{\pgfqpoint{1.184259in}{1.303392in}}%
\pgfpathclose%
\pgfusepath{fill}%
\end{pgfscope}%
\begin{pgfscope}%
\pgfpathrectangle{\pgfqpoint{0.100000in}{0.100000in}}{\pgfqpoint{3.007045in}{1.925000in}}%
\pgfusepath{clip}%
\pgfsetbuttcap%
\pgfsetmiterjoin%
\definecolor{currentfill}{rgb}{0.548235,0.752972,0.866959}%
\pgfsetfillcolor{currentfill}%
\pgfsetlinewidth{0.000000pt}%
\definecolor{currentstroke}{rgb}{0.000000,0.000000,0.000000}%
\pgfsetstrokecolor{currentstroke}%
\pgfsetstrokeopacity{0.000000}%
\pgfsetdash{}{0pt}%
\pgfpathmoveto{\pgfqpoint{2.362681in}{1.482619in}}%
\pgfpathlineto{\pgfqpoint{2.338611in}{1.478397in}}%
\pgfpathlineto{\pgfqpoint{2.342275in}{1.455683in}}%
\pgfpathlineto{\pgfqpoint{2.330370in}{1.453796in}}%
\pgfpathlineto{\pgfqpoint{2.331232in}{1.448101in}}%
\pgfpathlineto{\pgfqpoint{2.326140in}{1.446645in}}%
\pgfpathlineto{\pgfqpoint{2.303651in}{1.443639in}}%
\pgfpathlineto{\pgfqpoint{2.304504in}{1.437926in}}%
\pgfpathlineto{\pgfqpoint{2.283546in}{1.435118in}}%
\pgfpathlineto{\pgfqpoint{2.280842in}{1.457266in}}%
\pgfpathlineto{\pgfqpoint{2.269523in}{1.455942in}}%
\pgfpathlineto{\pgfqpoint{2.264263in}{1.501550in}}%
\pgfpathlineto{\pgfqpoint{2.298344in}{1.505689in}}%
\pgfpathlineto{\pgfqpoint{2.313368in}{1.507664in}}%
\pgfpathlineto{\pgfqpoint{2.313689in}{1.503876in}}%
\pgfpathlineto{\pgfqpoint{2.305190in}{1.494752in}}%
\pgfpathlineto{\pgfqpoint{2.302323in}{1.494776in}}%
\pgfpathlineto{\pgfqpoint{2.298978in}{1.488134in}}%
\pgfpathlineto{\pgfqpoint{2.298537in}{1.477920in}}%
\pgfpathlineto{\pgfqpoint{2.300991in}{1.473532in}}%
\pgfpathlineto{\pgfqpoint{2.313383in}{1.469662in}}%
\pgfpathlineto{\pgfqpoint{2.329237in}{1.497116in}}%
\pgfpathlineto{\pgfqpoint{2.339139in}{1.501127in}}%
\pgfpathlineto{\pgfqpoint{2.344111in}{1.505414in}}%
\pgfpathlineto{\pgfqpoint{2.353364in}{1.501784in}}%
\pgfpathlineto{\pgfqpoint{2.360186in}{1.491680in}}%
\pgfpathlineto{\pgfqpoint{2.362681in}{1.482619in}}%
\pgfpathclose%
\pgfusepath{fill}%
\end{pgfscope}%
\begin{pgfscope}%
\pgfpathrectangle{\pgfqpoint{0.100000in}{0.100000in}}{\pgfqpoint{3.007045in}{1.925000in}}%
\pgfusepath{clip}%
\pgfsetbuttcap%
\pgfsetmiterjoin%
\definecolor{currentfill}{rgb}{0.523137,0.739193,0.861546}%
\pgfsetfillcolor{currentfill}%
\pgfsetlinewidth{0.000000pt}%
\definecolor{currentstroke}{rgb}{0.000000,0.000000,0.000000}%
\pgfsetstrokecolor{currentstroke}%
\pgfsetstrokeopacity{0.000000}%
\pgfsetdash{}{0pt}%
\pgfpathmoveto{\pgfqpoint{2.558875in}{1.172717in}}%
\pgfpathlineto{\pgfqpoint{2.558583in}{1.161971in}}%
\pgfpathlineto{\pgfqpoint{2.553841in}{1.152723in}}%
\pgfpathlineto{\pgfqpoint{2.555307in}{1.150523in}}%
\pgfpathlineto{\pgfqpoint{2.549531in}{1.143698in}}%
\pgfpathlineto{\pgfqpoint{2.548827in}{1.135761in}}%
\pgfpathlineto{\pgfqpoint{2.538421in}{1.132193in}}%
\pgfpathlineto{\pgfqpoint{2.533411in}{1.132073in}}%
\pgfpathlineto{\pgfqpoint{2.527646in}{1.135652in}}%
\pgfpathlineto{\pgfqpoint{2.526634in}{1.142930in}}%
\pgfpathlineto{\pgfqpoint{2.526995in}{1.150325in}}%
\pgfpathlineto{\pgfqpoint{2.530752in}{1.154455in}}%
\pgfpathlineto{\pgfqpoint{2.532425in}{1.163870in}}%
\pgfpathlineto{\pgfqpoint{2.525737in}{1.174087in}}%
\pgfpathlineto{\pgfqpoint{2.527409in}{1.177624in}}%
\pgfpathlineto{\pgfqpoint{2.532999in}{1.178683in}}%
\pgfpathlineto{\pgfqpoint{2.532223in}{1.188646in}}%
\pgfpathlineto{\pgfqpoint{2.534247in}{1.194214in}}%
\pgfpathlineto{\pgfqpoint{2.523965in}{1.201691in}}%
\pgfpathlineto{\pgfqpoint{2.525644in}{1.210611in}}%
\pgfpathlineto{\pgfqpoint{2.532780in}{1.212166in}}%
\pgfpathlineto{\pgfqpoint{2.538153in}{1.216728in}}%
\pgfpathlineto{\pgfqpoint{2.543719in}{1.212838in}}%
\pgfpathlineto{\pgfqpoint{2.549426in}{1.216536in}}%
\pgfpathlineto{\pgfqpoint{2.560339in}{1.213292in}}%
\pgfpathlineto{\pgfqpoint{2.567300in}{1.215781in}}%
\pgfpathlineto{\pgfqpoint{2.569988in}{1.214538in}}%
\pgfpathlineto{\pgfqpoint{2.568471in}{1.207938in}}%
\pgfpathlineto{\pgfqpoint{2.571223in}{1.207045in}}%
\pgfpathlineto{\pgfqpoint{2.569924in}{1.199083in}}%
\pgfpathlineto{\pgfqpoint{2.562770in}{1.192250in}}%
\pgfpathlineto{\pgfqpoint{2.563877in}{1.187087in}}%
\pgfpathlineto{\pgfqpoint{2.558875in}{1.172717in}}%
\pgfpathclose%
\pgfusepath{fill}%
\end{pgfscope}%
\begin{pgfscope}%
\pgfpathrectangle{\pgfqpoint{0.100000in}{0.100000in}}{\pgfqpoint{3.007045in}{1.925000in}}%
\pgfusepath{clip}%
\pgfsetbuttcap%
\pgfsetmiterjoin%
\definecolor{currentfill}{rgb}{0.541961,0.749527,0.865606}%
\pgfsetfillcolor{currentfill}%
\pgfsetlinewidth{0.000000pt}%
\definecolor{currentstroke}{rgb}{0.000000,0.000000,0.000000}%
\pgfsetstrokecolor{currentstroke}%
\pgfsetstrokeopacity{0.000000}%
\pgfsetdash{}{0pt}%
\pgfpathmoveto{\pgfqpoint{2.697553in}{1.455914in}}%
\pgfpathlineto{\pgfqpoint{2.679109in}{1.451464in}}%
\pgfpathlineto{\pgfqpoint{2.675608in}{1.465295in}}%
\pgfpathlineto{\pgfqpoint{2.656543in}{1.461343in}}%
\pgfpathlineto{\pgfqpoint{2.652395in}{1.466521in}}%
\pgfpathlineto{\pgfqpoint{2.648634in}{1.483094in}}%
\pgfpathlineto{\pgfqpoint{2.648812in}{1.490538in}}%
\pgfpathlineto{\pgfqpoint{2.644415in}{1.507385in}}%
\pgfpathlineto{\pgfqpoint{2.647780in}{1.512793in}}%
\pgfpathlineto{\pgfqpoint{2.656259in}{1.521763in}}%
\pgfpathlineto{\pgfqpoint{2.664991in}{1.524305in}}%
\pgfpathlineto{\pgfqpoint{2.664388in}{1.534234in}}%
\pgfpathlineto{\pgfqpoint{2.672204in}{1.537694in}}%
\pgfpathlineto{\pgfqpoint{2.684161in}{1.539019in}}%
\pgfpathlineto{\pgfqpoint{2.688054in}{1.525089in}}%
\pgfpathlineto{\pgfqpoint{2.699199in}{1.524078in}}%
\pgfpathlineto{\pgfqpoint{2.716233in}{1.541048in}}%
\pgfpathlineto{\pgfqpoint{2.706539in}{1.571514in}}%
\pgfpathlineto{\pgfqpoint{2.712187in}{1.569616in}}%
\pgfpathlineto{\pgfqpoint{2.721573in}{1.573055in}}%
\pgfpathlineto{\pgfqpoint{2.733723in}{1.536267in}}%
\pgfpathlineto{\pgfqpoint{2.731549in}{1.525832in}}%
\pgfpathlineto{\pgfqpoint{2.739601in}{1.524017in}}%
\pgfpathlineto{\pgfqpoint{2.742037in}{1.517084in}}%
\pgfpathlineto{\pgfqpoint{2.739690in}{1.510307in}}%
\pgfpathlineto{\pgfqpoint{2.742514in}{1.505188in}}%
\pgfpathlineto{\pgfqpoint{2.743391in}{1.496243in}}%
\pgfpathlineto{\pgfqpoint{2.737431in}{1.497269in}}%
\pgfpathlineto{\pgfqpoint{2.737175in}{1.492249in}}%
\pgfpathlineto{\pgfqpoint{2.726801in}{1.495572in}}%
\pgfpathlineto{\pgfqpoint{2.722365in}{1.491084in}}%
\pgfpathlineto{\pgfqpoint{2.719947in}{1.482951in}}%
\pgfpathlineto{\pgfqpoint{2.692078in}{1.475538in}}%
\pgfpathlineto{\pgfqpoint{2.697553in}{1.455914in}}%
\pgfpathclose%
\pgfusepath{fill}%
\end{pgfscope}%
\begin{pgfscope}%
\pgfpathrectangle{\pgfqpoint{0.100000in}{0.100000in}}{\pgfqpoint{3.007045in}{1.925000in}}%
\pgfusepath{clip}%
\pgfsetbuttcap%
\pgfsetmiterjoin%
\definecolor{currentfill}{rgb}{0.321246,0.615179,0.800830}%
\pgfsetfillcolor{currentfill}%
\pgfsetlinewidth{0.000000pt}%
\definecolor{currentstroke}{rgb}{0.000000,0.000000,0.000000}%
\pgfsetstrokecolor{currentstroke}%
\pgfsetstrokeopacity{0.000000}%
\pgfsetdash{}{0pt}%
\pgfpathmoveto{\pgfqpoint{1.779241in}{1.344556in}}%
\pgfpathlineto{\pgfqpoint{1.813145in}{1.344757in}}%
\pgfpathlineto{\pgfqpoint{1.813376in}{1.321894in}}%
\pgfpathlineto{\pgfqpoint{1.784860in}{1.321668in}}%
\pgfpathlineto{\pgfqpoint{1.744960in}{1.321802in}}%
\pgfpathlineto{\pgfqpoint{1.739290in}{1.321824in}}%
\pgfpathlineto{\pgfqpoint{1.739541in}{1.344779in}}%
\pgfpathlineto{\pgfqpoint{1.779241in}{1.344556in}}%
\pgfpathclose%
\pgfusepath{fill}%
\end{pgfscope}%
\begin{pgfscope}%
\pgfpathrectangle{\pgfqpoint{0.100000in}{0.100000in}}{\pgfqpoint{3.007045in}{1.925000in}}%
\pgfusepath{clip}%
\pgfsetbuttcap%
\pgfsetmiterjoin%
\definecolor{currentfill}{rgb}{0.516863,0.735748,0.860192}%
\pgfsetfillcolor{currentfill}%
\pgfsetlinewidth{0.000000pt}%
\definecolor{currentstroke}{rgb}{0.000000,0.000000,0.000000}%
\pgfsetstrokecolor{currentstroke}%
\pgfsetstrokeopacity{0.000000}%
\pgfsetdash{}{0pt}%
\pgfpathmoveto{\pgfqpoint{2.222669in}{0.734886in}}%
\pgfpathlineto{\pgfqpoint{2.223953in}{0.726994in}}%
\pgfpathlineto{\pgfqpoint{2.230267in}{0.714046in}}%
\pgfpathlineto{\pgfqpoint{2.225729in}{0.705930in}}%
\pgfpathlineto{\pgfqpoint{2.226858in}{0.694238in}}%
\pgfpathlineto{\pgfqpoint{2.211673in}{0.692874in}}%
\pgfpathlineto{\pgfqpoint{2.202255in}{0.702095in}}%
\pgfpathlineto{\pgfqpoint{2.193997in}{0.705851in}}%
\pgfpathlineto{\pgfqpoint{2.190855in}{0.708411in}}%
\pgfpathlineto{\pgfqpoint{2.189892in}{0.719957in}}%
\pgfpathlineto{\pgfqpoint{2.178513in}{0.718958in}}%
\pgfpathlineto{\pgfqpoint{2.173613in}{0.721470in}}%
\pgfpathlineto{\pgfqpoint{2.174689in}{0.726922in}}%
\pgfpathlineto{\pgfqpoint{2.171759in}{0.737367in}}%
\pgfpathlineto{\pgfqpoint{2.172078in}{0.746133in}}%
\pgfpathlineto{\pgfqpoint{2.175967in}{0.749888in}}%
\pgfpathlineto{\pgfqpoint{2.176652in}{0.754493in}}%
\pgfpathlineto{\pgfqpoint{2.169944in}{0.753927in}}%
\pgfpathlineto{\pgfqpoint{2.168941in}{0.762907in}}%
\pgfpathlineto{\pgfqpoint{2.166832in}{0.787387in}}%
\pgfpathlineto{\pgfqpoint{2.176418in}{0.787936in}}%
\pgfpathlineto{\pgfqpoint{2.177866in}{0.793913in}}%
\pgfpathlineto{\pgfqpoint{2.192221in}{0.794560in}}%
\pgfpathlineto{\pgfqpoint{2.198442in}{0.789193in}}%
\pgfpathlineto{\pgfqpoint{2.197889in}{0.784340in}}%
\pgfpathlineto{\pgfqpoint{2.204331in}{0.776956in}}%
\pgfpathlineto{\pgfqpoint{2.213183in}{0.772767in}}%
\pgfpathlineto{\pgfqpoint{2.213720in}{0.767453in}}%
\pgfpathlineto{\pgfqpoint{2.217222in}{0.763282in}}%
\pgfpathlineto{\pgfqpoint{2.222615in}{0.760540in}}%
\pgfpathlineto{\pgfqpoint{2.223920in}{0.746448in}}%
\pgfpathlineto{\pgfqpoint{2.215985in}{0.745794in}}%
\pgfpathlineto{\pgfqpoint{2.217078in}{0.734332in}}%
\pgfpathlineto{\pgfqpoint{2.222669in}{0.734886in}}%
\pgfpathclose%
\pgfusepath{fill}%
\end{pgfscope}%
\begin{pgfscope}%
\pgfpathrectangle{\pgfqpoint{0.100000in}{0.100000in}}{\pgfqpoint{3.007045in}{1.925000in}}%
\pgfusepath{clip}%
\pgfsetbuttcap%
\pgfsetmiterjoin%
\definecolor{currentfill}{rgb}{0.435294,0.690965,0.842599}%
\pgfsetfillcolor{currentfill}%
\pgfsetlinewidth{0.000000pt}%
\definecolor{currentstroke}{rgb}{0.000000,0.000000,0.000000}%
\pgfsetstrokecolor{currentstroke}%
\pgfsetstrokeopacity{0.000000}%
\pgfsetdash{}{0pt}%
\pgfpathmoveto{\pgfqpoint{1.869963in}{1.345597in}}%
\pgfpathlineto{\pgfqpoint{1.858723in}{1.345342in}}%
\pgfpathlineto{\pgfqpoint{1.847419in}{1.345161in}}%
\pgfpathlineto{\pgfqpoint{1.847146in}{1.362415in}}%
\pgfpathlineto{\pgfqpoint{1.845372in}{1.362394in}}%
\pgfpathlineto{\pgfqpoint{1.844917in}{1.391266in}}%
\pgfpathlineto{\pgfqpoint{1.867818in}{1.391658in}}%
\pgfpathlineto{\pgfqpoint{1.868339in}{1.362827in}}%
\pgfpathlineto{\pgfqpoint{1.869522in}{1.362845in}}%
\pgfpathlineto{\pgfqpoint{1.869963in}{1.345597in}}%
\pgfpathclose%
\pgfusepath{fill}%
\end{pgfscope}%
\begin{pgfscope}%
\pgfpathrectangle{\pgfqpoint{0.100000in}{0.100000in}}{\pgfqpoint{3.007045in}{1.925000in}}%
\pgfusepath{clip}%
\pgfsetbuttcap%
\pgfsetmiterjoin%
\definecolor{currentfill}{rgb}{0.676817,0.816471,0.902376}%
\pgfsetfillcolor{currentfill}%
\pgfsetlinewidth{0.000000pt}%
\definecolor{currentstroke}{rgb}{0.000000,0.000000,0.000000}%
\pgfsetstrokecolor{currentstroke}%
\pgfsetstrokeopacity{0.000000}%
\pgfsetdash{}{0pt}%
\pgfpathmoveto{\pgfqpoint{2.283546in}{1.435118in}}%
\pgfpathlineto{\pgfqpoint{2.283625in}{1.434436in}}%
\pgfpathlineto{\pgfqpoint{2.261163in}{1.431916in}}%
\pgfpathlineto{\pgfqpoint{2.258163in}{1.454659in}}%
\pgfpathlineto{\pgfqpoint{2.246778in}{1.453408in}}%
\pgfpathlineto{\pgfqpoint{2.241654in}{1.499264in}}%
\pgfpathlineto{\pgfqpoint{2.252766in}{1.500298in}}%
\pgfpathlineto{\pgfqpoint{2.264263in}{1.501550in}}%
\pgfpathlineto{\pgfqpoint{2.269523in}{1.455942in}}%
\pgfpathlineto{\pgfqpoint{2.280842in}{1.457266in}}%
\pgfpathlineto{\pgfqpoint{2.283546in}{1.435118in}}%
\pgfpathclose%
\pgfusepath{fill}%
\end{pgfscope}%
\begin{pgfscope}%
\pgfpathrectangle{\pgfqpoint{0.100000in}{0.100000in}}{\pgfqpoint{3.007045in}{1.925000in}}%
\pgfusepath{clip}%
\pgfsetbuttcap%
\pgfsetmiterjoin%
\definecolor{currentfill}{rgb}{0.730950,0.839477,0.921323}%
\pgfsetfillcolor{currentfill}%
\pgfsetlinewidth{0.000000pt}%
\definecolor{currentstroke}{rgb}{0.000000,0.000000,0.000000}%
\pgfsetstrokecolor{currentstroke}%
\pgfsetstrokeopacity{0.000000}%
\pgfsetdash{}{0pt}%
\pgfpathmoveto{\pgfqpoint{2.418790in}{0.913875in}}%
\pgfpathlineto{\pgfqpoint{2.416370in}{0.915795in}}%
\pgfpathlineto{\pgfqpoint{2.421150in}{0.927599in}}%
\pgfpathlineto{\pgfqpoint{2.420735in}{0.932481in}}%
\pgfpathlineto{\pgfqpoint{2.405701in}{0.945045in}}%
\pgfpathlineto{\pgfqpoint{2.403998in}{0.954773in}}%
\pgfpathlineto{\pgfqpoint{2.399943in}{0.955887in}}%
\pgfpathlineto{\pgfqpoint{2.404400in}{0.961065in}}%
\pgfpathlineto{\pgfqpoint{2.414244in}{0.964014in}}%
\pgfpathlineto{\pgfqpoint{2.416656in}{0.974681in}}%
\pgfpathlineto{\pgfqpoint{2.429667in}{0.984471in}}%
\pgfpathlineto{\pgfqpoint{2.431676in}{0.978202in}}%
\pgfpathlineto{\pgfqpoint{2.437163in}{0.979641in}}%
\pgfpathlineto{\pgfqpoint{2.437960in}{0.976572in}}%
\pgfpathlineto{\pgfqpoint{2.446969in}{0.969925in}}%
\pgfpathlineto{\pgfqpoint{2.454474in}{0.954214in}}%
\pgfpathlineto{\pgfqpoint{2.459120in}{0.952751in}}%
\pgfpathlineto{\pgfqpoint{2.455786in}{0.951722in}}%
\pgfpathlineto{\pgfqpoint{2.453980in}{0.946316in}}%
\pgfpathlineto{\pgfqpoint{2.455447in}{0.943242in}}%
\pgfpathlineto{\pgfqpoint{2.451924in}{0.935651in}}%
\pgfpathlineto{\pgfqpoint{2.452389in}{0.929337in}}%
\pgfpathlineto{\pgfqpoint{2.435961in}{0.922177in}}%
\pgfpathlineto{\pgfqpoint{2.418790in}{0.913875in}}%
\pgfpathclose%
\pgfusepath{fill}%
\end{pgfscope}%
\begin{pgfscope}%
\pgfpathrectangle{\pgfqpoint{0.100000in}{0.100000in}}{\pgfqpoint{3.007045in}{1.925000in}}%
\pgfusepath{clip}%
\pgfsetbuttcap%
\pgfsetmiterjoin%
\definecolor{currentfill}{rgb}{0.529412,0.742637,0.862899}%
\pgfsetfillcolor{currentfill}%
\pgfsetlinewidth{0.000000pt}%
\definecolor{currentstroke}{rgb}{0.000000,0.000000,0.000000}%
\pgfsetstrokecolor{currentstroke}%
\pgfsetstrokeopacity{0.000000}%
\pgfsetdash{}{0pt}%
\pgfpathmoveto{\pgfqpoint{1.910330in}{0.935711in}}%
\pgfpathlineto{\pgfqpoint{1.910553in}{0.923358in}}%
\pgfpathlineto{\pgfqpoint{1.899149in}{0.923195in}}%
\pgfpathlineto{\pgfqpoint{1.899200in}{0.918414in}}%
\pgfpathlineto{\pgfqpoint{1.889632in}{0.918440in}}%
\pgfpathlineto{\pgfqpoint{1.851607in}{0.918465in}}%
\pgfpathlineto{\pgfqpoint{1.851596in}{0.920385in}}%
\pgfpathlineto{\pgfqpoint{1.851232in}{0.928170in}}%
\pgfpathlineto{\pgfqpoint{1.855032in}{0.933900in}}%
\pgfpathlineto{\pgfqpoint{1.855176in}{0.941374in}}%
\pgfpathlineto{\pgfqpoint{1.849514in}{0.944275in}}%
\pgfpathlineto{\pgfqpoint{1.845759in}{0.949170in}}%
\pgfpathlineto{\pgfqpoint{1.843659in}{0.955855in}}%
\pgfpathlineto{\pgfqpoint{1.832421in}{0.955990in}}%
\pgfpathlineto{\pgfqpoint{1.832387in}{0.968604in}}%
\pgfpathlineto{\pgfqpoint{1.886046in}{0.969526in}}%
\pgfpathlineto{\pgfqpoint{1.884426in}{0.969017in}}%
\pgfpathlineto{\pgfqpoint{1.884698in}{0.944172in}}%
\pgfpathlineto{\pgfqpoint{1.887586in}{0.941336in}}%
\pgfpathlineto{\pgfqpoint{1.910317in}{0.941444in}}%
\pgfpathlineto{\pgfqpoint{1.910330in}{0.935711in}}%
\pgfpathclose%
\pgfusepath{fill}%
\end{pgfscope}%
\begin{pgfscope}%
\pgfpathrectangle{\pgfqpoint{0.100000in}{0.100000in}}{\pgfqpoint{3.007045in}{1.925000in}}%
\pgfusepath{clip}%
\pgfsetbuttcap%
\pgfsetmiterjoin%
\definecolor{currentfill}{rgb}{0.435294,0.690965,0.842599}%
\pgfsetfillcolor{currentfill}%
\pgfsetlinewidth{0.000000pt}%
\definecolor{currentstroke}{rgb}{0.000000,0.000000,0.000000}%
\pgfsetstrokecolor{currentstroke}%
\pgfsetstrokeopacity{0.000000}%
\pgfsetdash{}{0pt}%
\pgfpathmoveto{\pgfqpoint{2.122132in}{1.377026in}}%
\pgfpathlineto{\pgfqpoint{2.078237in}{1.373879in}}%
\pgfpathlineto{\pgfqpoint{2.046233in}{1.372218in}}%
\pgfpathlineto{\pgfqpoint{2.044679in}{1.394919in}}%
\pgfpathlineto{\pgfqpoint{2.061841in}{1.396171in}}%
\pgfpathlineto{\pgfqpoint{2.073275in}{1.396646in}}%
\pgfpathlineto{\pgfqpoint{2.119150in}{1.400006in}}%
\pgfpathlineto{\pgfqpoint{2.122786in}{1.396116in}}%
\pgfpathlineto{\pgfqpoint{2.120693in}{1.385146in}}%
\pgfpathlineto{\pgfqpoint{2.122132in}{1.377026in}}%
\pgfpathclose%
\pgfusepath{fill}%
\end{pgfscope}%
\begin{pgfscope}%
\pgfpathrectangle{\pgfqpoint{0.100000in}{0.100000in}}{\pgfqpoint{3.007045in}{1.925000in}}%
\pgfusepath{clip}%
\pgfsetbuttcap%
\pgfsetmiterjoin%
\definecolor{currentfill}{rgb}{0.138547,0.452272,0.714510}%
\pgfsetfillcolor{currentfill}%
\pgfsetlinewidth{0.000000pt}%
\definecolor{currentstroke}{rgb}{0.000000,0.000000,0.000000}%
\pgfsetstrokecolor{currentstroke}%
\pgfsetstrokeopacity{0.000000}%
\pgfsetdash{}{0pt}%
\pgfpathmoveto{\pgfqpoint{1.495136in}{0.800124in}}%
\pgfpathlineto{\pgfqpoint{1.469347in}{0.801604in}}%
\pgfpathlineto{\pgfqpoint{1.471341in}{0.833174in}}%
\pgfpathlineto{\pgfqpoint{1.476764in}{0.832809in}}%
\pgfpathlineto{\pgfqpoint{1.478574in}{0.861428in}}%
\pgfpathlineto{\pgfqpoint{1.501084in}{0.860066in}}%
\pgfpathlineto{\pgfqpoint{1.500565in}{0.847796in}}%
\pgfpathlineto{\pgfqpoint{1.499337in}{0.825611in}}%
\pgfpathlineto{\pgfqpoint{1.496660in}{0.826118in}}%
\pgfpathlineto{\pgfqpoint{1.495136in}{0.800124in}}%
\pgfpathclose%
\pgfusepath{fill}%
\end{pgfscope}%
\begin{pgfscope}%
\pgfpathrectangle{\pgfqpoint{0.100000in}{0.100000in}}{\pgfqpoint{3.007045in}{1.925000in}}%
\pgfusepath{clip}%
\pgfsetbuttcap%
\pgfsetmiterjoin%
\definecolor{currentfill}{rgb}{0.356555,0.639293,0.814610}%
\pgfsetfillcolor{currentfill}%
\pgfsetlinewidth{0.000000pt}%
\definecolor{currentstroke}{rgb}{0.000000,0.000000,0.000000}%
\pgfsetstrokecolor{currentstroke}%
\pgfsetstrokeopacity{0.000000}%
\pgfsetdash{}{0pt}%
\pgfpathmoveto{\pgfqpoint{1.965268in}{0.666642in}}%
\pgfpathlineto{\pgfqpoint{1.965330in}{0.661044in}}%
\pgfpathlineto{\pgfqpoint{1.962071in}{0.656269in}}%
\pgfpathlineto{\pgfqpoint{1.959044in}{0.656761in}}%
\pgfpathlineto{\pgfqpoint{1.954469in}{0.643065in}}%
\pgfpathlineto{\pgfqpoint{1.949979in}{0.638422in}}%
\pgfpathlineto{\pgfqpoint{1.952131in}{0.625210in}}%
\pgfpathlineto{\pgfqpoint{1.942513in}{0.628295in}}%
\pgfpathlineto{\pgfqpoint{1.933963in}{0.633867in}}%
\pgfpathlineto{\pgfqpoint{1.933198in}{0.637846in}}%
\pgfpathlineto{\pgfqpoint{1.926639in}{0.641691in}}%
\pgfpathlineto{\pgfqpoint{1.922010in}{0.647995in}}%
\pgfpathlineto{\pgfqpoint{1.921110in}{0.659462in}}%
\pgfpathlineto{\pgfqpoint{1.923673in}{0.668219in}}%
\pgfpathlineto{\pgfqpoint{1.946656in}{0.668963in}}%
\pgfpathlineto{\pgfqpoint{1.949500in}{0.671101in}}%
\pgfpathlineto{\pgfqpoint{1.953769in}{0.665525in}}%
\pgfpathlineto{\pgfqpoint{1.963857in}{0.672381in}}%
\pgfpathlineto{\pgfqpoint{1.965268in}{0.666642in}}%
\pgfpathclose%
\pgfusepath{fill}%
\end{pgfscope}%
\begin{pgfscope}%
\pgfpathrectangle{\pgfqpoint{0.100000in}{0.100000in}}{\pgfqpoint{3.007045in}{1.925000in}}%
\pgfusepath{clip}%
\pgfsetbuttcap%
\pgfsetmiterjoin%
\definecolor{currentfill}{rgb}{0.510588,0.732303,0.858839}%
\pgfsetfillcolor{currentfill}%
\pgfsetlinewidth{0.000000pt}%
\definecolor{currentstroke}{rgb}{0.000000,0.000000,0.000000}%
\pgfsetstrokecolor{currentstroke}%
\pgfsetstrokeopacity{0.000000}%
\pgfsetdash{}{0pt}%
\pgfpathmoveto{\pgfqpoint{2.198442in}{0.789193in}}%
\pgfpathlineto{\pgfqpoint{2.192221in}{0.794560in}}%
\pgfpathlineto{\pgfqpoint{2.177866in}{0.793913in}}%
\pgfpathlineto{\pgfqpoint{2.176418in}{0.787936in}}%
\pgfpathlineto{\pgfqpoint{2.166832in}{0.787387in}}%
\pgfpathlineto{\pgfqpoint{2.142921in}{0.786166in}}%
\pgfpathlineto{\pgfqpoint{2.143299in}{0.800190in}}%
\pgfpathlineto{\pgfqpoint{2.144570in}{0.838481in}}%
\pgfpathlineto{\pgfqpoint{2.179696in}{0.840159in}}%
\pgfpathlineto{\pgfqpoint{2.202564in}{0.841744in}}%
\pgfpathlineto{\pgfqpoint{2.204539in}{0.819210in}}%
\pgfpathlineto{\pgfqpoint{2.206348in}{0.816022in}}%
\pgfpathlineto{\pgfqpoint{2.213109in}{0.813406in}}%
\pgfpathlineto{\pgfqpoint{2.210327in}{0.807475in}}%
\pgfpathlineto{\pgfqpoint{2.209929in}{0.801633in}}%
\pgfpathlineto{\pgfqpoint{2.207356in}{0.797594in}}%
\pgfpathlineto{\pgfqpoint{2.202677in}{0.796254in}}%
\pgfpathlineto{\pgfqpoint{2.198442in}{0.789193in}}%
\pgfpathclose%
\pgfusepath{fill}%
\end{pgfscope}%
\begin{pgfscope}%
\pgfpathrectangle{\pgfqpoint{0.100000in}{0.100000in}}{\pgfqpoint{3.007045in}{1.925000in}}%
\pgfusepath{clip}%
\pgfsetbuttcap%
\pgfsetmiterjoin%
\definecolor{currentfill}{rgb}{0.735871,0.841569,0.923045}%
\pgfsetfillcolor{currentfill}%
\pgfsetlinewidth{0.000000pt}%
\definecolor{currentstroke}{rgb}{0.000000,0.000000,0.000000}%
\pgfsetstrokecolor{currentstroke}%
\pgfsetstrokeopacity{0.000000}%
\pgfsetdash{}{0pt}%
\pgfpathmoveto{\pgfqpoint{2.474213in}{0.983871in}}%
\pgfpathlineto{\pgfqpoint{2.480170in}{0.978949in}}%
\pgfpathlineto{\pgfqpoint{2.489274in}{0.973776in}}%
\pgfpathlineto{\pgfqpoint{2.494999in}{0.975158in}}%
\pgfpathlineto{\pgfqpoint{2.498902in}{0.973883in}}%
\pgfpathlineto{\pgfqpoint{2.492448in}{0.960338in}}%
\pgfpathlineto{\pgfqpoint{2.486422in}{0.958962in}}%
\pgfpathlineto{\pgfqpoint{2.480332in}{0.961447in}}%
\pgfpathlineto{\pgfqpoint{2.476332in}{0.956183in}}%
\pgfpathlineto{\pgfqpoint{2.469542in}{0.954025in}}%
\pgfpathlineto{\pgfqpoint{2.459120in}{0.952751in}}%
\pgfpathlineto{\pgfqpoint{2.454474in}{0.954214in}}%
\pgfpathlineto{\pgfqpoint{2.446969in}{0.969925in}}%
\pgfpathlineto{\pgfqpoint{2.437960in}{0.976572in}}%
\pgfpathlineto{\pgfqpoint{2.437163in}{0.979641in}}%
\pgfpathlineto{\pgfqpoint{2.441498in}{0.987564in}}%
\pgfpathlineto{\pgfqpoint{2.450485in}{0.993442in}}%
\pgfpathlineto{\pgfqpoint{2.458453in}{0.990986in}}%
\pgfpathlineto{\pgfqpoint{2.462402in}{0.980829in}}%
\pgfpathlineto{\pgfqpoint{2.465463in}{0.979232in}}%
\pgfpathlineto{\pgfqpoint{2.468967in}{0.985644in}}%
\pgfpathlineto{\pgfqpoint{2.474213in}{0.983871in}}%
\pgfpathclose%
\pgfusepath{fill}%
\end{pgfscope}%
\begin{pgfscope}%
\pgfpathrectangle{\pgfqpoint{0.100000in}{0.100000in}}{\pgfqpoint{3.007045in}{1.925000in}}%
\pgfusepath{clip}%
\pgfsetbuttcap%
\pgfsetmiterjoin%
\definecolor{currentfill}{rgb}{0.808212,0.879985,0.947835}%
\pgfsetfillcolor{currentfill}%
\pgfsetlinewidth{0.000000pt}%
\definecolor{currentstroke}{rgb}{0.000000,0.000000,0.000000}%
\pgfsetstrokecolor{currentstroke}%
\pgfsetstrokeopacity{0.000000}%
\pgfsetdash{}{0pt}%
\pgfpathmoveto{\pgfqpoint{0.558722in}{1.943515in}}%
\pgfpathlineto{\pgfqpoint{0.550057in}{1.949005in}}%
\pgfpathlineto{\pgfqpoint{0.546058in}{1.955336in}}%
\pgfpathlineto{\pgfqpoint{0.547224in}{1.961309in}}%
\pgfpathlineto{\pgfqpoint{0.552633in}{1.957537in}}%
\pgfpathlineto{\pgfqpoint{0.559107in}{1.963705in}}%
\pgfpathlineto{\pgfqpoint{0.566457in}{1.957886in}}%
\pgfpathlineto{\pgfqpoint{0.558722in}{1.943515in}}%
\pgfpathclose%
\pgfusepath{fill}%
\end{pgfscope}%
\begin{pgfscope}%
\pgfpathrectangle{\pgfqpoint{0.100000in}{0.100000in}}{\pgfqpoint{3.007045in}{1.925000in}}%
\pgfusepath{clip}%
\pgfsetbuttcap%
\pgfsetmiterjoin%
\definecolor{currentfill}{rgb}{0.447843,0.697855,0.845306}%
\pgfsetfillcolor{currentfill}%
\pgfsetlinewidth{0.000000pt}%
\definecolor{currentstroke}{rgb}{0.000000,0.000000,0.000000}%
\pgfsetstrokecolor{currentstroke}%
\pgfsetstrokeopacity{0.000000}%
\pgfsetdash{}{0pt}%
\pgfpathmoveto{\pgfqpoint{1.386813in}{0.624056in}}%
\pgfpathlineto{\pgfqpoint{1.415290in}{0.621980in}}%
\pgfpathlineto{\pgfqpoint{1.433065in}{0.620988in}}%
\pgfpathlineto{\pgfqpoint{1.431514in}{0.595996in}}%
\pgfpathlineto{\pgfqpoint{1.429685in}{0.568203in}}%
\pgfpathlineto{\pgfqpoint{1.383908in}{0.571430in}}%
\pgfpathlineto{\pgfqpoint{1.389923in}{0.575693in}}%
\pgfpathlineto{\pgfqpoint{1.385588in}{0.580234in}}%
\pgfpathlineto{\pgfqpoint{1.391278in}{0.588321in}}%
\pgfpathlineto{\pgfqpoint{1.390664in}{0.594316in}}%
\pgfpathlineto{\pgfqpoint{1.386971in}{0.595647in}}%
\pgfpathlineto{\pgfqpoint{1.364071in}{0.597523in}}%
\pgfpathlineto{\pgfqpoint{1.363729in}{0.593646in}}%
\pgfpathlineto{\pgfqpoint{1.352050in}{0.594662in}}%
\pgfpathlineto{\pgfqpoint{1.350354in}{0.573776in}}%
\pgfpathlineto{\pgfqpoint{1.337515in}{0.574775in}}%
\pgfpathlineto{\pgfqpoint{1.336243in}{0.559459in}}%
\pgfpathlineto{\pgfqpoint{1.282209in}{0.612003in}}%
\pgfpathlineto{\pgfqpoint{1.318428in}{0.649073in}}%
\pgfpathlineto{\pgfqpoint{1.322677in}{0.647278in}}%
\pgfpathlineto{\pgfqpoint{1.328596in}{0.641304in}}%
\pgfpathlineto{\pgfqpoint{1.331766in}{0.642740in}}%
\pgfpathlineto{\pgfqpoint{1.337100in}{0.644492in}}%
\pgfpathlineto{\pgfqpoint{1.349725in}{0.635825in}}%
\pgfpathlineto{\pgfqpoint{1.351088in}{0.627381in}}%
\pgfpathlineto{\pgfqpoint{1.386813in}{0.624056in}}%
\pgfpathclose%
\pgfusepath{fill}%
\end{pgfscope}%
\begin{pgfscope}%
\pgfpathrectangle{\pgfqpoint{0.100000in}{0.100000in}}{\pgfqpoint{3.007045in}{1.925000in}}%
\pgfusepath{clip}%
\pgfsetbuttcap%
\pgfsetmiterjoin%
\definecolor{currentfill}{rgb}{0.321246,0.615179,0.800830}%
\pgfsetfillcolor{currentfill}%
\pgfsetlinewidth{0.000000pt}%
\definecolor{currentstroke}{rgb}{0.000000,0.000000,0.000000}%
\pgfsetstrokecolor{currentstroke}%
\pgfsetstrokeopacity{0.000000}%
\pgfsetdash{}{0pt}%
\pgfpathmoveto{\pgfqpoint{1.930951in}{0.750384in}}%
\pgfpathlineto{\pgfqpoint{1.932837in}{0.741766in}}%
\pgfpathlineto{\pgfqpoint{1.935072in}{0.740260in}}%
\pgfpathlineto{\pgfqpoint{1.898651in}{0.739626in}}%
\pgfpathlineto{\pgfqpoint{1.884006in}{0.739484in}}%
\pgfpathlineto{\pgfqpoint{1.884045in}{0.763316in}}%
\pgfpathlineto{\pgfqpoint{1.876274in}{0.763349in}}%
\pgfpathlineto{\pgfqpoint{1.876321in}{0.768157in}}%
\pgfpathlineto{\pgfqpoint{1.876541in}{0.789626in}}%
\pgfpathlineto{\pgfqpoint{1.884971in}{0.791103in}}%
\pgfpathlineto{\pgfqpoint{1.887657in}{0.795131in}}%
\pgfpathlineto{\pgfqpoint{1.887940in}{0.814920in}}%
\pgfpathlineto{\pgfqpoint{1.899320in}{0.814784in}}%
\pgfpathlineto{\pgfqpoint{1.909437in}{0.814697in}}%
\pgfpathlineto{\pgfqpoint{1.918185in}{0.809351in}}%
\pgfpathlineto{\pgfqpoint{1.910795in}{0.808969in}}%
\pgfpathlineto{\pgfqpoint{1.911801in}{0.803402in}}%
\pgfpathlineto{\pgfqpoint{1.918485in}{0.794539in}}%
\pgfpathlineto{\pgfqpoint{1.920677in}{0.772374in}}%
\pgfpathlineto{\pgfqpoint{1.917305in}{0.764130in}}%
\pgfpathlineto{\pgfqpoint{1.918474in}{0.757862in}}%
\pgfpathlineto{\pgfqpoint{1.921911in}{0.757915in}}%
\pgfpathlineto{\pgfqpoint{1.930951in}{0.750384in}}%
\pgfpathclose%
\pgfusepath{fill}%
\end{pgfscope}%
\begin{pgfscope}%
\pgfpathrectangle{\pgfqpoint{0.100000in}{0.100000in}}{\pgfqpoint{3.007045in}{1.925000in}}%
\pgfusepath{clip}%
\pgfsetbuttcap%
\pgfsetmiterjoin%
\definecolor{currentfill}{rgb}{0.290980,0.594510,0.789020}%
\pgfsetfillcolor{currentfill}%
\pgfsetlinewidth{0.000000pt}%
\definecolor{currentstroke}{rgb}{0.000000,0.000000,0.000000}%
\pgfsetstrokecolor{currentstroke}%
\pgfsetstrokeopacity{0.000000}%
\pgfsetdash{}{0pt}%
\pgfpathmoveto{\pgfqpoint{1.512812in}{1.275327in}}%
\pgfpathlineto{\pgfqpoint{1.511547in}{1.252543in}}%
\pgfpathlineto{\pgfqpoint{1.483897in}{1.254208in}}%
\pgfpathlineto{\pgfqpoint{1.460453in}{1.255605in}}%
\pgfpathlineto{\pgfqpoint{1.461934in}{1.278597in}}%
\pgfpathlineto{\pgfqpoint{1.460924in}{1.278668in}}%
\pgfpathlineto{\pgfqpoint{1.462537in}{1.301517in}}%
\pgfpathlineto{\pgfqpoint{1.455754in}{1.302023in}}%
\pgfpathlineto{\pgfqpoint{1.457408in}{1.324892in}}%
\pgfpathlineto{\pgfqpoint{1.485080in}{1.322839in}}%
\pgfpathlineto{\pgfqpoint{1.514112in}{1.321020in}}%
\pgfpathlineto{\pgfqpoint{1.512812in}{1.275327in}}%
\pgfpathclose%
\pgfusepath{fill}%
\end{pgfscope}%
\begin{pgfscope}%
\pgfpathrectangle{\pgfqpoint{0.100000in}{0.100000in}}{\pgfqpoint{3.007045in}{1.925000in}}%
\pgfusepath{clip}%
\pgfsetbuttcap%
\pgfsetmiterjoin%
\definecolor{currentfill}{rgb}{0.412042,0.677186,0.836263}%
\pgfsetfillcolor{currentfill}%
\pgfsetlinewidth{0.000000pt}%
\definecolor{currentstroke}{rgb}{0.000000,0.000000,0.000000}%
\pgfsetstrokecolor{currentstroke}%
\pgfsetstrokeopacity{0.000000}%
\pgfsetdash{}{0pt}%
\pgfpathmoveto{\pgfqpoint{2.345734in}{1.216099in}}%
\pgfpathlineto{\pgfqpoint{2.346250in}{1.205103in}}%
\pgfpathlineto{\pgfqpoint{2.330635in}{1.204428in}}%
\pgfpathlineto{\pgfqpoint{2.314937in}{1.203868in}}%
\pgfpathlineto{\pgfqpoint{2.305334in}{1.201390in}}%
\pgfpathlineto{\pgfqpoint{2.288393in}{1.199350in}}%
\pgfpathlineto{\pgfqpoint{2.285375in}{1.228022in}}%
\pgfpathlineto{\pgfqpoint{2.283471in}{1.248041in}}%
\pgfpathlineto{\pgfqpoint{2.283146in}{1.250837in}}%
\pgfpathlineto{\pgfqpoint{2.301562in}{1.253044in}}%
\pgfpathlineto{\pgfqpoint{2.306160in}{1.255224in}}%
\pgfpathlineto{\pgfqpoint{2.305349in}{1.261913in}}%
\pgfpathlineto{\pgfqpoint{2.322524in}{1.264076in}}%
\pgfpathlineto{\pgfqpoint{2.322114in}{1.267499in}}%
\pgfpathlineto{\pgfqpoint{2.327732in}{1.268420in}}%
\pgfpathlineto{\pgfqpoint{2.345898in}{1.268387in}}%
\pgfpathlineto{\pgfqpoint{2.346617in}{1.250238in}}%
\pgfpathlineto{\pgfqpoint{2.349506in}{1.250356in}}%
\pgfpathlineto{\pgfqpoint{2.350594in}{1.236213in}}%
\pgfpathlineto{\pgfqpoint{2.348588in}{1.219904in}}%
\pgfpathlineto{\pgfqpoint{2.345734in}{1.216099in}}%
\pgfpathclose%
\pgfusepath{fill}%
\end{pgfscope}%
\begin{pgfscope}%
\pgfpathrectangle{\pgfqpoint{0.100000in}{0.100000in}}{\pgfqpoint{3.007045in}{1.925000in}}%
\pgfusepath{clip}%
\pgfsetbuttcap%
\pgfsetmiterjoin%
\definecolor{currentfill}{rgb}{0.391865,0.663406,0.828389}%
\pgfsetfillcolor{currentfill}%
\pgfsetlinewidth{0.000000pt}%
\definecolor{currentstroke}{rgb}{0.000000,0.000000,0.000000}%
\pgfsetstrokecolor{currentstroke}%
\pgfsetstrokeopacity{0.000000}%
\pgfsetdash{}{0pt}%
\pgfpathmoveto{\pgfqpoint{1.925598in}{1.218425in}}%
\pgfpathlineto{\pgfqpoint{1.926181in}{1.198155in}}%
\pgfpathlineto{\pgfqpoint{1.906383in}{1.197788in}}%
\pgfpathlineto{\pgfqpoint{1.906198in}{1.203318in}}%
\pgfpathlineto{\pgfqpoint{1.880879in}{1.202651in}}%
\pgfpathlineto{\pgfqpoint{1.880439in}{1.202637in}}%
\pgfpathlineto{\pgfqpoint{1.879993in}{1.222750in}}%
\pgfpathlineto{\pgfqpoint{1.888599in}{1.223016in}}%
\pgfpathlineto{\pgfqpoint{1.889157in}{1.227613in}}%
\pgfpathlineto{\pgfqpoint{1.886676in}{1.239179in}}%
\pgfpathlineto{\pgfqpoint{1.890498in}{1.239370in}}%
\pgfpathlineto{\pgfqpoint{1.925211in}{1.241392in}}%
\pgfpathlineto{\pgfqpoint{1.925598in}{1.218425in}}%
\pgfpathclose%
\pgfusepath{fill}%
\end{pgfscope}%
\begin{pgfscope}%
\pgfpathrectangle{\pgfqpoint{0.100000in}{0.100000in}}{\pgfqpoint{3.007045in}{1.925000in}}%
\pgfusepath{clip}%
\pgfsetbuttcap%
\pgfsetmiterjoin%
\definecolor{currentfill}{rgb}{0.316201,0.611734,0.798862}%
\pgfsetfillcolor{currentfill}%
\pgfsetlinewidth{0.000000pt}%
\definecolor{currentstroke}{rgb}{0.000000,0.000000,0.000000}%
\pgfsetstrokecolor{currentstroke}%
\pgfsetstrokeopacity{0.000000}%
\pgfsetdash{}{0pt}%
\pgfpathmoveto{\pgfqpoint{1.842321in}{0.853020in}}%
\pgfpathlineto{\pgfqpoint{1.833670in}{0.850751in}}%
\pgfpathlineto{\pgfqpoint{1.830433in}{0.847601in}}%
\pgfpathlineto{\pgfqpoint{1.825651in}{0.849322in}}%
\pgfpathlineto{\pgfqpoint{1.809383in}{0.849499in}}%
\pgfpathlineto{\pgfqpoint{1.801936in}{0.851532in}}%
\pgfpathlineto{\pgfqpoint{1.801665in}{0.836897in}}%
\pgfpathlineto{\pgfqpoint{1.775725in}{0.836774in}}%
\pgfpathlineto{\pgfqpoint{1.775709in}{0.848268in}}%
\pgfpathlineto{\pgfqpoint{1.744379in}{0.848364in}}%
\pgfpathlineto{\pgfqpoint{1.744358in}{0.842618in}}%
\pgfpathlineto{\pgfqpoint{1.735784in}{0.842654in}}%
\pgfpathlineto{\pgfqpoint{1.724353in}{0.842730in}}%
\pgfpathlineto{\pgfqpoint{1.724413in}{0.848474in}}%
\pgfpathlineto{\pgfqpoint{1.713003in}{0.848605in}}%
\pgfpathlineto{\pgfqpoint{1.713476in}{0.872989in}}%
\pgfpathlineto{\pgfqpoint{1.719231in}{0.879638in}}%
\pgfpathlineto{\pgfqpoint{1.725397in}{0.881888in}}%
\pgfpathlineto{\pgfqpoint{1.732284in}{0.882217in}}%
\pgfpathlineto{\pgfqpoint{1.741580in}{0.885648in}}%
\pgfpathlineto{\pgfqpoint{1.751555in}{0.890916in}}%
\pgfpathlineto{\pgfqpoint{1.759980in}{0.886994in}}%
\pgfpathlineto{\pgfqpoint{1.763675in}{0.893735in}}%
\pgfpathlineto{\pgfqpoint{1.768002in}{0.897857in}}%
\pgfpathlineto{\pgfqpoint{1.765215in}{0.904084in}}%
\pgfpathlineto{\pgfqpoint{1.765492in}{0.911520in}}%
\pgfpathlineto{\pgfqpoint{1.800606in}{0.911576in}}%
\pgfpathlineto{\pgfqpoint{1.799463in}{0.919527in}}%
\pgfpathlineto{\pgfqpoint{1.819183in}{0.919101in}}%
\pgfpathlineto{\pgfqpoint{1.830635in}{0.919840in}}%
\pgfpathlineto{\pgfqpoint{1.826781in}{0.913265in}}%
\pgfpathlineto{\pgfqpoint{1.821996in}{0.913307in}}%
\pgfpathlineto{\pgfqpoint{1.821906in}{0.907584in}}%
\pgfpathlineto{\pgfqpoint{1.824480in}{0.903095in}}%
\pgfpathlineto{\pgfqpoint{1.822064in}{0.899000in}}%
\pgfpathlineto{\pgfqpoint{1.821794in}{0.883783in}}%
\pgfpathlineto{\pgfqpoint{1.818790in}{0.879030in}}%
\pgfpathlineto{\pgfqpoint{1.818712in}{0.873240in}}%
\pgfpathlineto{\pgfqpoint{1.822375in}{0.871354in}}%
\pgfpathlineto{\pgfqpoint{1.842357in}{0.871145in}}%
\pgfpathlineto{\pgfqpoint{1.842321in}{0.853020in}}%
\pgfpathclose%
\pgfusepath{fill}%
\end{pgfscope}%
\begin{pgfscope}%
\pgfpathrectangle{\pgfqpoint{0.100000in}{0.100000in}}{\pgfqpoint{3.007045in}{1.925000in}}%
\pgfusepath{clip}%
\pgfsetbuttcap%
\pgfsetmiterjoin%
\definecolor{currentfill}{rgb}{0.781638,0.862268,0.938977}%
\pgfsetfillcolor{currentfill}%
\pgfsetlinewidth{0.000000pt}%
\definecolor{currentstroke}{rgb}{0.000000,0.000000,0.000000}%
\pgfsetstrokecolor{currentstroke}%
\pgfsetstrokeopacity{0.000000}%
\pgfsetdash{}{0pt}%
\pgfpathmoveto{\pgfqpoint{2.994174in}{1.647565in}}%
\pgfpathlineto{\pgfqpoint{3.000400in}{1.646818in}}%
\pgfpathlineto{\pgfqpoint{2.997849in}{1.641810in}}%
\pgfpathlineto{\pgfqpoint{2.994174in}{1.647565in}}%
\pgfpathclose%
\pgfusepath{fill}%
\end{pgfscope}%
\begin{pgfscope}%
\pgfpathrectangle{\pgfqpoint{0.100000in}{0.100000in}}{\pgfqpoint{3.007045in}{1.925000in}}%
\pgfusepath{clip}%
\pgfsetbuttcap%
\pgfsetmiterjoin%
\definecolor{currentfill}{rgb}{0.781638,0.862268,0.938977}%
\pgfsetfillcolor{currentfill}%
\pgfsetlinewidth{0.000000pt}%
\definecolor{currentstroke}{rgb}{0.000000,0.000000,0.000000}%
\pgfsetstrokecolor{currentstroke}%
\pgfsetstrokeopacity{0.000000}%
\pgfsetdash{}{0pt}%
\pgfpathmoveto{\pgfqpoint{2.975595in}{1.630839in}}%
\pgfpathlineto{\pgfqpoint{2.978270in}{1.638556in}}%
\pgfpathlineto{\pgfqpoint{2.972759in}{1.647244in}}%
\pgfpathlineto{\pgfqpoint{2.969629in}{1.646428in}}%
\pgfpathlineto{\pgfqpoint{2.965671in}{1.652396in}}%
\pgfpathlineto{\pgfqpoint{2.962617in}{1.653501in}}%
\pgfpathlineto{\pgfqpoint{2.963183in}{1.661995in}}%
\pgfpathlineto{\pgfqpoint{2.962312in}{1.673903in}}%
\pgfpathlineto{\pgfqpoint{2.960139in}{1.680504in}}%
\pgfpathlineto{\pgfqpoint{2.966515in}{1.680747in}}%
\pgfpathlineto{\pgfqpoint{2.956012in}{1.701938in}}%
\pgfpathlineto{\pgfqpoint{2.945232in}{1.694510in}}%
\pgfpathlineto{\pgfqpoint{2.928393in}{1.726358in}}%
\pgfpathlineto{\pgfqpoint{2.930414in}{1.732814in}}%
\pgfpathlineto{\pgfqpoint{2.925595in}{1.735765in}}%
\pgfpathlineto{\pgfqpoint{2.926062in}{1.747145in}}%
\pgfpathlineto{\pgfqpoint{2.922774in}{1.754622in}}%
\pgfpathlineto{\pgfqpoint{2.915056in}{1.781252in}}%
\pgfpathlineto{\pgfqpoint{2.912260in}{1.792925in}}%
\pgfpathlineto{\pgfqpoint{2.951724in}{1.804308in}}%
\pgfpathlineto{\pgfqpoint{2.955175in}{1.793081in}}%
\pgfpathlineto{\pgfqpoint{2.972249in}{1.797222in}}%
\pgfpathlineto{\pgfqpoint{2.981003in}{1.769545in}}%
\pgfpathlineto{\pgfqpoint{2.987728in}{1.746399in}}%
\pgfpathlineto{\pgfqpoint{2.988875in}{1.746178in}}%
\pgfpathlineto{\pgfqpoint{3.003662in}{1.755298in}}%
\pgfpathlineto{\pgfqpoint{3.015744in}{1.733166in}}%
\pgfpathlineto{\pgfqpoint{3.010773in}{1.730803in}}%
\pgfpathlineto{\pgfqpoint{3.019840in}{1.712826in}}%
\pgfpathlineto{\pgfqpoint{3.014639in}{1.710091in}}%
\pgfpathlineto{\pgfqpoint{3.025401in}{1.689136in}}%
\pgfpathlineto{\pgfqpoint{3.031696in}{1.678779in}}%
\pgfpathlineto{\pgfqpoint{3.027762in}{1.672640in}}%
\pgfpathlineto{\pgfqpoint{3.022621in}{1.680721in}}%
\pgfpathlineto{\pgfqpoint{3.016503in}{1.677395in}}%
\pgfpathlineto{\pgfqpoint{3.023371in}{1.670491in}}%
\pgfpathlineto{\pgfqpoint{3.017940in}{1.661449in}}%
\pgfpathlineto{\pgfqpoint{3.012026in}{1.665983in}}%
\pgfpathlineto{\pgfqpoint{3.012283in}{1.670432in}}%
\pgfpathlineto{\pgfqpoint{3.008113in}{1.674442in}}%
\pgfpathlineto{\pgfqpoint{3.005865in}{1.669939in}}%
\pgfpathlineto{\pgfqpoint{3.005023in}{1.661197in}}%
\pgfpathlineto{\pgfqpoint{2.996716in}{1.662802in}}%
\pgfpathlineto{\pgfqpoint{2.986928in}{1.668315in}}%
\pgfpathlineto{\pgfqpoint{2.984644in}{1.665649in}}%
\pgfpathlineto{\pgfqpoint{2.987952in}{1.658494in}}%
\pgfpathlineto{\pgfqpoint{2.985414in}{1.647330in}}%
\pgfpathlineto{\pgfqpoint{2.987840in}{1.640394in}}%
\pgfpathlineto{\pgfqpoint{2.984073in}{1.632614in}}%
\pgfpathlineto{\pgfqpoint{2.975595in}{1.630839in}}%
\pgfpathclose%
\pgfusepath{fill}%
\end{pgfscope}%
\begin{pgfscope}%
\pgfpathrectangle{\pgfqpoint{0.100000in}{0.100000in}}{\pgfqpoint{3.007045in}{1.925000in}}%
\pgfusepath{clip}%
\pgfsetbuttcap%
\pgfsetmiterjoin%
\definecolor{currentfill}{rgb}{0.516863,0.735748,0.860192}%
\pgfsetfillcolor{currentfill}%
\pgfsetlinewidth{0.000000pt}%
\definecolor{currentstroke}{rgb}{0.000000,0.000000,0.000000}%
\pgfsetstrokecolor{currentstroke}%
\pgfsetstrokeopacity{0.000000}%
\pgfsetdash{}{0pt}%
\pgfpathmoveto{\pgfqpoint{1.566327in}{0.575160in}}%
\pgfpathlineto{\pgfqpoint{1.544968in}{0.575956in}}%
\pgfpathlineto{\pgfqpoint{1.546174in}{0.604074in}}%
\pgfpathlineto{\pgfqpoint{1.575754in}{0.602872in}}%
\pgfpathlineto{\pgfqpoint{1.576374in}{0.610043in}}%
\pgfpathlineto{\pgfqpoint{1.606454in}{0.609350in}}%
\pgfpathlineto{\pgfqpoint{1.610932in}{0.600948in}}%
\pgfpathlineto{\pgfqpoint{1.602998in}{0.592878in}}%
\pgfpathlineto{\pgfqpoint{1.593576in}{0.574175in}}%
\pgfpathlineto{\pgfqpoint{1.592881in}{0.569267in}}%
\pgfpathlineto{\pgfqpoint{1.576662in}{0.574910in}}%
\pgfpathlineto{\pgfqpoint{1.566327in}{0.575160in}}%
\pgfpathclose%
\pgfusepath{fill}%
\end{pgfscope}%
\begin{pgfscope}%
\pgfpathrectangle{\pgfqpoint{0.100000in}{0.100000in}}{\pgfqpoint{3.007045in}{1.925000in}}%
\pgfusepath{clip}%
\pgfsetbuttcap%
\pgfsetmiterjoin%
\definecolor{currentfill}{rgb}{0.351511,0.635848,0.812641}%
\pgfsetfillcolor{currentfill}%
\pgfsetlinewidth{0.000000pt}%
\definecolor{currentstroke}{rgb}{0.000000,0.000000,0.000000}%
\pgfsetstrokecolor{currentstroke}%
\pgfsetstrokeopacity{0.000000}%
\pgfsetdash{}{0pt}%
\pgfpathmoveto{\pgfqpoint{1.985200in}{0.741858in}}%
\pgfpathlineto{\pgfqpoint{1.989679in}{0.740203in}}%
\pgfpathlineto{\pgfqpoint{1.990829in}{0.734444in}}%
\pgfpathlineto{\pgfqpoint{1.986138in}{0.729844in}}%
\pgfpathlineto{\pgfqpoint{1.985938in}{0.725835in}}%
\pgfpathlineto{\pgfqpoint{1.991838in}{0.723800in}}%
\pgfpathlineto{\pgfqpoint{1.986961in}{0.717838in}}%
\pgfpathlineto{\pgfqpoint{1.988496in}{0.714303in}}%
\pgfpathlineto{\pgfqpoint{1.993233in}{0.713829in}}%
\pgfpathlineto{\pgfqpoint{1.988311in}{0.711075in}}%
\pgfpathlineto{\pgfqpoint{1.970224in}{0.710328in}}%
\pgfpathlineto{\pgfqpoint{1.970799in}{0.713279in}}%
\pgfpathlineto{\pgfqpoint{1.962387in}{0.712972in}}%
\pgfpathlineto{\pgfqpoint{1.959700in}{0.720849in}}%
\pgfpathlineto{\pgfqpoint{1.964704in}{0.732579in}}%
\pgfpathlineto{\pgfqpoint{1.967255in}{0.733621in}}%
\pgfpathlineto{\pgfqpoint{1.970228in}{0.741380in}}%
\pgfpathlineto{\pgfqpoint{1.985200in}{0.741858in}}%
\pgfpathclose%
\pgfusepath{fill}%
\end{pgfscope}%
\begin{pgfscope}%
\pgfpathrectangle{\pgfqpoint{0.100000in}{0.100000in}}{\pgfqpoint{3.007045in}{1.925000in}}%
\pgfusepath{clip}%
\pgfsetbuttcap%
\pgfsetmiterjoin%
\definecolor{currentfill}{rgb}{0.523137,0.739193,0.861546}%
\pgfsetfillcolor{currentfill}%
\pgfsetlinewidth{0.000000pt}%
\definecolor{currentstroke}{rgb}{0.000000,0.000000,0.000000}%
\pgfsetstrokecolor{currentstroke}%
\pgfsetstrokeopacity{0.000000}%
\pgfsetdash{}{0pt}%
\pgfpathmoveto{\pgfqpoint{1.760890in}{0.592713in}}%
\pgfpathlineto{\pgfqpoint{1.757986in}{0.598171in}}%
\pgfpathlineto{\pgfqpoint{1.753626in}{0.595060in}}%
\pgfpathlineto{\pgfqpoint{1.751572in}{0.571734in}}%
\pgfpathlineto{\pgfqpoint{1.737969in}{0.571808in}}%
\pgfpathlineto{\pgfqpoint{1.724776in}{0.579988in}}%
\pgfpathlineto{\pgfqpoint{1.723044in}{0.595566in}}%
\pgfpathlineto{\pgfqpoint{1.705582in}{0.593025in}}%
\pgfpathlineto{\pgfqpoint{1.701562in}{0.603092in}}%
\pgfpathlineto{\pgfqpoint{1.716755in}{0.610772in}}%
\pgfpathlineto{\pgfqpoint{1.728806in}{0.610818in}}%
\pgfpathlineto{\pgfqpoint{1.728210in}{0.613654in}}%
\pgfpathlineto{\pgfqpoint{1.731113in}{0.622526in}}%
\pgfpathlineto{\pgfqpoint{1.734214in}{0.624182in}}%
\pgfpathlineto{\pgfqpoint{1.730423in}{0.638033in}}%
\pgfpathlineto{\pgfqpoint{1.735433in}{0.640512in}}%
\pgfpathlineto{\pgfqpoint{1.756860in}{0.643794in}}%
\pgfpathlineto{\pgfqpoint{1.763851in}{0.642226in}}%
\pgfpathlineto{\pgfqpoint{1.766321in}{0.636474in}}%
\pgfpathlineto{\pgfqpoint{1.772144in}{0.632693in}}%
\pgfpathlineto{\pgfqpoint{1.779925in}{0.639399in}}%
\pgfpathlineto{\pgfqpoint{1.787706in}{0.634802in}}%
\pgfpathlineto{\pgfqpoint{1.800480in}{0.631756in}}%
\pgfpathlineto{\pgfqpoint{1.804231in}{0.625251in}}%
\pgfpathlineto{\pgfqpoint{1.810103in}{0.619771in}}%
\pgfpathlineto{\pgfqpoint{1.821877in}{0.611209in}}%
\pgfpathlineto{\pgfqpoint{1.803223in}{0.606687in}}%
\pgfpathlineto{\pgfqpoint{1.793823in}{0.610228in}}%
\pgfpathlineto{\pgfqpoint{1.786479in}{0.611427in}}%
\pgfpathlineto{\pgfqpoint{1.781271in}{0.614134in}}%
\pgfpathlineto{\pgfqpoint{1.775795in}{0.607167in}}%
\pgfpathlineto{\pgfqpoint{1.760890in}{0.592713in}}%
\pgfpathclose%
\pgfusepath{fill}%
\end{pgfscope}%
\begin{pgfscope}%
\pgfpathrectangle{\pgfqpoint{0.100000in}{0.100000in}}{\pgfqpoint{3.007045in}{1.925000in}}%
\pgfusepath{clip}%
\pgfsetbuttcap%
\pgfsetmiterjoin%
\definecolor{currentfill}{rgb}{0.681738,0.818562,0.904098}%
\pgfsetfillcolor{currentfill}%
\pgfsetlinewidth{0.000000pt}%
\definecolor{currentstroke}{rgb}{0.000000,0.000000,0.000000}%
\pgfsetstrokecolor{currentstroke}%
\pgfsetstrokeopacity{0.000000}%
\pgfsetdash{}{0pt}%
\pgfpathmoveto{\pgfqpoint{2.437163in}{0.979641in}}%
\pgfpathlineto{\pgfqpoint{2.431676in}{0.978202in}}%
\pgfpathlineto{\pgfqpoint{2.429667in}{0.984471in}}%
\pgfpathlineto{\pgfqpoint{2.416656in}{0.974681in}}%
\pgfpathlineto{\pgfqpoint{2.406942in}{0.983024in}}%
\pgfpathlineto{\pgfqpoint{2.400322in}{0.986669in}}%
\pgfpathlineto{\pgfqpoint{2.404364in}{0.992821in}}%
\pgfpathlineto{\pgfqpoint{2.399497in}{0.998207in}}%
\pgfpathlineto{\pgfqpoint{2.394844in}{0.994687in}}%
\pgfpathlineto{\pgfqpoint{2.392615in}{1.001420in}}%
\pgfpathlineto{\pgfqpoint{2.396925in}{1.003949in}}%
\pgfpathlineto{\pgfqpoint{2.399781in}{1.010130in}}%
\pgfpathlineto{\pgfqpoint{2.406434in}{1.016470in}}%
\pgfpathlineto{\pgfqpoint{2.407690in}{1.022099in}}%
\pgfpathlineto{\pgfqpoint{2.413702in}{1.026726in}}%
\pgfpathlineto{\pgfqpoint{2.416735in}{1.033753in}}%
\pgfpathlineto{\pgfqpoint{2.422280in}{1.037551in}}%
\pgfpathlineto{\pgfqpoint{2.434200in}{1.039055in}}%
\pgfpathlineto{\pgfqpoint{2.439622in}{1.028698in}}%
\pgfpathlineto{\pgfqpoint{2.450725in}{1.036221in}}%
\pgfpathlineto{\pgfqpoint{2.451559in}{1.040086in}}%
\pgfpathlineto{\pgfqpoint{2.461325in}{1.043806in}}%
\pgfpathlineto{\pgfqpoint{2.464192in}{1.048074in}}%
\pgfpathlineto{\pgfqpoint{2.469632in}{1.050830in}}%
\pgfpathlineto{\pgfqpoint{2.483475in}{1.055101in}}%
\pgfpathlineto{\pgfqpoint{2.490316in}{1.048764in}}%
\pgfpathlineto{\pgfqpoint{2.495062in}{1.040390in}}%
\pgfpathlineto{\pgfqpoint{2.481841in}{1.034709in}}%
\pgfpathlineto{\pgfqpoint{2.481176in}{1.031217in}}%
\pgfpathlineto{\pgfqpoint{2.476402in}{1.027509in}}%
\pgfpathlineto{\pgfqpoint{2.466964in}{1.026275in}}%
\pgfpathlineto{\pgfqpoint{2.462097in}{1.022483in}}%
\pgfpathlineto{\pgfqpoint{2.461449in}{1.016280in}}%
\pgfpathlineto{\pgfqpoint{2.458161in}{1.009811in}}%
\pgfpathlineto{\pgfqpoint{2.463756in}{1.002325in}}%
\pgfpathlineto{\pgfqpoint{2.460956in}{0.992711in}}%
\pgfpathlineto{\pgfqpoint{2.458453in}{0.990986in}}%
\pgfpathlineto{\pgfqpoint{2.450485in}{0.993442in}}%
\pgfpathlineto{\pgfqpoint{2.441498in}{0.987564in}}%
\pgfpathlineto{\pgfqpoint{2.437163in}{0.979641in}}%
\pgfpathclose%
\pgfusepath{fill}%
\end{pgfscope}%
\begin{pgfscope}%
\pgfpathrectangle{\pgfqpoint{0.100000in}{0.100000in}}{\pgfqpoint{3.007045in}{1.925000in}}%
\pgfusepath{clip}%
\pgfsetbuttcap%
\pgfsetmiterjoin%
\definecolor{currentfill}{rgb}{0.676817,0.816471,0.902376}%
\pgfsetfillcolor{currentfill}%
\pgfsetlinewidth{0.000000pt}%
\definecolor{currentstroke}{rgb}{0.000000,0.000000,0.000000}%
\pgfsetstrokecolor{currentstroke}%
\pgfsetstrokeopacity{0.000000}%
\pgfsetdash{}{0pt}%
\pgfpathmoveto{\pgfqpoint{2.131069in}{1.580709in}}%
\pgfpathlineto{\pgfqpoint{2.124022in}{1.562475in}}%
\pgfpathlineto{\pgfqpoint{2.118308in}{1.554107in}}%
\pgfpathlineto{\pgfqpoint{2.116985in}{1.541346in}}%
\pgfpathlineto{\pgfqpoint{2.111077in}{1.540159in}}%
\pgfpathlineto{\pgfqpoint{2.094036in}{1.542696in}}%
\pgfpathlineto{\pgfqpoint{2.093706in}{1.548428in}}%
\pgfpathlineto{\pgfqpoint{2.090120in}{1.553996in}}%
\pgfpathlineto{\pgfqpoint{2.084400in}{1.553797in}}%
\pgfpathlineto{\pgfqpoint{2.083707in}{1.565240in}}%
\pgfpathlineto{\pgfqpoint{2.078017in}{1.564986in}}%
\pgfpathlineto{\pgfqpoint{2.076489in}{1.587837in}}%
\pgfpathlineto{\pgfqpoint{2.064911in}{1.587042in}}%
\pgfpathlineto{\pgfqpoint{2.063792in}{1.604125in}}%
\pgfpathlineto{\pgfqpoint{2.048726in}{1.610846in}}%
\pgfpathlineto{\pgfqpoint{2.047301in}{1.632207in}}%
\pgfpathlineto{\pgfqpoint{2.087227in}{1.635040in}}%
\pgfpathlineto{\pgfqpoint{2.088086in}{1.623582in}}%
\pgfpathlineto{\pgfqpoint{2.111036in}{1.625398in}}%
\pgfpathlineto{\pgfqpoint{2.112389in}{1.608131in}}%
\pgfpathlineto{\pgfqpoint{2.123850in}{1.609038in}}%
\pgfpathlineto{\pgfqpoint{2.126263in}{1.603472in}}%
\pgfpathlineto{\pgfqpoint{2.128131in}{1.580457in}}%
\pgfpathlineto{\pgfqpoint{2.131069in}{1.580709in}}%
\pgfpathclose%
\pgfusepath{fill}%
\end{pgfscope}%
\begin{pgfscope}%
\pgfpathrectangle{\pgfqpoint{0.100000in}{0.100000in}}{\pgfqpoint{3.007045in}{1.925000in}}%
\pgfusepath{clip}%
\pgfsetbuttcap%
\pgfsetmiterjoin%
\definecolor{currentfill}{rgb}{0.396909,0.666851,0.830358}%
\pgfsetfillcolor{currentfill}%
\pgfsetlinewidth{0.000000pt}%
\definecolor{currentstroke}{rgb}{0.000000,0.000000,0.000000}%
\pgfsetstrokecolor{currentstroke}%
\pgfsetstrokeopacity{0.000000}%
\pgfsetdash{}{0pt}%
\pgfpathmoveto{\pgfqpoint{2.106178in}{0.880380in}}%
\pgfpathlineto{\pgfqpoint{2.085939in}{0.878995in}}%
\pgfpathlineto{\pgfqpoint{2.085014in}{0.905571in}}%
\pgfpathlineto{\pgfqpoint{2.090507in}{0.908195in}}%
\pgfpathlineto{\pgfqpoint{2.089884in}{0.925436in}}%
\pgfpathlineto{\pgfqpoint{2.074636in}{0.931335in}}%
\pgfpathlineto{\pgfqpoint{2.083454in}{0.947053in}}%
\pgfpathlineto{\pgfqpoint{2.083172in}{0.958841in}}%
\pgfpathlineto{\pgfqpoint{2.093449in}{0.959312in}}%
\pgfpathlineto{\pgfqpoint{2.098903in}{0.955455in}}%
\pgfpathlineto{\pgfqpoint{2.108338in}{0.951144in}}%
\pgfpathlineto{\pgfqpoint{2.108828in}{0.933214in}}%
\pgfpathlineto{\pgfqpoint{2.114171in}{0.933474in}}%
\pgfpathlineto{\pgfqpoint{2.114744in}{0.920210in}}%
\pgfpathlineto{\pgfqpoint{2.120897in}{0.915606in}}%
\pgfpathlineto{\pgfqpoint{2.125902in}{0.914003in}}%
\pgfpathlineto{\pgfqpoint{2.129095in}{0.910042in}}%
\pgfpathlineto{\pgfqpoint{2.129176in}{0.907553in}}%
\pgfpathlineto{\pgfqpoint{2.117945in}{0.906798in}}%
\pgfpathlineto{\pgfqpoint{2.114334in}{0.905115in}}%
\pgfpathlineto{\pgfqpoint{2.110613in}{0.897382in}}%
\pgfpathlineto{\pgfqpoint{2.107278in}{0.897155in}}%
\pgfpathlineto{\pgfqpoint{2.108156in}{0.880516in}}%
\pgfpathlineto{\pgfqpoint{2.106178in}{0.880380in}}%
\pgfpathclose%
\pgfusepath{fill}%
\end{pgfscope}%
\begin{pgfscope}%
\pgfpathrectangle{\pgfqpoint{0.100000in}{0.100000in}}{\pgfqpoint{3.007045in}{1.925000in}}%
\pgfusepath{clip}%
\pgfsetbuttcap%
\pgfsetmiterjoin%
\definecolor{currentfill}{rgb}{0.371688,0.649627,0.820515}%
\pgfsetfillcolor{currentfill}%
\pgfsetlinewidth{0.000000pt}%
\definecolor{currentstroke}{rgb}{0.000000,0.000000,0.000000}%
\pgfsetstrokecolor{currentstroke}%
\pgfsetstrokeopacity{0.000000}%
\pgfsetdash{}{0pt}%
\pgfpathmoveto{\pgfqpoint{1.650255in}{1.108902in}}%
\pgfpathlineto{\pgfqpoint{1.621825in}{1.109765in}}%
\pgfpathlineto{\pgfqpoint{1.622153in}{1.126974in}}%
\pgfpathlineto{\pgfqpoint{1.593668in}{1.127954in}}%
\pgfpathlineto{\pgfqpoint{1.594074in}{1.145204in}}%
\pgfpathlineto{\pgfqpoint{1.594289in}{1.150925in}}%
\pgfpathlineto{\pgfqpoint{1.622833in}{1.149886in}}%
\pgfpathlineto{\pgfqpoint{1.623055in}{1.155630in}}%
\pgfpathlineto{\pgfqpoint{1.651571in}{1.154782in}}%
\pgfpathlineto{\pgfqpoint{1.650255in}{1.108902in}}%
\pgfpathclose%
\pgfusepath{fill}%
\end{pgfscope}%
\begin{pgfscope}%
\pgfpathrectangle{\pgfqpoint{0.100000in}{0.100000in}}{\pgfqpoint{3.007045in}{1.925000in}}%
\pgfusepath{clip}%
\pgfsetbuttcap%
\pgfsetmiterjoin%
\definecolor{currentfill}{rgb}{0.154787,0.468512,0.722876}%
\pgfsetfillcolor{currentfill}%
\pgfsetlinewidth{0.000000pt}%
\definecolor{currentstroke}{rgb}{0.000000,0.000000,0.000000}%
\pgfsetstrokecolor{currentstroke}%
\pgfsetstrokeopacity{0.000000}%
\pgfsetdash{}{0pt}%
\pgfpathmoveto{\pgfqpoint{2.140155in}{0.982191in}}%
\pgfpathlineto{\pgfqpoint{2.141030in}{0.985137in}}%
\pgfpathlineto{\pgfqpoint{2.135867in}{0.998378in}}%
\pgfpathlineto{\pgfqpoint{2.132817in}{1.006098in}}%
\pgfpathlineto{\pgfqpoint{2.146362in}{1.013286in}}%
\pgfpathlineto{\pgfqpoint{2.152882in}{1.014762in}}%
\pgfpathlineto{\pgfqpoint{2.155476in}{1.018069in}}%
\pgfpathlineto{\pgfqpoint{2.156379in}{1.026602in}}%
\pgfpathlineto{\pgfqpoint{2.164956in}{1.024377in}}%
\pgfpathlineto{\pgfqpoint{2.174509in}{1.028592in}}%
\pgfpathlineto{\pgfqpoint{2.179536in}{1.021312in}}%
\pgfpathlineto{\pgfqpoint{2.189732in}{1.023554in}}%
\pgfpathlineto{\pgfqpoint{2.191739in}{0.996087in}}%
\pgfpathlineto{\pgfqpoint{2.188886in}{0.995787in}}%
\pgfpathlineto{\pgfqpoint{2.189194in}{0.983871in}}%
\pgfpathlineto{\pgfqpoint{2.183256in}{0.977798in}}%
\pgfpathlineto{\pgfqpoint{2.181638in}{0.973954in}}%
\pgfpathlineto{\pgfqpoint{2.169549in}{0.973787in}}%
\pgfpathlineto{\pgfqpoint{2.169480in}{0.968102in}}%
\pgfpathlineto{\pgfqpoint{2.166450in}{0.964466in}}%
\pgfpathlineto{\pgfqpoint{2.158552in}{0.966327in}}%
\pgfpathlineto{\pgfqpoint{2.146928in}{0.965917in}}%
\pgfpathlineto{\pgfqpoint{2.146159in}{0.970611in}}%
\pgfpathlineto{\pgfqpoint{2.140247in}{0.978276in}}%
\pgfpathlineto{\pgfqpoint{2.140155in}{0.982191in}}%
\pgfpathclose%
\pgfusepath{fill}%
\end{pgfscope}%
\begin{pgfscope}%
\pgfpathrectangle{\pgfqpoint{0.100000in}{0.100000in}}{\pgfqpoint{3.007045in}{1.925000in}}%
\pgfusepath{clip}%
\pgfsetbuttcap%
\pgfsetmiterjoin%
\definecolor{currentfill}{rgb}{0.466667,0.708189,0.849366}%
\pgfsetfillcolor{currentfill}%
\pgfsetlinewidth{0.000000pt}%
\definecolor{currentstroke}{rgb}{0.000000,0.000000,0.000000}%
\pgfsetstrokecolor{currentstroke}%
\pgfsetstrokeopacity{0.000000}%
\pgfsetdash{}{0pt}%
\pgfpathmoveto{\pgfqpoint{2.721216in}{1.279405in}}%
\pgfpathlineto{\pgfqpoint{2.739138in}{1.283230in}}%
\pgfpathlineto{\pgfqpoint{2.742511in}{1.290165in}}%
\pgfpathlineto{\pgfqpoint{2.746343in}{1.292558in}}%
\pgfpathlineto{\pgfqpoint{2.753598in}{1.293002in}}%
\pgfpathlineto{\pgfqpoint{2.753954in}{1.289095in}}%
\pgfpathlineto{\pgfqpoint{2.751124in}{1.279964in}}%
\pgfpathlineto{\pgfqpoint{2.754264in}{1.276110in}}%
\pgfpathlineto{\pgfqpoint{2.754436in}{1.271463in}}%
\pgfpathlineto{\pgfqpoint{2.751790in}{1.268575in}}%
\pgfpathlineto{\pgfqpoint{2.757050in}{1.263166in}}%
\pgfpathlineto{\pgfqpoint{2.751195in}{1.256851in}}%
\pgfpathlineto{\pgfqpoint{2.745705in}{1.255995in}}%
\pgfpathlineto{\pgfqpoint{2.744257in}{1.261114in}}%
\pgfpathlineto{\pgfqpoint{2.738533in}{1.258792in}}%
\pgfpathlineto{\pgfqpoint{2.730272in}{1.258328in}}%
\pgfpathlineto{\pgfqpoint{2.732034in}{1.262134in}}%
\pgfpathlineto{\pgfqpoint{2.731538in}{1.268601in}}%
\pgfpathlineto{\pgfqpoint{2.726638in}{1.268556in}}%
\pgfpathlineto{\pgfqpoint{2.716371in}{1.278367in}}%
\pgfpathlineto{\pgfqpoint{2.721216in}{1.279405in}}%
\pgfpathclose%
\pgfusepath{fill}%
\end{pgfscope}%
\begin{pgfscope}%
\pgfpathrectangle{\pgfqpoint{0.100000in}{0.100000in}}{\pgfqpoint{3.007045in}{1.925000in}}%
\pgfusepath{clip}%
\pgfsetbuttcap%
\pgfsetmiterjoin%
\definecolor{currentfill}{rgb}{0.351511,0.635848,0.812641}%
\pgfsetfillcolor{currentfill}%
\pgfsetlinewidth{0.000000pt}%
\definecolor{currentstroke}{rgb}{0.000000,0.000000,0.000000}%
\pgfsetstrokecolor{currentstroke}%
\pgfsetstrokeopacity{0.000000}%
\pgfsetdash{}{0pt}%
\pgfpathmoveto{\pgfqpoint{1.509458in}{1.149192in}}%
\pgfpathlineto{\pgfqpoint{1.508618in}{1.149222in}}%
\pgfpathlineto{\pgfqpoint{1.510385in}{1.177845in}}%
\pgfpathlineto{\pgfqpoint{1.537445in}{1.176396in}}%
\pgfpathlineto{\pgfqpoint{1.567058in}{1.175010in}}%
\pgfpathlineto{\pgfqpoint{1.565600in}{1.146405in}}%
\pgfpathlineto{\pgfqpoint{1.509458in}{1.149192in}}%
\pgfpathclose%
\pgfusepath{fill}%
\end{pgfscope}%
\begin{pgfscope}%
\pgfpathrectangle{\pgfqpoint{0.100000in}{0.100000in}}{\pgfqpoint{3.007045in}{1.925000in}}%
\pgfusepath{clip}%
\pgfsetbuttcap%
\pgfsetmiterjoin%
\definecolor{currentfill}{rgb}{0.244106,0.557832,0.768889}%
\pgfsetfillcolor{currentfill}%
\pgfsetlinewidth{0.000000pt}%
\definecolor{currentstroke}{rgb}{0.000000,0.000000,0.000000}%
\pgfsetstrokecolor{currentstroke}%
\pgfsetstrokeopacity{0.000000}%
\pgfsetdash{}{0pt}%
\pgfpathmoveto{\pgfqpoint{1.596180in}{1.004385in}}%
\pgfpathlineto{\pgfqpoint{1.596980in}{1.029821in}}%
\pgfpathlineto{\pgfqpoint{1.590957in}{1.030045in}}%
\pgfpathlineto{\pgfqpoint{1.591838in}{1.053004in}}%
\pgfpathlineto{\pgfqpoint{1.631737in}{1.051813in}}%
\pgfpathlineto{\pgfqpoint{1.631910in}{1.063521in}}%
\pgfpathlineto{\pgfqpoint{1.660467in}{1.062737in}}%
\pgfpathlineto{\pgfqpoint{1.660705in}{1.074262in}}%
\pgfpathlineto{\pgfqpoint{1.676880in}{1.073769in}}%
\pgfpathlineto{\pgfqpoint{1.693377in}{1.073520in}}%
\pgfpathlineto{\pgfqpoint{1.692552in}{1.033341in}}%
\pgfpathlineto{\pgfqpoint{1.659741in}{1.033948in}}%
\pgfpathlineto{\pgfqpoint{1.659290in}{1.002553in}}%
\pgfpathlineto{\pgfqpoint{1.642737in}{1.002945in}}%
\pgfpathlineto{\pgfqpoint{1.596180in}{1.004385in}}%
\pgfpathclose%
\pgfusepath{fill}%
\end{pgfscope}%
\begin{pgfscope}%
\pgfpathrectangle{\pgfqpoint{0.100000in}{0.100000in}}{\pgfqpoint{3.007045in}{1.925000in}}%
\pgfusepath{clip}%
\pgfsetbuttcap%
\pgfsetmiterjoin%
\definecolor{currentfill}{rgb}{0.306113,0.604844,0.794925}%
\pgfsetfillcolor{currentfill}%
\pgfsetlinewidth{0.000000pt}%
\definecolor{currentstroke}{rgb}{0.000000,0.000000,0.000000}%
\pgfsetstrokecolor{currentstroke}%
\pgfsetstrokeopacity{0.000000}%
\pgfsetdash{}{0pt}%
\pgfpathmoveto{\pgfqpoint{1.526290in}{0.791703in}}%
\pgfpathlineto{\pgfqpoint{1.525327in}{0.794045in}}%
\pgfpathlineto{\pgfqpoint{1.514372in}{0.798412in}}%
\pgfpathlineto{\pgfqpoint{1.495136in}{0.800124in}}%
\pgfpathlineto{\pgfqpoint{1.496660in}{0.826118in}}%
\pgfpathlineto{\pgfqpoint{1.499337in}{0.825611in}}%
\pgfpathlineto{\pgfqpoint{1.500565in}{0.847796in}}%
\pgfpathlineto{\pgfqpoint{1.504435in}{0.848612in}}%
\pgfpathlineto{\pgfqpoint{1.516379in}{0.834892in}}%
\pgfpathlineto{\pgfqpoint{1.521539in}{0.834387in}}%
\pgfpathlineto{\pgfqpoint{1.526280in}{0.836772in}}%
\pgfpathlineto{\pgfqpoint{1.531947in}{0.833744in}}%
\pgfpathlineto{\pgfqpoint{1.533732in}{0.839225in}}%
\pgfpathlineto{\pgfqpoint{1.542619in}{0.830859in}}%
\pgfpathlineto{\pgfqpoint{1.543414in}{0.822716in}}%
\pgfpathlineto{\pgfqpoint{1.556351in}{0.822025in}}%
\pgfpathlineto{\pgfqpoint{1.555230in}{0.797022in}}%
\pgfpathlineto{\pgfqpoint{1.526556in}{0.798349in}}%
\pgfpathlineto{\pgfqpoint{1.526290in}{0.791703in}}%
\pgfpathclose%
\pgfusepath{fill}%
\end{pgfscope}%
\begin{pgfscope}%
\pgfpathrectangle{\pgfqpoint{0.100000in}{0.100000in}}{\pgfqpoint{3.007045in}{1.925000in}}%
\pgfusepath{clip}%
\pgfsetbuttcap%
\pgfsetmiterjoin%
\definecolor{currentfill}{rgb}{0.285936,0.591065,0.787051}%
\pgfsetfillcolor{currentfill}%
\pgfsetlinewidth{0.000000pt}%
\definecolor{currentstroke}{rgb}{0.000000,0.000000,0.000000}%
\pgfsetstrokecolor{currentstroke}%
\pgfsetstrokeopacity{0.000000}%
\pgfsetdash{}{0pt}%
\pgfpathmoveto{\pgfqpoint{0.819433in}{0.344425in}}%
\pgfpathlineto{\pgfqpoint{0.816099in}{0.346069in}}%
\pgfpathlineto{\pgfqpoint{0.814545in}{0.348123in}}%
\pgfpathlineto{\pgfqpoint{0.812387in}{0.348661in}}%
\pgfpathlineto{\pgfqpoint{0.812849in}{0.351051in}}%
\pgfpathlineto{\pgfqpoint{0.812080in}{0.351861in}}%
\pgfpathlineto{\pgfqpoint{0.810038in}{0.350072in}}%
\pgfpathlineto{\pgfqpoint{0.808207in}{0.349391in}}%
\pgfpathlineto{\pgfqpoint{0.806954in}{0.350824in}}%
\pgfpathlineto{\pgfqpoint{0.805473in}{0.350277in}}%
\pgfpathlineto{\pgfqpoint{0.804388in}{0.347206in}}%
\pgfpathlineto{\pgfqpoint{0.802543in}{0.351436in}}%
\pgfpathlineto{\pgfqpoint{0.802043in}{0.348092in}}%
\pgfpathlineto{\pgfqpoint{0.798587in}{0.347926in}}%
\pgfpathlineto{\pgfqpoint{0.798087in}{0.346708in}}%
\pgfpathlineto{\pgfqpoint{0.795113in}{0.346802in}}%
\pgfpathlineto{\pgfqpoint{0.791694in}{0.347913in}}%
\pgfpathlineto{\pgfqpoint{0.786584in}{0.347123in}}%
\pgfpathlineto{\pgfqpoint{0.784356in}{0.348068in}}%
\pgfpathlineto{\pgfqpoint{0.783495in}{0.346619in}}%
\pgfpathlineto{\pgfqpoint{0.781504in}{0.347865in}}%
\pgfpathlineto{\pgfqpoint{0.781399in}{0.349966in}}%
\pgfpathlineto{\pgfqpoint{0.779304in}{0.349025in}}%
\pgfpathlineto{\pgfqpoint{0.776366in}{0.349945in}}%
\pgfpathlineto{\pgfqpoint{0.776139in}{0.354814in}}%
\pgfpathlineto{\pgfqpoint{0.778619in}{0.356270in}}%
\pgfpathlineto{\pgfqpoint{0.781930in}{0.355808in}}%
\pgfpathlineto{\pgfqpoint{0.783469in}{0.354208in}}%
\pgfpathlineto{\pgfqpoint{0.785254in}{0.354993in}}%
\pgfpathlineto{\pgfqpoint{0.787547in}{0.354040in}}%
\pgfpathlineto{\pgfqpoint{0.788336in}{0.355429in}}%
\pgfpathlineto{\pgfqpoint{0.793181in}{0.356662in}}%
\pgfpathlineto{\pgfqpoint{0.790991in}{0.357629in}}%
\pgfpathlineto{\pgfqpoint{0.788388in}{0.357382in}}%
\pgfpathlineto{\pgfqpoint{0.785835in}{0.356265in}}%
\pgfpathlineto{\pgfqpoint{0.783898in}{0.359777in}}%
\pgfpathlineto{\pgfqpoint{0.783523in}{0.362343in}}%
\pgfpathlineto{\pgfqpoint{0.789148in}{0.366853in}}%
\pgfpathlineto{\pgfqpoint{0.793181in}{0.368023in}}%
\pgfpathlineto{\pgfqpoint{0.797910in}{0.370814in}}%
\pgfpathlineto{\pgfqpoint{0.800516in}{0.373276in}}%
\pgfpathlineto{\pgfqpoint{0.801383in}{0.377874in}}%
\pgfpathlineto{\pgfqpoint{0.803675in}{0.378105in}}%
\pgfpathlineto{\pgfqpoint{0.805964in}{0.377296in}}%
\pgfpathlineto{\pgfqpoint{0.811186in}{0.377817in}}%
\pgfpathlineto{\pgfqpoint{0.816257in}{0.377436in}}%
\pgfpathlineto{\pgfqpoint{0.816031in}{0.374807in}}%
\pgfpathlineto{\pgfqpoint{0.817543in}{0.372687in}}%
\pgfpathlineto{\pgfqpoint{0.822055in}{0.372655in}}%
\pgfpathlineto{\pgfqpoint{0.820708in}{0.377871in}}%
\pgfpathlineto{\pgfqpoint{0.822622in}{0.379076in}}%
\pgfpathlineto{\pgfqpoint{0.818515in}{0.382126in}}%
\pgfpathlineto{\pgfqpoint{0.817100in}{0.384096in}}%
\pgfpathlineto{\pgfqpoint{0.812559in}{0.384716in}}%
\pgfpathlineto{\pgfqpoint{0.809430in}{0.383487in}}%
\pgfpathlineto{\pgfqpoint{0.804683in}{0.384619in}}%
\pgfpathlineto{\pgfqpoint{0.800098in}{0.383859in}}%
\pgfpathlineto{\pgfqpoint{0.798506in}{0.380328in}}%
\pgfpathlineto{\pgfqpoint{0.797580in}{0.381676in}}%
\pgfpathlineto{\pgfqpoint{0.795111in}{0.381646in}}%
\pgfpathlineto{\pgfqpoint{0.790374in}{0.380576in}}%
\pgfpathlineto{\pgfqpoint{0.787763in}{0.377740in}}%
\pgfpathlineto{\pgfqpoint{0.785721in}{0.376868in}}%
\pgfpathlineto{\pgfqpoint{0.783285in}{0.376730in}}%
\pgfpathlineto{\pgfqpoint{0.781683in}{0.377953in}}%
\pgfpathlineto{\pgfqpoint{0.781141in}{0.373732in}}%
\pgfpathlineto{\pgfqpoint{0.777419in}{0.371708in}}%
\pgfpathlineto{\pgfqpoint{0.775658in}{0.371988in}}%
\pgfpathlineto{\pgfqpoint{0.774292in}{0.373277in}}%
\pgfpathlineto{\pgfqpoint{0.771887in}{0.370547in}}%
\pgfpathlineto{\pgfqpoint{0.769799in}{0.370122in}}%
\pgfpathlineto{\pgfqpoint{0.764543in}{0.372414in}}%
\pgfpathlineto{\pgfqpoint{0.763348in}{0.371469in}}%
\pgfpathlineto{\pgfqpoint{0.761565in}{0.371604in}}%
\pgfpathlineto{\pgfqpoint{0.760574in}{0.369853in}}%
\pgfpathlineto{\pgfqpoint{0.756927in}{0.371271in}}%
\pgfpathlineto{\pgfqpoint{0.753478in}{0.368527in}}%
\pgfpathlineto{\pgfqpoint{0.751025in}{0.367939in}}%
\pgfpathlineto{\pgfqpoint{0.752269in}{0.365262in}}%
\pgfpathlineto{\pgfqpoint{0.755024in}{0.362772in}}%
\pgfpathlineto{\pgfqpoint{0.757010in}{0.357656in}}%
\pgfpathlineto{\pgfqpoint{0.757017in}{0.355194in}}%
\pgfpathlineto{\pgfqpoint{0.755001in}{0.356698in}}%
\pgfpathlineto{\pgfqpoint{0.753292in}{0.354357in}}%
\pgfpathlineto{\pgfqpoint{0.749665in}{0.357072in}}%
\pgfpathlineto{\pgfqpoint{0.748423in}{0.355398in}}%
\pgfpathlineto{\pgfqpoint{0.743549in}{0.359090in}}%
\pgfpathlineto{\pgfqpoint{0.740274in}{0.361603in}}%
\pgfpathlineto{\pgfqpoint{0.741011in}{0.363689in}}%
\pgfpathlineto{\pgfqpoint{0.746018in}{0.370169in}}%
\pgfpathlineto{\pgfqpoint{0.745512in}{0.370591in}}%
\pgfpathlineto{\pgfqpoint{0.748076in}{0.373959in}}%
\pgfpathlineto{\pgfqpoint{0.751397in}{0.371451in}}%
\pgfpathlineto{\pgfqpoint{0.753894in}{0.374816in}}%
\pgfpathlineto{\pgfqpoint{0.756787in}{0.372685in}}%
\pgfpathlineto{\pgfqpoint{0.758045in}{0.374400in}}%
\pgfpathlineto{\pgfqpoint{0.759707in}{0.373154in}}%
\pgfpathlineto{\pgfqpoint{0.762258in}{0.376547in}}%
\pgfpathlineto{\pgfqpoint{0.763938in}{0.375291in}}%
\pgfpathlineto{\pgfqpoint{0.765233in}{0.376990in}}%
\pgfpathlineto{\pgfqpoint{0.766473in}{0.376050in}}%
\pgfpathlineto{\pgfqpoint{0.771354in}{0.382565in}}%
\pgfpathlineto{\pgfqpoint{0.772579in}{0.381674in}}%
\pgfpathlineto{\pgfqpoint{0.777700in}{0.388433in}}%
\pgfpathlineto{\pgfqpoint{0.778986in}{0.387487in}}%
\pgfpathlineto{\pgfqpoint{0.784096in}{0.394320in}}%
\pgfpathlineto{\pgfqpoint{0.784950in}{0.396276in}}%
\pgfpathlineto{\pgfqpoint{0.788663in}{0.401187in}}%
\pgfpathlineto{\pgfqpoint{0.789615in}{0.403627in}}%
\pgfpathlineto{\pgfqpoint{0.792854in}{0.407923in}}%
\pgfpathlineto{\pgfqpoint{0.796295in}{0.405189in}}%
\pgfpathlineto{\pgfqpoint{0.796605in}{0.403352in}}%
\pgfpathlineto{\pgfqpoint{0.809857in}{0.421160in}}%
\pgfpathlineto{\pgfqpoint{0.816153in}{0.429462in}}%
\pgfpathlineto{\pgfqpoint{0.826027in}{0.421875in}}%
\pgfpathlineto{\pgfqpoint{0.827109in}{0.423257in}}%
\pgfpathlineto{\pgfqpoint{0.847814in}{0.419834in}}%
\pgfpathlineto{\pgfqpoint{0.856121in}{0.418510in}}%
\pgfpathlineto{\pgfqpoint{0.869886in}{0.409139in}}%
\pgfpathlineto{\pgfqpoint{0.871795in}{0.412013in}}%
\pgfpathlineto{\pgfqpoint{0.875033in}{0.409830in}}%
\pgfpathlineto{\pgfqpoint{0.886359in}{0.402751in}}%
\pgfpathlineto{\pgfqpoint{0.880633in}{0.393554in}}%
\pgfpathlineto{\pgfqpoint{0.878779in}{0.389457in}}%
\pgfpathlineto{\pgfqpoint{0.870796in}{0.376637in}}%
\pgfpathlineto{\pgfqpoint{0.865448in}{0.380135in}}%
\pgfpathlineto{\pgfqpoint{0.864498in}{0.378218in}}%
\pgfpathlineto{\pgfqpoint{0.855608in}{0.363995in}}%
\pgfpathlineto{\pgfqpoint{0.853061in}{0.365671in}}%
\pgfpathlineto{\pgfqpoint{0.852394in}{0.364648in}}%
\pgfpathlineto{\pgfqpoint{0.840140in}{0.372751in}}%
\pgfpathlineto{\pgfqpoint{0.836391in}{0.367548in}}%
\pgfpathlineto{\pgfqpoint{0.832191in}{0.361126in}}%
\pgfpathlineto{\pgfqpoint{0.829442in}{0.362834in}}%
\pgfpathlineto{\pgfqpoint{0.827922in}{0.360561in}}%
\pgfpathlineto{\pgfqpoint{0.828821in}{0.359977in}}%
\pgfpathlineto{\pgfqpoint{0.824494in}{0.353761in}}%
\pgfpathlineto{\pgfqpoint{0.825526in}{0.353057in}}%
\pgfpathlineto{\pgfqpoint{0.821535in}{0.347132in}}%
\pgfpathlineto{\pgfqpoint{0.819433in}{0.344425in}}%
\pgfpathclose%
\pgfusepath{fill}%
\end{pgfscope}%
\begin{pgfscope}%
\pgfpathrectangle{\pgfqpoint{0.100000in}{0.100000in}}{\pgfqpoint{3.007045in}{1.925000in}}%
\pgfusepath{clip}%
\pgfsetbuttcap%
\pgfsetmiterjoin%
\definecolor{currentfill}{rgb}{0.642368,0.801830,0.890319}%
\pgfsetfillcolor{currentfill}%
\pgfsetlinewidth{0.000000pt}%
\definecolor{currentstroke}{rgb}{0.000000,0.000000,0.000000}%
\pgfsetstrokecolor{currentstroke}%
\pgfsetstrokeopacity{0.000000}%
\pgfsetdash{}{0pt}%
\pgfpathmoveto{\pgfqpoint{2.434200in}{1.039055in}}%
\pgfpathlineto{\pgfqpoint{2.439575in}{1.041742in}}%
\pgfpathlineto{\pgfqpoint{2.437440in}{1.045801in}}%
\pgfpathlineto{\pgfqpoint{2.446669in}{1.052394in}}%
\pgfpathlineto{\pgfqpoint{2.440019in}{1.062167in}}%
\pgfpathlineto{\pgfqpoint{2.435154in}{1.067105in}}%
\pgfpathlineto{\pgfqpoint{2.450852in}{1.085445in}}%
\pgfpathlineto{\pgfqpoint{2.453185in}{1.084073in}}%
\pgfpathlineto{\pgfqpoint{2.452938in}{1.079250in}}%
\pgfpathlineto{\pgfqpoint{2.455254in}{1.073522in}}%
\pgfpathlineto{\pgfqpoint{2.462901in}{1.070185in}}%
\pgfpathlineto{\pgfqpoint{2.469086in}{1.065710in}}%
\pgfpathlineto{\pgfqpoint{2.475445in}{1.067081in}}%
\pgfpathlineto{\pgfqpoint{2.483702in}{1.075698in}}%
\pgfpathlineto{\pgfqpoint{2.487921in}{1.088350in}}%
\pgfpathlineto{\pgfqpoint{2.489686in}{1.089390in}}%
\pgfpathlineto{\pgfqpoint{2.490450in}{1.093853in}}%
\pgfpathlineto{\pgfqpoint{2.495503in}{1.095488in}}%
\pgfpathlineto{\pgfqpoint{2.509239in}{1.087027in}}%
\pgfpathlineto{\pgfqpoint{2.510605in}{1.081728in}}%
\pgfpathlineto{\pgfqpoint{2.502786in}{1.075737in}}%
\pgfpathlineto{\pgfqpoint{2.512374in}{1.068702in}}%
\pgfpathlineto{\pgfqpoint{2.508760in}{1.065976in}}%
\pgfpathlineto{\pgfqpoint{2.511594in}{1.060569in}}%
\pgfpathlineto{\pgfqpoint{2.521063in}{1.052157in}}%
\pgfpathlineto{\pgfqpoint{2.505363in}{1.044619in}}%
\pgfpathlineto{\pgfqpoint{2.503517in}{1.042080in}}%
\pgfpathlineto{\pgfqpoint{2.495062in}{1.040390in}}%
\pgfpathlineto{\pgfqpoint{2.490316in}{1.048764in}}%
\pgfpathlineto{\pgfqpoint{2.483475in}{1.055101in}}%
\pgfpathlineto{\pgfqpoint{2.469632in}{1.050830in}}%
\pgfpathlineto{\pgfqpoint{2.464192in}{1.048074in}}%
\pgfpathlineto{\pgfqpoint{2.461325in}{1.043806in}}%
\pgfpathlineto{\pgfqpoint{2.451559in}{1.040086in}}%
\pgfpathlineto{\pgfqpoint{2.450725in}{1.036221in}}%
\pgfpathlineto{\pgfqpoint{2.439622in}{1.028698in}}%
\pgfpathlineto{\pgfqpoint{2.434200in}{1.039055in}}%
\pgfpathclose%
\pgfusepath{fill}%
\end{pgfscope}%
\begin{pgfscope}%
\pgfpathrectangle{\pgfqpoint{0.100000in}{0.100000in}}{\pgfqpoint{3.007045in}{1.925000in}}%
\pgfusepath{clip}%
\pgfsetbuttcap%
\pgfsetmiterjoin%
\definecolor{currentfill}{rgb}{0.296025,0.597955,0.790988}%
\pgfsetfillcolor{currentfill}%
\pgfsetlinewidth{0.000000pt}%
\definecolor{currentstroke}{rgb}{0.000000,0.000000,0.000000}%
\pgfsetstrokecolor{currentstroke}%
\pgfsetstrokeopacity{0.000000}%
\pgfsetdash{}{0pt}%
\pgfpathmoveto{\pgfqpoint{1.570193in}{1.688446in}}%
\pgfpathlineto{\pgfqpoint{1.569042in}{1.665415in}}%
\pgfpathlineto{\pgfqpoint{1.570614in}{1.665351in}}%
\pgfpathlineto{\pgfqpoint{1.569331in}{1.642166in}}%
\pgfpathlineto{\pgfqpoint{1.540553in}{1.643789in}}%
\pgfpathlineto{\pgfqpoint{1.542129in}{1.666889in}}%
\pgfpathlineto{\pgfqpoint{1.540344in}{1.666994in}}%
\pgfpathlineto{\pgfqpoint{1.541684in}{1.689945in}}%
\pgfpathlineto{\pgfqpoint{1.570193in}{1.688446in}}%
\pgfpathclose%
\pgfusepath{fill}%
\end{pgfscope}%
\begin{pgfscope}%
\pgfpathrectangle{\pgfqpoint{0.100000in}{0.100000in}}{\pgfqpoint{3.007045in}{1.925000in}}%
\pgfusepath{clip}%
\pgfsetbuttcap%
\pgfsetmiterjoin%
\definecolor{currentfill}{rgb}{0.108651,0.416563,0.689043}%
\pgfsetfillcolor{currentfill}%
\pgfsetlinewidth{0.000000pt}%
\definecolor{currentstroke}{rgb}{0.000000,0.000000,0.000000}%
\pgfsetstrokecolor{currentstroke}%
\pgfsetstrokeopacity{0.000000}%
\pgfsetdash{}{0pt}%
\pgfpathmoveto{\pgfqpoint{1.456990in}{0.978623in}}%
\pgfpathlineto{\pgfqpoint{1.507272in}{0.975589in}}%
\pgfpathlineto{\pgfqpoint{1.505836in}{0.946347in}}%
\pgfpathlineto{\pgfqpoint{1.477103in}{0.948069in}}%
\pgfpathlineto{\pgfqpoint{1.419469in}{0.951860in}}%
\pgfpathlineto{\pgfqpoint{1.421577in}{0.981072in}}%
\pgfpathlineto{\pgfqpoint{1.456990in}{0.978623in}}%
\pgfpathclose%
\pgfusepath{fill}%
\end{pgfscope}%
\begin{pgfscope}%
\pgfpathrectangle{\pgfqpoint{0.100000in}{0.100000in}}{\pgfqpoint{3.007045in}{1.925000in}}%
\pgfusepath{clip}%
\pgfsetbuttcap%
\pgfsetmiterjoin%
\definecolor{currentfill}{rgb}{0.361599,0.642737,0.816578}%
\pgfsetfillcolor{currentfill}%
\pgfsetlinewidth{0.000000pt}%
\definecolor{currentstroke}{rgb}{0.000000,0.000000,0.000000}%
\pgfsetstrokecolor{currentstroke}%
\pgfsetstrokeopacity{0.000000}%
\pgfsetdash{}{0pt}%
\pgfpathmoveto{\pgfqpoint{2.196663in}{1.378800in}}%
\pgfpathlineto{\pgfqpoint{2.198102in}{1.398160in}}%
\pgfpathlineto{\pgfqpoint{2.196389in}{1.411349in}}%
\pgfpathlineto{\pgfqpoint{2.193000in}{1.423406in}}%
\pgfpathlineto{\pgfqpoint{2.183445in}{1.440623in}}%
\pgfpathlineto{\pgfqpoint{2.177072in}{1.454810in}}%
\pgfpathlineto{\pgfqpoint{2.176584in}{1.460340in}}%
\pgfpathlineto{\pgfqpoint{2.180375in}{1.470113in}}%
\pgfpathlineto{\pgfqpoint{2.199018in}{1.471645in}}%
\pgfpathlineto{\pgfqpoint{2.221626in}{1.473931in}}%
\pgfpathlineto{\pgfqpoint{2.224008in}{1.451122in}}%
\pgfpathlineto{\pgfqpoint{2.246778in}{1.453408in}}%
\pgfpathlineto{\pgfqpoint{2.258163in}{1.454659in}}%
\pgfpathlineto{\pgfqpoint{2.261163in}{1.431916in}}%
\pgfpathlineto{\pgfqpoint{2.263712in}{1.409069in}}%
\pgfpathlineto{\pgfqpoint{2.252299in}{1.407809in}}%
\pgfpathlineto{\pgfqpoint{2.229609in}{1.405210in}}%
\pgfpathlineto{\pgfqpoint{2.232094in}{1.382445in}}%
\pgfpathlineto{\pgfqpoint{2.196663in}{1.378800in}}%
\pgfpathclose%
\pgfusepath{fill}%
\end{pgfscope}%
\begin{pgfscope}%
\pgfpathrectangle{\pgfqpoint{0.100000in}{0.100000in}}{\pgfqpoint{3.007045in}{1.925000in}}%
\pgfusepath{clip}%
\pgfsetbuttcap%
\pgfsetmiterjoin%
\definecolor{currentfill}{rgb}{0.130427,0.444152,0.710327}%
\pgfsetfillcolor{currentfill}%
\pgfsetlinewidth{0.000000pt}%
\definecolor{currentstroke}{rgb}{0.000000,0.000000,0.000000}%
\pgfsetstrokecolor{currentstroke}%
\pgfsetstrokeopacity{0.000000}%
\pgfsetdash{}{0pt}%
\pgfpathmoveto{\pgfqpoint{0.973337in}{1.262333in}}%
\pgfpathlineto{\pgfqpoint{0.963825in}{1.264056in}}%
\pgfpathlineto{\pgfqpoint{0.964301in}{1.266613in}}%
\pgfpathlineto{\pgfqpoint{0.957377in}{1.275606in}}%
\pgfpathlineto{\pgfqpoint{0.959849in}{1.282582in}}%
\pgfpathlineto{\pgfqpoint{0.957665in}{1.289984in}}%
\pgfpathlineto{\pgfqpoint{0.950156in}{1.292847in}}%
\pgfpathlineto{\pgfqpoint{0.943668in}{1.304520in}}%
\pgfpathlineto{\pgfqpoint{0.947451in}{1.310370in}}%
\pgfpathlineto{\pgfqpoint{0.937610in}{1.310762in}}%
\pgfpathlineto{\pgfqpoint{0.925436in}{1.305547in}}%
\pgfpathlineto{\pgfqpoint{0.923286in}{1.309346in}}%
\pgfpathlineto{\pgfqpoint{0.914647in}{1.310630in}}%
\pgfpathlineto{\pgfqpoint{0.913219in}{1.297816in}}%
\pgfpathlineto{\pgfqpoint{0.914434in}{1.291420in}}%
\pgfpathlineto{\pgfqpoint{0.910706in}{1.281361in}}%
\pgfpathlineto{\pgfqpoint{0.901338in}{1.275958in}}%
\pgfpathlineto{\pgfqpoint{0.878712in}{1.280320in}}%
\pgfpathlineto{\pgfqpoint{0.816917in}{1.293210in}}%
\pgfpathlineto{\pgfqpoint{0.832155in}{1.363654in}}%
\pgfpathlineto{\pgfqpoint{0.892707in}{1.351230in}}%
\pgfpathlineto{\pgfqpoint{0.908601in}{1.353242in}}%
\pgfpathlineto{\pgfqpoint{0.924294in}{1.367661in}}%
\pgfpathlineto{\pgfqpoint{0.934476in}{1.365698in}}%
\pgfpathlineto{\pgfqpoint{0.942058in}{1.369750in}}%
\pgfpathlineto{\pgfqpoint{0.947822in}{1.365378in}}%
\pgfpathlineto{\pgfqpoint{0.951603in}{1.367889in}}%
\pgfpathlineto{\pgfqpoint{0.960688in}{1.366618in}}%
\pgfpathlineto{\pgfqpoint{0.963794in}{1.361956in}}%
\pgfpathlineto{\pgfqpoint{0.970085in}{1.358692in}}%
\pgfpathlineto{\pgfqpoint{0.969963in}{1.352958in}}%
\pgfpathlineto{\pgfqpoint{0.972493in}{1.347031in}}%
\pgfpathlineto{\pgfqpoint{0.981285in}{1.351407in}}%
\pgfpathlineto{\pgfqpoint{0.978339in}{1.334944in}}%
\pgfpathlineto{\pgfqpoint{1.005052in}{1.330045in}}%
\pgfpathlineto{\pgfqpoint{1.029601in}{1.326060in}}%
\pgfpathlineto{\pgfqpoint{1.027613in}{1.314125in}}%
\pgfpathlineto{\pgfqpoint{1.019386in}{1.315423in}}%
\pgfpathlineto{\pgfqpoint{1.015074in}{1.317685in}}%
\pgfpathlineto{\pgfqpoint{1.008923in}{1.315572in}}%
\pgfpathlineto{\pgfqpoint{0.994498in}{1.314874in}}%
\pgfpathlineto{\pgfqpoint{0.989953in}{1.316140in}}%
\pgfpathlineto{\pgfqpoint{0.981761in}{1.312933in}}%
\pgfpathlineto{\pgfqpoint{0.973337in}{1.262333in}}%
\pgfpathclose%
\pgfusepath{fill}%
\end{pgfscope}%
\begin{pgfscope}%
\pgfpathrectangle{\pgfqpoint{0.100000in}{0.100000in}}{\pgfqpoint{3.007045in}{1.925000in}}%
\pgfusepath{clip}%
\pgfsetbuttcap%
\pgfsetmiterjoin%
\definecolor{currentfill}{rgb}{0.215686,0.529412,0.754248}%
\pgfsetfillcolor{currentfill}%
\pgfsetlinewidth{0.000000pt}%
\definecolor{currentstroke}{rgb}{0.000000,0.000000,0.000000}%
\pgfsetstrokecolor{currentstroke}%
\pgfsetstrokeopacity{0.000000}%
\pgfsetdash{}{0pt}%
\pgfpathmoveto{\pgfqpoint{1.730711in}{1.250739in}}%
\pgfpathlineto{\pgfqpoint{1.699296in}{1.251149in}}%
\pgfpathlineto{\pgfqpoint{1.699583in}{1.268350in}}%
\pgfpathlineto{\pgfqpoint{1.677475in}{1.268823in}}%
\pgfpathlineto{\pgfqpoint{1.678558in}{1.314735in}}%
\pgfpathlineto{\pgfqpoint{1.701276in}{1.314251in}}%
\pgfpathlineto{\pgfqpoint{1.701191in}{1.310424in}}%
\pgfpathlineto{\pgfqpoint{1.717094in}{1.310143in}}%
\pgfpathlineto{\pgfqpoint{1.719990in}{1.317507in}}%
\pgfpathlineto{\pgfqpoint{1.716390in}{1.322247in}}%
\pgfpathlineto{\pgfqpoint{1.739290in}{1.321824in}}%
\pgfpathlineto{\pgfqpoint{1.744960in}{1.321802in}}%
\pgfpathlineto{\pgfqpoint{1.744815in}{1.304579in}}%
\pgfpathlineto{\pgfqpoint{1.747728in}{1.298216in}}%
\pgfpathlineto{\pgfqpoint{1.764683in}{1.298101in}}%
\pgfpathlineto{\pgfqpoint{1.764595in}{1.275306in}}%
\pgfpathlineto{\pgfqpoint{1.753309in}{1.275315in}}%
\pgfpathlineto{\pgfqpoint{1.753153in}{1.258349in}}%
\pgfpathlineto{\pgfqpoint{1.731811in}{1.258499in}}%
\pgfpathlineto{\pgfqpoint{1.730711in}{1.250739in}}%
\pgfpathclose%
\pgfusepath{fill}%
\end{pgfscope}%
\begin{pgfscope}%
\pgfpathrectangle{\pgfqpoint{0.100000in}{0.100000in}}{\pgfqpoint{3.007045in}{1.925000in}}%
\pgfusepath{clip}%
\pgfsetbuttcap%
\pgfsetmiterjoin%
\definecolor{currentfill}{rgb}{0.371688,0.649627,0.820515}%
\pgfsetfillcolor{currentfill}%
\pgfsetlinewidth{0.000000pt}%
\definecolor{currentstroke}{rgb}{0.000000,0.000000,0.000000}%
\pgfsetstrokecolor{currentstroke}%
\pgfsetstrokeopacity{0.000000}%
\pgfsetdash{}{0pt}%
\pgfpathmoveto{\pgfqpoint{1.928151in}{1.455464in}}%
\pgfpathlineto{\pgfqpoint{1.928977in}{1.432597in}}%
\pgfpathlineto{\pgfqpoint{1.894560in}{1.431480in}}%
\pgfpathlineto{\pgfqpoint{1.893937in}{1.453516in}}%
\pgfpathlineto{\pgfqpoint{1.882427in}{1.454160in}}%
\pgfpathlineto{\pgfqpoint{1.865473in}{1.453759in}}%
\pgfpathlineto{\pgfqpoint{1.864898in}{1.494975in}}%
\pgfpathlineto{\pgfqpoint{1.864831in}{1.497891in}}%
\pgfpathlineto{\pgfqpoint{1.876437in}{1.499997in}}%
\pgfpathlineto{\pgfqpoint{1.876275in}{1.505700in}}%
\pgfpathlineto{\pgfqpoint{1.884271in}{1.505208in}}%
\pgfpathlineto{\pgfqpoint{1.889296in}{1.502422in}}%
\pgfpathlineto{\pgfqpoint{1.898889in}{1.500432in}}%
\pgfpathlineto{\pgfqpoint{1.903859in}{1.493796in}}%
\pgfpathlineto{\pgfqpoint{1.911427in}{1.491667in}}%
\pgfpathlineto{\pgfqpoint{1.917765in}{1.487298in}}%
\pgfpathlineto{\pgfqpoint{1.921271in}{1.478116in}}%
\pgfpathlineto{\pgfqpoint{1.910791in}{1.477769in}}%
\pgfpathlineto{\pgfqpoint{1.911530in}{1.454907in}}%
\pgfpathlineto{\pgfqpoint{1.928151in}{1.455464in}}%
\pgfpathclose%
\pgfusepath{fill}%
\end{pgfscope}%
\begin{pgfscope}%
\pgfpathrectangle{\pgfqpoint{0.100000in}{0.100000in}}{\pgfqpoint{3.007045in}{1.925000in}}%
\pgfusepath{clip}%
\pgfsetbuttcap%
\pgfsetmiterjoin%
\definecolor{currentfill}{rgb}{0.290980,0.594510,0.789020}%
\pgfsetfillcolor{currentfill}%
\pgfsetlinewidth{0.000000pt}%
\definecolor{currentstroke}{rgb}{0.000000,0.000000,0.000000}%
\pgfsetstrokecolor{currentstroke}%
\pgfsetstrokeopacity{0.000000}%
\pgfsetdash{}{0pt}%
\pgfpathmoveto{\pgfqpoint{1.637239in}{1.477803in}}%
\pgfpathlineto{\pgfqpoint{1.682752in}{1.476661in}}%
\pgfpathlineto{\pgfqpoint{1.682306in}{1.453680in}}%
\pgfpathlineto{\pgfqpoint{1.648008in}{1.454576in}}%
\pgfpathlineto{\pgfqpoint{1.636440in}{1.454988in}}%
\pgfpathlineto{\pgfqpoint{1.637239in}{1.477803in}}%
\pgfpathclose%
\pgfusepath{fill}%
\end{pgfscope}%
\begin{pgfscope}%
\pgfpathrectangle{\pgfqpoint{0.100000in}{0.100000in}}{\pgfqpoint{3.007045in}{1.925000in}}%
\pgfusepath{clip}%
\pgfsetbuttcap%
\pgfsetmiterjoin%
\definecolor{currentfill}{rgb}{0.316201,0.611734,0.798862}%
\pgfsetfillcolor{currentfill}%
\pgfsetlinewidth{0.000000pt}%
\definecolor{currentstroke}{rgb}{0.000000,0.000000,0.000000}%
\pgfsetstrokecolor{currentstroke}%
\pgfsetstrokeopacity{0.000000}%
\pgfsetdash{}{0pt}%
\pgfpathmoveto{\pgfqpoint{1.755176in}{1.198941in}}%
\pgfpathlineto{\pgfqpoint{1.755095in}{1.176082in}}%
\pgfpathlineto{\pgfqpoint{1.743716in}{1.176139in}}%
\pgfpathlineto{\pgfqpoint{1.732315in}{1.176225in}}%
\pgfpathlineto{\pgfqpoint{1.732247in}{1.170499in}}%
\pgfpathlineto{\pgfqpoint{1.719727in}{1.170628in}}%
\pgfpathlineto{\pgfqpoint{1.709420in}{1.170767in}}%
\pgfpathlineto{\pgfqpoint{1.709799in}{1.199381in}}%
\pgfpathlineto{\pgfqpoint{1.755176in}{1.198941in}}%
\pgfpathclose%
\pgfusepath{fill}%
\end{pgfscope}%
\begin{pgfscope}%
\pgfpathrectangle{\pgfqpoint{0.100000in}{0.100000in}}{\pgfqpoint{3.007045in}{1.925000in}}%
\pgfusepath{clip}%
\pgfsetbuttcap%
\pgfsetmiterjoin%
\definecolor{currentfill}{rgb}{0.523137,0.739193,0.861546}%
\pgfsetfillcolor{currentfill}%
\pgfsetlinewidth{0.000000pt}%
\definecolor{currentstroke}{rgb}{0.000000,0.000000,0.000000}%
\pgfsetstrokecolor{currentstroke}%
\pgfsetstrokeopacity{0.000000}%
\pgfsetdash{}{0pt}%
\pgfpathmoveto{\pgfqpoint{2.029155in}{1.469240in}}%
\pgfpathlineto{\pgfqpoint{2.030394in}{1.446824in}}%
\pgfpathlineto{\pgfqpoint{2.021528in}{1.446204in}}%
\pgfpathlineto{\pgfqpoint{2.017195in}{1.455365in}}%
\pgfpathlineto{\pgfqpoint{2.012501in}{1.460181in}}%
\pgfpathlineto{\pgfqpoint{2.012462in}{1.466036in}}%
\pgfpathlineto{\pgfqpoint{2.008307in}{1.475250in}}%
\pgfpathlineto{\pgfqpoint{2.013699in}{1.486047in}}%
\pgfpathlineto{\pgfqpoint{1.994330in}{1.484956in}}%
\pgfpathlineto{\pgfqpoint{1.993541in}{1.496568in}}%
\pgfpathlineto{\pgfqpoint{1.970609in}{1.495307in}}%
\pgfpathlineto{\pgfqpoint{1.970355in}{1.501053in}}%
\pgfpathlineto{\pgfqpoint{1.964639in}{1.500815in}}%
\pgfpathlineto{\pgfqpoint{1.962214in}{1.552559in}}%
\pgfpathlineto{\pgfqpoint{1.961960in}{1.558341in}}%
\pgfpathlineto{\pgfqpoint{1.973425in}{1.558760in}}%
\pgfpathlineto{\pgfqpoint{2.003003in}{1.560543in}}%
\pgfpathlineto{\pgfqpoint{2.002328in}{1.571903in}}%
\pgfpathlineto{\pgfqpoint{2.030834in}{1.573678in}}%
\pgfpathlineto{\pgfqpoint{2.031351in}{1.567858in}}%
\pgfpathlineto{\pgfqpoint{2.032705in}{1.544832in}}%
\pgfpathlineto{\pgfqpoint{2.042126in}{1.545379in}}%
\pgfpathlineto{\pgfqpoint{2.043986in}{1.516491in}}%
\pgfpathlineto{\pgfqpoint{2.045773in}{1.487575in}}%
\pgfpathlineto{\pgfqpoint{2.028120in}{1.486633in}}%
\pgfpathlineto{\pgfqpoint{2.029155in}{1.469240in}}%
\pgfpathclose%
\pgfusepath{fill}%
\end{pgfscope}%
\begin{pgfscope}%
\pgfpathrectangle{\pgfqpoint{0.100000in}{0.100000in}}{\pgfqpoint{3.007045in}{1.925000in}}%
\pgfusepath{clip}%
\pgfsetbuttcap%
\pgfsetmiterjoin%
\definecolor{currentfill}{rgb}{0.386820,0.659962,0.826421}%
\pgfsetfillcolor{currentfill}%
\pgfsetlinewidth{0.000000pt}%
\definecolor{currentstroke}{rgb}{0.000000,0.000000,0.000000}%
\pgfsetstrokecolor{currentstroke}%
\pgfsetstrokeopacity{0.000000}%
\pgfsetdash{}{0pt}%
\pgfpathmoveto{\pgfqpoint{0.459905in}{1.401168in}}%
\pgfpathlineto{\pgfqpoint{0.451708in}{1.400095in}}%
\pgfpathlineto{\pgfqpoint{0.446584in}{1.393519in}}%
\pgfpathlineto{\pgfqpoint{0.443854in}{1.394476in}}%
\pgfpathlineto{\pgfqpoint{0.437071in}{1.389645in}}%
\pgfpathlineto{\pgfqpoint{0.425415in}{1.393166in}}%
\pgfpathlineto{\pgfqpoint{0.423677in}{1.387762in}}%
\pgfpathlineto{\pgfqpoint{0.380635in}{1.401231in}}%
\pgfpathlineto{\pgfqpoint{0.384397in}{1.412475in}}%
\pgfpathlineto{\pgfqpoint{0.385058in}{1.423808in}}%
\pgfpathlineto{\pgfqpoint{0.387316in}{1.430255in}}%
\pgfpathlineto{\pgfqpoint{0.384276in}{1.433228in}}%
\pgfpathlineto{\pgfqpoint{0.385422in}{1.437519in}}%
\pgfpathlineto{\pgfqpoint{0.390077in}{1.441086in}}%
\pgfpathlineto{\pgfqpoint{0.395246in}{1.441320in}}%
\pgfpathlineto{\pgfqpoint{0.399108in}{1.444237in}}%
\pgfpathlineto{\pgfqpoint{0.407911in}{1.446662in}}%
\pgfpathlineto{\pgfqpoint{0.407285in}{1.452726in}}%
\pgfpathlineto{\pgfqpoint{0.412778in}{1.458416in}}%
\pgfpathlineto{\pgfqpoint{0.420836in}{1.470570in}}%
\pgfpathlineto{\pgfqpoint{0.424324in}{1.471482in}}%
\pgfpathlineto{\pgfqpoint{0.429023in}{1.481634in}}%
\pgfpathlineto{\pgfqpoint{0.484393in}{1.464821in}}%
\pgfpathlineto{\pgfqpoint{0.479723in}{1.447096in}}%
\pgfpathlineto{\pgfqpoint{0.470680in}{1.418127in}}%
\pgfpathlineto{\pgfqpoint{0.462527in}{1.420581in}}%
\pgfpathlineto{\pgfqpoint{0.462117in}{1.414113in}}%
\pgfpathlineto{\pgfqpoint{0.466986in}{1.410310in}}%
\pgfpathlineto{\pgfqpoint{0.464449in}{1.404081in}}%
\pgfpathlineto{\pgfqpoint{0.459905in}{1.401168in}}%
\pgfpathclose%
\pgfusepath{fill}%
\end{pgfscope}%
\begin{pgfscope}%
\pgfpathrectangle{\pgfqpoint{0.100000in}{0.100000in}}{\pgfqpoint{3.007045in}{1.925000in}}%
\pgfusepath{clip}%
\pgfsetbuttcap%
\pgfsetmiterjoin%
\definecolor{currentfill}{rgb}{0.391865,0.663406,0.828389}%
\pgfsetfillcolor{currentfill}%
\pgfsetlinewidth{0.000000pt}%
\definecolor{currentstroke}{rgb}{0.000000,0.000000,0.000000}%
\pgfsetstrokecolor{currentstroke}%
\pgfsetstrokeopacity{0.000000}%
\pgfsetdash{}{0pt}%
\pgfpathmoveto{\pgfqpoint{2.256834in}{1.132262in}}%
\pgfpathlineto{\pgfqpoint{2.256292in}{1.137004in}}%
\pgfpathlineto{\pgfqpoint{2.250151in}{1.140216in}}%
\pgfpathlineto{\pgfqpoint{2.249568in}{1.145924in}}%
\pgfpathlineto{\pgfqpoint{2.253275in}{1.146329in}}%
\pgfpathlineto{\pgfqpoint{2.261118in}{1.153025in}}%
\pgfpathlineto{\pgfqpoint{2.273431in}{1.154396in}}%
\pgfpathlineto{\pgfqpoint{2.275093in}{1.139834in}}%
\pgfpathlineto{\pgfqpoint{2.283647in}{1.145438in}}%
\pgfpathlineto{\pgfqpoint{2.288955in}{1.139422in}}%
\pgfpathlineto{\pgfqpoint{2.282277in}{1.134362in}}%
\pgfpathlineto{\pgfqpoint{2.277544in}{1.133106in}}%
\pgfpathlineto{\pgfqpoint{2.272501in}{1.126660in}}%
\pgfpathlineto{\pgfqpoint{2.264410in}{1.127604in}}%
\pgfpathlineto{\pgfqpoint{2.264234in}{1.131721in}}%
\pgfpathlineto{\pgfqpoint{2.256834in}{1.132262in}}%
\pgfpathclose%
\pgfusepath{fill}%
\end{pgfscope}%
\begin{pgfscope}%
\pgfpathrectangle{\pgfqpoint{0.100000in}{0.100000in}}{\pgfqpoint{3.007045in}{1.925000in}}%
\pgfusepath{clip}%
\pgfsetbuttcap%
\pgfsetmiterjoin%
\definecolor{currentfill}{rgb}{0.179146,0.492872,0.735425}%
\pgfsetfillcolor{currentfill}%
\pgfsetlinewidth{0.000000pt}%
\definecolor{currentstroke}{rgb}{0.000000,0.000000,0.000000}%
\pgfsetstrokecolor{currentstroke}%
\pgfsetstrokeopacity{0.000000}%
\pgfsetdash{}{0pt}%
\pgfpathmoveto{\pgfqpoint{1.610215in}{1.316599in}}%
\pgfpathlineto{\pgfqpoint{1.609594in}{1.293697in}}%
\pgfpathlineto{\pgfqpoint{1.587220in}{1.294590in}}%
\pgfpathlineto{\pgfqpoint{1.564548in}{1.295579in}}%
\pgfpathlineto{\pgfqpoint{1.565174in}{1.318418in}}%
\pgfpathlineto{\pgfqpoint{1.541840in}{1.319542in}}%
\pgfpathlineto{\pgfqpoint{1.543010in}{1.342313in}}%
\pgfpathlineto{\pgfqpoint{1.565129in}{1.341345in}}%
\pgfpathlineto{\pgfqpoint{1.588274in}{1.340363in}}%
\pgfpathlineto{\pgfqpoint{1.587433in}{1.317426in}}%
\pgfpathlineto{\pgfqpoint{1.610215in}{1.316599in}}%
\pgfpathclose%
\pgfusepath{fill}%
\end{pgfscope}%
\begin{pgfscope}%
\pgfpathrectangle{\pgfqpoint{0.100000in}{0.100000in}}{\pgfqpoint{3.007045in}{1.925000in}}%
\pgfusepath{clip}%
\pgfsetbuttcap%
\pgfsetmiterjoin%
\definecolor{currentfill}{rgb}{0.199446,0.513172,0.745882}%
\pgfsetfillcolor{currentfill}%
\pgfsetlinewidth{0.000000pt}%
\definecolor{currentstroke}{rgb}{0.000000,0.000000,0.000000}%
\pgfsetstrokecolor{currentstroke}%
\pgfsetstrokeopacity{0.000000}%
\pgfsetdash{}{0pt}%
\pgfpathmoveto{\pgfqpoint{1.656900in}{1.355436in}}%
\pgfpathlineto{\pgfqpoint{1.673964in}{1.354981in}}%
\pgfpathlineto{\pgfqpoint{1.674656in}{1.382036in}}%
\pgfpathlineto{\pgfqpoint{1.684732in}{1.378050in}}%
\pgfpathlineto{\pgfqpoint{1.685257in}{1.403080in}}%
\pgfpathlineto{\pgfqpoint{1.702210in}{1.402774in}}%
\pgfpathlineto{\pgfqpoint{1.699433in}{1.400596in}}%
\pgfpathlineto{\pgfqpoint{1.699908in}{1.394517in}}%
\pgfpathlineto{\pgfqpoint{1.697951in}{1.391327in}}%
\pgfpathlineto{\pgfqpoint{1.730762in}{1.390911in}}%
\pgfpathlineto{\pgfqpoint{1.730520in}{1.367928in}}%
\pgfpathlineto{\pgfqpoint{1.736283in}{1.367899in}}%
\pgfpathlineto{\pgfqpoint{1.739732in}{1.362155in}}%
\pgfpathlineto{\pgfqpoint{1.739541in}{1.344779in}}%
\pgfpathlineto{\pgfqpoint{1.739290in}{1.321824in}}%
\pgfpathlineto{\pgfqpoint{1.716390in}{1.322247in}}%
\pgfpathlineto{\pgfqpoint{1.716907in}{1.329278in}}%
\pgfpathlineto{\pgfqpoint{1.713888in}{1.330082in}}%
\pgfpathlineto{\pgfqpoint{1.706612in}{1.341900in}}%
\pgfpathlineto{\pgfqpoint{1.706125in}{1.349454in}}%
\pgfpathlineto{\pgfqpoint{1.688055in}{1.349861in}}%
\pgfpathlineto{\pgfqpoint{1.683358in}{1.349049in}}%
\pgfpathlineto{\pgfqpoint{1.683057in}{1.337584in}}%
\pgfpathlineto{\pgfqpoint{1.673497in}{1.337786in}}%
\pgfpathlineto{\pgfqpoint{1.656422in}{1.338221in}}%
\pgfpathlineto{\pgfqpoint{1.656900in}{1.355436in}}%
\pgfpathclose%
\pgfusepath{fill}%
\end{pgfscope}%
\begin{pgfscope}%
\pgfpathrectangle{\pgfqpoint{0.100000in}{0.100000in}}{\pgfqpoint{3.007045in}{1.925000in}}%
\pgfusepath{clip}%
\pgfsetbuttcap%
\pgfsetmiterjoin%
\definecolor{currentfill}{rgb}{0.491765,0.721968,0.854779}%
\pgfsetfillcolor{currentfill}%
\pgfsetlinewidth{0.000000pt}%
\definecolor{currentstroke}{rgb}{0.000000,0.000000,0.000000}%
\pgfsetstrokecolor{currentstroke}%
\pgfsetstrokeopacity{0.000000}%
\pgfsetdash{}{0pt}%
\pgfpathmoveto{\pgfqpoint{2.559434in}{0.779009in}}%
\pgfpathlineto{\pgfqpoint{2.556087in}{0.772933in}}%
\pgfpathlineto{\pgfqpoint{2.548564in}{0.780580in}}%
\pgfpathlineto{\pgfqpoint{2.535487in}{0.765030in}}%
\pgfpathlineto{\pgfqpoint{2.529482in}{0.767093in}}%
\pgfpathlineto{\pgfqpoint{2.526324in}{0.773753in}}%
\pgfpathlineto{\pgfqpoint{2.525012in}{0.782861in}}%
\pgfpathlineto{\pgfqpoint{2.519904in}{0.788430in}}%
\pgfpathlineto{\pgfqpoint{2.519777in}{0.793317in}}%
\pgfpathlineto{\pgfqpoint{2.516660in}{0.795332in}}%
\pgfpathlineto{\pgfqpoint{2.507037in}{0.798645in}}%
\pgfpathlineto{\pgfqpoint{2.503376in}{0.803575in}}%
\pgfpathlineto{\pgfqpoint{2.521533in}{0.825837in}}%
\pgfpathlineto{\pgfqpoint{2.525782in}{0.825781in}}%
\pgfpathlineto{\pgfqpoint{2.532363in}{0.822168in}}%
\pgfpathlineto{\pgfqpoint{2.545370in}{0.817990in}}%
\pgfpathlineto{\pgfqpoint{2.556125in}{0.810180in}}%
\pgfpathlineto{\pgfqpoint{2.542055in}{0.797900in}}%
\pgfpathlineto{\pgfqpoint{2.548566in}{0.790531in}}%
\pgfpathlineto{\pgfqpoint{2.554010in}{0.786819in}}%
\pgfpathlineto{\pgfqpoint{2.559434in}{0.779009in}}%
\pgfpathclose%
\pgfusepath{fill}%
\end{pgfscope}%
\begin{pgfscope}%
\pgfpathrectangle{\pgfqpoint{0.100000in}{0.100000in}}{\pgfqpoint{3.007045in}{1.925000in}}%
\pgfusepath{clip}%
\pgfsetbuttcap%
\pgfsetmiterjoin%
\definecolor{currentfill}{rgb}{0.321246,0.615179,0.800830}%
\pgfsetfillcolor{currentfill}%
\pgfsetlinewidth{0.000000pt}%
\definecolor{currentstroke}{rgb}{0.000000,0.000000,0.000000}%
\pgfsetstrokecolor{currentstroke}%
\pgfsetstrokeopacity{0.000000}%
\pgfsetdash{}{0pt}%
\pgfpathmoveto{\pgfqpoint{0.936059in}{0.163351in}}%
\pgfpathlineto{\pgfqpoint{0.936182in}{0.166254in}}%
\pgfpathlineto{\pgfqpoint{0.938191in}{0.164676in}}%
\pgfpathlineto{\pgfqpoint{0.936059in}{0.163351in}}%
\pgfpathclose%
\pgfusepath{fill}%
\end{pgfscope}%
\begin{pgfscope}%
\pgfpathrectangle{\pgfqpoint{0.100000in}{0.100000in}}{\pgfqpoint{3.007045in}{1.925000in}}%
\pgfusepath{clip}%
\pgfsetbuttcap%
\pgfsetmiterjoin%
\definecolor{currentfill}{rgb}{0.321246,0.615179,0.800830}%
\pgfsetfillcolor{currentfill}%
\pgfsetlinewidth{0.000000pt}%
\definecolor{currentstroke}{rgb}{0.000000,0.000000,0.000000}%
\pgfsetstrokecolor{currentstroke}%
\pgfsetstrokeopacity{0.000000}%
\pgfsetdash{}{0pt}%
\pgfpathmoveto{\pgfqpoint{0.933022in}{0.152128in}}%
\pgfpathlineto{\pgfqpoint{0.933743in}{0.154327in}}%
\pgfpathlineto{\pgfqpoint{0.935121in}{0.154547in}}%
\pgfpathlineto{\pgfqpoint{0.934597in}{0.151513in}}%
\pgfpathlineto{\pgfqpoint{0.933022in}{0.152128in}}%
\pgfpathclose%
\pgfusepath{fill}%
\end{pgfscope}%
\begin{pgfscope}%
\pgfpathrectangle{\pgfqpoint{0.100000in}{0.100000in}}{\pgfqpoint{3.007045in}{1.925000in}}%
\pgfusepath{clip}%
\pgfsetbuttcap%
\pgfsetmiterjoin%
\definecolor{currentfill}{rgb}{0.321246,0.615179,0.800830}%
\pgfsetfillcolor{currentfill}%
\pgfsetlinewidth{0.000000pt}%
\definecolor{currentstroke}{rgb}{0.000000,0.000000,0.000000}%
\pgfsetstrokecolor{currentstroke}%
\pgfsetstrokeopacity{0.000000}%
\pgfsetdash{}{0pt}%
\pgfpathmoveto{\pgfqpoint{0.935935in}{0.137518in}}%
\pgfpathlineto{\pgfqpoint{0.934705in}{0.142381in}}%
\pgfpathlineto{\pgfqpoint{0.935452in}{0.143773in}}%
\pgfpathlineto{\pgfqpoint{0.933837in}{0.145038in}}%
\pgfpathlineto{\pgfqpoint{0.935285in}{0.148282in}}%
\pgfpathlineto{\pgfqpoint{0.935191in}{0.150312in}}%
\pgfpathlineto{\pgfqpoint{0.936976in}{0.151065in}}%
\pgfpathlineto{\pgfqpoint{0.937372in}{0.148133in}}%
\pgfpathlineto{\pgfqpoint{0.935770in}{0.146761in}}%
\pgfpathlineto{\pgfqpoint{0.936802in}{0.145789in}}%
\pgfpathlineto{\pgfqpoint{0.936825in}{0.141016in}}%
\pgfpathlineto{\pgfqpoint{0.939296in}{0.140402in}}%
\pgfpathlineto{\pgfqpoint{0.935935in}{0.137518in}}%
\pgfpathclose%
\pgfusepath{fill}%
\end{pgfscope}%
\begin{pgfscope}%
\pgfpathrectangle{\pgfqpoint{0.100000in}{0.100000in}}{\pgfqpoint{3.007045in}{1.925000in}}%
\pgfusepath{clip}%
\pgfsetbuttcap%
\pgfsetmiterjoin%
\definecolor{currentfill}{rgb}{0.321246,0.615179,0.800830}%
\pgfsetfillcolor{currentfill}%
\pgfsetlinewidth{0.000000pt}%
\definecolor{currentstroke}{rgb}{0.000000,0.000000,0.000000}%
\pgfsetstrokecolor{currentstroke}%
\pgfsetstrokeopacity{0.000000}%
\pgfsetdash{}{0pt}%
\pgfpathmoveto{\pgfqpoint{0.957491in}{0.142127in}}%
\pgfpathlineto{\pgfqpoint{0.957670in}{0.139655in}}%
\pgfpathlineto{\pgfqpoint{0.956604in}{0.137353in}}%
\pgfpathlineto{\pgfqpoint{0.955422in}{0.137331in}}%
\pgfpathlineto{\pgfqpoint{0.954352in}{0.139936in}}%
\pgfpathlineto{\pgfqpoint{0.955529in}{0.140546in}}%
\pgfpathlineto{\pgfqpoint{0.955668in}{0.143108in}}%
\pgfpathlineto{\pgfqpoint{0.957491in}{0.142127in}}%
\pgfpathclose%
\pgfusepath{fill}%
\end{pgfscope}%
\begin{pgfscope}%
\pgfpathrectangle{\pgfqpoint{0.100000in}{0.100000in}}{\pgfqpoint{3.007045in}{1.925000in}}%
\pgfusepath{clip}%
\pgfsetbuttcap%
\pgfsetmiterjoin%
\definecolor{currentfill}{rgb}{0.321246,0.615179,0.800830}%
\pgfsetfillcolor{currentfill}%
\pgfsetlinewidth{0.000000pt}%
\definecolor{currentstroke}{rgb}{0.000000,0.000000,0.000000}%
\pgfsetstrokecolor{currentstroke}%
\pgfsetstrokeopacity{0.000000}%
\pgfsetdash{}{0pt}%
\pgfpathmoveto{\pgfqpoint{0.952217in}{0.157810in}}%
\pgfpathlineto{\pgfqpoint{0.953820in}{0.154455in}}%
\pgfpathlineto{\pgfqpoint{0.955416in}{0.154694in}}%
\pgfpathlineto{\pgfqpoint{0.956678in}{0.156700in}}%
\pgfpathlineto{\pgfqpoint{0.956199in}{0.158645in}}%
\pgfpathlineto{\pgfqpoint{0.958436in}{0.159148in}}%
\pgfpathlineto{\pgfqpoint{0.958681in}{0.160262in}}%
\pgfpathlineto{\pgfqpoint{0.960550in}{0.161270in}}%
\pgfpathlineto{\pgfqpoint{0.962912in}{0.161570in}}%
\pgfpathlineto{\pgfqpoint{0.964181in}{0.158832in}}%
\pgfpathlineto{\pgfqpoint{0.966014in}{0.159668in}}%
\pgfpathlineto{\pgfqpoint{0.965820in}{0.161252in}}%
\pgfpathlineto{\pgfqpoint{0.968598in}{0.161733in}}%
\pgfpathlineto{\pgfqpoint{0.970864in}{0.160909in}}%
\pgfpathlineto{\pgfqpoint{0.972584in}{0.162951in}}%
\pgfpathlineto{\pgfqpoint{0.972690in}{0.165137in}}%
\pgfpathlineto{\pgfqpoint{0.975165in}{0.163335in}}%
\pgfpathlineto{\pgfqpoint{0.977660in}{0.159042in}}%
\pgfpathlineto{\pgfqpoint{0.978047in}{0.156183in}}%
\pgfpathlineto{\pgfqpoint{0.979828in}{0.154295in}}%
\pgfpathlineto{\pgfqpoint{0.981654in}{0.154152in}}%
\pgfpathlineto{\pgfqpoint{0.981934in}{0.151716in}}%
\pgfpathlineto{\pgfqpoint{0.981007in}{0.149518in}}%
\pgfpathlineto{\pgfqpoint{0.977742in}{0.146985in}}%
\pgfpathlineto{\pgfqpoint{0.977518in}{0.144790in}}%
\pgfpathlineto{\pgfqpoint{0.976466in}{0.142679in}}%
\pgfpathlineto{\pgfqpoint{0.976184in}{0.138931in}}%
\pgfpathlineto{\pgfqpoint{0.975258in}{0.135551in}}%
\pgfpathlineto{\pgfqpoint{0.973176in}{0.134289in}}%
\pgfpathlineto{\pgfqpoint{0.967906in}{0.129628in}}%
\pgfpathlineto{\pgfqpoint{0.965959in}{0.130065in}}%
\pgfpathlineto{\pgfqpoint{0.963807in}{0.128960in}}%
\pgfpathlineto{\pgfqpoint{0.961890in}{0.130440in}}%
\pgfpathlineto{\pgfqpoint{0.959954in}{0.130233in}}%
\pgfpathlineto{\pgfqpoint{0.961401in}{0.136798in}}%
\pgfpathlineto{\pgfqpoint{0.961289in}{0.138714in}}%
\pgfpathlineto{\pgfqpoint{0.964237in}{0.140766in}}%
\pgfpathlineto{\pgfqpoint{0.965836in}{0.140351in}}%
\pgfpathlineto{\pgfqpoint{0.965862in}{0.142940in}}%
\pgfpathlineto{\pgfqpoint{0.967505in}{0.145942in}}%
\pgfpathlineto{\pgfqpoint{0.969065in}{0.151851in}}%
\pgfpathlineto{\pgfqpoint{0.966959in}{0.151516in}}%
\pgfpathlineto{\pgfqpoint{0.967397in}{0.149343in}}%
\pgfpathlineto{\pgfqpoint{0.966518in}{0.146372in}}%
\pgfpathlineto{\pgfqpoint{0.962740in}{0.141294in}}%
\pgfpathlineto{\pgfqpoint{0.959748in}{0.140140in}}%
\pgfpathlineto{\pgfqpoint{0.959257in}{0.141993in}}%
\pgfpathlineto{\pgfqpoint{0.956453in}{0.144467in}}%
\pgfpathlineto{\pgfqpoint{0.953985in}{0.143893in}}%
\pgfpathlineto{\pgfqpoint{0.952414in}{0.142851in}}%
\pgfpathlineto{\pgfqpoint{0.953506in}{0.147090in}}%
\pgfpathlineto{\pgfqpoint{0.955124in}{0.149168in}}%
\pgfpathlineto{\pgfqpoint{0.956748in}{0.149574in}}%
\pgfpathlineto{\pgfqpoint{0.957816in}{0.151764in}}%
\pgfpathlineto{\pgfqpoint{0.959671in}{0.153160in}}%
\pgfpathlineto{\pgfqpoint{0.959717in}{0.154964in}}%
\pgfpathlineto{\pgfqpoint{0.962294in}{0.155551in}}%
\pgfpathlineto{\pgfqpoint{0.962264in}{0.157330in}}%
\pgfpathlineto{\pgfqpoint{0.959964in}{0.158575in}}%
\pgfpathlineto{\pgfqpoint{0.957588in}{0.157646in}}%
\pgfpathlineto{\pgfqpoint{0.958906in}{0.156128in}}%
\pgfpathlineto{\pgfqpoint{0.956807in}{0.153413in}}%
\pgfpathlineto{\pgfqpoint{0.953556in}{0.150672in}}%
\pgfpathlineto{\pgfqpoint{0.951866in}{0.154211in}}%
\pgfpathlineto{\pgfqpoint{0.952217in}{0.157810in}}%
\pgfpathclose%
\pgfusepath{fill}%
\end{pgfscope}%
\begin{pgfscope}%
\pgfpathrectangle{\pgfqpoint{0.100000in}{0.100000in}}{\pgfqpoint{3.007045in}{1.925000in}}%
\pgfusepath{clip}%
\pgfsetbuttcap%
\pgfsetmiterjoin%
\definecolor{currentfill}{rgb}{0.321246,0.615179,0.800830}%
\pgfsetfillcolor{currentfill}%
\pgfsetlinewidth{0.000000pt}%
\definecolor{currentstroke}{rgb}{0.000000,0.000000,0.000000}%
\pgfsetstrokecolor{currentstroke}%
\pgfsetstrokeopacity{0.000000}%
\pgfsetdash{}{0pt}%
\pgfpathmoveto{\pgfqpoint{0.932602in}{0.158192in}}%
\pgfpathlineto{\pgfqpoint{0.930583in}{0.159046in}}%
\pgfpathlineto{\pgfqpoint{0.933442in}{0.160711in}}%
\pgfpathlineto{\pgfqpoint{0.934281in}{0.158960in}}%
\pgfpathlineto{\pgfqpoint{0.937131in}{0.159694in}}%
\pgfpathlineto{\pgfqpoint{0.937907in}{0.158017in}}%
\pgfpathlineto{\pgfqpoint{0.935667in}{0.157983in}}%
\pgfpathlineto{\pgfqpoint{0.934472in}{0.156058in}}%
\pgfpathlineto{\pgfqpoint{0.932593in}{0.155358in}}%
\pgfpathlineto{\pgfqpoint{0.931238in}{0.156086in}}%
\pgfpathlineto{\pgfqpoint{0.932602in}{0.158192in}}%
\pgfpathclose%
\pgfusepath{fill}%
\end{pgfscope}%
\begin{pgfscope}%
\pgfpathrectangle{\pgfqpoint{0.100000in}{0.100000in}}{\pgfqpoint{3.007045in}{1.925000in}}%
\pgfusepath{clip}%
\pgfsetbuttcap%
\pgfsetmiterjoin%
\definecolor{currentfill}{rgb}{0.321246,0.615179,0.800830}%
\pgfsetfillcolor{currentfill}%
\pgfsetlinewidth{0.000000pt}%
\definecolor{currentstroke}{rgb}{0.000000,0.000000,0.000000}%
\pgfsetstrokecolor{currentstroke}%
\pgfsetstrokeopacity{0.000000}%
\pgfsetdash{}{0pt}%
\pgfpathmoveto{\pgfqpoint{0.940503in}{0.148050in}}%
\pgfpathlineto{\pgfqpoint{0.940580in}{0.150009in}}%
\pgfpathlineto{\pgfqpoint{0.936401in}{0.153144in}}%
\pgfpathlineto{\pgfqpoint{0.939865in}{0.154873in}}%
\pgfpathlineto{\pgfqpoint{0.939188in}{0.156547in}}%
\pgfpathlineto{\pgfqpoint{0.940021in}{0.158363in}}%
\pgfpathlineto{\pgfqpoint{0.938477in}{0.159030in}}%
\pgfpathlineto{\pgfqpoint{0.937509in}{0.161120in}}%
\pgfpathlineto{\pgfqpoint{0.940507in}{0.163189in}}%
\pgfpathlineto{\pgfqpoint{0.940853in}{0.165668in}}%
\pgfpathlineto{\pgfqpoint{0.942540in}{0.166461in}}%
\pgfpathlineto{\pgfqpoint{0.942490in}{0.168936in}}%
\pgfpathlineto{\pgfqpoint{0.940371in}{0.169760in}}%
\pgfpathlineto{\pgfqpoint{0.935534in}{0.169087in}}%
\pgfpathlineto{\pgfqpoint{0.937048in}{0.170932in}}%
\pgfpathlineto{\pgfqpoint{0.938908in}{0.171837in}}%
\pgfpathlineto{\pgfqpoint{0.941019in}{0.171940in}}%
\pgfpathlineto{\pgfqpoint{0.940657in}{0.175170in}}%
\pgfpathlineto{\pgfqpoint{0.941701in}{0.177230in}}%
\pgfpathlineto{\pgfqpoint{0.943306in}{0.176542in}}%
\pgfpathlineto{\pgfqpoint{0.946910in}{0.174004in}}%
\pgfpathlineto{\pgfqpoint{0.946946in}{0.170521in}}%
\pgfpathlineto{\pgfqpoint{0.945241in}{0.169179in}}%
\pgfpathlineto{\pgfqpoint{0.945972in}{0.166130in}}%
\pgfpathlineto{\pgfqpoint{0.947228in}{0.166360in}}%
\pgfpathlineto{\pgfqpoint{0.948782in}{0.164350in}}%
\pgfpathlineto{\pgfqpoint{0.949856in}{0.160592in}}%
\pgfpathlineto{\pgfqpoint{0.949616in}{0.158564in}}%
\pgfpathlineto{\pgfqpoint{0.948423in}{0.157492in}}%
\pgfpathlineto{\pgfqpoint{0.950053in}{0.156174in}}%
\pgfpathlineto{\pgfqpoint{0.949902in}{0.154095in}}%
\pgfpathlineto{\pgfqpoint{0.948375in}{0.153768in}}%
\pgfpathlineto{\pgfqpoint{0.947443in}{0.155763in}}%
\pgfpathlineto{\pgfqpoint{0.946239in}{0.153864in}}%
\pgfpathlineto{\pgfqpoint{0.948621in}{0.152037in}}%
\pgfpathlineto{\pgfqpoint{0.948626in}{0.150372in}}%
\pgfpathlineto{\pgfqpoint{0.949742in}{0.147800in}}%
\pgfpathlineto{\pgfqpoint{0.948428in}{0.144747in}}%
\pgfpathlineto{\pgfqpoint{0.950997in}{0.145602in}}%
\pgfpathlineto{\pgfqpoint{0.950534in}{0.143772in}}%
\pgfpathlineto{\pgfqpoint{0.948944in}{0.142658in}}%
\pgfpathlineto{\pgfqpoint{0.947671in}{0.143289in}}%
\pgfpathlineto{\pgfqpoint{0.946770in}{0.140379in}}%
\pgfpathlineto{\pgfqpoint{0.949016in}{0.140644in}}%
\pgfpathlineto{\pgfqpoint{0.947609in}{0.137730in}}%
\pgfpathlineto{\pgfqpoint{0.947150in}{0.135497in}}%
\pgfpathlineto{\pgfqpoint{0.945969in}{0.134095in}}%
\pgfpathlineto{\pgfqpoint{0.943530in}{0.134941in}}%
\pgfpathlineto{\pgfqpoint{0.943106in}{0.139361in}}%
\pgfpathlineto{\pgfqpoint{0.941552in}{0.142463in}}%
\pgfpathlineto{\pgfqpoint{0.941399in}{0.144418in}}%
\pgfpathlineto{\pgfqpoint{0.942947in}{0.147821in}}%
\pgfpathlineto{\pgfqpoint{0.940503in}{0.148050in}}%
\pgfpathclose%
\pgfusepath{fill}%
\end{pgfscope}%
\begin{pgfscope}%
\pgfpathrectangle{\pgfqpoint{0.100000in}{0.100000in}}{\pgfqpoint{3.007045in}{1.925000in}}%
\pgfusepath{clip}%
\pgfsetbuttcap%
\pgfsetmiterjoin%
\definecolor{currentfill}{rgb}{0.248166,0.561892,0.770980}%
\pgfsetfillcolor{currentfill}%
\pgfsetlinewidth{0.000000pt}%
\definecolor{currentstroke}{rgb}{0.000000,0.000000,0.000000}%
\pgfsetstrokecolor{currentstroke}%
\pgfsetstrokeopacity{0.000000}%
\pgfsetdash{}{0pt}%
\pgfpathmoveto{\pgfqpoint{0.746029in}{1.809361in}}%
\pgfpathlineto{\pgfqpoint{0.740338in}{1.787082in}}%
\pgfpathlineto{\pgfqpoint{0.735917in}{1.779747in}}%
\pgfpathlineto{\pgfqpoint{0.726344in}{1.778628in}}%
\pgfpathlineto{\pgfqpoint{0.702590in}{1.784664in}}%
\pgfpathlineto{\pgfqpoint{0.675279in}{1.792196in}}%
\pgfpathlineto{\pgfqpoint{0.670581in}{1.789487in}}%
\pgfpathlineto{\pgfqpoint{0.668689in}{1.793331in}}%
\pgfpathlineto{\pgfqpoint{0.662371in}{1.789397in}}%
\pgfpathlineto{\pgfqpoint{0.651182in}{1.791391in}}%
\pgfpathlineto{\pgfqpoint{0.648904in}{1.795115in}}%
\pgfpathlineto{\pgfqpoint{0.651860in}{1.802864in}}%
\pgfpathlineto{\pgfqpoint{0.652827in}{1.812664in}}%
\pgfpathlineto{\pgfqpoint{0.652061in}{1.822418in}}%
\pgfpathlineto{\pgfqpoint{0.656028in}{1.830896in}}%
\pgfpathlineto{\pgfqpoint{0.657936in}{1.837275in}}%
\pgfpathlineto{\pgfqpoint{0.662469in}{1.839071in}}%
\pgfpathlineto{\pgfqpoint{0.665649in}{1.843218in}}%
\pgfpathlineto{\pgfqpoint{0.680569in}{1.839333in}}%
\pgfpathlineto{\pgfqpoint{0.682133in}{1.844866in}}%
\pgfpathlineto{\pgfqpoint{0.689171in}{1.848788in}}%
\pgfpathlineto{\pgfqpoint{0.691919in}{1.848024in}}%
\pgfpathlineto{\pgfqpoint{0.695822in}{1.854965in}}%
\pgfpathlineto{\pgfqpoint{0.700822in}{1.858600in}}%
\pgfpathlineto{\pgfqpoint{0.702364in}{1.864190in}}%
\pgfpathlineto{\pgfqpoint{0.708324in}{1.867692in}}%
\pgfpathlineto{\pgfqpoint{0.713573in}{1.866288in}}%
\pgfpathlineto{\pgfqpoint{0.713636in}{1.864996in}}%
\pgfpathlineto{\pgfqpoint{0.701521in}{1.821189in}}%
\pgfpathlineto{\pgfqpoint{0.746029in}{1.809361in}}%
\pgfpathclose%
\pgfusepath{fill}%
\end{pgfscope}%
\begin{pgfscope}%
\pgfpathrectangle{\pgfqpoint{0.100000in}{0.100000in}}{\pgfqpoint{3.007045in}{1.925000in}}%
\pgfusepath{clip}%
\pgfsetbuttcap%
\pgfsetmiterjoin%
\definecolor{currentfill}{rgb}{0.632526,0.797647,0.886874}%
\pgfsetfillcolor{currentfill}%
\pgfsetlinewidth{0.000000pt}%
\definecolor{currentstroke}{rgb}{0.000000,0.000000,0.000000}%
\pgfsetstrokecolor{currentstroke}%
\pgfsetstrokeopacity{0.000000}%
\pgfsetdash{}{0pt}%
\pgfpathmoveto{\pgfqpoint{2.488733in}{1.208153in}}%
\pgfpathlineto{\pgfqpoint{2.485346in}{1.213775in}}%
\pgfpathlineto{\pgfqpoint{2.476467in}{1.219730in}}%
\pgfpathlineto{\pgfqpoint{2.478476in}{1.225076in}}%
\pgfpathlineto{\pgfqpoint{2.467867in}{1.223907in}}%
\pgfpathlineto{\pgfqpoint{2.465782in}{1.225372in}}%
\pgfpathlineto{\pgfqpoint{2.462970in}{1.231712in}}%
\pgfpathlineto{\pgfqpoint{2.461646in}{1.242312in}}%
\pgfpathlineto{\pgfqpoint{2.465474in}{1.242785in}}%
\pgfpathlineto{\pgfqpoint{2.463037in}{1.262554in}}%
\pgfpathlineto{\pgfqpoint{2.480209in}{1.264505in}}%
\pgfpathlineto{\pgfqpoint{2.498296in}{1.267386in}}%
\pgfpathlineto{\pgfqpoint{2.502865in}{1.238879in}}%
\pgfpathlineto{\pgfqpoint{2.507767in}{1.239655in}}%
\pgfpathlineto{\pgfqpoint{2.510163in}{1.234931in}}%
\pgfpathlineto{\pgfqpoint{2.505584in}{1.228741in}}%
\pgfpathlineto{\pgfqpoint{2.507751in}{1.225110in}}%
\pgfpathlineto{\pgfqpoint{2.505896in}{1.220258in}}%
\pgfpathlineto{\pgfqpoint{2.500713in}{1.220454in}}%
\pgfpathlineto{\pgfqpoint{2.488733in}{1.208153in}}%
\pgfpathclose%
\pgfusepath{fill}%
\end{pgfscope}%
\begin{pgfscope}%
\pgfpathrectangle{\pgfqpoint{0.100000in}{0.100000in}}{\pgfqpoint{3.007045in}{1.925000in}}%
\pgfusepath{clip}%
\pgfsetbuttcap%
\pgfsetmiterjoin%
\definecolor{currentfill}{rgb}{0.637447,0.799739,0.888597}%
\pgfsetfillcolor{currentfill}%
\pgfsetlinewidth{0.000000pt}%
\definecolor{currentstroke}{rgb}{0.000000,0.000000,0.000000}%
\pgfsetstrokecolor{currentstroke}%
\pgfsetstrokeopacity{0.000000}%
\pgfsetdash{}{0pt}%
\pgfpathmoveto{\pgfqpoint{2.905539in}{1.431014in}}%
\pgfpathlineto{\pgfqpoint{2.906901in}{1.437874in}}%
\pgfpathlineto{\pgfqpoint{2.903419in}{1.453280in}}%
\pgfpathlineto{\pgfqpoint{2.896973in}{1.475592in}}%
\pgfpathlineto{\pgfqpoint{2.916591in}{1.481492in}}%
\pgfpathlineto{\pgfqpoint{2.917977in}{1.479575in}}%
\pgfpathlineto{\pgfqpoint{2.929528in}{1.490206in}}%
\pgfpathlineto{\pgfqpoint{2.933228in}{1.482168in}}%
\pgfpathlineto{\pgfqpoint{2.944479in}{1.471321in}}%
\pgfpathlineto{\pgfqpoint{2.948844in}{1.463490in}}%
\pgfpathlineto{\pgfqpoint{2.944989in}{1.461097in}}%
\pgfpathlineto{\pgfqpoint{2.945676in}{1.456613in}}%
\pgfpathlineto{\pgfqpoint{2.941792in}{1.452515in}}%
\pgfpathlineto{\pgfqpoint{2.928813in}{1.447884in}}%
\pgfpathlineto{\pgfqpoint{2.928414in}{1.457369in}}%
\pgfpathlineto{\pgfqpoint{2.925276in}{1.463337in}}%
\pgfpathlineto{\pgfqpoint{2.921321in}{1.458455in}}%
\pgfpathlineto{\pgfqpoint{2.921311in}{1.451942in}}%
\pgfpathlineto{\pgfqpoint{2.924101in}{1.446233in}}%
\pgfpathlineto{\pgfqpoint{2.922682in}{1.438922in}}%
\pgfpathlineto{\pgfqpoint{2.919165in}{1.438139in}}%
\pgfpathlineto{\pgfqpoint{2.905539in}{1.431014in}}%
\pgfpathclose%
\pgfusepath{fill}%
\end{pgfscope}%
\begin{pgfscope}%
\pgfpathrectangle{\pgfqpoint{0.100000in}{0.100000in}}{\pgfqpoint{3.007045in}{1.925000in}}%
\pgfusepath{clip}%
\pgfsetbuttcap%
\pgfsetmiterjoin%
\definecolor{currentfill}{rgb}{0.244106,0.557832,0.768889}%
\pgfsetfillcolor{currentfill}%
\pgfsetlinewidth{0.000000pt}%
\definecolor{currentstroke}{rgb}{0.000000,0.000000,0.000000}%
\pgfsetstrokecolor{currentstroke}%
\pgfsetstrokeopacity{0.000000}%
\pgfsetdash{}{0pt}%
\pgfpathmoveto{\pgfqpoint{1.806790in}{1.547902in}}%
\pgfpathlineto{\pgfqpoint{1.801048in}{1.550695in}}%
\pgfpathlineto{\pgfqpoint{1.783351in}{1.550568in}}%
\pgfpathlineto{\pgfqpoint{1.783347in}{1.556314in}}%
\pgfpathlineto{\pgfqpoint{1.766199in}{1.556278in}}%
\pgfpathlineto{\pgfqpoint{1.765864in}{1.580217in}}%
\pgfpathlineto{\pgfqpoint{1.806019in}{1.580465in}}%
\pgfpathlineto{\pgfqpoint{1.802241in}{1.583541in}}%
\pgfpathlineto{\pgfqpoint{1.829469in}{1.583832in}}%
\pgfpathlineto{\pgfqpoint{1.829097in}{1.594472in}}%
\pgfpathlineto{\pgfqpoint{1.827186in}{1.594478in}}%
\pgfpathlineto{\pgfqpoint{1.827033in}{1.605997in}}%
\pgfpathlineto{\pgfqpoint{1.826933in}{1.611742in}}%
\pgfpathlineto{\pgfqpoint{1.844395in}{1.612207in}}%
\pgfpathlineto{\pgfqpoint{1.844441in}{1.606089in}}%
\pgfpathlineto{\pgfqpoint{1.844608in}{1.594634in}}%
\pgfpathlineto{\pgfqpoint{1.840662in}{1.594594in}}%
\pgfpathlineto{\pgfqpoint{1.841127in}{1.578186in}}%
\pgfpathlineto{\pgfqpoint{1.841735in}{1.545791in}}%
\pgfpathlineto{\pgfqpoint{1.838457in}{1.548803in}}%
\pgfpathlineto{\pgfqpoint{1.830410in}{1.548866in}}%
\pgfpathlineto{\pgfqpoint{1.816671in}{1.557157in}}%
\pgfpathlineto{\pgfqpoint{1.813546in}{1.550942in}}%
\pgfpathlineto{\pgfqpoint{1.806790in}{1.547902in}}%
\pgfpathclose%
\pgfusepath{fill}%
\end{pgfscope}%
\begin{pgfscope}%
\pgfpathrectangle{\pgfqpoint{0.100000in}{0.100000in}}{\pgfqpoint{3.007045in}{1.925000in}}%
\pgfusepath{clip}%
\pgfsetbuttcap%
\pgfsetmiterjoin%
\definecolor{currentfill}{rgb}{0.114802,0.424437,0.695194}%
\pgfsetfillcolor{currentfill}%
\pgfsetlinewidth{0.000000pt}%
\definecolor{currentstroke}{rgb}{0.000000,0.000000,0.000000}%
\pgfsetstrokecolor{currentstroke}%
\pgfsetstrokeopacity{0.000000}%
\pgfsetdash{}{0pt}%
\pgfpathmoveto{\pgfqpoint{0.967728in}{1.403968in}}%
\pgfpathlineto{\pgfqpoint{0.967940in}{1.408953in}}%
\pgfpathlineto{\pgfqpoint{0.964652in}{1.410841in}}%
\pgfpathlineto{\pgfqpoint{0.964506in}{1.418586in}}%
\pgfpathlineto{\pgfqpoint{0.968279in}{1.422237in}}%
\pgfpathlineto{\pgfqpoint{0.969777in}{1.427234in}}%
\pgfpathlineto{\pgfqpoint{0.968314in}{1.431784in}}%
\pgfpathlineto{\pgfqpoint{0.955057in}{1.434288in}}%
\pgfpathlineto{\pgfqpoint{0.955530in}{1.447686in}}%
\pgfpathlineto{\pgfqpoint{0.954259in}{1.454699in}}%
\pgfpathlineto{\pgfqpoint{0.947546in}{1.454706in}}%
\pgfpathlineto{\pgfqpoint{0.946788in}{1.460996in}}%
\pgfpathlineto{\pgfqpoint{0.947898in}{1.467983in}}%
\pgfpathlineto{\pgfqpoint{0.953549in}{1.475280in}}%
\pgfpathlineto{\pgfqpoint{1.001890in}{1.466220in}}%
\pgfpathlineto{\pgfqpoint{0.989963in}{1.400079in}}%
\pgfpathlineto{\pgfqpoint{0.967728in}{1.403968in}}%
\pgfpathclose%
\pgfusepath{fill}%
\end{pgfscope}%
\begin{pgfscope}%
\pgfpathrectangle{\pgfqpoint{0.100000in}{0.100000in}}{\pgfqpoint{3.007045in}{1.925000in}}%
\pgfusepath{clip}%
\pgfsetbuttcap%
\pgfsetmiterjoin%
\definecolor{currentfill}{rgb}{0.248166,0.561892,0.770980}%
\pgfsetfillcolor{currentfill}%
\pgfsetlinewidth{0.000000pt}%
\definecolor{currentstroke}{rgb}{0.000000,0.000000,0.000000}%
\pgfsetstrokecolor{currentstroke}%
\pgfsetstrokeopacity{0.000000}%
\pgfsetdash{}{0pt}%
\pgfpathmoveto{\pgfqpoint{1.844374in}{1.020395in}}%
\pgfpathlineto{\pgfqpoint{1.822007in}{1.020648in}}%
\pgfpathlineto{\pgfqpoint{1.820409in}{1.024532in}}%
\pgfpathlineto{\pgfqpoint{1.792358in}{1.025328in}}%
\pgfpathlineto{\pgfqpoint{1.792356in}{1.069644in}}%
\pgfpathlineto{\pgfqpoint{1.792345in}{1.071169in}}%
\pgfpathlineto{\pgfqpoint{1.811752in}{1.070677in}}%
\pgfpathlineto{\pgfqpoint{1.821187in}{1.069825in}}%
\pgfpathlineto{\pgfqpoint{1.820896in}{1.060894in}}%
\pgfpathlineto{\pgfqpoint{1.834274in}{1.060508in}}%
\pgfpathlineto{\pgfqpoint{1.834159in}{1.056680in}}%
\pgfpathlineto{\pgfqpoint{1.843696in}{1.056469in}}%
\pgfpathlineto{\pgfqpoint{1.844766in}{1.050707in}}%
\pgfpathlineto{\pgfqpoint{1.844374in}{1.020395in}}%
\pgfpathclose%
\pgfusepath{fill}%
\end{pgfscope}%
\begin{pgfscope}%
\pgfpathrectangle{\pgfqpoint{0.100000in}{0.100000in}}{\pgfqpoint{3.007045in}{1.925000in}}%
\pgfusepath{clip}%
\pgfsetbuttcap%
\pgfsetmiterjoin%
\definecolor{currentfill}{rgb}{0.460392,0.704744,0.848012}%
\pgfsetfillcolor{currentfill}%
\pgfsetlinewidth{0.000000pt}%
\definecolor{currentstroke}{rgb}{0.000000,0.000000,0.000000}%
\pgfsetstrokecolor{currentstroke}%
\pgfsetstrokeopacity{0.000000}%
\pgfsetdash{}{0pt}%
\pgfpathmoveto{\pgfqpoint{2.597898in}{0.973667in}}%
\pgfpathlineto{\pgfqpoint{2.564813in}{0.967383in}}%
\pgfpathlineto{\pgfqpoint{2.562648in}{0.972115in}}%
\pgfpathlineto{\pgfqpoint{2.555149in}{0.977464in}}%
\pgfpathlineto{\pgfqpoint{2.547659in}{0.980532in}}%
\pgfpathlineto{\pgfqpoint{2.545915in}{0.985956in}}%
\pgfpathlineto{\pgfqpoint{2.549802in}{0.987470in}}%
\pgfpathlineto{\pgfqpoint{2.549576in}{0.995557in}}%
\pgfpathlineto{\pgfqpoint{2.557574in}{1.001121in}}%
\pgfpathlineto{\pgfqpoint{2.566736in}{1.001545in}}%
\pgfpathlineto{\pgfqpoint{2.561819in}{1.036295in}}%
\pgfpathlineto{\pgfqpoint{2.578144in}{1.039001in}}%
\pgfpathlineto{\pgfqpoint{2.578534in}{1.060078in}}%
\pgfpathlineto{\pgfqpoint{2.578887in}{1.072582in}}%
\pgfpathlineto{\pgfqpoint{2.584592in}{1.075050in}}%
\pgfpathlineto{\pgfqpoint{2.591032in}{1.080451in}}%
\pgfpathlineto{\pgfqpoint{2.596091in}{1.081050in}}%
\pgfpathlineto{\pgfqpoint{2.599528in}{1.076992in}}%
\pgfpathlineto{\pgfqpoint{2.604673in}{1.078358in}}%
\pgfpathlineto{\pgfqpoint{2.604048in}{1.043452in}}%
\pgfpathlineto{\pgfqpoint{2.608257in}{1.044199in}}%
\pgfpathlineto{\pgfqpoint{2.610860in}{1.024597in}}%
\pgfpathlineto{\pgfqpoint{2.605433in}{1.023762in}}%
\pgfpathlineto{\pgfqpoint{2.608855in}{1.001837in}}%
\pgfpathlineto{\pgfqpoint{2.611008in}{0.997970in}}%
\pgfpathlineto{\pgfqpoint{2.594940in}{0.995133in}}%
\pgfpathlineto{\pgfqpoint{2.597898in}{0.973667in}}%
\pgfpathclose%
\pgfusepath{fill}%
\end{pgfscope}%
\begin{pgfscope}%
\pgfpathrectangle{\pgfqpoint{0.100000in}{0.100000in}}{\pgfqpoint{3.007045in}{1.925000in}}%
\pgfusepath{clip}%
\pgfsetbuttcap%
\pgfsetmiterjoin%
\definecolor{currentfill}{rgb}{0.485490,0.718524,0.853426}%
\pgfsetfillcolor{currentfill}%
\pgfsetlinewidth{0.000000pt}%
\definecolor{currentstroke}{rgb}{0.000000,0.000000,0.000000}%
\pgfsetstrokecolor{currentstroke}%
\pgfsetstrokeopacity{0.000000}%
\pgfsetdash{}{0pt}%
\pgfpathmoveto{\pgfqpoint{1.166919in}{0.389823in}}%
\pgfpathlineto{\pgfqpoint{1.169914in}{0.398030in}}%
\pgfpathlineto{\pgfqpoint{1.170066in}{0.403152in}}%
\pgfpathlineto{\pgfqpoint{1.177565in}{0.397908in}}%
\pgfpathlineto{\pgfqpoint{1.179080in}{0.390976in}}%
\pgfpathlineto{\pgfqpoint{1.174659in}{0.386184in}}%
\pgfpathlineto{\pgfqpoint{1.166919in}{0.389823in}}%
\pgfpathclose%
\pgfusepath{fill}%
\end{pgfscope}%
\begin{pgfscope}%
\pgfpathrectangle{\pgfqpoint{0.100000in}{0.100000in}}{\pgfqpoint{3.007045in}{1.925000in}}%
\pgfusepath{clip}%
\pgfsetbuttcap%
\pgfsetmiterjoin%
\definecolor{currentfill}{rgb}{0.485490,0.718524,0.853426}%
\pgfsetfillcolor{currentfill}%
\pgfsetlinewidth{0.000000pt}%
\definecolor{currentstroke}{rgb}{0.000000,0.000000,0.000000}%
\pgfsetstrokecolor{currentstroke}%
\pgfsetstrokeopacity{0.000000}%
\pgfsetdash{}{0pt}%
\pgfpathmoveto{\pgfqpoint{1.171506in}{0.365544in}}%
\pgfpathlineto{\pgfqpoint{1.172099in}{0.368432in}}%
\pgfpathlineto{\pgfqpoint{1.181292in}{0.366754in}}%
\pgfpathlineto{\pgfqpoint{1.179030in}{0.362563in}}%
\pgfpathlineto{\pgfqpoint{1.171506in}{0.365544in}}%
\pgfpathclose%
\pgfusepath{fill}%
\end{pgfscope}%
\begin{pgfscope}%
\pgfpathrectangle{\pgfqpoint{0.100000in}{0.100000in}}{\pgfqpoint{3.007045in}{1.925000in}}%
\pgfusepath{clip}%
\pgfsetbuttcap%
\pgfsetmiterjoin%
\definecolor{currentfill}{rgb}{0.485490,0.718524,0.853426}%
\pgfsetfillcolor{currentfill}%
\pgfsetlinewidth{0.000000pt}%
\definecolor{currentstroke}{rgb}{0.000000,0.000000,0.000000}%
\pgfsetstrokecolor{currentstroke}%
\pgfsetstrokeopacity{0.000000}%
\pgfsetdash{}{0pt}%
\pgfpathmoveto{\pgfqpoint{1.189967in}{0.357288in}}%
\pgfpathlineto{\pgfqpoint{1.187711in}{0.361758in}}%
\pgfpathlineto{\pgfqpoint{1.194261in}{0.372053in}}%
\pgfpathlineto{\pgfqpoint{1.187347in}{0.380076in}}%
\pgfpathlineto{\pgfqpoint{1.187339in}{0.386651in}}%
\pgfpathlineto{\pgfqpoint{1.193857in}{0.392205in}}%
\pgfpathlineto{\pgfqpoint{1.198739in}{0.390671in}}%
\pgfpathlineto{\pgfqpoint{1.199911in}{0.385707in}}%
\pgfpathlineto{\pgfqpoint{1.198482in}{0.379133in}}%
\pgfpathlineto{\pgfqpoint{1.210453in}{0.373714in}}%
\pgfpathlineto{\pgfqpoint{1.213763in}{0.360359in}}%
\pgfpathlineto{\pgfqpoint{1.218074in}{0.354132in}}%
\pgfpathlineto{\pgfqpoint{1.215157in}{0.349080in}}%
\pgfpathlineto{\pgfqpoint{1.209803in}{0.348364in}}%
\pgfpathlineto{\pgfqpoint{1.200882in}{0.353223in}}%
\pgfpathlineto{\pgfqpoint{1.189967in}{0.357288in}}%
\pgfpathclose%
\pgfusepath{fill}%
\end{pgfscope}%
\begin{pgfscope}%
\pgfpathrectangle{\pgfqpoint{0.100000in}{0.100000in}}{\pgfqpoint{3.007045in}{1.925000in}}%
\pgfusepath{clip}%
\pgfsetbuttcap%
\pgfsetmiterjoin%
\definecolor{currentfill}{rgb}{0.485490,0.718524,0.853426}%
\pgfsetfillcolor{currentfill}%
\pgfsetlinewidth{0.000000pt}%
\definecolor{currentstroke}{rgb}{0.000000,0.000000,0.000000}%
\pgfsetstrokecolor{currentstroke}%
\pgfsetstrokeopacity{0.000000}%
\pgfsetdash{}{0pt}%
\pgfpathmoveto{\pgfqpoint{1.188373in}{0.411114in}}%
\pgfpathlineto{\pgfqpoint{1.197345in}{0.404388in}}%
\pgfpathlineto{\pgfqpoint{1.194069in}{0.401026in}}%
\pgfpathlineto{\pgfqpoint{1.187346in}{0.401919in}}%
\pgfpathlineto{\pgfqpoint{1.182053in}{0.407040in}}%
\pgfpathlineto{\pgfqpoint{1.178085in}{0.413644in}}%
\pgfpathlineto{\pgfqpoint{1.165729in}{0.423161in}}%
\pgfpathlineto{\pgfqpoint{1.174372in}{0.427023in}}%
\pgfpathlineto{\pgfqpoint{1.188373in}{0.411114in}}%
\pgfpathclose%
\pgfusepath{fill}%
\end{pgfscope}%
\begin{pgfscope}%
\pgfpathrectangle{\pgfqpoint{0.100000in}{0.100000in}}{\pgfqpoint{3.007045in}{1.925000in}}%
\pgfusepath{clip}%
\pgfsetbuttcap%
\pgfsetmiterjoin%
\definecolor{currentfill}{rgb}{0.280892,0.587620,0.785083}%
\pgfsetfillcolor{currentfill}%
\pgfsetlinewidth{0.000000pt}%
\definecolor{currentstroke}{rgb}{0.000000,0.000000,0.000000}%
\pgfsetstrokecolor{currentstroke}%
\pgfsetstrokeopacity{0.000000}%
\pgfsetdash{}{0pt}%
\pgfpathmoveto{\pgfqpoint{1.330767in}{0.811238in}}%
\pgfpathlineto{\pgfqpoint{1.328991in}{0.794521in}}%
\pgfpathlineto{\pgfqpoint{1.303863in}{0.796794in}}%
\pgfpathlineto{\pgfqpoint{1.304431in}{0.802517in}}%
\pgfpathlineto{\pgfqpoint{1.292993in}{0.803535in}}%
\pgfpathlineto{\pgfqpoint{1.294107in}{0.814272in}}%
\pgfpathlineto{\pgfqpoint{1.287264in}{0.814931in}}%
\pgfpathlineto{\pgfqpoint{1.288966in}{0.832202in}}%
\pgfpathlineto{\pgfqpoint{1.283287in}{0.832832in}}%
\pgfpathlineto{\pgfqpoint{1.286524in}{0.867195in}}%
\pgfpathlineto{\pgfqpoint{1.299641in}{0.865888in}}%
\pgfpathlineto{\pgfqpoint{1.300241in}{0.871662in}}%
\pgfpathlineto{\pgfqpoint{1.311575in}{0.870525in}}%
\pgfpathlineto{\pgfqpoint{1.312137in}{0.876255in}}%
\pgfpathlineto{\pgfqpoint{1.317814in}{0.875699in}}%
\pgfpathlineto{\pgfqpoint{1.318390in}{0.881414in}}%
\pgfpathlineto{\pgfqpoint{1.324074in}{0.880866in}}%
\pgfpathlineto{\pgfqpoint{1.324670in}{0.886744in}}%
\pgfpathlineto{\pgfqpoint{1.337618in}{0.885562in}}%
\pgfpathlineto{\pgfqpoint{1.336407in}{0.871962in}}%
\pgfpathlineto{\pgfqpoint{1.364391in}{0.869531in}}%
\pgfpathlineto{\pgfqpoint{1.361996in}{0.840962in}}%
\pgfpathlineto{\pgfqpoint{1.333803in}{0.843365in}}%
\pgfpathlineto{\pgfqpoint{1.330767in}{0.811238in}}%
\pgfpathclose%
\pgfusepath{fill}%
\end{pgfscope}%
\begin{pgfscope}%
\pgfpathrectangle{\pgfqpoint{0.100000in}{0.100000in}}{\pgfqpoint{3.007045in}{1.925000in}}%
\pgfusepath{clip}%
\pgfsetbuttcap%
\pgfsetmiterjoin%
\definecolor{currentfill}{rgb}{0.356555,0.639293,0.814610}%
\pgfsetfillcolor{currentfill}%
\pgfsetlinewidth{0.000000pt}%
\definecolor{currentstroke}{rgb}{0.000000,0.000000,0.000000}%
\pgfsetstrokecolor{currentstroke}%
\pgfsetstrokeopacity{0.000000}%
\pgfsetdash{}{0pt}%
\pgfpathmoveto{\pgfqpoint{2.590755in}{0.837224in}}%
\pgfpathlineto{\pgfqpoint{2.584761in}{0.832670in}}%
\pgfpathlineto{\pgfqpoint{2.577826in}{0.830633in}}%
\pgfpathlineto{\pgfqpoint{2.571176in}{0.834427in}}%
\pgfpathlineto{\pgfqpoint{2.566054in}{0.842571in}}%
\pgfpathlineto{\pgfqpoint{2.559661in}{0.847879in}}%
\pgfpathlineto{\pgfqpoint{2.560497in}{0.852920in}}%
\pgfpathlineto{\pgfqpoint{2.558099in}{0.855702in}}%
\pgfpathlineto{\pgfqpoint{2.557919in}{0.868208in}}%
\pgfpathlineto{\pgfqpoint{2.555420in}{0.871742in}}%
\pgfpathlineto{\pgfqpoint{2.551190in}{0.869149in}}%
\pgfpathlineto{\pgfqpoint{2.541892in}{0.875274in}}%
\pgfpathlineto{\pgfqpoint{2.545276in}{0.888835in}}%
\pgfpathlineto{\pgfqpoint{2.538562in}{0.893175in}}%
\pgfpathlineto{\pgfqpoint{2.545027in}{0.899687in}}%
\pgfpathlineto{\pgfqpoint{2.550816in}{0.897393in}}%
\pgfpathlineto{\pgfqpoint{2.555983in}{0.898307in}}%
\pgfpathlineto{\pgfqpoint{2.555178in}{0.902655in}}%
\pgfpathlineto{\pgfqpoint{2.562238in}{0.907450in}}%
\pgfpathlineto{\pgfqpoint{2.567423in}{0.900852in}}%
\pgfpathlineto{\pgfqpoint{2.571314in}{0.892287in}}%
\pgfpathlineto{\pgfqpoint{2.577854in}{0.890686in}}%
\pgfpathlineto{\pgfqpoint{2.582316in}{0.884878in}}%
\pgfpathlineto{\pgfqpoint{2.580510in}{0.880092in}}%
\pgfpathlineto{\pgfqpoint{2.596705in}{0.871224in}}%
\pgfpathlineto{\pgfqpoint{2.592778in}{0.867711in}}%
\pgfpathlineto{\pgfqpoint{2.597384in}{0.859907in}}%
\pgfpathlineto{\pgfqpoint{2.593582in}{0.855868in}}%
\pgfpathlineto{\pgfqpoint{2.595239in}{0.853039in}}%
\pgfpathlineto{\pgfqpoint{2.590755in}{0.837224in}}%
\pgfpathclose%
\pgfusepath{fill}%
\end{pgfscope}%
\begin{pgfscope}%
\pgfpathrectangle{\pgfqpoint{0.100000in}{0.100000in}}{\pgfqpoint{3.007045in}{1.925000in}}%
\pgfusepath{clip}%
\pgfsetbuttcap%
\pgfsetmiterjoin%
\definecolor{currentfill}{rgb}{0.183206,0.496932,0.737516}%
\pgfsetfillcolor{currentfill}%
\pgfsetlinewidth{0.000000pt}%
\definecolor{currentstroke}{rgb}{0.000000,0.000000,0.000000}%
\pgfsetstrokecolor{currentstroke}%
\pgfsetstrokeopacity{0.000000}%
\pgfsetdash{}{0pt}%
\pgfpathmoveto{\pgfqpoint{1.318468in}{1.629024in}}%
\pgfpathlineto{\pgfqpoint{1.316018in}{1.629285in}}%
\pgfpathlineto{\pgfqpoint{1.313493in}{1.606222in}}%
\pgfpathlineto{\pgfqpoint{1.311526in}{1.606441in}}%
\pgfpathlineto{\pgfqpoint{1.255957in}{1.613050in}}%
\pgfpathlineto{\pgfqpoint{1.258817in}{1.635994in}}%
\pgfpathlineto{\pgfqpoint{1.260561in}{1.635790in}}%
\pgfpathlineto{\pgfqpoint{1.263286in}{1.658626in}}%
\pgfpathlineto{\pgfqpoint{1.264869in}{1.658439in}}%
\pgfpathlineto{\pgfqpoint{1.267668in}{1.681381in}}%
\pgfpathlineto{\pgfqpoint{1.269340in}{1.682158in}}%
\pgfpathlineto{\pgfqpoint{1.272031in}{1.704281in}}%
\pgfpathlineto{\pgfqpoint{1.277730in}{1.703596in}}%
\pgfpathlineto{\pgfqpoint{1.277023in}{1.697772in}}%
\pgfpathlineto{\pgfqpoint{1.282748in}{1.697095in}}%
\pgfpathlineto{\pgfqpoint{1.283442in}{1.702901in}}%
\pgfpathlineto{\pgfqpoint{1.302614in}{1.700688in}}%
\pgfpathlineto{\pgfqpoint{1.306273in}{1.700261in}}%
\pgfpathlineto{\pgfqpoint{1.304728in}{1.686756in}}%
\pgfpathlineto{\pgfqpoint{1.310508in}{1.686101in}}%
\pgfpathlineto{\pgfqpoint{1.310066in}{1.682269in}}%
\pgfpathlineto{\pgfqpoint{1.315821in}{1.681622in}}%
\pgfpathlineto{\pgfqpoint{1.317293in}{1.677601in}}%
\pgfpathlineto{\pgfqpoint{1.336444in}{1.675484in}}%
\pgfpathlineto{\pgfqpoint{1.335063in}{1.663866in}}%
\pgfpathlineto{\pgfqpoint{1.329154in}{1.662581in}}%
\pgfpathlineto{\pgfqpoint{1.327908in}{1.651084in}}%
\pgfpathlineto{\pgfqpoint{1.320048in}{1.651933in}}%
\pgfpathlineto{\pgfqpoint{1.318468in}{1.629024in}}%
\pgfpathclose%
\pgfusepath{fill}%
\end{pgfscope}%
\begin{pgfscope}%
\pgfpathrectangle{\pgfqpoint{0.100000in}{0.100000in}}{\pgfqpoint{3.007045in}{1.925000in}}%
\pgfusepath{clip}%
\pgfsetbuttcap%
\pgfsetmiterjoin%
\definecolor{currentfill}{rgb}{0.203506,0.517232,0.747974}%
\pgfsetfillcolor{currentfill}%
\pgfsetlinewidth{0.000000pt}%
\definecolor{currentstroke}{rgb}{0.000000,0.000000,0.000000}%
\pgfsetstrokecolor{currentstroke}%
\pgfsetstrokeopacity{0.000000}%
\pgfsetdash{}{0pt}%
\pgfpathmoveto{\pgfqpoint{1.694633in}{0.447645in}}%
\pgfpathlineto{\pgfqpoint{1.695745in}{0.469116in}}%
\pgfpathlineto{\pgfqpoint{1.676720in}{0.488443in}}%
\pgfpathlineto{\pgfqpoint{1.693878in}{0.498881in}}%
\pgfpathlineto{\pgfqpoint{1.696363in}{0.505939in}}%
\pgfpathlineto{\pgfqpoint{1.704139in}{0.513794in}}%
\pgfpathlineto{\pgfqpoint{1.710725in}{0.510845in}}%
\pgfpathlineto{\pgfqpoint{1.710733in}{0.503026in}}%
\pgfpathlineto{\pgfqpoint{1.714651in}{0.500322in}}%
\pgfpathlineto{\pgfqpoint{1.717041in}{0.493709in}}%
\pgfpathlineto{\pgfqpoint{1.722932in}{0.488802in}}%
\pgfpathlineto{\pgfqpoint{1.728488in}{0.492886in}}%
\pgfpathlineto{\pgfqpoint{1.734233in}{0.490553in}}%
\pgfpathlineto{\pgfqpoint{1.738650in}{0.493869in}}%
\pgfpathlineto{\pgfqpoint{1.740375in}{0.500409in}}%
\pgfpathlineto{\pgfqpoint{1.745433in}{0.500436in}}%
\pgfpathlineto{\pgfqpoint{1.747709in}{0.509837in}}%
\pgfpathlineto{\pgfqpoint{1.757144in}{0.510663in}}%
\pgfpathlineto{\pgfqpoint{1.759753in}{0.508170in}}%
\pgfpathlineto{\pgfqpoint{1.758777in}{0.502145in}}%
\pgfpathlineto{\pgfqpoint{1.769052in}{0.484340in}}%
\pgfpathlineto{\pgfqpoint{1.765360in}{0.476226in}}%
\pgfpathlineto{\pgfqpoint{1.749852in}{0.462071in}}%
\pgfpathlineto{\pgfqpoint{1.745282in}{0.461023in}}%
\pgfpathlineto{\pgfqpoint{1.724620in}{0.448821in}}%
\pgfpathlineto{\pgfqpoint{1.714126in}{0.445000in}}%
\pgfpathlineto{\pgfqpoint{1.711446in}{0.448095in}}%
\pgfpathlineto{\pgfqpoint{1.700720in}{0.443275in}}%
\pgfpathlineto{\pgfqpoint{1.700851in}{0.451264in}}%
\pgfpathlineto{\pgfqpoint{1.694633in}{0.447645in}}%
\pgfpathclose%
\pgfusepath{fill}%
\end{pgfscope}%
\begin{pgfscope}%
\pgfpathrectangle{\pgfqpoint{0.100000in}{0.100000in}}{\pgfqpoint{3.007045in}{1.925000in}}%
\pgfusepath{clip}%
\pgfsetbuttcap%
\pgfsetmiterjoin%
\definecolor{currentfill}{rgb}{0.657132,0.808105,0.895486}%
\pgfsetfillcolor{currentfill}%
\pgfsetlinewidth{0.000000pt}%
\definecolor{currentstroke}{rgb}{0.000000,0.000000,0.000000}%
\pgfsetstrokecolor{currentstroke}%
\pgfsetstrokeopacity{0.000000}%
\pgfsetdash{}{0pt}%
\pgfpathmoveto{\pgfqpoint{2.609375in}{1.212759in}}%
\pgfpathlineto{\pgfqpoint{2.603332in}{1.202130in}}%
\pgfpathlineto{\pgfqpoint{2.599182in}{1.199561in}}%
\pgfpathlineto{\pgfqpoint{2.596426in}{1.190754in}}%
\pgfpathlineto{\pgfqpoint{2.589115in}{1.195276in}}%
\pgfpathlineto{\pgfqpoint{2.586940in}{1.188956in}}%
\pgfpathlineto{\pgfqpoint{2.569924in}{1.199083in}}%
\pgfpathlineto{\pgfqpoint{2.571223in}{1.207045in}}%
\pgfpathlineto{\pgfqpoint{2.568471in}{1.207938in}}%
\pgfpathlineto{\pgfqpoint{2.569988in}{1.214538in}}%
\pgfpathlineto{\pgfqpoint{2.567300in}{1.215781in}}%
\pgfpathlineto{\pgfqpoint{2.560339in}{1.213292in}}%
\pgfpathlineto{\pgfqpoint{2.554976in}{1.247539in}}%
\pgfpathlineto{\pgfqpoint{2.559174in}{1.248298in}}%
\pgfpathlineto{\pgfqpoint{2.609634in}{1.257414in}}%
\pgfpathlineto{\pgfqpoint{2.613340in}{1.256035in}}%
\pgfpathlineto{\pgfqpoint{2.608755in}{1.249542in}}%
\pgfpathlineto{\pgfqpoint{2.607688in}{1.243264in}}%
\pgfpathlineto{\pgfqpoint{2.598024in}{1.242727in}}%
\pgfpathlineto{\pgfqpoint{2.592721in}{1.234363in}}%
\pgfpathlineto{\pgfqpoint{2.590743in}{1.227736in}}%
\pgfpathlineto{\pgfqpoint{2.585351in}{1.220532in}}%
\pgfpathlineto{\pgfqpoint{2.589864in}{1.218745in}}%
\pgfpathlineto{\pgfqpoint{2.604647in}{1.216037in}}%
\pgfpathlineto{\pgfqpoint{2.609375in}{1.212759in}}%
\pgfpathclose%
\pgfusepath{fill}%
\end{pgfscope}%
\begin{pgfscope}%
\pgfpathrectangle{\pgfqpoint{0.100000in}{0.100000in}}{\pgfqpoint{3.007045in}{1.925000in}}%
\pgfusepath{clip}%
\pgfsetbuttcap%
\pgfsetmiterjoin%
\definecolor{currentfill}{rgb}{0.485490,0.718524,0.853426}%
\pgfsetfillcolor{currentfill}%
\pgfsetlinewidth{0.000000pt}%
\definecolor{currentstroke}{rgb}{0.000000,0.000000,0.000000}%
\pgfsetstrokecolor{currentstroke}%
\pgfsetstrokeopacity{0.000000}%
\pgfsetdash{}{0pt}%
\pgfpathmoveto{\pgfqpoint{2.749311in}{1.495307in}}%
\pgfpathlineto{\pgfqpoint{2.743391in}{1.496243in}}%
\pgfpathlineto{\pgfqpoint{2.742514in}{1.505188in}}%
\pgfpathlineto{\pgfqpoint{2.739690in}{1.510307in}}%
\pgfpathlineto{\pgfqpoint{2.742037in}{1.517084in}}%
\pgfpathlineto{\pgfqpoint{2.739601in}{1.524017in}}%
\pgfpathlineto{\pgfqpoint{2.731549in}{1.525832in}}%
\pgfpathlineto{\pgfqpoint{2.733723in}{1.536267in}}%
\pgfpathlineto{\pgfqpoint{2.721573in}{1.573055in}}%
\pgfpathlineto{\pgfqpoint{2.735921in}{1.578293in}}%
\pgfpathlineto{\pgfqpoint{2.747323in}{1.582366in}}%
\pgfpathlineto{\pgfqpoint{2.750775in}{1.572935in}}%
\pgfpathlineto{\pgfqpoint{2.747693in}{1.569195in}}%
\pgfpathlineto{\pgfqpoint{2.755118in}{1.563156in}}%
\pgfpathlineto{\pgfqpoint{2.763075in}{1.564033in}}%
\pgfpathlineto{\pgfqpoint{2.763396in}{1.560589in}}%
\pgfpathlineto{\pgfqpoint{2.756305in}{1.557859in}}%
\pgfpathlineto{\pgfqpoint{2.774908in}{1.505824in}}%
\pgfpathlineto{\pgfqpoint{2.766923in}{1.496257in}}%
\pgfpathlineto{\pgfqpoint{2.759273in}{1.493756in}}%
\pgfpathlineto{\pgfqpoint{2.749311in}{1.495307in}}%
\pgfpathclose%
\pgfusepath{fill}%
\end{pgfscope}%
\begin{pgfscope}%
\pgfpathrectangle{\pgfqpoint{0.100000in}{0.100000in}}{\pgfqpoint{3.007045in}{1.925000in}}%
\pgfusepath{clip}%
\pgfsetbuttcap%
\pgfsetmiterjoin%
\definecolor{currentfill}{rgb}{0.585882,0.773641,0.875079}%
\pgfsetfillcolor{currentfill}%
\pgfsetlinewidth{0.000000pt}%
\definecolor{currentstroke}{rgb}{0.000000,0.000000,0.000000}%
\pgfsetstrokecolor{currentstroke}%
\pgfsetstrokeopacity{0.000000}%
\pgfsetdash{}{0pt}%
\pgfpathmoveto{\pgfqpoint{2.148381in}{1.158667in}}%
\pgfpathlineto{\pgfqpoint{2.142303in}{1.159619in}}%
\pgfpathlineto{\pgfqpoint{2.105428in}{1.156288in}}%
\pgfpathlineto{\pgfqpoint{2.105221in}{1.159153in}}%
\pgfpathlineto{\pgfqpoint{2.103104in}{1.187835in}}%
\pgfpathlineto{\pgfqpoint{2.123761in}{1.189439in}}%
\pgfpathlineto{\pgfqpoint{2.128488in}{1.192008in}}%
\pgfpathlineto{\pgfqpoint{2.128961in}{1.204856in}}%
\pgfpathlineto{\pgfqpoint{2.149319in}{1.206705in}}%
\pgfpathlineto{\pgfqpoint{2.152293in}{1.171574in}}%
\pgfpathlineto{\pgfqpoint{2.148720in}{1.168849in}}%
\pgfpathlineto{\pgfqpoint{2.150340in}{1.161666in}}%
\pgfpathlineto{\pgfqpoint{2.148381in}{1.158667in}}%
\pgfpathclose%
\pgfusepath{fill}%
\end{pgfscope}%
\begin{pgfscope}%
\pgfpathrectangle{\pgfqpoint{0.100000in}{0.100000in}}{\pgfqpoint{3.007045in}{1.925000in}}%
\pgfusepath{clip}%
\pgfsetbuttcap%
\pgfsetmiterjoin%
\definecolor{currentfill}{rgb}{0.270804,0.580730,0.781146}%
\pgfsetfillcolor{currentfill}%
\pgfsetlinewidth{0.000000pt}%
\definecolor{currentstroke}{rgb}{0.000000,0.000000,0.000000}%
\pgfsetstrokecolor{currentstroke}%
\pgfsetstrokeopacity{0.000000}%
\pgfsetdash{}{0pt}%
\pgfpathmoveto{\pgfqpoint{1.552729in}{1.550914in}}%
\pgfpathlineto{\pgfqpoint{1.527075in}{1.552396in}}%
\pgfpathlineto{\pgfqpoint{1.524722in}{1.557253in}}%
\pgfpathlineto{\pgfqpoint{1.526953in}{1.561321in}}%
\pgfpathlineto{\pgfqpoint{1.524166in}{1.567570in}}%
\pgfpathlineto{\pgfqpoint{1.518319in}{1.571035in}}%
\pgfpathlineto{\pgfqpoint{1.522749in}{1.580062in}}%
\pgfpathlineto{\pgfqpoint{1.526579in}{1.583044in}}%
\pgfpathlineto{\pgfqpoint{1.523621in}{1.593929in}}%
\pgfpathlineto{\pgfqpoint{1.518744in}{1.599173in}}%
\pgfpathlineto{\pgfqpoint{1.554716in}{1.596972in}}%
\pgfpathlineto{\pgfqpoint{1.552729in}{1.550914in}}%
\pgfpathclose%
\pgfusepath{fill}%
\end{pgfscope}%
\begin{pgfscope}%
\pgfpathrectangle{\pgfqpoint{0.100000in}{0.100000in}}{\pgfqpoint{3.007045in}{1.925000in}}%
\pgfusepath{clip}%
\pgfsetbuttcap%
\pgfsetmiterjoin%
\definecolor{currentfill}{rgb}{0.093272,0.396878,0.673664}%
\pgfsetfillcolor{currentfill}%
\pgfsetlinewidth{0.000000pt}%
\definecolor{currentstroke}{rgb}{0.000000,0.000000,0.000000}%
\pgfsetstrokecolor{currentstroke}%
\pgfsetstrokeopacity{0.000000}%
\pgfsetdash{}{0pt}%
\pgfpathmoveto{\pgfqpoint{0.720329in}{0.385044in}}%
\pgfpathlineto{\pgfqpoint{0.727658in}{0.379320in}}%
\pgfpathlineto{\pgfqpoint{0.723393in}{0.373987in}}%
\pgfpathlineto{\pgfqpoint{0.713956in}{0.381377in}}%
\pgfpathlineto{\pgfqpoint{0.717385in}{0.381921in}}%
\pgfpathlineto{\pgfqpoint{0.720189in}{0.383616in}}%
\pgfpathlineto{\pgfqpoint{0.720329in}{0.385044in}}%
\pgfpathclose%
\pgfusepath{fill}%
\end{pgfscope}%
\begin{pgfscope}%
\pgfpathrectangle{\pgfqpoint{0.100000in}{0.100000in}}{\pgfqpoint{3.007045in}{1.925000in}}%
\pgfusepath{clip}%
\pgfsetbuttcap%
\pgfsetmiterjoin%
\definecolor{currentfill}{rgb}{0.031373,0.285675,0.564291}%
\pgfsetfillcolor{currentfill}%
\pgfsetlinewidth{0.000000pt}%
\definecolor{currentstroke}{rgb}{0.000000,0.000000,0.000000}%
\pgfsetstrokecolor{currentstroke}%
\pgfsetstrokeopacity{0.000000}%
\pgfsetdash{}{0pt}%
\pgfpathmoveto{\pgfqpoint{1.483456in}{1.461953in}}%
\pgfpathlineto{\pgfqpoint{1.475452in}{1.459167in}}%
\pgfpathlineto{\pgfqpoint{1.473591in}{1.432656in}}%
\pgfpathlineto{\pgfqpoint{1.431343in}{1.435804in}}%
\pgfpathlineto{\pgfqpoint{1.430442in}{1.441625in}}%
\pgfpathlineto{\pgfqpoint{1.431599in}{1.456389in}}%
\pgfpathlineto{\pgfqpoint{1.437357in}{1.456409in}}%
\pgfpathlineto{\pgfqpoint{1.439374in}{1.475257in}}%
\pgfpathlineto{\pgfqpoint{1.442332in}{1.509325in}}%
\pgfpathlineto{\pgfqpoint{1.452730in}{1.510166in}}%
\pgfpathlineto{\pgfqpoint{1.458166in}{1.512912in}}%
\pgfpathlineto{\pgfqpoint{1.464688in}{1.511398in}}%
\pgfpathlineto{\pgfqpoint{1.473297in}{1.518824in}}%
\pgfpathlineto{\pgfqpoint{1.479161in}{1.518521in}}%
\pgfpathlineto{\pgfqpoint{1.484012in}{1.521937in}}%
\pgfpathlineto{\pgfqpoint{1.482943in}{1.518221in}}%
\pgfpathlineto{\pgfqpoint{1.481358in}{1.495327in}}%
\pgfpathlineto{\pgfqpoint{1.480026in}{1.483812in}}%
\pgfpathlineto{\pgfqpoint{1.485701in}{1.483457in}}%
\pgfpathlineto{\pgfqpoint{1.483456in}{1.461953in}}%
\pgfpathclose%
\pgfusepath{fill}%
\end{pgfscope}%
\begin{pgfscope}%
\pgfpathrectangle{\pgfqpoint{0.100000in}{0.100000in}}{\pgfqpoint{3.007045in}{1.925000in}}%
\pgfusepath{clip}%
\pgfsetbuttcap%
\pgfsetmiterjoin%
\definecolor{currentfill}{rgb}{0.381776,0.656517,0.824452}%
\pgfsetfillcolor{currentfill}%
\pgfsetlinewidth{0.000000pt}%
\definecolor{currentstroke}{rgb}{0.000000,0.000000,0.000000}%
\pgfsetstrokecolor{currentstroke}%
\pgfsetstrokeopacity{0.000000}%
\pgfsetdash{}{0pt}%
\pgfpathmoveto{\pgfqpoint{2.021691in}{1.270010in}}%
\pgfpathlineto{\pgfqpoint{2.021108in}{1.281507in}}%
\pgfpathlineto{\pgfqpoint{1.998647in}{1.280445in}}%
\pgfpathlineto{\pgfqpoint{1.998849in}{1.274698in}}%
\pgfpathlineto{\pgfqpoint{1.981668in}{1.274180in}}%
\pgfpathlineto{\pgfqpoint{1.973511in}{1.274071in}}%
\pgfpathlineto{\pgfqpoint{1.971270in}{1.279448in}}%
\pgfpathlineto{\pgfqpoint{1.964916in}{1.284730in}}%
\pgfpathlineto{\pgfqpoint{1.967742in}{1.296340in}}%
\pgfpathlineto{\pgfqpoint{1.980445in}{1.299472in}}%
\pgfpathlineto{\pgfqpoint{1.980128in}{1.308936in}}%
\pgfpathlineto{\pgfqpoint{1.974408in}{1.308715in}}%
\pgfpathlineto{\pgfqpoint{1.973965in}{1.320197in}}%
\pgfpathlineto{\pgfqpoint{1.996691in}{1.320570in}}%
\pgfpathlineto{\pgfqpoint{2.000310in}{1.318595in}}%
\pgfpathlineto{\pgfqpoint{2.006146in}{1.322530in}}%
\pgfpathlineto{\pgfqpoint{2.009626in}{1.309620in}}%
\pgfpathlineto{\pgfqpoint{2.025542in}{1.310434in}}%
\pgfpathlineto{\pgfqpoint{2.027069in}{1.287453in}}%
\pgfpathlineto{\pgfqpoint{2.037908in}{1.288046in}}%
\pgfpathlineto{\pgfqpoint{2.038925in}{1.270934in}}%
\pgfpathlineto{\pgfqpoint{2.021691in}{1.270010in}}%
\pgfpathclose%
\pgfusepath{fill}%
\end{pgfscope}%
\begin{pgfscope}%
\pgfpathrectangle{\pgfqpoint{0.100000in}{0.100000in}}{\pgfqpoint{3.007045in}{1.925000in}}%
\pgfusepath{clip}%
\pgfsetbuttcap%
\pgfsetmiterjoin%
\definecolor{currentfill}{rgb}{0.460392,0.704744,0.848012}%
\pgfsetfillcolor{currentfill}%
\pgfsetlinewidth{0.000000pt}%
\definecolor{currentstroke}{rgb}{0.000000,0.000000,0.000000}%
\pgfsetstrokecolor{currentstroke}%
\pgfsetstrokeopacity{0.000000}%
\pgfsetdash{}{0pt}%
\pgfpathmoveto{\pgfqpoint{2.273936in}{0.628314in}}%
\pgfpathlineto{\pgfqpoint{2.282576in}{0.629115in}}%
\pgfpathlineto{\pgfqpoint{2.284322in}{0.609983in}}%
\pgfpathlineto{\pgfqpoint{2.294588in}{0.610608in}}%
\pgfpathlineto{\pgfqpoint{2.291737in}{0.606164in}}%
\pgfpathlineto{\pgfqpoint{2.295880in}{0.596792in}}%
\pgfpathlineto{\pgfqpoint{2.288959in}{0.589362in}}%
\pgfpathlineto{\pgfqpoint{2.289539in}{0.581381in}}%
\pgfpathlineto{\pgfqpoint{2.275662in}{0.585133in}}%
\pgfpathlineto{\pgfqpoint{2.265886in}{0.586394in}}%
\pgfpathlineto{\pgfqpoint{2.252297in}{0.586303in}}%
\pgfpathlineto{\pgfqpoint{2.238850in}{0.583884in}}%
\pgfpathlineto{\pgfqpoint{2.221344in}{0.578728in}}%
\pgfpathlineto{\pgfqpoint{2.212863in}{0.577456in}}%
\pgfpathlineto{\pgfqpoint{2.202262in}{0.574181in}}%
\pgfpathlineto{\pgfqpoint{2.206102in}{0.581924in}}%
\pgfpathlineto{\pgfqpoint{2.210055in}{0.585128in}}%
\pgfpathlineto{\pgfqpoint{2.205014in}{0.590892in}}%
\pgfpathlineto{\pgfqpoint{2.207470in}{0.596991in}}%
\pgfpathlineto{\pgfqpoint{2.206456in}{0.600906in}}%
\pgfpathlineto{\pgfqpoint{2.199079in}{0.604683in}}%
\pgfpathlineto{\pgfqpoint{2.192424in}{0.612482in}}%
\pgfpathlineto{\pgfqpoint{2.193761in}{0.621372in}}%
\pgfpathlineto{\pgfqpoint{2.245494in}{0.625658in}}%
\pgfpathlineto{\pgfqpoint{2.273936in}{0.628314in}}%
\pgfpathclose%
\pgfusepath{fill}%
\end{pgfscope}%
\begin{pgfscope}%
\pgfpathrectangle{\pgfqpoint{0.100000in}{0.100000in}}{\pgfqpoint{3.007045in}{1.925000in}}%
\pgfusepath{clip}%
\pgfsetbuttcap%
\pgfsetmiterjoin%
\definecolor{currentfill}{rgb}{0.351511,0.635848,0.812641}%
\pgfsetfillcolor{currentfill}%
\pgfsetlinewidth{0.000000pt}%
\definecolor{currentstroke}{rgb}{0.000000,0.000000,0.000000}%
\pgfsetstrokecolor{currentstroke}%
\pgfsetstrokeopacity{0.000000}%
\pgfsetdash{}{0pt}%
\pgfpathmoveto{\pgfqpoint{1.539001in}{1.204970in}}%
\pgfpathlineto{\pgfqpoint{1.537445in}{1.176396in}}%
\pgfpathlineto{\pgfqpoint{1.510385in}{1.177845in}}%
\pgfpathlineto{\pgfqpoint{1.509436in}{1.177873in}}%
\pgfpathlineto{\pgfqpoint{1.511195in}{1.206454in}}%
\pgfpathlineto{\pgfqpoint{1.539001in}{1.204970in}}%
\pgfpathclose%
\pgfusepath{fill}%
\end{pgfscope}%
\begin{pgfscope}%
\pgfpathrectangle{\pgfqpoint{0.100000in}{0.100000in}}{\pgfqpoint{3.007045in}{1.925000in}}%
\pgfusepath{clip}%
\pgfsetbuttcap%
\pgfsetmiterjoin%
\definecolor{currentfill}{rgb}{0.523137,0.739193,0.861546}%
\pgfsetfillcolor{currentfill}%
\pgfsetlinewidth{0.000000pt}%
\definecolor{currentstroke}{rgb}{0.000000,0.000000,0.000000}%
\pgfsetstrokecolor{currentstroke}%
\pgfsetstrokeopacity{0.000000}%
\pgfsetdash{}{0pt}%
\pgfpathmoveto{\pgfqpoint{1.975423in}{0.932334in}}%
\pgfpathlineto{\pgfqpoint{1.969720in}{0.923885in}}%
\pgfpathlineto{\pgfqpoint{1.966993in}{0.915480in}}%
\pgfpathlineto{\pgfqpoint{1.967250in}{0.907985in}}%
\pgfpathlineto{\pgfqpoint{1.944456in}{0.907639in}}%
\pgfpathlineto{\pgfqpoint{1.944264in}{0.919360in}}%
\pgfpathlineto{\pgfqpoint{1.939696in}{0.925677in}}%
\pgfpathlineto{\pgfqpoint{1.940723in}{0.927923in}}%
\pgfpathlineto{\pgfqpoint{1.935018in}{0.933586in}}%
\pgfpathlineto{\pgfqpoint{1.927752in}{0.935378in}}%
\pgfpathlineto{\pgfqpoint{1.929258in}{0.940159in}}%
\pgfpathlineto{\pgfqpoint{1.921677in}{0.946416in}}%
\pgfpathlineto{\pgfqpoint{1.923472in}{0.954935in}}%
\pgfpathlineto{\pgfqpoint{1.954901in}{0.955275in}}%
\pgfpathlineto{\pgfqpoint{1.954694in}{0.961038in}}%
\pgfpathlineto{\pgfqpoint{1.960613in}{0.961138in}}%
\pgfpathlineto{\pgfqpoint{1.960431in}{0.971833in}}%
\pgfpathlineto{\pgfqpoint{1.962703in}{0.971887in}}%
\pgfpathlineto{\pgfqpoint{1.969967in}{0.964093in}}%
\pgfpathlineto{\pgfqpoint{1.971212in}{0.955492in}}%
\pgfpathlineto{\pgfqpoint{1.966442in}{0.954472in}}%
\pgfpathlineto{\pgfqpoint{1.966949in}{0.932026in}}%
\pgfpathlineto{\pgfqpoint{1.975423in}{0.932334in}}%
\pgfpathclose%
\pgfusepath{fill}%
\end{pgfscope}%
\begin{pgfscope}%
\pgfpathrectangle{\pgfqpoint{0.100000in}{0.100000in}}{\pgfqpoint{3.007045in}{1.925000in}}%
\pgfusepath{clip}%
\pgfsetbuttcap%
\pgfsetmiterjoin%
\definecolor{currentfill}{rgb}{0.504314,0.728858,0.857486}%
\pgfsetfillcolor{currentfill}%
\pgfsetlinewidth{0.000000pt}%
\definecolor{currentstroke}{rgb}{0.000000,0.000000,0.000000}%
\pgfsetstrokecolor{currentstroke}%
\pgfsetstrokeopacity{0.000000}%
\pgfsetdash{}{0pt}%
\pgfpathmoveto{\pgfqpoint{1.944978in}{1.116719in}}%
\pgfpathlineto{\pgfqpoint{1.945763in}{1.089375in}}%
\pgfpathlineto{\pgfqpoint{1.917287in}{1.088655in}}%
\pgfpathlineto{\pgfqpoint{1.911673in}{1.091410in}}%
\pgfpathlineto{\pgfqpoint{1.906201in}{1.091563in}}%
\pgfpathlineto{\pgfqpoint{1.906297in}{1.097100in}}%
\pgfpathlineto{\pgfqpoint{1.894852in}{1.097146in}}%
\pgfpathlineto{\pgfqpoint{1.883415in}{1.107800in}}%
\pgfpathlineto{\pgfqpoint{1.883385in}{1.113543in}}%
\pgfpathlineto{\pgfqpoint{1.872048in}{1.113887in}}%
\pgfpathlineto{\pgfqpoint{1.872212in}{1.129384in}}%
\pgfpathlineto{\pgfqpoint{1.872251in}{1.132295in}}%
\pgfpathlineto{\pgfqpoint{1.877944in}{1.138491in}}%
\pgfpathlineto{\pgfqpoint{1.877465in}{1.141048in}}%
\pgfpathlineto{\pgfqpoint{1.882047in}{1.149212in}}%
\pgfpathlineto{\pgfqpoint{1.884869in}{1.157037in}}%
\pgfpathlineto{\pgfqpoint{1.889132in}{1.155695in}}%
\pgfpathlineto{\pgfqpoint{1.903384in}{1.151278in}}%
\pgfpathlineto{\pgfqpoint{1.909292in}{1.151308in}}%
\pgfpathlineto{\pgfqpoint{1.920007in}{1.151193in}}%
\pgfpathlineto{\pgfqpoint{1.920082in}{1.139613in}}%
\pgfpathlineto{\pgfqpoint{1.944438in}{1.140114in}}%
\pgfpathlineto{\pgfqpoint{1.944978in}{1.116719in}}%
\pgfpathclose%
\pgfusepath{fill}%
\end{pgfscope}%
\begin{pgfscope}%
\pgfpathrectangle{\pgfqpoint{0.100000in}{0.100000in}}{\pgfqpoint{3.007045in}{1.925000in}}%
\pgfusepath{clip}%
\pgfsetbuttcap%
\pgfsetmiterjoin%
\definecolor{currentfill}{rgb}{0.150727,0.464452,0.720784}%
\pgfsetfillcolor{currentfill}%
\pgfsetlinewidth{0.000000pt}%
\definecolor{currentstroke}{rgb}{0.000000,0.000000,0.000000}%
\pgfsetstrokecolor{currentstroke}%
\pgfsetstrokeopacity{0.000000}%
\pgfsetdash{}{0pt}%
\pgfpathmoveto{\pgfqpoint{2.505019in}{0.681519in}}%
\pgfpathlineto{\pgfqpoint{2.503207in}{0.687158in}}%
\pgfpathlineto{\pgfqpoint{2.498935in}{0.687218in}}%
\pgfpathlineto{\pgfqpoint{2.496078in}{0.707079in}}%
\pgfpathlineto{\pgfqpoint{2.499221in}{0.711875in}}%
\pgfpathlineto{\pgfqpoint{2.496462in}{0.714530in}}%
\pgfpathlineto{\pgfqpoint{2.489650in}{0.715531in}}%
\pgfpathlineto{\pgfqpoint{2.489704in}{0.732476in}}%
\pgfpathlineto{\pgfqpoint{2.485433in}{0.742115in}}%
\pgfpathlineto{\pgfqpoint{2.509344in}{0.740112in}}%
\pgfpathlineto{\pgfqpoint{2.512057in}{0.734800in}}%
\pgfpathlineto{\pgfqpoint{2.513924in}{0.731882in}}%
\pgfpathlineto{\pgfqpoint{2.520931in}{0.729238in}}%
\pgfpathlineto{\pgfqpoint{2.526011in}{0.723942in}}%
\pgfpathlineto{\pgfqpoint{2.534070in}{0.723697in}}%
\pgfpathlineto{\pgfqpoint{2.537534in}{0.718203in}}%
\pgfpathlineto{\pgfqpoint{2.544812in}{0.713097in}}%
\pgfpathlineto{\pgfqpoint{2.549716in}{0.711373in}}%
\pgfpathlineto{\pgfqpoint{2.553497in}{0.705848in}}%
\pgfpathlineto{\pgfqpoint{2.550608in}{0.701257in}}%
\pgfpathlineto{\pgfqpoint{2.545964in}{0.706161in}}%
\pgfpathlineto{\pgfqpoint{2.536542in}{0.704088in}}%
\pgfpathlineto{\pgfqpoint{2.532623in}{0.707455in}}%
\pgfpathlineto{\pgfqpoint{2.529648in}{0.698799in}}%
\pgfpathlineto{\pgfqpoint{2.524353in}{0.695341in}}%
\pgfpathlineto{\pgfqpoint{2.528105in}{0.691406in}}%
\pgfpathlineto{\pgfqpoint{2.523188in}{0.681053in}}%
\pgfpathlineto{\pgfqpoint{2.518569in}{0.682840in}}%
\pgfpathlineto{\pgfqpoint{2.511887in}{0.680666in}}%
\pgfpathlineto{\pgfqpoint{2.505019in}{0.681519in}}%
\pgfpathclose%
\pgfusepath{fill}%
\end{pgfscope}%
\begin{pgfscope}%
\pgfpathrectangle{\pgfqpoint{0.100000in}{0.100000in}}{\pgfqpoint{3.007045in}{1.925000in}}%
\pgfusepath{clip}%
\pgfsetbuttcap%
\pgfsetmiterjoin%
\definecolor{currentfill}{rgb}{0.270804,0.580730,0.781146}%
\pgfsetfillcolor{currentfill}%
\pgfsetlinewidth{0.000000pt}%
\definecolor{currentstroke}{rgb}{0.000000,0.000000,0.000000}%
\pgfsetstrokecolor{currentstroke}%
\pgfsetstrokeopacity{0.000000}%
\pgfsetdash{}{0pt}%
\pgfpathmoveto{\pgfqpoint{1.704306in}{1.528001in}}%
\pgfpathlineto{\pgfqpoint{1.703956in}{1.505078in}}%
\pgfpathlineto{\pgfqpoint{1.703851in}{1.499277in}}%
\pgfpathlineto{\pgfqpoint{1.671967in}{1.499872in}}%
\pgfpathlineto{\pgfqpoint{1.637786in}{1.500890in}}%
\pgfpathlineto{\pgfqpoint{1.637825in}{1.506641in}}%
\pgfpathlineto{\pgfqpoint{1.632109in}{1.506861in}}%
\pgfpathlineto{\pgfqpoint{1.633052in}{1.541386in}}%
\pgfpathlineto{\pgfqpoint{1.668319in}{1.540277in}}%
\pgfpathlineto{\pgfqpoint{1.684336in}{1.539811in}}%
\pgfpathlineto{\pgfqpoint{1.684152in}{1.528327in}}%
\pgfpathlineto{\pgfqpoint{1.704306in}{1.528001in}}%
\pgfpathclose%
\pgfusepath{fill}%
\end{pgfscope}%
\begin{pgfscope}%
\pgfpathrectangle{\pgfqpoint{0.100000in}{0.100000in}}{\pgfqpoint{3.007045in}{1.925000in}}%
\pgfusepath{clip}%
\pgfsetbuttcap%
\pgfsetmiterjoin%
\definecolor{currentfill}{rgb}{0.491765,0.721968,0.854779}%
\pgfsetfillcolor{currentfill}%
\pgfsetlinewidth{0.000000pt}%
\definecolor{currentstroke}{rgb}{0.000000,0.000000,0.000000}%
\pgfsetstrokecolor{currentstroke}%
\pgfsetstrokeopacity{0.000000}%
\pgfsetdash{}{0pt}%
\pgfpathmoveto{\pgfqpoint{1.616016in}{0.651438in}}%
\pgfpathlineto{\pgfqpoint{1.586901in}{0.635387in}}%
\pgfpathlineto{\pgfqpoint{1.575122in}{0.656245in}}%
\pgfpathlineto{\pgfqpoint{1.565148in}{0.655103in}}%
\pgfpathlineto{\pgfqpoint{1.551784in}{0.680740in}}%
\pgfpathlineto{\pgfqpoint{1.577602in}{0.694365in}}%
\pgfpathlineto{\pgfqpoint{1.578119in}{0.708542in}}%
\pgfpathlineto{\pgfqpoint{1.600832in}{0.707630in}}%
\pgfpathlineto{\pgfqpoint{1.607101in}{0.688968in}}%
\pgfpathlineto{\pgfqpoint{1.611326in}{0.679134in}}%
\pgfpathlineto{\pgfqpoint{1.603281in}{0.674775in}}%
\pgfpathlineto{\pgfqpoint{1.616016in}{0.651438in}}%
\pgfpathclose%
\pgfusepath{fill}%
\end{pgfscope}%
\begin{pgfscope}%
\pgfpathrectangle{\pgfqpoint{0.100000in}{0.100000in}}{\pgfqpoint{3.007045in}{1.925000in}}%
\pgfusepath{clip}%
\pgfsetbuttcap%
\pgfsetmiterjoin%
\definecolor{currentfill}{rgb}{0.567059,0.763306,0.871019}%
\pgfsetfillcolor{currentfill}%
\pgfsetlinewidth{0.000000pt}%
\definecolor{currentstroke}{rgb}{0.000000,0.000000,0.000000}%
\pgfsetstrokecolor{currentstroke}%
\pgfsetstrokeopacity{0.000000}%
\pgfsetdash{}{0pt}%
\pgfpathmoveto{\pgfqpoint{2.609375in}{1.212759in}}%
\pgfpathlineto{\pgfqpoint{2.614276in}{1.211702in}}%
\pgfpathlineto{\pgfqpoint{2.616014in}{1.216485in}}%
\pgfpathlineto{\pgfqpoint{2.621536in}{1.211734in}}%
\pgfpathlineto{\pgfqpoint{2.622986in}{1.207692in}}%
\pgfpathlineto{\pgfqpoint{2.620005in}{1.199545in}}%
\pgfpathlineto{\pgfqpoint{2.625996in}{1.195904in}}%
\pgfpathlineto{\pgfqpoint{2.624882in}{1.186857in}}%
\pgfpathlineto{\pgfqpoint{2.621002in}{1.178691in}}%
\pgfpathlineto{\pgfqpoint{2.619686in}{1.171111in}}%
\pgfpathlineto{\pgfqpoint{2.613166in}{1.160149in}}%
\pgfpathlineto{\pgfqpoint{2.609006in}{1.155780in}}%
\pgfpathlineto{\pgfqpoint{2.581715in}{1.169068in}}%
\pgfpathlineto{\pgfqpoint{2.577995in}{1.164006in}}%
\pgfpathlineto{\pgfqpoint{2.569128in}{1.165501in}}%
\pgfpathlineto{\pgfqpoint{2.565022in}{1.171056in}}%
\pgfpathlineto{\pgfqpoint{2.558875in}{1.172717in}}%
\pgfpathlineto{\pgfqpoint{2.563877in}{1.187087in}}%
\pgfpathlineto{\pgfqpoint{2.562770in}{1.192250in}}%
\pgfpathlineto{\pgfqpoint{2.569924in}{1.199083in}}%
\pgfpathlineto{\pgfqpoint{2.586940in}{1.188956in}}%
\pgfpathlineto{\pgfqpoint{2.589115in}{1.195276in}}%
\pgfpathlineto{\pgfqpoint{2.596426in}{1.190754in}}%
\pgfpathlineto{\pgfqpoint{2.599182in}{1.199561in}}%
\pgfpathlineto{\pgfqpoint{2.603332in}{1.202130in}}%
\pgfpathlineto{\pgfqpoint{2.609375in}{1.212759in}}%
\pgfpathclose%
\pgfusepath{fill}%
\end{pgfscope}%
\begin{pgfscope}%
\pgfpathrectangle{\pgfqpoint{0.100000in}{0.100000in}}{\pgfqpoint{3.007045in}{1.925000in}}%
\pgfusepath{clip}%
\pgfsetbuttcap%
\pgfsetmiterjoin%
\definecolor{currentfill}{rgb}{0.285936,0.591065,0.787051}%
\pgfsetfillcolor{currentfill}%
\pgfsetlinewidth{0.000000pt}%
\definecolor{currentstroke}{rgb}{0.000000,0.000000,0.000000}%
\pgfsetstrokecolor{currentstroke}%
\pgfsetstrokeopacity{0.000000}%
\pgfsetdash{}{0pt}%
\pgfpathmoveto{\pgfqpoint{2.028875in}{0.983778in}}%
\pgfpathlineto{\pgfqpoint{2.025516in}{0.974832in}}%
\pgfpathlineto{\pgfqpoint{1.977561in}{0.972537in}}%
\pgfpathlineto{\pgfqpoint{1.977205in}{0.994031in}}%
\pgfpathlineto{\pgfqpoint{1.971458in}{0.993921in}}%
\pgfpathlineto{\pgfqpoint{1.971148in}{1.011094in}}%
\pgfpathlineto{\pgfqpoint{1.981594in}{1.012166in}}%
\pgfpathlineto{\pgfqpoint{1.984516in}{1.009519in}}%
\pgfpathlineto{\pgfqpoint{1.994242in}{1.009746in}}%
\pgfpathlineto{\pgfqpoint{1.996111in}{1.017492in}}%
\pgfpathlineto{\pgfqpoint{1.995876in}{1.024384in}}%
\pgfpathlineto{\pgfqpoint{2.005554in}{1.024952in}}%
\pgfpathlineto{\pgfqpoint{2.005455in}{1.027915in}}%
\pgfpathlineto{\pgfqpoint{2.022844in}{1.028625in}}%
\pgfpathlineto{\pgfqpoint{2.023343in}{1.013571in}}%
\pgfpathlineto{\pgfqpoint{2.029271in}{1.013848in}}%
\pgfpathlineto{\pgfqpoint{2.022047in}{1.002653in}}%
\pgfpathlineto{\pgfqpoint{2.028045in}{0.989826in}}%
\pgfpathlineto{\pgfqpoint{2.028875in}{0.983778in}}%
\pgfpathclose%
\pgfusepath{fill}%
\end{pgfscope}%
\begin{pgfscope}%
\pgfpathrectangle{\pgfqpoint{0.100000in}{0.100000in}}{\pgfqpoint{3.007045in}{1.925000in}}%
\pgfusepath{clip}%
\pgfsetbuttcap%
\pgfsetmiterjoin%
\definecolor{currentfill}{rgb}{0.183206,0.496932,0.737516}%
\pgfsetfillcolor{currentfill}%
\pgfsetlinewidth{0.000000pt}%
\definecolor{currentstroke}{rgb}{0.000000,0.000000,0.000000}%
\pgfsetstrokecolor{currentstroke}%
\pgfsetstrokeopacity{0.000000}%
\pgfsetdash{}{0pt}%
\pgfpathmoveto{\pgfqpoint{1.331685in}{1.819464in}}%
\pgfpathlineto{\pgfqpoint{1.375731in}{1.814864in}}%
\pgfpathlineto{\pgfqpoint{1.371755in}{1.774248in}}%
\pgfpathlineto{\pgfqpoint{1.346116in}{1.776879in}}%
\pgfpathlineto{\pgfqpoint{1.346744in}{1.782663in}}%
\pgfpathlineto{\pgfqpoint{1.341017in}{1.783290in}}%
\pgfpathlineto{\pgfqpoint{1.341649in}{1.789040in}}%
\pgfpathlineto{\pgfqpoint{1.332164in}{1.790080in}}%
\pgfpathlineto{\pgfqpoint{1.329938in}{1.796177in}}%
\pgfpathlineto{\pgfqpoint{1.331869in}{1.813477in}}%
\pgfpathlineto{\pgfqpoint{1.331685in}{1.819464in}}%
\pgfpathclose%
\pgfusepath{fill}%
\end{pgfscope}%
\begin{pgfscope}%
\pgfpathrectangle{\pgfqpoint{0.100000in}{0.100000in}}{\pgfqpoint{3.007045in}{1.925000in}}%
\pgfusepath{clip}%
\pgfsetbuttcap%
\pgfsetmiterjoin%
\definecolor{currentfill}{rgb}{0.401953,0.670296,0.832326}%
\pgfsetfillcolor{currentfill}%
\pgfsetlinewidth{0.000000pt}%
\definecolor{currentstroke}{rgb}{0.000000,0.000000,0.000000}%
\pgfsetstrokecolor{currentstroke}%
\pgfsetstrokeopacity{0.000000}%
\pgfsetdash{}{0pt}%
\pgfpathmoveto{\pgfqpoint{1.782023in}{1.648669in}}%
\pgfpathlineto{\pgfqpoint{1.764918in}{1.648716in}}%
\pgfpathlineto{\pgfqpoint{1.764876in}{1.642920in}}%
\pgfpathlineto{\pgfqpoint{1.718799in}{1.643280in}}%
\pgfpathlineto{\pgfqpoint{1.719048in}{1.660538in}}%
\pgfpathlineto{\pgfqpoint{1.718261in}{1.672097in}}%
\pgfpathlineto{\pgfqpoint{1.724037in}{1.672065in}}%
\pgfpathlineto{\pgfqpoint{1.724288in}{1.695202in}}%
\pgfpathlineto{\pgfqpoint{1.747387in}{1.695072in}}%
\pgfpathlineto{\pgfqpoint{1.746352in}{1.706621in}}%
\pgfpathlineto{\pgfqpoint{1.746314in}{1.723984in}}%
\pgfpathlineto{\pgfqpoint{1.745890in}{1.729762in}}%
\pgfpathlineto{\pgfqpoint{1.745680in}{1.764366in}}%
\pgfpathlineto{\pgfqpoint{1.762948in}{1.764365in}}%
\pgfpathlineto{\pgfqpoint{1.762966in}{1.752742in}}%
\pgfpathlineto{\pgfqpoint{1.797671in}{1.752902in}}%
\pgfpathlineto{\pgfqpoint{1.798346in}{1.718105in}}%
\pgfpathlineto{\pgfqpoint{1.798628in}{1.689186in}}%
\pgfpathlineto{\pgfqpoint{1.787190in}{1.689017in}}%
\pgfpathlineto{\pgfqpoint{1.787806in}{1.648648in}}%
\pgfpathlineto{\pgfqpoint{1.782023in}{1.648669in}}%
\pgfpathclose%
\pgfusepath{fill}%
\end{pgfscope}%
\begin{pgfscope}%
\pgfpathrectangle{\pgfqpoint{0.100000in}{0.100000in}}{\pgfqpoint{3.007045in}{1.925000in}}%
\pgfusepath{clip}%
\pgfsetbuttcap%
\pgfsetmiterjoin%
\definecolor{currentfill}{rgb}{0.223806,0.537532,0.758431}%
\pgfsetfillcolor{currentfill}%
\pgfsetlinewidth{0.000000pt}%
\definecolor{currentstroke}{rgb}{0.000000,0.000000,0.000000}%
\pgfsetstrokecolor{currentstroke}%
\pgfsetstrokeopacity{0.000000}%
\pgfsetdash{}{0pt}%
\pgfpathmoveto{\pgfqpoint{1.910648in}{0.511564in}}%
\pgfpathlineto{\pgfqpoint{1.909910in}{0.542133in}}%
\pgfpathlineto{\pgfqpoint{1.908086in}{0.571754in}}%
\pgfpathlineto{\pgfqpoint{1.910006in}{0.578352in}}%
\pgfpathlineto{\pgfqpoint{1.909431in}{0.599271in}}%
\pgfpathlineto{\pgfqpoint{1.915657in}{0.600227in}}%
\pgfpathlineto{\pgfqpoint{1.921122in}{0.606604in}}%
\pgfpathlineto{\pgfqpoint{1.927440in}{0.602431in}}%
\pgfpathlineto{\pgfqpoint{1.930070in}{0.596671in}}%
\pgfpathlineto{\pgfqpoint{1.954185in}{0.597454in}}%
\pgfpathlineto{\pgfqpoint{1.959192in}{0.587885in}}%
\pgfpathlineto{\pgfqpoint{1.958167in}{0.578691in}}%
\pgfpathlineto{\pgfqpoint{1.959526in}{0.575009in}}%
\pgfpathlineto{\pgfqpoint{1.965180in}{0.570849in}}%
\pgfpathlineto{\pgfqpoint{1.966602in}{0.562051in}}%
\pgfpathlineto{\pgfqpoint{1.970123in}{0.557651in}}%
\pgfpathlineto{\pgfqpoint{1.975186in}{0.557014in}}%
\pgfpathlineto{\pgfqpoint{1.976175in}{0.548583in}}%
\pgfpathlineto{\pgfqpoint{1.990185in}{0.543999in}}%
\pgfpathlineto{\pgfqpoint{1.988235in}{0.542241in}}%
\pgfpathlineto{\pgfqpoint{1.992985in}{0.532196in}}%
\pgfpathlineto{\pgfqpoint{1.997286in}{0.531031in}}%
\pgfpathlineto{\pgfqpoint{1.998738in}{0.526353in}}%
\pgfpathlineto{\pgfqpoint{1.993498in}{0.524333in}}%
\pgfpathlineto{\pgfqpoint{1.981476in}{0.527516in}}%
\pgfpathlineto{\pgfqpoint{1.982084in}{0.530926in}}%
\pgfpathlineto{\pgfqpoint{1.976086in}{0.538805in}}%
\pgfpathlineto{\pgfqpoint{1.967294in}{0.537756in}}%
\pgfpathlineto{\pgfqpoint{1.962543in}{0.530961in}}%
\pgfpathlineto{\pgfqpoint{1.953040in}{0.522060in}}%
\pgfpathlineto{\pgfqpoint{1.952704in}{0.529849in}}%
\pgfpathlineto{\pgfqpoint{1.943821in}{0.525996in}}%
\pgfpathlineto{\pgfqpoint{1.937749in}{0.521438in}}%
\pgfpathlineto{\pgfqpoint{1.939627in}{0.515520in}}%
\pgfpathlineto{\pgfqpoint{1.946540in}{0.514748in}}%
\pgfpathlineto{\pgfqpoint{1.955604in}{0.516241in}}%
\pgfpathlineto{\pgfqpoint{1.963363in}{0.511582in}}%
\pgfpathlineto{\pgfqpoint{1.957159in}{0.505997in}}%
\pgfpathlineto{\pgfqpoint{1.943062in}{0.512593in}}%
\pgfpathlineto{\pgfqpoint{1.937453in}{0.512510in}}%
\pgfpathlineto{\pgfqpoint{1.928770in}{0.508924in}}%
\pgfpathlineto{\pgfqpoint{1.910648in}{0.511564in}}%
\pgfpathclose%
\pgfusepath{fill}%
\end{pgfscope}%
\begin{pgfscope}%
\pgfpathrectangle{\pgfqpoint{0.100000in}{0.100000in}}{\pgfqpoint{3.007045in}{1.925000in}}%
\pgfusepath{clip}%
\pgfsetbuttcap%
\pgfsetmiterjoin%
\definecolor{currentfill}{rgb}{0.215686,0.529412,0.754248}%
\pgfsetfillcolor{currentfill}%
\pgfsetlinewidth{0.000000pt}%
\definecolor{currentstroke}{rgb}{0.000000,0.000000,0.000000}%
\pgfsetstrokecolor{currentstroke}%
\pgfsetstrokeopacity{0.000000}%
\pgfsetdash{}{0pt}%
\pgfpathmoveto{\pgfqpoint{1.547367in}{0.412624in}}%
\pgfpathlineto{\pgfqpoint{1.512666in}{0.412373in}}%
\pgfpathlineto{\pgfqpoint{1.512921in}{0.424052in}}%
\pgfpathlineto{\pgfqpoint{1.514242in}{0.453233in}}%
\pgfpathlineto{\pgfqpoint{1.513462in}{0.453257in}}%
\pgfpathlineto{\pgfqpoint{1.514641in}{0.483384in}}%
\pgfpathlineto{\pgfqpoint{1.550168in}{0.481716in}}%
\pgfpathlineto{\pgfqpoint{1.549167in}{0.452106in}}%
\pgfpathlineto{\pgfqpoint{1.547367in}{0.412624in}}%
\pgfpathclose%
\pgfusepath{fill}%
\end{pgfscope}%
\begin{pgfscope}%
\pgfpathrectangle{\pgfqpoint{0.100000in}{0.100000in}}{\pgfqpoint{3.007045in}{1.925000in}}%
\pgfusepath{clip}%
\pgfsetbuttcap%
\pgfsetmiterjoin%
\definecolor{currentfill}{rgb}{0.285936,0.591065,0.787051}%
\pgfsetfillcolor{currentfill}%
\pgfsetlinewidth{0.000000pt}%
\definecolor{currentstroke}{rgb}{0.000000,0.000000,0.000000}%
\pgfsetstrokecolor{currentstroke}%
\pgfsetstrokeopacity{0.000000}%
\pgfsetdash{}{0pt}%
\pgfpathmoveto{\pgfqpoint{2.488716in}{0.626214in}}%
\pgfpathlineto{\pgfqpoint{2.481583in}{0.625736in}}%
\pgfpathlineto{\pgfqpoint{2.479352in}{0.633603in}}%
\pgfpathlineto{\pgfqpoint{2.469936in}{0.633604in}}%
\pgfpathlineto{\pgfqpoint{2.464523in}{0.639748in}}%
\pgfpathlineto{\pgfqpoint{2.457148in}{0.641063in}}%
\pgfpathlineto{\pgfqpoint{2.454390in}{0.661837in}}%
\pgfpathlineto{\pgfqpoint{2.450133in}{0.661231in}}%
\pgfpathlineto{\pgfqpoint{2.450078in}{0.667343in}}%
\pgfpathlineto{\pgfqpoint{2.443165in}{0.674743in}}%
\pgfpathlineto{\pgfqpoint{2.442052in}{0.679496in}}%
\pgfpathlineto{\pgfqpoint{2.443946in}{0.680145in}}%
\pgfpathlineto{\pgfqpoint{2.445730in}{0.689304in}}%
\pgfpathlineto{\pgfqpoint{2.448668in}{0.692985in}}%
\pgfpathlineto{\pgfqpoint{2.447789in}{0.700910in}}%
\pgfpathlineto{\pgfqpoint{2.455203in}{0.702082in}}%
\pgfpathlineto{\pgfqpoint{2.457652in}{0.694920in}}%
\pgfpathlineto{\pgfqpoint{2.477917in}{0.698299in}}%
\pgfpathlineto{\pgfqpoint{2.492063in}{0.691931in}}%
\pgfpathlineto{\pgfqpoint{2.498935in}{0.687218in}}%
\pgfpathlineto{\pgfqpoint{2.503207in}{0.687158in}}%
\pgfpathlineto{\pgfqpoint{2.505019in}{0.681519in}}%
\pgfpathlineto{\pgfqpoint{2.508359in}{0.677599in}}%
\pgfpathlineto{\pgfqpoint{2.503124in}{0.674785in}}%
\pgfpathlineto{\pgfqpoint{2.497961in}{0.669319in}}%
\pgfpathlineto{\pgfqpoint{2.498304in}{0.661191in}}%
\pgfpathlineto{\pgfqpoint{2.503300in}{0.656930in}}%
\pgfpathlineto{\pgfqpoint{2.487249in}{0.654855in}}%
\pgfpathlineto{\pgfqpoint{2.488999in}{0.640420in}}%
\pgfpathlineto{\pgfqpoint{2.504409in}{0.641856in}}%
\pgfpathlineto{\pgfqpoint{2.502750in}{0.627127in}}%
\pgfpathlineto{\pgfqpoint{2.488716in}{0.626214in}}%
\pgfpathclose%
\pgfusepath{fill}%
\end{pgfscope}%
\begin{pgfscope}%
\pgfpathrectangle{\pgfqpoint{0.100000in}{0.100000in}}{\pgfqpoint{3.007045in}{1.925000in}}%
\pgfusepath{clip}%
\pgfsetbuttcap%
\pgfsetmiterjoin%
\definecolor{currentfill}{rgb}{0.585882,0.773641,0.875079}%
\pgfsetfillcolor{currentfill}%
\pgfsetlinewidth{0.000000pt}%
\definecolor{currentstroke}{rgb}{0.000000,0.000000,0.000000}%
\pgfsetstrokecolor{currentstroke}%
\pgfsetstrokeopacity{0.000000}%
\pgfsetdash{}{0pt}%
\pgfpathmoveto{\pgfqpoint{0.499509in}{1.627848in}}%
\pgfpathlineto{\pgfqpoint{0.492514in}{1.625981in}}%
\pgfpathlineto{\pgfqpoint{0.490196in}{1.618627in}}%
\pgfpathlineto{\pgfqpoint{0.462254in}{1.627100in}}%
\pgfpathlineto{\pgfqpoint{0.464341in}{1.633883in}}%
\pgfpathlineto{\pgfqpoint{0.447557in}{1.638893in}}%
\pgfpathlineto{\pgfqpoint{0.447535in}{1.643476in}}%
\pgfpathlineto{\pgfqpoint{0.451005in}{1.654418in}}%
\pgfpathlineto{\pgfqpoint{0.441475in}{1.657513in}}%
\pgfpathlineto{\pgfqpoint{0.434242in}{1.663279in}}%
\pgfpathlineto{\pgfqpoint{0.432010in}{1.670306in}}%
\pgfpathlineto{\pgfqpoint{0.423095in}{1.674955in}}%
\pgfpathlineto{\pgfqpoint{0.416059in}{1.671542in}}%
\pgfpathlineto{\pgfqpoint{0.406441in}{1.674762in}}%
\pgfpathlineto{\pgfqpoint{0.413049in}{1.689210in}}%
\pgfpathlineto{\pgfqpoint{0.416964in}{1.700001in}}%
\pgfpathlineto{\pgfqpoint{0.432427in}{1.695452in}}%
\pgfpathlineto{\pgfqpoint{0.436719in}{1.699407in}}%
\pgfpathlineto{\pgfqpoint{0.437975in}{1.704042in}}%
\pgfpathlineto{\pgfqpoint{0.443505in}{1.702226in}}%
\pgfpathlineto{\pgfqpoint{0.449171in}{1.720447in}}%
\pgfpathlineto{\pgfqpoint{0.444068in}{1.723357in}}%
\pgfpathlineto{\pgfqpoint{0.451037in}{1.744559in}}%
\pgfpathlineto{\pgfqpoint{0.448349in}{1.745454in}}%
\pgfpathlineto{\pgfqpoint{0.451249in}{1.754244in}}%
\pgfpathlineto{\pgfqpoint{0.465504in}{1.749616in}}%
\pgfpathlineto{\pgfqpoint{0.469874in}{1.763357in}}%
\pgfpathlineto{\pgfqpoint{0.484457in}{1.758626in}}%
\pgfpathlineto{\pgfqpoint{0.488462in}{1.755366in}}%
\pgfpathlineto{\pgfqpoint{0.490020in}{1.750804in}}%
\pgfpathlineto{\pgfqpoint{0.496816in}{1.743625in}}%
\pgfpathlineto{\pgfqpoint{0.496421in}{1.737333in}}%
\pgfpathlineto{\pgfqpoint{0.493784in}{1.734570in}}%
\pgfpathlineto{\pgfqpoint{0.496057in}{1.728942in}}%
\pgfpathlineto{\pgfqpoint{0.500179in}{1.724897in}}%
\pgfpathlineto{\pgfqpoint{0.501234in}{1.718853in}}%
\pgfpathlineto{\pgfqpoint{0.506524in}{1.713649in}}%
\pgfpathlineto{\pgfqpoint{0.536189in}{1.704696in}}%
\pgfpathlineto{\pgfqpoint{0.533873in}{1.701067in}}%
\pgfpathlineto{\pgfqpoint{0.530035in}{1.695577in}}%
\pgfpathlineto{\pgfqpoint{0.527922in}{1.686852in}}%
\pgfpathlineto{\pgfqpoint{0.523556in}{1.681083in}}%
\pgfpathlineto{\pgfqpoint{0.521156in}{1.673183in}}%
\pgfpathlineto{\pgfqpoint{0.521441in}{1.664987in}}%
\pgfpathlineto{\pgfqpoint{0.519173in}{1.654672in}}%
\pgfpathlineto{\pgfqpoint{0.512713in}{1.647749in}}%
\pgfpathlineto{\pgfqpoint{0.505425in}{1.642739in}}%
\pgfpathlineto{\pgfqpoint{0.504292in}{1.636653in}}%
\pgfpathlineto{\pgfqpoint{0.499509in}{1.627848in}}%
\pgfpathclose%
\pgfusepath{fill}%
\end{pgfscope}%
\begin{pgfscope}%
\pgfpathrectangle{\pgfqpoint{0.100000in}{0.100000in}}{\pgfqpoint{3.007045in}{1.925000in}}%
\pgfusepath{clip}%
\pgfsetbuttcap%
\pgfsetmiterjoin%
\definecolor{currentfill}{rgb}{0.381776,0.656517,0.824452}%
\pgfsetfillcolor{currentfill}%
\pgfsetlinewidth{0.000000pt}%
\definecolor{currentstroke}{rgb}{0.000000,0.000000,0.000000}%
\pgfsetstrokecolor{currentstroke}%
\pgfsetstrokeopacity{0.000000}%
\pgfsetdash{}{0pt}%
\pgfpathmoveto{\pgfqpoint{2.321001in}{0.750113in}}%
\pgfpathlineto{\pgfqpoint{2.312802in}{0.748173in}}%
\pgfpathlineto{\pgfqpoint{2.295738in}{0.746225in}}%
\pgfpathlineto{\pgfqpoint{2.293144in}{0.771183in}}%
\pgfpathlineto{\pgfqpoint{2.289813in}{0.770791in}}%
\pgfpathlineto{\pgfqpoint{2.287725in}{0.796480in}}%
\pgfpathlineto{\pgfqpoint{2.281013in}{0.795971in}}%
\pgfpathlineto{\pgfqpoint{2.278832in}{0.801530in}}%
\pgfpathlineto{\pgfqpoint{2.267838in}{0.800489in}}%
\pgfpathlineto{\pgfqpoint{2.264190in}{0.805986in}}%
\pgfpathlineto{\pgfqpoint{2.258979in}{0.805773in}}%
\pgfpathlineto{\pgfqpoint{2.263640in}{0.812852in}}%
\pgfpathlineto{\pgfqpoint{2.262319in}{0.816941in}}%
\pgfpathlineto{\pgfqpoint{2.271634in}{0.825054in}}%
\pgfpathlineto{\pgfqpoint{2.279304in}{0.827039in}}%
\pgfpathlineto{\pgfqpoint{2.290820in}{0.827401in}}%
\pgfpathlineto{\pgfqpoint{2.297886in}{0.828687in}}%
\pgfpathlineto{\pgfqpoint{2.298990in}{0.824653in}}%
\pgfpathlineto{\pgfqpoint{2.317320in}{0.826818in}}%
\pgfpathlineto{\pgfqpoint{2.318665in}{0.820809in}}%
\pgfpathlineto{\pgfqpoint{2.326336in}{0.819580in}}%
\pgfpathlineto{\pgfqpoint{2.327623in}{0.805924in}}%
\pgfpathlineto{\pgfqpoint{2.332358in}{0.802455in}}%
\pgfpathlineto{\pgfqpoint{2.343332in}{0.803548in}}%
\pgfpathlineto{\pgfqpoint{2.346769in}{0.793773in}}%
\pgfpathlineto{\pgfqpoint{2.352410in}{0.787703in}}%
\pgfpathlineto{\pgfqpoint{2.352426in}{0.785269in}}%
\pgfpathlineto{\pgfqpoint{2.333138in}{0.783244in}}%
\pgfpathlineto{\pgfqpoint{2.335233in}{0.759815in}}%
\pgfpathlineto{\pgfqpoint{2.317396in}{0.757985in}}%
\pgfpathlineto{\pgfqpoint{2.321001in}{0.750113in}}%
\pgfpathclose%
\pgfusepath{fill}%
\end{pgfscope}%
\begin{pgfscope}%
\pgfpathrectangle{\pgfqpoint{0.100000in}{0.100000in}}{\pgfqpoint{3.007045in}{1.925000in}}%
\pgfusepath{clip}%
\pgfsetbuttcap%
\pgfsetmiterjoin%
\definecolor{currentfill}{rgb}{0.592157,0.777086,0.876432}%
\pgfsetfillcolor{currentfill}%
\pgfsetlinewidth{0.000000pt}%
\definecolor{currentstroke}{rgb}{0.000000,0.000000,0.000000}%
\pgfsetstrokecolor{currentstroke}%
\pgfsetstrokeopacity{0.000000}%
\pgfsetdash{}{0pt}%
\pgfpathmoveto{\pgfqpoint{1.104848in}{1.099688in}}%
\pgfpathlineto{\pgfqpoint{1.098980in}{1.105934in}}%
\pgfpathlineto{\pgfqpoint{1.087280in}{1.105261in}}%
\pgfpathlineto{\pgfqpoint{1.085598in}{1.110401in}}%
\pgfpathlineto{\pgfqpoint{1.056798in}{1.114784in}}%
\pgfpathlineto{\pgfqpoint{1.045110in}{1.115632in}}%
\pgfpathlineto{\pgfqpoint{1.047972in}{1.134089in}}%
\pgfpathlineto{\pgfqpoint{1.048160in}{1.141447in}}%
\pgfpathlineto{\pgfqpoint{1.050444in}{1.156057in}}%
\pgfpathlineto{\pgfqpoint{1.059704in}{1.212390in}}%
\pgfpathlineto{\pgfqpoint{1.091973in}{1.207431in}}%
\pgfpathlineto{\pgfqpoint{1.141305in}{1.200383in}}%
\pgfpathlineto{\pgfqpoint{1.139110in}{1.197239in}}%
\pgfpathlineto{\pgfqpoint{1.142160in}{1.192960in}}%
\pgfpathlineto{\pgfqpoint{1.136431in}{1.191196in}}%
\pgfpathlineto{\pgfqpoint{1.128210in}{1.131426in}}%
\pgfpathlineto{\pgfqpoint{1.121297in}{1.132386in}}%
\pgfpathlineto{\pgfqpoint{1.123455in}{1.127357in}}%
\pgfpathlineto{\pgfqpoint{1.123423in}{1.117299in}}%
\pgfpathlineto{\pgfqpoint{1.121625in}{1.109932in}}%
\pgfpathlineto{\pgfqpoint{1.112363in}{1.107267in}}%
\pgfpathlineto{\pgfqpoint{1.104848in}{1.099688in}}%
\pgfpathclose%
\pgfusepath{fill}%
\end{pgfscope}%
\begin{pgfscope}%
\pgfpathrectangle{\pgfqpoint{0.100000in}{0.100000in}}{\pgfqpoint{3.007045in}{1.925000in}}%
\pgfusepath{clip}%
\pgfsetbuttcap%
\pgfsetmiterjoin%
\definecolor{currentfill}{rgb}{0.093272,0.396878,0.673664}%
\pgfsetfillcolor{currentfill}%
\pgfsetlinewidth{0.000000pt}%
\definecolor{currentstroke}{rgb}{0.000000,0.000000,0.000000}%
\pgfsetstrokecolor{currentstroke}%
\pgfsetstrokeopacity{0.000000}%
\pgfsetdash{}{0pt}%
\pgfpathmoveto{\pgfqpoint{1.001047in}{0.482991in}}%
\pgfpathlineto{\pgfqpoint{0.989639in}{0.489920in}}%
\pgfpathlineto{\pgfqpoint{0.966890in}{0.503916in}}%
\pgfpathlineto{\pgfqpoint{0.959476in}{0.508253in}}%
\pgfpathlineto{\pgfqpoint{0.952820in}{0.498342in}}%
\pgfpathlineto{\pgfqpoint{0.944682in}{0.503522in}}%
\pgfpathlineto{\pgfqpoint{0.941118in}{0.505304in}}%
\pgfpathlineto{\pgfqpoint{0.919823in}{0.519872in}}%
\pgfpathlineto{\pgfqpoint{0.895892in}{0.537283in}}%
\pgfpathlineto{\pgfqpoint{0.877896in}{0.551221in}}%
\pgfpathlineto{\pgfqpoint{0.878340in}{0.551804in}}%
\pgfpathlineto{\pgfqpoint{0.871135in}{0.557489in}}%
\pgfpathlineto{\pgfqpoint{0.872509in}{0.559210in}}%
\pgfpathlineto{\pgfqpoint{0.865552in}{0.564907in}}%
\pgfpathlineto{\pgfqpoint{0.867050in}{0.566744in}}%
\pgfpathlineto{\pgfqpoint{0.859713in}{0.572782in}}%
\pgfpathlineto{\pgfqpoint{0.847780in}{0.582922in}}%
\pgfpathlineto{\pgfqpoint{0.848948in}{0.584275in}}%
\pgfpathlineto{\pgfqpoint{0.845795in}{0.587009in}}%
\pgfpathlineto{\pgfqpoint{0.844619in}{0.585660in}}%
\pgfpathlineto{\pgfqpoint{0.837473in}{0.591959in}}%
\pgfpathlineto{\pgfqpoint{0.838673in}{0.593306in}}%
\pgfpathlineto{\pgfqpoint{0.835607in}{0.596048in}}%
\pgfpathlineto{\pgfqpoint{0.834401in}{0.594707in}}%
\pgfpathlineto{\pgfqpoint{0.830895in}{0.597873in}}%
\pgfpathlineto{\pgfqpoint{0.832109in}{0.599208in}}%
\pgfpathlineto{\pgfqpoint{0.826925in}{0.603954in}}%
\pgfpathlineto{\pgfqpoint{0.825701in}{0.602628in}}%
\pgfpathlineto{\pgfqpoint{0.813031in}{0.614554in}}%
\pgfpathlineto{\pgfqpoint{0.809438in}{0.610810in}}%
\pgfpathlineto{\pgfqpoint{0.803267in}{0.616806in}}%
\pgfpathlineto{\pgfqpoint{0.800883in}{0.622716in}}%
\pgfpathlineto{\pgfqpoint{0.801608in}{0.625470in}}%
\pgfpathlineto{\pgfqpoint{0.801127in}{0.628309in}}%
\pgfpathlineto{\pgfqpoint{0.799100in}{0.632308in}}%
\pgfpathlineto{\pgfqpoint{0.801794in}{0.631253in}}%
\pgfpathlineto{\pgfqpoint{0.804644in}{0.631085in}}%
\pgfpathlineto{\pgfqpoint{0.806937in}{0.632414in}}%
\pgfpathlineto{\pgfqpoint{0.812022in}{0.637296in}}%
\pgfpathlineto{\pgfqpoint{0.813808in}{0.636711in}}%
\pgfpathlineto{\pgfqpoint{0.817714in}{0.632498in}}%
\pgfpathlineto{\pgfqpoint{0.826485in}{0.626672in}}%
\pgfpathlineto{\pgfqpoint{0.829143in}{0.625614in}}%
\pgfpathlineto{\pgfqpoint{0.834897in}{0.625269in}}%
\pgfpathlineto{\pgfqpoint{0.842826in}{0.626872in}}%
\pgfpathlineto{\pgfqpoint{0.846807in}{0.629874in}}%
\pgfpathlineto{\pgfqpoint{0.847784in}{0.631457in}}%
\pgfpathlineto{\pgfqpoint{0.850593in}{0.633815in}}%
\pgfpathlineto{\pgfqpoint{0.852477in}{0.634580in}}%
\pgfpathlineto{\pgfqpoint{0.858398in}{0.634527in}}%
\pgfpathlineto{\pgfqpoint{0.860204in}{0.636304in}}%
\pgfpathlineto{\pgfqpoint{0.864117in}{0.636616in}}%
\pgfpathlineto{\pgfqpoint{0.868789in}{0.637634in}}%
\pgfpathlineto{\pgfqpoint{0.871460in}{0.633851in}}%
\pgfpathlineto{\pgfqpoint{0.876444in}{0.632688in}}%
\pgfpathlineto{\pgfqpoint{0.885282in}{0.633242in}}%
\pgfpathlineto{\pgfqpoint{0.894306in}{0.634804in}}%
\pgfpathlineto{\pgfqpoint{0.897261in}{0.633456in}}%
\pgfpathlineto{\pgfqpoint{0.897619in}{0.631397in}}%
\pgfpathlineto{\pgfqpoint{0.901143in}{0.628563in}}%
\pgfpathlineto{\pgfqpoint{0.904509in}{0.626669in}}%
\pgfpathlineto{\pgfqpoint{0.907341in}{0.625778in}}%
\pgfpathlineto{\pgfqpoint{0.913294in}{0.626230in}}%
\pgfpathlineto{\pgfqpoint{0.921174in}{0.629150in}}%
\pgfpathlineto{\pgfqpoint{0.923879in}{0.629562in}}%
\pgfpathlineto{\pgfqpoint{0.925211in}{0.625453in}}%
\pgfpathlineto{\pgfqpoint{0.924492in}{0.622843in}}%
\pgfpathlineto{\pgfqpoint{0.926548in}{0.622888in}}%
\pgfpathlineto{\pgfqpoint{0.927653in}{0.620890in}}%
\pgfpathlineto{\pgfqpoint{0.927665in}{0.618526in}}%
\pgfpathlineto{\pgfqpoint{0.923257in}{0.616818in}}%
\pgfpathlineto{\pgfqpoint{0.919869in}{0.618044in}}%
\pgfpathlineto{\pgfqpoint{0.920400in}{0.614171in}}%
\pgfpathlineto{\pgfqpoint{0.923838in}{0.612830in}}%
\pgfpathlineto{\pgfqpoint{0.925840in}{0.614446in}}%
\pgfpathlineto{\pgfqpoint{0.930393in}{0.614357in}}%
\pgfpathlineto{\pgfqpoint{0.932785in}{0.611562in}}%
\pgfpathlineto{\pgfqpoint{0.930092in}{0.608495in}}%
\pgfpathlineto{\pgfqpoint{0.929589in}{0.606710in}}%
\pgfpathlineto{\pgfqpoint{0.931595in}{0.605231in}}%
\pgfpathlineto{\pgfqpoint{0.931935in}{0.603180in}}%
\pgfpathlineto{\pgfqpoint{0.934772in}{0.603947in}}%
\pgfpathlineto{\pgfqpoint{0.938741in}{0.601961in}}%
\pgfpathlineto{\pgfqpoint{0.941530in}{0.601138in}}%
\pgfpathlineto{\pgfqpoint{0.943304in}{0.596728in}}%
\pgfpathlineto{\pgfqpoint{0.945069in}{0.596940in}}%
\pgfpathlineto{\pgfqpoint{0.946594in}{0.594438in}}%
\pgfpathlineto{\pgfqpoint{0.941495in}{0.590893in}}%
\pgfpathlineto{\pgfqpoint{0.939975in}{0.589037in}}%
\pgfpathlineto{\pgfqpoint{0.945374in}{0.585331in}}%
\pgfpathlineto{\pgfqpoint{0.944737in}{0.583986in}}%
\pgfpathlineto{\pgfqpoint{0.942095in}{0.583582in}}%
\pgfpathlineto{\pgfqpoint{0.946113in}{0.580534in}}%
\pgfpathlineto{\pgfqpoint{0.946529in}{0.578018in}}%
\pgfpathlineto{\pgfqpoint{0.949099in}{0.579130in}}%
\pgfpathlineto{\pgfqpoint{0.953212in}{0.577928in}}%
\pgfpathlineto{\pgfqpoint{0.953739in}{0.576380in}}%
\pgfpathlineto{\pgfqpoint{0.951749in}{0.574813in}}%
\pgfpathlineto{\pgfqpoint{0.955961in}{0.573632in}}%
\pgfpathlineto{\pgfqpoint{0.958806in}{0.574118in}}%
\pgfpathlineto{\pgfqpoint{0.962206in}{0.570550in}}%
\pgfpathlineto{\pgfqpoint{0.963220in}{0.570510in}}%
\pgfpathlineto{\pgfqpoint{0.964026in}{0.566308in}}%
\pgfpathlineto{\pgfqpoint{0.966643in}{0.564927in}}%
\pgfpathlineto{\pgfqpoint{0.966745in}{0.562715in}}%
\pgfpathlineto{\pgfqpoint{0.968656in}{0.561617in}}%
\pgfpathlineto{\pgfqpoint{0.969071in}{0.557591in}}%
\pgfpathlineto{\pgfqpoint{0.974627in}{0.550898in}}%
\pgfpathlineto{\pgfqpoint{0.976869in}{0.550826in}}%
\pgfpathlineto{\pgfqpoint{0.980019in}{0.548774in}}%
\pgfpathlineto{\pgfqpoint{0.981988in}{0.546062in}}%
\pgfpathlineto{\pgfqpoint{0.984408in}{0.544629in}}%
\pgfpathlineto{\pgfqpoint{0.986456in}{0.538416in}}%
\pgfpathlineto{\pgfqpoint{0.988132in}{0.536134in}}%
\pgfpathlineto{\pgfqpoint{0.992596in}{0.535000in}}%
\pgfpathlineto{\pgfqpoint{0.998469in}{0.532993in}}%
\pgfpathlineto{\pgfqpoint{1.000842in}{0.532798in}}%
\pgfpathlineto{\pgfqpoint{1.003539in}{0.531122in}}%
\pgfpathlineto{\pgfqpoint{1.006445in}{0.526960in}}%
\pgfpathlineto{\pgfqpoint{1.008246in}{0.522245in}}%
\pgfpathlineto{\pgfqpoint{1.008217in}{0.520872in}}%
\pgfpathlineto{\pgfqpoint{1.010008in}{0.517358in}}%
\pgfpathlineto{\pgfqpoint{1.011929in}{0.514984in}}%
\pgfpathlineto{\pgfqpoint{1.012633in}{0.512658in}}%
\pgfpathlineto{\pgfqpoint{1.015157in}{0.509681in}}%
\pgfpathlineto{\pgfqpoint{1.014103in}{0.506544in}}%
\pgfpathlineto{\pgfqpoint{1.001047in}{0.482991in}}%
\pgfpathclose%
\pgfusepath{fill}%
\end{pgfscope}%
\begin{pgfscope}%
\pgfpathrectangle{\pgfqpoint{0.100000in}{0.100000in}}{\pgfqpoint{3.007045in}{1.925000in}}%
\pgfusepath{clip}%
\pgfsetbuttcap%
\pgfsetmiterjoin%
\definecolor{currentfill}{rgb}{0.541961,0.749527,0.865606}%
\pgfsetfillcolor{currentfill}%
\pgfsetlinewidth{0.000000pt}%
\definecolor{currentstroke}{rgb}{0.000000,0.000000,0.000000}%
\pgfsetstrokecolor{currentstroke}%
\pgfsetstrokeopacity{0.000000}%
\pgfsetdash{}{0pt}%
\pgfpathmoveto{\pgfqpoint{2.119654in}{1.040653in}}%
\pgfpathlineto{\pgfqpoint{2.116603in}{1.041678in}}%
\pgfpathlineto{\pgfqpoint{2.115921in}{1.053268in}}%
\pgfpathlineto{\pgfqpoint{2.100489in}{1.052130in}}%
\pgfpathlineto{\pgfqpoint{2.099179in}{1.072352in}}%
\pgfpathlineto{\pgfqpoint{2.116398in}{1.073683in}}%
\pgfpathlineto{\pgfqpoint{2.128732in}{1.074544in}}%
\pgfpathlineto{\pgfqpoint{2.131121in}{1.073626in}}%
\pgfpathlineto{\pgfqpoint{2.134908in}{1.067855in}}%
\pgfpathlineto{\pgfqpoint{2.128798in}{1.058421in}}%
\pgfpathlineto{\pgfqpoint{2.133456in}{1.046059in}}%
\pgfpathlineto{\pgfqpoint{2.123429in}{1.044085in}}%
\pgfpathlineto{\pgfqpoint{2.119654in}{1.040653in}}%
\pgfpathclose%
\pgfusepath{fill}%
\end{pgfscope}%
\begin{pgfscope}%
\pgfpathrectangle{\pgfqpoint{0.100000in}{0.100000in}}{\pgfqpoint{3.007045in}{1.925000in}}%
\pgfusepath{clip}%
\pgfsetbuttcap%
\pgfsetmiterjoin%
\definecolor{currentfill}{rgb}{0.376732,0.653072,0.822484}%
\pgfsetfillcolor{currentfill}%
\pgfsetlinewidth{0.000000pt}%
\definecolor{currentstroke}{rgb}{0.000000,0.000000,0.000000}%
\pgfsetstrokecolor{currentstroke}%
\pgfsetstrokeopacity{0.000000}%
\pgfsetdash{}{0pt}%
\pgfpathmoveto{\pgfqpoint{2.073753in}{0.803617in}}%
\pgfpathlineto{\pgfqpoint{2.059401in}{0.802761in}}%
\pgfpathlineto{\pgfqpoint{2.050685in}{0.804294in}}%
\pgfpathlineto{\pgfqpoint{2.050947in}{0.799400in}}%
\pgfpathlineto{\pgfqpoint{2.039467in}{0.797809in}}%
\pgfpathlineto{\pgfqpoint{2.039761in}{0.792018in}}%
\pgfpathlineto{\pgfqpoint{2.030424in}{0.792915in}}%
\pgfpathlineto{\pgfqpoint{2.036938in}{0.797675in}}%
\pgfpathlineto{\pgfqpoint{2.022165in}{0.797050in}}%
\pgfpathlineto{\pgfqpoint{2.021640in}{0.808564in}}%
\pgfpathlineto{\pgfqpoint{2.010153in}{0.808128in}}%
\pgfpathlineto{\pgfqpoint{2.009788in}{0.816800in}}%
\pgfpathlineto{\pgfqpoint{1.993574in}{0.816135in}}%
\pgfpathlineto{\pgfqpoint{2.003512in}{0.827202in}}%
\pgfpathlineto{\pgfqpoint{2.003733in}{0.832951in}}%
\pgfpathlineto{\pgfqpoint{2.013158in}{0.836265in}}%
\pgfpathlineto{\pgfqpoint{2.013407in}{0.843839in}}%
\pgfpathlineto{\pgfqpoint{2.022978in}{0.837629in}}%
\pgfpathlineto{\pgfqpoint{2.028291in}{0.837950in}}%
\pgfpathlineto{\pgfqpoint{2.027867in}{0.843557in}}%
\pgfpathlineto{\pgfqpoint{2.033456in}{0.846814in}}%
\pgfpathlineto{\pgfqpoint{2.072679in}{0.849028in}}%
\pgfpathlineto{\pgfqpoint{2.085442in}{0.845945in}}%
\pgfpathlineto{\pgfqpoint{2.085917in}{0.838289in}}%
\pgfpathlineto{\pgfqpoint{2.086876in}{0.823888in}}%
\pgfpathlineto{\pgfqpoint{2.072446in}{0.823053in}}%
\pgfpathlineto{\pgfqpoint{2.073753in}{0.803617in}}%
\pgfpathclose%
\pgfusepath{fill}%
\end{pgfscope}%
\begin{pgfscope}%
\pgfpathrectangle{\pgfqpoint{0.100000in}{0.100000in}}{\pgfqpoint{3.007045in}{1.925000in}}%
\pgfusepath{clip}%
\pgfsetbuttcap%
\pgfsetmiterjoin%
\definecolor{currentfill}{rgb}{0.726028,0.837386,0.919600}%
\pgfsetfillcolor{currentfill}%
\pgfsetlinewidth{0.000000pt}%
\definecolor{currentstroke}{rgb}{0.000000,0.000000,0.000000}%
\pgfsetstrokecolor{currentstroke}%
\pgfsetstrokeopacity{0.000000}%
\pgfsetdash{}{0pt}%
\pgfpathmoveto{\pgfqpoint{2.938382in}{1.591168in}}%
\pgfpathlineto{\pgfqpoint{2.930725in}{1.593936in}}%
\pgfpathlineto{\pgfqpoint{2.931682in}{1.596631in}}%
\pgfpathlineto{\pgfqpoint{2.926133in}{1.600139in}}%
\pgfpathlineto{\pgfqpoint{2.921331in}{1.597973in}}%
\pgfpathlineto{\pgfqpoint{2.920147in}{1.603400in}}%
\pgfpathlineto{\pgfqpoint{2.912256in}{1.602890in}}%
\pgfpathlineto{\pgfqpoint{2.904580in}{1.599504in}}%
\pgfpathlineto{\pgfqpoint{2.895122in}{1.630668in}}%
\pgfpathlineto{\pgfqpoint{2.885278in}{1.661827in}}%
\pgfpathlineto{\pgfqpoint{2.874149in}{1.694991in}}%
\pgfpathlineto{\pgfqpoint{2.879401in}{1.698668in}}%
\pgfpathlineto{\pgfqpoint{2.883573in}{1.693339in}}%
\pgfpathlineto{\pgfqpoint{2.886335in}{1.697801in}}%
\pgfpathlineto{\pgfqpoint{2.885576in}{1.706733in}}%
\pgfpathlineto{\pgfqpoint{2.889390in}{1.704990in}}%
\pgfpathlineto{\pgfqpoint{2.892641in}{1.707989in}}%
\pgfpathlineto{\pgfqpoint{2.886793in}{1.712832in}}%
\pgfpathlineto{\pgfqpoint{2.888624in}{1.719661in}}%
\pgfpathlineto{\pgfqpoint{2.897841in}{1.730562in}}%
\pgfpathlineto{\pgfqpoint{2.896524in}{1.736677in}}%
\pgfpathlineto{\pgfqpoint{2.900020in}{1.741049in}}%
\pgfpathlineto{\pgfqpoint{2.900699in}{1.746177in}}%
\pgfpathlineto{\pgfqpoint{2.895909in}{1.750061in}}%
\pgfpathlineto{\pgfqpoint{2.897245in}{1.759098in}}%
\pgfpathlineto{\pgfqpoint{2.893988in}{1.761345in}}%
\pgfpathlineto{\pgfqpoint{2.895208in}{1.771201in}}%
\pgfpathlineto{\pgfqpoint{2.900371in}{1.778621in}}%
\pgfpathlineto{\pgfqpoint{2.898990in}{1.789167in}}%
\pgfpathlineto{\pgfqpoint{2.912260in}{1.792925in}}%
\pgfpathlineto{\pgfqpoint{2.915056in}{1.781252in}}%
\pgfpathlineto{\pgfqpoint{2.922774in}{1.754622in}}%
\pgfpathlineto{\pgfqpoint{2.926062in}{1.747145in}}%
\pgfpathlineto{\pgfqpoint{2.925595in}{1.735765in}}%
\pgfpathlineto{\pgfqpoint{2.930414in}{1.732814in}}%
\pgfpathlineto{\pgfqpoint{2.928393in}{1.726358in}}%
\pgfpathlineto{\pgfqpoint{2.945232in}{1.694510in}}%
\pgfpathlineto{\pgfqpoint{2.956012in}{1.701938in}}%
\pgfpathlineto{\pgfqpoint{2.966515in}{1.680747in}}%
\pgfpathlineto{\pgfqpoint{2.960139in}{1.680504in}}%
\pgfpathlineto{\pgfqpoint{2.962312in}{1.673903in}}%
\pgfpathlineto{\pgfqpoint{2.963183in}{1.661995in}}%
\pgfpathlineto{\pgfqpoint{2.962617in}{1.653501in}}%
\pgfpathlineto{\pgfqpoint{2.965671in}{1.652396in}}%
\pgfpathlineto{\pgfqpoint{2.969629in}{1.646428in}}%
\pgfpathlineto{\pgfqpoint{2.972759in}{1.647244in}}%
\pgfpathlineto{\pgfqpoint{2.978270in}{1.638556in}}%
\pgfpathlineto{\pgfqpoint{2.975595in}{1.630839in}}%
\pgfpathlineto{\pgfqpoint{2.972528in}{1.630382in}}%
\pgfpathlineto{\pgfqpoint{2.972069in}{1.622078in}}%
\pgfpathlineto{\pgfqpoint{2.963086in}{1.616945in}}%
\pgfpathlineto{\pgfqpoint{2.962287in}{1.612214in}}%
\pgfpathlineto{\pgfqpoint{2.958299in}{1.609229in}}%
\pgfpathlineto{\pgfqpoint{2.947252in}{1.614168in}}%
\pgfpathlineto{\pgfqpoint{2.941831in}{1.608526in}}%
\pgfpathlineto{\pgfqpoint{2.940412in}{1.601930in}}%
\pgfpathlineto{\pgfqpoint{2.945057in}{1.595018in}}%
\pgfpathlineto{\pgfqpoint{2.938382in}{1.591168in}}%
\pgfpathclose%
\pgfusepath{fill}%
\end{pgfscope}%
\begin{pgfscope}%
\pgfpathrectangle{\pgfqpoint{0.100000in}{0.100000in}}{\pgfqpoint{3.007045in}{1.925000in}}%
\pgfusepath{clip}%
\pgfsetbuttcap%
\pgfsetmiterjoin%
\definecolor{currentfill}{rgb}{0.498039,0.725413,0.856132}%
\pgfsetfillcolor{currentfill}%
\pgfsetlinewidth{0.000000pt}%
\definecolor{currentstroke}{rgb}{0.000000,0.000000,0.000000}%
\pgfsetstrokecolor{currentstroke}%
\pgfsetstrokeopacity{0.000000}%
\pgfsetdash{}{0pt}%
\pgfpathmoveto{\pgfqpoint{2.116053in}{0.743899in}}%
\pgfpathlineto{\pgfqpoint{2.087936in}{0.742436in}}%
\pgfpathlineto{\pgfqpoint{2.086624in}{0.765838in}}%
\pgfpathlineto{\pgfqpoint{2.079034in}{0.765390in}}%
\pgfpathlineto{\pgfqpoint{2.078426in}{0.775122in}}%
\pgfpathlineto{\pgfqpoint{2.075429in}{0.776700in}}%
\pgfpathlineto{\pgfqpoint{2.074344in}{0.794001in}}%
\pgfpathlineto{\pgfqpoint{2.091654in}{0.795102in}}%
\pgfpathlineto{\pgfqpoint{2.105804in}{0.798159in}}%
\pgfpathlineto{\pgfqpoint{2.105570in}{0.801993in}}%
\pgfpathlineto{\pgfqpoint{2.117152in}{0.802691in}}%
\pgfpathlineto{\pgfqpoint{2.117772in}{0.794063in}}%
\pgfpathlineto{\pgfqpoint{2.129211in}{0.792870in}}%
\pgfpathlineto{\pgfqpoint{2.131570in}{0.789934in}}%
\pgfpathlineto{\pgfqpoint{2.128772in}{0.785353in}}%
\pgfpathlineto{\pgfqpoint{2.121230in}{0.782758in}}%
\pgfpathlineto{\pgfqpoint{2.122360in}{0.768291in}}%
\pgfpathlineto{\pgfqpoint{2.114414in}{0.767709in}}%
\pgfpathlineto{\pgfqpoint{2.116053in}{0.743899in}}%
\pgfpathclose%
\pgfusepath{fill}%
\end{pgfscope}%
\begin{pgfscope}%
\pgfpathrectangle{\pgfqpoint{0.100000in}{0.100000in}}{\pgfqpoint{3.007045in}{1.925000in}}%
\pgfusepath{clip}%
\pgfsetbuttcap%
\pgfsetmiterjoin%
\definecolor{currentfill}{rgb}{0.567059,0.763306,0.871019}%
\pgfsetfillcolor{currentfill}%
\pgfsetlinewidth{0.000000pt}%
\definecolor{currentstroke}{rgb}{0.000000,0.000000,0.000000}%
\pgfsetstrokecolor{currentstroke}%
\pgfsetstrokeopacity{0.000000}%
\pgfsetdash{}{0pt}%
\pgfpathmoveto{\pgfqpoint{0.536189in}{1.704696in}}%
\pgfpathlineto{\pgfqpoint{0.506524in}{1.713649in}}%
\pgfpathlineto{\pgfqpoint{0.501234in}{1.718853in}}%
\pgfpathlineto{\pgfqpoint{0.500179in}{1.724897in}}%
\pgfpathlineto{\pgfqpoint{0.496057in}{1.728942in}}%
\pgfpathlineto{\pgfqpoint{0.493784in}{1.734570in}}%
\pgfpathlineto{\pgfqpoint{0.496421in}{1.737333in}}%
\pgfpathlineto{\pgfqpoint{0.496816in}{1.743625in}}%
\pgfpathlineto{\pgfqpoint{0.490020in}{1.750804in}}%
\pgfpathlineto{\pgfqpoint{0.488462in}{1.755366in}}%
\pgfpathlineto{\pgfqpoint{0.484457in}{1.758626in}}%
\pgfpathlineto{\pgfqpoint{0.469874in}{1.763357in}}%
\pgfpathlineto{\pgfqpoint{0.475014in}{1.769793in}}%
\pgfpathlineto{\pgfqpoint{0.475846in}{1.775523in}}%
\pgfpathlineto{\pgfqpoint{0.480074in}{1.779212in}}%
\pgfpathlineto{\pgfqpoint{0.481516in}{1.783693in}}%
\pgfpathlineto{\pgfqpoint{0.488870in}{1.806888in}}%
\pgfpathlineto{\pgfqpoint{0.498370in}{1.806808in}}%
\pgfpathlineto{\pgfqpoint{0.507744in}{1.796488in}}%
\pgfpathlineto{\pgfqpoint{0.510888in}{1.782623in}}%
\pgfpathlineto{\pgfqpoint{0.524572in}{1.784544in}}%
\pgfpathlineto{\pgfqpoint{0.529311in}{1.780896in}}%
\pgfpathlineto{\pgfqpoint{0.536092in}{1.785540in}}%
\pgfpathlineto{\pgfqpoint{0.542761in}{1.806508in}}%
\pgfpathlineto{\pgfqpoint{0.574289in}{1.797104in}}%
\pgfpathlineto{\pgfqpoint{0.567701in}{1.775221in}}%
\pgfpathlineto{\pgfqpoint{0.563789in}{1.776383in}}%
\pgfpathlineto{\pgfqpoint{0.558944in}{1.759738in}}%
\pgfpathlineto{\pgfqpoint{0.561232in}{1.755167in}}%
\pgfpathlineto{\pgfqpoint{0.551653in}{1.755806in}}%
\pgfpathlineto{\pgfqpoint{0.545786in}{1.757536in}}%
\pgfpathlineto{\pgfqpoint{0.542490in}{1.755480in}}%
\pgfpathlineto{\pgfqpoint{0.540688in}{1.747112in}}%
\pgfpathlineto{\pgfqpoint{0.547023in}{1.735645in}}%
\pgfpathlineto{\pgfqpoint{0.545552in}{1.726930in}}%
\pgfpathlineto{\pgfqpoint{0.541409in}{1.723741in}}%
\pgfpathlineto{\pgfqpoint{0.541696in}{1.715411in}}%
\pgfpathlineto{\pgfqpoint{0.535329in}{1.713270in}}%
\pgfpathlineto{\pgfqpoint{0.536189in}{1.704696in}}%
\pgfpathclose%
\pgfusepath{fill}%
\end{pgfscope}%
\begin{pgfscope}%
\pgfpathrectangle{\pgfqpoint{0.100000in}{0.100000in}}{\pgfqpoint{3.007045in}{1.925000in}}%
\pgfusepath{clip}%
\pgfsetbuttcap%
\pgfsetmiterjoin%
\definecolor{currentfill}{rgb}{0.627605,0.795556,0.885152}%
\pgfsetfillcolor{currentfill}%
\pgfsetlinewidth{0.000000pt}%
\definecolor{currentstroke}{rgb}{0.000000,0.000000,0.000000}%
\pgfsetstrokecolor{currentstroke}%
\pgfsetstrokeopacity{0.000000}%
\pgfsetdash{}{0pt}%
\pgfpathmoveto{\pgfqpoint{2.539586in}{0.435076in}}%
\pgfpathlineto{\pgfqpoint{2.510105in}{0.430828in}}%
\pgfpathlineto{\pgfqpoint{2.515004in}{0.442453in}}%
\pgfpathlineto{\pgfqpoint{2.510239in}{0.443614in}}%
\pgfpathlineto{\pgfqpoint{2.504846in}{0.439692in}}%
\pgfpathlineto{\pgfqpoint{2.503926in}{0.433696in}}%
\pgfpathlineto{\pgfqpoint{2.500319in}{0.432366in}}%
\pgfpathlineto{\pgfqpoint{2.490587in}{0.443682in}}%
\pgfpathlineto{\pgfqpoint{2.492902in}{0.456197in}}%
\pgfpathlineto{\pgfqpoint{2.490893in}{0.464809in}}%
\pgfpathlineto{\pgfqpoint{2.493043in}{0.468035in}}%
\pgfpathlineto{\pgfqpoint{2.496096in}{0.492116in}}%
\pgfpathlineto{\pgfqpoint{2.495655in}{0.499863in}}%
\pgfpathlineto{\pgfqpoint{2.508454in}{0.501648in}}%
\pgfpathlineto{\pgfqpoint{2.517736in}{0.501110in}}%
\pgfpathlineto{\pgfqpoint{2.521853in}{0.495228in}}%
\pgfpathlineto{\pgfqpoint{2.531241in}{0.493155in}}%
\pgfpathlineto{\pgfqpoint{2.533154in}{0.479245in}}%
\pgfpathlineto{\pgfqpoint{2.530722in}{0.475320in}}%
\pgfpathlineto{\pgfqpoint{2.534472in}{0.469934in}}%
\pgfpathlineto{\pgfqpoint{2.539586in}{0.435076in}}%
\pgfpathclose%
\pgfusepath{fill}%
\end{pgfscope}%
\begin{pgfscope}%
\pgfpathrectangle{\pgfqpoint{0.100000in}{0.100000in}}{\pgfqpoint{3.007045in}{1.925000in}}%
\pgfusepath{clip}%
\pgfsetbuttcap%
\pgfsetmiterjoin%
\definecolor{currentfill}{rgb}{0.701423,0.826928,0.910988}%
\pgfsetfillcolor{currentfill}%
\pgfsetlinewidth{0.000000pt}%
\definecolor{currentstroke}{rgb}{0.000000,0.000000,0.000000}%
\pgfsetstrokecolor{currentstroke}%
\pgfsetstrokeopacity{0.000000}%
\pgfsetdash{}{0pt}%
\pgfpathmoveto{\pgfqpoint{2.867151in}{1.544262in}}%
\pgfpathlineto{\pgfqpoint{2.867535in}{1.548390in}}%
\pgfpathlineto{\pgfqpoint{2.861539in}{1.556030in}}%
\pgfpathlineto{\pgfqpoint{2.863308in}{1.564293in}}%
\pgfpathlineto{\pgfqpoint{2.867917in}{1.577405in}}%
\pgfpathlineto{\pgfqpoint{2.870513in}{1.573950in}}%
\pgfpathlineto{\pgfqpoint{2.876946in}{1.580426in}}%
\pgfpathlineto{\pgfqpoint{2.875977in}{1.584608in}}%
\pgfpathlineto{\pgfqpoint{2.881617in}{1.586847in}}%
\pgfpathlineto{\pgfqpoint{2.876355in}{1.598852in}}%
\pgfpathlineto{\pgfqpoint{2.885308in}{1.602759in}}%
\pgfpathlineto{\pgfqpoint{2.882477in}{1.612503in}}%
\pgfpathlineto{\pgfqpoint{2.877796in}{1.621050in}}%
\pgfpathlineto{\pgfqpoint{2.884101in}{1.617605in}}%
\pgfpathlineto{\pgfqpoint{2.884731in}{1.625099in}}%
\pgfpathlineto{\pgfqpoint{2.894447in}{1.627269in}}%
\pgfpathlineto{\pgfqpoint{2.895122in}{1.630668in}}%
\pgfpathlineto{\pgfqpoint{2.904580in}{1.599504in}}%
\pgfpathlineto{\pgfqpoint{2.912256in}{1.602890in}}%
\pgfpathlineto{\pgfqpoint{2.920147in}{1.603400in}}%
\pgfpathlineto{\pgfqpoint{2.921331in}{1.597973in}}%
\pgfpathlineto{\pgfqpoint{2.926133in}{1.600139in}}%
\pgfpathlineto{\pgfqpoint{2.931682in}{1.596631in}}%
\pgfpathlineto{\pgfqpoint{2.930725in}{1.593936in}}%
\pgfpathlineto{\pgfqpoint{2.938382in}{1.591168in}}%
\pgfpathlineto{\pgfqpoint{2.939431in}{1.584564in}}%
\pgfpathlineto{\pgfqpoint{2.937676in}{1.578867in}}%
\pgfpathlineto{\pgfqpoint{2.933280in}{1.576054in}}%
\pgfpathlineto{\pgfqpoint{2.932692in}{1.561193in}}%
\pgfpathlineto{\pgfqpoint{2.928157in}{1.545679in}}%
\pgfpathlineto{\pgfqpoint{2.928476in}{1.542999in}}%
\pgfpathlineto{\pgfqpoint{2.922960in}{1.542355in}}%
\pgfpathlineto{\pgfqpoint{2.914625in}{1.535783in}}%
\pgfpathlineto{\pgfqpoint{2.909180in}{1.525824in}}%
\pgfpathlineto{\pgfqpoint{2.879160in}{1.518963in}}%
\pgfpathlineto{\pgfqpoint{2.875745in}{1.523338in}}%
\pgfpathlineto{\pgfqpoint{2.871705in}{1.532829in}}%
\pgfpathlineto{\pgfqpoint{2.869055in}{1.532209in}}%
\pgfpathlineto{\pgfqpoint{2.867151in}{1.544262in}}%
\pgfpathclose%
\pgfusepath{fill}%
\end{pgfscope}%
\begin{pgfscope}%
\pgfpathrectangle{\pgfqpoint{0.100000in}{0.100000in}}{\pgfqpoint{3.007045in}{1.925000in}}%
\pgfusepath{clip}%
\pgfsetbuttcap%
\pgfsetmiterjoin%
\definecolor{currentfill}{rgb}{0.460392,0.704744,0.848012}%
\pgfsetfillcolor{currentfill}%
\pgfsetlinewidth{0.000000pt}%
\definecolor{currentstroke}{rgb}{0.000000,0.000000,0.000000}%
\pgfsetstrokecolor{currentstroke}%
\pgfsetstrokeopacity{0.000000}%
\pgfsetdash{}{0pt}%
\pgfpathmoveto{\pgfqpoint{2.082341in}{0.997150in}}%
\pgfpathlineto{\pgfqpoint{2.081317in}{1.001182in}}%
\pgfpathlineto{\pgfqpoint{2.082880in}{1.008816in}}%
\pgfpathlineto{\pgfqpoint{2.081022in}{1.010070in}}%
\pgfpathlineto{\pgfqpoint{2.078167in}{1.013748in}}%
\pgfpathlineto{\pgfqpoint{2.085427in}{1.025478in}}%
\pgfpathlineto{\pgfqpoint{2.090698in}{1.026843in}}%
\pgfpathlineto{\pgfqpoint{2.091233in}{1.034083in}}%
\pgfpathlineto{\pgfqpoint{2.101574in}{1.034888in}}%
\pgfpathlineto{\pgfqpoint{2.100489in}{1.052130in}}%
\pgfpathlineto{\pgfqpoint{2.115921in}{1.053268in}}%
\pgfpathlineto{\pgfqpoint{2.116603in}{1.041678in}}%
\pgfpathlineto{\pgfqpoint{2.119654in}{1.040653in}}%
\pgfpathlineto{\pgfqpoint{2.127608in}{1.032537in}}%
\pgfpathlineto{\pgfqpoint{2.129826in}{1.024933in}}%
\pgfpathlineto{\pgfqpoint{2.127814in}{1.020206in}}%
\pgfpathlineto{\pgfqpoint{2.132817in}{1.006098in}}%
\pgfpathlineto{\pgfqpoint{2.135867in}{0.998378in}}%
\pgfpathlineto{\pgfqpoint{2.116101in}{0.997143in}}%
\pgfpathlineto{\pgfqpoint{2.115261in}{1.009765in}}%
\pgfpathlineto{\pgfqpoint{2.097982in}{1.008842in}}%
\pgfpathlineto{\pgfqpoint{2.098746in}{0.997437in}}%
\pgfpathlineto{\pgfqpoint{2.082341in}{0.997150in}}%
\pgfpathclose%
\pgfusepath{fill}%
\end{pgfscope}%
\begin{pgfscope}%
\pgfpathrectangle{\pgfqpoint{0.100000in}{0.100000in}}{\pgfqpoint{3.007045in}{1.925000in}}%
\pgfusepath{clip}%
\pgfsetbuttcap%
\pgfsetmiterjoin%
\definecolor{currentfill}{rgb}{0.541961,0.749527,0.865606}%
\pgfsetfillcolor{currentfill}%
\pgfsetlinewidth{0.000000pt}%
\definecolor{currentstroke}{rgb}{0.000000,0.000000,0.000000}%
\pgfsetstrokecolor{currentstroke}%
\pgfsetstrokeopacity{0.000000}%
\pgfsetdash{}{0pt}%
\pgfpathmoveto{\pgfqpoint{2.654171in}{1.354581in}}%
\pgfpathlineto{\pgfqpoint{2.645166in}{1.349038in}}%
\pgfpathlineto{\pgfqpoint{2.637385in}{1.345841in}}%
\pgfpathlineto{\pgfqpoint{2.627962in}{1.353837in}}%
\pgfpathlineto{\pgfqpoint{2.623988in}{1.352894in}}%
\pgfpathlineto{\pgfqpoint{2.621003in}{1.357074in}}%
\pgfpathlineto{\pgfqpoint{2.609177in}{1.353242in}}%
\pgfpathlineto{\pgfqpoint{2.605715in}{1.356835in}}%
\pgfpathlineto{\pgfqpoint{2.608964in}{1.367546in}}%
\pgfpathlineto{\pgfqpoint{2.607627in}{1.374496in}}%
\pgfpathlineto{\pgfqpoint{2.626550in}{1.378374in}}%
\pgfpathlineto{\pgfqpoint{2.625680in}{1.382482in}}%
\pgfpathlineto{\pgfqpoint{2.656505in}{1.389182in}}%
\pgfpathlineto{\pgfqpoint{2.660061in}{1.392938in}}%
\pgfpathlineto{\pgfqpoint{2.652249in}{1.418613in}}%
\pgfpathlineto{\pgfqpoint{2.669988in}{1.422124in}}%
\pgfpathlineto{\pgfqpoint{2.689741in}{1.426141in}}%
\pgfpathlineto{\pgfqpoint{2.695962in}{1.404095in}}%
\pgfpathlineto{\pgfqpoint{2.691983in}{1.402891in}}%
\pgfpathlineto{\pgfqpoint{2.692396in}{1.393752in}}%
\pgfpathlineto{\pgfqpoint{2.691171in}{1.379954in}}%
\pgfpathlineto{\pgfqpoint{2.686454in}{1.378941in}}%
\pgfpathlineto{\pgfqpoint{2.679506in}{1.367379in}}%
\pgfpathlineto{\pgfqpoint{2.677192in}{1.366767in}}%
\pgfpathlineto{\pgfqpoint{2.669738in}{1.366595in}}%
\pgfpathlineto{\pgfqpoint{2.661755in}{1.363168in}}%
\pgfpathlineto{\pgfqpoint{2.661650in}{1.359044in}}%
\pgfpathlineto{\pgfqpoint{2.654171in}{1.354581in}}%
\pgfpathclose%
\pgfusepath{fill}%
\end{pgfscope}%
\begin{pgfscope}%
\pgfpathrectangle{\pgfqpoint{0.100000in}{0.100000in}}{\pgfqpoint{3.007045in}{1.925000in}}%
\pgfusepath{clip}%
\pgfsetbuttcap%
\pgfsetmiterjoin%
\definecolor{currentfill}{rgb}{0.523137,0.739193,0.861546}%
\pgfsetfillcolor{currentfill}%
\pgfsetlinewidth{0.000000pt}%
\definecolor{currentstroke}{rgb}{0.000000,0.000000,0.000000}%
\pgfsetstrokecolor{currentstroke}%
\pgfsetstrokeopacity{0.000000}%
\pgfsetdash{}{0pt}%
\pgfpathmoveto{\pgfqpoint{2.245494in}{0.625658in}}%
\pgfpathlineto{\pgfqpoint{2.193761in}{0.621372in}}%
\pgfpathlineto{\pgfqpoint{2.191464in}{0.637662in}}%
\pgfpathlineto{\pgfqpoint{2.185584in}{0.641031in}}%
\pgfpathlineto{\pgfqpoint{2.184550in}{0.645723in}}%
\pgfpathlineto{\pgfqpoint{2.187693in}{0.649271in}}%
\pgfpathlineto{\pgfqpoint{2.192871in}{0.650628in}}%
\pgfpathlineto{\pgfqpoint{2.191786in}{0.667860in}}%
\pgfpathlineto{\pgfqpoint{2.194637in}{0.668142in}}%
\pgfpathlineto{\pgfqpoint{2.194774in}{0.676909in}}%
\pgfpathlineto{\pgfqpoint{2.228144in}{0.679871in}}%
\pgfpathlineto{\pgfqpoint{2.229320in}{0.666780in}}%
\pgfpathlineto{\pgfqpoint{2.233755in}{0.660011in}}%
\pgfpathlineto{\pgfqpoint{2.241536in}{0.660624in}}%
\pgfpathlineto{\pgfqpoint{2.244783in}{0.625988in}}%
\pgfpathlineto{\pgfqpoint{2.245494in}{0.625658in}}%
\pgfpathclose%
\pgfusepath{fill}%
\end{pgfscope}%
\begin{pgfscope}%
\pgfpathrectangle{\pgfqpoint{0.100000in}{0.100000in}}{\pgfqpoint{3.007045in}{1.925000in}}%
\pgfusepath{clip}%
\pgfsetbuttcap%
\pgfsetmiterjoin%
\definecolor{currentfill}{rgb}{0.485490,0.718524,0.853426}%
\pgfsetfillcolor{currentfill}%
\pgfsetlinewidth{0.000000pt}%
\definecolor{currentstroke}{rgb}{0.000000,0.000000,0.000000}%
\pgfsetstrokecolor{currentstroke}%
\pgfsetstrokeopacity{0.000000}%
\pgfsetdash{}{0pt}%
\pgfpathmoveto{\pgfqpoint{1.054598in}{0.801868in}}%
\pgfpathlineto{\pgfqpoint{0.997334in}{0.810949in}}%
\pgfpathlineto{\pgfqpoint{1.003150in}{0.848168in}}%
\pgfpathlineto{\pgfqpoint{1.011345in}{0.900484in}}%
\pgfpathlineto{\pgfqpoint{1.047397in}{0.894955in}}%
\pgfpathlineto{\pgfqpoint{1.082425in}{0.889855in}}%
\pgfpathlineto{\pgfqpoint{1.127835in}{0.883763in}}%
\pgfpathlineto{\pgfqpoint{1.133587in}{0.879798in}}%
\pgfpathlineto{\pgfqpoint{1.151795in}{0.871370in}}%
\pgfpathlineto{\pgfqpoint{1.150364in}{0.859555in}}%
\pgfpathlineto{\pgfqpoint{1.177035in}{0.856282in}}%
\pgfpathlineto{\pgfqpoint{1.173550in}{0.827766in}}%
\pgfpathlineto{\pgfqpoint{1.166565in}{0.828636in}}%
\pgfpathlineto{\pgfqpoint{1.165157in}{0.817109in}}%
\pgfpathlineto{\pgfqpoint{1.147690in}{0.819949in}}%
\pgfpathlineto{\pgfqpoint{1.146247in}{0.808107in}}%
\pgfpathlineto{\pgfqpoint{1.115027in}{0.811851in}}%
\pgfpathlineto{\pgfqpoint{1.060540in}{0.819258in}}%
\pgfpathlineto{\pgfqpoint{1.057206in}{0.819980in}}%
\pgfpathlineto{\pgfqpoint{1.054598in}{0.801868in}}%
\pgfpathclose%
\pgfusepath{fill}%
\end{pgfscope}%
\begin{pgfscope}%
\pgfpathrectangle{\pgfqpoint{0.100000in}{0.100000in}}{\pgfqpoint{3.007045in}{1.925000in}}%
\pgfusepath{clip}%
\pgfsetbuttcap%
\pgfsetmiterjoin%
\definecolor{currentfill}{rgb}{0.579608,0.770196,0.873725}%
\pgfsetfillcolor{currentfill}%
\pgfsetlinewidth{0.000000pt}%
\definecolor{currentstroke}{rgb}{0.000000,0.000000,0.000000}%
\pgfsetstrokecolor{currentstroke}%
\pgfsetstrokeopacity{0.000000}%
\pgfsetdash{}{0pt}%
\pgfpathmoveto{\pgfqpoint{2.246778in}{1.453408in}}%
\pgfpathlineto{\pgfqpoint{2.224008in}{1.451122in}}%
\pgfpathlineto{\pgfqpoint{2.221626in}{1.473931in}}%
\pgfpathlineto{\pgfqpoint{2.199018in}{1.471645in}}%
\pgfpathlineto{\pgfqpoint{2.196560in}{1.494814in}}%
\pgfpathlineto{\pgfqpoint{2.207015in}{1.495613in}}%
\pgfpathlineto{\pgfqpoint{2.204865in}{1.518579in}}%
\pgfpathlineto{\pgfqpoint{2.250233in}{1.523286in}}%
\pgfpathlineto{\pgfqpoint{2.252766in}{1.500298in}}%
\pgfpathlineto{\pgfqpoint{2.241654in}{1.499264in}}%
\pgfpathlineto{\pgfqpoint{2.246778in}{1.453408in}}%
\pgfpathclose%
\pgfusepath{fill}%
\end{pgfscope}%
\begin{pgfscope}%
\pgfpathrectangle{\pgfqpoint{0.100000in}{0.100000in}}{\pgfqpoint{3.007045in}{1.925000in}}%
\pgfusepath{clip}%
\pgfsetbuttcap%
\pgfsetmiterjoin%
\definecolor{currentfill}{rgb}{0.381776,0.656517,0.824452}%
\pgfsetfillcolor{currentfill}%
\pgfsetlinewidth{0.000000pt}%
\definecolor{currentstroke}{rgb}{0.000000,0.000000,0.000000}%
\pgfsetstrokecolor{currentstroke}%
\pgfsetstrokeopacity{0.000000}%
\pgfsetdash{}{0pt}%
\pgfpathmoveto{\pgfqpoint{2.163243in}{0.916924in}}%
\pgfpathlineto{\pgfqpoint{2.166375in}{0.924298in}}%
\pgfpathlineto{\pgfqpoint{2.160890in}{0.928692in}}%
\pgfpathlineto{\pgfqpoint{2.161409in}{0.940303in}}%
\pgfpathlineto{\pgfqpoint{2.161575in}{0.944740in}}%
\pgfpathlineto{\pgfqpoint{2.170276in}{0.951262in}}%
\pgfpathlineto{\pgfqpoint{2.188043in}{0.950490in}}%
\pgfpathlineto{\pgfqpoint{2.188078in}{0.943307in}}%
\pgfpathlineto{\pgfqpoint{2.210327in}{0.936644in}}%
\pgfpathlineto{\pgfqpoint{2.219736in}{0.935214in}}%
\pgfpathlineto{\pgfqpoint{2.220248in}{0.918653in}}%
\pgfpathlineto{\pgfqpoint{2.223823in}{0.914321in}}%
\pgfpathlineto{\pgfqpoint{2.223886in}{0.911812in}}%
\pgfpathlineto{\pgfqpoint{2.216592in}{0.906228in}}%
\pgfpathlineto{\pgfqpoint{2.212159in}{0.906542in}}%
\pgfpathlineto{\pgfqpoint{2.213336in}{0.888623in}}%
\pgfpathlineto{\pgfqpoint{2.193121in}{0.887325in}}%
\pgfpathlineto{\pgfqpoint{2.192088in}{0.890541in}}%
\pgfpathlineto{\pgfqpoint{2.190733in}{0.912451in}}%
\pgfpathlineto{\pgfqpoint{2.189759in}{0.915550in}}%
\pgfpathlineto{\pgfqpoint{2.184378in}{0.916884in}}%
\pgfpathlineto{\pgfqpoint{2.171317in}{0.911958in}}%
\pgfpathlineto{\pgfqpoint{2.167138in}{0.916537in}}%
\pgfpathlineto{\pgfqpoint{2.163243in}{0.916924in}}%
\pgfpathclose%
\pgfusepath{fill}%
\end{pgfscope}%
\begin{pgfscope}%
\pgfpathrectangle{\pgfqpoint{0.100000in}{0.100000in}}{\pgfqpoint{3.007045in}{1.925000in}}%
\pgfusepath{clip}%
\pgfsetbuttcap%
\pgfsetmiterjoin%
\definecolor{currentfill}{rgb}{0.604706,0.783975,0.879139}%
\pgfsetfillcolor{currentfill}%
\pgfsetlinewidth{0.000000pt}%
\definecolor{currentstroke}{rgb}{0.000000,0.000000,0.000000}%
\pgfsetstrokecolor{currentstroke}%
\pgfsetstrokeopacity{0.000000}%
\pgfsetdash{}{0pt}%
\pgfpathmoveto{\pgfqpoint{2.525737in}{1.174087in}}%
\pgfpathlineto{\pgfqpoint{2.532425in}{1.163870in}}%
\pgfpathlineto{\pgfqpoint{2.530752in}{1.154455in}}%
\pgfpathlineto{\pgfqpoint{2.526995in}{1.150325in}}%
\pgfpathlineto{\pgfqpoint{2.526634in}{1.142930in}}%
\pgfpathlineto{\pgfqpoint{2.522292in}{1.144995in}}%
\pgfpathlineto{\pgfqpoint{2.510167in}{1.130498in}}%
\pgfpathlineto{\pgfqpoint{2.501225in}{1.130708in}}%
\pgfpathlineto{\pgfqpoint{2.494710in}{1.137649in}}%
\pgfpathlineto{\pgfqpoint{2.490413in}{1.135206in}}%
\pgfpathlineto{\pgfqpoint{2.481701in}{1.138352in}}%
\pgfpathlineto{\pgfqpoint{2.496539in}{1.150737in}}%
\pgfpathlineto{\pgfqpoint{2.497293in}{1.156805in}}%
\pgfpathlineto{\pgfqpoint{2.494103in}{1.161327in}}%
\pgfpathlineto{\pgfqpoint{2.487851in}{1.166255in}}%
\pgfpathlineto{\pgfqpoint{2.489719in}{1.169881in}}%
\pgfpathlineto{\pgfqpoint{2.485316in}{1.174226in}}%
\pgfpathlineto{\pgfqpoint{2.487519in}{1.178106in}}%
\pgfpathlineto{\pgfqpoint{2.484306in}{1.188484in}}%
\pgfpathlineto{\pgfqpoint{2.489753in}{1.191584in}}%
\pgfpathlineto{\pgfqpoint{2.494442in}{1.196721in}}%
\pgfpathlineto{\pgfqpoint{2.498857in}{1.196503in}}%
\pgfpathlineto{\pgfqpoint{2.500337in}{1.190919in}}%
\pgfpathlineto{\pgfqpoint{2.507565in}{1.184564in}}%
\pgfpathlineto{\pgfqpoint{2.514455in}{1.181489in}}%
\pgfpathlineto{\pgfqpoint{2.517598in}{1.175060in}}%
\pgfpathlineto{\pgfqpoint{2.525737in}{1.174087in}}%
\pgfpathclose%
\pgfusepath{fill}%
\end{pgfscope}%
\begin{pgfscope}%
\pgfpathrectangle{\pgfqpoint{0.100000in}{0.100000in}}{\pgfqpoint{3.007045in}{1.925000in}}%
\pgfusepath{clip}%
\pgfsetbuttcap%
\pgfsetmiterjoin%
\definecolor{currentfill}{rgb}{0.231926,0.545652,0.762614}%
\pgfsetfillcolor{currentfill}%
\pgfsetlinewidth{0.000000pt}%
\definecolor{currentstroke}{rgb}{0.000000,0.000000,0.000000}%
\pgfsetstrokecolor{currentstroke}%
\pgfsetstrokeopacity{0.000000}%
\pgfsetdash{}{0pt}%
\pgfpathmoveto{\pgfqpoint{0.836260in}{1.093783in}}%
\pgfpathlineto{\pgfqpoint{0.776912in}{1.106019in}}%
\pgfpathlineto{\pgfqpoint{0.788118in}{1.158139in}}%
\pgfpathlineto{\pgfqpoint{0.798495in}{1.207353in}}%
\pgfpathlineto{\pgfqpoint{0.850288in}{1.196720in}}%
\pgfpathlineto{\pgfqpoint{0.875936in}{1.191519in}}%
\pgfpathlineto{\pgfqpoint{0.875518in}{1.183498in}}%
\pgfpathlineto{\pgfqpoint{0.878976in}{1.180671in}}%
\pgfpathlineto{\pgfqpoint{0.882037in}{1.174273in}}%
\pgfpathlineto{\pgfqpoint{0.877415in}{1.166653in}}%
\pgfpathlineto{\pgfqpoint{0.872431in}{1.163681in}}%
\pgfpathlineto{\pgfqpoint{0.869732in}{1.146966in}}%
\pgfpathlineto{\pgfqpoint{0.864002in}{1.147526in}}%
\pgfpathlineto{\pgfqpoint{0.863012in}{1.142510in}}%
\pgfpathlineto{\pgfqpoint{0.857375in}{1.143704in}}%
\pgfpathlineto{\pgfqpoint{0.854337in}{1.126666in}}%
\pgfpathlineto{\pgfqpoint{0.843170in}{1.128894in}}%
\pgfpathlineto{\pgfqpoint{0.836260in}{1.093783in}}%
\pgfpathclose%
\pgfusepath{fill}%
\end{pgfscope}%
\begin{pgfscope}%
\pgfpathrectangle{\pgfqpoint{0.100000in}{0.100000in}}{\pgfqpoint{3.007045in}{1.925000in}}%
\pgfusepath{clip}%
\pgfsetbuttcap%
\pgfsetmiterjoin%
\definecolor{currentfill}{rgb}{0.260715,0.573841,0.777209}%
\pgfsetfillcolor{currentfill}%
\pgfsetlinewidth{0.000000pt}%
\definecolor{currentstroke}{rgb}{0.000000,0.000000,0.000000}%
\pgfsetstrokecolor{currentstroke}%
\pgfsetstrokeopacity{0.000000}%
\pgfsetdash{}{0pt}%
\pgfpathmoveto{\pgfqpoint{2.253561in}{1.079316in}}%
\pgfpathlineto{\pgfqpoint{2.259515in}{1.072291in}}%
\pgfpathlineto{\pgfqpoint{2.261909in}{1.072347in}}%
\pgfpathlineto{\pgfqpoint{2.262113in}{1.062729in}}%
\pgfpathlineto{\pgfqpoint{2.267818in}{1.063028in}}%
\pgfpathlineto{\pgfqpoint{2.270021in}{1.058081in}}%
\pgfpathlineto{\pgfqpoint{2.261228in}{1.057564in}}%
\pgfpathlineto{\pgfqpoint{2.260372in}{1.054200in}}%
\pgfpathlineto{\pgfqpoint{2.239874in}{1.053973in}}%
\pgfpathlineto{\pgfqpoint{2.238225in}{1.047271in}}%
\pgfpathlineto{\pgfqpoint{2.234875in}{1.045920in}}%
\pgfpathlineto{\pgfqpoint{2.218727in}{1.043400in}}%
\pgfpathlineto{\pgfqpoint{2.217524in}{1.046026in}}%
\pgfpathlineto{\pgfqpoint{2.210809in}{1.047577in}}%
\pgfpathlineto{\pgfqpoint{2.207346in}{1.056620in}}%
\pgfpathlineto{\pgfqpoint{2.208703in}{1.064571in}}%
\pgfpathlineto{\pgfqpoint{2.204976in}{1.073343in}}%
\pgfpathlineto{\pgfqpoint{2.205852in}{1.076709in}}%
\pgfpathlineto{\pgfqpoint{2.208829in}{1.082188in}}%
\pgfpathlineto{\pgfqpoint{2.212960in}{1.083080in}}%
\pgfpathlineto{\pgfqpoint{2.211772in}{1.090499in}}%
\pgfpathlineto{\pgfqpoint{2.214133in}{1.095845in}}%
\pgfpathlineto{\pgfqpoint{2.217176in}{1.095188in}}%
\pgfpathlineto{\pgfqpoint{2.219675in}{1.101261in}}%
\pgfpathlineto{\pgfqpoint{2.223926in}{1.097599in}}%
\pgfpathlineto{\pgfqpoint{2.224931in}{1.092417in}}%
\pgfpathlineto{\pgfqpoint{2.229804in}{1.090057in}}%
\pgfpathlineto{\pgfqpoint{2.241377in}{1.090992in}}%
\pgfpathlineto{\pgfqpoint{2.246327in}{1.087565in}}%
\pgfpathlineto{\pgfqpoint{2.249127in}{1.081944in}}%
\pgfpathlineto{\pgfqpoint{2.253561in}{1.079316in}}%
\pgfpathclose%
\pgfusepath{fill}%
\end{pgfscope}%
\begin{pgfscope}%
\pgfpathrectangle{\pgfqpoint{0.100000in}{0.100000in}}{\pgfqpoint{3.007045in}{1.925000in}}%
\pgfusepath{clip}%
\pgfsetbuttcap%
\pgfsetmiterjoin%
\definecolor{currentfill}{rgb}{0.296025,0.597955,0.790988}%
\pgfsetfillcolor{currentfill}%
\pgfsetlinewidth{0.000000pt}%
\definecolor{currentstroke}{rgb}{0.000000,0.000000,0.000000}%
\pgfsetstrokecolor{currentstroke}%
\pgfsetstrokeopacity{0.000000}%
\pgfsetdash{}{0pt}%
\pgfpathmoveto{\pgfqpoint{1.596789in}{1.798977in}}%
\pgfpathlineto{\pgfqpoint{1.595955in}{1.780137in}}%
\pgfpathlineto{\pgfqpoint{1.596746in}{1.768470in}}%
\pgfpathlineto{\pgfqpoint{1.596215in}{1.756877in}}%
\pgfpathlineto{\pgfqpoint{1.573086in}{1.757987in}}%
\pgfpathlineto{\pgfqpoint{1.573646in}{1.769590in}}%
\pgfpathlineto{\pgfqpoint{1.544810in}{1.771147in}}%
\pgfpathlineto{\pgfqpoint{1.543838in}{1.782832in}}%
\pgfpathlineto{\pgfqpoint{1.544920in}{1.801546in}}%
\pgfpathlineto{\pgfqpoint{1.596789in}{1.798977in}}%
\pgfpathclose%
\pgfusepath{fill}%
\end{pgfscope}%
\begin{pgfscope}%
\pgfpathrectangle{\pgfqpoint{0.100000in}{0.100000in}}{\pgfqpoint{3.007045in}{1.925000in}}%
\pgfusepath{clip}%
\pgfsetbuttcap%
\pgfsetmiterjoin%
\definecolor{currentfill}{rgb}{0.290980,0.594510,0.789020}%
\pgfsetfillcolor{currentfill}%
\pgfsetlinewidth{0.000000pt}%
\definecolor{currentstroke}{rgb}{0.000000,0.000000,0.000000}%
\pgfsetstrokecolor{currentstroke}%
\pgfsetstrokeopacity{0.000000}%
\pgfsetdash{}{0pt}%
\pgfpathmoveto{\pgfqpoint{1.694674in}{0.889113in}}%
\pgfpathlineto{\pgfqpoint{1.694493in}{0.877607in}}%
\pgfpathlineto{\pgfqpoint{1.691843in}{0.877659in}}%
\pgfpathlineto{\pgfqpoint{1.691602in}{0.864102in}}%
\pgfpathlineto{\pgfqpoint{1.683133in}{0.863693in}}%
\pgfpathlineto{\pgfqpoint{1.678967in}{0.860772in}}%
\pgfpathlineto{\pgfqpoint{1.674047in}{0.866744in}}%
\pgfpathlineto{\pgfqpoint{1.670026in}{0.865689in}}%
\pgfpathlineto{\pgfqpoint{1.667600in}{0.860848in}}%
\pgfpathlineto{\pgfqpoint{1.627761in}{0.861901in}}%
\pgfpathlineto{\pgfqpoint{1.628500in}{0.893638in}}%
\pgfpathlineto{\pgfqpoint{1.619336in}{0.893761in}}%
\pgfpathlineto{\pgfqpoint{1.609800in}{0.896970in}}%
\pgfpathlineto{\pgfqpoint{1.594438in}{0.897482in}}%
\pgfpathlineto{\pgfqpoint{1.594488in}{0.908994in}}%
\pgfpathlineto{\pgfqpoint{1.594900in}{0.920448in}}%
\pgfpathlineto{\pgfqpoint{1.600584in}{0.920240in}}%
\pgfpathlineto{\pgfqpoint{1.601428in}{0.949236in}}%
\pgfpathlineto{\pgfqpoint{1.641329in}{0.947962in}}%
\pgfpathlineto{\pgfqpoint{1.647016in}{0.947800in}}%
\pgfpathlineto{\pgfqpoint{1.646710in}{0.935974in}}%
\pgfpathlineto{\pgfqpoint{1.658027in}{0.932849in}}%
\pgfpathlineto{\pgfqpoint{1.685760in}{0.932254in}}%
\pgfpathlineto{\pgfqpoint{1.685427in}{0.912299in}}%
\pgfpathlineto{\pgfqpoint{1.684983in}{0.896595in}}%
\pgfpathlineto{\pgfqpoint{1.691799in}{0.898304in}}%
\pgfpathlineto{\pgfqpoint{1.694674in}{0.889113in}}%
\pgfpathclose%
\pgfusepath{fill}%
\end{pgfscope}%
\begin{pgfscope}%
\pgfpathrectangle{\pgfqpoint{0.100000in}{0.100000in}}{\pgfqpoint{3.007045in}{1.925000in}}%
\pgfusepath{clip}%
\pgfsetbuttcap%
\pgfsetmiterjoin%
\definecolor{currentfill}{rgb}{0.412042,0.677186,0.836263}%
\pgfsetfillcolor{currentfill}%
\pgfsetlinewidth{0.000000pt}%
\definecolor{currentstroke}{rgb}{0.000000,0.000000,0.000000}%
\pgfsetstrokecolor{currentstroke}%
\pgfsetstrokeopacity{0.000000}%
\pgfsetdash{}{0pt}%
\pgfpathmoveto{\pgfqpoint{1.776405in}{1.429765in}}%
\pgfpathlineto{\pgfqpoint{1.779303in}{1.429779in}}%
\pgfpathlineto{\pgfqpoint{1.779266in}{1.452736in}}%
\pgfpathlineto{\pgfqpoint{1.831065in}{1.453126in}}%
\pgfpathlineto{\pgfqpoint{1.836747in}{1.453219in}}%
\pgfpathlineto{\pgfqpoint{1.837085in}{1.430186in}}%
\pgfpathlineto{\pgfqpoint{1.821649in}{1.429987in}}%
\pgfpathlineto{\pgfqpoint{1.799028in}{1.429855in}}%
\pgfpathlineto{\pgfqpoint{1.799116in}{1.413633in}}%
\pgfpathlineto{\pgfqpoint{1.776471in}{1.413556in}}%
\pgfpathlineto{\pgfqpoint{1.776405in}{1.429765in}}%
\pgfpathclose%
\pgfusepath{fill}%
\end{pgfscope}%
\begin{pgfscope}%
\pgfpathrectangle{\pgfqpoint{0.100000in}{0.100000in}}{\pgfqpoint{3.007045in}{1.925000in}}%
\pgfusepath{clip}%
\pgfsetbuttcap%
\pgfsetmiterjoin%
\definecolor{currentfill}{rgb}{0.391865,0.663406,0.828389}%
\pgfsetfillcolor{currentfill}%
\pgfsetlinewidth{0.000000pt}%
\definecolor{currentstroke}{rgb}{0.000000,0.000000,0.000000}%
\pgfsetstrokecolor{currentstroke}%
\pgfsetstrokeopacity{0.000000}%
\pgfsetdash{}{0pt}%
\pgfpathmoveto{\pgfqpoint{2.353818in}{1.009495in}}%
\pgfpathlineto{\pgfqpoint{2.341210in}{1.008147in}}%
\pgfpathlineto{\pgfqpoint{2.336554in}{1.012826in}}%
\pgfpathlineto{\pgfqpoint{2.332746in}{1.021710in}}%
\pgfpathlineto{\pgfqpoint{2.334697in}{1.028538in}}%
\pgfpathlineto{\pgfqpoint{2.332415in}{1.038665in}}%
\pgfpathlineto{\pgfqpoint{2.333634in}{1.044388in}}%
\pgfpathlineto{\pgfqpoint{2.334876in}{1.050418in}}%
\pgfpathlineto{\pgfqpoint{2.340592in}{1.057330in}}%
\pgfpathlineto{\pgfqpoint{2.347601in}{1.052815in}}%
\pgfpathlineto{\pgfqpoint{2.350874in}{1.053026in}}%
\pgfpathlineto{\pgfqpoint{2.358324in}{1.060460in}}%
\pgfpathlineto{\pgfqpoint{2.363870in}{1.060375in}}%
\pgfpathlineto{\pgfqpoint{2.372484in}{1.056245in}}%
\pgfpathlineto{\pgfqpoint{2.371954in}{1.050139in}}%
\pgfpathlineto{\pgfqpoint{2.376071in}{1.035433in}}%
\pgfpathlineto{\pgfqpoint{2.362675in}{1.024583in}}%
\pgfpathlineto{\pgfqpoint{2.360603in}{1.018528in}}%
\pgfpathlineto{\pgfqpoint{2.355024in}{1.013333in}}%
\pgfpathlineto{\pgfqpoint{2.353818in}{1.009495in}}%
\pgfpathclose%
\pgfusepath{fill}%
\end{pgfscope}%
\begin{pgfscope}%
\pgfpathrectangle{\pgfqpoint{0.100000in}{0.100000in}}{\pgfqpoint{3.007045in}{1.925000in}}%
\pgfusepath{clip}%
\pgfsetbuttcap%
\pgfsetmiterjoin%
\definecolor{currentfill}{rgb}{0.351511,0.635848,0.812641}%
\pgfsetfillcolor{currentfill}%
\pgfsetlinewidth{0.000000pt}%
\definecolor{currentstroke}{rgb}{0.000000,0.000000,0.000000}%
\pgfsetstrokecolor{currentstroke}%
\pgfsetstrokeopacity{0.000000}%
\pgfsetdash{}{0pt}%
\pgfpathmoveto{\pgfqpoint{2.321129in}{1.054650in}}%
\pgfpathlineto{\pgfqpoint{2.318137in}{1.056117in}}%
\pgfpathlineto{\pgfqpoint{2.310463in}{1.047437in}}%
\pgfpathlineto{\pgfqpoint{2.309458in}{1.055417in}}%
\pgfpathlineto{\pgfqpoint{2.302296in}{1.059040in}}%
\pgfpathlineto{\pgfqpoint{2.301227in}{1.065647in}}%
\pgfpathlineto{\pgfqpoint{2.291623in}{1.065687in}}%
\pgfpathlineto{\pgfqpoint{2.291416in}{1.086166in}}%
\pgfpathlineto{\pgfqpoint{2.289598in}{1.088435in}}%
\pgfpathlineto{\pgfqpoint{2.294182in}{1.093073in}}%
\pgfpathlineto{\pgfqpoint{2.300122in}{1.093676in}}%
\pgfpathlineto{\pgfqpoint{2.314503in}{1.080433in}}%
\pgfpathlineto{\pgfqpoint{2.316608in}{1.083386in}}%
\pgfpathlineto{\pgfqpoint{2.327579in}{1.069371in}}%
\pgfpathlineto{\pgfqpoint{2.322848in}{1.065287in}}%
\pgfpathlineto{\pgfqpoint{2.321129in}{1.054650in}}%
\pgfpathclose%
\pgfusepath{fill}%
\end{pgfscope}%
\begin{pgfscope}%
\pgfpathrectangle{\pgfqpoint{0.100000in}{0.100000in}}{\pgfqpoint{3.007045in}{1.925000in}}%
\pgfusepath{clip}%
\pgfsetbuttcap%
\pgfsetmiterjoin%
\definecolor{currentfill}{rgb}{0.401953,0.670296,0.832326}%
\pgfsetfillcolor{currentfill}%
\pgfsetlinewidth{0.000000pt}%
\definecolor{currentstroke}{rgb}{0.000000,0.000000,0.000000}%
\pgfsetstrokecolor{currentstroke}%
\pgfsetstrokeopacity{0.000000}%
\pgfsetdash{}{0pt}%
\pgfpathmoveto{\pgfqpoint{1.418635in}{1.235643in}}%
\pgfpathlineto{\pgfqpoint{1.414767in}{1.184726in}}%
\pgfpathlineto{\pgfqpoint{1.358836in}{1.188865in}}%
\pgfpathlineto{\pgfqpoint{1.330895in}{1.191503in}}%
\pgfpathlineto{\pgfqpoint{1.333705in}{1.220024in}}%
\pgfpathlineto{\pgfqpoint{1.345447in}{1.218910in}}%
\pgfpathlineto{\pgfqpoint{1.348876in}{1.253120in}}%
\pgfpathlineto{\pgfqpoint{1.343169in}{1.253667in}}%
\pgfpathlineto{\pgfqpoint{1.346503in}{1.285013in}}%
\pgfpathlineto{\pgfqpoint{1.392134in}{1.280968in}}%
\pgfpathlineto{\pgfqpoint{1.390844in}{1.264339in}}%
\pgfpathlineto{\pgfqpoint{1.420694in}{1.261942in}}%
\pgfpathlineto{\pgfqpoint{1.420426in}{1.258524in}}%
\pgfpathlineto{\pgfqpoint{1.418635in}{1.235643in}}%
\pgfpathclose%
\pgfusepath{fill}%
\end{pgfscope}%
\begin{pgfscope}%
\pgfpathrectangle{\pgfqpoint{0.100000in}{0.100000in}}{\pgfqpoint{3.007045in}{1.925000in}}%
\pgfusepath{clip}%
\pgfsetbuttcap%
\pgfsetmiterjoin%
\definecolor{currentfill}{rgb}{0.306113,0.604844,0.794925}%
\pgfsetfillcolor{currentfill}%
\pgfsetlinewidth{0.000000pt}%
\definecolor{currentstroke}{rgb}{0.000000,0.000000,0.000000}%
\pgfsetstrokecolor{currentstroke}%
\pgfsetstrokeopacity{0.000000}%
\pgfsetdash{}{0pt}%
\pgfpathmoveto{\pgfqpoint{1.555230in}{0.797022in}}%
\pgfpathlineto{\pgfqpoint{1.553928in}{0.768166in}}%
\pgfpathlineto{\pgfqpoint{1.525314in}{0.769593in}}%
\pgfpathlineto{\pgfqpoint{1.526290in}{0.791703in}}%
\pgfpathlineto{\pgfqpoint{1.526556in}{0.798349in}}%
\pgfpathlineto{\pgfqpoint{1.555230in}{0.797022in}}%
\pgfpathclose%
\pgfusepath{fill}%
\end{pgfscope}%
\begin{pgfscope}%
\pgfpathrectangle{\pgfqpoint{0.100000in}{0.100000in}}{\pgfqpoint{3.007045in}{1.925000in}}%
\pgfusepath{clip}%
\pgfsetbuttcap%
\pgfsetmiterjoin%
\definecolor{currentfill}{rgb}{0.366644,0.646182,0.818547}%
\pgfsetfillcolor{currentfill}%
\pgfsetlinewidth{0.000000pt}%
\definecolor{currentstroke}{rgb}{0.000000,0.000000,0.000000}%
\pgfsetstrokecolor{currentstroke}%
\pgfsetstrokeopacity{0.000000}%
\pgfsetdash{}{0pt}%
\pgfpathmoveto{\pgfqpoint{1.882047in}{1.149212in}}%
\pgfpathlineto{\pgfqpoint{1.877465in}{1.141048in}}%
\pgfpathlineto{\pgfqpoint{1.877944in}{1.138491in}}%
\pgfpathlineto{\pgfqpoint{1.872251in}{1.132295in}}%
\pgfpathlineto{\pgfqpoint{1.872212in}{1.129384in}}%
\pgfpathlineto{\pgfqpoint{1.849349in}{1.129915in}}%
\pgfpathlineto{\pgfqpoint{1.848856in}{1.148030in}}%
\pgfpathlineto{\pgfqpoint{1.840824in}{1.149749in}}%
\pgfpathlineto{\pgfqpoint{1.835689in}{1.147108in}}%
\pgfpathlineto{\pgfqpoint{1.835366in}{1.168017in}}%
\pgfpathlineto{\pgfqpoint{1.834849in}{1.196667in}}%
\pgfpathlineto{\pgfqpoint{1.855015in}{1.197513in}}%
\pgfpathlineto{\pgfqpoint{1.855189in}{1.180142in}}%
\pgfpathlineto{\pgfqpoint{1.859241in}{1.178993in}}%
\pgfpathlineto{\pgfqpoint{1.860786in}{1.170457in}}%
\pgfpathlineto{\pgfqpoint{1.859930in}{1.165889in}}%
\pgfpathlineto{\pgfqpoint{1.862920in}{1.162275in}}%
\pgfpathlineto{\pgfqpoint{1.867949in}{1.161555in}}%
\pgfpathlineto{\pgfqpoint{1.870359in}{1.156300in}}%
\pgfpathlineto{\pgfqpoint{1.875772in}{1.154964in}}%
\pgfpathlineto{\pgfqpoint{1.877817in}{1.149309in}}%
\pgfpathlineto{\pgfqpoint{1.882047in}{1.149212in}}%
\pgfpathclose%
\pgfusepath{fill}%
\end{pgfscope}%
\begin{pgfscope}%
\pgfpathrectangle{\pgfqpoint{0.100000in}{0.100000in}}{\pgfqpoint{3.007045in}{1.925000in}}%
\pgfusepath{clip}%
\pgfsetbuttcap%
\pgfsetmiterjoin%
\definecolor{currentfill}{rgb}{0.336378,0.625513,0.806736}%
\pgfsetfillcolor{currentfill}%
\pgfsetlinewidth{0.000000pt}%
\definecolor{currentstroke}{rgb}{0.000000,0.000000,0.000000}%
\pgfsetstrokecolor{currentstroke}%
\pgfsetstrokeopacity{0.000000}%
\pgfsetdash{}{0pt}%
\pgfpathmoveto{\pgfqpoint{2.422257in}{1.215672in}}%
\pgfpathlineto{\pgfqpoint{2.422680in}{1.209735in}}%
\pgfpathlineto{\pgfqpoint{2.416942in}{1.209349in}}%
\pgfpathlineto{\pgfqpoint{2.417290in}{1.203518in}}%
\pgfpathlineto{\pgfqpoint{2.405836in}{1.202691in}}%
\pgfpathlineto{\pgfqpoint{2.405917in}{1.200954in}}%
\pgfpathlineto{\pgfqpoint{2.394275in}{1.199273in}}%
\pgfpathlineto{\pgfqpoint{2.393794in}{1.205858in}}%
\pgfpathlineto{\pgfqpoint{2.381162in}{1.206723in}}%
\pgfpathlineto{\pgfqpoint{2.366865in}{1.205475in}}%
\pgfpathlineto{\pgfqpoint{2.366064in}{1.217273in}}%
\pgfpathlineto{\pgfqpoint{2.345734in}{1.216099in}}%
\pgfpathlineto{\pgfqpoint{2.348588in}{1.219904in}}%
\pgfpathlineto{\pgfqpoint{2.350594in}{1.236213in}}%
\pgfpathlineto{\pgfqpoint{2.349506in}{1.250356in}}%
\pgfpathlineto{\pgfqpoint{2.346617in}{1.250238in}}%
\pgfpathlineto{\pgfqpoint{2.345898in}{1.268387in}}%
\pgfpathlineto{\pgfqpoint{2.350943in}{1.269110in}}%
\pgfpathlineto{\pgfqpoint{2.359392in}{1.270321in}}%
\pgfpathlineto{\pgfqpoint{2.359900in}{1.266193in}}%
\pgfpathlineto{\pgfqpoint{2.376250in}{1.266421in}}%
\pgfpathlineto{\pgfqpoint{2.376485in}{1.262685in}}%
\pgfpathlineto{\pgfqpoint{2.390704in}{1.263733in}}%
\pgfpathlineto{\pgfqpoint{2.389048in}{1.277134in}}%
\pgfpathlineto{\pgfqpoint{2.402337in}{1.279030in}}%
\pgfpathlineto{\pgfqpoint{2.409849in}{1.281241in}}%
\pgfpathlineto{\pgfqpoint{2.411854in}{1.281854in}}%
\pgfpathlineto{\pgfqpoint{2.414093in}{1.259940in}}%
\pgfpathlineto{\pgfqpoint{2.415184in}{1.255293in}}%
\pgfpathlineto{\pgfqpoint{2.416625in}{1.241076in}}%
\pgfpathlineto{\pgfqpoint{2.415203in}{1.238420in}}%
\pgfpathlineto{\pgfqpoint{2.403598in}{1.237933in}}%
\pgfpathlineto{\pgfqpoint{2.403956in}{1.231787in}}%
\pgfpathlineto{\pgfqpoint{2.407820in}{1.232038in}}%
\pgfpathlineto{\pgfqpoint{2.410154in}{1.226352in}}%
\pgfpathlineto{\pgfqpoint{2.410789in}{1.216745in}}%
\pgfpathlineto{\pgfqpoint{2.422257in}{1.215672in}}%
\pgfpathclose%
\pgfusepath{fill}%
\end{pgfscope}%
\begin{pgfscope}%
\pgfpathrectangle{\pgfqpoint{0.100000in}{0.100000in}}{\pgfqpoint{3.007045in}{1.925000in}}%
\pgfusepath{clip}%
\pgfsetbuttcap%
\pgfsetmiterjoin%
\definecolor{currentfill}{rgb}{0.447843,0.697855,0.845306}%
\pgfsetfillcolor{currentfill}%
\pgfsetlinewidth{0.000000pt}%
\definecolor{currentstroke}{rgb}{0.000000,0.000000,0.000000}%
\pgfsetstrokecolor{currentstroke}%
\pgfsetstrokeopacity{0.000000}%
\pgfsetdash{}{0pt}%
\pgfpathmoveto{\pgfqpoint{1.792177in}{1.207263in}}%
\pgfpathlineto{\pgfqpoint{1.792476in}{1.224430in}}%
\pgfpathlineto{\pgfqpoint{1.790743in}{1.236578in}}%
\pgfpathlineto{\pgfqpoint{1.798792in}{1.236593in}}%
\pgfpathlineto{\pgfqpoint{1.833661in}{1.237372in}}%
\pgfpathlineto{\pgfqpoint{1.834576in}{1.208087in}}%
\pgfpathlineto{\pgfqpoint{1.811789in}{1.208062in}}%
\pgfpathlineto{\pgfqpoint{1.811728in}{1.201377in}}%
\pgfpathlineto{\pgfqpoint{1.792230in}{1.201536in}}%
\pgfpathlineto{\pgfqpoint{1.792177in}{1.207263in}}%
\pgfpathclose%
\pgfusepath{fill}%
\end{pgfscope}%
\begin{pgfscope}%
\pgfpathrectangle{\pgfqpoint{0.100000in}{0.100000in}}{\pgfqpoint{3.007045in}{1.925000in}}%
\pgfusepath{clip}%
\pgfsetbuttcap%
\pgfsetmiterjoin%
\definecolor{currentfill}{rgb}{0.447843,0.697855,0.845306}%
\pgfsetfillcolor{currentfill}%
\pgfsetlinewidth{0.000000pt}%
\definecolor{currentstroke}{rgb}{0.000000,0.000000,0.000000}%
\pgfsetstrokecolor{currentstroke}%
\pgfsetstrokeopacity{0.000000}%
\pgfsetdash{}{0pt}%
\pgfpathmoveto{\pgfqpoint{1.744270in}{1.498734in}}%
\pgfpathlineto{\pgfqpoint{1.743883in}{1.504541in}}%
\pgfpathlineto{\pgfqpoint{1.703956in}{1.505078in}}%
\pgfpathlineto{\pgfqpoint{1.704306in}{1.528001in}}%
\pgfpathlineto{\pgfqpoint{1.704575in}{1.547315in}}%
\pgfpathlineto{\pgfqpoint{1.712446in}{1.545741in}}%
\pgfpathlineto{\pgfqpoint{1.717578in}{1.542068in}}%
\pgfpathlineto{\pgfqpoint{1.720167in}{1.545159in}}%
\pgfpathlineto{\pgfqpoint{1.720360in}{1.556564in}}%
\pgfpathlineto{\pgfqpoint{1.766199in}{1.556278in}}%
\pgfpathlineto{\pgfqpoint{1.783347in}{1.556314in}}%
\pgfpathlineto{\pgfqpoint{1.783351in}{1.550568in}}%
\pgfpathlineto{\pgfqpoint{1.783670in}{1.521817in}}%
\pgfpathlineto{\pgfqpoint{1.749650in}{1.521778in}}%
\pgfpathlineto{\pgfqpoint{1.749628in}{1.512662in}}%
\pgfpathlineto{\pgfqpoint{1.755381in}{1.509019in}}%
\pgfpathlineto{\pgfqpoint{1.755330in}{1.498635in}}%
\pgfpathlineto{\pgfqpoint{1.744270in}{1.498734in}}%
\pgfpathclose%
\pgfusepath{fill}%
\end{pgfscope}%
\begin{pgfscope}%
\pgfpathrectangle{\pgfqpoint{0.100000in}{0.100000in}}{\pgfqpoint{3.007045in}{1.925000in}}%
\pgfusepath{clip}%
\pgfsetbuttcap%
\pgfsetmiterjoin%
\definecolor{currentfill}{rgb}{0.460392,0.704744,0.848012}%
\pgfsetfillcolor{currentfill}%
\pgfsetlinewidth{0.000000pt}%
\definecolor{currentstroke}{rgb}{0.000000,0.000000,0.000000}%
\pgfsetstrokecolor{currentstroke}%
\pgfsetstrokeopacity{0.000000}%
\pgfsetdash{}{0pt}%
\pgfpathmoveto{\pgfqpoint{1.890280in}{0.989939in}}%
\pgfpathlineto{\pgfqpoint{1.890382in}{0.969613in}}%
\pgfpathlineto{\pgfqpoint{1.886046in}{0.969526in}}%
\pgfpathlineto{\pgfqpoint{1.832387in}{0.968604in}}%
\pgfpathlineto{\pgfqpoint{1.821240in}{0.968478in}}%
\pgfpathlineto{\pgfqpoint{1.821509in}{0.984906in}}%
\pgfpathlineto{\pgfqpoint{1.822007in}{1.020648in}}%
\pgfpathlineto{\pgfqpoint{1.844374in}{1.020395in}}%
\pgfpathlineto{\pgfqpoint{1.845462in}{1.020387in}}%
\pgfpathlineto{\pgfqpoint{1.845487in}{1.001557in}}%
\pgfpathlineto{\pgfqpoint{1.859842in}{1.001580in}}%
\pgfpathlineto{\pgfqpoint{1.859810in}{0.990030in}}%
\pgfpathlineto{\pgfqpoint{1.890280in}{0.989939in}}%
\pgfpathclose%
\pgfusepath{fill}%
\end{pgfscope}%
\begin{pgfscope}%
\pgfpathrectangle{\pgfqpoint{0.100000in}{0.100000in}}{\pgfqpoint{3.007045in}{1.925000in}}%
\pgfusepath{clip}%
\pgfsetbuttcap%
\pgfsetmiterjoin%
\definecolor{currentfill}{rgb}{0.491765,0.721968,0.854779}%
\pgfsetfillcolor{currentfill}%
\pgfsetlinewidth{0.000000pt}%
\definecolor{currentstroke}{rgb}{0.000000,0.000000,0.000000}%
\pgfsetstrokecolor{currentstroke}%
\pgfsetstrokeopacity{0.000000}%
\pgfsetdash{}{0pt}%
\pgfpathmoveto{\pgfqpoint{2.285799in}{1.004151in}}%
\pgfpathlineto{\pgfqpoint{2.311789in}{1.005478in}}%
\pgfpathlineto{\pgfqpoint{2.315188in}{1.000641in}}%
\pgfpathlineto{\pgfqpoint{2.303650in}{1.003101in}}%
\pgfpathlineto{\pgfqpoint{2.298509in}{1.001042in}}%
\pgfpathlineto{\pgfqpoint{2.295566in}{0.990662in}}%
\pgfpathlineto{\pgfqpoint{2.296204in}{0.983928in}}%
\pgfpathlineto{\pgfqpoint{2.299083in}{0.977051in}}%
\pgfpathlineto{\pgfqpoint{2.296056in}{0.970309in}}%
\pgfpathlineto{\pgfqpoint{2.291608in}{0.970240in}}%
\pgfpathlineto{\pgfqpoint{2.291000in}{0.961775in}}%
\pgfpathlineto{\pgfqpoint{2.285629in}{0.965628in}}%
\pgfpathlineto{\pgfqpoint{2.277359in}{0.967152in}}%
\pgfpathlineto{\pgfqpoint{2.269975in}{0.962101in}}%
\pgfpathlineto{\pgfqpoint{2.267876in}{0.965337in}}%
\pgfpathlineto{\pgfqpoint{2.263169in}{0.965439in}}%
\pgfpathlineto{\pgfqpoint{2.261147in}{0.968764in}}%
\pgfpathlineto{\pgfqpoint{2.261071in}{0.973710in}}%
\pgfpathlineto{\pgfqpoint{2.258038in}{0.978829in}}%
\pgfpathlineto{\pgfqpoint{2.258356in}{0.992865in}}%
\pgfpathlineto{\pgfqpoint{2.263992in}{0.994959in}}%
\pgfpathlineto{\pgfqpoint{2.271083in}{0.992318in}}%
\pgfpathlineto{\pgfqpoint{2.275809in}{0.988285in}}%
\pgfpathlineto{\pgfqpoint{2.278118in}{0.995018in}}%
\pgfpathlineto{\pgfqpoint{2.285270in}{0.997505in}}%
\pgfpathlineto{\pgfqpoint{2.285799in}{1.004151in}}%
\pgfpathclose%
\pgfusepath{fill}%
\end{pgfscope}%
\begin{pgfscope}%
\pgfpathrectangle{\pgfqpoint{0.100000in}{0.100000in}}{\pgfqpoint{3.007045in}{1.925000in}}%
\pgfusepath{clip}%
\pgfsetbuttcap%
\pgfsetmiterjoin%
\definecolor{currentfill}{rgb}{0.341423,0.628958,0.808704}%
\pgfsetfillcolor{currentfill}%
\pgfsetlinewidth{0.000000pt}%
\definecolor{currentstroke}{rgb}{0.000000,0.000000,0.000000}%
\pgfsetstrokecolor{currentstroke}%
\pgfsetstrokeopacity{0.000000}%
\pgfsetdash{}{0pt}%
\pgfpathmoveto{\pgfqpoint{1.339144in}{1.451533in}}%
\pgfpathlineto{\pgfqpoint{1.298841in}{1.455523in}}%
\pgfpathlineto{\pgfqpoint{1.245609in}{1.461365in}}%
\pgfpathlineto{\pgfqpoint{1.251044in}{1.505633in}}%
\pgfpathlineto{\pgfqpoint{1.255504in}{1.539978in}}%
\pgfpathlineto{\pgfqpoint{1.257251in}{1.559802in}}%
\pgfpathlineto{\pgfqpoint{1.262489in}{1.559658in}}%
\pgfpathlineto{\pgfqpoint{1.303187in}{1.554906in}}%
\pgfpathlineto{\pgfqpoint{1.348840in}{1.549898in}}%
\pgfpathlineto{\pgfqpoint{1.343611in}{1.496079in}}%
\pgfpathlineto{\pgfqpoint{1.339144in}{1.451533in}}%
\pgfpathclose%
\pgfusepath{fill}%
\end{pgfscope}%
\begin{pgfscope}%
\pgfpathrectangle{\pgfqpoint{0.100000in}{0.100000in}}{\pgfqpoint{3.007045in}{1.925000in}}%
\pgfusepath{clip}%
\pgfsetbuttcap%
\pgfsetmiterjoin%
\definecolor{currentfill}{rgb}{0.485490,0.718524,0.853426}%
\pgfsetfillcolor{currentfill}%
\pgfsetlinewidth{0.000000pt}%
\definecolor{currentstroke}{rgb}{0.000000,0.000000,0.000000}%
\pgfsetstrokecolor{currentstroke}%
\pgfsetstrokeopacity{0.000000}%
\pgfsetdash{}{0pt}%
\pgfpathmoveto{\pgfqpoint{1.154637in}{0.238482in}}%
\pgfpathlineto{\pgfqpoint{1.153566in}{0.243456in}}%
\pgfpathlineto{\pgfqpoint{1.148403in}{0.253718in}}%
\pgfpathlineto{\pgfqpoint{1.150064in}{0.259799in}}%
\pgfpathlineto{\pgfqpoint{1.161104in}{0.270073in}}%
\pgfpathlineto{\pgfqpoint{1.164300in}{0.280924in}}%
\pgfpathlineto{\pgfqpoint{1.168171in}{0.288539in}}%
\pgfpathlineto{\pgfqpoint{1.167443in}{0.293757in}}%
\pgfpathlineto{\pgfqpoint{1.171591in}{0.300793in}}%
\pgfpathlineto{\pgfqpoint{1.191981in}{0.303197in}}%
\pgfpathlineto{\pgfqpoint{1.194052in}{0.305195in}}%
\pgfpathlineto{\pgfqpoint{1.194899in}{0.313558in}}%
\pgfpathlineto{\pgfqpoint{1.199467in}{0.319810in}}%
\pgfpathlineto{\pgfqpoint{1.203895in}{0.319010in}}%
\pgfpathlineto{\pgfqpoint{1.206973in}{0.312220in}}%
\pgfpathlineto{\pgfqpoint{1.209897in}{0.301367in}}%
\pgfpathlineto{\pgfqpoint{1.215982in}{0.294627in}}%
\pgfpathlineto{\pgfqpoint{1.221823in}{0.282731in}}%
\pgfpathlineto{\pgfqpoint{1.224745in}{0.266660in}}%
\pgfpathlineto{\pgfqpoint{1.219165in}{0.259374in}}%
\pgfpathlineto{\pgfqpoint{1.223841in}{0.256519in}}%
\pgfpathlineto{\pgfqpoint{1.220981in}{0.250601in}}%
\pgfpathlineto{\pgfqpoint{1.221882in}{0.243306in}}%
\pgfpathlineto{\pgfqpoint{1.224870in}{0.237321in}}%
\pgfpathlineto{\pgfqpoint{1.222256in}{0.234938in}}%
\pgfpathlineto{\pgfqpoint{1.208527in}{0.234063in}}%
\pgfpathlineto{\pgfqpoint{1.197467in}{0.236205in}}%
\pgfpathlineto{\pgfqpoint{1.188592in}{0.242263in}}%
\pgfpathlineto{\pgfqpoint{1.183559in}{0.241689in}}%
\pgfpathlineto{\pgfqpoint{1.172818in}{0.243600in}}%
\pgfpathlineto{\pgfqpoint{1.154637in}{0.238482in}}%
\pgfpathclose%
\pgfusepath{fill}%
\end{pgfscope}%
\begin{pgfscope}%
\pgfpathrectangle{\pgfqpoint{0.100000in}{0.100000in}}{\pgfqpoint{3.007045in}{1.925000in}}%
\pgfusepath{clip}%
\pgfsetbuttcap%
\pgfsetmiterjoin%
\definecolor{currentfill}{rgb}{0.441569,0.694410,0.843952}%
\pgfsetfillcolor{currentfill}%
\pgfsetlinewidth{0.000000pt}%
\definecolor{currentstroke}{rgb}{0.000000,0.000000,0.000000}%
\pgfsetstrokecolor{currentstroke}%
\pgfsetstrokeopacity{0.000000}%
\pgfsetdash{}{0pt}%
\pgfpathmoveto{\pgfqpoint{1.823694in}{1.499014in}}%
\pgfpathlineto{\pgfqpoint{1.822887in}{1.493279in}}%
\pgfpathlineto{\pgfqpoint{1.790073in}{1.492987in}}%
\pgfpathlineto{\pgfqpoint{1.789742in}{1.510236in}}%
\pgfpathlineto{\pgfqpoint{1.795936in}{1.510293in}}%
\pgfpathlineto{\pgfqpoint{1.795889in}{1.521880in}}%
\pgfpathlineto{\pgfqpoint{1.783670in}{1.521817in}}%
\pgfpathlineto{\pgfqpoint{1.783351in}{1.550568in}}%
\pgfpathlineto{\pgfqpoint{1.801048in}{1.550695in}}%
\pgfpathlineto{\pgfqpoint{1.806790in}{1.547902in}}%
\pgfpathlineto{\pgfqpoint{1.807185in}{1.527732in}}%
\pgfpathlineto{\pgfqpoint{1.818599in}{1.527791in}}%
\pgfpathlineto{\pgfqpoint{1.818879in}{1.510489in}}%
\pgfpathlineto{\pgfqpoint{1.824600in}{1.510542in}}%
\pgfpathlineto{\pgfqpoint{1.823694in}{1.499014in}}%
\pgfpathclose%
\pgfusepath{fill}%
\end{pgfscope}%
\begin{pgfscope}%
\pgfpathrectangle{\pgfqpoint{0.100000in}{0.100000in}}{\pgfqpoint{3.007045in}{1.925000in}}%
\pgfusepath{clip}%
\pgfsetbuttcap%
\pgfsetmiterjoin%
\definecolor{currentfill}{rgb}{0.412042,0.677186,0.836263}%
\pgfsetfillcolor{currentfill}%
\pgfsetlinewidth{0.000000pt}%
\definecolor{currentstroke}{rgb}{0.000000,0.000000,0.000000}%
\pgfsetstrokecolor{currentstroke}%
\pgfsetstrokeopacity{0.000000}%
\pgfsetdash{}{0pt}%
\pgfpathmoveto{\pgfqpoint{2.214133in}{1.095845in}}%
\pgfpathlineto{\pgfqpoint{2.211772in}{1.090499in}}%
\pgfpathlineto{\pgfqpoint{2.212960in}{1.083080in}}%
\pgfpathlineto{\pgfqpoint{2.208829in}{1.082188in}}%
\pgfpathlineto{\pgfqpoint{2.205852in}{1.076709in}}%
\pgfpathlineto{\pgfqpoint{2.205878in}{1.081086in}}%
\pgfpathlineto{\pgfqpoint{2.200736in}{1.080556in}}%
\pgfpathlineto{\pgfqpoint{2.197741in}{1.085673in}}%
\pgfpathlineto{\pgfqpoint{2.185322in}{1.077996in}}%
\pgfpathlineto{\pgfqpoint{2.184734in}{1.072536in}}%
\pgfpathlineto{\pgfqpoint{2.172010in}{1.077414in}}%
\pgfpathlineto{\pgfqpoint{2.165193in}{1.079905in}}%
\pgfpathlineto{\pgfqpoint{2.158994in}{1.078092in}}%
\pgfpathlineto{\pgfqpoint{2.156206in}{1.071747in}}%
\pgfpathlineto{\pgfqpoint{2.152373in}{1.076112in}}%
\pgfpathlineto{\pgfqpoint{2.144755in}{1.073731in}}%
\pgfpathlineto{\pgfqpoint{2.139094in}{1.074231in}}%
\pgfpathlineto{\pgfqpoint{2.141157in}{1.068895in}}%
\pgfpathlineto{\pgfqpoint{2.134908in}{1.067855in}}%
\pgfpathlineto{\pgfqpoint{2.131121in}{1.073626in}}%
\pgfpathlineto{\pgfqpoint{2.137521in}{1.089690in}}%
\pgfpathlineto{\pgfqpoint{2.134514in}{1.098251in}}%
\pgfpathlineto{\pgfqpoint{2.136503in}{1.101602in}}%
\pgfpathlineto{\pgfqpoint{2.134603in}{1.106062in}}%
\pgfpathlineto{\pgfqpoint{2.135045in}{1.113976in}}%
\pgfpathlineto{\pgfqpoint{2.136969in}{1.118978in}}%
\pgfpathlineto{\pgfqpoint{2.150356in}{1.119886in}}%
\pgfpathlineto{\pgfqpoint{2.154476in}{1.114634in}}%
\pgfpathlineto{\pgfqpoint{2.160227in}{1.118407in}}%
\pgfpathlineto{\pgfqpoint{2.171500in}{1.120112in}}%
\pgfpathlineto{\pgfqpoint{2.178440in}{1.120233in}}%
\pgfpathlineto{\pgfqpoint{2.190430in}{1.118463in}}%
\pgfpathlineto{\pgfqpoint{2.192366in}{1.120724in}}%
\pgfpathlineto{\pgfqpoint{2.200334in}{1.121433in}}%
\pgfpathlineto{\pgfqpoint{2.201173in}{1.112813in}}%
\pgfpathlineto{\pgfqpoint{2.202065in}{1.104184in}}%
\pgfpathlineto{\pgfqpoint{2.207678in}{1.104917in}}%
\pgfpathlineto{\pgfqpoint{2.207958in}{1.101054in}}%
\pgfpathlineto{\pgfqpoint{2.213695in}{1.101544in}}%
\pgfpathlineto{\pgfqpoint{2.214133in}{1.095845in}}%
\pgfpathclose%
\pgfusepath{fill}%
\end{pgfscope}%
\begin{pgfscope}%
\pgfpathrectangle{\pgfqpoint{0.100000in}{0.100000in}}{\pgfqpoint{3.007045in}{1.925000in}}%
\pgfusepath{clip}%
\pgfsetbuttcap%
\pgfsetmiterjoin%
\definecolor{currentfill}{rgb}{0.711265,0.831111,0.914433}%
\pgfsetfillcolor{currentfill}%
\pgfsetlinewidth{0.000000pt}%
\definecolor{currentstroke}{rgb}{0.000000,0.000000,0.000000}%
\pgfsetstrokecolor{currentstroke}%
\pgfsetstrokeopacity{0.000000}%
\pgfsetdash{}{0pt}%
\pgfpathmoveto{\pgfqpoint{1.929386in}{1.648471in}}%
\pgfpathlineto{\pgfqpoint{1.930881in}{1.608680in}}%
\pgfpathlineto{\pgfqpoint{1.907975in}{1.607918in}}%
\pgfpathlineto{\pgfqpoint{1.896759in}{1.607543in}}%
\pgfpathlineto{\pgfqpoint{1.896284in}{1.624795in}}%
\pgfpathlineto{\pgfqpoint{1.861397in}{1.624038in}}%
\pgfpathlineto{\pgfqpoint{1.860671in}{1.638052in}}%
\pgfpathlineto{\pgfqpoint{1.860483in}{1.664391in}}%
\pgfpathlineto{\pgfqpoint{1.859346in}{1.688805in}}%
\pgfpathlineto{\pgfqpoint{1.859255in}{1.710576in}}%
\pgfpathlineto{\pgfqpoint{1.858121in}{1.722001in}}%
\pgfpathlineto{\pgfqpoint{1.857189in}{1.733722in}}%
\pgfpathlineto{\pgfqpoint{1.856761in}{1.771102in}}%
\pgfpathlineto{\pgfqpoint{1.861336in}{1.770943in}}%
\pgfpathlineto{\pgfqpoint{1.872825in}{1.765570in}}%
\pgfpathlineto{\pgfqpoint{1.876937in}{1.765915in}}%
\pgfpathlineto{\pgfqpoint{1.876187in}{1.758800in}}%
\pgfpathlineto{\pgfqpoint{1.882756in}{1.759721in}}%
\pgfpathlineto{\pgfqpoint{1.889291in}{1.744713in}}%
\pgfpathlineto{\pgfqpoint{1.893197in}{1.746419in}}%
\pgfpathlineto{\pgfqpoint{1.892633in}{1.752996in}}%
\pgfpathlineto{\pgfqpoint{1.902968in}{1.754402in}}%
\pgfpathlineto{\pgfqpoint{1.905332in}{1.748213in}}%
\pgfpathlineto{\pgfqpoint{1.911782in}{1.744531in}}%
\pgfpathlineto{\pgfqpoint{1.918443in}{1.744191in}}%
\pgfpathlineto{\pgfqpoint{1.918786in}{1.738606in}}%
\pgfpathlineto{\pgfqpoint{1.931137in}{1.734731in}}%
\pgfpathlineto{\pgfqpoint{1.939426in}{1.737368in}}%
\pgfpathlineto{\pgfqpoint{1.948847in}{1.744739in}}%
\pgfpathlineto{\pgfqpoint{1.950215in}{1.724734in}}%
\pgfpathlineto{\pgfqpoint{1.951342in}{1.696561in}}%
\pgfpathlineto{\pgfqpoint{1.941485in}{1.683918in}}%
\pgfpathlineto{\pgfqpoint{1.931784in}{1.673108in}}%
\pgfpathlineto{\pgfqpoint{1.927592in}{1.670463in}}%
\pgfpathlineto{\pgfqpoint{1.918507in}{1.660575in}}%
\pgfpathlineto{\pgfqpoint{1.904812in}{1.649918in}}%
\pgfpathlineto{\pgfqpoint{1.907843in}{1.644693in}}%
\pgfpathlineto{\pgfqpoint{1.911711in}{1.642752in}}%
\pgfpathlineto{\pgfqpoint{1.918643in}{1.643971in}}%
\pgfpathlineto{\pgfqpoint{1.929386in}{1.648471in}}%
\pgfpathclose%
\pgfusepath{fill}%
\end{pgfscope}%
\begin{pgfscope}%
\pgfpathrectangle{\pgfqpoint{0.100000in}{0.100000in}}{\pgfqpoint{3.007045in}{1.925000in}}%
\pgfusepath{clip}%
\pgfsetbuttcap%
\pgfsetmiterjoin%
\definecolor{currentfill}{rgb}{0.371688,0.649627,0.820515}%
\pgfsetfillcolor{currentfill}%
\pgfsetlinewidth{0.000000pt}%
\definecolor{currentstroke}{rgb}{0.000000,0.000000,0.000000}%
\pgfsetstrokecolor{currentstroke}%
\pgfsetstrokeopacity{0.000000}%
\pgfsetdash{}{0pt}%
\pgfpathmoveto{\pgfqpoint{1.923061in}{1.300995in}}%
\pgfpathlineto{\pgfqpoint{1.905724in}{1.300373in}}%
\pgfpathlineto{\pgfqpoint{1.904114in}{1.352301in}}%
\pgfpathlineto{\pgfqpoint{1.938345in}{1.353373in}}%
\pgfpathlineto{\pgfqpoint{1.972422in}{1.354713in}}%
\pgfpathlineto{\pgfqpoint{1.973436in}{1.331705in}}%
\pgfpathlineto{\pgfqpoint{1.950534in}{1.330806in}}%
\pgfpathlineto{\pgfqpoint{1.950752in}{1.325057in}}%
\pgfpathlineto{\pgfqpoint{1.927902in}{1.324319in}}%
\pgfpathlineto{\pgfqpoint{1.928815in}{1.301237in}}%
\pgfpathlineto{\pgfqpoint{1.923061in}{1.300995in}}%
\pgfpathclose%
\pgfusepath{fill}%
\end{pgfscope}%
\begin{pgfscope}%
\pgfpathrectangle{\pgfqpoint{0.100000in}{0.100000in}}{\pgfqpoint{3.007045in}{1.925000in}}%
\pgfusepath{clip}%
\pgfsetbuttcap%
\pgfsetmiterjoin%
\definecolor{currentfill}{rgb}{0.554510,0.756417,0.868312}%
\pgfsetfillcolor{currentfill}%
\pgfsetlinewidth{0.000000pt}%
\definecolor{currentstroke}{rgb}{0.000000,0.000000,0.000000}%
\pgfsetstrokecolor{currentstroke}%
\pgfsetstrokeopacity{0.000000}%
\pgfsetdash{}{0pt}%
\pgfpathmoveto{\pgfqpoint{2.549716in}{0.711373in}}%
\pgfpathlineto{\pgfqpoint{2.544812in}{0.713097in}}%
\pgfpathlineto{\pgfqpoint{2.537534in}{0.718203in}}%
\pgfpathlineto{\pgfqpoint{2.534070in}{0.723697in}}%
\pgfpathlineto{\pgfqpoint{2.526011in}{0.723942in}}%
\pgfpathlineto{\pgfqpoint{2.520931in}{0.729238in}}%
\pgfpathlineto{\pgfqpoint{2.513924in}{0.731882in}}%
\pgfpathlineto{\pgfqpoint{2.512057in}{0.734800in}}%
\pgfpathlineto{\pgfqpoint{2.530406in}{0.743490in}}%
\pgfpathlineto{\pgfqpoint{2.528796in}{0.747623in}}%
\pgfpathlineto{\pgfqpoint{2.523508in}{0.752566in}}%
\pgfpathlineto{\pgfqpoint{2.522083in}{0.760192in}}%
\pgfpathlineto{\pgfqpoint{2.529482in}{0.767093in}}%
\pgfpathlineto{\pgfqpoint{2.535487in}{0.765030in}}%
\pgfpathlineto{\pgfqpoint{2.548564in}{0.780580in}}%
\pgfpathlineto{\pgfqpoint{2.556087in}{0.772933in}}%
\pgfpathlineto{\pgfqpoint{2.559434in}{0.779009in}}%
\pgfpathlineto{\pgfqpoint{2.567570in}{0.778188in}}%
\pgfpathlineto{\pgfqpoint{2.572033in}{0.772743in}}%
\pgfpathlineto{\pgfqpoint{2.576021in}{0.772345in}}%
\pgfpathlineto{\pgfqpoint{2.581007in}{0.767715in}}%
\pgfpathlineto{\pgfqpoint{2.584609in}{0.761996in}}%
\pgfpathlineto{\pgfqpoint{2.584009in}{0.757368in}}%
\pgfpathlineto{\pgfqpoint{2.574779in}{0.751322in}}%
\pgfpathlineto{\pgfqpoint{2.570706in}{0.743972in}}%
\pgfpathlineto{\pgfqpoint{2.562709in}{0.737347in}}%
\pgfpathlineto{\pgfqpoint{2.564979in}{0.730555in}}%
\pgfpathlineto{\pgfqpoint{2.552611in}{0.712060in}}%
\pgfpathlineto{\pgfqpoint{2.549716in}{0.711373in}}%
\pgfpathclose%
\pgfusepath{fill}%
\end{pgfscope}%
\begin{pgfscope}%
\pgfpathrectangle{\pgfqpoint{0.100000in}{0.100000in}}{\pgfqpoint{3.007045in}{1.925000in}}%
\pgfusepath{clip}%
\pgfsetbuttcap%
\pgfsetmiterjoin%
\definecolor{currentfill}{rgb}{0.396909,0.666851,0.830358}%
\pgfsetfillcolor{currentfill}%
\pgfsetlinewidth{0.000000pt}%
\definecolor{currentstroke}{rgb}{0.000000,0.000000,0.000000}%
\pgfsetstrokecolor{currentstroke}%
\pgfsetstrokeopacity{0.000000}%
\pgfsetdash{}{0pt}%
\pgfpathmoveto{\pgfqpoint{0.613279in}{1.618283in}}%
\pgfpathlineto{\pgfqpoint{0.651898in}{1.607795in}}%
\pgfpathlineto{\pgfqpoint{0.653418in}{1.613391in}}%
\pgfpathlineto{\pgfqpoint{0.680340in}{1.605736in}}%
\pgfpathlineto{\pgfqpoint{0.675770in}{1.589097in}}%
\pgfpathlineto{\pgfqpoint{0.661467in}{1.533821in}}%
\pgfpathlineto{\pgfqpoint{0.662044in}{1.533668in}}%
\pgfpathlineto{\pgfqpoint{0.651488in}{1.492874in}}%
\pgfpathlineto{\pgfqpoint{0.647669in}{1.474848in}}%
\pgfpathlineto{\pgfqpoint{0.623923in}{1.480760in}}%
\pgfpathlineto{\pgfqpoint{0.594261in}{1.488948in}}%
\pgfpathlineto{\pgfqpoint{0.592563in}{1.489427in}}%
\pgfpathlineto{\pgfqpoint{0.605512in}{1.537505in}}%
\pgfpathlineto{\pgfqpoint{0.578445in}{1.544847in}}%
\pgfpathlineto{\pgfqpoint{0.586719in}{1.572193in}}%
\pgfpathlineto{\pgfqpoint{0.588372in}{1.571730in}}%
\pgfpathlineto{\pgfqpoint{0.596099in}{1.599150in}}%
\pgfpathlineto{\pgfqpoint{0.597590in}{1.604775in}}%
\pgfpathlineto{\pgfqpoint{0.603136in}{1.603204in}}%
\pgfpathlineto{\pgfqpoint{0.607867in}{1.619794in}}%
\pgfpathlineto{\pgfqpoint{0.613279in}{1.618283in}}%
\pgfpathclose%
\pgfusepath{fill}%
\end{pgfscope}%
\begin{pgfscope}%
\pgfpathrectangle{\pgfqpoint{0.100000in}{0.100000in}}{\pgfqpoint{3.007045in}{1.925000in}}%
\pgfusepath{clip}%
\pgfsetbuttcap%
\pgfsetmiterjoin%
\definecolor{currentfill}{rgb}{0.671895,0.814379,0.900654}%
\pgfsetfillcolor{currentfill}%
\pgfsetlinewidth{0.000000pt}%
\definecolor{currentstroke}{rgb}{0.000000,0.000000,0.000000}%
\pgfsetstrokecolor{currentstroke}%
\pgfsetstrokeopacity{0.000000}%
\pgfsetdash{}{0pt}%
\pgfpathmoveto{\pgfqpoint{2.623105in}{0.353228in}}%
\pgfpathlineto{\pgfqpoint{2.599929in}{0.349210in}}%
\pgfpathlineto{\pgfqpoint{2.597057in}{0.366782in}}%
\pgfpathlineto{\pgfqpoint{2.579681in}{0.363910in}}%
\pgfpathlineto{\pgfqpoint{2.577000in}{0.380965in}}%
\pgfpathlineto{\pgfqpoint{2.574578in}{0.398555in}}%
\pgfpathlineto{\pgfqpoint{2.592108in}{0.401133in}}%
\pgfpathlineto{\pgfqpoint{2.591178in}{0.407023in}}%
\pgfpathlineto{\pgfqpoint{2.597054in}{0.407964in}}%
\pgfpathlineto{\pgfqpoint{2.596119in}{0.413805in}}%
\pgfpathlineto{\pgfqpoint{2.607766in}{0.416135in}}%
\pgfpathlineto{\pgfqpoint{2.600616in}{0.426095in}}%
\pgfpathlineto{\pgfqpoint{2.597645in}{0.425944in}}%
\pgfpathlineto{\pgfqpoint{2.590385in}{0.434304in}}%
\pgfpathlineto{\pgfqpoint{2.593275in}{0.442838in}}%
\pgfpathlineto{\pgfqpoint{2.609094in}{0.445229in}}%
\pgfpathlineto{\pgfqpoint{2.614673in}{0.446167in}}%
\pgfpathlineto{\pgfqpoint{2.615609in}{0.440573in}}%
\pgfpathlineto{\pgfqpoint{2.642481in}{0.444772in}}%
\pgfpathlineto{\pgfqpoint{2.647094in}{0.435626in}}%
\pgfpathlineto{\pgfqpoint{2.655837in}{0.422084in}}%
\pgfpathlineto{\pgfqpoint{2.661202in}{0.412146in}}%
\pgfpathlineto{\pgfqpoint{2.667843in}{0.397101in}}%
\pgfpathlineto{\pgfqpoint{2.669777in}{0.383631in}}%
\pgfpathlineto{\pgfqpoint{2.670514in}{0.365060in}}%
\pgfpathlineto{\pgfqpoint{2.661292in}{0.364487in}}%
\pgfpathlineto{\pgfqpoint{2.622244in}{0.358169in}}%
\pgfpathlineto{\pgfqpoint{2.623105in}{0.353228in}}%
\pgfpathclose%
\pgfusepath{fill}%
\end{pgfscope}%
\begin{pgfscope}%
\pgfpathrectangle{\pgfqpoint{0.100000in}{0.100000in}}{\pgfqpoint{3.007045in}{1.925000in}}%
\pgfusepath{clip}%
\pgfsetbuttcap%
\pgfsetmiterjoin%
\definecolor{currentfill}{rgb}{0.585882,0.773641,0.875079}%
\pgfsetfillcolor{currentfill}%
\pgfsetlinewidth{0.000000pt}%
\definecolor{currentstroke}{rgb}{0.000000,0.000000,0.000000}%
\pgfsetstrokecolor{currentstroke}%
\pgfsetstrokeopacity{0.000000}%
\pgfsetdash{}{0pt}%
\pgfpathmoveto{\pgfqpoint{2.201173in}{1.112813in}}%
\pgfpathlineto{\pgfqpoint{2.200334in}{1.121433in}}%
\pgfpathlineto{\pgfqpoint{2.192366in}{1.120724in}}%
\pgfpathlineto{\pgfqpoint{2.190430in}{1.118463in}}%
\pgfpathlineto{\pgfqpoint{2.178440in}{1.120233in}}%
\pgfpathlineto{\pgfqpoint{2.171500in}{1.120112in}}%
\pgfpathlineto{\pgfqpoint{2.168758in}{1.129637in}}%
\pgfpathlineto{\pgfqpoint{2.170118in}{1.138155in}}%
\pgfpathlineto{\pgfqpoint{2.175011in}{1.140154in}}%
\pgfpathlineto{\pgfqpoint{2.175194in}{1.144388in}}%
\pgfpathlineto{\pgfqpoint{2.169528in}{1.143936in}}%
\pgfpathlineto{\pgfqpoint{2.168048in}{1.161256in}}%
\pgfpathlineto{\pgfqpoint{2.177539in}{1.161835in}}%
\pgfpathlineto{\pgfqpoint{2.176581in}{1.173291in}}%
\pgfpathlineto{\pgfqpoint{2.182227in}{1.173751in}}%
\pgfpathlineto{\pgfqpoint{2.181588in}{1.182377in}}%
\pgfpathlineto{\pgfqpoint{2.194423in}{1.183326in}}%
\pgfpathlineto{\pgfqpoint{2.197194in}{1.183592in}}%
\pgfpathlineto{\pgfqpoint{2.198043in}{1.175099in}}%
\pgfpathlineto{\pgfqpoint{2.210701in}{1.176219in}}%
\pgfpathlineto{\pgfqpoint{2.213096in}{1.157168in}}%
\pgfpathlineto{\pgfqpoint{2.215737in}{1.157449in}}%
\pgfpathlineto{\pgfqpoint{2.218022in}{1.153845in}}%
\pgfpathlineto{\pgfqpoint{2.219718in}{1.138972in}}%
\pgfpathlineto{\pgfqpoint{2.218133in}{1.136791in}}%
\pgfpathlineto{\pgfqpoint{2.220327in}{1.114564in}}%
\pgfpathlineto{\pgfqpoint{2.201173in}{1.112813in}}%
\pgfpathclose%
\pgfusepath{fill}%
\end{pgfscope}%
\begin{pgfscope}%
\pgfpathrectangle{\pgfqpoint{0.100000in}{0.100000in}}{\pgfqpoint{3.007045in}{1.925000in}}%
\pgfusepath{clip}%
\pgfsetbuttcap%
\pgfsetmiterjoin%
\definecolor{currentfill}{rgb}{0.381776,0.656517,0.824452}%
\pgfsetfillcolor{currentfill}%
\pgfsetlinewidth{0.000000pt}%
\definecolor{currentstroke}{rgb}{0.000000,0.000000,0.000000}%
\pgfsetstrokecolor{currentstroke}%
\pgfsetstrokeopacity{0.000000}%
\pgfsetdash{}{0pt}%
\pgfpathmoveto{\pgfqpoint{1.899049in}{0.722750in}}%
\pgfpathlineto{\pgfqpoint{1.906378in}{0.722687in}}%
\pgfpathlineto{\pgfqpoint{1.910744in}{0.717166in}}%
\pgfpathlineto{\pgfqpoint{1.916505in}{0.717286in}}%
\pgfpathlineto{\pgfqpoint{1.916979in}{0.699909in}}%
\pgfpathlineto{\pgfqpoint{1.922946in}{0.694307in}}%
\pgfpathlineto{\pgfqpoint{1.923673in}{0.668219in}}%
\pgfpathlineto{\pgfqpoint{1.921110in}{0.659462in}}%
\pgfpathlineto{\pgfqpoint{1.906637in}{0.659096in}}%
\pgfpathlineto{\pgfqpoint{1.906746in}{0.653283in}}%
\pgfpathlineto{\pgfqpoint{1.886783in}{0.652756in}}%
\pgfpathlineto{\pgfqpoint{1.888359in}{0.662427in}}%
\pgfpathlineto{\pgfqpoint{1.891446in}{0.668175in}}%
\pgfpathlineto{\pgfqpoint{1.888095in}{0.681930in}}%
\pgfpathlineto{\pgfqpoint{1.895056in}{0.682024in}}%
\pgfpathlineto{\pgfqpoint{1.897220in}{0.687919in}}%
\pgfpathlineto{\pgfqpoint{1.896662in}{0.702398in}}%
\pgfpathlineto{\pgfqpoint{1.890952in}{0.702305in}}%
\pgfpathlineto{\pgfqpoint{1.890590in}{0.711973in}}%
\pgfpathlineto{\pgfqpoint{1.893306in}{0.722562in}}%
\pgfpathlineto{\pgfqpoint{1.899049in}{0.722750in}}%
\pgfpathclose%
\pgfusepath{fill}%
\end{pgfscope}%
\begin{pgfscope}%
\pgfpathrectangle{\pgfqpoint{0.100000in}{0.100000in}}{\pgfqpoint{3.007045in}{1.925000in}}%
\pgfusepath{clip}%
\pgfsetbuttcap%
\pgfsetmiterjoin%
\definecolor{currentfill}{rgb}{0.223806,0.537532,0.758431}%
\pgfsetfillcolor{currentfill}%
\pgfsetlinewidth{0.000000pt}%
\definecolor{currentstroke}{rgb}{0.000000,0.000000,0.000000}%
\pgfsetstrokecolor{currentstroke}%
\pgfsetstrokeopacity{0.000000}%
\pgfsetdash{}{0pt}%
\pgfpathmoveto{\pgfqpoint{1.092509in}{1.504079in}}%
\pgfpathlineto{\pgfqpoint{1.098004in}{1.495582in}}%
\pgfpathlineto{\pgfqpoint{1.097179in}{1.490393in}}%
\pgfpathlineto{\pgfqpoint{1.108398in}{1.488604in}}%
\pgfpathlineto{\pgfqpoint{1.119290in}{1.484034in}}%
\pgfpathlineto{\pgfqpoint{1.118840in}{1.481195in}}%
\pgfpathlineto{\pgfqpoint{1.135060in}{1.473922in}}%
\pgfpathlineto{\pgfqpoint{1.143778in}{1.473368in}}%
\pgfpathlineto{\pgfqpoint{1.170426in}{1.469676in}}%
\pgfpathlineto{\pgfqpoint{1.173551in}{1.471182in}}%
\pgfpathlineto{\pgfqpoint{1.174139in}{1.469186in}}%
\pgfpathlineto{\pgfqpoint{1.171716in}{1.446544in}}%
\pgfpathlineto{\pgfqpoint{1.168621in}{1.423872in}}%
\pgfpathlineto{\pgfqpoint{1.166601in}{1.424158in}}%
\pgfpathlineto{\pgfqpoint{1.163522in}{1.401453in}}%
\pgfpathlineto{\pgfqpoint{1.164443in}{1.401325in}}%
\pgfpathlineto{\pgfqpoint{1.162889in}{1.390045in}}%
\pgfpathlineto{\pgfqpoint{1.139789in}{1.393330in}}%
\pgfpathlineto{\pgfqpoint{1.089437in}{1.400986in}}%
\pgfpathlineto{\pgfqpoint{1.089929in}{1.412247in}}%
\pgfpathlineto{\pgfqpoint{1.092414in}{1.429227in}}%
\pgfpathlineto{\pgfqpoint{1.085388in}{1.435891in}}%
\pgfpathlineto{\pgfqpoint{1.079954in}{1.447589in}}%
\pgfpathlineto{\pgfqpoint{1.075159in}{1.453148in}}%
\pgfpathlineto{\pgfqpoint{1.070708in}{1.461628in}}%
\pgfpathlineto{\pgfqpoint{1.069165in}{1.468405in}}%
\pgfpathlineto{\pgfqpoint{1.069796in}{1.477757in}}%
\pgfpathlineto{\pgfqpoint{1.067969in}{1.484605in}}%
\pgfpathlineto{\pgfqpoint{1.053944in}{1.487032in}}%
\pgfpathlineto{\pgfqpoint{1.059788in}{1.522455in}}%
\pgfpathlineto{\pgfqpoint{1.062108in}{1.518116in}}%
\pgfpathlineto{\pgfqpoint{1.067064in}{1.517262in}}%
\pgfpathlineto{\pgfqpoint{1.070042in}{1.506895in}}%
\pgfpathlineto{\pgfqpoint{1.076162in}{1.510378in}}%
\pgfpathlineto{\pgfqpoint{1.078184in}{1.513711in}}%
\pgfpathlineto{\pgfqpoint{1.083266in}{1.515471in}}%
\pgfpathlineto{\pgfqpoint{1.087640in}{1.512697in}}%
\pgfpathlineto{\pgfqpoint{1.092509in}{1.504079in}}%
\pgfpathclose%
\pgfusepath{fill}%
\end{pgfscope}%
\begin{pgfscope}%
\pgfpathrectangle{\pgfqpoint{0.100000in}{0.100000in}}{\pgfqpoint{3.007045in}{1.925000in}}%
\pgfusepath{clip}%
\pgfsetbuttcap%
\pgfsetmiterjoin%
\definecolor{currentfill}{rgb}{0.765398,0.854118,0.933379}%
\pgfsetfillcolor{currentfill}%
\pgfsetlinewidth{0.000000pt}%
\definecolor{currentstroke}{rgb}{0.000000,0.000000,0.000000}%
\pgfsetstrokecolor{currentstroke}%
\pgfsetstrokeopacity{0.000000}%
\pgfsetdash{}{0pt}%
\pgfpathmoveto{\pgfqpoint{2.031351in}{1.567858in}}%
\pgfpathlineto{\pgfqpoint{2.030834in}{1.573678in}}%
\pgfpathlineto{\pgfqpoint{2.002328in}{1.571903in}}%
\pgfpathlineto{\pgfqpoint{2.000756in}{1.600279in}}%
\pgfpathlineto{\pgfqpoint{2.006057in}{1.600572in}}%
\pgfpathlineto{\pgfqpoint{2.004866in}{1.621636in}}%
\pgfpathlineto{\pgfqpoint{2.043917in}{1.613272in}}%
\pgfpathlineto{\pgfqpoint{2.048726in}{1.610846in}}%
\pgfpathlineto{\pgfqpoint{2.063792in}{1.604125in}}%
\pgfpathlineto{\pgfqpoint{2.064911in}{1.587042in}}%
\pgfpathlineto{\pgfqpoint{2.076489in}{1.587837in}}%
\pgfpathlineto{\pgfqpoint{2.078017in}{1.564986in}}%
\pgfpathlineto{\pgfqpoint{2.066412in}{1.564272in}}%
\pgfpathlineto{\pgfqpoint{2.054884in}{1.563482in}}%
\pgfpathlineto{\pgfqpoint{2.054639in}{1.569212in}}%
\pgfpathlineto{\pgfqpoint{2.031351in}{1.567858in}}%
\pgfpathclose%
\pgfusepath{fill}%
\end{pgfscope}%
\begin{pgfscope}%
\pgfpathrectangle{\pgfqpoint{0.100000in}{0.100000in}}{\pgfqpoint{3.007045in}{1.925000in}}%
\pgfusepath{clip}%
\pgfsetbuttcap%
\pgfsetmiterjoin%
\definecolor{currentfill}{rgb}{0.102499,0.408689,0.682891}%
\pgfsetfillcolor{currentfill}%
\pgfsetlinewidth{0.000000pt}%
\definecolor{currentstroke}{rgb}{0.000000,0.000000,0.000000}%
\pgfsetstrokecolor{currentstroke}%
\pgfsetstrokeopacity{0.000000}%
\pgfsetdash{}{0pt}%
\pgfpathmoveto{\pgfqpoint{0.955057in}{1.434288in}}%
\pgfpathlineto{\pgfqpoint{0.952072in}{1.424129in}}%
\pgfpathlineto{\pgfqpoint{0.948249in}{1.426793in}}%
\pgfpathlineto{\pgfqpoint{0.941439in}{1.428069in}}%
\pgfpathlineto{\pgfqpoint{0.934955in}{1.436848in}}%
\pgfpathlineto{\pgfqpoint{0.931789in}{1.444541in}}%
\pgfpathlineto{\pgfqpoint{0.926367in}{1.445605in}}%
\pgfpathlineto{\pgfqpoint{0.925285in}{1.439963in}}%
\pgfpathlineto{\pgfqpoint{0.917819in}{1.441425in}}%
\pgfpathlineto{\pgfqpoint{0.916709in}{1.435786in}}%
\pgfpathlineto{\pgfqpoint{0.900100in}{1.438993in}}%
\pgfpathlineto{\pgfqpoint{0.903388in}{1.455868in}}%
\pgfpathlineto{\pgfqpoint{0.895172in}{1.457614in}}%
\pgfpathlineto{\pgfqpoint{0.892728in}{1.460530in}}%
\pgfpathlineto{\pgfqpoint{0.896143in}{1.463222in}}%
\pgfpathlineto{\pgfqpoint{0.894533in}{1.469396in}}%
\pgfpathlineto{\pgfqpoint{0.898811in}{1.492021in}}%
\pgfpathlineto{\pgfqpoint{0.909954in}{1.489763in}}%
\pgfpathlineto{\pgfqpoint{0.912230in}{1.501011in}}%
\pgfpathlineto{\pgfqpoint{0.895520in}{1.504390in}}%
\pgfpathlineto{\pgfqpoint{0.896614in}{1.509734in}}%
\pgfpathlineto{\pgfqpoint{0.883755in}{1.512410in}}%
\pgfpathlineto{\pgfqpoint{0.884911in}{1.517908in}}%
\pgfpathlineto{\pgfqpoint{0.882414in}{1.527693in}}%
\pgfpathlineto{\pgfqpoint{0.878844in}{1.527019in}}%
\pgfpathlineto{\pgfqpoint{0.881395in}{1.530161in}}%
\pgfpathlineto{\pgfqpoint{0.891068in}{1.533017in}}%
\pgfpathlineto{\pgfqpoint{0.897650in}{1.536577in}}%
\pgfpathlineto{\pgfqpoint{0.900875in}{1.550400in}}%
\pgfpathlineto{\pgfqpoint{0.903257in}{1.555054in}}%
\pgfpathlineto{\pgfqpoint{0.905554in}{1.566329in}}%
\pgfpathlineto{\pgfqpoint{0.933521in}{1.560492in}}%
\pgfpathlineto{\pgfqpoint{0.935411in}{1.569991in}}%
\pgfpathlineto{\pgfqpoint{0.938309in}{1.576703in}}%
\pgfpathlineto{\pgfqpoint{0.952519in}{1.573584in}}%
\pgfpathlineto{\pgfqpoint{0.956111in}{1.570703in}}%
\pgfpathlineto{\pgfqpoint{0.958428in}{1.575728in}}%
\pgfpathlineto{\pgfqpoint{0.962156in}{1.577626in}}%
\pgfpathlineto{\pgfqpoint{0.970525in}{1.573027in}}%
\pgfpathlineto{\pgfqpoint{0.981565in}{1.573774in}}%
\pgfpathlineto{\pgfqpoint{0.983095in}{1.569716in}}%
\pgfpathlineto{\pgfqpoint{0.989150in}{1.572014in}}%
\pgfpathlineto{\pgfqpoint{0.998004in}{1.571906in}}%
\pgfpathlineto{\pgfqpoint{1.001653in}{1.579735in}}%
\pgfpathlineto{\pgfqpoint{1.006564in}{1.581699in}}%
\pgfpathlineto{\pgfqpoint{1.009746in}{1.577792in}}%
\pgfpathlineto{\pgfqpoint{1.015015in}{1.563234in}}%
\pgfpathlineto{\pgfqpoint{1.018543in}{1.560865in}}%
\pgfpathlineto{\pgfqpoint{1.012932in}{1.528900in}}%
\pgfpathlineto{\pgfqpoint{1.005109in}{1.526977in}}%
\pgfpathlineto{\pgfqpoint{0.995747in}{1.527955in}}%
\pgfpathlineto{\pgfqpoint{0.992126in}{1.508384in}}%
\pgfpathlineto{\pgfqpoint{0.999466in}{1.507030in}}%
\pgfpathlineto{\pgfqpoint{1.000614in}{1.501315in}}%
\pgfpathlineto{\pgfqpoint{1.007407in}{1.497511in}}%
\pgfpathlineto{\pgfqpoint{1.005297in}{1.485446in}}%
\pgfpathlineto{\pgfqpoint{1.001890in}{1.466220in}}%
\pgfpathlineto{\pgfqpoint{0.953549in}{1.475280in}}%
\pgfpathlineto{\pgfqpoint{0.947898in}{1.467983in}}%
\pgfpathlineto{\pgfqpoint{0.946788in}{1.460996in}}%
\pgfpathlineto{\pgfqpoint{0.947546in}{1.454706in}}%
\pgfpathlineto{\pgfqpoint{0.954259in}{1.454699in}}%
\pgfpathlineto{\pgfqpoint{0.955530in}{1.447686in}}%
\pgfpathlineto{\pgfqpoint{0.955057in}{1.434288in}}%
\pgfpathclose%
\pgfusepath{fill}%
\end{pgfscope}%
\begin{pgfscope}%
\pgfpathrectangle{\pgfqpoint{0.100000in}{0.100000in}}{\pgfqpoint{3.007045in}{1.925000in}}%
\pgfusepath{clip}%
\pgfsetbuttcap%
\pgfsetmiterjoin%
\definecolor{currentfill}{rgb}{0.056363,0.349635,0.636755}%
\pgfsetfillcolor{currentfill}%
\pgfsetlinewidth{0.000000pt}%
\definecolor{currentstroke}{rgb}{0.000000,0.000000,0.000000}%
\pgfsetstrokecolor{currentstroke}%
\pgfsetstrokeopacity{0.000000}%
\pgfsetdash{}{0pt}%
\pgfpathmoveto{\pgfqpoint{0.615087in}{1.868865in}}%
\pgfpathlineto{\pgfqpoint{0.618659in}{1.862354in}}%
\pgfpathlineto{\pgfqpoint{0.620178in}{1.855997in}}%
\pgfpathlineto{\pgfqpoint{0.627702in}{1.851025in}}%
\pgfpathlineto{\pgfqpoint{0.629717in}{1.846964in}}%
\pgfpathlineto{\pgfqpoint{0.635771in}{1.843623in}}%
\pgfpathlineto{\pgfqpoint{0.640168in}{1.838234in}}%
\pgfpathlineto{\pgfqpoint{0.652989in}{1.834668in}}%
\pgfpathlineto{\pgfqpoint{0.656028in}{1.830896in}}%
\pgfpathlineto{\pgfqpoint{0.652061in}{1.822418in}}%
\pgfpathlineto{\pgfqpoint{0.652827in}{1.812664in}}%
\pgfpathlineto{\pgfqpoint{0.651860in}{1.802864in}}%
\pgfpathlineto{\pgfqpoint{0.648904in}{1.795115in}}%
\pgfpathlineto{\pgfqpoint{0.651182in}{1.791391in}}%
\pgfpathlineto{\pgfqpoint{0.640972in}{1.753699in}}%
\pgfpathlineto{\pgfqpoint{0.616044in}{1.760740in}}%
\pgfpathlineto{\pgfqpoint{0.567701in}{1.775221in}}%
\pgfpathlineto{\pgfqpoint{0.574289in}{1.797104in}}%
\pgfpathlineto{\pgfqpoint{0.579595in}{1.795500in}}%
\pgfpathlineto{\pgfqpoint{0.579759in}{1.803398in}}%
\pgfpathlineto{\pgfqpoint{0.585013in}{1.811079in}}%
\pgfpathlineto{\pgfqpoint{0.585552in}{1.821819in}}%
\pgfpathlineto{\pgfqpoint{0.583582in}{1.827974in}}%
\pgfpathlineto{\pgfqpoint{0.597616in}{1.841914in}}%
\pgfpathlineto{\pgfqpoint{0.599282in}{1.851636in}}%
\pgfpathlineto{\pgfqpoint{0.595548in}{1.853584in}}%
\pgfpathlineto{\pgfqpoint{0.595348in}{1.857817in}}%
\pgfpathlineto{\pgfqpoint{0.600903in}{1.863151in}}%
\pgfpathlineto{\pgfqpoint{0.610134in}{1.868414in}}%
\pgfpathlineto{\pgfqpoint{0.615087in}{1.868865in}}%
\pgfpathclose%
\pgfusepath{fill}%
\end{pgfscope}%
\begin{pgfscope}%
\pgfpathrectangle{\pgfqpoint{0.100000in}{0.100000in}}{\pgfqpoint{3.007045in}{1.925000in}}%
\pgfusepath{clip}%
\pgfsetbuttcap%
\pgfsetmiterjoin%
\definecolor{currentfill}{rgb}{0.351511,0.635848,0.812641}%
\pgfsetfillcolor{currentfill}%
\pgfsetlinewidth{0.000000pt}%
\definecolor{currentstroke}{rgb}{0.000000,0.000000,0.000000}%
\pgfsetstrokecolor{currentstroke}%
\pgfsetstrokeopacity{0.000000}%
\pgfsetdash{}{0pt}%
\pgfpathmoveto{\pgfqpoint{0.465761in}{1.347164in}}%
\pgfpathlineto{\pgfqpoint{0.460387in}{1.347970in}}%
\pgfpathlineto{\pgfqpoint{0.450179in}{1.340731in}}%
\pgfpathlineto{\pgfqpoint{0.445384in}{1.329930in}}%
\pgfpathlineto{\pgfqpoint{0.439260in}{1.328328in}}%
\pgfpathlineto{\pgfqpoint{0.435140in}{1.324290in}}%
\pgfpathlineto{\pgfqpoint{0.431828in}{1.313269in}}%
\pgfpathlineto{\pgfqpoint{0.424973in}{1.314267in}}%
\pgfpathlineto{\pgfqpoint{0.420345in}{1.320339in}}%
\pgfpathlineto{\pgfqpoint{0.421047in}{1.323928in}}%
\pgfpathlineto{\pgfqpoint{0.417135in}{1.329622in}}%
\pgfpathlineto{\pgfqpoint{0.392426in}{1.337198in}}%
\pgfpathlineto{\pgfqpoint{0.390745in}{1.343728in}}%
\pgfpathlineto{\pgfqpoint{0.387533in}{1.347636in}}%
\pgfpathlineto{\pgfqpoint{0.389453in}{1.357782in}}%
\pgfpathlineto{\pgfqpoint{0.383879in}{1.360469in}}%
\pgfpathlineto{\pgfqpoint{0.379169in}{1.366194in}}%
\pgfpathlineto{\pgfqpoint{0.383653in}{1.376276in}}%
\pgfpathlineto{\pgfqpoint{0.386077in}{1.384484in}}%
\pgfpathlineto{\pgfqpoint{0.378819in}{1.386738in}}%
\pgfpathlineto{\pgfqpoint{0.381027in}{1.394938in}}%
\pgfpathlineto{\pgfqpoint{0.380635in}{1.401231in}}%
\pgfpathlineto{\pgfqpoint{0.423677in}{1.387762in}}%
\pgfpathlineto{\pgfqpoint{0.425415in}{1.393166in}}%
\pgfpathlineto{\pgfqpoint{0.437071in}{1.389645in}}%
\pgfpathlineto{\pgfqpoint{0.443854in}{1.394476in}}%
\pgfpathlineto{\pgfqpoint{0.446584in}{1.393519in}}%
\pgfpathlineto{\pgfqpoint{0.451708in}{1.400095in}}%
\pgfpathlineto{\pgfqpoint{0.459905in}{1.401168in}}%
\pgfpathlineto{\pgfqpoint{0.461995in}{1.397568in}}%
\pgfpathlineto{\pgfqpoint{0.458164in}{1.392328in}}%
\pgfpathlineto{\pgfqpoint{0.455804in}{1.383109in}}%
\pgfpathlineto{\pgfqpoint{0.457960in}{1.381368in}}%
\pgfpathlineto{\pgfqpoint{0.467458in}{1.355133in}}%
\pgfpathlineto{\pgfqpoint{0.465761in}{1.347164in}}%
\pgfpathclose%
\pgfusepath{fill}%
\end{pgfscope}%
\begin{pgfscope}%
\pgfpathrectangle{\pgfqpoint{0.100000in}{0.100000in}}{\pgfqpoint{3.007045in}{1.925000in}}%
\pgfusepath{clip}%
\pgfsetbuttcap%
\pgfsetmiterjoin%
\definecolor{currentfill}{rgb}{0.341423,0.628958,0.808704}%
\pgfsetfillcolor{currentfill}%
\pgfsetlinewidth{0.000000pt}%
\definecolor{currentstroke}{rgb}{0.000000,0.000000,0.000000}%
\pgfsetstrokecolor{currentstroke}%
\pgfsetstrokeopacity{0.000000}%
\pgfsetdash{}{0pt}%
\pgfpathmoveto{\pgfqpoint{1.746105in}{1.095854in}}%
\pgfpathlineto{\pgfqpoint{1.745387in}{1.069764in}}%
\pgfpathlineto{\pgfqpoint{1.722556in}{1.070090in}}%
\pgfpathlineto{\pgfqpoint{1.722658in}{1.078703in}}%
\pgfpathlineto{\pgfqpoint{1.702042in}{1.079070in}}%
\pgfpathlineto{\pgfqpoint{1.701914in}{1.073348in}}%
\pgfpathlineto{\pgfqpoint{1.693377in}{1.073520in}}%
\pgfpathlineto{\pgfqpoint{1.676880in}{1.073769in}}%
\pgfpathlineto{\pgfqpoint{1.677388in}{1.091095in}}%
\pgfpathlineto{\pgfqpoint{1.678559in}{1.102524in}}%
\pgfpathlineto{\pgfqpoint{1.702564in}{1.102052in}}%
\pgfpathlineto{\pgfqpoint{1.702839in}{1.116383in}}%
\pgfpathlineto{\pgfqpoint{1.723731in}{1.116097in}}%
\pgfpathlineto{\pgfqpoint{1.723295in}{1.096039in}}%
\pgfpathlineto{\pgfqpoint{1.746105in}{1.095854in}}%
\pgfpathclose%
\pgfusepath{fill}%
\end{pgfscope}%
\begin{pgfscope}%
\pgfpathrectangle{\pgfqpoint{0.100000in}{0.100000in}}{\pgfqpoint{3.007045in}{1.925000in}}%
\pgfusepath{clip}%
\pgfsetbuttcap%
\pgfsetmiterjoin%
\definecolor{currentfill}{rgb}{0.429020,0.687520,0.841246}%
\pgfsetfillcolor{currentfill}%
\pgfsetlinewidth{0.000000pt}%
\definecolor{currentstroke}{rgb}{0.000000,0.000000,0.000000}%
\pgfsetstrokecolor{currentstroke}%
\pgfsetstrokeopacity{0.000000}%
\pgfsetdash{}{0pt}%
\pgfpathmoveto{\pgfqpoint{1.102971in}{1.695747in}}%
\pgfpathlineto{\pgfqpoint{1.099659in}{1.696257in}}%
\pgfpathlineto{\pgfqpoint{1.093536in}{1.701272in}}%
\pgfpathlineto{\pgfqpoint{1.070809in}{1.700125in}}%
\pgfpathlineto{\pgfqpoint{1.066660in}{1.704559in}}%
\pgfpathlineto{\pgfqpoint{1.059670in}{1.716732in}}%
\pgfpathlineto{\pgfqpoint{1.053818in}{1.723481in}}%
\pgfpathlineto{\pgfqpoint{1.043222in}{1.727442in}}%
\pgfpathlineto{\pgfqpoint{1.026102in}{1.729804in}}%
\pgfpathlineto{\pgfqpoint{1.025046in}{1.724117in}}%
\pgfpathlineto{\pgfqpoint{1.019738in}{1.725093in}}%
\pgfpathlineto{\pgfqpoint{1.013851in}{1.726227in}}%
\pgfpathlineto{\pgfqpoint{1.016513in}{1.740285in}}%
\pgfpathlineto{\pgfqpoint{1.010013in}{1.742368in}}%
\pgfpathlineto{\pgfqpoint{1.005933in}{1.746597in}}%
\pgfpathlineto{\pgfqpoint{1.009841in}{1.767764in}}%
\pgfpathlineto{\pgfqpoint{0.999503in}{1.769057in}}%
\pgfpathlineto{\pgfqpoint{0.994442in}{1.773321in}}%
\pgfpathlineto{\pgfqpoint{0.994400in}{1.776622in}}%
\pgfpathlineto{\pgfqpoint{0.987897in}{1.779953in}}%
\pgfpathlineto{\pgfqpoint{0.976971in}{1.780376in}}%
\pgfpathlineto{\pgfqpoint{0.973893in}{1.786912in}}%
\pgfpathlineto{\pgfqpoint{0.975711in}{1.795488in}}%
\pgfpathlineto{\pgfqpoint{0.974188in}{1.804466in}}%
\pgfpathlineto{\pgfqpoint{0.978555in}{1.808720in}}%
\pgfpathlineto{\pgfqpoint{0.975124in}{1.816423in}}%
\pgfpathlineto{\pgfqpoint{1.011784in}{1.809131in}}%
\pgfpathlineto{\pgfqpoint{1.016942in}{1.805156in}}%
\pgfpathlineto{\pgfqpoint{1.018457in}{1.797935in}}%
\pgfpathlineto{\pgfqpoint{1.043735in}{1.793392in}}%
\pgfpathlineto{\pgfqpoint{1.045450in}{1.802919in}}%
\pgfpathlineto{\pgfqpoint{1.070974in}{1.798414in}}%
\pgfpathlineto{\pgfqpoint{1.071976in}{1.804099in}}%
\pgfpathlineto{\pgfqpoint{1.080804in}{1.802553in}}%
\pgfpathlineto{\pgfqpoint{1.081806in}{1.808264in}}%
\pgfpathlineto{\pgfqpoint{1.115936in}{1.802454in}}%
\pgfpathlineto{\pgfqpoint{1.115258in}{1.796858in}}%
\pgfpathlineto{\pgfqpoint{1.120232in}{1.796028in}}%
\pgfpathlineto{\pgfqpoint{1.119270in}{1.790180in}}%
\pgfpathlineto{\pgfqpoint{1.127771in}{1.788966in}}%
\pgfpathlineto{\pgfqpoint{1.123379in}{1.762908in}}%
\pgfpathlineto{\pgfqpoint{1.110393in}{1.763495in}}%
\pgfpathlineto{\pgfqpoint{1.107591in}{1.752201in}}%
\pgfpathlineto{\pgfqpoint{1.094250in}{1.746968in}}%
\pgfpathlineto{\pgfqpoint{1.092154in}{1.736531in}}%
\pgfpathlineto{\pgfqpoint{1.097369in}{1.732746in}}%
\pgfpathlineto{\pgfqpoint{1.103070in}{1.731808in}}%
\pgfpathlineto{\pgfqpoint{1.107755in}{1.728110in}}%
\pgfpathlineto{\pgfqpoint{1.104710in}{1.709101in}}%
\pgfpathlineto{\pgfqpoint{1.102514in}{1.707496in}}%
\pgfpathlineto{\pgfqpoint{1.102971in}{1.695747in}}%
\pgfpathclose%
\pgfusepath{fill}%
\end{pgfscope}%
\begin{pgfscope}%
\pgfpathrectangle{\pgfqpoint{0.100000in}{0.100000in}}{\pgfqpoint{3.007045in}{1.925000in}}%
\pgfusepath{clip}%
\pgfsetbuttcap%
\pgfsetmiterjoin%
\definecolor{currentfill}{rgb}{0.285936,0.591065,0.787051}%
\pgfsetfillcolor{currentfill}%
\pgfsetlinewidth{0.000000pt}%
\definecolor{currentstroke}{rgb}{0.000000,0.000000,0.000000}%
\pgfsetstrokecolor{currentstroke}%
\pgfsetstrokeopacity{0.000000}%
\pgfsetdash{}{0pt}%
\pgfpathmoveto{\pgfqpoint{1.522557in}{1.725806in}}%
\pgfpathlineto{\pgfqpoint{1.539856in}{1.724787in}}%
\pgfpathlineto{\pgfqpoint{1.538782in}{1.736495in}}%
\pgfpathlineto{\pgfqpoint{1.540108in}{1.759776in}}%
\pgfpathlineto{\pgfqpoint{1.539033in}{1.771449in}}%
\pgfpathlineto{\pgfqpoint{1.544810in}{1.771147in}}%
\pgfpathlineto{\pgfqpoint{1.573646in}{1.769590in}}%
\pgfpathlineto{\pgfqpoint{1.573086in}{1.757987in}}%
\pgfpathlineto{\pgfqpoint{1.557416in}{1.758816in}}%
\pgfpathlineto{\pgfqpoint{1.556105in}{1.735501in}}%
\pgfpathlineto{\pgfqpoint{1.557103in}{1.723835in}}%
\pgfpathlineto{\pgfqpoint{1.580095in}{1.722634in}}%
\pgfpathlineto{\pgfqpoint{1.580956in}{1.710973in}}%
\pgfpathlineto{\pgfqpoint{1.579862in}{1.687980in}}%
\pgfpathlineto{\pgfqpoint{1.570193in}{1.688446in}}%
\pgfpathlineto{\pgfqpoint{1.541684in}{1.689945in}}%
\pgfpathlineto{\pgfqpoint{1.516455in}{1.691466in}}%
\pgfpathlineto{\pgfqpoint{1.517941in}{1.714507in}}%
\pgfpathlineto{\pgfqpoint{1.521849in}{1.714257in}}%
\pgfpathlineto{\pgfqpoint{1.522557in}{1.725806in}}%
\pgfpathclose%
\pgfusepath{fill}%
\end{pgfscope}%
\begin{pgfscope}%
\pgfpathrectangle{\pgfqpoint{0.100000in}{0.100000in}}{\pgfqpoint{3.007045in}{1.925000in}}%
\pgfusepath{clip}%
\pgfsetbuttcap%
\pgfsetmiterjoin%
\definecolor{currentfill}{rgb}{0.406997,0.673741,0.834295}%
\pgfsetfillcolor{currentfill}%
\pgfsetlinewidth{0.000000pt}%
\definecolor{currentstroke}{rgb}{0.000000,0.000000,0.000000}%
\pgfsetstrokecolor{currentstroke}%
\pgfsetstrokeopacity{0.000000}%
\pgfsetdash{}{0pt}%
\pgfpathmoveto{\pgfqpoint{2.024381in}{1.212989in}}%
\pgfpathlineto{\pgfqpoint{2.014487in}{1.208775in}}%
\pgfpathlineto{\pgfqpoint{2.009474in}{1.209217in}}%
\pgfpathlineto{\pgfqpoint{2.008276in}{1.212623in}}%
\pgfpathlineto{\pgfqpoint{2.013813in}{1.217404in}}%
\pgfpathlineto{\pgfqpoint{2.017601in}{1.220842in}}%
\pgfpathlineto{\pgfqpoint{2.021248in}{1.230411in}}%
\pgfpathlineto{\pgfqpoint{2.027473in}{1.235391in}}%
\pgfpathlineto{\pgfqpoint{2.029027in}{1.239956in}}%
\pgfpathlineto{\pgfqpoint{2.028561in}{1.247285in}}%
\pgfpathlineto{\pgfqpoint{2.022785in}{1.247040in}}%
\pgfpathlineto{\pgfqpoint{2.021691in}{1.270010in}}%
\pgfpathlineto{\pgfqpoint{2.038925in}{1.270934in}}%
\pgfpathlineto{\pgfqpoint{2.037908in}{1.288046in}}%
\pgfpathlineto{\pgfqpoint{2.027069in}{1.287453in}}%
\pgfpathlineto{\pgfqpoint{2.025542in}{1.310434in}}%
\pgfpathlineto{\pgfqpoint{2.059792in}{1.312587in}}%
\pgfpathlineto{\pgfqpoint{2.059573in}{1.315459in}}%
\pgfpathlineto{\pgfqpoint{2.087331in}{1.317506in}}%
\pgfpathlineto{\pgfqpoint{2.090562in}{1.283138in}}%
\pgfpathlineto{\pgfqpoint{2.073487in}{1.281826in}}%
\pgfpathlineto{\pgfqpoint{2.075117in}{1.258653in}}%
\pgfpathlineto{\pgfqpoint{2.072749in}{1.252662in}}%
\pgfpathlineto{\pgfqpoint{2.065605in}{1.247690in}}%
\pgfpathlineto{\pgfqpoint{2.058847in}{1.247095in}}%
\pgfpathlineto{\pgfqpoint{2.060239in}{1.229418in}}%
\pgfpathlineto{\pgfqpoint{2.043262in}{1.228036in}}%
\pgfpathlineto{\pgfqpoint{2.044100in}{1.215046in}}%
\pgfpathlineto{\pgfqpoint{2.039739in}{1.216849in}}%
\pgfpathlineto{\pgfqpoint{2.032407in}{1.214497in}}%
\pgfpathlineto{\pgfqpoint{2.025473in}{1.215274in}}%
\pgfpathlineto{\pgfqpoint{2.024381in}{1.212989in}}%
\pgfpathclose%
\pgfusepath{fill}%
\end{pgfscope}%
\begin{pgfscope}%
\pgfpathrectangle{\pgfqpoint{0.100000in}{0.100000in}}{\pgfqpoint{3.007045in}{1.925000in}}%
\pgfusepath{clip}%
\pgfsetbuttcap%
\pgfsetmiterjoin%
\definecolor{currentfill}{rgb}{0.195386,0.509112,0.743791}%
\pgfsetfillcolor{currentfill}%
\pgfsetlinewidth{0.000000pt}%
\definecolor{currentstroke}{rgb}{0.000000,0.000000,0.000000}%
\pgfsetstrokecolor{currentstroke}%
\pgfsetstrokeopacity{0.000000}%
\pgfsetdash{}{0pt}%
\pgfpathmoveto{\pgfqpoint{1.553171in}{1.527805in}}%
\pgfpathlineto{\pgfqpoint{1.557925in}{1.527548in}}%
\pgfpathlineto{\pgfqpoint{1.598571in}{1.525674in}}%
\pgfpathlineto{\pgfqpoint{1.597846in}{1.508282in}}%
\pgfpathlineto{\pgfqpoint{1.596903in}{1.479359in}}%
\pgfpathlineto{\pgfqpoint{1.586215in}{1.479799in}}%
\pgfpathlineto{\pgfqpoint{1.555401in}{1.480969in}}%
\pgfpathlineto{\pgfqpoint{1.551371in}{1.482830in}}%
\pgfpathlineto{\pgfqpoint{1.551887in}{1.504699in}}%
\pgfpathlineto{\pgfqpoint{1.553171in}{1.527805in}}%
\pgfpathclose%
\pgfusepath{fill}%
\end{pgfscope}%
\begin{pgfscope}%
\pgfpathrectangle{\pgfqpoint{0.100000in}{0.100000in}}{\pgfqpoint{3.007045in}{1.925000in}}%
\pgfusepath{clip}%
\pgfsetbuttcap%
\pgfsetmiterjoin%
\definecolor{currentfill}{rgb}{0.321246,0.615179,0.800830}%
\pgfsetfillcolor{currentfill}%
\pgfsetlinewidth{0.000000pt}%
\definecolor{currentstroke}{rgb}{0.000000,0.000000,0.000000}%
\pgfsetstrokecolor{currentstroke}%
\pgfsetstrokeopacity{0.000000}%
\pgfsetdash{}{0pt}%
\pgfpathmoveto{\pgfqpoint{1.963991in}{0.607879in}}%
\pgfpathlineto{\pgfqpoint{1.967968in}{0.612514in}}%
\pgfpathlineto{\pgfqpoint{1.964384in}{0.616831in}}%
\pgfpathlineto{\pgfqpoint{1.967749in}{0.623734in}}%
\pgfpathlineto{\pgfqpoint{1.971855in}{0.623291in}}%
\pgfpathlineto{\pgfqpoint{1.972763in}{0.629511in}}%
\pgfpathlineto{\pgfqpoint{1.985698in}{0.633433in}}%
\pgfpathlineto{\pgfqpoint{1.993008in}{0.630363in}}%
\pgfpathlineto{\pgfqpoint{1.998230in}{0.632422in}}%
\pgfpathlineto{\pgfqpoint{2.020031in}{0.633536in}}%
\pgfpathlineto{\pgfqpoint{2.049764in}{0.635164in}}%
\pgfpathlineto{\pgfqpoint{2.053466in}{0.638280in}}%
\pgfpathlineto{\pgfqpoint{2.058100in}{0.638257in}}%
\pgfpathlineto{\pgfqpoint{2.065780in}{0.638658in}}%
\pgfpathlineto{\pgfqpoint{2.065550in}{0.641578in}}%
\pgfpathlineto{\pgfqpoint{2.075185in}{0.642117in}}%
\pgfpathlineto{\pgfqpoint{2.076881in}{0.613497in}}%
\pgfpathlineto{\pgfqpoint{2.072666in}{0.613237in}}%
\pgfpathlineto{\pgfqpoint{2.037435in}{0.611183in}}%
\pgfpathlineto{\pgfqpoint{2.010177in}{0.609791in}}%
\pgfpathlineto{\pgfqpoint{1.963991in}{0.607879in}}%
\pgfpathclose%
\pgfusepath{fill}%
\end{pgfscope}%
\begin{pgfscope}%
\pgfpathrectangle{\pgfqpoint{0.100000in}{0.100000in}}{\pgfqpoint{3.007045in}{1.925000in}}%
\pgfusepath{clip}%
\pgfsetbuttcap%
\pgfsetmiterjoin%
\definecolor{currentfill}{rgb}{0.401953,0.670296,0.832326}%
\pgfsetfillcolor{currentfill}%
\pgfsetlinewidth{0.000000pt}%
\definecolor{currentstroke}{rgb}{0.000000,0.000000,0.000000}%
\pgfsetstrokecolor{currentstroke}%
\pgfsetstrokeopacity{0.000000}%
\pgfsetdash{}{0pt}%
\pgfpathmoveto{\pgfqpoint{1.867818in}{1.391658in}}%
\pgfpathlineto{\pgfqpoint{1.866993in}{1.430735in}}%
\pgfpathlineto{\pgfqpoint{1.894560in}{1.431480in}}%
\pgfpathlineto{\pgfqpoint{1.928977in}{1.432597in}}%
\pgfpathlineto{\pgfqpoint{1.934699in}{1.432796in}}%
\pgfpathlineto{\pgfqpoint{1.936000in}{1.405159in}}%
\pgfpathlineto{\pgfqpoint{1.913064in}{1.404438in}}%
\pgfpathlineto{\pgfqpoint{1.913408in}{1.392837in}}%
\pgfpathlineto{\pgfqpoint{1.867818in}{1.391658in}}%
\pgfpathclose%
\pgfusepath{fill}%
\end{pgfscope}%
\begin{pgfscope}%
\pgfpathrectangle{\pgfqpoint{0.100000in}{0.100000in}}{\pgfqpoint{3.007045in}{1.925000in}}%
\pgfusepath{clip}%
\pgfsetbuttcap%
\pgfsetmiterjoin%
\definecolor{currentfill}{rgb}{0.592157,0.777086,0.876432}%
\pgfsetfillcolor{currentfill}%
\pgfsetlinewidth{0.000000pt}%
\definecolor{currentstroke}{rgb}{0.000000,0.000000,0.000000}%
\pgfsetstrokecolor{currentstroke}%
\pgfsetstrokeopacity{0.000000}%
\pgfsetdash{}{0pt}%
\pgfpathmoveto{\pgfqpoint{2.608126in}{1.441492in}}%
\pgfpathlineto{\pgfqpoint{2.601922in}{1.443430in}}%
\pgfpathlineto{\pgfqpoint{2.572079in}{1.437927in}}%
\pgfpathlineto{\pgfqpoint{2.570755in}{1.438725in}}%
\pgfpathlineto{\pgfqpoint{2.557784in}{1.431044in}}%
\pgfpathlineto{\pgfqpoint{2.552432in}{1.429022in}}%
\pgfpathlineto{\pgfqpoint{2.546681in}{1.434064in}}%
\pgfpathlineto{\pgfqpoint{2.543334in}{1.433990in}}%
\pgfpathlineto{\pgfqpoint{2.539326in}{1.435418in}}%
\pgfpathlineto{\pgfqpoint{2.542194in}{1.444103in}}%
\pgfpathlineto{\pgfqpoint{2.547820in}{1.448105in}}%
\pgfpathlineto{\pgfqpoint{2.550364in}{1.452200in}}%
\pgfpathlineto{\pgfqpoint{2.544628in}{1.462243in}}%
\pgfpathlineto{\pgfqpoint{2.540521in}{1.463483in}}%
\pgfpathlineto{\pgfqpoint{2.539800in}{1.468753in}}%
\pgfpathlineto{\pgfqpoint{2.535305in}{1.480591in}}%
\pgfpathlineto{\pgfqpoint{2.544988in}{1.486580in}}%
\pgfpathlineto{\pgfqpoint{2.553954in}{1.490859in}}%
\pgfpathlineto{\pgfqpoint{2.565267in}{1.494184in}}%
\pgfpathlineto{\pgfqpoint{2.584950in}{1.497319in}}%
\pgfpathlineto{\pgfqpoint{2.595524in}{1.497454in}}%
\pgfpathlineto{\pgfqpoint{2.605383in}{1.492846in}}%
\pgfpathlineto{\pgfqpoint{2.614787in}{1.497022in}}%
\pgfpathlineto{\pgfqpoint{2.634861in}{1.500680in}}%
\pgfpathlineto{\pgfqpoint{2.644415in}{1.507385in}}%
\pgfpathlineto{\pgfqpoint{2.648812in}{1.490538in}}%
\pgfpathlineto{\pgfqpoint{2.648634in}{1.483094in}}%
\pgfpathlineto{\pgfqpoint{2.652395in}{1.466521in}}%
\pgfpathlineto{\pgfqpoint{2.656543in}{1.461343in}}%
\pgfpathlineto{\pgfqpoint{2.661352in}{1.457374in}}%
\pgfpathlineto{\pgfqpoint{2.646021in}{1.453686in}}%
\pgfpathlineto{\pgfqpoint{2.647379in}{1.448629in}}%
\pgfpathlineto{\pgfqpoint{2.635144in}{1.448720in}}%
\pgfpathlineto{\pgfqpoint{2.634427in}{1.453744in}}%
\pgfpathlineto{\pgfqpoint{2.609773in}{1.449142in}}%
\pgfpathlineto{\pgfqpoint{2.608126in}{1.441492in}}%
\pgfpathclose%
\pgfusepath{fill}%
\end{pgfscope}%
\begin{pgfscope}%
\pgfpathrectangle{\pgfqpoint{0.100000in}{0.100000in}}{\pgfqpoint{3.007045in}{1.925000in}}%
\pgfusepath{clip}%
\pgfsetbuttcap%
\pgfsetmiterjoin%
\definecolor{currentfill}{rgb}{0.676817,0.816471,0.902376}%
\pgfsetfillcolor{currentfill}%
\pgfsetlinewidth{0.000000pt}%
\definecolor{currentstroke}{rgb}{0.000000,0.000000,0.000000}%
\pgfsetstrokecolor{currentstroke}%
\pgfsetstrokeopacity{0.000000}%
\pgfsetdash{}{0pt}%
\pgfpathmoveto{\pgfqpoint{2.507767in}{1.239655in}}%
\pgfpathlineto{\pgfqpoint{2.502865in}{1.238879in}}%
\pgfpathlineto{\pgfqpoint{2.498296in}{1.267386in}}%
\pgfpathlineto{\pgfqpoint{2.493292in}{1.298554in}}%
\pgfpathlineto{\pgfqpoint{2.491069in}{1.312382in}}%
\pgfpathlineto{\pgfqpoint{2.508763in}{1.315516in}}%
\pgfpathlineto{\pgfqpoint{2.506874in}{1.324911in}}%
\pgfpathlineto{\pgfqpoint{2.513140in}{1.337398in}}%
\pgfpathlineto{\pgfqpoint{2.528039in}{1.339978in}}%
\pgfpathlineto{\pgfqpoint{2.533624in}{1.328350in}}%
\pgfpathlineto{\pgfqpoint{2.540085in}{1.329055in}}%
\pgfpathlineto{\pgfqpoint{2.550029in}{1.332847in}}%
\pgfpathlineto{\pgfqpoint{2.552895in}{1.336088in}}%
\pgfpathlineto{\pgfqpoint{2.554469in}{1.327047in}}%
\pgfpathlineto{\pgfqpoint{2.574642in}{1.330275in}}%
\pgfpathlineto{\pgfqpoint{2.576932in}{1.318538in}}%
\pgfpathlineto{\pgfqpoint{2.574047in}{1.302973in}}%
\pgfpathlineto{\pgfqpoint{2.567920in}{1.283023in}}%
\pgfpathlineto{\pgfqpoint{2.561814in}{1.273283in}}%
\pgfpathlineto{\pgfqpoint{2.556847in}{1.259865in}}%
\pgfpathlineto{\pgfqpoint{2.559394in}{1.254953in}}%
\pgfpathlineto{\pgfqpoint{2.559174in}{1.248298in}}%
\pgfpathlineto{\pgfqpoint{2.554976in}{1.247539in}}%
\pgfpathlineto{\pgfqpoint{2.507767in}{1.239655in}}%
\pgfpathclose%
\pgfusepath{fill}%
\end{pgfscope}%
\begin{pgfscope}%
\pgfpathrectangle{\pgfqpoint{0.100000in}{0.100000in}}{\pgfqpoint{3.007045in}{1.925000in}}%
\pgfusepath{clip}%
\pgfsetbuttcap%
\pgfsetmiterjoin%
\definecolor{currentfill}{rgb}{0.610980,0.787420,0.880492}%
\pgfsetfillcolor{currentfill}%
\pgfsetlinewidth{0.000000pt}%
\definecolor{currentstroke}{rgb}{0.000000,0.000000,0.000000}%
\pgfsetstrokecolor{currentstroke}%
\pgfsetstrokeopacity{0.000000}%
\pgfsetdash{}{0pt}%
\pgfpathmoveto{\pgfqpoint{1.506378in}{0.538351in}}%
\pgfpathlineto{\pgfqpoint{1.501272in}{0.538689in}}%
\pgfpathlineto{\pgfqpoint{1.501924in}{0.549696in}}%
\pgfpathlineto{\pgfqpoint{1.498020in}{0.549906in}}%
\pgfpathlineto{\pgfqpoint{1.498977in}{0.564304in}}%
\pgfpathlineto{\pgfqpoint{1.524943in}{0.562745in}}%
\pgfpathlineto{\pgfqpoint{1.525500in}{0.576932in}}%
\pgfpathlineto{\pgfqpoint{1.544968in}{0.575956in}}%
\pgfpathlineto{\pgfqpoint{1.566327in}{0.575160in}}%
\pgfpathlineto{\pgfqpoint{1.565594in}{0.551126in}}%
\pgfpathlineto{\pgfqpoint{1.546459in}{0.551872in}}%
\pgfpathlineto{\pgfqpoint{1.545579in}{0.528102in}}%
\pgfpathlineto{\pgfqpoint{1.530916in}{0.536323in}}%
\pgfpathlineto{\pgfqpoint{1.506378in}{0.538351in}}%
\pgfpathclose%
\pgfusepath{fill}%
\end{pgfscope}%
\begin{pgfscope}%
\pgfpathrectangle{\pgfqpoint{0.100000in}{0.100000in}}{\pgfqpoint{3.007045in}{1.925000in}}%
\pgfusepath{clip}%
\pgfsetbuttcap%
\pgfsetmiterjoin%
\definecolor{currentfill}{rgb}{0.386820,0.659962,0.826421}%
\pgfsetfillcolor{currentfill}%
\pgfsetlinewidth{0.000000pt}%
\definecolor{currentstroke}{rgb}{0.000000,0.000000,0.000000}%
\pgfsetstrokecolor{currentstroke}%
\pgfsetstrokeopacity{0.000000}%
\pgfsetdash{}{0pt}%
\pgfpathmoveto{\pgfqpoint{2.408663in}{0.823968in}}%
\pgfpathlineto{\pgfqpoint{2.414961in}{0.829005in}}%
\pgfpathlineto{\pgfqpoint{2.424340in}{0.826019in}}%
\pgfpathlineto{\pgfqpoint{2.430059in}{0.826283in}}%
\pgfpathlineto{\pgfqpoint{2.430794in}{0.832351in}}%
\pgfpathlineto{\pgfqpoint{2.435742in}{0.842463in}}%
\pgfpathlineto{\pgfqpoint{2.440442in}{0.846333in}}%
\pgfpathlineto{\pgfqpoint{2.447000in}{0.848563in}}%
\pgfpathlineto{\pgfqpoint{2.452246in}{0.846939in}}%
\pgfpathlineto{\pgfqpoint{2.468037in}{0.838774in}}%
\pgfpathlineto{\pgfqpoint{2.473536in}{0.831913in}}%
\pgfpathlineto{\pgfqpoint{2.468621in}{0.832178in}}%
\pgfpathlineto{\pgfqpoint{2.460350in}{0.826839in}}%
\pgfpathlineto{\pgfqpoint{2.452424in}{0.825463in}}%
\pgfpathlineto{\pgfqpoint{2.449159in}{0.822767in}}%
\pgfpathlineto{\pgfqpoint{2.446193in}{0.816402in}}%
\pgfpathlineto{\pgfqpoint{2.442122in}{0.815272in}}%
\pgfpathlineto{\pgfqpoint{2.441136in}{0.811210in}}%
\pgfpathlineto{\pgfqpoint{2.432103in}{0.811727in}}%
\pgfpathlineto{\pgfqpoint{2.424800in}{0.803148in}}%
\pgfpathlineto{\pgfqpoint{2.417379in}{0.810753in}}%
\pgfpathlineto{\pgfqpoint{2.408663in}{0.823968in}}%
\pgfpathclose%
\pgfusepath{fill}%
\end{pgfscope}%
\begin{pgfscope}%
\pgfpathrectangle{\pgfqpoint{0.100000in}{0.100000in}}{\pgfqpoint{3.007045in}{1.925000in}}%
\pgfusepath{clip}%
\pgfsetbuttcap%
\pgfsetmiterjoin%
\definecolor{currentfill}{rgb}{0.447843,0.697855,0.845306}%
\pgfsetfillcolor{currentfill}%
\pgfsetlinewidth{0.000000pt}%
\definecolor{currentstroke}{rgb}{0.000000,0.000000,0.000000}%
\pgfsetstrokecolor{currentstroke}%
\pgfsetstrokeopacity{0.000000}%
\pgfsetdash{}{0pt}%
\pgfpathmoveto{\pgfqpoint{1.854945in}{1.201835in}}%
\pgfpathlineto{\pgfqpoint{1.855015in}{1.197513in}}%
\pgfpathlineto{\pgfqpoint{1.834849in}{1.196667in}}%
\pgfpathlineto{\pgfqpoint{1.834576in}{1.208087in}}%
\pgfpathlineto{\pgfqpoint{1.833661in}{1.237372in}}%
\pgfpathlineto{\pgfqpoint{1.844539in}{1.237709in}}%
\pgfpathlineto{\pgfqpoint{1.853670in}{1.237862in}}%
\pgfpathlineto{\pgfqpoint{1.854945in}{1.201835in}}%
\pgfpathclose%
\pgfusepath{fill}%
\end{pgfscope}%
\begin{pgfscope}%
\pgfpathrectangle{\pgfqpoint{0.100000in}{0.100000in}}{\pgfqpoint{3.007045in}{1.925000in}}%
\pgfusepath{clip}%
\pgfsetbuttcap%
\pgfsetmiterjoin%
\definecolor{currentfill}{rgb}{0.472941,0.711634,0.850719}%
\pgfsetfillcolor{currentfill}%
\pgfsetlinewidth{0.000000pt}%
\definecolor{currentstroke}{rgb}{0.000000,0.000000,0.000000}%
\pgfsetstrokecolor{currentstroke}%
\pgfsetstrokeopacity{0.000000}%
\pgfsetdash{}{0pt}%
\pgfpathmoveto{\pgfqpoint{1.538952in}{0.605702in}}%
\pgfpathlineto{\pgfqpoint{1.516589in}{0.606731in}}%
\pgfpathlineto{\pgfqpoint{1.509816in}{0.607078in}}%
\pgfpathlineto{\pgfqpoint{1.510333in}{0.616812in}}%
\pgfpathlineto{\pgfqpoint{1.481225in}{0.618416in}}%
\pgfpathlineto{\pgfqpoint{1.483290in}{0.651066in}}%
\pgfpathlineto{\pgfqpoint{1.476325in}{0.651539in}}%
\pgfpathlineto{\pgfqpoint{1.478210in}{0.684741in}}%
\pgfpathlineto{\pgfqpoint{1.482882in}{0.684491in}}%
\pgfpathlineto{\pgfqpoint{1.507461in}{0.683119in}}%
\pgfpathlineto{\pgfqpoint{1.505217in}{0.648714in}}%
\pgfpathlineto{\pgfqpoint{1.510361in}{0.642115in}}%
\pgfpathlineto{\pgfqpoint{1.520806in}{0.642095in}}%
\pgfpathlineto{\pgfqpoint{1.527773in}{0.637409in}}%
\pgfpathlineto{\pgfqpoint{1.532433in}{0.642045in}}%
\pgfpathlineto{\pgfqpoint{1.540604in}{0.640218in}}%
\pgfpathlineto{\pgfqpoint{1.538952in}{0.605702in}}%
\pgfpathclose%
\pgfusepath{fill}%
\end{pgfscope}%
\begin{pgfscope}%
\pgfpathrectangle{\pgfqpoint{0.100000in}{0.100000in}}{\pgfqpoint{3.007045in}{1.925000in}}%
\pgfusepath{clip}%
\pgfsetbuttcap%
\pgfsetmiterjoin%
\definecolor{currentfill}{rgb}{0.265759,0.577286,0.779177}%
\pgfsetfillcolor{currentfill}%
\pgfsetlinewidth{0.000000pt}%
\definecolor{currentstroke}{rgb}{0.000000,0.000000,0.000000}%
\pgfsetstrokecolor{currentstroke}%
\pgfsetstrokeopacity{0.000000}%
\pgfsetdash{}{0pt}%
\pgfpathmoveto{\pgfqpoint{2.076881in}{0.613497in}}%
\pgfpathlineto{\pgfqpoint{2.075185in}{0.642117in}}%
\pgfpathlineto{\pgfqpoint{2.065550in}{0.641578in}}%
\pgfpathlineto{\pgfqpoint{2.065780in}{0.638658in}}%
\pgfpathlineto{\pgfqpoint{2.058100in}{0.638257in}}%
\pgfpathlineto{\pgfqpoint{2.056646in}{0.646958in}}%
\pgfpathlineto{\pgfqpoint{2.055851in}{0.662822in}}%
\pgfpathlineto{\pgfqpoint{2.073893in}{0.665116in}}%
\pgfpathlineto{\pgfqpoint{2.088016in}{0.667046in}}%
\pgfpathlineto{\pgfqpoint{2.089615in}{0.643002in}}%
\pgfpathlineto{\pgfqpoint{2.121151in}{0.645024in}}%
\pgfpathlineto{\pgfqpoint{2.123918in}{0.610337in}}%
\pgfpathlineto{\pgfqpoint{2.095084in}{0.608430in}}%
\pgfpathlineto{\pgfqpoint{2.094715in}{0.614481in}}%
\pgfpathlineto{\pgfqpoint{2.085520in}{0.614773in}}%
\pgfpathlineto{\pgfqpoint{2.076881in}{0.613497in}}%
\pgfpathclose%
\pgfusepath{fill}%
\end{pgfscope}%
\begin{pgfscope}%
\pgfpathrectangle{\pgfqpoint{0.100000in}{0.100000in}}{\pgfqpoint{3.007045in}{1.925000in}}%
\pgfusepath{clip}%
\pgfsetbuttcap%
\pgfsetmiterjoin%
\definecolor{currentfill}{rgb}{0.376732,0.653072,0.822484}%
\pgfsetfillcolor{currentfill}%
\pgfsetlinewidth{0.000000pt}%
\definecolor{currentstroke}{rgb}{0.000000,0.000000,0.000000}%
\pgfsetstrokecolor{currentstroke}%
\pgfsetstrokeopacity{0.000000}%
\pgfsetdash{}{0pt}%
\pgfpathmoveto{\pgfqpoint{2.624882in}{1.186857in}}%
\pgfpathlineto{\pgfqpoint{2.625996in}{1.195904in}}%
\pgfpathlineto{\pgfqpoint{2.620005in}{1.199545in}}%
\pgfpathlineto{\pgfqpoint{2.622986in}{1.207692in}}%
\pgfpathlineto{\pgfqpoint{2.621536in}{1.211734in}}%
\pgfpathlineto{\pgfqpoint{2.634232in}{1.214856in}}%
\pgfpathlineto{\pgfqpoint{2.637469in}{1.212817in}}%
\pgfpathlineto{\pgfqpoint{2.643347in}{1.220746in}}%
\pgfpathlineto{\pgfqpoint{2.644431in}{1.224390in}}%
\pgfpathlineto{\pgfqpoint{2.646870in}{1.236936in}}%
\pgfpathlineto{\pgfqpoint{2.649600in}{1.238259in}}%
\pgfpathlineto{\pgfqpoint{2.651264in}{1.258406in}}%
\pgfpathlineto{\pgfqpoint{2.655502in}{1.265959in}}%
\pgfpathlineto{\pgfqpoint{2.667563in}{1.268330in}}%
\pgfpathlineto{\pgfqpoint{2.664785in}{1.260556in}}%
\pgfpathlineto{\pgfqpoint{2.672090in}{1.258055in}}%
\pgfpathlineto{\pgfqpoint{2.675901in}{1.254838in}}%
\pgfpathlineto{\pgfqpoint{2.673952in}{1.243959in}}%
\pgfpathlineto{\pgfqpoint{2.677394in}{1.240185in}}%
\pgfpathlineto{\pgfqpoint{2.684505in}{1.237247in}}%
\pgfpathlineto{\pgfqpoint{2.686960in}{1.233890in}}%
\pgfpathlineto{\pgfqpoint{2.703023in}{1.225995in}}%
\pgfpathlineto{\pgfqpoint{2.705509in}{1.221124in}}%
\pgfpathlineto{\pgfqpoint{2.704557in}{1.214940in}}%
\pgfpathlineto{\pgfqpoint{2.706153in}{1.210930in}}%
\pgfpathlineto{\pgfqpoint{2.715119in}{1.210272in}}%
\pgfpathlineto{\pgfqpoint{2.717999in}{1.199175in}}%
\pgfpathlineto{\pgfqpoint{2.726182in}{1.190843in}}%
\pgfpathlineto{\pgfqpoint{2.729054in}{1.179712in}}%
\pgfpathlineto{\pgfqpoint{2.733376in}{1.175355in}}%
\pgfpathlineto{\pgfqpoint{2.732254in}{1.169909in}}%
\pgfpathlineto{\pgfqpoint{2.727131in}{1.173840in}}%
\pgfpathlineto{\pgfqpoint{2.714351in}{1.177897in}}%
\pgfpathlineto{\pgfqpoint{2.708831in}{1.176381in}}%
\pgfpathlineto{\pgfqpoint{2.695566in}{1.180728in}}%
\pgfpathlineto{\pgfqpoint{2.691552in}{1.182331in}}%
\pgfpathlineto{\pgfqpoint{2.682867in}{1.178238in}}%
\pgfpathlineto{\pgfqpoint{2.680316in}{1.180573in}}%
\pgfpathlineto{\pgfqpoint{2.678549in}{1.188294in}}%
\pgfpathlineto{\pgfqpoint{2.668861in}{1.192884in}}%
\pgfpathlineto{\pgfqpoint{2.662595in}{1.183012in}}%
\pgfpathlineto{\pgfqpoint{2.663143in}{1.179525in}}%
\pgfpathlineto{\pgfqpoint{2.658001in}{1.178784in}}%
\pgfpathlineto{\pgfqpoint{2.648785in}{1.188481in}}%
\pgfpathlineto{\pgfqpoint{2.644402in}{1.195106in}}%
\pgfpathlineto{\pgfqpoint{2.634662in}{1.181683in}}%
\pgfpathlineto{\pgfqpoint{2.631429in}{1.181706in}}%
\pgfpathlineto{\pgfqpoint{2.624882in}{1.186857in}}%
\pgfpathclose%
\pgfusepath{fill}%
\end{pgfscope}%
\begin{pgfscope}%
\pgfpathrectangle{\pgfqpoint{0.100000in}{0.100000in}}{\pgfqpoint{3.007045in}{1.925000in}}%
\pgfusepath{clip}%
\pgfsetbuttcap%
\pgfsetmiterjoin%
\definecolor{currentfill}{rgb}{0.681738,0.818562,0.904098}%
\pgfsetfillcolor{currentfill}%
\pgfsetlinewidth{0.000000pt}%
\definecolor{currentstroke}{rgb}{0.000000,0.000000,0.000000}%
\pgfsetstrokecolor{currentstroke}%
\pgfsetstrokeopacity{0.000000}%
\pgfsetdash{}{0pt}%
\pgfpathmoveto{\pgfqpoint{0.521615in}{1.353758in}}%
\pgfpathlineto{\pgfqpoint{0.516565in}{1.336204in}}%
\pgfpathlineto{\pgfqpoint{0.492015in}{1.343273in}}%
\pgfpathlineto{\pgfqpoint{0.489895in}{1.349175in}}%
\pgfpathlineto{\pgfqpoint{0.486181in}{1.350484in}}%
\pgfpathlineto{\pgfqpoint{0.480726in}{1.347606in}}%
\pgfpathlineto{\pgfqpoint{0.465761in}{1.347164in}}%
\pgfpathlineto{\pgfqpoint{0.467458in}{1.355133in}}%
\pgfpathlineto{\pgfqpoint{0.457960in}{1.381368in}}%
\pgfpathlineto{\pgfqpoint{0.455804in}{1.383109in}}%
\pgfpathlineto{\pgfqpoint{0.458164in}{1.392328in}}%
\pgfpathlineto{\pgfqpoint{0.461995in}{1.397568in}}%
\pgfpathlineto{\pgfqpoint{0.459905in}{1.401168in}}%
\pgfpathlineto{\pgfqpoint{0.464449in}{1.404081in}}%
\pgfpathlineto{\pgfqpoint{0.466986in}{1.410310in}}%
\pgfpathlineto{\pgfqpoint{0.462117in}{1.414113in}}%
\pgfpathlineto{\pgfqpoint{0.462527in}{1.420581in}}%
\pgfpathlineto{\pgfqpoint{0.470680in}{1.418127in}}%
\pgfpathlineto{\pgfqpoint{0.483507in}{1.414404in}}%
\pgfpathlineto{\pgfqpoint{0.479961in}{1.402391in}}%
\pgfpathlineto{\pgfqpoint{0.482087in}{1.399330in}}%
\pgfpathlineto{\pgfqpoint{0.488838in}{1.398768in}}%
\pgfpathlineto{\pgfqpoint{0.495392in}{1.401958in}}%
\pgfpathlineto{\pgfqpoint{0.503882in}{1.397397in}}%
\pgfpathlineto{\pgfqpoint{0.512166in}{1.383346in}}%
\pgfpathlineto{\pgfqpoint{0.518020in}{1.379672in}}%
\pgfpathlineto{\pgfqpoint{0.517118in}{1.374961in}}%
\pgfpathlineto{\pgfqpoint{0.520278in}{1.368998in}}%
\pgfpathlineto{\pgfqpoint{0.517066in}{1.358013in}}%
\pgfpathlineto{\pgfqpoint{0.521615in}{1.353758in}}%
\pgfpathclose%
\pgfusepath{fill}%
\end{pgfscope}%
\begin{pgfscope}%
\pgfpathrectangle{\pgfqpoint{0.100000in}{0.100000in}}{\pgfqpoint{3.007045in}{1.925000in}}%
\pgfusepath{clip}%
\pgfsetbuttcap%
\pgfsetmiterjoin%
\definecolor{currentfill}{rgb}{0.290980,0.594510,0.789020}%
\pgfsetfillcolor{currentfill}%
\pgfsetlinewidth{0.000000pt}%
\definecolor{currentstroke}{rgb}{0.000000,0.000000,0.000000}%
\pgfsetstrokecolor{currentstroke}%
\pgfsetstrokeopacity{0.000000}%
\pgfsetdash{}{0pt}%
\pgfpathmoveto{\pgfqpoint{1.590530in}{1.433489in}}%
\pgfpathlineto{\pgfqpoint{1.591111in}{1.462290in}}%
\pgfpathlineto{\pgfqpoint{1.613769in}{1.461576in}}%
\pgfpathlineto{\pgfqpoint{1.614383in}{1.478679in}}%
\pgfpathlineto{\pgfqpoint{1.637041in}{1.477807in}}%
\pgfpathlineto{\pgfqpoint{1.637239in}{1.477803in}}%
\pgfpathlineto{\pgfqpoint{1.636440in}{1.454988in}}%
\pgfpathlineto{\pgfqpoint{1.648008in}{1.454576in}}%
\pgfpathlineto{\pgfqpoint{1.647409in}{1.431503in}}%
\pgfpathlineto{\pgfqpoint{1.590530in}{1.433489in}}%
\pgfpathclose%
\pgfusepath{fill}%
\end{pgfscope}%
\begin{pgfscope}%
\pgfpathrectangle{\pgfqpoint{0.100000in}{0.100000in}}{\pgfqpoint{3.007045in}{1.925000in}}%
\pgfusepath{clip}%
\pgfsetbuttcap%
\pgfsetmiterjoin%
\definecolor{currentfill}{rgb}{0.472941,0.711634,0.850719}%
\pgfsetfillcolor{currentfill}%
\pgfsetlinewidth{0.000000pt}%
\definecolor{currentstroke}{rgb}{0.000000,0.000000,0.000000}%
\pgfsetstrokecolor{currentstroke}%
\pgfsetstrokeopacity{0.000000}%
\pgfsetdash{}{0pt}%
\pgfpathmoveto{\pgfqpoint{0.797625in}{1.916024in}}%
\pgfpathlineto{\pgfqpoint{0.789229in}{1.906864in}}%
\pgfpathlineto{\pgfqpoint{0.787685in}{1.900755in}}%
\pgfpathlineto{\pgfqpoint{0.783449in}{1.895856in}}%
\pgfpathlineto{\pgfqpoint{0.789000in}{1.894475in}}%
\pgfpathlineto{\pgfqpoint{0.780353in}{1.860745in}}%
\pgfpathlineto{\pgfqpoint{0.777516in}{1.855475in}}%
\pgfpathlineto{\pgfqpoint{0.773351in}{1.839140in}}%
\pgfpathlineto{\pgfqpoint{0.768647in}{1.844176in}}%
\pgfpathlineto{\pgfqpoint{0.762060in}{1.843258in}}%
\pgfpathlineto{\pgfqpoint{0.759267in}{1.846349in}}%
\pgfpathlineto{\pgfqpoint{0.749934in}{1.845470in}}%
\pgfpathlineto{\pgfqpoint{0.748542in}{1.851631in}}%
\pgfpathlineto{\pgfqpoint{0.744909in}{1.856652in}}%
\pgfpathlineto{\pgfqpoint{0.735184in}{1.852084in}}%
\pgfpathlineto{\pgfqpoint{0.730840in}{1.859613in}}%
\pgfpathlineto{\pgfqpoint{0.719784in}{1.864006in}}%
\pgfpathlineto{\pgfqpoint{0.713636in}{1.864996in}}%
\pgfpathlineto{\pgfqpoint{0.713573in}{1.866288in}}%
\pgfpathlineto{\pgfqpoint{0.713541in}{1.878149in}}%
\pgfpathlineto{\pgfqpoint{0.710511in}{1.880181in}}%
\pgfpathlineto{\pgfqpoint{0.699664in}{1.876125in}}%
\pgfpathlineto{\pgfqpoint{0.695862in}{1.879153in}}%
\pgfpathlineto{\pgfqpoint{0.685273in}{1.876771in}}%
\pgfpathlineto{\pgfqpoint{0.683348in}{1.884220in}}%
\pgfpathlineto{\pgfqpoint{0.676269in}{1.883139in}}%
\pgfpathlineto{\pgfqpoint{0.675310in}{1.876711in}}%
\pgfpathlineto{\pgfqpoint{0.667442in}{1.879016in}}%
\pgfpathlineto{\pgfqpoint{0.666800in}{1.883212in}}%
\pgfpathlineto{\pgfqpoint{0.657890in}{1.895448in}}%
\pgfpathlineto{\pgfqpoint{0.659239in}{1.897832in}}%
\pgfpathlineto{\pgfqpoint{0.654234in}{1.907098in}}%
\pgfpathlineto{\pgfqpoint{0.650154in}{1.917880in}}%
\pgfpathlineto{\pgfqpoint{0.650253in}{1.923532in}}%
\pgfpathlineto{\pgfqpoint{0.648297in}{1.930642in}}%
\pgfpathlineto{\pgfqpoint{0.650462in}{1.932235in}}%
\pgfpathlineto{\pgfqpoint{0.654305in}{1.933747in}}%
\pgfpathlineto{\pgfqpoint{0.652693in}{1.955469in}}%
\pgfpathlineto{\pgfqpoint{0.684552in}{1.946256in}}%
\pgfpathlineto{\pgfqpoint{0.723267in}{1.935529in}}%
\pgfpathlineto{\pgfqpoint{0.752077in}{1.927793in}}%
\pgfpathlineto{\pgfqpoint{0.797625in}{1.916024in}}%
\pgfpathclose%
\pgfusepath{fill}%
\end{pgfscope}%
\begin{pgfscope}%
\pgfpathrectangle{\pgfqpoint{0.100000in}{0.100000in}}{\pgfqpoint{3.007045in}{1.925000in}}%
\pgfusepath{clip}%
\pgfsetbuttcap%
\pgfsetmiterjoin%
\definecolor{currentfill}{rgb}{0.396909,0.666851,0.830358}%
\pgfsetfillcolor{currentfill}%
\pgfsetlinewidth{0.000000pt}%
\definecolor{currentstroke}{rgb}{0.000000,0.000000,0.000000}%
\pgfsetstrokecolor{currentstroke}%
\pgfsetstrokeopacity{0.000000}%
\pgfsetdash{}{0pt}%
\pgfpathmoveto{\pgfqpoint{2.367977in}{0.755519in}}%
\pgfpathlineto{\pgfqpoint{2.361889in}{0.757434in}}%
\pgfpathlineto{\pgfqpoint{2.354579in}{0.763090in}}%
\pgfpathlineto{\pgfqpoint{2.352032in}{0.771233in}}%
\pgfpathlineto{\pgfqpoint{2.352426in}{0.785269in}}%
\pgfpathlineto{\pgfqpoint{2.352410in}{0.787703in}}%
\pgfpathlineto{\pgfqpoint{2.355976in}{0.788053in}}%
\pgfpathlineto{\pgfqpoint{2.357695in}{0.794657in}}%
\pgfpathlineto{\pgfqpoint{2.365423in}{0.795571in}}%
\pgfpathlineto{\pgfqpoint{2.375369in}{0.793380in}}%
\pgfpathlineto{\pgfqpoint{2.382395in}{0.804196in}}%
\pgfpathlineto{\pgfqpoint{2.386212in}{0.799739in}}%
\pgfpathlineto{\pgfqpoint{2.390080in}{0.791606in}}%
\pgfpathlineto{\pgfqpoint{2.390153in}{0.787078in}}%
\pgfpathlineto{\pgfqpoint{2.387448in}{0.788188in}}%
\pgfpathlineto{\pgfqpoint{2.377909in}{0.787069in}}%
\pgfpathlineto{\pgfqpoint{2.379793in}{0.770306in}}%
\pgfpathlineto{\pgfqpoint{2.375546in}{0.768744in}}%
\pgfpathlineto{\pgfqpoint{2.375145in}{0.759957in}}%
\pgfpathlineto{\pgfqpoint{2.373071in}{0.752288in}}%
\pgfpathlineto{\pgfqpoint{2.367977in}{0.755519in}}%
\pgfpathclose%
\pgfusepath{fill}%
\end{pgfscope}%
\begin{pgfscope}%
\pgfpathrectangle{\pgfqpoint{0.100000in}{0.100000in}}{\pgfqpoint{3.007045in}{1.925000in}}%
\pgfusepath{clip}%
\pgfsetbuttcap%
\pgfsetmiterjoin%
\definecolor{currentfill}{rgb}{0.417086,0.680631,0.838231}%
\pgfsetfillcolor{currentfill}%
\pgfsetlinewidth{0.000000pt}%
\definecolor{currentstroke}{rgb}{0.000000,0.000000,0.000000}%
\pgfsetstrokecolor{currentstroke}%
\pgfsetstrokeopacity{0.000000}%
\pgfsetdash{}{0pt}%
\pgfpathmoveto{\pgfqpoint{2.715119in}{1.210272in}}%
\pgfpathlineto{\pgfqpoint{2.706153in}{1.210930in}}%
\pgfpathlineto{\pgfqpoint{2.704557in}{1.214940in}}%
\pgfpathlineto{\pgfqpoint{2.705509in}{1.221124in}}%
\pgfpathlineto{\pgfqpoint{2.703023in}{1.225995in}}%
\pgfpathlineto{\pgfqpoint{2.686960in}{1.233890in}}%
\pgfpathlineto{\pgfqpoint{2.684505in}{1.237247in}}%
\pgfpathlineto{\pgfqpoint{2.677394in}{1.240185in}}%
\pgfpathlineto{\pgfqpoint{2.673952in}{1.243959in}}%
\pgfpathlineto{\pgfqpoint{2.675901in}{1.254838in}}%
\pgfpathlineto{\pgfqpoint{2.672090in}{1.258055in}}%
\pgfpathlineto{\pgfqpoint{2.664785in}{1.260556in}}%
\pgfpathlineto{\pgfqpoint{2.667563in}{1.268330in}}%
\pgfpathlineto{\pgfqpoint{2.716037in}{1.278304in}}%
\pgfpathlineto{\pgfqpoint{2.716371in}{1.278367in}}%
\pgfpathlineto{\pgfqpoint{2.726638in}{1.268556in}}%
\pgfpathlineto{\pgfqpoint{2.726766in}{1.261212in}}%
\pgfpathlineto{\pgfqpoint{2.723696in}{1.256764in}}%
\pgfpathlineto{\pgfqpoint{2.715251in}{1.250311in}}%
\pgfpathlineto{\pgfqpoint{2.713010in}{1.242278in}}%
\pgfpathlineto{\pgfqpoint{2.715503in}{1.232898in}}%
\pgfpathlineto{\pgfqpoint{2.715994in}{1.223495in}}%
\pgfpathlineto{\pgfqpoint{2.713299in}{1.211895in}}%
\pgfpathlineto{\pgfqpoint{2.715119in}{1.210272in}}%
\pgfpathclose%
\pgfusepath{fill}%
\end{pgfscope}%
\begin{pgfscope}%
\pgfpathrectangle{\pgfqpoint{0.100000in}{0.100000in}}{\pgfqpoint{3.007045in}{1.925000in}}%
\pgfusepath{clip}%
\pgfsetbuttcap%
\pgfsetmiterjoin%
\definecolor{currentfill}{rgb}{0.301069,0.601399,0.792957}%
\pgfsetfillcolor{currentfill}%
\pgfsetlinewidth{0.000000pt}%
\definecolor{currentstroke}{rgb}{0.000000,0.000000,0.000000}%
\pgfsetstrokecolor{currentstroke}%
\pgfsetstrokeopacity{0.000000}%
\pgfsetdash{}{0pt}%
\pgfpathmoveto{\pgfqpoint{1.534859in}{1.055660in}}%
\pgfpathlineto{\pgfqpoint{1.534600in}{1.061649in}}%
\pgfpathlineto{\pgfqpoint{1.536034in}{1.090382in}}%
\pgfpathlineto{\pgfqpoint{1.535329in}{1.090418in}}%
\pgfpathlineto{\pgfqpoint{1.536726in}{1.118994in}}%
\pgfpathlineto{\pgfqpoint{1.593165in}{1.116517in}}%
\pgfpathlineto{\pgfqpoint{1.591676in}{1.059068in}}%
\pgfpathlineto{\pgfqpoint{1.563552in}{1.060256in}}%
\pgfpathlineto{\pgfqpoint{1.563272in}{1.054193in}}%
\pgfpathlineto{\pgfqpoint{1.534859in}{1.055660in}}%
\pgfpathclose%
\pgfusepath{fill}%
\end{pgfscope}%
\begin{pgfscope}%
\pgfpathrectangle{\pgfqpoint{0.100000in}{0.100000in}}{\pgfqpoint{3.007045in}{1.925000in}}%
\pgfusepath{clip}%
\pgfsetbuttcap%
\pgfsetmiterjoin%
\definecolor{currentfill}{rgb}{0.711265,0.831111,0.914433}%
\pgfsetfillcolor{currentfill}%
\pgfsetlinewidth{0.000000pt}%
\definecolor{currentstroke}{rgb}{0.000000,0.000000,0.000000}%
\pgfsetstrokecolor{currentstroke}%
\pgfsetstrokeopacity{0.000000}%
\pgfsetdash{}{0pt}%
\pgfpathmoveto{\pgfqpoint{2.204865in}{1.518579in}}%
\pgfpathlineto{\pgfqpoint{2.185259in}{1.517003in}}%
\pgfpathlineto{\pgfqpoint{2.183604in}{1.529014in}}%
\pgfpathlineto{\pgfqpoint{2.190595in}{1.532298in}}%
\pgfpathlineto{\pgfqpoint{2.190444in}{1.542079in}}%
\pgfpathlineto{\pgfqpoint{2.194250in}{1.543292in}}%
\pgfpathlineto{\pgfqpoint{2.196243in}{1.547670in}}%
\pgfpathlineto{\pgfqpoint{2.202517in}{1.547510in}}%
\pgfpathlineto{\pgfqpoint{2.204610in}{1.554457in}}%
\pgfpathlineto{\pgfqpoint{2.210201in}{1.559918in}}%
\pgfpathlineto{\pgfqpoint{2.211986in}{1.556493in}}%
\pgfpathlineto{\pgfqpoint{2.212190in}{1.546743in}}%
\pgfpathlineto{\pgfqpoint{2.210346in}{1.543780in}}%
\pgfpathlineto{\pgfqpoint{2.211388in}{1.536695in}}%
\pgfpathlineto{\pgfqpoint{2.216884in}{1.535425in}}%
\pgfpathlineto{\pgfqpoint{2.221340in}{1.546582in}}%
\pgfpathlineto{\pgfqpoint{2.222009in}{1.560126in}}%
\pgfpathlineto{\pgfqpoint{2.220169in}{1.566322in}}%
\pgfpathlineto{\pgfqpoint{2.227849in}{1.567100in}}%
\pgfpathlineto{\pgfqpoint{2.228428in}{1.561303in}}%
\pgfpathlineto{\pgfqpoint{2.245496in}{1.563087in}}%
\pgfpathlineto{\pgfqpoint{2.247824in}{1.546163in}}%
\pgfpathlineto{\pgfqpoint{2.250233in}{1.523286in}}%
\pgfpathlineto{\pgfqpoint{2.204865in}{1.518579in}}%
\pgfpathclose%
\pgfusepath{fill}%
\end{pgfscope}%
\begin{pgfscope}%
\pgfpathrectangle{\pgfqpoint{0.100000in}{0.100000in}}{\pgfqpoint{3.007045in}{1.925000in}}%
\pgfusepath{clip}%
\pgfsetbuttcap%
\pgfsetmiterjoin%
\definecolor{currentfill}{rgb}{0.316201,0.611734,0.798862}%
\pgfsetfillcolor{currentfill}%
\pgfsetlinewidth{0.000000pt}%
\definecolor{currentstroke}{rgb}{0.000000,0.000000,0.000000}%
\pgfsetstrokecolor{currentstroke}%
\pgfsetstrokeopacity{0.000000}%
\pgfsetdash{}{0pt}%
\pgfpathmoveto{\pgfqpoint{1.541684in}{1.689945in}}%
\pgfpathlineto{\pgfqpoint{1.540344in}{1.666994in}}%
\pgfpathlineto{\pgfqpoint{1.542129in}{1.666889in}}%
\pgfpathlineto{\pgfqpoint{1.540553in}{1.643789in}}%
\pgfpathlineto{\pgfqpoint{1.514103in}{1.645425in}}%
\pgfpathlineto{\pgfqpoint{1.518209in}{1.642479in}}%
\pgfpathlineto{\pgfqpoint{1.519039in}{1.638417in}}%
\pgfpathlineto{\pgfqpoint{1.516525in}{1.631523in}}%
\pgfpathlineto{\pgfqpoint{1.506325in}{1.629275in}}%
\pgfpathlineto{\pgfqpoint{1.503158in}{1.631150in}}%
\pgfpathlineto{\pgfqpoint{1.495994in}{1.623275in}}%
\pgfpathlineto{\pgfqpoint{1.495332in}{1.629121in}}%
\pgfpathlineto{\pgfqpoint{1.483880in}{1.629904in}}%
\pgfpathlineto{\pgfqpoint{1.485073in}{1.647065in}}%
\pgfpathlineto{\pgfqpoint{1.465987in}{1.648477in}}%
\pgfpathlineto{\pgfqpoint{1.466420in}{1.654245in}}%
\pgfpathlineto{\pgfqpoint{1.449237in}{1.655533in}}%
\pgfpathlineto{\pgfqpoint{1.450582in}{1.673001in}}%
\pgfpathlineto{\pgfqpoint{1.448415in}{1.673164in}}%
\pgfpathlineto{\pgfqpoint{1.450238in}{1.696191in}}%
\pgfpathlineto{\pgfqpoint{1.447485in}{1.696400in}}%
\pgfpathlineto{\pgfqpoint{1.448796in}{1.712766in}}%
\pgfpathlineto{\pgfqpoint{1.446373in}{1.716181in}}%
\pgfpathlineto{\pgfqpoint{1.447405in}{1.726710in}}%
\pgfpathlineto{\pgfqpoint{1.441706in}{1.725605in}}%
\pgfpathlineto{\pgfqpoint{1.442195in}{1.731563in}}%
\pgfpathlineto{\pgfqpoint{1.465127in}{1.729765in}}%
\pgfpathlineto{\pgfqpoint{1.522557in}{1.725806in}}%
\pgfpathlineto{\pgfqpoint{1.521849in}{1.714257in}}%
\pgfpathlineto{\pgfqpoint{1.517941in}{1.714507in}}%
\pgfpathlineto{\pgfqpoint{1.516455in}{1.691466in}}%
\pgfpathlineto{\pgfqpoint{1.541684in}{1.689945in}}%
\pgfpathclose%
\pgfusepath{fill}%
\end{pgfscope}%
\begin{pgfscope}%
\pgfpathrectangle{\pgfqpoint{0.100000in}{0.100000in}}{\pgfqpoint{3.007045in}{1.925000in}}%
\pgfusepath{clip}%
\pgfsetbuttcap%
\pgfsetmiterjoin%
\definecolor{currentfill}{rgb}{0.479216,0.715079,0.852072}%
\pgfsetfillcolor{currentfill}%
\pgfsetlinewidth{0.000000pt}%
\definecolor{currentstroke}{rgb}{0.000000,0.000000,0.000000}%
\pgfsetstrokecolor{currentstroke}%
\pgfsetstrokeopacity{0.000000}%
\pgfsetdash{}{0pt}%
\pgfpathmoveto{\pgfqpoint{2.754436in}{1.271463in}}%
\pgfpathlineto{\pgfqpoint{2.754264in}{1.276110in}}%
\pgfpathlineto{\pgfqpoint{2.751124in}{1.279964in}}%
\pgfpathlineto{\pgfqpoint{2.753954in}{1.289095in}}%
\pgfpathlineto{\pgfqpoint{2.753598in}{1.293002in}}%
\pgfpathlineto{\pgfqpoint{2.746343in}{1.292558in}}%
\pgfpathlineto{\pgfqpoint{2.742511in}{1.290165in}}%
\pgfpathlineto{\pgfqpoint{2.739138in}{1.283230in}}%
\pgfpathlineto{\pgfqpoint{2.721216in}{1.279405in}}%
\pgfpathlineto{\pgfqpoint{2.724100in}{1.282597in}}%
\pgfpathlineto{\pgfqpoint{2.725778in}{1.289195in}}%
\pgfpathlineto{\pgfqpoint{2.725049in}{1.295828in}}%
\pgfpathlineto{\pgfqpoint{2.726798in}{1.302007in}}%
\pgfpathlineto{\pgfqpoint{2.725428in}{1.306338in}}%
\pgfpathlineto{\pgfqpoint{2.735728in}{1.317531in}}%
\pgfpathlineto{\pgfqpoint{2.741068in}{1.332531in}}%
\pgfpathlineto{\pgfqpoint{2.746311in}{1.336376in}}%
\pgfpathlineto{\pgfqpoint{2.755066in}{1.346497in}}%
\pgfpathlineto{\pgfqpoint{2.762277in}{1.343620in}}%
\pgfpathlineto{\pgfqpoint{2.764535in}{1.335742in}}%
\pgfpathlineto{\pgfqpoint{2.769299in}{1.335490in}}%
\pgfpathlineto{\pgfqpoint{2.777561in}{1.327350in}}%
\pgfpathlineto{\pgfqpoint{2.791685in}{1.322963in}}%
\pgfpathlineto{\pgfqpoint{2.806809in}{1.301783in}}%
\pgfpathlineto{\pgfqpoint{2.811323in}{1.284246in}}%
\pgfpathlineto{\pgfqpoint{2.815570in}{1.283815in}}%
\pgfpathlineto{\pgfqpoint{2.811785in}{1.275012in}}%
\pgfpathlineto{\pgfqpoint{2.803631in}{1.265614in}}%
\pgfpathlineto{\pgfqpoint{2.798237in}{1.247006in}}%
\pgfpathlineto{\pgfqpoint{2.795440in}{1.242918in}}%
\pgfpathlineto{\pgfqpoint{2.790081in}{1.241725in}}%
\pgfpathlineto{\pgfqpoint{2.790953in}{1.257015in}}%
\pgfpathlineto{\pgfqpoint{2.775479in}{1.258863in}}%
\pgfpathlineto{\pgfqpoint{2.767772in}{1.261597in}}%
\pgfpathlineto{\pgfqpoint{2.757096in}{1.268170in}}%
\pgfpathlineto{\pgfqpoint{2.754436in}{1.271463in}}%
\pgfpathclose%
\pgfusepath{fill}%
\end{pgfscope}%
\begin{pgfscope}%
\pgfpathrectangle{\pgfqpoint{0.100000in}{0.100000in}}{\pgfqpoint{3.007045in}{1.925000in}}%
\pgfusepath{clip}%
\pgfsetbuttcap%
\pgfsetmiterjoin%
\definecolor{currentfill}{rgb}{0.056363,0.349635,0.636755}%
\pgfsetfillcolor{currentfill}%
\pgfsetlinewidth{0.000000pt}%
\definecolor{currentstroke}{rgb}{0.000000,0.000000,0.000000}%
\pgfsetstrokecolor{currentstroke}%
\pgfsetstrokeopacity{0.000000}%
\pgfsetdash{}{0pt}%
\pgfpathmoveto{\pgfqpoint{1.454880in}{1.037434in}}%
\pgfpathlineto{\pgfqpoint{1.453227in}{1.011793in}}%
\pgfpathlineto{\pgfqpoint{1.427563in}{1.013412in}}%
\pgfpathlineto{\pgfqpoint{1.429358in}{1.039255in}}%
\pgfpathlineto{\pgfqpoint{1.404039in}{1.041199in}}%
\pgfpathlineto{\pgfqpoint{1.405343in}{1.057957in}}%
\pgfpathlineto{\pgfqpoint{1.405795in}{1.064151in}}%
\pgfpathlineto{\pgfqpoint{1.455281in}{1.060433in}}%
\pgfpathlineto{\pgfqpoint{1.453744in}{1.037504in}}%
\pgfpathlineto{\pgfqpoint{1.454880in}{1.037434in}}%
\pgfpathclose%
\pgfusepath{fill}%
\end{pgfscope}%
\begin{pgfscope}%
\pgfpathrectangle{\pgfqpoint{0.100000in}{0.100000in}}{\pgfqpoint{3.007045in}{1.925000in}}%
\pgfusepath{clip}%
\pgfsetbuttcap%
\pgfsetmiterjoin%
\definecolor{currentfill}{rgb}{0.270804,0.580730,0.781146}%
\pgfsetfillcolor{currentfill}%
\pgfsetlinewidth{0.000000pt}%
\definecolor{currentstroke}{rgb}{0.000000,0.000000,0.000000}%
\pgfsetstrokecolor{currentstroke}%
\pgfsetstrokeopacity{0.000000}%
\pgfsetdash{}{0pt}%
\pgfpathmoveto{\pgfqpoint{2.397002in}{0.647345in}}%
\pgfpathlineto{\pgfqpoint{2.397277in}{0.644995in}}%
\pgfpathlineto{\pgfqpoint{2.412441in}{0.646597in}}%
\pgfpathlineto{\pgfqpoint{2.415134in}{0.621502in}}%
\pgfpathlineto{\pgfqpoint{2.378692in}{0.619214in}}%
\pgfpathlineto{\pgfqpoint{2.375927in}{0.644931in}}%
\pgfpathlineto{\pgfqpoint{2.397002in}{0.647345in}}%
\pgfpathclose%
\pgfusepath{fill}%
\end{pgfscope}%
\begin{pgfscope}%
\pgfpathrectangle{\pgfqpoint{0.100000in}{0.100000in}}{\pgfqpoint{3.007045in}{1.925000in}}%
\pgfusepath{clip}%
\pgfsetbuttcap%
\pgfsetmiterjoin%
\definecolor{currentfill}{rgb}{0.523137,0.739193,0.861546}%
\pgfsetfillcolor{currentfill}%
\pgfsetlinewidth{0.000000pt}%
\definecolor{currentstroke}{rgb}{0.000000,0.000000,0.000000}%
\pgfsetstrokecolor{currentstroke}%
\pgfsetstrokeopacity{0.000000}%
\pgfsetdash{}{0pt}%
\pgfpathmoveto{\pgfqpoint{2.116398in}{1.073683in}}%
\pgfpathlineto{\pgfqpoint{2.099179in}{1.072352in}}%
\pgfpathlineto{\pgfqpoint{2.098280in}{1.086724in}}%
\pgfpathlineto{\pgfqpoint{2.075282in}{1.085153in}}%
\pgfpathlineto{\pgfqpoint{2.075029in}{1.090942in}}%
\pgfpathlineto{\pgfqpoint{2.046290in}{1.089601in}}%
\pgfpathlineto{\pgfqpoint{2.044424in}{1.118250in}}%
\pgfpathlineto{\pgfqpoint{2.050157in}{1.118638in}}%
\pgfpathlineto{\pgfqpoint{2.049708in}{1.124375in}}%
\pgfpathlineto{\pgfqpoint{2.067309in}{1.125374in}}%
\pgfpathlineto{\pgfqpoint{2.073272in}{1.125381in}}%
\pgfpathlineto{\pgfqpoint{2.072910in}{1.131159in}}%
\pgfpathlineto{\pgfqpoint{2.095572in}{1.132813in}}%
\pgfpathlineto{\pgfqpoint{2.096418in}{1.118342in}}%
\pgfpathlineto{\pgfqpoint{2.097811in}{1.095356in}}%
\pgfpathlineto{\pgfqpoint{2.114940in}{1.096476in}}%
\pgfpathlineto{\pgfqpoint{2.116398in}{1.073683in}}%
\pgfpathclose%
\pgfusepath{fill}%
\end{pgfscope}%
\begin{pgfscope}%
\pgfpathrectangle{\pgfqpoint{0.100000in}{0.100000in}}{\pgfqpoint{3.007045in}{1.925000in}}%
\pgfusepath{clip}%
\pgfsetbuttcap%
\pgfsetmiterjoin%
\definecolor{currentfill}{rgb}{0.331334,0.622068,0.804767}%
\pgfsetfillcolor{currentfill}%
\pgfsetlinewidth{0.000000pt}%
\definecolor{currentstroke}{rgb}{0.000000,0.000000,0.000000}%
\pgfsetstrokecolor{currentstroke}%
\pgfsetstrokeopacity{0.000000}%
\pgfsetdash{}{0pt}%
\pgfpathmoveto{\pgfqpoint{2.241536in}{0.660624in}}%
\pgfpathlineto{\pgfqpoint{2.233755in}{0.660011in}}%
\pgfpathlineto{\pgfqpoint{2.229320in}{0.666780in}}%
\pgfpathlineto{\pgfqpoint{2.228144in}{0.679871in}}%
\pgfpathlineto{\pgfqpoint{2.226858in}{0.694238in}}%
\pgfpathlineto{\pgfqpoint{2.225729in}{0.705930in}}%
\pgfpathlineto{\pgfqpoint{2.230267in}{0.714046in}}%
\pgfpathlineto{\pgfqpoint{2.223953in}{0.726994in}}%
\pgfpathlineto{\pgfqpoint{2.222669in}{0.734886in}}%
\pgfpathlineto{\pgfqpoint{2.233994in}{0.735748in}}%
\pgfpathlineto{\pgfqpoint{2.233682in}{0.738645in}}%
\pgfpathlineto{\pgfqpoint{2.250388in}{0.740313in}}%
\pgfpathlineto{\pgfqpoint{2.252087in}{0.743598in}}%
\pgfpathlineto{\pgfqpoint{2.279667in}{0.746313in}}%
\pgfpathlineto{\pgfqpoint{2.281053in}{0.729035in}}%
\pgfpathlineto{\pgfqpoint{2.274136in}{0.722813in}}%
\pgfpathlineto{\pgfqpoint{2.276404in}{0.718014in}}%
\pgfpathlineto{\pgfqpoint{2.277794in}{0.699248in}}%
\pgfpathlineto{\pgfqpoint{2.278425in}{0.693743in}}%
\pgfpathlineto{\pgfqpoint{2.261248in}{0.691899in}}%
\pgfpathlineto{\pgfqpoint{2.260621in}{0.697611in}}%
\pgfpathlineto{\pgfqpoint{2.254880in}{0.697001in}}%
\pgfpathlineto{\pgfqpoint{2.253069in}{0.691091in}}%
\pgfpathlineto{\pgfqpoint{2.254972in}{0.670665in}}%
\pgfpathlineto{\pgfqpoint{2.252110in}{0.670372in}}%
\pgfpathlineto{\pgfqpoint{2.252934in}{0.661782in}}%
\pgfpathlineto{\pgfqpoint{2.241536in}{0.660624in}}%
\pgfpathclose%
\pgfusepath{fill}%
\end{pgfscope}%
\begin{pgfscope}%
\pgfpathrectangle{\pgfqpoint{0.100000in}{0.100000in}}{\pgfqpoint{3.007045in}{1.925000in}}%
\pgfusepath{clip}%
\pgfsetbuttcap%
\pgfsetmiterjoin%
\definecolor{currentfill}{rgb}{0.529412,0.742637,0.862899}%
\pgfsetfillcolor{currentfill}%
\pgfsetlinewidth{0.000000pt}%
\definecolor{currentstroke}{rgb}{0.000000,0.000000,0.000000}%
\pgfsetstrokecolor{currentstroke}%
\pgfsetstrokeopacity{0.000000}%
\pgfsetdash{}{0pt}%
\pgfpathmoveto{\pgfqpoint{0.565403in}{1.940383in}}%
\pgfpathlineto{\pgfqpoint{0.569196in}{1.933206in}}%
\pgfpathlineto{\pgfqpoint{0.564241in}{1.932920in}}%
\pgfpathlineto{\pgfqpoint{0.561680in}{1.929053in}}%
\pgfpathlineto{\pgfqpoint{0.565762in}{1.918633in}}%
\pgfpathlineto{\pgfqpoint{0.569487in}{1.913159in}}%
\pgfpathlineto{\pgfqpoint{0.569117in}{1.907951in}}%
\pgfpathlineto{\pgfqpoint{0.564737in}{1.906268in}}%
\pgfpathlineto{\pgfqpoint{0.563966in}{1.911840in}}%
\pgfpathlineto{\pgfqpoint{0.559562in}{1.916054in}}%
\pgfpathlineto{\pgfqpoint{0.561416in}{1.920783in}}%
\pgfpathlineto{\pgfqpoint{0.556562in}{1.930703in}}%
\pgfpathlineto{\pgfqpoint{0.565403in}{1.940383in}}%
\pgfpathclose%
\pgfusepath{fill}%
\end{pgfscope}%
\begin{pgfscope}%
\pgfpathrectangle{\pgfqpoint{0.100000in}{0.100000in}}{\pgfqpoint{3.007045in}{1.925000in}}%
\pgfusepath{clip}%
\pgfsetbuttcap%
\pgfsetmiterjoin%
\definecolor{currentfill}{rgb}{0.529412,0.742637,0.862899}%
\pgfsetfillcolor{currentfill}%
\pgfsetlinewidth{0.000000pt}%
\definecolor{currentstroke}{rgb}{0.000000,0.000000,0.000000}%
\pgfsetstrokecolor{currentstroke}%
\pgfsetstrokeopacity{0.000000}%
\pgfsetdash{}{0pt}%
\pgfpathmoveto{\pgfqpoint{0.574289in}{1.797104in}}%
\pgfpathlineto{\pgfqpoint{0.542761in}{1.806508in}}%
\pgfpathlineto{\pgfqpoint{0.493956in}{1.821847in}}%
\pgfpathlineto{\pgfqpoint{0.501743in}{1.847839in}}%
\pgfpathlineto{\pgfqpoint{0.510872in}{1.844976in}}%
\pgfpathlineto{\pgfqpoint{0.515078in}{1.857792in}}%
\pgfpathlineto{\pgfqpoint{0.514994in}{1.864013in}}%
\pgfpathlineto{\pgfqpoint{0.502572in}{1.867867in}}%
\pgfpathlineto{\pgfqpoint{0.505442in}{1.879195in}}%
\pgfpathlineto{\pgfqpoint{0.510848in}{1.895622in}}%
\pgfpathlineto{\pgfqpoint{0.512685in}{1.901288in}}%
\pgfpathlineto{\pgfqpoint{0.536363in}{1.893635in}}%
\pgfpathlineto{\pgfqpoint{0.546025in}{1.896043in}}%
\pgfpathlineto{\pgfqpoint{0.549783in}{1.902072in}}%
\pgfpathlineto{\pgfqpoint{0.556850in}{1.906039in}}%
\pgfpathlineto{\pgfqpoint{0.556363in}{1.910254in}}%
\pgfpathlineto{\pgfqpoint{0.560656in}{1.907168in}}%
\pgfpathlineto{\pgfqpoint{0.559298in}{1.890060in}}%
\pgfpathlineto{\pgfqpoint{0.563615in}{1.893690in}}%
\pgfpathlineto{\pgfqpoint{0.564306in}{1.899159in}}%
\pgfpathlineto{\pgfqpoint{0.570797in}{1.906935in}}%
\pgfpathlineto{\pgfqpoint{0.575526in}{1.910604in}}%
\pgfpathlineto{\pgfqpoint{0.572185in}{1.916693in}}%
\pgfpathlineto{\pgfqpoint{0.573093in}{1.922337in}}%
\pgfpathlineto{\pgfqpoint{0.567141in}{1.928935in}}%
\pgfpathlineto{\pgfqpoint{0.574554in}{1.930262in}}%
\pgfpathlineto{\pgfqpoint{0.568642in}{1.941046in}}%
\pgfpathlineto{\pgfqpoint{0.578403in}{1.950226in}}%
\pgfpathlineto{\pgfqpoint{0.576681in}{1.953612in}}%
\pgfpathlineto{\pgfqpoint{0.623985in}{1.938935in}}%
\pgfpathlineto{\pgfqpoint{0.643432in}{1.933133in}}%
\pgfpathlineto{\pgfqpoint{0.650462in}{1.932235in}}%
\pgfpathlineto{\pgfqpoint{0.648297in}{1.930642in}}%
\pgfpathlineto{\pgfqpoint{0.650253in}{1.923532in}}%
\pgfpathlineto{\pgfqpoint{0.643966in}{1.926557in}}%
\pgfpathlineto{\pgfqpoint{0.633546in}{1.923434in}}%
\pgfpathlineto{\pgfqpoint{0.630805in}{1.915337in}}%
\pgfpathlineto{\pgfqpoint{0.634619in}{1.904375in}}%
\pgfpathlineto{\pgfqpoint{0.631834in}{1.900038in}}%
\pgfpathlineto{\pgfqpoint{0.622765in}{1.897664in}}%
\pgfpathlineto{\pgfqpoint{0.617908in}{1.888997in}}%
\pgfpathlineto{\pgfqpoint{0.619147in}{1.875248in}}%
\pgfpathlineto{\pgfqpoint{0.615087in}{1.868865in}}%
\pgfpathlineto{\pgfqpoint{0.610134in}{1.868414in}}%
\pgfpathlineto{\pgfqpoint{0.600903in}{1.863151in}}%
\pgfpathlineto{\pgfqpoint{0.595348in}{1.857817in}}%
\pgfpathlineto{\pgfqpoint{0.595548in}{1.853584in}}%
\pgfpathlineto{\pgfqpoint{0.599282in}{1.851636in}}%
\pgfpathlineto{\pgfqpoint{0.597616in}{1.841914in}}%
\pgfpathlineto{\pgfqpoint{0.583582in}{1.827974in}}%
\pgfpathlineto{\pgfqpoint{0.585552in}{1.821819in}}%
\pgfpathlineto{\pgfqpoint{0.585013in}{1.811079in}}%
\pgfpathlineto{\pgfqpoint{0.579759in}{1.803398in}}%
\pgfpathlineto{\pgfqpoint{0.579595in}{1.795500in}}%
\pgfpathlineto{\pgfqpoint{0.574289in}{1.797104in}}%
\pgfpathclose%
\pgfusepath{fill}%
\end{pgfscope}%
\begin{pgfscope}%
\pgfpathrectangle{\pgfqpoint{0.100000in}{0.100000in}}{\pgfqpoint{3.007045in}{1.925000in}}%
\pgfusepath{clip}%
\pgfsetbuttcap%
\pgfsetmiterjoin%
\definecolor{currentfill}{rgb}{0.435294,0.690965,0.842599}%
\pgfsetfillcolor{currentfill}%
\pgfsetlinewidth{0.000000pt}%
\definecolor{currentstroke}{rgb}{0.000000,0.000000,0.000000}%
\pgfsetstrokecolor{currentstroke}%
\pgfsetstrokeopacity{0.000000}%
\pgfsetdash{}{0pt}%
\pgfpathmoveto{\pgfqpoint{2.196663in}{1.378800in}}%
\pgfpathlineto{\pgfqpoint{2.232094in}{1.382445in}}%
\pgfpathlineto{\pgfqpoint{2.229609in}{1.405210in}}%
\pgfpathlineto{\pgfqpoint{2.252299in}{1.407809in}}%
\pgfpathlineto{\pgfqpoint{2.254932in}{1.384910in}}%
\pgfpathlineto{\pgfqpoint{2.272022in}{1.386856in}}%
\pgfpathlineto{\pgfqpoint{2.275068in}{1.363856in}}%
\pgfpathlineto{\pgfqpoint{2.269373in}{1.363358in}}%
\pgfpathlineto{\pgfqpoint{2.223806in}{1.358262in}}%
\pgfpathlineto{\pgfqpoint{2.201376in}{1.356174in}}%
\pgfpathlineto{\pgfqpoint{2.200252in}{1.367456in}}%
\pgfpathlineto{\pgfqpoint{2.193342in}{1.366778in}}%
\pgfpathlineto{\pgfqpoint{2.196663in}{1.378800in}}%
\pgfpathclose%
\pgfusepath{fill}%
\end{pgfscope}%
\begin{pgfscope}%
\pgfpathrectangle{\pgfqpoint{0.100000in}{0.100000in}}{\pgfqpoint{3.007045in}{1.925000in}}%
\pgfusepath{clip}%
\pgfsetbuttcap%
\pgfsetmiterjoin%
\definecolor{currentfill}{rgb}{0.460392,0.704744,0.848012}%
\pgfsetfillcolor{currentfill}%
\pgfsetlinewidth{0.000000pt}%
\definecolor{currentstroke}{rgb}{0.000000,0.000000,0.000000}%
\pgfsetstrokecolor{currentstroke}%
\pgfsetstrokeopacity{0.000000}%
\pgfsetdash{}{0pt}%
\pgfpathmoveto{\pgfqpoint{1.102971in}{1.695747in}}%
\pgfpathlineto{\pgfqpoint{1.119016in}{1.693144in}}%
\pgfpathlineto{\pgfqpoint{1.113428in}{1.662137in}}%
\pgfpathlineto{\pgfqpoint{1.102024in}{1.663968in}}%
\pgfpathlineto{\pgfqpoint{1.073596in}{1.668769in}}%
\pgfpathlineto{\pgfqpoint{1.067509in}{1.667877in}}%
\pgfpathlineto{\pgfqpoint{1.037941in}{1.673119in}}%
\pgfpathlineto{\pgfqpoint{1.036887in}{1.681674in}}%
\pgfpathlineto{\pgfqpoint{1.040593in}{1.686492in}}%
\pgfpathlineto{\pgfqpoint{1.037077in}{1.690684in}}%
\pgfpathlineto{\pgfqpoint{1.033247in}{1.689867in}}%
\pgfpathlineto{\pgfqpoint{1.030764in}{1.699834in}}%
\pgfpathlineto{\pgfqpoint{1.027516in}{1.705647in}}%
\pgfpathlineto{\pgfqpoint{1.025028in}{1.717016in}}%
\pgfpathlineto{\pgfqpoint{1.019449in}{1.720433in}}%
\pgfpathlineto{\pgfqpoint{1.019738in}{1.725093in}}%
\pgfpathlineto{\pgfqpoint{1.025046in}{1.724117in}}%
\pgfpathlineto{\pgfqpoint{1.026102in}{1.729804in}}%
\pgfpathlineto{\pgfqpoint{1.043222in}{1.727442in}}%
\pgfpathlineto{\pgfqpoint{1.053818in}{1.723481in}}%
\pgfpathlineto{\pgfqpoint{1.059670in}{1.716732in}}%
\pgfpathlineto{\pgfqpoint{1.066660in}{1.704559in}}%
\pgfpathlineto{\pgfqpoint{1.070809in}{1.700125in}}%
\pgfpathlineto{\pgfqpoint{1.093536in}{1.701272in}}%
\pgfpathlineto{\pgfqpoint{1.099659in}{1.696257in}}%
\pgfpathlineto{\pgfqpoint{1.102971in}{1.695747in}}%
\pgfpathclose%
\pgfusepath{fill}%
\end{pgfscope}%
\begin{pgfscope}%
\pgfpathrectangle{\pgfqpoint{0.100000in}{0.100000in}}{\pgfqpoint{3.007045in}{1.925000in}}%
\pgfusepath{clip}%
\pgfsetbuttcap%
\pgfsetmiterjoin%
\definecolor{currentfill}{rgb}{0.417086,0.680631,0.838231}%
\pgfsetfillcolor{currentfill}%
\pgfsetlinewidth{0.000000pt}%
\definecolor{currentstroke}{rgb}{0.000000,0.000000,0.000000}%
\pgfsetstrokecolor{currentstroke}%
\pgfsetstrokeopacity{0.000000}%
\pgfsetdash{}{0pt}%
\pgfpathmoveto{\pgfqpoint{2.802639in}{1.349500in}}%
\pgfpathlineto{\pgfqpoint{2.804907in}{1.347247in}}%
\pgfpathlineto{\pgfqpoint{2.799417in}{1.334517in}}%
\pgfpathlineto{\pgfqpoint{2.789776in}{1.325534in}}%
\pgfpathlineto{\pgfqpoint{2.791685in}{1.322963in}}%
\pgfpathlineto{\pgfqpoint{2.777561in}{1.327350in}}%
\pgfpathlineto{\pgfqpoint{2.769299in}{1.335490in}}%
\pgfpathlineto{\pgfqpoint{2.764535in}{1.335742in}}%
\pgfpathlineto{\pgfqpoint{2.762277in}{1.343620in}}%
\pgfpathlineto{\pgfqpoint{2.755066in}{1.346497in}}%
\pgfpathlineto{\pgfqpoint{2.753885in}{1.357653in}}%
\pgfpathlineto{\pgfqpoint{2.756767in}{1.359194in}}%
\pgfpathlineto{\pgfqpoint{2.758253in}{1.365099in}}%
\pgfpathlineto{\pgfqpoint{2.753651in}{1.370464in}}%
\pgfpathlineto{\pgfqpoint{2.762147in}{1.386718in}}%
\pgfpathlineto{\pgfqpoint{2.763125in}{1.394125in}}%
\pgfpathlineto{\pgfqpoint{2.768627in}{1.400117in}}%
\pgfpathlineto{\pgfqpoint{2.794033in}{1.391447in}}%
\pgfpathlineto{\pgfqpoint{2.803517in}{1.405975in}}%
\pgfpathlineto{\pgfqpoint{2.805188in}{1.401167in}}%
\pgfpathlineto{\pgfqpoint{2.810869in}{1.395086in}}%
\pgfpathlineto{\pgfqpoint{2.812809in}{1.384940in}}%
\pgfpathlineto{\pgfqpoint{2.810062in}{1.361931in}}%
\pgfpathlineto{\pgfqpoint{2.803807in}{1.360108in}}%
\pgfpathlineto{\pgfqpoint{2.802639in}{1.349500in}}%
\pgfpathclose%
\pgfusepath{fill}%
\end{pgfscope}%
\begin{pgfscope}%
\pgfpathrectangle{\pgfqpoint{0.100000in}{0.100000in}}{\pgfqpoint{3.007045in}{1.925000in}}%
\pgfusepath{clip}%
\pgfsetbuttcap%
\pgfsetmiterjoin%
\definecolor{currentfill}{rgb}{0.396909,0.666851,0.830358}%
\pgfsetfillcolor{currentfill}%
\pgfsetlinewidth{0.000000pt}%
\definecolor{currentstroke}{rgb}{0.000000,0.000000,0.000000}%
\pgfsetstrokecolor{currentstroke}%
\pgfsetstrokeopacity{0.000000}%
\pgfsetdash{}{0pt}%
\pgfpathmoveto{\pgfqpoint{2.427990in}{0.752079in}}%
\pgfpathlineto{\pgfqpoint{2.441826in}{0.762195in}}%
\pgfpathlineto{\pgfqpoint{2.441164in}{0.768340in}}%
\pgfpathlineto{\pgfqpoint{2.436633in}{0.771428in}}%
\pgfpathlineto{\pgfqpoint{2.433317in}{0.776989in}}%
\pgfpathlineto{\pgfqpoint{2.432968in}{0.782406in}}%
\pgfpathlineto{\pgfqpoint{2.437607in}{0.788428in}}%
\pgfpathlineto{\pgfqpoint{2.441806in}{0.790953in}}%
\pgfpathlineto{\pgfqpoint{2.443204in}{0.795028in}}%
\pgfpathlineto{\pgfqpoint{2.448640in}{0.798533in}}%
\pgfpathlineto{\pgfqpoint{2.449695in}{0.794403in}}%
\pgfpathlineto{\pgfqpoint{2.454423in}{0.792007in}}%
\pgfpathlineto{\pgfqpoint{2.461072in}{0.785247in}}%
\pgfpathlineto{\pgfqpoint{2.464577in}{0.779288in}}%
\pgfpathlineto{\pgfqpoint{2.464508in}{0.773648in}}%
\pgfpathlineto{\pgfqpoint{2.470222in}{0.769671in}}%
\pgfpathlineto{\pgfqpoint{2.467299in}{0.764175in}}%
\pgfpathlineto{\pgfqpoint{2.469397in}{0.761860in}}%
\pgfpathlineto{\pgfqpoint{2.463033in}{0.756433in}}%
\pgfpathlineto{\pgfqpoint{2.460691in}{0.751658in}}%
\pgfpathlineto{\pgfqpoint{2.457386in}{0.739684in}}%
\pgfpathlineto{\pgfqpoint{2.458380in}{0.737732in}}%
\pgfpathlineto{\pgfqpoint{2.451822in}{0.729156in}}%
\pgfpathlineto{\pgfqpoint{2.444853in}{0.725072in}}%
\pgfpathlineto{\pgfqpoint{2.434230in}{0.742118in}}%
\pgfpathlineto{\pgfqpoint{2.427990in}{0.752079in}}%
\pgfpathclose%
\pgfusepath{fill}%
\end{pgfscope}%
\begin{pgfscope}%
\pgfpathrectangle{\pgfqpoint{0.100000in}{0.100000in}}{\pgfqpoint{3.007045in}{1.925000in}}%
\pgfusepath{clip}%
\pgfsetbuttcap%
\pgfsetmiterjoin%
\definecolor{currentfill}{rgb}{0.280892,0.587620,0.785083}%
\pgfsetfillcolor{currentfill}%
\pgfsetlinewidth{0.000000pt}%
\definecolor{currentstroke}{rgb}{0.000000,0.000000,0.000000}%
\pgfsetstrokecolor{currentstroke}%
\pgfsetstrokeopacity{0.000000}%
\pgfsetdash{}{0pt}%
\pgfpathmoveto{\pgfqpoint{2.113075in}{0.848655in}}%
\pgfpathlineto{\pgfqpoint{2.110261in}{0.847530in}}%
\pgfpathlineto{\pgfqpoint{2.110699in}{0.840860in}}%
\pgfpathlineto{\pgfqpoint{2.103193in}{0.839402in}}%
\pgfpathlineto{\pgfqpoint{2.085917in}{0.838289in}}%
\pgfpathlineto{\pgfqpoint{2.085442in}{0.845945in}}%
\pgfpathlineto{\pgfqpoint{2.082243in}{0.851523in}}%
\pgfpathlineto{\pgfqpoint{2.081112in}{0.868758in}}%
\pgfpathlineto{\pgfqpoint{2.078219in}{0.868579in}}%
\pgfpathlineto{\pgfqpoint{2.077608in}{0.878428in}}%
\pgfpathlineto{\pgfqpoint{2.085939in}{0.878995in}}%
\pgfpathlineto{\pgfqpoint{2.106178in}{0.880380in}}%
\pgfpathlineto{\pgfqpoint{2.106779in}{0.871379in}}%
\pgfpathlineto{\pgfqpoint{2.110591in}{0.871650in}}%
\pgfpathlineto{\pgfqpoint{2.112919in}{0.865058in}}%
\pgfpathlineto{\pgfqpoint{2.113075in}{0.848655in}}%
\pgfpathclose%
\pgfusepath{fill}%
\end{pgfscope}%
\begin{pgfscope}%
\pgfpathrectangle{\pgfqpoint{0.100000in}{0.100000in}}{\pgfqpoint{3.007045in}{1.925000in}}%
\pgfusepath{clip}%
\pgfsetbuttcap%
\pgfsetmiterjoin%
\definecolor{currentfill}{rgb}{0.460392,0.704744,0.848012}%
\pgfsetfillcolor{currentfill}%
\pgfsetlinewidth{0.000000pt}%
\definecolor{currentstroke}{rgb}{0.000000,0.000000,0.000000}%
\pgfsetstrokecolor{currentstroke}%
\pgfsetstrokeopacity{0.000000}%
\pgfsetdash{}{0pt}%
\pgfpathmoveto{\pgfqpoint{2.179696in}{0.840159in}}%
\pgfpathlineto{\pgfqpoint{2.144570in}{0.838481in}}%
\pgfpathlineto{\pgfqpoint{2.144807in}{0.847920in}}%
\pgfpathlineto{\pgfqpoint{2.145830in}{0.876417in}}%
\pgfpathlineto{\pgfqpoint{2.139583in}{0.883616in}}%
\pgfpathlineto{\pgfqpoint{2.151351in}{0.884373in}}%
\pgfpathlineto{\pgfqpoint{2.150055in}{0.903459in}}%
\pgfpathlineto{\pgfqpoint{2.147248in}{0.909640in}}%
\pgfpathlineto{\pgfqpoint{2.149635in}{0.914318in}}%
\pgfpathlineto{\pgfqpoint{2.163243in}{0.916924in}}%
\pgfpathlineto{\pgfqpoint{2.167138in}{0.916537in}}%
\pgfpathlineto{\pgfqpoint{2.171317in}{0.911958in}}%
\pgfpathlineto{\pgfqpoint{2.184378in}{0.916884in}}%
\pgfpathlineto{\pgfqpoint{2.189759in}{0.915550in}}%
\pgfpathlineto{\pgfqpoint{2.190733in}{0.912451in}}%
\pgfpathlineto{\pgfqpoint{2.192088in}{0.890541in}}%
\pgfpathlineto{\pgfqpoint{2.193121in}{0.887325in}}%
\pgfpathlineto{\pgfqpoint{2.193931in}{0.875265in}}%
\pgfpathlineto{\pgfqpoint{2.189352in}{0.871261in}}%
\pgfpathlineto{\pgfqpoint{2.182587in}{0.873244in}}%
\pgfpathlineto{\pgfqpoint{2.183086in}{0.867245in}}%
\pgfpathlineto{\pgfqpoint{2.178278in}{0.857446in}}%
\pgfpathlineto{\pgfqpoint{2.179696in}{0.840159in}}%
\pgfpathclose%
\pgfusepath{fill}%
\end{pgfscope}%
\begin{pgfscope}%
\pgfpathrectangle{\pgfqpoint{0.100000in}{0.100000in}}{\pgfqpoint{3.007045in}{1.925000in}}%
\pgfusepath{clip}%
\pgfsetbuttcap%
\pgfsetmiterjoin%
\definecolor{currentfill}{rgb}{0.311157,0.608289,0.796894}%
\pgfsetfillcolor{currentfill}%
\pgfsetlinewidth{0.000000pt}%
\definecolor{currentstroke}{rgb}{0.000000,0.000000,0.000000}%
\pgfsetstrokecolor{currentstroke}%
\pgfsetstrokeopacity{0.000000}%
\pgfsetdash{}{0pt}%
\pgfpathmoveto{\pgfqpoint{1.662731in}{0.409312in}}%
\pgfpathlineto{\pgfqpoint{1.664160in}{0.412781in}}%
\pgfpathlineto{\pgfqpoint{1.665098in}{0.416805in}}%
\pgfpathlineto{\pgfqpoint{1.670154in}{0.417855in}}%
\pgfpathlineto{\pgfqpoint{1.678639in}{0.424954in}}%
\pgfpathlineto{\pgfqpoint{1.680357in}{0.422649in}}%
\pgfpathlineto{\pgfqpoint{1.670606in}{0.415983in}}%
\pgfpathlineto{\pgfqpoint{1.662731in}{0.409312in}}%
\pgfpathclose%
\pgfusepath{fill}%
\end{pgfscope}%
\begin{pgfscope}%
\pgfpathrectangle{\pgfqpoint{0.100000in}{0.100000in}}{\pgfqpoint{3.007045in}{1.925000in}}%
\pgfusepath{clip}%
\pgfsetbuttcap%
\pgfsetmiterjoin%
\definecolor{currentfill}{rgb}{0.311157,0.608289,0.796894}%
\pgfsetfillcolor{currentfill}%
\pgfsetlinewidth{0.000000pt}%
\definecolor{currentstroke}{rgb}{0.000000,0.000000,0.000000}%
\pgfsetstrokecolor{currentstroke}%
\pgfsetstrokeopacity{0.000000}%
\pgfsetdash{}{0pt}%
\pgfpathmoveto{\pgfqpoint{1.694633in}{0.447645in}}%
\pgfpathlineto{\pgfqpoint{1.685694in}{0.443169in}}%
\pgfpathlineto{\pgfqpoint{1.680452in}{0.447923in}}%
\pgfpathlineto{\pgfqpoint{1.679297in}{0.453438in}}%
\pgfpathlineto{\pgfqpoint{1.674645in}{0.451462in}}%
\pgfpathlineto{\pgfqpoint{1.678471in}{0.443997in}}%
\pgfpathlineto{\pgfqpoint{1.690125in}{0.434156in}}%
\pgfpathlineto{\pgfqpoint{1.679523in}{0.428823in}}%
\pgfpathlineto{\pgfqpoint{1.676716in}{0.426224in}}%
\pgfpathlineto{\pgfqpoint{1.667536in}{0.431953in}}%
\pgfpathlineto{\pgfqpoint{1.661661in}{0.438978in}}%
\pgfpathlineto{\pgfqpoint{1.655495in}{0.438831in}}%
\pgfpathlineto{\pgfqpoint{1.651860in}{0.442260in}}%
\pgfpathlineto{\pgfqpoint{1.645296in}{0.442654in}}%
\pgfpathlineto{\pgfqpoint{1.632363in}{0.431928in}}%
\pgfpathlineto{\pgfqpoint{1.622101in}{0.436398in}}%
\pgfpathlineto{\pgfqpoint{1.621822in}{0.440642in}}%
\pgfpathlineto{\pgfqpoint{1.613769in}{0.442713in}}%
\pgfpathlineto{\pgfqpoint{1.609200in}{0.451305in}}%
\pgfpathlineto{\pgfqpoint{1.621451in}{0.460671in}}%
\pgfpathlineto{\pgfqpoint{1.611272in}{0.473936in}}%
\pgfpathlineto{\pgfqpoint{1.619767in}{0.480550in}}%
\pgfpathlineto{\pgfqpoint{1.641992in}{0.498317in}}%
\pgfpathlineto{\pgfqpoint{1.644405in}{0.511356in}}%
\pgfpathlineto{\pgfqpoint{1.648009in}{0.514439in}}%
\pgfpathlineto{\pgfqpoint{1.663618in}{0.514403in}}%
\pgfpathlineto{\pgfqpoint{1.668257in}{0.510605in}}%
\pgfpathlineto{\pgfqpoint{1.681469in}{0.494214in}}%
\pgfpathlineto{\pgfqpoint{1.676720in}{0.488443in}}%
\pgfpathlineto{\pgfqpoint{1.695745in}{0.469116in}}%
\pgfpathlineto{\pgfqpoint{1.694633in}{0.447645in}}%
\pgfpathclose%
\pgfusepath{fill}%
\end{pgfscope}%
\begin{pgfscope}%
\pgfpathrectangle{\pgfqpoint{0.100000in}{0.100000in}}{\pgfqpoint{3.007045in}{1.925000in}}%
\pgfusepath{clip}%
\pgfsetbuttcap%
\pgfsetmiterjoin%
\definecolor{currentfill}{rgb}{0.396909,0.666851,0.830358}%
\pgfsetfillcolor{currentfill}%
\pgfsetlinewidth{0.000000pt}%
\definecolor{currentstroke}{rgb}{0.000000,0.000000,0.000000}%
\pgfsetstrokecolor{currentstroke}%
\pgfsetstrokeopacity{0.000000}%
\pgfsetdash{}{0pt}%
\pgfpathmoveto{\pgfqpoint{1.985200in}{0.741858in}}%
\pgfpathlineto{\pgfqpoint{1.970228in}{0.741380in}}%
\pgfpathlineto{\pgfqpoint{1.968862in}{0.741315in}}%
\pgfpathlineto{\pgfqpoint{1.968042in}{0.766663in}}%
\pgfpathlineto{\pgfqpoint{1.938599in}{0.766214in}}%
\pgfpathlineto{\pgfqpoint{1.939283in}{0.759647in}}%
\pgfpathlineto{\pgfqpoint{1.931371in}{0.753155in}}%
\pgfpathlineto{\pgfqpoint{1.930951in}{0.750384in}}%
\pgfpathlineto{\pgfqpoint{1.921911in}{0.757915in}}%
\pgfpathlineto{\pgfqpoint{1.918474in}{0.757862in}}%
\pgfpathlineto{\pgfqpoint{1.917305in}{0.764130in}}%
\pgfpathlineto{\pgfqpoint{1.920677in}{0.772374in}}%
\pgfpathlineto{\pgfqpoint{1.918485in}{0.794539in}}%
\pgfpathlineto{\pgfqpoint{1.911801in}{0.803402in}}%
\pgfpathlineto{\pgfqpoint{1.910795in}{0.808969in}}%
\pgfpathlineto{\pgfqpoint{1.918185in}{0.809351in}}%
\pgfpathlineto{\pgfqpoint{1.923831in}{0.809696in}}%
\pgfpathlineto{\pgfqpoint{1.924611in}{0.820832in}}%
\pgfpathlineto{\pgfqpoint{1.924362in}{0.838176in}}%
\pgfpathlineto{\pgfqpoint{1.951736in}{0.838426in}}%
\pgfpathlineto{\pgfqpoint{1.951971in}{0.821985in}}%
\pgfpathlineto{\pgfqpoint{1.961337in}{0.821725in}}%
\pgfpathlineto{\pgfqpoint{1.962924in}{0.816378in}}%
\pgfpathlineto{\pgfqpoint{1.966388in}{0.815925in}}%
\pgfpathlineto{\pgfqpoint{1.967563in}{0.808897in}}%
\pgfpathlineto{\pgfqpoint{1.974688in}{0.804739in}}%
\pgfpathlineto{\pgfqpoint{1.984377in}{0.804898in}}%
\pgfpathlineto{\pgfqpoint{1.980611in}{0.811051in}}%
\pgfpathlineto{\pgfqpoint{1.984723in}{0.815682in}}%
\pgfpathlineto{\pgfqpoint{1.993621in}{0.816052in}}%
\pgfpathlineto{\pgfqpoint{1.998299in}{0.813452in}}%
\pgfpathlineto{\pgfqpoint{1.996950in}{0.809911in}}%
\pgfpathlineto{\pgfqpoint{1.986772in}{0.806255in}}%
\pgfpathlineto{\pgfqpoint{1.991139in}{0.803150in}}%
\pgfpathlineto{\pgfqpoint{1.988015in}{0.799119in}}%
\pgfpathlineto{\pgfqpoint{1.992182in}{0.792843in}}%
\pgfpathlineto{\pgfqpoint{1.984613in}{0.793228in}}%
\pgfpathlineto{\pgfqpoint{1.980623in}{0.785127in}}%
\pgfpathlineto{\pgfqpoint{1.985316in}{0.780979in}}%
\pgfpathlineto{\pgfqpoint{1.980431in}{0.778978in}}%
\pgfpathlineto{\pgfqpoint{1.985442in}{0.766767in}}%
\pgfpathlineto{\pgfqpoint{1.985739in}{0.761415in}}%
\pgfpathlineto{\pgfqpoint{1.988710in}{0.756370in}}%
\pgfpathlineto{\pgfqpoint{1.988976in}{0.750911in}}%
\pgfpathlineto{\pgfqpoint{1.985200in}{0.741858in}}%
\pgfpathclose%
\pgfusepath{fill}%
\end{pgfscope}%
\begin{pgfscope}%
\pgfpathrectangle{\pgfqpoint{0.100000in}{0.100000in}}{\pgfqpoint{3.007045in}{1.925000in}}%
\pgfusepath{clip}%
\pgfsetbuttcap%
\pgfsetmiterjoin%
\definecolor{currentfill}{rgb}{0.745713,0.845752,0.926490}%
\pgfsetfillcolor{currentfill}%
\pgfsetlinewidth{0.000000pt}%
\definecolor{currentstroke}{rgb}{0.000000,0.000000,0.000000}%
\pgfsetstrokecolor{currentstroke}%
\pgfsetstrokeopacity{0.000000}%
\pgfsetdash{}{0pt}%
\pgfpathmoveto{\pgfqpoint{2.793875in}{1.572201in}}%
\pgfpathlineto{\pgfqpoint{2.763396in}{1.560589in}}%
\pgfpathlineto{\pgfqpoint{2.763075in}{1.564033in}}%
\pgfpathlineto{\pgfqpoint{2.755118in}{1.563156in}}%
\pgfpathlineto{\pgfqpoint{2.747693in}{1.569195in}}%
\pgfpathlineto{\pgfqpoint{2.750775in}{1.572935in}}%
\pgfpathlineto{\pgfqpoint{2.747323in}{1.582366in}}%
\pgfpathlineto{\pgfqpoint{2.735921in}{1.578293in}}%
\pgfpathlineto{\pgfqpoint{2.735305in}{1.582994in}}%
\pgfpathlineto{\pgfqpoint{2.718437in}{1.632224in}}%
\pgfpathlineto{\pgfqpoint{2.713963in}{1.634075in}}%
\pgfpathlineto{\pgfqpoint{2.740305in}{1.639901in}}%
\pgfpathlineto{\pgfqpoint{2.776771in}{1.649915in}}%
\pgfpathlineto{\pgfqpoint{2.778407in}{1.644079in}}%
\pgfpathlineto{\pgfqpoint{2.777816in}{1.638392in}}%
\pgfpathlineto{\pgfqpoint{2.780372in}{1.636703in}}%
\pgfpathlineto{\pgfqpoint{2.780749in}{1.624141in}}%
\pgfpathlineto{\pgfqpoint{2.786419in}{1.617481in}}%
\pgfpathlineto{\pgfqpoint{2.787188in}{1.608949in}}%
\pgfpathlineto{\pgfqpoint{2.789757in}{1.602899in}}%
\pgfpathlineto{\pgfqpoint{2.787394in}{1.596652in}}%
\pgfpathlineto{\pgfqpoint{2.787488in}{1.586749in}}%
\pgfpathlineto{\pgfqpoint{2.793875in}{1.572201in}}%
\pgfpathclose%
\pgfusepath{fill}%
\end{pgfscope}%
\begin{pgfscope}%
\pgfpathrectangle{\pgfqpoint{0.100000in}{0.100000in}}{\pgfqpoint{3.007045in}{1.925000in}}%
\pgfusepath{clip}%
\pgfsetbuttcap%
\pgfsetmiterjoin%
\definecolor{currentfill}{rgb}{0.447843,0.697855,0.845306}%
\pgfsetfillcolor{currentfill}%
\pgfsetlinewidth{0.000000pt}%
\definecolor{currentstroke}{rgb}{0.000000,0.000000,0.000000}%
\pgfsetstrokecolor{currentstroke}%
\pgfsetstrokeopacity{0.000000}%
\pgfsetdash{}{0pt}%
\pgfpathmoveto{\pgfqpoint{1.415290in}{0.621980in}}%
\pgfpathlineto{\pgfqpoint{1.386813in}{0.624056in}}%
\pgfpathlineto{\pgfqpoint{1.391806in}{0.690752in}}%
\pgfpathlineto{\pgfqpoint{1.396350in}{0.690465in}}%
\pgfpathlineto{\pgfqpoint{1.398836in}{0.719387in}}%
\pgfpathlineto{\pgfqpoint{1.370186in}{0.721452in}}%
\pgfpathlineto{\pgfqpoint{1.372150in}{0.750256in}}%
\pgfpathlineto{\pgfqpoint{1.379483in}{0.749728in}}%
\pgfpathlineto{\pgfqpoint{1.429539in}{0.746331in}}%
\pgfpathlineto{\pgfqpoint{1.427496in}{0.717282in}}%
\pgfpathlineto{\pgfqpoint{1.425000in}{0.688410in}}%
\pgfpathlineto{\pgfqpoint{1.420501in}{0.688714in}}%
\pgfpathlineto{\pgfqpoint{1.417736in}{0.651743in}}%
\pgfpathlineto{\pgfqpoint{1.415290in}{0.621980in}}%
\pgfpathclose%
\pgfusepath{fill}%
\end{pgfscope}%
\begin{pgfscope}%
\pgfpathrectangle{\pgfqpoint{0.100000in}{0.100000in}}{\pgfqpoint{3.007045in}{1.925000in}}%
\pgfusepath{clip}%
\pgfsetbuttcap%
\pgfsetmiterjoin%
\definecolor{currentfill}{rgb}{0.093272,0.396878,0.673664}%
\pgfsetfillcolor{currentfill}%
\pgfsetlinewidth{0.000000pt}%
\definecolor{currentstroke}{rgb}{0.000000,0.000000,0.000000}%
\pgfsetstrokecolor{currentstroke}%
\pgfsetstrokeopacity{0.000000}%
\pgfsetdash{}{0pt}%
\pgfpathmoveto{\pgfqpoint{0.713739in}{0.318670in}}%
\pgfpathlineto{\pgfqpoint{0.713522in}{0.320731in}}%
\pgfpathlineto{\pgfqpoint{0.715798in}{0.320350in}}%
\pgfpathlineto{\pgfqpoint{0.717849in}{0.316514in}}%
\pgfpathlineto{\pgfqpoint{0.713739in}{0.318670in}}%
\pgfpathclose%
\pgfusepath{fill}%
\end{pgfscope}%
\begin{pgfscope}%
\pgfpathrectangle{\pgfqpoint{0.100000in}{0.100000in}}{\pgfqpoint{3.007045in}{1.925000in}}%
\pgfusepath{clip}%
\pgfsetbuttcap%
\pgfsetmiterjoin%
\definecolor{currentfill}{rgb}{0.093272,0.396878,0.673664}%
\pgfsetfillcolor{currentfill}%
\pgfsetlinewidth{0.000000pt}%
\definecolor{currentstroke}{rgb}{0.000000,0.000000,0.000000}%
\pgfsetstrokecolor{currentstroke}%
\pgfsetstrokeopacity{0.000000}%
\pgfsetdash{}{0pt}%
\pgfpathmoveto{\pgfqpoint{0.707133in}{0.321121in}}%
\pgfpathlineto{\pgfqpoint{0.708642in}{0.322633in}}%
\pgfpathlineto{\pgfqpoint{0.712152in}{0.322886in}}%
\pgfpathlineto{\pgfqpoint{0.711451in}{0.321404in}}%
\pgfpathlineto{\pgfqpoint{0.707133in}{0.321121in}}%
\pgfpathclose%
\pgfusepath{fill}%
\end{pgfscope}%
\begin{pgfscope}%
\pgfpathrectangle{\pgfqpoint{0.100000in}{0.100000in}}{\pgfqpoint{3.007045in}{1.925000in}}%
\pgfusepath{clip}%
\pgfsetbuttcap%
\pgfsetmiterjoin%
\definecolor{currentfill}{rgb}{0.093272,0.396878,0.673664}%
\pgfsetfillcolor{currentfill}%
\pgfsetlinewidth{0.000000pt}%
\definecolor{currentstroke}{rgb}{0.000000,0.000000,0.000000}%
\pgfsetstrokecolor{currentstroke}%
\pgfsetstrokeopacity{0.000000}%
\pgfsetdash{}{0pt}%
\pgfpathmoveto{\pgfqpoint{0.701239in}{0.348319in}}%
\pgfpathlineto{\pgfqpoint{0.697922in}{0.351049in}}%
\pgfpathlineto{\pgfqpoint{0.698993in}{0.352386in}}%
\pgfpathlineto{\pgfqpoint{0.701818in}{0.353160in}}%
\pgfpathlineto{\pgfqpoint{0.703476in}{0.352730in}}%
\pgfpathlineto{\pgfqpoint{0.704495in}{0.354085in}}%
\pgfpathlineto{\pgfqpoint{0.707132in}{0.354637in}}%
\pgfpathlineto{\pgfqpoint{0.708693in}{0.353375in}}%
\pgfpathlineto{\pgfqpoint{0.712359in}{0.352736in}}%
\pgfpathlineto{\pgfqpoint{0.711910in}{0.351797in}}%
\pgfpathlineto{\pgfqpoint{0.708260in}{0.352473in}}%
\pgfpathlineto{\pgfqpoint{0.704182in}{0.352191in}}%
\pgfpathlineto{\pgfqpoint{0.703882in}{0.349772in}}%
\pgfpathlineto{\pgfqpoint{0.701239in}{0.348319in}}%
\pgfpathclose%
\pgfusepath{fill}%
\end{pgfscope}%
\begin{pgfscope}%
\pgfpathrectangle{\pgfqpoint{0.100000in}{0.100000in}}{\pgfqpoint{3.007045in}{1.925000in}}%
\pgfusepath{clip}%
\pgfsetbuttcap%
\pgfsetmiterjoin%
\definecolor{currentfill}{rgb}{0.093272,0.396878,0.673664}%
\pgfsetfillcolor{currentfill}%
\pgfsetlinewidth{0.000000pt}%
\definecolor{currentstroke}{rgb}{0.000000,0.000000,0.000000}%
\pgfsetstrokecolor{currentstroke}%
\pgfsetstrokeopacity{0.000000}%
\pgfsetdash{}{0pt}%
\pgfpathmoveto{\pgfqpoint{0.758615in}{0.341736in}}%
\pgfpathlineto{\pgfqpoint{0.757407in}{0.339850in}}%
\pgfpathlineto{\pgfqpoint{0.759513in}{0.339536in}}%
\pgfpathlineto{\pgfqpoint{0.759838in}{0.338128in}}%
\pgfpathlineto{\pgfqpoint{0.761603in}{0.337668in}}%
\pgfpathlineto{\pgfqpoint{0.762040in}{0.334943in}}%
\pgfpathlineto{\pgfqpoint{0.760463in}{0.333422in}}%
\pgfpathlineto{\pgfqpoint{0.758461in}{0.333128in}}%
\pgfpathlineto{\pgfqpoint{0.753516in}{0.335984in}}%
\pgfpathlineto{\pgfqpoint{0.749366in}{0.335186in}}%
\pgfpathlineto{\pgfqpoint{0.746236in}{0.338305in}}%
\pgfpathlineto{\pgfqpoint{0.745834in}{0.340586in}}%
\pgfpathlineto{\pgfqpoint{0.747476in}{0.341739in}}%
\pgfpathlineto{\pgfqpoint{0.750424in}{0.341197in}}%
\pgfpathlineto{\pgfqpoint{0.750868in}{0.342943in}}%
\pgfpathlineto{\pgfqpoint{0.753106in}{0.339546in}}%
\pgfpathlineto{\pgfqpoint{0.754901in}{0.343295in}}%
\pgfpathlineto{\pgfqpoint{0.758150in}{0.343053in}}%
\pgfpathlineto{\pgfqpoint{0.759750in}{0.345315in}}%
\pgfpathlineto{\pgfqpoint{0.760957in}{0.344040in}}%
\pgfpathlineto{\pgfqpoint{0.760988in}{0.342549in}}%
\pgfpathlineto{\pgfqpoint{0.758615in}{0.341736in}}%
\pgfpathclose%
\pgfusepath{fill}%
\end{pgfscope}%
\begin{pgfscope}%
\pgfpathrectangle{\pgfqpoint{0.100000in}{0.100000in}}{\pgfqpoint{3.007045in}{1.925000in}}%
\pgfusepath{clip}%
\pgfsetbuttcap%
\pgfsetmiterjoin%
\definecolor{currentfill}{rgb}{0.093272,0.396878,0.673664}%
\pgfsetfillcolor{currentfill}%
\pgfsetlinewidth{0.000000pt}%
\definecolor{currentstroke}{rgb}{0.000000,0.000000,0.000000}%
\pgfsetstrokecolor{currentstroke}%
\pgfsetstrokeopacity{0.000000}%
\pgfsetdash{}{0pt}%
\pgfpathmoveto{\pgfqpoint{0.751166in}{0.331121in}}%
\pgfpathlineto{\pgfqpoint{0.751716in}{0.328502in}}%
\pgfpathlineto{\pgfqpoint{0.749924in}{0.328523in}}%
\pgfpathlineto{\pgfqpoint{0.747420in}{0.325126in}}%
\pgfpathlineto{\pgfqpoint{0.750723in}{0.323377in}}%
\pgfpathlineto{\pgfqpoint{0.746585in}{0.321465in}}%
\pgfpathlineto{\pgfqpoint{0.744840in}{0.322182in}}%
\pgfpathlineto{\pgfqpoint{0.743772in}{0.324350in}}%
\pgfpathlineto{\pgfqpoint{0.742472in}{0.323661in}}%
\pgfpathlineto{\pgfqpoint{0.742188in}{0.321974in}}%
\pgfpathlineto{\pgfqpoint{0.740508in}{0.321472in}}%
\pgfpathlineto{\pgfqpoint{0.734035in}{0.323447in}}%
\pgfpathlineto{\pgfqpoint{0.734346in}{0.321507in}}%
\pgfpathlineto{\pgfqpoint{0.732208in}{0.320942in}}%
\pgfpathlineto{\pgfqpoint{0.729299in}{0.323471in}}%
\pgfpathlineto{\pgfqpoint{0.727715in}{0.323803in}}%
\pgfpathlineto{\pgfqpoint{0.723753in}{0.321601in}}%
\pgfpathlineto{\pgfqpoint{0.719755in}{0.321321in}}%
\pgfpathlineto{\pgfqpoint{0.720327in}{0.323273in}}%
\pgfpathlineto{\pgfqpoint{0.724725in}{0.324272in}}%
\pgfpathlineto{\pgfqpoint{0.724952in}{0.327121in}}%
\pgfpathlineto{\pgfqpoint{0.722436in}{0.327416in}}%
\pgfpathlineto{\pgfqpoint{0.719318in}{0.325777in}}%
\pgfpathlineto{\pgfqpoint{0.718473in}{0.330811in}}%
\pgfpathlineto{\pgfqpoint{0.720682in}{0.334691in}}%
\pgfpathlineto{\pgfqpoint{0.719629in}{0.337498in}}%
\pgfpathlineto{\pgfqpoint{0.724542in}{0.340506in}}%
\pgfpathlineto{\pgfqpoint{0.728409in}{0.340545in}}%
\pgfpathlineto{\pgfqpoint{0.730761in}{0.339863in}}%
\pgfpathlineto{\pgfqpoint{0.732294in}{0.338142in}}%
\pgfpathlineto{\pgfqpoint{0.730507in}{0.333695in}}%
\pgfpathlineto{\pgfqpoint{0.733490in}{0.335333in}}%
\pgfpathlineto{\pgfqpoint{0.734831in}{0.339974in}}%
\pgfpathlineto{\pgfqpoint{0.736842in}{0.341018in}}%
\pgfpathlineto{\pgfqpoint{0.738720in}{0.340689in}}%
\pgfpathlineto{\pgfqpoint{0.738975in}{0.337460in}}%
\pgfpathlineto{\pgfqpoint{0.740912in}{0.337016in}}%
\pgfpathlineto{\pgfqpoint{0.741053in}{0.339842in}}%
\pgfpathlineto{\pgfqpoint{0.745035in}{0.337912in}}%
\pgfpathlineto{\pgfqpoint{0.746149in}{0.335843in}}%
\pgfpathlineto{\pgfqpoint{0.748088in}{0.333888in}}%
\pgfpathlineto{\pgfqpoint{0.746213in}{0.332218in}}%
\pgfpathlineto{\pgfqpoint{0.751350in}{0.331886in}}%
\pgfpathlineto{\pgfqpoint{0.751166in}{0.331121in}}%
\pgfpathclose%
\pgfusepath{fill}%
\end{pgfscope}%
\begin{pgfscope}%
\pgfpathrectangle{\pgfqpoint{0.100000in}{0.100000in}}{\pgfqpoint{3.007045in}{1.925000in}}%
\pgfusepath{clip}%
\pgfsetbuttcap%
\pgfsetmiterjoin%
\definecolor{currentfill}{rgb}{0.093272,0.396878,0.673664}%
\pgfsetfillcolor{currentfill}%
\pgfsetlinewidth{0.000000pt}%
\definecolor{currentstroke}{rgb}{0.000000,0.000000,0.000000}%
\pgfsetstrokecolor{currentstroke}%
\pgfsetstrokeopacity{0.000000}%
\pgfsetdash{}{0pt}%
\pgfpathmoveto{\pgfqpoint{0.757017in}{0.355194in}}%
\pgfpathlineto{\pgfqpoint{0.755520in}{0.355499in}}%
\pgfpathlineto{\pgfqpoint{0.753460in}{0.353984in}}%
\pgfpathlineto{\pgfqpoint{0.748969in}{0.353910in}}%
\pgfpathlineto{\pgfqpoint{0.747754in}{0.354827in}}%
\pgfpathlineto{\pgfqpoint{0.745252in}{0.353754in}}%
\pgfpathlineto{\pgfqpoint{0.742936in}{0.353492in}}%
\pgfpathlineto{\pgfqpoint{0.741809in}{0.351678in}}%
\pgfpathlineto{\pgfqpoint{0.740287in}{0.352714in}}%
\pgfpathlineto{\pgfqpoint{0.739684in}{0.350603in}}%
\pgfpathlineto{\pgfqpoint{0.733752in}{0.349525in}}%
\pgfpathlineto{\pgfqpoint{0.731314in}{0.349960in}}%
\pgfpathlineto{\pgfqpoint{0.728847in}{0.351743in}}%
\pgfpathlineto{\pgfqpoint{0.725995in}{0.352347in}}%
\pgfpathlineto{\pgfqpoint{0.725319in}{0.351090in}}%
\pgfpathlineto{\pgfqpoint{0.721782in}{0.351198in}}%
\pgfpathlineto{\pgfqpoint{0.719594in}{0.350212in}}%
\pgfpathlineto{\pgfqpoint{0.719062in}{0.352698in}}%
\pgfpathlineto{\pgfqpoint{0.717117in}{0.351179in}}%
\pgfpathlineto{\pgfqpoint{0.714893in}{0.351866in}}%
\pgfpathlineto{\pgfqpoint{0.713962in}{0.350368in}}%
\pgfpathlineto{\pgfqpoint{0.712047in}{0.351741in}}%
\pgfpathlineto{\pgfqpoint{0.715201in}{0.352581in}}%
\pgfpathlineto{\pgfqpoint{0.717812in}{0.354377in}}%
\pgfpathlineto{\pgfqpoint{0.719341in}{0.353168in}}%
\pgfpathlineto{\pgfqpoint{0.722429in}{0.352828in}}%
\pgfpathlineto{\pgfqpoint{0.726759in}{0.358371in}}%
\pgfpathlineto{\pgfqpoint{0.728487in}{0.358327in}}%
\pgfpathlineto{\pgfqpoint{0.730107in}{0.357059in}}%
\pgfpathlineto{\pgfqpoint{0.731400in}{0.357708in}}%
\pgfpathlineto{\pgfqpoint{0.734357in}{0.355834in}}%
\pgfpathlineto{\pgfqpoint{0.737813in}{0.354884in}}%
\pgfpathlineto{\pgfqpoint{0.738983in}{0.355722in}}%
\pgfpathlineto{\pgfqpoint{0.740865in}{0.355528in}}%
\pgfpathlineto{\pgfqpoint{0.743549in}{0.359090in}}%
\pgfpathlineto{\pgfqpoint{0.748423in}{0.355398in}}%
\pgfpathlineto{\pgfqpoint{0.749665in}{0.357072in}}%
\pgfpathlineto{\pgfqpoint{0.753292in}{0.354357in}}%
\pgfpathlineto{\pgfqpoint{0.755001in}{0.356698in}}%
\pgfpathlineto{\pgfqpoint{0.757017in}{0.355194in}}%
\pgfpathclose%
\pgfusepath{fill}%
\end{pgfscope}%
\begin{pgfscope}%
\pgfpathrectangle{\pgfqpoint{0.100000in}{0.100000in}}{\pgfqpoint{3.007045in}{1.925000in}}%
\pgfusepath{clip}%
\pgfsetbuttcap%
\pgfsetmiterjoin%
\definecolor{currentfill}{rgb}{0.652211,0.806013,0.893764}%
\pgfsetfillcolor{currentfill}%
\pgfsetlinewidth{0.000000pt}%
\definecolor{currentstroke}{rgb}{0.000000,0.000000,0.000000}%
\pgfsetstrokecolor{currentstroke}%
\pgfsetstrokeopacity{0.000000}%
\pgfsetdash{}{0pt}%
\pgfpathmoveto{\pgfqpoint{1.020483in}{1.573439in}}%
\pgfpathlineto{\pgfqpoint{1.018543in}{1.560865in}}%
\pgfpathlineto{\pgfqpoint{1.015015in}{1.563234in}}%
\pgfpathlineto{\pgfqpoint{1.009746in}{1.577792in}}%
\pgfpathlineto{\pgfqpoint{1.006564in}{1.581699in}}%
\pgfpathlineto{\pgfqpoint{1.001653in}{1.579735in}}%
\pgfpathlineto{\pgfqpoint{0.987365in}{1.582354in}}%
\pgfpathlineto{\pgfqpoint{0.988238in}{1.587058in}}%
\pgfpathlineto{\pgfqpoint{0.975154in}{1.589701in}}%
\pgfpathlineto{\pgfqpoint{0.971016in}{1.593385in}}%
\pgfpathlineto{\pgfqpoint{0.972074in}{1.604894in}}%
\pgfpathlineto{\pgfqpoint{0.965491in}{1.606172in}}%
\pgfpathlineto{\pgfqpoint{0.966618in}{1.611886in}}%
\pgfpathlineto{\pgfqpoint{0.960968in}{1.612983in}}%
\pgfpathlineto{\pgfqpoint{0.964235in}{1.629990in}}%
\pgfpathlineto{\pgfqpoint{0.963238in}{1.636470in}}%
\pgfpathlineto{\pgfqpoint{0.959812in}{1.636263in}}%
\pgfpathlineto{\pgfqpoint{0.955356in}{1.640352in}}%
\pgfpathlineto{\pgfqpoint{0.957302in}{1.649942in}}%
\pgfpathlineto{\pgfqpoint{0.964271in}{1.652566in}}%
\pgfpathlineto{\pgfqpoint{0.966686in}{1.656395in}}%
\pgfpathlineto{\pgfqpoint{0.981446in}{1.653571in}}%
\pgfpathlineto{\pgfqpoint{0.988214in}{1.659319in}}%
\pgfpathlineto{\pgfqpoint{0.994682in}{1.658385in}}%
\pgfpathlineto{\pgfqpoint{0.997011in}{1.655494in}}%
\pgfpathlineto{\pgfqpoint{1.004464in}{1.652753in}}%
\pgfpathlineto{\pgfqpoint{1.010507in}{1.658210in}}%
\pgfpathlineto{\pgfqpoint{1.016570in}{1.660701in}}%
\pgfpathlineto{\pgfqpoint{1.019996in}{1.667445in}}%
\pgfpathlineto{\pgfqpoint{1.024131in}{1.670887in}}%
\pgfpathlineto{\pgfqpoint{1.025044in}{1.675119in}}%
\pgfpathlineto{\pgfqpoint{1.037941in}{1.673119in}}%
\pgfpathlineto{\pgfqpoint{1.067509in}{1.667877in}}%
\pgfpathlineto{\pgfqpoint{1.073596in}{1.668769in}}%
\pgfpathlineto{\pgfqpoint{1.102024in}{1.663968in}}%
\pgfpathlineto{\pgfqpoint{1.101102in}{1.658254in}}%
\pgfpathlineto{\pgfqpoint{1.102412in}{1.652202in}}%
\pgfpathlineto{\pgfqpoint{1.110938in}{1.650768in}}%
\pgfpathlineto{\pgfqpoint{1.109783in}{1.645289in}}%
\pgfpathlineto{\pgfqpoint{1.105971in}{1.645829in}}%
\pgfpathlineto{\pgfqpoint{1.101489in}{1.634717in}}%
\pgfpathlineto{\pgfqpoint{1.099750in}{1.623315in}}%
\pgfpathlineto{\pgfqpoint{1.094101in}{1.624217in}}%
\pgfpathlineto{\pgfqpoint{1.081871in}{1.620429in}}%
\pgfpathlineto{\pgfqpoint{1.079975in}{1.609049in}}%
\pgfpathlineto{\pgfqpoint{1.073953in}{1.610065in}}%
\pgfpathlineto{\pgfqpoint{1.072000in}{1.598524in}}%
\pgfpathlineto{\pgfqpoint{1.084238in}{1.596483in}}%
\pgfpathlineto{\pgfqpoint{1.082392in}{1.585392in}}%
\pgfpathlineto{\pgfqpoint{1.040422in}{1.591840in}}%
\pgfpathlineto{\pgfqpoint{1.024388in}{1.595276in}}%
\pgfpathlineto{\pgfqpoint{1.020483in}{1.573439in}}%
\pgfpathclose%
\pgfusepath{fill}%
\end{pgfscope}%
\begin{pgfscope}%
\pgfpathrectangle{\pgfqpoint{0.100000in}{0.100000in}}{\pgfqpoint{3.007045in}{1.925000in}}%
\pgfusepath{clip}%
\pgfsetbuttcap%
\pgfsetmiterjoin%
\definecolor{currentfill}{rgb}{0.417086,0.680631,0.838231}%
\pgfsetfillcolor{currentfill}%
\pgfsetlinewidth{0.000000pt}%
\definecolor{currentstroke}{rgb}{0.000000,0.000000,0.000000}%
\pgfsetstrokecolor{currentstroke}%
\pgfsetstrokeopacity{0.000000}%
\pgfsetdash{}{0pt}%
\pgfpathmoveto{\pgfqpoint{1.607157in}{1.202049in}}%
\pgfpathlineto{\pgfqpoint{1.595536in}{1.202468in}}%
\pgfpathlineto{\pgfqpoint{1.584352in}{1.202908in}}%
\pgfpathlineto{\pgfqpoint{1.586476in}{1.248770in}}%
\pgfpathlineto{\pgfqpoint{1.631190in}{1.247136in}}%
\pgfpathlineto{\pgfqpoint{1.630558in}{1.224197in}}%
\pgfpathlineto{\pgfqpoint{1.608008in}{1.224954in}}%
\pgfpathlineto{\pgfqpoint{1.607157in}{1.202049in}}%
\pgfpathclose%
\pgfusepath{fill}%
\end{pgfscope}%
\begin{pgfscope}%
\pgfpathrectangle{\pgfqpoint{0.100000in}{0.100000in}}{\pgfqpoint{3.007045in}{1.925000in}}%
\pgfusepath{clip}%
\pgfsetbuttcap%
\pgfsetmiterjoin%
\definecolor{currentfill}{rgb}{0.622684,0.793464,0.883429}%
\pgfsetfillcolor{currentfill}%
\pgfsetlinewidth{0.000000pt}%
\definecolor{currentstroke}{rgb}{0.000000,0.000000,0.000000}%
\pgfsetstrokecolor{currentstroke}%
\pgfsetstrokeopacity{0.000000}%
\pgfsetdash{}{0pt}%
\pgfpathmoveto{\pgfqpoint{0.596099in}{1.599150in}}%
\pgfpathlineto{\pgfqpoint{0.573938in}{1.605496in}}%
\pgfpathlineto{\pgfqpoint{0.530203in}{1.618670in}}%
\pgfpathlineto{\pgfqpoint{0.499509in}{1.627848in}}%
\pgfpathlineto{\pgfqpoint{0.504292in}{1.636653in}}%
\pgfpathlineto{\pgfqpoint{0.505425in}{1.642739in}}%
\pgfpathlineto{\pgfqpoint{0.512713in}{1.647749in}}%
\pgfpathlineto{\pgfqpoint{0.519173in}{1.654672in}}%
\pgfpathlineto{\pgfqpoint{0.521441in}{1.664987in}}%
\pgfpathlineto{\pgfqpoint{0.521156in}{1.673183in}}%
\pgfpathlineto{\pgfqpoint{0.523556in}{1.681083in}}%
\pgfpathlineto{\pgfqpoint{0.527922in}{1.686852in}}%
\pgfpathlineto{\pgfqpoint{0.530035in}{1.695577in}}%
\pgfpathlineto{\pgfqpoint{0.533873in}{1.701067in}}%
\pgfpathlineto{\pgfqpoint{0.563962in}{1.691905in}}%
\pgfpathlineto{\pgfqpoint{0.596303in}{1.682378in}}%
\pgfpathlineto{\pgfqpoint{0.594595in}{1.678719in}}%
\pgfpathlineto{\pgfqpoint{0.588656in}{1.658125in}}%
\pgfpathlineto{\pgfqpoint{0.605211in}{1.653481in}}%
\pgfpathlineto{\pgfqpoint{0.609918in}{1.648785in}}%
\pgfpathlineto{\pgfqpoint{0.608444in}{1.643504in}}%
\pgfpathlineto{\pgfqpoint{0.619536in}{1.640423in}}%
\pgfpathlineto{\pgfqpoint{0.613279in}{1.618283in}}%
\pgfpathlineto{\pgfqpoint{0.607867in}{1.619794in}}%
\pgfpathlineto{\pgfqpoint{0.603136in}{1.603204in}}%
\pgfpathlineto{\pgfqpoint{0.597590in}{1.604775in}}%
\pgfpathlineto{\pgfqpoint{0.596099in}{1.599150in}}%
\pgfpathclose%
\pgfusepath{fill}%
\end{pgfscope}%
\begin{pgfscope}%
\pgfpathrectangle{\pgfqpoint{0.100000in}{0.100000in}}{\pgfqpoint{3.007045in}{1.925000in}}%
\pgfusepath{clip}%
\pgfsetbuttcap%
\pgfsetmiterjoin%
\definecolor{currentfill}{rgb}{0.326290,0.618624,0.802799}%
\pgfsetfillcolor{currentfill}%
\pgfsetlinewidth{0.000000pt}%
\definecolor{currentstroke}{rgb}{0.000000,0.000000,0.000000}%
\pgfsetstrokecolor{currentstroke}%
\pgfsetstrokeopacity{0.000000}%
\pgfsetdash{}{0pt}%
\pgfpathmoveto{\pgfqpoint{1.446935in}{1.181960in}}%
\pgfpathlineto{\pgfqpoint{1.414739in}{1.184340in}}%
\pgfpathlineto{\pgfqpoint{1.414767in}{1.184726in}}%
\pgfpathlineto{\pgfqpoint{1.418635in}{1.235643in}}%
\pgfpathlineto{\pgfqpoint{1.455104in}{1.233038in}}%
\pgfpathlineto{\pgfqpoint{1.453477in}{1.210184in}}%
\pgfpathlineto{\pgfqpoint{1.449096in}{1.210470in}}%
\pgfpathlineto{\pgfqpoint{1.446935in}{1.181960in}}%
\pgfpathclose%
\pgfusepath{fill}%
\end{pgfscope}%
\begin{pgfscope}%
\pgfpathrectangle{\pgfqpoint{0.100000in}{0.100000in}}{\pgfqpoint{3.007045in}{1.925000in}}%
\pgfusepath{clip}%
\pgfsetbuttcap%
\pgfsetmiterjoin%
\definecolor{currentfill}{rgb}{0.784591,0.864237,0.939962}%
\pgfsetfillcolor{currentfill}%
\pgfsetlinewidth{0.000000pt}%
\definecolor{currentstroke}{rgb}{0.000000,0.000000,0.000000}%
\pgfsetstrokecolor{currentstroke}%
\pgfsetstrokeopacity{0.000000}%
\pgfsetdash{}{0pt}%
\pgfpathmoveto{\pgfqpoint{0.774157in}{0.904247in}}%
\pgfpathlineto{\pgfqpoint{0.778392in}{0.924882in}}%
\pgfpathlineto{\pgfqpoint{0.794371in}{1.003188in}}%
\pgfpathlineto{\pgfqpoint{0.798014in}{1.020963in}}%
\pgfpathlineto{\pgfqpoint{0.800968in}{1.030983in}}%
\pgfpathlineto{\pgfqpoint{0.800098in}{1.036517in}}%
\pgfpathlineto{\pgfqpoint{0.803014in}{1.039936in}}%
\pgfpathlineto{\pgfqpoint{0.807519in}{1.038421in}}%
\pgfpathlineto{\pgfqpoint{0.809522in}{1.041947in}}%
\pgfpathlineto{\pgfqpoint{0.818051in}{1.044246in}}%
\pgfpathlineto{\pgfqpoint{0.832280in}{1.046410in}}%
\pgfpathlineto{\pgfqpoint{0.839109in}{1.048947in}}%
\pgfpathlineto{\pgfqpoint{0.842909in}{1.052579in}}%
\pgfpathlineto{\pgfqpoint{0.842643in}{1.059879in}}%
\pgfpathlineto{\pgfqpoint{0.845589in}{1.070038in}}%
\pgfpathlineto{\pgfqpoint{0.854802in}{1.090140in}}%
\pgfpathlineto{\pgfqpoint{0.836260in}{1.093783in}}%
\pgfpathlineto{\pgfqpoint{0.843170in}{1.128894in}}%
\pgfpathlineto{\pgfqpoint{0.854337in}{1.126666in}}%
\pgfpathlineto{\pgfqpoint{0.870763in}{1.123009in}}%
\pgfpathlineto{\pgfqpoint{0.937370in}{1.110856in}}%
\pgfpathlineto{\pgfqpoint{0.958933in}{1.107125in}}%
\pgfpathlineto{\pgfqpoint{0.951769in}{1.095219in}}%
\pgfpathlineto{\pgfqpoint{0.944689in}{1.092843in}}%
\pgfpathlineto{\pgfqpoint{0.941624in}{1.085677in}}%
\pgfpathlineto{\pgfqpoint{0.932657in}{1.082202in}}%
\pgfpathlineto{\pgfqpoint{0.926450in}{1.083530in}}%
\pgfpathlineto{\pgfqpoint{0.923032in}{1.079857in}}%
\pgfpathlineto{\pgfqpoint{0.915481in}{1.078785in}}%
\pgfpathlineto{\pgfqpoint{0.947453in}{1.073147in}}%
\pgfpathlineto{\pgfqpoint{0.916273in}{0.895048in}}%
\pgfpathlineto{\pgfqpoint{0.911088in}{0.896114in}}%
\pgfpathlineto{\pgfqpoint{0.903247in}{0.903940in}}%
\pgfpathlineto{\pgfqpoint{0.896277in}{0.906751in}}%
\pgfpathlineto{\pgfqpoint{0.890889in}{0.911616in}}%
\pgfpathlineto{\pgfqpoint{0.883290in}{0.911640in}}%
\pgfpathlineto{\pgfqpoint{0.876178in}{0.916207in}}%
\pgfpathlineto{\pgfqpoint{0.868902in}{0.911853in}}%
\pgfpathlineto{\pgfqpoint{0.863233in}{0.897098in}}%
\pgfpathlineto{\pgfqpoint{0.875838in}{0.894759in}}%
\pgfpathlineto{\pgfqpoint{0.873121in}{0.885088in}}%
\pgfpathlineto{\pgfqpoint{0.860696in}{0.887387in}}%
\pgfpathlineto{\pgfqpoint{0.837638in}{0.894986in}}%
\pgfpathlineto{\pgfqpoint{0.829566in}{0.885369in}}%
\pgfpathlineto{\pgfqpoint{0.805861in}{0.897916in}}%
\pgfpathlineto{\pgfqpoint{0.774157in}{0.904247in}}%
\pgfpathclose%
\pgfusepath{fill}%
\end{pgfscope}%
\begin{pgfscope}%
\pgfpathrectangle{\pgfqpoint{0.100000in}{0.100000in}}{\pgfqpoint{3.007045in}{1.925000in}}%
\pgfusepath{clip}%
\pgfsetbuttcap%
\pgfsetmiterjoin%
\definecolor{currentfill}{rgb}{0.529412,0.742637,0.862899}%
\pgfsetfillcolor{currentfill}%
\pgfsetlinewidth{0.000000pt}%
\definecolor{currentstroke}{rgb}{0.000000,0.000000,0.000000}%
\pgfsetstrokecolor{currentstroke}%
\pgfsetstrokeopacity{0.000000}%
\pgfsetdash{}{0pt}%
\pgfpathmoveto{\pgfqpoint{1.330895in}{1.191503in}}%
\pgfpathlineto{\pgfqpoint{1.325845in}{1.145704in}}%
\pgfpathlineto{\pgfqpoint{1.308723in}{1.147516in}}%
\pgfpathlineto{\pgfqpoint{1.310620in}{1.164502in}}%
\pgfpathlineto{\pgfqpoint{1.245658in}{1.171547in}}%
\pgfpathlineto{\pgfqpoint{1.242494in}{1.143237in}}%
\pgfpathlineto{\pgfqpoint{1.209760in}{1.146848in}}%
\pgfpathlineto{\pgfqpoint{1.213877in}{1.153875in}}%
\pgfpathlineto{\pgfqpoint{1.213448in}{1.158253in}}%
\pgfpathlineto{\pgfqpoint{1.208767in}{1.163511in}}%
\pgfpathlineto{\pgfqpoint{1.200990in}{1.169461in}}%
\pgfpathlineto{\pgfqpoint{1.201403in}{1.171929in}}%
\pgfpathlineto{\pgfqpoint{1.202987in}{1.177221in}}%
\pgfpathlineto{\pgfqpoint{1.203807in}{1.188442in}}%
\pgfpathlineto{\pgfqpoint{1.216380in}{1.193870in}}%
\pgfpathlineto{\pgfqpoint{1.226581in}{1.206947in}}%
\pgfpathlineto{\pgfqpoint{1.220004in}{1.212317in}}%
\pgfpathlineto{\pgfqpoint{1.225797in}{1.217368in}}%
\pgfpathlineto{\pgfqpoint{1.232976in}{1.220918in}}%
\pgfpathlineto{\pgfqpoint{1.234437in}{1.230208in}}%
\pgfpathlineto{\pgfqpoint{1.237002in}{1.232833in}}%
\pgfpathlineto{\pgfqpoint{1.237983in}{1.247422in}}%
\pgfpathlineto{\pgfqpoint{1.267871in}{1.244110in}}%
\pgfpathlineto{\pgfqpoint{1.265951in}{1.226957in}}%
\pgfpathlineto{\pgfqpoint{1.311373in}{1.222190in}}%
\pgfpathlineto{\pgfqpoint{1.333705in}{1.220024in}}%
\pgfpathlineto{\pgfqpoint{1.330895in}{1.191503in}}%
\pgfpathclose%
\pgfusepath{fill}%
\end{pgfscope}%
\begin{pgfscope}%
\pgfpathrectangle{\pgfqpoint{0.100000in}{0.100000in}}{\pgfqpoint{3.007045in}{1.925000in}}%
\pgfusepath{clip}%
\pgfsetbuttcap%
\pgfsetmiterjoin%
\definecolor{currentfill}{rgb}{0.080969,0.381130,0.661361}%
\pgfsetfillcolor{currentfill}%
\pgfsetlinewidth{0.000000pt}%
\definecolor{currentstroke}{rgb}{0.000000,0.000000,0.000000}%
\pgfsetstrokecolor{currentstroke}%
\pgfsetstrokeopacity{0.000000}%
\pgfsetdash{}{0pt}%
\pgfpathmoveto{\pgfqpoint{1.514641in}{0.483384in}}%
\pgfpathlineto{\pgfqpoint{1.473940in}{0.485183in}}%
\pgfpathlineto{\pgfqpoint{1.475915in}{0.520991in}}%
\pgfpathlineto{\pgfqpoint{1.441878in}{0.523060in}}%
\pgfpathlineto{\pgfqpoint{1.444627in}{0.567276in}}%
\pgfpathlineto{\pgfqpoint{1.429685in}{0.568203in}}%
\pgfpathlineto{\pgfqpoint{1.431514in}{0.595996in}}%
\pgfpathlineto{\pgfqpoint{1.479753in}{0.593339in}}%
\pgfpathlineto{\pgfqpoint{1.478167in}{0.565413in}}%
\pgfpathlineto{\pgfqpoint{1.498977in}{0.564304in}}%
\pgfpathlineto{\pgfqpoint{1.498020in}{0.549906in}}%
\pgfpathlineto{\pgfqpoint{1.501924in}{0.549696in}}%
\pgfpathlineto{\pgfqpoint{1.501272in}{0.538689in}}%
\pgfpathlineto{\pgfqpoint{1.506378in}{0.538351in}}%
\pgfpathlineto{\pgfqpoint{1.505414in}{0.519674in}}%
\pgfpathlineto{\pgfqpoint{1.516499in}{0.519146in}}%
\pgfpathlineto{\pgfqpoint{1.514641in}{0.483384in}}%
\pgfpathclose%
\pgfusepath{fill}%
\end{pgfscope}%
\begin{pgfscope}%
\pgfpathrectangle{\pgfqpoint{0.100000in}{0.100000in}}{\pgfqpoint{3.007045in}{1.925000in}}%
\pgfusepath{clip}%
\pgfsetbuttcap%
\pgfsetmiterjoin%
\definecolor{currentfill}{rgb}{0.485490,0.718524,0.853426}%
\pgfsetfillcolor{currentfill}%
\pgfsetlinewidth{0.000000pt}%
\definecolor{currentstroke}{rgb}{0.000000,0.000000,0.000000}%
\pgfsetstrokecolor{currentstroke}%
\pgfsetstrokeopacity{0.000000}%
\pgfsetdash{}{0pt}%
\pgfpathmoveto{\pgfqpoint{0.989331in}{0.759721in}}%
\pgfpathlineto{\pgfqpoint{0.997334in}{0.810949in}}%
\pgfpathlineto{\pgfqpoint{1.054598in}{0.801868in}}%
\pgfpathlineto{\pgfqpoint{1.061462in}{0.800953in}}%
\pgfpathlineto{\pgfqpoint{1.062578in}{0.792196in}}%
\pgfpathlineto{\pgfqpoint{1.061362in}{0.785689in}}%
\pgfpathlineto{\pgfqpoint{1.064806in}{0.782977in}}%
\pgfpathlineto{\pgfqpoint{1.063832in}{0.777370in}}%
\pgfpathlineto{\pgfqpoint{1.065981in}{0.771962in}}%
\pgfpathlineto{\pgfqpoint{1.063758in}{0.762298in}}%
\pgfpathlineto{\pgfqpoint{1.070654in}{0.759703in}}%
\pgfpathlineto{\pgfqpoint{1.087735in}{0.757358in}}%
\pgfpathlineto{\pgfqpoint{1.080529in}{0.703340in}}%
\pgfpathlineto{\pgfqpoint{1.029591in}{0.710510in}}%
\pgfpathlineto{\pgfqpoint{1.025251in}{0.680923in}}%
\pgfpathlineto{\pgfqpoint{0.978016in}{0.688001in}}%
\pgfpathlineto{\pgfqpoint{0.989331in}{0.759721in}}%
\pgfpathclose%
\pgfusepath{fill}%
\end{pgfscope}%
\begin{pgfscope}%
\pgfpathrectangle{\pgfqpoint{0.100000in}{0.100000in}}{\pgfqpoint{3.007045in}{1.925000in}}%
\pgfusepath{clip}%
\pgfsetbuttcap%
\pgfsetmiterjoin%
\definecolor{currentfill}{rgb}{0.296025,0.597955,0.790988}%
\pgfsetfillcolor{currentfill}%
\pgfsetlinewidth{0.000000pt}%
\definecolor{currentstroke}{rgb}{0.000000,0.000000,0.000000}%
\pgfsetstrokecolor{currentstroke}%
\pgfsetstrokeopacity{0.000000}%
\pgfsetdash{}{0pt}%
\pgfpathmoveto{\pgfqpoint{1.669507in}{1.592201in}}%
\pgfpathlineto{\pgfqpoint{1.634891in}{1.593241in}}%
\pgfpathlineto{\pgfqpoint{1.587550in}{1.595331in}}%
\pgfpathlineto{\pgfqpoint{1.588665in}{1.618076in}}%
\pgfpathlineto{\pgfqpoint{1.633089in}{1.616328in}}%
\pgfpathlineto{\pgfqpoint{1.633892in}{1.639443in}}%
\pgfpathlineto{\pgfqpoint{1.649938in}{1.638890in}}%
\pgfpathlineto{\pgfqpoint{1.668280in}{1.638340in}}%
\pgfpathlineto{\pgfqpoint{1.667675in}{1.615294in}}%
\pgfpathlineto{\pgfqpoint{1.668146in}{1.596346in}}%
\pgfpathlineto{\pgfqpoint{1.669507in}{1.592201in}}%
\pgfpathclose%
\pgfusepath{fill}%
\end{pgfscope}%
\begin{pgfscope}%
\pgfpathrectangle{\pgfqpoint{0.100000in}{0.100000in}}{\pgfqpoint{3.007045in}{1.925000in}}%
\pgfusepath{clip}%
\pgfsetbuttcap%
\pgfsetmiterjoin%
\definecolor{currentfill}{rgb}{0.696501,0.824837,0.909266}%
\pgfsetfillcolor{currentfill}%
\pgfsetlinewidth{0.000000pt}%
\definecolor{currentstroke}{rgb}{0.000000,0.000000,0.000000}%
\pgfsetstrokecolor{currentstroke}%
\pgfsetstrokeopacity{0.000000}%
\pgfsetdash{}{0pt}%
\pgfpathmoveto{\pgfqpoint{2.826283in}{1.627777in}}%
\pgfpathlineto{\pgfqpoint{2.827912in}{1.634106in}}%
\pgfpathlineto{\pgfqpoint{2.822013in}{1.635422in}}%
\pgfpathlineto{\pgfqpoint{2.823531in}{1.641882in}}%
\pgfpathlineto{\pgfqpoint{2.815249in}{1.643795in}}%
\pgfpathlineto{\pgfqpoint{2.816481in}{1.647652in}}%
\pgfpathlineto{\pgfqpoint{2.812663in}{1.658624in}}%
\pgfpathlineto{\pgfqpoint{2.823785in}{1.661144in}}%
\pgfpathlineto{\pgfqpoint{2.860264in}{1.671267in}}%
\pgfpathlineto{\pgfqpoint{2.862575in}{1.687976in}}%
\pgfpathlineto{\pgfqpoint{2.873035in}{1.690583in}}%
\pgfpathlineto{\pgfqpoint{2.874149in}{1.694991in}}%
\pgfpathlineto{\pgfqpoint{2.885278in}{1.661827in}}%
\pgfpathlineto{\pgfqpoint{2.895122in}{1.630668in}}%
\pgfpathlineto{\pgfqpoint{2.894447in}{1.627269in}}%
\pgfpathlineto{\pgfqpoint{2.884731in}{1.625099in}}%
\pgfpathlineto{\pgfqpoint{2.884101in}{1.617605in}}%
\pgfpathlineto{\pgfqpoint{2.877796in}{1.621050in}}%
\pgfpathlineto{\pgfqpoint{2.875904in}{1.623095in}}%
\pgfpathlineto{\pgfqpoint{2.869747in}{1.621836in}}%
\pgfpathlineto{\pgfqpoint{2.868383in}{1.624903in}}%
\pgfpathlineto{\pgfqpoint{2.862905in}{1.625347in}}%
\pgfpathlineto{\pgfqpoint{2.858546in}{1.629287in}}%
\pgfpathlineto{\pgfqpoint{2.854814in}{1.623529in}}%
\pgfpathlineto{\pgfqpoint{2.849656in}{1.622164in}}%
\pgfpathlineto{\pgfqpoint{2.846880in}{1.617175in}}%
\pgfpathlineto{\pgfqpoint{2.848338in}{1.611817in}}%
\pgfpathlineto{\pgfqpoint{2.837477in}{1.608579in}}%
\pgfpathlineto{\pgfqpoint{2.834224in}{1.609178in}}%
\pgfpathlineto{\pgfqpoint{2.837240in}{1.625008in}}%
\pgfpathlineto{\pgfqpoint{2.826283in}{1.627777in}}%
\pgfpathclose%
\pgfusepath{fill}%
\end{pgfscope}%
\begin{pgfscope}%
\pgfpathrectangle{\pgfqpoint{0.100000in}{0.100000in}}{\pgfqpoint{3.007045in}{1.925000in}}%
\pgfusepath{clip}%
\pgfsetbuttcap%
\pgfsetmiterjoin%
\definecolor{currentfill}{rgb}{0.162907,0.476632,0.727059}%
\pgfsetfillcolor{currentfill}%
\pgfsetlinewidth{0.000000pt}%
\definecolor{currentstroke}{rgb}{0.000000,0.000000,0.000000}%
\pgfsetstrokecolor{currentstroke}%
\pgfsetstrokeopacity{0.000000}%
\pgfsetdash{}{0pt}%
\pgfpathmoveto{\pgfqpoint{1.610638in}{1.328045in}}%
\pgfpathlineto{\pgfqpoint{1.611668in}{1.362462in}}%
\pgfpathlineto{\pgfqpoint{1.657059in}{1.361181in}}%
\pgfpathlineto{\pgfqpoint{1.656900in}{1.355436in}}%
\pgfpathlineto{\pgfqpoint{1.656422in}{1.338221in}}%
\pgfpathlineto{\pgfqpoint{1.673497in}{1.337786in}}%
\pgfpathlineto{\pgfqpoint{1.672961in}{1.314877in}}%
\pgfpathlineto{\pgfqpoint{1.633082in}{1.315887in}}%
\pgfpathlineto{\pgfqpoint{1.633318in}{1.327373in}}%
\pgfpathlineto{\pgfqpoint{1.610638in}{1.328045in}}%
\pgfpathclose%
\pgfusepath{fill}%
\end{pgfscope}%
\begin{pgfscope}%
\pgfpathrectangle{\pgfqpoint{0.100000in}{0.100000in}}{\pgfqpoint{3.007045in}{1.925000in}}%
\pgfusepath{clip}%
\pgfsetbuttcap%
\pgfsetmiterjoin%
\definecolor{currentfill}{rgb}{0.441569,0.694410,0.843952}%
\pgfsetfillcolor{currentfill}%
\pgfsetlinewidth{0.000000pt}%
\definecolor{currentstroke}{rgb}{0.000000,0.000000,0.000000}%
\pgfsetstrokecolor{currentstroke}%
\pgfsetstrokeopacity{0.000000}%
\pgfsetdash{}{0pt}%
\pgfpathmoveto{\pgfqpoint{1.333705in}{1.220024in}}%
\pgfpathlineto{\pgfqpoint{1.311373in}{1.222190in}}%
\pgfpathlineto{\pgfqpoint{1.265951in}{1.226957in}}%
\pgfpathlineto{\pgfqpoint{1.267871in}{1.244110in}}%
\pgfpathlineto{\pgfqpoint{1.237983in}{1.247422in}}%
\pgfpathlineto{\pgfqpoint{1.233045in}{1.254595in}}%
\pgfpathlineto{\pgfqpoint{1.229691in}{1.263388in}}%
\pgfpathlineto{\pgfqpoint{1.222250in}{1.286251in}}%
\pgfpathlineto{\pgfqpoint{1.216879in}{1.294670in}}%
\pgfpathlineto{\pgfqpoint{1.217135in}{1.298862in}}%
\pgfpathlineto{\pgfqpoint{1.262254in}{1.293547in}}%
\pgfpathlineto{\pgfqpoint{1.294492in}{1.290219in}}%
\pgfpathlineto{\pgfqpoint{1.346503in}{1.285013in}}%
\pgfpathlineto{\pgfqpoint{1.343169in}{1.253667in}}%
\pgfpathlineto{\pgfqpoint{1.348876in}{1.253120in}}%
\pgfpathlineto{\pgfqpoint{1.345447in}{1.218910in}}%
\pgfpathlineto{\pgfqpoint{1.333705in}{1.220024in}}%
\pgfpathclose%
\pgfusepath{fill}%
\end{pgfscope}%
\begin{pgfscope}%
\pgfpathrectangle{\pgfqpoint{0.100000in}{0.100000in}}{\pgfqpoint{3.007045in}{1.925000in}}%
\pgfusepath{clip}%
\pgfsetbuttcap%
\pgfsetmiterjoin%
\definecolor{currentfill}{rgb}{0.491765,0.721968,0.854779}%
\pgfsetfillcolor{currentfill}%
\pgfsetlinewidth{0.000000pt}%
\definecolor{currentstroke}{rgb}{0.000000,0.000000,0.000000}%
\pgfsetstrokecolor{currentstroke}%
\pgfsetstrokeopacity{0.000000}%
\pgfsetdash{}{0pt}%
\pgfpathmoveto{\pgfqpoint{0.741011in}{1.724091in}}%
\pgfpathlineto{\pgfqpoint{0.746177in}{1.746042in}}%
\pgfpathlineto{\pgfqpoint{0.743350in}{1.746702in}}%
\pgfpathlineto{\pgfqpoint{0.742787in}{1.756014in}}%
\pgfpathlineto{\pgfqpoint{0.737197in}{1.757418in}}%
\pgfpathlineto{\pgfqpoint{0.740246in}{1.767188in}}%
\pgfpathlineto{\pgfqpoint{0.748808in}{1.770376in}}%
\pgfpathlineto{\pgfqpoint{0.752113in}{1.767443in}}%
\pgfpathlineto{\pgfqpoint{0.758923in}{1.767006in}}%
\pgfpathlineto{\pgfqpoint{0.760765in}{1.761473in}}%
\pgfpathlineto{\pgfqpoint{0.765426in}{1.755324in}}%
\pgfpathlineto{\pgfqpoint{0.765843in}{1.746677in}}%
\pgfpathlineto{\pgfqpoint{0.773003in}{1.745268in}}%
\pgfpathlineto{\pgfqpoint{0.774865in}{1.752771in}}%
\pgfpathlineto{\pgfqpoint{0.789165in}{1.749329in}}%
\pgfpathlineto{\pgfqpoint{0.794807in}{1.753888in}}%
\pgfpathlineto{\pgfqpoint{0.807745in}{1.750662in}}%
\pgfpathlineto{\pgfqpoint{0.812473in}{1.770492in}}%
\pgfpathlineto{\pgfqpoint{0.844253in}{1.762934in}}%
\pgfpathlineto{\pgfqpoint{0.873084in}{1.756371in}}%
\pgfpathlineto{\pgfqpoint{0.873285in}{1.749159in}}%
\pgfpathlineto{\pgfqpoint{0.884881in}{1.738405in}}%
\pgfpathlineto{\pgfqpoint{0.884196in}{1.733912in}}%
\pgfpathlineto{\pgfqpoint{0.895843in}{1.733369in}}%
\pgfpathlineto{\pgfqpoint{0.897388in}{1.731726in}}%
\pgfpathlineto{\pgfqpoint{0.891175in}{1.716853in}}%
\pgfpathlineto{\pgfqpoint{0.885687in}{1.708694in}}%
\pgfpathlineto{\pgfqpoint{0.885040in}{1.701682in}}%
\pgfpathlineto{\pgfqpoint{0.881292in}{1.694653in}}%
\pgfpathlineto{\pgfqpoint{0.883784in}{1.683111in}}%
\pgfpathlineto{\pgfqpoint{0.877655in}{1.681396in}}%
\pgfpathlineto{\pgfqpoint{0.874072in}{1.677313in}}%
\pgfpathlineto{\pgfqpoint{0.875599in}{1.669863in}}%
\pgfpathlineto{\pgfqpoint{0.871070in}{1.663297in}}%
\pgfpathlineto{\pgfqpoint{0.865364in}{1.658676in}}%
\pgfpathlineto{\pgfqpoint{0.859318in}{1.660383in}}%
\pgfpathlineto{\pgfqpoint{0.865574in}{1.646714in}}%
\pgfpathlineto{\pgfqpoint{0.860069in}{1.641241in}}%
\pgfpathlineto{\pgfqpoint{0.802276in}{1.654267in}}%
\pgfpathlineto{\pgfqpoint{0.792690in}{1.650559in}}%
\pgfpathlineto{\pgfqpoint{0.786641in}{1.653470in}}%
\pgfpathlineto{\pgfqpoint{0.786144in}{1.663045in}}%
\pgfpathlineto{\pgfqpoint{0.770539in}{1.666797in}}%
\pgfpathlineto{\pgfqpoint{0.780091in}{1.680597in}}%
\pgfpathlineto{\pgfqpoint{0.786067in}{1.686768in}}%
\pgfpathlineto{\pgfqpoint{0.784656in}{1.695189in}}%
\pgfpathlineto{\pgfqpoint{0.778737in}{1.702487in}}%
\pgfpathlineto{\pgfqpoint{0.774980in}{1.703762in}}%
\pgfpathlineto{\pgfqpoint{0.771689in}{1.716185in}}%
\pgfpathlineto{\pgfqpoint{0.741011in}{1.724091in}}%
\pgfpathclose%
\pgfusepath{fill}%
\end{pgfscope}%
\begin{pgfscope}%
\pgfpathrectangle{\pgfqpoint{0.100000in}{0.100000in}}{\pgfqpoint{3.007045in}{1.925000in}}%
\pgfusepath{clip}%
\pgfsetbuttcap%
\pgfsetmiterjoin%
\definecolor{currentfill}{rgb}{0.671895,0.814379,0.900654}%
\pgfsetfillcolor{currentfill}%
\pgfsetlinewidth{0.000000pt}%
\definecolor{currentstroke}{rgb}{0.000000,0.000000,0.000000}%
\pgfsetstrokecolor{currentstroke}%
\pgfsetstrokeopacity{0.000000}%
\pgfsetdash{}{0pt}%
\pgfpathmoveto{\pgfqpoint{2.492884in}{0.596677in}}%
\pgfpathlineto{\pgfqpoint{2.516161in}{0.600413in}}%
\pgfpathlineto{\pgfqpoint{2.520191in}{0.572396in}}%
\pgfpathlineto{\pgfqpoint{2.522149in}{0.556021in}}%
\pgfpathlineto{\pgfqpoint{2.513609in}{0.551975in}}%
\pgfpathlineto{\pgfqpoint{2.513059in}{0.555566in}}%
\pgfpathlineto{\pgfqpoint{2.501841in}{0.554053in}}%
\pgfpathlineto{\pgfqpoint{2.504526in}{0.536259in}}%
\pgfpathlineto{\pgfqpoint{2.496864in}{0.535136in}}%
\pgfpathlineto{\pgfqpoint{2.498437in}{0.523892in}}%
\pgfpathlineto{\pgfqpoint{2.494299in}{0.520921in}}%
\pgfpathlineto{\pgfqpoint{2.488304in}{0.521528in}}%
\pgfpathlineto{\pgfqpoint{2.486048in}{0.518807in}}%
\pgfpathlineto{\pgfqpoint{2.485328in}{0.522899in}}%
\pgfpathlineto{\pgfqpoint{2.480931in}{0.529259in}}%
\pgfpathlineto{\pgfqpoint{2.469105in}{0.528163in}}%
\pgfpathlineto{\pgfqpoint{2.464882in}{0.534167in}}%
\pgfpathlineto{\pgfqpoint{2.458365in}{0.538964in}}%
\pgfpathlineto{\pgfqpoint{2.455703in}{0.543435in}}%
\pgfpathlineto{\pgfqpoint{2.450937in}{0.544451in}}%
\pgfpathlineto{\pgfqpoint{2.444679in}{0.548491in}}%
\pgfpathlineto{\pgfqpoint{2.442205in}{0.558231in}}%
\pgfpathlineto{\pgfqpoint{2.444618in}{0.558566in}}%
\pgfpathlineto{\pgfqpoint{2.446412in}{0.569336in}}%
\pgfpathlineto{\pgfqpoint{2.470809in}{0.572873in}}%
\pgfpathlineto{\pgfqpoint{2.471783in}{0.578648in}}%
\pgfpathlineto{\pgfqpoint{2.476107in}{0.580682in}}%
\pgfpathlineto{\pgfqpoint{2.479736in}{0.576043in}}%
\pgfpathlineto{\pgfqpoint{2.485231in}{0.574838in}}%
\pgfpathlineto{\pgfqpoint{2.487697in}{0.577489in}}%
\pgfpathlineto{\pgfqpoint{2.488025in}{0.585155in}}%
\pgfpathlineto{\pgfqpoint{2.492884in}{0.596677in}}%
\pgfpathclose%
\pgfusepath{fill}%
\end{pgfscope}%
\begin{pgfscope}%
\pgfpathrectangle{\pgfqpoint{0.100000in}{0.100000in}}{\pgfqpoint{3.007045in}{1.925000in}}%
\pgfusepath{clip}%
\pgfsetbuttcap%
\pgfsetmiterjoin%
\definecolor{currentfill}{rgb}{0.346467,0.632403,0.810673}%
\pgfsetfillcolor{currentfill}%
\pgfsetlinewidth{0.000000pt}%
\definecolor{currentstroke}{rgb}{0.000000,0.000000,0.000000}%
\pgfsetstrokecolor{currentstroke}%
\pgfsetstrokeopacity{0.000000}%
\pgfsetdash{}{0pt}%
\pgfpathmoveto{\pgfqpoint{2.364743in}{1.174580in}}%
\pgfpathlineto{\pgfqpoint{2.370900in}{1.172765in}}%
\pgfpathlineto{\pgfqpoint{2.374431in}{1.146809in}}%
\pgfpathlineto{\pgfqpoint{2.360800in}{1.150692in}}%
\pgfpathlineto{\pgfqpoint{2.356171in}{1.148869in}}%
\pgfpathlineto{\pgfqpoint{2.354011in}{1.144847in}}%
\pgfpathlineto{\pgfqpoint{2.348520in}{1.145975in}}%
\pgfpathlineto{\pgfqpoint{2.342358in}{1.152328in}}%
\pgfpathlineto{\pgfqpoint{2.336776in}{1.153283in}}%
\pgfpathlineto{\pgfqpoint{2.331133in}{1.151744in}}%
\pgfpathlineto{\pgfqpoint{2.323384in}{1.154002in}}%
\pgfpathlineto{\pgfqpoint{2.326824in}{1.138484in}}%
\pgfpathlineto{\pgfqpoint{2.317016in}{1.135906in}}%
\pgfpathlineto{\pgfqpoint{2.308461in}{1.129091in}}%
\pgfpathlineto{\pgfqpoint{2.297453in}{1.137486in}}%
\pgfpathlineto{\pgfqpoint{2.296120in}{1.144067in}}%
\pgfpathlineto{\pgfqpoint{2.288955in}{1.139422in}}%
\pgfpathlineto{\pgfqpoint{2.283647in}{1.145438in}}%
\pgfpathlineto{\pgfqpoint{2.290373in}{1.148367in}}%
\pgfpathlineto{\pgfqpoint{2.294975in}{1.154537in}}%
\pgfpathlineto{\pgfqpoint{2.289964in}{1.156159in}}%
\pgfpathlineto{\pgfqpoint{2.291496in}{1.161362in}}%
\pgfpathlineto{\pgfqpoint{2.288006in}{1.165183in}}%
\pgfpathlineto{\pgfqpoint{2.291530in}{1.169095in}}%
\pgfpathlineto{\pgfqpoint{2.288706in}{1.196382in}}%
\pgfpathlineto{\pgfqpoint{2.288393in}{1.199350in}}%
\pgfpathlineto{\pgfqpoint{2.305334in}{1.201390in}}%
\pgfpathlineto{\pgfqpoint{2.314937in}{1.203868in}}%
\pgfpathlineto{\pgfqpoint{2.330635in}{1.204428in}}%
\pgfpathlineto{\pgfqpoint{2.331627in}{1.183748in}}%
\pgfpathlineto{\pgfqpoint{2.338836in}{1.184076in}}%
\pgfpathlineto{\pgfqpoint{2.340324in}{1.169258in}}%
\pgfpathlineto{\pgfqpoint{2.353578in}{1.170749in}}%
\pgfpathlineto{\pgfqpoint{2.364743in}{1.174580in}}%
\pgfpathclose%
\pgfusepath{fill}%
\end{pgfscope}%
\begin{pgfscope}%
\pgfpathrectangle{\pgfqpoint{0.100000in}{0.100000in}}{\pgfqpoint{3.007045in}{1.925000in}}%
\pgfusepath{clip}%
\pgfsetbuttcap%
\pgfsetmiterjoin%
\definecolor{currentfill}{rgb}{0.422745,0.684075,0.839892}%
\pgfsetfillcolor{currentfill}%
\pgfsetlinewidth{0.000000pt}%
\definecolor{currentstroke}{rgb}{0.000000,0.000000,0.000000}%
\pgfsetstrokecolor{currentstroke}%
\pgfsetstrokeopacity{0.000000}%
\pgfsetdash{}{0pt}%
\pgfpathmoveto{\pgfqpoint{2.439065in}{1.211016in}}%
\pgfpathlineto{\pgfqpoint{2.428341in}{1.210131in}}%
\pgfpathlineto{\pgfqpoint{2.427935in}{1.216104in}}%
\pgfpathlineto{\pgfqpoint{2.422257in}{1.215672in}}%
\pgfpathlineto{\pgfqpoint{2.410789in}{1.216745in}}%
\pgfpathlineto{\pgfqpoint{2.410154in}{1.226352in}}%
\pgfpathlineto{\pgfqpoint{2.407820in}{1.232038in}}%
\pgfpathlineto{\pgfqpoint{2.403956in}{1.231787in}}%
\pgfpathlineto{\pgfqpoint{2.403598in}{1.237933in}}%
\pgfpathlineto{\pgfqpoint{2.415203in}{1.238420in}}%
\pgfpathlineto{\pgfqpoint{2.416625in}{1.241076in}}%
\pgfpathlineto{\pgfqpoint{2.415184in}{1.255293in}}%
\pgfpathlineto{\pgfqpoint{2.441070in}{1.257982in}}%
\pgfpathlineto{\pgfqpoint{2.440547in}{1.262665in}}%
\pgfpathlineto{\pgfqpoint{2.457021in}{1.264583in}}%
\pgfpathlineto{\pgfqpoint{2.463037in}{1.262554in}}%
\pgfpathlineto{\pgfqpoint{2.465474in}{1.242785in}}%
\pgfpathlineto{\pgfqpoint{2.461646in}{1.242312in}}%
\pgfpathlineto{\pgfqpoint{2.462970in}{1.231712in}}%
\pgfpathlineto{\pgfqpoint{2.465782in}{1.225372in}}%
\pgfpathlineto{\pgfqpoint{2.456449in}{1.222465in}}%
\pgfpathlineto{\pgfqpoint{2.445563in}{1.221325in}}%
\pgfpathlineto{\pgfqpoint{2.445630in}{1.214029in}}%
\pgfpathlineto{\pgfqpoint{2.439738in}{1.214050in}}%
\pgfpathlineto{\pgfqpoint{2.439065in}{1.211016in}}%
\pgfpathclose%
\pgfusepath{fill}%
\end{pgfscope}%
\begin{pgfscope}%
\pgfpathrectangle{\pgfqpoint{0.100000in}{0.100000in}}{\pgfqpoint{3.007045in}{1.925000in}}%
\pgfusepath{clip}%
\pgfsetbuttcap%
\pgfsetmiterjoin%
\definecolor{currentfill}{rgb}{0.366644,0.646182,0.818547}%
\pgfsetfillcolor{currentfill}%
\pgfsetlinewidth{0.000000pt}%
\definecolor{currentstroke}{rgb}{0.000000,0.000000,0.000000}%
\pgfsetstrokecolor{currentstroke}%
\pgfsetstrokeopacity{0.000000}%
\pgfsetdash{}{0pt}%
\pgfpathmoveto{\pgfqpoint{1.704306in}{1.528001in}}%
\pgfpathlineto{\pgfqpoint{1.684152in}{1.528327in}}%
\pgfpathlineto{\pgfqpoint{1.684336in}{1.539811in}}%
\pgfpathlineto{\pgfqpoint{1.668319in}{1.540277in}}%
\pgfpathlineto{\pgfqpoint{1.668554in}{1.549944in}}%
\pgfpathlineto{\pgfqpoint{1.678758in}{1.549649in}}%
\pgfpathlineto{\pgfqpoint{1.679482in}{1.551623in}}%
\pgfpathlineto{\pgfqpoint{1.703819in}{1.551137in}}%
\pgfpathlineto{\pgfqpoint{1.701484in}{1.554430in}}%
\pgfpathlineto{\pgfqpoint{1.694140in}{1.556870in}}%
\pgfpathlineto{\pgfqpoint{1.686157in}{1.569986in}}%
\pgfpathlineto{\pgfqpoint{1.687099in}{1.572723in}}%
\pgfpathlineto{\pgfqpoint{1.695859in}{1.579141in}}%
\pgfpathlineto{\pgfqpoint{1.699144in}{1.583940in}}%
\pgfpathlineto{\pgfqpoint{1.700191in}{1.591568in}}%
\pgfpathlineto{\pgfqpoint{1.699675in}{1.597325in}}%
\pgfpathlineto{\pgfqpoint{1.714002in}{1.597084in}}%
\pgfpathlineto{\pgfqpoint{1.714476in}{1.591318in}}%
\pgfpathlineto{\pgfqpoint{1.714182in}{1.568146in}}%
\pgfpathlineto{\pgfqpoint{1.714573in}{1.556685in}}%
\pgfpathlineto{\pgfqpoint{1.720360in}{1.556564in}}%
\pgfpathlineto{\pgfqpoint{1.720167in}{1.545159in}}%
\pgfpathlineto{\pgfqpoint{1.717578in}{1.542068in}}%
\pgfpathlineto{\pgfqpoint{1.712446in}{1.545741in}}%
\pgfpathlineto{\pgfqpoint{1.704575in}{1.547315in}}%
\pgfpathlineto{\pgfqpoint{1.704306in}{1.528001in}}%
\pgfpathclose%
\pgfusepath{fill}%
\end{pgfscope}%
\begin{pgfscope}%
\pgfpathrectangle{\pgfqpoint{0.100000in}{0.100000in}}{\pgfqpoint{3.007045in}{1.925000in}}%
\pgfusepath{clip}%
\pgfsetbuttcap%
\pgfsetmiterjoin%
\definecolor{currentfill}{rgb}{0.175087,0.488812,0.733333}%
\pgfsetfillcolor{currentfill}%
\pgfsetlinewidth{0.000000pt}%
\definecolor{currentstroke}{rgb}{0.000000,0.000000,0.000000}%
\pgfsetstrokecolor{currentstroke}%
\pgfsetstrokeopacity{0.000000}%
\pgfsetdash{}{0pt}%
\pgfpathmoveto{\pgfqpoint{2.234453in}{1.308012in}}%
\pgfpathlineto{\pgfqpoint{2.233410in}{1.317320in}}%
\pgfpathlineto{\pgfqpoint{2.213486in}{1.315239in}}%
\pgfpathlineto{\pgfqpoint{2.211201in}{1.336530in}}%
\pgfpathlineto{\pgfqpoint{2.224488in}{1.337759in}}%
\pgfpathlineto{\pgfqpoint{2.225748in}{1.340539in}}%
\pgfpathlineto{\pgfqpoint{2.223806in}{1.358262in}}%
\pgfpathlineto{\pgfqpoint{2.269373in}{1.363358in}}%
\pgfpathlineto{\pgfqpoint{2.271732in}{1.342879in}}%
\pgfpathlineto{\pgfqpoint{2.253575in}{1.340858in}}%
\pgfpathlineto{\pgfqpoint{2.257315in}{1.308409in}}%
\pgfpathlineto{\pgfqpoint{2.240301in}{1.306723in}}%
\pgfpathlineto{\pgfqpoint{2.234453in}{1.308012in}}%
\pgfpathclose%
\pgfusepath{fill}%
\end{pgfscope}%
\begin{pgfscope}%
\pgfpathrectangle{\pgfqpoint{0.100000in}{0.100000in}}{\pgfqpoint{3.007045in}{1.925000in}}%
\pgfusepath{clip}%
\pgfsetbuttcap%
\pgfsetmiterjoin%
\definecolor{currentfill}{rgb}{0.760477,0.852026,0.931657}%
\pgfsetfillcolor{currentfill}%
\pgfsetlinewidth{0.000000pt}%
\definecolor{currentstroke}{rgb}{0.000000,0.000000,0.000000}%
\pgfsetstrokecolor{currentstroke}%
\pgfsetstrokeopacity{0.000000}%
\pgfsetdash{}{0pt}%
\pgfpathmoveto{\pgfqpoint{1.229691in}{1.263388in}}%
\pgfpathlineto{\pgfqpoint{1.233045in}{1.254595in}}%
\pgfpathlineto{\pgfqpoint{1.237983in}{1.247422in}}%
\pgfpathlineto{\pgfqpoint{1.237002in}{1.232833in}}%
\pgfpathlineto{\pgfqpoint{1.234437in}{1.230208in}}%
\pgfpathlineto{\pgfqpoint{1.232976in}{1.220918in}}%
\pgfpathlineto{\pgfqpoint{1.225797in}{1.217368in}}%
\pgfpathlineto{\pgfqpoint{1.220004in}{1.212317in}}%
\pgfpathlineto{\pgfqpoint{1.214794in}{1.212322in}}%
\pgfpathlineto{\pgfqpoint{1.212843in}{1.220333in}}%
\pgfpathlineto{\pgfqpoint{1.205530in}{1.228275in}}%
\pgfpathlineto{\pgfqpoint{1.186601in}{1.231472in}}%
\pgfpathlineto{\pgfqpoint{1.186728in}{1.236591in}}%
\pgfpathlineto{\pgfqpoint{1.189821in}{1.258740in}}%
\pgfpathlineto{\pgfqpoint{1.189706in}{1.265642in}}%
\pgfpathlineto{\pgfqpoint{1.192494in}{1.261100in}}%
\pgfpathlineto{\pgfqpoint{1.205670in}{1.255974in}}%
\pgfpathlineto{\pgfqpoint{1.213800in}{1.255082in}}%
\pgfpathlineto{\pgfqpoint{1.220648in}{1.256986in}}%
\pgfpathlineto{\pgfqpoint{1.222855in}{1.255180in}}%
\pgfpathlineto{\pgfqpoint{1.229691in}{1.263388in}}%
\pgfpathclose%
\pgfusepath{fill}%
\end{pgfscope}%
\begin{pgfscope}%
\pgfpathrectangle{\pgfqpoint{0.100000in}{0.100000in}}{\pgfqpoint{3.007045in}{1.925000in}}%
\pgfusepath{clip}%
\pgfsetbuttcap%
\pgfsetmiterjoin%
\definecolor{currentfill}{rgb}{0.256286,0.570012,0.775163}%
\pgfsetfillcolor{currentfill}%
\pgfsetlinewidth{0.000000pt}%
\definecolor{currentstroke}{rgb}{0.000000,0.000000,0.000000}%
\pgfsetstrokecolor{currentstroke}%
\pgfsetstrokeopacity{0.000000}%
\pgfsetdash{}{0pt}%
\pgfpathmoveto{\pgfqpoint{2.215737in}{1.157449in}}%
\pgfpathlineto{\pgfqpoint{2.213096in}{1.157168in}}%
\pgfpathlineto{\pgfqpoint{2.210701in}{1.176219in}}%
\pgfpathlineto{\pgfqpoint{2.198043in}{1.175099in}}%
\pgfpathlineto{\pgfqpoint{2.197194in}{1.183592in}}%
\pgfpathlineto{\pgfqpoint{2.194423in}{1.183326in}}%
\pgfpathlineto{\pgfqpoint{2.195596in}{1.188844in}}%
\pgfpathlineto{\pgfqpoint{2.193431in}{1.193828in}}%
\pgfpathlineto{\pgfqpoint{2.189712in}{1.229760in}}%
\pgfpathlineto{\pgfqpoint{2.212424in}{1.232038in}}%
\pgfpathlineto{\pgfqpoint{2.212181in}{1.234336in}}%
\pgfpathlineto{\pgfqpoint{2.231222in}{1.236425in}}%
\pgfpathlineto{\pgfqpoint{2.233016in}{1.218375in}}%
\pgfpathlineto{\pgfqpoint{2.247375in}{1.219993in}}%
\pgfpathlineto{\pgfqpoint{2.247429in}{1.209448in}}%
\pgfpathlineto{\pgfqpoint{2.262289in}{1.211080in}}%
\pgfpathlineto{\pgfqpoint{2.264298in}{1.193927in}}%
\pgfpathlineto{\pgfqpoint{2.264834in}{1.189201in}}%
\pgfpathlineto{\pgfqpoint{2.248039in}{1.187423in}}%
\pgfpathlineto{\pgfqpoint{2.245949in}{1.180376in}}%
\pgfpathlineto{\pgfqpoint{2.225671in}{1.177961in}}%
\pgfpathlineto{\pgfqpoint{2.227838in}{1.158721in}}%
\pgfpathlineto{\pgfqpoint{2.215737in}{1.157449in}}%
\pgfpathclose%
\pgfusepath{fill}%
\end{pgfscope}%
\begin{pgfscope}%
\pgfpathrectangle{\pgfqpoint{0.100000in}{0.100000in}}{\pgfqpoint{3.007045in}{1.925000in}}%
\pgfusepath{clip}%
\pgfsetbuttcap%
\pgfsetmiterjoin%
\definecolor{currentfill}{rgb}{0.154787,0.468512,0.722876}%
\pgfsetfillcolor{currentfill}%
\pgfsetlinewidth{0.000000pt}%
\definecolor{currentstroke}{rgb}{0.000000,0.000000,0.000000}%
\pgfsetstrokecolor{currentstroke}%
\pgfsetstrokeopacity{0.000000}%
\pgfsetdash{}{0pt}%
\pgfpathmoveto{\pgfqpoint{0.765905in}{0.863949in}}%
\pgfpathlineto{\pgfqpoint{0.754327in}{0.807329in}}%
\pgfpathlineto{\pgfqpoint{0.748143in}{0.777058in}}%
\pgfpathlineto{\pgfqpoint{0.673269in}{0.823873in}}%
\pgfpathlineto{\pgfqpoint{0.675440in}{0.831946in}}%
\pgfpathlineto{\pgfqpoint{0.681603in}{0.837269in}}%
\pgfpathlineto{\pgfqpoint{0.604721in}{0.848024in}}%
\pgfpathlineto{\pgfqpoint{0.611833in}{0.877391in}}%
\pgfpathlineto{\pgfqpoint{0.613032in}{0.877119in}}%
\pgfpathlineto{\pgfqpoint{0.618157in}{0.899792in}}%
\pgfpathlineto{\pgfqpoint{0.673193in}{0.887262in}}%
\pgfpathlineto{\pgfqpoint{0.691055in}{0.881658in}}%
\pgfpathlineto{\pgfqpoint{0.691629in}{0.872688in}}%
\pgfpathlineto{\pgfqpoint{0.687845in}{0.862115in}}%
\pgfpathlineto{\pgfqpoint{0.689447in}{0.856914in}}%
\pgfpathlineto{\pgfqpoint{0.710499in}{0.852063in}}%
\pgfpathlineto{\pgfqpoint{0.716592in}{0.880125in}}%
\pgfpathlineto{\pgfqpoint{0.733331in}{0.876625in}}%
\pgfpathlineto{\pgfqpoint{0.732137in}{0.870987in}}%
\pgfpathlineto{\pgfqpoint{0.765905in}{0.863949in}}%
\pgfpathclose%
\pgfusepath{fill}%
\end{pgfscope}%
\begin{pgfscope}%
\pgfpathrectangle{\pgfqpoint{0.100000in}{0.100000in}}{\pgfqpoint{3.007045in}{1.925000in}}%
\pgfusepath{clip}%
\pgfsetbuttcap%
\pgfsetroundjoin%
\pgfsetlinewidth{0.050187pt}%
\definecolor{currentstroke}{rgb}{1.000000,1.000000,1.000000}%
\pgfsetstrokecolor{currentstroke}%
\pgfsetdash{}{0pt}%
\pgfusepath{stroke}%
\end{pgfscope}%
\begin{pgfscope}%
\pgfpathrectangle{\pgfqpoint{0.100000in}{0.100000in}}{\pgfqpoint{3.007045in}{1.925000in}}%
\pgfusepath{clip}%
\pgfsetbuttcap%
\pgfsetroundjoin%
\pgfsetlinewidth{0.050187pt}%
\definecolor{currentstroke}{rgb}{1.000000,1.000000,1.000000}%
\pgfsetstrokecolor{currentstroke}%
\pgfsetdash{}{0pt}%
\pgfusepath{stroke}%
\end{pgfscope}%
\begin{pgfscope}%
\pgfpathrectangle{\pgfqpoint{0.100000in}{0.100000in}}{\pgfqpoint{3.007045in}{1.925000in}}%
\pgfusepath{clip}%
\pgfsetbuttcap%
\pgfsetroundjoin%
\pgfsetlinewidth{0.050187pt}%
\definecolor{currentstroke}{rgb}{1.000000,1.000000,1.000000}%
\pgfsetstrokecolor{currentstroke}%
\pgfsetdash{}{0pt}%
\pgfusepath{stroke}%
\end{pgfscope}%
\begin{pgfscope}%
\pgfpathrectangle{\pgfqpoint{0.100000in}{0.100000in}}{\pgfqpoint{3.007045in}{1.925000in}}%
\pgfusepath{clip}%
\pgfsetbuttcap%
\pgfsetroundjoin%
\pgfsetlinewidth{0.050187pt}%
\definecolor{currentstroke}{rgb}{1.000000,1.000000,1.000000}%
\pgfsetstrokecolor{currentstroke}%
\pgfsetdash{}{0pt}%
\pgfusepath{stroke}%
\end{pgfscope}%
\begin{pgfscope}%
\pgfpathrectangle{\pgfqpoint{0.100000in}{0.100000in}}{\pgfqpoint{3.007045in}{1.925000in}}%
\pgfusepath{clip}%
\pgfsetbuttcap%
\pgfsetroundjoin%
\pgfsetlinewidth{0.050187pt}%
\definecolor{currentstroke}{rgb}{1.000000,1.000000,1.000000}%
\pgfsetstrokecolor{currentstroke}%
\pgfsetdash{}{0pt}%
\pgfusepath{stroke}%
\end{pgfscope}%
\begin{pgfscope}%
\pgfpathrectangle{\pgfqpoint{0.100000in}{0.100000in}}{\pgfqpoint{3.007045in}{1.925000in}}%
\pgfusepath{clip}%
\pgfsetbuttcap%
\pgfsetroundjoin%
\pgfsetlinewidth{0.050187pt}%
\definecolor{currentstroke}{rgb}{1.000000,1.000000,1.000000}%
\pgfsetstrokecolor{currentstroke}%
\pgfsetdash{}{0pt}%
\pgfusepath{stroke}%
\end{pgfscope}%
\begin{pgfscope}%
\pgfpathrectangle{\pgfqpoint{0.100000in}{0.100000in}}{\pgfqpoint{3.007045in}{1.925000in}}%
\pgfusepath{clip}%
\pgfsetbuttcap%
\pgfsetroundjoin%
\pgfsetlinewidth{0.050187pt}%
\definecolor{currentstroke}{rgb}{1.000000,1.000000,1.000000}%
\pgfsetstrokecolor{currentstroke}%
\pgfsetdash{}{0pt}%
\pgfpathmoveto{\pgfqpoint{0.814503in}{1.911787in}}%
\pgfpathlineto{\pgfqpoint{0.798909in}{1.849997in}}%
\pgfpathlineto{\pgfqpoint{0.787821in}{1.805724in}}%
\pgfpathlineto{\pgfqpoint{0.774637in}{1.752465in}}%
\pgfpathlineto{\pgfqpoint{0.772049in}{1.739338in}}%
\pgfpathlineto{\pgfqpoint{0.773739in}{1.727298in}}%
\pgfpathlineto{\pgfqpoint{0.771499in}{1.716207in}}%
\pgfpathlineto{\pgfqpoint{0.679470in}{1.740222in}}%
\pgfpathlineto{\pgfqpoint{0.671157in}{1.737433in}}%
\pgfpathlineto{\pgfqpoint{0.664042in}{1.739853in}}%
\pgfpathlineto{\pgfqpoint{0.631104in}{1.740645in}}%
\pgfpathlineto{\pgfqpoint{0.619963in}{1.737504in}}%
\pgfpathlineto{\pgfqpoint{0.608900in}{1.738565in}}%
\pgfpathlineto{\pgfqpoint{0.604211in}{1.743481in}}%
\pgfpathlineto{\pgfqpoint{0.574722in}{1.742409in}}%
\pgfpathlineto{\pgfqpoint{0.570064in}{1.750443in}}%
\pgfpathlineto{\pgfqpoint{0.561223in}{1.754676in}}%
\pgfpathlineto{\pgfqpoint{0.548360in}{1.757255in}}%
\pgfpathlineto{\pgfqpoint{0.526157in}{1.753479in}}%
\pgfpathlineto{\pgfqpoint{0.517960in}{1.757179in}}%
\pgfpathlineto{\pgfqpoint{0.505323in}{1.767019in}}%
\pgfpathlineto{\pgfqpoint{0.509122in}{1.786306in}}%
\pgfpathlineto{\pgfqpoint{0.507790in}{1.796155in}}%
\pgfpathlineto{\pgfqpoint{0.497775in}{1.806612in}}%
\pgfpathlineto{\pgfqpoint{0.491359in}{1.805963in}}%
\pgfpathlineto{\pgfqpoint{0.486757in}{1.816540in}}%
\pgfpathlineto{\pgfqpoint{0.475856in}{1.820815in}}%
\pgfpathlineto{\pgfqpoint{0.467931in}{1.820270in}}%
\pgfpathlineto{\pgfqpoint{0.465368in}{1.831133in}}%
\pgfpathlineto{\pgfqpoint{0.473262in}{1.829977in}}%
\pgfpathlineto{\pgfqpoint{0.472635in}{1.845075in}}%
\pgfpathlineto{\pgfqpoint{0.469194in}{1.854040in}}%
\pgfpathlineto{\pgfqpoint{0.470711in}{1.865514in}}%
\pgfpathlineto{\pgfqpoint{0.474841in}{1.874144in}}%
\pgfpathlineto{\pgfqpoint{0.474608in}{1.890521in}}%
\pgfpathlineto{\pgfqpoint{0.472404in}{1.896463in}}%
\pgfpathlineto{\pgfqpoint{0.476212in}{1.915526in}}%
\pgfpathlineto{\pgfqpoint{0.475112in}{1.927812in}}%
\pgfpathlineto{\pgfqpoint{0.471305in}{1.933703in}}%
\pgfpathlineto{\pgfqpoint{0.471886in}{1.952963in}}%
\pgfpathlineto{\pgfqpoint{0.477424in}{1.967128in}}%
\pgfpathlineto{\pgfqpoint{0.483447in}{1.963661in}}%
\pgfpathlineto{\pgfqpoint{0.503389in}{1.943061in}}%
\pgfpathlineto{\pgfqpoint{0.527555in}{1.931787in}}%
\pgfpathlineto{\pgfqpoint{0.539928in}{1.930453in}}%
\pgfpathlineto{\pgfqpoint{0.551585in}{1.925654in}}%
\pgfpathlineto{\pgfqpoint{0.554894in}{1.909503in}}%
\pgfpathlineto{\pgfqpoint{0.544678in}{1.906443in}}%
\pgfpathlineto{\pgfqpoint{0.537908in}{1.893715in}}%
\pgfpathlineto{\pgfqpoint{0.545857in}{1.894422in}}%
\pgfpathlineto{\pgfqpoint{0.549062in}{1.900135in}}%
\pgfpathlineto{\pgfqpoint{0.560272in}{1.907280in}}%
\pgfpathlineto{\pgfqpoint{0.559693in}{1.896670in}}%
\pgfpathlineto{\pgfqpoint{0.552208in}{1.894959in}}%
\pgfpathlineto{\pgfqpoint{0.553451in}{1.881322in}}%
\pgfpathlineto{\pgfqpoint{0.544337in}{1.867661in}}%
\pgfpathlineto{\pgfqpoint{0.538503in}{1.874157in}}%
\pgfpathlineto{\pgfqpoint{0.521168in}{1.870625in}}%
\pgfpathlineto{\pgfqpoint{0.520256in}{1.862677in}}%
\pgfpathlineto{\pgfqpoint{0.526232in}{1.857836in}}%
\pgfpathlineto{\pgfqpoint{0.535349in}{1.857360in}}%
\pgfpathlineto{\pgfqpoint{0.547955in}{1.867706in}}%
\pgfpathlineto{\pgfqpoint{0.557932in}{1.868565in}}%
\pgfpathlineto{\pgfqpoint{0.558315in}{1.880017in}}%
\pgfpathlineto{\pgfqpoint{0.563457in}{1.896838in}}%
\pgfpathlineto{\pgfqpoint{0.575643in}{1.909348in}}%
\pgfpathlineto{\pgfqpoint{0.571580in}{1.919037in}}%
\pgfpathlineto{\pgfqpoint{0.574355in}{1.929469in}}%
\pgfpathlineto{\pgfqpoint{0.568949in}{1.939866in}}%
\pgfpathlineto{\pgfqpoint{0.576654in}{1.953161in}}%
\pgfpathlineto{\pgfqpoint{0.577808in}{1.961160in}}%
\pgfpathlineto{\pgfqpoint{0.571119in}{1.966417in}}%
\pgfpathlineto{\pgfqpoint{0.572232in}{1.979864in}}%
\pgfpathlineto{\pgfqpoint{0.652382in}{1.955513in}}%
\pgfpathlineto{\pgfqpoint{0.737492in}{1.931646in}}%
\pgfpathlineto{\pgfqpoint{0.814503in}{1.911787in}}%
\pgfusepath{stroke}%
\end{pgfscope}%
\begin{pgfscope}%
\pgfpathrectangle{\pgfqpoint{0.100000in}{0.100000in}}{\pgfqpoint{3.007045in}{1.925000in}}%
\pgfusepath{clip}%
\pgfsetbuttcap%
\pgfsetroundjoin%
\pgfsetlinewidth{0.050187pt}%
\definecolor{currentstroke}{rgb}{1.000000,1.000000,1.000000}%
\pgfsetstrokecolor{currentstroke}%
\pgfsetdash{}{0pt}%
\pgfpathmoveto{\pgfqpoint{0.559048in}{1.933344in}}%
\pgfpathlineto{\pgfqpoint{0.562994in}{1.918128in}}%
\pgfpathlineto{\pgfqpoint{0.569183in}{1.913074in}}%
\pgfpathlineto{\pgfqpoint{0.564245in}{1.906641in}}%
\pgfpathlineto{\pgfqpoint{0.559432in}{1.916068in}}%
\pgfpathlineto{\pgfqpoint{0.559048in}{1.933344in}}%
\pgfusepath{stroke}%
\end{pgfscope}%
\begin{pgfscope}%
\pgfpathrectangle{\pgfqpoint{0.100000in}{0.100000in}}{\pgfqpoint{3.007045in}{1.925000in}}%
\pgfusepath{clip}%
\pgfsetbuttcap%
\pgfsetroundjoin%
\pgfsetlinewidth{0.050187pt}%
\definecolor{currentstroke}{rgb}{1.000000,1.000000,1.000000}%
\pgfsetstrokecolor{currentstroke}%
\pgfsetdash{}{0pt}%
\pgfpathmoveto{\pgfqpoint{0.855834in}{1.901800in}}%
\pgfpathlineto{\pgfqpoint{0.941509in}{1.882574in}}%
\pgfpathlineto{\pgfqpoint{1.022196in}{1.866264in}}%
\pgfpathlineto{\pgfqpoint{1.084275in}{1.854890in}}%
\pgfpathlineto{\pgfqpoint{1.138397in}{1.845797in}}%
\pgfpathlineto{\pgfqpoint{1.192639in}{1.837451in}}%
\pgfpathlineto{\pgfqpoint{1.238831in}{1.830944in}}%
\pgfpathlineto{\pgfqpoint{1.285096in}{1.824977in}}%
\pgfpathlineto{\pgfqpoint{1.331427in}{1.819551in}}%
\pgfpathlineto{\pgfqpoint{1.375088in}{1.814940in}}%
\pgfpathlineto{\pgfqpoint{1.369083in}{1.748399in}}%
\pgfpathlineto{\pgfqpoint{1.360146in}{1.658478in}}%
\pgfpathlineto{\pgfqpoint{1.355458in}{1.612230in}}%
\pgfpathlineto{\pgfqpoint{1.348717in}{1.549903in}}%
\pgfpathlineto{\pgfqpoint{1.300990in}{1.555114in}}%
\pgfpathlineto{\pgfqpoint{1.246348in}{1.561317in}}%
\pgfpathlineto{\pgfqpoint{1.170478in}{1.571700in}}%
\pgfpathlineto{\pgfqpoint{1.136592in}{1.576534in}}%
\pgfpathlineto{\pgfqpoint{1.053106in}{1.589639in}}%
\pgfpathlineto{\pgfqpoint{1.024367in}{1.594897in}}%
\pgfpathlineto{\pgfqpoint{1.018377in}{1.560840in}}%
\pgfpathlineto{\pgfqpoint{1.015103in}{1.563270in}}%
\pgfpathlineto{\pgfqpoint{1.008949in}{1.579653in}}%
\pgfpathlineto{\pgfqpoint{1.001038in}{1.578644in}}%
\pgfpathlineto{\pgfqpoint{0.995375in}{1.570362in}}%
\pgfpathlineto{\pgfqpoint{0.983625in}{1.569254in}}%
\pgfpathlineto{\pgfqpoint{0.981264in}{1.573217in}}%
\pgfpathlineto{\pgfqpoint{0.970258in}{1.572737in}}%
\pgfpathlineto{\pgfqpoint{0.964724in}{1.576574in}}%
\pgfpathlineto{\pgfqpoint{0.956987in}{1.570638in}}%
\pgfpathlineto{\pgfqpoint{0.938194in}{1.575980in}}%
\pgfpathlineto{\pgfqpoint{0.929993in}{1.573111in}}%
\pgfpathlineto{\pgfqpoint{0.925581in}{1.585841in}}%
\pgfpathlineto{\pgfqpoint{0.925374in}{1.597953in}}%
\pgfpathlineto{\pgfqpoint{0.912420in}{1.606664in}}%
\pgfpathlineto{\pgfqpoint{0.914532in}{1.619218in}}%
\pgfpathlineto{\pgfqpoint{0.905308in}{1.639956in}}%
\pgfpathlineto{\pgfqpoint{0.905930in}{1.657584in}}%
\pgfpathlineto{\pgfqpoint{0.898002in}{1.666211in}}%
\pgfpathlineto{\pgfqpoint{0.890564in}{1.658578in}}%
\pgfpathlineto{\pgfqpoint{0.880439in}{1.654450in}}%
\pgfpathlineto{\pgfqpoint{0.872732in}{1.662959in}}%
\pgfpathlineto{\pgfqpoint{0.873985in}{1.676547in}}%
\pgfpathlineto{\pgfqpoint{0.883228in}{1.681901in}}%
\pgfpathlineto{\pgfqpoint{0.880771in}{1.690510in}}%
\pgfpathlineto{\pgfqpoint{0.896838in}{1.732179in}}%
\pgfpathlineto{\pgfqpoint{0.884169in}{1.733223in}}%
\pgfpathlineto{\pgfqpoint{0.882867in}{1.740700in}}%
\pgfpathlineto{\pgfqpoint{0.873638in}{1.747149in}}%
\pgfpathlineto{\pgfqpoint{0.874228in}{1.754357in}}%
\pgfpathlineto{\pgfqpoint{0.869360in}{1.759957in}}%
\pgfpathlineto{\pgfqpoint{0.861436in}{1.780289in}}%
\pgfpathlineto{\pgfqpoint{0.854194in}{1.784436in}}%
\pgfpathlineto{\pgfqpoint{0.849009in}{1.801984in}}%
\pgfpathlineto{\pgfqpoint{0.850232in}{1.813647in}}%
\pgfpathlineto{\pgfqpoint{0.840559in}{1.835130in}}%
\pgfpathlineto{\pgfqpoint{0.855834in}{1.901800in}}%
\pgfusepath{stroke}%
\end{pgfscope}%
\begin{pgfscope}%
\pgfpathrectangle{\pgfqpoint{0.100000in}{0.100000in}}{\pgfqpoint{3.007045in}{1.925000in}}%
\pgfusepath{clip}%
\pgfsetbuttcap%
\pgfsetroundjoin%
\pgfsetlinewidth{0.050187pt}%
\definecolor{currentstroke}{rgb}{1.000000,1.000000,1.000000}%
\pgfsetstrokecolor{currentstroke}%
\pgfsetdash{}{0pt}%
\pgfpathmoveto{\pgfqpoint{2.924097in}{1.558905in}}%
\pgfpathlineto{\pgfqpoint{2.922382in}{1.566196in}}%
\pgfpathlineto{\pgfqpoint{2.912852in}{1.572583in}}%
\pgfpathlineto{\pgfqpoint{2.887667in}{1.654962in}}%
\pgfpathlineto{\pgfqpoint{2.874093in}{1.694717in}}%
\pgfpathlineto{\pgfqpoint{2.885006in}{1.706170in}}%
\pgfpathlineto{\pgfqpoint{2.895614in}{1.734199in}}%
\pgfpathlineto{\pgfqpoint{2.900905in}{1.743156in}}%
\pgfpathlineto{\pgfqpoint{2.897157in}{1.746918in}}%
\pgfpathlineto{\pgfqpoint{2.895902in}{1.771799in}}%
\pgfpathlineto{\pgfqpoint{2.900635in}{1.779426in}}%
\pgfpathlineto{\pgfqpoint{2.898560in}{1.797155in}}%
\pgfpathlineto{\pgfqpoint{2.917451in}{1.855321in}}%
\pgfpathlineto{\pgfqpoint{2.925963in}{1.855639in}}%
\pgfpathlineto{\pgfqpoint{2.929497in}{1.845249in}}%
\pgfpathlineto{\pgfqpoint{2.937061in}{1.842266in}}%
\pgfpathlineto{\pgfqpoint{2.951351in}{1.854487in}}%
\pgfpathlineto{\pgfqpoint{2.962555in}{1.861705in}}%
\pgfpathlineto{\pgfqpoint{2.987243in}{1.848994in}}%
\pgfpathlineto{\pgfqpoint{3.009547in}{1.778417in}}%
\pgfpathlineto{\pgfqpoint{3.013786in}{1.761053in}}%
\pgfpathlineto{\pgfqpoint{3.023556in}{1.759017in}}%
\pgfpathlineto{\pgfqpoint{3.036806in}{1.747224in}}%
\pgfpathlineto{\pgfqpoint{3.036056in}{1.740353in}}%
\pgfpathlineto{\pgfqpoint{3.045081in}{1.732208in}}%
\pgfpathlineto{\pgfqpoint{3.052413in}{1.736881in}}%
\pgfpathlineto{\pgfqpoint{3.067756in}{1.718987in}}%
\pgfpathlineto{\pgfqpoint{3.060899in}{1.704668in}}%
\pgfpathlineto{\pgfqpoint{3.051704in}{1.704384in}}%
\pgfpathlineto{\pgfqpoint{3.044326in}{1.691601in}}%
\pgfpathlineto{\pgfqpoint{3.035390in}{1.689788in}}%
\pgfpathlineto{\pgfqpoint{3.028829in}{1.682982in}}%
\pgfpathlineto{\pgfqpoint{3.009234in}{1.675682in}}%
\pgfpathlineto{\pgfqpoint{2.997374in}{1.663918in}}%
\pgfpathlineto{\pgfqpoint{2.991350in}{1.672362in}}%
\pgfpathlineto{\pgfqpoint{2.985875in}{1.666310in}}%
\pgfpathlineto{\pgfqpoint{2.987374in}{1.641845in}}%
\pgfpathlineto{\pgfqpoint{2.983040in}{1.632155in}}%
\pgfpathlineto{\pgfqpoint{2.973585in}{1.634802in}}%
\pgfpathlineto{\pgfqpoint{2.972059in}{1.624866in}}%
\pgfpathlineto{\pgfqpoint{2.967973in}{1.620781in}}%
\pgfpathlineto{\pgfqpoint{2.959797in}{1.621777in}}%
\pgfpathlineto{\pgfqpoint{2.960312in}{1.612488in}}%
\pgfpathlineto{\pgfqpoint{2.947919in}{1.615126in}}%
\pgfpathlineto{\pgfqpoint{2.941151in}{1.602256in}}%
\pgfpathlineto{\pgfqpoint{2.943707in}{1.595522in}}%
\pgfpathlineto{\pgfqpoint{2.939681in}{1.584322in}}%
\pgfpathlineto{\pgfqpoint{2.933343in}{1.576084in}}%
\pgfpathlineto{\pgfqpoint{2.931748in}{1.558887in}}%
\pgfpathlineto{\pgfqpoint{2.924097in}{1.558905in}}%
\pgfusepath{stroke}%
\end{pgfscope}%
\begin{pgfscope}%
\pgfpathrectangle{\pgfqpoint{0.100000in}{0.100000in}}{\pgfqpoint{3.007045in}{1.925000in}}%
\pgfusepath{clip}%
\pgfsetbuttcap%
\pgfsetroundjoin%
\pgfsetlinewidth{0.050187pt}%
\definecolor{currentstroke}{rgb}{1.000000,1.000000,1.000000}%
\pgfsetstrokecolor{currentstroke}%
\pgfsetdash{}{0pt}%
\pgfpathmoveto{\pgfqpoint{3.012754in}{1.670658in}}%
\pgfpathlineto{\pgfqpoint{3.018349in}{1.676519in}}%
\pgfpathlineto{\pgfqpoint{3.023676in}{1.671010in}}%
\pgfpathlineto{\pgfqpoint{3.014108in}{1.663712in}}%
\pgfpathlineto{\pgfqpoint{3.012754in}{1.670658in}}%
\pgfusepath{stroke}%
\end{pgfscope}%
\begin{pgfscope}%
\pgfpathrectangle{\pgfqpoint{0.100000in}{0.100000in}}{\pgfqpoint{3.007045in}{1.925000in}}%
\pgfusepath{clip}%
\pgfsetbuttcap%
\pgfsetroundjoin%
\pgfsetlinewidth{0.050187pt}%
\definecolor{currentstroke}{rgb}{1.000000,1.000000,1.000000}%
\pgfsetstrokecolor{currentstroke}%
\pgfsetdash{}{0pt}%
\pgfpathmoveto{\pgfqpoint{1.355458in}{1.612230in}}%
\pgfpathlineto{\pgfqpoint{1.360146in}{1.658478in}}%
\pgfpathlineto{\pgfqpoint{1.369083in}{1.748399in}}%
\pgfpathlineto{\pgfqpoint{1.375088in}{1.814940in}}%
\pgfpathlineto{\pgfqpoint{1.424263in}{1.810326in}}%
\pgfpathlineto{\pgfqpoint{1.487174in}{1.805318in}}%
\pgfpathlineto{\pgfqpoint{1.544677in}{1.801614in}}%
\pgfpathlineto{\pgfqpoint{1.596745in}{1.798979in}}%
\pgfpathlineto{\pgfqpoint{1.674447in}{1.796314in}}%
\pgfpathlineto{\pgfqpoint{1.679750in}{1.774414in}}%
\pgfpathlineto{\pgfqpoint{1.677852in}{1.763948in}}%
\pgfpathlineto{\pgfqpoint{1.677227in}{1.742437in}}%
\pgfpathlineto{\pgfqpoint{1.680835in}{1.726326in}}%
\pgfpathlineto{\pgfqpoint{1.689100in}{1.702558in}}%
\pgfpathlineto{\pgfqpoint{1.689081in}{1.672635in}}%
\pgfpathlineto{\pgfqpoint{1.690611in}{1.637873in}}%
\pgfpathlineto{\pgfqpoint{1.692711in}{1.628504in}}%
\pgfpathlineto{\pgfqpoint{1.698854in}{1.618219in}}%
\pgfpathlineto{\pgfqpoint{1.700880in}{1.602193in}}%
\pgfpathlineto{\pgfqpoint{1.700005in}{1.591493in}}%
\pgfpathlineto{\pgfqpoint{1.634865in}{1.592905in}}%
\pgfpathlineto{\pgfqpoint{1.587476in}{1.595307in}}%
\pgfpathlineto{\pgfqpoint{1.518008in}{1.599005in}}%
\pgfpathlineto{\pgfqpoint{1.449501in}{1.603871in}}%
\pgfpathlineto{\pgfqpoint{1.403872in}{1.607552in}}%
\pgfpathlineto{\pgfqpoint{1.355458in}{1.612230in}}%
\pgfusepath{stroke}%
\end{pgfscope}%
\begin{pgfscope}%
\pgfpathrectangle{\pgfqpoint{0.100000in}{0.100000in}}{\pgfqpoint{3.007045in}{1.925000in}}%
\pgfusepath{clip}%
\pgfsetbuttcap%
\pgfsetroundjoin%
\pgfsetlinewidth{0.050187pt}%
\definecolor{currentstroke}{rgb}{1.000000,1.000000,1.000000}%
\pgfsetstrokecolor{currentstroke}%
\pgfsetdash{}{0pt}%
\pgfpathmoveto{\pgfqpoint{1.335768in}{1.418681in}}%
\pgfpathlineto{\pgfqpoint{1.343289in}{1.496206in}}%
\pgfpathlineto{\pgfqpoint{1.348717in}{1.549903in}}%
\pgfpathlineto{\pgfqpoint{1.355458in}{1.612230in}}%
\pgfpathlineto{\pgfqpoint{1.403872in}{1.607552in}}%
\pgfpathlineto{\pgfqpoint{1.449501in}{1.603871in}}%
\pgfpathlineto{\pgfqpoint{1.518008in}{1.599005in}}%
\pgfpathlineto{\pgfqpoint{1.587476in}{1.595307in}}%
\pgfpathlineto{\pgfqpoint{1.634865in}{1.592905in}}%
\pgfpathlineto{\pgfqpoint{1.700005in}{1.591493in}}%
\pgfpathlineto{\pgfqpoint{1.695595in}{1.578623in}}%
\pgfpathlineto{\pgfqpoint{1.686792in}{1.568520in}}%
\pgfpathlineto{\pgfqpoint{1.693533in}{1.556888in}}%
\pgfpathlineto{\pgfqpoint{1.700966in}{1.554402in}}%
\pgfpathlineto{\pgfqpoint{1.704494in}{1.547720in}}%
\pgfpathlineto{\pgfqpoint{1.703684in}{1.498934in}}%
\pgfpathlineto{\pgfqpoint{1.702328in}{1.430275in}}%
\pgfpathlineto{\pgfqpoint{1.697315in}{1.412242in}}%
\pgfpathlineto{\pgfqpoint{1.701799in}{1.402266in}}%
\pgfpathlineto{\pgfqpoint{1.692796in}{1.380831in}}%
\pgfpathlineto{\pgfqpoint{1.702281in}{1.363552in}}%
\pgfpathlineto{\pgfqpoint{1.694222in}{1.364874in}}%
\pgfpathlineto{\pgfqpoint{1.688724in}{1.375625in}}%
\pgfpathlineto{\pgfqpoint{1.669102in}{1.382994in}}%
\pgfpathlineto{\pgfqpoint{1.656727in}{1.389477in}}%
\pgfpathlineto{\pgfqpoint{1.635962in}{1.390016in}}%
\pgfpathlineto{\pgfqpoint{1.628754in}{1.384110in}}%
\pgfpathlineto{\pgfqpoint{1.605250in}{1.395738in}}%
\pgfpathlineto{\pgfqpoint{1.603447in}{1.399416in}}%
\pgfpathlineto{\pgfqpoint{1.521419in}{1.403298in}}%
\pgfpathlineto{\pgfqpoint{1.471583in}{1.406155in}}%
\pgfpathlineto{\pgfqpoint{1.430419in}{1.409379in}}%
\pgfpathlineto{\pgfqpoint{1.362400in}{1.415805in}}%
\pgfpathlineto{\pgfqpoint{1.335768in}{1.418681in}}%
\pgfusepath{stroke}%
\end{pgfscope}%
\begin{pgfscope}%
\pgfpathrectangle{\pgfqpoint{0.100000in}{0.100000in}}{\pgfqpoint{3.007045in}{1.925000in}}%
\pgfusepath{clip}%
\pgfsetbuttcap%
\pgfsetroundjoin%
\pgfsetlinewidth{0.050187pt}%
\definecolor{currentstroke}{rgb}{1.000000,1.000000,1.000000}%
\pgfsetstrokecolor{currentstroke}%
\pgfsetdash{}{0pt}%
\pgfpathmoveto{\pgfqpoint{1.322867in}{1.287410in}}%
\pgfpathlineto{\pgfqpoint{1.279137in}{1.291403in}}%
\pgfpathlineto{\pgfqpoint{1.183815in}{1.303168in}}%
\pgfpathlineto{\pgfqpoint{1.131961in}{1.310611in}}%
\pgfpathlineto{\pgfqpoint{1.076345in}{1.318594in}}%
\pgfpathlineto{\pgfqpoint{1.029497in}{1.326091in}}%
\pgfpathlineto{\pgfqpoint{0.978076in}{1.334895in}}%
\pgfpathlineto{\pgfqpoint{0.989767in}{1.399751in}}%
\pgfpathlineto{\pgfqpoint{1.001610in}{1.466254in}}%
\pgfpathlineto{\pgfqpoint{1.018377in}{1.560840in}}%
\pgfpathlineto{\pgfqpoint{1.024367in}{1.594897in}}%
\pgfpathlineto{\pgfqpoint{1.053106in}{1.589639in}}%
\pgfpathlineto{\pgfqpoint{1.136592in}{1.576534in}}%
\pgfpathlineto{\pgfqpoint{1.170478in}{1.571700in}}%
\pgfpathlineto{\pgfqpoint{1.246348in}{1.561317in}}%
\pgfpathlineto{\pgfqpoint{1.300990in}{1.555114in}}%
\pgfpathlineto{\pgfqpoint{1.348717in}{1.549903in}}%
\pgfpathlineto{\pgfqpoint{1.343289in}{1.496206in}}%
\pgfpathlineto{\pgfqpoint{1.335768in}{1.418681in}}%
\pgfpathlineto{\pgfqpoint{1.329313in}{1.352796in}}%
\pgfpathlineto{\pgfqpoint{1.322867in}{1.287410in}}%
\pgfusepath{stroke}%
\end{pgfscope}%
\begin{pgfscope}%
\pgfpathrectangle{\pgfqpoint{0.100000in}{0.100000in}}{\pgfqpoint{3.007045in}{1.925000in}}%
\pgfusepath{clip}%
\pgfsetbuttcap%
\pgfsetroundjoin%
\pgfsetlinewidth{0.050187pt}%
\definecolor{currentstroke}{rgb}{1.000000,1.000000,1.000000}%
\pgfsetstrokecolor{currentstroke}%
\pgfsetdash{}{0pt}%
\pgfpathmoveto{\pgfqpoint{2.122408in}{1.376872in}}%
\pgfpathlineto{\pgfqpoint{2.066992in}{1.372935in}}%
\pgfpathlineto{\pgfqpoint{1.984391in}{1.369415in}}%
\pgfpathlineto{\pgfqpoint{1.981252in}{1.377759in}}%
\pgfpathlineto{\pgfqpoint{1.962935in}{1.383996in}}%
\pgfpathlineto{\pgfqpoint{1.958894in}{1.395780in}}%
\pgfpathlineto{\pgfqpoint{1.957202in}{1.410359in}}%
\pgfpathlineto{\pgfqpoint{1.961330in}{1.417827in}}%
\pgfpathlineto{\pgfqpoint{1.954811in}{1.424995in}}%
\pgfpathlineto{\pgfqpoint{1.953238in}{1.433544in}}%
\pgfpathlineto{\pgfqpoint{1.951137in}{1.452456in}}%
\pgfpathlineto{\pgfqpoint{1.944889in}{1.462729in}}%
\pgfpathlineto{\pgfqpoint{1.933793in}{1.468495in}}%
\pgfpathlineto{\pgfqpoint{1.921718in}{1.478009in}}%
\pgfpathlineto{\pgfqpoint{1.915477in}{1.489266in}}%
\pgfpathlineto{\pgfqpoint{1.904277in}{1.493800in}}%
\pgfpathlineto{\pgfqpoint{1.897693in}{1.501178in}}%
\pgfpathlineto{\pgfqpoint{1.889718in}{1.502430in}}%
\pgfpathlineto{\pgfqpoint{1.875478in}{1.513385in}}%
\pgfpathlineto{\pgfqpoint{1.877791in}{1.525991in}}%
\pgfpathlineto{\pgfqpoint{1.877345in}{1.549965in}}%
\pgfpathlineto{\pgfqpoint{1.881737in}{1.556564in}}%
\pgfpathlineto{\pgfqpoint{1.877791in}{1.566528in}}%
\pgfpathlineto{\pgfqpoint{1.870840in}{1.568455in}}%
\pgfpathlineto{\pgfqpoint{1.871415in}{1.577207in}}%
\pgfpathlineto{\pgfqpoint{1.880038in}{1.591032in}}%
\pgfpathlineto{\pgfqpoint{1.897116in}{1.601964in}}%
\pgfpathlineto{\pgfqpoint{1.896052in}{1.640844in}}%
\pgfpathlineto{\pgfqpoint{1.904597in}{1.646684in}}%
\pgfpathlineto{\pgfqpoint{1.912684in}{1.642793in}}%
\pgfpathlineto{\pgfqpoint{1.929157in}{1.648489in}}%
\pgfpathlineto{\pgfqpoint{1.960149in}{1.662803in}}%
\pgfpathlineto{\pgfqpoint{1.964187in}{1.658371in}}%
\pgfpathlineto{\pgfqpoint{1.958317in}{1.638296in}}%
\pgfpathlineto{\pgfqpoint{1.967050in}{1.642702in}}%
\pgfpathlineto{\pgfqpoint{1.982007in}{1.638302in}}%
\pgfpathlineto{\pgfqpoint{1.991199in}{1.634631in}}%
\pgfpathlineto{\pgfqpoint{1.996352in}{1.623871in}}%
\pgfpathlineto{\pgfqpoint{2.043506in}{1.613716in}}%
\pgfpathlineto{\pgfqpoint{2.057600in}{1.606750in}}%
\pgfpathlineto{\pgfqpoint{2.071931in}{1.606832in}}%
\pgfpathlineto{\pgfqpoint{2.086656in}{1.603958in}}%
\pgfpathlineto{\pgfqpoint{2.096212in}{1.594142in}}%
\pgfpathlineto{\pgfqpoint{2.105344in}{1.589291in}}%
\pgfpathlineto{\pgfqpoint{2.107021in}{1.575301in}}%
\pgfpathlineto{\pgfqpoint{2.104324in}{1.566517in}}%
\pgfpathlineto{\pgfqpoint{2.114501in}{1.565867in}}%
\pgfpathlineto{\pgfqpoint{2.111072in}{1.555679in}}%
\pgfpathlineto{\pgfqpoint{2.114334in}{1.552060in}}%
\pgfpathlineto{\pgfqpoint{2.117577in}{1.542443in}}%
\pgfpathlineto{\pgfqpoint{2.107659in}{1.537351in}}%
\pgfpathlineto{\pgfqpoint{2.101900in}{1.523164in}}%
\pgfpathlineto{\pgfqpoint{2.100093in}{1.513132in}}%
\pgfpathlineto{\pgfqpoint{2.105618in}{1.511415in}}%
\pgfpathlineto{\pgfqpoint{2.112691in}{1.518938in}}%
\pgfpathlineto{\pgfqpoint{2.118705in}{1.531974in}}%
\pgfpathlineto{\pgfqpoint{2.126842in}{1.536502in}}%
\pgfpathlineto{\pgfqpoint{2.132957in}{1.530613in}}%
\pgfpathlineto{\pgfqpoint{2.126938in}{1.512787in}}%
\pgfpathlineto{\pgfqpoint{2.125050in}{1.498949in}}%
\pgfpathlineto{\pgfqpoint{2.126828in}{1.489001in}}%
\pgfpathlineto{\pgfqpoint{2.121239in}{1.483374in}}%
\pgfpathlineto{\pgfqpoint{2.118447in}{1.469622in}}%
\pgfpathlineto{\pgfqpoint{2.120728in}{1.455170in}}%
\pgfpathlineto{\pgfqpoint{2.114132in}{1.433796in}}%
\pgfpathlineto{\pgfqpoint{2.114271in}{1.423115in}}%
\pgfpathlineto{\pgfqpoint{2.119483in}{1.399971in}}%
\pgfpathlineto{\pgfqpoint{2.122861in}{1.396001in}}%
\pgfpathlineto{\pgfqpoint{2.122408in}{1.376872in}}%
\pgfusepath{stroke}%
\end{pgfscope}%
\begin{pgfscope}%
\pgfpathrectangle{\pgfqpoint{0.100000in}{0.100000in}}{\pgfqpoint{3.007045in}{1.925000in}}%
\pgfusepath{clip}%
\pgfsetbuttcap%
\pgfsetroundjoin%
\pgfsetlinewidth{0.050187pt}%
\definecolor{currentstroke}{rgb}{1.000000,1.000000,1.000000}%
\pgfsetstrokecolor{currentstroke}%
\pgfsetdash{}{0pt}%
\pgfpathmoveto{\pgfqpoint{2.143186in}{1.564449in}}%
\pgfpathlineto{\pgfqpoint{2.144423in}{1.555218in}}%
\pgfpathlineto{\pgfqpoint{2.133091in}{1.530893in}}%
\pgfpathlineto{\pgfqpoint{2.128055in}{1.537939in}}%
\pgfpathlineto{\pgfqpoint{2.143186in}{1.564449in}}%
\pgfusepath{stroke}%
\end{pgfscope}%
\begin{pgfscope}%
\pgfpathrectangle{\pgfqpoint{0.100000in}{0.100000in}}{\pgfqpoint{3.007045in}{1.925000in}}%
\pgfusepath{clip}%
\pgfsetbuttcap%
\pgfsetroundjoin%
\pgfsetlinewidth{0.050187pt}%
\definecolor{currentstroke}{rgb}{1.000000,1.000000,1.000000}%
\pgfsetstrokecolor{currentstroke}%
\pgfsetdash{}{0pt}%
\pgfpathmoveto{\pgfqpoint{0.771499in}{1.716207in}}%
\pgfpathlineto{\pgfqpoint{0.773739in}{1.727298in}}%
\pgfpathlineto{\pgfqpoint{0.772049in}{1.739338in}}%
\pgfpathlineto{\pgfqpoint{0.774637in}{1.752465in}}%
\pgfpathlineto{\pgfqpoint{0.787821in}{1.805724in}}%
\pgfpathlineto{\pgfqpoint{0.798909in}{1.849997in}}%
\pgfpathlineto{\pgfqpoint{0.814503in}{1.911787in}}%
\pgfpathlineto{\pgfqpoint{0.855834in}{1.901800in}}%
\pgfpathlineto{\pgfqpoint{0.840559in}{1.835130in}}%
\pgfpathlineto{\pgfqpoint{0.850232in}{1.813647in}}%
\pgfpathlineto{\pgfqpoint{0.849009in}{1.801984in}}%
\pgfpathlineto{\pgfqpoint{0.854194in}{1.784436in}}%
\pgfpathlineto{\pgfqpoint{0.861436in}{1.780289in}}%
\pgfpathlineto{\pgfqpoint{0.869360in}{1.759957in}}%
\pgfpathlineto{\pgfqpoint{0.874228in}{1.754357in}}%
\pgfpathlineto{\pgfqpoint{0.873638in}{1.747149in}}%
\pgfpathlineto{\pgfqpoint{0.882867in}{1.740700in}}%
\pgfpathlineto{\pgfqpoint{0.884169in}{1.733223in}}%
\pgfpathlineto{\pgfqpoint{0.896838in}{1.732179in}}%
\pgfpathlineto{\pgfqpoint{0.880771in}{1.690510in}}%
\pgfpathlineto{\pgfqpoint{0.883228in}{1.681901in}}%
\pgfpathlineto{\pgfqpoint{0.873985in}{1.676547in}}%
\pgfpathlineto{\pgfqpoint{0.872732in}{1.662959in}}%
\pgfpathlineto{\pgfqpoint{0.880439in}{1.654450in}}%
\pgfpathlineto{\pgfqpoint{0.890564in}{1.658578in}}%
\pgfpathlineto{\pgfqpoint{0.898002in}{1.666211in}}%
\pgfpathlineto{\pgfqpoint{0.905930in}{1.657584in}}%
\pgfpathlineto{\pgfqpoint{0.905308in}{1.639956in}}%
\pgfpathlineto{\pgfqpoint{0.914532in}{1.619218in}}%
\pgfpathlineto{\pgfqpoint{0.912420in}{1.606664in}}%
\pgfpathlineto{\pgfqpoint{0.925374in}{1.597953in}}%
\pgfpathlineto{\pgfqpoint{0.925581in}{1.585841in}}%
\pgfpathlineto{\pgfqpoint{0.929993in}{1.573111in}}%
\pgfpathlineto{\pgfqpoint{0.938194in}{1.575980in}}%
\pgfpathlineto{\pgfqpoint{0.956987in}{1.570638in}}%
\pgfpathlineto{\pgfqpoint{0.964724in}{1.576574in}}%
\pgfpathlineto{\pgfqpoint{0.970258in}{1.572737in}}%
\pgfpathlineto{\pgfqpoint{0.981264in}{1.573217in}}%
\pgfpathlineto{\pgfqpoint{0.983625in}{1.569254in}}%
\pgfpathlineto{\pgfqpoint{0.995375in}{1.570362in}}%
\pgfpathlineto{\pgfqpoint{1.001038in}{1.578644in}}%
\pgfpathlineto{\pgfqpoint{1.008949in}{1.579653in}}%
\pgfpathlineto{\pgfqpoint{1.015103in}{1.563270in}}%
\pgfpathlineto{\pgfqpoint{1.018377in}{1.560840in}}%
\pgfpathlineto{\pgfqpoint{1.001610in}{1.466254in}}%
\pgfpathlineto{\pgfqpoint{0.989767in}{1.399751in}}%
\pgfpathlineto{\pgfqpoint{0.896364in}{1.417782in}}%
\pgfpathlineto{\pgfqpoint{0.845912in}{1.427818in}}%
\pgfpathlineto{\pgfqpoint{0.798720in}{1.438195in}}%
\pgfpathlineto{\pgfqpoint{0.754611in}{1.448181in}}%
\pgfpathlineto{\pgfqpoint{0.703563in}{1.460492in}}%
\pgfpathlineto{\pgfqpoint{0.729893in}{1.568511in}}%
\pgfpathlineto{\pgfqpoint{0.731313in}{1.576405in}}%
\pgfpathlineto{\pgfqpoint{0.743022in}{1.597090in}}%
\pgfpathlineto{\pgfqpoint{0.730898in}{1.609522in}}%
\pgfpathlineto{\pgfqpoint{0.733416in}{1.621731in}}%
\pgfpathlineto{\pgfqpoint{0.738031in}{1.624787in}}%
\pgfpathlineto{\pgfqpoint{0.746274in}{1.637367in}}%
\pgfpathlineto{\pgfqpoint{0.758289in}{1.646070in}}%
\pgfpathlineto{\pgfqpoint{0.758946in}{1.652510in}}%
\pgfpathlineto{\pgfqpoint{0.766198in}{1.658944in}}%
\pgfpathlineto{\pgfqpoint{0.772219in}{1.670988in}}%
\pgfpathlineto{\pgfqpoint{0.784570in}{1.683729in}}%
\pgfpathlineto{\pgfqpoint{0.783693in}{1.696314in}}%
\pgfpathlineto{\pgfqpoint{0.775249in}{1.703304in}}%
\pgfpathlineto{\pgfqpoint{0.771499in}{1.716207in}}%
\pgfusepath{stroke}%
\end{pgfscope}%
\begin{pgfscope}%
\pgfpathrectangle{\pgfqpoint{0.100000in}{0.100000in}}{\pgfqpoint{3.007045in}{1.925000in}}%
\pgfusepath{clip}%
\pgfsetbuttcap%
\pgfsetroundjoin%
\pgfsetlinewidth{0.050187pt}%
\definecolor{currentstroke}{rgb}{1.000000,1.000000,1.000000}%
\pgfsetstrokecolor{currentstroke}%
\pgfsetdash{}{0pt}%
\pgfpathmoveto{\pgfqpoint{2.816139in}{1.505410in}}%
\pgfpathlineto{\pgfqpoint{2.812258in}{1.517661in}}%
\pgfpathlineto{\pgfqpoint{2.805056in}{1.554825in}}%
\pgfpathlineto{\pgfqpoint{2.795745in}{1.566258in}}%
\pgfpathlineto{\pgfqpoint{2.787577in}{1.586820in}}%
\pgfpathlineto{\pgfqpoint{2.790259in}{1.602064in}}%
\pgfpathlineto{\pgfqpoint{2.788108in}{1.613300in}}%
\pgfpathlineto{\pgfqpoint{2.781130in}{1.624335in}}%
\pgfpathlineto{\pgfqpoint{2.780843in}{1.636492in}}%
\pgfpathlineto{\pgfqpoint{2.776791in}{1.649602in}}%
\pgfpathlineto{\pgfqpoint{2.813053in}{1.658527in}}%
\pgfpathlineto{\pgfqpoint{2.860154in}{1.671241in}}%
\pgfpathlineto{\pgfqpoint{2.862034in}{1.663950in}}%
\pgfpathlineto{\pgfqpoint{2.859035in}{1.652344in}}%
\pgfpathlineto{\pgfqpoint{2.866075in}{1.643057in}}%
\pgfpathlineto{\pgfqpoint{2.862336in}{1.631290in}}%
\pgfpathlineto{\pgfqpoint{2.847477in}{1.616498in}}%
\pgfpathlineto{\pgfqpoint{2.851562in}{1.605386in}}%
\pgfpathlineto{\pgfqpoint{2.848535in}{1.581900in}}%
\pgfpathlineto{\pgfqpoint{2.844297in}{1.568658in}}%
\pgfpathlineto{\pgfqpoint{2.849703in}{1.531000in}}%
\pgfpathlineto{\pgfqpoint{2.848975in}{1.517751in}}%
\pgfpathlineto{\pgfqpoint{2.854206in}{1.513499in}}%
\pgfpathlineto{\pgfqpoint{2.816139in}{1.505410in}}%
\pgfusepath{stroke}%
\end{pgfscope}%
\begin{pgfscope}%
\pgfpathrectangle{\pgfqpoint{0.100000in}{0.100000in}}{\pgfqpoint{3.007045in}{1.925000in}}%
\pgfusepath{clip}%
\pgfsetbuttcap%
\pgfsetroundjoin%
\pgfsetlinewidth{0.050187pt}%
\definecolor{currentstroke}{rgb}{1.000000,1.000000,1.000000}%
\pgfsetstrokecolor{currentstroke}%
\pgfsetdash{}{0pt}%
\pgfpathmoveto{\pgfqpoint{1.702328in}{1.430275in}}%
\pgfpathlineto{\pgfqpoint{1.703684in}{1.498934in}}%
\pgfpathlineto{\pgfqpoint{1.704494in}{1.547720in}}%
\pgfpathlineto{\pgfqpoint{1.700966in}{1.554402in}}%
\pgfpathlineto{\pgfqpoint{1.693533in}{1.556888in}}%
\pgfpathlineto{\pgfqpoint{1.686792in}{1.568520in}}%
\pgfpathlineto{\pgfqpoint{1.695595in}{1.578623in}}%
\pgfpathlineto{\pgfqpoint{1.700005in}{1.591493in}}%
\pgfpathlineto{\pgfqpoint{1.700880in}{1.602193in}}%
\pgfpathlineto{\pgfqpoint{1.698854in}{1.618219in}}%
\pgfpathlineto{\pgfqpoint{1.692711in}{1.628504in}}%
\pgfpathlineto{\pgfqpoint{1.690611in}{1.637873in}}%
\pgfpathlineto{\pgfqpoint{1.689081in}{1.672635in}}%
\pgfpathlineto{\pgfqpoint{1.689100in}{1.702558in}}%
\pgfpathlineto{\pgfqpoint{1.680835in}{1.726326in}}%
\pgfpathlineto{\pgfqpoint{1.677227in}{1.742437in}}%
\pgfpathlineto{\pgfqpoint{1.677852in}{1.763948in}}%
\pgfpathlineto{\pgfqpoint{1.679750in}{1.774414in}}%
\pgfpathlineto{\pgfqpoint{1.674447in}{1.796314in}}%
\pgfpathlineto{\pgfqpoint{1.710550in}{1.795591in}}%
\pgfpathlineto{\pgfqpoint{1.765389in}{1.795120in}}%
\pgfpathlineto{\pgfqpoint{1.765689in}{1.820011in}}%
\pgfpathlineto{\pgfqpoint{1.779648in}{1.817270in}}%
\pgfpathlineto{\pgfqpoint{1.786338in}{1.786920in}}%
\pgfpathlineto{\pgfqpoint{1.791270in}{1.776008in}}%
\pgfpathlineto{\pgfqpoint{1.803535in}{1.775685in}}%
\pgfpathlineto{\pgfqpoint{1.806280in}{1.771983in}}%
\pgfpathlineto{\pgfqpoint{1.823388in}{1.770342in}}%
\pgfpathlineto{\pgfqpoint{1.826264in}{1.762818in}}%
\pgfpathlineto{\pgfqpoint{1.838054in}{1.764509in}}%
\pgfpathlineto{\pgfqpoint{1.847211in}{1.771547in}}%
\pgfpathlineto{\pgfqpoint{1.863003in}{1.771284in}}%
\pgfpathlineto{\pgfqpoint{1.872775in}{1.765624in}}%
\pgfpathlineto{\pgfqpoint{1.873898in}{1.760316in}}%
\pgfpathlineto{\pgfqpoint{1.883196in}{1.759205in}}%
\pgfpathlineto{\pgfqpoint{1.889264in}{1.744727in}}%
\pgfpathlineto{\pgfqpoint{1.893179in}{1.753629in}}%
\pgfpathlineto{\pgfqpoint{1.903859in}{1.754177in}}%
\pgfpathlineto{\pgfqpoint{1.906561in}{1.747244in}}%
\pgfpathlineto{\pgfqpoint{1.918579in}{1.744079in}}%
\pgfpathlineto{\pgfqpoint{1.925206in}{1.734093in}}%
\pgfpathlineto{\pgfqpoint{1.939904in}{1.737194in}}%
\pgfpathlineto{\pgfqpoint{1.956078in}{1.749448in}}%
\pgfpathlineto{\pgfqpoint{1.961975in}{1.738646in}}%
\pgfpathlineto{\pgfqpoint{1.988513in}{1.741602in}}%
\pgfpathlineto{\pgfqpoint{1.999857in}{1.734188in}}%
\pgfpathlineto{\pgfqpoint{2.006460in}{1.736835in}}%
\pgfpathlineto{\pgfqpoint{2.011755in}{1.732661in}}%
\pgfpathlineto{\pgfqpoint{1.996048in}{1.722759in}}%
\pgfpathlineto{\pgfqpoint{1.973627in}{1.713935in}}%
\pgfpathlineto{\pgfqpoint{1.951422in}{1.696299in}}%
\pgfpathlineto{\pgfqpoint{1.932171in}{1.673105in}}%
\pgfpathlineto{\pgfqpoint{1.917597in}{1.659394in}}%
\pgfpathlineto{\pgfqpoint{1.904832in}{1.649969in}}%
\pgfpathlineto{\pgfqpoint{1.896052in}{1.640844in}}%
\pgfpathlineto{\pgfqpoint{1.897116in}{1.601964in}}%
\pgfpathlineto{\pgfqpoint{1.880038in}{1.591032in}}%
\pgfpathlineto{\pgfqpoint{1.871415in}{1.577207in}}%
\pgfpathlineto{\pgfqpoint{1.870840in}{1.568455in}}%
\pgfpathlineto{\pgfqpoint{1.877791in}{1.566528in}}%
\pgfpathlineto{\pgfqpoint{1.881737in}{1.556564in}}%
\pgfpathlineto{\pgfqpoint{1.877345in}{1.549965in}}%
\pgfpathlineto{\pgfqpoint{1.877791in}{1.525991in}}%
\pgfpathlineto{\pgfqpoint{1.875478in}{1.513385in}}%
\pgfpathlineto{\pgfqpoint{1.889718in}{1.502430in}}%
\pgfpathlineto{\pgfqpoint{1.897693in}{1.501178in}}%
\pgfpathlineto{\pgfqpoint{1.904277in}{1.493800in}}%
\pgfpathlineto{\pgfqpoint{1.915477in}{1.489266in}}%
\pgfpathlineto{\pgfqpoint{1.921718in}{1.478009in}}%
\pgfpathlineto{\pgfqpoint{1.933793in}{1.468495in}}%
\pgfpathlineto{\pgfqpoint{1.944889in}{1.462729in}}%
\pgfpathlineto{\pgfqpoint{1.951137in}{1.452456in}}%
\pgfpathlineto{\pgfqpoint{1.953238in}{1.433544in}}%
\pgfpathlineto{\pgfqpoint{1.894352in}{1.431406in}}%
\pgfpathlineto{\pgfqpoint{1.844156in}{1.430356in}}%
\pgfpathlineto{\pgfqpoint{1.798422in}{1.429684in}}%
\pgfpathlineto{\pgfqpoint{1.750041in}{1.429758in}}%
\pgfpathlineto{\pgfqpoint{1.702328in}{1.430275in}}%
\pgfusepath{stroke}%
\end{pgfscope}%
\begin{pgfscope}%
\pgfpathrectangle{\pgfqpoint{0.100000in}{0.100000in}}{\pgfqpoint{3.007045in}{1.925000in}}%
\pgfusepath{clip}%
\pgfsetbuttcap%
\pgfsetroundjoin%
\pgfsetlinewidth{0.050187pt}%
\definecolor{currentstroke}{rgb}{1.000000,1.000000,1.000000}%
\pgfsetstrokecolor{currentstroke}%
\pgfsetdash{}{0pt}%
\pgfpathmoveto{\pgfqpoint{0.365302in}{1.558630in}}%
\pgfpathlineto{\pgfqpoint{0.360639in}{1.567212in}}%
\pgfpathlineto{\pgfqpoint{0.363617in}{1.589239in}}%
\pgfpathlineto{\pgfqpoint{0.369364in}{1.600750in}}%
\pgfpathlineto{\pgfqpoint{0.366491in}{1.616323in}}%
\pgfpathlineto{\pgfqpoint{0.372487in}{1.622871in}}%
\pgfpathlineto{\pgfqpoint{0.383425in}{1.640522in}}%
\pgfpathlineto{\pgfqpoint{0.392700in}{1.651178in}}%
\pgfpathlineto{\pgfqpoint{0.406257in}{1.674384in}}%
\pgfpathlineto{\pgfqpoint{0.416655in}{1.699664in}}%
\pgfpathlineto{\pgfqpoint{0.427694in}{1.723373in}}%
\pgfpathlineto{\pgfqpoint{0.429925in}{1.733275in}}%
\pgfpathlineto{\pgfqpoint{0.445153in}{1.761603in}}%
\pgfpathlineto{\pgfqpoint{0.448085in}{1.774045in}}%
\pgfpathlineto{\pgfqpoint{0.454568in}{1.787110in}}%
\pgfpathlineto{\pgfqpoint{0.456886in}{1.803047in}}%
\pgfpathlineto{\pgfqpoint{0.469263in}{1.810819in}}%
\pgfpathlineto{\pgfqpoint{0.483940in}{1.814736in}}%
\pgfpathlineto{\pgfqpoint{0.491359in}{1.805963in}}%
\pgfpathlineto{\pgfqpoint{0.497775in}{1.806612in}}%
\pgfpathlineto{\pgfqpoint{0.507790in}{1.796155in}}%
\pgfpathlineto{\pgfqpoint{0.509122in}{1.786306in}}%
\pgfpathlineto{\pgfqpoint{0.505323in}{1.767019in}}%
\pgfpathlineto{\pgfqpoint{0.517960in}{1.757179in}}%
\pgfpathlineto{\pgfqpoint{0.526157in}{1.753479in}}%
\pgfpathlineto{\pgfqpoint{0.548360in}{1.757255in}}%
\pgfpathlineto{\pgfqpoint{0.561223in}{1.754676in}}%
\pgfpathlineto{\pgfqpoint{0.570064in}{1.750443in}}%
\pgfpathlineto{\pgfqpoint{0.574722in}{1.742409in}}%
\pgfpathlineto{\pgfqpoint{0.604211in}{1.743481in}}%
\pgfpathlineto{\pgfqpoint{0.608900in}{1.738565in}}%
\pgfpathlineto{\pgfqpoint{0.619963in}{1.737504in}}%
\pgfpathlineto{\pgfqpoint{0.631104in}{1.740645in}}%
\pgfpathlineto{\pgfqpoint{0.664042in}{1.739853in}}%
\pgfpathlineto{\pgfqpoint{0.671157in}{1.737433in}}%
\pgfpathlineto{\pgfqpoint{0.679470in}{1.740222in}}%
\pgfpathlineto{\pgfqpoint{0.771499in}{1.716207in}}%
\pgfpathlineto{\pgfqpoint{0.775249in}{1.703304in}}%
\pgfpathlineto{\pgfqpoint{0.783693in}{1.696314in}}%
\pgfpathlineto{\pgfqpoint{0.784570in}{1.683729in}}%
\pgfpathlineto{\pgfqpoint{0.772219in}{1.670988in}}%
\pgfpathlineto{\pgfqpoint{0.766198in}{1.658944in}}%
\pgfpathlineto{\pgfqpoint{0.758946in}{1.652510in}}%
\pgfpathlineto{\pgfqpoint{0.758289in}{1.646070in}}%
\pgfpathlineto{\pgfqpoint{0.746274in}{1.637367in}}%
\pgfpathlineto{\pgfqpoint{0.738031in}{1.624787in}}%
\pgfpathlineto{\pgfqpoint{0.733416in}{1.621731in}}%
\pgfpathlineto{\pgfqpoint{0.730898in}{1.609522in}}%
\pgfpathlineto{\pgfqpoint{0.743022in}{1.597090in}}%
\pgfpathlineto{\pgfqpoint{0.731313in}{1.576405in}}%
\pgfpathlineto{\pgfqpoint{0.729893in}{1.568511in}}%
\pgfpathlineto{\pgfqpoint{0.703563in}{1.460492in}}%
\pgfpathlineto{\pgfqpoint{0.648178in}{1.474684in}}%
\pgfpathlineto{\pgfqpoint{0.594748in}{1.488461in}}%
\pgfpathlineto{\pgfqpoint{0.562516in}{1.497430in}}%
\pgfpathlineto{\pgfqpoint{0.521101in}{1.509228in}}%
\pgfpathlineto{\pgfqpoint{0.455041in}{1.530020in}}%
\pgfpathlineto{\pgfqpoint{0.383230in}{1.552378in}}%
\pgfpathlineto{\pgfqpoint{0.365302in}{1.558630in}}%
\pgfusepath{stroke}%
\end{pgfscope}%
\begin{pgfscope}%
\pgfpathrectangle{\pgfqpoint{0.100000in}{0.100000in}}{\pgfqpoint{3.007045in}{1.925000in}}%
\pgfusepath{clip}%
\pgfsetbuttcap%
\pgfsetroundjoin%
\pgfsetlinewidth{0.050187pt}%
\definecolor{currentstroke}{rgb}{1.000000,1.000000,1.000000}%
\pgfsetstrokecolor{currentstroke}%
\pgfsetdash{}{0pt}%
\pgfpathmoveto{\pgfqpoint{2.854206in}{1.513499in}}%
\pgfpathlineto{\pgfqpoint{2.848975in}{1.517751in}}%
\pgfpathlineto{\pgfqpoint{2.849703in}{1.531000in}}%
\pgfpathlineto{\pgfqpoint{2.844297in}{1.568658in}}%
\pgfpathlineto{\pgfqpoint{2.848535in}{1.581900in}}%
\pgfpathlineto{\pgfqpoint{2.851562in}{1.605386in}}%
\pgfpathlineto{\pgfqpoint{2.847477in}{1.616498in}}%
\pgfpathlineto{\pgfqpoint{2.862336in}{1.631290in}}%
\pgfpathlineto{\pgfqpoint{2.866075in}{1.643057in}}%
\pgfpathlineto{\pgfqpoint{2.859035in}{1.652344in}}%
\pgfpathlineto{\pgfqpoint{2.862034in}{1.663950in}}%
\pgfpathlineto{\pgfqpoint{2.860154in}{1.671241in}}%
\pgfpathlineto{\pgfqpoint{2.861768in}{1.686854in}}%
\pgfpathlineto{\pgfqpoint{2.864787in}{1.691674in}}%
\pgfpathlineto{\pgfqpoint{2.874093in}{1.694717in}}%
\pgfpathlineto{\pgfqpoint{2.887667in}{1.654962in}}%
\pgfpathlineto{\pgfqpoint{2.912852in}{1.572583in}}%
\pgfpathlineto{\pgfqpoint{2.922382in}{1.566196in}}%
\pgfpathlineto{\pgfqpoint{2.924097in}{1.558905in}}%
\pgfpathlineto{\pgfqpoint{2.929126in}{1.555958in}}%
\pgfpathlineto{\pgfqpoint{2.928744in}{1.542735in}}%
\pgfpathlineto{\pgfqpoint{2.923415in}{1.542522in}}%
\pgfpathlineto{\pgfqpoint{2.912617in}{1.534288in}}%
\pgfpathlineto{\pgfqpoint{2.909503in}{1.526026in}}%
\pgfpathlineto{\pgfqpoint{2.880603in}{1.518876in}}%
\pgfpathlineto{\pgfqpoint{2.854206in}{1.513499in}}%
\pgfusepath{stroke}%
\end{pgfscope}%
\begin{pgfscope}%
\pgfpathrectangle{\pgfqpoint{0.100000in}{0.100000in}}{\pgfqpoint{3.007045in}{1.925000in}}%
\pgfusepath{clip}%
\pgfsetbuttcap%
\pgfsetroundjoin%
\pgfsetlinewidth{0.050187pt}%
\definecolor{currentstroke}{rgb}{1.000000,1.000000,1.000000}%
\pgfsetstrokecolor{currentstroke}%
\pgfsetdash{}{0pt}%
\pgfpathmoveto{\pgfqpoint{1.950519in}{1.226895in}}%
\pgfpathlineto{\pgfqpoint{1.935262in}{1.242004in}}%
\pgfpathlineto{\pgfqpoint{1.886493in}{1.239191in}}%
\pgfpathlineto{\pgfqpoint{1.810432in}{1.236685in}}%
\pgfpathlineto{\pgfqpoint{1.733914in}{1.237878in}}%
\pgfpathlineto{\pgfqpoint{1.728543in}{1.247244in}}%
\pgfpathlineto{\pgfqpoint{1.730736in}{1.256440in}}%
\pgfpathlineto{\pgfqpoint{1.729695in}{1.272194in}}%
\pgfpathlineto{\pgfqpoint{1.725651in}{1.287453in}}%
\pgfpathlineto{\pgfqpoint{1.726121in}{1.295511in}}%
\pgfpathlineto{\pgfqpoint{1.717525in}{1.304616in}}%
\pgfpathlineto{\pgfqpoint{1.719395in}{1.317284in}}%
\pgfpathlineto{\pgfqpoint{1.706209in}{1.342308in}}%
\pgfpathlineto{\pgfqpoint{1.702281in}{1.363552in}}%
\pgfpathlineto{\pgfqpoint{1.692796in}{1.380831in}}%
\pgfpathlineto{\pgfqpoint{1.701799in}{1.402266in}}%
\pgfpathlineto{\pgfqpoint{1.697315in}{1.412242in}}%
\pgfpathlineto{\pgfqpoint{1.702328in}{1.430275in}}%
\pgfpathlineto{\pgfqpoint{1.750041in}{1.429758in}}%
\pgfpathlineto{\pgfqpoint{1.798422in}{1.429684in}}%
\pgfpathlineto{\pgfqpoint{1.844156in}{1.430356in}}%
\pgfpathlineto{\pgfqpoint{1.894352in}{1.431406in}}%
\pgfpathlineto{\pgfqpoint{1.953238in}{1.433544in}}%
\pgfpathlineto{\pgfqpoint{1.954811in}{1.424995in}}%
\pgfpathlineto{\pgfqpoint{1.961330in}{1.417827in}}%
\pgfpathlineto{\pgfqpoint{1.957202in}{1.410359in}}%
\pgfpathlineto{\pgfqpoint{1.958894in}{1.395780in}}%
\pgfpathlineto{\pgfqpoint{1.962935in}{1.383996in}}%
\pgfpathlineto{\pgfqpoint{1.981252in}{1.377759in}}%
\pgfpathlineto{\pgfqpoint{1.984391in}{1.369415in}}%
\pgfpathlineto{\pgfqpoint{1.994442in}{1.360047in}}%
\pgfpathlineto{\pgfqpoint{1.998543in}{1.350356in}}%
\pgfpathlineto{\pgfqpoint{2.008727in}{1.343855in}}%
\pgfpathlineto{\pgfqpoint{2.010321in}{1.336029in}}%
\pgfpathlineto{\pgfqpoint{2.008338in}{1.324180in}}%
\pgfpathlineto{\pgfqpoint{2.003154in}{1.320629in}}%
\pgfpathlineto{\pgfqpoint{2.001591in}{1.309350in}}%
\pgfpathlineto{\pgfqpoint{1.986693in}{1.300392in}}%
\pgfpathlineto{\pgfqpoint{1.967275in}{1.295485in}}%
\pgfpathlineto{\pgfqpoint{1.965496in}{1.284201in}}%
\pgfpathlineto{\pgfqpoint{1.972987in}{1.276140in}}%
\pgfpathlineto{\pgfqpoint{1.973292in}{1.266004in}}%
\pgfpathlineto{\pgfqpoint{1.967249in}{1.258037in}}%
\pgfpathlineto{\pgfqpoint{1.964076in}{1.246200in}}%
\pgfpathlineto{\pgfqpoint{1.953571in}{1.242282in}}%
\pgfpathlineto{\pgfqpoint{1.954240in}{1.229091in}}%
\pgfpathlineto{\pgfqpoint{1.950519in}{1.226895in}}%
\pgfusepath{stroke}%
\end{pgfscope}%
\begin{pgfscope}%
\pgfpathrectangle{\pgfqpoint{0.100000in}{0.100000in}}{\pgfqpoint{3.007045in}{1.925000in}}%
\pgfusepath{clip}%
\pgfsetbuttcap%
\pgfsetroundjoin%
\pgfsetlinewidth{0.050187pt}%
\definecolor{currentstroke}{rgb}{1.000000,1.000000,1.000000}%
\pgfsetstrokecolor{currentstroke}%
\pgfsetdash{}{0pt}%
\pgfpathmoveto{\pgfqpoint{2.897124in}{1.475380in}}%
\pgfpathlineto{\pgfqpoint{2.882673in}{1.473100in}}%
\pgfpathlineto{\pgfqpoint{2.816297in}{1.458017in}}%
\pgfpathlineto{\pgfqpoint{2.815137in}{1.459774in}}%
\pgfpathlineto{\pgfqpoint{2.816139in}{1.505410in}}%
\pgfpathlineto{\pgfqpoint{2.854206in}{1.513499in}}%
\pgfpathlineto{\pgfqpoint{2.880603in}{1.518876in}}%
\pgfpathlineto{\pgfqpoint{2.909503in}{1.526026in}}%
\pgfpathlineto{\pgfqpoint{2.912617in}{1.534288in}}%
\pgfpathlineto{\pgfqpoint{2.923415in}{1.542522in}}%
\pgfpathlineto{\pgfqpoint{2.928744in}{1.542735in}}%
\pgfpathlineto{\pgfqpoint{2.935751in}{1.530721in}}%
\pgfpathlineto{\pgfqpoint{2.929393in}{1.513184in}}%
\pgfpathlineto{\pgfqpoint{2.928460in}{1.502901in}}%
\pgfpathlineto{\pgfqpoint{2.941328in}{1.503867in}}%
\pgfpathlineto{\pgfqpoint{2.947164in}{1.498934in}}%
\pgfpathlineto{\pgfqpoint{2.960252in}{1.478771in}}%
\pgfpathlineto{\pgfqpoint{2.966756in}{1.476318in}}%
\pgfpathlineto{\pgfqpoint{2.977652in}{1.477228in}}%
\pgfpathlineto{\pgfqpoint{2.985234in}{1.484080in}}%
\pgfpathlineto{\pgfqpoint{2.990274in}{1.477965in}}%
\pgfpathlineto{\pgfqpoint{2.958579in}{1.461239in}}%
\pgfpathlineto{\pgfqpoint{2.957589in}{1.473241in}}%
\pgfpathlineto{\pgfqpoint{2.943191in}{1.454588in}}%
\pgfpathlineto{\pgfqpoint{2.938129in}{1.451375in}}%
\pgfpathlineto{\pgfqpoint{2.931068in}{1.462137in}}%
\pgfpathlineto{\pgfqpoint{2.929135in}{1.463612in}}%
\pgfpathlineto{\pgfqpoint{2.922562in}{1.467094in}}%
\pgfpathlineto{\pgfqpoint{2.916830in}{1.481216in}}%
\pgfpathlineto{\pgfqpoint{2.897124in}{1.475380in}}%
\pgfusepath{stroke}%
\end{pgfscope}%
\begin{pgfscope}%
\pgfpathrectangle{\pgfqpoint{0.100000in}{0.100000in}}{\pgfqpoint{3.007045in}{1.925000in}}%
\pgfusepath{clip}%
\pgfsetbuttcap%
\pgfsetroundjoin%
\pgfsetlinewidth{0.050187pt}%
\definecolor{currentstroke}{rgb}{1.000000,1.000000,1.000000}%
\pgfsetstrokecolor{currentstroke}%
\pgfsetdash{}{0pt}%
\pgfpathmoveto{\pgfqpoint{1.416825in}{1.212637in}}%
\pgfpathlineto{\pgfqpoint{1.422117in}{1.278255in}}%
\pgfpathlineto{\pgfqpoint{1.392146in}{1.280688in}}%
\pgfpathlineto{\pgfqpoint{1.322867in}{1.287410in}}%
\pgfpathlineto{\pgfqpoint{1.329313in}{1.352796in}}%
\pgfpathlineto{\pgfqpoint{1.335768in}{1.418681in}}%
\pgfpathlineto{\pgfqpoint{1.362400in}{1.415805in}}%
\pgfpathlineto{\pgfqpoint{1.430419in}{1.409379in}}%
\pgfpathlineto{\pgfqpoint{1.471583in}{1.406155in}}%
\pgfpathlineto{\pgfqpoint{1.521419in}{1.403298in}}%
\pgfpathlineto{\pgfqpoint{1.603447in}{1.399416in}}%
\pgfpathlineto{\pgfqpoint{1.605250in}{1.395738in}}%
\pgfpathlineto{\pgfqpoint{1.628754in}{1.384110in}}%
\pgfpathlineto{\pgfqpoint{1.635962in}{1.390016in}}%
\pgfpathlineto{\pgfqpoint{1.656727in}{1.389477in}}%
\pgfpathlineto{\pgfqpoint{1.669102in}{1.382994in}}%
\pgfpathlineto{\pgfqpoint{1.688724in}{1.375625in}}%
\pgfpathlineto{\pgfqpoint{1.694222in}{1.364874in}}%
\pgfpathlineto{\pgfqpoint{1.702281in}{1.363552in}}%
\pgfpathlineto{\pgfqpoint{1.706209in}{1.342308in}}%
\pgfpathlineto{\pgfqpoint{1.719395in}{1.317284in}}%
\pgfpathlineto{\pgfqpoint{1.717525in}{1.304616in}}%
\pgfpathlineto{\pgfqpoint{1.726121in}{1.295511in}}%
\pgfpathlineto{\pgfqpoint{1.725651in}{1.287453in}}%
\pgfpathlineto{\pgfqpoint{1.729695in}{1.272194in}}%
\pgfpathlineto{\pgfqpoint{1.730736in}{1.256440in}}%
\pgfpathlineto{\pgfqpoint{1.728543in}{1.247244in}}%
\pgfpathlineto{\pgfqpoint{1.733914in}{1.237878in}}%
\pgfpathlineto{\pgfqpoint{1.741281in}{1.220848in}}%
\pgfpathlineto{\pgfqpoint{1.748332in}{1.213916in}}%
\pgfpathlineto{\pgfqpoint{1.756737in}{1.198896in}}%
\pgfpathlineto{\pgfqpoint{1.732915in}{1.198649in}}%
\pgfpathlineto{\pgfqpoint{1.681409in}{1.199446in}}%
\pgfpathlineto{\pgfqpoint{1.624502in}{1.201190in}}%
\pgfpathlineto{\pgfqpoint{1.567259in}{1.203386in}}%
\pgfpathlineto{\pgfqpoint{1.483096in}{1.207983in}}%
\pgfpathlineto{\pgfqpoint{1.416825in}{1.212637in}}%
\pgfusepath{stroke}%
\end{pgfscope}%
\begin{pgfscope}%
\pgfpathrectangle{\pgfqpoint{0.100000in}{0.100000in}}{\pgfqpoint{3.007045in}{1.925000in}}%
\pgfusepath{clip}%
\pgfsetbuttcap%
\pgfsetroundjoin%
\pgfsetlinewidth{0.050187pt}%
\definecolor{currentstroke}{rgb}{1.000000,1.000000,1.000000}%
\pgfsetstrokecolor{currentstroke}%
\pgfsetdash{}{0pt}%
\pgfpathmoveto{\pgfqpoint{2.512669in}{1.410567in}}%
\pgfpathlineto{\pgfqpoint{2.539023in}{1.435695in}}%
\pgfpathlineto{\pgfqpoint{2.542299in}{1.444627in}}%
\pgfpathlineto{\pgfqpoint{2.550021in}{1.452274in}}%
\pgfpathlineto{\pgfqpoint{2.544226in}{1.463407in}}%
\pgfpathlineto{\pgfqpoint{2.536962in}{1.469919in}}%
\pgfpathlineto{\pgfqpoint{2.534862in}{1.481459in}}%
\pgfpathlineto{\pgfqpoint{2.561903in}{1.493308in}}%
\pgfpathlineto{\pgfqpoint{2.584291in}{1.497056in}}%
\pgfpathlineto{\pgfqpoint{2.596321in}{1.497306in}}%
\pgfpathlineto{\pgfqpoint{2.605485in}{1.492779in}}%
\pgfpathlineto{\pgfqpoint{2.614418in}{1.496822in}}%
\pgfpathlineto{\pgfqpoint{2.636206in}{1.501353in}}%
\pgfpathlineto{\pgfqpoint{2.643734in}{1.507201in}}%
\pgfpathlineto{\pgfqpoint{2.654897in}{1.520147in}}%
\pgfpathlineto{\pgfqpoint{2.665052in}{1.525868in}}%
\pgfpathlineto{\pgfqpoint{2.665769in}{1.531355in}}%
\pgfpathlineto{\pgfqpoint{2.660439in}{1.543873in}}%
\pgfpathlineto{\pgfqpoint{2.664292in}{1.551240in}}%
\pgfpathlineto{\pgfqpoint{2.659104in}{1.559161in}}%
\pgfpathlineto{\pgfqpoint{2.651159in}{1.559714in}}%
\pgfpathlineto{\pgfqpoint{2.671043in}{1.583641in}}%
\pgfpathlineto{\pgfqpoint{2.673405in}{1.592759in}}%
\pgfpathlineto{\pgfqpoint{2.689054in}{1.616015in}}%
\pgfpathlineto{\pgfqpoint{2.703581in}{1.628599in}}%
\pgfpathlineto{\pgfqpoint{2.713557in}{1.633854in}}%
\pgfpathlineto{\pgfqpoint{2.746198in}{1.641250in}}%
\pgfpathlineto{\pgfqpoint{2.776791in}{1.649602in}}%
\pgfpathlineto{\pgfqpoint{2.780843in}{1.636492in}}%
\pgfpathlineto{\pgfqpoint{2.781130in}{1.624335in}}%
\pgfpathlineto{\pgfqpoint{2.788108in}{1.613300in}}%
\pgfpathlineto{\pgfqpoint{2.790259in}{1.602064in}}%
\pgfpathlineto{\pgfqpoint{2.787577in}{1.586820in}}%
\pgfpathlineto{\pgfqpoint{2.795745in}{1.566258in}}%
\pgfpathlineto{\pgfqpoint{2.805056in}{1.554825in}}%
\pgfpathlineto{\pgfqpoint{2.812258in}{1.517661in}}%
\pgfpathlineto{\pgfqpoint{2.816139in}{1.505410in}}%
\pgfpathlineto{\pgfqpoint{2.815137in}{1.459774in}}%
\pgfpathlineto{\pgfqpoint{2.816297in}{1.458017in}}%
\pgfpathlineto{\pgfqpoint{2.824776in}{1.408935in}}%
\pgfpathlineto{\pgfqpoint{2.829531in}{1.404463in}}%
\pgfpathlineto{\pgfqpoint{2.819303in}{1.394526in}}%
\pgfpathlineto{\pgfqpoint{2.824350in}{1.388822in}}%
\pgfpathlineto{\pgfqpoint{2.819914in}{1.380196in}}%
\pgfpathlineto{\pgfqpoint{2.819952in}{1.376524in}}%
\pgfpathlineto{\pgfqpoint{2.814406in}{1.373210in}}%
\pgfpathlineto{\pgfqpoint{2.811706in}{1.365884in}}%
\pgfpathlineto{\pgfqpoint{2.812563in}{1.386028in}}%
\pgfpathlineto{\pgfqpoint{2.795375in}{1.390462in}}%
\pgfpathlineto{\pgfqpoint{2.768496in}{1.399631in}}%
\pgfpathlineto{\pgfqpoint{2.764695in}{1.404147in}}%
\pgfpathlineto{\pgfqpoint{2.753463in}{1.405229in}}%
\pgfpathlineto{\pgfqpoint{2.747002in}{1.411951in}}%
\pgfpathlineto{\pgfqpoint{2.743516in}{1.425338in}}%
\pgfpathlineto{\pgfqpoint{2.734357in}{1.427003in}}%
\pgfpathlineto{\pgfqpoint{2.728231in}{1.434026in}}%
\pgfpathlineto{\pgfqpoint{2.650253in}{1.418305in}}%
\pgfpathlineto{\pgfqpoint{2.612985in}{1.410597in}}%
\pgfpathlineto{\pgfqpoint{2.556401in}{1.400270in}}%
\pgfpathlineto{\pgfqpoint{2.515654in}{1.393396in}}%
\pgfpathlineto{\pgfqpoint{2.512669in}{1.410567in}}%
\pgfusepath{stroke}%
\end{pgfscope}%
\begin{pgfscope}%
\pgfpathrectangle{\pgfqpoint{0.100000in}{0.100000in}}{\pgfqpoint{3.007045in}{1.925000in}}%
\pgfusepath{clip}%
\pgfsetbuttcap%
\pgfsetroundjoin%
\pgfsetlinewidth{0.050187pt}%
\definecolor{currentstroke}{rgb}{1.000000,1.000000,1.000000}%
\pgfsetstrokecolor{currentstroke}%
\pgfsetdash{}{0pt}%
\pgfpathmoveto{\pgfqpoint{2.825722in}{1.361793in}}%
\pgfpathlineto{\pgfqpoint{2.813662in}{1.358036in}}%
\pgfpathlineto{\pgfqpoint{2.811634in}{1.361490in}}%
\pgfpathlineto{\pgfqpoint{2.815505in}{1.373084in}}%
\pgfpathlineto{\pgfqpoint{2.822673in}{1.374226in}}%
\pgfpathlineto{\pgfqpoint{2.822032in}{1.377876in}}%
\pgfpathlineto{\pgfqpoint{2.828469in}{1.383355in}}%
\pgfpathlineto{\pgfqpoint{2.847078in}{1.387719in}}%
\pgfpathlineto{\pgfqpoint{2.855363in}{1.394327in}}%
\pgfpathlineto{\pgfqpoint{2.873974in}{1.399833in}}%
\pgfpathlineto{\pgfqpoint{2.882462in}{1.397818in}}%
\pgfpathlineto{\pgfqpoint{2.889602in}{1.406695in}}%
\pgfpathlineto{\pgfqpoint{2.900394in}{1.407868in}}%
\pgfpathlineto{\pgfqpoint{2.881983in}{1.390553in}}%
\pgfpathlineto{\pgfqpoint{2.825722in}{1.361793in}}%
\pgfusepath{stroke}%
\end{pgfscope}%
\begin{pgfscope}%
\pgfpathrectangle{\pgfqpoint{0.100000in}{0.100000in}}{\pgfqpoint{3.007045in}{1.925000in}}%
\pgfusepath{clip}%
\pgfsetbuttcap%
\pgfsetroundjoin%
\pgfsetlinewidth{0.050187pt}%
\definecolor{currentstroke}{rgb}{1.000000,1.000000,1.000000}%
\pgfsetstrokecolor{currentstroke}%
\pgfsetdash{}{0pt}%
\pgfpathmoveto{\pgfqpoint{2.554781in}{1.247456in}}%
\pgfpathlineto{\pgfqpoint{2.502630in}{1.238834in}}%
\pgfpathlineto{\pgfqpoint{2.493170in}{1.298432in}}%
\pgfpathlineto{\pgfqpoint{2.479125in}{1.386284in}}%
\pgfpathlineto{\pgfqpoint{2.512669in}{1.410567in}}%
\pgfpathlineto{\pgfqpoint{2.515654in}{1.393396in}}%
\pgfpathlineto{\pgfqpoint{2.556401in}{1.400270in}}%
\pgfpathlineto{\pgfqpoint{2.612985in}{1.410597in}}%
\pgfpathlineto{\pgfqpoint{2.650253in}{1.418305in}}%
\pgfpathlineto{\pgfqpoint{2.728231in}{1.434026in}}%
\pgfpathlineto{\pgfqpoint{2.734357in}{1.427003in}}%
\pgfpathlineto{\pgfqpoint{2.743516in}{1.425338in}}%
\pgfpathlineto{\pgfqpoint{2.747002in}{1.411951in}}%
\pgfpathlineto{\pgfqpoint{2.753463in}{1.405229in}}%
\pgfpathlineto{\pgfqpoint{2.764695in}{1.404147in}}%
\pgfpathlineto{\pgfqpoint{2.768496in}{1.399631in}}%
\pgfpathlineto{\pgfqpoint{2.764637in}{1.396146in}}%
\pgfpathlineto{\pgfqpoint{2.761160in}{1.383818in}}%
\pgfpathlineto{\pgfqpoint{2.753012in}{1.369966in}}%
\pgfpathlineto{\pgfqpoint{2.758481in}{1.363942in}}%
\pgfpathlineto{\pgfqpoint{2.754006in}{1.357488in}}%
\pgfpathlineto{\pgfqpoint{2.756531in}{1.343334in}}%
\pgfpathlineto{\pgfqpoint{2.764590in}{1.335904in}}%
\pgfpathlineto{\pgfqpoint{2.783714in}{1.323833in}}%
\pgfpathlineto{\pgfqpoint{2.768314in}{1.306801in}}%
\pgfpathlineto{\pgfqpoint{2.768098in}{1.300337in}}%
\pgfpathlineto{\pgfqpoint{2.755560in}{1.292000in}}%
\pgfpathlineto{\pgfqpoint{2.741696in}{1.290435in}}%
\pgfpathlineto{\pgfqpoint{2.738266in}{1.283189in}}%
\pgfpathlineto{\pgfqpoint{2.699704in}{1.274845in}}%
\pgfpathlineto{\pgfqpoint{2.654706in}{1.265782in}}%
\pgfpathlineto{\pgfqpoint{2.623776in}{1.260244in}}%
\pgfpathlineto{\pgfqpoint{2.554781in}{1.247456in}}%
\pgfusepath{stroke}%
\end{pgfscope}%
\begin{pgfscope}%
\pgfpathrectangle{\pgfqpoint{0.100000in}{0.100000in}}{\pgfqpoint{3.007045in}{1.925000in}}%
\pgfusepath{clip}%
\pgfsetbuttcap%
\pgfsetroundjoin%
\pgfsetlinewidth{0.050187pt}%
\definecolor{currentstroke}{rgb}{1.000000,1.000000,1.000000}%
\pgfsetstrokecolor{currentstroke}%
\pgfsetdash{}{0pt}%
\pgfpathmoveto{\pgfqpoint{2.816297in}{1.458017in}}%
\pgfpathlineto{\pgfqpoint{2.882673in}{1.473100in}}%
\pgfpathlineto{\pgfqpoint{2.897124in}{1.475380in}}%
\pgfpathlineto{\pgfqpoint{2.906688in}{1.437776in}}%
\pgfpathlineto{\pgfqpoint{2.905174in}{1.431042in}}%
\pgfpathlineto{\pgfqpoint{2.874451in}{1.419148in}}%
\pgfpathlineto{\pgfqpoint{2.856110in}{1.414991in}}%
\pgfpathlineto{\pgfqpoint{2.848316in}{1.405665in}}%
\pgfpathlineto{\pgfqpoint{2.824350in}{1.388822in}}%
\pgfpathlineto{\pgfqpoint{2.819303in}{1.394526in}}%
\pgfpathlineto{\pgfqpoint{2.829531in}{1.404463in}}%
\pgfpathlineto{\pgfqpoint{2.824776in}{1.408935in}}%
\pgfpathlineto{\pgfqpoint{2.816297in}{1.458017in}}%
\pgfusepath{stroke}%
\end{pgfscope}%
\begin{pgfscope}%
\pgfpathrectangle{\pgfqpoint{0.100000in}{0.100000in}}{\pgfqpoint{3.007045in}{1.925000in}}%
\pgfusepath{clip}%
\pgfsetbuttcap%
\pgfsetroundjoin%
\pgfsetlinewidth{0.050187pt}%
\definecolor{currentstroke}{rgb}{1.000000,1.000000,1.000000}%
\pgfsetstrokecolor{currentstroke}%
\pgfsetdash{}{0pt}%
\pgfpathmoveto{\pgfqpoint{2.897124in}{1.475380in}}%
\pgfpathlineto{\pgfqpoint{2.916830in}{1.481216in}}%
\pgfpathlineto{\pgfqpoint{2.922562in}{1.467094in}}%
\pgfpathlineto{\pgfqpoint{2.929135in}{1.463612in}}%
\pgfpathlineto{\pgfqpoint{2.921024in}{1.457664in}}%
\pgfpathlineto{\pgfqpoint{2.922047in}{1.440194in}}%
\pgfpathlineto{\pgfqpoint{2.905174in}{1.431042in}}%
\pgfpathlineto{\pgfqpoint{2.906688in}{1.437776in}}%
\pgfpathlineto{\pgfqpoint{2.897124in}{1.475380in}}%
\pgfusepath{stroke}%
\end{pgfscope}%
\begin{pgfscope}%
\pgfpathrectangle{\pgfqpoint{0.100000in}{0.100000in}}{\pgfqpoint{3.007045in}{1.925000in}}%
\pgfusepath{clip}%
\pgfsetbuttcap%
\pgfsetroundjoin%
\pgfsetlinewidth{0.050187pt}%
\definecolor{currentstroke}{rgb}{1.000000,1.000000,1.000000}%
\pgfsetstrokecolor{currentstroke}%
\pgfsetdash{}{0pt}%
\pgfpathmoveto{\pgfqpoint{2.753338in}{1.285842in}}%
\pgfpathlineto{\pgfqpoint{2.755560in}{1.292000in}}%
\pgfpathlineto{\pgfqpoint{2.768098in}{1.300337in}}%
\pgfpathlineto{\pgfqpoint{2.768314in}{1.306801in}}%
\pgfpathlineto{\pgfqpoint{2.783714in}{1.323833in}}%
\pgfpathlineto{\pgfqpoint{2.764590in}{1.335904in}}%
\pgfpathlineto{\pgfqpoint{2.756531in}{1.343334in}}%
\pgfpathlineto{\pgfqpoint{2.754006in}{1.357488in}}%
\pgfpathlineto{\pgfqpoint{2.758481in}{1.363942in}}%
\pgfpathlineto{\pgfqpoint{2.753012in}{1.369966in}}%
\pgfpathlineto{\pgfqpoint{2.761160in}{1.383818in}}%
\pgfpathlineto{\pgfqpoint{2.764637in}{1.396146in}}%
\pgfpathlineto{\pgfqpoint{2.768496in}{1.399631in}}%
\pgfpathlineto{\pgfqpoint{2.795375in}{1.390462in}}%
\pgfpathlineto{\pgfqpoint{2.812563in}{1.386028in}}%
\pgfpathlineto{\pgfqpoint{2.811706in}{1.365884in}}%
\pgfpathlineto{\pgfqpoint{2.801265in}{1.350609in}}%
\pgfpathlineto{\pgfqpoint{2.809869in}{1.348362in}}%
\pgfpathlineto{\pgfqpoint{2.818801in}{1.341807in}}%
\pgfpathlineto{\pgfqpoint{2.819328in}{1.323878in}}%
\pgfpathlineto{\pgfqpoint{2.816624in}{1.311186in}}%
\pgfpathlineto{\pgfqpoint{2.818423in}{1.300760in}}%
\pgfpathlineto{\pgfqpoint{2.802896in}{1.265581in}}%
\pgfpathlineto{\pgfqpoint{2.797316in}{1.249154in}}%
\pgfpathlineto{\pgfqpoint{2.789524in}{1.257138in}}%
\pgfpathlineto{\pgfqpoint{2.779192in}{1.255778in}}%
\pgfpathlineto{\pgfqpoint{2.753333in}{1.270718in}}%
\pgfpathlineto{\pgfqpoint{2.750687in}{1.278719in}}%
\pgfpathlineto{\pgfqpoint{2.753338in}{1.285842in}}%
\pgfusepath{stroke}%
\end{pgfscope}%
\begin{pgfscope}%
\pgfpathrectangle{\pgfqpoint{0.100000in}{0.100000in}}{\pgfqpoint{3.007045in}{1.925000in}}%
\pgfusepath{clip}%
\pgfsetbuttcap%
\pgfsetroundjoin%
\pgfsetlinewidth{0.050187pt}%
\definecolor{currentstroke}{rgb}{1.000000,1.000000,1.000000}%
\pgfsetstrokecolor{currentstroke}%
\pgfsetdash{}{0pt}%
\pgfpathmoveto{\pgfqpoint{2.134480in}{1.068255in}}%
\pgfpathlineto{\pgfqpoint{2.133271in}{1.084575in}}%
\pgfpathlineto{\pgfqpoint{2.137970in}{1.092672in}}%
\pgfpathlineto{\pgfqpoint{2.134873in}{1.096640in}}%
\pgfpathlineto{\pgfqpoint{2.146167in}{1.109630in}}%
\pgfpathlineto{\pgfqpoint{2.145520in}{1.112760in}}%
\pgfpathlineto{\pgfqpoint{2.156492in}{1.135423in}}%
\pgfpathlineto{\pgfqpoint{2.154275in}{1.146352in}}%
\pgfpathlineto{\pgfqpoint{2.146318in}{1.157795in}}%
\pgfpathlineto{\pgfqpoint{2.151841in}{1.171716in}}%
\pgfpathlineto{\pgfqpoint{2.144682in}{1.263329in}}%
\pgfpathlineto{\pgfqpoint{2.139527in}{1.327598in}}%
\pgfpathlineto{\pgfqpoint{2.146648in}{1.322273in}}%
\pgfpathlineto{\pgfqpoint{2.154593in}{1.322417in}}%
\pgfpathlineto{\pgfqpoint{2.173359in}{1.333287in}}%
\pgfpathlineto{\pgfqpoint{2.230918in}{1.338658in}}%
\pgfpathlineto{\pgfqpoint{2.273522in}{1.343148in}}%
\pgfpathlineto{\pgfqpoint{2.273898in}{1.338979in}}%
\pgfpathlineto{\pgfqpoint{2.283623in}{1.250907in}}%
\pgfpathlineto{\pgfqpoint{2.291994in}{1.168959in}}%
\pgfpathlineto{\pgfqpoint{2.288386in}{1.165112in}}%
\pgfpathlineto{\pgfqpoint{2.293928in}{1.155563in}}%
\pgfpathlineto{\pgfqpoint{2.293900in}{1.148681in}}%
\pgfpathlineto{\pgfqpoint{2.285991in}{1.146952in}}%
\pgfpathlineto{\pgfqpoint{2.277140in}{1.140317in}}%
\pgfpathlineto{\pgfqpoint{2.271148in}{1.142929in}}%
\pgfpathlineto{\pgfqpoint{2.262180in}{1.138682in}}%
\pgfpathlineto{\pgfqpoint{2.264955in}{1.130153in}}%
\pgfpathlineto{\pgfqpoint{2.255736in}{1.121584in}}%
\pgfpathlineto{\pgfqpoint{2.253191in}{1.111671in}}%
\pgfpathlineto{\pgfqpoint{2.246853in}{1.110045in}}%
\pgfpathlineto{\pgfqpoint{2.242123in}{1.102538in}}%
\pgfpathlineto{\pgfqpoint{2.242742in}{1.094978in}}%
\pgfpathlineto{\pgfqpoint{2.237177in}{1.089662in}}%
\pgfpathlineto{\pgfqpoint{2.228803in}{1.090482in}}%
\pgfpathlineto{\pgfqpoint{2.222438in}{1.098640in}}%
\pgfpathlineto{\pgfqpoint{2.211653in}{1.090784in}}%
\pgfpathlineto{\pgfqpoint{2.212410in}{1.083920in}}%
\pgfpathlineto{\pgfqpoint{2.200414in}{1.079910in}}%
\pgfpathlineto{\pgfqpoint{2.197401in}{1.084958in}}%
\pgfpathlineto{\pgfqpoint{2.185723in}{1.079234in}}%
\pgfpathlineto{\pgfqpoint{2.181457in}{1.071040in}}%
\pgfpathlineto{\pgfqpoint{2.167393in}{1.079442in}}%
\pgfpathlineto{\pgfqpoint{2.145103in}{1.073876in}}%
\pgfpathlineto{\pgfqpoint{2.134480in}{1.068255in}}%
\pgfusepath{stroke}%
\end{pgfscope}%
\begin{pgfscope}%
\pgfpathrectangle{\pgfqpoint{0.100000in}{0.100000in}}{\pgfqpoint{3.007045in}{1.925000in}}%
\pgfusepath{clip}%
\pgfsetbuttcap%
\pgfsetroundjoin%
\pgfsetlinewidth{0.050187pt}%
\definecolor{currentstroke}{rgb}{1.000000,1.000000,1.000000}%
\pgfsetstrokecolor{currentstroke}%
\pgfsetdash{}{0pt}%
\pgfpathmoveto{\pgfqpoint{0.562516in}{1.497430in}}%
\pgfpathlineto{\pgfqpoint{0.594748in}{1.488461in}}%
\pgfpathlineto{\pgfqpoint{0.648178in}{1.474684in}}%
\pgfpathlineto{\pgfqpoint{0.703563in}{1.460492in}}%
\pgfpathlineto{\pgfqpoint{0.754611in}{1.448181in}}%
\pgfpathlineto{\pgfqpoint{0.798720in}{1.438195in}}%
\pgfpathlineto{\pgfqpoint{0.845912in}{1.427818in}}%
\pgfpathlineto{\pgfqpoint{0.832277in}{1.363472in}}%
\pgfpathlineto{\pgfqpoint{0.820130in}{1.306329in}}%
\pgfpathlineto{\pgfqpoint{0.800203in}{1.214130in}}%
\pgfpathlineto{\pgfqpoint{0.785243in}{1.144539in}}%
\pgfpathlineto{\pgfqpoint{0.777160in}{1.105704in}}%
\pgfpathlineto{\pgfqpoint{0.766797in}{1.055300in}}%
\pgfpathlineto{\pgfqpoint{0.759627in}{1.045075in}}%
\pgfpathlineto{\pgfqpoint{0.753855in}{1.044733in}}%
\pgfpathlineto{\pgfqpoint{0.749739in}{1.053659in}}%
\pgfpathlineto{\pgfqpoint{0.740288in}{1.056900in}}%
\pgfpathlineto{\pgfqpoint{0.730103in}{1.055764in}}%
\pgfpathlineto{\pgfqpoint{0.727212in}{1.048463in}}%
\pgfpathlineto{\pgfqpoint{0.726911in}{1.023110in}}%
\pgfpathlineto{\pgfqpoint{0.723861in}{1.017321in}}%
\pgfpathlineto{\pgfqpoint{0.725604in}{0.996967in}}%
\pgfpathlineto{\pgfqpoint{0.719235in}{0.983407in}}%
\pgfpathlineto{\pgfqpoint{0.667589in}{1.062934in}}%
\pgfpathlineto{\pgfqpoint{0.616726in}{1.140213in}}%
\pgfpathlineto{\pgfqpoint{0.590250in}{1.180736in}}%
\pgfpathlineto{\pgfqpoint{0.568163in}{1.215713in}}%
\pgfpathlineto{\pgfqpoint{0.540333in}{1.258980in}}%
\pgfpathlineto{\pgfqpoint{0.508830in}{1.307514in}}%
\pgfpathlineto{\pgfqpoint{0.521782in}{1.353580in}}%
\pgfpathlineto{\pgfqpoint{0.547847in}{1.445968in}}%
\pgfpathlineto{\pgfqpoint{0.562516in}{1.497430in}}%
\pgfusepath{stroke}%
\end{pgfscope}%
\begin{pgfscope}%
\pgfpathrectangle{\pgfqpoint{0.100000in}{0.100000in}}{\pgfqpoint{3.007045in}{1.925000in}}%
\pgfusepath{clip}%
\pgfsetbuttcap%
\pgfsetroundjoin%
\pgfsetlinewidth{0.050187pt}%
\definecolor{currentstroke}{rgb}{1.000000,1.000000,1.000000}%
\pgfsetstrokecolor{currentstroke}%
\pgfsetdash{}{0pt}%
\pgfpathmoveto{\pgfqpoint{0.845912in}{1.427818in}}%
\pgfpathlineto{\pgfqpoint{0.896364in}{1.417782in}}%
\pgfpathlineto{\pgfqpoint{0.989767in}{1.399751in}}%
\pgfpathlineto{\pgfqpoint{0.978076in}{1.334895in}}%
\pgfpathlineto{\pgfqpoint{1.029497in}{1.326091in}}%
\pgfpathlineto{\pgfqpoint{1.076345in}{1.318594in}}%
\pgfpathlineto{\pgfqpoint{1.068205in}{1.267320in}}%
\pgfpathlineto{\pgfqpoint{1.059579in}{1.212028in}}%
\pgfpathlineto{\pgfqpoint{1.048031in}{1.139428in}}%
\pgfpathlineto{\pgfqpoint{1.047732in}{1.133342in}}%
\pgfpathlineto{\pgfqpoint{1.035744in}{1.058103in}}%
\pgfpathlineto{\pgfqpoint{0.986397in}{1.065751in}}%
\pgfpathlineto{\pgfqpoint{0.961251in}{1.070801in}}%
\pgfpathlineto{\pgfqpoint{0.870361in}{1.086778in}}%
\pgfpathlineto{\pgfqpoint{0.836134in}{1.093528in}}%
\pgfpathlineto{\pgfqpoint{0.777160in}{1.105704in}}%
\pgfpathlineto{\pgfqpoint{0.785243in}{1.144539in}}%
\pgfpathlineto{\pgfqpoint{0.800203in}{1.214130in}}%
\pgfpathlineto{\pgfqpoint{0.820130in}{1.306329in}}%
\pgfpathlineto{\pgfqpoint{0.832277in}{1.363472in}}%
\pgfpathlineto{\pgfqpoint{0.845912in}{1.427818in}}%
\pgfusepath{stroke}%
\end{pgfscope}%
\begin{pgfscope}%
\pgfpathrectangle{\pgfqpoint{0.100000in}{0.100000in}}{\pgfqpoint{3.007045in}{1.925000in}}%
\pgfusepath{clip}%
\pgfsetbuttcap%
\pgfsetroundjoin%
\pgfsetlinewidth{0.050187pt}%
\definecolor{currentstroke}{rgb}{1.000000,1.000000,1.000000}%
\pgfsetstrokecolor{currentstroke}%
\pgfsetdash{}{0pt}%
\pgfpathmoveto{\pgfqpoint{0.365302in}{1.558630in}}%
\pgfpathlineto{\pgfqpoint{0.383230in}{1.552378in}}%
\pgfpathlineto{\pgfqpoint{0.455041in}{1.530020in}}%
\pgfpathlineto{\pgfqpoint{0.521101in}{1.509228in}}%
\pgfpathlineto{\pgfqpoint{0.562516in}{1.497430in}}%
\pgfpathlineto{\pgfqpoint{0.547847in}{1.445968in}}%
\pgfpathlineto{\pgfqpoint{0.521782in}{1.353580in}}%
\pgfpathlineto{\pgfqpoint{0.508830in}{1.307514in}}%
\pgfpathlineto{\pgfqpoint{0.540333in}{1.258980in}}%
\pgfpathlineto{\pgfqpoint{0.568163in}{1.215713in}}%
\pgfpathlineto{\pgfqpoint{0.590250in}{1.180736in}}%
\pgfpathlineto{\pgfqpoint{0.616726in}{1.140213in}}%
\pgfpathlineto{\pgfqpoint{0.667589in}{1.062934in}}%
\pgfpathlineto{\pgfqpoint{0.719235in}{0.983407in}}%
\pgfpathlineto{\pgfqpoint{0.717162in}{0.975520in}}%
\pgfpathlineto{\pgfqpoint{0.723391in}{0.962960in}}%
\pgfpathlineto{\pgfqpoint{0.724623in}{0.945778in}}%
\pgfpathlineto{\pgfqpoint{0.733671in}{0.937444in}}%
\pgfpathlineto{\pgfqpoint{0.734029in}{0.930728in}}%
\pgfpathlineto{\pgfqpoint{0.717829in}{0.923106in}}%
\pgfpathlineto{\pgfqpoint{0.710111in}{0.915475in}}%
\pgfpathlineto{\pgfqpoint{0.707704in}{0.898624in}}%
\pgfpathlineto{\pgfqpoint{0.703798in}{0.889434in}}%
\pgfpathlineto{\pgfqpoint{0.695567in}{0.881688in}}%
\pgfpathlineto{\pgfqpoint{0.687380in}{0.861545in}}%
\pgfpathlineto{\pgfqpoint{0.698891in}{0.851052in}}%
\pgfpathlineto{\pgfqpoint{0.697405in}{0.842403in}}%
\pgfpathlineto{\pgfqpoint{0.687967in}{0.836384in}}%
\pgfpathlineto{\pgfqpoint{0.681459in}{0.837446in}}%
\pgfpathlineto{\pgfqpoint{0.604646in}{0.848106in}}%
\pgfpathlineto{\pgfqpoint{0.547913in}{0.856154in}}%
\pgfpathlineto{\pgfqpoint{0.550394in}{0.865314in}}%
\pgfpathlineto{\pgfqpoint{0.546700in}{0.880521in}}%
\pgfpathlineto{\pgfqpoint{0.546328in}{0.895860in}}%
\pgfpathlineto{\pgfqpoint{0.543914in}{0.904837in}}%
\pgfpathlineto{\pgfqpoint{0.536479in}{0.917671in}}%
\pgfpathlineto{\pgfqpoint{0.515092in}{0.947282in}}%
\pgfpathlineto{\pgfqpoint{0.508088in}{0.950910in}}%
\pgfpathlineto{\pgfqpoint{0.499063in}{0.950853in}}%
\pgfpathlineto{\pgfqpoint{0.501124in}{0.960194in}}%
\pgfpathlineto{\pgfqpoint{0.496851in}{0.971875in}}%
\pgfpathlineto{\pgfqpoint{0.475860in}{0.977681in}}%
\pgfpathlineto{\pgfqpoint{0.463064in}{0.988432in}}%
\pgfpathlineto{\pgfqpoint{0.462003in}{0.995009in}}%
\pgfpathlineto{\pgfqpoint{0.447270in}{1.011295in}}%
\pgfpathlineto{\pgfqpoint{0.433245in}{1.014411in}}%
\pgfpathlineto{\pgfqpoint{0.420251in}{1.022687in}}%
\pgfpathlineto{\pgfqpoint{0.403161in}{1.025551in}}%
\pgfpathlineto{\pgfqpoint{0.395859in}{1.036589in}}%
\pgfpathlineto{\pgfqpoint{0.402794in}{1.054061in}}%
\pgfpathlineto{\pgfqpoint{0.400673in}{1.057990in}}%
\pgfpathlineto{\pgfqpoint{0.406447in}{1.072555in}}%
\pgfpathlineto{\pgfqpoint{0.396175in}{1.080306in}}%
\pgfpathlineto{\pgfqpoint{0.399508in}{1.094359in}}%
\pgfpathlineto{\pgfqpoint{0.394021in}{1.097951in}}%
\pgfpathlineto{\pgfqpoint{0.389283in}{1.111233in}}%
\pgfpathlineto{\pgfqpoint{0.383566in}{1.115284in}}%
\pgfpathlineto{\pgfqpoint{0.383136in}{1.124877in}}%
\pgfpathlineto{\pgfqpoint{0.378657in}{1.131642in}}%
\pgfpathlineto{\pgfqpoint{0.371908in}{1.154426in}}%
\pgfpathlineto{\pgfqpoint{0.364523in}{1.165339in}}%
\pgfpathlineto{\pgfqpoint{0.366128in}{1.183864in}}%
\pgfpathlineto{\pgfqpoint{0.374812in}{1.185728in}}%
\pgfpathlineto{\pgfqpoint{0.380467in}{1.195776in}}%
\pgfpathlineto{\pgfqpoint{0.377050in}{1.206674in}}%
\pgfpathlineto{\pgfqpoint{0.367814in}{1.208493in}}%
\pgfpathlineto{\pgfqpoint{0.363206in}{1.213591in}}%
\pgfpathlineto{\pgfqpoint{0.355744in}{1.232340in}}%
\pgfpathlineto{\pgfqpoint{0.359231in}{1.239073in}}%
\pgfpathlineto{\pgfqpoint{0.356757in}{1.251611in}}%
\pgfpathlineto{\pgfqpoint{0.362235in}{1.267849in}}%
\pgfpathlineto{\pgfqpoint{0.368652in}{1.261869in}}%
\pgfpathlineto{\pgfqpoint{0.365740in}{1.254809in}}%
\pgfpathlineto{\pgfqpoint{0.376861in}{1.243618in}}%
\pgfpathlineto{\pgfqpoint{0.376151in}{1.260237in}}%
\pgfpathlineto{\pgfqpoint{0.371386in}{1.264694in}}%
\pgfpathlineto{\pgfqpoint{0.376833in}{1.279306in}}%
\pgfpathlineto{\pgfqpoint{0.397434in}{1.276979in}}%
\pgfpathlineto{\pgfqpoint{0.394662in}{1.282413in}}%
\pgfpathlineto{\pgfqpoint{0.381065in}{1.281878in}}%
\pgfpathlineto{\pgfqpoint{0.374596in}{1.290108in}}%
\pgfpathlineto{\pgfqpoint{0.366455in}{1.282800in}}%
\pgfpathlineto{\pgfqpoint{0.362134in}{1.270584in}}%
\pgfpathlineto{\pgfqpoint{0.350603in}{1.287022in}}%
\pgfpathlineto{\pgfqpoint{0.346162in}{1.290014in}}%
\pgfpathlineto{\pgfqpoint{0.347854in}{1.307904in}}%
\pgfpathlineto{\pgfqpoint{0.344340in}{1.318454in}}%
\pgfpathlineto{\pgfqpoint{0.337968in}{1.328366in}}%
\pgfpathlineto{\pgfqpoint{0.324921in}{1.358735in}}%
\pgfpathlineto{\pgfqpoint{0.329187in}{1.365451in}}%
\pgfpathlineto{\pgfqpoint{0.329137in}{1.386675in}}%
\pgfpathlineto{\pgfqpoint{0.336178in}{1.398512in}}%
\pgfpathlineto{\pgfqpoint{0.337793in}{1.417013in}}%
\pgfpathlineto{\pgfqpoint{0.331153in}{1.438209in}}%
\pgfpathlineto{\pgfqpoint{0.322344in}{1.451701in}}%
\pgfpathlineto{\pgfqpoint{0.323915in}{1.463876in}}%
\pgfpathlineto{\pgfqpoint{0.348653in}{1.493342in}}%
\pgfpathlineto{\pgfqpoint{0.349881in}{1.503394in}}%
\pgfpathlineto{\pgfqpoint{0.361028in}{1.522568in}}%
\pgfpathlineto{\pgfqpoint{0.362575in}{1.540742in}}%
\pgfpathlineto{\pgfqpoint{0.358991in}{1.545379in}}%
\pgfpathlineto{\pgfqpoint{0.365302in}{1.558630in}}%
\pgfusepath{stroke}%
\end{pgfscope}%
\begin{pgfscope}%
\pgfpathrectangle{\pgfqpoint{0.100000in}{0.100000in}}{\pgfqpoint{3.007045in}{1.925000in}}%
\pgfusepath{clip}%
\pgfsetbuttcap%
\pgfsetroundjoin%
\pgfsetlinewidth{0.050187pt}%
\definecolor{currentstroke}{rgb}{1.000000,1.000000,1.000000}%
\pgfsetstrokecolor{currentstroke}%
\pgfsetdash{}{0pt}%
\pgfpathmoveto{\pgfqpoint{2.291994in}{1.168959in}}%
\pgfpathlineto{\pgfqpoint{2.283623in}{1.250907in}}%
\pgfpathlineto{\pgfqpoint{2.273898in}{1.338979in}}%
\pgfpathlineto{\pgfqpoint{2.337587in}{1.348464in}}%
\pgfpathlineto{\pgfqpoint{2.354491in}{1.344065in}}%
\pgfpathlineto{\pgfqpoint{2.362598in}{1.339289in}}%
\pgfpathlineto{\pgfqpoint{2.372757in}{1.340614in}}%
\pgfpathlineto{\pgfqpoint{2.386151in}{1.332707in}}%
\pgfpathlineto{\pgfqpoint{2.411100in}{1.344475in}}%
\pgfpathlineto{\pgfqpoint{2.424868in}{1.344868in}}%
\pgfpathlineto{\pgfqpoint{2.440946in}{1.362805in}}%
\pgfpathlineto{\pgfqpoint{2.457305in}{1.373716in}}%
\pgfpathlineto{\pgfqpoint{2.479125in}{1.386284in}}%
\pgfpathlineto{\pgfqpoint{2.493170in}{1.298432in}}%
\pgfpathlineto{\pgfqpoint{2.486644in}{1.292790in}}%
\pgfpathlineto{\pgfqpoint{2.490858in}{1.287608in}}%
\pgfpathlineto{\pgfqpoint{2.492536in}{1.276247in}}%
\pgfpathlineto{\pgfqpoint{2.489267in}{1.265577in}}%
\pgfpathlineto{\pgfqpoint{2.488835in}{1.249980in}}%
\pgfpathlineto{\pgfqpoint{2.485764in}{1.229680in}}%
\pgfpathlineto{\pgfqpoint{2.470230in}{1.211575in}}%
\pgfpathlineto{\pgfqpoint{2.463690in}{1.207732in}}%
\pgfpathlineto{\pgfqpoint{2.458605in}{1.211041in}}%
\pgfpathlineto{\pgfqpoint{2.446020in}{1.193784in}}%
\pgfpathlineto{\pgfqpoint{2.447199in}{1.177168in}}%
\pgfpathlineto{\pgfqpoint{2.433183in}{1.181164in}}%
\pgfpathlineto{\pgfqpoint{2.426546in}{1.164570in}}%
\pgfpathlineto{\pgfqpoint{2.429901in}{1.152811in}}%
\pgfpathlineto{\pgfqpoint{2.424649in}{1.151074in}}%
\pgfpathlineto{\pgfqpoint{2.423914in}{1.141782in}}%
\pgfpathlineto{\pgfqpoint{2.411023in}{1.138034in}}%
\pgfpathlineto{\pgfqpoint{2.404325in}{1.145536in}}%
\pgfpathlineto{\pgfqpoint{2.395706in}{1.148445in}}%
\pgfpathlineto{\pgfqpoint{2.393680in}{1.156052in}}%
\pgfpathlineto{\pgfqpoint{2.385887in}{1.154714in}}%
\pgfpathlineto{\pgfqpoint{2.380782in}{1.147717in}}%
\pgfpathlineto{\pgfqpoint{2.373478in}{1.145257in}}%
\pgfpathlineto{\pgfqpoint{2.360576in}{1.150208in}}%
\pgfpathlineto{\pgfqpoint{2.353440in}{1.144317in}}%
\pgfpathlineto{\pgfqpoint{2.343298in}{1.151287in}}%
\pgfpathlineto{\pgfqpoint{2.323842in}{1.153426in}}%
\pgfpathlineto{\pgfqpoint{2.317973in}{1.166088in}}%
\pgfpathlineto{\pgfqpoint{2.310538in}{1.171698in}}%
\pgfpathlineto{\pgfqpoint{2.303330in}{1.168107in}}%
\pgfpathlineto{\pgfqpoint{2.291994in}{1.168959in}}%
\pgfusepath{stroke}%
\end{pgfscope}%
\begin{pgfscope}%
\pgfpathrectangle{\pgfqpoint{0.100000in}{0.100000in}}{\pgfqpoint{3.007045in}{1.925000in}}%
\pgfusepath{clip}%
\pgfsetbuttcap%
\pgfsetroundjoin%
\pgfsetlinewidth{0.050187pt}%
\definecolor{currentstroke}{rgb}{1.000000,1.000000,1.000000}%
\pgfsetstrokecolor{currentstroke}%
\pgfsetdash{}{0pt}%
\pgfpathmoveto{\pgfqpoint{2.081137in}{1.010475in}}%
\pgfpathlineto{\pgfqpoint{2.073675in}{1.016537in}}%
\pgfpathlineto{\pgfqpoint{2.067590in}{1.013653in}}%
\pgfpathlineto{\pgfqpoint{2.059812in}{1.028177in}}%
\pgfpathlineto{\pgfqpoint{2.063786in}{1.037307in}}%
\pgfpathlineto{\pgfqpoint{2.058058in}{1.047582in}}%
\pgfpathlineto{\pgfqpoint{2.057751in}{1.055672in}}%
\pgfpathlineto{\pgfqpoint{2.050012in}{1.058555in}}%
\pgfpathlineto{\pgfqpoint{2.046423in}{1.064650in}}%
\pgfpathlineto{\pgfqpoint{2.031291in}{1.072263in}}%
\pgfpathlineto{\pgfqpoint{2.018150in}{1.081642in}}%
\pgfpathlineto{\pgfqpoint{2.012044in}{1.088723in}}%
\pgfpathlineto{\pgfqpoint{2.011918in}{1.097359in}}%
\pgfpathlineto{\pgfqpoint{2.020088in}{1.113940in}}%
\pgfpathlineto{\pgfqpoint{2.022602in}{1.126623in}}%
\pgfpathlineto{\pgfqpoint{2.015945in}{1.133789in}}%
\pgfpathlineto{\pgfqpoint{2.007129in}{1.136489in}}%
\pgfpathlineto{\pgfqpoint{1.996433in}{1.130585in}}%
\pgfpathlineto{\pgfqpoint{1.991906in}{1.140721in}}%
\pgfpathlineto{\pgfqpoint{1.989621in}{1.154462in}}%
\pgfpathlineto{\pgfqpoint{1.973858in}{1.166711in}}%
\pgfpathlineto{\pgfqpoint{1.970708in}{1.172146in}}%
\pgfpathlineto{\pgfqpoint{1.956311in}{1.184451in}}%
\pgfpathlineto{\pgfqpoint{1.951798in}{1.193397in}}%
\pgfpathlineto{\pgfqpoint{1.947730in}{1.211138in}}%
\pgfpathlineto{\pgfqpoint{1.950519in}{1.226895in}}%
\pgfpathlineto{\pgfqpoint{1.954240in}{1.229091in}}%
\pgfpathlineto{\pgfqpoint{1.953571in}{1.242282in}}%
\pgfpathlineto{\pgfqpoint{1.964076in}{1.246200in}}%
\pgfpathlineto{\pgfqpoint{1.967249in}{1.258037in}}%
\pgfpathlineto{\pgfqpoint{1.973292in}{1.266004in}}%
\pgfpathlineto{\pgfqpoint{1.972987in}{1.276140in}}%
\pgfpathlineto{\pgfqpoint{1.965496in}{1.284201in}}%
\pgfpathlineto{\pgfqpoint{1.967275in}{1.295485in}}%
\pgfpathlineto{\pgfqpoint{1.986693in}{1.300392in}}%
\pgfpathlineto{\pgfqpoint{2.001591in}{1.309350in}}%
\pgfpathlineto{\pgfqpoint{2.003154in}{1.320629in}}%
\pgfpathlineto{\pgfqpoint{2.008338in}{1.324180in}}%
\pgfpathlineto{\pgfqpoint{2.010321in}{1.336029in}}%
\pgfpathlineto{\pgfqpoint{2.008727in}{1.343855in}}%
\pgfpathlineto{\pgfqpoint{1.998543in}{1.350356in}}%
\pgfpathlineto{\pgfqpoint{1.994442in}{1.360047in}}%
\pgfpathlineto{\pgfqpoint{1.984391in}{1.369415in}}%
\pgfpathlineto{\pgfqpoint{2.066992in}{1.372935in}}%
\pgfpathlineto{\pgfqpoint{2.122408in}{1.376872in}}%
\pgfpathlineto{\pgfqpoint{2.121395in}{1.365215in}}%
\pgfpathlineto{\pgfqpoint{2.130839in}{1.349136in}}%
\pgfpathlineto{\pgfqpoint{2.134804in}{1.335398in}}%
\pgfpathlineto{\pgfqpoint{2.139527in}{1.327598in}}%
\pgfpathlineto{\pgfqpoint{2.144682in}{1.263329in}}%
\pgfpathlineto{\pgfqpoint{2.151841in}{1.171716in}}%
\pgfpathlineto{\pgfqpoint{2.146318in}{1.157795in}}%
\pgfpathlineto{\pgfqpoint{2.154275in}{1.146352in}}%
\pgfpathlineto{\pgfqpoint{2.156492in}{1.135423in}}%
\pgfpathlineto{\pgfqpoint{2.145520in}{1.112760in}}%
\pgfpathlineto{\pgfqpoint{2.146167in}{1.109630in}}%
\pgfpathlineto{\pgfqpoint{2.134873in}{1.096640in}}%
\pgfpathlineto{\pgfqpoint{2.137970in}{1.092672in}}%
\pgfpathlineto{\pgfqpoint{2.133271in}{1.084575in}}%
\pgfpathlineto{\pgfqpoint{2.134480in}{1.068255in}}%
\pgfpathlineto{\pgfqpoint{2.128773in}{1.058240in}}%
\pgfpathlineto{\pgfqpoint{2.133415in}{1.046407in}}%
\pgfpathlineto{\pgfqpoint{2.113967in}{1.039976in}}%
\pgfpathlineto{\pgfqpoint{2.112175in}{1.032983in}}%
\pgfpathlineto{\pgfqpoint{2.117486in}{1.024119in}}%
\pgfpathlineto{\pgfqpoint{2.115042in}{1.018342in}}%
\pgfpathlineto{\pgfqpoint{2.094189in}{1.025479in}}%
\pgfpathlineto{\pgfqpoint{2.087308in}{1.026197in}}%
\pgfpathlineto{\pgfqpoint{2.078730in}{1.015347in}}%
\pgfpathlineto{\pgfqpoint{2.081137in}{1.010475in}}%
\pgfusepath{stroke}%
\end{pgfscope}%
\begin{pgfscope}%
\pgfpathrectangle{\pgfqpoint{0.100000in}{0.100000in}}{\pgfqpoint{3.007045in}{1.925000in}}%
\pgfusepath{clip}%
\pgfsetbuttcap%
\pgfsetroundjoin%
\pgfsetlinewidth{0.050187pt}%
\definecolor{currentstroke}{rgb}{1.000000,1.000000,1.000000}%
\pgfsetstrokecolor{currentstroke}%
\pgfsetdash{}{0pt}%
\pgfpathmoveto{\pgfqpoint{2.688153in}{1.209860in}}%
\pgfpathlineto{\pgfqpoint{2.682414in}{1.218382in}}%
\pgfpathlineto{\pgfqpoint{2.685652in}{1.223151in}}%
\pgfpathlineto{\pgfqpoint{2.693583in}{1.217795in}}%
\pgfpathlineto{\pgfqpoint{2.688153in}{1.209860in}}%
\pgfusepath{stroke}%
\end{pgfscope}%
\begin{pgfscope}%
\pgfpathrectangle{\pgfqpoint{0.100000in}{0.100000in}}{\pgfqpoint{3.007045in}{1.925000in}}%
\pgfusepath{clip}%
\pgfsetbuttcap%
\pgfsetroundjoin%
\pgfsetlinewidth{0.050187pt}%
\definecolor{currentstroke}{rgb}{1.000000,1.000000,1.000000}%
\pgfsetstrokecolor{currentstroke}%
\pgfsetdash{}{0pt}%
\pgfpathmoveto{\pgfqpoint{2.738266in}{1.283189in}}%
\pgfpathlineto{\pgfqpoint{2.741696in}{1.290435in}}%
\pgfpathlineto{\pgfqpoint{2.755560in}{1.292000in}}%
\pgfpathlineto{\pgfqpoint{2.753338in}{1.285842in}}%
\pgfpathlineto{\pgfqpoint{2.748764in}{1.277976in}}%
\pgfpathlineto{\pgfqpoint{2.751863in}{1.268604in}}%
\pgfpathlineto{\pgfqpoint{2.764095in}{1.257374in}}%
\pgfpathlineto{\pgfqpoint{2.766934in}{1.245549in}}%
\pgfpathlineto{\pgfqpoint{2.781027in}{1.230823in}}%
\pgfpathlineto{\pgfqpoint{2.786556in}{1.231453in}}%
\pgfpathlineto{\pgfqpoint{2.793443in}{1.209347in}}%
\pgfpathlineto{\pgfqpoint{2.792313in}{1.209127in}}%
\pgfpathlineto{\pgfqpoint{2.791057in}{1.208880in}}%
\pgfpathlineto{\pgfqpoint{2.760342in}{1.203000in}}%
\pgfpathlineto{\pgfqpoint{2.743908in}{1.261445in}}%
\pgfpathlineto{\pgfqpoint{2.738266in}{1.283189in}}%
\pgfusepath{stroke}%
\end{pgfscope}%
\begin{pgfscope}%
\pgfpathrectangle{\pgfqpoint{0.100000in}{0.100000in}}{\pgfqpoint{3.007045in}{1.925000in}}%
\pgfusepath{clip}%
\pgfsetbuttcap%
\pgfsetroundjoin%
\pgfsetlinewidth{0.050187pt}%
\definecolor{currentstroke}{rgb}{1.000000,1.000000,1.000000}%
\pgfsetstrokecolor{currentstroke}%
\pgfsetdash{}{0pt}%
\pgfpathmoveto{\pgfqpoint{2.451373in}{1.085086in}}%
\pgfpathlineto{\pgfqpoint{2.441667in}{1.085443in}}%
\pgfpathlineto{\pgfqpoint{2.432723in}{1.091604in}}%
\pgfpathlineto{\pgfqpoint{2.420641in}{1.110328in}}%
\pgfpathlineto{\pgfqpoint{2.410356in}{1.120244in}}%
\pgfpathlineto{\pgfqpoint{2.413044in}{1.127903in}}%
\pgfpathlineto{\pgfqpoint{2.411023in}{1.138034in}}%
\pgfpathlineto{\pgfqpoint{2.423914in}{1.141782in}}%
\pgfpathlineto{\pgfqpoint{2.424649in}{1.151074in}}%
\pgfpathlineto{\pgfqpoint{2.429901in}{1.152811in}}%
\pgfpathlineto{\pgfqpoint{2.426546in}{1.164570in}}%
\pgfpathlineto{\pgfqpoint{2.433183in}{1.181164in}}%
\pgfpathlineto{\pgfqpoint{2.447199in}{1.177168in}}%
\pgfpathlineto{\pgfqpoint{2.446020in}{1.193784in}}%
\pgfpathlineto{\pgfqpoint{2.458605in}{1.211041in}}%
\pgfpathlineto{\pgfqpoint{2.463690in}{1.207732in}}%
\pgfpathlineto{\pgfqpoint{2.470230in}{1.211575in}}%
\pgfpathlineto{\pgfqpoint{2.485764in}{1.229680in}}%
\pgfpathlineto{\pgfqpoint{2.488835in}{1.249980in}}%
\pgfpathlineto{\pgfqpoint{2.489267in}{1.265577in}}%
\pgfpathlineto{\pgfqpoint{2.492536in}{1.276247in}}%
\pgfpathlineto{\pgfqpoint{2.490858in}{1.287608in}}%
\pgfpathlineto{\pgfqpoint{2.486644in}{1.292790in}}%
\pgfpathlineto{\pgfqpoint{2.493170in}{1.298432in}}%
\pgfpathlineto{\pgfqpoint{2.502630in}{1.238834in}}%
\pgfpathlineto{\pgfqpoint{2.554781in}{1.247456in}}%
\pgfpathlineto{\pgfqpoint{2.560186in}{1.213439in}}%
\pgfpathlineto{\pgfqpoint{2.579067in}{1.235892in}}%
\pgfpathlineto{\pgfqpoint{2.583515in}{1.233653in}}%
\pgfpathlineto{\pgfqpoint{2.590060in}{1.246663in}}%
\pgfpathlineto{\pgfqpoint{2.599126in}{1.242584in}}%
\pgfpathlineto{\pgfqpoint{2.607031in}{1.243360in}}%
\pgfpathlineto{\pgfqpoint{2.612246in}{1.252404in}}%
\pgfpathlineto{\pgfqpoint{2.619817in}{1.257435in}}%
\pgfpathlineto{\pgfqpoint{2.630342in}{1.253004in}}%
\pgfpathlineto{\pgfqpoint{2.637267in}{1.254535in}}%
\pgfpathlineto{\pgfqpoint{2.647188in}{1.237365in}}%
\pgfpathlineto{\pgfqpoint{2.644319in}{1.224366in}}%
\pgfpathlineto{\pgfqpoint{2.618406in}{1.239002in}}%
\pgfpathlineto{\pgfqpoint{2.613553in}{1.226987in}}%
\pgfpathlineto{\pgfqpoint{2.615154in}{1.221459in}}%
\pgfpathlineto{\pgfqpoint{2.604125in}{1.202772in}}%
\pgfpathlineto{\pgfqpoint{2.598920in}{1.199055in}}%
\pgfpathlineto{\pgfqpoint{2.596570in}{1.190815in}}%
\pgfpathlineto{\pgfqpoint{2.587696in}{1.191677in}}%
\pgfpathlineto{\pgfqpoint{2.581324in}{1.169185in}}%
\pgfpathlineto{\pgfqpoint{2.577757in}{1.164024in}}%
\pgfpathlineto{\pgfqpoint{2.568590in}{1.165745in}}%
\pgfpathlineto{\pgfqpoint{2.559215in}{1.172839in}}%
\pgfpathlineto{\pgfqpoint{2.558890in}{1.161951in}}%
\pgfpathlineto{\pgfqpoint{2.549814in}{1.143653in}}%
\pgfpathlineto{\pgfqpoint{2.547556in}{1.130626in}}%
\pgfpathlineto{\pgfqpoint{2.540593in}{1.121950in}}%
\pgfpathlineto{\pgfqpoint{2.535313in}{1.108096in}}%
\pgfpathlineto{\pgfqpoint{2.535053in}{1.095270in}}%
\pgfpathlineto{\pgfqpoint{2.526900in}{1.092874in}}%
\pgfpathlineto{\pgfqpoint{2.517650in}{1.085614in}}%
\pgfpathlineto{\pgfqpoint{2.508744in}{1.084231in}}%
\pgfpathlineto{\pgfqpoint{2.506692in}{1.078086in}}%
\pgfpathlineto{\pgfqpoint{2.492341in}{1.071792in}}%
\pgfpathlineto{\pgfqpoint{2.484322in}{1.077155in}}%
\pgfpathlineto{\pgfqpoint{2.475358in}{1.066973in}}%
\pgfpathlineto{\pgfqpoint{2.459950in}{1.069974in}}%
\pgfpathlineto{\pgfqpoint{2.450500in}{1.080661in}}%
\pgfpathlineto{\pgfqpoint{2.451373in}{1.085086in}}%
\pgfusepath{stroke}%
\end{pgfscope}%
\begin{pgfscope}%
\pgfpathrectangle{\pgfqpoint{0.100000in}{0.100000in}}{\pgfqpoint{3.007045in}{1.925000in}}%
\pgfusepath{clip}%
\pgfsetbuttcap%
\pgfsetroundjoin%
\pgfsetlinewidth{0.050187pt}%
\definecolor{currentstroke}{rgb}{1.000000,1.000000,1.000000}%
\pgfsetstrokecolor{currentstroke}%
\pgfsetdash{}{0pt}%
\pgfpathmoveto{\pgfqpoint{2.791057in}{1.208880in}}%
\pgfpathlineto{\pgfqpoint{2.790664in}{1.196870in}}%
\pgfpathlineto{\pgfqpoint{2.786042in}{1.190966in}}%
\pgfpathlineto{\pgfqpoint{2.783092in}{1.177859in}}%
\pgfpathlineto{\pgfqpoint{2.769772in}{1.171817in}}%
\pgfpathlineto{\pgfqpoint{2.758597in}{1.170057in}}%
\pgfpathlineto{\pgfqpoint{2.761853in}{1.178677in}}%
\pgfpathlineto{\pgfqpoint{2.744646in}{1.185900in}}%
\pgfpathlineto{\pgfqpoint{2.730667in}{1.194944in}}%
\pgfpathlineto{\pgfqpoint{2.739321in}{1.208462in}}%
\pgfpathlineto{\pgfqpoint{2.732385in}{1.218971in}}%
\pgfpathlineto{\pgfqpoint{2.733202in}{1.228035in}}%
\pgfpathlineto{\pgfqpoint{2.728187in}{1.230522in}}%
\pgfpathlineto{\pgfqpoint{2.724117in}{1.245248in}}%
\pgfpathlineto{\pgfqpoint{2.733741in}{1.265210in}}%
\pgfpathlineto{\pgfqpoint{2.726510in}{1.268470in}}%
\pgfpathlineto{\pgfqpoint{2.724637in}{1.258622in}}%
\pgfpathlineto{\pgfqpoint{2.714315in}{1.255880in}}%
\pgfpathlineto{\pgfqpoint{2.715401in}{1.223478in}}%
\pgfpathlineto{\pgfqpoint{2.713513in}{1.213050in}}%
\pgfpathlineto{\pgfqpoint{2.718690in}{1.198187in}}%
\pgfpathlineto{\pgfqpoint{2.726654in}{1.191036in}}%
\pgfpathlineto{\pgfqpoint{2.723043in}{1.186546in}}%
\pgfpathlineto{\pgfqpoint{2.731172in}{1.179985in}}%
\pgfpathlineto{\pgfqpoint{2.734111in}{1.169336in}}%
\pgfpathlineto{\pgfqpoint{2.719210in}{1.178155in}}%
\pgfpathlineto{\pgfqpoint{2.709780in}{1.177009in}}%
\pgfpathlineto{\pgfqpoint{2.702458in}{1.186080in}}%
\pgfpathlineto{\pgfqpoint{2.695002in}{1.186964in}}%
\pgfpathlineto{\pgfqpoint{2.684407in}{1.182420in}}%
\pgfpathlineto{\pgfqpoint{2.680297in}{1.188082in}}%
\pgfpathlineto{\pgfqpoint{2.685698in}{1.199967in}}%
\pgfpathlineto{\pgfqpoint{2.688153in}{1.209860in}}%
\pgfpathlineto{\pgfqpoint{2.693583in}{1.217795in}}%
\pgfpathlineto{\pgfqpoint{2.685652in}{1.223151in}}%
\pgfpathlineto{\pgfqpoint{2.682414in}{1.218382in}}%
\pgfpathlineto{\pgfqpoint{2.674491in}{1.223222in}}%
\pgfpathlineto{\pgfqpoint{2.660462in}{1.226448in}}%
\pgfpathlineto{\pgfqpoint{2.661734in}{1.233539in}}%
\pgfpathlineto{\pgfqpoint{2.655375in}{1.237653in}}%
\pgfpathlineto{\pgfqpoint{2.647188in}{1.237365in}}%
\pgfpathlineto{\pgfqpoint{2.637267in}{1.254535in}}%
\pgfpathlineto{\pgfqpoint{2.630342in}{1.253004in}}%
\pgfpathlineto{\pgfqpoint{2.619817in}{1.257435in}}%
\pgfpathlineto{\pgfqpoint{2.612246in}{1.252404in}}%
\pgfpathlineto{\pgfqpoint{2.607031in}{1.243360in}}%
\pgfpathlineto{\pgfqpoint{2.599126in}{1.242584in}}%
\pgfpathlineto{\pgfqpoint{2.590060in}{1.246663in}}%
\pgfpathlineto{\pgfqpoint{2.583515in}{1.233653in}}%
\pgfpathlineto{\pgfqpoint{2.579067in}{1.235892in}}%
\pgfpathlineto{\pgfqpoint{2.560186in}{1.213439in}}%
\pgfpathlineto{\pgfqpoint{2.554781in}{1.247456in}}%
\pgfpathlineto{\pgfqpoint{2.623776in}{1.260244in}}%
\pgfpathlineto{\pgfqpoint{2.654706in}{1.265782in}}%
\pgfpathlineto{\pgfqpoint{2.699704in}{1.274845in}}%
\pgfpathlineto{\pgfqpoint{2.738266in}{1.283189in}}%
\pgfpathlineto{\pgfqpoint{2.743908in}{1.261445in}}%
\pgfpathlineto{\pgfqpoint{2.760342in}{1.203000in}}%
\pgfpathlineto{\pgfqpoint{2.791057in}{1.208880in}}%
\pgfusepath{stroke}%
\end{pgfscope}%
\begin{pgfscope}%
\pgfpathrectangle{\pgfqpoint{0.100000in}{0.100000in}}{\pgfqpoint{3.007045in}{1.925000in}}%
\pgfusepath{clip}%
\pgfsetbuttcap%
\pgfsetroundjoin%
\pgfsetlinewidth{0.050187pt}%
\definecolor{currentstroke}{rgb}{1.000000,1.000000,1.000000}%
\pgfsetstrokecolor{currentstroke}%
\pgfsetdash{}{0pt}%
\pgfpathmoveto{\pgfqpoint{1.402207in}{1.014889in}}%
\pgfpathlineto{\pgfqpoint{1.351869in}{1.019693in}}%
\pgfpathlineto{\pgfqpoint{1.299648in}{1.024321in}}%
\pgfpathlineto{\pgfqpoint{1.239303in}{1.030616in}}%
\pgfpathlineto{\pgfqpoint{1.149648in}{1.041279in}}%
\pgfpathlineto{\pgfqpoint{1.117842in}{1.046179in}}%
\pgfpathlineto{\pgfqpoint{1.035744in}{1.058103in}}%
\pgfpathlineto{\pgfqpoint{1.047732in}{1.133342in}}%
\pgfpathlineto{\pgfqpoint{1.048031in}{1.139428in}}%
\pgfpathlineto{\pgfqpoint{1.059579in}{1.212028in}}%
\pgfpathlineto{\pgfqpoint{1.068205in}{1.267320in}}%
\pgfpathlineto{\pgfqpoint{1.076345in}{1.318594in}}%
\pgfpathlineto{\pgfqpoint{1.131961in}{1.310611in}}%
\pgfpathlineto{\pgfqpoint{1.183815in}{1.303168in}}%
\pgfpathlineto{\pgfqpoint{1.279137in}{1.291403in}}%
\pgfpathlineto{\pgfqpoint{1.322867in}{1.287410in}}%
\pgfpathlineto{\pgfqpoint{1.392146in}{1.280688in}}%
\pgfpathlineto{\pgfqpoint{1.422117in}{1.278255in}}%
\pgfpathlineto{\pgfqpoint{1.416825in}{1.212637in}}%
\pgfpathlineto{\pgfqpoint{1.412045in}{1.149457in}}%
\pgfpathlineto{\pgfqpoint{1.408197in}{1.098013in}}%
\pgfpathlineto{\pgfqpoint{1.402207in}{1.014889in}}%
\pgfusepath{stroke}%
\end{pgfscope}%
\begin{pgfscope}%
\pgfpathrectangle{\pgfqpoint{0.100000in}{0.100000in}}{\pgfqpoint{3.007045in}{1.925000in}}%
\pgfusepath{clip}%
\pgfsetbuttcap%
\pgfsetroundjoin%
\pgfsetlinewidth{0.050187pt}%
\definecolor{currentstroke}{rgb}{1.000000,1.000000,1.000000}%
\pgfsetstrokecolor{currentstroke}%
\pgfsetdash{}{0pt}%
\pgfpathmoveto{\pgfqpoint{2.068136in}{0.977587in}}%
\pgfpathlineto{\pgfqpoint{2.069833in}{0.985194in}}%
\pgfpathlineto{\pgfqpoint{2.078628in}{0.983482in}}%
\pgfpathlineto{\pgfqpoint{2.081137in}{1.010475in}}%
\pgfpathlineto{\pgfqpoint{2.078730in}{1.015347in}}%
\pgfpathlineto{\pgfqpoint{2.087308in}{1.026197in}}%
\pgfpathlineto{\pgfqpoint{2.094189in}{1.025479in}}%
\pgfpathlineto{\pgfqpoint{2.115042in}{1.018342in}}%
\pgfpathlineto{\pgfqpoint{2.117486in}{1.024119in}}%
\pgfpathlineto{\pgfqpoint{2.112175in}{1.032983in}}%
\pgfpathlineto{\pgfqpoint{2.113967in}{1.039976in}}%
\pgfpathlineto{\pgfqpoint{2.133415in}{1.046407in}}%
\pgfpathlineto{\pgfqpoint{2.128773in}{1.058240in}}%
\pgfpathlineto{\pgfqpoint{2.134480in}{1.068255in}}%
\pgfpathlineto{\pgfqpoint{2.145103in}{1.073876in}}%
\pgfpathlineto{\pgfqpoint{2.167393in}{1.079442in}}%
\pgfpathlineto{\pgfqpoint{2.181457in}{1.071040in}}%
\pgfpathlineto{\pgfqpoint{2.185723in}{1.079234in}}%
\pgfpathlineto{\pgfqpoint{2.197401in}{1.084958in}}%
\pgfpathlineto{\pgfqpoint{2.200414in}{1.079910in}}%
\pgfpathlineto{\pgfqpoint{2.212410in}{1.083920in}}%
\pgfpathlineto{\pgfqpoint{2.211653in}{1.090784in}}%
\pgfpathlineto{\pgfqpoint{2.222438in}{1.098640in}}%
\pgfpathlineto{\pgfqpoint{2.228803in}{1.090482in}}%
\pgfpathlineto{\pgfqpoint{2.237177in}{1.089662in}}%
\pgfpathlineto{\pgfqpoint{2.242742in}{1.094978in}}%
\pgfpathlineto{\pgfqpoint{2.242123in}{1.102538in}}%
\pgfpathlineto{\pgfqpoint{2.246853in}{1.110045in}}%
\pgfpathlineto{\pgfqpoint{2.253191in}{1.111671in}}%
\pgfpathlineto{\pgfqpoint{2.255736in}{1.121584in}}%
\pgfpathlineto{\pgfqpoint{2.264955in}{1.130153in}}%
\pgfpathlineto{\pgfqpoint{2.262180in}{1.138682in}}%
\pgfpathlineto{\pgfqpoint{2.271148in}{1.142929in}}%
\pgfpathlineto{\pgfqpoint{2.277140in}{1.140317in}}%
\pgfpathlineto{\pgfqpoint{2.285991in}{1.146952in}}%
\pgfpathlineto{\pgfqpoint{2.293900in}{1.148681in}}%
\pgfpathlineto{\pgfqpoint{2.293928in}{1.155563in}}%
\pgfpathlineto{\pgfqpoint{2.288386in}{1.165112in}}%
\pgfpathlineto{\pgfqpoint{2.291994in}{1.168959in}}%
\pgfpathlineto{\pgfqpoint{2.303330in}{1.168107in}}%
\pgfpathlineto{\pgfqpoint{2.310538in}{1.171698in}}%
\pgfpathlineto{\pgfqpoint{2.317973in}{1.166088in}}%
\pgfpathlineto{\pgfqpoint{2.323842in}{1.153426in}}%
\pgfpathlineto{\pgfqpoint{2.343298in}{1.151287in}}%
\pgfpathlineto{\pgfqpoint{2.353440in}{1.144317in}}%
\pgfpathlineto{\pgfqpoint{2.360576in}{1.150208in}}%
\pgfpathlineto{\pgfqpoint{2.373478in}{1.145257in}}%
\pgfpathlineto{\pgfqpoint{2.380782in}{1.147717in}}%
\pgfpathlineto{\pgfqpoint{2.385887in}{1.154714in}}%
\pgfpathlineto{\pgfqpoint{2.393680in}{1.156052in}}%
\pgfpathlineto{\pgfqpoint{2.395706in}{1.148445in}}%
\pgfpathlineto{\pgfqpoint{2.404325in}{1.145536in}}%
\pgfpathlineto{\pgfqpoint{2.411023in}{1.138034in}}%
\pgfpathlineto{\pgfqpoint{2.413044in}{1.127903in}}%
\pgfpathlineto{\pgfqpoint{2.410356in}{1.120244in}}%
\pgfpathlineto{\pgfqpoint{2.420641in}{1.110328in}}%
\pgfpathlineto{\pgfqpoint{2.432723in}{1.091604in}}%
\pgfpathlineto{\pgfqpoint{2.441667in}{1.085443in}}%
\pgfpathlineto{\pgfqpoint{2.451373in}{1.085086in}}%
\pgfpathlineto{\pgfqpoint{2.433451in}{1.064522in}}%
\pgfpathlineto{\pgfqpoint{2.415810in}{1.052075in}}%
\pgfpathlineto{\pgfqpoint{2.409317in}{1.042191in}}%
\pgfpathlineto{\pgfqpoint{2.409427in}{1.036826in}}%
\pgfpathlineto{\pgfqpoint{2.399872in}{1.032715in}}%
\pgfpathlineto{\pgfqpoint{2.397134in}{1.024978in}}%
\pgfpathlineto{\pgfqpoint{2.377229in}{1.017186in}}%
\pgfpathlineto{\pgfqpoint{2.370173in}{1.012125in}}%
\pgfpathlineto{\pgfqpoint{2.369218in}{1.011045in}}%
\pgfpathlineto{\pgfqpoint{2.311961in}{1.005616in}}%
\pgfpathlineto{\pgfqpoint{2.277384in}{1.002708in}}%
\pgfpathlineto{\pgfqpoint{2.220644in}{0.999497in}}%
\pgfpathlineto{\pgfqpoint{2.149956in}{0.992504in}}%
\pgfpathlineto{\pgfqpoint{2.138278in}{0.994121in}}%
\pgfpathlineto{\pgfqpoint{2.140708in}{0.982200in}}%
\pgfpathlineto{\pgfqpoint{2.068136in}{0.977587in}}%
\pgfusepath{stroke}%
\end{pgfscope}%
\begin{pgfscope}%
\pgfpathrectangle{\pgfqpoint{0.100000in}{0.100000in}}{\pgfqpoint{3.007045in}{1.925000in}}%
\pgfusepath{clip}%
\pgfsetbuttcap%
\pgfsetroundjoin%
\pgfsetlinewidth{0.050187pt}%
\definecolor{currentstroke}{rgb}{1.000000,1.000000,1.000000}%
\pgfsetstrokecolor{currentstroke}%
\pgfsetdash{}{0pt}%
\pgfpathmoveto{\pgfqpoint{1.402207in}{1.014889in}}%
\pgfpathlineto{\pgfqpoint{1.408197in}{1.098013in}}%
\pgfpathlineto{\pgfqpoint{1.412045in}{1.149457in}}%
\pgfpathlineto{\pgfqpoint{1.416825in}{1.212637in}}%
\pgfpathlineto{\pgfqpoint{1.483096in}{1.207983in}}%
\pgfpathlineto{\pgfqpoint{1.567259in}{1.203386in}}%
\pgfpathlineto{\pgfqpoint{1.624502in}{1.201190in}}%
\pgfpathlineto{\pgfqpoint{1.681409in}{1.199446in}}%
\pgfpathlineto{\pgfqpoint{1.732915in}{1.198649in}}%
\pgfpathlineto{\pgfqpoint{1.756737in}{1.198896in}}%
\pgfpathlineto{\pgfqpoint{1.767223in}{1.190339in}}%
\pgfpathlineto{\pgfqpoint{1.775436in}{1.192065in}}%
\pgfpathlineto{\pgfqpoint{1.775693in}{1.184601in}}%
\pgfpathlineto{\pgfqpoint{1.769593in}{1.171690in}}%
\pgfpathlineto{\pgfqpoint{1.770255in}{1.163533in}}%
\pgfpathlineto{\pgfqpoint{1.776095in}{1.158154in}}%
\pgfpathlineto{\pgfqpoint{1.781459in}{1.146947in}}%
\pgfpathlineto{\pgfqpoint{1.792341in}{1.140513in}}%
\pgfpathlineto{\pgfqpoint{1.791978in}{1.098269in}}%
\pgfpathlineto{\pgfqpoint{1.792298in}{1.001125in}}%
\pgfpathlineto{\pgfqpoint{1.719359in}{1.001461in}}%
\pgfpathlineto{\pgfqpoint{1.642551in}{1.002805in}}%
\pgfpathlineto{\pgfqpoint{1.586009in}{1.004800in}}%
\pgfpathlineto{\pgfqpoint{1.538812in}{1.006621in}}%
\pgfpathlineto{\pgfqpoint{1.459299in}{1.011309in}}%
\pgfpathlineto{\pgfqpoint{1.402207in}{1.014889in}}%
\pgfusepath{stroke}%
\end{pgfscope}%
\begin{pgfscope}%
\pgfpathrectangle{\pgfqpoint{0.100000in}{0.100000in}}{\pgfqpoint{3.007045in}{1.925000in}}%
\pgfusepath{clip}%
\pgfsetbuttcap%
\pgfsetroundjoin%
\pgfsetlinewidth{0.050187pt}%
\definecolor{currentstroke}{rgb}{1.000000,1.000000,1.000000}%
\pgfsetstrokecolor{currentstroke}%
\pgfsetdash{}{0pt}%
\pgfpathmoveto{\pgfqpoint{2.475422in}{1.025896in}}%
\pgfpathlineto{\pgfqpoint{2.394637in}{1.014547in}}%
\pgfpathlineto{\pgfqpoint{2.370173in}{1.012125in}}%
\pgfpathlineto{\pgfqpoint{2.377229in}{1.017186in}}%
\pgfpathlineto{\pgfqpoint{2.397134in}{1.024978in}}%
\pgfpathlineto{\pgfqpoint{2.399872in}{1.032715in}}%
\pgfpathlineto{\pgfqpoint{2.409427in}{1.036826in}}%
\pgfpathlineto{\pgfqpoint{2.409317in}{1.042191in}}%
\pgfpathlineto{\pgfqpoint{2.415810in}{1.052075in}}%
\pgfpathlineto{\pgfqpoint{2.433451in}{1.064522in}}%
\pgfpathlineto{\pgfqpoint{2.451373in}{1.085086in}}%
\pgfpathlineto{\pgfqpoint{2.450500in}{1.080661in}}%
\pgfpathlineto{\pgfqpoint{2.459950in}{1.069974in}}%
\pgfpathlineto{\pgfqpoint{2.475358in}{1.066973in}}%
\pgfpathlineto{\pgfqpoint{2.484322in}{1.077155in}}%
\pgfpathlineto{\pgfqpoint{2.492341in}{1.071792in}}%
\pgfpathlineto{\pgfqpoint{2.506692in}{1.078086in}}%
\pgfpathlineto{\pgfqpoint{2.508744in}{1.084231in}}%
\pgfpathlineto{\pgfqpoint{2.517650in}{1.085614in}}%
\pgfpathlineto{\pgfqpoint{2.526900in}{1.092874in}}%
\pgfpathlineto{\pgfqpoint{2.535053in}{1.095270in}}%
\pgfpathlineto{\pgfqpoint{2.535313in}{1.108096in}}%
\pgfpathlineto{\pgfqpoint{2.540593in}{1.121950in}}%
\pgfpathlineto{\pgfqpoint{2.547556in}{1.130626in}}%
\pgfpathlineto{\pgfqpoint{2.549814in}{1.143653in}}%
\pgfpathlineto{\pgfqpoint{2.558890in}{1.161951in}}%
\pgfpathlineto{\pgfqpoint{2.559215in}{1.172839in}}%
\pgfpathlineto{\pgfqpoint{2.568590in}{1.165745in}}%
\pgfpathlineto{\pgfqpoint{2.577757in}{1.164024in}}%
\pgfpathlineto{\pgfqpoint{2.581324in}{1.169185in}}%
\pgfpathlineto{\pgfqpoint{2.587696in}{1.191677in}}%
\pgfpathlineto{\pgfqpoint{2.596570in}{1.190815in}}%
\pgfpathlineto{\pgfqpoint{2.598920in}{1.199055in}}%
\pgfpathlineto{\pgfqpoint{2.604125in}{1.202772in}}%
\pgfpathlineto{\pgfqpoint{2.615154in}{1.221459in}}%
\pgfpathlineto{\pgfqpoint{2.613553in}{1.226987in}}%
\pgfpathlineto{\pgfqpoint{2.618406in}{1.239002in}}%
\pgfpathlineto{\pgfqpoint{2.644319in}{1.224366in}}%
\pgfpathlineto{\pgfqpoint{2.647188in}{1.237365in}}%
\pgfpathlineto{\pgfqpoint{2.655375in}{1.237653in}}%
\pgfpathlineto{\pgfqpoint{2.661734in}{1.233539in}}%
\pgfpathlineto{\pgfqpoint{2.660462in}{1.226448in}}%
\pgfpathlineto{\pgfqpoint{2.674491in}{1.223222in}}%
\pgfpathlineto{\pgfqpoint{2.682414in}{1.218382in}}%
\pgfpathlineto{\pgfqpoint{2.688153in}{1.209860in}}%
\pgfpathlineto{\pgfqpoint{2.685698in}{1.199967in}}%
\pgfpathlineto{\pgfqpoint{2.680740in}{1.199156in}}%
\pgfpathlineto{\pgfqpoint{2.677865in}{1.184228in}}%
\pgfpathlineto{\pgfqpoint{2.684163in}{1.178391in}}%
\pgfpathlineto{\pgfqpoint{2.693030in}{1.183109in}}%
\pgfpathlineto{\pgfqpoint{2.701257in}{1.173146in}}%
\pgfpathlineto{\pgfqpoint{2.719638in}{1.171355in}}%
\pgfpathlineto{\pgfqpoint{2.722805in}{1.165624in}}%
\pgfpathlineto{\pgfqpoint{2.739815in}{1.160050in}}%
\pgfpathlineto{\pgfqpoint{2.737716in}{1.153473in}}%
\pgfpathlineto{\pgfqpoint{2.739929in}{1.134884in}}%
\pgfpathlineto{\pgfqpoint{2.746795in}{1.127865in}}%
\pgfpathlineto{\pgfqpoint{2.736664in}{1.124902in}}%
\pgfpathlineto{\pgfqpoint{2.738007in}{1.116967in}}%
\pgfpathlineto{\pgfqpoint{2.748815in}{1.110238in}}%
\pgfpathlineto{\pgfqpoint{2.749789in}{1.103601in}}%
\pgfpathlineto{\pgfqpoint{2.743687in}{1.098614in}}%
\pgfpathlineto{\pgfqpoint{2.731375in}{1.110429in}}%
\pgfpathlineto{\pgfqpoint{2.730158in}{1.101808in}}%
\pgfpathlineto{\pgfqpoint{2.740465in}{1.097709in}}%
\pgfpathlineto{\pgfqpoint{2.741496in}{1.093469in}}%
\pgfpathlineto{\pgfqpoint{2.755633in}{1.099072in}}%
\pgfpathlineto{\pgfqpoint{2.766466in}{1.100555in}}%
\pgfpathlineto{\pgfqpoint{2.777563in}{1.078165in}}%
\pgfpathlineto{\pgfqpoint{2.776325in}{1.077922in}}%
\pgfpathlineto{\pgfqpoint{2.771311in}{1.076886in}}%
\pgfpathlineto{\pgfqpoint{2.769833in}{1.076577in}}%
\pgfpathlineto{\pgfqpoint{2.768856in}{1.076386in}}%
\pgfpathlineto{\pgfqpoint{2.702785in}{1.062681in}}%
\pgfpathlineto{\pgfqpoint{2.643688in}{1.050585in}}%
\pgfpathlineto{\pgfqpoint{2.561981in}{1.036510in}}%
\pgfpathlineto{\pgfqpoint{2.492587in}{1.027345in}}%
\pgfpathlineto{\pgfqpoint{2.475422in}{1.025896in}}%
\pgfusepath{stroke}%
\end{pgfscope}%
\begin{pgfscope}%
\pgfpathrectangle{\pgfqpoint{0.100000in}{0.100000in}}{\pgfqpoint{3.007045in}{1.925000in}}%
\pgfusepath{clip}%
\pgfsetbuttcap%
\pgfsetroundjoin%
\pgfsetlinewidth{0.050187pt}%
\definecolor{currentstroke}{rgb}{1.000000,1.000000,1.000000}%
\pgfsetstrokecolor{currentstroke}%
\pgfsetdash{}{0pt}%
\pgfpathmoveto{\pgfqpoint{2.769772in}{1.171817in}}%
\pgfpathlineto{\pgfqpoint{2.783092in}{1.177859in}}%
\pgfpathlineto{\pgfqpoint{2.772455in}{1.146743in}}%
\pgfpathlineto{\pgfqpoint{2.765421in}{1.130265in}}%
\pgfpathlineto{\pgfqpoint{2.760353in}{1.139590in}}%
\pgfpathlineto{\pgfqpoint{2.769365in}{1.161915in}}%
\pgfpathlineto{\pgfqpoint{2.769772in}{1.171817in}}%
\pgfusepath{stroke}%
\end{pgfscope}%
\begin{pgfscope}%
\pgfpathrectangle{\pgfqpoint{0.100000in}{0.100000in}}{\pgfqpoint{3.007045in}{1.925000in}}%
\pgfusepath{clip}%
\pgfsetbuttcap%
\pgfsetroundjoin%
\pgfsetlinewidth{0.050187pt}%
\definecolor{currentstroke}{rgb}{1.000000,1.000000,1.000000}%
\pgfsetstrokecolor{currentstroke}%
\pgfsetdash{}{0pt}%
\pgfpathmoveto{\pgfqpoint{2.068136in}{0.977587in}}%
\pgfpathlineto{\pgfqpoint{2.064919in}{0.977122in}}%
\pgfpathlineto{\pgfqpoint{2.061886in}{0.976909in}}%
\pgfpathlineto{\pgfqpoint{2.063184in}{0.967596in}}%
\pgfpathlineto{\pgfqpoint{2.055358in}{0.958208in}}%
\pgfpathlineto{\pgfqpoint{2.053816in}{0.943518in}}%
\pgfpathlineto{\pgfqpoint{2.018831in}{0.940931in}}%
\pgfpathlineto{\pgfqpoint{2.021881in}{0.947828in}}%
\pgfpathlineto{\pgfqpoint{2.034495in}{0.960424in}}%
\pgfpathlineto{\pgfqpoint{2.034844in}{0.967726in}}%
\pgfpathlineto{\pgfqpoint{2.029259in}{0.974639in}}%
\pgfpathlineto{\pgfqpoint{1.977188in}{0.971886in}}%
\pgfpathlineto{\pgfqpoint{1.886150in}{0.968996in}}%
\pgfpathlineto{\pgfqpoint{1.832874in}{0.968025in}}%
\pgfpathlineto{\pgfqpoint{1.792603in}{0.967662in}}%
\pgfpathlineto{\pgfqpoint{1.792298in}{1.001125in}}%
\pgfpathlineto{\pgfqpoint{1.791978in}{1.098269in}}%
\pgfpathlineto{\pgfqpoint{1.792341in}{1.140513in}}%
\pgfpathlineto{\pgfqpoint{1.781459in}{1.146947in}}%
\pgfpathlineto{\pgfqpoint{1.776095in}{1.158154in}}%
\pgfpathlineto{\pgfqpoint{1.770255in}{1.163533in}}%
\pgfpathlineto{\pgfqpoint{1.769593in}{1.171690in}}%
\pgfpathlineto{\pgfqpoint{1.775693in}{1.184601in}}%
\pgfpathlineto{\pgfqpoint{1.775436in}{1.192065in}}%
\pgfpathlineto{\pgfqpoint{1.767223in}{1.190339in}}%
\pgfpathlineto{\pgfqpoint{1.756737in}{1.198896in}}%
\pgfpathlineto{\pgfqpoint{1.748332in}{1.213916in}}%
\pgfpathlineto{\pgfqpoint{1.741281in}{1.220848in}}%
\pgfpathlineto{\pgfqpoint{1.733914in}{1.237878in}}%
\pgfpathlineto{\pgfqpoint{1.810432in}{1.236685in}}%
\pgfpathlineto{\pgfqpoint{1.886493in}{1.239191in}}%
\pgfpathlineto{\pgfqpoint{1.935262in}{1.242004in}}%
\pgfpathlineto{\pgfqpoint{1.950519in}{1.226895in}}%
\pgfpathlineto{\pgfqpoint{1.947730in}{1.211138in}}%
\pgfpathlineto{\pgfqpoint{1.951798in}{1.193397in}}%
\pgfpathlineto{\pgfqpoint{1.956311in}{1.184451in}}%
\pgfpathlineto{\pgfqpoint{1.970708in}{1.172146in}}%
\pgfpathlineto{\pgfqpoint{1.973858in}{1.166711in}}%
\pgfpathlineto{\pgfqpoint{1.989621in}{1.154462in}}%
\pgfpathlineto{\pgfqpoint{1.991906in}{1.140721in}}%
\pgfpathlineto{\pgfqpoint{1.996433in}{1.130585in}}%
\pgfpathlineto{\pgfqpoint{2.007129in}{1.136489in}}%
\pgfpathlineto{\pgfqpoint{2.015945in}{1.133789in}}%
\pgfpathlineto{\pgfqpoint{2.022602in}{1.126623in}}%
\pgfpathlineto{\pgfqpoint{2.020088in}{1.113940in}}%
\pgfpathlineto{\pgfqpoint{2.011918in}{1.097359in}}%
\pgfpathlineto{\pgfqpoint{2.012044in}{1.088723in}}%
\pgfpathlineto{\pgfqpoint{2.018150in}{1.081642in}}%
\pgfpathlineto{\pgfqpoint{2.031291in}{1.072263in}}%
\pgfpathlineto{\pgfqpoint{2.046423in}{1.064650in}}%
\pgfpathlineto{\pgfqpoint{2.050012in}{1.058555in}}%
\pgfpathlineto{\pgfqpoint{2.057751in}{1.055672in}}%
\pgfpathlineto{\pgfqpoint{2.058058in}{1.047582in}}%
\pgfpathlineto{\pgfqpoint{2.063786in}{1.037307in}}%
\pgfpathlineto{\pgfqpoint{2.059812in}{1.028177in}}%
\pgfpathlineto{\pgfqpoint{2.067590in}{1.013653in}}%
\pgfpathlineto{\pgfqpoint{2.073675in}{1.016537in}}%
\pgfpathlineto{\pgfqpoint{2.081137in}{1.010475in}}%
\pgfpathlineto{\pgfqpoint{2.078628in}{0.983482in}}%
\pgfpathlineto{\pgfqpoint{2.069833in}{0.985194in}}%
\pgfpathlineto{\pgfqpoint{2.068136in}{0.977587in}}%
\pgfusepath{stroke}%
\end{pgfscope}%
\begin{pgfscope}%
\pgfpathrectangle{\pgfqpoint{0.100000in}{0.100000in}}{\pgfqpoint{3.007045in}{1.925000in}}%
\pgfusepath{clip}%
\pgfsetbuttcap%
\pgfsetroundjoin%
\pgfsetlinewidth{0.050187pt}%
\definecolor{currentstroke}{rgb}{1.000000,1.000000,1.000000}%
\pgfsetstrokecolor{currentstroke}%
\pgfsetdash{}{0pt}%
\pgfpathmoveto{\pgfqpoint{0.719235in}{0.983407in}}%
\pgfpathlineto{\pgfqpoint{0.725604in}{0.996967in}}%
\pgfpathlineto{\pgfqpoint{0.723861in}{1.017321in}}%
\pgfpathlineto{\pgfqpoint{0.726911in}{1.023110in}}%
\pgfpathlineto{\pgfqpoint{0.727212in}{1.048463in}}%
\pgfpathlineto{\pgfqpoint{0.730103in}{1.055764in}}%
\pgfpathlineto{\pgfqpoint{0.740288in}{1.056900in}}%
\pgfpathlineto{\pgfqpoint{0.749739in}{1.053659in}}%
\pgfpathlineto{\pgfqpoint{0.753855in}{1.044733in}}%
\pgfpathlineto{\pgfqpoint{0.759627in}{1.045075in}}%
\pgfpathlineto{\pgfqpoint{0.766797in}{1.055300in}}%
\pgfpathlineto{\pgfqpoint{0.777160in}{1.105704in}}%
\pgfpathlineto{\pgfqpoint{0.836134in}{1.093528in}}%
\pgfpathlineto{\pgfqpoint{0.870361in}{1.086778in}}%
\pgfpathlineto{\pgfqpoint{0.961251in}{1.070801in}}%
\pgfpathlineto{\pgfqpoint{0.986397in}{1.065751in}}%
\pgfpathlineto{\pgfqpoint{1.035744in}{1.058103in}}%
\pgfpathlineto{\pgfqpoint{1.025626in}{0.992952in}}%
\pgfpathlineto{\pgfqpoint{1.015102in}{0.924993in}}%
\pgfpathlineto{\pgfqpoint{1.002983in}{0.848537in}}%
\pgfpathlineto{\pgfqpoint{0.989353in}{0.760774in}}%
\pgfpathlineto{\pgfqpoint{0.978343in}{0.688693in}}%
\pgfpathlineto{\pgfqpoint{0.899477in}{0.701217in}}%
\pgfpathlineto{\pgfqpoint{0.864826in}{0.707151in}}%
\pgfpathlineto{\pgfqpoint{0.849347in}{0.716418in}}%
\pgfpathlineto{\pgfqpoint{0.748420in}{0.777313in}}%
\pgfpathlineto{\pgfqpoint{0.672680in}{0.823525in}}%
\pgfpathlineto{\pgfqpoint{0.675202in}{0.831715in}}%
\pgfpathlineto{\pgfqpoint{0.681459in}{0.837446in}}%
\pgfpathlineto{\pgfqpoint{0.687967in}{0.836384in}}%
\pgfpathlineto{\pgfqpoint{0.697405in}{0.842403in}}%
\pgfpathlineto{\pgfqpoint{0.698891in}{0.851052in}}%
\pgfpathlineto{\pgfqpoint{0.687380in}{0.861545in}}%
\pgfpathlineto{\pgfqpoint{0.695567in}{0.881688in}}%
\pgfpathlineto{\pgfqpoint{0.703798in}{0.889434in}}%
\pgfpathlineto{\pgfqpoint{0.707704in}{0.898624in}}%
\pgfpathlineto{\pgfqpoint{0.710111in}{0.915475in}}%
\pgfpathlineto{\pgfqpoint{0.717829in}{0.923106in}}%
\pgfpathlineto{\pgfqpoint{0.734029in}{0.930728in}}%
\pgfpathlineto{\pgfqpoint{0.733671in}{0.937444in}}%
\pgfpathlineto{\pgfqpoint{0.724623in}{0.945778in}}%
\pgfpathlineto{\pgfqpoint{0.723391in}{0.962960in}}%
\pgfpathlineto{\pgfqpoint{0.717162in}{0.975520in}}%
\pgfpathlineto{\pgfqpoint{0.719235in}{0.983407in}}%
\pgfusepath{stroke}%
\end{pgfscope}%
\begin{pgfscope}%
\pgfpathrectangle{\pgfqpoint{0.100000in}{0.100000in}}{\pgfqpoint{3.007045in}{1.925000in}}%
\pgfusepath{clip}%
\pgfsetbuttcap%
\pgfsetroundjoin%
\pgfsetlinewidth{0.050187pt}%
\definecolor{currentstroke}{rgb}{1.000000,1.000000,1.000000}%
\pgfsetstrokecolor{currentstroke}%
\pgfsetdash{}{0pt}%
\pgfpathmoveto{\pgfqpoint{1.801141in}{0.779025in}}%
\pgfpathlineto{\pgfqpoint{1.786521in}{0.783547in}}%
\pgfpathlineto{\pgfqpoint{1.767788in}{0.797893in}}%
\pgfpathlineto{\pgfqpoint{1.759467in}{0.800976in}}%
\pgfpathlineto{\pgfqpoint{1.754180in}{0.794780in}}%
\pgfpathlineto{\pgfqpoint{1.747503in}{0.794467in}}%
\pgfpathlineto{\pgfqpoint{1.739055in}{0.799728in}}%
\pgfpathlineto{\pgfqpoint{1.725799in}{0.792988in}}%
\pgfpathlineto{\pgfqpoint{1.721080in}{0.796307in}}%
\pgfpathlineto{\pgfqpoint{1.708036in}{0.788443in}}%
\pgfpathlineto{\pgfqpoint{1.694275in}{0.789894in}}%
\pgfpathlineto{\pgfqpoint{1.683661in}{0.795004in}}%
\pgfpathlineto{\pgfqpoint{1.676227in}{0.793081in}}%
\pgfpathlineto{\pgfqpoint{1.664329in}{0.800308in}}%
\pgfpathlineto{\pgfqpoint{1.657711in}{0.787562in}}%
\pgfpathlineto{\pgfqpoint{1.650942in}{0.798523in}}%
\pgfpathlineto{\pgfqpoint{1.637572in}{0.794271in}}%
\pgfpathlineto{\pgfqpoint{1.625879in}{0.804682in}}%
\pgfpathlineto{\pgfqpoint{1.615661in}{0.796290in}}%
\pgfpathlineto{\pgfqpoint{1.610544in}{0.804002in}}%
\pgfpathlineto{\pgfqpoint{1.603167in}{0.806507in}}%
\pgfpathlineto{\pgfqpoint{1.598675in}{0.813942in}}%
\pgfpathlineto{\pgfqpoint{1.589028in}{0.816058in}}%
\pgfpathlineto{\pgfqpoint{1.583461in}{0.810461in}}%
\pgfpathlineto{\pgfqpoint{1.574003in}{0.817714in}}%
\pgfpathlineto{\pgfqpoint{1.569594in}{0.816058in}}%
\pgfpathlineto{\pgfqpoint{1.553918in}{0.821935in}}%
\pgfpathlineto{\pgfqpoint{1.544100in}{0.822590in}}%
\pgfpathlineto{\pgfqpoint{1.539706in}{0.835071in}}%
\pgfpathlineto{\pgfqpoint{1.531836in}{0.833517in}}%
\pgfpathlineto{\pgfqpoint{1.522821in}{0.836596in}}%
\pgfpathlineto{\pgfqpoint{1.516885in}{0.834817in}}%
\pgfpathlineto{\pgfqpoint{1.504154in}{0.848834in}}%
\pgfpathlineto{\pgfqpoint{1.500607in}{0.847917in}}%
\pgfpathlineto{\pgfqpoint{1.503828in}{0.904762in}}%
\pgfpathlineto{\pgfqpoint{1.507338in}{0.975120in}}%
\pgfpathlineto{\pgfqpoint{1.449726in}{0.978343in}}%
\pgfpathlineto{\pgfqpoint{1.392875in}{0.982641in}}%
\pgfpathlineto{\pgfqpoint{1.348953in}{0.986455in}}%
\pgfpathlineto{\pgfqpoint{1.351869in}{1.019693in}}%
\pgfpathlineto{\pgfqpoint{1.402207in}{1.014889in}}%
\pgfpathlineto{\pgfqpoint{1.459299in}{1.011309in}}%
\pgfpathlineto{\pgfqpoint{1.538812in}{1.006621in}}%
\pgfpathlineto{\pgfqpoint{1.586009in}{1.004800in}}%
\pgfpathlineto{\pgfqpoint{1.642551in}{1.002805in}}%
\pgfpathlineto{\pgfqpoint{1.719359in}{1.001461in}}%
\pgfpathlineto{\pgfqpoint{1.792298in}{1.001125in}}%
\pgfpathlineto{\pgfqpoint{1.792603in}{0.967662in}}%
\pgfpathlineto{\pgfqpoint{1.796697in}{0.942452in}}%
\pgfpathlineto{\pgfqpoint{1.803058in}{0.895889in}}%
\pgfpathlineto{\pgfqpoint{1.801746in}{0.816372in}}%
\pgfpathlineto{\pgfqpoint{1.801141in}{0.779025in}}%
\pgfusepath{stroke}%
\end{pgfscope}%
\begin{pgfscope}%
\pgfpathrectangle{\pgfqpoint{0.100000in}{0.100000in}}{\pgfqpoint{3.007045in}{1.925000in}}%
\pgfusepath{clip}%
\pgfsetbuttcap%
\pgfsetroundjoin%
\pgfsetlinewidth{0.050187pt}%
\definecolor{currentstroke}{rgb}{1.000000,1.000000,1.000000}%
\pgfsetstrokecolor{currentstroke}%
\pgfsetdash{}{0pt}%
\pgfpathmoveto{\pgfqpoint{2.348542in}{0.902516in}}%
\pgfpathlineto{\pgfqpoint{2.348585in}{0.917250in}}%
\pgfpathlineto{\pgfqpoint{2.361388in}{0.922908in}}%
\pgfpathlineto{\pgfqpoint{2.361925in}{0.931952in}}%
\pgfpathlineto{\pgfqpoint{2.373389in}{0.943000in}}%
\pgfpathlineto{\pgfqpoint{2.387718in}{0.945263in}}%
\pgfpathlineto{\pgfqpoint{2.400316in}{0.957439in}}%
\pgfpathlineto{\pgfqpoint{2.413472in}{0.962881in}}%
\pgfpathlineto{\pgfqpoint{2.423346in}{0.979220in}}%
\pgfpathlineto{\pgfqpoint{2.432339in}{0.978033in}}%
\pgfpathlineto{\pgfqpoint{2.451345in}{0.992920in}}%
\pgfpathlineto{\pgfqpoint{2.461381in}{0.993200in}}%
\pgfpathlineto{\pgfqpoint{2.465620in}{1.004543in}}%
\pgfpathlineto{\pgfqpoint{2.473597in}{1.012440in}}%
\pgfpathlineto{\pgfqpoint{2.475422in}{1.025896in}}%
\pgfpathlineto{\pgfqpoint{2.492587in}{1.027345in}}%
\pgfpathlineto{\pgfqpoint{2.561981in}{1.036510in}}%
\pgfpathlineto{\pgfqpoint{2.643688in}{1.050585in}}%
\pgfpathlineto{\pgfqpoint{2.702785in}{1.062681in}}%
\pgfpathlineto{\pgfqpoint{2.768856in}{1.076386in}}%
\pgfpathlineto{\pgfqpoint{2.787736in}{1.050465in}}%
\pgfpathlineto{\pgfqpoint{2.777509in}{1.052120in}}%
\pgfpathlineto{\pgfqpoint{2.764382in}{1.048922in}}%
\pgfpathlineto{\pgfqpoint{2.751446in}{1.035712in}}%
\pgfpathlineto{\pgfqpoint{2.742149in}{1.036662in}}%
\pgfpathlineto{\pgfqpoint{2.740973in}{1.028816in}}%
\pgfpathlineto{\pgfqpoint{2.757777in}{1.035021in}}%
\pgfpathlineto{\pgfqpoint{2.774686in}{1.037579in}}%
\pgfpathlineto{\pgfqpoint{2.780965in}{1.034175in}}%
\pgfpathlineto{\pgfqpoint{2.789414in}{1.038056in}}%
\pgfpathlineto{\pgfqpoint{2.793774in}{1.035336in}}%
\pgfpathlineto{\pgfqpoint{2.797633in}{1.022389in}}%
\pgfpathlineto{\pgfqpoint{2.789602in}{1.018379in}}%
\pgfpathlineto{\pgfqpoint{2.784090in}{1.002591in}}%
\pgfpathlineto{\pgfqpoint{2.778311in}{0.996443in}}%
\pgfpathlineto{\pgfqpoint{2.760595in}{0.997788in}}%
\pgfpathlineto{\pgfqpoint{2.761547in}{1.007059in}}%
\pgfpathlineto{\pgfqpoint{2.751872in}{1.003009in}}%
\pgfpathlineto{\pgfqpoint{2.749767in}{0.995277in}}%
\pgfpathlineto{\pgfqpoint{2.756399in}{0.987268in}}%
\pgfpathlineto{\pgfqpoint{2.758652in}{0.973674in}}%
\pgfpathlineto{\pgfqpoint{2.747800in}{0.965929in}}%
\pgfpathlineto{\pgfqpoint{2.759537in}{0.963205in}}%
\pgfpathlineto{\pgfqpoint{2.764848in}{0.968900in}}%
\pgfpathlineto{\pgfqpoint{2.776591in}{0.969594in}}%
\pgfpathlineto{\pgfqpoint{2.770556in}{0.957295in}}%
\pgfpathlineto{\pgfqpoint{2.763165in}{0.951369in}}%
\pgfpathlineto{\pgfqpoint{2.740807in}{0.945436in}}%
\pgfpathlineto{\pgfqpoint{2.717851in}{0.924616in}}%
\pgfpathlineto{\pgfqpoint{2.703697in}{0.904102in}}%
\pgfpathlineto{\pgfqpoint{2.703613in}{0.895771in}}%
\pgfpathlineto{\pgfqpoint{2.697959in}{0.884275in}}%
\pgfpathlineto{\pgfqpoint{2.668941in}{0.876711in}}%
\pgfpathlineto{\pgfqpoint{2.599931in}{0.926180in}}%
\pgfpathlineto{\pgfqpoint{2.539158in}{0.917251in}}%
\pgfpathlineto{\pgfqpoint{2.538649in}{0.925496in}}%
\pgfpathlineto{\pgfqpoint{2.529411in}{0.934781in}}%
\pgfpathlineto{\pgfqpoint{2.522418in}{0.937054in}}%
\pgfpathlineto{\pgfqpoint{2.456388in}{0.930193in}}%
\pgfpathlineto{\pgfqpoint{2.441211in}{0.925045in}}%
\pgfpathlineto{\pgfqpoint{2.413807in}{0.911413in}}%
\pgfpathlineto{\pgfqpoint{2.390124in}{0.907640in}}%
\pgfpathlineto{\pgfqpoint{2.348542in}{0.902516in}}%
\pgfusepath{stroke}%
\end{pgfscope}%
\begin{pgfscope}%
\pgfpathrectangle{\pgfqpoint{0.100000in}{0.100000in}}{\pgfqpoint{3.007045in}{1.925000in}}%
\pgfusepath{clip}%
\pgfsetbuttcap%
\pgfsetroundjoin%
\pgfsetlinewidth{0.050187pt}%
\definecolor{currentstroke}{rgb}{1.000000,1.000000,1.000000}%
\pgfsetstrokecolor{currentstroke}%
\pgfsetdash{}{0pt}%
\pgfpathmoveto{\pgfqpoint{2.348542in}{0.902516in}}%
\pgfpathlineto{\pgfqpoint{2.279433in}{0.894926in}}%
\pgfpathlineto{\pgfqpoint{2.216213in}{0.889241in}}%
\pgfpathlineto{\pgfqpoint{2.151235in}{0.885039in}}%
\pgfpathlineto{\pgfqpoint{2.140081in}{0.883422in}}%
\pgfpathlineto{\pgfqpoint{2.096280in}{0.880053in}}%
\pgfpathlineto{\pgfqpoint{2.026124in}{0.875959in}}%
\pgfpathlineto{\pgfqpoint{2.038606in}{0.886320in}}%
\pgfpathlineto{\pgfqpoint{2.035758in}{0.897226in}}%
\pgfpathlineto{\pgfqpoint{2.038517in}{0.910446in}}%
\pgfpathlineto{\pgfqpoint{2.042732in}{0.916670in}}%
\pgfpathlineto{\pgfqpoint{2.042574in}{0.925329in}}%
\pgfpathlineto{\pgfqpoint{2.053797in}{0.930775in}}%
\pgfpathlineto{\pgfqpoint{2.053816in}{0.943518in}}%
\pgfpathlineto{\pgfqpoint{2.055358in}{0.958208in}}%
\pgfpathlineto{\pgfqpoint{2.063184in}{0.967596in}}%
\pgfpathlineto{\pgfqpoint{2.061886in}{0.976909in}}%
\pgfpathlineto{\pgfqpoint{2.064919in}{0.977122in}}%
\pgfpathlineto{\pgfqpoint{2.068136in}{0.977587in}}%
\pgfpathlineto{\pgfqpoint{2.140708in}{0.982200in}}%
\pgfpathlineto{\pgfqpoint{2.138278in}{0.994121in}}%
\pgfpathlineto{\pgfqpoint{2.149956in}{0.992504in}}%
\pgfpathlineto{\pgfqpoint{2.220644in}{0.999497in}}%
\pgfpathlineto{\pgfqpoint{2.277384in}{1.002708in}}%
\pgfpathlineto{\pgfqpoint{2.311961in}{1.005616in}}%
\pgfpathlineto{\pgfqpoint{2.369218in}{1.011045in}}%
\pgfpathlineto{\pgfqpoint{2.370173in}{1.012125in}}%
\pgfpathlineto{\pgfqpoint{2.394637in}{1.014547in}}%
\pgfpathlineto{\pgfqpoint{2.475422in}{1.025896in}}%
\pgfpathlineto{\pgfqpoint{2.473597in}{1.012440in}}%
\pgfpathlineto{\pgfqpoint{2.465620in}{1.004543in}}%
\pgfpathlineto{\pgfqpoint{2.461381in}{0.993200in}}%
\pgfpathlineto{\pgfqpoint{2.451345in}{0.992920in}}%
\pgfpathlineto{\pgfqpoint{2.432339in}{0.978033in}}%
\pgfpathlineto{\pgfqpoint{2.423346in}{0.979220in}}%
\pgfpathlineto{\pgfqpoint{2.413472in}{0.962881in}}%
\pgfpathlineto{\pgfqpoint{2.400316in}{0.957439in}}%
\pgfpathlineto{\pgfqpoint{2.387718in}{0.945263in}}%
\pgfpathlineto{\pgfqpoint{2.373389in}{0.943000in}}%
\pgfpathlineto{\pgfqpoint{2.361925in}{0.931952in}}%
\pgfpathlineto{\pgfqpoint{2.361388in}{0.922908in}}%
\pgfpathlineto{\pgfqpoint{2.348585in}{0.917250in}}%
\pgfpathlineto{\pgfqpoint{2.348542in}{0.902516in}}%
\pgfusepath{stroke}%
\end{pgfscope}%
\begin{pgfscope}%
\pgfpathrectangle{\pgfqpoint{0.100000in}{0.100000in}}{\pgfqpoint{3.007045in}{1.925000in}}%
\pgfusepath{clip}%
\pgfsetbuttcap%
\pgfsetroundjoin%
\pgfsetlinewidth{0.050187pt}%
\definecolor{currentstroke}{rgb}{1.000000,1.000000,1.000000}%
\pgfsetstrokecolor{currentstroke}%
\pgfsetdash{}{0pt}%
\pgfpathmoveto{\pgfqpoint{1.122887in}{0.697915in}}%
\pgfpathlineto{\pgfqpoint{1.117958in}{0.705835in}}%
\pgfpathlineto{\pgfqpoint{1.119998in}{0.712658in}}%
\pgfpathlineto{\pgfqpoint{1.219064in}{0.701043in}}%
\pgfpathlineto{\pgfqpoint{1.319441in}{0.691038in}}%
\pgfpathlineto{\pgfqpoint{1.322368in}{0.724970in}}%
\pgfpathlineto{\pgfqpoint{1.337455in}{0.871739in}}%
\pgfpathlineto{\pgfqpoint{1.347370in}{0.986543in}}%
\pgfpathlineto{\pgfqpoint{1.348953in}{0.986455in}}%
\pgfpathlineto{\pgfqpoint{1.449726in}{0.978343in}}%
\pgfpathlineto{\pgfqpoint{1.507338in}{0.975120in}}%
\pgfpathlineto{\pgfqpoint{1.500607in}{0.847917in}}%
\pgfpathlineto{\pgfqpoint{1.504154in}{0.848834in}}%
\pgfpathlineto{\pgfqpoint{1.516885in}{0.834817in}}%
\pgfpathlineto{\pgfqpoint{1.522821in}{0.836596in}}%
\pgfpathlineto{\pgfqpoint{1.531836in}{0.833517in}}%
\pgfpathlineto{\pgfqpoint{1.539706in}{0.835071in}}%
\pgfpathlineto{\pgfqpoint{1.544100in}{0.822590in}}%
\pgfpathlineto{\pgfqpoint{1.553918in}{0.821935in}}%
\pgfpathlineto{\pgfqpoint{1.569594in}{0.816058in}}%
\pgfpathlineto{\pgfqpoint{1.574003in}{0.817714in}}%
\pgfpathlineto{\pgfqpoint{1.583461in}{0.810461in}}%
\pgfpathlineto{\pgfqpoint{1.589028in}{0.816058in}}%
\pgfpathlineto{\pgfqpoint{1.598675in}{0.813942in}}%
\pgfpathlineto{\pgfqpoint{1.603167in}{0.806507in}}%
\pgfpathlineto{\pgfqpoint{1.610544in}{0.804002in}}%
\pgfpathlineto{\pgfqpoint{1.615661in}{0.796290in}}%
\pgfpathlineto{\pgfqpoint{1.625879in}{0.804682in}}%
\pgfpathlineto{\pgfqpoint{1.637572in}{0.794271in}}%
\pgfpathlineto{\pgfqpoint{1.650942in}{0.798523in}}%
\pgfpathlineto{\pgfqpoint{1.657711in}{0.787562in}}%
\pgfpathlineto{\pgfqpoint{1.664329in}{0.800308in}}%
\pgfpathlineto{\pgfqpoint{1.676227in}{0.793081in}}%
\pgfpathlineto{\pgfqpoint{1.683661in}{0.795004in}}%
\pgfpathlineto{\pgfqpoint{1.694275in}{0.789894in}}%
\pgfpathlineto{\pgfqpoint{1.708036in}{0.788443in}}%
\pgfpathlineto{\pgfqpoint{1.721080in}{0.796307in}}%
\pgfpathlineto{\pgfqpoint{1.725799in}{0.792988in}}%
\pgfpathlineto{\pgfqpoint{1.739055in}{0.799728in}}%
\pgfpathlineto{\pgfqpoint{1.747503in}{0.794467in}}%
\pgfpathlineto{\pgfqpoint{1.754180in}{0.794780in}}%
\pgfpathlineto{\pgfqpoint{1.759467in}{0.800976in}}%
\pgfpathlineto{\pgfqpoint{1.767788in}{0.797893in}}%
\pgfpathlineto{\pgfqpoint{1.786521in}{0.783547in}}%
\pgfpathlineto{\pgfqpoint{1.801141in}{0.779025in}}%
\pgfpathlineto{\pgfqpoint{1.807003in}{0.773486in}}%
\pgfpathlineto{\pgfqpoint{1.814342in}{0.776503in}}%
\pgfpathlineto{\pgfqpoint{1.825462in}{0.774193in}}%
\pgfpathlineto{\pgfqpoint{1.826608in}{0.670746in}}%
\pgfpathlineto{\pgfqpoint{1.834331in}{0.664193in}}%
\pgfpathlineto{\pgfqpoint{1.840521in}{0.652119in}}%
\pgfpathlineto{\pgfqpoint{1.838345in}{0.644041in}}%
\pgfpathlineto{\pgfqpoint{1.843206in}{0.640621in}}%
\pgfpathlineto{\pgfqpoint{1.847157in}{0.625724in}}%
\pgfpathlineto{\pgfqpoint{1.854722in}{0.617746in}}%
\pgfpathlineto{\pgfqpoint{1.856418in}{0.600809in}}%
\pgfpathlineto{\pgfqpoint{1.855102in}{0.593637in}}%
\pgfpathlineto{\pgfqpoint{1.844817in}{0.574659in}}%
\pgfpathlineto{\pgfqpoint{1.843635in}{0.561013in}}%
\pgfpathlineto{\pgfqpoint{1.847123in}{0.558165in}}%
\pgfpathlineto{\pgfqpoint{1.847427in}{0.546210in}}%
\pgfpathlineto{\pgfqpoint{1.843881in}{0.538697in}}%
\pgfpathlineto{\pgfqpoint{1.838366in}{0.535867in}}%
\pgfpathlineto{\pgfqpoint{1.832966in}{0.526048in}}%
\pgfpathlineto{\pgfqpoint{1.839848in}{0.516539in}}%
\pgfpathlineto{\pgfqpoint{1.826489in}{0.516351in}}%
\pgfpathlineto{\pgfqpoint{1.790766in}{0.500008in}}%
\pgfpathlineto{\pgfqpoint{1.797587in}{0.509789in}}%
\pgfpathlineto{\pgfqpoint{1.789328in}{0.515062in}}%
\pgfpathlineto{\pgfqpoint{1.787602in}{0.524028in}}%
\pgfpathlineto{\pgfqpoint{1.771483in}{0.508411in}}%
\pgfpathlineto{\pgfqpoint{1.778637in}{0.497737in}}%
\pgfpathlineto{\pgfqpoint{1.768432in}{0.484145in}}%
\pgfpathlineto{\pgfqpoint{1.762940in}{0.484430in}}%
\pgfpathlineto{\pgfqpoint{1.757774in}{0.469628in}}%
\pgfpathlineto{\pgfqpoint{1.741435in}{0.457985in}}%
\pgfpathlineto{\pgfqpoint{1.726171in}{0.453789in}}%
\pgfpathlineto{\pgfqpoint{1.699548in}{0.442873in}}%
\pgfpathlineto{\pgfqpoint{1.696789in}{0.448962in}}%
\pgfpathlineto{\pgfqpoint{1.680205in}{0.443357in}}%
\pgfpathlineto{\pgfqpoint{1.690405in}{0.433814in}}%
\pgfpathlineto{\pgfqpoint{1.674316in}{0.425524in}}%
\pgfpathlineto{\pgfqpoint{1.656962in}{0.413006in}}%
\pgfpathlineto{\pgfqpoint{1.652794in}{0.418825in}}%
\pgfpathlineto{\pgfqpoint{1.638592in}{0.410104in}}%
\pgfpathlineto{\pgfqpoint{1.652512in}{0.406792in}}%
\pgfpathlineto{\pgfqpoint{1.642016in}{0.393080in}}%
\pgfpathlineto{\pgfqpoint{1.636893in}{0.397161in}}%
\pgfpathlineto{\pgfqpoint{1.630012in}{0.390606in}}%
\pgfpathlineto{\pgfqpoint{1.638588in}{0.384881in}}%
\pgfpathlineto{\pgfqpoint{1.633518in}{0.376512in}}%
\pgfpathlineto{\pgfqpoint{1.627286in}{0.356702in}}%
\pgfpathlineto{\pgfqpoint{1.618297in}{0.337622in}}%
\pgfpathlineto{\pgfqpoint{1.622332in}{0.325116in}}%
\pgfpathlineto{\pgfqpoint{1.629192in}{0.295630in}}%
\pgfpathlineto{\pgfqpoint{1.629809in}{0.283697in}}%
\pgfpathlineto{\pgfqpoint{1.640407in}{0.268045in}}%
\pgfpathlineto{\pgfqpoint{1.632235in}{0.268959in}}%
\pgfpathlineto{\pgfqpoint{1.624297in}{0.261063in}}%
\pgfpathlineto{\pgfqpoint{1.615000in}{0.267269in}}%
\pgfpathlineto{\pgfqpoint{1.611663in}{0.273460in}}%
\pgfpathlineto{\pgfqpoint{1.598432in}{0.276340in}}%
\pgfpathlineto{\pgfqpoint{1.578220in}{0.276693in}}%
\pgfpathlineto{\pgfqpoint{1.563329in}{0.288418in}}%
\pgfpathlineto{\pgfqpoint{1.549807in}{0.290376in}}%
\pgfpathlineto{\pgfqpoint{1.541604in}{0.299694in}}%
\pgfpathlineto{\pgfqpoint{1.524424in}{0.303444in}}%
\pgfpathlineto{\pgfqpoint{1.515030in}{0.333403in}}%
\pgfpathlineto{\pgfqpoint{1.505425in}{0.345405in}}%
\pgfpathlineto{\pgfqpoint{1.507057in}{0.356809in}}%
\pgfpathlineto{\pgfqpoint{1.501101in}{0.365143in}}%
\pgfpathlineto{\pgfqpoint{1.504839in}{0.376534in}}%
\pgfpathlineto{\pgfqpoint{1.501755in}{0.384882in}}%
\pgfpathlineto{\pgfqpoint{1.492103in}{0.388664in}}%
\pgfpathlineto{\pgfqpoint{1.483078in}{0.398289in}}%
\pgfpathlineto{\pgfqpoint{1.479790in}{0.411177in}}%
\pgfpathlineto{\pgfqpoint{1.471250in}{0.422889in}}%
\pgfpathlineto{\pgfqpoint{1.459883in}{0.432001in}}%
\pgfpathlineto{\pgfqpoint{1.449644in}{0.458142in}}%
\pgfpathlineto{\pgfqpoint{1.441377in}{0.486731in}}%
\pgfpathlineto{\pgfqpoint{1.434594in}{0.498034in}}%
\pgfpathlineto{\pgfqpoint{1.422834in}{0.507551in}}%
\pgfpathlineto{\pgfqpoint{1.419895in}{0.514459in}}%
\pgfpathlineto{\pgfqpoint{1.408884in}{0.518748in}}%
\pgfpathlineto{\pgfqpoint{1.398016in}{0.537050in}}%
\pgfpathlineto{\pgfqpoint{1.388091in}{0.535645in}}%
\pgfpathlineto{\pgfqpoint{1.378014in}{0.540216in}}%
\pgfpathlineto{\pgfqpoint{1.363728in}{0.539332in}}%
\pgfpathlineto{\pgfqpoint{1.349216in}{0.546876in}}%
\pgfpathlineto{\pgfqpoint{1.345123in}{0.539702in}}%
\pgfpathlineto{\pgfqpoint{1.328163in}{0.539522in}}%
\pgfpathlineto{\pgfqpoint{1.319529in}{0.525922in}}%
\pgfpathlineto{\pgfqpoint{1.312019in}{0.509082in}}%
\pgfpathlineto{\pgfqpoint{1.296050in}{0.491006in}}%
\pgfpathlineto{\pgfqpoint{1.277957in}{0.498939in}}%
\pgfpathlineto{\pgfqpoint{1.275390in}{0.504181in}}%
\pgfpathlineto{\pgfqpoint{1.264395in}{0.508180in}}%
\pgfpathlineto{\pgfqpoint{1.261017in}{0.513650in}}%
\pgfpathlineto{\pgfqpoint{1.246421in}{0.519171in}}%
\pgfpathlineto{\pgfqpoint{1.238256in}{0.530458in}}%
\pgfpathlineto{\pgfqpoint{1.228720in}{0.535904in}}%
\pgfpathlineto{\pgfqpoint{1.220515in}{0.545406in}}%
\pgfpathlineto{\pgfqpoint{1.214124in}{0.561500in}}%
\pgfpathlineto{\pgfqpoint{1.214839in}{0.583492in}}%
\pgfpathlineto{\pgfqpoint{1.207341in}{0.594605in}}%
\pgfpathlineto{\pgfqpoint{1.206478in}{0.606641in}}%
\pgfpathlineto{\pgfqpoint{1.189842in}{0.624662in}}%
\pgfpathlineto{\pgfqpoint{1.180162in}{0.628509in}}%
\pgfpathlineto{\pgfqpoint{1.169890in}{0.645324in}}%
\pgfpathlineto{\pgfqpoint{1.161136in}{0.652005in}}%
\pgfpathlineto{\pgfqpoint{1.150014in}{0.668274in}}%
\pgfpathlineto{\pgfqpoint{1.138634in}{0.675314in}}%
\pgfpathlineto{\pgfqpoint{1.131182in}{0.693355in}}%
\pgfpathlineto{\pgfqpoint{1.122887in}{0.697915in}}%
\pgfpathlineto{\pgfqpoint{1.122887in}{0.697915in}}%
\pgfusepath{stroke}%
\end{pgfscope}%
\begin{pgfscope}%
\pgfpathrectangle{\pgfqpoint{0.100000in}{0.100000in}}{\pgfqpoint{3.007045in}{1.925000in}}%
\pgfusepath{clip}%
\pgfsetbuttcap%
\pgfsetroundjoin%
\pgfsetlinewidth{0.050187pt}%
\definecolor{currentstroke}{rgb}{1.000000,1.000000,1.000000}%
\pgfsetstrokecolor{currentstroke}%
\pgfsetdash{}{0pt}%
\pgfpathmoveto{\pgfqpoint{1.035744in}{1.058103in}}%
\pgfpathlineto{\pgfqpoint{1.117842in}{1.046179in}}%
\pgfpathlineto{\pgfqpoint{1.149648in}{1.041279in}}%
\pgfpathlineto{\pgfqpoint{1.239303in}{1.030616in}}%
\pgfpathlineto{\pgfqpoint{1.299648in}{1.024321in}}%
\pgfpathlineto{\pgfqpoint{1.351869in}{1.019693in}}%
\pgfpathlineto{\pgfqpoint{1.348953in}{0.986455in}}%
\pgfpathlineto{\pgfqpoint{1.347370in}{0.986543in}}%
\pgfpathlineto{\pgfqpoint{1.342587in}{0.929495in}}%
\pgfpathlineto{\pgfqpoint{1.337455in}{0.871739in}}%
\pgfpathlineto{\pgfqpoint{1.331519in}{0.811277in}}%
\pgfpathlineto{\pgfqpoint{1.322368in}{0.724970in}}%
\pgfpathlineto{\pgfqpoint{1.319441in}{0.691038in}}%
\pgfpathlineto{\pgfqpoint{1.265644in}{0.696466in}}%
\pgfpathlineto{\pgfqpoint{1.219064in}{0.701043in}}%
\pgfpathlineto{\pgfqpoint{1.154643in}{0.708360in}}%
\pgfpathlineto{\pgfqpoint{1.119998in}{0.712658in}}%
\pgfpathlineto{\pgfqpoint{1.117958in}{0.705835in}}%
\pgfpathlineto{\pgfqpoint{1.122887in}{0.697915in}}%
\pgfpathlineto{\pgfqpoint{1.081250in}{0.703323in}}%
\pgfpathlineto{\pgfqpoint{1.029878in}{0.710696in}}%
\pgfpathlineto{\pgfqpoint{1.025202in}{0.681640in}}%
\pgfpathlineto{\pgfqpoint{0.978343in}{0.688693in}}%
\pgfpathlineto{\pgfqpoint{0.989353in}{0.760774in}}%
\pgfpathlineto{\pgfqpoint{1.002983in}{0.848537in}}%
\pgfpathlineto{\pgfqpoint{1.015102in}{0.924993in}}%
\pgfpathlineto{\pgfqpoint{1.025626in}{0.992952in}}%
\pgfpathlineto{\pgfqpoint{1.035744in}{1.058103in}}%
\pgfusepath{stroke}%
\end{pgfscope}%
\begin{pgfscope}%
\pgfpathrectangle{\pgfqpoint{0.100000in}{0.100000in}}{\pgfqpoint{3.007045in}{1.925000in}}%
\pgfusepath{clip}%
\pgfsetbuttcap%
\pgfsetroundjoin%
\pgfsetlinewidth{0.050187pt}%
\definecolor{currentstroke}{rgb}{1.000000,1.000000,1.000000}%
\pgfsetstrokecolor{currentstroke}%
\pgfsetdash{}{0pt}%
\pgfpathmoveto{\pgfqpoint{2.341115in}{0.635787in}}%
\pgfpathlineto{\pgfqpoint{2.262793in}{0.627088in}}%
\pgfpathlineto{\pgfqpoint{2.193752in}{0.621728in}}%
\pgfpathlineto{\pgfqpoint{2.192889in}{0.613268in}}%
\pgfpathlineto{\pgfqpoint{2.199230in}{0.605215in}}%
\pgfpathlineto{\pgfqpoint{2.206971in}{0.600484in}}%
\pgfpathlineto{\pgfqpoint{2.206853in}{0.588027in}}%
\pgfpathlineto{\pgfqpoint{2.204803in}{0.579705in}}%
\pgfpathlineto{\pgfqpoint{2.197971in}{0.573707in}}%
\pgfpathlineto{\pgfqpoint{2.188449in}{0.574342in}}%
\pgfpathlineto{\pgfqpoint{2.179446in}{0.581774in}}%
\pgfpathlineto{\pgfqpoint{2.177842in}{0.594998in}}%
\pgfpathlineto{\pgfqpoint{2.171137in}{0.602689in}}%
\pgfpathlineto{\pgfqpoint{2.166573in}{0.575162in}}%
\pgfpathlineto{\pgfqpoint{2.151065in}{0.577778in}}%
\pgfpathlineto{\pgfqpoint{2.145672in}{0.625913in}}%
\pgfpathlineto{\pgfqpoint{2.139832in}{0.676640in}}%
\pgfpathlineto{\pgfqpoint{2.142001in}{0.749828in}}%
\pgfpathlineto{\pgfqpoint{2.143347in}{0.800051in}}%
\pgfpathlineto{\pgfqpoint{2.146210in}{0.876679in}}%
\pgfpathlineto{\pgfqpoint{2.140081in}{0.883422in}}%
\pgfpathlineto{\pgfqpoint{2.151235in}{0.885039in}}%
\pgfpathlineto{\pgfqpoint{2.216213in}{0.889241in}}%
\pgfpathlineto{\pgfqpoint{2.279433in}{0.894926in}}%
\pgfpathlineto{\pgfqpoint{2.296066in}{0.836665in}}%
\pgfpathlineto{\pgfqpoint{2.307410in}{0.793901in}}%
\pgfpathlineto{\pgfqpoint{2.317578in}{0.758099in}}%
\pgfpathlineto{\pgfqpoint{2.323459in}{0.743683in}}%
\pgfpathlineto{\pgfqpoint{2.332718in}{0.730267in}}%
\pgfpathlineto{\pgfqpoint{2.331222in}{0.723440in}}%
\pgfpathlineto{\pgfqpoint{2.337837in}{0.720114in}}%
\pgfpathlineto{\pgfqpoint{2.330020in}{0.709765in}}%
\pgfpathlineto{\pgfqpoint{2.327380in}{0.691335in}}%
\pgfpathlineto{\pgfqpoint{2.335019in}{0.669779in}}%
\pgfpathlineto{\pgfqpoint{2.333958in}{0.648089in}}%
\pgfpathlineto{\pgfqpoint{2.341115in}{0.635787in}}%
\pgfusepath{stroke}%
\end{pgfscope}%
\begin{pgfscope}%
\pgfpathrectangle{\pgfqpoint{0.100000in}{0.100000in}}{\pgfqpoint{3.007045in}{1.925000in}}%
\pgfusepath{clip}%
\pgfsetbuttcap%
\pgfsetroundjoin%
\pgfsetlinewidth{0.050187pt}%
\definecolor{currentstroke}{rgb}{1.000000,1.000000,1.000000}%
\pgfsetstrokecolor{currentstroke}%
\pgfsetdash{}{0pt}%
\pgfpathmoveto{\pgfqpoint{2.151065in}{0.577778in}}%
\pgfpathlineto{\pgfqpoint{2.135151in}{0.573224in}}%
\pgfpathlineto{\pgfqpoint{2.123289in}{0.576148in}}%
\pgfpathlineto{\pgfqpoint{2.101259in}{0.569141in}}%
\pgfpathlineto{\pgfqpoint{2.084645in}{0.560117in}}%
\pgfpathlineto{\pgfqpoint{2.077828in}{0.576424in}}%
\pgfpathlineto{\pgfqpoint{2.070801in}{0.583259in}}%
\pgfpathlineto{\pgfqpoint{2.067411in}{0.591438in}}%
\pgfpathlineto{\pgfqpoint{2.072380in}{0.613576in}}%
\pgfpathlineto{\pgfqpoint{2.025296in}{0.610677in}}%
\pgfpathlineto{\pgfqpoint{1.964244in}{0.608030in}}%
\pgfpathlineto{\pgfqpoint{1.967876in}{0.613540in}}%
\pgfpathlineto{\pgfqpoint{1.963386in}{0.623948in}}%
\pgfpathlineto{\pgfqpoint{1.967860in}{0.626069in}}%
\pgfpathlineto{\pgfqpoint{1.966867in}{0.636068in}}%
\pgfpathlineto{\pgfqpoint{1.974665in}{0.645765in}}%
\pgfpathlineto{\pgfqpoint{1.978931in}{0.664590in}}%
\pgfpathlineto{\pgfqpoint{1.993202in}{0.676993in}}%
\pgfpathlineto{\pgfqpoint{1.987913in}{0.689038in}}%
\pgfpathlineto{\pgfqpoint{1.998033in}{0.695971in}}%
\pgfpathlineto{\pgfqpoint{1.989303in}{0.708523in}}%
\pgfpathlineto{\pgfqpoint{1.982985in}{0.737257in}}%
\pgfpathlineto{\pgfqpoint{1.985381in}{0.742474in}}%
\pgfpathlineto{\pgfqpoint{1.989164in}{0.752482in}}%
\pgfpathlineto{\pgfqpoint{1.985643in}{0.762985in}}%
\pgfpathlineto{\pgfqpoint{1.987365in}{0.767763in}}%
\pgfpathlineto{\pgfqpoint{1.981171in}{0.785815in}}%
\pgfpathlineto{\pgfqpoint{1.999458in}{0.818229in}}%
\pgfpathlineto{\pgfqpoint{2.000308in}{0.826857in}}%
\pgfpathlineto{\pgfqpoint{2.009117in}{0.833130in}}%
\pgfpathlineto{\pgfqpoint{2.018516in}{0.853839in}}%
\pgfpathlineto{\pgfqpoint{2.018070in}{0.862257in}}%
\pgfpathlineto{\pgfqpoint{2.029777in}{0.870862in}}%
\pgfpathlineto{\pgfqpoint{2.026124in}{0.875959in}}%
\pgfpathlineto{\pgfqpoint{2.096280in}{0.880053in}}%
\pgfpathlineto{\pgfqpoint{2.140081in}{0.883422in}}%
\pgfpathlineto{\pgfqpoint{2.146210in}{0.876679in}}%
\pgfpathlineto{\pgfqpoint{2.143347in}{0.800051in}}%
\pgfpathlineto{\pgfqpoint{2.142001in}{0.749828in}}%
\pgfpathlineto{\pgfqpoint{2.139832in}{0.676640in}}%
\pgfpathlineto{\pgfqpoint{2.145672in}{0.625913in}}%
\pgfpathlineto{\pgfqpoint{2.151065in}{0.577778in}}%
\pgfusepath{stroke}%
\end{pgfscope}%
\begin{pgfscope}%
\pgfpathrectangle{\pgfqpoint{0.100000in}{0.100000in}}{\pgfqpoint{3.007045in}{1.925000in}}%
\pgfusepath{clip}%
\pgfsetbuttcap%
\pgfsetroundjoin%
\pgfsetlinewidth{0.050187pt}%
\definecolor{currentstroke}{rgb}{1.000000,1.000000,1.000000}%
\pgfsetstrokecolor{currentstroke}%
\pgfsetdash{}{0pt}%
\pgfpathmoveto{\pgfqpoint{2.279433in}{0.894926in}}%
\pgfpathlineto{\pgfqpoint{2.348542in}{0.902516in}}%
\pgfpathlineto{\pgfqpoint{2.390124in}{0.907640in}}%
\pgfpathlineto{\pgfqpoint{2.413807in}{0.911413in}}%
\pgfpathlineto{\pgfqpoint{2.403997in}{0.895940in}}%
\pgfpathlineto{\pgfqpoint{2.404003in}{0.888645in}}%
\pgfpathlineto{\pgfqpoint{2.414101in}{0.884711in}}%
\pgfpathlineto{\pgfqpoint{2.420968in}{0.878359in}}%
\pgfpathlineto{\pgfqpoint{2.431346in}{0.877574in}}%
\pgfpathlineto{\pgfqpoint{2.441046in}{0.859694in}}%
\pgfpathlineto{\pgfqpoint{2.451568in}{0.847085in}}%
\pgfpathlineto{\pgfqpoint{2.470255in}{0.836488in}}%
\pgfpathlineto{\pgfqpoint{2.475566in}{0.828036in}}%
\pgfpathlineto{\pgfqpoint{2.492230in}{0.818895in}}%
\pgfpathlineto{\pgfqpoint{2.493483in}{0.812266in}}%
\pgfpathlineto{\pgfqpoint{2.507001in}{0.799084in}}%
\pgfpathlineto{\pgfqpoint{2.517356in}{0.795331in}}%
\pgfpathlineto{\pgfqpoint{2.524689in}{0.782884in}}%
\pgfpathlineto{\pgfqpoint{2.527935in}{0.768894in}}%
\pgfpathlineto{\pgfqpoint{2.538690in}{0.763486in}}%
\pgfpathlineto{\pgfqpoint{2.547349in}{0.748448in}}%
\pgfpathlineto{\pgfqpoint{2.550141in}{0.737389in}}%
\pgfpathlineto{\pgfqpoint{2.562829in}{0.732681in}}%
\pgfpathlineto{\pgfqpoint{2.552182in}{0.712295in}}%
\pgfpathlineto{\pgfqpoint{2.548176in}{0.700248in}}%
\pgfpathlineto{\pgfqpoint{2.550790in}{0.694600in}}%
\pgfpathlineto{\pgfqpoint{2.546225in}{0.685215in}}%
\pgfpathlineto{\pgfqpoint{2.544285in}{0.672239in}}%
\pgfpathlineto{\pgfqpoint{2.536167in}{0.669822in}}%
\pgfpathlineto{\pgfqpoint{2.540454in}{0.657942in}}%
\pgfpathlineto{\pgfqpoint{2.540180in}{0.642869in}}%
\pgfpathlineto{\pgfqpoint{2.525751in}{0.643434in}}%
\pgfpathlineto{\pgfqpoint{2.515700in}{0.646184in}}%
\pgfpathlineto{\pgfqpoint{2.511350in}{0.641072in}}%
\pgfpathlineto{\pgfqpoint{2.514596in}{0.628996in}}%
\pgfpathlineto{\pgfqpoint{2.513890in}{0.614955in}}%
\pgfpathlineto{\pgfqpoint{2.507549in}{0.613881in}}%
\pgfpathlineto{\pgfqpoint{2.502419in}{0.626994in}}%
\pgfpathlineto{\pgfqpoint{2.450185in}{0.623513in}}%
\pgfpathlineto{\pgfqpoint{2.384322in}{0.619919in}}%
\pgfpathlineto{\pgfqpoint{2.351105in}{0.617581in}}%
\pgfpathlineto{\pgfqpoint{2.341115in}{0.635787in}}%
\pgfpathlineto{\pgfqpoint{2.333958in}{0.648089in}}%
\pgfpathlineto{\pgfqpoint{2.335019in}{0.669779in}}%
\pgfpathlineto{\pgfqpoint{2.327380in}{0.691335in}}%
\pgfpathlineto{\pgfqpoint{2.330020in}{0.709765in}}%
\pgfpathlineto{\pgfqpoint{2.337837in}{0.720114in}}%
\pgfpathlineto{\pgfqpoint{2.331222in}{0.723440in}}%
\pgfpathlineto{\pgfqpoint{2.332718in}{0.730267in}}%
\pgfpathlineto{\pgfqpoint{2.323459in}{0.743683in}}%
\pgfpathlineto{\pgfqpoint{2.317578in}{0.758099in}}%
\pgfpathlineto{\pgfqpoint{2.307410in}{0.793901in}}%
\pgfpathlineto{\pgfqpoint{2.296066in}{0.836665in}}%
\pgfpathlineto{\pgfqpoint{2.279433in}{0.894926in}}%
\pgfusepath{stroke}%
\end{pgfscope}%
\begin{pgfscope}%
\pgfpathrectangle{\pgfqpoint{0.100000in}{0.100000in}}{\pgfqpoint{3.007045in}{1.925000in}}%
\pgfusepath{clip}%
\pgfsetbuttcap%
\pgfsetroundjoin%
\pgfsetlinewidth{0.050187pt}%
\definecolor{currentstroke}{rgb}{1.000000,1.000000,1.000000}%
\pgfsetstrokecolor{currentstroke}%
\pgfsetdash{}{0pt}%
\pgfpathmoveto{\pgfqpoint{2.413807in}{0.911413in}}%
\pgfpathlineto{\pgfqpoint{2.441211in}{0.925045in}}%
\pgfpathlineto{\pgfqpoint{2.456388in}{0.930193in}}%
\pgfpathlineto{\pgfqpoint{2.522418in}{0.937054in}}%
\pgfpathlineto{\pgfqpoint{2.529411in}{0.934781in}}%
\pgfpathlineto{\pgfqpoint{2.538649in}{0.925496in}}%
\pgfpathlineto{\pgfqpoint{2.539158in}{0.917251in}}%
\pgfpathlineto{\pgfqpoint{2.599931in}{0.926180in}}%
\pgfpathlineto{\pgfqpoint{2.668941in}{0.876711in}}%
\pgfpathlineto{\pgfqpoint{2.655999in}{0.863222in}}%
\pgfpathlineto{\pgfqpoint{2.638241in}{0.831896in}}%
\pgfpathlineto{\pgfqpoint{2.643296in}{0.825164in}}%
\pgfpathlineto{\pgfqpoint{2.633823in}{0.812096in}}%
\pgfpathlineto{\pgfqpoint{2.624400in}{0.810617in}}%
\pgfpathlineto{\pgfqpoint{2.624365in}{0.802765in}}%
\pgfpathlineto{\pgfqpoint{2.617549in}{0.794548in}}%
\pgfpathlineto{\pgfqpoint{2.609066in}{0.792881in}}%
\pgfpathlineto{\pgfqpoint{2.610917in}{0.785581in}}%
\pgfpathlineto{\pgfqpoint{2.606183in}{0.779976in}}%
\pgfpathlineto{\pgfqpoint{2.594860in}{0.775134in}}%
\pgfpathlineto{\pgfqpoint{2.581058in}{0.764120in}}%
\pgfpathlineto{\pgfqpoint{2.583663in}{0.756995in}}%
\pgfpathlineto{\pgfqpoint{2.574977in}{0.752521in}}%
\pgfpathlineto{\pgfqpoint{2.567651in}{0.757225in}}%
\pgfpathlineto{\pgfqpoint{2.562294in}{0.736773in}}%
\pgfpathlineto{\pgfqpoint{2.550141in}{0.737389in}}%
\pgfpathlineto{\pgfqpoint{2.547349in}{0.748448in}}%
\pgfpathlineto{\pgfqpoint{2.538690in}{0.763486in}}%
\pgfpathlineto{\pgfqpoint{2.527935in}{0.768894in}}%
\pgfpathlineto{\pgfqpoint{2.524689in}{0.782884in}}%
\pgfpathlineto{\pgfqpoint{2.517356in}{0.795331in}}%
\pgfpathlineto{\pgfqpoint{2.507001in}{0.799084in}}%
\pgfpathlineto{\pgfqpoint{2.493483in}{0.812266in}}%
\pgfpathlineto{\pgfqpoint{2.492230in}{0.818895in}}%
\pgfpathlineto{\pgfqpoint{2.475566in}{0.828036in}}%
\pgfpathlineto{\pgfqpoint{2.470255in}{0.836488in}}%
\pgfpathlineto{\pgfqpoint{2.451568in}{0.847085in}}%
\pgfpathlineto{\pgfqpoint{2.441046in}{0.859694in}}%
\pgfpathlineto{\pgfqpoint{2.431346in}{0.877574in}}%
\pgfpathlineto{\pgfqpoint{2.420968in}{0.878359in}}%
\pgfpathlineto{\pgfqpoint{2.414101in}{0.884711in}}%
\pgfpathlineto{\pgfqpoint{2.404003in}{0.888645in}}%
\pgfpathlineto{\pgfqpoint{2.403997in}{0.895940in}}%
\pgfpathlineto{\pgfqpoint{2.413807in}{0.911413in}}%
\pgfusepath{stroke}%
\end{pgfscope}%
\begin{pgfscope}%
\pgfpathrectangle{\pgfqpoint{0.100000in}{0.100000in}}{\pgfqpoint{3.007045in}{1.925000in}}%
\pgfusepath{clip}%
\pgfsetbuttcap%
\pgfsetroundjoin%
\pgfsetlinewidth{0.050187pt}%
\definecolor{currentstroke}{rgb}{1.000000,1.000000,1.000000}%
\pgfsetstrokecolor{currentstroke}%
\pgfsetdash{}{0pt}%
\pgfpathmoveto{\pgfqpoint{1.792603in}{0.967662in}}%
\pgfpathlineto{\pgfqpoint{1.832874in}{0.968025in}}%
\pgfpathlineto{\pgfqpoint{1.886150in}{0.968996in}}%
\pgfpathlineto{\pgfqpoint{1.977188in}{0.971886in}}%
\pgfpathlineto{\pgfqpoint{2.029259in}{0.974639in}}%
\pgfpathlineto{\pgfqpoint{2.034844in}{0.967726in}}%
\pgfpathlineto{\pgfqpoint{2.034495in}{0.960424in}}%
\pgfpathlineto{\pgfqpoint{2.021881in}{0.947828in}}%
\pgfpathlineto{\pgfqpoint{2.018831in}{0.940931in}}%
\pgfpathlineto{\pgfqpoint{2.053816in}{0.943518in}}%
\pgfpathlineto{\pgfqpoint{2.053797in}{0.930775in}}%
\pgfpathlineto{\pgfqpoint{2.042574in}{0.925329in}}%
\pgfpathlineto{\pgfqpoint{2.042732in}{0.916670in}}%
\pgfpathlineto{\pgfqpoint{2.038517in}{0.910446in}}%
\pgfpathlineto{\pgfqpoint{2.035758in}{0.897226in}}%
\pgfpathlineto{\pgfqpoint{2.038606in}{0.886320in}}%
\pgfpathlineto{\pgfqpoint{2.026124in}{0.875959in}}%
\pgfpathlineto{\pgfqpoint{2.029777in}{0.870862in}}%
\pgfpathlineto{\pgfqpoint{2.018070in}{0.862257in}}%
\pgfpathlineto{\pgfqpoint{2.018516in}{0.853839in}}%
\pgfpathlineto{\pgfqpoint{2.009117in}{0.833130in}}%
\pgfpathlineto{\pgfqpoint{2.000308in}{0.826857in}}%
\pgfpathlineto{\pgfqpoint{1.999458in}{0.818229in}}%
\pgfpathlineto{\pgfqpoint{1.981171in}{0.785815in}}%
\pgfpathlineto{\pgfqpoint{1.987365in}{0.767763in}}%
\pgfpathlineto{\pgfqpoint{1.985643in}{0.762985in}}%
\pgfpathlineto{\pgfqpoint{1.989164in}{0.752482in}}%
\pgfpathlineto{\pgfqpoint{1.985381in}{0.742474in}}%
\pgfpathlineto{\pgfqpoint{1.935375in}{0.740410in}}%
\pgfpathlineto{\pgfqpoint{1.870456in}{0.739340in}}%
\pgfpathlineto{\pgfqpoint{1.825680in}{0.738936in}}%
\pgfpathlineto{\pgfqpoint{1.825462in}{0.774193in}}%
\pgfpathlineto{\pgfqpoint{1.814342in}{0.776503in}}%
\pgfpathlineto{\pgfqpoint{1.807003in}{0.773486in}}%
\pgfpathlineto{\pgfqpoint{1.801141in}{0.779025in}}%
\pgfpathlineto{\pgfqpoint{1.801746in}{0.816372in}}%
\pgfpathlineto{\pgfqpoint{1.803058in}{0.895889in}}%
\pgfpathlineto{\pgfqpoint{1.796697in}{0.942452in}}%
\pgfpathlineto{\pgfqpoint{1.792603in}{0.967662in}}%
\pgfusepath{stroke}%
\end{pgfscope}%
\begin{pgfscope}%
\pgfpathrectangle{\pgfqpoint{0.100000in}{0.100000in}}{\pgfqpoint{3.007045in}{1.925000in}}%
\pgfusepath{clip}%
\pgfsetbuttcap%
\pgfsetroundjoin%
\pgfsetlinewidth{0.050187pt}%
\definecolor{currentstroke}{rgb}{1.000000,1.000000,1.000000}%
\pgfsetstrokecolor{currentstroke}%
\pgfsetdash{}{0pt}%
\pgfpathmoveto{\pgfqpoint{1.825680in}{0.738936in}}%
\pgfpathlineto{\pgfqpoint{1.870456in}{0.739340in}}%
\pgfpathlineto{\pgfqpoint{1.935375in}{0.740410in}}%
\pgfpathlineto{\pgfqpoint{1.985381in}{0.742474in}}%
\pgfpathlineto{\pgfqpoint{1.982985in}{0.737257in}}%
\pgfpathlineto{\pgfqpoint{1.989303in}{0.708523in}}%
\pgfpathlineto{\pgfqpoint{1.998033in}{0.695971in}}%
\pgfpathlineto{\pgfqpoint{1.987913in}{0.689038in}}%
\pgfpathlineto{\pgfqpoint{1.993202in}{0.676993in}}%
\pgfpathlineto{\pgfqpoint{1.978931in}{0.664590in}}%
\pgfpathlineto{\pgfqpoint{1.974665in}{0.645765in}}%
\pgfpathlineto{\pgfqpoint{1.966867in}{0.636068in}}%
\pgfpathlineto{\pgfqpoint{1.967860in}{0.626069in}}%
\pgfpathlineto{\pgfqpoint{1.963386in}{0.623948in}}%
\pgfpathlineto{\pgfqpoint{1.967876in}{0.613540in}}%
\pgfpathlineto{\pgfqpoint{1.964244in}{0.608030in}}%
\pgfpathlineto{\pgfqpoint{2.025296in}{0.610677in}}%
\pgfpathlineto{\pgfqpoint{2.072380in}{0.613576in}}%
\pgfpathlineto{\pgfqpoint{2.067411in}{0.591438in}}%
\pgfpathlineto{\pgfqpoint{2.070801in}{0.583259in}}%
\pgfpathlineto{\pgfqpoint{2.077828in}{0.576424in}}%
\pgfpathlineto{\pgfqpoint{2.084645in}{0.560117in}}%
\pgfpathlineto{\pgfqpoint{2.063113in}{0.563874in}}%
\pgfpathlineto{\pgfqpoint{2.055170in}{0.570050in}}%
\pgfpathlineto{\pgfqpoint{2.045706in}{0.570344in}}%
\pgfpathlineto{\pgfqpoint{2.035757in}{0.556820in}}%
\pgfpathlineto{\pgfqpoint{2.037740in}{0.550661in}}%
\pgfpathlineto{\pgfqpoint{2.069524in}{0.546935in}}%
\pgfpathlineto{\pgfqpoint{2.077845in}{0.539855in}}%
\pgfpathlineto{\pgfqpoint{2.091386in}{0.536213in}}%
\pgfpathlineto{\pgfqpoint{2.085472in}{0.527832in}}%
\pgfpathlineto{\pgfqpoint{2.075509in}{0.520522in}}%
\pgfpathlineto{\pgfqpoint{2.101808in}{0.504085in}}%
\pgfpathlineto{\pgfqpoint{2.114058in}{0.501516in}}%
\pgfpathlineto{\pgfqpoint{2.119152in}{0.488080in}}%
\pgfpathlineto{\pgfqpoint{2.107618in}{0.485609in}}%
\pgfpathlineto{\pgfqpoint{2.085940in}{0.502507in}}%
\pgfpathlineto{\pgfqpoint{2.077465in}{0.504856in}}%
\pgfpathlineto{\pgfqpoint{2.073363in}{0.511532in}}%
\pgfpathlineto{\pgfqpoint{2.060791in}{0.508792in}}%
\pgfpathlineto{\pgfqpoint{2.056855in}{0.500190in}}%
\pgfpathlineto{\pgfqpoint{2.059348in}{0.490604in}}%
\pgfpathlineto{\pgfqpoint{2.050884in}{0.484948in}}%
\pgfpathlineto{\pgfqpoint{2.043160in}{0.498883in}}%
\pgfpathlineto{\pgfqpoint{2.029631in}{0.494728in}}%
\pgfpathlineto{\pgfqpoint{2.024562in}{0.486414in}}%
\pgfpathlineto{\pgfqpoint{2.014953in}{0.488779in}}%
\pgfpathlineto{\pgfqpoint{2.008810in}{0.499279in}}%
\pgfpathlineto{\pgfqpoint{1.992513in}{0.502746in}}%
\pgfpathlineto{\pgfqpoint{1.989436in}{0.508210in}}%
\pgfpathlineto{\pgfqpoint{1.972456in}{0.517727in}}%
\pgfpathlineto{\pgfqpoint{1.968241in}{0.526049in}}%
\pgfpathlineto{\pgfqpoint{1.954018in}{0.522648in}}%
\pgfpathlineto{\pgfqpoint{1.955825in}{0.530267in}}%
\pgfpathlineto{\pgfqpoint{1.938148in}{0.522451in}}%
\pgfpathlineto{\pgfqpoint{1.942896in}{0.514343in}}%
\pgfpathlineto{\pgfqpoint{1.929245in}{0.509548in}}%
\pgfpathlineto{\pgfqpoint{1.911163in}{0.512183in}}%
\pgfpathlineto{\pgfqpoint{1.874567in}{0.524719in}}%
\pgfpathlineto{\pgfqpoint{1.846328in}{0.522231in}}%
\pgfpathlineto{\pgfqpoint{1.843881in}{0.538697in}}%
\pgfpathlineto{\pgfqpoint{1.847427in}{0.546210in}}%
\pgfpathlineto{\pgfqpoint{1.847123in}{0.558165in}}%
\pgfpathlineto{\pgfqpoint{1.843635in}{0.561013in}}%
\pgfpathlineto{\pgfqpoint{1.844817in}{0.574659in}}%
\pgfpathlineto{\pgfqpoint{1.855102in}{0.593637in}}%
\pgfpathlineto{\pgfqpoint{1.856418in}{0.600809in}}%
\pgfpathlineto{\pgfqpoint{1.854722in}{0.617746in}}%
\pgfpathlineto{\pgfqpoint{1.847157in}{0.625724in}}%
\pgfpathlineto{\pgfqpoint{1.843206in}{0.640621in}}%
\pgfpathlineto{\pgfqpoint{1.838345in}{0.644041in}}%
\pgfpathlineto{\pgfqpoint{1.840521in}{0.652119in}}%
\pgfpathlineto{\pgfqpoint{1.834331in}{0.664193in}}%
\pgfpathlineto{\pgfqpoint{1.826608in}{0.670746in}}%
\pgfpathlineto{\pgfqpoint{1.825680in}{0.738936in}}%
\pgfusepath{stroke}%
\end{pgfscope}%
\begin{pgfscope}%
\pgfpathrectangle{\pgfqpoint{0.100000in}{0.100000in}}{\pgfqpoint{3.007045in}{1.925000in}}%
\pgfusepath{clip}%
\pgfsetbuttcap%
\pgfsetroundjoin%
\pgfsetlinewidth{0.050187pt}%
\definecolor{currentstroke}{rgb}{1.000000,1.000000,1.000000}%
\pgfsetstrokecolor{currentstroke}%
\pgfsetdash{}{0pt}%
\pgfpathmoveto{\pgfqpoint{1.945436in}{0.513722in}}%
\pgfpathlineto{\pgfqpoint{1.951913in}{0.517575in}}%
\pgfpathlineto{\pgfqpoint{1.959776in}{0.513019in}}%
\pgfpathlineto{\pgfqpoint{1.955391in}{0.506751in}}%
\pgfpathlineto{\pgfqpoint{1.945436in}{0.513722in}}%
\pgfusepath{stroke}%
\end{pgfscope}%
\begin{pgfscope}%
\pgfpathrectangle{\pgfqpoint{0.100000in}{0.100000in}}{\pgfqpoint{3.007045in}{1.925000in}}%
\pgfusepath{clip}%
\pgfsetbuttcap%
\pgfsetroundjoin%
\pgfsetlinewidth{0.050187pt}%
\definecolor{currentstroke}{rgb}{1.000000,1.000000,1.000000}%
\pgfsetstrokecolor{currentstroke}%
\pgfsetdash{}{0pt}%
\pgfpathmoveto{\pgfqpoint{2.206853in}{0.588027in}}%
\pgfpathlineto{\pgfqpoint{2.206971in}{0.600484in}}%
\pgfpathlineto{\pgfqpoint{2.199230in}{0.605215in}}%
\pgfpathlineto{\pgfqpoint{2.192889in}{0.613268in}}%
\pgfpathlineto{\pgfqpoint{2.193752in}{0.621728in}}%
\pgfpathlineto{\pgfqpoint{2.262793in}{0.627088in}}%
\pgfpathlineto{\pgfqpoint{2.341115in}{0.635787in}}%
\pgfpathlineto{\pgfqpoint{2.351105in}{0.617581in}}%
\pgfpathlineto{\pgfqpoint{2.384322in}{0.619919in}}%
\pgfpathlineto{\pgfqpoint{2.450185in}{0.623513in}}%
\pgfpathlineto{\pgfqpoint{2.502419in}{0.626994in}}%
\pgfpathlineto{\pgfqpoint{2.507549in}{0.613881in}}%
\pgfpathlineto{\pgfqpoint{2.513890in}{0.614955in}}%
\pgfpathlineto{\pgfqpoint{2.514596in}{0.628996in}}%
\pgfpathlineto{\pgfqpoint{2.511350in}{0.641072in}}%
\pgfpathlineto{\pgfqpoint{2.515700in}{0.646184in}}%
\pgfpathlineto{\pgfqpoint{2.525751in}{0.643434in}}%
\pgfpathlineto{\pgfqpoint{2.540180in}{0.642869in}}%
\pgfpathlineto{\pgfqpoint{2.542399in}{0.632080in}}%
\pgfpathlineto{\pgfqpoint{2.546837in}{0.625915in}}%
\pgfpathlineto{\pgfqpoint{2.550303in}{0.612426in}}%
\pgfpathlineto{\pgfqpoint{2.561042in}{0.591552in}}%
\pgfpathlineto{\pgfqpoint{2.561096in}{0.585906in}}%
\pgfpathlineto{\pgfqpoint{2.576968in}{0.561396in}}%
\pgfpathlineto{\pgfqpoint{2.578503in}{0.556332in}}%
\pgfpathlineto{\pgfqpoint{2.602314in}{0.521715in}}%
\pgfpathlineto{\pgfqpoint{2.598541in}{0.521156in}}%
\pgfpathlineto{\pgfqpoint{2.608579in}{0.496534in}}%
\pgfpathlineto{\pgfqpoint{2.636155in}{0.453729in}}%
\pgfpathlineto{\pgfqpoint{2.652165in}{0.422256in}}%
\pgfpathlineto{\pgfqpoint{2.657516in}{0.416746in}}%
\pgfpathlineto{\pgfqpoint{2.666685in}{0.396991in}}%
\pgfpathlineto{\pgfqpoint{2.669866in}{0.365359in}}%
\pgfpathlineto{\pgfqpoint{2.671132in}{0.341707in}}%
\pgfpathlineto{\pgfqpoint{2.669603in}{0.326563in}}%
\pgfpathlineto{\pgfqpoint{2.664687in}{0.315737in}}%
\pgfpathlineto{\pgfqpoint{2.666971in}{0.301540in}}%
\pgfpathlineto{\pgfqpoint{2.661678in}{0.290303in}}%
\pgfpathlineto{\pgfqpoint{2.645962in}{0.281088in}}%
\pgfpathlineto{\pgfqpoint{2.635740in}{0.281761in}}%
\pgfpathlineto{\pgfqpoint{2.629103in}{0.276944in}}%
\pgfpathlineto{\pgfqpoint{2.627201in}{0.289805in}}%
\pgfpathlineto{\pgfqpoint{2.616214in}{0.293200in}}%
\pgfpathlineto{\pgfqpoint{2.606364in}{0.311145in}}%
\pgfpathlineto{\pgfqpoint{2.605252in}{0.319317in}}%
\pgfpathlineto{\pgfqpoint{2.587624in}{0.324280in}}%
\pgfpathlineto{\pgfqpoint{2.576265in}{0.323205in}}%
\pgfpathlineto{\pgfqpoint{2.569838in}{0.335063in}}%
\pgfpathlineto{\pgfqpoint{2.562480in}{0.356375in}}%
\pgfpathlineto{\pgfqpoint{2.552268in}{0.360695in}}%
\pgfpathlineto{\pgfqpoint{2.546732in}{0.372908in}}%
\pgfpathlineto{\pgfqpoint{2.533193in}{0.380106in}}%
\pgfpathlineto{\pgfqpoint{2.525367in}{0.389110in}}%
\pgfpathlineto{\pgfqpoint{2.511867in}{0.413587in}}%
\pgfpathlineto{\pgfqpoint{2.510279in}{0.430815in}}%
\pgfpathlineto{\pgfqpoint{2.517636in}{0.441838in}}%
\pgfpathlineto{\pgfqpoint{2.516901in}{0.449497in}}%
\pgfpathlineto{\pgfqpoint{2.508399in}{0.450316in}}%
\pgfpathlineto{\pgfqpoint{2.501297in}{0.455644in}}%
\pgfpathlineto{\pgfqpoint{2.497407in}{0.449138in}}%
\pgfpathlineto{\pgfqpoint{2.504224in}{0.443837in}}%
\pgfpathlineto{\pgfqpoint{2.499262in}{0.434305in}}%
\pgfpathlineto{\pgfqpoint{2.491236in}{0.442217in}}%
\pgfpathlineto{\pgfqpoint{2.492167in}{0.464156in}}%
\pgfpathlineto{\pgfqpoint{2.496039in}{0.481956in}}%
\pgfpathlineto{\pgfqpoint{2.494048in}{0.512497in}}%
\pgfpathlineto{\pgfqpoint{2.486048in}{0.519772in}}%
\pgfpathlineto{\pgfqpoint{2.482024in}{0.529107in}}%
\pgfpathlineto{\pgfqpoint{2.468242in}{0.528899in}}%
\pgfpathlineto{\pgfqpoint{2.454618in}{0.544259in}}%
\pgfpathlineto{\pgfqpoint{2.445483in}{0.548863in}}%
\pgfpathlineto{\pgfqpoint{2.442774in}{0.558590in}}%
\pgfpathlineto{\pgfqpoint{2.433840in}{0.562006in}}%
\pgfpathlineto{\pgfqpoint{2.426435in}{0.572746in}}%
\pgfpathlineto{\pgfqpoint{2.406875in}{0.581513in}}%
\pgfpathlineto{\pgfqpoint{2.391668in}{0.581737in}}%
\pgfpathlineto{\pgfqpoint{2.385048in}{0.578370in}}%
\pgfpathlineto{\pgfqpoint{2.386693in}{0.567859in}}%
\pgfpathlineto{\pgfqpoint{2.379790in}{0.568357in}}%
\pgfpathlineto{\pgfqpoint{2.358555in}{0.553661in}}%
\pgfpathlineto{\pgfqpoint{2.333048in}{0.547836in}}%
\pgfpathlineto{\pgfqpoint{2.332629in}{0.555043in}}%
\pgfpathlineto{\pgfqpoint{2.326976in}{0.562087in}}%
\pgfpathlineto{\pgfqpoint{2.311810in}{0.571811in}}%
\pgfpathlineto{\pgfqpoint{2.290041in}{0.581781in}}%
\pgfpathlineto{\pgfqpoint{2.266395in}{0.587059in}}%
\pgfpathlineto{\pgfqpoint{2.261956in}{0.594236in}}%
\pgfpathlineto{\pgfqpoint{2.253431in}{0.588250in}}%
\pgfpathlineto{\pgfqpoint{2.243165in}{0.586930in}}%
\pgfpathlineto{\pgfqpoint{2.207429in}{0.577514in}}%
\pgfpathlineto{\pgfqpoint{2.206853in}{0.588027in}}%
\pgfusepath{stroke}%
\end{pgfscope}%
\begin{pgfscope}%
\pgfpathrectangle{\pgfqpoint{0.100000in}{0.100000in}}{\pgfqpoint{3.007045in}{1.925000in}}%
\pgfusepath{clip}%
\pgfsetbuttcap%
\pgfsetroundjoin%
\pgfsetlinewidth{0.050187pt}%
\definecolor{currentstroke}{rgb}{1.000000,1.000000,1.000000}%
\pgfsetstrokecolor{currentstroke}%
\pgfsetdash{}{0pt}%
\pgfpathmoveto{\pgfqpoint{2.605783in}{0.522207in}}%
\pgfpathlineto{\pgfqpoint{2.615029in}{0.511452in}}%
\pgfpathlineto{\pgfqpoint{2.604349in}{0.510775in}}%
\pgfpathlineto{\pgfqpoint{2.605783in}{0.522207in}}%
\pgfusepath{stroke}%
\end{pgfscope}%
\begin{pgfscope}%
\pgfpathrectangle{\pgfqpoint{0.100000in}{0.100000in}}{\pgfqpoint{3.007045in}{1.925000in}}%
\pgfusepath{clip}%
\pgfsetbuttcap%
\pgfsetroundjoin%
\pgfsetlinewidth{0.050187pt}%
\definecolor{currentstroke}{rgb}{1.000000,1.000000,1.000000}%
\pgfsetstrokecolor{currentstroke}%
\pgfsetdash{}{0pt}%
\pgfpathmoveto{\pgfqpoint{2.061492in}{1.750249in}}%
\pgfpathlineto{\pgfqpoint{2.056476in}{1.740485in}}%
\pgfpathlineto{\pgfqpoint{2.044508in}{1.734787in}}%
\pgfpathlineto{\pgfqpoint{2.039325in}{1.727117in}}%
\pgfpathlineto{\pgfqpoint{2.033253in}{1.732641in}}%
\pgfpathlineto{\pgfqpoint{2.061492in}{1.750249in}}%
\pgfusepath{stroke}%
\end{pgfscope}%
\begin{pgfscope}%
\pgfpathrectangle{\pgfqpoint{0.100000in}{0.100000in}}{\pgfqpoint{3.007045in}{1.925000in}}%
\pgfusepath{clip}%
\pgfsetbuttcap%
\pgfsetroundjoin%
\pgfsetlinewidth{0.050187pt}%
\definecolor{currentstroke}{rgb}{1.000000,1.000000,1.000000}%
\pgfsetstrokecolor{currentstroke}%
\pgfsetdash{}{0pt}%
\pgfpathmoveto{\pgfqpoint{2.065561in}{1.691489in}}%
\pgfpathlineto{\pgfqpoint{2.077807in}{1.702912in}}%
\pgfpathlineto{\pgfqpoint{2.096694in}{1.705904in}}%
\pgfpathlineto{\pgfqpoint{2.091492in}{1.697959in}}%
\pgfpathlineto{\pgfqpoint{2.078543in}{1.686470in}}%
\pgfpathlineto{\pgfqpoint{2.070973in}{1.671723in}}%
\pgfpathlineto{\pgfqpoint{2.067767in}{1.679722in}}%
\pgfpathlineto{\pgfqpoint{2.062075in}{1.680870in}}%
\pgfpathlineto{\pgfqpoint{2.061503in}{1.688091in}}%
\pgfpathlineto{\pgfqpoint{2.065561in}{1.691489in}}%
\pgfusepath{stroke}%
\end{pgfscope}%
\begin{pgfscope}%
\pgfpathrectangle{\pgfqpoint{0.100000in}{0.100000in}}{\pgfqpoint{3.007045in}{1.925000in}}%
\pgfusepath{clip}%
\pgfsetbuttcap%
\pgfsetroundjoin%
\pgfsetlinewidth{0.050187pt}%
\definecolor{currentstroke}{rgb}{1.000000,1.000000,1.000000}%
\pgfsetstrokecolor{currentstroke}%
\pgfsetdash{}{0pt}%
\pgfpathmoveto{\pgfqpoint{2.114334in}{1.552060in}}%
\pgfpathlineto{\pgfqpoint{2.111072in}{1.555679in}}%
\pgfpathlineto{\pgfqpoint{2.114501in}{1.565867in}}%
\pgfpathlineto{\pgfqpoint{2.104324in}{1.566517in}}%
\pgfpathlineto{\pgfqpoint{2.107021in}{1.575301in}}%
\pgfpathlineto{\pgfqpoint{2.105344in}{1.589291in}}%
\pgfpathlineto{\pgfqpoint{2.096212in}{1.594142in}}%
\pgfpathlineto{\pgfqpoint{2.086656in}{1.603958in}}%
\pgfpathlineto{\pgfqpoint{2.071931in}{1.606832in}}%
\pgfpathlineto{\pgfqpoint{2.057600in}{1.606750in}}%
\pgfpathlineto{\pgfqpoint{2.043506in}{1.613716in}}%
\pgfpathlineto{\pgfqpoint{1.996352in}{1.623871in}}%
\pgfpathlineto{\pgfqpoint{1.991199in}{1.634631in}}%
\pgfpathlineto{\pgfqpoint{1.982007in}{1.638302in}}%
\pgfpathlineto{\pgfqpoint{1.999369in}{1.646540in}}%
\pgfpathlineto{\pgfqpoint{2.009177in}{1.656818in}}%
\pgfpathlineto{\pgfqpoint{2.027446in}{1.659605in}}%
\pgfpathlineto{\pgfqpoint{2.038687in}{1.670074in}}%
\pgfpathlineto{\pgfqpoint{2.044585in}{1.670493in}}%
\pgfpathlineto{\pgfqpoint{2.049086in}{1.677958in}}%
\pgfpathlineto{\pgfqpoint{2.060059in}{1.686803in}}%
\pgfpathlineto{\pgfqpoint{2.061015in}{1.680552in}}%
\pgfpathlineto{\pgfqpoint{2.065959in}{1.679261in}}%
\pgfpathlineto{\pgfqpoint{2.069679in}{1.671805in}}%
\pgfpathlineto{\pgfqpoint{2.070349in}{1.659095in}}%
\pgfpathlineto{\pgfqpoint{2.081529in}{1.666691in}}%
\pgfpathlineto{\pgfqpoint{2.094557in}{1.668277in}}%
\pgfpathlineto{\pgfqpoint{2.105677in}{1.664298in}}%
\pgfpathlineto{\pgfqpoint{2.120759in}{1.643609in}}%
\pgfpathlineto{\pgfqpoint{2.137230in}{1.646915in}}%
\pgfpathlineto{\pgfqpoint{2.143923in}{1.641378in}}%
\pgfpathlineto{\pgfqpoint{2.154687in}{1.640884in}}%
\pgfpathlineto{\pgfqpoint{2.161839in}{1.650814in}}%
\pgfpathlineto{\pgfqpoint{2.175410in}{1.659592in}}%
\pgfpathlineto{\pgfqpoint{2.188452in}{1.662342in}}%
\pgfpathlineto{\pgfqpoint{2.204637in}{1.662628in}}%
\pgfpathlineto{\pgfqpoint{2.216463in}{1.669409in}}%
\pgfpathlineto{\pgfqpoint{2.226115in}{1.666285in}}%
\pgfpathlineto{\pgfqpoint{2.228171in}{1.651944in}}%
\pgfpathlineto{\pgfqpoint{2.248877in}{1.649698in}}%
\pgfpathlineto{\pgfqpoint{2.260125in}{1.656411in}}%
\pgfpathlineto{\pgfqpoint{2.267921in}{1.641288in}}%
\pgfpathlineto{\pgfqpoint{2.282798in}{1.623800in}}%
\pgfpathlineto{\pgfqpoint{2.261976in}{1.623901in}}%
\pgfpathlineto{\pgfqpoint{2.255397in}{1.621740in}}%
\pgfpathlineto{\pgfqpoint{2.246379in}{1.624549in}}%
\pgfpathlineto{\pgfqpoint{2.245772in}{1.612443in}}%
\pgfpathlineto{\pgfqpoint{2.229399in}{1.621904in}}%
\pgfpathlineto{\pgfqpoint{2.208341in}{1.624795in}}%
\pgfpathlineto{\pgfqpoint{2.202544in}{1.615581in}}%
\pgfpathlineto{\pgfqpoint{2.174992in}{1.611101in}}%
\pgfpathlineto{\pgfqpoint{2.171822in}{1.603295in}}%
\pgfpathlineto{\pgfqpoint{2.152650in}{1.600953in}}%
\pgfpathlineto{\pgfqpoint{2.146854in}{1.593019in}}%
\pgfpathlineto{\pgfqpoint{2.136713in}{1.590891in}}%
\pgfpathlineto{\pgfqpoint{2.128613in}{1.572076in}}%
\pgfpathlineto{\pgfqpoint{2.118357in}{1.553854in}}%
\pgfpathlineto{\pgfqpoint{2.114334in}{1.552060in}}%
\pgfusepath{stroke}%
\end{pgfscope}%
\begin{pgfscope}%
\pgfpathrectangle{\pgfqpoint{0.100000in}{0.100000in}}{\pgfqpoint{3.007045in}{1.925000in}}%
\pgfusepath{clip}%
\pgfsetbuttcap%
\pgfsetroundjoin%
\pgfsetlinewidth{0.050187pt}%
\definecolor{currentstroke}{rgb}{1.000000,1.000000,1.000000}%
\pgfsetstrokecolor{currentstroke}%
\pgfsetdash{}{0pt}%
\pgfpathmoveto{\pgfqpoint{2.173359in}{1.333287in}}%
\pgfpathlineto{\pgfqpoint{2.183137in}{1.343579in}}%
\pgfpathlineto{\pgfqpoint{2.187591in}{1.358500in}}%
\pgfpathlineto{\pgfqpoint{2.192889in}{1.367149in}}%
\pgfpathlineto{\pgfqpoint{2.196137in}{1.378919in}}%
\pgfpathlineto{\pgfqpoint{2.197147in}{1.402387in}}%
\pgfpathlineto{\pgfqpoint{2.192256in}{1.424878in}}%
\pgfpathlineto{\pgfqpoint{2.176094in}{1.459319in}}%
\pgfpathlineto{\pgfqpoint{2.180471in}{1.470147in}}%
\pgfpathlineto{\pgfqpoint{2.174769in}{1.485115in}}%
\pgfpathlineto{\pgfqpoint{2.184550in}{1.505816in}}%
\pgfpathlineto{\pgfqpoint{2.182956in}{1.528897in}}%
\pgfpathlineto{\pgfqpoint{2.189769in}{1.531802in}}%
\pgfpathlineto{\pgfqpoint{2.190625in}{1.542784in}}%
\pgfpathlineto{\pgfqpoint{2.202718in}{1.549833in}}%
\pgfpathlineto{\pgfqpoint{2.209558in}{1.548695in}}%
\pgfpathlineto{\pgfqpoint{2.211480in}{1.536917in}}%
\pgfpathlineto{\pgfqpoint{2.216819in}{1.536455in}}%
\pgfpathlineto{\pgfqpoint{2.221689in}{1.553444in}}%
\pgfpathlineto{\pgfqpoint{2.220034in}{1.566657in}}%
\pgfpathlineto{\pgfqpoint{2.223201in}{1.574245in}}%
\pgfpathlineto{\pgfqpoint{2.240322in}{1.582083in}}%
\pgfpathlineto{\pgfqpoint{2.232520in}{1.584872in}}%
\pgfpathlineto{\pgfqpoint{2.229979in}{1.591605in}}%
\pgfpathlineto{\pgfqpoint{2.235599in}{1.603430in}}%
\pgfpathlineto{\pgfqpoint{2.246683in}{1.607514in}}%
\pgfpathlineto{\pgfqpoint{2.259554in}{1.600537in}}%
\pgfpathlineto{\pgfqpoint{2.271689in}{1.600450in}}%
\pgfpathlineto{\pgfqpoint{2.277325in}{1.592301in}}%
\pgfpathlineto{\pgfqpoint{2.285823in}{1.592867in}}%
\pgfpathlineto{\pgfqpoint{2.312058in}{1.581538in}}%
\pgfpathlineto{\pgfqpoint{2.317300in}{1.570571in}}%
\pgfpathlineto{\pgfqpoint{2.311575in}{1.566773in}}%
\pgfpathlineto{\pgfqpoint{2.313339in}{1.558555in}}%
\pgfpathlineto{\pgfqpoint{2.318994in}{1.554899in}}%
\pgfpathlineto{\pgfqpoint{2.322146in}{1.544796in}}%
\pgfpathlineto{\pgfqpoint{2.321705in}{1.520181in}}%
\pgfpathlineto{\pgfqpoint{2.314238in}{1.514301in}}%
\pgfpathlineto{\pgfqpoint{2.312562in}{1.501372in}}%
\pgfpathlineto{\pgfqpoint{2.298712in}{1.489423in}}%
\pgfpathlineto{\pgfqpoint{2.299534in}{1.474960in}}%
\pgfpathlineto{\pgfqpoint{2.311673in}{1.469893in}}%
\pgfpathlineto{\pgfqpoint{2.325368in}{1.487939in}}%
\pgfpathlineto{\pgfqpoint{2.326500in}{1.494484in}}%
\pgfpathlineto{\pgfqpoint{2.343579in}{1.505363in}}%
\pgfpathlineto{\pgfqpoint{2.354440in}{1.500324in}}%
\pgfpathlineto{\pgfqpoint{2.361256in}{1.488936in}}%
\pgfpathlineto{\pgfqpoint{2.372251in}{1.449395in}}%
\pgfpathlineto{\pgfqpoint{2.378077in}{1.436886in}}%
\pgfpathlineto{\pgfqpoint{2.376253in}{1.431715in}}%
\pgfpathlineto{\pgfqpoint{2.376430in}{1.414106in}}%
\pgfpathlineto{\pgfqpoint{2.365832in}{1.415796in}}%
\pgfpathlineto{\pgfqpoint{2.359845in}{1.402637in}}%
\pgfpathlineto{\pgfqpoint{2.359015in}{1.393692in}}%
\pgfpathlineto{\pgfqpoint{2.351009in}{1.387949in}}%
\pgfpathlineto{\pgfqpoint{2.349230in}{1.370503in}}%
\pgfpathlineto{\pgfqpoint{2.337587in}{1.348464in}}%
\pgfpathlineto{\pgfqpoint{2.273898in}{1.338979in}}%
\pgfpathlineto{\pgfqpoint{2.273522in}{1.343148in}}%
\pgfpathlineto{\pgfqpoint{2.230918in}{1.338658in}}%
\pgfpathlineto{\pgfqpoint{2.173359in}{1.333287in}}%
\pgfusepath{stroke}%
\end{pgfscope}%
\begin{pgfscope}%
\pgfpathrectangle{\pgfqpoint{0.100000in}{0.100000in}}{\pgfqpoint{3.007045in}{1.925000in}}%
\pgfusepath{clip}%
\pgfsetbuttcap%
\pgfsetroundjoin%
\pgfsetlinewidth{0.050187pt}%
\definecolor{currentstroke}{rgb}{1.000000,1.000000,1.000000}%
\pgfsetstrokecolor{currentstroke}%
\pgfsetdash{}{0pt}%
\pgfpathmoveto{\pgfqpoint{0.000000in}{0.000000in}}%
\pgfusepath{stroke}%
\end{pgfscope}%
\begin{pgfscope}%
\pgfpathrectangle{\pgfqpoint{0.100000in}{0.100000in}}{\pgfqpoint{3.007045in}{1.925000in}}%
\pgfusepath{clip}%
\pgfsetbuttcap%
\pgfsetroundjoin%
\pgfsetlinewidth{0.050187pt}%
\definecolor{currentstroke}{rgb}{1.000000,1.000000,1.000000}%
\pgfsetstrokecolor{currentstroke}%
\pgfsetdash{}{0pt}%
\pgfusepath{stroke}%
\end{pgfscope}%
\begin{pgfscope}%
\pgfpathrectangle{\pgfqpoint{0.100000in}{0.100000in}}{\pgfqpoint{3.007045in}{1.925000in}}%
\pgfusepath{clip}%
\pgfsetbuttcap%
\pgfsetroundjoin%
\pgfsetlinewidth{0.050187pt}%
\definecolor{currentstroke}{rgb}{1.000000,1.000000,1.000000}%
\pgfsetstrokecolor{currentstroke}%
\pgfsetdash{}{0pt}%
\pgfusepath{stroke}%
\end{pgfscope}%
\begin{pgfscope}%
\pgfpathrectangle{\pgfqpoint{0.100000in}{0.100000in}}{\pgfqpoint{3.007045in}{1.925000in}}%
\pgfusepath{clip}%
\pgfsetbuttcap%
\pgfsetroundjoin%
\pgfsetlinewidth{0.050187pt}%
\definecolor{currentstroke}{rgb}{1.000000,1.000000,1.000000}%
\pgfsetstrokecolor{currentstroke}%
\pgfsetdash{}{0pt}%
\pgfusepath{stroke}%
\end{pgfscope}%
\begin{pgfscope}%
\pgfpathrectangle{\pgfqpoint{0.100000in}{0.100000in}}{\pgfqpoint{3.007045in}{1.925000in}}%
\pgfusepath{clip}%
\pgfsetbuttcap%
\pgfsetroundjoin%
\pgfsetlinewidth{0.050187pt}%
\definecolor{currentstroke}{rgb}{1.000000,1.000000,1.000000}%
\pgfsetstrokecolor{currentstroke}%
\pgfsetdash{}{0pt}%
\pgfusepath{stroke}%
\end{pgfscope}%
\begin{pgfscope}%
\pgfpathrectangle{\pgfqpoint{0.100000in}{0.100000in}}{\pgfqpoint{3.007045in}{1.925000in}}%
\pgfusepath{clip}%
\pgfsetbuttcap%
\pgfsetroundjoin%
\pgfsetlinewidth{0.050187pt}%
\definecolor{currentstroke}{rgb}{1.000000,1.000000,1.000000}%
\pgfsetstrokecolor{currentstroke}%
\pgfsetdash{}{0pt}%
\pgfusepath{stroke}%
\end{pgfscope}%
\begin{pgfscope}%
\pgfpathrectangle{\pgfqpoint{0.100000in}{0.100000in}}{\pgfqpoint{3.007045in}{1.925000in}}%
\pgfusepath{clip}%
\pgfsetbuttcap%
\pgfsetroundjoin%
\pgfsetlinewidth{0.050187pt}%
\definecolor{currentstroke}{rgb}{1.000000,1.000000,1.000000}%
\pgfsetstrokecolor{currentstroke}%
\pgfsetdash{}{0pt}%
\pgfusepath{stroke}%
\end{pgfscope}%
\begin{pgfscope}%
\pgfpathrectangle{\pgfqpoint{0.100000in}{0.100000in}}{\pgfqpoint{3.007045in}{1.925000in}}%
\pgfusepath{clip}%
\pgfsetbuttcap%
\pgfsetroundjoin%
\pgfsetlinewidth{0.050187pt}%
\definecolor{currentstroke}{rgb}{1.000000,1.000000,1.000000}%
\pgfsetstrokecolor{currentstroke}%
\pgfsetdash{}{0pt}%
\pgfusepath{stroke}%
\end{pgfscope}%
\begin{pgfscope}%
\pgfpathrectangle{\pgfqpoint{0.100000in}{0.100000in}}{\pgfqpoint{3.007045in}{1.925000in}}%
\pgfusepath{clip}%
\pgfsetbuttcap%
\pgfsetroundjoin%
\pgfsetlinewidth{0.050187pt}%
\definecolor{currentstroke}{rgb}{1.000000,1.000000,1.000000}%
\pgfsetstrokecolor{currentstroke}%
\pgfsetdash{}{0pt}%
\pgfusepath{stroke}%
\end{pgfscope}%
\begin{pgfscope}%
\pgfpathrectangle{\pgfqpoint{0.100000in}{0.100000in}}{\pgfqpoint{3.007045in}{1.925000in}}%
\pgfusepath{clip}%
\pgfsetbuttcap%
\pgfsetroundjoin%
\pgfsetlinewidth{0.050187pt}%
\definecolor{currentstroke}{rgb}{1.000000,1.000000,1.000000}%
\pgfsetstrokecolor{currentstroke}%
\pgfsetdash{}{0pt}%
\pgfusepath{stroke}%
\end{pgfscope}%
\begin{pgfscope}%
\pgfpathrectangle{\pgfqpoint{0.100000in}{0.100000in}}{\pgfqpoint{3.007045in}{1.925000in}}%
\pgfusepath{clip}%
\pgfsetbuttcap%
\pgfsetroundjoin%
\pgfsetlinewidth{0.050187pt}%
\definecolor{currentstroke}{rgb}{1.000000,1.000000,1.000000}%
\pgfsetstrokecolor{currentstroke}%
\pgfsetdash{}{0pt}%
\pgfusepath{stroke}%
\end{pgfscope}%
\begin{pgfscope}%
\pgfpathrectangle{\pgfqpoint{0.100000in}{0.100000in}}{\pgfqpoint{3.007045in}{1.925000in}}%
\pgfusepath{clip}%
\pgfsetbuttcap%
\pgfsetroundjoin%
\pgfsetlinewidth{0.050187pt}%
\definecolor{currentstroke}{rgb}{1.000000,1.000000,1.000000}%
\pgfsetstrokecolor{currentstroke}%
\pgfsetdash{}{0pt}%
\pgfusepath{stroke}%
\end{pgfscope}%
\begin{pgfscope}%
\pgfpathrectangle{\pgfqpoint{0.100000in}{0.100000in}}{\pgfqpoint{3.007045in}{1.925000in}}%
\pgfusepath{clip}%
\pgfsetbuttcap%
\pgfsetroundjoin%
\pgfsetlinewidth{0.050187pt}%
\definecolor{currentstroke}{rgb}{1.000000,1.000000,1.000000}%
\pgfsetstrokecolor{currentstroke}%
\pgfsetdash{}{0pt}%
\pgfusepath{stroke}%
\end{pgfscope}%
\begin{pgfscope}%
\pgfpathrectangle{\pgfqpoint{0.100000in}{0.100000in}}{\pgfqpoint{3.007045in}{1.925000in}}%
\pgfusepath{clip}%
\pgfsetbuttcap%
\pgfsetroundjoin%
\pgfsetlinewidth{0.050187pt}%
\definecolor{currentstroke}{rgb}{1.000000,1.000000,1.000000}%
\pgfsetstrokecolor{currentstroke}%
\pgfsetdash{}{0pt}%
\pgfusepath{stroke}%
\end{pgfscope}%
\begin{pgfscope}%
\pgfpathrectangle{\pgfqpoint{0.100000in}{0.100000in}}{\pgfqpoint{3.007045in}{1.925000in}}%
\pgfusepath{clip}%
\pgfsetbuttcap%
\pgfsetroundjoin%
\pgfsetlinewidth{0.050187pt}%
\definecolor{currentstroke}{rgb}{1.000000,1.000000,1.000000}%
\pgfsetstrokecolor{currentstroke}%
\pgfsetdash{}{0pt}%
\pgfusepath{stroke}%
\end{pgfscope}%
\begin{pgfscope}%
\pgfpathrectangle{\pgfqpoint{0.100000in}{0.100000in}}{\pgfqpoint{3.007045in}{1.925000in}}%
\pgfusepath{clip}%
\pgfsetbuttcap%
\pgfsetroundjoin%
\pgfsetlinewidth{0.050187pt}%
\definecolor{currentstroke}{rgb}{1.000000,1.000000,1.000000}%
\pgfsetstrokecolor{currentstroke}%
\pgfsetdash{}{0pt}%
\pgfusepath{stroke}%
\end{pgfscope}%
\begin{pgfscope}%
\pgfpathrectangle{\pgfqpoint{0.100000in}{0.100000in}}{\pgfqpoint{3.007045in}{1.925000in}}%
\pgfusepath{clip}%
\pgfsetbuttcap%
\pgfsetroundjoin%
\pgfsetlinewidth{0.050187pt}%
\definecolor{currentstroke}{rgb}{1.000000,1.000000,1.000000}%
\pgfsetstrokecolor{currentstroke}%
\pgfsetdash{}{0pt}%
\pgfusepath{stroke}%
\end{pgfscope}%
\begin{pgfscope}%
\pgfpathrectangle{\pgfqpoint{0.100000in}{0.100000in}}{\pgfqpoint{3.007045in}{1.925000in}}%
\pgfusepath{clip}%
\pgfsetbuttcap%
\pgfsetroundjoin%
\pgfsetlinewidth{0.050187pt}%
\definecolor{currentstroke}{rgb}{1.000000,1.000000,1.000000}%
\pgfsetstrokecolor{currentstroke}%
\pgfsetdash{}{0pt}%
\pgfusepath{stroke}%
\end{pgfscope}%
\begin{pgfscope}%
\pgfpathrectangle{\pgfqpoint{0.100000in}{0.100000in}}{\pgfqpoint{3.007045in}{1.925000in}}%
\pgfusepath{clip}%
\pgfsetbuttcap%
\pgfsetroundjoin%
\pgfsetlinewidth{0.050187pt}%
\definecolor{currentstroke}{rgb}{1.000000,1.000000,1.000000}%
\pgfsetstrokecolor{currentstroke}%
\pgfsetdash{}{0pt}%
\pgfusepath{stroke}%
\end{pgfscope}%
\begin{pgfscope}%
\pgfpathrectangle{\pgfqpoint{0.100000in}{0.100000in}}{\pgfqpoint{3.007045in}{1.925000in}}%
\pgfusepath{clip}%
\pgfsetbuttcap%
\pgfsetroundjoin%
\pgfsetlinewidth{0.050187pt}%
\definecolor{currentstroke}{rgb}{1.000000,1.000000,1.000000}%
\pgfsetstrokecolor{currentstroke}%
\pgfsetdash{}{0pt}%
\pgfusepath{stroke}%
\end{pgfscope}%
\begin{pgfscope}%
\pgfpathrectangle{\pgfqpoint{0.100000in}{0.100000in}}{\pgfqpoint{3.007045in}{1.925000in}}%
\pgfusepath{clip}%
\pgfsetbuttcap%
\pgfsetroundjoin%
\pgfsetlinewidth{0.050187pt}%
\definecolor{currentstroke}{rgb}{1.000000,1.000000,1.000000}%
\pgfsetstrokecolor{currentstroke}%
\pgfsetdash{}{0pt}%
\pgfusepath{stroke}%
\end{pgfscope}%
\begin{pgfscope}%
\pgfpathrectangle{\pgfqpoint{0.100000in}{0.100000in}}{\pgfqpoint{3.007045in}{1.925000in}}%
\pgfusepath{clip}%
\pgfsetbuttcap%
\pgfsetroundjoin%
\pgfsetlinewidth{0.050187pt}%
\definecolor{currentstroke}{rgb}{1.000000,1.000000,1.000000}%
\pgfsetstrokecolor{currentstroke}%
\pgfsetdash{}{0pt}%
\pgfusepath{stroke}%
\end{pgfscope}%
\begin{pgfscope}%
\pgfpathrectangle{\pgfqpoint{0.100000in}{0.100000in}}{\pgfqpoint{3.007045in}{1.925000in}}%
\pgfusepath{clip}%
\pgfsetbuttcap%
\pgfsetroundjoin%
\pgfsetlinewidth{0.050187pt}%
\definecolor{currentstroke}{rgb}{1.000000,1.000000,1.000000}%
\pgfsetstrokecolor{currentstroke}%
\pgfsetdash{}{0pt}%
\pgfusepath{stroke}%
\end{pgfscope}%
\begin{pgfscope}%
\pgfpathrectangle{\pgfqpoint{0.100000in}{0.100000in}}{\pgfqpoint{3.007045in}{1.925000in}}%
\pgfusepath{clip}%
\pgfsetbuttcap%
\pgfsetroundjoin%
\pgfsetlinewidth{0.050187pt}%
\definecolor{currentstroke}{rgb}{1.000000,1.000000,1.000000}%
\pgfsetstrokecolor{currentstroke}%
\pgfsetdash{}{0pt}%
\pgfusepath{stroke}%
\end{pgfscope}%
\begin{pgfscope}%
\pgfpathrectangle{\pgfqpoint{0.100000in}{0.100000in}}{\pgfqpoint{3.007045in}{1.925000in}}%
\pgfusepath{clip}%
\pgfsetbuttcap%
\pgfsetroundjoin%
\pgfsetlinewidth{0.050187pt}%
\definecolor{currentstroke}{rgb}{1.000000,1.000000,1.000000}%
\pgfsetstrokecolor{currentstroke}%
\pgfsetdash{}{0pt}%
\pgfusepath{stroke}%
\end{pgfscope}%
\begin{pgfscope}%
\pgfpathrectangle{\pgfqpoint{0.100000in}{0.100000in}}{\pgfqpoint{3.007045in}{1.925000in}}%
\pgfusepath{clip}%
\pgfsetbuttcap%
\pgfsetroundjoin%
\pgfsetlinewidth{0.050187pt}%
\definecolor{currentstroke}{rgb}{1.000000,1.000000,1.000000}%
\pgfsetstrokecolor{currentstroke}%
\pgfsetdash{}{0pt}%
\pgfusepath{stroke}%
\end{pgfscope}%
\begin{pgfscope}%
\pgfpathrectangle{\pgfqpoint{0.100000in}{0.100000in}}{\pgfqpoint{3.007045in}{1.925000in}}%
\pgfusepath{clip}%
\pgfsetbuttcap%
\pgfsetroundjoin%
\pgfsetlinewidth{0.050187pt}%
\definecolor{currentstroke}{rgb}{1.000000,1.000000,1.000000}%
\pgfsetstrokecolor{currentstroke}%
\pgfsetdash{}{0pt}%
\pgfusepath{stroke}%
\end{pgfscope}%
\begin{pgfscope}%
\pgfpathrectangle{\pgfqpoint{0.100000in}{0.100000in}}{\pgfqpoint{3.007045in}{1.925000in}}%
\pgfusepath{clip}%
\pgfsetbuttcap%
\pgfsetroundjoin%
\pgfsetlinewidth{0.050187pt}%
\definecolor{currentstroke}{rgb}{1.000000,1.000000,1.000000}%
\pgfsetstrokecolor{currentstroke}%
\pgfsetdash{}{0pt}%
\pgfusepath{stroke}%
\end{pgfscope}%
\begin{pgfscope}%
\pgfpathrectangle{\pgfqpoint{2.857045in}{0.350000in}}{\pgfqpoint{0.300705in}{0.962500in}}%
\pgfusepath{clip}%
\pgfsetbuttcap%
\pgfsetmiterjoin%
\definecolor{currentfill}{rgb}{0.790496,0.868174,0.941930}%
\pgfsetfillcolor{currentfill}%
\pgfsetlinewidth{0.000000pt}%
\definecolor{currentstroke}{rgb}{0.000000,0.000000,0.000000}%
\pgfsetstrokecolor{currentstroke}%
\pgfsetstrokeopacity{0.000000}%
\pgfsetdash{}{0pt}%
\pgfpathmoveto{\pgfqpoint{3.107615in}{0.350000in}}%
\pgfpathlineto{\pgfqpoint{3.157750in}{0.350000in}}%
\pgfpathlineto{\pgfqpoint{3.157750in}{0.446250in}}%
\pgfpathlineto{\pgfqpoint{3.107615in}{0.446250in}}%
\pgfpathlineto{\pgfqpoint{3.107615in}{0.350000in}}%
\pgfpathclose%
\pgfusepath{fill}%
\end{pgfscope}%
\begin{pgfscope}%
\pgfpathrectangle{\pgfqpoint{2.857045in}{0.350000in}}{\pgfqpoint{0.300705in}{0.962500in}}%
\pgfusepath{clip}%
\pgfsetbuttcap%
\pgfsetmiterjoin%
\definecolor{currentfill}{rgb}{0.760477,0.852026,0.931657}%
\pgfsetfillcolor{currentfill}%
\pgfsetlinewidth{0.000000pt}%
\definecolor{currentstroke}{rgb}{0.000000,0.000000,0.000000}%
\pgfsetstrokecolor{currentstroke}%
\pgfsetstrokeopacity{0.000000}%
\pgfsetdash{}{0pt}%
\pgfpathmoveto{\pgfqpoint{3.105780in}{0.359625in}}%
\pgfpathlineto{\pgfqpoint{3.157750in}{0.359625in}}%
\pgfpathlineto{\pgfqpoint{3.157750in}{0.455875in}}%
\pgfpathlineto{\pgfqpoint{3.105780in}{0.455875in}}%
\pgfpathlineto{\pgfqpoint{3.105780in}{0.359625in}}%
\pgfpathclose%
\pgfusepath{fill}%
\end{pgfscope}%
\begin{pgfscope}%
\pgfpathrectangle{\pgfqpoint{2.857045in}{0.350000in}}{\pgfqpoint{0.300705in}{0.962500in}}%
\pgfusepath{clip}%
\pgfsetbuttcap%
\pgfsetmiterjoin%
\definecolor{currentfill}{rgb}{0.721107,0.835294,0.917878}%
\pgfsetfillcolor{currentfill}%
\pgfsetlinewidth{0.000000pt}%
\definecolor{currentstroke}{rgb}{0.000000,0.000000,0.000000}%
\pgfsetstrokecolor{currentstroke}%
\pgfsetstrokeopacity{0.000000}%
\pgfsetdash{}{0pt}%
\pgfpathmoveto{\pgfqpoint{3.103965in}{0.369250in}}%
\pgfpathlineto{\pgfqpoint{3.157750in}{0.369250in}}%
\pgfpathlineto{\pgfqpoint{3.157750in}{0.465500in}}%
\pgfpathlineto{\pgfqpoint{3.103965in}{0.465500in}}%
\pgfpathlineto{\pgfqpoint{3.103965in}{0.369250in}}%
\pgfpathclose%
\pgfusepath{fill}%
\end{pgfscope}%
\begin{pgfscope}%
\pgfpathrectangle{\pgfqpoint{2.857045in}{0.350000in}}{\pgfqpoint{0.300705in}{0.962500in}}%
\pgfusepath{clip}%
\pgfsetbuttcap%
\pgfsetmiterjoin%
\definecolor{currentfill}{rgb}{0.691580,0.822745,0.907543}%
\pgfsetfillcolor{currentfill}%
\pgfsetlinewidth{0.000000pt}%
\definecolor{currentstroke}{rgb}{0.000000,0.000000,0.000000}%
\pgfsetstrokecolor{currentstroke}%
\pgfsetstrokeopacity{0.000000}%
\pgfsetdash{}{0pt}%
\pgfpathmoveto{\pgfqpoint{3.102756in}{0.378875in}}%
\pgfpathlineto{\pgfqpoint{3.157750in}{0.378875in}}%
\pgfpathlineto{\pgfqpoint{3.157750in}{0.475125in}}%
\pgfpathlineto{\pgfqpoint{3.102756in}{0.475125in}}%
\pgfpathlineto{\pgfqpoint{3.102756in}{0.378875in}}%
\pgfpathclose%
\pgfusepath{fill}%
\end{pgfscope}%
\begin{pgfscope}%
\pgfpathrectangle{\pgfqpoint{2.857045in}{0.350000in}}{\pgfqpoint{0.300705in}{0.962500in}}%
\pgfusepath{clip}%
\pgfsetbuttcap%
\pgfsetmiterjoin%
\definecolor{currentfill}{rgb}{0.676817,0.816471,0.902376}%
\pgfsetfillcolor{currentfill}%
\pgfsetlinewidth{0.000000pt}%
\definecolor{currentstroke}{rgb}{0.000000,0.000000,0.000000}%
\pgfsetstrokecolor{currentstroke}%
\pgfsetstrokeopacity{0.000000}%
\pgfsetdash{}{0pt}%
\pgfpathmoveto{\pgfqpoint{3.102170in}{0.388500in}}%
\pgfpathlineto{\pgfqpoint{3.157750in}{0.388500in}}%
\pgfpathlineto{\pgfqpoint{3.157750in}{0.484750in}}%
\pgfpathlineto{\pgfqpoint{3.102170in}{0.484750in}}%
\pgfpathlineto{\pgfqpoint{3.102170in}{0.388500in}}%
\pgfpathclose%
\pgfusepath{fill}%
\end{pgfscope}%
\begin{pgfscope}%
\pgfpathrectangle{\pgfqpoint{2.857045in}{0.350000in}}{\pgfqpoint{0.300705in}{0.962500in}}%
\pgfusepath{clip}%
\pgfsetbuttcap%
\pgfsetmiterjoin%
\definecolor{currentfill}{rgb}{0.671895,0.814379,0.900654}%
\pgfsetfillcolor{currentfill}%
\pgfsetlinewidth{0.000000pt}%
\definecolor{currentstroke}{rgb}{0.000000,0.000000,0.000000}%
\pgfsetstrokecolor{currentstroke}%
\pgfsetstrokeopacity{0.000000}%
\pgfsetdash{}{0pt}%
\pgfpathmoveto{\pgfqpoint{3.101868in}{0.398125in}}%
\pgfpathlineto{\pgfqpoint{3.157750in}{0.398125in}}%
\pgfpathlineto{\pgfqpoint{3.157750in}{0.494375in}}%
\pgfpathlineto{\pgfqpoint{3.101868in}{0.494375in}}%
\pgfpathlineto{\pgfqpoint{3.101868in}{0.398125in}}%
\pgfpathclose%
\pgfusepath{fill}%
\end{pgfscope}%
\begin{pgfscope}%
\pgfpathrectangle{\pgfqpoint{2.857045in}{0.350000in}}{\pgfqpoint{0.300705in}{0.962500in}}%
\pgfusepath{clip}%
\pgfsetbuttcap%
\pgfsetmiterjoin%
\definecolor{currentfill}{rgb}{0.657132,0.808105,0.895486}%
\pgfsetfillcolor{currentfill}%
\pgfsetlinewidth{0.000000pt}%
\definecolor{currentstroke}{rgb}{0.000000,0.000000,0.000000}%
\pgfsetstrokecolor{currentstroke}%
\pgfsetstrokeopacity{0.000000}%
\pgfsetdash{}{0pt}%
\pgfpathmoveto{\pgfqpoint{3.101144in}{0.407750in}}%
\pgfpathlineto{\pgfqpoint{3.157750in}{0.407750in}}%
\pgfpathlineto{\pgfqpoint{3.157750in}{0.504000in}}%
\pgfpathlineto{\pgfqpoint{3.101144in}{0.504000in}}%
\pgfpathlineto{\pgfqpoint{3.101144in}{0.407750in}}%
\pgfpathclose%
\pgfusepath{fill}%
\end{pgfscope}%
\begin{pgfscope}%
\pgfpathrectangle{\pgfqpoint{2.857045in}{0.350000in}}{\pgfqpoint{0.300705in}{0.962500in}}%
\pgfusepath{clip}%
\pgfsetbuttcap%
\pgfsetmiterjoin%
\definecolor{currentfill}{rgb}{0.637447,0.799739,0.888597}%
\pgfsetfillcolor{currentfill}%
\pgfsetlinewidth{0.000000pt}%
\definecolor{currentstroke}{rgb}{0.000000,0.000000,0.000000}%
\pgfsetstrokecolor{currentstroke}%
\pgfsetstrokeopacity{0.000000}%
\pgfsetdash{}{0pt}%
\pgfpathmoveto{\pgfqpoint{3.100283in}{0.417375in}}%
\pgfpathlineto{\pgfqpoint{3.157750in}{0.417375in}}%
\pgfpathlineto{\pgfqpoint{3.157750in}{0.513625in}}%
\pgfpathlineto{\pgfqpoint{3.100283in}{0.513625in}}%
\pgfpathlineto{\pgfqpoint{3.100283in}{0.417375in}}%
\pgfpathclose%
\pgfusepath{fill}%
\end{pgfscope}%
\begin{pgfscope}%
\pgfpathrectangle{\pgfqpoint{2.857045in}{0.350000in}}{\pgfqpoint{0.300705in}{0.962500in}}%
\pgfusepath{clip}%
\pgfsetbuttcap%
\pgfsetmiterjoin%
\definecolor{currentfill}{rgb}{0.632526,0.797647,0.886874}%
\pgfsetfillcolor{currentfill}%
\pgfsetlinewidth{0.000000pt}%
\definecolor{currentstroke}{rgb}{0.000000,0.000000,0.000000}%
\pgfsetstrokecolor{currentstroke}%
\pgfsetstrokeopacity{0.000000}%
\pgfsetdash{}{0pt}%
\pgfpathmoveto{\pgfqpoint{3.100079in}{0.427000in}}%
\pgfpathlineto{\pgfqpoint{3.157750in}{0.427000in}}%
\pgfpathlineto{\pgfqpoint{3.157750in}{0.523250in}}%
\pgfpathlineto{\pgfqpoint{3.100079in}{0.523250in}}%
\pgfpathlineto{\pgfqpoint{3.100079in}{0.427000in}}%
\pgfpathclose%
\pgfusepath{fill}%
\end{pgfscope}%
\begin{pgfscope}%
\pgfpathrectangle{\pgfqpoint{2.857045in}{0.350000in}}{\pgfqpoint{0.300705in}{0.962500in}}%
\pgfusepath{clip}%
\pgfsetbuttcap%
\pgfsetmiterjoin%
\definecolor{currentfill}{rgb}{0.632526,0.797647,0.886874}%
\pgfsetfillcolor{currentfill}%
\pgfsetlinewidth{0.000000pt}%
\definecolor{currentstroke}{rgb}{0.000000,0.000000,0.000000}%
\pgfsetstrokecolor{currentstroke}%
\pgfsetstrokeopacity{0.000000}%
\pgfsetdash{}{0pt}%
\pgfpathmoveto{\pgfqpoint{3.100079in}{0.436625in}}%
\pgfpathlineto{\pgfqpoint{3.157750in}{0.436625in}}%
\pgfpathlineto{\pgfqpoint{3.157750in}{0.532875in}}%
\pgfpathlineto{\pgfqpoint{3.100079in}{0.532875in}}%
\pgfpathlineto{\pgfqpoint{3.100079in}{0.436625in}}%
\pgfpathclose%
\pgfusepath{fill}%
\end{pgfscope}%
\begin{pgfscope}%
\pgfpathrectangle{\pgfqpoint{2.857045in}{0.350000in}}{\pgfqpoint{0.300705in}{0.962500in}}%
\pgfusepath{clip}%
\pgfsetbuttcap%
\pgfsetmiterjoin%
\definecolor{currentfill}{rgb}{0.632526,0.797647,0.886874}%
\pgfsetfillcolor{currentfill}%
\pgfsetlinewidth{0.000000pt}%
\definecolor{currentstroke}{rgb}{0.000000,0.000000,0.000000}%
\pgfsetstrokecolor{currentstroke}%
\pgfsetstrokeopacity{0.000000}%
\pgfsetdash{}{0pt}%
\pgfpathmoveto{\pgfqpoint{3.099996in}{0.446250in}}%
\pgfpathlineto{\pgfqpoint{3.157750in}{0.446250in}}%
\pgfpathlineto{\pgfqpoint{3.157750in}{0.542500in}}%
\pgfpathlineto{\pgfqpoint{3.099996in}{0.542500in}}%
\pgfpathlineto{\pgfqpoint{3.099996in}{0.446250in}}%
\pgfpathclose%
\pgfusepath{fill}%
\end{pgfscope}%
\begin{pgfscope}%
\pgfpathrectangle{\pgfqpoint{2.857045in}{0.350000in}}{\pgfqpoint{0.300705in}{0.962500in}}%
\pgfusepath{clip}%
\pgfsetbuttcap%
\pgfsetmiterjoin%
\definecolor{currentfill}{rgb}{0.632526,0.797647,0.886874}%
\pgfsetfillcolor{currentfill}%
\pgfsetlinewidth{0.000000pt}%
\definecolor{currentstroke}{rgb}{0.000000,0.000000,0.000000}%
\pgfsetstrokecolor{currentstroke}%
\pgfsetstrokeopacity{0.000000}%
\pgfsetdash{}{0pt}%
\pgfpathmoveto{\pgfqpoint{3.099996in}{0.455875in}}%
\pgfpathlineto{\pgfqpoint{3.157750in}{0.455875in}}%
\pgfpathlineto{\pgfqpoint{3.157750in}{0.552125in}}%
\pgfpathlineto{\pgfqpoint{3.099996in}{0.552125in}}%
\pgfpathlineto{\pgfqpoint{3.099996in}{0.455875in}}%
\pgfpathclose%
\pgfusepath{fill}%
\end{pgfscope}%
\begin{pgfscope}%
\pgfpathrectangle{\pgfqpoint{2.857045in}{0.350000in}}{\pgfqpoint{0.300705in}{0.962500in}}%
\pgfusepath{clip}%
\pgfsetbuttcap%
\pgfsetmiterjoin%
\definecolor{currentfill}{rgb}{0.627605,0.795556,0.885152}%
\pgfsetfillcolor{currentfill}%
\pgfsetlinewidth{0.000000pt}%
\definecolor{currentstroke}{rgb}{0.000000,0.000000,0.000000}%
\pgfsetstrokecolor{currentstroke}%
\pgfsetstrokeopacity{0.000000}%
\pgfsetdash{}{0pt}%
\pgfpathmoveto{\pgfqpoint{3.099947in}{0.465500in}}%
\pgfpathlineto{\pgfqpoint{3.157750in}{0.465500in}}%
\pgfpathlineto{\pgfqpoint{3.157750in}{0.561750in}}%
\pgfpathlineto{\pgfqpoint{3.099947in}{0.561750in}}%
\pgfpathlineto{\pgfqpoint{3.099947in}{0.465500in}}%
\pgfpathclose%
\pgfusepath{fill}%
\end{pgfscope}%
\begin{pgfscope}%
\pgfpathrectangle{\pgfqpoint{2.857045in}{0.350000in}}{\pgfqpoint{0.300705in}{0.962500in}}%
\pgfusepath{clip}%
\pgfsetbuttcap%
\pgfsetmiterjoin%
\definecolor{currentfill}{rgb}{0.592157,0.777086,0.876432}%
\pgfsetfillcolor{currentfill}%
\pgfsetlinewidth{0.000000pt}%
\definecolor{currentstroke}{rgb}{0.000000,0.000000,0.000000}%
\pgfsetstrokecolor{currentstroke}%
\pgfsetstrokeopacity{0.000000}%
\pgfsetdash{}{0pt}%
\pgfpathmoveto{\pgfqpoint{3.098467in}{0.475125in}}%
\pgfpathlineto{\pgfqpoint{3.157750in}{0.475125in}}%
\pgfpathlineto{\pgfqpoint{3.157750in}{0.571375in}}%
\pgfpathlineto{\pgfqpoint{3.098467in}{0.571375in}}%
\pgfpathlineto{\pgfqpoint{3.098467in}{0.475125in}}%
\pgfpathclose%
\pgfusepath{fill}%
\end{pgfscope}%
\begin{pgfscope}%
\pgfpathrectangle{\pgfqpoint{2.857045in}{0.350000in}}{\pgfqpoint{0.300705in}{0.962500in}}%
\pgfusepath{clip}%
\pgfsetbuttcap%
\pgfsetmiterjoin%
\definecolor{currentfill}{rgb}{0.585882,0.773641,0.875079}%
\pgfsetfillcolor{currentfill}%
\pgfsetlinewidth{0.000000pt}%
\definecolor{currentstroke}{rgb}{0.000000,0.000000,0.000000}%
\pgfsetstrokecolor{currentstroke}%
\pgfsetstrokeopacity{0.000000}%
\pgfsetdash{}{0pt}%
\pgfpathmoveto{\pgfqpoint{3.098376in}{0.484750in}}%
\pgfpathlineto{\pgfqpoint{3.157750in}{0.484750in}}%
\pgfpathlineto{\pgfqpoint{3.157750in}{0.581000in}}%
\pgfpathlineto{\pgfqpoint{3.098376in}{0.581000in}}%
\pgfpathlineto{\pgfqpoint{3.098376in}{0.484750in}}%
\pgfpathclose%
\pgfusepath{fill}%
\end{pgfscope}%
\begin{pgfscope}%
\pgfpathrectangle{\pgfqpoint{2.857045in}{0.350000in}}{\pgfqpoint{0.300705in}{0.962500in}}%
\pgfusepath{clip}%
\pgfsetbuttcap%
\pgfsetmiterjoin%
\definecolor{currentfill}{rgb}{0.579608,0.770196,0.873725}%
\pgfsetfillcolor{currentfill}%
\pgfsetlinewidth{0.000000pt}%
\definecolor{currentstroke}{rgb}{0.000000,0.000000,0.000000}%
\pgfsetstrokecolor{currentstroke}%
\pgfsetstrokeopacity{0.000000}%
\pgfsetdash{}{0pt}%
\pgfpathmoveto{\pgfqpoint{3.098186in}{0.494375in}}%
\pgfpathlineto{\pgfqpoint{3.157750in}{0.494375in}}%
\pgfpathlineto{\pgfqpoint{3.157750in}{0.590625in}}%
\pgfpathlineto{\pgfqpoint{3.098186in}{0.590625in}}%
\pgfpathlineto{\pgfqpoint{3.098186in}{0.494375in}}%
\pgfpathclose%
\pgfusepath{fill}%
\end{pgfscope}%
\begin{pgfscope}%
\pgfpathrectangle{\pgfqpoint{2.857045in}{0.350000in}}{\pgfqpoint{0.300705in}{0.962500in}}%
\pgfusepath{clip}%
\pgfsetbuttcap%
\pgfsetmiterjoin%
\definecolor{currentfill}{rgb}{0.567059,0.763306,0.871019}%
\pgfsetfillcolor{currentfill}%
\pgfsetlinewidth{0.000000pt}%
\definecolor{currentstroke}{rgb}{0.000000,0.000000,0.000000}%
\pgfsetstrokecolor{currentstroke}%
\pgfsetstrokeopacity{0.000000}%
\pgfsetdash{}{0pt}%
\pgfpathmoveto{\pgfqpoint{3.097756in}{0.504000in}}%
\pgfpathlineto{\pgfqpoint{3.157750in}{0.504000in}}%
\pgfpathlineto{\pgfqpoint{3.157750in}{0.600250in}}%
\pgfpathlineto{\pgfqpoint{3.097756in}{0.600250in}}%
\pgfpathlineto{\pgfqpoint{3.097756in}{0.504000in}}%
\pgfpathclose%
\pgfusepath{fill}%
\end{pgfscope}%
\begin{pgfscope}%
\pgfpathrectangle{\pgfqpoint{2.857045in}{0.350000in}}{\pgfqpoint{0.300705in}{0.962500in}}%
\pgfusepath{clip}%
\pgfsetbuttcap%
\pgfsetmiterjoin%
\definecolor{currentfill}{rgb}{0.567059,0.763306,0.871019}%
\pgfsetfillcolor{currentfill}%
\pgfsetlinewidth{0.000000pt}%
\definecolor{currentstroke}{rgb}{0.000000,0.000000,0.000000}%
\pgfsetstrokecolor{currentstroke}%
\pgfsetstrokeopacity{0.000000}%
\pgfsetdash{}{0pt}%
\pgfpathmoveto{\pgfqpoint{3.097609in}{0.513625in}}%
\pgfpathlineto{\pgfqpoint{3.157750in}{0.513625in}}%
\pgfpathlineto{\pgfqpoint{3.157750in}{0.609875in}}%
\pgfpathlineto{\pgfqpoint{3.097609in}{0.609875in}}%
\pgfpathlineto{\pgfqpoint{3.097609in}{0.513625in}}%
\pgfpathclose%
\pgfusepath{fill}%
\end{pgfscope}%
\begin{pgfscope}%
\pgfpathrectangle{\pgfqpoint{2.857045in}{0.350000in}}{\pgfqpoint{0.300705in}{0.962500in}}%
\pgfusepath{clip}%
\pgfsetbuttcap%
\pgfsetmiterjoin%
\definecolor{currentfill}{rgb}{0.554510,0.756417,0.868312}%
\pgfsetfillcolor{currentfill}%
\pgfsetlinewidth{0.000000pt}%
\definecolor{currentstroke}{rgb}{0.000000,0.000000,0.000000}%
\pgfsetstrokecolor{currentstroke}%
\pgfsetstrokeopacity{0.000000}%
\pgfsetdash{}{0pt}%
\pgfpathmoveto{\pgfqpoint{3.097293in}{0.523250in}}%
\pgfpathlineto{\pgfqpoint{3.157750in}{0.523250in}}%
\pgfpathlineto{\pgfqpoint{3.157750in}{0.619500in}}%
\pgfpathlineto{\pgfqpoint{3.097293in}{0.619500in}}%
\pgfpathlineto{\pgfqpoint{3.097293in}{0.523250in}}%
\pgfpathclose%
\pgfusepath{fill}%
\end{pgfscope}%
\begin{pgfscope}%
\pgfpathrectangle{\pgfqpoint{2.857045in}{0.350000in}}{\pgfqpoint{0.300705in}{0.962500in}}%
\pgfusepath{clip}%
\pgfsetbuttcap%
\pgfsetmiterjoin%
\definecolor{currentfill}{rgb}{0.554510,0.756417,0.868312}%
\pgfsetfillcolor{currentfill}%
\pgfsetlinewidth{0.000000pt}%
\definecolor{currentstroke}{rgb}{0.000000,0.000000,0.000000}%
\pgfsetstrokecolor{currentstroke}%
\pgfsetstrokeopacity{0.000000}%
\pgfsetdash{}{0pt}%
\pgfpathmoveto{\pgfqpoint{3.097268in}{0.532875in}}%
\pgfpathlineto{\pgfqpoint{3.157750in}{0.532875in}}%
\pgfpathlineto{\pgfqpoint{3.157750in}{0.629125in}}%
\pgfpathlineto{\pgfqpoint{3.097268in}{0.629125in}}%
\pgfpathlineto{\pgfqpoint{3.097268in}{0.532875in}}%
\pgfpathclose%
\pgfusepath{fill}%
\end{pgfscope}%
\begin{pgfscope}%
\pgfpathrectangle{\pgfqpoint{2.857045in}{0.350000in}}{\pgfqpoint{0.300705in}{0.962500in}}%
\pgfusepath{clip}%
\pgfsetbuttcap%
\pgfsetmiterjoin%
\definecolor{currentfill}{rgb}{0.548235,0.752972,0.866959}%
\pgfsetfillcolor{currentfill}%
\pgfsetlinewidth{0.000000pt}%
\definecolor{currentstroke}{rgb}{0.000000,0.000000,0.000000}%
\pgfsetstrokecolor{currentstroke}%
\pgfsetstrokeopacity{0.000000}%
\pgfsetdash{}{0pt}%
\pgfpathmoveto{\pgfqpoint{3.097104in}{0.542500in}}%
\pgfpathlineto{\pgfqpoint{3.157750in}{0.542500in}}%
\pgfpathlineto{\pgfqpoint{3.157750in}{0.638750in}}%
\pgfpathlineto{\pgfqpoint{3.097104in}{0.638750in}}%
\pgfpathlineto{\pgfqpoint{3.097104in}{0.542500in}}%
\pgfpathclose%
\pgfusepath{fill}%
\end{pgfscope}%
\begin{pgfscope}%
\pgfpathrectangle{\pgfqpoint{2.857045in}{0.350000in}}{\pgfqpoint{0.300705in}{0.962500in}}%
\pgfusepath{clip}%
\pgfsetbuttcap%
\pgfsetmiterjoin%
\definecolor{currentfill}{rgb}{0.541961,0.749527,0.865606}%
\pgfsetfillcolor{currentfill}%
\pgfsetlinewidth{0.000000pt}%
\definecolor{currentstroke}{rgb}{0.000000,0.000000,0.000000}%
\pgfsetstrokecolor{currentstroke}%
\pgfsetstrokeopacity{0.000000}%
\pgfsetdash{}{0pt}%
\pgfpathmoveto{\pgfqpoint{3.096865in}{0.552125in}}%
\pgfpathlineto{\pgfqpoint{3.157750in}{0.552125in}}%
\pgfpathlineto{\pgfqpoint{3.157750in}{0.648375in}}%
\pgfpathlineto{\pgfqpoint{3.096865in}{0.648375in}}%
\pgfpathlineto{\pgfqpoint{3.096865in}{0.552125in}}%
\pgfpathclose%
\pgfusepath{fill}%
\end{pgfscope}%
\begin{pgfscope}%
\pgfpathrectangle{\pgfqpoint{2.857045in}{0.350000in}}{\pgfqpoint{0.300705in}{0.962500in}}%
\pgfusepath{clip}%
\pgfsetbuttcap%
\pgfsetmiterjoin%
\definecolor{currentfill}{rgb}{0.535686,0.746082,0.864252}%
\pgfsetfillcolor{currentfill}%
\pgfsetlinewidth{0.000000pt}%
\definecolor{currentstroke}{rgb}{0.000000,0.000000,0.000000}%
\pgfsetstrokecolor{currentstroke}%
\pgfsetstrokeopacity{0.000000}%
\pgfsetdash{}{0pt}%
\pgfpathmoveto{\pgfqpoint{3.096667in}{0.561750in}}%
\pgfpathlineto{\pgfqpoint{3.157750in}{0.561750in}}%
\pgfpathlineto{\pgfqpoint{3.157750in}{0.658000in}}%
\pgfpathlineto{\pgfqpoint{3.096667in}{0.658000in}}%
\pgfpathlineto{\pgfqpoint{3.096667in}{0.561750in}}%
\pgfpathclose%
\pgfusepath{fill}%
\end{pgfscope}%
\begin{pgfscope}%
\pgfpathrectangle{\pgfqpoint{2.857045in}{0.350000in}}{\pgfqpoint{0.300705in}{0.962500in}}%
\pgfusepath{clip}%
\pgfsetbuttcap%
\pgfsetmiterjoin%
\definecolor{currentfill}{rgb}{0.535686,0.746082,0.864252}%
\pgfsetfillcolor{currentfill}%
\pgfsetlinewidth{0.000000pt}%
\definecolor{currentstroke}{rgb}{0.000000,0.000000,0.000000}%
\pgfsetstrokecolor{currentstroke}%
\pgfsetstrokeopacity{0.000000}%
\pgfsetdash{}{0pt}%
\pgfpathmoveto{\pgfqpoint{3.096667in}{0.571375in}}%
\pgfpathlineto{\pgfqpoint{3.157750in}{0.571375in}}%
\pgfpathlineto{\pgfqpoint{3.157750in}{0.667625in}}%
\pgfpathlineto{\pgfqpoint{3.096667in}{0.667625in}}%
\pgfpathlineto{\pgfqpoint{3.096667in}{0.571375in}}%
\pgfpathclose%
\pgfusepath{fill}%
\end{pgfscope}%
\begin{pgfscope}%
\pgfpathrectangle{\pgfqpoint{2.857045in}{0.350000in}}{\pgfqpoint{0.300705in}{0.962500in}}%
\pgfusepath{clip}%
\pgfsetbuttcap%
\pgfsetmiterjoin%
\definecolor{currentfill}{rgb}{0.535686,0.746082,0.864252}%
\pgfsetfillcolor{currentfill}%
\pgfsetlinewidth{0.000000pt}%
\definecolor{currentstroke}{rgb}{0.000000,0.000000,0.000000}%
\pgfsetstrokecolor{currentstroke}%
\pgfsetstrokeopacity{0.000000}%
\pgfsetdash{}{0pt}%
\pgfpathmoveto{\pgfqpoint{3.096667in}{0.581000in}}%
\pgfpathlineto{\pgfqpoint{3.157750in}{0.581000in}}%
\pgfpathlineto{\pgfqpoint{3.157750in}{0.677250in}}%
\pgfpathlineto{\pgfqpoint{3.096667in}{0.677250in}}%
\pgfpathlineto{\pgfqpoint{3.096667in}{0.581000in}}%
\pgfpathclose%
\pgfusepath{fill}%
\end{pgfscope}%
\begin{pgfscope}%
\pgfpathrectangle{\pgfqpoint{2.857045in}{0.350000in}}{\pgfqpoint{0.300705in}{0.962500in}}%
\pgfusepath{clip}%
\pgfsetbuttcap%
\pgfsetmiterjoin%
\definecolor{currentfill}{rgb}{0.535686,0.746082,0.864252}%
\pgfsetfillcolor{currentfill}%
\pgfsetlinewidth{0.000000pt}%
\definecolor{currentstroke}{rgb}{0.000000,0.000000,0.000000}%
\pgfsetstrokecolor{currentstroke}%
\pgfsetstrokeopacity{0.000000}%
\pgfsetdash{}{0pt}%
\pgfpathmoveto{\pgfqpoint{3.096667in}{0.590625in}}%
\pgfpathlineto{\pgfqpoint{3.157750in}{0.590625in}}%
\pgfpathlineto{\pgfqpoint{3.157750in}{0.686875in}}%
\pgfpathlineto{\pgfqpoint{3.096667in}{0.686875in}}%
\pgfpathlineto{\pgfqpoint{3.096667in}{0.590625in}}%
\pgfpathclose%
\pgfusepath{fill}%
\end{pgfscope}%
\begin{pgfscope}%
\pgfpathrectangle{\pgfqpoint{2.857045in}{0.350000in}}{\pgfqpoint{0.300705in}{0.962500in}}%
\pgfusepath{clip}%
\pgfsetbuttcap%
\pgfsetmiterjoin%
\definecolor{currentfill}{rgb}{0.535686,0.746082,0.864252}%
\pgfsetfillcolor{currentfill}%
\pgfsetlinewidth{0.000000pt}%
\definecolor{currentstroke}{rgb}{0.000000,0.000000,0.000000}%
\pgfsetstrokecolor{currentstroke}%
\pgfsetstrokeopacity{0.000000}%
\pgfsetdash{}{0pt}%
\pgfpathmoveto{\pgfqpoint{3.096513in}{0.600250in}}%
\pgfpathlineto{\pgfqpoint{3.157750in}{0.600250in}}%
\pgfpathlineto{\pgfqpoint{3.157750in}{0.696500in}}%
\pgfpathlineto{\pgfqpoint{3.096513in}{0.696500in}}%
\pgfpathlineto{\pgfqpoint{3.096513in}{0.600250in}}%
\pgfpathclose%
\pgfusepath{fill}%
\end{pgfscope}%
\begin{pgfscope}%
\pgfpathrectangle{\pgfqpoint{2.857045in}{0.350000in}}{\pgfqpoint{0.300705in}{0.962500in}}%
\pgfusepath{clip}%
\pgfsetbuttcap%
\pgfsetmiterjoin%
\definecolor{currentfill}{rgb}{0.529412,0.742637,0.862899}%
\pgfsetfillcolor{currentfill}%
\pgfsetlinewidth{0.000000pt}%
\definecolor{currentstroke}{rgb}{0.000000,0.000000,0.000000}%
\pgfsetstrokecolor{currentstroke}%
\pgfsetstrokeopacity{0.000000}%
\pgfsetdash{}{0pt}%
\pgfpathmoveto{\pgfqpoint{3.096381in}{0.609875in}}%
\pgfpathlineto{\pgfqpoint{3.157750in}{0.609875in}}%
\pgfpathlineto{\pgfqpoint{3.157750in}{0.706125in}}%
\pgfpathlineto{\pgfqpoint{3.096381in}{0.706125in}}%
\pgfpathlineto{\pgfqpoint{3.096381in}{0.609875in}}%
\pgfpathclose%
\pgfusepath{fill}%
\end{pgfscope}%
\begin{pgfscope}%
\pgfpathrectangle{\pgfqpoint{2.857045in}{0.350000in}}{\pgfqpoint{0.300705in}{0.962500in}}%
\pgfusepath{clip}%
\pgfsetbuttcap%
\pgfsetmiterjoin%
\definecolor{currentfill}{rgb}{0.529412,0.742637,0.862899}%
\pgfsetfillcolor{currentfill}%
\pgfsetlinewidth{0.000000pt}%
\definecolor{currentstroke}{rgb}{0.000000,0.000000,0.000000}%
\pgfsetstrokecolor{currentstroke}%
\pgfsetstrokeopacity{0.000000}%
\pgfsetdash{}{0pt}%
\pgfpathmoveto{\pgfqpoint{3.096347in}{0.619500in}}%
\pgfpathlineto{\pgfqpoint{3.157750in}{0.619500in}}%
\pgfpathlineto{\pgfqpoint{3.157750in}{0.715750in}}%
\pgfpathlineto{\pgfqpoint{3.096347in}{0.715750in}}%
\pgfpathlineto{\pgfqpoint{3.096347in}{0.619500in}}%
\pgfpathclose%
\pgfusepath{fill}%
\end{pgfscope}%
\begin{pgfscope}%
\pgfpathrectangle{\pgfqpoint{2.857045in}{0.350000in}}{\pgfqpoint{0.300705in}{0.962500in}}%
\pgfusepath{clip}%
\pgfsetbuttcap%
\pgfsetmiterjoin%
\definecolor{currentfill}{rgb}{0.529412,0.742637,0.862899}%
\pgfsetfillcolor{currentfill}%
\pgfsetlinewidth{0.000000pt}%
\definecolor{currentstroke}{rgb}{0.000000,0.000000,0.000000}%
\pgfsetstrokecolor{currentstroke}%
\pgfsetstrokeopacity{0.000000}%
\pgfsetdash{}{0pt}%
\pgfpathmoveto{\pgfqpoint{3.096347in}{0.629125in}}%
\pgfpathlineto{\pgfqpoint{3.157750in}{0.629125in}}%
\pgfpathlineto{\pgfqpoint{3.157750in}{0.725375in}}%
\pgfpathlineto{\pgfqpoint{3.096347in}{0.725375in}}%
\pgfpathlineto{\pgfqpoint{3.096347in}{0.629125in}}%
\pgfpathclose%
\pgfusepath{fill}%
\end{pgfscope}%
\begin{pgfscope}%
\pgfpathrectangle{\pgfqpoint{2.857045in}{0.350000in}}{\pgfqpoint{0.300705in}{0.962500in}}%
\pgfusepath{clip}%
\pgfsetbuttcap%
\pgfsetmiterjoin%
\definecolor{currentfill}{rgb}{0.516863,0.735748,0.860192}%
\pgfsetfillcolor{currentfill}%
\pgfsetlinewidth{0.000000pt}%
\definecolor{currentstroke}{rgb}{0.000000,0.000000,0.000000}%
\pgfsetstrokecolor{currentstroke}%
\pgfsetstrokeopacity{0.000000}%
\pgfsetdash{}{0pt}%
\pgfpathmoveto{\pgfqpoint{3.095995in}{0.638750in}}%
\pgfpathlineto{\pgfqpoint{3.157750in}{0.638750in}}%
\pgfpathlineto{\pgfqpoint{3.157750in}{0.735000in}}%
\pgfpathlineto{\pgfqpoint{3.095995in}{0.735000in}}%
\pgfpathlineto{\pgfqpoint{3.095995in}{0.638750in}}%
\pgfpathclose%
\pgfusepath{fill}%
\end{pgfscope}%
\begin{pgfscope}%
\pgfpathrectangle{\pgfqpoint{2.857045in}{0.350000in}}{\pgfqpoint{0.300705in}{0.962500in}}%
\pgfusepath{clip}%
\pgfsetbuttcap%
\pgfsetmiterjoin%
\definecolor{currentfill}{rgb}{0.510588,0.732303,0.858839}%
\pgfsetfillcolor{currentfill}%
\pgfsetlinewidth{0.000000pt}%
\definecolor{currentstroke}{rgb}{0.000000,0.000000,0.000000}%
\pgfsetstrokecolor{currentstroke}%
\pgfsetstrokeopacity{0.000000}%
\pgfsetdash{}{0pt}%
\pgfpathmoveto{\pgfqpoint{3.095716in}{0.648375in}}%
\pgfpathlineto{\pgfqpoint{3.157750in}{0.648375in}}%
\pgfpathlineto{\pgfqpoint{3.157750in}{0.744625in}}%
\pgfpathlineto{\pgfqpoint{3.095716in}{0.744625in}}%
\pgfpathlineto{\pgfqpoint{3.095716in}{0.648375in}}%
\pgfpathclose%
\pgfusepath{fill}%
\end{pgfscope}%
\begin{pgfscope}%
\pgfpathrectangle{\pgfqpoint{2.857045in}{0.350000in}}{\pgfqpoint{0.300705in}{0.962500in}}%
\pgfusepath{clip}%
\pgfsetbuttcap%
\pgfsetmiterjoin%
\definecolor{currentfill}{rgb}{0.504314,0.728858,0.857486}%
\pgfsetfillcolor{currentfill}%
\pgfsetlinewidth{0.000000pt}%
\definecolor{currentstroke}{rgb}{0.000000,0.000000,0.000000}%
\pgfsetstrokecolor{currentstroke}%
\pgfsetstrokeopacity{0.000000}%
\pgfsetdash{}{0pt}%
\pgfpathmoveto{\pgfqpoint{3.095387in}{0.658000in}}%
\pgfpathlineto{\pgfqpoint{3.157750in}{0.658000in}}%
\pgfpathlineto{\pgfqpoint{3.157750in}{0.754250in}}%
\pgfpathlineto{\pgfqpoint{3.095387in}{0.754250in}}%
\pgfpathlineto{\pgfqpoint{3.095387in}{0.658000in}}%
\pgfpathclose%
\pgfusepath{fill}%
\end{pgfscope}%
\begin{pgfscope}%
\pgfpathrectangle{\pgfqpoint{2.857045in}{0.350000in}}{\pgfqpoint{0.300705in}{0.962500in}}%
\pgfusepath{clip}%
\pgfsetbuttcap%
\pgfsetmiterjoin%
\definecolor{currentfill}{rgb}{0.491765,0.721968,0.854779}%
\pgfsetfillcolor{currentfill}%
\pgfsetlinewidth{0.000000pt}%
\definecolor{currentstroke}{rgb}{0.000000,0.000000,0.000000}%
\pgfsetstrokecolor{currentstroke}%
\pgfsetstrokeopacity{0.000000}%
\pgfsetdash{}{0pt}%
\pgfpathmoveto{\pgfqpoint{3.095060in}{0.667625in}}%
\pgfpathlineto{\pgfqpoint{3.157750in}{0.667625in}}%
\pgfpathlineto{\pgfqpoint{3.157750in}{0.763875in}}%
\pgfpathlineto{\pgfqpoint{3.095060in}{0.763875in}}%
\pgfpathlineto{\pgfqpoint{3.095060in}{0.667625in}}%
\pgfpathclose%
\pgfusepath{fill}%
\end{pgfscope}%
\begin{pgfscope}%
\pgfpathrectangle{\pgfqpoint{2.857045in}{0.350000in}}{\pgfqpoint{0.300705in}{0.962500in}}%
\pgfusepath{clip}%
\pgfsetbuttcap%
\pgfsetmiterjoin%
\definecolor{currentfill}{rgb}{0.485490,0.718524,0.853426}%
\pgfsetfillcolor{currentfill}%
\pgfsetlinewidth{0.000000pt}%
\definecolor{currentstroke}{rgb}{0.000000,0.000000,0.000000}%
\pgfsetstrokecolor{currentstroke}%
\pgfsetstrokeopacity{0.000000}%
\pgfsetdash{}{0pt}%
\pgfpathmoveto{\pgfqpoint{3.094912in}{0.677250in}}%
\pgfpathlineto{\pgfqpoint{3.157750in}{0.677250in}}%
\pgfpathlineto{\pgfqpoint{3.157750in}{0.773500in}}%
\pgfpathlineto{\pgfqpoint{3.094912in}{0.773500in}}%
\pgfpathlineto{\pgfqpoint{3.094912in}{0.677250in}}%
\pgfpathclose%
\pgfusepath{fill}%
\end{pgfscope}%
\begin{pgfscope}%
\pgfpathrectangle{\pgfqpoint{2.857045in}{0.350000in}}{\pgfqpoint{0.300705in}{0.962500in}}%
\pgfusepath{clip}%
\pgfsetbuttcap%
\pgfsetmiterjoin%
\definecolor{currentfill}{rgb}{0.485490,0.718524,0.853426}%
\pgfsetfillcolor{currentfill}%
\pgfsetlinewidth{0.000000pt}%
\definecolor{currentstroke}{rgb}{0.000000,0.000000,0.000000}%
\pgfsetstrokecolor{currentstroke}%
\pgfsetstrokeopacity{0.000000}%
\pgfsetdash{}{0pt}%
\pgfpathmoveto{\pgfqpoint{3.094713in}{0.686875in}}%
\pgfpathlineto{\pgfqpoint{3.157750in}{0.686875in}}%
\pgfpathlineto{\pgfqpoint{3.157750in}{0.783125in}}%
\pgfpathlineto{\pgfqpoint{3.094713in}{0.783125in}}%
\pgfpathlineto{\pgfqpoint{3.094713in}{0.686875in}}%
\pgfpathclose%
\pgfusepath{fill}%
\end{pgfscope}%
\begin{pgfscope}%
\pgfpathrectangle{\pgfqpoint{2.857045in}{0.350000in}}{\pgfqpoint{0.300705in}{0.962500in}}%
\pgfusepath{clip}%
\pgfsetbuttcap%
\pgfsetmiterjoin%
\definecolor{currentfill}{rgb}{0.479216,0.715079,0.852072}%
\pgfsetfillcolor{currentfill}%
\pgfsetlinewidth{0.000000pt}%
\definecolor{currentstroke}{rgb}{0.000000,0.000000,0.000000}%
\pgfsetstrokecolor{currentstroke}%
\pgfsetstrokeopacity{0.000000}%
\pgfsetdash{}{0pt}%
\pgfpathmoveto{\pgfqpoint{3.094672in}{0.696500in}}%
\pgfpathlineto{\pgfqpoint{3.157750in}{0.696500in}}%
\pgfpathlineto{\pgfqpoint{3.157750in}{0.792750in}}%
\pgfpathlineto{\pgfqpoint{3.094672in}{0.792750in}}%
\pgfpathlineto{\pgfqpoint{3.094672in}{0.696500in}}%
\pgfpathclose%
\pgfusepath{fill}%
\end{pgfscope}%
\begin{pgfscope}%
\pgfpathrectangle{\pgfqpoint{2.857045in}{0.350000in}}{\pgfqpoint{0.300705in}{0.962500in}}%
\pgfusepath{clip}%
\pgfsetbuttcap%
\pgfsetmiterjoin%
\definecolor{currentfill}{rgb}{0.479216,0.715079,0.852072}%
\pgfsetfillcolor{currentfill}%
\pgfsetlinewidth{0.000000pt}%
\definecolor{currentstroke}{rgb}{0.000000,0.000000,0.000000}%
\pgfsetstrokecolor{currentstroke}%
\pgfsetstrokeopacity{0.000000}%
\pgfsetdash{}{0pt}%
\pgfpathmoveto{\pgfqpoint{3.094672in}{0.706125in}}%
\pgfpathlineto{\pgfqpoint{3.157750in}{0.706125in}}%
\pgfpathlineto{\pgfqpoint{3.157750in}{0.802375in}}%
\pgfpathlineto{\pgfqpoint{3.094672in}{0.802375in}}%
\pgfpathlineto{\pgfqpoint{3.094672in}{0.706125in}}%
\pgfpathclose%
\pgfusepath{fill}%
\end{pgfscope}%
\begin{pgfscope}%
\pgfpathrectangle{\pgfqpoint{2.857045in}{0.350000in}}{\pgfqpoint{0.300705in}{0.962500in}}%
\pgfusepath{clip}%
\pgfsetbuttcap%
\pgfsetmiterjoin%
\definecolor{currentfill}{rgb}{0.460392,0.704744,0.848012}%
\pgfsetfillcolor{currentfill}%
\pgfsetlinewidth{0.000000pt}%
\definecolor{currentstroke}{rgb}{0.000000,0.000000,0.000000}%
\pgfsetstrokecolor{currentstroke}%
\pgfsetstrokeopacity{0.000000}%
\pgfsetdash{}{0pt}%
\pgfpathmoveto{\pgfqpoint{3.094015in}{0.715750in}}%
\pgfpathlineto{\pgfqpoint{3.157750in}{0.715750in}}%
\pgfpathlineto{\pgfqpoint{3.157750in}{0.812000in}}%
\pgfpathlineto{\pgfqpoint{3.094015in}{0.812000in}}%
\pgfpathlineto{\pgfqpoint{3.094015in}{0.715750in}}%
\pgfpathclose%
\pgfusepath{fill}%
\end{pgfscope}%
\begin{pgfscope}%
\pgfpathrectangle{\pgfqpoint{2.857045in}{0.350000in}}{\pgfqpoint{0.300705in}{0.962500in}}%
\pgfusepath{clip}%
\pgfsetbuttcap%
\pgfsetmiterjoin%
\definecolor{currentfill}{rgb}{0.460392,0.704744,0.848012}%
\pgfsetfillcolor{currentfill}%
\pgfsetlinewidth{0.000000pt}%
\definecolor{currentstroke}{rgb}{0.000000,0.000000,0.000000}%
\pgfsetstrokecolor{currentstroke}%
\pgfsetstrokeopacity{0.000000}%
\pgfsetdash{}{0pt}%
\pgfpathmoveto{\pgfqpoint{3.093973in}{0.725375in}}%
\pgfpathlineto{\pgfqpoint{3.157750in}{0.725375in}}%
\pgfpathlineto{\pgfqpoint{3.157750in}{0.821625in}}%
\pgfpathlineto{\pgfqpoint{3.093973in}{0.821625in}}%
\pgfpathlineto{\pgfqpoint{3.093973in}{0.725375in}}%
\pgfpathclose%
\pgfusepath{fill}%
\end{pgfscope}%
\begin{pgfscope}%
\pgfpathrectangle{\pgfqpoint{2.857045in}{0.350000in}}{\pgfqpoint{0.300705in}{0.962500in}}%
\pgfusepath{clip}%
\pgfsetbuttcap%
\pgfsetmiterjoin%
\definecolor{currentfill}{rgb}{0.460392,0.704744,0.848012}%
\pgfsetfillcolor{currentfill}%
\pgfsetlinewidth{0.000000pt}%
\definecolor{currentstroke}{rgb}{0.000000,0.000000,0.000000}%
\pgfsetstrokecolor{currentstroke}%
\pgfsetstrokeopacity{0.000000}%
\pgfsetdash{}{0pt}%
\pgfpathmoveto{\pgfqpoint{3.093972in}{0.735000in}}%
\pgfpathlineto{\pgfqpoint{3.157750in}{0.735000in}}%
\pgfpathlineto{\pgfqpoint{3.157750in}{0.831250in}}%
\pgfpathlineto{\pgfqpoint{3.093972in}{0.831250in}}%
\pgfpathlineto{\pgfqpoint{3.093972in}{0.735000in}}%
\pgfpathclose%
\pgfusepath{fill}%
\end{pgfscope}%
\begin{pgfscope}%
\pgfpathrectangle{\pgfqpoint{2.857045in}{0.350000in}}{\pgfqpoint{0.300705in}{0.962500in}}%
\pgfusepath{clip}%
\pgfsetbuttcap%
\pgfsetmiterjoin%
\definecolor{currentfill}{rgb}{0.460392,0.704744,0.848012}%
\pgfsetfillcolor{currentfill}%
\pgfsetlinewidth{0.000000pt}%
\definecolor{currentstroke}{rgb}{0.000000,0.000000,0.000000}%
\pgfsetstrokecolor{currentstroke}%
\pgfsetstrokeopacity{0.000000}%
\pgfsetdash{}{0pt}%
\pgfpathmoveto{\pgfqpoint{3.093896in}{0.744625in}}%
\pgfpathlineto{\pgfqpoint{3.157750in}{0.744625in}}%
\pgfpathlineto{\pgfqpoint{3.157750in}{0.840875in}}%
\pgfpathlineto{\pgfqpoint{3.093896in}{0.840875in}}%
\pgfpathlineto{\pgfqpoint{3.093896in}{0.744625in}}%
\pgfpathclose%
\pgfusepath{fill}%
\end{pgfscope}%
\begin{pgfscope}%
\pgfpathrectangle{\pgfqpoint{2.857045in}{0.350000in}}{\pgfqpoint{0.300705in}{0.962500in}}%
\pgfusepath{clip}%
\pgfsetbuttcap%
\pgfsetmiterjoin%
\definecolor{currentfill}{rgb}{0.441569,0.694410,0.843952}%
\pgfsetfillcolor{currentfill}%
\pgfsetlinewidth{0.000000pt}%
\definecolor{currentstroke}{rgb}{0.000000,0.000000,0.000000}%
\pgfsetstrokecolor{currentstroke}%
\pgfsetstrokeopacity{0.000000}%
\pgfsetdash{}{0pt}%
\pgfpathmoveto{\pgfqpoint{3.093315in}{0.754250in}}%
\pgfpathlineto{\pgfqpoint{3.157750in}{0.754250in}}%
\pgfpathlineto{\pgfqpoint{3.157750in}{0.850500in}}%
\pgfpathlineto{\pgfqpoint{3.093315in}{0.850500in}}%
\pgfpathlineto{\pgfqpoint{3.093315in}{0.754250in}}%
\pgfpathclose%
\pgfusepath{fill}%
\end{pgfscope}%
\begin{pgfscope}%
\pgfpathrectangle{\pgfqpoint{2.857045in}{0.350000in}}{\pgfqpoint{0.300705in}{0.962500in}}%
\pgfusepath{clip}%
\pgfsetbuttcap%
\pgfsetmiterjoin%
\definecolor{currentfill}{rgb}{0.435294,0.690965,0.842599}%
\pgfsetfillcolor{currentfill}%
\pgfsetlinewidth{0.000000pt}%
\definecolor{currentstroke}{rgb}{0.000000,0.000000,0.000000}%
\pgfsetstrokecolor{currentstroke}%
\pgfsetstrokeopacity{0.000000}%
\pgfsetdash{}{0pt}%
\pgfpathmoveto{\pgfqpoint{3.093140in}{0.763875in}}%
\pgfpathlineto{\pgfqpoint{3.157750in}{0.763875in}}%
\pgfpathlineto{\pgfqpoint{3.157750in}{0.860125in}}%
\pgfpathlineto{\pgfqpoint{3.093140in}{0.860125in}}%
\pgfpathlineto{\pgfqpoint{3.093140in}{0.763875in}}%
\pgfpathclose%
\pgfusepath{fill}%
\end{pgfscope}%
\begin{pgfscope}%
\pgfpathrectangle{\pgfqpoint{2.857045in}{0.350000in}}{\pgfqpoint{0.300705in}{0.962500in}}%
\pgfusepath{clip}%
\pgfsetbuttcap%
\pgfsetmiterjoin%
\definecolor{currentfill}{rgb}{0.435294,0.690965,0.842599}%
\pgfsetfillcolor{currentfill}%
\pgfsetlinewidth{0.000000pt}%
\definecolor{currentstroke}{rgb}{0.000000,0.000000,0.000000}%
\pgfsetstrokecolor{currentstroke}%
\pgfsetstrokeopacity{0.000000}%
\pgfsetdash{}{0pt}%
\pgfpathmoveto{\pgfqpoint{3.093114in}{0.773500in}}%
\pgfpathlineto{\pgfqpoint{3.157750in}{0.773500in}}%
\pgfpathlineto{\pgfqpoint{3.157750in}{0.869750in}}%
\pgfpathlineto{\pgfqpoint{3.093114in}{0.869750in}}%
\pgfpathlineto{\pgfqpoint{3.093114in}{0.773500in}}%
\pgfpathclose%
\pgfusepath{fill}%
\end{pgfscope}%
\begin{pgfscope}%
\pgfpathrectangle{\pgfqpoint{2.857045in}{0.350000in}}{\pgfqpoint{0.300705in}{0.962500in}}%
\pgfusepath{clip}%
\pgfsetbuttcap%
\pgfsetmiterjoin%
\definecolor{currentfill}{rgb}{0.429020,0.687520,0.841246}%
\pgfsetfillcolor{currentfill}%
\pgfsetlinewidth{0.000000pt}%
\definecolor{currentstroke}{rgb}{0.000000,0.000000,0.000000}%
\pgfsetstrokecolor{currentstroke}%
\pgfsetstrokeopacity{0.000000}%
\pgfsetdash{}{0pt}%
\pgfpathmoveto{\pgfqpoint{3.092877in}{0.783125in}}%
\pgfpathlineto{\pgfqpoint{3.157750in}{0.783125in}}%
\pgfpathlineto{\pgfqpoint{3.157750in}{0.879375in}}%
\pgfpathlineto{\pgfqpoint{3.092877in}{0.879375in}}%
\pgfpathlineto{\pgfqpoint{3.092877in}{0.783125in}}%
\pgfpathclose%
\pgfusepath{fill}%
\end{pgfscope}%
\begin{pgfscope}%
\pgfpathrectangle{\pgfqpoint{2.857045in}{0.350000in}}{\pgfqpoint{0.300705in}{0.962500in}}%
\pgfusepath{clip}%
\pgfsetbuttcap%
\pgfsetmiterjoin%
\definecolor{currentfill}{rgb}{0.429020,0.687520,0.841246}%
\pgfsetfillcolor{currentfill}%
\pgfsetlinewidth{0.000000pt}%
\definecolor{currentstroke}{rgb}{0.000000,0.000000,0.000000}%
\pgfsetstrokecolor{currentstroke}%
\pgfsetstrokeopacity{0.000000}%
\pgfsetdash{}{0pt}%
\pgfpathmoveto{\pgfqpoint{3.092877in}{0.792750in}}%
\pgfpathlineto{\pgfqpoint{3.157750in}{0.792750in}}%
\pgfpathlineto{\pgfqpoint{3.157750in}{0.889000in}}%
\pgfpathlineto{\pgfqpoint{3.092877in}{0.889000in}}%
\pgfpathlineto{\pgfqpoint{3.092877in}{0.792750in}}%
\pgfpathclose%
\pgfusepath{fill}%
\end{pgfscope}%
\begin{pgfscope}%
\pgfpathrectangle{\pgfqpoint{2.857045in}{0.350000in}}{\pgfqpoint{0.300705in}{0.962500in}}%
\pgfusepath{clip}%
\pgfsetbuttcap%
\pgfsetmiterjoin%
\definecolor{currentfill}{rgb}{0.429020,0.687520,0.841246}%
\pgfsetfillcolor{currentfill}%
\pgfsetlinewidth{0.000000pt}%
\definecolor{currentstroke}{rgb}{0.000000,0.000000,0.000000}%
\pgfsetstrokecolor{currentstroke}%
\pgfsetstrokeopacity{0.000000}%
\pgfsetdash{}{0pt}%
\pgfpathmoveto{\pgfqpoint{3.092877in}{0.802375in}}%
\pgfpathlineto{\pgfqpoint{3.157750in}{0.802375in}}%
\pgfpathlineto{\pgfqpoint{3.157750in}{0.898625in}}%
\pgfpathlineto{\pgfqpoint{3.092877in}{0.898625in}}%
\pgfpathlineto{\pgfqpoint{3.092877in}{0.802375in}}%
\pgfpathclose%
\pgfusepath{fill}%
\end{pgfscope}%
\begin{pgfscope}%
\pgfpathrectangle{\pgfqpoint{2.857045in}{0.350000in}}{\pgfqpoint{0.300705in}{0.962500in}}%
\pgfusepath{clip}%
\pgfsetbuttcap%
\pgfsetmiterjoin%
\definecolor{currentfill}{rgb}{0.429020,0.687520,0.841246}%
\pgfsetfillcolor{currentfill}%
\pgfsetlinewidth{0.000000pt}%
\definecolor{currentstroke}{rgb}{0.000000,0.000000,0.000000}%
\pgfsetstrokecolor{currentstroke}%
\pgfsetstrokeopacity{0.000000}%
\pgfsetdash{}{0pt}%
\pgfpathmoveto{\pgfqpoint{3.092877in}{0.812000in}}%
\pgfpathlineto{\pgfqpoint{3.157750in}{0.812000in}}%
\pgfpathlineto{\pgfqpoint{3.157750in}{0.908250in}}%
\pgfpathlineto{\pgfqpoint{3.092877in}{0.908250in}}%
\pgfpathlineto{\pgfqpoint{3.092877in}{0.812000in}}%
\pgfpathclose%
\pgfusepath{fill}%
\end{pgfscope}%
\begin{pgfscope}%
\pgfpathrectangle{\pgfqpoint{2.857045in}{0.350000in}}{\pgfqpoint{0.300705in}{0.962500in}}%
\pgfusepath{clip}%
\pgfsetbuttcap%
\pgfsetmiterjoin%
\definecolor{currentfill}{rgb}{0.429020,0.687520,0.841246}%
\pgfsetfillcolor{currentfill}%
\pgfsetlinewidth{0.000000pt}%
\definecolor{currentstroke}{rgb}{0.000000,0.000000,0.000000}%
\pgfsetstrokecolor{currentstroke}%
\pgfsetstrokeopacity{0.000000}%
\pgfsetdash{}{0pt}%
\pgfpathmoveto{\pgfqpoint{3.092877in}{0.821625in}}%
\pgfpathlineto{\pgfqpoint{3.157750in}{0.821625in}}%
\pgfpathlineto{\pgfqpoint{3.157750in}{0.917875in}}%
\pgfpathlineto{\pgfqpoint{3.092877in}{0.917875in}}%
\pgfpathlineto{\pgfqpoint{3.092877in}{0.821625in}}%
\pgfpathclose%
\pgfusepath{fill}%
\end{pgfscope}%
\begin{pgfscope}%
\pgfpathrectangle{\pgfqpoint{2.857045in}{0.350000in}}{\pgfqpoint{0.300705in}{0.962500in}}%
\pgfusepath{clip}%
\pgfsetbuttcap%
\pgfsetmiterjoin%
\definecolor{currentfill}{rgb}{0.429020,0.687520,0.841246}%
\pgfsetfillcolor{currentfill}%
\pgfsetlinewidth{0.000000pt}%
\definecolor{currentstroke}{rgb}{0.000000,0.000000,0.000000}%
\pgfsetstrokecolor{currentstroke}%
\pgfsetstrokeopacity{0.000000}%
\pgfsetdash{}{0pt}%
\pgfpathmoveto{\pgfqpoint{3.092877in}{0.831250in}}%
\pgfpathlineto{\pgfqpoint{3.157750in}{0.831250in}}%
\pgfpathlineto{\pgfqpoint{3.157750in}{0.927500in}}%
\pgfpathlineto{\pgfqpoint{3.092877in}{0.927500in}}%
\pgfpathlineto{\pgfqpoint{3.092877in}{0.831250in}}%
\pgfpathclose%
\pgfusepath{fill}%
\end{pgfscope}%
\begin{pgfscope}%
\pgfpathrectangle{\pgfqpoint{2.857045in}{0.350000in}}{\pgfqpoint{0.300705in}{0.962500in}}%
\pgfusepath{clip}%
\pgfsetbuttcap%
\pgfsetmiterjoin%
\definecolor{currentfill}{rgb}{0.429020,0.687520,0.841246}%
\pgfsetfillcolor{currentfill}%
\pgfsetlinewidth{0.000000pt}%
\definecolor{currentstroke}{rgb}{0.000000,0.000000,0.000000}%
\pgfsetstrokecolor{currentstroke}%
\pgfsetstrokeopacity{0.000000}%
\pgfsetdash{}{0pt}%
\pgfpathmoveto{\pgfqpoint{3.092809in}{0.840875in}}%
\pgfpathlineto{\pgfqpoint{3.157750in}{0.840875in}}%
\pgfpathlineto{\pgfqpoint{3.157750in}{0.937125in}}%
\pgfpathlineto{\pgfqpoint{3.092809in}{0.937125in}}%
\pgfpathlineto{\pgfqpoint{3.092809in}{0.840875in}}%
\pgfpathclose%
\pgfusepath{fill}%
\end{pgfscope}%
\begin{pgfscope}%
\pgfpathrectangle{\pgfqpoint{2.857045in}{0.350000in}}{\pgfqpoint{0.300705in}{0.962500in}}%
\pgfusepath{clip}%
\pgfsetbuttcap%
\pgfsetmiterjoin%
\definecolor{currentfill}{rgb}{0.422745,0.684075,0.839892}%
\pgfsetfillcolor{currentfill}%
\pgfsetlinewidth{0.000000pt}%
\definecolor{currentstroke}{rgb}{0.000000,0.000000,0.000000}%
\pgfsetstrokecolor{currentstroke}%
\pgfsetstrokeopacity{0.000000}%
\pgfsetdash{}{0pt}%
\pgfpathmoveto{\pgfqpoint{3.092707in}{0.850500in}}%
\pgfpathlineto{\pgfqpoint{3.157750in}{0.850500in}}%
\pgfpathlineto{\pgfqpoint{3.157750in}{0.946750in}}%
\pgfpathlineto{\pgfqpoint{3.092707in}{0.946750in}}%
\pgfpathlineto{\pgfqpoint{3.092707in}{0.850500in}}%
\pgfpathclose%
\pgfusepath{fill}%
\end{pgfscope}%
\begin{pgfscope}%
\pgfpathrectangle{\pgfqpoint{2.857045in}{0.350000in}}{\pgfqpoint{0.300705in}{0.962500in}}%
\pgfusepath{clip}%
\pgfsetbuttcap%
\pgfsetmiterjoin%
\definecolor{currentfill}{rgb}{0.422745,0.684075,0.839892}%
\pgfsetfillcolor{currentfill}%
\pgfsetlinewidth{0.000000pt}%
\definecolor{currentstroke}{rgb}{0.000000,0.000000,0.000000}%
\pgfsetstrokecolor{currentstroke}%
\pgfsetstrokeopacity{0.000000}%
\pgfsetdash{}{0pt}%
\pgfpathmoveto{\pgfqpoint{3.092541in}{0.860125in}}%
\pgfpathlineto{\pgfqpoint{3.157750in}{0.860125in}}%
\pgfpathlineto{\pgfqpoint{3.157750in}{0.956375in}}%
\pgfpathlineto{\pgfqpoint{3.092541in}{0.956375in}}%
\pgfpathlineto{\pgfqpoint{3.092541in}{0.860125in}}%
\pgfpathclose%
\pgfusepath{fill}%
\end{pgfscope}%
\begin{pgfscope}%
\pgfpathrectangle{\pgfqpoint{2.857045in}{0.350000in}}{\pgfqpoint{0.300705in}{0.962500in}}%
\pgfusepath{clip}%
\pgfsetbuttcap%
\pgfsetmiterjoin%
\definecolor{currentfill}{rgb}{0.422745,0.684075,0.839892}%
\pgfsetfillcolor{currentfill}%
\pgfsetlinewidth{0.000000pt}%
\definecolor{currentstroke}{rgb}{0.000000,0.000000,0.000000}%
\pgfsetstrokecolor{currentstroke}%
\pgfsetstrokeopacity{0.000000}%
\pgfsetdash{}{0pt}%
\pgfpathmoveto{\pgfqpoint{3.092541in}{0.869750in}}%
\pgfpathlineto{\pgfqpoint{3.157750in}{0.869750in}}%
\pgfpathlineto{\pgfqpoint{3.157750in}{0.966000in}}%
\pgfpathlineto{\pgfqpoint{3.092541in}{0.966000in}}%
\pgfpathlineto{\pgfqpoint{3.092541in}{0.869750in}}%
\pgfpathclose%
\pgfusepath{fill}%
\end{pgfscope}%
\begin{pgfscope}%
\pgfpathrectangle{\pgfqpoint{2.857045in}{0.350000in}}{\pgfqpoint{0.300705in}{0.962500in}}%
\pgfusepath{clip}%
\pgfsetbuttcap%
\pgfsetmiterjoin%
\definecolor{currentfill}{rgb}{0.417086,0.680631,0.838231}%
\pgfsetfillcolor{currentfill}%
\pgfsetlinewidth{0.000000pt}%
\definecolor{currentstroke}{rgb}{0.000000,0.000000,0.000000}%
\pgfsetstrokecolor{currentstroke}%
\pgfsetstrokeopacity{0.000000}%
\pgfsetdash{}{0pt}%
\pgfpathmoveto{\pgfqpoint{3.092485in}{0.879375in}}%
\pgfpathlineto{\pgfqpoint{3.157750in}{0.879375in}}%
\pgfpathlineto{\pgfqpoint{3.157750in}{0.975625in}}%
\pgfpathlineto{\pgfqpoint{3.092485in}{0.975625in}}%
\pgfpathlineto{\pgfqpoint{3.092485in}{0.879375in}}%
\pgfpathclose%
\pgfusepath{fill}%
\end{pgfscope}%
\begin{pgfscope}%
\pgfpathrectangle{\pgfqpoint{2.857045in}{0.350000in}}{\pgfqpoint{0.300705in}{0.962500in}}%
\pgfusepath{clip}%
\pgfsetbuttcap%
\pgfsetmiterjoin%
\definecolor{currentfill}{rgb}{0.417086,0.680631,0.838231}%
\pgfsetfillcolor{currentfill}%
\pgfsetlinewidth{0.000000pt}%
\definecolor{currentstroke}{rgb}{0.000000,0.000000,0.000000}%
\pgfsetstrokecolor{currentstroke}%
\pgfsetstrokeopacity{0.000000}%
\pgfsetdash{}{0pt}%
\pgfpathmoveto{\pgfqpoint{3.092350in}{0.889000in}}%
\pgfpathlineto{\pgfqpoint{3.157750in}{0.889000in}}%
\pgfpathlineto{\pgfqpoint{3.157750in}{0.985250in}}%
\pgfpathlineto{\pgfqpoint{3.092350in}{0.985250in}}%
\pgfpathlineto{\pgfqpoint{3.092350in}{0.889000in}}%
\pgfpathclose%
\pgfusepath{fill}%
\end{pgfscope}%
\begin{pgfscope}%
\pgfpathrectangle{\pgfqpoint{2.857045in}{0.350000in}}{\pgfqpoint{0.300705in}{0.962500in}}%
\pgfusepath{clip}%
\pgfsetbuttcap%
\pgfsetmiterjoin%
\definecolor{currentfill}{rgb}{0.417086,0.680631,0.838231}%
\pgfsetfillcolor{currentfill}%
\pgfsetlinewidth{0.000000pt}%
\definecolor{currentstroke}{rgb}{0.000000,0.000000,0.000000}%
\pgfsetstrokecolor{currentstroke}%
\pgfsetstrokeopacity{0.000000}%
\pgfsetdash{}{0pt}%
\pgfpathmoveto{\pgfqpoint{3.092295in}{0.898625in}}%
\pgfpathlineto{\pgfqpoint{3.157750in}{0.898625in}}%
\pgfpathlineto{\pgfqpoint{3.157750in}{0.994875in}}%
\pgfpathlineto{\pgfqpoint{3.092295in}{0.994875in}}%
\pgfpathlineto{\pgfqpoint{3.092295in}{0.898625in}}%
\pgfpathclose%
\pgfusepath{fill}%
\end{pgfscope}%
\begin{pgfscope}%
\pgfpathrectangle{\pgfqpoint{2.857045in}{0.350000in}}{\pgfqpoint{0.300705in}{0.962500in}}%
\pgfusepath{clip}%
\pgfsetbuttcap%
\pgfsetmiterjoin%
\definecolor{currentfill}{rgb}{0.417086,0.680631,0.838231}%
\pgfsetfillcolor{currentfill}%
\pgfsetlinewidth{0.000000pt}%
\definecolor{currentstroke}{rgb}{0.000000,0.000000,0.000000}%
\pgfsetstrokecolor{currentstroke}%
\pgfsetstrokeopacity{0.000000}%
\pgfsetdash{}{0pt}%
\pgfpathmoveto{\pgfqpoint{3.092295in}{0.908250in}}%
\pgfpathlineto{\pgfqpoint{3.157750in}{0.908250in}}%
\pgfpathlineto{\pgfqpoint{3.157750in}{1.004500in}}%
\pgfpathlineto{\pgfqpoint{3.092295in}{1.004500in}}%
\pgfpathlineto{\pgfqpoint{3.092295in}{0.908250in}}%
\pgfpathclose%
\pgfusepath{fill}%
\end{pgfscope}%
\begin{pgfscope}%
\pgfpathrectangle{\pgfqpoint{2.857045in}{0.350000in}}{\pgfqpoint{0.300705in}{0.962500in}}%
\pgfusepath{clip}%
\pgfsetbuttcap%
\pgfsetmiterjoin%
\definecolor{currentfill}{rgb}{0.406997,0.673741,0.834295}%
\pgfsetfillcolor{currentfill}%
\pgfsetlinewidth{0.000000pt}%
\definecolor{currentstroke}{rgb}{0.000000,0.000000,0.000000}%
\pgfsetstrokecolor{currentstroke}%
\pgfsetstrokeopacity{0.000000}%
\pgfsetdash{}{0pt}%
\pgfpathmoveto{\pgfqpoint{3.092038in}{0.917875in}}%
\pgfpathlineto{\pgfqpoint{3.157750in}{0.917875in}}%
\pgfpathlineto{\pgfqpoint{3.157750in}{1.014125in}}%
\pgfpathlineto{\pgfqpoint{3.092038in}{1.014125in}}%
\pgfpathlineto{\pgfqpoint{3.092038in}{0.917875in}}%
\pgfpathclose%
\pgfusepath{fill}%
\end{pgfscope}%
\begin{pgfscope}%
\pgfpathrectangle{\pgfqpoint{2.857045in}{0.350000in}}{\pgfqpoint{0.300705in}{0.962500in}}%
\pgfusepath{clip}%
\pgfsetbuttcap%
\pgfsetmiterjoin%
\definecolor{currentfill}{rgb}{0.406997,0.673741,0.834295}%
\pgfsetfillcolor{currentfill}%
\pgfsetlinewidth{0.000000pt}%
\definecolor{currentstroke}{rgb}{0.000000,0.000000,0.000000}%
\pgfsetstrokecolor{currentstroke}%
\pgfsetstrokeopacity{0.000000}%
\pgfsetdash{}{0pt}%
\pgfpathmoveto{\pgfqpoint{3.091863in}{0.927500in}}%
\pgfpathlineto{\pgfqpoint{3.157750in}{0.927500in}}%
\pgfpathlineto{\pgfqpoint{3.157750in}{1.023750in}}%
\pgfpathlineto{\pgfqpoint{3.091863in}{1.023750in}}%
\pgfpathlineto{\pgfqpoint{3.091863in}{0.927500in}}%
\pgfpathclose%
\pgfusepath{fill}%
\end{pgfscope}%
\begin{pgfscope}%
\pgfpathrectangle{\pgfqpoint{2.857045in}{0.350000in}}{\pgfqpoint{0.300705in}{0.962500in}}%
\pgfusepath{clip}%
\pgfsetbuttcap%
\pgfsetmiterjoin%
\definecolor{currentfill}{rgb}{0.401953,0.670296,0.832326}%
\pgfsetfillcolor{currentfill}%
\pgfsetlinewidth{0.000000pt}%
\definecolor{currentstroke}{rgb}{0.000000,0.000000,0.000000}%
\pgfsetstrokecolor{currentstroke}%
\pgfsetstrokeopacity{0.000000}%
\pgfsetdash{}{0pt}%
\pgfpathmoveto{\pgfqpoint{3.091782in}{0.937125in}}%
\pgfpathlineto{\pgfqpoint{3.157750in}{0.937125in}}%
\pgfpathlineto{\pgfqpoint{3.157750in}{1.033375in}}%
\pgfpathlineto{\pgfqpoint{3.091782in}{1.033375in}}%
\pgfpathlineto{\pgfqpoint{3.091782in}{0.937125in}}%
\pgfpathclose%
\pgfusepath{fill}%
\end{pgfscope}%
\begin{pgfscope}%
\pgfpathrectangle{\pgfqpoint{2.857045in}{0.350000in}}{\pgfqpoint{0.300705in}{0.962500in}}%
\pgfusepath{clip}%
\pgfsetbuttcap%
\pgfsetmiterjoin%
\definecolor{currentfill}{rgb}{0.396909,0.666851,0.830358}%
\pgfsetfillcolor{currentfill}%
\pgfsetlinewidth{0.000000pt}%
\definecolor{currentstroke}{rgb}{0.000000,0.000000,0.000000}%
\pgfsetstrokecolor{currentstroke}%
\pgfsetstrokeopacity{0.000000}%
\pgfsetdash{}{0pt}%
\pgfpathmoveto{\pgfqpoint{3.091583in}{0.946750in}}%
\pgfpathlineto{\pgfqpoint{3.157750in}{0.946750in}}%
\pgfpathlineto{\pgfqpoint{3.157750in}{1.043000in}}%
\pgfpathlineto{\pgfqpoint{3.091583in}{1.043000in}}%
\pgfpathlineto{\pgfqpoint{3.091583in}{0.946750in}}%
\pgfpathclose%
\pgfusepath{fill}%
\end{pgfscope}%
\begin{pgfscope}%
\pgfpathrectangle{\pgfqpoint{2.857045in}{0.350000in}}{\pgfqpoint{0.300705in}{0.962500in}}%
\pgfusepath{clip}%
\pgfsetbuttcap%
\pgfsetmiterjoin%
\definecolor{currentfill}{rgb}{0.396909,0.666851,0.830358}%
\pgfsetfillcolor{currentfill}%
\pgfsetlinewidth{0.000000pt}%
\definecolor{currentstroke}{rgb}{0.000000,0.000000,0.000000}%
\pgfsetstrokecolor{currentstroke}%
\pgfsetstrokeopacity{0.000000}%
\pgfsetdash{}{0pt}%
\pgfpathmoveto{\pgfqpoint{3.091572in}{0.956375in}}%
\pgfpathlineto{\pgfqpoint{3.157750in}{0.956375in}}%
\pgfpathlineto{\pgfqpoint{3.157750in}{1.052625in}}%
\pgfpathlineto{\pgfqpoint{3.091572in}{1.052625in}}%
\pgfpathlineto{\pgfqpoint{3.091572in}{0.956375in}}%
\pgfpathclose%
\pgfusepath{fill}%
\end{pgfscope}%
\begin{pgfscope}%
\pgfpathrectangle{\pgfqpoint{2.857045in}{0.350000in}}{\pgfqpoint{0.300705in}{0.962500in}}%
\pgfusepath{clip}%
\pgfsetbuttcap%
\pgfsetmiterjoin%
\definecolor{currentfill}{rgb}{0.391865,0.663406,0.828389}%
\pgfsetfillcolor{currentfill}%
\pgfsetlinewidth{0.000000pt}%
\definecolor{currentstroke}{rgb}{0.000000,0.000000,0.000000}%
\pgfsetstrokecolor{currentstroke}%
\pgfsetstrokeopacity{0.000000}%
\pgfsetdash{}{0pt}%
\pgfpathmoveto{\pgfqpoint{3.091289in}{0.966000in}}%
\pgfpathlineto{\pgfqpoint{3.157750in}{0.966000in}}%
\pgfpathlineto{\pgfqpoint{3.157750in}{1.062250in}}%
\pgfpathlineto{\pgfqpoint{3.091289in}{1.062250in}}%
\pgfpathlineto{\pgfqpoint{3.091289in}{0.966000in}}%
\pgfpathclose%
\pgfusepath{fill}%
\end{pgfscope}%
\begin{pgfscope}%
\pgfpathrectangle{\pgfqpoint{2.857045in}{0.350000in}}{\pgfqpoint{0.300705in}{0.962500in}}%
\pgfusepath{clip}%
\pgfsetbuttcap%
\pgfsetmiterjoin%
\definecolor{currentfill}{rgb}{0.381776,0.656517,0.824452}%
\pgfsetfillcolor{currentfill}%
\pgfsetlinewidth{0.000000pt}%
\definecolor{currentstroke}{rgb}{0.000000,0.000000,0.000000}%
\pgfsetstrokecolor{currentstroke}%
\pgfsetstrokeopacity{0.000000}%
\pgfsetdash{}{0pt}%
\pgfpathmoveto{\pgfqpoint{3.090894in}{0.975625in}}%
\pgfpathlineto{\pgfqpoint{3.157750in}{0.975625in}}%
\pgfpathlineto{\pgfqpoint{3.157750in}{1.071875in}}%
\pgfpathlineto{\pgfqpoint{3.090894in}{1.071875in}}%
\pgfpathlineto{\pgfqpoint{3.090894in}{0.975625in}}%
\pgfpathclose%
\pgfusepath{fill}%
\end{pgfscope}%
\begin{pgfscope}%
\pgfpathrectangle{\pgfqpoint{2.857045in}{0.350000in}}{\pgfqpoint{0.300705in}{0.962500in}}%
\pgfusepath{clip}%
\pgfsetbuttcap%
\pgfsetmiterjoin%
\definecolor{currentfill}{rgb}{0.381776,0.656517,0.824452}%
\pgfsetfillcolor{currentfill}%
\pgfsetlinewidth{0.000000pt}%
\definecolor{currentstroke}{rgb}{0.000000,0.000000,0.000000}%
\pgfsetstrokecolor{currentstroke}%
\pgfsetstrokeopacity{0.000000}%
\pgfsetdash{}{0pt}%
\pgfpathmoveto{\pgfqpoint{3.090744in}{0.985250in}}%
\pgfpathlineto{\pgfqpoint{3.157750in}{0.985250in}}%
\pgfpathlineto{\pgfqpoint{3.157750in}{1.081500in}}%
\pgfpathlineto{\pgfqpoint{3.090744in}{1.081500in}}%
\pgfpathlineto{\pgfqpoint{3.090744in}{0.985250in}}%
\pgfpathclose%
\pgfusepath{fill}%
\end{pgfscope}%
\begin{pgfscope}%
\pgfpathrectangle{\pgfqpoint{2.857045in}{0.350000in}}{\pgfqpoint{0.300705in}{0.962500in}}%
\pgfusepath{clip}%
\pgfsetbuttcap%
\pgfsetmiterjoin%
\definecolor{currentfill}{rgb}{0.376732,0.653072,0.822484}%
\pgfsetfillcolor{currentfill}%
\pgfsetlinewidth{0.000000pt}%
\definecolor{currentstroke}{rgb}{0.000000,0.000000,0.000000}%
\pgfsetstrokecolor{currentstroke}%
\pgfsetstrokeopacity{0.000000}%
\pgfsetdash{}{0pt}%
\pgfpathmoveto{\pgfqpoint{3.090572in}{0.994875in}}%
\pgfpathlineto{\pgfqpoint{3.157750in}{0.994875in}}%
\pgfpathlineto{\pgfqpoint{3.157750in}{1.091125in}}%
\pgfpathlineto{\pgfqpoint{3.090572in}{1.091125in}}%
\pgfpathlineto{\pgfqpoint{3.090572in}{0.994875in}}%
\pgfpathclose%
\pgfusepath{fill}%
\end{pgfscope}%
\begin{pgfscope}%
\pgfpathrectangle{\pgfqpoint{2.857045in}{0.350000in}}{\pgfqpoint{0.300705in}{0.962500in}}%
\pgfusepath{clip}%
\pgfsetbuttcap%
\pgfsetmiterjoin%
\definecolor{currentfill}{rgb}{0.376732,0.653072,0.822484}%
\pgfsetfillcolor{currentfill}%
\pgfsetlinewidth{0.000000pt}%
\definecolor{currentstroke}{rgb}{0.000000,0.000000,0.000000}%
\pgfsetstrokecolor{currentstroke}%
\pgfsetstrokeopacity{0.000000}%
\pgfsetdash{}{0pt}%
\pgfpathmoveto{\pgfqpoint{3.090572in}{1.004500in}}%
\pgfpathlineto{\pgfqpoint{3.157750in}{1.004500in}}%
\pgfpathlineto{\pgfqpoint{3.157750in}{1.100750in}}%
\pgfpathlineto{\pgfqpoint{3.090572in}{1.100750in}}%
\pgfpathlineto{\pgfqpoint{3.090572in}{1.004500in}}%
\pgfpathclose%
\pgfusepath{fill}%
\end{pgfscope}%
\begin{pgfscope}%
\pgfpathrectangle{\pgfqpoint{2.857045in}{0.350000in}}{\pgfqpoint{0.300705in}{0.962500in}}%
\pgfusepath{clip}%
\pgfsetbuttcap%
\pgfsetmiterjoin%
\definecolor{currentfill}{rgb}{0.371688,0.649627,0.820515}%
\pgfsetfillcolor{currentfill}%
\pgfsetlinewidth{0.000000pt}%
\definecolor{currentstroke}{rgb}{0.000000,0.000000,0.000000}%
\pgfsetstrokecolor{currentstroke}%
\pgfsetstrokeopacity{0.000000}%
\pgfsetdash{}{0pt}%
\pgfpathmoveto{\pgfqpoint{3.090370in}{1.014125in}}%
\pgfpathlineto{\pgfqpoint{3.157750in}{1.014125in}}%
\pgfpathlineto{\pgfqpoint{3.157750in}{1.110375in}}%
\pgfpathlineto{\pgfqpoint{3.090370in}{1.110375in}}%
\pgfpathlineto{\pgfqpoint{3.090370in}{1.014125in}}%
\pgfpathclose%
\pgfusepath{fill}%
\end{pgfscope}%
\begin{pgfscope}%
\pgfpathrectangle{\pgfqpoint{2.857045in}{0.350000in}}{\pgfqpoint{0.300705in}{0.962500in}}%
\pgfusepath{clip}%
\pgfsetbuttcap%
\pgfsetmiterjoin%
\definecolor{currentfill}{rgb}{0.371688,0.649627,0.820515}%
\pgfsetfillcolor{currentfill}%
\pgfsetlinewidth{0.000000pt}%
\definecolor{currentstroke}{rgb}{0.000000,0.000000,0.000000}%
\pgfsetstrokecolor{currentstroke}%
\pgfsetstrokeopacity{0.000000}%
\pgfsetdash{}{0pt}%
\pgfpathmoveto{\pgfqpoint{3.090370in}{1.023750in}}%
\pgfpathlineto{\pgfqpoint{3.157750in}{1.023750in}}%
\pgfpathlineto{\pgfqpoint{3.157750in}{1.120000in}}%
\pgfpathlineto{\pgfqpoint{3.090370in}{1.120000in}}%
\pgfpathlineto{\pgfqpoint{3.090370in}{1.023750in}}%
\pgfpathclose%
\pgfusepath{fill}%
\end{pgfscope}%
\begin{pgfscope}%
\pgfpathrectangle{\pgfqpoint{2.857045in}{0.350000in}}{\pgfqpoint{0.300705in}{0.962500in}}%
\pgfusepath{clip}%
\pgfsetbuttcap%
\pgfsetmiterjoin%
\definecolor{currentfill}{rgb}{0.366644,0.646182,0.818547}%
\pgfsetfillcolor{currentfill}%
\pgfsetlinewidth{0.000000pt}%
\definecolor{currentstroke}{rgb}{0.000000,0.000000,0.000000}%
\pgfsetstrokecolor{currentstroke}%
\pgfsetstrokeopacity{0.000000}%
\pgfsetdash{}{0pt}%
\pgfpathmoveto{\pgfqpoint{3.090080in}{1.033375in}}%
\pgfpathlineto{\pgfqpoint{3.157750in}{1.033375in}}%
\pgfpathlineto{\pgfqpoint{3.157750in}{1.129625in}}%
\pgfpathlineto{\pgfqpoint{3.090080in}{1.129625in}}%
\pgfpathlineto{\pgfqpoint{3.090080in}{1.033375in}}%
\pgfpathclose%
\pgfusepath{fill}%
\end{pgfscope}%
\begin{pgfscope}%
\pgfpathrectangle{\pgfqpoint{2.857045in}{0.350000in}}{\pgfqpoint{0.300705in}{0.962500in}}%
\pgfusepath{clip}%
\pgfsetbuttcap%
\pgfsetmiterjoin%
\definecolor{currentfill}{rgb}{0.361599,0.642737,0.816578}%
\pgfsetfillcolor{currentfill}%
\pgfsetlinewidth{0.000000pt}%
\definecolor{currentstroke}{rgb}{0.000000,0.000000,0.000000}%
\pgfsetstrokecolor{currentstroke}%
\pgfsetstrokeopacity{0.000000}%
\pgfsetdash{}{0pt}%
\pgfpathmoveto{\pgfqpoint{3.089945in}{1.043000in}}%
\pgfpathlineto{\pgfqpoint{3.157750in}{1.043000in}}%
\pgfpathlineto{\pgfqpoint{3.157750in}{1.139250in}}%
\pgfpathlineto{\pgfqpoint{3.089945in}{1.139250in}}%
\pgfpathlineto{\pgfqpoint{3.089945in}{1.043000in}}%
\pgfpathclose%
\pgfusepath{fill}%
\end{pgfscope}%
\begin{pgfscope}%
\pgfpathrectangle{\pgfqpoint{2.857045in}{0.350000in}}{\pgfqpoint{0.300705in}{0.962500in}}%
\pgfusepath{clip}%
\pgfsetbuttcap%
\pgfsetmiterjoin%
\definecolor{currentfill}{rgb}{0.361599,0.642737,0.816578}%
\pgfsetfillcolor{currentfill}%
\pgfsetlinewidth{0.000000pt}%
\definecolor{currentstroke}{rgb}{0.000000,0.000000,0.000000}%
\pgfsetstrokecolor{currentstroke}%
\pgfsetstrokeopacity{0.000000}%
\pgfsetdash{}{0pt}%
\pgfpathmoveto{\pgfqpoint{3.089932in}{1.052625in}}%
\pgfpathlineto{\pgfqpoint{3.157750in}{1.052625in}}%
\pgfpathlineto{\pgfqpoint{3.157750in}{1.148875in}}%
\pgfpathlineto{\pgfqpoint{3.089932in}{1.148875in}}%
\pgfpathlineto{\pgfqpoint{3.089932in}{1.052625in}}%
\pgfpathclose%
\pgfusepath{fill}%
\end{pgfscope}%
\begin{pgfscope}%
\pgfpathrectangle{\pgfqpoint{2.857045in}{0.350000in}}{\pgfqpoint{0.300705in}{0.962500in}}%
\pgfusepath{clip}%
\pgfsetbuttcap%
\pgfsetmiterjoin%
\definecolor{currentfill}{rgb}{0.351511,0.635848,0.812641}%
\pgfsetfillcolor{currentfill}%
\pgfsetlinewidth{0.000000pt}%
\definecolor{currentstroke}{rgb}{0.000000,0.000000,0.000000}%
\pgfsetstrokecolor{currentstroke}%
\pgfsetstrokeopacity{0.000000}%
\pgfsetdash{}{0pt}%
\pgfpathmoveto{\pgfqpoint{3.089626in}{1.062250in}}%
\pgfpathlineto{\pgfqpoint{3.157750in}{1.062250in}}%
\pgfpathlineto{\pgfqpoint{3.157750in}{1.158500in}}%
\pgfpathlineto{\pgfqpoint{3.089626in}{1.158500in}}%
\pgfpathlineto{\pgfqpoint{3.089626in}{1.062250in}}%
\pgfpathclose%
\pgfusepath{fill}%
\end{pgfscope}%
\begin{pgfscope}%
\pgfpathrectangle{\pgfqpoint{2.857045in}{0.350000in}}{\pgfqpoint{0.300705in}{0.962500in}}%
\pgfusepath{clip}%
\pgfsetbuttcap%
\pgfsetmiterjoin%
\definecolor{currentfill}{rgb}{0.351511,0.635848,0.812641}%
\pgfsetfillcolor{currentfill}%
\pgfsetlinewidth{0.000000pt}%
\definecolor{currentstroke}{rgb}{0.000000,0.000000,0.000000}%
\pgfsetstrokecolor{currentstroke}%
\pgfsetstrokeopacity{0.000000}%
\pgfsetdash{}{0pt}%
\pgfpathmoveto{\pgfqpoint{3.089486in}{1.071875in}}%
\pgfpathlineto{\pgfqpoint{3.157750in}{1.071875in}}%
\pgfpathlineto{\pgfqpoint{3.157750in}{1.168125in}}%
\pgfpathlineto{\pgfqpoint{3.089486in}{1.168125in}}%
\pgfpathlineto{\pgfqpoint{3.089486in}{1.071875in}}%
\pgfpathclose%
\pgfusepath{fill}%
\end{pgfscope}%
\begin{pgfscope}%
\pgfpathrectangle{\pgfqpoint{2.857045in}{0.350000in}}{\pgfqpoint{0.300705in}{0.962500in}}%
\pgfusepath{clip}%
\pgfsetbuttcap%
\pgfsetmiterjoin%
\definecolor{currentfill}{rgb}{0.346467,0.632403,0.810673}%
\pgfsetfillcolor{currentfill}%
\pgfsetlinewidth{0.000000pt}%
\definecolor{currentstroke}{rgb}{0.000000,0.000000,0.000000}%
\pgfsetstrokecolor{currentstroke}%
\pgfsetstrokeopacity{0.000000}%
\pgfsetdash{}{0pt}%
\pgfpathmoveto{\pgfqpoint{3.089241in}{1.081500in}}%
\pgfpathlineto{\pgfqpoint{3.157750in}{1.081500in}}%
\pgfpathlineto{\pgfqpoint{3.157750in}{1.177750in}}%
\pgfpathlineto{\pgfqpoint{3.089241in}{1.177750in}}%
\pgfpathlineto{\pgfqpoint{3.089241in}{1.081500in}}%
\pgfpathclose%
\pgfusepath{fill}%
\end{pgfscope}%
\begin{pgfscope}%
\pgfpathrectangle{\pgfqpoint{2.857045in}{0.350000in}}{\pgfqpoint{0.300705in}{0.962500in}}%
\pgfusepath{clip}%
\pgfsetbuttcap%
\pgfsetmiterjoin%
\definecolor{currentfill}{rgb}{0.336378,0.625513,0.806736}%
\pgfsetfillcolor{currentfill}%
\pgfsetlinewidth{0.000000pt}%
\definecolor{currentstroke}{rgb}{0.000000,0.000000,0.000000}%
\pgfsetstrokecolor{currentstroke}%
\pgfsetstrokeopacity{0.000000}%
\pgfsetdash{}{0pt}%
\pgfpathmoveto{\pgfqpoint{3.088864in}{1.091125in}}%
\pgfpathlineto{\pgfqpoint{3.157750in}{1.091125in}}%
\pgfpathlineto{\pgfqpoint{3.157750in}{1.187375in}}%
\pgfpathlineto{\pgfqpoint{3.088864in}{1.187375in}}%
\pgfpathlineto{\pgfqpoint{3.088864in}{1.091125in}}%
\pgfpathclose%
\pgfusepath{fill}%
\end{pgfscope}%
\begin{pgfscope}%
\pgfpathrectangle{\pgfqpoint{2.857045in}{0.350000in}}{\pgfqpoint{0.300705in}{0.962500in}}%
\pgfusepath{clip}%
\pgfsetbuttcap%
\pgfsetmiterjoin%
\definecolor{currentfill}{rgb}{0.336378,0.625513,0.806736}%
\pgfsetfillcolor{currentfill}%
\pgfsetlinewidth{0.000000pt}%
\definecolor{currentstroke}{rgb}{0.000000,0.000000,0.000000}%
\pgfsetstrokecolor{currentstroke}%
\pgfsetstrokeopacity{0.000000}%
\pgfsetdash{}{0pt}%
\pgfpathmoveto{\pgfqpoint{3.088766in}{1.100750in}}%
\pgfpathlineto{\pgfqpoint{3.157750in}{1.100750in}}%
\pgfpathlineto{\pgfqpoint{3.157750in}{1.197000in}}%
\pgfpathlineto{\pgfqpoint{3.088766in}{1.197000in}}%
\pgfpathlineto{\pgfqpoint{3.088766in}{1.100750in}}%
\pgfpathclose%
\pgfusepath{fill}%
\end{pgfscope}%
\begin{pgfscope}%
\pgfpathrectangle{\pgfqpoint{2.857045in}{0.350000in}}{\pgfqpoint{0.300705in}{0.962500in}}%
\pgfusepath{clip}%
\pgfsetbuttcap%
\pgfsetmiterjoin%
\definecolor{currentfill}{rgb}{0.326290,0.618624,0.802799}%
\pgfsetfillcolor{currentfill}%
\pgfsetlinewidth{0.000000pt}%
\definecolor{currentstroke}{rgb}{0.000000,0.000000,0.000000}%
\pgfsetstrokecolor{currentstroke}%
\pgfsetstrokeopacity{0.000000}%
\pgfsetdash{}{0pt}%
\pgfpathmoveto{\pgfqpoint{3.088423in}{1.110375in}}%
\pgfpathlineto{\pgfqpoint{3.157750in}{1.110375in}}%
\pgfpathlineto{\pgfqpoint{3.157750in}{1.206625in}}%
\pgfpathlineto{\pgfqpoint{3.088423in}{1.206625in}}%
\pgfpathlineto{\pgfqpoint{3.088423in}{1.110375in}}%
\pgfpathclose%
\pgfusepath{fill}%
\end{pgfscope}%
\begin{pgfscope}%
\pgfpathrectangle{\pgfqpoint{2.857045in}{0.350000in}}{\pgfqpoint{0.300705in}{0.962500in}}%
\pgfusepath{clip}%
\pgfsetbuttcap%
\pgfsetmiterjoin%
\definecolor{currentfill}{rgb}{0.311157,0.608289,0.796894}%
\pgfsetfillcolor{currentfill}%
\pgfsetlinewidth{0.000000pt}%
\definecolor{currentstroke}{rgb}{0.000000,0.000000,0.000000}%
\pgfsetstrokecolor{currentstroke}%
\pgfsetstrokeopacity{0.000000}%
\pgfsetdash{}{0pt}%
\pgfpathmoveto{\pgfqpoint{3.087813in}{1.120000in}}%
\pgfpathlineto{\pgfqpoint{3.157750in}{1.120000in}}%
\pgfpathlineto{\pgfqpoint{3.157750in}{1.216250in}}%
\pgfpathlineto{\pgfqpoint{3.087813in}{1.216250in}}%
\pgfpathlineto{\pgfqpoint{3.087813in}{1.120000in}}%
\pgfpathclose%
\pgfusepath{fill}%
\end{pgfscope}%
\begin{pgfscope}%
\pgfpathrectangle{\pgfqpoint{2.857045in}{0.350000in}}{\pgfqpoint{0.300705in}{0.962500in}}%
\pgfusepath{clip}%
\pgfsetbuttcap%
\pgfsetmiterjoin%
\definecolor{currentfill}{rgb}{0.311157,0.608289,0.796894}%
\pgfsetfillcolor{currentfill}%
\pgfsetlinewidth{0.000000pt}%
\definecolor{currentstroke}{rgb}{0.000000,0.000000,0.000000}%
\pgfsetstrokecolor{currentstroke}%
\pgfsetstrokeopacity{0.000000}%
\pgfsetdash{}{0pt}%
\pgfpathmoveto{\pgfqpoint{3.087661in}{1.129625in}}%
\pgfpathlineto{\pgfqpoint{3.157750in}{1.129625in}}%
\pgfpathlineto{\pgfqpoint{3.157750in}{1.225875in}}%
\pgfpathlineto{\pgfqpoint{3.087661in}{1.225875in}}%
\pgfpathlineto{\pgfqpoint{3.087661in}{1.129625in}}%
\pgfpathclose%
\pgfusepath{fill}%
\end{pgfscope}%
\begin{pgfscope}%
\pgfpathrectangle{\pgfqpoint{2.857045in}{0.350000in}}{\pgfqpoint{0.300705in}{0.962500in}}%
\pgfusepath{clip}%
\pgfsetbuttcap%
\pgfsetmiterjoin%
\definecolor{currentfill}{rgb}{0.311157,0.608289,0.796894}%
\pgfsetfillcolor{currentfill}%
\pgfsetlinewidth{0.000000pt}%
\definecolor{currentstroke}{rgb}{0.000000,0.000000,0.000000}%
\pgfsetstrokecolor{currentstroke}%
\pgfsetstrokeopacity{0.000000}%
\pgfsetdash{}{0pt}%
\pgfpathmoveto{\pgfqpoint{3.087661in}{1.139250in}}%
\pgfpathlineto{\pgfqpoint{3.157750in}{1.139250in}}%
\pgfpathlineto{\pgfqpoint{3.157750in}{1.235500in}}%
\pgfpathlineto{\pgfqpoint{3.087661in}{1.235500in}}%
\pgfpathlineto{\pgfqpoint{3.087661in}{1.139250in}}%
\pgfpathclose%
\pgfusepath{fill}%
\end{pgfscope}%
\begin{pgfscope}%
\pgfpathrectangle{\pgfqpoint{2.857045in}{0.350000in}}{\pgfqpoint{0.300705in}{0.962500in}}%
\pgfusepath{clip}%
\pgfsetbuttcap%
\pgfsetmiterjoin%
\definecolor{currentfill}{rgb}{0.296025,0.597955,0.790988}%
\pgfsetfillcolor{currentfill}%
\pgfsetlinewidth{0.000000pt}%
\definecolor{currentstroke}{rgb}{0.000000,0.000000,0.000000}%
\pgfsetstrokecolor{currentstroke}%
\pgfsetstrokeopacity{0.000000}%
\pgfsetdash{}{0pt}%
\pgfpathmoveto{\pgfqpoint{3.087211in}{1.148875in}}%
\pgfpathlineto{\pgfqpoint{3.157750in}{1.148875in}}%
\pgfpathlineto{\pgfqpoint{3.157750in}{1.245125in}}%
\pgfpathlineto{\pgfqpoint{3.087211in}{1.245125in}}%
\pgfpathlineto{\pgfqpoint{3.087211in}{1.148875in}}%
\pgfpathclose%
\pgfusepath{fill}%
\end{pgfscope}%
\begin{pgfscope}%
\pgfpathrectangle{\pgfqpoint{2.857045in}{0.350000in}}{\pgfqpoint{0.300705in}{0.962500in}}%
\pgfusepath{clip}%
\pgfsetbuttcap%
\pgfsetmiterjoin%
\definecolor{currentfill}{rgb}{0.290980,0.594510,0.789020}%
\pgfsetfillcolor{currentfill}%
\pgfsetlinewidth{0.000000pt}%
\definecolor{currentstroke}{rgb}{0.000000,0.000000,0.000000}%
\pgfsetstrokecolor{currentstroke}%
\pgfsetstrokeopacity{0.000000}%
\pgfsetdash{}{0pt}%
\pgfpathmoveto{\pgfqpoint{3.086798in}{1.158500in}}%
\pgfpathlineto{\pgfqpoint{3.157750in}{1.158500in}}%
\pgfpathlineto{\pgfqpoint{3.157750in}{1.254750in}}%
\pgfpathlineto{\pgfqpoint{3.086798in}{1.254750in}}%
\pgfpathlineto{\pgfqpoint{3.086798in}{1.158500in}}%
\pgfpathclose%
\pgfusepath{fill}%
\end{pgfscope}%
\begin{pgfscope}%
\pgfpathrectangle{\pgfqpoint{2.857045in}{0.350000in}}{\pgfqpoint{0.300705in}{0.962500in}}%
\pgfusepath{clip}%
\pgfsetbuttcap%
\pgfsetmiterjoin%
\definecolor{currentfill}{rgb}{0.270804,0.580730,0.781146}%
\pgfsetfillcolor{currentfill}%
\pgfsetlinewidth{0.000000pt}%
\definecolor{currentstroke}{rgb}{0.000000,0.000000,0.000000}%
\pgfsetstrokecolor{currentstroke}%
\pgfsetstrokeopacity{0.000000}%
\pgfsetdash{}{0pt}%
\pgfpathmoveto{\pgfqpoint{3.085933in}{1.168125in}}%
\pgfpathlineto{\pgfqpoint{3.157750in}{1.168125in}}%
\pgfpathlineto{\pgfqpoint{3.157750in}{1.264375in}}%
\pgfpathlineto{\pgfqpoint{3.085933in}{1.264375in}}%
\pgfpathlineto{\pgfqpoint{3.085933in}{1.168125in}}%
\pgfpathclose%
\pgfusepath{fill}%
\end{pgfscope}%
\begin{pgfscope}%
\pgfpathrectangle{\pgfqpoint{2.857045in}{0.350000in}}{\pgfqpoint{0.300705in}{0.962500in}}%
\pgfusepath{clip}%
\pgfsetbuttcap%
\pgfsetmiterjoin%
\definecolor{currentfill}{rgb}{0.256286,0.570012,0.775163}%
\pgfsetfillcolor{currentfill}%
\pgfsetlinewidth{0.000000pt}%
\definecolor{currentstroke}{rgb}{0.000000,0.000000,0.000000}%
\pgfsetstrokecolor{currentstroke}%
\pgfsetstrokeopacity{0.000000}%
\pgfsetdash{}{0pt}%
\pgfpathmoveto{\pgfqpoint{3.085385in}{1.177750in}}%
\pgfpathlineto{\pgfqpoint{3.157750in}{1.177750in}}%
\pgfpathlineto{\pgfqpoint{3.157750in}{1.274000in}}%
\pgfpathlineto{\pgfqpoint{3.085385in}{1.274000in}}%
\pgfpathlineto{\pgfqpoint{3.085385in}{1.177750in}}%
\pgfpathclose%
\pgfusepath{fill}%
\end{pgfscope}%
\begin{pgfscope}%
\pgfpathrectangle{\pgfqpoint{2.857045in}{0.350000in}}{\pgfqpoint{0.300705in}{0.962500in}}%
\pgfusepath{clip}%
\pgfsetbuttcap%
\pgfsetmiterjoin%
\definecolor{currentfill}{rgb}{0.256286,0.570012,0.775163}%
\pgfsetfillcolor{currentfill}%
\pgfsetlinewidth{0.000000pt}%
\definecolor{currentstroke}{rgb}{0.000000,0.000000,0.000000}%
\pgfsetstrokecolor{currentstroke}%
\pgfsetstrokeopacity{0.000000}%
\pgfsetdash{}{0pt}%
\pgfpathmoveto{\pgfqpoint{3.085357in}{1.187375in}}%
\pgfpathlineto{\pgfqpoint{3.157750in}{1.187375in}}%
\pgfpathlineto{\pgfqpoint{3.157750in}{1.283625in}}%
\pgfpathlineto{\pgfqpoint{3.085357in}{1.283625in}}%
\pgfpathlineto{\pgfqpoint{3.085357in}{1.187375in}}%
\pgfpathclose%
\pgfusepath{fill}%
\end{pgfscope}%
\begin{pgfscope}%
\pgfpathrectangle{\pgfqpoint{2.857045in}{0.350000in}}{\pgfqpoint{0.300705in}{0.962500in}}%
\pgfusepath{clip}%
\pgfsetbuttcap%
\pgfsetmiterjoin%
\definecolor{currentfill}{rgb}{0.244106,0.557832,0.768889}%
\pgfsetfillcolor{currentfill}%
\pgfsetlinewidth{0.000000pt}%
\definecolor{currentstroke}{rgb}{0.000000,0.000000,0.000000}%
\pgfsetstrokecolor{currentstroke}%
\pgfsetstrokeopacity{0.000000}%
\pgfsetdash{}{0pt}%
\pgfpathmoveto{\pgfqpoint{3.084778in}{1.197000in}}%
\pgfpathlineto{\pgfqpoint{3.157750in}{1.197000in}}%
\pgfpathlineto{\pgfqpoint{3.157750in}{1.293250in}}%
\pgfpathlineto{\pgfqpoint{3.084778in}{1.293250in}}%
\pgfpathlineto{\pgfqpoint{3.084778in}{1.197000in}}%
\pgfpathclose%
\pgfusepath{fill}%
\end{pgfscope}%
\begin{pgfscope}%
\pgfpathrectangle{\pgfqpoint{2.857045in}{0.350000in}}{\pgfqpoint{0.300705in}{0.962500in}}%
\pgfusepath{clip}%
\pgfsetbuttcap%
\pgfsetmiterjoin%
\definecolor{currentfill}{rgb}{0.223806,0.537532,0.758431}%
\pgfsetfillcolor{currentfill}%
\pgfsetlinewidth{0.000000pt}%
\definecolor{currentstroke}{rgb}{0.000000,0.000000,0.000000}%
\pgfsetstrokecolor{currentstroke}%
\pgfsetstrokeopacity{0.000000}%
\pgfsetdash{}{0pt}%
\pgfpathmoveto{\pgfqpoint{3.083688in}{1.206625in}}%
\pgfpathlineto{\pgfqpoint{3.157750in}{1.206625in}}%
\pgfpathlineto{\pgfqpoint{3.157750in}{1.302875in}}%
\pgfpathlineto{\pgfqpoint{3.083688in}{1.302875in}}%
\pgfpathlineto{\pgfqpoint{3.083688in}{1.206625in}}%
\pgfpathclose%
\pgfusepath{fill}%
\end{pgfscope}%
\begin{pgfscope}%
\pgfpathrectangle{\pgfqpoint{2.857045in}{0.350000in}}{\pgfqpoint{0.300705in}{0.962500in}}%
\pgfusepath{clip}%
\pgfsetbuttcap%
\pgfsetmiterjoin%
\definecolor{currentfill}{rgb}{0.223806,0.537532,0.758431}%
\pgfsetfillcolor{currentfill}%
\pgfsetlinewidth{0.000000pt}%
\definecolor{currentstroke}{rgb}{0.000000,0.000000,0.000000}%
\pgfsetstrokecolor{currentstroke}%
\pgfsetstrokeopacity{0.000000}%
\pgfsetdash{}{0pt}%
\pgfpathmoveto{\pgfqpoint{3.083619in}{1.216250in}}%
\pgfpathlineto{\pgfqpoint{3.157750in}{1.216250in}}%
\pgfpathlineto{\pgfqpoint{3.157750in}{1.312500in}}%
\pgfpathlineto{\pgfqpoint{3.083619in}{1.312500in}}%
\pgfpathlineto{\pgfqpoint{3.083619in}{1.216250in}}%
\pgfpathclose%
\pgfusepath{fill}%
\end{pgfscope}%
\begin{pgfscope}%
\pgfpathrectangle{\pgfqpoint{2.857045in}{0.350000in}}{\pgfqpoint{0.300705in}{0.962500in}}%
\pgfusepath{clip}%
\pgfsetbuttcap%
\pgfsetmiterjoin%
\definecolor{currentfill}{rgb}{0.215686,0.529412,0.754248}%
\pgfsetfillcolor{currentfill}%
\pgfsetlinewidth{0.000000pt}%
\definecolor{currentstroke}{rgb}{0.000000,0.000000,0.000000}%
\pgfsetstrokecolor{currentstroke}%
\pgfsetstrokeopacity{0.000000}%
\pgfsetdash{}{0pt}%
\pgfpathmoveto{\pgfqpoint{3.083105in}{1.225875in}}%
\pgfpathlineto{\pgfqpoint{3.157750in}{1.225875in}}%
\pgfpathlineto{\pgfqpoint{3.157750in}{1.322125in}}%
\pgfpathlineto{\pgfqpoint{3.083105in}{1.322125in}}%
\pgfpathlineto{\pgfqpoint{3.083105in}{1.225875in}}%
\pgfpathclose%
\pgfusepath{fill}%
\end{pgfscope}%
\begin{pgfscope}%
\pgfpathrectangle{\pgfqpoint{2.857045in}{0.350000in}}{\pgfqpoint{0.300705in}{0.962500in}}%
\pgfusepath{clip}%
\pgfsetbuttcap%
\pgfsetmiterjoin%
\definecolor{currentfill}{rgb}{0.207566,0.521292,0.750065}%
\pgfsetfillcolor{currentfill}%
\pgfsetlinewidth{0.000000pt}%
\definecolor{currentstroke}{rgb}{0.000000,0.000000,0.000000}%
\pgfsetstrokecolor{currentstroke}%
\pgfsetstrokeopacity{0.000000}%
\pgfsetdash{}{0pt}%
\pgfpathmoveto{\pgfqpoint{3.082791in}{1.235500in}}%
\pgfpathlineto{\pgfqpoint{3.157750in}{1.235500in}}%
\pgfpathlineto{\pgfqpoint{3.157750in}{1.331750in}}%
\pgfpathlineto{\pgfqpoint{3.082791in}{1.331750in}}%
\pgfpathlineto{\pgfqpoint{3.082791in}{1.235500in}}%
\pgfpathclose%
\pgfusepath{fill}%
\end{pgfscope}%
\begin{pgfscope}%
\pgfpathrectangle{\pgfqpoint{2.857045in}{0.350000in}}{\pgfqpoint{0.300705in}{0.962500in}}%
\pgfusepath{clip}%
\pgfsetbuttcap%
\pgfsetmiterjoin%
\definecolor{currentfill}{rgb}{0.187266,0.500992,0.739608}%
\pgfsetfillcolor{currentfill}%
\pgfsetlinewidth{0.000000pt}%
\definecolor{currentstroke}{rgb}{0.000000,0.000000,0.000000}%
\pgfsetstrokecolor{currentstroke}%
\pgfsetstrokeopacity{0.000000}%
\pgfsetdash{}{0pt}%
\pgfpathmoveto{\pgfqpoint{3.081559in}{1.245125in}}%
\pgfpathlineto{\pgfqpoint{3.157750in}{1.245125in}}%
\pgfpathlineto{\pgfqpoint{3.157750in}{1.341375in}}%
\pgfpathlineto{\pgfqpoint{3.081559in}{1.341375in}}%
\pgfpathlineto{\pgfqpoint{3.081559in}{1.245125in}}%
\pgfpathclose%
\pgfusepath{fill}%
\end{pgfscope}%
\begin{pgfscope}%
\pgfpathrectangle{\pgfqpoint{2.857045in}{0.350000in}}{\pgfqpoint{0.300705in}{0.962500in}}%
\pgfusepath{clip}%
\pgfsetbuttcap%
\pgfsetmiterjoin%
\definecolor{currentfill}{rgb}{0.166967,0.480692,0.729150}%
\pgfsetfillcolor{currentfill}%
\pgfsetlinewidth{0.000000pt}%
\definecolor{currentstroke}{rgb}{0.000000,0.000000,0.000000}%
\pgfsetstrokecolor{currentstroke}%
\pgfsetstrokeopacity{0.000000}%
\pgfsetdash{}{0pt}%
\pgfpathmoveto{\pgfqpoint{3.080544in}{1.254750in}}%
\pgfpathlineto{\pgfqpoint{3.157750in}{1.254750in}}%
\pgfpathlineto{\pgfqpoint{3.157750in}{1.351000in}}%
\pgfpathlineto{\pgfqpoint{3.080544in}{1.351000in}}%
\pgfpathlineto{\pgfqpoint{3.080544in}{1.254750in}}%
\pgfpathclose%
\pgfusepath{fill}%
\end{pgfscope}%
\begin{pgfscope}%
\pgfpathrectangle{\pgfqpoint{2.857045in}{0.350000in}}{\pgfqpoint{0.300705in}{0.962500in}}%
\pgfusepath{clip}%
\pgfsetbuttcap%
\pgfsetmiterjoin%
\definecolor{currentfill}{rgb}{0.166967,0.480692,0.729150}%
\pgfsetfillcolor{currentfill}%
\pgfsetlinewidth{0.000000pt}%
\definecolor{currentstroke}{rgb}{0.000000,0.000000,0.000000}%
\pgfsetstrokecolor{currentstroke}%
\pgfsetstrokeopacity{0.000000}%
\pgfsetdash{}{0pt}%
\pgfpathmoveto{\pgfqpoint{3.080544in}{1.264375in}}%
\pgfpathlineto{\pgfqpoint{3.157750in}{1.264375in}}%
\pgfpathlineto{\pgfqpoint{3.157750in}{1.360625in}}%
\pgfpathlineto{\pgfqpoint{3.080544in}{1.360625in}}%
\pgfpathlineto{\pgfqpoint{3.080544in}{1.264375in}}%
\pgfpathclose%
\pgfusepath{fill}%
\end{pgfscope}%
\begin{pgfscope}%
\pgfpathrectangle{\pgfqpoint{2.857045in}{0.350000in}}{\pgfqpoint{0.300705in}{0.962500in}}%
\pgfusepath{clip}%
\pgfsetbuttcap%
\pgfsetmiterjoin%
\definecolor{currentfill}{rgb}{0.130427,0.444152,0.710327}%
\pgfsetfillcolor{currentfill}%
\pgfsetlinewidth{0.000000pt}%
\definecolor{currentstroke}{rgb}{0.000000,0.000000,0.000000}%
\pgfsetstrokecolor{currentstroke}%
\pgfsetstrokeopacity{0.000000}%
\pgfsetdash{}{0pt}%
\pgfpathmoveto{\pgfqpoint{3.078510in}{1.274000in}}%
\pgfpathlineto{\pgfqpoint{3.157750in}{1.274000in}}%
\pgfpathlineto{\pgfqpoint{3.157750in}{1.370250in}}%
\pgfpathlineto{\pgfqpoint{3.078510in}{1.370250in}}%
\pgfpathlineto{\pgfqpoint{3.078510in}{1.274000in}}%
\pgfpathclose%
\pgfusepath{fill}%
\end{pgfscope}%
\begin{pgfscope}%
\pgfpathrectangle{\pgfqpoint{2.857045in}{0.350000in}}{\pgfqpoint{0.300705in}{0.962500in}}%
\pgfusepath{clip}%
\pgfsetbuttcap%
\pgfsetmiterjoin%
\definecolor{currentfill}{rgb}{0.087120,0.389004,0.667512}%
\pgfsetfillcolor{currentfill}%
\pgfsetlinewidth{0.000000pt}%
\definecolor{currentstroke}{rgb}{0.000000,0.000000,0.000000}%
\pgfsetstrokecolor{currentstroke}%
\pgfsetstrokeopacity{0.000000}%
\pgfsetdash{}{0pt}%
\pgfpathmoveto{\pgfqpoint{3.075409in}{1.283625in}}%
\pgfpathlineto{\pgfqpoint{3.157750in}{1.283625in}}%
\pgfpathlineto{\pgfqpoint{3.157750in}{1.379875in}}%
\pgfpathlineto{\pgfqpoint{3.075409in}{1.379875in}}%
\pgfpathlineto{\pgfqpoint{3.075409in}{1.283625in}}%
\pgfpathclose%
\pgfusepath{fill}%
\end{pgfscope}%
\begin{pgfscope}%
\pgfpathrectangle{\pgfqpoint{2.857045in}{0.350000in}}{\pgfqpoint{0.300705in}{0.962500in}}%
\pgfusepath{clip}%
\pgfsetbuttcap%
\pgfsetmiterjoin%
\definecolor{currentfill}{rgb}{0.040984,0.329950,0.621376}%
\pgfsetfillcolor{currentfill}%
\pgfsetlinewidth{0.000000pt}%
\definecolor{currentstroke}{rgb}{0.000000,0.000000,0.000000}%
\pgfsetstrokecolor{currentstroke}%
\pgfsetstrokeopacity{0.000000}%
\pgfsetdash{}{0pt}%
\pgfpathmoveto{\pgfqpoint{3.072243in}{1.293250in}}%
\pgfpathlineto{\pgfqpoint{3.157750in}{1.293250in}}%
\pgfpathlineto{\pgfqpoint{3.157750in}{1.389500in}}%
\pgfpathlineto{\pgfqpoint{3.072243in}{1.389500in}}%
\pgfpathlineto{\pgfqpoint{3.072243in}{1.293250in}}%
\pgfpathclose%
\pgfusepath{fill}%
\end{pgfscope}%
\begin{pgfscope}%
\pgfpathrectangle{\pgfqpoint{2.857045in}{0.350000in}}{\pgfqpoint{0.300705in}{0.962500in}}%
\pgfusepath{clip}%
\pgfsetbuttcap%
\pgfsetmiterjoin%
\definecolor{currentfill}{rgb}{0.031373,0.188235,0.419608}%
\pgfsetfillcolor{currentfill}%
\pgfsetlinewidth{0.000000pt}%
\definecolor{currentstroke}{rgb}{0.000000,0.000000,0.000000}%
\pgfsetstrokecolor{currentstroke}%
\pgfsetstrokeopacity{0.000000}%
\pgfsetdash{}{0pt}%
\pgfpathmoveto{\pgfqpoint{3.064343in}{1.302875in}}%
\pgfpathlineto{\pgfqpoint{3.157750in}{1.302875in}}%
\pgfpathlineto{\pgfqpoint{3.157750in}{1.399125in}}%
\pgfpathlineto{\pgfqpoint{3.064343in}{1.399125in}}%
\pgfpathlineto{\pgfqpoint{3.064343in}{1.302875in}}%
\pgfpathclose%
\pgfusepath{fill}%
\end{pgfscope}%
\begin{pgfscope}%
\definecolor{textcolor}{rgb}{0.000000,0.000000,0.000000}%
\pgfsetstrokecolor{textcolor}%
\pgfsetfillcolor{textcolor}%
\pgftext[x=2.977327in,y=1.360625in,left,base]{\color{textcolor}\setmainfont{Lato}\rmfamily\fontsize{7.000000}{8.400000}\selectfont 31.1\%}%
\end{pgfscope}%
\begin{pgfscope}%
\definecolor{textcolor}{rgb}{0.000000,0.000000,0.000000}%
\pgfsetstrokecolor{textcolor}%
\pgfsetfillcolor{textcolor}%
\pgftext[x=2.977327in,y=0.234500in,left,base]{\color{textcolor}\setmainfont{Lato}\rmfamily\fontsize{7.000000}{8.400000}\selectfont 12.3\%}%
\end{pgfscope}%
\end{pgfpicture}%
\makeatother%
\endgroup%
 \hspace{-3mm} %% Creator: Matplotlib, PGF backend
%%
%% To include the figure in your LaTeX document, write
%%   \input{<filename>.pgf}
%%
%% Make sure the required packages are loaded in your preamble
%%   \usepackage{pgf}
%%
%% Also ensure that all the required font packages are loaded; for instance,
%% the lmodern package is sometimes necessary when using math font.
%%   \usepackage{lmodern}
%%
%% Figures using additional raster images can only be included by \input if
%% they are in the same directory as the main LaTeX file. For loading figures
%% from other directories you can use the `import` package
%%   \usepackage{import}
%%
%% and then include the figures with
%%   \import{<path to file>}{<filename>.pgf}
%%
%% Matplotlib used the following preamble
%%   
%%   \usepackage{fontspec}
%%   \setmainfont{DejaVuSerif.ttf}[Path=\detokenize{/home/brian/miniconda3/lib/python3.8/site-packages/matplotlib/mpl-data/fonts/ttf/}]
%%   \setsansfont{DejaVuSans.ttf}[Path=\detokenize{/home/brian/miniconda3/lib/python3.8/site-packages/matplotlib/mpl-data/fonts/ttf/}]
%%   \setmonofont{DejaVuSansMono.ttf}[Path=\detokenize{/home/brian/miniconda3/lib/python3.8/site-packages/matplotlib/mpl-data/fonts/ttf/}]
%%   \makeatletter\@ifpackageloaded{underscore}{}{\usepackage[strings]{underscore}}\makeatother
%%
\begingroup%
\makeatletter%
\begin{pgfpicture}%
\pgfpathrectangle{\pgfpointorigin}{\pgfqpoint{3.347362in}{2.125000in}}%
\pgfusepath{use as bounding box, clip}%
\begin{pgfscope}%
\pgfsetbuttcap%
\pgfsetmiterjoin%
\pgfsetlinewidth{0.000000pt}%
\definecolor{currentstroke}{rgb}{0.000000,0.000000,0.000000}%
\pgfsetstrokecolor{currentstroke}%
\pgfsetstrokeopacity{0.000000}%
\pgfsetdash{}{0pt}%
\pgfpathmoveto{\pgfqpoint{0.000000in}{0.000000in}}%
\pgfpathlineto{\pgfqpoint{3.347362in}{0.000000in}}%
\pgfpathlineto{\pgfqpoint{3.347362in}{2.125000in}}%
\pgfpathlineto{\pgfqpoint{0.000000in}{2.125000in}}%
\pgfpathlineto{\pgfqpoint{0.000000in}{0.000000in}}%
\pgfpathclose%
\pgfusepath{}%
\end{pgfscope}%
\begin{pgfscope}%
\pgfpathrectangle{\pgfqpoint{0.100000in}{0.100000in}}{\pgfqpoint{3.007045in}{1.925000in}}%
\pgfusepath{clip}%
\pgfsetbuttcap%
\pgfsetmiterjoin%
\pgfsetlinewidth{0.000000pt}%
\definecolor{currentstroke}{rgb}{0.000000,0.000000,0.000000}%
\pgfsetstrokecolor{currentstroke}%
\pgfsetstrokeopacity{0.000000}%
\pgfsetdash{}{0pt}%
\pgfpathmoveto{\pgfqpoint{0.100000in}{0.100000in}}%
\pgfpathlineto{\pgfqpoint{3.107045in}{0.100000in}}%
\pgfpathlineto{\pgfqpoint{3.107045in}{2.025000in}}%
\pgfpathlineto{\pgfqpoint{0.100000in}{2.025000in}}%
\pgfpathlineto{\pgfqpoint{0.100000in}{0.100000in}}%
\pgfpathclose%
\pgfusepath{}%
\end{pgfscope}%
\begin{pgfscope}%
\pgfpathrectangle{\pgfqpoint{0.100000in}{0.100000in}}{\pgfqpoint{3.007045in}{1.925000in}}%
\pgfusepath{clip}%
\pgfsetbuttcap%
\pgfsetmiterjoin%
\definecolor{currentfill}{rgb}{0.513033,0.795755,0.510542}%
\pgfsetfillcolor{currentfill}%
\pgfsetlinewidth{0.000000pt}%
\definecolor{currentstroke}{rgb}{0.000000,0.000000,0.000000}%
\pgfsetstrokecolor{currentstroke}%
\pgfsetstrokeopacity{0.000000}%
\pgfsetdash{}{0pt}%
\pgfpathmoveto{\pgfqpoint{1.787185in}{1.275241in}}%
\pgfpathlineto{\pgfqpoint{1.764595in}{1.275306in}}%
\pgfpathlineto{\pgfqpoint{1.764683in}{1.298101in}}%
\pgfpathlineto{\pgfqpoint{1.747728in}{1.298216in}}%
\pgfpathlineto{\pgfqpoint{1.744815in}{1.304579in}}%
\pgfpathlineto{\pgfqpoint{1.744960in}{1.321802in}}%
\pgfpathlineto{\pgfqpoint{1.784860in}{1.321668in}}%
\pgfpathlineto{\pgfqpoint{1.784942in}{1.304507in}}%
\pgfpathlineto{\pgfqpoint{1.787084in}{1.304515in}}%
\pgfpathlineto{\pgfqpoint{1.787185in}{1.275241in}}%
\pgfpathclose%
\pgfusepath{fill}%
\end{pgfscope}%
\begin{pgfscope}%
\pgfpathrectangle{\pgfqpoint{0.100000in}{0.100000in}}{\pgfqpoint{3.007045in}{1.925000in}}%
\pgfusepath{clip}%
\pgfsetbuttcap%
\pgfsetmiterjoin%
\definecolor{currentfill}{rgb}{0.585006,0.829343,0.569719}%
\pgfsetfillcolor{currentfill}%
\pgfsetlinewidth{0.000000pt}%
\definecolor{currentstroke}{rgb}{0.000000,0.000000,0.000000}%
\pgfsetstrokecolor{currentstroke}%
\pgfsetstrokeopacity{0.000000}%
\pgfsetdash{}{0pt}%
\pgfpathmoveto{\pgfqpoint{2.287383in}{1.019436in}}%
\pgfpathlineto{\pgfqpoint{2.292665in}{1.018214in}}%
\pgfpathlineto{\pgfqpoint{2.306270in}{1.030998in}}%
\pgfpathlineto{\pgfqpoint{2.302206in}{1.033847in}}%
\pgfpathlineto{\pgfqpoint{2.301947in}{1.038398in}}%
\pgfpathlineto{\pgfqpoint{2.310463in}{1.047437in}}%
\pgfpathlineto{\pgfqpoint{2.318137in}{1.056117in}}%
\pgfpathlineto{\pgfqpoint{2.321129in}{1.054650in}}%
\pgfpathlineto{\pgfqpoint{2.333634in}{1.044388in}}%
\pgfpathlineto{\pgfqpoint{2.332415in}{1.038665in}}%
\pgfpathlineto{\pgfqpoint{2.334697in}{1.028538in}}%
\pgfpathlineto{\pgfqpoint{2.332746in}{1.021710in}}%
\pgfpathlineto{\pgfqpoint{2.336554in}{1.012826in}}%
\pgfpathlineto{\pgfqpoint{2.341210in}{1.008147in}}%
\pgfpathlineto{\pgfqpoint{2.340163in}{1.003243in}}%
\pgfpathlineto{\pgfqpoint{2.336026in}{1.001374in}}%
\pgfpathlineto{\pgfqpoint{2.338194in}{0.987373in}}%
\pgfpathlineto{\pgfqpoint{2.334580in}{0.979199in}}%
\pgfpathlineto{\pgfqpoint{2.325950in}{0.982785in}}%
\pgfpathlineto{\pgfqpoint{2.317969in}{0.990753in}}%
\pgfpathlineto{\pgfqpoint{2.319984in}{0.992600in}}%
\pgfpathlineto{\pgfqpoint{2.315188in}{1.000641in}}%
\pgfpathlineto{\pgfqpoint{2.311789in}{1.005478in}}%
\pgfpathlineto{\pgfqpoint{2.285799in}{1.004151in}}%
\pgfpathlineto{\pgfqpoint{2.287383in}{1.019436in}}%
\pgfpathclose%
\pgfusepath{fill}%
\end{pgfscope}%
\begin{pgfscope}%
\pgfpathrectangle{\pgfqpoint{0.100000in}{0.100000in}}{\pgfqpoint{3.007045in}{1.925000in}}%
\pgfusepath{clip}%
\pgfsetbuttcap%
\pgfsetmiterjoin%
\definecolor{currentfill}{rgb}{0.790542,0.917785,0.763768}%
\pgfsetfillcolor{currentfill}%
\pgfsetlinewidth{0.000000pt}%
\definecolor{currentstroke}{rgb}{0.000000,0.000000,0.000000}%
\pgfsetstrokecolor{currentstroke}%
\pgfsetstrokeopacity{0.000000}%
\pgfsetdash{}{0pt}%
\pgfpathmoveto{\pgfqpoint{1.471341in}{0.833174in}}%
\pgfpathlineto{\pgfqpoint{1.469347in}{0.801604in}}%
\pgfpathlineto{\pgfqpoint{1.440696in}{0.803294in}}%
\pgfpathlineto{\pgfqpoint{1.442781in}{0.834883in}}%
\pgfpathlineto{\pgfqpoint{1.471341in}{0.833174in}}%
\pgfpathclose%
\pgfusepath{fill}%
\end{pgfscope}%
\begin{pgfscope}%
\pgfpathrectangle{\pgfqpoint{0.100000in}{0.100000in}}{\pgfqpoint{3.007045in}{1.925000in}}%
\pgfusepath{clip}%
\pgfsetbuttcap%
\pgfsetmiterjoin%
\definecolor{currentfill}{rgb}{0.557324,0.816424,0.546959}%
\pgfsetfillcolor{currentfill}%
\pgfsetlinewidth{0.000000pt}%
\definecolor{currentstroke}{rgb}{0.000000,0.000000,0.000000}%
\pgfsetstrokecolor{currentstroke}%
\pgfsetstrokeopacity{0.000000}%
\pgfsetdash{}{0pt}%
\pgfpathmoveto{\pgfqpoint{1.408374in}{1.098989in}}%
\pgfpathlineto{\pgfqpoint{1.408351in}{1.098572in}}%
\pgfpathlineto{\pgfqpoint{1.405795in}{1.064151in}}%
\pgfpathlineto{\pgfqpoint{1.405343in}{1.057957in}}%
\pgfpathlineto{\pgfqpoint{1.351564in}{1.062383in}}%
\pgfpathlineto{\pgfqpoint{1.300552in}{1.067248in}}%
\pgfpathlineto{\pgfqpoint{1.301086in}{1.073162in}}%
\pgfpathlineto{\pgfqpoint{1.271819in}{1.087203in}}%
\pgfpathlineto{\pgfqpoint{1.252696in}{1.088003in}}%
\pgfpathlineto{\pgfqpoint{1.251074in}{1.090450in}}%
\pgfpathlineto{\pgfqpoint{1.253600in}{1.112871in}}%
\pgfpathlineto{\pgfqpoint{1.259199in}{1.112273in}}%
\pgfpathlineto{\pgfqpoint{1.261017in}{1.129389in}}%
\pgfpathlineto{\pgfqpoint{1.306529in}{1.124802in}}%
\pgfpathlineto{\pgfqpoint{1.334664in}{1.121680in}}%
\pgfpathlineto{\pgfqpoint{1.333267in}{1.105178in}}%
\pgfpathlineto{\pgfqpoint{1.372398in}{1.101802in}}%
\pgfpathlineto{\pgfqpoint{1.408374in}{1.098989in}}%
\pgfpathclose%
\pgfusepath{fill}%
\end{pgfscope}%
\begin{pgfscope}%
\pgfpathrectangle{\pgfqpoint{0.100000in}{0.100000in}}{\pgfqpoint{3.007045in}{1.925000in}}%
\pgfusepath{clip}%
\pgfsetbuttcap%
\pgfsetmiterjoin%
\definecolor{currentfill}{rgb}{0.648320,0.858116,0.624344}%
\pgfsetfillcolor{currentfill}%
\pgfsetlinewidth{0.000000pt}%
\definecolor{currentstroke}{rgb}{0.000000,0.000000,0.000000}%
\pgfsetstrokecolor{currentstroke}%
\pgfsetstrokeopacity{0.000000}%
\pgfsetdash{}{0pt}%
\pgfpathmoveto{\pgfqpoint{1.769446in}{1.131864in}}%
\pgfpathlineto{\pgfqpoint{1.764032in}{1.131806in}}%
\pgfpathlineto{\pgfqpoint{1.763130in}{1.141541in}}%
\pgfpathlineto{\pgfqpoint{1.763181in}{1.160646in}}%
\pgfpathlineto{\pgfqpoint{1.743298in}{1.160738in}}%
\pgfpathlineto{\pgfqpoint{1.743716in}{1.176139in}}%
\pgfpathlineto{\pgfqpoint{1.755095in}{1.176082in}}%
\pgfpathlineto{\pgfqpoint{1.766486in}{1.176026in}}%
\pgfpathlineto{\pgfqpoint{1.769538in}{1.174111in}}%
\pgfpathlineto{\pgfqpoint{1.767101in}{1.168173in}}%
\pgfpathlineto{\pgfqpoint{1.772963in}{1.162326in}}%
\pgfpathlineto{\pgfqpoint{1.780164in}{1.150731in}}%
\pgfpathlineto{\pgfqpoint{1.777411in}{1.146354in}}%
\pgfpathlineto{\pgfqpoint{1.776978in}{1.132450in}}%
\pgfpathlineto{\pgfqpoint{1.769446in}{1.131864in}}%
\pgfpathclose%
\pgfusepath{fill}%
\end{pgfscope}%
\begin{pgfscope}%
\pgfpathrectangle{\pgfqpoint{0.100000in}{0.100000in}}{\pgfqpoint{3.007045in}{1.925000in}}%
\pgfusepath{clip}%
\pgfsetbuttcap%
\pgfsetmiterjoin%
\definecolor{currentfill}{rgb}{0.685721,0.873864,0.660761}%
\pgfsetfillcolor{currentfill}%
\pgfsetlinewidth{0.000000pt}%
\definecolor{currentstroke}{rgb}{0.000000,0.000000,0.000000}%
\pgfsetstrokecolor{currentstroke}%
\pgfsetstrokeopacity{0.000000}%
\pgfsetdash{}{0pt}%
\pgfpathmoveto{\pgfqpoint{2.317320in}{0.826818in}}%
\pgfpathlineto{\pgfqpoint{2.298990in}{0.824653in}}%
\pgfpathlineto{\pgfqpoint{2.297886in}{0.828687in}}%
\pgfpathlineto{\pgfqpoint{2.290820in}{0.827401in}}%
\pgfpathlineto{\pgfqpoint{2.279304in}{0.827039in}}%
\pgfpathlineto{\pgfqpoint{2.271751in}{0.844592in}}%
\pgfpathlineto{\pgfqpoint{2.275437in}{0.845959in}}%
\pgfpathlineto{\pgfqpoint{2.282179in}{0.853795in}}%
\pgfpathlineto{\pgfqpoint{2.284621in}{0.861875in}}%
\pgfpathlineto{\pgfqpoint{2.287764in}{0.864843in}}%
\pgfpathlineto{\pgfqpoint{2.286587in}{0.869005in}}%
\pgfpathlineto{\pgfqpoint{2.311838in}{0.871613in}}%
\pgfpathlineto{\pgfqpoint{2.311402in}{0.880390in}}%
\pgfpathlineto{\pgfqpoint{2.305031in}{0.879714in}}%
\pgfpathlineto{\pgfqpoint{2.306322in}{0.886459in}}%
\pgfpathlineto{\pgfqpoint{2.310101in}{0.886859in}}%
\pgfpathlineto{\pgfqpoint{2.313442in}{0.898402in}}%
\pgfpathlineto{\pgfqpoint{2.332547in}{0.900599in}}%
\pgfpathlineto{\pgfqpoint{2.333466in}{0.891984in}}%
\pgfpathlineto{\pgfqpoint{2.335741in}{0.890096in}}%
\pgfpathlineto{\pgfqpoint{2.332579in}{0.883358in}}%
\pgfpathlineto{\pgfqpoint{2.334077in}{0.871563in}}%
\pgfpathlineto{\pgfqpoint{2.337348in}{0.840608in}}%
\pgfpathlineto{\pgfqpoint{2.322977in}{0.839287in}}%
\pgfpathlineto{\pgfqpoint{2.320886in}{0.830332in}}%
\pgfpathlineto{\pgfqpoint{2.317320in}{0.826818in}}%
\pgfpathclose%
\pgfusepath{fill}%
\end{pgfscope}%
\begin{pgfscope}%
\pgfpathrectangle{\pgfqpoint{0.100000in}{0.100000in}}{\pgfqpoint{3.007045in}{1.925000in}}%
\pgfusepath{clip}%
\pgfsetbuttcap%
\pgfsetmiterjoin%
\definecolor{currentfill}{rgb}{0.681046,0.871895,0.656209}%
\pgfsetfillcolor{currentfill}%
\pgfsetlinewidth{0.000000pt}%
\definecolor{currentstroke}{rgb}{0.000000,0.000000,0.000000}%
\pgfsetstrokecolor{currentstroke}%
\pgfsetstrokeopacity{0.000000}%
\pgfsetdash{}{0pt}%
\pgfpathmoveto{\pgfqpoint{1.447246in}{1.572779in}}%
\pgfpathlineto{\pgfqpoint{1.445034in}{1.544102in}}%
\pgfpathlineto{\pgfqpoint{1.400379in}{1.547840in}}%
\pgfpathlineto{\pgfqpoint{1.401374in}{1.559284in}}%
\pgfpathlineto{\pgfqpoint{1.406341in}{1.607594in}}%
\pgfpathlineto{\pgfqpoint{1.403892in}{1.607820in}}%
\pgfpathlineto{\pgfqpoint{1.405844in}{1.630008in}}%
\pgfpathlineto{\pgfqpoint{1.409091in}{1.629777in}}%
\pgfpathlineto{\pgfqpoint{1.410950in}{1.652859in}}%
\pgfpathlineto{\pgfqpoint{1.448770in}{1.649792in}}%
\pgfpathlineto{\pgfqpoint{1.450846in}{1.649617in}}%
\pgfpathlineto{\pgfqpoint{1.449255in}{1.626605in}}%
\pgfpathlineto{\pgfqpoint{1.451489in}{1.626431in}}%
\pgfpathlineto{\pgfqpoint{1.450280in}{1.611196in}}%
\pgfpathlineto{\pgfqpoint{1.447246in}{1.572779in}}%
\pgfpathclose%
\pgfusepath{fill}%
\end{pgfscope}%
\begin{pgfscope}%
\pgfpathrectangle{\pgfqpoint{0.100000in}{0.100000in}}{\pgfqpoint{3.007045in}{1.925000in}}%
\pgfusepath{clip}%
\pgfsetbuttcap%
\pgfsetmiterjoin%
\definecolor{currentfill}{rgb}{0.167705,0.577578,0.294948}%
\pgfsetfillcolor{currentfill}%
\pgfsetlinewidth{0.000000pt}%
\definecolor{currentstroke}{rgb}{0.000000,0.000000,0.000000}%
\pgfsetstrokecolor{currentstroke}%
\pgfsetstrokeopacity{0.000000}%
\pgfsetdash{}{0pt}%
\pgfpathmoveto{\pgfqpoint{2.765291in}{1.169198in}}%
\pgfpathlineto{\pgfqpoint{2.770616in}{1.173606in}}%
\pgfpathlineto{\pgfqpoint{2.789514in}{1.179938in}}%
\pgfpathlineto{\pgfqpoint{2.785420in}{1.167710in}}%
\pgfpathlineto{\pgfqpoint{2.782226in}{1.168112in}}%
\pgfpathlineto{\pgfqpoint{2.778730in}{1.161873in}}%
\pgfpathlineto{\pgfqpoint{2.776737in}{1.148749in}}%
\pgfpathlineto{\pgfqpoint{2.778073in}{1.146484in}}%
\pgfpathlineto{\pgfqpoint{2.774225in}{1.130181in}}%
\pgfpathlineto{\pgfqpoint{2.772699in}{1.120440in}}%
\pgfpathlineto{\pgfqpoint{2.768773in}{1.114106in}}%
\pgfpathlineto{\pgfqpoint{2.765038in}{1.113269in}}%
\pgfpathlineto{\pgfqpoint{2.760275in}{1.121718in}}%
\pgfpathlineto{\pgfqpoint{2.760413in}{1.142327in}}%
\pgfpathlineto{\pgfqpoint{2.763321in}{1.150894in}}%
\pgfpathlineto{\pgfqpoint{2.763523in}{1.158416in}}%
\pgfpathlineto{\pgfqpoint{2.768180in}{1.159798in}}%
\pgfpathlineto{\pgfqpoint{2.769212in}{1.165623in}}%
\pgfpathlineto{\pgfqpoint{2.765291in}{1.169198in}}%
\pgfpathclose%
\pgfusepath{fill}%
\end{pgfscope}%
\begin{pgfscope}%
\pgfpathrectangle{\pgfqpoint{0.100000in}{0.100000in}}{\pgfqpoint{3.007045in}{1.925000in}}%
\pgfusepath{clip}%
\pgfsetbuttcap%
\pgfsetmiterjoin%
\definecolor{currentfill}{rgb}{0.276078,0.680969,0.375087}%
\pgfsetfillcolor{currentfill}%
\pgfsetlinewidth{0.000000pt}%
\definecolor{currentstroke}{rgb}{0.000000,0.000000,0.000000}%
\pgfsetstrokecolor{currentstroke}%
\pgfsetstrokeopacity{0.000000}%
\pgfsetdash{}{0pt}%
\pgfpathmoveto{\pgfqpoint{2.090405in}{1.513286in}}%
\pgfpathlineto{\pgfqpoint{2.073413in}{1.512435in}}%
\pgfpathlineto{\pgfqpoint{2.073048in}{1.518236in}}%
\pgfpathlineto{\pgfqpoint{2.043986in}{1.516491in}}%
\pgfpathlineto{\pgfqpoint{2.042126in}{1.545379in}}%
\pgfpathlineto{\pgfqpoint{2.032705in}{1.544832in}}%
\pgfpathlineto{\pgfqpoint{2.031351in}{1.567858in}}%
\pgfpathlineto{\pgfqpoint{2.054639in}{1.569212in}}%
\pgfpathlineto{\pgfqpoint{2.054884in}{1.563482in}}%
\pgfpathlineto{\pgfqpoint{2.066412in}{1.564272in}}%
\pgfpathlineto{\pgfqpoint{2.067108in}{1.552783in}}%
\pgfpathlineto{\pgfqpoint{2.069326in}{1.547117in}}%
\pgfpathlineto{\pgfqpoint{2.076668in}{1.547607in}}%
\pgfpathlineto{\pgfqpoint{2.077656in}{1.530311in}}%
\pgfpathlineto{\pgfqpoint{2.088935in}{1.530951in}}%
\pgfpathlineto{\pgfqpoint{2.090405in}{1.513286in}}%
\pgfpathclose%
\pgfusepath{fill}%
\end{pgfscope}%
\begin{pgfscope}%
\pgfpathrectangle{\pgfqpoint{0.100000in}{0.100000in}}{\pgfqpoint{3.007045in}{1.925000in}}%
\pgfusepath{clip}%
\pgfsetbuttcap%
\pgfsetmiterjoin%
\definecolor{currentfill}{rgb}{0.257255,0.671742,0.365859}%
\pgfsetfillcolor{currentfill}%
\pgfsetlinewidth{0.000000pt}%
\definecolor{currentstroke}{rgb}{0.000000,0.000000,0.000000}%
\pgfsetstrokecolor{currentstroke}%
\pgfsetstrokeopacity{0.000000}%
\pgfsetdash{}{0pt}%
\pgfpathmoveto{\pgfqpoint{2.836904in}{1.462520in}}%
\pgfpathlineto{\pgfqpoint{2.816097in}{1.458112in}}%
\pgfpathlineto{\pgfqpoint{2.815641in}{1.458005in}}%
\pgfpathlineto{\pgfqpoint{2.814539in}{1.460221in}}%
\pgfpathlineto{\pgfqpoint{2.815796in}{1.505313in}}%
\pgfpathlineto{\pgfqpoint{2.813736in}{1.511035in}}%
\pgfpathlineto{\pgfqpoint{2.807383in}{1.541913in}}%
\pgfpathlineto{\pgfqpoint{2.825651in}{1.545463in}}%
\pgfpathlineto{\pgfqpoint{2.828637in}{1.543127in}}%
\pgfpathlineto{\pgfqpoint{2.828599in}{1.533431in}}%
\pgfpathlineto{\pgfqpoint{2.822507in}{1.532326in}}%
\pgfpathlineto{\pgfqpoint{2.824976in}{1.520081in}}%
\pgfpathlineto{\pgfqpoint{2.829097in}{1.520876in}}%
\pgfpathlineto{\pgfqpoint{2.831656in}{1.508759in}}%
\pgfpathlineto{\pgfqpoint{2.827848in}{1.505323in}}%
\pgfpathlineto{\pgfqpoint{2.832224in}{1.502258in}}%
\pgfpathlineto{\pgfqpoint{2.833508in}{1.484777in}}%
\pgfpathlineto{\pgfqpoint{2.836904in}{1.462520in}}%
\pgfpathclose%
\pgfusepath{fill}%
\end{pgfscope}%
\begin{pgfscope}%
\pgfpathrectangle{\pgfqpoint{0.100000in}{0.100000in}}{\pgfqpoint{3.007045in}{1.925000in}}%
\pgfusepath{clip}%
\pgfsetbuttcap%
\pgfsetmiterjoin%
\definecolor{currentfill}{rgb}{0.501961,0.790588,0.501438}%
\pgfsetfillcolor{currentfill}%
\pgfsetlinewidth{0.000000pt}%
\definecolor{currentstroke}{rgb}{0.000000,0.000000,0.000000}%
\pgfsetstrokecolor{currentstroke}%
\pgfsetstrokeopacity{0.000000}%
\pgfsetdash{}{0pt}%
\pgfpathmoveto{\pgfqpoint{2.429643in}{1.109432in}}%
\pgfpathlineto{\pgfqpoint{2.425369in}{1.103463in}}%
\pgfpathlineto{\pgfqpoint{2.420050in}{1.106990in}}%
\pgfpathlineto{\pgfqpoint{2.421148in}{1.110912in}}%
\pgfpathlineto{\pgfqpoint{2.410869in}{1.119825in}}%
\pgfpathlineto{\pgfqpoint{2.411679in}{1.127319in}}%
\pgfpathlineto{\pgfqpoint{2.401483in}{1.125255in}}%
\pgfpathlineto{\pgfqpoint{2.396061in}{1.120463in}}%
\pgfpathlineto{\pgfqpoint{2.398193in}{1.116290in}}%
\pgfpathlineto{\pgfqpoint{2.392249in}{1.108544in}}%
\pgfpathlineto{\pgfqpoint{2.383396in}{1.107611in}}%
\pgfpathlineto{\pgfqpoint{2.379152in}{1.110286in}}%
\pgfpathlineto{\pgfqpoint{2.370833in}{1.107548in}}%
\pgfpathlineto{\pgfqpoint{2.359204in}{1.116262in}}%
\pgfpathlineto{\pgfqpoint{2.365037in}{1.122682in}}%
\pgfpathlineto{\pgfqpoint{2.364836in}{1.127691in}}%
\pgfpathlineto{\pgfqpoint{2.368680in}{1.131320in}}%
\pgfpathlineto{\pgfqpoint{2.372393in}{1.126861in}}%
\pgfpathlineto{\pgfqpoint{2.378692in}{1.128851in}}%
\pgfpathlineto{\pgfqpoint{2.377693in}{1.133822in}}%
\pgfpathlineto{\pgfqpoint{2.379735in}{1.139390in}}%
\pgfpathlineto{\pgfqpoint{2.384672in}{1.146612in}}%
\pgfpathlineto{\pgfqpoint{2.385587in}{1.155405in}}%
\pgfpathlineto{\pgfqpoint{2.392277in}{1.158355in}}%
\pgfpathlineto{\pgfqpoint{2.395723in}{1.148881in}}%
\pgfpathlineto{\pgfqpoint{2.401519in}{1.149007in}}%
\pgfpathlineto{\pgfqpoint{2.404127in}{1.159387in}}%
\pgfpathlineto{\pgfqpoint{2.403876in}{1.166074in}}%
\pgfpathlineto{\pgfqpoint{2.407704in}{1.166293in}}%
\pgfpathlineto{\pgfqpoint{2.407930in}{1.161976in}}%
\pgfpathlineto{\pgfqpoint{2.413032in}{1.162228in}}%
\pgfpathlineto{\pgfqpoint{2.414411in}{1.156435in}}%
\pgfpathlineto{\pgfqpoint{2.420478in}{1.156867in}}%
\pgfpathlineto{\pgfqpoint{2.420941in}{1.150887in}}%
\pgfpathlineto{\pgfqpoint{2.428157in}{1.152335in}}%
\pgfpathlineto{\pgfqpoint{2.437602in}{1.145887in}}%
\pgfpathlineto{\pgfqpoint{2.438948in}{1.139413in}}%
\pgfpathlineto{\pgfqpoint{2.431774in}{1.134957in}}%
\pgfpathlineto{\pgfqpoint{2.428704in}{1.126732in}}%
\pgfpathlineto{\pgfqpoint{2.428858in}{1.122661in}}%
\pgfpathlineto{\pgfqpoint{2.433994in}{1.115719in}}%
\pgfpathlineto{\pgfqpoint{2.429643in}{1.109432in}}%
\pgfpathclose%
\pgfusepath{fill}%
\end{pgfscope}%
\begin{pgfscope}%
\pgfpathrectangle{\pgfqpoint{0.100000in}{0.100000in}}{\pgfqpoint{3.007045in}{1.925000in}}%
\pgfusepath{clip}%
\pgfsetbuttcap%
\pgfsetmiterjoin%
\definecolor{currentfill}{rgb}{0.211995,0.624821,0.330381}%
\pgfsetfillcolor{currentfill}%
\pgfsetlinewidth{0.000000pt}%
\definecolor{currentstroke}{rgb}{0.000000,0.000000,0.000000}%
\pgfsetstrokecolor{currentstroke}%
\pgfsetstrokeopacity{0.000000}%
\pgfsetdash{}{0pt}%
\pgfpathmoveto{\pgfqpoint{1.194470in}{0.906193in}}%
\pgfpathlineto{\pgfqpoint{1.166043in}{0.909544in}}%
\pgfpathlineto{\pgfqpoint{1.171643in}{0.956432in}}%
\pgfpathlineto{\pgfqpoint{1.164202in}{0.962160in}}%
\pgfpathlineto{\pgfqpoint{1.165289in}{0.970829in}}%
\pgfpathlineto{\pgfqpoint{1.173700in}{0.972532in}}%
\pgfpathlineto{\pgfqpoint{1.139908in}{0.976686in}}%
\pgfpathlineto{\pgfqpoint{1.141795in}{0.990914in}}%
\pgfpathlineto{\pgfqpoint{1.102849in}{0.996413in}}%
\pgfpathlineto{\pgfqpoint{1.108498in}{1.034565in}}%
\pgfpathlineto{\pgfqpoint{1.118519in}{1.040536in}}%
\pgfpathlineto{\pgfqpoint{1.120573in}{1.045885in}}%
\pgfpathlineto{\pgfqpoint{1.169817in}{1.038957in}}%
\pgfpathlineto{\pgfqpoint{1.235453in}{1.031177in}}%
\pgfpathlineto{\pgfqpoint{1.235393in}{1.026018in}}%
\pgfpathlineto{\pgfqpoint{1.231839in}{1.013157in}}%
\pgfpathlineto{\pgfqpoint{1.225227in}{1.010762in}}%
\pgfpathlineto{\pgfqpoint{1.225457in}{0.993376in}}%
\pgfpathlineto{\pgfqpoint{1.223675in}{0.983885in}}%
\pgfpathlineto{\pgfqpoint{1.221151in}{0.983550in}}%
\pgfpathlineto{\pgfqpoint{1.218512in}{0.975812in}}%
\pgfpathlineto{\pgfqpoint{1.209230in}{0.966315in}}%
\pgfpathlineto{\pgfqpoint{1.201592in}{0.967633in}}%
\pgfpathlineto{\pgfqpoint{1.194470in}{0.906193in}}%
\pgfpathclose%
\pgfusepath{fill}%
\end{pgfscope}%
\begin{pgfscope}%
\pgfpathrectangle{\pgfqpoint{0.100000in}{0.100000in}}{\pgfqpoint{3.007045in}{1.925000in}}%
\pgfusepath{clip}%
\pgfsetbuttcap%
\pgfsetmiterjoin%
\definecolor{currentfill}{rgb}{0.618224,0.844844,0.597032}%
\pgfsetfillcolor{currentfill}%
\pgfsetlinewidth{0.000000pt}%
\definecolor{currentstroke}{rgb}{0.000000,0.000000,0.000000}%
\pgfsetstrokecolor{currentstroke}%
\pgfsetstrokeopacity{0.000000}%
\pgfsetdash{}{0pt}%
\pgfpathmoveto{\pgfqpoint{2.022844in}{1.028625in}}%
\pgfpathlineto{\pgfqpoint{2.005455in}{1.027915in}}%
\pgfpathlineto{\pgfqpoint{2.005554in}{1.024952in}}%
\pgfpathlineto{\pgfqpoint{1.995876in}{1.024384in}}%
\pgfpathlineto{\pgfqpoint{1.993598in}{1.031047in}}%
\pgfpathlineto{\pgfqpoint{1.993117in}{1.046153in}}%
\pgfpathlineto{\pgfqpoint{1.981728in}{1.045692in}}%
\pgfpathlineto{\pgfqpoint{1.973010in}{1.044275in}}%
\pgfpathlineto{\pgfqpoint{1.972796in}{1.051471in}}%
\pgfpathlineto{\pgfqpoint{1.975453in}{1.054398in}}%
\pgfpathlineto{\pgfqpoint{1.974361in}{1.084952in}}%
\pgfpathlineto{\pgfqpoint{1.990680in}{1.085683in}}%
\pgfpathlineto{\pgfqpoint{2.001501in}{1.072967in}}%
\pgfpathlineto{\pgfqpoint{2.010046in}{1.075951in}}%
\pgfpathlineto{\pgfqpoint{2.014367in}{1.079611in}}%
\pgfpathlineto{\pgfqpoint{2.020734in}{1.079515in}}%
\pgfpathlineto{\pgfqpoint{2.024754in}{1.078038in}}%
\pgfpathlineto{\pgfqpoint{2.031875in}{1.071974in}}%
\pgfpathlineto{\pgfqpoint{2.034604in}{1.066814in}}%
\pgfpathlineto{\pgfqpoint{2.040262in}{1.068526in}}%
\pgfpathlineto{\pgfqpoint{2.045818in}{1.065000in}}%
\pgfpathlineto{\pgfqpoint{2.054473in}{1.056687in}}%
\pgfpathlineto{\pgfqpoint{2.058159in}{1.055368in}}%
\pgfpathlineto{\pgfqpoint{2.060389in}{1.049493in}}%
\pgfpathlineto{\pgfqpoint{2.058254in}{1.047208in}}%
\pgfpathlineto{\pgfqpoint{2.053627in}{1.049022in}}%
\pgfpathlineto{\pgfqpoint{2.025561in}{1.047408in}}%
\pgfpathlineto{\pgfqpoint{2.026620in}{1.028636in}}%
\pgfpathlineto{\pgfqpoint{2.022844in}{1.028625in}}%
\pgfpathclose%
\pgfusepath{fill}%
\end{pgfscope}%
\begin{pgfscope}%
\pgfpathrectangle{\pgfqpoint{0.100000in}{0.100000in}}{\pgfqpoint{3.007045in}{1.925000in}}%
\pgfusepath{clip}%
\pgfsetbuttcap%
\pgfsetmiterjoin%
\definecolor{currentfill}{rgb}{0.794233,0.919262,0.767705}%
\pgfsetfillcolor{currentfill}%
\pgfsetlinewidth{0.000000pt}%
\definecolor{currentstroke}{rgb}{0.000000,0.000000,0.000000}%
\pgfsetstrokecolor{currentstroke}%
\pgfsetstrokeopacity{0.000000}%
\pgfsetdash{}{0pt}%
\pgfpathmoveto{\pgfqpoint{2.309942in}{0.708924in}}%
\pgfpathlineto{\pgfqpoint{2.309841in}{0.700126in}}%
\pgfpathlineto{\pgfqpoint{2.301073in}{0.698033in}}%
\pgfpathlineto{\pgfqpoint{2.297933in}{0.689896in}}%
\pgfpathlineto{\pgfqpoint{2.290454in}{0.689136in}}%
\pgfpathlineto{\pgfqpoint{2.289927in}{0.694879in}}%
\pgfpathlineto{\pgfqpoint{2.284645in}{0.694345in}}%
\pgfpathlineto{\pgfqpoint{2.283592in}{0.699581in}}%
\pgfpathlineto{\pgfqpoint{2.277794in}{0.699248in}}%
\pgfpathlineto{\pgfqpoint{2.276404in}{0.718014in}}%
\pgfpathlineto{\pgfqpoint{2.274136in}{0.722813in}}%
\pgfpathlineto{\pgfqpoint{2.281053in}{0.729035in}}%
\pgfpathlineto{\pgfqpoint{2.285996in}{0.729613in}}%
\pgfpathlineto{\pgfqpoint{2.285327in}{0.735328in}}%
\pgfpathlineto{\pgfqpoint{2.290941in}{0.736895in}}%
\pgfpathlineto{\pgfqpoint{2.290204in}{0.743581in}}%
\pgfpathlineto{\pgfqpoint{2.295738in}{0.746225in}}%
\pgfpathlineto{\pgfqpoint{2.312802in}{0.748173in}}%
\pgfpathlineto{\pgfqpoint{2.321001in}{0.750113in}}%
\pgfpathlineto{\pgfqpoint{2.323388in}{0.743508in}}%
\pgfpathlineto{\pgfqpoint{2.330141in}{0.735377in}}%
\pgfpathlineto{\pgfqpoint{2.327177in}{0.732545in}}%
\pgfpathlineto{\pgfqpoint{2.311941in}{0.730610in}}%
\pgfpathlineto{\pgfqpoint{2.312554in}{0.726822in}}%
\pgfpathlineto{\pgfqpoint{2.306778in}{0.726142in}}%
\pgfpathlineto{\pgfqpoint{2.308027in}{0.714591in}}%
\pgfpathlineto{\pgfqpoint{2.309942in}{0.708924in}}%
\pgfpathclose%
\pgfusepath{fill}%
\end{pgfscope}%
\begin{pgfscope}%
\pgfpathrectangle{\pgfqpoint{0.100000in}{0.100000in}}{\pgfqpoint{3.007045in}{1.925000in}}%
\pgfusepath{clip}%
\pgfsetbuttcap%
\pgfsetmiterjoin%
\definecolor{currentfill}{rgb}{0.732472,0.893549,0.706282}%
\pgfsetfillcolor{currentfill}%
\pgfsetlinewidth{0.000000pt}%
\definecolor{currentstroke}{rgb}{0.000000,0.000000,0.000000}%
\pgfsetstrokecolor{currentstroke}%
\pgfsetstrokeopacity{0.000000}%
\pgfsetdash{}{0pt}%
\pgfpathmoveto{\pgfqpoint{0.829274in}{1.562729in}}%
\pgfpathlineto{\pgfqpoint{0.828248in}{1.558524in}}%
\pgfpathlineto{\pgfqpoint{0.823175in}{1.550976in}}%
\pgfpathlineto{\pgfqpoint{0.824049in}{1.547593in}}%
\pgfpathlineto{\pgfqpoint{0.819243in}{1.542361in}}%
\pgfpathlineto{\pgfqpoint{0.809080in}{1.497950in}}%
\pgfpathlineto{\pgfqpoint{0.811385in}{1.497272in}}%
\pgfpathlineto{\pgfqpoint{0.798116in}{1.438303in}}%
\pgfpathlineto{\pgfqpoint{0.734607in}{1.453046in}}%
\pgfpathlineto{\pgfqpoint{0.703710in}{1.460774in}}%
\pgfpathlineto{\pgfqpoint{0.703346in}{1.460924in}}%
\pgfpathlineto{\pgfqpoint{0.729959in}{1.568628in}}%
\pgfpathlineto{\pgfqpoint{0.733534in}{1.580107in}}%
\pgfpathlineto{\pgfqpoint{0.736446in}{1.580574in}}%
\pgfpathlineto{\pgfqpoint{0.739760in}{1.573820in}}%
\pgfpathlineto{\pgfqpoint{0.746496in}{1.573138in}}%
\pgfpathlineto{\pgfqpoint{0.749229in}{1.584343in}}%
\pgfpathlineto{\pgfqpoint{0.755702in}{1.582709in}}%
\pgfpathlineto{\pgfqpoint{0.758886in}{1.587771in}}%
\pgfpathlineto{\pgfqpoint{0.762552in}{1.586835in}}%
\pgfpathlineto{\pgfqpoint{0.763908in}{1.592366in}}%
\pgfpathlineto{\pgfqpoint{0.768216in}{1.591309in}}%
\pgfpathlineto{\pgfqpoint{0.771340in}{1.602215in}}%
\pgfpathlineto{\pgfqpoint{0.775430in}{1.609514in}}%
\pgfpathlineto{\pgfqpoint{0.783037in}{1.611433in}}%
\pgfpathlineto{\pgfqpoint{0.780379in}{1.600195in}}%
\pgfpathlineto{\pgfqpoint{0.777625in}{1.600841in}}%
\pgfpathlineto{\pgfqpoint{0.774977in}{1.589715in}}%
\pgfpathlineto{\pgfqpoint{0.779546in}{1.588394in}}%
\pgfpathlineto{\pgfqpoint{0.781320in}{1.594007in}}%
\pgfpathlineto{\pgfqpoint{0.808189in}{1.587646in}}%
\pgfpathlineto{\pgfqpoint{0.811192in}{1.587476in}}%
\pgfpathlineto{\pgfqpoint{0.820142in}{1.591987in}}%
\pgfpathlineto{\pgfqpoint{0.824845in}{1.588003in}}%
\pgfpathlineto{\pgfqpoint{0.824028in}{1.581837in}}%
\pgfpathlineto{\pgfqpoint{0.829905in}{1.577513in}}%
\pgfpathlineto{\pgfqpoint{0.827488in}{1.571131in}}%
\pgfpathlineto{\pgfqpoint{0.829274in}{1.562729in}}%
\pgfpathclose%
\pgfusepath{fill}%
\end{pgfscope}%
\begin{pgfscope}%
\pgfpathrectangle{\pgfqpoint{0.100000in}{0.100000in}}{\pgfqpoint{3.007045in}{1.925000in}}%
\pgfusepath{clip}%
\pgfsetbuttcap%
\pgfsetmiterjoin%
\definecolor{currentfill}{rgb}{0.718447,0.887643,0.692626}%
\pgfsetfillcolor{currentfill}%
\pgfsetlinewidth{0.000000pt}%
\definecolor{currentstroke}{rgb}{0.000000,0.000000,0.000000}%
\pgfsetstrokecolor{currentstroke}%
\pgfsetstrokeopacity{0.000000}%
\pgfsetdash{}{0pt}%
\pgfpathmoveto{\pgfqpoint{1.910648in}{0.511564in}}%
\pgfpathlineto{\pgfqpoint{1.903450in}{0.513741in}}%
\pgfpathlineto{\pgfqpoint{1.888369in}{0.520432in}}%
\pgfpathlineto{\pgfqpoint{1.876777in}{0.523740in}}%
\pgfpathlineto{\pgfqpoint{1.868009in}{0.522526in}}%
\pgfpathlineto{\pgfqpoint{1.859158in}{0.523013in}}%
\pgfpathlineto{\pgfqpoint{1.843789in}{0.520220in}}%
\pgfpathlineto{\pgfqpoint{1.839662in}{0.517608in}}%
\pgfpathlineto{\pgfqpoint{1.834496in}{0.525796in}}%
\pgfpathlineto{\pgfqpoint{1.840585in}{0.532547in}}%
\pgfpathlineto{\pgfqpoint{1.842229in}{0.537436in}}%
\pgfpathlineto{\pgfqpoint{1.846781in}{0.541896in}}%
\pgfpathlineto{\pgfqpoint{1.846835in}{0.557616in}}%
\pgfpathlineto{\pgfqpoint{1.843280in}{0.560440in}}%
\pgfpathlineto{\pgfqpoint{1.847088in}{0.567676in}}%
\pgfpathlineto{\pgfqpoint{1.844939in}{0.573686in}}%
\pgfpathlineto{\pgfqpoint{1.847738in}{0.581014in}}%
\pgfpathlineto{\pgfqpoint{1.850658in}{0.583534in}}%
\pgfpathlineto{\pgfqpoint{1.853934in}{0.596710in}}%
\pgfpathlineto{\pgfqpoint{1.856414in}{0.600398in}}%
\pgfpathlineto{\pgfqpoint{1.854154in}{0.604333in}}%
\pgfpathlineto{\pgfqpoint{1.854631in}{0.617364in}}%
\pgfpathlineto{\pgfqpoint{1.855514in}{0.623416in}}%
\pgfpathlineto{\pgfqpoint{1.860787in}{0.623525in}}%
\pgfpathlineto{\pgfqpoint{1.863641in}{0.629386in}}%
\pgfpathlineto{\pgfqpoint{1.872399in}{0.629573in}}%
\pgfpathlineto{\pgfqpoint{1.883960in}{0.629796in}}%
\pgfpathlineto{\pgfqpoint{1.886870in}{0.628687in}}%
\pgfpathlineto{\pgfqpoint{1.892685in}{0.626874in}}%
\pgfpathlineto{\pgfqpoint{1.895632in}{0.622616in}}%
\pgfpathlineto{\pgfqpoint{1.896550in}{0.598540in}}%
\pgfpathlineto{\pgfqpoint{1.909431in}{0.599271in}}%
\pgfpathlineto{\pgfqpoint{1.910006in}{0.578352in}}%
\pgfpathlineto{\pgfqpoint{1.908086in}{0.571754in}}%
\pgfpathlineto{\pgfqpoint{1.909910in}{0.542133in}}%
\pgfpathlineto{\pgfqpoint{1.910648in}{0.511564in}}%
\pgfpathclose%
\pgfusepath{fill}%
\end{pgfscope}%
\begin{pgfscope}%
\pgfpathrectangle{\pgfqpoint{0.100000in}{0.100000in}}{\pgfqpoint{3.007045in}{1.925000in}}%
\pgfusepath{clip}%
\pgfsetbuttcap%
\pgfsetmiterjoin%
\definecolor{currentfill}{rgb}{0.200923,0.613010,0.321522}%
\pgfsetfillcolor{currentfill}%
\pgfsetlinewidth{0.000000pt}%
\definecolor{currentstroke}{rgb}{0.000000,0.000000,0.000000}%
\pgfsetstrokecolor{currentstroke}%
\pgfsetstrokeopacity{0.000000}%
\pgfsetdash{}{0pt}%
\pgfpathmoveto{\pgfqpoint{0.677176in}{1.047687in}}%
\pgfpathlineto{\pgfqpoint{0.672324in}{1.047770in}}%
\pgfpathlineto{\pgfqpoint{0.654898in}{1.051840in}}%
\pgfpathlineto{\pgfqpoint{0.593532in}{1.066855in}}%
\pgfpathlineto{\pgfqpoint{0.573744in}{1.071996in}}%
\pgfpathlineto{\pgfqpoint{0.554172in}{1.076466in}}%
\pgfpathlineto{\pgfqpoint{0.557260in}{1.082859in}}%
\pgfpathlineto{\pgfqpoint{0.556142in}{1.096690in}}%
\pgfpathlineto{\pgfqpoint{0.557824in}{1.101480in}}%
\pgfpathlineto{\pgfqpoint{0.556496in}{1.110688in}}%
\pgfpathlineto{\pgfqpoint{0.558630in}{1.113262in}}%
\pgfpathlineto{\pgfqpoint{0.553369in}{1.124234in}}%
\pgfpathlineto{\pgfqpoint{0.554104in}{1.127797in}}%
\pgfpathlineto{\pgfqpoint{0.551392in}{1.137382in}}%
\pgfpathlineto{\pgfqpoint{0.551954in}{1.148285in}}%
\pgfpathlineto{\pgfqpoint{0.554321in}{1.151293in}}%
\pgfpathlineto{\pgfqpoint{0.552868in}{1.161985in}}%
\pgfpathlineto{\pgfqpoint{0.546742in}{1.170366in}}%
\pgfpathlineto{\pgfqpoint{0.543660in}{1.171437in}}%
\pgfpathlineto{\pgfqpoint{0.543712in}{1.184135in}}%
\pgfpathlineto{\pgfqpoint{0.540581in}{1.185427in}}%
\pgfpathlineto{\pgfqpoint{0.543469in}{1.191912in}}%
\pgfpathlineto{\pgfqpoint{0.539193in}{1.195591in}}%
\pgfpathlineto{\pgfqpoint{0.537345in}{1.200997in}}%
\pgfpathlineto{\pgfqpoint{0.532644in}{1.204682in}}%
\pgfpathlineto{\pgfqpoint{0.532191in}{1.211551in}}%
\pgfpathlineto{\pgfqpoint{0.530069in}{1.215499in}}%
\pgfpathlineto{\pgfqpoint{0.522889in}{1.217818in}}%
\pgfpathlineto{\pgfqpoint{0.527157in}{1.220456in}}%
\pgfpathlineto{\pgfqpoint{0.527593in}{1.228218in}}%
\pgfpathlineto{\pgfqpoint{0.524524in}{1.231551in}}%
\pgfpathlineto{\pgfqpoint{0.525021in}{1.240725in}}%
\pgfpathlineto{\pgfqpoint{0.514807in}{1.248637in}}%
\pgfpathlineto{\pgfqpoint{0.513336in}{1.257851in}}%
\pgfpathlineto{\pgfqpoint{0.514663in}{1.260246in}}%
\pgfpathlineto{\pgfqpoint{0.519485in}{1.264180in}}%
\pgfpathlineto{\pgfqpoint{0.522497in}{1.269886in}}%
\pgfpathlineto{\pgfqpoint{0.520542in}{1.277526in}}%
\pgfpathlineto{\pgfqpoint{0.524302in}{1.283464in}}%
\pgfpathlineto{\pgfqpoint{0.533766in}{1.269179in}}%
\pgfpathlineto{\pgfqpoint{0.540209in}{1.259221in}}%
\pgfpathlineto{\pgfqpoint{0.567833in}{1.216455in}}%
\pgfpathlineto{\pgfqpoint{0.590725in}{1.181138in}}%
\pgfpathlineto{\pgfqpoint{0.616671in}{1.141013in}}%
\pgfpathlineto{\pgfqpoint{0.667024in}{1.063268in}}%
\pgfpathlineto{\pgfqpoint{0.677176in}{1.047687in}}%
\pgfpathclose%
\pgfusepath{fill}%
\end{pgfscope}%
\begin{pgfscope}%
\pgfpathrectangle{\pgfqpoint{0.100000in}{0.100000in}}{\pgfqpoint{3.007045in}{1.925000in}}%
\pgfusepath{clip}%
\pgfsetbuttcap%
\pgfsetmiterjoin%
\definecolor{currentfill}{rgb}{0.623760,0.847428,0.601584}%
\pgfsetfillcolor{currentfill}%
\pgfsetlinewidth{0.000000pt}%
\definecolor{currentstroke}{rgb}{0.000000,0.000000,0.000000}%
\pgfsetstrokecolor{currentstroke}%
\pgfsetstrokeopacity{0.000000}%
\pgfsetdash{}{0pt}%
\pgfpathmoveto{\pgfqpoint{2.442052in}{0.679496in}}%
\pgfpathlineto{\pgfqpoint{2.428996in}{0.678445in}}%
\pgfpathlineto{\pgfqpoint{2.425943in}{0.682862in}}%
\pgfpathlineto{\pgfqpoint{2.421279in}{0.683375in}}%
\pgfpathlineto{\pgfqpoint{2.420802in}{0.691745in}}%
\pgfpathlineto{\pgfqpoint{2.422816in}{0.692785in}}%
\pgfpathlineto{\pgfqpoint{2.420123in}{0.701779in}}%
\pgfpathlineto{\pgfqpoint{2.437033in}{0.704008in}}%
\pgfpathlineto{\pgfqpoint{2.447789in}{0.700910in}}%
\pgfpathlineto{\pgfqpoint{2.448668in}{0.692985in}}%
\pgfpathlineto{\pgfqpoint{2.445730in}{0.689304in}}%
\pgfpathlineto{\pgfqpoint{2.443946in}{0.680145in}}%
\pgfpathlineto{\pgfqpoint{2.442052in}{0.679496in}}%
\pgfpathclose%
\pgfusepath{fill}%
\end{pgfscope}%
\begin{pgfscope}%
\pgfpathrectangle{\pgfqpoint{0.100000in}{0.100000in}}{\pgfqpoint{3.007045in}{1.925000in}}%
\pgfusepath{clip}%
\pgfsetbuttcap%
\pgfsetmiterjoin%
\definecolor{currentfill}{rgb}{0.252595,0.668128,0.362860}%
\pgfsetfillcolor{currentfill}%
\pgfsetlinewidth{0.000000pt}%
\definecolor{currentstroke}{rgb}{0.000000,0.000000,0.000000}%
\pgfsetstrokecolor{currentstroke}%
\pgfsetstrokeopacity{0.000000}%
\pgfsetdash{}{0pt}%
\pgfpathmoveto{\pgfqpoint{0.404299in}{1.600395in}}%
\pgfpathlineto{\pgfqpoint{0.415190in}{1.591253in}}%
\pgfpathlineto{\pgfqpoint{0.419385in}{1.591902in}}%
\pgfpathlineto{\pgfqpoint{0.426249in}{1.588654in}}%
\pgfpathlineto{\pgfqpoint{0.429770in}{1.590236in}}%
\pgfpathlineto{\pgfqpoint{0.436255in}{1.587889in}}%
\pgfpathlineto{\pgfqpoint{0.446705in}{1.586464in}}%
\pgfpathlineto{\pgfqpoint{0.460693in}{1.593122in}}%
\pgfpathlineto{\pgfqpoint{0.465330in}{1.591720in}}%
\pgfpathlineto{\pgfqpoint{0.469189in}{1.594432in}}%
\pgfpathlineto{\pgfqpoint{0.474662in}{1.592743in}}%
\pgfpathlineto{\pgfqpoint{0.455007in}{1.530523in}}%
\pgfpathlineto{\pgfqpoint{0.397486in}{1.548285in}}%
\pgfpathlineto{\pgfqpoint{0.383232in}{1.552564in}}%
\pgfpathlineto{\pgfqpoint{0.384587in}{1.562231in}}%
\pgfpathlineto{\pgfqpoint{0.389591in}{1.567203in}}%
\pgfpathlineto{\pgfqpoint{0.387396in}{1.575981in}}%
\pgfpathlineto{\pgfqpoint{0.381248in}{1.578474in}}%
\pgfpathlineto{\pgfqpoint{0.384929in}{1.586076in}}%
\pgfpathlineto{\pgfqpoint{0.390435in}{1.585257in}}%
\pgfpathlineto{\pgfqpoint{0.395865in}{1.592654in}}%
\pgfpathlineto{\pgfqpoint{0.404299in}{1.600395in}}%
\pgfpathclose%
\pgfusepath{fill}%
\end{pgfscope}%
\begin{pgfscope}%
\pgfpathrectangle{\pgfqpoint{0.100000in}{0.100000in}}{\pgfqpoint{3.007045in}{1.925000in}}%
\pgfusepath{clip}%
\pgfsetbuttcap%
\pgfsetmiterjoin%
\definecolor{currentfill}{rgb}{0.507497,0.793172,0.505990}%
\pgfsetfillcolor{currentfill}%
\pgfsetlinewidth{0.000000pt}%
\definecolor{currentstroke}{rgb}{0.000000,0.000000,0.000000}%
\pgfsetstrokecolor{currentstroke}%
\pgfsetstrokeopacity{0.000000}%
\pgfsetdash{}{0pt}%
\pgfpathmoveto{\pgfqpoint{2.295738in}{0.746225in}}%
\pgfpathlineto{\pgfqpoint{2.290204in}{0.743581in}}%
\pgfpathlineto{\pgfqpoint{2.290941in}{0.736895in}}%
\pgfpathlineto{\pgfqpoint{2.285327in}{0.735328in}}%
\pgfpathlineto{\pgfqpoint{2.285996in}{0.729613in}}%
\pgfpathlineto{\pgfqpoint{2.281053in}{0.729035in}}%
\pgfpathlineto{\pgfqpoint{2.279667in}{0.746313in}}%
\pgfpathlineto{\pgfqpoint{2.252087in}{0.743598in}}%
\pgfpathlineto{\pgfqpoint{2.247289in}{0.747109in}}%
\pgfpathlineto{\pgfqpoint{2.243973in}{0.753410in}}%
\pgfpathlineto{\pgfqpoint{2.242492in}{0.770671in}}%
\pgfpathlineto{\pgfqpoint{2.244576in}{0.775293in}}%
\pgfpathlineto{\pgfqpoint{2.249772in}{0.779385in}}%
\pgfpathlineto{\pgfqpoint{2.248013in}{0.785562in}}%
\pgfpathlineto{\pgfqpoint{2.257385in}{0.796619in}}%
\pgfpathlineto{\pgfqpoint{2.255427in}{0.799367in}}%
\pgfpathlineto{\pgfqpoint{2.258979in}{0.805773in}}%
\pgfpathlineto{\pgfqpoint{2.264190in}{0.805986in}}%
\pgfpathlineto{\pgfqpoint{2.267838in}{0.800489in}}%
\pgfpathlineto{\pgfqpoint{2.278832in}{0.801530in}}%
\pgfpathlineto{\pgfqpoint{2.281013in}{0.795971in}}%
\pgfpathlineto{\pgfqpoint{2.287725in}{0.796480in}}%
\pgfpathlineto{\pgfqpoint{2.289813in}{0.770791in}}%
\pgfpathlineto{\pgfqpoint{2.293144in}{0.771183in}}%
\pgfpathlineto{\pgfqpoint{2.295738in}{0.746225in}}%
\pgfpathclose%
\pgfusepath{fill}%
\end{pgfscope}%
\begin{pgfscope}%
\pgfpathrectangle{\pgfqpoint{0.100000in}{0.100000in}}{\pgfqpoint{3.007045in}{1.925000in}}%
\pgfusepath{clip}%
\pgfsetbuttcap%
\pgfsetmiterjoin%
\definecolor{currentfill}{rgb}{0.737147,0.895517,0.710834}%
\pgfsetfillcolor{currentfill}%
\pgfsetlinewidth{0.000000pt}%
\definecolor{currentstroke}{rgb}{0.000000,0.000000,0.000000}%
\pgfsetstrokecolor{currentstroke}%
\pgfsetstrokeopacity{0.000000}%
\pgfsetdash{}{0pt}%
\pgfpathmoveto{\pgfqpoint{2.067386in}{1.211625in}}%
\pgfpathlineto{\pgfqpoint{2.066190in}{1.226950in}}%
\pgfpathlineto{\pgfqpoint{2.060239in}{1.229418in}}%
\pgfpathlineto{\pgfqpoint{2.058847in}{1.247095in}}%
\pgfpathlineto{\pgfqpoint{2.065605in}{1.247690in}}%
\pgfpathlineto{\pgfqpoint{2.072749in}{1.252662in}}%
\pgfpathlineto{\pgfqpoint{2.075117in}{1.258653in}}%
\pgfpathlineto{\pgfqpoint{2.073487in}{1.281826in}}%
\pgfpathlineto{\pgfqpoint{2.090562in}{1.283138in}}%
\pgfpathlineto{\pgfqpoint{2.107101in}{1.284740in}}%
\pgfpathlineto{\pgfqpoint{2.107905in}{1.276923in}}%
\pgfpathlineto{\pgfqpoint{2.110362in}{1.252216in}}%
\pgfpathlineto{\pgfqpoint{2.099165in}{1.251343in}}%
\pgfpathlineto{\pgfqpoint{2.100744in}{1.229280in}}%
\pgfpathlineto{\pgfqpoint{2.094960in}{1.228863in}}%
\pgfpathlineto{\pgfqpoint{2.087435in}{1.213369in}}%
\pgfpathlineto{\pgfqpoint{2.067386in}{1.211625in}}%
\pgfpathclose%
\pgfusepath{fill}%
\end{pgfscope}%
\begin{pgfscope}%
\pgfpathrectangle{\pgfqpoint{0.100000in}{0.100000in}}{\pgfqpoint{3.007045in}{1.925000in}}%
\pgfusepath{clip}%
\pgfsetbuttcap%
\pgfsetmiterjoin%
\definecolor{currentfill}{rgb}{0.709097,0.883706,0.683522}%
\pgfsetfillcolor{currentfill}%
\pgfsetlinewidth{0.000000pt}%
\definecolor{currentstroke}{rgb}{0.000000,0.000000,0.000000}%
\pgfsetstrokecolor{currentstroke}%
\pgfsetstrokeopacity{0.000000}%
\pgfsetdash{}{0pt}%
\pgfpathmoveto{\pgfqpoint{1.952131in}{0.625210in}}%
\pgfpathlineto{\pgfqpoint{1.961233in}{0.620328in}}%
\pgfpathlineto{\pgfqpoint{1.958991in}{0.616692in}}%
\pgfpathlineto{\pgfqpoint{1.958935in}{0.610678in}}%
\pgfpathlineto{\pgfqpoint{1.954440in}{0.605869in}}%
\pgfpathlineto{\pgfqpoint{1.954185in}{0.597454in}}%
\pgfpathlineto{\pgfqpoint{1.930070in}{0.596671in}}%
\pgfpathlineto{\pgfqpoint{1.927440in}{0.602431in}}%
\pgfpathlineto{\pgfqpoint{1.921122in}{0.606604in}}%
\pgfpathlineto{\pgfqpoint{1.915657in}{0.600227in}}%
\pgfpathlineto{\pgfqpoint{1.909431in}{0.599271in}}%
\pgfpathlineto{\pgfqpoint{1.896550in}{0.598540in}}%
\pgfpathlineto{\pgfqpoint{1.895632in}{0.622616in}}%
\pgfpathlineto{\pgfqpoint{1.892685in}{0.626874in}}%
\pgfpathlineto{\pgfqpoint{1.886870in}{0.628687in}}%
\pgfpathlineto{\pgfqpoint{1.890720in}{0.630209in}}%
\pgfpathlineto{\pgfqpoint{1.901476in}{0.640393in}}%
\pgfpathlineto{\pgfqpoint{1.896682in}{0.645552in}}%
\pgfpathlineto{\pgfqpoint{1.890709in}{0.647179in}}%
\pgfpathlineto{\pgfqpoint{1.886783in}{0.652756in}}%
\pgfpathlineto{\pgfqpoint{1.906746in}{0.653283in}}%
\pgfpathlineto{\pgfqpoint{1.906637in}{0.659096in}}%
\pgfpathlineto{\pgfqpoint{1.921110in}{0.659462in}}%
\pgfpathlineto{\pgfqpoint{1.922010in}{0.647995in}}%
\pgfpathlineto{\pgfqpoint{1.926639in}{0.641691in}}%
\pgfpathlineto{\pgfqpoint{1.933198in}{0.637846in}}%
\pgfpathlineto{\pgfqpoint{1.933963in}{0.633867in}}%
\pgfpathlineto{\pgfqpoint{1.942513in}{0.628295in}}%
\pgfpathlineto{\pgfqpoint{1.952131in}{0.625210in}}%
\pgfpathclose%
\pgfusepath{fill}%
\end{pgfscope}%
\begin{pgfscope}%
\pgfpathrectangle{\pgfqpoint{0.100000in}{0.100000in}}{\pgfqpoint{3.007045in}{1.925000in}}%
\pgfusepath{clip}%
\pgfsetbuttcap%
\pgfsetmiterjoin%
\definecolor{currentfill}{rgb}{0.468743,0.775087,0.474125}%
\pgfsetfillcolor{currentfill}%
\pgfsetlinewidth{0.000000pt}%
\definecolor{currentstroke}{rgb}{0.000000,0.000000,0.000000}%
\pgfsetstrokecolor{currentstroke}%
\pgfsetstrokeopacity{0.000000}%
\pgfsetdash{}{0pt}%
\pgfpathmoveto{\pgfqpoint{2.131121in}{1.073626in}}%
\pgfpathlineto{\pgfqpoint{2.128732in}{1.074544in}}%
\pgfpathlineto{\pgfqpoint{2.116398in}{1.073683in}}%
\pgfpathlineto{\pgfqpoint{2.114940in}{1.096476in}}%
\pgfpathlineto{\pgfqpoint{2.097811in}{1.095356in}}%
\pgfpathlineto{\pgfqpoint{2.096418in}{1.118342in}}%
\pgfpathlineto{\pgfqpoint{2.119277in}{1.119542in}}%
\pgfpathlineto{\pgfqpoint{2.118010in}{1.123238in}}%
\pgfpathlineto{\pgfqpoint{2.117823in}{1.135815in}}%
\pgfpathlineto{\pgfqpoint{2.135746in}{1.137366in}}%
\pgfpathlineto{\pgfqpoint{2.136969in}{1.118978in}}%
\pgfpathlineto{\pgfqpoint{2.135045in}{1.113976in}}%
\pgfpathlineto{\pgfqpoint{2.134603in}{1.106062in}}%
\pgfpathlineto{\pgfqpoint{2.136503in}{1.101602in}}%
\pgfpathlineto{\pgfqpoint{2.134514in}{1.098251in}}%
\pgfpathlineto{\pgfqpoint{2.137521in}{1.089690in}}%
\pgfpathlineto{\pgfqpoint{2.131121in}{1.073626in}}%
\pgfpathclose%
\pgfusepath{fill}%
\end{pgfscope}%
\begin{pgfscope}%
\pgfpathrectangle{\pgfqpoint{0.100000in}{0.100000in}}{\pgfqpoint{3.007045in}{1.925000in}}%
\pgfusepath{clip}%
\pgfsetbuttcap%
\pgfsetmiterjoin%
\definecolor{currentfill}{rgb}{0.568397,0.821592,0.556063}%
\pgfsetfillcolor{currentfill}%
\pgfsetlinewidth{0.000000pt}%
\definecolor{currentstroke}{rgb}{0.000000,0.000000,0.000000}%
\pgfsetstrokecolor{currentstroke}%
\pgfsetstrokeopacity{0.000000}%
\pgfsetdash{}{0pt}%
\pgfpathmoveto{\pgfqpoint{2.450037in}{1.161421in}}%
\pgfpathlineto{\pgfqpoint{2.440950in}{1.158535in}}%
\pgfpathlineto{\pgfqpoint{2.435873in}{1.154519in}}%
\pgfpathlineto{\pgfqpoint{2.437602in}{1.145887in}}%
\pgfpathlineto{\pgfqpoint{2.428157in}{1.152335in}}%
\pgfpathlineto{\pgfqpoint{2.420941in}{1.150887in}}%
\pgfpathlineto{\pgfqpoint{2.420478in}{1.156867in}}%
\pgfpathlineto{\pgfqpoint{2.414411in}{1.156435in}}%
\pgfpathlineto{\pgfqpoint{2.413032in}{1.162228in}}%
\pgfpathlineto{\pgfqpoint{2.407930in}{1.161976in}}%
\pgfpathlineto{\pgfqpoint{2.407704in}{1.166293in}}%
\pgfpathlineto{\pgfqpoint{2.413821in}{1.167145in}}%
\pgfpathlineto{\pgfqpoint{2.413102in}{1.179720in}}%
\pgfpathlineto{\pgfqpoint{2.418866in}{1.180027in}}%
\pgfpathlineto{\pgfqpoint{2.418057in}{1.191742in}}%
\pgfpathlineto{\pgfqpoint{2.419602in}{1.197771in}}%
\pgfpathlineto{\pgfqpoint{2.417290in}{1.203518in}}%
\pgfpathlineto{\pgfqpoint{2.416942in}{1.209349in}}%
\pgfpathlineto{\pgfqpoint{2.422680in}{1.209735in}}%
\pgfpathlineto{\pgfqpoint{2.422257in}{1.215672in}}%
\pgfpathlineto{\pgfqpoint{2.427935in}{1.216104in}}%
\pgfpathlineto{\pgfqpoint{2.428341in}{1.210131in}}%
\pgfpathlineto{\pgfqpoint{2.439065in}{1.211016in}}%
\pgfpathlineto{\pgfqpoint{2.439887in}{1.202342in}}%
\pgfpathlineto{\pgfqpoint{2.446995in}{1.200159in}}%
\pgfpathlineto{\pgfqpoint{2.446120in}{1.194088in}}%
\pgfpathlineto{\pgfqpoint{2.447277in}{1.183280in}}%
\pgfpathlineto{\pgfqpoint{2.441696in}{1.172351in}}%
\pgfpathlineto{\pgfqpoint{2.450037in}{1.161421in}}%
\pgfpathclose%
\pgfusepath{fill}%
\end{pgfscope}%
\begin{pgfscope}%
\pgfpathrectangle{\pgfqpoint{0.100000in}{0.100000in}}{\pgfqpoint{3.007045in}{1.925000in}}%
\pgfusepath{clip}%
\pgfsetbuttcap%
\pgfsetmiterjoin%
\definecolor{currentfill}{rgb}{0.407843,0.745559,0.439677}%
\pgfsetfillcolor{currentfill}%
\pgfsetlinewidth{0.000000pt}%
\definecolor{currentstroke}{rgb}{0.000000,0.000000,0.000000}%
\pgfsetstrokecolor{currentstroke}%
\pgfsetstrokeopacity{0.000000}%
\pgfsetdash{}{0pt}%
\pgfpathmoveto{\pgfqpoint{1.594392in}{1.173857in}}%
\pgfpathlineto{\pgfqpoint{1.623649in}{1.172778in}}%
\pgfpathlineto{\pgfqpoint{1.623055in}{1.155630in}}%
\pgfpathlineto{\pgfqpoint{1.622833in}{1.149886in}}%
\pgfpathlineto{\pgfqpoint{1.594289in}{1.150925in}}%
\pgfpathlineto{\pgfqpoint{1.594074in}{1.145204in}}%
\pgfpathlineto{\pgfqpoint{1.565600in}{1.146405in}}%
\pgfpathlineto{\pgfqpoint{1.567058in}{1.175010in}}%
\pgfpathlineto{\pgfqpoint{1.594392in}{1.173857in}}%
\pgfpathclose%
\pgfusepath{fill}%
\end{pgfscope}%
\begin{pgfscope}%
\pgfpathrectangle{\pgfqpoint{0.100000in}{0.100000in}}{\pgfqpoint{3.007045in}{1.925000in}}%
\pgfusepath{clip}%
\pgfsetbuttcap%
\pgfsetmiterjoin%
\definecolor{currentfill}{rgb}{0.699746,0.879769,0.674418}%
\pgfsetfillcolor{currentfill}%
\pgfsetlinewidth{0.000000pt}%
\definecolor{currentstroke}{rgb}{0.000000,0.000000,0.000000}%
\pgfsetstrokecolor{currentstroke}%
\pgfsetstrokeopacity{0.000000}%
\pgfsetdash{}{0pt}%
\pgfpathmoveto{\pgfqpoint{2.695089in}{1.071478in}}%
\pgfpathlineto{\pgfqpoint{2.691815in}{1.071736in}}%
\pgfpathlineto{\pgfqpoint{2.691645in}{1.081236in}}%
\pgfpathlineto{\pgfqpoint{2.685692in}{1.079312in}}%
\pgfpathlineto{\pgfqpoint{2.680930in}{1.081273in}}%
\pgfpathlineto{\pgfqpoint{2.674370in}{1.085389in}}%
\pgfpathlineto{\pgfqpoint{2.667768in}{1.085077in}}%
\pgfpathlineto{\pgfqpoint{2.665311in}{1.095015in}}%
\pgfpathlineto{\pgfqpoint{2.670092in}{1.099194in}}%
\pgfpathlineto{\pgfqpoint{2.646388in}{1.101662in}}%
\pgfpathlineto{\pgfqpoint{2.645322in}{1.106286in}}%
\pgfpathlineto{\pgfqpoint{2.649572in}{1.112889in}}%
\pgfpathlineto{\pgfqpoint{2.649818in}{1.126222in}}%
\pgfpathlineto{\pgfqpoint{2.644603in}{1.131634in}}%
\pgfpathlineto{\pgfqpoint{2.647609in}{1.142718in}}%
\pgfpathlineto{\pgfqpoint{2.654207in}{1.139998in}}%
\pgfpathlineto{\pgfqpoint{2.658319in}{1.135059in}}%
\pgfpathlineto{\pgfqpoint{2.663369in}{1.133914in}}%
\pgfpathlineto{\pgfqpoint{2.665522in}{1.153025in}}%
\pgfpathlineto{\pgfqpoint{2.679314in}{1.148103in}}%
\pgfpathlineto{\pgfqpoint{2.681074in}{1.143862in}}%
\pgfpathlineto{\pgfqpoint{2.685472in}{1.142255in}}%
\pgfpathlineto{\pgfqpoint{2.692040in}{1.154397in}}%
\pgfpathlineto{\pgfqpoint{2.697515in}{1.156406in}}%
\pgfpathlineto{\pgfqpoint{2.701679in}{1.149182in}}%
\pgfpathlineto{\pgfqpoint{2.706818in}{1.146620in}}%
\pgfpathlineto{\pgfqpoint{2.713398in}{1.148167in}}%
\pgfpathlineto{\pgfqpoint{2.716860in}{1.144435in}}%
\pgfpathlineto{\pgfqpoint{2.723470in}{1.137350in}}%
\pgfpathlineto{\pgfqpoint{2.724698in}{1.127666in}}%
\pgfpathlineto{\pgfqpoint{2.713382in}{1.123521in}}%
\pgfpathlineto{\pgfqpoint{2.716713in}{1.110416in}}%
\pgfpathlineto{\pgfqpoint{2.711789in}{1.111134in}}%
\pgfpathlineto{\pgfqpoint{2.704080in}{1.100551in}}%
\pgfpathlineto{\pgfqpoint{2.713742in}{1.098160in}}%
\pgfpathlineto{\pgfqpoint{2.716714in}{1.091684in}}%
\pgfpathlineto{\pgfqpoint{2.695089in}{1.071478in}}%
\pgfpathclose%
\pgfusepath{fill}%
\end{pgfscope}%
\begin{pgfscope}%
\pgfpathrectangle{\pgfqpoint{0.100000in}{0.100000in}}{\pgfqpoint{3.007045in}{1.925000in}}%
\pgfusepath{clip}%
\pgfsetbuttcap%
\pgfsetmiterjoin%
\definecolor{currentfill}{rgb}{0.713772,0.885675,0.688074}%
\pgfsetfillcolor{currentfill}%
\pgfsetlinewidth{0.000000pt}%
\definecolor{currentstroke}{rgb}{0.000000,0.000000,0.000000}%
\pgfsetstrokecolor{currentstroke}%
\pgfsetstrokeopacity{0.000000}%
\pgfsetdash{}{0pt}%
\pgfpathmoveto{\pgfqpoint{1.681563in}{0.619854in}}%
\pgfpathlineto{\pgfqpoint{1.668209in}{0.612458in}}%
\pgfpathlineto{\pgfqpoint{1.654308in}{0.604797in}}%
\pgfpathlineto{\pgfqpoint{1.642936in}{0.624588in}}%
\pgfpathlineto{\pgfqpoint{1.639177in}{0.622324in}}%
\pgfpathlineto{\pgfqpoint{1.635041in}{0.627484in}}%
\pgfpathlineto{\pgfqpoint{1.620418in}{0.653867in}}%
\pgfpathlineto{\pgfqpoint{1.616016in}{0.651438in}}%
\pgfpathlineto{\pgfqpoint{1.603281in}{0.674775in}}%
\pgfpathlineto{\pgfqpoint{1.611326in}{0.679134in}}%
\pgfpathlineto{\pgfqpoint{1.624705in}{0.686669in}}%
\pgfpathlineto{\pgfqpoint{1.628904in}{0.682713in}}%
\pgfpathlineto{\pgfqpoint{1.633343in}{0.684748in}}%
\pgfpathlineto{\pgfqpoint{1.655346in}{0.689754in}}%
\pgfpathlineto{\pgfqpoint{1.663204in}{0.675422in}}%
\pgfpathlineto{\pgfqpoint{1.665733in}{0.676821in}}%
\pgfpathlineto{\pgfqpoint{1.675361in}{0.659431in}}%
\pgfpathlineto{\pgfqpoint{1.663193in}{0.652637in}}%
\pgfpathlineto{\pgfqpoint{1.681563in}{0.619854in}}%
\pgfpathclose%
\pgfusepath{fill}%
\end{pgfscope}%
\begin{pgfscope}%
\pgfpathrectangle{\pgfqpoint{0.100000in}{0.100000in}}{\pgfqpoint{3.007045in}{1.925000in}}%
\pgfusepath{clip}%
\pgfsetbuttcap%
\pgfsetmiterjoin%
\definecolor{currentfill}{rgb}{0.634295,0.852211,0.610688}%
\pgfsetfillcolor{currentfill}%
\pgfsetlinewidth{0.000000pt}%
\definecolor{currentstroke}{rgb}{0.000000,0.000000,0.000000}%
\pgfsetstrokecolor{currentstroke}%
\pgfsetstrokeopacity{0.000000}%
\pgfsetdash{}{0pt}%
\pgfpathmoveto{\pgfqpoint{2.498857in}{1.196503in}}%
\pgfpathlineto{\pgfqpoint{2.494442in}{1.196721in}}%
\pgfpathlineto{\pgfqpoint{2.487686in}{1.206566in}}%
\pgfpathlineto{\pgfqpoint{2.488733in}{1.208153in}}%
\pgfpathlineto{\pgfqpoint{2.500713in}{1.220454in}}%
\pgfpathlineto{\pgfqpoint{2.505896in}{1.220258in}}%
\pgfpathlineto{\pgfqpoint{2.507751in}{1.225110in}}%
\pgfpathlineto{\pgfqpoint{2.505584in}{1.228741in}}%
\pgfpathlineto{\pgfqpoint{2.510163in}{1.234931in}}%
\pgfpathlineto{\pgfqpoint{2.507767in}{1.239655in}}%
\pgfpathlineto{\pgfqpoint{2.554976in}{1.247539in}}%
\pgfpathlineto{\pgfqpoint{2.560339in}{1.213292in}}%
\pgfpathlineto{\pgfqpoint{2.549426in}{1.216536in}}%
\pgfpathlineto{\pgfqpoint{2.543719in}{1.212838in}}%
\pgfpathlineto{\pgfqpoint{2.538153in}{1.216728in}}%
\pgfpathlineto{\pgfqpoint{2.532780in}{1.212166in}}%
\pgfpathlineto{\pgfqpoint{2.525644in}{1.210611in}}%
\pgfpathlineto{\pgfqpoint{2.523965in}{1.201691in}}%
\pgfpathlineto{\pgfqpoint{2.520381in}{1.200511in}}%
\pgfpathlineto{\pgfqpoint{2.505727in}{1.202402in}}%
\pgfpathlineto{\pgfqpoint{2.498857in}{1.196503in}}%
\pgfpathclose%
\pgfusepath{fill}%
\end{pgfscope}%
\begin{pgfscope}%
\pgfpathrectangle{\pgfqpoint{0.100000in}{0.100000in}}{\pgfqpoint{3.007045in}{1.925000in}}%
\pgfusepath{clip}%
\pgfsetbuttcap%
\pgfsetmiterjoin%
\definecolor{currentfill}{rgb}{0.439216,0.760938,0.455056}%
\pgfsetfillcolor{currentfill}%
\pgfsetlinewidth{0.000000pt}%
\definecolor{currentstroke}{rgb}{0.000000,0.000000,0.000000}%
\pgfsetstrokecolor{currentstroke}%
\pgfsetstrokeopacity{0.000000}%
\pgfsetdash{}{0pt}%
\pgfpathmoveto{\pgfqpoint{2.677192in}{1.366767in}}%
\pgfpathlineto{\pgfqpoint{2.679466in}{1.361047in}}%
\pgfpathlineto{\pgfqpoint{2.686221in}{1.354519in}}%
\pgfpathlineto{\pgfqpoint{2.688600in}{1.346784in}}%
\pgfpathlineto{\pgfqpoint{2.692918in}{1.348382in}}%
\pgfpathlineto{\pgfqpoint{2.694956in}{1.344848in}}%
\pgfpathlineto{\pgfqpoint{2.680799in}{1.334029in}}%
\pgfpathlineto{\pgfqpoint{2.670437in}{1.328641in}}%
\pgfpathlineto{\pgfqpoint{2.664036in}{1.332084in}}%
\pgfpathlineto{\pgfqpoint{2.648248in}{1.330436in}}%
\pgfpathlineto{\pgfqpoint{2.645950in}{1.339764in}}%
\pgfpathlineto{\pgfqpoint{2.654171in}{1.354581in}}%
\pgfpathlineto{\pgfqpoint{2.661650in}{1.359044in}}%
\pgfpathlineto{\pgfqpoint{2.661755in}{1.363168in}}%
\pgfpathlineto{\pgfqpoint{2.669738in}{1.366595in}}%
\pgfpathlineto{\pgfqpoint{2.677192in}{1.366767in}}%
\pgfpathclose%
\pgfusepath{fill}%
\end{pgfscope}%
\begin{pgfscope}%
\pgfpathrectangle{\pgfqpoint{0.100000in}{0.100000in}}{\pgfqpoint{3.007045in}{1.925000in}}%
\pgfusepath{clip}%
\pgfsetbuttcap%
\pgfsetmiterjoin%
\definecolor{currentfill}{rgb}{0.439216,0.760938,0.455056}%
\pgfsetfillcolor{currentfill}%
\pgfsetlinewidth{0.000000pt}%
\definecolor{currentstroke}{rgb}{0.000000,0.000000,0.000000}%
\pgfsetstrokecolor{currentstroke}%
\pgfsetstrokeopacity{0.000000}%
\pgfsetdash{}{0pt}%
\pgfpathmoveto{\pgfqpoint{2.141224in}{0.664018in}}%
\pgfpathlineto{\pgfqpoint{2.139795in}{0.676959in}}%
\pgfpathlineto{\pgfqpoint{2.141241in}{0.722568in}}%
\pgfpathlineto{\pgfqpoint{2.141820in}{0.745927in}}%
\pgfpathlineto{\pgfqpoint{2.141932in}{0.750061in}}%
\pgfpathlineto{\pgfqpoint{2.151248in}{0.751063in}}%
\pgfpathlineto{\pgfqpoint{2.154101in}{0.754354in}}%
\pgfpathlineto{\pgfqpoint{2.158791in}{0.754738in}}%
\pgfpathlineto{\pgfqpoint{2.163129in}{0.759429in}}%
\pgfpathlineto{\pgfqpoint{2.168941in}{0.762907in}}%
\pgfpathlineto{\pgfqpoint{2.169944in}{0.753927in}}%
\pgfpathlineto{\pgfqpoint{2.176652in}{0.754493in}}%
\pgfpathlineto{\pgfqpoint{2.175967in}{0.749888in}}%
\pgfpathlineto{\pgfqpoint{2.172078in}{0.746133in}}%
\pgfpathlineto{\pgfqpoint{2.171759in}{0.737367in}}%
\pgfpathlineto{\pgfqpoint{2.174689in}{0.726922in}}%
\pgfpathlineto{\pgfqpoint{2.173613in}{0.721470in}}%
\pgfpathlineto{\pgfqpoint{2.178513in}{0.718958in}}%
\pgfpathlineto{\pgfqpoint{2.189892in}{0.719957in}}%
\pgfpathlineto{\pgfqpoint{2.190855in}{0.708411in}}%
\pgfpathlineto{\pgfqpoint{2.193997in}{0.705851in}}%
\pgfpathlineto{\pgfqpoint{2.191190in}{0.704948in}}%
\pgfpathlineto{\pgfqpoint{2.191937in}{0.696888in}}%
\pgfpathlineto{\pgfqpoint{2.186226in}{0.696389in}}%
\pgfpathlineto{\pgfqpoint{2.186978in}{0.688005in}}%
\pgfpathlineto{\pgfqpoint{2.184529in}{0.686866in}}%
\pgfpathlineto{\pgfqpoint{2.161771in}{0.684999in}}%
\pgfpathlineto{\pgfqpoint{2.156603in}{0.673169in}}%
\pgfpathlineto{\pgfqpoint{2.161161in}{0.669683in}}%
\pgfpathlineto{\pgfqpoint{2.162408in}{0.665683in}}%
\pgfpathlineto{\pgfqpoint{2.141224in}{0.664018in}}%
\pgfpathclose%
\pgfusepath{fill}%
\end{pgfscope}%
\begin{pgfscope}%
\pgfpathrectangle{\pgfqpoint{0.100000in}{0.100000in}}{\pgfqpoint{3.007045in}{1.925000in}}%
\pgfusepath{clip}%
\pgfsetbuttcap%
\pgfsetmiterjoin%
\definecolor{currentfill}{rgb}{0.607151,0.839677,0.587928}%
\pgfsetfillcolor{currentfill}%
\pgfsetlinewidth{0.000000pt}%
\definecolor{currentstroke}{rgb}{0.000000,0.000000,0.000000}%
\pgfsetstrokecolor{currentstroke}%
\pgfsetstrokeopacity{0.000000}%
\pgfsetdash{}{0pt}%
\pgfpathmoveto{\pgfqpoint{1.984723in}{0.815682in}}%
\pgfpathlineto{\pgfqpoint{1.980611in}{0.811051in}}%
\pgfpathlineto{\pgfqpoint{1.984377in}{0.804898in}}%
\pgfpathlineto{\pgfqpoint{1.974688in}{0.804739in}}%
\pgfpathlineto{\pgfqpoint{1.967563in}{0.808897in}}%
\pgfpathlineto{\pgfqpoint{1.966388in}{0.815925in}}%
\pgfpathlineto{\pgfqpoint{1.962924in}{0.816378in}}%
\pgfpathlineto{\pgfqpoint{1.961337in}{0.821725in}}%
\pgfpathlineto{\pgfqpoint{1.951971in}{0.821985in}}%
\pgfpathlineto{\pgfqpoint{1.951736in}{0.838426in}}%
\pgfpathlineto{\pgfqpoint{1.955232in}{0.842505in}}%
\pgfpathlineto{\pgfqpoint{1.955013in}{0.847952in}}%
\pgfpathlineto{\pgfqpoint{1.951519in}{0.850943in}}%
\pgfpathlineto{\pgfqpoint{1.951312in}{0.867300in}}%
\pgfpathlineto{\pgfqpoint{1.945345in}{0.867212in}}%
\pgfpathlineto{\pgfqpoint{1.945167in}{0.874395in}}%
\pgfpathlineto{\pgfqpoint{1.950043in}{0.876695in}}%
\pgfpathlineto{\pgfqpoint{1.956926in}{0.874466in}}%
\pgfpathlineto{\pgfqpoint{1.956817in}{0.878848in}}%
\pgfpathlineto{\pgfqpoint{1.963244in}{0.878965in}}%
\pgfpathlineto{\pgfqpoint{1.968811in}{0.876703in}}%
\pgfpathlineto{\pgfqpoint{1.972810in}{0.872450in}}%
\pgfpathlineto{\pgfqpoint{1.980505in}{0.873918in}}%
\pgfpathlineto{\pgfqpoint{1.980621in}{0.867935in}}%
\pgfpathlineto{\pgfqpoint{1.983593in}{0.865113in}}%
\pgfpathlineto{\pgfqpoint{1.983824in}{0.856375in}}%
\pgfpathlineto{\pgfqpoint{1.986708in}{0.856447in}}%
\pgfpathlineto{\pgfqpoint{1.987166in}{0.839058in}}%
\pgfpathlineto{\pgfqpoint{1.990126in}{0.839097in}}%
\pgfpathlineto{\pgfqpoint{1.990401in}{0.830295in}}%
\pgfpathlineto{\pgfqpoint{1.987469in}{0.830190in}}%
\pgfpathlineto{\pgfqpoint{1.986497in}{0.818543in}}%
\pgfpathlineto{\pgfqpoint{1.984723in}{0.815682in}}%
\pgfpathclose%
\pgfusepath{fill}%
\end{pgfscope}%
\begin{pgfscope}%
\pgfpathrectangle{\pgfqpoint{0.100000in}{0.100000in}}{\pgfqpoint{3.007045in}{1.925000in}}%
\pgfusepath{clip}%
\pgfsetbuttcap%
\pgfsetmiterjoin%
\definecolor{currentfill}{rgb}{0.601615,0.837093,0.583376}%
\pgfsetfillcolor{currentfill}%
\pgfsetlinewidth{0.000000pt}%
\definecolor{currentstroke}{rgb}{0.000000,0.000000,0.000000}%
\pgfsetstrokecolor{currentstroke}%
\pgfsetstrokeopacity{0.000000}%
\pgfsetdash{}{0pt}%
\pgfpathmoveto{\pgfqpoint{1.408374in}{1.098989in}}%
\pgfpathlineto{\pgfqpoint{1.372398in}{1.101802in}}%
\pgfpathlineto{\pgfqpoint{1.333267in}{1.105178in}}%
\pgfpathlineto{\pgfqpoint{1.334664in}{1.121680in}}%
\pgfpathlineto{\pgfqpoint{1.351745in}{1.120688in}}%
\pgfpathlineto{\pgfqpoint{1.353166in}{1.131470in}}%
\pgfpathlineto{\pgfqpoint{1.355807in}{1.160158in}}%
\pgfpathlineto{\pgfqpoint{1.358836in}{1.188865in}}%
\pgfpathlineto{\pgfqpoint{1.414767in}{1.184726in}}%
\pgfpathlineto{\pgfqpoint{1.414739in}{1.184340in}}%
\pgfpathlineto{\pgfqpoint{1.410528in}{1.127149in}}%
\pgfpathlineto{\pgfqpoint{1.408374in}{1.098989in}}%
\pgfpathclose%
\pgfusepath{fill}%
\end{pgfscope}%
\begin{pgfscope}%
\pgfpathrectangle{\pgfqpoint{0.100000in}{0.100000in}}{\pgfqpoint{3.007045in}{1.925000in}}%
\pgfusepath{clip}%
\pgfsetbuttcap%
\pgfsetmiterjoin%
\definecolor{currentfill}{rgb}{0.681046,0.871895,0.656209}%
\pgfsetfillcolor{currentfill}%
\pgfsetlinewidth{0.000000pt}%
\definecolor{currentstroke}{rgb}{0.000000,0.000000,0.000000}%
\pgfsetstrokecolor{currentstroke}%
\pgfsetstrokeopacity{0.000000}%
\pgfsetdash{}{0pt}%
\pgfpathmoveto{\pgfqpoint{2.013407in}{0.843839in}}%
\pgfpathlineto{\pgfqpoint{2.013158in}{0.836265in}}%
\pgfpathlineto{\pgfqpoint{2.003733in}{0.832951in}}%
\pgfpathlineto{\pgfqpoint{2.003512in}{0.827202in}}%
\pgfpathlineto{\pgfqpoint{1.993574in}{0.816135in}}%
\pgfpathlineto{\pgfqpoint{1.993621in}{0.816052in}}%
\pgfpathlineto{\pgfqpoint{1.984723in}{0.815682in}}%
\pgfpathlineto{\pgfqpoint{1.986497in}{0.818543in}}%
\pgfpathlineto{\pgfqpoint{1.987469in}{0.830190in}}%
\pgfpathlineto{\pgfqpoint{1.990401in}{0.830295in}}%
\pgfpathlineto{\pgfqpoint{1.990126in}{0.839097in}}%
\pgfpathlineto{\pgfqpoint{1.987166in}{0.839058in}}%
\pgfpathlineto{\pgfqpoint{1.986708in}{0.856447in}}%
\pgfpathlineto{\pgfqpoint{1.983824in}{0.856375in}}%
\pgfpathlineto{\pgfqpoint{1.983593in}{0.865113in}}%
\pgfpathlineto{\pgfqpoint{1.980621in}{0.867935in}}%
\pgfpathlineto{\pgfqpoint{1.980505in}{0.873918in}}%
\pgfpathlineto{\pgfqpoint{1.986161in}{0.874058in}}%
\pgfpathlineto{\pgfqpoint{1.985207in}{0.903228in}}%
\pgfpathlineto{\pgfqpoint{2.008210in}{0.904668in}}%
\pgfpathlineto{\pgfqpoint{2.025623in}{0.904867in}}%
\pgfpathlineto{\pgfqpoint{2.033514in}{0.905154in}}%
\pgfpathlineto{\pgfqpoint{2.037206in}{0.901885in}}%
\pgfpathlineto{\pgfqpoint{2.034157in}{0.893004in}}%
\pgfpathlineto{\pgfqpoint{2.037839in}{0.890956in}}%
\pgfpathlineto{\pgfqpoint{2.036782in}{0.884180in}}%
\pgfpathlineto{\pgfqpoint{2.032840in}{0.884060in}}%
\pgfpathlineto{\pgfqpoint{2.030962in}{0.877859in}}%
\pgfpathlineto{\pgfqpoint{2.026701in}{0.874367in}}%
\pgfpathlineto{\pgfqpoint{2.029442in}{0.869945in}}%
\pgfpathlineto{\pgfqpoint{2.020371in}{0.864609in}}%
\pgfpathlineto{\pgfqpoint{2.018227in}{0.853487in}}%
\pgfpathlineto{\pgfqpoint{2.012143in}{0.849788in}}%
\pgfpathlineto{\pgfqpoint{2.013407in}{0.843839in}}%
\pgfpathclose%
\pgfusepath{fill}%
\end{pgfscope}%
\begin{pgfscope}%
\pgfpathrectangle{\pgfqpoint{0.100000in}{0.100000in}}{\pgfqpoint{3.007045in}{1.925000in}}%
\pgfusepath{clip}%
\pgfsetbuttcap%
\pgfsetmiterjoin%
\definecolor{currentfill}{rgb}{0.623760,0.847428,0.601584}%
\pgfsetfillcolor{currentfill}%
\pgfsetlinewidth{0.000000pt}%
\definecolor{currentstroke}{rgb}{0.000000,0.000000,0.000000}%
\pgfsetstrokecolor{currentstroke}%
\pgfsetstrokeopacity{0.000000}%
\pgfsetdash{}{0pt}%
\pgfpathmoveto{\pgfqpoint{2.283647in}{1.145438in}}%
\pgfpathlineto{\pgfqpoint{2.275093in}{1.139834in}}%
\pgfpathlineto{\pgfqpoint{2.273431in}{1.154396in}}%
\pgfpathlineto{\pgfqpoint{2.261118in}{1.153025in}}%
\pgfpathlineto{\pgfqpoint{2.259410in}{1.171566in}}%
\pgfpathlineto{\pgfqpoint{2.253429in}{1.166775in}}%
\pgfpathlineto{\pgfqpoint{2.247318in}{1.166002in}}%
\pgfpathlineto{\pgfqpoint{2.245949in}{1.180376in}}%
\pgfpathlineto{\pgfqpoint{2.248039in}{1.187423in}}%
\pgfpathlineto{\pgfqpoint{2.264834in}{1.189201in}}%
\pgfpathlineto{\pgfqpoint{2.264298in}{1.193927in}}%
\pgfpathlineto{\pgfqpoint{2.288706in}{1.196382in}}%
\pgfpathlineto{\pgfqpoint{2.291530in}{1.169095in}}%
\pgfpathlineto{\pgfqpoint{2.288006in}{1.165183in}}%
\pgfpathlineto{\pgfqpoint{2.291496in}{1.161362in}}%
\pgfpathlineto{\pgfqpoint{2.289964in}{1.156159in}}%
\pgfpathlineto{\pgfqpoint{2.294975in}{1.154537in}}%
\pgfpathlineto{\pgfqpoint{2.290373in}{1.148367in}}%
\pgfpathlineto{\pgfqpoint{2.283647in}{1.145438in}}%
\pgfpathclose%
\pgfusepath{fill}%
\end{pgfscope}%
\begin{pgfscope}%
\pgfpathrectangle{\pgfqpoint{0.100000in}{0.100000in}}{\pgfqpoint{3.007045in}{1.925000in}}%
\pgfusepath{clip}%
\pgfsetbuttcap%
\pgfsetmiterjoin%
\definecolor{currentfill}{rgb}{0.607151,0.839677,0.587928}%
\pgfsetfillcolor{currentfill}%
\pgfsetlinewidth{0.000000pt}%
\definecolor{currentstroke}{rgb}{0.000000,0.000000,0.000000}%
\pgfsetstrokecolor{currentstroke}%
\pgfsetstrokeopacity{0.000000}%
\pgfsetdash{}{0pt}%
\pgfpathmoveto{\pgfqpoint{1.440696in}{0.803294in}}%
\pgfpathlineto{\pgfqpoint{1.469347in}{0.801604in}}%
\pgfpathlineto{\pgfqpoint{1.495136in}{0.800124in}}%
\pgfpathlineto{\pgfqpoint{1.514372in}{0.798412in}}%
\pgfpathlineto{\pgfqpoint{1.525327in}{0.794045in}}%
\pgfpathlineto{\pgfqpoint{1.526290in}{0.791703in}}%
\pgfpathlineto{\pgfqpoint{1.525314in}{0.769593in}}%
\pgfpathlineto{\pgfqpoint{1.523941in}{0.740360in}}%
\pgfpathlineto{\pgfqpoint{1.486605in}{0.742533in}}%
\pgfpathlineto{\pgfqpoint{1.465814in}{0.743950in}}%
\pgfpathlineto{\pgfqpoint{1.467658in}{0.772689in}}%
\pgfpathlineto{\pgfqpoint{1.438908in}{0.774476in}}%
\pgfpathlineto{\pgfqpoint{1.440696in}{0.803294in}}%
\pgfpathclose%
\pgfusepath{fill}%
\end{pgfscope}%
\begin{pgfscope}%
\pgfpathrectangle{\pgfqpoint{0.100000in}{0.100000in}}{\pgfqpoint{3.007045in}{1.925000in}}%
\pgfusepath{clip}%
\pgfsetbuttcap%
\pgfsetmiterjoin%
\definecolor{currentfill}{rgb}{0.585006,0.829343,0.569719}%
\pgfsetfillcolor{currentfill}%
\pgfsetlinewidth{0.000000pt}%
\definecolor{currentstroke}{rgb}{0.000000,0.000000,0.000000}%
\pgfsetstrokecolor{currentstroke}%
\pgfsetstrokeopacity{0.000000}%
\pgfsetdash{}{0pt}%
\pgfpathmoveto{\pgfqpoint{1.963991in}{0.607879in}}%
\pgfpathlineto{\pgfqpoint{1.958935in}{0.610678in}}%
\pgfpathlineto{\pgfqpoint{1.958991in}{0.616692in}}%
\pgfpathlineto{\pgfqpoint{1.961233in}{0.620328in}}%
\pgfpathlineto{\pgfqpoint{1.952131in}{0.625210in}}%
\pgfpathlineto{\pgfqpoint{1.949979in}{0.638422in}}%
\pgfpathlineto{\pgfqpoint{1.954469in}{0.643065in}}%
\pgfpathlineto{\pgfqpoint{1.959044in}{0.656761in}}%
\pgfpathlineto{\pgfqpoint{1.962071in}{0.656269in}}%
\pgfpathlineto{\pgfqpoint{1.965330in}{0.661044in}}%
\pgfpathlineto{\pgfqpoint{1.965268in}{0.666642in}}%
\pgfpathlineto{\pgfqpoint{1.969490in}{0.679240in}}%
\pgfpathlineto{\pgfqpoint{1.969121in}{0.688091in}}%
\pgfpathlineto{\pgfqpoint{1.978768in}{0.688413in}}%
\pgfpathlineto{\pgfqpoint{1.989120in}{0.691955in}}%
\pgfpathlineto{\pgfqpoint{1.987667in}{0.684255in}}%
\pgfpathlineto{\pgfqpoint{1.995075in}{0.682326in}}%
\pgfpathlineto{\pgfqpoint{1.993029in}{0.676683in}}%
\pgfpathlineto{\pgfqpoint{1.986876in}{0.672737in}}%
\pgfpathlineto{\pgfqpoint{1.986197in}{0.669162in}}%
\pgfpathlineto{\pgfqpoint{1.978327in}{0.664514in}}%
\pgfpathlineto{\pgfqpoint{1.981608in}{0.653144in}}%
\pgfpathlineto{\pgfqpoint{1.989789in}{0.649517in}}%
\pgfpathlineto{\pgfqpoint{2.013321in}{0.650638in}}%
\pgfpathlineto{\pgfqpoint{2.019201in}{0.650937in}}%
\pgfpathlineto{\pgfqpoint{2.020031in}{0.633536in}}%
\pgfpathlineto{\pgfqpoint{1.998230in}{0.632422in}}%
\pgfpathlineto{\pgfqpoint{1.993008in}{0.630363in}}%
\pgfpathlineto{\pgfqpoint{1.985698in}{0.633433in}}%
\pgfpathlineto{\pgfqpoint{1.972763in}{0.629511in}}%
\pgfpathlineto{\pgfqpoint{1.971855in}{0.623291in}}%
\pgfpathlineto{\pgfqpoint{1.967749in}{0.623734in}}%
\pgfpathlineto{\pgfqpoint{1.964384in}{0.616831in}}%
\pgfpathlineto{\pgfqpoint{1.967968in}{0.612514in}}%
\pgfpathlineto{\pgfqpoint{1.963991in}{0.607879in}}%
\pgfpathclose%
\pgfusepath{fill}%
\end{pgfscope}%
\begin{pgfscope}%
\pgfpathrectangle{\pgfqpoint{0.100000in}{0.100000in}}{\pgfqpoint{3.007045in}{1.925000in}}%
\pgfusepath{clip}%
\pgfsetbuttcap%
\pgfsetmiterjoin%
\definecolor{currentfill}{rgb}{0.518570,0.798339,0.515094}%
\pgfsetfillcolor{currentfill}%
\pgfsetlinewidth{0.000000pt}%
\definecolor{currentstroke}{rgb}{0.000000,0.000000,0.000000}%
\pgfsetstrokecolor{currentstroke}%
\pgfsetstrokeopacity{0.000000}%
\pgfsetdash{}{0pt}%
\pgfpathmoveto{\pgfqpoint{1.990535in}{1.152936in}}%
\pgfpathlineto{\pgfqpoint{1.967023in}{1.152094in}}%
\pgfpathlineto{\pgfqpoint{1.967135in}{1.146275in}}%
\pgfpathlineto{\pgfqpoint{1.963320in}{1.146179in}}%
\pgfpathlineto{\pgfqpoint{1.955717in}{1.146050in}}%
\pgfpathlineto{\pgfqpoint{1.953772in}{1.157587in}}%
\pgfpathlineto{\pgfqpoint{1.939480in}{1.158394in}}%
\pgfpathlineto{\pgfqpoint{1.938875in}{1.179366in}}%
\pgfpathlineto{\pgfqpoint{1.932193in}{1.179188in}}%
\pgfpathlineto{\pgfqpoint{1.931854in}{1.198270in}}%
\pgfpathlineto{\pgfqpoint{1.926181in}{1.198155in}}%
\pgfpathlineto{\pgfqpoint{1.925598in}{1.218425in}}%
\pgfpathlineto{\pgfqpoint{1.948406in}{1.218688in}}%
\pgfpathlineto{\pgfqpoint{1.948122in}{1.215494in}}%
\pgfpathlineto{\pgfqpoint{1.977971in}{1.216245in}}%
\pgfpathlineto{\pgfqpoint{1.978770in}{1.193340in}}%
\pgfpathlineto{\pgfqpoint{1.996210in}{1.193778in}}%
\pgfpathlineto{\pgfqpoint{1.992899in}{1.184317in}}%
\pgfpathlineto{\pgfqpoint{1.997146in}{1.174029in}}%
\pgfpathlineto{\pgfqpoint{1.995501in}{1.164493in}}%
\pgfpathlineto{\pgfqpoint{1.979066in}{1.163974in}}%
\pgfpathlineto{\pgfqpoint{1.986624in}{1.158384in}}%
\pgfpathlineto{\pgfqpoint{1.990535in}{1.152936in}}%
\pgfpathclose%
\pgfusepath{fill}%
\end{pgfscope}%
\begin{pgfscope}%
\pgfpathrectangle{\pgfqpoint{0.100000in}{0.100000in}}{\pgfqpoint{3.007045in}{1.925000in}}%
\pgfusepath{clip}%
\pgfsetbuttcap%
\pgfsetmiterjoin%
\definecolor{currentfill}{rgb}{0.357647,0.720953,0.415071}%
\pgfsetfillcolor{currentfill}%
\pgfsetlinewidth{0.000000pt}%
\definecolor{currentstroke}{rgb}{0.000000,0.000000,0.000000}%
\pgfsetstrokecolor{currentstroke}%
\pgfsetstrokeopacity{0.000000}%
\pgfsetdash{}{0pt}%
\pgfpathmoveto{\pgfqpoint{0.856506in}{1.901711in}}%
\pgfpathlineto{\pgfqpoint{0.848850in}{1.869247in}}%
\pgfpathlineto{\pgfqpoint{0.817161in}{1.876758in}}%
\pgfpathlineto{\pgfqpoint{0.822653in}{1.899281in}}%
\pgfpathlineto{\pgfqpoint{0.812070in}{1.901824in}}%
\pgfpathlineto{\pgfqpoint{0.814525in}{1.911735in}}%
\pgfpathlineto{\pgfqpoint{0.856506in}{1.901711in}}%
\pgfpathclose%
\pgfusepath{fill}%
\end{pgfscope}%
\begin{pgfscope}%
\pgfpathrectangle{\pgfqpoint{0.100000in}{0.100000in}}{\pgfqpoint{3.007045in}{1.925000in}}%
\pgfusepath{clip}%
\pgfsetbuttcap%
\pgfsetmiterjoin%
\definecolor{currentfill}{rgb}{0.657670,0.862053,0.633449}%
\pgfsetfillcolor{currentfill}%
\pgfsetlinewidth{0.000000pt}%
\definecolor{currentstroke}{rgb}{0.000000,0.000000,0.000000}%
\pgfsetstrokecolor{currentstroke}%
\pgfsetstrokeopacity{0.000000}%
\pgfsetdash{}{0pt}%
\pgfpathmoveto{\pgfqpoint{2.349537in}{1.370183in}}%
\pgfpathlineto{\pgfqpoint{2.343612in}{1.373779in}}%
\pgfpathlineto{\pgfqpoint{2.320460in}{1.370056in}}%
\pgfpathlineto{\pgfqpoint{2.303185in}{1.367262in}}%
\pgfpathlineto{\pgfqpoint{2.299946in}{1.390294in}}%
\pgfpathlineto{\pgfqpoint{2.296340in}{1.413257in}}%
\pgfpathlineto{\pgfqpoint{2.286455in}{1.412044in}}%
\pgfpathlineto{\pgfqpoint{2.283625in}{1.434436in}}%
\pgfpathlineto{\pgfqpoint{2.283546in}{1.435118in}}%
\pgfpathlineto{\pgfqpoint{2.304504in}{1.437926in}}%
\pgfpathlineto{\pgfqpoint{2.303651in}{1.443639in}}%
\pgfpathlineto{\pgfqpoint{2.326140in}{1.446645in}}%
\pgfpathlineto{\pgfqpoint{2.331232in}{1.448101in}}%
\pgfpathlineto{\pgfqpoint{2.330370in}{1.453796in}}%
\pgfpathlineto{\pgfqpoint{2.342275in}{1.455683in}}%
\pgfpathlineto{\pgfqpoint{2.338611in}{1.478397in}}%
\pgfpathlineto{\pgfqpoint{2.362681in}{1.482619in}}%
\pgfpathlineto{\pgfqpoint{2.364694in}{1.474272in}}%
\pgfpathlineto{\pgfqpoint{2.368291in}{1.465466in}}%
\pgfpathlineto{\pgfqpoint{2.371281in}{1.451948in}}%
\pgfpathlineto{\pgfqpoint{2.373711in}{1.444991in}}%
\pgfpathlineto{\pgfqpoint{2.377887in}{1.438966in}}%
\pgfpathlineto{\pgfqpoint{2.376582in}{1.425757in}}%
\pgfpathlineto{\pgfqpoint{2.377704in}{1.422799in}}%
\pgfpathlineto{\pgfqpoint{2.376418in}{1.412332in}}%
\pgfpathlineto{\pgfqpoint{2.367395in}{1.416634in}}%
\pgfpathlineto{\pgfqpoint{2.363048in}{1.413679in}}%
\pgfpathlineto{\pgfqpoint{2.360076in}{1.403017in}}%
\pgfpathlineto{\pgfqpoint{2.361096in}{1.399839in}}%
\pgfpathlineto{\pgfqpoint{2.359012in}{1.392852in}}%
\pgfpathlineto{\pgfqpoint{2.352893in}{1.389860in}}%
\pgfpathlineto{\pgfqpoint{2.350480in}{1.384293in}}%
\pgfpathlineto{\pgfqpoint{2.351493in}{1.374397in}}%
\pgfpathlineto{\pgfqpoint{2.349537in}{1.370183in}}%
\pgfpathclose%
\pgfusepath{fill}%
\end{pgfscope}%
\begin{pgfscope}%
\pgfpathrectangle{\pgfqpoint{0.100000in}{0.100000in}}{\pgfqpoint{3.007045in}{1.925000in}}%
\pgfusepath{clip}%
\pgfsetbuttcap%
\pgfsetmiterjoin%
\definecolor{currentfill}{rgb}{0.579469,0.826759,0.565167}%
\pgfsetfillcolor{currentfill}%
\pgfsetlinewidth{0.000000pt}%
\definecolor{currentstroke}{rgb}{0.000000,0.000000,0.000000}%
\pgfsetstrokecolor{currentstroke}%
\pgfsetstrokeopacity{0.000000}%
\pgfsetdash{}{0pt}%
\pgfpathmoveto{\pgfqpoint{0.652693in}{1.955469in}}%
\pgfpathlineto{\pgfqpoint{0.654305in}{1.933747in}}%
\pgfpathlineto{\pgfqpoint{0.650462in}{1.932235in}}%
\pgfpathlineto{\pgfqpoint{0.643432in}{1.933133in}}%
\pgfpathlineto{\pgfqpoint{0.623985in}{1.938935in}}%
\pgfpathlineto{\pgfqpoint{0.576681in}{1.953612in}}%
\pgfpathlineto{\pgfqpoint{0.577191in}{1.962609in}}%
\pgfpathlineto{\pgfqpoint{0.572821in}{1.961633in}}%
\pgfpathlineto{\pgfqpoint{0.569082in}{1.972907in}}%
\pgfpathlineto{\pgfqpoint{0.572639in}{1.979931in}}%
\pgfpathlineto{\pgfqpoint{0.614766in}{1.966556in}}%
\pgfpathlineto{\pgfqpoint{0.652693in}{1.955469in}}%
\pgfpathclose%
\pgfusepath{fill}%
\end{pgfscope}%
\begin{pgfscope}%
\pgfpathrectangle{\pgfqpoint{0.100000in}{0.100000in}}{\pgfqpoint{3.007045in}{1.925000in}}%
\pgfusepath{clip}%
\pgfsetbuttcap%
\pgfsetmiterjoin%
\definecolor{currentfill}{rgb}{0.676371,0.869927,0.651657}%
\pgfsetfillcolor{currentfill}%
\pgfsetlinewidth{0.000000pt}%
\definecolor{currentstroke}{rgb}{0.000000,0.000000,0.000000}%
\pgfsetstrokecolor{currentstroke}%
\pgfsetstrokeopacity{0.000000}%
\pgfsetdash{}{0pt}%
\pgfpathmoveto{\pgfqpoint{1.035972in}{1.058230in}}%
\pgfpathlineto{\pgfqpoint{1.025816in}{0.993300in}}%
\pgfpathlineto{\pgfqpoint{1.100725in}{0.982093in}}%
\pgfpathlineto{\pgfqpoint{1.117536in}{0.979664in}}%
\pgfpathlineto{\pgfqpoint{1.111260in}{0.934336in}}%
\pgfpathlineto{\pgfqpoint{1.049496in}{0.943259in}}%
\pgfpathlineto{\pgfqpoint{1.046069in}{0.920555in}}%
\pgfpathlineto{\pgfqpoint{1.015230in}{0.925309in}}%
\pgfpathlineto{\pgfqpoint{1.011345in}{0.900484in}}%
\pgfpathlineto{\pgfqpoint{1.003150in}{0.848168in}}%
\pgfpathlineto{\pgfqpoint{0.986809in}{0.850694in}}%
\pgfpathlineto{\pgfqpoint{0.978685in}{0.845988in}}%
\pgfpathlineto{\pgfqpoint{0.975090in}{0.841991in}}%
\pgfpathlineto{\pgfqpoint{0.971108in}{0.842140in}}%
\pgfpathlineto{\pgfqpoint{0.960834in}{0.834656in}}%
\pgfpathlineto{\pgfqpoint{0.954625in}{0.842378in}}%
\pgfpathlineto{\pgfqpoint{0.949199in}{0.843379in}}%
\pgfpathlineto{\pgfqpoint{0.953797in}{0.870982in}}%
\pgfpathlineto{\pgfqpoint{0.913272in}{0.877904in}}%
\pgfpathlineto{\pgfqpoint{0.916273in}{0.895048in}}%
\pgfpathlineto{\pgfqpoint{0.947453in}{1.073147in}}%
\pgfpathlineto{\pgfqpoint{0.986309in}{1.066160in}}%
\pgfpathlineto{\pgfqpoint{1.035972in}{1.058230in}}%
\pgfpathclose%
\pgfusepath{fill}%
\end{pgfscope}%
\begin{pgfscope}%
\pgfpathrectangle{\pgfqpoint{0.100000in}{0.100000in}}{\pgfqpoint{3.007045in}{1.925000in}}%
\pgfusepath{clip}%
\pgfsetbuttcap%
\pgfsetmiterjoin%
\definecolor{currentfill}{rgb}{0.320000,0.702499,0.396617}%
\pgfsetfillcolor{currentfill}%
\pgfsetlinewidth{0.000000pt}%
\definecolor{currentstroke}{rgb}{0.000000,0.000000,0.000000}%
\pgfsetstrokecolor{currentstroke}%
\pgfsetstrokeopacity{0.000000}%
\pgfsetdash{}{0pt}%
\pgfpathmoveto{\pgfqpoint{1.907975in}{1.607918in}}%
\pgfpathlineto{\pgfqpoint{1.930881in}{1.608680in}}%
\pgfpathlineto{\pgfqpoint{1.959683in}{1.609718in}}%
\pgfpathlineto{\pgfqpoint{1.960172in}{1.598201in}}%
\pgfpathlineto{\pgfqpoint{1.971604in}{1.598763in}}%
\pgfpathlineto{\pgfqpoint{1.973425in}{1.558760in}}%
\pgfpathlineto{\pgfqpoint{1.961960in}{1.558341in}}%
\pgfpathlineto{\pgfqpoint{1.962214in}{1.552559in}}%
\pgfpathlineto{\pgfqpoint{1.933521in}{1.551363in}}%
\pgfpathlineto{\pgfqpoint{1.933694in}{1.545701in}}%
\pgfpathlineto{\pgfqpoint{1.905032in}{1.544903in}}%
\pgfpathlineto{\pgfqpoint{1.904215in}{1.573411in}}%
\pgfpathlineto{\pgfqpoint{1.909946in}{1.573607in}}%
\pgfpathlineto{\pgfqpoint{1.909099in}{1.596409in}}%
\pgfpathlineto{\pgfqpoint{1.907975in}{1.607918in}}%
\pgfpathclose%
\pgfusepath{fill}%
\end{pgfscope}%
\begin{pgfscope}%
\pgfpathrectangle{\pgfqpoint{0.100000in}{0.100000in}}{\pgfqpoint{3.007045in}{1.925000in}}%
\pgfusepath{clip}%
\pgfsetbuttcap%
\pgfsetmiterjoin%
\definecolor{currentfill}{rgb}{0.562860,0.819008,0.551511}%
\pgfsetfillcolor{currentfill}%
\pgfsetlinewidth{0.000000pt}%
\definecolor{currentstroke}{rgb}{0.000000,0.000000,0.000000}%
\pgfsetstrokecolor{currentstroke}%
\pgfsetstrokeopacity{0.000000}%
\pgfsetdash{}{0pt}%
\pgfpathmoveto{\pgfqpoint{1.454083in}{0.686391in}}%
\pgfpathlineto{\pgfqpoint{1.425000in}{0.688410in}}%
\pgfpathlineto{\pgfqpoint{1.427496in}{0.717282in}}%
\pgfpathlineto{\pgfqpoint{1.429539in}{0.746331in}}%
\pgfpathlineto{\pgfqpoint{1.437038in}{0.746267in}}%
\pgfpathlineto{\pgfqpoint{1.438908in}{0.774476in}}%
\pgfpathlineto{\pgfqpoint{1.467658in}{0.772689in}}%
\pgfpathlineto{\pgfqpoint{1.465814in}{0.743950in}}%
\pgfpathlineto{\pgfqpoint{1.458235in}{0.744453in}}%
\pgfpathlineto{\pgfqpoint{1.454083in}{0.686391in}}%
\pgfpathclose%
\pgfusepath{fill}%
\end{pgfscope}%
\begin{pgfscope}%
\pgfpathrectangle{\pgfqpoint{0.100000in}{0.100000in}}{\pgfqpoint{3.007045in}{1.925000in}}%
\pgfusepath{clip}%
\pgfsetbuttcap%
\pgfsetmiterjoin%
\definecolor{currentfill}{rgb}{0.573933,0.824175,0.560615}%
\pgfsetfillcolor{currentfill}%
\pgfsetlinewidth{0.000000pt}%
\definecolor{currentstroke}{rgb}{0.000000,0.000000,0.000000}%
\pgfsetstrokecolor{currentstroke}%
\pgfsetstrokeopacity{0.000000}%
\pgfsetdash{}{0pt}%
\pgfpathmoveto{\pgfqpoint{2.079034in}{0.765390in}}%
\pgfpathlineto{\pgfqpoint{2.086624in}{0.765838in}}%
\pgfpathlineto{\pgfqpoint{2.087936in}{0.742436in}}%
\pgfpathlineto{\pgfqpoint{2.065136in}{0.740946in}}%
\pgfpathlineto{\pgfqpoint{2.065317in}{0.737960in}}%
\pgfpathlineto{\pgfqpoint{2.052237in}{0.736873in}}%
\pgfpathlineto{\pgfqpoint{2.031102in}{0.744404in}}%
\pgfpathlineto{\pgfqpoint{2.025625in}{0.749814in}}%
\pgfpathlineto{\pgfqpoint{2.031123in}{0.752696in}}%
\pgfpathlineto{\pgfqpoint{2.030481in}{0.763778in}}%
\pgfpathlineto{\pgfqpoint{2.033226in}{0.766351in}}%
\pgfpathlineto{\pgfqpoint{2.039047in}{0.766251in}}%
\pgfpathlineto{\pgfqpoint{2.062975in}{0.759697in}}%
\pgfpathlineto{\pgfqpoint{2.065349in}{0.764544in}}%
\pgfpathlineto{\pgfqpoint{2.079034in}{0.765390in}}%
\pgfpathclose%
\pgfusepath{fill}%
\end{pgfscope}%
\begin{pgfscope}%
\pgfpathrectangle{\pgfqpoint{0.100000in}{0.100000in}}{\pgfqpoint{3.007045in}{1.925000in}}%
\pgfusepath{clip}%
\pgfsetbuttcap%
\pgfsetmiterjoin%
\definecolor{currentfill}{rgb}{0.723122,0.889612,0.697178}%
\pgfsetfillcolor{currentfill}%
\pgfsetlinewidth{0.000000pt}%
\definecolor{currentstroke}{rgb}{0.000000,0.000000,0.000000}%
\pgfsetstrokecolor{currentstroke}%
\pgfsetstrokeopacity{0.000000}%
\pgfsetdash{}{0pt}%
\pgfpathmoveto{\pgfqpoint{1.973511in}{1.274071in}}%
\pgfpathlineto{\pgfqpoint{1.952591in}{1.273282in}}%
\pgfpathlineto{\pgfqpoint{1.952380in}{1.279020in}}%
\pgfpathlineto{\pgfqpoint{1.935244in}{1.278471in}}%
\pgfpathlineto{\pgfqpoint{1.923843in}{1.278117in}}%
\pgfpathlineto{\pgfqpoint{1.923061in}{1.300995in}}%
\pgfpathlineto{\pgfqpoint{1.928815in}{1.301237in}}%
\pgfpathlineto{\pgfqpoint{1.927902in}{1.324319in}}%
\pgfpathlineto{\pgfqpoint{1.950752in}{1.325057in}}%
\pgfpathlineto{\pgfqpoint{1.950534in}{1.330806in}}%
\pgfpathlineto{\pgfqpoint{1.973436in}{1.331705in}}%
\pgfpathlineto{\pgfqpoint{1.973965in}{1.320197in}}%
\pgfpathlineto{\pgfqpoint{1.974408in}{1.308715in}}%
\pgfpathlineto{\pgfqpoint{1.980128in}{1.308936in}}%
\pgfpathlineto{\pgfqpoint{1.980445in}{1.299472in}}%
\pgfpathlineto{\pgfqpoint{1.967742in}{1.296340in}}%
\pgfpathlineto{\pgfqpoint{1.964916in}{1.284730in}}%
\pgfpathlineto{\pgfqpoint{1.971270in}{1.279448in}}%
\pgfpathlineto{\pgfqpoint{1.973511in}{1.274071in}}%
\pgfpathclose%
\pgfusepath{fill}%
\end{pgfscope}%
\begin{pgfscope}%
\pgfpathrectangle{\pgfqpoint{0.100000in}{0.100000in}}{\pgfqpoint{3.007045in}{1.925000in}}%
\pgfusepath{clip}%
\pgfsetbuttcap%
\pgfsetmiterjoin%
\definecolor{currentfill}{rgb}{0.439216,0.760938,0.455056}%
\pgfsetfillcolor{currentfill}%
\pgfsetlinewidth{0.000000pt}%
\definecolor{currentstroke}{rgb}{0.000000,0.000000,0.000000}%
\pgfsetstrokecolor{currentstroke}%
\pgfsetstrokeopacity{0.000000}%
\pgfsetdash{}{0pt}%
\pgfpathmoveto{\pgfqpoint{1.765864in}{1.580217in}}%
\pgfpathlineto{\pgfqpoint{1.766199in}{1.556278in}}%
\pgfpathlineto{\pgfqpoint{1.720360in}{1.556564in}}%
\pgfpathlineto{\pgfqpoint{1.714573in}{1.556685in}}%
\pgfpathlineto{\pgfqpoint{1.714182in}{1.568146in}}%
\pgfpathlineto{\pgfqpoint{1.714476in}{1.591318in}}%
\pgfpathlineto{\pgfqpoint{1.714002in}{1.597084in}}%
\pgfpathlineto{\pgfqpoint{1.713584in}{1.614400in}}%
\pgfpathlineto{\pgfqpoint{1.713868in}{1.637553in}}%
\pgfpathlineto{\pgfqpoint{1.718828in}{1.637478in}}%
\pgfpathlineto{\pgfqpoint{1.718799in}{1.643280in}}%
\pgfpathlineto{\pgfqpoint{1.764876in}{1.642920in}}%
\pgfpathlineto{\pgfqpoint{1.765646in}{1.613993in}}%
\pgfpathlineto{\pgfqpoint{1.765864in}{1.580217in}}%
\pgfpathclose%
\pgfusepath{fill}%
\end{pgfscope}%
\begin{pgfscope}%
\pgfpathrectangle{\pgfqpoint{0.100000in}{0.100000in}}{\pgfqpoint{3.007045in}{1.925000in}}%
\pgfusepath{clip}%
\pgfsetbuttcap%
\pgfsetmiterjoin%
\definecolor{currentfill}{rgb}{0.612687,0.842261,0.592480}%
\pgfsetfillcolor{currentfill}%
\pgfsetlinewidth{0.000000pt}%
\definecolor{currentstroke}{rgb}{0.000000,0.000000,0.000000}%
\pgfsetstrokecolor{currentstroke}%
\pgfsetstrokeopacity{0.000000}%
\pgfsetdash{}{0pt}%
\pgfpathmoveto{\pgfqpoint{2.364743in}{1.174580in}}%
\pgfpathlineto{\pgfqpoint{2.353578in}{1.170749in}}%
\pgfpathlineto{\pgfqpoint{2.340324in}{1.169258in}}%
\pgfpathlineto{\pgfqpoint{2.338836in}{1.184076in}}%
\pgfpathlineto{\pgfqpoint{2.331627in}{1.183748in}}%
\pgfpathlineto{\pgfqpoint{2.330635in}{1.204428in}}%
\pgfpathlineto{\pgfqpoint{2.346250in}{1.205103in}}%
\pgfpathlineto{\pgfqpoint{2.345734in}{1.216099in}}%
\pgfpathlineto{\pgfqpoint{2.366064in}{1.217273in}}%
\pgfpathlineto{\pgfqpoint{2.366865in}{1.205475in}}%
\pgfpathlineto{\pgfqpoint{2.362689in}{1.195710in}}%
\pgfpathlineto{\pgfqpoint{2.363339in}{1.187311in}}%
\pgfpathlineto{\pgfqpoint{2.365382in}{1.186478in}}%
\pgfpathlineto{\pgfqpoint{2.364743in}{1.174580in}}%
\pgfpathclose%
\pgfusepath{fill}%
\end{pgfscope}%
\begin{pgfscope}%
\pgfpathrectangle{\pgfqpoint{0.100000in}{0.100000in}}{\pgfqpoint{3.007045in}{1.925000in}}%
\pgfusepath{clip}%
\pgfsetbuttcap%
\pgfsetmiterjoin%
\definecolor{currentfill}{rgb}{0.134025,0.542330,0.268281}%
\pgfsetfillcolor{currentfill}%
\pgfsetlinewidth{0.000000pt}%
\definecolor{currentstroke}{rgb}{0.000000,0.000000,0.000000}%
\pgfsetstrokecolor{currentstroke}%
\pgfsetstrokeopacity{0.000000}%
\pgfsetdash{}{0pt}%
\pgfpathmoveto{\pgfqpoint{3.013596in}{1.761178in}}%
\pgfpathlineto{\pgfqpoint{3.033217in}{1.758207in}}%
\pgfpathlineto{\pgfqpoint{3.030957in}{1.753391in}}%
\pgfpathlineto{\pgfqpoint{3.036928in}{1.745059in}}%
\pgfpathlineto{\pgfqpoint{3.036155in}{1.740067in}}%
\pgfpathlineto{\pgfqpoint{3.045228in}{1.732473in}}%
\pgfpathlineto{\pgfqpoint{3.046987in}{1.737585in}}%
\pgfpathlineto{\pgfqpoint{3.052581in}{1.737319in}}%
\pgfpathlineto{\pgfqpoint{3.064704in}{1.724060in}}%
\pgfpathlineto{\pgfqpoint{3.067422in}{1.717798in}}%
\pgfpathlineto{\pgfqpoint{3.063897in}{1.711796in}}%
\pgfpathlineto{\pgfqpoint{3.061530in}{1.704211in}}%
\pgfpathlineto{\pgfqpoint{3.052328in}{1.703378in}}%
\pgfpathlineto{\pgfqpoint{3.052800in}{1.697836in}}%
\pgfpathlineto{\pgfqpoint{3.045823in}{1.697647in}}%
\pgfpathlineto{\pgfqpoint{3.044375in}{1.690587in}}%
\pgfpathlineto{\pgfqpoint{3.036932in}{1.689276in}}%
\pgfpathlineto{\pgfqpoint{3.034646in}{1.679524in}}%
\pgfpathlineto{\pgfqpoint{3.031696in}{1.678779in}}%
\pgfpathlineto{\pgfqpoint{3.025401in}{1.689136in}}%
\pgfpathlineto{\pgfqpoint{3.014639in}{1.710091in}}%
\pgfpathlineto{\pgfqpoint{3.019840in}{1.712826in}}%
\pgfpathlineto{\pgfqpoint{3.010773in}{1.730803in}}%
\pgfpathlineto{\pgfqpoint{3.015744in}{1.733166in}}%
\pgfpathlineto{\pgfqpoint{3.003662in}{1.755298in}}%
\pgfpathlineto{\pgfqpoint{3.013596in}{1.761178in}}%
\pgfpathclose%
\pgfusepath{fill}%
\end{pgfscope}%
\begin{pgfscope}%
\pgfpathrectangle{\pgfqpoint{0.100000in}{0.100000in}}{\pgfqpoint{3.007045in}{1.925000in}}%
\pgfusepath{clip}%
\pgfsetbuttcap%
\pgfsetmiterjoin%
\definecolor{currentfill}{rgb}{0.618224,0.844844,0.597032}%
\pgfsetfillcolor{currentfill}%
\pgfsetlinewidth{0.000000pt}%
\definecolor{currentstroke}{rgb}{0.000000,0.000000,0.000000}%
\pgfsetstrokecolor{currentstroke}%
\pgfsetstrokeopacity{0.000000}%
\pgfsetdash{}{0pt}%
\pgfpathmoveto{\pgfqpoint{2.378309in}{1.338718in}}%
\pgfpathlineto{\pgfqpoint{2.365754in}{1.334376in}}%
\pgfpathlineto{\pgfqpoint{2.346907in}{1.331869in}}%
\pgfpathlineto{\pgfqpoint{2.342797in}{1.334180in}}%
\pgfpathlineto{\pgfqpoint{2.345273in}{1.312326in}}%
\pgfpathlineto{\pgfqpoint{2.322615in}{1.309553in}}%
\pgfpathlineto{\pgfqpoint{2.305531in}{1.307364in}}%
\pgfpathlineto{\pgfqpoint{2.303481in}{1.324508in}}%
\pgfpathlineto{\pgfqpoint{2.297909in}{1.323832in}}%
\pgfpathlineto{\pgfqpoint{2.295342in}{1.329269in}}%
\pgfpathlineto{\pgfqpoint{2.292953in}{1.341716in}}%
\pgfpathlineto{\pgfqpoint{2.323909in}{1.346635in}}%
\pgfpathlineto{\pgfqpoint{2.320460in}{1.370056in}}%
\pgfpathlineto{\pgfqpoint{2.343612in}{1.373779in}}%
\pgfpathlineto{\pgfqpoint{2.349537in}{1.370183in}}%
\pgfpathlineto{\pgfqpoint{2.343956in}{1.363335in}}%
\pgfpathlineto{\pgfqpoint{2.338987in}{1.354529in}}%
\pgfpathlineto{\pgfqpoint{2.338422in}{1.346684in}}%
\pgfpathlineto{\pgfqpoint{2.344381in}{1.346977in}}%
\pgfpathlineto{\pgfqpoint{2.358766in}{1.342811in}}%
\pgfpathlineto{\pgfqpoint{2.364318in}{1.338530in}}%
\pgfpathlineto{\pgfqpoint{2.375956in}{1.341494in}}%
\pgfpathlineto{\pgfqpoint{2.378309in}{1.338718in}}%
\pgfpathclose%
\pgfusepath{fill}%
\end{pgfscope}%
\begin{pgfscope}%
\pgfpathrectangle{\pgfqpoint{0.100000in}{0.100000in}}{\pgfqpoint{3.007045in}{1.925000in}}%
\pgfusepath{clip}%
\pgfsetbuttcap%
\pgfsetmiterjoin%
\definecolor{currentfill}{rgb}{0.226759,0.640569,0.342191}%
\pgfsetfillcolor{currentfill}%
\pgfsetlinewidth{0.000000pt}%
\definecolor{currentstroke}{rgb}{0.000000,0.000000,0.000000}%
\pgfsetstrokecolor{currentstroke}%
\pgfsetstrokeopacity{0.000000}%
\pgfsetdash{}{0pt}%
\pgfpathmoveto{\pgfqpoint{2.000756in}{1.600279in}}%
\pgfpathlineto{\pgfqpoint{2.002328in}{1.571903in}}%
\pgfpathlineto{\pgfqpoint{2.003003in}{1.560543in}}%
\pgfpathlineto{\pgfqpoint{1.973425in}{1.558760in}}%
\pgfpathlineto{\pgfqpoint{1.971604in}{1.598763in}}%
\pgfpathlineto{\pgfqpoint{1.960172in}{1.598201in}}%
\pgfpathlineto{\pgfqpoint{1.959683in}{1.609718in}}%
\pgfpathlineto{\pgfqpoint{1.930881in}{1.608680in}}%
\pgfpathlineto{\pgfqpoint{1.929386in}{1.648471in}}%
\pgfpathlineto{\pgfqpoint{1.937520in}{1.651178in}}%
\pgfpathlineto{\pgfqpoint{1.944689in}{1.656422in}}%
\pgfpathlineto{\pgfqpoint{1.949266in}{1.655981in}}%
\pgfpathlineto{\pgfqpoint{1.955129in}{1.661825in}}%
\pgfpathlineto{\pgfqpoint{1.960258in}{1.663453in}}%
\pgfpathlineto{\pgfqpoint{1.965406in}{1.659054in}}%
\pgfpathlineto{\pgfqpoint{1.959719in}{1.649522in}}%
\pgfpathlineto{\pgfqpoint{1.959203in}{1.638150in}}%
\pgfpathlineto{\pgfqpoint{1.968638in}{1.644005in}}%
\pgfpathlineto{\pgfqpoint{1.975507in}{1.638992in}}%
\pgfpathlineto{\pgfqpoint{1.976511in}{1.616232in}}%
\pgfpathlineto{\pgfqpoint{1.982301in}{1.616494in}}%
\pgfpathlineto{\pgfqpoint{1.982576in}{1.610763in}}%
\pgfpathlineto{\pgfqpoint{1.988231in}{1.611105in}}%
\pgfpathlineto{\pgfqpoint{1.988873in}{1.599607in}}%
\pgfpathlineto{\pgfqpoint{2.000756in}{1.600279in}}%
\pgfpathclose%
\pgfusepath{fill}%
\end{pgfscope}%
\begin{pgfscope}%
\pgfpathrectangle{\pgfqpoint{0.100000in}{0.100000in}}{\pgfqpoint{3.007045in}{1.925000in}}%
\pgfusepath{clip}%
\pgfsetbuttcap%
\pgfsetmiterjoin%
\definecolor{currentfill}{rgb}{0.723122,0.889612,0.697178}%
\pgfsetfillcolor{currentfill}%
\pgfsetlinewidth{0.000000pt}%
\definecolor{currentstroke}{rgb}{0.000000,0.000000,0.000000}%
\pgfsetstrokecolor{currentstroke}%
\pgfsetstrokeopacity{0.000000}%
\pgfsetdash{}{0pt}%
\pgfpathmoveto{\pgfqpoint{2.212181in}{1.234336in}}%
\pgfpathlineto{\pgfqpoint{2.212424in}{1.232038in}}%
\pgfpathlineto{\pgfqpoint{2.189712in}{1.229760in}}%
\pgfpathlineto{\pgfqpoint{2.189499in}{1.232103in}}%
\pgfpathlineto{\pgfqpoint{2.169632in}{1.230355in}}%
\pgfpathlineto{\pgfqpoint{2.171191in}{1.213179in}}%
\pgfpathlineto{\pgfqpoint{2.154631in}{1.211728in}}%
\pgfpathlineto{\pgfqpoint{2.154351in}{1.223276in}}%
\pgfpathlineto{\pgfqpoint{2.147976in}{1.224103in}}%
\pgfpathlineto{\pgfqpoint{2.149319in}{1.206705in}}%
\pgfpathlineto{\pgfqpoint{2.128961in}{1.204856in}}%
\pgfpathlineto{\pgfqpoint{2.126956in}{1.227545in}}%
\pgfpathlineto{\pgfqpoint{2.125964in}{1.244692in}}%
\pgfpathlineto{\pgfqpoint{2.146371in}{1.246681in}}%
\pgfpathlineto{\pgfqpoint{2.145056in}{1.262796in}}%
\pgfpathlineto{\pgfqpoint{2.166328in}{1.264572in}}%
\pgfpathlineto{\pgfqpoint{2.165690in}{1.271203in}}%
\pgfpathlineto{\pgfqpoint{2.171318in}{1.271716in}}%
\pgfpathlineto{\pgfqpoint{2.173650in}{1.276857in}}%
\pgfpathlineto{\pgfqpoint{2.190977in}{1.278330in}}%
\pgfpathlineto{\pgfqpoint{2.192034in}{1.266835in}}%
\pgfpathlineto{\pgfqpoint{2.195116in}{1.264251in}}%
\pgfpathlineto{\pgfqpoint{2.202654in}{1.264977in}}%
\pgfpathlineto{\pgfqpoint{2.204170in}{1.247864in}}%
\pgfpathlineto{\pgfqpoint{2.211166in}{1.244697in}}%
\pgfpathlineto{\pgfqpoint{2.212181in}{1.234336in}}%
\pgfpathclose%
\pgfusepath{fill}%
\end{pgfscope}%
\begin{pgfscope}%
\pgfpathrectangle{\pgfqpoint{0.100000in}{0.100000in}}{\pgfqpoint{3.007045in}{1.925000in}}%
\pgfusepath{clip}%
\pgfsetbuttcap%
\pgfsetmiterjoin%
\definecolor{currentfill}{rgb}{0.601615,0.837093,0.583376}%
\pgfsetfillcolor{currentfill}%
\pgfsetlinewidth{0.000000pt}%
\definecolor{currentstroke}{rgb}{0.000000,0.000000,0.000000}%
\pgfsetstrokecolor{currentstroke}%
\pgfsetstrokeopacity{0.000000}%
\pgfsetdash{}{0pt}%
\pgfpathmoveto{\pgfqpoint{1.851931in}{0.616696in}}%
\pgfpathlineto{\pgfqpoint{1.834269in}{0.615251in}}%
\pgfpathlineto{\pgfqpoint{1.821877in}{0.611209in}}%
\pgfpathlineto{\pgfqpoint{1.810103in}{0.619771in}}%
\pgfpathlineto{\pgfqpoint{1.804231in}{0.625251in}}%
\pgfpathlineto{\pgfqpoint{1.800480in}{0.631756in}}%
\pgfpathlineto{\pgfqpoint{1.787706in}{0.634802in}}%
\pgfpathlineto{\pgfqpoint{1.779925in}{0.639399in}}%
\pgfpathlineto{\pgfqpoint{1.776581in}{0.643123in}}%
\pgfpathlineto{\pgfqpoint{1.773764in}{0.656330in}}%
\pgfpathlineto{\pgfqpoint{1.775875in}{0.660550in}}%
\pgfpathlineto{\pgfqpoint{1.803211in}{0.660558in}}%
\pgfpathlineto{\pgfqpoint{1.799830in}{0.669148in}}%
\pgfpathlineto{\pgfqpoint{1.827733in}{0.669764in}}%
\pgfpathlineto{\pgfqpoint{1.835531in}{0.660856in}}%
\pgfpathlineto{\pgfqpoint{1.838775in}{0.656201in}}%
\pgfpathlineto{\pgfqpoint{1.839877in}{0.650274in}}%
\pgfpathlineto{\pgfqpoint{1.837976in}{0.644603in}}%
\pgfpathlineto{\pgfqpoint{1.843180in}{0.636011in}}%
\pgfpathlineto{\pgfqpoint{1.847519in}{0.631354in}}%
\pgfpathlineto{\pgfqpoint{1.846842in}{0.625571in}}%
\pgfpathlineto{\pgfqpoint{1.851931in}{0.616696in}}%
\pgfpathclose%
\pgfusepath{fill}%
\end{pgfscope}%
\begin{pgfscope}%
\pgfpathrectangle{\pgfqpoint{0.100000in}{0.100000in}}{\pgfqpoint{3.007045in}{1.925000in}}%
\pgfusepath{clip}%
\pgfsetbuttcap%
\pgfsetmiterjoin%
\definecolor{currentfill}{rgb}{0.307451,0.696348,0.390465}%
\pgfsetfillcolor{currentfill}%
\pgfsetlinewidth{0.000000pt}%
\definecolor{currentstroke}{rgb}{0.000000,0.000000,0.000000}%
\pgfsetstrokecolor{currentstroke}%
\pgfsetstrokeopacity{0.000000}%
\pgfsetdash{}{0pt}%
\pgfpathmoveto{\pgfqpoint{2.639205in}{0.819156in}}%
\pgfpathlineto{\pgfqpoint{2.628686in}{0.826071in}}%
\pgfpathlineto{\pgfqpoint{2.612678in}{0.829037in}}%
\pgfpathlineto{\pgfqpoint{2.597720in}{0.838651in}}%
\pgfpathlineto{\pgfqpoint{2.590755in}{0.837224in}}%
\pgfpathlineto{\pgfqpoint{2.595239in}{0.853039in}}%
\pgfpathlineto{\pgfqpoint{2.593582in}{0.855868in}}%
\pgfpathlineto{\pgfqpoint{2.597384in}{0.859907in}}%
\pgfpathlineto{\pgfqpoint{2.592778in}{0.867711in}}%
\pgfpathlineto{\pgfqpoint{2.596705in}{0.871224in}}%
\pgfpathlineto{\pgfqpoint{2.580510in}{0.880092in}}%
\pgfpathlineto{\pgfqpoint{2.582316in}{0.884878in}}%
\pgfpathlineto{\pgfqpoint{2.577854in}{0.890686in}}%
\pgfpathlineto{\pgfqpoint{2.571314in}{0.892287in}}%
\pgfpathlineto{\pgfqpoint{2.582994in}{0.902658in}}%
\pgfpathlineto{\pgfqpoint{2.595973in}{0.904310in}}%
\pgfpathlineto{\pgfqpoint{2.607323in}{0.893730in}}%
\pgfpathlineto{\pgfqpoint{2.613555in}{0.916458in}}%
\pgfpathlineto{\pgfqpoint{2.637551in}{0.899069in}}%
\pgfpathlineto{\pgfqpoint{2.671279in}{0.875019in}}%
\pgfpathlineto{\pgfqpoint{2.663098in}{0.870424in}}%
\pgfpathlineto{\pgfqpoint{2.657985in}{0.865431in}}%
\pgfpathlineto{\pgfqpoint{2.651920in}{0.856646in}}%
\pgfpathlineto{\pgfqpoint{2.646223in}{0.845821in}}%
\pgfpathlineto{\pgfqpoint{2.644013in}{0.838187in}}%
\pgfpathlineto{\pgfqpoint{2.644164in}{0.824590in}}%
\pgfpathlineto{\pgfqpoint{2.639205in}{0.819156in}}%
\pgfpathclose%
\pgfusepath{fill}%
\end{pgfscope}%
\begin{pgfscope}%
\pgfpathrectangle{\pgfqpoint{0.100000in}{0.100000in}}{\pgfqpoint{3.007045in}{1.925000in}}%
\pgfusepath{clip}%
\pgfsetbuttcap%
\pgfsetmiterjoin%
\definecolor{currentfill}{rgb}{0.241522,0.656317,0.354002}%
\pgfsetfillcolor{currentfill}%
\pgfsetlinewidth{0.000000pt}%
\definecolor{currentstroke}{rgb}{0.000000,0.000000,0.000000}%
\pgfsetstrokecolor{currentstroke}%
\pgfsetstrokeopacity{0.000000}%
\pgfsetdash{}{0pt}%
\pgfpathmoveto{\pgfqpoint{2.192935in}{1.614287in}}%
\pgfpathlineto{\pgfqpoint{2.188229in}{1.662007in}}%
\pgfpathlineto{\pgfqpoint{2.204739in}{1.662788in}}%
\pgfpathlineto{\pgfqpoint{2.213687in}{1.668241in}}%
\pgfpathlineto{\pgfqpoint{2.228170in}{1.671708in}}%
\pgfpathlineto{\pgfqpoint{2.226016in}{1.665581in}}%
\pgfpathlineto{\pgfqpoint{2.227291in}{1.652489in}}%
\pgfpathlineto{\pgfqpoint{2.236977in}{1.650715in}}%
\pgfpathlineto{\pgfqpoint{2.243546in}{1.654372in}}%
\pgfpathlineto{\pgfqpoint{2.247528in}{1.650275in}}%
\pgfpathlineto{\pgfqpoint{2.255094in}{1.656649in}}%
\pgfpathlineto{\pgfqpoint{2.269105in}{1.658515in}}%
\pgfpathlineto{\pgfqpoint{2.268196in}{1.652678in}}%
\pgfpathlineto{\pgfqpoint{2.270778in}{1.646529in}}%
\pgfpathlineto{\pgfqpoint{2.269357in}{1.637766in}}%
\pgfpathlineto{\pgfqpoint{2.275952in}{1.634432in}}%
\pgfpathlineto{\pgfqpoint{2.278409in}{1.628086in}}%
\pgfpathlineto{\pgfqpoint{2.287170in}{1.626728in}}%
\pgfpathlineto{\pgfqpoint{2.296641in}{1.634077in}}%
\pgfpathlineto{\pgfqpoint{2.302457in}{1.627925in}}%
\pgfpathlineto{\pgfqpoint{2.301033in}{1.623548in}}%
\pgfpathlineto{\pgfqpoint{2.294668in}{1.624039in}}%
\pgfpathlineto{\pgfqpoint{2.288149in}{1.622604in}}%
\pgfpathlineto{\pgfqpoint{2.283389in}{1.624709in}}%
\pgfpathlineto{\pgfqpoint{2.272380in}{1.623457in}}%
\pgfpathlineto{\pgfqpoint{2.261434in}{1.619276in}}%
\pgfpathlineto{\pgfqpoint{2.254060in}{1.620885in}}%
\pgfpathlineto{\pgfqpoint{2.251911in}{1.624956in}}%
\pgfpathlineto{\pgfqpoint{2.246315in}{1.624264in}}%
\pgfpathlineto{\pgfqpoint{2.244624in}{1.611374in}}%
\pgfpathlineto{\pgfqpoint{2.236991in}{1.615863in}}%
\pgfpathlineto{\pgfqpoint{2.232484in}{1.620710in}}%
\pgfpathlineto{\pgfqpoint{2.225249in}{1.622997in}}%
\pgfpathlineto{\pgfqpoint{2.213158in}{1.624960in}}%
\pgfpathlineto{\pgfqpoint{2.207852in}{1.624039in}}%
\pgfpathlineto{\pgfqpoint{2.201875in}{1.614996in}}%
\pgfpathlineto{\pgfqpoint{2.192935in}{1.614287in}}%
\pgfpathclose%
\pgfusepath{fill}%
\end{pgfscope}%
\begin{pgfscope}%
\pgfpathrectangle{\pgfqpoint{0.100000in}{0.100000in}}{\pgfqpoint{3.007045in}{1.925000in}}%
\pgfusepath{clip}%
\pgfsetbuttcap%
\pgfsetmiterjoin%
\definecolor{currentfill}{rgb}{0.618224,0.844844,0.597032}%
\pgfsetfillcolor{currentfill}%
\pgfsetlinewidth{0.000000pt}%
\definecolor{currentstroke}{rgb}{0.000000,0.000000,0.000000}%
\pgfsetstrokecolor{currentstroke}%
\pgfsetstrokeopacity{0.000000}%
\pgfsetdash{}{0pt}%
\pgfpathmoveto{\pgfqpoint{2.085014in}{0.905571in}}%
\pgfpathlineto{\pgfqpoint{2.069407in}{0.904987in}}%
\pgfpathlineto{\pgfqpoint{2.068214in}{0.916226in}}%
\pgfpathlineto{\pgfqpoint{2.061506in}{0.919581in}}%
\pgfpathlineto{\pgfqpoint{2.051337in}{0.919577in}}%
\pgfpathlineto{\pgfqpoint{2.045514in}{0.912506in}}%
\pgfpathlineto{\pgfqpoint{2.043371in}{0.916596in}}%
\pgfpathlineto{\pgfqpoint{2.042608in}{0.925600in}}%
\pgfpathlineto{\pgfqpoint{2.049440in}{0.927189in}}%
\pgfpathlineto{\pgfqpoint{2.052220in}{0.930837in}}%
\pgfpathlineto{\pgfqpoint{2.051897in}{0.936223in}}%
\pgfpathlineto{\pgfqpoint{2.054350in}{0.941105in}}%
\pgfpathlineto{\pgfqpoint{2.055729in}{0.949079in}}%
\pgfpathlineto{\pgfqpoint{2.060132in}{0.952292in}}%
\pgfpathlineto{\pgfqpoint{2.054510in}{0.958420in}}%
\pgfpathlineto{\pgfqpoint{2.062788in}{0.960749in}}%
\pgfpathlineto{\pgfqpoint{2.061897in}{0.966247in}}%
\pgfpathlineto{\pgfqpoint{2.061582in}{0.976889in}}%
\pgfpathlineto{\pgfqpoint{2.064440in}{0.977028in}}%
\pgfpathlineto{\pgfqpoint{2.068022in}{0.977336in}}%
\pgfpathlineto{\pgfqpoint{2.071414in}{0.971616in}}%
\pgfpathlineto{\pgfqpoint{2.068381in}{0.961543in}}%
\pgfpathlineto{\pgfqpoint{2.065703in}{0.958280in}}%
\pgfpathlineto{\pgfqpoint{2.083172in}{0.958841in}}%
\pgfpathlineto{\pgfqpoint{2.083454in}{0.947053in}}%
\pgfpathlineto{\pgfqpoint{2.074636in}{0.931335in}}%
\pgfpathlineto{\pgfqpoint{2.089884in}{0.925436in}}%
\pgfpathlineto{\pgfqpoint{2.090507in}{0.908195in}}%
\pgfpathlineto{\pgfqpoint{2.085014in}{0.905571in}}%
\pgfpathclose%
\pgfusepath{fill}%
\end{pgfscope}%
\begin{pgfscope}%
\pgfpathrectangle{\pgfqpoint{0.100000in}{0.100000in}}{\pgfqpoint{3.007045in}{1.925000in}}%
\pgfusepath{clip}%
\pgfsetbuttcap%
\pgfsetmiterjoin%
\definecolor{currentfill}{rgb}{0.596078,0.834510,0.578824}%
\pgfsetfillcolor{currentfill}%
\pgfsetlinewidth{0.000000pt}%
\definecolor{currentstroke}{rgb}{0.000000,0.000000,0.000000}%
\pgfsetstrokecolor{currentstroke}%
\pgfsetstrokeopacity{0.000000}%
\pgfsetdash{}{0pt}%
\pgfpathmoveto{\pgfqpoint{2.492448in}{0.960338in}}%
\pgfpathlineto{\pgfqpoint{2.498902in}{0.973883in}}%
\pgfpathlineto{\pgfqpoint{2.494999in}{0.975158in}}%
\pgfpathlineto{\pgfqpoint{2.489274in}{0.973776in}}%
\pgfpathlineto{\pgfqpoint{2.480170in}{0.978949in}}%
\pgfpathlineto{\pgfqpoint{2.474213in}{0.983871in}}%
\pgfpathlineto{\pgfqpoint{2.479657in}{0.995365in}}%
\pgfpathlineto{\pgfqpoint{2.486521in}{0.996045in}}%
\pgfpathlineto{\pgfqpoint{2.494094in}{0.994158in}}%
\pgfpathlineto{\pgfqpoint{2.499152in}{0.989884in}}%
\pgfpathlineto{\pgfqpoint{2.506001in}{0.993623in}}%
\pgfpathlineto{\pgfqpoint{2.510978in}{0.993811in}}%
\pgfpathlineto{\pgfqpoint{2.522081in}{0.997367in}}%
\pgfpathlineto{\pgfqpoint{2.532109in}{0.998638in}}%
\pgfpathlineto{\pgfqpoint{2.533363in}{0.985590in}}%
\pgfpathlineto{\pgfqpoint{2.531826in}{0.973469in}}%
\pgfpathlineto{\pgfqpoint{2.535380in}{0.962820in}}%
\pgfpathlineto{\pgfqpoint{2.523665in}{0.961730in}}%
\pgfpathlineto{\pgfqpoint{2.523136in}{0.963616in}}%
\pgfpathlineto{\pgfqpoint{2.492448in}{0.960338in}}%
\pgfpathclose%
\pgfusepath{fill}%
\end{pgfscope}%
\begin{pgfscope}%
\pgfpathrectangle{\pgfqpoint{0.100000in}{0.100000in}}{\pgfqpoint{3.007045in}{1.925000in}}%
\pgfusepath{clip}%
\pgfsetbuttcap%
\pgfsetmiterjoin%
\definecolor{currentfill}{rgb}{0.535179,0.806090,0.528750}%
\pgfsetfillcolor{currentfill}%
\pgfsetlinewidth{0.000000pt}%
\definecolor{currentstroke}{rgb}{0.000000,0.000000,0.000000}%
\pgfsetstrokecolor{currentstroke}%
\pgfsetstrokeopacity{0.000000}%
\pgfsetdash{}{0pt}%
\pgfpathmoveto{\pgfqpoint{2.350943in}{1.269110in}}%
\pgfpathlineto{\pgfqpoint{2.345898in}{1.268387in}}%
\pgfpathlineto{\pgfqpoint{2.327732in}{1.268420in}}%
\pgfpathlineto{\pgfqpoint{2.322615in}{1.309553in}}%
\pgfpathlineto{\pgfqpoint{2.345273in}{1.312326in}}%
\pgfpathlineto{\pgfqpoint{2.342797in}{1.334180in}}%
\pgfpathlineto{\pgfqpoint{2.346907in}{1.331869in}}%
\pgfpathlineto{\pgfqpoint{2.365754in}{1.334376in}}%
\pgfpathlineto{\pgfqpoint{2.371168in}{1.333253in}}%
\pgfpathlineto{\pgfqpoint{2.375866in}{1.305055in}}%
\pgfpathlineto{\pgfqpoint{2.361935in}{1.302986in}}%
\pgfpathlineto{\pgfqpoint{2.364496in}{1.284008in}}%
\pgfpathlineto{\pgfqpoint{2.349350in}{1.280970in}}%
\pgfpathlineto{\pgfqpoint{2.350943in}{1.269110in}}%
\pgfpathclose%
\pgfusepath{fill}%
\end{pgfscope}%
\begin{pgfscope}%
\pgfpathrectangle{\pgfqpoint{0.100000in}{0.100000in}}{\pgfqpoint{3.007045in}{1.925000in}}%
\pgfusepath{clip}%
\pgfsetbuttcap%
\pgfsetmiterjoin%
\definecolor{currentfill}{rgb}{0.579469,0.826759,0.565167}%
\pgfsetfillcolor{currentfill}%
\pgfsetlinewidth{0.000000pt}%
\definecolor{currentstroke}{rgb}{0.000000,0.000000,0.000000}%
\pgfsetstrokecolor{currentstroke}%
\pgfsetstrokeopacity{0.000000}%
\pgfsetdash{}{0pt}%
\pgfpathmoveto{\pgfqpoint{2.880647in}{1.472875in}}%
\pgfpathlineto{\pgfqpoint{2.836904in}{1.462520in}}%
\pgfpathlineto{\pgfqpoint{2.833508in}{1.484777in}}%
\pgfpathlineto{\pgfqpoint{2.832224in}{1.502258in}}%
\pgfpathlineto{\pgfqpoint{2.827848in}{1.505323in}}%
\pgfpathlineto{\pgfqpoint{2.831656in}{1.508759in}}%
\pgfpathlineto{\pgfqpoint{2.862374in}{1.515277in}}%
\pgfpathlineto{\pgfqpoint{2.865846in}{1.512168in}}%
\pgfpathlineto{\pgfqpoint{2.865008in}{1.506125in}}%
\pgfpathlineto{\pgfqpoint{2.867051in}{1.490726in}}%
\pgfpathlineto{\pgfqpoint{2.873175in}{1.485821in}}%
\pgfpathlineto{\pgfqpoint{2.872055in}{1.481134in}}%
\pgfpathlineto{\pgfqpoint{2.878530in}{1.481285in}}%
\pgfpathlineto{\pgfqpoint{2.880647in}{1.472875in}}%
\pgfpathclose%
\pgfusepath{fill}%
\end{pgfscope}%
\begin{pgfscope}%
\pgfpathrectangle{\pgfqpoint{0.100000in}{0.100000in}}{\pgfqpoint{3.007045in}{1.925000in}}%
\pgfusepath{clip}%
\pgfsetbuttcap%
\pgfsetmiterjoin%
\definecolor{currentfill}{rgb}{0.596078,0.834510,0.578824}%
\pgfsetfillcolor{currentfill}%
\pgfsetlinewidth{0.000000pt}%
\definecolor{currentstroke}{rgb}{0.000000,0.000000,0.000000}%
\pgfsetstrokecolor{currentstroke}%
\pgfsetstrokeopacity{0.000000}%
\pgfsetdash{}{0pt}%
\pgfpathmoveto{\pgfqpoint{1.631910in}{1.063521in}}%
\pgfpathlineto{\pgfqpoint{1.632419in}{1.080689in}}%
\pgfpathlineto{\pgfqpoint{1.621015in}{1.081045in}}%
\pgfpathlineto{\pgfqpoint{1.621825in}{1.109765in}}%
\pgfpathlineto{\pgfqpoint{1.650255in}{1.108902in}}%
\pgfpathlineto{\pgfqpoint{1.672984in}{1.108364in}}%
\pgfpathlineto{\pgfqpoint{1.672876in}{1.102642in}}%
\pgfpathlineto{\pgfqpoint{1.678559in}{1.102524in}}%
\pgfpathlineto{\pgfqpoint{1.677388in}{1.091095in}}%
\pgfpathlineto{\pgfqpoint{1.676880in}{1.073769in}}%
\pgfpathlineto{\pgfqpoint{1.660705in}{1.074262in}}%
\pgfpathlineto{\pgfqpoint{1.660467in}{1.062737in}}%
\pgfpathlineto{\pgfqpoint{1.631910in}{1.063521in}}%
\pgfpathclose%
\pgfusepath{fill}%
\end{pgfscope}%
\begin{pgfscope}%
\pgfpathrectangle{\pgfqpoint{0.100000in}{0.100000in}}{\pgfqpoint{3.007045in}{1.925000in}}%
\pgfusepath{clip}%
\pgfsetbuttcap%
\pgfsetmiterjoin%
\definecolor{currentfill}{rgb}{0.794233,0.919262,0.767705}%
\pgfsetfillcolor{currentfill}%
\pgfsetlinewidth{0.000000pt}%
\definecolor{currentstroke}{rgb}{0.000000,0.000000,0.000000}%
\pgfsetstrokecolor{currentstroke}%
\pgfsetstrokeopacity{0.000000}%
\pgfsetdash{}{0pt}%
\pgfpathmoveto{\pgfqpoint{1.449237in}{1.655533in}}%
\pgfpathlineto{\pgfqpoint{1.448770in}{1.649792in}}%
\pgfpathlineto{\pgfqpoint{1.410950in}{1.652859in}}%
\pgfpathlineto{\pgfqpoint{1.379987in}{1.655666in}}%
\pgfpathlineto{\pgfqpoint{1.382174in}{1.678842in}}%
\pgfpathlineto{\pgfqpoint{1.379578in}{1.679082in}}%
\pgfpathlineto{\pgfqpoint{1.381806in}{1.702240in}}%
\pgfpathlineto{\pgfqpoint{1.407305in}{1.699840in}}%
\pgfpathlineto{\pgfqpoint{1.409324in}{1.722778in}}%
\pgfpathlineto{\pgfqpoint{1.429757in}{1.720971in}}%
\pgfpathlineto{\pgfqpoint{1.433959in}{1.727722in}}%
\pgfpathlineto{\pgfqpoint{1.441706in}{1.725605in}}%
\pgfpathlineto{\pgfqpoint{1.447405in}{1.726710in}}%
\pgfpathlineto{\pgfqpoint{1.446373in}{1.716181in}}%
\pgfpathlineto{\pgfqpoint{1.448796in}{1.712766in}}%
\pgfpathlineto{\pgfqpoint{1.447485in}{1.696400in}}%
\pgfpathlineto{\pgfqpoint{1.450238in}{1.696191in}}%
\pgfpathlineto{\pgfqpoint{1.448415in}{1.673164in}}%
\pgfpathlineto{\pgfqpoint{1.450582in}{1.673001in}}%
\pgfpathlineto{\pgfqpoint{1.449237in}{1.655533in}}%
\pgfpathclose%
\pgfusepath{fill}%
\end{pgfscope}%
\begin{pgfscope}%
\pgfpathrectangle{\pgfqpoint{0.100000in}{0.100000in}}{\pgfqpoint{3.007045in}{1.925000in}}%
\pgfusepath{clip}%
\pgfsetbuttcap%
\pgfsetmiterjoin%
\definecolor{currentfill}{rgb}{0.845905,0.939931,0.822822}%
\pgfsetfillcolor{currentfill}%
\pgfsetlinewidth{0.000000pt}%
\definecolor{currentstroke}{rgb}{0.000000,0.000000,0.000000}%
\pgfsetstrokecolor{currentstroke}%
\pgfsetstrokeopacity{0.000000}%
\pgfsetdash{}{0pt}%
\pgfpathmoveto{\pgfqpoint{0.842950in}{0.338119in}}%
\pgfpathlineto{\pgfqpoint{0.844097in}{0.340724in}}%
\pgfpathlineto{\pgfqpoint{0.845519in}{0.341453in}}%
\pgfpathlineto{\pgfqpoint{0.848054in}{0.340022in}}%
\pgfpathlineto{\pgfqpoint{0.847662in}{0.338845in}}%
\pgfpathlineto{\pgfqpoint{0.850122in}{0.336629in}}%
\pgfpathlineto{\pgfqpoint{0.844580in}{0.336908in}}%
\pgfpathlineto{\pgfqpoint{0.842950in}{0.338119in}}%
\pgfpathclose%
\pgfusepath{fill}%
\end{pgfscope}%
\begin{pgfscope}%
\pgfpathrectangle{\pgfqpoint{0.100000in}{0.100000in}}{\pgfqpoint{3.007045in}{1.925000in}}%
\pgfusepath{clip}%
\pgfsetbuttcap%
\pgfsetmiterjoin%
\definecolor{currentfill}{rgb}{0.845905,0.939931,0.822822}%
\pgfsetfillcolor{currentfill}%
\pgfsetlinewidth{0.000000pt}%
\definecolor{currentstroke}{rgb}{0.000000,0.000000,0.000000}%
\pgfsetstrokecolor{currentstroke}%
\pgfsetstrokeopacity{0.000000}%
\pgfsetdash{}{0pt}%
\pgfpathmoveto{\pgfqpoint{0.824088in}{0.336340in}}%
\pgfpathlineto{\pgfqpoint{0.826449in}{0.337990in}}%
\pgfpathlineto{\pgfqpoint{0.828677in}{0.338112in}}%
\pgfpathlineto{\pgfqpoint{0.837136in}{0.340612in}}%
\pgfpathlineto{\pgfqpoint{0.837987in}{0.342496in}}%
\pgfpathlineto{\pgfqpoint{0.839837in}{0.341361in}}%
\pgfpathlineto{\pgfqpoint{0.839821in}{0.339455in}}%
\pgfpathlineto{\pgfqpoint{0.834944in}{0.338315in}}%
\pgfpathlineto{\pgfqpoint{0.831383in}{0.335875in}}%
\pgfpathlineto{\pgfqpoint{0.829597in}{0.336005in}}%
\pgfpathlineto{\pgfqpoint{0.828843in}{0.334684in}}%
\pgfpathlineto{\pgfqpoint{0.825115in}{0.335470in}}%
\pgfpathlineto{\pgfqpoint{0.824088in}{0.336340in}}%
\pgfpathclose%
\pgfusepath{fill}%
\end{pgfscope}%
\begin{pgfscope}%
\pgfpathrectangle{\pgfqpoint{0.100000in}{0.100000in}}{\pgfqpoint{3.007045in}{1.925000in}}%
\pgfusepath{clip}%
\pgfsetbuttcap%
\pgfsetmiterjoin%
\definecolor{currentfill}{rgb}{0.845905,0.939931,0.822822}%
\pgfsetfillcolor{currentfill}%
\pgfsetlinewidth{0.000000pt}%
\definecolor{currentstroke}{rgb}{0.000000,0.000000,0.000000}%
\pgfsetstrokecolor{currentstroke}%
\pgfsetstrokeopacity{0.000000}%
\pgfsetdash{}{0pt}%
\pgfpathmoveto{\pgfqpoint{0.829895in}{0.342488in}}%
\pgfpathlineto{\pgfqpoint{0.832685in}{0.348119in}}%
\pgfpathlineto{\pgfqpoint{0.833319in}{0.346252in}}%
\pgfpathlineto{\pgfqpoint{0.829895in}{0.342488in}}%
\pgfpathclose%
\pgfusepath{fill}%
\end{pgfscope}%
\begin{pgfscope}%
\pgfpathrectangle{\pgfqpoint{0.100000in}{0.100000in}}{\pgfqpoint{3.007045in}{1.925000in}}%
\pgfusepath{clip}%
\pgfsetbuttcap%
\pgfsetmiterjoin%
\definecolor{currentfill}{rgb}{0.845905,0.939931,0.822822}%
\pgfsetfillcolor{currentfill}%
\pgfsetlinewidth{0.000000pt}%
\definecolor{currentstroke}{rgb}{0.000000,0.000000,0.000000}%
\pgfsetstrokecolor{currentstroke}%
\pgfsetstrokeopacity{0.000000}%
\pgfsetdash{}{0pt}%
\pgfpathmoveto{\pgfqpoint{0.880633in}{0.393554in}}%
\pgfpathlineto{\pgfqpoint{0.884245in}{0.391268in}}%
\pgfpathlineto{\pgfqpoint{0.886512in}{0.394965in}}%
\pgfpathlineto{\pgfqpoint{0.895631in}{0.389258in}}%
\pgfpathlineto{\pgfqpoint{0.894498in}{0.387400in}}%
\pgfpathlineto{\pgfqpoint{0.899657in}{0.384205in}}%
\pgfpathlineto{\pgfqpoint{0.914264in}{0.375762in}}%
\pgfpathlineto{\pgfqpoint{0.912562in}{0.370845in}}%
\pgfpathlineto{\pgfqpoint{0.909515in}{0.369615in}}%
\pgfpathlineto{\pgfqpoint{0.910203in}{0.367541in}}%
\pgfpathlineto{\pgfqpoint{0.907829in}{0.364423in}}%
\pgfpathlineto{\pgfqpoint{0.910185in}{0.362649in}}%
\pgfpathlineto{\pgfqpoint{0.912886in}{0.364295in}}%
\pgfpathlineto{\pgfqpoint{0.913444in}{0.363209in}}%
\pgfpathlineto{\pgfqpoint{0.916564in}{0.361368in}}%
\pgfpathlineto{\pgfqpoint{0.915470in}{0.359574in}}%
\pgfpathlineto{\pgfqpoint{0.917368in}{0.358439in}}%
\pgfpathlineto{\pgfqpoint{0.916310in}{0.356673in}}%
\pgfpathlineto{\pgfqpoint{0.917983in}{0.355693in}}%
\pgfpathlineto{\pgfqpoint{0.913691in}{0.348399in}}%
\pgfpathlineto{\pgfqpoint{0.914753in}{0.346694in}}%
\pgfpathlineto{\pgfqpoint{0.912069in}{0.341923in}}%
\pgfpathlineto{\pgfqpoint{0.920301in}{0.337334in}}%
\pgfpathlineto{\pgfqpoint{0.911029in}{0.320607in}}%
\pgfpathlineto{\pgfqpoint{0.901489in}{0.303399in}}%
\pgfpathlineto{\pgfqpoint{0.895195in}{0.292046in}}%
\pgfpathlineto{\pgfqpoint{0.891908in}{0.297356in}}%
\pgfpathlineto{\pgfqpoint{0.895042in}{0.297561in}}%
\pgfpathlineto{\pgfqpoint{0.895865in}{0.299933in}}%
\pgfpathlineto{\pgfqpoint{0.894005in}{0.301147in}}%
\pgfpathlineto{\pgfqpoint{0.893616in}{0.299471in}}%
\pgfpathlineto{\pgfqpoint{0.889618in}{0.300158in}}%
\pgfpathlineto{\pgfqpoint{0.888432in}{0.302470in}}%
\pgfpathlineto{\pgfqpoint{0.880446in}{0.309836in}}%
\pgfpathlineto{\pgfqpoint{0.876880in}{0.310671in}}%
\pgfpathlineto{\pgfqpoint{0.872039in}{0.312853in}}%
\pgfpathlineto{\pgfqpoint{0.867319in}{0.314295in}}%
\pgfpathlineto{\pgfqpoint{0.866304in}{0.315660in}}%
\pgfpathlineto{\pgfqpoint{0.866224in}{0.318437in}}%
\pgfpathlineto{\pgfqpoint{0.865182in}{0.320626in}}%
\pgfpathlineto{\pgfqpoint{0.862537in}{0.322620in}}%
\pgfpathlineto{\pgfqpoint{0.862785in}{0.325009in}}%
\pgfpathlineto{\pgfqpoint{0.861746in}{0.326681in}}%
\pgfpathlineto{\pgfqpoint{0.864218in}{0.329401in}}%
\pgfpathlineto{\pgfqpoint{0.864631in}{0.330890in}}%
\pgfpathlineto{\pgfqpoint{0.863184in}{0.332010in}}%
\pgfpathlineto{\pgfqpoint{0.861073in}{0.330560in}}%
\pgfpathlineto{\pgfqpoint{0.858822in}{0.331645in}}%
\pgfpathlineto{\pgfqpoint{0.857771in}{0.330884in}}%
\pgfpathlineto{\pgfqpoint{0.855435in}{0.335097in}}%
\pgfpathlineto{\pgfqpoint{0.853673in}{0.335957in}}%
\pgfpathlineto{\pgfqpoint{0.855410in}{0.337624in}}%
\pgfpathlineto{\pgfqpoint{0.855863in}{0.341580in}}%
\pgfpathlineto{\pgfqpoint{0.855648in}{0.344279in}}%
\pgfpathlineto{\pgfqpoint{0.852144in}{0.345090in}}%
\pgfpathlineto{\pgfqpoint{0.853807in}{0.346777in}}%
\pgfpathlineto{\pgfqpoint{0.850106in}{0.348448in}}%
\pgfpathlineto{\pgfqpoint{0.850605in}{0.352065in}}%
\pgfpathlineto{\pgfqpoint{0.848839in}{0.353027in}}%
\pgfpathlineto{\pgfqpoint{0.848400in}{0.354633in}}%
\pgfpathlineto{\pgfqpoint{0.846238in}{0.354335in}}%
\pgfpathlineto{\pgfqpoint{0.845449in}{0.352516in}}%
\pgfpathlineto{\pgfqpoint{0.842638in}{0.355946in}}%
\pgfpathlineto{\pgfqpoint{0.843920in}{0.358581in}}%
\pgfpathlineto{\pgfqpoint{0.839850in}{0.357610in}}%
\pgfpathlineto{\pgfqpoint{0.835041in}{0.357359in}}%
\pgfpathlineto{\pgfqpoint{0.836568in}{0.359811in}}%
\pgfpathlineto{\pgfqpoint{0.838051in}{0.358886in}}%
\pgfpathlineto{\pgfqpoint{0.838232in}{0.362073in}}%
\pgfpathlineto{\pgfqpoint{0.836781in}{0.362125in}}%
\pgfpathlineto{\pgfqpoint{0.832732in}{0.359491in}}%
\pgfpathlineto{\pgfqpoint{0.829761in}{0.356685in}}%
\pgfpathlineto{\pgfqpoint{0.830722in}{0.356227in}}%
\pgfpathlineto{\pgfqpoint{0.832645in}{0.357725in}}%
\pgfpathlineto{\pgfqpoint{0.834243in}{0.356277in}}%
\pgfpathlineto{\pgfqpoint{0.832157in}{0.354595in}}%
\pgfpathlineto{\pgfqpoint{0.830411in}{0.354246in}}%
\pgfpathlineto{\pgfqpoint{0.827722in}{0.355289in}}%
\pgfpathlineto{\pgfqpoint{0.828994in}{0.351862in}}%
\pgfpathlineto{\pgfqpoint{0.830872in}{0.350681in}}%
\pgfpathlineto{\pgfqpoint{0.825792in}{0.348726in}}%
\pgfpathlineto{\pgfqpoint{0.825045in}{0.347737in}}%
\pgfpathlineto{\pgfqpoint{0.828029in}{0.344167in}}%
\pgfpathlineto{\pgfqpoint{0.827087in}{0.340105in}}%
\pgfpathlineto{\pgfqpoint{0.825414in}{0.340100in}}%
\pgfpathlineto{\pgfqpoint{0.823966in}{0.342216in}}%
\pgfpathlineto{\pgfqpoint{0.824700in}{0.344679in}}%
\pgfpathlineto{\pgfqpoint{0.823443in}{0.345830in}}%
\pgfpathlineto{\pgfqpoint{0.821225in}{0.343325in}}%
\pgfpathlineto{\pgfqpoint{0.819433in}{0.344425in}}%
\pgfpathlineto{\pgfqpoint{0.821535in}{0.347132in}}%
\pgfpathlineto{\pgfqpoint{0.825526in}{0.353057in}}%
\pgfpathlineto{\pgfqpoint{0.824494in}{0.353761in}}%
\pgfpathlineto{\pgfqpoint{0.828821in}{0.359977in}}%
\pgfpathlineto{\pgfqpoint{0.827922in}{0.360561in}}%
\pgfpathlineto{\pgfqpoint{0.829442in}{0.362834in}}%
\pgfpathlineto{\pgfqpoint{0.832191in}{0.361126in}}%
\pgfpathlineto{\pgfqpoint{0.836391in}{0.367548in}}%
\pgfpathlineto{\pgfqpoint{0.840140in}{0.372751in}}%
\pgfpathlineto{\pgfqpoint{0.852394in}{0.364648in}}%
\pgfpathlineto{\pgfqpoint{0.853061in}{0.365671in}}%
\pgfpathlineto{\pgfqpoint{0.855608in}{0.363995in}}%
\pgfpathlineto{\pgfqpoint{0.864498in}{0.378218in}}%
\pgfpathlineto{\pgfqpoint{0.865448in}{0.380135in}}%
\pgfpathlineto{\pgfqpoint{0.870796in}{0.376637in}}%
\pgfpathlineto{\pgfqpoint{0.878779in}{0.389457in}}%
\pgfpathlineto{\pgfqpoint{0.880633in}{0.393554in}}%
\pgfpathclose%
\pgfusepath{fill}%
\end{pgfscope}%
\begin{pgfscope}%
\pgfpathrectangle{\pgfqpoint{0.100000in}{0.100000in}}{\pgfqpoint{3.007045in}{1.925000in}}%
\pgfusepath{clip}%
\pgfsetbuttcap%
\pgfsetmiterjoin%
\definecolor{currentfill}{rgb}{0.395294,0.739408,0.433526}%
\pgfsetfillcolor{currentfill}%
\pgfsetlinewidth{0.000000pt}%
\definecolor{currentstroke}{rgb}{0.000000,0.000000,0.000000}%
\pgfsetstrokecolor{currentstroke}%
\pgfsetstrokeopacity{0.000000}%
\pgfsetdash{}{0pt}%
\pgfpathmoveto{\pgfqpoint{2.453185in}{1.084073in}}%
\pgfpathlineto{\pgfqpoint{2.456556in}{1.087012in}}%
\pgfpathlineto{\pgfqpoint{2.458232in}{1.094688in}}%
\pgfpathlineto{\pgfqpoint{2.463061in}{1.098120in}}%
\pgfpathlineto{\pgfqpoint{2.459342in}{1.100923in}}%
\pgfpathlineto{\pgfqpoint{2.471855in}{1.105407in}}%
\pgfpathlineto{\pgfqpoint{2.476705in}{1.103396in}}%
\pgfpathlineto{\pgfqpoint{2.484474in}{1.093009in}}%
\pgfpathlineto{\pgfqpoint{2.489686in}{1.089390in}}%
\pgfpathlineto{\pgfqpoint{2.487921in}{1.088350in}}%
\pgfpathlineto{\pgfqpoint{2.483702in}{1.075698in}}%
\pgfpathlineto{\pgfqpoint{2.475445in}{1.067081in}}%
\pgfpathlineto{\pgfqpoint{2.469086in}{1.065710in}}%
\pgfpathlineto{\pgfqpoint{2.462901in}{1.070185in}}%
\pgfpathlineto{\pgfqpoint{2.455254in}{1.073522in}}%
\pgfpathlineto{\pgfqpoint{2.452938in}{1.079250in}}%
\pgfpathlineto{\pgfqpoint{2.453185in}{1.084073in}}%
\pgfpathclose%
\pgfusepath{fill}%
\end{pgfscope}%
\begin{pgfscope}%
\pgfpathrectangle{\pgfqpoint{0.100000in}{0.100000in}}{\pgfqpoint{3.007045in}{1.925000in}}%
\pgfusepath{clip}%
\pgfsetbuttcap%
\pgfsetmiterjoin%
\definecolor{currentfill}{rgb}{0.648320,0.858116,0.624344}%
\pgfsetfillcolor{currentfill}%
\pgfsetlinewidth{0.000000pt}%
\definecolor{currentstroke}{rgb}{0.000000,0.000000,0.000000}%
\pgfsetstrokecolor{currentstroke}%
\pgfsetstrokeopacity{0.000000}%
\pgfsetdash{}{0pt}%
\pgfpathmoveto{\pgfqpoint{2.533154in}{0.479245in}}%
\pgfpathlineto{\pgfqpoint{2.531241in}{0.493155in}}%
\pgfpathlineto{\pgfqpoint{2.521853in}{0.495228in}}%
\pgfpathlineto{\pgfqpoint{2.517736in}{0.501110in}}%
\pgfpathlineto{\pgfqpoint{2.522009in}{0.510090in}}%
\pgfpathlineto{\pgfqpoint{2.512188in}{0.520102in}}%
\pgfpathlineto{\pgfqpoint{2.549990in}{0.525560in}}%
\pgfpathlineto{\pgfqpoint{2.550293in}{0.531506in}}%
\pgfpathlineto{\pgfqpoint{2.547864in}{0.546632in}}%
\pgfpathlineto{\pgfqpoint{2.550313in}{0.542005in}}%
\pgfpathlineto{\pgfqpoint{2.555178in}{0.540975in}}%
\pgfpathlineto{\pgfqpoint{2.557842in}{0.535907in}}%
\pgfpathlineto{\pgfqpoint{2.567244in}{0.529939in}}%
\pgfpathlineto{\pgfqpoint{2.570842in}{0.520187in}}%
\pgfpathlineto{\pgfqpoint{2.576314in}{0.520862in}}%
\pgfpathlineto{\pgfqpoint{2.581927in}{0.518998in}}%
\pgfpathlineto{\pgfqpoint{2.584158in}{0.521537in}}%
\pgfpathlineto{\pgfqpoint{2.592396in}{0.508481in}}%
\pgfpathlineto{\pgfqpoint{2.599396in}{0.502567in}}%
\pgfpathlineto{\pgfqpoint{2.599254in}{0.497364in}}%
\pgfpathlineto{\pgfqpoint{2.602369in}{0.492044in}}%
\pgfpathlineto{\pgfqpoint{2.609094in}{0.445229in}}%
\pgfpathlineto{\pgfqpoint{2.593275in}{0.442838in}}%
\pgfpathlineto{\pgfqpoint{2.592017in}{0.447697in}}%
\pgfpathlineto{\pgfqpoint{2.587084in}{0.454555in}}%
\pgfpathlineto{\pgfqpoint{2.581515in}{0.455951in}}%
\pgfpathlineto{\pgfqpoint{2.580449in}{0.459760in}}%
\pgfpathlineto{\pgfqpoint{2.575198in}{0.465448in}}%
\pgfpathlineto{\pgfqpoint{2.565859in}{0.472579in}}%
\pgfpathlineto{\pgfqpoint{2.562765in}{0.480015in}}%
\pgfpathlineto{\pgfqpoint{2.556944in}{0.479141in}}%
\pgfpathlineto{\pgfqpoint{2.556078in}{0.484967in}}%
\pgfpathlineto{\pgfqpoint{2.538576in}{0.482276in}}%
\pgfpathlineto{\pgfqpoint{2.533154in}{0.479245in}}%
\pgfpathclose%
\pgfusepath{fill}%
\end{pgfscope}%
\begin{pgfscope}%
\pgfpathrectangle{\pgfqpoint{0.100000in}{0.100000in}}{\pgfqpoint{3.007045in}{1.925000in}}%
\pgfusepath{clip}%
\pgfsetbuttcap%
\pgfsetmiterjoin%
\definecolor{currentfill}{rgb}{0.667020,0.865990,0.642553}%
\pgfsetfillcolor{currentfill}%
\pgfsetlinewidth{0.000000pt}%
\definecolor{currentstroke}{rgb}{0.000000,0.000000,0.000000}%
\pgfsetstrokecolor{currentstroke}%
\pgfsetstrokeopacity{0.000000}%
\pgfsetdash{}{0pt}%
\pgfpathmoveto{\pgfqpoint{1.453227in}{1.011793in}}%
\pgfpathlineto{\pgfqpoint{1.459593in}{1.011387in}}%
\pgfpathlineto{\pgfqpoint{1.456990in}{0.978623in}}%
\pgfpathlineto{\pgfqpoint{1.421577in}{0.981072in}}%
\pgfpathlineto{\pgfqpoint{1.419469in}{0.951860in}}%
\pgfpathlineto{\pgfqpoint{1.390781in}{0.954030in}}%
\pgfpathlineto{\pgfqpoint{1.393120in}{0.983287in}}%
\pgfpathlineto{\pgfqpoint{1.348759in}{0.987002in}}%
\pgfpathlineto{\pgfqpoint{1.351668in}{1.019814in}}%
\pgfpathlineto{\pgfqpoint{1.367585in}{1.018108in}}%
\pgfpathlineto{\pgfqpoint{1.401987in}{1.015175in}}%
\pgfpathlineto{\pgfqpoint{1.404039in}{1.041199in}}%
\pgfpathlineto{\pgfqpoint{1.429358in}{1.039255in}}%
\pgfpathlineto{\pgfqpoint{1.427563in}{1.013412in}}%
\pgfpathlineto{\pgfqpoint{1.453227in}{1.011793in}}%
\pgfpathclose%
\pgfusepath{fill}%
\end{pgfscope}%
\begin{pgfscope}%
\pgfpathrectangle{\pgfqpoint{0.100000in}{0.100000in}}{\pgfqpoint{3.007045in}{1.925000in}}%
\pgfusepath{clip}%
\pgfsetbuttcap%
\pgfsetmiterjoin%
\definecolor{currentfill}{rgb}{0.623760,0.847428,0.601584}%
\pgfsetfillcolor{currentfill}%
\pgfsetlinewidth{0.000000pt}%
\definecolor{currentstroke}{rgb}{0.000000,0.000000,0.000000}%
\pgfsetstrokecolor{currentstroke}%
\pgfsetstrokeopacity{0.000000}%
\pgfsetdash{}{0pt}%
\pgfpathmoveto{\pgfqpoint{2.218727in}{1.043400in}}%
\pgfpathlineto{\pgfqpoint{2.223300in}{1.033829in}}%
\pgfpathlineto{\pgfqpoint{2.216520in}{1.033633in}}%
\pgfpathlineto{\pgfqpoint{2.212201in}{1.031659in}}%
\pgfpathlineto{\pgfqpoint{2.209941in}{1.021313in}}%
\pgfpathlineto{\pgfqpoint{2.200660in}{1.024227in}}%
\pgfpathlineto{\pgfqpoint{2.189732in}{1.023554in}}%
\pgfpathlineto{\pgfqpoint{2.179536in}{1.021312in}}%
\pgfpathlineto{\pgfqpoint{2.174509in}{1.028592in}}%
\pgfpathlineto{\pgfqpoint{2.171148in}{1.034806in}}%
\pgfpathlineto{\pgfqpoint{2.174985in}{1.043028in}}%
\pgfpathlineto{\pgfqpoint{2.172018in}{1.046217in}}%
\pgfpathlineto{\pgfqpoint{2.170940in}{1.054187in}}%
\pgfpathlineto{\pgfqpoint{2.164342in}{1.057050in}}%
\pgfpathlineto{\pgfqpoint{2.171424in}{1.072069in}}%
\pgfpathlineto{\pgfqpoint{2.172010in}{1.077414in}}%
\pgfpathlineto{\pgfqpoint{2.184734in}{1.072536in}}%
\pgfpathlineto{\pgfqpoint{2.185322in}{1.077996in}}%
\pgfpathlineto{\pgfqpoint{2.197741in}{1.085673in}}%
\pgfpathlineto{\pgfqpoint{2.200736in}{1.080556in}}%
\pgfpathlineto{\pgfqpoint{2.205878in}{1.081086in}}%
\pgfpathlineto{\pgfqpoint{2.205852in}{1.076709in}}%
\pgfpathlineto{\pgfqpoint{2.204976in}{1.073343in}}%
\pgfpathlineto{\pgfqpoint{2.208703in}{1.064571in}}%
\pgfpathlineto{\pgfqpoint{2.207346in}{1.056620in}}%
\pgfpathlineto{\pgfqpoint{2.210809in}{1.047577in}}%
\pgfpathlineto{\pgfqpoint{2.217524in}{1.046026in}}%
\pgfpathlineto{\pgfqpoint{2.218727in}{1.043400in}}%
\pgfpathclose%
\pgfusepath{fill}%
\end{pgfscope}%
\begin{pgfscope}%
\pgfpathrectangle{\pgfqpoint{0.100000in}{0.100000in}}{\pgfqpoint{3.007045in}{1.925000in}}%
\pgfusepath{clip}%
\pgfsetbuttcap%
\pgfsetmiterjoin%
\definecolor{currentfill}{rgb}{0.474279,0.777670,0.478677}%
\pgfsetfillcolor{currentfill}%
\pgfsetlinewidth{0.000000pt}%
\definecolor{currentstroke}{rgb}{0.000000,0.000000,0.000000}%
\pgfsetstrokecolor{currentstroke}%
\pgfsetstrokeopacity{0.000000}%
\pgfsetdash{}{0pt}%
\pgfpathmoveto{\pgfqpoint{1.865473in}{1.453759in}}%
\pgfpathlineto{\pgfqpoint{1.882427in}{1.454160in}}%
\pgfpathlineto{\pgfqpoint{1.893937in}{1.453516in}}%
\pgfpathlineto{\pgfqpoint{1.894560in}{1.431480in}}%
\pgfpathlineto{\pgfqpoint{1.866993in}{1.430735in}}%
\pgfpathlineto{\pgfqpoint{1.837085in}{1.430186in}}%
\pgfpathlineto{\pgfqpoint{1.836747in}{1.453219in}}%
\pgfpathlineto{\pgfqpoint{1.865473in}{1.453759in}}%
\pgfpathclose%
\pgfusepath{fill}%
\end{pgfscope}%
\begin{pgfscope}%
\pgfpathrectangle{\pgfqpoint{0.100000in}{0.100000in}}{\pgfqpoint{3.007045in}{1.925000in}}%
\pgfusepath{clip}%
\pgfsetbuttcap%
\pgfsetmiterjoin%
\definecolor{currentfill}{rgb}{0.769873,0.909296,0.742699}%
\pgfsetfillcolor{currentfill}%
\pgfsetlinewidth{0.000000pt}%
\definecolor{currentstroke}{rgb}{0.000000,0.000000,0.000000}%
\pgfsetstrokecolor{currentstroke}%
\pgfsetstrokeopacity{0.000000}%
\pgfsetdash{}{0pt}%
\pgfpathmoveto{\pgfqpoint{1.674656in}{1.382036in}}%
\pgfpathlineto{\pgfqpoint{1.667711in}{1.384730in}}%
\pgfpathlineto{\pgfqpoint{1.668294in}{1.409172in}}%
\pgfpathlineto{\pgfqpoint{1.656789in}{1.409421in}}%
\pgfpathlineto{\pgfqpoint{1.657257in}{1.431226in}}%
\pgfpathlineto{\pgfqpoint{1.647409in}{1.431503in}}%
\pgfpathlineto{\pgfqpoint{1.648008in}{1.454576in}}%
\pgfpathlineto{\pgfqpoint{1.682306in}{1.453680in}}%
\pgfpathlineto{\pgfqpoint{1.703058in}{1.453389in}}%
\pgfpathlineto{\pgfqpoint{1.702699in}{1.430317in}}%
\pgfpathlineto{\pgfqpoint{1.695709in}{1.430437in}}%
\pgfpathlineto{\pgfqpoint{1.695898in}{1.425999in}}%
\pgfpathlineto{\pgfqpoint{1.699167in}{1.423436in}}%
\pgfpathlineto{\pgfqpoint{1.698770in}{1.417142in}}%
\pgfpathlineto{\pgfqpoint{1.701304in}{1.411915in}}%
\pgfpathlineto{\pgfqpoint{1.702210in}{1.402774in}}%
\pgfpathlineto{\pgfqpoint{1.685257in}{1.403080in}}%
\pgfpathlineto{\pgfqpoint{1.684732in}{1.378050in}}%
\pgfpathlineto{\pgfqpoint{1.674656in}{1.382036in}}%
\pgfpathclose%
\pgfusepath{fill}%
\end{pgfscope}%
\begin{pgfscope}%
\pgfpathrectangle{\pgfqpoint{0.100000in}{0.100000in}}{\pgfqpoint{3.007045in}{1.925000in}}%
\pgfusepath{clip}%
\pgfsetbuttcap%
\pgfsetmiterjoin%
\definecolor{currentfill}{rgb}{0.376471,0.730181,0.424298}%
\pgfsetfillcolor{currentfill}%
\pgfsetlinewidth{0.000000pt}%
\definecolor{currentstroke}{rgb}{0.000000,0.000000,0.000000}%
\pgfsetstrokecolor{currentstroke}%
\pgfsetstrokeopacity{0.000000}%
\pgfsetdash{}{0pt}%
\pgfpathmoveto{\pgfqpoint{1.698475in}{1.199567in}}%
\pgfpathlineto{\pgfqpoint{1.709799in}{1.199381in}}%
\pgfpathlineto{\pgfqpoint{1.709420in}{1.170767in}}%
\pgfpathlineto{\pgfqpoint{1.652076in}{1.171950in}}%
\pgfpathlineto{\pgfqpoint{1.652796in}{1.200602in}}%
\pgfpathlineto{\pgfqpoint{1.698475in}{1.199567in}}%
\pgfpathclose%
\pgfusepath{fill}%
\end{pgfscope}%
\begin{pgfscope}%
\pgfpathrectangle{\pgfqpoint{0.100000in}{0.100000in}}{\pgfqpoint{3.007045in}{1.925000in}}%
\pgfusepath{clip}%
\pgfsetbuttcap%
\pgfsetmiterjoin%
\definecolor{currentfill}{rgb}{0.445490,0.764014,0.458131}%
\pgfsetfillcolor{currentfill}%
\pgfsetlinewidth{0.000000pt}%
\definecolor{currentstroke}{rgb}{0.000000,0.000000,0.000000}%
\pgfsetstrokecolor{currentstroke}%
\pgfsetstrokeopacity{0.000000}%
\pgfsetdash{}{0pt}%
\pgfpathmoveto{\pgfqpoint{2.013813in}{1.217404in}}%
\pgfpathlineto{\pgfqpoint{2.001133in}{1.217077in}}%
\pgfpathlineto{\pgfqpoint{2.000904in}{1.222841in}}%
\pgfpathlineto{\pgfqpoint{1.977807in}{1.222266in}}%
\pgfpathlineto{\pgfqpoint{1.977034in}{1.245635in}}%
\pgfpathlineto{\pgfqpoint{1.982776in}{1.245666in}}%
\pgfpathlineto{\pgfqpoint{1.981668in}{1.274180in}}%
\pgfpathlineto{\pgfqpoint{1.998849in}{1.274698in}}%
\pgfpathlineto{\pgfqpoint{1.998647in}{1.280445in}}%
\pgfpathlineto{\pgfqpoint{2.021108in}{1.281507in}}%
\pgfpathlineto{\pgfqpoint{2.021691in}{1.270010in}}%
\pgfpathlineto{\pgfqpoint{2.022785in}{1.247040in}}%
\pgfpathlineto{\pgfqpoint{2.028561in}{1.247285in}}%
\pgfpathlineto{\pgfqpoint{2.029027in}{1.239956in}}%
\pgfpathlineto{\pgfqpoint{2.027473in}{1.235391in}}%
\pgfpathlineto{\pgfqpoint{2.021248in}{1.230411in}}%
\pgfpathlineto{\pgfqpoint{2.017601in}{1.220842in}}%
\pgfpathlineto{\pgfqpoint{2.013813in}{1.217404in}}%
\pgfpathclose%
\pgfusepath{fill}%
\end{pgfscope}%
\begin{pgfscope}%
\pgfpathrectangle{\pgfqpoint{0.100000in}{0.100000in}}{\pgfqpoint{3.007045in}{1.925000in}}%
\pgfusepath{clip}%
\pgfsetbuttcap%
\pgfsetmiterjoin%
\definecolor{currentfill}{rgb}{0.681046,0.871895,0.656209}%
\pgfsetfillcolor{currentfill}%
\pgfsetlinewidth{0.000000pt}%
\definecolor{currentstroke}{rgb}{0.000000,0.000000,0.000000}%
\pgfsetstrokecolor{currentstroke}%
\pgfsetstrokeopacity{0.000000}%
\pgfsetdash{}{0pt}%
\pgfpathmoveto{\pgfqpoint{2.291623in}{1.065687in}}%
\pgfpathlineto{\pgfqpoint{2.292317in}{1.056941in}}%
\pgfpathlineto{\pgfqpoint{2.285344in}{1.059690in}}%
\pgfpathlineto{\pgfqpoint{2.281792in}{1.056352in}}%
\pgfpathlineto{\pgfqpoint{2.276076in}{1.059145in}}%
\pgfpathlineto{\pgfqpoint{2.270021in}{1.058081in}}%
\pgfpathlineto{\pgfqpoint{2.267818in}{1.063028in}}%
\pgfpathlineto{\pgfqpoint{2.262113in}{1.062729in}}%
\pgfpathlineto{\pgfqpoint{2.261909in}{1.072347in}}%
\pgfpathlineto{\pgfqpoint{2.259515in}{1.072291in}}%
\pgfpathlineto{\pgfqpoint{2.253561in}{1.079316in}}%
\pgfpathlineto{\pgfqpoint{2.258362in}{1.081298in}}%
\pgfpathlineto{\pgfqpoint{2.265212in}{1.092380in}}%
\pgfpathlineto{\pgfqpoint{2.270164in}{1.089523in}}%
\pgfpathlineto{\pgfqpoint{2.281936in}{1.092901in}}%
\pgfpathlineto{\pgfqpoint{2.283241in}{1.088142in}}%
\pgfpathlineto{\pgfqpoint{2.289598in}{1.088435in}}%
\pgfpathlineto{\pgfqpoint{2.291416in}{1.086166in}}%
\pgfpathlineto{\pgfqpoint{2.291623in}{1.065687in}}%
\pgfpathclose%
\pgfusepath{fill}%
\end{pgfscope}%
\begin{pgfscope}%
\pgfpathrectangle{\pgfqpoint{0.100000in}{0.100000in}}{\pgfqpoint{3.007045in}{1.925000in}}%
\pgfusepath{clip}%
\pgfsetbuttcap%
\pgfsetmiterjoin%
\definecolor{currentfill}{rgb}{0.779223,0.913233,0.751803}%
\pgfsetfillcolor{currentfill}%
\pgfsetlinewidth{0.000000pt}%
\definecolor{currentstroke}{rgb}{0.000000,0.000000,0.000000}%
\pgfsetstrokecolor{currentstroke}%
\pgfsetstrokeopacity{0.000000}%
\pgfsetdash{}{0pt}%
\pgfpathmoveto{\pgfqpoint{2.406992in}{0.858796in}}%
\pgfpathlineto{\pgfqpoint{2.409663in}{0.861811in}}%
\pgfpathlineto{\pgfqpoint{2.422010in}{0.860947in}}%
\pgfpathlineto{\pgfqpoint{2.421120in}{0.857223in}}%
\pgfpathlineto{\pgfqpoint{2.431219in}{0.848011in}}%
\pgfpathlineto{\pgfqpoint{2.440442in}{0.846333in}}%
\pgfpathlineto{\pgfqpoint{2.435742in}{0.842463in}}%
\pgfpathlineto{\pgfqpoint{2.430794in}{0.832351in}}%
\pgfpathlineto{\pgfqpoint{2.430059in}{0.826283in}}%
\pgfpathlineto{\pgfqpoint{2.424340in}{0.826019in}}%
\pgfpathlineto{\pgfqpoint{2.414961in}{0.829005in}}%
\pgfpathlineto{\pgfqpoint{2.408663in}{0.823968in}}%
\pgfpathlineto{\pgfqpoint{2.403550in}{0.830476in}}%
\pgfpathlineto{\pgfqpoint{2.392091in}{0.836488in}}%
\pgfpathlineto{\pgfqpoint{2.387448in}{0.834640in}}%
\pgfpathlineto{\pgfqpoint{2.380905in}{0.840878in}}%
\pgfpathlineto{\pgfqpoint{2.382703in}{0.849302in}}%
\pgfpathlineto{\pgfqpoint{2.386057in}{0.845572in}}%
\pgfpathlineto{\pgfqpoint{2.397316in}{0.844754in}}%
\pgfpathlineto{\pgfqpoint{2.399267in}{0.840607in}}%
\pgfpathlineto{\pgfqpoint{2.406734in}{0.847180in}}%
\pgfpathlineto{\pgfqpoint{2.407803in}{0.851715in}}%
\pgfpathlineto{\pgfqpoint{2.404649in}{0.856728in}}%
\pgfpathlineto{\pgfqpoint{2.406992in}{0.858796in}}%
\pgfpathclose%
\pgfusepath{fill}%
\end{pgfscope}%
\begin{pgfscope}%
\pgfpathrectangle{\pgfqpoint{0.100000in}{0.100000in}}{\pgfqpoint{3.007045in}{1.925000in}}%
\pgfusepath{clip}%
\pgfsetbuttcap%
\pgfsetmiterjoin%
\definecolor{currentfill}{rgb}{0.389020,0.736332,0.430450}%
\pgfsetfillcolor{currentfill}%
\pgfsetlinewidth{0.000000pt}%
\definecolor{currentstroke}{rgb}{0.000000,0.000000,0.000000}%
\pgfsetstrokecolor{currentstroke}%
\pgfsetstrokeopacity{0.000000}%
\pgfsetdash{}{0pt}%
\pgfpathmoveto{\pgfqpoint{1.565129in}{1.341345in}}%
\pgfpathlineto{\pgfqpoint{1.543010in}{1.342313in}}%
\pgfpathlineto{\pgfqpoint{1.519516in}{1.343574in}}%
\pgfpathlineto{\pgfqpoint{1.520062in}{1.366510in}}%
\pgfpathlineto{\pgfqpoint{1.521365in}{1.389410in}}%
\pgfpathlineto{\pgfqpoint{1.521452in}{1.403789in}}%
\pgfpathlineto{\pgfqpoint{1.566978in}{1.401393in}}%
\pgfpathlineto{\pgfqpoint{1.565129in}{1.341345in}}%
\pgfpathclose%
\pgfusepath{fill}%
\end{pgfscope}%
\begin{pgfscope}%
\pgfpathrectangle{\pgfqpoint{0.100000in}{0.100000in}}{\pgfqpoint{3.007045in}{1.925000in}}%
\pgfusepath{clip}%
\pgfsetbuttcap%
\pgfsetmiterjoin%
\definecolor{currentfill}{rgb}{0.607151,0.839677,0.587928}%
\pgfsetfillcolor{currentfill}%
\pgfsetlinewidth{0.000000pt}%
\definecolor{currentstroke}{rgb}{0.000000,0.000000,0.000000}%
\pgfsetstrokecolor{currentstroke}%
\pgfsetstrokeopacity{0.000000}%
\pgfsetdash{}{0pt}%
\pgfpathmoveto{\pgfqpoint{1.716390in}{1.322247in}}%
\pgfpathlineto{\pgfqpoint{1.719990in}{1.317507in}}%
\pgfpathlineto{\pgfqpoint{1.717094in}{1.310143in}}%
\pgfpathlineto{\pgfqpoint{1.701191in}{1.310424in}}%
\pgfpathlineto{\pgfqpoint{1.701276in}{1.314251in}}%
\pgfpathlineto{\pgfqpoint{1.678558in}{1.314735in}}%
\pgfpathlineto{\pgfqpoint{1.672961in}{1.314877in}}%
\pgfpathlineto{\pgfqpoint{1.673497in}{1.337786in}}%
\pgfpathlineto{\pgfqpoint{1.683057in}{1.337584in}}%
\pgfpathlineto{\pgfqpoint{1.683358in}{1.349049in}}%
\pgfpathlineto{\pgfqpoint{1.688055in}{1.349861in}}%
\pgfpathlineto{\pgfqpoint{1.706125in}{1.349454in}}%
\pgfpathlineto{\pgfqpoint{1.706612in}{1.341900in}}%
\pgfpathlineto{\pgfqpoint{1.713888in}{1.330082in}}%
\pgfpathlineto{\pgfqpoint{1.716907in}{1.329278in}}%
\pgfpathlineto{\pgfqpoint{1.716390in}{1.322247in}}%
\pgfpathclose%
\pgfusepath{fill}%
\end{pgfscope}%
\begin{pgfscope}%
\pgfpathrectangle{\pgfqpoint{0.100000in}{0.100000in}}{\pgfqpoint{3.007045in}{1.925000in}}%
\pgfusepath{clip}%
\pgfsetbuttcap%
\pgfsetmiterjoin%
\definecolor{currentfill}{rgb}{0.755848,0.903391,0.729043}%
\pgfsetfillcolor{currentfill}%
\pgfsetlinewidth{0.000000pt}%
\definecolor{currentstroke}{rgb}{0.000000,0.000000,0.000000}%
\pgfsetstrokecolor{currentstroke}%
\pgfsetstrokeopacity{0.000000}%
\pgfsetdash{}{0pt}%
\pgfpathmoveto{\pgfqpoint{1.898889in}{1.500432in}}%
\pgfpathlineto{\pgfqpoint{1.889296in}{1.502422in}}%
\pgfpathlineto{\pgfqpoint{1.884271in}{1.505208in}}%
\pgfpathlineto{\pgfqpoint{1.876275in}{1.505700in}}%
\pgfpathlineto{\pgfqpoint{1.876437in}{1.499997in}}%
\pgfpathlineto{\pgfqpoint{1.864831in}{1.497891in}}%
\pgfpathlineto{\pgfqpoint{1.864898in}{1.494975in}}%
\pgfpathlineto{\pgfqpoint{1.853444in}{1.494748in}}%
\pgfpathlineto{\pgfqpoint{1.853375in}{1.499519in}}%
\pgfpathlineto{\pgfqpoint{1.841921in}{1.499291in}}%
\pgfpathlineto{\pgfqpoint{1.823694in}{1.499014in}}%
\pgfpathlineto{\pgfqpoint{1.824600in}{1.510542in}}%
\pgfpathlineto{\pgfqpoint{1.818879in}{1.510489in}}%
\pgfpathlineto{\pgfqpoint{1.818599in}{1.527791in}}%
\pgfpathlineto{\pgfqpoint{1.807185in}{1.527732in}}%
\pgfpathlineto{\pgfqpoint{1.806790in}{1.547902in}}%
\pgfpathlineto{\pgfqpoint{1.813546in}{1.550942in}}%
\pgfpathlineto{\pgfqpoint{1.816671in}{1.557157in}}%
\pgfpathlineto{\pgfqpoint{1.830410in}{1.548866in}}%
\pgfpathlineto{\pgfqpoint{1.838457in}{1.548803in}}%
\pgfpathlineto{\pgfqpoint{1.841735in}{1.545791in}}%
\pgfpathlineto{\pgfqpoint{1.841127in}{1.578186in}}%
\pgfpathlineto{\pgfqpoint{1.872272in}{1.578521in}}%
\pgfpathlineto{\pgfqpoint{1.870534in}{1.568228in}}%
\pgfpathlineto{\pgfqpoint{1.875816in}{1.567788in}}%
\pgfpathlineto{\pgfqpoint{1.881731in}{1.559757in}}%
\pgfpathlineto{\pgfqpoint{1.881663in}{1.556766in}}%
\pgfpathlineto{\pgfqpoint{1.876801in}{1.549513in}}%
\pgfpathlineto{\pgfqpoint{1.876967in}{1.544122in}}%
\pgfpathlineto{\pgfqpoint{1.905032in}{1.544903in}}%
\pgfpathlineto{\pgfqpoint{1.906167in}{1.539095in}}%
\pgfpathlineto{\pgfqpoint{1.907419in}{1.500598in}}%
\pgfpathlineto{\pgfqpoint{1.898889in}{1.500432in}}%
\pgfpathclose%
\pgfusepath{fill}%
\end{pgfscope}%
\begin{pgfscope}%
\pgfpathrectangle{\pgfqpoint{0.100000in}{0.100000in}}{\pgfqpoint{3.007045in}{1.925000in}}%
\pgfusepath{clip}%
\pgfsetbuttcap%
\pgfsetmiterjoin%
\definecolor{currentfill}{rgb}{0.760523,0.905359,0.733595}%
\pgfsetfillcolor{currentfill}%
\pgfsetlinewidth{0.000000pt}%
\definecolor{currentstroke}{rgb}{0.000000,0.000000,0.000000}%
\pgfsetstrokecolor{currentstroke}%
\pgfsetstrokeopacity{0.000000}%
\pgfsetdash{}{0pt}%
\pgfpathmoveto{\pgfqpoint{2.037435in}{0.611183in}}%
\pgfpathlineto{\pgfqpoint{2.043709in}{0.592924in}}%
\pgfpathlineto{\pgfqpoint{2.046065in}{0.559905in}}%
\pgfpathlineto{\pgfqpoint{2.042685in}{0.564352in}}%
\pgfpathlineto{\pgfqpoint{2.032903in}{0.562233in}}%
\pgfpathlineto{\pgfqpoint{2.029617in}{0.558203in}}%
\pgfpathlineto{\pgfqpoint{2.023323in}{0.554899in}}%
\pgfpathlineto{\pgfqpoint{2.007733in}{0.550626in}}%
\pgfpathlineto{\pgfqpoint{2.005096in}{0.547382in}}%
\pgfpathlineto{\pgfqpoint{2.002634in}{0.548120in}}%
\pgfpathlineto{\pgfqpoint{1.990185in}{0.543999in}}%
\pgfpathlineto{\pgfqpoint{1.976175in}{0.548583in}}%
\pgfpathlineto{\pgfqpoint{1.975186in}{0.557014in}}%
\pgfpathlineto{\pgfqpoint{1.970123in}{0.557651in}}%
\pgfpathlineto{\pgfqpoint{1.966602in}{0.562051in}}%
\pgfpathlineto{\pgfqpoint{1.965180in}{0.570849in}}%
\pgfpathlineto{\pgfqpoint{1.959526in}{0.575009in}}%
\pgfpathlineto{\pgfqpoint{1.958167in}{0.578691in}}%
\pgfpathlineto{\pgfqpoint{1.959192in}{0.587885in}}%
\pgfpathlineto{\pgfqpoint{1.954185in}{0.597454in}}%
\pgfpathlineto{\pgfqpoint{1.954440in}{0.605869in}}%
\pgfpathlineto{\pgfqpoint{1.958935in}{0.610678in}}%
\pgfpathlineto{\pgfqpoint{1.963991in}{0.607879in}}%
\pgfpathlineto{\pgfqpoint{2.010177in}{0.609791in}}%
\pgfpathlineto{\pgfqpoint{2.037435in}{0.611183in}}%
\pgfpathclose%
\pgfusepath{fill}%
\end{pgfscope}%
\begin{pgfscope}%
\pgfpathrectangle{\pgfqpoint{0.100000in}{0.100000in}}{\pgfqpoint{3.007045in}{1.925000in}}%
\pgfusepath{clip}%
\pgfsetbuttcap%
\pgfsetmiterjoin%
\definecolor{currentfill}{rgb}{0.507497,0.793172,0.505990}%
\pgfsetfillcolor{currentfill}%
\pgfsetlinewidth{0.000000pt}%
\definecolor{currentstroke}{rgb}{0.000000,0.000000,0.000000}%
\pgfsetstrokecolor{currentstroke}%
\pgfsetstrokeopacity{0.000000}%
\pgfsetdash{}{0pt}%
\pgfpathmoveto{\pgfqpoint{1.596180in}{1.004385in}}%
\pgfpathlineto{\pgfqpoint{1.561807in}{1.005883in}}%
\pgfpathlineto{\pgfqpoint{1.562899in}{1.031213in}}%
\pgfpathlineto{\pgfqpoint{1.534548in}{1.032337in}}%
\pgfpathlineto{\pgfqpoint{1.534859in}{1.055660in}}%
\pgfpathlineto{\pgfqpoint{1.563272in}{1.054193in}}%
\pgfpathlineto{\pgfqpoint{1.563552in}{1.060256in}}%
\pgfpathlineto{\pgfqpoint{1.591676in}{1.059068in}}%
\pgfpathlineto{\pgfqpoint{1.591838in}{1.053004in}}%
\pgfpathlineto{\pgfqpoint{1.590957in}{1.030045in}}%
\pgfpathlineto{\pgfqpoint{1.596980in}{1.029821in}}%
\pgfpathlineto{\pgfqpoint{1.596180in}{1.004385in}}%
\pgfpathclose%
\pgfusepath{fill}%
\end{pgfscope}%
\begin{pgfscope}%
\pgfpathrectangle{\pgfqpoint{0.100000in}{0.100000in}}{\pgfqpoint{3.007045in}{1.925000in}}%
\pgfusepath{clip}%
\pgfsetbuttcap%
\pgfsetmiterjoin%
\definecolor{currentfill}{rgb}{0.607151,0.839677,0.587928}%
\pgfsetfillcolor{currentfill}%
\pgfsetlinewidth{0.000000pt}%
\definecolor{currentstroke}{rgb}{0.000000,0.000000,0.000000}%
\pgfsetstrokecolor{currentstroke}%
\pgfsetstrokeopacity{0.000000}%
\pgfsetdash{}{0pt}%
\pgfpathmoveto{\pgfqpoint{2.099536in}{1.444895in}}%
\pgfpathlineto{\pgfqpoint{2.088107in}{1.444073in}}%
\pgfpathlineto{\pgfqpoint{2.087655in}{1.449833in}}%
\pgfpathlineto{\pgfqpoint{2.058743in}{1.447959in}}%
\pgfpathlineto{\pgfqpoint{2.047314in}{1.447880in}}%
\pgfpathlineto{\pgfqpoint{2.046797in}{1.456021in}}%
\pgfpathlineto{\pgfqpoint{2.050014in}{1.463525in}}%
\pgfpathlineto{\pgfqpoint{2.049555in}{1.470541in}}%
\pgfpathlineto{\pgfqpoint{2.029155in}{1.469240in}}%
\pgfpathlineto{\pgfqpoint{2.028120in}{1.486633in}}%
\pgfpathlineto{\pgfqpoint{2.045773in}{1.487575in}}%
\pgfpathlineto{\pgfqpoint{2.043986in}{1.516491in}}%
\pgfpathlineto{\pgfqpoint{2.073048in}{1.518236in}}%
\pgfpathlineto{\pgfqpoint{2.073413in}{1.512435in}}%
\pgfpathlineto{\pgfqpoint{2.090405in}{1.513286in}}%
\pgfpathlineto{\pgfqpoint{2.092961in}{1.513610in}}%
\pgfpathlineto{\pgfqpoint{2.094574in}{1.490855in}}%
\pgfpathlineto{\pgfqpoint{2.101634in}{1.491321in}}%
\pgfpathlineto{\pgfqpoint{2.103472in}{1.468301in}}%
\pgfpathlineto{\pgfqpoint{2.097758in}{1.467853in}}%
\pgfpathlineto{\pgfqpoint{2.099536in}{1.444895in}}%
\pgfpathclose%
\pgfusepath{fill}%
\end{pgfscope}%
\begin{pgfscope}%
\pgfpathrectangle{\pgfqpoint{0.100000in}{0.100000in}}{\pgfqpoint{3.007045in}{1.925000in}}%
\pgfusepath{clip}%
\pgfsetbuttcap%
\pgfsetmiterjoin%
\definecolor{currentfill}{rgb}{0.699746,0.879769,0.674418}%
\pgfsetfillcolor{currentfill}%
\pgfsetlinewidth{0.000000pt}%
\definecolor{currentstroke}{rgb}{0.000000,0.000000,0.000000}%
\pgfsetstrokecolor{currentstroke}%
\pgfsetstrokeopacity{0.000000}%
\pgfsetdash{}{0pt}%
\pgfpathmoveto{\pgfqpoint{2.639205in}{0.819156in}}%
\pgfpathlineto{\pgfqpoint{2.636694in}{0.812102in}}%
\pgfpathlineto{\pgfqpoint{2.628381in}{0.812335in}}%
\pgfpathlineto{\pgfqpoint{2.623221in}{0.806040in}}%
\pgfpathlineto{\pgfqpoint{2.626021in}{0.803536in}}%
\pgfpathlineto{\pgfqpoint{2.619999in}{0.798689in}}%
\pgfpathlineto{\pgfqpoint{2.618917in}{0.795608in}}%
\pgfpathlineto{\pgfqpoint{2.611802in}{0.789793in}}%
\pgfpathlineto{\pgfqpoint{2.611478in}{0.786111in}}%
\pgfpathlineto{\pgfqpoint{2.605521in}{0.780009in}}%
\pgfpathlineto{\pgfqpoint{2.599364in}{0.777779in}}%
\pgfpathlineto{\pgfqpoint{2.588994in}{0.768511in}}%
\pgfpathlineto{\pgfqpoint{2.581007in}{0.767715in}}%
\pgfpathlineto{\pgfqpoint{2.576021in}{0.772345in}}%
\pgfpathlineto{\pgfqpoint{2.572033in}{0.772743in}}%
\pgfpathlineto{\pgfqpoint{2.567570in}{0.778188in}}%
\pgfpathlineto{\pgfqpoint{2.559434in}{0.779009in}}%
\pgfpathlineto{\pgfqpoint{2.554010in}{0.786819in}}%
\pgfpathlineto{\pgfqpoint{2.548566in}{0.790531in}}%
\pgfpathlineto{\pgfqpoint{2.542055in}{0.797900in}}%
\pgfpathlineto{\pgfqpoint{2.556125in}{0.810180in}}%
\pgfpathlineto{\pgfqpoint{2.570619in}{0.823049in}}%
\pgfpathlineto{\pgfqpoint{2.575563in}{0.818701in}}%
\pgfpathlineto{\pgfqpoint{2.584582in}{0.822960in}}%
\pgfpathlineto{\pgfqpoint{2.584761in}{0.832670in}}%
\pgfpathlineto{\pgfqpoint{2.590755in}{0.837224in}}%
\pgfpathlineto{\pgfqpoint{2.597720in}{0.838651in}}%
\pgfpathlineto{\pgfqpoint{2.612678in}{0.829037in}}%
\pgfpathlineto{\pgfqpoint{2.628686in}{0.826071in}}%
\pgfpathlineto{\pgfqpoint{2.639205in}{0.819156in}}%
\pgfpathclose%
\pgfusepath{fill}%
\end{pgfscope}%
\begin{pgfscope}%
\pgfpathrectangle{\pgfqpoint{0.100000in}{0.100000in}}{\pgfqpoint{3.007045in}{1.925000in}}%
\pgfusepath{clip}%
\pgfsetbuttcap%
\pgfsetmiterjoin%
\definecolor{currentfill}{rgb}{0.794233,0.919262,0.767705}%
\pgfsetfillcolor{currentfill}%
\pgfsetlinewidth{0.000000pt}%
\definecolor{currentstroke}{rgb}{0.000000,0.000000,0.000000}%
\pgfsetstrokecolor{currentstroke}%
\pgfsetstrokeopacity{0.000000}%
\pgfsetdash{}{0pt}%
\pgfpathmoveto{\pgfqpoint{2.675874in}{0.992417in}}%
\pgfpathlineto{\pgfqpoint{2.673153in}{0.980638in}}%
\pgfpathlineto{\pgfqpoint{2.674050in}{0.976293in}}%
\pgfpathlineto{\pgfqpoint{2.667375in}{0.972671in}}%
\pgfpathlineto{\pgfqpoint{2.661577in}{0.967658in}}%
\pgfpathlineto{\pgfqpoint{2.657195in}{0.967322in}}%
\pgfpathlineto{\pgfqpoint{2.653930in}{0.970205in}}%
\pgfpathlineto{\pgfqpoint{2.650737in}{0.964955in}}%
\pgfpathlineto{\pgfqpoint{2.641093in}{0.964571in}}%
\pgfpathlineto{\pgfqpoint{2.625940in}{0.956912in}}%
\pgfpathlineto{\pgfqpoint{2.618329in}{0.960663in}}%
\pgfpathlineto{\pgfqpoint{2.608908in}{0.972810in}}%
\pgfpathlineto{\pgfqpoint{2.608759in}{0.975733in}}%
\pgfpathlineto{\pgfqpoint{2.597898in}{0.973667in}}%
\pgfpathlineto{\pgfqpoint{2.594940in}{0.995133in}}%
\pgfpathlineto{\pgfqpoint{2.611008in}{0.997970in}}%
\pgfpathlineto{\pgfqpoint{2.608855in}{1.001837in}}%
\pgfpathlineto{\pgfqpoint{2.605433in}{1.023762in}}%
\pgfpathlineto{\pgfqpoint{2.610860in}{1.024597in}}%
\pgfpathlineto{\pgfqpoint{2.608257in}{1.044199in}}%
\pgfpathlineto{\pgfqpoint{2.643747in}{1.050644in}}%
\pgfpathlineto{\pgfqpoint{2.641967in}{1.044405in}}%
\pgfpathlineto{\pgfqpoint{2.645335in}{1.026823in}}%
\pgfpathlineto{\pgfqpoint{2.650549in}{1.026900in}}%
\pgfpathlineto{\pgfqpoint{2.654897in}{1.034173in}}%
\pgfpathlineto{\pgfqpoint{2.664233in}{1.034590in}}%
\pgfpathlineto{\pgfqpoint{2.671365in}{1.033005in}}%
\pgfpathlineto{\pgfqpoint{2.663029in}{1.005596in}}%
\pgfpathlineto{\pgfqpoint{2.668059in}{1.000440in}}%
\pgfpathlineto{\pgfqpoint{2.672444in}{0.992955in}}%
\pgfpathlineto{\pgfqpoint{2.675874in}{0.992417in}}%
\pgfpathclose%
\pgfusepath{fill}%
\end{pgfscope}%
\begin{pgfscope}%
\pgfpathrectangle{\pgfqpoint{0.100000in}{0.100000in}}{\pgfqpoint{3.007045in}{1.925000in}}%
\pgfusepath{clip}%
\pgfsetbuttcap%
\pgfsetmiterjoin%
\definecolor{currentfill}{rgb}{0.551788,0.813841,0.542407}%
\pgfsetfillcolor{currentfill}%
\pgfsetlinewidth{0.000000pt}%
\definecolor{currentstroke}{rgb}{0.000000,0.000000,0.000000}%
\pgfsetstrokecolor{currentstroke}%
\pgfsetstrokeopacity{0.000000}%
\pgfsetdash{}{0pt}%
\pgfpathmoveto{\pgfqpoint{1.584352in}{1.202908in}}%
\pgfpathlineto{\pgfqpoint{1.561531in}{1.203886in}}%
\pgfpathlineto{\pgfqpoint{1.563506in}{1.247230in}}%
\pgfpathlineto{\pgfqpoint{1.551658in}{1.248467in}}%
\pgfpathlineto{\pgfqpoint{1.552431in}{1.273213in}}%
\pgfpathlineto{\pgfqpoint{1.563271in}{1.272731in}}%
\pgfpathlineto{\pgfqpoint{1.564548in}{1.295579in}}%
\pgfpathlineto{\pgfqpoint{1.587220in}{1.294590in}}%
\pgfpathlineto{\pgfqpoint{1.586068in}{1.271711in}}%
\pgfpathlineto{\pgfqpoint{1.587386in}{1.271651in}}%
\pgfpathlineto{\pgfqpoint{1.586476in}{1.248770in}}%
\pgfpathlineto{\pgfqpoint{1.584352in}{1.202908in}}%
\pgfpathclose%
\pgfusepath{fill}%
\end{pgfscope}%
\begin{pgfscope}%
\pgfpathrectangle{\pgfqpoint{0.100000in}{0.100000in}}{\pgfqpoint{3.007045in}{1.925000in}}%
\pgfusepath{clip}%
\pgfsetbuttcap%
\pgfsetmiterjoin%
\definecolor{currentfill}{rgb}{0.513033,0.795755,0.510542}%
\pgfsetfillcolor{currentfill}%
\pgfsetlinewidth{0.000000pt}%
\definecolor{currentstroke}{rgb}{0.000000,0.000000,0.000000}%
\pgfsetstrokecolor{currentstroke}%
\pgfsetstrokeopacity{0.000000}%
\pgfsetdash{}{0pt}%
\pgfpathmoveto{\pgfqpoint{0.713573in}{1.866288in}}%
\pgfpathlineto{\pgfqpoint{0.708324in}{1.867692in}}%
\pgfpathlineto{\pgfqpoint{0.702364in}{1.864190in}}%
\pgfpathlineto{\pgfqpoint{0.700822in}{1.858600in}}%
\pgfpathlineto{\pgfqpoint{0.695822in}{1.854965in}}%
\pgfpathlineto{\pgfqpoint{0.691919in}{1.848024in}}%
\pgfpathlineto{\pgfqpoint{0.689171in}{1.848788in}}%
\pgfpathlineto{\pgfqpoint{0.682133in}{1.844866in}}%
\pgfpathlineto{\pgfqpoint{0.680569in}{1.839333in}}%
\pgfpathlineto{\pgfqpoint{0.665649in}{1.843218in}}%
\pgfpathlineto{\pgfqpoint{0.662469in}{1.839071in}}%
\pgfpathlineto{\pgfqpoint{0.657936in}{1.837275in}}%
\pgfpathlineto{\pgfqpoint{0.656028in}{1.830896in}}%
\pgfpathlineto{\pgfqpoint{0.652989in}{1.834668in}}%
\pgfpathlineto{\pgfqpoint{0.640168in}{1.838234in}}%
\pgfpathlineto{\pgfqpoint{0.635771in}{1.843623in}}%
\pgfpathlineto{\pgfqpoint{0.629717in}{1.846964in}}%
\pgfpathlineto{\pgfqpoint{0.627702in}{1.851025in}}%
\pgfpathlineto{\pgfqpoint{0.620178in}{1.855997in}}%
\pgfpathlineto{\pgfqpoint{0.618659in}{1.862354in}}%
\pgfpathlineto{\pgfqpoint{0.615087in}{1.868865in}}%
\pgfpathlineto{\pgfqpoint{0.619147in}{1.875248in}}%
\pgfpathlineto{\pgfqpoint{0.617908in}{1.888997in}}%
\pgfpathlineto{\pgfqpoint{0.622765in}{1.897664in}}%
\pgfpathlineto{\pgfqpoint{0.631834in}{1.900038in}}%
\pgfpathlineto{\pgfqpoint{0.634619in}{1.904375in}}%
\pgfpathlineto{\pgfqpoint{0.630805in}{1.915337in}}%
\pgfpathlineto{\pgfqpoint{0.633546in}{1.923434in}}%
\pgfpathlineto{\pgfqpoint{0.643966in}{1.926557in}}%
\pgfpathlineto{\pgfqpoint{0.650253in}{1.923532in}}%
\pgfpathlineto{\pgfqpoint{0.650154in}{1.917880in}}%
\pgfpathlineto{\pgfqpoint{0.654234in}{1.907098in}}%
\pgfpathlineto{\pgfqpoint{0.659239in}{1.897832in}}%
\pgfpathlineto{\pgfqpoint{0.657890in}{1.895448in}}%
\pgfpathlineto{\pgfqpoint{0.666800in}{1.883212in}}%
\pgfpathlineto{\pgfqpoint{0.667442in}{1.879016in}}%
\pgfpathlineto{\pgfqpoint{0.675310in}{1.876711in}}%
\pgfpathlineto{\pgfqpoint{0.676269in}{1.883139in}}%
\pgfpathlineto{\pgfqpoint{0.683348in}{1.884220in}}%
\pgfpathlineto{\pgfqpoint{0.685273in}{1.876771in}}%
\pgfpathlineto{\pgfqpoint{0.695862in}{1.879153in}}%
\pgfpathlineto{\pgfqpoint{0.699664in}{1.876125in}}%
\pgfpathlineto{\pgfqpoint{0.710511in}{1.880181in}}%
\pgfpathlineto{\pgfqpoint{0.713541in}{1.878149in}}%
\pgfpathlineto{\pgfqpoint{0.713573in}{1.866288in}}%
\pgfpathclose%
\pgfusepath{fill}%
\end{pgfscope}%
\begin{pgfscope}%
\pgfpathrectangle{\pgfqpoint{0.100000in}{0.100000in}}{\pgfqpoint{3.007045in}{1.925000in}}%
\pgfusepath{clip}%
\pgfsetbuttcap%
\pgfsetmiterjoin%
\definecolor{currentfill}{rgb}{0.638970,0.854179,0.615240}%
\pgfsetfillcolor{currentfill}%
\pgfsetlinewidth{0.000000pt}%
\definecolor{currentstroke}{rgb}{0.000000,0.000000,0.000000}%
\pgfsetstrokecolor{currentstroke}%
\pgfsetstrokeopacity{0.000000}%
\pgfsetdash{}{0pt}%
\pgfpathmoveto{\pgfqpoint{2.087534in}{0.559447in}}%
\pgfpathlineto{\pgfqpoint{2.082124in}{0.561840in}}%
\pgfpathlineto{\pgfqpoint{2.077352in}{0.576813in}}%
\pgfpathlineto{\pgfqpoint{2.072286in}{0.580627in}}%
\pgfpathlineto{\pgfqpoint{2.067276in}{0.590315in}}%
\pgfpathlineto{\pgfqpoint{2.067754in}{0.598346in}}%
\pgfpathlineto{\pgfqpoint{2.070931in}{0.604028in}}%
\pgfpathlineto{\pgfqpoint{2.072666in}{0.613237in}}%
\pgfpathlineto{\pgfqpoint{2.076881in}{0.613497in}}%
\pgfpathlineto{\pgfqpoint{2.085520in}{0.614773in}}%
\pgfpathlineto{\pgfqpoint{2.094715in}{0.614481in}}%
\pgfpathlineto{\pgfqpoint{2.095084in}{0.608430in}}%
\pgfpathlineto{\pgfqpoint{2.123918in}{0.610337in}}%
\pgfpathlineto{\pgfqpoint{2.121151in}{0.645024in}}%
\pgfpathlineto{\pgfqpoint{2.143322in}{0.646719in}}%
\pgfpathlineto{\pgfqpoint{2.146781in}{0.617810in}}%
\pgfpathlineto{\pgfqpoint{2.151467in}{0.577078in}}%
\pgfpathlineto{\pgfqpoint{2.147381in}{0.572588in}}%
\pgfpathlineto{\pgfqpoint{2.139247in}{0.575561in}}%
\pgfpathlineto{\pgfqpoint{2.134440in}{0.573344in}}%
\pgfpathlineto{\pgfqpoint{2.121286in}{0.575705in}}%
\pgfpathlineto{\pgfqpoint{2.113119in}{0.573784in}}%
\pgfpathlineto{\pgfqpoint{2.098215in}{0.568019in}}%
\pgfpathlineto{\pgfqpoint{2.087534in}{0.559447in}}%
\pgfpathclose%
\pgfusepath{fill}%
\end{pgfscope}%
\begin{pgfscope}%
\pgfpathrectangle{\pgfqpoint{0.100000in}{0.100000in}}{\pgfqpoint{3.007045in}{1.925000in}}%
\pgfusepath{clip}%
\pgfsetbuttcap%
\pgfsetmiterjoin%
\definecolor{currentfill}{rgb}{0.000000,0.266667,0.105882}%
\pgfsetfillcolor{currentfill}%
\pgfsetlinewidth{0.000000pt}%
\definecolor{currentstroke}{rgb}{0.000000,0.000000,0.000000}%
\pgfsetstrokecolor{currentstroke}%
\pgfsetstrokeopacity{0.000000}%
\pgfsetdash{}{0pt}%
\pgfpathmoveto{\pgfqpoint{2.990796in}{1.458726in}}%
\pgfpathlineto{\pgfqpoint{2.997029in}{1.452912in}}%
\pgfpathlineto{\pgfqpoint{2.995205in}{1.449674in}}%
\pgfpathlineto{\pgfqpoint{2.988017in}{1.448533in}}%
\pgfpathlineto{\pgfqpoint{2.992375in}{1.455283in}}%
\pgfpathlineto{\pgfqpoint{2.990796in}{1.458726in}}%
\pgfpathclose%
\pgfusepath{fill}%
\end{pgfscope}%
\begin{pgfscope}%
\pgfpathrectangle{\pgfqpoint{0.100000in}{0.100000in}}{\pgfqpoint{3.007045in}{1.925000in}}%
\pgfusepath{clip}%
\pgfsetbuttcap%
\pgfsetmiterjoin%
\definecolor{currentfill}{rgb}{0.540715,0.808674,0.533303}%
\pgfsetfillcolor{currentfill}%
\pgfsetlinewidth{0.000000pt}%
\definecolor{currentstroke}{rgb}{0.000000,0.000000,0.000000}%
\pgfsetstrokecolor{currentstroke}%
\pgfsetstrokeopacity{0.000000}%
\pgfsetdash{}{0pt}%
\pgfpathmoveto{\pgfqpoint{2.537236in}{1.043373in}}%
\pgfpathlineto{\pgfqpoint{2.526060in}{1.056473in}}%
\pgfpathlineto{\pgfqpoint{2.521063in}{1.052157in}}%
\pgfpathlineto{\pgfqpoint{2.511594in}{1.060569in}}%
\pgfpathlineto{\pgfqpoint{2.508760in}{1.065976in}}%
\pgfpathlineto{\pgfqpoint{2.512374in}{1.068702in}}%
\pgfpathlineto{\pgfqpoint{2.502786in}{1.075737in}}%
\pgfpathlineto{\pgfqpoint{2.510605in}{1.081728in}}%
\pgfpathlineto{\pgfqpoint{2.509239in}{1.087027in}}%
\pgfpathlineto{\pgfqpoint{2.516736in}{1.108073in}}%
\pgfpathlineto{\pgfqpoint{2.527642in}{1.106923in}}%
\pgfpathlineto{\pgfqpoint{2.535565in}{1.108737in}}%
\pgfpathlineto{\pgfqpoint{2.544389in}{1.103816in}}%
\pgfpathlineto{\pgfqpoint{2.554928in}{1.112611in}}%
\pgfpathlineto{\pgfqpoint{2.559027in}{1.119923in}}%
\pgfpathlineto{\pgfqpoint{2.566770in}{1.118700in}}%
\pgfpathlineto{\pgfqpoint{2.567398in}{1.112031in}}%
\pgfpathlineto{\pgfqpoint{2.578351in}{1.105297in}}%
\pgfpathlineto{\pgfqpoint{2.573073in}{1.099388in}}%
\pgfpathlineto{\pgfqpoint{2.570652in}{1.100774in}}%
\pgfpathlineto{\pgfqpoint{2.564677in}{1.094671in}}%
\pgfpathlineto{\pgfqpoint{2.562917in}{1.088318in}}%
\pgfpathlineto{\pgfqpoint{2.564592in}{1.081776in}}%
\pgfpathlineto{\pgfqpoint{2.574988in}{1.077569in}}%
\pgfpathlineto{\pgfqpoint{2.578887in}{1.072582in}}%
\pgfpathlineto{\pgfqpoint{2.578534in}{1.060078in}}%
\pgfpathlineto{\pgfqpoint{2.572984in}{1.054602in}}%
\pgfpathlineto{\pgfqpoint{2.547424in}{1.055944in}}%
\pgfpathlineto{\pgfqpoint{2.543457in}{1.053394in}}%
\pgfpathlineto{\pgfqpoint{2.537236in}{1.043373in}}%
\pgfpathclose%
\pgfusepath{fill}%
\end{pgfscope}%
\begin{pgfscope}%
\pgfpathrectangle{\pgfqpoint{0.100000in}{0.100000in}}{\pgfqpoint{3.007045in}{1.925000in}}%
\pgfusepath{clip}%
\pgfsetbuttcap%
\pgfsetmiterjoin%
\definecolor{currentfill}{rgb}{0.579469,0.826759,0.565167}%
\pgfsetfillcolor{currentfill}%
\pgfsetlinewidth{0.000000pt}%
\definecolor{currentstroke}{rgb}{0.000000,0.000000,0.000000}%
\pgfsetstrokecolor{currentstroke}%
\pgfsetstrokeopacity{0.000000}%
\pgfsetdash{}{0pt}%
\pgfpathmoveto{\pgfqpoint{1.072883in}{1.296139in}}%
\pgfpathlineto{\pgfqpoint{1.067620in}{1.298977in}}%
\pgfpathlineto{\pgfqpoint{1.068540in}{1.305641in}}%
\pgfpathlineto{\pgfqpoint{1.067332in}{1.310685in}}%
\pgfpathlineto{\pgfqpoint{1.057913in}{1.312141in}}%
\pgfpathlineto{\pgfqpoint{1.057220in}{1.307768in}}%
\pgfpathlineto{\pgfqpoint{1.052269in}{1.305131in}}%
\pgfpathlineto{\pgfqpoint{1.050961in}{1.311223in}}%
\pgfpathlineto{\pgfqpoint{1.033610in}{1.309798in}}%
\pgfpathlineto{\pgfqpoint{1.027613in}{1.314125in}}%
\pgfpathlineto{\pgfqpoint{1.029601in}{1.326060in}}%
\pgfpathlineto{\pgfqpoint{1.005052in}{1.330045in}}%
\pgfpathlineto{\pgfqpoint{0.978339in}{1.334944in}}%
\pgfpathlineto{\pgfqpoint{0.981285in}{1.351407in}}%
\pgfpathlineto{\pgfqpoint{0.972493in}{1.347031in}}%
\pgfpathlineto{\pgfqpoint{0.969963in}{1.352958in}}%
\pgfpathlineto{\pgfqpoint{0.970085in}{1.358692in}}%
\pgfpathlineto{\pgfqpoint{0.963794in}{1.361956in}}%
\pgfpathlineto{\pgfqpoint{0.960688in}{1.366618in}}%
\pgfpathlineto{\pgfqpoint{0.965270in}{1.373029in}}%
\pgfpathlineto{\pgfqpoint{0.968734in}{1.380626in}}%
\pgfpathlineto{\pgfqpoint{0.967550in}{1.392277in}}%
\pgfpathlineto{\pgfqpoint{0.965948in}{1.394443in}}%
\pgfpathlineto{\pgfqpoint{0.968744in}{1.400234in}}%
\pgfpathlineto{\pgfqpoint{0.967728in}{1.403968in}}%
\pgfpathlineto{\pgfqpoint{0.989963in}{1.400079in}}%
\pgfpathlineto{\pgfqpoint{1.001890in}{1.466220in}}%
\pgfpathlineto{\pgfqpoint{1.005297in}{1.485446in}}%
\pgfpathlineto{\pgfqpoint{1.016022in}{1.482462in}}%
\pgfpathlineto{\pgfqpoint{1.015288in}{1.478274in}}%
\pgfpathlineto{\pgfqpoint{1.026443in}{1.476331in}}%
\pgfpathlineto{\pgfqpoint{1.027092in}{1.480000in}}%
\pgfpathlineto{\pgfqpoint{1.038193in}{1.478061in}}%
\pgfpathlineto{\pgfqpoint{1.039144in}{1.483711in}}%
\pgfpathlineto{\pgfqpoint{1.052724in}{1.481408in}}%
\pgfpathlineto{\pgfqpoint{1.053944in}{1.487032in}}%
\pgfpathlineto{\pgfqpoint{1.067969in}{1.484605in}}%
\pgfpathlineto{\pgfqpoint{1.069796in}{1.477757in}}%
\pgfpathlineto{\pgfqpoint{1.069165in}{1.468405in}}%
\pgfpathlineto{\pgfqpoint{1.070708in}{1.461628in}}%
\pgfpathlineto{\pgfqpoint{1.075159in}{1.453148in}}%
\pgfpathlineto{\pgfqpoint{1.079954in}{1.447589in}}%
\pgfpathlineto{\pgfqpoint{1.085388in}{1.435891in}}%
\pgfpathlineto{\pgfqpoint{1.092414in}{1.429227in}}%
\pgfpathlineto{\pgfqpoint{1.089929in}{1.412247in}}%
\pgfpathlineto{\pgfqpoint{1.089437in}{1.400986in}}%
\pgfpathlineto{\pgfqpoint{1.139789in}{1.393330in}}%
\pgfpathlineto{\pgfqpoint{1.162889in}{1.390045in}}%
\pgfpathlineto{\pgfqpoint{1.162230in}{1.378562in}}%
\pgfpathlineto{\pgfqpoint{1.158048in}{1.350480in}}%
\pgfpathlineto{\pgfqpoint{1.137581in}{1.353507in}}%
\pgfpathlineto{\pgfqpoint{1.135759in}{1.336378in}}%
\pgfpathlineto{\pgfqpoint{1.131992in}{1.310519in}}%
\pgfpathlineto{\pgfqpoint{1.102069in}{1.314741in}}%
\pgfpathlineto{\pgfqpoint{1.076315in}{1.318750in}}%
\pgfpathlineto{\pgfqpoint{1.072883in}{1.296139in}}%
\pgfpathclose%
\pgfusepath{fill}%
\end{pgfscope}%
\begin{pgfscope}%
\pgfpathrectangle{\pgfqpoint{0.100000in}{0.100000in}}{\pgfqpoint{3.007045in}{1.925000in}}%
\pgfusepath{clip}%
\pgfsetbuttcap%
\pgfsetmiterjoin%
\definecolor{currentfill}{rgb}{0.457670,0.769919,0.465021}%
\pgfsetfillcolor{currentfill}%
\pgfsetlinewidth{0.000000pt}%
\definecolor{currentstroke}{rgb}{0.000000,0.000000,0.000000}%
\pgfsetstrokecolor{currentstroke}%
\pgfsetstrokeopacity{0.000000}%
\pgfsetdash{}{0pt}%
\pgfpathmoveto{\pgfqpoint{2.117083in}{1.446231in}}%
\pgfpathlineto{\pgfqpoint{2.099536in}{1.444895in}}%
\pgfpathlineto{\pgfqpoint{2.097758in}{1.467853in}}%
\pgfpathlineto{\pgfqpoint{2.103472in}{1.468301in}}%
\pgfpathlineto{\pgfqpoint{2.101634in}{1.491321in}}%
\pgfpathlineto{\pgfqpoint{2.108968in}{1.491853in}}%
\pgfpathlineto{\pgfqpoint{2.108536in}{1.497608in}}%
\pgfpathlineto{\pgfqpoint{2.124853in}{1.498938in}}%
\pgfpathlineto{\pgfqpoint{2.126881in}{1.489860in}}%
\pgfpathlineto{\pgfqpoint{2.121101in}{1.483818in}}%
\pgfpathlineto{\pgfqpoint{2.118059in}{1.468764in}}%
\pgfpathlineto{\pgfqpoint{2.120658in}{1.456120in}}%
\pgfpathlineto{\pgfqpoint{2.117083in}{1.446231in}}%
\pgfpathclose%
\pgfusepath{fill}%
\end{pgfscope}%
\begin{pgfscope}%
\pgfpathrectangle{\pgfqpoint{0.100000in}{0.100000in}}{\pgfqpoint{3.007045in}{1.925000in}}%
\pgfusepath{clip}%
\pgfsetbuttcap%
\pgfsetmiterjoin%
\definecolor{currentfill}{rgb}{0.779223,0.913233,0.751803}%
\pgfsetfillcolor{currentfill}%
\pgfsetlinewidth{0.000000pt}%
\definecolor{currentstroke}{rgb}{0.000000,0.000000,0.000000}%
\pgfsetstrokecolor{currentstroke}%
\pgfsetstrokeopacity{0.000000}%
\pgfsetdash{}{0pt}%
\pgfpathmoveto{\pgfqpoint{1.330767in}{0.811238in}}%
\pgfpathlineto{\pgfqpoint{1.333803in}{0.843365in}}%
\pgfpathlineto{\pgfqpoint{1.361996in}{0.840962in}}%
\pgfpathlineto{\pgfqpoint{1.364391in}{0.869531in}}%
\pgfpathlineto{\pgfqpoint{1.392891in}{0.867332in}}%
\pgfpathlineto{\pgfqpoint{1.390704in}{0.838664in}}%
\pgfpathlineto{\pgfqpoint{1.385676in}{0.839061in}}%
\pgfpathlineto{\pgfqpoint{1.383417in}{0.806844in}}%
\pgfpathlineto{\pgfqpoint{1.330767in}{0.811238in}}%
\pgfpathclose%
\pgfusepath{fill}%
\end{pgfscope}%
\begin{pgfscope}%
\pgfpathrectangle{\pgfqpoint{0.100000in}{0.100000in}}{\pgfqpoint{3.007045in}{1.925000in}}%
\pgfusepath{clip}%
\pgfsetbuttcap%
\pgfsetmiterjoin%
\definecolor{currentfill}{rgb}{0.363922,0.724029,0.418147}%
\pgfsetfillcolor{currentfill}%
\pgfsetlinewidth{0.000000pt}%
\definecolor{currentstroke}{rgb}{0.000000,0.000000,0.000000}%
\pgfsetstrokecolor{currentstroke}%
\pgfsetstrokeopacity{0.000000}%
\pgfsetdash{}{0pt}%
\pgfpathmoveto{\pgfqpoint{0.481516in}{1.783693in}}%
\pgfpathlineto{\pgfqpoint{0.465616in}{1.788403in}}%
\pgfpathlineto{\pgfqpoint{0.454831in}{1.792649in}}%
\pgfpathlineto{\pgfqpoint{0.462406in}{1.810296in}}%
\pgfpathlineto{\pgfqpoint{0.462308in}{1.820138in}}%
\pgfpathlineto{\pgfqpoint{0.467265in}{1.814731in}}%
\pgfpathlineto{\pgfqpoint{0.473024in}{1.816689in}}%
\pgfpathlineto{\pgfqpoint{0.479649in}{1.816372in}}%
\pgfpathlineto{\pgfqpoint{0.476951in}{1.821800in}}%
\pgfpathlineto{\pgfqpoint{0.468427in}{1.820044in}}%
\pgfpathlineto{\pgfqpoint{0.462073in}{1.828534in}}%
\pgfpathlineto{\pgfqpoint{0.465398in}{1.837277in}}%
\pgfpathlineto{\pgfqpoint{0.473406in}{1.837959in}}%
\pgfpathlineto{\pgfqpoint{0.472606in}{1.850993in}}%
\pgfpathlineto{\pgfqpoint{0.469231in}{1.854603in}}%
\pgfpathlineto{\pgfqpoint{0.471018in}{1.865773in}}%
\pgfpathlineto{\pgfqpoint{0.474450in}{1.863703in}}%
\pgfpathlineto{\pgfqpoint{0.482840in}{1.866326in}}%
\pgfpathlineto{\pgfqpoint{0.472793in}{1.873445in}}%
\pgfpathlineto{\pgfqpoint{0.474678in}{1.891400in}}%
\pgfpathlineto{\pgfqpoint{0.472885in}{1.900270in}}%
\pgfpathlineto{\pgfqpoint{0.474818in}{1.908630in}}%
\pgfpathlineto{\pgfqpoint{0.487655in}{1.903257in}}%
\pgfpathlineto{\pgfqpoint{0.510848in}{1.895622in}}%
\pgfpathlineto{\pgfqpoint{0.505442in}{1.879195in}}%
\pgfpathlineto{\pgfqpoint{0.502572in}{1.867867in}}%
\pgfpathlineto{\pgfqpoint{0.514994in}{1.864013in}}%
\pgfpathlineto{\pgfqpoint{0.515078in}{1.857792in}}%
\pgfpathlineto{\pgfqpoint{0.510872in}{1.844976in}}%
\pgfpathlineto{\pgfqpoint{0.501743in}{1.847839in}}%
\pgfpathlineto{\pgfqpoint{0.493956in}{1.821847in}}%
\pgfpathlineto{\pgfqpoint{0.542761in}{1.806508in}}%
\pgfpathlineto{\pgfqpoint{0.536092in}{1.785540in}}%
\pgfpathlineto{\pgfqpoint{0.529311in}{1.780896in}}%
\pgfpathlineto{\pgfqpoint{0.524572in}{1.784544in}}%
\pgfpathlineto{\pgfqpoint{0.510888in}{1.782623in}}%
\pgfpathlineto{\pgfqpoint{0.507744in}{1.796488in}}%
\pgfpathlineto{\pgfqpoint{0.498370in}{1.806808in}}%
\pgfpathlineto{\pgfqpoint{0.488870in}{1.806888in}}%
\pgfpathlineto{\pgfqpoint{0.481516in}{1.783693in}}%
\pgfpathclose%
\pgfusepath{fill}%
\end{pgfscope}%
\begin{pgfscope}%
\pgfpathrectangle{\pgfqpoint{0.100000in}{0.100000in}}{\pgfqpoint{3.007045in}{1.925000in}}%
\pgfusepath{clip}%
\pgfsetbuttcap%
\pgfsetmiterjoin%
\definecolor{currentfill}{rgb}{0.540715,0.808674,0.533303}%
\pgfsetfillcolor{currentfill}%
\pgfsetlinewidth{0.000000pt}%
\definecolor{currentstroke}{rgb}{0.000000,0.000000,0.000000}%
\pgfsetstrokecolor{currentstroke}%
\pgfsetstrokeopacity{0.000000}%
\pgfsetdash{}{0pt}%
\pgfpathmoveto{\pgfqpoint{0.812473in}{1.770492in}}%
\pgfpathlineto{\pgfqpoint{0.807745in}{1.750662in}}%
\pgfpathlineto{\pgfqpoint{0.794807in}{1.753888in}}%
\pgfpathlineto{\pgfqpoint{0.789165in}{1.749329in}}%
\pgfpathlineto{\pgfqpoint{0.774865in}{1.752771in}}%
\pgfpathlineto{\pgfqpoint{0.773003in}{1.745268in}}%
\pgfpathlineto{\pgfqpoint{0.765843in}{1.746677in}}%
\pgfpathlineto{\pgfqpoint{0.765426in}{1.755324in}}%
\pgfpathlineto{\pgfqpoint{0.760765in}{1.761473in}}%
\pgfpathlineto{\pgfqpoint{0.758923in}{1.767006in}}%
\pgfpathlineto{\pgfqpoint{0.752113in}{1.767443in}}%
\pgfpathlineto{\pgfqpoint{0.748808in}{1.770376in}}%
\pgfpathlineto{\pgfqpoint{0.740246in}{1.767188in}}%
\pgfpathlineto{\pgfqpoint{0.724939in}{1.766494in}}%
\pgfpathlineto{\pgfqpoint{0.724517in}{1.778735in}}%
\pgfpathlineto{\pgfqpoint{0.726344in}{1.778628in}}%
\pgfpathlineto{\pgfqpoint{0.735917in}{1.779747in}}%
\pgfpathlineto{\pgfqpoint{0.740338in}{1.787082in}}%
\pgfpathlineto{\pgfqpoint{0.746029in}{1.809361in}}%
\pgfpathlineto{\pgfqpoint{0.786311in}{1.799097in}}%
\pgfpathlineto{\pgfqpoint{0.784213in}{1.790565in}}%
\pgfpathlineto{\pgfqpoint{0.794601in}{1.781859in}}%
\pgfpathlineto{\pgfqpoint{0.808364in}{1.778465in}}%
\pgfpathlineto{\pgfqpoint{0.813833in}{1.776182in}}%
\pgfpathlineto{\pgfqpoint{0.812473in}{1.770492in}}%
\pgfpathclose%
\pgfusepath{fill}%
\end{pgfscope}%
\begin{pgfscope}%
\pgfpathrectangle{\pgfqpoint{0.100000in}{0.100000in}}{\pgfqpoint{3.007045in}{1.925000in}}%
\pgfusepath{clip}%
\pgfsetbuttcap%
\pgfsetmiterjoin%
\definecolor{currentfill}{rgb}{0.451765,0.767090,0.461207}%
\pgfsetfillcolor{currentfill}%
\pgfsetlinewidth{0.000000pt}%
\definecolor{currentstroke}{rgb}{0.000000,0.000000,0.000000}%
\pgfsetstrokecolor{currentstroke}%
\pgfsetstrokeopacity{0.000000}%
\pgfsetdash{}{0pt}%
\pgfpathmoveto{\pgfqpoint{0.748143in}{0.777058in}}%
\pgfpathlineto{\pgfqpoint{0.754327in}{0.807329in}}%
\pgfpathlineto{\pgfqpoint{0.795773in}{0.799049in}}%
\pgfpathlineto{\pgfqpoint{0.851184in}{0.788621in}}%
\pgfpathlineto{\pgfqpoint{0.912796in}{0.778126in}}%
\pgfpathlineto{\pgfqpoint{0.911810in}{0.772480in}}%
\pgfpathlineto{\pgfqpoint{0.927312in}{0.769778in}}%
\pgfpathlineto{\pgfqpoint{0.989331in}{0.759721in}}%
\pgfpathlineto{\pgfqpoint{0.978016in}{0.688001in}}%
\pgfpathlineto{\pgfqpoint{0.924756in}{0.696630in}}%
\pgfpathlineto{\pgfqpoint{0.864661in}{0.706945in}}%
\pgfpathlineto{\pgfqpoint{0.849482in}{0.716019in}}%
\pgfpathlineto{\pgfqpoint{0.766721in}{0.765622in}}%
\pgfpathlineto{\pgfqpoint{0.748143in}{0.777058in}}%
\pgfpathclose%
\pgfusepath{fill}%
\end{pgfscope}%
\begin{pgfscope}%
\pgfpathrectangle{\pgfqpoint{0.100000in}{0.100000in}}{\pgfqpoint{3.007045in}{1.925000in}}%
\pgfusepath{clip}%
\pgfsetbuttcap%
\pgfsetmiterjoin%
\definecolor{currentfill}{rgb}{0.568397,0.821592,0.556063}%
\pgfsetfillcolor{currentfill}%
\pgfsetlinewidth{0.000000pt}%
\definecolor{currentstroke}{rgb}{0.000000,0.000000,0.000000}%
\pgfsetstrokecolor{currentstroke}%
\pgfsetstrokeopacity{0.000000}%
\pgfsetdash{}{0pt}%
\pgfpathmoveto{\pgfqpoint{2.644431in}{1.224390in}}%
\pgfpathlineto{\pgfqpoint{2.643347in}{1.220746in}}%
\pgfpathlineto{\pgfqpoint{2.637469in}{1.212817in}}%
\pgfpathlineto{\pgfqpoint{2.634232in}{1.214856in}}%
\pgfpathlineto{\pgfqpoint{2.621536in}{1.211734in}}%
\pgfpathlineto{\pgfqpoint{2.616014in}{1.216485in}}%
\pgfpathlineto{\pgfqpoint{2.614276in}{1.211702in}}%
\pgfpathlineto{\pgfqpoint{2.609375in}{1.212759in}}%
\pgfpathlineto{\pgfqpoint{2.604647in}{1.216037in}}%
\pgfpathlineto{\pgfqpoint{2.589864in}{1.218745in}}%
\pgfpathlineto{\pgfqpoint{2.585351in}{1.220532in}}%
\pgfpathlineto{\pgfqpoint{2.590743in}{1.227736in}}%
\pgfpathlineto{\pgfqpoint{2.592721in}{1.234363in}}%
\pgfpathlineto{\pgfqpoint{2.598024in}{1.242727in}}%
\pgfpathlineto{\pgfqpoint{2.607688in}{1.243264in}}%
\pgfpathlineto{\pgfqpoint{2.614371in}{1.241096in}}%
\pgfpathlineto{\pgfqpoint{2.644431in}{1.224390in}}%
\pgfpathclose%
\pgfusepath{fill}%
\end{pgfscope}%
\begin{pgfscope}%
\pgfpathrectangle{\pgfqpoint{0.100000in}{0.100000in}}{\pgfqpoint{3.007045in}{1.925000in}}%
\pgfusepath{clip}%
\pgfsetbuttcap%
\pgfsetmiterjoin%
\definecolor{currentfill}{rgb}{0.834833,0.935502,0.811011}%
\pgfsetfillcolor{currentfill}%
\pgfsetlinewidth{0.000000pt}%
\definecolor{currentstroke}{rgb}{0.000000,0.000000,0.000000}%
\pgfsetstrokecolor{currentstroke}%
\pgfsetstrokeopacity{0.000000}%
\pgfsetdash{}{0pt}%
\pgfpathmoveto{\pgfqpoint{1.654308in}{0.604797in}}%
\pgfpathlineto{\pgfqpoint{1.643382in}{0.598629in}}%
\pgfpathlineto{\pgfqpoint{1.639929in}{0.589583in}}%
\pgfpathlineto{\pgfqpoint{1.622469in}{0.597909in}}%
\pgfpathlineto{\pgfqpoint{1.610932in}{0.600948in}}%
\pgfpathlineto{\pgfqpoint{1.606454in}{0.609350in}}%
\pgfpathlineto{\pgfqpoint{1.576374in}{0.610043in}}%
\pgfpathlineto{\pgfqpoint{1.571686in}{0.614167in}}%
\pgfpathlineto{\pgfqpoint{1.569891in}{0.623670in}}%
\pgfpathlineto{\pgfqpoint{1.586901in}{0.635387in}}%
\pgfpathlineto{\pgfqpoint{1.616016in}{0.651438in}}%
\pgfpathlineto{\pgfqpoint{1.620418in}{0.653867in}}%
\pgfpathlineto{\pgfqpoint{1.635041in}{0.627484in}}%
\pgfpathlineto{\pgfqpoint{1.639177in}{0.622324in}}%
\pgfpathlineto{\pgfqpoint{1.642936in}{0.624588in}}%
\pgfpathlineto{\pgfqpoint{1.654308in}{0.604797in}}%
\pgfpathclose%
\pgfusepath{fill}%
\end{pgfscope}%
\begin{pgfscope}%
\pgfpathrectangle{\pgfqpoint{0.100000in}{0.100000in}}{\pgfqpoint{3.007045in}{1.925000in}}%
\pgfusepath{clip}%
\pgfsetbuttcap%
\pgfsetmiterjoin%
\definecolor{currentfill}{rgb}{0.871742,0.950265,0.850381}%
\pgfsetfillcolor{currentfill}%
\pgfsetlinewidth{0.000000pt}%
\definecolor{currentstroke}{rgb}{0.000000,0.000000,0.000000}%
\pgfsetstrokecolor{currentstroke}%
\pgfsetstrokeopacity{0.000000}%
\pgfsetdash{}{0pt}%
\pgfpathmoveto{\pgfqpoint{1.521148in}{0.313814in}}%
\pgfpathlineto{\pgfqpoint{1.530987in}{0.320341in}}%
\pgfpathlineto{\pgfqpoint{1.534731in}{0.327636in}}%
\pgfpathlineto{\pgfqpoint{1.566652in}{0.326194in}}%
\pgfpathlineto{\pgfqpoint{1.592904in}{0.325092in}}%
\pgfpathlineto{\pgfqpoint{1.592550in}{0.313976in}}%
\pgfpathlineto{\pgfqpoint{1.599953in}{0.312537in}}%
\pgfpathlineto{\pgfqpoint{1.625397in}{0.311924in}}%
\pgfpathlineto{\pgfqpoint{1.627117in}{0.305150in}}%
\pgfpathlineto{\pgfqpoint{1.625131in}{0.302682in}}%
\pgfpathlineto{\pgfqpoint{1.630801in}{0.295513in}}%
\pgfpathlineto{\pgfqpoint{1.630560in}{0.289662in}}%
\pgfpathlineto{\pgfqpoint{1.633409in}{0.284357in}}%
\pgfpathlineto{\pgfqpoint{1.633735in}{0.277987in}}%
\pgfpathlineto{\pgfqpoint{1.641756in}{0.274374in}}%
\pgfpathlineto{\pgfqpoint{1.642081in}{0.269076in}}%
\pgfpathlineto{\pgfqpoint{1.628506in}{0.265666in}}%
\pgfpathlineto{\pgfqpoint{1.624844in}{0.260718in}}%
\pgfpathlineto{\pgfqpoint{1.619737in}{0.263995in}}%
\pgfpathlineto{\pgfqpoint{1.609852in}{0.273649in}}%
\pgfpathlineto{\pgfqpoint{1.597319in}{0.276746in}}%
\pgfpathlineto{\pgfqpoint{1.586571in}{0.275707in}}%
\pgfpathlineto{\pgfqpoint{1.575367in}{0.277979in}}%
\pgfpathlineto{\pgfqpoint{1.562903in}{0.288739in}}%
\pgfpathlineto{\pgfqpoint{1.550740in}{0.289846in}}%
\pgfpathlineto{\pgfqpoint{1.542376in}{0.299233in}}%
\pgfpathlineto{\pgfqpoint{1.540217in}{0.298637in}}%
\pgfpathlineto{\pgfqpoint{1.524231in}{0.303862in}}%
\pgfpathlineto{\pgfqpoint{1.525543in}{0.307224in}}%
\pgfpathlineto{\pgfqpoint{1.521148in}{0.313814in}}%
\pgfpathclose%
\pgfusepath{fill}%
\end{pgfscope}%
\begin{pgfscope}%
\pgfpathrectangle{\pgfqpoint{0.100000in}{0.100000in}}{\pgfqpoint{3.007045in}{1.925000in}}%
\pgfusepath{clip}%
\pgfsetbuttcap%
\pgfsetmiterjoin%
\definecolor{currentfill}{rgb}{0.529642,0.803506,0.524198}%
\pgfsetfillcolor{currentfill}%
\pgfsetlinewidth{0.000000pt}%
\definecolor{currentstroke}{rgb}{0.000000,0.000000,0.000000}%
\pgfsetstrokecolor{currentstroke}%
\pgfsetstrokeopacity{0.000000}%
\pgfsetdash{}{0pt}%
\pgfpathmoveto{\pgfqpoint{2.096418in}{1.118342in}}%
\pgfpathlineto{\pgfqpoint{2.095572in}{1.132813in}}%
\pgfpathlineto{\pgfqpoint{2.072910in}{1.131159in}}%
\pgfpathlineto{\pgfqpoint{2.073272in}{1.125381in}}%
\pgfpathlineto{\pgfqpoint{2.067309in}{1.125374in}}%
\pgfpathlineto{\pgfqpoint{2.065520in}{1.156638in}}%
\pgfpathlineto{\pgfqpoint{2.071143in}{1.157025in}}%
\pgfpathlineto{\pgfqpoint{2.105221in}{1.159153in}}%
\pgfpathlineto{\pgfqpoint{2.105428in}{1.156288in}}%
\pgfpathlineto{\pgfqpoint{2.142303in}{1.159619in}}%
\pgfpathlineto{\pgfqpoint{2.148381in}{1.158667in}}%
\pgfpathlineto{\pgfqpoint{2.151777in}{1.152326in}}%
\pgfpathlineto{\pgfqpoint{2.152006in}{1.147645in}}%
\pgfpathlineto{\pgfqpoint{2.155147in}{1.142366in}}%
\pgfpathlineto{\pgfqpoint{2.154997in}{1.139057in}}%
\pgfpathlineto{\pgfqpoint{2.135746in}{1.137366in}}%
\pgfpathlineto{\pgfqpoint{2.117823in}{1.135815in}}%
\pgfpathlineto{\pgfqpoint{2.118010in}{1.123238in}}%
\pgfpathlineto{\pgfqpoint{2.119277in}{1.119542in}}%
\pgfpathlineto{\pgfqpoint{2.096418in}{1.118342in}}%
\pgfpathclose%
\pgfusepath{fill}%
\end{pgfscope}%
\begin{pgfscope}%
\pgfpathrectangle{\pgfqpoint{0.100000in}{0.100000in}}{\pgfqpoint{3.007045in}{1.925000in}}%
\pgfusepath{clip}%
\pgfsetbuttcap%
\pgfsetmiterjoin%
\definecolor{currentfill}{rgb}{0.607151,0.839677,0.587928}%
\pgfsetfillcolor{currentfill}%
\pgfsetlinewidth{0.000000pt}%
\definecolor{currentstroke}{rgb}{0.000000,0.000000,0.000000}%
\pgfsetstrokecolor{currentstroke}%
\pgfsetstrokeopacity{0.000000}%
\pgfsetdash{}{0pt}%
\pgfpathmoveto{\pgfqpoint{2.547659in}{0.980532in}}%
\pgfpathlineto{\pgfqpoint{2.537459in}{0.985766in}}%
\pgfpathlineto{\pgfqpoint{2.533363in}{0.985590in}}%
\pgfpathlineto{\pgfqpoint{2.532109in}{0.998638in}}%
\pgfpathlineto{\pgfqpoint{2.522081in}{0.997367in}}%
\pgfpathlineto{\pgfqpoint{2.520017in}{1.015115in}}%
\pgfpathlineto{\pgfqpoint{2.514084in}{1.019211in}}%
\pgfpathlineto{\pgfqpoint{2.513186in}{1.021977in}}%
\pgfpathlineto{\pgfqpoint{2.515896in}{1.030173in}}%
\pgfpathlineto{\pgfqpoint{2.531083in}{1.032312in}}%
\pgfpathlineto{\pgfqpoint{2.547798in}{1.034084in}}%
\pgfpathlineto{\pgfqpoint{2.561819in}{1.036295in}}%
\pgfpathlineto{\pgfqpoint{2.566736in}{1.001545in}}%
\pgfpathlineto{\pgfqpoint{2.557574in}{1.001121in}}%
\pgfpathlineto{\pgfqpoint{2.549576in}{0.995557in}}%
\pgfpathlineto{\pgfqpoint{2.549802in}{0.987470in}}%
\pgfpathlineto{\pgfqpoint{2.545915in}{0.985956in}}%
\pgfpathlineto{\pgfqpoint{2.547659in}{0.980532in}}%
\pgfpathclose%
\pgfusepath{fill}%
\end{pgfscope}%
\begin{pgfscope}%
\pgfpathrectangle{\pgfqpoint{0.100000in}{0.100000in}}{\pgfqpoint{3.007045in}{1.925000in}}%
\pgfusepath{clip}%
\pgfsetbuttcap%
\pgfsetmiterjoin%
\definecolor{currentfill}{rgb}{0.485352,0.782837,0.487782}%
\pgfsetfillcolor{currentfill}%
\pgfsetlinewidth{0.000000pt}%
\definecolor{currentstroke}{rgb}{0.000000,0.000000,0.000000}%
\pgfsetstrokecolor{currentstroke}%
\pgfsetstrokeopacity{0.000000}%
\pgfsetdash{}{0pt}%
\pgfpathmoveto{\pgfqpoint{2.815570in}{1.283815in}}%
\pgfpathlineto{\pgfqpoint{2.811323in}{1.284246in}}%
\pgfpathlineto{\pgfqpoint{2.806809in}{1.301783in}}%
\pgfpathlineto{\pgfqpoint{2.791685in}{1.322963in}}%
\pgfpathlineto{\pgfqpoint{2.789776in}{1.325534in}}%
\pgfpathlineto{\pgfqpoint{2.799417in}{1.334517in}}%
\pgfpathlineto{\pgfqpoint{2.804907in}{1.347247in}}%
\pgfpathlineto{\pgfqpoint{2.809088in}{1.348480in}}%
\pgfpathlineto{\pgfqpoint{2.816490in}{1.347154in}}%
\pgfpathlineto{\pgfqpoint{2.819196in}{1.340247in}}%
\pgfpathlineto{\pgfqpoint{2.819265in}{1.328442in}}%
\pgfpathlineto{\pgfqpoint{2.821303in}{1.305330in}}%
\pgfpathlineto{\pgfqpoint{2.820085in}{1.298464in}}%
\pgfpathlineto{\pgfqpoint{2.815570in}{1.283815in}}%
\pgfpathclose%
\pgfusepath{fill}%
\end{pgfscope}%
\begin{pgfscope}%
\pgfpathrectangle{\pgfqpoint{0.100000in}{0.100000in}}{\pgfqpoint{3.007045in}{1.925000in}}%
\pgfusepath{clip}%
\pgfsetbuttcap%
\pgfsetmiterjoin%
\definecolor{currentfill}{rgb}{0.474279,0.777670,0.478677}%
\pgfsetfillcolor{currentfill}%
\pgfsetlinewidth{0.000000pt}%
\definecolor{currentstroke}{rgb}{0.000000,0.000000,0.000000}%
\pgfsetstrokecolor{currentstroke}%
\pgfsetstrokeopacity{0.000000}%
\pgfsetdash{}{0pt}%
\pgfpathmoveto{\pgfqpoint{1.633052in}{1.541386in}}%
\pgfpathlineto{\pgfqpoint{1.632109in}{1.506861in}}%
\pgfpathlineto{\pgfqpoint{1.597846in}{1.508282in}}%
\pgfpathlineto{\pgfqpoint{1.598571in}{1.525674in}}%
\pgfpathlineto{\pgfqpoint{1.557925in}{1.527548in}}%
\pgfpathlineto{\pgfqpoint{1.559045in}{1.550570in}}%
\pgfpathlineto{\pgfqpoint{1.552729in}{1.550914in}}%
\pgfpathlineto{\pgfqpoint{1.554716in}{1.596972in}}%
\pgfpathlineto{\pgfqpoint{1.587550in}{1.595331in}}%
\pgfpathlineto{\pgfqpoint{1.634891in}{1.593241in}}%
\pgfpathlineto{\pgfqpoint{1.633052in}{1.541386in}}%
\pgfpathclose%
\pgfusepath{fill}%
\end{pgfscope}%
\begin{pgfscope}%
\pgfpathrectangle{\pgfqpoint{0.100000in}{0.100000in}}{\pgfqpoint{3.007045in}{1.925000in}}%
\pgfusepath{clip}%
\pgfsetbuttcap%
\pgfsetmiterjoin%
\definecolor{currentfill}{rgb}{0.671696,0.867958,0.647105}%
\pgfsetfillcolor{currentfill}%
\pgfsetlinewidth{0.000000pt}%
\definecolor{currentstroke}{rgb}{0.000000,0.000000,0.000000}%
\pgfsetstrokecolor{currentstroke}%
\pgfsetstrokeopacity{0.000000}%
\pgfsetdash{}{0pt}%
\pgfpathmoveto{\pgfqpoint{1.827733in}{0.669764in}}%
\pgfpathlineto{\pgfqpoint{1.799830in}{0.669148in}}%
\pgfpathlineto{\pgfqpoint{1.803211in}{0.660558in}}%
\pgfpathlineto{\pgfqpoint{1.775875in}{0.660550in}}%
\pgfpathlineto{\pgfqpoint{1.773171in}{0.660592in}}%
\pgfpathlineto{\pgfqpoint{1.773020in}{0.706433in}}%
\pgfpathlineto{\pgfqpoint{1.763758in}{0.708620in}}%
\pgfpathlineto{\pgfqpoint{1.763868in}{0.730641in}}%
\pgfpathlineto{\pgfqpoint{1.787896in}{0.730793in}}%
\pgfpathlineto{\pgfqpoint{1.788689in}{0.729118in}}%
\pgfpathlineto{\pgfqpoint{1.825570in}{0.729490in}}%
\pgfpathlineto{\pgfqpoint{1.826203in}{0.670688in}}%
\pgfpathlineto{\pgfqpoint{1.827733in}{0.669764in}}%
\pgfpathclose%
\pgfusepath{fill}%
\end{pgfscope}%
\begin{pgfscope}%
\pgfpathrectangle{\pgfqpoint{0.100000in}{0.100000in}}{\pgfqpoint{3.007045in}{1.925000in}}%
\pgfusepath{clip}%
\pgfsetbuttcap%
\pgfsetmiterjoin%
\definecolor{currentfill}{rgb}{0.886505,0.956171,0.866128}%
\pgfsetfillcolor{currentfill}%
\pgfsetlinewidth{0.000000pt}%
\definecolor{currentstroke}{rgb}{0.000000,0.000000,0.000000}%
\pgfsetstrokecolor{currentstroke}%
\pgfsetstrokeopacity{0.000000}%
\pgfsetdash{}{0pt}%
\pgfpathmoveto{\pgfqpoint{0.877896in}{0.551221in}}%
\pgfpathlineto{\pgfqpoint{0.873853in}{0.546928in}}%
\pgfpathlineto{\pgfqpoint{0.875594in}{0.545533in}}%
\pgfpathlineto{\pgfqpoint{0.874170in}{0.543737in}}%
\pgfpathlineto{\pgfqpoint{0.876010in}{0.542324in}}%
\pgfpathlineto{\pgfqpoint{0.874603in}{0.540599in}}%
\pgfpathlineto{\pgfqpoint{0.876407in}{0.539190in}}%
\pgfpathlineto{\pgfqpoint{0.874516in}{0.537747in}}%
\pgfpathlineto{\pgfqpoint{0.870380in}{0.532422in}}%
\pgfpathlineto{\pgfqpoint{0.873947in}{0.529693in}}%
\pgfpathlineto{\pgfqpoint{0.872574in}{0.527836in}}%
\pgfpathlineto{\pgfqpoint{0.873845in}{0.526915in}}%
\pgfpathlineto{\pgfqpoint{0.867902in}{0.519376in}}%
\pgfpathlineto{\pgfqpoint{0.865978in}{0.517880in}}%
\pgfpathlineto{\pgfqpoint{0.860718in}{0.522144in}}%
\pgfpathlineto{\pgfqpoint{0.857919in}{0.518605in}}%
\pgfpathlineto{\pgfqpoint{0.852921in}{0.522581in}}%
\pgfpathlineto{\pgfqpoint{0.848387in}{0.517598in}}%
\pgfpathlineto{\pgfqpoint{0.845093in}{0.520174in}}%
\pgfpathlineto{\pgfqpoint{0.846505in}{0.521977in}}%
\pgfpathlineto{\pgfqpoint{0.844770in}{0.523398in}}%
\pgfpathlineto{\pgfqpoint{0.846127in}{0.525116in}}%
\pgfpathlineto{\pgfqpoint{0.842673in}{0.527910in}}%
\pgfpathlineto{\pgfqpoint{0.839868in}{0.524417in}}%
\pgfpathlineto{\pgfqpoint{0.836381in}{0.527215in}}%
\pgfpathlineto{\pgfqpoint{0.839219in}{0.530767in}}%
\pgfpathlineto{\pgfqpoint{0.832267in}{0.536448in}}%
\pgfpathlineto{\pgfqpoint{0.826641in}{0.529420in}}%
\pgfpathlineto{\pgfqpoint{0.818104in}{0.536633in}}%
\pgfpathlineto{\pgfqpoint{0.815410in}{0.533363in}}%
\pgfpathlineto{\pgfqpoint{0.810396in}{0.537663in}}%
\pgfpathlineto{\pgfqpoint{0.803129in}{0.529197in}}%
\pgfpathlineto{\pgfqpoint{0.801466in}{0.530632in}}%
\pgfpathlineto{\pgfqpoint{0.799988in}{0.528978in}}%
\pgfpathlineto{\pgfqpoint{0.790016in}{0.537796in}}%
\pgfpathlineto{\pgfqpoint{0.775661in}{0.550783in}}%
\pgfpathlineto{\pgfqpoint{0.768863in}{0.557118in}}%
\pgfpathlineto{\pgfqpoint{0.771871in}{0.560318in}}%
\pgfpathlineto{\pgfqpoint{0.770248in}{0.561865in}}%
\pgfpathlineto{\pgfqpoint{0.773326in}{0.565108in}}%
\pgfpathlineto{\pgfqpoint{0.771216in}{0.567098in}}%
\pgfpathlineto{\pgfqpoint{0.777383in}{0.573526in}}%
\pgfpathlineto{\pgfqpoint{0.776856in}{0.574080in}}%
\pgfpathlineto{\pgfqpoint{0.782998in}{0.580535in}}%
\pgfpathlineto{\pgfqpoint{0.784268in}{0.583082in}}%
\pgfpathlineto{\pgfqpoint{0.786955in}{0.580900in}}%
\pgfpathlineto{\pgfqpoint{0.788384in}{0.577765in}}%
\pgfpathlineto{\pgfqpoint{0.786641in}{0.576148in}}%
\pgfpathlineto{\pgfqpoint{0.784316in}{0.575077in}}%
\pgfpathlineto{\pgfqpoint{0.782913in}{0.573171in}}%
\pgfpathlineto{\pgfqpoint{0.780593in}{0.573219in}}%
\pgfpathlineto{\pgfqpoint{0.779806in}{0.568879in}}%
\pgfpathlineto{\pgfqpoint{0.781934in}{0.567331in}}%
\pgfpathlineto{\pgfqpoint{0.785364in}{0.563844in}}%
\pgfpathlineto{\pgfqpoint{0.788254in}{0.562327in}}%
\pgfpathlineto{\pgfqpoint{0.787949in}{0.560680in}}%
\pgfpathlineto{\pgfqpoint{0.789850in}{0.558307in}}%
\pgfpathlineto{\pgfqpoint{0.792509in}{0.557623in}}%
\pgfpathlineto{\pgfqpoint{0.793630in}{0.555844in}}%
\pgfpathlineto{\pgfqpoint{0.792534in}{0.553561in}}%
\pgfpathlineto{\pgfqpoint{0.797356in}{0.555466in}}%
\pgfpathlineto{\pgfqpoint{0.800166in}{0.557163in}}%
\pgfpathlineto{\pgfqpoint{0.801439in}{0.554540in}}%
\pgfpathlineto{\pgfqpoint{0.804446in}{0.553565in}}%
\pgfpathlineto{\pgfqpoint{0.804535in}{0.556258in}}%
\pgfpathlineto{\pgfqpoint{0.802504in}{0.559954in}}%
\pgfpathlineto{\pgfqpoint{0.799473in}{0.561628in}}%
\pgfpathlineto{\pgfqpoint{0.801889in}{0.563583in}}%
\pgfpathlineto{\pgfqpoint{0.801986in}{0.567136in}}%
\pgfpathlineto{\pgfqpoint{0.802817in}{0.569907in}}%
\pgfpathlineto{\pgfqpoint{0.801188in}{0.572131in}}%
\pgfpathlineto{\pgfqpoint{0.802671in}{0.576011in}}%
\pgfpathlineto{\pgfqpoint{0.804683in}{0.576334in}}%
\pgfpathlineto{\pgfqpoint{0.806090in}{0.575510in}}%
\pgfpathlineto{\pgfqpoint{0.805398in}{0.570181in}}%
\pgfpathlineto{\pgfqpoint{0.803112in}{0.567234in}}%
\pgfpathlineto{\pgfqpoint{0.803364in}{0.561395in}}%
\pgfpathlineto{\pgfqpoint{0.807010in}{0.559720in}}%
\pgfpathlineto{\pgfqpoint{0.807723in}{0.558079in}}%
\pgfpathlineto{\pgfqpoint{0.807544in}{0.554762in}}%
\pgfpathlineto{\pgfqpoint{0.812317in}{0.550291in}}%
\pgfpathlineto{\pgfqpoint{0.814525in}{0.552070in}}%
\pgfpathlineto{\pgfqpoint{0.815449in}{0.555639in}}%
\pgfpathlineto{\pgfqpoint{0.812956in}{0.557139in}}%
\pgfpathlineto{\pgfqpoint{0.812165in}{0.560243in}}%
\pgfpathlineto{\pgfqpoint{0.809060in}{0.561120in}}%
\pgfpathlineto{\pgfqpoint{0.807259in}{0.560638in}}%
\pgfpathlineto{\pgfqpoint{0.805657in}{0.562677in}}%
\pgfpathlineto{\pgfqpoint{0.805494in}{0.568104in}}%
\pgfpathlineto{\pgfqpoint{0.807400in}{0.569549in}}%
\pgfpathlineto{\pgfqpoint{0.811973in}{0.571987in}}%
\pgfpathlineto{\pgfqpoint{0.811391in}{0.574218in}}%
\pgfpathlineto{\pgfqpoint{0.806259in}{0.576839in}}%
\pgfpathlineto{\pgfqpoint{0.807373in}{0.578266in}}%
\pgfpathlineto{\pgfqpoint{0.803872in}{0.581047in}}%
\pgfpathlineto{\pgfqpoint{0.801672in}{0.583533in}}%
\pgfpathlineto{\pgfqpoint{0.799005in}{0.588521in}}%
\pgfpathlineto{\pgfqpoint{0.802634in}{0.593263in}}%
\pgfpathlineto{\pgfqpoint{0.804432in}{0.598840in}}%
\pgfpathlineto{\pgfqpoint{0.804689in}{0.601267in}}%
\pgfpathlineto{\pgfqpoint{0.803744in}{0.605764in}}%
\pgfpathlineto{\pgfqpoint{0.803822in}{0.608129in}}%
\pgfpathlineto{\pgfqpoint{0.803267in}{0.616806in}}%
\pgfpathlineto{\pgfqpoint{0.809438in}{0.610810in}}%
\pgfpathlineto{\pgfqpoint{0.813031in}{0.614554in}}%
\pgfpathlineto{\pgfqpoint{0.825701in}{0.602628in}}%
\pgfpathlineto{\pgfqpoint{0.826925in}{0.603954in}}%
\pgfpathlineto{\pgfqpoint{0.832109in}{0.599208in}}%
\pgfpathlineto{\pgfqpoint{0.830895in}{0.597873in}}%
\pgfpathlineto{\pgfqpoint{0.834401in}{0.594707in}}%
\pgfpathlineto{\pgfqpoint{0.835607in}{0.596048in}}%
\pgfpathlineto{\pgfqpoint{0.838673in}{0.593306in}}%
\pgfpathlineto{\pgfqpoint{0.837473in}{0.591959in}}%
\pgfpathlineto{\pgfqpoint{0.844619in}{0.585660in}}%
\pgfpathlineto{\pgfqpoint{0.845795in}{0.587009in}}%
\pgfpathlineto{\pgfqpoint{0.848948in}{0.584275in}}%
\pgfpathlineto{\pgfqpoint{0.847780in}{0.582922in}}%
\pgfpathlineto{\pgfqpoint{0.859713in}{0.572782in}}%
\pgfpathlineto{\pgfqpoint{0.867050in}{0.566744in}}%
\pgfpathlineto{\pgfqpoint{0.865552in}{0.564907in}}%
\pgfpathlineto{\pgfqpoint{0.872509in}{0.559210in}}%
\pgfpathlineto{\pgfqpoint{0.871135in}{0.557489in}}%
\pgfpathlineto{\pgfqpoint{0.878340in}{0.551804in}}%
\pgfpathlineto{\pgfqpoint{0.877896in}{0.551221in}}%
\pgfpathclose%
\pgfusepath{fill}%
\end{pgfscope}%
\begin{pgfscope}%
\pgfpathrectangle{\pgfqpoint{0.100000in}{0.100000in}}{\pgfqpoint{3.007045in}{1.925000in}}%
\pgfusepath{clip}%
\pgfsetbuttcap%
\pgfsetmiterjoin%
\definecolor{currentfill}{rgb}{0.727797,0.891580,0.701730}%
\pgfsetfillcolor{currentfill}%
\pgfsetlinewidth{0.000000pt}%
\definecolor{currentstroke}{rgb}{0.000000,0.000000,0.000000}%
\pgfsetstrokecolor{currentstroke}%
\pgfsetstrokeopacity{0.000000}%
\pgfsetdash{}{0pt}%
\pgfpathmoveto{\pgfqpoint{1.517880in}{1.481436in}}%
\pgfpathlineto{\pgfqpoint{1.516674in}{1.460211in}}%
\pgfpathlineto{\pgfqpoint{1.517464in}{1.451554in}}%
\pgfpathlineto{\pgfqpoint{1.511045in}{1.452183in}}%
\pgfpathlineto{\pgfqpoint{1.506479in}{1.454584in}}%
\pgfpathlineto{\pgfqpoint{1.496174in}{1.453745in}}%
\pgfpathlineto{\pgfqpoint{1.490267in}{1.460196in}}%
\pgfpathlineto{\pgfqpoint{1.483456in}{1.461953in}}%
\pgfpathlineto{\pgfqpoint{1.485701in}{1.483457in}}%
\pgfpathlineto{\pgfqpoint{1.517880in}{1.481436in}}%
\pgfpathclose%
\pgfusepath{fill}%
\end{pgfscope}%
\begin{pgfscope}%
\pgfpathrectangle{\pgfqpoint{0.100000in}{0.100000in}}{\pgfqpoint{3.007045in}{1.925000in}}%
\pgfusepath{clip}%
\pgfsetbuttcap%
\pgfsetmiterjoin%
\definecolor{currentfill}{rgb}{0.524106,0.800923,0.519646}%
\pgfsetfillcolor{currentfill}%
\pgfsetlinewidth{0.000000pt}%
\definecolor{currentstroke}{rgb}{0.000000,0.000000,0.000000}%
\pgfsetstrokecolor{currentstroke}%
\pgfsetstrokeopacity{0.000000}%
\pgfsetdash{}{0pt}%
\pgfpathmoveto{\pgfqpoint{2.006146in}{1.322530in}}%
\pgfpathlineto{\pgfqpoint{2.000310in}{1.318595in}}%
\pgfpathlineto{\pgfqpoint{1.996691in}{1.320570in}}%
\pgfpathlineto{\pgfqpoint{1.973965in}{1.320197in}}%
\pgfpathlineto{\pgfqpoint{1.973436in}{1.331705in}}%
\pgfpathlineto{\pgfqpoint{1.972422in}{1.354713in}}%
\pgfpathlineto{\pgfqpoint{1.983756in}{1.355205in}}%
\pgfpathlineto{\pgfqpoint{1.983506in}{1.360972in}}%
\pgfpathlineto{\pgfqpoint{1.992817in}{1.361377in}}%
\pgfpathlineto{\pgfqpoint{1.995792in}{1.357804in}}%
\pgfpathlineto{\pgfqpoint{1.997258in}{1.351410in}}%
\pgfpathlineto{\pgfqpoint{2.001078in}{1.349385in}}%
\pgfpathlineto{\pgfqpoint{2.008917in}{1.344598in}}%
\pgfpathlineto{\pgfqpoint{2.010393in}{1.337303in}}%
\pgfpathlineto{\pgfqpoint{2.009028in}{1.324274in}}%
\pgfpathlineto{\pgfqpoint{2.006146in}{1.322530in}}%
\pgfpathclose%
\pgfusepath{fill}%
\end{pgfscope}%
\begin{pgfscope}%
\pgfpathrectangle{\pgfqpoint{0.100000in}{0.100000in}}{\pgfqpoint{3.007045in}{1.925000in}}%
\pgfusepath{clip}%
\pgfsetbuttcap%
\pgfsetmiterjoin%
\definecolor{currentfill}{rgb}{0.451765,0.767090,0.461207}%
\pgfsetfillcolor{currentfill}%
\pgfsetlinewidth{0.000000pt}%
\definecolor{currentstroke}{rgb}{0.000000,0.000000,0.000000}%
\pgfsetstrokecolor{currentstroke}%
\pgfsetstrokeopacity{0.000000}%
\pgfsetdash{}{0pt}%
\pgfpathmoveto{\pgfqpoint{1.561531in}{1.203886in}}%
\pgfpathlineto{\pgfqpoint{1.584352in}{1.202908in}}%
\pgfpathlineto{\pgfqpoint{1.595536in}{1.202468in}}%
\pgfpathlineto{\pgfqpoint{1.594392in}{1.173857in}}%
\pgfpathlineto{\pgfqpoint{1.567058in}{1.175010in}}%
\pgfpathlineto{\pgfqpoint{1.537445in}{1.176396in}}%
\pgfpathlineto{\pgfqpoint{1.539001in}{1.204970in}}%
\pgfpathlineto{\pgfqpoint{1.561531in}{1.203886in}}%
\pgfpathclose%
\pgfusepath{fill}%
\end{pgfscope}%
\begin{pgfscope}%
\pgfpathrectangle{\pgfqpoint{0.100000in}{0.100000in}}{\pgfqpoint{3.007045in}{1.925000in}}%
\pgfusepath{clip}%
\pgfsetbuttcap%
\pgfsetmiterjoin%
\definecolor{currentfill}{rgb}{0.468743,0.775087,0.474125}%
\pgfsetfillcolor{currentfill}%
\pgfsetlinewidth{0.000000pt}%
\definecolor{currentstroke}{rgb}{0.000000,0.000000,0.000000}%
\pgfsetstrokecolor{currentstroke}%
\pgfsetstrokeopacity{0.000000}%
\pgfsetdash{}{0pt}%
\pgfpathmoveto{\pgfqpoint{1.995501in}{1.164493in}}%
\pgfpathlineto{\pgfqpoint{1.997146in}{1.174029in}}%
\pgfpathlineto{\pgfqpoint{1.992899in}{1.184317in}}%
\pgfpathlineto{\pgfqpoint{1.996210in}{1.193778in}}%
\pgfpathlineto{\pgfqpoint{1.978770in}{1.193340in}}%
\pgfpathlineto{\pgfqpoint{1.977971in}{1.216245in}}%
\pgfpathlineto{\pgfqpoint{1.977807in}{1.222266in}}%
\pgfpathlineto{\pgfqpoint{2.000904in}{1.222841in}}%
\pgfpathlineto{\pgfqpoint{2.001133in}{1.217077in}}%
\pgfpathlineto{\pgfqpoint{2.013813in}{1.217404in}}%
\pgfpathlineto{\pgfqpoint{2.008276in}{1.212623in}}%
\pgfpathlineto{\pgfqpoint{2.009474in}{1.209217in}}%
\pgfpathlineto{\pgfqpoint{2.014487in}{1.208775in}}%
\pgfpathlineto{\pgfqpoint{2.024381in}{1.212989in}}%
\pgfpathlineto{\pgfqpoint{2.025601in}{1.191743in}}%
\pgfpathlineto{\pgfqpoint{2.029984in}{1.174625in}}%
\pgfpathlineto{\pgfqpoint{2.018445in}{1.173868in}}%
\pgfpathlineto{\pgfqpoint{2.019642in}{1.156899in}}%
\pgfpathlineto{\pgfqpoint{2.011513in}{1.150677in}}%
\pgfpathlineto{\pgfqpoint{2.002639in}{1.150292in}}%
\pgfpathlineto{\pgfqpoint{1.997266in}{1.151743in}}%
\pgfpathlineto{\pgfqpoint{1.995501in}{1.164493in}}%
\pgfpathclose%
\pgfusepath{fill}%
\end{pgfscope}%
\begin{pgfscope}%
\pgfpathrectangle{\pgfqpoint{0.100000in}{0.100000in}}{\pgfqpoint{3.007045in}{1.925000in}}%
\pgfusepath{clip}%
\pgfsetbuttcap%
\pgfsetmiterjoin%
\definecolor{currentfill}{rgb}{0.671696,0.867958,0.647105}%
\pgfsetfillcolor{currentfill}%
\pgfsetlinewidth{0.000000pt}%
\definecolor{currentstroke}{rgb}{0.000000,0.000000,0.000000}%
\pgfsetstrokecolor{currentstroke}%
\pgfsetstrokeopacity{0.000000}%
\pgfsetdash{}{0pt}%
\pgfpathmoveto{\pgfqpoint{2.664388in}{1.534234in}}%
\pgfpathlineto{\pgfqpoint{2.661128in}{1.541996in}}%
\pgfpathlineto{\pgfqpoint{2.664794in}{1.548806in}}%
\pgfpathlineto{\pgfqpoint{2.662584in}{1.552711in}}%
\pgfpathlineto{\pgfqpoint{2.655181in}{1.554501in}}%
\pgfpathlineto{\pgfqpoint{2.651982in}{1.562017in}}%
\pgfpathlineto{\pgfqpoint{2.659262in}{1.569668in}}%
\pgfpathlineto{\pgfqpoint{2.658006in}{1.573468in}}%
\pgfpathlineto{\pgfqpoint{2.668332in}{1.581558in}}%
\pgfpathlineto{\pgfqpoint{2.671705in}{1.586619in}}%
\pgfpathlineto{\pgfqpoint{2.673065in}{1.592568in}}%
\pgfpathlineto{\pgfqpoint{2.682395in}{1.607388in}}%
\pgfpathlineto{\pgfqpoint{2.696457in}{1.624732in}}%
\pgfpathlineto{\pgfqpoint{2.702880in}{1.630799in}}%
\pgfpathlineto{\pgfqpoint{2.709073in}{1.634394in}}%
\pgfpathlineto{\pgfqpoint{2.713963in}{1.634075in}}%
\pgfpathlineto{\pgfqpoint{2.718437in}{1.632224in}}%
\pgfpathlineto{\pgfqpoint{2.735305in}{1.582994in}}%
\pgfpathlineto{\pgfqpoint{2.735921in}{1.578293in}}%
\pgfpathlineto{\pgfqpoint{2.721573in}{1.573055in}}%
\pgfpathlineto{\pgfqpoint{2.712187in}{1.569616in}}%
\pgfpathlineto{\pgfqpoint{2.706539in}{1.571514in}}%
\pgfpathlineto{\pgfqpoint{2.716233in}{1.541048in}}%
\pgfpathlineto{\pgfqpoint{2.699199in}{1.524078in}}%
\pgfpathlineto{\pgfqpoint{2.688054in}{1.525089in}}%
\pgfpathlineto{\pgfqpoint{2.684161in}{1.539019in}}%
\pgfpathlineto{\pgfqpoint{2.672204in}{1.537694in}}%
\pgfpathlineto{\pgfqpoint{2.664388in}{1.534234in}}%
\pgfpathclose%
\pgfusepath{fill}%
\end{pgfscope}%
\begin{pgfscope}%
\pgfpathrectangle{\pgfqpoint{0.100000in}{0.100000in}}{\pgfqpoint{3.007045in}{1.925000in}}%
\pgfusepath{clip}%
\pgfsetbuttcap%
\pgfsetmiterjoin%
\definecolor{currentfill}{rgb}{0.345098,0.714802,0.408920}%
\pgfsetfillcolor{currentfill}%
\pgfsetlinewidth{0.000000pt}%
\definecolor{currentstroke}{rgb}{0.000000,0.000000,0.000000}%
\pgfsetstrokecolor{currentstroke}%
\pgfsetstrokeopacity{0.000000}%
\pgfsetdash{}{0pt}%
\pgfpathmoveto{\pgfqpoint{1.907064in}{1.059811in}}%
\pgfpathlineto{\pgfqpoint{1.900128in}{1.055674in}}%
\pgfpathlineto{\pgfqpoint{1.896036in}{1.061694in}}%
\pgfpathlineto{\pgfqpoint{1.872518in}{1.061815in}}%
\pgfpathlineto{\pgfqpoint{1.872670in}{1.072352in}}%
\pgfpathlineto{\pgfqpoint{1.871618in}{1.080502in}}%
\pgfpathlineto{\pgfqpoint{1.872048in}{1.113887in}}%
\pgfpathlineto{\pgfqpoint{1.883385in}{1.113543in}}%
\pgfpathlineto{\pgfqpoint{1.883415in}{1.107800in}}%
\pgfpathlineto{\pgfqpoint{1.894852in}{1.097146in}}%
\pgfpathlineto{\pgfqpoint{1.906297in}{1.097100in}}%
\pgfpathlineto{\pgfqpoint{1.906201in}{1.091563in}}%
\pgfpathlineto{\pgfqpoint{1.911673in}{1.091410in}}%
\pgfpathlineto{\pgfqpoint{1.917287in}{1.088655in}}%
\pgfpathlineto{\pgfqpoint{1.918382in}{1.070506in}}%
\pgfpathlineto{\pgfqpoint{1.906913in}{1.070484in}}%
\pgfpathlineto{\pgfqpoint{1.907064in}{1.059811in}}%
\pgfpathclose%
\pgfusepath{fill}%
\end{pgfscope}%
\begin{pgfscope}%
\pgfpathrectangle{\pgfqpoint{0.100000in}{0.100000in}}{\pgfqpoint{3.007045in}{1.925000in}}%
\pgfusepath{clip}%
\pgfsetbuttcap%
\pgfsetmiterjoin%
\definecolor{currentfill}{rgb}{0.638970,0.854179,0.615240}%
\pgfsetfillcolor{currentfill}%
\pgfsetlinewidth{0.000000pt}%
\definecolor{currentstroke}{rgb}{0.000000,0.000000,0.000000}%
\pgfsetstrokecolor{currentstroke}%
\pgfsetstrokeopacity{0.000000}%
\pgfsetdash{}{0pt}%
\pgfpathmoveto{\pgfqpoint{2.045514in}{0.912506in}}%
\pgfpathlineto{\pgfqpoint{2.038540in}{0.912838in}}%
\pgfpathlineto{\pgfqpoint{2.039932in}{0.907365in}}%
\pgfpathlineto{\pgfqpoint{2.033514in}{0.905154in}}%
\pgfpathlineto{\pgfqpoint{2.025623in}{0.904867in}}%
\pgfpathlineto{\pgfqpoint{2.023624in}{0.941619in}}%
\pgfpathlineto{\pgfqpoint{2.021688in}{0.947698in}}%
\pgfpathlineto{\pgfqpoint{2.026615in}{0.954215in}}%
\pgfpathlineto{\pgfqpoint{2.033532in}{0.960342in}}%
\pgfpathlineto{\pgfqpoint{2.034251in}{0.967797in}}%
\pgfpathlineto{\pgfqpoint{2.025516in}{0.974832in}}%
\pgfpathlineto{\pgfqpoint{2.028875in}{0.983778in}}%
\pgfpathlineto{\pgfqpoint{2.038888in}{0.984228in}}%
\pgfpathlineto{\pgfqpoint{2.039044in}{0.968405in}}%
\pgfpathlineto{\pgfqpoint{2.050739in}{0.968865in}}%
\pgfpathlineto{\pgfqpoint{2.059156in}{0.970494in}}%
\pgfpathlineto{\pgfqpoint{2.061897in}{0.966247in}}%
\pgfpathlineto{\pgfqpoint{2.062788in}{0.960749in}}%
\pgfpathlineto{\pgfqpoint{2.054510in}{0.958420in}}%
\pgfpathlineto{\pgfqpoint{2.060132in}{0.952292in}}%
\pgfpathlineto{\pgfqpoint{2.055729in}{0.949079in}}%
\pgfpathlineto{\pgfqpoint{2.054350in}{0.941105in}}%
\pgfpathlineto{\pgfqpoint{2.051897in}{0.936223in}}%
\pgfpathlineto{\pgfqpoint{2.052220in}{0.930837in}}%
\pgfpathlineto{\pgfqpoint{2.049440in}{0.927189in}}%
\pgfpathlineto{\pgfqpoint{2.042608in}{0.925600in}}%
\pgfpathlineto{\pgfqpoint{2.043371in}{0.916596in}}%
\pgfpathlineto{\pgfqpoint{2.045514in}{0.912506in}}%
\pgfpathclose%
\pgfusepath{fill}%
\end{pgfscope}%
\begin{pgfscope}%
\pgfpathrectangle{\pgfqpoint{0.100000in}{0.100000in}}{\pgfqpoint{3.007045in}{1.925000in}}%
\pgfusepath{clip}%
\pgfsetbuttcap%
\pgfsetmiterjoin%
\definecolor{currentfill}{rgb}{0.801615,0.922215,0.775579}%
\pgfsetfillcolor{currentfill}%
\pgfsetlinewidth{0.000000pt}%
\definecolor{currentstroke}{rgb}{0.000000,0.000000,0.000000}%
\pgfsetstrokecolor{currentstroke}%
\pgfsetstrokeopacity{0.000000}%
\pgfsetdash{}{0pt}%
\pgfpathmoveto{\pgfqpoint{2.512057in}{0.734800in}}%
\pgfpathlineto{\pgfqpoint{2.509344in}{0.740112in}}%
\pgfpathlineto{\pgfqpoint{2.485433in}{0.742115in}}%
\pgfpathlineto{\pgfqpoint{2.478955in}{0.739426in}}%
\pgfpathlineto{\pgfqpoint{2.475118in}{0.743801in}}%
\pgfpathlineto{\pgfqpoint{2.476049in}{0.748944in}}%
\pgfpathlineto{\pgfqpoint{2.469526in}{0.751792in}}%
\pgfpathlineto{\pgfqpoint{2.460691in}{0.751658in}}%
\pgfpathlineto{\pgfqpoint{2.463033in}{0.756433in}}%
\pgfpathlineto{\pgfqpoint{2.469397in}{0.761860in}}%
\pgfpathlineto{\pgfqpoint{2.467299in}{0.764175in}}%
\pgfpathlineto{\pgfqpoint{2.470222in}{0.769671in}}%
\pgfpathlineto{\pgfqpoint{2.475814in}{0.775674in}}%
\pgfpathlineto{\pgfqpoint{2.482589in}{0.774200in}}%
\pgfpathlineto{\pgfqpoint{2.485475in}{0.775665in}}%
\pgfpathlineto{\pgfqpoint{2.488225in}{0.782556in}}%
\pgfpathlineto{\pgfqpoint{2.500186in}{0.786765in}}%
\pgfpathlineto{\pgfqpoint{2.505569in}{0.784566in}}%
\pgfpathlineto{\pgfqpoint{2.516660in}{0.795332in}}%
\pgfpathlineto{\pgfqpoint{2.519777in}{0.793317in}}%
\pgfpathlineto{\pgfqpoint{2.519904in}{0.788430in}}%
\pgfpathlineto{\pgfqpoint{2.525012in}{0.782861in}}%
\pgfpathlineto{\pgfqpoint{2.526324in}{0.773753in}}%
\pgfpathlineto{\pgfqpoint{2.529482in}{0.767093in}}%
\pgfpathlineto{\pgfqpoint{2.522083in}{0.760192in}}%
\pgfpathlineto{\pgfqpoint{2.523508in}{0.752566in}}%
\pgfpathlineto{\pgfqpoint{2.528796in}{0.747623in}}%
\pgfpathlineto{\pgfqpoint{2.530406in}{0.743490in}}%
\pgfpathlineto{\pgfqpoint{2.512057in}{0.734800in}}%
\pgfpathclose%
\pgfusepath{fill}%
\end{pgfscope}%
\begin{pgfscope}%
\pgfpathrectangle{\pgfqpoint{0.100000in}{0.100000in}}{\pgfqpoint{3.007045in}{1.925000in}}%
\pgfusepath{clip}%
\pgfsetbuttcap%
\pgfsetmiterjoin%
\definecolor{currentfill}{rgb}{0.823760,0.931073,0.799200}%
\pgfsetfillcolor{currentfill}%
\pgfsetlinewidth{0.000000pt}%
\definecolor{currentstroke}{rgb}{0.000000,0.000000,0.000000}%
\pgfsetstrokecolor{currentstroke}%
\pgfsetstrokeopacity{0.000000}%
\pgfsetdash{}{0pt}%
\pgfpathmoveto{\pgfqpoint{1.215312in}{0.701261in}}%
\pgfpathlineto{\pgfqpoint{1.212900in}{0.676457in}}%
\pgfpathlineto{\pgfqpoint{1.207983in}{0.628535in}}%
\pgfpathlineto{\pgfqpoint{1.205698in}{0.612940in}}%
\pgfpathlineto{\pgfqpoint{1.201869in}{0.611048in}}%
\pgfpathlineto{\pgfqpoint{1.194762in}{0.619908in}}%
\pgfpathlineto{\pgfqpoint{1.189798in}{0.624237in}}%
\pgfpathlineto{\pgfqpoint{1.180001in}{0.628883in}}%
\pgfpathlineto{\pgfqpoint{1.179940in}{0.630968in}}%
\pgfpathlineto{\pgfqpoint{1.171824in}{0.638589in}}%
\pgfpathlineto{\pgfqpoint{1.169910in}{0.645050in}}%
\pgfpathlineto{\pgfqpoint{1.160977in}{0.651771in}}%
\pgfpathlineto{\pgfqpoint{1.152306in}{0.666060in}}%
\pgfpathlineto{\pgfqpoint{1.145301in}{0.669338in}}%
\pgfpathlineto{\pgfqpoint{1.138853in}{0.674660in}}%
\pgfpathlineto{\pgfqpoint{1.131339in}{0.693210in}}%
\pgfpathlineto{\pgfqpoint{1.123511in}{0.697668in}}%
\pgfpathlineto{\pgfqpoint{1.080529in}{0.703340in}}%
\pgfpathlineto{\pgfqpoint{1.087735in}{0.757358in}}%
\pgfpathlineto{\pgfqpoint{1.089302in}{0.768831in}}%
\pgfpathlineto{\pgfqpoint{1.112021in}{0.765725in}}%
\pgfpathlineto{\pgfqpoint{1.112390in}{0.768586in}}%
\pgfpathlineto{\pgfqpoint{1.144512in}{0.779815in}}%
\pgfpathlineto{\pgfqpoint{1.141737in}{0.773832in}}%
\pgfpathlineto{\pgfqpoint{1.133803in}{0.710900in}}%
\pgfpathlineto{\pgfqpoint{1.186737in}{0.704479in}}%
\pgfpathlineto{\pgfqpoint{1.215312in}{0.701261in}}%
\pgfpathclose%
\pgfusepath{fill}%
\end{pgfscope}%
\begin{pgfscope}%
\pgfpathrectangle{\pgfqpoint{0.100000in}{0.100000in}}{\pgfqpoint{3.007045in}{1.925000in}}%
\pgfusepath{clip}%
\pgfsetbuttcap%
\pgfsetmiterjoin%
\definecolor{currentfill}{rgb}{0.779223,0.913233,0.751803}%
\pgfsetfillcolor{currentfill}%
\pgfsetlinewidth{0.000000pt}%
\definecolor{currentstroke}{rgb}{0.000000,0.000000,0.000000}%
\pgfsetstrokecolor{currentstroke}%
\pgfsetstrokeopacity{0.000000}%
\pgfsetdash{}{0pt}%
\pgfpathmoveto{\pgfqpoint{2.321001in}{0.750113in}}%
\pgfpathlineto{\pgfqpoint{2.317396in}{0.757985in}}%
\pgfpathlineto{\pgfqpoint{2.335233in}{0.759815in}}%
\pgfpathlineto{\pgfqpoint{2.333138in}{0.783244in}}%
\pgfpathlineto{\pgfqpoint{2.352426in}{0.785269in}}%
\pgfpathlineto{\pgfqpoint{2.352032in}{0.771233in}}%
\pgfpathlineto{\pgfqpoint{2.354579in}{0.763090in}}%
\pgfpathlineto{\pgfqpoint{2.361889in}{0.757434in}}%
\pgfpathlineto{\pgfqpoint{2.367977in}{0.755519in}}%
\pgfpathlineto{\pgfqpoint{2.360665in}{0.742287in}}%
\pgfpathlineto{\pgfqpoint{2.350162in}{0.739256in}}%
\pgfpathlineto{\pgfqpoint{2.349352in}{0.736500in}}%
\pgfpathlineto{\pgfqpoint{2.351237in}{0.719223in}}%
\pgfpathlineto{\pgfqpoint{2.336528in}{0.717430in}}%
\pgfpathlineto{\pgfqpoint{2.329526in}{0.710196in}}%
\pgfpathlineto{\pgfqpoint{2.330327in}{0.705545in}}%
\pgfpathlineto{\pgfqpoint{2.323124in}{0.704669in}}%
\pgfpathlineto{\pgfqpoint{2.318447in}{0.709939in}}%
\pgfpathlineto{\pgfqpoint{2.309942in}{0.708924in}}%
\pgfpathlineto{\pgfqpoint{2.308027in}{0.714591in}}%
\pgfpathlineto{\pgfqpoint{2.306778in}{0.726142in}}%
\pgfpathlineto{\pgfqpoint{2.312554in}{0.726822in}}%
\pgfpathlineto{\pgfqpoint{2.311941in}{0.730610in}}%
\pgfpathlineto{\pgfqpoint{2.327177in}{0.732545in}}%
\pgfpathlineto{\pgfqpoint{2.330141in}{0.735377in}}%
\pgfpathlineto{\pgfqpoint{2.323388in}{0.743508in}}%
\pgfpathlineto{\pgfqpoint{2.321001in}{0.750113in}}%
\pgfpathclose%
\pgfusepath{fill}%
\end{pgfscope}%
\begin{pgfscope}%
\pgfpathrectangle{\pgfqpoint{0.100000in}{0.100000in}}{\pgfqpoint{3.007045in}{1.925000in}}%
\pgfusepath{clip}%
\pgfsetbuttcap%
\pgfsetmiterjoin%
\definecolor{currentfill}{rgb}{0.671696,0.867958,0.647105}%
\pgfsetfillcolor{currentfill}%
\pgfsetlinewidth{0.000000pt}%
\definecolor{currentstroke}{rgb}{0.000000,0.000000,0.000000}%
\pgfsetstrokecolor{currentstroke}%
\pgfsetstrokeopacity{0.000000}%
\pgfsetdash{}{0pt}%
\pgfpathmoveto{\pgfqpoint{2.252087in}{0.743598in}}%
\pgfpathlineto{\pgfqpoint{2.250388in}{0.740313in}}%
\pgfpathlineto{\pgfqpoint{2.233682in}{0.738645in}}%
\pgfpathlineto{\pgfqpoint{2.233994in}{0.735748in}}%
\pgfpathlineto{\pgfqpoint{2.222669in}{0.734886in}}%
\pgfpathlineto{\pgfqpoint{2.217078in}{0.734332in}}%
\pgfpathlineto{\pgfqpoint{2.215985in}{0.745794in}}%
\pgfpathlineto{\pgfqpoint{2.223920in}{0.746448in}}%
\pgfpathlineto{\pgfqpoint{2.222615in}{0.760540in}}%
\pgfpathlineto{\pgfqpoint{2.217222in}{0.763282in}}%
\pgfpathlineto{\pgfqpoint{2.213720in}{0.767453in}}%
\pgfpathlineto{\pgfqpoint{2.213183in}{0.772767in}}%
\pgfpathlineto{\pgfqpoint{2.204331in}{0.776956in}}%
\pgfpathlineto{\pgfqpoint{2.197889in}{0.784340in}}%
\pgfpathlineto{\pgfqpoint{2.198442in}{0.789193in}}%
\pgfpathlineto{\pgfqpoint{2.202677in}{0.796254in}}%
\pgfpathlineto{\pgfqpoint{2.207356in}{0.797594in}}%
\pgfpathlineto{\pgfqpoint{2.209929in}{0.801633in}}%
\pgfpathlineto{\pgfqpoint{2.210327in}{0.807475in}}%
\pgfpathlineto{\pgfqpoint{2.213109in}{0.813406in}}%
\pgfpathlineto{\pgfqpoint{2.206348in}{0.816022in}}%
\pgfpathlineto{\pgfqpoint{2.204539in}{0.819210in}}%
\pgfpathlineto{\pgfqpoint{2.202564in}{0.841744in}}%
\pgfpathlineto{\pgfqpoint{2.202469in}{0.842698in}}%
\pgfpathlineto{\pgfqpoint{2.231240in}{0.844683in}}%
\pgfpathlineto{\pgfqpoint{2.236928in}{0.845086in}}%
\pgfpathlineto{\pgfqpoint{2.247672in}{0.832584in}}%
\pgfpathlineto{\pgfqpoint{2.246429in}{0.828970in}}%
\pgfpathlineto{\pgfqpoint{2.249130in}{0.825060in}}%
\pgfpathlineto{\pgfqpoint{2.254073in}{0.825883in}}%
\pgfpathlineto{\pgfqpoint{2.257720in}{0.820289in}}%
\pgfpathlineto{\pgfqpoint{2.262319in}{0.816941in}}%
\pgfpathlineto{\pgfqpoint{2.263640in}{0.812852in}}%
\pgfpathlineto{\pgfqpoint{2.258979in}{0.805773in}}%
\pgfpathlineto{\pgfqpoint{2.255427in}{0.799367in}}%
\pgfpathlineto{\pgfqpoint{2.257385in}{0.796619in}}%
\pgfpathlineto{\pgfqpoint{2.248013in}{0.785562in}}%
\pgfpathlineto{\pgfqpoint{2.249772in}{0.779385in}}%
\pgfpathlineto{\pgfqpoint{2.244576in}{0.775293in}}%
\pgfpathlineto{\pgfqpoint{2.242492in}{0.770671in}}%
\pgfpathlineto{\pgfqpoint{2.243973in}{0.753410in}}%
\pgfpathlineto{\pgfqpoint{2.247289in}{0.747109in}}%
\pgfpathlineto{\pgfqpoint{2.252087in}{0.743598in}}%
\pgfpathclose%
\pgfusepath{fill}%
\end{pgfscope}%
\begin{pgfscope}%
\pgfpathrectangle{\pgfqpoint{0.100000in}{0.100000in}}{\pgfqpoint{3.007045in}{1.925000in}}%
\pgfusepath{clip}%
\pgfsetbuttcap%
\pgfsetmiterjoin%
\definecolor{currentfill}{rgb}{0.629296,0.850012,0.606136}%
\pgfsetfillcolor{currentfill}%
\pgfsetlinewidth{0.000000pt}%
\definecolor{currentstroke}{rgb}{0.000000,0.000000,0.000000}%
\pgfsetstrokecolor{currentstroke}%
\pgfsetstrokeopacity{0.000000}%
\pgfsetdash{}{0pt}%
\pgfpathmoveto{\pgfqpoint{1.710630in}{1.795591in}}%
\pgfpathlineto{\pgfqpoint{1.710357in}{1.776679in}}%
\pgfpathlineto{\pgfqpoint{1.710937in}{1.765107in}}%
\pgfpathlineto{\pgfqpoint{1.676648in}{1.765752in}}%
\pgfpathlineto{\pgfqpoint{1.677958in}{1.760345in}}%
\pgfpathlineto{\pgfqpoint{1.676638in}{1.755236in}}%
\pgfpathlineto{\pgfqpoint{1.678294in}{1.748130in}}%
\pgfpathlineto{\pgfqpoint{1.677065in}{1.742344in}}%
\pgfpathlineto{\pgfqpoint{1.643166in}{1.743370in}}%
\pgfpathlineto{\pgfqpoint{1.625869in}{1.743967in}}%
\pgfpathlineto{\pgfqpoint{1.626307in}{1.755616in}}%
\pgfpathlineto{\pgfqpoint{1.625579in}{1.767286in}}%
\pgfpathlineto{\pgfqpoint{1.596746in}{1.768470in}}%
\pgfpathlineto{\pgfqpoint{1.595955in}{1.780137in}}%
\pgfpathlineto{\pgfqpoint{1.596789in}{1.798977in}}%
\pgfpathlineto{\pgfqpoint{1.638516in}{1.797389in}}%
\pgfpathlineto{\pgfqpoint{1.674481in}{1.796361in}}%
\pgfpathlineto{\pgfqpoint{1.710630in}{1.795591in}}%
\pgfpathclose%
\pgfusepath{fill}%
\end{pgfscope}%
\begin{pgfscope}%
\pgfpathrectangle{\pgfqpoint{0.100000in}{0.100000in}}{\pgfqpoint{3.007045in}{1.925000in}}%
\pgfusepath{clip}%
\pgfsetbuttcap%
\pgfsetmiterjoin%
\definecolor{currentfill}{rgb}{0.930980,0.973595,0.916863}%
\pgfsetfillcolor{currentfill}%
\pgfsetlinewidth{0.000000pt}%
\definecolor{currentstroke}{rgb}{0.000000,0.000000,0.000000}%
\pgfsetstrokecolor{currentstroke}%
\pgfsetstrokeopacity{0.000000}%
\pgfsetdash{}{0pt}%
\pgfpathmoveto{\pgfqpoint{1.521148in}{0.313814in}}%
\pgfpathlineto{\pgfqpoint{1.519248in}{0.324225in}}%
\pgfpathlineto{\pgfqpoint{1.516036in}{0.332359in}}%
\pgfpathlineto{\pgfqpoint{1.509870in}{0.338738in}}%
\pgfpathlineto{\pgfqpoint{1.505834in}{0.347871in}}%
\pgfpathlineto{\pgfqpoint{1.507664in}{0.355223in}}%
\pgfpathlineto{\pgfqpoint{1.503531in}{0.366333in}}%
\pgfpathlineto{\pgfqpoint{1.505470in}{0.376130in}}%
\pgfpathlineto{\pgfqpoint{1.502792in}{0.376996in}}%
\pgfpathlineto{\pgfqpoint{1.502216in}{0.384359in}}%
\pgfpathlineto{\pgfqpoint{1.491739in}{0.388953in}}%
\pgfpathlineto{\pgfqpoint{1.486685in}{0.396463in}}%
\pgfpathlineto{\pgfqpoint{1.482967in}{0.398437in}}%
\pgfpathlineto{\pgfqpoint{1.482372in}{0.405110in}}%
\pgfpathlineto{\pgfqpoint{1.476970in}{0.411788in}}%
\pgfpathlineto{\pgfqpoint{1.472047in}{0.422119in}}%
\pgfpathlineto{\pgfqpoint{1.464650in}{0.426083in}}%
\pgfpathlineto{\pgfqpoint{1.470461in}{0.425819in}}%
\pgfpathlineto{\pgfqpoint{1.512921in}{0.424052in}}%
\pgfpathlineto{\pgfqpoint{1.512666in}{0.412373in}}%
\pgfpathlineto{\pgfqpoint{1.547367in}{0.412624in}}%
\pgfpathlineto{\pgfqpoint{1.545690in}{0.365433in}}%
\pgfpathlineto{\pgfqpoint{1.560356in}{0.365181in}}%
\pgfpathlineto{\pgfqpoint{1.563594in}{0.356706in}}%
\pgfpathlineto{\pgfqpoint{1.564709in}{0.344532in}}%
\pgfpathlineto{\pgfqpoint{1.567623in}{0.344460in}}%
\pgfpathlineto{\pgfqpoint{1.566652in}{0.326194in}}%
\pgfpathlineto{\pgfqpoint{1.534731in}{0.327636in}}%
\pgfpathlineto{\pgfqpoint{1.530987in}{0.320341in}}%
\pgfpathlineto{\pgfqpoint{1.521148in}{0.313814in}}%
\pgfpathclose%
\pgfusepath{fill}%
\end{pgfscope}%
\begin{pgfscope}%
\pgfpathrectangle{\pgfqpoint{0.100000in}{0.100000in}}{\pgfqpoint{3.007045in}{1.925000in}}%
\pgfusepath{clip}%
\pgfsetbuttcap%
\pgfsetmiterjoin%
\definecolor{currentfill}{rgb}{0.535179,0.806090,0.528750}%
\pgfsetfillcolor{currentfill}%
\pgfsetlinewidth{0.000000pt}%
\definecolor{currentstroke}{rgb}{0.000000,0.000000,0.000000}%
\pgfsetstrokecolor{currentstroke}%
\pgfsetstrokeopacity{0.000000}%
\pgfsetdash{}{0pt}%
\pgfpathmoveto{\pgfqpoint{2.503300in}{0.656930in}}%
\pgfpathlineto{\pgfqpoint{2.498304in}{0.661191in}}%
\pgfpathlineto{\pgfqpoint{2.497961in}{0.669319in}}%
\pgfpathlineto{\pgfqpoint{2.503124in}{0.674785in}}%
\pgfpathlineto{\pgfqpoint{2.508359in}{0.677599in}}%
\pgfpathlineto{\pgfqpoint{2.505019in}{0.681519in}}%
\pgfpathlineto{\pgfqpoint{2.511887in}{0.680666in}}%
\pgfpathlineto{\pgfqpoint{2.518569in}{0.682840in}}%
\pgfpathlineto{\pgfqpoint{2.523188in}{0.681053in}}%
\pgfpathlineto{\pgfqpoint{2.528105in}{0.691406in}}%
\pgfpathlineto{\pgfqpoint{2.524353in}{0.695341in}}%
\pgfpathlineto{\pgfqpoint{2.529648in}{0.698799in}}%
\pgfpathlineto{\pgfqpoint{2.532623in}{0.707455in}}%
\pgfpathlineto{\pgfqpoint{2.536542in}{0.704088in}}%
\pgfpathlineto{\pgfqpoint{2.545964in}{0.706161in}}%
\pgfpathlineto{\pgfqpoint{2.550608in}{0.701257in}}%
\pgfpathlineto{\pgfqpoint{2.551853in}{0.698652in}}%
\pgfpathlineto{\pgfqpoint{2.548410in}{0.688938in}}%
\pgfpathlineto{\pgfqpoint{2.550045in}{0.680517in}}%
\pgfpathlineto{\pgfqpoint{2.543234in}{0.669808in}}%
\pgfpathlineto{\pgfqpoint{2.542300in}{0.662555in}}%
\pgfpathlineto{\pgfqpoint{2.533990in}{0.667284in}}%
\pgfpathlineto{\pgfqpoint{2.521337in}{0.669031in}}%
\pgfpathlineto{\pgfqpoint{2.514003in}{0.660961in}}%
\pgfpathlineto{\pgfqpoint{2.506695in}{0.661778in}}%
\pgfpathlineto{\pgfqpoint{2.503300in}{0.656930in}}%
\pgfpathclose%
\pgfusepath{fill}%
\end{pgfscope}%
\begin{pgfscope}%
\pgfpathrectangle{\pgfqpoint{0.100000in}{0.100000in}}{\pgfqpoint{3.007045in}{1.925000in}}%
\pgfusepath{clip}%
\pgfsetbuttcap%
\pgfsetmiterjoin%
\definecolor{currentfill}{rgb}{0.389020,0.736332,0.430450}%
\pgfsetfillcolor{currentfill}%
\pgfsetlinewidth{0.000000pt}%
\definecolor{currentstroke}{rgb}{0.000000,0.000000,0.000000}%
\pgfsetstrokecolor{currentstroke}%
\pgfsetstrokeopacity{0.000000}%
\pgfsetdash{}{0pt}%
\pgfpathmoveto{\pgfqpoint{1.326529in}{1.324213in}}%
\pgfpathlineto{\pgfqpoint{1.333387in}{1.392952in}}%
\pgfpathlineto{\pgfqpoint{1.335923in}{1.418503in}}%
\pgfpathlineto{\pgfqpoint{1.386507in}{1.413753in}}%
\pgfpathlineto{\pgfqpoint{1.396545in}{1.412828in}}%
\pgfpathlineto{\pgfqpoint{1.396236in}{1.398628in}}%
\pgfpathlineto{\pgfqpoint{1.394290in}{1.375922in}}%
\pgfpathlineto{\pgfqpoint{1.395452in}{1.375827in}}%
\pgfpathlineto{\pgfqpoint{1.393848in}{1.352944in}}%
\pgfpathlineto{\pgfqpoint{1.396518in}{1.346955in}}%
\pgfpathlineto{\pgfqpoint{1.395435in}{1.334971in}}%
\pgfpathlineto{\pgfqpoint{1.396817in}{1.327688in}}%
\pgfpathlineto{\pgfqpoint{1.395381in}{1.309582in}}%
\pgfpathlineto{\pgfqpoint{1.359322in}{1.312674in}}%
\pgfpathlineto{\pgfqpoint{1.358996in}{1.309848in}}%
\pgfpathlineto{\pgfqpoint{1.325420in}{1.312996in}}%
\pgfpathlineto{\pgfqpoint{1.326529in}{1.324213in}}%
\pgfpathclose%
\pgfusepath{fill}%
\end{pgfscope}%
\begin{pgfscope}%
\pgfpathrectangle{\pgfqpoint{0.100000in}{0.100000in}}{\pgfqpoint{3.007045in}{1.925000in}}%
\pgfusepath{clip}%
\pgfsetbuttcap%
\pgfsetmiterjoin%
\definecolor{currentfill}{rgb}{0.629296,0.850012,0.606136}%
\pgfsetfillcolor{currentfill}%
\pgfsetlinewidth{0.000000pt}%
\definecolor{currentstroke}{rgb}{0.000000,0.000000,0.000000}%
\pgfsetstrokecolor{currentstroke}%
\pgfsetstrokeopacity{0.000000}%
\pgfsetdash{}{0pt}%
\pgfpathmoveto{\pgfqpoint{2.011147in}{0.717467in}}%
\pgfpathlineto{\pgfqpoint{2.010976in}{0.714602in}}%
\pgfpathlineto{\pgfqpoint{2.003446in}{0.714212in}}%
\pgfpathlineto{\pgfqpoint{2.003803in}{0.705398in}}%
\pgfpathlineto{\pgfqpoint{1.998667in}{0.705092in}}%
\pgfpathlineto{\pgfqpoint{1.997671in}{0.713980in}}%
\pgfpathlineto{\pgfqpoint{1.993233in}{0.713829in}}%
\pgfpathlineto{\pgfqpoint{1.988496in}{0.714303in}}%
\pgfpathlineto{\pgfqpoint{1.986961in}{0.717838in}}%
\pgfpathlineto{\pgfqpoint{1.991838in}{0.723800in}}%
\pgfpathlineto{\pgfqpoint{1.985938in}{0.725835in}}%
\pgfpathlineto{\pgfqpoint{1.986138in}{0.729844in}}%
\pgfpathlineto{\pgfqpoint{1.990829in}{0.734444in}}%
\pgfpathlineto{\pgfqpoint{1.989679in}{0.740203in}}%
\pgfpathlineto{\pgfqpoint{1.985200in}{0.741858in}}%
\pgfpathlineto{\pgfqpoint{1.988976in}{0.750911in}}%
\pgfpathlineto{\pgfqpoint{1.988710in}{0.756370in}}%
\pgfpathlineto{\pgfqpoint{1.985739in}{0.761415in}}%
\pgfpathlineto{\pgfqpoint{1.985442in}{0.766767in}}%
\pgfpathlineto{\pgfqpoint{1.980431in}{0.778978in}}%
\pgfpathlineto{\pgfqpoint{1.985316in}{0.780979in}}%
\pgfpathlineto{\pgfqpoint{1.980623in}{0.785127in}}%
\pgfpathlineto{\pgfqpoint{1.984613in}{0.793228in}}%
\pgfpathlineto{\pgfqpoint{1.992182in}{0.792843in}}%
\pgfpathlineto{\pgfqpoint{1.988015in}{0.799119in}}%
\pgfpathlineto{\pgfqpoint{1.991139in}{0.803150in}}%
\pgfpathlineto{\pgfqpoint{1.986772in}{0.806255in}}%
\pgfpathlineto{\pgfqpoint{1.996950in}{0.809911in}}%
\pgfpathlineto{\pgfqpoint{1.998299in}{0.813452in}}%
\pgfpathlineto{\pgfqpoint{1.993621in}{0.816052in}}%
\pgfpathlineto{\pgfqpoint{1.993574in}{0.816135in}}%
\pgfpathlineto{\pgfqpoint{2.009788in}{0.816800in}}%
\pgfpathlineto{\pgfqpoint{2.010153in}{0.808128in}}%
\pgfpathlineto{\pgfqpoint{2.021640in}{0.808564in}}%
\pgfpathlineto{\pgfqpoint{2.022165in}{0.797050in}}%
\pgfpathlineto{\pgfqpoint{2.023674in}{0.765200in}}%
\pgfpathlineto{\pgfqpoint{2.018021in}{0.764899in}}%
\pgfpathlineto{\pgfqpoint{2.018223in}{0.761023in}}%
\pgfpathlineto{\pgfqpoint{2.009327in}{0.760562in}}%
\pgfpathlineto{\pgfqpoint{2.013096in}{0.749239in}}%
\pgfpathlineto{\pgfqpoint{2.013725in}{0.737732in}}%
\pgfpathlineto{\pgfqpoint{2.008261in}{0.731716in}}%
\pgfpathlineto{\pgfqpoint{2.011096in}{0.729886in}}%
\pgfpathlineto{\pgfqpoint{2.011147in}{0.717467in}}%
\pgfpathclose%
\pgfusepath{fill}%
\end{pgfscope}%
\begin{pgfscope}%
\pgfpathrectangle{\pgfqpoint{0.100000in}{0.100000in}}{\pgfqpoint{3.007045in}{1.925000in}}%
\pgfusepath{clip}%
\pgfsetbuttcap%
\pgfsetmiterjoin%
\definecolor{currentfill}{rgb}{0.326275,0.705575,0.399692}%
\pgfsetfillcolor{currentfill}%
\pgfsetlinewidth{0.000000pt}%
\definecolor{currentstroke}{rgb}{0.000000,0.000000,0.000000}%
\pgfsetstrokecolor{currentstroke}%
\pgfsetstrokeopacity{0.000000}%
\pgfsetdash{}{0pt}%
\pgfpathmoveto{\pgfqpoint{0.416964in}{1.700001in}}%
\pgfpathlineto{\pgfqpoint{0.427465in}{1.723088in}}%
\pgfpathlineto{\pgfqpoint{0.429108in}{1.730577in}}%
\pgfpathlineto{\pgfqpoint{0.441718in}{1.753288in}}%
\pgfpathlineto{\pgfqpoint{0.445758in}{1.765737in}}%
\pgfpathlineto{\pgfqpoint{0.453664in}{1.784658in}}%
\pgfpathlineto{\pgfqpoint{0.454831in}{1.792649in}}%
\pgfpathlineto{\pgfqpoint{0.465616in}{1.788403in}}%
\pgfpathlineto{\pgfqpoint{0.481516in}{1.783693in}}%
\pgfpathlineto{\pgfqpoint{0.480074in}{1.779212in}}%
\pgfpathlineto{\pgfqpoint{0.475846in}{1.775523in}}%
\pgfpathlineto{\pgfqpoint{0.475014in}{1.769793in}}%
\pgfpathlineto{\pgfqpoint{0.469874in}{1.763357in}}%
\pgfpathlineto{\pgfqpoint{0.465504in}{1.749616in}}%
\pgfpathlineto{\pgfqpoint{0.451249in}{1.754244in}}%
\pgfpathlineto{\pgfqpoint{0.448349in}{1.745454in}}%
\pgfpathlineto{\pgfqpoint{0.451037in}{1.744559in}}%
\pgfpathlineto{\pgfqpoint{0.444068in}{1.723357in}}%
\pgfpathlineto{\pgfqpoint{0.449171in}{1.720447in}}%
\pgfpathlineto{\pgfqpoint{0.443505in}{1.702226in}}%
\pgfpathlineto{\pgfqpoint{0.437975in}{1.704042in}}%
\pgfpathlineto{\pgfqpoint{0.436719in}{1.699407in}}%
\pgfpathlineto{\pgfqpoint{0.432427in}{1.695452in}}%
\pgfpathlineto{\pgfqpoint{0.416964in}{1.700001in}}%
\pgfpathclose%
\pgfusepath{fill}%
\end{pgfscope}%
\begin{pgfscope}%
\pgfpathrectangle{\pgfqpoint{0.100000in}{0.100000in}}{\pgfqpoint{3.007045in}{1.925000in}}%
\pgfusepath{clip}%
\pgfsetbuttcap%
\pgfsetmiterjoin%
\definecolor{currentfill}{rgb}{0.886505,0.956171,0.866128}%
\pgfsetfillcolor{currentfill}%
\pgfsetlinewidth{0.000000pt}%
\definecolor{currentstroke}{rgb}{0.000000,0.000000,0.000000}%
\pgfsetstrokecolor{currentstroke}%
\pgfsetstrokeopacity{0.000000}%
\pgfsetdash{}{0pt}%
\pgfpathmoveto{\pgfqpoint{0.377147in}{0.493245in}}%
\pgfpathlineto{\pgfqpoint{0.376520in}{0.497403in}}%
\pgfpathlineto{\pgfqpoint{0.377477in}{0.497571in}}%
\pgfpathlineto{\pgfqpoint{0.377147in}{0.493245in}}%
\pgfpathclose%
\pgfusepath{fill}%
\end{pgfscope}%
\begin{pgfscope}%
\pgfpathrectangle{\pgfqpoint{0.100000in}{0.100000in}}{\pgfqpoint{3.007045in}{1.925000in}}%
\pgfusepath{clip}%
\pgfsetbuttcap%
\pgfsetmiterjoin%
\definecolor{currentfill}{rgb}{0.886505,0.956171,0.866128}%
\pgfsetfillcolor{currentfill}%
\pgfsetlinewidth{0.000000pt}%
\definecolor{currentstroke}{rgb}{0.000000,0.000000,0.000000}%
\pgfsetstrokecolor{currentstroke}%
\pgfsetstrokeopacity{0.000000}%
\pgfsetdash{}{0pt}%
\pgfpathmoveto{\pgfqpoint{0.405385in}{0.465085in}}%
\pgfpathlineto{\pgfqpoint{0.405017in}{0.468510in}}%
\pgfpathlineto{\pgfqpoint{0.406749in}{0.467215in}}%
\pgfpathlineto{\pgfqpoint{0.407796in}{0.467652in}}%
\pgfpathlineto{\pgfqpoint{0.407185in}{0.472036in}}%
\pgfpathlineto{\pgfqpoint{0.408305in}{0.472399in}}%
\pgfpathlineto{\pgfqpoint{0.410270in}{0.469591in}}%
\pgfpathlineto{\pgfqpoint{0.409867in}{0.468623in}}%
\pgfpathlineto{\pgfqpoint{0.410756in}{0.466052in}}%
\pgfpathlineto{\pgfqpoint{0.406364in}{0.465975in}}%
\pgfpathlineto{\pgfqpoint{0.405385in}{0.465085in}}%
\pgfpathclose%
\pgfusepath{fill}%
\end{pgfscope}%
\begin{pgfscope}%
\pgfpathrectangle{\pgfqpoint{0.100000in}{0.100000in}}{\pgfqpoint{3.007045in}{1.925000in}}%
\pgfusepath{clip}%
\pgfsetbuttcap%
\pgfsetmiterjoin%
\definecolor{currentfill}{rgb}{0.886505,0.956171,0.866128}%
\pgfsetfillcolor{currentfill}%
\pgfsetlinewidth{0.000000pt}%
\definecolor{currentstroke}{rgb}{0.000000,0.000000,0.000000}%
\pgfsetstrokecolor{currentstroke}%
\pgfsetstrokeopacity{0.000000}%
\pgfsetdash{}{0pt}%
\pgfpathmoveto{\pgfqpoint{0.415072in}{0.452560in}}%
\pgfpathlineto{\pgfqpoint{0.414645in}{0.455401in}}%
\pgfpathlineto{\pgfqpoint{0.419140in}{0.455415in}}%
\pgfpathlineto{\pgfqpoint{0.420979in}{0.457297in}}%
\pgfpathlineto{\pgfqpoint{0.422131in}{0.456956in}}%
\pgfpathlineto{\pgfqpoint{0.423763in}{0.455028in}}%
\pgfpathlineto{\pgfqpoint{0.422007in}{0.454672in}}%
\pgfpathlineto{\pgfqpoint{0.421045in}{0.453407in}}%
\pgfpathlineto{\pgfqpoint{0.422643in}{0.451021in}}%
\pgfpathlineto{\pgfqpoint{0.423045in}{0.449284in}}%
\pgfpathlineto{\pgfqpoint{0.421244in}{0.449179in}}%
\pgfpathlineto{\pgfqpoint{0.420180in}{0.450490in}}%
\pgfpathlineto{\pgfqpoint{0.417429in}{0.452432in}}%
\pgfpathlineto{\pgfqpoint{0.415072in}{0.452560in}}%
\pgfpathclose%
\pgfusepath{fill}%
\end{pgfscope}%
\begin{pgfscope}%
\pgfpathrectangle{\pgfqpoint{0.100000in}{0.100000in}}{\pgfqpoint{3.007045in}{1.925000in}}%
\pgfusepath{clip}%
\pgfsetbuttcap%
\pgfsetmiterjoin%
\definecolor{currentfill}{rgb}{0.886505,0.956171,0.866128}%
\pgfsetfillcolor{currentfill}%
\pgfsetlinewidth{0.000000pt}%
\definecolor{currentstroke}{rgb}{0.000000,0.000000,0.000000}%
\pgfsetstrokecolor{currentstroke}%
\pgfsetstrokeopacity{0.000000}%
\pgfsetdash{}{0pt}%
\pgfpathmoveto{\pgfqpoint{0.412280in}{0.458704in}}%
\pgfpathlineto{\pgfqpoint{0.412274in}{0.460723in}}%
\pgfpathlineto{\pgfqpoint{0.408425in}{0.462631in}}%
\pgfpathlineto{\pgfqpoint{0.410054in}{0.463833in}}%
\pgfpathlineto{\pgfqpoint{0.412335in}{0.461427in}}%
\pgfpathlineto{\pgfqpoint{0.415525in}{0.460193in}}%
\pgfpathlineto{\pgfqpoint{0.417642in}{0.461549in}}%
\pgfpathlineto{\pgfqpoint{0.418112in}{0.459372in}}%
\pgfpathlineto{\pgfqpoint{0.415856in}{0.459267in}}%
\pgfpathlineto{\pgfqpoint{0.414135in}{0.457815in}}%
\pgfpathlineto{\pgfqpoint{0.412280in}{0.458704in}}%
\pgfpathclose%
\pgfusepath{fill}%
\end{pgfscope}%
\begin{pgfscope}%
\pgfpathrectangle{\pgfqpoint{0.100000in}{0.100000in}}{\pgfqpoint{3.007045in}{1.925000in}}%
\pgfusepath{clip}%
\pgfsetbuttcap%
\pgfsetmiterjoin%
\definecolor{currentfill}{rgb}{0.886505,0.956171,0.866128}%
\pgfsetfillcolor{currentfill}%
\pgfsetlinewidth{0.000000pt}%
\definecolor{currentstroke}{rgb}{0.000000,0.000000,0.000000}%
\pgfsetstrokecolor{currentstroke}%
\pgfsetstrokeopacity{0.000000}%
\pgfsetdash{}{0pt}%
\pgfpathmoveto{\pgfqpoint{0.389271in}{0.496402in}}%
\pgfpathlineto{\pgfqpoint{0.388843in}{0.498609in}}%
\pgfpathlineto{\pgfqpoint{0.391517in}{0.499040in}}%
\pgfpathlineto{\pgfqpoint{0.392195in}{0.498040in}}%
\pgfpathlineto{\pgfqpoint{0.391561in}{0.496110in}}%
\pgfpathlineto{\pgfqpoint{0.389271in}{0.496402in}}%
\pgfpathclose%
\pgfusepath{fill}%
\end{pgfscope}%
\begin{pgfscope}%
\pgfpathrectangle{\pgfqpoint{0.100000in}{0.100000in}}{\pgfqpoint{3.007045in}{1.925000in}}%
\pgfusepath{clip}%
\pgfsetbuttcap%
\pgfsetmiterjoin%
\definecolor{currentfill}{rgb}{0.886505,0.956171,0.866128}%
\pgfsetfillcolor{currentfill}%
\pgfsetlinewidth{0.000000pt}%
\definecolor{currentstroke}{rgb}{0.000000,0.000000,0.000000}%
\pgfsetstrokecolor{currentstroke}%
\pgfsetstrokeopacity{0.000000}%
\pgfsetdash{}{0pt}%
\pgfpathmoveto{\pgfqpoint{0.424629in}{0.447465in}}%
\pgfpathlineto{\pgfqpoint{0.424048in}{0.449531in}}%
\pgfpathlineto{\pgfqpoint{0.427830in}{0.447738in}}%
\pgfpathlineto{\pgfqpoint{0.426724in}{0.446647in}}%
\pgfpathlineto{\pgfqpoint{0.424629in}{0.447465in}}%
\pgfpathclose%
\pgfusepath{fill}%
\end{pgfscope}%
\begin{pgfscope}%
\pgfpathrectangle{\pgfqpoint{0.100000in}{0.100000in}}{\pgfqpoint{3.007045in}{1.925000in}}%
\pgfusepath{clip}%
\pgfsetbuttcap%
\pgfsetmiterjoin%
\definecolor{currentfill}{rgb}{0.886505,0.956171,0.866128}%
\pgfsetfillcolor{currentfill}%
\pgfsetlinewidth{0.000000pt}%
\definecolor{currentstroke}{rgb}{0.000000,0.000000,0.000000}%
\pgfsetstrokecolor{currentstroke}%
\pgfsetstrokeopacity{0.000000}%
\pgfsetdash{}{0pt}%
\pgfpathmoveto{\pgfqpoint{0.380099in}{0.509690in}}%
\pgfpathlineto{\pgfqpoint{0.380270in}{0.511273in}}%
\pgfpathlineto{\pgfqpoint{0.382192in}{0.510622in}}%
\pgfpathlineto{\pgfqpoint{0.381722in}{0.509141in}}%
\pgfpathlineto{\pgfqpoint{0.380099in}{0.509690in}}%
\pgfpathclose%
\pgfusepath{fill}%
\end{pgfscope}%
\begin{pgfscope}%
\pgfpathrectangle{\pgfqpoint{0.100000in}{0.100000in}}{\pgfqpoint{3.007045in}{1.925000in}}%
\pgfusepath{clip}%
\pgfsetbuttcap%
\pgfsetmiterjoin%
\definecolor{currentfill}{rgb}{0.886505,0.956171,0.866128}%
\pgfsetfillcolor{currentfill}%
\pgfsetlinewidth{0.000000pt}%
\definecolor{currentstroke}{rgb}{0.000000,0.000000,0.000000}%
\pgfsetstrokecolor{currentstroke}%
\pgfsetstrokeopacity{0.000000}%
\pgfsetdash{}{0pt}%
\pgfpathmoveto{\pgfqpoint{0.428107in}{0.449417in}}%
\pgfpathlineto{\pgfqpoint{0.427960in}{0.451570in}}%
\pgfpathlineto{\pgfqpoint{0.429777in}{0.452594in}}%
\pgfpathlineto{\pgfqpoint{0.430464in}{0.450596in}}%
\pgfpathlineto{\pgfqpoint{0.428107in}{0.449417in}}%
\pgfpathclose%
\pgfusepath{fill}%
\end{pgfscope}%
\begin{pgfscope}%
\pgfpathrectangle{\pgfqpoint{0.100000in}{0.100000in}}{\pgfqpoint{3.007045in}{1.925000in}}%
\pgfusepath{clip}%
\pgfsetbuttcap%
\pgfsetmiterjoin%
\definecolor{currentfill}{rgb}{0.886505,0.956171,0.866128}%
\pgfsetfillcolor{currentfill}%
\pgfsetlinewidth{0.000000pt}%
\definecolor{currentstroke}{rgb}{0.000000,0.000000,0.000000}%
\pgfsetstrokecolor{currentstroke}%
\pgfsetstrokeopacity{0.000000}%
\pgfsetdash{}{0pt}%
\pgfpathmoveto{\pgfqpoint{0.369259in}{0.522705in}}%
\pgfpathlineto{\pgfqpoint{0.369319in}{0.524685in}}%
\pgfpathlineto{\pgfqpoint{0.372098in}{0.524246in}}%
\pgfpathlineto{\pgfqpoint{0.373941in}{0.523135in}}%
\pgfpathlineto{\pgfqpoint{0.376137in}{0.523717in}}%
\pgfpathlineto{\pgfqpoint{0.377311in}{0.522297in}}%
\pgfpathlineto{\pgfqpoint{0.374074in}{0.522045in}}%
\pgfpathlineto{\pgfqpoint{0.373276in}{0.520714in}}%
\pgfpathlineto{\pgfqpoint{0.371566in}{0.523261in}}%
\pgfpathlineto{\pgfqpoint{0.369259in}{0.522705in}}%
\pgfpathclose%
\pgfusepath{fill}%
\end{pgfscope}%
\begin{pgfscope}%
\pgfpathrectangle{\pgfqpoint{0.100000in}{0.100000in}}{\pgfqpoint{3.007045in}{1.925000in}}%
\pgfusepath{clip}%
\pgfsetbuttcap%
\pgfsetmiterjoin%
\definecolor{currentfill}{rgb}{0.886505,0.956171,0.866128}%
\pgfsetfillcolor{currentfill}%
\pgfsetlinewidth{0.000000pt}%
\definecolor{currentstroke}{rgb}{0.000000,0.000000,0.000000}%
\pgfsetstrokecolor{currentstroke}%
\pgfsetstrokeopacity{0.000000}%
\pgfsetdash{}{0pt}%
\pgfpathmoveto{\pgfqpoint{0.465052in}{0.415402in}}%
\pgfpathlineto{\pgfqpoint{0.464568in}{0.416834in}}%
\pgfpathlineto{\pgfqpoint{0.468037in}{0.416734in}}%
\pgfpathlineto{\pgfqpoint{0.469080in}{0.414904in}}%
\pgfpathlineto{\pgfqpoint{0.468333in}{0.414126in}}%
\pgfpathlineto{\pgfqpoint{0.465052in}{0.415402in}}%
\pgfpathclose%
\pgfusepath{fill}%
\end{pgfscope}%
\begin{pgfscope}%
\pgfpathrectangle{\pgfqpoint{0.100000in}{0.100000in}}{\pgfqpoint{3.007045in}{1.925000in}}%
\pgfusepath{clip}%
\pgfsetbuttcap%
\pgfsetmiterjoin%
\definecolor{currentfill}{rgb}{0.886505,0.956171,0.866128}%
\pgfsetfillcolor{currentfill}%
\pgfsetlinewidth{0.000000pt}%
\definecolor{currentstroke}{rgb}{0.000000,0.000000,0.000000}%
\pgfsetstrokecolor{currentstroke}%
\pgfsetstrokeopacity{0.000000}%
\pgfsetdash{}{0pt}%
\pgfpathmoveto{\pgfqpoint{0.351227in}{0.571842in}}%
\pgfpathlineto{\pgfqpoint{0.350829in}{0.574370in}}%
\pgfpathlineto{\pgfqpoint{0.351503in}{0.575341in}}%
\pgfpathlineto{\pgfqpoint{0.353352in}{0.574845in}}%
\pgfpathlineto{\pgfqpoint{0.353069in}{0.573567in}}%
\pgfpathlineto{\pgfqpoint{0.351227in}{0.571842in}}%
\pgfpathclose%
\pgfusepath{fill}%
\end{pgfscope}%
\begin{pgfscope}%
\pgfpathrectangle{\pgfqpoint{0.100000in}{0.100000in}}{\pgfqpoint{3.007045in}{1.925000in}}%
\pgfusepath{clip}%
\pgfsetbuttcap%
\pgfsetmiterjoin%
\definecolor{currentfill}{rgb}{0.886505,0.956171,0.866128}%
\pgfsetfillcolor{currentfill}%
\pgfsetlinewidth{0.000000pt}%
\definecolor{currentstroke}{rgb}{0.000000,0.000000,0.000000}%
\pgfsetstrokecolor{currentstroke}%
\pgfsetstrokeopacity{0.000000}%
\pgfsetdash{}{0pt}%
\pgfpathmoveto{\pgfqpoint{0.486654in}{0.401495in}}%
\pgfpathlineto{\pgfqpoint{0.487380in}{0.403172in}}%
\pgfpathlineto{\pgfqpoint{0.490544in}{0.402448in}}%
\pgfpathlineto{\pgfqpoint{0.489751in}{0.400770in}}%
\pgfpathlineto{\pgfqpoint{0.486654in}{0.401495in}}%
\pgfpathclose%
\pgfusepath{fill}%
\end{pgfscope}%
\begin{pgfscope}%
\pgfpathrectangle{\pgfqpoint{0.100000in}{0.100000in}}{\pgfqpoint{3.007045in}{1.925000in}}%
\pgfusepath{clip}%
\pgfsetbuttcap%
\pgfsetmiterjoin%
\definecolor{currentfill}{rgb}{0.886505,0.956171,0.866128}%
\pgfsetfillcolor{currentfill}%
\pgfsetlinewidth{0.000000pt}%
\definecolor{currentstroke}{rgb}{0.000000,0.000000,0.000000}%
\pgfsetstrokecolor{currentstroke}%
\pgfsetstrokeopacity{0.000000}%
\pgfsetdash{}{0pt}%
\pgfpathmoveto{\pgfqpoint{0.355532in}{0.583863in}}%
\pgfpathlineto{\pgfqpoint{0.355346in}{0.585775in}}%
\pgfpathlineto{\pgfqpoint{0.352345in}{0.588784in}}%
\pgfpathlineto{\pgfqpoint{0.354611in}{0.589330in}}%
\pgfpathlineto{\pgfqpoint{0.353122in}{0.591113in}}%
\pgfpathlineto{\pgfqpoint{0.354013in}{0.592136in}}%
\pgfpathlineto{\pgfqpoint{0.356001in}{0.592468in}}%
\pgfpathlineto{\pgfqpoint{0.357361in}{0.589332in}}%
\pgfpathlineto{\pgfqpoint{0.359106in}{0.586517in}}%
\pgfpathlineto{\pgfqpoint{0.359311in}{0.584053in}}%
\pgfpathlineto{\pgfqpoint{0.357819in}{0.583219in}}%
\pgfpathlineto{\pgfqpoint{0.355532in}{0.583863in}}%
\pgfpathclose%
\pgfusepath{fill}%
\end{pgfscope}%
\begin{pgfscope}%
\pgfpathrectangle{\pgfqpoint{0.100000in}{0.100000in}}{\pgfqpoint{3.007045in}{1.925000in}}%
\pgfusepath{clip}%
\pgfsetbuttcap%
\pgfsetmiterjoin%
\definecolor{currentfill}{rgb}{0.886505,0.956171,0.866128}%
\pgfsetfillcolor{currentfill}%
\pgfsetlinewidth{0.000000pt}%
\definecolor{currentstroke}{rgb}{0.000000,0.000000,0.000000}%
\pgfsetstrokecolor{currentstroke}%
\pgfsetstrokeopacity{0.000000}%
\pgfsetdash{}{0pt}%
\pgfpathmoveto{\pgfqpoint{0.528105in}{0.380782in}}%
\pgfpathlineto{\pgfqpoint{0.527637in}{0.383157in}}%
\pgfpathlineto{\pgfqpoint{0.528964in}{0.383518in}}%
\pgfpathlineto{\pgfqpoint{0.532084in}{0.382526in}}%
\pgfpathlineto{\pgfqpoint{0.533585in}{0.381267in}}%
\pgfpathlineto{\pgfqpoint{0.535312in}{0.380961in}}%
\pgfpathlineto{\pgfqpoint{0.536723in}{0.379546in}}%
\pgfpathlineto{\pgfqpoint{0.538639in}{0.381047in}}%
\pgfpathlineto{\pgfqpoint{0.541713in}{0.382292in}}%
\pgfpathlineto{\pgfqpoint{0.544259in}{0.379787in}}%
\pgfpathlineto{\pgfqpoint{0.542790in}{0.385194in}}%
\pgfpathlineto{\pgfqpoint{0.544767in}{0.386107in}}%
\pgfpathlineto{\pgfqpoint{0.546961in}{0.385637in}}%
\pgfpathlineto{\pgfqpoint{0.550504in}{0.383713in}}%
\pgfpathlineto{\pgfqpoint{0.549689in}{0.382189in}}%
\pgfpathlineto{\pgfqpoint{0.551450in}{0.380178in}}%
\pgfpathlineto{\pgfqpoint{0.553278in}{0.380675in}}%
\pgfpathlineto{\pgfqpoint{0.552656in}{0.378168in}}%
\pgfpathlineto{\pgfqpoint{0.549658in}{0.377895in}}%
\pgfpathlineto{\pgfqpoint{0.548439in}{0.374718in}}%
\pgfpathlineto{\pgfqpoint{0.543340in}{0.376546in}}%
\pgfpathlineto{\pgfqpoint{0.541469in}{0.376728in}}%
\pgfpathlineto{\pgfqpoint{0.540243in}{0.375113in}}%
\pgfpathlineto{\pgfqpoint{0.538796in}{0.377500in}}%
\pgfpathlineto{\pgfqpoint{0.537382in}{0.378201in}}%
\pgfpathlineto{\pgfqpoint{0.533254in}{0.378923in}}%
\pgfpathlineto{\pgfqpoint{0.531395in}{0.379938in}}%
\pgfpathlineto{\pgfqpoint{0.528105in}{0.380782in}}%
\pgfpathclose%
\pgfusepath{fill}%
\end{pgfscope}%
\begin{pgfscope}%
\pgfpathrectangle{\pgfqpoint{0.100000in}{0.100000in}}{\pgfqpoint{3.007045in}{1.925000in}}%
\pgfusepath{clip}%
\pgfsetbuttcap%
\pgfsetmiterjoin%
\definecolor{currentfill}{rgb}{0.886505,0.956171,0.866128}%
\pgfsetfillcolor{currentfill}%
\pgfsetlinewidth{0.000000pt}%
\definecolor{currentstroke}{rgb}{0.000000,0.000000,0.000000}%
\pgfsetstrokecolor{currentstroke}%
\pgfsetstrokeopacity{0.000000}%
\pgfsetdash{}{0pt}%
\pgfpathmoveto{\pgfqpoint{0.500539in}{0.394151in}}%
\pgfpathlineto{\pgfqpoint{0.500162in}{0.396419in}}%
\pgfpathlineto{\pgfqpoint{0.502003in}{0.396647in}}%
\pgfpathlineto{\pgfqpoint{0.502049in}{0.394805in}}%
\pgfpathlineto{\pgfqpoint{0.500539in}{0.394151in}}%
\pgfpathclose%
\pgfusepath{fill}%
\end{pgfscope}%
\begin{pgfscope}%
\pgfpathrectangle{\pgfqpoint{0.100000in}{0.100000in}}{\pgfqpoint{3.007045in}{1.925000in}}%
\pgfusepath{clip}%
\pgfsetbuttcap%
\pgfsetmiterjoin%
\definecolor{currentfill}{rgb}{0.886505,0.956171,0.866128}%
\pgfsetfillcolor{currentfill}%
\pgfsetlinewidth{0.000000pt}%
\definecolor{currentstroke}{rgb}{0.000000,0.000000,0.000000}%
\pgfsetstrokecolor{currentstroke}%
\pgfsetstrokeopacity{0.000000}%
\pgfsetdash{}{0pt}%
\pgfpathmoveto{\pgfqpoint{0.507877in}{0.388268in}}%
\pgfpathlineto{\pgfqpoint{0.508331in}{0.389076in}}%
\pgfpathlineto{\pgfqpoint{0.513005in}{0.389208in}}%
\pgfpathlineto{\pgfqpoint{0.515777in}{0.391018in}}%
\pgfpathlineto{\pgfqpoint{0.519331in}{0.390800in}}%
\pgfpathlineto{\pgfqpoint{0.521540in}{0.387996in}}%
\pgfpathlineto{\pgfqpoint{0.523561in}{0.390645in}}%
\pgfpathlineto{\pgfqpoint{0.525758in}{0.391376in}}%
\pgfpathlineto{\pgfqpoint{0.528197in}{0.390858in}}%
\pgfpathlineto{\pgfqpoint{0.530050in}{0.389563in}}%
\pgfpathlineto{\pgfqpoint{0.530693in}{0.387435in}}%
\pgfpathlineto{\pgfqpoint{0.530317in}{0.386066in}}%
\pgfpathlineto{\pgfqpoint{0.528322in}{0.384961in}}%
\pgfpathlineto{\pgfqpoint{0.521632in}{0.386913in}}%
\pgfpathlineto{\pgfqpoint{0.517279in}{0.385660in}}%
\pgfpathlineto{\pgfqpoint{0.515618in}{0.386416in}}%
\pgfpathlineto{\pgfqpoint{0.507877in}{0.388268in}}%
\pgfpathclose%
\pgfusepath{fill}%
\end{pgfscope}%
\begin{pgfscope}%
\pgfpathrectangle{\pgfqpoint{0.100000in}{0.100000in}}{\pgfqpoint{3.007045in}{1.925000in}}%
\pgfusepath{clip}%
\pgfsetbuttcap%
\pgfsetmiterjoin%
\definecolor{currentfill}{rgb}{0.886505,0.956171,0.866128}%
\pgfsetfillcolor{currentfill}%
\pgfsetlinewidth{0.000000pt}%
\definecolor{currentstroke}{rgb}{0.000000,0.000000,0.000000}%
\pgfsetstrokecolor{currentstroke}%
\pgfsetstrokeopacity{0.000000}%
\pgfsetdash{}{0pt}%
\pgfpathmoveto{\pgfqpoint{0.453614in}{0.421300in}}%
\pgfpathlineto{\pgfqpoint{0.451899in}{0.423133in}}%
\pgfpathlineto{\pgfqpoint{0.450880in}{0.425283in}}%
\pgfpathlineto{\pgfqpoint{0.449241in}{0.426901in}}%
\pgfpathlineto{\pgfqpoint{0.449633in}{0.429133in}}%
\pgfpathlineto{\pgfqpoint{0.453091in}{0.425783in}}%
\pgfpathlineto{\pgfqpoint{0.455129in}{0.421602in}}%
\pgfpathlineto{\pgfqpoint{0.453614in}{0.421300in}}%
\pgfpathclose%
\pgfusepath{fill}%
\end{pgfscope}%
\begin{pgfscope}%
\pgfpathrectangle{\pgfqpoint{0.100000in}{0.100000in}}{\pgfqpoint{3.007045in}{1.925000in}}%
\pgfusepath{clip}%
\pgfsetbuttcap%
\pgfsetmiterjoin%
\definecolor{currentfill}{rgb}{0.886505,0.956171,0.866128}%
\pgfsetfillcolor{currentfill}%
\pgfsetlinewidth{0.000000pt}%
\definecolor{currentstroke}{rgb}{0.000000,0.000000,0.000000}%
\pgfsetstrokecolor{currentstroke}%
\pgfsetstrokeopacity{0.000000}%
\pgfsetdash{}{0pt}%
\pgfpathmoveto{\pgfqpoint{0.444349in}{0.432536in}}%
\pgfpathlineto{\pgfqpoint{0.443210in}{0.434576in}}%
\pgfpathlineto{\pgfqpoint{0.440911in}{0.435309in}}%
\pgfpathlineto{\pgfqpoint{0.440187in}{0.438580in}}%
\pgfpathlineto{\pgfqpoint{0.442735in}{0.437137in}}%
\pgfpathlineto{\pgfqpoint{0.443790in}{0.435233in}}%
\pgfpathlineto{\pgfqpoint{0.445554in}{0.435638in}}%
\pgfpathlineto{\pgfqpoint{0.448969in}{0.435304in}}%
\pgfpathlineto{\pgfqpoint{0.450814in}{0.436419in}}%
\pgfpathlineto{\pgfqpoint{0.452553in}{0.435695in}}%
\pgfpathlineto{\pgfqpoint{0.453462in}{0.433672in}}%
\pgfpathlineto{\pgfqpoint{0.453143in}{0.432430in}}%
\pgfpathlineto{\pgfqpoint{0.451239in}{0.431442in}}%
\pgfpathlineto{\pgfqpoint{0.449605in}{0.432844in}}%
\pgfpathlineto{\pgfqpoint{0.448738in}{0.430447in}}%
\pgfpathlineto{\pgfqpoint{0.447526in}{0.430880in}}%
\pgfpathlineto{\pgfqpoint{0.445671in}{0.433246in}}%
\pgfpathlineto{\pgfqpoint{0.444349in}{0.432536in}}%
\pgfpathclose%
\pgfusepath{fill}%
\end{pgfscope}%
\begin{pgfscope}%
\pgfpathrectangle{\pgfqpoint{0.100000in}{0.100000in}}{\pgfqpoint{3.007045in}{1.925000in}}%
\pgfusepath{clip}%
\pgfsetbuttcap%
\pgfsetmiterjoin%
\definecolor{currentfill}{rgb}{0.389020,0.736332,0.430450}%
\pgfsetfillcolor{currentfill}%
\pgfsetlinewidth{0.000000pt}%
\definecolor{currentstroke}{rgb}{0.000000,0.000000,0.000000}%
\pgfsetstrokecolor{currentstroke}%
\pgfsetstrokeopacity{0.000000}%
\pgfsetdash{}{0pt}%
\pgfpathmoveto{\pgfqpoint{1.876321in}{0.768157in}}%
\pgfpathlineto{\pgfqpoint{1.870545in}{0.767215in}}%
\pgfpathlineto{\pgfqpoint{1.856034in}{0.767401in}}%
\pgfpathlineto{\pgfqpoint{1.856051in}{0.769311in}}%
\pgfpathlineto{\pgfqpoint{1.842812in}{0.769497in}}%
\pgfpathlineto{\pgfqpoint{1.836570in}{0.777705in}}%
\pgfpathlineto{\pgfqpoint{1.829258in}{0.781767in}}%
\pgfpathlineto{\pgfqpoint{1.829586in}{0.787129in}}%
\pgfpathlineto{\pgfqpoint{1.828440in}{0.791464in}}%
\pgfpathlineto{\pgfqpoint{1.824314in}{0.796016in}}%
\pgfpathlineto{\pgfqpoint{1.822337in}{0.805756in}}%
\pgfpathlineto{\pgfqpoint{1.821492in}{0.816167in}}%
\pgfpathlineto{\pgfqpoint{1.813675in}{0.816167in}}%
\pgfpathlineto{\pgfqpoint{1.812989in}{0.826934in}}%
\pgfpathlineto{\pgfqpoint{1.830412in}{0.826752in}}%
\pgfpathlineto{\pgfqpoint{1.830433in}{0.847601in}}%
\pgfpathlineto{\pgfqpoint{1.833670in}{0.850751in}}%
\pgfpathlineto{\pgfqpoint{1.842321in}{0.853020in}}%
\pgfpathlineto{\pgfqpoint{1.859472in}{0.853168in}}%
\pgfpathlineto{\pgfqpoint{1.865214in}{0.855241in}}%
\pgfpathlineto{\pgfqpoint{1.882494in}{0.855536in}}%
\pgfpathlineto{\pgfqpoint{1.882114in}{0.846605in}}%
\pgfpathlineto{\pgfqpoint{1.887792in}{0.846596in}}%
\pgfpathlineto{\pgfqpoint{1.891528in}{0.843729in}}%
\pgfpathlineto{\pgfqpoint{1.892466in}{0.838016in}}%
\pgfpathlineto{\pgfqpoint{1.898412in}{0.836142in}}%
\pgfpathlineto{\pgfqpoint{1.899400in}{0.832281in}}%
\pgfpathlineto{\pgfqpoint{1.899320in}{0.814784in}}%
\pgfpathlineto{\pgfqpoint{1.887940in}{0.814920in}}%
\pgfpathlineto{\pgfqpoint{1.887657in}{0.795131in}}%
\pgfpathlineto{\pgfqpoint{1.884971in}{0.791103in}}%
\pgfpathlineto{\pgfqpoint{1.876541in}{0.789626in}}%
\pgfpathlineto{\pgfqpoint{1.876321in}{0.768157in}}%
\pgfpathclose%
\pgfusepath{fill}%
\end{pgfscope}%
\begin{pgfscope}%
\pgfpathrectangle{\pgfqpoint{0.100000in}{0.100000in}}{\pgfqpoint{3.007045in}{1.925000in}}%
\pgfusepath{clip}%
\pgfsetbuttcap%
\pgfsetmiterjoin%
\definecolor{currentfill}{rgb}{0.507497,0.793172,0.505990}%
\pgfsetfillcolor{currentfill}%
\pgfsetlinewidth{0.000000pt}%
\definecolor{currentstroke}{rgb}{0.000000,0.000000,0.000000}%
\pgfsetstrokecolor{currentstroke}%
\pgfsetstrokeopacity{0.000000}%
\pgfsetdash{}{0pt}%
\pgfpathmoveto{\pgfqpoint{0.596303in}{1.682378in}}%
\pgfpathlineto{\pgfqpoint{0.563962in}{1.691905in}}%
\pgfpathlineto{\pgfqpoint{0.533873in}{1.701067in}}%
\pgfpathlineto{\pgfqpoint{0.536189in}{1.704696in}}%
\pgfpathlineto{\pgfqpoint{0.535329in}{1.713270in}}%
\pgfpathlineto{\pgfqpoint{0.541696in}{1.715411in}}%
\pgfpathlineto{\pgfqpoint{0.541409in}{1.723741in}}%
\pgfpathlineto{\pgfqpoint{0.545552in}{1.726930in}}%
\pgfpathlineto{\pgfqpoint{0.547023in}{1.735645in}}%
\pgfpathlineto{\pgfqpoint{0.540688in}{1.747112in}}%
\pgfpathlineto{\pgfqpoint{0.542490in}{1.755480in}}%
\pgfpathlineto{\pgfqpoint{0.545786in}{1.757536in}}%
\pgfpathlineto{\pgfqpoint{0.551653in}{1.755806in}}%
\pgfpathlineto{\pgfqpoint{0.561232in}{1.755167in}}%
\pgfpathlineto{\pgfqpoint{0.558944in}{1.759738in}}%
\pgfpathlineto{\pgfqpoint{0.563789in}{1.776383in}}%
\pgfpathlineto{\pgfqpoint{0.567701in}{1.775221in}}%
\pgfpathlineto{\pgfqpoint{0.616044in}{1.760740in}}%
\pgfpathlineto{\pgfqpoint{0.640972in}{1.753699in}}%
\pgfpathlineto{\pgfqpoint{0.637219in}{1.740644in}}%
\pgfpathlineto{\pgfqpoint{0.630951in}{1.740790in}}%
\pgfpathlineto{\pgfqpoint{0.619982in}{1.737934in}}%
\pgfpathlineto{\pgfqpoint{0.610892in}{1.738691in}}%
\pgfpathlineto{\pgfqpoint{0.606293in}{1.736988in}}%
\pgfpathlineto{\pgfqpoint{0.608690in}{1.732151in}}%
\pgfpathlineto{\pgfqpoint{0.608895in}{1.725281in}}%
\pgfpathlineto{\pgfqpoint{0.602494in}{1.724955in}}%
\pgfpathlineto{\pgfqpoint{0.595993in}{1.712179in}}%
\pgfpathlineto{\pgfqpoint{0.598195in}{1.705394in}}%
\pgfpathlineto{\pgfqpoint{0.593379in}{1.692150in}}%
\pgfpathlineto{\pgfqpoint{0.593360in}{1.686222in}}%
\pgfpathlineto{\pgfqpoint{0.596303in}{1.682378in}}%
\pgfpathclose%
\pgfusepath{fill}%
\end{pgfscope}%
\begin{pgfscope}%
\pgfpathrectangle{\pgfqpoint{0.100000in}{0.100000in}}{\pgfqpoint{3.007045in}{1.925000in}}%
\pgfusepath{clip}%
\pgfsetbuttcap%
\pgfsetmiterjoin%
\definecolor{currentfill}{rgb}{0.751173,0.901423,0.724491}%
\pgfsetfillcolor{currentfill}%
\pgfsetlinewidth{0.000000pt}%
\definecolor{currentstroke}{rgb}{0.000000,0.000000,0.000000}%
\pgfsetstrokecolor{currentstroke}%
\pgfsetstrokeopacity{0.000000}%
\pgfsetdash{}{0pt}%
\pgfpathmoveto{\pgfqpoint{1.473659in}{0.890543in}}%
\pgfpathlineto{\pgfqpoint{1.444186in}{0.892364in}}%
\pgfpathlineto{\pgfqpoint{1.415306in}{0.894451in}}%
\pgfpathlineto{\pgfqpoint{1.417413in}{0.923222in}}%
\pgfpathlineto{\pgfqpoint{1.419469in}{0.951860in}}%
\pgfpathlineto{\pgfqpoint{1.477103in}{0.948069in}}%
\pgfpathlineto{\pgfqpoint{1.505836in}{0.946347in}}%
\pgfpathlineto{\pgfqpoint{1.505202in}{0.934838in}}%
\pgfpathlineto{\pgfqpoint{1.504252in}{0.917589in}}%
\pgfpathlineto{\pgfqpoint{1.475346in}{0.919280in}}%
\pgfpathlineto{\pgfqpoint{1.473659in}{0.890543in}}%
\pgfpathclose%
\pgfusepath{fill}%
\end{pgfscope}%
\begin{pgfscope}%
\pgfpathrectangle{\pgfqpoint{0.100000in}{0.100000in}}{\pgfqpoint{3.007045in}{1.925000in}}%
\pgfusepath{clip}%
\pgfsetbuttcap%
\pgfsetmiterjoin%
\definecolor{currentfill}{rgb}{0.562860,0.819008,0.551511}%
\pgfsetfillcolor{currentfill}%
\pgfsetlinewidth{0.000000pt}%
\definecolor{currentstroke}{rgb}{0.000000,0.000000,0.000000}%
\pgfsetstrokecolor{currentstroke}%
\pgfsetstrokeopacity{0.000000}%
\pgfsetdash{}{0pt}%
\pgfpathmoveto{\pgfqpoint{1.586901in}{0.635387in}}%
\pgfpathlineto{\pgfqpoint{1.569891in}{0.623670in}}%
\pgfpathlineto{\pgfqpoint{1.571686in}{0.614167in}}%
\pgfpathlineto{\pgfqpoint{1.576374in}{0.610043in}}%
\pgfpathlineto{\pgfqpoint{1.575754in}{0.602872in}}%
\pgfpathlineto{\pgfqpoint{1.546174in}{0.604074in}}%
\pgfpathlineto{\pgfqpoint{1.538952in}{0.605702in}}%
\pgfpathlineto{\pgfqpoint{1.540604in}{0.640218in}}%
\pgfpathlineto{\pgfqpoint{1.532433in}{0.642045in}}%
\pgfpathlineto{\pgfqpoint{1.527773in}{0.637409in}}%
\pgfpathlineto{\pgfqpoint{1.520806in}{0.642095in}}%
\pgfpathlineto{\pgfqpoint{1.510361in}{0.642115in}}%
\pgfpathlineto{\pgfqpoint{1.505217in}{0.648714in}}%
\pgfpathlineto{\pgfqpoint{1.507461in}{0.683119in}}%
\pgfpathlineto{\pgfqpoint{1.551784in}{0.680740in}}%
\pgfpathlineto{\pgfqpoint{1.565148in}{0.655103in}}%
\pgfpathlineto{\pgfqpoint{1.575122in}{0.656245in}}%
\pgfpathlineto{\pgfqpoint{1.586901in}{0.635387in}}%
\pgfpathclose%
\pgfusepath{fill}%
\end{pgfscope}%
\begin{pgfscope}%
\pgfpathrectangle{\pgfqpoint{0.100000in}{0.100000in}}{\pgfqpoint{3.007045in}{1.925000in}}%
\pgfusepath{clip}%
\pgfsetbuttcap%
\pgfsetmiterjoin%
\definecolor{currentfill}{rgb}{0.401569,0.742484,0.436601}%
\pgfsetfillcolor{currentfill}%
\pgfsetlinewidth{0.000000pt}%
\definecolor{currentstroke}{rgb}{0.000000,0.000000,0.000000}%
\pgfsetstrokecolor{currentstroke}%
\pgfsetstrokeopacity{0.000000}%
\pgfsetdash{}{0pt}%
\pgfpathmoveto{\pgfqpoint{2.442205in}{0.558231in}}%
\pgfpathlineto{\pgfqpoint{2.433773in}{0.561843in}}%
\pgfpathlineto{\pgfqpoint{2.427536in}{0.571164in}}%
\pgfpathlineto{\pgfqpoint{2.419054in}{0.576060in}}%
\pgfpathlineto{\pgfqpoint{2.409749in}{0.579100in}}%
\pgfpathlineto{\pgfqpoint{2.405641in}{0.582074in}}%
\pgfpathlineto{\pgfqpoint{2.408711in}{0.592974in}}%
\pgfpathlineto{\pgfqpoint{2.415159in}{0.598741in}}%
\pgfpathlineto{\pgfqpoint{2.413904in}{0.601638in}}%
\pgfpathlineto{\pgfqpoint{2.418230in}{0.606839in}}%
\pgfpathlineto{\pgfqpoint{2.416778in}{0.611640in}}%
\pgfpathlineto{\pgfqpoint{2.423706in}{0.617445in}}%
\pgfpathlineto{\pgfqpoint{2.422704in}{0.621971in}}%
\pgfpathlineto{\pgfqpoint{2.481583in}{0.625736in}}%
\pgfpathlineto{\pgfqpoint{2.488716in}{0.626214in}}%
\pgfpathlineto{\pgfqpoint{2.492884in}{0.596677in}}%
\pgfpathlineto{\pgfqpoint{2.488025in}{0.585155in}}%
\pgfpathlineto{\pgfqpoint{2.487697in}{0.577489in}}%
\pgfpathlineto{\pgfqpoint{2.485231in}{0.574838in}}%
\pgfpathlineto{\pgfqpoint{2.479736in}{0.576043in}}%
\pgfpathlineto{\pgfqpoint{2.476107in}{0.580682in}}%
\pgfpathlineto{\pgfqpoint{2.471783in}{0.578648in}}%
\pgfpathlineto{\pgfqpoint{2.470809in}{0.572873in}}%
\pgfpathlineto{\pgfqpoint{2.446412in}{0.569336in}}%
\pgfpathlineto{\pgfqpoint{2.444618in}{0.558566in}}%
\pgfpathlineto{\pgfqpoint{2.442205in}{0.558231in}}%
\pgfpathclose%
\pgfusepath{fill}%
\end{pgfscope}%
\begin{pgfscope}%
\pgfpathrectangle{\pgfqpoint{0.100000in}{0.100000in}}{\pgfqpoint{3.007045in}{1.925000in}}%
\pgfusepath{clip}%
\pgfsetbuttcap%
\pgfsetmiterjoin%
\definecolor{currentfill}{rgb}{0.535179,0.806090,0.528750}%
\pgfsetfillcolor{currentfill}%
\pgfsetlinewidth{0.000000pt}%
\definecolor{currentstroke}{rgb}{0.000000,0.000000,0.000000}%
\pgfsetstrokecolor{currentstroke}%
\pgfsetstrokeopacity{0.000000}%
\pgfsetdash{}{0pt}%
\pgfpathmoveto{\pgfqpoint{2.135867in}{0.998378in}}%
\pgfpathlineto{\pgfqpoint{2.141030in}{0.985137in}}%
\pgfpathlineto{\pgfqpoint{2.140155in}{0.982191in}}%
\pgfpathlineto{\pgfqpoint{2.117091in}{0.980740in}}%
\pgfpathlineto{\pgfqpoint{2.116101in}{0.997143in}}%
\pgfpathlineto{\pgfqpoint{2.135867in}{0.998378in}}%
\pgfpathclose%
\pgfusepath{fill}%
\end{pgfscope}%
\begin{pgfscope}%
\pgfpathrectangle{\pgfqpoint{0.100000in}{0.100000in}}{\pgfqpoint{3.007045in}{1.925000in}}%
\pgfusepath{clip}%
\pgfsetbuttcap%
\pgfsetmiterjoin%
\definecolor{currentfill}{rgb}{0.313725,0.699423,0.393541}%
\pgfsetfillcolor{currentfill}%
\pgfsetlinewidth{0.000000pt}%
\definecolor{currentstroke}{rgb}{0.000000,0.000000,0.000000}%
\pgfsetstrokecolor{currentstroke}%
\pgfsetstrokeopacity{0.000000}%
\pgfsetdash{}{0pt}%
\pgfpathmoveto{\pgfqpoint{1.861397in}{1.624038in}}%
\pgfpathlineto{\pgfqpoint{1.861813in}{1.606700in}}%
\pgfpathlineto{\pgfqpoint{1.844441in}{1.606089in}}%
\pgfpathlineto{\pgfqpoint{1.844395in}{1.612207in}}%
\pgfpathlineto{\pgfqpoint{1.826933in}{1.611742in}}%
\pgfpathlineto{\pgfqpoint{1.827033in}{1.605997in}}%
\pgfpathlineto{\pgfqpoint{1.801088in}{1.605771in}}%
\pgfpathlineto{\pgfqpoint{1.802129in}{1.614963in}}%
\pgfpathlineto{\pgfqpoint{1.798627in}{1.617392in}}%
\pgfpathlineto{\pgfqpoint{1.793675in}{1.615402in}}%
\pgfpathlineto{\pgfqpoint{1.782448in}{1.621431in}}%
\pgfpathlineto{\pgfqpoint{1.782023in}{1.648669in}}%
\pgfpathlineto{\pgfqpoint{1.787806in}{1.648648in}}%
\pgfpathlineto{\pgfqpoint{1.787190in}{1.689017in}}%
\pgfpathlineto{\pgfqpoint{1.798628in}{1.689186in}}%
\pgfpathlineto{\pgfqpoint{1.798346in}{1.718105in}}%
\pgfpathlineto{\pgfqpoint{1.827086in}{1.718431in}}%
\pgfpathlineto{\pgfqpoint{1.827059in}{1.721935in}}%
\pgfpathlineto{\pgfqpoint{1.858121in}{1.722001in}}%
\pgfpathlineto{\pgfqpoint{1.859255in}{1.710576in}}%
\pgfpathlineto{\pgfqpoint{1.859346in}{1.688805in}}%
\pgfpathlineto{\pgfqpoint{1.860483in}{1.664391in}}%
\pgfpathlineto{\pgfqpoint{1.860671in}{1.638052in}}%
\pgfpathlineto{\pgfqpoint{1.861397in}{1.624038in}}%
\pgfpathclose%
\pgfusepath{fill}%
\end{pgfscope}%
\begin{pgfscope}%
\pgfpathrectangle{\pgfqpoint{0.100000in}{0.100000in}}{\pgfqpoint{3.007045in}{1.925000in}}%
\pgfusepath{clip}%
\pgfsetbuttcap%
\pgfsetmiterjoin%
\definecolor{currentfill}{rgb}{0.535179,0.806090,0.528750}%
\pgfsetfillcolor{currentfill}%
\pgfsetlinewidth{0.000000pt}%
\definecolor{currentstroke}{rgb}{0.000000,0.000000,0.000000}%
\pgfsetstrokecolor{currentstroke}%
\pgfsetstrokeopacity{0.000000}%
\pgfsetdash{}{0pt}%
\pgfpathmoveto{\pgfqpoint{1.348840in}{1.549898in}}%
\pgfpathlineto{\pgfqpoint{1.303187in}{1.554906in}}%
\pgfpathlineto{\pgfqpoint{1.305764in}{1.578064in}}%
\pgfpathlineto{\pgfqpoint{1.308128in}{1.577801in}}%
\pgfpathlineto{\pgfqpoint{1.311526in}{1.606441in}}%
\pgfpathlineto{\pgfqpoint{1.313493in}{1.606222in}}%
\pgfpathlineto{\pgfqpoint{1.316018in}{1.629285in}}%
\pgfpathlineto{\pgfqpoint{1.318468in}{1.629024in}}%
\pgfpathlineto{\pgfqpoint{1.335389in}{1.627188in}}%
\pgfpathlineto{\pgfqpoint{1.334802in}{1.621794in}}%
\pgfpathlineto{\pgfqpoint{1.340576in}{1.621215in}}%
\pgfpathlineto{\pgfqpoint{1.339970in}{1.615629in}}%
\pgfpathlineto{\pgfqpoint{1.351714in}{1.614337in}}%
\pgfpathlineto{\pgfqpoint{1.355655in}{1.612365in}}%
\pgfpathlineto{\pgfqpoint{1.403892in}{1.607820in}}%
\pgfpathlineto{\pgfqpoint{1.406341in}{1.607594in}}%
\pgfpathlineto{\pgfqpoint{1.401374in}{1.559284in}}%
\pgfpathlineto{\pgfqpoint{1.351084in}{1.564019in}}%
\pgfpathlineto{\pgfqpoint{1.348840in}{1.549898in}}%
\pgfpathclose%
\pgfusepath{fill}%
\end{pgfscope}%
\begin{pgfscope}%
\pgfpathrectangle{\pgfqpoint{0.100000in}{0.100000in}}{\pgfqpoint{3.007045in}{1.925000in}}%
\pgfusepath{clip}%
\pgfsetbuttcap%
\pgfsetmiterjoin%
\definecolor{currentfill}{rgb}{0.294902,0.690196,0.384314}%
\pgfsetfillcolor{currentfill}%
\pgfsetlinewidth{0.000000pt}%
\definecolor{currentstroke}{rgb}{0.000000,0.000000,0.000000}%
\pgfsetstrokecolor{currentstroke}%
\pgfsetstrokeopacity{0.000000}%
\pgfsetdash{}{0pt}%
\pgfpathmoveto{\pgfqpoint{2.199018in}{1.471645in}}%
\pgfpathlineto{\pgfqpoint{2.180375in}{1.470113in}}%
\pgfpathlineto{\pgfqpoint{2.177932in}{1.479844in}}%
\pgfpathlineto{\pgfqpoint{2.175073in}{1.484492in}}%
\pgfpathlineto{\pgfqpoint{2.178913in}{1.490543in}}%
\pgfpathlineto{\pgfqpoint{2.184775in}{1.505418in}}%
\pgfpathlineto{\pgfqpoint{2.185259in}{1.517003in}}%
\pgfpathlineto{\pgfqpoint{2.204865in}{1.518579in}}%
\pgfpathlineto{\pgfqpoint{2.207015in}{1.495613in}}%
\pgfpathlineto{\pgfqpoint{2.196560in}{1.494814in}}%
\pgfpathlineto{\pgfqpoint{2.199018in}{1.471645in}}%
\pgfpathclose%
\pgfusepath{fill}%
\end{pgfscope}%
\begin{pgfscope}%
\pgfpathrectangle{\pgfqpoint{0.100000in}{0.100000in}}{\pgfqpoint{3.007045in}{1.925000in}}%
\pgfusepath{clip}%
\pgfsetbuttcap%
\pgfsetmiterjoin%
\definecolor{currentfill}{rgb}{0.755848,0.903391,0.729043}%
\pgfsetfillcolor{currentfill}%
\pgfsetlinewidth{0.000000pt}%
\definecolor{currentstroke}{rgb}{0.000000,0.000000,0.000000}%
\pgfsetstrokecolor{currentstroke}%
\pgfsetstrokeopacity{0.000000}%
\pgfsetdash{}{0pt}%
\pgfpathmoveto{\pgfqpoint{1.011802in}{1.226459in}}%
\pgfpathlineto{\pgfqpoint{1.010713in}{1.219317in}}%
\pgfpathlineto{\pgfqpoint{1.007174in}{1.215093in}}%
\pgfpathlineto{\pgfqpoint{1.007785in}{1.211607in}}%
\pgfpathlineto{\pgfqpoint{1.001745in}{1.203071in}}%
\pgfpathlineto{\pgfqpoint{0.998322in}{1.188795in}}%
\pgfpathlineto{\pgfqpoint{1.000947in}{1.182123in}}%
\pgfpathlineto{\pgfqpoint{0.999346in}{1.179674in}}%
\pgfpathlineto{\pgfqpoint{1.002802in}{1.166666in}}%
\pgfpathlineto{\pgfqpoint{1.001312in}{1.163987in}}%
\pgfpathlineto{\pgfqpoint{0.968213in}{1.169632in}}%
\pgfpathlineto{\pgfqpoint{0.936286in}{1.175309in}}%
\pgfpathlineto{\pgfqpoint{0.936407in}{1.175969in}}%
\pgfpathlineto{\pgfqpoint{0.942870in}{1.209747in}}%
\pgfpathlineto{\pgfqpoint{0.914996in}{1.214937in}}%
\pgfpathlineto{\pgfqpoint{0.907070in}{1.217261in}}%
\pgfpathlineto{\pgfqpoint{0.910572in}{1.236135in}}%
\pgfpathlineto{\pgfqpoint{0.915780in}{1.238145in}}%
\pgfpathlineto{\pgfqpoint{0.924573in}{1.236589in}}%
\pgfpathlineto{\pgfqpoint{0.927198in}{1.241278in}}%
\pgfpathlineto{\pgfqpoint{0.929754in}{1.256515in}}%
\pgfpathlineto{\pgfqpoint{0.937487in}{1.261054in}}%
\pgfpathlineto{\pgfqpoint{0.934985in}{1.263476in}}%
\pgfpathlineto{\pgfqpoint{0.974041in}{1.256437in}}%
\pgfpathlineto{\pgfqpoint{0.990893in}{1.253051in}}%
\pgfpathlineto{\pgfqpoint{1.022492in}{1.247657in}}%
\pgfpathlineto{\pgfqpoint{1.016300in}{1.244056in}}%
\pgfpathlineto{\pgfqpoint{1.015989in}{1.239279in}}%
\pgfpathlineto{\pgfqpoint{1.012248in}{1.233395in}}%
\pgfpathlineto{\pgfqpoint{1.011802in}{1.226459in}}%
\pgfpathclose%
\pgfusepath{fill}%
\end{pgfscope}%
\begin{pgfscope}%
\pgfpathrectangle{\pgfqpoint{0.100000in}{0.100000in}}{\pgfqpoint{3.007045in}{1.925000in}}%
\pgfusepath{clip}%
\pgfsetbuttcap%
\pgfsetmiterjoin%
\definecolor{currentfill}{rgb}{0.164014,0.573641,0.291995}%
\pgfsetfillcolor{currentfill}%
\pgfsetlinewidth{0.000000pt}%
\definecolor{currentstroke}{rgb}{0.000000,0.000000,0.000000}%
\pgfsetstrokecolor{currentstroke}%
\pgfsetstrokeopacity{0.000000}%
\pgfsetdash{}{0pt}%
\pgfpathmoveto{\pgfqpoint{0.397486in}{1.548285in}}%
\pgfpathlineto{\pgfqpoint{0.393613in}{1.542963in}}%
\pgfpathlineto{\pgfqpoint{0.385318in}{1.540362in}}%
\pgfpathlineto{\pgfqpoint{0.384988in}{1.532462in}}%
\pgfpathlineto{\pgfqpoint{0.380582in}{1.527032in}}%
\pgfpathlineto{\pgfqpoint{0.381897in}{1.521054in}}%
\pgfpathlineto{\pgfqpoint{0.381028in}{1.514238in}}%
\pgfpathlineto{\pgfqpoint{0.378145in}{1.511673in}}%
\pgfpathlineto{\pgfqpoint{0.372999in}{1.513265in}}%
\pgfpathlineto{\pgfqpoint{0.374697in}{1.518495in}}%
\pgfpathlineto{\pgfqpoint{0.360884in}{1.523094in}}%
\pgfpathlineto{\pgfqpoint{0.362297in}{1.539826in}}%
\pgfpathlineto{\pgfqpoint{0.358579in}{1.545963in}}%
\pgfpathlineto{\pgfqpoint{0.364222in}{1.554694in}}%
\pgfpathlineto{\pgfqpoint{0.365111in}{1.559248in}}%
\pgfpathlineto{\pgfqpoint{0.360771in}{1.567349in}}%
\pgfpathlineto{\pgfqpoint{0.362017in}{1.572568in}}%
\pgfpathlineto{\pgfqpoint{0.361686in}{1.582534in}}%
\pgfpathlineto{\pgfqpoint{0.365712in}{1.592867in}}%
\pgfpathlineto{\pgfqpoint{0.368947in}{1.597569in}}%
\pgfpathlineto{\pgfqpoint{0.369577in}{1.603351in}}%
\pgfpathlineto{\pgfqpoint{0.366697in}{1.609297in}}%
\pgfpathlineto{\pgfqpoint{0.367089in}{1.616757in}}%
\pgfpathlineto{\pgfqpoint{0.372778in}{1.622728in}}%
\pgfpathlineto{\pgfqpoint{0.382988in}{1.619317in}}%
\pgfpathlineto{\pgfqpoint{0.386071in}{1.611691in}}%
\pgfpathlineto{\pgfqpoint{0.382499in}{1.599727in}}%
\pgfpathlineto{\pgfqpoint{0.389344in}{1.598489in}}%
\pgfpathlineto{\pgfqpoint{0.397064in}{1.602137in}}%
\pgfpathlineto{\pgfqpoint{0.404299in}{1.600395in}}%
\pgfpathlineto{\pgfqpoint{0.395865in}{1.592654in}}%
\pgfpathlineto{\pgfqpoint{0.390435in}{1.585257in}}%
\pgfpathlineto{\pgfqpoint{0.384929in}{1.586076in}}%
\pgfpathlineto{\pgfqpoint{0.381248in}{1.578474in}}%
\pgfpathlineto{\pgfqpoint{0.387396in}{1.575981in}}%
\pgfpathlineto{\pgfqpoint{0.389591in}{1.567203in}}%
\pgfpathlineto{\pgfqpoint{0.384587in}{1.562231in}}%
\pgfpathlineto{\pgfqpoint{0.383232in}{1.552564in}}%
\pgfpathlineto{\pgfqpoint{0.397486in}{1.548285in}}%
\pgfpathclose%
\pgfusepath{fill}%
\end{pgfscope}%
\begin{pgfscope}%
\pgfpathrectangle{\pgfqpoint{0.100000in}{0.100000in}}{\pgfqpoint{3.007045in}{1.925000in}}%
\pgfusepath{clip}%
\pgfsetbuttcap%
\pgfsetmiterjoin%
\definecolor{currentfill}{rgb}{0.585006,0.829343,0.569719}%
\pgfsetfillcolor{currentfill}%
\pgfsetlinewidth{0.000000pt}%
\definecolor{currentstroke}{rgb}{0.000000,0.000000,0.000000}%
\pgfsetstrokecolor{currentstroke}%
\pgfsetstrokeopacity{0.000000}%
\pgfsetdash{}{0pt}%
\pgfpathmoveto{\pgfqpoint{2.283471in}{1.248041in}}%
\pgfpathlineto{\pgfqpoint{2.285375in}{1.228022in}}%
\pgfpathlineto{\pgfqpoint{2.265786in}{1.225833in}}%
\pgfpathlineto{\pgfqpoint{2.266331in}{1.211577in}}%
\pgfpathlineto{\pgfqpoint{2.262289in}{1.211080in}}%
\pgfpathlineto{\pgfqpoint{2.247429in}{1.209448in}}%
\pgfpathlineto{\pgfqpoint{2.247375in}{1.219993in}}%
\pgfpathlineto{\pgfqpoint{2.233016in}{1.218375in}}%
\pgfpathlineto{\pgfqpoint{2.231222in}{1.236425in}}%
\pgfpathlineto{\pgfqpoint{2.228217in}{1.264785in}}%
\pgfpathlineto{\pgfqpoint{2.239459in}{1.266008in}}%
\pgfpathlineto{\pgfqpoint{2.248890in}{1.267034in}}%
\pgfpathlineto{\pgfqpoint{2.249581in}{1.261349in}}%
\pgfpathlineto{\pgfqpoint{2.261811in}{1.262682in}}%
\pgfpathlineto{\pgfqpoint{2.262780in}{1.245532in}}%
\pgfpathlineto{\pgfqpoint{2.283471in}{1.248041in}}%
\pgfpathclose%
\pgfusepath{fill}%
\end{pgfscope}%
\begin{pgfscope}%
\pgfpathrectangle{\pgfqpoint{0.100000in}{0.100000in}}{\pgfqpoint{3.007045in}{1.925000in}}%
\pgfusepath{clip}%
\pgfsetbuttcap%
\pgfsetmiterjoin%
\definecolor{currentfill}{rgb}{0.704421,0.881738,0.678970}%
\pgfsetfillcolor{currentfill}%
\pgfsetlinewidth{0.000000pt}%
\definecolor{currentstroke}{rgb}{0.000000,0.000000,0.000000}%
\pgfsetstrokecolor{currentstroke}%
\pgfsetstrokeopacity{0.000000}%
\pgfsetdash{}{0pt}%
\pgfpathmoveto{\pgfqpoint{2.358324in}{1.060460in}}%
\pgfpathlineto{\pgfqpoint{2.350874in}{1.053026in}}%
\pgfpathlineto{\pgfqpoint{2.347601in}{1.052815in}}%
\pgfpathlineto{\pgfqpoint{2.340592in}{1.057330in}}%
\pgfpathlineto{\pgfqpoint{2.334876in}{1.050418in}}%
\pgfpathlineto{\pgfqpoint{2.333634in}{1.044388in}}%
\pgfpathlineto{\pgfqpoint{2.321129in}{1.054650in}}%
\pgfpathlineto{\pgfqpoint{2.322848in}{1.065287in}}%
\pgfpathlineto{\pgfqpoint{2.327579in}{1.069371in}}%
\pgfpathlineto{\pgfqpoint{2.316608in}{1.083386in}}%
\pgfpathlineto{\pgfqpoint{2.328598in}{1.094602in}}%
\pgfpathlineto{\pgfqpoint{2.342975in}{1.090758in}}%
\pgfpathlineto{\pgfqpoint{2.350883in}{1.086320in}}%
\pgfpathlineto{\pgfqpoint{2.358589in}{1.084899in}}%
\pgfpathlineto{\pgfqpoint{2.353799in}{1.080602in}}%
\pgfpathlineto{\pgfqpoint{2.350488in}{1.072384in}}%
\pgfpathlineto{\pgfqpoint{2.352145in}{1.068229in}}%
\pgfpathlineto{\pgfqpoint{2.358324in}{1.060460in}}%
\pgfpathclose%
\pgfusepath{fill}%
\end{pgfscope}%
\begin{pgfscope}%
\pgfpathrectangle{\pgfqpoint{0.100000in}{0.100000in}}{\pgfqpoint{3.007045in}{1.925000in}}%
\pgfusepath{clip}%
\pgfsetbuttcap%
\pgfsetmiterjoin%
\definecolor{currentfill}{rgb}{0.652995,0.860085,0.628897}%
\pgfsetfillcolor{currentfill}%
\pgfsetlinewidth{0.000000pt}%
\definecolor{currentstroke}{rgb}{0.000000,0.000000,0.000000}%
\pgfsetstrokecolor{currentstroke}%
\pgfsetstrokeopacity{0.000000}%
\pgfsetdash{}{0pt}%
\pgfpathmoveto{\pgfqpoint{1.835531in}{0.660856in}}%
\pgfpathlineto{\pgfqpoint{1.860218in}{0.661283in}}%
\pgfpathlineto{\pgfqpoint{1.860335in}{0.652602in}}%
\pgfpathlineto{\pgfqpoint{1.866139in}{0.652691in}}%
\pgfpathlineto{\pgfqpoint{1.866438in}{0.640996in}}%
\pgfpathlineto{\pgfqpoint{1.872143in}{0.641119in}}%
\pgfpathlineto{\pgfqpoint{1.872399in}{0.629573in}}%
\pgfpathlineto{\pgfqpoint{1.863641in}{0.629386in}}%
\pgfpathlineto{\pgfqpoint{1.860787in}{0.623525in}}%
\pgfpathlineto{\pgfqpoint{1.855514in}{0.623416in}}%
\pgfpathlineto{\pgfqpoint{1.854631in}{0.617364in}}%
\pgfpathlineto{\pgfqpoint{1.851931in}{0.616696in}}%
\pgfpathlineto{\pgfqpoint{1.846842in}{0.625571in}}%
\pgfpathlineto{\pgfqpoint{1.847519in}{0.631354in}}%
\pgfpathlineto{\pgfqpoint{1.843180in}{0.636011in}}%
\pgfpathlineto{\pgfqpoint{1.837976in}{0.644603in}}%
\pgfpathlineto{\pgfqpoint{1.839877in}{0.650274in}}%
\pgfpathlineto{\pgfqpoint{1.838775in}{0.656201in}}%
\pgfpathlineto{\pgfqpoint{1.835531in}{0.660856in}}%
\pgfpathclose%
\pgfusepath{fill}%
\end{pgfscope}%
\begin{pgfscope}%
\pgfpathrectangle{\pgfqpoint{0.100000in}{0.100000in}}{\pgfqpoint{3.007045in}{1.925000in}}%
\pgfusepath{clip}%
\pgfsetbuttcap%
\pgfsetmiterjoin%
\definecolor{currentfill}{rgb}{0.439216,0.760938,0.455056}%
\pgfsetfillcolor{currentfill}%
\pgfsetlinewidth{0.000000pt}%
\definecolor{currentstroke}{rgb}{0.000000,0.000000,0.000000}%
\pgfsetstrokecolor{currentstroke}%
\pgfsetstrokeopacity{0.000000}%
\pgfsetdash{}{0pt}%
\pgfpathmoveto{\pgfqpoint{1.460453in}{1.255605in}}%
\pgfpathlineto{\pgfqpoint{1.420426in}{1.258524in}}%
\pgfpathlineto{\pgfqpoint{1.420694in}{1.261942in}}%
\pgfpathlineto{\pgfqpoint{1.390844in}{1.264339in}}%
\pgfpathlineto{\pgfqpoint{1.392134in}{1.280968in}}%
\pgfpathlineto{\pgfqpoint{1.394532in}{1.283801in}}%
\pgfpathlineto{\pgfqpoint{1.395479in}{1.295197in}}%
\pgfpathlineto{\pgfqpoint{1.422895in}{1.292969in}}%
\pgfpathlineto{\pgfqpoint{1.423796in}{1.304378in}}%
\pgfpathlineto{\pgfqpoint{1.427276in}{1.304079in}}%
\pgfpathlineto{\pgfqpoint{1.455754in}{1.302023in}}%
\pgfpathlineto{\pgfqpoint{1.462537in}{1.301517in}}%
\pgfpathlineto{\pgfqpoint{1.460924in}{1.278668in}}%
\pgfpathlineto{\pgfqpoint{1.461934in}{1.278597in}}%
\pgfpathlineto{\pgfqpoint{1.460453in}{1.255605in}}%
\pgfpathclose%
\pgfusepath{fill}%
\end{pgfscope}%
\begin{pgfscope}%
\pgfpathrectangle{\pgfqpoint{0.100000in}{0.100000in}}{\pgfqpoint{3.007045in}{1.925000in}}%
\pgfusepath{clip}%
\pgfsetbuttcap%
\pgfsetmiterjoin%
\definecolor{currentfill}{rgb}{0.769873,0.909296,0.742699}%
\pgfsetfillcolor{currentfill}%
\pgfsetlinewidth{0.000000pt}%
\definecolor{currentstroke}{rgb}{0.000000,0.000000,0.000000}%
\pgfsetstrokecolor{currentstroke}%
\pgfsetstrokeopacity{0.000000}%
\pgfsetdash{}{0pt}%
\pgfpathmoveto{\pgfqpoint{2.191739in}{0.996087in}}%
\pgfpathlineto{\pgfqpoint{2.189732in}{1.023554in}}%
\pgfpathlineto{\pgfqpoint{2.200660in}{1.024227in}}%
\pgfpathlineto{\pgfqpoint{2.209941in}{1.021313in}}%
\pgfpathlineto{\pgfqpoint{2.212201in}{1.031659in}}%
\pgfpathlineto{\pgfqpoint{2.216520in}{1.033633in}}%
\pgfpathlineto{\pgfqpoint{2.223300in}{1.033829in}}%
\pgfpathlineto{\pgfqpoint{2.218727in}{1.043400in}}%
\pgfpathlineto{\pgfqpoint{2.234875in}{1.045920in}}%
\pgfpathlineto{\pgfqpoint{2.241309in}{1.038598in}}%
\pgfpathlineto{\pgfqpoint{2.241044in}{1.027825in}}%
\pgfpathlineto{\pgfqpoint{2.236965in}{1.019520in}}%
\pgfpathlineto{\pgfqpoint{2.247870in}{1.007837in}}%
\pgfpathlineto{\pgfqpoint{2.248936in}{1.000483in}}%
\pgfpathlineto{\pgfqpoint{2.230102in}{0.999971in}}%
\pgfpathlineto{\pgfqpoint{2.216487in}{0.998935in}}%
\pgfpathlineto{\pgfqpoint{2.191739in}{0.996087in}}%
\pgfpathclose%
\pgfusepath{fill}%
\end{pgfscope}%
\begin{pgfscope}%
\pgfpathrectangle{\pgfqpoint{0.100000in}{0.100000in}}{\pgfqpoint{3.007045in}{1.925000in}}%
\pgfusepath{clip}%
\pgfsetbuttcap%
\pgfsetmiterjoin%
\definecolor{currentfill}{rgb}{0.746498,0.899454,0.719938}%
\pgfsetfillcolor{currentfill}%
\pgfsetlinewidth{0.000000pt}%
\definecolor{currentstroke}{rgb}{0.000000,0.000000,0.000000}%
\pgfsetstrokecolor{currentstroke}%
\pgfsetstrokeopacity{0.000000}%
\pgfsetdash{}{0pt}%
\pgfpathmoveto{\pgfqpoint{0.829274in}{1.562729in}}%
\pgfpathlineto{\pgfqpoint{0.834605in}{1.565249in}}%
\pgfpathlineto{\pgfqpoint{0.836443in}{1.560436in}}%
\pgfpathlineto{\pgfqpoint{0.842609in}{1.559070in}}%
\pgfpathlineto{\pgfqpoint{0.849873in}{1.555251in}}%
\pgfpathlineto{\pgfqpoint{0.860657in}{1.551817in}}%
\pgfpathlineto{\pgfqpoint{0.861425in}{1.547079in}}%
\pgfpathlineto{\pgfqpoint{0.872520in}{1.537698in}}%
\pgfpathlineto{\pgfqpoint{0.878844in}{1.527019in}}%
\pgfpathlineto{\pgfqpoint{0.882414in}{1.527693in}}%
\pgfpathlineto{\pgfqpoint{0.884911in}{1.517908in}}%
\pgfpathlineto{\pgfqpoint{0.883755in}{1.512410in}}%
\pgfpathlineto{\pgfqpoint{0.896614in}{1.509734in}}%
\pgfpathlineto{\pgfqpoint{0.895520in}{1.504390in}}%
\pgfpathlineto{\pgfqpoint{0.912230in}{1.501011in}}%
\pgfpathlineto{\pgfqpoint{0.909954in}{1.489763in}}%
\pgfpathlineto{\pgfqpoint{0.898811in}{1.492021in}}%
\pgfpathlineto{\pgfqpoint{0.894533in}{1.469396in}}%
\pgfpathlineto{\pgfqpoint{0.896143in}{1.463222in}}%
\pgfpathlineto{\pgfqpoint{0.892728in}{1.460530in}}%
\pgfpathlineto{\pgfqpoint{0.887062in}{1.465893in}}%
\pgfpathlineto{\pgfqpoint{0.882125in}{1.465664in}}%
\pgfpathlineto{\pgfqpoint{0.884523in}{1.477298in}}%
\pgfpathlineto{\pgfqpoint{0.887326in}{1.476723in}}%
\pgfpathlineto{\pgfqpoint{0.891952in}{1.499400in}}%
\pgfpathlineto{\pgfqpoint{0.877767in}{1.502355in}}%
\pgfpathlineto{\pgfqpoint{0.873074in}{1.479709in}}%
\pgfpathlineto{\pgfqpoint{0.869518in}{1.474687in}}%
\pgfpathlineto{\pgfqpoint{0.861547in}{1.476417in}}%
\pgfpathlineto{\pgfqpoint{0.858463in}{1.461570in}}%
\pgfpathlineto{\pgfqpoint{0.851744in}{1.461596in}}%
\pgfpathlineto{\pgfqpoint{0.850305in}{1.455334in}}%
\pgfpathlineto{\pgfqpoint{0.840009in}{1.457550in}}%
\pgfpathlineto{\pgfqpoint{0.834301in}{1.430229in}}%
\pgfpathlineto{\pgfqpoint{0.808116in}{1.436493in}}%
\pgfpathlineto{\pgfqpoint{0.798116in}{1.438303in}}%
\pgfpathlineto{\pgfqpoint{0.811385in}{1.497272in}}%
\pgfpathlineto{\pgfqpoint{0.809080in}{1.497950in}}%
\pgfpathlineto{\pgfqpoint{0.819243in}{1.542361in}}%
\pgfpathlineto{\pgfqpoint{0.824049in}{1.547593in}}%
\pgfpathlineto{\pgfqpoint{0.823175in}{1.550976in}}%
\pgfpathlineto{\pgfqpoint{0.828248in}{1.558524in}}%
\pgfpathlineto{\pgfqpoint{0.829274in}{1.562729in}}%
\pgfpathclose%
\pgfusepath{fill}%
\end{pgfscope}%
\begin{pgfscope}%
\pgfpathrectangle{\pgfqpoint{0.100000in}{0.100000in}}{\pgfqpoint{3.007045in}{1.925000in}}%
\pgfusepath{clip}%
\pgfsetbuttcap%
\pgfsetmiterjoin%
\definecolor{currentfill}{rgb}{0.535179,0.806090,0.528750}%
\pgfsetfillcolor{currentfill}%
\pgfsetlinewidth{0.000000pt}%
\definecolor{currentstroke}{rgb}{0.000000,0.000000,0.000000}%
\pgfsetstrokecolor{currentstroke}%
\pgfsetstrokeopacity{0.000000}%
\pgfsetdash{}{0pt}%
\pgfpathmoveto{\pgfqpoint{0.429023in}{1.481634in}}%
\pgfpathlineto{\pgfqpoint{0.424324in}{1.471482in}}%
\pgfpathlineto{\pgfqpoint{0.420836in}{1.470570in}}%
\pgfpathlineto{\pgfqpoint{0.412778in}{1.458416in}}%
\pgfpathlineto{\pgfqpoint{0.407285in}{1.452726in}}%
\pgfpathlineto{\pgfqpoint{0.407911in}{1.446662in}}%
\pgfpathlineto{\pgfqpoint{0.399108in}{1.444237in}}%
\pgfpathlineto{\pgfqpoint{0.395246in}{1.441320in}}%
\pgfpathlineto{\pgfqpoint{0.390077in}{1.441086in}}%
\pgfpathlineto{\pgfqpoint{0.385422in}{1.437519in}}%
\pgfpathlineto{\pgfqpoint{0.384276in}{1.433228in}}%
\pgfpathlineto{\pgfqpoint{0.387316in}{1.430255in}}%
\pgfpathlineto{\pgfqpoint{0.385058in}{1.423808in}}%
\pgfpathlineto{\pgfqpoint{0.384397in}{1.412475in}}%
\pgfpathlineto{\pgfqpoint{0.355032in}{1.421828in}}%
\pgfpathlineto{\pgfqpoint{0.355543in}{1.423387in}}%
\pgfpathlineto{\pgfqpoint{0.332528in}{1.430977in}}%
\pgfpathlineto{\pgfqpoint{0.330534in}{1.438279in}}%
\pgfpathlineto{\pgfqpoint{0.327431in}{1.441471in}}%
\pgfpathlineto{\pgfqpoint{0.321640in}{1.452224in}}%
\pgfpathlineto{\pgfqpoint{0.323491in}{1.455583in}}%
\pgfpathlineto{\pgfqpoint{0.323252in}{1.464314in}}%
\pgfpathlineto{\pgfqpoint{0.332049in}{1.475470in}}%
\pgfpathlineto{\pgfqpoint{0.345804in}{1.489825in}}%
\pgfpathlineto{\pgfqpoint{0.349534in}{1.496541in}}%
\pgfpathlineto{\pgfqpoint{0.348214in}{1.499832in}}%
\pgfpathlineto{\pgfqpoint{0.355538in}{1.511740in}}%
\pgfpathlineto{\pgfqpoint{0.360884in}{1.523094in}}%
\pgfpathlineto{\pgfqpoint{0.374697in}{1.518495in}}%
\pgfpathlineto{\pgfqpoint{0.372999in}{1.513265in}}%
\pgfpathlineto{\pgfqpoint{0.378145in}{1.511673in}}%
\pgfpathlineto{\pgfqpoint{0.385689in}{1.509253in}}%
\pgfpathlineto{\pgfqpoint{0.386307in}{1.493716in}}%
\pgfpathlineto{\pgfqpoint{0.391169in}{1.489496in}}%
\pgfpathlineto{\pgfqpoint{0.391596in}{1.486250in}}%
\pgfpathlineto{\pgfqpoint{0.397986in}{1.484025in}}%
\pgfpathlineto{\pgfqpoint{0.399966in}{1.478464in}}%
\pgfpathlineto{\pgfqpoint{0.407088in}{1.477930in}}%
\pgfpathlineto{\pgfqpoint{0.405051in}{1.484139in}}%
\pgfpathlineto{\pgfqpoint{0.408114in}{1.487902in}}%
\pgfpathlineto{\pgfqpoint{0.414605in}{1.487528in}}%
\pgfpathlineto{\pgfqpoint{0.428385in}{1.494011in}}%
\pgfpathlineto{\pgfqpoint{0.431912in}{1.491802in}}%
\pgfpathlineto{\pgfqpoint{0.429023in}{1.481634in}}%
\pgfpathclose%
\pgfusepath{fill}%
\end{pgfscope}%
\begin{pgfscope}%
\pgfpathrectangle{\pgfqpoint{0.100000in}{0.100000in}}{\pgfqpoint{3.007045in}{1.925000in}}%
\pgfusepath{clip}%
\pgfsetbuttcap%
\pgfsetmiterjoin%
\definecolor{currentfill}{rgb}{0.535179,0.806090,0.528750}%
\pgfsetfillcolor{currentfill}%
\pgfsetlinewidth{0.000000pt}%
\definecolor{currentstroke}{rgb}{0.000000,0.000000,0.000000}%
\pgfsetstrokecolor{currentstroke}%
\pgfsetstrokeopacity{0.000000}%
\pgfsetdash{}{0pt}%
\pgfpathmoveto{\pgfqpoint{1.889132in}{1.155695in}}%
\pgfpathlineto{\pgfqpoint{1.884869in}{1.157037in}}%
\pgfpathlineto{\pgfqpoint{1.882047in}{1.149212in}}%
\pgfpathlineto{\pgfqpoint{1.877817in}{1.149309in}}%
\pgfpathlineto{\pgfqpoint{1.875772in}{1.154964in}}%
\pgfpathlineto{\pgfqpoint{1.870359in}{1.156300in}}%
\pgfpathlineto{\pgfqpoint{1.867949in}{1.161555in}}%
\pgfpathlineto{\pgfqpoint{1.862920in}{1.162275in}}%
\pgfpathlineto{\pgfqpoint{1.859930in}{1.165889in}}%
\pgfpathlineto{\pgfqpoint{1.860786in}{1.170457in}}%
\pgfpathlineto{\pgfqpoint{1.859241in}{1.178993in}}%
\pgfpathlineto{\pgfqpoint{1.855189in}{1.180142in}}%
\pgfpathlineto{\pgfqpoint{1.855015in}{1.197513in}}%
\pgfpathlineto{\pgfqpoint{1.854945in}{1.201835in}}%
\pgfpathlineto{\pgfqpoint{1.880439in}{1.202637in}}%
\pgfpathlineto{\pgfqpoint{1.880879in}{1.202651in}}%
\pgfpathlineto{\pgfqpoint{1.880876in}{1.180425in}}%
\pgfpathlineto{\pgfqpoint{1.889421in}{1.180509in}}%
\pgfpathlineto{\pgfqpoint{1.889132in}{1.155695in}}%
\pgfpathclose%
\pgfusepath{fill}%
\end{pgfscope}%
\begin{pgfscope}%
\pgfpathrectangle{\pgfqpoint{0.100000in}{0.100000in}}{\pgfqpoint{3.007045in}{1.925000in}}%
\pgfusepath{clip}%
\pgfsetbuttcap%
\pgfsetmiterjoin%
\definecolor{currentfill}{rgb}{0.727797,0.891580,0.701730}%
\pgfsetfillcolor{currentfill}%
\pgfsetlinewidth{0.000000pt}%
\definecolor{currentstroke}{rgb}{0.000000,0.000000,0.000000}%
\pgfsetstrokecolor{currentstroke}%
\pgfsetstrokeopacity{0.000000}%
\pgfsetdash{}{0pt}%
\pgfpathmoveto{\pgfqpoint{1.564310in}{0.846924in}}%
\pgfpathlineto{\pgfqpoint{1.564075in}{0.841182in}}%
\pgfpathlineto{\pgfqpoint{1.573014in}{0.840820in}}%
\pgfpathlineto{\pgfqpoint{1.572553in}{0.829339in}}%
\pgfpathlineto{\pgfqpoint{1.575401in}{0.829239in}}%
\pgfpathlineto{\pgfqpoint{1.574901in}{0.817578in}}%
\pgfpathlineto{\pgfqpoint{1.570469in}{0.816172in}}%
\pgfpathlineto{\pgfqpoint{1.560430in}{0.818079in}}%
\pgfpathlineto{\pgfqpoint{1.556351in}{0.822025in}}%
\pgfpathlineto{\pgfqpoint{1.543414in}{0.822716in}}%
\pgfpathlineto{\pgfqpoint{1.542619in}{0.830859in}}%
\pgfpathlineto{\pgfqpoint{1.542518in}{0.841248in}}%
\pgfpathlineto{\pgfqpoint{1.545976in}{0.847321in}}%
\pgfpathlineto{\pgfqpoint{1.564310in}{0.846924in}}%
\pgfpathclose%
\pgfusepath{fill}%
\end{pgfscope}%
\begin{pgfscope}%
\pgfpathrectangle{\pgfqpoint{0.100000in}{0.100000in}}{\pgfqpoint{3.007045in}{1.925000in}}%
\pgfusepath{clip}%
\pgfsetbuttcap%
\pgfsetmiterjoin%
\definecolor{currentfill}{rgb}{0.513033,0.795755,0.510542}%
\pgfsetfillcolor{currentfill}%
\pgfsetlinewidth{0.000000pt}%
\definecolor{currentstroke}{rgb}{0.000000,0.000000,0.000000}%
\pgfsetstrokecolor{currentstroke}%
\pgfsetstrokeopacity{0.000000}%
\pgfsetdash{}{0pt}%
\pgfpathmoveto{\pgfqpoint{1.721870in}{1.430045in}}%
\pgfpathlineto{\pgfqpoint{1.702699in}{1.430317in}}%
\pgfpathlineto{\pgfqpoint{1.703058in}{1.453389in}}%
\pgfpathlineto{\pgfqpoint{1.703851in}{1.499277in}}%
\pgfpathlineto{\pgfqpoint{1.703956in}{1.505078in}}%
\pgfpathlineto{\pgfqpoint{1.743883in}{1.504541in}}%
\pgfpathlineto{\pgfqpoint{1.744270in}{1.498734in}}%
\pgfpathlineto{\pgfqpoint{1.744188in}{1.475811in}}%
\pgfpathlineto{\pgfqpoint{1.721890in}{1.476069in}}%
\pgfpathlineto{\pgfqpoint{1.722176in}{1.453096in}}%
\pgfpathlineto{\pgfqpoint{1.721870in}{1.430045in}}%
\pgfpathclose%
\pgfusepath{fill}%
\end{pgfscope}%
\begin{pgfscope}%
\pgfpathrectangle{\pgfqpoint{0.100000in}{0.100000in}}{\pgfqpoint{3.007045in}{1.925000in}}%
\pgfusepath{clip}%
\pgfsetbuttcap%
\pgfsetmiterjoin%
\definecolor{currentfill}{rgb}{0.727797,0.891580,0.701730}%
\pgfsetfillcolor{currentfill}%
\pgfsetlinewidth{0.000000pt}%
\definecolor{currentstroke}{rgb}{0.000000,0.000000,0.000000}%
\pgfsetstrokecolor{currentstroke}%
\pgfsetstrokeopacity{0.000000}%
\pgfsetdash{}{0pt}%
\pgfpathmoveto{\pgfqpoint{1.234424in}{1.392114in}}%
\pgfpathlineto{\pgfqpoint{1.231571in}{1.369429in}}%
\pgfpathlineto{\pgfqpoint{1.226347in}{1.324123in}}%
\pgfpathlineto{\pgfqpoint{1.213860in}{1.325677in}}%
\pgfpathlineto{\pgfqpoint{1.210701in}{1.299766in}}%
\pgfpathlineto{\pgfqpoint{1.184259in}{1.303392in}}%
\pgfpathlineto{\pgfqpoint{1.131992in}{1.310519in}}%
\pgfpathlineto{\pgfqpoint{1.135759in}{1.336378in}}%
\pgfpathlineto{\pgfqpoint{1.137581in}{1.353507in}}%
\pgfpathlineto{\pgfqpoint{1.158048in}{1.350480in}}%
\pgfpathlineto{\pgfqpoint{1.162230in}{1.378562in}}%
\pgfpathlineto{\pgfqpoint{1.162889in}{1.390045in}}%
\pgfpathlineto{\pgfqpoint{1.164443in}{1.401325in}}%
\pgfpathlineto{\pgfqpoint{1.207752in}{1.395280in}}%
\pgfpathlineto{\pgfqpoint{1.234424in}{1.392114in}}%
\pgfpathclose%
\pgfusepath{fill}%
\end{pgfscope}%
\begin{pgfscope}%
\pgfpathrectangle{\pgfqpoint{0.100000in}{0.100000in}}{\pgfqpoint{3.007045in}{1.925000in}}%
\pgfusepath{clip}%
\pgfsetbuttcap%
\pgfsetmiterjoin%
\definecolor{currentfill}{rgb}{0.737147,0.895517,0.710834}%
\pgfsetfillcolor{currentfill}%
\pgfsetlinewidth{0.000000pt}%
\definecolor{currentstroke}{rgb}{0.000000,0.000000,0.000000}%
\pgfsetstrokecolor{currentstroke}%
\pgfsetstrokeopacity{0.000000}%
\pgfsetdash{}{0pt}%
\pgfpathmoveto{\pgfqpoint{1.652796in}{1.200602in}}%
\pgfpathlineto{\pgfqpoint{1.653431in}{1.223538in}}%
\pgfpathlineto{\pgfqpoint{1.654636in}{1.269418in}}%
\pgfpathlineto{\pgfqpoint{1.677475in}{1.268823in}}%
\pgfpathlineto{\pgfqpoint{1.699583in}{1.268350in}}%
\pgfpathlineto{\pgfqpoint{1.699296in}{1.251149in}}%
\pgfpathlineto{\pgfqpoint{1.698475in}{1.199567in}}%
\pgfpathlineto{\pgfqpoint{1.652796in}{1.200602in}}%
\pgfpathclose%
\pgfusepath{fill}%
\end{pgfscope}%
\begin{pgfscope}%
\pgfpathrectangle{\pgfqpoint{0.100000in}{0.100000in}}{\pgfqpoint{3.007045in}{1.925000in}}%
\pgfusepath{clip}%
\pgfsetbuttcap%
\pgfsetmiterjoin%
\definecolor{currentfill}{rgb}{0.618224,0.844844,0.597032}%
\pgfsetfillcolor{currentfill}%
\pgfsetlinewidth{0.000000pt}%
\definecolor{currentstroke}{rgb}{0.000000,0.000000,0.000000}%
\pgfsetstrokecolor{currentstroke}%
\pgfsetstrokeopacity{0.000000}%
\pgfsetdash{}{0pt}%
\pgfpathmoveto{\pgfqpoint{1.650048in}{1.679481in}}%
\pgfpathlineto{\pgfqpoint{1.638529in}{1.679874in}}%
\pgfpathlineto{\pgfqpoint{1.637567in}{1.685608in}}%
\pgfpathlineto{\pgfqpoint{1.638404in}{1.708645in}}%
\pgfpathlineto{\pgfqpoint{1.642987in}{1.708503in}}%
\pgfpathlineto{\pgfqpoint{1.643745in}{1.731622in}}%
\pgfpathlineto{\pgfqpoint{1.643166in}{1.743370in}}%
\pgfpathlineto{\pgfqpoint{1.677065in}{1.742344in}}%
\pgfpathlineto{\pgfqpoint{1.676789in}{1.740990in}}%
\pgfpathlineto{\pgfqpoint{1.705532in}{1.740473in}}%
\pgfpathlineto{\pgfqpoint{1.706015in}{1.718669in}}%
\pgfpathlineto{\pgfqpoint{1.711793in}{1.718564in}}%
\pgfpathlineto{\pgfqpoint{1.711752in}{1.712779in}}%
\pgfpathlineto{\pgfqpoint{1.734818in}{1.712519in}}%
\pgfpathlineto{\pgfqpoint{1.734882in}{1.718354in}}%
\pgfpathlineto{\pgfqpoint{1.740597in}{1.718328in}}%
\pgfpathlineto{\pgfqpoint{1.740624in}{1.724147in}}%
\pgfpathlineto{\pgfqpoint{1.746314in}{1.723984in}}%
\pgfpathlineto{\pgfqpoint{1.746352in}{1.706621in}}%
\pgfpathlineto{\pgfqpoint{1.747387in}{1.695072in}}%
\pgfpathlineto{\pgfqpoint{1.724288in}{1.695202in}}%
\pgfpathlineto{\pgfqpoint{1.689129in}{1.695752in}}%
\pgfpathlineto{\pgfqpoint{1.689494in}{1.678334in}}%
\pgfpathlineto{\pgfqpoint{1.650048in}{1.679481in}}%
\pgfpathclose%
\pgfusepath{fill}%
\end{pgfscope}%
\begin{pgfscope}%
\pgfpathrectangle{\pgfqpoint{0.100000in}{0.100000in}}{\pgfqpoint{3.007045in}{1.925000in}}%
\pgfusepath{clip}%
\pgfsetbuttcap%
\pgfsetmiterjoin%
\definecolor{currentfill}{rgb}{0.496424,0.788005,0.496886}%
\pgfsetfillcolor{currentfill}%
\pgfsetlinewidth{0.000000pt}%
\definecolor{currentstroke}{rgb}{0.000000,0.000000,0.000000}%
\pgfsetstrokecolor{currentstroke}%
\pgfsetstrokeopacity{0.000000}%
\pgfsetdash{}{0pt}%
\pgfpathmoveto{\pgfqpoint{1.767169in}{1.475679in}}%
\pgfpathlineto{\pgfqpoint{1.744188in}{1.475811in}}%
\pgfpathlineto{\pgfqpoint{1.744270in}{1.498734in}}%
\pgfpathlineto{\pgfqpoint{1.755330in}{1.498635in}}%
\pgfpathlineto{\pgfqpoint{1.755381in}{1.509019in}}%
\pgfpathlineto{\pgfqpoint{1.749628in}{1.512662in}}%
\pgfpathlineto{\pgfqpoint{1.749650in}{1.521778in}}%
\pgfpathlineto{\pgfqpoint{1.783670in}{1.521817in}}%
\pgfpathlineto{\pgfqpoint{1.795889in}{1.521880in}}%
\pgfpathlineto{\pgfqpoint{1.795936in}{1.510293in}}%
\pgfpathlineto{\pgfqpoint{1.789742in}{1.510236in}}%
\pgfpathlineto{\pgfqpoint{1.790073in}{1.492987in}}%
\pgfpathlineto{\pgfqpoint{1.782199in}{1.492987in}}%
\pgfpathlineto{\pgfqpoint{1.778660in}{1.495730in}}%
\pgfpathlineto{\pgfqpoint{1.778661in}{1.481448in}}%
\pgfpathlineto{\pgfqpoint{1.767240in}{1.481415in}}%
\pgfpathlineto{\pgfqpoint{1.767169in}{1.475679in}}%
\pgfpathclose%
\pgfusepath{fill}%
\end{pgfscope}%
\begin{pgfscope}%
\pgfpathrectangle{\pgfqpoint{0.100000in}{0.100000in}}{\pgfqpoint{3.007045in}{1.925000in}}%
\pgfusepath{clip}%
\pgfsetbuttcap%
\pgfsetmiterjoin%
\definecolor{currentfill}{rgb}{0.237832,0.652380,0.351050}%
\pgfsetfillcolor{currentfill}%
\pgfsetlinewidth{0.000000pt}%
\definecolor{currentstroke}{rgb}{0.000000,0.000000,0.000000}%
\pgfsetstrokecolor{currentstroke}%
\pgfsetstrokeopacity{0.000000}%
\pgfsetdash{}{0pt}%
\pgfpathmoveto{\pgfqpoint{1.948847in}{1.744739in}}%
\pgfpathlineto{\pgfqpoint{1.956874in}{1.748840in}}%
\pgfpathlineto{\pgfqpoint{1.961605in}{1.738804in}}%
\pgfpathlineto{\pgfqpoint{1.978390in}{1.739623in}}%
\pgfpathlineto{\pgfqpoint{1.988931in}{1.741585in}}%
\pgfpathlineto{\pgfqpoint{1.994032in}{1.740034in}}%
\pgfpathlineto{\pgfqpoint{1.995567in}{1.736342in}}%
\pgfpathlineto{\pgfqpoint{2.001096in}{1.733792in}}%
\pgfpathlineto{\pgfqpoint{2.005613in}{1.736504in}}%
\pgfpathlineto{\pgfqpoint{2.012380in}{1.736151in}}%
\pgfpathlineto{\pgfqpoint{2.011685in}{1.732318in}}%
\pgfpathlineto{\pgfqpoint{1.997324in}{1.723369in}}%
\pgfpathlineto{\pgfqpoint{1.987813in}{1.719063in}}%
\pgfpathlineto{\pgfqpoint{1.972465in}{1.713418in}}%
\pgfpathlineto{\pgfqpoint{1.963898in}{1.707804in}}%
\pgfpathlineto{\pgfqpoint{1.951342in}{1.696561in}}%
\pgfpathlineto{\pgfqpoint{1.950215in}{1.724734in}}%
\pgfpathlineto{\pgfqpoint{1.948847in}{1.744739in}}%
\pgfpathclose%
\pgfusepath{fill}%
\end{pgfscope}%
\begin{pgfscope}%
\pgfpathrectangle{\pgfqpoint{0.100000in}{0.100000in}}{\pgfqpoint{3.007045in}{1.925000in}}%
\pgfusepath{clip}%
\pgfsetbuttcap%
\pgfsetmiterjoin%
\definecolor{currentfill}{rgb}{0.634295,0.852211,0.610688}%
\pgfsetfillcolor{currentfill}%
\pgfsetlinewidth{0.000000pt}%
\definecolor{currentstroke}{rgb}{0.000000,0.000000,0.000000}%
\pgfsetstrokecolor{currentstroke}%
\pgfsetstrokeopacity{0.000000}%
\pgfsetdash{}{0pt}%
\pgfpathmoveto{\pgfqpoint{2.385587in}{1.155405in}}%
\pgfpathlineto{\pgfqpoint{2.380064in}{1.147699in}}%
\pgfpathlineto{\pgfqpoint{2.374431in}{1.146809in}}%
\pgfpathlineto{\pgfqpoint{2.370900in}{1.172765in}}%
\pgfpathlineto{\pgfqpoint{2.364743in}{1.174580in}}%
\pgfpathlineto{\pgfqpoint{2.365382in}{1.186478in}}%
\pgfpathlineto{\pgfqpoint{2.363339in}{1.187311in}}%
\pgfpathlineto{\pgfqpoint{2.362689in}{1.195710in}}%
\pgfpathlineto{\pgfqpoint{2.366865in}{1.205475in}}%
\pgfpathlineto{\pgfqpoint{2.381162in}{1.206723in}}%
\pgfpathlineto{\pgfqpoint{2.393794in}{1.205858in}}%
\pgfpathlineto{\pgfqpoint{2.394275in}{1.199273in}}%
\pgfpathlineto{\pgfqpoint{2.405917in}{1.200954in}}%
\pgfpathlineto{\pgfqpoint{2.405836in}{1.202691in}}%
\pgfpathlineto{\pgfqpoint{2.417290in}{1.203518in}}%
\pgfpathlineto{\pgfqpoint{2.419602in}{1.197771in}}%
\pgfpathlineto{\pgfqpoint{2.418057in}{1.191742in}}%
\pgfpathlineto{\pgfqpoint{2.418866in}{1.180027in}}%
\pgfpathlineto{\pgfqpoint{2.413102in}{1.179720in}}%
\pgfpathlineto{\pgfqpoint{2.413821in}{1.167145in}}%
\pgfpathlineto{\pgfqpoint{2.407704in}{1.166293in}}%
\pgfpathlineto{\pgfqpoint{2.403876in}{1.166074in}}%
\pgfpathlineto{\pgfqpoint{2.404127in}{1.159387in}}%
\pgfpathlineto{\pgfqpoint{2.401519in}{1.149007in}}%
\pgfpathlineto{\pgfqpoint{2.395723in}{1.148881in}}%
\pgfpathlineto{\pgfqpoint{2.392277in}{1.158355in}}%
\pgfpathlineto{\pgfqpoint{2.385587in}{1.155405in}}%
\pgfpathclose%
\pgfusepath{fill}%
\end{pgfscope}%
\begin{pgfscope}%
\pgfpathrectangle{\pgfqpoint{0.100000in}{0.100000in}}{\pgfqpoint{3.007045in}{1.925000in}}%
\pgfusepath{clip}%
\pgfsetbuttcap%
\pgfsetmiterjoin%
\definecolor{currentfill}{rgb}{0.501961,0.790588,0.501438}%
\pgfsetfillcolor{currentfill}%
\pgfsetlinewidth{0.000000pt}%
\definecolor{currentstroke}{rgb}{0.000000,0.000000,0.000000}%
\pgfsetstrokecolor{currentstroke}%
\pgfsetstrokeopacity{0.000000}%
\pgfsetdash{}{0pt}%
\pgfpathmoveto{\pgfqpoint{2.018995in}{1.370976in}}%
\pgfpathlineto{\pgfqpoint{2.020468in}{1.350608in}}%
\pgfpathlineto{\pgfqpoint{2.001078in}{1.349385in}}%
\pgfpathlineto{\pgfqpoint{1.997258in}{1.351410in}}%
\pgfpathlineto{\pgfqpoint{1.995792in}{1.357804in}}%
\pgfpathlineto{\pgfqpoint{1.992817in}{1.361377in}}%
\pgfpathlineto{\pgfqpoint{1.983506in}{1.360972in}}%
\pgfpathlineto{\pgfqpoint{1.983756in}{1.355205in}}%
\pgfpathlineto{\pgfqpoint{1.972422in}{1.354713in}}%
\pgfpathlineto{\pgfqpoint{1.938345in}{1.353373in}}%
\pgfpathlineto{\pgfqpoint{1.937355in}{1.364751in}}%
\pgfpathlineto{\pgfqpoint{1.936000in}{1.405159in}}%
\pgfpathlineto{\pgfqpoint{1.934699in}{1.432796in}}%
\pgfpathlineto{\pgfqpoint{1.953526in}{1.433541in}}%
\pgfpathlineto{\pgfqpoint{1.954322in}{1.428443in}}%
\pgfpathlineto{\pgfqpoint{1.954703in}{1.423577in}}%
\pgfpathlineto{\pgfqpoint{1.961879in}{1.417685in}}%
\pgfpathlineto{\pgfqpoint{1.956575in}{1.409480in}}%
\pgfpathlineto{\pgfqpoint{1.958754in}{1.394373in}}%
\pgfpathlineto{\pgfqpoint{1.960884in}{1.393037in}}%
\pgfpathlineto{\pgfqpoint{1.963562in}{1.383583in}}%
\pgfpathlineto{\pgfqpoint{1.969155in}{1.380236in}}%
\pgfpathlineto{\pgfqpoint{1.980017in}{1.377888in}}%
\pgfpathlineto{\pgfqpoint{1.984280in}{1.369339in}}%
\pgfpathlineto{\pgfqpoint{2.018995in}{1.370976in}}%
\pgfpathclose%
\pgfusepath{fill}%
\end{pgfscope}%
\begin{pgfscope}%
\pgfpathrectangle{\pgfqpoint{0.100000in}{0.100000in}}{\pgfqpoint{3.007045in}{1.925000in}}%
\pgfusepath{clip}%
\pgfsetbuttcap%
\pgfsetmiterjoin%
\definecolor{currentfill}{rgb}{0.551788,0.813841,0.542407}%
\pgfsetfillcolor{currentfill}%
\pgfsetlinewidth{0.000000pt}%
\definecolor{currentstroke}{rgb}{0.000000,0.000000,0.000000}%
\pgfsetstrokecolor{currentstroke}%
\pgfsetstrokeopacity{0.000000}%
\pgfsetdash{}{0pt}%
\pgfpathmoveto{\pgfqpoint{2.604673in}{1.078358in}}%
\pgfpathlineto{\pgfqpoint{2.599528in}{1.076992in}}%
\pgfpathlineto{\pgfqpoint{2.596091in}{1.081050in}}%
\pgfpathlineto{\pgfqpoint{2.591032in}{1.080451in}}%
\pgfpathlineto{\pgfqpoint{2.584592in}{1.075050in}}%
\pgfpathlineto{\pgfqpoint{2.578887in}{1.072582in}}%
\pgfpathlineto{\pgfqpoint{2.574988in}{1.077569in}}%
\pgfpathlineto{\pgfqpoint{2.564592in}{1.081776in}}%
\pgfpathlineto{\pgfqpoint{2.562917in}{1.088318in}}%
\pgfpathlineto{\pgfqpoint{2.564677in}{1.094671in}}%
\pgfpathlineto{\pgfqpoint{2.570652in}{1.100774in}}%
\pgfpathlineto{\pgfqpoint{2.573073in}{1.099388in}}%
\pgfpathlineto{\pgfqpoint{2.578351in}{1.105297in}}%
\pgfpathlineto{\pgfqpoint{2.579736in}{1.109975in}}%
\pgfpathlineto{\pgfqpoint{2.584797in}{1.114957in}}%
\pgfpathlineto{\pgfqpoint{2.587434in}{1.124379in}}%
\pgfpathlineto{\pgfqpoint{2.593692in}{1.126779in}}%
\pgfpathlineto{\pgfqpoint{2.598994in}{1.120903in}}%
\pgfpathlineto{\pgfqpoint{2.602817in}{1.120583in}}%
\pgfpathlineto{\pgfqpoint{2.609095in}{1.111445in}}%
\pgfpathlineto{\pgfqpoint{2.612770in}{1.112672in}}%
\pgfpathlineto{\pgfqpoint{2.621502in}{1.105990in}}%
\pgfpathlineto{\pgfqpoint{2.626460in}{1.104802in}}%
\pgfpathlineto{\pgfqpoint{2.623729in}{1.094309in}}%
\pgfpathlineto{\pgfqpoint{2.616873in}{1.090165in}}%
\pgfpathlineto{\pgfqpoint{2.614846in}{1.077543in}}%
\pgfpathlineto{\pgfqpoint{2.604673in}{1.078358in}}%
\pgfpathclose%
\pgfusepath{fill}%
\end{pgfscope}%
\begin{pgfscope}%
\pgfpathrectangle{\pgfqpoint{0.100000in}{0.100000in}}{\pgfqpoint{3.007045in}{1.925000in}}%
\pgfusepath{clip}%
\pgfsetbuttcap%
\pgfsetmiterjoin%
\definecolor{currentfill}{rgb}{0.507497,0.793172,0.505990}%
\pgfsetfillcolor{currentfill}%
\pgfsetlinewidth{0.000000pt}%
\definecolor{currentstroke}{rgb}{0.000000,0.000000,0.000000}%
\pgfsetstrokecolor{currentstroke}%
\pgfsetstrokeopacity{0.000000}%
\pgfsetdash{}{0pt}%
\pgfpathmoveto{\pgfqpoint{0.530203in}{1.618670in}}%
\pgfpathlineto{\pgfqpoint{0.573938in}{1.605496in}}%
\pgfpathlineto{\pgfqpoint{0.596099in}{1.599150in}}%
\pgfpathlineto{\pgfqpoint{0.588372in}{1.571730in}}%
\pgfpathlineto{\pgfqpoint{0.586719in}{1.572193in}}%
\pgfpathlineto{\pgfqpoint{0.578445in}{1.544847in}}%
\pgfpathlineto{\pgfqpoint{0.605512in}{1.537505in}}%
\pgfpathlineto{\pgfqpoint{0.592563in}{1.489427in}}%
\pgfpathlineto{\pgfqpoint{0.562393in}{1.497838in}}%
\pgfpathlineto{\pgfqpoint{0.520911in}{1.509711in}}%
\pgfpathlineto{\pgfqpoint{0.534635in}{1.557261in}}%
\pgfpathlineto{\pgfqpoint{0.513001in}{1.563845in}}%
\pgfpathlineto{\pgfqpoint{0.524578in}{1.602426in}}%
\pgfpathlineto{\pgfqpoint{0.530203in}{1.618670in}}%
\pgfpathclose%
\pgfusepath{fill}%
\end{pgfscope}%
\begin{pgfscope}%
\pgfpathrectangle{\pgfqpoint{0.100000in}{0.100000in}}{\pgfqpoint{3.007045in}{1.925000in}}%
\pgfusepath{clip}%
\pgfsetbuttcap%
\pgfsetmiterjoin%
\definecolor{currentfill}{rgb}{0.797924,0.920738,0.771642}%
\pgfsetfillcolor{currentfill}%
\pgfsetlinewidth{0.000000pt}%
\definecolor{currentstroke}{rgb}{0.000000,0.000000,0.000000}%
\pgfsetstrokecolor{currentstroke}%
\pgfsetstrokeopacity{0.000000}%
\pgfsetdash{}{0pt}%
\pgfpathmoveto{\pgfqpoint{1.799463in}{0.919527in}}%
\pgfpathlineto{\pgfqpoint{1.795605in}{0.946063in}}%
\pgfpathlineto{\pgfqpoint{1.771725in}{0.946010in}}%
\pgfpathlineto{\pgfqpoint{1.772350in}{0.979565in}}%
\pgfpathlineto{\pgfqpoint{1.792522in}{0.979422in}}%
\pgfpathlineto{\pgfqpoint{1.792483in}{0.985928in}}%
\pgfpathlineto{\pgfqpoint{1.821509in}{0.984906in}}%
\pgfpathlineto{\pgfqpoint{1.821240in}{0.968478in}}%
\pgfpathlineto{\pgfqpoint{1.832387in}{0.968604in}}%
\pgfpathlineto{\pgfqpoint{1.832421in}{0.955990in}}%
\pgfpathlineto{\pgfqpoint{1.843659in}{0.955855in}}%
\pgfpathlineto{\pgfqpoint{1.845759in}{0.949170in}}%
\pgfpathlineto{\pgfqpoint{1.849514in}{0.944275in}}%
\pgfpathlineto{\pgfqpoint{1.855176in}{0.941374in}}%
\pgfpathlineto{\pgfqpoint{1.855032in}{0.933900in}}%
\pgfpathlineto{\pgfqpoint{1.851232in}{0.928170in}}%
\pgfpathlineto{\pgfqpoint{1.851596in}{0.920385in}}%
\pgfpathlineto{\pgfqpoint{1.830635in}{0.919840in}}%
\pgfpathlineto{\pgfqpoint{1.819183in}{0.919101in}}%
\pgfpathlineto{\pgfqpoint{1.799463in}{0.919527in}}%
\pgfpathclose%
\pgfusepath{fill}%
\end{pgfscope}%
\begin{pgfscope}%
\pgfpathrectangle{\pgfqpoint{0.100000in}{0.100000in}}{\pgfqpoint{3.007045in}{1.925000in}}%
\pgfusepath{clip}%
\pgfsetbuttcap%
\pgfsetmiterjoin%
\definecolor{currentfill}{rgb}{0.709097,0.883706,0.683522}%
\pgfsetfillcolor{currentfill}%
\pgfsetlinewidth{0.000000pt}%
\definecolor{currentstroke}{rgb}{0.000000,0.000000,0.000000}%
\pgfsetstrokecolor{currentstroke}%
\pgfsetstrokeopacity{0.000000}%
\pgfsetdash{}{0pt}%
\pgfpathmoveto{\pgfqpoint{2.584761in}{0.832670in}}%
\pgfpathlineto{\pgfqpoint{2.584582in}{0.822960in}}%
\pgfpathlineto{\pgfqpoint{2.575563in}{0.818701in}}%
\pgfpathlineto{\pgfqpoint{2.570619in}{0.823049in}}%
\pgfpathlineto{\pgfqpoint{2.556125in}{0.810180in}}%
\pgfpathlineto{\pgfqpoint{2.545370in}{0.817990in}}%
\pgfpathlineto{\pgfqpoint{2.532363in}{0.822168in}}%
\pgfpathlineto{\pgfqpoint{2.525782in}{0.825781in}}%
\pgfpathlineto{\pgfqpoint{2.521533in}{0.825837in}}%
\pgfpathlineto{\pgfqpoint{2.530038in}{0.837997in}}%
\pgfpathlineto{\pgfqpoint{2.516115in}{0.841918in}}%
\pgfpathlineto{\pgfqpoint{2.506929in}{0.849528in}}%
\pgfpathlineto{\pgfqpoint{2.503173in}{0.844767in}}%
\pgfpathlineto{\pgfqpoint{2.492621in}{0.846691in}}%
\pgfpathlineto{\pgfqpoint{2.488445in}{0.853368in}}%
\pgfpathlineto{\pgfqpoint{2.482424in}{0.851583in}}%
\pgfpathlineto{\pgfqpoint{2.486659in}{0.864555in}}%
\pgfpathlineto{\pgfqpoint{2.484386in}{0.869099in}}%
\pgfpathlineto{\pgfqpoint{2.485317in}{0.874629in}}%
\pgfpathlineto{\pgfqpoint{2.490797in}{0.879764in}}%
\pgfpathlineto{\pgfqpoint{2.496666in}{0.891862in}}%
\pgfpathlineto{\pgfqpoint{2.502098in}{0.888878in}}%
\pgfpathlineto{\pgfqpoint{2.508933in}{0.891663in}}%
\pgfpathlineto{\pgfqpoint{2.508172in}{0.896146in}}%
\pgfpathlineto{\pgfqpoint{2.537783in}{0.898762in}}%
\pgfpathlineto{\pgfqpoint{2.538562in}{0.893175in}}%
\pgfpathlineto{\pgfqpoint{2.545276in}{0.888835in}}%
\pgfpathlineto{\pgfqpoint{2.541892in}{0.875274in}}%
\pgfpathlineto{\pgfqpoint{2.551190in}{0.869149in}}%
\pgfpathlineto{\pgfqpoint{2.555420in}{0.871742in}}%
\pgfpathlineto{\pgfqpoint{2.557919in}{0.868208in}}%
\pgfpathlineto{\pgfqpoint{2.558099in}{0.855702in}}%
\pgfpathlineto{\pgfqpoint{2.560497in}{0.852920in}}%
\pgfpathlineto{\pgfqpoint{2.559661in}{0.847879in}}%
\pgfpathlineto{\pgfqpoint{2.566054in}{0.842571in}}%
\pgfpathlineto{\pgfqpoint{2.571176in}{0.834427in}}%
\pgfpathlineto{\pgfqpoint{2.577826in}{0.830633in}}%
\pgfpathlineto{\pgfqpoint{2.584761in}{0.832670in}}%
\pgfpathclose%
\pgfusepath{fill}%
\end{pgfscope}%
\begin{pgfscope}%
\pgfpathrectangle{\pgfqpoint{0.100000in}{0.100000in}}{\pgfqpoint{3.007045in}{1.925000in}}%
\pgfusepath{clip}%
\pgfsetbuttcap%
\pgfsetmiterjoin%
\definecolor{currentfill}{rgb}{0.178777,0.589389,0.303806}%
\pgfsetfillcolor{currentfill}%
\pgfsetlinewidth{0.000000pt}%
\definecolor{currentstroke}{rgb}{0.000000,0.000000,0.000000}%
\pgfsetstrokecolor{currentstroke}%
\pgfsetstrokeopacity{0.000000}%
\pgfsetdash{}{0pt}%
\pgfpathmoveto{\pgfqpoint{2.413670in}{0.911416in}}%
\pgfpathlineto{\pgfqpoint{2.413500in}{0.907511in}}%
\pgfpathlineto{\pgfqpoint{2.408265in}{0.902647in}}%
\pgfpathlineto{\pgfqpoint{2.403903in}{0.896059in}}%
\pgfpathlineto{\pgfqpoint{2.402990in}{0.891059in}}%
\pgfpathlineto{\pgfqpoint{2.391995in}{0.891918in}}%
\pgfpathlineto{\pgfqpoint{2.389462in}{0.896298in}}%
\pgfpathlineto{\pgfqpoint{2.385967in}{0.896018in}}%
\pgfpathlineto{\pgfqpoint{2.385392in}{0.899513in}}%
\pgfpathlineto{\pgfqpoint{2.391763in}{0.908652in}}%
\pgfpathlineto{\pgfqpoint{2.383505in}{0.917672in}}%
\pgfpathlineto{\pgfqpoint{2.378628in}{0.917255in}}%
\pgfpathlineto{\pgfqpoint{2.379773in}{0.923563in}}%
\pgfpathlineto{\pgfqpoint{2.381105in}{0.925986in}}%
\pgfpathlineto{\pgfqpoint{2.392521in}{0.928792in}}%
\pgfpathlineto{\pgfqpoint{2.398481in}{0.931422in}}%
\pgfpathlineto{\pgfqpoint{2.406157in}{0.920864in}}%
\pgfpathlineto{\pgfqpoint{2.410454in}{0.917486in}}%
\pgfpathlineto{\pgfqpoint{2.413670in}{0.911416in}}%
\pgfpathclose%
\pgfusepath{fill}%
\end{pgfscope}%
\begin{pgfscope}%
\pgfpathrectangle{\pgfqpoint{0.100000in}{0.100000in}}{\pgfqpoint{3.007045in}{1.925000in}}%
\pgfusepath{clip}%
\pgfsetbuttcap%
\pgfsetmiterjoin%
\definecolor{currentfill}{rgb}{0.282353,0.684045,0.378162}%
\pgfsetfillcolor{currentfill}%
\pgfsetlinewidth{0.000000pt}%
\definecolor{currentstroke}{rgb}{0.000000,0.000000,0.000000}%
\pgfsetstrokecolor{currentstroke}%
\pgfsetstrokeopacity{0.000000}%
\pgfsetdash{}{0pt}%
\pgfpathmoveto{\pgfqpoint{1.923472in}{0.954935in}}%
\pgfpathlineto{\pgfqpoint{1.921677in}{0.946416in}}%
\pgfpathlineto{\pgfqpoint{1.921721in}{0.941554in}}%
\pgfpathlineto{\pgfqpoint{1.915999in}{0.941499in}}%
\pgfpathlineto{\pgfqpoint{1.916049in}{0.935756in}}%
\pgfpathlineto{\pgfqpoint{1.910330in}{0.935711in}}%
\pgfpathlineto{\pgfqpoint{1.910317in}{0.941444in}}%
\pgfpathlineto{\pgfqpoint{1.887586in}{0.941336in}}%
\pgfpathlineto{\pgfqpoint{1.884698in}{0.944172in}}%
\pgfpathlineto{\pgfqpoint{1.884426in}{0.969017in}}%
\pgfpathlineto{\pgfqpoint{1.886046in}{0.969526in}}%
\pgfpathlineto{\pgfqpoint{1.890382in}{0.969613in}}%
\pgfpathlineto{\pgfqpoint{1.890280in}{0.989939in}}%
\pgfpathlineto{\pgfqpoint{1.924728in}{0.990132in}}%
\pgfpathlineto{\pgfqpoint{1.924925in}{0.970656in}}%
\pgfpathlineto{\pgfqpoint{1.923336in}{0.970590in}}%
\pgfpathlineto{\pgfqpoint{1.923472in}{0.954935in}}%
\pgfpathclose%
\pgfusepath{fill}%
\end{pgfscope}%
\begin{pgfscope}%
\pgfpathrectangle{\pgfqpoint{0.100000in}{0.100000in}}{\pgfqpoint{3.007045in}{1.925000in}}%
\pgfusepath{clip}%
\pgfsetbuttcap%
\pgfsetmiterjoin%
\definecolor{currentfill}{rgb}{0.382745,0.733256,0.427374}%
\pgfsetfillcolor{currentfill}%
\pgfsetlinewidth{0.000000pt}%
\definecolor{currentstroke}{rgb}{0.000000,0.000000,0.000000}%
\pgfsetstrokecolor{currentstroke}%
\pgfsetstrokeopacity{0.000000}%
\pgfsetdash{}{0pt}%
\pgfpathmoveto{\pgfqpoint{2.680930in}{1.081273in}}%
\pgfpathlineto{\pgfqpoint{2.683306in}{1.069582in}}%
\pgfpathlineto{\pgfqpoint{2.679621in}{1.057585in}}%
\pgfpathlineto{\pgfqpoint{2.672729in}{1.056217in}}%
\pgfpathlineto{\pgfqpoint{2.643747in}{1.050644in}}%
\pgfpathlineto{\pgfqpoint{2.608257in}{1.044199in}}%
\pgfpathlineto{\pgfqpoint{2.604048in}{1.043452in}}%
\pgfpathlineto{\pgfqpoint{2.604673in}{1.078358in}}%
\pgfpathlineto{\pgfqpoint{2.614846in}{1.077543in}}%
\pgfpathlineto{\pgfqpoint{2.616873in}{1.090165in}}%
\pgfpathlineto{\pgfqpoint{2.623729in}{1.094309in}}%
\pgfpathlineto{\pgfqpoint{2.626460in}{1.104802in}}%
\pgfpathlineto{\pgfqpoint{2.633816in}{1.102244in}}%
\pgfpathlineto{\pgfqpoint{2.641373in}{1.123470in}}%
\pgfpathlineto{\pgfqpoint{2.641185in}{1.127114in}}%
\pgfpathlineto{\pgfqpoint{2.644603in}{1.131634in}}%
\pgfpathlineto{\pgfqpoint{2.649818in}{1.126222in}}%
\pgfpathlineto{\pgfqpoint{2.649572in}{1.112889in}}%
\pgfpathlineto{\pgfqpoint{2.645322in}{1.106286in}}%
\pgfpathlineto{\pgfqpoint{2.646388in}{1.101662in}}%
\pgfpathlineto{\pgfqpoint{2.670092in}{1.099194in}}%
\pgfpathlineto{\pgfqpoint{2.665311in}{1.095015in}}%
\pgfpathlineto{\pgfqpoint{2.667768in}{1.085077in}}%
\pgfpathlineto{\pgfqpoint{2.674370in}{1.085389in}}%
\pgfpathlineto{\pgfqpoint{2.680930in}{1.081273in}}%
\pgfpathclose%
\pgfusepath{fill}%
\end{pgfscope}%
\begin{pgfscope}%
\pgfpathrectangle{\pgfqpoint{0.100000in}{0.100000in}}{\pgfqpoint{3.007045in}{1.925000in}}%
\pgfusepath{clip}%
\pgfsetbuttcap%
\pgfsetmiterjoin%
\definecolor{currentfill}{rgb}{0.252595,0.668128,0.362860}%
\pgfsetfillcolor{currentfill}%
\pgfsetlinewidth{0.000000pt}%
\definecolor{currentstroke}{rgb}{0.000000,0.000000,0.000000}%
\pgfsetstrokecolor{currentstroke}%
\pgfsetstrokeopacity{0.000000}%
\pgfsetdash{}{0pt}%
\pgfpathmoveto{\pgfqpoint{0.878844in}{1.527019in}}%
\pgfpathlineto{\pgfqpoint{0.872520in}{1.537698in}}%
\pgfpathlineto{\pgfqpoint{0.861425in}{1.547079in}}%
\pgfpathlineto{\pgfqpoint{0.860657in}{1.551817in}}%
\pgfpathlineto{\pgfqpoint{0.849873in}{1.555251in}}%
\pgfpathlineto{\pgfqpoint{0.842609in}{1.559070in}}%
\pgfpathlineto{\pgfqpoint{0.836443in}{1.560436in}}%
\pgfpathlineto{\pgfqpoint{0.834605in}{1.565249in}}%
\pgfpathlineto{\pgfqpoint{0.829274in}{1.562729in}}%
\pgfpathlineto{\pgfqpoint{0.827488in}{1.571131in}}%
\pgfpathlineto{\pgfqpoint{0.829905in}{1.577513in}}%
\pgfpathlineto{\pgfqpoint{0.824028in}{1.581837in}}%
\pgfpathlineto{\pgfqpoint{0.824845in}{1.588003in}}%
\pgfpathlineto{\pgfqpoint{0.820142in}{1.591987in}}%
\pgfpathlineto{\pgfqpoint{0.824171in}{1.596209in}}%
\pgfpathlineto{\pgfqpoint{0.825412in}{1.600994in}}%
\pgfpathlineto{\pgfqpoint{0.822960in}{1.604421in}}%
\pgfpathlineto{\pgfqpoint{0.824810in}{1.609423in}}%
\pgfpathlineto{\pgfqpoint{0.830622in}{1.610723in}}%
\pgfpathlineto{\pgfqpoint{0.834363in}{1.617027in}}%
\pgfpathlineto{\pgfqpoint{0.841346in}{1.613085in}}%
\pgfpathlineto{\pgfqpoint{0.853704in}{1.621066in}}%
\pgfpathlineto{\pgfqpoint{0.857696in}{1.638685in}}%
\pgfpathlineto{\pgfqpoint{0.860069in}{1.641241in}}%
\pgfpathlineto{\pgfqpoint{0.865574in}{1.646714in}}%
\pgfpathlineto{\pgfqpoint{0.859318in}{1.660383in}}%
\pgfpathlineto{\pgfqpoint{0.865364in}{1.658676in}}%
\pgfpathlineto{\pgfqpoint{0.871070in}{1.663297in}}%
\pgfpathlineto{\pgfqpoint{0.873352in}{1.663504in}}%
\pgfpathlineto{\pgfqpoint{0.879064in}{1.657147in}}%
\pgfpathlineto{\pgfqpoint{0.890559in}{1.659065in}}%
\pgfpathlineto{\pgfqpoint{0.894600in}{1.662965in}}%
\pgfpathlineto{\pgfqpoint{0.901414in}{1.666132in}}%
\pgfpathlineto{\pgfqpoint{0.906088in}{1.658804in}}%
\pgfpathlineto{\pgfqpoint{0.903714in}{1.653787in}}%
\pgfpathlineto{\pgfqpoint{0.906555in}{1.644348in}}%
\pgfpathlineto{\pgfqpoint{0.905525in}{1.640271in}}%
\pgfpathlineto{\pgfqpoint{0.910956in}{1.624666in}}%
\pgfpathlineto{\pgfqpoint{0.914945in}{1.620213in}}%
\pgfpathlineto{\pgfqpoint{0.912228in}{1.607641in}}%
\pgfpathlineto{\pgfqpoint{0.916159in}{1.603777in}}%
\pgfpathlineto{\pgfqpoint{0.923418in}{1.601071in}}%
\pgfpathlineto{\pgfqpoint{0.926655in}{1.595511in}}%
\pgfpathlineto{\pgfqpoint{0.928393in}{1.581446in}}%
\pgfpathlineto{\pgfqpoint{0.927577in}{1.576545in}}%
\pgfpathlineto{\pgfqpoint{0.933731in}{1.569068in}}%
\pgfpathlineto{\pgfqpoint{0.935411in}{1.569991in}}%
\pgfpathlineto{\pgfqpoint{0.933521in}{1.560492in}}%
\pgfpathlineto{\pgfqpoint{0.905554in}{1.566329in}}%
\pgfpathlineto{\pgfqpoint{0.903257in}{1.555054in}}%
\pgfpathlineto{\pgfqpoint{0.900875in}{1.550400in}}%
\pgfpathlineto{\pgfqpoint{0.897650in}{1.536577in}}%
\pgfpathlineto{\pgfqpoint{0.891068in}{1.533017in}}%
\pgfpathlineto{\pgfqpoint{0.881395in}{1.530161in}}%
\pgfpathlineto{\pgfqpoint{0.878844in}{1.527019in}}%
\pgfpathclose%
\pgfusepath{fill}%
\end{pgfscope}%
\begin{pgfscope}%
\pgfpathrectangle{\pgfqpoint{0.100000in}{0.100000in}}{\pgfqpoint{3.007045in}{1.925000in}}%
\pgfusepath{clip}%
\pgfsetbuttcap%
\pgfsetmiterjoin%
\definecolor{currentfill}{rgb}{0.741822,0.897486,0.715386}%
\pgfsetfillcolor{currentfill}%
\pgfsetlinewidth{0.000000pt}%
\definecolor{currentstroke}{rgb}{0.000000,0.000000,0.000000}%
\pgfsetstrokecolor{currentstroke}%
\pgfsetstrokeopacity{0.000000}%
\pgfsetdash{}{0pt}%
\pgfpathmoveto{\pgfqpoint{1.002522in}{1.140541in}}%
\pgfpathlineto{\pgfqpoint{0.997732in}{1.138462in}}%
\pgfpathlineto{\pgfqpoint{0.992394in}{1.131729in}}%
\pgfpathlineto{\pgfqpoint{0.990253in}{1.132805in}}%
\pgfpathlineto{\pgfqpoint{0.982614in}{1.127004in}}%
\pgfpathlineto{\pgfqpoint{0.975147in}{1.126984in}}%
\pgfpathlineto{\pgfqpoint{0.967830in}{1.113017in}}%
\pgfpathlineto{\pgfqpoint{0.961951in}{1.110649in}}%
\pgfpathlineto{\pgfqpoint{0.958933in}{1.107125in}}%
\pgfpathlineto{\pgfqpoint{0.937370in}{1.110856in}}%
\pgfpathlineto{\pgfqpoint{0.870763in}{1.123009in}}%
\pgfpathlineto{\pgfqpoint{0.854337in}{1.126666in}}%
\pgfpathlineto{\pgfqpoint{0.857375in}{1.143704in}}%
\pgfpathlineto{\pgfqpoint{0.863012in}{1.142510in}}%
\pgfpathlineto{\pgfqpoint{0.864002in}{1.147526in}}%
\pgfpathlineto{\pgfqpoint{0.869732in}{1.146966in}}%
\pgfpathlineto{\pgfqpoint{0.872431in}{1.163681in}}%
\pgfpathlineto{\pgfqpoint{0.877415in}{1.166653in}}%
\pgfpathlineto{\pgfqpoint{0.882037in}{1.174273in}}%
\pgfpathlineto{\pgfqpoint{0.878976in}{1.180671in}}%
\pgfpathlineto{\pgfqpoint{0.875518in}{1.183498in}}%
\pgfpathlineto{\pgfqpoint{0.875936in}{1.191519in}}%
\pgfpathlineto{\pgfqpoint{0.880713in}{1.197892in}}%
\pgfpathlineto{\pgfqpoint{0.887613in}{1.196226in}}%
\pgfpathlineto{\pgfqpoint{0.892855in}{1.198668in}}%
\pgfpathlineto{\pgfqpoint{0.894249in}{1.207621in}}%
\pgfpathlineto{\pgfqpoint{0.900045in}{1.212653in}}%
\pgfpathlineto{\pgfqpoint{0.903422in}{1.212075in}}%
\pgfpathlineto{\pgfqpoint{0.907070in}{1.217261in}}%
\pgfpathlineto{\pgfqpoint{0.914996in}{1.214937in}}%
\pgfpathlineto{\pgfqpoint{0.942870in}{1.209747in}}%
\pgfpathlineto{\pgfqpoint{0.936407in}{1.175969in}}%
\pgfpathlineto{\pgfqpoint{0.913573in}{1.180156in}}%
\pgfpathlineto{\pgfqpoint{0.907847in}{1.175524in}}%
\pgfpathlineto{\pgfqpoint{0.904834in}{1.157752in}}%
\pgfpathlineto{\pgfqpoint{0.963729in}{1.147104in}}%
\pgfpathlineto{\pgfqpoint{1.002522in}{1.140541in}}%
\pgfpathclose%
\pgfusepath{fill}%
\end{pgfscope}%
\begin{pgfscope}%
\pgfpathrectangle{\pgfqpoint{0.100000in}{0.100000in}}{\pgfqpoint{3.007045in}{1.925000in}}%
\pgfusepath{clip}%
\pgfsetbuttcap%
\pgfsetmiterjoin%
\definecolor{currentfill}{rgb}{0.524106,0.800923,0.519646}%
\pgfsetfillcolor{currentfill}%
\pgfsetlinewidth{0.000000pt}%
\definecolor{currentstroke}{rgb}{0.000000,0.000000,0.000000}%
\pgfsetstrokecolor{currentstroke}%
\pgfsetstrokeopacity{0.000000}%
\pgfsetdash{}{0pt}%
\pgfpathmoveto{\pgfqpoint{1.768458in}{1.001242in}}%
\pgfpathlineto{\pgfqpoint{1.750868in}{1.001271in}}%
\pgfpathlineto{\pgfqpoint{1.749585in}{0.997555in}}%
\pgfpathlineto{\pgfqpoint{1.749489in}{0.974797in}}%
\pgfpathlineto{\pgfqpoint{1.729327in}{0.974933in}}%
\pgfpathlineto{\pgfqpoint{1.729236in}{0.963455in}}%
\pgfpathlineto{\pgfqpoint{1.719212in}{0.963573in}}%
\pgfpathlineto{\pgfqpoint{1.719649in}{1.001488in}}%
\pgfpathlineto{\pgfqpoint{1.721571in}{1.001483in}}%
\pgfpathlineto{\pgfqpoint{1.722263in}{1.049928in}}%
\pgfpathlineto{\pgfqpoint{1.767740in}{1.049552in}}%
\pgfpathlineto{\pgfqpoint{1.768402in}{1.026569in}}%
\pgfpathlineto{\pgfqpoint{1.768458in}{1.001242in}}%
\pgfpathclose%
\pgfusepath{fill}%
\end{pgfscope}%
\begin{pgfscope}%
\pgfpathrectangle{\pgfqpoint{0.100000in}{0.100000in}}{\pgfqpoint{3.007045in}{1.925000in}}%
\pgfusepath{clip}%
\pgfsetbuttcap%
\pgfsetmiterjoin%
\definecolor{currentfill}{rgb}{0.357647,0.720953,0.415071}%
\pgfsetfillcolor{currentfill}%
\pgfsetlinewidth{0.000000pt}%
\definecolor{currentstroke}{rgb}{0.000000,0.000000,0.000000}%
\pgfsetstrokecolor{currentstroke}%
\pgfsetstrokeopacity{0.000000}%
\pgfsetdash{}{0pt}%
\pgfpathmoveto{\pgfqpoint{1.192875in}{1.732905in}}%
\pgfpathlineto{\pgfqpoint{1.188345in}{1.720102in}}%
\pgfpathlineto{\pgfqpoint{1.188167in}{1.711464in}}%
\pgfpathlineto{\pgfqpoint{1.189559in}{1.703789in}}%
\pgfpathlineto{\pgfqpoint{1.186086in}{1.695738in}}%
\pgfpathlineto{\pgfqpoint{1.187986in}{1.693328in}}%
\pgfpathlineto{\pgfqpoint{1.189998in}{1.686616in}}%
\pgfpathlineto{\pgfqpoint{1.153729in}{1.691500in}}%
\pgfpathlineto{\pgfqpoint{1.119627in}{1.696907in}}%
\pgfpathlineto{\pgfqpoint{1.119016in}{1.693144in}}%
\pgfpathlineto{\pgfqpoint{1.102971in}{1.695747in}}%
\pgfpathlineto{\pgfqpoint{1.102514in}{1.707496in}}%
\pgfpathlineto{\pgfqpoint{1.104710in}{1.709101in}}%
\pgfpathlineto{\pgfqpoint{1.107755in}{1.728110in}}%
\pgfpathlineto{\pgfqpoint{1.103070in}{1.731808in}}%
\pgfpathlineto{\pgfqpoint{1.097369in}{1.732746in}}%
\pgfpathlineto{\pgfqpoint{1.092154in}{1.736531in}}%
\pgfpathlineto{\pgfqpoint{1.094250in}{1.746968in}}%
\pgfpathlineto{\pgfqpoint{1.107591in}{1.752201in}}%
\pgfpathlineto{\pgfqpoint{1.110393in}{1.763495in}}%
\pgfpathlineto{\pgfqpoint{1.123379in}{1.762908in}}%
\pgfpathlineto{\pgfqpoint{1.126744in}{1.761331in}}%
\pgfpathlineto{\pgfqpoint{1.135562in}{1.764609in}}%
\pgfpathlineto{\pgfqpoint{1.141833in}{1.762846in}}%
\pgfpathlineto{\pgfqpoint{1.148828in}{1.763186in}}%
\pgfpathlineto{\pgfqpoint{1.152002in}{1.758091in}}%
\pgfpathlineto{\pgfqpoint{1.156612in}{1.751283in}}%
\pgfpathlineto{\pgfqpoint{1.164883in}{1.747708in}}%
\pgfpathlineto{\pgfqpoint{1.170837in}{1.747018in}}%
\pgfpathlineto{\pgfqpoint{1.174368in}{1.744062in}}%
\pgfpathlineto{\pgfqpoint{1.187367in}{1.743646in}}%
\pgfpathlineto{\pgfqpoint{1.193102in}{1.740744in}}%
\pgfpathlineto{\pgfqpoint{1.192875in}{1.732905in}}%
\pgfpathclose%
\pgfusepath{fill}%
\end{pgfscope}%
\begin{pgfscope}%
\pgfpathrectangle{\pgfqpoint{0.100000in}{0.100000in}}{\pgfqpoint{3.007045in}{1.925000in}}%
\pgfusepath{clip}%
\pgfsetbuttcap%
\pgfsetmiterjoin%
\definecolor{currentfill}{rgb}{0.699746,0.879769,0.674418}%
\pgfsetfillcolor{currentfill}%
\pgfsetlinewidth{0.000000pt}%
\definecolor{currentstroke}{rgb}{0.000000,0.000000,0.000000}%
\pgfsetstrokecolor{currentstroke}%
\pgfsetstrokeopacity{0.000000}%
\pgfsetdash{}{0pt}%
\pgfpathmoveto{\pgfqpoint{1.685427in}{0.912299in}}%
\pgfpathlineto{\pgfqpoint{1.685760in}{0.932254in}}%
\pgfpathlineto{\pgfqpoint{1.685998in}{0.946658in}}%
\pgfpathlineto{\pgfqpoint{1.675476in}{0.946791in}}%
\pgfpathlineto{\pgfqpoint{1.675597in}{0.952567in}}%
\pgfpathlineto{\pgfqpoint{1.669953in}{0.952684in}}%
\pgfpathlineto{\pgfqpoint{1.670084in}{0.958402in}}%
\pgfpathlineto{\pgfqpoint{1.664393in}{0.958534in}}%
\pgfpathlineto{\pgfqpoint{1.664644in}{0.969981in}}%
\pgfpathlineto{\pgfqpoint{1.671642in}{0.972779in}}%
\pgfpathlineto{\pgfqpoint{1.670693in}{0.976025in}}%
\pgfpathlineto{\pgfqpoint{1.642047in}{0.976241in}}%
\pgfpathlineto{\pgfqpoint{1.642737in}{1.002945in}}%
\pgfpathlineto{\pgfqpoint{1.659290in}{1.002553in}}%
\pgfpathlineto{\pgfqpoint{1.719649in}{1.001488in}}%
\pgfpathlineto{\pgfqpoint{1.719212in}{0.963573in}}%
\pgfpathlineto{\pgfqpoint{1.729236in}{0.963455in}}%
\pgfpathlineto{\pgfqpoint{1.729327in}{0.974933in}}%
\pgfpathlineto{\pgfqpoint{1.749489in}{0.974797in}}%
\pgfpathlineto{\pgfqpoint{1.749585in}{0.997555in}}%
\pgfpathlineto{\pgfqpoint{1.750868in}{1.001271in}}%
\pgfpathlineto{\pgfqpoint{1.768458in}{1.001242in}}%
\pgfpathlineto{\pgfqpoint{1.772364in}{0.997534in}}%
\pgfpathlineto{\pgfqpoint{1.772350in}{0.979565in}}%
\pgfpathlineto{\pgfqpoint{1.771725in}{0.946010in}}%
\pgfpathlineto{\pgfqpoint{1.766020in}{0.946003in}}%
\pgfpathlineto{\pgfqpoint{1.766027in}{0.940252in}}%
\pgfpathlineto{\pgfqpoint{1.761242in}{0.940278in}}%
\pgfpathlineto{\pgfqpoint{1.757448in}{0.932993in}}%
\pgfpathlineto{\pgfqpoint{1.758114in}{0.923040in}}%
\pgfpathlineto{\pgfqpoint{1.746777in}{0.922489in}}%
\pgfpathlineto{\pgfqpoint{1.740818in}{0.919900in}}%
\pgfpathlineto{\pgfqpoint{1.737553in}{0.926013in}}%
\pgfpathlineto{\pgfqpoint{1.731364in}{0.926038in}}%
\pgfpathlineto{\pgfqpoint{1.708573in}{0.926325in}}%
\pgfpathlineto{\pgfqpoint{1.708380in}{0.911948in}}%
\pgfpathlineto{\pgfqpoint{1.685427in}{0.912299in}}%
\pgfpathclose%
\pgfusepath{fill}%
\end{pgfscope}%
\begin{pgfscope}%
\pgfpathrectangle{\pgfqpoint{0.100000in}{0.100000in}}{\pgfqpoint{3.007045in}{1.925000in}}%
\pgfusepath{clip}%
\pgfsetbuttcap%
\pgfsetmiterjoin%
\definecolor{currentfill}{rgb}{0.727797,0.891580,0.701730}%
\pgfsetfillcolor{currentfill}%
\pgfsetlinewidth{0.000000pt}%
\definecolor{currentstroke}{rgb}{0.000000,0.000000,0.000000}%
\pgfsetstrokecolor{currentstroke}%
\pgfsetstrokeopacity{0.000000}%
\pgfsetdash{}{0pt}%
\pgfpathmoveto{\pgfqpoint{1.551371in}{1.482830in}}%
\pgfpathlineto{\pgfqpoint{1.552579in}{1.474853in}}%
\pgfpathlineto{\pgfqpoint{1.546378in}{1.477641in}}%
\pgfpathlineto{\pgfqpoint{1.542308in}{1.477127in}}%
\pgfpathlineto{\pgfqpoint{1.538292in}{1.482046in}}%
\pgfpathlineto{\pgfqpoint{1.517980in}{1.483159in}}%
\pgfpathlineto{\pgfqpoint{1.517880in}{1.481436in}}%
\pgfpathlineto{\pgfqpoint{1.485701in}{1.483457in}}%
\pgfpathlineto{\pgfqpoint{1.480026in}{1.483812in}}%
\pgfpathlineto{\pgfqpoint{1.481358in}{1.495327in}}%
\pgfpathlineto{\pgfqpoint{1.482943in}{1.518221in}}%
\pgfpathlineto{\pgfqpoint{1.484012in}{1.521937in}}%
\pgfpathlineto{\pgfqpoint{1.489275in}{1.520056in}}%
\pgfpathlineto{\pgfqpoint{1.498577in}{1.523374in}}%
\pgfpathlineto{\pgfqpoint{1.501256in}{1.521545in}}%
\pgfpathlineto{\pgfqpoint{1.506367in}{1.525118in}}%
\pgfpathlineto{\pgfqpoint{1.512562in}{1.521360in}}%
\pgfpathlineto{\pgfqpoint{1.519350in}{1.527670in}}%
\pgfpathlineto{\pgfqpoint{1.518970in}{1.529708in}}%
\pgfpathlineto{\pgfqpoint{1.553171in}{1.527805in}}%
\pgfpathlineto{\pgfqpoint{1.551887in}{1.504699in}}%
\pgfpathlineto{\pgfqpoint{1.551371in}{1.482830in}}%
\pgfpathclose%
\pgfusepath{fill}%
\end{pgfscope}%
\begin{pgfscope}%
\pgfpathrectangle{\pgfqpoint{0.100000in}{0.100000in}}{\pgfqpoint{3.007045in}{1.925000in}}%
\pgfusepath{clip}%
\pgfsetbuttcap%
\pgfsetmiterjoin%
\definecolor{currentfill}{rgb}{0.546251,0.811257,0.537855}%
\pgfsetfillcolor{currentfill}%
\pgfsetlinewidth{0.000000pt}%
\definecolor{currentstroke}{rgb}{0.000000,0.000000,0.000000}%
\pgfsetstrokecolor{currentstroke}%
\pgfsetstrokeopacity{0.000000}%
\pgfsetdash{}{0pt}%
\pgfpathmoveto{\pgfqpoint{2.020031in}{0.633536in}}%
\pgfpathlineto{\pgfqpoint{2.019201in}{0.650937in}}%
\pgfpathlineto{\pgfqpoint{2.013321in}{0.650638in}}%
\pgfpathlineto{\pgfqpoint{2.013078in}{0.656435in}}%
\pgfpathlineto{\pgfqpoint{2.038852in}{0.657801in}}%
\pgfpathlineto{\pgfqpoint{2.047327in}{0.661747in}}%
\pgfpathlineto{\pgfqpoint{2.055851in}{0.662822in}}%
\pgfpathlineto{\pgfqpoint{2.056646in}{0.646958in}}%
\pgfpathlineto{\pgfqpoint{2.058100in}{0.638257in}}%
\pgfpathlineto{\pgfqpoint{2.053466in}{0.638280in}}%
\pgfpathlineto{\pgfqpoint{2.049764in}{0.635164in}}%
\pgfpathlineto{\pgfqpoint{2.020031in}{0.633536in}}%
\pgfpathclose%
\pgfusepath{fill}%
\end{pgfscope}%
\begin{pgfscope}%
\pgfpathrectangle{\pgfqpoint{0.100000in}{0.100000in}}{\pgfqpoint{3.007045in}{1.925000in}}%
\pgfusepath{clip}%
\pgfsetbuttcap%
\pgfsetmiterjoin%
\definecolor{currentfill}{rgb}{0.917693,0.968428,0.901361}%
\pgfsetfillcolor{currentfill}%
\pgfsetlinewidth{0.000000pt}%
\definecolor{currentstroke}{rgb}{0.000000,0.000000,0.000000}%
\pgfsetstrokecolor{currentstroke}%
\pgfsetstrokeopacity{0.000000}%
\pgfsetdash{}{0pt}%
\pgfpathmoveto{\pgfqpoint{1.386813in}{0.624056in}}%
\pgfpathlineto{\pgfqpoint{1.351088in}{0.627381in}}%
\pgfpathlineto{\pgfqpoint{1.349725in}{0.635825in}}%
\pgfpathlineto{\pgfqpoint{1.337100in}{0.644492in}}%
\pgfpathlineto{\pgfqpoint{1.331766in}{0.642740in}}%
\pgfpathlineto{\pgfqpoint{1.333802in}{0.666433in}}%
\pgfpathlineto{\pgfqpoint{1.332017in}{0.666591in}}%
\pgfpathlineto{\pgfqpoint{1.334479in}{0.695294in}}%
\pgfpathlineto{\pgfqpoint{1.319633in}{0.696676in}}%
\pgfpathlineto{\pgfqpoint{1.322173in}{0.725422in}}%
\pgfpathlineto{\pgfqpoint{1.370186in}{0.721452in}}%
\pgfpathlineto{\pgfqpoint{1.398836in}{0.719387in}}%
\pgfpathlineto{\pgfqpoint{1.396350in}{0.690465in}}%
\pgfpathlineto{\pgfqpoint{1.391806in}{0.690752in}}%
\pgfpathlineto{\pgfqpoint{1.386813in}{0.624056in}}%
\pgfpathclose%
\pgfusepath{fill}%
\end{pgfscope}%
\begin{pgfscope}%
\pgfpathrectangle{\pgfqpoint{0.100000in}{0.100000in}}{\pgfqpoint{3.007045in}{1.925000in}}%
\pgfusepath{clip}%
\pgfsetbuttcap%
\pgfsetmiterjoin%
\definecolor{currentfill}{rgb}{0.573933,0.824175,0.560615}%
\pgfsetfillcolor{currentfill}%
\pgfsetlinewidth{0.000000pt}%
\definecolor{currentstroke}{rgb}{0.000000,0.000000,0.000000}%
\pgfsetstrokecolor{currentstroke}%
\pgfsetstrokeopacity{0.000000}%
\pgfsetdash{}{0pt}%
\pgfpathmoveto{\pgfqpoint{2.044424in}{1.118250in}}%
\pgfpathlineto{\pgfqpoint{2.020049in}{1.117235in}}%
\pgfpathlineto{\pgfqpoint{2.018572in}{1.121893in}}%
\pgfpathlineto{\pgfqpoint{2.022871in}{1.129850in}}%
\pgfpathlineto{\pgfqpoint{2.015419in}{1.134149in}}%
\pgfpathlineto{\pgfqpoint{2.005193in}{1.136744in}}%
\pgfpathlineto{\pgfqpoint{2.000724in}{1.130408in}}%
\pgfpathlineto{\pgfqpoint{1.994783in}{1.133102in}}%
\pgfpathlineto{\pgfqpoint{1.991482in}{1.141762in}}%
\pgfpathlineto{\pgfqpoint{1.992952in}{1.144098in}}%
\pgfpathlineto{\pgfqpoint{1.990535in}{1.152936in}}%
\pgfpathlineto{\pgfqpoint{1.986624in}{1.158384in}}%
\pgfpathlineto{\pgfqpoint{1.979066in}{1.163974in}}%
\pgfpathlineto{\pgfqpoint{1.995501in}{1.164493in}}%
\pgfpathlineto{\pgfqpoint{1.997266in}{1.151743in}}%
\pgfpathlineto{\pgfqpoint{2.002639in}{1.150292in}}%
\pgfpathlineto{\pgfqpoint{2.011513in}{1.150677in}}%
\pgfpathlineto{\pgfqpoint{2.019642in}{1.156899in}}%
\pgfpathlineto{\pgfqpoint{2.018445in}{1.173868in}}%
\pgfpathlineto{\pgfqpoint{2.029984in}{1.174625in}}%
\pgfpathlineto{\pgfqpoint{2.049876in}{1.175940in}}%
\pgfpathlineto{\pgfqpoint{2.050719in}{1.164389in}}%
\pgfpathlineto{\pgfqpoint{2.070604in}{1.165647in}}%
\pgfpathlineto{\pgfqpoint{2.071143in}{1.157025in}}%
\pgfpathlineto{\pgfqpoint{2.065520in}{1.156638in}}%
\pgfpathlineto{\pgfqpoint{2.067309in}{1.125374in}}%
\pgfpathlineto{\pgfqpoint{2.049708in}{1.124375in}}%
\pgfpathlineto{\pgfqpoint{2.050157in}{1.118638in}}%
\pgfpathlineto{\pgfqpoint{2.044424in}{1.118250in}}%
\pgfpathclose%
\pgfusepath{fill}%
\end{pgfscope}%
\begin{pgfscope}%
\pgfpathrectangle{\pgfqpoint{0.100000in}{0.100000in}}{\pgfqpoint{3.007045in}{1.925000in}}%
\pgfusepath{clip}%
\pgfsetbuttcap%
\pgfsetmiterjoin%
\definecolor{currentfill}{rgb}{0.557324,0.816424,0.546959}%
\pgfsetfillcolor{currentfill}%
\pgfsetlinewidth{0.000000pt}%
\definecolor{currentstroke}{rgb}{0.000000,0.000000,0.000000}%
\pgfsetstrokecolor{currentstroke}%
\pgfsetstrokeopacity{0.000000}%
\pgfsetdash{}{0pt}%
\pgfpathmoveto{\pgfqpoint{1.968862in}{0.741315in}}%
\pgfpathlineto{\pgfqpoint{1.935072in}{0.740260in}}%
\pgfpathlineto{\pgfqpoint{1.932837in}{0.741766in}}%
\pgfpathlineto{\pgfqpoint{1.930951in}{0.750384in}}%
\pgfpathlineto{\pgfqpoint{1.931371in}{0.753155in}}%
\pgfpathlineto{\pgfqpoint{1.939283in}{0.759647in}}%
\pgfpathlineto{\pgfqpoint{1.938599in}{0.766214in}}%
\pgfpathlineto{\pgfqpoint{1.968042in}{0.766663in}}%
\pgfpathlineto{\pgfqpoint{1.968862in}{0.741315in}}%
\pgfpathclose%
\pgfusepath{fill}%
\end{pgfscope}%
\begin{pgfscope}%
\pgfpathrectangle{\pgfqpoint{0.100000in}{0.100000in}}{\pgfqpoint{3.007045in}{1.925000in}}%
\pgfusepath{clip}%
\pgfsetbuttcap%
\pgfsetmiterjoin%
\definecolor{currentfill}{rgb}{0.395294,0.739408,0.433526}%
\pgfsetfillcolor{currentfill}%
\pgfsetlinewidth{0.000000pt}%
\definecolor{currentstroke}{rgb}{0.000000,0.000000,0.000000}%
\pgfsetstrokecolor{currentstroke}%
\pgfsetstrokeopacity{0.000000}%
\pgfsetdash{}{0pt}%
\pgfpathmoveto{\pgfqpoint{1.455754in}{1.302023in}}%
\pgfpathlineto{\pgfqpoint{1.427276in}{1.304079in}}%
\pgfpathlineto{\pgfqpoint{1.428992in}{1.326930in}}%
\pgfpathlineto{\pgfqpoint{1.424949in}{1.327230in}}%
\pgfpathlineto{\pgfqpoint{1.426380in}{1.344786in}}%
\pgfpathlineto{\pgfqpoint{1.428115in}{1.350443in}}%
\pgfpathlineto{\pgfqpoint{1.429987in}{1.373229in}}%
\pgfpathlineto{\pgfqpoint{1.428602in}{1.373346in}}%
\pgfpathlineto{\pgfqpoint{1.430345in}{1.395850in}}%
\pgfpathlineto{\pgfqpoint{1.430688in}{1.409985in}}%
\pgfpathlineto{\pgfqpoint{1.471839in}{1.406868in}}%
\pgfpathlineto{\pgfqpoint{1.521452in}{1.403789in}}%
\pgfpathlineto{\pgfqpoint{1.521365in}{1.389410in}}%
\pgfpathlineto{\pgfqpoint{1.520062in}{1.366510in}}%
\pgfpathlineto{\pgfqpoint{1.519516in}{1.343574in}}%
\pgfpathlineto{\pgfqpoint{1.486374in}{1.345738in}}%
\pgfpathlineto{\pgfqpoint{1.485080in}{1.322839in}}%
\pgfpathlineto{\pgfqpoint{1.457408in}{1.324892in}}%
\pgfpathlineto{\pgfqpoint{1.455754in}{1.302023in}}%
\pgfpathclose%
\pgfusepath{fill}%
\end{pgfscope}%
\begin{pgfscope}%
\pgfpathrectangle{\pgfqpoint{0.100000in}{0.100000in}}{\pgfqpoint{3.007045in}{1.925000in}}%
\pgfusepath{clip}%
\pgfsetbuttcap%
\pgfsetmiterjoin%
\definecolor{currentfill}{rgb}{0.551788,0.813841,0.542407}%
\pgfsetfillcolor{currentfill}%
\pgfsetlinewidth{0.000000pt}%
\definecolor{currentstroke}{rgb}{0.000000,0.000000,0.000000}%
\pgfsetstrokecolor{currentstroke}%
\pgfsetstrokeopacity{0.000000}%
\pgfsetdash{}{0pt}%
\pgfpathmoveto{\pgfqpoint{2.292953in}{1.341716in}}%
\pgfpathlineto{\pgfqpoint{2.273138in}{1.338797in}}%
\pgfpathlineto{\pgfqpoint{2.271732in}{1.342879in}}%
\pgfpathlineto{\pgfqpoint{2.269373in}{1.363358in}}%
\pgfpathlineto{\pgfqpoint{2.275068in}{1.363856in}}%
\pgfpathlineto{\pgfqpoint{2.272022in}{1.386856in}}%
\pgfpathlineto{\pgfqpoint{2.299946in}{1.390294in}}%
\pgfpathlineto{\pgfqpoint{2.303185in}{1.367262in}}%
\pgfpathlineto{\pgfqpoint{2.320460in}{1.370056in}}%
\pgfpathlineto{\pgfqpoint{2.323909in}{1.346635in}}%
\pgfpathlineto{\pgfqpoint{2.292953in}{1.341716in}}%
\pgfpathclose%
\pgfusepath{fill}%
\end{pgfscope}%
\begin{pgfscope}%
\pgfpathrectangle{\pgfqpoint{0.100000in}{0.100000in}}{\pgfqpoint{3.007045in}{1.925000in}}%
\pgfusepath{clip}%
\pgfsetbuttcap%
\pgfsetmiterjoin%
\definecolor{currentfill}{rgb}{0.579469,0.826759,0.565167}%
\pgfsetfillcolor{currentfill}%
\pgfsetlinewidth{0.000000pt}%
\definecolor{currentstroke}{rgb}{0.000000,0.000000,0.000000}%
\pgfsetstrokecolor{currentstroke}%
\pgfsetstrokeopacity{0.000000}%
\pgfsetdash{}{0pt}%
\pgfpathmoveto{\pgfqpoint{2.741068in}{1.332531in}}%
\pgfpathlineto{\pgfqpoint{2.720167in}{1.343542in}}%
\pgfpathlineto{\pgfqpoint{2.725895in}{1.348689in}}%
\pgfpathlineto{\pgfqpoint{2.711768in}{1.357635in}}%
\pgfpathlineto{\pgfqpoint{2.723365in}{1.366527in}}%
\pgfpathlineto{\pgfqpoint{2.720526in}{1.370514in}}%
\pgfpathlineto{\pgfqpoint{2.723613in}{1.374911in}}%
\pgfpathlineto{\pgfqpoint{2.727132in}{1.374300in}}%
\pgfpathlineto{\pgfqpoint{2.735363in}{1.367506in}}%
\pgfpathlineto{\pgfqpoint{2.732190in}{1.364248in}}%
\pgfpathlineto{\pgfqpoint{2.738618in}{1.356791in}}%
\pgfpathlineto{\pgfqpoint{2.746998in}{1.361954in}}%
\pgfpathlineto{\pgfqpoint{2.753651in}{1.370464in}}%
\pgfpathlineto{\pgfqpoint{2.758253in}{1.365099in}}%
\pgfpathlineto{\pgfqpoint{2.756767in}{1.359194in}}%
\pgfpathlineto{\pgfqpoint{2.753885in}{1.357653in}}%
\pgfpathlineto{\pgfqpoint{2.755066in}{1.346497in}}%
\pgfpathlineto{\pgfqpoint{2.746311in}{1.336376in}}%
\pgfpathlineto{\pgfqpoint{2.741068in}{1.332531in}}%
\pgfpathclose%
\pgfusepath{fill}%
\end{pgfscope}%
\begin{pgfscope}%
\pgfpathrectangle{\pgfqpoint{0.100000in}{0.100000in}}{\pgfqpoint{3.007045in}{1.925000in}}%
\pgfusepath{clip}%
\pgfsetbuttcap%
\pgfsetmiterjoin%
\definecolor{currentfill}{rgb}{0.618224,0.844844,0.597032}%
\pgfsetfillcolor{currentfill}%
\pgfsetlinewidth{0.000000pt}%
\definecolor{currentstroke}{rgb}{0.000000,0.000000,0.000000}%
\pgfsetstrokecolor{currentstroke}%
\pgfsetstrokeopacity{0.000000}%
\pgfsetdash{}{0pt}%
\pgfpathmoveto{\pgfqpoint{1.538952in}{0.605702in}}%
\pgfpathlineto{\pgfqpoint{1.546174in}{0.604074in}}%
\pgfpathlineto{\pgfqpoint{1.544968in}{0.575956in}}%
\pgfpathlineto{\pgfqpoint{1.525500in}{0.576932in}}%
\pgfpathlineto{\pgfqpoint{1.515156in}{0.577420in}}%
\pgfpathlineto{\pgfqpoint{1.516589in}{0.606731in}}%
\pgfpathlineto{\pgfqpoint{1.538952in}{0.605702in}}%
\pgfpathclose%
\pgfusepath{fill}%
\end{pgfscope}%
\begin{pgfscope}%
\pgfpathrectangle{\pgfqpoint{0.100000in}{0.100000in}}{\pgfqpoint{3.007045in}{1.925000in}}%
\pgfusepath{clip}%
\pgfsetbuttcap%
\pgfsetmiterjoin%
\definecolor{currentfill}{rgb}{0.474279,0.777670,0.478677}%
\pgfsetfillcolor{currentfill}%
\pgfsetlinewidth{0.000000pt}%
\definecolor{currentstroke}{rgb}{0.000000,0.000000,0.000000}%
\pgfsetstrokecolor{currentstroke}%
\pgfsetstrokeopacity{0.000000}%
\pgfsetdash{}{0pt}%
\pgfpathmoveto{\pgfqpoint{1.880439in}{1.202637in}}%
\pgfpathlineto{\pgfqpoint{1.854945in}{1.201835in}}%
\pgfpathlineto{\pgfqpoint{1.853670in}{1.237862in}}%
\pgfpathlineto{\pgfqpoint{1.886676in}{1.239179in}}%
\pgfpathlineto{\pgfqpoint{1.889157in}{1.227613in}}%
\pgfpathlineto{\pgfqpoint{1.888599in}{1.223016in}}%
\pgfpathlineto{\pgfqpoint{1.879993in}{1.222750in}}%
\pgfpathlineto{\pgfqpoint{1.880439in}{1.202637in}}%
\pgfpathclose%
\pgfusepath{fill}%
\end{pgfscope}%
\begin{pgfscope}%
\pgfpathrectangle{\pgfqpoint{0.100000in}{0.100000in}}{\pgfqpoint{3.007045in}{1.925000in}}%
\pgfusepath{clip}%
\pgfsetbuttcap%
\pgfsetmiterjoin%
\definecolor{currentfill}{rgb}{0.395294,0.739408,0.433526}%
\pgfsetfillcolor{currentfill}%
\pgfsetlinewidth{0.000000pt}%
\definecolor{currentstroke}{rgb}{0.000000,0.000000,0.000000}%
\pgfsetstrokecolor{currentstroke}%
\pgfsetstrokeopacity{0.000000}%
\pgfsetdash{}{0pt}%
\pgfpathmoveto{\pgfqpoint{2.612770in}{1.112672in}}%
\pgfpathlineto{\pgfqpoint{2.609095in}{1.111445in}}%
\pgfpathlineto{\pgfqpoint{2.602817in}{1.120583in}}%
\pgfpathlineto{\pgfqpoint{2.598994in}{1.120903in}}%
\pgfpathlineto{\pgfqpoint{2.593692in}{1.126779in}}%
\pgfpathlineto{\pgfqpoint{2.587434in}{1.124379in}}%
\pgfpathlineto{\pgfqpoint{2.584797in}{1.114957in}}%
\pgfpathlineto{\pgfqpoint{2.579736in}{1.109975in}}%
\pgfpathlineto{\pgfqpoint{2.578351in}{1.105297in}}%
\pgfpathlineto{\pgfqpoint{2.567398in}{1.112031in}}%
\pgfpathlineto{\pgfqpoint{2.566770in}{1.118700in}}%
\pgfpathlineto{\pgfqpoint{2.559027in}{1.119923in}}%
\pgfpathlineto{\pgfqpoint{2.554928in}{1.112611in}}%
\pgfpathlineto{\pgfqpoint{2.544389in}{1.103816in}}%
\pgfpathlineto{\pgfqpoint{2.535565in}{1.108737in}}%
\pgfpathlineto{\pgfqpoint{2.538295in}{1.117062in}}%
\pgfpathlineto{\pgfqpoint{2.547549in}{1.130953in}}%
\pgfpathlineto{\pgfqpoint{2.548827in}{1.135761in}}%
\pgfpathlineto{\pgfqpoint{2.549531in}{1.143698in}}%
\pgfpathlineto{\pgfqpoint{2.555307in}{1.150523in}}%
\pgfpathlineto{\pgfqpoint{2.553841in}{1.152723in}}%
\pgfpathlineto{\pgfqpoint{2.558583in}{1.161971in}}%
\pgfpathlineto{\pgfqpoint{2.558875in}{1.172717in}}%
\pgfpathlineto{\pgfqpoint{2.565022in}{1.171056in}}%
\pgfpathlineto{\pgfqpoint{2.569128in}{1.165501in}}%
\pgfpathlineto{\pgfqpoint{2.577995in}{1.164006in}}%
\pgfpathlineto{\pgfqpoint{2.581715in}{1.169068in}}%
\pgfpathlineto{\pgfqpoint{2.609006in}{1.155780in}}%
\pgfpathlineto{\pgfqpoint{2.608923in}{1.148127in}}%
\pgfpathlineto{\pgfqpoint{2.606282in}{1.144642in}}%
\pgfpathlineto{\pgfqpoint{2.616773in}{1.129743in}}%
\pgfpathlineto{\pgfqpoint{2.617784in}{1.122931in}}%
\pgfpathlineto{\pgfqpoint{2.613630in}{1.118818in}}%
\pgfpathlineto{\pgfqpoint{2.612770in}{1.112672in}}%
\pgfpathclose%
\pgfusepath{fill}%
\end{pgfscope}%
\begin{pgfscope}%
\pgfpathrectangle{\pgfqpoint{0.100000in}{0.100000in}}{\pgfqpoint{3.007045in}{1.925000in}}%
\pgfusepath{clip}%
\pgfsetbuttcap%
\pgfsetmiterjoin%
\definecolor{currentfill}{rgb}{0.618224,0.844844,0.597032}%
\pgfsetfillcolor{currentfill}%
\pgfsetlinewidth{0.000000pt}%
\definecolor{currentstroke}{rgb}{0.000000,0.000000,0.000000}%
\pgfsetstrokecolor{currentstroke}%
\pgfsetstrokeopacity{0.000000}%
\pgfsetdash{}{0pt}%
\pgfpathmoveto{\pgfqpoint{1.525500in}{0.576932in}}%
\pgfpathlineto{\pgfqpoint{1.524943in}{0.562745in}}%
\pgfpathlineto{\pgfqpoint{1.498977in}{0.564304in}}%
\pgfpathlineto{\pgfqpoint{1.478167in}{0.565413in}}%
\pgfpathlineto{\pgfqpoint{1.479753in}{0.593339in}}%
\pgfpathlineto{\pgfqpoint{1.481225in}{0.618416in}}%
\pgfpathlineto{\pgfqpoint{1.510333in}{0.616812in}}%
\pgfpathlineto{\pgfqpoint{1.509816in}{0.607078in}}%
\pgfpathlineto{\pgfqpoint{1.516589in}{0.606731in}}%
\pgfpathlineto{\pgfqpoint{1.515156in}{0.577420in}}%
\pgfpathlineto{\pgfqpoint{1.525500in}{0.576932in}}%
\pgfpathclose%
\pgfusepath{fill}%
\end{pgfscope}%
\begin{pgfscope}%
\pgfpathrectangle{\pgfqpoint{0.100000in}{0.100000in}}{\pgfqpoint{3.007045in}{1.925000in}}%
\pgfusepath{clip}%
\pgfsetbuttcap%
\pgfsetmiterjoin%
\definecolor{currentfill}{rgb}{0.351373,0.717878,0.411995}%
\pgfsetfillcolor{currentfill}%
\pgfsetlinewidth{0.000000pt}%
\definecolor{currentstroke}{rgb}{0.000000,0.000000,0.000000}%
\pgfsetstrokecolor{currentstroke}%
\pgfsetstrokeopacity{0.000000}%
\pgfsetdash{}{0pt}%
\pgfpathmoveto{\pgfqpoint{2.671279in}{0.875019in}}%
\pgfpathlineto{\pgfqpoint{2.637551in}{0.899069in}}%
\pgfpathlineto{\pgfqpoint{2.641884in}{0.909457in}}%
\pgfpathlineto{\pgfqpoint{2.647369in}{0.914791in}}%
\pgfpathlineto{\pgfqpoint{2.647207in}{0.927059in}}%
\pgfpathlineto{\pgfqpoint{2.640458in}{0.935526in}}%
\pgfpathlineto{\pgfqpoint{2.645656in}{0.937482in}}%
\pgfpathlineto{\pgfqpoint{2.661874in}{0.940850in}}%
\pgfpathlineto{\pgfqpoint{2.668672in}{0.934954in}}%
\pgfpathlineto{\pgfqpoint{2.678361in}{0.924174in}}%
\pgfpathlineto{\pgfqpoint{2.683843in}{0.935913in}}%
\pgfpathlineto{\pgfqpoint{2.706932in}{0.940219in}}%
\pgfpathlineto{\pgfqpoint{2.714846in}{0.926915in}}%
\pgfpathlineto{\pgfqpoint{2.719169in}{0.923665in}}%
\pgfpathlineto{\pgfqpoint{2.709222in}{0.910276in}}%
\pgfpathlineto{\pgfqpoint{2.705650in}{0.901507in}}%
\pgfpathlineto{\pgfqpoint{2.702584in}{0.882702in}}%
\pgfpathlineto{\pgfqpoint{2.691067in}{0.882886in}}%
\pgfpathlineto{\pgfqpoint{2.680626in}{0.880256in}}%
\pgfpathlineto{\pgfqpoint{2.671279in}{0.875019in}}%
\pgfpathclose%
\pgfusepath{fill}%
\end{pgfscope}%
\begin{pgfscope}%
\pgfpathrectangle{\pgfqpoint{0.100000in}{0.100000in}}{\pgfqpoint{3.007045in}{1.925000in}}%
\pgfusepath{clip}%
\pgfsetbuttcap%
\pgfsetmiterjoin%
\definecolor{currentfill}{rgb}{0.573933,0.824175,0.560615}%
\pgfsetfillcolor{currentfill}%
\pgfsetlinewidth{0.000000pt}%
\definecolor{currentstroke}{rgb}{0.000000,0.000000,0.000000}%
\pgfsetstrokecolor{currentstroke}%
\pgfsetstrokeopacity{0.000000}%
\pgfsetdash{}{0pt}%
\pgfpathmoveto{\pgfqpoint{0.540209in}{1.259221in}}%
\pgfpathlineto{\pgfqpoint{0.533766in}{1.269179in}}%
\pgfpathlineto{\pgfqpoint{0.536090in}{1.281820in}}%
\pgfpathlineto{\pgfqpoint{0.533368in}{1.282853in}}%
\pgfpathlineto{\pgfqpoint{0.536518in}{1.297107in}}%
\pgfpathlineto{\pgfqpoint{0.542583in}{1.297900in}}%
\pgfpathlineto{\pgfqpoint{0.544151in}{1.303815in}}%
\pgfpathlineto{\pgfqpoint{0.522474in}{1.309901in}}%
\pgfpathlineto{\pgfqpoint{0.522542in}{1.311873in}}%
\pgfpathlineto{\pgfqpoint{0.510601in}{1.315136in}}%
\pgfpathlineto{\pgfqpoint{0.516565in}{1.336204in}}%
\pgfpathlineto{\pgfqpoint{0.521615in}{1.353758in}}%
\pgfpathlineto{\pgfqpoint{0.535789in}{1.402977in}}%
\pgfpathlineto{\pgfqpoint{0.547844in}{1.446415in}}%
\pgfpathlineto{\pgfqpoint{0.562393in}{1.497838in}}%
\pgfpathlineto{\pgfqpoint{0.592563in}{1.489427in}}%
\pgfpathlineto{\pgfqpoint{0.594261in}{1.488948in}}%
\pgfpathlineto{\pgfqpoint{0.584193in}{1.452083in}}%
\pgfpathlineto{\pgfqpoint{0.575970in}{1.423296in}}%
\pgfpathlineto{\pgfqpoint{0.576936in}{1.423028in}}%
\pgfpathlineto{\pgfqpoint{0.569070in}{1.394980in}}%
\pgfpathlineto{\pgfqpoint{0.567831in}{1.395284in}}%
\pgfpathlineto{\pgfqpoint{0.558741in}{1.362309in}}%
\pgfpathlineto{\pgfqpoint{0.594471in}{1.352638in}}%
\pgfpathlineto{\pgfqpoint{0.646486in}{1.339327in}}%
\pgfpathlineto{\pgfqpoint{0.648824in}{1.333373in}}%
\pgfpathlineto{\pgfqpoint{0.644679in}{1.322619in}}%
\pgfpathlineto{\pgfqpoint{0.644313in}{1.314171in}}%
\pgfpathlineto{\pgfqpoint{0.641368in}{1.309018in}}%
\pgfpathlineto{\pgfqpoint{0.632081in}{1.306373in}}%
\pgfpathlineto{\pgfqpoint{0.625081in}{1.300716in}}%
\pgfpathlineto{\pgfqpoint{0.625302in}{1.292738in}}%
\pgfpathlineto{\pgfqpoint{0.621234in}{1.289668in}}%
\pgfpathlineto{\pgfqpoint{0.620230in}{1.284320in}}%
\pgfpathlineto{\pgfqpoint{0.615468in}{1.284198in}}%
\pgfpathlineto{\pgfqpoint{0.562826in}{1.298061in}}%
\pgfpathlineto{\pgfqpoint{0.556383in}{1.290907in}}%
\pgfpathlineto{\pgfqpoint{0.554814in}{1.285052in}}%
\pgfpathlineto{\pgfqpoint{0.560395in}{1.283604in}}%
\pgfpathlineto{\pgfqpoint{0.552667in}{1.255842in}}%
\pgfpathlineto{\pgfqpoint{0.540209in}{1.259221in}}%
\pgfpathclose%
\pgfusepath{fill}%
\end{pgfscope}%
\begin{pgfscope}%
\pgfpathrectangle{\pgfqpoint{0.100000in}{0.100000in}}{\pgfqpoint{3.007045in}{1.925000in}}%
\pgfusepath{clip}%
\pgfsetbuttcap%
\pgfsetmiterjoin%
\definecolor{currentfill}{rgb}{0.535179,0.806090,0.528750}%
\pgfsetfillcolor{currentfill}%
\pgfsetlinewidth{0.000000pt}%
\definecolor{currentstroke}{rgb}{0.000000,0.000000,0.000000}%
\pgfsetstrokecolor{currentstroke}%
\pgfsetstrokeopacity{0.000000}%
\pgfsetdash{}{0pt}%
\pgfpathmoveto{\pgfqpoint{1.591111in}{1.462290in}}%
\pgfpathlineto{\pgfqpoint{1.585428in}{1.462527in}}%
\pgfpathlineto{\pgfqpoint{1.586215in}{1.479799in}}%
\pgfpathlineto{\pgfqpoint{1.596903in}{1.479359in}}%
\pgfpathlineto{\pgfqpoint{1.597846in}{1.508282in}}%
\pgfpathlineto{\pgfqpoint{1.632109in}{1.506861in}}%
\pgfpathlineto{\pgfqpoint{1.637825in}{1.506641in}}%
\pgfpathlineto{\pgfqpoint{1.637786in}{1.500890in}}%
\pgfpathlineto{\pgfqpoint{1.637041in}{1.477807in}}%
\pgfpathlineto{\pgfqpoint{1.614383in}{1.478679in}}%
\pgfpathlineto{\pgfqpoint{1.613769in}{1.461576in}}%
\pgfpathlineto{\pgfqpoint{1.591111in}{1.462290in}}%
\pgfpathclose%
\pgfusepath{fill}%
\end{pgfscope}%
\begin{pgfscope}%
\pgfpathrectangle{\pgfqpoint{0.100000in}{0.100000in}}{\pgfqpoint{3.007045in}{1.925000in}}%
\pgfusepath{clip}%
\pgfsetbuttcap%
\pgfsetmiterjoin%
\definecolor{currentfill}{rgb}{0.167705,0.577578,0.294948}%
\pgfsetfillcolor{currentfill}%
\pgfsetlinewidth{0.000000pt}%
\definecolor{currentstroke}{rgb}{0.000000,0.000000,0.000000}%
\pgfsetstrokecolor{currentstroke}%
\pgfsetstrokeopacity{0.000000}%
\pgfsetdash{}{0pt}%
\pgfpathmoveto{\pgfqpoint{2.735476in}{1.140922in}}%
\pgfpathlineto{\pgfqpoint{2.729649in}{1.140417in}}%
\pgfpathlineto{\pgfqpoint{2.725113in}{1.146693in}}%
\pgfpathlineto{\pgfqpoint{2.720436in}{1.148513in}}%
\pgfpathlineto{\pgfqpoint{2.716860in}{1.144435in}}%
\pgfpathlineto{\pgfqpoint{2.713398in}{1.148167in}}%
\pgfpathlineto{\pgfqpoint{2.706818in}{1.146620in}}%
\pgfpathlineto{\pgfqpoint{2.701679in}{1.149182in}}%
\pgfpathlineto{\pgfqpoint{2.697515in}{1.156406in}}%
\pgfpathlineto{\pgfqpoint{2.690990in}{1.162746in}}%
\pgfpathlineto{\pgfqpoint{2.695805in}{1.170685in}}%
\pgfpathlineto{\pgfqpoint{2.697143in}{1.177380in}}%
\pgfpathlineto{\pgfqpoint{2.704729in}{1.171795in}}%
\pgfpathlineto{\pgfqpoint{2.718443in}{1.173001in}}%
\pgfpathlineto{\pgfqpoint{2.728305in}{1.164174in}}%
\pgfpathlineto{\pgfqpoint{2.740946in}{1.160176in}}%
\pgfpathlineto{\pgfqpoint{2.739016in}{1.152520in}}%
\pgfpathlineto{\pgfqpoint{2.740280in}{1.146999in}}%
\pgfpathlineto{\pgfqpoint{2.738357in}{1.140814in}}%
\pgfpathlineto{\pgfqpoint{2.735476in}{1.140922in}}%
\pgfpathclose%
\pgfusepath{fill}%
\end{pgfscope}%
\begin{pgfscope}%
\pgfpathrectangle{\pgfqpoint{0.100000in}{0.100000in}}{\pgfqpoint{3.007045in}{1.925000in}}%
\pgfusepath{clip}%
\pgfsetbuttcap%
\pgfsetmiterjoin%
\definecolor{currentfill}{rgb}{0.518570,0.798339,0.515094}%
\pgfsetfillcolor{currentfill}%
\pgfsetlinewidth{0.000000pt}%
\definecolor{currentstroke}{rgb}{0.000000,0.000000,0.000000}%
\pgfsetstrokecolor{currentstroke}%
\pgfsetstrokeopacity{0.000000}%
\pgfsetdash{}{0pt}%
\pgfpathmoveto{\pgfqpoint{1.631910in}{1.063521in}}%
\pgfpathlineto{\pgfqpoint{1.631737in}{1.051813in}}%
\pgfpathlineto{\pgfqpoint{1.591838in}{1.053004in}}%
\pgfpathlineto{\pgfqpoint{1.591676in}{1.059068in}}%
\pgfpathlineto{\pgfqpoint{1.593165in}{1.116517in}}%
\pgfpathlineto{\pgfqpoint{1.593668in}{1.127954in}}%
\pgfpathlineto{\pgfqpoint{1.622153in}{1.126974in}}%
\pgfpathlineto{\pgfqpoint{1.621825in}{1.109765in}}%
\pgfpathlineto{\pgfqpoint{1.621015in}{1.081045in}}%
\pgfpathlineto{\pgfqpoint{1.632419in}{1.080689in}}%
\pgfpathlineto{\pgfqpoint{1.631910in}{1.063521in}}%
\pgfpathclose%
\pgfusepath{fill}%
\end{pgfscope}%
\begin{pgfscope}%
\pgfpathrectangle{\pgfqpoint{0.100000in}{0.100000in}}{\pgfqpoint{3.007045in}{1.925000in}}%
\pgfusepath{clip}%
\pgfsetbuttcap%
\pgfsetmiterjoin%
\definecolor{currentfill}{rgb}{0.704421,0.881738,0.678970}%
\pgfsetfillcolor{currentfill}%
\pgfsetlinewidth{0.000000pt}%
\definecolor{currentstroke}{rgb}{0.000000,0.000000,0.000000}%
\pgfsetstrokecolor{currentstroke}%
\pgfsetstrokeopacity{0.000000}%
\pgfsetdash{}{0pt}%
\pgfpathmoveto{\pgfqpoint{2.397002in}{0.647345in}}%
\pgfpathlineto{\pgfqpoint{2.375927in}{0.644931in}}%
\pgfpathlineto{\pgfqpoint{2.366514in}{0.643850in}}%
\pgfpathlineto{\pgfqpoint{2.365478in}{0.655570in}}%
\pgfpathlineto{\pgfqpoint{2.359552in}{0.655091in}}%
\pgfpathlineto{\pgfqpoint{2.358460in}{0.666668in}}%
\pgfpathlineto{\pgfqpoint{2.366771in}{0.667222in}}%
\pgfpathlineto{\pgfqpoint{2.370955in}{0.669671in}}%
\pgfpathlineto{\pgfqpoint{2.367869in}{0.675418in}}%
\pgfpathlineto{\pgfqpoint{2.367559in}{0.680287in}}%
\pgfpathlineto{\pgfqpoint{2.362322in}{0.683671in}}%
\pgfpathlineto{\pgfqpoint{2.357840in}{0.689197in}}%
\pgfpathlineto{\pgfqpoint{2.356903in}{0.698990in}}%
\pgfpathlineto{\pgfqpoint{2.394856in}{0.702917in}}%
\pgfpathlineto{\pgfqpoint{2.395249in}{0.698742in}}%
\pgfpathlineto{\pgfqpoint{2.402438in}{0.696626in}}%
\pgfpathlineto{\pgfqpoint{2.404240in}{0.684768in}}%
\pgfpathlineto{\pgfqpoint{2.408654in}{0.685257in}}%
\pgfpathlineto{\pgfqpoint{2.412962in}{0.682194in}}%
\pgfpathlineto{\pgfqpoint{2.414666in}{0.666517in}}%
\pgfpathlineto{\pgfqpoint{2.395188in}{0.664380in}}%
\pgfpathlineto{\pgfqpoint{2.397002in}{0.647345in}}%
\pgfpathclose%
\pgfusepath{fill}%
\end{pgfscope}%
\begin{pgfscope}%
\pgfpathrectangle{\pgfqpoint{0.100000in}{0.100000in}}{\pgfqpoint{3.007045in}{1.925000in}}%
\pgfusepath{clip}%
\pgfsetbuttcap%
\pgfsetmiterjoin%
\definecolor{currentfill}{rgb}{0.237832,0.652380,0.351050}%
\pgfsetfillcolor{currentfill}%
\pgfsetlinewidth{0.000000pt}%
\definecolor{currentstroke}{rgb}{0.000000,0.000000,0.000000}%
\pgfsetstrokecolor{currentstroke}%
\pgfsetstrokeopacity{0.000000}%
\pgfsetdash{}{0pt}%
\pgfpathmoveto{\pgfqpoint{2.209634in}{1.589574in}}%
\pgfpathlineto{\pgfqpoint{2.206269in}{1.590978in}}%
\pgfpathlineto{\pgfqpoint{2.208079in}{1.601957in}}%
\pgfpathlineto{\pgfqpoint{2.211439in}{1.600059in}}%
\pgfpathlineto{\pgfqpoint{2.212615in}{1.592362in}}%
\pgfpathlineto{\pgfqpoint{2.209634in}{1.589574in}}%
\pgfpathclose%
\pgfusepath{fill}%
\end{pgfscope}%
\begin{pgfscope}%
\pgfpathrectangle{\pgfqpoint{0.100000in}{0.100000in}}{\pgfqpoint{3.007045in}{1.925000in}}%
\pgfusepath{clip}%
\pgfsetbuttcap%
\pgfsetmiterjoin%
\definecolor{currentfill}{rgb}{0.237832,0.652380,0.351050}%
\pgfsetfillcolor{currentfill}%
\pgfsetlinewidth{0.000000pt}%
\definecolor{currentstroke}{rgb}{0.000000,0.000000,0.000000}%
\pgfsetstrokecolor{currentstroke}%
\pgfsetstrokeopacity{0.000000}%
\pgfsetdash{}{0pt}%
\pgfpathmoveto{\pgfqpoint{2.271744in}{1.600141in}}%
\pgfpathlineto{\pgfqpoint{2.269589in}{1.599880in}}%
\pgfpathlineto{\pgfqpoint{2.273175in}{1.571785in}}%
\pgfpathlineto{\pgfqpoint{2.250552in}{1.569368in}}%
\pgfpathlineto{\pgfqpoint{2.251176in}{1.563619in}}%
\pgfpathlineto{\pgfqpoint{2.245496in}{1.563087in}}%
\pgfpathlineto{\pgfqpoint{2.228428in}{1.561303in}}%
\pgfpathlineto{\pgfqpoint{2.227849in}{1.567100in}}%
\pgfpathlineto{\pgfqpoint{2.220169in}{1.566322in}}%
\pgfpathlineto{\pgfqpoint{2.223475in}{1.574310in}}%
\pgfpathlineto{\pgfqpoint{2.230581in}{1.578556in}}%
\pgfpathlineto{\pgfqpoint{2.234669in}{1.583276in}}%
\pgfpathlineto{\pgfqpoint{2.231732in}{1.586161in}}%
\pgfpathlineto{\pgfqpoint{2.230119in}{1.592502in}}%
\pgfpathlineto{\pgfqpoint{2.237120in}{1.602372in}}%
\pgfpathlineto{\pgfqpoint{2.243327in}{1.604920in}}%
\pgfpathlineto{\pgfqpoint{2.244332in}{1.607881in}}%
\pgfpathlineto{\pgfqpoint{2.259869in}{1.600439in}}%
\pgfpathlineto{\pgfqpoint{2.265685in}{1.601992in}}%
\pgfpathlineto{\pgfqpoint{2.271744in}{1.600141in}}%
\pgfpathclose%
\pgfusepath{fill}%
\end{pgfscope}%
\begin{pgfscope}%
\pgfpathrectangle{\pgfqpoint{0.100000in}{0.100000in}}{\pgfqpoint{3.007045in}{1.925000in}}%
\pgfusepath{clip}%
\pgfsetbuttcap%
\pgfsetmiterjoin%
\definecolor{currentfill}{rgb}{0.699746,0.879769,0.674418}%
\pgfsetfillcolor{currentfill}%
\pgfsetlinewidth{0.000000pt}%
\definecolor{currentstroke}{rgb}{0.000000,0.000000,0.000000}%
\pgfsetstrokecolor{currentstroke}%
\pgfsetstrokeopacity{0.000000}%
\pgfsetdash{}{0pt}%
\pgfpathmoveto{\pgfqpoint{2.415134in}{0.621502in}}%
\pgfpathlineto{\pgfqpoint{2.422704in}{0.621971in}}%
\pgfpathlineto{\pgfqpoint{2.423706in}{0.617445in}}%
\pgfpathlineto{\pgfqpoint{2.416778in}{0.611640in}}%
\pgfpathlineto{\pgfqpoint{2.418230in}{0.606839in}}%
\pgfpathlineto{\pgfqpoint{2.413904in}{0.601638in}}%
\pgfpathlineto{\pgfqpoint{2.415159in}{0.598741in}}%
\pgfpathlineto{\pgfqpoint{2.408711in}{0.592974in}}%
\pgfpathlineto{\pgfqpoint{2.405641in}{0.582074in}}%
\pgfpathlineto{\pgfqpoint{2.395790in}{0.579657in}}%
\pgfpathlineto{\pgfqpoint{2.389487in}{0.580324in}}%
\pgfpathlineto{\pgfqpoint{2.384731in}{0.575233in}}%
\pgfpathlineto{\pgfqpoint{2.386694in}{0.571228in}}%
\pgfpathlineto{\pgfqpoint{2.374564in}{0.573201in}}%
\pgfpathlineto{\pgfqpoint{2.347979in}{0.570392in}}%
\pgfpathlineto{\pgfqpoint{2.340428in}{0.571556in}}%
\pgfpathlineto{\pgfqpoint{2.339480in}{0.580648in}}%
\pgfpathlineto{\pgfqpoint{2.324742in}{0.580445in}}%
\pgfpathlineto{\pgfqpoint{2.322486in}{0.604682in}}%
\pgfpathlineto{\pgfqpoint{2.316671in}{0.604111in}}%
\pgfpathlineto{\pgfqpoint{2.317482in}{0.595459in}}%
\pgfpathlineto{\pgfqpoint{2.296561in}{0.593554in}}%
\pgfpathlineto{\pgfqpoint{2.288959in}{0.589362in}}%
\pgfpathlineto{\pgfqpoint{2.295880in}{0.596792in}}%
\pgfpathlineto{\pgfqpoint{2.291737in}{0.606164in}}%
\pgfpathlineto{\pgfqpoint{2.294588in}{0.610608in}}%
\pgfpathlineto{\pgfqpoint{2.284322in}{0.609983in}}%
\pgfpathlineto{\pgfqpoint{2.282576in}{0.629115in}}%
\pgfpathlineto{\pgfqpoint{2.341073in}{0.635732in}}%
\pgfpathlineto{\pgfqpoint{2.345288in}{0.628691in}}%
\pgfpathlineto{\pgfqpoint{2.346074in}{0.624253in}}%
\pgfpathlineto{\pgfqpoint{2.351014in}{0.617483in}}%
\pgfpathlineto{\pgfqpoint{2.378692in}{0.619214in}}%
\pgfpathlineto{\pgfqpoint{2.415134in}{0.621502in}}%
\pgfpathclose%
\pgfusepath{fill}%
\end{pgfscope}%
\begin{pgfscope}%
\pgfpathrectangle{\pgfqpoint{0.100000in}{0.100000in}}{\pgfqpoint{3.007045in}{1.925000in}}%
\pgfusepath{clip}%
\pgfsetbuttcap%
\pgfsetmiterjoin%
\definecolor{currentfill}{rgb}{0.794233,0.919262,0.767705}%
\pgfsetfillcolor{currentfill}%
\pgfsetlinewidth{0.000000pt}%
\definecolor{currentstroke}{rgb}{0.000000,0.000000,0.000000}%
\pgfsetstrokecolor{currentstroke}%
\pgfsetstrokeopacity{0.000000}%
\pgfsetdash{}{0pt}%
\pgfpathmoveto{\pgfqpoint{1.448770in}{1.649792in}}%
\pgfpathlineto{\pgfqpoint{1.449237in}{1.655533in}}%
\pgfpathlineto{\pgfqpoint{1.466420in}{1.654245in}}%
\pgfpathlineto{\pgfqpoint{1.465987in}{1.648477in}}%
\pgfpathlineto{\pgfqpoint{1.485073in}{1.647065in}}%
\pgfpathlineto{\pgfqpoint{1.483880in}{1.629904in}}%
\pgfpathlineto{\pgfqpoint{1.495332in}{1.629121in}}%
\pgfpathlineto{\pgfqpoint{1.495994in}{1.623275in}}%
\pgfpathlineto{\pgfqpoint{1.482806in}{1.614036in}}%
\pgfpathlineto{\pgfqpoint{1.480364in}{1.610710in}}%
\pgfpathlineto{\pgfqpoint{1.470594in}{1.606845in}}%
\pgfpathlineto{\pgfqpoint{1.459052in}{1.608999in}}%
\pgfpathlineto{\pgfqpoint{1.455103in}{1.611872in}}%
\pgfpathlineto{\pgfqpoint{1.450280in}{1.611196in}}%
\pgfpathlineto{\pgfqpoint{1.451489in}{1.626431in}}%
\pgfpathlineto{\pgfqpoint{1.449255in}{1.626605in}}%
\pgfpathlineto{\pgfqpoint{1.450846in}{1.649617in}}%
\pgfpathlineto{\pgfqpoint{1.448770in}{1.649792in}}%
\pgfpathclose%
\pgfusepath{fill}%
\end{pgfscope}%
\begin{pgfscope}%
\pgfpathrectangle{\pgfqpoint{0.100000in}{0.100000in}}{\pgfqpoint{3.007045in}{1.925000in}}%
\pgfusepath{clip}%
\pgfsetbuttcap%
\pgfsetmiterjoin%
\definecolor{currentfill}{rgb}{0.451765,0.767090,0.461207}%
\pgfsetfillcolor{currentfill}%
\pgfsetlinewidth{0.000000pt}%
\definecolor{currentstroke}{rgb}{0.000000,0.000000,0.000000}%
\pgfsetstrokecolor{currentstroke}%
\pgfsetstrokeopacity{0.000000}%
\pgfsetdash{}{0pt}%
\pgfpathmoveto{\pgfqpoint{1.408351in}{1.098572in}}%
\pgfpathlineto{\pgfqpoint{1.408374in}{1.098989in}}%
\pgfpathlineto{\pgfqpoint{1.410528in}{1.127149in}}%
\pgfpathlineto{\pgfqpoint{1.439278in}{1.125169in}}%
\pgfpathlineto{\pgfqpoint{1.480135in}{1.122427in}}%
\pgfpathlineto{\pgfqpoint{1.478511in}{1.093767in}}%
\pgfpathlineto{\pgfqpoint{1.432949in}{1.096770in}}%
\pgfpathlineto{\pgfqpoint{1.408351in}{1.098572in}}%
\pgfpathclose%
\pgfusepath{fill}%
\end{pgfscope}%
\begin{pgfscope}%
\pgfpathrectangle{\pgfqpoint{0.100000in}{0.100000in}}{\pgfqpoint{3.007045in}{1.925000in}}%
\pgfusepath{clip}%
\pgfsetbuttcap%
\pgfsetmiterjoin%
\definecolor{currentfill}{rgb}{0.451765,0.767090,0.461207}%
\pgfsetfillcolor{currentfill}%
\pgfsetlinewidth{0.000000pt}%
\definecolor{currentstroke}{rgb}{0.000000,0.000000,0.000000}%
\pgfsetstrokecolor{currentstroke}%
\pgfsetstrokeopacity{0.000000}%
\pgfsetdash{}{0pt}%
\pgfpathmoveto{\pgfqpoint{1.753660in}{1.413597in}}%
\pgfpathlineto{\pgfqpoint{1.753560in}{1.390798in}}%
\pgfpathlineto{\pgfqpoint{1.730762in}{1.390911in}}%
\pgfpathlineto{\pgfqpoint{1.697951in}{1.391327in}}%
\pgfpathlineto{\pgfqpoint{1.699908in}{1.394517in}}%
\pgfpathlineto{\pgfqpoint{1.699433in}{1.400596in}}%
\pgfpathlineto{\pgfqpoint{1.702210in}{1.402774in}}%
\pgfpathlineto{\pgfqpoint{1.701304in}{1.411915in}}%
\pgfpathlineto{\pgfqpoint{1.698770in}{1.417142in}}%
\pgfpathlineto{\pgfqpoint{1.699167in}{1.423436in}}%
\pgfpathlineto{\pgfqpoint{1.695898in}{1.425999in}}%
\pgfpathlineto{\pgfqpoint{1.695709in}{1.430437in}}%
\pgfpathlineto{\pgfqpoint{1.702699in}{1.430317in}}%
\pgfpathlineto{\pgfqpoint{1.721870in}{1.430045in}}%
\pgfpathlineto{\pgfqpoint{1.731061in}{1.429926in}}%
\pgfpathlineto{\pgfqpoint{1.730861in}{1.413920in}}%
\pgfpathlineto{\pgfqpoint{1.753660in}{1.413597in}}%
\pgfpathclose%
\pgfusepath{fill}%
\end{pgfscope}%
\begin{pgfscope}%
\pgfpathrectangle{\pgfqpoint{0.100000in}{0.100000in}}{\pgfqpoint{3.007045in}{1.925000in}}%
\pgfusepath{clip}%
\pgfsetbuttcap%
\pgfsetmiterjoin%
\definecolor{currentfill}{rgb}{0.407843,0.745559,0.439677}%
\pgfsetfillcolor{currentfill}%
\pgfsetlinewidth{0.000000pt}%
\definecolor{currentstroke}{rgb}{0.000000,0.000000,0.000000}%
\pgfsetstrokecolor{currentstroke}%
\pgfsetstrokeopacity{0.000000}%
\pgfsetdash{}{0pt}%
\pgfpathmoveto{\pgfqpoint{2.867151in}{1.544262in}}%
\pgfpathlineto{\pgfqpoint{2.869055in}{1.532209in}}%
\pgfpathlineto{\pgfqpoint{2.871705in}{1.532829in}}%
\pgfpathlineto{\pgfqpoint{2.875745in}{1.523338in}}%
\pgfpathlineto{\pgfqpoint{2.879160in}{1.518963in}}%
\pgfpathlineto{\pgfqpoint{2.862374in}{1.515277in}}%
\pgfpathlineto{\pgfqpoint{2.831656in}{1.508759in}}%
\pgfpathlineto{\pgfqpoint{2.829097in}{1.520876in}}%
\pgfpathlineto{\pgfqpoint{2.824976in}{1.520081in}}%
\pgfpathlineto{\pgfqpoint{2.822507in}{1.532326in}}%
\pgfpathlineto{\pgfqpoint{2.828599in}{1.533431in}}%
\pgfpathlineto{\pgfqpoint{2.828637in}{1.543127in}}%
\pgfpathlineto{\pgfqpoint{2.835404in}{1.542568in}}%
\pgfpathlineto{\pgfqpoint{2.846989in}{1.546225in}}%
\pgfpathlineto{\pgfqpoint{2.850891in}{1.542069in}}%
\pgfpathlineto{\pgfqpoint{2.860753in}{1.546240in}}%
\pgfpathlineto{\pgfqpoint{2.867151in}{1.544262in}}%
\pgfpathclose%
\pgfusepath{fill}%
\end{pgfscope}%
\begin{pgfscope}%
\pgfpathrectangle{\pgfqpoint{0.100000in}{0.100000in}}{\pgfqpoint{3.007045in}{1.925000in}}%
\pgfusepath{clip}%
\pgfsetbuttcap%
\pgfsetmiterjoin%
\definecolor{currentfill}{rgb}{0.557324,0.816424,0.546959}%
\pgfsetfillcolor{currentfill}%
\pgfsetlinewidth{0.000000pt}%
\definecolor{currentstroke}{rgb}{0.000000,0.000000,0.000000}%
\pgfsetstrokecolor{currentstroke}%
\pgfsetstrokeopacity{0.000000}%
\pgfsetdash{}{0pt}%
\pgfpathmoveto{\pgfqpoint{1.147775in}{0.452101in}}%
\pgfpathlineto{\pgfqpoint{1.141912in}{0.462482in}}%
\pgfpathlineto{\pgfqpoint{1.134776in}{0.466628in}}%
\pgfpathlineto{\pgfqpoint{1.136483in}{0.471729in}}%
\pgfpathlineto{\pgfqpoint{1.136838in}{0.481546in}}%
\pgfpathlineto{\pgfqpoint{1.140040in}{0.487692in}}%
\pgfpathlineto{\pgfqpoint{1.147213in}{0.482417in}}%
\pgfpathlineto{\pgfqpoint{1.157196in}{0.484131in}}%
\pgfpathlineto{\pgfqpoint{1.159892in}{0.482747in}}%
\pgfpathlineto{\pgfqpoint{1.158835in}{0.470012in}}%
\pgfpathlineto{\pgfqpoint{1.155158in}{0.464696in}}%
\pgfpathlineto{\pgfqpoint{1.158256in}{0.454687in}}%
\pgfpathlineto{\pgfqpoint{1.157522in}{0.448171in}}%
\pgfpathlineto{\pgfqpoint{1.147775in}{0.452101in}}%
\pgfpathclose%
\pgfusepath{fill}%
\end{pgfscope}%
\begin{pgfscope}%
\pgfpathrectangle{\pgfqpoint{0.100000in}{0.100000in}}{\pgfqpoint{3.007045in}{1.925000in}}%
\pgfusepath{clip}%
\pgfsetbuttcap%
\pgfsetmiterjoin%
\definecolor{currentfill}{rgb}{0.769873,0.909296,0.742699}%
\pgfsetfillcolor{currentfill}%
\pgfsetlinewidth{0.000000pt}%
\definecolor{currentstroke}{rgb}{0.000000,0.000000,0.000000}%
\pgfsetstrokecolor{currentstroke}%
\pgfsetstrokeopacity{0.000000}%
\pgfsetdash{}{0pt}%
\pgfpathmoveto{\pgfqpoint{0.495869in}{0.921526in}}%
\pgfpathlineto{\pgfqpoint{0.488883in}{0.925335in}}%
\pgfpathlineto{\pgfqpoint{0.489033in}{0.933786in}}%
\pgfpathlineto{\pgfqpoint{0.495364in}{0.929362in}}%
\pgfpathlineto{\pgfqpoint{0.495869in}{0.921526in}}%
\pgfpathclose%
\pgfusepath{fill}%
\end{pgfscope}%
\begin{pgfscope}%
\pgfpathrectangle{\pgfqpoint{0.100000in}{0.100000in}}{\pgfqpoint{3.007045in}{1.925000in}}%
\pgfusepath{clip}%
\pgfsetbuttcap%
\pgfsetmiterjoin%
\definecolor{currentfill}{rgb}{0.769873,0.909296,0.742699}%
\pgfsetfillcolor{currentfill}%
\pgfsetlinewidth{0.000000pt}%
\definecolor{currentstroke}{rgb}{0.000000,0.000000,0.000000}%
\pgfsetstrokecolor{currentstroke}%
\pgfsetstrokeopacity{0.000000}%
\pgfsetdash{}{0pt}%
\pgfpathmoveto{\pgfqpoint{0.536458in}{0.917204in}}%
\pgfpathlineto{\pgfqpoint{0.528897in}{0.929641in}}%
\pgfpathlineto{\pgfqpoint{0.522315in}{0.935716in}}%
\pgfpathlineto{\pgfqpoint{0.515735in}{0.945831in}}%
\pgfpathlineto{\pgfqpoint{0.511433in}{0.949293in}}%
\pgfpathlineto{\pgfqpoint{0.504379in}{0.947346in}}%
\pgfpathlineto{\pgfqpoint{0.498832in}{0.951061in}}%
\pgfpathlineto{\pgfqpoint{0.501532in}{0.957039in}}%
\pgfpathlineto{\pgfqpoint{0.499515in}{0.967768in}}%
\pgfpathlineto{\pgfqpoint{0.496988in}{0.971830in}}%
\pgfpathlineto{\pgfqpoint{0.486075in}{0.974392in}}%
\pgfpathlineto{\pgfqpoint{0.481940in}{0.973863in}}%
\pgfpathlineto{\pgfqpoint{0.463162in}{0.988331in}}%
\pgfpathlineto{\pgfqpoint{0.461865in}{0.997033in}}%
\pgfpathlineto{\pgfqpoint{0.453370in}{1.006642in}}%
\pgfpathlineto{\pgfqpoint{0.456601in}{1.011141in}}%
\pgfpathlineto{\pgfqpoint{0.464444in}{1.039777in}}%
\pgfpathlineto{\pgfqpoint{0.472706in}{1.036029in}}%
\pgfpathlineto{\pgfqpoint{0.473306in}{1.031380in}}%
\pgfpathlineto{\pgfqpoint{0.515049in}{1.020412in}}%
\pgfpathlineto{\pgfqpoint{0.557933in}{1.009610in}}%
\pgfpathlineto{\pgfqpoint{0.573744in}{1.071996in}}%
\pgfpathlineto{\pgfqpoint{0.593532in}{1.066855in}}%
\pgfpathlineto{\pgfqpoint{0.654898in}{1.051840in}}%
\pgfpathlineto{\pgfqpoint{0.672324in}{1.047770in}}%
\pgfpathlineto{\pgfqpoint{0.677176in}{1.047687in}}%
\pgfpathlineto{\pgfqpoint{0.718691in}{0.983740in}}%
\pgfpathlineto{\pgfqpoint{0.716895in}{0.976496in}}%
\pgfpathlineto{\pgfqpoint{0.719619in}{0.967635in}}%
\pgfpathlineto{\pgfqpoint{0.723182in}{0.963163in}}%
\pgfpathlineto{\pgfqpoint{0.723368in}{0.955271in}}%
\pgfpathlineto{\pgfqpoint{0.726621in}{0.944753in}}%
\pgfpathlineto{\pgfqpoint{0.735192in}{0.932928in}}%
\pgfpathlineto{\pgfqpoint{0.740331in}{0.928918in}}%
\pgfpathlineto{\pgfqpoint{0.762651in}{0.922152in}}%
\pgfpathlineto{\pgfqpoint{0.767196in}{0.926377in}}%
\pgfpathlineto{\pgfqpoint{0.778392in}{0.924882in}}%
\pgfpathlineto{\pgfqpoint{0.774157in}{0.904247in}}%
\pgfpathlineto{\pgfqpoint{0.765905in}{0.863949in}}%
\pgfpathlineto{\pgfqpoint{0.732137in}{0.870987in}}%
\pgfpathlineto{\pgfqpoint{0.733331in}{0.876625in}}%
\pgfpathlineto{\pgfqpoint{0.716592in}{0.880125in}}%
\pgfpathlineto{\pgfqpoint{0.710499in}{0.852063in}}%
\pgfpathlineto{\pgfqpoint{0.689447in}{0.856914in}}%
\pgfpathlineto{\pgfqpoint{0.687845in}{0.862115in}}%
\pgfpathlineto{\pgfqpoint{0.691629in}{0.872688in}}%
\pgfpathlineto{\pgfqpoint{0.691055in}{0.881658in}}%
\pgfpathlineto{\pgfqpoint{0.673193in}{0.887262in}}%
\pgfpathlineto{\pgfqpoint{0.618157in}{0.899792in}}%
\pgfpathlineto{\pgfqpoint{0.567402in}{0.912130in}}%
\pgfpathlineto{\pgfqpoint{0.556193in}{0.915257in}}%
\pgfpathlineto{\pgfqpoint{0.550214in}{0.920766in}}%
\pgfpathlineto{\pgfqpoint{0.543010in}{0.923557in}}%
\pgfpathlineto{\pgfqpoint{0.538502in}{0.921211in}}%
\pgfpathlineto{\pgfqpoint{0.536458in}{0.917204in}}%
\pgfpathclose%
\pgfusepath{fill}%
\end{pgfscope}%
\begin{pgfscope}%
\pgfpathrectangle{\pgfqpoint{0.100000in}{0.100000in}}{\pgfqpoint{3.007045in}{1.925000in}}%
\pgfusepath{clip}%
\pgfsetbuttcap%
\pgfsetmiterjoin%
\definecolor{currentfill}{rgb}{0.557324,0.816424,0.546959}%
\pgfsetfillcolor{currentfill}%
\pgfsetlinewidth{0.000000pt}%
\definecolor{currentstroke}{rgb}{0.000000,0.000000,0.000000}%
\pgfsetstrokecolor{currentstroke}%
\pgfsetstrokeopacity{0.000000}%
\pgfsetdash{}{0pt}%
\pgfpathmoveto{\pgfqpoint{2.241536in}{0.660624in}}%
\pgfpathlineto{\pgfqpoint{2.252934in}{0.661782in}}%
\pgfpathlineto{\pgfqpoint{2.252110in}{0.670372in}}%
\pgfpathlineto{\pgfqpoint{2.254972in}{0.670665in}}%
\pgfpathlineto{\pgfqpoint{2.253069in}{0.691091in}}%
\pgfpathlineto{\pgfqpoint{2.254880in}{0.697001in}}%
\pgfpathlineto{\pgfqpoint{2.260621in}{0.697611in}}%
\pgfpathlineto{\pgfqpoint{2.261248in}{0.691899in}}%
\pgfpathlineto{\pgfqpoint{2.278425in}{0.693743in}}%
\pgfpathlineto{\pgfqpoint{2.277794in}{0.699248in}}%
\pgfpathlineto{\pgfqpoint{2.283592in}{0.699581in}}%
\pgfpathlineto{\pgfqpoint{2.284645in}{0.694345in}}%
\pgfpathlineto{\pgfqpoint{2.289927in}{0.694879in}}%
\pgfpathlineto{\pgfqpoint{2.290454in}{0.689136in}}%
\pgfpathlineto{\pgfqpoint{2.297933in}{0.689896in}}%
\pgfpathlineto{\pgfqpoint{2.298312in}{0.683587in}}%
\pgfpathlineto{\pgfqpoint{2.294624in}{0.672055in}}%
\pgfpathlineto{\pgfqpoint{2.272250in}{0.669798in}}%
\pgfpathlineto{\pgfqpoint{2.270062in}{0.663726in}}%
\pgfpathlineto{\pgfqpoint{2.273936in}{0.628314in}}%
\pgfpathlineto{\pgfqpoint{2.245494in}{0.625658in}}%
\pgfpathlineto{\pgfqpoint{2.244783in}{0.625988in}}%
\pgfpathlineto{\pgfqpoint{2.241536in}{0.660624in}}%
\pgfpathclose%
\pgfusepath{fill}%
\end{pgfscope}%
\begin{pgfscope}%
\pgfpathrectangle{\pgfqpoint{0.100000in}{0.100000in}}{\pgfqpoint{3.007045in}{1.925000in}}%
\pgfusepath{clip}%
\pgfsetbuttcap%
\pgfsetmiterjoin%
\definecolor{currentfill}{rgb}{0.760523,0.905359,0.733595}%
\pgfsetfillcolor{currentfill}%
\pgfsetlinewidth{0.000000pt}%
\definecolor{currentstroke}{rgb}{0.000000,0.000000,0.000000}%
\pgfsetstrokecolor{currentstroke}%
\pgfsetstrokeopacity{0.000000}%
\pgfsetdash{}{0pt}%
\pgfpathmoveto{\pgfqpoint{2.033514in}{0.905154in}}%
\pgfpathlineto{\pgfqpoint{2.039932in}{0.907365in}}%
\pgfpathlineto{\pgfqpoint{2.038540in}{0.912838in}}%
\pgfpathlineto{\pgfqpoint{2.045514in}{0.912506in}}%
\pgfpathlineto{\pgfqpoint{2.051337in}{0.919577in}}%
\pgfpathlineto{\pgfqpoint{2.061506in}{0.919581in}}%
\pgfpathlineto{\pgfqpoint{2.068214in}{0.916226in}}%
\pgfpathlineto{\pgfqpoint{2.069407in}{0.904987in}}%
\pgfpathlineto{\pgfqpoint{2.085014in}{0.905571in}}%
\pgfpathlineto{\pgfqpoint{2.085939in}{0.878995in}}%
\pgfpathlineto{\pgfqpoint{2.077608in}{0.878428in}}%
\pgfpathlineto{\pgfqpoint{2.078219in}{0.868579in}}%
\pgfpathlineto{\pgfqpoint{2.081112in}{0.868758in}}%
\pgfpathlineto{\pgfqpoint{2.082243in}{0.851523in}}%
\pgfpathlineto{\pgfqpoint{2.085442in}{0.845945in}}%
\pgfpathlineto{\pgfqpoint{2.072679in}{0.849028in}}%
\pgfpathlineto{\pgfqpoint{2.033456in}{0.846814in}}%
\pgfpathlineto{\pgfqpoint{2.027867in}{0.843557in}}%
\pgfpathlineto{\pgfqpoint{2.028291in}{0.837950in}}%
\pgfpathlineto{\pgfqpoint{2.022978in}{0.837629in}}%
\pgfpathlineto{\pgfqpoint{2.013407in}{0.843839in}}%
\pgfpathlineto{\pgfqpoint{2.012143in}{0.849788in}}%
\pgfpathlineto{\pgfqpoint{2.018227in}{0.853487in}}%
\pgfpathlineto{\pgfqpoint{2.020371in}{0.864609in}}%
\pgfpathlineto{\pgfqpoint{2.029442in}{0.869945in}}%
\pgfpathlineto{\pgfqpoint{2.026701in}{0.874367in}}%
\pgfpathlineto{\pgfqpoint{2.030962in}{0.877859in}}%
\pgfpathlineto{\pgfqpoint{2.032840in}{0.884060in}}%
\pgfpathlineto{\pgfqpoint{2.036782in}{0.884180in}}%
\pgfpathlineto{\pgfqpoint{2.037839in}{0.890956in}}%
\pgfpathlineto{\pgfqpoint{2.034157in}{0.893004in}}%
\pgfpathlineto{\pgfqpoint{2.037206in}{0.901885in}}%
\pgfpathlineto{\pgfqpoint{2.033514in}{0.905154in}}%
\pgfpathclose%
\pgfusepath{fill}%
\end{pgfscope}%
\begin{pgfscope}%
\pgfpathrectangle{\pgfqpoint{0.100000in}{0.100000in}}{\pgfqpoint{3.007045in}{1.925000in}}%
\pgfusepath{clip}%
\pgfsetbuttcap%
\pgfsetmiterjoin%
\definecolor{currentfill}{rgb}{0.562860,0.819008,0.551511}%
\pgfsetfillcolor{currentfill}%
\pgfsetlinewidth{0.000000pt}%
\definecolor{currentstroke}{rgb}{0.000000,0.000000,0.000000}%
\pgfsetstrokecolor{currentstroke}%
\pgfsetstrokeopacity{0.000000}%
\pgfsetdash{}{0pt}%
\pgfpathmoveto{\pgfqpoint{2.372484in}{1.056245in}}%
\pgfpathlineto{\pgfqpoint{2.363870in}{1.060375in}}%
\pgfpathlineto{\pgfqpoint{2.358324in}{1.060460in}}%
\pgfpathlineto{\pgfqpoint{2.352145in}{1.068229in}}%
\pgfpathlineto{\pgfqpoint{2.350488in}{1.072384in}}%
\pgfpathlineto{\pgfqpoint{2.353799in}{1.080602in}}%
\pgfpathlineto{\pgfqpoint{2.358589in}{1.084899in}}%
\pgfpathlineto{\pgfqpoint{2.362396in}{1.092621in}}%
\pgfpathlineto{\pgfqpoint{2.367300in}{1.093259in}}%
\pgfpathlineto{\pgfqpoint{2.371064in}{1.101491in}}%
\pgfpathlineto{\pgfqpoint{2.370833in}{1.107548in}}%
\pgfpathlineto{\pgfqpoint{2.379152in}{1.110286in}}%
\pgfpathlineto{\pgfqpoint{2.383396in}{1.107611in}}%
\pgfpathlineto{\pgfqpoint{2.392249in}{1.108544in}}%
\pgfpathlineto{\pgfqpoint{2.396141in}{1.102157in}}%
\pgfpathlineto{\pgfqpoint{2.389841in}{1.100458in}}%
\pgfpathlineto{\pgfqpoint{2.381551in}{1.091127in}}%
\pgfpathlineto{\pgfqpoint{2.382872in}{1.085587in}}%
\pgfpathlineto{\pgfqpoint{2.391489in}{1.082969in}}%
\pgfpathlineto{\pgfqpoint{2.400437in}{1.076137in}}%
\pgfpathlineto{\pgfqpoint{2.392967in}{1.076431in}}%
\pgfpathlineto{\pgfqpoint{2.392156in}{1.068534in}}%
\pgfpathlineto{\pgfqpoint{2.384647in}{1.067115in}}%
\pgfpathlineto{\pgfqpoint{2.378157in}{1.062296in}}%
\pgfpathlineto{\pgfqpoint{2.375741in}{1.064118in}}%
\pgfpathlineto{\pgfqpoint{2.370868in}{1.060399in}}%
\pgfpathlineto{\pgfqpoint{2.372484in}{1.056245in}}%
\pgfpathclose%
\pgfusepath{fill}%
\end{pgfscope}%
\begin{pgfscope}%
\pgfpathrectangle{\pgfqpoint{0.100000in}{0.100000in}}{\pgfqpoint{3.007045in}{1.925000in}}%
\pgfusepath{clip}%
\pgfsetbuttcap%
\pgfsetmiterjoin%
\definecolor{currentfill}{rgb}{0.039293,0.461130,0.200615}%
\pgfsetfillcolor{currentfill}%
\pgfsetlinewidth{0.000000pt}%
\definecolor{currentstroke}{rgb}{0.000000,0.000000,0.000000}%
\pgfsetstrokecolor{currentstroke}%
\pgfsetstrokeopacity{0.000000}%
\pgfsetdash{}{0pt}%
\pgfpathmoveto{\pgfqpoint{2.623105in}{0.353228in}}%
\pgfpathlineto{\pgfqpoint{2.628250in}{0.323017in}}%
\pgfpathlineto{\pgfqpoint{2.599950in}{0.318474in}}%
\pgfpathlineto{\pgfqpoint{2.594969in}{0.321852in}}%
\pgfpathlineto{\pgfqpoint{2.580427in}{0.321862in}}%
\pgfpathlineto{\pgfqpoint{2.573972in}{0.324804in}}%
\pgfpathlineto{\pgfqpoint{2.569745in}{0.333328in}}%
\pgfpathlineto{\pgfqpoint{2.565657in}{0.347157in}}%
\pgfpathlineto{\pgfqpoint{2.562738in}{0.352174in}}%
\pgfpathlineto{\pgfqpoint{2.555247in}{0.358065in}}%
\pgfpathlineto{\pgfqpoint{2.547708in}{0.358122in}}%
\pgfpathlineto{\pgfqpoint{2.543685in}{0.367836in}}%
\pgfpathlineto{\pgfqpoint{2.546700in}{0.369561in}}%
\pgfpathlineto{\pgfqpoint{2.547508in}{0.376728in}}%
\pgfpathlineto{\pgfqpoint{2.577000in}{0.380965in}}%
\pgfpathlineto{\pgfqpoint{2.579681in}{0.363910in}}%
\pgfpathlineto{\pgfqpoint{2.597057in}{0.366782in}}%
\pgfpathlineto{\pgfqpoint{2.599929in}{0.349210in}}%
\pgfpathlineto{\pgfqpoint{2.623105in}{0.353228in}}%
\pgfpathclose%
\pgfusepath{fill}%
\end{pgfscope}%
\begin{pgfscope}%
\pgfpathrectangle{\pgfqpoint{0.100000in}{0.100000in}}{\pgfqpoint{3.007045in}{1.925000in}}%
\pgfusepath{clip}%
\pgfsetbuttcap%
\pgfsetmiterjoin%
\definecolor{currentfill}{rgb}{0.760523,0.905359,0.733595}%
\pgfsetfillcolor{currentfill}%
\pgfsetlinewidth{0.000000pt}%
\definecolor{currentstroke}{rgb}{0.000000,0.000000,0.000000}%
\pgfsetstrokecolor{currentstroke}%
\pgfsetstrokeopacity{0.000000}%
\pgfsetdash{}{0pt}%
\pgfpathmoveto{\pgfqpoint{0.440258in}{1.107424in}}%
\pgfpathlineto{\pgfqpoint{0.406868in}{1.117231in}}%
\pgfpathlineto{\pgfqpoint{0.382061in}{1.124859in}}%
\pgfpathlineto{\pgfqpoint{0.377822in}{1.132382in}}%
\pgfpathlineto{\pgfqpoint{0.377753in}{1.140606in}}%
\pgfpathlineto{\pgfqpoint{0.375122in}{1.142276in}}%
\pgfpathlineto{\pgfqpoint{0.370888in}{1.155493in}}%
\pgfpathlineto{\pgfqpoint{0.366684in}{1.159674in}}%
\pgfpathlineto{\pgfqpoint{0.363995in}{1.165891in}}%
\pgfpathlineto{\pgfqpoint{0.365227in}{1.184145in}}%
\pgfpathlineto{\pgfqpoint{0.371166in}{1.184004in}}%
\pgfpathlineto{\pgfqpoint{0.374775in}{1.187299in}}%
\pgfpathlineto{\pgfqpoint{0.378870in}{1.195364in}}%
\pgfpathlineto{\pgfqpoint{0.377503in}{1.204628in}}%
\pgfpathlineto{\pgfqpoint{0.365785in}{1.209702in}}%
\pgfpathlineto{\pgfqpoint{0.361242in}{1.215875in}}%
\pgfpathlineto{\pgfqpoint{0.359301in}{1.222111in}}%
\pgfpathlineto{\pgfqpoint{0.359639in}{1.227533in}}%
\pgfpathlineto{\pgfqpoint{0.368463in}{1.226772in}}%
\pgfpathlineto{\pgfqpoint{0.368846in}{1.236860in}}%
\pgfpathlineto{\pgfqpoint{0.370754in}{1.240933in}}%
\pgfpathlineto{\pgfqpoint{0.377116in}{1.242201in}}%
\pgfpathlineto{\pgfqpoint{0.384479in}{1.238341in}}%
\pgfpathlineto{\pgfqpoint{0.388058in}{1.239323in}}%
\pgfpathlineto{\pgfqpoint{0.407661in}{1.233230in}}%
\pgfpathlineto{\pgfqpoint{0.408920in}{1.226474in}}%
\pgfpathlineto{\pgfqpoint{0.404848in}{1.220444in}}%
\pgfpathlineto{\pgfqpoint{0.404869in}{1.211825in}}%
\pgfpathlineto{\pgfqpoint{0.413366in}{1.209180in}}%
\pgfpathlineto{\pgfqpoint{0.413081in}{1.203149in}}%
\pgfpathlineto{\pgfqpoint{0.409170in}{1.194540in}}%
\pgfpathlineto{\pgfqpoint{0.412185in}{1.187494in}}%
\pgfpathlineto{\pgfqpoint{0.421673in}{1.178083in}}%
\pgfpathlineto{\pgfqpoint{0.433028in}{1.157462in}}%
\pgfpathlineto{\pgfqpoint{0.430464in}{1.147237in}}%
\pgfpathlineto{\pgfqpoint{0.426809in}{1.147328in}}%
\pgfpathlineto{\pgfqpoint{0.423734in}{1.137658in}}%
\pgfpathlineto{\pgfqpoint{0.428309in}{1.126977in}}%
\pgfpathlineto{\pgfqpoint{0.432260in}{1.122040in}}%
\pgfpathlineto{\pgfqpoint{0.435906in}{1.120684in}}%
\pgfpathlineto{\pgfqpoint{0.440258in}{1.107424in}}%
\pgfpathclose%
\pgfusepath{fill}%
\end{pgfscope}%
\begin{pgfscope}%
\pgfpathrectangle{\pgfqpoint{0.100000in}{0.100000in}}{\pgfqpoint{3.007045in}{1.925000in}}%
\pgfusepath{clip}%
\pgfsetbuttcap%
\pgfsetmiterjoin%
\definecolor{currentfill}{rgb}{0.845905,0.939931,0.822822}%
\pgfsetfillcolor{currentfill}%
\pgfsetlinewidth{0.000000pt}%
\definecolor{currentstroke}{rgb}{0.000000,0.000000,0.000000}%
\pgfsetstrokecolor{currentstroke}%
\pgfsetstrokeopacity{0.000000}%
\pgfsetdash{}{0pt}%
\pgfpathmoveto{\pgfqpoint{0.966052in}{0.192399in}}%
\pgfpathlineto{\pgfqpoint{0.964721in}{0.193135in}}%
\pgfpathlineto{\pgfqpoint{0.963493in}{0.195565in}}%
\pgfpathlineto{\pgfqpoint{0.958655in}{0.195219in}}%
\pgfpathlineto{\pgfqpoint{0.957247in}{0.193413in}}%
\pgfpathlineto{\pgfqpoint{0.955692in}{0.192812in}}%
\pgfpathlineto{\pgfqpoint{0.954245in}{0.194809in}}%
\pgfpathlineto{\pgfqpoint{0.951030in}{0.194631in}}%
\pgfpathlineto{\pgfqpoint{0.952673in}{0.196998in}}%
\pgfpathlineto{\pgfqpoint{0.954685in}{0.197089in}}%
\pgfpathlineto{\pgfqpoint{0.954041in}{0.199019in}}%
\pgfpathlineto{\pgfqpoint{0.954053in}{0.201532in}}%
\pgfpathlineto{\pgfqpoint{0.954782in}{0.203274in}}%
\pgfpathlineto{\pgfqpoint{0.953594in}{0.205042in}}%
\pgfpathlineto{\pgfqpoint{0.954325in}{0.206175in}}%
\pgfpathlineto{\pgfqpoint{0.956732in}{0.206941in}}%
\pgfpathlineto{\pgfqpoint{0.957233in}{0.209806in}}%
\pgfpathlineto{\pgfqpoint{0.954751in}{0.209244in}}%
\pgfpathlineto{\pgfqpoint{0.954872in}{0.213244in}}%
\pgfpathlineto{\pgfqpoint{0.953569in}{0.215767in}}%
\pgfpathlineto{\pgfqpoint{0.954140in}{0.218651in}}%
\pgfpathlineto{\pgfqpoint{0.953788in}{0.220128in}}%
\pgfpathlineto{\pgfqpoint{0.952256in}{0.220622in}}%
\pgfpathlineto{\pgfqpoint{0.949173in}{0.223004in}}%
\pgfpathlineto{\pgfqpoint{0.948744in}{0.225034in}}%
\pgfpathlineto{\pgfqpoint{0.949609in}{0.226010in}}%
\pgfpathlineto{\pgfqpoint{0.948660in}{0.227534in}}%
\pgfpathlineto{\pgfqpoint{0.949321in}{0.229307in}}%
\pgfpathlineto{\pgfqpoint{0.949075in}{0.234015in}}%
\pgfpathlineto{\pgfqpoint{0.949908in}{0.234459in}}%
\pgfpathlineto{\pgfqpoint{0.950927in}{0.237714in}}%
\pgfpathlineto{\pgfqpoint{0.949190in}{0.237528in}}%
\pgfpathlineto{\pgfqpoint{0.949382in}{0.239911in}}%
\pgfpathlineto{\pgfqpoint{0.951032in}{0.247085in}}%
\pgfpathlineto{\pgfqpoint{0.951672in}{0.251760in}}%
\pgfpathlineto{\pgfqpoint{0.949767in}{0.251172in}}%
\pgfpathlineto{\pgfqpoint{0.950024in}{0.249097in}}%
\pgfpathlineto{\pgfqpoint{0.948795in}{0.248793in}}%
\pgfpathlineto{\pgfqpoint{0.948847in}{0.246567in}}%
\pgfpathlineto{\pgfqpoint{0.947286in}{0.241526in}}%
\pgfpathlineto{\pgfqpoint{0.946946in}{0.236340in}}%
\pgfpathlineto{\pgfqpoint{0.945783in}{0.232618in}}%
\pgfpathlineto{\pgfqpoint{0.944897in}{0.227958in}}%
\pgfpathlineto{\pgfqpoint{0.942843in}{0.225828in}}%
\pgfpathlineto{\pgfqpoint{0.941004in}{0.228956in}}%
\pgfpathlineto{\pgfqpoint{0.941264in}{0.231409in}}%
\pgfpathlineto{\pgfqpoint{0.939615in}{0.232446in}}%
\pgfpathlineto{\pgfqpoint{0.935598in}{0.233694in}}%
\pgfpathlineto{\pgfqpoint{0.937893in}{0.236489in}}%
\pgfpathlineto{\pgfqpoint{0.938389in}{0.240874in}}%
\pgfpathlineto{\pgfqpoint{0.938282in}{0.243236in}}%
\pgfpathlineto{\pgfqpoint{0.935953in}{0.243285in}}%
\pgfpathlineto{\pgfqpoint{0.934974in}{0.245815in}}%
\pgfpathlineto{\pgfqpoint{0.933942in}{0.246488in}}%
\pgfpathlineto{\pgfqpoint{0.934379in}{0.249543in}}%
\pgfpathlineto{\pgfqpoint{0.931764in}{0.250440in}}%
\pgfpathlineto{\pgfqpoint{0.930132in}{0.252504in}}%
\pgfpathlineto{\pgfqpoint{0.929218in}{0.251104in}}%
\pgfpathlineto{\pgfqpoint{0.931463in}{0.248759in}}%
\pgfpathlineto{\pgfqpoint{0.931593in}{0.245629in}}%
\pgfpathlineto{\pgfqpoint{0.933934in}{0.242455in}}%
\pgfpathlineto{\pgfqpoint{0.931815in}{0.241739in}}%
\pgfpathlineto{\pgfqpoint{0.932820in}{0.240335in}}%
\pgfpathlineto{\pgfqpoint{0.934512in}{0.240731in}}%
\pgfpathlineto{\pgfqpoint{0.934225in}{0.234374in}}%
\pgfpathlineto{\pgfqpoint{0.930700in}{0.234187in}}%
\pgfpathlineto{\pgfqpoint{0.931316in}{0.236252in}}%
\pgfpathlineto{\pgfqpoint{0.928812in}{0.235626in}}%
\pgfpathlineto{\pgfqpoint{0.925421in}{0.233892in}}%
\pgfpathlineto{\pgfqpoint{0.925428in}{0.237038in}}%
\pgfpathlineto{\pgfqpoint{0.924460in}{0.239424in}}%
\pgfpathlineto{\pgfqpoint{0.922320in}{0.240051in}}%
\pgfpathlineto{\pgfqpoint{0.919294in}{0.247052in}}%
\pgfpathlineto{\pgfqpoint{0.919630in}{0.248516in}}%
\pgfpathlineto{\pgfqpoint{0.917767in}{0.252248in}}%
\pgfpathlineto{\pgfqpoint{0.917999in}{0.256789in}}%
\pgfpathlineto{\pgfqpoint{0.916349in}{0.261630in}}%
\pgfpathlineto{\pgfqpoint{0.913473in}{0.265334in}}%
\pgfpathlineto{\pgfqpoint{0.912193in}{0.268097in}}%
\pgfpathlineto{\pgfqpoint{0.908339in}{0.273326in}}%
\pgfpathlineto{\pgfqpoint{0.904581in}{0.280446in}}%
\pgfpathlineto{\pgfqpoint{0.907013in}{0.280057in}}%
\pgfpathlineto{\pgfqpoint{0.910459in}{0.281718in}}%
\pgfpathlineto{\pgfqpoint{0.910823in}{0.286283in}}%
\pgfpathlineto{\pgfqpoint{0.911674in}{0.287962in}}%
\pgfpathlineto{\pgfqpoint{0.907295in}{0.286400in}}%
\pgfpathlineto{\pgfqpoint{0.905838in}{0.288835in}}%
\pgfpathlineto{\pgfqpoint{0.902783in}{0.288080in}}%
\pgfpathlineto{\pgfqpoint{0.901946in}{0.286374in}}%
\pgfpathlineto{\pgfqpoint{0.897692in}{0.289155in}}%
\pgfpathlineto{\pgfqpoint{0.895195in}{0.292046in}}%
\pgfpathlineto{\pgfqpoint{0.901489in}{0.303399in}}%
\pgfpathlineto{\pgfqpoint{0.905283in}{0.298997in}}%
\pgfpathlineto{\pgfqpoint{0.907695in}{0.299493in}}%
\pgfpathlineto{\pgfqpoint{0.910640in}{0.295087in}}%
\pgfpathlineto{\pgfqpoint{0.915744in}{0.296619in}}%
\pgfpathlineto{\pgfqpoint{0.920918in}{0.294241in}}%
\pgfpathlineto{\pgfqpoint{0.921920in}{0.292840in}}%
\pgfpathlineto{\pgfqpoint{0.917937in}{0.288586in}}%
\pgfpathlineto{\pgfqpoint{0.918519in}{0.285751in}}%
\pgfpathlineto{\pgfqpoint{0.921178in}{0.282019in}}%
\pgfpathlineto{\pgfqpoint{0.920601in}{0.278454in}}%
\pgfpathlineto{\pgfqpoint{0.925898in}{0.261651in}}%
\pgfpathlineto{\pgfqpoint{0.924680in}{0.256502in}}%
\pgfpathlineto{\pgfqpoint{0.923189in}{0.253738in}}%
\pgfpathlineto{\pgfqpoint{0.924096in}{0.253355in}}%
\pgfpathlineto{\pgfqpoint{0.927167in}{0.254345in}}%
\pgfpathlineto{\pgfqpoint{0.932410in}{0.254987in}}%
\pgfpathlineto{\pgfqpoint{0.936447in}{0.254101in}}%
\pgfpathlineto{\pgfqpoint{0.938989in}{0.256067in}}%
\pgfpathlineto{\pgfqpoint{0.940309in}{0.258917in}}%
\pgfpathlineto{\pgfqpoint{0.942720in}{0.259083in}}%
\pgfpathlineto{\pgfqpoint{0.946116in}{0.262219in}}%
\pgfpathlineto{\pgfqpoint{0.948733in}{0.261499in}}%
\pgfpathlineto{\pgfqpoint{0.954529in}{0.261913in}}%
\pgfpathlineto{\pgfqpoint{0.955655in}{0.261214in}}%
\pgfpathlineto{\pgfqpoint{0.957563in}{0.255964in}}%
\pgfpathlineto{\pgfqpoint{0.957437in}{0.253580in}}%
\pgfpathlineto{\pgfqpoint{0.955678in}{0.251817in}}%
\pgfpathlineto{\pgfqpoint{0.955644in}{0.247945in}}%
\pgfpathlineto{\pgfqpoint{0.958149in}{0.245787in}}%
\pgfpathlineto{\pgfqpoint{0.958180in}{0.243117in}}%
\pgfpathlineto{\pgfqpoint{0.959385in}{0.241184in}}%
\pgfpathlineto{\pgfqpoint{0.958932in}{0.238807in}}%
\pgfpathlineto{\pgfqpoint{0.959232in}{0.235056in}}%
\pgfpathlineto{\pgfqpoint{0.962530in}{0.230044in}}%
\pgfpathlineto{\pgfqpoint{0.962897in}{0.225216in}}%
\pgfpathlineto{\pgfqpoint{0.964318in}{0.222393in}}%
\pgfpathlineto{\pgfqpoint{0.963412in}{0.220631in}}%
\pgfpathlineto{\pgfqpoint{0.964300in}{0.213106in}}%
\pgfpathlineto{\pgfqpoint{0.964057in}{0.210474in}}%
\pgfpathlineto{\pgfqpoint{0.964942in}{0.205541in}}%
\pgfpathlineto{\pgfqpoint{0.965301in}{0.196768in}}%
\pgfpathlineto{\pgfqpoint{0.966052in}{0.192399in}}%
\pgfpathclose%
\pgfusepath{fill}%
\end{pgfscope}%
\begin{pgfscope}%
\pgfpathrectangle{\pgfqpoint{0.100000in}{0.100000in}}{\pgfqpoint{3.007045in}{1.925000in}}%
\pgfusepath{clip}%
\pgfsetbuttcap%
\pgfsetmiterjoin%
\definecolor{currentfill}{rgb}{0.845905,0.939931,0.822822}%
\pgfsetfillcolor{currentfill}%
\pgfsetlinewidth{0.000000pt}%
\definecolor{currentstroke}{rgb}{0.000000,0.000000,0.000000}%
\pgfsetstrokecolor{currentstroke}%
\pgfsetstrokeopacity{0.000000}%
\pgfsetdash{}{0pt}%
\pgfpathmoveto{\pgfqpoint{0.936626in}{0.197068in}}%
\pgfpathlineto{\pgfqpoint{0.938877in}{0.201656in}}%
\pgfpathlineto{\pgfqpoint{0.940558in}{0.202970in}}%
\pgfpathlineto{\pgfqpoint{0.942011in}{0.205039in}}%
\pgfpathlineto{\pgfqpoint{0.942374in}{0.211894in}}%
\pgfpathlineto{\pgfqpoint{0.944008in}{0.216194in}}%
\pgfpathlineto{\pgfqpoint{0.944696in}{0.219988in}}%
\pgfpathlineto{\pgfqpoint{0.945984in}{0.221447in}}%
\pgfpathlineto{\pgfqpoint{0.945425in}{0.223366in}}%
\pgfpathlineto{\pgfqpoint{0.945878in}{0.227745in}}%
\pgfpathlineto{\pgfqpoint{0.947540in}{0.225877in}}%
\pgfpathlineto{\pgfqpoint{0.947168in}{0.222342in}}%
\pgfpathlineto{\pgfqpoint{0.948688in}{0.221616in}}%
\pgfpathlineto{\pgfqpoint{0.952311in}{0.218681in}}%
\pgfpathlineto{\pgfqpoint{0.952653in}{0.214058in}}%
\pgfpathlineto{\pgfqpoint{0.952346in}{0.208455in}}%
\pgfpathlineto{\pgfqpoint{0.950987in}{0.209391in}}%
\pgfpathlineto{\pgfqpoint{0.951266in}{0.216399in}}%
\pgfpathlineto{\pgfqpoint{0.950054in}{0.216109in}}%
\pgfpathlineto{\pgfqpoint{0.948419in}{0.213095in}}%
\pgfpathlineto{\pgfqpoint{0.949973in}{0.212546in}}%
\pgfpathlineto{\pgfqpoint{0.949906in}{0.210661in}}%
\pgfpathlineto{\pgfqpoint{0.948880in}{0.207820in}}%
\pgfpathlineto{\pgfqpoint{0.950125in}{0.207070in}}%
\pgfpathlineto{\pgfqpoint{0.949683in}{0.203304in}}%
\pgfpathlineto{\pgfqpoint{0.947074in}{0.204157in}}%
\pgfpathlineto{\pgfqpoint{0.948879in}{0.200841in}}%
\pgfpathlineto{\pgfqpoint{0.947064in}{0.199809in}}%
\pgfpathlineto{\pgfqpoint{0.946264in}{0.200919in}}%
\pgfpathlineto{\pgfqpoint{0.944697in}{0.199414in}}%
\pgfpathlineto{\pgfqpoint{0.942282in}{0.199077in}}%
\pgfpathlineto{\pgfqpoint{0.939613in}{0.197250in}}%
\pgfpathlineto{\pgfqpoint{0.936626in}{0.197068in}}%
\pgfpathclose%
\pgfusepath{fill}%
\end{pgfscope}%
\begin{pgfscope}%
\pgfpathrectangle{\pgfqpoint{0.100000in}{0.100000in}}{\pgfqpoint{3.007045in}{1.925000in}}%
\pgfusepath{clip}%
\pgfsetbuttcap%
\pgfsetmiterjoin%
\definecolor{currentfill}{rgb}{0.845905,0.939931,0.822822}%
\pgfsetfillcolor{currentfill}%
\pgfsetlinewidth{0.000000pt}%
\definecolor{currentstroke}{rgb}{0.000000,0.000000,0.000000}%
\pgfsetstrokecolor{currentstroke}%
\pgfsetstrokeopacity{0.000000}%
\pgfsetdash{}{0pt}%
\pgfpathmoveto{\pgfqpoint{0.924230in}{0.224566in}}%
\pgfpathlineto{\pgfqpoint{0.924247in}{0.226972in}}%
\pgfpathlineto{\pgfqpoint{0.925304in}{0.230453in}}%
\pgfpathlineto{\pgfqpoint{0.928832in}{0.230375in}}%
\pgfpathlineto{\pgfqpoint{0.930194in}{0.229696in}}%
\pgfpathlineto{\pgfqpoint{0.931247in}{0.231217in}}%
\pgfpathlineto{\pgfqpoint{0.932446in}{0.229785in}}%
\pgfpathlineto{\pgfqpoint{0.935561in}{0.230884in}}%
\pgfpathlineto{\pgfqpoint{0.937650in}{0.227795in}}%
\pgfpathlineto{\pgfqpoint{0.938131in}{0.224744in}}%
\pgfpathlineto{\pgfqpoint{0.941809in}{0.222536in}}%
\pgfpathlineto{\pgfqpoint{0.943001in}{0.220916in}}%
\pgfpathlineto{\pgfqpoint{0.942398in}{0.218706in}}%
\pgfpathlineto{\pgfqpoint{0.941010in}{0.217751in}}%
\pgfpathlineto{\pgfqpoint{0.939636in}{0.215477in}}%
\pgfpathlineto{\pgfqpoint{0.936886in}{0.216547in}}%
\pgfpathlineto{\pgfqpoint{0.934778in}{0.220682in}}%
\pgfpathlineto{\pgfqpoint{0.933152in}{0.222753in}}%
\pgfpathlineto{\pgfqpoint{0.931898in}{0.225478in}}%
\pgfpathlineto{\pgfqpoint{0.928775in}{0.222468in}}%
\pgfpathlineto{\pgfqpoint{0.924230in}{0.224566in}}%
\pgfpathclose%
\pgfusepath{fill}%
\end{pgfscope}%
\begin{pgfscope}%
\pgfpathrectangle{\pgfqpoint{0.100000in}{0.100000in}}{\pgfqpoint{3.007045in}{1.925000in}}%
\pgfusepath{clip}%
\pgfsetbuttcap%
\pgfsetmiterjoin%
\definecolor{currentfill}{rgb}{0.681046,0.871895,0.656209}%
\pgfsetfillcolor{currentfill}%
\pgfsetlinewidth{0.000000pt}%
\definecolor{currentstroke}{rgb}{0.000000,0.000000,0.000000}%
\pgfsetstrokecolor{currentstroke}%
\pgfsetstrokeopacity{0.000000}%
\pgfsetdash{}{0pt}%
\pgfpathmoveto{\pgfqpoint{2.803517in}{1.405975in}}%
\pgfpathlineto{\pgfqpoint{2.803333in}{1.417011in}}%
\pgfpathlineto{\pgfqpoint{2.822125in}{1.424128in}}%
\pgfpathlineto{\pgfqpoint{2.824695in}{1.409031in}}%
\pgfpathlineto{\pgfqpoint{2.829259in}{1.404526in}}%
\pgfpathlineto{\pgfqpoint{2.819156in}{1.394508in}}%
\pgfpathlineto{\pgfqpoint{2.824324in}{1.387968in}}%
\pgfpathlineto{\pgfqpoint{2.826877in}{1.382803in}}%
\pgfpathlineto{\pgfqpoint{2.833070in}{1.385813in}}%
\pgfpathlineto{\pgfqpoint{2.846367in}{1.387920in}}%
\pgfpathlineto{\pgfqpoint{2.853875in}{1.393696in}}%
\pgfpathlineto{\pgfqpoint{2.867459in}{1.397196in}}%
\pgfpathlineto{\pgfqpoint{2.873958in}{1.399810in}}%
\pgfpathlineto{\pgfqpoint{2.883054in}{1.409613in}}%
\pgfpathlineto{\pgfqpoint{2.884754in}{1.413273in}}%
\pgfpathlineto{\pgfqpoint{2.890760in}{1.408191in}}%
\pgfpathlineto{\pgfqpoint{2.895504in}{1.410186in}}%
\pgfpathlineto{\pgfqpoint{2.899044in}{1.407001in}}%
\pgfpathlineto{\pgfqpoint{2.904883in}{1.413872in}}%
\pgfpathlineto{\pgfqpoint{2.909280in}{1.413649in}}%
\pgfpathlineto{\pgfqpoint{2.882440in}{1.391754in}}%
\pgfpathlineto{\pgfqpoint{2.865813in}{1.380465in}}%
\pgfpathlineto{\pgfqpoint{2.859459in}{1.375218in}}%
\pgfpathlineto{\pgfqpoint{2.849223in}{1.369251in}}%
\pgfpathlineto{\pgfqpoint{2.845081in}{1.369000in}}%
\pgfpathlineto{\pgfqpoint{2.835181in}{1.363297in}}%
\pgfpathlineto{\pgfqpoint{2.824782in}{1.361319in}}%
\pgfpathlineto{\pgfqpoint{2.818078in}{1.357032in}}%
\pgfpathlineto{\pgfqpoint{2.810932in}{1.358841in}}%
\pgfpathlineto{\pgfqpoint{2.808914in}{1.354582in}}%
\pgfpathlineto{\pgfqpoint{2.802639in}{1.349500in}}%
\pgfpathlineto{\pgfqpoint{2.803807in}{1.360108in}}%
\pgfpathlineto{\pgfqpoint{2.810062in}{1.361931in}}%
\pgfpathlineto{\pgfqpoint{2.812809in}{1.384940in}}%
\pgfpathlineto{\pgfqpoint{2.810869in}{1.395086in}}%
\pgfpathlineto{\pgfqpoint{2.805188in}{1.401167in}}%
\pgfpathlineto{\pgfqpoint{2.803517in}{1.405975in}}%
\pgfpathclose%
\pgfusepath{fill}%
\end{pgfscope}%
\begin{pgfscope}%
\pgfpathrectangle{\pgfqpoint{0.100000in}{0.100000in}}{\pgfqpoint{3.007045in}{1.925000in}}%
\pgfusepath{clip}%
\pgfsetbuttcap%
\pgfsetmiterjoin%
\definecolor{currentfill}{rgb}{0.513033,0.795755,0.510542}%
\pgfsetfillcolor{currentfill}%
\pgfsetlinewidth{0.000000pt}%
\definecolor{currentstroke}{rgb}{0.000000,0.000000,0.000000}%
\pgfsetstrokecolor{currentstroke}%
\pgfsetstrokeopacity{0.000000}%
\pgfsetdash{}{0pt}%
\pgfpathmoveto{\pgfqpoint{0.981446in}{1.653571in}}%
\pgfpathlineto{\pgfqpoint{0.972309in}{1.661350in}}%
\pgfpathlineto{\pgfqpoint{0.973053in}{1.678906in}}%
\pgfpathlineto{\pgfqpoint{0.969357in}{1.685011in}}%
\pgfpathlineto{\pgfqpoint{0.971079in}{1.692732in}}%
\pgfpathlineto{\pgfqpoint{0.973346in}{1.694790in}}%
\pgfpathlineto{\pgfqpoint{0.982127in}{1.696489in}}%
\pgfpathlineto{\pgfqpoint{0.985045in}{1.700348in}}%
\pgfpathlineto{\pgfqpoint{0.987484in}{1.710382in}}%
\pgfpathlineto{\pgfqpoint{0.983438in}{1.717518in}}%
\pgfpathlineto{\pgfqpoint{0.977851in}{1.718526in}}%
\pgfpathlineto{\pgfqpoint{0.979570in}{1.727393in}}%
\pgfpathlineto{\pgfqpoint{0.967944in}{1.729505in}}%
\pgfpathlineto{\pgfqpoint{0.972359in}{1.752059in}}%
\pgfpathlineto{\pgfqpoint{0.960752in}{1.754525in}}%
\pgfpathlineto{\pgfqpoint{0.964785in}{1.774681in}}%
\pgfpathlineto{\pgfqpoint{0.962444in}{1.782513in}}%
\pgfpathlineto{\pgfqpoint{0.963166in}{1.790823in}}%
\pgfpathlineto{\pgfqpoint{0.967284in}{1.791973in}}%
\pgfpathlineto{\pgfqpoint{0.969863in}{1.800673in}}%
\pgfpathlineto{\pgfqpoint{0.974188in}{1.804466in}}%
\pgfpathlineto{\pgfqpoint{0.975711in}{1.795488in}}%
\pgfpathlineto{\pgfqpoint{0.973893in}{1.786912in}}%
\pgfpathlineto{\pgfqpoint{0.976971in}{1.780376in}}%
\pgfpathlineto{\pgfqpoint{0.987897in}{1.779953in}}%
\pgfpathlineto{\pgfqpoint{0.994400in}{1.776622in}}%
\pgfpathlineto{\pgfqpoint{0.994442in}{1.773321in}}%
\pgfpathlineto{\pgfqpoint{0.999503in}{1.769057in}}%
\pgfpathlineto{\pgfqpoint{1.009841in}{1.767764in}}%
\pgfpathlineto{\pgfqpoint{1.005933in}{1.746597in}}%
\pgfpathlineto{\pgfqpoint{1.010013in}{1.742368in}}%
\pgfpathlineto{\pgfqpoint{1.016513in}{1.740285in}}%
\pgfpathlineto{\pgfqpoint{1.013851in}{1.726227in}}%
\pgfpathlineto{\pgfqpoint{1.019738in}{1.725093in}}%
\pgfpathlineto{\pgfqpoint{1.019449in}{1.720433in}}%
\pgfpathlineto{\pgfqpoint{1.025028in}{1.717016in}}%
\pgfpathlineto{\pgfqpoint{1.027516in}{1.705647in}}%
\pgfpathlineto{\pgfqpoint{1.030764in}{1.699834in}}%
\pgfpathlineto{\pgfqpoint{1.033247in}{1.689867in}}%
\pgfpathlineto{\pgfqpoint{1.037077in}{1.690684in}}%
\pgfpathlineto{\pgfqpoint{1.040593in}{1.686492in}}%
\pgfpathlineto{\pgfqpoint{1.036887in}{1.681674in}}%
\pgfpathlineto{\pgfqpoint{1.037941in}{1.673119in}}%
\pgfpathlineto{\pgfqpoint{1.025044in}{1.675119in}}%
\pgfpathlineto{\pgfqpoint{1.024131in}{1.670887in}}%
\pgfpathlineto{\pgfqpoint{1.019996in}{1.667445in}}%
\pgfpathlineto{\pgfqpoint{1.016570in}{1.660701in}}%
\pgfpathlineto{\pgfqpoint{1.010507in}{1.658210in}}%
\pgfpathlineto{\pgfqpoint{1.004464in}{1.652753in}}%
\pgfpathlineto{\pgfqpoint{0.997011in}{1.655494in}}%
\pgfpathlineto{\pgfqpoint{0.994682in}{1.658385in}}%
\pgfpathlineto{\pgfqpoint{0.988214in}{1.659319in}}%
\pgfpathlineto{\pgfqpoint{0.981446in}{1.653571in}}%
\pgfpathclose%
\pgfusepath{fill}%
\end{pgfscope}%
\begin{pgfscope}%
\pgfpathrectangle{\pgfqpoint{0.100000in}{0.100000in}}{\pgfqpoint{3.007045in}{1.925000in}}%
\pgfusepath{clip}%
\pgfsetbuttcap%
\pgfsetmiterjoin%
\definecolor{currentfill}{rgb}{0.496424,0.788005,0.496886}%
\pgfsetfillcolor{currentfill}%
\pgfsetlinewidth{0.000000pt}%
\definecolor{currentstroke}{rgb}{0.000000,0.000000,0.000000}%
\pgfsetstrokecolor{currentstroke}%
\pgfsetstrokeopacity{0.000000}%
\pgfsetdash{}{0pt}%
\pgfpathmoveto{\pgfqpoint{1.653431in}{1.223538in}}%
\pgfpathlineto{\pgfqpoint{1.630558in}{1.224197in}}%
\pgfpathlineto{\pgfqpoint{1.631190in}{1.247136in}}%
\pgfpathlineto{\pgfqpoint{1.632031in}{1.278421in}}%
\pgfpathlineto{\pgfqpoint{1.635249in}{1.282362in}}%
\pgfpathlineto{\pgfqpoint{1.643733in}{1.288581in}}%
\pgfpathlineto{\pgfqpoint{1.655235in}{1.292354in}}%
\pgfpathlineto{\pgfqpoint{1.654636in}{1.269418in}}%
\pgfpathlineto{\pgfqpoint{1.653431in}{1.223538in}}%
\pgfpathclose%
\pgfusepath{fill}%
\end{pgfscope}%
\begin{pgfscope}%
\pgfpathrectangle{\pgfqpoint{0.100000in}{0.100000in}}{\pgfqpoint{3.007045in}{1.925000in}}%
\pgfusepath{clip}%
\pgfsetbuttcap%
\pgfsetmiterjoin%
\definecolor{currentfill}{rgb}{0.568397,0.821592,0.556063}%
\pgfsetfillcolor{currentfill}%
\pgfsetlinewidth{0.000000pt}%
\definecolor{currentstroke}{rgb}{0.000000,0.000000,0.000000}%
\pgfsetstrokecolor{currentstroke}%
\pgfsetstrokeopacity{0.000000}%
\pgfsetdash{}{0pt}%
\pgfpathmoveto{\pgfqpoint{0.630991in}{1.684466in}}%
\pgfpathlineto{\pgfqpoint{0.628849in}{1.673291in}}%
\pgfpathlineto{\pgfqpoint{0.619536in}{1.640423in}}%
\pgfpathlineto{\pgfqpoint{0.608444in}{1.643504in}}%
\pgfpathlineto{\pgfqpoint{0.609918in}{1.648785in}}%
\pgfpathlineto{\pgfqpoint{0.605211in}{1.653481in}}%
\pgfpathlineto{\pgfqpoint{0.588656in}{1.658125in}}%
\pgfpathlineto{\pgfqpoint{0.594595in}{1.678719in}}%
\pgfpathlineto{\pgfqpoint{0.596303in}{1.682378in}}%
\pgfpathlineto{\pgfqpoint{0.593360in}{1.686222in}}%
\pgfpathlineto{\pgfqpoint{0.593379in}{1.692150in}}%
\pgfpathlineto{\pgfqpoint{0.598195in}{1.705394in}}%
\pgfpathlineto{\pgfqpoint{0.595993in}{1.712179in}}%
\pgfpathlineto{\pgfqpoint{0.602494in}{1.724955in}}%
\pgfpathlineto{\pgfqpoint{0.608895in}{1.725281in}}%
\pgfpathlineto{\pgfqpoint{0.608690in}{1.732151in}}%
\pgfpathlineto{\pgfqpoint{0.606293in}{1.736988in}}%
\pgfpathlineto{\pgfqpoint{0.610892in}{1.738691in}}%
\pgfpathlineto{\pgfqpoint{0.619982in}{1.737934in}}%
\pgfpathlineto{\pgfqpoint{0.630951in}{1.740790in}}%
\pgfpathlineto{\pgfqpoint{0.620630in}{1.705467in}}%
\pgfpathlineto{\pgfqpoint{0.626202in}{1.703795in}}%
\pgfpathlineto{\pgfqpoint{0.624611in}{1.698303in}}%
\pgfpathlineto{\pgfqpoint{0.630102in}{1.696622in}}%
\pgfpathlineto{\pgfqpoint{0.625624in}{1.685979in}}%
\pgfpathlineto{\pgfqpoint{0.630991in}{1.684466in}}%
\pgfpathclose%
\pgfusepath{fill}%
\end{pgfscope}%
\begin{pgfscope}%
\pgfpathrectangle{\pgfqpoint{0.100000in}{0.100000in}}{\pgfqpoint{3.007045in}{1.925000in}}%
\pgfusepath{clip}%
\pgfsetbuttcap%
\pgfsetmiterjoin%
\definecolor{currentfill}{rgb}{0.414118,0.748635,0.442753}%
\pgfsetfillcolor{currentfill}%
\pgfsetlinewidth{0.000000pt}%
\definecolor{currentstroke}{rgb}{0.000000,0.000000,0.000000}%
\pgfsetstrokecolor{currentstroke}%
\pgfsetstrokeopacity{0.000000}%
\pgfsetdash{}{0pt}%
\pgfpathmoveto{\pgfqpoint{1.510401in}{1.206498in}}%
\pgfpathlineto{\pgfqpoint{1.453477in}{1.210184in}}%
\pgfpathlineto{\pgfqpoint{1.455104in}{1.233038in}}%
\pgfpathlineto{\pgfqpoint{1.418635in}{1.235643in}}%
\pgfpathlineto{\pgfqpoint{1.420426in}{1.258524in}}%
\pgfpathlineto{\pgfqpoint{1.460453in}{1.255605in}}%
\pgfpathlineto{\pgfqpoint{1.483897in}{1.254208in}}%
\pgfpathlineto{\pgfqpoint{1.482455in}{1.231145in}}%
\pgfpathlineto{\pgfqpoint{1.511500in}{1.229428in}}%
\pgfpathlineto{\pgfqpoint{1.510401in}{1.206498in}}%
\pgfpathclose%
\pgfusepath{fill}%
\end{pgfscope}%
\begin{pgfscope}%
\pgfpathrectangle{\pgfqpoint{0.100000in}{0.100000in}}{\pgfqpoint{3.007045in}{1.925000in}}%
\pgfusepath{clip}%
\pgfsetbuttcap%
\pgfsetmiterjoin%
\definecolor{currentfill}{rgb}{0.568397,0.821592,0.556063}%
\pgfsetfillcolor{currentfill}%
\pgfsetlinewidth{0.000000pt}%
\definecolor{currentstroke}{rgb}{0.000000,0.000000,0.000000}%
\pgfsetstrokecolor{currentstroke}%
\pgfsetstrokeopacity{0.000000}%
\pgfsetdash{}{0pt}%
\pgfpathmoveto{\pgfqpoint{1.035972in}{1.058230in}}%
\pgfpathlineto{\pgfqpoint{0.986309in}{1.066160in}}%
\pgfpathlineto{\pgfqpoint{0.947453in}{1.073147in}}%
\pgfpathlineto{\pgfqpoint{0.915481in}{1.078785in}}%
\pgfpathlineto{\pgfqpoint{0.923032in}{1.079857in}}%
\pgfpathlineto{\pgfqpoint{0.926450in}{1.083530in}}%
\pgfpathlineto{\pgfqpoint{0.932657in}{1.082202in}}%
\pgfpathlineto{\pgfqpoint{0.941624in}{1.085677in}}%
\pgfpathlineto{\pgfqpoint{0.944689in}{1.092843in}}%
\pgfpathlineto{\pgfqpoint{0.951769in}{1.095219in}}%
\pgfpathlineto{\pgfqpoint{0.958933in}{1.107125in}}%
\pgfpathlineto{\pgfqpoint{0.961951in}{1.110649in}}%
\pgfpathlineto{\pgfqpoint{0.967830in}{1.113017in}}%
\pgfpathlineto{\pgfqpoint{0.975147in}{1.126984in}}%
\pgfpathlineto{\pgfqpoint{0.982614in}{1.127004in}}%
\pgfpathlineto{\pgfqpoint{0.990253in}{1.132805in}}%
\pgfpathlineto{\pgfqpoint{0.992394in}{1.131729in}}%
\pgfpathlineto{\pgfqpoint{0.997732in}{1.138462in}}%
\pgfpathlineto{\pgfqpoint{1.002522in}{1.140541in}}%
\pgfpathlineto{\pgfqpoint{1.005017in}{1.143547in}}%
\pgfpathlineto{\pgfqpoint{1.002930in}{1.149720in}}%
\pgfpathlineto{\pgfqpoint{1.001312in}{1.163987in}}%
\pgfpathlineto{\pgfqpoint{1.050444in}{1.156057in}}%
\pgfpathlineto{\pgfqpoint{1.048160in}{1.141447in}}%
\pgfpathlineto{\pgfqpoint{1.047972in}{1.134089in}}%
\pgfpathlineto{\pgfqpoint{1.045110in}{1.115632in}}%
\pgfpathlineto{\pgfqpoint{1.056798in}{1.114784in}}%
\pgfpathlineto{\pgfqpoint{1.085598in}{1.110401in}}%
\pgfpathlineto{\pgfqpoint{1.087280in}{1.105261in}}%
\pgfpathlineto{\pgfqpoint{1.098980in}{1.105934in}}%
\pgfpathlineto{\pgfqpoint{1.104848in}{1.099688in}}%
\pgfpathlineto{\pgfqpoint{1.103247in}{1.096115in}}%
\pgfpathlineto{\pgfqpoint{1.094867in}{1.088434in}}%
\pgfpathlineto{\pgfqpoint{1.093042in}{1.080708in}}%
\pgfpathlineto{\pgfqpoint{1.088472in}{1.075836in}}%
\pgfpathlineto{\pgfqpoint{1.083485in}{1.074787in}}%
\pgfpathlineto{\pgfqpoint{1.077407in}{1.066910in}}%
\pgfpathlineto{\pgfqpoint{1.073504in}{1.056252in}}%
\pgfpathlineto{\pgfqpoint{1.070635in}{1.052989in}}%
\pgfpathlineto{\pgfqpoint{1.035972in}{1.058230in}}%
\pgfpathclose%
\pgfusepath{fill}%
\end{pgfscope}%
\begin{pgfscope}%
\pgfpathrectangle{\pgfqpoint{0.100000in}{0.100000in}}{\pgfqpoint{3.007045in}{1.925000in}}%
\pgfusepath{clip}%
\pgfsetbuttcap%
\pgfsetmiterjoin%
\definecolor{currentfill}{rgb}{0.704421,0.881738,0.678970}%
\pgfsetfillcolor{currentfill}%
\pgfsetlinewidth{0.000000pt}%
\definecolor{currentstroke}{rgb}{0.000000,0.000000,0.000000}%
\pgfsetstrokecolor{currentstroke}%
\pgfsetstrokeopacity{0.000000}%
\pgfsetdash{}{0pt}%
\pgfpathmoveto{\pgfqpoint{2.241623in}{0.891209in}}%
\pgfpathlineto{\pgfqpoint{2.240204in}{0.883283in}}%
\pgfpathlineto{\pgfqpoint{2.241813in}{0.880528in}}%
\pgfpathlineto{\pgfqpoint{2.241299in}{0.867393in}}%
\pgfpathlineto{\pgfqpoint{2.244466in}{0.861498in}}%
\pgfpathlineto{\pgfqpoint{2.238722in}{0.856984in}}%
\pgfpathlineto{\pgfqpoint{2.232658in}{0.857963in}}%
\pgfpathlineto{\pgfqpoint{2.231240in}{0.844683in}}%
\pgfpathlineto{\pgfqpoint{2.202469in}{0.842698in}}%
\pgfpathlineto{\pgfqpoint{2.202564in}{0.841744in}}%
\pgfpathlineto{\pgfqpoint{2.179696in}{0.840159in}}%
\pgfpathlineto{\pgfqpoint{2.178278in}{0.857446in}}%
\pgfpathlineto{\pgfqpoint{2.183086in}{0.867245in}}%
\pgfpathlineto{\pgfqpoint{2.182587in}{0.873244in}}%
\pgfpathlineto{\pgfqpoint{2.189352in}{0.871261in}}%
\pgfpathlineto{\pgfqpoint{2.193931in}{0.875265in}}%
\pgfpathlineto{\pgfqpoint{2.193121in}{0.887325in}}%
\pgfpathlineto{\pgfqpoint{2.213336in}{0.888623in}}%
\pgfpathlineto{\pgfqpoint{2.212159in}{0.906542in}}%
\pgfpathlineto{\pgfqpoint{2.216592in}{0.906228in}}%
\pgfpathlineto{\pgfqpoint{2.223886in}{0.911812in}}%
\pgfpathlineto{\pgfqpoint{2.223823in}{0.914321in}}%
\pgfpathlineto{\pgfqpoint{2.235740in}{0.906849in}}%
\pgfpathlineto{\pgfqpoint{2.237551in}{0.900340in}}%
\pgfpathlineto{\pgfqpoint{2.240394in}{0.900117in}}%
\pgfpathlineto{\pgfqpoint{2.241623in}{0.891209in}}%
\pgfpathclose%
\pgfusepath{fill}%
\end{pgfscope}%
\begin{pgfscope}%
\pgfpathrectangle{\pgfqpoint{0.100000in}{0.100000in}}{\pgfqpoint{3.007045in}{1.925000in}}%
\pgfusepath{clip}%
\pgfsetbuttcap%
\pgfsetmiterjoin%
\definecolor{currentfill}{rgb}{0.723122,0.889612,0.697178}%
\pgfsetfillcolor{currentfill}%
\pgfsetlinewidth{0.000000pt}%
\definecolor{currentstroke}{rgb}{0.000000,0.000000,0.000000}%
\pgfsetstrokecolor{currentstroke}%
\pgfsetstrokeopacity{0.000000}%
\pgfsetdash{}{0pt}%
\pgfpathmoveto{\pgfqpoint{1.062663in}{1.231515in}}%
\pgfpathlineto{\pgfqpoint{1.072883in}{1.296139in}}%
\pgfpathlineto{\pgfqpoint{1.076315in}{1.318750in}}%
\pgfpathlineto{\pgfqpoint{1.102069in}{1.314741in}}%
\pgfpathlineto{\pgfqpoint{1.131992in}{1.310519in}}%
\pgfpathlineto{\pgfqpoint{1.184259in}{1.303392in}}%
\pgfpathlineto{\pgfqpoint{1.183883in}{1.298271in}}%
\pgfpathlineto{\pgfqpoint{1.192988in}{1.292052in}}%
\pgfpathlineto{\pgfqpoint{1.193426in}{1.288196in}}%
\pgfpathlineto{\pgfqpoint{1.188332in}{1.274642in}}%
\pgfpathlineto{\pgfqpoint{1.189706in}{1.265642in}}%
\pgfpathlineto{\pgfqpoint{1.189821in}{1.258740in}}%
\pgfpathlineto{\pgfqpoint{1.186728in}{1.236591in}}%
\pgfpathlineto{\pgfqpoint{1.186601in}{1.231472in}}%
\pgfpathlineto{\pgfqpoint{1.166255in}{1.233746in}}%
\pgfpathlineto{\pgfqpoint{1.167476in}{1.245055in}}%
\pgfpathlineto{\pgfqpoint{1.153503in}{1.246887in}}%
\pgfpathlineto{\pgfqpoint{1.151850in}{1.235425in}}%
\pgfpathlineto{\pgfqpoint{1.146223in}{1.236379in}}%
\pgfpathlineto{\pgfqpoint{1.145415in}{1.230549in}}%
\pgfpathlineto{\pgfqpoint{1.120000in}{1.233895in}}%
\pgfpathlineto{\pgfqpoint{1.118778in}{1.225384in}}%
\pgfpathlineto{\pgfqpoint{1.087872in}{1.229867in}}%
\pgfpathlineto{\pgfqpoint{1.087451in}{1.227034in}}%
\pgfpathlineto{\pgfqpoint{1.062663in}{1.231515in}}%
\pgfpathclose%
\pgfusepath{fill}%
\end{pgfscope}%
\begin{pgfscope}%
\pgfpathrectangle{\pgfqpoint{0.100000in}{0.100000in}}{\pgfqpoint{3.007045in}{1.925000in}}%
\pgfusepath{clip}%
\pgfsetbuttcap%
\pgfsetmiterjoin%
\definecolor{currentfill}{rgb}{0.774548,0.911265,0.747251}%
\pgfsetfillcolor{currentfill}%
\pgfsetlinewidth{0.000000pt}%
\definecolor{currentstroke}{rgb}{0.000000,0.000000,0.000000}%
\pgfsetstrokecolor{currentstroke}%
\pgfsetstrokeopacity{0.000000}%
\pgfsetdash{}{0pt}%
\pgfpathmoveto{\pgfqpoint{1.441506in}{0.486973in}}%
\pgfpathlineto{\pgfqpoint{1.473940in}{0.485183in}}%
\pgfpathlineto{\pgfqpoint{1.470461in}{0.425819in}}%
\pgfpathlineto{\pgfqpoint{1.464650in}{0.426083in}}%
\pgfpathlineto{\pgfqpoint{1.460228in}{0.431872in}}%
\pgfpathlineto{\pgfqpoint{1.458257in}{0.442114in}}%
\pgfpathlineto{\pgfqpoint{1.454336in}{0.450453in}}%
\pgfpathlineto{\pgfqpoint{1.455112in}{0.452787in}}%
\pgfpathlineto{\pgfqpoint{1.449609in}{0.458112in}}%
\pgfpathlineto{\pgfqpoint{1.447489in}{0.469258in}}%
\pgfpathlineto{\pgfqpoint{1.441856in}{0.477516in}}%
\pgfpathlineto{\pgfqpoint{1.441506in}{0.486973in}}%
\pgfpathclose%
\pgfusepath{fill}%
\end{pgfscope}%
\begin{pgfscope}%
\pgfpathrectangle{\pgfqpoint{0.100000in}{0.100000in}}{\pgfqpoint{3.007045in}{1.925000in}}%
\pgfusepath{clip}%
\pgfsetbuttcap%
\pgfsetmiterjoin%
\definecolor{currentfill}{rgb}{0.590542,0.831926,0.574271}%
\pgfsetfillcolor{currentfill}%
\pgfsetlinewidth{0.000000pt}%
\definecolor{currentstroke}{rgb}{0.000000,0.000000,0.000000}%
\pgfsetstrokecolor{currentstroke}%
\pgfsetstrokeopacity{0.000000}%
\pgfsetdash{}{0pt}%
\pgfpathmoveto{\pgfqpoint{1.775875in}{0.660550in}}%
\pgfpathlineto{\pgfqpoint{1.773764in}{0.656330in}}%
\pgfpathlineto{\pgfqpoint{1.776581in}{0.643123in}}%
\pgfpathlineto{\pgfqpoint{1.779925in}{0.639399in}}%
\pgfpathlineto{\pgfqpoint{1.772144in}{0.632693in}}%
\pgfpathlineto{\pgfqpoint{1.766321in}{0.636474in}}%
\pgfpathlineto{\pgfqpoint{1.763851in}{0.642226in}}%
\pgfpathlineto{\pgfqpoint{1.756860in}{0.643794in}}%
\pgfpathlineto{\pgfqpoint{1.735433in}{0.640512in}}%
\pgfpathlineto{\pgfqpoint{1.730423in}{0.638033in}}%
\pgfpathlineto{\pgfqpoint{1.732103in}{0.644712in}}%
\pgfpathlineto{\pgfqpoint{1.723459in}{0.650585in}}%
\pgfpathlineto{\pgfqpoint{1.722872in}{0.654828in}}%
\pgfpathlineto{\pgfqpoint{1.718592in}{0.656205in}}%
\pgfpathlineto{\pgfqpoint{1.712375in}{0.673683in}}%
\pgfpathlineto{\pgfqpoint{1.703043in}{0.687922in}}%
\pgfpathlineto{\pgfqpoint{1.694814in}{0.693189in}}%
\pgfpathlineto{\pgfqpoint{1.712066in}{0.694873in}}%
\pgfpathlineto{\pgfqpoint{1.712393in}{0.726751in}}%
\pgfpathlineto{\pgfqpoint{1.720316in}{0.726579in}}%
\pgfpathlineto{\pgfqpoint{1.725321in}{0.722021in}}%
\pgfpathlineto{\pgfqpoint{1.728058in}{0.722452in}}%
\pgfpathlineto{\pgfqpoint{1.736915in}{0.718690in}}%
\pgfpathlineto{\pgfqpoint{1.735350in}{0.734610in}}%
\pgfpathlineto{\pgfqpoint{1.755158in}{0.734645in}}%
\pgfpathlineto{\pgfqpoint{1.762705in}{0.734578in}}%
\pgfpathlineto{\pgfqpoint{1.763868in}{0.730641in}}%
\pgfpathlineto{\pgfqpoint{1.763758in}{0.708620in}}%
\pgfpathlineto{\pgfqpoint{1.773020in}{0.706433in}}%
\pgfpathlineto{\pgfqpoint{1.773171in}{0.660592in}}%
\pgfpathlineto{\pgfqpoint{1.775875in}{0.660550in}}%
\pgfpathclose%
\pgfusepath{fill}%
\end{pgfscope}%
\begin{pgfscope}%
\pgfpathrectangle{\pgfqpoint{0.100000in}{0.100000in}}{\pgfqpoint{3.007045in}{1.925000in}}%
\pgfusepath{clip}%
\pgfsetbuttcap%
\pgfsetmiterjoin%
\definecolor{currentfill}{rgb}{0.699746,0.879769,0.674418}%
\pgfsetfillcolor{currentfill}%
\pgfsetlinewidth{0.000000pt}%
\definecolor{currentstroke}{rgb}{0.000000,0.000000,0.000000}%
\pgfsetstrokecolor{currentstroke}%
\pgfsetstrokeopacity{0.000000}%
\pgfsetdash{}{0pt}%
\pgfpathmoveto{\pgfqpoint{0.407690in}{1.233195in}}%
\pgfpathlineto{\pgfqpoint{0.404604in}{1.238315in}}%
\pgfpathlineto{\pgfqpoint{0.409870in}{1.255826in}}%
\pgfpathlineto{\pgfqpoint{0.409743in}{1.259767in}}%
\pgfpathlineto{\pgfqpoint{0.413844in}{1.273461in}}%
\pgfpathlineto{\pgfqpoint{0.409876in}{1.274771in}}%
\pgfpathlineto{\pgfqpoint{0.409917in}{1.279165in}}%
\pgfpathlineto{\pgfqpoint{0.414282in}{1.280552in}}%
\pgfpathlineto{\pgfqpoint{0.417431in}{1.287454in}}%
\pgfpathlineto{\pgfqpoint{0.412425in}{1.289011in}}%
\pgfpathlineto{\pgfqpoint{0.416471in}{1.302466in}}%
\pgfpathlineto{\pgfqpoint{0.404388in}{1.306594in}}%
\pgfpathlineto{\pgfqpoint{0.400092in}{1.304827in}}%
\pgfpathlineto{\pgfqpoint{0.396070in}{1.307741in}}%
\pgfpathlineto{\pgfqpoint{0.395623in}{1.317666in}}%
\pgfpathlineto{\pgfqpoint{0.393792in}{1.321038in}}%
\pgfpathlineto{\pgfqpoint{0.393337in}{1.331101in}}%
\pgfpathlineto{\pgfqpoint{0.388567in}{1.334258in}}%
\pgfpathlineto{\pgfqpoint{0.392426in}{1.337198in}}%
\pgfpathlineto{\pgfqpoint{0.417135in}{1.329622in}}%
\pgfpathlineto{\pgfqpoint{0.421047in}{1.323928in}}%
\pgfpathlineto{\pgfqpoint{0.420345in}{1.320339in}}%
\pgfpathlineto{\pgfqpoint{0.424973in}{1.314267in}}%
\pgfpathlineto{\pgfqpoint{0.431828in}{1.313269in}}%
\pgfpathlineto{\pgfqpoint{0.435140in}{1.324290in}}%
\pgfpathlineto{\pgfqpoint{0.439260in}{1.328328in}}%
\pgfpathlineto{\pgfqpoint{0.445384in}{1.329930in}}%
\pgfpathlineto{\pgfqpoint{0.450179in}{1.340731in}}%
\pgfpathlineto{\pgfqpoint{0.460387in}{1.347970in}}%
\pgfpathlineto{\pgfqpoint{0.465761in}{1.347164in}}%
\pgfpathlineto{\pgfqpoint{0.480726in}{1.347606in}}%
\pgfpathlineto{\pgfqpoint{0.486181in}{1.350484in}}%
\pgfpathlineto{\pgfqpoint{0.489895in}{1.349175in}}%
\pgfpathlineto{\pgfqpoint{0.492015in}{1.343273in}}%
\pgfpathlineto{\pgfqpoint{0.516565in}{1.336204in}}%
\pgfpathlineto{\pgfqpoint{0.510601in}{1.315136in}}%
\pgfpathlineto{\pgfqpoint{0.522542in}{1.311873in}}%
\pgfpathlineto{\pgfqpoint{0.522474in}{1.309901in}}%
\pgfpathlineto{\pgfqpoint{0.544151in}{1.303815in}}%
\pgfpathlineto{\pgfqpoint{0.542583in}{1.297900in}}%
\pgfpathlineto{\pgfqpoint{0.536518in}{1.297107in}}%
\pgfpathlineto{\pgfqpoint{0.533368in}{1.282853in}}%
\pgfpathlineto{\pgfqpoint{0.536090in}{1.281820in}}%
\pgfpathlineto{\pgfqpoint{0.533766in}{1.269179in}}%
\pgfpathlineto{\pgfqpoint{0.524302in}{1.283464in}}%
\pgfpathlineto{\pgfqpoint{0.520542in}{1.277526in}}%
\pgfpathlineto{\pgfqpoint{0.522497in}{1.269886in}}%
\pgfpathlineto{\pgfqpoint{0.519485in}{1.264180in}}%
\pgfpathlineto{\pgfqpoint{0.514663in}{1.260246in}}%
\pgfpathlineto{\pgfqpoint{0.510697in}{1.267522in}}%
\pgfpathlineto{\pgfqpoint{0.502964in}{1.265519in}}%
\pgfpathlineto{\pgfqpoint{0.497623in}{1.272306in}}%
\pgfpathlineto{\pgfqpoint{0.488781in}{1.270732in}}%
\pgfpathlineto{\pgfqpoint{0.475080in}{1.259356in}}%
\pgfpathlineto{\pgfqpoint{0.464303in}{1.248933in}}%
\pgfpathlineto{\pgfqpoint{0.455151in}{1.243275in}}%
\pgfpathlineto{\pgfqpoint{0.446079in}{1.262743in}}%
\pgfpathlineto{\pgfqpoint{0.440085in}{1.241248in}}%
\pgfpathlineto{\pgfqpoint{0.436885in}{1.243705in}}%
\pgfpathlineto{\pgfqpoint{0.424991in}{1.242670in}}%
\pgfpathlineto{\pgfqpoint{0.407690in}{1.233195in}}%
\pgfpathclose%
\pgfusepath{fill}%
\end{pgfscope}%
\begin{pgfscope}%
\pgfpathrectangle{\pgfqpoint{0.100000in}{0.100000in}}{\pgfqpoint{3.007045in}{1.925000in}}%
\pgfusepath{clip}%
\pgfsetbuttcap%
\pgfsetmiterjoin%
\definecolor{currentfill}{rgb}{0.420392,0.751711,0.445829}%
\pgfsetfillcolor{currentfill}%
\pgfsetlinewidth{0.000000pt}%
\definecolor{currentstroke}{rgb}{0.000000,0.000000,0.000000}%
\pgfsetstrokecolor{currentstroke}%
\pgfsetstrokeopacity{0.000000}%
\pgfsetdash{}{0pt}%
\pgfpathmoveto{\pgfqpoint{1.588665in}{1.618076in}}%
\pgfpathlineto{\pgfqpoint{1.587099in}{1.618157in}}%
\pgfpathlineto{\pgfqpoint{1.588163in}{1.641202in}}%
\pgfpathlineto{\pgfqpoint{1.569331in}{1.642166in}}%
\pgfpathlineto{\pgfqpoint{1.570614in}{1.665351in}}%
\pgfpathlineto{\pgfqpoint{1.569042in}{1.665415in}}%
\pgfpathlineto{\pgfqpoint{1.570193in}{1.688446in}}%
\pgfpathlineto{\pgfqpoint{1.579862in}{1.687980in}}%
\pgfpathlineto{\pgfqpoint{1.614445in}{1.686458in}}%
\pgfpathlineto{\pgfqpoint{1.615660in}{1.680679in}}%
\pgfpathlineto{\pgfqpoint{1.638529in}{1.679874in}}%
\pgfpathlineto{\pgfqpoint{1.650048in}{1.679481in}}%
\pgfpathlineto{\pgfqpoint{1.650635in}{1.662065in}}%
\pgfpathlineto{\pgfqpoint{1.649938in}{1.638890in}}%
\pgfpathlineto{\pgfqpoint{1.633892in}{1.639443in}}%
\pgfpathlineto{\pgfqpoint{1.633089in}{1.616328in}}%
\pgfpathlineto{\pgfqpoint{1.588665in}{1.618076in}}%
\pgfpathclose%
\pgfusepath{fill}%
\end{pgfscope}%
\begin{pgfscope}%
\pgfpathrectangle{\pgfqpoint{0.100000in}{0.100000in}}{\pgfqpoint{3.007045in}{1.925000in}}%
\pgfusepath{clip}%
\pgfsetbuttcap%
\pgfsetmiterjoin%
\definecolor{currentfill}{rgb}{0.320000,0.702499,0.396617}%
\pgfsetfillcolor{currentfill}%
\pgfsetlinewidth{0.000000pt}%
\definecolor{currentstroke}{rgb}{0.000000,0.000000,0.000000}%
\pgfsetstrokecolor{currentstroke}%
\pgfsetstrokeopacity{0.000000}%
\pgfsetdash{}{0pt}%
\pgfpathmoveto{\pgfqpoint{2.867151in}{1.544262in}}%
\pgfpathlineto{\pgfqpoint{2.860753in}{1.546240in}}%
\pgfpathlineto{\pgfqpoint{2.850891in}{1.542069in}}%
\pgfpathlineto{\pgfqpoint{2.846989in}{1.546225in}}%
\pgfpathlineto{\pgfqpoint{2.835404in}{1.542568in}}%
\pgfpathlineto{\pgfqpoint{2.828637in}{1.543127in}}%
\pgfpathlineto{\pgfqpoint{2.825651in}{1.545463in}}%
\pgfpathlineto{\pgfqpoint{2.829267in}{1.550294in}}%
\pgfpathlineto{\pgfqpoint{2.829522in}{1.557864in}}%
\pgfpathlineto{\pgfqpoint{2.826130in}{1.558806in}}%
\pgfpathlineto{\pgfqpoint{2.827305in}{1.571893in}}%
\pgfpathlineto{\pgfqpoint{2.821240in}{1.572694in}}%
\pgfpathlineto{\pgfqpoint{2.821564in}{1.578723in}}%
\pgfpathlineto{\pgfqpoint{2.811802in}{1.580671in}}%
\pgfpathlineto{\pgfqpoint{2.815683in}{1.583923in}}%
\pgfpathlineto{\pgfqpoint{2.819866in}{1.593654in}}%
\pgfpathlineto{\pgfqpoint{2.811115in}{1.594372in}}%
\pgfpathlineto{\pgfqpoint{2.808838in}{1.600051in}}%
\pgfpathlineto{\pgfqpoint{2.808017in}{1.607896in}}%
\pgfpathlineto{\pgfqpoint{2.810189in}{1.620019in}}%
\pgfpathlineto{\pgfqpoint{2.817950in}{1.618314in}}%
\pgfpathlineto{\pgfqpoint{2.819281in}{1.624179in}}%
\pgfpathlineto{\pgfqpoint{2.825014in}{1.621968in}}%
\pgfpathlineto{\pgfqpoint{2.826283in}{1.627777in}}%
\pgfpathlineto{\pgfqpoint{2.837240in}{1.625008in}}%
\pgfpathlineto{\pgfqpoint{2.834224in}{1.609178in}}%
\pgfpathlineto{\pgfqpoint{2.837477in}{1.608579in}}%
\pgfpathlineto{\pgfqpoint{2.848338in}{1.611817in}}%
\pgfpathlineto{\pgfqpoint{2.846880in}{1.617175in}}%
\pgfpathlineto{\pgfqpoint{2.849656in}{1.622164in}}%
\pgfpathlineto{\pgfqpoint{2.854814in}{1.623529in}}%
\pgfpathlineto{\pgfqpoint{2.858546in}{1.629287in}}%
\pgfpathlineto{\pgfqpoint{2.862905in}{1.625347in}}%
\pgfpathlineto{\pgfqpoint{2.868383in}{1.624903in}}%
\pgfpathlineto{\pgfqpoint{2.869747in}{1.621836in}}%
\pgfpathlineto{\pgfqpoint{2.875904in}{1.623095in}}%
\pgfpathlineto{\pgfqpoint{2.877796in}{1.621050in}}%
\pgfpathlineto{\pgfqpoint{2.882477in}{1.612503in}}%
\pgfpathlineto{\pgfqpoint{2.885308in}{1.602759in}}%
\pgfpathlineto{\pgfqpoint{2.876355in}{1.598852in}}%
\pgfpathlineto{\pgfqpoint{2.881617in}{1.586847in}}%
\pgfpathlineto{\pgfqpoint{2.875977in}{1.584608in}}%
\pgfpathlineto{\pgfqpoint{2.876946in}{1.580426in}}%
\pgfpathlineto{\pgfqpoint{2.870513in}{1.573950in}}%
\pgfpathlineto{\pgfqpoint{2.867917in}{1.577405in}}%
\pgfpathlineto{\pgfqpoint{2.863308in}{1.564293in}}%
\pgfpathlineto{\pgfqpoint{2.861539in}{1.556030in}}%
\pgfpathlineto{\pgfqpoint{2.867535in}{1.548390in}}%
\pgfpathlineto{\pgfqpoint{2.867151in}{1.544262in}}%
\pgfpathclose%
\pgfusepath{fill}%
\end{pgfscope}%
\begin{pgfscope}%
\pgfpathrectangle{\pgfqpoint{0.100000in}{0.100000in}}{\pgfqpoint{3.007045in}{1.925000in}}%
\pgfusepath{clip}%
\pgfsetbuttcap%
\pgfsetmiterjoin%
\definecolor{currentfill}{rgb}{0.468743,0.775087,0.474125}%
\pgfsetfillcolor{currentfill}%
\pgfsetlinewidth{0.000000pt}%
\definecolor{currentstroke}{rgb}{0.000000,0.000000,0.000000}%
\pgfsetstrokecolor{currentstroke}%
\pgfsetstrokeopacity{0.000000}%
\pgfsetdash{}{0pt}%
\pgfpathmoveto{\pgfqpoint{2.552895in}{1.336088in}}%
\pgfpathlineto{\pgfqpoint{2.550029in}{1.332847in}}%
\pgfpathlineto{\pgfqpoint{2.540085in}{1.329055in}}%
\pgfpathlineto{\pgfqpoint{2.533624in}{1.328350in}}%
\pgfpathlineto{\pgfqpoint{2.528039in}{1.339978in}}%
\pgfpathlineto{\pgfqpoint{2.513140in}{1.337398in}}%
\pgfpathlineto{\pgfqpoint{2.509694in}{1.358087in}}%
\pgfpathlineto{\pgfqpoint{2.503675in}{1.356445in}}%
\pgfpathlineto{\pgfqpoint{2.484276in}{1.354588in}}%
\pgfpathlineto{\pgfqpoint{2.479266in}{1.385678in}}%
\pgfpathlineto{\pgfqpoint{2.494271in}{1.396072in}}%
\pgfpathlineto{\pgfqpoint{2.507781in}{1.407380in}}%
\pgfpathlineto{\pgfqpoint{2.516225in}{1.413233in}}%
\pgfpathlineto{\pgfqpoint{2.526196in}{1.422773in}}%
\pgfpathlineto{\pgfqpoint{2.534317in}{1.431974in}}%
\pgfpathlineto{\pgfqpoint{2.539326in}{1.435418in}}%
\pgfpathlineto{\pgfqpoint{2.543334in}{1.433990in}}%
\pgfpathlineto{\pgfqpoint{2.549510in}{1.399000in}}%
\pgfpathlineto{\pgfqpoint{2.556395in}{1.400149in}}%
\pgfpathlineto{\pgfqpoint{2.558369in}{1.389969in}}%
\pgfpathlineto{\pgfqpoint{2.556898in}{1.389007in}}%
\pgfpathlineto{\pgfqpoint{2.558934in}{1.375533in}}%
\pgfpathlineto{\pgfqpoint{2.559925in}{1.364184in}}%
\pgfpathlineto{\pgfqpoint{2.554563in}{1.361326in}}%
\pgfpathlineto{\pgfqpoint{2.555437in}{1.355961in}}%
\pgfpathlineto{\pgfqpoint{2.549986in}{1.354398in}}%
\pgfpathlineto{\pgfqpoint{2.552895in}{1.336088in}}%
\pgfpathclose%
\pgfusepath{fill}%
\end{pgfscope}%
\begin{pgfscope}%
\pgfpathrectangle{\pgfqpoint{0.100000in}{0.100000in}}{\pgfqpoint{3.007045in}{1.925000in}}%
\pgfusepath{clip}%
\pgfsetbuttcap%
\pgfsetmiterjoin%
\definecolor{currentfill}{rgb}{0.590542,0.831926,0.574271}%
\pgfsetfillcolor{currentfill}%
\pgfsetlinewidth{0.000000pt}%
\definecolor{currentstroke}{rgb}{0.000000,0.000000,0.000000}%
\pgfsetstrokecolor{currentstroke}%
\pgfsetstrokeopacity{0.000000}%
\pgfsetdash{}{0pt}%
\pgfpathmoveto{\pgfqpoint{1.677475in}{1.268823in}}%
\pgfpathlineto{\pgfqpoint{1.654636in}{1.269418in}}%
\pgfpathlineto{\pgfqpoint{1.655235in}{1.292354in}}%
\pgfpathlineto{\pgfqpoint{1.643733in}{1.288581in}}%
\pgfpathlineto{\pgfqpoint{1.643851in}{1.292672in}}%
\pgfpathlineto{\pgfqpoint{1.638646in}{1.292816in}}%
\pgfpathlineto{\pgfqpoint{1.638392in}{1.287167in}}%
\pgfpathlineto{\pgfqpoint{1.621155in}{1.286648in}}%
\pgfpathlineto{\pgfqpoint{1.618684in}{1.284817in}}%
\pgfpathlineto{\pgfqpoint{1.610083in}{1.285141in}}%
\pgfpathlineto{\pgfqpoint{1.609594in}{1.293697in}}%
\pgfpathlineto{\pgfqpoint{1.610215in}{1.316599in}}%
\pgfpathlineto{\pgfqpoint{1.610638in}{1.328045in}}%
\pgfpathlineto{\pgfqpoint{1.633318in}{1.327373in}}%
\pgfpathlineto{\pgfqpoint{1.633082in}{1.315887in}}%
\pgfpathlineto{\pgfqpoint{1.672961in}{1.314877in}}%
\pgfpathlineto{\pgfqpoint{1.678558in}{1.314735in}}%
\pgfpathlineto{\pgfqpoint{1.677475in}{1.268823in}}%
\pgfpathclose%
\pgfusepath{fill}%
\end{pgfscope}%
\begin{pgfscope}%
\pgfpathrectangle{\pgfqpoint{0.100000in}{0.100000in}}{\pgfqpoint{3.007045in}{1.925000in}}%
\pgfusepath{clip}%
\pgfsetbuttcap%
\pgfsetmiterjoin%
\definecolor{currentfill}{rgb}{0.652995,0.860085,0.628897}%
\pgfsetfillcolor{currentfill}%
\pgfsetlinewidth{0.000000pt}%
\definecolor{currentstroke}{rgb}{0.000000,0.000000,0.000000}%
\pgfsetstrokecolor{currentstroke}%
\pgfsetstrokeopacity{0.000000}%
\pgfsetdash{}{0pt}%
\pgfpathmoveto{\pgfqpoint{1.779266in}{1.452736in}}%
\pgfpathlineto{\pgfqpoint{1.778980in}{1.469921in}}%
\pgfpathlineto{\pgfqpoint{1.767569in}{1.469929in}}%
\pgfpathlineto{\pgfqpoint{1.767169in}{1.475679in}}%
\pgfpathlineto{\pgfqpoint{1.767240in}{1.481415in}}%
\pgfpathlineto{\pgfqpoint{1.778661in}{1.481448in}}%
\pgfpathlineto{\pgfqpoint{1.778660in}{1.495730in}}%
\pgfpathlineto{\pgfqpoint{1.782199in}{1.492987in}}%
\pgfpathlineto{\pgfqpoint{1.790073in}{1.492987in}}%
\pgfpathlineto{\pgfqpoint{1.822887in}{1.493279in}}%
\pgfpathlineto{\pgfqpoint{1.823694in}{1.499014in}}%
\pgfpathlineto{\pgfqpoint{1.841921in}{1.499291in}}%
\pgfpathlineto{\pgfqpoint{1.842251in}{1.476305in}}%
\pgfpathlineto{\pgfqpoint{1.830743in}{1.476124in}}%
\pgfpathlineto{\pgfqpoint{1.831065in}{1.453126in}}%
\pgfpathlineto{\pgfqpoint{1.779266in}{1.452736in}}%
\pgfpathclose%
\pgfusepath{fill}%
\end{pgfscope}%
\begin{pgfscope}%
\pgfpathrectangle{\pgfqpoint{0.100000in}{0.100000in}}{\pgfqpoint{3.007045in}{1.925000in}}%
\pgfusepath{clip}%
\pgfsetbuttcap%
\pgfsetmiterjoin%
\definecolor{currentfill}{rgb}{0.507497,0.793172,0.505990}%
\pgfsetfillcolor{currentfill}%
\pgfsetlinewidth{0.000000pt}%
\definecolor{currentstroke}{rgb}{0.000000,0.000000,0.000000}%
\pgfsetstrokecolor{currentstroke}%
\pgfsetstrokeopacity{0.000000}%
\pgfsetdash{}{0pt}%
\pgfpathmoveto{\pgfqpoint{1.890498in}{1.239370in}}%
\pgfpathlineto{\pgfqpoint{1.889907in}{1.259660in}}%
\pgfpathlineto{\pgfqpoint{1.867080in}{1.259096in}}%
\pgfpathlineto{\pgfqpoint{1.866638in}{1.276378in}}%
\pgfpathlineto{\pgfqpoint{1.855298in}{1.276147in}}%
\pgfpathlineto{\pgfqpoint{1.854836in}{1.299029in}}%
\pgfpathlineto{\pgfqpoint{1.905724in}{1.300373in}}%
\pgfpathlineto{\pgfqpoint{1.923061in}{1.300995in}}%
\pgfpathlineto{\pgfqpoint{1.923843in}{1.278117in}}%
\pgfpathlineto{\pgfqpoint{1.935244in}{1.278471in}}%
\pgfpathlineto{\pgfqpoint{1.935891in}{1.242279in}}%
\pgfpathlineto{\pgfqpoint{1.925211in}{1.241392in}}%
\pgfpathlineto{\pgfqpoint{1.890498in}{1.239370in}}%
\pgfpathclose%
\pgfusepath{fill}%
\end{pgfscope}%
\begin{pgfscope}%
\pgfpathrectangle{\pgfqpoint{0.100000in}{0.100000in}}{\pgfqpoint{3.007045in}{1.925000in}}%
\pgfusepath{clip}%
\pgfsetbuttcap%
\pgfsetmiterjoin%
\definecolor{currentfill}{rgb}{0.906621,0.964121,0.888443}%
\pgfsetfillcolor{currentfill}%
\pgfsetlinewidth{0.000000pt}%
\definecolor{currentstroke}{rgb}{0.000000,0.000000,0.000000}%
\pgfsetstrokecolor{currentstroke}%
\pgfsetstrokeopacity{0.000000}%
\pgfsetdash{}{0pt}%
\pgfpathmoveto{\pgfqpoint{0.967728in}{1.403968in}}%
\pgfpathlineto{\pgfqpoint{0.968744in}{1.400234in}}%
\pgfpathlineto{\pgfqpoint{0.965948in}{1.394443in}}%
\pgfpathlineto{\pgfqpoint{0.967550in}{1.392277in}}%
\pgfpathlineto{\pgfqpoint{0.968734in}{1.380626in}}%
\pgfpathlineto{\pgfqpoint{0.965270in}{1.373029in}}%
\pgfpathlineto{\pgfqpoint{0.960688in}{1.366618in}}%
\pgfpathlineto{\pgfqpoint{0.951603in}{1.367889in}}%
\pgfpathlineto{\pgfqpoint{0.947822in}{1.365378in}}%
\pgfpathlineto{\pgfqpoint{0.942058in}{1.369750in}}%
\pgfpathlineto{\pgfqpoint{0.934476in}{1.365698in}}%
\pgfpathlineto{\pgfqpoint{0.924294in}{1.367661in}}%
\pgfpathlineto{\pgfqpoint{0.908601in}{1.353242in}}%
\pgfpathlineto{\pgfqpoint{0.892707in}{1.351230in}}%
\pgfpathlineto{\pgfqpoint{0.832155in}{1.363654in}}%
\pgfpathlineto{\pgfqpoint{0.845804in}{1.427734in}}%
\pgfpathlineto{\pgfqpoint{0.853680in}{1.425681in}}%
\pgfpathlineto{\pgfqpoint{0.895826in}{1.417699in}}%
\pgfpathlineto{\pgfqpoint{0.900100in}{1.438993in}}%
\pgfpathlineto{\pgfqpoint{0.916709in}{1.435786in}}%
\pgfpathlineto{\pgfqpoint{0.917819in}{1.441425in}}%
\pgfpathlineto{\pgfqpoint{0.925285in}{1.439963in}}%
\pgfpathlineto{\pgfqpoint{0.926367in}{1.445605in}}%
\pgfpathlineto{\pgfqpoint{0.931789in}{1.444541in}}%
\pgfpathlineto{\pgfqpoint{0.934955in}{1.436848in}}%
\pgfpathlineto{\pgfqpoint{0.941439in}{1.428069in}}%
\pgfpathlineto{\pgfqpoint{0.948249in}{1.426793in}}%
\pgfpathlineto{\pgfqpoint{0.952072in}{1.424129in}}%
\pgfpathlineto{\pgfqpoint{0.955057in}{1.434288in}}%
\pgfpathlineto{\pgfqpoint{0.968314in}{1.431784in}}%
\pgfpathlineto{\pgfqpoint{0.969777in}{1.427234in}}%
\pgfpathlineto{\pgfqpoint{0.968279in}{1.422237in}}%
\pgfpathlineto{\pgfqpoint{0.964506in}{1.418586in}}%
\pgfpathlineto{\pgfqpoint{0.964652in}{1.410841in}}%
\pgfpathlineto{\pgfqpoint{0.967940in}{1.408953in}}%
\pgfpathlineto{\pgfqpoint{0.967728in}{1.403968in}}%
\pgfpathclose%
\pgfusepath{fill}%
\end{pgfscope}%
\begin{pgfscope}%
\pgfpathrectangle{\pgfqpoint{0.100000in}{0.100000in}}{\pgfqpoint{3.007045in}{1.925000in}}%
\pgfusepath{clip}%
\pgfsetbuttcap%
\pgfsetmiterjoin%
\definecolor{currentfill}{rgb}{0.657670,0.862053,0.633449}%
\pgfsetfillcolor{currentfill}%
\pgfsetlinewidth{0.000000pt}%
\definecolor{currentstroke}{rgb}{0.000000,0.000000,0.000000}%
\pgfsetstrokecolor{currentstroke}%
\pgfsetstrokeopacity{0.000000}%
\pgfsetdash{}{0pt}%
\pgfpathmoveto{\pgfqpoint{2.359204in}{1.116262in}}%
\pgfpathlineto{\pgfqpoint{2.353384in}{1.121802in}}%
\pgfpathlineto{\pgfqpoint{2.347373in}{1.122044in}}%
\pgfpathlineto{\pgfqpoint{2.339247in}{1.131363in}}%
\pgfpathlineto{\pgfqpoint{2.341607in}{1.134315in}}%
\pgfpathlineto{\pgfqpoint{2.338933in}{1.139975in}}%
\pgfpathlineto{\pgfqpoint{2.332965in}{1.140303in}}%
\pgfpathlineto{\pgfqpoint{2.326824in}{1.138484in}}%
\pgfpathlineto{\pgfqpoint{2.323384in}{1.154002in}}%
\pgfpathlineto{\pgfqpoint{2.331133in}{1.151744in}}%
\pgfpathlineto{\pgfqpoint{2.336776in}{1.153283in}}%
\pgfpathlineto{\pgfqpoint{2.342358in}{1.152328in}}%
\pgfpathlineto{\pgfqpoint{2.348520in}{1.145975in}}%
\pgfpathlineto{\pgfqpoint{2.354011in}{1.144847in}}%
\pgfpathlineto{\pgfqpoint{2.356171in}{1.148869in}}%
\pgfpathlineto{\pgfqpoint{2.360800in}{1.150692in}}%
\pgfpathlineto{\pgfqpoint{2.374431in}{1.146809in}}%
\pgfpathlineto{\pgfqpoint{2.380064in}{1.147699in}}%
\pgfpathlineto{\pgfqpoint{2.385587in}{1.155405in}}%
\pgfpathlineto{\pgfqpoint{2.384672in}{1.146612in}}%
\pgfpathlineto{\pgfqpoint{2.379735in}{1.139390in}}%
\pgfpathlineto{\pgfqpoint{2.377693in}{1.133822in}}%
\pgfpathlineto{\pgfqpoint{2.378692in}{1.128851in}}%
\pgfpathlineto{\pgfqpoint{2.372393in}{1.126861in}}%
\pgfpathlineto{\pgfqpoint{2.368680in}{1.131320in}}%
\pgfpathlineto{\pgfqpoint{2.364836in}{1.127691in}}%
\pgfpathlineto{\pgfqpoint{2.365037in}{1.122682in}}%
\pgfpathlineto{\pgfqpoint{2.359204in}{1.116262in}}%
\pgfpathclose%
\pgfusepath{fill}%
\end{pgfscope}%
\begin{pgfscope}%
\pgfpathrectangle{\pgfqpoint{0.100000in}{0.100000in}}{\pgfqpoint{3.007045in}{1.925000in}}%
\pgfusepath{clip}%
\pgfsetbuttcap%
\pgfsetmiterjoin%
\definecolor{currentfill}{rgb}{0.685721,0.873864,0.660761}%
\pgfsetfillcolor{currentfill}%
\pgfsetlinewidth{0.000000pt}%
\definecolor{currentstroke}{rgb}{0.000000,0.000000,0.000000}%
\pgfsetstrokecolor{currentstroke}%
\pgfsetstrokeopacity{0.000000}%
\pgfsetdash{}{0pt}%
\pgfpathmoveto{\pgfqpoint{2.382677in}{0.742228in}}%
\pgfpathlineto{\pgfqpoint{2.378977in}{0.740714in}}%
\pgfpathlineto{\pgfqpoint{2.375605in}{0.733287in}}%
\pgfpathlineto{\pgfqpoint{2.368826in}{0.730586in}}%
\pgfpathlineto{\pgfqpoint{2.363559in}{0.734763in}}%
\pgfpathlineto{\pgfqpoint{2.360665in}{0.742287in}}%
\pgfpathlineto{\pgfqpoint{2.367977in}{0.755519in}}%
\pgfpathlineto{\pgfqpoint{2.373071in}{0.752288in}}%
\pgfpathlineto{\pgfqpoint{2.375145in}{0.759957in}}%
\pgfpathlineto{\pgfqpoint{2.375546in}{0.768744in}}%
\pgfpathlineto{\pgfqpoint{2.379793in}{0.770306in}}%
\pgfpathlineto{\pgfqpoint{2.377909in}{0.787069in}}%
\pgfpathlineto{\pgfqpoint{2.387448in}{0.788188in}}%
\pgfpathlineto{\pgfqpoint{2.390153in}{0.787078in}}%
\pgfpathlineto{\pgfqpoint{2.390888in}{0.783936in}}%
\pgfpathlineto{\pgfqpoint{2.405437in}{0.788426in}}%
\pgfpathlineto{\pgfqpoint{2.411719in}{0.790134in}}%
\pgfpathlineto{\pgfqpoint{2.417835in}{0.773607in}}%
\pgfpathlineto{\pgfqpoint{2.415340in}{0.771423in}}%
\pgfpathlineto{\pgfqpoint{2.427990in}{0.752079in}}%
\pgfpathlineto{\pgfqpoint{2.434230in}{0.742118in}}%
\pgfpathlineto{\pgfqpoint{2.432076in}{0.743806in}}%
\pgfpathlineto{\pgfqpoint{2.425722in}{0.735586in}}%
\pgfpathlineto{\pgfqpoint{2.422656in}{0.729020in}}%
\pgfpathlineto{\pgfqpoint{2.427863in}{0.724959in}}%
\pgfpathlineto{\pgfqpoint{2.425917in}{0.719822in}}%
\pgfpathlineto{\pgfqpoint{2.410702in}{0.718714in}}%
\pgfpathlineto{\pgfqpoint{2.409404in}{0.729924in}}%
\pgfpathlineto{\pgfqpoint{2.395922in}{0.728401in}}%
\pgfpathlineto{\pgfqpoint{2.395370in}{0.738234in}}%
\pgfpathlineto{\pgfqpoint{2.392692in}{0.741268in}}%
\pgfpathlineto{\pgfqpoint{2.382677in}{0.742228in}}%
\pgfpathclose%
\pgfusepath{fill}%
\end{pgfscope}%
\begin{pgfscope}%
\pgfpathrectangle{\pgfqpoint{0.100000in}{0.100000in}}{\pgfqpoint{3.007045in}{1.925000in}}%
\pgfusepath{clip}%
\pgfsetbuttcap%
\pgfsetmiterjoin%
\definecolor{currentfill}{rgb}{0.474279,0.777670,0.478677}%
\pgfsetfillcolor{currentfill}%
\pgfsetlinewidth{0.000000pt}%
\definecolor{currentstroke}{rgb}{0.000000,0.000000,0.000000}%
\pgfsetstrokecolor{currentstroke}%
\pgfsetstrokeopacity{0.000000}%
\pgfsetdash{}{0pt}%
\pgfpathmoveto{\pgfqpoint{1.668319in}{1.540277in}}%
\pgfpathlineto{\pgfqpoint{1.633052in}{1.541386in}}%
\pgfpathlineto{\pgfqpoint{1.634891in}{1.593241in}}%
\pgfpathlineto{\pgfqpoint{1.669507in}{1.592201in}}%
\pgfpathlineto{\pgfqpoint{1.700191in}{1.591568in}}%
\pgfpathlineto{\pgfqpoint{1.699144in}{1.583940in}}%
\pgfpathlineto{\pgfqpoint{1.695859in}{1.579141in}}%
\pgfpathlineto{\pgfqpoint{1.687099in}{1.572723in}}%
\pgfpathlineto{\pgfqpoint{1.686157in}{1.569986in}}%
\pgfpathlineto{\pgfqpoint{1.694140in}{1.556870in}}%
\pgfpathlineto{\pgfqpoint{1.701484in}{1.554430in}}%
\pgfpathlineto{\pgfqpoint{1.703819in}{1.551137in}}%
\pgfpathlineto{\pgfqpoint{1.679482in}{1.551623in}}%
\pgfpathlineto{\pgfqpoint{1.678758in}{1.549649in}}%
\pgfpathlineto{\pgfqpoint{1.668554in}{1.549944in}}%
\pgfpathlineto{\pgfqpoint{1.668319in}{1.540277in}}%
\pgfpathclose%
\pgfusepath{fill}%
\end{pgfscope}%
\begin{pgfscope}%
\pgfpathrectangle{\pgfqpoint{0.100000in}{0.100000in}}{\pgfqpoint{3.007045in}{1.925000in}}%
\pgfusepath{clip}%
\pgfsetbuttcap%
\pgfsetmiterjoin%
\definecolor{currentfill}{rgb}{0.401569,0.742484,0.436601}%
\pgfsetfillcolor{currentfill}%
\pgfsetlinewidth{0.000000pt}%
\definecolor{currentstroke}{rgb}{0.000000,0.000000,0.000000}%
\pgfsetstrokecolor{currentstroke}%
\pgfsetstrokeopacity{0.000000}%
\pgfsetdash{}{0pt}%
\pgfpathmoveto{\pgfqpoint{2.493292in}{1.298554in}}%
\pgfpathlineto{\pgfqpoint{2.486521in}{1.293715in}}%
\pgfpathlineto{\pgfqpoint{2.476757in}{1.293301in}}%
\pgfpathlineto{\pgfqpoint{2.473535in}{1.295735in}}%
\pgfpathlineto{\pgfqpoint{2.472733in}{1.301176in}}%
\pgfpathlineto{\pgfqpoint{2.464330in}{1.299934in}}%
\pgfpathlineto{\pgfqpoint{2.461738in}{1.316890in}}%
\pgfpathlineto{\pgfqpoint{2.466018in}{1.317526in}}%
\pgfpathlineto{\pgfqpoint{2.460651in}{1.350974in}}%
\pgfpathlineto{\pgfqpoint{2.484276in}{1.354588in}}%
\pgfpathlineto{\pgfqpoint{2.503675in}{1.356445in}}%
\pgfpathlineto{\pgfqpoint{2.509694in}{1.358087in}}%
\pgfpathlineto{\pgfqpoint{2.513140in}{1.337398in}}%
\pgfpathlineto{\pgfqpoint{2.506874in}{1.324911in}}%
\pgfpathlineto{\pgfqpoint{2.508763in}{1.315516in}}%
\pgfpathlineto{\pgfqpoint{2.491069in}{1.312382in}}%
\pgfpathlineto{\pgfqpoint{2.493292in}{1.298554in}}%
\pgfpathclose%
\pgfusepath{fill}%
\end{pgfscope}%
\begin{pgfscope}%
\pgfpathrectangle{\pgfqpoint{0.100000in}{0.100000in}}{\pgfqpoint{3.007045in}{1.925000in}}%
\pgfusepath{clip}%
\pgfsetbuttcap%
\pgfsetmiterjoin%
\definecolor{currentfill}{rgb}{0.490888,0.785421,0.492334}%
\pgfsetfillcolor{currentfill}%
\pgfsetlinewidth{0.000000pt}%
\definecolor{currentstroke}{rgb}{0.000000,0.000000,0.000000}%
\pgfsetstrokecolor{currentstroke}%
\pgfsetstrokeopacity{0.000000}%
\pgfsetdash{}{0pt}%
\pgfpathmoveto{\pgfqpoint{1.386507in}{1.413753in}}%
\pgfpathlineto{\pgfqpoint{1.335923in}{1.418503in}}%
\pgfpathlineto{\pgfqpoint{1.339144in}{1.451533in}}%
\pgfpathlineto{\pgfqpoint{1.343611in}{1.496079in}}%
\pgfpathlineto{\pgfqpoint{1.348840in}{1.549898in}}%
\pgfpathlineto{\pgfqpoint{1.351084in}{1.564019in}}%
\pgfpathlineto{\pgfqpoint{1.401374in}{1.559284in}}%
\pgfpathlineto{\pgfqpoint{1.400379in}{1.547840in}}%
\pgfpathlineto{\pgfqpoint{1.445034in}{1.544102in}}%
\pgfpathlineto{\pgfqpoint{1.442332in}{1.509325in}}%
\pgfpathlineto{\pgfqpoint{1.439374in}{1.475257in}}%
\pgfpathlineto{\pgfqpoint{1.437357in}{1.456409in}}%
\pgfpathlineto{\pgfqpoint{1.431599in}{1.456389in}}%
\pgfpathlineto{\pgfqpoint{1.429680in}{1.455658in}}%
\pgfpathlineto{\pgfqpoint{1.399572in}{1.458189in}}%
\pgfpathlineto{\pgfqpoint{1.395146in}{1.457695in}}%
\pgfpathlineto{\pgfqpoint{1.390013in}{1.453572in}}%
\pgfpathlineto{\pgfqpoint{1.386507in}{1.413753in}}%
\pgfpathclose%
\pgfusepath{fill}%
\end{pgfscope}%
\begin{pgfscope}%
\pgfpathrectangle{\pgfqpoint{0.100000in}{0.100000in}}{\pgfqpoint{3.007045in}{1.925000in}}%
\pgfusepath{clip}%
\pgfsetbuttcap%
\pgfsetmiterjoin%
\definecolor{currentfill}{rgb}{0.623760,0.847428,0.601584}%
\pgfsetfillcolor{currentfill}%
\pgfsetlinewidth{0.000000pt}%
\definecolor{currentstroke}{rgb}{0.000000,0.000000,0.000000}%
\pgfsetstrokecolor{currentstroke}%
\pgfsetstrokeopacity{0.000000}%
\pgfsetdash{}{0pt}%
\pgfpathmoveto{\pgfqpoint{2.406434in}{1.016470in}}%
\pgfpathlineto{\pgfqpoint{2.399781in}{1.010130in}}%
\pgfpathlineto{\pgfqpoint{2.396925in}{1.003949in}}%
\pgfpathlineto{\pgfqpoint{2.392615in}{1.001420in}}%
\pgfpathlineto{\pgfqpoint{2.388640in}{1.002575in}}%
\pgfpathlineto{\pgfqpoint{2.380497in}{0.997923in}}%
\pgfpathlineto{\pgfqpoint{2.372662in}{0.995532in}}%
\pgfpathlineto{\pgfqpoint{2.361719in}{1.000172in}}%
\pgfpathlineto{\pgfqpoint{2.358542in}{0.999323in}}%
\pgfpathlineto{\pgfqpoint{2.353818in}{1.009495in}}%
\pgfpathlineto{\pgfqpoint{2.355024in}{1.013333in}}%
\pgfpathlineto{\pgfqpoint{2.360603in}{1.018528in}}%
\pgfpathlineto{\pgfqpoint{2.362675in}{1.024583in}}%
\pgfpathlineto{\pgfqpoint{2.376071in}{1.035433in}}%
\pgfpathlineto{\pgfqpoint{2.378235in}{1.029769in}}%
\pgfpathlineto{\pgfqpoint{2.378204in}{1.022084in}}%
\pgfpathlineto{\pgfqpoint{2.380873in}{1.017870in}}%
\pgfpathlineto{\pgfqpoint{2.373385in}{1.015073in}}%
\pgfpathlineto{\pgfqpoint{2.372965in}{1.012422in}}%
\pgfpathlineto{\pgfqpoint{2.406434in}{1.016470in}}%
\pgfpathclose%
\pgfusepath{fill}%
\end{pgfscope}%
\begin{pgfscope}%
\pgfpathrectangle{\pgfqpoint{0.100000in}{0.100000in}}{\pgfqpoint{3.007045in}{1.925000in}}%
\pgfusepath{clip}%
\pgfsetbuttcap%
\pgfsetmiterjoin%
\definecolor{currentfill}{rgb}{0.193541,0.605136,0.315617}%
\pgfsetfillcolor{currentfill}%
\pgfsetlinewidth{0.000000pt}%
\definecolor{currentstroke}{rgb}{0.000000,0.000000,0.000000}%
\pgfsetstrokecolor{currentstroke}%
\pgfsetstrokeopacity{0.000000}%
\pgfsetdash{}{0pt}%
\pgfpathmoveto{\pgfqpoint{2.385967in}{0.896018in}}%
\pgfpathlineto{\pgfqpoint{2.377707in}{0.892532in}}%
\pgfpathlineto{\pgfqpoint{2.375790in}{0.888043in}}%
\pgfpathlineto{\pgfqpoint{2.371189in}{0.888722in}}%
\pgfpathlineto{\pgfqpoint{2.366727in}{0.881594in}}%
\pgfpathlineto{\pgfqpoint{2.360142in}{0.881259in}}%
\pgfpathlineto{\pgfqpoint{2.362404in}{0.891696in}}%
\pgfpathlineto{\pgfqpoint{2.356919in}{0.901553in}}%
\pgfpathlineto{\pgfqpoint{2.348646in}{0.902478in}}%
\pgfpathlineto{\pgfqpoint{2.348504in}{0.918229in}}%
\pgfpathlineto{\pgfqpoint{2.351726in}{0.921503in}}%
\pgfpathlineto{\pgfqpoint{2.356471in}{0.920502in}}%
\pgfpathlineto{\pgfqpoint{2.361956in}{0.924269in}}%
\pgfpathlineto{\pgfqpoint{2.366318in}{0.919795in}}%
\pgfpathlineto{\pgfqpoint{2.372864in}{0.923192in}}%
\pgfpathlineto{\pgfqpoint{2.379773in}{0.923563in}}%
\pgfpathlineto{\pgfqpoint{2.378628in}{0.917255in}}%
\pgfpathlineto{\pgfqpoint{2.383505in}{0.917672in}}%
\pgfpathlineto{\pgfqpoint{2.391763in}{0.908652in}}%
\pgfpathlineto{\pgfqpoint{2.385392in}{0.899513in}}%
\pgfpathlineto{\pgfqpoint{2.385967in}{0.896018in}}%
\pgfpathclose%
\pgfusepath{fill}%
\end{pgfscope}%
\begin{pgfscope}%
\pgfpathrectangle{\pgfqpoint{0.100000in}{0.100000in}}{\pgfqpoint{3.007045in}{1.925000in}}%
\pgfusepath{clip}%
\pgfsetbuttcap%
\pgfsetmiterjoin%
\definecolor{currentfill}{rgb}{0.546251,0.811257,0.537855}%
\pgfsetfillcolor{currentfill}%
\pgfsetlinewidth{0.000000pt}%
\definecolor{currentstroke}{rgb}{0.000000,0.000000,0.000000}%
\pgfsetstrokecolor{currentstroke}%
\pgfsetstrokeopacity{0.000000}%
\pgfsetdash{}{0pt}%
\pgfpathmoveto{\pgfqpoint{2.410399in}{1.309530in}}%
\pgfpathlineto{\pgfqpoint{2.407666in}{1.313899in}}%
\pgfpathlineto{\pgfqpoint{2.407064in}{1.318707in}}%
\pgfpathlineto{\pgfqpoint{2.411263in}{1.323487in}}%
\pgfpathlineto{\pgfqpoint{2.416075in}{1.324178in}}%
\pgfpathlineto{\pgfqpoint{2.415460in}{1.329120in}}%
\pgfpathlineto{\pgfqpoint{2.420098in}{1.329773in}}%
\pgfpathlineto{\pgfqpoint{2.419430in}{1.334702in}}%
\pgfpathlineto{\pgfqpoint{2.414804in}{1.334085in}}%
\pgfpathlineto{\pgfqpoint{2.413516in}{1.344212in}}%
\pgfpathlineto{\pgfqpoint{2.424873in}{1.344692in}}%
\pgfpathlineto{\pgfqpoint{2.430966in}{1.350397in}}%
\pgfpathlineto{\pgfqpoint{2.437543in}{1.358914in}}%
\pgfpathlineto{\pgfqpoint{2.456731in}{1.373827in}}%
\pgfpathlineto{\pgfqpoint{2.461753in}{1.375576in}}%
\pgfpathlineto{\pgfqpoint{2.479266in}{1.385678in}}%
\pgfpathlineto{\pgfqpoint{2.484276in}{1.354588in}}%
\pgfpathlineto{\pgfqpoint{2.460651in}{1.350974in}}%
\pgfpathlineto{\pgfqpoint{2.466018in}{1.317526in}}%
\pgfpathlineto{\pgfqpoint{2.461738in}{1.316890in}}%
\pgfpathlineto{\pgfqpoint{2.445530in}{1.314439in}}%
\pgfpathlineto{\pgfqpoint{2.446146in}{1.309080in}}%
\pgfpathlineto{\pgfqpoint{2.434855in}{1.307908in}}%
\pgfpathlineto{\pgfqpoint{2.434169in}{1.312710in}}%
\pgfpathlineto{\pgfqpoint{2.410399in}{1.309530in}}%
\pgfpathclose%
\pgfusepath{fill}%
\end{pgfscope}%
\begin{pgfscope}%
\pgfpathrectangle{\pgfqpoint{0.100000in}{0.100000in}}{\pgfqpoint{3.007045in}{1.925000in}}%
\pgfusepath{clip}%
\pgfsetbuttcap%
\pgfsetmiterjoin%
\definecolor{currentfill}{rgb}{0.723122,0.889612,0.697178}%
\pgfsetfillcolor{currentfill}%
\pgfsetlinewidth{0.000000pt}%
\definecolor{currentstroke}{rgb}{0.000000,0.000000,0.000000}%
\pgfsetstrokecolor{currentstroke}%
\pgfsetstrokeopacity{0.000000}%
\pgfsetdash{}{0pt}%
\pgfpathmoveto{\pgfqpoint{0.630991in}{1.684466in}}%
\pgfpathlineto{\pgfqpoint{0.625624in}{1.685979in}}%
\pgfpathlineto{\pgfqpoint{0.630102in}{1.696622in}}%
\pgfpathlineto{\pgfqpoint{0.624611in}{1.698303in}}%
\pgfpathlineto{\pgfqpoint{0.626202in}{1.703795in}}%
\pgfpathlineto{\pgfqpoint{0.620630in}{1.705467in}}%
\pgfpathlineto{\pgfqpoint{0.630951in}{1.740790in}}%
\pgfpathlineto{\pgfqpoint{0.637219in}{1.740644in}}%
\pgfpathlineto{\pgfqpoint{0.640972in}{1.753699in}}%
\pgfpathlineto{\pgfqpoint{0.651182in}{1.791391in}}%
\pgfpathlineto{\pgfqpoint{0.662371in}{1.789397in}}%
\pgfpathlineto{\pgfqpoint{0.668689in}{1.793331in}}%
\pgfpathlineto{\pgfqpoint{0.670581in}{1.789487in}}%
\pgfpathlineto{\pgfqpoint{0.675279in}{1.792196in}}%
\pgfpathlineto{\pgfqpoint{0.702590in}{1.784664in}}%
\pgfpathlineto{\pgfqpoint{0.726344in}{1.778628in}}%
\pgfpathlineto{\pgfqpoint{0.724517in}{1.778735in}}%
\pgfpathlineto{\pgfqpoint{0.724939in}{1.766494in}}%
\pgfpathlineto{\pgfqpoint{0.740246in}{1.767188in}}%
\pgfpathlineto{\pgfqpoint{0.737197in}{1.757418in}}%
\pgfpathlineto{\pgfqpoint{0.742787in}{1.756014in}}%
\pgfpathlineto{\pgfqpoint{0.743350in}{1.746702in}}%
\pgfpathlineto{\pgfqpoint{0.746177in}{1.746042in}}%
\pgfpathlineto{\pgfqpoint{0.741011in}{1.724091in}}%
\pgfpathlineto{\pgfqpoint{0.724313in}{1.728496in}}%
\pgfpathlineto{\pgfqpoint{0.722204in}{1.719436in}}%
\pgfpathlineto{\pgfqpoint{0.718256in}{1.717477in}}%
\pgfpathlineto{\pgfqpoint{0.715135in}{1.709465in}}%
\pgfpathlineto{\pgfqpoint{0.712881in}{1.710060in}}%
\pgfpathlineto{\pgfqpoint{0.709253in}{1.696086in}}%
\pgfpathlineto{\pgfqpoint{0.704883in}{1.694237in}}%
\pgfpathlineto{\pgfqpoint{0.694542in}{1.697073in}}%
\pgfpathlineto{\pgfqpoint{0.693304in}{1.692341in}}%
\pgfpathlineto{\pgfqpoint{0.681006in}{1.694917in}}%
\pgfpathlineto{\pgfqpoint{0.680362in}{1.684919in}}%
\pgfpathlineto{\pgfqpoint{0.685286in}{1.683617in}}%
\pgfpathlineto{\pgfqpoint{0.683125in}{1.670416in}}%
\pgfpathlineto{\pgfqpoint{0.630991in}{1.684466in}}%
\pgfpathclose%
\pgfusepath{fill}%
\end{pgfscope}%
\begin{pgfscope}%
\pgfpathrectangle{\pgfqpoint{0.100000in}{0.100000in}}{\pgfqpoint{3.007045in}{1.925000in}}%
\pgfusepath{clip}%
\pgfsetbuttcap%
\pgfsetmiterjoin%
\definecolor{currentfill}{rgb}{0.401569,0.742484,0.436601}%
\pgfsetfillcolor{currentfill}%
\pgfsetlinewidth{0.000000pt}%
\definecolor{currentstroke}{rgb}{0.000000,0.000000,0.000000}%
\pgfsetstrokecolor{currentstroke}%
\pgfsetstrokeopacity{0.000000}%
\pgfsetdash{}{0pt}%
\pgfpathmoveto{\pgfqpoint{1.351564in}{1.062383in}}%
\pgfpathlineto{\pgfqpoint{1.347270in}{1.020215in}}%
\pgfpathlineto{\pgfqpoint{1.298942in}{1.024507in}}%
\pgfpathlineto{\pgfqpoint{1.277959in}{1.026457in}}%
\pgfpathlineto{\pgfqpoint{1.238883in}{1.030797in}}%
\pgfpathlineto{\pgfqpoint{1.241118in}{1.050301in}}%
\pgfpathlineto{\pgfqpoint{1.242720in}{1.057254in}}%
\pgfpathlineto{\pgfqpoint{1.241493in}{1.067690in}}%
\pgfpathlineto{\pgfqpoint{1.236457in}{1.074727in}}%
\pgfpathlineto{\pgfqpoint{1.225618in}{1.073383in}}%
\pgfpathlineto{\pgfqpoint{1.231022in}{1.088493in}}%
\pgfpathlineto{\pgfqpoint{1.229068in}{1.091653in}}%
\pgfpathlineto{\pgfqpoint{1.232103in}{1.090942in}}%
\pgfpathlineto{\pgfqpoint{1.241682in}{1.094643in}}%
\pgfpathlineto{\pgfqpoint{1.245702in}{1.097910in}}%
\pgfpathlineto{\pgfqpoint{1.251074in}{1.090450in}}%
\pgfpathlineto{\pgfqpoint{1.252696in}{1.088003in}}%
\pgfpathlineto{\pgfqpoint{1.271819in}{1.087203in}}%
\pgfpathlineto{\pgfqpoint{1.301086in}{1.073162in}}%
\pgfpathlineto{\pgfqpoint{1.300552in}{1.067248in}}%
\pgfpathlineto{\pgfqpoint{1.351564in}{1.062383in}}%
\pgfpathclose%
\pgfusepath{fill}%
\end{pgfscope}%
\begin{pgfscope}%
\pgfpathrectangle{\pgfqpoint{0.100000in}{0.100000in}}{\pgfqpoint{3.007045in}{1.925000in}}%
\pgfusepath{clip}%
\pgfsetbuttcap%
\pgfsetmiterjoin%
\definecolor{currentfill}{rgb}{0.643645,0.856148,0.619792}%
\pgfsetfillcolor{currentfill}%
\pgfsetlinewidth{0.000000pt}%
\definecolor{currentstroke}{rgb}{0.000000,0.000000,0.000000}%
\pgfsetstrokecolor{currentstroke}%
\pgfsetstrokeopacity{0.000000}%
\pgfsetdash{}{0pt}%
\pgfpathmoveto{\pgfqpoint{1.972796in}{1.051471in}}%
\pgfpathlineto{\pgfqpoint{1.964595in}{1.051309in}}%
\pgfpathlineto{\pgfqpoint{1.964484in}{1.057050in}}%
\pgfpathlineto{\pgfqpoint{1.953009in}{1.056717in}}%
\pgfpathlineto{\pgfqpoint{1.952290in}{1.080650in}}%
\pgfpathlineto{\pgfqpoint{1.946082in}{1.080499in}}%
\pgfpathlineto{\pgfqpoint{1.945763in}{1.089375in}}%
\pgfpathlineto{\pgfqpoint{1.944978in}{1.116719in}}%
\pgfpathlineto{\pgfqpoint{1.949655in}{1.115080in}}%
\pgfpathlineto{\pgfqpoint{1.956351in}{1.117548in}}%
\pgfpathlineto{\pgfqpoint{1.956143in}{1.126474in}}%
\pgfpathlineto{\pgfqpoint{1.963877in}{1.126688in}}%
\pgfpathlineto{\pgfqpoint{1.963320in}{1.146179in}}%
\pgfpathlineto{\pgfqpoint{1.967135in}{1.146275in}}%
\pgfpathlineto{\pgfqpoint{1.967023in}{1.152094in}}%
\pgfpathlineto{\pgfqpoint{1.990535in}{1.152936in}}%
\pgfpathlineto{\pgfqpoint{1.992952in}{1.144098in}}%
\pgfpathlineto{\pgfqpoint{1.991482in}{1.141762in}}%
\pgfpathlineto{\pgfqpoint{1.994783in}{1.133102in}}%
\pgfpathlineto{\pgfqpoint{2.000724in}{1.130408in}}%
\pgfpathlineto{\pgfqpoint{2.005193in}{1.136744in}}%
\pgfpathlineto{\pgfqpoint{2.015419in}{1.134149in}}%
\pgfpathlineto{\pgfqpoint{2.022871in}{1.129850in}}%
\pgfpathlineto{\pgfqpoint{2.018572in}{1.121893in}}%
\pgfpathlineto{\pgfqpoint{2.020049in}{1.117235in}}%
\pgfpathlineto{\pgfqpoint{2.044424in}{1.118250in}}%
\pgfpathlineto{\pgfqpoint{2.046290in}{1.089601in}}%
\pgfpathlineto{\pgfqpoint{2.029064in}{1.088879in}}%
\pgfpathlineto{\pgfqpoint{2.029399in}{1.083124in}}%
\pgfpathlineto{\pgfqpoint{2.020734in}{1.079515in}}%
\pgfpathlineto{\pgfqpoint{2.014367in}{1.079611in}}%
\pgfpathlineto{\pgfqpoint{2.010046in}{1.075951in}}%
\pgfpathlineto{\pgfqpoint{2.001501in}{1.072967in}}%
\pgfpathlineto{\pgfqpoint{1.990680in}{1.085683in}}%
\pgfpathlineto{\pgfqpoint{1.974361in}{1.084952in}}%
\pgfpathlineto{\pgfqpoint{1.975453in}{1.054398in}}%
\pgfpathlineto{\pgfqpoint{1.972796in}{1.051471in}}%
\pgfpathclose%
\pgfusepath{fill}%
\end{pgfscope}%
\begin{pgfscope}%
\pgfpathrectangle{\pgfqpoint{0.100000in}{0.100000in}}{\pgfqpoint{3.007045in}{1.925000in}}%
\pgfusepath{clip}%
\pgfsetbuttcap%
\pgfsetmiterjoin%
\definecolor{currentfill}{rgb}{0.732472,0.893549,0.706282}%
\pgfsetfillcolor{currentfill}%
\pgfsetlinewidth{0.000000pt}%
\definecolor{currentstroke}{rgb}{0.000000,0.000000,0.000000}%
\pgfsetstrokecolor{currentstroke}%
\pgfsetstrokeopacity{0.000000}%
\pgfsetdash{}{0pt}%
\pgfpathmoveto{\pgfqpoint{1.792356in}{1.069644in}}%
\pgfpathlineto{\pgfqpoint{1.768296in}{1.069654in}}%
\pgfpathlineto{\pgfqpoint{1.768960in}{1.084214in}}%
\pgfpathlineto{\pgfqpoint{1.768947in}{1.112944in}}%
\pgfpathlineto{\pgfqpoint{1.769426in}{1.115817in}}%
\pgfpathlineto{\pgfqpoint{1.769446in}{1.131864in}}%
\pgfpathlineto{\pgfqpoint{1.776978in}{1.132450in}}%
\pgfpathlineto{\pgfqpoint{1.777411in}{1.146354in}}%
\pgfpathlineto{\pgfqpoint{1.780164in}{1.150731in}}%
\pgfpathlineto{\pgfqpoint{1.772963in}{1.162326in}}%
\pgfpathlineto{\pgfqpoint{1.767101in}{1.168173in}}%
\pgfpathlineto{\pgfqpoint{1.792583in}{1.168136in}}%
\pgfpathlineto{\pgfqpoint{1.792627in}{1.163111in}}%
\pgfpathlineto{\pgfqpoint{1.812471in}{1.163163in}}%
\pgfpathlineto{\pgfqpoint{1.812537in}{1.167939in}}%
\pgfpathlineto{\pgfqpoint{1.835366in}{1.168017in}}%
\pgfpathlineto{\pgfqpoint{1.835689in}{1.147108in}}%
\pgfpathlineto{\pgfqpoint{1.840824in}{1.149749in}}%
\pgfpathlineto{\pgfqpoint{1.848856in}{1.148030in}}%
\pgfpathlineto{\pgfqpoint{1.849349in}{1.129915in}}%
\pgfpathlineto{\pgfqpoint{1.849007in}{1.104433in}}%
\pgfpathlineto{\pgfqpoint{1.820514in}{1.104792in}}%
\pgfpathlineto{\pgfqpoint{1.820190in}{1.081544in}}%
\pgfpathlineto{\pgfqpoint{1.821187in}{1.069825in}}%
\pgfpathlineto{\pgfqpoint{1.811752in}{1.070677in}}%
\pgfpathlineto{\pgfqpoint{1.792345in}{1.071169in}}%
\pgfpathlineto{\pgfqpoint{1.792356in}{1.069644in}}%
\pgfpathclose%
\pgfusepath{fill}%
\end{pgfscope}%
\begin{pgfscope}%
\pgfpathrectangle{\pgfqpoint{0.100000in}{0.100000in}}{\pgfqpoint{3.007045in}{1.925000in}}%
\pgfusepath{clip}%
\pgfsetbuttcap%
\pgfsetmiterjoin%
\definecolor{currentfill}{rgb}{0.585006,0.829343,0.569719}%
\pgfsetfillcolor{currentfill}%
\pgfsetlinewidth{0.000000pt}%
\definecolor{currentstroke}{rgb}{0.000000,0.000000,0.000000}%
\pgfsetstrokecolor{currentstroke}%
\pgfsetstrokeopacity{0.000000}%
\pgfsetdash{}{0pt}%
\pgfpathmoveto{\pgfqpoint{1.285700in}{1.824835in}}%
\pgfpathlineto{\pgfqpoint{1.331685in}{1.819464in}}%
\pgfpathlineto{\pgfqpoint{1.331869in}{1.813477in}}%
\pgfpathlineto{\pgfqpoint{1.329938in}{1.796177in}}%
\pgfpathlineto{\pgfqpoint{1.332164in}{1.790080in}}%
\pgfpathlineto{\pgfqpoint{1.295653in}{1.794279in}}%
\pgfpathlineto{\pgfqpoint{1.286131in}{1.795421in}}%
\pgfpathlineto{\pgfqpoint{1.288237in}{1.812721in}}%
\pgfpathlineto{\pgfqpoint{1.282530in}{1.813419in}}%
\pgfpathlineto{\pgfqpoint{1.285700in}{1.824835in}}%
\pgfpathclose%
\pgfusepath{fill}%
\end{pgfscope}%
\begin{pgfscope}%
\pgfpathrectangle{\pgfqpoint{0.100000in}{0.100000in}}{\pgfqpoint{3.007045in}{1.925000in}}%
\pgfusepath{clip}%
\pgfsetbuttcap%
\pgfsetmiterjoin%
\definecolor{currentfill}{rgb}{0.676371,0.869927,0.651657}%
\pgfsetfillcolor{currentfill}%
\pgfsetlinewidth{0.000000pt}%
\definecolor{currentstroke}{rgb}{0.000000,0.000000,0.000000}%
\pgfsetstrokecolor{currentstroke}%
\pgfsetstrokeopacity{0.000000}%
\pgfsetdash{}{0pt}%
\pgfpathmoveto{\pgfqpoint{2.568608in}{0.576327in}}%
\pgfpathlineto{\pgfqpoint{2.562689in}{0.572385in}}%
\pgfpathlineto{\pgfqpoint{2.551269in}{0.570455in}}%
\pgfpathlineto{\pgfqpoint{2.549830in}{0.579481in}}%
\pgfpathlineto{\pgfqpoint{2.545862in}{0.584302in}}%
\pgfpathlineto{\pgfqpoint{2.532632in}{0.582139in}}%
\pgfpathlineto{\pgfqpoint{2.526180in}{0.575201in}}%
\pgfpathlineto{\pgfqpoint{2.520191in}{0.572396in}}%
\pgfpathlineto{\pgfqpoint{2.516161in}{0.600413in}}%
\pgfpathlineto{\pgfqpoint{2.492884in}{0.596677in}}%
\pgfpathlineto{\pgfqpoint{2.488716in}{0.626214in}}%
\pgfpathlineto{\pgfqpoint{2.502750in}{0.627127in}}%
\pgfpathlineto{\pgfqpoint{2.504409in}{0.641856in}}%
\pgfpathlineto{\pgfqpoint{2.488999in}{0.640420in}}%
\pgfpathlineto{\pgfqpoint{2.487249in}{0.654855in}}%
\pgfpathlineto{\pgfqpoint{2.503300in}{0.656930in}}%
\pgfpathlineto{\pgfqpoint{2.506695in}{0.661778in}}%
\pgfpathlineto{\pgfqpoint{2.514003in}{0.660961in}}%
\pgfpathlineto{\pgfqpoint{2.521337in}{0.669031in}}%
\pgfpathlineto{\pgfqpoint{2.533990in}{0.667284in}}%
\pgfpathlineto{\pgfqpoint{2.542300in}{0.662555in}}%
\pgfpathlineto{\pgfqpoint{2.545263in}{0.656150in}}%
\pgfpathlineto{\pgfqpoint{2.543774in}{0.645180in}}%
\pgfpathlineto{\pgfqpoint{2.546129in}{0.642161in}}%
\pgfpathlineto{\pgfqpoint{2.546297in}{0.634841in}}%
\pgfpathlineto{\pgfqpoint{2.551061in}{0.622627in}}%
\pgfpathlineto{\pgfqpoint{2.552074in}{0.616494in}}%
\pgfpathlineto{\pgfqpoint{2.560429in}{0.594452in}}%
\pgfpathlineto{\pgfqpoint{2.568608in}{0.576327in}}%
\pgfpathclose%
\pgfusepath{fill}%
\end{pgfscope}%
\begin{pgfscope}%
\pgfpathrectangle{\pgfqpoint{0.100000in}{0.100000in}}{\pgfqpoint{3.007045in}{1.925000in}}%
\pgfusepath{clip}%
\pgfsetbuttcap%
\pgfsetmiterjoin%
\definecolor{currentfill}{rgb}{0.618224,0.844844,0.597032}%
\pgfsetfillcolor{currentfill}%
\pgfsetlinewidth{0.000000pt}%
\definecolor{currentstroke}{rgb}{0.000000,0.000000,0.000000}%
\pgfsetstrokecolor{currentstroke}%
\pgfsetstrokeopacity{0.000000}%
\pgfsetdash{}{0pt}%
\pgfpathmoveto{\pgfqpoint{1.342266in}{0.937180in}}%
\pgfpathlineto{\pgfqpoint{1.341578in}{0.929505in}}%
\pgfpathlineto{\pgfqpoint{1.337618in}{0.885562in}}%
\pgfpathlineto{\pgfqpoint{1.324670in}{0.886744in}}%
\pgfpathlineto{\pgfqpoint{1.324074in}{0.880866in}}%
\pgfpathlineto{\pgfqpoint{1.318390in}{0.881414in}}%
\pgfpathlineto{\pgfqpoint{1.317814in}{0.875699in}}%
\pgfpathlineto{\pgfqpoint{1.312137in}{0.876255in}}%
\pgfpathlineto{\pgfqpoint{1.311575in}{0.870525in}}%
\pgfpathlineto{\pgfqpoint{1.300241in}{0.871662in}}%
\pgfpathlineto{\pgfqpoint{1.299641in}{0.865888in}}%
\pgfpathlineto{\pgfqpoint{1.286524in}{0.867195in}}%
\pgfpathlineto{\pgfqpoint{1.276783in}{0.868180in}}%
\pgfpathlineto{\pgfqpoint{1.280556in}{0.903442in}}%
\pgfpathlineto{\pgfqpoint{1.295559in}{0.908588in}}%
\pgfpathlineto{\pgfqpoint{1.307382in}{0.907359in}}%
\pgfpathlineto{\pgfqpoint{1.308509in}{0.919262in}}%
\pgfpathlineto{\pgfqpoint{1.303832in}{0.919724in}}%
\pgfpathlineto{\pgfqpoint{1.292797in}{0.945962in}}%
\pgfpathlineto{\pgfqpoint{1.290299in}{0.944870in}}%
\pgfpathlineto{\pgfqpoint{1.271875in}{0.946736in}}%
\pgfpathlineto{\pgfqpoint{1.273425in}{0.954257in}}%
\pgfpathlineto{\pgfqpoint{1.273752in}{0.963960in}}%
\pgfpathlineto{\pgfqpoint{1.271190in}{0.975810in}}%
\pgfpathlineto{\pgfqpoint{1.293808in}{0.973457in}}%
\pgfpathlineto{\pgfqpoint{1.293521in}{0.970602in}}%
\pgfpathlineto{\pgfqpoint{1.304909in}{0.969450in}}%
\pgfpathlineto{\pgfqpoint{1.304327in}{0.963661in}}%
\pgfpathlineto{\pgfqpoint{1.327114in}{0.961542in}}%
\pgfpathlineto{\pgfqpoint{1.324480in}{0.938840in}}%
\pgfpathlineto{\pgfqpoint{1.342266in}{0.937180in}}%
\pgfpathclose%
\pgfusepath{fill}%
\end{pgfscope}%
\begin{pgfscope}%
\pgfpathrectangle{\pgfqpoint{0.100000in}{0.100000in}}{\pgfqpoint{3.007045in}{1.925000in}}%
\pgfusepath{clip}%
\pgfsetbuttcap%
\pgfsetmiterjoin%
\definecolor{currentfill}{rgb}{0.540715,0.808674,0.533303}%
\pgfsetfillcolor{currentfill}%
\pgfsetlinewidth{0.000000pt}%
\definecolor{currentstroke}{rgb}{0.000000,0.000000,0.000000}%
\pgfsetstrokecolor{currentstroke}%
\pgfsetstrokeopacity{0.000000}%
\pgfsetdash{}{0pt}%
\pgfpathmoveto{\pgfqpoint{2.625680in}{1.382482in}}%
\pgfpathlineto{\pgfqpoint{2.619453in}{1.411977in}}%
\pgfpathlineto{\pgfqpoint{2.612717in}{1.410639in}}%
\pgfpathlineto{\pgfqpoint{2.608126in}{1.441492in}}%
\pgfpathlineto{\pgfqpoint{2.609773in}{1.449142in}}%
\pgfpathlineto{\pgfqpoint{2.634427in}{1.453744in}}%
\pgfpathlineto{\pgfqpoint{2.635144in}{1.448720in}}%
\pgfpathlineto{\pgfqpoint{2.647379in}{1.448629in}}%
\pgfpathlineto{\pgfqpoint{2.646021in}{1.453686in}}%
\pgfpathlineto{\pgfqpoint{2.661352in}{1.457374in}}%
\pgfpathlineto{\pgfqpoint{2.656543in}{1.461343in}}%
\pgfpathlineto{\pgfqpoint{2.675608in}{1.465295in}}%
\pgfpathlineto{\pgfqpoint{2.679109in}{1.451464in}}%
\pgfpathlineto{\pgfqpoint{2.680763in}{1.444336in}}%
\pgfpathlineto{\pgfqpoint{2.672580in}{1.444093in}}%
\pgfpathlineto{\pgfqpoint{2.667190in}{1.440512in}}%
\pgfpathlineto{\pgfqpoint{2.669988in}{1.422124in}}%
\pgfpathlineto{\pgfqpoint{2.652249in}{1.418613in}}%
\pgfpathlineto{\pgfqpoint{2.660061in}{1.392938in}}%
\pgfpathlineto{\pgfqpoint{2.656505in}{1.389182in}}%
\pgfpathlineto{\pgfqpoint{2.625680in}{1.382482in}}%
\pgfpathclose%
\pgfusepath{fill}%
\end{pgfscope}%
\begin{pgfscope}%
\pgfpathrectangle{\pgfqpoint{0.100000in}{0.100000in}}{\pgfqpoint{3.007045in}{1.925000in}}%
\pgfusepath{clip}%
\pgfsetbuttcap%
\pgfsetmiterjoin%
\definecolor{currentfill}{rgb}{0.468743,0.775087,0.474125}%
\pgfsetfillcolor{currentfill}%
\pgfsetlinewidth{0.000000pt}%
\definecolor{currentstroke}{rgb}{0.000000,0.000000,0.000000}%
\pgfsetstrokecolor{currentstroke}%
\pgfsetstrokeopacity{0.000000}%
\pgfsetdash{}{0pt}%
\pgfpathmoveto{\pgfqpoint{0.683738in}{1.619619in}}%
\pgfpathlineto{\pgfqpoint{0.687625in}{1.619603in}}%
\pgfpathlineto{\pgfqpoint{0.695577in}{1.627414in}}%
\pgfpathlineto{\pgfqpoint{0.698781in}{1.628508in}}%
\pgfpathlineto{\pgfqpoint{0.716216in}{1.624121in}}%
\pgfpathlineto{\pgfqpoint{0.720108in}{1.619262in}}%
\pgfpathlineto{\pgfqpoint{0.718698in}{1.613732in}}%
\pgfpathlineto{\pgfqpoint{0.730914in}{1.610662in}}%
\pgfpathlineto{\pgfqpoint{0.733470in}{1.622014in}}%
\pgfpathlineto{\pgfqpoint{0.738581in}{1.626973in}}%
\pgfpathlineto{\pgfqpoint{0.746366in}{1.637288in}}%
\pgfpathlineto{\pgfqpoint{0.754151in}{1.641671in}}%
\pgfpathlineto{\pgfqpoint{0.766659in}{1.638492in}}%
\pgfpathlineto{\pgfqpoint{0.763942in}{1.627183in}}%
\pgfpathlineto{\pgfqpoint{0.767582in}{1.623021in}}%
\pgfpathlineto{\pgfqpoint{0.765709in}{1.615018in}}%
\pgfpathlineto{\pgfqpoint{0.770478in}{1.610681in}}%
\pgfpathlineto{\pgfqpoint{0.775430in}{1.609514in}}%
\pgfpathlineto{\pgfqpoint{0.771340in}{1.602215in}}%
\pgfpathlineto{\pgfqpoint{0.768216in}{1.591309in}}%
\pgfpathlineto{\pgfqpoint{0.763908in}{1.592366in}}%
\pgfpathlineto{\pgfqpoint{0.762552in}{1.586835in}}%
\pgfpathlineto{\pgfqpoint{0.758886in}{1.587771in}}%
\pgfpathlineto{\pgfqpoint{0.755702in}{1.582709in}}%
\pgfpathlineto{\pgfqpoint{0.749229in}{1.584343in}}%
\pgfpathlineto{\pgfqpoint{0.746496in}{1.573138in}}%
\pgfpathlineto{\pgfqpoint{0.739760in}{1.573820in}}%
\pgfpathlineto{\pgfqpoint{0.736446in}{1.580574in}}%
\pgfpathlineto{\pgfqpoint{0.733534in}{1.580107in}}%
\pgfpathlineto{\pgfqpoint{0.729959in}{1.568628in}}%
\pgfpathlineto{\pgfqpoint{0.703346in}{1.460924in}}%
\pgfpathlineto{\pgfqpoint{0.647669in}{1.474848in}}%
\pgfpathlineto{\pgfqpoint{0.651488in}{1.492874in}}%
\pgfpathlineto{\pgfqpoint{0.662044in}{1.533668in}}%
\pgfpathlineto{\pgfqpoint{0.661467in}{1.533821in}}%
\pgfpathlineto{\pgfqpoint{0.675770in}{1.589097in}}%
\pgfpathlineto{\pgfqpoint{0.680340in}{1.605736in}}%
\pgfpathlineto{\pgfqpoint{0.683738in}{1.619619in}}%
\pgfpathclose%
\pgfusepath{fill}%
\end{pgfscope}%
\begin{pgfscope}%
\pgfpathrectangle{\pgfqpoint{0.100000in}{0.100000in}}{\pgfqpoint{3.007045in}{1.925000in}}%
\pgfusepath{clip}%
\pgfsetbuttcap%
\pgfsetmiterjoin%
\definecolor{currentfill}{rgb}{0.629296,0.850012,0.606136}%
\pgfsetfillcolor{currentfill}%
\pgfsetlinewidth{0.000000pt}%
\definecolor{currentstroke}{rgb}{0.000000,0.000000,0.000000}%
\pgfsetstrokecolor{currentstroke}%
\pgfsetstrokeopacity{0.000000}%
\pgfsetdash{}{0pt}%
\pgfpathmoveto{\pgfqpoint{2.402438in}{0.696626in}}%
\pgfpathlineto{\pgfqpoint{2.395249in}{0.698742in}}%
\pgfpathlineto{\pgfqpoint{2.394856in}{0.702917in}}%
\pgfpathlineto{\pgfqpoint{2.389600in}{0.714570in}}%
\pgfpathlineto{\pgfqpoint{2.386745in}{0.718055in}}%
\pgfpathlineto{\pgfqpoint{2.389930in}{0.724977in}}%
\pgfpathlineto{\pgfqpoint{2.395922in}{0.728401in}}%
\pgfpathlineto{\pgfqpoint{2.409404in}{0.729924in}}%
\pgfpathlineto{\pgfqpoint{2.410702in}{0.718714in}}%
\pgfpathlineto{\pgfqpoint{2.425917in}{0.719822in}}%
\pgfpathlineto{\pgfqpoint{2.428514in}{0.718344in}}%
\pgfpathlineto{\pgfqpoint{2.429949in}{0.711648in}}%
\pgfpathlineto{\pgfqpoint{2.433519in}{0.709598in}}%
\pgfpathlineto{\pgfqpoint{2.437033in}{0.704008in}}%
\pgfpathlineto{\pgfqpoint{2.420123in}{0.701779in}}%
\pgfpathlineto{\pgfqpoint{2.412679in}{0.701290in}}%
\pgfpathlineto{\pgfqpoint{2.413072in}{0.697996in}}%
\pgfpathlineto{\pgfqpoint{2.402438in}{0.696626in}}%
\pgfpathclose%
\pgfusepath{fill}%
\end{pgfscope}%
\begin{pgfscope}%
\pgfpathrectangle{\pgfqpoint{0.100000in}{0.100000in}}{\pgfqpoint{3.007045in}{1.925000in}}%
\pgfusepath{clip}%
\pgfsetbuttcap%
\pgfsetmiterjoin%
\definecolor{currentfill}{rgb}{0.451765,0.767090,0.461207}%
\pgfsetfillcolor{currentfill}%
\pgfsetlinewidth{0.000000pt}%
\definecolor{currentstroke}{rgb}{0.000000,0.000000,0.000000}%
\pgfsetstrokecolor{currentstroke}%
\pgfsetstrokeopacity{0.000000}%
\pgfsetdash{}{0pt}%
\pgfpathmoveto{\pgfqpoint{1.511195in}{1.206454in}}%
\pgfpathlineto{\pgfqpoint{1.510401in}{1.206498in}}%
\pgfpathlineto{\pgfqpoint{1.511500in}{1.229428in}}%
\pgfpathlineto{\pgfqpoint{1.482455in}{1.231145in}}%
\pgfpathlineto{\pgfqpoint{1.483897in}{1.254208in}}%
\pgfpathlineto{\pgfqpoint{1.511547in}{1.252543in}}%
\pgfpathlineto{\pgfqpoint{1.512812in}{1.275327in}}%
\pgfpathlineto{\pgfqpoint{1.552431in}{1.273213in}}%
\pgfpathlineto{\pgfqpoint{1.551658in}{1.248467in}}%
\pgfpathlineto{\pgfqpoint{1.563506in}{1.247230in}}%
\pgfpathlineto{\pgfqpoint{1.561531in}{1.203886in}}%
\pgfpathlineto{\pgfqpoint{1.539001in}{1.204970in}}%
\pgfpathlineto{\pgfqpoint{1.511195in}{1.206454in}}%
\pgfpathclose%
\pgfusepath{fill}%
\end{pgfscope}%
\begin{pgfscope}%
\pgfpathrectangle{\pgfqpoint{0.100000in}{0.100000in}}{\pgfqpoint{3.007045in}{1.925000in}}%
\pgfusepath{clip}%
\pgfsetbuttcap%
\pgfsetmiterjoin%
\definecolor{currentfill}{rgb}{0.551788,0.813841,0.542407}%
\pgfsetfillcolor{currentfill}%
\pgfsetlinewidth{0.000000pt}%
\definecolor{currentstroke}{rgb}{0.000000,0.000000,0.000000}%
\pgfsetstrokecolor{currentstroke}%
\pgfsetstrokeopacity{0.000000}%
\pgfsetdash{}{0pt}%
\pgfpathmoveto{\pgfqpoint{1.298841in}{1.455523in}}%
\pgfpathlineto{\pgfqpoint{1.339144in}{1.451533in}}%
\pgfpathlineto{\pgfqpoint{1.335923in}{1.418503in}}%
\pgfpathlineto{\pgfqpoint{1.333387in}{1.392952in}}%
\pgfpathlineto{\pgfqpoint{1.292810in}{1.397021in}}%
\pgfpathlineto{\pgfqpoint{1.293878in}{1.409890in}}%
\pgfpathlineto{\pgfqpoint{1.298841in}{1.455523in}}%
\pgfpathclose%
\pgfusepath{fill}%
\end{pgfscope}%
\begin{pgfscope}%
\pgfpathrectangle{\pgfqpoint{0.100000in}{0.100000in}}{\pgfqpoint{3.007045in}{1.925000in}}%
\pgfusepath{clip}%
\pgfsetbuttcap%
\pgfsetmiterjoin%
\definecolor{currentfill}{rgb}{0.662345,0.864022,0.638001}%
\pgfsetfillcolor{currentfill}%
\pgfsetlinewidth{0.000000pt}%
\definecolor{currentstroke}{rgb}{0.000000,0.000000,0.000000}%
\pgfsetstrokecolor{currentstroke}%
\pgfsetstrokeopacity{0.000000}%
\pgfsetdash{}{0pt}%
\pgfpathmoveto{\pgfqpoint{0.374195in}{1.289386in}}%
\pgfpathlineto{\pgfqpoint{0.377156in}{1.290848in}}%
\pgfpathlineto{\pgfqpoint{0.377641in}{1.299051in}}%
\pgfpathlineto{\pgfqpoint{0.376145in}{1.309067in}}%
\pgfpathlineto{\pgfqpoint{0.371022in}{1.321407in}}%
\pgfpathlineto{\pgfqpoint{0.373315in}{1.325445in}}%
\pgfpathlineto{\pgfqpoint{0.382079in}{1.325306in}}%
\pgfpathlineto{\pgfqpoint{0.388700in}{1.330958in}}%
\pgfpathlineto{\pgfqpoint{0.388567in}{1.334258in}}%
\pgfpathlineto{\pgfqpoint{0.393337in}{1.331101in}}%
\pgfpathlineto{\pgfqpoint{0.393792in}{1.321038in}}%
\pgfpathlineto{\pgfqpoint{0.395623in}{1.317666in}}%
\pgfpathlineto{\pgfqpoint{0.396070in}{1.307741in}}%
\pgfpathlineto{\pgfqpoint{0.400092in}{1.304827in}}%
\pgfpathlineto{\pgfqpoint{0.404388in}{1.306594in}}%
\pgfpathlineto{\pgfqpoint{0.416471in}{1.302466in}}%
\pgfpathlineto{\pgfqpoint{0.412425in}{1.289011in}}%
\pgfpathlineto{\pgfqpoint{0.417431in}{1.287454in}}%
\pgfpathlineto{\pgfqpoint{0.414282in}{1.280552in}}%
\pgfpathlineto{\pgfqpoint{0.409917in}{1.279165in}}%
\pgfpathlineto{\pgfqpoint{0.409876in}{1.274771in}}%
\pgfpathlineto{\pgfqpoint{0.413844in}{1.273461in}}%
\pgfpathlineto{\pgfqpoint{0.409743in}{1.259767in}}%
\pgfpathlineto{\pgfqpoint{0.409870in}{1.255826in}}%
\pgfpathlineto{\pgfqpoint{0.404604in}{1.238315in}}%
\pgfpathlineto{\pgfqpoint{0.407690in}{1.233195in}}%
\pgfpathlineto{\pgfqpoint{0.407661in}{1.233230in}}%
\pgfpathlineto{\pgfqpoint{0.388058in}{1.239323in}}%
\pgfpathlineto{\pgfqpoint{0.384479in}{1.238341in}}%
\pgfpathlineto{\pgfqpoint{0.377116in}{1.242201in}}%
\pgfpathlineto{\pgfqpoint{0.370754in}{1.240933in}}%
\pgfpathlineto{\pgfqpoint{0.368846in}{1.236860in}}%
\pgfpathlineto{\pgfqpoint{0.368463in}{1.226772in}}%
\pgfpathlineto{\pgfqpoint{0.359639in}{1.227533in}}%
\pgfpathlineto{\pgfqpoint{0.359301in}{1.222111in}}%
\pgfpathlineto{\pgfqpoint{0.355662in}{1.228408in}}%
\pgfpathlineto{\pgfqpoint{0.355631in}{1.232502in}}%
\pgfpathlineto{\pgfqpoint{0.358662in}{1.239132in}}%
\pgfpathlineto{\pgfqpoint{0.358656in}{1.248068in}}%
\pgfpathlineto{\pgfqpoint{0.356187in}{1.251923in}}%
\pgfpathlineto{\pgfqpoint{0.359786in}{1.258991in}}%
\pgfpathlineto{\pgfqpoint{0.361404in}{1.267963in}}%
\pgfpathlineto{\pgfqpoint{0.367391in}{1.268171in}}%
\pgfpathlineto{\pgfqpoint{0.365394in}{1.253654in}}%
\pgfpathlineto{\pgfqpoint{0.369838in}{1.250965in}}%
\pgfpathlineto{\pgfqpoint{0.373837in}{1.244770in}}%
\pgfpathlineto{\pgfqpoint{0.376610in}{1.245915in}}%
\pgfpathlineto{\pgfqpoint{0.376740in}{1.254309in}}%
\pgfpathlineto{\pgfqpoint{0.370828in}{1.266174in}}%
\pgfpathlineto{\pgfqpoint{0.373731in}{1.272170in}}%
\pgfpathlineto{\pgfqpoint{0.370644in}{1.277375in}}%
\pgfpathlineto{\pgfqpoint{0.379239in}{1.281168in}}%
\pgfpathlineto{\pgfqpoint{0.378012in}{1.285740in}}%
\pgfpathlineto{\pgfqpoint{0.374195in}{1.289386in}}%
\pgfpathclose%
\pgfusepath{fill}%
\end{pgfscope}%
\begin{pgfscope}%
\pgfpathrectangle{\pgfqpoint{0.100000in}{0.100000in}}{\pgfqpoint{3.007045in}{1.925000in}}%
\pgfusepath{clip}%
\pgfsetbuttcap%
\pgfsetmiterjoin%
\definecolor{currentfill}{rgb}{0.662345,0.864022,0.638001}%
\pgfsetfillcolor{currentfill}%
\pgfsetlinewidth{0.000000pt}%
\definecolor{currentstroke}{rgb}{0.000000,0.000000,0.000000}%
\pgfsetstrokecolor{currentstroke}%
\pgfsetstrokeopacity{0.000000}%
\pgfsetdash{}{0pt}%
\pgfpathmoveto{\pgfqpoint{0.369150in}{1.288189in}}%
\pgfpathlineto{\pgfqpoint{0.366091in}{1.274818in}}%
\pgfpathlineto{\pgfqpoint{0.361852in}{1.270865in}}%
\pgfpathlineto{\pgfqpoint{0.357053in}{1.277694in}}%
\pgfpathlineto{\pgfqpoint{0.354298in}{1.277874in}}%
\pgfpathlineto{\pgfqpoint{0.349286in}{1.287909in}}%
\pgfpathlineto{\pgfqpoint{0.344216in}{1.290431in}}%
\pgfpathlineto{\pgfqpoint{0.347199in}{1.297896in}}%
\pgfpathlineto{\pgfqpoint{0.347482in}{1.307924in}}%
\pgfpathlineto{\pgfqpoint{0.352980in}{1.307630in}}%
\pgfpathlineto{\pgfqpoint{0.358679in}{1.298249in}}%
\pgfpathlineto{\pgfqpoint{0.366648in}{1.294372in}}%
\pgfpathlineto{\pgfqpoint{0.369150in}{1.288189in}}%
\pgfpathclose%
\pgfusepath{fill}%
\end{pgfscope}%
\begin{pgfscope}%
\pgfpathrectangle{\pgfqpoint{0.100000in}{0.100000in}}{\pgfqpoint{3.007045in}{1.925000in}}%
\pgfusepath{clip}%
\pgfsetbuttcap%
\pgfsetmiterjoin%
\definecolor{currentfill}{rgb}{0.579469,0.826759,0.565167}%
\pgfsetfillcolor{currentfill}%
\pgfsetlinewidth{0.000000pt}%
\definecolor{currentstroke}{rgb}{0.000000,0.000000,0.000000}%
\pgfsetstrokecolor{currentstroke}%
\pgfsetstrokeopacity{0.000000}%
\pgfsetdash{}{0pt}%
\pgfpathmoveto{\pgfqpoint{2.025542in}{1.310434in}}%
\pgfpathlineto{\pgfqpoint{2.009626in}{1.309620in}}%
\pgfpathlineto{\pgfqpoint{2.006146in}{1.322530in}}%
\pgfpathlineto{\pgfqpoint{2.009028in}{1.324274in}}%
\pgfpathlineto{\pgfqpoint{2.010393in}{1.337303in}}%
\pgfpathlineto{\pgfqpoint{2.008917in}{1.344598in}}%
\pgfpathlineto{\pgfqpoint{2.001078in}{1.349385in}}%
\pgfpathlineto{\pgfqpoint{2.020468in}{1.350608in}}%
\pgfpathlineto{\pgfqpoint{2.018995in}{1.370976in}}%
\pgfpathlineto{\pgfqpoint{2.046233in}{1.372218in}}%
\pgfpathlineto{\pgfqpoint{2.078237in}{1.373879in}}%
\pgfpathlineto{\pgfqpoint{2.079874in}{1.351505in}}%
\pgfpathlineto{\pgfqpoint{2.085588in}{1.351908in}}%
\pgfpathlineto{\pgfqpoint{2.085344in}{1.346129in}}%
\pgfpathlineto{\pgfqpoint{2.087331in}{1.317506in}}%
\pgfpathlineto{\pgfqpoint{2.059573in}{1.315459in}}%
\pgfpathlineto{\pgfqpoint{2.059792in}{1.312587in}}%
\pgfpathlineto{\pgfqpoint{2.025542in}{1.310434in}}%
\pgfpathclose%
\pgfusepath{fill}%
\end{pgfscope}%
\begin{pgfscope}%
\pgfpathrectangle{\pgfqpoint{0.100000in}{0.100000in}}{\pgfqpoint{3.007045in}{1.925000in}}%
\pgfusepath{clip}%
\pgfsetbuttcap%
\pgfsetmiterjoin%
\definecolor{currentfill}{rgb}{0.546251,0.811257,0.537855}%
\pgfsetfillcolor{currentfill}%
\pgfsetlinewidth{0.000000pt}%
\definecolor{currentstroke}{rgb}{0.000000,0.000000,0.000000}%
\pgfsetstrokecolor{currentstroke}%
\pgfsetstrokeopacity{0.000000}%
\pgfsetdash{}{0pt}%
\pgfpathmoveto{\pgfqpoint{2.169528in}{1.143936in}}%
\pgfpathlineto{\pgfqpoint{2.175194in}{1.144388in}}%
\pgfpathlineto{\pgfqpoint{2.175011in}{1.140154in}}%
\pgfpathlineto{\pgfqpoint{2.170118in}{1.138155in}}%
\pgfpathlineto{\pgfqpoint{2.168758in}{1.129637in}}%
\pgfpathlineto{\pgfqpoint{2.171500in}{1.120112in}}%
\pgfpathlineto{\pgfqpoint{2.160227in}{1.118407in}}%
\pgfpathlineto{\pgfqpoint{2.154476in}{1.114634in}}%
\pgfpathlineto{\pgfqpoint{2.150356in}{1.119886in}}%
\pgfpathlineto{\pgfqpoint{2.136969in}{1.118978in}}%
\pgfpathlineto{\pgfqpoint{2.135746in}{1.137366in}}%
\pgfpathlineto{\pgfqpoint{2.154997in}{1.139057in}}%
\pgfpathlineto{\pgfqpoint{2.155147in}{1.142366in}}%
\pgfpathlineto{\pgfqpoint{2.169528in}{1.143936in}}%
\pgfpathclose%
\pgfusepath{fill}%
\end{pgfscope}%
\begin{pgfscope}%
\pgfpathrectangle{\pgfqpoint{0.100000in}{0.100000in}}{\pgfqpoint{3.007045in}{1.925000in}}%
\pgfusepath{clip}%
\pgfsetbuttcap%
\pgfsetmiterjoin%
\definecolor{currentfill}{rgb}{0.573933,0.824175,0.560615}%
\pgfsetfillcolor{currentfill}%
\pgfsetlinewidth{0.000000pt}%
\definecolor{currentstroke}{rgb}{0.000000,0.000000,0.000000}%
\pgfsetstrokecolor{currentstroke}%
\pgfsetstrokeopacity{0.000000}%
\pgfsetdash{}{0pt}%
\pgfpathmoveto{\pgfqpoint{1.735784in}{0.842654in}}%
\pgfpathlineto{\pgfqpoint{1.735756in}{0.836902in}}%
\pgfpathlineto{\pgfqpoint{1.729891in}{0.836933in}}%
\pgfpathlineto{\pgfqpoint{1.729695in}{0.813812in}}%
\pgfpathlineto{\pgfqpoint{1.718089in}{0.813933in}}%
\pgfpathlineto{\pgfqpoint{1.695332in}{0.814267in}}%
\pgfpathlineto{\pgfqpoint{1.695616in}{0.831530in}}%
\pgfpathlineto{\pgfqpoint{1.689815in}{0.831695in}}%
\pgfpathlineto{\pgfqpoint{1.689906in}{0.837361in}}%
\pgfpathlineto{\pgfqpoint{1.672844in}{0.837714in}}%
\pgfpathlineto{\pgfqpoint{1.672960in}{0.843473in}}%
\pgfpathlineto{\pgfqpoint{1.667244in}{0.843587in}}%
\pgfpathlineto{\pgfqpoint{1.667600in}{0.860848in}}%
\pgfpathlineto{\pgfqpoint{1.670026in}{0.865689in}}%
\pgfpathlineto{\pgfqpoint{1.674047in}{0.866744in}}%
\pgfpathlineto{\pgfqpoint{1.678967in}{0.860772in}}%
\pgfpathlineto{\pgfqpoint{1.683133in}{0.863693in}}%
\pgfpathlineto{\pgfqpoint{1.691602in}{0.864102in}}%
\pgfpathlineto{\pgfqpoint{1.691843in}{0.877659in}}%
\pgfpathlineto{\pgfqpoint{1.694493in}{0.877607in}}%
\pgfpathlineto{\pgfqpoint{1.694674in}{0.889113in}}%
\pgfpathlineto{\pgfqpoint{1.719456in}{0.888768in}}%
\pgfpathlineto{\pgfqpoint{1.719231in}{0.879638in}}%
\pgfpathlineto{\pgfqpoint{1.713476in}{0.872989in}}%
\pgfpathlineto{\pgfqpoint{1.713003in}{0.848605in}}%
\pgfpathlineto{\pgfqpoint{1.724413in}{0.848474in}}%
\pgfpathlineto{\pgfqpoint{1.724353in}{0.842730in}}%
\pgfpathlineto{\pgfqpoint{1.735784in}{0.842654in}}%
\pgfpathclose%
\pgfusepath{fill}%
\end{pgfscope}%
\begin{pgfscope}%
\pgfpathrectangle{\pgfqpoint{0.100000in}{0.100000in}}{\pgfqpoint{3.007045in}{1.925000in}}%
\pgfusepath{clip}%
\pgfsetbuttcap%
\pgfsetmiterjoin%
\definecolor{currentfill}{rgb}{0.551788,0.813841,0.542407}%
\pgfsetfillcolor{currentfill}%
\pgfsetlinewidth{0.000000pt}%
\definecolor{currentstroke}{rgb}{0.000000,0.000000,0.000000}%
\pgfsetstrokecolor{currentstroke}%
\pgfsetstrokeopacity{0.000000}%
\pgfsetdash{}{0pt}%
\pgfpathmoveto{\pgfqpoint{2.643747in}{1.050644in}}%
\pgfpathlineto{\pgfqpoint{2.672729in}{1.056217in}}%
\pgfpathlineto{\pgfqpoint{2.674093in}{1.045577in}}%
\pgfpathlineto{\pgfqpoint{2.671365in}{1.033005in}}%
\pgfpathlineto{\pgfqpoint{2.664233in}{1.034590in}}%
\pgfpathlineto{\pgfqpoint{2.654897in}{1.034173in}}%
\pgfpathlineto{\pgfqpoint{2.650549in}{1.026900in}}%
\pgfpathlineto{\pgfqpoint{2.645335in}{1.026823in}}%
\pgfpathlineto{\pgfqpoint{2.641967in}{1.044405in}}%
\pgfpathlineto{\pgfqpoint{2.643747in}{1.050644in}}%
\pgfpathclose%
\pgfusepath{fill}%
\end{pgfscope}%
\begin{pgfscope}%
\pgfpathrectangle{\pgfqpoint{0.100000in}{0.100000in}}{\pgfqpoint{3.007045in}{1.925000in}}%
\pgfusepath{clip}%
\pgfsetbuttcap%
\pgfsetmiterjoin%
\definecolor{currentfill}{rgb}{0.351373,0.717878,0.411995}%
\pgfsetfillcolor{currentfill}%
\pgfsetlinewidth{0.000000pt}%
\definecolor{currentstroke}{rgb}{0.000000,0.000000,0.000000}%
\pgfsetstrokecolor{currentstroke}%
\pgfsetstrokeopacity{0.000000}%
\pgfsetdash{}{0pt}%
\pgfpathmoveto{\pgfqpoint{1.003150in}{0.848168in}}%
\pgfpathlineto{\pgfqpoint{0.997334in}{0.810949in}}%
\pgfpathlineto{\pgfqpoint{0.989331in}{0.759721in}}%
\pgfpathlineto{\pgfqpoint{0.927312in}{0.769778in}}%
\pgfpathlineto{\pgfqpoint{0.911810in}{0.772480in}}%
\pgfpathlineto{\pgfqpoint{0.912796in}{0.778126in}}%
\pgfpathlineto{\pgfqpoint{0.920548in}{0.822552in}}%
\pgfpathlineto{\pgfqpoint{0.923055in}{0.834676in}}%
\pgfpathlineto{\pgfqpoint{0.930194in}{0.836921in}}%
\pgfpathlineto{\pgfqpoint{0.938796in}{0.837657in}}%
\pgfpathlineto{\pgfqpoint{0.947992in}{0.836130in}}%
\pgfpathlineto{\pgfqpoint{0.949199in}{0.843379in}}%
\pgfpathlineto{\pgfqpoint{0.954625in}{0.842378in}}%
\pgfpathlineto{\pgfqpoint{0.960834in}{0.834656in}}%
\pgfpathlineto{\pgfqpoint{0.971108in}{0.842140in}}%
\pgfpathlineto{\pgfqpoint{0.975090in}{0.841991in}}%
\pgfpathlineto{\pgfqpoint{0.978685in}{0.845988in}}%
\pgfpathlineto{\pgfqpoint{0.986809in}{0.850694in}}%
\pgfpathlineto{\pgfqpoint{1.003150in}{0.848168in}}%
\pgfpathclose%
\pgfusepath{fill}%
\end{pgfscope}%
\begin{pgfscope}%
\pgfpathrectangle{\pgfqpoint{0.100000in}{0.100000in}}{\pgfqpoint{3.007045in}{1.925000in}}%
\pgfusepath{clip}%
\pgfsetbuttcap%
\pgfsetmiterjoin%
\definecolor{currentfill}{rgb}{0.513033,0.795755,0.510542}%
\pgfsetfillcolor{currentfill}%
\pgfsetlinewidth{0.000000pt}%
\definecolor{currentstroke}{rgb}{0.000000,0.000000,0.000000}%
\pgfsetstrokecolor{currentstroke}%
\pgfsetstrokeopacity{0.000000}%
\pgfsetdash{}{0pt}%
\pgfpathmoveto{\pgfqpoint{1.730423in}{0.638033in}}%
\pgfpathlineto{\pgfqpoint{1.734214in}{0.624182in}}%
\pgfpathlineto{\pgfqpoint{1.731113in}{0.622526in}}%
\pgfpathlineto{\pgfqpoint{1.728210in}{0.613654in}}%
\pgfpathlineto{\pgfqpoint{1.728806in}{0.610818in}}%
\pgfpathlineto{\pgfqpoint{1.716755in}{0.610772in}}%
\pgfpathlineto{\pgfqpoint{1.701562in}{0.603092in}}%
\pgfpathlineto{\pgfqpoint{1.696915in}{0.623980in}}%
\pgfpathlineto{\pgfqpoint{1.697477in}{0.628654in}}%
\pgfpathlineto{\pgfqpoint{1.681563in}{0.619854in}}%
\pgfpathlineto{\pgfqpoint{1.663193in}{0.652637in}}%
\pgfpathlineto{\pgfqpoint{1.675361in}{0.659431in}}%
\pgfpathlineto{\pgfqpoint{1.665733in}{0.676821in}}%
\pgfpathlineto{\pgfqpoint{1.694814in}{0.693189in}}%
\pgfpathlineto{\pgfqpoint{1.703043in}{0.687922in}}%
\pgfpathlineto{\pgfqpoint{1.712375in}{0.673683in}}%
\pgfpathlineto{\pgfqpoint{1.718592in}{0.656205in}}%
\pgfpathlineto{\pgfqpoint{1.722872in}{0.654828in}}%
\pgfpathlineto{\pgfqpoint{1.723459in}{0.650585in}}%
\pgfpathlineto{\pgfqpoint{1.732103in}{0.644712in}}%
\pgfpathlineto{\pgfqpoint{1.730423in}{0.638033in}}%
\pgfpathclose%
\pgfusepath{fill}%
\end{pgfscope}%
\begin{pgfscope}%
\pgfpathrectangle{\pgfqpoint{0.100000in}{0.100000in}}{\pgfqpoint{3.007045in}{1.925000in}}%
\pgfusepath{clip}%
\pgfsetbuttcap%
\pgfsetmiterjoin%
\definecolor{currentfill}{rgb}{0.590542,0.831926,0.574271}%
\pgfsetfillcolor{currentfill}%
\pgfsetlinewidth{0.000000pt}%
\definecolor{currentstroke}{rgb}{0.000000,0.000000,0.000000}%
\pgfsetstrokecolor{currentstroke}%
\pgfsetstrokeopacity{0.000000}%
\pgfsetdash{}{0pt}%
\pgfpathmoveto{\pgfqpoint{1.643166in}{1.743370in}}%
\pgfpathlineto{\pgfqpoint{1.643745in}{1.731622in}}%
\pgfpathlineto{\pgfqpoint{1.642987in}{1.708503in}}%
\pgfpathlineto{\pgfqpoint{1.638404in}{1.708645in}}%
\pgfpathlineto{\pgfqpoint{1.615296in}{1.709477in}}%
\pgfpathlineto{\pgfqpoint{1.614607in}{1.721113in}}%
\pgfpathlineto{\pgfqpoint{1.580095in}{1.722634in}}%
\pgfpathlineto{\pgfqpoint{1.557103in}{1.723835in}}%
\pgfpathlineto{\pgfqpoint{1.556105in}{1.735501in}}%
\pgfpathlineto{\pgfqpoint{1.557416in}{1.758816in}}%
\pgfpathlineto{\pgfqpoint{1.573086in}{1.757987in}}%
\pgfpathlineto{\pgfqpoint{1.596215in}{1.756877in}}%
\pgfpathlineto{\pgfqpoint{1.596746in}{1.768470in}}%
\pgfpathlineto{\pgfqpoint{1.625579in}{1.767286in}}%
\pgfpathlineto{\pgfqpoint{1.626307in}{1.755616in}}%
\pgfpathlineto{\pgfqpoint{1.625869in}{1.743967in}}%
\pgfpathlineto{\pgfqpoint{1.643166in}{1.743370in}}%
\pgfpathclose%
\pgfusepath{fill}%
\end{pgfscope}%
\begin{pgfscope}%
\pgfpathrectangle{\pgfqpoint{0.100000in}{0.100000in}}{\pgfqpoint{3.007045in}{1.925000in}}%
\pgfusepath{clip}%
\pgfsetbuttcap%
\pgfsetmiterjoin%
\definecolor{currentfill}{rgb}{0.671696,0.867958,0.647105}%
\pgfsetfillcolor{currentfill}%
\pgfsetlinewidth{0.000000pt}%
\definecolor{currentstroke}{rgb}{0.000000,0.000000,0.000000}%
\pgfsetstrokecolor{currentstroke}%
\pgfsetstrokeopacity{0.000000}%
\pgfsetdash{}{0pt}%
\pgfpathmoveto{\pgfqpoint{2.060711in}{0.508116in}}%
\pgfpathlineto{\pgfqpoint{2.060425in}{0.500127in}}%
\pgfpathlineto{\pgfqpoint{2.056411in}{0.498585in}}%
\pgfpathlineto{\pgfqpoint{2.060074in}{0.492771in}}%
\pgfpathlineto{\pgfqpoint{2.058740in}{0.489589in}}%
\pgfpathlineto{\pgfqpoint{2.051172in}{0.484014in}}%
\pgfpathlineto{\pgfqpoint{2.048181in}{0.488038in}}%
\pgfpathlineto{\pgfqpoint{2.046171in}{0.495142in}}%
\pgfpathlineto{\pgfqpoint{2.042744in}{0.497419in}}%
\pgfpathlineto{\pgfqpoint{2.040294in}{0.492004in}}%
\pgfpathlineto{\pgfqpoint{2.034473in}{0.496696in}}%
\pgfpathlineto{\pgfqpoint{2.026621in}{0.490352in}}%
\pgfpathlineto{\pgfqpoint{2.026344in}{0.485680in}}%
\pgfpathlineto{\pgfqpoint{2.021444in}{0.485001in}}%
\pgfpathlineto{\pgfqpoint{2.014806in}{0.481223in}}%
\pgfpathlineto{\pgfqpoint{2.007886in}{0.488255in}}%
\pgfpathlineto{\pgfqpoint{2.000041in}{0.488340in}}%
\pgfpathlineto{\pgfqpoint{1.989313in}{0.491975in}}%
\pgfpathlineto{\pgfqpoint{1.985959in}{0.494974in}}%
\pgfpathlineto{\pgfqpoint{1.991281in}{0.500512in}}%
\pgfpathlineto{\pgfqpoint{1.989147in}{0.507119in}}%
\pgfpathlineto{\pgfqpoint{1.973270in}{0.512709in}}%
\pgfpathlineto{\pgfqpoint{1.968012in}{0.516450in}}%
\pgfpathlineto{\pgfqpoint{1.968031in}{0.523710in}}%
\pgfpathlineto{\pgfqpoint{1.958964in}{0.523855in}}%
\pgfpathlineto{\pgfqpoint{1.953040in}{0.522060in}}%
\pgfpathlineto{\pgfqpoint{1.962543in}{0.530961in}}%
\pgfpathlineto{\pgfqpoint{1.967294in}{0.537756in}}%
\pgfpathlineto{\pgfqpoint{1.976086in}{0.538805in}}%
\pgfpathlineto{\pgfqpoint{1.982084in}{0.530926in}}%
\pgfpathlineto{\pgfqpoint{1.981476in}{0.527516in}}%
\pgfpathlineto{\pgfqpoint{1.993498in}{0.524333in}}%
\pgfpathlineto{\pgfqpoint{1.998738in}{0.526353in}}%
\pgfpathlineto{\pgfqpoint{1.997286in}{0.531031in}}%
\pgfpathlineto{\pgfqpoint{1.992985in}{0.532196in}}%
\pgfpathlineto{\pgfqpoint{1.988235in}{0.542241in}}%
\pgfpathlineto{\pgfqpoint{1.990185in}{0.543999in}}%
\pgfpathlineto{\pgfqpoint{2.002634in}{0.548120in}}%
\pgfpathlineto{\pgfqpoint{2.005096in}{0.547382in}}%
\pgfpathlineto{\pgfqpoint{2.005579in}{0.541500in}}%
\pgfpathlineto{\pgfqpoint{2.010082in}{0.536853in}}%
\pgfpathlineto{\pgfqpoint{2.030059in}{0.536911in}}%
\pgfpathlineto{\pgfqpoint{2.034075in}{0.531818in}}%
\pgfpathlineto{\pgfqpoint{2.040150in}{0.528704in}}%
\pgfpathlineto{\pgfqpoint{2.042354in}{0.524234in}}%
\pgfpathlineto{\pgfqpoint{2.050269in}{0.524553in}}%
\pgfpathlineto{\pgfqpoint{2.053131in}{0.522603in}}%
\pgfpathlineto{\pgfqpoint{2.051671in}{0.516053in}}%
\pgfpathlineto{\pgfqpoint{2.054795in}{0.510464in}}%
\pgfpathlineto{\pgfqpoint{2.060711in}{0.508116in}}%
\pgfpathclose%
\pgfusepath{fill}%
\end{pgfscope}%
\begin{pgfscope}%
\pgfpathrectangle{\pgfqpoint{0.100000in}{0.100000in}}{\pgfqpoint{3.007045in}{1.925000in}}%
\pgfusepath{clip}%
\pgfsetbuttcap%
\pgfsetmiterjoin%
\definecolor{currentfill}{rgb}{0.723122,0.889612,0.697178}%
\pgfsetfillcolor{currentfill}%
\pgfsetlinewidth{0.000000pt}%
\definecolor{currentstroke}{rgb}{0.000000,0.000000,0.000000}%
\pgfsetstrokecolor{currentstroke}%
\pgfsetstrokeopacity{0.000000}%
\pgfsetdash{}{0pt}%
\pgfpathmoveto{\pgfqpoint{1.502115in}{0.878788in}}%
\pgfpathlineto{\pgfqpoint{1.504252in}{0.917589in}}%
\pgfpathlineto{\pgfqpoint{1.505202in}{0.934838in}}%
\pgfpathlineto{\pgfqpoint{1.510237in}{0.939913in}}%
\pgfpathlineto{\pgfqpoint{1.515519in}{0.937448in}}%
\pgfpathlineto{\pgfqpoint{1.520249in}{0.932634in}}%
\pgfpathlineto{\pgfqpoint{1.528739in}{0.932611in}}%
\pgfpathlineto{\pgfqpoint{1.532268in}{0.939536in}}%
\pgfpathlineto{\pgfqpoint{1.538682in}{0.941826in}}%
\pgfpathlineto{\pgfqpoint{1.538310in}{0.928586in}}%
\pgfpathlineto{\pgfqpoint{1.578070in}{0.926846in}}%
\pgfpathlineto{\pgfqpoint{1.577863in}{0.909634in}}%
\pgfpathlineto{\pgfqpoint{1.576781in}{0.879670in}}%
\pgfpathlineto{\pgfqpoint{1.569929in}{0.881210in}}%
\pgfpathlineto{\pgfqpoint{1.514160in}{0.883851in}}%
\pgfpathlineto{\pgfqpoint{1.513852in}{0.878138in}}%
\pgfpathlineto{\pgfqpoint{1.502115in}{0.878788in}}%
\pgfpathclose%
\pgfusepath{fill}%
\end{pgfscope}%
\begin{pgfscope}%
\pgfpathrectangle{\pgfqpoint{0.100000in}{0.100000in}}{\pgfqpoint{3.007045in}{1.925000in}}%
\pgfusepath{clip}%
\pgfsetbuttcap%
\pgfsetmiterjoin%
\definecolor{currentfill}{rgb}{0.474279,0.777670,0.478677}%
\pgfsetfillcolor{currentfill}%
\pgfsetlinewidth{0.000000pt}%
\definecolor{currentstroke}{rgb}{0.000000,0.000000,0.000000}%
\pgfsetstrokecolor{currentstroke}%
\pgfsetstrokeopacity{0.000000}%
\pgfsetdash{}{0pt}%
\pgfpathmoveto{\pgfqpoint{2.523965in}{1.201691in}}%
\pgfpathlineto{\pgfqpoint{2.534247in}{1.194214in}}%
\pgfpathlineto{\pgfqpoint{2.532223in}{1.188646in}}%
\pgfpathlineto{\pgfqpoint{2.532999in}{1.178683in}}%
\pgfpathlineto{\pgfqpoint{2.527409in}{1.177624in}}%
\pgfpathlineto{\pgfqpoint{2.525737in}{1.174087in}}%
\pgfpathlineto{\pgfqpoint{2.517598in}{1.175060in}}%
\pgfpathlineto{\pgfqpoint{2.514455in}{1.181489in}}%
\pgfpathlineto{\pgfqpoint{2.507565in}{1.184564in}}%
\pgfpathlineto{\pgfqpoint{2.500337in}{1.190919in}}%
\pgfpathlineto{\pgfqpoint{2.498857in}{1.196503in}}%
\pgfpathlineto{\pgfqpoint{2.505727in}{1.202402in}}%
\pgfpathlineto{\pgfqpoint{2.520381in}{1.200511in}}%
\pgfpathlineto{\pgfqpoint{2.523965in}{1.201691in}}%
\pgfpathclose%
\pgfusepath{fill}%
\end{pgfscope}%
\begin{pgfscope}%
\pgfpathrectangle{\pgfqpoint{0.100000in}{0.100000in}}{\pgfqpoint{3.007045in}{1.925000in}}%
\pgfusepath{clip}%
\pgfsetbuttcap%
\pgfsetmiterjoin%
\definecolor{currentfill}{rgb}{0.718447,0.887643,0.692626}%
\pgfsetfillcolor{currentfill}%
\pgfsetlinewidth{0.000000pt}%
\definecolor{currentstroke}{rgb}{0.000000,0.000000,0.000000}%
\pgfsetstrokecolor{currentstroke}%
\pgfsetstrokeopacity{0.000000}%
\pgfsetdash{}{0pt}%
\pgfpathmoveto{\pgfqpoint{2.362396in}{1.092621in}}%
\pgfpathlineto{\pgfqpoint{2.358589in}{1.084899in}}%
\pgfpathlineto{\pgfqpoint{2.350883in}{1.086320in}}%
\pgfpathlineto{\pgfqpoint{2.342975in}{1.090758in}}%
\pgfpathlineto{\pgfqpoint{2.328598in}{1.094602in}}%
\pgfpathlineto{\pgfqpoint{2.316608in}{1.083386in}}%
\pgfpathlineto{\pgfqpoint{2.314503in}{1.080433in}}%
\pgfpathlineto{\pgfqpoint{2.300122in}{1.093676in}}%
\pgfpathlineto{\pgfqpoint{2.294182in}{1.093073in}}%
\pgfpathlineto{\pgfqpoint{2.289598in}{1.088435in}}%
\pgfpathlineto{\pgfqpoint{2.283241in}{1.088142in}}%
\pgfpathlineto{\pgfqpoint{2.281936in}{1.092901in}}%
\pgfpathlineto{\pgfqpoint{2.287143in}{1.100325in}}%
\pgfpathlineto{\pgfqpoint{2.288667in}{1.106429in}}%
\pgfpathlineto{\pgfqpoint{2.288467in}{1.113713in}}%
\pgfpathlineto{\pgfqpoint{2.292494in}{1.120777in}}%
\pgfpathlineto{\pgfqpoint{2.287382in}{1.127701in}}%
\pgfpathlineto{\pgfqpoint{2.285748in}{1.132441in}}%
\pgfpathlineto{\pgfqpoint{2.282277in}{1.134362in}}%
\pgfpathlineto{\pgfqpoint{2.288955in}{1.139422in}}%
\pgfpathlineto{\pgfqpoint{2.296120in}{1.144067in}}%
\pgfpathlineto{\pgfqpoint{2.297453in}{1.137486in}}%
\pgfpathlineto{\pgfqpoint{2.308461in}{1.129091in}}%
\pgfpathlineto{\pgfqpoint{2.317016in}{1.135906in}}%
\pgfpathlineto{\pgfqpoint{2.326824in}{1.138484in}}%
\pgfpathlineto{\pgfqpoint{2.332965in}{1.140303in}}%
\pgfpathlineto{\pgfqpoint{2.338933in}{1.139975in}}%
\pgfpathlineto{\pgfqpoint{2.341607in}{1.134315in}}%
\pgfpathlineto{\pgfqpoint{2.339247in}{1.131363in}}%
\pgfpathlineto{\pgfqpoint{2.347373in}{1.122044in}}%
\pgfpathlineto{\pgfqpoint{2.336897in}{1.108708in}}%
\pgfpathlineto{\pgfqpoint{2.344556in}{1.097617in}}%
\pgfpathlineto{\pgfqpoint{2.355335in}{1.096669in}}%
\pgfpathlineto{\pgfqpoint{2.362396in}{1.092621in}}%
\pgfpathclose%
\pgfusepath{fill}%
\end{pgfscope}%
\begin{pgfscope}%
\pgfpathrectangle{\pgfqpoint{0.100000in}{0.100000in}}{\pgfqpoint{3.007045in}{1.925000in}}%
\pgfusepath{clip}%
\pgfsetbuttcap%
\pgfsetmiterjoin%
\definecolor{currentfill}{rgb}{0.529642,0.803506,0.524198}%
\pgfsetfillcolor{currentfill}%
\pgfsetlinewidth{0.000000pt}%
\definecolor{currentstroke}{rgb}{0.000000,0.000000,0.000000}%
\pgfsetstrokecolor{currentstroke}%
\pgfsetstrokeopacity{0.000000}%
\pgfsetdash{}{0pt}%
\pgfpathmoveto{\pgfqpoint{2.231222in}{1.236425in}}%
\pgfpathlineto{\pgfqpoint{2.212181in}{1.234336in}}%
\pgfpathlineto{\pgfqpoint{2.211166in}{1.244697in}}%
\pgfpathlineto{\pgfqpoint{2.204170in}{1.247864in}}%
\pgfpathlineto{\pgfqpoint{2.202654in}{1.264977in}}%
\pgfpathlineto{\pgfqpoint{2.195116in}{1.264251in}}%
\pgfpathlineto{\pgfqpoint{2.192034in}{1.266835in}}%
\pgfpathlineto{\pgfqpoint{2.190977in}{1.278330in}}%
\pgfpathlineto{\pgfqpoint{2.173650in}{1.276857in}}%
\pgfpathlineto{\pgfqpoint{2.172150in}{1.293918in}}%
\pgfpathlineto{\pgfqpoint{2.195043in}{1.295965in}}%
\pgfpathlineto{\pgfqpoint{2.196588in}{1.278782in}}%
\pgfpathlineto{\pgfqpoint{2.211433in}{1.280191in}}%
\pgfpathlineto{\pgfqpoint{2.210917in}{1.285878in}}%
\pgfpathlineto{\pgfqpoint{2.221882in}{1.287153in}}%
\pgfpathlineto{\pgfqpoint{2.224496in}{1.264365in}}%
\pgfpathlineto{\pgfqpoint{2.228217in}{1.264785in}}%
\pgfpathlineto{\pgfqpoint{2.231222in}{1.236425in}}%
\pgfpathclose%
\pgfusepath{fill}%
\end{pgfscope}%
\begin{pgfscope}%
\pgfpathrectangle{\pgfqpoint{0.100000in}{0.100000in}}{\pgfqpoint{3.007045in}{1.925000in}}%
\pgfusepath{clip}%
\pgfsetbuttcap%
\pgfsetmiterjoin%
\definecolor{currentfill}{rgb}{0.662345,0.864022,0.638001}%
\pgfsetfillcolor{currentfill}%
\pgfsetlinewidth{0.000000pt}%
\definecolor{currentstroke}{rgb}{0.000000,0.000000,0.000000}%
\pgfsetstrokecolor{currentstroke}%
\pgfsetstrokeopacity{0.000000}%
\pgfsetdash{}{0pt}%
\pgfpathmoveto{\pgfqpoint{1.980505in}{0.873918in}}%
\pgfpathlineto{\pgfqpoint{1.972810in}{0.872450in}}%
\pgfpathlineto{\pgfqpoint{1.968811in}{0.876703in}}%
\pgfpathlineto{\pgfqpoint{1.963244in}{0.878965in}}%
\pgfpathlineto{\pgfqpoint{1.956817in}{0.878848in}}%
\pgfpathlineto{\pgfqpoint{1.956926in}{0.874466in}}%
\pgfpathlineto{\pgfqpoint{1.950043in}{0.876695in}}%
\pgfpathlineto{\pgfqpoint{1.945167in}{0.874395in}}%
\pgfpathlineto{\pgfqpoint{1.941240in}{0.873795in}}%
\pgfpathlineto{\pgfqpoint{1.933967in}{0.876483in}}%
\pgfpathlineto{\pgfqpoint{1.927950in}{0.876198in}}%
\pgfpathlineto{\pgfqpoint{1.927775in}{0.895762in}}%
\pgfpathlineto{\pgfqpoint{1.941745in}{0.895933in}}%
\pgfpathlineto{\pgfqpoint{1.944537in}{0.902457in}}%
\pgfpathlineto{\pgfqpoint{1.944456in}{0.907639in}}%
\pgfpathlineto{\pgfqpoint{1.967250in}{0.907985in}}%
\pgfpathlineto{\pgfqpoint{1.966993in}{0.915480in}}%
\pgfpathlineto{\pgfqpoint{1.969720in}{0.923885in}}%
\pgfpathlineto{\pgfqpoint{1.975423in}{0.932334in}}%
\pgfpathlineto{\pgfqpoint{1.984303in}{0.932387in}}%
\pgfpathlineto{\pgfqpoint{1.985207in}{0.903228in}}%
\pgfpathlineto{\pgfqpoint{1.986161in}{0.874058in}}%
\pgfpathlineto{\pgfqpoint{1.980505in}{0.873918in}}%
\pgfpathclose%
\pgfusepath{fill}%
\end{pgfscope}%
\begin{pgfscope}%
\pgfpathrectangle{\pgfqpoint{0.100000in}{0.100000in}}{\pgfqpoint{3.007045in}{1.925000in}}%
\pgfusepath{clip}%
\pgfsetbuttcap%
\pgfsetmiterjoin%
\definecolor{currentfill}{rgb}{0.751173,0.901423,0.724491}%
\pgfsetfillcolor{currentfill}%
\pgfsetlinewidth{0.000000pt}%
\definecolor{currentstroke}{rgb}{0.000000,0.000000,0.000000}%
\pgfsetstrokecolor{currentstroke}%
\pgfsetstrokeopacity{0.000000}%
\pgfsetdash{}{0pt}%
\pgfpathmoveto{\pgfqpoint{1.553928in}{0.768166in}}%
\pgfpathlineto{\pgfqpoint{1.555230in}{0.797022in}}%
\pgfpathlineto{\pgfqpoint{1.556351in}{0.822025in}}%
\pgfpathlineto{\pgfqpoint{1.560430in}{0.818079in}}%
\pgfpathlineto{\pgfqpoint{1.570469in}{0.816172in}}%
\pgfpathlineto{\pgfqpoint{1.574901in}{0.817578in}}%
\pgfpathlineto{\pgfqpoint{1.581430in}{0.811070in}}%
\pgfpathlineto{\pgfqpoint{1.586260in}{0.813360in}}%
\pgfpathlineto{\pgfqpoint{1.588268in}{0.817057in}}%
\pgfpathlineto{\pgfqpoint{1.598818in}{0.813845in}}%
\pgfpathlineto{\pgfqpoint{1.605264in}{0.805706in}}%
\pgfpathlineto{\pgfqpoint{1.610850in}{0.805245in}}%
\pgfpathlineto{\pgfqpoint{1.608887in}{0.798669in}}%
\pgfpathlineto{\pgfqpoint{1.607913in}{0.770715in}}%
\pgfpathlineto{\pgfqpoint{1.583500in}{0.771553in}}%
\pgfpathlineto{\pgfqpoint{1.583365in}{0.766882in}}%
\pgfpathlineto{\pgfqpoint{1.553928in}{0.768166in}}%
\pgfpathclose%
\pgfusepath{fill}%
\end{pgfscope}%
\begin{pgfscope}%
\pgfpathrectangle{\pgfqpoint{0.100000in}{0.100000in}}{\pgfqpoint{3.007045in}{1.925000in}}%
\pgfusepath{clip}%
\pgfsetbuttcap%
\pgfsetmiterjoin%
\definecolor{currentfill}{rgb}{0.524106,0.800923,0.519646}%
\pgfsetfillcolor{currentfill}%
\pgfsetlinewidth{0.000000pt}%
\definecolor{currentstroke}{rgb}{0.000000,0.000000,0.000000}%
\pgfsetstrokecolor{currentstroke}%
\pgfsetstrokeopacity{0.000000}%
\pgfsetdash{}{0pt}%
\pgfpathmoveto{\pgfqpoint{1.872670in}{1.072352in}}%
\pgfpathlineto{\pgfqpoint{1.849904in}{1.072618in}}%
\pgfpathlineto{\pgfqpoint{1.849656in}{1.062106in}}%
\pgfpathlineto{\pgfqpoint{1.846715in}{1.062147in}}%
\pgfpathlineto{\pgfqpoint{1.843696in}{1.056469in}}%
\pgfpathlineto{\pgfqpoint{1.834159in}{1.056680in}}%
\pgfpathlineto{\pgfqpoint{1.834274in}{1.060508in}}%
\pgfpathlineto{\pgfqpoint{1.820896in}{1.060894in}}%
\pgfpathlineto{\pgfqpoint{1.821187in}{1.069825in}}%
\pgfpathlineto{\pgfqpoint{1.820190in}{1.081544in}}%
\pgfpathlineto{\pgfqpoint{1.820514in}{1.104792in}}%
\pgfpathlineto{\pgfqpoint{1.849007in}{1.104433in}}%
\pgfpathlineto{\pgfqpoint{1.849349in}{1.129915in}}%
\pgfpathlineto{\pgfqpoint{1.872212in}{1.129384in}}%
\pgfpathlineto{\pgfqpoint{1.872048in}{1.113887in}}%
\pgfpathlineto{\pgfqpoint{1.871618in}{1.080502in}}%
\pgfpathlineto{\pgfqpoint{1.872670in}{1.072352in}}%
\pgfpathclose%
\pgfusepath{fill}%
\end{pgfscope}%
\begin{pgfscope}%
\pgfpathrectangle{\pgfqpoint{0.100000in}{0.100000in}}{\pgfqpoint{3.007045in}{1.925000in}}%
\pgfusepath{clip}%
\pgfsetbuttcap%
\pgfsetmiterjoin%
\definecolor{currentfill}{rgb}{0.607151,0.839677,0.587928}%
\pgfsetfillcolor{currentfill}%
\pgfsetlinewidth{0.000000pt}%
\definecolor{currentstroke}{rgb}{0.000000,0.000000,0.000000}%
\pgfsetstrokecolor{currentstroke}%
\pgfsetstrokeopacity{0.000000}%
\pgfsetdash{}{0pt}%
\pgfpathmoveto{\pgfqpoint{0.433096in}{0.983854in}}%
\pgfpathlineto{\pgfqpoint{0.425437in}{0.987474in}}%
\pgfpathlineto{\pgfqpoint{0.427947in}{0.992997in}}%
\pgfpathlineto{\pgfqpoint{0.436545in}{0.987083in}}%
\pgfpathlineto{\pgfqpoint{0.433096in}{0.983854in}}%
\pgfpathclose%
\pgfusepath{fill}%
\end{pgfscope}%
\begin{pgfscope}%
\pgfpathrectangle{\pgfqpoint{0.100000in}{0.100000in}}{\pgfqpoint{3.007045in}{1.925000in}}%
\pgfusepath{clip}%
\pgfsetbuttcap%
\pgfsetmiterjoin%
\definecolor{currentfill}{rgb}{0.607151,0.839677,0.587928}%
\pgfsetfillcolor{currentfill}%
\pgfsetlinewidth{0.000000pt}%
\definecolor{currentstroke}{rgb}{0.000000,0.000000,0.000000}%
\pgfsetstrokecolor{currentstroke}%
\pgfsetstrokeopacity{0.000000}%
\pgfsetdash{}{0pt}%
\pgfpathmoveto{\pgfqpoint{0.411502in}{0.985534in}}%
\pgfpathlineto{\pgfqpoint{0.407141in}{0.993168in}}%
\pgfpathlineto{\pgfqpoint{0.412033in}{0.994414in}}%
\pgfpathlineto{\pgfqpoint{0.419874in}{0.989359in}}%
\pgfpathlineto{\pgfqpoint{0.418366in}{0.986829in}}%
\pgfpathlineto{\pgfqpoint{0.411502in}{0.985534in}}%
\pgfpathclose%
\pgfusepath{fill}%
\end{pgfscope}%
\begin{pgfscope}%
\pgfpathrectangle{\pgfqpoint{0.100000in}{0.100000in}}{\pgfqpoint{3.007045in}{1.925000in}}%
\pgfusepath{clip}%
\pgfsetbuttcap%
\pgfsetmiterjoin%
\definecolor{currentfill}{rgb}{0.607151,0.839677,0.587928}%
\pgfsetfillcolor{currentfill}%
\pgfsetlinewidth{0.000000pt}%
\definecolor{currentstroke}{rgb}{0.000000,0.000000,0.000000}%
\pgfsetstrokecolor{currentstroke}%
\pgfsetstrokeopacity{0.000000}%
\pgfsetdash{}{0pt}%
\pgfpathmoveto{\pgfqpoint{0.464444in}{1.039777in}}%
\pgfpathlineto{\pgfqpoint{0.456601in}{1.011141in}}%
\pgfpathlineto{\pgfqpoint{0.453370in}{1.006642in}}%
\pgfpathlineto{\pgfqpoint{0.447005in}{1.011573in}}%
\pgfpathlineto{\pgfqpoint{0.440321in}{1.011780in}}%
\pgfpathlineto{\pgfqpoint{0.433003in}{1.014607in}}%
\pgfpathlineto{\pgfqpoint{0.427568in}{1.019544in}}%
\pgfpathlineto{\pgfqpoint{0.416032in}{1.023886in}}%
\pgfpathlineto{\pgfqpoint{0.403356in}{1.025371in}}%
\pgfpathlineto{\pgfqpoint{0.401817in}{1.031469in}}%
\pgfpathlineto{\pgfqpoint{0.395689in}{1.036947in}}%
\pgfpathlineto{\pgfqpoint{0.400361in}{1.044762in}}%
\pgfpathlineto{\pgfqpoint{0.399420in}{1.048155in}}%
\pgfpathlineto{\pgfqpoint{0.402613in}{1.053373in}}%
\pgfpathlineto{\pgfqpoint{0.400409in}{1.057912in}}%
\pgfpathlineto{\pgfqpoint{0.403882in}{1.063816in}}%
\pgfpathlineto{\pgfqpoint{0.405983in}{1.073145in}}%
\pgfpathlineto{\pgfqpoint{0.400206in}{1.075730in}}%
\pgfpathlineto{\pgfqpoint{0.395182in}{1.083867in}}%
\pgfpathlineto{\pgfqpoint{0.399136in}{1.090944in}}%
\pgfpathlineto{\pgfqpoint{0.398305in}{1.096127in}}%
\pgfpathlineto{\pgfqpoint{0.393442in}{1.098296in}}%
\pgfpathlineto{\pgfqpoint{0.388518in}{1.111365in}}%
\pgfpathlineto{\pgfqpoint{0.382876in}{1.115594in}}%
\pgfpathlineto{\pgfqpoint{0.382061in}{1.124859in}}%
\pgfpathlineto{\pgfqpoint{0.406868in}{1.117231in}}%
\pgfpathlineto{\pgfqpoint{0.440258in}{1.107424in}}%
\pgfpathlineto{\pgfqpoint{0.441277in}{1.107129in}}%
\pgfpathlineto{\pgfqpoint{0.438131in}{1.096030in}}%
\pgfpathlineto{\pgfqpoint{0.443707in}{1.094459in}}%
\pgfpathlineto{\pgfqpoint{0.442087in}{1.088884in}}%
\pgfpathlineto{\pgfqpoint{0.445664in}{1.083905in}}%
\pgfpathlineto{\pgfqpoint{0.451207in}{1.080317in}}%
\pgfpathlineto{\pgfqpoint{0.449631in}{1.074733in}}%
\pgfpathlineto{\pgfqpoint{0.453296in}{1.073681in}}%
\pgfpathlineto{\pgfqpoint{0.451737in}{1.068135in}}%
\pgfpathlineto{\pgfqpoint{0.459131in}{1.065999in}}%
\pgfpathlineto{\pgfqpoint{0.457558in}{1.060425in}}%
\pgfpathlineto{\pgfqpoint{0.463546in}{1.059117in}}%
\pgfpathlineto{\pgfqpoint{0.461538in}{1.053350in}}%
\pgfpathlineto{\pgfqpoint{0.465944in}{1.051393in}}%
\pgfpathlineto{\pgfqpoint{0.464444in}{1.039777in}}%
\pgfpathclose%
\pgfusepath{fill}%
\end{pgfscope}%
\begin{pgfscope}%
\pgfpathrectangle{\pgfqpoint{0.100000in}{0.100000in}}{\pgfqpoint{3.007045in}{1.925000in}}%
\pgfusepath{clip}%
\pgfsetbuttcap%
\pgfsetmiterjoin%
\definecolor{currentfill}{rgb}{0.546251,0.811257,0.537855}%
\pgfsetfillcolor{currentfill}%
\pgfsetlinewidth{0.000000pt}%
\definecolor{currentstroke}{rgb}{0.000000,0.000000,0.000000}%
\pgfsetstrokecolor{currentstroke}%
\pgfsetstrokeopacity{0.000000}%
\pgfsetdash{}{0pt}%
\pgfpathmoveto{\pgfqpoint{2.825651in}{1.545463in}}%
\pgfpathlineto{\pgfqpoint{2.807383in}{1.541913in}}%
\pgfpathlineto{\pgfqpoint{2.803917in}{1.557332in}}%
\pgfpathlineto{\pgfqpoint{2.794553in}{1.564410in}}%
\pgfpathlineto{\pgfqpoint{2.793875in}{1.572201in}}%
\pgfpathlineto{\pgfqpoint{2.787488in}{1.586749in}}%
\pgfpathlineto{\pgfqpoint{2.787394in}{1.596652in}}%
\pgfpathlineto{\pgfqpoint{2.789757in}{1.602899in}}%
\pgfpathlineto{\pgfqpoint{2.787188in}{1.608949in}}%
\pgfpathlineto{\pgfqpoint{2.786419in}{1.617481in}}%
\pgfpathlineto{\pgfqpoint{2.780749in}{1.624141in}}%
\pgfpathlineto{\pgfqpoint{2.780372in}{1.636703in}}%
\pgfpathlineto{\pgfqpoint{2.777816in}{1.638392in}}%
\pgfpathlineto{\pgfqpoint{2.778407in}{1.644079in}}%
\pgfpathlineto{\pgfqpoint{2.776771in}{1.649915in}}%
\pgfpathlineto{\pgfqpoint{2.809280in}{1.658194in}}%
\pgfpathlineto{\pgfqpoint{2.812663in}{1.658624in}}%
\pgfpathlineto{\pgfqpoint{2.816481in}{1.647652in}}%
\pgfpathlineto{\pgfqpoint{2.815249in}{1.643795in}}%
\pgfpathlineto{\pgfqpoint{2.823531in}{1.641882in}}%
\pgfpathlineto{\pgfqpoint{2.822013in}{1.635422in}}%
\pgfpathlineto{\pgfqpoint{2.827912in}{1.634106in}}%
\pgfpathlineto{\pgfqpoint{2.826283in}{1.627777in}}%
\pgfpathlineto{\pgfqpoint{2.825014in}{1.621968in}}%
\pgfpathlineto{\pgfqpoint{2.819281in}{1.624179in}}%
\pgfpathlineto{\pgfqpoint{2.817950in}{1.618314in}}%
\pgfpathlineto{\pgfqpoint{2.810189in}{1.620019in}}%
\pgfpathlineto{\pgfqpoint{2.808017in}{1.607896in}}%
\pgfpathlineto{\pgfqpoint{2.808838in}{1.600051in}}%
\pgfpathlineto{\pgfqpoint{2.811115in}{1.594372in}}%
\pgfpathlineto{\pgfqpoint{2.819866in}{1.593654in}}%
\pgfpathlineto{\pgfqpoint{2.815683in}{1.583923in}}%
\pgfpathlineto{\pgfqpoint{2.811802in}{1.580671in}}%
\pgfpathlineto{\pgfqpoint{2.821564in}{1.578723in}}%
\pgfpathlineto{\pgfqpoint{2.821240in}{1.572694in}}%
\pgfpathlineto{\pgfqpoint{2.827305in}{1.571893in}}%
\pgfpathlineto{\pgfqpoint{2.826130in}{1.558806in}}%
\pgfpathlineto{\pgfqpoint{2.829522in}{1.557864in}}%
\pgfpathlineto{\pgfqpoint{2.829267in}{1.550294in}}%
\pgfpathlineto{\pgfqpoint{2.825651in}{1.545463in}}%
\pgfpathclose%
\pgfusepath{fill}%
\end{pgfscope}%
\begin{pgfscope}%
\pgfpathrectangle{\pgfqpoint{0.100000in}{0.100000in}}{\pgfqpoint{3.007045in}{1.925000in}}%
\pgfusepath{clip}%
\pgfsetbuttcap%
\pgfsetmiterjoin%
\definecolor{currentfill}{rgb}{0.551788,0.813841,0.542407}%
\pgfsetfillcolor{currentfill}%
\pgfsetlinewidth{0.000000pt}%
\definecolor{currentstroke}{rgb}{0.000000,0.000000,0.000000}%
\pgfsetstrokecolor{currentstroke}%
\pgfsetstrokeopacity{0.000000}%
\pgfsetdash{}{0pt}%
\pgfpathmoveto{\pgfqpoint{1.969121in}{0.688091in}}%
\pgfpathlineto{\pgfqpoint{1.966832in}{0.693993in}}%
\pgfpathlineto{\pgfqpoint{1.970011in}{0.695703in}}%
\pgfpathlineto{\pgfqpoint{1.968808in}{0.707431in}}%
\pgfpathlineto{\pgfqpoint{1.970224in}{0.710328in}}%
\pgfpathlineto{\pgfqpoint{1.988311in}{0.711075in}}%
\pgfpathlineto{\pgfqpoint{1.993233in}{0.713829in}}%
\pgfpathlineto{\pgfqpoint{1.997671in}{0.713980in}}%
\pgfpathlineto{\pgfqpoint{1.998667in}{0.705092in}}%
\pgfpathlineto{\pgfqpoint{2.003803in}{0.705398in}}%
\pgfpathlineto{\pgfqpoint{2.003446in}{0.714212in}}%
\pgfpathlineto{\pgfqpoint{2.010976in}{0.714602in}}%
\pgfpathlineto{\pgfqpoint{2.011147in}{0.717467in}}%
\pgfpathlineto{\pgfqpoint{2.020779in}{0.710545in}}%
\pgfpathlineto{\pgfqpoint{2.018058in}{0.708669in}}%
\pgfpathlineto{\pgfqpoint{2.016613in}{0.701553in}}%
\pgfpathlineto{\pgfqpoint{2.012854in}{0.699525in}}%
\pgfpathlineto{\pgfqpoint{2.011898in}{0.691368in}}%
\pgfpathlineto{\pgfqpoint{2.013470in}{0.675004in}}%
\pgfpathlineto{\pgfqpoint{2.013078in}{0.656435in}}%
\pgfpathlineto{\pgfqpoint{2.013321in}{0.650638in}}%
\pgfpathlineto{\pgfqpoint{1.989789in}{0.649517in}}%
\pgfpathlineto{\pgfqpoint{1.981608in}{0.653144in}}%
\pgfpathlineto{\pgfqpoint{1.978327in}{0.664514in}}%
\pgfpathlineto{\pgfqpoint{1.986197in}{0.669162in}}%
\pgfpathlineto{\pgfqpoint{1.986876in}{0.672737in}}%
\pgfpathlineto{\pgfqpoint{1.993029in}{0.676683in}}%
\pgfpathlineto{\pgfqpoint{1.995075in}{0.682326in}}%
\pgfpathlineto{\pgfqpoint{1.987667in}{0.684255in}}%
\pgfpathlineto{\pgfqpoint{1.989120in}{0.691955in}}%
\pgfpathlineto{\pgfqpoint{1.978768in}{0.688413in}}%
\pgfpathlineto{\pgfqpoint{1.969121in}{0.688091in}}%
\pgfpathclose%
\pgfusepath{fill}%
\end{pgfscope}%
\begin{pgfscope}%
\pgfpathrectangle{\pgfqpoint{0.100000in}{0.100000in}}{\pgfqpoint{3.007045in}{1.925000in}}%
\pgfusepath{clip}%
\pgfsetbuttcap%
\pgfsetmiterjoin%
\definecolor{currentfill}{rgb}{0.868051,0.948789,0.846444}%
\pgfsetfillcolor{currentfill}%
\pgfsetlinewidth{0.000000pt}%
\definecolor{currentstroke}{rgb}{0.000000,0.000000,0.000000}%
\pgfsetstrokecolor{currentstroke}%
\pgfsetstrokeopacity{0.000000}%
\pgfsetdash{}{0pt}%
\pgfpathmoveto{\pgfqpoint{0.441277in}{1.107129in}}%
\pgfpathlineto{\pgfqpoint{0.440258in}{1.107424in}}%
\pgfpathlineto{\pgfqpoint{0.435906in}{1.120684in}}%
\pgfpathlineto{\pgfqpoint{0.432260in}{1.122040in}}%
\pgfpathlineto{\pgfqpoint{0.428309in}{1.126977in}}%
\pgfpathlineto{\pgfqpoint{0.423734in}{1.137658in}}%
\pgfpathlineto{\pgfqpoint{0.426809in}{1.147328in}}%
\pgfpathlineto{\pgfqpoint{0.430464in}{1.147237in}}%
\pgfpathlineto{\pgfqpoint{0.433028in}{1.157462in}}%
\pgfpathlineto{\pgfqpoint{0.421673in}{1.178083in}}%
\pgfpathlineto{\pgfqpoint{0.450589in}{1.194118in}}%
\pgfpathlineto{\pgfqpoint{0.461047in}{1.194974in}}%
\pgfpathlineto{\pgfqpoint{0.473609in}{1.193492in}}%
\pgfpathlineto{\pgfqpoint{0.492430in}{1.204219in}}%
\pgfpathlineto{\pgfqpoint{0.498018in}{1.202710in}}%
\pgfpathlineto{\pgfqpoint{0.502749in}{1.206653in}}%
\pgfpathlineto{\pgfqpoint{0.503936in}{1.210812in}}%
\pgfpathlineto{\pgfqpoint{0.521553in}{1.220832in}}%
\pgfpathlineto{\pgfqpoint{0.522889in}{1.217818in}}%
\pgfpathlineto{\pgfqpoint{0.530069in}{1.215499in}}%
\pgfpathlineto{\pgfqpoint{0.532191in}{1.211551in}}%
\pgfpathlineto{\pgfqpoint{0.532644in}{1.204682in}}%
\pgfpathlineto{\pgfqpoint{0.537345in}{1.200997in}}%
\pgfpathlineto{\pgfqpoint{0.539193in}{1.195591in}}%
\pgfpathlineto{\pgfqpoint{0.543469in}{1.191912in}}%
\pgfpathlineto{\pgfqpoint{0.540581in}{1.185427in}}%
\pgfpathlineto{\pgfqpoint{0.543712in}{1.184135in}}%
\pgfpathlineto{\pgfqpoint{0.543660in}{1.171437in}}%
\pgfpathlineto{\pgfqpoint{0.546742in}{1.170366in}}%
\pgfpathlineto{\pgfqpoint{0.552868in}{1.161985in}}%
\pgfpathlineto{\pgfqpoint{0.554321in}{1.151293in}}%
\pgfpathlineto{\pgfqpoint{0.551954in}{1.148285in}}%
\pgfpathlineto{\pgfqpoint{0.551392in}{1.137382in}}%
\pgfpathlineto{\pgfqpoint{0.554104in}{1.127797in}}%
\pgfpathlineto{\pgfqpoint{0.553369in}{1.124234in}}%
\pgfpathlineto{\pgfqpoint{0.558630in}{1.113262in}}%
\pgfpathlineto{\pgfqpoint{0.556496in}{1.110688in}}%
\pgfpathlineto{\pgfqpoint{0.557824in}{1.101480in}}%
\pgfpathlineto{\pgfqpoint{0.556142in}{1.096690in}}%
\pgfpathlineto{\pgfqpoint{0.557260in}{1.082859in}}%
\pgfpathlineto{\pgfqpoint{0.554172in}{1.076466in}}%
\pgfpathlineto{\pgfqpoint{0.522151in}{1.084892in}}%
\pgfpathlineto{\pgfqpoint{0.475084in}{1.097667in}}%
\pgfpathlineto{\pgfqpoint{0.441277in}{1.107129in}}%
\pgfpathclose%
\pgfusepath{fill}%
\end{pgfscope}%
\begin{pgfscope}%
\pgfpathrectangle{\pgfqpoint{0.100000in}{0.100000in}}{\pgfqpoint{3.007045in}{1.925000in}}%
\pgfusepath{clip}%
\pgfsetbuttcap%
\pgfsetmiterjoin%
\definecolor{currentfill}{rgb}{0.474279,0.777670,0.478677}%
\pgfsetfillcolor{currentfill}%
\pgfsetlinewidth{0.000000pt}%
\definecolor{currentstroke}{rgb}{0.000000,0.000000,0.000000}%
\pgfsetstrokecolor{currentstroke}%
\pgfsetstrokeopacity{0.000000}%
\pgfsetdash{}{0pt}%
\pgfpathmoveto{\pgfqpoint{0.856506in}{1.901711in}}%
\pgfpathlineto{\pgfqpoint{0.892466in}{1.893289in}}%
\pgfpathlineto{\pgfqpoint{0.913025in}{1.888772in}}%
\pgfpathlineto{\pgfqpoint{0.912328in}{1.878097in}}%
\pgfpathlineto{\pgfqpoint{0.913309in}{1.873552in}}%
\pgfpathlineto{\pgfqpoint{0.910046in}{1.867809in}}%
\pgfpathlineto{\pgfqpoint{0.901194in}{1.867977in}}%
\pgfpathlineto{\pgfqpoint{0.901644in}{1.862062in}}%
\pgfpathlineto{\pgfqpoint{0.896692in}{1.839394in}}%
\pgfpathlineto{\pgfqpoint{0.889217in}{1.841046in}}%
\pgfpathlineto{\pgfqpoint{0.886583in}{1.827387in}}%
\pgfpathlineto{\pgfqpoint{0.880330in}{1.828801in}}%
\pgfpathlineto{\pgfqpoint{0.877958in}{1.823506in}}%
\pgfpathlineto{\pgfqpoint{0.873246in}{1.821840in}}%
\pgfpathlineto{\pgfqpoint{0.862566in}{1.825406in}}%
\pgfpathlineto{\pgfqpoint{0.863113in}{1.830232in}}%
\pgfpathlineto{\pgfqpoint{0.860078in}{1.837613in}}%
\pgfpathlineto{\pgfqpoint{0.859696in}{1.849285in}}%
\pgfpathlineto{\pgfqpoint{0.845385in}{1.846522in}}%
\pgfpathlineto{\pgfqpoint{0.844425in}{1.850601in}}%
\pgfpathlineto{\pgfqpoint{0.848850in}{1.869247in}}%
\pgfpathlineto{\pgfqpoint{0.856506in}{1.901711in}}%
\pgfpathclose%
\pgfusepath{fill}%
\end{pgfscope}%
\begin{pgfscope}%
\pgfpathrectangle{\pgfqpoint{0.100000in}{0.100000in}}{\pgfqpoint{3.007045in}{1.925000in}}%
\pgfusepath{clip}%
\pgfsetbuttcap%
\pgfsetmiterjoin%
\definecolor{currentfill}{rgb}{0.557324,0.816424,0.546959}%
\pgfsetfillcolor{currentfill}%
\pgfsetlinewidth{0.000000pt}%
\definecolor{currentstroke}{rgb}{0.000000,0.000000,0.000000}%
\pgfsetstrokecolor{currentstroke}%
\pgfsetstrokeopacity{0.000000}%
\pgfsetdash{}{0pt}%
\pgfpathmoveto{\pgfqpoint{2.113075in}{0.848655in}}%
\pgfpathlineto{\pgfqpoint{2.112919in}{0.865058in}}%
\pgfpathlineto{\pgfqpoint{2.110591in}{0.871650in}}%
\pgfpathlineto{\pgfqpoint{2.106779in}{0.871379in}}%
\pgfpathlineto{\pgfqpoint{2.106178in}{0.880380in}}%
\pgfpathlineto{\pgfqpoint{2.108156in}{0.880516in}}%
\pgfpathlineto{\pgfqpoint{2.107278in}{0.897155in}}%
\pgfpathlineto{\pgfqpoint{2.110613in}{0.897382in}}%
\pgfpathlineto{\pgfqpoint{2.114334in}{0.905115in}}%
\pgfpathlineto{\pgfqpoint{2.117945in}{0.906798in}}%
\pgfpathlineto{\pgfqpoint{2.129176in}{0.907553in}}%
\pgfpathlineto{\pgfqpoint{2.129095in}{0.910042in}}%
\pgfpathlineto{\pgfqpoint{2.142348in}{0.908782in}}%
\pgfpathlineto{\pgfqpoint{2.147248in}{0.909640in}}%
\pgfpathlineto{\pgfqpoint{2.150055in}{0.903459in}}%
\pgfpathlineto{\pgfqpoint{2.151351in}{0.884373in}}%
\pgfpathlineto{\pgfqpoint{2.139583in}{0.883616in}}%
\pgfpathlineto{\pgfqpoint{2.145830in}{0.876417in}}%
\pgfpathlineto{\pgfqpoint{2.144807in}{0.847920in}}%
\pgfpathlineto{\pgfqpoint{2.123839in}{0.846519in}}%
\pgfpathlineto{\pgfqpoint{2.123649in}{0.849414in}}%
\pgfpathlineto{\pgfqpoint{2.113075in}{0.848655in}}%
\pgfpathclose%
\pgfusepath{fill}%
\end{pgfscope}%
\begin{pgfscope}%
\pgfpathrectangle{\pgfqpoint{0.100000in}{0.100000in}}{\pgfqpoint{3.007045in}{1.925000in}}%
\pgfusepath{clip}%
\pgfsetbuttcap%
\pgfsetmiterjoin%
\definecolor{currentfill}{rgb}{0.000000,0.271711,0.107974}%
\pgfsetfillcolor{currentfill}%
\pgfsetlinewidth{0.000000pt}%
\definecolor{currentstroke}{rgb}{0.000000,0.000000,0.000000}%
\pgfsetstrokecolor{currentstroke}%
\pgfsetstrokeopacity{0.000000}%
\pgfsetdash{}{0pt}%
\pgfpathmoveto{\pgfqpoint{0.556363in}{1.910254in}}%
\pgfpathlineto{\pgfqpoint{0.556850in}{1.906039in}}%
\pgfpathlineto{\pgfqpoint{0.549783in}{1.902072in}}%
\pgfpathlineto{\pgfqpoint{0.546025in}{1.896043in}}%
\pgfpathlineto{\pgfqpoint{0.536363in}{1.893635in}}%
\pgfpathlineto{\pgfqpoint{0.512685in}{1.901288in}}%
\pgfpathlineto{\pgfqpoint{0.510848in}{1.895622in}}%
\pgfpathlineto{\pgfqpoint{0.487655in}{1.903257in}}%
\pgfpathlineto{\pgfqpoint{0.474818in}{1.908630in}}%
\pgfpathlineto{\pgfqpoint{0.476195in}{1.922659in}}%
\pgfpathlineto{\pgfqpoint{0.474995in}{1.928796in}}%
\pgfpathlineto{\pgfqpoint{0.470655in}{1.935298in}}%
\pgfpathlineto{\pgfqpoint{0.470537in}{1.940585in}}%
\pgfpathlineto{\pgfqpoint{0.473848in}{1.954693in}}%
\pgfpathlineto{\pgfqpoint{0.481028in}{1.966614in}}%
\pgfpathlineto{\pgfqpoint{0.489160in}{1.956885in}}%
\pgfpathlineto{\pgfqpoint{0.494976in}{1.953083in}}%
\pgfpathlineto{\pgfqpoint{0.502419in}{1.944351in}}%
\pgfpathlineto{\pgfqpoint{0.509754in}{1.940761in}}%
\pgfpathlineto{\pgfqpoint{0.515790in}{1.939825in}}%
\pgfpathlineto{\pgfqpoint{0.520125in}{1.935808in}}%
\pgfpathlineto{\pgfqpoint{0.533931in}{1.930174in}}%
\pgfpathlineto{\pgfqpoint{0.539348in}{1.932466in}}%
\pgfpathlineto{\pgfqpoint{0.542272in}{1.925230in}}%
\pgfpathlineto{\pgfqpoint{0.555512in}{1.925149in}}%
\pgfpathlineto{\pgfqpoint{0.557364in}{1.921639in}}%
\pgfpathlineto{\pgfqpoint{0.555607in}{1.915548in}}%
\pgfpathlineto{\pgfqpoint{0.556363in}{1.910254in}}%
\pgfpathclose%
\pgfusepath{fill}%
\end{pgfscope}%
\begin{pgfscope}%
\pgfpathrectangle{\pgfqpoint{0.100000in}{0.100000in}}{\pgfqpoint{3.007045in}{1.925000in}}%
\pgfusepath{clip}%
\pgfsetbuttcap%
\pgfsetmiterjoin%
\definecolor{currentfill}{rgb}{0.524106,0.800923,0.519646}%
\pgfsetfillcolor{currentfill}%
\pgfsetlinewidth{0.000000pt}%
\definecolor{currentstroke}{rgb}{0.000000,0.000000,0.000000}%
\pgfsetstrokecolor{currentstroke}%
\pgfsetstrokeopacity{0.000000}%
\pgfsetdash{}{0pt}%
\pgfpathmoveto{\pgfqpoint{2.288393in}{1.199350in}}%
\pgfpathlineto{\pgfqpoint{2.288706in}{1.196382in}}%
\pgfpathlineto{\pgfqpoint{2.264298in}{1.193927in}}%
\pgfpathlineto{\pgfqpoint{2.262289in}{1.211080in}}%
\pgfpathlineto{\pgfqpoint{2.266331in}{1.211577in}}%
\pgfpathlineto{\pgfqpoint{2.265786in}{1.225833in}}%
\pgfpathlineto{\pgfqpoint{2.285375in}{1.228022in}}%
\pgfpathlineto{\pgfqpoint{2.288393in}{1.199350in}}%
\pgfpathclose%
\pgfusepath{fill}%
\end{pgfscope}%
\begin{pgfscope}%
\pgfpathrectangle{\pgfqpoint{0.100000in}{0.100000in}}{\pgfqpoint{3.007045in}{1.925000in}}%
\pgfusepath{clip}%
\pgfsetbuttcap%
\pgfsetmiterjoin%
\definecolor{currentfill}{rgb}{0.376471,0.730181,0.424298}%
\pgfsetfillcolor{currentfill}%
\pgfsetlinewidth{0.000000pt}%
\definecolor{currentstroke}{rgb}{0.000000,0.000000,0.000000}%
\pgfsetstrokecolor{currentstroke}%
\pgfsetstrokeopacity{0.000000}%
\pgfsetdash{}{0pt}%
\pgfpathmoveto{\pgfqpoint{1.905032in}{1.544903in}}%
\pgfpathlineto{\pgfqpoint{1.876967in}{1.544122in}}%
\pgfpathlineto{\pgfqpoint{1.876801in}{1.549513in}}%
\pgfpathlineto{\pgfqpoint{1.881663in}{1.556766in}}%
\pgfpathlineto{\pgfqpoint{1.881731in}{1.559757in}}%
\pgfpathlineto{\pgfqpoint{1.875816in}{1.567788in}}%
\pgfpathlineto{\pgfqpoint{1.870534in}{1.568228in}}%
\pgfpathlineto{\pgfqpoint{1.872272in}{1.578521in}}%
\pgfpathlineto{\pgfqpoint{1.874776in}{1.580802in}}%
\pgfpathlineto{\pgfqpoint{1.877917in}{1.589345in}}%
\pgfpathlineto{\pgfqpoint{1.881222in}{1.592179in}}%
\pgfpathlineto{\pgfqpoint{1.894425in}{1.598113in}}%
\pgfpathlineto{\pgfqpoint{1.896908in}{1.602038in}}%
\pgfpathlineto{\pgfqpoint{1.896759in}{1.607543in}}%
\pgfpathlineto{\pgfqpoint{1.907975in}{1.607918in}}%
\pgfpathlineto{\pgfqpoint{1.909099in}{1.596409in}}%
\pgfpathlineto{\pgfqpoint{1.909946in}{1.573607in}}%
\pgfpathlineto{\pgfqpoint{1.904215in}{1.573411in}}%
\pgfpathlineto{\pgfqpoint{1.905032in}{1.544903in}}%
\pgfpathclose%
\pgfusepath{fill}%
\end{pgfscope}%
\begin{pgfscope}%
\pgfpathrectangle{\pgfqpoint{0.100000in}{0.100000in}}{\pgfqpoint{3.007045in}{1.925000in}}%
\pgfusepath{clip}%
\pgfsetbuttcap%
\pgfsetmiterjoin%
\definecolor{currentfill}{rgb}{0.727797,0.891580,0.701730}%
\pgfsetfillcolor{currentfill}%
\pgfsetlinewidth{0.000000pt}%
\definecolor{currentstroke}{rgb}{0.000000,0.000000,0.000000}%
\pgfsetstrokecolor{currentstroke}%
\pgfsetstrokeopacity{0.000000}%
\pgfsetdash{}{0pt}%
\pgfpathmoveto{\pgfqpoint{2.052237in}{0.736873in}}%
\pgfpathlineto{\pgfqpoint{2.048959in}{0.728755in}}%
\pgfpathlineto{\pgfqpoint{2.040123in}{0.721921in}}%
\pgfpathlineto{\pgfqpoint{2.027190in}{0.716066in}}%
\pgfpathlineto{\pgfqpoint{2.020779in}{0.710545in}}%
\pgfpathlineto{\pgfqpoint{2.011147in}{0.717467in}}%
\pgfpathlineto{\pgfqpoint{2.011096in}{0.729886in}}%
\pgfpathlineto{\pgfqpoint{2.008261in}{0.731716in}}%
\pgfpathlineto{\pgfqpoint{2.013725in}{0.737732in}}%
\pgfpathlineto{\pgfqpoint{2.013096in}{0.749239in}}%
\pgfpathlineto{\pgfqpoint{2.009327in}{0.760562in}}%
\pgfpathlineto{\pgfqpoint{2.018223in}{0.761023in}}%
\pgfpathlineto{\pgfqpoint{2.018021in}{0.764899in}}%
\pgfpathlineto{\pgfqpoint{2.023674in}{0.765200in}}%
\pgfpathlineto{\pgfqpoint{2.023847in}{0.761087in}}%
\pgfpathlineto{\pgfqpoint{2.030481in}{0.763778in}}%
\pgfpathlineto{\pgfqpoint{2.031123in}{0.752696in}}%
\pgfpathlineto{\pgfqpoint{2.025625in}{0.749814in}}%
\pgfpathlineto{\pgfqpoint{2.031102in}{0.744404in}}%
\pgfpathlineto{\pgfqpoint{2.052237in}{0.736873in}}%
\pgfpathclose%
\pgfusepath{fill}%
\end{pgfscope}%
\begin{pgfscope}%
\pgfpathrectangle{\pgfqpoint{0.100000in}{0.100000in}}{\pgfqpoint{3.007045in}{1.925000in}}%
\pgfusepath{clip}%
\pgfsetbuttcap%
\pgfsetmiterjoin%
\definecolor{currentfill}{rgb}{0.363922,0.724029,0.418147}%
\pgfsetfillcolor{currentfill}%
\pgfsetlinewidth{0.000000pt}%
\definecolor{currentstroke}{rgb}{0.000000,0.000000,0.000000}%
\pgfsetstrokecolor{currentstroke}%
\pgfsetstrokeopacity{0.000000}%
\pgfsetdash{}{0pt}%
\pgfpathmoveto{\pgfqpoint{1.827033in}{1.605997in}}%
\pgfpathlineto{\pgfqpoint{1.827186in}{1.594478in}}%
\pgfpathlineto{\pgfqpoint{1.829097in}{1.594472in}}%
\pgfpathlineto{\pgfqpoint{1.829469in}{1.583832in}}%
\pgfpathlineto{\pgfqpoint{1.802241in}{1.583541in}}%
\pgfpathlineto{\pgfqpoint{1.806019in}{1.580465in}}%
\pgfpathlineto{\pgfqpoint{1.765864in}{1.580217in}}%
\pgfpathlineto{\pgfqpoint{1.765646in}{1.613993in}}%
\pgfpathlineto{\pgfqpoint{1.764876in}{1.642920in}}%
\pgfpathlineto{\pgfqpoint{1.764918in}{1.648716in}}%
\pgfpathlineto{\pgfqpoint{1.782023in}{1.648669in}}%
\pgfpathlineto{\pgfqpoint{1.782448in}{1.621431in}}%
\pgfpathlineto{\pgfqpoint{1.793675in}{1.615402in}}%
\pgfpathlineto{\pgfqpoint{1.798627in}{1.617392in}}%
\pgfpathlineto{\pgfqpoint{1.802129in}{1.614963in}}%
\pgfpathlineto{\pgfqpoint{1.801088in}{1.605771in}}%
\pgfpathlineto{\pgfqpoint{1.827033in}{1.605997in}}%
\pgfpathclose%
\pgfusepath{fill}%
\end{pgfscope}%
\begin{pgfscope}%
\pgfpathrectangle{\pgfqpoint{0.100000in}{0.100000in}}{\pgfqpoint{3.007045in}{1.925000in}}%
\pgfusepath{clip}%
\pgfsetbuttcap%
\pgfsetmiterjoin%
\definecolor{currentfill}{rgb}{0.590542,0.831926,0.574271}%
\pgfsetfillcolor{currentfill}%
\pgfsetlinewidth{0.000000pt}%
\definecolor{currentstroke}{rgb}{0.000000,0.000000,0.000000}%
\pgfsetstrokecolor{currentstroke}%
\pgfsetstrokeopacity{0.000000}%
\pgfsetdash{}{0pt}%
\pgfpathmoveto{\pgfqpoint{1.876321in}{0.768157in}}%
\pgfpathlineto{\pgfqpoint{1.876274in}{0.763349in}}%
\pgfpathlineto{\pgfqpoint{1.884045in}{0.763316in}}%
\pgfpathlineto{\pgfqpoint{1.884006in}{0.739484in}}%
\pgfpathlineto{\pgfqpoint{1.838686in}{0.738812in}}%
\pgfpathlineto{\pgfqpoint{1.835563in}{0.743537in}}%
\pgfpathlineto{\pgfqpoint{1.836073in}{0.747122in}}%
\pgfpathlineto{\pgfqpoint{1.841164in}{0.748497in}}%
\pgfpathlineto{\pgfqpoint{1.846121in}{0.759629in}}%
\pgfpathlineto{\pgfqpoint{1.842406in}{0.765728in}}%
\pgfpathlineto{\pgfqpoint{1.842812in}{0.769497in}}%
\pgfpathlineto{\pgfqpoint{1.856051in}{0.769311in}}%
\pgfpathlineto{\pgfqpoint{1.856034in}{0.767401in}}%
\pgfpathlineto{\pgfqpoint{1.870545in}{0.767215in}}%
\pgfpathlineto{\pgfqpoint{1.876321in}{0.768157in}}%
\pgfpathclose%
\pgfusepath{fill}%
\end{pgfscope}%
\begin{pgfscope}%
\pgfpathrectangle{\pgfqpoint{0.100000in}{0.100000in}}{\pgfqpoint{3.007045in}{1.925000in}}%
\pgfusepath{clip}%
\pgfsetbuttcap%
\pgfsetmiterjoin%
\definecolor{currentfill}{rgb}{0.723122,0.889612,0.697178}%
\pgfsetfillcolor{currentfill}%
\pgfsetlinewidth{0.000000pt}%
\definecolor{currentstroke}{rgb}{0.000000,0.000000,0.000000}%
\pgfsetstrokecolor{currentstroke}%
\pgfsetstrokeopacity{0.000000}%
\pgfsetdash{}{0pt}%
\pgfpathmoveto{\pgfqpoint{2.061841in}{1.396171in}}%
\pgfpathlineto{\pgfqpoint{2.044679in}{1.394919in}}%
\pgfpathlineto{\pgfqpoint{2.021967in}{1.394387in}}%
\pgfpathlineto{\pgfqpoint{2.022119in}{1.391513in}}%
\pgfpathlineto{\pgfqpoint{1.993663in}{1.389926in}}%
\pgfpathlineto{\pgfqpoint{1.992238in}{1.415509in}}%
\pgfpathlineto{\pgfqpoint{1.986647in}{1.415927in}}%
\pgfpathlineto{\pgfqpoint{1.980936in}{1.413027in}}%
\pgfpathlineto{\pgfqpoint{1.980060in}{1.429597in}}%
\pgfpathlineto{\pgfqpoint{1.979503in}{1.438168in}}%
\pgfpathlineto{\pgfqpoint{1.996750in}{1.439088in}}%
\pgfpathlineto{\pgfqpoint{1.996384in}{1.444829in}}%
\pgfpathlineto{\pgfqpoint{2.021528in}{1.446204in}}%
\pgfpathlineto{\pgfqpoint{2.030394in}{1.446824in}}%
\pgfpathlineto{\pgfqpoint{2.029155in}{1.469240in}}%
\pgfpathlineto{\pgfqpoint{2.049555in}{1.470541in}}%
\pgfpathlineto{\pgfqpoint{2.050014in}{1.463525in}}%
\pgfpathlineto{\pgfqpoint{2.046797in}{1.456021in}}%
\pgfpathlineto{\pgfqpoint{2.047314in}{1.447880in}}%
\pgfpathlineto{\pgfqpoint{2.058743in}{1.447959in}}%
\pgfpathlineto{\pgfqpoint{2.061841in}{1.396171in}}%
\pgfpathclose%
\pgfusepath{fill}%
\end{pgfscope}%
\begin{pgfscope}%
\pgfpathrectangle{\pgfqpoint{0.100000in}{0.100000in}}{\pgfqpoint{3.007045in}{1.925000in}}%
\pgfusepath{clip}%
\pgfsetbuttcap%
\pgfsetmiterjoin%
\definecolor{currentfill}{rgb}{0.485352,0.782837,0.487782}%
\pgfsetfillcolor{currentfill}%
\pgfsetlinewidth{0.000000pt}%
\definecolor{currentstroke}{rgb}{0.000000,0.000000,0.000000}%
\pgfsetstrokecolor{currentstroke}%
\pgfsetstrokeopacity{0.000000}%
\pgfsetdash{}{0pt}%
\pgfpathmoveto{\pgfqpoint{1.082392in}{1.585392in}}%
\pgfpathlineto{\pgfqpoint{1.114575in}{1.580480in}}%
\pgfpathlineto{\pgfqpoint{1.115456in}{1.579920in}}%
\pgfpathlineto{\pgfqpoint{1.154054in}{1.574061in}}%
\pgfpathlineto{\pgfqpoint{1.169718in}{1.571914in}}%
\pgfpathlineto{\pgfqpoint{1.170691in}{1.563235in}}%
\pgfpathlineto{\pgfqpoint{1.174953in}{1.552958in}}%
\pgfpathlineto{\pgfqpoint{1.180687in}{1.550302in}}%
\pgfpathlineto{\pgfqpoint{1.188802in}{1.543168in}}%
\pgfpathlineto{\pgfqpoint{1.193646in}{1.534813in}}%
\pgfpathlineto{\pgfqpoint{1.197692in}{1.530943in}}%
\pgfpathlineto{\pgfqpoint{1.199309in}{1.522550in}}%
\pgfpathlineto{\pgfqpoint{1.197829in}{1.512182in}}%
\pgfpathlineto{\pgfqpoint{1.131931in}{1.521748in}}%
\pgfpathlineto{\pgfqpoint{1.131047in}{1.515926in}}%
\pgfpathlineto{\pgfqpoint{1.119707in}{1.517659in}}%
\pgfpathlineto{\pgfqpoint{1.118845in}{1.511867in}}%
\pgfpathlineto{\pgfqpoint{1.113003in}{1.512751in}}%
\pgfpathlineto{\pgfqpoint{1.112286in}{1.507103in}}%
\pgfpathlineto{\pgfqpoint{1.103807in}{1.508400in}}%
\pgfpathlineto{\pgfqpoint{1.102940in}{1.502630in}}%
\pgfpathlineto{\pgfqpoint{1.092509in}{1.504079in}}%
\pgfpathlineto{\pgfqpoint{1.087640in}{1.512697in}}%
\pgfpathlineto{\pgfqpoint{1.083266in}{1.515471in}}%
\pgfpathlineto{\pgfqpoint{1.078184in}{1.513711in}}%
\pgfpathlineto{\pgfqpoint{1.076162in}{1.510378in}}%
\pgfpathlineto{\pgfqpoint{1.070042in}{1.506895in}}%
\pgfpathlineto{\pgfqpoint{1.067064in}{1.517262in}}%
\pgfpathlineto{\pgfqpoint{1.062108in}{1.518116in}}%
\pgfpathlineto{\pgfqpoint{1.059788in}{1.522455in}}%
\pgfpathlineto{\pgfqpoint{1.061138in}{1.530583in}}%
\pgfpathlineto{\pgfqpoint{1.058870in}{1.534775in}}%
\pgfpathlineto{\pgfqpoint{1.057765in}{1.542345in}}%
\pgfpathlineto{\pgfqpoint{1.053058in}{1.551890in}}%
\pgfpathlineto{\pgfqpoint{1.054900in}{1.557314in}}%
\pgfpathlineto{\pgfqpoint{1.051087in}{1.562499in}}%
\pgfpathlineto{\pgfqpoint{1.037496in}{1.564829in}}%
\pgfpathlineto{\pgfqpoint{1.038445in}{1.570273in}}%
\pgfpathlineto{\pgfqpoint{1.020483in}{1.573439in}}%
\pgfpathlineto{\pgfqpoint{1.024388in}{1.595276in}}%
\pgfpathlineto{\pgfqpoint{1.040422in}{1.591840in}}%
\pgfpathlineto{\pgfqpoint{1.082392in}{1.585392in}}%
\pgfpathclose%
\pgfusepath{fill}%
\end{pgfscope}%
\begin{pgfscope}%
\pgfpathrectangle{\pgfqpoint{0.100000in}{0.100000in}}{\pgfqpoint{3.007045in}{1.925000in}}%
\pgfusepath{clip}%
\pgfsetbuttcap%
\pgfsetmiterjoin%
\definecolor{currentfill}{rgb}{0.695071,0.877801,0.669865}%
\pgfsetfillcolor{currentfill}%
\pgfsetlinewidth{0.000000pt}%
\definecolor{currentstroke}{rgb}{0.000000,0.000000,0.000000}%
\pgfsetstrokecolor{currentstroke}%
\pgfsetstrokeopacity{0.000000}%
\pgfsetdash{}{0pt}%
\pgfpathmoveto{\pgfqpoint{2.143299in}{0.800190in}}%
\pgfpathlineto{\pgfqpoint{2.139121in}{0.799917in}}%
\pgfpathlineto{\pgfqpoint{2.136914in}{0.795797in}}%
\pgfpathlineto{\pgfqpoint{2.129211in}{0.792870in}}%
\pgfpathlineto{\pgfqpoint{2.117772in}{0.794063in}}%
\pgfpathlineto{\pgfqpoint{2.117152in}{0.802691in}}%
\pgfpathlineto{\pgfqpoint{2.105570in}{0.801993in}}%
\pgfpathlineto{\pgfqpoint{2.105804in}{0.798159in}}%
\pgfpathlineto{\pgfqpoint{2.091654in}{0.795102in}}%
\pgfpathlineto{\pgfqpoint{2.074344in}{0.794001in}}%
\pgfpathlineto{\pgfqpoint{2.073753in}{0.803617in}}%
\pgfpathlineto{\pgfqpoint{2.072446in}{0.823053in}}%
\pgfpathlineto{\pgfqpoint{2.086876in}{0.823888in}}%
\pgfpathlineto{\pgfqpoint{2.085917in}{0.838289in}}%
\pgfpathlineto{\pgfqpoint{2.103193in}{0.839402in}}%
\pgfpathlineto{\pgfqpoint{2.110699in}{0.840860in}}%
\pgfpathlineto{\pgfqpoint{2.110261in}{0.847530in}}%
\pgfpathlineto{\pgfqpoint{2.113075in}{0.848655in}}%
\pgfpathlineto{\pgfqpoint{2.123649in}{0.849414in}}%
\pgfpathlineto{\pgfqpoint{2.123839in}{0.846519in}}%
\pgfpathlineto{\pgfqpoint{2.144807in}{0.847920in}}%
\pgfpathlineto{\pgfqpoint{2.144570in}{0.838481in}}%
\pgfpathlineto{\pgfqpoint{2.143299in}{0.800190in}}%
\pgfpathclose%
\pgfusepath{fill}%
\end{pgfscope}%
\begin{pgfscope}%
\pgfpathrectangle{\pgfqpoint{0.100000in}{0.100000in}}{\pgfqpoint{3.007045in}{1.925000in}}%
\pgfusepath{clip}%
\pgfsetbuttcap%
\pgfsetmiterjoin%
\definecolor{currentfill}{rgb}{0.134025,0.542330,0.268281}%
\pgfsetfillcolor{currentfill}%
\pgfsetlinewidth{0.000000pt}%
\definecolor{currentstroke}{rgb}{0.000000,0.000000,0.000000}%
\pgfsetstrokecolor{currentstroke}%
\pgfsetstrokeopacity{0.000000}%
\pgfsetdash{}{0pt}%
\pgfpathmoveto{\pgfqpoint{2.898990in}{1.789167in}}%
\pgfpathlineto{\pgfqpoint{2.897996in}{1.797410in}}%
\pgfpathlineto{\pgfqpoint{2.917566in}{1.855793in}}%
\pgfpathlineto{\pgfqpoint{2.926613in}{1.855322in}}%
\pgfpathlineto{\pgfqpoint{2.928868in}{1.845060in}}%
\pgfpathlineto{\pgfqpoint{2.936844in}{1.841940in}}%
\pgfpathlineto{\pgfqpoint{2.948719in}{1.853007in}}%
\pgfpathlineto{\pgfqpoint{2.957139in}{1.855503in}}%
\pgfpathlineto{\pgfqpoint{2.956602in}{1.859933in}}%
\pgfpathlineto{\pgfqpoint{2.962545in}{1.861779in}}%
\pgfpathlineto{\pgfqpoint{2.977427in}{1.855560in}}%
\pgfpathlineto{\pgfqpoint{2.982323in}{1.850596in}}%
\pgfpathlineto{\pgfqpoint{2.987134in}{1.849352in}}%
\pgfpathlineto{\pgfqpoint{2.996362in}{1.819764in}}%
\pgfpathlineto{\pgfqpoint{3.009567in}{1.778140in}}%
\pgfpathlineto{\pgfqpoint{3.009737in}{1.774049in}}%
\pgfpathlineto{\pgfqpoint{3.013596in}{1.761178in}}%
\pgfpathlineto{\pgfqpoint{3.003662in}{1.755298in}}%
\pgfpathlineto{\pgfqpoint{2.988875in}{1.746178in}}%
\pgfpathlineto{\pgfqpoint{2.987728in}{1.746399in}}%
\pgfpathlineto{\pgfqpoint{2.981003in}{1.769545in}}%
\pgfpathlineto{\pgfqpoint{2.972249in}{1.797222in}}%
\pgfpathlineto{\pgfqpoint{2.955175in}{1.793081in}}%
\pgfpathlineto{\pgfqpoint{2.951724in}{1.804308in}}%
\pgfpathlineto{\pgfqpoint{2.912260in}{1.792925in}}%
\pgfpathlineto{\pgfqpoint{2.898990in}{1.789167in}}%
\pgfpathclose%
\pgfusepath{fill}%
\end{pgfscope}%
\begin{pgfscope}%
\pgfpathrectangle{\pgfqpoint{0.100000in}{0.100000in}}{\pgfqpoint{3.007045in}{1.925000in}}%
\pgfusepath{clip}%
\pgfsetbuttcap%
\pgfsetmiterjoin%
\definecolor{currentfill}{rgb}{0.535179,0.806090,0.528750}%
\pgfsetfillcolor{currentfill}%
\pgfsetlinewidth{0.000000pt}%
\definecolor{currentstroke}{rgb}{0.000000,0.000000,0.000000}%
\pgfsetstrokecolor{currentstroke}%
\pgfsetstrokeopacity{0.000000}%
\pgfsetdash{}{0pt}%
\pgfpathmoveto{\pgfqpoint{2.278491in}{1.292509in}}%
\pgfpathlineto{\pgfqpoint{2.273138in}{1.338797in}}%
\pgfpathlineto{\pgfqpoint{2.292953in}{1.341716in}}%
\pgfpathlineto{\pgfqpoint{2.295342in}{1.329269in}}%
\pgfpathlineto{\pgfqpoint{2.297909in}{1.323832in}}%
\pgfpathlineto{\pgfqpoint{2.303481in}{1.324508in}}%
\pgfpathlineto{\pgfqpoint{2.305531in}{1.307364in}}%
\pgfpathlineto{\pgfqpoint{2.299961in}{1.306677in}}%
\pgfpathlineto{\pgfqpoint{2.301283in}{1.295237in}}%
\pgfpathlineto{\pgfqpoint{2.278491in}{1.292509in}}%
\pgfpathclose%
\pgfusepath{fill}%
\end{pgfscope}%
\begin{pgfscope}%
\pgfpathrectangle{\pgfqpoint{0.100000in}{0.100000in}}{\pgfqpoint{3.007045in}{1.925000in}}%
\pgfusepath{clip}%
\pgfsetbuttcap%
\pgfsetmiterjoin%
\definecolor{currentfill}{rgb}{0.618224,0.844844,0.597032}%
\pgfsetfillcolor{currentfill}%
\pgfsetlinewidth{0.000000pt}%
\definecolor{currentstroke}{rgb}{0.000000,0.000000,0.000000}%
\pgfsetstrokecolor{currentstroke}%
\pgfsetstrokeopacity{0.000000}%
\pgfsetdash{}{0pt}%
\pgfpathmoveto{\pgfqpoint{1.102849in}{0.996413in}}%
\pgfpathlineto{\pgfqpoint{1.100725in}{0.982093in}}%
\pgfpathlineto{\pgfqpoint{1.025816in}{0.993300in}}%
\pgfpathlineto{\pgfqpoint{1.035972in}{1.058230in}}%
\pgfpathlineto{\pgfqpoint{1.070635in}{1.052989in}}%
\pgfpathlineto{\pgfqpoint{1.073504in}{1.056252in}}%
\pgfpathlineto{\pgfqpoint{1.077407in}{1.066910in}}%
\pgfpathlineto{\pgfqpoint{1.083485in}{1.074787in}}%
\pgfpathlineto{\pgfqpoint{1.088472in}{1.075836in}}%
\pgfpathlineto{\pgfqpoint{1.093042in}{1.080708in}}%
\pgfpathlineto{\pgfqpoint{1.094867in}{1.088434in}}%
\pgfpathlineto{\pgfqpoint{1.103247in}{1.096115in}}%
\pgfpathlineto{\pgfqpoint{1.104848in}{1.099688in}}%
\pgfpathlineto{\pgfqpoint{1.112363in}{1.107267in}}%
\pgfpathlineto{\pgfqpoint{1.121625in}{1.109932in}}%
\pgfpathlineto{\pgfqpoint{1.123300in}{1.108295in}}%
\pgfpathlineto{\pgfqpoint{1.124305in}{1.096464in}}%
\pgfpathlineto{\pgfqpoint{1.121195in}{1.073917in}}%
\pgfpathlineto{\pgfqpoint{1.139523in}{1.071438in}}%
\pgfpathlineto{\pgfqpoint{1.139253in}{1.069444in}}%
\pgfpathlineto{\pgfqpoint{1.162663in}{1.066636in}}%
\pgfpathlineto{\pgfqpoint{1.161246in}{1.055657in}}%
\pgfpathlineto{\pgfqpoint{1.169817in}{1.038957in}}%
\pgfpathlineto{\pgfqpoint{1.120573in}{1.045885in}}%
\pgfpathlineto{\pgfqpoint{1.118519in}{1.040536in}}%
\pgfpathlineto{\pgfqpoint{1.108498in}{1.034565in}}%
\pgfpathlineto{\pgfqpoint{1.102849in}{0.996413in}}%
\pgfpathclose%
\pgfusepath{fill}%
\end{pgfscope}%
\begin{pgfscope}%
\pgfpathrectangle{\pgfqpoint{0.100000in}{0.100000in}}{\pgfqpoint{3.007045in}{1.925000in}}%
\pgfusepath{clip}%
\pgfsetbuttcap%
\pgfsetmiterjoin%
\definecolor{currentfill}{rgb}{0.376471,0.730181,0.424298}%
\pgfsetfillcolor{currentfill}%
\pgfsetlinewidth{0.000000pt}%
\definecolor{currentstroke}{rgb}{0.000000,0.000000,0.000000}%
\pgfsetstrokecolor{currentstroke}%
\pgfsetstrokeopacity{0.000000}%
\pgfsetdash{}{0pt}%
\pgfpathmoveto{\pgfqpoint{2.033429in}{1.724806in}}%
\pgfpathlineto{\pgfqpoint{2.028764in}{1.727403in}}%
\pgfpathlineto{\pgfqpoint{2.031261in}{1.731523in}}%
\pgfpathlineto{\pgfqpoint{2.049089in}{1.741926in}}%
\pgfpathlineto{\pgfqpoint{2.055441in}{1.748010in}}%
\pgfpathlineto{\pgfqpoint{2.063105in}{1.751149in}}%
\pgfpathlineto{\pgfqpoint{2.058606in}{1.742131in}}%
\pgfpathlineto{\pgfqpoint{2.053691in}{1.738050in}}%
\pgfpathlineto{\pgfqpoint{2.042387in}{1.732402in}}%
\pgfpathlineto{\pgfqpoint{2.044192in}{1.729855in}}%
\pgfpathlineto{\pgfqpoint{2.033429in}{1.724806in}}%
\pgfpathclose%
\pgfusepath{fill}%
\end{pgfscope}%
\begin{pgfscope}%
\pgfpathrectangle{\pgfqpoint{0.100000in}{0.100000in}}{\pgfqpoint{3.007045in}{1.925000in}}%
\pgfusepath{clip}%
\pgfsetbuttcap%
\pgfsetmiterjoin%
\definecolor{currentfill}{rgb}{0.376471,0.730181,0.424298}%
\pgfsetfillcolor{currentfill}%
\pgfsetlinewidth{0.000000pt}%
\definecolor{currentstroke}{rgb}{0.000000,0.000000,0.000000}%
\pgfsetstrokecolor{currentstroke}%
\pgfsetstrokeopacity{0.000000}%
\pgfsetdash{}{0pt}%
\pgfpathmoveto{\pgfqpoint{2.088047in}{1.667928in}}%
\pgfpathlineto{\pgfqpoint{2.088768in}{1.658202in}}%
\pgfpathlineto{\pgfqpoint{2.091627in}{1.652681in}}%
\pgfpathlineto{\pgfqpoint{2.085981in}{1.652184in}}%
\pgfpathlineto{\pgfqpoint{2.087227in}{1.635040in}}%
\pgfpathlineto{\pgfqpoint{2.047301in}{1.632207in}}%
\pgfpathlineto{\pgfqpoint{2.046234in}{1.649481in}}%
\pgfpathlineto{\pgfqpoint{2.051916in}{1.649810in}}%
\pgfpathlineto{\pgfqpoint{2.048453in}{1.655274in}}%
\pgfpathlineto{\pgfqpoint{2.047218in}{1.673072in}}%
\pgfpathlineto{\pgfqpoint{2.048913in}{1.677724in}}%
\pgfpathlineto{\pgfqpoint{2.058859in}{1.686684in}}%
\pgfpathlineto{\pgfqpoint{2.062270in}{1.688002in}}%
\pgfpathlineto{\pgfqpoint{2.068764in}{1.697183in}}%
\pgfpathlineto{\pgfqpoint{2.077462in}{1.703060in}}%
\pgfpathlineto{\pgfqpoint{2.089417in}{1.705852in}}%
\pgfpathlineto{\pgfqpoint{2.096001in}{1.705831in}}%
\pgfpathlineto{\pgfqpoint{2.100493in}{1.701354in}}%
\pgfpathlineto{\pgfqpoint{2.091665in}{1.700318in}}%
\pgfpathlineto{\pgfqpoint{2.090480in}{1.696264in}}%
\pgfpathlineto{\pgfqpoint{2.085473in}{1.693327in}}%
\pgfpathlineto{\pgfqpoint{2.078290in}{1.686386in}}%
\pgfpathlineto{\pgfqpoint{2.078268in}{1.682642in}}%
\pgfpathlineto{\pgfqpoint{2.069709in}{1.670334in}}%
\pgfpathlineto{\pgfqpoint{2.068597in}{1.661571in}}%
\pgfpathlineto{\pgfqpoint{2.071052in}{1.658913in}}%
\pgfpathlineto{\pgfqpoint{2.078864in}{1.668386in}}%
\pgfpathlineto{\pgfqpoint{2.082700in}{1.671038in}}%
\pgfpathlineto{\pgfqpoint{2.088047in}{1.667928in}}%
\pgfpathclose%
\pgfusepath{fill}%
\end{pgfscope}%
\begin{pgfscope}%
\pgfpathrectangle{\pgfqpoint{0.100000in}{0.100000in}}{\pgfqpoint{3.007045in}{1.925000in}}%
\pgfusepath{clip}%
\pgfsetbuttcap%
\pgfsetmiterjoin%
\definecolor{currentfill}{rgb}{0.769873,0.909296,0.742699}%
\pgfsetfillcolor{currentfill}%
\pgfsetlinewidth{0.000000pt}%
\definecolor{currentstroke}{rgb}{0.000000,0.000000,0.000000}%
\pgfsetstrokecolor{currentstroke}%
\pgfsetstrokeopacity{0.000000}%
\pgfsetdash{}{0pt}%
\pgfpathmoveto{\pgfqpoint{1.251074in}{1.090450in}}%
\pgfpathlineto{\pgfqpoint{1.245702in}{1.097910in}}%
\pgfpathlineto{\pgfqpoint{1.241682in}{1.094643in}}%
\pgfpathlineto{\pgfqpoint{1.232103in}{1.090942in}}%
\pgfpathlineto{\pgfqpoint{1.229068in}{1.091653in}}%
\pgfpathlineto{\pgfqpoint{1.223896in}{1.096979in}}%
\pgfpathlineto{\pgfqpoint{1.220511in}{1.109227in}}%
\pgfpathlineto{\pgfqpoint{1.205700in}{1.130953in}}%
\pgfpathlineto{\pgfqpoint{1.213054in}{1.137702in}}%
\pgfpathlineto{\pgfqpoint{1.213461in}{1.141563in}}%
\pgfpathlineto{\pgfqpoint{1.209760in}{1.146848in}}%
\pgfpathlineto{\pgfqpoint{1.242494in}{1.143237in}}%
\pgfpathlineto{\pgfqpoint{1.245658in}{1.171547in}}%
\pgfpathlineto{\pgfqpoint{1.310620in}{1.164502in}}%
\pgfpathlineto{\pgfqpoint{1.308723in}{1.147516in}}%
\pgfpathlineto{\pgfqpoint{1.306529in}{1.124802in}}%
\pgfpathlineto{\pgfqpoint{1.261017in}{1.129389in}}%
\pgfpathlineto{\pgfqpoint{1.259199in}{1.112273in}}%
\pgfpathlineto{\pgfqpoint{1.253600in}{1.112871in}}%
\pgfpathlineto{\pgfqpoint{1.251074in}{1.090450in}}%
\pgfpathclose%
\pgfusepath{fill}%
\end{pgfscope}%
\begin{pgfscope}%
\pgfpathrectangle{\pgfqpoint{0.100000in}{0.100000in}}{\pgfqpoint{3.007045in}{1.925000in}}%
\pgfusepath{clip}%
\pgfsetbuttcap%
\pgfsetmiterjoin%
\definecolor{currentfill}{rgb}{0.370196,0.727105,0.421223}%
\pgfsetfillcolor{currentfill}%
\pgfsetlinewidth{0.000000pt}%
\definecolor{currentstroke}{rgb}{0.000000,0.000000,0.000000}%
\pgfsetstrokecolor{currentstroke}%
\pgfsetstrokeopacity{0.000000}%
\pgfsetdash{}{0pt}%
\pgfpathmoveto{\pgfqpoint{2.498296in}{1.267386in}}%
\pgfpathlineto{\pgfqpoint{2.480209in}{1.264505in}}%
\pgfpathlineto{\pgfqpoint{2.478385in}{1.281786in}}%
\pgfpathlineto{\pgfqpoint{2.474602in}{1.281312in}}%
\pgfpathlineto{\pgfqpoint{2.474283in}{1.290040in}}%
\pgfpathlineto{\pgfqpoint{2.476757in}{1.293301in}}%
\pgfpathlineto{\pgfqpoint{2.486521in}{1.293715in}}%
\pgfpathlineto{\pgfqpoint{2.493292in}{1.298554in}}%
\pgfpathlineto{\pgfqpoint{2.498296in}{1.267386in}}%
\pgfpathclose%
\pgfusepath{fill}%
\end{pgfscope}%
\begin{pgfscope}%
\pgfpathrectangle{\pgfqpoint{0.100000in}{0.100000in}}{\pgfqpoint{3.007045in}{1.925000in}}%
\pgfusepath{clip}%
\pgfsetbuttcap%
\pgfsetmiterjoin%
\definecolor{currentfill}{rgb}{0.474279,0.777670,0.478677}%
\pgfsetfillcolor{currentfill}%
\pgfsetlinewidth{0.000000pt}%
\definecolor{currentstroke}{rgb}{0.000000,0.000000,0.000000}%
\pgfsetstrokecolor{currentstroke}%
\pgfsetstrokeopacity{0.000000}%
\pgfsetdash{}{0pt}%
\pgfpathmoveto{\pgfqpoint{2.161409in}{0.940303in}}%
\pgfpathlineto{\pgfqpoint{2.160890in}{0.928692in}}%
\pgfpathlineto{\pgfqpoint{2.166375in}{0.924298in}}%
\pgfpathlineto{\pgfqpoint{2.163243in}{0.916924in}}%
\pgfpathlineto{\pgfqpoint{2.149635in}{0.914318in}}%
\pgfpathlineto{\pgfqpoint{2.147248in}{0.909640in}}%
\pgfpathlineto{\pgfqpoint{2.142348in}{0.908782in}}%
\pgfpathlineto{\pgfqpoint{2.129095in}{0.910042in}}%
\pgfpathlineto{\pgfqpoint{2.125902in}{0.914003in}}%
\pgfpathlineto{\pgfqpoint{2.120897in}{0.915606in}}%
\pgfpathlineto{\pgfqpoint{2.114744in}{0.920210in}}%
\pgfpathlineto{\pgfqpoint{2.114171in}{0.933474in}}%
\pgfpathlineto{\pgfqpoint{2.116000in}{0.935950in}}%
\pgfpathlineto{\pgfqpoint{2.136863in}{0.936941in}}%
\pgfpathlineto{\pgfqpoint{2.143491in}{0.939640in}}%
\pgfpathlineto{\pgfqpoint{2.144945in}{0.937174in}}%
\pgfpathlineto{\pgfqpoint{2.151093in}{0.939458in}}%
\pgfpathlineto{\pgfqpoint{2.161409in}{0.940303in}}%
\pgfpathclose%
\pgfusepath{fill}%
\end{pgfscope}%
\begin{pgfscope}%
\pgfpathrectangle{\pgfqpoint{0.100000in}{0.100000in}}{\pgfqpoint{3.007045in}{1.925000in}}%
\pgfusepath{clip}%
\pgfsetbuttcap%
\pgfsetmiterjoin%
\definecolor{currentfill}{rgb}{0.732472,0.893549,0.706282}%
\pgfsetfillcolor{currentfill}%
\pgfsetlinewidth{0.000000pt}%
\definecolor{currentstroke}{rgb}{0.000000,0.000000,0.000000}%
\pgfsetstrokecolor{currentstroke}%
\pgfsetstrokeopacity{0.000000}%
\pgfsetdash{}{0pt}%
\pgfpathmoveto{\pgfqpoint{1.282225in}{0.694354in}}%
\pgfpathlineto{\pgfqpoint{1.265329in}{0.695996in}}%
\pgfpathlineto{\pgfqpoint{1.219256in}{0.700832in}}%
\pgfpathlineto{\pgfqpoint{1.222778in}{0.735139in}}%
\pgfpathlineto{\pgfqpoint{1.194858in}{0.738140in}}%
\pgfpathlineto{\pgfqpoint{1.198272in}{0.767309in}}%
\pgfpathlineto{\pgfqpoint{1.200282in}{0.767082in}}%
\pgfpathlineto{\pgfqpoint{1.201518in}{0.778286in}}%
\pgfpathlineto{\pgfqpoint{1.224219in}{0.776176in}}%
\pgfpathlineto{\pgfqpoint{1.226820in}{0.793144in}}%
\pgfpathlineto{\pgfqpoint{1.231742in}{0.838561in}}%
\pgfpathlineto{\pgfqpoint{1.237556in}{0.837894in}}%
\pgfpathlineto{\pgfqpoint{1.236948in}{0.832126in}}%
\pgfpathlineto{\pgfqpoint{1.271248in}{0.828230in}}%
\pgfpathlineto{\pgfqpoint{1.271855in}{0.834010in}}%
\pgfpathlineto{\pgfqpoint{1.283287in}{0.832832in}}%
\pgfpathlineto{\pgfqpoint{1.288966in}{0.832202in}}%
\pgfpathlineto{\pgfqpoint{1.287264in}{0.814931in}}%
\pgfpathlineto{\pgfqpoint{1.294107in}{0.814272in}}%
\pgfpathlineto{\pgfqpoint{1.292993in}{0.803535in}}%
\pgfpathlineto{\pgfqpoint{1.304431in}{0.802517in}}%
\pgfpathlineto{\pgfqpoint{1.303863in}{0.796794in}}%
\pgfpathlineto{\pgfqpoint{1.292309in}{0.797855in}}%
\pgfpathlineto{\pgfqpoint{1.291172in}{0.786411in}}%
\pgfpathlineto{\pgfqpoint{1.288691in}{0.786623in}}%
\pgfpathlineto{\pgfqpoint{1.285953in}{0.758267in}}%
\pgfpathlineto{\pgfqpoint{1.283283in}{0.758525in}}%
\pgfpathlineto{\pgfqpoint{1.280477in}{0.729323in}}%
\pgfpathlineto{\pgfqpoint{1.285541in}{0.728841in}}%
\pgfpathlineto{\pgfqpoint{1.282225in}{0.694354in}}%
\pgfpathclose%
\pgfusepath{fill}%
\end{pgfscope}%
\begin{pgfscope}%
\pgfpathrectangle{\pgfqpoint{0.100000in}{0.100000in}}{\pgfqpoint{3.007045in}{1.925000in}}%
\pgfusepath{clip}%
\pgfsetbuttcap%
\pgfsetmiterjoin%
\definecolor{currentfill}{rgb}{0.357647,0.720953,0.415071}%
\pgfsetfillcolor{currentfill}%
\pgfsetlinewidth{0.000000pt}%
\definecolor{currentstroke}{rgb}{0.000000,0.000000,0.000000}%
\pgfsetstrokecolor{currentstroke}%
\pgfsetstrokeopacity{0.000000}%
\pgfsetdash{}{0pt}%
\pgfpathmoveto{\pgfqpoint{1.193280in}{1.837329in}}%
\pgfpathlineto{\pgfqpoint{1.189784in}{1.820439in}}%
\pgfpathlineto{\pgfqpoint{1.187177in}{1.820830in}}%
\pgfpathlineto{\pgfqpoint{1.185494in}{1.809331in}}%
\pgfpathlineto{\pgfqpoint{1.182880in}{1.799848in}}%
\pgfpathlineto{\pgfqpoint{1.180237in}{1.801764in}}%
\pgfpathlineto{\pgfqpoint{1.177030in}{1.787023in}}%
\pgfpathlineto{\pgfqpoint{1.174638in}{1.770993in}}%
\pgfpathlineto{\pgfqpoint{1.167096in}{1.773103in}}%
\pgfpathlineto{\pgfqpoint{1.165616in}{1.768587in}}%
\pgfpathlineto{\pgfqpoint{1.154039in}{1.770483in}}%
\pgfpathlineto{\pgfqpoint{1.152002in}{1.758091in}}%
\pgfpathlineto{\pgfqpoint{1.148828in}{1.763186in}}%
\pgfpathlineto{\pgfqpoint{1.141833in}{1.762846in}}%
\pgfpathlineto{\pgfqpoint{1.135562in}{1.764609in}}%
\pgfpathlineto{\pgfqpoint{1.126744in}{1.761331in}}%
\pgfpathlineto{\pgfqpoint{1.123379in}{1.762908in}}%
\pgfpathlineto{\pgfqpoint{1.127771in}{1.788966in}}%
\pgfpathlineto{\pgfqpoint{1.119270in}{1.790180in}}%
\pgfpathlineto{\pgfqpoint{1.120232in}{1.796028in}}%
\pgfpathlineto{\pgfqpoint{1.115258in}{1.796858in}}%
\pgfpathlineto{\pgfqpoint{1.115936in}{1.802454in}}%
\pgfpathlineto{\pgfqpoint{1.081806in}{1.808264in}}%
\pgfpathlineto{\pgfqpoint{1.080804in}{1.802553in}}%
\pgfpathlineto{\pgfqpoint{1.071976in}{1.804099in}}%
\pgfpathlineto{\pgfqpoint{1.070974in}{1.798414in}}%
\pgfpathlineto{\pgfqpoint{1.045450in}{1.802919in}}%
\pgfpathlineto{\pgfqpoint{1.046502in}{1.808652in}}%
\pgfpathlineto{\pgfqpoint{1.052373in}{1.807583in}}%
\pgfpathlineto{\pgfqpoint{1.061807in}{1.858712in}}%
\pgfpathlineto{\pgfqpoint{1.107359in}{1.850859in}}%
\pgfpathlineto{\pgfqpoint{1.160277in}{1.842278in}}%
\pgfpathlineto{\pgfqpoint{1.193280in}{1.837329in}}%
\pgfpathclose%
\pgfusepath{fill}%
\end{pgfscope}%
\begin{pgfscope}%
\pgfpathrectangle{\pgfqpoint{0.100000in}{0.100000in}}{\pgfqpoint{3.007045in}{1.925000in}}%
\pgfusepath{clip}%
\pgfsetbuttcap%
\pgfsetmiterjoin%
\definecolor{currentfill}{rgb}{0.208304,0.620884,0.327428}%
\pgfsetfillcolor{currentfill}%
\pgfsetlinewidth{0.000000pt}%
\definecolor{currentstroke}{rgb}{0.000000,0.000000,0.000000}%
\pgfsetstrokecolor{currentstroke}%
\pgfsetstrokeopacity{0.000000}%
\pgfsetdash{}{0pt}%
\pgfpathmoveto{\pgfqpoint{2.755737in}{0.991248in}}%
\pgfpathlineto{\pgfqpoint{2.755915in}{0.984607in}}%
\pgfpathlineto{\pgfqpoint{2.744917in}{0.981075in}}%
\pgfpathlineto{\pgfqpoint{2.737785in}{0.985361in}}%
\pgfpathlineto{\pgfqpoint{2.724069in}{0.990555in}}%
\pgfpathlineto{\pgfqpoint{2.723536in}{0.997192in}}%
\pgfpathlineto{\pgfqpoint{2.725670in}{0.999941in}}%
\pgfpathlineto{\pgfqpoint{2.722156in}{1.004752in}}%
\pgfpathlineto{\pgfqpoint{2.720462in}{1.011491in}}%
\pgfpathlineto{\pgfqpoint{2.710520in}{1.014800in}}%
\pgfpathlineto{\pgfqpoint{2.709850in}{1.020460in}}%
\pgfpathlineto{\pgfqpoint{2.705428in}{1.026285in}}%
\pgfpathlineto{\pgfqpoint{2.709028in}{1.030469in}}%
\pgfpathlineto{\pgfqpoint{2.715659in}{1.028449in}}%
\pgfpathlineto{\pgfqpoint{2.716657in}{1.024638in}}%
\pgfpathlineto{\pgfqpoint{2.722671in}{1.024298in}}%
\pgfpathlineto{\pgfqpoint{2.728149in}{1.020939in}}%
\pgfpathlineto{\pgfqpoint{2.731963in}{1.024371in}}%
\pgfpathlineto{\pgfqpoint{2.735129in}{1.019199in}}%
\pgfpathlineto{\pgfqpoint{2.738228in}{1.024682in}}%
\pgfpathlineto{\pgfqpoint{2.741385in}{1.024328in}}%
\pgfpathlineto{\pgfqpoint{2.743345in}{1.029859in}}%
\pgfpathlineto{\pgfqpoint{2.751811in}{1.031450in}}%
\pgfpathlineto{\pgfqpoint{2.758023in}{1.035375in}}%
\pgfpathlineto{\pgfqpoint{2.761386in}{1.033527in}}%
\pgfpathlineto{\pgfqpoint{2.767606in}{1.037752in}}%
\pgfpathlineto{\pgfqpoint{2.775648in}{1.039823in}}%
\pgfpathlineto{\pgfqpoint{2.782087in}{1.034841in}}%
\pgfpathlineto{\pgfqpoint{2.784638in}{1.040494in}}%
\pgfpathlineto{\pgfqpoint{2.789579in}{1.040321in}}%
\pgfpathlineto{\pgfqpoint{2.793501in}{1.036335in}}%
\pgfpathlineto{\pgfqpoint{2.797982in}{1.024799in}}%
\pgfpathlineto{\pgfqpoint{2.796768in}{1.016487in}}%
\pgfpathlineto{\pgfqpoint{2.790807in}{1.014831in}}%
\pgfpathlineto{\pgfqpoint{2.786449in}{1.006708in}}%
\pgfpathlineto{\pgfqpoint{2.785854in}{1.002606in}}%
\pgfpathlineto{\pgfqpoint{2.779981in}{0.995686in}}%
\pgfpathlineto{\pgfqpoint{2.772764in}{0.995498in}}%
\pgfpathlineto{\pgfqpoint{2.764411in}{0.997827in}}%
\pgfpathlineto{\pgfqpoint{2.762533in}{0.994980in}}%
\pgfpathlineto{\pgfqpoint{2.754523in}{0.995316in}}%
\pgfpathlineto{\pgfqpoint{2.755737in}{0.991248in}}%
\pgfpathclose%
\pgfusepath{fill}%
\end{pgfscope}%
\begin{pgfscope}%
\pgfpathrectangle{\pgfqpoint{0.100000in}{0.100000in}}{\pgfqpoint{3.007045in}{1.925000in}}%
\pgfusepath{clip}%
\pgfsetbuttcap%
\pgfsetmiterjoin%
\definecolor{currentfill}{rgb}{0.208304,0.620884,0.327428}%
\pgfsetfillcolor{currentfill}%
\pgfsetlinewidth{0.000000pt}%
\definecolor{currentstroke}{rgb}{0.000000,0.000000,0.000000}%
\pgfsetstrokecolor{currentstroke}%
\pgfsetstrokeopacity{0.000000}%
\pgfsetdash{}{0pt}%
\pgfpathmoveto{\pgfqpoint{2.787444in}{1.058403in}}%
\pgfpathlineto{\pgfqpoint{2.798434in}{1.043326in}}%
\pgfpathlineto{\pgfqpoint{2.793014in}{1.044103in}}%
\pgfpathlineto{\pgfqpoint{2.786756in}{1.058061in}}%
\pgfpathlineto{\pgfqpoint{2.787444in}{1.058403in}}%
\pgfpathclose%
\pgfusepath{fill}%
\end{pgfscope}%
\begin{pgfscope}%
\pgfpathrectangle{\pgfqpoint{0.100000in}{0.100000in}}{\pgfqpoint{3.007045in}{1.925000in}}%
\pgfusepath{clip}%
\pgfsetbuttcap%
\pgfsetmiterjoin%
\definecolor{currentfill}{rgb}{0.557324,0.816424,0.546959}%
\pgfsetfillcolor{currentfill}%
\pgfsetlinewidth{0.000000pt}%
\definecolor{currentstroke}{rgb}{0.000000,0.000000,0.000000}%
\pgfsetstrokecolor{currentstroke}%
\pgfsetstrokeopacity{0.000000}%
\pgfsetdash{}{0pt}%
\pgfpathmoveto{\pgfqpoint{1.619645in}{0.516644in}}%
\pgfpathlineto{\pgfqpoint{1.621659in}{0.515303in}}%
\pgfpathlineto{\pgfqpoint{1.638148in}{0.525132in}}%
\pgfpathlineto{\pgfqpoint{1.648009in}{0.514439in}}%
\pgfpathlineto{\pgfqpoint{1.644405in}{0.511356in}}%
\pgfpathlineto{\pgfqpoint{1.641992in}{0.498317in}}%
\pgfpathlineto{\pgfqpoint{1.619767in}{0.480550in}}%
\pgfpathlineto{\pgfqpoint{1.613261in}{0.488228in}}%
\pgfpathlineto{\pgfqpoint{1.605940in}{0.497218in}}%
\pgfpathlineto{\pgfqpoint{1.619645in}{0.516644in}}%
\pgfpathclose%
\pgfusepath{fill}%
\end{pgfscope}%
\begin{pgfscope}%
\pgfpathrectangle{\pgfqpoint{0.100000in}{0.100000in}}{\pgfqpoint{3.007045in}{1.925000in}}%
\pgfusepath{clip}%
\pgfsetbuttcap%
\pgfsetmiterjoin%
\definecolor{currentfill}{rgb}{0.634295,0.852211,0.610688}%
\pgfsetfillcolor{currentfill}%
\pgfsetlinewidth{0.000000pt}%
\definecolor{currentstroke}{rgb}{0.000000,0.000000,0.000000}%
\pgfsetstrokecolor{currentstroke}%
\pgfsetstrokeopacity{0.000000}%
\pgfsetdash{}{0pt}%
\pgfpathmoveto{\pgfqpoint{1.590530in}{1.433489in}}%
\pgfpathlineto{\pgfqpoint{1.647409in}{1.431503in}}%
\pgfpathlineto{\pgfqpoint{1.657257in}{1.431226in}}%
\pgfpathlineto{\pgfqpoint{1.656789in}{1.409421in}}%
\pgfpathlineto{\pgfqpoint{1.624011in}{1.410336in}}%
\pgfpathlineto{\pgfqpoint{1.622340in}{1.404124in}}%
\pgfpathlineto{\pgfqpoint{1.622963in}{1.397529in}}%
\pgfpathlineto{\pgfqpoint{1.621472in}{1.391045in}}%
\pgfpathlineto{\pgfqpoint{1.618866in}{1.388633in}}%
\pgfpathlineto{\pgfqpoint{1.612271in}{1.391873in}}%
\pgfpathlineto{\pgfqpoint{1.604815in}{1.396442in}}%
\pgfpathlineto{\pgfqpoint{1.603438in}{1.399865in}}%
\pgfpathlineto{\pgfqpoint{1.602167in}{1.402735in}}%
\pgfpathlineto{\pgfqpoint{1.594388in}{1.405377in}}%
\pgfpathlineto{\pgfqpoint{1.590580in}{1.410117in}}%
\pgfpathlineto{\pgfqpoint{1.586747in}{1.410853in}}%
\pgfpathlineto{\pgfqpoint{1.584652in}{1.417090in}}%
\pgfpathlineto{\pgfqpoint{1.576855in}{1.422263in}}%
\pgfpathlineto{\pgfqpoint{1.572614in}{1.429955in}}%
\pgfpathlineto{\pgfqpoint{1.568472in}{1.430306in}}%
\pgfpathlineto{\pgfqpoint{1.566436in}{1.434578in}}%
\pgfpathlineto{\pgfqpoint{1.590530in}{1.433489in}}%
\pgfpathclose%
\pgfusepath{fill}%
\end{pgfscope}%
\begin{pgfscope}%
\pgfpathrectangle{\pgfqpoint{0.100000in}{0.100000in}}{\pgfqpoint{3.007045in}{1.925000in}}%
\pgfusepath{clip}%
\pgfsetbuttcap%
\pgfsetmiterjoin%
\definecolor{currentfill}{rgb}{0.662345,0.864022,0.638001}%
\pgfsetfillcolor{currentfill}%
\pgfsetlinewidth{0.000000pt}%
\definecolor{currentstroke}{rgb}{0.000000,0.000000,0.000000}%
\pgfsetstrokecolor{currentstroke}%
\pgfsetstrokeopacity{0.000000}%
\pgfsetdash{}{0pt}%
\pgfpathmoveto{\pgfqpoint{2.256834in}{1.132262in}}%
\pgfpathlineto{\pgfqpoint{2.245396in}{1.130991in}}%
\pgfpathlineto{\pgfqpoint{2.240875in}{1.128586in}}%
\pgfpathlineto{\pgfqpoint{2.239682in}{1.138991in}}%
\pgfpathlineto{\pgfqpoint{2.235786in}{1.140573in}}%
\pgfpathlineto{\pgfqpoint{2.219718in}{1.138972in}}%
\pgfpathlineto{\pgfqpoint{2.218022in}{1.153845in}}%
\pgfpathlineto{\pgfqpoint{2.215737in}{1.157449in}}%
\pgfpathlineto{\pgfqpoint{2.227838in}{1.158721in}}%
\pgfpathlineto{\pgfqpoint{2.225671in}{1.177961in}}%
\pgfpathlineto{\pgfqpoint{2.245949in}{1.180376in}}%
\pgfpathlineto{\pgfqpoint{2.247318in}{1.166002in}}%
\pgfpathlineto{\pgfqpoint{2.253429in}{1.166775in}}%
\pgfpathlineto{\pgfqpoint{2.259410in}{1.171566in}}%
\pgfpathlineto{\pgfqpoint{2.261118in}{1.153025in}}%
\pgfpathlineto{\pgfqpoint{2.253275in}{1.146329in}}%
\pgfpathlineto{\pgfqpoint{2.249568in}{1.145924in}}%
\pgfpathlineto{\pgfqpoint{2.250151in}{1.140216in}}%
\pgfpathlineto{\pgfqpoint{2.256292in}{1.137004in}}%
\pgfpathlineto{\pgfqpoint{2.256834in}{1.132262in}}%
\pgfpathclose%
\pgfusepath{fill}%
\end{pgfscope}%
\begin{pgfscope}%
\pgfpathrectangle{\pgfqpoint{0.100000in}{0.100000in}}{\pgfqpoint{3.007045in}{1.925000in}}%
\pgfusepath{clip}%
\pgfsetbuttcap%
\pgfsetmiterjoin%
\definecolor{currentfill}{rgb}{0.507497,0.793172,0.505990}%
\pgfsetfillcolor{currentfill}%
\pgfsetlinewidth{0.000000pt}%
\definecolor{currentstroke}{rgb}{0.000000,0.000000,0.000000}%
\pgfsetstrokecolor{currentstroke}%
\pgfsetstrokeopacity{0.000000}%
\pgfsetdash{}{0pt}%
\pgfpathmoveto{\pgfqpoint{1.719649in}{1.001488in}}%
\pgfpathlineto{\pgfqpoint{1.659290in}{1.002553in}}%
\pgfpathlineto{\pgfqpoint{1.659741in}{1.033948in}}%
\pgfpathlineto{\pgfqpoint{1.692552in}{1.033341in}}%
\pgfpathlineto{\pgfqpoint{1.693377in}{1.073520in}}%
\pgfpathlineto{\pgfqpoint{1.701914in}{1.073348in}}%
\pgfpathlineto{\pgfqpoint{1.702042in}{1.079070in}}%
\pgfpathlineto{\pgfqpoint{1.722658in}{1.078703in}}%
\pgfpathlineto{\pgfqpoint{1.722556in}{1.070090in}}%
\pgfpathlineto{\pgfqpoint{1.722263in}{1.049928in}}%
\pgfpathlineto{\pgfqpoint{1.721571in}{1.001483in}}%
\pgfpathlineto{\pgfqpoint{1.719649in}{1.001488in}}%
\pgfpathclose%
\pgfusepath{fill}%
\end{pgfscope}%
\begin{pgfscope}%
\pgfpathrectangle{\pgfqpoint{0.100000in}{0.100000in}}{\pgfqpoint{3.007045in}{1.925000in}}%
\pgfusepath{clip}%
\pgfsetbuttcap%
\pgfsetmiterjoin%
\definecolor{currentfill}{rgb}{0.351373,0.717878,0.411995}%
\pgfsetfillcolor{currentfill}%
\pgfsetlinewidth{0.000000pt}%
\definecolor{currentstroke}{rgb}{0.000000,0.000000,0.000000}%
\pgfsetstrokecolor{currentstroke}%
\pgfsetstrokeopacity{0.000000}%
\pgfsetdash{}{0pt}%
\pgfpathmoveto{\pgfqpoint{2.406434in}{1.016470in}}%
\pgfpathlineto{\pgfqpoint{2.372965in}{1.012422in}}%
\pgfpathlineto{\pgfqpoint{2.373385in}{1.015073in}}%
\pgfpathlineto{\pgfqpoint{2.380873in}{1.017870in}}%
\pgfpathlineto{\pgfqpoint{2.397238in}{1.025171in}}%
\pgfpathlineto{\pgfqpoint{2.399471in}{1.032629in}}%
\pgfpathlineto{\pgfqpoint{2.409387in}{1.036508in}}%
\pgfpathlineto{\pgfqpoint{2.409125in}{1.041720in}}%
\pgfpathlineto{\pgfqpoint{2.415182in}{1.047525in}}%
\pgfpathlineto{\pgfqpoint{2.415558in}{1.052709in}}%
\pgfpathlineto{\pgfqpoint{2.423037in}{1.058784in}}%
\pgfpathlineto{\pgfqpoint{2.435154in}{1.067105in}}%
\pgfpathlineto{\pgfqpoint{2.440019in}{1.062167in}}%
\pgfpathlineto{\pgfqpoint{2.446669in}{1.052394in}}%
\pgfpathlineto{\pgfqpoint{2.437440in}{1.045801in}}%
\pgfpathlineto{\pgfqpoint{2.439575in}{1.041742in}}%
\pgfpathlineto{\pgfqpoint{2.434200in}{1.039055in}}%
\pgfpathlineto{\pgfqpoint{2.422280in}{1.037551in}}%
\pgfpathlineto{\pgfqpoint{2.416735in}{1.033753in}}%
\pgfpathlineto{\pgfqpoint{2.413702in}{1.026726in}}%
\pgfpathlineto{\pgfqpoint{2.407690in}{1.022099in}}%
\pgfpathlineto{\pgfqpoint{2.406434in}{1.016470in}}%
\pgfpathclose%
\pgfusepath{fill}%
\end{pgfscope}%
\begin{pgfscope}%
\pgfpathrectangle{\pgfqpoint{0.100000in}{0.100000in}}{\pgfqpoint{3.007045in}{1.925000in}}%
\pgfusepath{clip}%
\pgfsetbuttcap%
\pgfsetmiterjoin%
\definecolor{currentfill}{rgb}{0.751173,0.901423,0.724491}%
\pgfsetfillcolor{currentfill}%
\pgfsetlinewidth{0.000000pt}%
\definecolor{currentstroke}{rgb}{0.000000,0.000000,0.000000}%
\pgfsetstrokecolor{currentstroke}%
\pgfsetstrokeopacity{0.000000}%
\pgfsetdash{}{0pt}%
\pgfpathmoveto{\pgfqpoint{1.419469in}{0.951860in}}%
\pgfpathlineto{\pgfqpoint{1.417413in}{0.923222in}}%
\pgfpathlineto{\pgfqpoint{1.388570in}{0.925898in}}%
\pgfpathlineto{\pgfqpoint{1.341578in}{0.929505in}}%
\pgfpathlineto{\pgfqpoint{1.342266in}{0.937180in}}%
\pgfpathlineto{\pgfqpoint{1.346674in}{0.987189in}}%
\pgfpathlineto{\pgfqpoint{1.348759in}{0.987002in}}%
\pgfpathlineto{\pgfqpoint{1.393120in}{0.983287in}}%
\pgfpathlineto{\pgfqpoint{1.390781in}{0.954030in}}%
\pgfpathlineto{\pgfqpoint{1.419469in}{0.951860in}}%
\pgfpathclose%
\pgfusepath{fill}%
\end{pgfscope}%
\begin{pgfscope}%
\pgfpathrectangle{\pgfqpoint{0.100000in}{0.100000in}}{\pgfqpoint{3.007045in}{1.925000in}}%
\pgfusepath{clip}%
\pgfsetbuttcap%
\pgfsetmiterjoin%
\definecolor{currentfill}{rgb}{0.820069,0.929596,0.795263}%
\pgfsetfillcolor{currentfill}%
\pgfsetlinewidth{0.000000pt}%
\definecolor{currentstroke}{rgb}{0.000000,0.000000,0.000000}%
\pgfsetstrokecolor{currentstroke}%
\pgfsetstrokeopacity{0.000000}%
\pgfsetdash{}{0pt}%
\pgfpathmoveto{\pgfqpoint{1.605884in}{0.418493in}}%
\pgfpathlineto{\pgfqpoint{1.600656in}{0.414463in}}%
\pgfpathlineto{\pgfqpoint{1.601772in}{0.410582in}}%
\pgfpathlineto{\pgfqpoint{1.547367in}{0.412624in}}%
\pgfpathlineto{\pgfqpoint{1.549167in}{0.452106in}}%
\pgfpathlineto{\pgfqpoint{1.550168in}{0.481716in}}%
\pgfpathlineto{\pgfqpoint{1.514641in}{0.483384in}}%
\pgfpathlineto{\pgfqpoint{1.516499in}{0.519146in}}%
\pgfpathlineto{\pgfqpoint{1.505414in}{0.519674in}}%
\pgfpathlineto{\pgfqpoint{1.506378in}{0.538351in}}%
\pgfpathlineto{\pgfqpoint{1.530916in}{0.536323in}}%
\pgfpathlineto{\pgfqpoint{1.545579in}{0.528102in}}%
\pgfpathlineto{\pgfqpoint{1.546459in}{0.551872in}}%
\pgfpathlineto{\pgfqpoint{1.565594in}{0.551126in}}%
\pgfpathlineto{\pgfqpoint{1.575104in}{0.537324in}}%
\pgfpathlineto{\pgfqpoint{1.582069in}{0.543767in}}%
\pgfpathlineto{\pgfqpoint{1.597051in}{0.530605in}}%
\pgfpathlineto{\pgfqpoint{1.598646in}{0.524138in}}%
\pgfpathlineto{\pgfqpoint{1.606065in}{0.530962in}}%
\pgfpathlineto{\pgfqpoint{1.615672in}{0.518019in}}%
\pgfpathlineto{\pgfqpoint{1.619645in}{0.516644in}}%
\pgfpathlineto{\pgfqpoint{1.605940in}{0.497218in}}%
\pgfpathlineto{\pgfqpoint{1.613261in}{0.488228in}}%
\pgfpathlineto{\pgfqpoint{1.585508in}{0.466431in}}%
\pgfpathlineto{\pgfqpoint{1.595954in}{0.453216in}}%
\pgfpathlineto{\pgfqpoint{1.590934in}{0.451560in}}%
\pgfpathlineto{\pgfqpoint{1.606432in}{0.418922in}}%
\pgfpathlineto{\pgfqpoint{1.605884in}{0.418493in}}%
\pgfpathclose%
\pgfusepath{fill}%
\end{pgfscope}%
\begin{pgfscope}%
\pgfpathrectangle{\pgfqpoint{0.100000in}{0.100000in}}{\pgfqpoint{3.007045in}{1.925000in}}%
\pgfusepath{clip}%
\pgfsetbuttcap%
\pgfsetmiterjoin%
\definecolor{currentfill}{rgb}{0.727797,0.891580,0.701730}%
\pgfsetfillcolor{currentfill}%
\pgfsetlinewidth{0.000000pt}%
\definecolor{currentstroke}{rgb}{0.000000,0.000000,0.000000}%
\pgfsetstrokecolor{currentstroke}%
\pgfsetstrokeopacity{0.000000}%
\pgfsetdash{}{0pt}%
\pgfpathmoveto{\pgfqpoint{1.542619in}{0.830859in}}%
\pgfpathlineto{\pgfqpoint{1.533732in}{0.839225in}}%
\pgfpathlineto{\pgfqpoint{1.531947in}{0.833744in}}%
\pgfpathlineto{\pgfqpoint{1.526280in}{0.836772in}}%
\pgfpathlineto{\pgfqpoint{1.521539in}{0.834387in}}%
\pgfpathlineto{\pgfqpoint{1.516379in}{0.834892in}}%
\pgfpathlineto{\pgfqpoint{1.504435in}{0.848612in}}%
\pgfpathlineto{\pgfqpoint{1.500565in}{0.847796in}}%
\pgfpathlineto{\pgfqpoint{1.501084in}{0.860066in}}%
\pgfpathlineto{\pgfqpoint{1.502115in}{0.878788in}}%
\pgfpathlineto{\pgfqpoint{1.513852in}{0.878138in}}%
\pgfpathlineto{\pgfqpoint{1.514160in}{0.883851in}}%
\pgfpathlineto{\pgfqpoint{1.569929in}{0.881210in}}%
\pgfpathlineto{\pgfqpoint{1.576781in}{0.879670in}}%
\pgfpathlineto{\pgfqpoint{1.576210in}{0.863673in}}%
\pgfpathlineto{\pgfqpoint{1.565031in}{0.864162in}}%
\pgfpathlineto{\pgfqpoint{1.564310in}{0.846924in}}%
\pgfpathlineto{\pgfqpoint{1.545976in}{0.847321in}}%
\pgfpathlineto{\pgfqpoint{1.542518in}{0.841248in}}%
\pgfpathlineto{\pgfqpoint{1.542619in}{0.830859in}}%
\pgfpathclose%
\pgfusepath{fill}%
\end{pgfscope}%
\begin{pgfscope}%
\pgfpathrectangle{\pgfqpoint{0.100000in}{0.100000in}}{\pgfqpoint{3.007045in}{1.925000in}}%
\pgfusepath{clip}%
\pgfsetbuttcap%
\pgfsetmiterjoin%
\definecolor{currentfill}{rgb}{0.842215,0.938454,0.818885}%
\pgfsetfillcolor{currentfill}%
\pgfsetlinewidth{0.000000pt}%
\definecolor{currentstroke}{rgb}{0.000000,0.000000,0.000000}%
\pgfsetstrokecolor{currentstroke}%
\pgfsetstrokeopacity{0.000000}%
\pgfsetdash{}{0pt}%
\pgfpathmoveto{\pgfqpoint{1.702839in}{1.116383in}}%
\pgfpathlineto{\pgfqpoint{1.702564in}{1.102052in}}%
\pgfpathlineto{\pgfqpoint{1.678559in}{1.102524in}}%
\pgfpathlineto{\pgfqpoint{1.672876in}{1.102642in}}%
\pgfpathlineto{\pgfqpoint{1.672984in}{1.108364in}}%
\pgfpathlineto{\pgfqpoint{1.650255in}{1.108902in}}%
\pgfpathlineto{\pgfqpoint{1.651571in}{1.154782in}}%
\pgfpathlineto{\pgfqpoint{1.652076in}{1.171950in}}%
\pgfpathlineto{\pgfqpoint{1.709420in}{1.170767in}}%
\pgfpathlineto{\pgfqpoint{1.719727in}{1.170628in}}%
\pgfpathlineto{\pgfqpoint{1.719484in}{1.147603in}}%
\pgfpathlineto{\pgfqpoint{1.712256in}{1.145853in}}%
\pgfpathlineto{\pgfqpoint{1.708830in}{1.147175in}}%
\pgfpathlineto{\pgfqpoint{1.701412in}{1.144959in}}%
\pgfpathlineto{\pgfqpoint{1.701186in}{1.136436in}}%
\pgfpathlineto{\pgfqpoint{1.695539in}{1.136550in}}%
\pgfpathlineto{\pgfqpoint{1.695279in}{1.122232in}}%
\pgfpathlineto{\pgfqpoint{1.700977in}{1.122113in}}%
\pgfpathlineto{\pgfqpoint{1.702839in}{1.116383in}}%
\pgfpathclose%
\pgfusepath{fill}%
\end{pgfscope}%
\begin{pgfscope}%
\pgfpathrectangle{\pgfqpoint{0.100000in}{0.100000in}}{\pgfqpoint{3.007045in}{1.925000in}}%
\pgfusepath{clip}%
\pgfsetbuttcap%
\pgfsetmiterjoin%
\definecolor{currentfill}{rgb}{0.590542,0.831926,0.574271}%
\pgfsetfillcolor{currentfill}%
\pgfsetlinewidth{0.000000pt}%
\definecolor{currentstroke}{rgb}{0.000000,0.000000,0.000000}%
\pgfsetstrokecolor{currentstroke}%
\pgfsetstrokeopacity{0.000000}%
\pgfsetdash{}{0pt}%
\pgfpathmoveto{\pgfqpoint{2.505332in}{0.935315in}}%
\pgfpathlineto{\pgfqpoint{2.472948in}{0.932025in}}%
\pgfpathlineto{\pgfqpoint{2.472311in}{0.936614in}}%
\pgfpathlineto{\pgfqpoint{2.460723in}{0.945065in}}%
\pgfpathlineto{\pgfqpoint{2.455447in}{0.943242in}}%
\pgfpathlineto{\pgfqpoint{2.453980in}{0.946316in}}%
\pgfpathlineto{\pgfqpoint{2.455786in}{0.951722in}}%
\pgfpathlineto{\pgfqpoint{2.459120in}{0.952751in}}%
\pgfpathlineto{\pgfqpoint{2.469542in}{0.954025in}}%
\pgfpathlineto{\pgfqpoint{2.476332in}{0.956183in}}%
\pgfpathlineto{\pgfqpoint{2.480332in}{0.961447in}}%
\pgfpathlineto{\pgfqpoint{2.486422in}{0.958962in}}%
\pgfpathlineto{\pgfqpoint{2.492448in}{0.960338in}}%
\pgfpathlineto{\pgfqpoint{2.523136in}{0.963616in}}%
\pgfpathlineto{\pgfqpoint{2.523665in}{0.961730in}}%
\pgfpathlineto{\pgfqpoint{2.524892in}{0.952370in}}%
\pgfpathlineto{\pgfqpoint{2.523035in}{0.944991in}}%
\pgfpathlineto{\pgfqpoint{2.523780in}{0.937038in}}%
\pgfpathlineto{\pgfqpoint{2.505332in}{0.935315in}}%
\pgfpathclose%
\pgfusepath{fill}%
\end{pgfscope}%
\begin{pgfscope}%
\pgfpathrectangle{\pgfqpoint{0.100000in}{0.100000in}}{\pgfqpoint{3.007045in}{1.925000in}}%
\pgfusepath{clip}%
\pgfsetbuttcap%
\pgfsetmiterjoin%
\definecolor{currentfill}{rgb}{0.579469,0.826759,0.565167}%
\pgfsetfillcolor{currentfill}%
\pgfsetlinewidth{0.000000pt}%
\definecolor{currentstroke}{rgb}{0.000000,0.000000,0.000000}%
\pgfsetstrokecolor{currentstroke}%
\pgfsetstrokeopacity{0.000000}%
\pgfsetdash{}{0pt}%
\pgfpathmoveto{\pgfqpoint{2.452246in}{0.846939in}}%
\pgfpathlineto{\pgfqpoint{2.447000in}{0.848563in}}%
\pgfpathlineto{\pgfqpoint{2.440442in}{0.846333in}}%
\pgfpathlineto{\pgfqpoint{2.431219in}{0.848011in}}%
\pgfpathlineto{\pgfqpoint{2.421120in}{0.857223in}}%
\pgfpathlineto{\pgfqpoint{2.422010in}{0.860947in}}%
\pgfpathlineto{\pgfqpoint{2.409663in}{0.861811in}}%
\pgfpathlineto{\pgfqpoint{2.406992in}{0.858796in}}%
\pgfpathlineto{\pgfqpoint{2.404051in}{0.865135in}}%
\pgfpathlineto{\pgfqpoint{2.403763in}{0.872065in}}%
\pgfpathlineto{\pgfqpoint{2.399169in}{0.874931in}}%
\pgfpathlineto{\pgfqpoint{2.402644in}{0.888304in}}%
\pgfpathlineto{\pgfqpoint{2.403995in}{0.889257in}}%
\pgfpathlineto{\pgfqpoint{2.410336in}{0.885056in}}%
\pgfpathlineto{\pgfqpoint{2.413604in}{0.885230in}}%
\pgfpathlineto{\pgfqpoint{2.420077in}{0.879039in}}%
\pgfpathlineto{\pgfqpoint{2.423854in}{0.877561in}}%
\pgfpathlineto{\pgfqpoint{2.429221in}{0.879173in}}%
\pgfpathlineto{\pgfqpoint{2.439934in}{0.864954in}}%
\pgfpathlineto{\pgfqpoint{2.442202in}{0.858557in}}%
\pgfpathlineto{\pgfqpoint{2.449961in}{0.851491in}}%
\pgfpathlineto{\pgfqpoint{2.452246in}{0.846939in}}%
\pgfpathclose%
\pgfusepath{fill}%
\end{pgfscope}%
\begin{pgfscope}%
\pgfpathrectangle{\pgfqpoint{0.100000in}{0.100000in}}{\pgfqpoint{3.007045in}{1.925000in}}%
\pgfusepath{clip}%
\pgfsetbuttcap%
\pgfsetmiterjoin%
\definecolor{currentfill}{rgb}{0.463206,0.772503,0.469573}%
\pgfsetfillcolor{currentfill}%
\pgfsetlinewidth{0.000000pt}%
\definecolor{currentstroke}{rgb}{0.000000,0.000000,0.000000}%
\pgfsetstrokecolor{currentstroke}%
\pgfsetstrokeopacity{0.000000}%
\pgfsetdash{}{0pt}%
\pgfpathmoveto{\pgfqpoint{0.816917in}{1.293210in}}%
\pgfpathlineto{\pgfqpoint{0.811909in}{1.269809in}}%
\pgfpathlineto{\pgfqpoint{0.799942in}{1.214078in}}%
\pgfpathlineto{\pgfqpoint{0.752166in}{1.224502in}}%
\pgfpathlineto{\pgfqpoint{0.714021in}{1.265975in}}%
\pgfpathlineto{\pgfqpoint{0.679545in}{1.274162in}}%
\pgfpathlineto{\pgfqpoint{0.643335in}{1.283223in}}%
\pgfpathlineto{\pgfqpoint{0.620230in}{1.284320in}}%
\pgfpathlineto{\pgfqpoint{0.621234in}{1.289668in}}%
\pgfpathlineto{\pgfqpoint{0.625302in}{1.292738in}}%
\pgfpathlineto{\pgfqpoint{0.625081in}{1.300716in}}%
\pgfpathlineto{\pgfqpoint{0.632081in}{1.306373in}}%
\pgfpathlineto{\pgfqpoint{0.641368in}{1.309018in}}%
\pgfpathlineto{\pgfqpoint{0.644313in}{1.314171in}}%
\pgfpathlineto{\pgfqpoint{0.644679in}{1.322619in}}%
\pgfpathlineto{\pgfqpoint{0.648824in}{1.333373in}}%
\pgfpathlineto{\pgfqpoint{0.646486in}{1.339327in}}%
\pgfpathlineto{\pgfqpoint{0.671162in}{1.376694in}}%
\pgfpathlineto{\pgfqpoint{0.682240in}{1.373981in}}%
\pgfpathlineto{\pgfqpoint{0.687900in}{1.396831in}}%
\pgfpathlineto{\pgfqpoint{0.703710in}{1.460774in}}%
\pgfpathlineto{\pgfqpoint{0.734607in}{1.453046in}}%
\pgfpathlineto{\pgfqpoint{0.798116in}{1.438303in}}%
\pgfpathlineto{\pgfqpoint{0.808116in}{1.436493in}}%
\pgfpathlineto{\pgfqpoint{0.834301in}{1.430229in}}%
\pgfpathlineto{\pgfqpoint{0.845804in}{1.427734in}}%
\pgfpathlineto{\pgfqpoint{0.832155in}{1.363654in}}%
\pgfpathlineto{\pgfqpoint{0.816917in}{1.293210in}}%
\pgfpathclose%
\pgfusepath{fill}%
\end{pgfscope}%
\begin{pgfscope}%
\pgfpathrectangle{\pgfqpoint{0.100000in}{0.100000in}}{\pgfqpoint{3.007045in}{1.925000in}}%
\pgfusepath{clip}%
\pgfsetbuttcap%
\pgfsetmiterjoin%
\definecolor{currentfill}{rgb}{0.727797,0.891580,0.701730}%
\pgfsetfillcolor{currentfill}%
\pgfsetlinewidth{0.000000pt}%
\definecolor{currentstroke}{rgb}{0.000000,0.000000,0.000000}%
\pgfsetstrokecolor{currentstroke}%
\pgfsetstrokeopacity{0.000000}%
\pgfsetdash{}{0pt}%
\pgfpathmoveto{\pgfqpoint{1.566436in}{1.434578in}}%
\pgfpathlineto{\pgfqpoint{1.568472in}{1.430306in}}%
\pgfpathlineto{\pgfqpoint{1.572614in}{1.429955in}}%
\pgfpathlineto{\pgfqpoint{1.576855in}{1.422263in}}%
\pgfpathlineto{\pgfqpoint{1.584652in}{1.417090in}}%
\pgfpathlineto{\pgfqpoint{1.586747in}{1.410853in}}%
\pgfpathlineto{\pgfqpoint{1.590580in}{1.410117in}}%
\pgfpathlineto{\pgfqpoint{1.594388in}{1.405377in}}%
\pgfpathlineto{\pgfqpoint{1.602167in}{1.402735in}}%
\pgfpathlineto{\pgfqpoint{1.603438in}{1.399865in}}%
\pgfpathlineto{\pgfqpoint{1.566978in}{1.401393in}}%
\pgfpathlineto{\pgfqpoint{1.521452in}{1.403789in}}%
\pgfpathlineto{\pgfqpoint{1.521204in}{1.412401in}}%
\pgfpathlineto{\pgfqpoint{1.522496in}{1.435436in}}%
\pgfpathlineto{\pgfqpoint{1.522608in}{1.451050in}}%
\pgfpathlineto{\pgfqpoint{1.531144in}{1.449134in}}%
\pgfpathlineto{\pgfqpoint{1.541165in}{1.448356in}}%
\pgfpathlineto{\pgfqpoint{1.551671in}{1.451429in}}%
\pgfpathlineto{\pgfqpoint{1.550851in}{1.435352in}}%
\pgfpathlineto{\pgfqpoint{1.566436in}{1.434578in}}%
\pgfpathclose%
\pgfusepath{fill}%
\end{pgfscope}%
\begin{pgfscope}%
\pgfpathrectangle{\pgfqpoint{0.100000in}{0.100000in}}{\pgfqpoint{3.007045in}{1.925000in}}%
\pgfusepath{clip}%
\pgfsetbuttcap%
\pgfsetmiterjoin%
\definecolor{currentfill}{rgb}{0.568397,0.821592,0.556063}%
\pgfsetfillcolor{currentfill}%
\pgfsetlinewidth{0.000000pt}%
\definecolor{currentstroke}{rgb}{0.000000,0.000000,0.000000}%
\pgfsetstrokecolor{currentstroke}%
\pgfsetstrokeopacity{0.000000}%
\pgfsetdash{}{0pt}%
\pgfpathmoveto{\pgfqpoint{2.359552in}{0.655091in}}%
\pgfpathlineto{\pgfqpoint{2.343997in}{0.653135in}}%
\pgfpathlineto{\pgfqpoint{2.345074in}{0.640987in}}%
\pgfpathlineto{\pgfqpoint{2.339044in}{0.640515in}}%
\pgfpathlineto{\pgfqpoint{2.334338in}{0.646016in}}%
\pgfpathlineto{\pgfqpoint{2.332928in}{0.653635in}}%
\pgfpathlineto{\pgfqpoint{2.334926in}{0.671284in}}%
\pgfpathlineto{\pgfqpoint{2.333467in}{0.676210in}}%
\pgfpathlineto{\pgfqpoint{2.329091in}{0.680831in}}%
\pgfpathlineto{\pgfqpoint{2.312691in}{0.679850in}}%
\pgfpathlineto{\pgfqpoint{2.313278in}{0.674089in}}%
\pgfpathlineto{\pgfqpoint{2.294624in}{0.672055in}}%
\pgfpathlineto{\pgfqpoint{2.298312in}{0.683587in}}%
\pgfpathlineto{\pgfqpoint{2.297933in}{0.689896in}}%
\pgfpathlineto{\pgfqpoint{2.301073in}{0.698033in}}%
\pgfpathlineto{\pgfqpoint{2.309841in}{0.700126in}}%
\pgfpathlineto{\pgfqpoint{2.309942in}{0.708924in}}%
\pgfpathlineto{\pgfqpoint{2.318447in}{0.709939in}}%
\pgfpathlineto{\pgfqpoint{2.323124in}{0.704669in}}%
\pgfpathlineto{\pgfqpoint{2.330327in}{0.705545in}}%
\pgfpathlineto{\pgfqpoint{2.331743in}{0.699866in}}%
\pgfpathlineto{\pgfqpoint{2.339099in}{0.697275in}}%
\pgfpathlineto{\pgfqpoint{2.356903in}{0.698990in}}%
\pgfpathlineto{\pgfqpoint{2.357840in}{0.689197in}}%
\pgfpathlineto{\pgfqpoint{2.362322in}{0.683671in}}%
\pgfpathlineto{\pgfqpoint{2.367559in}{0.680287in}}%
\pgfpathlineto{\pgfqpoint{2.367869in}{0.675418in}}%
\pgfpathlineto{\pgfqpoint{2.370955in}{0.669671in}}%
\pgfpathlineto{\pgfqpoint{2.366771in}{0.667222in}}%
\pgfpathlineto{\pgfqpoint{2.358460in}{0.666668in}}%
\pgfpathlineto{\pgfqpoint{2.359552in}{0.655091in}}%
\pgfpathclose%
\pgfusepath{fill}%
\end{pgfscope}%
\begin{pgfscope}%
\pgfpathrectangle{\pgfqpoint{0.100000in}{0.100000in}}{\pgfqpoint{3.007045in}{1.925000in}}%
\pgfusepath{clip}%
\pgfsetbuttcap%
\pgfsetmiterjoin%
\definecolor{currentfill}{rgb}{0.479815,0.780254,0.483230}%
\pgfsetfillcolor{currentfill}%
\pgfsetlinewidth{0.000000pt}%
\definecolor{currentstroke}{rgb}{0.000000,0.000000,0.000000}%
\pgfsetstrokecolor{currentstroke}%
\pgfsetstrokeopacity{0.000000}%
\pgfsetdash{}{0pt}%
\pgfpathmoveto{\pgfqpoint{0.923792in}{1.677916in}}%
\pgfpathlineto{\pgfqpoint{0.921890in}{1.674121in}}%
\pgfpathlineto{\pgfqpoint{0.914135in}{1.673457in}}%
\pgfpathlineto{\pgfqpoint{0.908010in}{1.668930in}}%
\pgfpathlineto{\pgfqpoint{0.901414in}{1.666132in}}%
\pgfpathlineto{\pgfqpoint{0.894600in}{1.662965in}}%
\pgfpathlineto{\pgfqpoint{0.890559in}{1.659065in}}%
\pgfpathlineto{\pgfqpoint{0.879064in}{1.657147in}}%
\pgfpathlineto{\pgfqpoint{0.873352in}{1.663504in}}%
\pgfpathlineto{\pgfqpoint{0.871070in}{1.663297in}}%
\pgfpathlineto{\pgfqpoint{0.875599in}{1.669863in}}%
\pgfpathlineto{\pgfqpoint{0.874072in}{1.677313in}}%
\pgfpathlineto{\pgfqpoint{0.877655in}{1.681396in}}%
\pgfpathlineto{\pgfqpoint{0.883784in}{1.683111in}}%
\pgfpathlineto{\pgfqpoint{0.881292in}{1.694653in}}%
\pgfpathlineto{\pgfqpoint{0.885040in}{1.701682in}}%
\pgfpathlineto{\pgfqpoint{0.885687in}{1.708694in}}%
\pgfpathlineto{\pgfqpoint{0.891175in}{1.716853in}}%
\pgfpathlineto{\pgfqpoint{0.897388in}{1.731726in}}%
\pgfpathlineto{\pgfqpoint{0.895843in}{1.733369in}}%
\pgfpathlineto{\pgfqpoint{0.884196in}{1.733912in}}%
\pgfpathlineto{\pgfqpoint{0.884881in}{1.738405in}}%
\pgfpathlineto{\pgfqpoint{0.873285in}{1.749159in}}%
\pgfpathlineto{\pgfqpoint{0.873084in}{1.756371in}}%
\pgfpathlineto{\pgfqpoint{0.869191in}{1.761494in}}%
\pgfpathlineto{\pgfqpoint{0.861882in}{1.780983in}}%
\pgfpathlineto{\pgfqpoint{0.853297in}{1.785851in}}%
\pgfpathlineto{\pgfqpoint{0.853094in}{1.789354in}}%
\pgfpathlineto{\pgfqpoint{0.845810in}{1.797316in}}%
\pgfpathlineto{\pgfqpoint{0.851627in}{1.798562in}}%
\pgfpathlineto{\pgfqpoint{0.859543in}{1.795823in}}%
\pgfpathlineto{\pgfqpoint{0.866931in}{1.795554in}}%
\pgfpathlineto{\pgfqpoint{0.878913in}{1.786617in}}%
\pgfpathlineto{\pgfqpoint{0.878165in}{1.781128in}}%
\pgfpathlineto{\pgfqpoint{0.883118in}{1.777032in}}%
\pgfpathlineto{\pgfqpoint{0.890065in}{1.775890in}}%
\pgfpathlineto{\pgfqpoint{0.900754in}{1.769170in}}%
\pgfpathlineto{\pgfqpoint{0.906982in}{1.761729in}}%
\pgfpathlineto{\pgfqpoint{0.910196in}{1.762234in}}%
\pgfpathlineto{\pgfqpoint{0.921428in}{1.759848in}}%
\pgfpathlineto{\pgfqpoint{0.926970in}{1.761628in}}%
\pgfpathlineto{\pgfqpoint{0.926220in}{1.769071in}}%
\pgfpathlineto{\pgfqpoint{0.927695in}{1.775797in}}%
\pgfpathlineto{\pgfqpoint{0.925310in}{1.782011in}}%
\pgfpathlineto{\pgfqpoint{0.927953in}{1.789677in}}%
\pgfpathlineto{\pgfqpoint{0.948336in}{1.785640in}}%
\pgfpathlineto{\pgfqpoint{0.942694in}{1.758174in}}%
\pgfpathlineto{\pgfqpoint{0.949935in}{1.756722in}}%
\pgfpathlineto{\pgfqpoint{0.945308in}{1.734086in}}%
\pgfpathlineto{\pgfqpoint{0.940260in}{1.735106in}}%
\pgfpathlineto{\pgfqpoint{0.936258in}{1.730024in}}%
\pgfpathlineto{\pgfqpoint{0.927484in}{1.729880in}}%
\pgfpathlineto{\pgfqpoint{0.926725in}{1.726127in}}%
\pgfpathlineto{\pgfqpoint{0.921093in}{1.727298in}}%
\pgfpathlineto{\pgfqpoint{0.915399in}{1.718366in}}%
\pgfpathlineto{\pgfqpoint{0.914970in}{1.712349in}}%
\pgfpathlineto{\pgfqpoint{0.917545in}{1.708081in}}%
\pgfpathlineto{\pgfqpoint{0.917747in}{1.701623in}}%
\pgfpathlineto{\pgfqpoint{0.914111in}{1.694407in}}%
\pgfpathlineto{\pgfqpoint{0.913359in}{1.689243in}}%
\pgfpathlineto{\pgfqpoint{0.917372in}{1.685810in}}%
\pgfpathlineto{\pgfqpoint{0.918191in}{1.680868in}}%
\pgfpathlineto{\pgfqpoint{0.923792in}{1.677916in}}%
\pgfpathclose%
\pgfusepath{fill}%
\end{pgfscope}%
\begin{pgfscope}%
\pgfpathrectangle{\pgfqpoint{0.100000in}{0.100000in}}{\pgfqpoint{3.007045in}{1.925000in}}%
\pgfusepath{clip}%
\pgfsetbuttcap%
\pgfsetmiterjoin%
\definecolor{currentfill}{rgb}{0.420392,0.751711,0.445829}%
\pgfsetfillcolor{currentfill}%
\pgfsetlinewidth{0.000000pt}%
\definecolor{currentstroke}{rgb}{0.000000,0.000000,0.000000}%
\pgfsetstrokecolor{currentstroke}%
\pgfsetstrokeopacity{0.000000}%
\pgfsetdash{}{0pt}%
\pgfpathmoveto{\pgfqpoint{1.844441in}{1.606089in}}%
\pgfpathlineto{\pgfqpoint{1.861813in}{1.606700in}}%
\pgfpathlineto{\pgfqpoint{1.861397in}{1.624038in}}%
\pgfpathlineto{\pgfqpoint{1.896284in}{1.624795in}}%
\pgfpathlineto{\pgfqpoint{1.896759in}{1.607543in}}%
\pgfpathlineto{\pgfqpoint{1.896908in}{1.602038in}}%
\pgfpathlineto{\pgfqpoint{1.894425in}{1.598113in}}%
\pgfpathlineto{\pgfqpoint{1.881222in}{1.592179in}}%
\pgfpathlineto{\pgfqpoint{1.877917in}{1.589345in}}%
\pgfpathlineto{\pgfqpoint{1.874776in}{1.580802in}}%
\pgfpathlineto{\pgfqpoint{1.872272in}{1.578521in}}%
\pgfpathlineto{\pgfqpoint{1.841127in}{1.578186in}}%
\pgfpathlineto{\pgfqpoint{1.840662in}{1.594594in}}%
\pgfpathlineto{\pgfqpoint{1.844608in}{1.594634in}}%
\pgfpathlineto{\pgfqpoint{1.844441in}{1.606089in}}%
\pgfpathclose%
\pgfusepath{fill}%
\end{pgfscope}%
\begin{pgfscope}%
\pgfpathrectangle{\pgfqpoint{0.100000in}{0.100000in}}{\pgfqpoint{3.007045in}{1.925000in}}%
\pgfusepath{clip}%
\pgfsetbuttcap%
\pgfsetmiterjoin%
\definecolor{currentfill}{rgb}{0.507497,0.793172,0.505990}%
\pgfsetfillcolor{currentfill}%
\pgfsetlinewidth{0.000000pt}%
\definecolor{currentstroke}{rgb}{0.000000,0.000000,0.000000}%
\pgfsetstrokecolor{currentstroke}%
\pgfsetstrokeopacity{0.000000}%
\pgfsetdash{}{0pt}%
\pgfpathmoveto{\pgfqpoint{1.561807in}{1.005883in}}%
\pgfpathlineto{\pgfqpoint{1.537864in}{1.006990in}}%
\pgfpathlineto{\pgfqpoint{1.533400in}{1.007220in}}%
\pgfpathlineto{\pgfqpoint{1.534548in}{1.032337in}}%
\pgfpathlineto{\pgfqpoint{1.562899in}{1.031213in}}%
\pgfpathlineto{\pgfqpoint{1.561807in}{1.005883in}}%
\pgfpathclose%
\pgfusepath{fill}%
\end{pgfscope}%
\begin{pgfscope}%
\pgfpathrectangle{\pgfqpoint{0.100000in}{0.100000in}}{\pgfqpoint{3.007045in}{1.925000in}}%
\pgfusepath{clip}%
\pgfsetbuttcap%
\pgfsetmiterjoin%
\definecolor{currentfill}{rgb}{0.568397,0.821592,0.556063}%
\pgfsetfillcolor{currentfill}%
\pgfsetlinewidth{0.000000pt}%
\definecolor{currentstroke}{rgb}{0.000000,0.000000,0.000000}%
\pgfsetstrokecolor{currentstroke}%
\pgfsetstrokeopacity{0.000000}%
\pgfsetdash{}{0pt}%
\pgfpathmoveto{\pgfqpoint{2.174509in}{1.028592in}}%
\pgfpathlineto{\pgfqpoint{2.164956in}{1.024377in}}%
\pgfpathlineto{\pgfqpoint{2.156379in}{1.026602in}}%
\pgfpathlineto{\pgfqpoint{2.155476in}{1.018069in}}%
\pgfpathlineto{\pgfqpoint{2.152882in}{1.014762in}}%
\pgfpathlineto{\pgfqpoint{2.146362in}{1.013286in}}%
\pgfpathlineto{\pgfqpoint{2.132817in}{1.006098in}}%
\pgfpathlineto{\pgfqpoint{2.127814in}{1.020206in}}%
\pgfpathlineto{\pgfqpoint{2.129826in}{1.024933in}}%
\pgfpathlineto{\pgfqpoint{2.127608in}{1.032537in}}%
\pgfpathlineto{\pgfqpoint{2.119654in}{1.040653in}}%
\pgfpathlineto{\pgfqpoint{2.123429in}{1.044085in}}%
\pgfpathlineto{\pgfqpoint{2.133456in}{1.046059in}}%
\pgfpathlineto{\pgfqpoint{2.128798in}{1.058421in}}%
\pgfpathlineto{\pgfqpoint{2.134908in}{1.067855in}}%
\pgfpathlineto{\pgfqpoint{2.141157in}{1.068895in}}%
\pgfpathlineto{\pgfqpoint{2.139094in}{1.074231in}}%
\pgfpathlineto{\pgfqpoint{2.144755in}{1.073731in}}%
\pgfpathlineto{\pgfqpoint{2.152373in}{1.076112in}}%
\pgfpathlineto{\pgfqpoint{2.156206in}{1.071747in}}%
\pgfpathlineto{\pgfqpoint{2.158994in}{1.078092in}}%
\pgfpathlineto{\pgfqpoint{2.165193in}{1.079905in}}%
\pgfpathlineto{\pgfqpoint{2.172010in}{1.077414in}}%
\pgfpathlineto{\pgfqpoint{2.171424in}{1.072069in}}%
\pgfpathlineto{\pgfqpoint{2.164342in}{1.057050in}}%
\pgfpathlineto{\pgfqpoint{2.170940in}{1.054187in}}%
\pgfpathlineto{\pgfqpoint{2.172018in}{1.046217in}}%
\pgfpathlineto{\pgfqpoint{2.174985in}{1.043028in}}%
\pgfpathlineto{\pgfqpoint{2.171148in}{1.034806in}}%
\pgfpathlineto{\pgfqpoint{2.174509in}{1.028592in}}%
\pgfpathclose%
\pgfusepath{fill}%
\end{pgfscope}%
\begin{pgfscope}%
\pgfpathrectangle{\pgfqpoint{0.100000in}{0.100000in}}{\pgfqpoint{3.007045in}{1.925000in}}%
\pgfusepath{clip}%
\pgfsetbuttcap%
\pgfsetmiterjoin%
\definecolor{currentfill}{rgb}{0.618224,0.844844,0.597032}%
\pgfsetfillcolor{currentfill}%
\pgfsetlinewidth{0.000000pt}%
\definecolor{currentstroke}{rgb}{0.000000,0.000000,0.000000}%
\pgfsetstrokecolor{currentstroke}%
\pgfsetstrokeopacity{0.000000}%
\pgfsetdash{}{0pt}%
\pgfpathmoveto{\pgfqpoint{1.381806in}{1.702240in}}%
\pgfpathlineto{\pgfqpoint{1.379578in}{1.679082in}}%
\pgfpathlineto{\pgfqpoint{1.382174in}{1.678842in}}%
\pgfpathlineto{\pgfqpoint{1.379987in}{1.655666in}}%
\pgfpathlineto{\pgfqpoint{1.371277in}{1.656487in}}%
\pgfpathlineto{\pgfqpoint{1.370707in}{1.650612in}}%
\pgfpathlineto{\pgfqpoint{1.359600in}{1.651719in}}%
\pgfpathlineto{\pgfqpoint{1.360257in}{1.658366in}}%
\pgfpathlineto{\pgfqpoint{1.346246in}{1.659789in}}%
\pgfpathlineto{\pgfqpoint{1.343644in}{1.662962in}}%
\pgfpathlineto{\pgfqpoint{1.335063in}{1.663866in}}%
\pgfpathlineto{\pgfqpoint{1.336444in}{1.675484in}}%
\pgfpathlineto{\pgfqpoint{1.317293in}{1.677601in}}%
\pgfpathlineto{\pgfqpoint{1.315821in}{1.681622in}}%
\pgfpathlineto{\pgfqpoint{1.310066in}{1.682269in}}%
\pgfpathlineto{\pgfqpoint{1.310508in}{1.686101in}}%
\pgfpathlineto{\pgfqpoint{1.304728in}{1.686756in}}%
\pgfpathlineto{\pgfqpoint{1.306273in}{1.700261in}}%
\pgfpathlineto{\pgfqpoint{1.302614in}{1.700688in}}%
\pgfpathlineto{\pgfqpoint{1.305247in}{1.723687in}}%
\pgfpathlineto{\pgfqpoint{1.307173in}{1.723478in}}%
\pgfpathlineto{\pgfqpoint{1.309172in}{1.740759in}}%
\pgfpathlineto{\pgfqpoint{1.314888in}{1.740131in}}%
\pgfpathlineto{\pgfqpoint{1.314224in}{1.734361in}}%
\pgfpathlineto{\pgfqpoint{1.325673in}{1.733068in}}%
\pgfpathlineto{\pgfqpoint{1.325013in}{1.727299in}}%
\pgfpathlineto{\pgfqpoint{1.342199in}{1.725419in}}%
\pgfpathlineto{\pgfqpoint{1.341580in}{1.719656in}}%
\pgfpathlineto{\pgfqpoint{1.345202in}{1.719293in}}%
\pgfpathlineto{\pgfqpoint{1.343969in}{1.707734in}}%
\pgfpathlineto{\pgfqpoint{1.361096in}{1.705933in}}%
\pgfpathlineto{\pgfqpoint{1.364804in}{1.703953in}}%
\pgfpathlineto{\pgfqpoint{1.381806in}{1.702240in}}%
\pgfpathclose%
\pgfusepath{fill}%
\end{pgfscope}%
\begin{pgfscope}%
\pgfpathrectangle{\pgfqpoint{0.100000in}{0.100000in}}{\pgfqpoint{3.007045in}{1.925000in}}%
\pgfusepath{clip}%
\pgfsetbuttcap%
\pgfsetmiterjoin%
\definecolor{currentfill}{rgb}{0.662345,0.864022,0.638001}%
\pgfsetfillcolor{currentfill}%
\pgfsetlinewidth{0.000000pt}%
\definecolor{currentstroke}{rgb}{0.000000,0.000000,0.000000}%
\pgfsetstrokecolor{currentstroke}%
\pgfsetstrokeopacity{0.000000}%
\pgfsetdash{}{0pt}%
\pgfpathmoveto{\pgfqpoint{2.128961in}{1.204856in}}%
\pgfpathlineto{\pgfqpoint{2.128488in}{1.192008in}}%
\pgfpathlineto{\pgfqpoint{2.123761in}{1.189439in}}%
\pgfpathlineto{\pgfqpoint{2.103104in}{1.187835in}}%
\pgfpathlineto{\pgfqpoint{2.105221in}{1.159153in}}%
\pgfpathlineto{\pgfqpoint{2.071143in}{1.157025in}}%
\pgfpathlineto{\pgfqpoint{2.070604in}{1.165647in}}%
\pgfpathlineto{\pgfqpoint{2.076427in}{1.165836in}}%
\pgfpathlineto{\pgfqpoint{2.075079in}{1.186079in}}%
\pgfpathlineto{\pgfqpoint{2.069337in}{1.185778in}}%
\pgfpathlineto{\pgfqpoint{2.068581in}{1.195334in}}%
\pgfpathlineto{\pgfqpoint{2.064705in}{1.195908in}}%
\pgfpathlineto{\pgfqpoint{2.064257in}{1.202733in}}%
\pgfpathlineto{\pgfqpoint{2.068009in}{1.203031in}}%
\pgfpathlineto{\pgfqpoint{2.067386in}{1.211625in}}%
\pgfpathlineto{\pgfqpoint{2.087435in}{1.213369in}}%
\pgfpathlineto{\pgfqpoint{2.094960in}{1.228863in}}%
\pgfpathlineto{\pgfqpoint{2.100744in}{1.229280in}}%
\pgfpathlineto{\pgfqpoint{2.099165in}{1.251343in}}%
\pgfpathlineto{\pgfqpoint{2.110362in}{1.252216in}}%
\pgfpathlineto{\pgfqpoint{2.107905in}{1.276923in}}%
\pgfpathlineto{\pgfqpoint{2.113613in}{1.277542in}}%
\pgfpathlineto{\pgfqpoint{2.116838in}{1.244113in}}%
\pgfpathlineto{\pgfqpoint{2.125964in}{1.244692in}}%
\pgfpathlineto{\pgfqpoint{2.126956in}{1.227545in}}%
\pgfpathlineto{\pgfqpoint{2.128961in}{1.204856in}}%
\pgfpathclose%
\pgfusepath{fill}%
\end{pgfscope}%
\begin{pgfscope}%
\pgfpathrectangle{\pgfqpoint{0.100000in}{0.100000in}}{\pgfqpoint{3.007045in}{1.925000in}}%
\pgfusepath{clip}%
\pgfsetbuttcap%
\pgfsetmiterjoin%
\definecolor{currentfill}{rgb}{0.695071,0.877801,0.669865}%
\pgfsetfillcolor{currentfill}%
\pgfsetlinewidth{0.000000pt}%
\definecolor{currentstroke}{rgb}{0.000000,0.000000,0.000000}%
\pgfsetstrokecolor{currentstroke}%
\pgfsetstrokeopacity{0.000000}%
\pgfsetdash{}{0pt}%
\pgfpathmoveto{\pgfqpoint{1.381806in}{1.702240in}}%
\pgfpathlineto{\pgfqpoint{1.364804in}{1.703953in}}%
\pgfpathlineto{\pgfqpoint{1.361096in}{1.705933in}}%
\pgfpathlineto{\pgfqpoint{1.343969in}{1.707734in}}%
\pgfpathlineto{\pgfqpoint{1.345202in}{1.719293in}}%
\pgfpathlineto{\pgfqpoint{1.341580in}{1.719656in}}%
\pgfpathlineto{\pgfqpoint{1.342199in}{1.725419in}}%
\pgfpathlineto{\pgfqpoint{1.325013in}{1.727299in}}%
\pgfpathlineto{\pgfqpoint{1.325673in}{1.733068in}}%
\pgfpathlineto{\pgfqpoint{1.314224in}{1.734361in}}%
\pgfpathlineto{\pgfqpoint{1.314888in}{1.740131in}}%
\pgfpathlineto{\pgfqpoint{1.317232in}{1.745061in}}%
\pgfpathlineto{\pgfqpoint{1.318715in}{1.758150in}}%
\pgfpathlineto{\pgfqpoint{1.324957in}{1.756945in}}%
\pgfpathlineto{\pgfqpoint{1.335725in}{1.759869in}}%
\pgfpathlineto{\pgfqpoint{1.341818in}{1.758120in}}%
\pgfpathlineto{\pgfqpoint{1.351556in}{1.758299in}}%
\pgfpathlineto{\pgfqpoint{1.355902in}{1.752996in}}%
\pgfpathlineto{\pgfqpoint{1.366019in}{1.751221in}}%
\pgfpathlineto{\pgfqpoint{1.369270in}{1.748112in}}%
\pgfpathlineto{\pgfqpoint{1.376795in}{1.745675in}}%
\pgfpathlineto{\pgfqpoint{1.386114in}{1.754885in}}%
\pgfpathlineto{\pgfqpoint{1.389403in}{1.754089in}}%
\pgfpathlineto{\pgfqpoint{1.390611in}{1.748693in}}%
\pgfpathlineto{\pgfqpoint{1.394175in}{1.746414in}}%
\pgfpathlineto{\pgfqpoint{1.403722in}{1.748968in}}%
\pgfpathlineto{\pgfqpoint{1.407533in}{1.753794in}}%
\pgfpathlineto{\pgfqpoint{1.422411in}{1.752597in}}%
\pgfpathlineto{\pgfqpoint{1.425483in}{1.768138in}}%
\pgfpathlineto{\pgfqpoint{1.422948in}{1.768339in}}%
\pgfpathlineto{\pgfqpoint{1.423931in}{1.779899in}}%
\pgfpathlineto{\pgfqpoint{1.452658in}{1.777477in}}%
\pgfpathlineto{\pgfqpoint{1.464168in}{1.776657in}}%
\pgfpathlineto{\pgfqpoint{1.463268in}{1.765000in}}%
\pgfpathlineto{\pgfqpoint{1.465533in}{1.764842in}}%
\pgfpathlineto{\pgfqpoint{1.463751in}{1.741609in}}%
\pgfpathlineto{\pgfqpoint{1.466009in}{1.741430in}}%
\pgfpathlineto{\pgfqpoint{1.465127in}{1.729765in}}%
\pgfpathlineto{\pgfqpoint{1.442195in}{1.731563in}}%
\pgfpathlineto{\pgfqpoint{1.441706in}{1.725605in}}%
\pgfpathlineto{\pgfqpoint{1.433959in}{1.727722in}}%
\pgfpathlineto{\pgfqpoint{1.429757in}{1.720971in}}%
\pgfpathlineto{\pgfqpoint{1.409324in}{1.722778in}}%
\pgfpathlineto{\pgfqpoint{1.407305in}{1.699840in}}%
\pgfpathlineto{\pgfqpoint{1.381806in}{1.702240in}}%
\pgfpathclose%
\pgfusepath{fill}%
\end{pgfscope}%
\begin{pgfscope}%
\pgfpathrectangle{\pgfqpoint{0.100000in}{0.100000in}}{\pgfqpoint{3.007045in}{1.925000in}}%
\pgfusepath{clip}%
\pgfsetbuttcap%
\pgfsetmiterjoin%
\definecolor{currentfill}{rgb}{0.000000,0.266667,0.105882}%
\pgfsetfillcolor{currentfill}%
\pgfsetlinewidth{0.000000pt}%
\definecolor{currentstroke}{rgb}{0.000000,0.000000,0.000000}%
\pgfsetstrokecolor{currentstroke}%
\pgfsetstrokeopacity{0.000000}%
\pgfsetdash{}{0pt}%
\pgfpathmoveto{\pgfqpoint{2.962613in}{1.457141in}}%
\pgfpathlineto{\pgfqpoint{2.972406in}{1.450863in}}%
\pgfpathlineto{\pgfqpoint{2.960151in}{1.447007in}}%
\pgfpathlineto{\pgfqpoint{2.958968in}{1.452964in}}%
\pgfpathlineto{\pgfqpoint{2.962613in}{1.457141in}}%
\pgfpathclose%
\pgfusepath{fill}%
\end{pgfscope}%
\begin{pgfscope}%
\pgfpathrectangle{\pgfqpoint{0.100000in}{0.100000in}}{\pgfqpoint{3.007045in}{1.925000in}}%
\pgfusepath{clip}%
\pgfsetbuttcap%
\pgfsetmiterjoin%
\definecolor{currentfill}{rgb}{0.557324,0.816424,0.546959}%
\pgfsetfillcolor{currentfill}%
\pgfsetlinewidth{0.000000pt}%
\definecolor{currentstroke}{rgb}{0.000000,0.000000,0.000000}%
\pgfsetstrokecolor{currentstroke}%
\pgfsetstrokeopacity{0.000000}%
\pgfsetdash{}{0pt}%
\pgfpathmoveto{\pgfqpoint{2.644603in}{1.131634in}}%
\pgfpathlineto{\pgfqpoint{2.641185in}{1.127114in}}%
\pgfpathlineto{\pgfqpoint{2.641373in}{1.123470in}}%
\pgfpathlineto{\pgfqpoint{2.633816in}{1.102244in}}%
\pgfpathlineto{\pgfqpoint{2.626460in}{1.104802in}}%
\pgfpathlineto{\pgfqpoint{2.621502in}{1.105990in}}%
\pgfpathlineto{\pgfqpoint{2.612770in}{1.112672in}}%
\pgfpathlineto{\pgfqpoint{2.613630in}{1.118818in}}%
\pgfpathlineto{\pgfqpoint{2.617784in}{1.122931in}}%
\pgfpathlineto{\pgfqpoint{2.616773in}{1.129743in}}%
\pgfpathlineto{\pgfqpoint{2.606282in}{1.144642in}}%
\pgfpathlineto{\pgfqpoint{2.608923in}{1.148127in}}%
\pgfpathlineto{\pgfqpoint{2.609006in}{1.155780in}}%
\pgfpathlineto{\pgfqpoint{2.613166in}{1.160149in}}%
\pgfpathlineto{\pgfqpoint{2.619686in}{1.171111in}}%
\pgfpathlineto{\pgfqpoint{2.621002in}{1.178691in}}%
\pgfpathlineto{\pgfqpoint{2.624882in}{1.186857in}}%
\pgfpathlineto{\pgfqpoint{2.631429in}{1.181706in}}%
\pgfpathlineto{\pgfqpoint{2.634662in}{1.181683in}}%
\pgfpathlineto{\pgfqpoint{2.644402in}{1.195106in}}%
\pgfpathlineto{\pgfqpoint{2.648785in}{1.188481in}}%
\pgfpathlineto{\pgfqpoint{2.658001in}{1.178784in}}%
\pgfpathlineto{\pgfqpoint{2.663143in}{1.179525in}}%
\pgfpathlineto{\pgfqpoint{2.650904in}{1.157743in}}%
\pgfpathlineto{\pgfqpoint{2.656779in}{1.157595in}}%
\pgfpathlineto{\pgfqpoint{2.665522in}{1.153025in}}%
\pgfpathlineto{\pgfqpoint{2.663369in}{1.133914in}}%
\pgfpathlineto{\pgfqpoint{2.658319in}{1.135059in}}%
\pgfpathlineto{\pgfqpoint{2.654207in}{1.139998in}}%
\pgfpathlineto{\pgfqpoint{2.647609in}{1.142718in}}%
\pgfpathlineto{\pgfqpoint{2.644603in}{1.131634in}}%
\pgfpathclose%
\pgfusepath{fill}%
\end{pgfscope}%
\begin{pgfscope}%
\pgfpathrectangle{\pgfqpoint{0.100000in}{0.100000in}}{\pgfqpoint{3.007045in}{1.925000in}}%
\pgfusepath{clip}%
\pgfsetbuttcap%
\pgfsetmiterjoin%
\definecolor{currentfill}{rgb}{0.790542,0.917785,0.763768}%
\pgfsetfillcolor{currentfill}%
\pgfsetlinewidth{0.000000pt}%
\definecolor{currentstroke}{rgb}{0.000000,0.000000,0.000000}%
\pgfsetstrokecolor{currentstroke}%
\pgfsetstrokeopacity{0.000000}%
\pgfsetdash{}{0pt}%
\pgfpathmoveto{\pgfqpoint{1.444186in}{0.892364in}}%
\pgfpathlineto{\pgfqpoint{1.442027in}{0.863792in}}%
\pgfpathlineto{\pgfqpoint{1.421405in}{0.865168in}}%
\pgfpathlineto{\pgfqpoint{1.392891in}{0.867332in}}%
\pgfpathlineto{\pgfqpoint{1.364391in}{0.869531in}}%
\pgfpathlineto{\pgfqpoint{1.336407in}{0.871962in}}%
\pgfpathlineto{\pgfqpoint{1.337618in}{0.885562in}}%
\pgfpathlineto{\pgfqpoint{1.341578in}{0.929505in}}%
\pgfpathlineto{\pgfqpoint{1.388570in}{0.925898in}}%
\pgfpathlineto{\pgfqpoint{1.417413in}{0.923222in}}%
\pgfpathlineto{\pgfqpoint{1.415306in}{0.894451in}}%
\pgfpathlineto{\pgfqpoint{1.444186in}{0.892364in}}%
\pgfpathclose%
\pgfusepath{fill}%
\end{pgfscope}%
\begin{pgfscope}%
\pgfpathrectangle{\pgfqpoint{0.100000in}{0.100000in}}{\pgfqpoint{3.007045in}{1.925000in}}%
\pgfusepath{clip}%
\pgfsetbuttcap%
\pgfsetmiterjoin%
\definecolor{currentfill}{rgb}{0.709097,0.883706,0.683522}%
\pgfsetfillcolor{currentfill}%
\pgfsetlinewidth{0.000000pt}%
\definecolor{currentstroke}{rgb}{0.000000,0.000000,0.000000}%
\pgfsetstrokecolor{currentstroke}%
\pgfsetstrokeopacity{0.000000}%
\pgfsetdash{}{0pt}%
\pgfpathmoveto{\pgfqpoint{1.918185in}{0.809351in}}%
\pgfpathlineto{\pgfqpoint{1.909437in}{0.814697in}}%
\pgfpathlineto{\pgfqpoint{1.899320in}{0.814784in}}%
\pgfpathlineto{\pgfqpoint{1.899400in}{0.832281in}}%
\pgfpathlineto{\pgfqpoint{1.898412in}{0.836142in}}%
\pgfpathlineto{\pgfqpoint{1.892466in}{0.838016in}}%
\pgfpathlineto{\pgfqpoint{1.891528in}{0.843729in}}%
\pgfpathlineto{\pgfqpoint{1.887792in}{0.846596in}}%
\pgfpathlineto{\pgfqpoint{1.882114in}{0.846605in}}%
\pgfpathlineto{\pgfqpoint{1.882494in}{0.855536in}}%
\pgfpathlineto{\pgfqpoint{1.865214in}{0.855241in}}%
\pgfpathlineto{\pgfqpoint{1.864117in}{0.861985in}}%
\pgfpathlineto{\pgfqpoint{1.871864in}{0.868767in}}%
\pgfpathlineto{\pgfqpoint{1.884196in}{0.881637in}}%
\pgfpathlineto{\pgfqpoint{1.888066in}{0.882118in}}%
\pgfpathlineto{\pgfqpoint{1.887849in}{0.901272in}}%
\pgfpathlineto{\pgfqpoint{1.907800in}{0.901293in}}%
\pgfpathlineto{\pgfqpoint{1.907859in}{0.895543in}}%
\pgfpathlineto{\pgfqpoint{1.927775in}{0.895762in}}%
\pgfpathlineto{\pgfqpoint{1.927950in}{0.876198in}}%
\pgfpathlineto{\pgfqpoint{1.933967in}{0.876483in}}%
\pgfpathlineto{\pgfqpoint{1.941240in}{0.873795in}}%
\pgfpathlineto{\pgfqpoint{1.945167in}{0.874395in}}%
\pgfpathlineto{\pgfqpoint{1.945345in}{0.867212in}}%
\pgfpathlineto{\pgfqpoint{1.951312in}{0.867300in}}%
\pgfpathlineto{\pgfqpoint{1.951519in}{0.850943in}}%
\pgfpathlineto{\pgfqpoint{1.955013in}{0.847952in}}%
\pgfpathlineto{\pgfqpoint{1.955232in}{0.842505in}}%
\pgfpathlineto{\pgfqpoint{1.951736in}{0.838426in}}%
\pgfpathlineto{\pgfqpoint{1.924362in}{0.838176in}}%
\pgfpathlineto{\pgfqpoint{1.924611in}{0.820832in}}%
\pgfpathlineto{\pgfqpoint{1.923831in}{0.809696in}}%
\pgfpathlineto{\pgfqpoint{1.918185in}{0.809351in}}%
\pgfpathclose%
\pgfusepath{fill}%
\end{pgfscope}%
\begin{pgfscope}%
\pgfpathrectangle{\pgfqpoint{0.100000in}{0.100000in}}{\pgfqpoint{3.007045in}{1.925000in}}%
\pgfusepath{clip}%
\pgfsetbuttcap%
\pgfsetmiterjoin%
\definecolor{currentfill}{rgb}{0.585006,0.829343,0.569719}%
\pgfsetfillcolor{currentfill}%
\pgfsetlinewidth{0.000000pt}%
\definecolor{currentstroke}{rgb}{0.000000,0.000000,0.000000}%
\pgfsetstrokecolor{currentstroke}%
\pgfsetstrokeopacity{0.000000}%
\pgfsetdash{}{0pt}%
\pgfpathmoveto{\pgfqpoint{1.311526in}{1.606441in}}%
\pgfpathlineto{\pgfqpoint{1.308128in}{1.577801in}}%
\pgfpathlineto{\pgfqpoint{1.305764in}{1.578064in}}%
\pgfpathlineto{\pgfqpoint{1.303187in}{1.554906in}}%
\pgfpathlineto{\pgfqpoint{1.262489in}{1.559658in}}%
\pgfpathlineto{\pgfqpoint{1.257251in}{1.559802in}}%
\pgfpathlineto{\pgfqpoint{1.246154in}{1.561183in}}%
\pgfpathlineto{\pgfqpoint{1.247672in}{1.573410in}}%
\pgfpathlineto{\pgfqpoint{1.224309in}{1.576380in}}%
\pgfpathlineto{\pgfqpoint{1.227555in}{1.587451in}}%
\pgfpathlineto{\pgfqpoint{1.230291in}{1.608576in}}%
\pgfpathlineto{\pgfqpoint{1.221965in}{1.610407in}}%
\pgfpathlineto{\pgfqpoint{1.222383in}{1.622797in}}%
\pgfpathlineto{\pgfqpoint{1.212290in}{1.629948in}}%
\pgfpathlineto{\pgfqpoint{1.200926in}{1.631495in}}%
\pgfpathlineto{\pgfqpoint{1.201711in}{1.637190in}}%
\pgfpathlineto{\pgfqpoint{1.197757in}{1.637726in}}%
\pgfpathlineto{\pgfqpoint{1.201035in}{1.646271in}}%
\pgfpathlineto{\pgfqpoint{1.195568in}{1.655776in}}%
\pgfpathlineto{\pgfqpoint{1.192292in}{1.666923in}}%
\pgfpathlineto{\pgfqpoint{1.188964in}{1.669306in}}%
\pgfpathlineto{\pgfqpoint{1.190856in}{1.681412in}}%
\pgfpathlineto{\pgfqpoint{1.189998in}{1.686616in}}%
\pgfpathlineto{\pgfqpoint{1.187986in}{1.693328in}}%
\pgfpathlineto{\pgfqpoint{1.240870in}{1.686663in}}%
\pgfpathlineto{\pgfqpoint{1.240737in}{1.685718in}}%
\pgfpathlineto{\pgfqpoint{1.269340in}{1.682158in}}%
\pgfpathlineto{\pgfqpoint{1.267668in}{1.681381in}}%
\pgfpathlineto{\pgfqpoint{1.264869in}{1.658439in}}%
\pgfpathlineto{\pgfqpoint{1.263286in}{1.658626in}}%
\pgfpathlineto{\pgfqpoint{1.260561in}{1.635790in}}%
\pgfpathlineto{\pgfqpoint{1.258817in}{1.635994in}}%
\pgfpathlineto{\pgfqpoint{1.255957in}{1.613050in}}%
\pgfpathlineto{\pgfqpoint{1.311526in}{1.606441in}}%
\pgfpathclose%
\pgfusepath{fill}%
\end{pgfscope}%
\begin{pgfscope}%
\pgfpathrectangle{\pgfqpoint{0.100000in}{0.100000in}}{\pgfqpoint{3.007045in}{1.925000in}}%
\pgfusepath{clip}%
\pgfsetbuttcap%
\pgfsetmiterjoin%
\definecolor{currentfill}{rgb}{0.709097,0.883706,0.683522}%
\pgfsetfillcolor{currentfill}%
\pgfsetlinewidth{0.000000pt}%
\definecolor{currentstroke}{rgb}{0.000000,0.000000,0.000000}%
\pgfsetstrokecolor{currentstroke}%
\pgfsetstrokeopacity{0.000000}%
\pgfsetdash{}{0pt}%
\pgfpathmoveto{\pgfqpoint{2.498935in}{0.687218in}}%
\pgfpathlineto{\pgfqpoint{2.492063in}{0.691931in}}%
\pgfpathlineto{\pgfqpoint{2.477917in}{0.698299in}}%
\pgfpathlineto{\pgfqpoint{2.457652in}{0.694920in}}%
\pgfpathlineto{\pgfqpoint{2.455203in}{0.702082in}}%
\pgfpathlineto{\pgfqpoint{2.447789in}{0.700910in}}%
\pgfpathlineto{\pgfqpoint{2.437033in}{0.704008in}}%
\pgfpathlineto{\pgfqpoint{2.433519in}{0.709598in}}%
\pgfpathlineto{\pgfqpoint{2.429949in}{0.711648in}}%
\pgfpathlineto{\pgfqpoint{2.428514in}{0.718344in}}%
\pgfpathlineto{\pgfqpoint{2.425917in}{0.719822in}}%
\pgfpathlineto{\pgfqpoint{2.427863in}{0.724959in}}%
\pgfpathlineto{\pgfqpoint{2.422656in}{0.729020in}}%
\pgfpathlineto{\pgfqpoint{2.425722in}{0.735586in}}%
\pgfpathlineto{\pgfqpoint{2.432076in}{0.743806in}}%
\pgfpathlineto{\pgfqpoint{2.434230in}{0.742118in}}%
\pgfpathlineto{\pgfqpoint{2.444853in}{0.725072in}}%
\pgfpathlineto{\pgfqpoint{2.451822in}{0.729156in}}%
\pgfpathlineto{\pgfqpoint{2.458380in}{0.737732in}}%
\pgfpathlineto{\pgfqpoint{2.457386in}{0.739684in}}%
\pgfpathlineto{\pgfqpoint{2.460691in}{0.751658in}}%
\pgfpathlineto{\pgfqpoint{2.469526in}{0.751792in}}%
\pgfpathlineto{\pgfqpoint{2.476049in}{0.748944in}}%
\pgfpathlineto{\pgfqpoint{2.475118in}{0.743801in}}%
\pgfpathlineto{\pgfqpoint{2.478955in}{0.739426in}}%
\pgfpathlineto{\pgfqpoint{2.485433in}{0.742115in}}%
\pgfpathlineto{\pgfqpoint{2.489704in}{0.732476in}}%
\pgfpathlineto{\pgfqpoint{2.489650in}{0.715531in}}%
\pgfpathlineto{\pgfqpoint{2.496462in}{0.714530in}}%
\pgfpathlineto{\pgfqpoint{2.499221in}{0.711875in}}%
\pgfpathlineto{\pgfqpoint{2.496078in}{0.707079in}}%
\pgfpathlineto{\pgfqpoint{2.498935in}{0.687218in}}%
\pgfpathclose%
\pgfusepath{fill}%
\end{pgfscope}%
\begin{pgfscope}%
\pgfpathrectangle{\pgfqpoint{0.100000in}{0.100000in}}{\pgfqpoint{3.007045in}{1.925000in}}%
\pgfusepath{clip}%
\pgfsetbuttcap%
\pgfsetmiterjoin%
\definecolor{currentfill}{rgb}{0.629296,0.850012,0.606136}%
\pgfsetfillcolor{currentfill}%
\pgfsetlinewidth{0.000000pt}%
\definecolor{currentstroke}{rgb}{0.000000,0.000000,0.000000}%
\pgfsetstrokecolor{currentstroke}%
\pgfsetstrokeopacity{0.000000}%
\pgfsetdash{}{0pt}%
\pgfpathmoveto{\pgfqpoint{2.741244in}{1.095623in}}%
\pgfpathlineto{\pgfqpoint{2.738779in}{1.094235in}}%
\pgfpathlineto{\pgfqpoint{2.723519in}{1.072852in}}%
\pgfpathlineto{\pgfqpoint{2.723871in}{1.066185in}}%
\pgfpathlineto{\pgfqpoint{2.701345in}{1.062144in}}%
\pgfpathlineto{\pgfqpoint{2.695089in}{1.071478in}}%
\pgfpathlineto{\pgfqpoint{2.716714in}{1.091684in}}%
\pgfpathlineto{\pgfqpoint{2.713742in}{1.098160in}}%
\pgfpathlineto{\pgfqpoint{2.704080in}{1.100551in}}%
\pgfpathlineto{\pgfqpoint{2.711789in}{1.111134in}}%
\pgfpathlineto{\pgfqpoint{2.716713in}{1.110416in}}%
\pgfpathlineto{\pgfqpoint{2.713382in}{1.123521in}}%
\pgfpathlineto{\pgfqpoint{2.724698in}{1.127666in}}%
\pgfpathlineto{\pgfqpoint{2.723470in}{1.137350in}}%
\pgfpathlineto{\pgfqpoint{2.716860in}{1.144435in}}%
\pgfpathlineto{\pgfqpoint{2.720436in}{1.148513in}}%
\pgfpathlineto{\pgfqpoint{2.725113in}{1.146693in}}%
\pgfpathlineto{\pgfqpoint{2.729649in}{1.140417in}}%
\pgfpathlineto{\pgfqpoint{2.735476in}{1.140922in}}%
\pgfpathlineto{\pgfqpoint{2.739950in}{1.134391in}}%
\pgfpathlineto{\pgfqpoint{2.744726in}{1.133904in}}%
\pgfpathlineto{\pgfqpoint{2.747294in}{1.127106in}}%
\pgfpathlineto{\pgfqpoint{2.744962in}{1.123172in}}%
\pgfpathlineto{\pgfqpoint{2.740631in}{1.125917in}}%
\pgfpathlineto{\pgfqpoint{2.742495in}{1.113720in}}%
\pgfpathlineto{\pgfqpoint{2.750051in}{1.107970in}}%
\pgfpathlineto{\pgfqpoint{2.749079in}{1.102041in}}%
\pgfpathlineto{\pgfqpoint{2.745560in}{1.100962in}}%
\pgfpathlineto{\pgfqpoint{2.741244in}{1.095623in}}%
\pgfpathclose%
\pgfusepath{fill}%
\end{pgfscope}%
\begin{pgfscope}%
\pgfpathrectangle{\pgfqpoint{0.100000in}{0.100000in}}{\pgfqpoint{3.007045in}{1.925000in}}%
\pgfusepath{clip}%
\pgfsetbuttcap%
\pgfsetmiterjoin%
\definecolor{currentfill}{rgb}{0.629296,0.850012,0.606136}%
\pgfsetfillcolor{currentfill}%
\pgfsetlinewidth{0.000000pt}%
\definecolor{currentstroke}{rgb}{0.000000,0.000000,0.000000}%
\pgfsetstrokecolor{currentstroke}%
\pgfsetstrokeopacity{0.000000}%
\pgfsetdash{}{0pt}%
\pgfpathmoveto{\pgfqpoint{2.386694in}{0.571228in}}%
\pgfpathlineto{\pgfqpoint{2.386777in}{0.566973in}}%
\pgfpathlineto{\pgfqpoint{2.379860in}{0.568127in}}%
\pgfpathlineto{\pgfqpoint{2.372613in}{0.564198in}}%
\pgfpathlineto{\pgfqpoint{2.358211in}{0.553384in}}%
\pgfpathlineto{\pgfqpoint{2.349321in}{0.550333in}}%
\pgfpathlineto{\pgfqpoint{2.344406in}{0.550144in}}%
\pgfpathlineto{\pgfqpoint{2.336004in}{0.546671in}}%
\pgfpathlineto{\pgfqpoint{2.332928in}{0.547891in}}%
\pgfpathlineto{\pgfqpoint{2.332519in}{0.554903in}}%
\pgfpathlineto{\pgfqpoint{2.324536in}{0.563497in}}%
\pgfpathlineto{\pgfqpoint{2.320470in}{0.564006in}}%
\pgfpathlineto{\pgfqpoint{2.314549in}{0.569015in}}%
\pgfpathlineto{\pgfqpoint{2.302272in}{0.575567in}}%
\pgfpathlineto{\pgfqpoint{2.289539in}{0.581381in}}%
\pgfpathlineto{\pgfqpoint{2.288959in}{0.589362in}}%
\pgfpathlineto{\pgfqpoint{2.296561in}{0.593554in}}%
\pgfpathlineto{\pgfqpoint{2.317482in}{0.595459in}}%
\pgfpathlineto{\pgfqpoint{2.316671in}{0.604111in}}%
\pgfpathlineto{\pgfqpoint{2.322486in}{0.604682in}}%
\pgfpathlineto{\pgfqpoint{2.324742in}{0.580445in}}%
\pgfpathlineto{\pgfqpoint{2.339480in}{0.580648in}}%
\pgfpathlineto{\pgfqpoint{2.340428in}{0.571556in}}%
\pgfpathlineto{\pgfqpoint{2.347979in}{0.570392in}}%
\pgfpathlineto{\pgfqpoint{2.374564in}{0.573201in}}%
\pgfpathlineto{\pgfqpoint{2.386694in}{0.571228in}}%
\pgfpathclose%
\pgfusepath{fill}%
\end{pgfscope}%
\begin{pgfscope}%
\pgfpathrectangle{\pgfqpoint{0.100000in}{0.100000in}}{\pgfqpoint{3.007045in}{1.925000in}}%
\pgfusepath{clip}%
\pgfsetbuttcap%
\pgfsetmiterjoin%
\definecolor{currentfill}{rgb}{0.276078,0.680969,0.375087}%
\pgfsetfillcolor{currentfill}%
\pgfsetlinewidth{0.000000pt}%
\definecolor{currentstroke}{rgb}{0.000000,0.000000,0.000000}%
\pgfsetstrokecolor{currentstroke}%
\pgfsetstrokeopacity{0.000000}%
\pgfsetdash{}{0pt}%
\pgfpathmoveto{\pgfqpoint{2.578534in}{1.060078in}}%
\pgfpathlineto{\pgfqpoint{2.578144in}{1.039001in}}%
\pgfpathlineto{\pgfqpoint{2.561819in}{1.036295in}}%
\pgfpathlineto{\pgfqpoint{2.547798in}{1.034084in}}%
\pgfpathlineto{\pgfqpoint{2.531083in}{1.032312in}}%
\pgfpathlineto{\pgfqpoint{2.529927in}{1.036274in}}%
\pgfpathlineto{\pgfqpoint{2.537514in}{1.041038in}}%
\pgfpathlineto{\pgfqpoint{2.537236in}{1.043373in}}%
\pgfpathlineto{\pgfqpoint{2.543457in}{1.053394in}}%
\pgfpathlineto{\pgfqpoint{2.547424in}{1.055944in}}%
\pgfpathlineto{\pgfqpoint{2.572984in}{1.054602in}}%
\pgfpathlineto{\pgfqpoint{2.578534in}{1.060078in}}%
\pgfpathclose%
\pgfusepath{fill}%
\end{pgfscope}%
\begin{pgfscope}%
\pgfpathrectangle{\pgfqpoint{0.100000in}{0.100000in}}{\pgfqpoint{3.007045in}{1.925000in}}%
\pgfusepath{clip}%
\pgfsetbuttcap%
\pgfsetmiterjoin%
\definecolor{currentfill}{rgb}{0.618224,0.844844,0.597032}%
\pgfsetfillcolor{currentfill}%
\pgfsetlinewidth{0.000000pt}%
\definecolor{currentstroke}{rgb}{0.000000,0.000000,0.000000}%
\pgfsetstrokecolor{currentstroke}%
\pgfsetstrokeopacity{0.000000}%
\pgfsetdash{}{0pt}%
\pgfpathmoveto{\pgfqpoint{1.286524in}{0.867195in}}%
\pgfpathlineto{\pgfqpoint{1.283287in}{0.832832in}}%
\pgfpathlineto{\pgfqpoint{1.271855in}{0.834010in}}%
\pgfpathlineto{\pgfqpoint{1.271248in}{0.828230in}}%
\pgfpathlineto{\pgfqpoint{1.236948in}{0.832126in}}%
\pgfpathlineto{\pgfqpoint{1.237556in}{0.837894in}}%
\pgfpathlineto{\pgfqpoint{1.231742in}{0.838561in}}%
\pgfpathlineto{\pgfqpoint{1.233670in}{0.855551in}}%
\pgfpathlineto{\pgfqpoint{1.210832in}{0.858083in}}%
\pgfpathlineto{\pgfqpoint{1.212799in}{0.875172in}}%
\pgfpathlineto{\pgfqpoint{1.214018in}{0.875058in}}%
\pgfpathlineto{\pgfqpoint{1.217220in}{0.903588in}}%
\pgfpathlineto{\pgfqpoint{1.218495in}{0.915021in}}%
\pgfpathlineto{\pgfqpoint{1.281006in}{0.908273in}}%
\pgfpathlineto{\pgfqpoint{1.280556in}{0.903442in}}%
\pgfpathlineto{\pgfqpoint{1.276783in}{0.868180in}}%
\pgfpathlineto{\pgfqpoint{1.286524in}{0.867195in}}%
\pgfpathclose%
\pgfusepath{fill}%
\end{pgfscope}%
\begin{pgfscope}%
\pgfpathrectangle{\pgfqpoint{0.100000in}{0.100000in}}{\pgfqpoint{3.007045in}{1.925000in}}%
\pgfusepath{clip}%
\pgfsetbuttcap%
\pgfsetmiterjoin%
\definecolor{currentfill}{rgb}{0.662345,0.864022,0.638001}%
\pgfsetfillcolor{currentfill}%
\pgfsetlinewidth{0.000000pt}%
\definecolor{currentstroke}{rgb}{0.000000,0.000000,0.000000}%
\pgfsetstrokecolor{currentstroke}%
\pgfsetstrokeopacity{0.000000}%
\pgfsetdash{}{0pt}%
\pgfpathmoveto{\pgfqpoint{2.959007in}{1.458513in}}%
\pgfpathlineto{\pgfqpoint{2.957672in}{1.458775in}}%
\pgfpathlineto{\pgfqpoint{2.959346in}{1.462818in}}%
\pgfpathlineto{\pgfqpoint{2.955719in}{1.471405in}}%
\pgfpathlineto{\pgfqpoint{2.948844in}{1.463490in}}%
\pgfpathlineto{\pgfqpoint{2.944479in}{1.471321in}}%
\pgfpathlineto{\pgfqpoint{2.933228in}{1.482168in}}%
\pgfpathlineto{\pgfqpoint{2.929528in}{1.490206in}}%
\pgfpathlineto{\pgfqpoint{2.917977in}{1.479575in}}%
\pgfpathlineto{\pgfqpoint{2.916591in}{1.481492in}}%
\pgfpathlineto{\pgfqpoint{2.896973in}{1.475592in}}%
\pgfpathlineto{\pgfqpoint{2.880647in}{1.472875in}}%
\pgfpathlineto{\pgfqpoint{2.878530in}{1.481285in}}%
\pgfpathlineto{\pgfqpoint{2.872055in}{1.481134in}}%
\pgfpathlineto{\pgfqpoint{2.873175in}{1.485821in}}%
\pgfpathlineto{\pgfqpoint{2.867051in}{1.490726in}}%
\pgfpathlineto{\pgfqpoint{2.865008in}{1.506125in}}%
\pgfpathlineto{\pgfqpoint{2.865846in}{1.512168in}}%
\pgfpathlineto{\pgfqpoint{2.862374in}{1.515277in}}%
\pgfpathlineto{\pgfqpoint{2.879160in}{1.518963in}}%
\pgfpathlineto{\pgfqpoint{2.909180in}{1.525824in}}%
\pgfpathlineto{\pgfqpoint{2.914625in}{1.535783in}}%
\pgfpathlineto{\pgfqpoint{2.922960in}{1.542355in}}%
\pgfpathlineto{\pgfqpoint{2.928476in}{1.542999in}}%
\pgfpathlineto{\pgfqpoint{2.934550in}{1.531865in}}%
\pgfpathlineto{\pgfqpoint{2.942818in}{1.530677in}}%
\pgfpathlineto{\pgfqpoint{2.939653in}{1.526607in}}%
\pgfpathlineto{\pgfqpoint{2.932124in}{1.521534in}}%
\pgfpathlineto{\pgfqpoint{2.933679in}{1.519559in}}%
\pgfpathlineto{\pgfqpoint{2.928376in}{1.511191in}}%
\pgfpathlineto{\pgfqpoint{2.930098in}{1.505248in}}%
\pgfpathlineto{\pgfqpoint{2.935189in}{1.506115in}}%
\pgfpathlineto{\pgfqpoint{2.944880in}{1.501863in}}%
\pgfpathlineto{\pgfqpoint{2.950388in}{1.495521in}}%
\pgfpathlineto{\pgfqpoint{2.949426in}{1.488196in}}%
\pgfpathlineto{\pgfqpoint{2.958190in}{1.486046in}}%
\pgfpathlineto{\pgfqpoint{2.959875in}{1.479333in}}%
\pgfpathlineto{\pgfqpoint{2.966716in}{1.476431in}}%
\pgfpathlineto{\pgfqpoint{2.973137in}{1.477451in}}%
\pgfpathlineto{\pgfqpoint{2.988983in}{1.486082in}}%
\pgfpathlineto{\pgfqpoint{2.991233in}{1.478471in}}%
\pgfpathlineto{\pgfqpoint{2.968516in}{1.467337in}}%
\pgfpathlineto{\pgfqpoint{2.967427in}{1.463763in}}%
\pgfpathlineto{\pgfqpoint{2.959007in}{1.458513in}}%
\pgfpathclose%
\pgfusepath{fill}%
\end{pgfscope}%
\begin{pgfscope}%
\pgfpathrectangle{\pgfqpoint{0.100000in}{0.100000in}}{\pgfqpoint{3.007045in}{1.925000in}}%
\pgfusepath{clip}%
\pgfsetbuttcap%
\pgfsetmiterjoin%
\definecolor{currentfill}{rgb}{0.886505,0.956171,0.866128}%
\pgfsetfillcolor{currentfill}%
\pgfsetlinewidth{0.000000pt}%
\definecolor{currentstroke}{rgb}{0.000000,0.000000,0.000000}%
\pgfsetstrokecolor{currentstroke}%
\pgfsetstrokeopacity{0.000000}%
\pgfsetdash{}{0pt}%
\pgfpathmoveto{\pgfqpoint{0.678803in}{0.558557in}}%
\pgfpathlineto{\pgfqpoint{0.679997in}{0.561480in}}%
\pgfpathlineto{\pgfqpoint{0.680052in}{0.564426in}}%
\pgfpathlineto{\pgfqpoint{0.679031in}{0.568728in}}%
\pgfpathlineto{\pgfqpoint{0.679664in}{0.571354in}}%
\pgfpathlineto{\pgfqpoint{0.678267in}{0.576636in}}%
\pgfpathlineto{\pgfqpoint{0.676653in}{0.578838in}}%
\pgfpathlineto{\pgfqpoint{0.671826in}{0.579804in}}%
\pgfpathlineto{\pgfqpoint{0.671241in}{0.581751in}}%
\pgfpathlineto{\pgfqpoint{0.671704in}{0.585285in}}%
\pgfpathlineto{\pgfqpoint{0.673540in}{0.587189in}}%
\pgfpathlineto{\pgfqpoint{0.675922in}{0.587938in}}%
\pgfpathlineto{\pgfqpoint{0.678005in}{0.589551in}}%
\pgfpathlineto{\pgfqpoint{0.678705in}{0.588737in}}%
\pgfpathlineto{\pgfqpoint{0.677833in}{0.586473in}}%
\pgfpathlineto{\pgfqpoint{0.680118in}{0.580300in}}%
\pgfpathlineto{\pgfqpoint{0.681622in}{0.579560in}}%
\pgfpathlineto{\pgfqpoint{0.686508in}{0.578300in}}%
\pgfpathlineto{\pgfqpoint{0.687618in}{0.575930in}}%
\pgfpathlineto{\pgfqpoint{0.687530in}{0.573360in}}%
\pgfpathlineto{\pgfqpoint{0.685839in}{0.571449in}}%
\pgfpathlineto{\pgfqpoint{0.687044in}{0.567758in}}%
\pgfpathlineto{\pgfqpoint{0.687213in}{0.565076in}}%
\pgfpathlineto{\pgfqpoint{0.688961in}{0.561585in}}%
\pgfpathlineto{\pgfqpoint{0.691199in}{0.560100in}}%
\pgfpathlineto{\pgfqpoint{0.692475in}{0.557535in}}%
\pgfpathlineto{\pgfqpoint{0.691676in}{0.556369in}}%
\pgfpathlineto{\pgfqpoint{0.688983in}{0.555610in}}%
\pgfpathlineto{\pgfqpoint{0.687659in}{0.557811in}}%
\pgfpathlineto{\pgfqpoint{0.685611in}{0.559618in}}%
\pgfpathlineto{\pgfqpoint{0.682975in}{0.560083in}}%
\pgfpathlineto{\pgfqpoint{0.678803in}{0.558557in}}%
\pgfpathclose%
\pgfusepath{fill}%
\end{pgfscope}%
\begin{pgfscope}%
\pgfpathrectangle{\pgfqpoint{0.100000in}{0.100000in}}{\pgfqpoint{3.007045in}{1.925000in}}%
\pgfusepath{clip}%
\pgfsetbuttcap%
\pgfsetmiterjoin%
\definecolor{currentfill}{rgb}{0.886505,0.956171,0.866128}%
\pgfsetfillcolor{currentfill}%
\pgfsetlinewidth{0.000000pt}%
\definecolor{currentstroke}{rgb}{0.000000,0.000000,0.000000}%
\pgfsetstrokecolor{currentstroke}%
\pgfsetstrokeopacity{0.000000}%
\pgfsetdash{}{0pt}%
\pgfpathmoveto{\pgfqpoint{0.784268in}{0.583082in}}%
\pgfpathlineto{\pgfqpoint{0.782998in}{0.580535in}}%
\pgfpathlineto{\pgfqpoint{0.776856in}{0.574080in}}%
\pgfpathlineto{\pgfqpoint{0.777383in}{0.573526in}}%
\pgfpathlineto{\pgfqpoint{0.771216in}{0.567098in}}%
\pgfpathlineto{\pgfqpoint{0.773326in}{0.565108in}}%
\pgfpathlineto{\pgfqpoint{0.770248in}{0.561865in}}%
\pgfpathlineto{\pgfqpoint{0.771871in}{0.560318in}}%
\pgfpathlineto{\pgfqpoint{0.768863in}{0.557118in}}%
\pgfpathlineto{\pgfqpoint{0.775661in}{0.550783in}}%
\pgfpathlineto{\pgfqpoint{0.790016in}{0.537796in}}%
\pgfpathlineto{\pgfqpoint{0.799988in}{0.528978in}}%
\pgfpathlineto{\pgfqpoint{0.801466in}{0.530632in}}%
\pgfpathlineto{\pgfqpoint{0.803129in}{0.529197in}}%
\pgfpathlineto{\pgfqpoint{0.804798in}{0.527778in}}%
\pgfpathlineto{\pgfqpoint{0.800473in}{0.522699in}}%
\pgfpathlineto{\pgfqpoint{0.798922in}{0.524065in}}%
\pgfpathlineto{\pgfqpoint{0.793312in}{0.517519in}}%
\pgfpathlineto{\pgfqpoint{0.791670in}{0.518914in}}%
\pgfpathlineto{\pgfqpoint{0.788832in}{0.515575in}}%
\pgfpathlineto{\pgfqpoint{0.787134in}{0.516986in}}%
\pgfpathlineto{\pgfqpoint{0.775774in}{0.503541in}}%
\pgfpathlineto{\pgfqpoint{0.777669in}{0.501945in}}%
\pgfpathlineto{\pgfqpoint{0.773225in}{0.497185in}}%
\pgfpathlineto{\pgfqpoint{0.771756in}{0.498478in}}%
\pgfpathlineto{\pgfqpoint{0.770286in}{0.496757in}}%
\pgfpathlineto{\pgfqpoint{0.768805in}{0.498004in}}%
\pgfpathlineto{\pgfqpoint{0.767409in}{0.496302in}}%
\pgfpathlineto{\pgfqpoint{0.764086in}{0.499185in}}%
\pgfpathlineto{\pgfqpoint{0.762810in}{0.497757in}}%
\pgfpathlineto{\pgfqpoint{0.761153in}{0.499198in}}%
\pgfpathlineto{\pgfqpoint{0.756937in}{0.494067in}}%
\pgfpathlineto{\pgfqpoint{0.755431in}{0.495345in}}%
\pgfpathlineto{\pgfqpoint{0.751165in}{0.490590in}}%
\pgfpathlineto{\pgfqpoint{0.741587in}{0.499176in}}%
\pgfpathlineto{\pgfqpoint{0.742883in}{0.500799in}}%
\pgfpathlineto{\pgfqpoint{0.741434in}{0.502134in}}%
\pgfpathlineto{\pgfqpoint{0.742884in}{0.503844in}}%
\pgfpathlineto{\pgfqpoint{0.741312in}{0.505279in}}%
\pgfpathlineto{\pgfqpoint{0.742738in}{0.506896in}}%
\pgfpathlineto{\pgfqpoint{0.741050in}{0.508543in}}%
\pgfpathlineto{\pgfqpoint{0.746586in}{0.510216in}}%
\pgfpathlineto{\pgfqpoint{0.746712in}{0.512326in}}%
\pgfpathlineto{\pgfqpoint{0.745027in}{0.513729in}}%
\pgfpathlineto{\pgfqpoint{0.746433in}{0.515178in}}%
\pgfpathlineto{\pgfqpoint{0.748958in}{0.513680in}}%
\pgfpathlineto{\pgfqpoint{0.748481in}{0.510045in}}%
\pgfpathlineto{\pgfqpoint{0.750210in}{0.505718in}}%
\pgfpathlineto{\pgfqpoint{0.755351in}{0.502374in}}%
\pgfpathlineto{\pgfqpoint{0.757588in}{0.502208in}}%
\pgfpathlineto{\pgfqpoint{0.761221in}{0.502993in}}%
\pgfpathlineto{\pgfqpoint{0.764189in}{0.504114in}}%
\pgfpathlineto{\pgfqpoint{0.765417in}{0.505253in}}%
\pgfpathlineto{\pgfqpoint{0.767984in}{0.511139in}}%
\pgfpathlineto{\pgfqpoint{0.770424in}{0.514447in}}%
\pgfpathlineto{\pgfqpoint{0.771090in}{0.519688in}}%
\pgfpathlineto{\pgfqpoint{0.769265in}{0.521681in}}%
\pgfpathlineto{\pgfqpoint{0.771550in}{0.523004in}}%
\pgfpathlineto{\pgfqpoint{0.773952in}{0.519957in}}%
\pgfpathlineto{\pgfqpoint{0.778165in}{0.520246in}}%
\pgfpathlineto{\pgfqpoint{0.780075in}{0.522157in}}%
\pgfpathlineto{\pgfqpoint{0.781283in}{0.524482in}}%
\pgfpathlineto{\pgfqpoint{0.780153in}{0.528966in}}%
\pgfpathlineto{\pgfqpoint{0.776623in}{0.527750in}}%
\pgfpathlineto{\pgfqpoint{0.774378in}{0.527915in}}%
\pgfpathlineto{\pgfqpoint{0.773486in}{0.530158in}}%
\pgfpathlineto{\pgfqpoint{0.770430in}{0.530368in}}%
\pgfpathlineto{\pgfqpoint{0.767348in}{0.531211in}}%
\pgfpathlineto{\pgfqpoint{0.764853in}{0.531128in}}%
\pgfpathlineto{\pgfqpoint{0.757318in}{0.529341in}}%
\pgfpathlineto{\pgfqpoint{0.759189in}{0.532678in}}%
\pgfpathlineto{\pgfqpoint{0.759643in}{0.537554in}}%
\pgfpathlineto{\pgfqpoint{0.757195in}{0.537867in}}%
\pgfpathlineto{\pgfqpoint{0.757549in}{0.535521in}}%
\pgfpathlineto{\pgfqpoint{0.755714in}{0.534840in}}%
\pgfpathlineto{\pgfqpoint{0.755754in}{0.537318in}}%
\pgfpathlineto{\pgfqpoint{0.754836in}{0.539913in}}%
\pgfpathlineto{\pgfqpoint{0.749644in}{0.544946in}}%
\pgfpathlineto{\pgfqpoint{0.743760in}{0.546795in}}%
\pgfpathlineto{\pgfqpoint{0.742220in}{0.547611in}}%
\pgfpathlineto{\pgfqpoint{0.735804in}{0.558236in}}%
\pgfpathlineto{\pgfqpoint{0.734630in}{0.560668in}}%
\pgfpathlineto{\pgfqpoint{0.734596in}{0.562774in}}%
\pgfpathlineto{\pgfqpoint{0.736774in}{0.565731in}}%
\pgfpathlineto{\pgfqpoint{0.738729in}{0.567328in}}%
\pgfpathlineto{\pgfqpoint{0.739066in}{0.574040in}}%
\pgfpathlineto{\pgfqpoint{0.742356in}{0.572549in}}%
\pgfpathlineto{\pgfqpoint{0.746458in}{0.574359in}}%
\pgfpathlineto{\pgfqpoint{0.742466in}{0.581233in}}%
\pgfpathlineto{\pgfqpoint{0.740777in}{0.583420in}}%
\pgfpathlineto{\pgfqpoint{0.739815in}{0.588701in}}%
\pgfpathlineto{\pgfqpoint{0.738785in}{0.592111in}}%
\pgfpathlineto{\pgfqpoint{0.740864in}{0.594187in}}%
\pgfpathlineto{\pgfqpoint{0.750011in}{0.591581in}}%
\pgfpathlineto{\pgfqpoint{0.750589in}{0.589917in}}%
\pgfpathlineto{\pgfqpoint{0.753942in}{0.589343in}}%
\pgfpathlineto{\pgfqpoint{0.753820in}{0.591119in}}%
\pgfpathlineto{\pgfqpoint{0.762748in}{0.589868in}}%
\pgfpathlineto{\pgfqpoint{0.764937in}{0.585166in}}%
\pgfpathlineto{\pgfqpoint{0.766490in}{0.583615in}}%
\pgfpathlineto{\pgfqpoint{0.768297in}{0.583840in}}%
\pgfpathlineto{\pgfqpoint{0.767184in}{0.586211in}}%
\pgfpathlineto{\pgfqpoint{0.770001in}{0.588417in}}%
\pgfpathlineto{\pgfqpoint{0.777196in}{0.586185in}}%
\pgfpathlineto{\pgfqpoint{0.781761in}{0.584481in}}%
\pgfpathlineto{\pgfqpoint{0.784268in}{0.583082in}}%
\pgfpathclose%
\pgfusepath{fill}%
\end{pgfscope}%
\begin{pgfscope}%
\pgfpathrectangle{\pgfqpoint{0.100000in}{0.100000in}}{\pgfqpoint{3.007045in}{1.925000in}}%
\pgfusepath{clip}%
\pgfsetbuttcap%
\pgfsetmiterjoin%
\definecolor{currentfill}{rgb}{0.727797,0.891580,0.701730}%
\pgfsetfillcolor{currentfill}%
\pgfsetlinewidth{0.000000pt}%
\definecolor{currentstroke}{rgb}{0.000000,0.000000,0.000000}%
\pgfsetstrokecolor{currentstroke}%
\pgfsetstrokeopacity{0.000000}%
\pgfsetdash{}{0pt}%
\pgfpathmoveto{\pgfqpoint{1.521452in}{1.403789in}}%
\pgfpathlineto{\pgfqpoint{1.471839in}{1.406868in}}%
\pgfpathlineto{\pgfqpoint{1.473591in}{1.432656in}}%
\pgfpathlineto{\pgfqpoint{1.475452in}{1.459167in}}%
\pgfpathlineto{\pgfqpoint{1.483456in}{1.461953in}}%
\pgfpathlineto{\pgfqpoint{1.490267in}{1.460196in}}%
\pgfpathlineto{\pgfqpoint{1.496174in}{1.453745in}}%
\pgfpathlineto{\pgfqpoint{1.506479in}{1.454584in}}%
\pgfpathlineto{\pgfqpoint{1.511045in}{1.452183in}}%
\pgfpathlineto{\pgfqpoint{1.517464in}{1.451554in}}%
\pgfpathlineto{\pgfqpoint{1.522608in}{1.451050in}}%
\pgfpathlineto{\pgfqpoint{1.522496in}{1.435436in}}%
\pgfpathlineto{\pgfqpoint{1.521204in}{1.412401in}}%
\pgfpathlineto{\pgfqpoint{1.521452in}{1.403789in}}%
\pgfpathclose%
\pgfusepath{fill}%
\end{pgfscope}%
\begin{pgfscope}%
\pgfpathrectangle{\pgfqpoint{0.100000in}{0.100000in}}{\pgfqpoint{3.007045in}{1.925000in}}%
\pgfusepath{clip}%
\pgfsetbuttcap%
\pgfsetmiterjoin%
\definecolor{currentfill}{rgb}{0.479815,0.780254,0.483230}%
\pgfsetfillcolor{currentfill}%
\pgfsetlinewidth{0.000000pt}%
\definecolor{currentstroke}{rgb}{0.000000,0.000000,0.000000}%
\pgfsetstrokecolor{currentstroke}%
\pgfsetstrokeopacity{0.000000}%
\pgfsetdash{}{0pt}%
\pgfpathmoveto{\pgfqpoint{0.763180in}{1.655715in}}%
\pgfpathlineto{\pgfqpoint{0.759381in}{1.652797in}}%
\pgfpathlineto{\pgfqpoint{0.758646in}{1.647388in}}%
\pgfpathlineto{\pgfqpoint{0.754151in}{1.641671in}}%
\pgfpathlineto{\pgfqpoint{0.746366in}{1.637288in}}%
\pgfpathlineto{\pgfqpoint{0.738581in}{1.626973in}}%
\pgfpathlineto{\pgfqpoint{0.733470in}{1.622014in}}%
\pgfpathlineto{\pgfqpoint{0.730914in}{1.610662in}}%
\pgfpathlineto{\pgfqpoint{0.718698in}{1.613732in}}%
\pgfpathlineto{\pgfqpoint{0.720108in}{1.619262in}}%
\pgfpathlineto{\pgfqpoint{0.716216in}{1.624121in}}%
\pgfpathlineto{\pgfqpoint{0.698781in}{1.628508in}}%
\pgfpathlineto{\pgfqpoint{0.695577in}{1.627414in}}%
\pgfpathlineto{\pgfqpoint{0.687625in}{1.619603in}}%
\pgfpathlineto{\pgfqpoint{0.683738in}{1.619619in}}%
\pgfpathlineto{\pgfqpoint{0.671575in}{1.622748in}}%
\pgfpathlineto{\pgfqpoint{0.677402in}{1.629304in}}%
\pgfpathlineto{\pgfqpoint{0.678278in}{1.634204in}}%
\pgfpathlineto{\pgfqpoint{0.685898in}{1.641698in}}%
\pgfpathlineto{\pgfqpoint{0.684947in}{1.644999in}}%
\pgfpathlineto{\pgfqpoint{0.678193in}{1.649100in}}%
\pgfpathlineto{\pgfqpoint{0.678258in}{1.651609in}}%
\pgfpathlineto{\pgfqpoint{0.689599in}{1.651922in}}%
\pgfpathlineto{\pgfqpoint{0.689503in}{1.657799in}}%
\pgfpathlineto{\pgfqpoint{0.694447in}{1.659187in}}%
\pgfpathlineto{\pgfqpoint{0.695518in}{1.663103in}}%
\pgfpathlineto{\pgfqpoint{0.690132in}{1.668259in}}%
\pgfpathlineto{\pgfqpoint{0.683125in}{1.670416in}}%
\pgfpathlineto{\pgfqpoint{0.685286in}{1.683617in}}%
\pgfpathlineto{\pgfqpoint{0.680362in}{1.684919in}}%
\pgfpathlineto{\pgfqpoint{0.681006in}{1.694917in}}%
\pgfpathlineto{\pgfqpoint{0.693304in}{1.692341in}}%
\pgfpathlineto{\pgfqpoint{0.694542in}{1.697073in}}%
\pgfpathlineto{\pgfqpoint{0.704883in}{1.694237in}}%
\pgfpathlineto{\pgfqpoint{0.709253in}{1.696086in}}%
\pgfpathlineto{\pgfqpoint{0.712881in}{1.710060in}}%
\pgfpathlineto{\pgfqpoint{0.715135in}{1.709465in}}%
\pgfpathlineto{\pgfqpoint{0.718256in}{1.717477in}}%
\pgfpathlineto{\pgfqpoint{0.722204in}{1.719436in}}%
\pgfpathlineto{\pgfqpoint{0.732306in}{1.716868in}}%
\pgfpathlineto{\pgfqpoint{0.727679in}{1.706314in}}%
\pgfpathlineto{\pgfqpoint{0.729208in}{1.700974in}}%
\pgfpathlineto{\pgfqpoint{0.727519in}{1.694386in}}%
\pgfpathlineto{\pgfqpoint{0.728377in}{1.682025in}}%
\pgfpathlineto{\pgfqpoint{0.732751in}{1.676168in}}%
\pgfpathlineto{\pgfqpoint{0.733120in}{1.669073in}}%
\pgfpathlineto{\pgfqpoint{0.742633in}{1.666720in}}%
\pgfpathlineto{\pgfqpoint{0.741293in}{1.661247in}}%
\pgfpathlineto{\pgfqpoint{0.763180in}{1.655715in}}%
\pgfpathclose%
\pgfusepath{fill}%
\end{pgfscope}%
\begin{pgfscope}%
\pgfpathrectangle{\pgfqpoint{0.100000in}{0.100000in}}{\pgfqpoint{3.007045in}{1.925000in}}%
\pgfusepath{clip}%
\pgfsetbuttcap%
\pgfsetmiterjoin%
\definecolor{currentfill}{rgb}{0.451765,0.767090,0.461207}%
\pgfsetfillcolor{currentfill}%
\pgfsetlinewidth{0.000000pt}%
\definecolor{currentstroke}{rgb}{0.000000,0.000000,0.000000}%
\pgfsetstrokecolor{currentstroke}%
\pgfsetstrokeopacity{0.000000}%
\pgfsetdash{}{0pt}%
\pgfpathmoveto{\pgfqpoint{1.509458in}{1.149192in}}%
\pgfpathlineto{\pgfqpoint{1.507527in}{1.120611in}}%
\pgfpathlineto{\pgfqpoint{1.480135in}{1.122427in}}%
\pgfpathlineto{\pgfqpoint{1.439278in}{1.125169in}}%
\pgfpathlineto{\pgfqpoint{1.441646in}{1.153706in}}%
\pgfpathlineto{\pgfqpoint{1.446054in}{1.153418in}}%
\pgfpathlineto{\pgfqpoint{1.448195in}{1.181868in}}%
\pgfpathlineto{\pgfqpoint{1.509436in}{1.177873in}}%
\pgfpathlineto{\pgfqpoint{1.510385in}{1.177845in}}%
\pgfpathlineto{\pgfqpoint{1.508618in}{1.149222in}}%
\pgfpathlineto{\pgfqpoint{1.509458in}{1.149192in}}%
\pgfpathclose%
\pgfusepath{fill}%
\end{pgfscope}%
\begin{pgfscope}%
\pgfpathrectangle{\pgfqpoint{0.100000in}{0.100000in}}{\pgfqpoint{3.007045in}{1.925000in}}%
\pgfusepath{clip}%
\pgfsetbuttcap%
\pgfsetmiterjoin%
\definecolor{currentfill}{rgb}{0.551788,0.813841,0.542407}%
\pgfsetfillcolor{currentfill}%
\pgfsetlinewidth{0.000000pt}%
\definecolor{currentstroke}{rgb}{0.000000,0.000000,0.000000}%
\pgfsetstrokecolor{currentstroke}%
\pgfsetstrokeopacity{0.000000}%
\pgfsetdash{}{0pt}%
\pgfpathmoveto{\pgfqpoint{2.070604in}{1.165647in}}%
\pgfpathlineto{\pgfqpoint{2.050719in}{1.164389in}}%
\pgfpathlineto{\pgfqpoint{2.049876in}{1.175940in}}%
\pgfpathlineto{\pgfqpoint{2.029984in}{1.174625in}}%
\pgfpathlineto{\pgfqpoint{2.025601in}{1.191743in}}%
\pgfpathlineto{\pgfqpoint{2.024381in}{1.212989in}}%
\pgfpathlineto{\pgfqpoint{2.025473in}{1.215274in}}%
\pgfpathlineto{\pgfqpoint{2.032407in}{1.214497in}}%
\pgfpathlineto{\pgfqpoint{2.039739in}{1.216849in}}%
\pgfpathlineto{\pgfqpoint{2.044100in}{1.215046in}}%
\pgfpathlineto{\pgfqpoint{2.043262in}{1.228036in}}%
\pgfpathlineto{\pgfqpoint{2.060239in}{1.229418in}}%
\pgfpathlineto{\pgfqpoint{2.066190in}{1.226950in}}%
\pgfpathlineto{\pgfqpoint{2.067386in}{1.211625in}}%
\pgfpathlineto{\pgfqpoint{2.068009in}{1.203031in}}%
\pgfpathlineto{\pgfqpoint{2.064257in}{1.202733in}}%
\pgfpathlineto{\pgfqpoint{2.064705in}{1.195908in}}%
\pgfpathlineto{\pgfqpoint{2.068581in}{1.195334in}}%
\pgfpathlineto{\pgfqpoint{2.069337in}{1.185778in}}%
\pgfpathlineto{\pgfqpoint{2.075079in}{1.186079in}}%
\pgfpathlineto{\pgfqpoint{2.076427in}{1.165836in}}%
\pgfpathlineto{\pgfqpoint{2.070604in}{1.165647in}}%
\pgfpathclose%
\pgfusepath{fill}%
\end{pgfscope}%
\begin{pgfscope}%
\pgfpathrectangle{\pgfqpoint{0.100000in}{0.100000in}}{\pgfqpoint{3.007045in}{1.925000in}}%
\pgfusepath{clip}%
\pgfsetbuttcap%
\pgfsetmiterjoin%
\definecolor{currentfill}{rgb}{0.643645,0.856148,0.619792}%
\pgfsetfillcolor{currentfill}%
\pgfsetlinewidth{0.000000pt}%
\definecolor{currentstroke}{rgb}{0.000000,0.000000,0.000000}%
\pgfsetstrokecolor{currentstroke}%
\pgfsetstrokeopacity{0.000000}%
\pgfsetdash{}{0pt}%
\pgfpathmoveto{\pgfqpoint{2.803517in}{1.405975in}}%
\pgfpathlineto{\pgfqpoint{2.794033in}{1.391447in}}%
\pgfpathlineto{\pgfqpoint{2.768627in}{1.400117in}}%
\pgfpathlineto{\pgfqpoint{2.765754in}{1.403877in}}%
\pgfpathlineto{\pgfqpoint{2.760698in}{1.403338in}}%
\pgfpathlineto{\pgfqpoint{2.752320in}{1.406665in}}%
\pgfpathlineto{\pgfqpoint{2.747986in}{1.413389in}}%
\pgfpathlineto{\pgfqpoint{2.745672in}{1.421591in}}%
\pgfpathlineto{\pgfqpoint{2.739778in}{1.426983in}}%
\pgfpathlineto{\pgfqpoint{2.768356in}{1.454967in}}%
\pgfpathlineto{\pgfqpoint{2.776040in}{1.453025in}}%
\pgfpathlineto{\pgfqpoint{2.787399in}{1.454470in}}%
\pgfpathlineto{\pgfqpoint{2.787830in}{1.459571in}}%
\pgfpathlineto{\pgfqpoint{2.794361in}{1.458299in}}%
\pgfpathlineto{\pgfqpoint{2.795846in}{1.455210in}}%
\pgfpathlineto{\pgfqpoint{2.806124in}{1.452725in}}%
\pgfpathlineto{\pgfqpoint{2.815310in}{1.453061in}}%
\pgfpathlineto{\pgfqpoint{2.815641in}{1.458005in}}%
\pgfpathlineto{\pgfqpoint{2.816097in}{1.458112in}}%
\pgfpathlineto{\pgfqpoint{2.822125in}{1.424128in}}%
\pgfpathlineto{\pgfqpoint{2.803333in}{1.417011in}}%
\pgfpathlineto{\pgfqpoint{2.803517in}{1.405975in}}%
\pgfpathclose%
\pgfusepath{fill}%
\end{pgfscope}%
\begin{pgfscope}%
\pgfpathrectangle{\pgfqpoint{0.100000in}{0.100000in}}{\pgfqpoint{3.007045in}{1.925000in}}%
\pgfusepath{clip}%
\pgfsetbuttcap%
\pgfsetmiterjoin%
\definecolor{currentfill}{rgb}{0.363922,0.724029,0.418147}%
\pgfsetfillcolor{currentfill}%
\pgfsetlinewidth{0.000000pt}%
\definecolor{currentstroke}{rgb}{0.000000,0.000000,0.000000}%
\pgfsetstrokecolor{currentstroke}%
\pgfsetstrokeopacity{0.000000}%
\pgfsetdash{}{0pt}%
\pgfpathmoveto{\pgfqpoint{1.219256in}{0.700832in}}%
\pgfpathlineto{\pgfqpoint{1.215312in}{0.701261in}}%
\pgfpathlineto{\pgfqpoint{1.186737in}{0.704479in}}%
\pgfpathlineto{\pgfqpoint{1.133803in}{0.710900in}}%
\pgfpathlineto{\pgfqpoint{1.141737in}{0.773832in}}%
\pgfpathlineto{\pgfqpoint{1.144512in}{0.779815in}}%
\pgfpathlineto{\pgfqpoint{1.147002in}{0.802056in}}%
\pgfpathlineto{\pgfqpoint{1.146247in}{0.808107in}}%
\pgfpathlineto{\pgfqpoint{1.147690in}{0.819949in}}%
\pgfpathlineto{\pgfqpoint{1.165157in}{0.817109in}}%
\pgfpathlineto{\pgfqpoint{1.166565in}{0.828636in}}%
\pgfpathlineto{\pgfqpoint{1.173550in}{0.827766in}}%
\pgfpathlineto{\pgfqpoint{1.177035in}{0.856282in}}%
\pgfpathlineto{\pgfqpoint{1.210171in}{0.852394in}}%
\pgfpathlineto{\pgfqpoint{1.210832in}{0.858083in}}%
\pgfpathlineto{\pgfqpoint{1.233670in}{0.855551in}}%
\pgfpathlineto{\pgfqpoint{1.231742in}{0.838561in}}%
\pgfpathlineto{\pgfqpoint{1.226820in}{0.793144in}}%
\pgfpathlineto{\pgfqpoint{1.224219in}{0.776176in}}%
\pgfpathlineto{\pgfqpoint{1.201518in}{0.778286in}}%
\pgfpathlineto{\pgfqpoint{1.200282in}{0.767082in}}%
\pgfpathlineto{\pgfqpoint{1.198272in}{0.767309in}}%
\pgfpathlineto{\pgfqpoint{1.194858in}{0.738140in}}%
\pgfpathlineto{\pgfqpoint{1.222778in}{0.735139in}}%
\pgfpathlineto{\pgfqpoint{1.219256in}{0.700832in}}%
\pgfpathclose%
\pgfusepath{fill}%
\end{pgfscope}%
\begin{pgfscope}%
\pgfpathrectangle{\pgfqpoint{0.100000in}{0.100000in}}{\pgfqpoint{3.007045in}{1.925000in}}%
\pgfusepath{clip}%
\pgfsetbuttcap%
\pgfsetmiterjoin%
\definecolor{currentfill}{rgb}{0.629296,0.850012,0.606136}%
\pgfsetfillcolor{currentfill}%
\pgfsetlinewidth{0.000000pt}%
\definecolor{currentstroke}{rgb}{0.000000,0.000000,0.000000}%
\pgfsetstrokecolor{currentstroke}%
\pgfsetstrokeopacity{0.000000}%
\pgfsetdash{}{0pt}%
\pgfpathmoveto{\pgfqpoint{2.394856in}{0.702917in}}%
\pgfpathlineto{\pgfqpoint{2.356903in}{0.698990in}}%
\pgfpathlineto{\pgfqpoint{2.339099in}{0.697275in}}%
\pgfpathlineto{\pgfqpoint{2.331743in}{0.699866in}}%
\pgfpathlineto{\pgfqpoint{2.330327in}{0.705545in}}%
\pgfpathlineto{\pgfqpoint{2.329526in}{0.710196in}}%
\pgfpathlineto{\pgfqpoint{2.336528in}{0.717430in}}%
\pgfpathlineto{\pgfqpoint{2.351237in}{0.719223in}}%
\pgfpathlineto{\pgfqpoint{2.349352in}{0.736500in}}%
\pgfpathlineto{\pgfqpoint{2.350162in}{0.739256in}}%
\pgfpathlineto{\pgfqpoint{2.360665in}{0.742287in}}%
\pgfpathlineto{\pgfqpoint{2.363559in}{0.734763in}}%
\pgfpathlineto{\pgfqpoint{2.368826in}{0.730586in}}%
\pgfpathlineto{\pgfqpoint{2.375605in}{0.733287in}}%
\pgfpathlineto{\pgfqpoint{2.378977in}{0.740714in}}%
\pgfpathlineto{\pgfqpoint{2.382677in}{0.742228in}}%
\pgfpathlineto{\pgfqpoint{2.392692in}{0.741268in}}%
\pgfpathlineto{\pgfqpoint{2.395370in}{0.738234in}}%
\pgfpathlineto{\pgfqpoint{2.395922in}{0.728401in}}%
\pgfpathlineto{\pgfqpoint{2.389930in}{0.724977in}}%
\pgfpathlineto{\pgfqpoint{2.386745in}{0.718055in}}%
\pgfpathlineto{\pgfqpoint{2.389600in}{0.714570in}}%
\pgfpathlineto{\pgfqpoint{2.394856in}{0.702917in}}%
\pgfpathclose%
\pgfusepath{fill}%
\end{pgfscope}%
\begin{pgfscope}%
\pgfpathrectangle{\pgfqpoint{0.100000in}{0.100000in}}{\pgfqpoint{3.007045in}{1.925000in}}%
\pgfusepath{clip}%
\pgfsetbuttcap%
\pgfsetmiterjoin%
\definecolor{currentfill}{rgb}{0.112495,0.523875,0.252903}%
\pgfsetfillcolor{currentfill}%
\pgfsetlinewidth{0.000000pt}%
\definecolor{currentstroke}{rgb}{0.000000,0.000000,0.000000}%
\pgfsetstrokecolor{currentstroke}%
\pgfsetstrokeopacity{0.000000}%
\pgfsetdash{}{0pt}%
\pgfpathmoveto{\pgfqpoint{2.313368in}{1.507664in}}%
\pgfpathlineto{\pgfqpoint{2.298344in}{1.505689in}}%
\pgfpathlineto{\pgfqpoint{2.291082in}{1.562710in}}%
\pgfpathlineto{\pgfqpoint{2.290081in}{1.574045in}}%
\pgfpathlineto{\pgfqpoint{2.273175in}{1.571785in}}%
\pgfpathlineto{\pgfqpoint{2.269589in}{1.599880in}}%
\pgfpathlineto{\pgfqpoint{2.271744in}{1.600141in}}%
\pgfpathlineto{\pgfqpoint{2.277048in}{1.592666in}}%
\pgfpathlineto{\pgfqpoint{2.285970in}{1.592606in}}%
\pgfpathlineto{\pgfqpoint{2.292882in}{1.588599in}}%
\pgfpathlineto{\pgfqpoint{2.308293in}{1.585034in}}%
\pgfpathlineto{\pgfqpoint{2.311509in}{1.579156in}}%
\pgfpathlineto{\pgfqpoint{2.316771in}{1.573334in}}%
\pgfpathlineto{\pgfqpoint{2.320179in}{1.565731in}}%
\pgfpathlineto{\pgfqpoint{2.311941in}{1.567027in}}%
\pgfpathlineto{\pgfqpoint{2.313141in}{1.558957in}}%
\pgfpathlineto{\pgfqpoint{2.319127in}{1.556401in}}%
\pgfpathlineto{\pgfqpoint{2.322517in}{1.545363in}}%
\pgfpathlineto{\pgfqpoint{2.321374in}{1.538390in}}%
\pgfpathlineto{\pgfqpoint{2.322792in}{1.520480in}}%
\pgfpathlineto{\pgfqpoint{2.318985in}{1.515613in}}%
\pgfpathlineto{\pgfqpoint{2.314903in}{1.514937in}}%
\pgfpathlineto{\pgfqpoint{2.313368in}{1.507664in}}%
\pgfpathclose%
\pgfusepath{fill}%
\end{pgfscope}%
\begin{pgfscope}%
\pgfpathrectangle{\pgfqpoint{0.100000in}{0.100000in}}{\pgfqpoint{3.007045in}{1.925000in}}%
\pgfusepath{clip}%
\pgfsetbuttcap%
\pgfsetmiterjoin%
\definecolor{currentfill}{rgb}{0.573933,0.824175,0.560615}%
\pgfsetfillcolor{currentfill}%
\pgfsetlinewidth{0.000000pt}%
\definecolor{currentstroke}{rgb}{0.000000,0.000000,0.000000}%
\pgfsetstrokecolor{currentstroke}%
\pgfsetstrokeopacity{0.000000}%
\pgfsetdash{}{0pt}%
\pgfpathmoveto{\pgfqpoint{2.327732in}{1.268420in}}%
\pgfpathlineto{\pgfqpoint{2.322114in}{1.267499in}}%
\pgfpathlineto{\pgfqpoint{2.322524in}{1.264076in}}%
\pgfpathlineto{\pgfqpoint{2.305349in}{1.261913in}}%
\pgfpathlineto{\pgfqpoint{2.306160in}{1.255224in}}%
\pgfpathlineto{\pgfqpoint{2.301562in}{1.253044in}}%
\pgfpathlineto{\pgfqpoint{2.283146in}{1.250837in}}%
\pgfpathlineto{\pgfqpoint{2.278491in}{1.292509in}}%
\pgfpathlineto{\pgfqpoint{2.301283in}{1.295237in}}%
\pgfpathlineto{\pgfqpoint{2.299961in}{1.306677in}}%
\pgfpathlineto{\pgfqpoint{2.305531in}{1.307364in}}%
\pgfpathlineto{\pgfqpoint{2.322615in}{1.309553in}}%
\pgfpathlineto{\pgfqpoint{2.327732in}{1.268420in}}%
\pgfpathclose%
\pgfusepath{fill}%
\end{pgfscope}%
\begin{pgfscope}%
\pgfpathrectangle{\pgfqpoint{0.100000in}{0.100000in}}{\pgfqpoint{3.007045in}{1.925000in}}%
\pgfusepath{clip}%
\pgfsetbuttcap%
\pgfsetmiterjoin%
\definecolor{currentfill}{rgb}{0.463206,0.772503,0.469573}%
\pgfsetfillcolor{currentfill}%
\pgfsetlinewidth{0.000000pt}%
\definecolor{currentstroke}{rgb}{0.000000,0.000000,0.000000}%
\pgfsetstrokecolor{currentstroke}%
\pgfsetstrokeopacity{0.000000}%
\pgfsetdash{}{0pt}%
\pgfpathmoveto{\pgfqpoint{2.625680in}{1.382482in}}%
\pgfpathlineto{\pgfqpoint{2.626550in}{1.378374in}}%
\pgfpathlineto{\pgfqpoint{2.607627in}{1.374496in}}%
\pgfpathlineto{\pgfqpoint{2.604647in}{1.373975in}}%
\pgfpathlineto{\pgfqpoint{2.595492in}{1.381828in}}%
\pgfpathlineto{\pgfqpoint{2.581784in}{1.380075in}}%
\pgfpathlineto{\pgfqpoint{2.558934in}{1.375533in}}%
\pgfpathlineto{\pgfqpoint{2.556898in}{1.389007in}}%
\pgfpathlineto{\pgfqpoint{2.558369in}{1.389969in}}%
\pgfpathlineto{\pgfqpoint{2.556395in}{1.400149in}}%
\pgfpathlineto{\pgfqpoint{2.549510in}{1.399000in}}%
\pgfpathlineto{\pgfqpoint{2.543334in}{1.433990in}}%
\pgfpathlineto{\pgfqpoint{2.546681in}{1.434064in}}%
\pgfpathlineto{\pgfqpoint{2.552432in}{1.429022in}}%
\pgfpathlineto{\pgfqpoint{2.557784in}{1.431044in}}%
\pgfpathlineto{\pgfqpoint{2.570755in}{1.438725in}}%
\pgfpathlineto{\pgfqpoint{2.572079in}{1.437927in}}%
\pgfpathlineto{\pgfqpoint{2.601922in}{1.443430in}}%
\pgfpathlineto{\pgfqpoint{2.608126in}{1.441492in}}%
\pgfpathlineto{\pgfqpoint{2.612717in}{1.410639in}}%
\pgfpathlineto{\pgfqpoint{2.619453in}{1.411977in}}%
\pgfpathlineto{\pgfqpoint{2.625680in}{1.382482in}}%
\pgfpathclose%
\pgfusepath{fill}%
\end{pgfscope}%
\begin{pgfscope}%
\pgfpathrectangle{\pgfqpoint{0.100000in}{0.100000in}}{\pgfqpoint{3.007045in}{1.925000in}}%
\pgfusepath{clip}%
\pgfsetbuttcap%
\pgfsetmiterjoin%
\definecolor{currentfill}{rgb}{0.695071,0.877801,0.669865}%
\pgfsetfillcolor{currentfill}%
\pgfsetlinewidth{0.000000pt}%
\definecolor{currentstroke}{rgb}{0.000000,0.000000,0.000000}%
\pgfsetstrokecolor{currentstroke}%
\pgfsetstrokeopacity{0.000000}%
\pgfsetdash{}{0pt}%
\pgfpathmoveto{\pgfqpoint{2.252299in}{1.407809in}}%
\pgfpathlineto{\pgfqpoint{2.263712in}{1.409069in}}%
\pgfpathlineto{\pgfqpoint{2.261163in}{1.431916in}}%
\pgfpathlineto{\pgfqpoint{2.283625in}{1.434436in}}%
\pgfpathlineto{\pgfqpoint{2.286455in}{1.412044in}}%
\pgfpathlineto{\pgfqpoint{2.296340in}{1.413257in}}%
\pgfpathlineto{\pgfqpoint{2.299946in}{1.390294in}}%
\pgfpathlineto{\pgfqpoint{2.272022in}{1.386856in}}%
\pgfpathlineto{\pgfqpoint{2.254932in}{1.384910in}}%
\pgfpathlineto{\pgfqpoint{2.252299in}{1.407809in}}%
\pgfpathclose%
\pgfusepath{fill}%
\end{pgfscope}%
\begin{pgfscope}%
\pgfpathrectangle{\pgfqpoint{0.100000in}{0.100000in}}{\pgfqpoint{3.007045in}{1.925000in}}%
\pgfusepath{clip}%
\pgfsetbuttcap%
\pgfsetmiterjoin%
\definecolor{currentfill}{rgb}{0.513033,0.795755,0.510542}%
\pgfsetfillcolor{currentfill}%
\pgfsetlinewidth{0.000000pt}%
\definecolor{currentstroke}{rgb}{0.000000,0.000000,0.000000}%
\pgfsetstrokecolor{currentstroke}%
\pgfsetstrokeopacity{0.000000}%
\pgfsetdash{}{0pt}%
\pgfpathmoveto{\pgfqpoint{1.677065in}{1.742344in}}%
\pgfpathlineto{\pgfqpoint{1.678294in}{1.748130in}}%
\pgfpathlineto{\pgfqpoint{1.676638in}{1.755236in}}%
\pgfpathlineto{\pgfqpoint{1.677958in}{1.760345in}}%
\pgfpathlineto{\pgfqpoint{1.676648in}{1.765752in}}%
\pgfpathlineto{\pgfqpoint{1.710937in}{1.765107in}}%
\pgfpathlineto{\pgfqpoint{1.745680in}{1.764366in}}%
\pgfpathlineto{\pgfqpoint{1.745890in}{1.729762in}}%
\pgfpathlineto{\pgfqpoint{1.746314in}{1.723984in}}%
\pgfpathlineto{\pgfqpoint{1.740624in}{1.724147in}}%
\pgfpathlineto{\pgfqpoint{1.740597in}{1.718328in}}%
\pgfpathlineto{\pgfqpoint{1.734882in}{1.718354in}}%
\pgfpathlineto{\pgfqpoint{1.734818in}{1.712519in}}%
\pgfpathlineto{\pgfqpoint{1.711752in}{1.712779in}}%
\pgfpathlineto{\pgfqpoint{1.711793in}{1.718564in}}%
\pgfpathlineto{\pgfqpoint{1.706015in}{1.718669in}}%
\pgfpathlineto{\pgfqpoint{1.705532in}{1.740473in}}%
\pgfpathlineto{\pgfqpoint{1.676789in}{1.740990in}}%
\pgfpathlineto{\pgfqpoint{1.677065in}{1.742344in}}%
\pgfpathclose%
\pgfusepath{fill}%
\end{pgfscope}%
\begin{pgfscope}%
\pgfpathrectangle{\pgfqpoint{0.100000in}{0.100000in}}{\pgfqpoint{3.007045in}{1.925000in}}%
\pgfusepath{clip}%
\pgfsetbuttcap%
\pgfsetmiterjoin%
\definecolor{currentfill}{rgb}{0.590542,0.831926,0.574271}%
\pgfsetfillcolor{currentfill}%
\pgfsetlinewidth{0.000000pt}%
\definecolor{currentstroke}{rgb}{0.000000,0.000000,0.000000}%
\pgfsetstrokecolor{currentstroke}%
\pgfsetstrokeopacity{0.000000}%
\pgfsetdash{}{0pt}%
\pgfpathmoveto{\pgfqpoint{2.667563in}{1.268330in}}%
\pgfpathlineto{\pgfqpoint{2.655502in}{1.265959in}}%
\pgfpathlineto{\pgfqpoint{2.652060in}{1.280358in}}%
\pgfpathlineto{\pgfqpoint{2.638827in}{1.299653in}}%
\pgfpathlineto{\pgfqpoint{2.636603in}{1.298784in}}%
\pgfpathlineto{\pgfqpoint{2.632407in}{1.305014in}}%
\pgfpathlineto{\pgfqpoint{2.627758in}{1.303889in}}%
\pgfpathlineto{\pgfqpoint{2.625195in}{1.307053in}}%
\pgfpathlineto{\pgfqpoint{2.628110in}{1.312979in}}%
\pgfpathlineto{\pgfqpoint{2.626236in}{1.316064in}}%
\pgfpathlineto{\pgfqpoint{2.633083in}{1.330389in}}%
\pgfpathlineto{\pgfqpoint{2.645950in}{1.339764in}}%
\pgfpathlineto{\pgfqpoint{2.648248in}{1.330436in}}%
\pgfpathlineto{\pgfqpoint{2.664036in}{1.332084in}}%
\pgfpathlineto{\pgfqpoint{2.670437in}{1.328641in}}%
\pgfpathlineto{\pgfqpoint{2.680799in}{1.334029in}}%
\pgfpathlineto{\pgfqpoint{2.690299in}{1.329057in}}%
\pgfpathlineto{\pgfqpoint{2.684366in}{1.322430in}}%
\pgfpathlineto{\pgfqpoint{2.693385in}{1.305552in}}%
\pgfpathlineto{\pgfqpoint{2.687398in}{1.300784in}}%
\pgfpathlineto{\pgfqpoint{2.690460in}{1.296019in}}%
\pgfpathlineto{\pgfqpoint{2.697015in}{1.296543in}}%
\pgfpathlineto{\pgfqpoint{2.701481in}{1.290448in}}%
\pgfpathlineto{\pgfqpoint{2.706007in}{1.289043in}}%
\pgfpathlineto{\pgfqpoint{2.716037in}{1.278304in}}%
\pgfpathlineto{\pgfqpoint{2.667563in}{1.268330in}}%
\pgfpathclose%
\pgfusepath{fill}%
\end{pgfscope}%
\begin{pgfscope}%
\pgfpathrectangle{\pgfqpoint{0.100000in}{0.100000in}}{\pgfqpoint{3.007045in}{1.925000in}}%
\pgfusepath{clip}%
\pgfsetbuttcap%
\pgfsetmiterjoin%
\definecolor{currentfill}{rgb}{0.652995,0.860085,0.628897}%
\pgfsetfillcolor{currentfill}%
\pgfsetlinewidth{0.000000pt}%
\definecolor{currentstroke}{rgb}{0.000000,0.000000,0.000000}%
\pgfsetstrokecolor{currentstroke}%
\pgfsetstrokeopacity{0.000000}%
\pgfsetdash{}{0pt}%
\pgfpathmoveto{\pgfqpoint{1.395381in}{1.309582in}}%
\pgfpathlineto{\pgfqpoint{1.396484in}{1.306583in}}%
\pgfpathlineto{\pgfqpoint{1.395479in}{1.295197in}}%
\pgfpathlineto{\pgfqpoint{1.394532in}{1.283801in}}%
\pgfpathlineto{\pgfqpoint{1.392134in}{1.280968in}}%
\pgfpathlineto{\pgfqpoint{1.346503in}{1.285013in}}%
\pgfpathlineto{\pgfqpoint{1.294492in}{1.290219in}}%
\pgfpathlineto{\pgfqpoint{1.262254in}{1.293547in}}%
\pgfpathlineto{\pgfqpoint{1.267098in}{1.336682in}}%
\pgfpathlineto{\pgfqpoint{1.297671in}{1.333074in}}%
\pgfpathlineto{\pgfqpoint{1.297072in}{1.327304in}}%
\pgfpathlineto{\pgfqpoint{1.326529in}{1.324213in}}%
\pgfpathlineto{\pgfqpoint{1.325420in}{1.312996in}}%
\pgfpathlineto{\pgfqpoint{1.358996in}{1.309848in}}%
\pgfpathlineto{\pgfqpoint{1.359322in}{1.312674in}}%
\pgfpathlineto{\pgfqpoint{1.395381in}{1.309582in}}%
\pgfpathclose%
\pgfusepath{fill}%
\end{pgfscope}%
\begin{pgfscope}%
\pgfpathrectangle{\pgfqpoint{0.100000in}{0.100000in}}{\pgfqpoint{3.007045in}{1.925000in}}%
\pgfusepath{clip}%
\pgfsetbuttcap%
\pgfsetmiterjoin%
\definecolor{currentfill}{rgb}{0.307451,0.696348,0.390465}%
\pgfsetfillcolor{currentfill}%
\pgfsetlinewidth{0.000000pt}%
\definecolor{currentstroke}{rgb}{0.000000,0.000000,0.000000}%
\pgfsetstrokecolor{currentstroke}%
\pgfsetstrokeopacity{0.000000}%
\pgfsetdash{}{0pt}%
\pgfpathmoveto{\pgfqpoint{1.201403in}{1.171929in}}%
\pgfpathlineto{\pgfqpoint{1.200990in}{1.169461in}}%
\pgfpathlineto{\pgfqpoint{1.208767in}{1.163511in}}%
\pgfpathlineto{\pgfqpoint{1.213448in}{1.158253in}}%
\pgfpathlineto{\pgfqpoint{1.213877in}{1.153875in}}%
\pgfpathlineto{\pgfqpoint{1.209760in}{1.146848in}}%
\pgfpathlineto{\pgfqpoint{1.213461in}{1.141563in}}%
\pgfpathlineto{\pgfqpoint{1.213054in}{1.137702in}}%
\pgfpathlineto{\pgfqpoint{1.205700in}{1.130953in}}%
\pgfpathlineto{\pgfqpoint{1.202180in}{1.129872in}}%
\pgfpathlineto{\pgfqpoint{1.193414in}{1.130886in}}%
\pgfpathlineto{\pgfqpoint{1.189161in}{1.136735in}}%
\pgfpathlineto{\pgfqpoint{1.187185in}{1.151775in}}%
\pgfpathlineto{\pgfqpoint{1.188060in}{1.158199in}}%
\pgfpathlineto{\pgfqpoint{1.193507in}{1.162201in}}%
\pgfpathlineto{\pgfqpoint{1.186288in}{1.164150in}}%
\pgfpathlineto{\pgfqpoint{1.180201in}{1.170746in}}%
\pgfpathlineto{\pgfqpoint{1.181782in}{1.174522in}}%
\pgfpathlineto{\pgfqpoint{1.201403in}{1.171929in}}%
\pgfpathclose%
\pgfusepath{fill}%
\end{pgfscope}%
\begin{pgfscope}%
\pgfpathrectangle{\pgfqpoint{0.100000in}{0.100000in}}{\pgfqpoint{3.007045in}{1.925000in}}%
\pgfusepath{clip}%
\pgfsetbuttcap%
\pgfsetmiterjoin%
\definecolor{currentfill}{rgb}{0.529642,0.803506,0.524198}%
\pgfsetfillcolor{currentfill}%
\pgfsetlinewidth{0.000000pt}%
\definecolor{currentstroke}{rgb}{0.000000,0.000000,0.000000}%
\pgfsetstrokecolor{currentstroke}%
\pgfsetstrokeopacity{0.000000}%
\pgfsetdash{}{0pt}%
\pgfpathmoveto{\pgfqpoint{2.114171in}{0.933474in}}%
\pgfpathlineto{\pgfqpoint{2.108828in}{0.933214in}}%
\pgfpathlineto{\pgfqpoint{2.108338in}{0.951144in}}%
\pgfpathlineto{\pgfqpoint{2.110066in}{0.954874in}}%
\pgfpathlineto{\pgfqpoint{2.116626in}{0.955211in}}%
\pgfpathlineto{\pgfqpoint{2.115651in}{0.980661in}}%
\pgfpathlineto{\pgfqpoint{2.117091in}{0.980740in}}%
\pgfpathlineto{\pgfqpoint{2.140155in}{0.982191in}}%
\pgfpathlineto{\pgfqpoint{2.140247in}{0.978276in}}%
\pgfpathlineto{\pgfqpoint{2.146159in}{0.970611in}}%
\pgfpathlineto{\pgfqpoint{2.146928in}{0.965917in}}%
\pgfpathlineto{\pgfqpoint{2.158552in}{0.966327in}}%
\pgfpathlineto{\pgfqpoint{2.166450in}{0.964466in}}%
\pgfpathlineto{\pgfqpoint{2.169480in}{0.968102in}}%
\pgfpathlineto{\pgfqpoint{2.169549in}{0.973787in}}%
\pgfpathlineto{\pgfqpoint{2.181638in}{0.973954in}}%
\pgfpathlineto{\pgfqpoint{2.188596in}{0.973610in}}%
\pgfpathlineto{\pgfqpoint{2.187491in}{0.967104in}}%
\pgfpathlineto{\pgfqpoint{2.188978in}{0.963509in}}%
\pgfpathlineto{\pgfqpoint{2.188043in}{0.950490in}}%
\pgfpathlineto{\pgfqpoint{2.170276in}{0.951262in}}%
\pgfpathlineto{\pgfqpoint{2.161575in}{0.944740in}}%
\pgfpathlineto{\pgfqpoint{2.161409in}{0.940303in}}%
\pgfpathlineto{\pgfqpoint{2.151093in}{0.939458in}}%
\pgfpathlineto{\pgfqpoint{2.144945in}{0.937174in}}%
\pgfpathlineto{\pgfqpoint{2.143491in}{0.939640in}}%
\pgfpathlineto{\pgfqpoint{2.136863in}{0.936941in}}%
\pgfpathlineto{\pgfqpoint{2.116000in}{0.935950in}}%
\pgfpathlineto{\pgfqpoint{2.114171in}{0.933474in}}%
\pgfpathclose%
\pgfusepath{fill}%
\end{pgfscope}%
\begin{pgfscope}%
\pgfpathrectangle{\pgfqpoint{0.100000in}{0.100000in}}{\pgfqpoint{3.007045in}{1.925000in}}%
\pgfusepath{clip}%
\pgfsetbuttcap%
\pgfsetmiterjoin%
\definecolor{currentfill}{rgb}{0.845905,0.939931,0.822822}%
\pgfsetfillcolor{currentfill}%
\pgfsetlinewidth{0.000000pt}%
\definecolor{currentstroke}{rgb}{0.000000,0.000000,0.000000}%
\pgfsetstrokecolor{currentstroke}%
\pgfsetstrokeopacity{0.000000}%
\pgfsetdash{}{0pt}%
\pgfpathmoveto{\pgfqpoint{0.920301in}{0.337334in}}%
\pgfpathlineto{\pgfqpoint{0.912069in}{0.341923in}}%
\pgfpathlineto{\pgfqpoint{0.914753in}{0.346694in}}%
\pgfpathlineto{\pgfqpoint{0.913691in}{0.348399in}}%
\pgfpathlineto{\pgfqpoint{0.917983in}{0.355693in}}%
\pgfpathlineto{\pgfqpoint{0.916310in}{0.356673in}}%
\pgfpathlineto{\pgfqpoint{0.917368in}{0.358439in}}%
\pgfpathlineto{\pgfqpoint{0.915470in}{0.359574in}}%
\pgfpathlineto{\pgfqpoint{0.916564in}{0.361368in}}%
\pgfpathlineto{\pgfqpoint{0.913444in}{0.363209in}}%
\pgfpathlineto{\pgfqpoint{0.912886in}{0.364295in}}%
\pgfpathlineto{\pgfqpoint{0.910185in}{0.362649in}}%
\pgfpathlineto{\pgfqpoint{0.907829in}{0.364423in}}%
\pgfpathlineto{\pgfqpoint{0.910203in}{0.367541in}}%
\pgfpathlineto{\pgfqpoint{0.909515in}{0.369615in}}%
\pgfpathlineto{\pgfqpoint{0.912562in}{0.370845in}}%
\pgfpathlineto{\pgfqpoint{0.914264in}{0.375762in}}%
\pgfpathlineto{\pgfqpoint{0.899657in}{0.384205in}}%
\pgfpathlineto{\pgfqpoint{0.894498in}{0.387400in}}%
\pgfpathlineto{\pgfqpoint{0.895631in}{0.389258in}}%
\pgfpathlineto{\pgfqpoint{0.886512in}{0.394965in}}%
\pgfpathlineto{\pgfqpoint{0.884245in}{0.391268in}}%
\pgfpathlineto{\pgfqpoint{0.880633in}{0.393554in}}%
\pgfpathlineto{\pgfqpoint{0.886359in}{0.402751in}}%
\pgfpathlineto{\pgfqpoint{0.875033in}{0.409830in}}%
\pgfpathlineto{\pgfqpoint{0.871795in}{0.412013in}}%
\pgfpathlineto{\pgfqpoint{0.869886in}{0.409139in}}%
\pgfpathlineto{\pgfqpoint{0.856121in}{0.418510in}}%
\pgfpathlineto{\pgfqpoint{0.847814in}{0.419834in}}%
\pgfpathlineto{\pgfqpoint{0.827109in}{0.423257in}}%
\pgfpathlineto{\pgfqpoint{0.826027in}{0.421875in}}%
\pgfpathlineto{\pgfqpoint{0.816153in}{0.429462in}}%
\pgfpathlineto{\pgfqpoint{0.809857in}{0.421160in}}%
\pgfpathlineto{\pgfqpoint{0.807974in}{0.419972in}}%
\pgfpathlineto{\pgfqpoint{0.805001in}{0.422296in}}%
\pgfpathlineto{\pgfqpoint{0.803696in}{0.420629in}}%
\pgfpathlineto{\pgfqpoint{0.800378in}{0.423260in}}%
\pgfpathlineto{\pgfqpoint{0.798438in}{0.420763in}}%
\pgfpathlineto{\pgfqpoint{0.796550in}{0.422002in}}%
\pgfpathlineto{\pgfqpoint{0.787335in}{0.429069in}}%
\pgfpathlineto{\pgfqpoint{0.776021in}{0.437944in}}%
\pgfpathlineto{\pgfqpoint{0.768555in}{0.444177in}}%
\pgfpathlineto{\pgfqpoint{0.770089in}{0.446056in}}%
\pgfpathlineto{\pgfqpoint{0.757361in}{0.456738in}}%
\pgfpathlineto{\pgfqpoint{0.755991in}{0.455118in}}%
\pgfpathlineto{\pgfqpoint{0.749271in}{0.460813in}}%
\pgfpathlineto{\pgfqpoint{0.747517in}{0.458771in}}%
\pgfpathlineto{\pgfqpoint{0.741635in}{0.463832in}}%
\pgfpathlineto{\pgfqpoint{0.733232in}{0.471329in}}%
\pgfpathlineto{\pgfqpoint{0.737384in}{0.476047in}}%
\pgfpathlineto{\pgfqpoint{0.736206in}{0.477104in}}%
\pgfpathlineto{\pgfqpoint{0.741905in}{0.483472in}}%
\pgfpathlineto{\pgfqpoint{0.742309in}{0.483156in}}%
\pgfpathlineto{\pgfqpoint{0.748071in}{0.489521in}}%
\pgfpathlineto{\pgfqpoint{0.750170in}{0.491012in}}%
\pgfpathlineto{\pgfqpoint{0.751165in}{0.490590in}}%
\pgfpathlineto{\pgfqpoint{0.755431in}{0.495345in}}%
\pgfpathlineto{\pgfqpoint{0.756937in}{0.494067in}}%
\pgfpathlineto{\pgfqpoint{0.761153in}{0.499198in}}%
\pgfpathlineto{\pgfqpoint{0.762810in}{0.497757in}}%
\pgfpathlineto{\pgfqpoint{0.764086in}{0.499185in}}%
\pgfpathlineto{\pgfqpoint{0.767409in}{0.496302in}}%
\pgfpathlineto{\pgfqpoint{0.768805in}{0.498004in}}%
\pgfpathlineto{\pgfqpoint{0.770286in}{0.496757in}}%
\pgfpathlineto{\pgfqpoint{0.771756in}{0.498478in}}%
\pgfpathlineto{\pgfqpoint{0.773225in}{0.497185in}}%
\pgfpathlineto{\pgfqpoint{0.777669in}{0.501945in}}%
\pgfpathlineto{\pgfqpoint{0.775774in}{0.503541in}}%
\pgfpathlineto{\pgfqpoint{0.787134in}{0.516986in}}%
\pgfpathlineto{\pgfqpoint{0.788832in}{0.515575in}}%
\pgfpathlineto{\pgfqpoint{0.791670in}{0.518914in}}%
\pgfpathlineto{\pgfqpoint{0.793312in}{0.517519in}}%
\pgfpathlineto{\pgfqpoint{0.798922in}{0.524065in}}%
\pgfpathlineto{\pgfqpoint{0.800473in}{0.522699in}}%
\pgfpathlineto{\pgfqpoint{0.804798in}{0.527778in}}%
\pgfpathlineto{\pgfqpoint{0.803129in}{0.529197in}}%
\pgfpathlineto{\pgfqpoint{0.810396in}{0.537663in}}%
\pgfpathlineto{\pgfqpoint{0.815410in}{0.533363in}}%
\pgfpathlineto{\pgfqpoint{0.818104in}{0.536633in}}%
\pgfpathlineto{\pgfqpoint{0.826641in}{0.529420in}}%
\pgfpathlineto{\pgfqpoint{0.832267in}{0.536448in}}%
\pgfpathlineto{\pgfqpoint{0.839219in}{0.530767in}}%
\pgfpathlineto{\pgfqpoint{0.836381in}{0.527215in}}%
\pgfpathlineto{\pgfqpoint{0.839868in}{0.524417in}}%
\pgfpathlineto{\pgfqpoint{0.842673in}{0.527910in}}%
\pgfpathlineto{\pgfqpoint{0.846127in}{0.525116in}}%
\pgfpathlineto{\pgfqpoint{0.844770in}{0.523398in}}%
\pgfpathlineto{\pgfqpoint{0.846505in}{0.521977in}}%
\pgfpathlineto{\pgfqpoint{0.845093in}{0.520174in}}%
\pgfpathlineto{\pgfqpoint{0.848387in}{0.517598in}}%
\pgfpathlineto{\pgfqpoint{0.852921in}{0.522581in}}%
\pgfpathlineto{\pgfqpoint{0.857919in}{0.518605in}}%
\pgfpathlineto{\pgfqpoint{0.860718in}{0.522144in}}%
\pgfpathlineto{\pgfqpoint{0.865978in}{0.517880in}}%
\pgfpathlineto{\pgfqpoint{0.867902in}{0.519376in}}%
\pgfpathlineto{\pgfqpoint{0.873845in}{0.526915in}}%
\pgfpathlineto{\pgfqpoint{0.872574in}{0.527836in}}%
\pgfpathlineto{\pgfqpoint{0.873947in}{0.529693in}}%
\pgfpathlineto{\pgfqpoint{0.870380in}{0.532422in}}%
\pgfpathlineto{\pgfqpoint{0.874516in}{0.537747in}}%
\pgfpathlineto{\pgfqpoint{0.876407in}{0.539190in}}%
\pgfpathlineto{\pgfqpoint{0.874603in}{0.540599in}}%
\pgfpathlineto{\pgfqpoint{0.876010in}{0.542324in}}%
\pgfpathlineto{\pgfqpoint{0.874170in}{0.543737in}}%
\pgfpathlineto{\pgfqpoint{0.875594in}{0.545533in}}%
\pgfpathlineto{\pgfqpoint{0.873853in}{0.546928in}}%
\pgfpathlineto{\pgfqpoint{0.877896in}{0.551221in}}%
\pgfpathlineto{\pgfqpoint{0.895892in}{0.537283in}}%
\pgfpathlineto{\pgfqpoint{0.919823in}{0.519872in}}%
\pgfpathlineto{\pgfqpoint{0.941118in}{0.505304in}}%
\pgfpathlineto{\pgfqpoint{0.944682in}{0.503522in}}%
\pgfpathlineto{\pgfqpoint{0.952820in}{0.498342in}}%
\pgfpathlineto{\pgfqpoint{0.959476in}{0.508253in}}%
\pgfpathlineto{\pgfqpoint{0.966890in}{0.503916in}}%
\pgfpathlineto{\pgfqpoint{0.989639in}{0.489920in}}%
\pgfpathlineto{\pgfqpoint{1.001047in}{0.482991in}}%
\pgfpathlineto{\pgfqpoint{0.992914in}{0.468320in}}%
\pgfpathlineto{\pgfqpoint{0.981844in}{0.448351in}}%
\pgfpathlineto{\pgfqpoint{0.970915in}{0.428636in}}%
\pgfpathlineto{\pgfqpoint{0.960932in}{0.410628in}}%
\pgfpathlineto{\pgfqpoint{0.950240in}{0.391342in}}%
\pgfpathlineto{\pgfqpoint{0.939656in}{0.372249in}}%
\pgfpathlineto{\pgfqpoint{0.929170in}{0.353333in}}%
\pgfpathlineto{\pgfqpoint{0.920301in}{0.337334in}}%
\pgfpathclose%
\pgfusepath{fill}%
\end{pgfscope}%
\begin{pgfscope}%
\pgfpathrectangle{\pgfqpoint{0.100000in}{0.100000in}}{\pgfqpoint{3.007045in}{1.925000in}}%
\pgfusepath{clip}%
\pgfsetbuttcap%
\pgfsetmiterjoin%
\definecolor{currentfill}{rgb}{0.634295,0.852211,0.610688}%
\pgfsetfillcolor{currentfill}%
\pgfsetlinewidth{0.000000pt}%
\definecolor{currentstroke}{rgb}{0.000000,0.000000,0.000000}%
\pgfsetstrokecolor{currentstroke}%
\pgfsetstrokeopacity{0.000000}%
\pgfsetdash{}{0pt}%
\pgfpathmoveto{\pgfqpoint{1.924728in}{0.990132in}}%
\pgfpathlineto{\pgfqpoint{1.890280in}{0.989939in}}%
\pgfpathlineto{\pgfqpoint{1.859810in}{0.990030in}}%
\pgfpathlineto{\pgfqpoint{1.859842in}{1.001580in}}%
\pgfpathlineto{\pgfqpoint{1.845487in}{1.001557in}}%
\pgfpathlineto{\pgfqpoint{1.845462in}{1.020387in}}%
\pgfpathlineto{\pgfqpoint{1.844374in}{1.020395in}}%
\pgfpathlineto{\pgfqpoint{1.844766in}{1.050707in}}%
\pgfpathlineto{\pgfqpoint{1.843696in}{1.056469in}}%
\pgfpathlineto{\pgfqpoint{1.846715in}{1.062147in}}%
\pgfpathlineto{\pgfqpoint{1.849656in}{1.062106in}}%
\pgfpathlineto{\pgfqpoint{1.849904in}{1.072618in}}%
\pgfpathlineto{\pgfqpoint{1.872670in}{1.072352in}}%
\pgfpathlineto{\pgfqpoint{1.872518in}{1.061815in}}%
\pgfpathlineto{\pgfqpoint{1.896036in}{1.061694in}}%
\pgfpathlineto{\pgfqpoint{1.900128in}{1.055674in}}%
\pgfpathlineto{\pgfqpoint{1.907064in}{1.059811in}}%
\pgfpathlineto{\pgfqpoint{1.907105in}{1.050239in}}%
\pgfpathlineto{\pgfqpoint{1.915821in}{1.046178in}}%
\pgfpathlineto{\pgfqpoint{1.915878in}{1.043286in}}%
\pgfpathlineto{\pgfqpoint{1.916864in}{1.007392in}}%
\pgfpathlineto{\pgfqpoint{1.925420in}{1.007434in}}%
\pgfpathlineto{\pgfqpoint{1.924728in}{0.990132in}}%
\pgfpathclose%
\pgfusepath{fill}%
\end{pgfscope}%
\begin{pgfscope}%
\pgfpathrectangle{\pgfqpoint{0.100000in}{0.100000in}}{\pgfqpoint{3.007045in}{1.925000in}}%
\pgfusepath{clip}%
\pgfsetbuttcap%
\pgfsetmiterjoin%
\definecolor{currentfill}{rgb}{0.601615,0.837093,0.583376}%
\pgfsetfillcolor{currentfill}%
\pgfsetlinewidth{0.000000pt}%
\definecolor{currentstroke}{rgb}{0.000000,0.000000,0.000000}%
\pgfsetstrokecolor{currentstroke}%
\pgfsetstrokeopacity{0.000000}%
\pgfsetdash{}{0pt}%
\pgfpathmoveto{\pgfqpoint{1.829586in}{0.787129in}}%
\pgfpathlineto{\pgfqpoint{1.825054in}{0.789133in}}%
\pgfpathlineto{\pgfqpoint{1.814189in}{0.790353in}}%
\pgfpathlineto{\pgfqpoint{1.806094in}{0.799451in}}%
\pgfpathlineto{\pgfqpoint{1.801001in}{0.799445in}}%
\pgfpathlineto{\pgfqpoint{1.800644in}{0.779398in}}%
\pgfpathlineto{\pgfqpoint{1.786299in}{0.783647in}}%
\pgfpathlineto{\pgfqpoint{1.779488in}{0.786480in}}%
\pgfpathlineto{\pgfqpoint{1.774881in}{0.791133in}}%
\pgfpathlineto{\pgfqpoint{1.760297in}{0.800742in}}%
\pgfpathlineto{\pgfqpoint{1.756132in}{0.794865in}}%
\pgfpathlineto{\pgfqpoint{1.749934in}{0.794479in}}%
\pgfpathlineto{\pgfqpoint{1.742020in}{0.795983in}}%
\pgfpathlineto{\pgfqpoint{1.740181in}{0.799040in}}%
\pgfpathlineto{\pgfqpoint{1.730486in}{0.795042in}}%
\pgfpathlineto{\pgfqpoint{1.726173in}{0.795162in}}%
\pgfpathlineto{\pgfqpoint{1.720879in}{0.800114in}}%
\pgfpathlineto{\pgfqpoint{1.720987in}{0.811838in}}%
\pgfpathlineto{\pgfqpoint{1.718089in}{0.813933in}}%
\pgfpathlineto{\pgfqpoint{1.729695in}{0.813812in}}%
\pgfpathlineto{\pgfqpoint{1.729891in}{0.836933in}}%
\pgfpathlineto{\pgfqpoint{1.735756in}{0.836902in}}%
\pgfpathlineto{\pgfqpoint{1.735784in}{0.842654in}}%
\pgfpathlineto{\pgfqpoint{1.744358in}{0.842618in}}%
\pgfpathlineto{\pgfqpoint{1.744379in}{0.848364in}}%
\pgfpathlineto{\pgfqpoint{1.775709in}{0.848268in}}%
\pgfpathlineto{\pgfqpoint{1.775725in}{0.836774in}}%
\pgfpathlineto{\pgfqpoint{1.801665in}{0.836897in}}%
\pgfpathlineto{\pgfqpoint{1.801936in}{0.851532in}}%
\pgfpathlineto{\pgfqpoint{1.809383in}{0.849499in}}%
\pgfpathlineto{\pgfqpoint{1.825651in}{0.849322in}}%
\pgfpathlineto{\pgfqpoint{1.830433in}{0.847601in}}%
\pgfpathlineto{\pgfqpoint{1.830412in}{0.826752in}}%
\pgfpathlineto{\pgfqpoint{1.812989in}{0.826934in}}%
\pgfpathlineto{\pgfqpoint{1.813675in}{0.816167in}}%
\pgfpathlineto{\pgfqpoint{1.821492in}{0.816167in}}%
\pgfpathlineto{\pgfqpoint{1.822337in}{0.805756in}}%
\pgfpathlineto{\pgfqpoint{1.824314in}{0.796016in}}%
\pgfpathlineto{\pgfqpoint{1.828440in}{0.791464in}}%
\pgfpathlineto{\pgfqpoint{1.829586in}{0.787129in}}%
\pgfpathclose%
\pgfusepath{fill}%
\end{pgfscope}%
\begin{pgfscope}%
\pgfpathrectangle{\pgfqpoint{0.100000in}{0.100000in}}{\pgfqpoint{3.007045in}{1.925000in}}%
\pgfusepath{clip}%
\pgfsetbuttcap%
\pgfsetmiterjoin%
\definecolor{currentfill}{rgb}{0.797924,0.920738,0.771642}%
\pgfsetfillcolor{currentfill}%
\pgfsetlinewidth{0.000000pt}%
\definecolor{currentstroke}{rgb}{0.000000,0.000000,0.000000}%
\pgfsetstrokecolor{currentstroke}%
\pgfsetstrokeopacity{0.000000}%
\pgfsetdash{}{0pt}%
\pgfpathmoveto{\pgfqpoint{1.798597in}{1.275256in}}%
\pgfpathlineto{\pgfqpoint{1.787185in}{1.275241in}}%
\pgfpathlineto{\pgfqpoint{1.787084in}{1.304515in}}%
\pgfpathlineto{\pgfqpoint{1.784942in}{1.304507in}}%
\pgfpathlineto{\pgfqpoint{1.784860in}{1.321668in}}%
\pgfpathlineto{\pgfqpoint{1.813376in}{1.321894in}}%
\pgfpathlineto{\pgfqpoint{1.813145in}{1.344757in}}%
\pgfpathlineto{\pgfqpoint{1.847419in}{1.345161in}}%
\pgfpathlineto{\pgfqpoint{1.858723in}{1.345342in}}%
\pgfpathlineto{\pgfqpoint{1.859156in}{1.322514in}}%
\pgfpathlineto{\pgfqpoint{1.853457in}{1.322433in}}%
\pgfpathlineto{\pgfqpoint{1.853741in}{1.305138in}}%
\pgfpathlineto{\pgfqpoint{1.854836in}{1.299029in}}%
\pgfpathlineto{\pgfqpoint{1.855298in}{1.276147in}}%
\pgfpathlineto{\pgfqpoint{1.843908in}{1.275998in}}%
\pgfpathlineto{\pgfqpoint{1.826979in}{1.275432in}}%
\pgfpathlineto{\pgfqpoint{1.798597in}{1.275256in}}%
\pgfpathclose%
\pgfusepath{fill}%
\end{pgfscope}%
\begin{pgfscope}%
\pgfpathrectangle{\pgfqpoint{0.100000in}{0.100000in}}{\pgfqpoint{3.007045in}{1.925000in}}%
\pgfusepath{clip}%
\pgfsetbuttcap%
\pgfsetmiterjoin%
\definecolor{currentfill}{rgb}{0.751173,0.901423,0.724491}%
\pgfsetfillcolor{currentfill}%
\pgfsetlinewidth{0.000000pt}%
\definecolor{currentstroke}{rgb}{0.000000,0.000000,0.000000}%
\pgfsetstrokecolor{currentstroke}%
\pgfsetstrokeopacity{0.000000}%
\pgfsetdash{}{0pt}%
\pgfpathmoveto{\pgfqpoint{1.502115in}{0.878788in}}%
\pgfpathlineto{\pgfqpoint{1.501084in}{0.860066in}}%
\pgfpathlineto{\pgfqpoint{1.478574in}{0.861428in}}%
\pgfpathlineto{\pgfqpoint{1.471813in}{0.861849in}}%
\pgfpathlineto{\pgfqpoint{1.473659in}{0.890543in}}%
\pgfpathlineto{\pgfqpoint{1.475346in}{0.919280in}}%
\pgfpathlineto{\pgfqpoint{1.504252in}{0.917589in}}%
\pgfpathlineto{\pgfqpoint{1.502115in}{0.878788in}}%
\pgfpathclose%
\pgfusepath{fill}%
\end{pgfscope}%
\begin{pgfscope}%
\pgfpathrectangle{\pgfqpoint{0.100000in}{0.100000in}}{\pgfqpoint{3.007045in}{1.925000in}}%
\pgfusepath{clip}%
\pgfsetbuttcap%
\pgfsetmiterjoin%
\definecolor{currentfill}{rgb}{0.618224,0.844844,0.597032}%
\pgfsetfillcolor{currentfill}%
\pgfsetlinewidth{0.000000pt}%
\definecolor{currentstroke}{rgb}{0.000000,0.000000,0.000000}%
\pgfsetstrokecolor{currentstroke}%
\pgfsetstrokeopacity{0.000000}%
\pgfsetdash{}{0pt}%
\pgfpathmoveto{\pgfqpoint{1.594488in}{0.908994in}}%
\pgfpathlineto{\pgfqpoint{1.577863in}{0.909634in}}%
\pgfpathlineto{\pgfqpoint{1.578070in}{0.926846in}}%
\pgfpathlineto{\pgfqpoint{1.538310in}{0.928586in}}%
\pgfpathlineto{\pgfqpoint{1.538682in}{0.941826in}}%
\pgfpathlineto{\pgfqpoint{1.539102in}{0.951802in}}%
\pgfpathlineto{\pgfqpoint{1.561768in}{0.950581in}}%
\pgfpathlineto{\pgfqpoint{1.562805in}{0.979407in}}%
\pgfpathlineto{\pgfqpoint{1.553847in}{0.991951in}}%
\pgfpathlineto{\pgfqpoint{1.547875in}{0.992381in}}%
\pgfpathlineto{\pgfqpoint{1.541765in}{0.998932in}}%
\pgfpathlineto{\pgfqpoint{1.537864in}{1.006990in}}%
\pgfpathlineto{\pgfqpoint{1.561807in}{1.005883in}}%
\pgfpathlineto{\pgfqpoint{1.596180in}{1.004385in}}%
\pgfpathlineto{\pgfqpoint{1.642737in}{1.002945in}}%
\pgfpathlineto{\pgfqpoint{1.642047in}{0.976241in}}%
\pgfpathlineto{\pgfqpoint{1.641329in}{0.947962in}}%
\pgfpathlineto{\pgfqpoint{1.601428in}{0.949236in}}%
\pgfpathlineto{\pgfqpoint{1.600584in}{0.920240in}}%
\pgfpathlineto{\pgfqpoint{1.594900in}{0.920448in}}%
\pgfpathlineto{\pgfqpoint{1.594488in}{0.908994in}}%
\pgfpathclose%
\pgfusepath{fill}%
\end{pgfscope}%
\begin{pgfscope}%
\pgfpathrectangle{\pgfqpoint{0.100000in}{0.100000in}}{\pgfqpoint{3.007045in}{1.925000in}}%
\pgfusepath{clip}%
\pgfsetbuttcap%
\pgfsetmiterjoin%
\definecolor{currentfill}{rgb}{0.968627,0.988235,0.960784}%
\pgfsetfillcolor{currentfill}%
\pgfsetlinewidth{0.000000pt}%
\definecolor{currentstroke}{rgb}{0.000000,0.000000,0.000000}%
\pgfsetstrokecolor{currentstroke}%
\pgfsetstrokeopacity{0.000000}%
\pgfsetdash{}{0pt}%
\pgfpathmoveto{\pgfqpoint{0.907070in}{1.217261in}}%
\pgfpathlineto{\pgfqpoint{0.903422in}{1.212075in}}%
\pgfpathlineto{\pgfqpoint{0.900045in}{1.212653in}}%
\pgfpathlineto{\pgfqpoint{0.894249in}{1.207621in}}%
\pgfpathlineto{\pgfqpoint{0.892855in}{1.198668in}}%
\pgfpathlineto{\pgfqpoint{0.887613in}{1.196226in}}%
\pgfpathlineto{\pgfqpoint{0.880713in}{1.197892in}}%
\pgfpathlineto{\pgfqpoint{0.875936in}{1.191519in}}%
\pgfpathlineto{\pgfqpoint{0.850288in}{1.196720in}}%
\pgfpathlineto{\pgfqpoint{0.798495in}{1.207353in}}%
\pgfpathlineto{\pgfqpoint{0.799942in}{1.214078in}}%
\pgfpathlineto{\pgfqpoint{0.811909in}{1.269809in}}%
\pgfpathlineto{\pgfqpoint{0.816917in}{1.293210in}}%
\pgfpathlineto{\pgfqpoint{0.878712in}{1.280320in}}%
\pgfpathlineto{\pgfqpoint{0.901338in}{1.275958in}}%
\pgfpathlineto{\pgfqpoint{0.910706in}{1.281361in}}%
\pgfpathlineto{\pgfqpoint{0.914434in}{1.291420in}}%
\pgfpathlineto{\pgfqpoint{0.913219in}{1.297816in}}%
\pgfpathlineto{\pgfqpoint{0.914647in}{1.310630in}}%
\pgfpathlineto{\pgfqpoint{0.923286in}{1.309346in}}%
\pgfpathlineto{\pgfqpoint{0.925436in}{1.305547in}}%
\pgfpathlineto{\pgfqpoint{0.937610in}{1.310762in}}%
\pgfpathlineto{\pgfqpoint{0.947451in}{1.310370in}}%
\pgfpathlineto{\pgfqpoint{0.943668in}{1.304520in}}%
\pgfpathlineto{\pgfqpoint{0.950156in}{1.292847in}}%
\pgfpathlineto{\pgfqpoint{0.957665in}{1.289984in}}%
\pgfpathlineto{\pgfqpoint{0.959849in}{1.282582in}}%
\pgfpathlineto{\pgfqpoint{0.957377in}{1.275606in}}%
\pgfpathlineto{\pgfqpoint{0.964301in}{1.266613in}}%
\pgfpathlineto{\pgfqpoint{0.963825in}{1.264056in}}%
\pgfpathlineto{\pgfqpoint{0.973337in}{1.262333in}}%
\pgfpathlineto{\pgfqpoint{0.974041in}{1.256437in}}%
\pgfpathlineto{\pgfqpoint{0.934985in}{1.263476in}}%
\pgfpathlineto{\pgfqpoint{0.937487in}{1.261054in}}%
\pgfpathlineto{\pgfqpoint{0.929754in}{1.256515in}}%
\pgfpathlineto{\pgfqpoint{0.927198in}{1.241278in}}%
\pgfpathlineto{\pgfqpoint{0.924573in}{1.236589in}}%
\pgfpathlineto{\pgfqpoint{0.915780in}{1.238145in}}%
\pgfpathlineto{\pgfqpoint{0.910572in}{1.236135in}}%
\pgfpathlineto{\pgfqpoint{0.907070in}{1.217261in}}%
\pgfpathclose%
\pgfusepath{fill}%
\end{pgfscope}%
\begin{pgfscope}%
\pgfpathrectangle{\pgfqpoint{0.100000in}{0.100000in}}{\pgfqpoint{3.007045in}{1.925000in}}%
\pgfusepath{clip}%
\pgfsetbuttcap%
\pgfsetmiterjoin%
\definecolor{currentfill}{rgb}{0.535179,0.806090,0.528750}%
\pgfsetfillcolor{currentfill}%
\pgfsetlinewidth{0.000000pt}%
\definecolor{currentstroke}{rgb}{0.000000,0.000000,0.000000}%
\pgfsetstrokecolor{currentstroke}%
\pgfsetstrokeopacity{0.000000}%
\pgfsetdash{}{0pt}%
\pgfpathmoveto{\pgfqpoint{2.375866in}{1.305055in}}%
\pgfpathlineto{\pgfqpoint{2.371168in}{1.333253in}}%
\pgfpathlineto{\pgfqpoint{2.365754in}{1.334376in}}%
\pgfpathlineto{\pgfqpoint{2.378309in}{1.338718in}}%
\pgfpathlineto{\pgfqpoint{2.383725in}{1.333912in}}%
\pgfpathlineto{\pgfqpoint{2.389914in}{1.332610in}}%
\pgfpathlineto{\pgfqpoint{2.400156in}{1.337527in}}%
\pgfpathlineto{\pgfqpoint{2.411268in}{1.344518in}}%
\pgfpathlineto{\pgfqpoint{2.413516in}{1.344212in}}%
\pgfpathlineto{\pgfqpoint{2.414804in}{1.334085in}}%
\pgfpathlineto{\pgfqpoint{2.419430in}{1.334702in}}%
\pgfpathlineto{\pgfqpoint{2.420098in}{1.329773in}}%
\pgfpathlineto{\pgfqpoint{2.415460in}{1.329120in}}%
\pgfpathlineto{\pgfqpoint{2.416075in}{1.324178in}}%
\pgfpathlineto{\pgfqpoint{2.411263in}{1.323487in}}%
\pgfpathlineto{\pgfqpoint{2.407064in}{1.318707in}}%
\pgfpathlineto{\pgfqpoint{2.407666in}{1.313899in}}%
\pgfpathlineto{\pgfqpoint{2.394589in}{1.312199in}}%
\pgfpathlineto{\pgfqpoint{2.395446in}{1.307499in}}%
\pgfpathlineto{\pgfqpoint{2.375866in}{1.305055in}}%
\pgfpathclose%
\pgfusepath{fill}%
\end{pgfscope}%
\begin{pgfscope}%
\pgfpathrectangle{\pgfqpoint{0.100000in}{0.100000in}}{\pgfqpoint{3.007045in}{1.925000in}}%
\pgfusepath{clip}%
\pgfsetbuttcap%
\pgfsetmiterjoin%
\definecolor{currentfill}{rgb}{0.585006,0.829343,0.569719}%
\pgfsetfillcolor{currentfill}%
\pgfsetlinewidth{0.000000pt}%
\definecolor{currentstroke}{rgb}{0.000000,0.000000,0.000000}%
\pgfsetstrokecolor{currentstroke}%
\pgfsetstrokeopacity{0.000000}%
\pgfsetdash{}{0pt}%
\pgfpathmoveto{\pgfqpoint{1.359600in}{1.651719in}}%
\pgfpathlineto{\pgfqpoint{1.355655in}{1.612365in}}%
\pgfpathlineto{\pgfqpoint{1.351714in}{1.614337in}}%
\pgfpathlineto{\pgfqpoint{1.339970in}{1.615629in}}%
\pgfpathlineto{\pgfqpoint{1.340576in}{1.621215in}}%
\pgfpathlineto{\pgfqpoint{1.334802in}{1.621794in}}%
\pgfpathlineto{\pgfqpoint{1.335389in}{1.627188in}}%
\pgfpathlineto{\pgfqpoint{1.318468in}{1.629024in}}%
\pgfpathlineto{\pgfqpoint{1.320048in}{1.651933in}}%
\pgfpathlineto{\pgfqpoint{1.327908in}{1.651084in}}%
\pgfpathlineto{\pgfqpoint{1.329154in}{1.662581in}}%
\pgfpathlineto{\pgfqpoint{1.335063in}{1.663866in}}%
\pgfpathlineto{\pgfqpoint{1.343644in}{1.662962in}}%
\pgfpathlineto{\pgfqpoint{1.346246in}{1.659789in}}%
\pgfpathlineto{\pgfqpoint{1.360257in}{1.658366in}}%
\pgfpathlineto{\pgfqpoint{1.359600in}{1.651719in}}%
\pgfpathclose%
\pgfusepath{fill}%
\end{pgfscope}%
\begin{pgfscope}%
\pgfpathrectangle{\pgfqpoint{0.100000in}{0.100000in}}{\pgfqpoint{3.007045in}{1.925000in}}%
\pgfusepath{clip}%
\pgfsetbuttcap%
\pgfsetmiterjoin%
\definecolor{currentfill}{rgb}{0.294902,0.690196,0.384314}%
\pgfsetfillcolor{currentfill}%
\pgfsetlinewidth{0.000000pt}%
\definecolor{currentstroke}{rgb}{0.000000,0.000000,0.000000}%
\pgfsetstrokecolor{currentstroke}%
\pgfsetstrokeopacity{0.000000}%
\pgfsetdash{}{0pt}%
\pgfpathmoveto{\pgfqpoint{2.765291in}{1.169198in}}%
\pgfpathlineto{\pgfqpoint{2.759352in}{1.165916in}}%
\pgfpathlineto{\pgfqpoint{2.756954in}{1.169348in}}%
\pgfpathlineto{\pgfqpoint{2.759204in}{1.174214in}}%
\pgfpathlineto{\pgfqpoint{2.757805in}{1.180412in}}%
\pgfpathlineto{\pgfqpoint{2.751639in}{1.179119in}}%
\pgfpathlineto{\pgfqpoint{2.751955in}{1.186579in}}%
\pgfpathlineto{\pgfqpoint{2.751060in}{1.190437in}}%
\pgfpathlineto{\pgfqpoint{2.753781in}{1.202785in}}%
\pgfpathlineto{\pgfqpoint{2.758832in}{1.209240in}}%
\pgfpathlineto{\pgfqpoint{2.745705in}{1.255995in}}%
\pgfpathlineto{\pgfqpoint{2.751195in}{1.256851in}}%
\pgfpathlineto{\pgfqpoint{2.757050in}{1.263166in}}%
\pgfpathlineto{\pgfqpoint{2.761478in}{1.260531in}}%
\pgfpathlineto{\pgfqpoint{2.765354in}{1.253983in}}%
\pgfpathlineto{\pgfqpoint{2.766799in}{1.245106in}}%
\pgfpathlineto{\pgfqpoint{2.770462in}{1.242822in}}%
\pgfpathlineto{\pgfqpoint{2.773747in}{1.236358in}}%
\pgfpathlineto{\pgfqpoint{2.781281in}{1.230541in}}%
\pgfpathlineto{\pgfqpoint{2.786967in}{1.229663in}}%
\pgfpathlineto{\pgfqpoint{2.793235in}{1.209405in}}%
\pgfpathlineto{\pgfqpoint{2.791000in}{1.184936in}}%
\pgfpathlineto{\pgfqpoint{2.789514in}{1.179938in}}%
\pgfpathlineto{\pgfqpoint{2.770616in}{1.173606in}}%
\pgfpathlineto{\pgfqpoint{2.765291in}{1.169198in}}%
\pgfpathclose%
\pgfusepath{fill}%
\end{pgfscope}%
\begin{pgfscope}%
\pgfpathrectangle{\pgfqpoint{0.100000in}{0.100000in}}{\pgfqpoint{3.007045in}{1.925000in}}%
\pgfusepath{clip}%
\pgfsetbuttcap%
\pgfsetmiterjoin%
\definecolor{currentfill}{rgb}{0.395294,0.739408,0.433526}%
\pgfsetfillcolor{currentfill}%
\pgfsetlinewidth{0.000000pt}%
\definecolor{currentstroke}{rgb}{0.000000,0.000000,0.000000}%
\pgfsetstrokecolor{currentstroke}%
\pgfsetstrokeopacity{0.000000}%
\pgfsetdash{}{0pt}%
\pgfpathmoveto{\pgfqpoint{1.779241in}{1.344556in}}%
\pgfpathlineto{\pgfqpoint{1.739541in}{1.344779in}}%
\pgfpathlineto{\pgfqpoint{1.739732in}{1.362155in}}%
\pgfpathlineto{\pgfqpoint{1.736283in}{1.367899in}}%
\pgfpathlineto{\pgfqpoint{1.730520in}{1.367928in}}%
\pgfpathlineto{\pgfqpoint{1.730762in}{1.390911in}}%
\pgfpathlineto{\pgfqpoint{1.753560in}{1.390798in}}%
\pgfpathlineto{\pgfqpoint{1.776480in}{1.390760in}}%
\pgfpathlineto{\pgfqpoint{1.776477in}{1.362015in}}%
\pgfpathlineto{\pgfqpoint{1.779293in}{1.361999in}}%
\pgfpathlineto{\pgfqpoint{1.779241in}{1.344556in}}%
\pgfpathclose%
\pgfusepath{fill}%
\end{pgfscope}%
\begin{pgfscope}%
\pgfpathrectangle{\pgfqpoint{0.100000in}{0.100000in}}{\pgfqpoint{3.007045in}{1.925000in}}%
\pgfusepath{clip}%
\pgfsetbuttcap%
\pgfsetmiterjoin%
\definecolor{currentfill}{rgb}{0.765198,0.907328,0.738147}%
\pgfsetfillcolor{currentfill}%
\pgfsetlinewidth{0.000000pt}%
\definecolor{currentstroke}{rgb}{0.000000,0.000000,0.000000}%
\pgfsetstrokecolor{currentstroke}%
\pgfsetstrokeopacity{0.000000}%
\pgfsetdash{}{0pt}%
\pgfpathmoveto{\pgfqpoint{1.011802in}{1.226459in}}%
\pgfpathlineto{\pgfqpoint{1.057846in}{1.219034in}}%
\pgfpathlineto{\pgfqpoint{1.061002in}{1.220921in}}%
\pgfpathlineto{\pgfqpoint{1.059704in}{1.212390in}}%
\pgfpathlineto{\pgfqpoint{1.050444in}{1.156057in}}%
\pgfpathlineto{\pgfqpoint{1.001312in}{1.163987in}}%
\pgfpathlineto{\pgfqpoint{1.002802in}{1.166666in}}%
\pgfpathlineto{\pgfqpoint{0.999346in}{1.179674in}}%
\pgfpathlineto{\pgfqpoint{1.000947in}{1.182123in}}%
\pgfpathlineto{\pgfqpoint{0.998322in}{1.188795in}}%
\pgfpathlineto{\pgfqpoint{1.001745in}{1.203071in}}%
\pgfpathlineto{\pgfqpoint{1.007785in}{1.211607in}}%
\pgfpathlineto{\pgfqpoint{1.007174in}{1.215093in}}%
\pgfpathlineto{\pgfqpoint{1.010713in}{1.219317in}}%
\pgfpathlineto{\pgfqpoint{1.011802in}{1.226459in}}%
\pgfpathclose%
\pgfusepath{fill}%
\end{pgfscope}%
\begin{pgfscope}%
\pgfpathrectangle{\pgfqpoint{0.100000in}{0.100000in}}{\pgfqpoint{3.007045in}{1.925000in}}%
\pgfusepath{clip}%
\pgfsetbuttcap%
\pgfsetmiterjoin%
\definecolor{currentfill}{rgb}{0.507497,0.793172,0.505990}%
\pgfsetfillcolor{currentfill}%
\pgfsetlinewidth{0.000000pt}%
\definecolor{currentstroke}{rgb}{0.000000,0.000000,0.000000}%
\pgfsetstrokecolor{currentstroke}%
\pgfsetstrokeopacity{0.000000}%
\pgfsetdash{}{0pt}%
\pgfpathmoveto{\pgfqpoint{2.334077in}{0.871563in}}%
\pgfpathlineto{\pgfqpoint{2.332579in}{0.883358in}}%
\pgfpathlineto{\pgfqpoint{2.335741in}{0.890096in}}%
\pgfpathlineto{\pgfqpoint{2.333466in}{0.891984in}}%
\pgfpathlineto{\pgfqpoint{2.332547in}{0.900599in}}%
\pgfpathlineto{\pgfqpoint{2.348646in}{0.902478in}}%
\pgfpathlineto{\pgfqpoint{2.356919in}{0.901553in}}%
\pgfpathlineto{\pgfqpoint{2.362404in}{0.891696in}}%
\pgfpathlineto{\pgfqpoint{2.360142in}{0.881259in}}%
\pgfpathlineto{\pgfqpoint{2.358828in}{0.878027in}}%
\pgfpathlineto{\pgfqpoint{2.351405in}{0.875886in}}%
\pgfpathlineto{\pgfqpoint{2.349378in}{0.873307in}}%
\pgfpathlineto{\pgfqpoint{2.334077in}{0.871563in}}%
\pgfpathclose%
\pgfusepath{fill}%
\end{pgfscope}%
\begin{pgfscope}%
\pgfpathrectangle{\pgfqpoint{0.100000in}{0.100000in}}{\pgfqpoint{3.007045in}{1.925000in}}%
\pgfusepath{clip}%
\pgfsetbuttcap%
\pgfsetmiterjoin%
\definecolor{currentfill}{rgb}{0.160323,0.569704,0.289043}%
\pgfsetfillcolor{currentfill}%
\pgfsetlinewidth{0.000000pt}%
\definecolor{currentstroke}{rgb}{0.000000,0.000000,0.000000}%
\pgfsetstrokecolor{currentstroke}%
\pgfsetstrokeopacity{0.000000}%
\pgfsetdash{}{0pt}%
\pgfpathmoveto{\pgfqpoint{0.372778in}{1.622728in}}%
\pgfpathlineto{\pgfqpoint{0.383904in}{1.641105in}}%
\pgfpathlineto{\pgfqpoint{0.396803in}{1.656959in}}%
\pgfpathlineto{\pgfqpoint{0.406441in}{1.674762in}}%
\pgfpathlineto{\pgfqpoint{0.416059in}{1.671542in}}%
\pgfpathlineto{\pgfqpoint{0.423095in}{1.674955in}}%
\pgfpathlineto{\pgfqpoint{0.432010in}{1.670306in}}%
\pgfpathlineto{\pgfqpoint{0.434242in}{1.663279in}}%
\pgfpathlineto{\pgfqpoint{0.441475in}{1.657513in}}%
\pgfpathlineto{\pgfqpoint{0.451005in}{1.654418in}}%
\pgfpathlineto{\pgfqpoint{0.447535in}{1.643476in}}%
\pgfpathlineto{\pgfqpoint{0.447557in}{1.638893in}}%
\pgfpathlineto{\pgfqpoint{0.464341in}{1.633883in}}%
\pgfpathlineto{\pgfqpoint{0.462254in}{1.627100in}}%
\pgfpathlineto{\pgfqpoint{0.490196in}{1.618627in}}%
\pgfpathlineto{\pgfqpoint{0.493741in}{1.612121in}}%
\pgfpathlineto{\pgfqpoint{0.493766in}{1.605432in}}%
\pgfpathlineto{\pgfqpoint{0.487657in}{1.599535in}}%
\pgfpathlineto{\pgfqpoint{0.485002in}{1.594976in}}%
\pgfpathlineto{\pgfqpoint{0.476029in}{1.597225in}}%
\pgfpathlineto{\pgfqpoint{0.474662in}{1.592743in}}%
\pgfpathlineto{\pgfqpoint{0.469189in}{1.594432in}}%
\pgfpathlineto{\pgfqpoint{0.465330in}{1.591720in}}%
\pgfpathlineto{\pgfqpoint{0.460693in}{1.593122in}}%
\pgfpathlineto{\pgfqpoint{0.446705in}{1.586464in}}%
\pgfpathlineto{\pgfqpoint{0.436255in}{1.587889in}}%
\pgfpathlineto{\pgfqpoint{0.429770in}{1.590236in}}%
\pgfpathlineto{\pgfqpoint{0.426249in}{1.588654in}}%
\pgfpathlineto{\pgfqpoint{0.419385in}{1.591902in}}%
\pgfpathlineto{\pgfqpoint{0.415190in}{1.591253in}}%
\pgfpathlineto{\pgfqpoint{0.404299in}{1.600395in}}%
\pgfpathlineto{\pgfqpoint{0.397064in}{1.602137in}}%
\pgfpathlineto{\pgfqpoint{0.389344in}{1.598489in}}%
\pgfpathlineto{\pgfqpoint{0.382499in}{1.599727in}}%
\pgfpathlineto{\pgfqpoint{0.386071in}{1.611691in}}%
\pgfpathlineto{\pgfqpoint{0.382988in}{1.619317in}}%
\pgfpathlineto{\pgfqpoint{0.372778in}{1.622728in}}%
\pgfpathclose%
\pgfusepath{fill}%
\end{pgfscope}%
\begin{pgfscope}%
\pgfpathrectangle{\pgfqpoint{0.100000in}{0.100000in}}{\pgfqpoint{3.007045in}{1.925000in}}%
\pgfusepath{clip}%
\pgfsetbuttcap%
\pgfsetmiterjoin%
\definecolor{currentfill}{rgb}{0.197232,0.609073,0.318570}%
\pgfsetfillcolor{currentfill}%
\pgfsetlinewidth{0.000000pt}%
\definecolor{currentstroke}{rgb}{0.000000,0.000000,0.000000}%
\pgfsetstrokecolor{currentstroke}%
\pgfsetstrokeopacity{0.000000}%
\pgfsetdash{}{0pt}%
\pgfpathmoveto{\pgfqpoint{2.605371in}{0.522544in}}%
\pgfpathlineto{\pgfqpoint{2.591696in}{0.520372in}}%
\pgfpathlineto{\pgfqpoint{2.592396in}{0.508481in}}%
\pgfpathlineto{\pgfqpoint{2.584158in}{0.521537in}}%
\pgfpathlineto{\pgfqpoint{2.581927in}{0.518998in}}%
\pgfpathlineto{\pgfqpoint{2.576314in}{0.520862in}}%
\pgfpathlineto{\pgfqpoint{2.570842in}{0.520187in}}%
\pgfpathlineto{\pgfqpoint{2.567244in}{0.529939in}}%
\pgfpathlineto{\pgfqpoint{2.557842in}{0.535907in}}%
\pgfpathlineto{\pgfqpoint{2.555178in}{0.540975in}}%
\pgfpathlineto{\pgfqpoint{2.550313in}{0.542005in}}%
\pgfpathlineto{\pgfqpoint{2.547864in}{0.546632in}}%
\pgfpathlineto{\pgfqpoint{2.541214in}{0.552000in}}%
\pgfpathlineto{\pgfqpoint{2.538085in}{0.559398in}}%
\pgfpathlineto{\pgfqpoint{2.533920in}{0.561053in}}%
\pgfpathlineto{\pgfqpoint{2.522149in}{0.556021in}}%
\pgfpathlineto{\pgfqpoint{2.520191in}{0.572396in}}%
\pgfpathlineto{\pgfqpoint{2.526180in}{0.575201in}}%
\pgfpathlineto{\pgfqpoint{2.532632in}{0.582139in}}%
\pgfpathlineto{\pgfqpoint{2.545862in}{0.584302in}}%
\pgfpathlineto{\pgfqpoint{2.549830in}{0.579481in}}%
\pgfpathlineto{\pgfqpoint{2.551269in}{0.570455in}}%
\pgfpathlineto{\pgfqpoint{2.562689in}{0.572385in}}%
\pgfpathlineto{\pgfqpoint{2.568608in}{0.576327in}}%
\pgfpathlineto{\pgfqpoint{2.581840in}{0.553861in}}%
\pgfpathlineto{\pgfqpoint{2.588561in}{0.543384in}}%
\pgfpathlineto{\pgfqpoint{2.605371in}{0.522544in}}%
\pgfpathclose%
\pgfusepath{fill}%
\end{pgfscope}%
\begin{pgfscope}%
\pgfpathrectangle{\pgfqpoint{0.100000in}{0.100000in}}{\pgfqpoint{3.007045in}{1.925000in}}%
\pgfusepath{clip}%
\pgfsetbuttcap%
\pgfsetmiterjoin%
\definecolor{currentfill}{rgb}{0.676371,0.869927,0.651657}%
\pgfsetfillcolor{currentfill}%
\pgfsetlinewidth{0.000000pt}%
\definecolor{currentstroke}{rgb}{0.000000,0.000000,0.000000}%
\pgfsetstrokecolor{currentstroke}%
\pgfsetstrokeopacity{0.000000}%
\pgfsetdash{}{0pt}%
\pgfpathmoveto{\pgfqpoint{0.778392in}{0.924882in}}%
\pgfpathlineto{\pgfqpoint{0.767196in}{0.926377in}}%
\pgfpathlineto{\pgfqpoint{0.762651in}{0.922152in}}%
\pgfpathlineto{\pgfqpoint{0.740331in}{0.928918in}}%
\pgfpathlineto{\pgfqpoint{0.735192in}{0.932928in}}%
\pgfpathlineto{\pgfqpoint{0.726621in}{0.944753in}}%
\pgfpathlineto{\pgfqpoint{0.723368in}{0.955271in}}%
\pgfpathlineto{\pgfqpoint{0.723182in}{0.963163in}}%
\pgfpathlineto{\pgfqpoint{0.719619in}{0.967635in}}%
\pgfpathlineto{\pgfqpoint{0.716895in}{0.976496in}}%
\pgfpathlineto{\pgfqpoint{0.718691in}{0.983740in}}%
\pgfpathlineto{\pgfqpoint{0.677176in}{1.047687in}}%
\pgfpathlineto{\pgfqpoint{0.667024in}{1.063268in}}%
\pgfpathlineto{\pgfqpoint{0.616671in}{1.141013in}}%
\pgfpathlineto{\pgfqpoint{0.590725in}{1.181138in}}%
\pgfpathlineto{\pgfqpoint{0.567833in}{1.216455in}}%
\pgfpathlineto{\pgfqpoint{0.571645in}{1.215280in}}%
\pgfpathlineto{\pgfqpoint{0.614328in}{1.243704in}}%
\pgfpathlineto{\pgfqpoint{0.596506in}{1.278590in}}%
\pgfpathlineto{\pgfqpoint{0.597849in}{1.283744in}}%
\pgfpathlineto{\pgfqpoint{0.615468in}{1.284198in}}%
\pgfpathlineto{\pgfqpoint{0.620230in}{1.284320in}}%
\pgfpathlineto{\pgfqpoint{0.643335in}{1.283223in}}%
\pgfpathlineto{\pgfqpoint{0.679545in}{1.274162in}}%
\pgfpathlineto{\pgfqpoint{0.714021in}{1.265975in}}%
\pgfpathlineto{\pgfqpoint{0.752166in}{1.224502in}}%
\pgfpathlineto{\pgfqpoint{0.799942in}{1.214078in}}%
\pgfpathlineto{\pgfqpoint{0.798495in}{1.207353in}}%
\pgfpathlineto{\pgfqpoint{0.788118in}{1.158139in}}%
\pgfpathlineto{\pgfqpoint{0.776912in}{1.106019in}}%
\pgfpathlineto{\pgfqpoint{0.836260in}{1.093783in}}%
\pgfpathlineto{\pgfqpoint{0.854802in}{1.090140in}}%
\pgfpathlineto{\pgfqpoint{0.845589in}{1.070038in}}%
\pgfpathlineto{\pgfqpoint{0.842643in}{1.059879in}}%
\pgfpathlineto{\pgfqpoint{0.842909in}{1.052579in}}%
\pgfpathlineto{\pgfqpoint{0.839109in}{1.048947in}}%
\pgfpathlineto{\pgfqpoint{0.832280in}{1.046410in}}%
\pgfpathlineto{\pgfqpoint{0.818051in}{1.044246in}}%
\pgfpathlineto{\pgfqpoint{0.809522in}{1.041947in}}%
\pgfpathlineto{\pgfqpoint{0.807519in}{1.038421in}}%
\pgfpathlineto{\pgfqpoint{0.803014in}{1.039936in}}%
\pgfpathlineto{\pgfqpoint{0.800098in}{1.036517in}}%
\pgfpathlineto{\pgfqpoint{0.800968in}{1.030983in}}%
\pgfpathlineto{\pgfqpoint{0.798014in}{1.020963in}}%
\pgfpathlineto{\pgfqpoint{0.794371in}{1.003188in}}%
\pgfpathlineto{\pgfqpoint{0.778392in}{0.924882in}}%
\pgfpathclose%
\pgfusepath{fill}%
\end{pgfscope}%
\begin{pgfscope}%
\pgfpathrectangle{\pgfqpoint{0.100000in}{0.100000in}}{\pgfqpoint{3.007045in}{1.925000in}}%
\pgfusepath{clip}%
\pgfsetbuttcap%
\pgfsetmiterjoin%
\definecolor{currentfill}{rgb}{0.524106,0.800923,0.519646}%
\pgfsetfillcolor{currentfill}%
\pgfsetlinewidth{0.000000pt}%
\definecolor{currentstroke}{rgb}{0.000000,0.000000,0.000000}%
\pgfsetstrokecolor{currentstroke}%
\pgfsetstrokeopacity{0.000000}%
\pgfsetdash{}{0pt}%
\pgfpathmoveto{\pgfqpoint{2.480209in}{1.264505in}}%
\pgfpathlineto{\pgfqpoint{2.463037in}{1.262554in}}%
\pgfpathlineto{\pgfqpoint{2.457021in}{1.264583in}}%
\pgfpathlineto{\pgfqpoint{2.440547in}{1.262665in}}%
\pgfpathlineto{\pgfqpoint{2.441070in}{1.257982in}}%
\pgfpathlineto{\pgfqpoint{2.415184in}{1.255293in}}%
\pgfpathlineto{\pgfqpoint{2.414093in}{1.259940in}}%
\pgfpathlineto{\pgfqpoint{2.411854in}{1.281854in}}%
\pgfpathlineto{\pgfqpoint{2.409849in}{1.281241in}}%
\pgfpathlineto{\pgfqpoint{2.408904in}{1.287720in}}%
\pgfpathlineto{\pgfqpoint{2.413572in}{1.288426in}}%
\pgfpathlineto{\pgfqpoint{2.410399in}{1.309530in}}%
\pgfpathlineto{\pgfqpoint{2.434169in}{1.312710in}}%
\pgfpathlineto{\pgfqpoint{2.434855in}{1.307908in}}%
\pgfpathlineto{\pgfqpoint{2.446146in}{1.309080in}}%
\pgfpathlineto{\pgfqpoint{2.445530in}{1.314439in}}%
\pgfpathlineto{\pgfqpoint{2.461738in}{1.316890in}}%
\pgfpathlineto{\pgfqpoint{2.464330in}{1.299934in}}%
\pgfpathlineto{\pgfqpoint{2.472733in}{1.301176in}}%
\pgfpathlineto{\pgfqpoint{2.473535in}{1.295735in}}%
\pgfpathlineto{\pgfqpoint{2.476757in}{1.293301in}}%
\pgfpathlineto{\pgfqpoint{2.474283in}{1.290040in}}%
\pgfpathlineto{\pgfqpoint{2.474602in}{1.281312in}}%
\pgfpathlineto{\pgfqpoint{2.478385in}{1.281786in}}%
\pgfpathlineto{\pgfqpoint{2.480209in}{1.264505in}}%
\pgfpathclose%
\pgfusepath{fill}%
\end{pgfscope}%
\begin{pgfscope}%
\pgfpathrectangle{\pgfqpoint{0.100000in}{0.100000in}}{\pgfqpoint{3.007045in}{1.925000in}}%
\pgfusepath{clip}%
\pgfsetbuttcap%
\pgfsetmiterjoin%
\definecolor{currentfill}{rgb}{0.432941,0.757862,0.451980}%
\pgfsetfillcolor{currentfill}%
\pgfsetlinewidth{0.000000pt}%
\definecolor{currentstroke}{rgb}{0.000000,0.000000,0.000000}%
\pgfsetstrokecolor{currentstroke}%
\pgfsetstrokeopacity{0.000000}%
\pgfsetdash{}{0pt}%
\pgfpathmoveto{\pgfqpoint{2.514084in}{1.019211in}}%
\pgfpathlineto{\pgfqpoint{2.520017in}{1.015115in}}%
\pgfpathlineto{\pgfqpoint{2.522081in}{0.997367in}}%
\pgfpathlineto{\pgfqpoint{2.510978in}{0.993811in}}%
\pgfpathlineto{\pgfqpoint{2.506001in}{0.993623in}}%
\pgfpathlineto{\pgfqpoint{2.499152in}{0.989884in}}%
\pgfpathlineto{\pgfqpoint{2.494094in}{0.994158in}}%
\pgfpathlineto{\pgfqpoint{2.486521in}{0.996045in}}%
\pgfpathlineto{\pgfqpoint{2.490105in}{1.002416in}}%
\pgfpathlineto{\pgfqpoint{2.479391in}{1.010461in}}%
\pgfpathlineto{\pgfqpoint{2.473634in}{1.012683in}}%
\pgfpathlineto{\pgfqpoint{2.475252in}{1.017742in}}%
\pgfpathlineto{\pgfqpoint{2.475027in}{1.025720in}}%
\pgfpathlineto{\pgfqpoint{2.492133in}{1.027383in}}%
\pgfpathlineto{\pgfqpoint{2.492311in}{1.025615in}}%
\pgfpathlineto{\pgfqpoint{2.500787in}{1.015480in}}%
\pgfpathlineto{\pgfqpoint{2.506602in}{1.020237in}}%
\pgfpathlineto{\pgfqpoint{2.514084in}{1.019211in}}%
\pgfpathclose%
\pgfusepath{fill}%
\end{pgfscope}%
\begin{pgfscope}%
\pgfpathrectangle{\pgfqpoint{0.100000in}{0.100000in}}{\pgfqpoint{3.007045in}{1.925000in}}%
\pgfusepath{clip}%
\pgfsetbuttcap%
\pgfsetmiterjoin%
\definecolor{currentfill}{rgb}{0.389020,0.736332,0.430450}%
\pgfsetfillcolor{currentfill}%
\pgfsetlinewidth{0.000000pt}%
\definecolor{currentstroke}{rgb}{0.000000,0.000000,0.000000}%
\pgfsetstrokecolor{currentstroke}%
\pgfsetstrokeopacity{0.000000}%
\pgfsetdash{}{0pt}%
\pgfpathmoveto{\pgfqpoint{1.396518in}{1.346955in}}%
\pgfpathlineto{\pgfqpoint{1.426380in}{1.344786in}}%
\pgfpathlineto{\pgfqpoint{1.424949in}{1.327230in}}%
\pgfpathlineto{\pgfqpoint{1.428992in}{1.326930in}}%
\pgfpathlineto{\pgfqpoint{1.427276in}{1.304079in}}%
\pgfpathlineto{\pgfqpoint{1.423796in}{1.304378in}}%
\pgfpathlineto{\pgfqpoint{1.422895in}{1.292969in}}%
\pgfpathlineto{\pgfqpoint{1.395479in}{1.295197in}}%
\pgfpathlineto{\pgfqpoint{1.396484in}{1.306583in}}%
\pgfpathlineto{\pgfqpoint{1.395381in}{1.309582in}}%
\pgfpathlineto{\pgfqpoint{1.396817in}{1.327688in}}%
\pgfpathlineto{\pgfqpoint{1.395435in}{1.334971in}}%
\pgfpathlineto{\pgfqpoint{1.396518in}{1.346955in}}%
\pgfpathclose%
\pgfusepath{fill}%
\end{pgfscope}%
\begin{pgfscope}%
\pgfpathrectangle{\pgfqpoint{0.100000in}{0.100000in}}{\pgfqpoint{3.007045in}{1.925000in}}%
\pgfusepath{clip}%
\pgfsetbuttcap%
\pgfsetmiterjoin%
\definecolor{currentfill}{rgb}{0.662345,0.864022,0.638001}%
\pgfsetfillcolor{currentfill}%
\pgfsetlinewidth{0.000000pt}%
\definecolor{currentstroke}{rgb}{0.000000,0.000000,0.000000}%
\pgfsetstrokecolor{currentstroke}%
\pgfsetstrokeopacity{0.000000}%
\pgfsetdash{}{0pt}%
\pgfpathmoveto{\pgfqpoint{2.285799in}{1.004151in}}%
\pgfpathlineto{\pgfqpoint{2.285270in}{0.997505in}}%
\pgfpathlineto{\pgfqpoint{2.278118in}{0.995018in}}%
\pgfpathlineto{\pgfqpoint{2.275809in}{0.988285in}}%
\pgfpathlineto{\pgfqpoint{2.271083in}{0.992318in}}%
\pgfpathlineto{\pgfqpoint{2.263992in}{0.994959in}}%
\pgfpathlineto{\pgfqpoint{2.258356in}{0.992865in}}%
\pgfpathlineto{\pgfqpoint{2.258514in}{1.000440in}}%
\pgfpathlineto{\pgfqpoint{2.248936in}{1.000483in}}%
\pgfpathlineto{\pgfqpoint{2.247870in}{1.007837in}}%
\pgfpathlineto{\pgfqpoint{2.236965in}{1.019520in}}%
\pgfpathlineto{\pgfqpoint{2.241044in}{1.027825in}}%
\pgfpathlineto{\pgfqpoint{2.241309in}{1.038598in}}%
\pgfpathlineto{\pgfqpoint{2.234875in}{1.045920in}}%
\pgfpathlineto{\pgfqpoint{2.238225in}{1.047271in}}%
\pgfpathlineto{\pgfqpoint{2.239874in}{1.053973in}}%
\pgfpathlineto{\pgfqpoint{2.260372in}{1.054200in}}%
\pgfpathlineto{\pgfqpoint{2.259090in}{1.046090in}}%
\pgfpathlineto{\pgfqpoint{2.260460in}{1.038312in}}%
\pgfpathlineto{\pgfqpoint{2.263654in}{1.034861in}}%
\pgfpathlineto{\pgfqpoint{2.269307in}{1.034406in}}%
\pgfpathlineto{\pgfqpoint{2.271241in}{1.027836in}}%
\pgfpathlineto{\pgfqpoint{2.274407in}{1.023601in}}%
\pgfpathlineto{\pgfqpoint{2.286020in}{1.023971in}}%
\pgfpathlineto{\pgfqpoint{2.287383in}{1.019436in}}%
\pgfpathlineto{\pgfqpoint{2.285799in}{1.004151in}}%
\pgfpathclose%
\pgfusepath{fill}%
\end{pgfscope}%
\begin{pgfscope}%
\pgfpathrectangle{\pgfqpoint{0.100000in}{0.100000in}}{\pgfqpoint{3.007045in}{1.925000in}}%
\pgfusepath{clip}%
\pgfsetbuttcap%
\pgfsetmiterjoin%
\definecolor{currentfill}{rgb}{0.573933,0.824175,0.560615}%
\pgfsetfillcolor{currentfill}%
\pgfsetlinewidth{0.000000pt}%
\definecolor{currentstroke}{rgb}{0.000000,0.000000,0.000000}%
\pgfsetstrokecolor{currentstroke}%
\pgfsetstrokeopacity{0.000000}%
\pgfsetdash{}{0pt}%
\pgfpathmoveto{\pgfqpoint{2.688736in}{0.994006in}}%
\pgfpathlineto{\pgfqpoint{2.675874in}{0.992417in}}%
\pgfpathlineto{\pgfqpoint{2.672444in}{0.992955in}}%
\pgfpathlineto{\pgfqpoint{2.668059in}{1.000440in}}%
\pgfpathlineto{\pgfqpoint{2.663029in}{1.005596in}}%
\pgfpathlineto{\pgfqpoint{2.671365in}{1.033005in}}%
\pgfpathlineto{\pgfqpoint{2.679615in}{1.030971in}}%
\pgfpathlineto{\pgfqpoint{2.689119in}{1.033149in}}%
\pgfpathlineto{\pgfqpoint{2.697499in}{1.030502in}}%
\pgfpathlineto{\pgfqpoint{2.698517in}{1.027123in}}%
\pgfpathlineto{\pgfqpoint{2.705428in}{1.026285in}}%
\pgfpathlineto{\pgfqpoint{2.709850in}{1.020460in}}%
\pgfpathlineto{\pgfqpoint{2.710520in}{1.014800in}}%
\pgfpathlineto{\pgfqpoint{2.708269in}{1.015278in}}%
\pgfpathlineto{\pgfqpoint{2.688736in}{0.994006in}}%
\pgfpathclose%
\pgfusepath{fill}%
\end{pgfscope}%
\begin{pgfscope}%
\pgfpathrectangle{\pgfqpoint{0.100000in}{0.100000in}}{\pgfqpoint{3.007045in}{1.925000in}}%
\pgfusepath{clip}%
\pgfsetbuttcap%
\pgfsetmiterjoin%
\definecolor{currentfill}{rgb}{0.746498,0.899454,0.719938}%
\pgfsetfillcolor{currentfill}%
\pgfsetlinewidth{0.000000pt}%
\definecolor{currentstroke}{rgb}{0.000000,0.000000,0.000000}%
\pgfsetstrokecolor{currentstroke}%
\pgfsetstrokeopacity{0.000000}%
\pgfsetdash{}{0pt}%
\pgfpathmoveto{\pgfqpoint{0.892728in}{1.460530in}}%
\pgfpathlineto{\pgfqpoint{0.895172in}{1.457614in}}%
\pgfpathlineto{\pgfqpoint{0.903388in}{1.455868in}}%
\pgfpathlineto{\pgfqpoint{0.900100in}{1.438993in}}%
\pgfpathlineto{\pgfqpoint{0.895826in}{1.417699in}}%
\pgfpathlineto{\pgfqpoint{0.853680in}{1.425681in}}%
\pgfpathlineto{\pgfqpoint{0.845804in}{1.427734in}}%
\pgfpathlineto{\pgfqpoint{0.834301in}{1.430229in}}%
\pgfpathlineto{\pgfqpoint{0.840009in}{1.457550in}}%
\pgfpathlineto{\pgfqpoint{0.850305in}{1.455334in}}%
\pgfpathlineto{\pgfqpoint{0.851744in}{1.461596in}}%
\pgfpathlineto{\pgfqpoint{0.858463in}{1.461570in}}%
\pgfpathlineto{\pgfqpoint{0.861547in}{1.476417in}}%
\pgfpathlineto{\pgfqpoint{0.869518in}{1.474687in}}%
\pgfpathlineto{\pgfqpoint{0.873074in}{1.479709in}}%
\pgfpathlineto{\pgfqpoint{0.877767in}{1.502355in}}%
\pgfpathlineto{\pgfqpoint{0.891952in}{1.499400in}}%
\pgfpathlineto{\pgfqpoint{0.887326in}{1.476723in}}%
\pgfpathlineto{\pgfqpoint{0.884523in}{1.477298in}}%
\pgfpathlineto{\pgfqpoint{0.882125in}{1.465664in}}%
\pgfpathlineto{\pgfqpoint{0.887062in}{1.465893in}}%
\pgfpathlineto{\pgfqpoint{0.892728in}{1.460530in}}%
\pgfpathclose%
\pgfusepath{fill}%
\end{pgfscope}%
\begin{pgfscope}%
\pgfpathrectangle{\pgfqpoint{0.100000in}{0.100000in}}{\pgfqpoint{3.007045in}{1.925000in}}%
\pgfusepath{clip}%
\pgfsetbuttcap%
\pgfsetmiterjoin%
\definecolor{currentfill}{rgb}{0.596078,0.834510,0.578824}%
\pgfsetfillcolor{currentfill}%
\pgfsetlinewidth{0.000000pt}%
\definecolor{currentstroke}{rgb}{0.000000,0.000000,0.000000}%
\pgfsetstrokecolor{currentstroke}%
\pgfsetstrokeopacity{0.000000}%
\pgfsetdash{}{0pt}%
\pgfpathmoveto{\pgfqpoint{2.286587in}{0.869005in}}%
\pgfpathlineto{\pgfqpoint{2.279674in}{0.894582in}}%
\pgfpathlineto{\pgfqpoint{2.265651in}{0.893385in}}%
\pgfpathlineto{\pgfqpoint{2.264140in}{0.908633in}}%
\pgfpathlineto{\pgfqpoint{2.261069in}{0.912869in}}%
\pgfpathlineto{\pgfqpoint{2.262417in}{0.916982in}}%
\pgfpathlineto{\pgfqpoint{2.261477in}{0.928516in}}%
\pgfpathlineto{\pgfqpoint{2.264733in}{0.928140in}}%
\pgfpathlineto{\pgfqpoint{2.278483in}{0.930848in}}%
\pgfpathlineto{\pgfqpoint{2.285385in}{0.933858in}}%
\pgfpathlineto{\pgfqpoint{2.287005in}{0.929042in}}%
\pgfpathlineto{\pgfqpoint{2.297460in}{0.921002in}}%
\pgfpathlineto{\pgfqpoint{2.299822in}{0.927995in}}%
\pgfpathlineto{\pgfqpoint{2.308274in}{0.925835in}}%
\pgfpathlineto{\pgfqpoint{2.313145in}{0.917931in}}%
\pgfpathlineto{\pgfqpoint{2.309785in}{0.907283in}}%
\pgfpathlineto{\pgfqpoint{2.313442in}{0.898402in}}%
\pgfpathlineto{\pgfqpoint{2.310101in}{0.886859in}}%
\pgfpathlineto{\pgfqpoint{2.306322in}{0.886459in}}%
\pgfpathlineto{\pgfqpoint{2.305031in}{0.879714in}}%
\pgfpathlineto{\pgfqpoint{2.311402in}{0.880390in}}%
\pgfpathlineto{\pgfqpoint{2.311838in}{0.871613in}}%
\pgfpathlineto{\pgfqpoint{2.286587in}{0.869005in}}%
\pgfpathclose%
\pgfusepath{fill}%
\end{pgfscope}%
\begin{pgfscope}%
\pgfpathrectangle{\pgfqpoint{0.100000in}{0.100000in}}{\pgfqpoint{3.007045in}{1.925000in}}%
\pgfusepath{clip}%
\pgfsetbuttcap%
\pgfsetmiterjoin%
\definecolor{currentfill}{rgb}{0.590542,0.831926,0.574271}%
\pgfsetfillcolor{currentfill}%
\pgfsetlinewidth{0.000000pt}%
\definecolor{currentstroke}{rgb}{0.000000,0.000000,0.000000}%
\pgfsetstrokecolor{currentstroke}%
\pgfsetstrokeopacity{0.000000}%
\pgfsetdash{}{0pt}%
\pgfpathmoveto{\pgfqpoint{2.239459in}{1.266008in}}%
\pgfpathlineto{\pgfqpoint{2.228217in}{1.264785in}}%
\pgfpathlineto{\pgfqpoint{2.224496in}{1.264365in}}%
\pgfpathlineto{\pgfqpoint{2.221882in}{1.287153in}}%
\pgfpathlineto{\pgfqpoint{2.210917in}{1.285878in}}%
\pgfpathlineto{\pgfqpoint{2.211433in}{1.280191in}}%
\pgfpathlineto{\pgfqpoint{2.196588in}{1.278782in}}%
\pgfpathlineto{\pgfqpoint{2.195043in}{1.295965in}}%
\pgfpathlineto{\pgfqpoint{2.215428in}{1.298031in}}%
\pgfpathlineto{\pgfqpoint{2.213486in}{1.315239in}}%
\pgfpathlineto{\pgfqpoint{2.233410in}{1.317320in}}%
\pgfpathlineto{\pgfqpoint{2.234453in}{1.308012in}}%
\pgfpathlineto{\pgfqpoint{2.233530in}{1.300204in}}%
\pgfpathlineto{\pgfqpoint{2.234886in}{1.288646in}}%
\pgfpathlineto{\pgfqpoint{2.236831in}{1.288875in}}%
\pgfpathlineto{\pgfqpoint{2.239459in}{1.266008in}}%
\pgfpathclose%
\pgfusepath{fill}%
\end{pgfscope}%
\begin{pgfscope}%
\pgfpathrectangle{\pgfqpoint{0.100000in}{0.100000in}}{\pgfqpoint{3.007045in}{1.925000in}}%
\pgfusepath{clip}%
\pgfsetbuttcap%
\pgfsetmiterjoin%
\definecolor{currentfill}{rgb}{0.389020,0.736332,0.430450}%
\pgfsetfillcolor{currentfill}%
\pgfsetlinewidth{0.000000pt}%
\definecolor{currentstroke}{rgb}{0.000000,0.000000,0.000000}%
\pgfsetstrokecolor{currentstroke}%
\pgfsetstrokeopacity{0.000000}%
\pgfsetdash{}{0pt}%
\pgfpathmoveto{\pgfqpoint{0.374195in}{1.289386in}}%
\pgfpathlineto{\pgfqpoint{0.369150in}{1.288189in}}%
\pgfpathlineto{\pgfqpoint{0.366648in}{1.294372in}}%
\pgfpathlineto{\pgfqpoint{0.358679in}{1.298249in}}%
\pgfpathlineto{\pgfqpoint{0.352980in}{1.307630in}}%
\pgfpathlineto{\pgfqpoint{0.347482in}{1.307924in}}%
\pgfpathlineto{\pgfqpoint{0.344373in}{1.319665in}}%
\pgfpathlineto{\pgfqpoint{0.339643in}{1.325177in}}%
\pgfpathlineto{\pgfqpoint{0.335632in}{1.333256in}}%
\pgfpathlineto{\pgfqpoint{0.333488in}{1.341560in}}%
\pgfpathlineto{\pgfqpoint{0.327012in}{1.352333in}}%
\pgfpathlineto{\pgfqpoint{0.324629in}{1.358271in}}%
\pgfpathlineto{\pgfqpoint{0.328917in}{1.365788in}}%
\pgfpathlineto{\pgfqpoint{0.328154in}{1.376358in}}%
\pgfpathlineto{\pgfqpoint{0.328676in}{1.386959in}}%
\pgfpathlineto{\pgfqpoint{0.331174in}{1.392985in}}%
\pgfpathlineto{\pgfqpoint{0.335747in}{1.398988in}}%
\pgfpathlineto{\pgfqpoint{0.337136in}{1.407814in}}%
\pgfpathlineto{\pgfqpoint{0.337272in}{1.417649in}}%
\pgfpathlineto{\pgfqpoint{0.332528in}{1.430977in}}%
\pgfpathlineto{\pgfqpoint{0.355543in}{1.423387in}}%
\pgfpathlineto{\pgfqpoint{0.355032in}{1.421828in}}%
\pgfpathlineto{\pgfqpoint{0.384397in}{1.412475in}}%
\pgfpathlineto{\pgfqpoint{0.380635in}{1.401231in}}%
\pgfpathlineto{\pgfqpoint{0.381027in}{1.394938in}}%
\pgfpathlineto{\pgfqpoint{0.378819in}{1.386738in}}%
\pgfpathlineto{\pgfqpoint{0.386077in}{1.384484in}}%
\pgfpathlineto{\pgfqpoint{0.383653in}{1.376276in}}%
\pgfpathlineto{\pgfqpoint{0.379169in}{1.366194in}}%
\pgfpathlineto{\pgfqpoint{0.383879in}{1.360469in}}%
\pgfpathlineto{\pgfqpoint{0.389453in}{1.357782in}}%
\pgfpathlineto{\pgfqpoint{0.387533in}{1.347636in}}%
\pgfpathlineto{\pgfqpoint{0.390745in}{1.343728in}}%
\pgfpathlineto{\pgfqpoint{0.392426in}{1.337198in}}%
\pgfpathlineto{\pgfqpoint{0.388567in}{1.334258in}}%
\pgfpathlineto{\pgfqpoint{0.388700in}{1.330958in}}%
\pgfpathlineto{\pgfqpoint{0.382079in}{1.325306in}}%
\pgfpathlineto{\pgfqpoint{0.373315in}{1.325445in}}%
\pgfpathlineto{\pgfqpoint{0.371022in}{1.321407in}}%
\pgfpathlineto{\pgfqpoint{0.376145in}{1.309067in}}%
\pgfpathlineto{\pgfqpoint{0.377641in}{1.299051in}}%
\pgfpathlineto{\pgfqpoint{0.377156in}{1.290848in}}%
\pgfpathlineto{\pgfqpoint{0.374195in}{1.289386in}}%
\pgfpathclose%
\pgfusepath{fill}%
\end{pgfscope}%
\begin{pgfscope}%
\pgfpathrectangle{\pgfqpoint{0.100000in}{0.100000in}}{\pgfqpoint{3.007045in}{1.925000in}}%
\pgfusepath{clip}%
\pgfsetbuttcap%
\pgfsetmiterjoin%
\definecolor{currentfill}{rgb}{0.439216,0.760938,0.455056}%
\pgfsetfillcolor{currentfill}%
\pgfsetlinewidth{0.000000pt}%
\definecolor{currentstroke}{rgb}{0.000000,0.000000,0.000000}%
\pgfsetstrokecolor{currentstroke}%
\pgfsetstrokeopacity{0.000000}%
\pgfsetdash{}{0pt}%
\pgfpathmoveto{\pgfqpoint{2.458232in}{1.094688in}}%
\pgfpathlineto{\pgfqpoint{2.446820in}{1.093542in}}%
\pgfpathlineto{\pgfqpoint{2.441207in}{1.098026in}}%
\pgfpathlineto{\pgfqpoint{2.440031in}{1.105220in}}%
\pgfpathlineto{\pgfqpoint{2.435308in}{1.112304in}}%
\pgfpathlineto{\pgfqpoint{2.429643in}{1.109432in}}%
\pgfpathlineto{\pgfqpoint{2.433994in}{1.115719in}}%
\pgfpathlineto{\pgfqpoint{2.428858in}{1.122661in}}%
\pgfpathlineto{\pgfqpoint{2.428704in}{1.126732in}}%
\pgfpathlineto{\pgfqpoint{2.431774in}{1.134957in}}%
\pgfpathlineto{\pgfqpoint{2.438948in}{1.139413in}}%
\pgfpathlineto{\pgfqpoint{2.437602in}{1.145887in}}%
\pgfpathlineto{\pgfqpoint{2.435873in}{1.154519in}}%
\pgfpathlineto{\pgfqpoint{2.440950in}{1.158535in}}%
\pgfpathlineto{\pgfqpoint{2.450037in}{1.161421in}}%
\pgfpathlineto{\pgfqpoint{2.458391in}{1.154238in}}%
\pgfpathlineto{\pgfqpoint{2.463450in}{1.158871in}}%
\pgfpathlineto{\pgfqpoint{2.466955in}{1.155210in}}%
\pgfpathlineto{\pgfqpoint{2.480993in}{1.155854in}}%
\pgfpathlineto{\pgfqpoint{2.487851in}{1.166255in}}%
\pgfpathlineto{\pgfqpoint{2.494103in}{1.161327in}}%
\pgfpathlineto{\pgfqpoint{2.497293in}{1.156805in}}%
\pgfpathlineto{\pgfqpoint{2.496539in}{1.150737in}}%
\pgfpathlineto{\pgfqpoint{2.481701in}{1.138352in}}%
\pgfpathlineto{\pgfqpoint{2.477515in}{1.130920in}}%
\pgfpathlineto{\pgfqpoint{2.476974in}{1.118028in}}%
\pgfpathlineto{\pgfqpoint{2.470294in}{1.118309in}}%
\pgfpathlineto{\pgfqpoint{2.467580in}{1.113861in}}%
\pgfpathlineto{\pgfqpoint{2.471855in}{1.105407in}}%
\pgfpathlineto{\pgfqpoint{2.459342in}{1.100923in}}%
\pgfpathlineto{\pgfqpoint{2.463061in}{1.098120in}}%
\pgfpathlineto{\pgfqpoint{2.458232in}{1.094688in}}%
\pgfpathclose%
\pgfusepath{fill}%
\end{pgfscope}%
\begin{pgfscope}%
\pgfpathrectangle{\pgfqpoint{0.100000in}{0.100000in}}{\pgfqpoint{3.007045in}{1.925000in}}%
\pgfusepath{clip}%
\pgfsetbuttcap%
\pgfsetmiterjoin%
\definecolor{currentfill}{rgb}{0.727797,0.891580,0.701730}%
\pgfsetfillcolor{currentfill}%
\pgfsetlinewidth{0.000000pt}%
\definecolor{currentstroke}{rgb}{0.000000,0.000000,0.000000}%
\pgfsetstrokecolor{currentstroke}%
\pgfsetstrokeopacity{0.000000}%
\pgfsetdash{}{0pt}%
\pgfpathmoveto{\pgfqpoint{1.298841in}{1.455523in}}%
\pgfpathlineto{\pgfqpoint{1.293878in}{1.409890in}}%
\pgfpathlineto{\pgfqpoint{1.292810in}{1.397021in}}%
\pgfpathlineto{\pgfqpoint{1.273764in}{1.398937in}}%
\pgfpathlineto{\pgfqpoint{1.272590in}{1.387487in}}%
\pgfpathlineto{\pgfqpoint{1.270231in}{1.387736in}}%
\pgfpathlineto{\pgfqpoint{1.267095in}{1.378734in}}%
\pgfpathlineto{\pgfqpoint{1.255653in}{1.381012in}}%
\pgfpathlineto{\pgfqpoint{1.256175in}{1.385798in}}%
\pgfpathlineto{\pgfqpoint{1.259419in}{1.388949in}}%
\pgfpathlineto{\pgfqpoint{1.234424in}{1.392114in}}%
\pgfpathlineto{\pgfqpoint{1.207752in}{1.395280in}}%
\pgfpathlineto{\pgfqpoint{1.164443in}{1.401325in}}%
\pgfpathlineto{\pgfqpoint{1.163522in}{1.401453in}}%
\pgfpathlineto{\pgfqpoint{1.166601in}{1.424158in}}%
\pgfpathlineto{\pgfqpoint{1.168621in}{1.423872in}}%
\pgfpathlineto{\pgfqpoint{1.171716in}{1.446544in}}%
\pgfpathlineto{\pgfqpoint{1.174139in}{1.469186in}}%
\pgfpathlineto{\pgfqpoint{1.173551in}{1.471182in}}%
\pgfpathlineto{\pgfqpoint{1.193681in}{1.468364in}}%
\pgfpathlineto{\pgfqpoint{1.218116in}{1.464571in}}%
\pgfpathlineto{\pgfqpoint{1.245609in}{1.461365in}}%
\pgfpathlineto{\pgfqpoint{1.298841in}{1.455523in}}%
\pgfpathclose%
\pgfusepath{fill}%
\end{pgfscope}%
\begin{pgfscope}%
\pgfpathrectangle{\pgfqpoint{0.100000in}{0.100000in}}{\pgfqpoint{3.007045in}{1.925000in}}%
\pgfusepath{clip}%
\pgfsetbuttcap%
\pgfsetmiterjoin%
\definecolor{currentfill}{rgb}{0.496424,0.788005,0.496886}%
\pgfsetfillcolor{currentfill}%
\pgfsetlinewidth{0.000000pt}%
\definecolor{currentstroke}{rgb}{0.000000,0.000000,0.000000}%
\pgfsetstrokecolor{currentstroke}%
\pgfsetstrokeopacity{0.000000}%
\pgfsetdash{}{0pt}%
\pgfpathmoveto{\pgfqpoint{1.756771in}{1.198935in}}%
\pgfpathlineto{\pgfqpoint{1.755176in}{1.198941in}}%
\pgfpathlineto{\pgfqpoint{1.709799in}{1.199381in}}%
\pgfpathlineto{\pgfqpoint{1.698475in}{1.199567in}}%
\pgfpathlineto{\pgfqpoint{1.699296in}{1.251149in}}%
\pgfpathlineto{\pgfqpoint{1.730711in}{1.250739in}}%
\pgfpathlineto{\pgfqpoint{1.727958in}{1.247635in}}%
\pgfpathlineto{\pgfqpoint{1.732568in}{1.242722in}}%
\pgfpathlineto{\pgfqpoint{1.734017in}{1.237625in}}%
\pgfpathlineto{\pgfqpoint{1.741287in}{1.219682in}}%
\pgfpathlineto{\pgfqpoint{1.747673in}{1.215377in}}%
\pgfpathlineto{\pgfqpoint{1.748211in}{1.211203in}}%
\pgfpathlineto{\pgfqpoint{1.751959in}{1.207454in}}%
\pgfpathlineto{\pgfqpoint{1.751069in}{1.202842in}}%
\pgfpathlineto{\pgfqpoint{1.756771in}{1.198935in}}%
\pgfpathclose%
\pgfusepath{fill}%
\end{pgfscope}%
\begin{pgfscope}%
\pgfpathrectangle{\pgfqpoint{0.100000in}{0.100000in}}{\pgfqpoint{3.007045in}{1.925000in}}%
\pgfusepath{clip}%
\pgfsetbuttcap%
\pgfsetmiterjoin%
\definecolor{currentfill}{rgb}{0.827451,0.932549,0.803137}%
\pgfsetfillcolor{currentfill}%
\pgfsetlinewidth{0.000000pt}%
\definecolor{currentstroke}{rgb}{0.000000,0.000000,0.000000}%
\pgfsetstrokecolor{currentstroke}%
\pgfsetstrokeopacity{0.000000}%
\pgfsetdash{}{0pt}%
\pgfpathmoveto{\pgfqpoint{1.370186in}{0.721452in}}%
\pgfpathlineto{\pgfqpoint{1.322173in}{0.725422in}}%
\pgfpathlineto{\pgfqpoint{1.319633in}{0.696676in}}%
\pgfpathlineto{\pgfqpoint{1.319120in}{0.690958in}}%
\pgfpathlineto{\pgfqpoint{1.282225in}{0.694354in}}%
\pgfpathlineto{\pgfqpoint{1.285541in}{0.728841in}}%
\pgfpathlineto{\pgfqpoint{1.280477in}{0.729323in}}%
\pgfpathlineto{\pgfqpoint{1.283283in}{0.758525in}}%
\pgfpathlineto{\pgfqpoint{1.285953in}{0.758267in}}%
\pgfpathlineto{\pgfqpoint{1.288691in}{0.786623in}}%
\pgfpathlineto{\pgfqpoint{1.291172in}{0.786411in}}%
\pgfpathlineto{\pgfqpoint{1.292309in}{0.797855in}}%
\pgfpathlineto{\pgfqpoint{1.303863in}{0.796794in}}%
\pgfpathlineto{\pgfqpoint{1.328991in}{0.794521in}}%
\pgfpathlineto{\pgfqpoint{1.327701in}{0.782533in}}%
\pgfpathlineto{\pgfqpoint{1.381713in}{0.778104in}}%
\pgfpathlineto{\pgfqpoint{1.379483in}{0.749728in}}%
\pgfpathlineto{\pgfqpoint{1.372150in}{0.750256in}}%
\pgfpathlineto{\pgfqpoint{1.370186in}{0.721452in}}%
\pgfpathclose%
\pgfusepath{fill}%
\end{pgfscope}%
\begin{pgfscope}%
\pgfpathrectangle{\pgfqpoint{0.100000in}{0.100000in}}{\pgfqpoint{3.007045in}{1.925000in}}%
\pgfusepath{clip}%
\pgfsetbuttcap%
\pgfsetmiterjoin%
\definecolor{currentfill}{rgb}{0.732472,0.893549,0.706282}%
\pgfsetfillcolor{currentfill}%
\pgfsetlinewidth{0.000000pt}%
\definecolor{currentstroke}{rgb}{0.000000,0.000000,0.000000}%
\pgfsetstrokecolor{currentstroke}%
\pgfsetstrokeopacity{0.000000}%
\pgfsetdash{}{0pt}%
\pgfpathmoveto{\pgfqpoint{2.139901in}{1.326614in}}%
\pgfpathlineto{\pgfqpoint{2.141986in}{1.299655in}}%
\pgfpathlineto{\pgfqpoint{2.117970in}{1.297350in}}%
\pgfpathlineto{\pgfqpoint{2.118514in}{1.291662in}}%
\pgfpathlineto{\pgfqpoint{2.107035in}{1.290511in}}%
\pgfpathlineto{\pgfqpoint{2.107101in}{1.284740in}}%
\pgfpathlineto{\pgfqpoint{2.090562in}{1.283138in}}%
\pgfpathlineto{\pgfqpoint{2.087331in}{1.317506in}}%
\pgfpathlineto{\pgfqpoint{2.085344in}{1.346129in}}%
\pgfpathlineto{\pgfqpoint{2.085588in}{1.351908in}}%
\pgfpathlineto{\pgfqpoint{2.079874in}{1.351505in}}%
\pgfpathlineto{\pgfqpoint{2.078237in}{1.373879in}}%
\pgfpathlineto{\pgfqpoint{2.122132in}{1.377026in}}%
\pgfpathlineto{\pgfqpoint{2.121650in}{1.363916in}}%
\pgfpathlineto{\pgfqpoint{2.123749in}{1.358371in}}%
\pgfpathlineto{\pgfqpoint{2.130184in}{1.350131in}}%
\pgfpathlineto{\pgfqpoint{2.134976in}{1.335269in}}%
\pgfpathlineto{\pgfqpoint{2.139901in}{1.326614in}}%
\pgfpathclose%
\pgfusepath{fill}%
\end{pgfscope}%
\begin{pgfscope}%
\pgfpathrectangle{\pgfqpoint{0.100000in}{0.100000in}}{\pgfqpoint{3.007045in}{1.925000in}}%
\pgfusepath{clip}%
\pgfsetbuttcap%
\pgfsetmiterjoin%
\definecolor{currentfill}{rgb}{0.376471,0.730181,0.424298}%
\pgfsetfillcolor{currentfill}%
\pgfsetlinewidth{0.000000pt}%
\definecolor{currentstroke}{rgb}{0.000000,0.000000,0.000000}%
\pgfsetstrokecolor{currentstroke}%
\pgfsetstrokeopacity{0.000000}%
\pgfsetdash{}{0pt}%
\pgfpathmoveto{\pgfqpoint{1.652076in}{1.171950in}}%
\pgfpathlineto{\pgfqpoint{1.651571in}{1.154782in}}%
\pgfpathlineto{\pgfqpoint{1.623055in}{1.155630in}}%
\pgfpathlineto{\pgfqpoint{1.623649in}{1.172778in}}%
\pgfpathlineto{\pgfqpoint{1.594392in}{1.173857in}}%
\pgfpathlineto{\pgfqpoint{1.595536in}{1.202468in}}%
\pgfpathlineto{\pgfqpoint{1.607157in}{1.202049in}}%
\pgfpathlineto{\pgfqpoint{1.652796in}{1.200602in}}%
\pgfpathlineto{\pgfqpoint{1.652076in}{1.171950in}}%
\pgfpathclose%
\pgfusepath{fill}%
\end{pgfscope}%
\begin{pgfscope}%
\pgfpathrectangle{\pgfqpoint{0.100000in}{0.100000in}}{\pgfqpoint{3.007045in}{1.925000in}}%
\pgfusepath{clip}%
\pgfsetbuttcap%
\pgfsetmiterjoin%
\definecolor{currentfill}{rgb}{0.657670,0.862053,0.633449}%
\pgfsetfillcolor{currentfill}%
\pgfsetlinewidth{0.000000pt}%
\definecolor{currentstroke}{rgb}{0.000000,0.000000,0.000000}%
\pgfsetstrokecolor{currentstroke}%
\pgfsetstrokeopacity{0.000000}%
\pgfsetdash{}{0pt}%
\pgfpathmoveto{\pgfqpoint{1.426380in}{1.344786in}}%
\pgfpathlineto{\pgfqpoint{1.396518in}{1.346955in}}%
\pgfpathlineto{\pgfqpoint{1.393848in}{1.352944in}}%
\pgfpathlineto{\pgfqpoint{1.395452in}{1.375827in}}%
\pgfpathlineto{\pgfqpoint{1.394290in}{1.375922in}}%
\pgfpathlineto{\pgfqpoint{1.396236in}{1.398628in}}%
\pgfpathlineto{\pgfqpoint{1.396545in}{1.412828in}}%
\pgfpathlineto{\pgfqpoint{1.386507in}{1.413753in}}%
\pgfpathlineto{\pgfqpoint{1.390013in}{1.453572in}}%
\pgfpathlineto{\pgfqpoint{1.395146in}{1.457695in}}%
\pgfpathlineto{\pgfqpoint{1.399572in}{1.458189in}}%
\pgfpathlineto{\pgfqpoint{1.429680in}{1.455658in}}%
\pgfpathlineto{\pgfqpoint{1.431599in}{1.456389in}}%
\pgfpathlineto{\pgfqpoint{1.430442in}{1.441625in}}%
\pgfpathlineto{\pgfqpoint{1.431343in}{1.435804in}}%
\pgfpathlineto{\pgfqpoint{1.473591in}{1.432656in}}%
\pgfpathlineto{\pgfqpoint{1.471839in}{1.406868in}}%
\pgfpathlineto{\pgfqpoint{1.430688in}{1.409985in}}%
\pgfpathlineto{\pgfqpoint{1.430345in}{1.395850in}}%
\pgfpathlineto{\pgfqpoint{1.428602in}{1.373346in}}%
\pgfpathlineto{\pgfqpoint{1.429987in}{1.373229in}}%
\pgfpathlineto{\pgfqpoint{1.428115in}{1.350443in}}%
\pgfpathlineto{\pgfqpoint{1.426380in}{1.344786in}}%
\pgfpathclose%
\pgfusepath{fill}%
\end{pgfscope}%
\begin{pgfscope}%
\pgfpathrectangle{\pgfqpoint{0.100000in}{0.100000in}}{\pgfqpoint{3.007045in}{1.925000in}}%
\pgfusepath{clip}%
\pgfsetbuttcap%
\pgfsetmiterjoin%
\definecolor{currentfill}{rgb}{0.252595,0.668128,0.362860}%
\pgfsetfillcolor{currentfill}%
\pgfsetlinewidth{0.000000pt}%
\definecolor{currentstroke}{rgb}{0.000000,0.000000,0.000000}%
\pgfsetstrokecolor{currentstroke}%
\pgfsetstrokeopacity{0.000000}%
\pgfsetdash{}{0pt}%
\pgfpathmoveto{\pgfqpoint{1.797461in}{1.775190in}}%
\pgfpathlineto{\pgfqpoint{1.804817in}{1.775072in}}%
\pgfpathlineto{\pgfqpoint{1.805689in}{1.772077in}}%
\pgfpathlineto{\pgfqpoint{1.823819in}{1.770365in}}%
\pgfpathlineto{\pgfqpoint{1.825676in}{1.763110in}}%
\pgfpathlineto{\pgfqpoint{1.840111in}{1.765309in}}%
\pgfpathlineto{\pgfqpoint{1.840141in}{1.768138in}}%
\pgfpathlineto{\pgfqpoint{1.851487in}{1.771989in}}%
\pgfpathlineto{\pgfqpoint{1.856761in}{1.771102in}}%
\pgfpathlineto{\pgfqpoint{1.857189in}{1.733722in}}%
\pgfpathlineto{\pgfqpoint{1.858121in}{1.722001in}}%
\pgfpathlineto{\pgfqpoint{1.827059in}{1.721935in}}%
\pgfpathlineto{\pgfqpoint{1.827086in}{1.718431in}}%
\pgfpathlineto{\pgfqpoint{1.798346in}{1.718105in}}%
\pgfpathlineto{\pgfqpoint{1.797671in}{1.752902in}}%
\pgfpathlineto{\pgfqpoint{1.797461in}{1.775190in}}%
\pgfpathclose%
\pgfusepath{fill}%
\end{pgfscope}%
\begin{pgfscope}%
\pgfpathrectangle{\pgfqpoint{0.100000in}{0.100000in}}{\pgfqpoint{3.007045in}{1.925000in}}%
\pgfusepath{clip}%
\pgfsetbuttcap%
\pgfsetmiterjoin%
\definecolor{currentfill}{rgb}{0.667020,0.865990,0.642553}%
\pgfsetfillcolor{currentfill}%
\pgfsetlinewidth{0.000000pt}%
\definecolor{currentstroke}{rgb}{0.000000,0.000000,0.000000}%
\pgfsetstrokecolor{currentstroke}%
\pgfsetstrokeopacity{0.000000}%
\pgfsetdash{}{0pt}%
\pgfpathmoveto{\pgfqpoint{2.145056in}{1.262796in}}%
\pgfpathlineto{\pgfqpoint{2.146371in}{1.246681in}}%
\pgfpathlineto{\pgfqpoint{2.125964in}{1.244692in}}%
\pgfpathlineto{\pgfqpoint{2.116838in}{1.244113in}}%
\pgfpathlineto{\pgfqpoint{2.113613in}{1.277542in}}%
\pgfpathlineto{\pgfqpoint{2.107905in}{1.276923in}}%
\pgfpathlineto{\pgfqpoint{2.107101in}{1.284740in}}%
\pgfpathlineto{\pgfqpoint{2.107035in}{1.290511in}}%
\pgfpathlineto{\pgfqpoint{2.118514in}{1.291662in}}%
\pgfpathlineto{\pgfqpoint{2.117970in}{1.297350in}}%
\pgfpathlineto{\pgfqpoint{2.141986in}{1.299655in}}%
\pgfpathlineto{\pgfqpoint{2.145056in}{1.262796in}}%
\pgfpathclose%
\pgfusepath{fill}%
\end{pgfscope}%
\begin{pgfscope}%
\pgfpathrectangle{\pgfqpoint{0.100000in}{0.100000in}}{\pgfqpoint{3.007045in}{1.925000in}}%
\pgfusepath{clip}%
\pgfsetbuttcap%
\pgfsetmiterjoin%
\definecolor{currentfill}{rgb}{0.648320,0.858116,0.624344}%
\pgfsetfillcolor{currentfill}%
\pgfsetlinewidth{0.000000pt}%
\definecolor{currentstroke}{rgb}{0.000000,0.000000,0.000000}%
\pgfsetstrokecolor{currentstroke}%
\pgfsetstrokeopacity{0.000000}%
\pgfsetdash{}{0pt}%
\pgfpathmoveto{\pgfqpoint{2.060711in}{0.508116in}}%
\pgfpathlineto{\pgfqpoint{2.054795in}{0.510464in}}%
\pgfpathlineto{\pgfqpoint{2.051671in}{0.516053in}}%
\pgfpathlineto{\pgfqpoint{2.053131in}{0.522603in}}%
\pgfpathlineto{\pgfqpoint{2.050269in}{0.524553in}}%
\pgfpathlineto{\pgfqpoint{2.042354in}{0.524234in}}%
\pgfpathlineto{\pgfqpoint{2.040150in}{0.528704in}}%
\pgfpathlineto{\pgfqpoint{2.034075in}{0.531818in}}%
\pgfpathlineto{\pgfqpoint{2.030059in}{0.536911in}}%
\pgfpathlineto{\pgfqpoint{2.010082in}{0.536853in}}%
\pgfpathlineto{\pgfqpoint{2.005579in}{0.541500in}}%
\pgfpathlineto{\pgfqpoint{2.005096in}{0.547382in}}%
\pgfpathlineto{\pgfqpoint{2.007733in}{0.550626in}}%
\pgfpathlineto{\pgfqpoint{2.023323in}{0.554899in}}%
\pgfpathlineto{\pgfqpoint{2.029617in}{0.558203in}}%
\pgfpathlineto{\pgfqpoint{2.032903in}{0.562233in}}%
\pgfpathlineto{\pgfqpoint{2.042685in}{0.564352in}}%
\pgfpathlineto{\pgfqpoint{2.046065in}{0.559905in}}%
\pgfpathlineto{\pgfqpoint{2.043709in}{0.592924in}}%
\pgfpathlineto{\pgfqpoint{2.037435in}{0.611183in}}%
\pgfpathlineto{\pgfqpoint{2.072666in}{0.613237in}}%
\pgfpathlineto{\pgfqpoint{2.070931in}{0.604028in}}%
\pgfpathlineto{\pgfqpoint{2.067754in}{0.598346in}}%
\pgfpathlineto{\pgfqpoint{2.067276in}{0.590315in}}%
\pgfpathlineto{\pgfqpoint{2.072286in}{0.580627in}}%
\pgfpathlineto{\pgfqpoint{2.077352in}{0.576813in}}%
\pgfpathlineto{\pgfqpoint{2.082124in}{0.561840in}}%
\pgfpathlineto{\pgfqpoint{2.087534in}{0.559447in}}%
\pgfpathlineto{\pgfqpoint{2.082198in}{0.557399in}}%
\pgfpathlineto{\pgfqpoint{2.076836in}{0.548759in}}%
\pgfpathlineto{\pgfqpoint{2.071543in}{0.549265in}}%
\pgfpathlineto{\pgfqpoint{2.070214in}{0.542954in}}%
\pgfpathlineto{\pgfqpoint{2.075634in}{0.542956in}}%
\pgfpathlineto{\pgfqpoint{2.077930in}{0.538591in}}%
\pgfpathlineto{\pgfqpoint{2.084493in}{0.539170in}}%
\pgfpathlineto{\pgfqpoint{2.085010in}{0.546770in}}%
\pgfpathlineto{\pgfqpoint{2.089855in}{0.550151in}}%
\pgfpathlineto{\pgfqpoint{2.100634in}{0.542949in}}%
\pgfpathlineto{\pgfqpoint{2.094872in}{0.530452in}}%
\pgfpathlineto{\pgfqpoint{2.081945in}{0.530204in}}%
\pgfpathlineto{\pgfqpoint{2.084051in}{0.524154in}}%
\pgfpathlineto{\pgfqpoint{2.081293in}{0.516251in}}%
\pgfpathlineto{\pgfqpoint{2.087768in}{0.513497in}}%
\pgfpathlineto{\pgfqpoint{2.093764in}{0.507672in}}%
\pgfpathlineto{\pgfqpoint{2.109818in}{0.505073in}}%
\pgfpathlineto{\pgfqpoint{2.117129in}{0.495373in}}%
\pgfpathlineto{\pgfqpoint{2.122084in}{0.494513in}}%
\pgfpathlineto{\pgfqpoint{2.119160in}{0.486899in}}%
\pgfpathlineto{\pgfqpoint{2.113765in}{0.483500in}}%
\pgfpathlineto{\pgfqpoint{2.109844in}{0.485263in}}%
\pgfpathlineto{\pgfqpoint{2.102674in}{0.484267in}}%
\pgfpathlineto{\pgfqpoint{2.093049in}{0.495483in}}%
\pgfpathlineto{\pgfqpoint{2.079626in}{0.500114in}}%
\pgfpathlineto{\pgfqpoint{2.072262in}{0.500507in}}%
\pgfpathlineto{\pgfqpoint{2.073553in}{0.504658in}}%
\pgfpathlineto{\pgfqpoint{2.072286in}{0.510828in}}%
\pgfpathlineto{\pgfqpoint{2.060711in}{0.508116in}}%
\pgfpathclose%
\pgfusepath{fill}%
\end{pgfscope}%
\begin{pgfscope}%
\pgfpathrectangle{\pgfqpoint{0.100000in}{0.100000in}}{\pgfqpoint{3.007045in}{1.925000in}}%
\pgfusepath{clip}%
\pgfsetbuttcap%
\pgfsetmiterjoin%
\definecolor{currentfill}{rgb}{0.568397,0.821592,0.556063}%
\pgfsetfillcolor{currentfill}%
\pgfsetlinewidth{0.000000pt}%
\definecolor{currentstroke}{rgb}{0.000000,0.000000,0.000000}%
\pgfsetstrokecolor{currentstroke}%
\pgfsetstrokeopacity{0.000000}%
\pgfsetdash{}{0pt}%
\pgfpathmoveto{\pgfqpoint{1.938345in}{1.353373in}}%
\pgfpathlineto{\pgfqpoint{1.904114in}{1.352301in}}%
\pgfpathlineto{\pgfqpoint{1.881224in}{1.351626in}}%
\pgfpathlineto{\pgfqpoint{1.881395in}{1.345922in}}%
\pgfpathlineto{\pgfqpoint{1.869963in}{1.345597in}}%
\pgfpathlineto{\pgfqpoint{1.869522in}{1.362845in}}%
\pgfpathlineto{\pgfqpoint{1.868339in}{1.362827in}}%
\pgfpathlineto{\pgfqpoint{1.867818in}{1.391658in}}%
\pgfpathlineto{\pgfqpoint{1.913408in}{1.392837in}}%
\pgfpathlineto{\pgfqpoint{1.913064in}{1.404438in}}%
\pgfpathlineto{\pgfqpoint{1.936000in}{1.405159in}}%
\pgfpathlineto{\pgfqpoint{1.937355in}{1.364751in}}%
\pgfpathlineto{\pgfqpoint{1.938345in}{1.353373in}}%
\pgfpathclose%
\pgfusepath{fill}%
\end{pgfscope}%
\begin{pgfscope}%
\pgfpathrectangle{\pgfqpoint{0.100000in}{0.100000in}}{\pgfqpoint{3.007045in}{1.925000in}}%
\pgfusepath{clip}%
\pgfsetbuttcap%
\pgfsetmiterjoin%
\definecolor{currentfill}{rgb}{0.838524,0.936978,0.814948}%
\pgfsetfillcolor{currentfill}%
\pgfsetlinewidth{0.000000pt}%
\definecolor{currentstroke}{rgb}{0.000000,0.000000,0.000000}%
\pgfsetstrokecolor{currentstroke}%
\pgfsetstrokeopacity{0.000000}%
\pgfsetdash{}{0pt}%
\pgfpathmoveto{\pgfqpoint{2.405437in}{0.788426in}}%
\pgfpathlineto{\pgfqpoint{2.390888in}{0.783936in}}%
\pgfpathlineto{\pgfqpoint{2.390153in}{0.787078in}}%
\pgfpathlineto{\pgfqpoint{2.390080in}{0.791606in}}%
\pgfpathlineto{\pgfqpoint{2.386212in}{0.799739in}}%
\pgfpathlineto{\pgfqpoint{2.382395in}{0.804196in}}%
\pgfpathlineto{\pgfqpoint{2.375369in}{0.793380in}}%
\pgfpathlineto{\pgfqpoint{2.365423in}{0.795571in}}%
\pgfpathlineto{\pgfqpoint{2.357695in}{0.794657in}}%
\pgfpathlineto{\pgfqpoint{2.355976in}{0.788053in}}%
\pgfpathlineto{\pgfqpoint{2.352410in}{0.787703in}}%
\pgfpathlineto{\pgfqpoint{2.346769in}{0.793773in}}%
\pgfpathlineto{\pgfqpoint{2.343332in}{0.803548in}}%
\pgfpathlineto{\pgfqpoint{2.332358in}{0.802455in}}%
\pgfpathlineto{\pgfqpoint{2.327623in}{0.805924in}}%
\pgfpathlineto{\pgfqpoint{2.326336in}{0.819580in}}%
\pgfpathlineto{\pgfqpoint{2.318665in}{0.820809in}}%
\pgfpathlineto{\pgfqpoint{2.317320in}{0.826818in}}%
\pgfpathlineto{\pgfqpoint{2.320886in}{0.830332in}}%
\pgfpathlineto{\pgfqpoint{2.322977in}{0.839287in}}%
\pgfpathlineto{\pgfqpoint{2.337348in}{0.840608in}}%
\pgfpathlineto{\pgfqpoint{2.334077in}{0.871563in}}%
\pgfpathlineto{\pgfqpoint{2.349378in}{0.873307in}}%
\pgfpathlineto{\pgfqpoint{2.351405in}{0.875886in}}%
\pgfpathlineto{\pgfqpoint{2.358828in}{0.878027in}}%
\pgfpathlineto{\pgfqpoint{2.359050in}{0.873951in}}%
\pgfpathlineto{\pgfqpoint{2.365118in}{0.869291in}}%
\pgfpathlineto{\pgfqpoint{2.371582in}{0.867326in}}%
\pgfpathlineto{\pgfqpoint{2.375538in}{0.858536in}}%
\pgfpathlineto{\pgfqpoint{2.372820in}{0.852600in}}%
\pgfpathlineto{\pgfqpoint{2.380206in}{0.847045in}}%
\pgfpathlineto{\pgfqpoint{2.382703in}{0.849302in}}%
\pgfpathlineto{\pgfqpoint{2.380905in}{0.840878in}}%
\pgfpathlineto{\pgfqpoint{2.387448in}{0.834640in}}%
\pgfpathlineto{\pgfqpoint{2.392091in}{0.836488in}}%
\pgfpathlineto{\pgfqpoint{2.403550in}{0.830476in}}%
\pgfpathlineto{\pgfqpoint{2.408663in}{0.823968in}}%
\pgfpathlineto{\pgfqpoint{2.417379in}{0.810753in}}%
\pgfpathlineto{\pgfqpoint{2.403890in}{0.805755in}}%
\pgfpathlineto{\pgfqpoint{2.405437in}{0.788426in}}%
\pgfpathclose%
\pgfusepath{fill}%
\end{pgfscope}%
\begin{pgfscope}%
\pgfpathrectangle{\pgfqpoint{0.100000in}{0.100000in}}{\pgfqpoint{3.007045in}{1.925000in}}%
\pgfusepath{clip}%
\pgfsetbuttcap%
\pgfsetmiterjoin%
\definecolor{currentfill}{rgb}{0.189850,0.601200,0.312664}%
\pgfsetfillcolor{currentfill}%
\pgfsetlinewidth{0.000000pt}%
\definecolor{currentstroke}{rgb}{0.000000,0.000000,0.000000}%
\pgfsetstrokecolor{currentstroke}%
\pgfsetstrokeopacity{0.000000}%
\pgfsetdash{}{0pt}%
\pgfpathmoveto{\pgfqpoint{2.418790in}{0.913875in}}%
\pgfpathlineto{\pgfqpoint{2.413670in}{0.911416in}}%
\pgfpathlineto{\pgfqpoint{2.410454in}{0.917486in}}%
\pgfpathlineto{\pgfqpoint{2.406157in}{0.920864in}}%
\pgfpathlineto{\pgfqpoint{2.398481in}{0.931422in}}%
\pgfpathlineto{\pgfqpoint{2.392521in}{0.928792in}}%
\pgfpathlineto{\pgfqpoint{2.381105in}{0.925986in}}%
\pgfpathlineto{\pgfqpoint{2.379773in}{0.923563in}}%
\pgfpathlineto{\pgfqpoint{2.372864in}{0.923192in}}%
\pgfpathlineto{\pgfqpoint{2.366318in}{0.919795in}}%
\pgfpathlineto{\pgfqpoint{2.361956in}{0.924269in}}%
\pgfpathlineto{\pgfqpoint{2.361650in}{0.931386in}}%
\pgfpathlineto{\pgfqpoint{2.364223in}{0.935952in}}%
\pgfpathlineto{\pgfqpoint{2.373522in}{0.943638in}}%
\pgfpathlineto{\pgfqpoint{2.379259in}{0.944820in}}%
\pgfpathlineto{\pgfqpoint{2.388077in}{0.945524in}}%
\pgfpathlineto{\pgfqpoint{2.399943in}{0.955887in}}%
\pgfpathlineto{\pgfqpoint{2.403998in}{0.954773in}}%
\pgfpathlineto{\pgfqpoint{2.405701in}{0.945045in}}%
\pgfpathlineto{\pgfqpoint{2.420735in}{0.932481in}}%
\pgfpathlineto{\pgfqpoint{2.421150in}{0.927599in}}%
\pgfpathlineto{\pgfqpoint{2.416370in}{0.915795in}}%
\pgfpathlineto{\pgfqpoint{2.418790in}{0.913875in}}%
\pgfpathclose%
\pgfusepath{fill}%
\end{pgfscope}%
\begin{pgfscope}%
\pgfpathrectangle{\pgfqpoint{0.100000in}{0.100000in}}{\pgfqpoint{3.007045in}{1.925000in}}%
\pgfusepath{clip}%
\pgfsetbuttcap%
\pgfsetmiterjoin%
\definecolor{currentfill}{rgb}{0.501961,0.790588,0.501438}%
\pgfsetfillcolor{currentfill}%
\pgfsetlinewidth{0.000000pt}%
\definecolor{currentstroke}{rgb}{0.000000,0.000000,0.000000}%
\pgfsetstrokecolor{currentstroke}%
\pgfsetstrokeopacity{0.000000}%
\pgfsetdash{}{0pt}%
\pgfpathmoveto{\pgfqpoint{2.044679in}{1.394919in}}%
\pgfpathlineto{\pgfqpoint{2.046233in}{1.372218in}}%
\pgfpathlineto{\pgfqpoint{2.018995in}{1.370976in}}%
\pgfpathlineto{\pgfqpoint{1.984280in}{1.369339in}}%
\pgfpathlineto{\pgfqpoint{1.980017in}{1.377888in}}%
\pgfpathlineto{\pgfqpoint{1.969155in}{1.380236in}}%
\pgfpathlineto{\pgfqpoint{1.963562in}{1.383583in}}%
\pgfpathlineto{\pgfqpoint{1.960884in}{1.393037in}}%
\pgfpathlineto{\pgfqpoint{1.958754in}{1.394373in}}%
\pgfpathlineto{\pgfqpoint{1.956575in}{1.409480in}}%
\pgfpathlineto{\pgfqpoint{1.961879in}{1.417685in}}%
\pgfpathlineto{\pgfqpoint{1.954703in}{1.423577in}}%
\pgfpathlineto{\pgfqpoint{1.954322in}{1.428443in}}%
\pgfpathlineto{\pgfqpoint{1.980060in}{1.429597in}}%
\pgfpathlineto{\pgfqpoint{1.980936in}{1.413027in}}%
\pgfpathlineto{\pgfqpoint{1.986647in}{1.415927in}}%
\pgfpathlineto{\pgfqpoint{1.992238in}{1.415509in}}%
\pgfpathlineto{\pgfqpoint{1.993663in}{1.389926in}}%
\pgfpathlineto{\pgfqpoint{2.022119in}{1.391513in}}%
\pgfpathlineto{\pgfqpoint{2.021967in}{1.394387in}}%
\pgfpathlineto{\pgfqpoint{2.044679in}{1.394919in}}%
\pgfpathclose%
\pgfusepath{fill}%
\end{pgfscope}%
\begin{pgfscope}%
\pgfpathrectangle{\pgfqpoint{0.100000in}{0.100000in}}{\pgfqpoint{3.007045in}{1.925000in}}%
\pgfusepath{clip}%
\pgfsetbuttcap%
\pgfsetmiterjoin%
\definecolor{currentfill}{rgb}{0.451765,0.767090,0.461207}%
\pgfsetfillcolor{currentfill}%
\pgfsetlinewidth{0.000000pt}%
\definecolor{currentstroke}{rgb}{0.000000,0.000000,0.000000}%
\pgfsetstrokecolor{currentstroke}%
\pgfsetstrokeopacity{0.000000}%
\pgfsetdash{}{0pt}%
\pgfpathmoveto{\pgfqpoint{1.923472in}{0.954935in}}%
\pgfpathlineto{\pgfqpoint{1.923336in}{0.970590in}}%
\pgfpathlineto{\pgfqpoint{1.924925in}{0.970656in}}%
\pgfpathlineto{\pgfqpoint{1.924728in}{0.990132in}}%
\pgfpathlineto{\pgfqpoint{1.925420in}{1.007434in}}%
\pgfpathlineto{\pgfqpoint{1.916864in}{1.007392in}}%
\pgfpathlineto{\pgfqpoint{1.915878in}{1.043286in}}%
\pgfpathlineto{\pgfqpoint{1.941697in}{1.043763in}}%
\pgfpathlineto{\pgfqpoint{1.942096in}{1.032256in}}%
\pgfpathlineto{\pgfqpoint{1.972430in}{1.032823in}}%
\pgfpathlineto{\pgfqpoint{1.972674in}{1.026092in}}%
\pgfpathlineto{\pgfqpoint{1.975616in}{1.018755in}}%
\pgfpathlineto{\pgfqpoint{1.981501in}{1.015319in}}%
\pgfpathlineto{\pgfqpoint{1.981594in}{1.012166in}}%
\pgfpathlineto{\pgfqpoint{1.971148in}{1.011094in}}%
\pgfpathlineto{\pgfqpoint{1.971458in}{0.993921in}}%
\pgfpathlineto{\pgfqpoint{1.977205in}{0.994031in}}%
\pgfpathlineto{\pgfqpoint{1.977561in}{0.972537in}}%
\pgfpathlineto{\pgfqpoint{1.962703in}{0.971887in}}%
\pgfpathlineto{\pgfqpoint{1.960431in}{0.971833in}}%
\pgfpathlineto{\pgfqpoint{1.960613in}{0.961138in}}%
\pgfpathlineto{\pgfqpoint{1.954694in}{0.961038in}}%
\pgfpathlineto{\pgfqpoint{1.954901in}{0.955275in}}%
\pgfpathlineto{\pgfqpoint{1.923472in}{0.954935in}}%
\pgfpathclose%
\pgfusepath{fill}%
\end{pgfscope}%
\begin{pgfscope}%
\pgfpathrectangle{\pgfqpoint{0.100000in}{0.100000in}}{\pgfqpoint{3.007045in}{1.925000in}}%
\pgfusepath{clip}%
\pgfsetbuttcap%
\pgfsetmiterjoin%
\definecolor{currentfill}{rgb}{0.657670,0.862053,0.633449}%
\pgfsetfillcolor{currentfill}%
\pgfsetlinewidth{0.000000pt}%
\definecolor{currentstroke}{rgb}{0.000000,0.000000,0.000000}%
\pgfsetstrokecolor{currentstroke}%
\pgfsetstrokeopacity{0.000000}%
\pgfsetdash{}{0pt}%
\pgfpathmoveto{\pgfqpoint{1.018543in}{1.560865in}}%
\pgfpathlineto{\pgfqpoint{1.020483in}{1.573439in}}%
\pgfpathlineto{\pgfqpoint{1.038445in}{1.570273in}}%
\pgfpathlineto{\pgfqpoint{1.037496in}{1.564829in}}%
\pgfpathlineto{\pgfqpoint{1.051087in}{1.562499in}}%
\pgfpathlineto{\pgfqpoint{1.054900in}{1.557314in}}%
\pgfpathlineto{\pgfqpoint{1.053058in}{1.551890in}}%
\pgfpathlineto{\pgfqpoint{1.057765in}{1.542345in}}%
\pgfpathlineto{\pgfqpoint{1.058870in}{1.534775in}}%
\pgfpathlineto{\pgfqpoint{1.061138in}{1.530583in}}%
\pgfpathlineto{\pgfqpoint{1.059788in}{1.522455in}}%
\pgfpathlineto{\pgfqpoint{1.053944in}{1.487032in}}%
\pgfpathlineto{\pgfqpoint{1.052724in}{1.481408in}}%
\pgfpathlineto{\pgfqpoint{1.039144in}{1.483711in}}%
\pgfpathlineto{\pgfqpoint{1.038193in}{1.478061in}}%
\pgfpathlineto{\pgfqpoint{1.027092in}{1.480000in}}%
\pgfpathlineto{\pgfqpoint{1.026443in}{1.476331in}}%
\pgfpathlineto{\pgfqpoint{1.015288in}{1.478274in}}%
\pgfpathlineto{\pgfqpoint{1.016022in}{1.482462in}}%
\pgfpathlineto{\pgfqpoint{1.005297in}{1.485446in}}%
\pgfpathlineto{\pgfqpoint{1.007407in}{1.497511in}}%
\pgfpathlineto{\pgfqpoint{1.000614in}{1.501315in}}%
\pgfpathlineto{\pgfqpoint{0.999466in}{1.507030in}}%
\pgfpathlineto{\pgfqpoint{0.992126in}{1.508384in}}%
\pgfpathlineto{\pgfqpoint{0.995747in}{1.527955in}}%
\pgfpathlineto{\pgfqpoint{1.005109in}{1.526977in}}%
\pgfpathlineto{\pgfqpoint{1.012932in}{1.528900in}}%
\pgfpathlineto{\pgfqpoint{1.018543in}{1.560865in}}%
\pgfpathclose%
\pgfusepath{fill}%
\end{pgfscope}%
\begin{pgfscope}%
\pgfpathrectangle{\pgfqpoint{0.100000in}{0.100000in}}{\pgfqpoint{3.007045in}{1.925000in}}%
\pgfusepath{clip}%
\pgfsetbuttcap%
\pgfsetmiterjoin%
\definecolor{currentfill}{rgb}{0.723122,0.889612,0.697178}%
\pgfsetfillcolor{currentfill}%
\pgfsetlinewidth{0.000000pt}%
\definecolor{currentstroke}{rgb}{0.000000,0.000000,0.000000}%
\pgfsetstrokecolor{currentstroke}%
\pgfsetstrokeopacity{0.000000}%
\pgfsetdash{}{0pt}%
\pgfpathmoveto{\pgfqpoint{1.578119in}{0.708542in}}%
\pgfpathlineto{\pgfqpoint{1.577602in}{0.694365in}}%
\pgfpathlineto{\pgfqpoint{1.551784in}{0.680740in}}%
\pgfpathlineto{\pgfqpoint{1.507461in}{0.683119in}}%
\pgfpathlineto{\pgfqpoint{1.482882in}{0.684491in}}%
\pgfpathlineto{\pgfqpoint{1.486605in}{0.742533in}}%
\pgfpathlineto{\pgfqpoint{1.523941in}{0.740360in}}%
\pgfpathlineto{\pgfqpoint{1.573659in}{0.737980in}}%
\pgfpathlineto{\pgfqpoint{1.572479in}{0.708902in}}%
\pgfpathlineto{\pgfqpoint{1.578119in}{0.708542in}}%
\pgfpathclose%
\pgfusepath{fill}%
\end{pgfscope}%
\begin{pgfscope}%
\pgfpathrectangle{\pgfqpoint{0.100000in}{0.100000in}}{\pgfqpoint{3.007045in}{1.925000in}}%
\pgfusepath{clip}%
\pgfsetbuttcap%
\pgfsetmiterjoin%
\definecolor{currentfill}{rgb}{0.741822,0.897486,0.715386}%
\pgfsetfillcolor{currentfill}%
\pgfsetlinewidth{0.000000pt}%
\definecolor{currentstroke}{rgb}{0.000000,0.000000,0.000000}%
\pgfsetstrokecolor{currentstroke}%
\pgfsetstrokeopacity{0.000000}%
\pgfsetdash{}{0pt}%
\pgfpathmoveto{\pgfqpoint{0.936407in}{1.175969in}}%
\pgfpathlineto{\pgfqpoint{0.936286in}{1.175309in}}%
\pgfpathlineto{\pgfqpoint{0.968213in}{1.169632in}}%
\pgfpathlineto{\pgfqpoint{1.001312in}{1.163987in}}%
\pgfpathlineto{\pgfqpoint{1.002930in}{1.149720in}}%
\pgfpathlineto{\pgfqpoint{1.005017in}{1.143547in}}%
\pgfpathlineto{\pgfqpoint{1.002522in}{1.140541in}}%
\pgfpathlineto{\pgfqpoint{0.963729in}{1.147104in}}%
\pgfpathlineto{\pgfqpoint{0.904834in}{1.157752in}}%
\pgfpathlineto{\pgfqpoint{0.907847in}{1.175524in}}%
\pgfpathlineto{\pgfqpoint{0.913573in}{1.180156in}}%
\pgfpathlineto{\pgfqpoint{0.936407in}{1.175969in}}%
\pgfpathclose%
\pgfusepath{fill}%
\end{pgfscope}%
\begin{pgfscope}%
\pgfpathrectangle{\pgfqpoint{0.100000in}{0.100000in}}{\pgfqpoint{3.007045in}{1.925000in}}%
\pgfusepath{clip}%
\pgfsetbuttcap%
\pgfsetmiterjoin%
\definecolor{currentfill}{rgb}{0.562860,0.819008,0.551511}%
\pgfsetfillcolor{currentfill}%
\pgfsetlinewidth{0.000000pt}%
\definecolor{currentstroke}{rgb}{0.000000,0.000000,0.000000}%
\pgfsetstrokecolor{currentstroke}%
\pgfsetstrokeopacity{0.000000}%
\pgfsetdash{}{0pt}%
\pgfpathmoveto{\pgfqpoint{1.786299in}{0.783647in}}%
\pgfpathlineto{\pgfqpoint{1.800644in}{0.779398in}}%
\pgfpathlineto{\pgfqpoint{1.801001in}{0.799445in}}%
\pgfpathlineto{\pgfqpoint{1.806094in}{0.799451in}}%
\pgfpathlineto{\pgfqpoint{1.814189in}{0.790353in}}%
\pgfpathlineto{\pgfqpoint{1.825054in}{0.789133in}}%
\pgfpathlineto{\pgfqpoint{1.829586in}{0.787129in}}%
\pgfpathlineto{\pgfqpoint{1.829258in}{0.781767in}}%
\pgfpathlineto{\pgfqpoint{1.836570in}{0.777705in}}%
\pgfpathlineto{\pgfqpoint{1.842812in}{0.769497in}}%
\pgfpathlineto{\pgfqpoint{1.842406in}{0.765728in}}%
\pgfpathlineto{\pgfqpoint{1.846121in}{0.759629in}}%
\pgfpathlineto{\pgfqpoint{1.841164in}{0.748497in}}%
\pgfpathlineto{\pgfqpoint{1.836073in}{0.747122in}}%
\pgfpathlineto{\pgfqpoint{1.835563in}{0.743537in}}%
\pgfpathlineto{\pgfqpoint{1.838686in}{0.738812in}}%
\pgfpathlineto{\pgfqpoint{1.825469in}{0.738647in}}%
\pgfpathlineto{\pgfqpoint{1.825570in}{0.729490in}}%
\pgfpathlineto{\pgfqpoint{1.788689in}{0.729118in}}%
\pgfpathlineto{\pgfqpoint{1.787896in}{0.730793in}}%
\pgfpathlineto{\pgfqpoint{1.763868in}{0.730641in}}%
\pgfpathlineto{\pgfqpoint{1.762705in}{0.734578in}}%
\pgfpathlineto{\pgfqpoint{1.755158in}{0.734645in}}%
\pgfpathlineto{\pgfqpoint{1.755274in}{0.762070in}}%
\pgfpathlineto{\pgfqpoint{1.767169in}{0.763420in}}%
\pgfpathlineto{\pgfqpoint{1.779169in}{0.763215in}}%
\pgfpathlineto{\pgfqpoint{1.786323in}{0.758894in}}%
\pgfpathlineto{\pgfqpoint{1.786299in}{0.783647in}}%
\pgfpathclose%
\pgfusepath{fill}%
\end{pgfscope}%
\begin{pgfscope}%
\pgfpathrectangle{\pgfqpoint{0.100000in}{0.100000in}}{\pgfqpoint{3.007045in}{1.925000in}}%
\pgfusepath{clip}%
\pgfsetbuttcap%
\pgfsetmiterjoin%
\definecolor{currentfill}{rgb}{0.013456,0.438985,0.182161}%
\pgfsetfillcolor{currentfill}%
\pgfsetlinewidth{0.000000pt}%
\definecolor{currentstroke}{rgb}{0.000000,0.000000,0.000000}%
\pgfsetstrokecolor{currentstroke}%
\pgfsetstrokeopacity{0.000000}%
\pgfsetdash{}{0pt}%
\pgfpathmoveto{\pgfqpoint{2.517736in}{0.501110in}}%
\pgfpathlineto{\pgfqpoint{2.508454in}{0.501648in}}%
\pgfpathlineto{\pgfqpoint{2.495655in}{0.499863in}}%
\pgfpathlineto{\pgfqpoint{2.494027in}{0.512423in}}%
\pgfpathlineto{\pgfqpoint{2.486048in}{0.518807in}}%
\pgfpathlineto{\pgfqpoint{2.488304in}{0.521528in}}%
\pgfpathlineto{\pgfqpoint{2.494299in}{0.520921in}}%
\pgfpathlineto{\pgfqpoint{2.498437in}{0.523892in}}%
\pgfpathlineto{\pgfqpoint{2.496864in}{0.535136in}}%
\pgfpathlineto{\pgfqpoint{2.504526in}{0.536259in}}%
\pgfpathlineto{\pgfqpoint{2.501841in}{0.554053in}}%
\pgfpathlineto{\pgfqpoint{2.513059in}{0.555566in}}%
\pgfpathlineto{\pgfqpoint{2.513609in}{0.551975in}}%
\pgfpathlineto{\pgfqpoint{2.522149in}{0.556021in}}%
\pgfpathlineto{\pgfqpoint{2.533920in}{0.561053in}}%
\pgfpathlineto{\pgfqpoint{2.538085in}{0.559398in}}%
\pgfpathlineto{\pgfqpoint{2.541214in}{0.552000in}}%
\pgfpathlineto{\pgfqpoint{2.547864in}{0.546632in}}%
\pgfpathlineto{\pgfqpoint{2.550293in}{0.531506in}}%
\pgfpathlineto{\pgfqpoint{2.549990in}{0.525560in}}%
\pgfpathlineto{\pgfqpoint{2.512188in}{0.520102in}}%
\pgfpathlineto{\pgfqpoint{2.522009in}{0.510090in}}%
\pgfpathlineto{\pgfqpoint{2.517736in}{0.501110in}}%
\pgfpathclose%
\pgfusepath{fill}%
\end{pgfscope}%
\begin{pgfscope}%
\pgfpathrectangle{\pgfqpoint{0.100000in}{0.100000in}}{\pgfqpoint{3.007045in}{1.925000in}}%
\pgfusepath{clip}%
\pgfsetbuttcap%
\pgfsetmiterjoin%
\definecolor{currentfill}{rgb}{0.612687,0.842261,0.592480}%
\pgfsetfillcolor{currentfill}%
\pgfsetlinewidth{0.000000pt}%
\definecolor{currentstroke}{rgb}{0.000000,0.000000,0.000000}%
\pgfsetstrokecolor{currentstroke}%
\pgfsetstrokeopacity{0.000000}%
\pgfsetdash{}{0pt}%
\pgfpathmoveto{\pgfqpoint{1.685427in}{0.912299in}}%
\pgfpathlineto{\pgfqpoint{1.708380in}{0.911948in}}%
\pgfpathlineto{\pgfqpoint{1.708573in}{0.926325in}}%
\pgfpathlineto{\pgfqpoint{1.731364in}{0.926038in}}%
\pgfpathlineto{\pgfqpoint{1.731290in}{0.917416in}}%
\pgfpathlineto{\pgfqpoint{1.734135in}{0.917404in}}%
\pgfpathlineto{\pgfqpoint{1.734052in}{0.905918in}}%
\pgfpathlineto{\pgfqpoint{1.728181in}{0.905962in}}%
\pgfpathlineto{\pgfqpoint{1.728070in}{0.900172in}}%
\pgfpathlineto{\pgfqpoint{1.725208in}{0.894438in}}%
\pgfpathlineto{\pgfqpoint{1.719521in}{0.894505in}}%
\pgfpathlineto{\pgfqpoint{1.719456in}{0.888768in}}%
\pgfpathlineto{\pgfqpoint{1.694674in}{0.889113in}}%
\pgfpathlineto{\pgfqpoint{1.691799in}{0.898304in}}%
\pgfpathlineto{\pgfqpoint{1.684983in}{0.896595in}}%
\pgfpathlineto{\pgfqpoint{1.685427in}{0.912299in}}%
\pgfpathclose%
\pgfusepath{fill}%
\end{pgfscope}%
\begin{pgfscope}%
\pgfpathrectangle{\pgfqpoint{0.100000in}{0.100000in}}{\pgfqpoint{3.007045in}{1.925000in}}%
\pgfusepath{clip}%
\pgfsetbuttcap%
\pgfsetmiterjoin%
\definecolor{currentfill}{rgb}{0.420392,0.751711,0.445829}%
\pgfsetfillcolor{currentfill}%
\pgfsetlinewidth{0.000000pt}%
\definecolor{currentstroke}{rgb}{0.000000,0.000000,0.000000}%
\pgfsetstrokecolor{currentstroke}%
\pgfsetstrokeopacity{0.000000}%
\pgfsetdash{}{0pt}%
\pgfpathmoveto{\pgfqpoint{2.474213in}{0.983871in}}%
\pgfpathlineto{\pgfqpoint{2.468967in}{0.985644in}}%
\pgfpathlineto{\pgfqpoint{2.465463in}{0.979232in}}%
\pgfpathlineto{\pgfqpoint{2.462402in}{0.980829in}}%
\pgfpathlineto{\pgfqpoint{2.458453in}{0.990986in}}%
\pgfpathlineto{\pgfqpoint{2.460956in}{0.992711in}}%
\pgfpathlineto{\pgfqpoint{2.463756in}{1.002325in}}%
\pgfpathlineto{\pgfqpoint{2.458161in}{1.009811in}}%
\pgfpathlineto{\pgfqpoint{2.461449in}{1.016280in}}%
\pgfpathlineto{\pgfqpoint{2.462097in}{1.022483in}}%
\pgfpathlineto{\pgfqpoint{2.466964in}{1.026275in}}%
\pgfpathlineto{\pgfqpoint{2.476402in}{1.027509in}}%
\pgfpathlineto{\pgfqpoint{2.475027in}{1.025720in}}%
\pgfpathlineto{\pgfqpoint{2.475252in}{1.017742in}}%
\pgfpathlineto{\pgfqpoint{2.473634in}{1.012683in}}%
\pgfpathlineto{\pgfqpoint{2.479391in}{1.010461in}}%
\pgfpathlineto{\pgfqpoint{2.490105in}{1.002416in}}%
\pgfpathlineto{\pgfqpoint{2.486521in}{0.996045in}}%
\pgfpathlineto{\pgfqpoint{2.479657in}{0.995365in}}%
\pgfpathlineto{\pgfqpoint{2.474213in}{0.983871in}}%
\pgfpathclose%
\pgfusepath{fill}%
\end{pgfscope}%
\begin{pgfscope}%
\pgfpathrectangle{\pgfqpoint{0.100000in}{0.100000in}}{\pgfqpoint{3.007045in}{1.925000in}}%
\pgfusepath{clip}%
\pgfsetbuttcap%
\pgfsetmiterjoin%
\definecolor{currentfill}{rgb}{0.474279,0.777670,0.478677}%
\pgfsetfillcolor{currentfill}%
\pgfsetlinewidth{0.000000pt}%
\definecolor{currentstroke}{rgb}{0.000000,0.000000,0.000000}%
\pgfsetstrokecolor{currentstroke}%
\pgfsetstrokeopacity{0.000000}%
\pgfsetdash{}{0pt}%
\pgfpathmoveto{\pgfqpoint{1.527075in}{1.552396in}}%
\pgfpathlineto{\pgfqpoint{1.552729in}{1.550914in}}%
\pgfpathlineto{\pgfqpoint{1.559045in}{1.550570in}}%
\pgfpathlineto{\pgfqpoint{1.557925in}{1.527548in}}%
\pgfpathlineto{\pgfqpoint{1.553171in}{1.527805in}}%
\pgfpathlineto{\pgfqpoint{1.518970in}{1.529708in}}%
\pgfpathlineto{\pgfqpoint{1.518468in}{1.535998in}}%
\pgfpathlineto{\pgfqpoint{1.524818in}{1.537110in}}%
\pgfpathlineto{\pgfqpoint{1.523533in}{1.542968in}}%
\pgfpathlineto{\pgfqpoint{1.527075in}{1.552396in}}%
\pgfpathclose%
\pgfusepath{fill}%
\end{pgfscope}%
\begin{pgfscope}%
\pgfpathrectangle{\pgfqpoint{0.100000in}{0.100000in}}{\pgfqpoint{3.007045in}{1.925000in}}%
\pgfusepath{clip}%
\pgfsetbuttcap%
\pgfsetmiterjoin%
\definecolor{currentfill}{rgb}{0.816378,0.928120,0.791326}%
\pgfsetfillcolor{currentfill}%
\pgfsetlinewidth{0.000000pt}%
\definecolor{currentstroke}{rgb}{0.000000,0.000000,0.000000}%
\pgfsetstrokecolor{currentstroke}%
\pgfsetstrokeopacity{0.000000}%
\pgfsetdash{}{0pt}%
\pgfpathmoveto{\pgfqpoint{1.265329in}{0.695996in}}%
\pgfpathlineto{\pgfqpoint{1.282225in}{0.694354in}}%
\pgfpathlineto{\pgfqpoint{1.319120in}{0.690958in}}%
\pgfpathlineto{\pgfqpoint{1.319633in}{0.696676in}}%
\pgfpathlineto{\pgfqpoint{1.334479in}{0.695294in}}%
\pgfpathlineto{\pgfqpoint{1.332017in}{0.666591in}}%
\pgfpathlineto{\pgfqpoint{1.333802in}{0.666433in}}%
\pgfpathlineto{\pgfqpoint{1.331766in}{0.642740in}}%
\pgfpathlineto{\pgfqpoint{1.328596in}{0.641304in}}%
\pgfpathlineto{\pgfqpoint{1.322677in}{0.647278in}}%
\pgfpathlineto{\pgfqpoint{1.318428in}{0.649073in}}%
\pgfpathlineto{\pgfqpoint{1.282209in}{0.612003in}}%
\pgfpathlineto{\pgfqpoint{1.255027in}{0.637292in}}%
\pgfpathlineto{\pgfqpoint{1.265329in}{0.695996in}}%
\pgfpathclose%
\pgfusepath{fill}%
\end{pgfscope}%
\begin{pgfscope}%
\pgfpathrectangle{\pgfqpoint{0.100000in}{0.100000in}}{\pgfqpoint{3.007045in}{1.925000in}}%
\pgfusepath{clip}%
\pgfsetbuttcap%
\pgfsetmiterjoin%
\definecolor{currentfill}{rgb}{0.751173,0.901423,0.724491}%
\pgfsetfillcolor{currentfill}%
\pgfsetlinewidth{0.000000pt}%
\definecolor{currentstroke}{rgb}{0.000000,0.000000,0.000000}%
\pgfsetstrokecolor{currentstroke}%
\pgfsetstrokeopacity{0.000000}%
\pgfsetdash{}{0pt}%
\pgfpathmoveto{\pgfqpoint{2.742724in}{1.035580in}}%
\pgfpathlineto{\pgfqpoint{2.738199in}{1.043269in}}%
\pgfpathlineto{\pgfqpoint{2.738529in}{1.052500in}}%
\pgfpathlineto{\pgfqpoint{2.733124in}{1.055886in}}%
\pgfpathlineto{\pgfqpoint{2.724911in}{1.056794in}}%
\pgfpathlineto{\pgfqpoint{2.723871in}{1.066185in}}%
\pgfpathlineto{\pgfqpoint{2.723519in}{1.072852in}}%
\pgfpathlineto{\pgfqpoint{2.738779in}{1.094235in}}%
\pgfpathlineto{\pgfqpoint{2.741244in}{1.095623in}}%
\pgfpathlineto{\pgfqpoint{2.746605in}{1.094479in}}%
\pgfpathlineto{\pgfqpoint{2.750900in}{1.100328in}}%
\pgfpathlineto{\pgfqpoint{2.761478in}{1.098464in}}%
\pgfpathlineto{\pgfqpoint{2.766403in}{1.100448in}}%
\pgfpathlineto{\pgfqpoint{2.771017in}{1.090586in}}%
\pgfpathlineto{\pgfqpoint{2.777353in}{1.079986in}}%
\pgfpathlineto{\pgfqpoint{2.787444in}{1.058403in}}%
\pgfpathlineto{\pgfqpoint{2.786756in}{1.058061in}}%
\pgfpathlineto{\pgfqpoint{2.780444in}{1.069950in}}%
\pgfpathlineto{\pgfqpoint{2.777839in}{1.065735in}}%
\pgfpathlineto{\pgfqpoint{2.788743in}{1.048436in}}%
\pgfpathlineto{\pgfqpoint{2.782567in}{1.052153in}}%
\pgfpathlineto{\pgfqpoint{2.776592in}{1.052297in}}%
\pgfpathlineto{\pgfqpoint{2.766420in}{1.044865in}}%
\pgfpathlineto{\pgfqpoint{2.757806in}{1.041934in}}%
\pgfpathlineto{\pgfqpoint{2.752708in}{1.036350in}}%
\pgfpathlineto{\pgfqpoint{2.742724in}{1.035580in}}%
\pgfpathclose%
\pgfusepath{fill}%
\end{pgfscope}%
\begin{pgfscope}%
\pgfpathrectangle{\pgfqpoint{0.100000in}{0.100000in}}{\pgfqpoint{3.007045in}{1.925000in}}%
\pgfusepath{clip}%
\pgfsetbuttcap%
\pgfsetmiterjoin%
\definecolor{currentfill}{rgb}{0.751173,0.901423,0.724491}%
\pgfsetfillcolor{currentfill}%
\pgfsetlinewidth{0.000000pt}%
\definecolor{currentstroke}{rgb}{0.000000,0.000000,0.000000}%
\pgfsetstrokecolor{currentstroke}%
\pgfsetstrokeopacity{0.000000}%
\pgfsetdash{}{0pt}%
\pgfpathmoveto{\pgfqpoint{1.685760in}{0.932254in}}%
\pgfpathlineto{\pgfqpoint{1.658027in}{0.932849in}}%
\pgfpathlineto{\pgfqpoint{1.646710in}{0.935974in}}%
\pgfpathlineto{\pgfqpoint{1.647016in}{0.947800in}}%
\pgfpathlineto{\pgfqpoint{1.641329in}{0.947962in}}%
\pgfpathlineto{\pgfqpoint{1.642047in}{0.976241in}}%
\pgfpathlineto{\pgfqpoint{1.670693in}{0.976025in}}%
\pgfpathlineto{\pgfqpoint{1.671642in}{0.972779in}}%
\pgfpathlineto{\pgfqpoint{1.664644in}{0.969981in}}%
\pgfpathlineto{\pgfqpoint{1.664393in}{0.958534in}}%
\pgfpathlineto{\pgfqpoint{1.670084in}{0.958402in}}%
\pgfpathlineto{\pgfqpoint{1.669953in}{0.952684in}}%
\pgfpathlineto{\pgfqpoint{1.675597in}{0.952567in}}%
\pgfpathlineto{\pgfqpoint{1.675476in}{0.946791in}}%
\pgfpathlineto{\pgfqpoint{1.685998in}{0.946658in}}%
\pgfpathlineto{\pgfqpoint{1.685760in}{0.932254in}}%
\pgfpathclose%
\pgfusepath{fill}%
\end{pgfscope}%
\begin{pgfscope}%
\pgfpathrectangle{\pgfqpoint{0.100000in}{0.100000in}}{\pgfqpoint{3.007045in}{1.925000in}}%
\pgfusepath{clip}%
\pgfsetbuttcap%
\pgfsetmiterjoin%
\definecolor{currentfill}{rgb}{0.596078,0.834510,0.578824}%
\pgfsetfillcolor{currentfill}%
\pgfsetlinewidth{0.000000pt}%
\definecolor{currentstroke}{rgb}{0.000000,0.000000,0.000000}%
\pgfsetstrokecolor{currentstroke}%
\pgfsetstrokeopacity{0.000000}%
\pgfsetdash{}{0pt}%
\pgfpathmoveto{\pgfqpoint{1.564548in}{1.295579in}}%
\pgfpathlineto{\pgfqpoint{1.563271in}{1.272731in}}%
\pgfpathlineto{\pgfqpoint{1.552431in}{1.273213in}}%
\pgfpathlineto{\pgfqpoint{1.512812in}{1.275327in}}%
\pgfpathlineto{\pgfqpoint{1.514112in}{1.321020in}}%
\pgfpathlineto{\pgfqpoint{1.485080in}{1.322839in}}%
\pgfpathlineto{\pgfqpoint{1.486374in}{1.345738in}}%
\pgfpathlineto{\pgfqpoint{1.519516in}{1.343574in}}%
\pgfpathlineto{\pgfqpoint{1.543010in}{1.342313in}}%
\pgfpathlineto{\pgfqpoint{1.541840in}{1.319542in}}%
\pgfpathlineto{\pgfqpoint{1.565174in}{1.318418in}}%
\pgfpathlineto{\pgfqpoint{1.564548in}{1.295579in}}%
\pgfpathclose%
\pgfusepath{fill}%
\end{pgfscope}%
\begin{pgfscope}%
\pgfpathrectangle{\pgfqpoint{0.100000in}{0.100000in}}{\pgfqpoint{3.007045in}{1.925000in}}%
\pgfusepath{clip}%
\pgfsetbuttcap%
\pgfsetmiterjoin%
\definecolor{currentfill}{rgb}{0.709097,0.883706,0.683522}%
\pgfsetfillcolor{currentfill}%
\pgfsetlinewidth{0.000000pt}%
\definecolor{currentstroke}{rgb}{0.000000,0.000000,0.000000}%
\pgfsetstrokecolor{currentstroke}%
\pgfsetstrokeopacity{0.000000}%
\pgfsetdash{}{0pt}%
\pgfpathmoveto{\pgfqpoint{2.025623in}{0.904867in}}%
\pgfpathlineto{\pgfqpoint{2.008210in}{0.904668in}}%
\pgfpathlineto{\pgfqpoint{1.985207in}{0.903228in}}%
\pgfpathlineto{\pgfqpoint{1.984303in}{0.932387in}}%
\pgfpathlineto{\pgfqpoint{1.975423in}{0.932334in}}%
\pgfpathlineto{\pgfqpoint{1.966949in}{0.932026in}}%
\pgfpathlineto{\pgfqpoint{1.966442in}{0.954472in}}%
\pgfpathlineto{\pgfqpoint{1.971212in}{0.955492in}}%
\pgfpathlineto{\pgfqpoint{1.969967in}{0.964093in}}%
\pgfpathlineto{\pgfqpoint{1.962703in}{0.971887in}}%
\pgfpathlineto{\pgfqpoint{1.977561in}{0.972537in}}%
\pgfpathlineto{\pgfqpoint{2.025516in}{0.974832in}}%
\pgfpathlineto{\pgfqpoint{2.034251in}{0.967797in}}%
\pgfpathlineto{\pgfqpoint{2.033532in}{0.960342in}}%
\pgfpathlineto{\pgfqpoint{2.026615in}{0.954215in}}%
\pgfpathlineto{\pgfqpoint{2.021688in}{0.947698in}}%
\pgfpathlineto{\pgfqpoint{2.023624in}{0.941619in}}%
\pgfpathlineto{\pgfqpoint{2.025623in}{0.904867in}}%
\pgfpathclose%
\pgfusepath{fill}%
\end{pgfscope}%
\begin{pgfscope}%
\pgfpathrectangle{\pgfqpoint{0.100000in}{0.100000in}}{\pgfqpoint{3.007045in}{1.925000in}}%
\pgfusepath{clip}%
\pgfsetbuttcap%
\pgfsetmiterjoin%
\definecolor{currentfill}{rgb}{0.451765,0.767090,0.461207}%
\pgfsetfillcolor{currentfill}%
\pgfsetlinewidth{0.000000pt}%
\definecolor{currentstroke}{rgb}{0.000000,0.000000,0.000000}%
\pgfsetstrokecolor{currentstroke}%
\pgfsetstrokeopacity{0.000000}%
\pgfsetdash{}{0pt}%
\pgfpathmoveto{\pgfqpoint{1.169817in}{1.038957in}}%
\pgfpathlineto{\pgfqpoint{1.161246in}{1.055657in}}%
\pgfpathlineto{\pgfqpoint{1.162663in}{1.066636in}}%
\pgfpathlineto{\pgfqpoint{1.139253in}{1.069444in}}%
\pgfpathlineto{\pgfqpoint{1.139523in}{1.071438in}}%
\pgfpathlineto{\pgfqpoint{1.140887in}{1.088045in}}%
\pgfpathlineto{\pgfqpoint{1.143557in}{1.105192in}}%
\pgfpathlineto{\pgfqpoint{1.150792in}{1.105369in}}%
\pgfpathlineto{\pgfqpoint{1.154865in}{1.136033in}}%
\pgfpathlineto{\pgfqpoint{1.193414in}{1.130886in}}%
\pgfpathlineto{\pgfqpoint{1.202180in}{1.129872in}}%
\pgfpathlineto{\pgfqpoint{1.205700in}{1.130953in}}%
\pgfpathlineto{\pgfqpoint{1.220511in}{1.109227in}}%
\pgfpathlineto{\pgfqpoint{1.223896in}{1.096979in}}%
\pgfpathlineto{\pgfqpoint{1.229068in}{1.091653in}}%
\pgfpathlineto{\pgfqpoint{1.231022in}{1.088493in}}%
\pgfpathlineto{\pgfqpoint{1.225618in}{1.073383in}}%
\pgfpathlineto{\pgfqpoint{1.236457in}{1.074727in}}%
\pgfpathlineto{\pgfqpoint{1.241493in}{1.067690in}}%
\pgfpathlineto{\pgfqpoint{1.242720in}{1.057254in}}%
\pgfpathlineto{\pgfqpoint{1.241118in}{1.050301in}}%
\pgfpathlineto{\pgfqpoint{1.238883in}{1.030797in}}%
\pgfpathlineto{\pgfqpoint{1.235453in}{1.031177in}}%
\pgfpathlineto{\pgfqpoint{1.169817in}{1.038957in}}%
\pgfpathclose%
\pgfusepath{fill}%
\end{pgfscope}%
\begin{pgfscope}%
\pgfpathrectangle{\pgfqpoint{0.100000in}{0.100000in}}{\pgfqpoint{3.007045in}{1.925000in}}%
\pgfusepath{clip}%
\pgfsetbuttcap%
\pgfsetmiterjoin%
\definecolor{currentfill}{rgb}{0.345098,0.714802,0.408920}%
\pgfsetfillcolor{currentfill}%
\pgfsetlinewidth{0.000000pt}%
\definecolor{currentstroke}{rgb}{0.000000,0.000000,0.000000}%
\pgfsetstrokecolor{currentstroke}%
\pgfsetstrokeopacity{0.000000}%
\pgfsetdash{}{0pt}%
\pgfpathmoveto{\pgfqpoint{1.853670in}{1.237862in}}%
\pgfpathlineto{\pgfqpoint{1.844539in}{1.237709in}}%
\pgfpathlineto{\pgfqpoint{1.843908in}{1.275998in}}%
\pgfpathlineto{\pgfqpoint{1.855298in}{1.276147in}}%
\pgfpathlineto{\pgfqpoint{1.866638in}{1.276378in}}%
\pgfpathlineto{\pgfqpoint{1.867080in}{1.259096in}}%
\pgfpathlineto{\pgfqpoint{1.889907in}{1.259660in}}%
\pgfpathlineto{\pgfqpoint{1.890498in}{1.239370in}}%
\pgfpathlineto{\pgfqpoint{1.886676in}{1.239179in}}%
\pgfpathlineto{\pgfqpoint{1.853670in}{1.237862in}}%
\pgfpathclose%
\pgfusepath{fill}%
\end{pgfscope}%
\begin{pgfscope}%
\pgfpathrectangle{\pgfqpoint{0.100000in}{0.100000in}}{\pgfqpoint{3.007045in}{1.925000in}}%
\pgfusepath{clip}%
\pgfsetbuttcap%
\pgfsetmiterjoin%
\definecolor{currentfill}{rgb}{0.875433,0.951742,0.854318}%
\pgfsetfillcolor{currentfill}%
\pgfsetlinewidth{0.000000pt}%
\definecolor{currentstroke}{rgb}{0.000000,0.000000,0.000000}%
\pgfsetstrokecolor{currentstroke}%
\pgfsetstrokeopacity{0.000000}%
\pgfsetdash{}{0pt}%
\pgfpathmoveto{\pgfqpoint{1.638148in}{0.525132in}}%
\pgfpathlineto{\pgfqpoint{1.621659in}{0.515303in}}%
\pgfpathlineto{\pgfqpoint{1.619645in}{0.516644in}}%
\pgfpathlineto{\pgfqpoint{1.615672in}{0.518019in}}%
\pgfpathlineto{\pgfqpoint{1.606065in}{0.530962in}}%
\pgfpathlineto{\pgfqpoint{1.598646in}{0.524138in}}%
\pgfpathlineto{\pgfqpoint{1.597051in}{0.530605in}}%
\pgfpathlineto{\pgfqpoint{1.582069in}{0.543767in}}%
\pgfpathlineto{\pgfqpoint{1.575104in}{0.537324in}}%
\pgfpathlineto{\pgfqpoint{1.565594in}{0.551126in}}%
\pgfpathlineto{\pgfqpoint{1.566327in}{0.575160in}}%
\pgfpathlineto{\pgfqpoint{1.576662in}{0.574910in}}%
\pgfpathlineto{\pgfqpoint{1.592881in}{0.569267in}}%
\pgfpathlineto{\pgfqpoint{1.593576in}{0.574175in}}%
\pgfpathlineto{\pgfqpoint{1.602998in}{0.592878in}}%
\pgfpathlineto{\pgfqpoint{1.610932in}{0.600948in}}%
\pgfpathlineto{\pgfqpoint{1.622469in}{0.597909in}}%
\pgfpathlineto{\pgfqpoint{1.639929in}{0.589583in}}%
\pgfpathlineto{\pgfqpoint{1.643382in}{0.598629in}}%
\pgfpathlineto{\pgfqpoint{1.654308in}{0.604797in}}%
\pgfpathlineto{\pgfqpoint{1.668209in}{0.612458in}}%
\pgfpathlineto{\pgfqpoint{1.670600in}{0.609460in}}%
\pgfpathlineto{\pgfqpoint{1.672091in}{0.600546in}}%
\pgfpathlineto{\pgfqpoint{1.678179in}{0.593593in}}%
\pgfpathlineto{\pgfqpoint{1.679583in}{0.587280in}}%
\pgfpathlineto{\pgfqpoint{1.659758in}{0.576115in}}%
\pgfpathlineto{\pgfqpoint{1.662849in}{0.573743in}}%
\pgfpathlineto{\pgfqpoint{1.665066in}{0.567789in}}%
\pgfpathlineto{\pgfqpoint{1.672046in}{0.560070in}}%
\pgfpathlineto{\pgfqpoint{1.677967in}{0.558409in}}%
\pgfpathlineto{\pgfqpoint{1.672597in}{0.554943in}}%
\pgfpathlineto{\pgfqpoint{1.668934in}{0.549504in}}%
\pgfpathlineto{\pgfqpoint{1.665918in}{0.546597in}}%
\pgfpathlineto{\pgfqpoint{1.655528in}{0.542512in}}%
\pgfpathlineto{\pgfqpoint{1.638148in}{0.525132in}}%
\pgfpathclose%
\pgfusepath{fill}%
\end{pgfscope}%
\begin{pgfscope}%
\pgfpathrectangle{\pgfqpoint{0.100000in}{0.100000in}}{\pgfqpoint{3.007045in}{1.925000in}}%
\pgfusepath{clip}%
\pgfsetbuttcap%
\pgfsetmiterjoin%
\definecolor{currentfill}{rgb}{0.652995,0.860085,0.628897}%
\pgfsetfillcolor{currentfill}%
\pgfsetlinewidth{0.000000pt}%
\definecolor{currentstroke}{rgb}{0.000000,0.000000,0.000000}%
\pgfsetstrokecolor{currentstroke}%
\pgfsetstrokeopacity{0.000000}%
\pgfsetdash{}{0pt}%
\pgfpathmoveto{\pgfqpoint{1.923673in}{0.668219in}}%
\pgfpathlineto{\pgfqpoint{1.922946in}{0.694307in}}%
\pgfpathlineto{\pgfqpoint{1.916979in}{0.699909in}}%
\pgfpathlineto{\pgfqpoint{1.916505in}{0.717286in}}%
\pgfpathlineto{\pgfqpoint{1.910744in}{0.717166in}}%
\pgfpathlineto{\pgfqpoint{1.906378in}{0.722687in}}%
\pgfpathlineto{\pgfqpoint{1.899049in}{0.722750in}}%
\pgfpathlineto{\pgfqpoint{1.898651in}{0.739626in}}%
\pgfpathlineto{\pgfqpoint{1.935072in}{0.740260in}}%
\pgfpathlineto{\pgfqpoint{1.968862in}{0.741315in}}%
\pgfpathlineto{\pgfqpoint{1.970228in}{0.741380in}}%
\pgfpathlineto{\pgfqpoint{1.967255in}{0.733621in}}%
\pgfpathlineto{\pgfqpoint{1.964704in}{0.732579in}}%
\pgfpathlineto{\pgfqpoint{1.959700in}{0.720849in}}%
\pgfpathlineto{\pgfqpoint{1.962387in}{0.712972in}}%
\pgfpathlineto{\pgfqpoint{1.970799in}{0.713279in}}%
\pgfpathlineto{\pgfqpoint{1.970224in}{0.710328in}}%
\pgfpathlineto{\pgfqpoint{1.968808in}{0.707431in}}%
\pgfpathlineto{\pgfqpoint{1.970011in}{0.695703in}}%
\pgfpathlineto{\pgfqpoint{1.966832in}{0.693993in}}%
\pgfpathlineto{\pgfqpoint{1.969121in}{0.688091in}}%
\pgfpathlineto{\pgfqpoint{1.969490in}{0.679240in}}%
\pgfpathlineto{\pgfqpoint{1.965268in}{0.666642in}}%
\pgfpathlineto{\pgfqpoint{1.963857in}{0.672381in}}%
\pgfpathlineto{\pgfqpoint{1.953769in}{0.665525in}}%
\pgfpathlineto{\pgfqpoint{1.949500in}{0.671101in}}%
\pgfpathlineto{\pgfqpoint{1.946656in}{0.668963in}}%
\pgfpathlineto{\pgfqpoint{1.923673in}{0.668219in}}%
\pgfpathclose%
\pgfusepath{fill}%
\end{pgfscope}%
\begin{pgfscope}%
\pgfpathrectangle{\pgfqpoint{0.100000in}{0.100000in}}{\pgfqpoint{3.007045in}{1.925000in}}%
\pgfusepath{clip}%
\pgfsetbuttcap%
\pgfsetmiterjoin%
\definecolor{currentfill}{rgb}{0.313725,0.699423,0.393541}%
\pgfsetfillcolor{currentfill}%
\pgfsetlinewidth{0.000000pt}%
\definecolor{currentstroke}{rgb}{0.000000,0.000000,0.000000}%
\pgfsetstrokecolor{currentstroke}%
\pgfsetstrokeopacity{0.000000}%
\pgfsetdash{}{0pt}%
\pgfpathmoveto{\pgfqpoint{2.192935in}{1.614287in}}%
\pgfpathlineto{\pgfqpoint{2.174903in}{1.611101in}}%
\pgfpathlineto{\pgfqpoint{2.172417in}{1.608214in}}%
\pgfpathlineto{\pgfqpoint{2.171847in}{1.600827in}}%
\pgfpathlineto{\pgfqpoint{2.163982in}{1.596790in}}%
\pgfpathlineto{\pgfqpoint{2.160048in}{1.590165in}}%
\pgfpathlineto{\pgfqpoint{2.155554in}{1.591476in}}%
\pgfpathlineto{\pgfqpoint{2.160765in}{1.598791in}}%
\pgfpathlineto{\pgfqpoint{2.163198in}{1.604703in}}%
\pgfpathlineto{\pgfqpoint{2.152294in}{1.599865in}}%
\pgfpathlineto{\pgfqpoint{2.149656in}{1.593559in}}%
\pgfpathlineto{\pgfqpoint{2.144154in}{1.589747in}}%
\pgfpathlineto{\pgfqpoint{2.139521in}{1.593745in}}%
\pgfpathlineto{\pgfqpoint{2.134646in}{1.588277in}}%
\pgfpathlineto{\pgfqpoint{2.131069in}{1.580709in}}%
\pgfpathlineto{\pgfqpoint{2.128131in}{1.580457in}}%
\pgfpathlineto{\pgfqpoint{2.126263in}{1.603472in}}%
\pgfpathlineto{\pgfqpoint{2.123850in}{1.609038in}}%
\pgfpathlineto{\pgfqpoint{2.112389in}{1.608131in}}%
\pgfpathlineto{\pgfqpoint{2.111036in}{1.625398in}}%
\pgfpathlineto{\pgfqpoint{2.088086in}{1.623582in}}%
\pgfpathlineto{\pgfqpoint{2.087227in}{1.635040in}}%
\pgfpathlineto{\pgfqpoint{2.085981in}{1.652184in}}%
\pgfpathlineto{\pgfqpoint{2.091627in}{1.652681in}}%
\pgfpathlineto{\pgfqpoint{2.088768in}{1.658202in}}%
\pgfpathlineto{\pgfqpoint{2.088047in}{1.667928in}}%
\pgfpathlineto{\pgfqpoint{2.094537in}{1.668479in}}%
\pgfpathlineto{\pgfqpoint{2.107303in}{1.662856in}}%
\pgfpathlineto{\pgfqpoint{2.109669in}{1.658235in}}%
\pgfpathlineto{\pgfqpoint{2.122391in}{1.643026in}}%
\pgfpathlineto{\pgfqpoint{2.132464in}{1.643757in}}%
\pgfpathlineto{\pgfqpoint{2.136667in}{1.646574in}}%
\pgfpathlineto{\pgfqpoint{2.143748in}{1.640811in}}%
\pgfpathlineto{\pgfqpoint{2.153377in}{1.649798in}}%
\pgfpathlineto{\pgfqpoint{2.156972in}{1.644308in}}%
\pgfpathlineto{\pgfqpoint{2.161790in}{1.650391in}}%
\pgfpathlineto{\pgfqpoint{2.176831in}{1.658426in}}%
\pgfpathlineto{\pgfqpoint{2.188229in}{1.662007in}}%
\pgfpathlineto{\pgfqpoint{2.192935in}{1.614287in}}%
\pgfpathclose%
\pgfusepath{fill}%
\end{pgfscope}%
\begin{pgfscope}%
\pgfpathrectangle{\pgfqpoint{0.100000in}{0.100000in}}{\pgfqpoint{3.007045in}{1.925000in}}%
\pgfusepath{clip}%
\pgfsetbuttcap%
\pgfsetmiterjoin%
\definecolor{currentfill}{rgb}{0.501961,0.790588,0.501438}%
\pgfsetfillcolor{currentfill}%
\pgfsetlinewidth{0.000000pt}%
\definecolor{currentstroke}{rgb}{0.000000,0.000000,0.000000}%
\pgfsetstrokecolor{currentstroke}%
\pgfsetstrokeopacity{0.000000}%
\pgfsetdash{}{0pt}%
\pgfpathmoveto{\pgfqpoint{2.753651in}{1.370464in}}%
\pgfpathlineto{\pgfqpoint{2.746998in}{1.361954in}}%
\pgfpathlineto{\pgfqpoint{2.738618in}{1.356791in}}%
\pgfpathlineto{\pgfqpoint{2.732190in}{1.364248in}}%
\pgfpathlineto{\pgfqpoint{2.735363in}{1.367506in}}%
\pgfpathlineto{\pgfqpoint{2.727132in}{1.374300in}}%
\pgfpathlineto{\pgfqpoint{2.723613in}{1.374911in}}%
\pgfpathlineto{\pgfqpoint{2.720526in}{1.370514in}}%
\pgfpathlineto{\pgfqpoint{2.723365in}{1.366527in}}%
\pgfpathlineto{\pgfqpoint{2.711768in}{1.357635in}}%
\pgfpathlineto{\pgfqpoint{2.701015in}{1.357842in}}%
\pgfpathlineto{\pgfqpoint{2.698167in}{1.352791in}}%
\pgfpathlineto{\pgfqpoint{2.698167in}{1.347287in}}%
\pgfpathlineto{\pgfqpoint{2.694956in}{1.344848in}}%
\pgfpathlineto{\pgfqpoint{2.692918in}{1.348382in}}%
\pgfpathlineto{\pgfqpoint{2.688600in}{1.346784in}}%
\pgfpathlineto{\pgfqpoint{2.686221in}{1.354519in}}%
\pgfpathlineto{\pgfqpoint{2.679466in}{1.361047in}}%
\pgfpathlineto{\pgfqpoint{2.677192in}{1.366767in}}%
\pgfpathlineto{\pgfqpoint{2.679506in}{1.367379in}}%
\pgfpathlineto{\pgfqpoint{2.686454in}{1.378941in}}%
\pgfpathlineto{\pgfqpoint{2.691171in}{1.379954in}}%
\pgfpathlineto{\pgfqpoint{2.692396in}{1.393752in}}%
\pgfpathlineto{\pgfqpoint{2.691983in}{1.402891in}}%
\pgfpathlineto{\pgfqpoint{2.695962in}{1.404095in}}%
\pgfpathlineto{\pgfqpoint{2.727566in}{1.410162in}}%
\pgfpathlineto{\pgfqpoint{2.721527in}{1.432985in}}%
\pgfpathlineto{\pgfqpoint{2.727454in}{1.434295in}}%
\pgfpathlineto{\pgfqpoint{2.732166in}{1.431244in}}%
\pgfpathlineto{\pgfqpoint{2.733867in}{1.427016in}}%
\pgfpathlineto{\pgfqpoint{2.739778in}{1.426983in}}%
\pgfpathlineto{\pgfqpoint{2.745672in}{1.421591in}}%
\pgfpathlineto{\pgfqpoint{2.747986in}{1.413389in}}%
\pgfpathlineto{\pgfqpoint{2.752320in}{1.406665in}}%
\pgfpathlineto{\pgfqpoint{2.760698in}{1.403338in}}%
\pgfpathlineto{\pgfqpoint{2.765754in}{1.403877in}}%
\pgfpathlineto{\pgfqpoint{2.768627in}{1.400117in}}%
\pgfpathlineto{\pgfqpoint{2.763125in}{1.394125in}}%
\pgfpathlineto{\pgfqpoint{2.762147in}{1.386718in}}%
\pgfpathlineto{\pgfqpoint{2.753651in}{1.370464in}}%
\pgfpathclose%
\pgfusepath{fill}%
\end{pgfscope}%
\begin{pgfscope}%
\pgfpathrectangle{\pgfqpoint{0.100000in}{0.100000in}}{\pgfqpoint{3.007045in}{1.925000in}}%
\pgfusepath{clip}%
\pgfsetbuttcap%
\pgfsetmiterjoin%
\definecolor{currentfill}{rgb}{0.451765,0.767090,0.461207}%
\pgfsetfillcolor{currentfill}%
\pgfsetlinewidth{0.000000pt}%
\definecolor{currentstroke}{rgb}{0.000000,0.000000,0.000000}%
\pgfsetstrokecolor{currentstroke}%
\pgfsetstrokeopacity{0.000000}%
\pgfsetdash{}{0pt}%
\pgfpathmoveto{\pgfqpoint{1.565600in}{1.146405in}}%
\pgfpathlineto{\pgfqpoint{1.594074in}{1.145204in}}%
\pgfpathlineto{\pgfqpoint{1.593668in}{1.127954in}}%
\pgfpathlineto{\pgfqpoint{1.593165in}{1.116517in}}%
\pgfpathlineto{\pgfqpoint{1.536726in}{1.118994in}}%
\pgfpathlineto{\pgfqpoint{1.507527in}{1.120611in}}%
\pgfpathlineto{\pgfqpoint{1.509458in}{1.149192in}}%
\pgfpathlineto{\pgfqpoint{1.565600in}{1.146405in}}%
\pgfpathclose%
\pgfusepath{fill}%
\end{pgfscope}%
\begin{pgfscope}%
\pgfpathrectangle{\pgfqpoint{0.100000in}{0.100000in}}{\pgfqpoint{3.007045in}{1.925000in}}%
\pgfusepath{clip}%
\pgfsetbuttcap%
\pgfsetmiterjoin%
\definecolor{currentfill}{rgb}{0.886505,0.956171,0.866128}%
\pgfsetfillcolor{currentfill}%
\pgfsetlinewidth{0.000000pt}%
\definecolor{currentstroke}{rgb}{0.000000,0.000000,0.000000}%
\pgfsetstrokecolor{currentstroke}%
\pgfsetstrokeopacity{0.000000}%
\pgfsetdash{}{0pt}%
\pgfpathmoveto{\pgfqpoint{0.558440in}{0.376370in}}%
\pgfpathlineto{\pgfqpoint{0.557442in}{0.379760in}}%
\pgfpathlineto{\pgfqpoint{0.559412in}{0.380419in}}%
\pgfpathlineto{\pgfqpoint{0.561194in}{0.379545in}}%
\pgfpathlineto{\pgfqpoint{0.561660in}{0.376234in}}%
\pgfpathlineto{\pgfqpoint{0.560353in}{0.375728in}}%
\pgfpathlineto{\pgfqpoint{0.558440in}{0.376370in}}%
\pgfpathclose%
\pgfusepath{fill}%
\end{pgfscope}%
\begin{pgfscope}%
\pgfpathrectangle{\pgfqpoint{0.100000in}{0.100000in}}{\pgfqpoint{3.007045in}{1.925000in}}%
\pgfusepath{clip}%
\pgfsetbuttcap%
\pgfsetmiterjoin%
\definecolor{currentfill}{rgb}{0.886505,0.956171,0.866128}%
\pgfsetfillcolor{currentfill}%
\pgfsetlinewidth{0.000000pt}%
\definecolor{currentstroke}{rgb}{0.000000,0.000000,0.000000}%
\pgfsetstrokecolor{currentstroke}%
\pgfsetstrokeopacity{0.000000}%
\pgfsetdash{}{0pt}%
\pgfpathmoveto{\pgfqpoint{0.562662in}{0.374546in}}%
\pgfpathlineto{\pgfqpoint{0.563801in}{0.376835in}}%
\pgfpathlineto{\pgfqpoint{0.565086in}{0.376399in}}%
\pgfpathlineto{\pgfqpoint{0.565958in}{0.374725in}}%
\pgfpathlineto{\pgfqpoint{0.562662in}{0.374546in}}%
\pgfpathclose%
\pgfusepath{fill}%
\end{pgfscope}%
\begin{pgfscope}%
\pgfpathrectangle{\pgfqpoint{0.100000in}{0.100000in}}{\pgfqpoint{3.007045in}{1.925000in}}%
\pgfusepath{clip}%
\pgfsetbuttcap%
\pgfsetmiterjoin%
\definecolor{currentfill}{rgb}{0.886505,0.956171,0.866128}%
\pgfsetfillcolor{currentfill}%
\pgfsetlinewidth{0.000000pt}%
\definecolor{currentstroke}{rgb}{0.000000,0.000000,0.000000}%
\pgfsetstrokecolor{currentstroke}%
\pgfsetstrokeopacity{0.000000}%
\pgfsetdash{}{0pt}%
\pgfpathmoveto{\pgfqpoint{0.630180in}{0.347609in}}%
\pgfpathlineto{\pgfqpoint{0.628737in}{0.345721in}}%
\pgfpathlineto{\pgfqpoint{0.627895in}{0.348703in}}%
\pgfpathlineto{\pgfqpoint{0.625581in}{0.348130in}}%
\pgfpathlineto{\pgfqpoint{0.628285in}{0.352021in}}%
\pgfpathlineto{\pgfqpoint{0.631284in}{0.352146in}}%
\pgfpathlineto{\pgfqpoint{0.630506in}{0.350620in}}%
\pgfpathlineto{\pgfqpoint{0.631605in}{0.349217in}}%
\pgfpathlineto{\pgfqpoint{0.634087in}{0.348313in}}%
\pgfpathlineto{\pgfqpoint{0.632370in}{0.345968in}}%
\pgfpathlineto{\pgfqpoint{0.631280in}{0.347967in}}%
\pgfpathlineto{\pgfqpoint{0.630180in}{0.347609in}}%
\pgfpathclose%
\pgfusepath{fill}%
\end{pgfscope}%
\begin{pgfscope}%
\pgfpathrectangle{\pgfqpoint{0.100000in}{0.100000in}}{\pgfqpoint{3.007045in}{1.925000in}}%
\pgfusepath{clip}%
\pgfsetbuttcap%
\pgfsetmiterjoin%
\definecolor{currentfill}{rgb}{0.886505,0.956171,0.866128}%
\pgfsetfillcolor{currentfill}%
\pgfsetlinewidth{0.000000pt}%
\definecolor{currentstroke}{rgb}{0.000000,0.000000,0.000000}%
\pgfsetstrokecolor{currentstroke}%
\pgfsetstrokeopacity{0.000000}%
\pgfsetdash{}{0pt}%
\pgfpathmoveto{\pgfqpoint{0.606094in}{0.356546in}}%
\pgfpathlineto{\pgfqpoint{0.606285in}{0.359081in}}%
\pgfpathlineto{\pgfqpoint{0.608596in}{0.359253in}}%
\pgfpathlineto{\pgfqpoint{0.608292in}{0.357179in}}%
\pgfpathlineto{\pgfqpoint{0.606094in}{0.356546in}}%
\pgfpathclose%
\pgfusepath{fill}%
\end{pgfscope}%
\begin{pgfscope}%
\pgfpathrectangle{\pgfqpoint{0.100000in}{0.100000in}}{\pgfqpoint{3.007045in}{1.925000in}}%
\pgfusepath{clip}%
\pgfsetbuttcap%
\pgfsetmiterjoin%
\definecolor{currentfill}{rgb}{0.886505,0.956171,0.866128}%
\pgfsetfillcolor{currentfill}%
\pgfsetlinewidth{0.000000pt}%
\definecolor{currentstroke}{rgb}{0.000000,0.000000,0.000000}%
\pgfsetstrokecolor{currentstroke}%
\pgfsetstrokeopacity{0.000000}%
\pgfsetdash{}{0pt}%
\pgfpathmoveto{\pgfqpoint{0.628222in}{0.338692in}}%
\pgfpathlineto{\pgfqpoint{0.634226in}{0.342115in}}%
\pgfpathlineto{\pgfqpoint{0.635742in}{0.340232in}}%
\pgfpathlineto{\pgfqpoint{0.633794in}{0.339431in}}%
\pgfpathlineto{\pgfqpoint{0.628222in}{0.338692in}}%
\pgfpathclose%
\pgfusepath{fill}%
\end{pgfscope}%
\begin{pgfscope}%
\pgfpathrectangle{\pgfqpoint{0.100000in}{0.100000in}}{\pgfqpoint{3.007045in}{1.925000in}}%
\pgfusepath{clip}%
\pgfsetbuttcap%
\pgfsetmiterjoin%
\definecolor{currentfill}{rgb}{0.886505,0.956171,0.866128}%
\pgfsetfillcolor{currentfill}%
\pgfsetlinewidth{0.000000pt}%
\definecolor{currentstroke}{rgb}{0.000000,0.000000,0.000000}%
\pgfsetstrokecolor{currentstroke}%
\pgfsetstrokeopacity{0.000000}%
\pgfsetdash{}{0pt}%
\pgfpathmoveto{\pgfqpoint{0.646196in}{0.346313in}}%
\pgfpathlineto{\pgfqpoint{0.645147in}{0.347754in}}%
\pgfpathlineto{\pgfqpoint{0.648533in}{0.349423in}}%
\pgfpathlineto{\pgfqpoint{0.647621in}{0.351703in}}%
\pgfpathlineto{\pgfqpoint{0.645804in}{0.351697in}}%
\pgfpathlineto{\pgfqpoint{0.644687in}{0.352622in}}%
\pgfpathlineto{\pgfqpoint{0.640184in}{0.353124in}}%
\pgfpathlineto{\pgfqpoint{0.636431in}{0.353197in}}%
\pgfpathlineto{\pgfqpoint{0.634892in}{0.354431in}}%
\pgfpathlineto{\pgfqpoint{0.633377in}{0.353497in}}%
\pgfpathlineto{\pgfqpoint{0.631296in}{0.353562in}}%
\pgfpathlineto{\pgfqpoint{0.631571in}{0.355560in}}%
\pgfpathlineto{\pgfqpoint{0.627061in}{0.355667in}}%
\pgfpathlineto{\pgfqpoint{0.624673in}{0.356334in}}%
\pgfpathlineto{\pgfqpoint{0.622641in}{0.358315in}}%
\pgfpathlineto{\pgfqpoint{0.624401in}{0.359859in}}%
\pgfpathlineto{\pgfqpoint{0.627993in}{0.361739in}}%
\pgfpathlineto{\pgfqpoint{0.626212in}{0.363517in}}%
\pgfpathlineto{\pgfqpoint{0.624339in}{0.363616in}}%
\pgfpathlineto{\pgfqpoint{0.620521in}{0.360581in}}%
\pgfpathlineto{\pgfqpoint{0.617999in}{0.360096in}}%
\pgfpathlineto{\pgfqpoint{0.615778in}{0.360562in}}%
\pgfpathlineto{\pgfqpoint{0.614484in}{0.358255in}}%
\pgfpathlineto{\pgfqpoint{0.612634in}{0.358313in}}%
\pgfpathlineto{\pgfqpoint{0.608448in}{0.361014in}}%
\pgfpathlineto{\pgfqpoint{0.608279in}{0.362605in}}%
\pgfpathlineto{\pgfqpoint{0.609968in}{0.362762in}}%
\pgfpathlineto{\pgfqpoint{0.610410in}{0.364673in}}%
\pgfpathlineto{\pgfqpoint{0.610099in}{0.367347in}}%
\pgfpathlineto{\pgfqpoint{0.606743in}{0.363415in}}%
\pgfpathlineto{\pgfqpoint{0.606556in}{0.361958in}}%
\pgfpathlineto{\pgfqpoint{0.603910in}{0.362284in}}%
\pgfpathlineto{\pgfqpoint{0.602151in}{0.363752in}}%
\pgfpathlineto{\pgfqpoint{0.603264in}{0.366602in}}%
\pgfpathlineto{\pgfqpoint{0.602235in}{0.369208in}}%
\pgfpathlineto{\pgfqpoint{0.601172in}{0.368904in}}%
\pgfpathlineto{\pgfqpoint{0.600989in}{0.365044in}}%
\pgfpathlineto{\pgfqpoint{0.595049in}{0.365340in}}%
\pgfpathlineto{\pgfqpoint{0.596365in}{0.363673in}}%
\pgfpathlineto{\pgfqpoint{0.594450in}{0.363174in}}%
\pgfpathlineto{\pgfqpoint{0.593766in}{0.364861in}}%
\pgfpathlineto{\pgfqpoint{0.592021in}{0.364091in}}%
\pgfpathlineto{\pgfqpoint{0.589288in}{0.365861in}}%
\pgfpathlineto{\pgfqpoint{0.585892in}{0.369053in}}%
\pgfpathlineto{\pgfqpoint{0.581946in}{0.370113in}}%
\pgfpathlineto{\pgfqpoint{0.578510in}{0.369321in}}%
\pgfpathlineto{\pgfqpoint{0.576281in}{0.370552in}}%
\pgfpathlineto{\pgfqpoint{0.575149in}{0.374376in}}%
\pgfpathlineto{\pgfqpoint{0.576135in}{0.376449in}}%
\pgfpathlineto{\pgfqpoint{0.580707in}{0.375972in}}%
\pgfpathlineto{\pgfqpoint{0.586276in}{0.378070in}}%
\pgfpathlineto{\pgfqpoint{0.587415in}{0.376136in}}%
\pgfpathlineto{\pgfqpoint{0.589357in}{0.375160in}}%
\pgfpathlineto{\pgfqpoint{0.590661in}{0.375506in}}%
\pgfpathlineto{\pgfqpoint{0.595051in}{0.373915in}}%
\pgfpathlineto{\pgfqpoint{0.597649in}{0.371586in}}%
\pgfpathlineto{\pgfqpoint{0.596738in}{0.368224in}}%
\pgfpathlineto{\pgfqpoint{0.598337in}{0.367723in}}%
\pgfpathlineto{\pgfqpoint{0.600669in}{0.370774in}}%
\pgfpathlineto{\pgfqpoint{0.604276in}{0.369985in}}%
\pgfpathlineto{\pgfqpoint{0.605343in}{0.367909in}}%
\pgfpathlineto{\pgfqpoint{0.607337in}{0.368055in}}%
\pgfpathlineto{\pgfqpoint{0.620695in}{0.370601in}}%
\pgfpathlineto{\pgfqpoint{0.625712in}{0.370166in}}%
\pgfpathlineto{\pgfqpoint{0.628012in}{0.370317in}}%
\pgfpathlineto{\pgfqpoint{0.633395in}{0.367875in}}%
\pgfpathlineto{\pgfqpoint{0.636602in}{0.364776in}}%
\pgfpathlineto{\pgfqpoint{0.635557in}{0.363488in}}%
\pgfpathlineto{\pgfqpoint{0.635523in}{0.359748in}}%
\pgfpathlineto{\pgfqpoint{0.637405in}{0.360249in}}%
\pgfpathlineto{\pgfqpoint{0.640275in}{0.358605in}}%
\pgfpathlineto{\pgfqpoint{0.639530in}{0.356749in}}%
\pgfpathlineto{\pgfqpoint{0.640742in}{0.355326in}}%
\pgfpathlineto{\pgfqpoint{0.642306in}{0.356377in}}%
\pgfpathlineto{\pgfqpoint{0.641680in}{0.358719in}}%
\pgfpathlineto{\pgfqpoint{0.640362in}{0.360039in}}%
\pgfpathlineto{\pgfqpoint{0.641503in}{0.361170in}}%
\pgfpathlineto{\pgfqpoint{0.644772in}{0.362563in}}%
\pgfpathlineto{\pgfqpoint{0.646911in}{0.363995in}}%
\pgfpathlineto{\pgfqpoint{0.651915in}{0.364863in}}%
\pgfpathlineto{\pgfqpoint{0.654420in}{0.364388in}}%
\pgfpathlineto{\pgfqpoint{0.656995in}{0.364705in}}%
\pgfpathlineto{\pgfqpoint{0.659517in}{0.364305in}}%
\pgfpathlineto{\pgfqpoint{0.665903in}{0.362144in}}%
\pgfpathlineto{\pgfqpoint{0.665627in}{0.363168in}}%
\pgfpathlineto{\pgfqpoint{0.669120in}{0.362795in}}%
\pgfpathlineto{\pgfqpoint{0.670704in}{0.361448in}}%
\pgfpathlineto{\pgfqpoint{0.668605in}{0.361721in}}%
\pgfpathlineto{\pgfqpoint{0.667310in}{0.360198in}}%
\pgfpathlineto{\pgfqpoint{0.664908in}{0.362318in}}%
\pgfpathlineto{\pgfqpoint{0.663515in}{0.360688in}}%
\pgfpathlineto{\pgfqpoint{0.660386in}{0.363372in}}%
\pgfpathlineto{\pgfqpoint{0.659049in}{0.361833in}}%
\pgfpathlineto{\pgfqpoint{0.655995in}{0.364392in}}%
\pgfpathlineto{\pgfqpoint{0.653338in}{0.361411in}}%
\pgfpathlineto{\pgfqpoint{0.654252in}{0.360591in}}%
\pgfpathlineto{\pgfqpoint{0.649075in}{0.354518in}}%
\pgfpathlineto{\pgfqpoint{0.647002in}{0.353387in}}%
\pgfpathlineto{\pgfqpoint{0.650028in}{0.350774in}}%
\pgfpathlineto{\pgfqpoint{0.646196in}{0.346313in}}%
\pgfpathclose%
\pgfusepath{fill}%
\end{pgfscope}%
\begin{pgfscope}%
\pgfpathrectangle{\pgfqpoint{0.100000in}{0.100000in}}{\pgfqpoint{3.007045in}{1.925000in}}%
\pgfusepath{clip}%
\pgfsetbuttcap%
\pgfsetmiterjoin%
\definecolor{currentfill}{rgb}{0.420392,0.751711,0.445829}%
\pgfsetfillcolor{currentfill}%
\pgfsetlinewidth{0.000000pt}%
\definecolor{currentstroke}{rgb}{0.000000,0.000000,0.000000}%
\pgfsetstrokecolor{currentstroke}%
\pgfsetstrokeopacity{0.000000}%
\pgfsetdash{}{0pt}%
\pgfpathmoveto{\pgfqpoint{1.588665in}{1.618076in}}%
\pgfpathlineto{\pgfqpoint{1.587550in}{1.595331in}}%
\pgfpathlineto{\pgfqpoint{1.554716in}{1.596972in}}%
\pgfpathlineto{\pgfqpoint{1.518744in}{1.599173in}}%
\pgfpathlineto{\pgfqpoint{1.514131in}{1.607335in}}%
\pgfpathlineto{\pgfqpoint{1.514093in}{1.614101in}}%
\pgfpathlineto{\pgfqpoint{1.517258in}{1.619082in}}%
\pgfpathlineto{\pgfqpoint{1.518046in}{1.627144in}}%
\pgfpathlineto{\pgfqpoint{1.516525in}{1.631523in}}%
\pgfpathlineto{\pgfqpoint{1.519039in}{1.638417in}}%
\pgfpathlineto{\pgfqpoint{1.518209in}{1.642479in}}%
\pgfpathlineto{\pgfqpoint{1.514103in}{1.645425in}}%
\pgfpathlineto{\pgfqpoint{1.540553in}{1.643789in}}%
\pgfpathlineto{\pgfqpoint{1.569331in}{1.642166in}}%
\pgfpathlineto{\pgfqpoint{1.588163in}{1.641202in}}%
\pgfpathlineto{\pgfqpoint{1.587099in}{1.618157in}}%
\pgfpathlineto{\pgfqpoint{1.588665in}{1.618076in}}%
\pgfpathclose%
\pgfusepath{fill}%
\end{pgfscope}%
\begin{pgfscope}%
\pgfpathrectangle{\pgfqpoint{0.100000in}{0.100000in}}{\pgfqpoint{3.007045in}{1.925000in}}%
\pgfusepath{clip}%
\pgfsetbuttcap%
\pgfsetmiterjoin%
\definecolor{currentfill}{rgb}{0.801615,0.922215,0.775579}%
\pgfsetfillcolor{currentfill}%
\pgfsetlinewidth{0.000000pt}%
\definecolor{currentstroke}{rgb}{0.000000,0.000000,0.000000}%
\pgfsetstrokecolor{currentstroke}%
\pgfsetstrokeopacity{0.000000}%
\pgfsetdash{}{0pt}%
\pgfpathmoveto{\pgfqpoint{1.509153in}{1.008633in}}%
\pgfpathlineto{\pgfqpoint{1.507272in}{0.975589in}}%
\pgfpathlineto{\pgfqpoint{1.456990in}{0.978623in}}%
\pgfpathlineto{\pgfqpoint{1.459593in}{1.011387in}}%
\pgfpathlineto{\pgfqpoint{1.453227in}{1.011793in}}%
\pgfpathlineto{\pgfqpoint{1.454880in}{1.037434in}}%
\pgfpathlineto{\pgfqpoint{1.476575in}{1.036016in}}%
\pgfpathlineto{\pgfqpoint{1.476919in}{1.041763in}}%
\pgfpathlineto{\pgfqpoint{1.505407in}{1.040070in}}%
\pgfpathlineto{\pgfqpoint{1.505902in}{1.034236in}}%
\pgfpathlineto{\pgfqpoint{1.504587in}{1.008926in}}%
\pgfpathlineto{\pgfqpoint{1.509153in}{1.008633in}}%
\pgfpathclose%
\pgfusepath{fill}%
\end{pgfscope}%
\begin{pgfscope}%
\pgfpathrectangle{\pgfqpoint{0.100000in}{0.100000in}}{\pgfqpoint{3.007045in}{1.925000in}}%
\pgfusepath{clip}%
\pgfsetbuttcap%
\pgfsetmiterjoin%
\definecolor{currentfill}{rgb}{0.507497,0.793172,0.505990}%
\pgfsetfillcolor{currentfill}%
\pgfsetlinewidth{0.000000pt}%
\definecolor{currentstroke}{rgb}{0.000000,0.000000,0.000000}%
\pgfsetstrokecolor{currentstroke}%
\pgfsetstrokeopacity{0.000000}%
\pgfsetdash{}{0pt}%
\pgfpathmoveto{\pgfqpoint{1.753660in}{1.413597in}}%
\pgfpathlineto{\pgfqpoint{1.730861in}{1.413920in}}%
\pgfpathlineto{\pgfqpoint{1.731061in}{1.429926in}}%
\pgfpathlineto{\pgfqpoint{1.721870in}{1.430045in}}%
\pgfpathlineto{\pgfqpoint{1.722176in}{1.453096in}}%
\pgfpathlineto{\pgfqpoint{1.721890in}{1.476069in}}%
\pgfpathlineto{\pgfqpoint{1.744188in}{1.475811in}}%
\pgfpathlineto{\pgfqpoint{1.767169in}{1.475679in}}%
\pgfpathlineto{\pgfqpoint{1.767569in}{1.469929in}}%
\pgfpathlineto{\pgfqpoint{1.778980in}{1.469921in}}%
\pgfpathlineto{\pgfqpoint{1.779266in}{1.452736in}}%
\pgfpathlineto{\pgfqpoint{1.779303in}{1.429779in}}%
\pgfpathlineto{\pgfqpoint{1.776405in}{1.429765in}}%
\pgfpathlineto{\pgfqpoint{1.753741in}{1.429789in}}%
\pgfpathlineto{\pgfqpoint{1.753660in}{1.413597in}}%
\pgfpathclose%
\pgfusepath{fill}%
\end{pgfscope}%
\begin{pgfscope}%
\pgfpathrectangle{\pgfqpoint{0.100000in}{0.100000in}}{\pgfqpoint{3.007045in}{1.925000in}}%
\pgfusepath{clip}%
\pgfsetbuttcap%
\pgfsetmiterjoin%
\definecolor{currentfill}{rgb}{0.345098,0.714802,0.408920}%
\pgfsetfillcolor{currentfill}%
\pgfsetlinewidth{0.000000pt}%
\definecolor{currentstroke}{rgb}{0.000000,0.000000,0.000000}%
\pgfsetstrokecolor{currentstroke}%
\pgfsetstrokeopacity{0.000000}%
\pgfsetdash{}{0pt}%
\pgfpathmoveto{\pgfqpoint{1.798792in}{1.236593in}}%
\pgfpathlineto{\pgfqpoint{1.790743in}{1.236578in}}%
\pgfpathlineto{\pgfqpoint{1.734017in}{1.237625in}}%
\pgfpathlineto{\pgfqpoint{1.732568in}{1.242722in}}%
\pgfpathlineto{\pgfqpoint{1.727958in}{1.247635in}}%
\pgfpathlineto{\pgfqpoint{1.730711in}{1.250739in}}%
\pgfpathlineto{\pgfqpoint{1.731811in}{1.258499in}}%
\pgfpathlineto{\pgfqpoint{1.753153in}{1.258349in}}%
\pgfpathlineto{\pgfqpoint{1.753309in}{1.275315in}}%
\pgfpathlineto{\pgfqpoint{1.764595in}{1.275306in}}%
\pgfpathlineto{\pgfqpoint{1.787185in}{1.275241in}}%
\pgfpathlineto{\pgfqpoint{1.798597in}{1.275256in}}%
\pgfpathlineto{\pgfqpoint{1.798792in}{1.236593in}}%
\pgfpathclose%
\pgfusepath{fill}%
\end{pgfscope}%
\begin{pgfscope}%
\pgfpathrectangle{\pgfqpoint{0.100000in}{0.100000in}}{\pgfqpoint{3.007045in}{1.925000in}}%
\pgfusepath{clip}%
\pgfsetbuttcap%
\pgfsetmiterjoin%
\definecolor{currentfill}{rgb}{0.737147,0.895517,0.710834}%
\pgfsetfillcolor{currentfill}%
\pgfsetlinewidth{0.000000pt}%
\definecolor{currentstroke}{rgb}{0.000000,0.000000,0.000000}%
\pgfsetstrokecolor{currentstroke}%
\pgfsetstrokeopacity{0.000000}%
\pgfsetdash{}{0pt}%
\pgfpathmoveto{\pgfqpoint{1.634574in}{0.311704in}}%
\pgfpathlineto{\pgfqpoint{1.632618in}{0.311751in}}%
\pgfpathlineto{\pgfqpoint{1.628886in}{0.326991in}}%
\pgfpathlineto{\pgfqpoint{1.629356in}{0.349314in}}%
\pgfpathlineto{\pgfqpoint{1.631677in}{0.363791in}}%
\pgfpathlineto{\pgfqpoint{1.647149in}{0.392795in}}%
\pgfpathlineto{\pgfqpoint{1.651127in}{0.395908in}}%
\pgfpathlineto{\pgfqpoint{1.655482in}{0.405989in}}%
\pgfpathlineto{\pgfqpoint{1.661560in}{0.414117in}}%
\pgfpathlineto{\pgfqpoint{1.664160in}{0.412781in}}%
\pgfpathlineto{\pgfqpoint{1.662731in}{0.409312in}}%
\pgfpathlineto{\pgfqpoint{1.652282in}{0.397408in}}%
\pgfpathlineto{\pgfqpoint{1.643553in}{0.383920in}}%
\pgfpathlineto{\pgfqpoint{1.637879in}{0.372867in}}%
\pgfpathlineto{\pgfqpoint{1.633473in}{0.362042in}}%
\pgfpathlineto{\pgfqpoint{1.630879in}{0.352084in}}%
\pgfpathlineto{\pgfqpoint{1.629805in}{0.336635in}}%
\pgfpathlineto{\pgfqpoint{1.631675in}{0.322573in}}%
\pgfpathlineto{\pgfqpoint{1.634574in}{0.311704in}}%
\pgfpathclose%
\pgfusepath{fill}%
\end{pgfscope}%
\begin{pgfscope}%
\pgfpathrectangle{\pgfqpoint{0.100000in}{0.100000in}}{\pgfqpoint{3.007045in}{1.925000in}}%
\pgfusepath{clip}%
\pgfsetbuttcap%
\pgfsetmiterjoin%
\definecolor{currentfill}{rgb}{0.737147,0.895517,0.710834}%
\pgfsetfillcolor{currentfill}%
\pgfsetlinewidth{0.000000pt}%
\definecolor{currentstroke}{rgb}{0.000000,0.000000,0.000000}%
\pgfsetstrokecolor{currentstroke}%
\pgfsetstrokeopacity{0.000000}%
\pgfsetdash{}{0pt}%
\pgfpathmoveto{\pgfqpoint{1.625397in}{0.311924in}}%
\pgfpathlineto{\pgfqpoint{1.599953in}{0.312537in}}%
\pgfpathlineto{\pgfqpoint{1.592550in}{0.313976in}}%
\pgfpathlineto{\pgfqpoint{1.592904in}{0.325092in}}%
\pgfpathlineto{\pgfqpoint{1.566652in}{0.326194in}}%
\pgfpathlineto{\pgfqpoint{1.567623in}{0.344460in}}%
\pgfpathlineto{\pgfqpoint{1.564709in}{0.344532in}}%
\pgfpathlineto{\pgfqpoint{1.563594in}{0.356706in}}%
\pgfpathlineto{\pgfqpoint{1.560356in}{0.365181in}}%
\pgfpathlineto{\pgfqpoint{1.545690in}{0.365433in}}%
\pgfpathlineto{\pgfqpoint{1.547367in}{0.412624in}}%
\pgfpathlineto{\pgfqpoint{1.601772in}{0.410582in}}%
\pgfpathlineto{\pgfqpoint{1.600656in}{0.414463in}}%
\pgfpathlineto{\pgfqpoint{1.605884in}{0.418493in}}%
\pgfpathlineto{\pgfqpoint{1.620569in}{0.414900in}}%
\pgfpathlineto{\pgfqpoint{1.632363in}{0.431928in}}%
\pgfpathlineto{\pgfqpoint{1.645296in}{0.442654in}}%
\pgfpathlineto{\pgfqpoint{1.651860in}{0.442260in}}%
\pgfpathlineto{\pgfqpoint{1.655495in}{0.438831in}}%
\pgfpathlineto{\pgfqpoint{1.661661in}{0.438978in}}%
\pgfpathlineto{\pgfqpoint{1.667536in}{0.431953in}}%
\pgfpathlineto{\pgfqpoint{1.666217in}{0.423659in}}%
\pgfpathlineto{\pgfqpoint{1.665504in}{0.419888in}}%
\pgfpathlineto{\pgfqpoint{1.658131in}{0.413570in}}%
\pgfpathlineto{\pgfqpoint{1.654267in}{0.413508in}}%
\pgfpathlineto{\pgfqpoint{1.651083in}{0.406992in}}%
\pgfpathlineto{\pgfqpoint{1.642502in}{0.393804in}}%
\pgfpathlineto{\pgfqpoint{1.637586in}{0.397797in}}%
\pgfpathlineto{\pgfqpoint{1.631800in}{0.395290in}}%
\pgfpathlineto{\pgfqpoint{1.632028in}{0.388063in}}%
\pgfpathlineto{\pgfqpoint{1.638379in}{0.385411in}}%
\pgfpathlineto{\pgfqpoint{1.634179in}{0.376996in}}%
\pgfpathlineto{\pgfqpoint{1.628882in}{0.361108in}}%
\pgfpathlineto{\pgfqpoint{1.626077in}{0.346812in}}%
\pgfpathlineto{\pgfqpoint{1.619237in}{0.338192in}}%
\pgfpathlineto{\pgfqpoint{1.618817in}{0.328087in}}%
\pgfpathlineto{\pgfqpoint{1.623729in}{0.322672in}}%
\pgfpathlineto{\pgfqpoint{1.625397in}{0.311924in}}%
\pgfpathclose%
\pgfusepath{fill}%
\end{pgfscope}%
\begin{pgfscope}%
\pgfpathrectangle{\pgfqpoint{0.100000in}{0.100000in}}{\pgfqpoint{3.007045in}{1.925000in}}%
\pgfusepath{clip}%
\pgfsetbuttcap%
\pgfsetmiterjoin%
\definecolor{currentfill}{rgb}{0.535179,0.806090,0.528750}%
\pgfsetfillcolor{currentfill}%
\pgfsetlinewidth{0.000000pt}%
\definecolor{currentstroke}{rgb}{0.000000,0.000000,0.000000}%
\pgfsetstrokecolor{currentstroke}%
\pgfsetstrokeopacity{0.000000}%
\pgfsetdash{}{0pt}%
\pgfpathmoveto{\pgfqpoint{1.703058in}{1.453389in}}%
\pgfpathlineto{\pgfqpoint{1.682306in}{1.453680in}}%
\pgfpathlineto{\pgfqpoint{1.682752in}{1.476661in}}%
\pgfpathlineto{\pgfqpoint{1.637239in}{1.477803in}}%
\pgfpathlineto{\pgfqpoint{1.637041in}{1.477807in}}%
\pgfpathlineto{\pgfqpoint{1.637786in}{1.500890in}}%
\pgfpathlineto{\pgfqpoint{1.671967in}{1.499872in}}%
\pgfpathlineto{\pgfqpoint{1.703851in}{1.499277in}}%
\pgfpathlineto{\pgfqpoint{1.703058in}{1.453389in}}%
\pgfpathclose%
\pgfusepath{fill}%
\end{pgfscope}%
\begin{pgfscope}%
\pgfpathrectangle{\pgfqpoint{0.100000in}{0.100000in}}{\pgfqpoint{3.007045in}{1.925000in}}%
\pgfusepath{clip}%
\pgfsetbuttcap%
\pgfsetmiterjoin%
\definecolor{currentfill}{rgb}{0.282353,0.684045,0.378162}%
\pgfsetfillcolor{currentfill}%
\pgfsetlinewidth{0.000000pt}%
\definecolor{currentstroke}{rgb}{0.000000,0.000000,0.000000}%
\pgfsetstrokecolor{currentstroke}%
\pgfsetstrokeopacity{0.000000}%
\pgfsetdash{}{0pt}%
\pgfpathmoveto{\pgfqpoint{1.944456in}{0.907639in}}%
\pgfpathlineto{\pgfqpoint{1.944537in}{0.902457in}}%
\pgfpathlineto{\pgfqpoint{1.941745in}{0.895933in}}%
\pgfpathlineto{\pgfqpoint{1.927775in}{0.895762in}}%
\pgfpathlineto{\pgfqpoint{1.907859in}{0.895543in}}%
\pgfpathlineto{\pgfqpoint{1.907800in}{0.901293in}}%
\pgfpathlineto{\pgfqpoint{1.887849in}{0.901272in}}%
\pgfpathlineto{\pgfqpoint{1.889712in}{0.906997in}}%
\pgfpathlineto{\pgfqpoint{1.889632in}{0.918440in}}%
\pgfpathlineto{\pgfqpoint{1.899200in}{0.918414in}}%
\pgfpathlineto{\pgfqpoint{1.899149in}{0.923195in}}%
\pgfpathlineto{\pgfqpoint{1.910553in}{0.923358in}}%
\pgfpathlineto{\pgfqpoint{1.910330in}{0.935711in}}%
\pgfpathlineto{\pgfqpoint{1.916049in}{0.935756in}}%
\pgfpathlineto{\pgfqpoint{1.915999in}{0.941499in}}%
\pgfpathlineto{\pgfqpoint{1.921721in}{0.941554in}}%
\pgfpathlineto{\pgfqpoint{1.921677in}{0.946416in}}%
\pgfpathlineto{\pgfqpoint{1.929258in}{0.940159in}}%
\pgfpathlineto{\pgfqpoint{1.927752in}{0.935378in}}%
\pgfpathlineto{\pgfqpoint{1.935018in}{0.933586in}}%
\pgfpathlineto{\pgfqpoint{1.940723in}{0.927923in}}%
\pgfpathlineto{\pgfqpoint{1.939696in}{0.925677in}}%
\pgfpathlineto{\pgfqpoint{1.944264in}{0.919360in}}%
\pgfpathlineto{\pgfqpoint{1.944456in}{0.907639in}}%
\pgfpathclose%
\pgfusepath{fill}%
\end{pgfscope}%
\begin{pgfscope}%
\pgfpathrectangle{\pgfqpoint{0.100000in}{0.100000in}}{\pgfqpoint{3.007045in}{1.925000in}}%
\pgfusepath{clip}%
\pgfsetbuttcap%
\pgfsetmiterjoin%
\definecolor{currentfill}{rgb}{0.623760,0.847428,0.601584}%
\pgfsetfillcolor{currentfill}%
\pgfsetlinewidth{0.000000pt}%
\definecolor{currentstroke}{rgb}{0.000000,0.000000,0.000000}%
\pgfsetstrokecolor{currentstroke}%
\pgfsetstrokeopacity{0.000000}%
\pgfsetdash{}{0pt}%
\pgfpathmoveto{\pgfqpoint{2.482424in}{0.851583in}}%
\pgfpathlineto{\pgfqpoint{2.474703in}{0.848311in}}%
\pgfpathlineto{\pgfqpoint{2.477899in}{0.842105in}}%
\pgfpathlineto{\pgfqpoint{2.475134in}{0.834677in}}%
\pgfpathlineto{\pgfqpoint{2.480732in}{0.829446in}}%
\pgfpathlineto{\pgfqpoint{2.479069in}{0.826299in}}%
\pgfpathlineto{\pgfqpoint{2.473536in}{0.831913in}}%
\pgfpathlineto{\pgfqpoint{2.468037in}{0.838774in}}%
\pgfpathlineto{\pgfqpoint{2.452246in}{0.846939in}}%
\pgfpathlineto{\pgfqpoint{2.449961in}{0.851491in}}%
\pgfpathlineto{\pgfqpoint{2.442202in}{0.858557in}}%
\pgfpathlineto{\pgfqpoint{2.439934in}{0.864954in}}%
\pgfpathlineto{\pgfqpoint{2.429221in}{0.879173in}}%
\pgfpathlineto{\pgfqpoint{2.423854in}{0.877561in}}%
\pgfpathlineto{\pgfqpoint{2.420077in}{0.879039in}}%
\pgfpathlineto{\pgfqpoint{2.413604in}{0.885230in}}%
\pgfpathlineto{\pgfqpoint{2.410336in}{0.885056in}}%
\pgfpathlineto{\pgfqpoint{2.403995in}{0.889257in}}%
\pgfpathlineto{\pgfqpoint{2.402990in}{0.891059in}}%
\pgfpathlineto{\pgfqpoint{2.403903in}{0.896059in}}%
\pgfpathlineto{\pgfqpoint{2.408265in}{0.902647in}}%
\pgfpathlineto{\pgfqpoint{2.413500in}{0.907511in}}%
\pgfpathlineto{\pgfqpoint{2.413670in}{0.911416in}}%
\pgfpathlineto{\pgfqpoint{2.418790in}{0.913875in}}%
\pgfpathlineto{\pgfqpoint{2.435961in}{0.922177in}}%
\pgfpathlineto{\pgfqpoint{2.452389in}{0.929337in}}%
\pgfpathlineto{\pgfqpoint{2.459680in}{0.930712in}}%
\pgfpathlineto{\pgfqpoint{2.462329in}{0.908061in}}%
\pgfpathlineto{\pgfqpoint{2.476040in}{0.897437in}}%
\pgfpathlineto{\pgfqpoint{2.481602in}{0.894314in}}%
\pgfpathlineto{\pgfqpoint{2.496666in}{0.891862in}}%
\pgfpathlineto{\pgfqpoint{2.490797in}{0.879764in}}%
\pgfpathlineto{\pgfqpoint{2.485317in}{0.874629in}}%
\pgfpathlineto{\pgfqpoint{2.484386in}{0.869099in}}%
\pgfpathlineto{\pgfqpoint{2.486659in}{0.864555in}}%
\pgfpathlineto{\pgfqpoint{2.482424in}{0.851583in}}%
\pgfpathclose%
\pgfusepath{fill}%
\end{pgfscope}%
\begin{pgfscope}%
\pgfpathrectangle{\pgfqpoint{0.100000in}{0.100000in}}{\pgfqpoint{3.007045in}{1.925000in}}%
\pgfusepath{clip}%
\pgfsetbuttcap%
\pgfsetmiterjoin%
\definecolor{currentfill}{rgb}{0.676371,0.869927,0.651657}%
\pgfsetfillcolor{currentfill}%
\pgfsetlinewidth{0.000000pt}%
\definecolor{currentstroke}{rgb}{0.000000,0.000000,0.000000}%
\pgfsetstrokecolor{currentstroke}%
\pgfsetstrokeopacity{0.000000}%
\pgfsetdash{}{0pt}%
\pgfpathmoveto{\pgfqpoint{2.253561in}{1.079316in}}%
\pgfpathlineto{\pgfqpoint{2.249127in}{1.081944in}}%
\pgfpathlineto{\pgfqpoint{2.246327in}{1.087565in}}%
\pgfpathlineto{\pgfqpoint{2.241377in}{1.090992in}}%
\pgfpathlineto{\pgfqpoint{2.229804in}{1.090057in}}%
\pgfpathlineto{\pgfqpoint{2.224931in}{1.092417in}}%
\pgfpathlineto{\pgfqpoint{2.223926in}{1.097599in}}%
\pgfpathlineto{\pgfqpoint{2.219675in}{1.101261in}}%
\pgfpathlineto{\pgfqpoint{2.217176in}{1.095188in}}%
\pgfpathlineto{\pgfqpoint{2.214133in}{1.095845in}}%
\pgfpathlineto{\pgfqpoint{2.213695in}{1.101544in}}%
\pgfpathlineto{\pgfqpoint{2.207958in}{1.101054in}}%
\pgfpathlineto{\pgfqpoint{2.207678in}{1.104917in}}%
\pgfpathlineto{\pgfqpoint{2.202065in}{1.104184in}}%
\pgfpathlineto{\pgfqpoint{2.201173in}{1.112813in}}%
\pgfpathlineto{\pgfqpoint{2.220327in}{1.114564in}}%
\pgfpathlineto{\pgfqpoint{2.218133in}{1.136791in}}%
\pgfpathlineto{\pgfqpoint{2.219718in}{1.138972in}}%
\pgfpathlineto{\pgfqpoint{2.235786in}{1.140573in}}%
\pgfpathlineto{\pgfqpoint{2.239682in}{1.138991in}}%
\pgfpathlineto{\pgfqpoint{2.240875in}{1.128586in}}%
\pgfpathlineto{\pgfqpoint{2.245396in}{1.130991in}}%
\pgfpathlineto{\pgfqpoint{2.256834in}{1.132262in}}%
\pgfpathlineto{\pgfqpoint{2.264234in}{1.131721in}}%
\pgfpathlineto{\pgfqpoint{2.264410in}{1.127604in}}%
\pgfpathlineto{\pgfqpoint{2.272501in}{1.126660in}}%
\pgfpathlineto{\pgfqpoint{2.277544in}{1.133106in}}%
\pgfpathlineto{\pgfqpoint{2.282277in}{1.134362in}}%
\pgfpathlineto{\pgfqpoint{2.285748in}{1.132441in}}%
\pgfpathlineto{\pgfqpoint{2.287382in}{1.127701in}}%
\pgfpathlineto{\pgfqpoint{2.292494in}{1.120777in}}%
\pgfpathlineto{\pgfqpoint{2.288467in}{1.113713in}}%
\pgfpathlineto{\pgfqpoint{2.288667in}{1.106429in}}%
\pgfpathlineto{\pgfqpoint{2.287143in}{1.100325in}}%
\pgfpathlineto{\pgfqpoint{2.281936in}{1.092901in}}%
\pgfpathlineto{\pgfqpoint{2.270164in}{1.089523in}}%
\pgfpathlineto{\pgfqpoint{2.265212in}{1.092380in}}%
\pgfpathlineto{\pgfqpoint{2.258362in}{1.081298in}}%
\pgfpathlineto{\pgfqpoint{2.253561in}{1.079316in}}%
\pgfpathclose%
\pgfusepath{fill}%
\end{pgfscope}%
\begin{pgfscope}%
\pgfpathrectangle{\pgfqpoint{0.100000in}{0.100000in}}{\pgfqpoint{3.007045in}{1.925000in}}%
\pgfusepath{clip}%
\pgfsetbuttcap%
\pgfsetmiterjoin%
\definecolor{currentfill}{rgb}{0.401569,0.742484,0.436601}%
\pgfsetfillcolor{currentfill}%
\pgfsetlinewidth{0.000000pt}%
\definecolor{currentstroke}{rgb}{0.000000,0.000000,0.000000}%
\pgfsetstrokecolor{currentstroke}%
\pgfsetstrokeopacity{0.000000}%
\pgfsetdash{}{0pt}%
\pgfpathmoveto{\pgfqpoint{1.610638in}{1.328045in}}%
\pgfpathlineto{\pgfqpoint{1.610215in}{1.316599in}}%
\pgfpathlineto{\pgfqpoint{1.587433in}{1.317426in}}%
\pgfpathlineto{\pgfqpoint{1.588274in}{1.340363in}}%
\pgfpathlineto{\pgfqpoint{1.565129in}{1.341345in}}%
\pgfpathlineto{\pgfqpoint{1.566978in}{1.401393in}}%
\pgfpathlineto{\pgfqpoint{1.603438in}{1.399865in}}%
\pgfpathlineto{\pgfqpoint{1.604815in}{1.396442in}}%
\pgfpathlineto{\pgfqpoint{1.612271in}{1.391873in}}%
\pgfpathlineto{\pgfqpoint{1.611668in}{1.362462in}}%
\pgfpathlineto{\pgfqpoint{1.610638in}{1.328045in}}%
\pgfpathclose%
\pgfusepath{fill}%
\end{pgfscope}%
\begin{pgfscope}%
\pgfpathrectangle{\pgfqpoint{0.100000in}{0.100000in}}{\pgfqpoint{3.007045in}{1.925000in}}%
\pgfusepath{clip}%
\pgfsetbuttcap%
\pgfsetmiterjoin%
\definecolor{currentfill}{rgb}{0.709097,0.883706,0.683522}%
\pgfsetfillcolor{currentfill}%
\pgfsetlinewidth{0.000000pt}%
\definecolor{currentstroke}{rgb}{0.000000,0.000000,0.000000}%
\pgfsetstrokecolor{currentstroke}%
\pgfsetstrokeopacity{0.000000}%
\pgfsetdash{}{0pt}%
\pgfpathmoveto{\pgfqpoint{2.283146in}{1.250837in}}%
\pgfpathlineto{\pgfqpoint{2.283471in}{1.248041in}}%
\pgfpathlineto{\pgfqpoint{2.262780in}{1.245532in}}%
\pgfpathlineto{\pgfqpoint{2.261811in}{1.262682in}}%
\pgfpathlineto{\pgfqpoint{2.249581in}{1.261349in}}%
\pgfpathlineto{\pgfqpoint{2.248890in}{1.267034in}}%
\pgfpathlineto{\pgfqpoint{2.239459in}{1.266008in}}%
\pgfpathlineto{\pgfqpoint{2.236831in}{1.288875in}}%
\pgfpathlineto{\pgfqpoint{2.234886in}{1.288646in}}%
\pgfpathlineto{\pgfqpoint{2.233530in}{1.300204in}}%
\pgfpathlineto{\pgfqpoint{2.234453in}{1.308012in}}%
\pgfpathlineto{\pgfqpoint{2.240301in}{1.306723in}}%
\pgfpathlineto{\pgfqpoint{2.257315in}{1.308409in}}%
\pgfpathlineto{\pgfqpoint{2.253575in}{1.340858in}}%
\pgfpathlineto{\pgfqpoint{2.271732in}{1.342879in}}%
\pgfpathlineto{\pgfqpoint{2.273138in}{1.338797in}}%
\pgfpathlineto{\pgfqpoint{2.278491in}{1.292509in}}%
\pgfpathlineto{\pgfqpoint{2.283146in}{1.250837in}}%
\pgfpathclose%
\pgfusepath{fill}%
\end{pgfscope}%
\begin{pgfscope}%
\pgfpathrectangle{\pgfqpoint{0.100000in}{0.100000in}}{\pgfqpoint{3.007045in}{1.925000in}}%
\pgfusepath{clip}%
\pgfsetbuttcap%
\pgfsetmiterjoin%
\definecolor{currentfill}{rgb}{0.557324,0.816424,0.546959}%
\pgfsetfillcolor{currentfill}%
\pgfsetlinewidth{0.000000pt}%
\definecolor{currentstroke}{rgb}{0.000000,0.000000,0.000000}%
\pgfsetstrokecolor{currentstroke}%
\pgfsetstrokeopacity{0.000000}%
\pgfsetdash{}{0pt}%
\pgfpathmoveto{\pgfqpoint{2.380873in}{1.017870in}}%
\pgfpathlineto{\pgfqpoint{2.378204in}{1.022084in}}%
\pgfpathlineto{\pgfqpoint{2.378235in}{1.029769in}}%
\pgfpathlineto{\pgfqpoint{2.376071in}{1.035433in}}%
\pgfpathlineto{\pgfqpoint{2.371954in}{1.050139in}}%
\pgfpathlineto{\pgfqpoint{2.372484in}{1.056245in}}%
\pgfpathlineto{\pgfqpoint{2.370868in}{1.060399in}}%
\pgfpathlineto{\pgfqpoint{2.375741in}{1.064118in}}%
\pgfpathlineto{\pgfqpoint{2.378157in}{1.062296in}}%
\pgfpathlineto{\pgfqpoint{2.384647in}{1.067115in}}%
\pgfpathlineto{\pgfqpoint{2.392156in}{1.068534in}}%
\pgfpathlineto{\pgfqpoint{2.392967in}{1.076431in}}%
\pgfpathlineto{\pgfqpoint{2.400437in}{1.076137in}}%
\pgfpathlineto{\pgfqpoint{2.410452in}{1.072460in}}%
\pgfpathlineto{\pgfqpoint{2.412270in}{1.064935in}}%
\pgfpathlineto{\pgfqpoint{2.417226in}{1.060057in}}%
\pgfpathlineto{\pgfqpoint{2.423037in}{1.058784in}}%
\pgfpathlineto{\pgfqpoint{2.415558in}{1.052709in}}%
\pgfpathlineto{\pgfqpoint{2.415182in}{1.047525in}}%
\pgfpathlineto{\pgfqpoint{2.409125in}{1.041720in}}%
\pgfpathlineto{\pgfqpoint{2.409387in}{1.036508in}}%
\pgfpathlineto{\pgfqpoint{2.399471in}{1.032629in}}%
\pgfpathlineto{\pgfqpoint{2.397238in}{1.025171in}}%
\pgfpathlineto{\pgfqpoint{2.380873in}{1.017870in}}%
\pgfpathclose%
\pgfusepath{fill}%
\end{pgfscope}%
\begin{pgfscope}%
\pgfpathrectangle{\pgfqpoint{0.100000in}{0.100000in}}{\pgfqpoint{3.007045in}{1.925000in}}%
\pgfusepath{clip}%
\pgfsetbuttcap%
\pgfsetmiterjoin%
\definecolor{currentfill}{rgb}{0.676371,0.869927,0.651657}%
\pgfsetfillcolor{currentfill}%
\pgfsetlinewidth{0.000000pt}%
\definecolor{currentstroke}{rgb}{0.000000,0.000000,0.000000}%
\pgfsetstrokecolor{currentstroke}%
\pgfsetstrokeopacity{0.000000}%
\pgfsetdash{}{0pt}%
\pgfpathmoveto{\pgfqpoint{2.508172in}{0.896146in}}%
\pgfpathlineto{\pgfqpoint{2.508933in}{0.891663in}}%
\pgfpathlineto{\pgfqpoint{2.502098in}{0.888878in}}%
\pgfpathlineto{\pgfqpoint{2.496666in}{0.891862in}}%
\pgfpathlineto{\pgfqpoint{2.481602in}{0.894314in}}%
\pgfpathlineto{\pgfqpoint{2.476040in}{0.897437in}}%
\pgfpathlineto{\pgfqpoint{2.462329in}{0.908061in}}%
\pgfpathlineto{\pgfqpoint{2.459680in}{0.930712in}}%
\pgfpathlineto{\pgfqpoint{2.452389in}{0.929337in}}%
\pgfpathlineto{\pgfqpoint{2.451924in}{0.935651in}}%
\pgfpathlineto{\pgfqpoint{2.455447in}{0.943242in}}%
\pgfpathlineto{\pgfqpoint{2.460723in}{0.945065in}}%
\pgfpathlineto{\pgfqpoint{2.472311in}{0.936614in}}%
\pgfpathlineto{\pgfqpoint{2.472948in}{0.932025in}}%
\pgfpathlineto{\pgfqpoint{2.505332in}{0.935315in}}%
\pgfpathlineto{\pgfqpoint{2.504209in}{0.927233in}}%
\pgfpathlineto{\pgfqpoint{2.500183in}{0.925858in}}%
\pgfpathlineto{\pgfqpoint{2.503825in}{0.913860in}}%
\pgfpathlineto{\pgfqpoint{2.505466in}{0.902631in}}%
\pgfpathlineto{\pgfqpoint{2.508172in}{0.896146in}}%
\pgfpathclose%
\pgfusepath{fill}%
\end{pgfscope}%
\begin{pgfscope}%
\pgfpathrectangle{\pgfqpoint{0.100000in}{0.100000in}}{\pgfqpoint{3.007045in}{1.925000in}}%
\pgfusepath{clip}%
\pgfsetbuttcap%
\pgfsetmiterjoin%
\definecolor{currentfill}{rgb}{0.886505,0.956171,0.866128}%
\pgfsetfillcolor{currentfill}%
\pgfsetlinewidth{0.000000pt}%
\definecolor{currentstroke}{rgb}{0.000000,0.000000,0.000000}%
\pgfsetstrokecolor{currentstroke}%
\pgfsetstrokeopacity{0.000000}%
\pgfsetdash{}{0pt}%
\pgfpathmoveto{\pgfqpoint{0.928770in}{0.169774in}}%
\pgfpathlineto{\pgfqpoint{0.929087in}{0.172237in}}%
\pgfpathlineto{\pgfqpoint{0.930780in}{0.170565in}}%
\pgfpathlineto{\pgfqpoint{0.928770in}{0.169774in}}%
\pgfpathclose%
\pgfusepath{fill}%
\end{pgfscope}%
\begin{pgfscope}%
\pgfpathrectangle{\pgfqpoint{0.100000in}{0.100000in}}{\pgfqpoint{3.007045in}{1.925000in}}%
\pgfusepath{clip}%
\pgfsetbuttcap%
\pgfsetmiterjoin%
\definecolor{currentfill}{rgb}{0.886505,0.956171,0.866128}%
\pgfsetfillcolor{currentfill}%
\pgfsetlinewidth{0.000000pt}%
\definecolor{currentstroke}{rgb}{0.000000,0.000000,0.000000}%
\pgfsetstrokecolor{currentstroke}%
\pgfsetstrokeopacity{0.000000}%
\pgfsetdash{}{0pt}%
\pgfpathmoveto{\pgfqpoint{0.927632in}{0.179344in}}%
\pgfpathlineto{\pgfqpoint{0.926956in}{0.181069in}}%
\pgfpathlineto{\pgfqpoint{0.926534in}{0.189086in}}%
\pgfpathlineto{\pgfqpoint{0.927077in}{0.189993in}}%
\pgfpathlineto{\pgfqpoint{0.926777in}{0.193425in}}%
\pgfpathlineto{\pgfqpoint{0.927516in}{0.194491in}}%
\pgfpathlineto{\pgfqpoint{0.926137in}{0.195868in}}%
\pgfpathlineto{\pgfqpoint{0.925991in}{0.197780in}}%
\pgfpathlineto{\pgfqpoint{0.928814in}{0.198830in}}%
\pgfpathlineto{\pgfqpoint{0.929984in}{0.200875in}}%
\pgfpathlineto{\pgfqpoint{0.928490in}{0.201725in}}%
\pgfpathlineto{\pgfqpoint{0.928312in}{0.203305in}}%
\pgfpathlineto{\pgfqpoint{0.929750in}{0.205372in}}%
\pgfpathlineto{\pgfqpoint{0.928020in}{0.209653in}}%
\pgfpathlineto{\pgfqpoint{0.932348in}{0.213098in}}%
\pgfpathlineto{\pgfqpoint{0.934167in}{0.209712in}}%
\pgfpathlineto{\pgfqpoint{0.937918in}{0.207074in}}%
\pgfpathlineto{\pgfqpoint{0.935283in}{0.206082in}}%
\pgfpathlineto{\pgfqpoint{0.936138in}{0.203890in}}%
\pgfpathlineto{\pgfqpoint{0.935915in}{0.201464in}}%
\pgfpathlineto{\pgfqpoint{0.934812in}{0.198205in}}%
\pgfpathlineto{\pgfqpoint{0.934411in}{0.195163in}}%
\pgfpathlineto{\pgfqpoint{0.933099in}{0.192388in}}%
\pgfpathlineto{\pgfqpoint{0.933499in}{0.191353in}}%
\pgfpathlineto{\pgfqpoint{0.932518in}{0.188498in}}%
\pgfpathlineto{\pgfqpoint{0.930488in}{0.187744in}}%
\pgfpathlineto{\pgfqpoint{0.931046in}{0.186350in}}%
\pgfpathlineto{\pgfqpoint{0.929012in}{0.184009in}}%
\pgfpathlineto{\pgfqpoint{0.927632in}{0.179344in}}%
\pgfpathclose%
\pgfusepath{fill}%
\end{pgfscope}%
\begin{pgfscope}%
\pgfpathrectangle{\pgfqpoint{0.100000in}{0.100000in}}{\pgfqpoint{3.007045in}{1.925000in}}%
\pgfusepath{clip}%
\pgfsetbuttcap%
\pgfsetmiterjoin%
\definecolor{currentfill}{rgb}{0.886505,0.956171,0.866128}%
\pgfsetfillcolor{currentfill}%
\pgfsetlinewidth{0.000000pt}%
\definecolor{currentstroke}{rgb}{0.000000,0.000000,0.000000}%
\pgfsetstrokecolor{currentstroke}%
\pgfsetstrokeopacity{0.000000}%
\pgfsetdash{}{0pt}%
\pgfpathmoveto{\pgfqpoint{0.921995in}{0.203350in}}%
\pgfpathlineto{\pgfqpoint{0.924523in}{0.205630in}}%
\pgfpathlineto{\pgfqpoint{0.924613in}{0.208652in}}%
\pgfpathlineto{\pgfqpoint{0.925424in}{0.210159in}}%
\pgfpathlineto{\pgfqpoint{0.927073in}{0.209363in}}%
\pgfpathlineto{\pgfqpoint{0.926551in}{0.205362in}}%
\pgfpathlineto{\pgfqpoint{0.925876in}{0.203745in}}%
\pgfpathlineto{\pgfqpoint{0.924259in}{0.202862in}}%
\pgfpathlineto{\pgfqpoint{0.921995in}{0.203350in}}%
\pgfpathclose%
\pgfusepath{fill}%
\end{pgfscope}%
\begin{pgfscope}%
\pgfpathrectangle{\pgfqpoint{0.100000in}{0.100000in}}{\pgfqpoint{3.007045in}{1.925000in}}%
\pgfusepath{clip}%
\pgfsetbuttcap%
\pgfsetmiterjoin%
\definecolor{currentfill}{rgb}{0.886505,0.956171,0.866128}%
\pgfsetfillcolor{currentfill}%
\pgfsetlinewidth{0.000000pt}%
\definecolor{currentstroke}{rgb}{0.000000,0.000000,0.000000}%
\pgfsetstrokecolor{currentstroke}%
\pgfsetstrokeopacity{0.000000}%
\pgfsetdash{}{0pt}%
\pgfpathmoveto{\pgfqpoint{0.924230in}{0.224566in}}%
\pgfpathlineto{\pgfqpoint{0.928775in}{0.222468in}}%
\pgfpathlineto{\pgfqpoint{0.931898in}{0.225478in}}%
\pgfpathlineto{\pgfqpoint{0.933152in}{0.222753in}}%
\pgfpathlineto{\pgfqpoint{0.934778in}{0.220682in}}%
\pgfpathlineto{\pgfqpoint{0.936886in}{0.216547in}}%
\pgfpathlineto{\pgfqpoint{0.939636in}{0.215477in}}%
\pgfpathlineto{\pgfqpoint{0.940303in}{0.214612in}}%
\pgfpathlineto{\pgfqpoint{0.938629in}{0.208282in}}%
\pgfpathlineto{\pgfqpoint{0.936276in}{0.208683in}}%
\pgfpathlineto{\pgfqpoint{0.931068in}{0.218130in}}%
\pgfpathlineto{\pgfqpoint{0.930637in}{0.216360in}}%
\pgfpathlineto{\pgfqpoint{0.931418in}{0.214692in}}%
\pgfpathlineto{\pgfqpoint{0.930183in}{0.212277in}}%
\pgfpathlineto{\pgfqpoint{0.927651in}{0.210591in}}%
\pgfpathlineto{\pgfqpoint{0.926290in}{0.211782in}}%
\pgfpathlineto{\pgfqpoint{0.925901in}{0.214793in}}%
\pgfpathlineto{\pgfqpoint{0.926598in}{0.218053in}}%
\pgfpathlineto{\pgfqpoint{0.924230in}{0.224566in}}%
\pgfpathclose%
\pgfusepath{fill}%
\end{pgfscope}%
\begin{pgfscope}%
\pgfpathrectangle{\pgfqpoint{0.100000in}{0.100000in}}{\pgfqpoint{3.007045in}{1.925000in}}%
\pgfusepath{clip}%
\pgfsetbuttcap%
\pgfsetmiterjoin%
\definecolor{currentfill}{rgb}{0.886505,0.956171,0.866128}%
\pgfsetfillcolor{currentfill}%
\pgfsetlinewidth{0.000000pt}%
\definecolor{currentstroke}{rgb}{0.000000,0.000000,0.000000}%
\pgfsetstrokecolor{currentstroke}%
\pgfsetstrokeopacity{0.000000}%
\pgfsetdash{}{0pt}%
\pgfpathmoveto{\pgfqpoint{0.949543in}{0.170923in}}%
\pgfpathlineto{\pgfqpoint{0.948341in}{0.174468in}}%
\pgfpathlineto{\pgfqpoint{0.951022in}{0.176053in}}%
\pgfpathlineto{\pgfqpoint{0.953034in}{0.175105in}}%
\pgfpathlineto{\pgfqpoint{0.953826in}{0.173319in}}%
\pgfpathlineto{\pgfqpoint{0.951936in}{0.170704in}}%
\pgfpathlineto{\pgfqpoint{0.949543in}{0.170923in}}%
\pgfpathclose%
\pgfusepath{fill}%
\end{pgfscope}%
\begin{pgfscope}%
\pgfpathrectangle{\pgfqpoint{0.100000in}{0.100000in}}{\pgfqpoint{3.007045in}{1.925000in}}%
\pgfusepath{clip}%
\pgfsetbuttcap%
\pgfsetmiterjoin%
\definecolor{currentfill}{rgb}{0.886505,0.956171,0.866128}%
\pgfsetfillcolor{currentfill}%
\pgfsetlinewidth{0.000000pt}%
\definecolor{currentstroke}{rgb}{0.000000,0.000000,0.000000}%
\pgfsetstrokecolor{currentstroke}%
\pgfsetstrokeopacity{0.000000}%
\pgfsetdash{}{0pt}%
\pgfpathmoveto{\pgfqpoint{0.941866in}{0.183097in}}%
\pgfpathlineto{\pgfqpoint{0.939091in}{0.182574in}}%
\pgfpathlineto{\pgfqpoint{0.938900in}{0.178993in}}%
\pgfpathlineto{\pgfqpoint{0.936985in}{0.178377in}}%
\pgfpathlineto{\pgfqpoint{0.937239in}{0.175986in}}%
\pgfpathlineto{\pgfqpoint{0.935359in}{0.173870in}}%
\pgfpathlineto{\pgfqpoint{0.934641in}{0.176663in}}%
\pgfpathlineto{\pgfqpoint{0.933054in}{0.175528in}}%
\pgfpathlineto{\pgfqpoint{0.932146in}{0.176239in}}%
\pgfpathlineto{\pgfqpoint{0.933479in}{0.179062in}}%
\pgfpathlineto{\pgfqpoint{0.935593in}{0.180734in}}%
\pgfpathlineto{\pgfqpoint{0.937053in}{0.180678in}}%
\pgfpathlineto{\pgfqpoint{0.938116in}{0.183219in}}%
\pgfpathlineto{\pgfqpoint{0.936870in}{0.183669in}}%
\pgfpathlineto{\pgfqpoint{0.935994in}{0.185420in}}%
\pgfpathlineto{\pgfqpoint{0.936562in}{0.187485in}}%
\pgfpathlineto{\pgfqpoint{0.936865in}{0.191439in}}%
\pgfpathlineto{\pgfqpoint{0.939368in}{0.192607in}}%
\pgfpathlineto{\pgfqpoint{0.942659in}{0.187128in}}%
\pgfpathlineto{\pgfqpoint{0.943420in}{0.184913in}}%
\pgfpathlineto{\pgfqpoint{0.941866in}{0.183097in}}%
\pgfpathclose%
\pgfusepath{fill}%
\end{pgfscope}%
\begin{pgfscope}%
\pgfpathrectangle{\pgfqpoint{0.100000in}{0.100000in}}{\pgfqpoint{3.007045in}{1.925000in}}%
\pgfusepath{clip}%
\pgfsetbuttcap%
\pgfsetmiterjoin%
\definecolor{currentfill}{rgb}{0.886505,0.956171,0.866128}%
\pgfsetfillcolor{currentfill}%
\pgfsetlinewidth{0.000000pt}%
\definecolor{currentstroke}{rgb}{0.000000,0.000000,0.000000}%
\pgfsetstrokecolor{currentstroke}%
\pgfsetstrokeopacity{0.000000}%
\pgfsetdash{}{0pt}%
\pgfpathmoveto{\pgfqpoint{0.951030in}{0.194631in}}%
\pgfpathlineto{\pgfqpoint{0.954245in}{0.194809in}}%
\pgfpathlineto{\pgfqpoint{0.955692in}{0.192812in}}%
\pgfpathlineto{\pgfqpoint{0.957247in}{0.193413in}}%
\pgfpathlineto{\pgfqpoint{0.958655in}{0.195219in}}%
\pgfpathlineto{\pgfqpoint{0.963493in}{0.195565in}}%
\pgfpathlineto{\pgfqpoint{0.964721in}{0.193135in}}%
\pgfpathlineto{\pgfqpoint{0.966052in}{0.192399in}}%
\pgfpathlineto{\pgfqpoint{0.966135in}{0.189796in}}%
\pgfpathlineto{\pgfqpoint{0.963633in}{0.186915in}}%
\pgfpathlineto{\pgfqpoint{0.966486in}{0.184667in}}%
\pgfpathlineto{\pgfqpoint{0.964468in}{0.180678in}}%
\pgfpathlineto{\pgfqpoint{0.966318in}{0.178588in}}%
\pgfpathlineto{\pgfqpoint{0.965355in}{0.173338in}}%
\pgfpathlineto{\pgfqpoint{0.967092in}{0.172775in}}%
\pgfpathlineto{\pgfqpoint{0.969300in}{0.170288in}}%
\pgfpathlineto{\pgfqpoint{0.972690in}{0.165137in}}%
\pgfpathlineto{\pgfqpoint{0.972584in}{0.162951in}}%
\pgfpathlineto{\pgfqpoint{0.970864in}{0.160909in}}%
\pgfpathlineto{\pgfqpoint{0.968598in}{0.161733in}}%
\pgfpathlineto{\pgfqpoint{0.965820in}{0.161252in}}%
\pgfpathlineto{\pgfqpoint{0.966014in}{0.159668in}}%
\pgfpathlineto{\pgfqpoint{0.964181in}{0.158832in}}%
\pgfpathlineto{\pgfqpoint{0.962912in}{0.161570in}}%
\pgfpathlineto{\pgfqpoint{0.960550in}{0.161270in}}%
\pgfpathlineto{\pgfqpoint{0.958681in}{0.160262in}}%
\pgfpathlineto{\pgfqpoint{0.958436in}{0.159148in}}%
\pgfpathlineto{\pgfqpoint{0.956199in}{0.158645in}}%
\pgfpathlineto{\pgfqpoint{0.956678in}{0.156700in}}%
\pgfpathlineto{\pgfqpoint{0.955416in}{0.154694in}}%
\pgfpathlineto{\pgfqpoint{0.953820in}{0.154455in}}%
\pgfpathlineto{\pgfqpoint{0.952217in}{0.157810in}}%
\pgfpathlineto{\pgfqpoint{0.955289in}{0.158119in}}%
\pgfpathlineto{\pgfqpoint{0.956244in}{0.159613in}}%
\pgfpathlineto{\pgfqpoint{0.958348in}{0.160762in}}%
\pgfpathlineto{\pgfqpoint{0.959317in}{0.163914in}}%
\pgfpathlineto{\pgfqpoint{0.960954in}{0.164908in}}%
\pgfpathlineto{\pgfqpoint{0.959218in}{0.166639in}}%
\pgfpathlineto{\pgfqpoint{0.958611in}{0.164391in}}%
\pgfpathlineto{\pgfqpoint{0.955800in}{0.164056in}}%
\pgfpathlineto{\pgfqpoint{0.954820in}{0.161400in}}%
\pgfpathlineto{\pgfqpoint{0.952564in}{0.162138in}}%
\pgfpathlineto{\pgfqpoint{0.952336in}{0.166328in}}%
\pgfpathlineto{\pgfqpoint{0.950396in}{0.166275in}}%
\pgfpathlineto{\pgfqpoint{0.951336in}{0.169331in}}%
\pgfpathlineto{\pgfqpoint{0.955934in}{0.171121in}}%
\pgfpathlineto{\pgfqpoint{0.955518in}{0.168806in}}%
\pgfpathlineto{\pgfqpoint{0.956516in}{0.167155in}}%
\pgfpathlineto{\pgfqpoint{0.957594in}{0.173312in}}%
\pgfpathlineto{\pgfqpoint{0.958602in}{0.175079in}}%
\pgfpathlineto{\pgfqpoint{0.961488in}{0.177541in}}%
\pgfpathlineto{\pgfqpoint{0.957917in}{0.179285in}}%
\pgfpathlineto{\pgfqpoint{0.958031in}{0.180854in}}%
\pgfpathlineto{\pgfqpoint{0.956711in}{0.183035in}}%
\pgfpathlineto{\pgfqpoint{0.956083in}{0.186705in}}%
\pgfpathlineto{\pgfqpoint{0.957478in}{0.186958in}}%
\pgfpathlineto{\pgfqpoint{0.957475in}{0.188952in}}%
\pgfpathlineto{\pgfqpoint{0.955586in}{0.189502in}}%
\pgfpathlineto{\pgfqpoint{0.954193in}{0.191220in}}%
\pgfpathlineto{\pgfqpoint{0.953928in}{0.193026in}}%
\pgfpathlineto{\pgfqpoint{0.952297in}{0.192656in}}%
\pgfpathlineto{\pgfqpoint{0.951030in}{0.194631in}}%
\pgfpathclose%
\pgfusepath{fill}%
\end{pgfscope}%
\begin{pgfscope}%
\pgfpathrectangle{\pgfqpoint{0.100000in}{0.100000in}}{\pgfqpoint{3.007045in}{1.925000in}}%
\pgfusepath{clip}%
\pgfsetbuttcap%
\pgfsetmiterjoin%
\definecolor{currentfill}{rgb}{0.886505,0.956171,0.866128}%
\pgfsetfillcolor{currentfill}%
\pgfsetlinewidth{0.000000pt}%
\definecolor{currentstroke}{rgb}{0.000000,0.000000,0.000000}%
\pgfsetstrokecolor{currentstroke}%
\pgfsetstrokeopacity{0.000000}%
\pgfsetdash{}{0pt}%
\pgfpathmoveto{\pgfqpoint{0.942241in}{0.179713in}}%
\pgfpathlineto{\pgfqpoint{0.944515in}{0.186094in}}%
\pgfpathlineto{\pgfqpoint{0.944156in}{0.190021in}}%
\pgfpathlineto{\pgfqpoint{0.944652in}{0.190685in}}%
\pgfpathlineto{\pgfqpoint{0.943856in}{0.194420in}}%
\pgfpathlineto{\pgfqpoint{0.945632in}{0.194990in}}%
\pgfpathlineto{\pgfqpoint{0.951121in}{0.190206in}}%
\pgfpathlineto{\pgfqpoint{0.954271in}{0.188027in}}%
\pgfpathlineto{\pgfqpoint{0.954701in}{0.185800in}}%
\pgfpathlineto{\pgfqpoint{0.953794in}{0.184014in}}%
\pgfpathlineto{\pgfqpoint{0.955507in}{0.182225in}}%
\pgfpathlineto{\pgfqpoint{0.956828in}{0.176810in}}%
\pgfpathlineto{\pgfqpoint{0.953089in}{0.176262in}}%
\pgfpathlineto{\pgfqpoint{0.951273in}{0.177255in}}%
\pgfpathlineto{\pgfqpoint{0.953068in}{0.180485in}}%
\pgfpathlineto{\pgfqpoint{0.951111in}{0.180223in}}%
\pgfpathlineto{\pgfqpoint{0.949807in}{0.178518in}}%
\pgfpathlineto{\pgfqpoint{0.947328in}{0.177601in}}%
\pgfpathlineto{\pgfqpoint{0.945308in}{0.179613in}}%
\pgfpathlineto{\pgfqpoint{0.944254in}{0.178896in}}%
\pgfpathlineto{\pgfqpoint{0.942241in}{0.179713in}}%
\pgfpathclose%
\pgfusepath{fill}%
\end{pgfscope}%
\begin{pgfscope}%
\pgfpathrectangle{\pgfqpoint{0.100000in}{0.100000in}}{\pgfqpoint{3.007045in}{1.925000in}}%
\pgfusepath{clip}%
\pgfsetbuttcap%
\pgfsetmiterjoin%
\definecolor{currentfill}{rgb}{0.288627,0.687120,0.381238}%
\pgfsetfillcolor{currentfill}%
\pgfsetlinewidth{0.000000pt}%
\definecolor{currentstroke}{rgb}{0.000000,0.000000,0.000000}%
\pgfsetstrokecolor{currentstroke}%
\pgfsetstrokeopacity{0.000000}%
\pgfsetdash{}{0pt}%
\pgfpathmoveto{\pgfqpoint{0.981446in}{1.653571in}}%
\pgfpathlineto{\pgfqpoint{0.966686in}{1.656395in}}%
\pgfpathlineto{\pgfqpoint{0.964271in}{1.652566in}}%
\pgfpathlineto{\pgfqpoint{0.957302in}{1.649942in}}%
\pgfpathlineto{\pgfqpoint{0.955593in}{1.658524in}}%
\pgfpathlineto{\pgfqpoint{0.944841in}{1.665499in}}%
\pgfpathlineto{\pgfqpoint{0.941856in}{1.670351in}}%
\pgfpathlineto{\pgfqpoint{0.936043in}{1.670714in}}%
\pgfpathlineto{\pgfqpoint{0.926926in}{1.665686in}}%
\pgfpathlineto{\pgfqpoint{0.923792in}{1.677916in}}%
\pgfpathlineto{\pgfqpoint{0.918191in}{1.680868in}}%
\pgfpathlineto{\pgfqpoint{0.917372in}{1.685810in}}%
\pgfpathlineto{\pgfqpoint{0.913359in}{1.689243in}}%
\pgfpathlineto{\pgfqpoint{0.914111in}{1.694407in}}%
\pgfpathlineto{\pgfqpoint{0.917747in}{1.701623in}}%
\pgfpathlineto{\pgfqpoint{0.917545in}{1.708081in}}%
\pgfpathlineto{\pgfqpoint{0.914970in}{1.712349in}}%
\pgfpathlineto{\pgfqpoint{0.915399in}{1.718366in}}%
\pgfpathlineto{\pgfqpoint{0.921093in}{1.727298in}}%
\pgfpathlineto{\pgfqpoint{0.926725in}{1.726127in}}%
\pgfpathlineto{\pgfqpoint{0.927484in}{1.729880in}}%
\pgfpathlineto{\pgfqpoint{0.936258in}{1.730024in}}%
\pgfpathlineto{\pgfqpoint{0.940260in}{1.735106in}}%
\pgfpathlineto{\pgfqpoint{0.945308in}{1.734086in}}%
\pgfpathlineto{\pgfqpoint{0.949935in}{1.756722in}}%
\pgfpathlineto{\pgfqpoint{0.942694in}{1.758174in}}%
\pgfpathlineto{\pgfqpoint{0.948336in}{1.785640in}}%
\pgfpathlineto{\pgfqpoint{0.962444in}{1.782513in}}%
\pgfpathlineto{\pgfqpoint{0.964785in}{1.774681in}}%
\pgfpathlineto{\pgfqpoint{0.960752in}{1.754525in}}%
\pgfpathlineto{\pgfqpoint{0.972359in}{1.752059in}}%
\pgfpathlineto{\pgfqpoint{0.967944in}{1.729505in}}%
\pgfpathlineto{\pgfqpoint{0.979570in}{1.727393in}}%
\pgfpathlineto{\pgfqpoint{0.977851in}{1.718526in}}%
\pgfpathlineto{\pgfqpoint{0.983438in}{1.717518in}}%
\pgfpathlineto{\pgfqpoint{0.987484in}{1.710382in}}%
\pgfpathlineto{\pgfqpoint{0.985045in}{1.700348in}}%
\pgfpathlineto{\pgfqpoint{0.982127in}{1.696489in}}%
\pgfpathlineto{\pgfqpoint{0.973346in}{1.694790in}}%
\pgfpathlineto{\pgfqpoint{0.971079in}{1.692732in}}%
\pgfpathlineto{\pgfqpoint{0.969357in}{1.685011in}}%
\pgfpathlineto{\pgfqpoint{0.973053in}{1.678906in}}%
\pgfpathlineto{\pgfqpoint{0.972309in}{1.661350in}}%
\pgfpathlineto{\pgfqpoint{0.981446in}{1.653571in}}%
\pgfpathclose%
\pgfusepath{fill}%
\end{pgfscope}%
\begin{pgfscope}%
\pgfpathrectangle{\pgfqpoint{0.100000in}{0.100000in}}{\pgfqpoint{3.007045in}{1.925000in}}%
\pgfusepath{clip}%
\pgfsetbuttcap%
\pgfsetmiterjoin%
\definecolor{currentfill}{rgb}{0.607151,0.839677,0.587928}%
\pgfsetfillcolor{currentfill}%
\pgfsetlinewidth{0.000000pt}%
\definecolor{currentstroke}{rgb}{0.000000,0.000000,0.000000}%
\pgfsetstrokecolor{currentstroke}%
\pgfsetstrokeopacity{0.000000}%
\pgfsetdash{}{0pt}%
\pgfpathmoveto{\pgfqpoint{1.631190in}{1.247136in}}%
\pgfpathlineto{\pgfqpoint{1.586476in}{1.248770in}}%
\pgfpathlineto{\pgfqpoint{1.587386in}{1.271651in}}%
\pgfpathlineto{\pgfqpoint{1.586068in}{1.271711in}}%
\pgfpathlineto{\pgfqpoint{1.587220in}{1.294590in}}%
\pgfpathlineto{\pgfqpoint{1.609594in}{1.293697in}}%
\pgfpathlineto{\pgfqpoint{1.610083in}{1.285141in}}%
\pgfpathlineto{\pgfqpoint{1.618684in}{1.284817in}}%
\pgfpathlineto{\pgfqpoint{1.621155in}{1.286648in}}%
\pgfpathlineto{\pgfqpoint{1.638392in}{1.287167in}}%
\pgfpathlineto{\pgfqpoint{1.638646in}{1.292816in}}%
\pgfpathlineto{\pgfqpoint{1.643851in}{1.292672in}}%
\pgfpathlineto{\pgfqpoint{1.643733in}{1.288581in}}%
\pgfpathlineto{\pgfqpoint{1.635249in}{1.282362in}}%
\pgfpathlineto{\pgfqpoint{1.632031in}{1.278421in}}%
\pgfpathlineto{\pgfqpoint{1.631190in}{1.247136in}}%
\pgfpathclose%
\pgfusepath{fill}%
\end{pgfscope}%
\begin{pgfscope}%
\pgfpathrectangle{\pgfqpoint{0.100000in}{0.100000in}}{\pgfqpoint{3.007045in}{1.925000in}}%
\pgfusepath{clip}%
\pgfsetbuttcap%
\pgfsetmiterjoin%
\definecolor{currentfill}{rgb}{0.568397,0.821592,0.556063}%
\pgfsetfillcolor{currentfill}%
\pgfsetlinewidth{0.000000pt}%
\definecolor{currentstroke}{rgb}{0.000000,0.000000,0.000000}%
\pgfsetstrokecolor{currentstroke}%
\pgfsetstrokeopacity{0.000000}%
\pgfsetdash{}{0pt}%
\pgfpathmoveto{\pgfqpoint{0.683125in}{1.670416in}}%
\pgfpathlineto{\pgfqpoint{0.690132in}{1.668259in}}%
\pgfpathlineto{\pgfqpoint{0.695518in}{1.663103in}}%
\pgfpathlineto{\pgfqpoint{0.694447in}{1.659187in}}%
\pgfpathlineto{\pgfqpoint{0.689503in}{1.657799in}}%
\pgfpathlineto{\pgfqpoint{0.689599in}{1.651922in}}%
\pgfpathlineto{\pgfqpoint{0.678258in}{1.651609in}}%
\pgfpathlineto{\pgfqpoint{0.678193in}{1.649100in}}%
\pgfpathlineto{\pgfqpoint{0.684947in}{1.644999in}}%
\pgfpathlineto{\pgfqpoint{0.685898in}{1.641698in}}%
\pgfpathlineto{\pgfqpoint{0.678278in}{1.634204in}}%
\pgfpathlineto{\pgfqpoint{0.677402in}{1.629304in}}%
\pgfpathlineto{\pgfqpoint{0.671575in}{1.622748in}}%
\pgfpathlineto{\pgfqpoint{0.683738in}{1.619619in}}%
\pgfpathlineto{\pgfqpoint{0.680340in}{1.605736in}}%
\pgfpathlineto{\pgfqpoint{0.653418in}{1.613391in}}%
\pgfpathlineto{\pgfqpoint{0.651898in}{1.607795in}}%
\pgfpathlineto{\pgfqpoint{0.613279in}{1.618283in}}%
\pgfpathlineto{\pgfqpoint{0.619536in}{1.640423in}}%
\pgfpathlineto{\pgfqpoint{0.628849in}{1.673291in}}%
\pgfpathlineto{\pgfqpoint{0.630991in}{1.684466in}}%
\pgfpathlineto{\pgfqpoint{0.683125in}{1.670416in}}%
\pgfpathclose%
\pgfusepath{fill}%
\end{pgfscope}%
\begin{pgfscope}%
\pgfpathrectangle{\pgfqpoint{0.100000in}{0.100000in}}{\pgfqpoint{3.007045in}{1.925000in}}%
\pgfusepath{clip}%
\pgfsetbuttcap%
\pgfsetmiterjoin%
\definecolor{currentfill}{rgb}{0.765198,0.907328,0.738147}%
\pgfsetfillcolor{currentfill}%
\pgfsetlinewidth{0.000000pt}%
\definecolor{currentstroke}{rgb}{0.000000,0.000000,0.000000}%
\pgfsetstrokecolor{currentstroke}%
\pgfsetstrokeopacity{0.000000}%
\pgfsetdash{}{0pt}%
\pgfpathmoveto{\pgfqpoint{1.072883in}{1.296139in}}%
\pgfpathlineto{\pgfqpoint{1.062663in}{1.231515in}}%
\pgfpathlineto{\pgfqpoint{1.061002in}{1.220921in}}%
\pgfpathlineto{\pgfqpoint{1.057846in}{1.219034in}}%
\pgfpathlineto{\pgfqpoint{1.011802in}{1.226459in}}%
\pgfpathlineto{\pgfqpoint{1.012248in}{1.233395in}}%
\pgfpathlineto{\pgfqpoint{1.015989in}{1.239279in}}%
\pgfpathlineto{\pgfqpoint{1.016300in}{1.244056in}}%
\pgfpathlineto{\pgfqpoint{1.022492in}{1.247657in}}%
\pgfpathlineto{\pgfqpoint{0.990893in}{1.253051in}}%
\pgfpathlineto{\pgfqpoint{0.974041in}{1.256437in}}%
\pgfpathlineto{\pgfqpoint{0.973337in}{1.262333in}}%
\pgfpathlineto{\pgfqpoint{0.981761in}{1.312933in}}%
\pgfpathlineto{\pgfqpoint{0.989953in}{1.316140in}}%
\pgfpathlineto{\pgfqpoint{0.994498in}{1.314874in}}%
\pgfpathlineto{\pgfqpoint{1.008923in}{1.315572in}}%
\pgfpathlineto{\pgfqpoint{1.015074in}{1.317685in}}%
\pgfpathlineto{\pgfqpoint{1.019386in}{1.315423in}}%
\pgfpathlineto{\pgfqpoint{1.027613in}{1.314125in}}%
\pgfpathlineto{\pgfqpoint{1.033610in}{1.309798in}}%
\pgfpathlineto{\pgfqpoint{1.050961in}{1.311223in}}%
\pgfpathlineto{\pgfqpoint{1.052269in}{1.305131in}}%
\pgfpathlineto{\pgfqpoint{1.057220in}{1.307768in}}%
\pgfpathlineto{\pgfqpoint{1.057913in}{1.312141in}}%
\pgfpathlineto{\pgfqpoint{1.067332in}{1.310685in}}%
\pgfpathlineto{\pgfqpoint{1.068540in}{1.305641in}}%
\pgfpathlineto{\pgfqpoint{1.067620in}{1.298977in}}%
\pgfpathlineto{\pgfqpoint{1.072883in}{1.296139in}}%
\pgfpathclose%
\pgfusepath{fill}%
\end{pgfscope}%
\begin{pgfscope}%
\pgfpathrectangle{\pgfqpoint{0.100000in}{0.100000in}}{\pgfqpoint{3.007045in}{1.925000in}}%
\pgfusepath{clip}%
\pgfsetbuttcap%
\pgfsetmiterjoin%
\definecolor{currentfill}{rgb}{0.638970,0.854179,0.615240}%
\pgfsetfillcolor{currentfill}%
\pgfsetlinewidth{0.000000pt}%
\definecolor{currentstroke}{rgb}{0.000000,0.000000,0.000000}%
\pgfsetstrokecolor{currentstroke}%
\pgfsetstrokeopacity{0.000000}%
\pgfsetdash{}{0pt}%
\pgfpathmoveto{\pgfqpoint{1.772350in}{0.979565in}}%
\pgfpathlineto{\pgfqpoint{1.772364in}{0.997534in}}%
\pgfpathlineto{\pgfqpoint{1.768458in}{1.001242in}}%
\pgfpathlineto{\pgfqpoint{1.768402in}{1.026569in}}%
\pgfpathlineto{\pgfqpoint{1.767740in}{1.049552in}}%
\pgfpathlineto{\pgfqpoint{1.768296in}{1.069654in}}%
\pgfpathlineto{\pgfqpoint{1.792356in}{1.069644in}}%
\pgfpathlineto{\pgfqpoint{1.792358in}{1.025328in}}%
\pgfpathlineto{\pgfqpoint{1.820409in}{1.024532in}}%
\pgfpathlineto{\pgfqpoint{1.822007in}{1.020648in}}%
\pgfpathlineto{\pgfqpoint{1.821509in}{0.984906in}}%
\pgfpathlineto{\pgfqpoint{1.792483in}{0.985928in}}%
\pgfpathlineto{\pgfqpoint{1.792522in}{0.979422in}}%
\pgfpathlineto{\pgfqpoint{1.772350in}{0.979565in}}%
\pgfpathclose%
\pgfusepath{fill}%
\end{pgfscope}%
\begin{pgfscope}%
\pgfpathrectangle{\pgfqpoint{0.100000in}{0.100000in}}{\pgfqpoint{3.007045in}{1.925000in}}%
\pgfusepath{clip}%
\pgfsetbuttcap%
\pgfsetmiterjoin%
\definecolor{currentfill}{rgb}{0.657670,0.862053,0.633449}%
\pgfsetfillcolor{currentfill}%
\pgfsetlinewidth{0.000000pt}%
\definecolor{currentstroke}{rgb}{0.000000,0.000000,0.000000}%
\pgfsetstrokecolor{currentstroke}%
\pgfsetstrokeopacity{0.000000}%
\pgfsetdash{}{0pt}%
\pgfpathmoveto{\pgfqpoint{2.584120in}{0.229934in}}%
\pgfpathlineto{\pgfqpoint{2.584920in}{0.232373in}}%
\pgfpathlineto{\pgfqpoint{2.594250in}{0.236772in}}%
\pgfpathlineto{\pgfqpoint{2.597851in}{0.242148in}}%
\pgfpathlineto{\pgfqpoint{2.603807in}{0.246137in}}%
\pgfpathlineto{\pgfqpoint{2.612413in}{0.245063in}}%
\pgfpathlineto{\pgfqpoint{2.616307in}{0.242603in}}%
\pgfpathlineto{\pgfqpoint{2.608919in}{0.238771in}}%
\pgfpathlineto{\pgfqpoint{2.602623in}{0.240537in}}%
\pgfpathlineto{\pgfqpoint{2.601276in}{0.237478in}}%
\pgfpathlineto{\pgfqpoint{2.591656in}{0.231982in}}%
\pgfpathlineto{\pgfqpoint{2.584120in}{0.229934in}}%
\pgfpathclose%
\pgfusepath{fill}%
\end{pgfscope}%
\begin{pgfscope}%
\pgfpathrectangle{\pgfqpoint{0.100000in}{0.100000in}}{\pgfqpoint{3.007045in}{1.925000in}}%
\pgfusepath{clip}%
\pgfsetbuttcap%
\pgfsetmiterjoin%
\definecolor{currentfill}{rgb}{0.657670,0.862053,0.633449}%
\pgfsetfillcolor{currentfill}%
\pgfsetlinewidth{0.000000pt}%
\definecolor{currentstroke}{rgb}{0.000000,0.000000,0.000000}%
\pgfsetstrokecolor{currentstroke}%
\pgfsetstrokeopacity{0.000000}%
\pgfsetdash{}{0pt}%
\pgfpathmoveto{\pgfqpoint{2.599950in}{0.318474in}}%
\pgfpathlineto{\pgfqpoint{2.628250in}{0.323017in}}%
\pgfpathlineto{\pgfqpoint{2.623105in}{0.353228in}}%
\pgfpathlineto{\pgfqpoint{2.622244in}{0.358169in}}%
\pgfpathlineto{\pgfqpoint{2.661292in}{0.364487in}}%
\pgfpathlineto{\pgfqpoint{2.670514in}{0.365060in}}%
\pgfpathlineto{\pgfqpoint{2.671006in}{0.349163in}}%
\pgfpathlineto{\pgfqpoint{2.673290in}{0.330949in}}%
\pgfpathlineto{\pgfqpoint{2.672635in}{0.325576in}}%
\pgfpathlineto{\pgfqpoint{2.666787in}{0.322924in}}%
\pgfpathlineto{\pgfqpoint{2.663898in}{0.305786in}}%
\pgfpathlineto{\pgfqpoint{2.665452in}{0.296790in}}%
\pgfpathlineto{\pgfqpoint{2.656966in}{0.286274in}}%
\pgfpathlineto{\pgfqpoint{2.652560in}{0.288018in}}%
\pgfpathlineto{\pgfqpoint{2.650535in}{0.283393in}}%
\pgfpathlineto{\pgfqpoint{2.645495in}{0.280171in}}%
\pgfpathlineto{\pgfqpoint{2.636303in}{0.281237in}}%
\pgfpathlineto{\pgfqpoint{2.633609in}{0.278212in}}%
\pgfpathlineto{\pgfqpoint{2.622499in}{0.274880in}}%
\pgfpathlineto{\pgfqpoint{2.616272in}{0.281159in}}%
\pgfpathlineto{\pgfqpoint{2.616427in}{0.288354in}}%
\pgfpathlineto{\pgfqpoint{2.614021in}{0.297409in}}%
\pgfpathlineto{\pgfqpoint{2.611003in}{0.299907in}}%
\pgfpathlineto{\pgfqpoint{2.607596in}{0.307893in}}%
\pgfpathlineto{\pgfqpoint{2.601344in}{0.312735in}}%
\pgfpathlineto{\pgfqpoint{2.599950in}{0.318474in}}%
\pgfpathclose%
\pgfusepath{fill}%
\end{pgfscope}%
\begin{pgfscope}%
\pgfpathrectangle{\pgfqpoint{0.100000in}{0.100000in}}{\pgfqpoint{3.007045in}{1.925000in}}%
\pgfusepath{clip}%
\pgfsetbuttcap%
\pgfsetmiterjoin%
\definecolor{currentfill}{rgb}{0.551788,0.813841,0.542407}%
\pgfsetfillcolor{currentfill}%
\pgfsetlinewidth{0.000000pt}%
\definecolor{currentstroke}{rgb}{0.000000,0.000000,0.000000}%
\pgfsetstrokecolor{currentstroke}%
\pgfsetstrokeopacity{0.000000}%
\pgfsetdash{}{0pt}%
\pgfpathmoveto{\pgfqpoint{2.139795in}{0.676959in}}%
\pgfpathlineto{\pgfqpoint{2.130566in}{0.673645in}}%
\pgfpathlineto{\pgfqpoint{2.115516in}{0.670799in}}%
\pgfpathlineto{\pgfqpoint{2.113560in}{0.697106in}}%
\pgfpathlineto{\pgfqpoint{2.090671in}{0.695575in}}%
\pgfpathlineto{\pgfqpoint{2.087936in}{0.742436in}}%
\pgfpathlineto{\pgfqpoint{2.116053in}{0.743899in}}%
\pgfpathlineto{\pgfqpoint{2.141820in}{0.745927in}}%
\pgfpathlineto{\pgfqpoint{2.141241in}{0.722568in}}%
\pgfpathlineto{\pgfqpoint{2.139795in}{0.676959in}}%
\pgfpathclose%
\pgfusepath{fill}%
\end{pgfscope}%
\begin{pgfscope}%
\pgfpathrectangle{\pgfqpoint{0.100000in}{0.100000in}}{\pgfqpoint{3.007045in}{1.925000in}}%
\pgfusepath{clip}%
\pgfsetbuttcap%
\pgfsetmiterjoin%
\definecolor{currentfill}{rgb}{0.468743,0.775087,0.474125}%
\pgfsetfillcolor{currentfill}%
\pgfsetlinewidth{0.000000pt}%
\definecolor{currentstroke}{rgb}{0.000000,0.000000,0.000000}%
\pgfsetstrokecolor{currentstroke}%
\pgfsetstrokeopacity{0.000000}%
\pgfsetdash{}{0pt}%
\pgfpathmoveto{\pgfqpoint{1.054085in}{0.578918in}}%
\pgfpathlineto{\pgfqpoint{1.057990in}{0.585618in}}%
\pgfpathlineto{\pgfqpoint{1.069716in}{0.583261in}}%
\pgfpathlineto{\pgfqpoint{1.065530in}{0.579554in}}%
\pgfpathlineto{\pgfqpoint{1.060064in}{0.581478in}}%
\pgfpathlineto{\pgfqpoint{1.054085in}{0.578918in}}%
\pgfpathclose%
\pgfusepath{fill}%
\end{pgfscope}%
\begin{pgfscope}%
\pgfpathrectangle{\pgfqpoint{0.100000in}{0.100000in}}{\pgfqpoint{3.007045in}{1.925000in}}%
\pgfusepath{clip}%
\pgfsetbuttcap%
\pgfsetmiterjoin%
\definecolor{currentfill}{rgb}{0.468743,0.775087,0.474125}%
\pgfsetfillcolor{currentfill}%
\pgfsetlinewidth{0.000000pt}%
\definecolor{currentstroke}{rgb}{0.000000,0.000000,0.000000}%
\pgfsetstrokecolor{currentstroke}%
\pgfsetstrokeopacity{0.000000}%
\pgfsetdash{}{0pt}%
\pgfpathmoveto{\pgfqpoint{1.094825in}{0.552023in}}%
\pgfpathlineto{\pgfqpoint{1.087738in}{0.562041in}}%
\pgfpathlineto{\pgfqpoint{1.084656in}{0.570735in}}%
\pgfpathlineto{\pgfqpoint{1.085398in}{0.574624in}}%
\pgfpathlineto{\pgfqpoint{1.094010in}{0.578238in}}%
\pgfpathlineto{\pgfqpoint{1.104098in}{0.576218in}}%
\pgfpathlineto{\pgfqpoint{1.107468in}{0.572007in}}%
\pgfpathlineto{\pgfqpoint{1.113289in}{0.568890in}}%
\pgfpathlineto{\pgfqpoint{1.115269in}{0.562992in}}%
\pgfpathlineto{\pgfqpoint{1.112904in}{0.558446in}}%
\pgfpathlineto{\pgfqpoint{1.108126in}{0.556840in}}%
\pgfpathlineto{\pgfqpoint{1.104455in}{0.552095in}}%
\pgfpathlineto{\pgfqpoint{1.094825in}{0.552023in}}%
\pgfpathclose%
\pgfusepath{fill}%
\end{pgfscope}%
\begin{pgfscope}%
\pgfpathrectangle{\pgfqpoint{0.100000in}{0.100000in}}{\pgfqpoint{3.007045in}{1.925000in}}%
\pgfusepath{clip}%
\pgfsetbuttcap%
\pgfsetmiterjoin%
\definecolor{currentfill}{rgb}{0.420392,0.751711,0.445829}%
\pgfsetfillcolor{currentfill}%
\pgfsetlinewidth{0.000000pt}%
\definecolor{currentstroke}{rgb}{0.000000,0.000000,0.000000}%
\pgfsetstrokecolor{currentstroke}%
\pgfsetstrokeopacity{0.000000}%
\pgfsetdash{}{0pt}%
\pgfpathmoveto{\pgfqpoint{1.638404in}{1.708645in}}%
\pgfpathlineto{\pgfqpoint{1.637567in}{1.685608in}}%
\pgfpathlineto{\pgfqpoint{1.638529in}{1.679874in}}%
\pgfpathlineto{\pgfqpoint{1.615660in}{1.680679in}}%
\pgfpathlineto{\pgfqpoint{1.614445in}{1.686458in}}%
\pgfpathlineto{\pgfqpoint{1.615296in}{1.709477in}}%
\pgfpathlineto{\pgfqpoint{1.638404in}{1.708645in}}%
\pgfpathclose%
\pgfusepath{fill}%
\end{pgfscope}%
\begin{pgfscope}%
\pgfpathrectangle{\pgfqpoint{0.100000in}{0.100000in}}{\pgfqpoint{3.007045in}{1.925000in}}%
\pgfusepath{clip}%
\pgfsetbuttcap%
\pgfsetmiterjoin%
\definecolor{currentfill}{rgb}{0.524106,0.800923,0.519646}%
\pgfsetfillcolor{currentfill}%
\pgfsetlinewidth{0.000000pt}%
\definecolor{currentstroke}{rgb}{0.000000,0.000000,0.000000}%
\pgfsetstrokecolor{currentstroke}%
\pgfsetstrokeopacity{0.000000}%
\pgfsetdash{}{0pt}%
\pgfpathmoveto{\pgfqpoint{1.245609in}{1.461365in}}%
\pgfpathlineto{\pgfqpoint{1.218116in}{1.464571in}}%
\pgfpathlineto{\pgfqpoint{1.193681in}{1.468364in}}%
\pgfpathlineto{\pgfqpoint{1.197829in}{1.512182in}}%
\pgfpathlineto{\pgfqpoint{1.199309in}{1.522550in}}%
\pgfpathlineto{\pgfqpoint{1.197692in}{1.530943in}}%
\pgfpathlineto{\pgfqpoint{1.193646in}{1.534813in}}%
\pgfpathlineto{\pgfqpoint{1.188802in}{1.543168in}}%
\pgfpathlineto{\pgfqpoint{1.180687in}{1.550302in}}%
\pgfpathlineto{\pgfqpoint{1.174953in}{1.552958in}}%
\pgfpathlineto{\pgfqpoint{1.170691in}{1.563235in}}%
\pgfpathlineto{\pgfqpoint{1.169718in}{1.571914in}}%
\pgfpathlineto{\pgfqpoint{1.154054in}{1.574061in}}%
\pgfpathlineto{\pgfqpoint{1.160708in}{1.586055in}}%
\pgfpathlineto{\pgfqpoint{1.163161in}{1.587525in}}%
\pgfpathlineto{\pgfqpoint{1.135376in}{1.591485in}}%
\pgfpathlineto{\pgfqpoint{1.138319in}{1.605460in}}%
\pgfpathlineto{\pgfqpoint{1.140335in}{1.607233in}}%
\pgfpathlineto{\pgfqpoint{1.155409in}{1.604740in}}%
\pgfpathlineto{\pgfqpoint{1.167306in}{1.606868in}}%
\pgfpathlineto{\pgfqpoint{1.169858in}{1.624279in}}%
\pgfpathlineto{\pgfqpoint{1.171947in}{1.631734in}}%
\pgfpathlineto{\pgfqpoint{1.179523in}{1.630622in}}%
\pgfpathlineto{\pgfqpoint{1.184148in}{1.635775in}}%
\pgfpathlineto{\pgfqpoint{1.189792in}{1.634979in}}%
\pgfpathlineto{\pgfqpoint{1.190334in}{1.638786in}}%
\pgfpathlineto{\pgfqpoint{1.197757in}{1.637726in}}%
\pgfpathlineto{\pgfqpoint{1.201711in}{1.637190in}}%
\pgfpathlineto{\pgfqpoint{1.200926in}{1.631495in}}%
\pgfpathlineto{\pgfqpoint{1.212290in}{1.629948in}}%
\pgfpathlineto{\pgfqpoint{1.222383in}{1.622797in}}%
\pgfpathlineto{\pgfqpoint{1.221965in}{1.610407in}}%
\pgfpathlineto{\pgfqpoint{1.230291in}{1.608576in}}%
\pgfpathlineto{\pgfqpoint{1.227555in}{1.587451in}}%
\pgfpathlineto{\pgfqpoint{1.224309in}{1.576380in}}%
\pgfpathlineto{\pgfqpoint{1.247672in}{1.573410in}}%
\pgfpathlineto{\pgfqpoint{1.246154in}{1.561183in}}%
\pgfpathlineto{\pgfqpoint{1.257251in}{1.559802in}}%
\pgfpathlineto{\pgfqpoint{1.255504in}{1.539978in}}%
\pgfpathlineto{\pgfqpoint{1.251044in}{1.505633in}}%
\pgfpathlineto{\pgfqpoint{1.245609in}{1.461365in}}%
\pgfpathclose%
\pgfusepath{fill}%
\end{pgfscope}%
\begin{pgfscope}%
\pgfpathrectangle{\pgfqpoint{0.100000in}{0.100000in}}{\pgfqpoint{3.007045in}{1.925000in}}%
\pgfusepath{clip}%
\pgfsetbuttcap%
\pgfsetmiterjoin%
\definecolor{currentfill}{rgb}{0.648320,0.858116,0.624344}%
\pgfsetfillcolor{currentfill}%
\pgfsetlinewidth{0.000000pt}%
\definecolor{currentstroke}{rgb}{0.000000,0.000000,0.000000}%
\pgfsetstrokecolor{currentstroke}%
\pgfsetstrokeopacity{0.000000}%
\pgfsetdash{}{0pt}%
\pgfpathmoveto{\pgfqpoint{1.836747in}{1.453219in}}%
\pgfpathlineto{\pgfqpoint{1.831065in}{1.453126in}}%
\pgfpathlineto{\pgfqpoint{1.830743in}{1.476124in}}%
\pgfpathlineto{\pgfqpoint{1.842251in}{1.476305in}}%
\pgfpathlineto{\pgfqpoint{1.841921in}{1.499291in}}%
\pgfpathlineto{\pgfqpoint{1.853375in}{1.499519in}}%
\pgfpathlineto{\pgfqpoint{1.853444in}{1.494748in}}%
\pgfpathlineto{\pgfqpoint{1.864898in}{1.494975in}}%
\pgfpathlineto{\pgfqpoint{1.865473in}{1.453759in}}%
\pgfpathlineto{\pgfqpoint{1.836747in}{1.453219in}}%
\pgfpathclose%
\pgfusepath{fill}%
\end{pgfscope}%
\begin{pgfscope}%
\pgfpathrectangle{\pgfqpoint{0.100000in}{0.100000in}}{\pgfqpoint{3.007045in}{1.925000in}}%
\pgfusepath{clip}%
\pgfsetbuttcap%
\pgfsetmiterjoin%
\definecolor{currentfill}{rgb}{0.451765,0.767090,0.461207}%
\pgfsetfillcolor{currentfill}%
\pgfsetlinewidth{0.000000pt}%
\definecolor{currentstroke}{rgb}{0.000000,0.000000,0.000000}%
\pgfsetstrokecolor{currentstroke}%
\pgfsetstrokeopacity{0.000000}%
\pgfsetdash{}{0pt}%
\pgfpathmoveto{\pgfqpoint{1.507527in}{1.120611in}}%
\pgfpathlineto{\pgfqpoint{1.536726in}{1.118994in}}%
\pgfpathlineto{\pgfqpoint{1.535329in}{1.090418in}}%
\pgfpathlineto{\pgfqpoint{1.502124in}{1.092214in}}%
\pgfpathlineto{\pgfqpoint{1.478511in}{1.093767in}}%
\pgfpathlineto{\pgfqpoint{1.480135in}{1.122427in}}%
\pgfpathlineto{\pgfqpoint{1.507527in}{1.120611in}}%
\pgfpathclose%
\pgfusepath{fill}%
\end{pgfscope}%
\begin{pgfscope}%
\pgfpathrectangle{\pgfqpoint{0.100000in}{0.100000in}}{\pgfqpoint{3.007045in}{1.925000in}}%
\pgfusepath{clip}%
\pgfsetbuttcap%
\pgfsetmiterjoin%
\definecolor{currentfill}{rgb}{0.760523,0.905359,0.733595}%
\pgfsetfillcolor{currentfill}%
\pgfsetlinewidth{0.000000pt}%
\definecolor{currentstroke}{rgb}{0.000000,0.000000,0.000000}%
\pgfsetstrokecolor{currentstroke}%
\pgfsetstrokeopacity{0.000000}%
\pgfsetdash{}{0pt}%
\pgfpathmoveto{\pgfqpoint{1.450280in}{1.611196in}}%
\pgfpathlineto{\pgfqpoint{1.455103in}{1.611872in}}%
\pgfpathlineto{\pgfqpoint{1.459052in}{1.608999in}}%
\pgfpathlineto{\pgfqpoint{1.470594in}{1.606845in}}%
\pgfpathlineto{\pgfqpoint{1.480364in}{1.610710in}}%
\pgfpathlineto{\pgfqpoint{1.482806in}{1.614036in}}%
\pgfpathlineto{\pgfqpoint{1.495994in}{1.623275in}}%
\pgfpathlineto{\pgfqpoint{1.503158in}{1.631150in}}%
\pgfpathlineto{\pgfqpoint{1.506325in}{1.629275in}}%
\pgfpathlineto{\pgfqpoint{1.516525in}{1.631523in}}%
\pgfpathlineto{\pgfqpoint{1.518046in}{1.627144in}}%
\pgfpathlineto{\pgfqpoint{1.517258in}{1.619082in}}%
\pgfpathlineto{\pgfqpoint{1.514093in}{1.614101in}}%
\pgfpathlineto{\pgfqpoint{1.514131in}{1.607335in}}%
\pgfpathlineto{\pgfqpoint{1.518744in}{1.599173in}}%
\pgfpathlineto{\pgfqpoint{1.523621in}{1.593929in}}%
\pgfpathlineto{\pgfqpoint{1.526579in}{1.583044in}}%
\pgfpathlineto{\pgfqpoint{1.522749in}{1.580062in}}%
\pgfpathlineto{\pgfqpoint{1.518319in}{1.571035in}}%
\pgfpathlineto{\pgfqpoint{1.524166in}{1.567570in}}%
\pgfpathlineto{\pgfqpoint{1.494665in}{1.569370in}}%
\pgfpathlineto{\pgfqpoint{1.447246in}{1.572779in}}%
\pgfpathlineto{\pgfqpoint{1.450280in}{1.611196in}}%
\pgfpathclose%
\pgfusepath{fill}%
\end{pgfscope}%
\begin{pgfscope}%
\pgfpathrectangle{\pgfqpoint{0.100000in}{0.100000in}}{\pgfqpoint{3.007045in}{1.925000in}}%
\pgfusepath{clip}%
\pgfsetbuttcap%
\pgfsetmiterjoin%
\definecolor{currentfill}{rgb}{0.445490,0.764014,0.458131}%
\pgfsetfillcolor{currentfill}%
\pgfsetlinewidth{0.000000pt}%
\definecolor{currentstroke}{rgb}{0.000000,0.000000,0.000000}%
\pgfsetstrokecolor{currentstroke}%
\pgfsetstrokeopacity{0.000000}%
\pgfsetdash{}{0pt}%
\pgfpathmoveto{\pgfqpoint{1.139523in}{1.071438in}}%
\pgfpathlineto{\pgfqpoint{1.121195in}{1.073917in}}%
\pgfpathlineto{\pgfqpoint{1.124305in}{1.096464in}}%
\pgfpathlineto{\pgfqpoint{1.123300in}{1.108295in}}%
\pgfpathlineto{\pgfqpoint{1.121625in}{1.109932in}}%
\pgfpathlineto{\pgfqpoint{1.123423in}{1.117299in}}%
\pgfpathlineto{\pgfqpoint{1.123455in}{1.127357in}}%
\pgfpathlineto{\pgfqpoint{1.121297in}{1.132386in}}%
\pgfpathlineto{\pgfqpoint{1.128210in}{1.131426in}}%
\pgfpathlineto{\pgfqpoint{1.136431in}{1.191196in}}%
\pgfpathlineto{\pgfqpoint{1.142160in}{1.192960in}}%
\pgfpathlineto{\pgfqpoint{1.143492in}{1.187243in}}%
\pgfpathlineto{\pgfqpoint{1.146599in}{1.183208in}}%
\pgfpathlineto{\pgfqpoint{1.157547in}{1.181738in}}%
\pgfpathlineto{\pgfqpoint{1.164471in}{1.172432in}}%
\pgfpathlineto{\pgfqpoint{1.169868in}{1.170865in}}%
\pgfpathlineto{\pgfqpoint{1.172957in}{1.174640in}}%
\pgfpathlineto{\pgfqpoint{1.177735in}{1.174146in}}%
\pgfpathlineto{\pgfqpoint{1.180201in}{1.170746in}}%
\pgfpathlineto{\pgfqpoint{1.186288in}{1.164150in}}%
\pgfpathlineto{\pgfqpoint{1.193507in}{1.162201in}}%
\pgfpathlineto{\pgfqpoint{1.188060in}{1.158199in}}%
\pgfpathlineto{\pgfqpoint{1.187185in}{1.151775in}}%
\pgfpathlineto{\pgfqpoint{1.189161in}{1.136735in}}%
\pgfpathlineto{\pgfqpoint{1.193414in}{1.130886in}}%
\pgfpathlineto{\pgfqpoint{1.154865in}{1.136033in}}%
\pgfpathlineto{\pgfqpoint{1.150792in}{1.105369in}}%
\pgfpathlineto{\pgfqpoint{1.143557in}{1.105192in}}%
\pgfpathlineto{\pgfqpoint{1.140887in}{1.088045in}}%
\pgfpathlineto{\pgfqpoint{1.139523in}{1.071438in}}%
\pgfpathclose%
\pgfusepath{fill}%
\end{pgfscope}%
\begin{pgfscope}%
\pgfpathrectangle{\pgfqpoint{0.100000in}{0.100000in}}{\pgfqpoint{3.007045in}{1.925000in}}%
\pgfusepath{clip}%
\pgfsetbuttcap%
\pgfsetmiterjoin%
\definecolor{currentfill}{rgb}{0.727797,0.891580,0.701730}%
\pgfsetfillcolor{currentfill}%
\pgfsetlinewidth{0.000000pt}%
\definecolor{currentstroke}{rgb}{0.000000,0.000000,0.000000}%
\pgfsetstrokecolor{currentstroke}%
\pgfsetstrokeopacity{0.000000}%
\pgfsetdash{}{0pt}%
\pgfpathmoveto{\pgfqpoint{1.447246in}{1.572779in}}%
\pgfpathlineto{\pgfqpoint{1.494665in}{1.569370in}}%
\pgfpathlineto{\pgfqpoint{1.524166in}{1.567570in}}%
\pgfpathlineto{\pgfqpoint{1.526953in}{1.561321in}}%
\pgfpathlineto{\pgfqpoint{1.524722in}{1.557253in}}%
\pgfpathlineto{\pgfqpoint{1.527075in}{1.552396in}}%
\pgfpathlineto{\pgfqpoint{1.523533in}{1.542968in}}%
\pgfpathlineto{\pgfqpoint{1.524818in}{1.537110in}}%
\pgfpathlineto{\pgfqpoint{1.518468in}{1.535998in}}%
\pgfpathlineto{\pgfqpoint{1.518970in}{1.529708in}}%
\pgfpathlineto{\pgfqpoint{1.519350in}{1.527670in}}%
\pgfpathlineto{\pgfqpoint{1.512562in}{1.521360in}}%
\pgfpathlineto{\pgfqpoint{1.506367in}{1.525118in}}%
\pgfpathlineto{\pgfqpoint{1.501256in}{1.521545in}}%
\pgfpathlineto{\pgfqpoint{1.498577in}{1.523374in}}%
\pgfpathlineto{\pgfqpoint{1.489275in}{1.520056in}}%
\pgfpathlineto{\pgfqpoint{1.484012in}{1.521937in}}%
\pgfpathlineto{\pgfqpoint{1.479161in}{1.518521in}}%
\pgfpathlineto{\pgfqpoint{1.473297in}{1.518824in}}%
\pgfpathlineto{\pgfqpoint{1.464688in}{1.511398in}}%
\pgfpathlineto{\pgfqpoint{1.458166in}{1.512912in}}%
\pgfpathlineto{\pgfqpoint{1.452730in}{1.510166in}}%
\pgfpathlineto{\pgfqpoint{1.442332in}{1.509325in}}%
\pgfpathlineto{\pgfqpoint{1.445034in}{1.544102in}}%
\pgfpathlineto{\pgfqpoint{1.447246in}{1.572779in}}%
\pgfpathclose%
\pgfusepath{fill}%
\end{pgfscope}%
\begin{pgfscope}%
\pgfpathrectangle{\pgfqpoint{0.100000in}{0.100000in}}{\pgfqpoint{3.007045in}{1.925000in}}%
\pgfusepath{clip}%
\pgfsetbuttcap%
\pgfsetmiterjoin%
\definecolor{currentfill}{rgb}{0.612687,0.842261,0.592480}%
\pgfsetfillcolor{currentfill}%
\pgfsetlinewidth{0.000000pt}%
\definecolor{currentstroke}{rgb}{0.000000,0.000000,0.000000}%
\pgfsetstrokecolor{currentstroke}%
\pgfsetstrokeopacity{0.000000}%
\pgfsetdash{}{0pt}%
\pgfpathmoveto{\pgfqpoint{2.379259in}{0.944820in}}%
\pgfpathlineto{\pgfqpoint{2.373522in}{0.943638in}}%
\pgfpathlineto{\pgfqpoint{2.364223in}{0.935952in}}%
\pgfpathlineto{\pgfqpoint{2.354830in}{0.942630in}}%
\pgfpathlineto{\pgfqpoint{2.341668in}{0.946890in}}%
\pgfpathlineto{\pgfqpoint{2.329679in}{0.943109in}}%
\pgfpathlineto{\pgfqpoint{2.327024in}{0.948030in}}%
\pgfpathlineto{\pgfqpoint{2.321351in}{0.950113in}}%
\pgfpathlineto{\pgfqpoint{2.317579in}{0.954327in}}%
\pgfpathlineto{\pgfqpoint{2.322488in}{0.960571in}}%
\pgfpathlineto{\pgfqpoint{2.319716in}{0.965996in}}%
\pgfpathlineto{\pgfqpoint{2.308631in}{0.975465in}}%
\pgfpathlineto{\pgfqpoint{2.309261in}{0.983763in}}%
\pgfpathlineto{\pgfqpoint{2.317969in}{0.990753in}}%
\pgfpathlineto{\pgfqpoint{2.325950in}{0.982785in}}%
\pgfpathlineto{\pgfqpoint{2.334580in}{0.979199in}}%
\pgfpathlineto{\pgfqpoint{2.338194in}{0.987373in}}%
\pgfpathlineto{\pgfqpoint{2.336026in}{1.001374in}}%
\pgfpathlineto{\pgfqpoint{2.340163in}{1.003243in}}%
\pgfpathlineto{\pgfqpoint{2.341210in}{1.008147in}}%
\pgfpathlineto{\pgfqpoint{2.353818in}{1.009495in}}%
\pgfpathlineto{\pgfqpoint{2.358542in}{0.999323in}}%
\pgfpathlineto{\pgfqpoint{2.361719in}{1.000172in}}%
\pgfpathlineto{\pgfqpoint{2.372662in}{0.995532in}}%
\pgfpathlineto{\pgfqpoint{2.370660in}{0.984801in}}%
\pgfpathlineto{\pgfqpoint{2.376395in}{0.971244in}}%
\pgfpathlineto{\pgfqpoint{2.369682in}{0.964768in}}%
\pgfpathlineto{\pgfqpoint{2.374820in}{0.956140in}}%
\pgfpathlineto{\pgfqpoint{2.378270in}{0.952775in}}%
\pgfpathlineto{\pgfqpoint{2.379259in}{0.944820in}}%
\pgfpathclose%
\pgfusepath{fill}%
\end{pgfscope}%
\begin{pgfscope}%
\pgfpathrectangle{\pgfqpoint{0.100000in}{0.100000in}}{\pgfqpoint{3.007045in}{1.925000in}}%
\pgfusepath{clip}%
\pgfsetbuttcap%
\pgfsetmiterjoin%
\definecolor{currentfill}{rgb}{0.451765,0.767090,0.461207}%
\pgfsetfillcolor{currentfill}%
\pgfsetlinewidth{0.000000pt}%
\definecolor{currentstroke}{rgb}{0.000000,0.000000,0.000000}%
\pgfsetstrokecolor{currentstroke}%
\pgfsetstrokeopacity{0.000000}%
\pgfsetdash{}{0pt}%
\pgfpathmoveto{\pgfqpoint{1.511195in}{1.206454in}}%
\pgfpathlineto{\pgfqpoint{1.509436in}{1.177873in}}%
\pgfpathlineto{\pgfqpoint{1.448195in}{1.181868in}}%
\pgfpathlineto{\pgfqpoint{1.446935in}{1.181960in}}%
\pgfpathlineto{\pgfqpoint{1.449096in}{1.210470in}}%
\pgfpathlineto{\pgfqpoint{1.453477in}{1.210184in}}%
\pgfpathlineto{\pgfqpoint{1.510401in}{1.206498in}}%
\pgfpathlineto{\pgfqpoint{1.511195in}{1.206454in}}%
\pgfpathclose%
\pgfusepath{fill}%
\end{pgfscope}%
\begin{pgfscope}%
\pgfpathrectangle{\pgfqpoint{0.100000in}{0.100000in}}{\pgfqpoint{3.007045in}{1.925000in}}%
\pgfusepath{clip}%
\pgfsetbuttcap%
\pgfsetmiterjoin%
\definecolor{currentfill}{rgb}{0.529642,0.803506,0.524198}%
\pgfsetfillcolor{currentfill}%
\pgfsetlinewidth{0.000000pt}%
\definecolor{currentstroke}{rgb}{0.000000,0.000000,0.000000}%
\pgfsetstrokecolor{currentstroke}%
\pgfsetstrokeopacity{0.000000}%
\pgfsetdash{}{0pt}%
\pgfpathmoveto{\pgfqpoint{2.585937in}{0.924191in}}%
\pgfpathlineto{\pgfqpoint{2.599449in}{0.926351in}}%
\pgfpathlineto{\pgfqpoint{2.613555in}{0.916458in}}%
\pgfpathlineto{\pgfqpoint{2.607323in}{0.893730in}}%
\pgfpathlineto{\pgfqpoint{2.595973in}{0.904310in}}%
\pgfpathlineto{\pgfqpoint{2.582994in}{0.902658in}}%
\pgfpathlineto{\pgfqpoint{2.571314in}{0.892287in}}%
\pgfpathlineto{\pgfqpoint{2.567423in}{0.900852in}}%
\pgfpathlineto{\pgfqpoint{2.562238in}{0.907450in}}%
\pgfpathlineto{\pgfqpoint{2.551891in}{0.919345in}}%
\pgfpathlineto{\pgfqpoint{2.585937in}{0.924191in}}%
\pgfpathclose%
\pgfusepath{fill}%
\end{pgfscope}%
\begin{pgfscope}%
\pgfpathrectangle{\pgfqpoint{0.100000in}{0.100000in}}{\pgfqpoint{3.007045in}{1.925000in}}%
\pgfusepath{clip}%
\pgfsetbuttcap%
\pgfsetmiterjoin%
\definecolor{currentfill}{rgb}{0.769873,0.909296,0.742699}%
\pgfsetfillcolor{currentfill}%
\pgfsetlinewidth{0.000000pt}%
\definecolor{currentstroke}{rgb}{0.000000,0.000000,0.000000}%
\pgfsetstrokecolor{currentstroke}%
\pgfsetstrokeopacity{0.000000}%
\pgfsetdash{}{0pt}%
\pgfpathmoveto{\pgfqpoint{1.473940in}{0.485183in}}%
\pgfpathlineto{\pgfqpoint{1.441506in}{0.486973in}}%
\pgfpathlineto{\pgfqpoint{1.435567in}{0.493311in}}%
\pgfpathlineto{\pgfqpoint{1.434480in}{0.498439in}}%
\pgfpathlineto{\pgfqpoint{1.422748in}{0.507253in}}%
\pgfpathlineto{\pgfqpoint{1.418771in}{0.514227in}}%
\pgfpathlineto{\pgfqpoint{1.415118in}{0.514942in}}%
\pgfpathlineto{\pgfqpoint{1.410285in}{0.521202in}}%
\pgfpathlineto{\pgfqpoint{1.406272in}{0.522948in}}%
\pgfpathlineto{\pgfqpoint{1.401030in}{0.533989in}}%
\pgfpathlineto{\pgfqpoint{1.394375in}{0.537885in}}%
\pgfpathlineto{\pgfqpoint{1.390298in}{0.536623in}}%
\pgfpathlineto{\pgfqpoint{1.369726in}{0.539824in}}%
\pgfpathlineto{\pgfqpoint{1.364524in}{0.539321in}}%
\pgfpathlineto{\pgfqpoint{1.349461in}{0.546748in}}%
\pgfpathlineto{\pgfqpoint{1.336243in}{0.559459in}}%
\pgfpathlineto{\pgfqpoint{1.337515in}{0.574775in}}%
\pgfpathlineto{\pgfqpoint{1.350354in}{0.573776in}}%
\pgfpathlineto{\pgfqpoint{1.352050in}{0.594662in}}%
\pgfpathlineto{\pgfqpoint{1.363729in}{0.593646in}}%
\pgfpathlineto{\pgfqpoint{1.364071in}{0.597523in}}%
\pgfpathlineto{\pgfqpoint{1.386971in}{0.595647in}}%
\pgfpathlineto{\pgfqpoint{1.390664in}{0.594316in}}%
\pgfpathlineto{\pgfqpoint{1.391278in}{0.588321in}}%
\pgfpathlineto{\pgfqpoint{1.385588in}{0.580234in}}%
\pgfpathlineto{\pgfqpoint{1.389923in}{0.575693in}}%
\pgfpathlineto{\pgfqpoint{1.383908in}{0.571430in}}%
\pgfpathlineto{\pgfqpoint{1.429685in}{0.568203in}}%
\pgfpathlineto{\pgfqpoint{1.444627in}{0.567276in}}%
\pgfpathlineto{\pgfqpoint{1.441878in}{0.523060in}}%
\pgfpathlineto{\pgfqpoint{1.475915in}{0.520991in}}%
\pgfpathlineto{\pgfqpoint{1.473940in}{0.485183in}}%
\pgfpathclose%
\pgfusepath{fill}%
\end{pgfscope}%
\begin{pgfscope}%
\pgfpathrectangle{\pgfqpoint{0.100000in}{0.100000in}}{\pgfqpoint{3.007045in}{1.925000in}}%
\pgfusepath{clip}%
\pgfsetbuttcap%
\pgfsetmiterjoin%
\definecolor{currentfill}{rgb}{0.345098,0.714802,0.408920}%
\pgfsetfillcolor{currentfill}%
\pgfsetlinewidth{0.000000pt}%
\definecolor{currentstroke}{rgb}{0.000000,0.000000,0.000000}%
\pgfsetstrokecolor{currentstroke}%
\pgfsetstrokeopacity{0.000000}%
\pgfsetdash{}{0pt}%
\pgfpathmoveto{\pgfqpoint{2.559174in}{1.248298in}}%
\pgfpathlineto{\pgfqpoint{2.559394in}{1.254953in}}%
\pgfpathlineto{\pgfqpoint{2.556847in}{1.259865in}}%
\pgfpathlineto{\pgfqpoint{2.561814in}{1.273283in}}%
\pgfpathlineto{\pgfqpoint{2.567920in}{1.283023in}}%
\pgfpathlineto{\pgfqpoint{2.574047in}{1.302973in}}%
\pgfpathlineto{\pgfqpoint{2.576932in}{1.318538in}}%
\pgfpathlineto{\pgfqpoint{2.599096in}{1.322608in}}%
\pgfpathlineto{\pgfqpoint{2.606011in}{1.320523in}}%
\pgfpathlineto{\pgfqpoint{2.609389in}{1.317404in}}%
\pgfpathlineto{\pgfqpoint{2.608640in}{1.313232in}}%
\pgfpathlineto{\pgfqpoint{2.612866in}{1.308899in}}%
\pgfpathlineto{\pgfqpoint{2.608820in}{1.295800in}}%
\pgfpathlineto{\pgfqpoint{2.610752in}{1.291249in}}%
\pgfpathlineto{\pgfqpoint{2.616493in}{1.288353in}}%
\pgfpathlineto{\pgfqpoint{2.613885in}{1.282571in}}%
\pgfpathlineto{\pgfqpoint{2.614370in}{1.276431in}}%
\pgfpathlineto{\pgfqpoint{2.612460in}{1.264832in}}%
\pgfpathlineto{\pgfqpoint{2.609634in}{1.257414in}}%
\pgfpathlineto{\pgfqpoint{2.559174in}{1.248298in}}%
\pgfpathclose%
\pgfusepath{fill}%
\end{pgfscope}%
\begin{pgfscope}%
\pgfpathrectangle{\pgfqpoint{0.100000in}{0.100000in}}{\pgfqpoint{3.007045in}{1.925000in}}%
\pgfusepath{clip}%
\pgfsetbuttcap%
\pgfsetmiterjoin%
\definecolor{currentfill}{rgb}{0.671696,0.867958,0.647105}%
\pgfsetfillcolor{currentfill}%
\pgfsetlinewidth{0.000000pt}%
\definecolor{currentstroke}{rgb}{0.000000,0.000000,0.000000}%
\pgfsetstrokecolor{currentstroke}%
\pgfsetstrokeopacity{0.000000}%
\pgfsetdash{}{0pt}%
\pgfpathmoveto{\pgfqpoint{1.963320in}{1.146179in}}%
\pgfpathlineto{\pgfqpoint{1.963877in}{1.126688in}}%
\pgfpathlineto{\pgfqpoint{1.956143in}{1.126474in}}%
\pgfpathlineto{\pgfqpoint{1.956351in}{1.117548in}}%
\pgfpathlineto{\pgfqpoint{1.949655in}{1.115080in}}%
\pgfpathlineto{\pgfqpoint{1.944978in}{1.116719in}}%
\pgfpathlineto{\pgfqpoint{1.944438in}{1.140114in}}%
\pgfpathlineto{\pgfqpoint{1.920082in}{1.139613in}}%
\pgfpathlineto{\pgfqpoint{1.920007in}{1.151193in}}%
\pgfpathlineto{\pgfqpoint{1.909292in}{1.151308in}}%
\pgfpathlineto{\pgfqpoint{1.909303in}{1.157973in}}%
\pgfpathlineto{\pgfqpoint{1.939480in}{1.158394in}}%
\pgfpathlineto{\pgfqpoint{1.953772in}{1.157587in}}%
\pgfpathlineto{\pgfqpoint{1.955717in}{1.146050in}}%
\pgfpathlineto{\pgfqpoint{1.963320in}{1.146179in}}%
\pgfpathclose%
\pgfusepath{fill}%
\end{pgfscope}%
\begin{pgfscope}%
\pgfpathrectangle{\pgfqpoint{0.100000in}{0.100000in}}{\pgfqpoint{3.007045in}{1.925000in}}%
\pgfusepath{clip}%
\pgfsetbuttcap%
\pgfsetmiterjoin%
\definecolor{currentfill}{rgb}{0.612687,0.842261,0.592480}%
\pgfsetfillcolor{currentfill}%
\pgfsetlinewidth{0.000000pt}%
\definecolor{currentstroke}{rgb}{0.000000,0.000000,0.000000}%
\pgfsetstrokecolor{currentstroke}%
\pgfsetstrokeopacity{0.000000}%
\pgfsetdash{}{0pt}%
\pgfpathmoveto{\pgfqpoint{1.835366in}{1.168017in}}%
\pgfpathlineto{\pgfqpoint{1.812537in}{1.167939in}}%
\pgfpathlineto{\pgfqpoint{1.812471in}{1.163163in}}%
\pgfpathlineto{\pgfqpoint{1.792627in}{1.163111in}}%
\pgfpathlineto{\pgfqpoint{1.792583in}{1.168136in}}%
\pgfpathlineto{\pgfqpoint{1.767101in}{1.168173in}}%
\pgfpathlineto{\pgfqpoint{1.769538in}{1.174111in}}%
\pgfpathlineto{\pgfqpoint{1.766486in}{1.176026in}}%
\pgfpathlineto{\pgfqpoint{1.755095in}{1.176082in}}%
\pgfpathlineto{\pgfqpoint{1.755176in}{1.198941in}}%
\pgfpathlineto{\pgfqpoint{1.756771in}{1.198935in}}%
\pgfpathlineto{\pgfqpoint{1.762123in}{1.192600in}}%
\pgfpathlineto{\pgfqpoint{1.768029in}{1.189813in}}%
\pgfpathlineto{\pgfqpoint{1.772661in}{1.193312in}}%
\pgfpathlineto{\pgfqpoint{1.770162in}{1.207494in}}%
\pgfpathlineto{\pgfqpoint{1.792177in}{1.207263in}}%
\pgfpathlineto{\pgfqpoint{1.792230in}{1.201536in}}%
\pgfpathlineto{\pgfqpoint{1.811728in}{1.201377in}}%
\pgfpathlineto{\pgfqpoint{1.811789in}{1.208062in}}%
\pgfpathlineto{\pgfqpoint{1.834576in}{1.208087in}}%
\pgfpathlineto{\pgfqpoint{1.834849in}{1.196667in}}%
\pgfpathlineto{\pgfqpoint{1.835366in}{1.168017in}}%
\pgfpathclose%
\pgfusepath{fill}%
\end{pgfscope}%
\begin{pgfscope}%
\pgfpathrectangle{\pgfqpoint{0.100000in}{0.100000in}}{\pgfqpoint{3.007045in}{1.925000in}}%
\pgfusepath{clip}%
\pgfsetbuttcap%
\pgfsetmiterjoin%
\definecolor{currentfill}{rgb}{0.623760,0.847428,0.601584}%
\pgfsetfillcolor{currentfill}%
\pgfsetlinewidth{0.000000pt}%
\definecolor{currentstroke}{rgb}{0.000000,0.000000,0.000000}%
\pgfsetstrokecolor{currentstroke}%
\pgfsetstrokeopacity{0.000000}%
\pgfsetdash{}{0pt}%
\pgfpathmoveto{\pgfqpoint{2.068022in}{0.977336in}}%
\pgfpathlineto{\pgfqpoint{2.064440in}{0.977028in}}%
\pgfpathlineto{\pgfqpoint{2.061582in}{0.976889in}}%
\pgfpathlineto{\pgfqpoint{2.061897in}{0.966247in}}%
\pgfpathlineto{\pgfqpoint{2.059156in}{0.970494in}}%
\pgfpathlineto{\pgfqpoint{2.050739in}{0.968865in}}%
\pgfpathlineto{\pgfqpoint{2.039044in}{0.968405in}}%
\pgfpathlineto{\pgfqpoint{2.038888in}{0.984228in}}%
\pgfpathlineto{\pgfqpoint{2.028875in}{0.983778in}}%
\pgfpathlineto{\pgfqpoint{2.028045in}{0.989826in}}%
\pgfpathlineto{\pgfqpoint{2.022047in}{1.002653in}}%
\pgfpathlineto{\pgfqpoint{2.029271in}{1.013848in}}%
\pgfpathlineto{\pgfqpoint{2.023343in}{1.013571in}}%
\pgfpathlineto{\pgfqpoint{2.022844in}{1.028625in}}%
\pgfpathlineto{\pgfqpoint{2.026620in}{1.028636in}}%
\pgfpathlineto{\pgfqpoint{2.025561in}{1.047408in}}%
\pgfpathlineto{\pgfqpoint{2.053627in}{1.049022in}}%
\pgfpathlineto{\pgfqpoint{2.058254in}{1.047208in}}%
\pgfpathlineto{\pgfqpoint{2.064198in}{1.036075in}}%
\pgfpathlineto{\pgfqpoint{2.059538in}{1.028875in}}%
\pgfpathlineto{\pgfqpoint{2.067443in}{1.014097in}}%
\pgfpathlineto{\pgfqpoint{2.072603in}{1.010193in}}%
\pgfpathlineto{\pgfqpoint{2.077300in}{1.012074in}}%
\pgfpathlineto{\pgfqpoint{2.081022in}{1.010070in}}%
\pgfpathlineto{\pgfqpoint{2.082880in}{1.008816in}}%
\pgfpathlineto{\pgfqpoint{2.081317in}{1.001182in}}%
\pgfpathlineto{\pgfqpoint{2.082341in}{0.997150in}}%
\pgfpathlineto{\pgfqpoint{2.078631in}{0.992336in}}%
\pgfpathlineto{\pgfqpoint{2.080971in}{0.989204in}}%
\pgfpathlineto{\pgfqpoint{2.077777in}{0.982593in}}%
\pgfpathlineto{\pgfqpoint{2.070320in}{0.985858in}}%
\pgfpathlineto{\pgfqpoint{2.068022in}{0.977336in}}%
\pgfpathclose%
\pgfusepath{fill}%
\end{pgfscope}%
\begin{pgfscope}%
\pgfpathrectangle{\pgfqpoint{0.100000in}{0.100000in}}{\pgfqpoint{3.007045in}{1.925000in}}%
\pgfusepath{clip}%
\pgfsetbuttcap%
\pgfsetmiterjoin%
\definecolor{currentfill}{rgb}{0.676371,0.869927,0.651657}%
\pgfsetfillcolor{currentfill}%
\pgfsetlinewidth{0.000000pt}%
\definecolor{currentstroke}{rgb}{0.000000,0.000000,0.000000}%
\pgfsetstrokecolor{currentstroke}%
\pgfsetstrokeopacity{0.000000}%
\pgfsetdash{}{0pt}%
\pgfpathmoveto{\pgfqpoint{2.139795in}{0.676959in}}%
\pgfpathlineto{\pgfqpoint{2.141224in}{0.664018in}}%
\pgfpathlineto{\pgfqpoint{2.143322in}{0.646719in}}%
\pgfpathlineto{\pgfqpoint{2.121151in}{0.645024in}}%
\pgfpathlineto{\pgfqpoint{2.089615in}{0.643002in}}%
\pgfpathlineto{\pgfqpoint{2.088016in}{0.667046in}}%
\pgfpathlineto{\pgfqpoint{2.073893in}{0.665116in}}%
\pgfpathlineto{\pgfqpoint{2.072910in}{0.680552in}}%
\pgfpathlineto{\pgfqpoint{2.068533in}{0.682516in}}%
\pgfpathlineto{\pgfqpoint{2.067842in}{0.694087in}}%
\pgfpathlineto{\pgfqpoint{2.090671in}{0.695575in}}%
\pgfpathlineto{\pgfqpoint{2.113560in}{0.697106in}}%
\pgfpathlineto{\pgfqpoint{2.115516in}{0.670799in}}%
\pgfpathlineto{\pgfqpoint{2.130566in}{0.673645in}}%
\pgfpathlineto{\pgfqpoint{2.139795in}{0.676959in}}%
\pgfpathclose%
\pgfusepath{fill}%
\end{pgfscope}%
\begin{pgfscope}%
\pgfpathrectangle{\pgfqpoint{0.100000in}{0.100000in}}{\pgfqpoint{3.007045in}{1.925000in}}%
\pgfusepath{clip}%
\pgfsetbuttcap%
\pgfsetmiterjoin%
\definecolor{currentfill}{rgb}{0.585006,0.829343,0.569719}%
\pgfsetfillcolor{currentfill}%
\pgfsetlinewidth{0.000000pt}%
\definecolor{currentstroke}{rgb}{0.000000,0.000000,0.000000}%
\pgfsetstrokecolor{currentstroke}%
\pgfsetstrokeopacity{0.000000}%
\pgfsetdash{}{0pt}%
\pgfpathmoveto{\pgfqpoint{1.314888in}{1.740131in}}%
\pgfpathlineto{\pgfqpoint{1.309172in}{1.740759in}}%
\pgfpathlineto{\pgfqpoint{1.307173in}{1.723478in}}%
\pgfpathlineto{\pgfqpoint{1.305247in}{1.723687in}}%
\pgfpathlineto{\pgfqpoint{1.302614in}{1.700688in}}%
\pgfpathlineto{\pgfqpoint{1.283442in}{1.702901in}}%
\pgfpathlineto{\pgfqpoint{1.282748in}{1.697095in}}%
\pgfpathlineto{\pgfqpoint{1.277023in}{1.697772in}}%
\pgfpathlineto{\pgfqpoint{1.277730in}{1.703596in}}%
\pgfpathlineto{\pgfqpoint{1.272031in}{1.704281in}}%
\pgfpathlineto{\pgfqpoint{1.268284in}{1.704741in}}%
\pgfpathlineto{\pgfqpoint{1.269690in}{1.716224in}}%
\pgfpathlineto{\pgfqpoint{1.263981in}{1.716912in}}%
\pgfpathlineto{\pgfqpoint{1.267101in}{1.728242in}}%
\pgfpathlineto{\pgfqpoint{1.269873in}{1.750687in}}%
\pgfpathlineto{\pgfqpoint{1.265560in}{1.751233in}}%
\pgfpathlineto{\pgfqpoint{1.263629in}{1.757459in}}%
\pgfpathlineto{\pgfqpoint{1.265605in}{1.764425in}}%
\pgfpathlineto{\pgfqpoint{1.272584in}{1.760976in}}%
\pgfpathlineto{\pgfqpoint{1.289529in}{1.757826in}}%
\pgfpathlineto{\pgfqpoint{1.291202in}{1.771653in}}%
\pgfpathlineto{\pgfqpoint{1.292930in}{1.771456in}}%
\pgfpathlineto{\pgfqpoint{1.295653in}{1.794279in}}%
\pgfpathlineto{\pgfqpoint{1.332164in}{1.790080in}}%
\pgfpathlineto{\pgfqpoint{1.341649in}{1.789040in}}%
\pgfpathlineto{\pgfqpoint{1.341017in}{1.783290in}}%
\pgfpathlineto{\pgfqpoint{1.346744in}{1.782663in}}%
\pgfpathlineto{\pgfqpoint{1.346116in}{1.776879in}}%
\pgfpathlineto{\pgfqpoint{1.371755in}{1.774248in}}%
\pgfpathlineto{\pgfqpoint{1.369270in}{1.748112in}}%
\pgfpathlineto{\pgfqpoint{1.366019in}{1.751221in}}%
\pgfpathlineto{\pgfqpoint{1.355902in}{1.752996in}}%
\pgfpathlineto{\pgfqpoint{1.351556in}{1.758299in}}%
\pgfpathlineto{\pgfqpoint{1.341818in}{1.758120in}}%
\pgfpathlineto{\pgfqpoint{1.335725in}{1.759869in}}%
\pgfpathlineto{\pgfqpoint{1.324957in}{1.756945in}}%
\pgfpathlineto{\pgfqpoint{1.318715in}{1.758150in}}%
\pgfpathlineto{\pgfqpoint{1.317232in}{1.745061in}}%
\pgfpathlineto{\pgfqpoint{1.314888in}{1.740131in}}%
\pgfpathclose%
\pgfusepath{fill}%
\end{pgfscope}%
\begin{pgfscope}%
\pgfpathrectangle{\pgfqpoint{0.100000in}{0.100000in}}{\pgfqpoint{3.007045in}{1.925000in}}%
\pgfusepath{clip}%
\pgfsetbuttcap%
\pgfsetmiterjoin%
\definecolor{currentfill}{rgb}{0.812687,0.926644,0.787389}%
\pgfsetfillcolor{currentfill}%
\pgfsetlinewidth{0.000000pt}%
\definecolor{currentstroke}{rgb}{0.000000,0.000000,0.000000}%
\pgfsetstrokecolor{currentstroke}%
\pgfsetstrokeopacity{0.000000}%
\pgfsetdash{}{0pt}%
\pgfpathmoveto{\pgfqpoint{1.512921in}{0.424052in}}%
\pgfpathlineto{\pgfqpoint{1.470461in}{0.425819in}}%
\pgfpathlineto{\pgfqpoint{1.473940in}{0.485183in}}%
\pgfpathlineto{\pgfqpoint{1.514641in}{0.483384in}}%
\pgfpathlineto{\pgfqpoint{1.513462in}{0.453257in}}%
\pgfpathlineto{\pgfqpoint{1.514242in}{0.453233in}}%
\pgfpathlineto{\pgfqpoint{1.512921in}{0.424052in}}%
\pgfpathclose%
\pgfusepath{fill}%
\end{pgfscope}%
\begin{pgfscope}%
\pgfpathrectangle{\pgfqpoint{0.100000in}{0.100000in}}{\pgfqpoint{3.007045in}{1.925000in}}%
\pgfusepath{clip}%
\pgfsetbuttcap%
\pgfsetmiterjoin%
\definecolor{currentfill}{rgb}{0.816378,0.928120,0.791326}%
\pgfsetfillcolor{currentfill}%
\pgfsetlinewidth{0.000000pt}%
\definecolor{currentstroke}{rgb}{0.000000,0.000000,0.000000}%
\pgfsetstrokecolor{currentstroke}%
\pgfsetstrokeopacity{0.000000}%
\pgfsetdash{}{0pt}%
\pgfpathmoveto{\pgfqpoint{1.219256in}{0.700832in}}%
\pgfpathlineto{\pgfqpoint{1.265329in}{0.695996in}}%
\pgfpathlineto{\pgfqpoint{1.255027in}{0.637292in}}%
\pgfpathlineto{\pgfqpoint{1.205698in}{0.612940in}}%
\pgfpathlineto{\pgfqpoint{1.207983in}{0.628535in}}%
\pgfpathlineto{\pgfqpoint{1.212900in}{0.676457in}}%
\pgfpathlineto{\pgfqpoint{1.215312in}{0.701261in}}%
\pgfpathlineto{\pgfqpoint{1.219256in}{0.700832in}}%
\pgfpathclose%
\pgfusepath{fill}%
\end{pgfscope}%
\begin{pgfscope}%
\pgfpathrectangle{\pgfqpoint{0.100000in}{0.100000in}}{\pgfqpoint{3.007045in}{1.925000in}}%
\pgfusepath{clip}%
\pgfsetbuttcap%
\pgfsetmiterjoin%
\definecolor{currentfill}{rgb}{0.618224,0.844844,0.597032}%
\pgfsetfillcolor{currentfill}%
\pgfsetlinewidth{0.000000pt}%
\definecolor{currentstroke}{rgb}{0.000000,0.000000,0.000000}%
\pgfsetstrokecolor{currentstroke}%
\pgfsetstrokeopacity{0.000000}%
\pgfsetdash{}{0pt}%
\pgfpathmoveto{\pgfqpoint{1.347270in}{1.020215in}}%
\pgfpathlineto{\pgfqpoint{1.351668in}{1.019814in}}%
\pgfpathlineto{\pgfqpoint{1.348759in}{0.987002in}}%
\pgfpathlineto{\pgfqpoint{1.346674in}{0.987189in}}%
\pgfpathlineto{\pgfqpoint{1.342266in}{0.937180in}}%
\pgfpathlineto{\pgfqpoint{1.324480in}{0.938840in}}%
\pgfpathlineto{\pgfqpoint{1.327114in}{0.961542in}}%
\pgfpathlineto{\pgfqpoint{1.304327in}{0.963661in}}%
\pgfpathlineto{\pgfqpoint{1.304909in}{0.969450in}}%
\pgfpathlineto{\pgfqpoint{1.293521in}{0.970602in}}%
\pgfpathlineto{\pgfqpoint{1.293808in}{0.973457in}}%
\pgfpathlineto{\pgfqpoint{1.298942in}{1.024507in}}%
\pgfpathlineto{\pgfqpoint{1.347270in}{1.020215in}}%
\pgfpathclose%
\pgfusepath{fill}%
\end{pgfscope}%
\begin{pgfscope}%
\pgfpathrectangle{\pgfqpoint{0.100000in}{0.100000in}}{\pgfqpoint{3.007045in}{1.925000in}}%
\pgfusepath{clip}%
\pgfsetbuttcap%
\pgfsetmiterjoin%
\definecolor{currentfill}{rgb}{0.886505,0.956171,0.866128}%
\pgfsetfillcolor{currentfill}%
\pgfsetlinewidth{0.000000pt}%
\definecolor{currentstroke}{rgb}{0.000000,0.000000,0.000000}%
\pgfsetstrokecolor{currentstroke}%
\pgfsetstrokeopacity{0.000000}%
\pgfsetdash{}{0pt}%
\pgfpathmoveto{\pgfqpoint{0.610456in}{0.538425in}}%
\pgfpathlineto{\pgfqpoint{0.607882in}{0.541289in}}%
\pgfpathlineto{\pgfqpoint{0.608049in}{0.546486in}}%
\pgfpathlineto{\pgfqpoint{0.607286in}{0.548897in}}%
\pgfpathlineto{\pgfqpoint{0.610143in}{0.549575in}}%
\pgfpathlineto{\pgfqpoint{0.608906in}{0.546215in}}%
\pgfpathlineto{\pgfqpoint{0.609173in}{0.543568in}}%
\pgfpathlineto{\pgfqpoint{0.611044in}{0.540017in}}%
\pgfpathlineto{\pgfqpoint{0.610456in}{0.538425in}}%
\pgfpathclose%
\pgfusepath{fill}%
\end{pgfscope}%
\begin{pgfscope}%
\pgfpathrectangle{\pgfqpoint{0.100000in}{0.100000in}}{\pgfqpoint{3.007045in}{1.925000in}}%
\pgfusepath{clip}%
\pgfsetbuttcap%
\pgfsetmiterjoin%
\definecolor{currentfill}{rgb}{0.886505,0.956171,0.866128}%
\pgfsetfillcolor{currentfill}%
\pgfsetlinewidth{0.000000pt}%
\definecolor{currentstroke}{rgb}{0.000000,0.000000,0.000000}%
\pgfsetstrokecolor{currentstroke}%
\pgfsetstrokeopacity{0.000000}%
\pgfsetdash{}{0pt}%
\pgfpathmoveto{\pgfqpoint{0.652105in}{0.476446in}}%
\pgfpathlineto{\pgfqpoint{0.651830in}{0.477944in}}%
\pgfpathlineto{\pgfqpoint{0.649670in}{0.479666in}}%
\pgfpathlineto{\pgfqpoint{0.648834in}{0.482194in}}%
\pgfpathlineto{\pgfqpoint{0.648697in}{0.484544in}}%
\pgfpathlineto{\pgfqpoint{0.647047in}{0.487579in}}%
\pgfpathlineto{\pgfqpoint{0.647342in}{0.491406in}}%
\pgfpathlineto{\pgfqpoint{0.648359in}{0.493435in}}%
\pgfpathlineto{\pgfqpoint{0.649744in}{0.493109in}}%
\pgfpathlineto{\pgfqpoint{0.652856in}{0.488673in}}%
\pgfpathlineto{\pgfqpoint{0.656744in}{0.489158in}}%
\pgfpathlineto{\pgfqpoint{0.659223in}{0.488519in}}%
\pgfpathlineto{\pgfqpoint{0.659628in}{0.487236in}}%
\pgfpathlineto{\pgfqpoint{0.661768in}{0.486250in}}%
\pgfpathlineto{\pgfqpoint{0.661612in}{0.484269in}}%
\pgfpathlineto{\pgfqpoint{0.663795in}{0.482269in}}%
\pgfpathlineto{\pgfqpoint{0.663710in}{0.479661in}}%
\pgfpathlineto{\pgfqpoint{0.661944in}{0.477630in}}%
\pgfpathlineto{\pgfqpoint{0.660195in}{0.476581in}}%
\pgfpathlineto{\pgfqpoint{0.659577in}{0.472416in}}%
\pgfpathlineto{\pgfqpoint{0.657595in}{0.473029in}}%
\pgfpathlineto{\pgfqpoint{0.655097in}{0.475183in}}%
\pgfpathlineto{\pgfqpoint{0.652105in}{0.476446in}}%
\pgfpathclose%
\pgfusepath{fill}%
\end{pgfscope}%
\begin{pgfscope}%
\pgfpathrectangle{\pgfqpoint{0.100000in}{0.100000in}}{\pgfqpoint{3.007045in}{1.925000in}}%
\pgfusepath{clip}%
\pgfsetbuttcap%
\pgfsetmiterjoin%
\definecolor{currentfill}{rgb}{0.886505,0.956171,0.866128}%
\pgfsetfillcolor{currentfill}%
\pgfsetlinewidth{0.000000pt}%
\definecolor{currentstroke}{rgb}{0.000000,0.000000,0.000000}%
\pgfsetstrokecolor{currentstroke}%
\pgfsetstrokeopacity{0.000000}%
\pgfsetdash{}{0pt}%
\pgfpathmoveto{\pgfqpoint{0.809857in}{0.421160in}}%
\pgfpathlineto{\pgfqpoint{0.796605in}{0.403352in}}%
\pgfpathlineto{\pgfqpoint{0.796295in}{0.405189in}}%
\pgfpathlineto{\pgfqpoint{0.792854in}{0.407923in}}%
\pgfpathlineto{\pgfqpoint{0.789615in}{0.403627in}}%
\pgfpathlineto{\pgfqpoint{0.788663in}{0.401187in}}%
\pgfpathlineto{\pgfqpoint{0.784950in}{0.396276in}}%
\pgfpathlineto{\pgfqpoint{0.775612in}{0.403305in}}%
\pgfpathlineto{\pgfqpoint{0.761232in}{0.414538in}}%
\pgfpathlineto{\pgfqpoint{0.750064in}{0.423440in}}%
\pgfpathlineto{\pgfqpoint{0.751507in}{0.425233in}}%
\pgfpathlineto{\pgfqpoint{0.748296in}{0.427861in}}%
\pgfpathlineto{\pgfqpoint{0.746839in}{0.426097in}}%
\pgfpathlineto{\pgfqpoint{0.745138in}{0.427348in}}%
\pgfpathlineto{\pgfqpoint{0.743807in}{0.425686in}}%
\pgfpathlineto{\pgfqpoint{0.742145in}{0.427074in}}%
\pgfpathlineto{\pgfqpoint{0.743447in}{0.428714in}}%
\pgfpathlineto{\pgfqpoint{0.734378in}{0.436156in}}%
\pgfpathlineto{\pgfqpoint{0.733009in}{0.434561in}}%
\pgfpathlineto{\pgfqpoint{0.731767in}{0.435552in}}%
\pgfpathlineto{\pgfqpoint{0.730325in}{0.433863in}}%
\pgfpathlineto{\pgfqpoint{0.728779in}{0.435116in}}%
\pgfpathlineto{\pgfqpoint{0.726974in}{0.433888in}}%
\pgfpathlineto{\pgfqpoint{0.724234in}{0.430706in}}%
\pgfpathlineto{\pgfqpoint{0.722451in}{0.432254in}}%
\pgfpathlineto{\pgfqpoint{0.721076in}{0.430661in}}%
\pgfpathlineto{\pgfqpoint{0.719488in}{0.432032in}}%
\pgfpathlineto{\pgfqpoint{0.716770in}{0.428913in}}%
\pgfpathlineto{\pgfqpoint{0.715184in}{0.430242in}}%
\pgfpathlineto{\pgfqpoint{0.713230in}{0.429162in}}%
\pgfpathlineto{\pgfqpoint{0.710582in}{0.426191in}}%
\pgfpathlineto{\pgfqpoint{0.709062in}{0.427491in}}%
\pgfpathlineto{\pgfqpoint{0.706317in}{0.424353in}}%
\pgfpathlineto{\pgfqpoint{0.704186in}{0.426137in}}%
\pgfpathlineto{\pgfqpoint{0.701434in}{0.423011in}}%
\pgfpathlineto{\pgfqpoint{0.699846in}{0.424401in}}%
\pgfpathlineto{\pgfqpoint{0.697072in}{0.421245in}}%
\pgfpathlineto{\pgfqpoint{0.695145in}{0.420096in}}%
\pgfpathlineto{\pgfqpoint{0.693623in}{0.421487in}}%
\pgfpathlineto{\pgfqpoint{0.691710in}{0.420379in}}%
\pgfpathlineto{\pgfqpoint{0.688938in}{0.417253in}}%
\pgfpathlineto{\pgfqpoint{0.687384in}{0.418587in}}%
\pgfpathlineto{\pgfqpoint{0.686127in}{0.417159in}}%
\pgfpathlineto{\pgfqpoint{0.684613in}{0.418474in}}%
\pgfpathlineto{\pgfqpoint{0.681584in}{0.415593in}}%
\pgfpathlineto{\pgfqpoint{0.681431in}{0.415610in}}%
\pgfpathlineto{\pgfqpoint{0.678710in}{0.418051in}}%
\pgfpathlineto{\pgfqpoint{0.677618in}{0.416243in}}%
\pgfpathlineto{\pgfqpoint{0.676726in}{0.415156in}}%
\pgfpathlineto{\pgfqpoint{0.670768in}{0.416525in}}%
\pgfpathlineto{\pgfqpoint{0.671401in}{0.419366in}}%
\pgfpathlineto{\pgfqpoint{0.674143in}{0.420400in}}%
\pgfpathlineto{\pgfqpoint{0.677137in}{0.423980in}}%
\pgfpathlineto{\pgfqpoint{0.678234in}{0.426008in}}%
\pgfpathlineto{\pgfqpoint{0.678253in}{0.428853in}}%
\pgfpathlineto{\pgfqpoint{0.679315in}{0.431241in}}%
\pgfpathlineto{\pgfqpoint{0.680432in}{0.430540in}}%
\pgfpathlineto{\pgfqpoint{0.681892in}{0.432397in}}%
\pgfpathlineto{\pgfqpoint{0.685663in}{0.432497in}}%
\pgfpathlineto{\pgfqpoint{0.686784in}{0.434521in}}%
\pgfpathlineto{\pgfqpoint{0.688419in}{0.439211in}}%
\pgfpathlineto{\pgfqpoint{0.688889in}{0.441645in}}%
\pgfpathlineto{\pgfqpoint{0.690056in}{0.443550in}}%
\pgfpathlineto{\pgfqpoint{0.690186in}{0.446179in}}%
\pgfpathlineto{\pgfqpoint{0.692680in}{0.447617in}}%
\pgfpathlineto{\pgfqpoint{0.692600in}{0.450868in}}%
\pgfpathlineto{\pgfqpoint{0.690991in}{0.451746in}}%
\pgfpathlineto{\pgfqpoint{0.688110in}{0.448749in}}%
\pgfpathlineto{\pgfqpoint{0.685848in}{0.448554in}}%
\pgfpathlineto{\pgfqpoint{0.684889in}{0.449901in}}%
\pgfpathlineto{\pgfqpoint{0.680145in}{0.450430in}}%
\pgfpathlineto{\pgfqpoint{0.675467in}{0.452772in}}%
\pgfpathlineto{\pgfqpoint{0.673367in}{0.454382in}}%
\pgfpathlineto{\pgfqpoint{0.670995in}{0.457386in}}%
\pgfpathlineto{\pgfqpoint{0.670409in}{0.459079in}}%
\pgfpathlineto{\pgfqpoint{0.671394in}{0.461167in}}%
\pgfpathlineto{\pgfqpoint{0.672846in}{0.461844in}}%
\pgfpathlineto{\pgfqpoint{0.672427in}{0.464293in}}%
\pgfpathlineto{\pgfqpoint{0.672655in}{0.467035in}}%
\pgfpathlineto{\pgfqpoint{0.673992in}{0.467451in}}%
\pgfpathlineto{\pgfqpoint{0.672827in}{0.471052in}}%
\pgfpathlineto{\pgfqpoint{0.673952in}{0.471769in}}%
\pgfpathlineto{\pgfqpoint{0.671748in}{0.473544in}}%
\pgfpathlineto{\pgfqpoint{0.671417in}{0.476039in}}%
\pgfpathlineto{\pgfqpoint{0.672034in}{0.478136in}}%
\pgfpathlineto{\pgfqpoint{0.673606in}{0.477783in}}%
\pgfpathlineto{\pgfqpoint{0.674154in}{0.479291in}}%
\pgfpathlineto{\pgfqpoint{0.671865in}{0.479944in}}%
\pgfpathlineto{\pgfqpoint{0.671055in}{0.482481in}}%
\pgfpathlineto{\pgfqpoint{0.674291in}{0.481708in}}%
\pgfpathlineto{\pgfqpoint{0.678236in}{0.483360in}}%
\pgfpathlineto{\pgfqpoint{0.680888in}{0.483413in}}%
\pgfpathlineto{\pgfqpoint{0.679590in}{0.486507in}}%
\pgfpathlineto{\pgfqpoint{0.680253in}{0.488047in}}%
\pgfpathlineto{\pgfqpoint{0.682631in}{0.488869in}}%
\pgfpathlineto{\pgfqpoint{0.682850in}{0.492468in}}%
\pgfpathlineto{\pgfqpoint{0.680782in}{0.491150in}}%
\pgfpathlineto{\pgfqpoint{0.679720in}{0.494316in}}%
\pgfpathlineto{\pgfqpoint{0.682387in}{0.496907in}}%
\pgfpathlineto{\pgfqpoint{0.680696in}{0.499993in}}%
\pgfpathlineto{\pgfqpoint{0.681929in}{0.501414in}}%
\pgfpathlineto{\pgfqpoint{0.683631in}{0.501161in}}%
\pgfpathlineto{\pgfqpoint{0.684549in}{0.502655in}}%
\pgfpathlineto{\pgfqpoint{0.684007in}{0.504296in}}%
\pgfpathlineto{\pgfqpoint{0.681235in}{0.504565in}}%
\pgfpathlineto{\pgfqpoint{0.682519in}{0.506916in}}%
\pgfpathlineto{\pgfqpoint{0.686294in}{0.507512in}}%
\pgfpathlineto{\pgfqpoint{0.687166in}{0.510160in}}%
\pgfpathlineto{\pgfqpoint{0.689933in}{0.507661in}}%
\pgfpathlineto{\pgfqpoint{0.692175in}{0.508261in}}%
\pgfpathlineto{\pgfqpoint{0.693238in}{0.510668in}}%
\pgfpathlineto{\pgfqpoint{0.695178in}{0.511813in}}%
\pgfpathlineto{\pgfqpoint{0.704163in}{0.514029in}}%
\pgfpathlineto{\pgfqpoint{0.708310in}{0.514388in}}%
\pgfpathlineto{\pgfqpoint{0.709700in}{0.512813in}}%
\pgfpathlineto{\pgfqpoint{0.712038in}{0.513756in}}%
\pgfpathlineto{\pgfqpoint{0.712689in}{0.516016in}}%
\pgfpathlineto{\pgfqpoint{0.717386in}{0.519817in}}%
\pgfpathlineto{\pgfqpoint{0.721882in}{0.521512in}}%
\pgfpathlineto{\pgfqpoint{0.726704in}{0.521249in}}%
\pgfpathlineto{\pgfqpoint{0.729612in}{0.518760in}}%
\pgfpathlineto{\pgfqpoint{0.730444in}{0.514976in}}%
\pgfpathlineto{\pgfqpoint{0.730775in}{0.511205in}}%
\pgfpathlineto{\pgfqpoint{0.731879in}{0.508992in}}%
\pgfpathlineto{\pgfqpoint{0.734192in}{0.507876in}}%
\pgfpathlineto{\pgfqpoint{0.738215in}{0.508943in}}%
\pgfpathlineto{\pgfqpoint{0.741050in}{0.508543in}}%
\pgfpathlineto{\pgfqpoint{0.742738in}{0.506896in}}%
\pgfpathlineto{\pgfqpoint{0.741312in}{0.505279in}}%
\pgfpathlineto{\pgfqpoint{0.742884in}{0.503844in}}%
\pgfpathlineto{\pgfqpoint{0.741434in}{0.502134in}}%
\pgfpathlineto{\pgfqpoint{0.742883in}{0.500799in}}%
\pgfpathlineto{\pgfqpoint{0.741587in}{0.499176in}}%
\pgfpathlineto{\pgfqpoint{0.751165in}{0.490590in}}%
\pgfpathlineto{\pgfqpoint{0.750170in}{0.491012in}}%
\pgfpathlineto{\pgfqpoint{0.748071in}{0.489521in}}%
\pgfpathlineto{\pgfqpoint{0.742309in}{0.483156in}}%
\pgfpathlineto{\pgfqpoint{0.741905in}{0.483472in}}%
\pgfpathlineto{\pgfqpoint{0.736206in}{0.477104in}}%
\pgfpathlineto{\pgfqpoint{0.737384in}{0.476047in}}%
\pgfpathlineto{\pgfqpoint{0.733232in}{0.471329in}}%
\pgfpathlineto{\pgfqpoint{0.741635in}{0.463832in}}%
\pgfpathlineto{\pgfqpoint{0.747517in}{0.458771in}}%
\pgfpathlineto{\pgfqpoint{0.749271in}{0.460813in}}%
\pgfpathlineto{\pgfqpoint{0.755991in}{0.455118in}}%
\pgfpathlineto{\pgfqpoint{0.757361in}{0.456738in}}%
\pgfpathlineto{\pgfqpoint{0.770089in}{0.446056in}}%
\pgfpathlineto{\pgfqpoint{0.768555in}{0.444177in}}%
\pgfpathlineto{\pgfqpoint{0.776021in}{0.437944in}}%
\pgfpathlineto{\pgfqpoint{0.787335in}{0.429069in}}%
\pgfpathlineto{\pgfqpoint{0.796550in}{0.422002in}}%
\pgfpathlineto{\pgfqpoint{0.798438in}{0.420763in}}%
\pgfpathlineto{\pgfqpoint{0.800378in}{0.423260in}}%
\pgfpathlineto{\pgfqpoint{0.803696in}{0.420629in}}%
\pgfpathlineto{\pgfqpoint{0.805001in}{0.422296in}}%
\pgfpathlineto{\pgfqpoint{0.807974in}{0.419972in}}%
\pgfpathlineto{\pgfqpoint{0.809857in}{0.421160in}}%
\pgfpathclose%
\pgfusepath{fill}%
\end{pgfscope}%
\begin{pgfscope}%
\pgfpathrectangle{\pgfqpoint{0.100000in}{0.100000in}}{\pgfqpoint{3.007045in}{1.925000in}}%
\pgfusepath{clip}%
\pgfsetbuttcap%
\pgfsetmiterjoin%
\definecolor{currentfill}{rgb}{0.252595,0.668128,0.362860}%
\pgfsetfillcolor{currentfill}%
\pgfsetlinewidth{0.000000pt}%
\definecolor{currentstroke}{rgb}{0.000000,0.000000,0.000000}%
\pgfsetstrokecolor{currentstroke}%
\pgfsetstrokeopacity{0.000000}%
\pgfsetdash{}{0pt}%
\pgfpathmoveto{\pgfqpoint{0.860069in}{1.641241in}}%
\pgfpathlineto{\pgfqpoint{0.857696in}{1.638685in}}%
\pgfpathlineto{\pgfqpoint{0.853704in}{1.621066in}}%
\pgfpathlineto{\pgfqpoint{0.841346in}{1.613085in}}%
\pgfpathlineto{\pgfqpoint{0.834363in}{1.617027in}}%
\pgfpathlineto{\pgfqpoint{0.830622in}{1.610723in}}%
\pgfpathlineto{\pgfqpoint{0.824810in}{1.609423in}}%
\pgfpathlineto{\pgfqpoint{0.822960in}{1.604421in}}%
\pgfpathlineto{\pgfqpoint{0.825412in}{1.600994in}}%
\pgfpathlineto{\pgfqpoint{0.824171in}{1.596209in}}%
\pgfpathlineto{\pgfqpoint{0.820142in}{1.591987in}}%
\pgfpathlineto{\pgfqpoint{0.811192in}{1.587476in}}%
\pgfpathlineto{\pgfqpoint{0.808189in}{1.587646in}}%
\pgfpathlineto{\pgfqpoint{0.781320in}{1.594007in}}%
\pgfpathlineto{\pgfqpoint{0.779546in}{1.588394in}}%
\pgfpathlineto{\pgfqpoint{0.774977in}{1.589715in}}%
\pgfpathlineto{\pgfqpoint{0.777625in}{1.600841in}}%
\pgfpathlineto{\pgfqpoint{0.780379in}{1.600195in}}%
\pgfpathlineto{\pgfqpoint{0.783037in}{1.611433in}}%
\pgfpathlineto{\pgfqpoint{0.775430in}{1.609514in}}%
\pgfpathlineto{\pgfqpoint{0.770478in}{1.610681in}}%
\pgfpathlineto{\pgfqpoint{0.765709in}{1.615018in}}%
\pgfpathlineto{\pgfqpoint{0.767582in}{1.623021in}}%
\pgfpathlineto{\pgfqpoint{0.763942in}{1.627183in}}%
\pgfpathlineto{\pgfqpoint{0.766659in}{1.638492in}}%
\pgfpathlineto{\pgfqpoint{0.754151in}{1.641671in}}%
\pgfpathlineto{\pgfqpoint{0.758646in}{1.647388in}}%
\pgfpathlineto{\pgfqpoint{0.759381in}{1.652797in}}%
\pgfpathlineto{\pgfqpoint{0.763180in}{1.655715in}}%
\pgfpathlineto{\pgfqpoint{0.766636in}{1.659171in}}%
\pgfpathlineto{\pgfqpoint{0.770539in}{1.666797in}}%
\pgfpathlineto{\pgfqpoint{0.786144in}{1.663045in}}%
\pgfpathlineto{\pgfqpoint{0.786641in}{1.653470in}}%
\pgfpathlineto{\pgfqpoint{0.792690in}{1.650559in}}%
\pgfpathlineto{\pgfqpoint{0.802276in}{1.654267in}}%
\pgfpathlineto{\pgfqpoint{0.860069in}{1.641241in}}%
\pgfpathclose%
\pgfusepath{fill}%
\end{pgfscope}%
\begin{pgfscope}%
\pgfpathrectangle{\pgfqpoint{0.100000in}{0.100000in}}{\pgfqpoint{3.007045in}{1.925000in}}%
\pgfusepath{clip}%
\pgfsetbuttcap%
\pgfsetmiterjoin%
\definecolor{currentfill}{rgb}{0.546251,0.811257,0.537855}%
\pgfsetfillcolor{currentfill}%
\pgfsetlinewidth{0.000000pt}%
\definecolor{currentstroke}{rgb}{0.000000,0.000000,0.000000}%
\pgfsetstrokecolor{currentstroke}%
\pgfsetstrokeopacity{0.000000}%
\pgfsetdash{}{0pt}%
\pgfpathmoveto{\pgfqpoint{2.081022in}{1.010070in}}%
\pgfpathlineto{\pgfqpoint{2.077300in}{1.012074in}}%
\pgfpathlineto{\pgfqpoint{2.072603in}{1.010193in}}%
\pgfpathlineto{\pgfqpoint{2.067443in}{1.014097in}}%
\pgfpathlineto{\pgfqpoint{2.059538in}{1.028875in}}%
\pgfpathlineto{\pgfqpoint{2.064198in}{1.036075in}}%
\pgfpathlineto{\pgfqpoint{2.058254in}{1.047208in}}%
\pgfpathlineto{\pgfqpoint{2.060389in}{1.049493in}}%
\pgfpathlineto{\pgfqpoint{2.058159in}{1.055368in}}%
\pgfpathlineto{\pgfqpoint{2.054473in}{1.056687in}}%
\pgfpathlineto{\pgfqpoint{2.045818in}{1.065000in}}%
\pgfpathlineto{\pgfqpoint{2.040262in}{1.068526in}}%
\pgfpathlineto{\pgfqpoint{2.034604in}{1.066814in}}%
\pgfpathlineto{\pgfqpoint{2.031875in}{1.071974in}}%
\pgfpathlineto{\pgfqpoint{2.024754in}{1.078038in}}%
\pgfpathlineto{\pgfqpoint{2.020734in}{1.079515in}}%
\pgfpathlineto{\pgfqpoint{2.029399in}{1.083124in}}%
\pgfpathlineto{\pgfqpoint{2.029064in}{1.088879in}}%
\pgfpathlineto{\pgfqpoint{2.046290in}{1.089601in}}%
\pgfpathlineto{\pgfqpoint{2.075029in}{1.090942in}}%
\pgfpathlineto{\pgfqpoint{2.075282in}{1.085153in}}%
\pgfpathlineto{\pgfqpoint{2.098280in}{1.086724in}}%
\pgfpathlineto{\pgfqpoint{2.099179in}{1.072352in}}%
\pgfpathlineto{\pgfqpoint{2.100489in}{1.052130in}}%
\pgfpathlineto{\pgfqpoint{2.101574in}{1.034888in}}%
\pgfpathlineto{\pgfqpoint{2.091233in}{1.034083in}}%
\pgfpathlineto{\pgfqpoint{2.090698in}{1.026843in}}%
\pgfpathlineto{\pgfqpoint{2.085427in}{1.025478in}}%
\pgfpathlineto{\pgfqpoint{2.078167in}{1.013748in}}%
\pgfpathlineto{\pgfqpoint{2.081022in}{1.010070in}}%
\pgfpathclose%
\pgfusepath{fill}%
\end{pgfscope}%
\begin{pgfscope}%
\pgfpathrectangle{\pgfqpoint{0.100000in}{0.100000in}}{\pgfqpoint{3.007045in}{1.925000in}}%
\pgfusepath{clip}%
\pgfsetbuttcap%
\pgfsetmiterjoin%
\definecolor{currentfill}{rgb}{0.568397,0.821592,0.556063}%
\pgfsetfillcolor{currentfill}%
\pgfsetlinewidth{0.000000pt}%
\definecolor{currentstroke}{rgb}{0.000000,0.000000,0.000000}%
\pgfsetstrokecolor{currentstroke}%
\pgfsetstrokeopacity{0.000000}%
\pgfsetdash{}{0pt}%
\pgfpathmoveto{\pgfqpoint{2.721216in}{1.279405in}}%
\pgfpathlineto{\pgfqpoint{2.716371in}{1.278367in}}%
\pgfpathlineto{\pgfqpoint{2.716037in}{1.278304in}}%
\pgfpathlineto{\pgfqpoint{2.706007in}{1.289043in}}%
\pgfpathlineto{\pgfqpoint{2.701481in}{1.290448in}}%
\pgfpathlineto{\pgfqpoint{2.697015in}{1.296543in}}%
\pgfpathlineto{\pgfqpoint{2.690460in}{1.296019in}}%
\pgfpathlineto{\pgfqpoint{2.687398in}{1.300784in}}%
\pgfpathlineto{\pgfqpoint{2.693385in}{1.305552in}}%
\pgfpathlineto{\pgfqpoint{2.684366in}{1.322430in}}%
\pgfpathlineto{\pgfqpoint{2.690299in}{1.329057in}}%
\pgfpathlineto{\pgfqpoint{2.680799in}{1.334029in}}%
\pgfpathlineto{\pgfqpoint{2.694956in}{1.344848in}}%
\pgfpathlineto{\pgfqpoint{2.698167in}{1.347287in}}%
\pgfpathlineto{\pgfqpoint{2.698167in}{1.352791in}}%
\pgfpathlineto{\pgfqpoint{2.701015in}{1.357842in}}%
\pgfpathlineto{\pgfqpoint{2.711768in}{1.357635in}}%
\pgfpathlineto{\pgfqpoint{2.725895in}{1.348689in}}%
\pgfpathlineto{\pgfqpoint{2.720167in}{1.343542in}}%
\pgfpathlineto{\pgfqpoint{2.741068in}{1.332531in}}%
\pgfpathlineto{\pgfqpoint{2.735728in}{1.317531in}}%
\pgfpathlineto{\pgfqpoint{2.725428in}{1.306338in}}%
\pgfpathlineto{\pgfqpoint{2.726798in}{1.302007in}}%
\pgfpathlineto{\pgfqpoint{2.725049in}{1.295828in}}%
\pgfpathlineto{\pgfqpoint{2.725778in}{1.289195in}}%
\pgfpathlineto{\pgfqpoint{2.724100in}{1.282597in}}%
\pgfpathlineto{\pgfqpoint{2.721216in}{1.279405in}}%
\pgfpathclose%
\pgfusepath{fill}%
\end{pgfscope}%
\begin{pgfscope}%
\pgfpathrectangle{\pgfqpoint{0.100000in}{0.100000in}}{\pgfqpoint{3.007045in}{1.925000in}}%
\pgfusepath{clip}%
\pgfsetbuttcap%
\pgfsetmiterjoin%
\definecolor{currentfill}{rgb}{0.732472,0.893549,0.706282}%
\pgfsetfillcolor{currentfill}%
\pgfsetlinewidth{0.000000pt}%
\definecolor{currentstroke}{rgb}{0.000000,0.000000,0.000000}%
\pgfsetstrokecolor{currentstroke}%
\pgfsetstrokeopacity{0.000000}%
\pgfsetdash{}{0pt}%
\pgfpathmoveto{\pgfqpoint{2.422704in}{0.621971in}}%
\pgfpathlineto{\pgfqpoint{2.415134in}{0.621502in}}%
\pgfpathlineto{\pgfqpoint{2.412441in}{0.646597in}}%
\pgfpathlineto{\pgfqpoint{2.397277in}{0.644995in}}%
\pgfpathlineto{\pgfqpoint{2.397002in}{0.647345in}}%
\pgfpathlineto{\pgfqpoint{2.395188in}{0.664380in}}%
\pgfpathlineto{\pgfqpoint{2.414666in}{0.666517in}}%
\pgfpathlineto{\pgfqpoint{2.412962in}{0.682194in}}%
\pgfpathlineto{\pgfqpoint{2.408654in}{0.685257in}}%
\pgfpathlineto{\pgfqpoint{2.404240in}{0.684768in}}%
\pgfpathlineto{\pgfqpoint{2.402438in}{0.696626in}}%
\pgfpathlineto{\pgfqpoint{2.413072in}{0.697996in}}%
\pgfpathlineto{\pgfqpoint{2.412679in}{0.701290in}}%
\pgfpathlineto{\pgfqpoint{2.420123in}{0.701779in}}%
\pgfpathlineto{\pgfqpoint{2.422816in}{0.692785in}}%
\pgfpathlineto{\pgfqpoint{2.420802in}{0.691745in}}%
\pgfpathlineto{\pgfqpoint{2.421279in}{0.683375in}}%
\pgfpathlineto{\pgfqpoint{2.425943in}{0.682862in}}%
\pgfpathlineto{\pgfqpoint{2.428996in}{0.678445in}}%
\pgfpathlineto{\pgfqpoint{2.442052in}{0.679496in}}%
\pgfpathlineto{\pgfqpoint{2.443165in}{0.674743in}}%
\pgfpathlineto{\pgfqpoint{2.450078in}{0.667343in}}%
\pgfpathlineto{\pgfqpoint{2.450133in}{0.661231in}}%
\pgfpathlineto{\pgfqpoint{2.454390in}{0.661837in}}%
\pgfpathlineto{\pgfqpoint{2.457148in}{0.641063in}}%
\pgfpathlineto{\pgfqpoint{2.464523in}{0.639748in}}%
\pgfpathlineto{\pgfqpoint{2.469936in}{0.633604in}}%
\pgfpathlineto{\pgfqpoint{2.479352in}{0.633603in}}%
\pgfpathlineto{\pgfqpoint{2.481583in}{0.625736in}}%
\pgfpathlineto{\pgfqpoint{2.422704in}{0.621971in}}%
\pgfpathclose%
\pgfusepath{fill}%
\end{pgfscope}%
\begin{pgfscope}%
\pgfpathrectangle{\pgfqpoint{0.100000in}{0.100000in}}{\pgfqpoint{3.007045in}{1.925000in}}%
\pgfusepath{clip}%
\pgfsetbuttcap%
\pgfsetmiterjoin%
\definecolor{currentfill}{rgb}{0.573933,0.824175,0.560615}%
\pgfsetfillcolor{currentfill}%
\pgfsetlinewidth{0.000000pt}%
\definecolor{currentstroke}{rgb}{0.000000,0.000000,0.000000}%
\pgfsetstrokecolor{currentstroke}%
\pgfsetstrokeopacity{0.000000}%
\pgfsetdash{}{0pt}%
\pgfpathmoveto{\pgfqpoint{2.030481in}{0.763778in}}%
\pgfpathlineto{\pgfqpoint{2.023847in}{0.761087in}}%
\pgfpathlineto{\pgfqpoint{2.023674in}{0.765200in}}%
\pgfpathlineto{\pgfqpoint{2.022165in}{0.797050in}}%
\pgfpathlineto{\pgfqpoint{2.036938in}{0.797675in}}%
\pgfpathlineto{\pgfqpoint{2.030424in}{0.792915in}}%
\pgfpathlineto{\pgfqpoint{2.039761in}{0.792018in}}%
\pgfpathlineto{\pgfqpoint{2.039467in}{0.797809in}}%
\pgfpathlineto{\pgfqpoint{2.050947in}{0.799400in}}%
\pgfpathlineto{\pgfqpoint{2.050685in}{0.804294in}}%
\pgfpathlineto{\pgfqpoint{2.059401in}{0.802761in}}%
\pgfpathlineto{\pgfqpoint{2.073753in}{0.803617in}}%
\pgfpathlineto{\pgfqpoint{2.074344in}{0.794001in}}%
\pgfpathlineto{\pgfqpoint{2.075429in}{0.776700in}}%
\pgfpathlineto{\pgfqpoint{2.078426in}{0.775122in}}%
\pgfpathlineto{\pgfqpoint{2.079034in}{0.765390in}}%
\pgfpathlineto{\pgfqpoint{2.065349in}{0.764544in}}%
\pgfpathlineto{\pgfqpoint{2.062975in}{0.759697in}}%
\pgfpathlineto{\pgfqpoint{2.039047in}{0.766251in}}%
\pgfpathlineto{\pgfqpoint{2.033226in}{0.766351in}}%
\pgfpathlineto{\pgfqpoint{2.030481in}{0.763778in}}%
\pgfpathclose%
\pgfusepath{fill}%
\end{pgfscope}%
\begin{pgfscope}%
\pgfpathrectangle{\pgfqpoint{0.100000in}{0.100000in}}{\pgfqpoint{3.007045in}{1.925000in}}%
\pgfusepath{clip}%
\pgfsetbuttcap%
\pgfsetmiterjoin%
\definecolor{currentfill}{rgb}{0.842215,0.938454,0.818885}%
\pgfsetfillcolor{currentfill}%
\pgfsetlinewidth{0.000000pt}%
\definecolor{currentstroke}{rgb}{0.000000,0.000000,0.000000}%
\pgfsetstrokecolor{currentstroke}%
\pgfsetstrokeopacity{0.000000}%
\pgfsetdash{}{0pt}%
\pgfpathmoveto{\pgfqpoint{1.611326in}{0.679134in}}%
\pgfpathlineto{\pgfqpoint{1.607101in}{0.688968in}}%
\pgfpathlineto{\pgfqpoint{1.600832in}{0.707630in}}%
\pgfpathlineto{\pgfqpoint{1.578119in}{0.708542in}}%
\pgfpathlineto{\pgfqpoint{1.572479in}{0.708902in}}%
\pgfpathlineto{\pgfqpoint{1.573659in}{0.737980in}}%
\pgfpathlineto{\pgfqpoint{1.581896in}{0.737526in}}%
\pgfpathlineto{\pgfqpoint{1.582071in}{0.741206in}}%
\pgfpathlineto{\pgfqpoint{1.610081in}{0.739788in}}%
\pgfpathlineto{\pgfqpoint{1.611193in}{0.768407in}}%
\pgfpathlineto{\pgfqpoint{1.607913in}{0.770715in}}%
\pgfpathlineto{\pgfqpoint{1.608887in}{0.798669in}}%
\pgfpathlineto{\pgfqpoint{1.614953in}{0.795791in}}%
\pgfpathlineto{\pgfqpoint{1.625906in}{0.804863in}}%
\pgfpathlineto{\pgfqpoint{1.631785in}{0.798613in}}%
\pgfpathlineto{\pgfqpoint{1.636022in}{0.799536in}}%
\pgfpathlineto{\pgfqpoint{1.634992in}{0.767696in}}%
\pgfpathlineto{\pgfqpoint{1.640734in}{0.767480in}}%
\pgfpathlineto{\pgfqpoint{1.639106in}{0.738203in}}%
\pgfpathlineto{\pgfqpoint{1.659429in}{0.737338in}}%
\pgfpathlineto{\pgfqpoint{1.658448in}{0.708465in}}%
\pgfpathlineto{\pgfqpoint{1.655743in}{0.708579in}}%
\pgfpathlineto{\pgfqpoint{1.655346in}{0.689754in}}%
\pgfpathlineto{\pgfqpoint{1.633343in}{0.684748in}}%
\pgfpathlineto{\pgfqpoint{1.628904in}{0.682713in}}%
\pgfpathlineto{\pgfqpoint{1.624705in}{0.686669in}}%
\pgfpathlineto{\pgfqpoint{1.611326in}{0.679134in}}%
\pgfpathclose%
\pgfusepath{fill}%
\end{pgfscope}%
\begin{pgfscope}%
\pgfpathrectangle{\pgfqpoint{0.100000in}{0.100000in}}{\pgfqpoint{3.007045in}{1.925000in}}%
\pgfusepath{clip}%
\pgfsetbuttcap%
\pgfsetmiterjoin%
\definecolor{currentfill}{rgb}{0.607151,0.839677,0.587928}%
\pgfsetfillcolor{currentfill}%
\pgfsetlinewidth{0.000000pt}%
\definecolor{currentstroke}{rgb}{0.000000,0.000000,0.000000}%
\pgfsetstrokecolor{currentstroke}%
\pgfsetstrokeopacity{0.000000}%
\pgfsetdash{}{0pt}%
\pgfpathmoveto{\pgfqpoint{2.265651in}{0.893385in}}%
\pgfpathlineto{\pgfqpoint{2.241623in}{0.891209in}}%
\pgfpathlineto{\pgfqpoint{2.240394in}{0.900117in}}%
\pgfpathlineto{\pgfqpoint{2.237551in}{0.900340in}}%
\pgfpathlineto{\pgfqpoint{2.235740in}{0.906849in}}%
\pgfpathlineto{\pgfqpoint{2.223823in}{0.914321in}}%
\pgfpathlineto{\pgfqpoint{2.220248in}{0.918653in}}%
\pgfpathlineto{\pgfqpoint{2.219736in}{0.935214in}}%
\pgfpathlineto{\pgfqpoint{2.226333in}{0.935993in}}%
\pgfpathlineto{\pgfqpoint{2.228851in}{0.933436in}}%
\pgfpathlineto{\pgfqpoint{2.240739in}{0.937396in}}%
\pgfpathlineto{\pgfqpoint{2.251280in}{0.935449in}}%
\pgfpathlineto{\pgfqpoint{2.254813in}{0.936868in}}%
\pgfpathlineto{\pgfqpoint{2.257324in}{0.931388in}}%
\pgfpathlineto{\pgfqpoint{2.261477in}{0.928516in}}%
\pgfpathlineto{\pgfqpoint{2.262417in}{0.916982in}}%
\pgfpathlineto{\pgfqpoint{2.261069in}{0.912869in}}%
\pgfpathlineto{\pgfqpoint{2.264140in}{0.908633in}}%
\pgfpathlineto{\pgfqpoint{2.265651in}{0.893385in}}%
\pgfpathclose%
\pgfusepath{fill}%
\end{pgfscope}%
\begin{pgfscope}%
\pgfpathrectangle{\pgfqpoint{0.100000in}{0.100000in}}{\pgfqpoint{3.007045in}{1.925000in}}%
\pgfusepath{clip}%
\pgfsetbuttcap%
\pgfsetmiterjoin%
\definecolor{currentfill}{rgb}{0.643645,0.856148,0.619792}%
\pgfsetfillcolor{currentfill}%
\pgfsetlinewidth{0.000000pt}%
\definecolor{currentstroke}{rgb}{0.000000,0.000000,0.000000}%
\pgfsetstrokecolor{currentstroke}%
\pgfsetstrokeopacity{0.000000}%
\pgfsetdash{}{0pt}%
\pgfpathmoveto{\pgfqpoint{2.141820in}{0.745927in}}%
\pgfpathlineto{\pgfqpoint{2.116053in}{0.743899in}}%
\pgfpathlineto{\pgfqpoint{2.114414in}{0.767709in}}%
\pgfpathlineto{\pgfqpoint{2.122360in}{0.768291in}}%
\pgfpathlineto{\pgfqpoint{2.121230in}{0.782758in}}%
\pgfpathlineto{\pgfqpoint{2.128772in}{0.785353in}}%
\pgfpathlineto{\pgfqpoint{2.131570in}{0.789934in}}%
\pgfpathlineto{\pgfqpoint{2.129211in}{0.792870in}}%
\pgfpathlineto{\pgfqpoint{2.136914in}{0.795797in}}%
\pgfpathlineto{\pgfqpoint{2.139121in}{0.799917in}}%
\pgfpathlineto{\pgfqpoint{2.143299in}{0.800190in}}%
\pgfpathlineto{\pgfqpoint{2.142921in}{0.786166in}}%
\pgfpathlineto{\pgfqpoint{2.166832in}{0.787387in}}%
\pgfpathlineto{\pgfqpoint{2.168941in}{0.762907in}}%
\pgfpathlineto{\pgfqpoint{2.163129in}{0.759429in}}%
\pgfpathlineto{\pgfqpoint{2.158791in}{0.754738in}}%
\pgfpathlineto{\pgfqpoint{2.154101in}{0.754354in}}%
\pgfpathlineto{\pgfqpoint{2.151248in}{0.751063in}}%
\pgfpathlineto{\pgfqpoint{2.141932in}{0.750061in}}%
\pgfpathlineto{\pgfqpoint{2.141820in}{0.745927in}}%
\pgfpathclose%
\pgfusepath{fill}%
\end{pgfscope}%
\begin{pgfscope}%
\pgfpathrectangle{\pgfqpoint{0.100000in}{0.100000in}}{\pgfqpoint{3.007045in}{1.925000in}}%
\pgfusepath{clip}%
\pgfsetbuttcap%
\pgfsetmiterjoin%
\definecolor{currentfill}{rgb}{0.474279,0.777670,0.478677}%
\pgfsetfillcolor{currentfill}%
\pgfsetlinewidth{0.000000pt}%
\definecolor{currentstroke}{rgb}{0.000000,0.000000,0.000000}%
\pgfsetstrokecolor{currentstroke}%
\pgfsetstrokeopacity{0.000000}%
\pgfsetdash{}{0pt}%
\pgfpathmoveto{\pgfqpoint{2.695962in}{1.404095in}}%
\pgfpathlineto{\pgfqpoint{2.689741in}{1.426141in}}%
\pgfpathlineto{\pgfqpoint{2.669988in}{1.422124in}}%
\pgfpathlineto{\pgfqpoint{2.667190in}{1.440512in}}%
\pgfpathlineto{\pgfqpoint{2.672580in}{1.444093in}}%
\pgfpathlineto{\pgfqpoint{2.680763in}{1.444336in}}%
\pgfpathlineto{\pgfqpoint{2.679109in}{1.451464in}}%
\pgfpathlineto{\pgfqpoint{2.697553in}{1.455914in}}%
\pgfpathlineto{\pgfqpoint{2.700669in}{1.446056in}}%
\pgfpathlineto{\pgfqpoint{2.710631in}{1.447455in}}%
\pgfpathlineto{\pgfqpoint{2.711428in}{1.444046in}}%
\pgfpathlineto{\pgfqpoint{2.721879in}{1.446280in}}%
\pgfpathlineto{\pgfqpoint{2.723883in}{1.436359in}}%
\pgfpathlineto{\pgfqpoint{2.727454in}{1.434295in}}%
\pgfpathlineto{\pgfqpoint{2.721527in}{1.432985in}}%
\pgfpathlineto{\pgfqpoint{2.727566in}{1.410162in}}%
\pgfpathlineto{\pgfqpoint{2.695962in}{1.404095in}}%
\pgfpathclose%
\pgfusepath{fill}%
\end{pgfscope}%
\begin{pgfscope}%
\pgfpathrectangle{\pgfqpoint{0.100000in}{0.100000in}}{\pgfqpoint{3.007045in}{1.925000in}}%
\pgfusepath{clip}%
\pgfsetbuttcap%
\pgfsetmiterjoin%
\definecolor{currentfill}{rgb}{0.618224,0.844844,0.597032}%
\pgfsetfillcolor{currentfill}%
\pgfsetlinewidth{0.000000pt}%
\definecolor{currentstroke}{rgb}{0.000000,0.000000,0.000000}%
\pgfsetstrokecolor{currentstroke}%
\pgfsetstrokeopacity{0.000000}%
\pgfsetdash{}{0pt}%
\pgfpathmoveto{\pgfqpoint{1.700709in}{0.760907in}}%
\pgfpathlineto{\pgfqpoint{1.695830in}{0.764016in}}%
\pgfpathlineto{\pgfqpoint{1.670614in}{0.764978in}}%
\pgfpathlineto{\pgfqpoint{1.664958in}{0.765812in}}%
\pgfpathlineto{\pgfqpoint{1.665670in}{0.801052in}}%
\pgfpathlineto{\pgfqpoint{1.668937in}{0.795613in}}%
\pgfpathlineto{\pgfqpoint{1.678257in}{0.793022in}}%
\pgfpathlineto{\pgfqpoint{1.683354in}{0.804224in}}%
\pgfpathlineto{\pgfqpoint{1.695332in}{0.814267in}}%
\pgfpathlineto{\pgfqpoint{1.718089in}{0.813933in}}%
\pgfpathlineto{\pgfqpoint{1.720987in}{0.811838in}}%
\pgfpathlineto{\pgfqpoint{1.720879in}{0.800114in}}%
\pgfpathlineto{\pgfqpoint{1.726173in}{0.795162in}}%
\pgfpathlineto{\pgfqpoint{1.730486in}{0.795042in}}%
\pgfpathlineto{\pgfqpoint{1.725337in}{0.792070in}}%
\pgfpathlineto{\pgfqpoint{1.724887in}{0.764449in}}%
\pgfpathlineto{\pgfqpoint{1.700709in}{0.760907in}}%
\pgfpathclose%
\pgfusepath{fill}%
\end{pgfscope}%
\begin{pgfscope}%
\pgfpathrectangle{\pgfqpoint{0.100000in}{0.100000in}}{\pgfqpoint{3.007045in}{1.925000in}}%
\pgfusepath{clip}%
\pgfsetbuttcap%
\pgfsetmiterjoin%
\definecolor{currentfill}{rgb}{0.579469,0.826759,0.565167}%
\pgfsetfillcolor{currentfill}%
\pgfsetlinewidth{0.000000pt}%
\definecolor{currentstroke}{rgb}{0.000000,0.000000,0.000000}%
\pgfsetstrokecolor{currentstroke}%
\pgfsetstrokeopacity{0.000000}%
\pgfsetdash{}{0pt}%
\pgfpathmoveto{\pgfqpoint{2.151467in}{0.577078in}}%
\pgfpathlineto{\pgfqpoint{2.146781in}{0.617810in}}%
\pgfpathlineto{\pgfqpoint{2.143322in}{0.646719in}}%
\pgfpathlineto{\pgfqpoint{2.141224in}{0.664018in}}%
\pgfpathlineto{\pgfqpoint{2.162408in}{0.665683in}}%
\pgfpathlineto{\pgfqpoint{2.161161in}{0.669683in}}%
\pgfpathlineto{\pgfqpoint{2.156603in}{0.673169in}}%
\pgfpathlineto{\pgfqpoint{2.161771in}{0.684999in}}%
\pgfpathlineto{\pgfqpoint{2.184529in}{0.686866in}}%
\pgfpathlineto{\pgfqpoint{2.186978in}{0.688005in}}%
\pgfpathlineto{\pgfqpoint{2.186226in}{0.696389in}}%
\pgfpathlineto{\pgfqpoint{2.191937in}{0.696888in}}%
\pgfpathlineto{\pgfqpoint{2.191190in}{0.704948in}}%
\pgfpathlineto{\pgfqpoint{2.193997in}{0.705851in}}%
\pgfpathlineto{\pgfqpoint{2.202255in}{0.702095in}}%
\pgfpathlineto{\pgfqpoint{2.211673in}{0.692874in}}%
\pgfpathlineto{\pgfqpoint{2.226858in}{0.694238in}}%
\pgfpathlineto{\pgfqpoint{2.228144in}{0.679871in}}%
\pgfpathlineto{\pgfqpoint{2.194774in}{0.676909in}}%
\pgfpathlineto{\pgfqpoint{2.194637in}{0.668142in}}%
\pgfpathlineto{\pgfqpoint{2.191786in}{0.667860in}}%
\pgfpathlineto{\pgfqpoint{2.192871in}{0.650628in}}%
\pgfpathlineto{\pgfqpoint{2.187693in}{0.649271in}}%
\pgfpathlineto{\pgfqpoint{2.184550in}{0.645723in}}%
\pgfpathlineto{\pgfqpoint{2.185584in}{0.641031in}}%
\pgfpathlineto{\pgfqpoint{2.191464in}{0.637662in}}%
\pgfpathlineto{\pgfqpoint{2.193761in}{0.621372in}}%
\pgfpathlineto{\pgfqpoint{2.192424in}{0.612482in}}%
\pgfpathlineto{\pgfqpoint{2.199079in}{0.604683in}}%
\pgfpathlineto{\pgfqpoint{2.206456in}{0.600906in}}%
\pgfpathlineto{\pgfqpoint{2.207470in}{0.596991in}}%
\pgfpathlineto{\pgfqpoint{2.205014in}{0.590892in}}%
\pgfpathlineto{\pgfqpoint{2.210055in}{0.585128in}}%
\pgfpathlineto{\pgfqpoint{2.206102in}{0.581924in}}%
\pgfpathlineto{\pgfqpoint{2.202262in}{0.574181in}}%
\pgfpathlineto{\pgfqpoint{2.194379in}{0.571471in}}%
\pgfpathlineto{\pgfqpoint{2.188678in}{0.573144in}}%
\pgfpathlineto{\pgfqpoint{2.179172in}{0.581268in}}%
\pgfpathlineto{\pgfqpoint{2.177184in}{0.585669in}}%
\pgfpathlineto{\pgfqpoint{2.178736in}{0.589471in}}%
\pgfpathlineto{\pgfqpoint{2.176721in}{0.596910in}}%
\pgfpathlineto{\pgfqpoint{2.173439in}{0.598944in}}%
\pgfpathlineto{\pgfqpoint{2.169431in}{0.595956in}}%
\pgfpathlineto{\pgfqpoint{2.167717in}{0.586842in}}%
\pgfpathlineto{\pgfqpoint{2.168005in}{0.577878in}}%
\pgfpathlineto{\pgfqpoint{2.166722in}{0.574427in}}%
\pgfpathlineto{\pgfqpoint{2.154103in}{0.578631in}}%
\pgfpathlineto{\pgfqpoint{2.151467in}{0.577078in}}%
\pgfpathclose%
\pgfusepath{fill}%
\end{pgfscope}%
\begin{pgfscope}%
\pgfpathrectangle{\pgfqpoint{0.100000in}{0.100000in}}{\pgfqpoint{3.007045in}{1.925000in}}%
\pgfusepath{clip}%
\pgfsetbuttcap%
\pgfsetmiterjoin%
\definecolor{currentfill}{rgb}{0.623760,0.847428,0.601584}%
\pgfsetfillcolor{currentfill}%
\pgfsetlinewidth{0.000000pt}%
\definecolor{currentstroke}{rgb}{0.000000,0.000000,0.000000}%
\pgfsetstrokecolor{currentstroke}%
\pgfsetstrokeopacity{0.000000}%
\pgfsetdash{}{0pt}%
\pgfpathmoveto{\pgfqpoint{1.217220in}{0.903588in}}%
\pgfpathlineto{\pgfqpoint{1.214018in}{0.875058in}}%
\pgfpathlineto{\pgfqpoint{1.212799in}{0.875172in}}%
\pgfpathlineto{\pgfqpoint{1.210832in}{0.858083in}}%
\pgfpathlineto{\pgfqpoint{1.210171in}{0.852394in}}%
\pgfpathlineto{\pgfqpoint{1.177035in}{0.856282in}}%
\pgfpathlineto{\pgfqpoint{1.150364in}{0.859555in}}%
\pgfpathlineto{\pgfqpoint{1.151795in}{0.871370in}}%
\pgfpathlineto{\pgfqpoint{1.133587in}{0.879798in}}%
\pgfpathlineto{\pgfqpoint{1.127835in}{0.883763in}}%
\pgfpathlineto{\pgfqpoint{1.082425in}{0.889855in}}%
\pgfpathlineto{\pgfqpoint{1.047397in}{0.894955in}}%
\pgfpathlineto{\pgfqpoint{1.011345in}{0.900484in}}%
\pgfpathlineto{\pgfqpoint{1.015230in}{0.925309in}}%
\pgfpathlineto{\pgfqpoint{1.046069in}{0.920555in}}%
\pgfpathlineto{\pgfqpoint{1.049496in}{0.943259in}}%
\pgfpathlineto{\pgfqpoint{1.111260in}{0.934336in}}%
\pgfpathlineto{\pgfqpoint{1.117536in}{0.979664in}}%
\pgfpathlineto{\pgfqpoint{1.100725in}{0.982093in}}%
\pgfpathlineto{\pgfqpoint{1.102849in}{0.996413in}}%
\pgfpathlineto{\pgfqpoint{1.141795in}{0.990914in}}%
\pgfpathlineto{\pgfqpoint{1.139908in}{0.976686in}}%
\pgfpathlineto{\pgfqpoint{1.173700in}{0.972532in}}%
\pgfpathlineto{\pgfqpoint{1.165289in}{0.970829in}}%
\pgfpathlineto{\pgfqpoint{1.164202in}{0.962160in}}%
\pgfpathlineto{\pgfqpoint{1.171643in}{0.956432in}}%
\pgfpathlineto{\pgfqpoint{1.166043in}{0.909544in}}%
\pgfpathlineto{\pgfqpoint{1.194470in}{0.906193in}}%
\pgfpathlineto{\pgfqpoint{1.217220in}{0.903588in}}%
\pgfpathclose%
\pgfusepath{fill}%
\end{pgfscope}%
\begin{pgfscope}%
\pgfpathrectangle{\pgfqpoint{0.100000in}{0.100000in}}{\pgfqpoint{3.007045in}{1.925000in}}%
\pgfusepath{clip}%
\pgfsetbuttcap%
\pgfsetmiterjoin%
\definecolor{currentfill}{rgb}{0.868051,0.948789,0.846444}%
\pgfsetfillcolor{currentfill}%
\pgfsetlinewidth{0.000000pt}%
\definecolor{currentstroke}{rgb}{0.000000,0.000000,0.000000}%
\pgfsetstrokecolor{currentstroke}%
\pgfsetstrokeopacity{0.000000}%
\pgfsetdash{}{0pt}%
\pgfpathmoveto{\pgfqpoint{1.677967in}{0.558409in}}%
\pgfpathlineto{\pgfqpoint{1.672046in}{0.560070in}}%
\pgfpathlineto{\pgfqpoint{1.665066in}{0.567789in}}%
\pgfpathlineto{\pgfqpoint{1.662849in}{0.573743in}}%
\pgfpathlineto{\pgfqpoint{1.659758in}{0.576115in}}%
\pgfpathlineto{\pgfqpoint{1.679583in}{0.587280in}}%
\pgfpathlineto{\pgfqpoint{1.678179in}{0.593593in}}%
\pgfpathlineto{\pgfqpoint{1.672091in}{0.600546in}}%
\pgfpathlineto{\pgfqpoint{1.670600in}{0.609460in}}%
\pgfpathlineto{\pgfqpoint{1.668209in}{0.612458in}}%
\pgfpathlineto{\pgfqpoint{1.681563in}{0.619854in}}%
\pgfpathlineto{\pgfqpoint{1.697477in}{0.628654in}}%
\pgfpathlineto{\pgfqpoint{1.696915in}{0.623980in}}%
\pgfpathlineto{\pgfqpoint{1.701562in}{0.603092in}}%
\pgfpathlineto{\pgfqpoint{1.705582in}{0.593025in}}%
\pgfpathlineto{\pgfqpoint{1.723044in}{0.595566in}}%
\pgfpathlineto{\pgfqpoint{1.724776in}{0.579988in}}%
\pgfpathlineto{\pgfqpoint{1.726035in}{0.554362in}}%
\pgfpathlineto{\pgfqpoint{1.709639in}{0.553402in}}%
\pgfpathlineto{\pgfqpoint{1.709428in}{0.556522in}}%
\pgfpathlineto{\pgfqpoint{1.704636in}{0.564588in}}%
\pgfpathlineto{\pgfqpoint{1.696897in}{0.563937in}}%
\pgfpathlineto{\pgfqpoint{1.690544in}{0.561282in}}%
\pgfpathlineto{\pgfqpoint{1.677967in}{0.558409in}}%
\pgfpathclose%
\pgfusepath{fill}%
\end{pgfscope}%
\begin{pgfscope}%
\pgfpathrectangle{\pgfqpoint{0.100000in}{0.100000in}}{\pgfqpoint{3.007045in}{1.925000in}}%
\pgfusepath{clip}%
\pgfsetbuttcap%
\pgfsetmiterjoin%
\definecolor{currentfill}{rgb}{0.351373,0.717878,0.411995}%
\pgfsetfillcolor{currentfill}%
\pgfsetlinewidth{0.000000pt}%
\definecolor{currentstroke}{rgb}{0.000000,0.000000,0.000000}%
\pgfsetstrokecolor{currentstroke}%
\pgfsetstrokeopacity{0.000000}%
\pgfsetdash{}{0pt}%
\pgfpathmoveto{\pgfqpoint{2.509239in}{1.087027in}}%
\pgfpathlineto{\pgfqpoint{2.495503in}{1.095488in}}%
\pgfpathlineto{\pgfqpoint{2.490450in}{1.093853in}}%
\pgfpathlineto{\pgfqpoint{2.489686in}{1.089390in}}%
\pgfpathlineto{\pgfqpoint{2.484474in}{1.093009in}}%
\pgfpathlineto{\pgfqpoint{2.476705in}{1.103396in}}%
\pgfpathlineto{\pgfqpoint{2.471855in}{1.105407in}}%
\pgfpathlineto{\pgfqpoint{2.467580in}{1.113861in}}%
\pgfpathlineto{\pgfqpoint{2.470294in}{1.118309in}}%
\pgfpathlineto{\pgfqpoint{2.476974in}{1.118028in}}%
\pgfpathlineto{\pgfqpoint{2.477515in}{1.130920in}}%
\pgfpathlineto{\pgfqpoint{2.481701in}{1.138352in}}%
\pgfpathlineto{\pgfqpoint{2.490413in}{1.135206in}}%
\pgfpathlineto{\pgfqpoint{2.494710in}{1.137649in}}%
\pgfpathlineto{\pgfqpoint{2.501225in}{1.130708in}}%
\pgfpathlineto{\pgfqpoint{2.510167in}{1.130498in}}%
\pgfpathlineto{\pgfqpoint{2.522292in}{1.144995in}}%
\pgfpathlineto{\pgfqpoint{2.526634in}{1.142930in}}%
\pgfpathlineto{\pgfqpoint{2.527646in}{1.135652in}}%
\pgfpathlineto{\pgfqpoint{2.533411in}{1.132073in}}%
\pgfpathlineto{\pgfqpoint{2.538421in}{1.132193in}}%
\pgfpathlineto{\pgfqpoint{2.548827in}{1.135761in}}%
\pgfpathlineto{\pgfqpoint{2.547549in}{1.130953in}}%
\pgfpathlineto{\pgfqpoint{2.538295in}{1.117062in}}%
\pgfpathlineto{\pgfqpoint{2.535565in}{1.108737in}}%
\pgfpathlineto{\pgfqpoint{2.527642in}{1.106923in}}%
\pgfpathlineto{\pgfqpoint{2.516736in}{1.108073in}}%
\pgfpathlineto{\pgfqpoint{2.509239in}{1.087027in}}%
\pgfpathclose%
\pgfusepath{fill}%
\end{pgfscope}%
\begin{pgfscope}%
\pgfpathrectangle{\pgfqpoint{0.100000in}{0.100000in}}{\pgfqpoint{3.007045in}{1.925000in}}%
\pgfusepath{clip}%
\pgfsetbuttcap%
\pgfsetmiterjoin%
\definecolor{currentfill}{rgb}{0.332549,0.708651,0.402768}%
\pgfsetfillcolor{currentfill}%
\pgfsetlinewidth{0.000000pt}%
\definecolor{currentstroke}{rgb}{0.000000,0.000000,0.000000}%
\pgfsetstrokecolor{currentstroke}%
\pgfsetstrokeopacity{0.000000}%
\pgfsetdash{}{0pt}%
\pgfpathmoveto{\pgfqpoint{2.768356in}{1.454967in}}%
\pgfpathlineto{\pgfqpoint{2.739778in}{1.426983in}}%
\pgfpathlineto{\pgfqpoint{2.733867in}{1.427016in}}%
\pgfpathlineto{\pgfqpoint{2.732166in}{1.431244in}}%
\pgfpathlineto{\pgfqpoint{2.727454in}{1.434295in}}%
\pgfpathlineto{\pgfqpoint{2.723883in}{1.436359in}}%
\pgfpathlineto{\pgfqpoint{2.721879in}{1.446280in}}%
\pgfpathlineto{\pgfqpoint{2.711428in}{1.444046in}}%
\pgfpathlineto{\pgfqpoint{2.710631in}{1.447455in}}%
\pgfpathlineto{\pgfqpoint{2.700669in}{1.446056in}}%
\pgfpathlineto{\pgfqpoint{2.697553in}{1.455914in}}%
\pgfpathlineto{\pgfqpoint{2.692078in}{1.475538in}}%
\pgfpathlineto{\pgfqpoint{2.719947in}{1.482951in}}%
\pgfpathlineto{\pgfqpoint{2.722365in}{1.491084in}}%
\pgfpathlineto{\pgfqpoint{2.726801in}{1.495572in}}%
\pgfpathlineto{\pgfqpoint{2.737175in}{1.492249in}}%
\pgfpathlineto{\pgfqpoint{2.737431in}{1.497269in}}%
\pgfpathlineto{\pgfqpoint{2.743391in}{1.496243in}}%
\pgfpathlineto{\pgfqpoint{2.749311in}{1.495307in}}%
\pgfpathlineto{\pgfqpoint{2.749558in}{1.490026in}}%
\pgfpathlineto{\pgfqpoint{2.753121in}{1.482442in}}%
\pgfpathlineto{\pgfqpoint{2.750879in}{1.474563in}}%
\pgfpathlineto{\pgfqpoint{2.756645in}{1.469554in}}%
\pgfpathlineto{\pgfqpoint{2.766813in}{1.466787in}}%
\pgfpathlineto{\pgfqpoint{2.763804in}{1.456101in}}%
\pgfpathlineto{\pgfqpoint{2.768356in}{1.454967in}}%
\pgfpathclose%
\pgfusepath{fill}%
\end{pgfscope}%
\begin{pgfscope}%
\pgfpathrectangle{\pgfqpoint{0.100000in}{0.100000in}}{\pgfqpoint{3.007045in}{1.925000in}}%
\pgfusepath{clip}%
\pgfsetbuttcap%
\pgfsetmiterjoin%
\definecolor{currentfill}{rgb}{0.420392,0.751711,0.445829}%
\pgfsetfillcolor{currentfill}%
\pgfsetlinewidth{0.000000pt}%
\definecolor{currentstroke}{rgb}{0.000000,0.000000,0.000000}%
\pgfsetstrokecolor{currentstroke}%
\pgfsetstrokeopacity{0.000000}%
\pgfsetdash{}{0pt}%
\pgfpathmoveto{\pgfqpoint{1.614445in}{1.686458in}}%
\pgfpathlineto{\pgfqpoint{1.579862in}{1.687980in}}%
\pgfpathlineto{\pgfqpoint{1.580956in}{1.710973in}}%
\pgfpathlineto{\pgfqpoint{1.580095in}{1.722634in}}%
\pgfpathlineto{\pgfqpoint{1.614607in}{1.721113in}}%
\pgfpathlineto{\pgfqpoint{1.615296in}{1.709477in}}%
\pgfpathlineto{\pgfqpoint{1.614445in}{1.686458in}}%
\pgfpathclose%
\pgfusepath{fill}%
\end{pgfscope}%
\begin{pgfscope}%
\pgfpathrectangle{\pgfqpoint{0.100000in}{0.100000in}}{\pgfqpoint{3.007045in}{1.925000in}}%
\pgfusepath{clip}%
\pgfsetbuttcap%
\pgfsetmiterjoin%
\definecolor{currentfill}{rgb}{0.699746,0.879769,0.674418}%
\pgfsetfillcolor{currentfill}%
\pgfsetlinewidth{0.000000pt}%
\definecolor{currentstroke}{rgb}{0.000000,0.000000,0.000000}%
\pgfsetstrokecolor{currentstroke}%
\pgfsetstrokeopacity{0.000000}%
\pgfsetdash{}{0pt}%
\pgfpathmoveto{\pgfqpoint{1.627761in}{0.861901in}}%
\pgfpathlineto{\pgfqpoint{1.627401in}{0.850427in}}%
\pgfpathlineto{\pgfqpoint{1.604562in}{0.851135in}}%
\pgfpathlineto{\pgfqpoint{1.604926in}{0.862611in}}%
\pgfpathlineto{\pgfqpoint{1.576210in}{0.863673in}}%
\pgfpathlineto{\pgfqpoint{1.576781in}{0.879670in}}%
\pgfpathlineto{\pgfqpoint{1.577863in}{0.909634in}}%
\pgfpathlineto{\pgfqpoint{1.594488in}{0.908994in}}%
\pgfpathlineto{\pgfqpoint{1.594438in}{0.897482in}}%
\pgfpathlineto{\pgfqpoint{1.609800in}{0.896970in}}%
\pgfpathlineto{\pgfqpoint{1.619336in}{0.893761in}}%
\pgfpathlineto{\pgfqpoint{1.628500in}{0.893638in}}%
\pgfpathlineto{\pgfqpoint{1.627761in}{0.861901in}}%
\pgfpathclose%
\pgfusepath{fill}%
\end{pgfscope}%
\begin{pgfscope}%
\pgfpathrectangle{\pgfqpoint{0.100000in}{0.100000in}}{\pgfqpoint{3.007045in}{1.925000in}}%
\pgfusepath{clip}%
\pgfsetbuttcap%
\pgfsetmiterjoin%
\definecolor{currentfill}{rgb}{0.407843,0.745559,0.439677}%
\pgfsetfillcolor{currentfill}%
\pgfsetlinewidth{0.000000pt}%
\definecolor{currentstroke}{rgb}{0.000000,0.000000,0.000000}%
\pgfsetstrokecolor{currentstroke}%
\pgfsetstrokeopacity{0.000000}%
\pgfsetdash{}{0pt}%
\pgfpathmoveto{\pgfqpoint{1.821649in}{1.429987in}}%
\pgfpathlineto{\pgfqpoint{1.822052in}{1.390899in}}%
\pgfpathlineto{\pgfqpoint{1.776480in}{1.390760in}}%
\pgfpathlineto{\pgfqpoint{1.753560in}{1.390798in}}%
\pgfpathlineto{\pgfqpoint{1.753660in}{1.413597in}}%
\pgfpathlineto{\pgfqpoint{1.753741in}{1.429789in}}%
\pgfpathlineto{\pgfqpoint{1.776405in}{1.429765in}}%
\pgfpathlineto{\pgfqpoint{1.776471in}{1.413556in}}%
\pgfpathlineto{\pgfqpoint{1.799116in}{1.413633in}}%
\pgfpathlineto{\pgfqpoint{1.799028in}{1.429855in}}%
\pgfpathlineto{\pgfqpoint{1.821649in}{1.429987in}}%
\pgfpathclose%
\pgfusepath{fill}%
\end{pgfscope}%
\begin{pgfscope}%
\pgfpathrectangle{\pgfqpoint{0.100000in}{0.100000in}}{\pgfqpoint{3.007045in}{1.925000in}}%
\pgfusepath{clip}%
\pgfsetbuttcap%
\pgfsetmiterjoin%
\definecolor{currentfill}{rgb}{0.657670,0.862053,0.633449}%
\pgfsetfillcolor{currentfill}%
\pgfsetlinewidth{0.000000pt}%
\definecolor{currentstroke}{rgb}{0.000000,0.000000,0.000000}%
\pgfsetstrokecolor{currentstroke}%
\pgfsetstrokeopacity{0.000000}%
\pgfsetdash{}{0pt}%
\pgfpathmoveto{\pgfqpoint{2.370833in}{1.107548in}}%
\pgfpathlineto{\pgfqpoint{2.371064in}{1.101491in}}%
\pgfpathlineto{\pgfqpoint{2.367300in}{1.093259in}}%
\pgfpathlineto{\pgfqpoint{2.362396in}{1.092621in}}%
\pgfpathlineto{\pgfqpoint{2.355335in}{1.096669in}}%
\pgfpathlineto{\pgfqpoint{2.344556in}{1.097617in}}%
\pgfpathlineto{\pgfqpoint{2.336897in}{1.108708in}}%
\pgfpathlineto{\pgfqpoint{2.347373in}{1.122044in}}%
\pgfpathlineto{\pgfqpoint{2.353384in}{1.121802in}}%
\pgfpathlineto{\pgfqpoint{2.359204in}{1.116262in}}%
\pgfpathlineto{\pgfqpoint{2.370833in}{1.107548in}}%
\pgfpathclose%
\pgfusepath{fill}%
\end{pgfscope}%
\begin{pgfscope}%
\pgfpathrectangle{\pgfqpoint{0.100000in}{0.100000in}}{\pgfqpoint{3.007045in}{1.925000in}}%
\pgfusepath{clip}%
\pgfsetbuttcap%
\pgfsetmiterjoin%
\definecolor{currentfill}{rgb}{0.643645,0.856148,0.619792}%
\pgfsetfillcolor{currentfill}%
\pgfsetlinewidth{0.000000pt}%
\definecolor{currentstroke}{rgb}{0.000000,0.000000,0.000000}%
\pgfsetstrokecolor{currentstroke}%
\pgfsetstrokeopacity{0.000000}%
\pgfsetdash{}{0pt}%
\pgfpathmoveto{\pgfqpoint{2.172150in}{1.293918in}}%
\pgfpathlineto{\pgfqpoint{2.173650in}{1.276857in}}%
\pgfpathlineto{\pgfqpoint{2.171318in}{1.271716in}}%
\pgfpathlineto{\pgfqpoint{2.165690in}{1.271203in}}%
\pgfpathlineto{\pgfqpoint{2.166328in}{1.264572in}}%
\pgfpathlineto{\pgfqpoint{2.145056in}{1.262796in}}%
\pgfpathlineto{\pgfqpoint{2.141986in}{1.299655in}}%
\pgfpathlineto{\pgfqpoint{2.139901in}{1.326614in}}%
\pgfpathlineto{\pgfqpoint{2.145664in}{1.322773in}}%
\pgfpathlineto{\pgfqpoint{2.155565in}{1.322406in}}%
\pgfpathlineto{\pgfqpoint{2.168897in}{1.329206in}}%
\pgfpathlineto{\pgfqpoint{2.172150in}{1.293918in}}%
\pgfpathclose%
\pgfusepath{fill}%
\end{pgfscope}%
\begin{pgfscope}%
\pgfpathrectangle{\pgfqpoint{0.100000in}{0.100000in}}{\pgfqpoint{3.007045in}{1.925000in}}%
\pgfusepath{clip}%
\pgfsetbuttcap%
\pgfsetmiterjoin%
\definecolor{currentfill}{rgb}{0.601615,0.837093,0.583376}%
\pgfsetfillcolor{currentfill}%
\pgfsetlinewidth{0.000000pt}%
\definecolor{currentstroke}{rgb}{0.000000,0.000000,0.000000}%
\pgfsetstrokecolor{currentstroke}%
\pgfsetstrokeopacity{0.000000}%
\pgfsetdash{}{0pt}%
\pgfpathmoveto{\pgfqpoint{1.553928in}{0.768166in}}%
\pgfpathlineto{\pgfqpoint{1.583365in}{0.766882in}}%
\pgfpathlineto{\pgfqpoint{1.583500in}{0.771553in}}%
\pgfpathlineto{\pgfqpoint{1.607913in}{0.770715in}}%
\pgfpathlineto{\pgfqpoint{1.611193in}{0.768407in}}%
\pgfpathlineto{\pgfqpoint{1.610081in}{0.739788in}}%
\pgfpathlineto{\pgfqpoint{1.582071in}{0.741206in}}%
\pgfpathlineto{\pgfqpoint{1.581896in}{0.737526in}}%
\pgfpathlineto{\pgfqpoint{1.573659in}{0.737980in}}%
\pgfpathlineto{\pgfqpoint{1.523941in}{0.740360in}}%
\pgfpathlineto{\pgfqpoint{1.525314in}{0.769593in}}%
\pgfpathlineto{\pgfqpoint{1.553928in}{0.768166in}}%
\pgfpathclose%
\pgfusepath{fill}%
\end{pgfscope}%
\begin{pgfscope}%
\pgfpathrectangle{\pgfqpoint{0.100000in}{0.100000in}}{\pgfqpoint{3.007045in}{1.925000in}}%
\pgfusepath{clip}%
\pgfsetbuttcap%
\pgfsetmiterjoin%
\definecolor{currentfill}{rgb}{0.573933,0.824175,0.560615}%
\pgfsetfillcolor{currentfill}%
\pgfsetlinewidth{0.000000pt}%
\definecolor{currentstroke}{rgb}{0.000000,0.000000,0.000000}%
\pgfsetstrokecolor{currentstroke}%
\pgfsetstrokeopacity{0.000000}%
\pgfsetdash{}{0pt}%
\pgfpathmoveto{\pgfqpoint{1.939480in}{1.158394in}}%
\pgfpathlineto{\pgfqpoint{1.909303in}{1.157973in}}%
\pgfpathlineto{\pgfqpoint{1.909292in}{1.151308in}}%
\pgfpathlineto{\pgfqpoint{1.903384in}{1.151278in}}%
\pgfpathlineto{\pgfqpoint{1.889132in}{1.155695in}}%
\pgfpathlineto{\pgfqpoint{1.889421in}{1.180509in}}%
\pgfpathlineto{\pgfqpoint{1.880876in}{1.180425in}}%
\pgfpathlineto{\pgfqpoint{1.880879in}{1.202651in}}%
\pgfpathlineto{\pgfqpoint{1.906198in}{1.203318in}}%
\pgfpathlineto{\pgfqpoint{1.906383in}{1.197788in}}%
\pgfpathlineto{\pgfqpoint{1.926181in}{1.198155in}}%
\pgfpathlineto{\pgfqpoint{1.931854in}{1.198270in}}%
\pgfpathlineto{\pgfqpoint{1.932193in}{1.179188in}}%
\pgfpathlineto{\pgfqpoint{1.938875in}{1.179366in}}%
\pgfpathlineto{\pgfqpoint{1.939480in}{1.158394in}}%
\pgfpathclose%
\pgfusepath{fill}%
\end{pgfscope}%
\begin{pgfscope}%
\pgfpathrectangle{\pgfqpoint{0.100000in}{0.100000in}}{\pgfqpoint{3.007045in}{1.925000in}}%
\pgfusepath{clip}%
\pgfsetbuttcap%
\pgfsetmiterjoin%
\definecolor{currentfill}{rgb}{0.657670,0.862053,0.633449}%
\pgfsetfillcolor{currentfill}%
\pgfsetlinewidth{0.000000pt}%
\definecolor{currentstroke}{rgb}{0.000000,0.000000,0.000000}%
\pgfsetstrokecolor{currentstroke}%
\pgfsetstrokeopacity{0.000000}%
\pgfsetdash{}{0pt}%
\pgfpathmoveto{\pgfqpoint{1.719231in}{0.879638in}}%
\pgfpathlineto{\pgfqpoint{1.719456in}{0.888768in}}%
\pgfpathlineto{\pgfqpoint{1.719521in}{0.894505in}}%
\pgfpathlineto{\pgfqpoint{1.725208in}{0.894438in}}%
\pgfpathlineto{\pgfqpoint{1.728070in}{0.900172in}}%
\pgfpathlineto{\pgfqpoint{1.728181in}{0.905962in}}%
\pgfpathlineto{\pgfqpoint{1.734052in}{0.905918in}}%
\pgfpathlineto{\pgfqpoint{1.734135in}{0.917404in}}%
\pgfpathlineto{\pgfqpoint{1.731290in}{0.917416in}}%
\pgfpathlineto{\pgfqpoint{1.731364in}{0.926038in}}%
\pgfpathlineto{\pgfqpoint{1.737553in}{0.926013in}}%
\pgfpathlineto{\pgfqpoint{1.740818in}{0.919900in}}%
\pgfpathlineto{\pgfqpoint{1.746777in}{0.922489in}}%
\pgfpathlineto{\pgfqpoint{1.758114in}{0.923040in}}%
\pgfpathlineto{\pgfqpoint{1.757448in}{0.932993in}}%
\pgfpathlineto{\pgfqpoint{1.761242in}{0.940278in}}%
\pgfpathlineto{\pgfqpoint{1.766027in}{0.940252in}}%
\pgfpathlineto{\pgfqpoint{1.766020in}{0.946003in}}%
\pgfpathlineto{\pgfqpoint{1.771725in}{0.946010in}}%
\pgfpathlineto{\pgfqpoint{1.795605in}{0.946063in}}%
\pgfpathlineto{\pgfqpoint{1.799463in}{0.919527in}}%
\pgfpathlineto{\pgfqpoint{1.800606in}{0.911576in}}%
\pgfpathlineto{\pgfqpoint{1.765492in}{0.911520in}}%
\pgfpathlineto{\pgfqpoint{1.765215in}{0.904084in}}%
\pgfpathlineto{\pgfqpoint{1.768002in}{0.897857in}}%
\pgfpathlineto{\pgfqpoint{1.763675in}{0.893735in}}%
\pgfpathlineto{\pgfqpoint{1.759980in}{0.886994in}}%
\pgfpathlineto{\pgfqpoint{1.751555in}{0.890916in}}%
\pgfpathlineto{\pgfqpoint{1.741580in}{0.885648in}}%
\pgfpathlineto{\pgfqpoint{1.732284in}{0.882217in}}%
\pgfpathlineto{\pgfqpoint{1.725397in}{0.881888in}}%
\pgfpathlineto{\pgfqpoint{1.719231in}{0.879638in}}%
\pgfpathclose%
\pgfusepath{fill}%
\end{pgfscope}%
\begin{pgfscope}%
\pgfpathrectangle{\pgfqpoint{0.100000in}{0.100000in}}{\pgfqpoint{3.007045in}{1.925000in}}%
\pgfusepath{clip}%
\pgfsetbuttcap%
\pgfsetmiterjoin%
\definecolor{currentfill}{rgb}{0.474279,0.777670,0.478677}%
\pgfsetfillcolor{currentfill}%
\pgfsetlinewidth{0.000000pt}%
\definecolor{currentstroke}{rgb}{0.000000,0.000000,0.000000}%
\pgfsetstrokecolor{currentstroke}%
\pgfsetstrokeopacity{0.000000}%
\pgfsetdash{}{0pt}%
\pgfpathmoveto{\pgfqpoint{1.756771in}{1.198935in}}%
\pgfpathlineto{\pgfqpoint{1.751069in}{1.202842in}}%
\pgfpathlineto{\pgfqpoint{1.751959in}{1.207454in}}%
\pgfpathlineto{\pgfqpoint{1.748211in}{1.211203in}}%
\pgfpathlineto{\pgfqpoint{1.747673in}{1.215377in}}%
\pgfpathlineto{\pgfqpoint{1.741287in}{1.219682in}}%
\pgfpathlineto{\pgfqpoint{1.734017in}{1.237625in}}%
\pgfpathlineto{\pgfqpoint{1.790743in}{1.236578in}}%
\pgfpathlineto{\pgfqpoint{1.792476in}{1.224430in}}%
\pgfpathlineto{\pgfqpoint{1.792177in}{1.207263in}}%
\pgfpathlineto{\pgfqpoint{1.770162in}{1.207494in}}%
\pgfpathlineto{\pgfqpoint{1.772661in}{1.193312in}}%
\pgfpathlineto{\pgfqpoint{1.768029in}{1.189813in}}%
\pgfpathlineto{\pgfqpoint{1.762123in}{1.192600in}}%
\pgfpathlineto{\pgfqpoint{1.756771in}{1.198935in}}%
\pgfpathclose%
\pgfusepath{fill}%
\end{pgfscope}%
\begin{pgfscope}%
\pgfpathrectangle{\pgfqpoint{0.100000in}{0.100000in}}{\pgfqpoint{3.007045in}{1.925000in}}%
\pgfusepath{clip}%
\pgfsetbuttcap%
\pgfsetmiterjoin%
\definecolor{currentfill}{rgb}{0.401569,0.742484,0.436601}%
\pgfsetfillcolor{currentfill}%
\pgfsetlinewidth{0.000000pt}%
\definecolor{currentstroke}{rgb}{0.000000,0.000000,0.000000}%
\pgfsetstrokecolor{currentstroke}%
\pgfsetstrokeopacity{0.000000}%
\pgfsetdash{}{0pt}%
\pgfpathmoveto{\pgfqpoint{2.285385in}{0.933858in}}%
\pgfpathlineto{\pgfqpoint{2.278483in}{0.930848in}}%
\pgfpathlineto{\pgfqpoint{2.264733in}{0.928140in}}%
\pgfpathlineto{\pgfqpoint{2.261477in}{0.928516in}}%
\pgfpathlineto{\pgfqpoint{2.257324in}{0.931388in}}%
\pgfpathlineto{\pgfqpoint{2.254813in}{0.936868in}}%
\pgfpathlineto{\pgfqpoint{2.254335in}{0.941757in}}%
\pgfpathlineto{\pgfqpoint{2.257521in}{0.951222in}}%
\pgfpathlineto{\pgfqpoint{2.251263in}{0.956468in}}%
\pgfpathlineto{\pgfqpoint{2.248417in}{0.962697in}}%
\pgfpathlineto{\pgfqpoint{2.253187in}{0.966596in}}%
\pgfpathlineto{\pgfqpoint{2.258371in}{0.965969in}}%
\pgfpathlineto{\pgfqpoint{2.261147in}{0.968764in}}%
\pgfpathlineto{\pgfqpoint{2.263169in}{0.965439in}}%
\pgfpathlineto{\pgfqpoint{2.267876in}{0.965337in}}%
\pgfpathlineto{\pgfqpoint{2.269975in}{0.962101in}}%
\pgfpathlineto{\pgfqpoint{2.277359in}{0.967152in}}%
\pgfpathlineto{\pgfqpoint{2.285629in}{0.965628in}}%
\pgfpathlineto{\pgfqpoint{2.291000in}{0.961775in}}%
\pgfpathlineto{\pgfqpoint{2.293939in}{0.955321in}}%
\pgfpathlineto{\pgfqpoint{2.292762in}{0.944860in}}%
\pgfpathlineto{\pgfqpoint{2.285385in}{0.933858in}}%
\pgfpathclose%
\pgfusepath{fill}%
\end{pgfscope}%
\begin{pgfscope}%
\pgfpathrectangle{\pgfqpoint{0.100000in}{0.100000in}}{\pgfqpoint{3.007045in}{1.925000in}}%
\pgfusepath{clip}%
\pgfsetbuttcap%
\pgfsetmiterjoin%
\definecolor{currentfill}{rgb}{0.507497,0.793172,0.505990}%
\pgfsetfillcolor{currentfill}%
\pgfsetlinewidth{0.000000pt}%
\definecolor{currentstroke}{rgb}{0.000000,0.000000,0.000000}%
\pgfsetstrokecolor{currentstroke}%
\pgfsetstrokeopacity{0.000000}%
\pgfsetdash{}{0pt}%
\pgfpathmoveto{\pgfqpoint{1.193280in}{1.837329in}}%
\pgfpathlineto{\pgfqpoint{1.244920in}{1.830121in}}%
\pgfpathlineto{\pgfqpoint{1.285700in}{1.824835in}}%
\pgfpathlineto{\pgfqpoint{1.282530in}{1.813419in}}%
\pgfpathlineto{\pgfqpoint{1.288237in}{1.812721in}}%
\pgfpathlineto{\pgfqpoint{1.286131in}{1.795421in}}%
\pgfpathlineto{\pgfqpoint{1.295653in}{1.794279in}}%
\pgfpathlineto{\pgfqpoint{1.292930in}{1.771456in}}%
\pgfpathlineto{\pgfqpoint{1.291202in}{1.771653in}}%
\pgfpathlineto{\pgfqpoint{1.289529in}{1.757826in}}%
\pgfpathlineto{\pgfqpoint{1.272584in}{1.760976in}}%
\pgfpathlineto{\pgfqpoint{1.265605in}{1.764425in}}%
\pgfpathlineto{\pgfqpoint{1.263629in}{1.757459in}}%
\pgfpathlineto{\pgfqpoint{1.260272in}{1.755961in}}%
\pgfpathlineto{\pgfqpoint{1.255859in}{1.748834in}}%
\pgfpathlineto{\pgfqpoint{1.246943in}{1.744290in}}%
\pgfpathlineto{\pgfqpoint{1.241819in}{1.745330in}}%
\pgfpathlineto{\pgfqpoint{1.240243in}{1.741597in}}%
\pgfpathlineto{\pgfqpoint{1.224257in}{1.742711in}}%
\pgfpathlineto{\pgfqpoint{1.217097in}{1.745725in}}%
\pgfpathlineto{\pgfqpoint{1.214896in}{1.741346in}}%
\pgfpathlineto{\pgfqpoint{1.206785in}{1.744306in}}%
\pgfpathlineto{\pgfqpoint{1.201819in}{1.739713in}}%
\pgfpathlineto{\pgfqpoint{1.195585in}{1.736909in}}%
\pgfpathlineto{\pgfqpoint{1.192875in}{1.732905in}}%
\pgfpathlineto{\pgfqpoint{1.193102in}{1.740744in}}%
\pgfpathlineto{\pgfqpoint{1.187367in}{1.743646in}}%
\pgfpathlineto{\pgfqpoint{1.174368in}{1.744062in}}%
\pgfpathlineto{\pgfqpoint{1.170837in}{1.747018in}}%
\pgfpathlineto{\pgfqpoint{1.164883in}{1.747708in}}%
\pgfpathlineto{\pgfqpoint{1.156612in}{1.751283in}}%
\pgfpathlineto{\pgfqpoint{1.152002in}{1.758091in}}%
\pgfpathlineto{\pgfqpoint{1.154039in}{1.770483in}}%
\pgfpathlineto{\pgfqpoint{1.165616in}{1.768587in}}%
\pgfpathlineto{\pgfqpoint{1.167096in}{1.773103in}}%
\pgfpathlineto{\pgfqpoint{1.174638in}{1.770993in}}%
\pgfpathlineto{\pgfqpoint{1.177030in}{1.787023in}}%
\pgfpathlineto{\pgfqpoint{1.180237in}{1.801764in}}%
\pgfpathlineto{\pgfqpoint{1.182880in}{1.799848in}}%
\pgfpathlineto{\pgfqpoint{1.185494in}{1.809331in}}%
\pgfpathlineto{\pgfqpoint{1.187177in}{1.820830in}}%
\pgfpathlineto{\pgfqpoint{1.189784in}{1.820439in}}%
\pgfpathlineto{\pgfqpoint{1.193280in}{1.837329in}}%
\pgfpathclose%
\pgfusepath{fill}%
\end{pgfscope}%
\begin{pgfscope}%
\pgfpathrectangle{\pgfqpoint{0.100000in}{0.100000in}}{\pgfqpoint{3.007045in}{1.925000in}}%
\pgfusepath{clip}%
\pgfsetbuttcap%
\pgfsetmiterjoin%
\definecolor{currentfill}{rgb}{0.182468,0.593326,0.306759}%
\pgfsetfillcolor{currentfill}%
\pgfsetlinewidth{0.000000pt}%
\definecolor{currentstroke}{rgb}{0.000000,0.000000,0.000000}%
\pgfsetstrokecolor{currentstroke}%
\pgfsetstrokeopacity{0.000000}%
\pgfsetdash{}{0pt}%
\pgfpathmoveto{\pgfqpoint{2.642481in}{0.444772in}}%
\pgfpathlineto{\pgfqpoint{2.615609in}{0.440573in}}%
\pgfpathlineto{\pgfqpoint{2.614673in}{0.446167in}}%
\pgfpathlineto{\pgfqpoint{2.609094in}{0.445229in}}%
\pgfpathlineto{\pgfqpoint{2.602369in}{0.492044in}}%
\pgfpathlineto{\pgfqpoint{2.599254in}{0.497364in}}%
\pgfpathlineto{\pgfqpoint{2.599396in}{0.502567in}}%
\pgfpathlineto{\pgfqpoint{2.592396in}{0.508481in}}%
\pgfpathlineto{\pgfqpoint{2.591696in}{0.520372in}}%
\pgfpathlineto{\pgfqpoint{2.605371in}{0.522544in}}%
\pgfpathlineto{\pgfqpoint{2.616569in}{0.510498in}}%
\pgfpathlineto{\pgfqpoint{2.620884in}{0.502421in}}%
\pgfpathlineto{\pgfqpoint{2.617307in}{0.495527in}}%
\pgfpathlineto{\pgfqpoint{2.618107in}{0.489509in}}%
\pgfpathlineto{\pgfqpoint{2.622871in}{0.477008in}}%
\pgfpathlineto{\pgfqpoint{2.636790in}{0.456290in}}%
\pgfpathlineto{\pgfqpoint{2.642481in}{0.444772in}}%
\pgfpathclose%
\pgfusepath{fill}%
\end{pgfscope}%
\begin{pgfscope}%
\pgfpathrectangle{\pgfqpoint{0.100000in}{0.100000in}}{\pgfqpoint{3.007045in}{1.925000in}}%
\pgfusepath{clip}%
\pgfsetbuttcap%
\pgfsetmiterjoin%
\definecolor{currentfill}{rgb}{0.657670,0.862053,0.633449}%
\pgfsetfillcolor{currentfill}%
\pgfsetlinewidth{0.000000pt}%
\definecolor{currentstroke}{rgb}{0.000000,0.000000,0.000000}%
\pgfsetstrokecolor{currentstroke}%
\pgfsetstrokeopacity{0.000000}%
\pgfsetdash{}{0pt}%
\pgfpathmoveto{\pgfqpoint{1.865214in}{0.855241in}}%
\pgfpathlineto{\pgfqpoint{1.859472in}{0.853168in}}%
\pgfpathlineto{\pgfqpoint{1.842321in}{0.853020in}}%
\pgfpathlineto{\pgfqpoint{1.842357in}{0.871145in}}%
\pgfpathlineto{\pgfqpoint{1.822375in}{0.871354in}}%
\pgfpathlineto{\pgfqpoint{1.818712in}{0.873240in}}%
\pgfpathlineto{\pgfqpoint{1.818790in}{0.879030in}}%
\pgfpathlineto{\pgfqpoint{1.821794in}{0.883783in}}%
\pgfpathlineto{\pgfqpoint{1.822064in}{0.899000in}}%
\pgfpathlineto{\pgfqpoint{1.824480in}{0.903095in}}%
\pgfpathlineto{\pgfqpoint{1.821906in}{0.907584in}}%
\pgfpathlineto{\pgfqpoint{1.821996in}{0.913307in}}%
\pgfpathlineto{\pgfqpoint{1.826781in}{0.913265in}}%
\pgfpathlineto{\pgfqpoint{1.830635in}{0.919840in}}%
\pgfpathlineto{\pgfqpoint{1.851596in}{0.920385in}}%
\pgfpathlineto{\pgfqpoint{1.851607in}{0.918465in}}%
\pgfpathlineto{\pgfqpoint{1.889632in}{0.918440in}}%
\pgfpathlineto{\pgfqpoint{1.889712in}{0.906997in}}%
\pgfpathlineto{\pgfqpoint{1.887849in}{0.901272in}}%
\pgfpathlineto{\pgfqpoint{1.888066in}{0.882118in}}%
\pgfpathlineto{\pgfqpoint{1.884196in}{0.881637in}}%
\pgfpathlineto{\pgfqpoint{1.871864in}{0.868767in}}%
\pgfpathlineto{\pgfqpoint{1.864117in}{0.861985in}}%
\pgfpathlineto{\pgfqpoint{1.865214in}{0.855241in}}%
\pgfpathclose%
\pgfusepath{fill}%
\end{pgfscope}%
\begin{pgfscope}%
\pgfpathrectangle{\pgfqpoint{0.100000in}{0.100000in}}{\pgfqpoint{3.007045in}{1.925000in}}%
\pgfusepath{clip}%
\pgfsetbuttcap%
\pgfsetmiterjoin%
\definecolor{currentfill}{rgb}{0.704421,0.881738,0.678970}%
\pgfsetfillcolor{currentfill}%
\pgfsetlinewidth{0.000000pt}%
\definecolor{currentstroke}{rgb}{0.000000,0.000000,0.000000}%
\pgfsetstrokecolor{currentstroke}%
\pgfsetstrokeopacity{0.000000}%
\pgfsetdash{}{0pt}%
\pgfpathmoveto{\pgfqpoint{0.912796in}{0.778126in}}%
\pgfpathlineto{\pgfqpoint{0.851184in}{0.788621in}}%
\pgfpathlineto{\pgfqpoint{0.795773in}{0.799049in}}%
\pgfpathlineto{\pgfqpoint{0.754327in}{0.807329in}}%
\pgfpathlineto{\pgfqpoint{0.765905in}{0.863949in}}%
\pgfpathlineto{\pgfqpoint{0.774157in}{0.904247in}}%
\pgfpathlineto{\pgfqpoint{0.805861in}{0.897916in}}%
\pgfpathlineto{\pgfqpoint{0.829566in}{0.885369in}}%
\pgfpathlineto{\pgfqpoint{0.837638in}{0.894986in}}%
\pgfpathlineto{\pgfqpoint{0.860696in}{0.887387in}}%
\pgfpathlineto{\pgfqpoint{0.873121in}{0.885088in}}%
\pgfpathlineto{\pgfqpoint{0.875838in}{0.894759in}}%
\pgfpathlineto{\pgfqpoint{0.863233in}{0.897098in}}%
\pgfpathlineto{\pgfqpoint{0.868902in}{0.911853in}}%
\pgfpathlineto{\pgfqpoint{0.876178in}{0.916207in}}%
\pgfpathlineto{\pgfqpoint{0.883290in}{0.911640in}}%
\pgfpathlineto{\pgfqpoint{0.890889in}{0.911616in}}%
\pgfpathlineto{\pgfqpoint{0.896277in}{0.906751in}}%
\pgfpathlineto{\pgfqpoint{0.903247in}{0.903940in}}%
\pgfpathlineto{\pgfqpoint{0.911088in}{0.896114in}}%
\pgfpathlineto{\pgfqpoint{0.916273in}{0.895048in}}%
\pgfpathlineto{\pgfqpoint{0.913272in}{0.877904in}}%
\pgfpathlineto{\pgfqpoint{0.953797in}{0.870982in}}%
\pgfpathlineto{\pgfqpoint{0.949199in}{0.843379in}}%
\pgfpathlineto{\pgfqpoint{0.947992in}{0.836130in}}%
\pgfpathlineto{\pgfqpoint{0.938796in}{0.837657in}}%
\pgfpathlineto{\pgfqpoint{0.930194in}{0.836921in}}%
\pgfpathlineto{\pgfqpoint{0.923055in}{0.834676in}}%
\pgfpathlineto{\pgfqpoint{0.920548in}{0.822552in}}%
\pgfpathlineto{\pgfqpoint{0.912796in}{0.778126in}}%
\pgfpathclose%
\pgfusepath{fill}%
\end{pgfscope}%
\begin{pgfscope}%
\pgfpathrectangle{\pgfqpoint{0.100000in}{0.100000in}}{\pgfqpoint{3.007045in}{1.925000in}}%
\pgfusepath{clip}%
\pgfsetbuttcap%
\pgfsetmiterjoin%
\definecolor{currentfill}{rgb}{0.513033,0.795755,0.510542}%
\pgfsetfillcolor{currentfill}%
\pgfsetlinewidth{0.000000pt}%
\definecolor{currentstroke}{rgb}{0.000000,0.000000,0.000000}%
\pgfsetstrokecolor{currentstroke}%
\pgfsetstrokeopacity{0.000000}%
\pgfsetdash{}{0pt}%
\pgfpathmoveto{\pgfqpoint{1.710630in}{1.795591in}}%
\pgfpathlineto{\pgfqpoint{1.765577in}{1.795047in}}%
\pgfpathlineto{\pgfqpoint{1.765638in}{1.820857in}}%
\pgfpathlineto{\pgfqpoint{1.769776in}{1.818763in}}%
\pgfpathlineto{\pgfqpoint{1.774181in}{1.819903in}}%
\pgfpathlineto{\pgfqpoint{1.779963in}{1.814824in}}%
\pgfpathlineto{\pgfqpoint{1.786263in}{1.787384in}}%
\pgfpathlineto{\pgfqpoint{1.785945in}{1.780293in}}%
\pgfpathlineto{\pgfqpoint{1.790535in}{1.776268in}}%
\pgfpathlineto{\pgfqpoint{1.797461in}{1.775190in}}%
\pgfpathlineto{\pgfqpoint{1.797671in}{1.752902in}}%
\pgfpathlineto{\pgfqpoint{1.762966in}{1.752742in}}%
\pgfpathlineto{\pgfqpoint{1.762948in}{1.764365in}}%
\pgfpathlineto{\pgfqpoint{1.745680in}{1.764366in}}%
\pgfpathlineto{\pgfqpoint{1.710937in}{1.765107in}}%
\pgfpathlineto{\pgfqpoint{1.710357in}{1.776679in}}%
\pgfpathlineto{\pgfqpoint{1.710630in}{1.795591in}}%
\pgfpathclose%
\pgfusepath{fill}%
\end{pgfscope}%
\begin{pgfscope}%
\pgfpathrectangle{\pgfqpoint{0.100000in}{0.100000in}}{\pgfqpoint{3.007045in}{1.925000in}}%
\pgfusepath{clip}%
\pgfsetbuttcap%
\pgfsetmiterjoin%
\definecolor{currentfill}{rgb}{0.490888,0.785421,0.492334}%
\pgfsetfillcolor{currentfill}%
\pgfsetlinewidth{0.000000pt}%
\definecolor{currentstroke}{rgb}{0.000000,0.000000,0.000000}%
\pgfsetstrokecolor{currentstroke}%
\pgfsetstrokeopacity{0.000000}%
\pgfsetdash{}{0pt}%
\pgfpathmoveto{\pgfqpoint{0.851627in}{1.798562in}}%
\pgfpathlineto{\pgfqpoint{0.847085in}{1.804318in}}%
\pgfpathlineto{\pgfqpoint{0.850775in}{1.806624in}}%
\pgfpathlineto{\pgfqpoint{0.850754in}{1.813504in}}%
\pgfpathlineto{\pgfqpoint{0.847161in}{1.818241in}}%
\pgfpathlineto{\pgfqpoint{0.840821in}{1.835141in}}%
\pgfpathlineto{\pgfqpoint{0.844425in}{1.850601in}}%
\pgfpathlineto{\pgfqpoint{0.845385in}{1.846522in}}%
\pgfpathlineto{\pgfqpoint{0.859696in}{1.849285in}}%
\pgfpathlineto{\pgfqpoint{0.860078in}{1.837613in}}%
\pgfpathlineto{\pgfqpoint{0.863113in}{1.830232in}}%
\pgfpathlineto{\pgfqpoint{0.862566in}{1.825406in}}%
\pgfpathlineto{\pgfqpoint{0.873246in}{1.821840in}}%
\pgfpathlineto{\pgfqpoint{0.877958in}{1.823506in}}%
\pgfpathlineto{\pgfqpoint{0.880330in}{1.828801in}}%
\pgfpathlineto{\pgfqpoint{0.886583in}{1.827387in}}%
\pgfpathlineto{\pgfqpoint{0.889217in}{1.841046in}}%
\pgfpathlineto{\pgfqpoint{0.896692in}{1.839394in}}%
\pgfpathlineto{\pgfqpoint{0.901644in}{1.862062in}}%
\pgfpathlineto{\pgfqpoint{0.901194in}{1.867977in}}%
\pgfpathlineto{\pgfqpoint{0.910046in}{1.867809in}}%
\pgfpathlineto{\pgfqpoint{0.913309in}{1.873552in}}%
\pgfpathlineto{\pgfqpoint{0.912328in}{1.878097in}}%
\pgfpathlineto{\pgfqpoint{0.913025in}{1.888772in}}%
\pgfpathlineto{\pgfqpoint{0.941334in}{1.882540in}}%
\pgfpathlineto{\pgfqpoint{0.941225in}{1.874500in}}%
\pgfpathlineto{\pgfqpoint{0.943633in}{1.869658in}}%
\pgfpathlineto{\pgfqpoint{0.949121in}{1.871112in}}%
\pgfpathlineto{\pgfqpoint{0.953218in}{1.860951in}}%
\pgfpathlineto{\pgfqpoint{0.949618in}{1.854808in}}%
\pgfpathlineto{\pgfqpoint{0.961059in}{1.847574in}}%
\pgfpathlineto{\pgfqpoint{0.959285in}{1.841945in}}%
\pgfpathlineto{\pgfqpoint{0.964152in}{1.834807in}}%
\pgfpathlineto{\pgfqpoint{0.961768in}{1.833418in}}%
\pgfpathlineto{\pgfqpoint{0.967054in}{1.825877in}}%
\pgfpathlineto{\pgfqpoint{0.966264in}{1.820814in}}%
\pgfpathlineto{\pgfqpoint{0.975124in}{1.816423in}}%
\pgfpathlineto{\pgfqpoint{0.978555in}{1.808720in}}%
\pgfpathlineto{\pgfqpoint{0.974188in}{1.804466in}}%
\pgfpathlineto{\pgfqpoint{0.969863in}{1.800673in}}%
\pgfpathlineto{\pgfqpoint{0.967284in}{1.791973in}}%
\pgfpathlineto{\pgfqpoint{0.963166in}{1.790823in}}%
\pgfpathlineto{\pgfqpoint{0.962444in}{1.782513in}}%
\pgfpathlineto{\pgfqpoint{0.948336in}{1.785640in}}%
\pgfpathlineto{\pgfqpoint{0.927953in}{1.789677in}}%
\pgfpathlineto{\pgfqpoint{0.925310in}{1.782011in}}%
\pgfpathlineto{\pgfqpoint{0.927695in}{1.775797in}}%
\pgfpathlineto{\pgfqpoint{0.926220in}{1.769071in}}%
\pgfpathlineto{\pgfqpoint{0.926970in}{1.761628in}}%
\pgfpathlineto{\pgfqpoint{0.921428in}{1.759848in}}%
\pgfpathlineto{\pgfqpoint{0.910196in}{1.762234in}}%
\pgfpathlineto{\pgfqpoint{0.906982in}{1.761729in}}%
\pgfpathlineto{\pgfqpoint{0.900754in}{1.769170in}}%
\pgfpathlineto{\pgfqpoint{0.890065in}{1.775890in}}%
\pgfpathlineto{\pgfqpoint{0.883118in}{1.777032in}}%
\pgfpathlineto{\pgfqpoint{0.878165in}{1.781128in}}%
\pgfpathlineto{\pgfqpoint{0.878913in}{1.786617in}}%
\pgfpathlineto{\pgfqpoint{0.866931in}{1.795554in}}%
\pgfpathlineto{\pgfqpoint{0.859543in}{1.795823in}}%
\pgfpathlineto{\pgfqpoint{0.851627in}{1.798562in}}%
\pgfpathclose%
\pgfusepath{fill}%
\end{pgfscope}%
\begin{pgfscope}%
\pgfpathrectangle{\pgfqpoint{0.100000in}{0.100000in}}{\pgfqpoint{3.007045in}{1.925000in}}%
\pgfusepath{clip}%
\pgfsetbuttcap%
\pgfsetmiterjoin%
\definecolor{currentfill}{rgb}{0.816378,0.928120,0.791326}%
\pgfsetfillcolor{currentfill}%
\pgfsetlinewidth{0.000000pt}%
\definecolor{currentstroke}{rgb}{0.000000,0.000000,0.000000}%
\pgfsetstrokecolor{currentstroke}%
\pgfsetstrokeopacity{0.000000}%
\pgfsetdash{}{0pt}%
\pgfpathmoveto{\pgfqpoint{1.336243in}{0.559459in}}%
\pgfpathlineto{\pgfqpoint{1.349461in}{0.546748in}}%
\pgfpathlineto{\pgfqpoint{1.344620in}{0.539783in}}%
\pgfpathlineto{\pgfqpoint{1.328189in}{0.539358in}}%
\pgfpathlineto{\pgfqpoint{1.323294in}{0.529791in}}%
\pgfpathlineto{\pgfqpoint{1.319302in}{0.525390in}}%
\pgfpathlineto{\pgfqpoint{1.316511in}{0.514705in}}%
\pgfpathlineto{\pgfqpoint{1.305003in}{0.500040in}}%
\pgfpathlineto{\pgfqpoint{1.299683in}{0.496151in}}%
\pgfpathlineto{\pgfqpoint{1.298360in}{0.491273in}}%
\pgfpathlineto{\pgfqpoint{1.291128in}{0.491622in}}%
\pgfpathlineto{\pgfqpoint{1.278806in}{0.498569in}}%
\pgfpathlineto{\pgfqpoint{1.273825in}{0.505053in}}%
\pgfpathlineto{\pgfqpoint{1.264636in}{0.507489in}}%
\pgfpathlineto{\pgfqpoint{1.261257in}{0.513470in}}%
\pgfpathlineto{\pgfqpoint{1.250316in}{0.516616in}}%
\pgfpathlineto{\pgfqpoint{1.241232in}{0.523169in}}%
\pgfpathlineto{\pgfqpoint{1.236587in}{0.531465in}}%
\pgfpathlineto{\pgfqpoint{1.231030in}{0.533550in}}%
\pgfpathlineto{\pgfqpoint{1.221419in}{0.542714in}}%
\pgfpathlineto{\pgfqpoint{1.219216in}{0.552154in}}%
\pgfpathlineto{\pgfqpoint{1.214219in}{0.561569in}}%
\pgfpathlineto{\pgfqpoint{1.213608in}{0.571913in}}%
\pgfpathlineto{\pgfqpoint{1.214655in}{0.583241in}}%
\pgfpathlineto{\pgfqpoint{1.207041in}{0.594453in}}%
\pgfpathlineto{\pgfqpoint{1.207348in}{0.601413in}}%
\pgfpathlineto{\pgfqpoint{1.204717in}{0.609116in}}%
\pgfpathlineto{\pgfqpoint{1.201869in}{0.611048in}}%
\pgfpathlineto{\pgfqpoint{1.205698in}{0.612940in}}%
\pgfpathlineto{\pgfqpoint{1.255027in}{0.637292in}}%
\pgfpathlineto{\pgfqpoint{1.282209in}{0.612003in}}%
\pgfpathlineto{\pgfqpoint{1.336243in}{0.559459in}}%
\pgfpathclose%
\pgfusepath{fill}%
\end{pgfscope}%
\begin{pgfscope}%
\pgfpathrectangle{\pgfqpoint{0.100000in}{0.100000in}}{\pgfqpoint{3.007045in}{1.925000in}}%
\pgfusepath{clip}%
\pgfsetbuttcap%
\pgfsetmiterjoin%
\definecolor{currentfill}{rgb}{0.590542,0.831926,0.574271}%
\pgfsetfillcolor{currentfill}%
\pgfsetlinewidth{0.000000pt}%
\definecolor{currentstroke}{rgb}{0.000000,0.000000,0.000000}%
\pgfsetstrokecolor{currentstroke}%
\pgfsetstrokeopacity{0.000000}%
\pgfsetdash{}{0pt}%
\pgfpathmoveto{\pgfqpoint{2.279304in}{0.827039in}}%
\pgfpathlineto{\pgfqpoint{2.271634in}{0.825054in}}%
\pgfpathlineto{\pgfqpoint{2.262319in}{0.816941in}}%
\pgfpathlineto{\pgfqpoint{2.257720in}{0.820289in}}%
\pgfpathlineto{\pgfqpoint{2.254073in}{0.825883in}}%
\pgfpathlineto{\pgfqpoint{2.249130in}{0.825060in}}%
\pgfpathlineto{\pgfqpoint{2.246429in}{0.828970in}}%
\pgfpathlineto{\pgfqpoint{2.247672in}{0.832584in}}%
\pgfpathlineto{\pgfqpoint{2.236928in}{0.845086in}}%
\pgfpathlineto{\pgfqpoint{2.231240in}{0.844683in}}%
\pgfpathlineto{\pgfqpoint{2.232658in}{0.857963in}}%
\pgfpathlineto{\pgfqpoint{2.238722in}{0.856984in}}%
\pgfpathlineto{\pgfqpoint{2.244466in}{0.861498in}}%
\pgfpathlineto{\pgfqpoint{2.241299in}{0.867393in}}%
\pgfpathlineto{\pgfqpoint{2.241813in}{0.880528in}}%
\pgfpathlineto{\pgfqpoint{2.240204in}{0.883283in}}%
\pgfpathlineto{\pgfqpoint{2.241623in}{0.891209in}}%
\pgfpathlineto{\pgfqpoint{2.265651in}{0.893385in}}%
\pgfpathlineto{\pgfqpoint{2.279674in}{0.894582in}}%
\pgfpathlineto{\pgfqpoint{2.286587in}{0.869005in}}%
\pgfpathlineto{\pgfqpoint{2.287764in}{0.864843in}}%
\pgfpathlineto{\pgfqpoint{2.284621in}{0.861875in}}%
\pgfpathlineto{\pgfqpoint{2.282179in}{0.853795in}}%
\pgfpathlineto{\pgfqpoint{2.275437in}{0.845959in}}%
\pgfpathlineto{\pgfqpoint{2.271751in}{0.844592in}}%
\pgfpathlineto{\pgfqpoint{2.279304in}{0.827039in}}%
\pgfpathclose%
\pgfusepath{fill}%
\end{pgfscope}%
\begin{pgfscope}%
\pgfpathrectangle{\pgfqpoint{0.100000in}{0.100000in}}{\pgfqpoint{3.007045in}{1.925000in}}%
\pgfusepath{clip}%
\pgfsetbuttcap%
\pgfsetmiterjoin%
\definecolor{currentfill}{rgb}{0.338824,0.711726,0.405844}%
\pgfsetfillcolor{currentfill}%
\pgfsetlinewidth{0.000000pt}%
\definecolor{currentstroke}{rgb}{0.000000,0.000000,0.000000}%
\pgfsetstrokecolor{currentstroke}%
\pgfsetstrokeopacity{0.000000}%
\pgfsetdash{}{0pt}%
\pgfpathmoveto{\pgfqpoint{0.521615in}{1.353758in}}%
\pgfpathlineto{\pgfqpoint{0.517066in}{1.358013in}}%
\pgfpathlineto{\pgfqpoint{0.520278in}{1.368998in}}%
\pgfpathlineto{\pgfqpoint{0.517118in}{1.374961in}}%
\pgfpathlineto{\pgfqpoint{0.518020in}{1.379672in}}%
\pgfpathlineto{\pgfqpoint{0.512166in}{1.383346in}}%
\pgfpathlineto{\pgfqpoint{0.503882in}{1.397397in}}%
\pgfpathlineto{\pgfqpoint{0.495392in}{1.401958in}}%
\pgfpathlineto{\pgfqpoint{0.488838in}{1.398768in}}%
\pgfpathlineto{\pgfqpoint{0.482087in}{1.399330in}}%
\pgfpathlineto{\pgfqpoint{0.479961in}{1.402391in}}%
\pgfpathlineto{\pgfqpoint{0.483507in}{1.414404in}}%
\pgfpathlineto{\pgfqpoint{0.470680in}{1.418127in}}%
\pgfpathlineto{\pgfqpoint{0.479723in}{1.447096in}}%
\pgfpathlineto{\pgfqpoint{0.484393in}{1.464821in}}%
\pgfpathlineto{\pgfqpoint{0.429023in}{1.481634in}}%
\pgfpathlineto{\pgfqpoint{0.431912in}{1.491802in}}%
\pgfpathlineto{\pgfqpoint{0.428385in}{1.494011in}}%
\pgfpathlineto{\pgfqpoint{0.414605in}{1.487528in}}%
\pgfpathlineto{\pgfqpoint{0.408114in}{1.487902in}}%
\pgfpathlineto{\pgfqpoint{0.405051in}{1.484139in}}%
\pgfpathlineto{\pgfqpoint{0.407088in}{1.477930in}}%
\pgfpathlineto{\pgfqpoint{0.399966in}{1.478464in}}%
\pgfpathlineto{\pgfqpoint{0.397986in}{1.484025in}}%
\pgfpathlineto{\pgfqpoint{0.391596in}{1.486250in}}%
\pgfpathlineto{\pgfqpoint{0.391169in}{1.489496in}}%
\pgfpathlineto{\pgfqpoint{0.386307in}{1.493716in}}%
\pgfpathlineto{\pgfqpoint{0.385689in}{1.509253in}}%
\pgfpathlineto{\pgfqpoint{0.378145in}{1.511673in}}%
\pgfpathlineto{\pgfqpoint{0.381028in}{1.514238in}}%
\pgfpathlineto{\pgfqpoint{0.381897in}{1.521054in}}%
\pgfpathlineto{\pgfqpoint{0.380582in}{1.527032in}}%
\pgfpathlineto{\pgfqpoint{0.384988in}{1.532462in}}%
\pgfpathlineto{\pgfqpoint{0.385318in}{1.540362in}}%
\pgfpathlineto{\pgfqpoint{0.393613in}{1.542963in}}%
\pgfpathlineto{\pgfqpoint{0.397486in}{1.548285in}}%
\pgfpathlineto{\pgfqpoint{0.455007in}{1.530523in}}%
\pgfpathlineto{\pgfqpoint{0.474662in}{1.592743in}}%
\pgfpathlineto{\pgfqpoint{0.476029in}{1.597225in}}%
\pgfpathlineto{\pgfqpoint{0.485002in}{1.594976in}}%
\pgfpathlineto{\pgfqpoint{0.487657in}{1.599535in}}%
\pgfpathlineto{\pgfqpoint{0.493766in}{1.605432in}}%
\pgfpathlineto{\pgfqpoint{0.493741in}{1.612121in}}%
\pgfpathlineto{\pgfqpoint{0.490196in}{1.618627in}}%
\pgfpathlineto{\pgfqpoint{0.492514in}{1.625981in}}%
\pgfpathlineto{\pgfqpoint{0.499509in}{1.627848in}}%
\pgfpathlineto{\pgfqpoint{0.530203in}{1.618670in}}%
\pgfpathlineto{\pgfqpoint{0.524578in}{1.602426in}}%
\pgfpathlineto{\pgfqpoint{0.513001in}{1.563845in}}%
\pgfpathlineto{\pgfqpoint{0.534635in}{1.557261in}}%
\pgfpathlineto{\pgfqpoint{0.520911in}{1.509711in}}%
\pgfpathlineto{\pgfqpoint{0.562393in}{1.497838in}}%
\pgfpathlineto{\pgfqpoint{0.547844in}{1.446415in}}%
\pgfpathlineto{\pgfqpoint{0.535789in}{1.402977in}}%
\pgfpathlineto{\pgfqpoint{0.521615in}{1.353758in}}%
\pgfpathclose%
\pgfusepath{fill}%
\end{pgfscope}%
\begin{pgfscope}%
\pgfpathrectangle{\pgfqpoint{0.100000in}{0.100000in}}{\pgfqpoint{3.007045in}{1.925000in}}%
\pgfusepath{clip}%
\pgfsetbuttcap%
\pgfsetmiterjoin%
\definecolor{currentfill}{rgb}{0.596078,0.834510,0.578824}%
\pgfsetfillcolor{currentfill}%
\pgfsetlinewidth{0.000000pt}%
\definecolor{currentstroke}{rgb}{0.000000,0.000000,0.000000}%
\pgfsetstrokecolor{currentstroke}%
\pgfsetstrokeopacity{0.000000}%
\pgfsetdash{}{0pt}%
\pgfpathmoveto{\pgfqpoint{2.213486in}{1.315239in}}%
\pgfpathlineto{\pgfqpoint{2.215428in}{1.298031in}}%
\pgfpathlineto{\pgfqpoint{2.195043in}{1.295965in}}%
\pgfpathlineto{\pgfqpoint{2.172150in}{1.293918in}}%
\pgfpathlineto{\pgfqpoint{2.168897in}{1.329206in}}%
\pgfpathlineto{\pgfqpoint{2.180297in}{1.338873in}}%
\pgfpathlineto{\pgfqpoint{2.184660in}{1.346048in}}%
\pgfpathlineto{\pgfqpoint{2.188143in}{1.357873in}}%
\pgfpathlineto{\pgfqpoint{2.193342in}{1.366778in}}%
\pgfpathlineto{\pgfqpoint{2.200252in}{1.367456in}}%
\pgfpathlineto{\pgfqpoint{2.201376in}{1.356174in}}%
\pgfpathlineto{\pgfqpoint{2.223806in}{1.358262in}}%
\pgfpathlineto{\pgfqpoint{2.225748in}{1.340539in}}%
\pgfpathlineto{\pgfqpoint{2.224488in}{1.337759in}}%
\pgfpathlineto{\pgfqpoint{2.211201in}{1.336530in}}%
\pgfpathlineto{\pgfqpoint{2.213486in}{1.315239in}}%
\pgfpathclose%
\pgfusepath{fill}%
\end{pgfscope}%
\begin{pgfscope}%
\pgfpathrectangle{\pgfqpoint{0.100000in}{0.100000in}}{\pgfqpoint{3.007045in}{1.925000in}}%
\pgfusepath{clip}%
\pgfsetbuttcap%
\pgfsetmiterjoin%
\definecolor{currentfill}{rgb}{0.690396,0.875832,0.665313}%
\pgfsetfillcolor{currentfill}%
\pgfsetlinewidth{0.000000pt}%
\definecolor{currentstroke}{rgb}{0.000000,0.000000,0.000000}%
\pgfsetstrokecolor{currentstroke}%
\pgfsetstrokeopacity{0.000000}%
\pgfsetdash{}{0pt}%
\pgfpathmoveto{\pgfqpoint{2.013078in}{0.656435in}}%
\pgfpathlineto{\pgfqpoint{2.013470in}{0.675004in}}%
\pgfpathlineto{\pgfqpoint{2.011898in}{0.691368in}}%
\pgfpathlineto{\pgfqpoint{2.012854in}{0.699525in}}%
\pgfpathlineto{\pgfqpoint{2.016613in}{0.701553in}}%
\pgfpathlineto{\pgfqpoint{2.018058in}{0.708669in}}%
\pgfpathlineto{\pgfqpoint{2.020779in}{0.710545in}}%
\pgfpathlineto{\pgfqpoint{2.027190in}{0.716066in}}%
\pgfpathlineto{\pgfqpoint{2.040123in}{0.721921in}}%
\pgfpathlineto{\pgfqpoint{2.048959in}{0.728755in}}%
\pgfpathlineto{\pgfqpoint{2.052237in}{0.736873in}}%
\pgfpathlineto{\pgfqpoint{2.065317in}{0.737960in}}%
\pgfpathlineto{\pgfqpoint{2.065136in}{0.740946in}}%
\pgfpathlineto{\pgfqpoint{2.087936in}{0.742436in}}%
\pgfpathlineto{\pgfqpoint{2.090671in}{0.695575in}}%
\pgfpathlineto{\pgfqpoint{2.067842in}{0.694087in}}%
\pgfpathlineto{\pgfqpoint{2.068533in}{0.682516in}}%
\pgfpathlineto{\pgfqpoint{2.072910in}{0.680552in}}%
\pgfpathlineto{\pgfqpoint{2.073893in}{0.665116in}}%
\pgfpathlineto{\pgfqpoint{2.055851in}{0.662822in}}%
\pgfpathlineto{\pgfqpoint{2.047327in}{0.661747in}}%
\pgfpathlineto{\pgfqpoint{2.038852in}{0.657801in}}%
\pgfpathlineto{\pgfqpoint{2.013078in}{0.656435in}}%
\pgfpathclose%
\pgfusepath{fill}%
\end{pgfscope}%
\begin{pgfscope}%
\pgfpathrectangle{\pgfqpoint{0.100000in}{0.100000in}}{\pgfqpoint{3.007045in}{1.925000in}}%
\pgfusepath{clip}%
\pgfsetbuttcap%
\pgfsetmiterjoin%
\definecolor{currentfill}{rgb}{0.643645,0.856148,0.619792}%
\pgfsetfillcolor{currentfill}%
\pgfsetlinewidth{0.000000pt}%
\definecolor{currentstroke}{rgb}{0.000000,0.000000,0.000000}%
\pgfsetstrokecolor{currentstroke}%
\pgfsetstrokeopacity{0.000000}%
\pgfsetdash{}{0pt}%
\pgfpathmoveto{\pgfqpoint{1.886783in}{0.652756in}}%
\pgfpathlineto{\pgfqpoint{1.890709in}{0.647179in}}%
\pgfpathlineto{\pgfqpoint{1.896682in}{0.645552in}}%
\pgfpathlineto{\pgfqpoint{1.901476in}{0.640393in}}%
\pgfpathlineto{\pgfqpoint{1.890720in}{0.630209in}}%
\pgfpathlineto{\pgfqpoint{1.886870in}{0.628687in}}%
\pgfpathlineto{\pgfqpoint{1.883960in}{0.629796in}}%
\pgfpathlineto{\pgfqpoint{1.872399in}{0.629573in}}%
\pgfpathlineto{\pgfqpoint{1.872143in}{0.641119in}}%
\pgfpathlineto{\pgfqpoint{1.866438in}{0.640996in}}%
\pgfpathlineto{\pgfqpoint{1.866139in}{0.652691in}}%
\pgfpathlineto{\pgfqpoint{1.860335in}{0.652602in}}%
\pgfpathlineto{\pgfqpoint{1.860218in}{0.661283in}}%
\pgfpathlineto{\pgfqpoint{1.835531in}{0.660856in}}%
\pgfpathlineto{\pgfqpoint{1.827733in}{0.669764in}}%
\pgfpathlineto{\pgfqpoint{1.826203in}{0.670688in}}%
\pgfpathlineto{\pgfqpoint{1.825570in}{0.729490in}}%
\pgfpathlineto{\pgfqpoint{1.825469in}{0.738647in}}%
\pgfpathlineto{\pgfqpoint{1.838686in}{0.738812in}}%
\pgfpathlineto{\pgfqpoint{1.884006in}{0.739484in}}%
\pgfpathlineto{\pgfqpoint{1.898651in}{0.739626in}}%
\pgfpathlineto{\pgfqpoint{1.899049in}{0.722750in}}%
\pgfpathlineto{\pgfqpoint{1.893306in}{0.722562in}}%
\pgfpathlineto{\pgfqpoint{1.890590in}{0.711973in}}%
\pgfpathlineto{\pgfqpoint{1.890952in}{0.702305in}}%
\pgfpathlineto{\pgfqpoint{1.896662in}{0.702398in}}%
\pgfpathlineto{\pgfqpoint{1.897220in}{0.687919in}}%
\pgfpathlineto{\pgfqpoint{1.895056in}{0.682024in}}%
\pgfpathlineto{\pgfqpoint{1.888095in}{0.681930in}}%
\pgfpathlineto{\pgfqpoint{1.891446in}{0.668175in}}%
\pgfpathlineto{\pgfqpoint{1.888359in}{0.662427in}}%
\pgfpathlineto{\pgfqpoint{1.886783in}{0.652756in}}%
\pgfpathclose%
\pgfusepath{fill}%
\end{pgfscope}%
\begin{pgfscope}%
\pgfpathrectangle{\pgfqpoint{0.100000in}{0.100000in}}{\pgfqpoint{3.007045in}{1.925000in}}%
\pgfusepath{clip}%
\pgfsetbuttcap%
\pgfsetmiterjoin%
\definecolor{currentfill}{rgb}{0.529642,0.803506,0.524198}%
\pgfsetfillcolor{currentfill}%
\pgfsetlinewidth{0.000000pt}%
\definecolor{currentstroke}{rgb}{0.000000,0.000000,0.000000}%
\pgfsetstrokecolor{currentstroke}%
\pgfsetstrokeopacity{0.000000}%
\pgfsetdash{}{0pt}%
\pgfpathmoveto{\pgfqpoint{2.399943in}{0.955887in}}%
\pgfpathlineto{\pgfqpoint{2.388077in}{0.945524in}}%
\pgfpathlineto{\pgfqpoint{2.379259in}{0.944820in}}%
\pgfpathlineto{\pgfqpoint{2.378270in}{0.952775in}}%
\pgfpathlineto{\pgfqpoint{2.374820in}{0.956140in}}%
\pgfpathlineto{\pgfqpoint{2.369682in}{0.964768in}}%
\pgfpathlineto{\pgfqpoint{2.376395in}{0.971244in}}%
\pgfpathlineto{\pgfqpoint{2.370660in}{0.984801in}}%
\pgfpathlineto{\pgfqpoint{2.372662in}{0.995532in}}%
\pgfpathlineto{\pgfqpoint{2.380497in}{0.997923in}}%
\pgfpathlineto{\pgfqpoint{2.388640in}{1.002575in}}%
\pgfpathlineto{\pgfqpoint{2.392615in}{1.001420in}}%
\pgfpathlineto{\pgfqpoint{2.394844in}{0.994687in}}%
\pgfpathlineto{\pgfqpoint{2.399497in}{0.998207in}}%
\pgfpathlineto{\pgfqpoint{2.404364in}{0.992821in}}%
\pgfpathlineto{\pgfqpoint{2.400322in}{0.986669in}}%
\pgfpathlineto{\pgfqpoint{2.406942in}{0.983024in}}%
\pgfpathlineto{\pgfqpoint{2.416656in}{0.974681in}}%
\pgfpathlineto{\pgfqpoint{2.414244in}{0.964014in}}%
\pgfpathlineto{\pgfqpoint{2.404400in}{0.961065in}}%
\pgfpathlineto{\pgfqpoint{2.399943in}{0.955887in}}%
\pgfpathclose%
\pgfusepath{fill}%
\end{pgfscope}%
\begin{pgfscope}%
\pgfpathrectangle{\pgfqpoint{0.100000in}{0.100000in}}{\pgfqpoint{3.007045in}{1.925000in}}%
\pgfusepath{clip}%
\pgfsetbuttcap%
\pgfsetmiterjoin%
\definecolor{currentfill}{rgb}{0.562860,0.819008,0.551511}%
\pgfsetfillcolor{currentfill}%
\pgfsetlinewidth{0.000000pt}%
\definecolor{currentstroke}{rgb}{0.000000,0.000000,0.000000}%
\pgfsetstrokecolor{currentstroke}%
\pgfsetstrokeopacity{0.000000}%
\pgfsetdash{}{0pt}%
\pgfpathmoveto{\pgfqpoint{1.092509in}{1.504079in}}%
\pgfpathlineto{\pgfqpoint{1.102940in}{1.502630in}}%
\pgfpathlineto{\pgfqpoint{1.103807in}{1.508400in}}%
\pgfpathlineto{\pgfqpoint{1.112286in}{1.507103in}}%
\pgfpathlineto{\pgfqpoint{1.113003in}{1.512751in}}%
\pgfpathlineto{\pgfqpoint{1.118845in}{1.511867in}}%
\pgfpathlineto{\pgfqpoint{1.119707in}{1.517659in}}%
\pgfpathlineto{\pgfqpoint{1.131047in}{1.515926in}}%
\pgfpathlineto{\pgfqpoint{1.131931in}{1.521748in}}%
\pgfpathlineto{\pgfqpoint{1.197829in}{1.512182in}}%
\pgfpathlineto{\pgfqpoint{1.193681in}{1.468364in}}%
\pgfpathlineto{\pgfqpoint{1.173551in}{1.471182in}}%
\pgfpathlineto{\pgfqpoint{1.170426in}{1.469676in}}%
\pgfpathlineto{\pgfqpoint{1.143778in}{1.473368in}}%
\pgfpathlineto{\pgfqpoint{1.135060in}{1.473922in}}%
\pgfpathlineto{\pgfqpoint{1.118840in}{1.481195in}}%
\pgfpathlineto{\pgfqpoint{1.119290in}{1.484034in}}%
\pgfpathlineto{\pgfqpoint{1.108398in}{1.488604in}}%
\pgfpathlineto{\pgfqpoint{1.097179in}{1.490393in}}%
\pgfpathlineto{\pgfqpoint{1.098004in}{1.495582in}}%
\pgfpathlineto{\pgfqpoint{1.092509in}{1.504079in}}%
\pgfpathclose%
\pgfusepath{fill}%
\end{pgfscope}%
\begin{pgfscope}%
\pgfpathrectangle{\pgfqpoint{0.100000in}{0.100000in}}{\pgfqpoint{3.007045in}{1.925000in}}%
\pgfusepath{clip}%
\pgfsetbuttcap%
\pgfsetmiterjoin%
\definecolor{currentfill}{rgb}{0.671696,0.867958,0.647105}%
\pgfsetfillcolor{currentfill}%
\pgfsetlinewidth{0.000000pt}%
\definecolor{currentstroke}{rgb}{0.000000,0.000000,0.000000}%
\pgfsetstrokecolor{currentstroke}%
\pgfsetstrokeopacity{0.000000}%
\pgfsetdash{}{0pt}%
\pgfpathmoveto{\pgfqpoint{2.119150in}{1.400006in}}%
\pgfpathlineto{\pgfqpoint{2.073275in}{1.396646in}}%
\pgfpathlineto{\pgfqpoint{2.061841in}{1.396171in}}%
\pgfpathlineto{\pgfqpoint{2.058743in}{1.447959in}}%
\pgfpathlineto{\pgfqpoint{2.087655in}{1.449833in}}%
\pgfpathlineto{\pgfqpoint{2.088107in}{1.444073in}}%
\pgfpathlineto{\pgfqpoint{2.099536in}{1.444895in}}%
\pgfpathlineto{\pgfqpoint{2.117083in}{1.446231in}}%
\pgfpathlineto{\pgfqpoint{2.116939in}{1.440991in}}%
\pgfpathlineto{\pgfqpoint{2.114427in}{1.436064in}}%
\pgfpathlineto{\pgfqpoint{2.113005in}{1.425983in}}%
\pgfpathlineto{\pgfqpoint{2.119150in}{1.400006in}}%
\pgfpathclose%
\pgfusepath{fill}%
\end{pgfscope}%
\begin{pgfscope}%
\pgfpathrectangle{\pgfqpoint{0.100000in}{0.100000in}}{\pgfqpoint{3.007045in}{1.925000in}}%
\pgfusepath{clip}%
\pgfsetbuttcap%
\pgfsetmiterjoin%
\definecolor{currentfill}{rgb}{0.760523,0.905359,0.733595}%
\pgfsetfillcolor{currentfill}%
\pgfsetlinewidth{0.000000pt}%
\definecolor{currentstroke}{rgb}{0.000000,0.000000,0.000000}%
\pgfsetstrokecolor{currentstroke}%
\pgfsetstrokeopacity{0.000000}%
\pgfsetdash{}{0pt}%
\pgfpathmoveto{\pgfqpoint{0.618157in}{0.899792in}}%
\pgfpathlineto{\pgfqpoint{0.613032in}{0.877119in}}%
\pgfpathlineto{\pgfqpoint{0.611833in}{0.877391in}}%
\pgfpathlineto{\pgfqpoint{0.604721in}{0.848024in}}%
\pgfpathlineto{\pgfqpoint{0.548111in}{0.856043in}}%
\pgfpathlineto{\pgfqpoint{0.548819in}{0.861290in}}%
\pgfpathlineto{\pgfqpoint{0.543625in}{0.866151in}}%
\pgfpathlineto{\pgfqpoint{0.546518in}{0.879093in}}%
\pgfpathlineto{\pgfqpoint{0.547290in}{0.888052in}}%
\pgfpathlineto{\pgfqpoint{0.545659in}{0.899091in}}%
\pgfpathlineto{\pgfqpoint{0.540654in}{0.912060in}}%
\pgfpathlineto{\pgfqpoint{0.536458in}{0.917204in}}%
\pgfpathlineto{\pgfqpoint{0.538502in}{0.921211in}}%
\pgfpathlineto{\pgfqpoint{0.543010in}{0.923557in}}%
\pgfpathlineto{\pgfqpoint{0.550214in}{0.920766in}}%
\pgfpathlineto{\pgfqpoint{0.556193in}{0.915257in}}%
\pgfpathlineto{\pgfqpoint{0.567402in}{0.912130in}}%
\pgfpathlineto{\pgfqpoint{0.618157in}{0.899792in}}%
\pgfpathclose%
\pgfusepath{fill}%
\end{pgfscope}%
\begin{pgfscope}%
\pgfpathrectangle{\pgfqpoint{0.100000in}{0.100000in}}{\pgfqpoint{3.007045in}{1.925000in}}%
\pgfusepath{clip}%
\pgfsetbuttcap%
\pgfsetmiterjoin%
\definecolor{currentfill}{rgb}{0.338824,0.711726,0.405844}%
\pgfsetfillcolor{currentfill}%
\pgfsetlinewidth{0.000000pt}%
\definecolor{currentstroke}{rgb}{0.000000,0.000000,0.000000}%
\pgfsetstrokecolor{currentstroke}%
\pgfsetstrokeopacity{0.000000}%
\pgfsetdash{}{0pt}%
\pgfpathmoveto{\pgfqpoint{2.000756in}{1.600279in}}%
\pgfpathlineto{\pgfqpoint{1.988873in}{1.599607in}}%
\pgfpathlineto{\pgfqpoint{1.988231in}{1.611105in}}%
\pgfpathlineto{\pgfqpoint{1.982576in}{1.610763in}}%
\pgfpathlineto{\pgfqpoint{1.982301in}{1.616494in}}%
\pgfpathlineto{\pgfqpoint{1.976511in}{1.616232in}}%
\pgfpathlineto{\pgfqpoint{1.975507in}{1.638992in}}%
\pgfpathlineto{\pgfqpoint{1.980609in}{1.637782in}}%
\pgfpathlineto{\pgfqpoint{1.984352in}{1.640321in}}%
\pgfpathlineto{\pgfqpoint{1.998858in}{1.646114in}}%
\pgfpathlineto{\pgfqpoint{2.009289in}{1.656368in}}%
\pgfpathlineto{\pgfqpoint{2.026300in}{1.659095in}}%
\pgfpathlineto{\pgfqpoint{2.033497in}{1.663365in}}%
\pgfpathlineto{\pgfqpoint{2.037806in}{1.669257in}}%
\pgfpathlineto{\pgfqpoint{2.044910in}{1.670501in}}%
\pgfpathlineto{\pgfqpoint{2.047218in}{1.673072in}}%
\pgfpathlineto{\pgfqpoint{2.048453in}{1.655274in}}%
\pgfpathlineto{\pgfqpoint{2.051916in}{1.649810in}}%
\pgfpathlineto{\pgfqpoint{2.046234in}{1.649481in}}%
\pgfpathlineto{\pgfqpoint{2.047301in}{1.632207in}}%
\pgfpathlineto{\pgfqpoint{2.048726in}{1.610846in}}%
\pgfpathlineto{\pgfqpoint{2.043917in}{1.613272in}}%
\pgfpathlineto{\pgfqpoint{2.004866in}{1.621636in}}%
\pgfpathlineto{\pgfqpoint{2.006057in}{1.600572in}}%
\pgfpathlineto{\pgfqpoint{2.000756in}{1.600279in}}%
\pgfpathclose%
\pgfusepath{fill}%
\end{pgfscope}%
\begin{pgfscope}%
\pgfpathrectangle{\pgfqpoint{0.100000in}{0.100000in}}{\pgfqpoint{3.007045in}{1.925000in}}%
\pgfusepath{clip}%
\pgfsetbuttcap%
\pgfsetmiterjoin%
\definecolor{currentfill}{rgb}{0.601615,0.837093,0.583376}%
\pgfsetfillcolor{currentfill}%
\pgfsetlinewidth{0.000000pt}%
\definecolor{currentstroke}{rgb}{0.000000,0.000000,0.000000}%
\pgfsetstrokecolor{currentstroke}%
\pgfsetstrokeopacity{0.000000}%
\pgfsetdash{}{0pt}%
\pgfpathmoveto{\pgfqpoint{1.330895in}{1.191503in}}%
\pgfpathlineto{\pgfqpoint{1.358836in}{1.188865in}}%
\pgfpathlineto{\pgfqpoint{1.355807in}{1.160158in}}%
\pgfpathlineto{\pgfqpoint{1.353166in}{1.131470in}}%
\pgfpathlineto{\pgfqpoint{1.351745in}{1.120688in}}%
\pgfpathlineto{\pgfqpoint{1.334664in}{1.121680in}}%
\pgfpathlineto{\pgfqpoint{1.306529in}{1.124802in}}%
\pgfpathlineto{\pgfqpoint{1.308723in}{1.147516in}}%
\pgfpathlineto{\pgfqpoint{1.325845in}{1.145704in}}%
\pgfpathlineto{\pgfqpoint{1.330895in}{1.191503in}}%
\pgfpathclose%
\pgfusepath{fill}%
\end{pgfscope}%
\begin{pgfscope}%
\pgfpathrectangle{\pgfqpoint{0.100000in}{0.100000in}}{\pgfqpoint{3.007045in}{1.925000in}}%
\pgfusepath{clip}%
\pgfsetbuttcap%
\pgfsetmiterjoin%
\definecolor{currentfill}{rgb}{0.573933,0.824175,0.560615}%
\pgfsetfillcolor{currentfill}%
\pgfsetlinewidth{0.000000pt}%
\definecolor{currentstroke}{rgb}{0.000000,0.000000,0.000000}%
\pgfsetstrokecolor{currentstroke}%
\pgfsetstrokeopacity{0.000000}%
\pgfsetdash{}{0pt}%
\pgfpathmoveto{\pgfqpoint{1.768296in}{1.069654in}}%
\pgfpathlineto{\pgfqpoint{1.767740in}{1.049552in}}%
\pgfpathlineto{\pgfqpoint{1.722263in}{1.049928in}}%
\pgfpathlineto{\pgfqpoint{1.722556in}{1.070090in}}%
\pgfpathlineto{\pgfqpoint{1.745387in}{1.069764in}}%
\pgfpathlineto{\pgfqpoint{1.746105in}{1.095854in}}%
\pgfpathlineto{\pgfqpoint{1.746586in}{1.115888in}}%
\pgfpathlineto{\pgfqpoint{1.769426in}{1.115817in}}%
\pgfpathlineto{\pgfqpoint{1.768947in}{1.112944in}}%
\pgfpathlineto{\pgfqpoint{1.768960in}{1.084214in}}%
\pgfpathlineto{\pgfqpoint{1.768296in}{1.069654in}}%
\pgfpathclose%
\pgfusepath{fill}%
\end{pgfscope}%
\begin{pgfscope}%
\pgfpathrectangle{\pgfqpoint{0.100000in}{0.100000in}}{\pgfqpoint{3.007045in}{1.925000in}}%
\pgfusepath{clip}%
\pgfsetbuttcap%
\pgfsetmiterjoin%
\definecolor{currentfill}{rgb}{0.562860,0.819008,0.551511}%
\pgfsetfillcolor{currentfill}%
\pgfsetlinewidth{0.000000pt}%
\definecolor{currentstroke}{rgb}{0.000000,0.000000,0.000000}%
\pgfsetstrokecolor{currentstroke}%
\pgfsetstrokeopacity{0.000000}%
\pgfsetdash{}{0pt}%
\pgfpathmoveto{\pgfqpoint{1.980060in}{1.429597in}}%
\pgfpathlineto{\pgfqpoint{1.954322in}{1.428443in}}%
\pgfpathlineto{\pgfqpoint{1.953526in}{1.433541in}}%
\pgfpathlineto{\pgfqpoint{1.934699in}{1.432796in}}%
\pgfpathlineto{\pgfqpoint{1.928977in}{1.432597in}}%
\pgfpathlineto{\pgfqpoint{1.928151in}{1.455464in}}%
\pgfpathlineto{\pgfqpoint{1.949409in}{1.456290in}}%
\pgfpathlineto{\pgfqpoint{1.942334in}{1.465087in}}%
\pgfpathlineto{\pgfqpoint{1.947232in}{1.465292in}}%
\pgfpathlineto{\pgfqpoint{1.947712in}{1.470634in}}%
\pgfpathlineto{\pgfqpoint{1.963568in}{1.471696in}}%
\pgfpathlineto{\pgfqpoint{1.966485in}{1.477622in}}%
\pgfpathlineto{\pgfqpoint{1.994624in}{1.478781in}}%
\pgfpathlineto{\pgfqpoint{1.994330in}{1.484956in}}%
\pgfpathlineto{\pgfqpoint{2.013699in}{1.486047in}}%
\pgfpathlineto{\pgfqpoint{2.008307in}{1.475250in}}%
\pgfpathlineto{\pgfqpoint{2.012462in}{1.466036in}}%
\pgfpathlineto{\pgfqpoint{2.012501in}{1.460181in}}%
\pgfpathlineto{\pgfqpoint{2.017195in}{1.455365in}}%
\pgfpathlineto{\pgfqpoint{2.021528in}{1.446204in}}%
\pgfpathlineto{\pgfqpoint{1.996384in}{1.444829in}}%
\pgfpathlineto{\pgfqpoint{1.996750in}{1.439088in}}%
\pgfpathlineto{\pgfqpoint{1.979503in}{1.438168in}}%
\pgfpathlineto{\pgfqpoint{1.980060in}{1.429597in}}%
\pgfpathclose%
\pgfusepath{fill}%
\end{pgfscope}%
\begin{pgfscope}%
\pgfpathrectangle{\pgfqpoint{0.100000in}{0.100000in}}{\pgfqpoint{3.007045in}{1.925000in}}%
\pgfusepath{clip}%
\pgfsetbuttcap%
\pgfsetmiterjoin%
\definecolor{currentfill}{rgb}{0.426667,0.754787,0.448904}%
\pgfsetfillcolor{currentfill}%
\pgfsetlinewidth{0.000000pt}%
\definecolor{currentstroke}{rgb}{0.000000,0.000000,0.000000}%
\pgfsetstrokecolor{currentstroke}%
\pgfsetstrokeopacity{0.000000}%
\pgfsetdash{}{0pt}%
\pgfpathmoveto{\pgfqpoint{1.045450in}{1.802919in}}%
\pgfpathlineto{\pgfqpoint{1.043735in}{1.793392in}}%
\pgfpathlineto{\pgfqpoint{1.018457in}{1.797935in}}%
\pgfpathlineto{\pgfqpoint{1.016942in}{1.805156in}}%
\pgfpathlineto{\pgfqpoint{1.011784in}{1.809131in}}%
\pgfpathlineto{\pgfqpoint{0.975124in}{1.816423in}}%
\pgfpathlineto{\pgfqpoint{0.966264in}{1.820814in}}%
\pgfpathlineto{\pgfqpoint{0.967054in}{1.825877in}}%
\pgfpathlineto{\pgfqpoint{0.961768in}{1.833418in}}%
\pgfpathlineto{\pgfqpoint{0.964152in}{1.834807in}}%
\pgfpathlineto{\pgfqpoint{0.959285in}{1.841945in}}%
\pgfpathlineto{\pgfqpoint{0.961059in}{1.847574in}}%
\pgfpathlineto{\pgfqpoint{0.949618in}{1.854808in}}%
\pgfpathlineto{\pgfqpoint{0.953218in}{1.860951in}}%
\pgfpathlineto{\pgfqpoint{0.949121in}{1.871112in}}%
\pgfpathlineto{\pgfqpoint{0.943633in}{1.869658in}}%
\pgfpathlineto{\pgfqpoint{0.941225in}{1.874500in}}%
\pgfpathlineto{\pgfqpoint{0.941334in}{1.882540in}}%
\pgfpathlineto{\pgfqpoint{0.981051in}{1.874255in}}%
\pgfpathlineto{\pgfqpoint{1.010982in}{1.868345in}}%
\pgfpathlineto{\pgfqpoint{1.061807in}{1.858712in}}%
\pgfpathlineto{\pgfqpoint{1.052373in}{1.807583in}}%
\pgfpathlineto{\pgfqpoint{1.046502in}{1.808652in}}%
\pgfpathlineto{\pgfqpoint{1.045450in}{1.802919in}}%
\pgfpathclose%
\pgfusepath{fill}%
\end{pgfscope}%
\begin{pgfscope}%
\pgfpathrectangle{\pgfqpoint{0.100000in}{0.100000in}}{\pgfqpoint{3.007045in}{1.925000in}}%
\pgfusepath{clip}%
\pgfsetbuttcap%
\pgfsetmiterjoin%
\definecolor{currentfill}{rgb}{0.671696,0.867958,0.647105}%
\pgfsetfillcolor{currentfill}%
\pgfsetlinewidth{0.000000pt}%
\definecolor{currentstroke}{rgb}{0.000000,0.000000,0.000000}%
\pgfsetstrokecolor{currentstroke}%
\pgfsetstrokeopacity{0.000000}%
\pgfsetdash{}{0pt}%
\pgfpathmoveto{\pgfqpoint{1.522557in}{1.725806in}}%
\pgfpathlineto{\pgfqpoint{1.465127in}{1.729765in}}%
\pgfpathlineto{\pgfqpoint{1.466009in}{1.741430in}}%
\pgfpathlineto{\pgfqpoint{1.463751in}{1.741609in}}%
\pgfpathlineto{\pgfqpoint{1.465533in}{1.764842in}}%
\pgfpathlineto{\pgfqpoint{1.463268in}{1.765000in}}%
\pgfpathlineto{\pgfqpoint{1.464168in}{1.776657in}}%
\pgfpathlineto{\pgfqpoint{1.452658in}{1.777477in}}%
\pgfpathlineto{\pgfqpoint{1.453609in}{1.789081in}}%
\pgfpathlineto{\pgfqpoint{1.457171in}{1.788804in}}%
\pgfpathlineto{\pgfqpoint{1.457650in}{1.794593in}}%
\pgfpathlineto{\pgfqpoint{1.463375in}{1.794138in}}%
\pgfpathlineto{\pgfqpoint{1.464391in}{1.806947in}}%
\pgfpathlineto{\pgfqpoint{1.508366in}{1.803801in}}%
\pgfpathlineto{\pgfqpoint{1.544920in}{1.801546in}}%
\pgfpathlineto{\pgfqpoint{1.543838in}{1.782832in}}%
\pgfpathlineto{\pgfqpoint{1.544810in}{1.771147in}}%
\pgfpathlineto{\pgfqpoint{1.539033in}{1.771449in}}%
\pgfpathlineto{\pgfqpoint{1.540108in}{1.759776in}}%
\pgfpathlineto{\pgfqpoint{1.538782in}{1.736495in}}%
\pgfpathlineto{\pgfqpoint{1.539856in}{1.724787in}}%
\pgfpathlineto{\pgfqpoint{1.522557in}{1.725806in}}%
\pgfpathclose%
\pgfusepath{fill}%
\end{pgfscope}%
\begin{pgfscope}%
\pgfpathrectangle{\pgfqpoint{0.100000in}{0.100000in}}{\pgfqpoint{3.007045in}{1.925000in}}%
\pgfusepath{clip}%
\pgfsetbuttcap%
\pgfsetmiterjoin%
\definecolor{currentfill}{rgb}{0.643645,0.856148,0.619792}%
\pgfsetfillcolor{currentfill}%
\pgfsetlinewidth{0.000000pt}%
\definecolor{currentstroke}{rgb}{0.000000,0.000000,0.000000}%
\pgfsetstrokecolor{currentstroke}%
\pgfsetstrokeopacity{0.000000}%
\pgfsetdash{}{0pt}%
\pgfpathmoveto{\pgfqpoint{1.591111in}{1.462290in}}%
\pgfpathlineto{\pgfqpoint{1.590530in}{1.433489in}}%
\pgfpathlineto{\pgfqpoint{1.566436in}{1.434578in}}%
\pgfpathlineto{\pgfqpoint{1.550851in}{1.435352in}}%
\pgfpathlineto{\pgfqpoint{1.551671in}{1.451429in}}%
\pgfpathlineto{\pgfqpoint{1.541165in}{1.448356in}}%
\pgfpathlineto{\pgfqpoint{1.531144in}{1.449134in}}%
\pgfpathlineto{\pgfqpoint{1.522608in}{1.451050in}}%
\pgfpathlineto{\pgfqpoint{1.517464in}{1.451554in}}%
\pgfpathlineto{\pgfqpoint{1.516674in}{1.460211in}}%
\pgfpathlineto{\pgfqpoint{1.517880in}{1.481436in}}%
\pgfpathlineto{\pgfqpoint{1.517980in}{1.483159in}}%
\pgfpathlineto{\pgfqpoint{1.538292in}{1.482046in}}%
\pgfpathlineto{\pgfqpoint{1.542308in}{1.477127in}}%
\pgfpathlineto{\pgfqpoint{1.546378in}{1.477641in}}%
\pgfpathlineto{\pgfqpoint{1.552579in}{1.474853in}}%
\pgfpathlineto{\pgfqpoint{1.551371in}{1.482830in}}%
\pgfpathlineto{\pgfqpoint{1.555401in}{1.480969in}}%
\pgfpathlineto{\pgfqpoint{1.586215in}{1.479799in}}%
\pgfpathlineto{\pgfqpoint{1.585428in}{1.462527in}}%
\pgfpathlineto{\pgfqpoint{1.591111in}{1.462290in}}%
\pgfpathclose%
\pgfusepath{fill}%
\end{pgfscope}%
\begin{pgfscope}%
\pgfpathrectangle{\pgfqpoint{0.100000in}{0.100000in}}{\pgfqpoint{3.007045in}{1.925000in}}%
\pgfusepath{clip}%
\pgfsetbuttcap%
\pgfsetmiterjoin%
\definecolor{currentfill}{rgb}{0.568397,0.821592,0.556063}%
\pgfsetfillcolor{currentfill}%
\pgfsetlinewidth{0.000000pt}%
\definecolor{currentstroke}{rgb}{0.000000,0.000000,0.000000}%
\pgfsetstrokecolor{currentstroke}%
\pgfsetstrokeopacity{0.000000}%
\pgfsetdash{}{0pt}%
\pgfpathmoveto{\pgfqpoint{2.429643in}{1.109432in}}%
\pgfpathlineto{\pgfqpoint{2.435308in}{1.112304in}}%
\pgfpathlineto{\pgfqpoint{2.440031in}{1.105220in}}%
\pgfpathlineto{\pgfqpoint{2.441207in}{1.098026in}}%
\pgfpathlineto{\pgfqpoint{2.446820in}{1.093542in}}%
\pgfpathlineto{\pgfqpoint{2.458232in}{1.094688in}}%
\pgfpathlineto{\pgfqpoint{2.456556in}{1.087012in}}%
\pgfpathlineto{\pgfqpoint{2.453185in}{1.084073in}}%
\pgfpathlineto{\pgfqpoint{2.450852in}{1.085445in}}%
\pgfpathlineto{\pgfqpoint{2.435154in}{1.067105in}}%
\pgfpathlineto{\pgfqpoint{2.423037in}{1.058784in}}%
\pgfpathlineto{\pgfqpoint{2.417226in}{1.060057in}}%
\pgfpathlineto{\pgfqpoint{2.412270in}{1.064935in}}%
\pgfpathlineto{\pgfqpoint{2.410452in}{1.072460in}}%
\pgfpathlineto{\pgfqpoint{2.400437in}{1.076137in}}%
\pgfpathlineto{\pgfqpoint{2.391489in}{1.082969in}}%
\pgfpathlineto{\pgfqpoint{2.382872in}{1.085587in}}%
\pgfpathlineto{\pgfqpoint{2.381551in}{1.091127in}}%
\pgfpathlineto{\pgfqpoint{2.389841in}{1.100458in}}%
\pgfpathlineto{\pgfqpoint{2.396141in}{1.102157in}}%
\pgfpathlineto{\pgfqpoint{2.392249in}{1.108544in}}%
\pgfpathlineto{\pgfqpoint{2.398193in}{1.116290in}}%
\pgfpathlineto{\pgfqpoint{2.396061in}{1.120463in}}%
\pgfpathlineto{\pgfqpoint{2.401483in}{1.125255in}}%
\pgfpathlineto{\pgfqpoint{2.411679in}{1.127319in}}%
\pgfpathlineto{\pgfqpoint{2.410869in}{1.119825in}}%
\pgfpathlineto{\pgfqpoint{2.421148in}{1.110912in}}%
\pgfpathlineto{\pgfqpoint{2.420050in}{1.106990in}}%
\pgfpathlineto{\pgfqpoint{2.425369in}{1.103463in}}%
\pgfpathlineto{\pgfqpoint{2.429643in}{1.109432in}}%
\pgfpathclose%
\pgfusepath{fill}%
\end{pgfscope}%
\begin{pgfscope}%
\pgfpathrectangle{\pgfqpoint{0.100000in}{0.100000in}}{\pgfqpoint{3.007045in}{1.925000in}}%
\pgfusepath{clip}%
\pgfsetbuttcap%
\pgfsetmiterjoin%
\definecolor{currentfill}{rgb}{0.288627,0.687120,0.381238}%
\pgfsetfillcolor{currentfill}%
\pgfsetlinewidth{0.000000pt}%
\definecolor{currentstroke}{rgb}{0.000000,0.000000,0.000000}%
\pgfsetstrokecolor{currentstroke}%
\pgfsetstrokeopacity{0.000000}%
\pgfsetdash{}{0pt}%
\pgfpathmoveto{\pgfqpoint{1.844917in}{1.391266in}}%
\pgfpathlineto{\pgfqpoint{1.822052in}{1.390899in}}%
\pgfpathlineto{\pgfqpoint{1.821649in}{1.429987in}}%
\pgfpathlineto{\pgfqpoint{1.837085in}{1.430186in}}%
\pgfpathlineto{\pgfqpoint{1.866993in}{1.430735in}}%
\pgfpathlineto{\pgfqpoint{1.867818in}{1.391658in}}%
\pgfpathlineto{\pgfqpoint{1.844917in}{1.391266in}}%
\pgfpathclose%
\pgfusepath{fill}%
\end{pgfscope}%
\begin{pgfscope}%
\pgfpathrectangle{\pgfqpoint{0.100000in}{0.100000in}}{\pgfqpoint{3.007045in}{1.925000in}}%
\pgfusepath{clip}%
\pgfsetbuttcap%
\pgfsetmiterjoin%
\definecolor{currentfill}{rgb}{0.479815,0.780254,0.483230}%
\pgfsetfillcolor{currentfill}%
\pgfsetlinewidth{0.000000pt}%
\definecolor{currentstroke}{rgb}{0.000000,0.000000,0.000000}%
\pgfsetstrokecolor{currentstroke}%
\pgfsetstrokeopacity{0.000000}%
\pgfsetdash{}{0pt}%
\pgfpathmoveto{\pgfqpoint{2.463450in}{1.158871in}}%
\pgfpathlineto{\pgfqpoint{2.458391in}{1.154238in}}%
\pgfpathlineto{\pgfqpoint{2.450037in}{1.161421in}}%
\pgfpathlineto{\pgfqpoint{2.441696in}{1.172351in}}%
\pgfpathlineto{\pgfqpoint{2.447277in}{1.183280in}}%
\pgfpathlineto{\pgfqpoint{2.446120in}{1.194088in}}%
\pgfpathlineto{\pgfqpoint{2.446995in}{1.200159in}}%
\pgfpathlineto{\pgfqpoint{2.439887in}{1.202342in}}%
\pgfpathlineto{\pgfqpoint{2.439065in}{1.211016in}}%
\pgfpathlineto{\pgfqpoint{2.439738in}{1.214050in}}%
\pgfpathlineto{\pgfqpoint{2.445630in}{1.214029in}}%
\pgfpathlineto{\pgfqpoint{2.445563in}{1.221325in}}%
\pgfpathlineto{\pgfqpoint{2.456449in}{1.222465in}}%
\pgfpathlineto{\pgfqpoint{2.465782in}{1.225372in}}%
\pgfpathlineto{\pgfqpoint{2.467867in}{1.223907in}}%
\pgfpathlineto{\pgfqpoint{2.478476in}{1.225076in}}%
\pgfpathlineto{\pgfqpoint{2.476467in}{1.219730in}}%
\pgfpathlineto{\pgfqpoint{2.485346in}{1.213775in}}%
\pgfpathlineto{\pgfqpoint{2.488733in}{1.208153in}}%
\pgfpathlineto{\pgfqpoint{2.487686in}{1.206566in}}%
\pgfpathlineto{\pgfqpoint{2.494442in}{1.196721in}}%
\pgfpathlineto{\pgfqpoint{2.489753in}{1.191584in}}%
\pgfpathlineto{\pgfqpoint{2.484306in}{1.188484in}}%
\pgfpathlineto{\pgfqpoint{2.478749in}{1.189995in}}%
\pgfpathlineto{\pgfqpoint{2.469733in}{1.178659in}}%
\pgfpathlineto{\pgfqpoint{2.463999in}{1.180859in}}%
\pgfpathlineto{\pgfqpoint{2.460301in}{1.177194in}}%
\pgfpathlineto{\pgfqpoint{2.461626in}{1.162571in}}%
\pgfpathlineto{\pgfqpoint{2.463450in}{1.158871in}}%
\pgfpathclose%
\pgfusepath{fill}%
\end{pgfscope}%
\begin{pgfscope}%
\pgfpathrectangle{\pgfqpoint{0.100000in}{0.100000in}}{\pgfqpoint{3.007045in}{1.925000in}}%
\pgfusepath{clip}%
\pgfsetbuttcap%
\pgfsetmiterjoin%
\definecolor{currentfill}{rgb}{0.612687,0.842261,0.592480}%
\pgfsetfillcolor{currentfill}%
\pgfsetlinewidth{0.000000pt}%
\definecolor{currentstroke}{rgb}{0.000000,0.000000,0.000000}%
\pgfsetstrokecolor{currentstroke}%
\pgfsetstrokeopacity{0.000000}%
\pgfsetdash{}{0pt}%
\pgfpathmoveto{\pgfqpoint{1.898889in}{1.500432in}}%
\pgfpathlineto{\pgfqpoint{1.907419in}{1.500598in}}%
\pgfpathlineto{\pgfqpoint{1.906167in}{1.539095in}}%
\pgfpathlineto{\pgfqpoint{1.905032in}{1.544903in}}%
\pgfpathlineto{\pgfqpoint{1.933694in}{1.545701in}}%
\pgfpathlineto{\pgfqpoint{1.933521in}{1.551363in}}%
\pgfpathlineto{\pgfqpoint{1.962214in}{1.552559in}}%
\pgfpathlineto{\pgfqpoint{1.964639in}{1.500815in}}%
\pgfpathlineto{\pgfqpoint{1.970355in}{1.501053in}}%
\pgfpathlineto{\pgfqpoint{1.970609in}{1.495307in}}%
\pgfpathlineto{\pgfqpoint{1.993541in}{1.496568in}}%
\pgfpathlineto{\pgfqpoint{1.994330in}{1.484956in}}%
\pgfpathlineto{\pgfqpoint{1.994624in}{1.478781in}}%
\pgfpathlineto{\pgfqpoint{1.966485in}{1.477622in}}%
\pgfpathlineto{\pgfqpoint{1.963568in}{1.471696in}}%
\pgfpathlineto{\pgfqpoint{1.947712in}{1.470634in}}%
\pgfpathlineto{\pgfqpoint{1.947232in}{1.465292in}}%
\pgfpathlineto{\pgfqpoint{1.942334in}{1.465087in}}%
\pgfpathlineto{\pgfqpoint{1.949409in}{1.456290in}}%
\pgfpathlineto{\pgfqpoint{1.928151in}{1.455464in}}%
\pgfpathlineto{\pgfqpoint{1.911530in}{1.454907in}}%
\pgfpathlineto{\pgfqpoint{1.910791in}{1.477769in}}%
\pgfpathlineto{\pgfqpoint{1.921271in}{1.478116in}}%
\pgfpathlineto{\pgfqpoint{1.917765in}{1.487298in}}%
\pgfpathlineto{\pgfqpoint{1.911427in}{1.491667in}}%
\pgfpathlineto{\pgfqpoint{1.903859in}{1.493796in}}%
\pgfpathlineto{\pgfqpoint{1.898889in}{1.500432in}}%
\pgfpathclose%
\pgfusepath{fill}%
\end{pgfscope}%
\begin{pgfscope}%
\pgfpathrectangle{\pgfqpoint{0.100000in}{0.100000in}}{\pgfqpoint{3.007045in}{1.925000in}}%
\pgfusepath{clip}%
\pgfsetbuttcap%
\pgfsetmiterjoin%
\definecolor{currentfill}{rgb}{0.652995,0.860085,0.628897}%
\pgfsetfillcolor{currentfill}%
\pgfsetlinewidth{0.000000pt}%
\definecolor{currentstroke}{rgb}{0.000000,0.000000,0.000000}%
\pgfsetstrokecolor{currentstroke}%
\pgfsetstrokeopacity{0.000000}%
\pgfsetdash{}{0pt}%
\pgfpathmoveto{\pgfqpoint{1.632363in}{0.431928in}}%
\pgfpathlineto{\pgfqpoint{1.620569in}{0.414900in}}%
\pgfpathlineto{\pgfqpoint{1.605884in}{0.418493in}}%
\pgfpathlineto{\pgfqpoint{1.606432in}{0.418922in}}%
\pgfpathlineto{\pgfqpoint{1.590934in}{0.451560in}}%
\pgfpathlineto{\pgfqpoint{1.595954in}{0.453216in}}%
\pgfpathlineto{\pgfqpoint{1.585508in}{0.466431in}}%
\pgfpathlineto{\pgfqpoint{1.613261in}{0.488228in}}%
\pgfpathlineto{\pgfqpoint{1.619767in}{0.480550in}}%
\pgfpathlineto{\pgfqpoint{1.611272in}{0.473936in}}%
\pgfpathlineto{\pgfqpoint{1.621451in}{0.460671in}}%
\pgfpathlineto{\pgfqpoint{1.609200in}{0.451305in}}%
\pgfpathlineto{\pgfqpoint{1.613769in}{0.442713in}}%
\pgfpathlineto{\pgfqpoint{1.621822in}{0.440642in}}%
\pgfpathlineto{\pgfqpoint{1.622101in}{0.436398in}}%
\pgfpathlineto{\pgfqpoint{1.632363in}{0.431928in}}%
\pgfpathclose%
\pgfusepath{fill}%
\end{pgfscope}%
\begin{pgfscope}%
\pgfpathrectangle{\pgfqpoint{0.100000in}{0.100000in}}{\pgfqpoint{3.007045in}{1.925000in}}%
\pgfusepath{clip}%
\pgfsetbuttcap%
\pgfsetmiterjoin%
\definecolor{currentfill}{rgb}{0.652995,0.860085,0.628897}%
\pgfsetfillcolor{currentfill}%
\pgfsetlinewidth{0.000000pt}%
\definecolor{currentstroke}{rgb}{0.000000,0.000000,0.000000}%
\pgfsetstrokecolor{currentstroke}%
\pgfsetstrokeopacity{0.000000}%
\pgfsetdash{}{0pt}%
\pgfpathmoveto{\pgfqpoint{2.516660in}{0.795332in}}%
\pgfpathlineto{\pgfqpoint{2.505569in}{0.784566in}}%
\pgfpathlineto{\pgfqpoint{2.500186in}{0.786765in}}%
\pgfpathlineto{\pgfqpoint{2.488225in}{0.782556in}}%
\pgfpathlineto{\pgfqpoint{2.485475in}{0.775665in}}%
\pgfpathlineto{\pgfqpoint{2.482589in}{0.774200in}}%
\pgfpathlineto{\pgfqpoint{2.475814in}{0.775674in}}%
\pgfpathlineto{\pgfqpoint{2.470222in}{0.769671in}}%
\pgfpathlineto{\pgfqpoint{2.464508in}{0.773648in}}%
\pgfpathlineto{\pgfqpoint{2.464577in}{0.779288in}}%
\pgfpathlineto{\pgfqpoint{2.461072in}{0.785247in}}%
\pgfpathlineto{\pgfqpoint{2.454423in}{0.792007in}}%
\pgfpathlineto{\pgfqpoint{2.449695in}{0.794403in}}%
\pgfpathlineto{\pgfqpoint{2.448640in}{0.798533in}}%
\pgfpathlineto{\pgfqpoint{2.441136in}{0.811210in}}%
\pgfpathlineto{\pgfqpoint{2.442122in}{0.815272in}}%
\pgfpathlineto{\pgfqpoint{2.446193in}{0.816402in}}%
\pgfpathlineto{\pgfqpoint{2.449159in}{0.822767in}}%
\pgfpathlineto{\pgfqpoint{2.452424in}{0.825463in}}%
\pgfpathlineto{\pgfqpoint{2.460350in}{0.826839in}}%
\pgfpathlineto{\pgfqpoint{2.468621in}{0.832178in}}%
\pgfpathlineto{\pgfqpoint{2.473536in}{0.831913in}}%
\pgfpathlineto{\pgfqpoint{2.479069in}{0.826299in}}%
\pgfpathlineto{\pgfqpoint{2.480732in}{0.829446in}}%
\pgfpathlineto{\pgfqpoint{2.475134in}{0.834677in}}%
\pgfpathlineto{\pgfqpoint{2.477899in}{0.842105in}}%
\pgfpathlineto{\pgfqpoint{2.474703in}{0.848311in}}%
\pgfpathlineto{\pgfqpoint{2.482424in}{0.851583in}}%
\pgfpathlineto{\pgfqpoint{2.488445in}{0.853368in}}%
\pgfpathlineto{\pgfqpoint{2.492621in}{0.846691in}}%
\pgfpathlineto{\pgfqpoint{2.503173in}{0.844767in}}%
\pgfpathlineto{\pgfqpoint{2.506929in}{0.849528in}}%
\pgfpathlineto{\pgfqpoint{2.516115in}{0.841918in}}%
\pgfpathlineto{\pgfqpoint{2.530038in}{0.837997in}}%
\pgfpathlineto{\pgfqpoint{2.521533in}{0.825837in}}%
\pgfpathlineto{\pgfqpoint{2.503376in}{0.803575in}}%
\pgfpathlineto{\pgfqpoint{2.507037in}{0.798645in}}%
\pgfpathlineto{\pgfqpoint{2.516660in}{0.795332in}}%
\pgfpathclose%
\pgfusepath{fill}%
\end{pgfscope}%
\begin{pgfscope}%
\pgfpathrectangle{\pgfqpoint{0.100000in}{0.100000in}}{\pgfqpoint{3.007045in}{1.925000in}}%
\pgfusepath{clip}%
\pgfsetbuttcap%
\pgfsetmiterjoin%
\definecolor{currentfill}{rgb}{0.699746,0.879769,0.674418}%
\pgfsetfillcolor{currentfill}%
\pgfsetlinewidth{0.000000pt}%
\definecolor{currentstroke}{rgb}{0.000000,0.000000,0.000000}%
\pgfsetstrokecolor{currentstroke}%
\pgfsetstrokeopacity{0.000000}%
\pgfsetdash{}{0pt}%
\pgfpathmoveto{\pgfqpoint{2.719169in}{0.923665in}}%
\pgfpathlineto{\pgfqpoint{2.714846in}{0.926915in}}%
\pgfpathlineto{\pgfqpoint{2.706932in}{0.940219in}}%
\pgfpathlineto{\pgfqpoint{2.705959in}{0.953701in}}%
\pgfpathlineto{\pgfqpoint{2.699683in}{0.959124in}}%
\pgfpathlineto{\pgfqpoint{2.696911in}{0.966874in}}%
\pgfpathlineto{\pgfqpoint{2.693069in}{0.968298in}}%
\pgfpathlineto{\pgfqpoint{2.688736in}{0.994006in}}%
\pgfpathlineto{\pgfqpoint{2.708269in}{1.015278in}}%
\pgfpathlineto{\pgfqpoint{2.710520in}{1.014800in}}%
\pgfpathlineto{\pgfqpoint{2.720462in}{1.011491in}}%
\pgfpathlineto{\pgfqpoint{2.722156in}{1.004752in}}%
\pgfpathlineto{\pgfqpoint{2.725670in}{0.999941in}}%
\pgfpathlineto{\pgfqpoint{2.723536in}{0.997192in}}%
\pgfpathlineto{\pgfqpoint{2.724069in}{0.990555in}}%
\pgfpathlineto{\pgfqpoint{2.737785in}{0.985361in}}%
\pgfpathlineto{\pgfqpoint{2.744917in}{0.981075in}}%
\pgfpathlineto{\pgfqpoint{2.755915in}{0.984607in}}%
\pgfpathlineto{\pgfqpoint{2.755737in}{0.991248in}}%
\pgfpathlineto{\pgfqpoint{2.762672in}{0.991320in}}%
\pgfpathlineto{\pgfqpoint{2.764291in}{0.988344in}}%
\pgfpathlineto{\pgfqpoint{2.761049in}{0.976325in}}%
\pgfpathlineto{\pgfqpoint{2.757169in}{0.971840in}}%
\pgfpathlineto{\pgfqpoint{2.760232in}{0.968674in}}%
\pgfpathlineto{\pgfqpoint{2.769233in}{0.972061in}}%
\pgfpathlineto{\pgfqpoint{2.773670in}{0.970838in}}%
\pgfpathlineto{\pgfqpoint{2.777250in}{0.973729in}}%
\pgfpathlineto{\pgfqpoint{2.783649in}{0.973537in}}%
\pgfpathlineto{\pgfqpoint{2.771520in}{0.950478in}}%
\pgfpathlineto{\pgfqpoint{2.767312in}{0.947670in}}%
\pgfpathlineto{\pgfqpoint{2.760927in}{0.948913in}}%
\pgfpathlineto{\pgfqpoint{2.754301in}{0.947904in}}%
\pgfpathlineto{\pgfqpoint{2.736174in}{0.939738in}}%
\pgfpathlineto{\pgfqpoint{2.719169in}{0.923665in}}%
\pgfpathclose%
\pgfusepath{fill}%
\end{pgfscope}%
\begin{pgfscope}%
\pgfpathrectangle{\pgfqpoint{0.100000in}{0.100000in}}{\pgfqpoint{3.007045in}{1.925000in}}%
\pgfusepath{clip}%
\pgfsetbuttcap%
\pgfsetmiterjoin%
\definecolor{currentfill}{rgb}{0.699746,0.879769,0.674418}%
\pgfsetfillcolor{currentfill}%
\pgfsetlinewidth{0.000000pt}%
\definecolor{currentstroke}{rgb}{0.000000,0.000000,0.000000}%
\pgfsetstrokecolor{currentstroke}%
\pgfsetstrokeopacity{0.000000}%
\pgfsetdash{}{0pt}%
\pgfpathmoveto{\pgfqpoint{1.481225in}{0.618416in}}%
\pgfpathlineto{\pgfqpoint{1.479753in}{0.593339in}}%
\pgfpathlineto{\pgfqpoint{1.431514in}{0.595996in}}%
\pgfpathlineto{\pgfqpoint{1.433065in}{0.620988in}}%
\pgfpathlineto{\pgfqpoint{1.415290in}{0.621980in}}%
\pgfpathlineto{\pgfqpoint{1.417736in}{0.651743in}}%
\pgfpathlineto{\pgfqpoint{1.420501in}{0.688714in}}%
\pgfpathlineto{\pgfqpoint{1.425000in}{0.688410in}}%
\pgfpathlineto{\pgfqpoint{1.454083in}{0.686391in}}%
\pgfpathlineto{\pgfqpoint{1.478210in}{0.684741in}}%
\pgfpathlineto{\pgfqpoint{1.476325in}{0.651539in}}%
\pgfpathlineto{\pgfqpoint{1.483290in}{0.651066in}}%
\pgfpathlineto{\pgfqpoint{1.481225in}{0.618416in}}%
\pgfpathclose%
\pgfusepath{fill}%
\end{pgfscope}%
\begin{pgfscope}%
\pgfpathrectangle{\pgfqpoint{0.100000in}{0.100000in}}{\pgfqpoint{3.007045in}{1.925000in}}%
\pgfusepath{clip}%
\pgfsetbuttcap%
\pgfsetmiterjoin%
\definecolor{currentfill}{rgb}{0.671696,0.867958,0.647105}%
\pgfsetfillcolor{currentfill}%
\pgfsetlinewidth{0.000000pt}%
\definecolor{currentstroke}{rgb}{0.000000,0.000000,0.000000}%
\pgfsetstrokecolor{currentstroke}%
\pgfsetstrokeopacity{0.000000}%
\pgfsetdash{}{0pt}%
\pgfpathmoveto{\pgfqpoint{2.382703in}{0.849302in}}%
\pgfpathlineto{\pgfqpoint{2.380206in}{0.847045in}}%
\pgfpathlineto{\pgfqpoint{2.372820in}{0.852600in}}%
\pgfpathlineto{\pgfqpoint{2.375538in}{0.858536in}}%
\pgfpathlineto{\pgfqpoint{2.371582in}{0.867326in}}%
\pgfpathlineto{\pgfqpoint{2.365118in}{0.869291in}}%
\pgfpathlineto{\pgfqpoint{2.359050in}{0.873951in}}%
\pgfpathlineto{\pgfqpoint{2.358828in}{0.878027in}}%
\pgfpathlineto{\pgfqpoint{2.360142in}{0.881259in}}%
\pgfpathlineto{\pgfqpoint{2.366727in}{0.881594in}}%
\pgfpathlineto{\pgfqpoint{2.371189in}{0.888722in}}%
\pgfpathlineto{\pgfqpoint{2.375790in}{0.888043in}}%
\pgfpathlineto{\pgfqpoint{2.377707in}{0.892532in}}%
\pgfpathlineto{\pgfqpoint{2.385967in}{0.896018in}}%
\pgfpathlineto{\pgfqpoint{2.389462in}{0.896298in}}%
\pgfpathlineto{\pgfqpoint{2.391995in}{0.891918in}}%
\pgfpathlineto{\pgfqpoint{2.402990in}{0.891059in}}%
\pgfpathlineto{\pgfqpoint{2.403995in}{0.889257in}}%
\pgfpathlineto{\pgfqpoint{2.402644in}{0.888304in}}%
\pgfpathlineto{\pgfqpoint{2.399169in}{0.874931in}}%
\pgfpathlineto{\pgfqpoint{2.403763in}{0.872065in}}%
\pgfpathlineto{\pgfqpoint{2.404051in}{0.865135in}}%
\pgfpathlineto{\pgfqpoint{2.406992in}{0.858796in}}%
\pgfpathlineto{\pgfqpoint{2.404649in}{0.856728in}}%
\pgfpathlineto{\pgfqpoint{2.407803in}{0.851715in}}%
\pgfpathlineto{\pgfqpoint{2.406734in}{0.847180in}}%
\pgfpathlineto{\pgfqpoint{2.399267in}{0.840607in}}%
\pgfpathlineto{\pgfqpoint{2.397316in}{0.844754in}}%
\pgfpathlineto{\pgfqpoint{2.386057in}{0.845572in}}%
\pgfpathlineto{\pgfqpoint{2.382703in}{0.849302in}}%
\pgfpathclose%
\pgfusepath{fill}%
\end{pgfscope}%
\begin{pgfscope}%
\pgfpathrectangle{\pgfqpoint{0.100000in}{0.100000in}}{\pgfqpoint{3.007045in}{1.925000in}}%
\pgfusepath{clip}%
\pgfsetbuttcap%
\pgfsetmiterjoin%
\definecolor{currentfill}{rgb}{0.648320,0.858116,0.624344}%
\pgfsetfillcolor{currentfill}%
\pgfsetlinewidth{0.000000pt}%
\definecolor{currentstroke}{rgb}{0.000000,0.000000,0.000000}%
\pgfsetstrokecolor{currentstroke}%
\pgfsetstrokeopacity{0.000000}%
\pgfsetdash{}{0pt}%
\pgfpathmoveto{\pgfqpoint{1.746105in}{1.095854in}}%
\pgfpathlineto{\pgfqpoint{1.723295in}{1.096039in}}%
\pgfpathlineto{\pgfqpoint{1.723731in}{1.116097in}}%
\pgfpathlineto{\pgfqpoint{1.702839in}{1.116383in}}%
\pgfpathlineto{\pgfqpoint{1.700977in}{1.122113in}}%
\pgfpathlineto{\pgfqpoint{1.695279in}{1.122232in}}%
\pgfpathlineto{\pgfqpoint{1.695539in}{1.136550in}}%
\pgfpathlineto{\pgfqpoint{1.701186in}{1.136436in}}%
\pgfpathlineto{\pgfqpoint{1.701412in}{1.144959in}}%
\pgfpathlineto{\pgfqpoint{1.708830in}{1.147175in}}%
\pgfpathlineto{\pgfqpoint{1.712256in}{1.145853in}}%
\pgfpathlineto{\pgfqpoint{1.719484in}{1.147603in}}%
\pgfpathlineto{\pgfqpoint{1.719727in}{1.170628in}}%
\pgfpathlineto{\pgfqpoint{1.732247in}{1.170499in}}%
\pgfpathlineto{\pgfqpoint{1.732315in}{1.176225in}}%
\pgfpathlineto{\pgfqpoint{1.743716in}{1.176139in}}%
\pgfpathlineto{\pgfqpoint{1.743298in}{1.160738in}}%
\pgfpathlineto{\pgfqpoint{1.763181in}{1.160646in}}%
\pgfpathlineto{\pgfqpoint{1.763130in}{1.141541in}}%
\pgfpathlineto{\pgfqpoint{1.764032in}{1.131806in}}%
\pgfpathlineto{\pgfqpoint{1.769446in}{1.131864in}}%
\pgfpathlineto{\pgfqpoint{1.769426in}{1.115817in}}%
\pgfpathlineto{\pgfqpoint{1.746586in}{1.115888in}}%
\pgfpathlineto{\pgfqpoint{1.746105in}{1.095854in}}%
\pgfpathclose%
\pgfusepath{fill}%
\end{pgfscope}%
\begin{pgfscope}%
\pgfpathrectangle{\pgfqpoint{0.100000in}{0.100000in}}{\pgfqpoint{3.007045in}{1.925000in}}%
\pgfusepath{clip}%
\pgfsetbuttcap%
\pgfsetmiterjoin%
\definecolor{currentfill}{rgb}{0.652995,0.860085,0.628897}%
\pgfsetfillcolor{currentfill}%
\pgfsetlinewidth{0.000000pt}%
\definecolor{currentstroke}{rgb}{0.000000,0.000000,0.000000}%
\pgfsetstrokecolor{currentstroke}%
\pgfsetstrokeopacity{0.000000}%
\pgfsetdash{}{0pt}%
\pgfpathmoveto{\pgfqpoint{1.704139in}{0.513794in}}%
\pgfpathlineto{\pgfqpoint{1.696363in}{0.505939in}}%
\pgfpathlineto{\pgfqpoint{1.693878in}{0.498881in}}%
\pgfpathlineto{\pgfqpoint{1.676720in}{0.488443in}}%
\pgfpathlineto{\pgfqpoint{1.681469in}{0.494214in}}%
\pgfpathlineto{\pgfqpoint{1.668257in}{0.510605in}}%
\pgfpathlineto{\pgfqpoint{1.663618in}{0.514403in}}%
\pgfpathlineto{\pgfqpoint{1.648009in}{0.514439in}}%
\pgfpathlineto{\pgfqpoint{1.638148in}{0.525132in}}%
\pgfpathlineto{\pgfqpoint{1.655528in}{0.542512in}}%
\pgfpathlineto{\pgfqpoint{1.665918in}{0.546597in}}%
\pgfpathlineto{\pgfqpoint{1.668934in}{0.549504in}}%
\pgfpathlineto{\pgfqpoint{1.675326in}{0.548748in}}%
\pgfpathlineto{\pgfqpoint{1.683561in}{0.532352in}}%
\pgfpathlineto{\pgfqpoint{1.690540in}{0.526743in}}%
\pgfpathlineto{\pgfqpoint{1.695319in}{0.526005in}}%
\pgfpathlineto{\pgfqpoint{1.699409in}{0.516089in}}%
\pgfpathlineto{\pgfqpoint{1.704139in}{0.513794in}}%
\pgfpathclose%
\pgfusepath{fill}%
\end{pgfscope}%
\begin{pgfscope}%
\pgfpathrectangle{\pgfqpoint{0.100000in}{0.100000in}}{\pgfqpoint{3.007045in}{1.925000in}}%
\pgfusepath{clip}%
\pgfsetbuttcap%
\pgfsetmiterjoin%
\definecolor{currentfill}{rgb}{0.463206,0.772503,0.469573}%
\pgfsetfillcolor{currentfill}%
\pgfsetlinewidth{0.000000pt}%
\definecolor{currentstroke}{rgb}{0.000000,0.000000,0.000000}%
\pgfsetstrokecolor{currentstroke}%
\pgfsetstrokeopacity{0.000000}%
\pgfsetdash{}{0pt}%
\pgfpathmoveto{\pgfqpoint{0.646486in}{1.339327in}}%
\pgfpathlineto{\pgfqpoint{0.594471in}{1.352638in}}%
\pgfpathlineto{\pgfqpoint{0.558741in}{1.362309in}}%
\pgfpathlineto{\pgfqpoint{0.567831in}{1.395284in}}%
\pgfpathlineto{\pgfqpoint{0.569070in}{1.394980in}}%
\pgfpathlineto{\pgfqpoint{0.576936in}{1.423028in}}%
\pgfpathlineto{\pgfqpoint{0.575970in}{1.423296in}}%
\pgfpathlineto{\pgfqpoint{0.584193in}{1.452083in}}%
\pgfpathlineto{\pgfqpoint{0.594261in}{1.488948in}}%
\pgfpathlineto{\pgfqpoint{0.623923in}{1.480760in}}%
\pgfpathlineto{\pgfqpoint{0.647669in}{1.474848in}}%
\pgfpathlineto{\pgfqpoint{0.703346in}{1.460924in}}%
\pgfpathlineto{\pgfqpoint{0.703710in}{1.460774in}}%
\pgfpathlineto{\pgfqpoint{0.687900in}{1.396831in}}%
\pgfpathlineto{\pgfqpoint{0.682240in}{1.373981in}}%
\pgfpathlineto{\pgfqpoint{0.671162in}{1.376694in}}%
\pgfpathlineto{\pgfqpoint{0.646486in}{1.339327in}}%
\pgfpathclose%
\pgfusepath{fill}%
\end{pgfscope}%
\begin{pgfscope}%
\pgfpathrectangle{\pgfqpoint{0.100000in}{0.100000in}}{\pgfqpoint{3.007045in}{1.925000in}}%
\pgfusepath{clip}%
\pgfsetbuttcap%
\pgfsetmiterjoin%
\definecolor{currentfill}{rgb}{0.685721,0.873864,0.660761}%
\pgfsetfillcolor{currentfill}%
\pgfsetlinewidth{0.000000pt}%
\definecolor{currentstroke}{rgb}{0.000000,0.000000,0.000000}%
\pgfsetstrokecolor{currentstroke}%
\pgfsetstrokeopacity{0.000000}%
\pgfsetdash{}{0pt}%
\pgfpathmoveto{\pgfqpoint{1.674656in}{1.382036in}}%
\pgfpathlineto{\pgfqpoint{1.673964in}{1.354981in}}%
\pgfpathlineto{\pgfqpoint{1.656900in}{1.355436in}}%
\pgfpathlineto{\pgfqpoint{1.657059in}{1.361181in}}%
\pgfpathlineto{\pgfqpoint{1.611668in}{1.362462in}}%
\pgfpathlineto{\pgfqpoint{1.612271in}{1.391873in}}%
\pgfpathlineto{\pgfqpoint{1.618866in}{1.388633in}}%
\pgfpathlineto{\pgfqpoint{1.621472in}{1.391045in}}%
\pgfpathlineto{\pgfqpoint{1.622963in}{1.397529in}}%
\pgfpathlineto{\pgfqpoint{1.622340in}{1.404124in}}%
\pgfpathlineto{\pgfqpoint{1.624011in}{1.410336in}}%
\pgfpathlineto{\pgfqpoint{1.656789in}{1.409421in}}%
\pgfpathlineto{\pgfqpoint{1.668294in}{1.409172in}}%
\pgfpathlineto{\pgfqpoint{1.667711in}{1.384730in}}%
\pgfpathlineto{\pgfqpoint{1.674656in}{1.382036in}}%
\pgfpathclose%
\pgfusepath{fill}%
\end{pgfscope}%
\begin{pgfscope}%
\pgfpathrectangle{\pgfqpoint{0.100000in}{0.100000in}}{\pgfqpoint{3.007045in}{1.925000in}}%
\pgfusepath{clip}%
\pgfsetbuttcap%
\pgfsetmiterjoin%
\definecolor{currentfill}{rgb}{0.648320,0.858116,0.624344}%
\pgfsetfillcolor{currentfill}%
\pgfsetlinewidth{0.000000pt}%
\definecolor{currentstroke}{rgb}{0.000000,0.000000,0.000000}%
\pgfsetstrokecolor{currentstroke}%
\pgfsetstrokeopacity{0.000000}%
\pgfsetdash{}{0pt}%
\pgfpathmoveto{\pgfqpoint{2.675874in}{0.992417in}}%
\pgfpathlineto{\pgfqpoint{2.688736in}{0.994006in}}%
\pgfpathlineto{\pgfqpoint{2.693069in}{0.968298in}}%
\pgfpathlineto{\pgfqpoint{2.696911in}{0.966874in}}%
\pgfpathlineto{\pgfqpoint{2.699683in}{0.959124in}}%
\pgfpathlineto{\pgfqpoint{2.705959in}{0.953701in}}%
\pgfpathlineto{\pgfqpoint{2.706932in}{0.940219in}}%
\pgfpathlineto{\pgfqpoint{2.683843in}{0.935913in}}%
\pgfpathlineto{\pgfqpoint{2.678361in}{0.924174in}}%
\pgfpathlineto{\pgfqpoint{2.668672in}{0.934954in}}%
\pgfpathlineto{\pgfqpoint{2.661874in}{0.940850in}}%
\pgfpathlineto{\pgfqpoint{2.652361in}{0.948046in}}%
\pgfpathlineto{\pgfqpoint{2.650566in}{0.955387in}}%
\pgfpathlineto{\pgfqpoint{2.650737in}{0.964955in}}%
\pgfpathlineto{\pgfqpoint{2.653930in}{0.970205in}}%
\pgfpathlineto{\pgfqpoint{2.657195in}{0.967322in}}%
\pgfpathlineto{\pgfqpoint{2.661577in}{0.967658in}}%
\pgfpathlineto{\pgfqpoint{2.667375in}{0.972671in}}%
\pgfpathlineto{\pgfqpoint{2.674050in}{0.976293in}}%
\pgfpathlineto{\pgfqpoint{2.673153in}{0.980638in}}%
\pgfpathlineto{\pgfqpoint{2.675874in}{0.992417in}}%
\pgfpathclose%
\pgfusepath{fill}%
\end{pgfscope}%
\begin{pgfscope}%
\pgfpathrectangle{\pgfqpoint{0.100000in}{0.100000in}}{\pgfqpoint{3.007045in}{1.925000in}}%
\pgfusepath{clip}%
\pgfsetbuttcap%
\pgfsetmiterjoin%
\definecolor{currentfill}{rgb}{0.601615,0.837093,0.583376}%
\pgfsetfillcolor{currentfill}%
\pgfsetlinewidth{0.000000pt}%
\definecolor{currentstroke}{rgb}{0.000000,0.000000,0.000000}%
\pgfsetstrokecolor{currentstroke}%
\pgfsetstrokeopacity{0.000000}%
\pgfsetdash{}{0pt}%
\pgfpathmoveto{\pgfqpoint{2.339044in}{0.640515in}}%
\pgfpathlineto{\pgfqpoint{2.341073in}{0.635732in}}%
\pgfpathlineto{\pgfqpoint{2.282576in}{0.629115in}}%
\pgfpathlineto{\pgfqpoint{2.273936in}{0.628314in}}%
\pgfpathlineto{\pgfqpoint{2.270062in}{0.663726in}}%
\pgfpathlineto{\pgfqpoint{2.272250in}{0.669798in}}%
\pgfpathlineto{\pgfqpoint{2.294624in}{0.672055in}}%
\pgfpathlineto{\pgfqpoint{2.313278in}{0.674089in}}%
\pgfpathlineto{\pgfqpoint{2.312691in}{0.679850in}}%
\pgfpathlineto{\pgfqpoint{2.329091in}{0.680831in}}%
\pgfpathlineto{\pgfqpoint{2.333467in}{0.676210in}}%
\pgfpathlineto{\pgfqpoint{2.334926in}{0.671284in}}%
\pgfpathlineto{\pgfqpoint{2.332928in}{0.653635in}}%
\pgfpathlineto{\pgfqpoint{2.334338in}{0.646016in}}%
\pgfpathlineto{\pgfqpoint{2.339044in}{0.640515in}}%
\pgfpathclose%
\pgfusepath{fill}%
\end{pgfscope}%
\begin{pgfscope}%
\pgfpathrectangle{\pgfqpoint{0.100000in}{0.100000in}}{\pgfqpoint{3.007045in}{1.925000in}}%
\pgfusepath{clip}%
\pgfsetbuttcap%
\pgfsetmiterjoin%
\definecolor{currentfill}{rgb}{0.612687,0.842261,0.592480}%
\pgfsetfillcolor{currentfill}%
\pgfsetlinewidth{0.000000pt}%
\definecolor{currentstroke}{rgb}{0.000000,0.000000,0.000000}%
\pgfsetstrokecolor{currentstroke}%
\pgfsetstrokeopacity{0.000000}%
\pgfsetdash{}{0pt}%
\pgfpathmoveto{\pgfqpoint{0.746029in}{1.809361in}}%
\pgfpathlineto{\pgfqpoint{0.701521in}{1.821189in}}%
\pgfpathlineto{\pgfqpoint{0.713636in}{1.864996in}}%
\pgfpathlineto{\pgfqpoint{0.719784in}{1.864006in}}%
\pgfpathlineto{\pgfqpoint{0.730840in}{1.859613in}}%
\pgfpathlineto{\pgfqpoint{0.735184in}{1.852084in}}%
\pgfpathlineto{\pgfqpoint{0.744909in}{1.856652in}}%
\pgfpathlineto{\pgfqpoint{0.748542in}{1.851631in}}%
\pgfpathlineto{\pgfqpoint{0.749934in}{1.845470in}}%
\pgfpathlineto{\pgfqpoint{0.759267in}{1.846349in}}%
\pgfpathlineto{\pgfqpoint{0.762060in}{1.843258in}}%
\pgfpathlineto{\pgfqpoint{0.768647in}{1.844176in}}%
\pgfpathlineto{\pgfqpoint{0.773351in}{1.839140in}}%
\pgfpathlineto{\pgfqpoint{0.777516in}{1.855475in}}%
\pgfpathlineto{\pgfqpoint{0.780353in}{1.860745in}}%
\pgfpathlineto{\pgfqpoint{0.789000in}{1.894475in}}%
\pgfpathlineto{\pgfqpoint{0.783449in}{1.895856in}}%
\pgfpathlineto{\pgfqpoint{0.787685in}{1.900755in}}%
\pgfpathlineto{\pgfqpoint{0.789229in}{1.906864in}}%
\pgfpathlineto{\pgfqpoint{0.797625in}{1.916024in}}%
\pgfpathlineto{\pgfqpoint{0.814525in}{1.911735in}}%
\pgfpathlineto{\pgfqpoint{0.812070in}{1.901824in}}%
\pgfpathlineto{\pgfqpoint{0.822653in}{1.899281in}}%
\pgfpathlineto{\pgfqpoint{0.817161in}{1.876758in}}%
\pgfpathlineto{\pgfqpoint{0.848850in}{1.869247in}}%
\pgfpathlineto{\pgfqpoint{0.844425in}{1.850601in}}%
\pgfpathlineto{\pgfqpoint{0.840821in}{1.835141in}}%
\pgfpathlineto{\pgfqpoint{0.847161in}{1.818241in}}%
\pgfpathlineto{\pgfqpoint{0.850754in}{1.813504in}}%
\pgfpathlineto{\pgfqpoint{0.850775in}{1.806624in}}%
\pgfpathlineto{\pgfqpoint{0.847085in}{1.804318in}}%
\pgfpathlineto{\pgfqpoint{0.851627in}{1.798562in}}%
\pgfpathlineto{\pgfqpoint{0.845810in}{1.797316in}}%
\pgfpathlineto{\pgfqpoint{0.853094in}{1.789354in}}%
\pgfpathlineto{\pgfqpoint{0.853297in}{1.785851in}}%
\pgfpathlineto{\pgfqpoint{0.861882in}{1.780983in}}%
\pgfpathlineto{\pgfqpoint{0.869191in}{1.761494in}}%
\pgfpathlineto{\pgfqpoint{0.873084in}{1.756371in}}%
\pgfpathlineto{\pgfqpoint{0.844253in}{1.762934in}}%
\pgfpathlineto{\pgfqpoint{0.812473in}{1.770492in}}%
\pgfpathlineto{\pgfqpoint{0.813833in}{1.776182in}}%
\pgfpathlineto{\pgfqpoint{0.808364in}{1.778465in}}%
\pgfpathlineto{\pgfqpoint{0.794601in}{1.781859in}}%
\pgfpathlineto{\pgfqpoint{0.784213in}{1.790565in}}%
\pgfpathlineto{\pgfqpoint{0.786311in}{1.799097in}}%
\pgfpathlineto{\pgfqpoint{0.746029in}{1.809361in}}%
\pgfpathclose%
\pgfusepath{fill}%
\end{pgfscope}%
\begin{pgfscope}%
\pgfpathrectangle{\pgfqpoint{0.100000in}{0.100000in}}{\pgfqpoint{3.007045in}{1.925000in}}%
\pgfusepath{clip}%
\pgfsetbuttcap%
\pgfsetmiterjoin%
\definecolor{currentfill}{rgb}{0.676371,0.869927,0.651657}%
\pgfsetfillcolor{currentfill}%
\pgfsetlinewidth{0.000000pt}%
\definecolor{currentstroke}{rgb}{0.000000,0.000000,0.000000}%
\pgfsetstrokecolor{currentstroke}%
\pgfsetstrokeopacity{0.000000}%
\pgfsetdash{}{0pt}%
\pgfpathmoveto{\pgfqpoint{1.839662in}{0.517608in}}%
\pgfpathlineto{\pgfqpoint{1.824849in}{0.515348in}}%
\pgfpathlineto{\pgfqpoint{1.809818in}{0.508665in}}%
\pgfpathlineto{\pgfqpoint{1.809447in}{0.530410in}}%
\pgfpathlineto{\pgfqpoint{1.804575in}{0.530516in}}%
\pgfpathlineto{\pgfqpoint{1.804317in}{0.545414in}}%
\pgfpathlineto{\pgfqpoint{1.795599in}{0.545307in}}%
\pgfpathlineto{\pgfqpoint{1.787662in}{0.570458in}}%
\pgfpathlineto{\pgfqpoint{1.774291in}{0.576196in}}%
\pgfpathlineto{\pgfqpoint{1.770975in}{0.581765in}}%
\pgfpathlineto{\pgfqpoint{1.765743in}{0.581820in}}%
\pgfpathlineto{\pgfqpoint{1.764790in}{0.589806in}}%
\pgfpathlineto{\pgfqpoint{1.760890in}{0.592713in}}%
\pgfpathlineto{\pgfqpoint{1.775795in}{0.607167in}}%
\pgfpathlineto{\pgfqpoint{1.781271in}{0.614134in}}%
\pgfpathlineto{\pgfqpoint{1.786479in}{0.611427in}}%
\pgfpathlineto{\pgfqpoint{1.793823in}{0.610228in}}%
\pgfpathlineto{\pgfqpoint{1.803223in}{0.606687in}}%
\pgfpathlineto{\pgfqpoint{1.821877in}{0.611209in}}%
\pgfpathlineto{\pgfqpoint{1.834269in}{0.615251in}}%
\pgfpathlineto{\pgfqpoint{1.851931in}{0.616696in}}%
\pgfpathlineto{\pgfqpoint{1.854631in}{0.617364in}}%
\pgfpathlineto{\pgfqpoint{1.854154in}{0.604333in}}%
\pgfpathlineto{\pgfqpoint{1.856414in}{0.600398in}}%
\pgfpathlineto{\pgfqpoint{1.853934in}{0.596710in}}%
\pgfpathlineto{\pgfqpoint{1.850658in}{0.583534in}}%
\pgfpathlineto{\pgfqpoint{1.847738in}{0.581014in}}%
\pgfpathlineto{\pgfqpoint{1.844939in}{0.573686in}}%
\pgfpathlineto{\pgfqpoint{1.847088in}{0.567676in}}%
\pgfpathlineto{\pgfqpoint{1.843280in}{0.560440in}}%
\pgfpathlineto{\pgfqpoint{1.846835in}{0.557616in}}%
\pgfpathlineto{\pgfqpoint{1.846781in}{0.541896in}}%
\pgfpathlineto{\pgfqpoint{1.842229in}{0.537436in}}%
\pgfpathlineto{\pgfqpoint{1.840585in}{0.532547in}}%
\pgfpathlineto{\pgfqpoint{1.834496in}{0.525796in}}%
\pgfpathlineto{\pgfqpoint{1.839662in}{0.517608in}}%
\pgfpathclose%
\pgfusepath{fill}%
\end{pgfscope}%
\begin{pgfscope}%
\pgfpathrectangle{\pgfqpoint{0.100000in}{0.100000in}}{\pgfqpoint{3.007045in}{1.925000in}}%
\pgfusepath{clip}%
\pgfsetbuttcap%
\pgfsetmiterjoin%
\definecolor{currentfill}{rgb}{0.823760,0.931073,0.799200}%
\pgfsetfillcolor{currentfill}%
\pgfsetlinewidth{0.000000pt}%
\definecolor{currentstroke}{rgb}{0.000000,0.000000,0.000000}%
\pgfsetstrokecolor{currentstroke}%
\pgfsetstrokeopacity{0.000000}%
\pgfsetdash{}{0pt}%
\pgfpathmoveto{\pgfqpoint{2.697515in}{1.156406in}}%
\pgfpathlineto{\pgfqpoint{2.692040in}{1.154397in}}%
\pgfpathlineto{\pgfqpoint{2.685472in}{1.142255in}}%
\pgfpathlineto{\pgfqpoint{2.681074in}{1.143862in}}%
\pgfpathlineto{\pgfqpoint{2.679314in}{1.148103in}}%
\pgfpathlineto{\pgfqpoint{2.665522in}{1.153025in}}%
\pgfpathlineto{\pgfqpoint{2.656779in}{1.157595in}}%
\pgfpathlineto{\pgfqpoint{2.650904in}{1.157743in}}%
\pgfpathlineto{\pgfqpoint{2.663143in}{1.179525in}}%
\pgfpathlineto{\pgfqpoint{2.662595in}{1.183012in}}%
\pgfpathlineto{\pgfqpoint{2.668861in}{1.192884in}}%
\pgfpathlineto{\pgfqpoint{2.678549in}{1.188294in}}%
\pgfpathlineto{\pgfqpoint{2.680316in}{1.180573in}}%
\pgfpathlineto{\pgfqpoint{2.682867in}{1.178238in}}%
\pgfpathlineto{\pgfqpoint{2.691552in}{1.182331in}}%
\pgfpathlineto{\pgfqpoint{2.695566in}{1.180728in}}%
\pgfpathlineto{\pgfqpoint{2.697143in}{1.177380in}}%
\pgfpathlineto{\pgfqpoint{2.695805in}{1.170685in}}%
\pgfpathlineto{\pgfqpoint{2.690990in}{1.162746in}}%
\pgfpathlineto{\pgfqpoint{2.697515in}{1.156406in}}%
\pgfpathclose%
\pgfusepath{fill}%
\end{pgfscope}%
\begin{pgfscope}%
\pgfpathrectangle{\pgfqpoint{0.100000in}{0.100000in}}{\pgfqpoint{3.007045in}{1.925000in}}%
\pgfusepath{clip}%
\pgfsetbuttcap%
\pgfsetmiterjoin%
\definecolor{currentfill}{rgb}{0.827451,0.932549,0.803137}%
\pgfsetfillcolor{currentfill}%
\pgfsetlinewidth{0.000000pt}%
\definecolor{currentstroke}{rgb}{0.000000,0.000000,0.000000}%
\pgfsetstrokecolor{currentstroke}%
\pgfsetstrokeopacity{0.000000}%
\pgfsetdash{}{0pt}%
\pgfpathmoveto{\pgfqpoint{1.534548in}{1.032337in}}%
\pgfpathlineto{\pgfqpoint{1.533400in}{1.007220in}}%
\pgfpathlineto{\pgfqpoint{1.509153in}{1.008633in}}%
\pgfpathlineto{\pgfqpoint{1.504587in}{1.008926in}}%
\pgfpathlineto{\pgfqpoint{1.505902in}{1.034236in}}%
\pgfpathlineto{\pgfqpoint{1.505407in}{1.040070in}}%
\pgfpathlineto{\pgfqpoint{1.476919in}{1.041763in}}%
\pgfpathlineto{\pgfqpoint{1.478439in}{1.065027in}}%
\pgfpathlineto{\pgfqpoint{1.478495in}{1.076500in}}%
\pgfpathlineto{\pgfqpoint{1.501163in}{1.075036in}}%
\pgfpathlineto{\pgfqpoint{1.502124in}{1.092214in}}%
\pgfpathlineto{\pgfqpoint{1.535329in}{1.090418in}}%
\pgfpathlineto{\pgfqpoint{1.536034in}{1.090382in}}%
\pgfpathlineto{\pgfqpoint{1.534600in}{1.061649in}}%
\pgfpathlineto{\pgfqpoint{1.534859in}{1.055660in}}%
\pgfpathlineto{\pgfqpoint{1.534548in}{1.032337in}}%
\pgfpathclose%
\pgfusepath{fill}%
\end{pgfscope}%
\begin{pgfscope}%
\pgfpathrectangle{\pgfqpoint{0.100000in}{0.100000in}}{\pgfqpoint{3.007045in}{1.925000in}}%
\pgfusepath{clip}%
\pgfsetbuttcap%
\pgfsetmiterjoin%
\definecolor{currentfill}{rgb}{0.562860,0.819008,0.551511}%
\pgfsetfillcolor{currentfill}%
\pgfsetlinewidth{0.000000pt}%
\definecolor{currentstroke}{rgb}{0.000000,0.000000,0.000000}%
\pgfsetstrokecolor{currentstroke}%
\pgfsetstrokeopacity{0.000000}%
\pgfsetdash{}{0pt}%
\pgfpathmoveto{\pgfqpoint{1.482882in}{0.684491in}}%
\pgfpathlineto{\pgfqpoint{1.478210in}{0.684741in}}%
\pgfpathlineto{\pgfqpoint{1.454083in}{0.686391in}}%
\pgfpathlineto{\pgfqpoint{1.458235in}{0.744453in}}%
\pgfpathlineto{\pgfqpoint{1.465814in}{0.743950in}}%
\pgfpathlineto{\pgfqpoint{1.486605in}{0.742533in}}%
\pgfpathlineto{\pgfqpoint{1.482882in}{0.684491in}}%
\pgfpathclose%
\pgfusepath{fill}%
\end{pgfscope}%
\begin{pgfscope}%
\pgfpathrectangle{\pgfqpoint{0.100000in}{0.100000in}}{\pgfqpoint{3.007045in}{1.925000in}}%
\pgfusepath{clip}%
\pgfsetbuttcap%
\pgfsetmiterjoin%
\definecolor{currentfill}{rgb}{0.827451,0.932549,0.803137}%
\pgfsetfillcolor{currentfill}%
\pgfsetlinewidth{0.000000pt}%
\definecolor{currentstroke}{rgb}{0.000000,0.000000,0.000000}%
\pgfsetstrokecolor{currentstroke}%
\pgfsetstrokeopacity{0.000000}%
\pgfsetdash{}{0pt}%
\pgfpathmoveto{\pgfqpoint{1.476919in}{1.041763in}}%
\pgfpathlineto{\pgfqpoint{1.476575in}{1.036016in}}%
\pgfpathlineto{\pgfqpoint{1.454880in}{1.037434in}}%
\pgfpathlineto{\pgfqpoint{1.453744in}{1.037504in}}%
\pgfpathlineto{\pgfqpoint{1.455281in}{1.060433in}}%
\pgfpathlineto{\pgfqpoint{1.405795in}{1.064151in}}%
\pgfpathlineto{\pgfqpoint{1.408351in}{1.098572in}}%
\pgfpathlineto{\pgfqpoint{1.432949in}{1.096770in}}%
\pgfpathlineto{\pgfqpoint{1.478511in}{1.093767in}}%
\pgfpathlineto{\pgfqpoint{1.502124in}{1.092214in}}%
\pgfpathlineto{\pgfqpoint{1.501163in}{1.075036in}}%
\pgfpathlineto{\pgfqpoint{1.478495in}{1.076500in}}%
\pgfpathlineto{\pgfqpoint{1.478439in}{1.065027in}}%
\pgfpathlineto{\pgfqpoint{1.476919in}{1.041763in}}%
\pgfpathclose%
\pgfusepath{fill}%
\end{pgfscope}%
\begin{pgfscope}%
\pgfpathrectangle{\pgfqpoint{0.100000in}{0.100000in}}{\pgfqpoint{3.007045in}{1.925000in}}%
\pgfusepath{clip}%
\pgfsetbuttcap%
\pgfsetmiterjoin%
\definecolor{currentfill}{rgb}{0.755848,0.903391,0.729043}%
\pgfsetfillcolor{currentfill}%
\pgfsetlinewidth{0.000000pt}%
\definecolor{currentstroke}{rgb}{0.000000,0.000000,0.000000}%
\pgfsetstrokecolor{currentstroke}%
\pgfsetstrokeopacity{0.000000}%
\pgfsetdash{}{0pt}%
\pgfpathmoveto{\pgfqpoint{1.917287in}{1.088655in}}%
\pgfpathlineto{\pgfqpoint{1.945763in}{1.089375in}}%
\pgfpathlineto{\pgfqpoint{1.946082in}{1.080499in}}%
\pgfpathlineto{\pgfqpoint{1.952290in}{1.080650in}}%
\pgfpathlineto{\pgfqpoint{1.953009in}{1.056717in}}%
\pgfpathlineto{\pgfqpoint{1.964484in}{1.057050in}}%
\pgfpathlineto{\pgfqpoint{1.964595in}{1.051309in}}%
\pgfpathlineto{\pgfqpoint{1.972796in}{1.051471in}}%
\pgfpathlineto{\pgfqpoint{1.973010in}{1.044275in}}%
\pgfpathlineto{\pgfqpoint{1.981728in}{1.045692in}}%
\pgfpathlineto{\pgfqpoint{1.993117in}{1.046153in}}%
\pgfpathlineto{\pgfqpoint{1.993598in}{1.031047in}}%
\pgfpathlineto{\pgfqpoint{1.995876in}{1.024384in}}%
\pgfpathlineto{\pgfqpoint{1.996111in}{1.017492in}}%
\pgfpathlineto{\pgfqpoint{1.994242in}{1.009746in}}%
\pgfpathlineto{\pgfqpoint{1.984516in}{1.009519in}}%
\pgfpathlineto{\pgfqpoint{1.981594in}{1.012166in}}%
\pgfpathlineto{\pgfqpoint{1.981501in}{1.015319in}}%
\pgfpathlineto{\pgfqpoint{1.975616in}{1.018755in}}%
\pgfpathlineto{\pgfqpoint{1.972674in}{1.026092in}}%
\pgfpathlineto{\pgfqpoint{1.972430in}{1.032823in}}%
\pgfpathlineto{\pgfqpoint{1.942096in}{1.032256in}}%
\pgfpathlineto{\pgfqpoint{1.941697in}{1.043763in}}%
\pgfpathlineto{\pgfqpoint{1.915878in}{1.043286in}}%
\pgfpathlineto{\pgfqpoint{1.915821in}{1.046178in}}%
\pgfpathlineto{\pgfqpoint{1.907105in}{1.050239in}}%
\pgfpathlineto{\pgfqpoint{1.907064in}{1.059811in}}%
\pgfpathlineto{\pgfqpoint{1.906913in}{1.070484in}}%
\pgfpathlineto{\pgfqpoint{1.918382in}{1.070506in}}%
\pgfpathlineto{\pgfqpoint{1.917287in}{1.088655in}}%
\pgfpathclose%
\pgfusepath{fill}%
\end{pgfscope}%
\begin{pgfscope}%
\pgfpathrectangle{\pgfqpoint{0.100000in}{0.100000in}}{\pgfqpoint{3.007045in}{1.925000in}}%
\pgfusepath{clip}%
\pgfsetbuttcap%
\pgfsetmiterjoin%
\definecolor{currentfill}{rgb}{0.540715,0.808674,0.533303}%
\pgfsetfillcolor{currentfill}%
\pgfsetlinewidth{0.000000pt}%
\definecolor{currentstroke}{rgb}{0.000000,0.000000,0.000000}%
\pgfsetstrokecolor{currentstroke}%
\pgfsetstrokeopacity{0.000000}%
\pgfsetdash{}{0pt}%
\pgfpathmoveto{\pgfqpoint{2.375927in}{0.644931in}}%
\pgfpathlineto{\pgfqpoint{2.378692in}{0.619214in}}%
\pgfpathlineto{\pgfqpoint{2.351014in}{0.617483in}}%
\pgfpathlineto{\pgfqpoint{2.346074in}{0.624253in}}%
\pgfpathlineto{\pgfqpoint{2.345288in}{0.628691in}}%
\pgfpathlineto{\pgfqpoint{2.341073in}{0.635732in}}%
\pgfpathlineto{\pgfqpoint{2.339044in}{0.640515in}}%
\pgfpathlineto{\pgfqpoint{2.345074in}{0.640987in}}%
\pgfpathlineto{\pgfqpoint{2.343997in}{0.653135in}}%
\pgfpathlineto{\pgfqpoint{2.359552in}{0.655091in}}%
\pgfpathlineto{\pgfqpoint{2.365478in}{0.655570in}}%
\pgfpathlineto{\pgfqpoint{2.366514in}{0.643850in}}%
\pgfpathlineto{\pgfqpoint{2.375927in}{0.644931in}}%
\pgfpathclose%
\pgfusepath{fill}%
\end{pgfscope}%
\begin{pgfscope}%
\pgfpathrectangle{\pgfqpoint{0.100000in}{0.100000in}}{\pgfqpoint{3.007045in}{1.925000in}}%
\pgfusepath{clip}%
\pgfsetbuttcap%
\pgfsetmiterjoin%
\definecolor{currentfill}{rgb}{0.871742,0.950265,0.850381}%
\pgfsetfillcolor{currentfill}%
\pgfsetlinewidth{0.000000pt}%
\definecolor{currentstroke}{rgb}{0.000000,0.000000,0.000000}%
\pgfsetstrokecolor{currentstroke}%
\pgfsetstrokeopacity{0.000000}%
\pgfsetdash{}{0pt}%
\pgfpathmoveto{\pgfqpoint{1.704139in}{0.513794in}}%
\pgfpathlineto{\pgfqpoint{1.699409in}{0.516089in}}%
\pgfpathlineto{\pgfqpoint{1.695319in}{0.526005in}}%
\pgfpathlineto{\pgfqpoint{1.690540in}{0.526743in}}%
\pgfpathlineto{\pgfqpoint{1.683561in}{0.532352in}}%
\pgfpathlineto{\pgfqpoint{1.675326in}{0.548748in}}%
\pgfpathlineto{\pgfqpoint{1.668934in}{0.549504in}}%
\pgfpathlineto{\pgfqpoint{1.672597in}{0.554943in}}%
\pgfpathlineto{\pgfqpoint{1.677967in}{0.558409in}}%
\pgfpathlineto{\pgfqpoint{1.690544in}{0.561282in}}%
\pgfpathlineto{\pgfqpoint{1.696897in}{0.563937in}}%
\pgfpathlineto{\pgfqpoint{1.704636in}{0.564588in}}%
\pgfpathlineto{\pgfqpoint{1.709428in}{0.556522in}}%
\pgfpathlineto{\pgfqpoint{1.709639in}{0.553402in}}%
\pgfpathlineto{\pgfqpoint{1.726035in}{0.554362in}}%
\pgfpathlineto{\pgfqpoint{1.724776in}{0.579988in}}%
\pgfpathlineto{\pgfqpoint{1.737969in}{0.571808in}}%
\pgfpathlineto{\pgfqpoint{1.751572in}{0.571734in}}%
\pgfpathlineto{\pgfqpoint{1.753626in}{0.595060in}}%
\pgfpathlineto{\pgfqpoint{1.757986in}{0.598171in}}%
\pgfpathlineto{\pgfqpoint{1.760890in}{0.592713in}}%
\pgfpathlineto{\pgfqpoint{1.764790in}{0.589806in}}%
\pgfpathlineto{\pgfqpoint{1.765743in}{0.581820in}}%
\pgfpathlineto{\pgfqpoint{1.770975in}{0.581765in}}%
\pgfpathlineto{\pgfqpoint{1.774291in}{0.576196in}}%
\pgfpathlineto{\pgfqpoint{1.787662in}{0.570458in}}%
\pgfpathlineto{\pgfqpoint{1.795599in}{0.545307in}}%
\pgfpathlineto{\pgfqpoint{1.804317in}{0.545414in}}%
\pgfpathlineto{\pgfqpoint{1.804575in}{0.530516in}}%
\pgfpathlineto{\pgfqpoint{1.809447in}{0.530410in}}%
\pgfpathlineto{\pgfqpoint{1.809818in}{0.508665in}}%
\pgfpathlineto{\pgfqpoint{1.794978in}{0.501878in}}%
\pgfpathlineto{\pgfqpoint{1.797904in}{0.509417in}}%
\pgfpathlineto{\pgfqpoint{1.787377in}{0.506124in}}%
\pgfpathlineto{\pgfqpoint{1.790301in}{0.517701in}}%
\pgfpathlineto{\pgfqpoint{1.787535in}{0.523539in}}%
\pgfpathlineto{\pgfqpoint{1.783015in}{0.521586in}}%
\pgfpathlineto{\pgfqpoint{1.779708in}{0.515842in}}%
\pgfpathlineto{\pgfqpoint{1.774796in}{0.517671in}}%
\pgfpathlineto{\pgfqpoint{1.771403in}{0.513633in}}%
\pgfpathlineto{\pgfqpoint{1.771383in}{0.507282in}}%
\pgfpathlineto{\pgfqpoint{1.777596in}{0.504206in}}%
\pgfpathlineto{\pgfqpoint{1.778463in}{0.493573in}}%
\pgfpathlineto{\pgfqpoint{1.788184in}{0.493368in}}%
\pgfpathlineto{\pgfqpoint{1.765360in}{0.476226in}}%
\pgfpathlineto{\pgfqpoint{1.769052in}{0.484340in}}%
\pgfpathlineto{\pgfqpoint{1.758777in}{0.502145in}}%
\pgfpathlineto{\pgfqpoint{1.759753in}{0.508170in}}%
\pgfpathlineto{\pgfqpoint{1.757144in}{0.510663in}}%
\pgfpathlineto{\pgfqpoint{1.747709in}{0.509837in}}%
\pgfpathlineto{\pgfqpoint{1.745433in}{0.500436in}}%
\pgfpathlineto{\pgfqpoint{1.740375in}{0.500409in}}%
\pgfpathlineto{\pgfqpoint{1.738650in}{0.493869in}}%
\pgfpathlineto{\pgfqpoint{1.734233in}{0.490553in}}%
\pgfpathlineto{\pgfqpoint{1.728488in}{0.492886in}}%
\pgfpathlineto{\pgfqpoint{1.722932in}{0.488802in}}%
\pgfpathlineto{\pgfqpoint{1.717041in}{0.493709in}}%
\pgfpathlineto{\pgfqpoint{1.714651in}{0.500322in}}%
\pgfpathlineto{\pgfqpoint{1.710733in}{0.503026in}}%
\pgfpathlineto{\pgfqpoint{1.710725in}{0.510845in}}%
\pgfpathlineto{\pgfqpoint{1.704139in}{0.513794in}}%
\pgfpathclose%
\pgfusepath{fill}%
\end{pgfscope}%
\begin{pgfscope}%
\pgfpathrectangle{\pgfqpoint{0.100000in}{0.100000in}}{\pgfqpoint{3.007045in}{1.925000in}}%
\pgfusepath{clip}%
\pgfsetbuttcap%
\pgfsetmiterjoin%
\definecolor{currentfill}{rgb}{0.774548,0.911265,0.747251}%
\pgfsetfillcolor{currentfill}%
\pgfsetlinewidth{0.000000pt}%
\definecolor{currentstroke}{rgb}{0.000000,0.000000,0.000000}%
\pgfsetstrokecolor{currentstroke}%
\pgfsetstrokeopacity{0.000000}%
\pgfsetdash{}{0pt}%
\pgfpathmoveto{\pgfqpoint{1.631785in}{0.798613in}}%
\pgfpathlineto{\pgfqpoint{1.625906in}{0.804863in}}%
\pgfpathlineto{\pgfqpoint{1.614953in}{0.795791in}}%
\pgfpathlineto{\pgfqpoint{1.608887in}{0.798669in}}%
\pgfpathlineto{\pgfqpoint{1.610850in}{0.805245in}}%
\pgfpathlineto{\pgfqpoint{1.605264in}{0.805706in}}%
\pgfpathlineto{\pgfqpoint{1.598818in}{0.813845in}}%
\pgfpathlineto{\pgfqpoint{1.588268in}{0.817057in}}%
\pgfpathlineto{\pgfqpoint{1.586260in}{0.813360in}}%
\pgfpathlineto{\pgfqpoint{1.581430in}{0.811070in}}%
\pgfpathlineto{\pgfqpoint{1.574901in}{0.817578in}}%
\pgfpathlineto{\pgfqpoint{1.575401in}{0.829239in}}%
\pgfpathlineto{\pgfqpoint{1.572553in}{0.829339in}}%
\pgfpathlineto{\pgfqpoint{1.573014in}{0.840820in}}%
\pgfpathlineto{\pgfqpoint{1.564075in}{0.841182in}}%
\pgfpathlineto{\pgfqpoint{1.564310in}{0.846924in}}%
\pgfpathlineto{\pgfqpoint{1.565031in}{0.864162in}}%
\pgfpathlineto{\pgfqpoint{1.576210in}{0.863673in}}%
\pgfpathlineto{\pgfqpoint{1.604926in}{0.862611in}}%
\pgfpathlineto{\pgfqpoint{1.604562in}{0.851135in}}%
\pgfpathlineto{\pgfqpoint{1.627401in}{0.850427in}}%
\pgfpathlineto{\pgfqpoint{1.633138in}{0.850249in}}%
\pgfpathlineto{\pgfqpoint{1.631785in}{0.798613in}}%
\pgfpathclose%
\pgfusepath{fill}%
\end{pgfscope}%
\begin{pgfscope}%
\pgfpathrectangle{\pgfqpoint{0.100000in}{0.100000in}}{\pgfqpoint{3.007045in}{1.925000in}}%
\pgfusepath{clip}%
\pgfsetbuttcap%
\pgfsetmiterjoin%
\definecolor{currentfill}{rgb}{0.535179,0.806090,0.528750}%
\pgfsetfillcolor{currentfill}%
\pgfsetlinewidth{0.000000pt}%
\definecolor{currentstroke}{rgb}{0.000000,0.000000,0.000000}%
\pgfsetstrokecolor{currentstroke}%
\pgfsetstrokeopacity{0.000000}%
\pgfsetdash{}{0pt}%
\pgfpathmoveto{\pgfqpoint{2.815641in}{1.458005in}}%
\pgfpathlineto{\pgfqpoint{2.815310in}{1.453061in}}%
\pgfpathlineto{\pgfqpoint{2.806124in}{1.452725in}}%
\pgfpathlineto{\pgfqpoint{2.795846in}{1.455210in}}%
\pgfpathlineto{\pgfqpoint{2.794361in}{1.458299in}}%
\pgfpathlineto{\pgfqpoint{2.787830in}{1.459571in}}%
\pgfpathlineto{\pgfqpoint{2.787399in}{1.454470in}}%
\pgfpathlineto{\pgfqpoint{2.776040in}{1.453025in}}%
\pgfpathlineto{\pgfqpoint{2.768356in}{1.454967in}}%
\pgfpathlineto{\pgfqpoint{2.763804in}{1.456101in}}%
\pgfpathlineto{\pgfqpoint{2.766813in}{1.466787in}}%
\pgfpathlineto{\pgfqpoint{2.756645in}{1.469554in}}%
\pgfpathlineto{\pgfqpoint{2.750879in}{1.474563in}}%
\pgfpathlineto{\pgfqpoint{2.753121in}{1.482442in}}%
\pgfpathlineto{\pgfqpoint{2.749558in}{1.490026in}}%
\pgfpathlineto{\pgfqpoint{2.749311in}{1.495307in}}%
\pgfpathlineto{\pgfqpoint{2.759273in}{1.493756in}}%
\pgfpathlineto{\pgfqpoint{2.766923in}{1.496257in}}%
\pgfpathlineto{\pgfqpoint{2.774908in}{1.505824in}}%
\pgfpathlineto{\pgfqpoint{2.756305in}{1.557859in}}%
\pgfpathlineto{\pgfqpoint{2.763396in}{1.560589in}}%
\pgfpathlineto{\pgfqpoint{2.793875in}{1.572201in}}%
\pgfpathlineto{\pgfqpoint{2.794553in}{1.564410in}}%
\pgfpathlineto{\pgfqpoint{2.803917in}{1.557332in}}%
\pgfpathlineto{\pgfqpoint{2.807383in}{1.541913in}}%
\pgfpathlineto{\pgfqpoint{2.813736in}{1.511035in}}%
\pgfpathlineto{\pgfqpoint{2.815796in}{1.505313in}}%
\pgfpathlineto{\pgfqpoint{2.814539in}{1.460221in}}%
\pgfpathlineto{\pgfqpoint{2.815641in}{1.458005in}}%
\pgfpathclose%
\pgfusepath{fill}%
\end{pgfscope}%
\begin{pgfscope}%
\pgfpathrectangle{\pgfqpoint{0.100000in}{0.100000in}}{\pgfqpoint{3.007045in}{1.925000in}}%
\pgfusepath{clip}%
\pgfsetbuttcap%
\pgfsetmiterjoin%
\definecolor{currentfill}{rgb}{0.579469,0.826759,0.565167}%
\pgfsetfillcolor{currentfill}%
\pgfsetlinewidth{0.000000pt}%
\definecolor{currentstroke}{rgb}{0.000000,0.000000,0.000000}%
\pgfsetstrokecolor{currentstroke}%
\pgfsetstrokeopacity{0.000000}%
\pgfsetdash{}{0pt}%
\pgfpathmoveto{\pgfqpoint{1.154054in}{1.574061in}}%
\pgfpathlineto{\pgfqpoint{1.115456in}{1.579920in}}%
\pgfpathlineto{\pgfqpoint{1.114575in}{1.580480in}}%
\pgfpathlineto{\pgfqpoint{1.082392in}{1.585392in}}%
\pgfpathlineto{\pgfqpoint{1.084238in}{1.596483in}}%
\pgfpathlineto{\pgfqpoint{1.072000in}{1.598524in}}%
\pgfpathlineto{\pgfqpoint{1.073953in}{1.610065in}}%
\pgfpathlineto{\pgfqpoint{1.079975in}{1.609049in}}%
\pgfpathlineto{\pgfqpoint{1.081871in}{1.620429in}}%
\pgfpathlineto{\pgfqpoint{1.094101in}{1.624217in}}%
\pgfpathlineto{\pgfqpoint{1.099750in}{1.623315in}}%
\pgfpathlineto{\pgfqpoint{1.101489in}{1.634717in}}%
\pgfpathlineto{\pgfqpoint{1.105971in}{1.645829in}}%
\pgfpathlineto{\pgfqpoint{1.109783in}{1.645289in}}%
\pgfpathlineto{\pgfqpoint{1.110938in}{1.650768in}}%
\pgfpathlineto{\pgfqpoint{1.102412in}{1.652202in}}%
\pgfpathlineto{\pgfqpoint{1.101102in}{1.658254in}}%
\pgfpathlineto{\pgfqpoint{1.102024in}{1.663968in}}%
\pgfpathlineto{\pgfqpoint{1.113428in}{1.662137in}}%
\pgfpathlineto{\pgfqpoint{1.119016in}{1.693144in}}%
\pgfpathlineto{\pgfqpoint{1.119627in}{1.696907in}}%
\pgfpathlineto{\pgfqpoint{1.153729in}{1.691500in}}%
\pgfpathlineto{\pgfqpoint{1.189998in}{1.686616in}}%
\pgfpathlineto{\pgfqpoint{1.190856in}{1.681412in}}%
\pgfpathlineto{\pgfqpoint{1.188964in}{1.669306in}}%
\pgfpathlineto{\pgfqpoint{1.192292in}{1.666923in}}%
\pgfpathlineto{\pgfqpoint{1.195568in}{1.655776in}}%
\pgfpathlineto{\pgfqpoint{1.201035in}{1.646271in}}%
\pgfpathlineto{\pgfqpoint{1.197757in}{1.637726in}}%
\pgfpathlineto{\pgfqpoint{1.190334in}{1.638786in}}%
\pgfpathlineto{\pgfqpoint{1.189792in}{1.634979in}}%
\pgfpathlineto{\pgfqpoint{1.184148in}{1.635775in}}%
\pgfpathlineto{\pgfqpoint{1.179523in}{1.630622in}}%
\pgfpathlineto{\pgfqpoint{1.171947in}{1.631734in}}%
\pgfpathlineto{\pgfqpoint{1.169858in}{1.624279in}}%
\pgfpathlineto{\pgfqpoint{1.167306in}{1.606868in}}%
\pgfpathlineto{\pgfqpoint{1.155409in}{1.604740in}}%
\pgfpathlineto{\pgfqpoint{1.140335in}{1.607233in}}%
\pgfpathlineto{\pgfqpoint{1.138319in}{1.605460in}}%
\pgfpathlineto{\pgfqpoint{1.135376in}{1.591485in}}%
\pgfpathlineto{\pgfqpoint{1.163161in}{1.587525in}}%
\pgfpathlineto{\pgfqpoint{1.160708in}{1.586055in}}%
\pgfpathlineto{\pgfqpoint{1.154054in}{1.574061in}}%
\pgfpathclose%
\pgfusepath{fill}%
\end{pgfscope}%
\begin{pgfscope}%
\pgfpathrectangle{\pgfqpoint{0.100000in}{0.100000in}}{\pgfqpoint{3.007045in}{1.925000in}}%
\pgfusepath{clip}%
\pgfsetbuttcap%
\pgfsetmiterjoin%
\definecolor{currentfill}{rgb}{0.507497,0.793172,0.505990}%
\pgfsetfillcolor{currentfill}%
\pgfsetlinewidth{0.000000pt}%
\definecolor{currentstroke}{rgb}{0.000000,0.000000,0.000000}%
\pgfsetstrokecolor{currentstroke}%
\pgfsetstrokeopacity{0.000000}%
\pgfsetdash{}{0pt}%
\pgfpathmoveto{\pgfqpoint{2.427990in}{0.752079in}}%
\pgfpathlineto{\pgfqpoint{2.415340in}{0.771423in}}%
\pgfpathlineto{\pgfqpoint{2.417835in}{0.773607in}}%
\pgfpathlineto{\pgfqpoint{2.411719in}{0.790134in}}%
\pgfpathlineto{\pgfqpoint{2.405437in}{0.788426in}}%
\pgfpathlineto{\pgfqpoint{2.403890in}{0.805755in}}%
\pgfpathlineto{\pgfqpoint{2.417379in}{0.810753in}}%
\pgfpathlineto{\pgfqpoint{2.424800in}{0.803148in}}%
\pgfpathlineto{\pgfqpoint{2.432103in}{0.811727in}}%
\pgfpathlineto{\pgfqpoint{2.441136in}{0.811210in}}%
\pgfpathlineto{\pgfqpoint{2.448640in}{0.798533in}}%
\pgfpathlineto{\pgfqpoint{2.443204in}{0.795028in}}%
\pgfpathlineto{\pgfqpoint{2.441806in}{0.790953in}}%
\pgfpathlineto{\pgfqpoint{2.437607in}{0.788428in}}%
\pgfpathlineto{\pgfqpoint{2.432968in}{0.782406in}}%
\pgfpathlineto{\pgfqpoint{2.433317in}{0.776989in}}%
\pgfpathlineto{\pgfqpoint{2.436633in}{0.771428in}}%
\pgfpathlineto{\pgfqpoint{2.441164in}{0.768340in}}%
\pgfpathlineto{\pgfqpoint{2.441826in}{0.762195in}}%
\pgfpathlineto{\pgfqpoint{2.427990in}{0.752079in}}%
\pgfpathclose%
\pgfusepath{fill}%
\end{pgfscope}%
\begin{pgfscope}%
\pgfpathrectangle{\pgfqpoint{0.100000in}{0.100000in}}{\pgfqpoint{3.007045in}{1.925000in}}%
\pgfusepath{clip}%
\pgfsetbuttcap%
\pgfsetmiterjoin%
\definecolor{currentfill}{rgb}{0.585006,0.829343,0.569719}%
\pgfsetfillcolor{currentfill}%
\pgfsetlinewidth{0.000000pt}%
\definecolor{currentstroke}{rgb}{0.000000,0.000000,0.000000}%
\pgfsetstrokecolor{currentstroke}%
\pgfsetstrokeopacity{0.000000}%
\pgfsetdash{}{0pt}%
\pgfpathmoveto{\pgfqpoint{1.263629in}{1.757459in}}%
\pgfpathlineto{\pgfqpoint{1.265560in}{1.751233in}}%
\pgfpathlineto{\pgfqpoint{1.269873in}{1.750687in}}%
\pgfpathlineto{\pgfqpoint{1.267101in}{1.728242in}}%
\pgfpathlineto{\pgfqpoint{1.263981in}{1.716912in}}%
\pgfpathlineto{\pgfqpoint{1.269690in}{1.716224in}}%
\pgfpathlineto{\pgfqpoint{1.268284in}{1.704741in}}%
\pgfpathlineto{\pgfqpoint{1.272031in}{1.704281in}}%
\pgfpathlineto{\pgfqpoint{1.269340in}{1.682158in}}%
\pgfpathlineto{\pgfqpoint{1.240737in}{1.685718in}}%
\pgfpathlineto{\pgfqpoint{1.240870in}{1.686663in}}%
\pgfpathlineto{\pgfqpoint{1.187986in}{1.693328in}}%
\pgfpathlineto{\pgfqpoint{1.186086in}{1.695738in}}%
\pgfpathlineto{\pgfqpoint{1.189559in}{1.703789in}}%
\pgfpathlineto{\pgfqpoint{1.188167in}{1.711464in}}%
\pgfpathlineto{\pgfqpoint{1.188345in}{1.720102in}}%
\pgfpathlineto{\pgfqpoint{1.192875in}{1.732905in}}%
\pgfpathlineto{\pgfqpoint{1.195585in}{1.736909in}}%
\pgfpathlineto{\pgfqpoint{1.201819in}{1.739713in}}%
\pgfpathlineto{\pgfqpoint{1.206785in}{1.744306in}}%
\pgfpathlineto{\pgfqpoint{1.214896in}{1.741346in}}%
\pgfpathlineto{\pgfqpoint{1.217097in}{1.745725in}}%
\pgfpathlineto{\pgfqpoint{1.224257in}{1.742711in}}%
\pgfpathlineto{\pgfqpoint{1.240243in}{1.741597in}}%
\pgfpathlineto{\pgfqpoint{1.241819in}{1.745330in}}%
\pgfpathlineto{\pgfqpoint{1.246943in}{1.744290in}}%
\pgfpathlineto{\pgfqpoint{1.255859in}{1.748834in}}%
\pgfpathlineto{\pgfqpoint{1.260272in}{1.755961in}}%
\pgfpathlineto{\pgfqpoint{1.263629in}{1.757459in}}%
\pgfpathclose%
\pgfusepath{fill}%
\end{pgfscope}%
\begin{pgfscope}%
\pgfpathrectangle{\pgfqpoint{0.100000in}{0.100000in}}{\pgfqpoint{3.007045in}{1.925000in}}%
\pgfusepath{clip}%
\pgfsetbuttcap%
\pgfsetmiterjoin%
\definecolor{currentfill}{rgb}{0.612687,0.842261,0.592480}%
\pgfsetfillcolor{currentfill}%
\pgfsetlinewidth{0.000000pt}%
\definecolor{currentstroke}{rgb}{0.000000,0.000000,0.000000}%
\pgfsetstrokecolor{currentstroke}%
\pgfsetstrokeopacity{0.000000}%
\pgfsetdash{}{0pt}%
\pgfpathmoveto{\pgfqpoint{2.644431in}{1.224390in}}%
\pgfpathlineto{\pgfqpoint{2.614371in}{1.241096in}}%
\pgfpathlineto{\pgfqpoint{2.607688in}{1.243264in}}%
\pgfpathlineto{\pgfqpoint{2.608755in}{1.249542in}}%
\pgfpathlineto{\pgfqpoint{2.613340in}{1.256035in}}%
\pgfpathlineto{\pgfqpoint{2.609634in}{1.257414in}}%
\pgfpathlineto{\pgfqpoint{2.612460in}{1.264832in}}%
\pgfpathlineto{\pgfqpoint{2.614370in}{1.276431in}}%
\pgfpathlineto{\pgfqpoint{2.613885in}{1.282571in}}%
\pgfpathlineto{\pgfqpoint{2.616493in}{1.288353in}}%
\pgfpathlineto{\pgfqpoint{2.627614in}{1.286206in}}%
\pgfpathlineto{\pgfqpoint{2.631148in}{1.284175in}}%
\pgfpathlineto{\pgfqpoint{2.636603in}{1.298784in}}%
\pgfpathlineto{\pgfqpoint{2.638827in}{1.299653in}}%
\pgfpathlineto{\pgfqpoint{2.652060in}{1.280358in}}%
\pgfpathlineto{\pgfqpoint{2.655502in}{1.265959in}}%
\pgfpathlineto{\pgfqpoint{2.651264in}{1.258406in}}%
\pgfpathlineto{\pgfqpoint{2.649600in}{1.238259in}}%
\pgfpathlineto{\pgfqpoint{2.646870in}{1.236936in}}%
\pgfpathlineto{\pgfqpoint{2.644431in}{1.224390in}}%
\pgfpathclose%
\pgfusepath{fill}%
\end{pgfscope}%
\begin{pgfscope}%
\pgfpathrectangle{\pgfqpoint{0.100000in}{0.100000in}}{\pgfqpoint{3.007045in}{1.925000in}}%
\pgfusepath{clip}%
\pgfsetbuttcap%
\pgfsetmiterjoin%
\definecolor{currentfill}{rgb}{0.886505,0.956171,0.866128}%
\pgfsetfillcolor{currentfill}%
\pgfsetlinewidth{0.000000pt}%
\definecolor{currentstroke}{rgb}{0.000000,0.000000,0.000000}%
\pgfsetstrokecolor{currentstroke}%
\pgfsetstrokeopacity{0.000000}%
\pgfsetdash{}{0pt}%
\pgfpathmoveto{\pgfqpoint{0.677092in}{0.410736in}}%
\pgfpathlineto{\pgfqpoint{0.679449in}{0.413309in}}%
\pgfpathlineto{\pgfqpoint{0.683920in}{0.412319in}}%
\pgfpathlineto{\pgfqpoint{0.681471in}{0.410739in}}%
\pgfpathlineto{\pgfqpoint{0.677092in}{0.410736in}}%
\pgfpathclose%
\pgfusepath{fill}%
\end{pgfscope}%
\begin{pgfscope}%
\pgfpathrectangle{\pgfqpoint{0.100000in}{0.100000in}}{\pgfqpoint{3.007045in}{1.925000in}}%
\pgfusepath{clip}%
\pgfsetbuttcap%
\pgfsetmiterjoin%
\definecolor{currentfill}{rgb}{0.886505,0.956171,0.866128}%
\pgfsetfillcolor{currentfill}%
\pgfsetlinewidth{0.000000pt}%
\definecolor{currentstroke}{rgb}{0.000000,0.000000,0.000000}%
\pgfsetstrokecolor{currentstroke}%
\pgfsetstrokeopacity{0.000000}%
\pgfsetdash{}{0pt}%
\pgfpathmoveto{\pgfqpoint{0.784950in}{0.396276in}}%
\pgfpathlineto{\pgfqpoint{0.784096in}{0.394320in}}%
\pgfpathlineto{\pgfqpoint{0.778986in}{0.387487in}}%
\pgfpathlineto{\pgfqpoint{0.777700in}{0.388433in}}%
\pgfpathlineto{\pgfqpoint{0.772579in}{0.381674in}}%
\pgfpathlineto{\pgfqpoint{0.771354in}{0.382565in}}%
\pgfpathlineto{\pgfqpoint{0.766473in}{0.376050in}}%
\pgfpathlineto{\pgfqpoint{0.765233in}{0.376990in}}%
\pgfpathlineto{\pgfqpoint{0.763938in}{0.375291in}}%
\pgfpathlineto{\pgfqpoint{0.762258in}{0.376547in}}%
\pgfpathlineto{\pgfqpoint{0.759707in}{0.373154in}}%
\pgfpathlineto{\pgfqpoint{0.758045in}{0.374400in}}%
\pgfpathlineto{\pgfqpoint{0.756787in}{0.372685in}}%
\pgfpathlineto{\pgfqpoint{0.753894in}{0.374816in}}%
\pgfpathlineto{\pgfqpoint{0.751397in}{0.371451in}}%
\pgfpathlineto{\pgfqpoint{0.748076in}{0.373959in}}%
\pgfpathlineto{\pgfqpoint{0.745512in}{0.370591in}}%
\pgfpathlineto{\pgfqpoint{0.746018in}{0.370169in}}%
\pgfpathlineto{\pgfqpoint{0.741011in}{0.363689in}}%
\pgfpathlineto{\pgfqpoint{0.740274in}{0.361603in}}%
\pgfpathlineto{\pgfqpoint{0.743549in}{0.359090in}}%
\pgfpathlineto{\pgfqpoint{0.740865in}{0.355528in}}%
\pgfpathlineto{\pgfqpoint{0.738983in}{0.355722in}}%
\pgfpathlineto{\pgfqpoint{0.737813in}{0.354884in}}%
\pgfpathlineto{\pgfqpoint{0.734357in}{0.355834in}}%
\pgfpathlineto{\pgfqpoint{0.731400in}{0.357708in}}%
\pgfpathlineto{\pgfqpoint{0.730107in}{0.357059in}}%
\pgfpathlineto{\pgfqpoint{0.728487in}{0.358327in}}%
\pgfpathlineto{\pgfqpoint{0.726759in}{0.358371in}}%
\pgfpathlineto{\pgfqpoint{0.722429in}{0.352828in}}%
\pgfpathlineto{\pgfqpoint{0.719341in}{0.353168in}}%
\pgfpathlineto{\pgfqpoint{0.717812in}{0.354377in}}%
\pgfpathlineto{\pgfqpoint{0.715201in}{0.352581in}}%
\pgfpathlineto{\pgfqpoint{0.712047in}{0.351741in}}%
\pgfpathlineto{\pgfqpoint{0.711910in}{0.351797in}}%
\pgfpathlineto{\pgfqpoint{0.712359in}{0.352736in}}%
\pgfpathlineto{\pgfqpoint{0.708693in}{0.353375in}}%
\pgfpathlineto{\pgfqpoint{0.707132in}{0.354637in}}%
\pgfpathlineto{\pgfqpoint{0.704495in}{0.354085in}}%
\pgfpathlineto{\pgfqpoint{0.703476in}{0.352730in}}%
\pgfpathlineto{\pgfqpoint{0.701818in}{0.353160in}}%
\pgfpathlineto{\pgfqpoint{0.698993in}{0.352386in}}%
\pgfpathlineto{\pgfqpoint{0.697922in}{0.351049in}}%
\pgfpathlineto{\pgfqpoint{0.701239in}{0.348319in}}%
\pgfpathlineto{\pgfqpoint{0.697491in}{0.346230in}}%
\pgfpathlineto{\pgfqpoint{0.696501in}{0.348097in}}%
\pgfpathlineto{\pgfqpoint{0.693475in}{0.346785in}}%
\pgfpathlineto{\pgfqpoint{0.692417in}{0.347973in}}%
\pgfpathlineto{\pgfqpoint{0.683563in}{0.349485in}}%
\pgfpathlineto{\pgfqpoint{0.682454in}{0.347098in}}%
\pgfpathlineto{\pgfqpoint{0.680426in}{0.348551in}}%
\pgfpathlineto{\pgfqpoint{0.680359in}{0.350534in}}%
\pgfpathlineto{\pgfqpoint{0.678074in}{0.351002in}}%
\pgfpathlineto{\pgfqpoint{0.677768in}{0.349141in}}%
\pgfpathlineto{\pgfqpoint{0.676340in}{0.347791in}}%
\pgfpathlineto{\pgfqpoint{0.672925in}{0.348867in}}%
\pgfpathlineto{\pgfqpoint{0.670574in}{0.351024in}}%
\pgfpathlineto{\pgfqpoint{0.667546in}{0.349762in}}%
\pgfpathlineto{\pgfqpoint{0.669799in}{0.346596in}}%
\pgfpathlineto{\pgfqpoint{0.664976in}{0.345492in}}%
\pgfpathlineto{\pgfqpoint{0.662946in}{0.343281in}}%
\pgfpathlineto{\pgfqpoint{0.662231in}{0.345505in}}%
\pgfpathlineto{\pgfqpoint{0.659923in}{0.345201in}}%
\pgfpathlineto{\pgfqpoint{0.658186in}{0.346446in}}%
\pgfpathlineto{\pgfqpoint{0.656689in}{0.346479in}}%
\pgfpathlineto{\pgfqpoint{0.652000in}{0.348708in}}%
\pgfpathlineto{\pgfqpoint{0.649972in}{0.347731in}}%
\pgfpathlineto{\pgfqpoint{0.648426in}{0.347772in}}%
\pgfpathlineto{\pgfqpoint{0.646196in}{0.346313in}}%
\pgfpathlineto{\pgfqpoint{0.650028in}{0.350774in}}%
\pgfpathlineto{\pgfqpoint{0.647002in}{0.353387in}}%
\pgfpathlineto{\pgfqpoint{0.649075in}{0.354518in}}%
\pgfpathlineto{\pgfqpoint{0.654252in}{0.360591in}}%
\pgfpathlineto{\pgfqpoint{0.653338in}{0.361411in}}%
\pgfpathlineto{\pgfqpoint{0.655995in}{0.364392in}}%
\pgfpathlineto{\pgfqpoint{0.659049in}{0.361833in}}%
\pgfpathlineto{\pgfqpoint{0.660386in}{0.363372in}}%
\pgfpathlineto{\pgfqpoint{0.663515in}{0.360688in}}%
\pgfpathlineto{\pgfqpoint{0.664908in}{0.362318in}}%
\pgfpathlineto{\pgfqpoint{0.667310in}{0.360198in}}%
\pgfpathlineto{\pgfqpoint{0.668605in}{0.361721in}}%
\pgfpathlineto{\pgfqpoint{0.670704in}{0.361448in}}%
\pgfpathlineto{\pgfqpoint{0.672992in}{0.359286in}}%
\pgfpathlineto{\pgfqpoint{0.675516in}{0.363060in}}%
\pgfpathlineto{\pgfqpoint{0.680318in}{0.364924in}}%
\pgfpathlineto{\pgfqpoint{0.685284in}{0.365748in}}%
\pgfpathlineto{\pgfqpoint{0.687937in}{0.364755in}}%
\pgfpathlineto{\pgfqpoint{0.690707in}{0.366122in}}%
\pgfpathlineto{\pgfqpoint{0.694569in}{0.365827in}}%
\pgfpathlineto{\pgfqpoint{0.695772in}{0.367709in}}%
\pgfpathlineto{\pgfqpoint{0.704688in}{0.375601in}}%
\pgfpathlineto{\pgfqpoint{0.706614in}{0.375694in}}%
\pgfpathlineto{\pgfqpoint{0.706494in}{0.377302in}}%
\pgfpathlineto{\pgfqpoint{0.709397in}{0.380342in}}%
\pgfpathlineto{\pgfqpoint{0.713956in}{0.381377in}}%
\pgfpathlineto{\pgfqpoint{0.723393in}{0.373987in}}%
\pgfpathlineto{\pgfqpoint{0.727658in}{0.379320in}}%
\pgfpathlineto{\pgfqpoint{0.720329in}{0.385044in}}%
\pgfpathlineto{\pgfqpoint{0.707689in}{0.387813in}}%
\pgfpathlineto{\pgfqpoint{0.705303in}{0.388793in}}%
\pgfpathlineto{\pgfqpoint{0.704791in}{0.390712in}}%
\pgfpathlineto{\pgfqpoint{0.705523in}{0.393336in}}%
\pgfpathlineto{\pgfqpoint{0.704478in}{0.394497in}}%
\pgfpathlineto{\pgfqpoint{0.707497in}{0.397391in}}%
\pgfpathlineto{\pgfqpoint{0.706468in}{0.397884in}}%
\pgfpathlineto{\pgfqpoint{0.700476in}{0.396199in}}%
\pgfpathlineto{\pgfqpoint{0.698526in}{0.391552in}}%
\pgfpathlineto{\pgfqpoint{0.697132in}{0.389984in}}%
\pgfpathlineto{\pgfqpoint{0.695255in}{0.390285in}}%
\pgfpathlineto{\pgfqpoint{0.694284in}{0.392577in}}%
\pgfpathlineto{\pgfqpoint{0.695108in}{0.393300in}}%
\pgfpathlineto{\pgfqpoint{0.695145in}{0.396106in}}%
\pgfpathlineto{\pgfqpoint{0.696329in}{0.402560in}}%
\pgfpathlineto{\pgfqpoint{0.695632in}{0.403858in}}%
\pgfpathlineto{\pgfqpoint{0.697277in}{0.406002in}}%
\pgfpathlineto{\pgfqpoint{0.695522in}{0.406874in}}%
\pgfpathlineto{\pgfqpoint{0.694538in}{0.405951in}}%
\pgfpathlineto{\pgfqpoint{0.691532in}{0.405166in}}%
\pgfpathlineto{\pgfqpoint{0.691789in}{0.407114in}}%
\pgfpathlineto{\pgfqpoint{0.690598in}{0.411229in}}%
\pgfpathlineto{\pgfqpoint{0.692288in}{0.413752in}}%
\pgfpathlineto{\pgfqpoint{0.690110in}{0.414243in}}%
\pgfpathlineto{\pgfqpoint{0.685710in}{0.414422in}}%
\pgfpathlineto{\pgfqpoint{0.681584in}{0.415593in}}%
\pgfpathlineto{\pgfqpoint{0.684613in}{0.418474in}}%
\pgfpathlineto{\pgfqpoint{0.686127in}{0.417159in}}%
\pgfpathlineto{\pgfqpoint{0.687384in}{0.418587in}}%
\pgfpathlineto{\pgfqpoint{0.688938in}{0.417253in}}%
\pgfpathlineto{\pgfqpoint{0.691710in}{0.420379in}}%
\pgfpathlineto{\pgfqpoint{0.693623in}{0.421487in}}%
\pgfpathlineto{\pgfqpoint{0.695145in}{0.420096in}}%
\pgfpathlineto{\pgfqpoint{0.697072in}{0.421245in}}%
\pgfpathlineto{\pgfqpoint{0.699846in}{0.424401in}}%
\pgfpathlineto{\pgfqpoint{0.701434in}{0.423011in}}%
\pgfpathlineto{\pgfqpoint{0.704186in}{0.426137in}}%
\pgfpathlineto{\pgfqpoint{0.706317in}{0.424353in}}%
\pgfpathlineto{\pgfqpoint{0.709062in}{0.427491in}}%
\pgfpathlineto{\pgfqpoint{0.710582in}{0.426191in}}%
\pgfpathlineto{\pgfqpoint{0.713230in}{0.429162in}}%
\pgfpathlineto{\pgfqpoint{0.715184in}{0.430242in}}%
\pgfpathlineto{\pgfqpoint{0.716770in}{0.428913in}}%
\pgfpathlineto{\pgfqpoint{0.719488in}{0.432032in}}%
\pgfpathlineto{\pgfqpoint{0.721076in}{0.430661in}}%
\pgfpathlineto{\pgfqpoint{0.722451in}{0.432254in}}%
\pgfpathlineto{\pgfqpoint{0.724234in}{0.430706in}}%
\pgfpathlineto{\pgfqpoint{0.726974in}{0.433888in}}%
\pgfpathlineto{\pgfqpoint{0.728779in}{0.435116in}}%
\pgfpathlineto{\pgfqpoint{0.730325in}{0.433863in}}%
\pgfpathlineto{\pgfqpoint{0.731767in}{0.435552in}}%
\pgfpathlineto{\pgfqpoint{0.733009in}{0.434561in}}%
\pgfpathlineto{\pgfqpoint{0.734378in}{0.436156in}}%
\pgfpathlineto{\pgfqpoint{0.743447in}{0.428714in}}%
\pgfpathlineto{\pgfqpoint{0.742145in}{0.427074in}}%
\pgfpathlineto{\pgfqpoint{0.743807in}{0.425686in}}%
\pgfpathlineto{\pgfqpoint{0.745138in}{0.427348in}}%
\pgfpathlineto{\pgfqpoint{0.746839in}{0.426097in}}%
\pgfpathlineto{\pgfqpoint{0.748296in}{0.427861in}}%
\pgfpathlineto{\pgfqpoint{0.751507in}{0.425233in}}%
\pgfpathlineto{\pgfqpoint{0.750064in}{0.423440in}}%
\pgfpathlineto{\pgfqpoint{0.761232in}{0.414538in}}%
\pgfpathlineto{\pgfqpoint{0.775612in}{0.403305in}}%
\pgfpathlineto{\pgfqpoint{0.784950in}{0.396276in}}%
\pgfpathclose%
\pgfusepath{fill}%
\end{pgfscope}%
\begin{pgfscope}%
\pgfpathrectangle{\pgfqpoint{0.100000in}{0.100000in}}{\pgfqpoint{3.007045in}{1.925000in}}%
\pgfusepath{clip}%
\pgfsetbuttcap%
\pgfsetmiterjoin%
\definecolor{currentfill}{rgb}{0.886505,0.956171,0.866128}%
\pgfsetfillcolor{currentfill}%
\pgfsetlinewidth{0.000000pt}%
\definecolor{currentstroke}{rgb}{0.000000,0.000000,0.000000}%
\pgfsetstrokecolor{currentstroke}%
\pgfsetstrokeopacity{0.000000}%
\pgfsetdash{}{0pt}%
\pgfpathmoveto{\pgfqpoint{0.677618in}{0.416243in}}%
\pgfpathlineto{\pgfqpoint{0.678710in}{0.418051in}}%
\pgfpathlineto{\pgfqpoint{0.681431in}{0.415610in}}%
\pgfpathlineto{\pgfqpoint{0.677618in}{0.416243in}}%
\pgfpathclose%
\pgfusepath{fill}%
\end{pgfscope}%
\begin{pgfscope}%
\pgfpathrectangle{\pgfqpoint{0.100000in}{0.100000in}}{\pgfqpoint{3.007045in}{1.925000in}}%
\pgfusepath{clip}%
\pgfsetbuttcap%
\pgfsetmiterjoin%
\definecolor{currentfill}{rgb}{0.288627,0.687120,0.381238}%
\pgfsetfillcolor{currentfill}%
\pgfsetlinewidth{0.000000pt}%
\definecolor{currentstroke}{rgb}{0.000000,0.000000,0.000000}%
\pgfsetstrokecolor{currentstroke}%
\pgfsetstrokeopacity{0.000000}%
\pgfsetdash{}{0pt}%
\pgfpathmoveto{\pgfqpoint{0.957302in}{1.649942in}}%
\pgfpathlineto{\pgfqpoint{0.955356in}{1.640352in}}%
\pgfpathlineto{\pgfqpoint{0.959812in}{1.636263in}}%
\pgfpathlineto{\pgfqpoint{0.963238in}{1.636470in}}%
\pgfpathlineto{\pgfqpoint{0.964235in}{1.629990in}}%
\pgfpathlineto{\pgfqpoint{0.960968in}{1.612983in}}%
\pgfpathlineto{\pgfqpoint{0.966618in}{1.611886in}}%
\pgfpathlineto{\pgfqpoint{0.965491in}{1.606172in}}%
\pgfpathlineto{\pgfqpoint{0.972074in}{1.604894in}}%
\pgfpathlineto{\pgfqpoint{0.971016in}{1.593385in}}%
\pgfpathlineto{\pgfqpoint{0.975154in}{1.589701in}}%
\pgfpathlineto{\pgfqpoint{0.988238in}{1.587058in}}%
\pgfpathlineto{\pgfqpoint{0.987365in}{1.582354in}}%
\pgfpathlineto{\pgfqpoint{1.001653in}{1.579735in}}%
\pgfpathlineto{\pgfqpoint{0.998004in}{1.571906in}}%
\pgfpathlineto{\pgfqpoint{0.989150in}{1.572014in}}%
\pgfpathlineto{\pgfqpoint{0.983095in}{1.569716in}}%
\pgfpathlineto{\pgfqpoint{0.981565in}{1.573774in}}%
\pgfpathlineto{\pgfqpoint{0.970525in}{1.573027in}}%
\pgfpathlineto{\pgfqpoint{0.962156in}{1.577626in}}%
\pgfpathlineto{\pgfqpoint{0.958428in}{1.575728in}}%
\pgfpathlineto{\pgfqpoint{0.956111in}{1.570703in}}%
\pgfpathlineto{\pgfqpoint{0.952519in}{1.573584in}}%
\pgfpathlineto{\pgfqpoint{0.938309in}{1.576703in}}%
\pgfpathlineto{\pgfqpoint{0.935411in}{1.569991in}}%
\pgfpathlineto{\pgfqpoint{0.933731in}{1.569068in}}%
\pgfpathlineto{\pgfqpoint{0.927577in}{1.576545in}}%
\pgfpathlineto{\pgfqpoint{0.928393in}{1.581446in}}%
\pgfpathlineto{\pgfqpoint{0.926655in}{1.595511in}}%
\pgfpathlineto{\pgfqpoint{0.923418in}{1.601071in}}%
\pgfpathlineto{\pgfqpoint{0.916159in}{1.603777in}}%
\pgfpathlineto{\pgfqpoint{0.912228in}{1.607641in}}%
\pgfpathlineto{\pgfqpoint{0.914945in}{1.620213in}}%
\pgfpathlineto{\pgfqpoint{0.910956in}{1.624666in}}%
\pgfpathlineto{\pgfqpoint{0.905525in}{1.640271in}}%
\pgfpathlineto{\pgfqpoint{0.906555in}{1.644348in}}%
\pgfpathlineto{\pgfqpoint{0.903714in}{1.653787in}}%
\pgfpathlineto{\pgfqpoint{0.906088in}{1.658804in}}%
\pgfpathlineto{\pgfqpoint{0.901414in}{1.666132in}}%
\pgfpathlineto{\pgfqpoint{0.908010in}{1.668930in}}%
\pgfpathlineto{\pgfqpoint{0.914135in}{1.673457in}}%
\pgfpathlineto{\pgfqpoint{0.921890in}{1.674121in}}%
\pgfpathlineto{\pgfqpoint{0.923792in}{1.677916in}}%
\pgfpathlineto{\pgfqpoint{0.926926in}{1.665686in}}%
\pgfpathlineto{\pgfqpoint{0.936043in}{1.670714in}}%
\pgfpathlineto{\pgfqpoint{0.941856in}{1.670351in}}%
\pgfpathlineto{\pgfqpoint{0.944841in}{1.665499in}}%
\pgfpathlineto{\pgfqpoint{0.955593in}{1.658524in}}%
\pgfpathlineto{\pgfqpoint{0.957302in}{1.649942in}}%
\pgfpathclose%
\pgfusepath{fill}%
\end{pgfscope}%
\begin{pgfscope}%
\pgfpathrectangle{\pgfqpoint{0.100000in}{0.100000in}}{\pgfqpoint{3.007045in}{1.925000in}}%
\pgfusepath{clip}%
\pgfsetbuttcap%
\pgfsetmiterjoin%
\definecolor{currentfill}{rgb}{0.479815,0.780254,0.483230}%
\pgfsetfillcolor{currentfill}%
\pgfsetlinewidth{0.000000pt}%
\definecolor{currentstroke}{rgb}{0.000000,0.000000,0.000000}%
\pgfsetstrokecolor{currentstroke}%
\pgfsetstrokeopacity{0.000000}%
\pgfsetdash{}{0pt}%
\pgfpathmoveto{\pgfqpoint{0.741011in}{1.724091in}}%
\pgfpathlineto{\pgfqpoint{0.771689in}{1.716185in}}%
\pgfpathlineto{\pgfqpoint{0.774980in}{1.703762in}}%
\pgfpathlineto{\pgfqpoint{0.778737in}{1.702487in}}%
\pgfpathlineto{\pgfqpoint{0.784656in}{1.695189in}}%
\pgfpathlineto{\pgfqpoint{0.786067in}{1.686768in}}%
\pgfpathlineto{\pgfqpoint{0.780091in}{1.680597in}}%
\pgfpathlineto{\pgfqpoint{0.770539in}{1.666797in}}%
\pgfpathlineto{\pgfqpoint{0.766636in}{1.659171in}}%
\pgfpathlineto{\pgfqpoint{0.763180in}{1.655715in}}%
\pgfpathlineto{\pgfqpoint{0.741293in}{1.661247in}}%
\pgfpathlineto{\pgfqpoint{0.742633in}{1.666720in}}%
\pgfpathlineto{\pgfqpoint{0.733120in}{1.669073in}}%
\pgfpathlineto{\pgfqpoint{0.732751in}{1.676168in}}%
\pgfpathlineto{\pgfqpoint{0.728377in}{1.682025in}}%
\pgfpathlineto{\pgfqpoint{0.727519in}{1.694386in}}%
\pgfpathlineto{\pgfqpoint{0.729208in}{1.700974in}}%
\pgfpathlineto{\pgfqpoint{0.727679in}{1.706314in}}%
\pgfpathlineto{\pgfqpoint{0.732306in}{1.716868in}}%
\pgfpathlineto{\pgfqpoint{0.722204in}{1.719436in}}%
\pgfpathlineto{\pgfqpoint{0.724313in}{1.728496in}}%
\pgfpathlineto{\pgfqpoint{0.741011in}{1.724091in}}%
\pgfpathclose%
\pgfusepath{fill}%
\end{pgfscope}%
\begin{pgfscope}%
\pgfpathrectangle{\pgfqpoint{0.100000in}{0.100000in}}{\pgfqpoint{3.007045in}{1.925000in}}%
\pgfusepath{clip}%
\pgfsetbuttcap%
\pgfsetmiterjoin%
\definecolor{currentfill}{rgb}{0.451765,0.767090,0.461207}%
\pgfsetfillcolor{currentfill}%
\pgfsetlinewidth{0.000000pt}%
\definecolor{currentstroke}{rgb}{0.000000,0.000000,0.000000}%
\pgfsetstrokecolor{currentstroke}%
\pgfsetstrokeopacity{0.000000}%
\pgfsetdash{}{0pt}%
\pgfpathmoveto{\pgfqpoint{1.439278in}{1.125169in}}%
\pgfpathlineto{\pgfqpoint{1.410528in}{1.127149in}}%
\pgfpathlineto{\pgfqpoint{1.414739in}{1.184340in}}%
\pgfpathlineto{\pgfqpoint{1.446935in}{1.181960in}}%
\pgfpathlineto{\pgfqpoint{1.448195in}{1.181868in}}%
\pgfpathlineto{\pgfqpoint{1.446054in}{1.153418in}}%
\pgfpathlineto{\pgfqpoint{1.441646in}{1.153706in}}%
\pgfpathlineto{\pgfqpoint{1.439278in}{1.125169in}}%
\pgfpathclose%
\pgfusepath{fill}%
\end{pgfscope}%
\begin{pgfscope}%
\pgfpathrectangle{\pgfqpoint{0.100000in}{0.100000in}}{\pgfqpoint{3.007045in}{1.925000in}}%
\pgfusepath{clip}%
\pgfsetbuttcap%
\pgfsetmiterjoin%
\definecolor{currentfill}{rgb}{0.893887,0.959123,0.874002}%
\pgfsetfillcolor{currentfill}%
\pgfsetlinewidth{0.000000pt}%
\definecolor{currentstroke}{rgb}{0.000000,0.000000,0.000000}%
\pgfsetstrokecolor{currentstroke}%
\pgfsetstrokeopacity{0.000000}%
\pgfsetdash{}{0pt}%
\pgfpathmoveto{\pgfqpoint{1.712393in}{0.726751in}}%
\pgfpathlineto{\pgfqpoint{1.712066in}{0.694873in}}%
\pgfpathlineto{\pgfqpoint{1.694814in}{0.693189in}}%
\pgfpathlineto{\pgfqpoint{1.665733in}{0.676821in}}%
\pgfpathlineto{\pgfqpoint{1.663204in}{0.675422in}}%
\pgfpathlineto{\pgfqpoint{1.655346in}{0.689754in}}%
\pgfpathlineto{\pgfqpoint{1.655743in}{0.708579in}}%
\pgfpathlineto{\pgfqpoint{1.658448in}{0.708465in}}%
\pgfpathlineto{\pgfqpoint{1.659429in}{0.737338in}}%
\pgfpathlineto{\pgfqpoint{1.639106in}{0.738203in}}%
\pgfpathlineto{\pgfqpoint{1.640734in}{0.767480in}}%
\pgfpathlineto{\pgfqpoint{1.634992in}{0.767696in}}%
\pgfpathlineto{\pgfqpoint{1.636022in}{0.799536in}}%
\pgfpathlineto{\pgfqpoint{1.638052in}{0.793421in}}%
\pgfpathlineto{\pgfqpoint{1.641958in}{0.793032in}}%
\pgfpathlineto{\pgfqpoint{1.649125in}{0.798383in}}%
\pgfpathlineto{\pgfqpoint{1.652538in}{0.797937in}}%
\pgfpathlineto{\pgfqpoint{1.651130in}{0.791612in}}%
\pgfpathlineto{\pgfqpoint{1.653822in}{0.786416in}}%
\pgfpathlineto{\pgfqpoint{1.656967in}{0.786599in}}%
\pgfpathlineto{\pgfqpoint{1.665670in}{0.801052in}}%
\pgfpathlineto{\pgfqpoint{1.664958in}{0.765812in}}%
\pgfpathlineto{\pgfqpoint{1.670614in}{0.764978in}}%
\pgfpathlineto{\pgfqpoint{1.695830in}{0.764016in}}%
\pgfpathlineto{\pgfqpoint{1.700709in}{0.760907in}}%
\pgfpathlineto{\pgfqpoint{1.700123in}{0.727126in}}%
\pgfpathlineto{\pgfqpoint{1.712393in}{0.726751in}}%
\pgfpathclose%
\pgfusepath{fill}%
\end{pgfscope}%
\begin{pgfscope}%
\pgfpathrectangle{\pgfqpoint{0.100000in}{0.100000in}}{\pgfqpoint{3.007045in}{1.925000in}}%
\pgfusepath{clip}%
\pgfsetbuttcap%
\pgfsetmiterjoin%
\definecolor{currentfill}{rgb}{0.623760,0.847428,0.601584}%
\pgfsetfillcolor{currentfill}%
\pgfsetlinewidth{0.000000pt}%
\definecolor{currentstroke}{rgb}{0.000000,0.000000,0.000000}%
\pgfsetstrokecolor{currentstroke}%
\pgfsetstrokeopacity{0.000000}%
\pgfsetdash{}{0pt}%
\pgfpathmoveto{\pgfqpoint{1.665670in}{0.801052in}}%
\pgfpathlineto{\pgfqpoint{1.656967in}{0.786599in}}%
\pgfpathlineto{\pgfqpoint{1.653822in}{0.786416in}}%
\pgfpathlineto{\pgfqpoint{1.651130in}{0.791612in}}%
\pgfpathlineto{\pgfqpoint{1.652538in}{0.797937in}}%
\pgfpathlineto{\pgfqpoint{1.649125in}{0.798383in}}%
\pgfpathlineto{\pgfqpoint{1.641958in}{0.793032in}}%
\pgfpathlineto{\pgfqpoint{1.638052in}{0.793421in}}%
\pgfpathlineto{\pgfqpoint{1.636022in}{0.799536in}}%
\pgfpathlineto{\pgfqpoint{1.631785in}{0.798613in}}%
\pgfpathlineto{\pgfqpoint{1.633138in}{0.850249in}}%
\pgfpathlineto{\pgfqpoint{1.627401in}{0.850427in}}%
\pgfpathlineto{\pgfqpoint{1.627761in}{0.861901in}}%
\pgfpathlineto{\pgfqpoint{1.667600in}{0.860848in}}%
\pgfpathlineto{\pgfqpoint{1.667244in}{0.843587in}}%
\pgfpathlineto{\pgfqpoint{1.672960in}{0.843473in}}%
\pgfpathlineto{\pgfqpoint{1.672844in}{0.837714in}}%
\pgfpathlineto{\pgfqpoint{1.689906in}{0.837361in}}%
\pgfpathlineto{\pgfqpoint{1.689815in}{0.831695in}}%
\pgfpathlineto{\pgfqpoint{1.695616in}{0.831530in}}%
\pgfpathlineto{\pgfqpoint{1.695332in}{0.814267in}}%
\pgfpathlineto{\pgfqpoint{1.683354in}{0.804224in}}%
\pgfpathlineto{\pgfqpoint{1.678257in}{0.793022in}}%
\pgfpathlineto{\pgfqpoint{1.668937in}{0.795613in}}%
\pgfpathlineto{\pgfqpoint{1.665670in}{0.801052in}}%
\pgfpathclose%
\pgfusepath{fill}%
\end{pgfscope}%
\begin{pgfscope}%
\pgfpathrectangle{\pgfqpoint{0.100000in}{0.100000in}}{\pgfqpoint{3.007045in}{1.925000in}}%
\pgfusepath{clip}%
\pgfsetbuttcap%
\pgfsetmiterjoin%
\definecolor{currentfill}{rgb}{0.363922,0.724029,0.418147}%
\pgfsetfillcolor{currentfill}%
\pgfsetlinewidth{0.000000pt}%
\definecolor{currentstroke}{rgb}{0.000000,0.000000,0.000000}%
\pgfsetstrokecolor{currentstroke}%
\pgfsetstrokeopacity{0.000000}%
\pgfsetdash{}{0pt}%
\pgfpathmoveto{\pgfqpoint{1.146247in}{0.808107in}}%
\pgfpathlineto{\pgfqpoint{1.147002in}{0.802056in}}%
\pgfpathlineto{\pgfqpoint{1.144512in}{0.779815in}}%
\pgfpathlineto{\pgfqpoint{1.112390in}{0.768586in}}%
\pgfpathlineto{\pgfqpoint{1.112021in}{0.765725in}}%
\pgfpathlineto{\pgfqpoint{1.089302in}{0.768831in}}%
\pgfpathlineto{\pgfqpoint{1.087735in}{0.757358in}}%
\pgfpathlineto{\pgfqpoint{1.070654in}{0.759703in}}%
\pgfpathlineto{\pgfqpoint{1.063758in}{0.762298in}}%
\pgfpathlineto{\pgfqpoint{1.065981in}{0.771962in}}%
\pgfpathlineto{\pgfqpoint{1.063832in}{0.777370in}}%
\pgfpathlineto{\pgfqpoint{1.064806in}{0.782977in}}%
\pgfpathlineto{\pgfqpoint{1.061362in}{0.785689in}}%
\pgfpathlineto{\pgfqpoint{1.062578in}{0.792196in}}%
\pgfpathlineto{\pgfqpoint{1.061462in}{0.800953in}}%
\pgfpathlineto{\pgfqpoint{1.054598in}{0.801868in}}%
\pgfpathlineto{\pgfqpoint{1.057206in}{0.819980in}}%
\pgfpathlineto{\pgfqpoint{1.060540in}{0.819258in}}%
\pgfpathlineto{\pgfqpoint{1.115027in}{0.811851in}}%
\pgfpathlineto{\pgfqpoint{1.146247in}{0.808107in}}%
\pgfpathclose%
\pgfusepath{fill}%
\end{pgfscope}%
\begin{pgfscope}%
\pgfpathrectangle{\pgfqpoint{0.100000in}{0.100000in}}{\pgfqpoint{3.007045in}{1.925000in}}%
\pgfusepath{clip}%
\pgfsetbuttcap%
\pgfsetmiterjoin%
\definecolor{currentfill}{rgb}{0.765198,0.907328,0.738147}%
\pgfsetfillcolor{currentfill}%
\pgfsetlinewidth{0.000000pt}%
\definecolor{currentstroke}{rgb}{0.000000,0.000000,0.000000}%
\pgfsetstrokecolor{currentstroke}%
\pgfsetstrokeopacity{0.000000}%
\pgfsetdash{}{0pt}%
\pgfpathmoveto{\pgfqpoint{1.718799in}{1.643280in}}%
\pgfpathlineto{\pgfqpoint{1.718828in}{1.637478in}}%
\pgfpathlineto{\pgfqpoint{1.713868in}{1.637553in}}%
\pgfpathlineto{\pgfqpoint{1.713584in}{1.614400in}}%
\pgfpathlineto{\pgfqpoint{1.714002in}{1.597084in}}%
\pgfpathlineto{\pgfqpoint{1.699675in}{1.597325in}}%
\pgfpathlineto{\pgfqpoint{1.700191in}{1.591568in}}%
\pgfpathlineto{\pgfqpoint{1.669507in}{1.592201in}}%
\pgfpathlineto{\pgfqpoint{1.668146in}{1.596346in}}%
\pgfpathlineto{\pgfqpoint{1.667675in}{1.615294in}}%
\pgfpathlineto{\pgfqpoint{1.668280in}{1.638340in}}%
\pgfpathlineto{\pgfqpoint{1.649938in}{1.638890in}}%
\pgfpathlineto{\pgfqpoint{1.650635in}{1.662065in}}%
\pgfpathlineto{\pgfqpoint{1.650048in}{1.679481in}}%
\pgfpathlineto{\pgfqpoint{1.689494in}{1.678334in}}%
\pgfpathlineto{\pgfqpoint{1.689129in}{1.695752in}}%
\pgfpathlineto{\pgfqpoint{1.724288in}{1.695202in}}%
\pgfpathlineto{\pgfqpoint{1.724037in}{1.672065in}}%
\pgfpathlineto{\pgfqpoint{1.718261in}{1.672097in}}%
\pgfpathlineto{\pgfqpoint{1.719048in}{1.660538in}}%
\pgfpathlineto{\pgfqpoint{1.718799in}{1.643280in}}%
\pgfpathclose%
\pgfusepath{fill}%
\end{pgfscope}%
\begin{pgfscope}%
\pgfpathrectangle{\pgfqpoint{0.100000in}{0.100000in}}{\pgfqpoint{3.007045in}{1.925000in}}%
\pgfusepath{clip}%
\pgfsetbuttcap%
\pgfsetmiterjoin%
\definecolor{currentfill}{rgb}{0.103883,0.516494,0.246751}%
\pgfsetfillcolor{currentfill}%
\pgfsetlinewidth{0.000000pt}%
\definecolor{currentstroke}{rgb}{0.000000,0.000000,0.000000}%
\pgfsetstrokecolor{currentstroke}%
\pgfsetstrokeopacity{0.000000}%
\pgfsetdash{}{0pt}%
\pgfpathmoveto{\pgfqpoint{2.298344in}{1.505689in}}%
\pgfpathlineto{\pgfqpoint{2.264263in}{1.501550in}}%
\pgfpathlineto{\pgfqpoint{2.252766in}{1.500298in}}%
\pgfpathlineto{\pgfqpoint{2.250233in}{1.523286in}}%
\pgfpathlineto{\pgfqpoint{2.247824in}{1.546163in}}%
\pgfpathlineto{\pgfqpoint{2.245496in}{1.563087in}}%
\pgfpathlineto{\pgfqpoint{2.251176in}{1.563619in}}%
\pgfpathlineto{\pgfqpoint{2.250552in}{1.569368in}}%
\pgfpathlineto{\pgfqpoint{2.273175in}{1.571785in}}%
\pgfpathlineto{\pgfqpoint{2.290081in}{1.574045in}}%
\pgfpathlineto{\pgfqpoint{2.291082in}{1.562710in}}%
\pgfpathlineto{\pgfqpoint{2.298344in}{1.505689in}}%
\pgfpathclose%
\pgfusepath{fill}%
\end{pgfscope}%
\begin{pgfscope}%
\pgfpathrectangle{\pgfqpoint{0.100000in}{0.100000in}}{\pgfqpoint{3.007045in}{1.925000in}}%
\pgfusepath{clip}%
\pgfsetbuttcap%
\pgfsetmiterjoin%
\definecolor{currentfill}{rgb}{0.000000,0.286844,0.114248}%
\pgfsetfillcolor{currentfill}%
\pgfsetlinewidth{0.000000pt}%
\definecolor{currentstroke}{rgb}{0.000000,0.000000,0.000000}%
\pgfsetstrokecolor{currentstroke}%
\pgfsetstrokeopacity{0.000000}%
\pgfsetdash{}{0pt}%
\pgfpathmoveto{\pgfqpoint{2.574578in}{0.398555in}}%
\pgfpathlineto{\pgfqpoint{2.577000in}{0.380965in}}%
\pgfpathlineto{\pgfqpoint{2.547508in}{0.376728in}}%
\pgfpathlineto{\pgfqpoint{2.547051in}{0.383388in}}%
\pgfpathlineto{\pgfqpoint{2.542240in}{0.387941in}}%
\pgfpathlineto{\pgfqpoint{2.539505in}{0.384621in}}%
\pgfpathlineto{\pgfqpoint{2.542207in}{0.377186in}}%
\pgfpathlineto{\pgfqpoint{2.536927in}{0.376168in}}%
\pgfpathlineto{\pgfqpoint{2.534833in}{0.376193in}}%
\pgfpathlineto{\pgfqpoint{2.521960in}{0.392832in}}%
\pgfpathlineto{\pgfqpoint{2.512567in}{0.407613in}}%
\pgfpathlineto{\pgfqpoint{2.503920in}{0.416388in}}%
\pgfpathlineto{\pgfqpoint{2.505714in}{0.422690in}}%
\pgfpathlineto{\pgfqpoint{2.510105in}{0.430828in}}%
\pgfpathlineto{\pgfqpoint{2.539586in}{0.435076in}}%
\pgfpathlineto{\pgfqpoint{2.542441in}{0.414576in}}%
\pgfpathlineto{\pgfqpoint{2.571480in}{0.418988in}}%
\pgfpathlineto{\pgfqpoint{2.574578in}{0.398555in}}%
\pgfpathclose%
\pgfusepath{fill}%
\end{pgfscope}%
\begin{pgfscope}%
\pgfpathrectangle{\pgfqpoint{0.100000in}{0.100000in}}{\pgfqpoint{3.007045in}{1.925000in}}%
\pgfusepath{clip}%
\pgfsetbuttcap%
\pgfsetmiterjoin%
\definecolor{currentfill}{rgb}{0.357647,0.720953,0.415071}%
\pgfsetfillcolor{currentfill}%
\pgfsetlinewidth{0.000000pt}%
\definecolor{currentstroke}{rgb}{0.000000,0.000000,0.000000}%
\pgfsetstrokecolor{currentstroke}%
\pgfsetstrokeopacity{0.000000}%
\pgfsetdash{}{0pt}%
\pgfpathmoveto{\pgfqpoint{2.555437in}{1.355961in}}%
\pgfpathlineto{\pgfqpoint{2.554563in}{1.361326in}}%
\pgfpathlineto{\pgfqpoint{2.559925in}{1.364184in}}%
\pgfpathlineto{\pgfqpoint{2.558934in}{1.375533in}}%
\pgfpathlineto{\pgfqpoint{2.581784in}{1.380075in}}%
\pgfpathlineto{\pgfqpoint{2.595492in}{1.381828in}}%
\pgfpathlineto{\pgfqpoint{2.604647in}{1.373975in}}%
\pgfpathlineto{\pgfqpoint{2.607627in}{1.374496in}}%
\pgfpathlineto{\pgfqpoint{2.608964in}{1.367546in}}%
\pgfpathlineto{\pgfqpoint{2.605715in}{1.356835in}}%
\pgfpathlineto{\pgfqpoint{2.573880in}{1.354042in}}%
\pgfpathlineto{\pgfqpoint{2.561808in}{1.358520in}}%
\pgfpathlineto{\pgfqpoint{2.555437in}{1.355961in}}%
\pgfpathclose%
\pgfusepath{fill}%
\end{pgfscope}%
\begin{pgfscope}%
\pgfpathrectangle{\pgfqpoint{0.100000in}{0.100000in}}{\pgfqpoint{3.007045in}{1.925000in}}%
\pgfusepath{clip}%
\pgfsetbuttcap%
\pgfsetmiterjoin%
\definecolor{currentfill}{rgb}{0.596078,0.834510,0.578824}%
\pgfsetfillcolor{currentfill}%
\pgfsetlinewidth{0.000000pt}%
\definecolor{currentstroke}{rgb}{0.000000,0.000000,0.000000}%
\pgfsetstrokecolor{currentstroke}%
\pgfsetstrokeopacity{0.000000}%
\pgfsetdash{}{0pt}%
\pgfpathmoveto{\pgfqpoint{2.169528in}{1.143936in}}%
\pgfpathlineto{\pgfqpoint{2.155147in}{1.142366in}}%
\pgfpathlineto{\pgfqpoint{2.152006in}{1.147645in}}%
\pgfpathlineto{\pgfqpoint{2.151777in}{1.152326in}}%
\pgfpathlineto{\pgfqpoint{2.148381in}{1.158667in}}%
\pgfpathlineto{\pgfqpoint{2.150340in}{1.161666in}}%
\pgfpathlineto{\pgfqpoint{2.148720in}{1.168849in}}%
\pgfpathlineto{\pgfqpoint{2.152293in}{1.171574in}}%
\pgfpathlineto{\pgfqpoint{2.149319in}{1.206705in}}%
\pgfpathlineto{\pgfqpoint{2.147976in}{1.224103in}}%
\pgfpathlineto{\pgfqpoint{2.154351in}{1.223276in}}%
\pgfpathlineto{\pgfqpoint{2.154631in}{1.211728in}}%
\pgfpathlineto{\pgfqpoint{2.171191in}{1.213179in}}%
\pgfpathlineto{\pgfqpoint{2.169632in}{1.230355in}}%
\pgfpathlineto{\pgfqpoint{2.189499in}{1.232103in}}%
\pgfpathlineto{\pgfqpoint{2.189712in}{1.229760in}}%
\pgfpathlineto{\pgfqpoint{2.193431in}{1.193828in}}%
\pgfpathlineto{\pgfqpoint{2.195596in}{1.188844in}}%
\pgfpathlineto{\pgfqpoint{2.194423in}{1.183326in}}%
\pgfpathlineto{\pgfqpoint{2.181588in}{1.182377in}}%
\pgfpathlineto{\pgfqpoint{2.182227in}{1.173751in}}%
\pgfpathlineto{\pgfqpoint{2.176581in}{1.173291in}}%
\pgfpathlineto{\pgfqpoint{2.177539in}{1.161835in}}%
\pgfpathlineto{\pgfqpoint{2.168048in}{1.161256in}}%
\pgfpathlineto{\pgfqpoint{2.169528in}{1.143936in}}%
\pgfpathclose%
\pgfusepath{fill}%
\end{pgfscope}%
\begin{pgfscope}%
\pgfpathrectangle{\pgfqpoint{0.100000in}{0.100000in}}{\pgfqpoint{3.007045in}{1.925000in}}%
\pgfusepath{clip}%
\pgfsetbuttcap%
\pgfsetmiterjoin%
\definecolor{currentfill}{rgb}{0.551788,0.813841,0.542407}%
\pgfsetfillcolor{currentfill}%
\pgfsetlinewidth{0.000000pt}%
\definecolor{currentstroke}{rgb}{0.000000,0.000000,0.000000}%
\pgfsetstrokecolor{currentstroke}%
\pgfsetstrokeopacity{0.000000}%
\pgfsetdash{}{0pt}%
\pgfpathmoveto{\pgfqpoint{1.333387in}{1.392952in}}%
\pgfpathlineto{\pgfqpoint{1.326529in}{1.324213in}}%
\pgfpathlineto{\pgfqpoint{1.297072in}{1.327304in}}%
\pgfpathlineto{\pgfqpoint{1.297671in}{1.333074in}}%
\pgfpathlineto{\pgfqpoint{1.267098in}{1.336682in}}%
\pgfpathlineto{\pgfqpoint{1.269866in}{1.358232in}}%
\pgfpathlineto{\pgfqpoint{1.272590in}{1.387487in}}%
\pgfpathlineto{\pgfqpoint{1.273764in}{1.398937in}}%
\pgfpathlineto{\pgfqpoint{1.292810in}{1.397021in}}%
\pgfpathlineto{\pgfqpoint{1.333387in}{1.392952in}}%
\pgfpathclose%
\pgfusepath{fill}%
\end{pgfscope}%
\begin{pgfscope}%
\pgfpathrectangle{\pgfqpoint{0.100000in}{0.100000in}}{\pgfqpoint{3.007045in}{1.925000in}}%
\pgfusepath{clip}%
\pgfsetbuttcap%
\pgfsetmiterjoin%
\definecolor{currentfill}{rgb}{0.783160,0.914833,0.755894}%
\pgfsetfillcolor{currentfill}%
\pgfsetlinewidth{0.000000pt}%
\definecolor{currentstroke}{rgb}{0.000000,0.000000,0.000000}%
\pgfsetstrokecolor{currentstroke}%
\pgfsetstrokeopacity{0.000000}%
\pgfsetdash{}{0pt}%
\pgfpathmoveto{\pgfqpoint{1.392891in}{0.867332in}}%
\pgfpathlineto{\pgfqpoint{1.421405in}{0.865168in}}%
\pgfpathlineto{\pgfqpoint{1.419358in}{0.836488in}}%
\pgfpathlineto{\pgfqpoint{1.442781in}{0.834883in}}%
\pgfpathlineto{\pgfqpoint{1.440696in}{0.803294in}}%
\pgfpathlineto{\pgfqpoint{1.383417in}{0.806844in}}%
\pgfpathlineto{\pgfqpoint{1.385676in}{0.839061in}}%
\pgfpathlineto{\pgfqpoint{1.390704in}{0.838664in}}%
\pgfpathlineto{\pgfqpoint{1.392891in}{0.867332in}}%
\pgfpathclose%
\pgfusepath{fill}%
\end{pgfscope}%
\begin{pgfscope}%
\pgfpathrectangle{\pgfqpoint{0.100000in}{0.100000in}}{\pgfqpoint{3.007045in}{1.925000in}}%
\pgfusepath{clip}%
\pgfsetbuttcap%
\pgfsetmiterjoin%
\definecolor{currentfill}{rgb}{0.695071,0.877801,0.669865}%
\pgfsetfillcolor{currentfill}%
\pgfsetlinewidth{0.000000pt}%
\definecolor{currentstroke}{rgb}{0.000000,0.000000,0.000000}%
\pgfsetstrokecolor{currentstroke}%
\pgfsetstrokeopacity{0.000000}%
\pgfsetdash{}{0pt}%
\pgfpathmoveto{\pgfqpoint{1.180201in}{1.170746in}}%
\pgfpathlineto{\pgfqpoint{1.177735in}{1.174146in}}%
\pgfpathlineto{\pgfqpoint{1.172957in}{1.174640in}}%
\pgfpathlineto{\pgfqpoint{1.169868in}{1.170865in}}%
\pgfpathlineto{\pgfqpoint{1.164471in}{1.172432in}}%
\pgfpathlineto{\pgfqpoint{1.157547in}{1.181738in}}%
\pgfpathlineto{\pgfqpoint{1.146599in}{1.183208in}}%
\pgfpathlineto{\pgfqpoint{1.143492in}{1.187243in}}%
\pgfpathlineto{\pgfqpoint{1.142160in}{1.192960in}}%
\pgfpathlineto{\pgfqpoint{1.139110in}{1.197239in}}%
\pgfpathlineto{\pgfqpoint{1.141305in}{1.200383in}}%
\pgfpathlineto{\pgfqpoint{1.091973in}{1.207431in}}%
\pgfpathlineto{\pgfqpoint{1.059704in}{1.212390in}}%
\pgfpathlineto{\pgfqpoint{1.061002in}{1.220921in}}%
\pgfpathlineto{\pgfqpoint{1.062663in}{1.231515in}}%
\pgfpathlineto{\pgfqpoint{1.087451in}{1.227034in}}%
\pgfpathlineto{\pgfqpoint{1.087872in}{1.229867in}}%
\pgfpathlineto{\pgfqpoint{1.118778in}{1.225384in}}%
\pgfpathlineto{\pgfqpoint{1.120000in}{1.233895in}}%
\pgfpathlineto{\pgfqpoint{1.145415in}{1.230549in}}%
\pgfpathlineto{\pgfqpoint{1.146223in}{1.236379in}}%
\pgfpathlineto{\pgfqpoint{1.151850in}{1.235425in}}%
\pgfpathlineto{\pgfqpoint{1.153503in}{1.246887in}}%
\pgfpathlineto{\pgfqpoint{1.167476in}{1.245055in}}%
\pgfpathlineto{\pgfqpoint{1.166255in}{1.233746in}}%
\pgfpathlineto{\pgfqpoint{1.186601in}{1.231472in}}%
\pgfpathlineto{\pgfqpoint{1.205530in}{1.228275in}}%
\pgfpathlineto{\pgfqpoint{1.212843in}{1.220333in}}%
\pgfpathlineto{\pgfqpoint{1.214794in}{1.212322in}}%
\pgfpathlineto{\pgfqpoint{1.220004in}{1.212317in}}%
\pgfpathlineto{\pgfqpoint{1.226581in}{1.206947in}}%
\pgfpathlineto{\pgfqpoint{1.216380in}{1.193870in}}%
\pgfpathlineto{\pgfqpoint{1.203807in}{1.188442in}}%
\pgfpathlineto{\pgfqpoint{1.202987in}{1.177221in}}%
\pgfpathlineto{\pgfqpoint{1.201403in}{1.171929in}}%
\pgfpathlineto{\pgfqpoint{1.181782in}{1.174522in}}%
\pgfpathlineto{\pgfqpoint{1.180201in}{1.170746in}}%
\pgfpathclose%
\pgfusepath{fill}%
\end{pgfscope}%
\begin{pgfscope}%
\pgfpathrectangle{\pgfqpoint{0.100000in}{0.100000in}}{\pgfqpoint{3.007045in}{1.925000in}}%
\pgfusepath{clip}%
\pgfsetbuttcap%
\pgfsetmiterjoin%
\definecolor{currentfill}{rgb}{0.501961,0.790588,0.501438}%
\pgfsetfillcolor{currentfill}%
\pgfsetlinewidth{0.000000pt}%
\definecolor{currentstroke}{rgb}{0.000000,0.000000,0.000000}%
\pgfsetstrokecolor{currentstroke}%
\pgfsetstrokeopacity{0.000000}%
\pgfsetdash{}{0pt}%
\pgfpathmoveto{\pgfqpoint{2.593275in}{0.442838in}}%
\pgfpathlineto{\pgfqpoint{2.590385in}{0.434304in}}%
\pgfpathlineto{\pgfqpoint{2.597645in}{0.425944in}}%
\pgfpathlineto{\pgfqpoint{2.600616in}{0.426095in}}%
\pgfpathlineto{\pgfqpoint{2.607766in}{0.416135in}}%
\pgfpathlineto{\pgfqpoint{2.596119in}{0.413805in}}%
\pgfpathlineto{\pgfqpoint{2.597054in}{0.407964in}}%
\pgfpathlineto{\pgfqpoint{2.591178in}{0.407023in}}%
\pgfpathlineto{\pgfqpoint{2.592108in}{0.401133in}}%
\pgfpathlineto{\pgfqpoint{2.574578in}{0.398555in}}%
\pgfpathlineto{\pgfqpoint{2.571480in}{0.418988in}}%
\pgfpathlineto{\pgfqpoint{2.542441in}{0.414576in}}%
\pgfpathlineto{\pgfqpoint{2.539586in}{0.435076in}}%
\pgfpathlineto{\pgfqpoint{2.534472in}{0.469934in}}%
\pgfpathlineto{\pgfqpoint{2.530722in}{0.475320in}}%
\pgfpathlineto{\pgfqpoint{2.533154in}{0.479245in}}%
\pgfpathlineto{\pgfqpoint{2.538576in}{0.482276in}}%
\pgfpathlineto{\pgfqpoint{2.556078in}{0.484967in}}%
\pgfpathlineto{\pgfqpoint{2.556944in}{0.479141in}}%
\pgfpathlineto{\pgfqpoint{2.562765in}{0.480015in}}%
\pgfpathlineto{\pgfqpoint{2.565859in}{0.472579in}}%
\pgfpathlineto{\pgfqpoint{2.575198in}{0.465448in}}%
\pgfpathlineto{\pgfqpoint{2.580449in}{0.459760in}}%
\pgfpathlineto{\pgfqpoint{2.581515in}{0.455951in}}%
\pgfpathlineto{\pgfqpoint{2.587084in}{0.454555in}}%
\pgfpathlineto{\pgfqpoint{2.592017in}{0.447697in}}%
\pgfpathlineto{\pgfqpoint{2.593275in}{0.442838in}}%
\pgfpathclose%
\pgfusepath{fill}%
\end{pgfscope}%
\begin{pgfscope}%
\pgfpathrectangle{\pgfqpoint{0.100000in}{0.100000in}}{\pgfqpoint{3.007045in}{1.925000in}}%
\pgfusepath{clip}%
\pgfsetbuttcap%
\pgfsetmiterjoin%
\definecolor{currentfill}{rgb}{0.893887,0.959123,0.874002}%
\pgfsetfillcolor{currentfill}%
\pgfsetlinewidth{0.000000pt}%
\definecolor{currentstroke}{rgb}{0.000000,0.000000,0.000000}%
\pgfsetstrokecolor{currentstroke}%
\pgfsetstrokeopacity{0.000000}%
\pgfsetdash{}{0pt}%
\pgfpathmoveto{\pgfqpoint{0.441277in}{1.107129in}}%
\pgfpathlineto{\pgfqpoint{0.475084in}{1.097667in}}%
\pgfpathlineto{\pgfqpoint{0.522151in}{1.084892in}}%
\pgfpathlineto{\pgfqpoint{0.554172in}{1.076466in}}%
\pgfpathlineto{\pgfqpoint{0.573744in}{1.071996in}}%
\pgfpathlineto{\pgfqpoint{0.557933in}{1.009610in}}%
\pgfpathlineto{\pgfqpoint{0.515049in}{1.020412in}}%
\pgfpathlineto{\pgfqpoint{0.473306in}{1.031380in}}%
\pgfpathlineto{\pgfqpoint{0.472706in}{1.036029in}}%
\pgfpathlineto{\pgfqpoint{0.464444in}{1.039777in}}%
\pgfpathlineto{\pgfqpoint{0.465944in}{1.051393in}}%
\pgfpathlineto{\pgfqpoint{0.461538in}{1.053350in}}%
\pgfpathlineto{\pgfqpoint{0.463546in}{1.059117in}}%
\pgfpathlineto{\pgfqpoint{0.457558in}{1.060425in}}%
\pgfpathlineto{\pgfqpoint{0.459131in}{1.065999in}}%
\pgfpathlineto{\pgfqpoint{0.451737in}{1.068135in}}%
\pgfpathlineto{\pgfqpoint{0.453296in}{1.073681in}}%
\pgfpathlineto{\pgfqpoint{0.449631in}{1.074733in}}%
\pgfpathlineto{\pgfqpoint{0.451207in}{1.080317in}}%
\pgfpathlineto{\pgfqpoint{0.445664in}{1.083905in}}%
\pgfpathlineto{\pgfqpoint{0.442087in}{1.088884in}}%
\pgfpathlineto{\pgfqpoint{0.443707in}{1.094459in}}%
\pgfpathlineto{\pgfqpoint{0.438131in}{1.096030in}}%
\pgfpathlineto{\pgfqpoint{0.441277in}{1.107129in}}%
\pgfpathclose%
\pgfusepath{fill}%
\end{pgfscope}%
\begin{pgfscope}%
\pgfpathrectangle{\pgfqpoint{0.100000in}{0.100000in}}{\pgfqpoint{3.007045in}{1.925000in}}%
\pgfusepath{clip}%
\pgfsetbuttcap%
\pgfsetmiterjoin%
\definecolor{currentfill}{rgb}{0.345098,0.714802,0.408920}%
\pgfsetfillcolor{currentfill}%
\pgfsetlinewidth{0.000000pt}%
\definecolor{currentstroke}{rgb}{0.000000,0.000000,0.000000}%
\pgfsetstrokecolor{currentstroke}%
\pgfsetstrokeopacity{0.000000}%
\pgfsetdash{}{0pt}%
\pgfpathmoveto{\pgfqpoint{1.844539in}{1.237709in}}%
\pgfpathlineto{\pgfqpoint{1.833661in}{1.237372in}}%
\pgfpathlineto{\pgfqpoint{1.798792in}{1.236593in}}%
\pgfpathlineto{\pgfqpoint{1.798597in}{1.275256in}}%
\pgfpathlineto{\pgfqpoint{1.826979in}{1.275432in}}%
\pgfpathlineto{\pgfqpoint{1.843908in}{1.275998in}}%
\pgfpathlineto{\pgfqpoint{1.844539in}{1.237709in}}%
\pgfpathclose%
\pgfusepath{fill}%
\end{pgfscope}%
\begin{pgfscope}%
\pgfpathrectangle{\pgfqpoint{0.100000in}{0.100000in}}{\pgfqpoint{3.007045in}{1.925000in}}%
\pgfusepath{clip}%
\pgfsetbuttcap%
\pgfsetmiterjoin%
\definecolor{currentfill}{rgb}{0.524106,0.800923,0.519646}%
\pgfsetfillcolor{currentfill}%
\pgfsetlinewidth{0.000000pt}%
\definecolor{currentstroke}{rgb}{0.000000,0.000000,0.000000}%
\pgfsetstrokecolor{currentstroke}%
\pgfsetstrokeopacity{0.000000}%
\pgfsetdash{}{0pt}%
\pgfpathmoveto{\pgfqpoint{2.287383in}{1.019436in}}%
\pgfpathlineto{\pgfqpoint{2.286020in}{1.023971in}}%
\pgfpathlineto{\pgfqpoint{2.274407in}{1.023601in}}%
\pgfpathlineto{\pgfqpoint{2.271241in}{1.027836in}}%
\pgfpathlineto{\pgfqpoint{2.269307in}{1.034406in}}%
\pgfpathlineto{\pgfqpoint{2.263654in}{1.034861in}}%
\pgfpathlineto{\pgfqpoint{2.260460in}{1.038312in}}%
\pgfpathlineto{\pgfqpoint{2.259090in}{1.046090in}}%
\pgfpathlineto{\pgfqpoint{2.260372in}{1.054200in}}%
\pgfpathlineto{\pgfqpoint{2.261228in}{1.057564in}}%
\pgfpathlineto{\pgfqpoint{2.270021in}{1.058081in}}%
\pgfpathlineto{\pgfqpoint{2.276076in}{1.059145in}}%
\pgfpathlineto{\pgfqpoint{2.281792in}{1.056352in}}%
\pgfpathlineto{\pgfqpoint{2.285344in}{1.059690in}}%
\pgfpathlineto{\pgfqpoint{2.292317in}{1.056941in}}%
\pgfpathlineto{\pgfqpoint{2.291623in}{1.065687in}}%
\pgfpathlineto{\pgfqpoint{2.301227in}{1.065647in}}%
\pgfpathlineto{\pgfqpoint{2.302296in}{1.059040in}}%
\pgfpathlineto{\pgfqpoint{2.309458in}{1.055417in}}%
\pgfpathlineto{\pgfqpoint{2.310463in}{1.047437in}}%
\pgfpathlineto{\pgfqpoint{2.301947in}{1.038398in}}%
\pgfpathlineto{\pgfqpoint{2.302206in}{1.033847in}}%
\pgfpathlineto{\pgfqpoint{2.306270in}{1.030998in}}%
\pgfpathlineto{\pgfqpoint{2.292665in}{1.018214in}}%
\pgfpathlineto{\pgfqpoint{2.287383in}{1.019436in}}%
\pgfpathclose%
\pgfusepath{fill}%
\end{pgfscope}%
\begin{pgfscope}%
\pgfpathrectangle{\pgfqpoint{0.100000in}{0.100000in}}{\pgfqpoint{3.007045in}{1.925000in}}%
\pgfusepath{clip}%
\pgfsetbuttcap%
\pgfsetmiterjoin%
\definecolor{currentfill}{rgb}{0.751173,0.901423,0.724491}%
\pgfsetfillcolor{currentfill}%
\pgfsetlinewidth{0.000000pt}%
\definecolor{currentstroke}{rgb}{0.000000,0.000000,0.000000}%
\pgfsetstrokecolor{currentstroke}%
\pgfsetstrokeopacity{0.000000}%
\pgfsetdash{}{0pt}%
\pgfpathmoveto{\pgfqpoint{1.442781in}{0.834883in}}%
\pgfpathlineto{\pgfqpoint{1.419358in}{0.836488in}}%
\pgfpathlineto{\pgfqpoint{1.421405in}{0.865168in}}%
\pgfpathlineto{\pgfqpoint{1.442027in}{0.863792in}}%
\pgfpathlineto{\pgfqpoint{1.444186in}{0.892364in}}%
\pgfpathlineto{\pgfqpoint{1.473659in}{0.890543in}}%
\pgfpathlineto{\pgfqpoint{1.471813in}{0.861849in}}%
\pgfpathlineto{\pgfqpoint{1.478574in}{0.861428in}}%
\pgfpathlineto{\pgfqpoint{1.476764in}{0.832809in}}%
\pgfpathlineto{\pgfqpoint{1.471341in}{0.833174in}}%
\pgfpathlineto{\pgfqpoint{1.442781in}{0.834883in}}%
\pgfpathclose%
\pgfusepath{fill}%
\end{pgfscope}%
\begin{pgfscope}%
\pgfpathrectangle{\pgfqpoint{0.100000in}{0.100000in}}{\pgfqpoint{3.007045in}{1.925000in}}%
\pgfusepath{clip}%
\pgfsetbuttcap%
\pgfsetmiterjoin%
\definecolor{currentfill}{rgb}{0.451765,0.767090,0.461207}%
\pgfsetfillcolor{currentfill}%
\pgfsetlinewidth{0.000000pt}%
\definecolor{currentstroke}{rgb}{0.000000,0.000000,0.000000}%
\pgfsetstrokecolor{currentstroke}%
\pgfsetstrokeopacity{0.000000}%
\pgfsetdash{}{0pt}%
\pgfpathmoveto{\pgfqpoint{1.405343in}{1.057957in}}%
\pgfpathlineto{\pgfqpoint{1.404039in}{1.041199in}}%
\pgfpathlineto{\pgfqpoint{1.401987in}{1.015175in}}%
\pgfpathlineto{\pgfqpoint{1.367585in}{1.018108in}}%
\pgfpathlineto{\pgfqpoint{1.351668in}{1.019814in}}%
\pgfpathlineto{\pgfqpoint{1.347270in}{1.020215in}}%
\pgfpathlineto{\pgfqpoint{1.351564in}{1.062383in}}%
\pgfpathlineto{\pgfqpoint{1.405343in}{1.057957in}}%
\pgfpathclose%
\pgfusepath{fill}%
\end{pgfscope}%
\begin{pgfscope}%
\pgfpathrectangle{\pgfqpoint{0.100000in}{0.100000in}}{\pgfqpoint{3.007045in}{1.925000in}}%
\pgfusepath{clip}%
\pgfsetbuttcap%
\pgfsetmiterjoin%
\definecolor{currentfill}{rgb}{0.812687,0.926644,0.787389}%
\pgfsetfillcolor{currentfill}%
\pgfsetlinewidth{0.000000pt}%
\definecolor{currentstroke}{rgb}{0.000000,0.000000,0.000000}%
\pgfsetstrokecolor{currentstroke}%
\pgfsetstrokeopacity{0.000000}%
\pgfsetdash{}{0pt}%
\pgfpathmoveto{\pgfqpoint{2.254813in}{0.936868in}}%
\pgfpathlineto{\pgfqpoint{2.251280in}{0.935449in}}%
\pgfpathlineto{\pgfqpoint{2.240739in}{0.937396in}}%
\pgfpathlineto{\pgfqpoint{2.228851in}{0.933436in}}%
\pgfpathlineto{\pgfqpoint{2.226333in}{0.935993in}}%
\pgfpathlineto{\pgfqpoint{2.219736in}{0.935214in}}%
\pgfpathlineto{\pgfqpoint{2.210327in}{0.936644in}}%
\pgfpathlineto{\pgfqpoint{2.188078in}{0.943307in}}%
\pgfpathlineto{\pgfqpoint{2.188043in}{0.950490in}}%
\pgfpathlineto{\pgfqpoint{2.188978in}{0.963509in}}%
\pgfpathlineto{\pgfqpoint{2.187491in}{0.967104in}}%
\pgfpathlineto{\pgfqpoint{2.188596in}{0.973610in}}%
\pgfpathlineto{\pgfqpoint{2.181638in}{0.973954in}}%
\pgfpathlineto{\pgfqpoint{2.183256in}{0.977798in}}%
\pgfpathlineto{\pgfqpoint{2.189194in}{0.983871in}}%
\pgfpathlineto{\pgfqpoint{2.188886in}{0.995787in}}%
\pgfpathlineto{\pgfqpoint{2.191739in}{0.996087in}}%
\pgfpathlineto{\pgfqpoint{2.216487in}{0.998935in}}%
\pgfpathlineto{\pgfqpoint{2.230102in}{0.999971in}}%
\pgfpathlineto{\pgfqpoint{2.248936in}{1.000483in}}%
\pgfpathlineto{\pgfqpoint{2.258514in}{1.000440in}}%
\pgfpathlineto{\pgfqpoint{2.258356in}{0.992865in}}%
\pgfpathlineto{\pgfqpoint{2.258038in}{0.978829in}}%
\pgfpathlineto{\pgfqpoint{2.261071in}{0.973710in}}%
\pgfpathlineto{\pgfqpoint{2.261147in}{0.968764in}}%
\pgfpathlineto{\pgfqpoint{2.258371in}{0.965969in}}%
\pgfpathlineto{\pgfqpoint{2.253187in}{0.966596in}}%
\pgfpathlineto{\pgfqpoint{2.248417in}{0.962697in}}%
\pgfpathlineto{\pgfqpoint{2.251263in}{0.956468in}}%
\pgfpathlineto{\pgfqpoint{2.257521in}{0.951222in}}%
\pgfpathlineto{\pgfqpoint{2.254335in}{0.941757in}}%
\pgfpathlineto{\pgfqpoint{2.254813in}{0.936868in}}%
\pgfpathclose%
\pgfusepath{fill}%
\end{pgfscope}%
\begin{pgfscope}%
\pgfpathrectangle{\pgfqpoint{0.100000in}{0.100000in}}{\pgfqpoint{3.007045in}{1.925000in}}%
\pgfusepath{clip}%
\pgfsetbuttcap%
\pgfsetmiterjoin%
\definecolor{currentfill}{rgb}{0.808997,0.925167,0.783453}%
\pgfsetfillcolor{currentfill}%
\pgfsetlinewidth{0.000000pt}%
\definecolor{currentstroke}{rgb}{0.000000,0.000000,0.000000}%
\pgfsetstrokecolor{currentstroke}%
\pgfsetstrokeopacity{0.000000}%
\pgfsetdash{}{0pt}%
\pgfpathmoveto{\pgfqpoint{0.421673in}{1.178083in}}%
\pgfpathlineto{\pgfqpoint{0.412185in}{1.187494in}}%
\pgfpathlineto{\pgfqpoint{0.409170in}{1.194540in}}%
\pgfpathlineto{\pgfqpoint{0.413081in}{1.203149in}}%
\pgfpathlineto{\pgfqpoint{0.413366in}{1.209180in}}%
\pgfpathlineto{\pgfqpoint{0.404869in}{1.211825in}}%
\pgfpathlineto{\pgfqpoint{0.404848in}{1.220444in}}%
\pgfpathlineto{\pgfqpoint{0.408920in}{1.226474in}}%
\pgfpathlineto{\pgfqpoint{0.407661in}{1.233230in}}%
\pgfpathlineto{\pgfqpoint{0.407690in}{1.233195in}}%
\pgfpathlineto{\pgfqpoint{0.424991in}{1.242670in}}%
\pgfpathlineto{\pgfqpoint{0.436885in}{1.243705in}}%
\pgfpathlineto{\pgfqpoint{0.440085in}{1.241248in}}%
\pgfpathlineto{\pgfqpoint{0.446079in}{1.262743in}}%
\pgfpathlineto{\pgfqpoint{0.455151in}{1.243275in}}%
\pgfpathlineto{\pgfqpoint{0.464303in}{1.248933in}}%
\pgfpathlineto{\pgfqpoint{0.475080in}{1.259356in}}%
\pgfpathlineto{\pgfqpoint{0.488781in}{1.270732in}}%
\pgfpathlineto{\pgfqpoint{0.497623in}{1.272306in}}%
\pgfpathlineto{\pgfqpoint{0.502964in}{1.265519in}}%
\pgfpathlineto{\pgfqpoint{0.510697in}{1.267522in}}%
\pgfpathlineto{\pgfqpoint{0.514663in}{1.260246in}}%
\pgfpathlineto{\pgfqpoint{0.513336in}{1.257851in}}%
\pgfpathlineto{\pgfqpoint{0.514807in}{1.248637in}}%
\pgfpathlineto{\pgfqpoint{0.525021in}{1.240725in}}%
\pgfpathlineto{\pgfqpoint{0.524524in}{1.231551in}}%
\pgfpathlineto{\pgfqpoint{0.527593in}{1.228218in}}%
\pgfpathlineto{\pgfqpoint{0.527157in}{1.220456in}}%
\pgfpathlineto{\pgfqpoint{0.522889in}{1.217818in}}%
\pgfpathlineto{\pgfqpoint{0.521553in}{1.220832in}}%
\pgfpathlineto{\pgfqpoint{0.503936in}{1.210812in}}%
\pgfpathlineto{\pgfqpoint{0.502749in}{1.206653in}}%
\pgfpathlineto{\pgfqpoint{0.498018in}{1.202710in}}%
\pgfpathlineto{\pgfqpoint{0.492430in}{1.204219in}}%
\pgfpathlineto{\pgfqpoint{0.473609in}{1.193492in}}%
\pgfpathlineto{\pgfqpoint{0.461047in}{1.194974in}}%
\pgfpathlineto{\pgfqpoint{0.450589in}{1.194118in}}%
\pgfpathlineto{\pgfqpoint{0.421673in}{1.178083in}}%
\pgfpathclose%
\pgfusepath{fill}%
\end{pgfscope}%
\begin{pgfscope}%
\pgfpathrectangle{\pgfqpoint{0.100000in}{0.100000in}}{\pgfqpoint{3.007045in}{1.925000in}}%
\pgfusepath{clip}%
\pgfsetbuttcap%
\pgfsetmiterjoin%
\definecolor{currentfill}{rgb}{0.794233,0.919262,0.767705}%
\pgfsetfillcolor{currentfill}%
\pgfsetlinewidth{0.000000pt}%
\definecolor{currentstroke}{rgb}{0.000000,0.000000,0.000000}%
\pgfsetstrokecolor{currentstroke}%
\pgfsetstrokeopacity{0.000000}%
\pgfsetdash{}{0pt}%
\pgfpathmoveto{\pgfqpoint{1.403892in}{1.607820in}}%
\pgfpathlineto{\pgfqpoint{1.355655in}{1.612365in}}%
\pgfpathlineto{\pgfqpoint{1.359600in}{1.651719in}}%
\pgfpathlineto{\pgfqpoint{1.370707in}{1.650612in}}%
\pgfpathlineto{\pgfqpoint{1.371277in}{1.656487in}}%
\pgfpathlineto{\pgfqpoint{1.379987in}{1.655666in}}%
\pgfpathlineto{\pgfqpoint{1.410950in}{1.652859in}}%
\pgfpathlineto{\pgfqpoint{1.409091in}{1.629777in}}%
\pgfpathlineto{\pgfqpoint{1.405844in}{1.630008in}}%
\pgfpathlineto{\pgfqpoint{1.403892in}{1.607820in}}%
\pgfpathclose%
\pgfusepath{fill}%
\end{pgfscope}%
\begin{pgfscope}%
\pgfpathrectangle{\pgfqpoint{0.100000in}{0.100000in}}{\pgfqpoint{3.007045in}{1.925000in}}%
\pgfusepath{clip}%
\pgfsetbuttcap%
\pgfsetmiterjoin%
\definecolor{currentfill}{rgb}{0.338824,0.711726,0.405844}%
\pgfsetfillcolor{currentfill}%
\pgfsetlinewidth{0.000000pt}%
\definecolor{currentstroke}{rgb}{0.000000,0.000000,0.000000}%
\pgfsetstrokecolor{currentstroke}%
\pgfsetstrokeopacity{0.000000}%
\pgfsetdash{}{0pt}%
\pgfpathmoveto{\pgfqpoint{2.751955in}{1.186579in}}%
\pgfpathlineto{\pgfqpoint{2.748588in}{1.187139in}}%
\pgfpathlineto{\pgfqpoint{2.739288in}{1.183920in}}%
\pgfpathlineto{\pgfqpoint{2.728190in}{1.196931in}}%
\pgfpathlineto{\pgfqpoint{2.729962in}{1.199364in}}%
\pgfpathlineto{\pgfqpoint{2.728781in}{1.207289in}}%
\pgfpathlineto{\pgfqpoint{2.734998in}{1.209843in}}%
\pgfpathlineto{\pgfqpoint{2.724066in}{1.216278in}}%
\pgfpathlineto{\pgfqpoint{2.728601in}{1.227398in}}%
\pgfpathlineto{\pgfqpoint{2.725297in}{1.231605in}}%
\pgfpathlineto{\pgfqpoint{2.725039in}{1.237269in}}%
\pgfpathlineto{\pgfqpoint{2.721783in}{1.241309in}}%
\pgfpathlineto{\pgfqpoint{2.724195in}{1.253074in}}%
\pgfpathlineto{\pgfqpoint{2.730272in}{1.258328in}}%
\pgfpathlineto{\pgfqpoint{2.738533in}{1.258792in}}%
\pgfpathlineto{\pgfqpoint{2.744257in}{1.261114in}}%
\pgfpathlineto{\pgfqpoint{2.745705in}{1.255995in}}%
\pgfpathlineto{\pgfqpoint{2.758832in}{1.209240in}}%
\pgfpathlineto{\pgfqpoint{2.753781in}{1.202785in}}%
\pgfpathlineto{\pgfqpoint{2.751060in}{1.190437in}}%
\pgfpathlineto{\pgfqpoint{2.751955in}{1.186579in}}%
\pgfpathclose%
\pgfusepath{fill}%
\end{pgfscope}%
\begin{pgfscope}%
\pgfpathrectangle{\pgfqpoint{0.100000in}{0.100000in}}{\pgfqpoint{3.007045in}{1.925000in}}%
\pgfusepath{clip}%
\pgfsetbuttcap%
\pgfsetmiterjoin%
\definecolor{currentfill}{rgb}{0.638970,0.854179,0.615240}%
\pgfsetfillcolor{currentfill}%
\pgfsetlinewidth{0.000000pt}%
\definecolor{currentstroke}{rgb}{0.000000,0.000000,0.000000}%
\pgfsetstrokecolor{currentstroke}%
\pgfsetstrokeopacity{0.000000}%
\pgfsetdash{}{0pt}%
\pgfpathmoveto{\pgfqpoint{2.124853in}{1.498938in}}%
\pgfpathlineto{\pgfqpoint{2.108536in}{1.497608in}}%
\pgfpathlineto{\pgfqpoint{2.108968in}{1.491853in}}%
\pgfpathlineto{\pgfqpoint{2.101634in}{1.491321in}}%
\pgfpathlineto{\pgfqpoint{2.094574in}{1.490855in}}%
\pgfpathlineto{\pgfqpoint{2.092961in}{1.513610in}}%
\pgfpathlineto{\pgfqpoint{2.090405in}{1.513286in}}%
\pgfpathlineto{\pgfqpoint{2.088935in}{1.530951in}}%
\pgfpathlineto{\pgfqpoint{2.077656in}{1.530311in}}%
\pgfpathlineto{\pgfqpoint{2.076668in}{1.547607in}}%
\pgfpathlineto{\pgfqpoint{2.069326in}{1.547117in}}%
\pgfpathlineto{\pgfqpoint{2.067108in}{1.552783in}}%
\pgfpathlineto{\pgfqpoint{2.066412in}{1.564272in}}%
\pgfpathlineto{\pgfqpoint{2.078017in}{1.564986in}}%
\pgfpathlineto{\pgfqpoint{2.083707in}{1.565240in}}%
\pgfpathlineto{\pgfqpoint{2.084400in}{1.553797in}}%
\pgfpathlineto{\pgfqpoint{2.090120in}{1.553996in}}%
\pgfpathlineto{\pgfqpoint{2.093706in}{1.548428in}}%
\pgfpathlineto{\pgfqpoint{2.094036in}{1.542696in}}%
\pgfpathlineto{\pgfqpoint{2.111077in}{1.540159in}}%
\pgfpathlineto{\pgfqpoint{2.102036in}{1.523219in}}%
\pgfpathlineto{\pgfqpoint{2.100053in}{1.512797in}}%
\pgfpathlineto{\pgfqpoint{2.104771in}{1.510691in}}%
\pgfpathlineto{\pgfqpoint{2.108185in}{1.516689in}}%
\pgfpathlineto{\pgfqpoint{2.112193in}{1.518509in}}%
\pgfpathlineto{\pgfqpoint{2.118815in}{1.532207in}}%
\pgfpathlineto{\pgfqpoint{2.124890in}{1.534323in}}%
\pgfpathlineto{\pgfqpoint{2.127996in}{1.537667in}}%
\pgfpathlineto{\pgfqpoint{2.133801in}{1.550138in}}%
\pgfpathlineto{\pgfqpoint{2.134382in}{1.555481in}}%
\pgfpathlineto{\pgfqpoint{2.143696in}{1.559230in}}%
\pgfpathlineto{\pgfqpoint{2.143649in}{1.552093in}}%
\pgfpathlineto{\pgfqpoint{2.137836in}{1.542472in}}%
\pgfpathlineto{\pgfqpoint{2.137692in}{1.536286in}}%
\pgfpathlineto{\pgfqpoint{2.134764in}{1.533978in}}%
\pgfpathlineto{\pgfqpoint{2.128245in}{1.517619in}}%
\pgfpathlineto{\pgfqpoint{2.124853in}{1.498938in}}%
\pgfpathclose%
\pgfusepath{fill}%
\end{pgfscope}%
\begin{pgfscope}%
\pgfpathrectangle{\pgfqpoint{0.100000in}{0.100000in}}{\pgfqpoint{3.007045in}{1.925000in}}%
\pgfusepath{clip}%
\pgfsetbuttcap%
\pgfsetmiterjoin%
\definecolor{currentfill}{rgb}{0.568397,0.821592,0.556063}%
\pgfsetfillcolor{currentfill}%
\pgfsetlinewidth{0.000000pt}%
\definecolor{currentstroke}{rgb}{0.000000,0.000000,0.000000}%
\pgfsetstrokecolor{currentstroke}%
\pgfsetstrokeopacity{0.000000}%
\pgfsetdash{}{0pt}%
\pgfpathmoveto{\pgfqpoint{2.636603in}{1.298784in}}%
\pgfpathlineto{\pgfqpoint{2.631148in}{1.284175in}}%
\pgfpathlineto{\pgfqpoint{2.627614in}{1.286206in}}%
\pgfpathlineto{\pgfqpoint{2.616493in}{1.288353in}}%
\pgfpathlineto{\pgfqpoint{2.610752in}{1.291249in}}%
\pgfpathlineto{\pgfqpoint{2.608820in}{1.295800in}}%
\pgfpathlineto{\pgfqpoint{2.612866in}{1.308899in}}%
\pgfpathlineto{\pgfqpoint{2.608640in}{1.313232in}}%
\pgfpathlineto{\pgfqpoint{2.609389in}{1.317404in}}%
\pgfpathlineto{\pgfqpoint{2.606011in}{1.320523in}}%
\pgfpathlineto{\pgfqpoint{2.599096in}{1.322608in}}%
\pgfpathlineto{\pgfqpoint{2.576932in}{1.318538in}}%
\pgfpathlineto{\pgfqpoint{2.574642in}{1.330275in}}%
\pgfpathlineto{\pgfqpoint{2.554469in}{1.327047in}}%
\pgfpathlineto{\pgfqpoint{2.552895in}{1.336088in}}%
\pgfpathlineto{\pgfqpoint{2.549986in}{1.354398in}}%
\pgfpathlineto{\pgfqpoint{2.555437in}{1.355961in}}%
\pgfpathlineto{\pgfqpoint{2.561808in}{1.358520in}}%
\pgfpathlineto{\pgfqpoint{2.573880in}{1.354042in}}%
\pgfpathlineto{\pgfqpoint{2.605715in}{1.356835in}}%
\pgfpathlineto{\pgfqpoint{2.609177in}{1.353242in}}%
\pgfpathlineto{\pgfqpoint{2.621003in}{1.357074in}}%
\pgfpathlineto{\pgfqpoint{2.623988in}{1.352894in}}%
\pgfpathlineto{\pgfqpoint{2.627962in}{1.353837in}}%
\pgfpathlineto{\pgfqpoint{2.637385in}{1.345841in}}%
\pgfpathlineto{\pgfqpoint{2.645166in}{1.349038in}}%
\pgfpathlineto{\pgfqpoint{2.654171in}{1.354581in}}%
\pgfpathlineto{\pgfqpoint{2.645950in}{1.339764in}}%
\pgfpathlineto{\pgfqpoint{2.633083in}{1.330389in}}%
\pgfpathlineto{\pgfqpoint{2.626236in}{1.316064in}}%
\pgfpathlineto{\pgfqpoint{2.628110in}{1.312979in}}%
\pgfpathlineto{\pgfqpoint{2.625195in}{1.307053in}}%
\pgfpathlineto{\pgfqpoint{2.627758in}{1.303889in}}%
\pgfpathlineto{\pgfqpoint{2.632407in}{1.305014in}}%
\pgfpathlineto{\pgfqpoint{2.636603in}{1.298784in}}%
\pgfpathclose%
\pgfusepath{fill}%
\end{pgfscope}%
\begin{pgfscope}%
\pgfpathrectangle{\pgfqpoint{0.100000in}{0.100000in}}{\pgfqpoint{3.007045in}{1.925000in}}%
\pgfusepath{clip}%
\pgfsetbuttcap%
\pgfsetmiterjoin%
\definecolor{currentfill}{rgb}{0.463206,0.772503,0.469573}%
\pgfsetfillcolor{currentfill}%
\pgfsetlinewidth{0.000000pt}%
\definecolor{currentstroke}{rgb}{0.000000,0.000000,0.000000}%
\pgfsetstrokecolor{currentstroke}%
\pgfsetstrokeopacity{0.000000}%
\pgfsetdash{}{0pt}%
\pgfpathmoveto{\pgfqpoint{1.847419in}{1.345161in}}%
\pgfpathlineto{\pgfqpoint{1.813145in}{1.344757in}}%
\pgfpathlineto{\pgfqpoint{1.779241in}{1.344556in}}%
\pgfpathlineto{\pgfqpoint{1.779293in}{1.361999in}}%
\pgfpathlineto{\pgfqpoint{1.776477in}{1.362015in}}%
\pgfpathlineto{\pgfqpoint{1.776480in}{1.390760in}}%
\pgfpathlineto{\pgfqpoint{1.822052in}{1.390899in}}%
\pgfpathlineto{\pgfqpoint{1.844917in}{1.391266in}}%
\pgfpathlineto{\pgfqpoint{1.845372in}{1.362394in}}%
\pgfpathlineto{\pgfqpoint{1.847146in}{1.362415in}}%
\pgfpathlineto{\pgfqpoint{1.847419in}{1.345161in}}%
\pgfpathclose%
\pgfusepath{fill}%
\end{pgfscope}%
\begin{pgfscope}%
\pgfpathrectangle{\pgfqpoint{0.100000in}{0.100000in}}{\pgfqpoint{3.007045in}{1.925000in}}%
\pgfusepath{clip}%
\pgfsetbuttcap%
\pgfsetmiterjoin%
\definecolor{currentfill}{rgb}{0.634295,0.852211,0.610688}%
\pgfsetfillcolor{currentfill}%
\pgfsetlinewidth{0.000000pt}%
\definecolor{currentstroke}{rgb}{0.000000,0.000000,0.000000}%
\pgfsetstrokecolor{currentstroke}%
\pgfsetstrokeopacity{0.000000}%
\pgfsetdash{}{0pt}%
\pgfpathmoveto{\pgfqpoint{2.905539in}{1.431014in}}%
\pgfpathlineto{\pgfqpoint{2.901522in}{1.431102in}}%
\pgfpathlineto{\pgfqpoint{2.893781in}{1.426472in}}%
\pgfpathlineto{\pgfqpoint{2.889192in}{1.426210in}}%
\pgfpathlineto{\pgfqpoint{2.881149in}{1.420623in}}%
\pgfpathlineto{\pgfqpoint{2.877803in}{1.421007in}}%
\pgfpathlineto{\pgfqpoint{2.855094in}{1.414428in}}%
\pgfpathlineto{\pgfqpoint{2.824324in}{1.387968in}}%
\pgfpathlineto{\pgfqpoint{2.819156in}{1.394508in}}%
\pgfpathlineto{\pgfqpoint{2.829259in}{1.404526in}}%
\pgfpathlineto{\pgfqpoint{2.824695in}{1.409031in}}%
\pgfpathlineto{\pgfqpoint{2.822125in}{1.424128in}}%
\pgfpathlineto{\pgfqpoint{2.816097in}{1.458112in}}%
\pgfpathlineto{\pgfqpoint{2.836904in}{1.462520in}}%
\pgfpathlineto{\pgfqpoint{2.880647in}{1.472875in}}%
\pgfpathlineto{\pgfqpoint{2.896973in}{1.475592in}}%
\pgfpathlineto{\pgfqpoint{2.903419in}{1.453280in}}%
\pgfpathlineto{\pgfqpoint{2.906901in}{1.437874in}}%
\pgfpathlineto{\pgfqpoint{2.905539in}{1.431014in}}%
\pgfpathclose%
\pgfusepath{fill}%
\end{pgfscope}%
\begin{pgfscope}%
\pgfpathrectangle{\pgfqpoint{0.100000in}{0.100000in}}{\pgfqpoint{3.007045in}{1.925000in}}%
\pgfusepath{clip}%
\pgfsetbuttcap%
\pgfsetmiterjoin%
\definecolor{currentfill}{rgb}{0.439216,0.760938,0.455056}%
\pgfsetfillcolor{currentfill}%
\pgfsetlinewidth{0.000000pt}%
\definecolor{currentstroke}{rgb}{0.000000,0.000000,0.000000}%
\pgfsetstrokecolor{currentstroke}%
\pgfsetstrokeopacity{0.000000}%
\pgfsetdash{}{0pt}%
\pgfpathmoveto{\pgfqpoint{1.977807in}{1.222266in}}%
\pgfpathlineto{\pgfqpoint{1.977971in}{1.216245in}}%
\pgfpathlineto{\pgfqpoint{1.948122in}{1.215494in}}%
\pgfpathlineto{\pgfqpoint{1.948406in}{1.218688in}}%
\pgfpathlineto{\pgfqpoint{1.925598in}{1.218425in}}%
\pgfpathlineto{\pgfqpoint{1.925211in}{1.241392in}}%
\pgfpathlineto{\pgfqpoint{1.935891in}{1.242279in}}%
\pgfpathlineto{\pgfqpoint{1.935244in}{1.278471in}}%
\pgfpathlineto{\pgfqpoint{1.952380in}{1.279020in}}%
\pgfpathlineto{\pgfqpoint{1.952591in}{1.273282in}}%
\pgfpathlineto{\pgfqpoint{1.973511in}{1.274071in}}%
\pgfpathlineto{\pgfqpoint{1.981668in}{1.274180in}}%
\pgfpathlineto{\pgfqpoint{1.982776in}{1.245666in}}%
\pgfpathlineto{\pgfqpoint{1.977034in}{1.245635in}}%
\pgfpathlineto{\pgfqpoint{1.977807in}{1.222266in}}%
\pgfpathclose%
\pgfusepath{fill}%
\end{pgfscope}%
\begin{pgfscope}%
\pgfpathrectangle{\pgfqpoint{0.100000in}{0.100000in}}{\pgfqpoint{3.007045in}{1.925000in}}%
\pgfusepath{clip}%
\pgfsetbuttcap%
\pgfsetmiterjoin%
\definecolor{currentfill}{rgb}{0.568397,0.821592,0.556063}%
\pgfsetfillcolor{currentfill}%
\pgfsetlinewidth{0.000000pt}%
\definecolor{currentstroke}{rgb}{0.000000,0.000000,0.000000}%
\pgfsetstrokecolor{currentstroke}%
\pgfsetstrokeopacity{0.000000}%
\pgfsetdash{}{0pt}%
\pgfpathmoveto{\pgfqpoint{2.348646in}{0.902478in}}%
\pgfpathlineto{\pgfqpoint{2.332547in}{0.900599in}}%
\pgfpathlineto{\pgfqpoint{2.313442in}{0.898402in}}%
\pgfpathlineto{\pgfqpoint{2.309785in}{0.907283in}}%
\pgfpathlineto{\pgfqpoint{2.313145in}{0.917931in}}%
\pgfpathlineto{\pgfqpoint{2.308274in}{0.925835in}}%
\pgfpathlineto{\pgfqpoint{2.313025in}{0.928226in}}%
\pgfpathlineto{\pgfqpoint{2.319657in}{0.939932in}}%
\pgfpathlineto{\pgfqpoint{2.318851in}{0.944884in}}%
\pgfpathlineto{\pgfqpoint{2.321351in}{0.950113in}}%
\pgfpathlineto{\pgfqpoint{2.327024in}{0.948030in}}%
\pgfpathlineto{\pgfqpoint{2.329679in}{0.943109in}}%
\pgfpathlineto{\pgfqpoint{2.341668in}{0.946890in}}%
\pgfpathlineto{\pgfqpoint{2.354830in}{0.942630in}}%
\pgfpathlineto{\pgfqpoint{2.364223in}{0.935952in}}%
\pgfpathlineto{\pgfqpoint{2.361650in}{0.931386in}}%
\pgfpathlineto{\pgfqpoint{2.361956in}{0.924269in}}%
\pgfpathlineto{\pgfqpoint{2.356471in}{0.920502in}}%
\pgfpathlineto{\pgfqpoint{2.351726in}{0.921503in}}%
\pgfpathlineto{\pgfqpoint{2.348504in}{0.918229in}}%
\pgfpathlineto{\pgfqpoint{2.348646in}{0.902478in}}%
\pgfpathclose%
\pgfusepath{fill}%
\end{pgfscope}%
\begin{pgfscope}%
\pgfpathrectangle{\pgfqpoint{0.100000in}{0.100000in}}{\pgfqpoint{3.007045in}{1.925000in}}%
\pgfusepath{clip}%
\pgfsetbuttcap%
\pgfsetmiterjoin%
\definecolor{currentfill}{rgb}{0.363922,0.724029,0.418147}%
\pgfsetfillcolor{currentfill}%
\pgfsetlinewidth{0.000000pt}%
\definecolor{currentstroke}{rgb}{0.000000,0.000000,0.000000}%
\pgfsetstrokecolor{currentstroke}%
\pgfsetstrokeopacity{0.000000}%
\pgfsetdash{}{0pt}%
\pgfpathmoveto{\pgfqpoint{1.280556in}{0.903442in}}%
\pgfpathlineto{\pgfqpoint{1.281006in}{0.908273in}}%
\pgfpathlineto{\pgfqpoint{1.218495in}{0.915021in}}%
\pgfpathlineto{\pgfqpoint{1.217220in}{0.903588in}}%
\pgfpathlineto{\pgfqpoint{1.194470in}{0.906193in}}%
\pgfpathlineto{\pgfqpoint{1.201592in}{0.967633in}}%
\pgfpathlineto{\pgfqpoint{1.209230in}{0.966315in}}%
\pgfpathlineto{\pgfqpoint{1.218512in}{0.975812in}}%
\pgfpathlineto{\pgfqpoint{1.221151in}{0.983550in}}%
\pgfpathlineto{\pgfqpoint{1.223675in}{0.983885in}}%
\pgfpathlineto{\pgfqpoint{1.225457in}{0.993376in}}%
\pgfpathlineto{\pgfqpoint{1.225227in}{1.010762in}}%
\pgfpathlineto{\pgfqpoint{1.231839in}{1.013157in}}%
\pgfpathlineto{\pgfqpoint{1.235393in}{1.026018in}}%
\pgfpathlineto{\pgfqpoint{1.235453in}{1.031177in}}%
\pgfpathlineto{\pgfqpoint{1.238883in}{1.030797in}}%
\pgfpathlineto{\pgfqpoint{1.277959in}{1.026457in}}%
\pgfpathlineto{\pgfqpoint{1.298942in}{1.024507in}}%
\pgfpathlineto{\pgfqpoint{1.293808in}{0.973457in}}%
\pgfpathlineto{\pgfqpoint{1.271190in}{0.975810in}}%
\pgfpathlineto{\pgfqpoint{1.273752in}{0.963960in}}%
\pgfpathlineto{\pgfqpoint{1.273425in}{0.954257in}}%
\pgfpathlineto{\pgfqpoint{1.271875in}{0.946736in}}%
\pgfpathlineto{\pgfqpoint{1.290299in}{0.944870in}}%
\pgfpathlineto{\pgfqpoint{1.292797in}{0.945962in}}%
\pgfpathlineto{\pgfqpoint{1.303832in}{0.919724in}}%
\pgfpathlineto{\pgfqpoint{1.308509in}{0.919262in}}%
\pgfpathlineto{\pgfqpoint{1.307382in}{0.907359in}}%
\pgfpathlineto{\pgfqpoint{1.295559in}{0.908588in}}%
\pgfpathlineto{\pgfqpoint{1.280556in}{0.903442in}}%
\pgfpathclose%
\pgfusepath{fill}%
\end{pgfscope}%
\begin{pgfscope}%
\pgfpathrectangle{\pgfqpoint{0.100000in}{0.100000in}}{\pgfqpoint{3.007045in}{1.925000in}}%
\pgfusepath{clip}%
\pgfsetbuttcap%
\pgfsetmiterjoin%
\definecolor{currentfill}{rgb}{0.607151,0.839677,0.587928}%
\pgfsetfillcolor{currentfill}%
\pgfsetlinewidth{0.000000pt}%
\definecolor{currentstroke}{rgb}{0.000000,0.000000,0.000000}%
\pgfsetstrokecolor{currentstroke}%
\pgfsetstrokeopacity{0.000000}%
\pgfsetdash{}{0pt}%
\pgfpathmoveto{\pgfqpoint{1.905724in}{1.300373in}}%
\pgfpathlineto{\pgfqpoint{1.854836in}{1.299029in}}%
\pgfpathlineto{\pgfqpoint{1.853741in}{1.305138in}}%
\pgfpathlineto{\pgfqpoint{1.853457in}{1.322433in}}%
\pgfpathlineto{\pgfqpoint{1.859156in}{1.322514in}}%
\pgfpathlineto{\pgfqpoint{1.858723in}{1.345342in}}%
\pgfpathlineto{\pgfqpoint{1.869963in}{1.345597in}}%
\pgfpathlineto{\pgfqpoint{1.881395in}{1.345922in}}%
\pgfpathlineto{\pgfqpoint{1.881224in}{1.351626in}}%
\pgfpathlineto{\pgfqpoint{1.904114in}{1.352301in}}%
\pgfpathlineto{\pgfqpoint{1.905724in}{1.300373in}}%
\pgfpathclose%
\pgfusepath{fill}%
\end{pgfscope}%
\begin{pgfscope}%
\pgfpathrectangle{\pgfqpoint{0.100000in}{0.100000in}}{\pgfqpoint{3.007045in}{1.925000in}}%
\pgfusepath{clip}%
\pgfsetbuttcap%
\pgfsetmiterjoin%
\definecolor{currentfill}{rgb}{0.463206,0.772503,0.469573}%
\pgfsetfillcolor{currentfill}%
\pgfsetlinewidth{0.000000pt}%
\definecolor{currentstroke}{rgb}{0.000000,0.000000,0.000000}%
\pgfsetstrokecolor{currentstroke}%
\pgfsetstrokeopacity{0.000000}%
\pgfsetdash{}{0pt}%
\pgfpathmoveto{\pgfqpoint{0.615468in}{1.284198in}}%
\pgfpathlineto{\pgfqpoint{0.597849in}{1.283744in}}%
\pgfpathlineto{\pgfqpoint{0.596506in}{1.278590in}}%
\pgfpathlineto{\pgfqpoint{0.614328in}{1.243704in}}%
\pgfpathlineto{\pgfqpoint{0.571645in}{1.215280in}}%
\pgfpathlineto{\pgfqpoint{0.567833in}{1.216455in}}%
\pgfpathlineto{\pgfqpoint{0.540209in}{1.259221in}}%
\pgfpathlineto{\pgfqpoint{0.552667in}{1.255842in}}%
\pgfpathlineto{\pgfqpoint{0.560395in}{1.283604in}}%
\pgfpathlineto{\pgfqpoint{0.554814in}{1.285052in}}%
\pgfpathlineto{\pgfqpoint{0.556383in}{1.290907in}}%
\pgfpathlineto{\pgfqpoint{0.562826in}{1.298061in}}%
\pgfpathlineto{\pgfqpoint{0.615468in}{1.284198in}}%
\pgfpathclose%
\pgfusepath{fill}%
\end{pgfscope}%
\begin{pgfscope}%
\pgfpathrectangle{\pgfqpoint{0.100000in}{0.100000in}}{\pgfqpoint{3.007045in}{1.925000in}}%
\pgfusepath{clip}%
\pgfsetbuttcap%
\pgfsetmiterjoin%
\definecolor{currentfill}{rgb}{0.263529,0.674817,0.368935}%
\pgfsetfillcolor{currentfill}%
\pgfsetlinewidth{0.000000pt}%
\definecolor{currentstroke}{rgb}{0.000000,0.000000,0.000000}%
\pgfsetstrokecolor{currentstroke}%
\pgfsetstrokeopacity{0.000000}%
\pgfsetdash{}{0pt}%
\pgfpathmoveto{\pgfqpoint{2.514084in}{1.019211in}}%
\pgfpathlineto{\pgfqpoint{2.506602in}{1.020237in}}%
\pgfpathlineto{\pgfqpoint{2.500787in}{1.015480in}}%
\pgfpathlineto{\pgfqpoint{2.492311in}{1.025615in}}%
\pgfpathlineto{\pgfqpoint{2.492133in}{1.027383in}}%
\pgfpathlineto{\pgfqpoint{2.475027in}{1.025720in}}%
\pgfpathlineto{\pgfqpoint{2.476402in}{1.027509in}}%
\pgfpathlineto{\pgfqpoint{2.481176in}{1.031217in}}%
\pgfpathlineto{\pgfqpoint{2.481841in}{1.034709in}}%
\pgfpathlineto{\pgfqpoint{2.495062in}{1.040390in}}%
\pgfpathlineto{\pgfqpoint{2.503517in}{1.042080in}}%
\pgfpathlineto{\pgfqpoint{2.505363in}{1.044619in}}%
\pgfpathlineto{\pgfqpoint{2.521063in}{1.052157in}}%
\pgfpathlineto{\pgfqpoint{2.526060in}{1.056473in}}%
\pgfpathlineto{\pgfqpoint{2.537236in}{1.043373in}}%
\pgfpathlineto{\pgfqpoint{2.537514in}{1.041038in}}%
\pgfpathlineto{\pgfqpoint{2.529927in}{1.036274in}}%
\pgfpathlineto{\pgfqpoint{2.531083in}{1.032312in}}%
\pgfpathlineto{\pgfqpoint{2.515896in}{1.030173in}}%
\pgfpathlineto{\pgfqpoint{2.513186in}{1.021977in}}%
\pgfpathlineto{\pgfqpoint{2.514084in}{1.019211in}}%
\pgfpathclose%
\pgfusepath{fill}%
\end{pgfscope}%
\begin{pgfscope}%
\pgfpathrectangle{\pgfqpoint{0.100000in}{0.100000in}}{\pgfqpoint{3.007045in}{1.925000in}}%
\pgfusepath{clip}%
\pgfsetbuttcap%
\pgfsetmiterjoin%
\definecolor{currentfill}{rgb}{0.562860,0.819008,0.551511}%
\pgfsetfillcolor{currentfill}%
\pgfsetlinewidth{0.000000pt}%
\definecolor{currentstroke}{rgb}{0.000000,0.000000,0.000000}%
\pgfsetstrokecolor{currentstroke}%
\pgfsetstrokeopacity{0.000000}%
\pgfsetdash{}{0pt}%
\pgfpathmoveto{\pgfqpoint{2.407666in}{1.313899in}}%
\pgfpathlineto{\pgfqpoint{2.410399in}{1.309530in}}%
\pgfpathlineto{\pgfqpoint{2.413572in}{1.288426in}}%
\pgfpathlineto{\pgfqpoint{2.408904in}{1.287720in}}%
\pgfpathlineto{\pgfqpoint{2.409849in}{1.281241in}}%
\pgfpathlineto{\pgfqpoint{2.402337in}{1.279030in}}%
\pgfpathlineto{\pgfqpoint{2.389048in}{1.277134in}}%
\pgfpathlineto{\pgfqpoint{2.390704in}{1.263733in}}%
\pgfpathlineto{\pgfqpoint{2.376485in}{1.262685in}}%
\pgfpathlineto{\pgfqpoint{2.376250in}{1.266421in}}%
\pgfpathlineto{\pgfqpoint{2.359900in}{1.266193in}}%
\pgfpathlineto{\pgfqpoint{2.359392in}{1.270321in}}%
\pgfpathlineto{\pgfqpoint{2.350943in}{1.269110in}}%
\pgfpathlineto{\pgfqpoint{2.349350in}{1.280970in}}%
\pgfpathlineto{\pgfqpoint{2.364496in}{1.284008in}}%
\pgfpathlineto{\pgfqpoint{2.361935in}{1.302986in}}%
\pgfpathlineto{\pgfqpoint{2.375866in}{1.305055in}}%
\pgfpathlineto{\pgfqpoint{2.395446in}{1.307499in}}%
\pgfpathlineto{\pgfqpoint{2.394589in}{1.312199in}}%
\pgfpathlineto{\pgfqpoint{2.407666in}{1.313899in}}%
\pgfpathclose%
\pgfusepath{fill}%
\end{pgfscope}%
\begin{pgfscope}%
\pgfpathrectangle{\pgfqpoint{0.100000in}{0.100000in}}{\pgfqpoint{3.007045in}{1.925000in}}%
\pgfusepath{clip}%
\pgfsetbuttcap%
\pgfsetmiterjoin%
\definecolor{currentfill}{rgb}{0.234141,0.648443,0.348097}%
\pgfsetfillcolor{currentfill}%
\pgfsetlinewidth{0.000000pt}%
\definecolor{currentstroke}{rgb}{0.000000,0.000000,0.000000}%
\pgfsetstrokecolor{currentstroke}%
\pgfsetstrokeopacity{0.000000}%
\pgfsetdash{}{0pt}%
\pgfpathmoveto{\pgfqpoint{2.285385in}{0.933858in}}%
\pgfpathlineto{\pgfqpoint{2.292762in}{0.944860in}}%
\pgfpathlineto{\pgfqpoint{2.293939in}{0.955321in}}%
\pgfpathlineto{\pgfqpoint{2.291000in}{0.961775in}}%
\pgfpathlineto{\pgfqpoint{2.291608in}{0.970240in}}%
\pgfpathlineto{\pgfqpoint{2.296056in}{0.970309in}}%
\pgfpathlineto{\pgfqpoint{2.299083in}{0.977051in}}%
\pgfpathlineto{\pgfqpoint{2.296204in}{0.983928in}}%
\pgfpathlineto{\pgfqpoint{2.295566in}{0.990662in}}%
\pgfpathlineto{\pgfqpoint{2.298509in}{1.001042in}}%
\pgfpathlineto{\pgfqpoint{2.303650in}{1.003101in}}%
\pgfpathlineto{\pgfqpoint{2.315188in}{1.000641in}}%
\pgfpathlineto{\pgfqpoint{2.319984in}{0.992600in}}%
\pgfpathlineto{\pgfqpoint{2.317969in}{0.990753in}}%
\pgfpathlineto{\pgfqpoint{2.309261in}{0.983763in}}%
\pgfpathlineto{\pgfqpoint{2.308631in}{0.975465in}}%
\pgfpathlineto{\pgfqpoint{2.319716in}{0.965996in}}%
\pgfpathlineto{\pgfqpoint{2.322488in}{0.960571in}}%
\pgfpathlineto{\pgfqpoint{2.317579in}{0.954327in}}%
\pgfpathlineto{\pgfqpoint{2.321351in}{0.950113in}}%
\pgfpathlineto{\pgfqpoint{2.318851in}{0.944884in}}%
\pgfpathlineto{\pgfqpoint{2.319657in}{0.939932in}}%
\pgfpathlineto{\pgfqpoint{2.313025in}{0.928226in}}%
\pgfpathlineto{\pgfqpoint{2.308274in}{0.925835in}}%
\pgfpathlineto{\pgfqpoint{2.299822in}{0.927995in}}%
\pgfpathlineto{\pgfqpoint{2.297460in}{0.921002in}}%
\pgfpathlineto{\pgfqpoint{2.287005in}{0.929042in}}%
\pgfpathlineto{\pgfqpoint{2.285385in}{0.933858in}}%
\pgfpathclose%
\pgfusepath{fill}%
\end{pgfscope}%
\begin{pgfscope}%
\pgfpathrectangle{\pgfqpoint{0.100000in}{0.100000in}}{\pgfqpoint{3.007045in}{1.925000in}}%
\pgfusepath{clip}%
\pgfsetbuttcap%
\pgfsetmiterjoin%
\definecolor{currentfill}{rgb}{0.357647,0.720953,0.415071}%
\pgfsetfillcolor{currentfill}%
\pgfsetlinewidth{0.000000pt}%
\definecolor{currentstroke}{rgb}{0.000000,0.000000,0.000000}%
\pgfsetstrokecolor{currentstroke}%
\pgfsetstrokeopacity{0.000000}%
\pgfsetdash{}{0pt}%
\pgfpathmoveto{\pgfqpoint{2.742724in}{1.035580in}}%
\pgfpathlineto{\pgfqpoint{2.743345in}{1.029859in}}%
\pgfpathlineto{\pgfqpoint{2.741385in}{1.024328in}}%
\pgfpathlineto{\pgfqpoint{2.738228in}{1.024682in}}%
\pgfpathlineto{\pgfqpoint{2.735129in}{1.019199in}}%
\pgfpathlineto{\pgfqpoint{2.731963in}{1.024371in}}%
\pgfpathlineto{\pgfqpoint{2.728149in}{1.020939in}}%
\pgfpathlineto{\pgfqpoint{2.722671in}{1.024298in}}%
\pgfpathlineto{\pgfqpoint{2.716657in}{1.024638in}}%
\pgfpathlineto{\pgfqpoint{2.715659in}{1.028449in}}%
\pgfpathlineto{\pgfqpoint{2.709028in}{1.030469in}}%
\pgfpathlineto{\pgfqpoint{2.705428in}{1.026285in}}%
\pgfpathlineto{\pgfqpoint{2.698517in}{1.027123in}}%
\pgfpathlineto{\pgfqpoint{2.697499in}{1.030502in}}%
\pgfpathlineto{\pgfqpoint{2.689119in}{1.033149in}}%
\pgfpathlineto{\pgfqpoint{2.679615in}{1.030971in}}%
\pgfpathlineto{\pgfqpoint{2.671365in}{1.033005in}}%
\pgfpathlineto{\pgfqpoint{2.674093in}{1.045577in}}%
\pgfpathlineto{\pgfqpoint{2.672729in}{1.056217in}}%
\pgfpathlineto{\pgfqpoint{2.679621in}{1.057585in}}%
\pgfpathlineto{\pgfqpoint{2.683306in}{1.069582in}}%
\pgfpathlineto{\pgfqpoint{2.680930in}{1.081273in}}%
\pgfpathlineto{\pgfqpoint{2.685692in}{1.079312in}}%
\pgfpathlineto{\pgfqpoint{2.691645in}{1.081236in}}%
\pgfpathlineto{\pgfqpoint{2.691815in}{1.071736in}}%
\pgfpathlineto{\pgfqpoint{2.695089in}{1.071478in}}%
\pgfpathlineto{\pgfqpoint{2.701345in}{1.062144in}}%
\pgfpathlineto{\pgfqpoint{2.723871in}{1.066185in}}%
\pgfpathlineto{\pgfqpoint{2.724911in}{1.056794in}}%
\pgfpathlineto{\pgfqpoint{2.733124in}{1.055886in}}%
\pgfpathlineto{\pgfqpoint{2.738529in}{1.052500in}}%
\pgfpathlineto{\pgfqpoint{2.738199in}{1.043269in}}%
\pgfpathlineto{\pgfqpoint{2.742724in}{1.035580in}}%
\pgfpathclose%
\pgfusepath{fill}%
\end{pgfscope}%
\begin{pgfscope}%
\pgfpathrectangle{\pgfqpoint{0.100000in}{0.100000in}}{\pgfqpoint{3.007045in}{1.925000in}}%
\pgfusepath{clip}%
\pgfsetbuttcap%
\pgfsetmiterjoin%
\definecolor{currentfill}{rgb}{0.638970,0.854179,0.615240}%
\pgfsetfillcolor{currentfill}%
\pgfsetlinewidth{0.000000pt}%
\definecolor{currentstroke}{rgb}{0.000000,0.000000,0.000000}%
\pgfsetstrokecolor{currentstroke}%
\pgfsetstrokeopacity{0.000000}%
\pgfsetdash{}{0pt}%
\pgfpathmoveto{\pgfqpoint{1.538682in}{0.941826in}}%
\pgfpathlineto{\pgfqpoint{1.532268in}{0.939536in}}%
\pgfpathlineto{\pgfqpoint{1.528739in}{0.932611in}}%
\pgfpathlineto{\pgfqpoint{1.520249in}{0.932634in}}%
\pgfpathlineto{\pgfqpoint{1.515519in}{0.937448in}}%
\pgfpathlineto{\pgfqpoint{1.510237in}{0.939913in}}%
\pgfpathlineto{\pgfqpoint{1.505202in}{0.934838in}}%
\pgfpathlineto{\pgfqpoint{1.505836in}{0.946347in}}%
\pgfpathlineto{\pgfqpoint{1.507272in}{0.975589in}}%
\pgfpathlineto{\pgfqpoint{1.509153in}{1.008633in}}%
\pgfpathlineto{\pgfqpoint{1.533400in}{1.007220in}}%
\pgfpathlineto{\pgfqpoint{1.537864in}{1.006990in}}%
\pgfpathlineto{\pgfqpoint{1.541765in}{0.998932in}}%
\pgfpathlineto{\pgfqpoint{1.547875in}{0.992381in}}%
\pgfpathlineto{\pgfqpoint{1.553847in}{0.991951in}}%
\pgfpathlineto{\pgfqpoint{1.562805in}{0.979407in}}%
\pgfpathlineto{\pgfqpoint{1.561768in}{0.950581in}}%
\pgfpathlineto{\pgfqpoint{1.539102in}{0.951802in}}%
\pgfpathlineto{\pgfqpoint{1.538682in}{0.941826in}}%
\pgfpathclose%
\pgfusepath{fill}%
\end{pgfscope}%
\begin{pgfscope}%
\pgfpathrectangle{\pgfqpoint{0.100000in}{0.100000in}}{\pgfqpoint{3.007045in}{1.925000in}}%
\pgfusepath{clip}%
\pgfsetbuttcap%
\pgfsetmiterjoin%
\definecolor{currentfill}{rgb}{0.812687,0.926644,0.787389}%
\pgfsetfillcolor{currentfill}%
\pgfsetlinewidth{0.000000pt}%
\definecolor{currentstroke}{rgb}{0.000000,0.000000,0.000000}%
\pgfsetstrokecolor{currentstroke}%
\pgfsetstrokeopacity{0.000000}%
\pgfsetdash{}{0pt}%
\pgfpathmoveto{\pgfqpoint{2.585937in}{0.924191in}}%
\pgfpathlineto{\pgfqpoint{2.551891in}{0.919345in}}%
\pgfpathlineto{\pgfqpoint{2.562238in}{0.907450in}}%
\pgfpathlineto{\pgfqpoint{2.555178in}{0.902655in}}%
\pgfpathlineto{\pgfqpoint{2.555983in}{0.898307in}}%
\pgfpathlineto{\pgfqpoint{2.550816in}{0.897393in}}%
\pgfpathlineto{\pgfqpoint{2.545027in}{0.899687in}}%
\pgfpathlineto{\pgfqpoint{2.538562in}{0.893175in}}%
\pgfpathlineto{\pgfqpoint{2.537783in}{0.898762in}}%
\pgfpathlineto{\pgfqpoint{2.508172in}{0.896146in}}%
\pgfpathlineto{\pgfqpoint{2.505466in}{0.902631in}}%
\pgfpathlineto{\pgfqpoint{2.503825in}{0.913860in}}%
\pgfpathlineto{\pgfqpoint{2.500183in}{0.925858in}}%
\pgfpathlineto{\pgfqpoint{2.504209in}{0.927233in}}%
\pgfpathlineto{\pgfqpoint{2.505332in}{0.935315in}}%
\pgfpathlineto{\pgfqpoint{2.523780in}{0.937038in}}%
\pgfpathlineto{\pgfqpoint{2.523035in}{0.944991in}}%
\pgfpathlineto{\pgfqpoint{2.524892in}{0.952370in}}%
\pgfpathlineto{\pgfqpoint{2.523665in}{0.961730in}}%
\pgfpathlineto{\pgfqpoint{2.535380in}{0.962820in}}%
\pgfpathlineto{\pgfqpoint{2.531826in}{0.973469in}}%
\pgfpathlineto{\pgfqpoint{2.533363in}{0.985590in}}%
\pgfpathlineto{\pgfqpoint{2.537459in}{0.985766in}}%
\pgfpathlineto{\pgfqpoint{2.547659in}{0.980532in}}%
\pgfpathlineto{\pgfqpoint{2.555149in}{0.977464in}}%
\pgfpathlineto{\pgfqpoint{2.562648in}{0.972115in}}%
\pgfpathlineto{\pgfqpoint{2.564813in}{0.967383in}}%
\pgfpathlineto{\pgfqpoint{2.569136in}{0.965197in}}%
\pgfpathlineto{\pgfqpoint{2.573216in}{0.959889in}}%
\pgfpathlineto{\pgfqpoint{2.571933in}{0.952085in}}%
\pgfpathlineto{\pgfqpoint{2.575191in}{0.947306in}}%
\pgfpathlineto{\pgfqpoint{2.575141in}{0.942834in}}%
\pgfpathlineto{\pgfqpoint{2.581758in}{0.944406in}}%
\pgfpathlineto{\pgfqpoint{2.586120in}{0.939582in}}%
\pgfpathlineto{\pgfqpoint{2.589091in}{0.931283in}}%
\pgfpathlineto{\pgfqpoint{2.585937in}{0.924191in}}%
\pgfpathclose%
\pgfusepath{fill}%
\end{pgfscope}%
\begin{pgfscope}%
\pgfpathrectangle{\pgfqpoint{0.100000in}{0.100000in}}{\pgfqpoint{3.007045in}{1.925000in}}%
\pgfusepath{clip}%
\pgfsetbuttcap%
\pgfsetmiterjoin%
\definecolor{currentfill}{rgb}{0.685721,0.873864,0.660761}%
\pgfsetfillcolor{currentfill}%
\pgfsetlinewidth{0.000000pt}%
\definecolor{currentstroke}{rgb}{0.000000,0.000000,0.000000}%
\pgfsetstrokecolor{currentstroke}%
\pgfsetstrokeopacity{0.000000}%
\pgfsetdash{}{0pt}%
\pgfpathmoveto{\pgfqpoint{1.712393in}{0.726751in}}%
\pgfpathlineto{\pgfqpoint{1.700123in}{0.727126in}}%
\pgfpathlineto{\pgfqpoint{1.700709in}{0.760907in}}%
\pgfpathlineto{\pgfqpoint{1.724887in}{0.764449in}}%
\pgfpathlineto{\pgfqpoint{1.725337in}{0.792070in}}%
\pgfpathlineto{\pgfqpoint{1.730486in}{0.795042in}}%
\pgfpathlineto{\pgfqpoint{1.740181in}{0.799040in}}%
\pgfpathlineto{\pgfqpoint{1.742020in}{0.795983in}}%
\pgfpathlineto{\pgfqpoint{1.749934in}{0.794479in}}%
\pgfpathlineto{\pgfqpoint{1.756132in}{0.794865in}}%
\pgfpathlineto{\pgfqpoint{1.760297in}{0.800742in}}%
\pgfpathlineto{\pgfqpoint{1.774881in}{0.791133in}}%
\pgfpathlineto{\pgfqpoint{1.779488in}{0.786480in}}%
\pgfpathlineto{\pgfqpoint{1.786299in}{0.783647in}}%
\pgfpathlineto{\pgfqpoint{1.786323in}{0.758894in}}%
\pgfpathlineto{\pgfqpoint{1.779169in}{0.763215in}}%
\pgfpathlineto{\pgfqpoint{1.767169in}{0.763420in}}%
\pgfpathlineto{\pgfqpoint{1.755274in}{0.762070in}}%
\pgfpathlineto{\pgfqpoint{1.755158in}{0.734645in}}%
\pgfpathlineto{\pgfqpoint{1.735350in}{0.734610in}}%
\pgfpathlineto{\pgfqpoint{1.736915in}{0.718690in}}%
\pgfpathlineto{\pgfqpoint{1.728058in}{0.722452in}}%
\pgfpathlineto{\pgfqpoint{1.725321in}{0.722021in}}%
\pgfpathlineto{\pgfqpoint{1.720316in}{0.726579in}}%
\pgfpathlineto{\pgfqpoint{1.712393in}{0.726751in}}%
\pgfpathclose%
\pgfusepath{fill}%
\end{pgfscope}%
\begin{pgfscope}%
\pgfpathrectangle{\pgfqpoint{0.100000in}{0.100000in}}{\pgfqpoint{3.007045in}{1.925000in}}%
\pgfusepath{clip}%
\pgfsetbuttcap%
\pgfsetmiterjoin%
\definecolor{currentfill}{rgb}{0.439216,0.760938,0.455056}%
\pgfsetfillcolor{currentfill}%
\pgfsetlinewidth{0.000000pt}%
\definecolor{currentstroke}{rgb}{0.000000,0.000000,0.000000}%
\pgfsetstrokecolor{currentstroke}%
\pgfsetstrokeopacity{0.000000}%
\pgfsetdash{}{0pt}%
\pgfpathmoveto{\pgfqpoint{2.484306in}{1.188484in}}%
\pgfpathlineto{\pgfqpoint{2.487519in}{1.178106in}}%
\pgfpathlineto{\pgfqpoint{2.485316in}{1.174226in}}%
\pgfpathlineto{\pgfqpoint{2.489719in}{1.169881in}}%
\pgfpathlineto{\pgfqpoint{2.487851in}{1.166255in}}%
\pgfpathlineto{\pgfqpoint{2.480993in}{1.155854in}}%
\pgfpathlineto{\pgfqpoint{2.466955in}{1.155210in}}%
\pgfpathlineto{\pgfqpoint{2.463450in}{1.158871in}}%
\pgfpathlineto{\pgfqpoint{2.461626in}{1.162571in}}%
\pgfpathlineto{\pgfqpoint{2.460301in}{1.177194in}}%
\pgfpathlineto{\pgfqpoint{2.463999in}{1.180859in}}%
\pgfpathlineto{\pgfqpoint{2.469733in}{1.178659in}}%
\pgfpathlineto{\pgfqpoint{2.478749in}{1.189995in}}%
\pgfpathlineto{\pgfqpoint{2.484306in}{1.188484in}}%
\pgfpathclose%
\pgfusepath{fill}%
\end{pgfscope}%
\begin{pgfscope}%
\pgfpathrectangle{\pgfqpoint{0.100000in}{0.100000in}}{\pgfqpoint{3.007045in}{1.925000in}}%
\pgfusepath{clip}%
\pgfsetbuttcap%
\pgfsetmiterjoin%
\definecolor{currentfill}{rgb}{0.790542,0.917785,0.763768}%
\pgfsetfillcolor{currentfill}%
\pgfsetlinewidth{0.000000pt}%
\definecolor{currentstroke}{rgb}{0.000000,0.000000,0.000000}%
\pgfsetstrokecolor{currentstroke}%
\pgfsetstrokeopacity{0.000000}%
\pgfsetdash{}{0pt}%
\pgfpathmoveto{\pgfqpoint{1.369270in}{1.748112in}}%
\pgfpathlineto{\pgfqpoint{1.371755in}{1.774248in}}%
\pgfpathlineto{\pgfqpoint{1.375731in}{1.814864in}}%
\pgfpathlineto{\pgfqpoint{1.424243in}{1.810279in}}%
\pgfpathlineto{\pgfqpoint{1.464391in}{1.806947in}}%
\pgfpathlineto{\pgfqpoint{1.463375in}{1.794138in}}%
\pgfpathlineto{\pgfqpoint{1.457650in}{1.794593in}}%
\pgfpathlineto{\pgfqpoint{1.457171in}{1.788804in}}%
\pgfpathlineto{\pgfqpoint{1.453609in}{1.789081in}}%
\pgfpathlineto{\pgfqpoint{1.452658in}{1.777477in}}%
\pgfpathlineto{\pgfqpoint{1.423931in}{1.779899in}}%
\pgfpathlineto{\pgfqpoint{1.422948in}{1.768339in}}%
\pgfpathlineto{\pgfqpoint{1.425483in}{1.768138in}}%
\pgfpathlineto{\pgfqpoint{1.422411in}{1.752597in}}%
\pgfpathlineto{\pgfqpoint{1.407533in}{1.753794in}}%
\pgfpathlineto{\pgfqpoint{1.403722in}{1.748968in}}%
\pgfpathlineto{\pgfqpoint{1.394175in}{1.746414in}}%
\pgfpathlineto{\pgfqpoint{1.390611in}{1.748693in}}%
\pgfpathlineto{\pgfqpoint{1.389403in}{1.754089in}}%
\pgfpathlineto{\pgfqpoint{1.386114in}{1.754885in}}%
\pgfpathlineto{\pgfqpoint{1.376795in}{1.745675in}}%
\pgfpathlineto{\pgfqpoint{1.369270in}{1.748112in}}%
\pgfpathclose%
\pgfusepath{fill}%
\end{pgfscope}%
\begin{pgfscope}%
\pgfpathrectangle{\pgfqpoint{0.100000in}{0.100000in}}{\pgfqpoint{3.007045in}{1.925000in}}%
\pgfusepath{clip}%
\pgfsetbuttcap%
\pgfsetmiterjoin%
\definecolor{currentfill}{rgb}{0.834833,0.935502,0.811011}%
\pgfsetfillcolor{currentfill}%
\pgfsetlinewidth{0.000000pt}%
\definecolor{currentstroke}{rgb}{0.000000,0.000000,0.000000}%
\pgfsetstrokecolor{currentstroke}%
\pgfsetstrokeopacity{0.000000}%
\pgfsetdash{}{0pt}%
\pgfpathmoveto{\pgfqpoint{1.330767in}{0.811238in}}%
\pgfpathlineto{\pgfqpoint{1.383417in}{0.806844in}}%
\pgfpathlineto{\pgfqpoint{1.440696in}{0.803294in}}%
\pgfpathlineto{\pgfqpoint{1.438908in}{0.774476in}}%
\pgfpathlineto{\pgfqpoint{1.437038in}{0.746267in}}%
\pgfpathlineto{\pgfqpoint{1.429539in}{0.746331in}}%
\pgfpathlineto{\pgfqpoint{1.379483in}{0.749728in}}%
\pgfpathlineto{\pgfqpoint{1.381713in}{0.778104in}}%
\pgfpathlineto{\pgfqpoint{1.327701in}{0.782533in}}%
\pgfpathlineto{\pgfqpoint{1.328991in}{0.794521in}}%
\pgfpathlineto{\pgfqpoint{1.330767in}{0.811238in}}%
\pgfpathclose%
\pgfusepath{fill}%
\end{pgfscope}%
\begin{pgfscope}%
\pgfpathrectangle{\pgfqpoint{0.100000in}{0.100000in}}{\pgfqpoint{3.007045in}{1.925000in}}%
\pgfusepath{clip}%
\pgfsetbuttcap%
\pgfsetmiterjoin%
\definecolor{currentfill}{rgb}{0.709097,0.883706,0.683522}%
\pgfsetfillcolor{currentfill}%
\pgfsetlinewidth{0.000000pt}%
\definecolor{currentstroke}{rgb}{0.000000,0.000000,0.000000}%
\pgfsetstrokecolor{currentstroke}%
\pgfsetstrokeopacity{0.000000}%
\pgfsetdash{}{0pt}%
\pgfpathmoveto{\pgfqpoint{2.637551in}{0.899069in}}%
\pgfpathlineto{\pgfqpoint{2.613555in}{0.916458in}}%
\pgfpathlineto{\pgfqpoint{2.599449in}{0.926351in}}%
\pgfpathlineto{\pgfqpoint{2.585937in}{0.924191in}}%
\pgfpathlineto{\pgfqpoint{2.589091in}{0.931283in}}%
\pgfpathlineto{\pgfqpoint{2.586120in}{0.939582in}}%
\pgfpathlineto{\pgfqpoint{2.581758in}{0.944406in}}%
\pgfpathlineto{\pgfqpoint{2.575141in}{0.942834in}}%
\pgfpathlineto{\pgfqpoint{2.575191in}{0.947306in}}%
\pgfpathlineto{\pgfqpoint{2.571933in}{0.952085in}}%
\pgfpathlineto{\pgfqpoint{2.573216in}{0.959889in}}%
\pgfpathlineto{\pgfqpoint{2.569136in}{0.965197in}}%
\pgfpathlineto{\pgfqpoint{2.564813in}{0.967383in}}%
\pgfpathlineto{\pgfqpoint{2.597898in}{0.973667in}}%
\pgfpathlineto{\pgfqpoint{2.608759in}{0.975733in}}%
\pgfpathlineto{\pgfqpoint{2.608908in}{0.972810in}}%
\pgfpathlineto{\pgfqpoint{2.618329in}{0.960663in}}%
\pgfpathlineto{\pgfqpoint{2.625940in}{0.956912in}}%
\pgfpathlineto{\pgfqpoint{2.641093in}{0.964571in}}%
\pgfpathlineto{\pgfqpoint{2.650737in}{0.964955in}}%
\pgfpathlineto{\pgfqpoint{2.650566in}{0.955387in}}%
\pgfpathlineto{\pgfqpoint{2.652361in}{0.948046in}}%
\pgfpathlineto{\pgfqpoint{2.661874in}{0.940850in}}%
\pgfpathlineto{\pgfqpoint{2.645656in}{0.937482in}}%
\pgfpathlineto{\pgfqpoint{2.640458in}{0.935526in}}%
\pgfpathlineto{\pgfqpoint{2.647207in}{0.927059in}}%
\pgfpathlineto{\pgfqpoint{2.647369in}{0.914791in}}%
\pgfpathlineto{\pgfqpoint{2.641884in}{0.909457in}}%
\pgfpathlineto{\pgfqpoint{2.637551in}{0.899069in}}%
\pgfpathclose%
\pgfusepath{fill}%
\end{pgfscope}%
\begin{pgfscope}%
\pgfpathrectangle{\pgfqpoint{0.100000in}{0.100000in}}{\pgfqpoint{3.007045in}{1.925000in}}%
\pgfusepath{clip}%
\pgfsetbuttcap%
\pgfsetmiterjoin%
\definecolor{currentfill}{rgb}{0.518570,0.798339,0.515094}%
\pgfsetfillcolor{currentfill}%
\pgfsetlinewidth{0.000000pt}%
\definecolor{currentstroke}{rgb}{0.000000,0.000000,0.000000}%
\pgfsetstrokecolor{currentstroke}%
\pgfsetstrokeopacity{0.000000}%
\pgfsetdash{}{0pt}%
\pgfpathmoveto{\pgfqpoint{1.653431in}{1.223538in}}%
\pgfpathlineto{\pgfqpoint{1.652796in}{1.200602in}}%
\pgfpathlineto{\pgfqpoint{1.607157in}{1.202049in}}%
\pgfpathlineto{\pgfqpoint{1.608008in}{1.224954in}}%
\pgfpathlineto{\pgfqpoint{1.630558in}{1.224197in}}%
\pgfpathlineto{\pgfqpoint{1.653431in}{1.223538in}}%
\pgfpathclose%
\pgfusepath{fill}%
\end{pgfscope}%
\begin{pgfscope}%
\pgfpathrectangle{\pgfqpoint{0.100000in}{0.100000in}}{\pgfqpoint{3.007045in}{1.925000in}}%
\pgfusepath{clip}%
\pgfsetbuttcap%
\pgfsetmiterjoin%
\definecolor{currentfill}{rgb}{0.585006,0.829343,0.569719}%
\pgfsetfillcolor{currentfill}%
\pgfsetlinewidth{0.000000pt}%
\definecolor{currentstroke}{rgb}{0.000000,0.000000,0.000000}%
\pgfsetstrokecolor{currentstroke}%
\pgfsetstrokeopacity{0.000000}%
\pgfsetdash{}{0pt}%
\pgfpathmoveto{\pgfqpoint{2.108338in}{0.951144in}}%
\pgfpathlineto{\pgfqpoint{2.098903in}{0.955455in}}%
\pgfpathlineto{\pgfqpoint{2.093449in}{0.959312in}}%
\pgfpathlineto{\pgfqpoint{2.083172in}{0.958841in}}%
\pgfpathlineto{\pgfqpoint{2.065703in}{0.958280in}}%
\pgfpathlineto{\pgfqpoint{2.068381in}{0.961543in}}%
\pgfpathlineto{\pgfqpoint{2.071414in}{0.971616in}}%
\pgfpathlineto{\pgfqpoint{2.068022in}{0.977336in}}%
\pgfpathlineto{\pgfqpoint{2.070320in}{0.985858in}}%
\pgfpathlineto{\pgfqpoint{2.077777in}{0.982593in}}%
\pgfpathlineto{\pgfqpoint{2.080971in}{0.989204in}}%
\pgfpathlineto{\pgfqpoint{2.078631in}{0.992336in}}%
\pgfpathlineto{\pgfqpoint{2.082341in}{0.997150in}}%
\pgfpathlineto{\pgfqpoint{2.098746in}{0.997437in}}%
\pgfpathlineto{\pgfqpoint{2.097982in}{1.008842in}}%
\pgfpathlineto{\pgfqpoint{2.115261in}{1.009765in}}%
\pgfpathlineto{\pgfqpoint{2.116101in}{0.997143in}}%
\pgfpathlineto{\pgfqpoint{2.117091in}{0.980740in}}%
\pgfpathlineto{\pgfqpoint{2.115651in}{0.980661in}}%
\pgfpathlineto{\pgfqpoint{2.116626in}{0.955211in}}%
\pgfpathlineto{\pgfqpoint{2.110066in}{0.954874in}}%
\pgfpathlineto{\pgfqpoint{2.108338in}{0.951144in}}%
\pgfpathclose%
\pgfusepath{fill}%
\end{pgfscope}%
\begin{pgfscope}%
\pgfpathrectangle{\pgfqpoint{0.100000in}{0.100000in}}{\pgfqpoint{3.007045in}{1.925000in}}%
\pgfusepath{clip}%
\pgfsetbuttcap%
\pgfsetmiterjoin%
\definecolor{currentfill}{rgb}{0.727797,0.891580,0.701730}%
\pgfsetfillcolor{currentfill}%
\pgfsetlinewidth{0.000000pt}%
\definecolor{currentstroke}{rgb}{0.000000,0.000000,0.000000}%
\pgfsetstrokecolor{currentstroke}%
\pgfsetstrokeopacity{0.000000}%
\pgfsetdash{}{0pt}%
\pgfpathmoveto{\pgfqpoint{1.184259in}{1.303392in}}%
\pgfpathlineto{\pgfqpoint{1.210701in}{1.299766in}}%
\pgfpathlineto{\pgfqpoint{1.213860in}{1.325677in}}%
\pgfpathlineto{\pgfqpoint{1.226347in}{1.324123in}}%
\pgfpathlineto{\pgfqpoint{1.231571in}{1.369429in}}%
\pgfpathlineto{\pgfqpoint{1.234424in}{1.392114in}}%
\pgfpathlineto{\pgfqpoint{1.259419in}{1.388949in}}%
\pgfpathlineto{\pgfqpoint{1.256175in}{1.385798in}}%
\pgfpathlineto{\pgfqpoint{1.255653in}{1.381012in}}%
\pgfpathlineto{\pgfqpoint{1.267095in}{1.378734in}}%
\pgfpathlineto{\pgfqpoint{1.270231in}{1.387736in}}%
\pgfpathlineto{\pgfqpoint{1.272590in}{1.387487in}}%
\pgfpathlineto{\pgfqpoint{1.269866in}{1.358232in}}%
\pgfpathlineto{\pgfqpoint{1.267098in}{1.336682in}}%
\pgfpathlineto{\pgfqpoint{1.262254in}{1.293547in}}%
\pgfpathlineto{\pgfqpoint{1.217135in}{1.298862in}}%
\pgfpathlineto{\pgfqpoint{1.216879in}{1.294670in}}%
\pgfpathlineto{\pgfqpoint{1.222250in}{1.286251in}}%
\pgfpathlineto{\pgfqpoint{1.229691in}{1.263388in}}%
\pgfpathlineto{\pgfqpoint{1.222855in}{1.255180in}}%
\pgfpathlineto{\pgfqpoint{1.220648in}{1.256986in}}%
\pgfpathlineto{\pgfqpoint{1.213800in}{1.255082in}}%
\pgfpathlineto{\pgfqpoint{1.205670in}{1.255974in}}%
\pgfpathlineto{\pgfqpoint{1.192494in}{1.261100in}}%
\pgfpathlineto{\pgfqpoint{1.189706in}{1.265642in}}%
\pgfpathlineto{\pgfqpoint{1.188332in}{1.274642in}}%
\pgfpathlineto{\pgfqpoint{1.193426in}{1.288196in}}%
\pgfpathlineto{\pgfqpoint{1.192988in}{1.292052in}}%
\pgfpathlineto{\pgfqpoint{1.183883in}{1.298271in}}%
\pgfpathlineto{\pgfqpoint{1.184259in}{1.303392in}}%
\pgfpathclose%
\pgfusepath{fill}%
\end{pgfscope}%
\begin{pgfscope}%
\pgfpathrectangle{\pgfqpoint{0.100000in}{0.100000in}}{\pgfqpoint{3.007045in}{1.925000in}}%
\pgfusepath{clip}%
\pgfsetbuttcap%
\pgfsetmiterjoin%
\definecolor{currentfill}{rgb}{0.426667,0.754787,0.448904}%
\pgfsetfillcolor{currentfill}%
\pgfsetlinewidth{0.000000pt}%
\definecolor{currentstroke}{rgb}{0.000000,0.000000,0.000000}%
\pgfsetstrokecolor{currentstroke}%
\pgfsetstrokeopacity{0.000000}%
\pgfsetdash{}{0pt}%
\pgfpathmoveto{\pgfqpoint{2.362681in}{1.482619in}}%
\pgfpathlineto{\pgfqpoint{2.338611in}{1.478397in}}%
\pgfpathlineto{\pgfqpoint{2.342275in}{1.455683in}}%
\pgfpathlineto{\pgfqpoint{2.330370in}{1.453796in}}%
\pgfpathlineto{\pgfqpoint{2.331232in}{1.448101in}}%
\pgfpathlineto{\pgfqpoint{2.326140in}{1.446645in}}%
\pgfpathlineto{\pgfqpoint{2.303651in}{1.443639in}}%
\pgfpathlineto{\pgfqpoint{2.304504in}{1.437926in}}%
\pgfpathlineto{\pgfqpoint{2.283546in}{1.435118in}}%
\pgfpathlineto{\pgfqpoint{2.280842in}{1.457266in}}%
\pgfpathlineto{\pgfqpoint{2.269523in}{1.455942in}}%
\pgfpathlineto{\pgfqpoint{2.264263in}{1.501550in}}%
\pgfpathlineto{\pgfqpoint{2.298344in}{1.505689in}}%
\pgfpathlineto{\pgfqpoint{2.313368in}{1.507664in}}%
\pgfpathlineto{\pgfqpoint{2.313689in}{1.503876in}}%
\pgfpathlineto{\pgfqpoint{2.305190in}{1.494752in}}%
\pgfpathlineto{\pgfqpoint{2.302323in}{1.494776in}}%
\pgfpathlineto{\pgfqpoint{2.298978in}{1.488134in}}%
\pgfpathlineto{\pgfqpoint{2.298537in}{1.477920in}}%
\pgfpathlineto{\pgfqpoint{2.300991in}{1.473532in}}%
\pgfpathlineto{\pgfqpoint{2.313383in}{1.469662in}}%
\pgfpathlineto{\pgfqpoint{2.329237in}{1.497116in}}%
\pgfpathlineto{\pgfqpoint{2.339139in}{1.501127in}}%
\pgfpathlineto{\pgfqpoint{2.344111in}{1.505414in}}%
\pgfpathlineto{\pgfqpoint{2.353364in}{1.501784in}}%
\pgfpathlineto{\pgfqpoint{2.360186in}{1.491680in}}%
\pgfpathlineto{\pgfqpoint{2.362681in}{1.482619in}}%
\pgfpathclose%
\pgfusepath{fill}%
\end{pgfscope}%
\begin{pgfscope}%
\pgfpathrectangle{\pgfqpoint{0.100000in}{0.100000in}}{\pgfqpoint{3.007045in}{1.925000in}}%
\pgfusepath{clip}%
\pgfsetbuttcap%
\pgfsetmiterjoin%
\definecolor{currentfill}{rgb}{0.426667,0.754787,0.448904}%
\pgfsetfillcolor{currentfill}%
\pgfsetlinewidth{0.000000pt}%
\definecolor{currentstroke}{rgb}{0.000000,0.000000,0.000000}%
\pgfsetstrokecolor{currentstroke}%
\pgfsetstrokeopacity{0.000000}%
\pgfsetdash{}{0pt}%
\pgfpathmoveto{\pgfqpoint{2.558875in}{1.172717in}}%
\pgfpathlineto{\pgfqpoint{2.558583in}{1.161971in}}%
\pgfpathlineto{\pgfqpoint{2.553841in}{1.152723in}}%
\pgfpathlineto{\pgfqpoint{2.555307in}{1.150523in}}%
\pgfpathlineto{\pgfqpoint{2.549531in}{1.143698in}}%
\pgfpathlineto{\pgfqpoint{2.548827in}{1.135761in}}%
\pgfpathlineto{\pgfqpoint{2.538421in}{1.132193in}}%
\pgfpathlineto{\pgfqpoint{2.533411in}{1.132073in}}%
\pgfpathlineto{\pgfqpoint{2.527646in}{1.135652in}}%
\pgfpathlineto{\pgfqpoint{2.526634in}{1.142930in}}%
\pgfpathlineto{\pgfqpoint{2.526995in}{1.150325in}}%
\pgfpathlineto{\pgfqpoint{2.530752in}{1.154455in}}%
\pgfpathlineto{\pgfqpoint{2.532425in}{1.163870in}}%
\pgfpathlineto{\pgfqpoint{2.525737in}{1.174087in}}%
\pgfpathlineto{\pgfqpoint{2.527409in}{1.177624in}}%
\pgfpathlineto{\pgfqpoint{2.532999in}{1.178683in}}%
\pgfpathlineto{\pgfqpoint{2.532223in}{1.188646in}}%
\pgfpathlineto{\pgfqpoint{2.534247in}{1.194214in}}%
\pgfpathlineto{\pgfqpoint{2.523965in}{1.201691in}}%
\pgfpathlineto{\pgfqpoint{2.525644in}{1.210611in}}%
\pgfpathlineto{\pgfqpoint{2.532780in}{1.212166in}}%
\pgfpathlineto{\pgfqpoint{2.538153in}{1.216728in}}%
\pgfpathlineto{\pgfqpoint{2.543719in}{1.212838in}}%
\pgfpathlineto{\pgfqpoint{2.549426in}{1.216536in}}%
\pgfpathlineto{\pgfqpoint{2.560339in}{1.213292in}}%
\pgfpathlineto{\pgfqpoint{2.567300in}{1.215781in}}%
\pgfpathlineto{\pgfqpoint{2.569988in}{1.214538in}}%
\pgfpathlineto{\pgfqpoint{2.568471in}{1.207938in}}%
\pgfpathlineto{\pgfqpoint{2.571223in}{1.207045in}}%
\pgfpathlineto{\pgfqpoint{2.569924in}{1.199083in}}%
\pgfpathlineto{\pgfqpoint{2.562770in}{1.192250in}}%
\pgfpathlineto{\pgfqpoint{2.563877in}{1.187087in}}%
\pgfpathlineto{\pgfqpoint{2.558875in}{1.172717in}}%
\pgfpathclose%
\pgfusepath{fill}%
\end{pgfscope}%
\begin{pgfscope}%
\pgfpathrectangle{\pgfqpoint{0.100000in}{0.100000in}}{\pgfqpoint{3.007045in}{1.925000in}}%
\pgfusepath{clip}%
\pgfsetbuttcap%
\pgfsetmiterjoin%
\definecolor{currentfill}{rgb}{0.557324,0.816424,0.546959}%
\pgfsetfillcolor{currentfill}%
\pgfsetlinewidth{0.000000pt}%
\definecolor{currentstroke}{rgb}{0.000000,0.000000,0.000000}%
\pgfsetstrokecolor{currentstroke}%
\pgfsetstrokeopacity{0.000000}%
\pgfsetdash{}{0pt}%
\pgfpathmoveto{\pgfqpoint{2.697553in}{1.455914in}}%
\pgfpathlineto{\pgfqpoint{2.679109in}{1.451464in}}%
\pgfpathlineto{\pgfqpoint{2.675608in}{1.465295in}}%
\pgfpathlineto{\pgfqpoint{2.656543in}{1.461343in}}%
\pgfpathlineto{\pgfqpoint{2.652395in}{1.466521in}}%
\pgfpathlineto{\pgfqpoint{2.648634in}{1.483094in}}%
\pgfpathlineto{\pgfqpoint{2.648812in}{1.490538in}}%
\pgfpathlineto{\pgfqpoint{2.644415in}{1.507385in}}%
\pgfpathlineto{\pgfqpoint{2.647780in}{1.512793in}}%
\pgfpathlineto{\pgfqpoint{2.656259in}{1.521763in}}%
\pgfpathlineto{\pgfqpoint{2.664991in}{1.524305in}}%
\pgfpathlineto{\pgfqpoint{2.664388in}{1.534234in}}%
\pgfpathlineto{\pgfqpoint{2.672204in}{1.537694in}}%
\pgfpathlineto{\pgfqpoint{2.684161in}{1.539019in}}%
\pgfpathlineto{\pgfqpoint{2.688054in}{1.525089in}}%
\pgfpathlineto{\pgfqpoint{2.699199in}{1.524078in}}%
\pgfpathlineto{\pgfqpoint{2.716233in}{1.541048in}}%
\pgfpathlineto{\pgfqpoint{2.706539in}{1.571514in}}%
\pgfpathlineto{\pgfqpoint{2.712187in}{1.569616in}}%
\pgfpathlineto{\pgfqpoint{2.721573in}{1.573055in}}%
\pgfpathlineto{\pgfqpoint{2.733723in}{1.536267in}}%
\pgfpathlineto{\pgfqpoint{2.731549in}{1.525832in}}%
\pgfpathlineto{\pgfqpoint{2.739601in}{1.524017in}}%
\pgfpathlineto{\pgfqpoint{2.742037in}{1.517084in}}%
\pgfpathlineto{\pgfqpoint{2.739690in}{1.510307in}}%
\pgfpathlineto{\pgfqpoint{2.742514in}{1.505188in}}%
\pgfpathlineto{\pgfqpoint{2.743391in}{1.496243in}}%
\pgfpathlineto{\pgfqpoint{2.737431in}{1.497269in}}%
\pgfpathlineto{\pgfqpoint{2.737175in}{1.492249in}}%
\pgfpathlineto{\pgfqpoint{2.726801in}{1.495572in}}%
\pgfpathlineto{\pgfqpoint{2.722365in}{1.491084in}}%
\pgfpathlineto{\pgfqpoint{2.719947in}{1.482951in}}%
\pgfpathlineto{\pgfqpoint{2.692078in}{1.475538in}}%
\pgfpathlineto{\pgfqpoint{2.697553in}{1.455914in}}%
\pgfpathclose%
\pgfusepath{fill}%
\end{pgfscope}%
\begin{pgfscope}%
\pgfpathrectangle{\pgfqpoint{0.100000in}{0.100000in}}{\pgfqpoint{3.007045in}{1.925000in}}%
\pgfusepath{clip}%
\pgfsetbuttcap%
\pgfsetmiterjoin%
\definecolor{currentfill}{rgb}{0.395294,0.739408,0.433526}%
\pgfsetfillcolor{currentfill}%
\pgfsetlinewidth{0.000000pt}%
\definecolor{currentstroke}{rgb}{0.000000,0.000000,0.000000}%
\pgfsetstrokecolor{currentstroke}%
\pgfsetstrokeopacity{0.000000}%
\pgfsetdash{}{0pt}%
\pgfpathmoveto{\pgfqpoint{1.779241in}{1.344556in}}%
\pgfpathlineto{\pgfqpoint{1.813145in}{1.344757in}}%
\pgfpathlineto{\pgfqpoint{1.813376in}{1.321894in}}%
\pgfpathlineto{\pgfqpoint{1.784860in}{1.321668in}}%
\pgfpathlineto{\pgfqpoint{1.744960in}{1.321802in}}%
\pgfpathlineto{\pgfqpoint{1.739290in}{1.321824in}}%
\pgfpathlineto{\pgfqpoint{1.739541in}{1.344779in}}%
\pgfpathlineto{\pgfqpoint{1.779241in}{1.344556in}}%
\pgfpathclose%
\pgfusepath{fill}%
\end{pgfscope}%
\begin{pgfscope}%
\pgfpathrectangle{\pgfqpoint{0.100000in}{0.100000in}}{\pgfqpoint{3.007045in}{1.925000in}}%
\pgfusepath{clip}%
\pgfsetbuttcap%
\pgfsetmiterjoin%
\definecolor{currentfill}{rgb}{0.755848,0.903391,0.729043}%
\pgfsetfillcolor{currentfill}%
\pgfsetlinewidth{0.000000pt}%
\definecolor{currentstroke}{rgb}{0.000000,0.000000,0.000000}%
\pgfsetstrokecolor{currentstroke}%
\pgfsetstrokeopacity{0.000000}%
\pgfsetdash{}{0pt}%
\pgfpathmoveto{\pgfqpoint{2.222669in}{0.734886in}}%
\pgfpathlineto{\pgfqpoint{2.223953in}{0.726994in}}%
\pgfpathlineto{\pgfqpoint{2.230267in}{0.714046in}}%
\pgfpathlineto{\pgfqpoint{2.225729in}{0.705930in}}%
\pgfpathlineto{\pgfqpoint{2.226858in}{0.694238in}}%
\pgfpathlineto{\pgfqpoint{2.211673in}{0.692874in}}%
\pgfpathlineto{\pgfqpoint{2.202255in}{0.702095in}}%
\pgfpathlineto{\pgfqpoint{2.193997in}{0.705851in}}%
\pgfpathlineto{\pgfqpoint{2.190855in}{0.708411in}}%
\pgfpathlineto{\pgfqpoint{2.189892in}{0.719957in}}%
\pgfpathlineto{\pgfqpoint{2.178513in}{0.718958in}}%
\pgfpathlineto{\pgfqpoint{2.173613in}{0.721470in}}%
\pgfpathlineto{\pgfqpoint{2.174689in}{0.726922in}}%
\pgfpathlineto{\pgfqpoint{2.171759in}{0.737367in}}%
\pgfpathlineto{\pgfqpoint{2.172078in}{0.746133in}}%
\pgfpathlineto{\pgfqpoint{2.175967in}{0.749888in}}%
\pgfpathlineto{\pgfqpoint{2.176652in}{0.754493in}}%
\pgfpathlineto{\pgfqpoint{2.169944in}{0.753927in}}%
\pgfpathlineto{\pgfqpoint{2.168941in}{0.762907in}}%
\pgfpathlineto{\pgfqpoint{2.166832in}{0.787387in}}%
\pgfpathlineto{\pgfqpoint{2.176418in}{0.787936in}}%
\pgfpathlineto{\pgfqpoint{2.177866in}{0.793913in}}%
\pgfpathlineto{\pgfqpoint{2.192221in}{0.794560in}}%
\pgfpathlineto{\pgfqpoint{2.198442in}{0.789193in}}%
\pgfpathlineto{\pgfqpoint{2.197889in}{0.784340in}}%
\pgfpathlineto{\pgfqpoint{2.204331in}{0.776956in}}%
\pgfpathlineto{\pgfqpoint{2.213183in}{0.772767in}}%
\pgfpathlineto{\pgfqpoint{2.213720in}{0.767453in}}%
\pgfpathlineto{\pgfqpoint{2.217222in}{0.763282in}}%
\pgfpathlineto{\pgfqpoint{2.222615in}{0.760540in}}%
\pgfpathlineto{\pgfqpoint{2.223920in}{0.746448in}}%
\pgfpathlineto{\pgfqpoint{2.215985in}{0.745794in}}%
\pgfpathlineto{\pgfqpoint{2.217078in}{0.734332in}}%
\pgfpathlineto{\pgfqpoint{2.222669in}{0.734886in}}%
\pgfpathclose%
\pgfusepath{fill}%
\end{pgfscope}%
\begin{pgfscope}%
\pgfpathrectangle{\pgfqpoint{0.100000in}{0.100000in}}{\pgfqpoint{3.007045in}{1.925000in}}%
\pgfusepath{clip}%
\pgfsetbuttcap%
\pgfsetmiterjoin%
\definecolor{currentfill}{rgb}{0.479815,0.780254,0.483230}%
\pgfsetfillcolor{currentfill}%
\pgfsetlinewidth{0.000000pt}%
\definecolor{currentstroke}{rgb}{0.000000,0.000000,0.000000}%
\pgfsetstrokecolor{currentstroke}%
\pgfsetstrokeopacity{0.000000}%
\pgfsetdash{}{0pt}%
\pgfpathmoveto{\pgfqpoint{1.869963in}{1.345597in}}%
\pgfpathlineto{\pgfqpoint{1.858723in}{1.345342in}}%
\pgfpathlineto{\pgfqpoint{1.847419in}{1.345161in}}%
\pgfpathlineto{\pgfqpoint{1.847146in}{1.362415in}}%
\pgfpathlineto{\pgfqpoint{1.845372in}{1.362394in}}%
\pgfpathlineto{\pgfqpoint{1.844917in}{1.391266in}}%
\pgfpathlineto{\pgfqpoint{1.867818in}{1.391658in}}%
\pgfpathlineto{\pgfqpoint{1.868339in}{1.362827in}}%
\pgfpathlineto{\pgfqpoint{1.869522in}{1.362845in}}%
\pgfpathlineto{\pgfqpoint{1.869963in}{1.345597in}}%
\pgfpathclose%
\pgfusepath{fill}%
\end{pgfscope}%
\begin{pgfscope}%
\pgfpathrectangle{\pgfqpoint{0.100000in}{0.100000in}}{\pgfqpoint{3.007045in}{1.925000in}}%
\pgfusepath{clip}%
\pgfsetbuttcap%
\pgfsetmiterjoin%
\definecolor{currentfill}{rgb}{0.662345,0.864022,0.638001}%
\pgfsetfillcolor{currentfill}%
\pgfsetlinewidth{0.000000pt}%
\definecolor{currentstroke}{rgb}{0.000000,0.000000,0.000000}%
\pgfsetstrokecolor{currentstroke}%
\pgfsetstrokeopacity{0.000000}%
\pgfsetdash{}{0pt}%
\pgfpathmoveto{\pgfqpoint{2.283546in}{1.435118in}}%
\pgfpathlineto{\pgfqpoint{2.283625in}{1.434436in}}%
\pgfpathlineto{\pgfqpoint{2.261163in}{1.431916in}}%
\pgfpathlineto{\pgfqpoint{2.258163in}{1.454659in}}%
\pgfpathlineto{\pgfqpoint{2.246778in}{1.453408in}}%
\pgfpathlineto{\pgfqpoint{2.241654in}{1.499264in}}%
\pgfpathlineto{\pgfqpoint{2.252766in}{1.500298in}}%
\pgfpathlineto{\pgfqpoint{2.264263in}{1.501550in}}%
\pgfpathlineto{\pgfqpoint{2.269523in}{1.455942in}}%
\pgfpathlineto{\pgfqpoint{2.280842in}{1.457266in}}%
\pgfpathlineto{\pgfqpoint{2.283546in}{1.435118in}}%
\pgfpathclose%
\pgfusepath{fill}%
\end{pgfscope}%
\begin{pgfscope}%
\pgfpathrectangle{\pgfqpoint{0.100000in}{0.100000in}}{\pgfqpoint{3.007045in}{1.925000in}}%
\pgfusepath{clip}%
\pgfsetbuttcap%
\pgfsetmiterjoin%
\definecolor{currentfill}{rgb}{0.320000,0.702499,0.396617}%
\pgfsetfillcolor{currentfill}%
\pgfsetlinewidth{0.000000pt}%
\definecolor{currentstroke}{rgb}{0.000000,0.000000,0.000000}%
\pgfsetstrokecolor{currentstroke}%
\pgfsetstrokeopacity{0.000000}%
\pgfsetdash{}{0pt}%
\pgfpathmoveto{\pgfqpoint{2.418790in}{0.913875in}}%
\pgfpathlineto{\pgfqpoint{2.416370in}{0.915795in}}%
\pgfpathlineto{\pgfqpoint{2.421150in}{0.927599in}}%
\pgfpathlineto{\pgfqpoint{2.420735in}{0.932481in}}%
\pgfpathlineto{\pgfqpoint{2.405701in}{0.945045in}}%
\pgfpathlineto{\pgfqpoint{2.403998in}{0.954773in}}%
\pgfpathlineto{\pgfqpoint{2.399943in}{0.955887in}}%
\pgfpathlineto{\pgfqpoint{2.404400in}{0.961065in}}%
\pgfpathlineto{\pgfqpoint{2.414244in}{0.964014in}}%
\pgfpathlineto{\pgfqpoint{2.416656in}{0.974681in}}%
\pgfpathlineto{\pgfqpoint{2.429667in}{0.984471in}}%
\pgfpathlineto{\pgfqpoint{2.431676in}{0.978202in}}%
\pgfpathlineto{\pgfqpoint{2.437163in}{0.979641in}}%
\pgfpathlineto{\pgfqpoint{2.437960in}{0.976572in}}%
\pgfpathlineto{\pgfqpoint{2.446969in}{0.969925in}}%
\pgfpathlineto{\pgfqpoint{2.454474in}{0.954214in}}%
\pgfpathlineto{\pgfqpoint{2.459120in}{0.952751in}}%
\pgfpathlineto{\pgfqpoint{2.455786in}{0.951722in}}%
\pgfpathlineto{\pgfqpoint{2.453980in}{0.946316in}}%
\pgfpathlineto{\pgfqpoint{2.455447in}{0.943242in}}%
\pgfpathlineto{\pgfqpoint{2.451924in}{0.935651in}}%
\pgfpathlineto{\pgfqpoint{2.452389in}{0.929337in}}%
\pgfpathlineto{\pgfqpoint{2.435961in}{0.922177in}}%
\pgfpathlineto{\pgfqpoint{2.418790in}{0.913875in}}%
\pgfpathclose%
\pgfusepath{fill}%
\end{pgfscope}%
\begin{pgfscope}%
\pgfpathrectangle{\pgfqpoint{0.100000in}{0.100000in}}{\pgfqpoint{3.007045in}{1.925000in}}%
\pgfusepath{clip}%
\pgfsetbuttcap%
\pgfsetmiterjoin%
\definecolor{currentfill}{rgb}{0.257255,0.671742,0.365859}%
\pgfsetfillcolor{currentfill}%
\pgfsetlinewidth{0.000000pt}%
\definecolor{currentstroke}{rgb}{0.000000,0.000000,0.000000}%
\pgfsetstrokecolor{currentstroke}%
\pgfsetstrokeopacity{0.000000}%
\pgfsetdash{}{0pt}%
\pgfpathmoveto{\pgfqpoint{1.910330in}{0.935711in}}%
\pgfpathlineto{\pgfqpoint{1.910553in}{0.923358in}}%
\pgfpathlineto{\pgfqpoint{1.899149in}{0.923195in}}%
\pgfpathlineto{\pgfqpoint{1.899200in}{0.918414in}}%
\pgfpathlineto{\pgfqpoint{1.889632in}{0.918440in}}%
\pgfpathlineto{\pgfqpoint{1.851607in}{0.918465in}}%
\pgfpathlineto{\pgfqpoint{1.851596in}{0.920385in}}%
\pgfpathlineto{\pgfqpoint{1.851232in}{0.928170in}}%
\pgfpathlineto{\pgfqpoint{1.855032in}{0.933900in}}%
\pgfpathlineto{\pgfqpoint{1.855176in}{0.941374in}}%
\pgfpathlineto{\pgfqpoint{1.849514in}{0.944275in}}%
\pgfpathlineto{\pgfqpoint{1.845759in}{0.949170in}}%
\pgfpathlineto{\pgfqpoint{1.843659in}{0.955855in}}%
\pgfpathlineto{\pgfqpoint{1.832421in}{0.955990in}}%
\pgfpathlineto{\pgfqpoint{1.832387in}{0.968604in}}%
\pgfpathlineto{\pgfqpoint{1.886046in}{0.969526in}}%
\pgfpathlineto{\pgfqpoint{1.884426in}{0.969017in}}%
\pgfpathlineto{\pgfqpoint{1.884698in}{0.944172in}}%
\pgfpathlineto{\pgfqpoint{1.887586in}{0.941336in}}%
\pgfpathlineto{\pgfqpoint{1.910317in}{0.941444in}}%
\pgfpathlineto{\pgfqpoint{1.910330in}{0.935711in}}%
\pgfpathclose%
\pgfusepath{fill}%
\end{pgfscope}%
\begin{pgfscope}%
\pgfpathrectangle{\pgfqpoint{0.100000in}{0.100000in}}{\pgfqpoint{3.007045in}{1.925000in}}%
\pgfusepath{clip}%
\pgfsetbuttcap%
\pgfsetmiterjoin%
\definecolor{currentfill}{rgb}{0.662345,0.864022,0.638001}%
\pgfsetfillcolor{currentfill}%
\pgfsetlinewidth{0.000000pt}%
\definecolor{currentstroke}{rgb}{0.000000,0.000000,0.000000}%
\pgfsetstrokecolor{currentstroke}%
\pgfsetstrokeopacity{0.000000}%
\pgfsetdash{}{0pt}%
\pgfpathmoveto{\pgfqpoint{2.122132in}{1.377026in}}%
\pgfpathlineto{\pgfqpoint{2.078237in}{1.373879in}}%
\pgfpathlineto{\pgfqpoint{2.046233in}{1.372218in}}%
\pgfpathlineto{\pgfqpoint{2.044679in}{1.394919in}}%
\pgfpathlineto{\pgfqpoint{2.061841in}{1.396171in}}%
\pgfpathlineto{\pgfqpoint{2.073275in}{1.396646in}}%
\pgfpathlineto{\pgfqpoint{2.119150in}{1.400006in}}%
\pgfpathlineto{\pgfqpoint{2.122786in}{1.396116in}}%
\pgfpathlineto{\pgfqpoint{2.120693in}{1.385146in}}%
\pgfpathlineto{\pgfqpoint{2.122132in}{1.377026in}}%
\pgfpathclose%
\pgfusepath{fill}%
\end{pgfscope}%
\begin{pgfscope}%
\pgfpathrectangle{\pgfqpoint{0.100000in}{0.100000in}}{\pgfqpoint{3.007045in}{1.925000in}}%
\pgfusepath{clip}%
\pgfsetbuttcap%
\pgfsetmiterjoin%
\definecolor{currentfill}{rgb}{0.723122,0.889612,0.697178}%
\pgfsetfillcolor{currentfill}%
\pgfsetlinewidth{0.000000pt}%
\definecolor{currentstroke}{rgb}{0.000000,0.000000,0.000000}%
\pgfsetstrokecolor{currentstroke}%
\pgfsetstrokeopacity{0.000000}%
\pgfsetdash{}{0pt}%
\pgfpathmoveto{\pgfqpoint{1.495136in}{0.800124in}}%
\pgfpathlineto{\pgfqpoint{1.469347in}{0.801604in}}%
\pgfpathlineto{\pgfqpoint{1.471341in}{0.833174in}}%
\pgfpathlineto{\pgfqpoint{1.476764in}{0.832809in}}%
\pgfpathlineto{\pgfqpoint{1.478574in}{0.861428in}}%
\pgfpathlineto{\pgfqpoint{1.501084in}{0.860066in}}%
\pgfpathlineto{\pgfqpoint{1.500565in}{0.847796in}}%
\pgfpathlineto{\pgfqpoint{1.499337in}{0.825611in}}%
\pgfpathlineto{\pgfqpoint{1.496660in}{0.826118in}}%
\pgfpathlineto{\pgfqpoint{1.495136in}{0.800124in}}%
\pgfpathclose%
\pgfusepath{fill}%
\end{pgfscope}%
\begin{pgfscope}%
\pgfpathrectangle{\pgfqpoint{0.100000in}{0.100000in}}{\pgfqpoint{3.007045in}{1.925000in}}%
\pgfusepath{clip}%
\pgfsetbuttcap%
\pgfsetmiterjoin%
\definecolor{currentfill}{rgb}{0.667020,0.865990,0.642553}%
\pgfsetfillcolor{currentfill}%
\pgfsetlinewidth{0.000000pt}%
\definecolor{currentstroke}{rgb}{0.000000,0.000000,0.000000}%
\pgfsetstrokecolor{currentstroke}%
\pgfsetstrokeopacity{0.000000}%
\pgfsetdash{}{0pt}%
\pgfpathmoveto{\pgfqpoint{1.965268in}{0.666642in}}%
\pgfpathlineto{\pgfqpoint{1.965330in}{0.661044in}}%
\pgfpathlineto{\pgfqpoint{1.962071in}{0.656269in}}%
\pgfpathlineto{\pgfqpoint{1.959044in}{0.656761in}}%
\pgfpathlineto{\pgfqpoint{1.954469in}{0.643065in}}%
\pgfpathlineto{\pgfqpoint{1.949979in}{0.638422in}}%
\pgfpathlineto{\pgfqpoint{1.952131in}{0.625210in}}%
\pgfpathlineto{\pgfqpoint{1.942513in}{0.628295in}}%
\pgfpathlineto{\pgfqpoint{1.933963in}{0.633867in}}%
\pgfpathlineto{\pgfqpoint{1.933198in}{0.637846in}}%
\pgfpathlineto{\pgfqpoint{1.926639in}{0.641691in}}%
\pgfpathlineto{\pgfqpoint{1.922010in}{0.647995in}}%
\pgfpathlineto{\pgfqpoint{1.921110in}{0.659462in}}%
\pgfpathlineto{\pgfqpoint{1.923673in}{0.668219in}}%
\pgfpathlineto{\pgfqpoint{1.946656in}{0.668963in}}%
\pgfpathlineto{\pgfqpoint{1.949500in}{0.671101in}}%
\pgfpathlineto{\pgfqpoint{1.953769in}{0.665525in}}%
\pgfpathlineto{\pgfqpoint{1.963857in}{0.672381in}}%
\pgfpathlineto{\pgfqpoint{1.965268in}{0.666642in}}%
\pgfpathclose%
\pgfusepath{fill}%
\end{pgfscope}%
\begin{pgfscope}%
\pgfpathrectangle{\pgfqpoint{0.100000in}{0.100000in}}{\pgfqpoint{3.007045in}{1.925000in}}%
\pgfusepath{clip}%
\pgfsetbuttcap%
\pgfsetmiterjoin%
\definecolor{currentfill}{rgb}{0.518570,0.798339,0.515094}%
\pgfsetfillcolor{currentfill}%
\pgfsetlinewidth{0.000000pt}%
\definecolor{currentstroke}{rgb}{0.000000,0.000000,0.000000}%
\pgfsetstrokecolor{currentstroke}%
\pgfsetstrokeopacity{0.000000}%
\pgfsetdash{}{0pt}%
\pgfpathmoveto{\pgfqpoint{2.198442in}{0.789193in}}%
\pgfpathlineto{\pgfqpoint{2.192221in}{0.794560in}}%
\pgfpathlineto{\pgfqpoint{2.177866in}{0.793913in}}%
\pgfpathlineto{\pgfqpoint{2.176418in}{0.787936in}}%
\pgfpathlineto{\pgfqpoint{2.166832in}{0.787387in}}%
\pgfpathlineto{\pgfqpoint{2.142921in}{0.786166in}}%
\pgfpathlineto{\pgfqpoint{2.143299in}{0.800190in}}%
\pgfpathlineto{\pgfqpoint{2.144570in}{0.838481in}}%
\pgfpathlineto{\pgfqpoint{2.179696in}{0.840159in}}%
\pgfpathlineto{\pgfqpoint{2.202564in}{0.841744in}}%
\pgfpathlineto{\pgfqpoint{2.204539in}{0.819210in}}%
\pgfpathlineto{\pgfqpoint{2.206348in}{0.816022in}}%
\pgfpathlineto{\pgfqpoint{2.213109in}{0.813406in}}%
\pgfpathlineto{\pgfqpoint{2.210327in}{0.807475in}}%
\pgfpathlineto{\pgfqpoint{2.209929in}{0.801633in}}%
\pgfpathlineto{\pgfqpoint{2.207356in}{0.797594in}}%
\pgfpathlineto{\pgfqpoint{2.202677in}{0.796254in}}%
\pgfpathlineto{\pgfqpoint{2.198442in}{0.789193in}}%
\pgfpathclose%
\pgfusepath{fill}%
\end{pgfscope}%
\begin{pgfscope}%
\pgfpathrectangle{\pgfqpoint{0.100000in}{0.100000in}}{\pgfqpoint{3.007045in}{1.925000in}}%
\pgfusepath{clip}%
\pgfsetbuttcap%
\pgfsetmiterjoin%
\definecolor{currentfill}{rgb}{0.457670,0.769919,0.465021}%
\pgfsetfillcolor{currentfill}%
\pgfsetlinewidth{0.000000pt}%
\definecolor{currentstroke}{rgb}{0.000000,0.000000,0.000000}%
\pgfsetstrokecolor{currentstroke}%
\pgfsetstrokeopacity{0.000000}%
\pgfsetdash{}{0pt}%
\pgfpathmoveto{\pgfqpoint{2.474213in}{0.983871in}}%
\pgfpathlineto{\pgfqpoint{2.480170in}{0.978949in}}%
\pgfpathlineto{\pgfqpoint{2.489274in}{0.973776in}}%
\pgfpathlineto{\pgfqpoint{2.494999in}{0.975158in}}%
\pgfpathlineto{\pgfqpoint{2.498902in}{0.973883in}}%
\pgfpathlineto{\pgfqpoint{2.492448in}{0.960338in}}%
\pgfpathlineto{\pgfqpoint{2.486422in}{0.958962in}}%
\pgfpathlineto{\pgfqpoint{2.480332in}{0.961447in}}%
\pgfpathlineto{\pgfqpoint{2.476332in}{0.956183in}}%
\pgfpathlineto{\pgfqpoint{2.469542in}{0.954025in}}%
\pgfpathlineto{\pgfqpoint{2.459120in}{0.952751in}}%
\pgfpathlineto{\pgfqpoint{2.454474in}{0.954214in}}%
\pgfpathlineto{\pgfqpoint{2.446969in}{0.969925in}}%
\pgfpathlineto{\pgfqpoint{2.437960in}{0.976572in}}%
\pgfpathlineto{\pgfqpoint{2.437163in}{0.979641in}}%
\pgfpathlineto{\pgfqpoint{2.441498in}{0.987564in}}%
\pgfpathlineto{\pgfqpoint{2.450485in}{0.993442in}}%
\pgfpathlineto{\pgfqpoint{2.458453in}{0.990986in}}%
\pgfpathlineto{\pgfqpoint{2.462402in}{0.980829in}}%
\pgfpathlineto{\pgfqpoint{2.465463in}{0.979232in}}%
\pgfpathlineto{\pgfqpoint{2.468967in}{0.985644in}}%
\pgfpathlineto{\pgfqpoint{2.474213in}{0.983871in}}%
\pgfpathclose%
\pgfusepath{fill}%
\end{pgfscope}%
\begin{pgfscope}%
\pgfpathrectangle{\pgfqpoint{0.100000in}{0.100000in}}{\pgfqpoint{3.007045in}{1.925000in}}%
\pgfusepath{clip}%
\pgfsetbuttcap%
\pgfsetmiterjoin%
\definecolor{currentfill}{rgb}{0.125413,0.534948,0.262130}%
\pgfsetfillcolor{currentfill}%
\pgfsetlinewidth{0.000000pt}%
\definecolor{currentstroke}{rgb}{0.000000,0.000000,0.000000}%
\pgfsetstrokecolor{currentstroke}%
\pgfsetstrokeopacity{0.000000}%
\pgfsetdash{}{0pt}%
\pgfpathmoveto{\pgfqpoint{0.558722in}{1.943515in}}%
\pgfpathlineto{\pgfqpoint{0.550057in}{1.949005in}}%
\pgfpathlineto{\pgfqpoint{0.546058in}{1.955336in}}%
\pgfpathlineto{\pgfqpoint{0.547224in}{1.961309in}}%
\pgfpathlineto{\pgfqpoint{0.552633in}{1.957537in}}%
\pgfpathlineto{\pgfqpoint{0.559107in}{1.963705in}}%
\pgfpathlineto{\pgfqpoint{0.566457in}{1.957886in}}%
\pgfpathlineto{\pgfqpoint{0.558722in}{1.943515in}}%
\pgfpathclose%
\pgfusepath{fill}%
\end{pgfscope}%
\begin{pgfscope}%
\pgfpathrectangle{\pgfqpoint{0.100000in}{0.100000in}}{\pgfqpoint{3.007045in}{1.925000in}}%
\pgfusepath{clip}%
\pgfsetbuttcap%
\pgfsetmiterjoin%
\definecolor{currentfill}{rgb}{0.618224,0.844844,0.597032}%
\pgfsetfillcolor{currentfill}%
\pgfsetlinewidth{0.000000pt}%
\definecolor{currentstroke}{rgb}{0.000000,0.000000,0.000000}%
\pgfsetstrokecolor{currentstroke}%
\pgfsetstrokeopacity{0.000000}%
\pgfsetdash{}{0pt}%
\pgfpathmoveto{\pgfqpoint{1.386813in}{0.624056in}}%
\pgfpathlineto{\pgfqpoint{1.415290in}{0.621980in}}%
\pgfpathlineto{\pgfqpoint{1.433065in}{0.620988in}}%
\pgfpathlineto{\pgfqpoint{1.431514in}{0.595996in}}%
\pgfpathlineto{\pgfqpoint{1.429685in}{0.568203in}}%
\pgfpathlineto{\pgfqpoint{1.383908in}{0.571430in}}%
\pgfpathlineto{\pgfqpoint{1.389923in}{0.575693in}}%
\pgfpathlineto{\pgfqpoint{1.385588in}{0.580234in}}%
\pgfpathlineto{\pgfqpoint{1.391278in}{0.588321in}}%
\pgfpathlineto{\pgfqpoint{1.390664in}{0.594316in}}%
\pgfpathlineto{\pgfqpoint{1.386971in}{0.595647in}}%
\pgfpathlineto{\pgfqpoint{1.364071in}{0.597523in}}%
\pgfpathlineto{\pgfqpoint{1.363729in}{0.593646in}}%
\pgfpathlineto{\pgfqpoint{1.352050in}{0.594662in}}%
\pgfpathlineto{\pgfqpoint{1.350354in}{0.573776in}}%
\pgfpathlineto{\pgfqpoint{1.337515in}{0.574775in}}%
\pgfpathlineto{\pgfqpoint{1.336243in}{0.559459in}}%
\pgfpathlineto{\pgfqpoint{1.282209in}{0.612003in}}%
\pgfpathlineto{\pgfqpoint{1.318428in}{0.649073in}}%
\pgfpathlineto{\pgfqpoint{1.322677in}{0.647278in}}%
\pgfpathlineto{\pgfqpoint{1.328596in}{0.641304in}}%
\pgfpathlineto{\pgfqpoint{1.331766in}{0.642740in}}%
\pgfpathlineto{\pgfqpoint{1.337100in}{0.644492in}}%
\pgfpathlineto{\pgfqpoint{1.349725in}{0.635825in}}%
\pgfpathlineto{\pgfqpoint{1.351088in}{0.627381in}}%
\pgfpathlineto{\pgfqpoint{1.386813in}{0.624056in}}%
\pgfpathclose%
\pgfusepath{fill}%
\end{pgfscope}%
\begin{pgfscope}%
\pgfpathrectangle{\pgfqpoint{0.100000in}{0.100000in}}{\pgfqpoint{3.007045in}{1.925000in}}%
\pgfusepath{clip}%
\pgfsetbuttcap%
\pgfsetmiterjoin%
\definecolor{currentfill}{rgb}{0.590542,0.831926,0.574271}%
\pgfsetfillcolor{currentfill}%
\pgfsetlinewidth{0.000000pt}%
\definecolor{currentstroke}{rgb}{0.000000,0.000000,0.000000}%
\pgfsetstrokecolor{currentstroke}%
\pgfsetstrokeopacity{0.000000}%
\pgfsetdash{}{0pt}%
\pgfpathmoveto{\pgfqpoint{1.930951in}{0.750384in}}%
\pgfpathlineto{\pgfqpoint{1.932837in}{0.741766in}}%
\pgfpathlineto{\pgfqpoint{1.935072in}{0.740260in}}%
\pgfpathlineto{\pgfqpoint{1.898651in}{0.739626in}}%
\pgfpathlineto{\pgfqpoint{1.884006in}{0.739484in}}%
\pgfpathlineto{\pgfqpoint{1.884045in}{0.763316in}}%
\pgfpathlineto{\pgfqpoint{1.876274in}{0.763349in}}%
\pgfpathlineto{\pgfqpoint{1.876321in}{0.768157in}}%
\pgfpathlineto{\pgfqpoint{1.876541in}{0.789626in}}%
\pgfpathlineto{\pgfqpoint{1.884971in}{0.791103in}}%
\pgfpathlineto{\pgfqpoint{1.887657in}{0.795131in}}%
\pgfpathlineto{\pgfqpoint{1.887940in}{0.814920in}}%
\pgfpathlineto{\pgfqpoint{1.899320in}{0.814784in}}%
\pgfpathlineto{\pgfqpoint{1.909437in}{0.814697in}}%
\pgfpathlineto{\pgfqpoint{1.918185in}{0.809351in}}%
\pgfpathlineto{\pgfqpoint{1.910795in}{0.808969in}}%
\pgfpathlineto{\pgfqpoint{1.911801in}{0.803402in}}%
\pgfpathlineto{\pgfqpoint{1.918485in}{0.794539in}}%
\pgfpathlineto{\pgfqpoint{1.920677in}{0.772374in}}%
\pgfpathlineto{\pgfqpoint{1.917305in}{0.764130in}}%
\pgfpathlineto{\pgfqpoint{1.918474in}{0.757862in}}%
\pgfpathlineto{\pgfqpoint{1.921911in}{0.757915in}}%
\pgfpathlineto{\pgfqpoint{1.930951in}{0.750384in}}%
\pgfpathclose%
\pgfusepath{fill}%
\end{pgfscope}%
\begin{pgfscope}%
\pgfpathrectangle{\pgfqpoint{0.100000in}{0.100000in}}{\pgfqpoint{3.007045in}{1.925000in}}%
\pgfusepath{clip}%
\pgfsetbuttcap%
\pgfsetmiterjoin%
\definecolor{currentfill}{rgb}{0.414118,0.748635,0.442753}%
\pgfsetfillcolor{currentfill}%
\pgfsetlinewidth{0.000000pt}%
\definecolor{currentstroke}{rgb}{0.000000,0.000000,0.000000}%
\pgfsetstrokecolor{currentstroke}%
\pgfsetstrokeopacity{0.000000}%
\pgfsetdash{}{0pt}%
\pgfpathmoveto{\pgfqpoint{1.512812in}{1.275327in}}%
\pgfpathlineto{\pgfqpoint{1.511547in}{1.252543in}}%
\pgfpathlineto{\pgfqpoint{1.483897in}{1.254208in}}%
\pgfpathlineto{\pgfqpoint{1.460453in}{1.255605in}}%
\pgfpathlineto{\pgfqpoint{1.461934in}{1.278597in}}%
\pgfpathlineto{\pgfqpoint{1.460924in}{1.278668in}}%
\pgfpathlineto{\pgfqpoint{1.462537in}{1.301517in}}%
\pgfpathlineto{\pgfqpoint{1.455754in}{1.302023in}}%
\pgfpathlineto{\pgfqpoint{1.457408in}{1.324892in}}%
\pgfpathlineto{\pgfqpoint{1.485080in}{1.322839in}}%
\pgfpathlineto{\pgfqpoint{1.514112in}{1.321020in}}%
\pgfpathlineto{\pgfqpoint{1.512812in}{1.275327in}}%
\pgfpathclose%
\pgfusepath{fill}%
\end{pgfscope}%
\begin{pgfscope}%
\pgfpathrectangle{\pgfqpoint{0.100000in}{0.100000in}}{\pgfqpoint{3.007045in}{1.925000in}}%
\pgfusepath{clip}%
\pgfsetbuttcap%
\pgfsetmiterjoin%
\definecolor{currentfill}{rgb}{0.568397,0.821592,0.556063}%
\pgfsetfillcolor{currentfill}%
\pgfsetlinewidth{0.000000pt}%
\definecolor{currentstroke}{rgb}{0.000000,0.000000,0.000000}%
\pgfsetstrokecolor{currentstroke}%
\pgfsetstrokeopacity{0.000000}%
\pgfsetdash{}{0pt}%
\pgfpathmoveto{\pgfqpoint{2.345734in}{1.216099in}}%
\pgfpathlineto{\pgfqpoint{2.346250in}{1.205103in}}%
\pgfpathlineto{\pgfqpoint{2.330635in}{1.204428in}}%
\pgfpathlineto{\pgfqpoint{2.314937in}{1.203868in}}%
\pgfpathlineto{\pgfqpoint{2.305334in}{1.201390in}}%
\pgfpathlineto{\pgfqpoint{2.288393in}{1.199350in}}%
\pgfpathlineto{\pgfqpoint{2.285375in}{1.228022in}}%
\pgfpathlineto{\pgfqpoint{2.283471in}{1.248041in}}%
\pgfpathlineto{\pgfqpoint{2.283146in}{1.250837in}}%
\pgfpathlineto{\pgfqpoint{2.301562in}{1.253044in}}%
\pgfpathlineto{\pgfqpoint{2.306160in}{1.255224in}}%
\pgfpathlineto{\pgfqpoint{2.305349in}{1.261913in}}%
\pgfpathlineto{\pgfqpoint{2.322524in}{1.264076in}}%
\pgfpathlineto{\pgfqpoint{2.322114in}{1.267499in}}%
\pgfpathlineto{\pgfqpoint{2.327732in}{1.268420in}}%
\pgfpathlineto{\pgfqpoint{2.345898in}{1.268387in}}%
\pgfpathlineto{\pgfqpoint{2.346617in}{1.250238in}}%
\pgfpathlineto{\pgfqpoint{2.349506in}{1.250356in}}%
\pgfpathlineto{\pgfqpoint{2.350594in}{1.236213in}}%
\pgfpathlineto{\pgfqpoint{2.348588in}{1.219904in}}%
\pgfpathlineto{\pgfqpoint{2.345734in}{1.216099in}}%
\pgfpathclose%
\pgfusepath{fill}%
\end{pgfscope}%
\begin{pgfscope}%
\pgfpathrectangle{\pgfqpoint{0.100000in}{0.100000in}}{\pgfqpoint{3.007045in}{1.925000in}}%
\pgfusepath{clip}%
\pgfsetbuttcap%
\pgfsetmiterjoin%
\definecolor{currentfill}{rgb}{0.518570,0.798339,0.515094}%
\pgfsetfillcolor{currentfill}%
\pgfsetlinewidth{0.000000pt}%
\definecolor{currentstroke}{rgb}{0.000000,0.000000,0.000000}%
\pgfsetstrokecolor{currentstroke}%
\pgfsetstrokeopacity{0.000000}%
\pgfsetdash{}{0pt}%
\pgfpathmoveto{\pgfqpoint{1.925598in}{1.218425in}}%
\pgfpathlineto{\pgfqpoint{1.926181in}{1.198155in}}%
\pgfpathlineto{\pgfqpoint{1.906383in}{1.197788in}}%
\pgfpathlineto{\pgfqpoint{1.906198in}{1.203318in}}%
\pgfpathlineto{\pgfqpoint{1.880879in}{1.202651in}}%
\pgfpathlineto{\pgfqpoint{1.880439in}{1.202637in}}%
\pgfpathlineto{\pgfqpoint{1.879993in}{1.222750in}}%
\pgfpathlineto{\pgfqpoint{1.888599in}{1.223016in}}%
\pgfpathlineto{\pgfqpoint{1.889157in}{1.227613in}}%
\pgfpathlineto{\pgfqpoint{1.886676in}{1.239179in}}%
\pgfpathlineto{\pgfqpoint{1.890498in}{1.239370in}}%
\pgfpathlineto{\pgfqpoint{1.925211in}{1.241392in}}%
\pgfpathlineto{\pgfqpoint{1.925598in}{1.218425in}}%
\pgfpathclose%
\pgfusepath{fill}%
\end{pgfscope}%
\begin{pgfscope}%
\pgfpathrectangle{\pgfqpoint{0.100000in}{0.100000in}}{\pgfqpoint{3.007045in}{1.925000in}}%
\pgfusepath{clip}%
\pgfsetbuttcap%
\pgfsetmiterjoin%
\definecolor{currentfill}{rgb}{0.638970,0.854179,0.615240}%
\pgfsetfillcolor{currentfill}%
\pgfsetlinewidth{0.000000pt}%
\definecolor{currentstroke}{rgb}{0.000000,0.000000,0.000000}%
\pgfsetstrokecolor{currentstroke}%
\pgfsetstrokeopacity{0.000000}%
\pgfsetdash{}{0pt}%
\pgfpathmoveto{\pgfqpoint{1.842321in}{0.853020in}}%
\pgfpathlineto{\pgfqpoint{1.833670in}{0.850751in}}%
\pgfpathlineto{\pgfqpoint{1.830433in}{0.847601in}}%
\pgfpathlineto{\pgfqpoint{1.825651in}{0.849322in}}%
\pgfpathlineto{\pgfqpoint{1.809383in}{0.849499in}}%
\pgfpathlineto{\pgfqpoint{1.801936in}{0.851532in}}%
\pgfpathlineto{\pgfqpoint{1.801665in}{0.836897in}}%
\pgfpathlineto{\pgfqpoint{1.775725in}{0.836774in}}%
\pgfpathlineto{\pgfqpoint{1.775709in}{0.848268in}}%
\pgfpathlineto{\pgfqpoint{1.744379in}{0.848364in}}%
\pgfpathlineto{\pgfqpoint{1.744358in}{0.842618in}}%
\pgfpathlineto{\pgfqpoint{1.735784in}{0.842654in}}%
\pgfpathlineto{\pgfqpoint{1.724353in}{0.842730in}}%
\pgfpathlineto{\pgfqpoint{1.724413in}{0.848474in}}%
\pgfpathlineto{\pgfqpoint{1.713003in}{0.848605in}}%
\pgfpathlineto{\pgfqpoint{1.713476in}{0.872989in}}%
\pgfpathlineto{\pgfqpoint{1.719231in}{0.879638in}}%
\pgfpathlineto{\pgfqpoint{1.725397in}{0.881888in}}%
\pgfpathlineto{\pgfqpoint{1.732284in}{0.882217in}}%
\pgfpathlineto{\pgfqpoint{1.741580in}{0.885648in}}%
\pgfpathlineto{\pgfqpoint{1.751555in}{0.890916in}}%
\pgfpathlineto{\pgfqpoint{1.759980in}{0.886994in}}%
\pgfpathlineto{\pgfqpoint{1.763675in}{0.893735in}}%
\pgfpathlineto{\pgfqpoint{1.768002in}{0.897857in}}%
\pgfpathlineto{\pgfqpoint{1.765215in}{0.904084in}}%
\pgfpathlineto{\pgfqpoint{1.765492in}{0.911520in}}%
\pgfpathlineto{\pgfqpoint{1.800606in}{0.911576in}}%
\pgfpathlineto{\pgfqpoint{1.799463in}{0.919527in}}%
\pgfpathlineto{\pgfqpoint{1.819183in}{0.919101in}}%
\pgfpathlineto{\pgfqpoint{1.830635in}{0.919840in}}%
\pgfpathlineto{\pgfqpoint{1.826781in}{0.913265in}}%
\pgfpathlineto{\pgfqpoint{1.821996in}{0.913307in}}%
\pgfpathlineto{\pgfqpoint{1.821906in}{0.907584in}}%
\pgfpathlineto{\pgfqpoint{1.824480in}{0.903095in}}%
\pgfpathlineto{\pgfqpoint{1.822064in}{0.899000in}}%
\pgfpathlineto{\pgfqpoint{1.821794in}{0.883783in}}%
\pgfpathlineto{\pgfqpoint{1.818790in}{0.879030in}}%
\pgfpathlineto{\pgfqpoint{1.818712in}{0.873240in}}%
\pgfpathlineto{\pgfqpoint{1.822375in}{0.871354in}}%
\pgfpathlineto{\pgfqpoint{1.842357in}{0.871145in}}%
\pgfpathlineto{\pgfqpoint{1.842321in}{0.853020in}}%
\pgfpathclose%
\pgfusepath{fill}%
\end{pgfscope}%
\begin{pgfscope}%
\pgfpathrectangle{\pgfqpoint{0.100000in}{0.100000in}}{\pgfqpoint{3.007045in}{1.925000in}}%
\pgfusepath{clip}%
\pgfsetbuttcap%
\pgfsetmiterjoin%
\definecolor{currentfill}{rgb}{0.357647,0.720953,0.415071}%
\pgfsetfillcolor{currentfill}%
\pgfsetlinewidth{0.000000pt}%
\definecolor{currentstroke}{rgb}{0.000000,0.000000,0.000000}%
\pgfsetstrokecolor{currentstroke}%
\pgfsetstrokeopacity{0.000000}%
\pgfsetdash{}{0pt}%
\pgfpathmoveto{\pgfqpoint{2.994174in}{1.647565in}}%
\pgfpathlineto{\pgfqpoint{3.000400in}{1.646818in}}%
\pgfpathlineto{\pgfqpoint{2.997849in}{1.641810in}}%
\pgfpathlineto{\pgfqpoint{2.994174in}{1.647565in}}%
\pgfpathclose%
\pgfusepath{fill}%
\end{pgfscope}%
\begin{pgfscope}%
\pgfpathrectangle{\pgfqpoint{0.100000in}{0.100000in}}{\pgfqpoint{3.007045in}{1.925000in}}%
\pgfusepath{clip}%
\pgfsetbuttcap%
\pgfsetmiterjoin%
\definecolor{currentfill}{rgb}{0.357647,0.720953,0.415071}%
\pgfsetfillcolor{currentfill}%
\pgfsetlinewidth{0.000000pt}%
\definecolor{currentstroke}{rgb}{0.000000,0.000000,0.000000}%
\pgfsetstrokecolor{currentstroke}%
\pgfsetstrokeopacity{0.000000}%
\pgfsetdash{}{0pt}%
\pgfpathmoveto{\pgfqpoint{2.975595in}{1.630839in}}%
\pgfpathlineto{\pgfqpoint{2.978270in}{1.638556in}}%
\pgfpathlineto{\pgfqpoint{2.972759in}{1.647244in}}%
\pgfpathlineto{\pgfqpoint{2.969629in}{1.646428in}}%
\pgfpathlineto{\pgfqpoint{2.965671in}{1.652396in}}%
\pgfpathlineto{\pgfqpoint{2.962617in}{1.653501in}}%
\pgfpathlineto{\pgfqpoint{2.963183in}{1.661995in}}%
\pgfpathlineto{\pgfqpoint{2.962312in}{1.673903in}}%
\pgfpathlineto{\pgfqpoint{2.960139in}{1.680504in}}%
\pgfpathlineto{\pgfqpoint{2.966515in}{1.680747in}}%
\pgfpathlineto{\pgfqpoint{2.956012in}{1.701938in}}%
\pgfpathlineto{\pgfqpoint{2.945232in}{1.694510in}}%
\pgfpathlineto{\pgfqpoint{2.928393in}{1.726358in}}%
\pgfpathlineto{\pgfqpoint{2.930414in}{1.732814in}}%
\pgfpathlineto{\pgfqpoint{2.925595in}{1.735765in}}%
\pgfpathlineto{\pgfqpoint{2.926062in}{1.747145in}}%
\pgfpathlineto{\pgfqpoint{2.922774in}{1.754622in}}%
\pgfpathlineto{\pgfqpoint{2.915056in}{1.781252in}}%
\pgfpathlineto{\pgfqpoint{2.912260in}{1.792925in}}%
\pgfpathlineto{\pgfqpoint{2.951724in}{1.804308in}}%
\pgfpathlineto{\pgfqpoint{2.955175in}{1.793081in}}%
\pgfpathlineto{\pgfqpoint{2.972249in}{1.797222in}}%
\pgfpathlineto{\pgfqpoint{2.981003in}{1.769545in}}%
\pgfpathlineto{\pgfqpoint{2.987728in}{1.746399in}}%
\pgfpathlineto{\pgfqpoint{2.988875in}{1.746178in}}%
\pgfpathlineto{\pgfqpoint{3.003662in}{1.755298in}}%
\pgfpathlineto{\pgfqpoint{3.015744in}{1.733166in}}%
\pgfpathlineto{\pgfqpoint{3.010773in}{1.730803in}}%
\pgfpathlineto{\pgfqpoint{3.019840in}{1.712826in}}%
\pgfpathlineto{\pgfqpoint{3.014639in}{1.710091in}}%
\pgfpathlineto{\pgfqpoint{3.025401in}{1.689136in}}%
\pgfpathlineto{\pgfqpoint{3.031696in}{1.678779in}}%
\pgfpathlineto{\pgfqpoint{3.027762in}{1.672640in}}%
\pgfpathlineto{\pgfqpoint{3.022621in}{1.680721in}}%
\pgfpathlineto{\pgfqpoint{3.016503in}{1.677395in}}%
\pgfpathlineto{\pgfqpoint{3.023371in}{1.670491in}}%
\pgfpathlineto{\pgfqpoint{3.017940in}{1.661449in}}%
\pgfpathlineto{\pgfqpoint{3.012026in}{1.665983in}}%
\pgfpathlineto{\pgfqpoint{3.012283in}{1.670432in}}%
\pgfpathlineto{\pgfqpoint{3.008113in}{1.674442in}}%
\pgfpathlineto{\pgfqpoint{3.005865in}{1.669939in}}%
\pgfpathlineto{\pgfqpoint{3.005023in}{1.661197in}}%
\pgfpathlineto{\pgfqpoint{2.996716in}{1.662802in}}%
\pgfpathlineto{\pgfqpoint{2.986928in}{1.668315in}}%
\pgfpathlineto{\pgfqpoint{2.984644in}{1.665649in}}%
\pgfpathlineto{\pgfqpoint{2.987952in}{1.658494in}}%
\pgfpathlineto{\pgfqpoint{2.985414in}{1.647330in}}%
\pgfpathlineto{\pgfqpoint{2.987840in}{1.640394in}}%
\pgfpathlineto{\pgfqpoint{2.984073in}{1.632614in}}%
\pgfpathlineto{\pgfqpoint{2.975595in}{1.630839in}}%
\pgfpathclose%
\pgfusepath{fill}%
\end{pgfscope}%
\begin{pgfscope}%
\pgfpathrectangle{\pgfqpoint{0.100000in}{0.100000in}}{\pgfqpoint{3.007045in}{1.925000in}}%
\pgfusepath{clip}%
\pgfsetbuttcap%
\pgfsetmiterjoin%
\definecolor{currentfill}{rgb}{0.562860,0.819008,0.551511}%
\pgfsetfillcolor{currentfill}%
\pgfsetlinewidth{0.000000pt}%
\definecolor{currentstroke}{rgb}{0.000000,0.000000,0.000000}%
\pgfsetstrokecolor{currentstroke}%
\pgfsetstrokeopacity{0.000000}%
\pgfsetdash{}{0pt}%
\pgfpathmoveto{\pgfqpoint{1.566327in}{0.575160in}}%
\pgfpathlineto{\pgfqpoint{1.544968in}{0.575956in}}%
\pgfpathlineto{\pgfqpoint{1.546174in}{0.604074in}}%
\pgfpathlineto{\pgfqpoint{1.575754in}{0.602872in}}%
\pgfpathlineto{\pgfqpoint{1.576374in}{0.610043in}}%
\pgfpathlineto{\pgfqpoint{1.606454in}{0.609350in}}%
\pgfpathlineto{\pgfqpoint{1.610932in}{0.600948in}}%
\pgfpathlineto{\pgfqpoint{1.602998in}{0.592878in}}%
\pgfpathlineto{\pgfqpoint{1.593576in}{0.574175in}}%
\pgfpathlineto{\pgfqpoint{1.592881in}{0.569267in}}%
\pgfpathlineto{\pgfqpoint{1.576662in}{0.574910in}}%
\pgfpathlineto{\pgfqpoint{1.566327in}{0.575160in}}%
\pgfpathclose%
\pgfusepath{fill}%
\end{pgfscope}%
\begin{pgfscope}%
\pgfpathrectangle{\pgfqpoint{0.100000in}{0.100000in}}{\pgfqpoint{3.007045in}{1.925000in}}%
\pgfusepath{clip}%
\pgfsetbuttcap%
\pgfsetmiterjoin%
\definecolor{currentfill}{rgb}{0.496424,0.788005,0.496886}%
\pgfsetfillcolor{currentfill}%
\pgfsetlinewidth{0.000000pt}%
\definecolor{currentstroke}{rgb}{0.000000,0.000000,0.000000}%
\pgfsetstrokecolor{currentstroke}%
\pgfsetstrokeopacity{0.000000}%
\pgfsetdash{}{0pt}%
\pgfpathmoveto{\pgfqpoint{1.985200in}{0.741858in}}%
\pgfpathlineto{\pgfqpoint{1.989679in}{0.740203in}}%
\pgfpathlineto{\pgfqpoint{1.990829in}{0.734444in}}%
\pgfpathlineto{\pgfqpoint{1.986138in}{0.729844in}}%
\pgfpathlineto{\pgfqpoint{1.985938in}{0.725835in}}%
\pgfpathlineto{\pgfqpoint{1.991838in}{0.723800in}}%
\pgfpathlineto{\pgfqpoint{1.986961in}{0.717838in}}%
\pgfpathlineto{\pgfqpoint{1.988496in}{0.714303in}}%
\pgfpathlineto{\pgfqpoint{1.993233in}{0.713829in}}%
\pgfpathlineto{\pgfqpoint{1.988311in}{0.711075in}}%
\pgfpathlineto{\pgfqpoint{1.970224in}{0.710328in}}%
\pgfpathlineto{\pgfqpoint{1.970799in}{0.713279in}}%
\pgfpathlineto{\pgfqpoint{1.962387in}{0.712972in}}%
\pgfpathlineto{\pgfqpoint{1.959700in}{0.720849in}}%
\pgfpathlineto{\pgfqpoint{1.964704in}{0.732579in}}%
\pgfpathlineto{\pgfqpoint{1.967255in}{0.733621in}}%
\pgfpathlineto{\pgfqpoint{1.970228in}{0.741380in}}%
\pgfpathlineto{\pgfqpoint{1.985200in}{0.741858in}}%
\pgfpathclose%
\pgfusepath{fill}%
\end{pgfscope}%
\begin{pgfscope}%
\pgfpathrectangle{\pgfqpoint{0.100000in}{0.100000in}}{\pgfqpoint{3.007045in}{1.925000in}}%
\pgfusepath{clip}%
\pgfsetbuttcap%
\pgfsetmiterjoin%
\definecolor{currentfill}{rgb}{0.638970,0.854179,0.615240}%
\pgfsetfillcolor{currentfill}%
\pgfsetlinewidth{0.000000pt}%
\definecolor{currentstroke}{rgb}{0.000000,0.000000,0.000000}%
\pgfsetstrokecolor{currentstroke}%
\pgfsetstrokeopacity{0.000000}%
\pgfsetdash{}{0pt}%
\pgfpathmoveto{\pgfqpoint{1.760890in}{0.592713in}}%
\pgfpathlineto{\pgfqpoint{1.757986in}{0.598171in}}%
\pgfpathlineto{\pgfqpoint{1.753626in}{0.595060in}}%
\pgfpathlineto{\pgfqpoint{1.751572in}{0.571734in}}%
\pgfpathlineto{\pgfqpoint{1.737969in}{0.571808in}}%
\pgfpathlineto{\pgfqpoint{1.724776in}{0.579988in}}%
\pgfpathlineto{\pgfqpoint{1.723044in}{0.595566in}}%
\pgfpathlineto{\pgfqpoint{1.705582in}{0.593025in}}%
\pgfpathlineto{\pgfqpoint{1.701562in}{0.603092in}}%
\pgfpathlineto{\pgfqpoint{1.716755in}{0.610772in}}%
\pgfpathlineto{\pgfqpoint{1.728806in}{0.610818in}}%
\pgfpathlineto{\pgfqpoint{1.728210in}{0.613654in}}%
\pgfpathlineto{\pgfqpoint{1.731113in}{0.622526in}}%
\pgfpathlineto{\pgfqpoint{1.734214in}{0.624182in}}%
\pgfpathlineto{\pgfqpoint{1.730423in}{0.638033in}}%
\pgfpathlineto{\pgfqpoint{1.735433in}{0.640512in}}%
\pgfpathlineto{\pgfqpoint{1.756860in}{0.643794in}}%
\pgfpathlineto{\pgfqpoint{1.763851in}{0.642226in}}%
\pgfpathlineto{\pgfqpoint{1.766321in}{0.636474in}}%
\pgfpathlineto{\pgfqpoint{1.772144in}{0.632693in}}%
\pgfpathlineto{\pgfqpoint{1.779925in}{0.639399in}}%
\pgfpathlineto{\pgfqpoint{1.787706in}{0.634802in}}%
\pgfpathlineto{\pgfqpoint{1.800480in}{0.631756in}}%
\pgfpathlineto{\pgfqpoint{1.804231in}{0.625251in}}%
\pgfpathlineto{\pgfqpoint{1.810103in}{0.619771in}}%
\pgfpathlineto{\pgfqpoint{1.821877in}{0.611209in}}%
\pgfpathlineto{\pgfqpoint{1.803223in}{0.606687in}}%
\pgfpathlineto{\pgfqpoint{1.793823in}{0.610228in}}%
\pgfpathlineto{\pgfqpoint{1.786479in}{0.611427in}}%
\pgfpathlineto{\pgfqpoint{1.781271in}{0.614134in}}%
\pgfpathlineto{\pgfqpoint{1.775795in}{0.607167in}}%
\pgfpathlineto{\pgfqpoint{1.760890in}{0.592713in}}%
\pgfpathclose%
\pgfusepath{fill}%
\end{pgfscope}%
\begin{pgfscope}%
\pgfpathrectangle{\pgfqpoint{0.100000in}{0.100000in}}{\pgfqpoint{3.007045in}{1.925000in}}%
\pgfusepath{clip}%
\pgfsetbuttcap%
\pgfsetmiterjoin%
\definecolor{currentfill}{rgb}{0.426667,0.754787,0.448904}%
\pgfsetfillcolor{currentfill}%
\pgfsetlinewidth{0.000000pt}%
\definecolor{currentstroke}{rgb}{0.000000,0.000000,0.000000}%
\pgfsetstrokecolor{currentstroke}%
\pgfsetstrokeopacity{0.000000}%
\pgfsetdash{}{0pt}%
\pgfpathmoveto{\pgfqpoint{2.437163in}{0.979641in}}%
\pgfpathlineto{\pgfqpoint{2.431676in}{0.978202in}}%
\pgfpathlineto{\pgfqpoint{2.429667in}{0.984471in}}%
\pgfpathlineto{\pgfqpoint{2.416656in}{0.974681in}}%
\pgfpathlineto{\pgfqpoint{2.406942in}{0.983024in}}%
\pgfpathlineto{\pgfqpoint{2.400322in}{0.986669in}}%
\pgfpathlineto{\pgfqpoint{2.404364in}{0.992821in}}%
\pgfpathlineto{\pgfqpoint{2.399497in}{0.998207in}}%
\pgfpathlineto{\pgfqpoint{2.394844in}{0.994687in}}%
\pgfpathlineto{\pgfqpoint{2.392615in}{1.001420in}}%
\pgfpathlineto{\pgfqpoint{2.396925in}{1.003949in}}%
\pgfpathlineto{\pgfqpoint{2.399781in}{1.010130in}}%
\pgfpathlineto{\pgfqpoint{2.406434in}{1.016470in}}%
\pgfpathlineto{\pgfqpoint{2.407690in}{1.022099in}}%
\pgfpathlineto{\pgfqpoint{2.413702in}{1.026726in}}%
\pgfpathlineto{\pgfqpoint{2.416735in}{1.033753in}}%
\pgfpathlineto{\pgfqpoint{2.422280in}{1.037551in}}%
\pgfpathlineto{\pgfqpoint{2.434200in}{1.039055in}}%
\pgfpathlineto{\pgfqpoint{2.439622in}{1.028698in}}%
\pgfpathlineto{\pgfqpoint{2.450725in}{1.036221in}}%
\pgfpathlineto{\pgfqpoint{2.451559in}{1.040086in}}%
\pgfpathlineto{\pgfqpoint{2.461325in}{1.043806in}}%
\pgfpathlineto{\pgfqpoint{2.464192in}{1.048074in}}%
\pgfpathlineto{\pgfqpoint{2.469632in}{1.050830in}}%
\pgfpathlineto{\pgfqpoint{2.483475in}{1.055101in}}%
\pgfpathlineto{\pgfqpoint{2.490316in}{1.048764in}}%
\pgfpathlineto{\pgfqpoint{2.495062in}{1.040390in}}%
\pgfpathlineto{\pgfqpoint{2.481841in}{1.034709in}}%
\pgfpathlineto{\pgfqpoint{2.481176in}{1.031217in}}%
\pgfpathlineto{\pgfqpoint{2.476402in}{1.027509in}}%
\pgfpathlineto{\pgfqpoint{2.466964in}{1.026275in}}%
\pgfpathlineto{\pgfqpoint{2.462097in}{1.022483in}}%
\pgfpathlineto{\pgfqpoint{2.461449in}{1.016280in}}%
\pgfpathlineto{\pgfqpoint{2.458161in}{1.009811in}}%
\pgfpathlineto{\pgfqpoint{2.463756in}{1.002325in}}%
\pgfpathlineto{\pgfqpoint{2.460956in}{0.992711in}}%
\pgfpathlineto{\pgfqpoint{2.458453in}{0.990986in}}%
\pgfpathlineto{\pgfqpoint{2.450485in}{0.993442in}}%
\pgfpathlineto{\pgfqpoint{2.441498in}{0.987564in}}%
\pgfpathlineto{\pgfqpoint{2.437163in}{0.979641in}}%
\pgfpathclose%
\pgfusepath{fill}%
\end{pgfscope}%
\begin{pgfscope}%
\pgfpathrectangle{\pgfqpoint{0.100000in}{0.100000in}}{\pgfqpoint{3.007045in}{1.925000in}}%
\pgfusepath{clip}%
\pgfsetbuttcap%
\pgfsetmiterjoin%
\definecolor{currentfill}{rgb}{0.332549,0.708651,0.402768}%
\pgfsetfillcolor{currentfill}%
\pgfsetlinewidth{0.000000pt}%
\definecolor{currentstroke}{rgb}{0.000000,0.000000,0.000000}%
\pgfsetstrokecolor{currentstroke}%
\pgfsetstrokeopacity{0.000000}%
\pgfsetdash{}{0pt}%
\pgfpathmoveto{\pgfqpoint{2.131069in}{1.580709in}}%
\pgfpathlineto{\pgfqpoint{2.124022in}{1.562475in}}%
\pgfpathlineto{\pgfqpoint{2.118308in}{1.554107in}}%
\pgfpathlineto{\pgfqpoint{2.116985in}{1.541346in}}%
\pgfpathlineto{\pgfqpoint{2.111077in}{1.540159in}}%
\pgfpathlineto{\pgfqpoint{2.094036in}{1.542696in}}%
\pgfpathlineto{\pgfqpoint{2.093706in}{1.548428in}}%
\pgfpathlineto{\pgfqpoint{2.090120in}{1.553996in}}%
\pgfpathlineto{\pgfqpoint{2.084400in}{1.553797in}}%
\pgfpathlineto{\pgfqpoint{2.083707in}{1.565240in}}%
\pgfpathlineto{\pgfqpoint{2.078017in}{1.564986in}}%
\pgfpathlineto{\pgfqpoint{2.076489in}{1.587837in}}%
\pgfpathlineto{\pgfqpoint{2.064911in}{1.587042in}}%
\pgfpathlineto{\pgfqpoint{2.063792in}{1.604125in}}%
\pgfpathlineto{\pgfqpoint{2.048726in}{1.610846in}}%
\pgfpathlineto{\pgfqpoint{2.047301in}{1.632207in}}%
\pgfpathlineto{\pgfqpoint{2.087227in}{1.635040in}}%
\pgfpathlineto{\pgfqpoint{2.088086in}{1.623582in}}%
\pgfpathlineto{\pgfqpoint{2.111036in}{1.625398in}}%
\pgfpathlineto{\pgfqpoint{2.112389in}{1.608131in}}%
\pgfpathlineto{\pgfqpoint{2.123850in}{1.609038in}}%
\pgfpathlineto{\pgfqpoint{2.126263in}{1.603472in}}%
\pgfpathlineto{\pgfqpoint{2.128131in}{1.580457in}}%
\pgfpathlineto{\pgfqpoint{2.131069in}{1.580709in}}%
\pgfpathclose%
\pgfusepath{fill}%
\end{pgfscope}%
\begin{pgfscope}%
\pgfpathrectangle{\pgfqpoint{0.100000in}{0.100000in}}{\pgfqpoint{3.007045in}{1.925000in}}%
\pgfusepath{clip}%
\pgfsetbuttcap%
\pgfsetmiterjoin%
\definecolor{currentfill}{rgb}{0.623760,0.847428,0.601584}%
\pgfsetfillcolor{currentfill}%
\pgfsetlinewidth{0.000000pt}%
\definecolor{currentstroke}{rgb}{0.000000,0.000000,0.000000}%
\pgfsetstrokecolor{currentstroke}%
\pgfsetstrokeopacity{0.000000}%
\pgfsetdash{}{0pt}%
\pgfpathmoveto{\pgfqpoint{2.106178in}{0.880380in}}%
\pgfpathlineto{\pgfqpoint{2.085939in}{0.878995in}}%
\pgfpathlineto{\pgfqpoint{2.085014in}{0.905571in}}%
\pgfpathlineto{\pgfqpoint{2.090507in}{0.908195in}}%
\pgfpathlineto{\pgfqpoint{2.089884in}{0.925436in}}%
\pgfpathlineto{\pgfqpoint{2.074636in}{0.931335in}}%
\pgfpathlineto{\pgfqpoint{2.083454in}{0.947053in}}%
\pgfpathlineto{\pgfqpoint{2.083172in}{0.958841in}}%
\pgfpathlineto{\pgfqpoint{2.093449in}{0.959312in}}%
\pgfpathlineto{\pgfqpoint{2.098903in}{0.955455in}}%
\pgfpathlineto{\pgfqpoint{2.108338in}{0.951144in}}%
\pgfpathlineto{\pgfqpoint{2.108828in}{0.933214in}}%
\pgfpathlineto{\pgfqpoint{2.114171in}{0.933474in}}%
\pgfpathlineto{\pgfqpoint{2.114744in}{0.920210in}}%
\pgfpathlineto{\pgfqpoint{2.120897in}{0.915606in}}%
\pgfpathlineto{\pgfqpoint{2.125902in}{0.914003in}}%
\pgfpathlineto{\pgfqpoint{2.129095in}{0.910042in}}%
\pgfpathlineto{\pgfqpoint{2.129176in}{0.907553in}}%
\pgfpathlineto{\pgfqpoint{2.117945in}{0.906798in}}%
\pgfpathlineto{\pgfqpoint{2.114334in}{0.905115in}}%
\pgfpathlineto{\pgfqpoint{2.110613in}{0.897382in}}%
\pgfpathlineto{\pgfqpoint{2.107278in}{0.897155in}}%
\pgfpathlineto{\pgfqpoint{2.108156in}{0.880516in}}%
\pgfpathlineto{\pgfqpoint{2.106178in}{0.880380in}}%
\pgfpathclose%
\pgfusepath{fill}%
\end{pgfscope}%
\begin{pgfscope}%
\pgfpathrectangle{\pgfqpoint{0.100000in}{0.100000in}}{\pgfqpoint{3.007045in}{1.925000in}}%
\pgfusepath{clip}%
\pgfsetbuttcap%
\pgfsetmiterjoin%
\definecolor{currentfill}{rgb}{0.376471,0.730181,0.424298}%
\pgfsetfillcolor{currentfill}%
\pgfsetlinewidth{0.000000pt}%
\definecolor{currentstroke}{rgb}{0.000000,0.000000,0.000000}%
\pgfsetstrokecolor{currentstroke}%
\pgfsetstrokeopacity{0.000000}%
\pgfsetdash{}{0pt}%
\pgfpathmoveto{\pgfqpoint{1.650255in}{1.108902in}}%
\pgfpathlineto{\pgfqpoint{1.621825in}{1.109765in}}%
\pgfpathlineto{\pgfqpoint{1.622153in}{1.126974in}}%
\pgfpathlineto{\pgfqpoint{1.593668in}{1.127954in}}%
\pgfpathlineto{\pgfqpoint{1.594074in}{1.145204in}}%
\pgfpathlineto{\pgfqpoint{1.594289in}{1.150925in}}%
\pgfpathlineto{\pgfqpoint{1.622833in}{1.149886in}}%
\pgfpathlineto{\pgfqpoint{1.623055in}{1.155630in}}%
\pgfpathlineto{\pgfqpoint{1.651571in}{1.154782in}}%
\pgfpathlineto{\pgfqpoint{1.650255in}{1.108902in}}%
\pgfpathclose%
\pgfusepath{fill}%
\end{pgfscope}%
\begin{pgfscope}%
\pgfpathrectangle{\pgfqpoint{0.100000in}{0.100000in}}{\pgfqpoint{3.007045in}{1.925000in}}%
\pgfusepath{clip}%
\pgfsetbuttcap%
\pgfsetmiterjoin%
\definecolor{currentfill}{rgb}{0.868051,0.948789,0.846444}%
\pgfsetfillcolor{currentfill}%
\pgfsetlinewidth{0.000000pt}%
\definecolor{currentstroke}{rgb}{0.000000,0.000000,0.000000}%
\pgfsetstrokecolor{currentstroke}%
\pgfsetstrokeopacity{0.000000}%
\pgfsetdash{}{0pt}%
\pgfpathmoveto{\pgfqpoint{2.140155in}{0.982191in}}%
\pgfpathlineto{\pgfqpoint{2.141030in}{0.985137in}}%
\pgfpathlineto{\pgfqpoint{2.135867in}{0.998378in}}%
\pgfpathlineto{\pgfqpoint{2.132817in}{1.006098in}}%
\pgfpathlineto{\pgfqpoint{2.146362in}{1.013286in}}%
\pgfpathlineto{\pgfqpoint{2.152882in}{1.014762in}}%
\pgfpathlineto{\pgfqpoint{2.155476in}{1.018069in}}%
\pgfpathlineto{\pgfqpoint{2.156379in}{1.026602in}}%
\pgfpathlineto{\pgfqpoint{2.164956in}{1.024377in}}%
\pgfpathlineto{\pgfqpoint{2.174509in}{1.028592in}}%
\pgfpathlineto{\pgfqpoint{2.179536in}{1.021312in}}%
\pgfpathlineto{\pgfqpoint{2.189732in}{1.023554in}}%
\pgfpathlineto{\pgfqpoint{2.191739in}{0.996087in}}%
\pgfpathlineto{\pgfqpoint{2.188886in}{0.995787in}}%
\pgfpathlineto{\pgfqpoint{2.189194in}{0.983871in}}%
\pgfpathlineto{\pgfqpoint{2.183256in}{0.977798in}}%
\pgfpathlineto{\pgfqpoint{2.181638in}{0.973954in}}%
\pgfpathlineto{\pgfqpoint{2.169549in}{0.973787in}}%
\pgfpathlineto{\pgfqpoint{2.169480in}{0.968102in}}%
\pgfpathlineto{\pgfqpoint{2.166450in}{0.964466in}}%
\pgfpathlineto{\pgfqpoint{2.158552in}{0.966327in}}%
\pgfpathlineto{\pgfqpoint{2.146928in}{0.965917in}}%
\pgfpathlineto{\pgfqpoint{2.146159in}{0.970611in}}%
\pgfpathlineto{\pgfqpoint{2.140247in}{0.978276in}}%
\pgfpathlineto{\pgfqpoint{2.140155in}{0.982191in}}%
\pgfpathclose%
\pgfusepath{fill}%
\end{pgfscope}%
\begin{pgfscope}%
\pgfpathrectangle{\pgfqpoint{0.100000in}{0.100000in}}{\pgfqpoint{3.007045in}{1.925000in}}%
\pgfusepath{clip}%
\pgfsetbuttcap%
\pgfsetmiterjoin%
\definecolor{currentfill}{rgb}{0.690396,0.875832,0.665313}%
\pgfsetfillcolor{currentfill}%
\pgfsetlinewidth{0.000000pt}%
\definecolor{currentstroke}{rgb}{0.000000,0.000000,0.000000}%
\pgfsetstrokecolor{currentstroke}%
\pgfsetstrokeopacity{0.000000}%
\pgfsetdash{}{0pt}%
\pgfpathmoveto{\pgfqpoint{2.721216in}{1.279405in}}%
\pgfpathlineto{\pgfqpoint{2.739138in}{1.283230in}}%
\pgfpathlineto{\pgfqpoint{2.742511in}{1.290165in}}%
\pgfpathlineto{\pgfqpoint{2.746343in}{1.292558in}}%
\pgfpathlineto{\pgfqpoint{2.753598in}{1.293002in}}%
\pgfpathlineto{\pgfqpoint{2.753954in}{1.289095in}}%
\pgfpathlineto{\pgfqpoint{2.751124in}{1.279964in}}%
\pgfpathlineto{\pgfqpoint{2.754264in}{1.276110in}}%
\pgfpathlineto{\pgfqpoint{2.754436in}{1.271463in}}%
\pgfpathlineto{\pgfqpoint{2.751790in}{1.268575in}}%
\pgfpathlineto{\pgfqpoint{2.757050in}{1.263166in}}%
\pgfpathlineto{\pgfqpoint{2.751195in}{1.256851in}}%
\pgfpathlineto{\pgfqpoint{2.745705in}{1.255995in}}%
\pgfpathlineto{\pgfqpoint{2.744257in}{1.261114in}}%
\pgfpathlineto{\pgfqpoint{2.738533in}{1.258792in}}%
\pgfpathlineto{\pgfqpoint{2.730272in}{1.258328in}}%
\pgfpathlineto{\pgfqpoint{2.732034in}{1.262134in}}%
\pgfpathlineto{\pgfqpoint{2.731538in}{1.268601in}}%
\pgfpathlineto{\pgfqpoint{2.726638in}{1.268556in}}%
\pgfpathlineto{\pgfqpoint{2.716371in}{1.278367in}}%
\pgfpathlineto{\pgfqpoint{2.721216in}{1.279405in}}%
\pgfpathclose%
\pgfusepath{fill}%
\end{pgfscope}%
\begin{pgfscope}%
\pgfpathrectangle{\pgfqpoint{0.100000in}{0.100000in}}{\pgfqpoint{3.007045in}{1.925000in}}%
\pgfusepath{clip}%
\pgfsetbuttcap%
\pgfsetmiterjoin%
\definecolor{currentfill}{rgb}{0.451765,0.767090,0.461207}%
\pgfsetfillcolor{currentfill}%
\pgfsetlinewidth{0.000000pt}%
\definecolor{currentstroke}{rgb}{0.000000,0.000000,0.000000}%
\pgfsetstrokecolor{currentstroke}%
\pgfsetstrokeopacity{0.000000}%
\pgfsetdash{}{0pt}%
\pgfpathmoveto{\pgfqpoint{1.509458in}{1.149192in}}%
\pgfpathlineto{\pgfqpoint{1.508618in}{1.149222in}}%
\pgfpathlineto{\pgfqpoint{1.510385in}{1.177845in}}%
\pgfpathlineto{\pgfqpoint{1.537445in}{1.176396in}}%
\pgfpathlineto{\pgfqpoint{1.567058in}{1.175010in}}%
\pgfpathlineto{\pgfqpoint{1.565600in}{1.146405in}}%
\pgfpathlineto{\pgfqpoint{1.509458in}{1.149192in}}%
\pgfpathclose%
\pgfusepath{fill}%
\end{pgfscope}%
\begin{pgfscope}%
\pgfpathrectangle{\pgfqpoint{0.100000in}{0.100000in}}{\pgfqpoint{3.007045in}{1.925000in}}%
\pgfusepath{clip}%
\pgfsetbuttcap%
\pgfsetmiterjoin%
\definecolor{currentfill}{rgb}{0.704421,0.881738,0.678970}%
\pgfsetfillcolor{currentfill}%
\pgfsetlinewidth{0.000000pt}%
\definecolor{currentstroke}{rgb}{0.000000,0.000000,0.000000}%
\pgfsetstrokecolor{currentstroke}%
\pgfsetstrokeopacity{0.000000}%
\pgfsetdash{}{0pt}%
\pgfpathmoveto{\pgfqpoint{1.596180in}{1.004385in}}%
\pgfpathlineto{\pgfqpoint{1.596980in}{1.029821in}}%
\pgfpathlineto{\pgfqpoint{1.590957in}{1.030045in}}%
\pgfpathlineto{\pgfqpoint{1.591838in}{1.053004in}}%
\pgfpathlineto{\pgfqpoint{1.631737in}{1.051813in}}%
\pgfpathlineto{\pgfqpoint{1.631910in}{1.063521in}}%
\pgfpathlineto{\pgfqpoint{1.660467in}{1.062737in}}%
\pgfpathlineto{\pgfqpoint{1.660705in}{1.074262in}}%
\pgfpathlineto{\pgfqpoint{1.676880in}{1.073769in}}%
\pgfpathlineto{\pgfqpoint{1.693377in}{1.073520in}}%
\pgfpathlineto{\pgfqpoint{1.692552in}{1.033341in}}%
\pgfpathlineto{\pgfqpoint{1.659741in}{1.033948in}}%
\pgfpathlineto{\pgfqpoint{1.659290in}{1.002553in}}%
\pgfpathlineto{\pgfqpoint{1.642737in}{1.002945in}}%
\pgfpathlineto{\pgfqpoint{1.596180in}{1.004385in}}%
\pgfpathclose%
\pgfusepath{fill}%
\end{pgfscope}%
\begin{pgfscope}%
\pgfpathrectangle{\pgfqpoint{0.100000in}{0.100000in}}{\pgfqpoint{3.007045in}{1.925000in}}%
\pgfusepath{clip}%
\pgfsetbuttcap%
\pgfsetmiterjoin%
\definecolor{currentfill}{rgb}{0.601615,0.837093,0.583376}%
\pgfsetfillcolor{currentfill}%
\pgfsetlinewidth{0.000000pt}%
\definecolor{currentstroke}{rgb}{0.000000,0.000000,0.000000}%
\pgfsetstrokecolor{currentstroke}%
\pgfsetstrokeopacity{0.000000}%
\pgfsetdash{}{0pt}%
\pgfpathmoveto{\pgfqpoint{1.526290in}{0.791703in}}%
\pgfpathlineto{\pgfqpoint{1.525327in}{0.794045in}}%
\pgfpathlineto{\pgfqpoint{1.514372in}{0.798412in}}%
\pgfpathlineto{\pgfqpoint{1.495136in}{0.800124in}}%
\pgfpathlineto{\pgfqpoint{1.496660in}{0.826118in}}%
\pgfpathlineto{\pgfqpoint{1.499337in}{0.825611in}}%
\pgfpathlineto{\pgfqpoint{1.500565in}{0.847796in}}%
\pgfpathlineto{\pgfqpoint{1.504435in}{0.848612in}}%
\pgfpathlineto{\pgfqpoint{1.516379in}{0.834892in}}%
\pgfpathlineto{\pgfqpoint{1.521539in}{0.834387in}}%
\pgfpathlineto{\pgfqpoint{1.526280in}{0.836772in}}%
\pgfpathlineto{\pgfqpoint{1.531947in}{0.833744in}}%
\pgfpathlineto{\pgfqpoint{1.533732in}{0.839225in}}%
\pgfpathlineto{\pgfqpoint{1.542619in}{0.830859in}}%
\pgfpathlineto{\pgfqpoint{1.543414in}{0.822716in}}%
\pgfpathlineto{\pgfqpoint{1.556351in}{0.822025in}}%
\pgfpathlineto{\pgfqpoint{1.555230in}{0.797022in}}%
\pgfpathlineto{\pgfqpoint{1.526556in}{0.798349in}}%
\pgfpathlineto{\pgfqpoint{1.526290in}{0.791703in}}%
\pgfpathclose%
\pgfusepath{fill}%
\end{pgfscope}%
\begin{pgfscope}%
\pgfpathrectangle{\pgfqpoint{0.100000in}{0.100000in}}{\pgfqpoint{3.007045in}{1.925000in}}%
\pgfusepath{clip}%
\pgfsetbuttcap%
\pgfsetmiterjoin%
\definecolor{currentfill}{rgb}{0.823760,0.931073,0.799200}%
\pgfsetfillcolor{currentfill}%
\pgfsetlinewidth{0.000000pt}%
\definecolor{currentstroke}{rgb}{0.000000,0.000000,0.000000}%
\pgfsetstrokecolor{currentstroke}%
\pgfsetstrokeopacity{0.000000}%
\pgfsetdash{}{0pt}%
\pgfpathmoveto{\pgfqpoint{0.819433in}{0.344425in}}%
\pgfpathlineto{\pgfqpoint{0.816099in}{0.346069in}}%
\pgfpathlineto{\pgfqpoint{0.814545in}{0.348123in}}%
\pgfpathlineto{\pgfqpoint{0.812387in}{0.348661in}}%
\pgfpathlineto{\pgfqpoint{0.812849in}{0.351051in}}%
\pgfpathlineto{\pgfqpoint{0.812080in}{0.351861in}}%
\pgfpathlineto{\pgfqpoint{0.810038in}{0.350072in}}%
\pgfpathlineto{\pgfqpoint{0.808207in}{0.349391in}}%
\pgfpathlineto{\pgfqpoint{0.806954in}{0.350824in}}%
\pgfpathlineto{\pgfqpoint{0.805473in}{0.350277in}}%
\pgfpathlineto{\pgfqpoint{0.804388in}{0.347206in}}%
\pgfpathlineto{\pgfqpoint{0.802543in}{0.351436in}}%
\pgfpathlineto{\pgfqpoint{0.802043in}{0.348092in}}%
\pgfpathlineto{\pgfqpoint{0.798587in}{0.347926in}}%
\pgfpathlineto{\pgfqpoint{0.798087in}{0.346708in}}%
\pgfpathlineto{\pgfqpoint{0.795113in}{0.346802in}}%
\pgfpathlineto{\pgfqpoint{0.791694in}{0.347913in}}%
\pgfpathlineto{\pgfqpoint{0.786584in}{0.347123in}}%
\pgfpathlineto{\pgfqpoint{0.784356in}{0.348068in}}%
\pgfpathlineto{\pgfqpoint{0.783495in}{0.346619in}}%
\pgfpathlineto{\pgfqpoint{0.781504in}{0.347865in}}%
\pgfpathlineto{\pgfqpoint{0.781399in}{0.349966in}}%
\pgfpathlineto{\pgfqpoint{0.779304in}{0.349025in}}%
\pgfpathlineto{\pgfqpoint{0.776366in}{0.349945in}}%
\pgfpathlineto{\pgfqpoint{0.776139in}{0.354814in}}%
\pgfpathlineto{\pgfqpoint{0.778619in}{0.356270in}}%
\pgfpathlineto{\pgfqpoint{0.781930in}{0.355808in}}%
\pgfpathlineto{\pgfqpoint{0.783469in}{0.354208in}}%
\pgfpathlineto{\pgfqpoint{0.785254in}{0.354993in}}%
\pgfpathlineto{\pgfqpoint{0.787547in}{0.354040in}}%
\pgfpathlineto{\pgfqpoint{0.788336in}{0.355429in}}%
\pgfpathlineto{\pgfqpoint{0.793181in}{0.356662in}}%
\pgfpathlineto{\pgfqpoint{0.790991in}{0.357629in}}%
\pgfpathlineto{\pgfqpoint{0.788388in}{0.357382in}}%
\pgfpathlineto{\pgfqpoint{0.785835in}{0.356265in}}%
\pgfpathlineto{\pgfqpoint{0.783898in}{0.359777in}}%
\pgfpathlineto{\pgfqpoint{0.783523in}{0.362343in}}%
\pgfpathlineto{\pgfqpoint{0.789148in}{0.366853in}}%
\pgfpathlineto{\pgfqpoint{0.793181in}{0.368023in}}%
\pgfpathlineto{\pgfqpoint{0.797910in}{0.370814in}}%
\pgfpathlineto{\pgfqpoint{0.800516in}{0.373276in}}%
\pgfpathlineto{\pgfqpoint{0.801383in}{0.377874in}}%
\pgfpathlineto{\pgfqpoint{0.803675in}{0.378105in}}%
\pgfpathlineto{\pgfqpoint{0.805964in}{0.377296in}}%
\pgfpathlineto{\pgfqpoint{0.811186in}{0.377817in}}%
\pgfpathlineto{\pgfqpoint{0.816257in}{0.377436in}}%
\pgfpathlineto{\pgfqpoint{0.816031in}{0.374807in}}%
\pgfpathlineto{\pgfqpoint{0.817543in}{0.372687in}}%
\pgfpathlineto{\pgfqpoint{0.822055in}{0.372655in}}%
\pgfpathlineto{\pgfqpoint{0.820708in}{0.377871in}}%
\pgfpathlineto{\pgfqpoint{0.822622in}{0.379076in}}%
\pgfpathlineto{\pgfqpoint{0.818515in}{0.382126in}}%
\pgfpathlineto{\pgfqpoint{0.817100in}{0.384096in}}%
\pgfpathlineto{\pgfqpoint{0.812559in}{0.384716in}}%
\pgfpathlineto{\pgfqpoint{0.809430in}{0.383487in}}%
\pgfpathlineto{\pgfqpoint{0.804683in}{0.384619in}}%
\pgfpathlineto{\pgfqpoint{0.800098in}{0.383859in}}%
\pgfpathlineto{\pgfqpoint{0.798506in}{0.380328in}}%
\pgfpathlineto{\pgfqpoint{0.797580in}{0.381676in}}%
\pgfpathlineto{\pgfqpoint{0.795111in}{0.381646in}}%
\pgfpathlineto{\pgfqpoint{0.790374in}{0.380576in}}%
\pgfpathlineto{\pgfqpoint{0.787763in}{0.377740in}}%
\pgfpathlineto{\pgfqpoint{0.785721in}{0.376868in}}%
\pgfpathlineto{\pgfqpoint{0.783285in}{0.376730in}}%
\pgfpathlineto{\pgfqpoint{0.781683in}{0.377953in}}%
\pgfpathlineto{\pgfqpoint{0.781141in}{0.373732in}}%
\pgfpathlineto{\pgfqpoint{0.777419in}{0.371708in}}%
\pgfpathlineto{\pgfqpoint{0.775658in}{0.371988in}}%
\pgfpathlineto{\pgfqpoint{0.774292in}{0.373277in}}%
\pgfpathlineto{\pgfqpoint{0.771887in}{0.370547in}}%
\pgfpathlineto{\pgfqpoint{0.769799in}{0.370122in}}%
\pgfpathlineto{\pgfqpoint{0.764543in}{0.372414in}}%
\pgfpathlineto{\pgfqpoint{0.763348in}{0.371469in}}%
\pgfpathlineto{\pgfqpoint{0.761565in}{0.371604in}}%
\pgfpathlineto{\pgfqpoint{0.760574in}{0.369853in}}%
\pgfpathlineto{\pgfqpoint{0.756927in}{0.371271in}}%
\pgfpathlineto{\pgfqpoint{0.753478in}{0.368527in}}%
\pgfpathlineto{\pgfqpoint{0.751025in}{0.367939in}}%
\pgfpathlineto{\pgfqpoint{0.752269in}{0.365262in}}%
\pgfpathlineto{\pgfqpoint{0.755024in}{0.362772in}}%
\pgfpathlineto{\pgfqpoint{0.757010in}{0.357656in}}%
\pgfpathlineto{\pgfqpoint{0.757017in}{0.355194in}}%
\pgfpathlineto{\pgfqpoint{0.755001in}{0.356698in}}%
\pgfpathlineto{\pgfqpoint{0.753292in}{0.354357in}}%
\pgfpathlineto{\pgfqpoint{0.749665in}{0.357072in}}%
\pgfpathlineto{\pgfqpoint{0.748423in}{0.355398in}}%
\pgfpathlineto{\pgfqpoint{0.743549in}{0.359090in}}%
\pgfpathlineto{\pgfqpoint{0.740274in}{0.361603in}}%
\pgfpathlineto{\pgfqpoint{0.741011in}{0.363689in}}%
\pgfpathlineto{\pgfqpoint{0.746018in}{0.370169in}}%
\pgfpathlineto{\pgfqpoint{0.745512in}{0.370591in}}%
\pgfpathlineto{\pgfqpoint{0.748076in}{0.373959in}}%
\pgfpathlineto{\pgfqpoint{0.751397in}{0.371451in}}%
\pgfpathlineto{\pgfqpoint{0.753894in}{0.374816in}}%
\pgfpathlineto{\pgfqpoint{0.756787in}{0.372685in}}%
\pgfpathlineto{\pgfqpoint{0.758045in}{0.374400in}}%
\pgfpathlineto{\pgfqpoint{0.759707in}{0.373154in}}%
\pgfpathlineto{\pgfqpoint{0.762258in}{0.376547in}}%
\pgfpathlineto{\pgfqpoint{0.763938in}{0.375291in}}%
\pgfpathlineto{\pgfqpoint{0.765233in}{0.376990in}}%
\pgfpathlineto{\pgfqpoint{0.766473in}{0.376050in}}%
\pgfpathlineto{\pgfqpoint{0.771354in}{0.382565in}}%
\pgfpathlineto{\pgfqpoint{0.772579in}{0.381674in}}%
\pgfpathlineto{\pgfqpoint{0.777700in}{0.388433in}}%
\pgfpathlineto{\pgfqpoint{0.778986in}{0.387487in}}%
\pgfpathlineto{\pgfqpoint{0.784096in}{0.394320in}}%
\pgfpathlineto{\pgfqpoint{0.784950in}{0.396276in}}%
\pgfpathlineto{\pgfqpoint{0.788663in}{0.401187in}}%
\pgfpathlineto{\pgfqpoint{0.789615in}{0.403627in}}%
\pgfpathlineto{\pgfqpoint{0.792854in}{0.407923in}}%
\pgfpathlineto{\pgfqpoint{0.796295in}{0.405189in}}%
\pgfpathlineto{\pgfqpoint{0.796605in}{0.403352in}}%
\pgfpathlineto{\pgfqpoint{0.809857in}{0.421160in}}%
\pgfpathlineto{\pgfqpoint{0.816153in}{0.429462in}}%
\pgfpathlineto{\pgfqpoint{0.826027in}{0.421875in}}%
\pgfpathlineto{\pgfqpoint{0.827109in}{0.423257in}}%
\pgfpathlineto{\pgfqpoint{0.847814in}{0.419834in}}%
\pgfpathlineto{\pgfqpoint{0.856121in}{0.418510in}}%
\pgfpathlineto{\pgfqpoint{0.869886in}{0.409139in}}%
\pgfpathlineto{\pgfqpoint{0.871795in}{0.412013in}}%
\pgfpathlineto{\pgfqpoint{0.875033in}{0.409830in}}%
\pgfpathlineto{\pgfqpoint{0.886359in}{0.402751in}}%
\pgfpathlineto{\pgfqpoint{0.880633in}{0.393554in}}%
\pgfpathlineto{\pgfqpoint{0.878779in}{0.389457in}}%
\pgfpathlineto{\pgfqpoint{0.870796in}{0.376637in}}%
\pgfpathlineto{\pgfqpoint{0.865448in}{0.380135in}}%
\pgfpathlineto{\pgfqpoint{0.864498in}{0.378218in}}%
\pgfpathlineto{\pgfqpoint{0.855608in}{0.363995in}}%
\pgfpathlineto{\pgfqpoint{0.853061in}{0.365671in}}%
\pgfpathlineto{\pgfqpoint{0.852394in}{0.364648in}}%
\pgfpathlineto{\pgfqpoint{0.840140in}{0.372751in}}%
\pgfpathlineto{\pgfqpoint{0.836391in}{0.367548in}}%
\pgfpathlineto{\pgfqpoint{0.832191in}{0.361126in}}%
\pgfpathlineto{\pgfqpoint{0.829442in}{0.362834in}}%
\pgfpathlineto{\pgfqpoint{0.827922in}{0.360561in}}%
\pgfpathlineto{\pgfqpoint{0.828821in}{0.359977in}}%
\pgfpathlineto{\pgfqpoint{0.824494in}{0.353761in}}%
\pgfpathlineto{\pgfqpoint{0.825526in}{0.353057in}}%
\pgfpathlineto{\pgfqpoint{0.821535in}{0.347132in}}%
\pgfpathlineto{\pgfqpoint{0.819433in}{0.344425in}}%
\pgfpathclose%
\pgfusepath{fill}%
\end{pgfscope}%
\begin{pgfscope}%
\pgfpathrectangle{\pgfqpoint{0.100000in}{0.100000in}}{\pgfqpoint{3.007045in}{1.925000in}}%
\pgfusepath{clip}%
\pgfsetbuttcap%
\pgfsetmiterjoin%
\definecolor{currentfill}{rgb}{0.307451,0.696348,0.390465}%
\pgfsetfillcolor{currentfill}%
\pgfsetlinewidth{0.000000pt}%
\definecolor{currentstroke}{rgb}{0.000000,0.000000,0.000000}%
\pgfsetstrokecolor{currentstroke}%
\pgfsetstrokeopacity{0.000000}%
\pgfsetdash{}{0pt}%
\pgfpathmoveto{\pgfqpoint{2.434200in}{1.039055in}}%
\pgfpathlineto{\pgfqpoint{2.439575in}{1.041742in}}%
\pgfpathlineto{\pgfqpoint{2.437440in}{1.045801in}}%
\pgfpathlineto{\pgfqpoint{2.446669in}{1.052394in}}%
\pgfpathlineto{\pgfqpoint{2.440019in}{1.062167in}}%
\pgfpathlineto{\pgfqpoint{2.435154in}{1.067105in}}%
\pgfpathlineto{\pgfqpoint{2.450852in}{1.085445in}}%
\pgfpathlineto{\pgfqpoint{2.453185in}{1.084073in}}%
\pgfpathlineto{\pgfqpoint{2.452938in}{1.079250in}}%
\pgfpathlineto{\pgfqpoint{2.455254in}{1.073522in}}%
\pgfpathlineto{\pgfqpoint{2.462901in}{1.070185in}}%
\pgfpathlineto{\pgfqpoint{2.469086in}{1.065710in}}%
\pgfpathlineto{\pgfqpoint{2.475445in}{1.067081in}}%
\pgfpathlineto{\pgfqpoint{2.483702in}{1.075698in}}%
\pgfpathlineto{\pgfqpoint{2.487921in}{1.088350in}}%
\pgfpathlineto{\pgfqpoint{2.489686in}{1.089390in}}%
\pgfpathlineto{\pgfqpoint{2.490450in}{1.093853in}}%
\pgfpathlineto{\pgfqpoint{2.495503in}{1.095488in}}%
\pgfpathlineto{\pgfqpoint{2.509239in}{1.087027in}}%
\pgfpathlineto{\pgfqpoint{2.510605in}{1.081728in}}%
\pgfpathlineto{\pgfqpoint{2.502786in}{1.075737in}}%
\pgfpathlineto{\pgfqpoint{2.512374in}{1.068702in}}%
\pgfpathlineto{\pgfqpoint{2.508760in}{1.065976in}}%
\pgfpathlineto{\pgfqpoint{2.511594in}{1.060569in}}%
\pgfpathlineto{\pgfqpoint{2.521063in}{1.052157in}}%
\pgfpathlineto{\pgfqpoint{2.505363in}{1.044619in}}%
\pgfpathlineto{\pgfqpoint{2.503517in}{1.042080in}}%
\pgfpathlineto{\pgfqpoint{2.495062in}{1.040390in}}%
\pgfpathlineto{\pgfqpoint{2.490316in}{1.048764in}}%
\pgfpathlineto{\pgfqpoint{2.483475in}{1.055101in}}%
\pgfpathlineto{\pgfqpoint{2.469632in}{1.050830in}}%
\pgfpathlineto{\pgfqpoint{2.464192in}{1.048074in}}%
\pgfpathlineto{\pgfqpoint{2.461325in}{1.043806in}}%
\pgfpathlineto{\pgfqpoint{2.451559in}{1.040086in}}%
\pgfpathlineto{\pgfqpoint{2.450725in}{1.036221in}}%
\pgfpathlineto{\pgfqpoint{2.439622in}{1.028698in}}%
\pgfpathlineto{\pgfqpoint{2.434200in}{1.039055in}}%
\pgfpathclose%
\pgfusepath{fill}%
\end{pgfscope}%
\begin{pgfscope}%
\pgfpathrectangle{\pgfqpoint{0.100000in}{0.100000in}}{\pgfqpoint{3.007045in}{1.925000in}}%
\pgfusepath{clip}%
\pgfsetbuttcap%
\pgfsetmiterjoin%
\definecolor{currentfill}{rgb}{0.420392,0.751711,0.445829}%
\pgfsetfillcolor{currentfill}%
\pgfsetlinewidth{0.000000pt}%
\definecolor{currentstroke}{rgb}{0.000000,0.000000,0.000000}%
\pgfsetstrokecolor{currentstroke}%
\pgfsetstrokeopacity{0.000000}%
\pgfsetdash{}{0pt}%
\pgfpathmoveto{\pgfqpoint{1.570193in}{1.688446in}}%
\pgfpathlineto{\pgfqpoint{1.569042in}{1.665415in}}%
\pgfpathlineto{\pgfqpoint{1.570614in}{1.665351in}}%
\pgfpathlineto{\pgfqpoint{1.569331in}{1.642166in}}%
\pgfpathlineto{\pgfqpoint{1.540553in}{1.643789in}}%
\pgfpathlineto{\pgfqpoint{1.542129in}{1.666889in}}%
\pgfpathlineto{\pgfqpoint{1.540344in}{1.666994in}}%
\pgfpathlineto{\pgfqpoint{1.541684in}{1.689945in}}%
\pgfpathlineto{\pgfqpoint{1.570193in}{1.688446in}}%
\pgfpathclose%
\pgfusepath{fill}%
\end{pgfscope}%
\begin{pgfscope}%
\pgfpathrectangle{\pgfqpoint{0.100000in}{0.100000in}}{\pgfqpoint{3.007045in}{1.925000in}}%
\pgfusepath{clip}%
\pgfsetbuttcap%
\pgfsetmiterjoin%
\definecolor{currentfill}{rgb}{0.751173,0.901423,0.724491}%
\pgfsetfillcolor{currentfill}%
\pgfsetlinewidth{0.000000pt}%
\definecolor{currentstroke}{rgb}{0.000000,0.000000,0.000000}%
\pgfsetstrokecolor{currentstroke}%
\pgfsetstrokeopacity{0.000000}%
\pgfsetdash{}{0pt}%
\pgfpathmoveto{\pgfqpoint{1.456990in}{0.978623in}}%
\pgfpathlineto{\pgfqpoint{1.507272in}{0.975589in}}%
\pgfpathlineto{\pgfqpoint{1.505836in}{0.946347in}}%
\pgfpathlineto{\pgfqpoint{1.477103in}{0.948069in}}%
\pgfpathlineto{\pgfqpoint{1.419469in}{0.951860in}}%
\pgfpathlineto{\pgfqpoint{1.421577in}{0.981072in}}%
\pgfpathlineto{\pgfqpoint{1.456990in}{0.978623in}}%
\pgfpathclose%
\pgfusepath{fill}%
\end{pgfscope}%
\begin{pgfscope}%
\pgfpathrectangle{\pgfqpoint{0.100000in}{0.100000in}}{\pgfqpoint{3.007045in}{1.925000in}}%
\pgfusepath{clip}%
\pgfsetbuttcap%
\pgfsetmiterjoin%
\definecolor{currentfill}{rgb}{0.704421,0.881738,0.678970}%
\pgfsetfillcolor{currentfill}%
\pgfsetlinewidth{0.000000pt}%
\definecolor{currentstroke}{rgb}{0.000000,0.000000,0.000000}%
\pgfsetstrokecolor{currentstroke}%
\pgfsetstrokeopacity{0.000000}%
\pgfsetdash{}{0pt}%
\pgfpathmoveto{\pgfqpoint{2.196663in}{1.378800in}}%
\pgfpathlineto{\pgfqpoint{2.198102in}{1.398160in}}%
\pgfpathlineto{\pgfqpoint{2.196389in}{1.411349in}}%
\pgfpathlineto{\pgfqpoint{2.193000in}{1.423406in}}%
\pgfpathlineto{\pgfqpoint{2.183445in}{1.440623in}}%
\pgfpathlineto{\pgfqpoint{2.177072in}{1.454810in}}%
\pgfpathlineto{\pgfqpoint{2.176584in}{1.460340in}}%
\pgfpathlineto{\pgfqpoint{2.180375in}{1.470113in}}%
\pgfpathlineto{\pgfqpoint{2.199018in}{1.471645in}}%
\pgfpathlineto{\pgfqpoint{2.221626in}{1.473931in}}%
\pgfpathlineto{\pgfqpoint{2.224008in}{1.451122in}}%
\pgfpathlineto{\pgfqpoint{2.246778in}{1.453408in}}%
\pgfpathlineto{\pgfqpoint{2.258163in}{1.454659in}}%
\pgfpathlineto{\pgfqpoint{2.261163in}{1.431916in}}%
\pgfpathlineto{\pgfqpoint{2.263712in}{1.409069in}}%
\pgfpathlineto{\pgfqpoint{2.252299in}{1.407809in}}%
\pgfpathlineto{\pgfqpoint{2.229609in}{1.405210in}}%
\pgfpathlineto{\pgfqpoint{2.232094in}{1.382445in}}%
\pgfpathlineto{\pgfqpoint{2.196663in}{1.378800in}}%
\pgfpathclose%
\pgfusepath{fill}%
\end{pgfscope}%
\begin{pgfscope}%
\pgfpathrectangle{\pgfqpoint{0.100000in}{0.100000in}}{\pgfqpoint{3.007045in}{1.925000in}}%
\pgfusepath{clip}%
\pgfsetbuttcap%
\pgfsetmiterjoin%
\definecolor{currentfill}{rgb}{0.890196,0.957647,0.870065}%
\pgfsetfillcolor{currentfill}%
\pgfsetlinewidth{0.000000pt}%
\definecolor{currentstroke}{rgb}{0.000000,0.000000,0.000000}%
\pgfsetstrokecolor{currentstroke}%
\pgfsetstrokeopacity{0.000000}%
\pgfsetdash{}{0pt}%
\pgfpathmoveto{\pgfqpoint{0.973337in}{1.262333in}}%
\pgfpathlineto{\pgfqpoint{0.963825in}{1.264056in}}%
\pgfpathlineto{\pgfqpoint{0.964301in}{1.266613in}}%
\pgfpathlineto{\pgfqpoint{0.957377in}{1.275606in}}%
\pgfpathlineto{\pgfqpoint{0.959849in}{1.282582in}}%
\pgfpathlineto{\pgfqpoint{0.957665in}{1.289984in}}%
\pgfpathlineto{\pgfqpoint{0.950156in}{1.292847in}}%
\pgfpathlineto{\pgfqpoint{0.943668in}{1.304520in}}%
\pgfpathlineto{\pgfqpoint{0.947451in}{1.310370in}}%
\pgfpathlineto{\pgfqpoint{0.937610in}{1.310762in}}%
\pgfpathlineto{\pgfqpoint{0.925436in}{1.305547in}}%
\pgfpathlineto{\pgfqpoint{0.923286in}{1.309346in}}%
\pgfpathlineto{\pgfqpoint{0.914647in}{1.310630in}}%
\pgfpathlineto{\pgfqpoint{0.913219in}{1.297816in}}%
\pgfpathlineto{\pgfqpoint{0.914434in}{1.291420in}}%
\pgfpathlineto{\pgfqpoint{0.910706in}{1.281361in}}%
\pgfpathlineto{\pgfqpoint{0.901338in}{1.275958in}}%
\pgfpathlineto{\pgfqpoint{0.878712in}{1.280320in}}%
\pgfpathlineto{\pgfqpoint{0.816917in}{1.293210in}}%
\pgfpathlineto{\pgfqpoint{0.832155in}{1.363654in}}%
\pgfpathlineto{\pgfqpoint{0.892707in}{1.351230in}}%
\pgfpathlineto{\pgfqpoint{0.908601in}{1.353242in}}%
\pgfpathlineto{\pgfqpoint{0.924294in}{1.367661in}}%
\pgfpathlineto{\pgfqpoint{0.934476in}{1.365698in}}%
\pgfpathlineto{\pgfqpoint{0.942058in}{1.369750in}}%
\pgfpathlineto{\pgfqpoint{0.947822in}{1.365378in}}%
\pgfpathlineto{\pgfqpoint{0.951603in}{1.367889in}}%
\pgfpathlineto{\pgfqpoint{0.960688in}{1.366618in}}%
\pgfpathlineto{\pgfqpoint{0.963794in}{1.361956in}}%
\pgfpathlineto{\pgfqpoint{0.970085in}{1.358692in}}%
\pgfpathlineto{\pgfqpoint{0.969963in}{1.352958in}}%
\pgfpathlineto{\pgfqpoint{0.972493in}{1.347031in}}%
\pgfpathlineto{\pgfqpoint{0.981285in}{1.351407in}}%
\pgfpathlineto{\pgfqpoint{0.978339in}{1.334944in}}%
\pgfpathlineto{\pgfqpoint{1.005052in}{1.330045in}}%
\pgfpathlineto{\pgfqpoint{1.029601in}{1.326060in}}%
\pgfpathlineto{\pgfqpoint{1.027613in}{1.314125in}}%
\pgfpathlineto{\pgfqpoint{1.019386in}{1.315423in}}%
\pgfpathlineto{\pgfqpoint{1.015074in}{1.317685in}}%
\pgfpathlineto{\pgfqpoint{1.008923in}{1.315572in}}%
\pgfpathlineto{\pgfqpoint{0.994498in}{1.314874in}}%
\pgfpathlineto{\pgfqpoint{0.989953in}{1.316140in}}%
\pgfpathlineto{\pgfqpoint{0.981761in}{1.312933in}}%
\pgfpathlineto{\pgfqpoint{0.973337in}{1.262333in}}%
\pgfpathclose%
\pgfusepath{fill}%
\end{pgfscope}%
\begin{pgfscope}%
\pgfpathrectangle{\pgfqpoint{0.100000in}{0.100000in}}{\pgfqpoint{3.007045in}{1.925000in}}%
\pgfusepath{clip}%
\pgfsetbuttcap%
\pgfsetmiterjoin%
\definecolor{currentfill}{rgb}{0.760523,0.905359,0.733595}%
\pgfsetfillcolor{currentfill}%
\pgfsetlinewidth{0.000000pt}%
\definecolor{currentstroke}{rgb}{0.000000,0.000000,0.000000}%
\pgfsetstrokecolor{currentstroke}%
\pgfsetstrokeopacity{0.000000}%
\pgfsetdash{}{0pt}%
\pgfpathmoveto{\pgfqpoint{1.730711in}{1.250739in}}%
\pgfpathlineto{\pgfqpoint{1.699296in}{1.251149in}}%
\pgfpathlineto{\pgfqpoint{1.699583in}{1.268350in}}%
\pgfpathlineto{\pgfqpoint{1.677475in}{1.268823in}}%
\pgfpathlineto{\pgfqpoint{1.678558in}{1.314735in}}%
\pgfpathlineto{\pgfqpoint{1.701276in}{1.314251in}}%
\pgfpathlineto{\pgfqpoint{1.701191in}{1.310424in}}%
\pgfpathlineto{\pgfqpoint{1.717094in}{1.310143in}}%
\pgfpathlineto{\pgfqpoint{1.719990in}{1.317507in}}%
\pgfpathlineto{\pgfqpoint{1.716390in}{1.322247in}}%
\pgfpathlineto{\pgfqpoint{1.739290in}{1.321824in}}%
\pgfpathlineto{\pgfqpoint{1.744960in}{1.321802in}}%
\pgfpathlineto{\pgfqpoint{1.744815in}{1.304579in}}%
\pgfpathlineto{\pgfqpoint{1.747728in}{1.298216in}}%
\pgfpathlineto{\pgfqpoint{1.764683in}{1.298101in}}%
\pgfpathlineto{\pgfqpoint{1.764595in}{1.275306in}}%
\pgfpathlineto{\pgfqpoint{1.753309in}{1.275315in}}%
\pgfpathlineto{\pgfqpoint{1.753153in}{1.258349in}}%
\pgfpathlineto{\pgfqpoint{1.731811in}{1.258499in}}%
\pgfpathlineto{\pgfqpoint{1.730711in}{1.250739in}}%
\pgfpathclose%
\pgfusepath{fill}%
\end{pgfscope}%
\begin{pgfscope}%
\pgfpathrectangle{\pgfqpoint{0.100000in}{0.100000in}}{\pgfqpoint{3.007045in}{1.925000in}}%
\pgfusepath{clip}%
\pgfsetbuttcap%
\pgfsetmiterjoin%
\definecolor{currentfill}{rgb}{0.634295,0.852211,0.610688}%
\pgfsetfillcolor{currentfill}%
\pgfsetlinewidth{0.000000pt}%
\definecolor{currentstroke}{rgb}{0.000000,0.000000,0.000000}%
\pgfsetstrokecolor{currentstroke}%
\pgfsetstrokeopacity{0.000000}%
\pgfsetdash{}{0pt}%
\pgfpathmoveto{\pgfqpoint{1.928151in}{1.455464in}}%
\pgfpathlineto{\pgfqpoint{1.928977in}{1.432597in}}%
\pgfpathlineto{\pgfqpoint{1.894560in}{1.431480in}}%
\pgfpathlineto{\pgfqpoint{1.893937in}{1.453516in}}%
\pgfpathlineto{\pgfqpoint{1.882427in}{1.454160in}}%
\pgfpathlineto{\pgfqpoint{1.865473in}{1.453759in}}%
\pgfpathlineto{\pgfqpoint{1.864898in}{1.494975in}}%
\pgfpathlineto{\pgfqpoint{1.864831in}{1.497891in}}%
\pgfpathlineto{\pgfqpoint{1.876437in}{1.499997in}}%
\pgfpathlineto{\pgfqpoint{1.876275in}{1.505700in}}%
\pgfpathlineto{\pgfqpoint{1.884271in}{1.505208in}}%
\pgfpathlineto{\pgfqpoint{1.889296in}{1.502422in}}%
\pgfpathlineto{\pgfqpoint{1.898889in}{1.500432in}}%
\pgfpathlineto{\pgfqpoint{1.903859in}{1.493796in}}%
\pgfpathlineto{\pgfqpoint{1.911427in}{1.491667in}}%
\pgfpathlineto{\pgfqpoint{1.917765in}{1.487298in}}%
\pgfpathlineto{\pgfqpoint{1.921271in}{1.478116in}}%
\pgfpathlineto{\pgfqpoint{1.910791in}{1.477769in}}%
\pgfpathlineto{\pgfqpoint{1.911530in}{1.454907in}}%
\pgfpathlineto{\pgfqpoint{1.928151in}{1.455464in}}%
\pgfpathclose%
\pgfusepath{fill}%
\end{pgfscope}%
\begin{pgfscope}%
\pgfpathrectangle{\pgfqpoint{0.100000in}{0.100000in}}{\pgfqpoint{3.007045in}{1.925000in}}%
\pgfusepath{clip}%
\pgfsetbuttcap%
\pgfsetmiterjoin%
\definecolor{currentfill}{rgb}{0.535179,0.806090,0.528750}%
\pgfsetfillcolor{currentfill}%
\pgfsetlinewidth{0.000000pt}%
\definecolor{currentstroke}{rgb}{0.000000,0.000000,0.000000}%
\pgfsetstrokecolor{currentstroke}%
\pgfsetstrokeopacity{0.000000}%
\pgfsetdash{}{0pt}%
\pgfpathmoveto{\pgfqpoint{1.637239in}{1.477803in}}%
\pgfpathlineto{\pgfqpoint{1.682752in}{1.476661in}}%
\pgfpathlineto{\pgfqpoint{1.682306in}{1.453680in}}%
\pgfpathlineto{\pgfqpoint{1.648008in}{1.454576in}}%
\pgfpathlineto{\pgfqpoint{1.636440in}{1.454988in}}%
\pgfpathlineto{\pgfqpoint{1.637239in}{1.477803in}}%
\pgfpathclose%
\pgfusepath{fill}%
\end{pgfscope}%
\begin{pgfscope}%
\pgfpathrectangle{\pgfqpoint{0.100000in}{0.100000in}}{\pgfqpoint{3.007045in}{1.925000in}}%
\pgfusepath{clip}%
\pgfsetbuttcap%
\pgfsetmiterjoin%
\definecolor{currentfill}{rgb}{0.474279,0.777670,0.478677}%
\pgfsetfillcolor{currentfill}%
\pgfsetlinewidth{0.000000pt}%
\definecolor{currentstroke}{rgb}{0.000000,0.000000,0.000000}%
\pgfsetstrokecolor{currentstroke}%
\pgfsetstrokeopacity{0.000000}%
\pgfsetdash{}{0pt}%
\pgfpathmoveto{\pgfqpoint{1.755176in}{1.198941in}}%
\pgfpathlineto{\pgfqpoint{1.755095in}{1.176082in}}%
\pgfpathlineto{\pgfqpoint{1.743716in}{1.176139in}}%
\pgfpathlineto{\pgfqpoint{1.732315in}{1.176225in}}%
\pgfpathlineto{\pgfqpoint{1.732247in}{1.170499in}}%
\pgfpathlineto{\pgfqpoint{1.719727in}{1.170628in}}%
\pgfpathlineto{\pgfqpoint{1.709420in}{1.170767in}}%
\pgfpathlineto{\pgfqpoint{1.709799in}{1.199381in}}%
\pgfpathlineto{\pgfqpoint{1.755176in}{1.198941in}}%
\pgfpathclose%
\pgfusepath{fill}%
\end{pgfscope}%
\begin{pgfscope}%
\pgfpathrectangle{\pgfqpoint{0.100000in}{0.100000in}}{\pgfqpoint{3.007045in}{1.925000in}}%
\pgfusepath{clip}%
\pgfsetbuttcap%
\pgfsetmiterjoin%
\definecolor{currentfill}{rgb}{0.445490,0.764014,0.458131}%
\pgfsetfillcolor{currentfill}%
\pgfsetlinewidth{0.000000pt}%
\definecolor{currentstroke}{rgb}{0.000000,0.000000,0.000000}%
\pgfsetstrokecolor{currentstroke}%
\pgfsetstrokeopacity{0.000000}%
\pgfsetdash{}{0pt}%
\pgfpathmoveto{\pgfqpoint{2.029155in}{1.469240in}}%
\pgfpathlineto{\pgfqpoint{2.030394in}{1.446824in}}%
\pgfpathlineto{\pgfqpoint{2.021528in}{1.446204in}}%
\pgfpathlineto{\pgfqpoint{2.017195in}{1.455365in}}%
\pgfpathlineto{\pgfqpoint{2.012501in}{1.460181in}}%
\pgfpathlineto{\pgfqpoint{2.012462in}{1.466036in}}%
\pgfpathlineto{\pgfqpoint{2.008307in}{1.475250in}}%
\pgfpathlineto{\pgfqpoint{2.013699in}{1.486047in}}%
\pgfpathlineto{\pgfqpoint{1.994330in}{1.484956in}}%
\pgfpathlineto{\pgfqpoint{1.993541in}{1.496568in}}%
\pgfpathlineto{\pgfqpoint{1.970609in}{1.495307in}}%
\pgfpathlineto{\pgfqpoint{1.970355in}{1.501053in}}%
\pgfpathlineto{\pgfqpoint{1.964639in}{1.500815in}}%
\pgfpathlineto{\pgfqpoint{1.962214in}{1.552559in}}%
\pgfpathlineto{\pgfqpoint{1.961960in}{1.558341in}}%
\pgfpathlineto{\pgfqpoint{1.973425in}{1.558760in}}%
\pgfpathlineto{\pgfqpoint{2.003003in}{1.560543in}}%
\pgfpathlineto{\pgfqpoint{2.002328in}{1.571903in}}%
\pgfpathlineto{\pgfqpoint{2.030834in}{1.573678in}}%
\pgfpathlineto{\pgfqpoint{2.031351in}{1.567858in}}%
\pgfpathlineto{\pgfqpoint{2.032705in}{1.544832in}}%
\pgfpathlineto{\pgfqpoint{2.042126in}{1.545379in}}%
\pgfpathlineto{\pgfqpoint{2.043986in}{1.516491in}}%
\pgfpathlineto{\pgfqpoint{2.045773in}{1.487575in}}%
\pgfpathlineto{\pgfqpoint{2.028120in}{1.486633in}}%
\pgfpathlineto{\pgfqpoint{2.029155in}{1.469240in}}%
\pgfpathclose%
\pgfusepath{fill}%
\end{pgfscope}%
\begin{pgfscope}%
\pgfpathrectangle{\pgfqpoint{0.100000in}{0.100000in}}{\pgfqpoint{3.007045in}{1.925000in}}%
\pgfusepath{clip}%
\pgfsetbuttcap%
\pgfsetmiterjoin%
\definecolor{currentfill}{rgb}{0.474279,0.777670,0.478677}%
\pgfsetfillcolor{currentfill}%
\pgfsetlinewidth{0.000000pt}%
\definecolor{currentstroke}{rgb}{0.000000,0.000000,0.000000}%
\pgfsetstrokecolor{currentstroke}%
\pgfsetstrokeopacity{0.000000}%
\pgfsetdash{}{0pt}%
\pgfpathmoveto{\pgfqpoint{0.459905in}{1.401168in}}%
\pgfpathlineto{\pgfqpoint{0.451708in}{1.400095in}}%
\pgfpathlineto{\pgfqpoint{0.446584in}{1.393519in}}%
\pgfpathlineto{\pgfqpoint{0.443854in}{1.394476in}}%
\pgfpathlineto{\pgfqpoint{0.437071in}{1.389645in}}%
\pgfpathlineto{\pgfqpoint{0.425415in}{1.393166in}}%
\pgfpathlineto{\pgfqpoint{0.423677in}{1.387762in}}%
\pgfpathlineto{\pgfqpoint{0.380635in}{1.401231in}}%
\pgfpathlineto{\pgfqpoint{0.384397in}{1.412475in}}%
\pgfpathlineto{\pgfqpoint{0.385058in}{1.423808in}}%
\pgfpathlineto{\pgfqpoint{0.387316in}{1.430255in}}%
\pgfpathlineto{\pgfqpoint{0.384276in}{1.433228in}}%
\pgfpathlineto{\pgfqpoint{0.385422in}{1.437519in}}%
\pgfpathlineto{\pgfqpoint{0.390077in}{1.441086in}}%
\pgfpathlineto{\pgfqpoint{0.395246in}{1.441320in}}%
\pgfpathlineto{\pgfqpoint{0.399108in}{1.444237in}}%
\pgfpathlineto{\pgfqpoint{0.407911in}{1.446662in}}%
\pgfpathlineto{\pgfqpoint{0.407285in}{1.452726in}}%
\pgfpathlineto{\pgfqpoint{0.412778in}{1.458416in}}%
\pgfpathlineto{\pgfqpoint{0.420836in}{1.470570in}}%
\pgfpathlineto{\pgfqpoint{0.424324in}{1.471482in}}%
\pgfpathlineto{\pgfqpoint{0.429023in}{1.481634in}}%
\pgfpathlineto{\pgfqpoint{0.484393in}{1.464821in}}%
\pgfpathlineto{\pgfqpoint{0.479723in}{1.447096in}}%
\pgfpathlineto{\pgfqpoint{0.470680in}{1.418127in}}%
\pgfpathlineto{\pgfqpoint{0.462527in}{1.420581in}}%
\pgfpathlineto{\pgfqpoint{0.462117in}{1.414113in}}%
\pgfpathlineto{\pgfqpoint{0.466986in}{1.410310in}}%
\pgfpathlineto{\pgfqpoint{0.464449in}{1.404081in}}%
\pgfpathlineto{\pgfqpoint{0.459905in}{1.401168in}}%
\pgfpathclose%
\pgfusepath{fill}%
\end{pgfscope}%
\begin{pgfscope}%
\pgfpathrectangle{\pgfqpoint{0.100000in}{0.100000in}}{\pgfqpoint{3.007045in}{1.925000in}}%
\pgfusepath{clip}%
\pgfsetbuttcap%
\pgfsetmiterjoin%
\definecolor{currentfill}{rgb}{0.657670,0.862053,0.633449}%
\pgfsetfillcolor{currentfill}%
\pgfsetlinewidth{0.000000pt}%
\definecolor{currentstroke}{rgb}{0.000000,0.000000,0.000000}%
\pgfsetstrokecolor{currentstroke}%
\pgfsetstrokeopacity{0.000000}%
\pgfsetdash{}{0pt}%
\pgfpathmoveto{\pgfqpoint{2.256834in}{1.132262in}}%
\pgfpathlineto{\pgfqpoint{2.256292in}{1.137004in}}%
\pgfpathlineto{\pgfqpoint{2.250151in}{1.140216in}}%
\pgfpathlineto{\pgfqpoint{2.249568in}{1.145924in}}%
\pgfpathlineto{\pgfqpoint{2.253275in}{1.146329in}}%
\pgfpathlineto{\pgfqpoint{2.261118in}{1.153025in}}%
\pgfpathlineto{\pgfqpoint{2.273431in}{1.154396in}}%
\pgfpathlineto{\pgfqpoint{2.275093in}{1.139834in}}%
\pgfpathlineto{\pgfqpoint{2.283647in}{1.145438in}}%
\pgfpathlineto{\pgfqpoint{2.288955in}{1.139422in}}%
\pgfpathlineto{\pgfqpoint{2.282277in}{1.134362in}}%
\pgfpathlineto{\pgfqpoint{2.277544in}{1.133106in}}%
\pgfpathlineto{\pgfqpoint{2.272501in}{1.126660in}}%
\pgfpathlineto{\pgfqpoint{2.264410in}{1.127604in}}%
\pgfpathlineto{\pgfqpoint{2.264234in}{1.131721in}}%
\pgfpathlineto{\pgfqpoint{2.256834in}{1.132262in}}%
\pgfpathclose%
\pgfusepath{fill}%
\end{pgfscope}%
\begin{pgfscope}%
\pgfpathrectangle{\pgfqpoint{0.100000in}{0.100000in}}{\pgfqpoint{3.007045in}{1.925000in}}%
\pgfusepath{clip}%
\pgfsetbuttcap%
\pgfsetmiterjoin%
\definecolor{currentfill}{rgb}{0.607151,0.839677,0.587928}%
\pgfsetfillcolor{currentfill}%
\pgfsetlinewidth{0.000000pt}%
\definecolor{currentstroke}{rgb}{0.000000,0.000000,0.000000}%
\pgfsetstrokecolor{currentstroke}%
\pgfsetstrokeopacity{0.000000}%
\pgfsetdash{}{0pt}%
\pgfpathmoveto{\pgfqpoint{1.610215in}{1.316599in}}%
\pgfpathlineto{\pgfqpoint{1.609594in}{1.293697in}}%
\pgfpathlineto{\pgfqpoint{1.587220in}{1.294590in}}%
\pgfpathlineto{\pgfqpoint{1.564548in}{1.295579in}}%
\pgfpathlineto{\pgfqpoint{1.565174in}{1.318418in}}%
\pgfpathlineto{\pgfqpoint{1.541840in}{1.319542in}}%
\pgfpathlineto{\pgfqpoint{1.543010in}{1.342313in}}%
\pgfpathlineto{\pgfqpoint{1.565129in}{1.341345in}}%
\pgfpathlineto{\pgfqpoint{1.588274in}{1.340363in}}%
\pgfpathlineto{\pgfqpoint{1.587433in}{1.317426in}}%
\pgfpathlineto{\pgfqpoint{1.610215in}{1.316599in}}%
\pgfpathclose%
\pgfusepath{fill}%
\end{pgfscope}%
\begin{pgfscope}%
\pgfpathrectangle{\pgfqpoint{0.100000in}{0.100000in}}{\pgfqpoint{3.007045in}{1.925000in}}%
\pgfusepath{clip}%
\pgfsetbuttcap%
\pgfsetmiterjoin%
\definecolor{currentfill}{rgb}{0.652995,0.860085,0.628897}%
\pgfsetfillcolor{currentfill}%
\pgfsetlinewidth{0.000000pt}%
\definecolor{currentstroke}{rgb}{0.000000,0.000000,0.000000}%
\pgfsetstrokecolor{currentstroke}%
\pgfsetstrokeopacity{0.000000}%
\pgfsetdash{}{0pt}%
\pgfpathmoveto{\pgfqpoint{1.656900in}{1.355436in}}%
\pgfpathlineto{\pgfqpoint{1.673964in}{1.354981in}}%
\pgfpathlineto{\pgfqpoint{1.674656in}{1.382036in}}%
\pgfpathlineto{\pgfqpoint{1.684732in}{1.378050in}}%
\pgfpathlineto{\pgfqpoint{1.685257in}{1.403080in}}%
\pgfpathlineto{\pgfqpoint{1.702210in}{1.402774in}}%
\pgfpathlineto{\pgfqpoint{1.699433in}{1.400596in}}%
\pgfpathlineto{\pgfqpoint{1.699908in}{1.394517in}}%
\pgfpathlineto{\pgfqpoint{1.697951in}{1.391327in}}%
\pgfpathlineto{\pgfqpoint{1.730762in}{1.390911in}}%
\pgfpathlineto{\pgfqpoint{1.730520in}{1.367928in}}%
\pgfpathlineto{\pgfqpoint{1.736283in}{1.367899in}}%
\pgfpathlineto{\pgfqpoint{1.739732in}{1.362155in}}%
\pgfpathlineto{\pgfqpoint{1.739541in}{1.344779in}}%
\pgfpathlineto{\pgfqpoint{1.739290in}{1.321824in}}%
\pgfpathlineto{\pgfqpoint{1.716390in}{1.322247in}}%
\pgfpathlineto{\pgfqpoint{1.716907in}{1.329278in}}%
\pgfpathlineto{\pgfqpoint{1.713888in}{1.330082in}}%
\pgfpathlineto{\pgfqpoint{1.706612in}{1.341900in}}%
\pgfpathlineto{\pgfqpoint{1.706125in}{1.349454in}}%
\pgfpathlineto{\pgfqpoint{1.688055in}{1.349861in}}%
\pgfpathlineto{\pgfqpoint{1.683358in}{1.349049in}}%
\pgfpathlineto{\pgfqpoint{1.683057in}{1.337584in}}%
\pgfpathlineto{\pgfqpoint{1.673497in}{1.337786in}}%
\pgfpathlineto{\pgfqpoint{1.656422in}{1.338221in}}%
\pgfpathlineto{\pgfqpoint{1.656900in}{1.355436in}}%
\pgfpathclose%
\pgfusepath{fill}%
\end{pgfscope}%
\begin{pgfscope}%
\pgfpathrectangle{\pgfqpoint{0.100000in}{0.100000in}}{\pgfqpoint{3.007045in}{1.925000in}}%
\pgfusepath{clip}%
\pgfsetbuttcap%
\pgfsetmiterjoin%
\definecolor{currentfill}{rgb}{0.518570,0.798339,0.515094}%
\pgfsetfillcolor{currentfill}%
\pgfsetlinewidth{0.000000pt}%
\definecolor{currentstroke}{rgb}{0.000000,0.000000,0.000000}%
\pgfsetstrokecolor{currentstroke}%
\pgfsetstrokeopacity{0.000000}%
\pgfsetdash{}{0pt}%
\pgfpathmoveto{\pgfqpoint{2.559434in}{0.779009in}}%
\pgfpathlineto{\pgfqpoint{2.556087in}{0.772933in}}%
\pgfpathlineto{\pgfqpoint{2.548564in}{0.780580in}}%
\pgfpathlineto{\pgfqpoint{2.535487in}{0.765030in}}%
\pgfpathlineto{\pgfqpoint{2.529482in}{0.767093in}}%
\pgfpathlineto{\pgfqpoint{2.526324in}{0.773753in}}%
\pgfpathlineto{\pgfqpoint{2.525012in}{0.782861in}}%
\pgfpathlineto{\pgfqpoint{2.519904in}{0.788430in}}%
\pgfpathlineto{\pgfqpoint{2.519777in}{0.793317in}}%
\pgfpathlineto{\pgfqpoint{2.516660in}{0.795332in}}%
\pgfpathlineto{\pgfqpoint{2.507037in}{0.798645in}}%
\pgfpathlineto{\pgfqpoint{2.503376in}{0.803575in}}%
\pgfpathlineto{\pgfqpoint{2.521533in}{0.825837in}}%
\pgfpathlineto{\pgfqpoint{2.525782in}{0.825781in}}%
\pgfpathlineto{\pgfqpoint{2.532363in}{0.822168in}}%
\pgfpathlineto{\pgfqpoint{2.545370in}{0.817990in}}%
\pgfpathlineto{\pgfqpoint{2.556125in}{0.810180in}}%
\pgfpathlineto{\pgfqpoint{2.542055in}{0.797900in}}%
\pgfpathlineto{\pgfqpoint{2.548566in}{0.790531in}}%
\pgfpathlineto{\pgfqpoint{2.554010in}{0.786819in}}%
\pgfpathlineto{\pgfqpoint{2.559434in}{0.779009in}}%
\pgfpathclose%
\pgfusepath{fill}%
\end{pgfscope}%
\begin{pgfscope}%
\pgfpathrectangle{\pgfqpoint{0.100000in}{0.100000in}}{\pgfqpoint{3.007045in}{1.925000in}}%
\pgfusepath{clip}%
\pgfsetbuttcap%
\pgfsetmiterjoin%
\definecolor{currentfill}{rgb}{0.856978,0.944360,0.834633}%
\pgfsetfillcolor{currentfill}%
\pgfsetlinewidth{0.000000pt}%
\definecolor{currentstroke}{rgb}{0.000000,0.000000,0.000000}%
\pgfsetstrokecolor{currentstroke}%
\pgfsetstrokeopacity{0.000000}%
\pgfsetdash{}{0pt}%
\pgfpathmoveto{\pgfqpoint{0.936059in}{0.163351in}}%
\pgfpathlineto{\pgfqpoint{0.936182in}{0.166254in}}%
\pgfpathlineto{\pgfqpoint{0.938191in}{0.164676in}}%
\pgfpathlineto{\pgfqpoint{0.936059in}{0.163351in}}%
\pgfpathclose%
\pgfusepath{fill}%
\end{pgfscope}%
\begin{pgfscope}%
\pgfpathrectangle{\pgfqpoint{0.100000in}{0.100000in}}{\pgfqpoint{3.007045in}{1.925000in}}%
\pgfusepath{clip}%
\pgfsetbuttcap%
\pgfsetmiterjoin%
\definecolor{currentfill}{rgb}{0.856978,0.944360,0.834633}%
\pgfsetfillcolor{currentfill}%
\pgfsetlinewidth{0.000000pt}%
\definecolor{currentstroke}{rgb}{0.000000,0.000000,0.000000}%
\pgfsetstrokecolor{currentstroke}%
\pgfsetstrokeopacity{0.000000}%
\pgfsetdash{}{0pt}%
\pgfpathmoveto{\pgfqpoint{0.933022in}{0.152128in}}%
\pgfpathlineto{\pgfqpoint{0.933743in}{0.154327in}}%
\pgfpathlineto{\pgfqpoint{0.935121in}{0.154547in}}%
\pgfpathlineto{\pgfqpoint{0.934597in}{0.151513in}}%
\pgfpathlineto{\pgfqpoint{0.933022in}{0.152128in}}%
\pgfpathclose%
\pgfusepath{fill}%
\end{pgfscope}%
\begin{pgfscope}%
\pgfpathrectangle{\pgfqpoint{0.100000in}{0.100000in}}{\pgfqpoint{3.007045in}{1.925000in}}%
\pgfusepath{clip}%
\pgfsetbuttcap%
\pgfsetmiterjoin%
\definecolor{currentfill}{rgb}{0.856978,0.944360,0.834633}%
\pgfsetfillcolor{currentfill}%
\pgfsetlinewidth{0.000000pt}%
\definecolor{currentstroke}{rgb}{0.000000,0.000000,0.000000}%
\pgfsetstrokecolor{currentstroke}%
\pgfsetstrokeopacity{0.000000}%
\pgfsetdash{}{0pt}%
\pgfpathmoveto{\pgfqpoint{0.935935in}{0.137518in}}%
\pgfpathlineto{\pgfqpoint{0.934705in}{0.142381in}}%
\pgfpathlineto{\pgfqpoint{0.935452in}{0.143773in}}%
\pgfpathlineto{\pgfqpoint{0.933837in}{0.145038in}}%
\pgfpathlineto{\pgfqpoint{0.935285in}{0.148282in}}%
\pgfpathlineto{\pgfqpoint{0.935191in}{0.150312in}}%
\pgfpathlineto{\pgfqpoint{0.936976in}{0.151065in}}%
\pgfpathlineto{\pgfqpoint{0.937372in}{0.148133in}}%
\pgfpathlineto{\pgfqpoint{0.935770in}{0.146761in}}%
\pgfpathlineto{\pgfqpoint{0.936802in}{0.145789in}}%
\pgfpathlineto{\pgfqpoint{0.936825in}{0.141016in}}%
\pgfpathlineto{\pgfqpoint{0.939296in}{0.140402in}}%
\pgfpathlineto{\pgfqpoint{0.935935in}{0.137518in}}%
\pgfpathclose%
\pgfusepath{fill}%
\end{pgfscope}%
\begin{pgfscope}%
\pgfpathrectangle{\pgfqpoint{0.100000in}{0.100000in}}{\pgfqpoint{3.007045in}{1.925000in}}%
\pgfusepath{clip}%
\pgfsetbuttcap%
\pgfsetmiterjoin%
\definecolor{currentfill}{rgb}{0.856978,0.944360,0.834633}%
\pgfsetfillcolor{currentfill}%
\pgfsetlinewidth{0.000000pt}%
\definecolor{currentstroke}{rgb}{0.000000,0.000000,0.000000}%
\pgfsetstrokecolor{currentstroke}%
\pgfsetstrokeopacity{0.000000}%
\pgfsetdash{}{0pt}%
\pgfpathmoveto{\pgfqpoint{0.957491in}{0.142127in}}%
\pgfpathlineto{\pgfqpoint{0.957670in}{0.139655in}}%
\pgfpathlineto{\pgfqpoint{0.956604in}{0.137353in}}%
\pgfpathlineto{\pgfqpoint{0.955422in}{0.137331in}}%
\pgfpathlineto{\pgfqpoint{0.954352in}{0.139936in}}%
\pgfpathlineto{\pgfqpoint{0.955529in}{0.140546in}}%
\pgfpathlineto{\pgfqpoint{0.955668in}{0.143108in}}%
\pgfpathlineto{\pgfqpoint{0.957491in}{0.142127in}}%
\pgfpathclose%
\pgfusepath{fill}%
\end{pgfscope}%
\begin{pgfscope}%
\pgfpathrectangle{\pgfqpoint{0.100000in}{0.100000in}}{\pgfqpoint{3.007045in}{1.925000in}}%
\pgfusepath{clip}%
\pgfsetbuttcap%
\pgfsetmiterjoin%
\definecolor{currentfill}{rgb}{0.856978,0.944360,0.834633}%
\pgfsetfillcolor{currentfill}%
\pgfsetlinewidth{0.000000pt}%
\definecolor{currentstroke}{rgb}{0.000000,0.000000,0.000000}%
\pgfsetstrokecolor{currentstroke}%
\pgfsetstrokeopacity{0.000000}%
\pgfsetdash{}{0pt}%
\pgfpathmoveto{\pgfqpoint{0.952217in}{0.157810in}}%
\pgfpathlineto{\pgfqpoint{0.953820in}{0.154455in}}%
\pgfpathlineto{\pgfqpoint{0.955416in}{0.154694in}}%
\pgfpathlineto{\pgfqpoint{0.956678in}{0.156700in}}%
\pgfpathlineto{\pgfqpoint{0.956199in}{0.158645in}}%
\pgfpathlineto{\pgfqpoint{0.958436in}{0.159148in}}%
\pgfpathlineto{\pgfqpoint{0.958681in}{0.160262in}}%
\pgfpathlineto{\pgfqpoint{0.960550in}{0.161270in}}%
\pgfpathlineto{\pgfqpoint{0.962912in}{0.161570in}}%
\pgfpathlineto{\pgfqpoint{0.964181in}{0.158832in}}%
\pgfpathlineto{\pgfqpoint{0.966014in}{0.159668in}}%
\pgfpathlineto{\pgfqpoint{0.965820in}{0.161252in}}%
\pgfpathlineto{\pgfqpoint{0.968598in}{0.161733in}}%
\pgfpathlineto{\pgfqpoint{0.970864in}{0.160909in}}%
\pgfpathlineto{\pgfqpoint{0.972584in}{0.162951in}}%
\pgfpathlineto{\pgfqpoint{0.972690in}{0.165137in}}%
\pgfpathlineto{\pgfqpoint{0.975165in}{0.163335in}}%
\pgfpathlineto{\pgfqpoint{0.977660in}{0.159042in}}%
\pgfpathlineto{\pgfqpoint{0.978047in}{0.156183in}}%
\pgfpathlineto{\pgfqpoint{0.979828in}{0.154295in}}%
\pgfpathlineto{\pgfqpoint{0.981654in}{0.154152in}}%
\pgfpathlineto{\pgfqpoint{0.981934in}{0.151716in}}%
\pgfpathlineto{\pgfqpoint{0.981007in}{0.149518in}}%
\pgfpathlineto{\pgfqpoint{0.977742in}{0.146985in}}%
\pgfpathlineto{\pgfqpoint{0.977518in}{0.144790in}}%
\pgfpathlineto{\pgfqpoint{0.976466in}{0.142679in}}%
\pgfpathlineto{\pgfqpoint{0.976184in}{0.138931in}}%
\pgfpathlineto{\pgfqpoint{0.975258in}{0.135551in}}%
\pgfpathlineto{\pgfqpoint{0.973176in}{0.134289in}}%
\pgfpathlineto{\pgfqpoint{0.967906in}{0.129628in}}%
\pgfpathlineto{\pgfqpoint{0.965959in}{0.130065in}}%
\pgfpathlineto{\pgfqpoint{0.963807in}{0.128960in}}%
\pgfpathlineto{\pgfqpoint{0.961890in}{0.130440in}}%
\pgfpathlineto{\pgfqpoint{0.959954in}{0.130233in}}%
\pgfpathlineto{\pgfqpoint{0.961401in}{0.136798in}}%
\pgfpathlineto{\pgfqpoint{0.961289in}{0.138714in}}%
\pgfpathlineto{\pgfqpoint{0.964237in}{0.140766in}}%
\pgfpathlineto{\pgfqpoint{0.965836in}{0.140351in}}%
\pgfpathlineto{\pgfqpoint{0.965862in}{0.142940in}}%
\pgfpathlineto{\pgfqpoint{0.967505in}{0.145942in}}%
\pgfpathlineto{\pgfqpoint{0.969065in}{0.151851in}}%
\pgfpathlineto{\pgfqpoint{0.966959in}{0.151516in}}%
\pgfpathlineto{\pgfqpoint{0.967397in}{0.149343in}}%
\pgfpathlineto{\pgfqpoint{0.966518in}{0.146372in}}%
\pgfpathlineto{\pgfqpoint{0.962740in}{0.141294in}}%
\pgfpathlineto{\pgfqpoint{0.959748in}{0.140140in}}%
\pgfpathlineto{\pgfqpoint{0.959257in}{0.141993in}}%
\pgfpathlineto{\pgfqpoint{0.956453in}{0.144467in}}%
\pgfpathlineto{\pgfqpoint{0.953985in}{0.143893in}}%
\pgfpathlineto{\pgfqpoint{0.952414in}{0.142851in}}%
\pgfpathlineto{\pgfqpoint{0.953506in}{0.147090in}}%
\pgfpathlineto{\pgfqpoint{0.955124in}{0.149168in}}%
\pgfpathlineto{\pgfqpoint{0.956748in}{0.149574in}}%
\pgfpathlineto{\pgfqpoint{0.957816in}{0.151764in}}%
\pgfpathlineto{\pgfqpoint{0.959671in}{0.153160in}}%
\pgfpathlineto{\pgfqpoint{0.959717in}{0.154964in}}%
\pgfpathlineto{\pgfqpoint{0.962294in}{0.155551in}}%
\pgfpathlineto{\pgfqpoint{0.962264in}{0.157330in}}%
\pgfpathlineto{\pgfqpoint{0.959964in}{0.158575in}}%
\pgfpathlineto{\pgfqpoint{0.957588in}{0.157646in}}%
\pgfpathlineto{\pgfqpoint{0.958906in}{0.156128in}}%
\pgfpathlineto{\pgfqpoint{0.956807in}{0.153413in}}%
\pgfpathlineto{\pgfqpoint{0.953556in}{0.150672in}}%
\pgfpathlineto{\pgfqpoint{0.951866in}{0.154211in}}%
\pgfpathlineto{\pgfqpoint{0.952217in}{0.157810in}}%
\pgfpathclose%
\pgfusepath{fill}%
\end{pgfscope}%
\begin{pgfscope}%
\pgfpathrectangle{\pgfqpoint{0.100000in}{0.100000in}}{\pgfqpoint{3.007045in}{1.925000in}}%
\pgfusepath{clip}%
\pgfsetbuttcap%
\pgfsetmiterjoin%
\definecolor{currentfill}{rgb}{0.856978,0.944360,0.834633}%
\pgfsetfillcolor{currentfill}%
\pgfsetlinewidth{0.000000pt}%
\definecolor{currentstroke}{rgb}{0.000000,0.000000,0.000000}%
\pgfsetstrokecolor{currentstroke}%
\pgfsetstrokeopacity{0.000000}%
\pgfsetdash{}{0pt}%
\pgfpathmoveto{\pgfqpoint{0.932602in}{0.158192in}}%
\pgfpathlineto{\pgfqpoint{0.930583in}{0.159046in}}%
\pgfpathlineto{\pgfqpoint{0.933442in}{0.160711in}}%
\pgfpathlineto{\pgfqpoint{0.934281in}{0.158960in}}%
\pgfpathlineto{\pgfqpoint{0.937131in}{0.159694in}}%
\pgfpathlineto{\pgfqpoint{0.937907in}{0.158017in}}%
\pgfpathlineto{\pgfqpoint{0.935667in}{0.157983in}}%
\pgfpathlineto{\pgfqpoint{0.934472in}{0.156058in}}%
\pgfpathlineto{\pgfqpoint{0.932593in}{0.155358in}}%
\pgfpathlineto{\pgfqpoint{0.931238in}{0.156086in}}%
\pgfpathlineto{\pgfqpoint{0.932602in}{0.158192in}}%
\pgfpathclose%
\pgfusepath{fill}%
\end{pgfscope}%
\begin{pgfscope}%
\pgfpathrectangle{\pgfqpoint{0.100000in}{0.100000in}}{\pgfqpoint{3.007045in}{1.925000in}}%
\pgfusepath{clip}%
\pgfsetbuttcap%
\pgfsetmiterjoin%
\definecolor{currentfill}{rgb}{0.856978,0.944360,0.834633}%
\pgfsetfillcolor{currentfill}%
\pgfsetlinewidth{0.000000pt}%
\definecolor{currentstroke}{rgb}{0.000000,0.000000,0.000000}%
\pgfsetstrokecolor{currentstroke}%
\pgfsetstrokeopacity{0.000000}%
\pgfsetdash{}{0pt}%
\pgfpathmoveto{\pgfqpoint{0.940503in}{0.148050in}}%
\pgfpathlineto{\pgfqpoint{0.940580in}{0.150009in}}%
\pgfpathlineto{\pgfqpoint{0.936401in}{0.153144in}}%
\pgfpathlineto{\pgfqpoint{0.939865in}{0.154873in}}%
\pgfpathlineto{\pgfqpoint{0.939188in}{0.156547in}}%
\pgfpathlineto{\pgfqpoint{0.940021in}{0.158363in}}%
\pgfpathlineto{\pgfqpoint{0.938477in}{0.159030in}}%
\pgfpathlineto{\pgfqpoint{0.937509in}{0.161120in}}%
\pgfpathlineto{\pgfqpoint{0.940507in}{0.163189in}}%
\pgfpathlineto{\pgfqpoint{0.940853in}{0.165668in}}%
\pgfpathlineto{\pgfqpoint{0.942540in}{0.166461in}}%
\pgfpathlineto{\pgfqpoint{0.942490in}{0.168936in}}%
\pgfpathlineto{\pgfqpoint{0.940371in}{0.169760in}}%
\pgfpathlineto{\pgfqpoint{0.935534in}{0.169087in}}%
\pgfpathlineto{\pgfqpoint{0.937048in}{0.170932in}}%
\pgfpathlineto{\pgfqpoint{0.938908in}{0.171837in}}%
\pgfpathlineto{\pgfqpoint{0.941019in}{0.171940in}}%
\pgfpathlineto{\pgfqpoint{0.940657in}{0.175170in}}%
\pgfpathlineto{\pgfqpoint{0.941701in}{0.177230in}}%
\pgfpathlineto{\pgfqpoint{0.943306in}{0.176542in}}%
\pgfpathlineto{\pgfqpoint{0.946910in}{0.174004in}}%
\pgfpathlineto{\pgfqpoint{0.946946in}{0.170521in}}%
\pgfpathlineto{\pgfqpoint{0.945241in}{0.169179in}}%
\pgfpathlineto{\pgfqpoint{0.945972in}{0.166130in}}%
\pgfpathlineto{\pgfqpoint{0.947228in}{0.166360in}}%
\pgfpathlineto{\pgfqpoint{0.948782in}{0.164350in}}%
\pgfpathlineto{\pgfqpoint{0.949856in}{0.160592in}}%
\pgfpathlineto{\pgfqpoint{0.949616in}{0.158564in}}%
\pgfpathlineto{\pgfqpoint{0.948423in}{0.157492in}}%
\pgfpathlineto{\pgfqpoint{0.950053in}{0.156174in}}%
\pgfpathlineto{\pgfqpoint{0.949902in}{0.154095in}}%
\pgfpathlineto{\pgfqpoint{0.948375in}{0.153768in}}%
\pgfpathlineto{\pgfqpoint{0.947443in}{0.155763in}}%
\pgfpathlineto{\pgfqpoint{0.946239in}{0.153864in}}%
\pgfpathlineto{\pgfqpoint{0.948621in}{0.152037in}}%
\pgfpathlineto{\pgfqpoint{0.948626in}{0.150372in}}%
\pgfpathlineto{\pgfqpoint{0.949742in}{0.147800in}}%
\pgfpathlineto{\pgfqpoint{0.948428in}{0.144747in}}%
\pgfpathlineto{\pgfqpoint{0.950997in}{0.145602in}}%
\pgfpathlineto{\pgfqpoint{0.950534in}{0.143772in}}%
\pgfpathlineto{\pgfqpoint{0.948944in}{0.142658in}}%
\pgfpathlineto{\pgfqpoint{0.947671in}{0.143289in}}%
\pgfpathlineto{\pgfqpoint{0.946770in}{0.140379in}}%
\pgfpathlineto{\pgfqpoint{0.949016in}{0.140644in}}%
\pgfpathlineto{\pgfqpoint{0.947609in}{0.137730in}}%
\pgfpathlineto{\pgfqpoint{0.947150in}{0.135497in}}%
\pgfpathlineto{\pgfqpoint{0.945969in}{0.134095in}}%
\pgfpathlineto{\pgfqpoint{0.943530in}{0.134941in}}%
\pgfpathlineto{\pgfqpoint{0.943106in}{0.139361in}}%
\pgfpathlineto{\pgfqpoint{0.941552in}{0.142463in}}%
\pgfpathlineto{\pgfqpoint{0.941399in}{0.144418in}}%
\pgfpathlineto{\pgfqpoint{0.942947in}{0.147821in}}%
\pgfpathlineto{\pgfqpoint{0.940503in}{0.148050in}}%
\pgfpathclose%
\pgfusepath{fill}%
\end{pgfscope}%
\begin{pgfscope}%
\pgfpathrectangle{\pgfqpoint{0.100000in}{0.100000in}}{\pgfqpoint{3.007045in}{1.925000in}}%
\pgfusepath{clip}%
\pgfsetbuttcap%
\pgfsetmiterjoin%
\definecolor{currentfill}{rgb}{0.751173,0.901423,0.724491}%
\pgfsetfillcolor{currentfill}%
\pgfsetlinewidth{0.000000pt}%
\definecolor{currentstroke}{rgb}{0.000000,0.000000,0.000000}%
\pgfsetstrokecolor{currentstroke}%
\pgfsetstrokeopacity{0.000000}%
\pgfsetdash{}{0pt}%
\pgfpathmoveto{\pgfqpoint{0.746029in}{1.809361in}}%
\pgfpathlineto{\pgfqpoint{0.740338in}{1.787082in}}%
\pgfpathlineto{\pgfqpoint{0.735917in}{1.779747in}}%
\pgfpathlineto{\pgfqpoint{0.726344in}{1.778628in}}%
\pgfpathlineto{\pgfqpoint{0.702590in}{1.784664in}}%
\pgfpathlineto{\pgfqpoint{0.675279in}{1.792196in}}%
\pgfpathlineto{\pgfqpoint{0.670581in}{1.789487in}}%
\pgfpathlineto{\pgfqpoint{0.668689in}{1.793331in}}%
\pgfpathlineto{\pgfqpoint{0.662371in}{1.789397in}}%
\pgfpathlineto{\pgfqpoint{0.651182in}{1.791391in}}%
\pgfpathlineto{\pgfqpoint{0.648904in}{1.795115in}}%
\pgfpathlineto{\pgfqpoint{0.651860in}{1.802864in}}%
\pgfpathlineto{\pgfqpoint{0.652827in}{1.812664in}}%
\pgfpathlineto{\pgfqpoint{0.652061in}{1.822418in}}%
\pgfpathlineto{\pgfqpoint{0.656028in}{1.830896in}}%
\pgfpathlineto{\pgfqpoint{0.657936in}{1.837275in}}%
\pgfpathlineto{\pgfqpoint{0.662469in}{1.839071in}}%
\pgfpathlineto{\pgfqpoint{0.665649in}{1.843218in}}%
\pgfpathlineto{\pgfqpoint{0.680569in}{1.839333in}}%
\pgfpathlineto{\pgfqpoint{0.682133in}{1.844866in}}%
\pgfpathlineto{\pgfqpoint{0.689171in}{1.848788in}}%
\pgfpathlineto{\pgfqpoint{0.691919in}{1.848024in}}%
\pgfpathlineto{\pgfqpoint{0.695822in}{1.854965in}}%
\pgfpathlineto{\pgfqpoint{0.700822in}{1.858600in}}%
\pgfpathlineto{\pgfqpoint{0.702364in}{1.864190in}}%
\pgfpathlineto{\pgfqpoint{0.708324in}{1.867692in}}%
\pgfpathlineto{\pgfqpoint{0.713573in}{1.866288in}}%
\pgfpathlineto{\pgfqpoint{0.713636in}{1.864996in}}%
\pgfpathlineto{\pgfqpoint{0.701521in}{1.821189in}}%
\pgfpathlineto{\pgfqpoint{0.746029in}{1.809361in}}%
\pgfpathclose%
\pgfusepath{fill}%
\end{pgfscope}%
\begin{pgfscope}%
\pgfpathrectangle{\pgfqpoint{0.100000in}{0.100000in}}{\pgfqpoint{3.007045in}{1.925000in}}%
\pgfusepath{clip}%
\pgfsetbuttcap%
\pgfsetmiterjoin%
\definecolor{currentfill}{rgb}{0.376471,0.730181,0.424298}%
\pgfsetfillcolor{currentfill}%
\pgfsetlinewidth{0.000000pt}%
\definecolor{currentstroke}{rgb}{0.000000,0.000000,0.000000}%
\pgfsetstrokecolor{currentstroke}%
\pgfsetstrokeopacity{0.000000}%
\pgfsetdash{}{0pt}%
\pgfpathmoveto{\pgfqpoint{2.488733in}{1.208153in}}%
\pgfpathlineto{\pgfqpoint{2.485346in}{1.213775in}}%
\pgfpathlineto{\pgfqpoint{2.476467in}{1.219730in}}%
\pgfpathlineto{\pgfqpoint{2.478476in}{1.225076in}}%
\pgfpathlineto{\pgfqpoint{2.467867in}{1.223907in}}%
\pgfpathlineto{\pgfqpoint{2.465782in}{1.225372in}}%
\pgfpathlineto{\pgfqpoint{2.462970in}{1.231712in}}%
\pgfpathlineto{\pgfqpoint{2.461646in}{1.242312in}}%
\pgfpathlineto{\pgfqpoint{2.465474in}{1.242785in}}%
\pgfpathlineto{\pgfqpoint{2.463037in}{1.262554in}}%
\pgfpathlineto{\pgfqpoint{2.480209in}{1.264505in}}%
\pgfpathlineto{\pgfqpoint{2.498296in}{1.267386in}}%
\pgfpathlineto{\pgfqpoint{2.502865in}{1.238879in}}%
\pgfpathlineto{\pgfqpoint{2.507767in}{1.239655in}}%
\pgfpathlineto{\pgfqpoint{2.510163in}{1.234931in}}%
\pgfpathlineto{\pgfqpoint{2.505584in}{1.228741in}}%
\pgfpathlineto{\pgfqpoint{2.507751in}{1.225110in}}%
\pgfpathlineto{\pgfqpoint{2.505896in}{1.220258in}}%
\pgfpathlineto{\pgfqpoint{2.500713in}{1.220454in}}%
\pgfpathlineto{\pgfqpoint{2.488733in}{1.208153in}}%
\pgfpathclose%
\pgfusepath{fill}%
\end{pgfscope}%
\begin{pgfscope}%
\pgfpathrectangle{\pgfqpoint{0.100000in}{0.100000in}}{\pgfqpoint{3.007045in}{1.925000in}}%
\pgfusepath{clip}%
\pgfsetbuttcap%
\pgfsetmiterjoin%
\definecolor{currentfill}{rgb}{0.623760,0.847428,0.601584}%
\pgfsetfillcolor{currentfill}%
\pgfsetlinewidth{0.000000pt}%
\definecolor{currentstroke}{rgb}{0.000000,0.000000,0.000000}%
\pgfsetstrokecolor{currentstroke}%
\pgfsetstrokeopacity{0.000000}%
\pgfsetdash{}{0pt}%
\pgfpathmoveto{\pgfqpoint{2.905539in}{1.431014in}}%
\pgfpathlineto{\pgfqpoint{2.906901in}{1.437874in}}%
\pgfpathlineto{\pgfqpoint{2.903419in}{1.453280in}}%
\pgfpathlineto{\pgfqpoint{2.896973in}{1.475592in}}%
\pgfpathlineto{\pgfqpoint{2.916591in}{1.481492in}}%
\pgfpathlineto{\pgfqpoint{2.917977in}{1.479575in}}%
\pgfpathlineto{\pgfqpoint{2.929528in}{1.490206in}}%
\pgfpathlineto{\pgfqpoint{2.933228in}{1.482168in}}%
\pgfpathlineto{\pgfqpoint{2.944479in}{1.471321in}}%
\pgfpathlineto{\pgfqpoint{2.948844in}{1.463490in}}%
\pgfpathlineto{\pgfqpoint{2.944989in}{1.461097in}}%
\pgfpathlineto{\pgfqpoint{2.945676in}{1.456613in}}%
\pgfpathlineto{\pgfqpoint{2.941792in}{1.452515in}}%
\pgfpathlineto{\pgfqpoint{2.928813in}{1.447884in}}%
\pgfpathlineto{\pgfqpoint{2.928414in}{1.457369in}}%
\pgfpathlineto{\pgfqpoint{2.925276in}{1.463337in}}%
\pgfpathlineto{\pgfqpoint{2.921321in}{1.458455in}}%
\pgfpathlineto{\pgfqpoint{2.921311in}{1.451942in}}%
\pgfpathlineto{\pgfqpoint{2.924101in}{1.446233in}}%
\pgfpathlineto{\pgfqpoint{2.922682in}{1.438922in}}%
\pgfpathlineto{\pgfqpoint{2.919165in}{1.438139in}}%
\pgfpathlineto{\pgfqpoint{2.905539in}{1.431014in}}%
\pgfpathclose%
\pgfusepath{fill}%
\end{pgfscope}%
\begin{pgfscope}%
\pgfpathrectangle{\pgfqpoint{0.100000in}{0.100000in}}{\pgfqpoint{3.007045in}{1.925000in}}%
\pgfusepath{clip}%
\pgfsetbuttcap%
\pgfsetmiterjoin%
\definecolor{currentfill}{rgb}{0.786851,0.916309,0.759831}%
\pgfsetfillcolor{currentfill}%
\pgfsetlinewidth{0.000000pt}%
\definecolor{currentstroke}{rgb}{0.000000,0.000000,0.000000}%
\pgfsetstrokecolor{currentstroke}%
\pgfsetstrokeopacity{0.000000}%
\pgfsetdash{}{0pt}%
\pgfpathmoveto{\pgfqpoint{1.806790in}{1.547902in}}%
\pgfpathlineto{\pgfqpoint{1.801048in}{1.550695in}}%
\pgfpathlineto{\pgfqpoint{1.783351in}{1.550568in}}%
\pgfpathlineto{\pgfqpoint{1.783347in}{1.556314in}}%
\pgfpathlineto{\pgfqpoint{1.766199in}{1.556278in}}%
\pgfpathlineto{\pgfqpoint{1.765864in}{1.580217in}}%
\pgfpathlineto{\pgfqpoint{1.806019in}{1.580465in}}%
\pgfpathlineto{\pgfqpoint{1.802241in}{1.583541in}}%
\pgfpathlineto{\pgfqpoint{1.829469in}{1.583832in}}%
\pgfpathlineto{\pgfqpoint{1.829097in}{1.594472in}}%
\pgfpathlineto{\pgfqpoint{1.827186in}{1.594478in}}%
\pgfpathlineto{\pgfqpoint{1.827033in}{1.605997in}}%
\pgfpathlineto{\pgfqpoint{1.826933in}{1.611742in}}%
\pgfpathlineto{\pgfqpoint{1.844395in}{1.612207in}}%
\pgfpathlineto{\pgfqpoint{1.844441in}{1.606089in}}%
\pgfpathlineto{\pgfqpoint{1.844608in}{1.594634in}}%
\pgfpathlineto{\pgfqpoint{1.840662in}{1.594594in}}%
\pgfpathlineto{\pgfqpoint{1.841127in}{1.578186in}}%
\pgfpathlineto{\pgfqpoint{1.841735in}{1.545791in}}%
\pgfpathlineto{\pgfqpoint{1.838457in}{1.548803in}}%
\pgfpathlineto{\pgfqpoint{1.830410in}{1.548866in}}%
\pgfpathlineto{\pgfqpoint{1.816671in}{1.557157in}}%
\pgfpathlineto{\pgfqpoint{1.813546in}{1.550942in}}%
\pgfpathlineto{\pgfqpoint{1.806790in}{1.547902in}}%
\pgfpathclose%
\pgfusepath{fill}%
\end{pgfscope}%
\begin{pgfscope}%
\pgfpathrectangle{\pgfqpoint{0.100000in}{0.100000in}}{\pgfqpoint{3.007045in}{1.925000in}}%
\pgfusepath{clip}%
\pgfsetbuttcap%
\pgfsetmiterjoin%
\definecolor{currentfill}{rgb}{0.713772,0.885675,0.688074}%
\pgfsetfillcolor{currentfill}%
\pgfsetlinewidth{0.000000pt}%
\definecolor{currentstroke}{rgb}{0.000000,0.000000,0.000000}%
\pgfsetstrokecolor{currentstroke}%
\pgfsetstrokeopacity{0.000000}%
\pgfsetdash{}{0pt}%
\pgfpathmoveto{\pgfqpoint{0.967728in}{1.403968in}}%
\pgfpathlineto{\pgfqpoint{0.967940in}{1.408953in}}%
\pgfpathlineto{\pgfqpoint{0.964652in}{1.410841in}}%
\pgfpathlineto{\pgfqpoint{0.964506in}{1.418586in}}%
\pgfpathlineto{\pgfqpoint{0.968279in}{1.422237in}}%
\pgfpathlineto{\pgfqpoint{0.969777in}{1.427234in}}%
\pgfpathlineto{\pgfqpoint{0.968314in}{1.431784in}}%
\pgfpathlineto{\pgfqpoint{0.955057in}{1.434288in}}%
\pgfpathlineto{\pgfqpoint{0.955530in}{1.447686in}}%
\pgfpathlineto{\pgfqpoint{0.954259in}{1.454699in}}%
\pgfpathlineto{\pgfqpoint{0.947546in}{1.454706in}}%
\pgfpathlineto{\pgfqpoint{0.946788in}{1.460996in}}%
\pgfpathlineto{\pgfqpoint{0.947898in}{1.467983in}}%
\pgfpathlineto{\pgfqpoint{0.953549in}{1.475280in}}%
\pgfpathlineto{\pgfqpoint{1.001890in}{1.466220in}}%
\pgfpathlineto{\pgfqpoint{0.989963in}{1.400079in}}%
\pgfpathlineto{\pgfqpoint{0.967728in}{1.403968in}}%
\pgfpathclose%
\pgfusepath{fill}%
\end{pgfscope}%
\begin{pgfscope}%
\pgfpathrectangle{\pgfqpoint{0.100000in}{0.100000in}}{\pgfqpoint{3.007045in}{1.925000in}}%
\pgfusepath{clip}%
\pgfsetbuttcap%
\pgfsetmiterjoin%
\definecolor{currentfill}{rgb}{0.468743,0.775087,0.474125}%
\pgfsetfillcolor{currentfill}%
\pgfsetlinewidth{0.000000pt}%
\definecolor{currentstroke}{rgb}{0.000000,0.000000,0.000000}%
\pgfsetstrokecolor{currentstroke}%
\pgfsetstrokeopacity{0.000000}%
\pgfsetdash{}{0pt}%
\pgfpathmoveto{\pgfqpoint{1.844374in}{1.020395in}}%
\pgfpathlineto{\pgfqpoint{1.822007in}{1.020648in}}%
\pgfpathlineto{\pgfqpoint{1.820409in}{1.024532in}}%
\pgfpathlineto{\pgfqpoint{1.792358in}{1.025328in}}%
\pgfpathlineto{\pgfqpoint{1.792356in}{1.069644in}}%
\pgfpathlineto{\pgfqpoint{1.792345in}{1.071169in}}%
\pgfpathlineto{\pgfqpoint{1.811752in}{1.070677in}}%
\pgfpathlineto{\pgfqpoint{1.821187in}{1.069825in}}%
\pgfpathlineto{\pgfqpoint{1.820896in}{1.060894in}}%
\pgfpathlineto{\pgfqpoint{1.834274in}{1.060508in}}%
\pgfpathlineto{\pgfqpoint{1.834159in}{1.056680in}}%
\pgfpathlineto{\pgfqpoint{1.843696in}{1.056469in}}%
\pgfpathlineto{\pgfqpoint{1.844766in}{1.050707in}}%
\pgfpathlineto{\pgfqpoint{1.844374in}{1.020395in}}%
\pgfpathclose%
\pgfusepath{fill}%
\end{pgfscope}%
\begin{pgfscope}%
\pgfpathrectangle{\pgfqpoint{0.100000in}{0.100000in}}{\pgfqpoint{3.007045in}{1.925000in}}%
\pgfusepath{clip}%
\pgfsetbuttcap%
\pgfsetmiterjoin%
\definecolor{currentfill}{rgb}{0.634295,0.852211,0.610688}%
\pgfsetfillcolor{currentfill}%
\pgfsetlinewidth{0.000000pt}%
\definecolor{currentstroke}{rgb}{0.000000,0.000000,0.000000}%
\pgfsetstrokecolor{currentstroke}%
\pgfsetstrokeopacity{0.000000}%
\pgfsetdash{}{0pt}%
\pgfpathmoveto{\pgfqpoint{2.597898in}{0.973667in}}%
\pgfpathlineto{\pgfqpoint{2.564813in}{0.967383in}}%
\pgfpathlineto{\pgfqpoint{2.562648in}{0.972115in}}%
\pgfpathlineto{\pgfqpoint{2.555149in}{0.977464in}}%
\pgfpathlineto{\pgfqpoint{2.547659in}{0.980532in}}%
\pgfpathlineto{\pgfqpoint{2.545915in}{0.985956in}}%
\pgfpathlineto{\pgfqpoint{2.549802in}{0.987470in}}%
\pgfpathlineto{\pgfqpoint{2.549576in}{0.995557in}}%
\pgfpathlineto{\pgfqpoint{2.557574in}{1.001121in}}%
\pgfpathlineto{\pgfqpoint{2.566736in}{1.001545in}}%
\pgfpathlineto{\pgfqpoint{2.561819in}{1.036295in}}%
\pgfpathlineto{\pgfqpoint{2.578144in}{1.039001in}}%
\pgfpathlineto{\pgfqpoint{2.578534in}{1.060078in}}%
\pgfpathlineto{\pgfqpoint{2.578887in}{1.072582in}}%
\pgfpathlineto{\pgfqpoint{2.584592in}{1.075050in}}%
\pgfpathlineto{\pgfqpoint{2.591032in}{1.080451in}}%
\pgfpathlineto{\pgfqpoint{2.596091in}{1.081050in}}%
\pgfpathlineto{\pgfqpoint{2.599528in}{1.076992in}}%
\pgfpathlineto{\pgfqpoint{2.604673in}{1.078358in}}%
\pgfpathlineto{\pgfqpoint{2.604048in}{1.043452in}}%
\pgfpathlineto{\pgfqpoint{2.608257in}{1.044199in}}%
\pgfpathlineto{\pgfqpoint{2.610860in}{1.024597in}}%
\pgfpathlineto{\pgfqpoint{2.605433in}{1.023762in}}%
\pgfpathlineto{\pgfqpoint{2.608855in}{1.001837in}}%
\pgfpathlineto{\pgfqpoint{2.611008in}{0.997970in}}%
\pgfpathlineto{\pgfqpoint{2.594940in}{0.995133in}}%
\pgfpathlineto{\pgfqpoint{2.597898in}{0.973667in}}%
\pgfpathclose%
\pgfusepath{fill}%
\end{pgfscope}%
\begin{pgfscope}%
\pgfpathrectangle{\pgfqpoint{0.100000in}{0.100000in}}{\pgfqpoint{3.007045in}{1.925000in}}%
\pgfusepath{clip}%
\pgfsetbuttcap%
\pgfsetmiterjoin%
\definecolor{currentfill}{rgb}{0.468743,0.775087,0.474125}%
\pgfsetfillcolor{currentfill}%
\pgfsetlinewidth{0.000000pt}%
\definecolor{currentstroke}{rgb}{0.000000,0.000000,0.000000}%
\pgfsetstrokecolor{currentstroke}%
\pgfsetstrokeopacity{0.000000}%
\pgfsetdash{}{0pt}%
\pgfpathmoveto{\pgfqpoint{1.166919in}{0.389823in}}%
\pgfpathlineto{\pgfqpoint{1.169914in}{0.398030in}}%
\pgfpathlineto{\pgfqpoint{1.170066in}{0.403152in}}%
\pgfpathlineto{\pgfqpoint{1.177565in}{0.397908in}}%
\pgfpathlineto{\pgfqpoint{1.179080in}{0.390976in}}%
\pgfpathlineto{\pgfqpoint{1.174659in}{0.386184in}}%
\pgfpathlineto{\pgfqpoint{1.166919in}{0.389823in}}%
\pgfpathclose%
\pgfusepath{fill}%
\end{pgfscope}%
\begin{pgfscope}%
\pgfpathrectangle{\pgfqpoint{0.100000in}{0.100000in}}{\pgfqpoint{3.007045in}{1.925000in}}%
\pgfusepath{clip}%
\pgfsetbuttcap%
\pgfsetmiterjoin%
\definecolor{currentfill}{rgb}{0.468743,0.775087,0.474125}%
\pgfsetfillcolor{currentfill}%
\pgfsetlinewidth{0.000000pt}%
\definecolor{currentstroke}{rgb}{0.000000,0.000000,0.000000}%
\pgfsetstrokecolor{currentstroke}%
\pgfsetstrokeopacity{0.000000}%
\pgfsetdash{}{0pt}%
\pgfpathmoveto{\pgfqpoint{1.171506in}{0.365544in}}%
\pgfpathlineto{\pgfqpoint{1.172099in}{0.368432in}}%
\pgfpathlineto{\pgfqpoint{1.181292in}{0.366754in}}%
\pgfpathlineto{\pgfqpoint{1.179030in}{0.362563in}}%
\pgfpathlineto{\pgfqpoint{1.171506in}{0.365544in}}%
\pgfpathclose%
\pgfusepath{fill}%
\end{pgfscope}%
\begin{pgfscope}%
\pgfpathrectangle{\pgfqpoint{0.100000in}{0.100000in}}{\pgfqpoint{3.007045in}{1.925000in}}%
\pgfusepath{clip}%
\pgfsetbuttcap%
\pgfsetmiterjoin%
\definecolor{currentfill}{rgb}{0.468743,0.775087,0.474125}%
\pgfsetfillcolor{currentfill}%
\pgfsetlinewidth{0.000000pt}%
\definecolor{currentstroke}{rgb}{0.000000,0.000000,0.000000}%
\pgfsetstrokecolor{currentstroke}%
\pgfsetstrokeopacity{0.000000}%
\pgfsetdash{}{0pt}%
\pgfpathmoveto{\pgfqpoint{1.189967in}{0.357288in}}%
\pgfpathlineto{\pgfqpoint{1.187711in}{0.361758in}}%
\pgfpathlineto{\pgfqpoint{1.194261in}{0.372053in}}%
\pgfpathlineto{\pgfqpoint{1.187347in}{0.380076in}}%
\pgfpathlineto{\pgfqpoint{1.187339in}{0.386651in}}%
\pgfpathlineto{\pgfqpoint{1.193857in}{0.392205in}}%
\pgfpathlineto{\pgfqpoint{1.198739in}{0.390671in}}%
\pgfpathlineto{\pgfqpoint{1.199911in}{0.385707in}}%
\pgfpathlineto{\pgfqpoint{1.198482in}{0.379133in}}%
\pgfpathlineto{\pgfqpoint{1.210453in}{0.373714in}}%
\pgfpathlineto{\pgfqpoint{1.213763in}{0.360359in}}%
\pgfpathlineto{\pgfqpoint{1.218074in}{0.354132in}}%
\pgfpathlineto{\pgfqpoint{1.215157in}{0.349080in}}%
\pgfpathlineto{\pgfqpoint{1.209803in}{0.348364in}}%
\pgfpathlineto{\pgfqpoint{1.200882in}{0.353223in}}%
\pgfpathlineto{\pgfqpoint{1.189967in}{0.357288in}}%
\pgfpathclose%
\pgfusepath{fill}%
\end{pgfscope}%
\begin{pgfscope}%
\pgfpathrectangle{\pgfqpoint{0.100000in}{0.100000in}}{\pgfqpoint{3.007045in}{1.925000in}}%
\pgfusepath{clip}%
\pgfsetbuttcap%
\pgfsetmiterjoin%
\definecolor{currentfill}{rgb}{0.468743,0.775087,0.474125}%
\pgfsetfillcolor{currentfill}%
\pgfsetlinewidth{0.000000pt}%
\definecolor{currentstroke}{rgb}{0.000000,0.000000,0.000000}%
\pgfsetstrokecolor{currentstroke}%
\pgfsetstrokeopacity{0.000000}%
\pgfsetdash{}{0pt}%
\pgfpathmoveto{\pgfqpoint{1.188373in}{0.411114in}}%
\pgfpathlineto{\pgfqpoint{1.197345in}{0.404388in}}%
\pgfpathlineto{\pgfqpoint{1.194069in}{0.401026in}}%
\pgfpathlineto{\pgfqpoint{1.187346in}{0.401919in}}%
\pgfpathlineto{\pgfqpoint{1.182053in}{0.407040in}}%
\pgfpathlineto{\pgfqpoint{1.178085in}{0.413644in}}%
\pgfpathlineto{\pgfqpoint{1.165729in}{0.423161in}}%
\pgfpathlineto{\pgfqpoint{1.174372in}{0.427023in}}%
\pgfpathlineto{\pgfqpoint{1.188373in}{0.411114in}}%
\pgfpathclose%
\pgfusepath{fill}%
\end{pgfscope}%
\begin{pgfscope}%
\pgfpathrectangle{\pgfqpoint{0.100000in}{0.100000in}}{\pgfqpoint{3.007045in}{1.925000in}}%
\pgfusepath{clip}%
\pgfsetbuttcap%
\pgfsetmiterjoin%
\definecolor{currentfill}{rgb}{0.634295,0.852211,0.610688}%
\pgfsetfillcolor{currentfill}%
\pgfsetlinewidth{0.000000pt}%
\definecolor{currentstroke}{rgb}{0.000000,0.000000,0.000000}%
\pgfsetstrokecolor{currentstroke}%
\pgfsetstrokeopacity{0.000000}%
\pgfsetdash{}{0pt}%
\pgfpathmoveto{\pgfqpoint{1.330767in}{0.811238in}}%
\pgfpathlineto{\pgfqpoint{1.328991in}{0.794521in}}%
\pgfpathlineto{\pgfqpoint{1.303863in}{0.796794in}}%
\pgfpathlineto{\pgfqpoint{1.304431in}{0.802517in}}%
\pgfpathlineto{\pgfqpoint{1.292993in}{0.803535in}}%
\pgfpathlineto{\pgfqpoint{1.294107in}{0.814272in}}%
\pgfpathlineto{\pgfqpoint{1.287264in}{0.814931in}}%
\pgfpathlineto{\pgfqpoint{1.288966in}{0.832202in}}%
\pgfpathlineto{\pgfqpoint{1.283287in}{0.832832in}}%
\pgfpathlineto{\pgfqpoint{1.286524in}{0.867195in}}%
\pgfpathlineto{\pgfqpoint{1.299641in}{0.865888in}}%
\pgfpathlineto{\pgfqpoint{1.300241in}{0.871662in}}%
\pgfpathlineto{\pgfqpoint{1.311575in}{0.870525in}}%
\pgfpathlineto{\pgfqpoint{1.312137in}{0.876255in}}%
\pgfpathlineto{\pgfqpoint{1.317814in}{0.875699in}}%
\pgfpathlineto{\pgfqpoint{1.318390in}{0.881414in}}%
\pgfpathlineto{\pgfqpoint{1.324074in}{0.880866in}}%
\pgfpathlineto{\pgfqpoint{1.324670in}{0.886744in}}%
\pgfpathlineto{\pgfqpoint{1.337618in}{0.885562in}}%
\pgfpathlineto{\pgfqpoint{1.336407in}{0.871962in}}%
\pgfpathlineto{\pgfqpoint{1.364391in}{0.869531in}}%
\pgfpathlineto{\pgfqpoint{1.361996in}{0.840962in}}%
\pgfpathlineto{\pgfqpoint{1.333803in}{0.843365in}}%
\pgfpathlineto{\pgfqpoint{1.330767in}{0.811238in}}%
\pgfpathclose%
\pgfusepath{fill}%
\end{pgfscope}%
\begin{pgfscope}%
\pgfpathrectangle{\pgfqpoint{0.100000in}{0.100000in}}{\pgfqpoint{3.007045in}{1.925000in}}%
\pgfusepath{clip}%
\pgfsetbuttcap%
\pgfsetmiterjoin%
\definecolor{currentfill}{rgb}{0.618224,0.844844,0.597032}%
\pgfsetfillcolor{currentfill}%
\pgfsetlinewidth{0.000000pt}%
\definecolor{currentstroke}{rgb}{0.000000,0.000000,0.000000}%
\pgfsetstrokecolor{currentstroke}%
\pgfsetstrokeopacity{0.000000}%
\pgfsetdash{}{0pt}%
\pgfpathmoveto{\pgfqpoint{2.590755in}{0.837224in}}%
\pgfpathlineto{\pgfqpoint{2.584761in}{0.832670in}}%
\pgfpathlineto{\pgfqpoint{2.577826in}{0.830633in}}%
\pgfpathlineto{\pgfqpoint{2.571176in}{0.834427in}}%
\pgfpathlineto{\pgfqpoint{2.566054in}{0.842571in}}%
\pgfpathlineto{\pgfqpoint{2.559661in}{0.847879in}}%
\pgfpathlineto{\pgfqpoint{2.560497in}{0.852920in}}%
\pgfpathlineto{\pgfqpoint{2.558099in}{0.855702in}}%
\pgfpathlineto{\pgfqpoint{2.557919in}{0.868208in}}%
\pgfpathlineto{\pgfqpoint{2.555420in}{0.871742in}}%
\pgfpathlineto{\pgfqpoint{2.551190in}{0.869149in}}%
\pgfpathlineto{\pgfqpoint{2.541892in}{0.875274in}}%
\pgfpathlineto{\pgfqpoint{2.545276in}{0.888835in}}%
\pgfpathlineto{\pgfqpoint{2.538562in}{0.893175in}}%
\pgfpathlineto{\pgfqpoint{2.545027in}{0.899687in}}%
\pgfpathlineto{\pgfqpoint{2.550816in}{0.897393in}}%
\pgfpathlineto{\pgfqpoint{2.555983in}{0.898307in}}%
\pgfpathlineto{\pgfqpoint{2.555178in}{0.902655in}}%
\pgfpathlineto{\pgfqpoint{2.562238in}{0.907450in}}%
\pgfpathlineto{\pgfqpoint{2.567423in}{0.900852in}}%
\pgfpathlineto{\pgfqpoint{2.571314in}{0.892287in}}%
\pgfpathlineto{\pgfqpoint{2.577854in}{0.890686in}}%
\pgfpathlineto{\pgfqpoint{2.582316in}{0.884878in}}%
\pgfpathlineto{\pgfqpoint{2.580510in}{0.880092in}}%
\pgfpathlineto{\pgfqpoint{2.596705in}{0.871224in}}%
\pgfpathlineto{\pgfqpoint{2.592778in}{0.867711in}}%
\pgfpathlineto{\pgfqpoint{2.597384in}{0.859907in}}%
\pgfpathlineto{\pgfqpoint{2.593582in}{0.855868in}}%
\pgfpathlineto{\pgfqpoint{2.595239in}{0.853039in}}%
\pgfpathlineto{\pgfqpoint{2.590755in}{0.837224in}}%
\pgfpathclose%
\pgfusepath{fill}%
\end{pgfscope}%
\begin{pgfscope}%
\pgfpathrectangle{\pgfqpoint{0.100000in}{0.100000in}}{\pgfqpoint{3.007045in}{1.925000in}}%
\pgfusepath{clip}%
\pgfsetbuttcap%
\pgfsetmiterjoin%
\definecolor{currentfill}{rgb}{0.585006,0.829343,0.569719}%
\pgfsetfillcolor{currentfill}%
\pgfsetlinewidth{0.000000pt}%
\definecolor{currentstroke}{rgb}{0.000000,0.000000,0.000000}%
\pgfsetstrokecolor{currentstroke}%
\pgfsetstrokeopacity{0.000000}%
\pgfsetdash{}{0pt}%
\pgfpathmoveto{\pgfqpoint{1.318468in}{1.629024in}}%
\pgfpathlineto{\pgfqpoint{1.316018in}{1.629285in}}%
\pgfpathlineto{\pgfqpoint{1.313493in}{1.606222in}}%
\pgfpathlineto{\pgfqpoint{1.311526in}{1.606441in}}%
\pgfpathlineto{\pgfqpoint{1.255957in}{1.613050in}}%
\pgfpathlineto{\pgfqpoint{1.258817in}{1.635994in}}%
\pgfpathlineto{\pgfqpoint{1.260561in}{1.635790in}}%
\pgfpathlineto{\pgfqpoint{1.263286in}{1.658626in}}%
\pgfpathlineto{\pgfqpoint{1.264869in}{1.658439in}}%
\pgfpathlineto{\pgfqpoint{1.267668in}{1.681381in}}%
\pgfpathlineto{\pgfqpoint{1.269340in}{1.682158in}}%
\pgfpathlineto{\pgfqpoint{1.272031in}{1.704281in}}%
\pgfpathlineto{\pgfqpoint{1.277730in}{1.703596in}}%
\pgfpathlineto{\pgfqpoint{1.277023in}{1.697772in}}%
\pgfpathlineto{\pgfqpoint{1.282748in}{1.697095in}}%
\pgfpathlineto{\pgfqpoint{1.283442in}{1.702901in}}%
\pgfpathlineto{\pgfqpoint{1.302614in}{1.700688in}}%
\pgfpathlineto{\pgfqpoint{1.306273in}{1.700261in}}%
\pgfpathlineto{\pgfqpoint{1.304728in}{1.686756in}}%
\pgfpathlineto{\pgfqpoint{1.310508in}{1.686101in}}%
\pgfpathlineto{\pgfqpoint{1.310066in}{1.682269in}}%
\pgfpathlineto{\pgfqpoint{1.315821in}{1.681622in}}%
\pgfpathlineto{\pgfqpoint{1.317293in}{1.677601in}}%
\pgfpathlineto{\pgfqpoint{1.336444in}{1.675484in}}%
\pgfpathlineto{\pgfqpoint{1.335063in}{1.663866in}}%
\pgfpathlineto{\pgfqpoint{1.329154in}{1.662581in}}%
\pgfpathlineto{\pgfqpoint{1.327908in}{1.651084in}}%
\pgfpathlineto{\pgfqpoint{1.320048in}{1.651933in}}%
\pgfpathlineto{\pgfqpoint{1.318468in}{1.629024in}}%
\pgfpathclose%
\pgfusepath{fill}%
\end{pgfscope}%
\begin{pgfscope}%
\pgfpathrectangle{\pgfqpoint{0.100000in}{0.100000in}}{\pgfqpoint{3.007045in}{1.925000in}}%
\pgfusepath{clip}%
\pgfsetbuttcap%
\pgfsetmiterjoin%
\definecolor{currentfill}{rgb}{0.834833,0.935502,0.811011}%
\pgfsetfillcolor{currentfill}%
\pgfsetlinewidth{0.000000pt}%
\definecolor{currentstroke}{rgb}{0.000000,0.000000,0.000000}%
\pgfsetstrokecolor{currentstroke}%
\pgfsetstrokeopacity{0.000000}%
\pgfsetdash{}{0pt}%
\pgfpathmoveto{\pgfqpoint{1.694633in}{0.447645in}}%
\pgfpathlineto{\pgfqpoint{1.695745in}{0.469116in}}%
\pgfpathlineto{\pgfqpoint{1.676720in}{0.488443in}}%
\pgfpathlineto{\pgfqpoint{1.693878in}{0.498881in}}%
\pgfpathlineto{\pgfqpoint{1.696363in}{0.505939in}}%
\pgfpathlineto{\pgfqpoint{1.704139in}{0.513794in}}%
\pgfpathlineto{\pgfqpoint{1.710725in}{0.510845in}}%
\pgfpathlineto{\pgfqpoint{1.710733in}{0.503026in}}%
\pgfpathlineto{\pgfqpoint{1.714651in}{0.500322in}}%
\pgfpathlineto{\pgfqpoint{1.717041in}{0.493709in}}%
\pgfpathlineto{\pgfqpoint{1.722932in}{0.488802in}}%
\pgfpathlineto{\pgfqpoint{1.728488in}{0.492886in}}%
\pgfpathlineto{\pgfqpoint{1.734233in}{0.490553in}}%
\pgfpathlineto{\pgfqpoint{1.738650in}{0.493869in}}%
\pgfpathlineto{\pgfqpoint{1.740375in}{0.500409in}}%
\pgfpathlineto{\pgfqpoint{1.745433in}{0.500436in}}%
\pgfpathlineto{\pgfqpoint{1.747709in}{0.509837in}}%
\pgfpathlineto{\pgfqpoint{1.757144in}{0.510663in}}%
\pgfpathlineto{\pgfqpoint{1.759753in}{0.508170in}}%
\pgfpathlineto{\pgfqpoint{1.758777in}{0.502145in}}%
\pgfpathlineto{\pgfqpoint{1.769052in}{0.484340in}}%
\pgfpathlineto{\pgfqpoint{1.765360in}{0.476226in}}%
\pgfpathlineto{\pgfqpoint{1.749852in}{0.462071in}}%
\pgfpathlineto{\pgfqpoint{1.745282in}{0.461023in}}%
\pgfpathlineto{\pgfqpoint{1.724620in}{0.448821in}}%
\pgfpathlineto{\pgfqpoint{1.714126in}{0.445000in}}%
\pgfpathlineto{\pgfqpoint{1.711446in}{0.448095in}}%
\pgfpathlineto{\pgfqpoint{1.700720in}{0.443275in}}%
\pgfpathlineto{\pgfqpoint{1.700851in}{0.451264in}}%
\pgfpathlineto{\pgfqpoint{1.694633in}{0.447645in}}%
\pgfpathclose%
\pgfusepath{fill}%
\end{pgfscope}%
\begin{pgfscope}%
\pgfpathrectangle{\pgfqpoint{0.100000in}{0.100000in}}{\pgfqpoint{3.007045in}{1.925000in}}%
\pgfusepath{clip}%
\pgfsetbuttcap%
\pgfsetmiterjoin%
\definecolor{currentfill}{rgb}{0.401569,0.742484,0.436601}%
\pgfsetfillcolor{currentfill}%
\pgfsetlinewidth{0.000000pt}%
\definecolor{currentstroke}{rgb}{0.000000,0.000000,0.000000}%
\pgfsetstrokecolor{currentstroke}%
\pgfsetstrokeopacity{0.000000}%
\pgfsetdash{}{0pt}%
\pgfpathmoveto{\pgfqpoint{2.609375in}{1.212759in}}%
\pgfpathlineto{\pgfqpoint{2.603332in}{1.202130in}}%
\pgfpathlineto{\pgfqpoint{2.599182in}{1.199561in}}%
\pgfpathlineto{\pgfqpoint{2.596426in}{1.190754in}}%
\pgfpathlineto{\pgfqpoint{2.589115in}{1.195276in}}%
\pgfpathlineto{\pgfqpoint{2.586940in}{1.188956in}}%
\pgfpathlineto{\pgfqpoint{2.569924in}{1.199083in}}%
\pgfpathlineto{\pgfqpoint{2.571223in}{1.207045in}}%
\pgfpathlineto{\pgfqpoint{2.568471in}{1.207938in}}%
\pgfpathlineto{\pgfqpoint{2.569988in}{1.214538in}}%
\pgfpathlineto{\pgfqpoint{2.567300in}{1.215781in}}%
\pgfpathlineto{\pgfqpoint{2.560339in}{1.213292in}}%
\pgfpathlineto{\pgfqpoint{2.554976in}{1.247539in}}%
\pgfpathlineto{\pgfqpoint{2.559174in}{1.248298in}}%
\pgfpathlineto{\pgfqpoint{2.609634in}{1.257414in}}%
\pgfpathlineto{\pgfqpoint{2.613340in}{1.256035in}}%
\pgfpathlineto{\pgfqpoint{2.608755in}{1.249542in}}%
\pgfpathlineto{\pgfqpoint{2.607688in}{1.243264in}}%
\pgfpathlineto{\pgfqpoint{2.598024in}{1.242727in}}%
\pgfpathlineto{\pgfqpoint{2.592721in}{1.234363in}}%
\pgfpathlineto{\pgfqpoint{2.590743in}{1.227736in}}%
\pgfpathlineto{\pgfqpoint{2.585351in}{1.220532in}}%
\pgfpathlineto{\pgfqpoint{2.589864in}{1.218745in}}%
\pgfpathlineto{\pgfqpoint{2.604647in}{1.216037in}}%
\pgfpathlineto{\pgfqpoint{2.609375in}{1.212759in}}%
\pgfpathclose%
\pgfusepath{fill}%
\end{pgfscope}%
\begin{pgfscope}%
\pgfpathrectangle{\pgfqpoint{0.100000in}{0.100000in}}{\pgfqpoint{3.007045in}{1.925000in}}%
\pgfusepath{clip}%
\pgfsetbuttcap%
\pgfsetmiterjoin%
\definecolor{currentfill}{rgb}{0.513033,0.795755,0.510542}%
\pgfsetfillcolor{currentfill}%
\pgfsetlinewidth{0.000000pt}%
\definecolor{currentstroke}{rgb}{0.000000,0.000000,0.000000}%
\pgfsetstrokecolor{currentstroke}%
\pgfsetstrokeopacity{0.000000}%
\pgfsetdash{}{0pt}%
\pgfpathmoveto{\pgfqpoint{2.749311in}{1.495307in}}%
\pgfpathlineto{\pgfqpoint{2.743391in}{1.496243in}}%
\pgfpathlineto{\pgfqpoint{2.742514in}{1.505188in}}%
\pgfpathlineto{\pgfqpoint{2.739690in}{1.510307in}}%
\pgfpathlineto{\pgfqpoint{2.742037in}{1.517084in}}%
\pgfpathlineto{\pgfqpoint{2.739601in}{1.524017in}}%
\pgfpathlineto{\pgfqpoint{2.731549in}{1.525832in}}%
\pgfpathlineto{\pgfqpoint{2.733723in}{1.536267in}}%
\pgfpathlineto{\pgfqpoint{2.721573in}{1.573055in}}%
\pgfpathlineto{\pgfqpoint{2.735921in}{1.578293in}}%
\pgfpathlineto{\pgfqpoint{2.747323in}{1.582366in}}%
\pgfpathlineto{\pgfqpoint{2.750775in}{1.572935in}}%
\pgfpathlineto{\pgfqpoint{2.747693in}{1.569195in}}%
\pgfpathlineto{\pgfqpoint{2.755118in}{1.563156in}}%
\pgfpathlineto{\pgfqpoint{2.763075in}{1.564033in}}%
\pgfpathlineto{\pgfqpoint{2.763396in}{1.560589in}}%
\pgfpathlineto{\pgfqpoint{2.756305in}{1.557859in}}%
\pgfpathlineto{\pgfqpoint{2.774908in}{1.505824in}}%
\pgfpathlineto{\pgfqpoint{2.766923in}{1.496257in}}%
\pgfpathlineto{\pgfqpoint{2.759273in}{1.493756in}}%
\pgfpathlineto{\pgfqpoint{2.749311in}{1.495307in}}%
\pgfpathclose%
\pgfusepath{fill}%
\end{pgfscope}%
\begin{pgfscope}%
\pgfpathrectangle{\pgfqpoint{0.100000in}{0.100000in}}{\pgfqpoint{3.007045in}{1.925000in}}%
\pgfusepath{clip}%
\pgfsetbuttcap%
\pgfsetmiterjoin%
\definecolor{currentfill}{rgb}{0.579469,0.826759,0.565167}%
\pgfsetfillcolor{currentfill}%
\pgfsetlinewidth{0.000000pt}%
\definecolor{currentstroke}{rgb}{0.000000,0.000000,0.000000}%
\pgfsetstrokecolor{currentstroke}%
\pgfsetstrokeopacity{0.000000}%
\pgfsetdash{}{0pt}%
\pgfpathmoveto{\pgfqpoint{2.148381in}{1.158667in}}%
\pgfpathlineto{\pgfqpoint{2.142303in}{1.159619in}}%
\pgfpathlineto{\pgfqpoint{2.105428in}{1.156288in}}%
\pgfpathlineto{\pgfqpoint{2.105221in}{1.159153in}}%
\pgfpathlineto{\pgfqpoint{2.103104in}{1.187835in}}%
\pgfpathlineto{\pgfqpoint{2.123761in}{1.189439in}}%
\pgfpathlineto{\pgfqpoint{2.128488in}{1.192008in}}%
\pgfpathlineto{\pgfqpoint{2.128961in}{1.204856in}}%
\pgfpathlineto{\pgfqpoint{2.149319in}{1.206705in}}%
\pgfpathlineto{\pgfqpoint{2.152293in}{1.171574in}}%
\pgfpathlineto{\pgfqpoint{2.148720in}{1.168849in}}%
\pgfpathlineto{\pgfqpoint{2.150340in}{1.161666in}}%
\pgfpathlineto{\pgfqpoint{2.148381in}{1.158667in}}%
\pgfpathclose%
\pgfusepath{fill}%
\end{pgfscope}%
\begin{pgfscope}%
\pgfpathrectangle{\pgfqpoint{0.100000in}{0.100000in}}{\pgfqpoint{3.007045in}{1.925000in}}%
\pgfusepath{clip}%
\pgfsetbuttcap%
\pgfsetmiterjoin%
\definecolor{currentfill}{rgb}{0.474279,0.777670,0.478677}%
\pgfsetfillcolor{currentfill}%
\pgfsetlinewidth{0.000000pt}%
\definecolor{currentstroke}{rgb}{0.000000,0.000000,0.000000}%
\pgfsetstrokecolor{currentstroke}%
\pgfsetstrokeopacity{0.000000}%
\pgfsetdash{}{0pt}%
\pgfpathmoveto{\pgfqpoint{1.552729in}{1.550914in}}%
\pgfpathlineto{\pgfqpoint{1.527075in}{1.552396in}}%
\pgfpathlineto{\pgfqpoint{1.524722in}{1.557253in}}%
\pgfpathlineto{\pgfqpoint{1.526953in}{1.561321in}}%
\pgfpathlineto{\pgfqpoint{1.524166in}{1.567570in}}%
\pgfpathlineto{\pgfqpoint{1.518319in}{1.571035in}}%
\pgfpathlineto{\pgfqpoint{1.522749in}{1.580062in}}%
\pgfpathlineto{\pgfqpoint{1.526579in}{1.583044in}}%
\pgfpathlineto{\pgfqpoint{1.523621in}{1.593929in}}%
\pgfpathlineto{\pgfqpoint{1.518744in}{1.599173in}}%
\pgfpathlineto{\pgfqpoint{1.554716in}{1.596972in}}%
\pgfpathlineto{\pgfqpoint{1.552729in}{1.550914in}}%
\pgfpathclose%
\pgfusepath{fill}%
\end{pgfscope}%
\begin{pgfscope}%
\pgfpathrectangle{\pgfqpoint{0.100000in}{0.100000in}}{\pgfqpoint{3.007045in}{1.925000in}}%
\pgfusepath{clip}%
\pgfsetbuttcap%
\pgfsetmiterjoin%
\definecolor{currentfill}{rgb}{0.886505,0.956171,0.866128}%
\pgfsetfillcolor{currentfill}%
\pgfsetlinewidth{0.000000pt}%
\definecolor{currentstroke}{rgb}{0.000000,0.000000,0.000000}%
\pgfsetstrokecolor{currentstroke}%
\pgfsetstrokeopacity{0.000000}%
\pgfsetdash{}{0pt}%
\pgfpathmoveto{\pgfqpoint{0.720329in}{0.385044in}}%
\pgfpathlineto{\pgfqpoint{0.727658in}{0.379320in}}%
\pgfpathlineto{\pgfqpoint{0.723393in}{0.373987in}}%
\pgfpathlineto{\pgfqpoint{0.713956in}{0.381377in}}%
\pgfpathlineto{\pgfqpoint{0.717385in}{0.381921in}}%
\pgfpathlineto{\pgfqpoint{0.720189in}{0.383616in}}%
\pgfpathlineto{\pgfqpoint{0.720329in}{0.385044in}}%
\pgfpathclose%
\pgfusepath{fill}%
\end{pgfscope}%
\begin{pgfscope}%
\pgfpathrectangle{\pgfqpoint{0.100000in}{0.100000in}}{\pgfqpoint{3.007045in}{1.925000in}}%
\pgfusepath{clip}%
\pgfsetbuttcap%
\pgfsetmiterjoin%
\definecolor{currentfill}{rgb}{0.727797,0.891580,0.701730}%
\pgfsetfillcolor{currentfill}%
\pgfsetlinewidth{0.000000pt}%
\definecolor{currentstroke}{rgb}{0.000000,0.000000,0.000000}%
\pgfsetstrokecolor{currentstroke}%
\pgfsetstrokeopacity{0.000000}%
\pgfsetdash{}{0pt}%
\pgfpathmoveto{\pgfqpoint{1.483456in}{1.461953in}}%
\pgfpathlineto{\pgfqpoint{1.475452in}{1.459167in}}%
\pgfpathlineto{\pgfqpoint{1.473591in}{1.432656in}}%
\pgfpathlineto{\pgfqpoint{1.431343in}{1.435804in}}%
\pgfpathlineto{\pgfqpoint{1.430442in}{1.441625in}}%
\pgfpathlineto{\pgfqpoint{1.431599in}{1.456389in}}%
\pgfpathlineto{\pgfqpoint{1.437357in}{1.456409in}}%
\pgfpathlineto{\pgfqpoint{1.439374in}{1.475257in}}%
\pgfpathlineto{\pgfqpoint{1.442332in}{1.509325in}}%
\pgfpathlineto{\pgfqpoint{1.452730in}{1.510166in}}%
\pgfpathlineto{\pgfqpoint{1.458166in}{1.512912in}}%
\pgfpathlineto{\pgfqpoint{1.464688in}{1.511398in}}%
\pgfpathlineto{\pgfqpoint{1.473297in}{1.518824in}}%
\pgfpathlineto{\pgfqpoint{1.479161in}{1.518521in}}%
\pgfpathlineto{\pgfqpoint{1.484012in}{1.521937in}}%
\pgfpathlineto{\pgfqpoint{1.482943in}{1.518221in}}%
\pgfpathlineto{\pgfqpoint{1.481358in}{1.495327in}}%
\pgfpathlineto{\pgfqpoint{1.480026in}{1.483812in}}%
\pgfpathlineto{\pgfqpoint{1.485701in}{1.483457in}}%
\pgfpathlineto{\pgfqpoint{1.483456in}{1.461953in}}%
\pgfpathclose%
\pgfusepath{fill}%
\end{pgfscope}%
\begin{pgfscope}%
\pgfpathrectangle{\pgfqpoint{0.100000in}{0.100000in}}{\pgfqpoint{3.007045in}{1.925000in}}%
\pgfusepath{clip}%
\pgfsetbuttcap%
\pgfsetmiterjoin%
\definecolor{currentfill}{rgb}{0.585006,0.829343,0.569719}%
\pgfsetfillcolor{currentfill}%
\pgfsetlinewidth{0.000000pt}%
\definecolor{currentstroke}{rgb}{0.000000,0.000000,0.000000}%
\pgfsetstrokecolor{currentstroke}%
\pgfsetstrokeopacity{0.000000}%
\pgfsetdash{}{0pt}%
\pgfpathmoveto{\pgfqpoint{2.021691in}{1.270010in}}%
\pgfpathlineto{\pgfqpoint{2.021108in}{1.281507in}}%
\pgfpathlineto{\pgfqpoint{1.998647in}{1.280445in}}%
\pgfpathlineto{\pgfqpoint{1.998849in}{1.274698in}}%
\pgfpathlineto{\pgfqpoint{1.981668in}{1.274180in}}%
\pgfpathlineto{\pgfqpoint{1.973511in}{1.274071in}}%
\pgfpathlineto{\pgfqpoint{1.971270in}{1.279448in}}%
\pgfpathlineto{\pgfqpoint{1.964916in}{1.284730in}}%
\pgfpathlineto{\pgfqpoint{1.967742in}{1.296340in}}%
\pgfpathlineto{\pgfqpoint{1.980445in}{1.299472in}}%
\pgfpathlineto{\pgfqpoint{1.980128in}{1.308936in}}%
\pgfpathlineto{\pgfqpoint{1.974408in}{1.308715in}}%
\pgfpathlineto{\pgfqpoint{1.973965in}{1.320197in}}%
\pgfpathlineto{\pgfqpoint{1.996691in}{1.320570in}}%
\pgfpathlineto{\pgfqpoint{2.000310in}{1.318595in}}%
\pgfpathlineto{\pgfqpoint{2.006146in}{1.322530in}}%
\pgfpathlineto{\pgfqpoint{2.009626in}{1.309620in}}%
\pgfpathlineto{\pgfqpoint{2.025542in}{1.310434in}}%
\pgfpathlineto{\pgfqpoint{2.027069in}{1.287453in}}%
\pgfpathlineto{\pgfqpoint{2.037908in}{1.288046in}}%
\pgfpathlineto{\pgfqpoint{2.038925in}{1.270934in}}%
\pgfpathlineto{\pgfqpoint{2.021691in}{1.270010in}}%
\pgfpathclose%
\pgfusepath{fill}%
\end{pgfscope}%
\begin{pgfscope}%
\pgfpathrectangle{\pgfqpoint{0.100000in}{0.100000in}}{\pgfqpoint{3.007045in}{1.925000in}}%
\pgfusepath{clip}%
\pgfsetbuttcap%
\pgfsetmiterjoin%
\definecolor{currentfill}{rgb}{0.676371,0.869927,0.651657}%
\pgfsetfillcolor{currentfill}%
\pgfsetlinewidth{0.000000pt}%
\definecolor{currentstroke}{rgb}{0.000000,0.000000,0.000000}%
\pgfsetstrokecolor{currentstroke}%
\pgfsetstrokeopacity{0.000000}%
\pgfsetdash{}{0pt}%
\pgfpathmoveto{\pgfqpoint{2.273936in}{0.628314in}}%
\pgfpathlineto{\pgfqpoint{2.282576in}{0.629115in}}%
\pgfpathlineto{\pgfqpoint{2.284322in}{0.609983in}}%
\pgfpathlineto{\pgfqpoint{2.294588in}{0.610608in}}%
\pgfpathlineto{\pgfqpoint{2.291737in}{0.606164in}}%
\pgfpathlineto{\pgfqpoint{2.295880in}{0.596792in}}%
\pgfpathlineto{\pgfqpoint{2.288959in}{0.589362in}}%
\pgfpathlineto{\pgfqpoint{2.289539in}{0.581381in}}%
\pgfpathlineto{\pgfqpoint{2.275662in}{0.585133in}}%
\pgfpathlineto{\pgfqpoint{2.265886in}{0.586394in}}%
\pgfpathlineto{\pgfqpoint{2.252297in}{0.586303in}}%
\pgfpathlineto{\pgfqpoint{2.238850in}{0.583884in}}%
\pgfpathlineto{\pgfqpoint{2.221344in}{0.578728in}}%
\pgfpathlineto{\pgfqpoint{2.212863in}{0.577456in}}%
\pgfpathlineto{\pgfqpoint{2.202262in}{0.574181in}}%
\pgfpathlineto{\pgfqpoint{2.206102in}{0.581924in}}%
\pgfpathlineto{\pgfqpoint{2.210055in}{0.585128in}}%
\pgfpathlineto{\pgfqpoint{2.205014in}{0.590892in}}%
\pgfpathlineto{\pgfqpoint{2.207470in}{0.596991in}}%
\pgfpathlineto{\pgfqpoint{2.206456in}{0.600906in}}%
\pgfpathlineto{\pgfqpoint{2.199079in}{0.604683in}}%
\pgfpathlineto{\pgfqpoint{2.192424in}{0.612482in}}%
\pgfpathlineto{\pgfqpoint{2.193761in}{0.621372in}}%
\pgfpathlineto{\pgfqpoint{2.245494in}{0.625658in}}%
\pgfpathlineto{\pgfqpoint{2.273936in}{0.628314in}}%
\pgfpathclose%
\pgfusepath{fill}%
\end{pgfscope}%
\begin{pgfscope}%
\pgfpathrectangle{\pgfqpoint{0.100000in}{0.100000in}}{\pgfqpoint{3.007045in}{1.925000in}}%
\pgfusepath{clip}%
\pgfsetbuttcap%
\pgfsetmiterjoin%
\definecolor{currentfill}{rgb}{0.451765,0.767090,0.461207}%
\pgfsetfillcolor{currentfill}%
\pgfsetlinewidth{0.000000pt}%
\definecolor{currentstroke}{rgb}{0.000000,0.000000,0.000000}%
\pgfsetstrokecolor{currentstroke}%
\pgfsetstrokeopacity{0.000000}%
\pgfsetdash{}{0pt}%
\pgfpathmoveto{\pgfqpoint{1.539001in}{1.204970in}}%
\pgfpathlineto{\pgfqpoint{1.537445in}{1.176396in}}%
\pgfpathlineto{\pgfqpoint{1.510385in}{1.177845in}}%
\pgfpathlineto{\pgfqpoint{1.509436in}{1.177873in}}%
\pgfpathlineto{\pgfqpoint{1.511195in}{1.206454in}}%
\pgfpathlineto{\pgfqpoint{1.539001in}{1.204970in}}%
\pgfpathclose%
\pgfusepath{fill}%
\end{pgfscope}%
\begin{pgfscope}%
\pgfpathrectangle{\pgfqpoint{0.100000in}{0.100000in}}{\pgfqpoint{3.007045in}{1.925000in}}%
\pgfusepath{clip}%
\pgfsetbuttcap%
\pgfsetmiterjoin%
\definecolor{currentfill}{rgb}{0.357647,0.720953,0.415071}%
\pgfsetfillcolor{currentfill}%
\pgfsetlinewidth{0.000000pt}%
\definecolor{currentstroke}{rgb}{0.000000,0.000000,0.000000}%
\pgfsetstrokecolor{currentstroke}%
\pgfsetstrokeopacity{0.000000}%
\pgfsetdash{}{0pt}%
\pgfpathmoveto{\pgfqpoint{1.975423in}{0.932334in}}%
\pgfpathlineto{\pgfqpoint{1.969720in}{0.923885in}}%
\pgfpathlineto{\pgfqpoint{1.966993in}{0.915480in}}%
\pgfpathlineto{\pgfqpoint{1.967250in}{0.907985in}}%
\pgfpathlineto{\pgfqpoint{1.944456in}{0.907639in}}%
\pgfpathlineto{\pgfqpoint{1.944264in}{0.919360in}}%
\pgfpathlineto{\pgfqpoint{1.939696in}{0.925677in}}%
\pgfpathlineto{\pgfqpoint{1.940723in}{0.927923in}}%
\pgfpathlineto{\pgfqpoint{1.935018in}{0.933586in}}%
\pgfpathlineto{\pgfqpoint{1.927752in}{0.935378in}}%
\pgfpathlineto{\pgfqpoint{1.929258in}{0.940159in}}%
\pgfpathlineto{\pgfqpoint{1.921677in}{0.946416in}}%
\pgfpathlineto{\pgfqpoint{1.923472in}{0.954935in}}%
\pgfpathlineto{\pgfqpoint{1.954901in}{0.955275in}}%
\pgfpathlineto{\pgfqpoint{1.954694in}{0.961038in}}%
\pgfpathlineto{\pgfqpoint{1.960613in}{0.961138in}}%
\pgfpathlineto{\pgfqpoint{1.960431in}{0.971833in}}%
\pgfpathlineto{\pgfqpoint{1.962703in}{0.971887in}}%
\pgfpathlineto{\pgfqpoint{1.969967in}{0.964093in}}%
\pgfpathlineto{\pgfqpoint{1.971212in}{0.955492in}}%
\pgfpathlineto{\pgfqpoint{1.966442in}{0.954472in}}%
\pgfpathlineto{\pgfqpoint{1.966949in}{0.932026in}}%
\pgfpathlineto{\pgfqpoint{1.975423in}{0.932334in}}%
\pgfpathclose%
\pgfusepath{fill}%
\end{pgfscope}%
\begin{pgfscope}%
\pgfpathrectangle{\pgfqpoint{0.100000in}{0.100000in}}{\pgfqpoint{3.007045in}{1.925000in}}%
\pgfusepath{clip}%
\pgfsetbuttcap%
\pgfsetmiterjoin%
\definecolor{currentfill}{rgb}{0.732472,0.893549,0.706282}%
\pgfsetfillcolor{currentfill}%
\pgfsetlinewidth{0.000000pt}%
\definecolor{currentstroke}{rgb}{0.000000,0.000000,0.000000}%
\pgfsetstrokecolor{currentstroke}%
\pgfsetstrokeopacity{0.000000}%
\pgfsetdash{}{0pt}%
\pgfpathmoveto{\pgfqpoint{1.944978in}{1.116719in}}%
\pgfpathlineto{\pgfqpoint{1.945763in}{1.089375in}}%
\pgfpathlineto{\pgfqpoint{1.917287in}{1.088655in}}%
\pgfpathlineto{\pgfqpoint{1.911673in}{1.091410in}}%
\pgfpathlineto{\pgfqpoint{1.906201in}{1.091563in}}%
\pgfpathlineto{\pgfqpoint{1.906297in}{1.097100in}}%
\pgfpathlineto{\pgfqpoint{1.894852in}{1.097146in}}%
\pgfpathlineto{\pgfqpoint{1.883415in}{1.107800in}}%
\pgfpathlineto{\pgfqpoint{1.883385in}{1.113543in}}%
\pgfpathlineto{\pgfqpoint{1.872048in}{1.113887in}}%
\pgfpathlineto{\pgfqpoint{1.872212in}{1.129384in}}%
\pgfpathlineto{\pgfqpoint{1.872251in}{1.132295in}}%
\pgfpathlineto{\pgfqpoint{1.877944in}{1.138491in}}%
\pgfpathlineto{\pgfqpoint{1.877465in}{1.141048in}}%
\pgfpathlineto{\pgfqpoint{1.882047in}{1.149212in}}%
\pgfpathlineto{\pgfqpoint{1.884869in}{1.157037in}}%
\pgfpathlineto{\pgfqpoint{1.889132in}{1.155695in}}%
\pgfpathlineto{\pgfqpoint{1.903384in}{1.151278in}}%
\pgfpathlineto{\pgfqpoint{1.909292in}{1.151308in}}%
\pgfpathlineto{\pgfqpoint{1.920007in}{1.151193in}}%
\pgfpathlineto{\pgfqpoint{1.920082in}{1.139613in}}%
\pgfpathlineto{\pgfqpoint{1.944438in}{1.140114in}}%
\pgfpathlineto{\pgfqpoint{1.944978in}{1.116719in}}%
\pgfpathclose%
\pgfusepath{fill}%
\end{pgfscope}%
\begin{pgfscope}%
\pgfpathrectangle{\pgfqpoint{0.100000in}{0.100000in}}{\pgfqpoint{3.007045in}{1.925000in}}%
\pgfusepath{clip}%
\pgfsetbuttcap%
\pgfsetmiterjoin%
\definecolor{currentfill}{rgb}{0.845905,0.939931,0.822822}%
\pgfsetfillcolor{currentfill}%
\pgfsetlinewidth{0.000000pt}%
\definecolor{currentstroke}{rgb}{0.000000,0.000000,0.000000}%
\pgfsetstrokecolor{currentstroke}%
\pgfsetstrokeopacity{0.000000}%
\pgfsetdash{}{0pt}%
\pgfpathmoveto{\pgfqpoint{2.505019in}{0.681519in}}%
\pgfpathlineto{\pgfqpoint{2.503207in}{0.687158in}}%
\pgfpathlineto{\pgfqpoint{2.498935in}{0.687218in}}%
\pgfpathlineto{\pgfqpoint{2.496078in}{0.707079in}}%
\pgfpathlineto{\pgfqpoint{2.499221in}{0.711875in}}%
\pgfpathlineto{\pgfqpoint{2.496462in}{0.714530in}}%
\pgfpathlineto{\pgfqpoint{2.489650in}{0.715531in}}%
\pgfpathlineto{\pgfqpoint{2.489704in}{0.732476in}}%
\pgfpathlineto{\pgfqpoint{2.485433in}{0.742115in}}%
\pgfpathlineto{\pgfqpoint{2.509344in}{0.740112in}}%
\pgfpathlineto{\pgfqpoint{2.512057in}{0.734800in}}%
\pgfpathlineto{\pgfqpoint{2.513924in}{0.731882in}}%
\pgfpathlineto{\pgfqpoint{2.520931in}{0.729238in}}%
\pgfpathlineto{\pgfqpoint{2.526011in}{0.723942in}}%
\pgfpathlineto{\pgfqpoint{2.534070in}{0.723697in}}%
\pgfpathlineto{\pgfqpoint{2.537534in}{0.718203in}}%
\pgfpathlineto{\pgfqpoint{2.544812in}{0.713097in}}%
\pgfpathlineto{\pgfqpoint{2.549716in}{0.711373in}}%
\pgfpathlineto{\pgfqpoint{2.553497in}{0.705848in}}%
\pgfpathlineto{\pgfqpoint{2.550608in}{0.701257in}}%
\pgfpathlineto{\pgfqpoint{2.545964in}{0.706161in}}%
\pgfpathlineto{\pgfqpoint{2.536542in}{0.704088in}}%
\pgfpathlineto{\pgfqpoint{2.532623in}{0.707455in}}%
\pgfpathlineto{\pgfqpoint{2.529648in}{0.698799in}}%
\pgfpathlineto{\pgfqpoint{2.524353in}{0.695341in}}%
\pgfpathlineto{\pgfqpoint{2.528105in}{0.691406in}}%
\pgfpathlineto{\pgfqpoint{2.523188in}{0.681053in}}%
\pgfpathlineto{\pgfqpoint{2.518569in}{0.682840in}}%
\pgfpathlineto{\pgfqpoint{2.511887in}{0.680666in}}%
\pgfpathlineto{\pgfqpoint{2.505019in}{0.681519in}}%
\pgfpathclose%
\pgfusepath{fill}%
\end{pgfscope}%
\begin{pgfscope}%
\pgfpathrectangle{\pgfqpoint{0.100000in}{0.100000in}}{\pgfqpoint{3.007045in}{1.925000in}}%
\pgfusepath{clip}%
\pgfsetbuttcap%
\pgfsetmiterjoin%
\definecolor{currentfill}{rgb}{0.474279,0.777670,0.478677}%
\pgfsetfillcolor{currentfill}%
\pgfsetlinewidth{0.000000pt}%
\definecolor{currentstroke}{rgb}{0.000000,0.000000,0.000000}%
\pgfsetstrokecolor{currentstroke}%
\pgfsetstrokeopacity{0.000000}%
\pgfsetdash{}{0pt}%
\pgfpathmoveto{\pgfqpoint{1.704306in}{1.528001in}}%
\pgfpathlineto{\pgfqpoint{1.703956in}{1.505078in}}%
\pgfpathlineto{\pgfqpoint{1.703851in}{1.499277in}}%
\pgfpathlineto{\pgfqpoint{1.671967in}{1.499872in}}%
\pgfpathlineto{\pgfqpoint{1.637786in}{1.500890in}}%
\pgfpathlineto{\pgfqpoint{1.637825in}{1.506641in}}%
\pgfpathlineto{\pgfqpoint{1.632109in}{1.506861in}}%
\pgfpathlineto{\pgfqpoint{1.633052in}{1.541386in}}%
\pgfpathlineto{\pgfqpoint{1.668319in}{1.540277in}}%
\pgfpathlineto{\pgfqpoint{1.684336in}{1.539811in}}%
\pgfpathlineto{\pgfqpoint{1.684152in}{1.528327in}}%
\pgfpathlineto{\pgfqpoint{1.704306in}{1.528001in}}%
\pgfpathclose%
\pgfusepath{fill}%
\end{pgfscope}%
\begin{pgfscope}%
\pgfpathrectangle{\pgfqpoint{0.100000in}{0.100000in}}{\pgfqpoint{3.007045in}{1.925000in}}%
\pgfusepath{clip}%
\pgfsetbuttcap%
\pgfsetmiterjoin%
\definecolor{currentfill}{rgb}{0.490888,0.785421,0.492334}%
\pgfsetfillcolor{currentfill}%
\pgfsetlinewidth{0.000000pt}%
\definecolor{currentstroke}{rgb}{0.000000,0.000000,0.000000}%
\pgfsetstrokecolor{currentstroke}%
\pgfsetstrokeopacity{0.000000}%
\pgfsetdash{}{0pt}%
\pgfpathmoveto{\pgfqpoint{1.616016in}{0.651438in}}%
\pgfpathlineto{\pgfqpoint{1.586901in}{0.635387in}}%
\pgfpathlineto{\pgfqpoint{1.575122in}{0.656245in}}%
\pgfpathlineto{\pgfqpoint{1.565148in}{0.655103in}}%
\pgfpathlineto{\pgfqpoint{1.551784in}{0.680740in}}%
\pgfpathlineto{\pgfqpoint{1.577602in}{0.694365in}}%
\pgfpathlineto{\pgfqpoint{1.578119in}{0.708542in}}%
\pgfpathlineto{\pgfqpoint{1.600832in}{0.707630in}}%
\pgfpathlineto{\pgfqpoint{1.607101in}{0.688968in}}%
\pgfpathlineto{\pgfqpoint{1.611326in}{0.679134in}}%
\pgfpathlineto{\pgfqpoint{1.603281in}{0.674775in}}%
\pgfpathlineto{\pgfqpoint{1.616016in}{0.651438in}}%
\pgfpathclose%
\pgfusepath{fill}%
\end{pgfscope}%
\begin{pgfscope}%
\pgfpathrectangle{\pgfqpoint{0.100000in}{0.100000in}}{\pgfqpoint{3.007045in}{1.925000in}}%
\pgfusepath{clip}%
\pgfsetbuttcap%
\pgfsetmiterjoin%
\definecolor{currentfill}{rgb}{0.638970,0.854179,0.615240}%
\pgfsetfillcolor{currentfill}%
\pgfsetlinewidth{0.000000pt}%
\definecolor{currentstroke}{rgb}{0.000000,0.000000,0.000000}%
\pgfsetstrokecolor{currentstroke}%
\pgfsetstrokeopacity{0.000000}%
\pgfsetdash{}{0pt}%
\pgfpathmoveto{\pgfqpoint{2.609375in}{1.212759in}}%
\pgfpathlineto{\pgfqpoint{2.614276in}{1.211702in}}%
\pgfpathlineto{\pgfqpoint{2.616014in}{1.216485in}}%
\pgfpathlineto{\pgfqpoint{2.621536in}{1.211734in}}%
\pgfpathlineto{\pgfqpoint{2.622986in}{1.207692in}}%
\pgfpathlineto{\pgfqpoint{2.620005in}{1.199545in}}%
\pgfpathlineto{\pgfqpoint{2.625996in}{1.195904in}}%
\pgfpathlineto{\pgfqpoint{2.624882in}{1.186857in}}%
\pgfpathlineto{\pgfqpoint{2.621002in}{1.178691in}}%
\pgfpathlineto{\pgfqpoint{2.619686in}{1.171111in}}%
\pgfpathlineto{\pgfqpoint{2.613166in}{1.160149in}}%
\pgfpathlineto{\pgfqpoint{2.609006in}{1.155780in}}%
\pgfpathlineto{\pgfqpoint{2.581715in}{1.169068in}}%
\pgfpathlineto{\pgfqpoint{2.577995in}{1.164006in}}%
\pgfpathlineto{\pgfqpoint{2.569128in}{1.165501in}}%
\pgfpathlineto{\pgfqpoint{2.565022in}{1.171056in}}%
\pgfpathlineto{\pgfqpoint{2.558875in}{1.172717in}}%
\pgfpathlineto{\pgfqpoint{2.563877in}{1.187087in}}%
\pgfpathlineto{\pgfqpoint{2.562770in}{1.192250in}}%
\pgfpathlineto{\pgfqpoint{2.569924in}{1.199083in}}%
\pgfpathlineto{\pgfqpoint{2.586940in}{1.188956in}}%
\pgfpathlineto{\pgfqpoint{2.589115in}{1.195276in}}%
\pgfpathlineto{\pgfqpoint{2.596426in}{1.190754in}}%
\pgfpathlineto{\pgfqpoint{2.599182in}{1.199561in}}%
\pgfpathlineto{\pgfqpoint{2.603332in}{1.202130in}}%
\pgfpathlineto{\pgfqpoint{2.609375in}{1.212759in}}%
\pgfpathclose%
\pgfusepath{fill}%
\end{pgfscope}%
\begin{pgfscope}%
\pgfpathrectangle{\pgfqpoint{0.100000in}{0.100000in}}{\pgfqpoint{3.007045in}{1.925000in}}%
\pgfusepath{clip}%
\pgfsetbuttcap%
\pgfsetmiterjoin%
\definecolor{currentfill}{rgb}{0.551788,0.813841,0.542407}%
\pgfsetfillcolor{currentfill}%
\pgfsetlinewidth{0.000000pt}%
\definecolor{currentstroke}{rgb}{0.000000,0.000000,0.000000}%
\pgfsetstrokecolor{currentstroke}%
\pgfsetstrokeopacity{0.000000}%
\pgfsetdash{}{0pt}%
\pgfpathmoveto{\pgfqpoint{2.028875in}{0.983778in}}%
\pgfpathlineto{\pgfqpoint{2.025516in}{0.974832in}}%
\pgfpathlineto{\pgfqpoint{1.977561in}{0.972537in}}%
\pgfpathlineto{\pgfqpoint{1.977205in}{0.994031in}}%
\pgfpathlineto{\pgfqpoint{1.971458in}{0.993921in}}%
\pgfpathlineto{\pgfqpoint{1.971148in}{1.011094in}}%
\pgfpathlineto{\pgfqpoint{1.981594in}{1.012166in}}%
\pgfpathlineto{\pgfqpoint{1.984516in}{1.009519in}}%
\pgfpathlineto{\pgfqpoint{1.994242in}{1.009746in}}%
\pgfpathlineto{\pgfqpoint{1.996111in}{1.017492in}}%
\pgfpathlineto{\pgfqpoint{1.995876in}{1.024384in}}%
\pgfpathlineto{\pgfqpoint{2.005554in}{1.024952in}}%
\pgfpathlineto{\pgfqpoint{2.005455in}{1.027915in}}%
\pgfpathlineto{\pgfqpoint{2.022844in}{1.028625in}}%
\pgfpathlineto{\pgfqpoint{2.023343in}{1.013571in}}%
\pgfpathlineto{\pgfqpoint{2.029271in}{1.013848in}}%
\pgfpathlineto{\pgfqpoint{2.022047in}{1.002653in}}%
\pgfpathlineto{\pgfqpoint{2.028045in}{0.989826in}}%
\pgfpathlineto{\pgfqpoint{2.028875in}{0.983778in}}%
\pgfpathclose%
\pgfusepath{fill}%
\end{pgfscope}%
\begin{pgfscope}%
\pgfpathrectangle{\pgfqpoint{0.100000in}{0.100000in}}{\pgfqpoint{3.007045in}{1.925000in}}%
\pgfusepath{clip}%
\pgfsetbuttcap%
\pgfsetmiterjoin%
\definecolor{currentfill}{rgb}{0.585006,0.829343,0.569719}%
\pgfsetfillcolor{currentfill}%
\pgfsetlinewidth{0.000000pt}%
\definecolor{currentstroke}{rgb}{0.000000,0.000000,0.000000}%
\pgfsetstrokecolor{currentstroke}%
\pgfsetstrokeopacity{0.000000}%
\pgfsetdash{}{0pt}%
\pgfpathmoveto{\pgfqpoint{1.331685in}{1.819464in}}%
\pgfpathlineto{\pgfqpoint{1.375731in}{1.814864in}}%
\pgfpathlineto{\pgfqpoint{1.371755in}{1.774248in}}%
\pgfpathlineto{\pgfqpoint{1.346116in}{1.776879in}}%
\pgfpathlineto{\pgfqpoint{1.346744in}{1.782663in}}%
\pgfpathlineto{\pgfqpoint{1.341017in}{1.783290in}}%
\pgfpathlineto{\pgfqpoint{1.341649in}{1.789040in}}%
\pgfpathlineto{\pgfqpoint{1.332164in}{1.790080in}}%
\pgfpathlineto{\pgfqpoint{1.329938in}{1.796177in}}%
\pgfpathlineto{\pgfqpoint{1.331869in}{1.813477in}}%
\pgfpathlineto{\pgfqpoint{1.331685in}{1.819464in}}%
\pgfpathclose%
\pgfusepath{fill}%
\end{pgfscope}%
\begin{pgfscope}%
\pgfpathrectangle{\pgfqpoint{0.100000in}{0.100000in}}{\pgfqpoint{3.007045in}{1.925000in}}%
\pgfusepath{clip}%
\pgfsetbuttcap%
\pgfsetmiterjoin%
\definecolor{currentfill}{rgb}{0.382745,0.733256,0.427374}%
\pgfsetfillcolor{currentfill}%
\pgfsetlinewidth{0.000000pt}%
\definecolor{currentstroke}{rgb}{0.000000,0.000000,0.000000}%
\pgfsetstrokecolor{currentstroke}%
\pgfsetstrokeopacity{0.000000}%
\pgfsetdash{}{0pt}%
\pgfpathmoveto{\pgfqpoint{1.782023in}{1.648669in}}%
\pgfpathlineto{\pgfqpoint{1.764918in}{1.648716in}}%
\pgfpathlineto{\pgfqpoint{1.764876in}{1.642920in}}%
\pgfpathlineto{\pgfqpoint{1.718799in}{1.643280in}}%
\pgfpathlineto{\pgfqpoint{1.719048in}{1.660538in}}%
\pgfpathlineto{\pgfqpoint{1.718261in}{1.672097in}}%
\pgfpathlineto{\pgfqpoint{1.724037in}{1.672065in}}%
\pgfpathlineto{\pgfqpoint{1.724288in}{1.695202in}}%
\pgfpathlineto{\pgfqpoint{1.747387in}{1.695072in}}%
\pgfpathlineto{\pgfqpoint{1.746352in}{1.706621in}}%
\pgfpathlineto{\pgfqpoint{1.746314in}{1.723984in}}%
\pgfpathlineto{\pgfqpoint{1.745890in}{1.729762in}}%
\pgfpathlineto{\pgfqpoint{1.745680in}{1.764366in}}%
\pgfpathlineto{\pgfqpoint{1.762948in}{1.764365in}}%
\pgfpathlineto{\pgfqpoint{1.762966in}{1.752742in}}%
\pgfpathlineto{\pgfqpoint{1.797671in}{1.752902in}}%
\pgfpathlineto{\pgfqpoint{1.798346in}{1.718105in}}%
\pgfpathlineto{\pgfqpoint{1.798628in}{1.689186in}}%
\pgfpathlineto{\pgfqpoint{1.787190in}{1.689017in}}%
\pgfpathlineto{\pgfqpoint{1.787806in}{1.648648in}}%
\pgfpathlineto{\pgfqpoint{1.782023in}{1.648669in}}%
\pgfpathclose%
\pgfusepath{fill}%
\end{pgfscope}%
\begin{pgfscope}%
\pgfpathrectangle{\pgfqpoint{0.100000in}{0.100000in}}{\pgfqpoint{3.007045in}{1.925000in}}%
\pgfusepath{clip}%
\pgfsetbuttcap%
\pgfsetmiterjoin%
\definecolor{currentfill}{rgb}{0.741822,0.897486,0.715386}%
\pgfsetfillcolor{currentfill}%
\pgfsetlinewidth{0.000000pt}%
\definecolor{currentstroke}{rgb}{0.000000,0.000000,0.000000}%
\pgfsetstrokecolor{currentstroke}%
\pgfsetstrokeopacity{0.000000}%
\pgfsetdash{}{0pt}%
\pgfpathmoveto{\pgfqpoint{1.910648in}{0.511564in}}%
\pgfpathlineto{\pgfqpoint{1.909910in}{0.542133in}}%
\pgfpathlineto{\pgfqpoint{1.908086in}{0.571754in}}%
\pgfpathlineto{\pgfqpoint{1.910006in}{0.578352in}}%
\pgfpathlineto{\pgfqpoint{1.909431in}{0.599271in}}%
\pgfpathlineto{\pgfqpoint{1.915657in}{0.600227in}}%
\pgfpathlineto{\pgfqpoint{1.921122in}{0.606604in}}%
\pgfpathlineto{\pgfqpoint{1.927440in}{0.602431in}}%
\pgfpathlineto{\pgfqpoint{1.930070in}{0.596671in}}%
\pgfpathlineto{\pgfqpoint{1.954185in}{0.597454in}}%
\pgfpathlineto{\pgfqpoint{1.959192in}{0.587885in}}%
\pgfpathlineto{\pgfqpoint{1.958167in}{0.578691in}}%
\pgfpathlineto{\pgfqpoint{1.959526in}{0.575009in}}%
\pgfpathlineto{\pgfqpoint{1.965180in}{0.570849in}}%
\pgfpathlineto{\pgfqpoint{1.966602in}{0.562051in}}%
\pgfpathlineto{\pgfqpoint{1.970123in}{0.557651in}}%
\pgfpathlineto{\pgfqpoint{1.975186in}{0.557014in}}%
\pgfpathlineto{\pgfqpoint{1.976175in}{0.548583in}}%
\pgfpathlineto{\pgfqpoint{1.990185in}{0.543999in}}%
\pgfpathlineto{\pgfqpoint{1.988235in}{0.542241in}}%
\pgfpathlineto{\pgfqpoint{1.992985in}{0.532196in}}%
\pgfpathlineto{\pgfqpoint{1.997286in}{0.531031in}}%
\pgfpathlineto{\pgfqpoint{1.998738in}{0.526353in}}%
\pgfpathlineto{\pgfqpoint{1.993498in}{0.524333in}}%
\pgfpathlineto{\pgfqpoint{1.981476in}{0.527516in}}%
\pgfpathlineto{\pgfqpoint{1.982084in}{0.530926in}}%
\pgfpathlineto{\pgfqpoint{1.976086in}{0.538805in}}%
\pgfpathlineto{\pgfqpoint{1.967294in}{0.537756in}}%
\pgfpathlineto{\pgfqpoint{1.962543in}{0.530961in}}%
\pgfpathlineto{\pgfqpoint{1.953040in}{0.522060in}}%
\pgfpathlineto{\pgfqpoint{1.952704in}{0.529849in}}%
\pgfpathlineto{\pgfqpoint{1.943821in}{0.525996in}}%
\pgfpathlineto{\pgfqpoint{1.937749in}{0.521438in}}%
\pgfpathlineto{\pgfqpoint{1.939627in}{0.515520in}}%
\pgfpathlineto{\pgfqpoint{1.946540in}{0.514748in}}%
\pgfpathlineto{\pgfqpoint{1.955604in}{0.516241in}}%
\pgfpathlineto{\pgfqpoint{1.963363in}{0.511582in}}%
\pgfpathlineto{\pgfqpoint{1.957159in}{0.505997in}}%
\pgfpathlineto{\pgfqpoint{1.943062in}{0.512593in}}%
\pgfpathlineto{\pgfqpoint{1.937453in}{0.512510in}}%
\pgfpathlineto{\pgfqpoint{1.928770in}{0.508924in}}%
\pgfpathlineto{\pgfqpoint{1.910648in}{0.511564in}}%
\pgfpathclose%
\pgfusepath{fill}%
\end{pgfscope}%
\begin{pgfscope}%
\pgfpathrectangle{\pgfqpoint{0.100000in}{0.100000in}}{\pgfqpoint{3.007045in}{1.925000in}}%
\pgfusepath{clip}%
\pgfsetbuttcap%
\pgfsetmiterjoin%
\definecolor{currentfill}{rgb}{0.713772,0.885675,0.688074}%
\pgfsetfillcolor{currentfill}%
\pgfsetlinewidth{0.000000pt}%
\definecolor{currentstroke}{rgb}{0.000000,0.000000,0.000000}%
\pgfsetstrokecolor{currentstroke}%
\pgfsetstrokeopacity{0.000000}%
\pgfsetdash{}{0pt}%
\pgfpathmoveto{\pgfqpoint{1.547367in}{0.412624in}}%
\pgfpathlineto{\pgfqpoint{1.512666in}{0.412373in}}%
\pgfpathlineto{\pgfqpoint{1.512921in}{0.424052in}}%
\pgfpathlineto{\pgfqpoint{1.514242in}{0.453233in}}%
\pgfpathlineto{\pgfqpoint{1.513462in}{0.453257in}}%
\pgfpathlineto{\pgfqpoint{1.514641in}{0.483384in}}%
\pgfpathlineto{\pgfqpoint{1.550168in}{0.481716in}}%
\pgfpathlineto{\pgfqpoint{1.549167in}{0.452106in}}%
\pgfpathlineto{\pgfqpoint{1.547367in}{0.412624in}}%
\pgfpathclose%
\pgfusepath{fill}%
\end{pgfscope}%
\begin{pgfscope}%
\pgfpathrectangle{\pgfqpoint{0.100000in}{0.100000in}}{\pgfqpoint{3.007045in}{1.925000in}}%
\pgfusepath{clip}%
\pgfsetbuttcap%
\pgfsetmiterjoin%
\definecolor{currentfill}{rgb}{0.727797,0.891580,0.701730}%
\pgfsetfillcolor{currentfill}%
\pgfsetlinewidth{0.000000pt}%
\definecolor{currentstroke}{rgb}{0.000000,0.000000,0.000000}%
\pgfsetstrokecolor{currentstroke}%
\pgfsetstrokeopacity{0.000000}%
\pgfsetdash{}{0pt}%
\pgfpathmoveto{\pgfqpoint{2.488716in}{0.626214in}}%
\pgfpathlineto{\pgfqpoint{2.481583in}{0.625736in}}%
\pgfpathlineto{\pgfqpoint{2.479352in}{0.633603in}}%
\pgfpathlineto{\pgfqpoint{2.469936in}{0.633604in}}%
\pgfpathlineto{\pgfqpoint{2.464523in}{0.639748in}}%
\pgfpathlineto{\pgfqpoint{2.457148in}{0.641063in}}%
\pgfpathlineto{\pgfqpoint{2.454390in}{0.661837in}}%
\pgfpathlineto{\pgfqpoint{2.450133in}{0.661231in}}%
\pgfpathlineto{\pgfqpoint{2.450078in}{0.667343in}}%
\pgfpathlineto{\pgfqpoint{2.443165in}{0.674743in}}%
\pgfpathlineto{\pgfqpoint{2.442052in}{0.679496in}}%
\pgfpathlineto{\pgfqpoint{2.443946in}{0.680145in}}%
\pgfpathlineto{\pgfqpoint{2.445730in}{0.689304in}}%
\pgfpathlineto{\pgfqpoint{2.448668in}{0.692985in}}%
\pgfpathlineto{\pgfqpoint{2.447789in}{0.700910in}}%
\pgfpathlineto{\pgfqpoint{2.455203in}{0.702082in}}%
\pgfpathlineto{\pgfqpoint{2.457652in}{0.694920in}}%
\pgfpathlineto{\pgfqpoint{2.477917in}{0.698299in}}%
\pgfpathlineto{\pgfqpoint{2.492063in}{0.691931in}}%
\pgfpathlineto{\pgfqpoint{2.498935in}{0.687218in}}%
\pgfpathlineto{\pgfqpoint{2.503207in}{0.687158in}}%
\pgfpathlineto{\pgfqpoint{2.505019in}{0.681519in}}%
\pgfpathlineto{\pgfqpoint{2.508359in}{0.677599in}}%
\pgfpathlineto{\pgfqpoint{2.503124in}{0.674785in}}%
\pgfpathlineto{\pgfqpoint{2.497961in}{0.669319in}}%
\pgfpathlineto{\pgfqpoint{2.498304in}{0.661191in}}%
\pgfpathlineto{\pgfqpoint{2.503300in}{0.656930in}}%
\pgfpathlineto{\pgfqpoint{2.487249in}{0.654855in}}%
\pgfpathlineto{\pgfqpoint{2.488999in}{0.640420in}}%
\pgfpathlineto{\pgfqpoint{2.504409in}{0.641856in}}%
\pgfpathlineto{\pgfqpoint{2.502750in}{0.627127in}}%
\pgfpathlineto{\pgfqpoint{2.488716in}{0.626214in}}%
\pgfpathclose%
\pgfusepath{fill}%
\end{pgfscope}%
\begin{pgfscope}%
\pgfpathrectangle{\pgfqpoint{0.100000in}{0.100000in}}{\pgfqpoint{3.007045in}{1.925000in}}%
\pgfusepath{clip}%
\pgfsetbuttcap%
\pgfsetmiterjoin%
\definecolor{currentfill}{rgb}{0.573933,0.824175,0.560615}%
\pgfsetfillcolor{currentfill}%
\pgfsetlinewidth{0.000000pt}%
\definecolor{currentstroke}{rgb}{0.000000,0.000000,0.000000}%
\pgfsetstrokecolor{currentstroke}%
\pgfsetstrokeopacity{0.000000}%
\pgfsetdash{}{0pt}%
\pgfpathmoveto{\pgfqpoint{0.499509in}{1.627848in}}%
\pgfpathlineto{\pgfqpoint{0.492514in}{1.625981in}}%
\pgfpathlineto{\pgfqpoint{0.490196in}{1.618627in}}%
\pgfpathlineto{\pgfqpoint{0.462254in}{1.627100in}}%
\pgfpathlineto{\pgfqpoint{0.464341in}{1.633883in}}%
\pgfpathlineto{\pgfqpoint{0.447557in}{1.638893in}}%
\pgfpathlineto{\pgfqpoint{0.447535in}{1.643476in}}%
\pgfpathlineto{\pgfqpoint{0.451005in}{1.654418in}}%
\pgfpathlineto{\pgfqpoint{0.441475in}{1.657513in}}%
\pgfpathlineto{\pgfqpoint{0.434242in}{1.663279in}}%
\pgfpathlineto{\pgfqpoint{0.432010in}{1.670306in}}%
\pgfpathlineto{\pgfqpoint{0.423095in}{1.674955in}}%
\pgfpathlineto{\pgfqpoint{0.416059in}{1.671542in}}%
\pgfpathlineto{\pgfqpoint{0.406441in}{1.674762in}}%
\pgfpathlineto{\pgfqpoint{0.413049in}{1.689210in}}%
\pgfpathlineto{\pgfqpoint{0.416964in}{1.700001in}}%
\pgfpathlineto{\pgfqpoint{0.432427in}{1.695452in}}%
\pgfpathlineto{\pgfqpoint{0.436719in}{1.699407in}}%
\pgfpathlineto{\pgfqpoint{0.437975in}{1.704042in}}%
\pgfpathlineto{\pgfqpoint{0.443505in}{1.702226in}}%
\pgfpathlineto{\pgfqpoint{0.449171in}{1.720447in}}%
\pgfpathlineto{\pgfqpoint{0.444068in}{1.723357in}}%
\pgfpathlineto{\pgfqpoint{0.451037in}{1.744559in}}%
\pgfpathlineto{\pgfqpoint{0.448349in}{1.745454in}}%
\pgfpathlineto{\pgfqpoint{0.451249in}{1.754244in}}%
\pgfpathlineto{\pgfqpoint{0.465504in}{1.749616in}}%
\pgfpathlineto{\pgfqpoint{0.469874in}{1.763357in}}%
\pgfpathlineto{\pgfqpoint{0.484457in}{1.758626in}}%
\pgfpathlineto{\pgfqpoint{0.488462in}{1.755366in}}%
\pgfpathlineto{\pgfqpoint{0.490020in}{1.750804in}}%
\pgfpathlineto{\pgfqpoint{0.496816in}{1.743625in}}%
\pgfpathlineto{\pgfqpoint{0.496421in}{1.737333in}}%
\pgfpathlineto{\pgfqpoint{0.493784in}{1.734570in}}%
\pgfpathlineto{\pgfqpoint{0.496057in}{1.728942in}}%
\pgfpathlineto{\pgfqpoint{0.500179in}{1.724897in}}%
\pgfpathlineto{\pgfqpoint{0.501234in}{1.718853in}}%
\pgfpathlineto{\pgfqpoint{0.506524in}{1.713649in}}%
\pgfpathlineto{\pgfqpoint{0.536189in}{1.704696in}}%
\pgfpathlineto{\pgfqpoint{0.533873in}{1.701067in}}%
\pgfpathlineto{\pgfqpoint{0.530035in}{1.695577in}}%
\pgfpathlineto{\pgfqpoint{0.527922in}{1.686852in}}%
\pgfpathlineto{\pgfqpoint{0.523556in}{1.681083in}}%
\pgfpathlineto{\pgfqpoint{0.521156in}{1.673183in}}%
\pgfpathlineto{\pgfqpoint{0.521441in}{1.664987in}}%
\pgfpathlineto{\pgfqpoint{0.519173in}{1.654672in}}%
\pgfpathlineto{\pgfqpoint{0.512713in}{1.647749in}}%
\pgfpathlineto{\pgfqpoint{0.505425in}{1.642739in}}%
\pgfpathlineto{\pgfqpoint{0.504292in}{1.636653in}}%
\pgfpathlineto{\pgfqpoint{0.499509in}{1.627848in}}%
\pgfpathclose%
\pgfusepath{fill}%
\end{pgfscope}%
\begin{pgfscope}%
\pgfpathrectangle{\pgfqpoint{0.100000in}{0.100000in}}{\pgfqpoint{3.007045in}{1.925000in}}%
\pgfusepath{clip}%
\pgfsetbuttcap%
\pgfsetmiterjoin%
\definecolor{currentfill}{rgb}{0.709097,0.883706,0.683522}%
\pgfsetfillcolor{currentfill}%
\pgfsetlinewidth{0.000000pt}%
\definecolor{currentstroke}{rgb}{0.000000,0.000000,0.000000}%
\pgfsetstrokecolor{currentstroke}%
\pgfsetstrokeopacity{0.000000}%
\pgfsetdash{}{0pt}%
\pgfpathmoveto{\pgfqpoint{2.321001in}{0.750113in}}%
\pgfpathlineto{\pgfqpoint{2.312802in}{0.748173in}}%
\pgfpathlineto{\pgfqpoint{2.295738in}{0.746225in}}%
\pgfpathlineto{\pgfqpoint{2.293144in}{0.771183in}}%
\pgfpathlineto{\pgfqpoint{2.289813in}{0.770791in}}%
\pgfpathlineto{\pgfqpoint{2.287725in}{0.796480in}}%
\pgfpathlineto{\pgfqpoint{2.281013in}{0.795971in}}%
\pgfpathlineto{\pgfqpoint{2.278832in}{0.801530in}}%
\pgfpathlineto{\pgfqpoint{2.267838in}{0.800489in}}%
\pgfpathlineto{\pgfqpoint{2.264190in}{0.805986in}}%
\pgfpathlineto{\pgfqpoint{2.258979in}{0.805773in}}%
\pgfpathlineto{\pgfqpoint{2.263640in}{0.812852in}}%
\pgfpathlineto{\pgfqpoint{2.262319in}{0.816941in}}%
\pgfpathlineto{\pgfqpoint{2.271634in}{0.825054in}}%
\pgfpathlineto{\pgfqpoint{2.279304in}{0.827039in}}%
\pgfpathlineto{\pgfqpoint{2.290820in}{0.827401in}}%
\pgfpathlineto{\pgfqpoint{2.297886in}{0.828687in}}%
\pgfpathlineto{\pgfqpoint{2.298990in}{0.824653in}}%
\pgfpathlineto{\pgfqpoint{2.317320in}{0.826818in}}%
\pgfpathlineto{\pgfqpoint{2.318665in}{0.820809in}}%
\pgfpathlineto{\pgfqpoint{2.326336in}{0.819580in}}%
\pgfpathlineto{\pgfqpoint{2.327623in}{0.805924in}}%
\pgfpathlineto{\pgfqpoint{2.332358in}{0.802455in}}%
\pgfpathlineto{\pgfqpoint{2.343332in}{0.803548in}}%
\pgfpathlineto{\pgfqpoint{2.346769in}{0.793773in}}%
\pgfpathlineto{\pgfqpoint{2.352410in}{0.787703in}}%
\pgfpathlineto{\pgfqpoint{2.352426in}{0.785269in}}%
\pgfpathlineto{\pgfqpoint{2.333138in}{0.783244in}}%
\pgfpathlineto{\pgfqpoint{2.335233in}{0.759815in}}%
\pgfpathlineto{\pgfqpoint{2.317396in}{0.757985in}}%
\pgfpathlineto{\pgfqpoint{2.321001in}{0.750113in}}%
\pgfpathclose%
\pgfusepath{fill}%
\end{pgfscope}%
\begin{pgfscope}%
\pgfpathrectangle{\pgfqpoint{0.100000in}{0.100000in}}{\pgfqpoint{3.007045in}{1.925000in}}%
\pgfusepath{clip}%
\pgfsetbuttcap%
\pgfsetmiterjoin%
\definecolor{currentfill}{rgb}{0.363922,0.724029,0.418147}%
\pgfsetfillcolor{currentfill}%
\pgfsetlinewidth{0.000000pt}%
\definecolor{currentstroke}{rgb}{0.000000,0.000000,0.000000}%
\pgfsetstrokecolor{currentstroke}%
\pgfsetstrokeopacity{0.000000}%
\pgfsetdash{}{0pt}%
\pgfpathmoveto{\pgfqpoint{1.104848in}{1.099688in}}%
\pgfpathlineto{\pgfqpoint{1.098980in}{1.105934in}}%
\pgfpathlineto{\pgfqpoint{1.087280in}{1.105261in}}%
\pgfpathlineto{\pgfqpoint{1.085598in}{1.110401in}}%
\pgfpathlineto{\pgfqpoint{1.056798in}{1.114784in}}%
\pgfpathlineto{\pgfqpoint{1.045110in}{1.115632in}}%
\pgfpathlineto{\pgfqpoint{1.047972in}{1.134089in}}%
\pgfpathlineto{\pgfqpoint{1.048160in}{1.141447in}}%
\pgfpathlineto{\pgfqpoint{1.050444in}{1.156057in}}%
\pgfpathlineto{\pgfqpoint{1.059704in}{1.212390in}}%
\pgfpathlineto{\pgfqpoint{1.091973in}{1.207431in}}%
\pgfpathlineto{\pgfqpoint{1.141305in}{1.200383in}}%
\pgfpathlineto{\pgfqpoint{1.139110in}{1.197239in}}%
\pgfpathlineto{\pgfqpoint{1.142160in}{1.192960in}}%
\pgfpathlineto{\pgfqpoint{1.136431in}{1.191196in}}%
\pgfpathlineto{\pgfqpoint{1.128210in}{1.131426in}}%
\pgfpathlineto{\pgfqpoint{1.121297in}{1.132386in}}%
\pgfpathlineto{\pgfqpoint{1.123455in}{1.127357in}}%
\pgfpathlineto{\pgfqpoint{1.123423in}{1.117299in}}%
\pgfpathlineto{\pgfqpoint{1.121625in}{1.109932in}}%
\pgfpathlineto{\pgfqpoint{1.112363in}{1.107267in}}%
\pgfpathlineto{\pgfqpoint{1.104848in}{1.099688in}}%
\pgfpathclose%
\pgfusepath{fill}%
\end{pgfscope}%
\begin{pgfscope}%
\pgfpathrectangle{\pgfqpoint{0.100000in}{0.100000in}}{\pgfqpoint{3.007045in}{1.925000in}}%
\pgfusepath{clip}%
\pgfsetbuttcap%
\pgfsetmiterjoin%
\definecolor{currentfill}{rgb}{0.886505,0.956171,0.866128}%
\pgfsetfillcolor{currentfill}%
\pgfsetlinewidth{0.000000pt}%
\definecolor{currentstroke}{rgb}{0.000000,0.000000,0.000000}%
\pgfsetstrokecolor{currentstroke}%
\pgfsetstrokeopacity{0.000000}%
\pgfsetdash{}{0pt}%
\pgfpathmoveto{\pgfqpoint{1.001047in}{0.482991in}}%
\pgfpathlineto{\pgfqpoint{0.989639in}{0.489920in}}%
\pgfpathlineto{\pgfqpoint{0.966890in}{0.503916in}}%
\pgfpathlineto{\pgfqpoint{0.959476in}{0.508253in}}%
\pgfpathlineto{\pgfqpoint{0.952820in}{0.498342in}}%
\pgfpathlineto{\pgfqpoint{0.944682in}{0.503522in}}%
\pgfpathlineto{\pgfqpoint{0.941118in}{0.505304in}}%
\pgfpathlineto{\pgfqpoint{0.919823in}{0.519872in}}%
\pgfpathlineto{\pgfqpoint{0.895892in}{0.537283in}}%
\pgfpathlineto{\pgfqpoint{0.877896in}{0.551221in}}%
\pgfpathlineto{\pgfqpoint{0.878340in}{0.551804in}}%
\pgfpathlineto{\pgfqpoint{0.871135in}{0.557489in}}%
\pgfpathlineto{\pgfqpoint{0.872509in}{0.559210in}}%
\pgfpathlineto{\pgfqpoint{0.865552in}{0.564907in}}%
\pgfpathlineto{\pgfqpoint{0.867050in}{0.566744in}}%
\pgfpathlineto{\pgfqpoint{0.859713in}{0.572782in}}%
\pgfpathlineto{\pgfqpoint{0.847780in}{0.582922in}}%
\pgfpathlineto{\pgfqpoint{0.848948in}{0.584275in}}%
\pgfpathlineto{\pgfqpoint{0.845795in}{0.587009in}}%
\pgfpathlineto{\pgfqpoint{0.844619in}{0.585660in}}%
\pgfpathlineto{\pgfqpoint{0.837473in}{0.591959in}}%
\pgfpathlineto{\pgfqpoint{0.838673in}{0.593306in}}%
\pgfpathlineto{\pgfqpoint{0.835607in}{0.596048in}}%
\pgfpathlineto{\pgfqpoint{0.834401in}{0.594707in}}%
\pgfpathlineto{\pgfqpoint{0.830895in}{0.597873in}}%
\pgfpathlineto{\pgfqpoint{0.832109in}{0.599208in}}%
\pgfpathlineto{\pgfqpoint{0.826925in}{0.603954in}}%
\pgfpathlineto{\pgfqpoint{0.825701in}{0.602628in}}%
\pgfpathlineto{\pgfqpoint{0.813031in}{0.614554in}}%
\pgfpathlineto{\pgfqpoint{0.809438in}{0.610810in}}%
\pgfpathlineto{\pgfqpoint{0.803267in}{0.616806in}}%
\pgfpathlineto{\pgfqpoint{0.800883in}{0.622716in}}%
\pgfpathlineto{\pgfqpoint{0.801608in}{0.625470in}}%
\pgfpathlineto{\pgfqpoint{0.801127in}{0.628309in}}%
\pgfpathlineto{\pgfqpoint{0.799100in}{0.632308in}}%
\pgfpathlineto{\pgfqpoint{0.801794in}{0.631253in}}%
\pgfpathlineto{\pgfqpoint{0.804644in}{0.631085in}}%
\pgfpathlineto{\pgfqpoint{0.806937in}{0.632414in}}%
\pgfpathlineto{\pgfqpoint{0.812022in}{0.637296in}}%
\pgfpathlineto{\pgfqpoint{0.813808in}{0.636711in}}%
\pgfpathlineto{\pgfqpoint{0.817714in}{0.632498in}}%
\pgfpathlineto{\pgfqpoint{0.826485in}{0.626672in}}%
\pgfpathlineto{\pgfqpoint{0.829143in}{0.625614in}}%
\pgfpathlineto{\pgfqpoint{0.834897in}{0.625269in}}%
\pgfpathlineto{\pgfqpoint{0.842826in}{0.626872in}}%
\pgfpathlineto{\pgfqpoint{0.846807in}{0.629874in}}%
\pgfpathlineto{\pgfqpoint{0.847784in}{0.631457in}}%
\pgfpathlineto{\pgfqpoint{0.850593in}{0.633815in}}%
\pgfpathlineto{\pgfqpoint{0.852477in}{0.634580in}}%
\pgfpathlineto{\pgfqpoint{0.858398in}{0.634527in}}%
\pgfpathlineto{\pgfqpoint{0.860204in}{0.636304in}}%
\pgfpathlineto{\pgfqpoint{0.864117in}{0.636616in}}%
\pgfpathlineto{\pgfqpoint{0.868789in}{0.637634in}}%
\pgfpathlineto{\pgfqpoint{0.871460in}{0.633851in}}%
\pgfpathlineto{\pgfqpoint{0.876444in}{0.632688in}}%
\pgfpathlineto{\pgfqpoint{0.885282in}{0.633242in}}%
\pgfpathlineto{\pgfqpoint{0.894306in}{0.634804in}}%
\pgfpathlineto{\pgfqpoint{0.897261in}{0.633456in}}%
\pgfpathlineto{\pgfqpoint{0.897619in}{0.631397in}}%
\pgfpathlineto{\pgfqpoint{0.901143in}{0.628563in}}%
\pgfpathlineto{\pgfqpoint{0.904509in}{0.626669in}}%
\pgfpathlineto{\pgfqpoint{0.907341in}{0.625778in}}%
\pgfpathlineto{\pgfqpoint{0.913294in}{0.626230in}}%
\pgfpathlineto{\pgfqpoint{0.921174in}{0.629150in}}%
\pgfpathlineto{\pgfqpoint{0.923879in}{0.629562in}}%
\pgfpathlineto{\pgfqpoint{0.925211in}{0.625453in}}%
\pgfpathlineto{\pgfqpoint{0.924492in}{0.622843in}}%
\pgfpathlineto{\pgfqpoint{0.926548in}{0.622888in}}%
\pgfpathlineto{\pgfqpoint{0.927653in}{0.620890in}}%
\pgfpathlineto{\pgfqpoint{0.927665in}{0.618526in}}%
\pgfpathlineto{\pgfqpoint{0.923257in}{0.616818in}}%
\pgfpathlineto{\pgfqpoint{0.919869in}{0.618044in}}%
\pgfpathlineto{\pgfqpoint{0.920400in}{0.614171in}}%
\pgfpathlineto{\pgfqpoint{0.923838in}{0.612830in}}%
\pgfpathlineto{\pgfqpoint{0.925840in}{0.614446in}}%
\pgfpathlineto{\pgfqpoint{0.930393in}{0.614357in}}%
\pgfpathlineto{\pgfqpoint{0.932785in}{0.611562in}}%
\pgfpathlineto{\pgfqpoint{0.930092in}{0.608495in}}%
\pgfpathlineto{\pgfqpoint{0.929589in}{0.606710in}}%
\pgfpathlineto{\pgfqpoint{0.931595in}{0.605231in}}%
\pgfpathlineto{\pgfqpoint{0.931935in}{0.603180in}}%
\pgfpathlineto{\pgfqpoint{0.934772in}{0.603947in}}%
\pgfpathlineto{\pgfqpoint{0.938741in}{0.601961in}}%
\pgfpathlineto{\pgfqpoint{0.941530in}{0.601138in}}%
\pgfpathlineto{\pgfqpoint{0.943304in}{0.596728in}}%
\pgfpathlineto{\pgfqpoint{0.945069in}{0.596940in}}%
\pgfpathlineto{\pgfqpoint{0.946594in}{0.594438in}}%
\pgfpathlineto{\pgfqpoint{0.941495in}{0.590893in}}%
\pgfpathlineto{\pgfqpoint{0.939975in}{0.589037in}}%
\pgfpathlineto{\pgfqpoint{0.945374in}{0.585331in}}%
\pgfpathlineto{\pgfqpoint{0.944737in}{0.583986in}}%
\pgfpathlineto{\pgfqpoint{0.942095in}{0.583582in}}%
\pgfpathlineto{\pgfqpoint{0.946113in}{0.580534in}}%
\pgfpathlineto{\pgfqpoint{0.946529in}{0.578018in}}%
\pgfpathlineto{\pgfqpoint{0.949099in}{0.579130in}}%
\pgfpathlineto{\pgfqpoint{0.953212in}{0.577928in}}%
\pgfpathlineto{\pgfqpoint{0.953739in}{0.576380in}}%
\pgfpathlineto{\pgfqpoint{0.951749in}{0.574813in}}%
\pgfpathlineto{\pgfqpoint{0.955961in}{0.573632in}}%
\pgfpathlineto{\pgfqpoint{0.958806in}{0.574118in}}%
\pgfpathlineto{\pgfqpoint{0.962206in}{0.570550in}}%
\pgfpathlineto{\pgfqpoint{0.963220in}{0.570510in}}%
\pgfpathlineto{\pgfqpoint{0.964026in}{0.566308in}}%
\pgfpathlineto{\pgfqpoint{0.966643in}{0.564927in}}%
\pgfpathlineto{\pgfqpoint{0.966745in}{0.562715in}}%
\pgfpathlineto{\pgfqpoint{0.968656in}{0.561617in}}%
\pgfpathlineto{\pgfqpoint{0.969071in}{0.557591in}}%
\pgfpathlineto{\pgfqpoint{0.974627in}{0.550898in}}%
\pgfpathlineto{\pgfqpoint{0.976869in}{0.550826in}}%
\pgfpathlineto{\pgfqpoint{0.980019in}{0.548774in}}%
\pgfpathlineto{\pgfqpoint{0.981988in}{0.546062in}}%
\pgfpathlineto{\pgfqpoint{0.984408in}{0.544629in}}%
\pgfpathlineto{\pgfqpoint{0.986456in}{0.538416in}}%
\pgfpathlineto{\pgfqpoint{0.988132in}{0.536134in}}%
\pgfpathlineto{\pgfqpoint{0.992596in}{0.535000in}}%
\pgfpathlineto{\pgfqpoint{0.998469in}{0.532993in}}%
\pgfpathlineto{\pgfqpoint{1.000842in}{0.532798in}}%
\pgfpathlineto{\pgfqpoint{1.003539in}{0.531122in}}%
\pgfpathlineto{\pgfqpoint{1.006445in}{0.526960in}}%
\pgfpathlineto{\pgfqpoint{1.008246in}{0.522245in}}%
\pgfpathlineto{\pgfqpoint{1.008217in}{0.520872in}}%
\pgfpathlineto{\pgfqpoint{1.010008in}{0.517358in}}%
\pgfpathlineto{\pgfqpoint{1.011929in}{0.514984in}}%
\pgfpathlineto{\pgfqpoint{1.012633in}{0.512658in}}%
\pgfpathlineto{\pgfqpoint{1.015157in}{0.509681in}}%
\pgfpathlineto{\pgfqpoint{1.014103in}{0.506544in}}%
\pgfpathlineto{\pgfqpoint{1.001047in}{0.482991in}}%
\pgfpathclose%
\pgfusepath{fill}%
\end{pgfscope}%
\begin{pgfscope}%
\pgfpathrectangle{\pgfqpoint{0.100000in}{0.100000in}}{\pgfqpoint{3.007045in}{1.925000in}}%
\pgfusepath{clip}%
\pgfsetbuttcap%
\pgfsetmiterjoin%
\definecolor{currentfill}{rgb}{0.376471,0.730181,0.424298}%
\pgfsetfillcolor{currentfill}%
\pgfsetlinewidth{0.000000pt}%
\definecolor{currentstroke}{rgb}{0.000000,0.000000,0.000000}%
\pgfsetstrokecolor{currentstroke}%
\pgfsetstrokeopacity{0.000000}%
\pgfsetdash{}{0pt}%
\pgfpathmoveto{\pgfqpoint{2.119654in}{1.040653in}}%
\pgfpathlineto{\pgfqpoint{2.116603in}{1.041678in}}%
\pgfpathlineto{\pgfqpoint{2.115921in}{1.053268in}}%
\pgfpathlineto{\pgfqpoint{2.100489in}{1.052130in}}%
\pgfpathlineto{\pgfqpoint{2.099179in}{1.072352in}}%
\pgfpathlineto{\pgfqpoint{2.116398in}{1.073683in}}%
\pgfpathlineto{\pgfqpoint{2.128732in}{1.074544in}}%
\pgfpathlineto{\pgfqpoint{2.131121in}{1.073626in}}%
\pgfpathlineto{\pgfqpoint{2.134908in}{1.067855in}}%
\pgfpathlineto{\pgfqpoint{2.128798in}{1.058421in}}%
\pgfpathlineto{\pgfqpoint{2.133456in}{1.046059in}}%
\pgfpathlineto{\pgfqpoint{2.123429in}{1.044085in}}%
\pgfpathlineto{\pgfqpoint{2.119654in}{1.040653in}}%
\pgfpathclose%
\pgfusepath{fill}%
\end{pgfscope}%
\begin{pgfscope}%
\pgfpathrectangle{\pgfqpoint{0.100000in}{0.100000in}}{\pgfqpoint{3.007045in}{1.925000in}}%
\pgfusepath{clip}%
\pgfsetbuttcap%
\pgfsetmiterjoin%
\definecolor{currentfill}{rgb}{0.676371,0.869927,0.651657}%
\pgfsetfillcolor{currentfill}%
\pgfsetlinewidth{0.000000pt}%
\definecolor{currentstroke}{rgb}{0.000000,0.000000,0.000000}%
\pgfsetstrokecolor{currentstroke}%
\pgfsetstrokeopacity{0.000000}%
\pgfsetdash{}{0pt}%
\pgfpathmoveto{\pgfqpoint{2.073753in}{0.803617in}}%
\pgfpathlineto{\pgfqpoint{2.059401in}{0.802761in}}%
\pgfpathlineto{\pgfqpoint{2.050685in}{0.804294in}}%
\pgfpathlineto{\pgfqpoint{2.050947in}{0.799400in}}%
\pgfpathlineto{\pgfqpoint{2.039467in}{0.797809in}}%
\pgfpathlineto{\pgfqpoint{2.039761in}{0.792018in}}%
\pgfpathlineto{\pgfqpoint{2.030424in}{0.792915in}}%
\pgfpathlineto{\pgfqpoint{2.036938in}{0.797675in}}%
\pgfpathlineto{\pgfqpoint{2.022165in}{0.797050in}}%
\pgfpathlineto{\pgfqpoint{2.021640in}{0.808564in}}%
\pgfpathlineto{\pgfqpoint{2.010153in}{0.808128in}}%
\pgfpathlineto{\pgfqpoint{2.009788in}{0.816800in}}%
\pgfpathlineto{\pgfqpoint{1.993574in}{0.816135in}}%
\pgfpathlineto{\pgfqpoint{2.003512in}{0.827202in}}%
\pgfpathlineto{\pgfqpoint{2.003733in}{0.832951in}}%
\pgfpathlineto{\pgfqpoint{2.013158in}{0.836265in}}%
\pgfpathlineto{\pgfqpoint{2.013407in}{0.843839in}}%
\pgfpathlineto{\pgfqpoint{2.022978in}{0.837629in}}%
\pgfpathlineto{\pgfqpoint{2.028291in}{0.837950in}}%
\pgfpathlineto{\pgfqpoint{2.027867in}{0.843557in}}%
\pgfpathlineto{\pgfqpoint{2.033456in}{0.846814in}}%
\pgfpathlineto{\pgfqpoint{2.072679in}{0.849028in}}%
\pgfpathlineto{\pgfqpoint{2.085442in}{0.845945in}}%
\pgfpathlineto{\pgfqpoint{2.085917in}{0.838289in}}%
\pgfpathlineto{\pgfqpoint{2.086876in}{0.823888in}}%
\pgfpathlineto{\pgfqpoint{2.072446in}{0.823053in}}%
\pgfpathlineto{\pgfqpoint{2.073753in}{0.803617in}}%
\pgfpathclose%
\pgfusepath{fill}%
\end{pgfscope}%
\begin{pgfscope}%
\pgfpathrectangle{\pgfqpoint{0.100000in}{0.100000in}}{\pgfqpoint{3.007045in}{1.925000in}}%
\pgfusepath{clip}%
\pgfsetbuttcap%
\pgfsetmiterjoin%
\definecolor{currentfill}{rgb}{0.457670,0.769919,0.465021}%
\pgfsetfillcolor{currentfill}%
\pgfsetlinewidth{0.000000pt}%
\definecolor{currentstroke}{rgb}{0.000000,0.000000,0.000000}%
\pgfsetstrokecolor{currentstroke}%
\pgfsetstrokeopacity{0.000000}%
\pgfsetdash{}{0pt}%
\pgfpathmoveto{\pgfqpoint{2.938382in}{1.591168in}}%
\pgfpathlineto{\pgfqpoint{2.930725in}{1.593936in}}%
\pgfpathlineto{\pgfqpoint{2.931682in}{1.596631in}}%
\pgfpathlineto{\pgfqpoint{2.926133in}{1.600139in}}%
\pgfpathlineto{\pgfqpoint{2.921331in}{1.597973in}}%
\pgfpathlineto{\pgfqpoint{2.920147in}{1.603400in}}%
\pgfpathlineto{\pgfqpoint{2.912256in}{1.602890in}}%
\pgfpathlineto{\pgfqpoint{2.904580in}{1.599504in}}%
\pgfpathlineto{\pgfqpoint{2.895122in}{1.630668in}}%
\pgfpathlineto{\pgfqpoint{2.885278in}{1.661827in}}%
\pgfpathlineto{\pgfqpoint{2.874149in}{1.694991in}}%
\pgfpathlineto{\pgfqpoint{2.879401in}{1.698668in}}%
\pgfpathlineto{\pgfqpoint{2.883573in}{1.693339in}}%
\pgfpathlineto{\pgfqpoint{2.886335in}{1.697801in}}%
\pgfpathlineto{\pgfqpoint{2.885576in}{1.706733in}}%
\pgfpathlineto{\pgfqpoint{2.889390in}{1.704990in}}%
\pgfpathlineto{\pgfqpoint{2.892641in}{1.707989in}}%
\pgfpathlineto{\pgfqpoint{2.886793in}{1.712832in}}%
\pgfpathlineto{\pgfqpoint{2.888624in}{1.719661in}}%
\pgfpathlineto{\pgfqpoint{2.897841in}{1.730562in}}%
\pgfpathlineto{\pgfqpoint{2.896524in}{1.736677in}}%
\pgfpathlineto{\pgfqpoint{2.900020in}{1.741049in}}%
\pgfpathlineto{\pgfqpoint{2.900699in}{1.746177in}}%
\pgfpathlineto{\pgfqpoint{2.895909in}{1.750061in}}%
\pgfpathlineto{\pgfqpoint{2.897245in}{1.759098in}}%
\pgfpathlineto{\pgfqpoint{2.893988in}{1.761345in}}%
\pgfpathlineto{\pgfqpoint{2.895208in}{1.771201in}}%
\pgfpathlineto{\pgfqpoint{2.900371in}{1.778621in}}%
\pgfpathlineto{\pgfqpoint{2.898990in}{1.789167in}}%
\pgfpathlineto{\pgfqpoint{2.912260in}{1.792925in}}%
\pgfpathlineto{\pgfqpoint{2.915056in}{1.781252in}}%
\pgfpathlineto{\pgfqpoint{2.922774in}{1.754622in}}%
\pgfpathlineto{\pgfqpoint{2.926062in}{1.747145in}}%
\pgfpathlineto{\pgfqpoint{2.925595in}{1.735765in}}%
\pgfpathlineto{\pgfqpoint{2.930414in}{1.732814in}}%
\pgfpathlineto{\pgfqpoint{2.928393in}{1.726358in}}%
\pgfpathlineto{\pgfqpoint{2.945232in}{1.694510in}}%
\pgfpathlineto{\pgfqpoint{2.956012in}{1.701938in}}%
\pgfpathlineto{\pgfqpoint{2.966515in}{1.680747in}}%
\pgfpathlineto{\pgfqpoint{2.960139in}{1.680504in}}%
\pgfpathlineto{\pgfqpoint{2.962312in}{1.673903in}}%
\pgfpathlineto{\pgfqpoint{2.963183in}{1.661995in}}%
\pgfpathlineto{\pgfqpoint{2.962617in}{1.653501in}}%
\pgfpathlineto{\pgfqpoint{2.965671in}{1.652396in}}%
\pgfpathlineto{\pgfqpoint{2.969629in}{1.646428in}}%
\pgfpathlineto{\pgfqpoint{2.972759in}{1.647244in}}%
\pgfpathlineto{\pgfqpoint{2.978270in}{1.638556in}}%
\pgfpathlineto{\pgfqpoint{2.975595in}{1.630839in}}%
\pgfpathlineto{\pgfqpoint{2.972528in}{1.630382in}}%
\pgfpathlineto{\pgfqpoint{2.972069in}{1.622078in}}%
\pgfpathlineto{\pgfqpoint{2.963086in}{1.616945in}}%
\pgfpathlineto{\pgfqpoint{2.962287in}{1.612214in}}%
\pgfpathlineto{\pgfqpoint{2.958299in}{1.609229in}}%
\pgfpathlineto{\pgfqpoint{2.947252in}{1.614168in}}%
\pgfpathlineto{\pgfqpoint{2.941831in}{1.608526in}}%
\pgfpathlineto{\pgfqpoint{2.940412in}{1.601930in}}%
\pgfpathlineto{\pgfqpoint{2.945057in}{1.595018in}}%
\pgfpathlineto{\pgfqpoint{2.938382in}{1.591168in}}%
\pgfpathclose%
\pgfusepath{fill}%
\end{pgfscope}%
\begin{pgfscope}%
\pgfpathrectangle{\pgfqpoint{0.100000in}{0.100000in}}{\pgfqpoint{3.007045in}{1.925000in}}%
\pgfusepath{clip}%
\pgfsetbuttcap%
\pgfsetmiterjoin%
\definecolor{currentfill}{rgb}{0.690396,0.875832,0.665313}%
\pgfsetfillcolor{currentfill}%
\pgfsetlinewidth{0.000000pt}%
\definecolor{currentstroke}{rgb}{0.000000,0.000000,0.000000}%
\pgfsetstrokecolor{currentstroke}%
\pgfsetstrokeopacity{0.000000}%
\pgfsetdash{}{0pt}%
\pgfpathmoveto{\pgfqpoint{2.116053in}{0.743899in}}%
\pgfpathlineto{\pgfqpoint{2.087936in}{0.742436in}}%
\pgfpathlineto{\pgfqpoint{2.086624in}{0.765838in}}%
\pgfpathlineto{\pgfqpoint{2.079034in}{0.765390in}}%
\pgfpathlineto{\pgfqpoint{2.078426in}{0.775122in}}%
\pgfpathlineto{\pgfqpoint{2.075429in}{0.776700in}}%
\pgfpathlineto{\pgfqpoint{2.074344in}{0.794001in}}%
\pgfpathlineto{\pgfqpoint{2.091654in}{0.795102in}}%
\pgfpathlineto{\pgfqpoint{2.105804in}{0.798159in}}%
\pgfpathlineto{\pgfqpoint{2.105570in}{0.801993in}}%
\pgfpathlineto{\pgfqpoint{2.117152in}{0.802691in}}%
\pgfpathlineto{\pgfqpoint{2.117772in}{0.794063in}}%
\pgfpathlineto{\pgfqpoint{2.129211in}{0.792870in}}%
\pgfpathlineto{\pgfqpoint{2.131570in}{0.789934in}}%
\pgfpathlineto{\pgfqpoint{2.128772in}{0.785353in}}%
\pgfpathlineto{\pgfqpoint{2.121230in}{0.782758in}}%
\pgfpathlineto{\pgfqpoint{2.122360in}{0.768291in}}%
\pgfpathlineto{\pgfqpoint{2.114414in}{0.767709in}}%
\pgfpathlineto{\pgfqpoint{2.116053in}{0.743899in}}%
\pgfpathclose%
\pgfusepath{fill}%
\end{pgfscope}%
\begin{pgfscope}%
\pgfpathrectangle{\pgfqpoint{0.100000in}{0.100000in}}{\pgfqpoint{3.007045in}{1.925000in}}%
\pgfusepath{clip}%
\pgfsetbuttcap%
\pgfsetmiterjoin%
\definecolor{currentfill}{rgb}{0.713772,0.885675,0.688074}%
\pgfsetfillcolor{currentfill}%
\pgfsetlinewidth{0.000000pt}%
\definecolor{currentstroke}{rgb}{0.000000,0.000000,0.000000}%
\pgfsetstrokecolor{currentstroke}%
\pgfsetstrokeopacity{0.000000}%
\pgfsetdash{}{0pt}%
\pgfpathmoveto{\pgfqpoint{0.536189in}{1.704696in}}%
\pgfpathlineto{\pgfqpoint{0.506524in}{1.713649in}}%
\pgfpathlineto{\pgfqpoint{0.501234in}{1.718853in}}%
\pgfpathlineto{\pgfqpoint{0.500179in}{1.724897in}}%
\pgfpathlineto{\pgfqpoint{0.496057in}{1.728942in}}%
\pgfpathlineto{\pgfqpoint{0.493784in}{1.734570in}}%
\pgfpathlineto{\pgfqpoint{0.496421in}{1.737333in}}%
\pgfpathlineto{\pgfqpoint{0.496816in}{1.743625in}}%
\pgfpathlineto{\pgfqpoint{0.490020in}{1.750804in}}%
\pgfpathlineto{\pgfqpoint{0.488462in}{1.755366in}}%
\pgfpathlineto{\pgfqpoint{0.484457in}{1.758626in}}%
\pgfpathlineto{\pgfqpoint{0.469874in}{1.763357in}}%
\pgfpathlineto{\pgfqpoint{0.475014in}{1.769793in}}%
\pgfpathlineto{\pgfqpoint{0.475846in}{1.775523in}}%
\pgfpathlineto{\pgfqpoint{0.480074in}{1.779212in}}%
\pgfpathlineto{\pgfqpoint{0.481516in}{1.783693in}}%
\pgfpathlineto{\pgfqpoint{0.488870in}{1.806888in}}%
\pgfpathlineto{\pgfqpoint{0.498370in}{1.806808in}}%
\pgfpathlineto{\pgfqpoint{0.507744in}{1.796488in}}%
\pgfpathlineto{\pgfqpoint{0.510888in}{1.782623in}}%
\pgfpathlineto{\pgfqpoint{0.524572in}{1.784544in}}%
\pgfpathlineto{\pgfqpoint{0.529311in}{1.780896in}}%
\pgfpathlineto{\pgfqpoint{0.536092in}{1.785540in}}%
\pgfpathlineto{\pgfqpoint{0.542761in}{1.806508in}}%
\pgfpathlineto{\pgfqpoint{0.574289in}{1.797104in}}%
\pgfpathlineto{\pgfqpoint{0.567701in}{1.775221in}}%
\pgfpathlineto{\pgfqpoint{0.563789in}{1.776383in}}%
\pgfpathlineto{\pgfqpoint{0.558944in}{1.759738in}}%
\pgfpathlineto{\pgfqpoint{0.561232in}{1.755167in}}%
\pgfpathlineto{\pgfqpoint{0.551653in}{1.755806in}}%
\pgfpathlineto{\pgfqpoint{0.545786in}{1.757536in}}%
\pgfpathlineto{\pgfqpoint{0.542490in}{1.755480in}}%
\pgfpathlineto{\pgfqpoint{0.540688in}{1.747112in}}%
\pgfpathlineto{\pgfqpoint{0.547023in}{1.735645in}}%
\pgfpathlineto{\pgfqpoint{0.545552in}{1.726930in}}%
\pgfpathlineto{\pgfqpoint{0.541409in}{1.723741in}}%
\pgfpathlineto{\pgfqpoint{0.541696in}{1.715411in}}%
\pgfpathlineto{\pgfqpoint{0.535329in}{1.713270in}}%
\pgfpathlineto{\pgfqpoint{0.536189in}{1.704696in}}%
\pgfpathclose%
\pgfusepath{fill}%
\end{pgfscope}%
\begin{pgfscope}%
\pgfpathrectangle{\pgfqpoint{0.100000in}{0.100000in}}{\pgfqpoint{3.007045in}{1.925000in}}%
\pgfusepath{clip}%
\pgfsetbuttcap%
\pgfsetmiterjoin%
\definecolor{currentfill}{rgb}{0.518570,0.798339,0.515094}%
\pgfsetfillcolor{currentfill}%
\pgfsetlinewidth{0.000000pt}%
\definecolor{currentstroke}{rgb}{0.000000,0.000000,0.000000}%
\pgfsetstrokecolor{currentstroke}%
\pgfsetstrokeopacity{0.000000}%
\pgfsetdash{}{0pt}%
\pgfpathmoveto{\pgfqpoint{2.539586in}{0.435076in}}%
\pgfpathlineto{\pgfqpoint{2.510105in}{0.430828in}}%
\pgfpathlineto{\pgfqpoint{2.515004in}{0.442453in}}%
\pgfpathlineto{\pgfqpoint{2.510239in}{0.443614in}}%
\pgfpathlineto{\pgfqpoint{2.504846in}{0.439692in}}%
\pgfpathlineto{\pgfqpoint{2.503926in}{0.433696in}}%
\pgfpathlineto{\pgfqpoint{2.500319in}{0.432366in}}%
\pgfpathlineto{\pgfqpoint{2.490587in}{0.443682in}}%
\pgfpathlineto{\pgfqpoint{2.492902in}{0.456197in}}%
\pgfpathlineto{\pgfqpoint{2.490893in}{0.464809in}}%
\pgfpathlineto{\pgfqpoint{2.493043in}{0.468035in}}%
\pgfpathlineto{\pgfqpoint{2.496096in}{0.492116in}}%
\pgfpathlineto{\pgfqpoint{2.495655in}{0.499863in}}%
\pgfpathlineto{\pgfqpoint{2.508454in}{0.501648in}}%
\pgfpathlineto{\pgfqpoint{2.517736in}{0.501110in}}%
\pgfpathlineto{\pgfqpoint{2.521853in}{0.495228in}}%
\pgfpathlineto{\pgfqpoint{2.531241in}{0.493155in}}%
\pgfpathlineto{\pgfqpoint{2.533154in}{0.479245in}}%
\pgfpathlineto{\pgfqpoint{2.530722in}{0.475320in}}%
\pgfpathlineto{\pgfqpoint{2.534472in}{0.469934in}}%
\pgfpathlineto{\pgfqpoint{2.539586in}{0.435076in}}%
\pgfpathclose%
\pgfusepath{fill}%
\end{pgfscope}%
\begin{pgfscope}%
\pgfpathrectangle{\pgfqpoint{0.100000in}{0.100000in}}{\pgfqpoint{3.007045in}{1.925000in}}%
\pgfusepath{clip}%
\pgfsetbuttcap%
\pgfsetmiterjoin%
\definecolor{currentfill}{rgb}{0.540715,0.808674,0.533303}%
\pgfsetfillcolor{currentfill}%
\pgfsetlinewidth{0.000000pt}%
\definecolor{currentstroke}{rgb}{0.000000,0.000000,0.000000}%
\pgfsetstrokecolor{currentstroke}%
\pgfsetstrokeopacity{0.000000}%
\pgfsetdash{}{0pt}%
\pgfpathmoveto{\pgfqpoint{2.867151in}{1.544262in}}%
\pgfpathlineto{\pgfqpoint{2.867535in}{1.548390in}}%
\pgfpathlineto{\pgfqpoint{2.861539in}{1.556030in}}%
\pgfpathlineto{\pgfqpoint{2.863308in}{1.564293in}}%
\pgfpathlineto{\pgfqpoint{2.867917in}{1.577405in}}%
\pgfpathlineto{\pgfqpoint{2.870513in}{1.573950in}}%
\pgfpathlineto{\pgfqpoint{2.876946in}{1.580426in}}%
\pgfpathlineto{\pgfqpoint{2.875977in}{1.584608in}}%
\pgfpathlineto{\pgfqpoint{2.881617in}{1.586847in}}%
\pgfpathlineto{\pgfqpoint{2.876355in}{1.598852in}}%
\pgfpathlineto{\pgfqpoint{2.885308in}{1.602759in}}%
\pgfpathlineto{\pgfqpoint{2.882477in}{1.612503in}}%
\pgfpathlineto{\pgfqpoint{2.877796in}{1.621050in}}%
\pgfpathlineto{\pgfqpoint{2.884101in}{1.617605in}}%
\pgfpathlineto{\pgfqpoint{2.884731in}{1.625099in}}%
\pgfpathlineto{\pgfqpoint{2.894447in}{1.627269in}}%
\pgfpathlineto{\pgfqpoint{2.895122in}{1.630668in}}%
\pgfpathlineto{\pgfqpoint{2.904580in}{1.599504in}}%
\pgfpathlineto{\pgfqpoint{2.912256in}{1.602890in}}%
\pgfpathlineto{\pgfqpoint{2.920147in}{1.603400in}}%
\pgfpathlineto{\pgfqpoint{2.921331in}{1.597973in}}%
\pgfpathlineto{\pgfqpoint{2.926133in}{1.600139in}}%
\pgfpathlineto{\pgfqpoint{2.931682in}{1.596631in}}%
\pgfpathlineto{\pgfqpoint{2.930725in}{1.593936in}}%
\pgfpathlineto{\pgfqpoint{2.938382in}{1.591168in}}%
\pgfpathlineto{\pgfqpoint{2.939431in}{1.584564in}}%
\pgfpathlineto{\pgfqpoint{2.937676in}{1.578867in}}%
\pgfpathlineto{\pgfqpoint{2.933280in}{1.576054in}}%
\pgfpathlineto{\pgfqpoint{2.932692in}{1.561193in}}%
\pgfpathlineto{\pgfqpoint{2.928157in}{1.545679in}}%
\pgfpathlineto{\pgfqpoint{2.928476in}{1.542999in}}%
\pgfpathlineto{\pgfqpoint{2.922960in}{1.542355in}}%
\pgfpathlineto{\pgfqpoint{2.914625in}{1.535783in}}%
\pgfpathlineto{\pgfqpoint{2.909180in}{1.525824in}}%
\pgfpathlineto{\pgfqpoint{2.879160in}{1.518963in}}%
\pgfpathlineto{\pgfqpoint{2.875745in}{1.523338in}}%
\pgfpathlineto{\pgfqpoint{2.871705in}{1.532829in}}%
\pgfpathlineto{\pgfqpoint{2.869055in}{1.532209in}}%
\pgfpathlineto{\pgfqpoint{2.867151in}{1.544262in}}%
\pgfpathclose%
\pgfusepath{fill}%
\end{pgfscope}%
\begin{pgfscope}%
\pgfpathrectangle{\pgfqpoint{0.100000in}{0.100000in}}{\pgfqpoint{3.007045in}{1.925000in}}%
\pgfusepath{clip}%
\pgfsetbuttcap%
\pgfsetmiterjoin%
\definecolor{currentfill}{rgb}{0.518570,0.798339,0.515094}%
\pgfsetfillcolor{currentfill}%
\pgfsetlinewidth{0.000000pt}%
\definecolor{currentstroke}{rgb}{0.000000,0.000000,0.000000}%
\pgfsetstrokecolor{currentstroke}%
\pgfsetstrokeopacity{0.000000}%
\pgfsetdash{}{0pt}%
\pgfpathmoveto{\pgfqpoint{2.082341in}{0.997150in}}%
\pgfpathlineto{\pgfqpoint{2.081317in}{1.001182in}}%
\pgfpathlineto{\pgfqpoint{2.082880in}{1.008816in}}%
\pgfpathlineto{\pgfqpoint{2.081022in}{1.010070in}}%
\pgfpathlineto{\pgfqpoint{2.078167in}{1.013748in}}%
\pgfpathlineto{\pgfqpoint{2.085427in}{1.025478in}}%
\pgfpathlineto{\pgfqpoint{2.090698in}{1.026843in}}%
\pgfpathlineto{\pgfqpoint{2.091233in}{1.034083in}}%
\pgfpathlineto{\pgfqpoint{2.101574in}{1.034888in}}%
\pgfpathlineto{\pgfqpoint{2.100489in}{1.052130in}}%
\pgfpathlineto{\pgfqpoint{2.115921in}{1.053268in}}%
\pgfpathlineto{\pgfqpoint{2.116603in}{1.041678in}}%
\pgfpathlineto{\pgfqpoint{2.119654in}{1.040653in}}%
\pgfpathlineto{\pgfqpoint{2.127608in}{1.032537in}}%
\pgfpathlineto{\pgfqpoint{2.129826in}{1.024933in}}%
\pgfpathlineto{\pgfqpoint{2.127814in}{1.020206in}}%
\pgfpathlineto{\pgfqpoint{2.132817in}{1.006098in}}%
\pgfpathlineto{\pgfqpoint{2.135867in}{0.998378in}}%
\pgfpathlineto{\pgfqpoint{2.116101in}{0.997143in}}%
\pgfpathlineto{\pgfqpoint{2.115261in}{1.009765in}}%
\pgfpathlineto{\pgfqpoint{2.097982in}{1.008842in}}%
\pgfpathlineto{\pgfqpoint{2.098746in}{0.997437in}}%
\pgfpathlineto{\pgfqpoint{2.082341in}{0.997150in}}%
\pgfpathclose%
\pgfusepath{fill}%
\end{pgfscope}%
\begin{pgfscope}%
\pgfpathrectangle{\pgfqpoint{0.100000in}{0.100000in}}{\pgfqpoint{3.007045in}{1.925000in}}%
\pgfusepath{clip}%
\pgfsetbuttcap%
\pgfsetmiterjoin%
\definecolor{currentfill}{rgb}{0.445490,0.764014,0.458131}%
\pgfsetfillcolor{currentfill}%
\pgfsetlinewidth{0.000000pt}%
\definecolor{currentstroke}{rgb}{0.000000,0.000000,0.000000}%
\pgfsetstrokecolor{currentstroke}%
\pgfsetstrokeopacity{0.000000}%
\pgfsetdash{}{0pt}%
\pgfpathmoveto{\pgfqpoint{2.654171in}{1.354581in}}%
\pgfpathlineto{\pgfqpoint{2.645166in}{1.349038in}}%
\pgfpathlineto{\pgfqpoint{2.637385in}{1.345841in}}%
\pgfpathlineto{\pgfqpoint{2.627962in}{1.353837in}}%
\pgfpathlineto{\pgfqpoint{2.623988in}{1.352894in}}%
\pgfpathlineto{\pgfqpoint{2.621003in}{1.357074in}}%
\pgfpathlineto{\pgfqpoint{2.609177in}{1.353242in}}%
\pgfpathlineto{\pgfqpoint{2.605715in}{1.356835in}}%
\pgfpathlineto{\pgfqpoint{2.608964in}{1.367546in}}%
\pgfpathlineto{\pgfqpoint{2.607627in}{1.374496in}}%
\pgfpathlineto{\pgfqpoint{2.626550in}{1.378374in}}%
\pgfpathlineto{\pgfqpoint{2.625680in}{1.382482in}}%
\pgfpathlineto{\pgfqpoint{2.656505in}{1.389182in}}%
\pgfpathlineto{\pgfqpoint{2.660061in}{1.392938in}}%
\pgfpathlineto{\pgfqpoint{2.652249in}{1.418613in}}%
\pgfpathlineto{\pgfqpoint{2.669988in}{1.422124in}}%
\pgfpathlineto{\pgfqpoint{2.689741in}{1.426141in}}%
\pgfpathlineto{\pgfqpoint{2.695962in}{1.404095in}}%
\pgfpathlineto{\pgfqpoint{2.691983in}{1.402891in}}%
\pgfpathlineto{\pgfqpoint{2.692396in}{1.393752in}}%
\pgfpathlineto{\pgfqpoint{2.691171in}{1.379954in}}%
\pgfpathlineto{\pgfqpoint{2.686454in}{1.378941in}}%
\pgfpathlineto{\pgfqpoint{2.679506in}{1.367379in}}%
\pgfpathlineto{\pgfqpoint{2.677192in}{1.366767in}}%
\pgfpathlineto{\pgfqpoint{2.669738in}{1.366595in}}%
\pgfpathlineto{\pgfqpoint{2.661755in}{1.363168in}}%
\pgfpathlineto{\pgfqpoint{2.661650in}{1.359044in}}%
\pgfpathlineto{\pgfqpoint{2.654171in}{1.354581in}}%
\pgfpathclose%
\pgfusepath{fill}%
\end{pgfscope}%
\begin{pgfscope}%
\pgfpathrectangle{\pgfqpoint{0.100000in}{0.100000in}}{\pgfqpoint{3.007045in}{1.925000in}}%
\pgfusepath{clip}%
\pgfsetbuttcap%
\pgfsetmiterjoin%
\definecolor{currentfill}{rgb}{0.439216,0.760938,0.455056}%
\pgfsetfillcolor{currentfill}%
\pgfsetlinewidth{0.000000pt}%
\definecolor{currentstroke}{rgb}{0.000000,0.000000,0.000000}%
\pgfsetstrokecolor{currentstroke}%
\pgfsetstrokeopacity{0.000000}%
\pgfsetdash{}{0pt}%
\pgfpathmoveto{\pgfqpoint{2.245494in}{0.625658in}}%
\pgfpathlineto{\pgfqpoint{2.193761in}{0.621372in}}%
\pgfpathlineto{\pgfqpoint{2.191464in}{0.637662in}}%
\pgfpathlineto{\pgfqpoint{2.185584in}{0.641031in}}%
\pgfpathlineto{\pgfqpoint{2.184550in}{0.645723in}}%
\pgfpathlineto{\pgfqpoint{2.187693in}{0.649271in}}%
\pgfpathlineto{\pgfqpoint{2.192871in}{0.650628in}}%
\pgfpathlineto{\pgfqpoint{2.191786in}{0.667860in}}%
\pgfpathlineto{\pgfqpoint{2.194637in}{0.668142in}}%
\pgfpathlineto{\pgfqpoint{2.194774in}{0.676909in}}%
\pgfpathlineto{\pgfqpoint{2.228144in}{0.679871in}}%
\pgfpathlineto{\pgfqpoint{2.229320in}{0.666780in}}%
\pgfpathlineto{\pgfqpoint{2.233755in}{0.660011in}}%
\pgfpathlineto{\pgfqpoint{2.241536in}{0.660624in}}%
\pgfpathlineto{\pgfqpoint{2.244783in}{0.625988in}}%
\pgfpathlineto{\pgfqpoint{2.245494in}{0.625658in}}%
\pgfpathclose%
\pgfusepath{fill}%
\end{pgfscope}%
\begin{pgfscope}%
\pgfpathrectangle{\pgfqpoint{0.100000in}{0.100000in}}{\pgfqpoint{3.007045in}{1.925000in}}%
\pgfusepath{clip}%
\pgfsetbuttcap%
\pgfsetmiterjoin%
\definecolor{currentfill}{rgb}{0.363922,0.724029,0.418147}%
\pgfsetfillcolor{currentfill}%
\pgfsetlinewidth{0.000000pt}%
\definecolor{currentstroke}{rgb}{0.000000,0.000000,0.000000}%
\pgfsetstrokecolor{currentstroke}%
\pgfsetstrokeopacity{0.000000}%
\pgfsetdash{}{0pt}%
\pgfpathmoveto{\pgfqpoint{1.054598in}{0.801868in}}%
\pgfpathlineto{\pgfqpoint{0.997334in}{0.810949in}}%
\pgfpathlineto{\pgfqpoint{1.003150in}{0.848168in}}%
\pgfpathlineto{\pgfqpoint{1.011345in}{0.900484in}}%
\pgfpathlineto{\pgfqpoint{1.047397in}{0.894955in}}%
\pgfpathlineto{\pgfqpoint{1.082425in}{0.889855in}}%
\pgfpathlineto{\pgfqpoint{1.127835in}{0.883763in}}%
\pgfpathlineto{\pgfqpoint{1.133587in}{0.879798in}}%
\pgfpathlineto{\pgfqpoint{1.151795in}{0.871370in}}%
\pgfpathlineto{\pgfqpoint{1.150364in}{0.859555in}}%
\pgfpathlineto{\pgfqpoint{1.177035in}{0.856282in}}%
\pgfpathlineto{\pgfqpoint{1.173550in}{0.827766in}}%
\pgfpathlineto{\pgfqpoint{1.166565in}{0.828636in}}%
\pgfpathlineto{\pgfqpoint{1.165157in}{0.817109in}}%
\pgfpathlineto{\pgfqpoint{1.147690in}{0.819949in}}%
\pgfpathlineto{\pgfqpoint{1.146247in}{0.808107in}}%
\pgfpathlineto{\pgfqpoint{1.115027in}{0.811851in}}%
\pgfpathlineto{\pgfqpoint{1.060540in}{0.819258in}}%
\pgfpathlineto{\pgfqpoint{1.057206in}{0.819980in}}%
\pgfpathlineto{\pgfqpoint{1.054598in}{0.801868in}}%
\pgfpathclose%
\pgfusepath{fill}%
\end{pgfscope}%
\begin{pgfscope}%
\pgfpathrectangle{\pgfqpoint{0.100000in}{0.100000in}}{\pgfqpoint{3.007045in}{1.925000in}}%
\pgfusepath{clip}%
\pgfsetbuttcap%
\pgfsetmiterjoin%
\definecolor{currentfill}{rgb}{0.468743,0.775087,0.474125}%
\pgfsetfillcolor{currentfill}%
\pgfsetlinewidth{0.000000pt}%
\definecolor{currentstroke}{rgb}{0.000000,0.000000,0.000000}%
\pgfsetstrokecolor{currentstroke}%
\pgfsetstrokeopacity{0.000000}%
\pgfsetdash{}{0pt}%
\pgfpathmoveto{\pgfqpoint{2.246778in}{1.453408in}}%
\pgfpathlineto{\pgfqpoint{2.224008in}{1.451122in}}%
\pgfpathlineto{\pgfqpoint{2.221626in}{1.473931in}}%
\pgfpathlineto{\pgfqpoint{2.199018in}{1.471645in}}%
\pgfpathlineto{\pgfqpoint{2.196560in}{1.494814in}}%
\pgfpathlineto{\pgfqpoint{2.207015in}{1.495613in}}%
\pgfpathlineto{\pgfqpoint{2.204865in}{1.518579in}}%
\pgfpathlineto{\pgfqpoint{2.250233in}{1.523286in}}%
\pgfpathlineto{\pgfqpoint{2.252766in}{1.500298in}}%
\pgfpathlineto{\pgfqpoint{2.241654in}{1.499264in}}%
\pgfpathlineto{\pgfqpoint{2.246778in}{1.453408in}}%
\pgfpathclose%
\pgfusepath{fill}%
\end{pgfscope}%
\begin{pgfscope}%
\pgfpathrectangle{\pgfqpoint{0.100000in}{0.100000in}}{\pgfqpoint{3.007045in}{1.925000in}}%
\pgfusepath{clip}%
\pgfsetbuttcap%
\pgfsetmiterjoin%
\definecolor{currentfill}{rgb}{0.643645,0.856148,0.619792}%
\pgfsetfillcolor{currentfill}%
\pgfsetlinewidth{0.000000pt}%
\definecolor{currentstroke}{rgb}{0.000000,0.000000,0.000000}%
\pgfsetstrokecolor{currentstroke}%
\pgfsetstrokeopacity{0.000000}%
\pgfsetdash{}{0pt}%
\pgfpathmoveto{\pgfqpoint{2.163243in}{0.916924in}}%
\pgfpathlineto{\pgfqpoint{2.166375in}{0.924298in}}%
\pgfpathlineto{\pgfqpoint{2.160890in}{0.928692in}}%
\pgfpathlineto{\pgfqpoint{2.161409in}{0.940303in}}%
\pgfpathlineto{\pgfqpoint{2.161575in}{0.944740in}}%
\pgfpathlineto{\pgfqpoint{2.170276in}{0.951262in}}%
\pgfpathlineto{\pgfqpoint{2.188043in}{0.950490in}}%
\pgfpathlineto{\pgfqpoint{2.188078in}{0.943307in}}%
\pgfpathlineto{\pgfqpoint{2.210327in}{0.936644in}}%
\pgfpathlineto{\pgfqpoint{2.219736in}{0.935214in}}%
\pgfpathlineto{\pgfqpoint{2.220248in}{0.918653in}}%
\pgfpathlineto{\pgfqpoint{2.223823in}{0.914321in}}%
\pgfpathlineto{\pgfqpoint{2.223886in}{0.911812in}}%
\pgfpathlineto{\pgfqpoint{2.216592in}{0.906228in}}%
\pgfpathlineto{\pgfqpoint{2.212159in}{0.906542in}}%
\pgfpathlineto{\pgfqpoint{2.213336in}{0.888623in}}%
\pgfpathlineto{\pgfqpoint{2.193121in}{0.887325in}}%
\pgfpathlineto{\pgfqpoint{2.192088in}{0.890541in}}%
\pgfpathlineto{\pgfqpoint{2.190733in}{0.912451in}}%
\pgfpathlineto{\pgfqpoint{2.189759in}{0.915550in}}%
\pgfpathlineto{\pgfqpoint{2.184378in}{0.916884in}}%
\pgfpathlineto{\pgfqpoint{2.171317in}{0.911958in}}%
\pgfpathlineto{\pgfqpoint{2.167138in}{0.916537in}}%
\pgfpathlineto{\pgfqpoint{2.163243in}{0.916924in}}%
\pgfpathclose%
\pgfusepath{fill}%
\end{pgfscope}%
\begin{pgfscope}%
\pgfpathrectangle{\pgfqpoint{0.100000in}{0.100000in}}{\pgfqpoint{3.007045in}{1.925000in}}%
\pgfusepath{clip}%
\pgfsetbuttcap%
\pgfsetmiterjoin%
\definecolor{currentfill}{rgb}{0.420392,0.751711,0.445829}%
\pgfsetfillcolor{currentfill}%
\pgfsetlinewidth{0.000000pt}%
\definecolor{currentstroke}{rgb}{0.000000,0.000000,0.000000}%
\pgfsetstrokecolor{currentstroke}%
\pgfsetstrokeopacity{0.000000}%
\pgfsetdash{}{0pt}%
\pgfpathmoveto{\pgfqpoint{2.525737in}{1.174087in}}%
\pgfpathlineto{\pgfqpoint{2.532425in}{1.163870in}}%
\pgfpathlineto{\pgfqpoint{2.530752in}{1.154455in}}%
\pgfpathlineto{\pgfqpoint{2.526995in}{1.150325in}}%
\pgfpathlineto{\pgfqpoint{2.526634in}{1.142930in}}%
\pgfpathlineto{\pgfqpoint{2.522292in}{1.144995in}}%
\pgfpathlineto{\pgfqpoint{2.510167in}{1.130498in}}%
\pgfpathlineto{\pgfqpoint{2.501225in}{1.130708in}}%
\pgfpathlineto{\pgfqpoint{2.494710in}{1.137649in}}%
\pgfpathlineto{\pgfqpoint{2.490413in}{1.135206in}}%
\pgfpathlineto{\pgfqpoint{2.481701in}{1.138352in}}%
\pgfpathlineto{\pgfqpoint{2.496539in}{1.150737in}}%
\pgfpathlineto{\pgfqpoint{2.497293in}{1.156805in}}%
\pgfpathlineto{\pgfqpoint{2.494103in}{1.161327in}}%
\pgfpathlineto{\pgfqpoint{2.487851in}{1.166255in}}%
\pgfpathlineto{\pgfqpoint{2.489719in}{1.169881in}}%
\pgfpathlineto{\pgfqpoint{2.485316in}{1.174226in}}%
\pgfpathlineto{\pgfqpoint{2.487519in}{1.178106in}}%
\pgfpathlineto{\pgfqpoint{2.484306in}{1.188484in}}%
\pgfpathlineto{\pgfqpoint{2.489753in}{1.191584in}}%
\pgfpathlineto{\pgfqpoint{2.494442in}{1.196721in}}%
\pgfpathlineto{\pgfqpoint{2.498857in}{1.196503in}}%
\pgfpathlineto{\pgfqpoint{2.500337in}{1.190919in}}%
\pgfpathlineto{\pgfqpoint{2.507565in}{1.184564in}}%
\pgfpathlineto{\pgfqpoint{2.514455in}{1.181489in}}%
\pgfpathlineto{\pgfqpoint{2.517598in}{1.175060in}}%
\pgfpathlineto{\pgfqpoint{2.525737in}{1.174087in}}%
\pgfpathclose%
\pgfusepath{fill}%
\end{pgfscope}%
\begin{pgfscope}%
\pgfpathrectangle{\pgfqpoint{0.100000in}{0.100000in}}{\pgfqpoint{3.007045in}{1.925000in}}%
\pgfusepath{clip}%
\pgfsetbuttcap%
\pgfsetmiterjoin%
\definecolor{currentfill}{rgb}{0.501961,0.790588,0.501438}%
\pgfsetfillcolor{currentfill}%
\pgfsetlinewidth{0.000000pt}%
\definecolor{currentstroke}{rgb}{0.000000,0.000000,0.000000}%
\pgfsetstrokecolor{currentstroke}%
\pgfsetstrokeopacity{0.000000}%
\pgfsetdash{}{0pt}%
\pgfpathmoveto{\pgfqpoint{0.836260in}{1.093783in}}%
\pgfpathlineto{\pgfqpoint{0.776912in}{1.106019in}}%
\pgfpathlineto{\pgfqpoint{0.788118in}{1.158139in}}%
\pgfpathlineto{\pgfqpoint{0.798495in}{1.207353in}}%
\pgfpathlineto{\pgfqpoint{0.850288in}{1.196720in}}%
\pgfpathlineto{\pgfqpoint{0.875936in}{1.191519in}}%
\pgfpathlineto{\pgfqpoint{0.875518in}{1.183498in}}%
\pgfpathlineto{\pgfqpoint{0.878976in}{1.180671in}}%
\pgfpathlineto{\pgfqpoint{0.882037in}{1.174273in}}%
\pgfpathlineto{\pgfqpoint{0.877415in}{1.166653in}}%
\pgfpathlineto{\pgfqpoint{0.872431in}{1.163681in}}%
\pgfpathlineto{\pgfqpoint{0.869732in}{1.146966in}}%
\pgfpathlineto{\pgfqpoint{0.864002in}{1.147526in}}%
\pgfpathlineto{\pgfqpoint{0.863012in}{1.142510in}}%
\pgfpathlineto{\pgfqpoint{0.857375in}{1.143704in}}%
\pgfpathlineto{\pgfqpoint{0.854337in}{1.126666in}}%
\pgfpathlineto{\pgfqpoint{0.843170in}{1.128894in}}%
\pgfpathlineto{\pgfqpoint{0.836260in}{1.093783in}}%
\pgfpathclose%
\pgfusepath{fill}%
\end{pgfscope}%
\begin{pgfscope}%
\pgfpathrectangle{\pgfqpoint{0.100000in}{0.100000in}}{\pgfqpoint{3.007045in}{1.925000in}}%
\pgfusepath{clip}%
\pgfsetbuttcap%
\pgfsetmiterjoin%
\definecolor{currentfill}{rgb}{0.732472,0.893549,0.706282}%
\pgfsetfillcolor{currentfill}%
\pgfsetlinewidth{0.000000pt}%
\definecolor{currentstroke}{rgb}{0.000000,0.000000,0.000000}%
\pgfsetstrokecolor{currentstroke}%
\pgfsetstrokeopacity{0.000000}%
\pgfsetdash{}{0pt}%
\pgfpathmoveto{\pgfqpoint{2.253561in}{1.079316in}}%
\pgfpathlineto{\pgfqpoint{2.259515in}{1.072291in}}%
\pgfpathlineto{\pgfqpoint{2.261909in}{1.072347in}}%
\pgfpathlineto{\pgfqpoint{2.262113in}{1.062729in}}%
\pgfpathlineto{\pgfqpoint{2.267818in}{1.063028in}}%
\pgfpathlineto{\pgfqpoint{2.270021in}{1.058081in}}%
\pgfpathlineto{\pgfqpoint{2.261228in}{1.057564in}}%
\pgfpathlineto{\pgfqpoint{2.260372in}{1.054200in}}%
\pgfpathlineto{\pgfqpoint{2.239874in}{1.053973in}}%
\pgfpathlineto{\pgfqpoint{2.238225in}{1.047271in}}%
\pgfpathlineto{\pgfqpoint{2.234875in}{1.045920in}}%
\pgfpathlineto{\pgfqpoint{2.218727in}{1.043400in}}%
\pgfpathlineto{\pgfqpoint{2.217524in}{1.046026in}}%
\pgfpathlineto{\pgfqpoint{2.210809in}{1.047577in}}%
\pgfpathlineto{\pgfqpoint{2.207346in}{1.056620in}}%
\pgfpathlineto{\pgfqpoint{2.208703in}{1.064571in}}%
\pgfpathlineto{\pgfqpoint{2.204976in}{1.073343in}}%
\pgfpathlineto{\pgfqpoint{2.205852in}{1.076709in}}%
\pgfpathlineto{\pgfqpoint{2.208829in}{1.082188in}}%
\pgfpathlineto{\pgfqpoint{2.212960in}{1.083080in}}%
\pgfpathlineto{\pgfqpoint{2.211772in}{1.090499in}}%
\pgfpathlineto{\pgfqpoint{2.214133in}{1.095845in}}%
\pgfpathlineto{\pgfqpoint{2.217176in}{1.095188in}}%
\pgfpathlineto{\pgfqpoint{2.219675in}{1.101261in}}%
\pgfpathlineto{\pgfqpoint{2.223926in}{1.097599in}}%
\pgfpathlineto{\pgfqpoint{2.224931in}{1.092417in}}%
\pgfpathlineto{\pgfqpoint{2.229804in}{1.090057in}}%
\pgfpathlineto{\pgfqpoint{2.241377in}{1.090992in}}%
\pgfpathlineto{\pgfqpoint{2.246327in}{1.087565in}}%
\pgfpathlineto{\pgfqpoint{2.249127in}{1.081944in}}%
\pgfpathlineto{\pgfqpoint{2.253561in}{1.079316in}}%
\pgfpathclose%
\pgfusepath{fill}%
\end{pgfscope}%
\begin{pgfscope}%
\pgfpathrectangle{\pgfqpoint{0.100000in}{0.100000in}}{\pgfqpoint{3.007045in}{1.925000in}}%
\pgfusepath{clip}%
\pgfsetbuttcap%
\pgfsetmiterjoin%
\definecolor{currentfill}{rgb}{0.420392,0.751711,0.445829}%
\pgfsetfillcolor{currentfill}%
\pgfsetlinewidth{0.000000pt}%
\definecolor{currentstroke}{rgb}{0.000000,0.000000,0.000000}%
\pgfsetstrokecolor{currentstroke}%
\pgfsetstrokeopacity{0.000000}%
\pgfsetdash{}{0pt}%
\pgfpathmoveto{\pgfqpoint{1.596789in}{1.798977in}}%
\pgfpathlineto{\pgfqpoint{1.595955in}{1.780137in}}%
\pgfpathlineto{\pgfqpoint{1.596746in}{1.768470in}}%
\pgfpathlineto{\pgfqpoint{1.596215in}{1.756877in}}%
\pgfpathlineto{\pgfqpoint{1.573086in}{1.757987in}}%
\pgfpathlineto{\pgfqpoint{1.573646in}{1.769590in}}%
\pgfpathlineto{\pgfqpoint{1.544810in}{1.771147in}}%
\pgfpathlineto{\pgfqpoint{1.543838in}{1.782832in}}%
\pgfpathlineto{\pgfqpoint{1.544920in}{1.801546in}}%
\pgfpathlineto{\pgfqpoint{1.596789in}{1.798977in}}%
\pgfpathclose%
\pgfusepath{fill}%
\end{pgfscope}%
\begin{pgfscope}%
\pgfpathrectangle{\pgfqpoint{0.100000in}{0.100000in}}{\pgfqpoint{3.007045in}{1.925000in}}%
\pgfusepath{clip}%
\pgfsetbuttcap%
\pgfsetmiterjoin%
\definecolor{currentfill}{rgb}{0.783160,0.914833,0.755894}%
\pgfsetfillcolor{currentfill}%
\pgfsetlinewidth{0.000000pt}%
\definecolor{currentstroke}{rgb}{0.000000,0.000000,0.000000}%
\pgfsetstrokecolor{currentstroke}%
\pgfsetstrokeopacity{0.000000}%
\pgfsetdash{}{0pt}%
\pgfpathmoveto{\pgfqpoint{1.694674in}{0.889113in}}%
\pgfpathlineto{\pgfqpoint{1.694493in}{0.877607in}}%
\pgfpathlineto{\pgfqpoint{1.691843in}{0.877659in}}%
\pgfpathlineto{\pgfqpoint{1.691602in}{0.864102in}}%
\pgfpathlineto{\pgfqpoint{1.683133in}{0.863693in}}%
\pgfpathlineto{\pgfqpoint{1.678967in}{0.860772in}}%
\pgfpathlineto{\pgfqpoint{1.674047in}{0.866744in}}%
\pgfpathlineto{\pgfqpoint{1.670026in}{0.865689in}}%
\pgfpathlineto{\pgfqpoint{1.667600in}{0.860848in}}%
\pgfpathlineto{\pgfqpoint{1.627761in}{0.861901in}}%
\pgfpathlineto{\pgfqpoint{1.628500in}{0.893638in}}%
\pgfpathlineto{\pgfqpoint{1.619336in}{0.893761in}}%
\pgfpathlineto{\pgfqpoint{1.609800in}{0.896970in}}%
\pgfpathlineto{\pgfqpoint{1.594438in}{0.897482in}}%
\pgfpathlineto{\pgfqpoint{1.594488in}{0.908994in}}%
\pgfpathlineto{\pgfqpoint{1.594900in}{0.920448in}}%
\pgfpathlineto{\pgfqpoint{1.600584in}{0.920240in}}%
\pgfpathlineto{\pgfqpoint{1.601428in}{0.949236in}}%
\pgfpathlineto{\pgfqpoint{1.641329in}{0.947962in}}%
\pgfpathlineto{\pgfqpoint{1.647016in}{0.947800in}}%
\pgfpathlineto{\pgfqpoint{1.646710in}{0.935974in}}%
\pgfpathlineto{\pgfqpoint{1.658027in}{0.932849in}}%
\pgfpathlineto{\pgfqpoint{1.685760in}{0.932254in}}%
\pgfpathlineto{\pgfqpoint{1.685427in}{0.912299in}}%
\pgfpathlineto{\pgfqpoint{1.684983in}{0.896595in}}%
\pgfpathlineto{\pgfqpoint{1.691799in}{0.898304in}}%
\pgfpathlineto{\pgfqpoint{1.694674in}{0.889113in}}%
\pgfpathclose%
\pgfusepath{fill}%
\end{pgfscope}%
\begin{pgfscope}%
\pgfpathrectangle{\pgfqpoint{0.100000in}{0.100000in}}{\pgfqpoint{3.007045in}{1.925000in}}%
\pgfusepath{clip}%
\pgfsetbuttcap%
\pgfsetmiterjoin%
\definecolor{currentfill}{rgb}{0.382745,0.733256,0.427374}%
\pgfsetfillcolor{currentfill}%
\pgfsetlinewidth{0.000000pt}%
\definecolor{currentstroke}{rgb}{0.000000,0.000000,0.000000}%
\pgfsetstrokecolor{currentstroke}%
\pgfsetstrokeopacity{0.000000}%
\pgfsetdash{}{0pt}%
\pgfpathmoveto{\pgfqpoint{1.776405in}{1.429765in}}%
\pgfpathlineto{\pgfqpoint{1.779303in}{1.429779in}}%
\pgfpathlineto{\pgfqpoint{1.779266in}{1.452736in}}%
\pgfpathlineto{\pgfqpoint{1.831065in}{1.453126in}}%
\pgfpathlineto{\pgfqpoint{1.836747in}{1.453219in}}%
\pgfpathlineto{\pgfqpoint{1.837085in}{1.430186in}}%
\pgfpathlineto{\pgfqpoint{1.821649in}{1.429987in}}%
\pgfpathlineto{\pgfqpoint{1.799028in}{1.429855in}}%
\pgfpathlineto{\pgfqpoint{1.799116in}{1.413633in}}%
\pgfpathlineto{\pgfqpoint{1.776471in}{1.413556in}}%
\pgfpathlineto{\pgfqpoint{1.776405in}{1.429765in}}%
\pgfpathclose%
\pgfusepath{fill}%
\end{pgfscope}%
\begin{pgfscope}%
\pgfpathrectangle{\pgfqpoint{0.100000in}{0.100000in}}{\pgfqpoint{3.007045in}{1.925000in}}%
\pgfusepath{clip}%
\pgfsetbuttcap%
\pgfsetmiterjoin%
\definecolor{currentfill}{rgb}{0.681046,0.871895,0.656209}%
\pgfsetfillcolor{currentfill}%
\pgfsetlinewidth{0.000000pt}%
\definecolor{currentstroke}{rgb}{0.000000,0.000000,0.000000}%
\pgfsetstrokecolor{currentstroke}%
\pgfsetstrokeopacity{0.000000}%
\pgfsetdash{}{0pt}%
\pgfpathmoveto{\pgfqpoint{2.353818in}{1.009495in}}%
\pgfpathlineto{\pgfqpoint{2.341210in}{1.008147in}}%
\pgfpathlineto{\pgfqpoint{2.336554in}{1.012826in}}%
\pgfpathlineto{\pgfqpoint{2.332746in}{1.021710in}}%
\pgfpathlineto{\pgfqpoint{2.334697in}{1.028538in}}%
\pgfpathlineto{\pgfqpoint{2.332415in}{1.038665in}}%
\pgfpathlineto{\pgfqpoint{2.333634in}{1.044388in}}%
\pgfpathlineto{\pgfqpoint{2.334876in}{1.050418in}}%
\pgfpathlineto{\pgfqpoint{2.340592in}{1.057330in}}%
\pgfpathlineto{\pgfqpoint{2.347601in}{1.052815in}}%
\pgfpathlineto{\pgfqpoint{2.350874in}{1.053026in}}%
\pgfpathlineto{\pgfqpoint{2.358324in}{1.060460in}}%
\pgfpathlineto{\pgfqpoint{2.363870in}{1.060375in}}%
\pgfpathlineto{\pgfqpoint{2.372484in}{1.056245in}}%
\pgfpathlineto{\pgfqpoint{2.371954in}{1.050139in}}%
\pgfpathlineto{\pgfqpoint{2.376071in}{1.035433in}}%
\pgfpathlineto{\pgfqpoint{2.362675in}{1.024583in}}%
\pgfpathlineto{\pgfqpoint{2.360603in}{1.018528in}}%
\pgfpathlineto{\pgfqpoint{2.355024in}{1.013333in}}%
\pgfpathlineto{\pgfqpoint{2.353818in}{1.009495in}}%
\pgfpathclose%
\pgfusepath{fill}%
\end{pgfscope}%
\begin{pgfscope}%
\pgfpathrectangle{\pgfqpoint{0.100000in}{0.100000in}}{\pgfqpoint{3.007045in}{1.925000in}}%
\pgfusepath{clip}%
\pgfsetbuttcap%
\pgfsetmiterjoin%
\definecolor{currentfill}{rgb}{0.657670,0.862053,0.633449}%
\pgfsetfillcolor{currentfill}%
\pgfsetlinewidth{0.000000pt}%
\definecolor{currentstroke}{rgb}{0.000000,0.000000,0.000000}%
\pgfsetstrokecolor{currentstroke}%
\pgfsetstrokeopacity{0.000000}%
\pgfsetdash{}{0pt}%
\pgfpathmoveto{\pgfqpoint{2.321129in}{1.054650in}}%
\pgfpathlineto{\pgfqpoint{2.318137in}{1.056117in}}%
\pgfpathlineto{\pgfqpoint{2.310463in}{1.047437in}}%
\pgfpathlineto{\pgfqpoint{2.309458in}{1.055417in}}%
\pgfpathlineto{\pgfqpoint{2.302296in}{1.059040in}}%
\pgfpathlineto{\pgfqpoint{2.301227in}{1.065647in}}%
\pgfpathlineto{\pgfqpoint{2.291623in}{1.065687in}}%
\pgfpathlineto{\pgfqpoint{2.291416in}{1.086166in}}%
\pgfpathlineto{\pgfqpoint{2.289598in}{1.088435in}}%
\pgfpathlineto{\pgfqpoint{2.294182in}{1.093073in}}%
\pgfpathlineto{\pgfqpoint{2.300122in}{1.093676in}}%
\pgfpathlineto{\pgfqpoint{2.314503in}{1.080433in}}%
\pgfpathlineto{\pgfqpoint{2.316608in}{1.083386in}}%
\pgfpathlineto{\pgfqpoint{2.327579in}{1.069371in}}%
\pgfpathlineto{\pgfqpoint{2.322848in}{1.065287in}}%
\pgfpathlineto{\pgfqpoint{2.321129in}{1.054650in}}%
\pgfpathclose%
\pgfusepath{fill}%
\end{pgfscope}%
\begin{pgfscope}%
\pgfpathrectangle{\pgfqpoint{0.100000in}{0.100000in}}{\pgfqpoint{3.007045in}{1.925000in}}%
\pgfusepath{clip}%
\pgfsetbuttcap%
\pgfsetmiterjoin%
\definecolor{currentfill}{rgb}{0.601615,0.837093,0.583376}%
\pgfsetfillcolor{currentfill}%
\pgfsetlinewidth{0.000000pt}%
\definecolor{currentstroke}{rgb}{0.000000,0.000000,0.000000}%
\pgfsetstrokecolor{currentstroke}%
\pgfsetstrokeopacity{0.000000}%
\pgfsetdash{}{0pt}%
\pgfpathmoveto{\pgfqpoint{1.418635in}{1.235643in}}%
\pgfpathlineto{\pgfqpoint{1.414767in}{1.184726in}}%
\pgfpathlineto{\pgfqpoint{1.358836in}{1.188865in}}%
\pgfpathlineto{\pgfqpoint{1.330895in}{1.191503in}}%
\pgfpathlineto{\pgfqpoint{1.333705in}{1.220024in}}%
\pgfpathlineto{\pgfqpoint{1.345447in}{1.218910in}}%
\pgfpathlineto{\pgfqpoint{1.348876in}{1.253120in}}%
\pgfpathlineto{\pgfqpoint{1.343169in}{1.253667in}}%
\pgfpathlineto{\pgfqpoint{1.346503in}{1.285013in}}%
\pgfpathlineto{\pgfqpoint{1.392134in}{1.280968in}}%
\pgfpathlineto{\pgfqpoint{1.390844in}{1.264339in}}%
\pgfpathlineto{\pgfqpoint{1.420694in}{1.261942in}}%
\pgfpathlineto{\pgfqpoint{1.420426in}{1.258524in}}%
\pgfpathlineto{\pgfqpoint{1.418635in}{1.235643in}}%
\pgfpathclose%
\pgfusepath{fill}%
\end{pgfscope}%
\begin{pgfscope}%
\pgfpathrectangle{\pgfqpoint{0.100000in}{0.100000in}}{\pgfqpoint{3.007045in}{1.925000in}}%
\pgfusepath{clip}%
\pgfsetbuttcap%
\pgfsetmiterjoin%
\definecolor{currentfill}{rgb}{0.601615,0.837093,0.583376}%
\pgfsetfillcolor{currentfill}%
\pgfsetlinewidth{0.000000pt}%
\definecolor{currentstroke}{rgb}{0.000000,0.000000,0.000000}%
\pgfsetstrokecolor{currentstroke}%
\pgfsetstrokeopacity{0.000000}%
\pgfsetdash{}{0pt}%
\pgfpathmoveto{\pgfqpoint{1.555230in}{0.797022in}}%
\pgfpathlineto{\pgfqpoint{1.553928in}{0.768166in}}%
\pgfpathlineto{\pgfqpoint{1.525314in}{0.769593in}}%
\pgfpathlineto{\pgfqpoint{1.526290in}{0.791703in}}%
\pgfpathlineto{\pgfqpoint{1.526556in}{0.798349in}}%
\pgfpathlineto{\pgfqpoint{1.555230in}{0.797022in}}%
\pgfpathclose%
\pgfusepath{fill}%
\end{pgfscope}%
\begin{pgfscope}%
\pgfpathrectangle{\pgfqpoint{0.100000in}{0.100000in}}{\pgfqpoint{3.007045in}{1.925000in}}%
\pgfusepath{clip}%
\pgfsetbuttcap%
\pgfsetmiterjoin%
\definecolor{currentfill}{rgb}{0.579469,0.826759,0.565167}%
\pgfsetfillcolor{currentfill}%
\pgfsetlinewidth{0.000000pt}%
\definecolor{currentstroke}{rgb}{0.000000,0.000000,0.000000}%
\pgfsetstrokecolor{currentstroke}%
\pgfsetstrokeopacity{0.000000}%
\pgfsetdash{}{0pt}%
\pgfpathmoveto{\pgfqpoint{1.882047in}{1.149212in}}%
\pgfpathlineto{\pgfqpoint{1.877465in}{1.141048in}}%
\pgfpathlineto{\pgfqpoint{1.877944in}{1.138491in}}%
\pgfpathlineto{\pgfqpoint{1.872251in}{1.132295in}}%
\pgfpathlineto{\pgfqpoint{1.872212in}{1.129384in}}%
\pgfpathlineto{\pgfqpoint{1.849349in}{1.129915in}}%
\pgfpathlineto{\pgfqpoint{1.848856in}{1.148030in}}%
\pgfpathlineto{\pgfqpoint{1.840824in}{1.149749in}}%
\pgfpathlineto{\pgfqpoint{1.835689in}{1.147108in}}%
\pgfpathlineto{\pgfqpoint{1.835366in}{1.168017in}}%
\pgfpathlineto{\pgfqpoint{1.834849in}{1.196667in}}%
\pgfpathlineto{\pgfqpoint{1.855015in}{1.197513in}}%
\pgfpathlineto{\pgfqpoint{1.855189in}{1.180142in}}%
\pgfpathlineto{\pgfqpoint{1.859241in}{1.178993in}}%
\pgfpathlineto{\pgfqpoint{1.860786in}{1.170457in}}%
\pgfpathlineto{\pgfqpoint{1.859930in}{1.165889in}}%
\pgfpathlineto{\pgfqpoint{1.862920in}{1.162275in}}%
\pgfpathlineto{\pgfqpoint{1.867949in}{1.161555in}}%
\pgfpathlineto{\pgfqpoint{1.870359in}{1.156300in}}%
\pgfpathlineto{\pgfqpoint{1.875772in}{1.154964in}}%
\pgfpathlineto{\pgfqpoint{1.877817in}{1.149309in}}%
\pgfpathlineto{\pgfqpoint{1.882047in}{1.149212in}}%
\pgfpathclose%
\pgfusepath{fill}%
\end{pgfscope}%
\begin{pgfscope}%
\pgfpathrectangle{\pgfqpoint{0.100000in}{0.100000in}}{\pgfqpoint{3.007045in}{1.925000in}}%
\pgfusepath{clip}%
\pgfsetbuttcap%
\pgfsetmiterjoin%
\definecolor{currentfill}{rgb}{0.794233,0.919262,0.767705}%
\pgfsetfillcolor{currentfill}%
\pgfsetlinewidth{0.000000pt}%
\definecolor{currentstroke}{rgb}{0.000000,0.000000,0.000000}%
\pgfsetstrokecolor{currentstroke}%
\pgfsetstrokeopacity{0.000000}%
\pgfsetdash{}{0pt}%
\pgfpathmoveto{\pgfqpoint{2.422257in}{1.215672in}}%
\pgfpathlineto{\pgfqpoint{2.422680in}{1.209735in}}%
\pgfpathlineto{\pgfqpoint{2.416942in}{1.209349in}}%
\pgfpathlineto{\pgfqpoint{2.417290in}{1.203518in}}%
\pgfpathlineto{\pgfqpoint{2.405836in}{1.202691in}}%
\pgfpathlineto{\pgfqpoint{2.405917in}{1.200954in}}%
\pgfpathlineto{\pgfqpoint{2.394275in}{1.199273in}}%
\pgfpathlineto{\pgfqpoint{2.393794in}{1.205858in}}%
\pgfpathlineto{\pgfqpoint{2.381162in}{1.206723in}}%
\pgfpathlineto{\pgfqpoint{2.366865in}{1.205475in}}%
\pgfpathlineto{\pgfqpoint{2.366064in}{1.217273in}}%
\pgfpathlineto{\pgfqpoint{2.345734in}{1.216099in}}%
\pgfpathlineto{\pgfqpoint{2.348588in}{1.219904in}}%
\pgfpathlineto{\pgfqpoint{2.350594in}{1.236213in}}%
\pgfpathlineto{\pgfqpoint{2.349506in}{1.250356in}}%
\pgfpathlineto{\pgfqpoint{2.346617in}{1.250238in}}%
\pgfpathlineto{\pgfqpoint{2.345898in}{1.268387in}}%
\pgfpathlineto{\pgfqpoint{2.350943in}{1.269110in}}%
\pgfpathlineto{\pgfqpoint{2.359392in}{1.270321in}}%
\pgfpathlineto{\pgfqpoint{2.359900in}{1.266193in}}%
\pgfpathlineto{\pgfqpoint{2.376250in}{1.266421in}}%
\pgfpathlineto{\pgfqpoint{2.376485in}{1.262685in}}%
\pgfpathlineto{\pgfqpoint{2.390704in}{1.263733in}}%
\pgfpathlineto{\pgfqpoint{2.389048in}{1.277134in}}%
\pgfpathlineto{\pgfqpoint{2.402337in}{1.279030in}}%
\pgfpathlineto{\pgfqpoint{2.409849in}{1.281241in}}%
\pgfpathlineto{\pgfqpoint{2.411854in}{1.281854in}}%
\pgfpathlineto{\pgfqpoint{2.414093in}{1.259940in}}%
\pgfpathlineto{\pgfqpoint{2.415184in}{1.255293in}}%
\pgfpathlineto{\pgfqpoint{2.416625in}{1.241076in}}%
\pgfpathlineto{\pgfqpoint{2.415203in}{1.238420in}}%
\pgfpathlineto{\pgfqpoint{2.403598in}{1.237933in}}%
\pgfpathlineto{\pgfqpoint{2.403956in}{1.231787in}}%
\pgfpathlineto{\pgfqpoint{2.407820in}{1.232038in}}%
\pgfpathlineto{\pgfqpoint{2.410154in}{1.226352in}}%
\pgfpathlineto{\pgfqpoint{2.410789in}{1.216745in}}%
\pgfpathlineto{\pgfqpoint{2.422257in}{1.215672in}}%
\pgfpathclose%
\pgfusepath{fill}%
\end{pgfscope}%
\begin{pgfscope}%
\pgfpathrectangle{\pgfqpoint{0.100000in}{0.100000in}}{\pgfqpoint{3.007045in}{1.925000in}}%
\pgfusepath{clip}%
\pgfsetbuttcap%
\pgfsetmiterjoin%
\definecolor{currentfill}{rgb}{0.474279,0.777670,0.478677}%
\pgfsetfillcolor{currentfill}%
\pgfsetlinewidth{0.000000pt}%
\definecolor{currentstroke}{rgb}{0.000000,0.000000,0.000000}%
\pgfsetstrokecolor{currentstroke}%
\pgfsetstrokeopacity{0.000000}%
\pgfsetdash{}{0pt}%
\pgfpathmoveto{\pgfqpoint{1.792177in}{1.207263in}}%
\pgfpathlineto{\pgfqpoint{1.792476in}{1.224430in}}%
\pgfpathlineto{\pgfqpoint{1.790743in}{1.236578in}}%
\pgfpathlineto{\pgfqpoint{1.798792in}{1.236593in}}%
\pgfpathlineto{\pgfqpoint{1.833661in}{1.237372in}}%
\pgfpathlineto{\pgfqpoint{1.834576in}{1.208087in}}%
\pgfpathlineto{\pgfqpoint{1.811789in}{1.208062in}}%
\pgfpathlineto{\pgfqpoint{1.811728in}{1.201377in}}%
\pgfpathlineto{\pgfqpoint{1.792230in}{1.201536in}}%
\pgfpathlineto{\pgfqpoint{1.792177in}{1.207263in}}%
\pgfpathclose%
\pgfusepath{fill}%
\end{pgfscope}%
\begin{pgfscope}%
\pgfpathrectangle{\pgfqpoint{0.100000in}{0.100000in}}{\pgfqpoint{3.007045in}{1.925000in}}%
\pgfusepath{clip}%
\pgfsetbuttcap%
\pgfsetmiterjoin%
\definecolor{currentfill}{rgb}{0.407843,0.745559,0.439677}%
\pgfsetfillcolor{currentfill}%
\pgfsetlinewidth{0.000000pt}%
\definecolor{currentstroke}{rgb}{0.000000,0.000000,0.000000}%
\pgfsetstrokecolor{currentstroke}%
\pgfsetstrokeopacity{0.000000}%
\pgfsetdash{}{0pt}%
\pgfpathmoveto{\pgfqpoint{1.744270in}{1.498734in}}%
\pgfpathlineto{\pgfqpoint{1.743883in}{1.504541in}}%
\pgfpathlineto{\pgfqpoint{1.703956in}{1.505078in}}%
\pgfpathlineto{\pgfqpoint{1.704306in}{1.528001in}}%
\pgfpathlineto{\pgfqpoint{1.704575in}{1.547315in}}%
\pgfpathlineto{\pgfqpoint{1.712446in}{1.545741in}}%
\pgfpathlineto{\pgfqpoint{1.717578in}{1.542068in}}%
\pgfpathlineto{\pgfqpoint{1.720167in}{1.545159in}}%
\pgfpathlineto{\pgfqpoint{1.720360in}{1.556564in}}%
\pgfpathlineto{\pgfqpoint{1.766199in}{1.556278in}}%
\pgfpathlineto{\pgfqpoint{1.783347in}{1.556314in}}%
\pgfpathlineto{\pgfqpoint{1.783351in}{1.550568in}}%
\pgfpathlineto{\pgfqpoint{1.783670in}{1.521817in}}%
\pgfpathlineto{\pgfqpoint{1.749650in}{1.521778in}}%
\pgfpathlineto{\pgfqpoint{1.749628in}{1.512662in}}%
\pgfpathlineto{\pgfqpoint{1.755381in}{1.509019in}}%
\pgfpathlineto{\pgfqpoint{1.755330in}{1.498635in}}%
\pgfpathlineto{\pgfqpoint{1.744270in}{1.498734in}}%
\pgfpathclose%
\pgfusepath{fill}%
\end{pgfscope}%
\begin{pgfscope}%
\pgfpathrectangle{\pgfqpoint{0.100000in}{0.100000in}}{\pgfqpoint{3.007045in}{1.925000in}}%
\pgfusepath{clip}%
\pgfsetbuttcap%
\pgfsetmiterjoin%
\definecolor{currentfill}{rgb}{0.389020,0.736332,0.430450}%
\pgfsetfillcolor{currentfill}%
\pgfsetlinewidth{0.000000pt}%
\definecolor{currentstroke}{rgb}{0.000000,0.000000,0.000000}%
\pgfsetstrokecolor{currentstroke}%
\pgfsetstrokeopacity{0.000000}%
\pgfsetdash{}{0pt}%
\pgfpathmoveto{\pgfqpoint{1.890280in}{0.989939in}}%
\pgfpathlineto{\pgfqpoint{1.890382in}{0.969613in}}%
\pgfpathlineto{\pgfqpoint{1.886046in}{0.969526in}}%
\pgfpathlineto{\pgfqpoint{1.832387in}{0.968604in}}%
\pgfpathlineto{\pgfqpoint{1.821240in}{0.968478in}}%
\pgfpathlineto{\pgfqpoint{1.821509in}{0.984906in}}%
\pgfpathlineto{\pgfqpoint{1.822007in}{1.020648in}}%
\pgfpathlineto{\pgfqpoint{1.844374in}{1.020395in}}%
\pgfpathlineto{\pgfqpoint{1.845462in}{1.020387in}}%
\pgfpathlineto{\pgfqpoint{1.845487in}{1.001557in}}%
\pgfpathlineto{\pgfqpoint{1.859842in}{1.001580in}}%
\pgfpathlineto{\pgfqpoint{1.859810in}{0.990030in}}%
\pgfpathlineto{\pgfqpoint{1.890280in}{0.989939in}}%
\pgfpathclose%
\pgfusepath{fill}%
\end{pgfscope}%
\begin{pgfscope}%
\pgfpathrectangle{\pgfqpoint{0.100000in}{0.100000in}}{\pgfqpoint{3.007045in}{1.925000in}}%
\pgfusepath{clip}%
\pgfsetbuttcap%
\pgfsetmiterjoin%
\definecolor{currentfill}{rgb}{0.648320,0.858116,0.624344}%
\pgfsetfillcolor{currentfill}%
\pgfsetlinewidth{0.000000pt}%
\definecolor{currentstroke}{rgb}{0.000000,0.000000,0.000000}%
\pgfsetstrokecolor{currentstroke}%
\pgfsetstrokeopacity{0.000000}%
\pgfsetdash{}{0pt}%
\pgfpathmoveto{\pgfqpoint{2.285799in}{1.004151in}}%
\pgfpathlineto{\pgfqpoint{2.311789in}{1.005478in}}%
\pgfpathlineto{\pgfqpoint{2.315188in}{1.000641in}}%
\pgfpathlineto{\pgfqpoint{2.303650in}{1.003101in}}%
\pgfpathlineto{\pgfqpoint{2.298509in}{1.001042in}}%
\pgfpathlineto{\pgfqpoint{2.295566in}{0.990662in}}%
\pgfpathlineto{\pgfqpoint{2.296204in}{0.983928in}}%
\pgfpathlineto{\pgfqpoint{2.299083in}{0.977051in}}%
\pgfpathlineto{\pgfqpoint{2.296056in}{0.970309in}}%
\pgfpathlineto{\pgfqpoint{2.291608in}{0.970240in}}%
\pgfpathlineto{\pgfqpoint{2.291000in}{0.961775in}}%
\pgfpathlineto{\pgfqpoint{2.285629in}{0.965628in}}%
\pgfpathlineto{\pgfqpoint{2.277359in}{0.967152in}}%
\pgfpathlineto{\pgfqpoint{2.269975in}{0.962101in}}%
\pgfpathlineto{\pgfqpoint{2.267876in}{0.965337in}}%
\pgfpathlineto{\pgfqpoint{2.263169in}{0.965439in}}%
\pgfpathlineto{\pgfqpoint{2.261147in}{0.968764in}}%
\pgfpathlineto{\pgfqpoint{2.261071in}{0.973710in}}%
\pgfpathlineto{\pgfqpoint{2.258038in}{0.978829in}}%
\pgfpathlineto{\pgfqpoint{2.258356in}{0.992865in}}%
\pgfpathlineto{\pgfqpoint{2.263992in}{0.994959in}}%
\pgfpathlineto{\pgfqpoint{2.271083in}{0.992318in}}%
\pgfpathlineto{\pgfqpoint{2.275809in}{0.988285in}}%
\pgfpathlineto{\pgfqpoint{2.278118in}{0.995018in}}%
\pgfpathlineto{\pgfqpoint{2.285270in}{0.997505in}}%
\pgfpathlineto{\pgfqpoint{2.285799in}{1.004151in}}%
\pgfpathclose%
\pgfusepath{fill}%
\end{pgfscope}%
\begin{pgfscope}%
\pgfpathrectangle{\pgfqpoint{0.100000in}{0.100000in}}{\pgfqpoint{3.007045in}{1.925000in}}%
\pgfusepath{clip}%
\pgfsetbuttcap%
\pgfsetmiterjoin%
\definecolor{currentfill}{rgb}{0.551788,0.813841,0.542407}%
\pgfsetfillcolor{currentfill}%
\pgfsetlinewidth{0.000000pt}%
\definecolor{currentstroke}{rgb}{0.000000,0.000000,0.000000}%
\pgfsetstrokecolor{currentstroke}%
\pgfsetstrokeopacity{0.000000}%
\pgfsetdash{}{0pt}%
\pgfpathmoveto{\pgfqpoint{1.339144in}{1.451533in}}%
\pgfpathlineto{\pgfqpoint{1.298841in}{1.455523in}}%
\pgfpathlineto{\pgfqpoint{1.245609in}{1.461365in}}%
\pgfpathlineto{\pgfqpoint{1.251044in}{1.505633in}}%
\pgfpathlineto{\pgfqpoint{1.255504in}{1.539978in}}%
\pgfpathlineto{\pgfqpoint{1.257251in}{1.559802in}}%
\pgfpathlineto{\pgfqpoint{1.262489in}{1.559658in}}%
\pgfpathlineto{\pgfqpoint{1.303187in}{1.554906in}}%
\pgfpathlineto{\pgfqpoint{1.348840in}{1.549898in}}%
\pgfpathlineto{\pgfqpoint{1.343611in}{1.496079in}}%
\pgfpathlineto{\pgfqpoint{1.339144in}{1.451533in}}%
\pgfpathclose%
\pgfusepath{fill}%
\end{pgfscope}%
\begin{pgfscope}%
\pgfpathrectangle{\pgfqpoint{0.100000in}{0.100000in}}{\pgfqpoint{3.007045in}{1.925000in}}%
\pgfusepath{clip}%
\pgfsetbuttcap%
\pgfsetmiterjoin%
\definecolor{currentfill}{rgb}{0.351373,0.717878,0.411995}%
\pgfsetfillcolor{currentfill}%
\pgfsetlinewidth{0.000000pt}%
\definecolor{currentstroke}{rgb}{0.000000,0.000000,0.000000}%
\pgfsetstrokecolor{currentstroke}%
\pgfsetstrokeopacity{0.000000}%
\pgfsetdash{}{0pt}%
\pgfpathmoveto{\pgfqpoint{1.154637in}{0.238482in}}%
\pgfpathlineto{\pgfqpoint{1.153566in}{0.243456in}}%
\pgfpathlineto{\pgfqpoint{1.148403in}{0.253718in}}%
\pgfpathlineto{\pgfqpoint{1.150064in}{0.259799in}}%
\pgfpathlineto{\pgfqpoint{1.161104in}{0.270073in}}%
\pgfpathlineto{\pgfqpoint{1.164300in}{0.280924in}}%
\pgfpathlineto{\pgfqpoint{1.168171in}{0.288539in}}%
\pgfpathlineto{\pgfqpoint{1.167443in}{0.293757in}}%
\pgfpathlineto{\pgfqpoint{1.171591in}{0.300793in}}%
\pgfpathlineto{\pgfqpoint{1.191981in}{0.303197in}}%
\pgfpathlineto{\pgfqpoint{1.194052in}{0.305195in}}%
\pgfpathlineto{\pgfqpoint{1.194899in}{0.313558in}}%
\pgfpathlineto{\pgfqpoint{1.199467in}{0.319810in}}%
\pgfpathlineto{\pgfqpoint{1.203895in}{0.319010in}}%
\pgfpathlineto{\pgfqpoint{1.206973in}{0.312220in}}%
\pgfpathlineto{\pgfqpoint{1.209897in}{0.301367in}}%
\pgfpathlineto{\pgfqpoint{1.215982in}{0.294627in}}%
\pgfpathlineto{\pgfqpoint{1.221823in}{0.282731in}}%
\pgfpathlineto{\pgfqpoint{1.224745in}{0.266660in}}%
\pgfpathlineto{\pgfqpoint{1.219165in}{0.259374in}}%
\pgfpathlineto{\pgfqpoint{1.223841in}{0.256519in}}%
\pgfpathlineto{\pgfqpoint{1.220981in}{0.250601in}}%
\pgfpathlineto{\pgfqpoint{1.221882in}{0.243306in}}%
\pgfpathlineto{\pgfqpoint{1.224870in}{0.237321in}}%
\pgfpathlineto{\pgfqpoint{1.222256in}{0.234938in}}%
\pgfpathlineto{\pgfqpoint{1.208527in}{0.234063in}}%
\pgfpathlineto{\pgfqpoint{1.197467in}{0.236205in}}%
\pgfpathlineto{\pgfqpoint{1.188592in}{0.242263in}}%
\pgfpathlineto{\pgfqpoint{1.183559in}{0.241689in}}%
\pgfpathlineto{\pgfqpoint{1.172818in}{0.243600in}}%
\pgfpathlineto{\pgfqpoint{1.154637in}{0.238482in}}%
\pgfpathclose%
\pgfusepath{fill}%
\end{pgfscope}%
\begin{pgfscope}%
\pgfpathrectangle{\pgfqpoint{0.100000in}{0.100000in}}{\pgfqpoint{3.007045in}{1.925000in}}%
\pgfusepath{clip}%
\pgfsetbuttcap%
\pgfsetmiterjoin%
\definecolor{currentfill}{rgb}{0.479815,0.780254,0.483230}%
\pgfsetfillcolor{currentfill}%
\pgfsetlinewidth{0.000000pt}%
\definecolor{currentstroke}{rgb}{0.000000,0.000000,0.000000}%
\pgfsetstrokecolor{currentstroke}%
\pgfsetstrokeopacity{0.000000}%
\pgfsetdash{}{0pt}%
\pgfpathmoveto{\pgfqpoint{1.823694in}{1.499014in}}%
\pgfpathlineto{\pgfqpoint{1.822887in}{1.493279in}}%
\pgfpathlineto{\pgfqpoint{1.790073in}{1.492987in}}%
\pgfpathlineto{\pgfqpoint{1.789742in}{1.510236in}}%
\pgfpathlineto{\pgfqpoint{1.795936in}{1.510293in}}%
\pgfpathlineto{\pgfqpoint{1.795889in}{1.521880in}}%
\pgfpathlineto{\pgfqpoint{1.783670in}{1.521817in}}%
\pgfpathlineto{\pgfqpoint{1.783351in}{1.550568in}}%
\pgfpathlineto{\pgfqpoint{1.801048in}{1.550695in}}%
\pgfpathlineto{\pgfqpoint{1.806790in}{1.547902in}}%
\pgfpathlineto{\pgfqpoint{1.807185in}{1.527732in}}%
\pgfpathlineto{\pgfqpoint{1.818599in}{1.527791in}}%
\pgfpathlineto{\pgfqpoint{1.818879in}{1.510489in}}%
\pgfpathlineto{\pgfqpoint{1.824600in}{1.510542in}}%
\pgfpathlineto{\pgfqpoint{1.823694in}{1.499014in}}%
\pgfpathclose%
\pgfusepath{fill}%
\end{pgfscope}%
\begin{pgfscope}%
\pgfpathrectangle{\pgfqpoint{0.100000in}{0.100000in}}{\pgfqpoint{3.007045in}{1.925000in}}%
\pgfusepath{clip}%
\pgfsetbuttcap%
\pgfsetmiterjoin%
\definecolor{currentfill}{rgb}{0.590542,0.831926,0.574271}%
\pgfsetfillcolor{currentfill}%
\pgfsetlinewidth{0.000000pt}%
\definecolor{currentstroke}{rgb}{0.000000,0.000000,0.000000}%
\pgfsetstrokecolor{currentstroke}%
\pgfsetstrokeopacity{0.000000}%
\pgfsetdash{}{0pt}%
\pgfpathmoveto{\pgfqpoint{2.214133in}{1.095845in}}%
\pgfpathlineto{\pgfqpoint{2.211772in}{1.090499in}}%
\pgfpathlineto{\pgfqpoint{2.212960in}{1.083080in}}%
\pgfpathlineto{\pgfqpoint{2.208829in}{1.082188in}}%
\pgfpathlineto{\pgfqpoint{2.205852in}{1.076709in}}%
\pgfpathlineto{\pgfqpoint{2.205878in}{1.081086in}}%
\pgfpathlineto{\pgfqpoint{2.200736in}{1.080556in}}%
\pgfpathlineto{\pgfqpoint{2.197741in}{1.085673in}}%
\pgfpathlineto{\pgfqpoint{2.185322in}{1.077996in}}%
\pgfpathlineto{\pgfqpoint{2.184734in}{1.072536in}}%
\pgfpathlineto{\pgfqpoint{2.172010in}{1.077414in}}%
\pgfpathlineto{\pgfqpoint{2.165193in}{1.079905in}}%
\pgfpathlineto{\pgfqpoint{2.158994in}{1.078092in}}%
\pgfpathlineto{\pgfqpoint{2.156206in}{1.071747in}}%
\pgfpathlineto{\pgfqpoint{2.152373in}{1.076112in}}%
\pgfpathlineto{\pgfqpoint{2.144755in}{1.073731in}}%
\pgfpathlineto{\pgfqpoint{2.139094in}{1.074231in}}%
\pgfpathlineto{\pgfqpoint{2.141157in}{1.068895in}}%
\pgfpathlineto{\pgfqpoint{2.134908in}{1.067855in}}%
\pgfpathlineto{\pgfqpoint{2.131121in}{1.073626in}}%
\pgfpathlineto{\pgfqpoint{2.137521in}{1.089690in}}%
\pgfpathlineto{\pgfqpoint{2.134514in}{1.098251in}}%
\pgfpathlineto{\pgfqpoint{2.136503in}{1.101602in}}%
\pgfpathlineto{\pgfqpoint{2.134603in}{1.106062in}}%
\pgfpathlineto{\pgfqpoint{2.135045in}{1.113976in}}%
\pgfpathlineto{\pgfqpoint{2.136969in}{1.118978in}}%
\pgfpathlineto{\pgfqpoint{2.150356in}{1.119886in}}%
\pgfpathlineto{\pgfqpoint{2.154476in}{1.114634in}}%
\pgfpathlineto{\pgfqpoint{2.160227in}{1.118407in}}%
\pgfpathlineto{\pgfqpoint{2.171500in}{1.120112in}}%
\pgfpathlineto{\pgfqpoint{2.178440in}{1.120233in}}%
\pgfpathlineto{\pgfqpoint{2.190430in}{1.118463in}}%
\pgfpathlineto{\pgfqpoint{2.192366in}{1.120724in}}%
\pgfpathlineto{\pgfqpoint{2.200334in}{1.121433in}}%
\pgfpathlineto{\pgfqpoint{2.201173in}{1.112813in}}%
\pgfpathlineto{\pgfqpoint{2.202065in}{1.104184in}}%
\pgfpathlineto{\pgfqpoint{2.207678in}{1.104917in}}%
\pgfpathlineto{\pgfqpoint{2.207958in}{1.101054in}}%
\pgfpathlineto{\pgfqpoint{2.213695in}{1.101544in}}%
\pgfpathlineto{\pgfqpoint{2.214133in}{1.095845in}}%
\pgfpathclose%
\pgfusepath{fill}%
\end{pgfscope}%
\begin{pgfscope}%
\pgfpathrectangle{\pgfqpoint{0.100000in}{0.100000in}}{\pgfqpoint{3.007045in}{1.925000in}}%
\pgfusepath{clip}%
\pgfsetbuttcap%
\pgfsetmiterjoin%
\definecolor{currentfill}{rgb}{0.395294,0.739408,0.433526}%
\pgfsetfillcolor{currentfill}%
\pgfsetlinewidth{0.000000pt}%
\definecolor{currentstroke}{rgb}{0.000000,0.000000,0.000000}%
\pgfsetstrokecolor{currentstroke}%
\pgfsetstrokeopacity{0.000000}%
\pgfsetdash{}{0pt}%
\pgfpathmoveto{\pgfqpoint{1.929386in}{1.648471in}}%
\pgfpathlineto{\pgfqpoint{1.930881in}{1.608680in}}%
\pgfpathlineto{\pgfqpoint{1.907975in}{1.607918in}}%
\pgfpathlineto{\pgfqpoint{1.896759in}{1.607543in}}%
\pgfpathlineto{\pgfqpoint{1.896284in}{1.624795in}}%
\pgfpathlineto{\pgfqpoint{1.861397in}{1.624038in}}%
\pgfpathlineto{\pgfqpoint{1.860671in}{1.638052in}}%
\pgfpathlineto{\pgfqpoint{1.860483in}{1.664391in}}%
\pgfpathlineto{\pgfqpoint{1.859346in}{1.688805in}}%
\pgfpathlineto{\pgfqpoint{1.859255in}{1.710576in}}%
\pgfpathlineto{\pgfqpoint{1.858121in}{1.722001in}}%
\pgfpathlineto{\pgfqpoint{1.857189in}{1.733722in}}%
\pgfpathlineto{\pgfqpoint{1.856761in}{1.771102in}}%
\pgfpathlineto{\pgfqpoint{1.861336in}{1.770943in}}%
\pgfpathlineto{\pgfqpoint{1.872825in}{1.765570in}}%
\pgfpathlineto{\pgfqpoint{1.876937in}{1.765915in}}%
\pgfpathlineto{\pgfqpoint{1.876187in}{1.758800in}}%
\pgfpathlineto{\pgfqpoint{1.882756in}{1.759721in}}%
\pgfpathlineto{\pgfqpoint{1.889291in}{1.744713in}}%
\pgfpathlineto{\pgfqpoint{1.893197in}{1.746419in}}%
\pgfpathlineto{\pgfqpoint{1.892633in}{1.752996in}}%
\pgfpathlineto{\pgfqpoint{1.902968in}{1.754402in}}%
\pgfpathlineto{\pgfqpoint{1.905332in}{1.748213in}}%
\pgfpathlineto{\pgfqpoint{1.911782in}{1.744531in}}%
\pgfpathlineto{\pgfqpoint{1.918443in}{1.744191in}}%
\pgfpathlineto{\pgfqpoint{1.918786in}{1.738606in}}%
\pgfpathlineto{\pgfqpoint{1.931137in}{1.734731in}}%
\pgfpathlineto{\pgfqpoint{1.939426in}{1.737368in}}%
\pgfpathlineto{\pgfqpoint{1.948847in}{1.744739in}}%
\pgfpathlineto{\pgfqpoint{1.950215in}{1.724734in}}%
\pgfpathlineto{\pgfqpoint{1.951342in}{1.696561in}}%
\pgfpathlineto{\pgfqpoint{1.941485in}{1.683918in}}%
\pgfpathlineto{\pgfqpoint{1.931784in}{1.673108in}}%
\pgfpathlineto{\pgfqpoint{1.927592in}{1.670463in}}%
\pgfpathlineto{\pgfqpoint{1.918507in}{1.660575in}}%
\pgfpathlineto{\pgfqpoint{1.904812in}{1.649918in}}%
\pgfpathlineto{\pgfqpoint{1.907843in}{1.644693in}}%
\pgfpathlineto{\pgfqpoint{1.911711in}{1.642752in}}%
\pgfpathlineto{\pgfqpoint{1.918643in}{1.643971in}}%
\pgfpathlineto{\pgfqpoint{1.929386in}{1.648471in}}%
\pgfpathclose%
\pgfusepath{fill}%
\end{pgfscope}%
\begin{pgfscope}%
\pgfpathrectangle{\pgfqpoint{0.100000in}{0.100000in}}{\pgfqpoint{3.007045in}{1.925000in}}%
\pgfusepath{clip}%
\pgfsetbuttcap%
\pgfsetmiterjoin%
\definecolor{currentfill}{rgb}{0.657670,0.862053,0.633449}%
\pgfsetfillcolor{currentfill}%
\pgfsetlinewidth{0.000000pt}%
\definecolor{currentstroke}{rgb}{0.000000,0.000000,0.000000}%
\pgfsetstrokecolor{currentstroke}%
\pgfsetstrokeopacity{0.000000}%
\pgfsetdash{}{0pt}%
\pgfpathmoveto{\pgfqpoint{1.923061in}{1.300995in}}%
\pgfpathlineto{\pgfqpoint{1.905724in}{1.300373in}}%
\pgfpathlineto{\pgfqpoint{1.904114in}{1.352301in}}%
\pgfpathlineto{\pgfqpoint{1.938345in}{1.353373in}}%
\pgfpathlineto{\pgfqpoint{1.972422in}{1.354713in}}%
\pgfpathlineto{\pgfqpoint{1.973436in}{1.331705in}}%
\pgfpathlineto{\pgfqpoint{1.950534in}{1.330806in}}%
\pgfpathlineto{\pgfqpoint{1.950752in}{1.325057in}}%
\pgfpathlineto{\pgfqpoint{1.927902in}{1.324319in}}%
\pgfpathlineto{\pgfqpoint{1.928815in}{1.301237in}}%
\pgfpathlineto{\pgfqpoint{1.923061in}{1.300995in}}%
\pgfpathclose%
\pgfusepath{fill}%
\end{pgfscope}%
\begin{pgfscope}%
\pgfpathrectangle{\pgfqpoint{0.100000in}{0.100000in}}{\pgfqpoint{3.007045in}{1.925000in}}%
\pgfusepath{clip}%
\pgfsetbuttcap%
\pgfsetmiterjoin%
\definecolor{currentfill}{rgb}{0.546251,0.811257,0.537855}%
\pgfsetfillcolor{currentfill}%
\pgfsetlinewidth{0.000000pt}%
\definecolor{currentstroke}{rgb}{0.000000,0.000000,0.000000}%
\pgfsetstrokecolor{currentstroke}%
\pgfsetstrokeopacity{0.000000}%
\pgfsetdash{}{0pt}%
\pgfpathmoveto{\pgfqpoint{2.549716in}{0.711373in}}%
\pgfpathlineto{\pgfqpoint{2.544812in}{0.713097in}}%
\pgfpathlineto{\pgfqpoint{2.537534in}{0.718203in}}%
\pgfpathlineto{\pgfqpoint{2.534070in}{0.723697in}}%
\pgfpathlineto{\pgfqpoint{2.526011in}{0.723942in}}%
\pgfpathlineto{\pgfqpoint{2.520931in}{0.729238in}}%
\pgfpathlineto{\pgfqpoint{2.513924in}{0.731882in}}%
\pgfpathlineto{\pgfqpoint{2.512057in}{0.734800in}}%
\pgfpathlineto{\pgfqpoint{2.530406in}{0.743490in}}%
\pgfpathlineto{\pgfqpoint{2.528796in}{0.747623in}}%
\pgfpathlineto{\pgfqpoint{2.523508in}{0.752566in}}%
\pgfpathlineto{\pgfqpoint{2.522083in}{0.760192in}}%
\pgfpathlineto{\pgfqpoint{2.529482in}{0.767093in}}%
\pgfpathlineto{\pgfqpoint{2.535487in}{0.765030in}}%
\pgfpathlineto{\pgfqpoint{2.548564in}{0.780580in}}%
\pgfpathlineto{\pgfqpoint{2.556087in}{0.772933in}}%
\pgfpathlineto{\pgfqpoint{2.559434in}{0.779009in}}%
\pgfpathlineto{\pgfqpoint{2.567570in}{0.778188in}}%
\pgfpathlineto{\pgfqpoint{2.572033in}{0.772743in}}%
\pgfpathlineto{\pgfqpoint{2.576021in}{0.772345in}}%
\pgfpathlineto{\pgfqpoint{2.581007in}{0.767715in}}%
\pgfpathlineto{\pgfqpoint{2.584609in}{0.761996in}}%
\pgfpathlineto{\pgfqpoint{2.584009in}{0.757368in}}%
\pgfpathlineto{\pgfqpoint{2.574779in}{0.751322in}}%
\pgfpathlineto{\pgfqpoint{2.570706in}{0.743972in}}%
\pgfpathlineto{\pgfqpoint{2.562709in}{0.737347in}}%
\pgfpathlineto{\pgfqpoint{2.564979in}{0.730555in}}%
\pgfpathlineto{\pgfqpoint{2.552611in}{0.712060in}}%
\pgfpathlineto{\pgfqpoint{2.549716in}{0.711373in}}%
\pgfpathclose%
\pgfusepath{fill}%
\end{pgfscope}%
\begin{pgfscope}%
\pgfpathrectangle{\pgfqpoint{0.100000in}{0.100000in}}{\pgfqpoint{3.007045in}{1.925000in}}%
\pgfusepath{clip}%
\pgfsetbuttcap%
\pgfsetmiterjoin%
\definecolor{currentfill}{rgb}{0.507497,0.793172,0.505990}%
\pgfsetfillcolor{currentfill}%
\pgfsetlinewidth{0.000000pt}%
\definecolor{currentstroke}{rgb}{0.000000,0.000000,0.000000}%
\pgfsetstrokecolor{currentstroke}%
\pgfsetstrokeopacity{0.000000}%
\pgfsetdash{}{0pt}%
\pgfpathmoveto{\pgfqpoint{0.613279in}{1.618283in}}%
\pgfpathlineto{\pgfqpoint{0.651898in}{1.607795in}}%
\pgfpathlineto{\pgfqpoint{0.653418in}{1.613391in}}%
\pgfpathlineto{\pgfqpoint{0.680340in}{1.605736in}}%
\pgfpathlineto{\pgfqpoint{0.675770in}{1.589097in}}%
\pgfpathlineto{\pgfqpoint{0.661467in}{1.533821in}}%
\pgfpathlineto{\pgfqpoint{0.662044in}{1.533668in}}%
\pgfpathlineto{\pgfqpoint{0.651488in}{1.492874in}}%
\pgfpathlineto{\pgfqpoint{0.647669in}{1.474848in}}%
\pgfpathlineto{\pgfqpoint{0.623923in}{1.480760in}}%
\pgfpathlineto{\pgfqpoint{0.594261in}{1.488948in}}%
\pgfpathlineto{\pgfqpoint{0.592563in}{1.489427in}}%
\pgfpathlineto{\pgfqpoint{0.605512in}{1.537505in}}%
\pgfpathlineto{\pgfqpoint{0.578445in}{1.544847in}}%
\pgfpathlineto{\pgfqpoint{0.586719in}{1.572193in}}%
\pgfpathlineto{\pgfqpoint{0.588372in}{1.571730in}}%
\pgfpathlineto{\pgfqpoint{0.596099in}{1.599150in}}%
\pgfpathlineto{\pgfqpoint{0.597590in}{1.604775in}}%
\pgfpathlineto{\pgfqpoint{0.603136in}{1.603204in}}%
\pgfpathlineto{\pgfqpoint{0.607867in}{1.619794in}}%
\pgfpathlineto{\pgfqpoint{0.613279in}{1.618283in}}%
\pgfpathclose%
\pgfusepath{fill}%
\end{pgfscope}%
\begin{pgfscope}%
\pgfpathrectangle{\pgfqpoint{0.100000in}{0.100000in}}{\pgfqpoint{3.007045in}{1.925000in}}%
\pgfusepath{clip}%
\pgfsetbuttcap%
\pgfsetmiterjoin%
\definecolor{currentfill}{rgb}{0.234141,0.648443,0.348097}%
\pgfsetfillcolor{currentfill}%
\pgfsetlinewidth{0.000000pt}%
\definecolor{currentstroke}{rgb}{0.000000,0.000000,0.000000}%
\pgfsetstrokecolor{currentstroke}%
\pgfsetstrokeopacity{0.000000}%
\pgfsetdash{}{0pt}%
\pgfpathmoveto{\pgfqpoint{2.623105in}{0.353228in}}%
\pgfpathlineto{\pgfqpoint{2.599929in}{0.349210in}}%
\pgfpathlineto{\pgfqpoint{2.597057in}{0.366782in}}%
\pgfpathlineto{\pgfqpoint{2.579681in}{0.363910in}}%
\pgfpathlineto{\pgfqpoint{2.577000in}{0.380965in}}%
\pgfpathlineto{\pgfqpoint{2.574578in}{0.398555in}}%
\pgfpathlineto{\pgfqpoint{2.592108in}{0.401133in}}%
\pgfpathlineto{\pgfqpoint{2.591178in}{0.407023in}}%
\pgfpathlineto{\pgfqpoint{2.597054in}{0.407964in}}%
\pgfpathlineto{\pgfqpoint{2.596119in}{0.413805in}}%
\pgfpathlineto{\pgfqpoint{2.607766in}{0.416135in}}%
\pgfpathlineto{\pgfqpoint{2.600616in}{0.426095in}}%
\pgfpathlineto{\pgfqpoint{2.597645in}{0.425944in}}%
\pgfpathlineto{\pgfqpoint{2.590385in}{0.434304in}}%
\pgfpathlineto{\pgfqpoint{2.593275in}{0.442838in}}%
\pgfpathlineto{\pgfqpoint{2.609094in}{0.445229in}}%
\pgfpathlineto{\pgfqpoint{2.614673in}{0.446167in}}%
\pgfpathlineto{\pgfqpoint{2.615609in}{0.440573in}}%
\pgfpathlineto{\pgfqpoint{2.642481in}{0.444772in}}%
\pgfpathlineto{\pgfqpoint{2.647094in}{0.435626in}}%
\pgfpathlineto{\pgfqpoint{2.655837in}{0.422084in}}%
\pgfpathlineto{\pgfqpoint{2.661202in}{0.412146in}}%
\pgfpathlineto{\pgfqpoint{2.667843in}{0.397101in}}%
\pgfpathlineto{\pgfqpoint{2.669777in}{0.383631in}}%
\pgfpathlineto{\pgfqpoint{2.670514in}{0.365060in}}%
\pgfpathlineto{\pgfqpoint{2.661292in}{0.364487in}}%
\pgfpathlineto{\pgfqpoint{2.622244in}{0.358169in}}%
\pgfpathlineto{\pgfqpoint{2.623105in}{0.353228in}}%
\pgfpathclose%
\pgfusepath{fill}%
\end{pgfscope}%
\begin{pgfscope}%
\pgfpathrectangle{\pgfqpoint{0.100000in}{0.100000in}}{\pgfqpoint{3.007045in}{1.925000in}}%
\pgfusepath{clip}%
\pgfsetbuttcap%
\pgfsetmiterjoin%
\definecolor{currentfill}{rgb}{0.681046,0.871895,0.656209}%
\pgfsetfillcolor{currentfill}%
\pgfsetlinewidth{0.000000pt}%
\definecolor{currentstroke}{rgb}{0.000000,0.000000,0.000000}%
\pgfsetstrokecolor{currentstroke}%
\pgfsetstrokeopacity{0.000000}%
\pgfsetdash{}{0pt}%
\pgfpathmoveto{\pgfqpoint{2.201173in}{1.112813in}}%
\pgfpathlineto{\pgfqpoint{2.200334in}{1.121433in}}%
\pgfpathlineto{\pgfqpoint{2.192366in}{1.120724in}}%
\pgfpathlineto{\pgfqpoint{2.190430in}{1.118463in}}%
\pgfpathlineto{\pgfqpoint{2.178440in}{1.120233in}}%
\pgfpathlineto{\pgfqpoint{2.171500in}{1.120112in}}%
\pgfpathlineto{\pgfqpoint{2.168758in}{1.129637in}}%
\pgfpathlineto{\pgfqpoint{2.170118in}{1.138155in}}%
\pgfpathlineto{\pgfqpoint{2.175011in}{1.140154in}}%
\pgfpathlineto{\pgfqpoint{2.175194in}{1.144388in}}%
\pgfpathlineto{\pgfqpoint{2.169528in}{1.143936in}}%
\pgfpathlineto{\pgfqpoint{2.168048in}{1.161256in}}%
\pgfpathlineto{\pgfqpoint{2.177539in}{1.161835in}}%
\pgfpathlineto{\pgfqpoint{2.176581in}{1.173291in}}%
\pgfpathlineto{\pgfqpoint{2.182227in}{1.173751in}}%
\pgfpathlineto{\pgfqpoint{2.181588in}{1.182377in}}%
\pgfpathlineto{\pgfqpoint{2.194423in}{1.183326in}}%
\pgfpathlineto{\pgfqpoint{2.197194in}{1.183592in}}%
\pgfpathlineto{\pgfqpoint{2.198043in}{1.175099in}}%
\pgfpathlineto{\pgfqpoint{2.210701in}{1.176219in}}%
\pgfpathlineto{\pgfqpoint{2.213096in}{1.157168in}}%
\pgfpathlineto{\pgfqpoint{2.215737in}{1.157449in}}%
\pgfpathlineto{\pgfqpoint{2.218022in}{1.153845in}}%
\pgfpathlineto{\pgfqpoint{2.219718in}{1.138972in}}%
\pgfpathlineto{\pgfqpoint{2.218133in}{1.136791in}}%
\pgfpathlineto{\pgfqpoint{2.220327in}{1.114564in}}%
\pgfpathlineto{\pgfqpoint{2.201173in}{1.112813in}}%
\pgfpathclose%
\pgfusepath{fill}%
\end{pgfscope}%
\begin{pgfscope}%
\pgfpathrectangle{\pgfqpoint{0.100000in}{0.100000in}}{\pgfqpoint{3.007045in}{1.925000in}}%
\pgfusepath{clip}%
\pgfsetbuttcap%
\pgfsetmiterjoin%
\definecolor{currentfill}{rgb}{0.629296,0.850012,0.606136}%
\pgfsetfillcolor{currentfill}%
\pgfsetlinewidth{0.000000pt}%
\definecolor{currentstroke}{rgb}{0.000000,0.000000,0.000000}%
\pgfsetstrokecolor{currentstroke}%
\pgfsetstrokeopacity{0.000000}%
\pgfsetdash{}{0pt}%
\pgfpathmoveto{\pgfqpoint{1.899049in}{0.722750in}}%
\pgfpathlineto{\pgfqpoint{1.906378in}{0.722687in}}%
\pgfpathlineto{\pgfqpoint{1.910744in}{0.717166in}}%
\pgfpathlineto{\pgfqpoint{1.916505in}{0.717286in}}%
\pgfpathlineto{\pgfqpoint{1.916979in}{0.699909in}}%
\pgfpathlineto{\pgfqpoint{1.922946in}{0.694307in}}%
\pgfpathlineto{\pgfqpoint{1.923673in}{0.668219in}}%
\pgfpathlineto{\pgfqpoint{1.921110in}{0.659462in}}%
\pgfpathlineto{\pgfqpoint{1.906637in}{0.659096in}}%
\pgfpathlineto{\pgfqpoint{1.906746in}{0.653283in}}%
\pgfpathlineto{\pgfqpoint{1.886783in}{0.652756in}}%
\pgfpathlineto{\pgfqpoint{1.888359in}{0.662427in}}%
\pgfpathlineto{\pgfqpoint{1.891446in}{0.668175in}}%
\pgfpathlineto{\pgfqpoint{1.888095in}{0.681930in}}%
\pgfpathlineto{\pgfqpoint{1.895056in}{0.682024in}}%
\pgfpathlineto{\pgfqpoint{1.897220in}{0.687919in}}%
\pgfpathlineto{\pgfqpoint{1.896662in}{0.702398in}}%
\pgfpathlineto{\pgfqpoint{1.890952in}{0.702305in}}%
\pgfpathlineto{\pgfqpoint{1.890590in}{0.711973in}}%
\pgfpathlineto{\pgfqpoint{1.893306in}{0.722562in}}%
\pgfpathlineto{\pgfqpoint{1.899049in}{0.722750in}}%
\pgfpathclose%
\pgfusepath{fill}%
\end{pgfscope}%
\begin{pgfscope}%
\pgfpathrectangle{\pgfqpoint{0.100000in}{0.100000in}}{\pgfqpoint{3.007045in}{1.925000in}}%
\pgfusepath{clip}%
\pgfsetbuttcap%
\pgfsetmiterjoin%
\definecolor{currentfill}{rgb}{0.585006,0.829343,0.569719}%
\pgfsetfillcolor{currentfill}%
\pgfsetlinewidth{0.000000pt}%
\definecolor{currentstroke}{rgb}{0.000000,0.000000,0.000000}%
\pgfsetstrokecolor{currentstroke}%
\pgfsetstrokeopacity{0.000000}%
\pgfsetdash{}{0pt}%
\pgfpathmoveto{\pgfqpoint{1.092509in}{1.504079in}}%
\pgfpathlineto{\pgfqpoint{1.098004in}{1.495582in}}%
\pgfpathlineto{\pgfqpoint{1.097179in}{1.490393in}}%
\pgfpathlineto{\pgfqpoint{1.108398in}{1.488604in}}%
\pgfpathlineto{\pgfqpoint{1.119290in}{1.484034in}}%
\pgfpathlineto{\pgfqpoint{1.118840in}{1.481195in}}%
\pgfpathlineto{\pgfqpoint{1.135060in}{1.473922in}}%
\pgfpathlineto{\pgfqpoint{1.143778in}{1.473368in}}%
\pgfpathlineto{\pgfqpoint{1.170426in}{1.469676in}}%
\pgfpathlineto{\pgfqpoint{1.173551in}{1.471182in}}%
\pgfpathlineto{\pgfqpoint{1.174139in}{1.469186in}}%
\pgfpathlineto{\pgfqpoint{1.171716in}{1.446544in}}%
\pgfpathlineto{\pgfqpoint{1.168621in}{1.423872in}}%
\pgfpathlineto{\pgfqpoint{1.166601in}{1.424158in}}%
\pgfpathlineto{\pgfqpoint{1.163522in}{1.401453in}}%
\pgfpathlineto{\pgfqpoint{1.164443in}{1.401325in}}%
\pgfpathlineto{\pgfqpoint{1.162889in}{1.390045in}}%
\pgfpathlineto{\pgfqpoint{1.139789in}{1.393330in}}%
\pgfpathlineto{\pgfqpoint{1.089437in}{1.400986in}}%
\pgfpathlineto{\pgfqpoint{1.089929in}{1.412247in}}%
\pgfpathlineto{\pgfqpoint{1.092414in}{1.429227in}}%
\pgfpathlineto{\pgfqpoint{1.085388in}{1.435891in}}%
\pgfpathlineto{\pgfqpoint{1.079954in}{1.447589in}}%
\pgfpathlineto{\pgfqpoint{1.075159in}{1.453148in}}%
\pgfpathlineto{\pgfqpoint{1.070708in}{1.461628in}}%
\pgfpathlineto{\pgfqpoint{1.069165in}{1.468405in}}%
\pgfpathlineto{\pgfqpoint{1.069796in}{1.477757in}}%
\pgfpathlineto{\pgfqpoint{1.067969in}{1.484605in}}%
\pgfpathlineto{\pgfqpoint{1.053944in}{1.487032in}}%
\pgfpathlineto{\pgfqpoint{1.059788in}{1.522455in}}%
\pgfpathlineto{\pgfqpoint{1.062108in}{1.518116in}}%
\pgfpathlineto{\pgfqpoint{1.067064in}{1.517262in}}%
\pgfpathlineto{\pgfqpoint{1.070042in}{1.506895in}}%
\pgfpathlineto{\pgfqpoint{1.076162in}{1.510378in}}%
\pgfpathlineto{\pgfqpoint{1.078184in}{1.513711in}}%
\pgfpathlineto{\pgfqpoint{1.083266in}{1.515471in}}%
\pgfpathlineto{\pgfqpoint{1.087640in}{1.512697in}}%
\pgfpathlineto{\pgfqpoint{1.092509in}{1.504079in}}%
\pgfpathclose%
\pgfusepath{fill}%
\end{pgfscope}%
\begin{pgfscope}%
\pgfpathrectangle{\pgfqpoint{0.100000in}{0.100000in}}{\pgfqpoint{3.007045in}{1.925000in}}%
\pgfusepath{clip}%
\pgfsetbuttcap%
\pgfsetmiterjoin%
\definecolor{currentfill}{rgb}{0.164014,0.573641,0.291995}%
\pgfsetfillcolor{currentfill}%
\pgfsetlinewidth{0.000000pt}%
\definecolor{currentstroke}{rgb}{0.000000,0.000000,0.000000}%
\pgfsetstrokecolor{currentstroke}%
\pgfsetstrokeopacity{0.000000}%
\pgfsetdash{}{0pt}%
\pgfpathmoveto{\pgfqpoint{2.031351in}{1.567858in}}%
\pgfpathlineto{\pgfqpoint{2.030834in}{1.573678in}}%
\pgfpathlineto{\pgfqpoint{2.002328in}{1.571903in}}%
\pgfpathlineto{\pgfqpoint{2.000756in}{1.600279in}}%
\pgfpathlineto{\pgfqpoint{2.006057in}{1.600572in}}%
\pgfpathlineto{\pgfqpoint{2.004866in}{1.621636in}}%
\pgfpathlineto{\pgfqpoint{2.043917in}{1.613272in}}%
\pgfpathlineto{\pgfqpoint{2.048726in}{1.610846in}}%
\pgfpathlineto{\pgfqpoint{2.063792in}{1.604125in}}%
\pgfpathlineto{\pgfqpoint{2.064911in}{1.587042in}}%
\pgfpathlineto{\pgfqpoint{2.076489in}{1.587837in}}%
\pgfpathlineto{\pgfqpoint{2.078017in}{1.564986in}}%
\pgfpathlineto{\pgfqpoint{2.066412in}{1.564272in}}%
\pgfpathlineto{\pgfqpoint{2.054884in}{1.563482in}}%
\pgfpathlineto{\pgfqpoint{2.054639in}{1.569212in}}%
\pgfpathlineto{\pgfqpoint{2.031351in}{1.567858in}}%
\pgfpathclose%
\pgfusepath{fill}%
\end{pgfscope}%
\begin{pgfscope}%
\pgfpathrectangle{\pgfqpoint{0.100000in}{0.100000in}}{\pgfqpoint{3.007045in}{1.925000in}}%
\pgfusepath{clip}%
\pgfsetbuttcap%
\pgfsetmiterjoin%
\definecolor{currentfill}{rgb}{0.831142,0.934025,0.807074}%
\pgfsetfillcolor{currentfill}%
\pgfsetlinewidth{0.000000pt}%
\definecolor{currentstroke}{rgb}{0.000000,0.000000,0.000000}%
\pgfsetstrokecolor{currentstroke}%
\pgfsetstrokeopacity{0.000000}%
\pgfsetdash{}{0pt}%
\pgfpathmoveto{\pgfqpoint{0.955057in}{1.434288in}}%
\pgfpathlineto{\pgfqpoint{0.952072in}{1.424129in}}%
\pgfpathlineto{\pgfqpoint{0.948249in}{1.426793in}}%
\pgfpathlineto{\pgfqpoint{0.941439in}{1.428069in}}%
\pgfpathlineto{\pgfqpoint{0.934955in}{1.436848in}}%
\pgfpathlineto{\pgfqpoint{0.931789in}{1.444541in}}%
\pgfpathlineto{\pgfqpoint{0.926367in}{1.445605in}}%
\pgfpathlineto{\pgfqpoint{0.925285in}{1.439963in}}%
\pgfpathlineto{\pgfqpoint{0.917819in}{1.441425in}}%
\pgfpathlineto{\pgfqpoint{0.916709in}{1.435786in}}%
\pgfpathlineto{\pgfqpoint{0.900100in}{1.438993in}}%
\pgfpathlineto{\pgfqpoint{0.903388in}{1.455868in}}%
\pgfpathlineto{\pgfqpoint{0.895172in}{1.457614in}}%
\pgfpathlineto{\pgfqpoint{0.892728in}{1.460530in}}%
\pgfpathlineto{\pgfqpoint{0.896143in}{1.463222in}}%
\pgfpathlineto{\pgfqpoint{0.894533in}{1.469396in}}%
\pgfpathlineto{\pgfqpoint{0.898811in}{1.492021in}}%
\pgfpathlineto{\pgfqpoint{0.909954in}{1.489763in}}%
\pgfpathlineto{\pgfqpoint{0.912230in}{1.501011in}}%
\pgfpathlineto{\pgfqpoint{0.895520in}{1.504390in}}%
\pgfpathlineto{\pgfqpoint{0.896614in}{1.509734in}}%
\pgfpathlineto{\pgfqpoint{0.883755in}{1.512410in}}%
\pgfpathlineto{\pgfqpoint{0.884911in}{1.517908in}}%
\pgfpathlineto{\pgfqpoint{0.882414in}{1.527693in}}%
\pgfpathlineto{\pgfqpoint{0.878844in}{1.527019in}}%
\pgfpathlineto{\pgfqpoint{0.881395in}{1.530161in}}%
\pgfpathlineto{\pgfqpoint{0.891068in}{1.533017in}}%
\pgfpathlineto{\pgfqpoint{0.897650in}{1.536577in}}%
\pgfpathlineto{\pgfqpoint{0.900875in}{1.550400in}}%
\pgfpathlineto{\pgfqpoint{0.903257in}{1.555054in}}%
\pgfpathlineto{\pgfqpoint{0.905554in}{1.566329in}}%
\pgfpathlineto{\pgfqpoint{0.933521in}{1.560492in}}%
\pgfpathlineto{\pgfqpoint{0.935411in}{1.569991in}}%
\pgfpathlineto{\pgfqpoint{0.938309in}{1.576703in}}%
\pgfpathlineto{\pgfqpoint{0.952519in}{1.573584in}}%
\pgfpathlineto{\pgfqpoint{0.956111in}{1.570703in}}%
\pgfpathlineto{\pgfqpoint{0.958428in}{1.575728in}}%
\pgfpathlineto{\pgfqpoint{0.962156in}{1.577626in}}%
\pgfpathlineto{\pgfqpoint{0.970525in}{1.573027in}}%
\pgfpathlineto{\pgfqpoint{0.981565in}{1.573774in}}%
\pgfpathlineto{\pgfqpoint{0.983095in}{1.569716in}}%
\pgfpathlineto{\pgfqpoint{0.989150in}{1.572014in}}%
\pgfpathlineto{\pgfqpoint{0.998004in}{1.571906in}}%
\pgfpathlineto{\pgfqpoint{1.001653in}{1.579735in}}%
\pgfpathlineto{\pgfqpoint{1.006564in}{1.581699in}}%
\pgfpathlineto{\pgfqpoint{1.009746in}{1.577792in}}%
\pgfpathlineto{\pgfqpoint{1.015015in}{1.563234in}}%
\pgfpathlineto{\pgfqpoint{1.018543in}{1.560865in}}%
\pgfpathlineto{\pgfqpoint{1.012932in}{1.528900in}}%
\pgfpathlineto{\pgfqpoint{1.005109in}{1.526977in}}%
\pgfpathlineto{\pgfqpoint{0.995747in}{1.527955in}}%
\pgfpathlineto{\pgfqpoint{0.992126in}{1.508384in}}%
\pgfpathlineto{\pgfqpoint{0.999466in}{1.507030in}}%
\pgfpathlineto{\pgfqpoint{1.000614in}{1.501315in}}%
\pgfpathlineto{\pgfqpoint{1.007407in}{1.497511in}}%
\pgfpathlineto{\pgfqpoint{1.005297in}{1.485446in}}%
\pgfpathlineto{\pgfqpoint{1.001890in}{1.466220in}}%
\pgfpathlineto{\pgfqpoint{0.953549in}{1.475280in}}%
\pgfpathlineto{\pgfqpoint{0.947898in}{1.467983in}}%
\pgfpathlineto{\pgfqpoint{0.946788in}{1.460996in}}%
\pgfpathlineto{\pgfqpoint{0.947546in}{1.454706in}}%
\pgfpathlineto{\pgfqpoint{0.954259in}{1.454699in}}%
\pgfpathlineto{\pgfqpoint{0.955530in}{1.447686in}}%
\pgfpathlineto{\pgfqpoint{0.955057in}{1.434288in}}%
\pgfpathclose%
\pgfusepath{fill}%
\end{pgfscope}%
\begin{pgfscope}%
\pgfpathrectangle{\pgfqpoint{0.100000in}{0.100000in}}{\pgfqpoint{3.007045in}{1.925000in}}%
\pgfusepath{clip}%
\pgfsetbuttcap%
\pgfsetmiterjoin%
\definecolor{currentfill}{rgb}{0.797924,0.920738,0.771642}%
\pgfsetfillcolor{currentfill}%
\pgfsetlinewidth{0.000000pt}%
\definecolor{currentstroke}{rgb}{0.000000,0.000000,0.000000}%
\pgfsetstrokecolor{currentstroke}%
\pgfsetstrokeopacity{0.000000}%
\pgfsetdash{}{0pt}%
\pgfpathmoveto{\pgfqpoint{0.615087in}{1.868865in}}%
\pgfpathlineto{\pgfqpoint{0.618659in}{1.862354in}}%
\pgfpathlineto{\pgfqpoint{0.620178in}{1.855997in}}%
\pgfpathlineto{\pgfqpoint{0.627702in}{1.851025in}}%
\pgfpathlineto{\pgfqpoint{0.629717in}{1.846964in}}%
\pgfpathlineto{\pgfqpoint{0.635771in}{1.843623in}}%
\pgfpathlineto{\pgfqpoint{0.640168in}{1.838234in}}%
\pgfpathlineto{\pgfqpoint{0.652989in}{1.834668in}}%
\pgfpathlineto{\pgfqpoint{0.656028in}{1.830896in}}%
\pgfpathlineto{\pgfqpoint{0.652061in}{1.822418in}}%
\pgfpathlineto{\pgfqpoint{0.652827in}{1.812664in}}%
\pgfpathlineto{\pgfqpoint{0.651860in}{1.802864in}}%
\pgfpathlineto{\pgfqpoint{0.648904in}{1.795115in}}%
\pgfpathlineto{\pgfqpoint{0.651182in}{1.791391in}}%
\pgfpathlineto{\pgfqpoint{0.640972in}{1.753699in}}%
\pgfpathlineto{\pgfqpoint{0.616044in}{1.760740in}}%
\pgfpathlineto{\pgfqpoint{0.567701in}{1.775221in}}%
\pgfpathlineto{\pgfqpoint{0.574289in}{1.797104in}}%
\pgfpathlineto{\pgfqpoint{0.579595in}{1.795500in}}%
\pgfpathlineto{\pgfqpoint{0.579759in}{1.803398in}}%
\pgfpathlineto{\pgfqpoint{0.585013in}{1.811079in}}%
\pgfpathlineto{\pgfqpoint{0.585552in}{1.821819in}}%
\pgfpathlineto{\pgfqpoint{0.583582in}{1.827974in}}%
\pgfpathlineto{\pgfqpoint{0.597616in}{1.841914in}}%
\pgfpathlineto{\pgfqpoint{0.599282in}{1.851636in}}%
\pgfpathlineto{\pgfqpoint{0.595548in}{1.853584in}}%
\pgfpathlineto{\pgfqpoint{0.595348in}{1.857817in}}%
\pgfpathlineto{\pgfqpoint{0.600903in}{1.863151in}}%
\pgfpathlineto{\pgfqpoint{0.610134in}{1.868414in}}%
\pgfpathlineto{\pgfqpoint{0.615087in}{1.868865in}}%
\pgfpathclose%
\pgfusepath{fill}%
\end{pgfscope}%
\begin{pgfscope}%
\pgfpathrectangle{\pgfqpoint{0.100000in}{0.100000in}}{\pgfqpoint{3.007045in}{1.925000in}}%
\pgfusepath{clip}%
\pgfsetbuttcap%
\pgfsetmiterjoin%
\definecolor{currentfill}{rgb}{0.681046,0.871895,0.656209}%
\pgfsetfillcolor{currentfill}%
\pgfsetlinewidth{0.000000pt}%
\definecolor{currentstroke}{rgb}{0.000000,0.000000,0.000000}%
\pgfsetstrokecolor{currentstroke}%
\pgfsetstrokeopacity{0.000000}%
\pgfsetdash{}{0pt}%
\pgfpathmoveto{\pgfqpoint{0.465761in}{1.347164in}}%
\pgfpathlineto{\pgfqpoint{0.460387in}{1.347970in}}%
\pgfpathlineto{\pgfqpoint{0.450179in}{1.340731in}}%
\pgfpathlineto{\pgfqpoint{0.445384in}{1.329930in}}%
\pgfpathlineto{\pgfqpoint{0.439260in}{1.328328in}}%
\pgfpathlineto{\pgfqpoint{0.435140in}{1.324290in}}%
\pgfpathlineto{\pgfqpoint{0.431828in}{1.313269in}}%
\pgfpathlineto{\pgfqpoint{0.424973in}{1.314267in}}%
\pgfpathlineto{\pgfqpoint{0.420345in}{1.320339in}}%
\pgfpathlineto{\pgfqpoint{0.421047in}{1.323928in}}%
\pgfpathlineto{\pgfqpoint{0.417135in}{1.329622in}}%
\pgfpathlineto{\pgfqpoint{0.392426in}{1.337198in}}%
\pgfpathlineto{\pgfqpoint{0.390745in}{1.343728in}}%
\pgfpathlineto{\pgfqpoint{0.387533in}{1.347636in}}%
\pgfpathlineto{\pgfqpoint{0.389453in}{1.357782in}}%
\pgfpathlineto{\pgfqpoint{0.383879in}{1.360469in}}%
\pgfpathlineto{\pgfqpoint{0.379169in}{1.366194in}}%
\pgfpathlineto{\pgfqpoint{0.383653in}{1.376276in}}%
\pgfpathlineto{\pgfqpoint{0.386077in}{1.384484in}}%
\pgfpathlineto{\pgfqpoint{0.378819in}{1.386738in}}%
\pgfpathlineto{\pgfqpoint{0.381027in}{1.394938in}}%
\pgfpathlineto{\pgfqpoint{0.380635in}{1.401231in}}%
\pgfpathlineto{\pgfqpoint{0.423677in}{1.387762in}}%
\pgfpathlineto{\pgfqpoint{0.425415in}{1.393166in}}%
\pgfpathlineto{\pgfqpoint{0.437071in}{1.389645in}}%
\pgfpathlineto{\pgfqpoint{0.443854in}{1.394476in}}%
\pgfpathlineto{\pgfqpoint{0.446584in}{1.393519in}}%
\pgfpathlineto{\pgfqpoint{0.451708in}{1.400095in}}%
\pgfpathlineto{\pgfqpoint{0.459905in}{1.401168in}}%
\pgfpathlineto{\pgfqpoint{0.461995in}{1.397568in}}%
\pgfpathlineto{\pgfqpoint{0.458164in}{1.392328in}}%
\pgfpathlineto{\pgfqpoint{0.455804in}{1.383109in}}%
\pgfpathlineto{\pgfqpoint{0.457960in}{1.381368in}}%
\pgfpathlineto{\pgfqpoint{0.467458in}{1.355133in}}%
\pgfpathlineto{\pgfqpoint{0.465761in}{1.347164in}}%
\pgfpathclose%
\pgfusepath{fill}%
\end{pgfscope}%
\begin{pgfscope}%
\pgfpathrectangle{\pgfqpoint{0.100000in}{0.100000in}}{\pgfqpoint{3.007045in}{1.925000in}}%
\pgfusepath{clip}%
\pgfsetbuttcap%
\pgfsetmiterjoin%
\definecolor{currentfill}{rgb}{0.507497,0.793172,0.505990}%
\pgfsetfillcolor{currentfill}%
\pgfsetlinewidth{0.000000pt}%
\definecolor{currentstroke}{rgb}{0.000000,0.000000,0.000000}%
\pgfsetstrokecolor{currentstroke}%
\pgfsetstrokeopacity{0.000000}%
\pgfsetdash{}{0pt}%
\pgfpathmoveto{\pgfqpoint{1.746105in}{1.095854in}}%
\pgfpathlineto{\pgfqpoint{1.745387in}{1.069764in}}%
\pgfpathlineto{\pgfqpoint{1.722556in}{1.070090in}}%
\pgfpathlineto{\pgfqpoint{1.722658in}{1.078703in}}%
\pgfpathlineto{\pgfqpoint{1.702042in}{1.079070in}}%
\pgfpathlineto{\pgfqpoint{1.701914in}{1.073348in}}%
\pgfpathlineto{\pgfqpoint{1.693377in}{1.073520in}}%
\pgfpathlineto{\pgfqpoint{1.676880in}{1.073769in}}%
\pgfpathlineto{\pgfqpoint{1.677388in}{1.091095in}}%
\pgfpathlineto{\pgfqpoint{1.678559in}{1.102524in}}%
\pgfpathlineto{\pgfqpoint{1.702564in}{1.102052in}}%
\pgfpathlineto{\pgfqpoint{1.702839in}{1.116383in}}%
\pgfpathlineto{\pgfqpoint{1.723731in}{1.116097in}}%
\pgfpathlineto{\pgfqpoint{1.723295in}{1.096039in}}%
\pgfpathlineto{\pgfqpoint{1.746105in}{1.095854in}}%
\pgfpathclose%
\pgfusepath{fill}%
\end{pgfscope}%
\begin{pgfscope}%
\pgfpathrectangle{\pgfqpoint{0.100000in}{0.100000in}}{\pgfqpoint{3.007045in}{1.925000in}}%
\pgfusepath{clip}%
\pgfsetbuttcap%
\pgfsetmiterjoin%
\definecolor{currentfill}{rgb}{0.524106,0.800923,0.519646}%
\pgfsetfillcolor{currentfill}%
\pgfsetlinewidth{0.000000pt}%
\definecolor{currentstroke}{rgb}{0.000000,0.000000,0.000000}%
\pgfsetstrokecolor{currentstroke}%
\pgfsetstrokeopacity{0.000000}%
\pgfsetdash{}{0pt}%
\pgfpathmoveto{\pgfqpoint{1.102971in}{1.695747in}}%
\pgfpathlineto{\pgfqpoint{1.099659in}{1.696257in}}%
\pgfpathlineto{\pgfqpoint{1.093536in}{1.701272in}}%
\pgfpathlineto{\pgfqpoint{1.070809in}{1.700125in}}%
\pgfpathlineto{\pgfqpoint{1.066660in}{1.704559in}}%
\pgfpathlineto{\pgfqpoint{1.059670in}{1.716732in}}%
\pgfpathlineto{\pgfqpoint{1.053818in}{1.723481in}}%
\pgfpathlineto{\pgfqpoint{1.043222in}{1.727442in}}%
\pgfpathlineto{\pgfqpoint{1.026102in}{1.729804in}}%
\pgfpathlineto{\pgfqpoint{1.025046in}{1.724117in}}%
\pgfpathlineto{\pgfqpoint{1.019738in}{1.725093in}}%
\pgfpathlineto{\pgfqpoint{1.013851in}{1.726227in}}%
\pgfpathlineto{\pgfqpoint{1.016513in}{1.740285in}}%
\pgfpathlineto{\pgfqpoint{1.010013in}{1.742368in}}%
\pgfpathlineto{\pgfqpoint{1.005933in}{1.746597in}}%
\pgfpathlineto{\pgfqpoint{1.009841in}{1.767764in}}%
\pgfpathlineto{\pgfqpoint{0.999503in}{1.769057in}}%
\pgfpathlineto{\pgfqpoint{0.994442in}{1.773321in}}%
\pgfpathlineto{\pgfqpoint{0.994400in}{1.776622in}}%
\pgfpathlineto{\pgfqpoint{0.987897in}{1.779953in}}%
\pgfpathlineto{\pgfqpoint{0.976971in}{1.780376in}}%
\pgfpathlineto{\pgfqpoint{0.973893in}{1.786912in}}%
\pgfpathlineto{\pgfqpoint{0.975711in}{1.795488in}}%
\pgfpathlineto{\pgfqpoint{0.974188in}{1.804466in}}%
\pgfpathlineto{\pgfqpoint{0.978555in}{1.808720in}}%
\pgfpathlineto{\pgfqpoint{0.975124in}{1.816423in}}%
\pgfpathlineto{\pgfqpoint{1.011784in}{1.809131in}}%
\pgfpathlineto{\pgfqpoint{1.016942in}{1.805156in}}%
\pgfpathlineto{\pgfqpoint{1.018457in}{1.797935in}}%
\pgfpathlineto{\pgfqpoint{1.043735in}{1.793392in}}%
\pgfpathlineto{\pgfqpoint{1.045450in}{1.802919in}}%
\pgfpathlineto{\pgfqpoint{1.070974in}{1.798414in}}%
\pgfpathlineto{\pgfqpoint{1.071976in}{1.804099in}}%
\pgfpathlineto{\pgfqpoint{1.080804in}{1.802553in}}%
\pgfpathlineto{\pgfqpoint{1.081806in}{1.808264in}}%
\pgfpathlineto{\pgfqpoint{1.115936in}{1.802454in}}%
\pgfpathlineto{\pgfqpoint{1.115258in}{1.796858in}}%
\pgfpathlineto{\pgfqpoint{1.120232in}{1.796028in}}%
\pgfpathlineto{\pgfqpoint{1.119270in}{1.790180in}}%
\pgfpathlineto{\pgfqpoint{1.127771in}{1.788966in}}%
\pgfpathlineto{\pgfqpoint{1.123379in}{1.762908in}}%
\pgfpathlineto{\pgfqpoint{1.110393in}{1.763495in}}%
\pgfpathlineto{\pgfqpoint{1.107591in}{1.752201in}}%
\pgfpathlineto{\pgfqpoint{1.094250in}{1.746968in}}%
\pgfpathlineto{\pgfqpoint{1.092154in}{1.736531in}}%
\pgfpathlineto{\pgfqpoint{1.097369in}{1.732746in}}%
\pgfpathlineto{\pgfqpoint{1.103070in}{1.731808in}}%
\pgfpathlineto{\pgfqpoint{1.107755in}{1.728110in}}%
\pgfpathlineto{\pgfqpoint{1.104710in}{1.709101in}}%
\pgfpathlineto{\pgfqpoint{1.102514in}{1.707496in}}%
\pgfpathlineto{\pgfqpoint{1.102971in}{1.695747in}}%
\pgfpathclose%
\pgfusepath{fill}%
\end{pgfscope}%
\begin{pgfscope}%
\pgfpathrectangle{\pgfqpoint{0.100000in}{0.100000in}}{\pgfqpoint{3.007045in}{1.925000in}}%
\pgfusepath{clip}%
\pgfsetbuttcap%
\pgfsetmiterjoin%
\definecolor{currentfill}{rgb}{0.457670,0.769919,0.465021}%
\pgfsetfillcolor{currentfill}%
\pgfsetlinewidth{0.000000pt}%
\definecolor{currentstroke}{rgb}{0.000000,0.000000,0.000000}%
\pgfsetstrokecolor{currentstroke}%
\pgfsetstrokeopacity{0.000000}%
\pgfsetdash{}{0pt}%
\pgfpathmoveto{\pgfqpoint{1.522557in}{1.725806in}}%
\pgfpathlineto{\pgfqpoint{1.539856in}{1.724787in}}%
\pgfpathlineto{\pgfqpoint{1.538782in}{1.736495in}}%
\pgfpathlineto{\pgfqpoint{1.540108in}{1.759776in}}%
\pgfpathlineto{\pgfqpoint{1.539033in}{1.771449in}}%
\pgfpathlineto{\pgfqpoint{1.544810in}{1.771147in}}%
\pgfpathlineto{\pgfqpoint{1.573646in}{1.769590in}}%
\pgfpathlineto{\pgfqpoint{1.573086in}{1.757987in}}%
\pgfpathlineto{\pgfqpoint{1.557416in}{1.758816in}}%
\pgfpathlineto{\pgfqpoint{1.556105in}{1.735501in}}%
\pgfpathlineto{\pgfqpoint{1.557103in}{1.723835in}}%
\pgfpathlineto{\pgfqpoint{1.580095in}{1.722634in}}%
\pgfpathlineto{\pgfqpoint{1.580956in}{1.710973in}}%
\pgfpathlineto{\pgfqpoint{1.579862in}{1.687980in}}%
\pgfpathlineto{\pgfqpoint{1.570193in}{1.688446in}}%
\pgfpathlineto{\pgfqpoint{1.541684in}{1.689945in}}%
\pgfpathlineto{\pgfqpoint{1.516455in}{1.691466in}}%
\pgfpathlineto{\pgfqpoint{1.517941in}{1.714507in}}%
\pgfpathlineto{\pgfqpoint{1.521849in}{1.714257in}}%
\pgfpathlineto{\pgfqpoint{1.522557in}{1.725806in}}%
\pgfpathclose%
\pgfusepath{fill}%
\end{pgfscope}%
\begin{pgfscope}%
\pgfpathrectangle{\pgfqpoint{0.100000in}{0.100000in}}{\pgfqpoint{3.007045in}{1.925000in}}%
\pgfusepath{clip}%
\pgfsetbuttcap%
\pgfsetmiterjoin%
\definecolor{currentfill}{rgb}{0.535179,0.806090,0.528750}%
\pgfsetfillcolor{currentfill}%
\pgfsetlinewidth{0.000000pt}%
\definecolor{currentstroke}{rgb}{0.000000,0.000000,0.000000}%
\pgfsetstrokecolor{currentstroke}%
\pgfsetstrokeopacity{0.000000}%
\pgfsetdash{}{0pt}%
\pgfpathmoveto{\pgfqpoint{2.024381in}{1.212989in}}%
\pgfpathlineto{\pgfqpoint{2.014487in}{1.208775in}}%
\pgfpathlineto{\pgfqpoint{2.009474in}{1.209217in}}%
\pgfpathlineto{\pgfqpoint{2.008276in}{1.212623in}}%
\pgfpathlineto{\pgfqpoint{2.013813in}{1.217404in}}%
\pgfpathlineto{\pgfqpoint{2.017601in}{1.220842in}}%
\pgfpathlineto{\pgfqpoint{2.021248in}{1.230411in}}%
\pgfpathlineto{\pgfqpoint{2.027473in}{1.235391in}}%
\pgfpathlineto{\pgfqpoint{2.029027in}{1.239956in}}%
\pgfpathlineto{\pgfqpoint{2.028561in}{1.247285in}}%
\pgfpathlineto{\pgfqpoint{2.022785in}{1.247040in}}%
\pgfpathlineto{\pgfqpoint{2.021691in}{1.270010in}}%
\pgfpathlineto{\pgfqpoint{2.038925in}{1.270934in}}%
\pgfpathlineto{\pgfqpoint{2.037908in}{1.288046in}}%
\pgfpathlineto{\pgfqpoint{2.027069in}{1.287453in}}%
\pgfpathlineto{\pgfqpoint{2.025542in}{1.310434in}}%
\pgfpathlineto{\pgfqpoint{2.059792in}{1.312587in}}%
\pgfpathlineto{\pgfqpoint{2.059573in}{1.315459in}}%
\pgfpathlineto{\pgfqpoint{2.087331in}{1.317506in}}%
\pgfpathlineto{\pgfqpoint{2.090562in}{1.283138in}}%
\pgfpathlineto{\pgfqpoint{2.073487in}{1.281826in}}%
\pgfpathlineto{\pgfqpoint{2.075117in}{1.258653in}}%
\pgfpathlineto{\pgfqpoint{2.072749in}{1.252662in}}%
\pgfpathlineto{\pgfqpoint{2.065605in}{1.247690in}}%
\pgfpathlineto{\pgfqpoint{2.058847in}{1.247095in}}%
\pgfpathlineto{\pgfqpoint{2.060239in}{1.229418in}}%
\pgfpathlineto{\pgfqpoint{2.043262in}{1.228036in}}%
\pgfpathlineto{\pgfqpoint{2.044100in}{1.215046in}}%
\pgfpathlineto{\pgfqpoint{2.039739in}{1.216849in}}%
\pgfpathlineto{\pgfqpoint{2.032407in}{1.214497in}}%
\pgfpathlineto{\pgfqpoint{2.025473in}{1.215274in}}%
\pgfpathlineto{\pgfqpoint{2.024381in}{1.212989in}}%
\pgfpathclose%
\pgfusepath{fill}%
\end{pgfscope}%
\begin{pgfscope}%
\pgfpathrectangle{\pgfqpoint{0.100000in}{0.100000in}}{\pgfqpoint{3.007045in}{1.925000in}}%
\pgfusepath{clip}%
\pgfsetbuttcap%
\pgfsetmiterjoin%
\definecolor{currentfill}{rgb}{0.596078,0.834510,0.578824}%
\pgfsetfillcolor{currentfill}%
\pgfsetlinewidth{0.000000pt}%
\definecolor{currentstroke}{rgb}{0.000000,0.000000,0.000000}%
\pgfsetstrokecolor{currentstroke}%
\pgfsetstrokeopacity{0.000000}%
\pgfsetdash{}{0pt}%
\pgfpathmoveto{\pgfqpoint{1.553171in}{1.527805in}}%
\pgfpathlineto{\pgfqpoint{1.557925in}{1.527548in}}%
\pgfpathlineto{\pgfqpoint{1.598571in}{1.525674in}}%
\pgfpathlineto{\pgfqpoint{1.597846in}{1.508282in}}%
\pgfpathlineto{\pgfqpoint{1.596903in}{1.479359in}}%
\pgfpathlineto{\pgfqpoint{1.586215in}{1.479799in}}%
\pgfpathlineto{\pgfqpoint{1.555401in}{1.480969in}}%
\pgfpathlineto{\pgfqpoint{1.551371in}{1.482830in}}%
\pgfpathlineto{\pgfqpoint{1.551887in}{1.504699in}}%
\pgfpathlineto{\pgfqpoint{1.553171in}{1.527805in}}%
\pgfpathclose%
\pgfusepath{fill}%
\end{pgfscope}%
\begin{pgfscope}%
\pgfpathrectangle{\pgfqpoint{0.100000in}{0.100000in}}{\pgfqpoint{3.007045in}{1.925000in}}%
\pgfusepath{clip}%
\pgfsetbuttcap%
\pgfsetmiterjoin%
\definecolor{currentfill}{rgb}{0.607151,0.839677,0.587928}%
\pgfsetfillcolor{currentfill}%
\pgfsetlinewidth{0.000000pt}%
\definecolor{currentstroke}{rgb}{0.000000,0.000000,0.000000}%
\pgfsetstrokecolor{currentstroke}%
\pgfsetstrokeopacity{0.000000}%
\pgfsetdash{}{0pt}%
\pgfpathmoveto{\pgfqpoint{1.963991in}{0.607879in}}%
\pgfpathlineto{\pgfqpoint{1.967968in}{0.612514in}}%
\pgfpathlineto{\pgfqpoint{1.964384in}{0.616831in}}%
\pgfpathlineto{\pgfqpoint{1.967749in}{0.623734in}}%
\pgfpathlineto{\pgfqpoint{1.971855in}{0.623291in}}%
\pgfpathlineto{\pgfqpoint{1.972763in}{0.629511in}}%
\pgfpathlineto{\pgfqpoint{1.985698in}{0.633433in}}%
\pgfpathlineto{\pgfqpoint{1.993008in}{0.630363in}}%
\pgfpathlineto{\pgfqpoint{1.998230in}{0.632422in}}%
\pgfpathlineto{\pgfqpoint{2.020031in}{0.633536in}}%
\pgfpathlineto{\pgfqpoint{2.049764in}{0.635164in}}%
\pgfpathlineto{\pgfqpoint{2.053466in}{0.638280in}}%
\pgfpathlineto{\pgfqpoint{2.058100in}{0.638257in}}%
\pgfpathlineto{\pgfqpoint{2.065780in}{0.638658in}}%
\pgfpathlineto{\pgfqpoint{2.065550in}{0.641578in}}%
\pgfpathlineto{\pgfqpoint{2.075185in}{0.642117in}}%
\pgfpathlineto{\pgfqpoint{2.076881in}{0.613497in}}%
\pgfpathlineto{\pgfqpoint{2.072666in}{0.613237in}}%
\pgfpathlineto{\pgfqpoint{2.037435in}{0.611183in}}%
\pgfpathlineto{\pgfqpoint{2.010177in}{0.609791in}}%
\pgfpathlineto{\pgfqpoint{1.963991in}{0.607879in}}%
\pgfpathclose%
\pgfusepath{fill}%
\end{pgfscope}%
\begin{pgfscope}%
\pgfpathrectangle{\pgfqpoint{0.100000in}{0.100000in}}{\pgfqpoint{3.007045in}{1.925000in}}%
\pgfusepath{clip}%
\pgfsetbuttcap%
\pgfsetmiterjoin%
\definecolor{currentfill}{rgb}{0.389020,0.736332,0.430450}%
\pgfsetfillcolor{currentfill}%
\pgfsetlinewidth{0.000000pt}%
\definecolor{currentstroke}{rgb}{0.000000,0.000000,0.000000}%
\pgfsetstrokecolor{currentstroke}%
\pgfsetstrokeopacity{0.000000}%
\pgfsetdash{}{0pt}%
\pgfpathmoveto{\pgfqpoint{1.867818in}{1.391658in}}%
\pgfpathlineto{\pgfqpoint{1.866993in}{1.430735in}}%
\pgfpathlineto{\pgfqpoint{1.894560in}{1.431480in}}%
\pgfpathlineto{\pgfqpoint{1.928977in}{1.432597in}}%
\pgfpathlineto{\pgfqpoint{1.934699in}{1.432796in}}%
\pgfpathlineto{\pgfqpoint{1.936000in}{1.405159in}}%
\pgfpathlineto{\pgfqpoint{1.913064in}{1.404438in}}%
\pgfpathlineto{\pgfqpoint{1.913408in}{1.392837in}}%
\pgfpathlineto{\pgfqpoint{1.867818in}{1.391658in}}%
\pgfpathclose%
\pgfusepath{fill}%
\end{pgfscope}%
\begin{pgfscope}%
\pgfpathrectangle{\pgfqpoint{0.100000in}{0.100000in}}{\pgfqpoint{3.007045in}{1.925000in}}%
\pgfusepath{clip}%
\pgfsetbuttcap%
\pgfsetmiterjoin%
\definecolor{currentfill}{rgb}{0.557324,0.816424,0.546959}%
\pgfsetfillcolor{currentfill}%
\pgfsetlinewidth{0.000000pt}%
\definecolor{currentstroke}{rgb}{0.000000,0.000000,0.000000}%
\pgfsetstrokecolor{currentstroke}%
\pgfsetstrokeopacity{0.000000}%
\pgfsetdash{}{0pt}%
\pgfpathmoveto{\pgfqpoint{2.608126in}{1.441492in}}%
\pgfpathlineto{\pgfqpoint{2.601922in}{1.443430in}}%
\pgfpathlineto{\pgfqpoint{2.572079in}{1.437927in}}%
\pgfpathlineto{\pgfqpoint{2.570755in}{1.438725in}}%
\pgfpathlineto{\pgfqpoint{2.557784in}{1.431044in}}%
\pgfpathlineto{\pgfqpoint{2.552432in}{1.429022in}}%
\pgfpathlineto{\pgfqpoint{2.546681in}{1.434064in}}%
\pgfpathlineto{\pgfqpoint{2.543334in}{1.433990in}}%
\pgfpathlineto{\pgfqpoint{2.539326in}{1.435418in}}%
\pgfpathlineto{\pgfqpoint{2.542194in}{1.444103in}}%
\pgfpathlineto{\pgfqpoint{2.547820in}{1.448105in}}%
\pgfpathlineto{\pgfqpoint{2.550364in}{1.452200in}}%
\pgfpathlineto{\pgfqpoint{2.544628in}{1.462243in}}%
\pgfpathlineto{\pgfqpoint{2.540521in}{1.463483in}}%
\pgfpathlineto{\pgfqpoint{2.539800in}{1.468753in}}%
\pgfpathlineto{\pgfqpoint{2.535305in}{1.480591in}}%
\pgfpathlineto{\pgfqpoint{2.544988in}{1.486580in}}%
\pgfpathlineto{\pgfqpoint{2.553954in}{1.490859in}}%
\pgfpathlineto{\pgfqpoint{2.565267in}{1.494184in}}%
\pgfpathlineto{\pgfqpoint{2.584950in}{1.497319in}}%
\pgfpathlineto{\pgfqpoint{2.595524in}{1.497454in}}%
\pgfpathlineto{\pgfqpoint{2.605383in}{1.492846in}}%
\pgfpathlineto{\pgfqpoint{2.614787in}{1.497022in}}%
\pgfpathlineto{\pgfqpoint{2.634861in}{1.500680in}}%
\pgfpathlineto{\pgfqpoint{2.644415in}{1.507385in}}%
\pgfpathlineto{\pgfqpoint{2.648812in}{1.490538in}}%
\pgfpathlineto{\pgfqpoint{2.648634in}{1.483094in}}%
\pgfpathlineto{\pgfqpoint{2.652395in}{1.466521in}}%
\pgfpathlineto{\pgfqpoint{2.656543in}{1.461343in}}%
\pgfpathlineto{\pgfqpoint{2.661352in}{1.457374in}}%
\pgfpathlineto{\pgfqpoint{2.646021in}{1.453686in}}%
\pgfpathlineto{\pgfqpoint{2.647379in}{1.448629in}}%
\pgfpathlineto{\pgfqpoint{2.635144in}{1.448720in}}%
\pgfpathlineto{\pgfqpoint{2.634427in}{1.453744in}}%
\pgfpathlineto{\pgfqpoint{2.609773in}{1.449142in}}%
\pgfpathlineto{\pgfqpoint{2.608126in}{1.441492in}}%
\pgfpathclose%
\pgfusepath{fill}%
\end{pgfscope}%
\begin{pgfscope}%
\pgfpathrectangle{\pgfqpoint{0.100000in}{0.100000in}}{\pgfqpoint{3.007045in}{1.925000in}}%
\pgfusepath{clip}%
\pgfsetbuttcap%
\pgfsetmiterjoin%
\definecolor{currentfill}{rgb}{0.457670,0.769919,0.465021}%
\pgfsetfillcolor{currentfill}%
\pgfsetlinewidth{0.000000pt}%
\definecolor{currentstroke}{rgb}{0.000000,0.000000,0.000000}%
\pgfsetstrokecolor{currentstroke}%
\pgfsetstrokeopacity{0.000000}%
\pgfsetdash{}{0pt}%
\pgfpathmoveto{\pgfqpoint{2.507767in}{1.239655in}}%
\pgfpathlineto{\pgfqpoint{2.502865in}{1.238879in}}%
\pgfpathlineto{\pgfqpoint{2.498296in}{1.267386in}}%
\pgfpathlineto{\pgfqpoint{2.493292in}{1.298554in}}%
\pgfpathlineto{\pgfqpoint{2.491069in}{1.312382in}}%
\pgfpathlineto{\pgfqpoint{2.508763in}{1.315516in}}%
\pgfpathlineto{\pgfqpoint{2.506874in}{1.324911in}}%
\pgfpathlineto{\pgfqpoint{2.513140in}{1.337398in}}%
\pgfpathlineto{\pgfqpoint{2.528039in}{1.339978in}}%
\pgfpathlineto{\pgfqpoint{2.533624in}{1.328350in}}%
\pgfpathlineto{\pgfqpoint{2.540085in}{1.329055in}}%
\pgfpathlineto{\pgfqpoint{2.550029in}{1.332847in}}%
\pgfpathlineto{\pgfqpoint{2.552895in}{1.336088in}}%
\pgfpathlineto{\pgfqpoint{2.554469in}{1.327047in}}%
\pgfpathlineto{\pgfqpoint{2.574642in}{1.330275in}}%
\pgfpathlineto{\pgfqpoint{2.576932in}{1.318538in}}%
\pgfpathlineto{\pgfqpoint{2.574047in}{1.302973in}}%
\pgfpathlineto{\pgfqpoint{2.567920in}{1.283023in}}%
\pgfpathlineto{\pgfqpoint{2.561814in}{1.273283in}}%
\pgfpathlineto{\pgfqpoint{2.556847in}{1.259865in}}%
\pgfpathlineto{\pgfqpoint{2.559394in}{1.254953in}}%
\pgfpathlineto{\pgfqpoint{2.559174in}{1.248298in}}%
\pgfpathlineto{\pgfqpoint{2.554976in}{1.247539in}}%
\pgfpathlineto{\pgfqpoint{2.507767in}{1.239655in}}%
\pgfpathclose%
\pgfusepath{fill}%
\end{pgfscope}%
\begin{pgfscope}%
\pgfpathrectangle{\pgfqpoint{0.100000in}{0.100000in}}{\pgfqpoint{3.007045in}{1.925000in}}%
\pgfusepath{clip}%
\pgfsetbuttcap%
\pgfsetmiterjoin%
\definecolor{currentfill}{rgb}{0.200923,0.613010,0.321522}%
\pgfsetfillcolor{currentfill}%
\pgfsetlinewidth{0.000000pt}%
\definecolor{currentstroke}{rgb}{0.000000,0.000000,0.000000}%
\pgfsetstrokecolor{currentstroke}%
\pgfsetstrokeopacity{0.000000}%
\pgfsetdash{}{0pt}%
\pgfpathmoveto{\pgfqpoint{1.506378in}{0.538351in}}%
\pgfpathlineto{\pgfqpoint{1.501272in}{0.538689in}}%
\pgfpathlineto{\pgfqpoint{1.501924in}{0.549696in}}%
\pgfpathlineto{\pgfqpoint{1.498020in}{0.549906in}}%
\pgfpathlineto{\pgfqpoint{1.498977in}{0.564304in}}%
\pgfpathlineto{\pgfqpoint{1.524943in}{0.562745in}}%
\pgfpathlineto{\pgfqpoint{1.525500in}{0.576932in}}%
\pgfpathlineto{\pgfqpoint{1.544968in}{0.575956in}}%
\pgfpathlineto{\pgfqpoint{1.566327in}{0.575160in}}%
\pgfpathlineto{\pgfqpoint{1.565594in}{0.551126in}}%
\pgfpathlineto{\pgfqpoint{1.546459in}{0.551872in}}%
\pgfpathlineto{\pgfqpoint{1.545579in}{0.528102in}}%
\pgfpathlineto{\pgfqpoint{1.530916in}{0.536323in}}%
\pgfpathlineto{\pgfqpoint{1.506378in}{0.538351in}}%
\pgfpathclose%
\pgfusepath{fill}%
\end{pgfscope}%
\begin{pgfscope}%
\pgfpathrectangle{\pgfqpoint{0.100000in}{0.100000in}}{\pgfqpoint{3.007045in}{1.925000in}}%
\pgfusepath{clip}%
\pgfsetbuttcap%
\pgfsetmiterjoin%
\definecolor{currentfill}{rgb}{0.513033,0.795755,0.510542}%
\pgfsetfillcolor{currentfill}%
\pgfsetlinewidth{0.000000pt}%
\definecolor{currentstroke}{rgb}{0.000000,0.000000,0.000000}%
\pgfsetstrokecolor{currentstroke}%
\pgfsetstrokeopacity{0.000000}%
\pgfsetdash{}{0pt}%
\pgfpathmoveto{\pgfqpoint{2.408663in}{0.823968in}}%
\pgfpathlineto{\pgfqpoint{2.414961in}{0.829005in}}%
\pgfpathlineto{\pgfqpoint{2.424340in}{0.826019in}}%
\pgfpathlineto{\pgfqpoint{2.430059in}{0.826283in}}%
\pgfpathlineto{\pgfqpoint{2.430794in}{0.832351in}}%
\pgfpathlineto{\pgfqpoint{2.435742in}{0.842463in}}%
\pgfpathlineto{\pgfqpoint{2.440442in}{0.846333in}}%
\pgfpathlineto{\pgfqpoint{2.447000in}{0.848563in}}%
\pgfpathlineto{\pgfqpoint{2.452246in}{0.846939in}}%
\pgfpathlineto{\pgfqpoint{2.468037in}{0.838774in}}%
\pgfpathlineto{\pgfqpoint{2.473536in}{0.831913in}}%
\pgfpathlineto{\pgfqpoint{2.468621in}{0.832178in}}%
\pgfpathlineto{\pgfqpoint{2.460350in}{0.826839in}}%
\pgfpathlineto{\pgfqpoint{2.452424in}{0.825463in}}%
\pgfpathlineto{\pgfqpoint{2.449159in}{0.822767in}}%
\pgfpathlineto{\pgfqpoint{2.446193in}{0.816402in}}%
\pgfpathlineto{\pgfqpoint{2.442122in}{0.815272in}}%
\pgfpathlineto{\pgfqpoint{2.441136in}{0.811210in}}%
\pgfpathlineto{\pgfqpoint{2.432103in}{0.811727in}}%
\pgfpathlineto{\pgfqpoint{2.424800in}{0.803148in}}%
\pgfpathlineto{\pgfqpoint{2.417379in}{0.810753in}}%
\pgfpathlineto{\pgfqpoint{2.408663in}{0.823968in}}%
\pgfpathclose%
\pgfusepath{fill}%
\end{pgfscope}%
\begin{pgfscope}%
\pgfpathrectangle{\pgfqpoint{0.100000in}{0.100000in}}{\pgfqpoint{3.007045in}{1.925000in}}%
\pgfusepath{clip}%
\pgfsetbuttcap%
\pgfsetmiterjoin%
\definecolor{currentfill}{rgb}{0.474279,0.777670,0.478677}%
\pgfsetfillcolor{currentfill}%
\pgfsetlinewidth{0.000000pt}%
\definecolor{currentstroke}{rgb}{0.000000,0.000000,0.000000}%
\pgfsetstrokecolor{currentstroke}%
\pgfsetstrokeopacity{0.000000}%
\pgfsetdash{}{0pt}%
\pgfpathmoveto{\pgfqpoint{1.854945in}{1.201835in}}%
\pgfpathlineto{\pgfqpoint{1.855015in}{1.197513in}}%
\pgfpathlineto{\pgfqpoint{1.834849in}{1.196667in}}%
\pgfpathlineto{\pgfqpoint{1.834576in}{1.208087in}}%
\pgfpathlineto{\pgfqpoint{1.833661in}{1.237372in}}%
\pgfpathlineto{\pgfqpoint{1.844539in}{1.237709in}}%
\pgfpathlineto{\pgfqpoint{1.853670in}{1.237862in}}%
\pgfpathlineto{\pgfqpoint{1.854945in}{1.201835in}}%
\pgfpathclose%
\pgfusepath{fill}%
\end{pgfscope}%
\begin{pgfscope}%
\pgfpathrectangle{\pgfqpoint{0.100000in}{0.100000in}}{\pgfqpoint{3.007045in}{1.925000in}}%
\pgfusepath{clip}%
\pgfsetbuttcap%
\pgfsetmiterjoin%
\definecolor{currentfill}{rgb}{0.590542,0.831926,0.574271}%
\pgfsetfillcolor{currentfill}%
\pgfsetlinewidth{0.000000pt}%
\definecolor{currentstroke}{rgb}{0.000000,0.000000,0.000000}%
\pgfsetstrokecolor{currentstroke}%
\pgfsetstrokeopacity{0.000000}%
\pgfsetdash{}{0pt}%
\pgfpathmoveto{\pgfqpoint{1.538952in}{0.605702in}}%
\pgfpathlineto{\pgfqpoint{1.516589in}{0.606731in}}%
\pgfpathlineto{\pgfqpoint{1.509816in}{0.607078in}}%
\pgfpathlineto{\pgfqpoint{1.510333in}{0.616812in}}%
\pgfpathlineto{\pgfqpoint{1.481225in}{0.618416in}}%
\pgfpathlineto{\pgfqpoint{1.483290in}{0.651066in}}%
\pgfpathlineto{\pgfqpoint{1.476325in}{0.651539in}}%
\pgfpathlineto{\pgfqpoint{1.478210in}{0.684741in}}%
\pgfpathlineto{\pgfqpoint{1.482882in}{0.684491in}}%
\pgfpathlineto{\pgfqpoint{1.507461in}{0.683119in}}%
\pgfpathlineto{\pgfqpoint{1.505217in}{0.648714in}}%
\pgfpathlineto{\pgfqpoint{1.510361in}{0.642115in}}%
\pgfpathlineto{\pgfqpoint{1.520806in}{0.642095in}}%
\pgfpathlineto{\pgfqpoint{1.527773in}{0.637409in}}%
\pgfpathlineto{\pgfqpoint{1.532433in}{0.642045in}}%
\pgfpathlineto{\pgfqpoint{1.540604in}{0.640218in}}%
\pgfpathlineto{\pgfqpoint{1.538952in}{0.605702in}}%
\pgfpathclose%
\pgfusepath{fill}%
\end{pgfscope}%
\begin{pgfscope}%
\pgfpathrectangle{\pgfqpoint{0.100000in}{0.100000in}}{\pgfqpoint{3.007045in}{1.925000in}}%
\pgfusepath{clip}%
\pgfsetbuttcap%
\pgfsetmiterjoin%
\definecolor{currentfill}{rgb}{0.746498,0.899454,0.719938}%
\pgfsetfillcolor{currentfill}%
\pgfsetlinewidth{0.000000pt}%
\definecolor{currentstroke}{rgb}{0.000000,0.000000,0.000000}%
\pgfsetstrokecolor{currentstroke}%
\pgfsetstrokeopacity{0.000000}%
\pgfsetdash{}{0pt}%
\pgfpathmoveto{\pgfqpoint{2.076881in}{0.613497in}}%
\pgfpathlineto{\pgfqpoint{2.075185in}{0.642117in}}%
\pgfpathlineto{\pgfqpoint{2.065550in}{0.641578in}}%
\pgfpathlineto{\pgfqpoint{2.065780in}{0.638658in}}%
\pgfpathlineto{\pgfqpoint{2.058100in}{0.638257in}}%
\pgfpathlineto{\pgfqpoint{2.056646in}{0.646958in}}%
\pgfpathlineto{\pgfqpoint{2.055851in}{0.662822in}}%
\pgfpathlineto{\pgfqpoint{2.073893in}{0.665116in}}%
\pgfpathlineto{\pgfqpoint{2.088016in}{0.667046in}}%
\pgfpathlineto{\pgfqpoint{2.089615in}{0.643002in}}%
\pgfpathlineto{\pgfqpoint{2.121151in}{0.645024in}}%
\pgfpathlineto{\pgfqpoint{2.123918in}{0.610337in}}%
\pgfpathlineto{\pgfqpoint{2.095084in}{0.608430in}}%
\pgfpathlineto{\pgfqpoint{2.094715in}{0.614481in}}%
\pgfpathlineto{\pgfqpoint{2.085520in}{0.614773in}}%
\pgfpathlineto{\pgfqpoint{2.076881in}{0.613497in}}%
\pgfpathclose%
\pgfusepath{fill}%
\end{pgfscope}%
\begin{pgfscope}%
\pgfpathrectangle{\pgfqpoint{0.100000in}{0.100000in}}{\pgfqpoint{3.007045in}{1.925000in}}%
\pgfusepath{clip}%
\pgfsetbuttcap%
\pgfsetmiterjoin%
\definecolor{currentfill}{rgb}{0.797924,0.920738,0.771642}%
\pgfsetfillcolor{currentfill}%
\pgfsetlinewidth{0.000000pt}%
\definecolor{currentstroke}{rgb}{0.000000,0.000000,0.000000}%
\pgfsetstrokecolor{currentstroke}%
\pgfsetstrokeopacity{0.000000}%
\pgfsetdash{}{0pt}%
\pgfpathmoveto{\pgfqpoint{2.624882in}{1.186857in}}%
\pgfpathlineto{\pgfqpoint{2.625996in}{1.195904in}}%
\pgfpathlineto{\pgfqpoint{2.620005in}{1.199545in}}%
\pgfpathlineto{\pgfqpoint{2.622986in}{1.207692in}}%
\pgfpathlineto{\pgfqpoint{2.621536in}{1.211734in}}%
\pgfpathlineto{\pgfqpoint{2.634232in}{1.214856in}}%
\pgfpathlineto{\pgfqpoint{2.637469in}{1.212817in}}%
\pgfpathlineto{\pgfqpoint{2.643347in}{1.220746in}}%
\pgfpathlineto{\pgfqpoint{2.644431in}{1.224390in}}%
\pgfpathlineto{\pgfqpoint{2.646870in}{1.236936in}}%
\pgfpathlineto{\pgfqpoint{2.649600in}{1.238259in}}%
\pgfpathlineto{\pgfqpoint{2.651264in}{1.258406in}}%
\pgfpathlineto{\pgfqpoint{2.655502in}{1.265959in}}%
\pgfpathlineto{\pgfqpoint{2.667563in}{1.268330in}}%
\pgfpathlineto{\pgfqpoint{2.664785in}{1.260556in}}%
\pgfpathlineto{\pgfqpoint{2.672090in}{1.258055in}}%
\pgfpathlineto{\pgfqpoint{2.675901in}{1.254838in}}%
\pgfpathlineto{\pgfqpoint{2.673952in}{1.243959in}}%
\pgfpathlineto{\pgfqpoint{2.677394in}{1.240185in}}%
\pgfpathlineto{\pgfqpoint{2.684505in}{1.237247in}}%
\pgfpathlineto{\pgfqpoint{2.686960in}{1.233890in}}%
\pgfpathlineto{\pgfqpoint{2.703023in}{1.225995in}}%
\pgfpathlineto{\pgfqpoint{2.705509in}{1.221124in}}%
\pgfpathlineto{\pgfqpoint{2.704557in}{1.214940in}}%
\pgfpathlineto{\pgfqpoint{2.706153in}{1.210930in}}%
\pgfpathlineto{\pgfqpoint{2.715119in}{1.210272in}}%
\pgfpathlineto{\pgfqpoint{2.717999in}{1.199175in}}%
\pgfpathlineto{\pgfqpoint{2.726182in}{1.190843in}}%
\pgfpathlineto{\pgfqpoint{2.729054in}{1.179712in}}%
\pgfpathlineto{\pgfqpoint{2.733376in}{1.175355in}}%
\pgfpathlineto{\pgfqpoint{2.732254in}{1.169909in}}%
\pgfpathlineto{\pgfqpoint{2.727131in}{1.173840in}}%
\pgfpathlineto{\pgfqpoint{2.714351in}{1.177897in}}%
\pgfpathlineto{\pgfqpoint{2.708831in}{1.176381in}}%
\pgfpathlineto{\pgfqpoint{2.695566in}{1.180728in}}%
\pgfpathlineto{\pgfqpoint{2.691552in}{1.182331in}}%
\pgfpathlineto{\pgfqpoint{2.682867in}{1.178238in}}%
\pgfpathlineto{\pgfqpoint{2.680316in}{1.180573in}}%
\pgfpathlineto{\pgfqpoint{2.678549in}{1.188294in}}%
\pgfpathlineto{\pgfqpoint{2.668861in}{1.192884in}}%
\pgfpathlineto{\pgfqpoint{2.662595in}{1.183012in}}%
\pgfpathlineto{\pgfqpoint{2.663143in}{1.179525in}}%
\pgfpathlineto{\pgfqpoint{2.658001in}{1.178784in}}%
\pgfpathlineto{\pgfqpoint{2.648785in}{1.188481in}}%
\pgfpathlineto{\pgfqpoint{2.644402in}{1.195106in}}%
\pgfpathlineto{\pgfqpoint{2.634662in}{1.181683in}}%
\pgfpathlineto{\pgfqpoint{2.631429in}{1.181706in}}%
\pgfpathlineto{\pgfqpoint{2.624882in}{1.186857in}}%
\pgfpathclose%
\pgfusepath{fill}%
\end{pgfscope}%
\begin{pgfscope}%
\pgfpathrectangle{\pgfqpoint{0.100000in}{0.100000in}}{\pgfqpoint{3.007045in}{1.925000in}}%
\pgfusepath{clip}%
\pgfsetbuttcap%
\pgfsetmiterjoin%
\definecolor{currentfill}{rgb}{0.193541,0.605136,0.315617}%
\pgfsetfillcolor{currentfill}%
\pgfsetlinewidth{0.000000pt}%
\definecolor{currentstroke}{rgb}{0.000000,0.000000,0.000000}%
\pgfsetstrokecolor{currentstroke}%
\pgfsetstrokeopacity{0.000000}%
\pgfsetdash{}{0pt}%
\pgfpathmoveto{\pgfqpoint{0.521615in}{1.353758in}}%
\pgfpathlineto{\pgfqpoint{0.516565in}{1.336204in}}%
\pgfpathlineto{\pgfqpoint{0.492015in}{1.343273in}}%
\pgfpathlineto{\pgfqpoint{0.489895in}{1.349175in}}%
\pgfpathlineto{\pgfqpoint{0.486181in}{1.350484in}}%
\pgfpathlineto{\pgfqpoint{0.480726in}{1.347606in}}%
\pgfpathlineto{\pgfqpoint{0.465761in}{1.347164in}}%
\pgfpathlineto{\pgfqpoint{0.467458in}{1.355133in}}%
\pgfpathlineto{\pgfqpoint{0.457960in}{1.381368in}}%
\pgfpathlineto{\pgfqpoint{0.455804in}{1.383109in}}%
\pgfpathlineto{\pgfqpoint{0.458164in}{1.392328in}}%
\pgfpathlineto{\pgfqpoint{0.461995in}{1.397568in}}%
\pgfpathlineto{\pgfqpoint{0.459905in}{1.401168in}}%
\pgfpathlineto{\pgfqpoint{0.464449in}{1.404081in}}%
\pgfpathlineto{\pgfqpoint{0.466986in}{1.410310in}}%
\pgfpathlineto{\pgfqpoint{0.462117in}{1.414113in}}%
\pgfpathlineto{\pgfqpoint{0.462527in}{1.420581in}}%
\pgfpathlineto{\pgfqpoint{0.470680in}{1.418127in}}%
\pgfpathlineto{\pgfqpoint{0.483507in}{1.414404in}}%
\pgfpathlineto{\pgfqpoint{0.479961in}{1.402391in}}%
\pgfpathlineto{\pgfqpoint{0.482087in}{1.399330in}}%
\pgfpathlineto{\pgfqpoint{0.488838in}{1.398768in}}%
\pgfpathlineto{\pgfqpoint{0.495392in}{1.401958in}}%
\pgfpathlineto{\pgfqpoint{0.503882in}{1.397397in}}%
\pgfpathlineto{\pgfqpoint{0.512166in}{1.383346in}}%
\pgfpathlineto{\pgfqpoint{0.518020in}{1.379672in}}%
\pgfpathlineto{\pgfqpoint{0.517118in}{1.374961in}}%
\pgfpathlineto{\pgfqpoint{0.520278in}{1.368998in}}%
\pgfpathlineto{\pgfqpoint{0.517066in}{1.358013in}}%
\pgfpathlineto{\pgfqpoint{0.521615in}{1.353758in}}%
\pgfpathclose%
\pgfusepath{fill}%
\end{pgfscope}%
\begin{pgfscope}%
\pgfpathrectangle{\pgfqpoint{0.100000in}{0.100000in}}{\pgfqpoint{3.007045in}{1.925000in}}%
\pgfusepath{clip}%
\pgfsetbuttcap%
\pgfsetmiterjoin%
\definecolor{currentfill}{rgb}{0.535179,0.806090,0.528750}%
\pgfsetfillcolor{currentfill}%
\pgfsetlinewidth{0.000000pt}%
\definecolor{currentstroke}{rgb}{0.000000,0.000000,0.000000}%
\pgfsetstrokecolor{currentstroke}%
\pgfsetstrokeopacity{0.000000}%
\pgfsetdash{}{0pt}%
\pgfpathmoveto{\pgfqpoint{1.590530in}{1.433489in}}%
\pgfpathlineto{\pgfqpoint{1.591111in}{1.462290in}}%
\pgfpathlineto{\pgfqpoint{1.613769in}{1.461576in}}%
\pgfpathlineto{\pgfqpoint{1.614383in}{1.478679in}}%
\pgfpathlineto{\pgfqpoint{1.637041in}{1.477807in}}%
\pgfpathlineto{\pgfqpoint{1.637239in}{1.477803in}}%
\pgfpathlineto{\pgfqpoint{1.636440in}{1.454988in}}%
\pgfpathlineto{\pgfqpoint{1.648008in}{1.454576in}}%
\pgfpathlineto{\pgfqpoint{1.647409in}{1.431503in}}%
\pgfpathlineto{\pgfqpoint{1.590530in}{1.433489in}}%
\pgfpathclose%
\pgfusepath{fill}%
\end{pgfscope}%
\begin{pgfscope}%
\pgfpathrectangle{\pgfqpoint{0.100000in}{0.100000in}}{\pgfqpoint{3.007045in}{1.925000in}}%
\pgfusepath{clip}%
\pgfsetbuttcap%
\pgfsetmiterjoin%
\definecolor{currentfill}{rgb}{0.320000,0.702499,0.396617}%
\pgfsetfillcolor{currentfill}%
\pgfsetlinewidth{0.000000pt}%
\definecolor{currentstroke}{rgb}{0.000000,0.000000,0.000000}%
\pgfsetstrokecolor{currentstroke}%
\pgfsetstrokeopacity{0.000000}%
\pgfsetdash{}{0pt}%
\pgfpathmoveto{\pgfqpoint{0.797625in}{1.916024in}}%
\pgfpathlineto{\pgfqpoint{0.789229in}{1.906864in}}%
\pgfpathlineto{\pgfqpoint{0.787685in}{1.900755in}}%
\pgfpathlineto{\pgfqpoint{0.783449in}{1.895856in}}%
\pgfpathlineto{\pgfqpoint{0.789000in}{1.894475in}}%
\pgfpathlineto{\pgfqpoint{0.780353in}{1.860745in}}%
\pgfpathlineto{\pgfqpoint{0.777516in}{1.855475in}}%
\pgfpathlineto{\pgfqpoint{0.773351in}{1.839140in}}%
\pgfpathlineto{\pgfqpoint{0.768647in}{1.844176in}}%
\pgfpathlineto{\pgfqpoint{0.762060in}{1.843258in}}%
\pgfpathlineto{\pgfqpoint{0.759267in}{1.846349in}}%
\pgfpathlineto{\pgfqpoint{0.749934in}{1.845470in}}%
\pgfpathlineto{\pgfqpoint{0.748542in}{1.851631in}}%
\pgfpathlineto{\pgfqpoint{0.744909in}{1.856652in}}%
\pgfpathlineto{\pgfqpoint{0.735184in}{1.852084in}}%
\pgfpathlineto{\pgfqpoint{0.730840in}{1.859613in}}%
\pgfpathlineto{\pgfqpoint{0.719784in}{1.864006in}}%
\pgfpathlineto{\pgfqpoint{0.713636in}{1.864996in}}%
\pgfpathlineto{\pgfqpoint{0.713573in}{1.866288in}}%
\pgfpathlineto{\pgfqpoint{0.713541in}{1.878149in}}%
\pgfpathlineto{\pgfqpoint{0.710511in}{1.880181in}}%
\pgfpathlineto{\pgfqpoint{0.699664in}{1.876125in}}%
\pgfpathlineto{\pgfqpoint{0.695862in}{1.879153in}}%
\pgfpathlineto{\pgfqpoint{0.685273in}{1.876771in}}%
\pgfpathlineto{\pgfqpoint{0.683348in}{1.884220in}}%
\pgfpathlineto{\pgfqpoint{0.676269in}{1.883139in}}%
\pgfpathlineto{\pgfqpoint{0.675310in}{1.876711in}}%
\pgfpathlineto{\pgfqpoint{0.667442in}{1.879016in}}%
\pgfpathlineto{\pgfqpoint{0.666800in}{1.883212in}}%
\pgfpathlineto{\pgfqpoint{0.657890in}{1.895448in}}%
\pgfpathlineto{\pgfqpoint{0.659239in}{1.897832in}}%
\pgfpathlineto{\pgfqpoint{0.654234in}{1.907098in}}%
\pgfpathlineto{\pgfqpoint{0.650154in}{1.917880in}}%
\pgfpathlineto{\pgfqpoint{0.650253in}{1.923532in}}%
\pgfpathlineto{\pgfqpoint{0.648297in}{1.930642in}}%
\pgfpathlineto{\pgfqpoint{0.650462in}{1.932235in}}%
\pgfpathlineto{\pgfqpoint{0.654305in}{1.933747in}}%
\pgfpathlineto{\pgfqpoint{0.652693in}{1.955469in}}%
\pgfpathlineto{\pgfqpoint{0.684552in}{1.946256in}}%
\pgfpathlineto{\pgfqpoint{0.723267in}{1.935529in}}%
\pgfpathlineto{\pgfqpoint{0.752077in}{1.927793in}}%
\pgfpathlineto{\pgfqpoint{0.797625in}{1.916024in}}%
\pgfpathclose%
\pgfusepath{fill}%
\end{pgfscope}%
\begin{pgfscope}%
\pgfpathrectangle{\pgfqpoint{0.100000in}{0.100000in}}{\pgfqpoint{3.007045in}{1.925000in}}%
\pgfusepath{clip}%
\pgfsetbuttcap%
\pgfsetmiterjoin%
\definecolor{currentfill}{rgb}{0.652995,0.860085,0.628897}%
\pgfsetfillcolor{currentfill}%
\pgfsetlinewidth{0.000000pt}%
\definecolor{currentstroke}{rgb}{0.000000,0.000000,0.000000}%
\pgfsetstrokecolor{currentstroke}%
\pgfsetstrokeopacity{0.000000}%
\pgfsetdash{}{0pt}%
\pgfpathmoveto{\pgfqpoint{2.367977in}{0.755519in}}%
\pgfpathlineto{\pgfqpoint{2.361889in}{0.757434in}}%
\pgfpathlineto{\pgfqpoint{2.354579in}{0.763090in}}%
\pgfpathlineto{\pgfqpoint{2.352032in}{0.771233in}}%
\pgfpathlineto{\pgfqpoint{2.352426in}{0.785269in}}%
\pgfpathlineto{\pgfqpoint{2.352410in}{0.787703in}}%
\pgfpathlineto{\pgfqpoint{2.355976in}{0.788053in}}%
\pgfpathlineto{\pgfqpoint{2.357695in}{0.794657in}}%
\pgfpathlineto{\pgfqpoint{2.365423in}{0.795571in}}%
\pgfpathlineto{\pgfqpoint{2.375369in}{0.793380in}}%
\pgfpathlineto{\pgfqpoint{2.382395in}{0.804196in}}%
\pgfpathlineto{\pgfqpoint{2.386212in}{0.799739in}}%
\pgfpathlineto{\pgfqpoint{2.390080in}{0.791606in}}%
\pgfpathlineto{\pgfqpoint{2.390153in}{0.787078in}}%
\pgfpathlineto{\pgfqpoint{2.387448in}{0.788188in}}%
\pgfpathlineto{\pgfqpoint{2.377909in}{0.787069in}}%
\pgfpathlineto{\pgfqpoint{2.379793in}{0.770306in}}%
\pgfpathlineto{\pgfqpoint{2.375546in}{0.768744in}}%
\pgfpathlineto{\pgfqpoint{2.375145in}{0.759957in}}%
\pgfpathlineto{\pgfqpoint{2.373071in}{0.752288in}}%
\pgfpathlineto{\pgfqpoint{2.367977in}{0.755519in}}%
\pgfpathclose%
\pgfusepath{fill}%
\end{pgfscope}%
\begin{pgfscope}%
\pgfpathrectangle{\pgfqpoint{0.100000in}{0.100000in}}{\pgfqpoint{3.007045in}{1.925000in}}%
\pgfusepath{clip}%
\pgfsetbuttcap%
\pgfsetmiterjoin%
\definecolor{currentfill}{rgb}{0.695071,0.877801,0.669865}%
\pgfsetfillcolor{currentfill}%
\pgfsetlinewidth{0.000000pt}%
\definecolor{currentstroke}{rgb}{0.000000,0.000000,0.000000}%
\pgfsetstrokecolor{currentstroke}%
\pgfsetstrokeopacity{0.000000}%
\pgfsetdash{}{0pt}%
\pgfpathmoveto{\pgfqpoint{2.715119in}{1.210272in}}%
\pgfpathlineto{\pgfqpoint{2.706153in}{1.210930in}}%
\pgfpathlineto{\pgfqpoint{2.704557in}{1.214940in}}%
\pgfpathlineto{\pgfqpoint{2.705509in}{1.221124in}}%
\pgfpathlineto{\pgfqpoint{2.703023in}{1.225995in}}%
\pgfpathlineto{\pgfqpoint{2.686960in}{1.233890in}}%
\pgfpathlineto{\pgfqpoint{2.684505in}{1.237247in}}%
\pgfpathlineto{\pgfqpoint{2.677394in}{1.240185in}}%
\pgfpathlineto{\pgfqpoint{2.673952in}{1.243959in}}%
\pgfpathlineto{\pgfqpoint{2.675901in}{1.254838in}}%
\pgfpathlineto{\pgfqpoint{2.672090in}{1.258055in}}%
\pgfpathlineto{\pgfqpoint{2.664785in}{1.260556in}}%
\pgfpathlineto{\pgfqpoint{2.667563in}{1.268330in}}%
\pgfpathlineto{\pgfqpoint{2.716037in}{1.278304in}}%
\pgfpathlineto{\pgfqpoint{2.716371in}{1.278367in}}%
\pgfpathlineto{\pgfqpoint{2.726638in}{1.268556in}}%
\pgfpathlineto{\pgfqpoint{2.726766in}{1.261212in}}%
\pgfpathlineto{\pgfqpoint{2.723696in}{1.256764in}}%
\pgfpathlineto{\pgfqpoint{2.715251in}{1.250311in}}%
\pgfpathlineto{\pgfqpoint{2.713010in}{1.242278in}}%
\pgfpathlineto{\pgfqpoint{2.715503in}{1.232898in}}%
\pgfpathlineto{\pgfqpoint{2.715994in}{1.223495in}}%
\pgfpathlineto{\pgfqpoint{2.713299in}{1.211895in}}%
\pgfpathlineto{\pgfqpoint{2.715119in}{1.210272in}}%
\pgfpathclose%
\pgfusepath{fill}%
\end{pgfscope}%
\begin{pgfscope}%
\pgfpathrectangle{\pgfqpoint{0.100000in}{0.100000in}}{\pgfqpoint{3.007045in}{1.925000in}}%
\pgfusepath{clip}%
\pgfsetbuttcap%
\pgfsetmiterjoin%
\definecolor{currentfill}{rgb}{0.507497,0.793172,0.505990}%
\pgfsetfillcolor{currentfill}%
\pgfsetlinewidth{0.000000pt}%
\definecolor{currentstroke}{rgb}{0.000000,0.000000,0.000000}%
\pgfsetstrokecolor{currentstroke}%
\pgfsetstrokeopacity{0.000000}%
\pgfsetdash{}{0pt}%
\pgfpathmoveto{\pgfqpoint{1.534859in}{1.055660in}}%
\pgfpathlineto{\pgfqpoint{1.534600in}{1.061649in}}%
\pgfpathlineto{\pgfqpoint{1.536034in}{1.090382in}}%
\pgfpathlineto{\pgfqpoint{1.535329in}{1.090418in}}%
\pgfpathlineto{\pgfqpoint{1.536726in}{1.118994in}}%
\pgfpathlineto{\pgfqpoint{1.593165in}{1.116517in}}%
\pgfpathlineto{\pgfqpoint{1.591676in}{1.059068in}}%
\pgfpathlineto{\pgfqpoint{1.563552in}{1.060256in}}%
\pgfpathlineto{\pgfqpoint{1.563272in}{1.054193in}}%
\pgfpathlineto{\pgfqpoint{1.534859in}{1.055660in}}%
\pgfpathclose%
\pgfusepath{fill}%
\end{pgfscope}%
\begin{pgfscope}%
\pgfpathrectangle{\pgfqpoint{0.100000in}{0.100000in}}{\pgfqpoint{3.007045in}{1.925000in}}%
\pgfusepath{clip}%
\pgfsetbuttcap%
\pgfsetmiterjoin%
\definecolor{currentfill}{rgb}{0.245213,0.660254,0.356955}%
\pgfsetfillcolor{currentfill}%
\pgfsetlinewidth{0.000000pt}%
\definecolor{currentstroke}{rgb}{0.000000,0.000000,0.000000}%
\pgfsetstrokecolor{currentstroke}%
\pgfsetstrokeopacity{0.000000}%
\pgfsetdash{}{0pt}%
\pgfpathmoveto{\pgfqpoint{2.204865in}{1.518579in}}%
\pgfpathlineto{\pgfqpoint{2.185259in}{1.517003in}}%
\pgfpathlineto{\pgfqpoint{2.183604in}{1.529014in}}%
\pgfpathlineto{\pgfqpoint{2.190595in}{1.532298in}}%
\pgfpathlineto{\pgfqpoint{2.190444in}{1.542079in}}%
\pgfpathlineto{\pgfqpoint{2.194250in}{1.543292in}}%
\pgfpathlineto{\pgfqpoint{2.196243in}{1.547670in}}%
\pgfpathlineto{\pgfqpoint{2.202517in}{1.547510in}}%
\pgfpathlineto{\pgfqpoint{2.204610in}{1.554457in}}%
\pgfpathlineto{\pgfqpoint{2.210201in}{1.559918in}}%
\pgfpathlineto{\pgfqpoint{2.211986in}{1.556493in}}%
\pgfpathlineto{\pgfqpoint{2.212190in}{1.546743in}}%
\pgfpathlineto{\pgfqpoint{2.210346in}{1.543780in}}%
\pgfpathlineto{\pgfqpoint{2.211388in}{1.536695in}}%
\pgfpathlineto{\pgfqpoint{2.216884in}{1.535425in}}%
\pgfpathlineto{\pgfqpoint{2.221340in}{1.546582in}}%
\pgfpathlineto{\pgfqpoint{2.222009in}{1.560126in}}%
\pgfpathlineto{\pgfqpoint{2.220169in}{1.566322in}}%
\pgfpathlineto{\pgfqpoint{2.227849in}{1.567100in}}%
\pgfpathlineto{\pgfqpoint{2.228428in}{1.561303in}}%
\pgfpathlineto{\pgfqpoint{2.245496in}{1.563087in}}%
\pgfpathlineto{\pgfqpoint{2.247824in}{1.546163in}}%
\pgfpathlineto{\pgfqpoint{2.250233in}{1.523286in}}%
\pgfpathlineto{\pgfqpoint{2.204865in}{1.518579in}}%
\pgfpathclose%
\pgfusepath{fill}%
\end{pgfscope}%
\begin{pgfscope}%
\pgfpathrectangle{\pgfqpoint{0.100000in}{0.100000in}}{\pgfqpoint{3.007045in}{1.925000in}}%
\pgfusepath{clip}%
\pgfsetbuttcap%
\pgfsetmiterjoin%
\definecolor{currentfill}{rgb}{0.662345,0.864022,0.638001}%
\pgfsetfillcolor{currentfill}%
\pgfsetlinewidth{0.000000pt}%
\definecolor{currentstroke}{rgb}{0.000000,0.000000,0.000000}%
\pgfsetstrokecolor{currentstroke}%
\pgfsetstrokeopacity{0.000000}%
\pgfsetdash{}{0pt}%
\pgfpathmoveto{\pgfqpoint{1.541684in}{1.689945in}}%
\pgfpathlineto{\pgfqpoint{1.540344in}{1.666994in}}%
\pgfpathlineto{\pgfqpoint{1.542129in}{1.666889in}}%
\pgfpathlineto{\pgfqpoint{1.540553in}{1.643789in}}%
\pgfpathlineto{\pgfqpoint{1.514103in}{1.645425in}}%
\pgfpathlineto{\pgfqpoint{1.518209in}{1.642479in}}%
\pgfpathlineto{\pgfqpoint{1.519039in}{1.638417in}}%
\pgfpathlineto{\pgfqpoint{1.516525in}{1.631523in}}%
\pgfpathlineto{\pgfqpoint{1.506325in}{1.629275in}}%
\pgfpathlineto{\pgfqpoint{1.503158in}{1.631150in}}%
\pgfpathlineto{\pgfqpoint{1.495994in}{1.623275in}}%
\pgfpathlineto{\pgfqpoint{1.495332in}{1.629121in}}%
\pgfpathlineto{\pgfqpoint{1.483880in}{1.629904in}}%
\pgfpathlineto{\pgfqpoint{1.485073in}{1.647065in}}%
\pgfpathlineto{\pgfqpoint{1.465987in}{1.648477in}}%
\pgfpathlineto{\pgfqpoint{1.466420in}{1.654245in}}%
\pgfpathlineto{\pgfqpoint{1.449237in}{1.655533in}}%
\pgfpathlineto{\pgfqpoint{1.450582in}{1.673001in}}%
\pgfpathlineto{\pgfqpoint{1.448415in}{1.673164in}}%
\pgfpathlineto{\pgfqpoint{1.450238in}{1.696191in}}%
\pgfpathlineto{\pgfqpoint{1.447485in}{1.696400in}}%
\pgfpathlineto{\pgfqpoint{1.448796in}{1.712766in}}%
\pgfpathlineto{\pgfqpoint{1.446373in}{1.716181in}}%
\pgfpathlineto{\pgfqpoint{1.447405in}{1.726710in}}%
\pgfpathlineto{\pgfqpoint{1.441706in}{1.725605in}}%
\pgfpathlineto{\pgfqpoint{1.442195in}{1.731563in}}%
\pgfpathlineto{\pgfqpoint{1.465127in}{1.729765in}}%
\pgfpathlineto{\pgfqpoint{1.522557in}{1.725806in}}%
\pgfpathlineto{\pgfqpoint{1.521849in}{1.714257in}}%
\pgfpathlineto{\pgfqpoint{1.517941in}{1.714507in}}%
\pgfpathlineto{\pgfqpoint{1.516455in}{1.691466in}}%
\pgfpathlineto{\pgfqpoint{1.541684in}{1.689945in}}%
\pgfpathclose%
\pgfusepath{fill}%
\end{pgfscope}%
\begin{pgfscope}%
\pgfpathrectangle{\pgfqpoint{0.100000in}{0.100000in}}{\pgfqpoint{3.007045in}{1.925000in}}%
\pgfusepath{clip}%
\pgfsetbuttcap%
\pgfsetmiterjoin%
\definecolor{currentfill}{rgb}{0.662345,0.864022,0.638001}%
\pgfsetfillcolor{currentfill}%
\pgfsetlinewidth{0.000000pt}%
\definecolor{currentstroke}{rgb}{0.000000,0.000000,0.000000}%
\pgfsetstrokecolor{currentstroke}%
\pgfsetstrokeopacity{0.000000}%
\pgfsetdash{}{0pt}%
\pgfpathmoveto{\pgfqpoint{2.754436in}{1.271463in}}%
\pgfpathlineto{\pgfqpoint{2.754264in}{1.276110in}}%
\pgfpathlineto{\pgfqpoint{2.751124in}{1.279964in}}%
\pgfpathlineto{\pgfqpoint{2.753954in}{1.289095in}}%
\pgfpathlineto{\pgfqpoint{2.753598in}{1.293002in}}%
\pgfpathlineto{\pgfqpoint{2.746343in}{1.292558in}}%
\pgfpathlineto{\pgfqpoint{2.742511in}{1.290165in}}%
\pgfpathlineto{\pgfqpoint{2.739138in}{1.283230in}}%
\pgfpathlineto{\pgfqpoint{2.721216in}{1.279405in}}%
\pgfpathlineto{\pgfqpoint{2.724100in}{1.282597in}}%
\pgfpathlineto{\pgfqpoint{2.725778in}{1.289195in}}%
\pgfpathlineto{\pgfqpoint{2.725049in}{1.295828in}}%
\pgfpathlineto{\pgfqpoint{2.726798in}{1.302007in}}%
\pgfpathlineto{\pgfqpoint{2.725428in}{1.306338in}}%
\pgfpathlineto{\pgfqpoint{2.735728in}{1.317531in}}%
\pgfpathlineto{\pgfqpoint{2.741068in}{1.332531in}}%
\pgfpathlineto{\pgfqpoint{2.746311in}{1.336376in}}%
\pgfpathlineto{\pgfqpoint{2.755066in}{1.346497in}}%
\pgfpathlineto{\pgfqpoint{2.762277in}{1.343620in}}%
\pgfpathlineto{\pgfqpoint{2.764535in}{1.335742in}}%
\pgfpathlineto{\pgfqpoint{2.769299in}{1.335490in}}%
\pgfpathlineto{\pgfqpoint{2.777561in}{1.327350in}}%
\pgfpathlineto{\pgfqpoint{2.791685in}{1.322963in}}%
\pgfpathlineto{\pgfqpoint{2.806809in}{1.301783in}}%
\pgfpathlineto{\pgfqpoint{2.811323in}{1.284246in}}%
\pgfpathlineto{\pgfqpoint{2.815570in}{1.283815in}}%
\pgfpathlineto{\pgfqpoint{2.811785in}{1.275012in}}%
\pgfpathlineto{\pgfqpoint{2.803631in}{1.265614in}}%
\pgfpathlineto{\pgfqpoint{2.798237in}{1.247006in}}%
\pgfpathlineto{\pgfqpoint{2.795440in}{1.242918in}}%
\pgfpathlineto{\pgfqpoint{2.790081in}{1.241725in}}%
\pgfpathlineto{\pgfqpoint{2.790953in}{1.257015in}}%
\pgfpathlineto{\pgfqpoint{2.775479in}{1.258863in}}%
\pgfpathlineto{\pgfqpoint{2.767772in}{1.261597in}}%
\pgfpathlineto{\pgfqpoint{2.757096in}{1.268170in}}%
\pgfpathlineto{\pgfqpoint{2.754436in}{1.271463in}}%
\pgfpathclose%
\pgfusepath{fill}%
\end{pgfscope}%
\begin{pgfscope}%
\pgfpathrectangle{\pgfqpoint{0.100000in}{0.100000in}}{\pgfqpoint{3.007045in}{1.925000in}}%
\pgfusepath{clip}%
\pgfsetbuttcap%
\pgfsetmiterjoin%
\definecolor{currentfill}{rgb}{0.827451,0.932549,0.803137}%
\pgfsetfillcolor{currentfill}%
\pgfsetlinewidth{0.000000pt}%
\definecolor{currentstroke}{rgb}{0.000000,0.000000,0.000000}%
\pgfsetstrokecolor{currentstroke}%
\pgfsetstrokeopacity{0.000000}%
\pgfsetdash{}{0pt}%
\pgfpathmoveto{\pgfqpoint{1.454880in}{1.037434in}}%
\pgfpathlineto{\pgfqpoint{1.453227in}{1.011793in}}%
\pgfpathlineto{\pgfqpoint{1.427563in}{1.013412in}}%
\pgfpathlineto{\pgfqpoint{1.429358in}{1.039255in}}%
\pgfpathlineto{\pgfqpoint{1.404039in}{1.041199in}}%
\pgfpathlineto{\pgfqpoint{1.405343in}{1.057957in}}%
\pgfpathlineto{\pgfqpoint{1.405795in}{1.064151in}}%
\pgfpathlineto{\pgfqpoint{1.455281in}{1.060433in}}%
\pgfpathlineto{\pgfqpoint{1.453744in}{1.037504in}}%
\pgfpathlineto{\pgfqpoint{1.454880in}{1.037434in}}%
\pgfpathclose%
\pgfusepath{fill}%
\end{pgfscope}%
\begin{pgfscope}%
\pgfpathrectangle{\pgfqpoint{0.100000in}{0.100000in}}{\pgfqpoint{3.007045in}{1.925000in}}%
\pgfusepath{clip}%
\pgfsetbuttcap%
\pgfsetmiterjoin%
\definecolor{currentfill}{rgb}{0.596078,0.834510,0.578824}%
\pgfsetfillcolor{currentfill}%
\pgfsetlinewidth{0.000000pt}%
\definecolor{currentstroke}{rgb}{0.000000,0.000000,0.000000}%
\pgfsetstrokecolor{currentstroke}%
\pgfsetstrokeopacity{0.000000}%
\pgfsetdash{}{0pt}%
\pgfpathmoveto{\pgfqpoint{2.397002in}{0.647345in}}%
\pgfpathlineto{\pgfqpoint{2.397277in}{0.644995in}}%
\pgfpathlineto{\pgfqpoint{2.412441in}{0.646597in}}%
\pgfpathlineto{\pgfqpoint{2.415134in}{0.621502in}}%
\pgfpathlineto{\pgfqpoint{2.378692in}{0.619214in}}%
\pgfpathlineto{\pgfqpoint{2.375927in}{0.644931in}}%
\pgfpathlineto{\pgfqpoint{2.397002in}{0.647345in}}%
\pgfpathclose%
\pgfusepath{fill}%
\end{pgfscope}%
\begin{pgfscope}%
\pgfpathrectangle{\pgfqpoint{0.100000in}{0.100000in}}{\pgfqpoint{3.007045in}{1.925000in}}%
\pgfusepath{clip}%
\pgfsetbuttcap%
\pgfsetmiterjoin%
\definecolor{currentfill}{rgb}{0.513033,0.795755,0.510542}%
\pgfsetfillcolor{currentfill}%
\pgfsetlinewidth{0.000000pt}%
\definecolor{currentstroke}{rgb}{0.000000,0.000000,0.000000}%
\pgfsetstrokecolor{currentstroke}%
\pgfsetstrokeopacity{0.000000}%
\pgfsetdash{}{0pt}%
\pgfpathmoveto{\pgfqpoint{2.116398in}{1.073683in}}%
\pgfpathlineto{\pgfqpoint{2.099179in}{1.072352in}}%
\pgfpathlineto{\pgfqpoint{2.098280in}{1.086724in}}%
\pgfpathlineto{\pgfqpoint{2.075282in}{1.085153in}}%
\pgfpathlineto{\pgfqpoint{2.075029in}{1.090942in}}%
\pgfpathlineto{\pgfqpoint{2.046290in}{1.089601in}}%
\pgfpathlineto{\pgfqpoint{2.044424in}{1.118250in}}%
\pgfpathlineto{\pgfqpoint{2.050157in}{1.118638in}}%
\pgfpathlineto{\pgfqpoint{2.049708in}{1.124375in}}%
\pgfpathlineto{\pgfqpoint{2.067309in}{1.125374in}}%
\pgfpathlineto{\pgfqpoint{2.073272in}{1.125381in}}%
\pgfpathlineto{\pgfqpoint{2.072910in}{1.131159in}}%
\pgfpathlineto{\pgfqpoint{2.095572in}{1.132813in}}%
\pgfpathlineto{\pgfqpoint{2.096418in}{1.118342in}}%
\pgfpathlineto{\pgfqpoint{2.097811in}{1.095356in}}%
\pgfpathlineto{\pgfqpoint{2.114940in}{1.096476in}}%
\pgfpathlineto{\pgfqpoint{2.116398in}{1.073683in}}%
\pgfpathclose%
\pgfusepath{fill}%
\end{pgfscope}%
\begin{pgfscope}%
\pgfpathrectangle{\pgfqpoint{0.100000in}{0.100000in}}{\pgfqpoint{3.007045in}{1.925000in}}%
\pgfusepath{clip}%
\pgfsetbuttcap%
\pgfsetmiterjoin%
\definecolor{currentfill}{rgb}{0.704421,0.881738,0.678970}%
\pgfsetfillcolor{currentfill}%
\pgfsetlinewidth{0.000000pt}%
\definecolor{currentstroke}{rgb}{0.000000,0.000000,0.000000}%
\pgfsetstrokecolor{currentstroke}%
\pgfsetstrokeopacity{0.000000}%
\pgfsetdash{}{0pt}%
\pgfpathmoveto{\pgfqpoint{2.241536in}{0.660624in}}%
\pgfpathlineto{\pgfqpoint{2.233755in}{0.660011in}}%
\pgfpathlineto{\pgfqpoint{2.229320in}{0.666780in}}%
\pgfpathlineto{\pgfqpoint{2.228144in}{0.679871in}}%
\pgfpathlineto{\pgfqpoint{2.226858in}{0.694238in}}%
\pgfpathlineto{\pgfqpoint{2.225729in}{0.705930in}}%
\pgfpathlineto{\pgfqpoint{2.230267in}{0.714046in}}%
\pgfpathlineto{\pgfqpoint{2.223953in}{0.726994in}}%
\pgfpathlineto{\pgfqpoint{2.222669in}{0.734886in}}%
\pgfpathlineto{\pgfqpoint{2.233994in}{0.735748in}}%
\pgfpathlineto{\pgfqpoint{2.233682in}{0.738645in}}%
\pgfpathlineto{\pgfqpoint{2.250388in}{0.740313in}}%
\pgfpathlineto{\pgfqpoint{2.252087in}{0.743598in}}%
\pgfpathlineto{\pgfqpoint{2.279667in}{0.746313in}}%
\pgfpathlineto{\pgfqpoint{2.281053in}{0.729035in}}%
\pgfpathlineto{\pgfqpoint{2.274136in}{0.722813in}}%
\pgfpathlineto{\pgfqpoint{2.276404in}{0.718014in}}%
\pgfpathlineto{\pgfqpoint{2.277794in}{0.699248in}}%
\pgfpathlineto{\pgfqpoint{2.278425in}{0.693743in}}%
\pgfpathlineto{\pgfqpoint{2.261248in}{0.691899in}}%
\pgfpathlineto{\pgfqpoint{2.260621in}{0.697611in}}%
\pgfpathlineto{\pgfqpoint{2.254880in}{0.697001in}}%
\pgfpathlineto{\pgfqpoint{2.253069in}{0.691091in}}%
\pgfpathlineto{\pgfqpoint{2.254972in}{0.670665in}}%
\pgfpathlineto{\pgfqpoint{2.252110in}{0.670372in}}%
\pgfpathlineto{\pgfqpoint{2.252934in}{0.661782in}}%
\pgfpathlineto{\pgfqpoint{2.241536in}{0.660624in}}%
\pgfpathclose%
\pgfusepath{fill}%
\end{pgfscope}%
\begin{pgfscope}%
\pgfpathrectangle{\pgfqpoint{0.100000in}{0.100000in}}{\pgfqpoint{3.007045in}{1.925000in}}%
\pgfusepath{clip}%
\pgfsetbuttcap%
\pgfsetmiterjoin%
\definecolor{currentfill}{rgb}{0.746498,0.899454,0.719938}%
\pgfsetfillcolor{currentfill}%
\pgfsetlinewidth{0.000000pt}%
\definecolor{currentstroke}{rgb}{0.000000,0.000000,0.000000}%
\pgfsetstrokecolor{currentstroke}%
\pgfsetstrokeopacity{0.000000}%
\pgfsetdash{}{0pt}%
\pgfpathmoveto{\pgfqpoint{0.565403in}{1.940383in}}%
\pgfpathlineto{\pgfqpoint{0.569196in}{1.933206in}}%
\pgfpathlineto{\pgfqpoint{0.564241in}{1.932920in}}%
\pgfpathlineto{\pgfqpoint{0.561680in}{1.929053in}}%
\pgfpathlineto{\pgfqpoint{0.565762in}{1.918633in}}%
\pgfpathlineto{\pgfqpoint{0.569487in}{1.913159in}}%
\pgfpathlineto{\pgfqpoint{0.569117in}{1.907951in}}%
\pgfpathlineto{\pgfqpoint{0.564737in}{1.906268in}}%
\pgfpathlineto{\pgfqpoint{0.563966in}{1.911840in}}%
\pgfpathlineto{\pgfqpoint{0.559562in}{1.916054in}}%
\pgfpathlineto{\pgfqpoint{0.561416in}{1.920783in}}%
\pgfpathlineto{\pgfqpoint{0.556562in}{1.930703in}}%
\pgfpathlineto{\pgfqpoint{0.565403in}{1.940383in}}%
\pgfpathclose%
\pgfusepath{fill}%
\end{pgfscope}%
\begin{pgfscope}%
\pgfpathrectangle{\pgfqpoint{0.100000in}{0.100000in}}{\pgfqpoint{3.007045in}{1.925000in}}%
\pgfusepath{clip}%
\pgfsetbuttcap%
\pgfsetmiterjoin%
\definecolor{currentfill}{rgb}{0.746498,0.899454,0.719938}%
\pgfsetfillcolor{currentfill}%
\pgfsetlinewidth{0.000000pt}%
\definecolor{currentstroke}{rgb}{0.000000,0.000000,0.000000}%
\pgfsetstrokecolor{currentstroke}%
\pgfsetstrokeopacity{0.000000}%
\pgfsetdash{}{0pt}%
\pgfpathmoveto{\pgfqpoint{0.574289in}{1.797104in}}%
\pgfpathlineto{\pgfqpoint{0.542761in}{1.806508in}}%
\pgfpathlineto{\pgfqpoint{0.493956in}{1.821847in}}%
\pgfpathlineto{\pgfqpoint{0.501743in}{1.847839in}}%
\pgfpathlineto{\pgfqpoint{0.510872in}{1.844976in}}%
\pgfpathlineto{\pgfqpoint{0.515078in}{1.857792in}}%
\pgfpathlineto{\pgfqpoint{0.514994in}{1.864013in}}%
\pgfpathlineto{\pgfqpoint{0.502572in}{1.867867in}}%
\pgfpathlineto{\pgfqpoint{0.505442in}{1.879195in}}%
\pgfpathlineto{\pgfqpoint{0.510848in}{1.895622in}}%
\pgfpathlineto{\pgfqpoint{0.512685in}{1.901288in}}%
\pgfpathlineto{\pgfqpoint{0.536363in}{1.893635in}}%
\pgfpathlineto{\pgfqpoint{0.546025in}{1.896043in}}%
\pgfpathlineto{\pgfqpoint{0.549783in}{1.902072in}}%
\pgfpathlineto{\pgfqpoint{0.556850in}{1.906039in}}%
\pgfpathlineto{\pgfqpoint{0.556363in}{1.910254in}}%
\pgfpathlineto{\pgfqpoint{0.560656in}{1.907168in}}%
\pgfpathlineto{\pgfqpoint{0.559298in}{1.890060in}}%
\pgfpathlineto{\pgfqpoint{0.563615in}{1.893690in}}%
\pgfpathlineto{\pgfqpoint{0.564306in}{1.899159in}}%
\pgfpathlineto{\pgfqpoint{0.570797in}{1.906935in}}%
\pgfpathlineto{\pgfqpoint{0.575526in}{1.910604in}}%
\pgfpathlineto{\pgfqpoint{0.572185in}{1.916693in}}%
\pgfpathlineto{\pgfqpoint{0.573093in}{1.922337in}}%
\pgfpathlineto{\pgfqpoint{0.567141in}{1.928935in}}%
\pgfpathlineto{\pgfqpoint{0.574554in}{1.930262in}}%
\pgfpathlineto{\pgfqpoint{0.568642in}{1.941046in}}%
\pgfpathlineto{\pgfqpoint{0.578403in}{1.950226in}}%
\pgfpathlineto{\pgfqpoint{0.576681in}{1.953612in}}%
\pgfpathlineto{\pgfqpoint{0.623985in}{1.938935in}}%
\pgfpathlineto{\pgfqpoint{0.643432in}{1.933133in}}%
\pgfpathlineto{\pgfqpoint{0.650462in}{1.932235in}}%
\pgfpathlineto{\pgfqpoint{0.648297in}{1.930642in}}%
\pgfpathlineto{\pgfqpoint{0.650253in}{1.923532in}}%
\pgfpathlineto{\pgfqpoint{0.643966in}{1.926557in}}%
\pgfpathlineto{\pgfqpoint{0.633546in}{1.923434in}}%
\pgfpathlineto{\pgfqpoint{0.630805in}{1.915337in}}%
\pgfpathlineto{\pgfqpoint{0.634619in}{1.904375in}}%
\pgfpathlineto{\pgfqpoint{0.631834in}{1.900038in}}%
\pgfpathlineto{\pgfqpoint{0.622765in}{1.897664in}}%
\pgfpathlineto{\pgfqpoint{0.617908in}{1.888997in}}%
\pgfpathlineto{\pgfqpoint{0.619147in}{1.875248in}}%
\pgfpathlineto{\pgfqpoint{0.615087in}{1.868865in}}%
\pgfpathlineto{\pgfqpoint{0.610134in}{1.868414in}}%
\pgfpathlineto{\pgfqpoint{0.600903in}{1.863151in}}%
\pgfpathlineto{\pgfqpoint{0.595348in}{1.857817in}}%
\pgfpathlineto{\pgfqpoint{0.595548in}{1.853584in}}%
\pgfpathlineto{\pgfqpoint{0.599282in}{1.851636in}}%
\pgfpathlineto{\pgfqpoint{0.597616in}{1.841914in}}%
\pgfpathlineto{\pgfqpoint{0.583582in}{1.827974in}}%
\pgfpathlineto{\pgfqpoint{0.585552in}{1.821819in}}%
\pgfpathlineto{\pgfqpoint{0.585013in}{1.811079in}}%
\pgfpathlineto{\pgfqpoint{0.579759in}{1.803398in}}%
\pgfpathlineto{\pgfqpoint{0.579595in}{1.795500in}}%
\pgfpathlineto{\pgfqpoint{0.574289in}{1.797104in}}%
\pgfpathclose%
\pgfusepath{fill}%
\end{pgfscope}%
\begin{pgfscope}%
\pgfpathrectangle{\pgfqpoint{0.100000in}{0.100000in}}{\pgfqpoint{3.007045in}{1.925000in}}%
\pgfusepath{clip}%
\pgfsetbuttcap%
\pgfsetmiterjoin%
\definecolor{currentfill}{rgb}{0.657670,0.862053,0.633449}%
\pgfsetfillcolor{currentfill}%
\pgfsetlinewidth{0.000000pt}%
\definecolor{currentstroke}{rgb}{0.000000,0.000000,0.000000}%
\pgfsetstrokecolor{currentstroke}%
\pgfsetstrokeopacity{0.000000}%
\pgfsetdash{}{0pt}%
\pgfpathmoveto{\pgfqpoint{2.196663in}{1.378800in}}%
\pgfpathlineto{\pgfqpoint{2.232094in}{1.382445in}}%
\pgfpathlineto{\pgfqpoint{2.229609in}{1.405210in}}%
\pgfpathlineto{\pgfqpoint{2.252299in}{1.407809in}}%
\pgfpathlineto{\pgfqpoint{2.254932in}{1.384910in}}%
\pgfpathlineto{\pgfqpoint{2.272022in}{1.386856in}}%
\pgfpathlineto{\pgfqpoint{2.275068in}{1.363856in}}%
\pgfpathlineto{\pgfqpoint{2.269373in}{1.363358in}}%
\pgfpathlineto{\pgfqpoint{2.223806in}{1.358262in}}%
\pgfpathlineto{\pgfqpoint{2.201376in}{1.356174in}}%
\pgfpathlineto{\pgfqpoint{2.200252in}{1.367456in}}%
\pgfpathlineto{\pgfqpoint{2.193342in}{1.366778in}}%
\pgfpathlineto{\pgfqpoint{2.196663in}{1.378800in}}%
\pgfpathclose%
\pgfusepath{fill}%
\end{pgfscope}%
\begin{pgfscope}%
\pgfpathrectangle{\pgfqpoint{0.100000in}{0.100000in}}{\pgfqpoint{3.007045in}{1.925000in}}%
\pgfusepath{clip}%
\pgfsetbuttcap%
\pgfsetmiterjoin%
\definecolor{currentfill}{rgb}{0.357647,0.720953,0.415071}%
\pgfsetfillcolor{currentfill}%
\pgfsetlinewidth{0.000000pt}%
\definecolor{currentstroke}{rgb}{0.000000,0.000000,0.000000}%
\pgfsetstrokecolor{currentstroke}%
\pgfsetstrokeopacity{0.000000}%
\pgfsetdash{}{0pt}%
\pgfpathmoveto{\pgfqpoint{1.102971in}{1.695747in}}%
\pgfpathlineto{\pgfqpoint{1.119016in}{1.693144in}}%
\pgfpathlineto{\pgfqpoint{1.113428in}{1.662137in}}%
\pgfpathlineto{\pgfqpoint{1.102024in}{1.663968in}}%
\pgfpathlineto{\pgfqpoint{1.073596in}{1.668769in}}%
\pgfpathlineto{\pgfqpoint{1.067509in}{1.667877in}}%
\pgfpathlineto{\pgfqpoint{1.037941in}{1.673119in}}%
\pgfpathlineto{\pgfqpoint{1.036887in}{1.681674in}}%
\pgfpathlineto{\pgfqpoint{1.040593in}{1.686492in}}%
\pgfpathlineto{\pgfqpoint{1.037077in}{1.690684in}}%
\pgfpathlineto{\pgfqpoint{1.033247in}{1.689867in}}%
\pgfpathlineto{\pgfqpoint{1.030764in}{1.699834in}}%
\pgfpathlineto{\pgfqpoint{1.027516in}{1.705647in}}%
\pgfpathlineto{\pgfqpoint{1.025028in}{1.717016in}}%
\pgfpathlineto{\pgfqpoint{1.019449in}{1.720433in}}%
\pgfpathlineto{\pgfqpoint{1.019738in}{1.725093in}}%
\pgfpathlineto{\pgfqpoint{1.025046in}{1.724117in}}%
\pgfpathlineto{\pgfqpoint{1.026102in}{1.729804in}}%
\pgfpathlineto{\pgfqpoint{1.043222in}{1.727442in}}%
\pgfpathlineto{\pgfqpoint{1.053818in}{1.723481in}}%
\pgfpathlineto{\pgfqpoint{1.059670in}{1.716732in}}%
\pgfpathlineto{\pgfqpoint{1.066660in}{1.704559in}}%
\pgfpathlineto{\pgfqpoint{1.070809in}{1.700125in}}%
\pgfpathlineto{\pgfqpoint{1.093536in}{1.701272in}}%
\pgfpathlineto{\pgfqpoint{1.099659in}{1.696257in}}%
\pgfpathlineto{\pgfqpoint{1.102971in}{1.695747in}}%
\pgfpathclose%
\pgfusepath{fill}%
\end{pgfscope}%
\begin{pgfscope}%
\pgfpathrectangle{\pgfqpoint{0.100000in}{0.100000in}}{\pgfqpoint{3.007045in}{1.925000in}}%
\pgfusepath{clip}%
\pgfsetbuttcap%
\pgfsetmiterjoin%
\definecolor{currentfill}{rgb}{0.723122,0.889612,0.697178}%
\pgfsetfillcolor{currentfill}%
\pgfsetlinewidth{0.000000pt}%
\definecolor{currentstroke}{rgb}{0.000000,0.000000,0.000000}%
\pgfsetstrokecolor{currentstroke}%
\pgfsetstrokeopacity{0.000000}%
\pgfsetdash{}{0pt}%
\pgfpathmoveto{\pgfqpoint{2.802639in}{1.349500in}}%
\pgfpathlineto{\pgfqpoint{2.804907in}{1.347247in}}%
\pgfpathlineto{\pgfqpoint{2.799417in}{1.334517in}}%
\pgfpathlineto{\pgfqpoint{2.789776in}{1.325534in}}%
\pgfpathlineto{\pgfqpoint{2.791685in}{1.322963in}}%
\pgfpathlineto{\pgfqpoint{2.777561in}{1.327350in}}%
\pgfpathlineto{\pgfqpoint{2.769299in}{1.335490in}}%
\pgfpathlineto{\pgfqpoint{2.764535in}{1.335742in}}%
\pgfpathlineto{\pgfqpoint{2.762277in}{1.343620in}}%
\pgfpathlineto{\pgfqpoint{2.755066in}{1.346497in}}%
\pgfpathlineto{\pgfqpoint{2.753885in}{1.357653in}}%
\pgfpathlineto{\pgfqpoint{2.756767in}{1.359194in}}%
\pgfpathlineto{\pgfqpoint{2.758253in}{1.365099in}}%
\pgfpathlineto{\pgfqpoint{2.753651in}{1.370464in}}%
\pgfpathlineto{\pgfqpoint{2.762147in}{1.386718in}}%
\pgfpathlineto{\pgfqpoint{2.763125in}{1.394125in}}%
\pgfpathlineto{\pgfqpoint{2.768627in}{1.400117in}}%
\pgfpathlineto{\pgfqpoint{2.794033in}{1.391447in}}%
\pgfpathlineto{\pgfqpoint{2.803517in}{1.405975in}}%
\pgfpathlineto{\pgfqpoint{2.805188in}{1.401167in}}%
\pgfpathlineto{\pgfqpoint{2.810869in}{1.395086in}}%
\pgfpathlineto{\pgfqpoint{2.812809in}{1.384940in}}%
\pgfpathlineto{\pgfqpoint{2.810062in}{1.361931in}}%
\pgfpathlineto{\pgfqpoint{2.803807in}{1.360108in}}%
\pgfpathlineto{\pgfqpoint{2.802639in}{1.349500in}}%
\pgfpathclose%
\pgfusepath{fill}%
\end{pgfscope}%
\begin{pgfscope}%
\pgfpathrectangle{\pgfqpoint{0.100000in}{0.100000in}}{\pgfqpoint{3.007045in}{1.925000in}}%
\pgfusepath{clip}%
\pgfsetbuttcap%
\pgfsetmiterjoin%
\definecolor{currentfill}{rgb}{0.601615,0.837093,0.583376}%
\pgfsetfillcolor{currentfill}%
\pgfsetlinewidth{0.000000pt}%
\definecolor{currentstroke}{rgb}{0.000000,0.000000,0.000000}%
\pgfsetstrokecolor{currentstroke}%
\pgfsetstrokeopacity{0.000000}%
\pgfsetdash{}{0pt}%
\pgfpathmoveto{\pgfqpoint{2.427990in}{0.752079in}}%
\pgfpathlineto{\pgfqpoint{2.441826in}{0.762195in}}%
\pgfpathlineto{\pgfqpoint{2.441164in}{0.768340in}}%
\pgfpathlineto{\pgfqpoint{2.436633in}{0.771428in}}%
\pgfpathlineto{\pgfqpoint{2.433317in}{0.776989in}}%
\pgfpathlineto{\pgfqpoint{2.432968in}{0.782406in}}%
\pgfpathlineto{\pgfqpoint{2.437607in}{0.788428in}}%
\pgfpathlineto{\pgfqpoint{2.441806in}{0.790953in}}%
\pgfpathlineto{\pgfqpoint{2.443204in}{0.795028in}}%
\pgfpathlineto{\pgfqpoint{2.448640in}{0.798533in}}%
\pgfpathlineto{\pgfqpoint{2.449695in}{0.794403in}}%
\pgfpathlineto{\pgfqpoint{2.454423in}{0.792007in}}%
\pgfpathlineto{\pgfqpoint{2.461072in}{0.785247in}}%
\pgfpathlineto{\pgfqpoint{2.464577in}{0.779288in}}%
\pgfpathlineto{\pgfqpoint{2.464508in}{0.773648in}}%
\pgfpathlineto{\pgfqpoint{2.470222in}{0.769671in}}%
\pgfpathlineto{\pgfqpoint{2.467299in}{0.764175in}}%
\pgfpathlineto{\pgfqpoint{2.469397in}{0.761860in}}%
\pgfpathlineto{\pgfqpoint{2.463033in}{0.756433in}}%
\pgfpathlineto{\pgfqpoint{2.460691in}{0.751658in}}%
\pgfpathlineto{\pgfqpoint{2.457386in}{0.739684in}}%
\pgfpathlineto{\pgfqpoint{2.458380in}{0.737732in}}%
\pgfpathlineto{\pgfqpoint{2.451822in}{0.729156in}}%
\pgfpathlineto{\pgfqpoint{2.444853in}{0.725072in}}%
\pgfpathlineto{\pgfqpoint{2.434230in}{0.742118in}}%
\pgfpathlineto{\pgfqpoint{2.427990in}{0.752079in}}%
\pgfpathclose%
\pgfusepath{fill}%
\end{pgfscope}%
\begin{pgfscope}%
\pgfpathrectangle{\pgfqpoint{0.100000in}{0.100000in}}{\pgfqpoint{3.007045in}{1.925000in}}%
\pgfusepath{clip}%
\pgfsetbuttcap%
\pgfsetmiterjoin%
\definecolor{currentfill}{rgb}{0.662345,0.864022,0.638001}%
\pgfsetfillcolor{currentfill}%
\pgfsetlinewidth{0.000000pt}%
\definecolor{currentstroke}{rgb}{0.000000,0.000000,0.000000}%
\pgfsetstrokecolor{currentstroke}%
\pgfsetstrokeopacity{0.000000}%
\pgfsetdash{}{0pt}%
\pgfpathmoveto{\pgfqpoint{2.113075in}{0.848655in}}%
\pgfpathlineto{\pgfqpoint{2.110261in}{0.847530in}}%
\pgfpathlineto{\pgfqpoint{2.110699in}{0.840860in}}%
\pgfpathlineto{\pgfqpoint{2.103193in}{0.839402in}}%
\pgfpathlineto{\pgfqpoint{2.085917in}{0.838289in}}%
\pgfpathlineto{\pgfqpoint{2.085442in}{0.845945in}}%
\pgfpathlineto{\pgfqpoint{2.082243in}{0.851523in}}%
\pgfpathlineto{\pgfqpoint{2.081112in}{0.868758in}}%
\pgfpathlineto{\pgfqpoint{2.078219in}{0.868579in}}%
\pgfpathlineto{\pgfqpoint{2.077608in}{0.878428in}}%
\pgfpathlineto{\pgfqpoint{2.085939in}{0.878995in}}%
\pgfpathlineto{\pgfqpoint{2.106178in}{0.880380in}}%
\pgfpathlineto{\pgfqpoint{2.106779in}{0.871379in}}%
\pgfpathlineto{\pgfqpoint{2.110591in}{0.871650in}}%
\pgfpathlineto{\pgfqpoint{2.112919in}{0.865058in}}%
\pgfpathlineto{\pgfqpoint{2.113075in}{0.848655in}}%
\pgfpathclose%
\pgfusepath{fill}%
\end{pgfscope}%
\begin{pgfscope}%
\pgfpathrectangle{\pgfqpoint{0.100000in}{0.100000in}}{\pgfqpoint{3.007045in}{1.925000in}}%
\pgfusepath{clip}%
\pgfsetbuttcap%
\pgfsetmiterjoin%
\definecolor{currentfill}{rgb}{0.546251,0.811257,0.537855}%
\pgfsetfillcolor{currentfill}%
\pgfsetlinewidth{0.000000pt}%
\definecolor{currentstroke}{rgb}{0.000000,0.000000,0.000000}%
\pgfsetstrokecolor{currentstroke}%
\pgfsetstrokeopacity{0.000000}%
\pgfsetdash{}{0pt}%
\pgfpathmoveto{\pgfqpoint{2.179696in}{0.840159in}}%
\pgfpathlineto{\pgfqpoint{2.144570in}{0.838481in}}%
\pgfpathlineto{\pgfqpoint{2.144807in}{0.847920in}}%
\pgfpathlineto{\pgfqpoint{2.145830in}{0.876417in}}%
\pgfpathlineto{\pgfqpoint{2.139583in}{0.883616in}}%
\pgfpathlineto{\pgfqpoint{2.151351in}{0.884373in}}%
\pgfpathlineto{\pgfqpoint{2.150055in}{0.903459in}}%
\pgfpathlineto{\pgfqpoint{2.147248in}{0.909640in}}%
\pgfpathlineto{\pgfqpoint{2.149635in}{0.914318in}}%
\pgfpathlineto{\pgfqpoint{2.163243in}{0.916924in}}%
\pgfpathlineto{\pgfqpoint{2.167138in}{0.916537in}}%
\pgfpathlineto{\pgfqpoint{2.171317in}{0.911958in}}%
\pgfpathlineto{\pgfqpoint{2.184378in}{0.916884in}}%
\pgfpathlineto{\pgfqpoint{2.189759in}{0.915550in}}%
\pgfpathlineto{\pgfqpoint{2.190733in}{0.912451in}}%
\pgfpathlineto{\pgfqpoint{2.192088in}{0.890541in}}%
\pgfpathlineto{\pgfqpoint{2.193121in}{0.887325in}}%
\pgfpathlineto{\pgfqpoint{2.193931in}{0.875265in}}%
\pgfpathlineto{\pgfqpoint{2.189352in}{0.871261in}}%
\pgfpathlineto{\pgfqpoint{2.182587in}{0.873244in}}%
\pgfpathlineto{\pgfqpoint{2.183086in}{0.867245in}}%
\pgfpathlineto{\pgfqpoint{2.178278in}{0.857446in}}%
\pgfpathlineto{\pgfqpoint{2.179696in}{0.840159in}}%
\pgfpathclose%
\pgfusepath{fill}%
\end{pgfscope}%
\begin{pgfscope}%
\pgfpathrectangle{\pgfqpoint{0.100000in}{0.100000in}}{\pgfqpoint{3.007045in}{1.925000in}}%
\pgfusepath{clip}%
\pgfsetbuttcap%
\pgfsetmiterjoin%
\definecolor{currentfill}{rgb}{0.634295,0.852211,0.610688}%
\pgfsetfillcolor{currentfill}%
\pgfsetlinewidth{0.000000pt}%
\definecolor{currentstroke}{rgb}{0.000000,0.000000,0.000000}%
\pgfsetstrokecolor{currentstroke}%
\pgfsetstrokeopacity{0.000000}%
\pgfsetdash{}{0pt}%
\pgfpathmoveto{\pgfqpoint{1.662731in}{0.409312in}}%
\pgfpathlineto{\pgfqpoint{1.664160in}{0.412781in}}%
\pgfpathlineto{\pgfqpoint{1.665098in}{0.416805in}}%
\pgfpathlineto{\pgfqpoint{1.670154in}{0.417855in}}%
\pgfpathlineto{\pgfqpoint{1.678639in}{0.424954in}}%
\pgfpathlineto{\pgfqpoint{1.680357in}{0.422649in}}%
\pgfpathlineto{\pgfqpoint{1.670606in}{0.415983in}}%
\pgfpathlineto{\pgfqpoint{1.662731in}{0.409312in}}%
\pgfpathclose%
\pgfusepath{fill}%
\end{pgfscope}%
\begin{pgfscope}%
\pgfpathrectangle{\pgfqpoint{0.100000in}{0.100000in}}{\pgfqpoint{3.007045in}{1.925000in}}%
\pgfusepath{clip}%
\pgfsetbuttcap%
\pgfsetmiterjoin%
\definecolor{currentfill}{rgb}{0.634295,0.852211,0.610688}%
\pgfsetfillcolor{currentfill}%
\pgfsetlinewidth{0.000000pt}%
\definecolor{currentstroke}{rgb}{0.000000,0.000000,0.000000}%
\pgfsetstrokecolor{currentstroke}%
\pgfsetstrokeopacity{0.000000}%
\pgfsetdash{}{0pt}%
\pgfpathmoveto{\pgfqpoint{1.694633in}{0.447645in}}%
\pgfpathlineto{\pgfqpoint{1.685694in}{0.443169in}}%
\pgfpathlineto{\pgfqpoint{1.680452in}{0.447923in}}%
\pgfpathlineto{\pgfqpoint{1.679297in}{0.453438in}}%
\pgfpathlineto{\pgfqpoint{1.674645in}{0.451462in}}%
\pgfpathlineto{\pgfqpoint{1.678471in}{0.443997in}}%
\pgfpathlineto{\pgfqpoint{1.690125in}{0.434156in}}%
\pgfpathlineto{\pgfqpoint{1.679523in}{0.428823in}}%
\pgfpathlineto{\pgfqpoint{1.676716in}{0.426224in}}%
\pgfpathlineto{\pgfqpoint{1.667536in}{0.431953in}}%
\pgfpathlineto{\pgfqpoint{1.661661in}{0.438978in}}%
\pgfpathlineto{\pgfqpoint{1.655495in}{0.438831in}}%
\pgfpathlineto{\pgfqpoint{1.651860in}{0.442260in}}%
\pgfpathlineto{\pgfqpoint{1.645296in}{0.442654in}}%
\pgfpathlineto{\pgfqpoint{1.632363in}{0.431928in}}%
\pgfpathlineto{\pgfqpoint{1.622101in}{0.436398in}}%
\pgfpathlineto{\pgfqpoint{1.621822in}{0.440642in}}%
\pgfpathlineto{\pgfqpoint{1.613769in}{0.442713in}}%
\pgfpathlineto{\pgfqpoint{1.609200in}{0.451305in}}%
\pgfpathlineto{\pgfqpoint{1.621451in}{0.460671in}}%
\pgfpathlineto{\pgfqpoint{1.611272in}{0.473936in}}%
\pgfpathlineto{\pgfqpoint{1.619767in}{0.480550in}}%
\pgfpathlineto{\pgfqpoint{1.641992in}{0.498317in}}%
\pgfpathlineto{\pgfqpoint{1.644405in}{0.511356in}}%
\pgfpathlineto{\pgfqpoint{1.648009in}{0.514439in}}%
\pgfpathlineto{\pgfqpoint{1.663618in}{0.514403in}}%
\pgfpathlineto{\pgfqpoint{1.668257in}{0.510605in}}%
\pgfpathlineto{\pgfqpoint{1.681469in}{0.494214in}}%
\pgfpathlineto{\pgfqpoint{1.676720in}{0.488443in}}%
\pgfpathlineto{\pgfqpoint{1.695745in}{0.469116in}}%
\pgfpathlineto{\pgfqpoint{1.694633in}{0.447645in}}%
\pgfpathclose%
\pgfusepath{fill}%
\end{pgfscope}%
\begin{pgfscope}%
\pgfpathrectangle{\pgfqpoint{0.100000in}{0.100000in}}{\pgfqpoint{3.007045in}{1.925000in}}%
\pgfusepath{clip}%
\pgfsetbuttcap%
\pgfsetmiterjoin%
\definecolor{currentfill}{rgb}{0.557324,0.816424,0.546959}%
\pgfsetfillcolor{currentfill}%
\pgfsetlinewidth{0.000000pt}%
\definecolor{currentstroke}{rgb}{0.000000,0.000000,0.000000}%
\pgfsetstrokecolor{currentstroke}%
\pgfsetstrokeopacity{0.000000}%
\pgfsetdash{}{0pt}%
\pgfpathmoveto{\pgfqpoint{1.985200in}{0.741858in}}%
\pgfpathlineto{\pgfqpoint{1.970228in}{0.741380in}}%
\pgfpathlineto{\pgfqpoint{1.968862in}{0.741315in}}%
\pgfpathlineto{\pgfqpoint{1.968042in}{0.766663in}}%
\pgfpathlineto{\pgfqpoint{1.938599in}{0.766214in}}%
\pgfpathlineto{\pgfqpoint{1.939283in}{0.759647in}}%
\pgfpathlineto{\pgfqpoint{1.931371in}{0.753155in}}%
\pgfpathlineto{\pgfqpoint{1.930951in}{0.750384in}}%
\pgfpathlineto{\pgfqpoint{1.921911in}{0.757915in}}%
\pgfpathlineto{\pgfqpoint{1.918474in}{0.757862in}}%
\pgfpathlineto{\pgfqpoint{1.917305in}{0.764130in}}%
\pgfpathlineto{\pgfqpoint{1.920677in}{0.772374in}}%
\pgfpathlineto{\pgfqpoint{1.918485in}{0.794539in}}%
\pgfpathlineto{\pgfqpoint{1.911801in}{0.803402in}}%
\pgfpathlineto{\pgfqpoint{1.910795in}{0.808969in}}%
\pgfpathlineto{\pgfqpoint{1.918185in}{0.809351in}}%
\pgfpathlineto{\pgfqpoint{1.923831in}{0.809696in}}%
\pgfpathlineto{\pgfqpoint{1.924611in}{0.820832in}}%
\pgfpathlineto{\pgfqpoint{1.924362in}{0.838176in}}%
\pgfpathlineto{\pgfqpoint{1.951736in}{0.838426in}}%
\pgfpathlineto{\pgfqpoint{1.951971in}{0.821985in}}%
\pgfpathlineto{\pgfqpoint{1.961337in}{0.821725in}}%
\pgfpathlineto{\pgfqpoint{1.962924in}{0.816378in}}%
\pgfpathlineto{\pgfqpoint{1.966388in}{0.815925in}}%
\pgfpathlineto{\pgfqpoint{1.967563in}{0.808897in}}%
\pgfpathlineto{\pgfqpoint{1.974688in}{0.804739in}}%
\pgfpathlineto{\pgfqpoint{1.984377in}{0.804898in}}%
\pgfpathlineto{\pgfqpoint{1.980611in}{0.811051in}}%
\pgfpathlineto{\pgfqpoint{1.984723in}{0.815682in}}%
\pgfpathlineto{\pgfqpoint{1.993621in}{0.816052in}}%
\pgfpathlineto{\pgfqpoint{1.998299in}{0.813452in}}%
\pgfpathlineto{\pgfqpoint{1.996950in}{0.809911in}}%
\pgfpathlineto{\pgfqpoint{1.986772in}{0.806255in}}%
\pgfpathlineto{\pgfqpoint{1.991139in}{0.803150in}}%
\pgfpathlineto{\pgfqpoint{1.988015in}{0.799119in}}%
\pgfpathlineto{\pgfqpoint{1.992182in}{0.792843in}}%
\pgfpathlineto{\pgfqpoint{1.984613in}{0.793228in}}%
\pgfpathlineto{\pgfqpoint{1.980623in}{0.785127in}}%
\pgfpathlineto{\pgfqpoint{1.985316in}{0.780979in}}%
\pgfpathlineto{\pgfqpoint{1.980431in}{0.778978in}}%
\pgfpathlineto{\pgfqpoint{1.985442in}{0.766767in}}%
\pgfpathlineto{\pgfqpoint{1.985739in}{0.761415in}}%
\pgfpathlineto{\pgfqpoint{1.988710in}{0.756370in}}%
\pgfpathlineto{\pgfqpoint{1.988976in}{0.750911in}}%
\pgfpathlineto{\pgfqpoint{1.985200in}{0.741858in}}%
\pgfpathclose%
\pgfusepath{fill}%
\end{pgfscope}%
\begin{pgfscope}%
\pgfpathrectangle{\pgfqpoint{0.100000in}{0.100000in}}{\pgfqpoint{3.007045in}{1.925000in}}%
\pgfusepath{clip}%
\pgfsetbuttcap%
\pgfsetmiterjoin%
\definecolor{currentfill}{rgb}{0.474279,0.777670,0.478677}%
\pgfsetfillcolor{currentfill}%
\pgfsetlinewidth{0.000000pt}%
\definecolor{currentstroke}{rgb}{0.000000,0.000000,0.000000}%
\pgfsetstrokecolor{currentstroke}%
\pgfsetstrokeopacity{0.000000}%
\pgfsetdash{}{0pt}%
\pgfpathmoveto{\pgfqpoint{2.793875in}{1.572201in}}%
\pgfpathlineto{\pgfqpoint{2.763396in}{1.560589in}}%
\pgfpathlineto{\pgfqpoint{2.763075in}{1.564033in}}%
\pgfpathlineto{\pgfqpoint{2.755118in}{1.563156in}}%
\pgfpathlineto{\pgfqpoint{2.747693in}{1.569195in}}%
\pgfpathlineto{\pgfqpoint{2.750775in}{1.572935in}}%
\pgfpathlineto{\pgfqpoint{2.747323in}{1.582366in}}%
\pgfpathlineto{\pgfqpoint{2.735921in}{1.578293in}}%
\pgfpathlineto{\pgfqpoint{2.735305in}{1.582994in}}%
\pgfpathlineto{\pgfqpoint{2.718437in}{1.632224in}}%
\pgfpathlineto{\pgfqpoint{2.713963in}{1.634075in}}%
\pgfpathlineto{\pgfqpoint{2.740305in}{1.639901in}}%
\pgfpathlineto{\pgfqpoint{2.776771in}{1.649915in}}%
\pgfpathlineto{\pgfqpoint{2.778407in}{1.644079in}}%
\pgfpathlineto{\pgfqpoint{2.777816in}{1.638392in}}%
\pgfpathlineto{\pgfqpoint{2.780372in}{1.636703in}}%
\pgfpathlineto{\pgfqpoint{2.780749in}{1.624141in}}%
\pgfpathlineto{\pgfqpoint{2.786419in}{1.617481in}}%
\pgfpathlineto{\pgfqpoint{2.787188in}{1.608949in}}%
\pgfpathlineto{\pgfqpoint{2.789757in}{1.602899in}}%
\pgfpathlineto{\pgfqpoint{2.787394in}{1.596652in}}%
\pgfpathlineto{\pgfqpoint{2.787488in}{1.586749in}}%
\pgfpathlineto{\pgfqpoint{2.793875in}{1.572201in}}%
\pgfpathclose%
\pgfusepath{fill}%
\end{pgfscope}%
\begin{pgfscope}%
\pgfpathrectangle{\pgfqpoint{0.100000in}{0.100000in}}{\pgfqpoint{3.007045in}{1.925000in}}%
\pgfusepath{clip}%
\pgfsetbuttcap%
\pgfsetmiterjoin%
\definecolor{currentfill}{rgb}{0.618224,0.844844,0.597032}%
\pgfsetfillcolor{currentfill}%
\pgfsetlinewidth{0.000000pt}%
\definecolor{currentstroke}{rgb}{0.000000,0.000000,0.000000}%
\pgfsetstrokecolor{currentstroke}%
\pgfsetstrokeopacity{0.000000}%
\pgfsetdash{}{0pt}%
\pgfpathmoveto{\pgfqpoint{1.415290in}{0.621980in}}%
\pgfpathlineto{\pgfqpoint{1.386813in}{0.624056in}}%
\pgfpathlineto{\pgfqpoint{1.391806in}{0.690752in}}%
\pgfpathlineto{\pgfqpoint{1.396350in}{0.690465in}}%
\pgfpathlineto{\pgfqpoint{1.398836in}{0.719387in}}%
\pgfpathlineto{\pgfqpoint{1.370186in}{0.721452in}}%
\pgfpathlineto{\pgfqpoint{1.372150in}{0.750256in}}%
\pgfpathlineto{\pgfqpoint{1.379483in}{0.749728in}}%
\pgfpathlineto{\pgfqpoint{1.429539in}{0.746331in}}%
\pgfpathlineto{\pgfqpoint{1.427496in}{0.717282in}}%
\pgfpathlineto{\pgfqpoint{1.425000in}{0.688410in}}%
\pgfpathlineto{\pgfqpoint{1.420501in}{0.688714in}}%
\pgfpathlineto{\pgfqpoint{1.417736in}{0.651743in}}%
\pgfpathlineto{\pgfqpoint{1.415290in}{0.621980in}}%
\pgfpathclose%
\pgfusepath{fill}%
\end{pgfscope}%
\begin{pgfscope}%
\pgfpathrectangle{\pgfqpoint{0.100000in}{0.100000in}}{\pgfqpoint{3.007045in}{1.925000in}}%
\pgfusepath{clip}%
\pgfsetbuttcap%
\pgfsetmiterjoin%
\definecolor{currentfill}{rgb}{0.886505,0.956171,0.866128}%
\pgfsetfillcolor{currentfill}%
\pgfsetlinewidth{0.000000pt}%
\definecolor{currentstroke}{rgb}{0.000000,0.000000,0.000000}%
\pgfsetstrokecolor{currentstroke}%
\pgfsetstrokeopacity{0.000000}%
\pgfsetdash{}{0pt}%
\pgfpathmoveto{\pgfqpoint{0.713739in}{0.318670in}}%
\pgfpathlineto{\pgfqpoint{0.713522in}{0.320731in}}%
\pgfpathlineto{\pgfqpoint{0.715798in}{0.320350in}}%
\pgfpathlineto{\pgfqpoint{0.717849in}{0.316514in}}%
\pgfpathlineto{\pgfqpoint{0.713739in}{0.318670in}}%
\pgfpathclose%
\pgfusepath{fill}%
\end{pgfscope}%
\begin{pgfscope}%
\pgfpathrectangle{\pgfqpoint{0.100000in}{0.100000in}}{\pgfqpoint{3.007045in}{1.925000in}}%
\pgfusepath{clip}%
\pgfsetbuttcap%
\pgfsetmiterjoin%
\definecolor{currentfill}{rgb}{0.886505,0.956171,0.866128}%
\pgfsetfillcolor{currentfill}%
\pgfsetlinewidth{0.000000pt}%
\definecolor{currentstroke}{rgb}{0.000000,0.000000,0.000000}%
\pgfsetstrokecolor{currentstroke}%
\pgfsetstrokeopacity{0.000000}%
\pgfsetdash{}{0pt}%
\pgfpathmoveto{\pgfqpoint{0.707133in}{0.321121in}}%
\pgfpathlineto{\pgfqpoint{0.708642in}{0.322633in}}%
\pgfpathlineto{\pgfqpoint{0.712152in}{0.322886in}}%
\pgfpathlineto{\pgfqpoint{0.711451in}{0.321404in}}%
\pgfpathlineto{\pgfqpoint{0.707133in}{0.321121in}}%
\pgfpathclose%
\pgfusepath{fill}%
\end{pgfscope}%
\begin{pgfscope}%
\pgfpathrectangle{\pgfqpoint{0.100000in}{0.100000in}}{\pgfqpoint{3.007045in}{1.925000in}}%
\pgfusepath{clip}%
\pgfsetbuttcap%
\pgfsetmiterjoin%
\definecolor{currentfill}{rgb}{0.886505,0.956171,0.866128}%
\pgfsetfillcolor{currentfill}%
\pgfsetlinewidth{0.000000pt}%
\definecolor{currentstroke}{rgb}{0.000000,0.000000,0.000000}%
\pgfsetstrokecolor{currentstroke}%
\pgfsetstrokeopacity{0.000000}%
\pgfsetdash{}{0pt}%
\pgfpathmoveto{\pgfqpoint{0.701239in}{0.348319in}}%
\pgfpathlineto{\pgfqpoint{0.697922in}{0.351049in}}%
\pgfpathlineto{\pgfqpoint{0.698993in}{0.352386in}}%
\pgfpathlineto{\pgfqpoint{0.701818in}{0.353160in}}%
\pgfpathlineto{\pgfqpoint{0.703476in}{0.352730in}}%
\pgfpathlineto{\pgfqpoint{0.704495in}{0.354085in}}%
\pgfpathlineto{\pgfqpoint{0.707132in}{0.354637in}}%
\pgfpathlineto{\pgfqpoint{0.708693in}{0.353375in}}%
\pgfpathlineto{\pgfqpoint{0.712359in}{0.352736in}}%
\pgfpathlineto{\pgfqpoint{0.711910in}{0.351797in}}%
\pgfpathlineto{\pgfqpoint{0.708260in}{0.352473in}}%
\pgfpathlineto{\pgfqpoint{0.704182in}{0.352191in}}%
\pgfpathlineto{\pgfqpoint{0.703882in}{0.349772in}}%
\pgfpathlineto{\pgfqpoint{0.701239in}{0.348319in}}%
\pgfpathclose%
\pgfusepath{fill}%
\end{pgfscope}%
\begin{pgfscope}%
\pgfpathrectangle{\pgfqpoint{0.100000in}{0.100000in}}{\pgfqpoint{3.007045in}{1.925000in}}%
\pgfusepath{clip}%
\pgfsetbuttcap%
\pgfsetmiterjoin%
\definecolor{currentfill}{rgb}{0.886505,0.956171,0.866128}%
\pgfsetfillcolor{currentfill}%
\pgfsetlinewidth{0.000000pt}%
\definecolor{currentstroke}{rgb}{0.000000,0.000000,0.000000}%
\pgfsetstrokecolor{currentstroke}%
\pgfsetstrokeopacity{0.000000}%
\pgfsetdash{}{0pt}%
\pgfpathmoveto{\pgfqpoint{0.758615in}{0.341736in}}%
\pgfpathlineto{\pgfqpoint{0.757407in}{0.339850in}}%
\pgfpathlineto{\pgfqpoint{0.759513in}{0.339536in}}%
\pgfpathlineto{\pgfqpoint{0.759838in}{0.338128in}}%
\pgfpathlineto{\pgfqpoint{0.761603in}{0.337668in}}%
\pgfpathlineto{\pgfqpoint{0.762040in}{0.334943in}}%
\pgfpathlineto{\pgfqpoint{0.760463in}{0.333422in}}%
\pgfpathlineto{\pgfqpoint{0.758461in}{0.333128in}}%
\pgfpathlineto{\pgfqpoint{0.753516in}{0.335984in}}%
\pgfpathlineto{\pgfqpoint{0.749366in}{0.335186in}}%
\pgfpathlineto{\pgfqpoint{0.746236in}{0.338305in}}%
\pgfpathlineto{\pgfqpoint{0.745834in}{0.340586in}}%
\pgfpathlineto{\pgfqpoint{0.747476in}{0.341739in}}%
\pgfpathlineto{\pgfqpoint{0.750424in}{0.341197in}}%
\pgfpathlineto{\pgfqpoint{0.750868in}{0.342943in}}%
\pgfpathlineto{\pgfqpoint{0.753106in}{0.339546in}}%
\pgfpathlineto{\pgfqpoint{0.754901in}{0.343295in}}%
\pgfpathlineto{\pgfqpoint{0.758150in}{0.343053in}}%
\pgfpathlineto{\pgfqpoint{0.759750in}{0.345315in}}%
\pgfpathlineto{\pgfqpoint{0.760957in}{0.344040in}}%
\pgfpathlineto{\pgfqpoint{0.760988in}{0.342549in}}%
\pgfpathlineto{\pgfqpoint{0.758615in}{0.341736in}}%
\pgfpathclose%
\pgfusepath{fill}%
\end{pgfscope}%
\begin{pgfscope}%
\pgfpathrectangle{\pgfqpoint{0.100000in}{0.100000in}}{\pgfqpoint{3.007045in}{1.925000in}}%
\pgfusepath{clip}%
\pgfsetbuttcap%
\pgfsetmiterjoin%
\definecolor{currentfill}{rgb}{0.886505,0.956171,0.866128}%
\pgfsetfillcolor{currentfill}%
\pgfsetlinewidth{0.000000pt}%
\definecolor{currentstroke}{rgb}{0.000000,0.000000,0.000000}%
\pgfsetstrokecolor{currentstroke}%
\pgfsetstrokeopacity{0.000000}%
\pgfsetdash{}{0pt}%
\pgfpathmoveto{\pgfqpoint{0.751166in}{0.331121in}}%
\pgfpathlineto{\pgfqpoint{0.751716in}{0.328502in}}%
\pgfpathlineto{\pgfqpoint{0.749924in}{0.328523in}}%
\pgfpathlineto{\pgfqpoint{0.747420in}{0.325126in}}%
\pgfpathlineto{\pgfqpoint{0.750723in}{0.323377in}}%
\pgfpathlineto{\pgfqpoint{0.746585in}{0.321465in}}%
\pgfpathlineto{\pgfqpoint{0.744840in}{0.322182in}}%
\pgfpathlineto{\pgfqpoint{0.743772in}{0.324350in}}%
\pgfpathlineto{\pgfqpoint{0.742472in}{0.323661in}}%
\pgfpathlineto{\pgfqpoint{0.742188in}{0.321974in}}%
\pgfpathlineto{\pgfqpoint{0.740508in}{0.321472in}}%
\pgfpathlineto{\pgfqpoint{0.734035in}{0.323447in}}%
\pgfpathlineto{\pgfqpoint{0.734346in}{0.321507in}}%
\pgfpathlineto{\pgfqpoint{0.732208in}{0.320942in}}%
\pgfpathlineto{\pgfqpoint{0.729299in}{0.323471in}}%
\pgfpathlineto{\pgfqpoint{0.727715in}{0.323803in}}%
\pgfpathlineto{\pgfqpoint{0.723753in}{0.321601in}}%
\pgfpathlineto{\pgfqpoint{0.719755in}{0.321321in}}%
\pgfpathlineto{\pgfqpoint{0.720327in}{0.323273in}}%
\pgfpathlineto{\pgfqpoint{0.724725in}{0.324272in}}%
\pgfpathlineto{\pgfqpoint{0.724952in}{0.327121in}}%
\pgfpathlineto{\pgfqpoint{0.722436in}{0.327416in}}%
\pgfpathlineto{\pgfqpoint{0.719318in}{0.325777in}}%
\pgfpathlineto{\pgfqpoint{0.718473in}{0.330811in}}%
\pgfpathlineto{\pgfqpoint{0.720682in}{0.334691in}}%
\pgfpathlineto{\pgfqpoint{0.719629in}{0.337498in}}%
\pgfpathlineto{\pgfqpoint{0.724542in}{0.340506in}}%
\pgfpathlineto{\pgfqpoint{0.728409in}{0.340545in}}%
\pgfpathlineto{\pgfqpoint{0.730761in}{0.339863in}}%
\pgfpathlineto{\pgfqpoint{0.732294in}{0.338142in}}%
\pgfpathlineto{\pgfqpoint{0.730507in}{0.333695in}}%
\pgfpathlineto{\pgfqpoint{0.733490in}{0.335333in}}%
\pgfpathlineto{\pgfqpoint{0.734831in}{0.339974in}}%
\pgfpathlineto{\pgfqpoint{0.736842in}{0.341018in}}%
\pgfpathlineto{\pgfqpoint{0.738720in}{0.340689in}}%
\pgfpathlineto{\pgfqpoint{0.738975in}{0.337460in}}%
\pgfpathlineto{\pgfqpoint{0.740912in}{0.337016in}}%
\pgfpathlineto{\pgfqpoint{0.741053in}{0.339842in}}%
\pgfpathlineto{\pgfqpoint{0.745035in}{0.337912in}}%
\pgfpathlineto{\pgfqpoint{0.746149in}{0.335843in}}%
\pgfpathlineto{\pgfqpoint{0.748088in}{0.333888in}}%
\pgfpathlineto{\pgfqpoint{0.746213in}{0.332218in}}%
\pgfpathlineto{\pgfqpoint{0.751350in}{0.331886in}}%
\pgfpathlineto{\pgfqpoint{0.751166in}{0.331121in}}%
\pgfpathclose%
\pgfusepath{fill}%
\end{pgfscope}%
\begin{pgfscope}%
\pgfpathrectangle{\pgfqpoint{0.100000in}{0.100000in}}{\pgfqpoint{3.007045in}{1.925000in}}%
\pgfusepath{clip}%
\pgfsetbuttcap%
\pgfsetmiterjoin%
\definecolor{currentfill}{rgb}{0.886505,0.956171,0.866128}%
\pgfsetfillcolor{currentfill}%
\pgfsetlinewidth{0.000000pt}%
\definecolor{currentstroke}{rgb}{0.000000,0.000000,0.000000}%
\pgfsetstrokecolor{currentstroke}%
\pgfsetstrokeopacity{0.000000}%
\pgfsetdash{}{0pt}%
\pgfpathmoveto{\pgfqpoint{0.757017in}{0.355194in}}%
\pgfpathlineto{\pgfqpoint{0.755520in}{0.355499in}}%
\pgfpathlineto{\pgfqpoint{0.753460in}{0.353984in}}%
\pgfpathlineto{\pgfqpoint{0.748969in}{0.353910in}}%
\pgfpathlineto{\pgfqpoint{0.747754in}{0.354827in}}%
\pgfpathlineto{\pgfqpoint{0.745252in}{0.353754in}}%
\pgfpathlineto{\pgfqpoint{0.742936in}{0.353492in}}%
\pgfpathlineto{\pgfqpoint{0.741809in}{0.351678in}}%
\pgfpathlineto{\pgfqpoint{0.740287in}{0.352714in}}%
\pgfpathlineto{\pgfqpoint{0.739684in}{0.350603in}}%
\pgfpathlineto{\pgfqpoint{0.733752in}{0.349525in}}%
\pgfpathlineto{\pgfqpoint{0.731314in}{0.349960in}}%
\pgfpathlineto{\pgfqpoint{0.728847in}{0.351743in}}%
\pgfpathlineto{\pgfqpoint{0.725995in}{0.352347in}}%
\pgfpathlineto{\pgfqpoint{0.725319in}{0.351090in}}%
\pgfpathlineto{\pgfqpoint{0.721782in}{0.351198in}}%
\pgfpathlineto{\pgfqpoint{0.719594in}{0.350212in}}%
\pgfpathlineto{\pgfqpoint{0.719062in}{0.352698in}}%
\pgfpathlineto{\pgfqpoint{0.717117in}{0.351179in}}%
\pgfpathlineto{\pgfqpoint{0.714893in}{0.351866in}}%
\pgfpathlineto{\pgfqpoint{0.713962in}{0.350368in}}%
\pgfpathlineto{\pgfqpoint{0.712047in}{0.351741in}}%
\pgfpathlineto{\pgfqpoint{0.715201in}{0.352581in}}%
\pgfpathlineto{\pgfqpoint{0.717812in}{0.354377in}}%
\pgfpathlineto{\pgfqpoint{0.719341in}{0.353168in}}%
\pgfpathlineto{\pgfqpoint{0.722429in}{0.352828in}}%
\pgfpathlineto{\pgfqpoint{0.726759in}{0.358371in}}%
\pgfpathlineto{\pgfqpoint{0.728487in}{0.358327in}}%
\pgfpathlineto{\pgfqpoint{0.730107in}{0.357059in}}%
\pgfpathlineto{\pgfqpoint{0.731400in}{0.357708in}}%
\pgfpathlineto{\pgfqpoint{0.734357in}{0.355834in}}%
\pgfpathlineto{\pgfqpoint{0.737813in}{0.354884in}}%
\pgfpathlineto{\pgfqpoint{0.738983in}{0.355722in}}%
\pgfpathlineto{\pgfqpoint{0.740865in}{0.355528in}}%
\pgfpathlineto{\pgfqpoint{0.743549in}{0.359090in}}%
\pgfpathlineto{\pgfqpoint{0.748423in}{0.355398in}}%
\pgfpathlineto{\pgfqpoint{0.749665in}{0.357072in}}%
\pgfpathlineto{\pgfqpoint{0.753292in}{0.354357in}}%
\pgfpathlineto{\pgfqpoint{0.755001in}{0.356698in}}%
\pgfpathlineto{\pgfqpoint{0.757017in}{0.355194in}}%
\pgfpathclose%
\pgfusepath{fill}%
\end{pgfscope}%
\begin{pgfscope}%
\pgfpathrectangle{\pgfqpoint{0.100000in}{0.100000in}}{\pgfqpoint{3.007045in}{1.925000in}}%
\pgfusepath{clip}%
\pgfsetbuttcap%
\pgfsetmiterjoin%
\definecolor{currentfill}{rgb}{0.746498,0.899454,0.719938}%
\pgfsetfillcolor{currentfill}%
\pgfsetlinewidth{0.000000pt}%
\definecolor{currentstroke}{rgb}{0.000000,0.000000,0.000000}%
\pgfsetstrokecolor{currentstroke}%
\pgfsetstrokeopacity{0.000000}%
\pgfsetdash{}{0pt}%
\pgfpathmoveto{\pgfqpoint{1.020483in}{1.573439in}}%
\pgfpathlineto{\pgfqpoint{1.018543in}{1.560865in}}%
\pgfpathlineto{\pgfqpoint{1.015015in}{1.563234in}}%
\pgfpathlineto{\pgfqpoint{1.009746in}{1.577792in}}%
\pgfpathlineto{\pgfqpoint{1.006564in}{1.581699in}}%
\pgfpathlineto{\pgfqpoint{1.001653in}{1.579735in}}%
\pgfpathlineto{\pgfqpoint{0.987365in}{1.582354in}}%
\pgfpathlineto{\pgfqpoint{0.988238in}{1.587058in}}%
\pgfpathlineto{\pgfqpoint{0.975154in}{1.589701in}}%
\pgfpathlineto{\pgfqpoint{0.971016in}{1.593385in}}%
\pgfpathlineto{\pgfqpoint{0.972074in}{1.604894in}}%
\pgfpathlineto{\pgfqpoint{0.965491in}{1.606172in}}%
\pgfpathlineto{\pgfqpoint{0.966618in}{1.611886in}}%
\pgfpathlineto{\pgfqpoint{0.960968in}{1.612983in}}%
\pgfpathlineto{\pgfqpoint{0.964235in}{1.629990in}}%
\pgfpathlineto{\pgfqpoint{0.963238in}{1.636470in}}%
\pgfpathlineto{\pgfqpoint{0.959812in}{1.636263in}}%
\pgfpathlineto{\pgfqpoint{0.955356in}{1.640352in}}%
\pgfpathlineto{\pgfqpoint{0.957302in}{1.649942in}}%
\pgfpathlineto{\pgfqpoint{0.964271in}{1.652566in}}%
\pgfpathlineto{\pgfqpoint{0.966686in}{1.656395in}}%
\pgfpathlineto{\pgfqpoint{0.981446in}{1.653571in}}%
\pgfpathlineto{\pgfqpoint{0.988214in}{1.659319in}}%
\pgfpathlineto{\pgfqpoint{0.994682in}{1.658385in}}%
\pgfpathlineto{\pgfqpoint{0.997011in}{1.655494in}}%
\pgfpathlineto{\pgfqpoint{1.004464in}{1.652753in}}%
\pgfpathlineto{\pgfqpoint{1.010507in}{1.658210in}}%
\pgfpathlineto{\pgfqpoint{1.016570in}{1.660701in}}%
\pgfpathlineto{\pgfqpoint{1.019996in}{1.667445in}}%
\pgfpathlineto{\pgfqpoint{1.024131in}{1.670887in}}%
\pgfpathlineto{\pgfqpoint{1.025044in}{1.675119in}}%
\pgfpathlineto{\pgfqpoint{1.037941in}{1.673119in}}%
\pgfpathlineto{\pgfqpoint{1.067509in}{1.667877in}}%
\pgfpathlineto{\pgfqpoint{1.073596in}{1.668769in}}%
\pgfpathlineto{\pgfqpoint{1.102024in}{1.663968in}}%
\pgfpathlineto{\pgfqpoint{1.101102in}{1.658254in}}%
\pgfpathlineto{\pgfqpoint{1.102412in}{1.652202in}}%
\pgfpathlineto{\pgfqpoint{1.110938in}{1.650768in}}%
\pgfpathlineto{\pgfqpoint{1.109783in}{1.645289in}}%
\pgfpathlineto{\pgfqpoint{1.105971in}{1.645829in}}%
\pgfpathlineto{\pgfqpoint{1.101489in}{1.634717in}}%
\pgfpathlineto{\pgfqpoint{1.099750in}{1.623315in}}%
\pgfpathlineto{\pgfqpoint{1.094101in}{1.624217in}}%
\pgfpathlineto{\pgfqpoint{1.081871in}{1.620429in}}%
\pgfpathlineto{\pgfqpoint{1.079975in}{1.609049in}}%
\pgfpathlineto{\pgfqpoint{1.073953in}{1.610065in}}%
\pgfpathlineto{\pgfqpoint{1.072000in}{1.598524in}}%
\pgfpathlineto{\pgfqpoint{1.084238in}{1.596483in}}%
\pgfpathlineto{\pgfqpoint{1.082392in}{1.585392in}}%
\pgfpathlineto{\pgfqpoint{1.040422in}{1.591840in}}%
\pgfpathlineto{\pgfqpoint{1.024388in}{1.595276in}}%
\pgfpathlineto{\pgfqpoint{1.020483in}{1.573439in}}%
\pgfpathclose%
\pgfusepath{fill}%
\end{pgfscope}%
\begin{pgfscope}%
\pgfpathrectangle{\pgfqpoint{0.100000in}{0.100000in}}{\pgfqpoint{3.007045in}{1.925000in}}%
\pgfusepath{clip}%
\pgfsetbuttcap%
\pgfsetmiterjoin%
\definecolor{currentfill}{rgb}{0.551788,0.813841,0.542407}%
\pgfsetfillcolor{currentfill}%
\pgfsetlinewidth{0.000000pt}%
\definecolor{currentstroke}{rgb}{0.000000,0.000000,0.000000}%
\pgfsetstrokecolor{currentstroke}%
\pgfsetstrokeopacity{0.000000}%
\pgfsetdash{}{0pt}%
\pgfpathmoveto{\pgfqpoint{1.607157in}{1.202049in}}%
\pgfpathlineto{\pgfqpoint{1.595536in}{1.202468in}}%
\pgfpathlineto{\pgfqpoint{1.584352in}{1.202908in}}%
\pgfpathlineto{\pgfqpoint{1.586476in}{1.248770in}}%
\pgfpathlineto{\pgfqpoint{1.631190in}{1.247136in}}%
\pgfpathlineto{\pgfqpoint{1.630558in}{1.224197in}}%
\pgfpathlineto{\pgfqpoint{1.608008in}{1.224954in}}%
\pgfpathlineto{\pgfqpoint{1.607157in}{1.202049in}}%
\pgfpathclose%
\pgfusepath{fill}%
\end{pgfscope}%
\begin{pgfscope}%
\pgfpathrectangle{\pgfqpoint{0.100000in}{0.100000in}}{\pgfqpoint{3.007045in}{1.925000in}}%
\pgfusepath{clip}%
\pgfsetbuttcap%
\pgfsetmiterjoin%
\definecolor{currentfill}{rgb}{0.445490,0.764014,0.458131}%
\pgfsetfillcolor{currentfill}%
\pgfsetlinewidth{0.000000pt}%
\definecolor{currentstroke}{rgb}{0.000000,0.000000,0.000000}%
\pgfsetstrokecolor{currentstroke}%
\pgfsetstrokeopacity{0.000000}%
\pgfsetdash{}{0pt}%
\pgfpathmoveto{\pgfqpoint{0.596099in}{1.599150in}}%
\pgfpathlineto{\pgfqpoint{0.573938in}{1.605496in}}%
\pgfpathlineto{\pgfqpoint{0.530203in}{1.618670in}}%
\pgfpathlineto{\pgfqpoint{0.499509in}{1.627848in}}%
\pgfpathlineto{\pgfqpoint{0.504292in}{1.636653in}}%
\pgfpathlineto{\pgfqpoint{0.505425in}{1.642739in}}%
\pgfpathlineto{\pgfqpoint{0.512713in}{1.647749in}}%
\pgfpathlineto{\pgfqpoint{0.519173in}{1.654672in}}%
\pgfpathlineto{\pgfqpoint{0.521441in}{1.664987in}}%
\pgfpathlineto{\pgfqpoint{0.521156in}{1.673183in}}%
\pgfpathlineto{\pgfqpoint{0.523556in}{1.681083in}}%
\pgfpathlineto{\pgfqpoint{0.527922in}{1.686852in}}%
\pgfpathlineto{\pgfqpoint{0.530035in}{1.695577in}}%
\pgfpathlineto{\pgfqpoint{0.533873in}{1.701067in}}%
\pgfpathlineto{\pgfqpoint{0.563962in}{1.691905in}}%
\pgfpathlineto{\pgfqpoint{0.596303in}{1.682378in}}%
\pgfpathlineto{\pgfqpoint{0.594595in}{1.678719in}}%
\pgfpathlineto{\pgfqpoint{0.588656in}{1.658125in}}%
\pgfpathlineto{\pgfqpoint{0.605211in}{1.653481in}}%
\pgfpathlineto{\pgfqpoint{0.609918in}{1.648785in}}%
\pgfpathlineto{\pgfqpoint{0.608444in}{1.643504in}}%
\pgfpathlineto{\pgfqpoint{0.619536in}{1.640423in}}%
\pgfpathlineto{\pgfqpoint{0.613279in}{1.618283in}}%
\pgfpathlineto{\pgfqpoint{0.607867in}{1.619794in}}%
\pgfpathlineto{\pgfqpoint{0.603136in}{1.603204in}}%
\pgfpathlineto{\pgfqpoint{0.597590in}{1.604775in}}%
\pgfpathlineto{\pgfqpoint{0.596099in}{1.599150in}}%
\pgfpathclose%
\pgfusepath{fill}%
\end{pgfscope}%
\begin{pgfscope}%
\pgfpathrectangle{\pgfqpoint{0.100000in}{0.100000in}}{\pgfqpoint{3.007045in}{1.925000in}}%
\pgfusepath{clip}%
\pgfsetbuttcap%
\pgfsetmiterjoin%
\definecolor{currentfill}{rgb}{0.439216,0.760938,0.455056}%
\pgfsetfillcolor{currentfill}%
\pgfsetlinewidth{0.000000pt}%
\definecolor{currentstroke}{rgb}{0.000000,0.000000,0.000000}%
\pgfsetstrokecolor{currentstroke}%
\pgfsetstrokeopacity{0.000000}%
\pgfsetdash{}{0pt}%
\pgfpathmoveto{\pgfqpoint{1.446935in}{1.181960in}}%
\pgfpathlineto{\pgfqpoint{1.414739in}{1.184340in}}%
\pgfpathlineto{\pgfqpoint{1.414767in}{1.184726in}}%
\pgfpathlineto{\pgfqpoint{1.418635in}{1.235643in}}%
\pgfpathlineto{\pgfqpoint{1.455104in}{1.233038in}}%
\pgfpathlineto{\pgfqpoint{1.453477in}{1.210184in}}%
\pgfpathlineto{\pgfqpoint{1.449096in}{1.210470in}}%
\pgfpathlineto{\pgfqpoint{1.446935in}{1.181960in}}%
\pgfpathclose%
\pgfusepath{fill}%
\end{pgfscope}%
\begin{pgfscope}%
\pgfpathrectangle{\pgfqpoint{0.100000in}{0.100000in}}{\pgfqpoint{3.007045in}{1.925000in}}%
\pgfusepath{clip}%
\pgfsetbuttcap%
\pgfsetmiterjoin%
\definecolor{currentfill}{rgb}{0.200923,0.613010,0.321522}%
\pgfsetfillcolor{currentfill}%
\pgfsetlinewidth{0.000000pt}%
\definecolor{currentstroke}{rgb}{0.000000,0.000000,0.000000}%
\pgfsetstrokecolor{currentstroke}%
\pgfsetstrokeopacity{0.000000}%
\pgfsetdash{}{0pt}%
\pgfpathmoveto{\pgfqpoint{0.774157in}{0.904247in}}%
\pgfpathlineto{\pgfqpoint{0.778392in}{0.924882in}}%
\pgfpathlineto{\pgfqpoint{0.794371in}{1.003188in}}%
\pgfpathlineto{\pgfqpoint{0.798014in}{1.020963in}}%
\pgfpathlineto{\pgfqpoint{0.800968in}{1.030983in}}%
\pgfpathlineto{\pgfqpoint{0.800098in}{1.036517in}}%
\pgfpathlineto{\pgfqpoint{0.803014in}{1.039936in}}%
\pgfpathlineto{\pgfqpoint{0.807519in}{1.038421in}}%
\pgfpathlineto{\pgfqpoint{0.809522in}{1.041947in}}%
\pgfpathlineto{\pgfqpoint{0.818051in}{1.044246in}}%
\pgfpathlineto{\pgfqpoint{0.832280in}{1.046410in}}%
\pgfpathlineto{\pgfqpoint{0.839109in}{1.048947in}}%
\pgfpathlineto{\pgfqpoint{0.842909in}{1.052579in}}%
\pgfpathlineto{\pgfqpoint{0.842643in}{1.059879in}}%
\pgfpathlineto{\pgfqpoint{0.845589in}{1.070038in}}%
\pgfpathlineto{\pgfqpoint{0.854802in}{1.090140in}}%
\pgfpathlineto{\pgfqpoint{0.836260in}{1.093783in}}%
\pgfpathlineto{\pgfqpoint{0.843170in}{1.128894in}}%
\pgfpathlineto{\pgfqpoint{0.854337in}{1.126666in}}%
\pgfpathlineto{\pgfqpoint{0.870763in}{1.123009in}}%
\pgfpathlineto{\pgfqpoint{0.937370in}{1.110856in}}%
\pgfpathlineto{\pgfqpoint{0.958933in}{1.107125in}}%
\pgfpathlineto{\pgfqpoint{0.951769in}{1.095219in}}%
\pgfpathlineto{\pgfqpoint{0.944689in}{1.092843in}}%
\pgfpathlineto{\pgfqpoint{0.941624in}{1.085677in}}%
\pgfpathlineto{\pgfqpoint{0.932657in}{1.082202in}}%
\pgfpathlineto{\pgfqpoint{0.926450in}{1.083530in}}%
\pgfpathlineto{\pgfqpoint{0.923032in}{1.079857in}}%
\pgfpathlineto{\pgfqpoint{0.915481in}{1.078785in}}%
\pgfpathlineto{\pgfqpoint{0.947453in}{1.073147in}}%
\pgfpathlineto{\pgfqpoint{0.916273in}{0.895048in}}%
\pgfpathlineto{\pgfqpoint{0.911088in}{0.896114in}}%
\pgfpathlineto{\pgfqpoint{0.903247in}{0.903940in}}%
\pgfpathlineto{\pgfqpoint{0.896277in}{0.906751in}}%
\pgfpathlineto{\pgfqpoint{0.890889in}{0.911616in}}%
\pgfpathlineto{\pgfqpoint{0.883290in}{0.911640in}}%
\pgfpathlineto{\pgfqpoint{0.876178in}{0.916207in}}%
\pgfpathlineto{\pgfqpoint{0.868902in}{0.911853in}}%
\pgfpathlineto{\pgfqpoint{0.863233in}{0.897098in}}%
\pgfpathlineto{\pgfqpoint{0.875838in}{0.894759in}}%
\pgfpathlineto{\pgfqpoint{0.873121in}{0.885088in}}%
\pgfpathlineto{\pgfqpoint{0.860696in}{0.887387in}}%
\pgfpathlineto{\pgfqpoint{0.837638in}{0.894986in}}%
\pgfpathlineto{\pgfqpoint{0.829566in}{0.885369in}}%
\pgfpathlineto{\pgfqpoint{0.805861in}{0.897916in}}%
\pgfpathlineto{\pgfqpoint{0.774157in}{0.904247in}}%
\pgfpathclose%
\pgfusepath{fill}%
\end{pgfscope}%
\begin{pgfscope}%
\pgfpathrectangle{\pgfqpoint{0.100000in}{0.100000in}}{\pgfqpoint{3.007045in}{1.925000in}}%
\pgfusepath{clip}%
\pgfsetbuttcap%
\pgfsetmiterjoin%
\definecolor{currentfill}{rgb}{0.797924,0.920738,0.771642}%
\pgfsetfillcolor{currentfill}%
\pgfsetlinewidth{0.000000pt}%
\definecolor{currentstroke}{rgb}{0.000000,0.000000,0.000000}%
\pgfsetstrokecolor{currentstroke}%
\pgfsetstrokeopacity{0.000000}%
\pgfsetdash{}{0pt}%
\pgfpathmoveto{\pgfqpoint{1.330895in}{1.191503in}}%
\pgfpathlineto{\pgfqpoint{1.325845in}{1.145704in}}%
\pgfpathlineto{\pgfqpoint{1.308723in}{1.147516in}}%
\pgfpathlineto{\pgfqpoint{1.310620in}{1.164502in}}%
\pgfpathlineto{\pgfqpoint{1.245658in}{1.171547in}}%
\pgfpathlineto{\pgfqpoint{1.242494in}{1.143237in}}%
\pgfpathlineto{\pgfqpoint{1.209760in}{1.146848in}}%
\pgfpathlineto{\pgfqpoint{1.213877in}{1.153875in}}%
\pgfpathlineto{\pgfqpoint{1.213448in}{1.158253in}}%
\pgfpathlineto{\pgfqpoint{1.208767in}{1.163511in}}%
\pgfpathlineto{\pgfqpoint{1.200990in}{1.169461in}}%
\pgfpathlineto{\pgfqpoint{1.201403in}{1.171929in}}%
\pgfpathlineto{\pgfqpoint{1.202987in}{1.177221in}}%
\pgfpathlineto{\pgfqpoint{1.203807in}{1.188442in}}%
\pgfpathlineto{\pgfqpoint{1.216380in}{1.193870in}}%
\pgfpathlineto{\pgfqpoint{1.226581in}{1.206947in}}%
\pgfpathlineto{\pgfqpoint{1.220004in}{1.212317in}}%
\pgfpathlineto{\pgfqpoint{1.225797in}{1.217368in}}%
\pgfpathlineto{\pgfqpoint{1.232976in}{1.220918in}}%
\pgfpathlineto{\pgfqpoint{1.234437in}{1.230208in}}%
\pgfpathlineto{\pgfqpoint{1.237002in}{1.232833in}}%
\pgfpathlineto{\pgfqpoint{1.237983in}{1.247422in}}%
\pgfpathlineto{\pgfqpoint{1.267871in}{1.244110in}}%
\pgfpathlineto{\pgfqpoint{1.265951in}{1.226957in}}%
\pgfpathlineto{\pgfqpoint{1.311373in}{1.222190in}}%
\pgfpathlineto{\pgfqpoint{1.333705in}{1.220024in}}%
\pgfpathlineto{\pgfqpoint{1.330895in}{1.191503in}}%
\pgfpathclose%
\pgfusepath{fill}%
\end{pgfscope}%
\begin{pgfscope}%
\pgfpathrectangle{\pgfqpoint{0.100000in}{0.100000in}}{\pgfqpoint{3.007045in}{1.925000in}}%
\pgfusepath{clip}%
\pgfsetbuttcap%
\pgfsetmiterjoin%
\definecolor{currentfill}{rgb}{0.755848,0.903391,0.729043}%
\pgfsetfillcolor{currentfill}%
\pgfsetlinewidth{0.000000pt}%
\definecolor{currentstroke}{rgb}{0.000000,0.000000,0.000000}%
\pgfsetstrokecolor{currentstroke}%
\pgfsetstrokeopacity{0.000000}%
\pgfsetdash{}{0pt}%
\pgfpathmoveto{\pgfqpoint{1.514641in}{0.483384in}}%
\pgfpathlineto{\pgfqpoint{1.473940in}{0.485183in}}%
\pgfpathlineto{\pgfqpoint{1.475915in}{0.520991in}}%
\pgfpathlineto{\pgfqpoint{1.441878in}{0.523060in}}%
\pgfpathlineto{\pgfqpoint{1.444627in}{0.567276in}}%
\pgfpathlineto{\pgfqpoint{1.429685in}{0.568203in}}%
\pgfpathlineto{\pgfqpoint{1.431514in}{0.595996in}}%
\pgfpathlineto{\pgfqpoint{1.479753in}{0.593339in}}%
\pgfpathlineto{\pgfqpoint{1.478167in}{0.565413in}}%
\pgfpathlineto{\pgfqpoint{1.498977in}{0.564304in}}%
\pgfpathlineto{\pgfqpoint{1.498020in}{0.549906in}}%
\pgfpathlineto{\pgfqpoint{1.501924in}{0.549696in}}%
\pgfpathlineto{\pgfqpoint{1.501272in}{0.538689in}}%
\pgfpathlineto{\pgfqpoint{1.506378in}{0.538351in}}%
\pgfpathlineto{\pgfqpoint{1.505414in}{0.519674in}}%
\pgfpathlineto{\pgfqpoint{1.516499in}{0.519146in}}%
\pgfpathlineto{\pgfqpoint{1.514641in}{0.483384in}}%
\pgfpathclose%
\pgfusepath{fill}%
\end{pgfscope}%
\begin{pgfscope}%
\pgfpathrectangle{\pgfqpoint{0.100000in}{0.100000in}}{\pgfqpoint{3.007045in}{1.925000in}}%
\pgfusepath{clip}%
\pgfsetbuttcap%
\pgfsetmiterjoin%
\definecolor{currentfill}{rgb}{0.363922,0.724029,0.418147}%
\pgfsetfillcolor{currentfill}%
\pgfsetlinewidth{0.000000pt}%
\definecolor{currentstroke}{rgb}{0.000000,0.000000,0.000000}%
\pgfsetstrokecolor{currentstroke}%
\pgfsetstrokeopacity{0.000000}%
\pgfsetdash{}{0pt}%
\pgfpathmoveto{\pgfqpoint{0.989331in}{0.759721in}}%
\pgfpathlineto{\pgfqpoint{0.997334in}{0.810949in}}%
\pgfpathlineto{\pgfqpoint{1.054598in}{0.801868in}}%
\pgfpathlineto{\pgfqpoint{1.061462in}{0.800953in}}%
\pgfpathlineto{\pgfqpoint{1.062578in}{0.792196in}}%
\pgfpathlineto{\pgfqpoint{1.061362in}{0.785689in}}%
\pgfpathlineto{\pgfqpoint{1.064806in}{0.782977in}}%
\pgfpathlineto{\pgfqpoint{1.063832in}{0.777370in}}%
\pgfpathlineto{\pgfqpoint{1.065981in}{0.771962in}}%
\pgfpathlineto{\pgfqpoint{1.063758in}{0.762298in}}%
\pgfpathlineto{\pgfqpoint{1.070654in}{0.759703in}}%
\pgfpathlineto{\pgfqpoint{1.087735in}{0.757358in}}%
\pgfpathlineto{\pgfqpoint{1.080529in}{0.703340in}}%
\pgfpathlineto{\pgfqpoint{1.029591in}{0.710510in}}%
\pgfpathlineto{\pgfqpoint{1.025251in}{0.680923in}}%
\pgfpathlineto{\pgfqpoint{0.978016in}{0.688001in}}%
\pgfpathlineto{\pgfqpoint{0.989331in}{0.759721in}}%
\pgfpathclose%
\pgfusepath{fill}%
\end{pgfscope}%
\begin{pgfscope}%
\pgfpathrectangle{\pgfqpoint{0.100000in}{0.100000in}}{\pgfqpoint{3.007045in}{1.925000in}}%
\pgfusepath{clip}%
\pgfsetbuttcap%
\pgfsetmiterjoin%
\definecolor{currentfill}{rgb}{0.420392,0.751711,0.445829}%
\pgfsetfillcolor{currentfill}%
\pgfsetlinewidth{0.000000pt}%
\definecolor{currentstroke}{rgb}{0.000000,0.000000,0.000000}%
\pgfsetstrokecolor{currentstroke}%
\pgfsetstrokeopacity{0.000000}%
\pgfsetdash{}{0pt}%
\pgfpathmoveto{\pgfqpoint{1.669507in}{1.592201in}}%
\pgfpathlineto{\pgfqpoint{1.634891in}{1.593241in}}%
\pgfpathlineto{\pgfqpoint{1.587550in}{1.595331in}}%
\pgfpathlineto{\pgfqpoint{1.588665in}{1.618076in}}%
\pgfpathlineto{\pgfqpoint{1.633089in}{1.616328in}}%
\pgfpathlineto{\pgfqpoint{1.633892in}{1.639443in}}%
\pgfpathlineto{\pgfqpoint{1.649938in}{1.638890in}}%
\pgfpathlineto{\pgfqpoint{1.668280in}{1.638340in}}%
\pgfpathlineto{\pgfqpoint{1.667675in}{1.615294in}}%
\pgfpathlineto{\pgfqpoint{1.668146in}{1.596346in}}%
\pgfpathlineto{\pgfqpoint{1.669507in}{1.592201in}}%
\pgfpathclose%
\pgfusepath{fill}%
\end{pgfscope}%
\begin{pgfscope}%
\pgfpathrectangle{\pgfqpoint{0.100000in}{0.100000in}}{\pgfqpoint{3.007045in}{1.925000in}}%
\pgfusepath{clip}%
\pgfsetbuttcap%
\pgfsetmiterjoin%
\definecolor{currentfill}{rgb}{0.426667,0.754787,0.448904}%
\pgfsetfillcolor{currentfill}%
\pgfsetlinewidth{0.000000pt}%
\definecolor{currentstroke}{rgb}{0.000000,0.000000,0.000000}%
\pgfsetstrokecolor{currentstroke}%
\pgfsetstrokeopacity{0.000000}%
\pgfsetdash{}{0pt}%
\pgfpathmoveto{\pgfqpoint{2.826283in}{1.627777in}}%
\pgfpathlineto{\pgfqpoint{2.827912in}{1.634106in}}%
\pgfpathlineto{\pgfqpoint{2.822013in}{1.635422in}}%
\pgfpathlineto{\pgfqpoint{2.823531in}{1.641882in}}%
\pgfpathlineto{\pgfqpoint{2.815249in}{1.643795in}}%
\pgfpathlineto{\pgfqpoint{2.816481in}{1.647652in}}%
\pgfpathlineto{\pgfqpoint{2.812663in}{1.658624in}}%
\pgfpathlineto{\pgfqpoint{2.823785in}{1.661144in}}%
\pgfpathlineto{\pgfqpoint{2.860264in}{1.671267in}}%
\pgfpathlineto{\pgfqpoint{2.862575in}{1.687976in}}%
\pgfpathlineto{\pgfqpoint{2.873035in}{1.690583in}}%
\pgfpathlineto{\pgfqpoint{2.874149in}{1.694991in}}%
\pgfpathlineto{\pgfqpoint{2.885278in}{1.661827in}}%
\pgfpathlineto{\pgfqpoint{2.895122in}{1.630668in}}%
\pgfpathlineto{\pgfqpoint{2.894447in}{1.627269in}}%
\pgfpathlineto{\pgfqpoint{2.884731in}{1.625099in}}%
\pgfpathlineto{\pgfqpoint{2.884101in}{1.617605in}}%
\pgfpathlineto{\pgfqpoint{2.877796in}{1.621050in}}%
\pgfpathlineto{\pgfqpoint{2.875904in}{1.623095in}}%
\pgfpathlineto{\pgfqpoint{2.869747in}{1.621836in}}%
\pgfpathlineto{\pgfqpoint{2.868383in}{1.624903in}}%
\pgfpathlineto{\pgfqpoint{2.862905in}{1.625347in}}%
\pgfpathlineto{\pgfqpoint{2.858546in}{1.629287in}}%
\pgfpathlineto{\pgfqpoint{2.854814in}{1.623529in}}%
\pgfpathlineto{\pgfqpoint{2.849656in}{1.622164in}}%
\pgfpathlineto{\pgfqpoint{2.846880in}{1.617175in}}%
\pgfpathlineto{\pgfqpoint{2.848338in}{1.611817in}}%
\pgfpathlineto{\pgfqpoint{2.837477in}{1.608579in}}%
\pgfpathlineto{\pgfqpoint{2.834224in}{1.609178in}}%
\pgfpathlineto{\pgfqpoint{2.837240in}{1.625008in}}%
\pgfpathlineto{\pgfqpoint{2.826283in}{1.627777in}}%
\pgfpathclose%
\pgfusepath{fill}%
\end{pgfscope}%
\begin{pgfscope}%
\pgfpathrectangle{\pgfqpoint{0.100000in}{0.100000in}}{\pgfqpoint{3.007045in}{1.925000in}}%
\pgfusepath{clip}%
\pgfsetbuttcap%
\pgfsetmiterjoin%
\definecolor{currentfill}{rgb}{0.607151,0.839677,0.587928}%
\pgfsetfillcolor{currentfill}%
\pgfsetlinewidth{0.000000pt}%
\definecolor{currentstroke}{rgb}{0.000000,0.000000,0.000000}%
\pgfsetstrokecolor{currentstroke}%
\pgfsetstrokeopacity{0.000000}%
\pgfsetdash{}{0pt}%
\pgfpathmoveto{\pgfqpoint{1.610638in}{1.328045in}}%
\pgfpathlineto{\pgfqpoint{1.611668in}{1.362462in}}%
\pgfpathlineto{\pgfqpoint{1.657059in}{1.361181in}}%
\pgfpathlineto{\pgfqpoint{1.656900in}{1.355436in}}%
\pgfpathlineto{\pgfqpoint{1.656422in}{1.338221in}}%
\pgfpathlineto{\pgfqpoint{1.673497in}{1.337786in}}%
\pgfpathlineto{\pgfqpoint{1.672961in}{1.314877in}}%
\pgfpathlineto{\pgfqpoint{1.633082in}{1.315887in}}%
\pgfpathlineto{\pgfqpoint{1.633318in}{1.327373in}}%
\pgfpathlineto{\pgfqpoint{1.610638in}{1.328045in}}%
\pgfpathclose%
\pgfusepath{fill}%
\end{pgfscope}%
\begin{pgfscope}%
\pgfpathrectangle{\pgfqpoint{0.100000in}{0.100000in}}{\pgfqpoint{3.007045in}{1.925000in}}%
\pgfusepath{clip}%
\pgfsetbuttcap%
\pgfsetmiterjoin%
\definecolor{currentfill}{rgb}{0.755848,0.903391,0.729043}%
\pgfsetfillcolor{currentfill}%
\pgfsetlinewidth{0.000000pt}%
\definecolor{currentstroke}{rgb}{0.000000,0.000000,0.000000}%
\pgfsetstrokecolor{currentstroke}%
\pgfsetstrokeopacity{0.000000}%
\pgfsetdash{}{0pt}%
\pgfpathmoveto{\pgfqpoint{1.333705in}{1.220024in}}%
\pgfpathlineto{\pgfqpoint{1.311373in}{1.222190in}}%
\pgfpathlineto{\pgfqpoint{1.265951in}{1.226957in}}%
\pgfpathlineto{\pgfqpoint{1.267871in}{1.244110in}}%
\pgfpathlineto{\pgfqpoint{1.237983in}{1.247422in}}%
\pgfpathlineto{\pgfqpoint{1.233045in}{1.254595in}}%
\pgfpathlineto{\pgfqpoint{1.229691in}{1.263388in}}%
\pgfpathlineto{\pgfqpoint{1.222250in}{1.286251in}}%
\pgfpathlineto{\pgfqpoint{1.216879in}{1.294670in}}%
\pgfpathlineto{\pgfqpoint{1.217135in}{1.298862in}}%
\pgfpathlineto{\pgfqpoint{1.262254in}{1.293547in}}%
\pgfpathlineto{\pgfqpoint{1.294492in}{1.290219in}}%
\pgfpathlineto{\pgfqpoint{1.346503in}{1.285013in}}%
\pgfpathlineto{\pgfqpoint{1.343169in}{1.253667in}}%
\pgfpathlineto{\pgfqpoint{1.348876in}{1.253120in}}%
\pgfpathlineto{\pgfqpoint{1.345447in}{1.218910in}}%
\pgfpathlineto{\pgfqpoint{1.333705in}{1.220024in}}%
\pgfpathclose%
\pgfusepath{fill}%
\end{pgfscope}%
\begin{pgfscope}%
\pgfpathrectangle{\pgfqpoint{0.100000in}{0.100000in}}{\pgfqpoint{3.007045in}{1.925000in}}%
\pgfusepath{clip}%
\pgfsetbuttcap%
\pgfsetmiterjoin%
\definecolor{currentfill}{rgb}{0.370196,0.727105,0.421223}%
\pgfsetfillcolor{currentfill}%
\pgfsetlinewidth{0.000000pt}%
\definecolor{currentstroke}{rgb}{0.000000,0.000000,0.000000}%
\pgfsetstrokecolor{currentstroke}%
\pgfsetstrokeopacity{0.000000}%
\pgfsetdash{}{0pt}%
\pgfpathmoveto{\pgfqpoint{0.741011in}{1.724091in}}%
\pgfpathlineto{\pgfqpoint{0.746177in}{1.746042in}}%
\pgfpathlineto{\pgfqpoint{0.743350in}{1.746702in}}%
\pgfpathlineto{\pgfqpoint{0.742787in}{1.756014in}}%
\pgfpathlineto{\pgfqpoint{0.737197in}{1.757418in}}%
\pgfpathlineto{\pgfqpoint{0.740246in}{1.767188in}}%
\pgfpathlineto{\pgfqpoint{0.748808in}{1.770376in}}%
\pgfpathlineto{\pgfqpoint{0.752113in}{1.767443in}}%
\pgfpathlineto{\pgfqpoint{0.758923in}{1.767006in}}%
\pgfpathlineto{\pgfqpoint{0.760765in}{1.761473in}}%
\pgfpathlineto{\pgfqpoint{0.765426in}{1.755324in}}%
\pgfpathlineto{\pgfqpoint{0.765843in}{1.746677in}}%
\pgfpathlineto{\pgfqpoint{0.773003in}{1.745268in}}%
\pgfpathlineto{\pgfqpoint{0.774865in}{1.752771in}}%
\pgfpathlineto{\pgfqpoint{0.789165in}{1.749329in}}%
\pgfpathlineto{\pgfqpoint{0.794807in}{1.753888in}}%
\pgfpathlineto{\pgfqpoint{0.807745in}{1.750662in}}%
\pgfpathlineto{\pgfqpoint{0.812473in}{1.770492in}}%
\pgfpathlineto{\pgfqpoint{0.844253in}{1.762934in}}%
\pgfpathlineto{\pgfqpoint{0.873084in}{1.756371in}}%
\pgfpathlineto{\pgfqpoint{0.873285in}{1.749159in}}%
\pgfpathlineto{\pgfqpoint{0.884881in}{1.738405in}}%
\pgfpathlineto{\pgfqpoint{0.884196in}{1.733912in}}%
\pgfpathlineto{\pgfqpoint{0.895843in}{1.733369in}}%
\pgfpathlineto{\pgfqpoint{0.897388in}{1.731726in}}%
\pgfpathlineto{\pgfqpoint{0.891175in}{1.716853in}}%
\pgfpathlineto{\pgfqpoint{0.885687in}{1.708694in}}%
\pgfpathlineto{\pgfqpoint{0.885040in}{1.701682in}}%
\pgfpathlineto{\pgfqpoint{0.881292in}{1.694653in}}%
\pgfpathlineto{\pgfqpoint{0.883784in}{1.683111in}}%
\pgfpathlineto{\pgfqpoint{0.877655in}{1.681396in}}%
\pgfpathlineto{\pgfqpoint{0.874072in}{1.677313in}}%
\pgfpathlineto{\pgfqpoint{0.875599in}{1.669863in}}%
\pgfpathlineto{\pgfqpoint{0.871070in}{1.663297in}}%
\pgfpathlineto{\pgfqpoint{0.865364in}{1.658676in}}%
\pgfpathlineto{\pgfqpoint{0.859318in}{1.660383in}}%
\pgfpathlineto{\pgfqpoint{0.865574in}{1.646714in}}%
\pgfpathlineto{\pgfqpoint{0.860069in}{1.641241in}}%
\pgfpathlineto{\pgfqpoint{0.802276in}{1.654267in}}%
\pgfpathlineto{\pgfqpoint{0.792690in}{1.650559in}}%
\pgfpathlineto{\pgfqpoint{0.786641in}{1.653470in}}%
\pgfpathlineto{\pgfqpoint{0.786144in}{1.663045in}}%
\pgfpathlineto{\pgfqpoint{0.770539in}{1.666797in}}%
\pgfpathlineto{\pgfqpoint{0.780091in}{1.680597in}}%
\pgfpathlineto{\pgfqpoint{0.786067in}{1.686768in}}%
\pgfpathlineto{\pgfqpoint{0.784656in}{1.695189in}}%
\pgfpathlineto{\pgfqpoint{0.778737in}{1.702487in}}%
\pgfpathlineto{\pgfqpoint{0.774980in}{1.703762in}}%
\pgfpathlineto{\pgfqpoint{0.771689in}{1.716185in}}%
\pgfpathlineto{\pgfqpoint{0.741011in}{1.724091in}}%
\pgfpathclose%
\pgfusepath{fill}%
\end{pgfscope}%
\begin{pgfscope}%
\pgfpathrectangle{\pgfqpoint{0.100000in}{0.100000in}}{\pgfqpoint{3.007045in}{1.925000in}}%
\pgfusepath{clip}%
\pgfsetbuttcap%
\pgfsetmiterjoin%
\definecolor{currentfill}{rgb}{0.671696,0.867958,0.647105}%
\pgfsetfillcolor{currentfill}%
\pgfsetlinewidth{0.000000pt}%
\definecolor{currentstroke}{rgb}{0.000000,0.000000,0.000000}%
\pgfsetstrokecolor{currentstroke}%
\pgfsetstrokeopacity{0.000000}%
\pgfsetdash{}{0pt}%
\pgfpathmoveto{\pgfqpoint{2.492884in}{0.596677in}}%
\pgfpathlineto{\pgfqpoint{2.516161in}{0.600413in}}%
\pgfpathlineto{\pgfqpoint{2.520191in}{0.572396in}}%
\pgfpathlineto{\pgfqpoint{2.522149in}{0.556021in}}%
\pgfpathlineto{\pgfqpoint{2.513609in}{0.551975in}}%
\pgfpathlineto{\pgfqpoint{2.513059in}{0.555566in}}%
\pgfpathlineto{\pgfqpoint{2.501841in}{0.554053in}}%
\pgfpathlineto{\pgfqpoint{2.504526in}{0.536259in}}%
\pgfpathlineto{\pgfqpoint{2.496864in}{0.535136in}}%
\pgfpathlineto{\pgfqpoint{2.498437in}{0.523892in}}%
\pgfpathlineto{\pgfqpoint{2.494299in}{0.520921in}}%
\pgfpathlineto{\pgfqpoint{2.488304in}{0.521528in}}%
\pgfpathlineto{\pgfqpoint{2.486048in}{0.518807in}}%
\pgfpathlineto{\pgfqpoint{2.485328in}{0.522899in}}%
\pgfpathlineto{\pgfqpoint{2.480931in}{0.529259in}}%
\pgfpathlineto{\pgfqpoint{2.469105in}{0.528163in}}%
\pgfpathlineto{\pgfqpoint{2.464882in}{0.534167in}}%
\pgfpathlineto{\pgfqpoint{2.458365in}{0.538964in}}%
\pgfpathlineto{\pgfqpoint{2.455703in}{0.543435in}}%
\pgfpathlineto{\pgfqpoint{2.450937in}{0.544451in}}%
\pgfpathlineto{\pgfqpoint{2.444679in}{0.548491in}}%
\pgfpathlineto{\pgfqpoint{2.442205in}{0.558231in}}%
\pgfpathlineto{\pgfqpoint{2.444618in}{0.558566in}}%
\pgfpathlineto{\pgfqpoint{2.446412in}{0.569336in}}%
\pgfpathlineto{\pgfqpoint{2.470809in}{0.572873in}}%
\pgfpathlineto{\pgfqpoint{2.471783in}{0.578648in}}%
\pgfpathlineto{\pgfqpoint{2.476107in}{0.580682in}}%
\pgfpathlineto{\pgfqpoint{2.479736in}{0.576043in}}%
\pgfpathlineto{\pgfqpoint{2.485231in}{0.574838in}}%
\pgfpathlineto{\pgfqpoint{2.487697in}{0.577489in}}%
\pgfpathlineto{\pgfqpoint{2.488025in}{0.585155in}}%
\pgfpathlineto{\pgfqpoint{2.492884in}{0.596677in}}%
\pgfpathclose%
\pgfusepath{fill}%
\end{pgfscope}%
\begin{pgfscope}%
\pgfpathrectangle{\pgfqpoint{0.100000in}{0.100000in}}{\pgfqpoint{3.007045in}{1.925000in}}%
\pgfusepath{clip}%
\pgfsetbuttcap%
\pgfsetmiterjoin%
\definecolor{currentfill}{rgb}{0.704421,0.881738,0.678970}%
\pgfsetfillcolor{currentfill}%
\pgfsetlinewidth{0.000000pt}%
\definecolor{currentstroke}{rgb}{0.000000,0.000000,0.000000}%
\pgfsetstrokecolor{currentstroke}%
\pgfsetstrokeopacity{0.000000}%
\pgfsetdash{}{0pt}%
\pgfpathmoveto{\pgfqpoint{2.364743in}{1.174580in}}%
\pgfpathlineto{\pgfqpoint{2.370900in}{1.172765in}}%
\pgfpathlineto{\pgfqpoint{2.374431in}{1.146809in}}%
\pgfpathlineto{\pgfqpoint{2.360800in}{1.150692in}}%
\pgfpathlineto{\pgfqpoint{2.356171in}{1.148869in}}%
\pgfpathlineto{\pgfqpoint{2.354011in}{1.144847in}}%
\pgfpathlineto{\pgfqpoint{2.348520in}{1.145975in}}%
\pgfpathlineto{\pgfqpoint{2.342358in}{1.152328in}}%
\pgfpathlineto{\pgfqpoint{2.336776in}{1.153283in}}%
\pgfpathlineto{\pgfqpoint{2.331133in}{1.151744in}}%
\pgfpathlineto{\pgfqpoint{2.323384in}{1.154002in}}%
\pgfpathlineto{\pgfqpoint{2.326824in}{1.138484in}}%
\pgfpathlineto{\pgfqpoint{2.317016in}{1.135906in}}%
\pgfpathlineto{\pgfqpoint{2.308461in}{1.129091in}}%
\pgfpathlineto{\pgfqpoint{2.297453in}{1.137486in}}%
\pgfpathlineto{\pgfqpoint{2.296120in}{1.144067in}}%
\pgfpathlineto{\pgfqpoint{2.288955in}{1.139422in}}%
\pgfpathlineto{\pgfqpoint{2.283647in}{1.145438in}}%
\pgfpathlineto{\pgfqpoint{2.290373in}{1.148367in}}%
\pgfpathlineto{\pgfqpoint{2.294975in}{1.154537in}}%
\pgfpathlineto{\pgfqpoint{2.289964in}{1.156159in}}%
\pgfpathlineto{\pgfqpoint{2.291496in}{1.161362in}}%
\pgfpathlineto{\pgfqpoint{2.288006in}{1.165183in}}%
\pgfpathlineto{\pgfqpoint{2.291530in}{1.169095in}}%
\pgfpathlineto{\pgfqpoint{2.288706in}{1.196382in}}%
\pgfpathlineto{\pgfqpoint{2.288393in}{1.199350in}}%
\pgfpathlineto{\pgfqpoint{2.305334in}{1.201390in}}%
\pgfpathlineto{\pgfqpoint{2.314937in}{1.203868in}}%
\pgfpathlineto{\pgfqpoint{2.330635in}{1.204428in}}%
\pgfpathlineto{\pgfqpoint{2.331627in}{1.183748in}}%
\pgfpathlineto{\pgfqpoint{2.338836in}{1.184076in}}%
\pgfpathlineto{\pgfqpoint{2.340324in}{1.169258in}}%
\pgfpathlineto{\pgfqpoint{2.353578in}{1.170749in}}%
\pgfpathlineto{\pgfqpoint{2.364743in}{1.174580in}}%
\pgfpathclose%
\pgfusepath{fill}%
\end{pgfscope}%
\begin{pgfscope}%
\pgfpathrectangle{\pgfqpoint{0.100000in}{0.100000in}}{\pgfqpoint{3.007045in}{1.925000in}}%
\pgfusepath{clip}%
\pgfsetbuttcap%
\pgfsetmiterjoin%
\definecolor{currentfill}{rgb}{0.540715,0.808674,0.533303}%
\pgfsetfillcolor{currentfill}%
\pgfsetlinewidth{0.000000pt}%
\definecolor{currentstroke}{rgb}{0.000000,0.000000,0.000000}%
\pgfsetstrokecolor{currentstroke}%
\pgfsetstrokeopacity{0.000000}%
\pgfsetdash{}{0pt}%
\pgfpathmoveto{\pgfqpoint{2.439065in}{1.211016in}}%
\pgfpathlineto{\pgfqpoint{2.428341in}{1.210131in}}%
\pgfpathlineto{\pgfqpoint{2.427935in}{1.216104in}}%
\pgfpathlineto{\pgfqpoint{2.422257in}{1.215672in}}%
\pgfpathlineto{\pgfqpoint{2.410789in}{1.216745in}}%
\pgfpathlineto{\pgfqpoint{2.410154in}{1.226352in}}%
\pgfpathlineto{\pgfqpoint{2.407820in}{1.232038in}}%
\pgfpathlineto{\pgfqpoint{2.403956in}{1.231787in}}%
\pgfpathlineto{\pgfqpoint{2.403598in}{1.237933in}}%
\pgfpathlineto{\pgfqpoint{2.415203in}{1.238420in}}%
\pgfpathlineto{\pgfqpoint{2.416625in}{1.241076in}}%
\pgfpathlineto{\pgfqpoint{2.415184in}{1.255293in}}%
\pgfpathlineto{\pgfqpoint{2.441070in}{1.257982in}}%
\pgfpathlineto{\pgfqpoint{2.440547in}{1.262665in}}%
\pgfpathlineto{\pgfqpoint{2.457021in}{1.264583in}}%
\pgfpathlineto{\pgfqpoint{2.463037in}{1.262554in}}%
\pgfpathlineto{\pgfqpoint{2.465474in}{1.242785in}}%
\pgfpathlineto{\pgfqpoint{2.461646in}{1.242312in}}%
\pgfpathlineto{\pgfqpoint{2.462970in}{1.231712in}}%
\pgfpathlineto{\pgfqpoint{2.465782in}{1.225372in}}%
\pgfpathlineto{\pgfqpoint{2.456449in}{1.222465in}}%
\pgfpathlineto{\pgfqpoint{2.445563in}{1.221325in}}%
\pgfpathlineto{\pgfqpoint{2.445630in}{1.214029in}}%
\pgfpathlineto{\pgfqpoint{2.439738in}{1.214050in}}%
\pgfpathlineto{\pgfqpoint{2.439065in}{1.211016in}}%
\pgfpathclose%
\pgfusepath{fill}%
\end{pgfscope}%
\begin{pgfscope}%
\pgfpathrectangle{\pgfqpoint{0.100000in}{0.100000in}}{\pgfqpoint{3.007045in}{1.925000in}}%
\pgfusepath{clip}%
\pgfsetbuttcap%
\pgfsetmiterjoin%
\definecolor{currentfill}{rgb}{0.389020,0.736332,0.430450}%
\pgfsetfillcolor{currentfill}%
\pgfsetlinewidth{0.000000pt}%
\definecolor{currentstroke}{rgb}{0.000000,0.000000,0.000000}%
\pgfsetstrokecolor{currentstroke}%
\pgfsetstrokeopacity{0.000000}%
\pgfsetdash{}{0pt}%
\pgfpathmoveto{\pgfqpoint{1.704306in}{1.528001in}}%
\pgfpathlineto{\pgfqpoint{1.684152in}{1.528327in}}%
\pgfpathlineto{\pgfqpoint{1.684336in}{1.539811in}}%
\pgfpathlineto{\pgfqpoint{1.668319in}{1.540277in}}%
\pgfpathlineto{\pgfqpoint{1.668554in}{1.549944in}}%
\pgfpathlineto{\pgfqpoint{1.678758in}{1.549649in}}%
\pgfpathlineto{\pgfqpoint{1.679482in}{1.551623in}}%
\pgfpathlineto{\pgfqpoint{1.703819in}{1.551137in}}%
\pgfpathlineto{\pgfqpoint{1.701484in}{1.554430in}}%
\pgfpathlineto{\pgfqpoint{1.694140in}{1.556870in}}%
\pgfpathlineto{\pgfqpoint{1.686157in}{1.569986in}}%
\pgfpathlineto{\pgfqpoint{1.687099in}{1.572723in}}%
\pgfpathlineto{\pgfqpoint{1.695859in}{1.579141in}}%
\pgfpathlineto{\pgfqpoint{1.699144in}{1.583940in}}%
\pgfpathlineto{\pgfqpoint{1.700191in}{1.591568in}}%
\pgfpathlineto{\pgfqpoint{1.699675in}{1.597325in}}%
\pgfpathlineto{\pgfqpoint{1.714002in}{1.597084in}}%
\pgfpathlineto{\pgfqpoint{1.714476in}{1.591318in}}%
\pgfpathlineto{\pgfqpoint{1.714182in}{1.568146in}}%
\pgfpathlineto{\pgfqpoint{1.714573in}{1.556685in}}%
\pgfpathlineto{\pgfqpoint{1.720360in}{1.556564in}}%
\pgfpathlineto{\pgfqpoint{1.720167in}{1.545159in}}%
\pgfpathlineto{\pgfqpoint{1.717578in}{1.542068in}}%
\pgfpathlineto{\pgfqpoint{1.712446in}{1.545741in}}%
\pgfpathlineto{\pgfqpoint{1.704575in}{1.547315in}}%
\pgfpathlineto{\pgfqpoint{1.704306in}{1.528001in}}%
\pgfpathclose%
\pgfusepath{fill}%
\end{pgfscope}%
\begin{pgfscope}%
\pgfpathrectangle{\pgfqpoint{0.100000in}{0.100000in}}{\pgfqpoint{3.007045in}{1.925000in}}%
\pgfusepath{clip}%
\pgfsetbuttcap%
\pgfsetmiterjoin%
\definecolor{currentfill}{rgb}{0.676371,0.869927,0.651657}%
\pgfsetfillcolor{currentfill}%
\pgfsetlinewidth{0.000000pt}%
\definecolor{currentstroke}{rgb}{0.000000,0.000000,0.000000}%
\pgfsetstrokecolor{currentstroke}%
\pgfsetstrokeopacity{0.000000}%
\pgfsetdash{}{0pt}%
\pgfpathmoveto{\pgfqpoint{2.234453in}{1.308012in}}%
\pgfpathlineto{\pgfqpoint{2.233410in}{1.317320in}}%
\pgfpathlineto{\pgfqpoint{2.213486in}{1.315239in}}%
\pgfpathlineto{\pgfqpoint{2.211201in}{1.336530in}}%
\pgfpathlineto{\pgfqpoint{2.224488in}{1.337759in}}%
\pgfpathlineto{\pgfqpoint{2.225748in}{1.340539in}}%
\pgfpathlineto{\pgfqpoint{2.223806in}{1.358262in}}%
\pgfpathlineto{\pgfqpoint{2.269373in}{1.363358in}}%
\pgfpathlineto{\pgfqpoint{2.271732in}{1.342879in}}%
\pgfpathlineto{\pgfqpoint{2.253575in}{1.340858in}}%
\pgfpathlineto{\pgfqpoint{2.257315in}{1.308409in}}%
\pgfpathlineto{\pgfqpoint{2.240301in}{1.306723in}}%
\pgfpathlineto{\pgfqpoint{2.234453in}{1.308012in}}%
\pgfpathclose%
\pgfusepath{fill}%
\end{pgfscope}%
\begin{pgfscope}%
\pgfpathrectangle{\pgfqpoint{0.100000in}{0.100000in}}{\pgfqpoint{3.007045in}{1.925000in}}%
\pgfusepath{clip}%
\pgfsetbuttcap%
\pgfsetmiterjoin%
\definecolor{currentfill}{rgb}{0.699746,0.879769,0.674418}%
\pgfsetfillcolor{currentfill}%
\pgfsetlinewidth{0.000000pt}%
\definecolor{currentstroke}{rgb}{0.000000,0.000000,0.000000}%
\pgfsetstrokecolor{currentstroke}%
\pgfsetstrokeopacity{0.000000}%
\pgfsetdash{}{0pt}%
\pgfpathmoveto{\pgfqpoint{1.229691in}{1.263388in}}%
\pgfpathlineto{\pgfqpoint{1.233045in}{1.254595in}}%
\pgfpathlineto{\pgfqpoint{1.237983in}{1.247422in}}%
\pgfpathlineto{\pgfqpoint{1.237002in}{1.232833in}}%
\pgfpathlineto{\pgfqpoint{1.234437in}{1.230208in}}%
\pgfpathlineto{\pgfqpoint{1.232976in}{1.220918in}}%
\pgfpathlineto{\pgfqpoint{1.225797in}{1.217368in}}%
\pgfpathlineto{\pgfqpoint{1.220004in}{1.212317in}}%
\pgfpathlineto{\pgfqpoint{1.214794in}{1.212322in}}%
\pgfpathlineto{\pgfqpoint{1.212843in}{1.220333in}}%
\pgfpathlineto{\pgfqpoint{1.205530in}{1.228275in}}%
\pgfpathlineto{\pgfqpoint{1.186601in}{1.231472in}}%
\pgfpathlineto{\pgfqpoint{1.186728in}{1.236591in}}%
\pgfpathlineto{\pgfqpoint{1.189821in}{1.258740in}}%
\pgfpathlineto{\pgfqpoint{1.189706in}{1.265642in}}%
\pgfpathlineto{\pgfqpoint{1.192494in}{1.261100in}}%
\pgfpathlineto{\pgfqpoint{1.205670in}{1.255974in}}%
\pgfpathlineto{\pgfqpoint{1.213800in}{1.255082in}}%
\pgfpathlineto{\pgfqpoint{1.220648in}{1.256986in}}%
\pgfpathlineto{\pgfqpoint{1.222855in}{1.255180in}}%
\pgfpathlineto{\pgfqpoint{1.229691in}{1.263388in}}%
\pgfpathclose%
\pgfusepath{fill}%
\end{pgfscope}%
\begin{pgfscope}%
\pgfpathrectangle{\pgfqpoint{0.100000in}{0.100000in}}{\pgfqpoint{3.007045in}{1.925000in}}%
\pgfusepath{clip}%
\pgfsetbuttcap%
\pgfsetmiterjoin%
\definecolor{currentfill}{rgb}{0.797924,0.920738,0.771642}%
\pgfsetfillcolor{currentfill}%
\pgfsetlinewidth{0.000000pt}%
\definecolor{currentstroke}{rgb}{0.000000,0.000000,0.000000}%
\pgfsetstrokecolor{currentstroke}%
\pgfsetstrokeopacity{0.000000}%
\pgfsetdash{}{0pt}%
\pgfpathmoveto{\pgfqpoint{2.215737in}{1.157449in}}%
\pgfpathlineto{\pgfqpoint{2.213096in}{1.157168in}}%
\pgfpathlineto{\pgfqpoint{2.210701in}{1.176219in}}%
\pgfpathlineto{\pgfqpoint{2.198043in}{1.175099in}}%
\pgfpathlineto{\pgfqpoint{2.197194in}{1.183592in}}%
\pgfpathlineto{\pgfqpoint{2.194423in}{1.183326in}}%
\pgfpathlineto{\pgfqpoint{2.195596in}{1.188844in}}%
\pgfpathlineto{\pgfqpoint{2.193431in}{1.193828in}}%
\pgfpathlineto{\pgfqpoint{2.189712in}{1.229760in}}%
\pgfpathlineto{\pgfqpoint{2.212424in}{1.232038in}}%
\pgfpathlineto{\pgfqpoint{2.212181in}{1.234336in}}%
\pgfpathlineto{\pgfqpoint{2.231222in}{1.236425in}}%
\pgfpathlineto{\pgfqpoint{2.233016in}{1.218375in}}%
\pgfpathlineto{\pgfqpoint{2.247375in}{1.219993in}}%
\pgfpathlineto{\pgfqpoint{2.247429in}{1.209448in}}%
\pgfpathlineto{\pgfqpoint{2.262289in}{1.211080in}}%
\pgfpathlineto{\pgfqpoint{2.264298in}{1.193927in}}%
\pgfpathlineto{\pgfqpoint{2.264834in}{1.189201in}}%
\pgfpathlineto{\pgfqpoint{2.248039in}{1.187423in}}%
\pgfpathlineto{\pgfqpoint{2.245949in}{1.180376in}}%
\pgfpathlineto{\pgfqpoint{2.225671in}{1.177961in}}%
\pgfpathlineto{\pgfqpoint{2.227838in}{1.158721in}}%
\pgfpathlineto{\pgfqpoint{2.215737in}{1.157449in}}%
\pgfpathclose%
\pgfusepath{fill}%
\end{pgfscope}%
\begin{pgfscope}%
\pgfpathrectangle{\pgfqpoint{0.100000in}{0.100000in}}{\pgfqpoint{3.007045in}{1.925000in}}%
\pgfusepath{clip}%
\pgfsetbuttcap%
\pgfsetmiterjoin%
\definecolor{currentfill}{rgb}{0.643645,0.856148,0.619792}%
\pgfsetfillcolor{currentfill}%
\pgfsetlinewidth{0.000000pt}%
\definecolor{currentstroke}{rgb}{0.000000,0.000000,0.000000}%
\pgfsetstrokecolor{currentstroke}%
\pgfsetstrokeopacity{0.000000}%
\pgfsetdash{}{0pt}%
\pgfpathmoveto{\pgfqpoint{0.765905in}{0.863949in}}%
\pgfpathlineto{\pgfqpoint{0.754327in}{0.807329in}}%
\pgfpathlineto{\pgfqpoint{0.748143in}{0.777058in}}%
\pgfpathlineto{\pgfqpoint{0.673269in}{0.823873in}}%
\pgfpathlineto{\pgfqpoint{0.675440in}{0.831946in}}%
\pgfpathlineto{\pgfqpoint{0.681603in}{0.837269in}}%
\pgfpathlineto{\pgfqpoint{0.604721in}{0.848024in}}%
\pgfpathlineto{\pgfqpoint{0.611833in}{0.877391in}}%
\pgfpathlineto{\pgfqpoint{0.613032in}{0.877119in}}%
\pgfpathlineto{\pgfqpoint{0.618157in}{0.899792in}}%
\pgfpathlineto{\pgfqpoint{0.673193in}{0.887262in}}%
\pgfpathlineto{\pgfqpoint{0.691055in}{0.881658in}}%
\pgfpathlineto{\pgfqpoint{0.691629in}{0.872688in}}%
\pgfpathlineto{\pgfqpoint{0.687845in}{0.862115in}}%
\pgfpathlineto{\pgfqpoint{0.689447in}{0.856914in}}%
\pgfpathlineto{\pgfqpoint{0.710499in}{0.852063in}}%
\pgfpathlineto{\pgfqpoint{0.716592in}{0.880125in}}%
\pgfpathlineto{\pgfqpoint{0.733331in}{0.876625in}}%
\pgfpathlineto{\pgfqpoint{0.732137in}{0.870987in}}%
\pgfpathlineto{\pgfqpoint{0.765905in}{0.863949in}}%
\pgfpathclose%
\pgfusepath{fill}%
\end{pgfscope}%
\begin{pgfscope}%
\pgfpathrectangle{\pgfqpoint{0.100000in}{0.100000in}}{\pgfqpoint{3.007045in}{1.925000in}}%
\pgfusepath{clip}%
\pgfsetbuttcap%
\pgfsetroundjoin%
\pgfsetlinewidth{0.050187pt}%
\definecolor{currentstroke}{rgb}{1.000000,1.000000,1.000000}%
\pgfsetstrokecolor{currentstroke}%
\pgfsetdash{}{0pt}%
\pgfusepath{stroke}%
\end{pgfscope}%
\begin{pgfscope}%
\pgfpathrectangle{\pgfqpoint{0.100000in}{0.100000in}}{\pgfqpoint{3.007045in}{1.925000in}}%
\pgfusepath{clip}%
\pgfsetbuttcap%
\pgfsetroundjoin%
\pgfsetlinewidth{0.050187pt}%
\definecolor{currentstroke}{rgb}{1.000000,1.000000,1.000000}%
\pgfsetstrokecolor{currentstroke}%
\pgfsetdash{}{0pt}%
\pgfusepath{stroke}%
\end{pgfscope}%
\begin{pgfscope}%
\pgfpathrectangle{\pgfqpoint{0.100000in}{0.100000in}}{\pgfqpoint{3.007045in}{1.925000in}}%
\pgfusepath{clip}%
\pgfsetbuttcap%
\pgfsetroundjoin%
\pgfsetlinewidth{0.050187pt}%
\definecolor{currentstroke}{rgb}{1.000000,1.000000,1.000000}%
\pgfsetstrokecolor{currentstroke}%
\pgfsetdash{}{0pt}%
\pgfusepath{stroke}%
\end{pgfscope}%
\begin{pgfscope}%
\pgfpathrectangle{\pgfqpoint{0.100000in}{0.100000in}}{\pgfqpoint{3.007045in}{1.925000in}}%
\pgfusepath{clip}%
\pgfsetbuttcap%
\pgfsetroundjoin%
\pgfsetlinewidth{0.050187pt}%
\definecolor{currentstroke}{rgb}{1.000000,1.000000,1.000000}%
\pgfsetstrokecolor{currentstroke}%
\pgfsetdash{}{0pt}%
\pgfusepath{stroke}%
\end{pgfscope}%
\begin{pgfscope}%
\pgfpathrectangle{\pgfqpoint{0.100000in}{0.100000in}}{\pgfqpoint{3.007045in}{1.925000in}}%
\pgfusepath{clip}%
\pgfsetbuttcap%
\pgfsetroundjoin%
\pgfsetlinewidth{0.050187pt}%
\definecolor{currentstroke}{rgb}{1.000000,1.000000,1.000000}%
\pgfsetstrokecolor{currentstroke}%
\pgfsetdash{}{0pt}%
\pgfusepath{stroke}%
\end{pgfscope}%
\begin{pgfscope}%
\pgfpathrectangle{\pgfqpoint{0.100000in}{0.100000in}}{\pgfqpoint{3.007045in}{1.925000in}}%
\pgfusepath{clip}%
\pgfsetbuttcap%
\pgfsetroundjoin%
\pgfsetlinewidth{0.050187pt}%
\definecolor{currentstroke}{rgb}{1.000000,1.000000,1.000000}%
\pgfsetstrokecolor{currentstroke}%
\pgfsetdash{}{0pt}%
\pgfusepath{stroke}%
\end{pgfscope}%
\begin{pgfscope}%
\pgfpathrectangle{\pgfqpoint{0.100000in}{0.100000in}}{\pgfqpoint{3.007045in}{1.925000in}}%
\pgfusepath{clip}%
\pgfsetbuttcap%
\pgfsetroundjoin%
\pgfsetlinewidth{0.050187pt}%
\definecolor{currentstroke}{rgb}{1.000000,1.000000,1.000000}%
\pgfsetstrokecolor{currentstroke}%
\pgfsetdash{}{0pt}%
\pgfpathmoveto{\pgfqpoint{0.814503in}{1.911787in}}%
\pgfpathlineto{\pgfqpoint{0.798909in}{1.849997in}}%
\pgfpathlineto{\pgfqpoint{0.787821in}{1.805724in}}%
\pgfpathlineto{\pgfqpoint{0.774637in}{1.752465in}}%
\pgfpathlineto{\pgfqpoint{0.772049in}{1.739338in}}%
\pgfpathlineto{\pgfqpoint{0.773739in}{1.727298in}}%
\pgfpathlineto{\pgfqpoint{0.771499in}{1.716207in}}%
\pgfpathlineto{\pgfqpoint{0.679470in}{1.740222in}}%
\pgfpathlineto{\pgfqpoint{0.671157in}{1.737433in}}%
\pgfpathlineto{\pgfqpoint{0.664042in}{1.739853in}}%
\pgfpathlineto{\pgfqpoint{0.631104in}{1.740645in}}%
\pgfpathlineto{\pgfqpoint{0.619963in}{1.737504in}}%
\pgfpathlineto{\pgfqpoint{0.608900in}{1.738565in}}%
\pgfpathlineto{\pgfqpoint{0.604211in}{1.743481in}}%
\pgfpathlineto{\pgfqpoint{0.574722in}{1.742409in}}%
\pgfpathlineto{\pgfqpoint{0.570064in}{1.750443in}}%
\pgfpathlineto{\pgfqpoint{0.561223in}{1.754676in}}%
\pgfpathlineto{\pgfqpoint{0.548360in}{1.757255in}}%
\pgfpathlineto{\pgfqpoint{0.526157in}{1.753479in}}%
\pgfpathlineto{\pgfqpoint{0.517960in}{1.757179in}}%
\pgfpathlineto{\pgfqpoint{0.505323in}{1.767019in}}%
\pgfpathlineto{\pgfqpoint{0.509122in}{1.786306in}}%
\pgfpathlineto{\pgfqpoint{0.507790in}{1.796155in}}%
\pgfpathlineto{\pgfqpoint{0.497775in}{1.806612in}}%
\pgfpathlineto{\pgfqpoint{0.491359in}{1.805963in}}%
\pgfpathlineto{\pgfqpoint{0.486757in}{1.816540in}}%
\pgfpathlineto{\pgfqpoint{0.475856in}{1.820815in}}%
\pgfpathlineto{\pgfqpoint{0.467931in}{1.820270in}}%
\pgfpathlineto{\pgfqpoint{0.465368in}{1.831133in}}%
\pgfpathlineto{\pgfqpoint{0.473262in}{1.829977in}}%
\pgfpathlineto{\pgfqpoint{0.472635in}{1.845075in}}%
\pgfpathlineto{\pgfqpoint{0.469194in}{1.854040in}}%
\pgfpathlineto{\pgfqpoint{0.470711in}{1.865514in}}%
\pgfpathlineto{\pgfqpoint{0.474841in}{1.874144in}}%
\pgfpathlineto{\pgfqpoint{0.474608in}{1.890521in}}%
\pgfpathlineto{\pgfqpoint{0.472404in}{1.896463in}}%
\pgfpathlineto{\pgfqpoint{0.476212in}{1.915526in}}%
\pgfpathlineto{\pgfqpoint{0.475112in}{1.927812in}}%
\pgfpathlineto{\pgfqpoint{0.471305in}{1.933703in}}%
\pgfpathlineto{\pgfqpoint{0.471886in}{1.952963in}}%
\pgfpathlineto{\pgfqpoint{0.477424in}{1.967128in}}%
\pgfpathlineto{\pgfqpoint{0.483447in}{1.963661in}}%
\pgfpathlineto{\pgfqpoint{0.503389in}{1.943061in}}%
\pgfpathlineto{\pgfqpoint{0.527555in}{1.931787in}}%
\pgfpathlineto{\pgfqpoint{0.539928in}{1.930453in}}%
\pgfpathlineto{\pgfqpoint{0.551585in}{1.925654in}}%
\pgfpathlineto{\pgfqpoint{0.554894in}{1.909503in}}%
\pgfpathlineto{\pgfqpoint{0.544678in}{1.906443in}}%
\pgfpathlineto{\pgfqpoint{0.537908in}{1.893715in}}%
\pgfpathlineto{\pgfqpoint{0.545857in}{1.894422in}}%
\pgfpathlineto{\pgfqpoint{0.549062in}{1.900135in}}%
\pgfpathlineto{\pgfqpoint{0.560272in}{1.907280in}}%
\pgfpathlineto{\pgfqpoint{0.559693in}{1.896670in}}%
\pgfpathlineto{\pgfqpoint{0.552208in}{1.894959in}}%
\pgfpathlineto{\pgfqpoint{0.553451in}{1.881322in}}%
\pgfpathlineto{\pgfqpoint{0.544337in}{1.867661in}}%
\pgfpathlineto{\pgfqpoint{0.538503in}{1.874157in}}%
\pgfpathlineto{\pgfqpoint{0.521168in}{1.870625in}}%
\pgfpathlineto{\pgfqpoint{0.520256in}{1.862677in}}%
\pgfpathlineto{\pgfqpoint{0.526232in}{1.857836in}}%
\pgfpathlineto{\pgfqpoint{0.535349in}{1.857360in}}%
\pgfpathlineto{\pgfqpoint{0.547955in}{1.867706in}}%
\pgfpathlineto{\pgfqpoint{0.557932in}{1.868565in}}%
\pgfpathlineto{\pgfqpoint{0.558315in}{1.880017in}}%
\pgfpathlineto{\pgfqpoint{0.563457in}{1.896838in}}%
\pgfpathlineto{\pgfqpoint{0.575643in}{1.909348in}}%
\pgfpathlineto{\pgfqpoint{0.571580in}{1.919037in}}%
\pgfpathlineto{\pgfqpoint{0.574355in}{1.929469in}}%
\pgfpathlineto{\pgfqpoint{0.568949in}{1.939866in}}%
\pgfpathlineto{\pgfqpoint{0.576654in}{1.953161in}}%
\pgfpathlineto{\pgfqpoint{0.577808in}{1.961160in}}%
\pgfpathlineto{\pgfqpoint{0.571119in}{1.966417in}}%
\pgfpathlineto{\pgfqpoint{0.572232in}{1.979864in}}%
\pgfpathlineto{\pgfqpoint{0.652382in}{1.955513in}}%
\pgfpathlineto{\pgfqpoint{0.737492in}{1.931646in}}%
\pgfpathlineto{\pgfqpoint{0.814503in}{1.911787in}}%
\pgfusepath{stroke}%
\end{pgfscope}%
\begin{pgfscope}%
\pgfpathrectangle{\pgfqpoint{0.100000in}{0.100000in}}{\pgfqpoint{3.007045in}{1.925000in}}%
\pgfusepath{clip}%
\pgfsetbuttcap%
\pgfsetroundjoin%
\pgfsetlinewidth{0.050187pt}%
\definecolor{currentstroke}{rgb}{1.000000,1.000000,1.000000}%
\pgfsetstrokecolor{currentstroke}%
\pgfsetdash{}{0pt}%
\pgfpathmoveto{\pgfqpoint{0.559048in}{1.933344in}}%
\pgfpathlineto{\pgfqpoint{0.562994in}{1.918128in}}%
\pgfpathlineto{\pgfqpoint{0.569183in}{1.913074in}}%
\pgfpathlineto{\pgfqpoint{0.564245in}{1.906641in}}%
\pgfpathlineto{\pgfqpoint{0.559432in}{1.916068in}}%
\pgfpathlineto{\pgfqpoint{0.559048in}{1.933344in}}%
\pgfusepath{stroke}%
\end{pgfscope}%
\begin{pgfscope}%
\pgfpathrectangle{\pgfqpoint{0.100000in}{0.100000in}}{\pgfqpoint{3.007045in}{1.925000in}}%
\pgfusepath{clip}%
\pgfsetbuttcap%
\pgfsetroundjoin%
\pgfsetlinewidth{0.050187pt}%
\definecolor{currentstroke}{rgb}{1.000000,1.000000,1.000000}%
\pgfsetstrokecolor{currentstroke}%
\pgfsetdash{}{0pt}%
\pgfpathmoveto{\pgfqpoint{0.855834in}{1.901800in}}%
\pgfpathlineto{\pgfqpoint{0.941509in}{1.882574in}}%
\pgfpathlineto{\pgfqpoint{1.022196in}{1.866264in}}%
\pgfpathlineto{\pgfqpoint{1.084275in}{1.854890in}}%
\pgfpathlineto{\pgfqpoint{1.138397in}{1.845797in}}%
\pgfpathlineto{\pgfqpoint{1.192639in}{1.837451in}}%
\pgfpathlineto{\pgfqpoint{1.238831in}{1.830944in}}%
\pgfpathlineto{\pgfqpoint{1.285096in}{1.824977in}}%
\pgfpathlineto{\pgfqpoint{1.331427in}{1.819551in}}%
\pgfpathlineto{\pgfqpoint{1.375088in}{1.814940in}}%
\pgfpathlineto{\pgfqpoint{1.369083in}{1.748399in}}%
\pgfpathlineto{\pgfqpoint{1.360146in}{1.658478in}}%
\pgfpathlineto{\pgfqpoint{1.355458in}{1.612230in}}%
\pgfpathlineto{\pgfqpoint{1.348717in}{1.549903in}}%
\pgfpathlineto{\pgfqpoint{1.300990in}{1.555114in}}%
\pgfpathlineto{\pgfqpoint{1.246348in}{1.561317in}}%
\pgfpathlineto{\pgfqpoint{1.170478in}{1.571700in}}%
\pgfpathlineto{\pgfqpoint{1.136592in}{1.576534in}}%
\pgfpathlineto{\pgfqpoint{1.053106in}{1.589639in}}%
\pgfpathlineto{\pgfqpoint{1.024367in}{1.594897in}}%
\pgfpathlineto{\pgfqpoint{1.018377in}{1.560840in}}%
\pgfpathlineto{\pgfqpoint{1.015103in}{1.563270in}}%
\pgfpathlineto{\pgfqpoint{1.008949in}{1.579653in}}%
\pgfpathlineto{\pgfqpoint{1.001038in}{1.578644in}}%
\pgfpathlineto{\pgfqpoint{0.995375in}{1.570362in}}%
\pgfpathlineto{\pgfqpoint{0.983625in}{1.569254in}}%
\pgfpathlineto{\pgfqpoint{0.981264in}{1.573217in}}%
\pgfpathlineto{\pgfqpoint{0.970258in}{1.572737in}}%
\pgfpathlineto{\pgfqpoint{0.964724in}{1.576574in}}%
\pgfpathlineto{\pgfqpoint{0.956987in}{1.570638in}}%
\pgfpathlineto{\pgfqpoint{0.938194in}{1.575980in}}%
\pgfpathlineto{\pgfqpoint{0.929993in}{1.573111in}}%
\pgfpathlineto{\pgfqpoint{0.925581in}{1.585841in}}%
\pgfpathlineto{\pgfqpoint{0.925374in}{1.597953in}}%
\pgfpathlineto{\pgfqpoint{0.912420in}{1.606664in}}%
\pgfpathlineto{\pgfqpoint{0.914532in}{1.619218in}}%
\pgfpathlineto{\pgfqpoint{0.905308in}{1.639956in}}%
\pgfpathlineto{\pgfqpoint{0.905930in}{1.657584in}}%
\pgfpathlineto{\pgfqpoint{0.898002in}{1.666211in}}%
\pgfpathlineto{\pgfqpoint{0.890564in}{1.658578in}}%
\pgfpathlineto{\pgfqpoint{0.880439in}{1.654450in}}%
\pgfpathlineto{\pgfqpoint{0.872732in}{1.662959in}}%
\pgfpathlineto{\pgfqpoint{0.873985in}{1.676547in}}%
\pgfpathlineto{\pgfqpoint{0.883228in}{1.681901in}}%
\pgfpathlineto{\pgfqpoint{0.880771in}{1.690510in}}%
\pgfpathlineto{\pgfqpoint{0.896838in}{1.732179in}}%
\pgfpathlineto{\pgfqpoint{0.884169in}{1.733223in}}%
\pgfpathlineto{\pgfqpoint{0.882867in}{1.740700in}}%
\pgfpathlineto{\pgfqpoint{0.873638in}{1.747149in}}%
\pgfpathlineto{\pgfqpoint{0.874228in}{1.754357in}}%
\pgfpathlineto{\pgfqpoint{0.869360in}{1.759957in}}%
\pgfpathlineto{\pgfqpoint{0.861436in}{1.780289in}}%
\pgfpathlineto{\pgfqpoint{0.854194in}{1.784436in}}%
\pgfpathlineto{\pgfqpoint{0.849009in}{1.801984in}}%
\pgfpathlineto{\pgfqpoint{0.850232in}{1.813647in}}%
\pgfpathlineto{\pgfqpoint{0.840559in}{1.835130in}}%
\pgfpathlineto{\pgfqpoint{0.855834in}{1.901800in}}%
\pgfusepath{stroke}%
\end{pgfscope}%
\begin{pgfscope}%
\pgfpathrectangle{\pgfqpoint{0.100000in}{0.100000in}}{\pgfqpoint{3.007045in}{1.925000in}}%
\pgfusepath{clip}%
\pgfsetbuttcap%
\pgfsetroundjoin%
\pgfsetlinewidth{0.050187pt}%
\definecolor{currentstroke}{rgb}{1.000000,1.000000,1.000000}%
\pgfsetstrokecolor{currentstroke}%
\pgfsetdash{}{0pt}%
\pgfpathmoveto{\pgfqpoint{2.924097in}{1.558905in}}%
\pgfpathlineto{\pgfqpoint{2.922382in}{1.566196in}}%
\pgfpathlineto{\pgfqpoint{2.912852in}{1.572583in}}%
\pgfpathlineto{\pgfqpoint{2.887667in}{1.654962in}}%
\pgfpathlineto{\pgfqpoint{2.874093in}{1.694717in}}%
\pgfpathlineto{\pgfqpoint{2.885006in}{1.706170in}}%
\pgfpathlineto{\pgfqpoint{2.895614in}{1.734199in}}%
\pgfpathlineto{\pgfqpoint{2.900905in}{1.743156in}}%
\pgfpathlineto{\pgfqpoint{2.897157in}{1.746918in}}%
\pgfpathlineto{\pgfqpoint{2.895902in}{1.771799in}}%
\pgfpathlineto{\pgfqpoint{2.900635in}{1.779426in}}%
\pgfpathlineto{\pgfqpoint{2.898560in}{1.797155in}}%
\pgfpathlineto{\pgfqpoint{2.917451in}{1.855321in}}%
\pgfpathlineto{\pgfqpoint{2.925963in}{1.855639in}}%
\pgfpathlineto{\pgfqpoint{2.929497in}{1.845249in}}%
\pgfpathlineto{\pgfqpoint{2.937061in}{1.842266in}}%
\pgfpathlineto{\pgfqpoint{2.951351in}{1.854487in}}%
\pgfpathlineto{\pgfqpoint{2.962555in}{1.861705in}}%
\pgfpathlineto{\pgfqpoint{2.987243in}{1.848994in}}%
\pgfpathlineto{\pgfqpoint{3.009547in}{1.778417in}}%
\pgfpathlineto{\pgfqpoint{3.013786in}{1.761053in}}%
\pgfpathlineto{\pgfqpoint{3.023556in}{1.759017in}}%
\pgfpathlineto{\pgfqpoint{3.036806in}{1.747224in}}%
\pgfpathlineto{\pgfqpoint{3.036056in}{1.740353in}}%
\pgfpathlineto{\pgfqpoint{3.045081in}{1.732208in}}%
\pgfpathlineto{\pgfqpoint{3.052413in}{1.736881in}}%
\pgfpathlineto{\pgfqpoint{3.067756in}{1.718987in}}%
\pgfpathlineto{\pgfqpoint{3.060899in}{1.704668in}}%
\pgfpathlineto{\pgfqpoint{3.051704in}{1.704384in}}%
\pgfpathlineto{\pgfqpoint{3.044326in}{1.691601in}}%
\pgfpathlineto{\pgfqpoint{3.035390in}{1.689788in}}%
\pgfpathlineto{\pgfqpoint{3.028829in}{1.682982in}}%
\pgfpathlineto{\pgfqpoint{3.009234in}{1.675682in}}%
\pgfpathlineto{\pgfqpoint{2.997374in}{1.663918in}}%
\pgfpathlineto{\pgfqpoint{2.991350in}{1.672362in}}%
\pgfpathlineto{\pgfqpoint{2.985875in}{1.666310in}}%
\pgfpathlineto{\pgfqpoint{2.987374in}{1.641845in}}%
\pgfpathlineto{\pgfqpoint{2.983040in}{1.632155in}}%
\pgfpathlineto{\pgfqpoint{2.973585in}{1.634802in}}%
\pgfpathlineto{\pgfqpoint{2.972059in}{1.624866in}}%
\pgfpathlineto{\pgfqpoint{2.967973in}{1.620781in}}%
\pgfpathlineto{\pgfqpoint{2.959797in}{1.621777in}}%
\pgfpathlineto{\pgfqpoint{2.960312in}{1.612488in}}%
\pgfpathlineto{\pgfqpoint{2.947919in}{1.615126in}}%
\pgfpathlineto{\pgfqpoint{2.941151in}{1.602256in}}%
\pgfpathlineto{\pgfqpoint{2.943707in}{1.595522in}}%
\pgfpathlineto{\pgfqpoint{2.939681in}{1.584322in}}%
\pgfpathlineto{\pgfqpoint{2.933343in}{1.576084in}}%
\pgfpathlineto{\pgfqpoint{2.931748in}{1.558887in}}%
\pgfpathlineto{\pgfqpoint{2.924097in}{1.558905in}}%
\pgfusepath{stroke}%
\end{pgfscope}%
\begin{pgfscope}%
\pgfpathrectangle{\pgfqpoint{0.100000in}{0.100000in}}{\pgfqpoint{3.007045in}{1.925000in}}%
\pgfusepath{clip}%
\pgfsetbuttcap%
\pgfsetroundjoin%
\pgfsetlinewidth{0.050187pt}%
\definecolor{currentstroke}{rgb}{1.000000,1.000000,1.000000}%
\pgfsetstrokecolor{currentstroke}%
\pgfsetdash{}{0pt}%
\pgfpathmoveto{\pgfqpoint{3.012754in}{1.670658in}}%
\pgfpathlineto{\pgfqpoint{3.018349in}{1.676519in}}%
\pgfpathlineto{\pgfqpoint{3.023676in}{1.671010in}}%
\pgfpathlineto{\pgfqpoint{3.014108in}{1.663712in}}%
\pgfpathlineto{\pgfqpoint{3.012754in}{1.670658in}}%
\pgfusepath{stroke}%
\end{pgfscope}%
\begin{pgfscope}%
\pgfpathrectangle{\pgfqpoint{0.100000in}{0.100000in}}{\pgfqpoint{3.007045in}{1.925000in}}%
\pgfusepath{clip}%
\pgfsetbuttcap%
\pgfsetroundjoin%
\pgfsetlinewidth{0.050187pt}%
\definecolor{currentstroke}{rgb}{1.000000,1.000000,1.000000}%
\pgfsetstrokecolor{currentstroke}%
\pgfsetdash{}{0pt}%
\pgfpathmoveto{\pgfqpoint{1.355458in}{1.612230in}}%
\pgfpathlineto{\pgfqpoint{1.360146in}{1.658478in}}%
\pgfpathlineto{\pgfqpoint{1.369083in}{1.748399in}}%
\pgfpathlineto{\pgfqpoint{1.375088in}{1.814940in}}%
\pgfpathlineto{\pgfqpoint{1.424263in}{1.810326in}}%
\pgfpathlineto{\pgfqpoint{1.487174in}{1.805318in}}%
\pgfpathlineto{\pgfqpoint{1.544677in}{1.801614in}}%
\pgfpathlineto{\pgfqpoint{1.596745in}{1.798979in}}%
\pgfpathlineto{\pgfqpoint{1.674447in}{1.796314in}}%
\pgfpathlineto{\pgfqpoint{1.679750in}{1.774414in}}%
\pgfpathlineto{\pgfqpoint{1.677852in}{1.763948in}}%
\pgfpathlineto{\pgfqpoint{1.677227in}{1.742437in}}%
\pgfpathlineto{\pgfqpoint{1.680835in}{1.726326in}}%
\pgfpathlineto{\pgfqpoint{1.689100in}{1.702558in}}%
\pgfpathlineto{\pgfqpoint{1.689081in}{1.672635in}}%
\pgfpathlineto{\pgfqpoint{1.690611in}{1.637873in}}%
\pgfpathlineto{\pgfqpoint{1.692711in}{1.628504in}}%
\pgfpathlineto{\pgfqpoint{1.698854in}{1.618219in}}%
\pgfpathlineto{\pgfqpoint{1.700880in}{1.602193in}}%
\pgfpathlineto{\pgfqpoint{1.700005in}{1.591493in}}%
\pgfpathlineto{\pgfqpoint{1.634865in}{1.592905in}}%
\pgfpathlineto{\pgfqpoint{1.587476in}{1.595307in}}%
\pgfpathlineto{\pgfqpoint{1.518008in}{1.599005in}}%
\pgfpathlineto{\pgfqpoint{1.449501in}{1.603871in}}%
\pgfpathlineto{\pgfqpoint{1.403872in}{1.607552in}}%
\pgfpathlineto{\pgfqpoint{1.355458in}{1.612230in}}%
\pgfusepath{stroke}%
\end{pgfscope}%
\begin{pgfscope}%
\pgfpathrectangle{\pgfqpoint{0.100000in}{0.100000in}}{\pgfqpoint{3.007045in}{1.925000in}}%
\pgfusepath{clip}%
\pgfsetbuttcap%
\pgfsetroundjoin%
\pgfsetlinewidth{0.050187pt}%
\definecolor{currentstroke}{rgb}{1.000000,1.000000,1.000000}%
\pgfsetstrokecolor{currentstroke}%
\pgfsetdash{}{0pt}%
\pgfpathmoveto{\pgfqpoint{1.335768in}{1.418681in}}%
\pgfpathlineto{\pgfqpoint{1.343289in}{1.496206in}}%
\pgfpathlineto{\pgfqpoint{1.348717in}{1.549903in}}%
\pgfpathlineto{\pgfqpoint{1.355458in}{1.612230in}}%
\pgfpathlineto{\pgfqpoint{1.403872in}{1.607552in}}%
\pgfpathlineto{\pgfqpoint{1.449501in}{1.603871in}}%
\pgfpathlineto{\pgfqpoint{1.518008in}{1.599005in}}%
\pgfpathlineto{\pgfqpoint{1.587476in}{1.595307in}}%
\pgfpathlineto{\pgfqpoint{1.634865in}{1.592905in}}%
\pgfpathlineto{\pgfqpoint{1.700005in}{1.591493in}}%
\pgfpathlineto{\pgfqpoint{1.695595in}{1.578623in}}%
\pgfpathlineto{\pgfqpoint{1.686792in}{1.568520in}}%
\pgfpathlineto{\pgfqpoint{1.693533in}{1.556888in}}%
\pgfpathlineto{\pgfqpoint{1.700966in}{1.554402in}}%
\pgfpathlineto{\pgfqpoint{1.704494in}{1.547720in}}%
\pgfpathlineto{\pgfqpoint{1.703684in}{1.498934in}}%
\pgfpathlineto{\pgfqpoint{1.702328in}{1.430275in}}%
\pgfpathlineto{\pgfqpoint{1.697315in}{1.412242in}}%
\pgfpathlineto{\pgfqpoint{1.701799in}{1.402266in}}%
\pgfpathlineto{\pgfqpoint{1.692796in}{1.380831in}}%
\pgfpathlineto{\pgfqpoint{1.702281in}{1.363552in}}%
\pgfpathlineto{\pgfqpoint{1.694222in}{1.364874in}}%
\pgfpathlineto{\pgfqpoint{1.688724in}{1.375625in}}%
\pgfpathlineto{\pgfqpoint{1.669102in}{1.382994in}}%
\pgfpathlineto{\pgfqpoint{1.656727in}{1.389477in}}%
\pgfpathlineto{\pgfqpoint{1.635962in}{1.390016in}}%
\pgfpathlineto{\pgfqpoint{1.628754in}{1.384110in}}%
\pgfpathlineto{\pgfqpoint{1.605250in}{1.395738in}}%
\pgfpathlineto{\pgfqpoint{1.603447in}{1.399416in}}%
\pgfpathlineto{\pgfqpoint{1.521419in}{1.403298in}}%
\pgfpathlineto{\pgfqpoint{1.471583in}{1.406155in}}%
\pgfpathlineto{\pgfqpoint{1.430419in}{1.409379in}}%
\pgfpathlineto{\pgfqpoint{1.362400in}{1.415805in}}%
\pgfpathlineto{\pgfqpoint{1.335768in}{1.418681in}}%
\pgfusepath{stroke}%
\end{pgfscope}%
\begin{pgfscope}%
\pgfpathrectangle{\pgfqpoint{0.100000in}{0.100000in}}{\pgfqpoint{3.007045in}{1.925000in}}%
\pgfusepath{clip}%
\pgfsetbuttcap%
\pgfsetroundjoin%
\pgfsetlinewidth{0.050187pt}%
\definecolor{currentstroke}{rgb}{1.000000,1.000000,1.000000}%
\pgfsetstrokecolor{currentstroke}%
\pgfsetdash{}{0pt}%
\pgfpathmoveto{\pgfqpoint{1.322867in}{1.287410in}}%
\pgfpathlineto{\pgfqpoint{1.279137in}{1.291403in}}%
\pgfpathlineto{\pgfqpoint{1.183815in}{1.303168in}}%
\pgfpathlineto{\pgfqpoint{1.131961in}{1.310611in}}%
\pgfpathlineto{\pgfqpoint{1.076345in}{1.318594in}}%
\pgfpathlineto{\pgfqpoint{1.029497in}{1.326091in}}%
\pgfpathlineto{\pgfqpoint{0.978076in}{1.334895in}}%
\pgfpathlineto{\pgfqpoint{0.989767in}{1.399751in}}%
\pgfpathlineto{\pgfqpoint{1.001610in}{1.466254in}}%
\pgfpathlineto{\pgfqpoint{1.018377in}{1.560840in}}%
\pgfpathlineto{\pgfqpoint{1.024367in}{1.594897in}}%
\pgfpathlineto{\pgfqpoint{1.053106in}{1.589639in}}%
\pgfpathlineto{\pgfqpoint{1.136592in}{1.576534in}}%
\pgfpathlineto{\pgfqpoint{1.170478in}{1.571700in}}%
\pgfpathlineto{\pgfqpoint{1.246348in}{1.561317in}}%
\pgfpathlineto{\pgfqpoint{1.300990in}{1.555114in}}%
\pgfpathlineto{\pgfqpoint{1.348717in}{1.549903in}}%
\pgfpathlineto{\pgfqpoint{1.343289in}{1.496206in}}%
\pgfpathlineto{\pgfqpoint{1.335768in}{1.418681in}}%
\pgfpathlineto{\pgfqpoint{1.329313in}{1.352796in}}%
\pgfpathlineto{\pgfqpoint{1.322867in}{1.287410in}}%
\pgfusepath{stroke}%
\end{pgfscope}%
\begin{pgfscope}%
\pgfpathrectangle{\pgfqpoint{0.100000in}{0.100000in}}{\pgfqpoint{3.007045in}{1.925000in}}%
\pgfusepath{clip}%
\pgfsetbuttcap%
\pgfsetroundjoin%
\pgfsetlinewidth{0.050187pt}%
\definecolor{currentstroke}{rgb}{1.000000,1.000000,1.000000}%
\pgfsetstrokecolor{currentstroke}%
\pgfsetdash{}{0pt}%
\pgfpathmoveto{\pgfqpoint{2.122408in}{1.376872in}}%
\pgfpathlineto{\pgfqpoint{2.066992in}{1.372935in}}%
\pgfpathlineto{\pgfqpoint{1.984391in}{1.369415in}}%
\pgfpathlineto{\pgfqpoint{1.981252in}{1.377759in}}%
\pgfpathlineto{\pgfqpoint{1.962935in}{1.383996in}}%
\pgfpathlineto{\pgfqpoint{1.958894in}{1.395780in}}%
\pgfpathlineto{\pgfqpoint{1.957202in}{1.410359in}}%
\pgfpathlineto{\pgfqpoint{1.961330in}{1.417827in}}%
\pgfpathlineto{\pgfqpoint{1.954811in}{1.424995in}}%
\pgfpathlineto{\pgfqpoint{1.953238in}{1.433544in}}%
\pgfpathlineto{\pgfqpoint{1.951137in}{1.452456in}}%
\pgfpathlineto{\pgfqpoint{1.944889in}{1.462729in}}%
\pgfpathlineto{\pgfqpoint{1.933793in}{1.468495in}}%
\pgfpathlineto{\pgfqpoint{1.921718in}{1.478009in}}%
\pgfpathlineto{\pgfqpoint{1.915477in}{1.489266in}}%
\pgfpathlineto{\pgfqpoint{1.904277in}{1.493800in}}%
\pgfpathlineto{\pgfqpoint{1.897693in}{1.501178in}}%
\pgfpathlineto{\pgfqpoint{1.889718in}{1.502430in}}%
\pgfpathlineto{\pgfqpoint{1.875478in}{1.513385in}}%
\pgfpathlineto{\pgfqpoint{1.877791in}{1.525991in}}%
\pgfpathlineto{\pgfqpoint{1.877345in}{1.549965in}}%
\pgfpathlineto{\pgfqpoint{1.881737in}{1.556564in}}%
\pgfpathlineto{\pgfqpoint{1.877791in}{1.566528in}}%
\pgfpathlineto{\pgfqpoint{1.870840in}{1.568455in}}%
\pgfpathlineto{\pgfqpoint{1.871415in}{1.577207in}}%
\pgfpathlineto{\pgfqpoint{1.880038in}{1.591032in}}%
\pgfpathlineto{\pgfqpoint{1.897116in}{1.601964in}}%
\pgfpathlineto{\pgfqpoint{1.896052in}{1.640844in}}%
\pgfpathlineto{\pgfqpoint{1.904597in}{1.646684in}}%
\pgfpathlineto{\pgfqpoint{1.912684in}{1.642793in}}%
\pgfpathlineto{\pgfqpoint{1.929157in}{1.648489in}}%
\pgfpathlineto{\pgfqpoint{1.960149in}{1.662803in}}%
\pgfpathlineto{\pgfqpoint{1.964187in}{1.658371in}}%
\pgfpathlineto{\pgfqpoint{1.958317in}{1.638296in}}%
\pgfpathlineto{\pgfqpoint{1.967050in}{1.642702in}}%
\pgfpathlineto{\pgfqpoint{1.982007in}{1.638302in}}%
\pgfpathlineto{\pgfqpoint{1.991199in}{1.634631in}}%
\pgfpathlineto{\pgfqpoint{1.996352in}{1.623871in}}%
\pgfpathlineto{\pgfqpoint{2.043506in}{1.613716in}}%
\pgfpathlineto{\pgfqpoint{2.057600in}{1.606750in}}%
\pgfpathlineto{\pgfqpoint{2.071931in}{1.606832in}}%
\pgfpathlineto{\pgfqpoint{2.086656in}{1.603958in}}%
\pgfpathlineto{\pgfqpoint{2.096212in}{1.594142in}}%
\pgfpathlineto{\pgfqpoint{2.105344in}{1.589291in}}%
\pgfpathlineto{\pgfqpoint{2.107021in}{1.575301in}}%
\pgfpathlineto{\pgfqpoint{2.104324in}{1.566517in}}%
\pgfpathlineto{\pgfqpoint{2.114501in}{1.565867in}}%
\pgfpathlineto{\pgfqpoint{2.111072in}{1.555679in}}%
\pgfpathlineto{\pgfqpoint{2.114334in}{1.552060in}}%
\pgfpathlineto{\pgfqpoint{2.117577in}{1.542443in}}%
\pgfpathlineto{\pgfqpoint{2.107659in}{1.537351in}}%
\pgfpathlineto{\pgfqpoint{2.101900in}{1.523164in}}%
\pgfpathlineto{\pgfqpoint{2.100093in}{1.513132in}}%
\pgfpathlineto{\pgfqpoint{2.105618in}{1.511415in}}%
\pgfpathlineto{\pgfqpoint{2.112691in}{1.518938in}}%
\pgfpathlineto{\pgfqpoint{2.118705in}{1.531974in}}%
\pgfpathlineto{\pgfqpoint{2.126842in}{1.536502in}}%
\pgfpathlineto{\pgfqpoint{2.132957in}{1.530613in}}%
\pgfpathlineto{\pgfqpoint{2.126938in}{1.512787in}}%
\pgfpathlineto{\pgfqpoint{2.125050in}{1.498949in}}%
\pgfpathlineto{\pgfqpoint{2.126828in}{1.489001in}}%
\pgfpathlineto{\pgfqpoint{2.121239in}{1.483374in}}%
\pgfpathlineto{\pgfqpoint{2.118447in}{1.469622in}}%
\pgfpathlineto{\pgfqpoint{2.120728in}{1.455170in}}%
\pgfpathlineto{\pgfqpoint{2.114132in}{1.433796in}}%
\pgfpathlineto{\pgfqpoint{2.114271in}{1.423115in}}%
\pgfpathlineto{\pgfqpoint{2.119483in}{1.399971in}}%
\pgfpathlineto{\pgfqpoint{2.122861in}{1.396001in}}%
\pgfpathlineto{\pgfqpoint{2.122408in}{1.376872in}}%
\pgfusepath{stroke}%
\end{pgfscope}%
\begin{pgfscope}%
\pgfpathrectangle{\pgfqpoint{0.100000in}{0.100000in}}{\pgfqpoint{3.007045in}{1.925000in}}%
\pgfusepath{clip}%
\pgfsetbuttcap%
\pgfsetroundjoin%
\pgfsetlinewidth{0.050187pt}%
\definecolor{currentstroke}{rgb}{1.000000,1.000000,1.000000}%
\pgfsetstrokecolor{currentstroke}%
\pgfsetdash{}{0pt}%
\pgfpathmoveto{\pgfqpoint{2.143186in}{1.564449in}}%
\pgfpathlineto{\pgfqpoint{2.144423in}{1.555218in}}%
\pgfpathlineto{\pgfqpoint{2.133091in}{1.530893in}}%
\pgfpathlineto{\pgfqpoint{2.128055in}{1.537939in}}%
\pgfpathlineto{\pgfqpoint{2.143186in}{1.564449in}}%
\pgfusepath{stroke}%
\end{pgfscope}%
\begin{pgfscope}%
\pgfpathrectangle{\pgfqpoint{0.100000in}{0.100000in}}{\pgfqpoint{3.007045in}{1.925000in}}%
\pgfusepath{clip}%
\pgfsetbuttcap%
\pgfsetroundjoin%
\pgfsetlinewidth{0.050187pt}%
\definecolor{currentstroke}{rgb}{1.000000,1.000000,1.000000}%
\pgfsetstrokecolor{currentstroke}%
\pgfsetdash{}{0pt}%
\pgfpathmoveto{\pgfqpoint{0.771499in}{1.716207in}}%
\pgfpathlineto{\pgfqpoint{0.773739in}{1.727298in}}%
\pgfpathlineto{\pgfqpoint{0.772049in}{1.739338in}}%
\pgfpathlineto{\pgfqpoint{0.774637in}{1.752465in}}%
\pgfpathlineto{\pgfqpoint{0.787821in}{1.805724in}}%
\pgfpathlineto{\pgfqpoint{0.798909in}{1.849997in}}%
\pgfpathlineto{\pgfqpoint{0.814503in}{1.911787in}}%
\pgfpathlineto{\pgfqpoint{0.855834in}{1.901800in}}%
\pgfpathlineto{\pgfqpoint{0.840559in}{1.835130in}}%
\pgfpathlineto{\pgfqpoint{0.850232in}{1.813647in}}%
\pgfpathlineto{\pgfqpoint{0.849009in}{1.801984in}}%
\pgfpathlineto{\pgfqpoint{0.854194in}{1.784436in}}%
\pgfpathlineto{\pgfqpoint{0.861436in}{1.780289in}}%
\pgfpathlineto{\pgfqpoint{0.869360in}{1.759957in}}%
\pgfpathlineto{\pgfqpoint{0.874228in}{1.754357in}}%
\pgfpathlineto{\pgfqpoint{0.873638in}{1.747149in}}%
\pgfpathlineto{\pgfqpoint{0.882867in}{1.740700in}}%
\pgfpathlineto{\pgfqpoint{0.884169in}{1.733223in}}%
\pgfpathlineto{\pgfqpoint{0.896838in}{1.732179in}}%
\pgfpathlineto{\pgfqpoint{0.880771in}{1.690510in}}%
\pgfpathlineto{\pgfqpoint{0.883228in}{1.681901in}}%
\pgfpathlineto{\pgfqpoint{0.873985in}{1.676547in}}%
\pgfpathlineto{\pgfqpoint{0.872732in}{1.662959in}}%
\pgfpathlineto{\pgfqpoint{0.880439in}{1.654450in}}%
\pgfpathlineto{\pgfqpoint{0.890564in}{1.658578in}}%
\pgfpathlineto{\pgfqpoint{0.898002in}{1.666211in}}%
\pgfpathlineto{\pgfqpoint{0.905930in}{1.657584in}}%
\pgfpathlineto{\pgfqpoint{0.905308in}{1.639956in}}%
\pgfpathlineto{\pgfqpoint{0.914532in}{1.619218in}}%
\pgfpathlineto{\pgfqpoint{0.912420in}{1.606664in}}%
\pgfpathlineto{\pgfqpoint{0.925374in}{1.597953in}}%
\pgfpathlineto{\pgfqpoint{0.925581in}{1.585841in}}%
\pgfpathlineto{\pgfqpoint{0.929993in}{1.573111in}}%
\pgfpathlineto{\pgfqpoint{0.938194in}{1.575980in}}%
\pgfpathlineto{\pgfqpoint{0.956987in}{1.570638in}}%
\pgfpathlineto{\pgfqpoint{0.964724in}{1.576574in}}%
\pgfpathlineto{\pgfqpoint{0.970258in}{1.572737in}}%
\pgfpathlineto{\pgfqpoint{0.981264in}{1.573217in}}%
\pgfpathlineto{\pgfqpoint{0.983625in}{1.569254in}}%
\pgfpathlineto{\pgfqpoint{0.995375in}{1.570362in}}%
\pgfpathlineto{\pgfqpoint{1.001038in}{1.578644in}}%
\pgfpathlineto{\pgfqpoint{1.008949in}{1.579653in}}%
\pgfpathlineto{\pgfqpoint{1.015103in}{1.563270in}}%
\pgfpathlineto{\pgfqpoint{1.018377in}{1.560840in}}%
\pgfpathlineto{\pgfqpoint{1.001610in}{1.466254in}}%
\pgfpathlineto{\pgfqpoint{0.989767in}{1.399751in}}%
\pgfpathlineto{\pgfqpoint{0.896364in}{1.417782in}}%
\pgfpathlineto{\pgfqpoint{0.845912in}{1.427818in}}%
\pgfpathlineto{\pgfqpoint{0.798720in}{1.438195in}}%
\pgfpathlineto{\pgfqpoint{0.754611in}{1.448181in}}%
\pgfpathlineto{\pgfqpoint{0.703563in}{1.460492in}}%
\pgfpathlineto{\pgfqpoint{0.729893in}{1.568511in}}%
\pgfpathlineto{\pgfqpoint{0.731313in}{1.576405in}}%
\pgfpathlineto{\pgfqpoint{0.743022in}{1.597090in}}%
\pgfpathlineto{\pgfqpoint{0.730898in}{1.609522in}}%
\pgfpathlineto{\pgfqpoint{0.733416in}{1.621731in}}%
\pgfpathlineto{\pgfqpoint{0.738031in}{1.624787in}}%
\pgfpathlineto{\pgfqpoint{0.746274in}{1.637367in}}%
\pgfpathlineto{\pgfqpoint{0.758289in}{1.646070in}}%
\pgfpathlineto{\pgfqpoint{0.758946in}{1.652510in}}%
\pgfpathlineto{\pgfqpoint{0.766198in}{1.658944in}}%
\pgfpathlineto{\pgfqpoint{0.772219in}{1.670988in}}%
\pgfpathlineto{\pgfqpoint{0.784570in}{1.683729in}}%
\pgfpathlineto{\pgfqpoint{0.783693in}{1.696314in}}%
\pgfpathlineto{\pgfqpoint{0.775249in}{1.703304in}}%
\pgfpathlineto{\pgfqpoint{0.771499in}{1.716207in}}%
\pgfusepath{stroke}%
\end{pgfscope}%
\begin{pgfscope}%
\pgfpathrectangle{\pgfqpoint{0.100000in}{0.100000in}}{\pgfqpoint{3.007045in}{1.925000in}}%
\pgfusepath{clip}%
\pgfsetbuttcap%
\pgfsetroundjoin%
\pgfsetlinewidth{0.050187pt}%
\definecolor{currentstroke}{rgb}{1.000000,1.000000,1.000000}%
\pgfsetstrokecolor{currentstroke}%
\pgfsetdash{}{0pt}%
\pgfpathmoveto{\pgfqpoint{2.816139in}{1.505410in}}%
\pgfpathlineto{\pgfqpoint{2.812258in}{1.517661in}}%
\pgfpathlineto{\pgfqpoint{2.805056in}{1.554825in}}%
\pgfpathlineto{\pgfqpoint{2.795745in}{1.566258in}}%
\pgfpathlineto{\pgfqpoint{2.787577in}{1.586820in}}%
\pgfpathlineto{\pgfqpoint{2.790259in}{1.602064in}}%
\pgfpathlineto{\pgfqpoint{2.788108in}{1.613300in}}%
\pgfpathlineto{\pgfqpoint{2.781130in}{1.624335in}}%
\pgfpathlineto{\pgfqpoint{2.780843in}{1.636492in}}%
\pgfpathlineto{\pgfqpoint{2.776791in}{1.649602in}}%
\pgfpathlineto{\pgfqpoint{2.813053in}{1.658527in}}%
\pgfpathlineto{\pgfqpoint{2.860154in}{1.671241in}}%
\pgfpathlineto{\pgfqpoint{2.862034in}{1.663950in}}%
\pgfpathlineto{\pgfqpoint{2.859035in}{1.652344in}}%
\pgfpathlineto{\pgfqpoint{2.866075in}{1.643057in}}%
\pgfpathlineto{\pgfqpoint{2.862336in}{1.631290in}}%
\pgfpathlineto{\pgfqpoint{2.847477in}{1.616498in}}%
\pgfpathlineto{\pgfqpoint{2.851562in}{1.605386in}}%
\pgfpathlineto{\pgfqpoint{2.848535in}{1.581900in}}%
\pgfpathlineto{\pgfqpoint{2.844297in}{1.568658in}}%
\pgfpathlineto{\pgfqpoint{2.849703in}{1.531000in}}%
\pgfpathlineto{\pgfqpoint{2.848975in}{1.517751in}}%
\pgfpathlineto{\pgfqpoint{2.854206in}{1.513499in}}%
\pgfpathlineto{\pgfqpoint{2.816139in}{1.505410in}}%
\pgfusepath{stroke}%
\end{pgfscope}%
\begin{pgfscope}%
\pgfpathrectangle{\pgfqpoint{0.100000in}{0.100000in}}{\pgfqpoint{3.007045in}{1.925000in}}%
\pgfusepath{clip}%
\pgfsetbuttcap%
\pgfsetroundjoin%
\pgfsetlinewidth{0.050187pt}%
\definecolor{currentstroke}{rgb}{1.000000,1.000000,1.000000}%
\pgfsetstrokecolor{currentstroke}%
\pgfsetdash{}{0pt}%
\pgfpathmoveto{\pgfqpoint{1.702328in}{1.430275in}}%
\pgfpathlineto{\pgfqpoint{1.703684in}{1.498934in}}%
\pgfpathlineto{\pgfqpoint{1.704494in}{1.547720in}}%
\pgfpathlineto{\pgfqpoint{1.700966in}{1.554402in}}%
\pgfpathlineto{\pgfqpoint{1.693533in}{1.556888in}}%
\pgfpathlineto{\pgfqpoint{1.686792in}{1.568520in}}%
\pgfpathlineto{\pgfqpoint{1.695595in}{1.578623in}}%
\pgfpathlineto{\pgfqpoint{1.700005in}{1.591493in}}%
\pgfpathlineto{\pgfqpoint{1.700880in}{1.602193in}}%
\pgfpathlineto{\pgfqpoint{1.698854in}{1.618219in}}%
\pgfpathlineto{\pgfqpoint{1.692711in}{1.628504in}}%
\pgfpathlineto{\pgfqpoint{1.690611in}{1.637873in}}%
\pgfpathlineto{\pgfqpoint{1.689081in}{1.672635in}}%
\pgfpathlineto{\pgfqpoint{1.689100in}{1.702558in}}%
\pgfpathlineto{\pgfqpoint{1.680835in}{1.726326in}}%
\pgfpathlineto{\pgfqpoint{1.677227in}{1.742437in}}%
\pgfpathlineto{\pgfqpoint{1.677852in}{1.763948in}}%
\pgfpathlineto{\pgfqpoint{1.679750in}{1.774414in}}%
\pgfpathlineto{\pgfqpoint{1.674447in}{1.796314in}}%
\pgfpathlineto{\pgfqpoint{1.710550in}{1.795591in}}%
\pgfpathlineto{\pgfqpoint{1.765389in}{1.795120in}}%
\pgfpathlineto{\pgfqpoint{1.765689in}{1.820011in}}%
\pgfpathlineto{\pgfqpoint{1.779648in}{1.817270in}}%
\pgfpathlineto{\pgfqpoint{1.786338in}{1.786920in}}%
\pgfpathlineto{\pgfqpoint{1.791270in}{1.776008in}}%
\pgfpathlineto{\pgfqpoint{1.803535in}{1.775685in}}%
\pgfpathlineto{\pgfqpoint{1.806280in}{1.771983in}}%
\pgfpathlineto{\pgfqpoint{1.823388in}{1.770342in}}%
\pgfpathlineto{\pgfqpoint{1.826264in}{1.762818in}}%
\pgfpathlineto{\pgfqpoint{1.838054in}{1.764509in}}%
\pgfpathlineto{\pgfqpoint{1.847211in}{1.771547in}}%
\pgfpathlineto{\pgfqpoint{1.863003in}{1.771284in}}%
\pgfpathlineto{\pgfqpoint{1.872775in}{1.765624in}}%
\pgfpathlineto{\pgfqpoint{1.873898in}{1.760316in}}%
\pgfpathlineto{\pgfqpoint{1.883196in}{1.759205in}}%
\pgfpathlineto{\pgfqpoint{1.889264in}{1.744727in}}%
\pgfpathlineto{\pgfqpoint{1.893179in}{1.753629in}}%
\pgfpathlineto{\pgfqpoint{1.903859in}{1.754177in}}%
\pgfpathlineto{\pgfqpoint{1.906561in}{1.747244in}}%
\pgfpathlineto{\pgfqpoint{1.918579in}{1.744079in}}%
\pgfpathlineto{\pgfqpoint{1.925206in}{1.734093in}}%
\pgfpathlineto{\pgfqpoint{1.939904in}{1.737194in}}%
\pgfpathlineto{\pgfqpoint{1.956078in}{1.749448in}}%
\pgfpathlineto{\pgfqpoint{1.961975in}{1.738646in}}%
\pgfpathlineto{\pgfqpoint{1.988513in}{1.741602in}}%
\pgfpathlineto{\pgfqpoint{1.999857in}{1.734188in}}%
\pgfpathlineto{\pgfqpoint{2.006460in}{1.736835in}}%
\pgfpathlineto{\pgfqpoint{2.011755in}{1.732661in}}%
\pgfpathlineto{\pgfqpoint{1.996048in}{1.722759in}}%
\pgfpathlineto{\pgfqpoint{1.973627in}{1.713935in}}%
\pgfpathlineto{\pgfqpoint{1.951422in}{1.696299in}}%
\pgfpathlineto{\pgfqpoint{1.932171in}{1.673105in}}%
\pgfpathlineto{\pgfqpoint{1.917597in}{1.659394in}}%
\pgfpathlineto{\pgfqpoint{1.904832in}{1.649969in}}%
\pgfpathlineto{\pgfqpoint{1.896052in}{1.640844in}}%
\pgfpathlineto{\pgfqpoint{1.897116in}{1.601964in}}%
\pgfpathlineto{\pgfqpoint{1.880038in}{1.591032in}}%
\pgfpathlineto{\pgfqpoint{1.871415in}{1.577207in}}%
\pgfpathlineto{\pgfqpoint{1.870840in}{1.568455in}}%
\pgfpathlineto{\pgfqpoint{1.877791in}{1.566528in}}%
\pgfpathlineto{\pgfqpoint{1.881737in}{1.556564in}}%
\pgfpathlineto{\pgfqpoint{1.877345in}{1.549965in}}%
\pgfpathlineto{\pgfqpoint{1.877791in}{1.525991in}}%
\pgfpathlineto{\pgfqpoint{1.875478in}{1.513385in}}%
\pgfpathlineto{\pgfqpoint{1.889718in}{1.502430in}}%
\pgfpathlineto{\pgfqpoint{1.897693in}{1.501178in}}%
\pgfpathlineto{\pgfqpoint{1.904277in}{1.493800in}}%
\pgfpathlineto{\pgfqpoint{1.915477in}{1.489266in}}%
\pgfpathlineto{\pgfqpoint{1.921718in}{1.478009in}}%
\pgfpathlineto{\pgfqpoint{1.933793in}{1.468495in}}%
\pgfpathlineto{\pgfqpoint{1.944889in}{1.462729in}}%
\pgfpathlineto{\pgfqpoint{1.951137in}{1.452456in}}%
\pgfpathlineto{\pgfqpoint{1.953238in}{1.433544in}}%
\pgfpathlineto{\pgfqpoint{1.894352in}{1.431406in}}%
\pgfpathlineto{\pgfqpoint{1.844156in}{1.430356in}}%
\pgfpathlineto{\pgfqpoint{1.798422in}{1.429684in}}%
\pgfpathlineto{\pgfqpoint{1.750041in}{1.429758in}}%
\pgfpathlineto{\pgfqpoint{1.702328in}{1.430275in}}%
\pgfusepath{stroke}%
\end{pgfscope}%
\begin{pgfscope}%
\pgfpathrectangle{\pgfqpoint{0.100000in}{0.100000in}}{\pgfqpoint{3.007045in}{1.925000in}}%
\pgfusepath{clip}%
\pgfsetbuttcap%
\pgfsetroundjoin%
\pgfsetlinewidth{0.050187pt}%
\definecolor{currentstroke}{rgb}{1.000000,1.000000,1.000000}%
\pgfsetstrokecolor{currentstroke}%
\pgfsetdash{}{0pt}%
\pgfpathmoveto{\pgfqpoint{0.365302in}{1.558630in}}%
\pgfpathlineto{\pgfqpoint{0.360639in}{1.567212in}}%
\pgfpathlineto{\pgfqpoint{0.363617in}{1.589239in}}%
\pgfpathlineto{\pgfqpoint{0.369364in}{1.600750in}}%
\pgfpathlineto{\pgfqpoint{0.366491in}{1.616323in}}%
\pgfpathlineto{\pgfqpoint{0.372487in}{1.622871in}}%
\pgfpathlineto{\pgfqpoint{0.383425in}{1.640522in}}%
\pgfpathlineto{\pgfqpoint{0.392700in}{1.651178in}}%
\pgfpathlineto{\pgfqpoint{0.406257in}{1.674384in}}%
\pgfpathlineto{\pgfqpoint{0.416655in}{1.699664in}}%
\pgfpathlineto{\pgfqpoint{0.427694in}{1.723373in}}%
\pgfpathlineto{\pgfqpoint{0.429925in}{1.733275in}}%
\pgfpathlineto{\pgfqpoint{0.445153in}{1.761603in}}%
\pgfpathlineto{\pgfqpoint{0.448085in}{1.774045in}}%
\pgfpathlineto{\pgfqpoint{0.454568in}{1.787110in}}%
\pgfpathlineto{\pgfqpoint{0.456886in}{1.803047in}}%
\pgfpathlineto{\pgfqpoint{0.469263in}{1.810819in}}%
\pgfpathlineto{\pgfqpoint{0.483940in}{1.814736in}}%
\pgfpathlineto{\pgfqpoint{0.491359in}{1.805963in}}%
\pgfpathlineto{\pgfqpoint{0.497775in}{1.806612in}}%
\pgfpathlineto{\pgfqpoint{0.507790in}{1.796155in}}%
\pgfpathlineto{\pgfqpoint{0.509122in}{1.786306in}}%
\pgfpathlineto{\pgfqpoint{0.505323in}{1.767019in}}%
\pgfpathlineto{\pgfqpoint{0.517960in}{1.757179in}}%
\pgfpathlineto{\pgfqpoint{0.526157in}{1.753479in}}%
\pgfpathlineto{\pgfqpoint{0.548360in}{1.757255in}}%
\pgfpathlineto{\pgfqpoint{0.561223in}{1.754676in}}%
\pgfpathlineto{\pgfqpoint{0.570064in}{1.750443in}}%
\pgfpathlineto{\pgfqpoint{0.574722in}{1.742409in}}%
\pgfpathlineto{\pgfqpoint{0.604211in}{1.743481in}}%
\pgfpathlineto{\pgfqpoint{0.608900in}{1.738565in}}%
\pgfpathlineto{\pgfqpoint{0.619963in}{1.737504in}}%
\pgfpathlineto{\pgfqpoint{0.631104in}{1.740645in}}%
\pgfpathlineto{\pgfqpoint{0.664042in}{1.739853in}}%
\pgfpathlineto{\pgfqpoint{0.671157in}{1.737433in}}%
\pgfpathlineto{\pgfqpoint{0.679470in}{1.740222in}}%
\pgfpathlineto{\pgfqpoint{0.771499in}{1.716207in}}%
\pgfpathlineto{\pgfqpoint{0.775249in}{1.703304in}}%
\pgfpathlineto{\pgfqpoint{0.783693in}{1.696314in}}%
\pgfpathlineto{\pgfqpoint{0.784570in}{1.683729in}}%
\pgfpathlineto{\pgfqpoint{0.772219in}{1.670988in}}%
\pgfpathlineto{\pgfqpoint{0.766198in}{1.658944in}}%
\pgfpathlineto{\pgfqpoint{0.758946in}{1.652510in}}%
\pgfpathlineto{\pgfqpoint{0.758289in}{1.646070in}}%
\pgfpathlineto{\pgfqpoint{0.746274in}{1.637367in}}%
\pgfpathlineto{\pgfqpoint{0.738031in}{1.624787in}}%
\pgfpathlineto{\pgfqpoint{0.733416in}{1.621731in}}%
\pgfpathlineto{\pgfqpoint{0.730898in}{1.609522in}}%
\pgfpathlineto{\pgfqpoint{0.743022in}{1.597090in}}%
\pgfpathlineto{\pgfqpoint{0.731313in}{1.576405in}}%
\pgfpathlineto{\pgfqpoint{0.729893in}{1.568511in}}%
\pgfpathlineto{\pgfqpoint{0.703563in}{1.460492in}}%
\pgfpathlineto{\pgfqpoint{0.648178in}{1.474684in}}%
\pgfpathlineto{\pgfqpoint{0.594748in}{1.488461in}}%
\pgfpathlineto{\pgfqpoint{0.562516in}{1.497430in}}%
\pgfpathlineto{\pgfqpoint{0.521101in}{1.509228in}}%
\pgfpathlineto{\pgfqpoint{0.455041in}{1.530020in}}%
\pgfpathlineto{\pgfqpoint{0.383230in}{1.552378in}}%
\pgfpathlineto{\pgfqpoint{0.365302in}{1.558630in}}%
\pgfusepath{stroke}%
\end{pgfscope}%
\begin{pgfscope}%
\pgfpathrectangle{\pgfqpoint{0.100000in}{0.100000in}}{\pgfqpoint{3.007045in}{1.925000in}}%
\pgfusepath{clip}%
\pgfsetbuttcap%
\pgfsetroundjoin%
\pgfsetlinewidth{0.050187pt}%
\definecolor{currentstroke}{rgb}{1.000000,1.000000,1.000000}%
\pgfsetstrokecolor{currentstroke}%
\pgfsetdash{}{0pt}%
\pgfpathmoveto{\pgfqpoint{2.854206in}{1.513499in}}%
\pgfpathlineto{\pgfqpoint{2.848975in}{1.517751in}}%
\pgfpathlineto{\pgfqpoint{2.849703in}{1.531000in}}%
\pgfpathlineto{\pgfqpoint{2.844297in}{1.568658in}}%
\pgfpathlineto{\pgfqpoint{2.848535in}{1.581900in}}%
\pgfpathlineto{\pgfqpoint{2.851562in}{1.605386in}}%
\pgfpathlineto{\pgfqpoint{2.847477in}{1.616498in}}%
\pgfpathlineto{\pgfqpoint{2.862336in}{1.631290in}}%
\pgfpathlineto{\pgfqpoint{2.866075in}{1.643057in}}%
\pgfpathlineto{\pgfqpoint{2.859035in}{1.652344in}}%
\pgfpathlineto{\pgfqpoint{2.862034in}{1.663950in}}%
\pgfpathlineto{\pgfqpoint{2.860154in}{1.671241in}}%
\pgfpathlineto{\pgfqpoint{2.861768in}{1.686854in}}%
\pgfpathlineto{\pgfqpoint{2.864787in}{1.691674in}}%
\pgfpathlineto{\pgfqpoint{2.874093in}{1.694717in}}%
\pgfpathlineto{\pgfqpoint{2.887667in}{1.654962in}}%
\pgfpathlineto{\pgfqpoint{2.912852in}{1.572583in}}%
\pgfpathlineto{\pgfqpoint{2.922382in}{1.566196in}}%
\pgfpathlineto{\pgfqpoint{2.924097in}{1.558905in}}%
\pgfpathlineto{\pgfqpoint{2.929126in}{1.555958in}}%
\pgfpathlineto{\pgfqpoint{2.928744in}{1.542735in}}%
\pgfpathlineto{\pgfqpoint{2.923415in}{1.542522in}}%
\pgfpathlineto{\pgfqpoint{2.912617in}{1.534288in}}%
\pgfpathlineto{\pgfqpoint{2.909503in}{1.526026in}}%
\pgfpathlineto{\pgfqpoint{2.880603in}{1.518876in}}%
\pgfpathlineto{\pgfqpoint{2.854206in}{1.513499in}}%
\pgfusepath{stroke}%
\end{pgfscope}%
\begin{pgfscope}%
\pgfpathrectangle{\pgfqpoint{0.100000in}{0.100000in}}{\pgfqpoint{3.007045in}{1.925000in}}%
\pgfusepath{clip}%
\pgfsetbuttcap%
\pgfsetroundjoin%
\pgfsetlinewidth{0.050187pt}%
\definecolor{currentstroke}{rgb}{1.000000,1.000000,1.000000}%
\pgfsetstrokecolor{currentstroke}%
\pgfsetdash{}{0pt}%
\pgfpathmoveto{\pgfqpoint{1.950519in}{1.226895in}}%
\pgfpathlineto{\pgfqpoint{1.935262in}{1.242004in}}%
\pgfpathlineto{\pgfqpoint{1.886493in}{1.239191in}}%
\pgfpathlineto{\pgfqpoint{1.810432in}{1.236685in}}%
\pgfpathlineto{\pgfqpoint{1.733914in}{1.237878in}}%
\pgfpathlineto{\pgfqpoint{1.728543in}{1.247244in}}%
\pgfpathlineto{\pgfqpoint{1.730736in}{1.256440in}}%
\pgfpathlineto{\pgfqpoint{1.729695in}{1.272194in}}%
\pgfpathlineto{\pgfqpoint{1.725651in}{1.287453in}}%
\pgfpathlineto{\pgfqpoint{1.726121in}{1.295511in}}%
\pgfpathlineto{\pgfqpoint{1.717525in}{1.304616in}}%
\pgfpathlineto{\pgfqpoint{1.719395in}{1.317284in}}%
\pgfpathlineto{\pgfqpoint{1.706209in}{1.342308in}}%
\pgfpathlineto{\pgfqpoint{1.702281in}{1.363552in}}%
\pgfpathlineto{\pgfqpoint{1.692796in}{1.380831in}}%
\pgfpathlineto{\pgfqpoint{1.701799in}{1.402266in}}%
\pgfpathlineto{\pgfqpoint{1.697315in}{1.412242in}}%
\pgfpathlineto{\pgfqpoint{1.702328in}{1.430275in}}%
\pgfpathlineto{\pgfqpoint{1.750041in}{1.429758in}}%
\pgfpathlineto{\pgfqpoint{1.798422in}{1.429684in}}%
\pgfpathlineto{\pgfqpoint{1.844156in}{1.430356in}}%
\pgfpathlineto{\pgfqpoint{1.894352in}{1.431406in}}%
\pgfpathlineto{\pgfqpoint{1.953238in}{1.433544in}}%
\pgfpathlineto{\pgfqpoint{1.954811in}{1.424995in}}%
\pgfpathlineto{\pgfqpoint{1.961330in}{1.417827in}}%
\pgfpathlineto{\pgfqpoint{1.957202in}{1.410359in}}%
\pgfpathlineto{\pgfqpoint{1.958894in}{1.395780in}}%
\pgfpathlineto{\pgfqpoint{1.962935in}{1.383996in}}%
\pgfpathlineto{\pgfqpoint{1.981252in}{1.377759in}}%
\pgfpathlineto{\pgfqpoint{1.984391in}{1.369415in}}%
\pgfpathlineto{\pgfqpoint{1.994442in}{1.360047in}}%
\pgfpathlineto{\pgfqpoint{1.998543in}{1.350356in}}%
\pgfpathlineto{\pgfqpoint{2.008727in}{1.343855in}}%
\pgfpathlineto{\pgfqpoint{2.010321in}{1.336029in}}%
\pgfpathlineto{\pgfqpoint{2.008338in}{1.324180in}}%
\pgfpathlineto{\pgfqpoint{2.003154in}{1.320629in}}%
\pgfpathlineto{\pgfqpoint{2.001591in}{1.309350in}}%
\pgfpathlineto{\pgfqpoint{1.986693in}{1.300392in}}%
\pgfpathlineto{\pgfqpoint{1.967275in}{1.295485in}}%
\pgfpathlineto{\pgfqpoint{1.965496in}{1.284201in}}%
\pgfpathlineto{\pgfqpoint{1.972987in}{1.276140in}}%
\pgfpathlineto{\pgfqpoint{1.973292in}{1.266004in}}%
\pgfpathlineto{\pgfqpoint{1.967249in}{1.258037in}}%
\pgfpathlineto{\pgfqpoint{1.964076in}{1.246200in}}%
\pgfpathlineto{\pgfqpoint{1.953571in}{1.242282in}}%
\pgfpathlineto{\pgfqpoint{1.954240in}{1.229091in}}%
\pgfpathlineto{\pgfqpoint{1.950519in}{1.226895in}}%
\pgfusepath{stroke}%
\end{pgfscope}%
\begin{pgfscope}%
\pgfpathrectangle{\pgfqpoint{0.100000in}{0.100000in}}{\pgfqpoint{3.007045in}{1.925000in}}%
\pgfusepath{clip}%
\pgfsetbuttcap%
\pgfsetroundjoin%
\pgfsetlinewidth{0.050187pt}%
\definecolor{currentstroke}{rgb}{1.000000,1.000000,1.000000}%
\pgfsetstrokecolor{currentstroke}%
\pgfsetdash{}{0pt}%
\pgfpathmoveto{\pgfqpoint{2.897124in}{1.475380in}}%
\pgfpathlineto{\pgfqpoint{2.882673in}{1.473100in}}%
\pgfpathlineto{\pgfqpoint{2.816297in}{1.458017in}}%
\pgfpathlineto{\pgfqpoint{2.815137in}{1.459774in}}%
\pgfpathlineto{\pgfqpoint{2.816139in}{1.505410in}}%
\pgfpathlineto{\pgfqpoint{2.854206in}{1.513499in}}%
\pgfpathlineto{\pgfqpoint{2.880603in}{1.518876in}}%
\pgfpathlineto{\pgfqpoint{2.909503in}{1.526026in}}%
\pgfpathlineto{\pgfqpoint{2.912617in}{1.534288in}}%
\pgfpathlineto{\pgfqpoint{2.923415in}{1.542522in}}%
\pgfpathlineto{\pgfqpoint{2.928744in}{1.542735in}}%
\pgfpathlineto{\pgfqpoint{2.935751in}{1.530721in}}%
\pgfpathlineto{\pgfqpoint{2.929393in}{1.513184in}}%
\pgfpathlineto{\pgfqpoint{2.928460in}{1.502901in}}%
\pgfpathlineto{\pgfqpoint{2.941328in}{1.503867in}}%
\pgfpathlineto{\pgfqpoint{2.947164in}{1.498934in}}%
\pgfpathlineto{\pgfqpoint{2.960252in}{1.478771in}}%
\pgfpathlineto{\pgfqpoint{2.966756in}{1.476318in}}%
\pgfpathlineto{\pgfqpoint{2.977652in}{1.477228in}}%
\pgfpathlineto{\pgfqpoint{2.985234in}{1.484080in}}%
\pgfpathlineto{\pgfqpoint{2.990274in}{1.477965in}}%
\pgfpathlineto{\pgfqpoint{2.958579in}{1.461239in}}%
\pgfpathlineto{\pgfqpoint{2.957589in}{1.473241in}}%
\pgfpathlineto{\pgfqpoint{2.943191in}{1.454588in}}%
\pgfpathlineto{\pgfqpoint{2.938129in}{1.451375in}}%
\pgfpathlineto{\pgfqpoint{2.931068in}{1.462137in}}%
\pgfpathlineto{\pgfqpoint{2.929135in}{1.463612in}}%
\pgfpathlineto{\pgfqpoint{2.922562in}{1.467094in}}%
\pgfpathlineto{\pgfqpoint{2.916830in}{1.481216in}}%
\pgfpathlineto{\pgfqpoint{2.897124in}{1.475380in}}%
\pgfusepath{stroke}%
\end{pgfscope}%
\begin{pgfscope}%
\pgfpathrectangle{\pgfqpoint{0.100000in}{0.100000in}}{\pgfqpoint{3.007045in}{1.925000in}}%
\pgfusepath{clip}%
\pgfsetbuttcap%
\pgfsetroundjoin%
\pgfsetlinewidth{0.050187pt}%
\definecolor{currentstroke}{rgb}{1.000000,1.000000,1.000000}%
\pgfsetstrokecolor{currentstroke}%
\pgfsetdash{}{0pt}%
\pgfpathmoveto{\pgfqpoint{1.416825in}{1.212637in}}%
\pgfpathlineto{\pgfqpoint{1.422117in}{1.278255in}}%
\pgfpathlineto{\pgfqpoint{1.392146in}{1.280688in}}%
\pgfpathlineto{\pgfqpoint{1.322867in}{1.287410in}}%
\pgfpathlineto{\pgfqpoint{1.329313in}{1.352796in}}%
\pgfpathlineto{\pgfqpoint{1.335768in}{1.418681in}}%
\pgfpathlineto{\pgfqpoint{1.362400in}{1.415805in}}%
\pgfpathlineto{\pgfqpoint{1.430419in}{1.409379in}}%
\pgfpathlineto{\pgfqpoint{1.471583in}{1.406155in}}%
\pgfpathlineto{\pgfqpoint{1.521419in}{1.403298in}}%
\pgfpathlineto{\pgfqpoint{1.603447in}{1.399416in}}%
\pgfpathlineto{\pgfqpoint{1.605250in}{1.395738in}}%
\pgfpathlineto{\pgfqpoint{1.628754in}{1.384110in}}%
\pgfpathlineto{\pgfqpoint{1.635962in}{1.390016in}}%
\pgfpathlineto{\pgfqpoint{1.656727in}{1.389477in}}%
\pgfpathlineto{\pgfqpoint{1.669102in}{1.382994in}}%
\pgfpathlineto{\pgfqpoint{1.688724in}{1.375625in}}%
\pgfpathlineto{\pgfqpoint{1.694222in}{1.364874in}}%
\pgfpathlineto{\pgfqpoint{1.702281in}{1.363552in}}%
\pgfpathlineto{\pgfqpoint{1.706209in}{1.342308in}}%
\pgfpathlineto{\pgfqpoint{1.719395in}{1.317284in}}%
\pgfpathlineto{\pgfqpoint{1.717525in}{1.304616in}}%
\pgfpathlineto{\pgfqpoint{1.726121in}{1.295511in}}%
\pgfpathlineto{\pgfqpoint{1.725651in}{1.287453in}}%
\pgfpathlineto{\pgfqpoint{1.729695in}{1.272194in}}%
\pgfpathlineto{\pgfqpoint{1.730736in}{1.256440in}}%
\pgfpathlineto{\pgfqpoint{1.728543in}{1.247244in}}%
\pgfpathlineto{\pgfqpoint{1.733914in}{1.237878in}}%
\pgfpathlineto{\pgfqpoint{1.741281in}{1.220848in}}%
\pgfpathlineto{\pgfqpoint{1.748332in}{1.213916in}}%
\pgfpathlineto{\pgfqpoint{1.756737in}{1.198896in}}%
\pgfpathlineto{\pgfqpoint{1.732915in}{1.198649in}}%
\pgfpathlineto{\pgfqpoint{1.681409in}{1.199446in}}%
\pgfpathlineto{\pgfqpoint{1.624502in}{1.201190in}}%
\pgfpathlineto{\pgfqpoint{1.567259in}{1.203386in}}%
\pgfpathlineto{\pgfqpoint{1.483096in}{1.207983in}}%
\pgfpathlineto{\pgfqpoint{1.416825in}{1.212637in}}%
\pgfusepath{stroke}%
\end{pgfscope}%
\begin{pgfscope}%
\pgfpathrectangle{\pgfqpoint{0.100000in}{0.100000in}}{\pgfqpoint{3.007045in}{1.925000in}}%
\pgfusepath{clip}%
\pgfsetbuttcap%
\pgfsetroundjoin%
\pgfsetlinewidth{0.050187pt}%
\definecolor{currentstroke}{rgb}{1.000000,1.000000,1.000000}%
\pgfsetstrokecolor{currentstroke}%
\pgfsetdash{}{0pt}%
\pgfpathmoveto{\pgfqpoint{2.512669in}{1.410567in}}%
\pgfpathlineto{\pgfqpoint{2.539023in}{1.435695in}}%
\pgfpathlineto{\pgfqpoint{2.542299in}{1.444627in}}%
\pgfpathlineto{\pgfqpoint{2.550021in}{1.452274in}}%
\pgfpathlineto{\pgfqpoint{2.544226in}{1.463407in}}%
\pgfpathlineto{\pgfqpoint{2.536962in}{1.469919in}}%
\pgfpathlineto{\pgfqpoint{2.534862in}{1.481459in}}%
\pgfpathlineto{\pgfqpoint{2.561903in}{1.493308in}}%
\pgfpathlineto{\pgfqpoint{2.584291in}{1.497056in}}%
\pgfpathlineto{\pgfqpoint{2.596321in}{1.497306in}}%
\pgfpathlineto{\pgfqpoint{2.605485in}{1.492779in}}%
\pgfpathlineto{\pgfqpoint{2.614418in}{1.496822in}}%
\pgfpathlineto{\pgfqpoint{2.636206in}{1.501353in}}%
\pgfpathlineto{\pgfqpoint{2.643734in}{1.507201in}}%
\pgfpathlineto{\pgfqpoint{2.654897in}{1.520147in}}%
\pgfpathlineto{\pgfqpoint{2.665052in}{1.525868in}}%
\pgfpathlineto{\pgfqpoint{2.665769in}{1.531355in}}%
\pgfpathlineto{\pgfqpoint{2.660439in}{1.543873in}}%
\pgfpathlineto{\pgfqpoint{2.664292in}{1.551240in}}%
\pgfpathlineto{\pgfqpoint{2.659104in}{1.559161in}}%
\pgfpathlineto{\pgfqpoint{2.651159in}{1.559714in}}%
\pgfpathlineto{\pgfqpoint{2.671043in}{1.583641in}}%
\pgfpathlineto{\pgfqpoint{2.673405in}{1.592759in}}%
\pgfpathlineto{\pgfqpoint{2.689054in}{1.616015in}}%
\pgfpathlineto{\pgfqpoint{2.703581in}{1.628599in}}%
\pgfpathlineto{\pgfqpoint{2.713557in}{1.633854in}}%
\pgfpathlineto{\pgfqpoint{2.746198in}{1.641250in}}%
\pgfpathlineto{\pgfqpoint{2.776791in}{1.649602in}}%
\pgfpathlineto{\pgfqpoint{2.780843in}{1.636492in}}%
\pgfpathlineto{\pgfqpoint{2.781130in}{1.624335in}}%
\pgfpathlineto{\pgfqpoint{2.788108in}{1.613300in}}%
\pgfpathlineto{\pgfqpoint{2.790259in}{1.602064in}}%
\pgfpathlineto{\pgfqpoint{2.787577in}{1.586820in}}%
\pgfpathlineto{\pgfqpoint{2.795745in}{1.566258in}}%
\pgfpathlineto{\pgfqpoint{2.805056in}{1.554825in}}%
\pgfpathlineto{\pgfqpoint{2.812258in}{1.517661in}}%
\pgfpathlineto{\pgfqpoint{2.816139in}{1.505410in}}%
\pgfpathlineto{\pgfqpoint{2.815137in}{1.459774in}}%
\pgfpathlineto{\pgfqpoint{2.816297in}{1.458017in}}%
\pgfpathlineto{\pgfqpoint{2.824776in}{1.408935in}}%
\pgfpathlineto{\pgfqpoint{2.829531in}{1.404463in}}%
\pgfpathlineto{\pgfqpoint{2.819303in}{1.394526in}}%
\pgfpathlineto{\pgfqpoint{2.824350in}{1.388822in}}%
\pgfpathlineto{\pgfqpoint{2.819914in}{1.380196in}}%
\pgfpathlineto{\pgfqpoint{2.819952in}{1.376524in}}%
\pgfpathlineto{\pgfqpoint{2.814406in}{1.373210in}}%
\pgfpathlineto{\pgfqpoint{2.811706in}{1.365884in}}%
\pgfpathlineto{\pgfqpoint{2.812563in}{1.386028in}}%
\pgfpathlineto{\pgfqpoint{2.795375in}{1.390462in}}%
\pgfpathlineto{\pgfqpoint{2.768496in}{1.399631in}}%
\pgfpathlineto{\pgfqpoint{2.764695in}{1.404147in}}%
\pgfpathlineto{\pgfqpoint{2.753463in}{1.405229in}}%
\pgfpathlineto{\pgfqpoint{2.747002in}{1.411951in}}%
\pgfpathlineto{\pgfqpoint{2.743516in}{1.425338in}}%
\pgfpathlineto{\pgfqpoint{2.734357in}{1.427003in}}%
\pgfpathlineto{\pgfqpoint{2.728231in}{1.434026in}}%
\pgfpathlineto{\pgfqpoint{2.650253in}{1.418305in}}%
\pgfpathlineto{\pgfqpoint{2.612985in}{1.410597in}}%
\pgfpathlineto{\pgfqpoint{2.556401in}{1.400270in}}%
\pgfpathlineto{\pgfqpoint{2.515654in}{1.393396in}}%
\pgfpathlineto{\pgfqpoint{2.512669in}{1.410567in}}%
\pgfusepath{stroke}%
\end{pgfscope}%
\begin{pgfscope}%
\pgfpathrectangle{\pgfqpoint{0.100000in}{0.100000in}}{\pgfqpoint{3.007045in}{1.925000in}}%
\pgfusepath{clip}%
\pgfsetbuttcap%
\pgfsetroundjoin%
\pgfsetlinewidth{0.050187pt}%
\definecolor{currentstroke}{rgb}{1.000000,1.000000,1.000000}%
\pgfsetstrokecolor{currentstroke}%
\pgfsetdash{}{0pt}%
\pgfpathmoveto{\pgfqpoint{2.825722in}{1.361793in}}%
\pgfpathlineto{\pgfqpoint{2.813662in}{1.358036in}}%
\pgfpathlineto{\pgfqpoint{2.811634in}{1.361490in}}%
\pgfpathlineto{\pgfqpoint{2.815505in}{1.373084in}}%
\pgfpathlineto{\pgfqpoint{2.822673in}{1.374226in}}%
\pgfpathlineto{\pgfqpoint{2.822032in}{1.377876in}}%
\pgfpathlineto{\pgfqpoint{2.828469in}{1.383355in}}%
\pgfpathlineto{\pgfqpoint{2.847078in}{1.387719in}}%
\pgfpathlineto{\pgfqpoint{2.855363in}{1.394327in}}%
\pgfpathlineto{\pgfqpoint{2.873974in}{1.399833in}}%
\pgfpathlineto{\pgfqpoint{2.882462in}{1.397818in}}%
\pgfpathlineto{\pgfqpoint{2.889602in}{1.406695in}}%
\pgfpathlineto{\pgfqpoint{2.900394in}{1.407868in}}%
\pgfpathlineto{\pgfqpoint{2.881983in}{1.390553in}}%
\pgfpathlineto{\pgfqpoint{2.825722in}{1.361793in}}%
\pgfusepath{stroke}%
\end{pgfscope}%
\begin{pgfscope}%
\pgfpathrectangle{\pgfqpoint{0.100000in}{0.100000in}}{\pgfqpoint{3.007045in}{1.925000in}}%
\pgfusepath{clip}%
\pgfsetbuttcap%
\pgfsetroundjoin%
\pgfsetlinewidth{0.050187pt}%
\definecolor{currentstroke}{rgb}{1.000000,1.000000,1.000000}%
\pgfsetstrokecolor{currentstroke}%
\pgfsetdash{}{0pt}%
\pgfpathmoveto{\pgfqpoint{2.554781in}{1.247456in}}%
\pgfpathlineto{\pgfqpoint{2.502630in}{1.238834in}}%
\pgfpathlineto{\pgfqpoint{2.493170in}{1.298432in}}%
\pgfpathlineto{\pgfqpoint{2.479125in}{1.386284in}}%
\pgfpathlineto{\pgfqpoint{2.512669in}{1.410567in}}%
\pgfpathlineto{\pgfqpoint{2.515654in}{1.393396in}}%
\pgfpathlineto{\pgfqpoint{2.556401in}{1.400270in}}%
\pgfpathlineto{\pgfqpoint{2.612985in}{1.410597in}}%
\pgfpathlineto{\pgfqpoint{2.650253in}{1.418305in}}%
\pgfpathlineto{\pgfqpoint{2.728231in}{1.434026in}}%
\pgfpathlineto{\pgfqpoint{2.734357in}{1.427003in}}%
\pgfpathlineto{\pgfqpoint{2.743516in}{1.425338in}}%
\pgfpathlineto{\pgfqpoint{2.747002in}{1.411951in}}%
\pgfpathlineto{\pgfqpoint{2.753463in}{1.405229in}}%
\pgfpathlineto{\pgfqpoint{2.764695in}{1.404147in}}%
\pgfpathlineto{\pgfqpoint{2.768496in}{1.399631in}}%
\pgfpathlineto{\pgfqpoint{2.764637in}{1.396146in}}%
\pgfpathlineto{\pgfqpoint{2.761160in}{1.383818in}}%
\pgfpathlineto{\pgfqpoint{2.753012in}{1.369966in}}%
\pgfpathlineto{\pgfqpoint{2.758481in}{1.363942in}}%
\pgfpathlineto{\pgfqpoint{2.754006in}{1.357488in}}%
\pgfpathlineto{\pgfqpoint{2.756531in}{1.343334in}}%
\pgfpathlineto{\pgfqpoint{2.764590in}{1.335904in}}%
\pgfpathlineto{\pgfqpoint{2.783714in}{1.323833in}}%
\pgfpathlineto{\pgfqpoint{2.768314in}{1.306801in}}%
\pgfpathlineto{\pgfqpoint{2.768098in}{1.300337in}}%
\pgfpathlineto{\pgfqpoint{2.755560in}{1.292000in}}%
\pgfpathlineto{\pgfqpoint{2.741696in}{1.290435in}}%
\pgfpathlineto{\pgfqpoint{2.738266in}{1.283189in}}%
\pgfpathlineto{\pgfqpoint{2.699704in}{1.274845in}}%
\pgfpathlineto{\pgfqpoint{2.654706in}{1.265782in}}%
\pgfpathlineto{\pgfqpoint{2.623776in}{1.260244in}}%
\pgfpathlineto{\pgfqpoint{2.554781in}{1.247456in}}%
\pgfusepath{stroke}%
\end{pgfscope}%
\begin{pgfscope}%
\pgfpathrectangle{\pgfqpoint{0.100000in}{0.100000in}}{\pgfqpoint{3.007045in}{1.925000in}}%
\pgfusepath{clip}%
\pgfsetbuttcap%
\pgfsetroundjoin%
\pgfsetlinewidth{0.050187pt}%
\definecolor{currentstroke}{rgb}{1.000000,1.000000,1.000000}%
\pgfsetstrokecolor{currentstroke}%
\pgfsetdash{}{0pt}%
\pgfpathmoveto{\pgfqpoint{2.816297in}{1.458017in}}%
\pgfpathlineto{\pgfqpoint{2.882673in}{1.473100in}}%
\pgfpathlineto{\pgfqpoint{2.897124in}{1.475380in}}%
\pgfpathlineto{\pgfqpoint{2.906688in}{1.437776in}}%
\pgfpathlineto{\pgfqpoint{2.905174in}{1.431042in}}%
\pgfpathlineto{\pgfqpoint{2.874451in}{1.419148in}}%
\pgfpathlineto{\pgfqpoint{2.856110in}{1.414991in}}%
\pgfpathlineto{\pgfqpoint{2.848316in}{1.405665in}}%
\pgfpathlineto{\pgfqpoint{2.824350in}{1.388822in}}%
\pgfpathlineto{\pgfqpoint{2.819303in}{1.394526in}}%
\pgfpathlineto{\pgfqpoint{2.829531in}{1.404463in}}%
\pgfpathlineto{\pgfqpoint{2.824776in}{1.408935in}}%
\pgfpathlineto{\pgfqpoint{2.816297in}{1.458017in}}%
\pgfusepath{stroke}%
\end{pgfscope}%
\begin{pgfscope}%
\pgfpathrectangle{\pgfqpoint{0.100000in}{0.100000in}}{\pgfqpoint{3.007045in}{1.925000in}}%
\pgfusepath{clip}%
\pgfsetbuttcap%
\pgfsetroundjoin%
\pgfsetlinewidth{0.050187pt}%
\definecolor{currentstroke}{rgb}{1.000000,1.000000,1.000000}%
\pgfsetstrokecolor{currentstroke}%
\pgfsetdash{}{0pt}%
\pgfpathmoveto{\pgfqpoint{2.897124in}{1.475380in}}%
\pgfpathlineto{\pgfqpoint{2.916830in}{1.481216in}}%
\pgfpathlineto{\pgfqpoint{2.922562in}{1.467094in}}%
\pgfpathlineto{\pgfqpoint{2.929135in}{1.463612in}}%
\pgfpathlineto{\pgfqpoint{2.921024in}{1.457664in}}%
\pgfpathlineto{\pgfqpoint{2.922047in}{1.440194in}}%
\pgfpathlineto{\pgfqpoint{2.905174in}{1.431042in}}%
\pgfpathlineto{\pgfqpoint{2.906688in}{1.437776in}}%
\pgfpathlineto{\pgfqpoint{2.897124in}{1.475380in}}%
\pgfusepath{stroke}%
\end{pgfscope}%
\begin{pgfscope}%
\pgfpathrectangle{\pgfqpoint{0.100000in}{0.100000in}}{\pgfqpoint{3.007045in}{1.925000in}}%
\pgfusepath{clip}%
\pgfsetbuttcap%
\pgfsetroundjoin%
\pgfsetlinewidth{0.050187pt}%
\definecolor{currentstroke}{rgb}{1.000000,1.000000,1.000000}%
\pgfsetstrokecolor{currentstroke}%
\pgfsetdash{}{0pt}%
\pgfpathmoveto{\pgfqpoint{2.753338in}{1.285842in}}%
\pgfpathlineto{\pgfqpoint{2.755560in}{1.292000in}}%
\pgfpathlineto{\pgfqpoint{2.768098in}{1.300337in}}%
\pgfpathlineto{\pgfqpoint{2.768314in}{1.306801in}}%
\pgfpathlineto{\pgfqpoint{2.783714in}{1.323833in}}%
\pgfpathlineto{\pgfqpoint{2.764590in}{1.335904in}}%
\pgfpathlineto{\pgfqpoint{2.756531in}{1.343334in}}%
\pgfpathlineto{\pgfqpoint{2.754006in}{1.357488in}}%
\pgfpathlineto{\pgfqpoint{2.758481in}{1.363942in}}%
\pgfpathlineto{\pgfqpoint{2.753012in}{1.369966in}}%
\pgfpathlineto{\pgfqpoint{2.761160in}{1.383818in}}%
\pgfpathlineto{\pgfqpoint{2.764637in}{1.396146in}}%
\pgfpathlineto{\pgfqpoint{2.768496in}{1.399631in}}%
\pgfpathlineto{\pgfqpoint{2.795375in}{1.390462in}}%
\pgfpathlineto{\pgfqpoint{2.812563in}{1.386028in}}%
\pgfpathlineto{\pgfqpoint{2.811706in}{1.365884in}}%
\pgfpathlineto{\pgfqpoint{2.801265in}{1.350609in}}%
\pgfpathlineto{\pgfqpoint{2.809869in}{1.348362in}}%
\pgfpathlineto{\pgfqpoint{2.818801in}{1.341807in}}%
\pgfpathlineto{\pgfqpoint{2.819328in}{1.323878in}}%
\pgfpathlineto{\pgfqpoint{2.816624in}{1.311186in}}%
\pgfpathlineto{\pgfqpoint{2.818423in}{1.300760in}}%
\pgfpathlineto{\pgfqpoint{2.802896in}{1.265581in}}%
\pgfpathlineto{\pgfqpoint{2.797316in}{1.249154in}}%
\pgfpathlineto{\pgfqpoint{2.789524in}{1.257138in}}%
\pgfpathlineto{\pgfqpoint{2.779192in}{1.255778in}}%
\pgfpathlineto{\pgfqpoint{2.753333in}{1.270718in}}%
\pgfpathlineto{\pgfqpoint{2.750687in}{1.278719in}}%
\pgfpathlineto{\pgfqpoint{2.753338in}{1.285842in}}%
\pgfusepath{stroke}%
\end{pgfscope}%
\begin{pgfscope}%
\pgfpathrectangle{\pgfqpoint{0.100000in}{0.100000in}}{\pgfqpoint{3.007045in}{1.925000in}}%
\pgfusepath{clip}%
\pgfsetbuttcap%
\pgfsetroundjoin%
\pgfsetlinewidth{0.050187pt}%
\definecolor{currentstroke}{rgb}{1.000000,1.000000,1.000000}%
\pgfsetstrokecolor{currentstroke}%
\pgfsetdash{}{0pt}%
\pgfpathmoveto{\pgfqpoint{2.134480in}{1.068255in}}%
\pgfpathlineto{\pgfqpoint{2.133271in}{1.084575in}}%
\pgfpathlineto{\pgfqpoint{2.137970in}{1.092672in}}%
\pgfpathlineto{\pgfqpoint{2.134873in}{1.096640in}}%
\pgfpathlineto{\pgfqpoint{2.146167in}{1.109630in}}%
\pgfpathlineto{\pgfqpoint{2.145520in}{1.112760in}}%
\pgfpathlineto{\pgfqpoint{2.156492in}{1.135423in}}%
\pgfpathlineto{\pgfqpoint{2.154275in}{1.146352in}}%
\pgfpathlineto{\pgfqpoint{2.146318in}{1.157795in}}%
\pgfpathlineto{\pgfqpoint{2.151841in}{1.171716in}}%
\pgfpathlineto{\pgfqpoint{2.144682in}{1.263329in}}%
\pgfpathlineto{\pgfqpoint{2.139527in}{1.327598in}}%
\pgfpathlineto{\pgfqpoint{2.146648in}{1.322273in}}%
\pgfpathlineto{\pgfqpoint{2.154593in}{1.322417in}}%
\pgfpathlineto{\pgfqpoint{2.173359in}{1.333287in}}%
\pgfpathlineto{\pgfqpoint{2.230918in}{1.338658in}}%
\pgfpathlineto{\pgfqpoint{2.273522in}{1.343148in}}%
\pgfpathlineto{\pgfqpoint{2.273898in}{1.338979in}}%
\pgfpathlineto{\pgfqpoint{2.283623in}{1.250907in}}%
\pgfpathlineto{\pgfqpoint{2.291994in}{1.168959in}}%
\pgfpathlineto{\pgfqpoint{2.288386in}{1.165112in}}%
\pgfpathlineto{\pgfqpoint{2.293928in}{1.155563in}}%
\pgfpathlineto{\pgfqpoint{2.293900in}{1.148681in}}%
\pgfpathlineto{\pgfqpoint{2.285991in}{1.146952in}}%
\pgfpathlineto{\pgfqpoint{2.277140in}{1.140317in}}%
\pgfpathlineto{\pgfqpoint{2.271148in}{1.142929in}}%
\pgfpathlineto{\pgfqpoint{2.262180in}{1.138682in}}%
\pgfpathlineto{\pgfqpoint{2.264955in}{1.130153in}}%
\pgfpathlineto{\pgfqpoint{2.255736in}{1.121584in}}%
\pgfpathlineto{\pgfqpoint{2.253191in}{1.111671in}}%
\pgfpathlineto{\pgfqpoint{2.246853in}{1.110045in}}%
\pgfpathlineto{\pgfqpoint{2.242123in}{1.102538in}}%
\pgfpathlineto{\pgfqpoint{2.242742in}{1.094978in}}%
\pgfpathlineto{\pgfqpoint{2.237177in}{1.089662in}}%
\pgfpathlineto{\pgfqpoint{2.228803in}{1.090482in}}%
\pgfpathlineto{\pgfqpoint{2.222438in}{1.098640in}}%
\pgfpathlineto{\pgfqpoint{2.211653in}{1.090784in}}%
\pgfpathlineto{\pgfqpoint{2.212410in}{1.083920in}}%
\pgfpathlineto{\pgfqpoint{2.200414in}{1.079910in}}%
\pgfpathlineto{\pgfqpoint{2.197401in}{1.084958in}}%
\pgfpathlineto{\pgfqpoint{2.185723in}{1.079234in}}%
\pgfpathlineto{\pgfqpoint{2.181457in}{1.071040in}}%
\pgfpathlineto{\pgfqpoint{2.167393in}{1.079442in}}%
\pgfpathlineto{\pgfqpoint{2.145103in}{1.073876in}}%
\pgfpathlineto{\pgfqpoint{2.134480in}{1.068255in}}%
\pgfusepath{stroke}%
\end{pgfscope}%
\begin{pgfscope}%
\pgfpathrectangle{\pgfqpoint{0.100000in}{0.100000in}}{\pgfqpoint{3.007045in}{1.925000in}}%
\pgfusepath{clip}%
\pgfsetbuttcap%
\pgfsetroundjoin%
\pgfsetlinewidth{0.050187pt}%
\definecolor{currentstroke}{rgb}{1.000000,1.000000,1.000000}%
\pgfsetstrokecolor{currentstroke}%
\pgfsetdash{}{0pt}%
\pgfpathmoveto{\pgfqpoint{0.562516in}{1.497430in}}%
\pgfpathlineto{\pgfqpoint{0.594748in}{1.488461in}}%
\pgfpathlineto{\pgfqpoint{0.648178in}{1.474684in}}%
\pgfpathlineto{\pgfqpoint{0.703563in}{1.460492in}}%
\pgfpathlineto{\pgfqpoint{0.754611in}{1.448181in}}%
\pgfpathlineto{\pgfqpoint{0.798720in}{1.438195in}}%
\pgfpathlineto{\pgfqpoint{0.845912in}{1.427818in}}%
\pgfpathlineto{\pgfqpoint{0.832277in}{1.363472in}}%
\pgfpathlineto{\pgfqpoint{0.820130in}{1.306329in}}%
\pgfpathlineto{\pgfqpoint{0.800203in}{1.214130in}}%
\pgfpathlineto{\pgfqpoint{0.785243in}{1.144539in}}%
\pgfpathlineto{\pgfqpoint{0.777160in}{1.105704in}}%
\pgfpathlineto{\pgfqpoint{0.766797in}{1.055300in}}%
\pgfpathlineto{\pgfqpoint{0.759627in}{1.045075in}}%
\pgfpathlineto{\pgfqpoint{0.753855in}{1.044733in}}%
\pgfpathlineto{\pgfqpoint{0.749739in}{1.053659in}}%
\pgfpathlineto{\pgfqpoint{0.740288in}{1.056900in}}%
\pgfpathlineto{\pgfqpoint{0.730103in}{1.055764in}}%
\pgfpathlineto{\pgfqpoint{0.727212in}{1.048463in}}%
\pgfpathlineto{\pgfqpoint{0.726911in}{1.023110in}}%
\pgfpathlineto{\pgfqpoint{0.723861in}{1.017321in}}%
\pgfpathlineto{\pgfqpoint{0.725604in}{0.996967in}}%
\pgfpathlineto{\pgfqpoint{0.719235in}{0.983407in}}%
\pgfpathlineto{\pgfqpoint{0.667589in}{1.062934in}}%
\pgfpathlineto{\pgfqpoint{0.616726in}{1.140213in}}%
\pgfpathlineto{\pgfqpoint{0.590250in}{1.180736in}}%
\pgfpathlineto{\pgfqpoint{0.568163in}{1.215713in}}%
\pgfpathlineto{\pgfqpoint{0.540333in}{1.258980in}}%
\pgfpathlineto{\pgfqpoint{0.508830in}{1.307514in}}%
\pgfpathlineto{\pgfqpoint{0.521782in}{1.353580in}}%
\pgfpathlineto{\pgfqpoint{0.547847in}{1.445968in}}%
\pgfpathlineto{\pgfqpoint{0.562516in}{1.497430in}}%
\pgfusepath{stroke}%
\end{pgfscope}%
\begin{pgfscope}%
\pgfpathrectangle{\pgfqpoint{0.100000in}{0.100000in}}{\pgfqpoint{3.007045in}{1.925000in}}%
\pgfusepath{clip}%
\pgfsetbuttcap%
\pgfsetroundjoin%
\pgfsetlinewidth{0.050187pt}%
\definecolor{currentstroke}{rgb}{1.000000,1.000000,1.000000}%
\pgfsetstrokecolor{currentstroke}%
\pgfsetdash{}{0pt}%
\pgfpathmoveto{\pgfqpoint{0.845912in}{1.427818in}}%
\pgfpathlineto{\pgfqpoint{0.896364in}{1.417782in}}%
\pgfpathlineto{\pgfqpoint{0.989767in}{1.399751in}}%
\pgfpathlineto{\pgfqpoint{0.978076in}{1.334895in}}%
\pgfpathlineto{\pgfqpoint{1.029497in}{1.326091in}}%
\pgfpathlineto{\pgfqpoint{1.076345in}{1.318594in}}%
\pgfpathlineto{\pgfqpoint{1.068205in}{1.267320in}}%
\pgfpathlineto{\pgfqpoint{1.059579in}{1.212028in}}%
\pgfpathlineto{\pgfqpoint{1.048031in}{1.139428in}}%
\pgfpathlineto{\pgfqpoint{1.047732in}{1.133342in}}%
\pgfpathlineto{\pgfqpoint{1.035744in}{1.058103in}}%
\pgfpathlineto{\pgfqpoint{0.986397in}{1.065751in}}%
\pgfpathlineto{\pgfqpoint{0.961251in}{1.070801in}}%
\pgfpathlineto{\pgfqpoint{0.870361in}{1.086778in}}%
\pgfpathlineto{\pgfqpoint{0.836134in}{1.093528in}}%
\pgfpathlineto{\pgfqpoint{0.777160in}{1.105704in}}%
\pgfpathlineto{\pgfqpoint{0.785243in}{1.144539in}}%
\pgfpathlineto{\pgfqpoint{0.800203in}{1.214130in}}%
\pgfpathlineto{\pgfqpoint{0.820130in}{1.306329in}}%
\pgfpathlineto{\pgfqpoint{0.832277in}{1.363472in}}%
\pgfpathlineto{\pgfqpoint{0.845912in}{1.427818in}}%
\pgfusepath{stroke}%
\end{pgfscope}%
\begin{pgfscope}%
\pgfpathrectangle{\pgfqpoint{0.100000in}{0.100000in}}{\pgfqpoint{3.007045in}{1.925000in}}%
\pgfusepath{clip}%
\pgfsetbuttcap%
\pgfsetroundjoin%
\pgfsetlinewidth{0.050187pt}%
\definecolor{currentstroke}{rgb}{1.000000,1.000000,1.000000}%
\pgfsetstrokecolor{currentstroke}%
\pgfsetdash{}{0pt}%
\pgfpathmoveto{\pgfqpoint{0.365302in}{1.558630in}}%
\pgfpathlineto{\pgfqpoint{0.383230in}{1.552378in}}%
\pgfpathlineto{\pgfqpoint{0.455041in}{1.530020in}}%
\pgfpathlineto{\pgfqpoint{0.521101in}{1.509228in}}%
\pgfpathlineto{\pgfqpoint{0.562516in}{1.497430in}}%
\pgfpathlineto{\pgfqpoint{0.547847in}{1.445968in}}%
\pgfpathlineto{\pgfqpoint{0.521782in}{1.353580in}}%
\pgfpathlineto{\pgfqpoint{0.508830in}{1.307514in}}%
\pgfpathlineto{\pgfqpoint{0.540333in}{1.258980in}}%
\pgfpathlineto{\pgfqpoint{0.568163in}{1.215713in}}%
\pgfpathlineto{\pgfqpoint{0.590250in}{1.180736in}}%
\pgfpathlineto{\pgfqpoint{0.616726in}{1.140213in}}%
\pgfpathlineto{\pgfqpoint{0.667589in}{1.062934in}}%
\pgfpathlineto{\pgfqpoint{0.719235in}{0.983407in}}%
\pgfpathlineto{\pgfqpoint{0.717162in}{0.975520in}}%
\pgfpathlineto{\pgfqpoint{0.723391in}{0.962960in}}%
\pgfpathlineto{\pgfqpoint{0.724623in}{0.945778in}}%
\pgfpathlineto{\pgfqpoint{0.733671in}{0.937444in}}%
\pgfpathlineto{\pgfqpoint{0.734029in}{0.930728in}}%
\pgfpathlineto{\pgfqpoint{0.717829in}{0.923106in}}%
\pgfpathlineto{\pgfqpoint{0.710111in}{0.915475in}}%
\pgfpathlineto{\pgfqpoint{0.707704in}{0.898624in}}%
\pgfpathlineto{\pgfqpoint{0.703798in}{0.889434in}}%
\pgfpathlineto{\pgfqpoint{0.695567in}{0.881688in}}%
\pgfpathlineto{\pgfqpoint{0.687380in}{0.861545in}}%
\pgfpathlineto{\pgfqpoint{0.698891in}{0.851052in}}%
\pgfpathlineto{\pgfqpoint{0.697405in}{0.842403in}}%
\pgfpathlineto{\pgfqpoint{0.687967in}{0.836384in}}%
\pgfpathlineto{\pgfqpoint{0.681459in}{0.837446in}}%
\pgfpathlineto{\pgfqpoint{0.604646in}{0.848106in}}%
\pgfpathlineto{\pgfqpoint{0.547913in}{0.856154in}}%
\pgfpathlineto{\pgfqpoint{0.550394in}{0.865314in}}%
\pgfpathlineto{\pgfqpoint{0.546700in}{0.880521in}}%
\pgfpathlineto{\pgfqpoint{0.546328in}{0.895860in}}%
\pgfpathlineto{\pgfqpoint{0.543914in}{0.904837in}}%
\pgfpathlineto{\pgfqpoint{0.536479in}{0.917671in}}%
\pgfpathlineto{\pgfqpoint{0.515092in}{0.947282in}}%
\pgfpathlineto{\pgfqpoint{0.508088in}{0.950910in}}%
\pgfpathlineto{\pgfqpoint{0.499063in}{0.950853in}}%
\pgfpathlineto{\pgfqpoint{0.501124in}{0.960194in}}%
\pgfpathlineto{\pgfqpoint{0.496851in}{0.971875in}}%
\pgfpathlineto{\pgfqpoint{0.475860in}{0.977681in}}%
\pgfpathlineto{\pgfqpoint{0.463064in}{0.988432in}}%
\pgfpathlineto{\pgfqpoint{0.462003in}{0.995009in}}%
\pgfpathlineto{\pgfqpoint{0.447270in}{1.011295in}}%
\pgfpathlineto{\pgfqpoint{0.433245in}{1.014411in}}%
\pgfpathlineto{\pgfqpoint{0.420251in}{1.022687in}}%
\pgfpathlineto{\pgfqpoint{0.403161in}{1.025551in}}%
\pgfpathlineto{\pgfqpoint{0.395859in}{1.036589in}}%
\pgfpathlineto{\pgfqpoint{0.402794in}{1.054061in}}%
\pgfpathlineto{\pgfqpoint{0.400673in}{1.057990in}}%
\pgfpathlineto{\pgfqpoint{0.406447in}{1.072555in}}%
\pgfpathlineto{\pgfqpoint{0.396175in}{1.080306in}}%
\pgfpathlineto{\pgfqpoint{0.399508in}{1.094359in}}%
\pgfpathlineto{\pgfqpoint{0.394021in}{1.097951in}}%
\pgfpathlineto{\pgfqpoint{0.389283in}{1.111233in}}%
\pgfpathlineto{\pgfqpoint{0.383566in}{1.115284in}}%
\pgfpathlineto{\pgfqpoint{0.383136in}{1.124877in}}%
\pgfpathlineto{\pgfqpoint{0.378657in}{1.131642in}}%
\pgfpathlineto{\pgfqpoint{0.371908in}{1.154426in}}%
\pgfpathlineto{\pgfqpoint{0.364523in}{1.165339in}}%
\pgfpathlineto{\pgfqpoint{0.366128in}{1.183864in}}%
\pgfpathlineto{\pgfqpoint{0.374812in}{1.185728in}}%
\pgfpathlineto{\pgfqpoint{0.380467in}{1.195776in}}%
\pgfpathlineto{\pgfqpoint{0.377050in}{1.206674in}}%
\pgfpathlineto{\pgfqpoint{0.367814in}{1.208493in}}%
\pgfpathlineto{\pgfqpoint{0.363206in}{1.213591in}}%
\pgfpathlineto{\pgfqpoint{0.355744in}{1.232340in}}%
\pgfpathlineto{\pgfqpoint{0.359231in}{1.239073in}}%
\pgfpathlineto{\pgfqpoint{0.356757in}{1.251611in}}%
\pgfpathlineto{\pgfqpoint{0.362235in}{1.267849in}}%
\pgfpathlineto{\pgfqpoint{0.368652in}{1.261869in}}%
\pgfpathlineto{\pgfqpoint{0.365740in}{1.254809in}}%
\pgfpathlineto{\pgfqpoint{0.376861in}{1.243618in}}%
\pgfpathlineto{\pgfqpoint{0.376151in}{1.260237in}}%
\pgfpathlineto{\pgfqpoint{0.371386in}{1.264694in}}%
\pgfpathlineto{\pgfqpoint{0.376833in}{1.279306in}}%
\pgfpathlineto{\pgfqpoint{0.397434in}{1.276979in}}%
\pgfpathlineto{\pgfqpoint{0.394662in}{1.282413in}}%
\pgfpathlineto{\pgfqpoint{0.381065in}{1.281878in}}%
\pgfpathlineto{\pgfqpoint{0.374596in}{1.290108in}}%
\pgfpathlineto{\pgfqpoint{0.366455in}{1.282800in}}%
\pgfpathlineto{\pgfqpoint{0.362134in}{1.270584in}}%
\pgfpathlineto{\pgfqpoint{0.350603in}{1.287022in}}%
\pgfpathlineto{\pgfqpoint{0.346162in}{1.290014in}}%
\pgfpathlineto{\pgfqpoint{0.347854in}{1.307904in}}%
\pgfpathlineto{\pgfqpoint{0.344340in}{1.318454in}}%
\pgfpathlineto{\pgfqpoint{0.337968in}{1.328366in}}%
\pgfpathlineto{\pgfqpoint{0.324921in}{1.358735in}}%
\pgfpathlineto{\pgfqpoint{0.329187in}{1.365451in}}%
\pgfpathlineto{\pgfqpoint{0.329137in}{1.386675in}}%
\pgfpathlineto{\pgfqpoint{0.336178in}{1.398512in}}%
\pgfpathlineto{\pgfqpoint{0.337793in}{1.417013in}}%
\pgfpathlineto{\pgfqpoint{0.331153in}{1.438209in}}%
\pgfpathlineto{\pgfqpoint{0.322344in}{1.451701in}}%
\pgfpathlineto{\pgfqpoint{0.323915in}{1.463876in}}%
\pgfpathlineto{\pgfqpoint{0.348653in}{1.493342in}}%
\pgfpathlineto{\pgfqpoint{0.349881in}{1.503394in}}%
\pgfpathlineto{\pgfqpoint{0.361028in}{1.522568in}}%
\pgfpathlineto{\pgfqpoint{0.362575in}{1.540742in}}%
\pgfpathlineto{\pgfqpoint{0.358991in}{1.545379in}}%
\pgfpathlineto{\pgfqpoint{0.365302in}{1.558630in}}%
\pgfusepath{stroke}%
\end{pgfscope}%
\begin{pgfscope}%
\pgfpathrectangle{\pgfqpoint{0.100000in}{0.100000in}}{\pgfqpoint{3.007045in}{1.925000in}}%
\pgfusepath{clip}%
\pgfsetbuttcap%
\pgfsetroundjoin%
\pgfsetlinewidth{0.050187pt}%
\definecolor{currentstroke}{rgb}{1.000000,1.000000,1.000000}%
\pgfsetstrokecolor{currentstroke}%
\pgfsetdash{}{0pt}%
\pgfpathmoveto{\pgfqpoint{2.291994in}{1.168959in}}%
\pgfpathlineto{\pgfqpoint{2.283623in}{1.250907in}}%
\pgfpathlineto{\pgfqpoint{2.273898in}{1.338979in}}%
\pgfpathlineto{\pgfqpoint{2.337587in}{1.348464in}}%
\pgfpathlineto{\pgfqpoint{2.354491in}{1.344065in}}%
\pgfpathlineto{\pgfqpoint{2.362598in}{1.339289in}}%
\pgfpathlineto{\pgfqpoint{2.372757in}{1.340614in}}%
\pgfpathlineto{\pgfqpoint{2.386151in}{1.332707in}}%
\pgfpathlineto{\pgfqpoint{2.411100in}{1.344475in}}%
\pgfpathlineto{\pgfqpoint{2.424868in}{1.344868in}}%
\pgfpathlineto{\pgfqpoint{2.440946in}{1.362805in}}%
\pgfpathlineto{\pgfqpoint{2.457305in}{1.373716in}}%
\pgfpathlineto{\pgfqpoint{2.479125in}{1.386284in}}%
\pgfpathlineto{\pgfqpoint{2.493170in}{1.298432in}}%
\pgfpathlineto{\pgfqpoint{2.486644in}{1.292790in}}%
\pgfpathlineto{\pgfqpoint{2.490858in}{1.287608in}}%
\pgfpathlineto{\pgfqpoint{2.492536in}{1.276247in}}%
\pgfpathlineto{\pgfqpoint{2.489267in}{1.265577in}}%
\pgfpathlineto{\pgfqpoint{2.488835in}{1.249980in}}%
\pgfpathlineto{\pgfqpoint{2.485764in}{1.229680in}}%
\pgfpathlineto{\pgfqpoint{2.470230in}{1.211575in}}%
\pgfpathlineto{\pgfqpoint{2.463690in}{1.207732in}}%
\pgfpathlineto{\pgfqpoint{2.458605in}{1.211041in}}%
\pgfpathlineto{\pgfqpoint{2.446020in}{1.193784in}}%
\pgfpathlineto{\pgfqpoint{2.447199in}{1.177168in}}%
\pgfpathlineto{\pgfqpoint{2.433183in}{1.181164in}}%
\pgfpathlineto{\pgfqpoint{2.426546in}{1.164570in}}%
\pgfpathlineto{\pgfqpoint{2.429901in}{1.152811in}}%
\pgfpathlineto{\pgfqpoint{2.424649in}{1.151074in}}%
\pgfpathlineto{\pgfqpoint{2.423914in}{1.141782in}}%
\pgfpathlineto{\pgfqpoint{2.411023in}{1.138034in}}%
\pgfpathlineto{\pgfqpoint{2.404325in}{1.145536in}}%
\pgfpathlineto{\pgfqpoint{2.395706in}{1.148445in}}%
\pgfpathlineto{\pgfqpoint{2.393680in}{1.156052in}}%
\pgfpathlineto{\pgfqpoint{2.385887in}{1.154714in}}%
\pgfpathlineto{\pgfqpoint{2.380782in}{1.147717in}}%
\pgfpathlineto{\pgfqpoint{2.373478in}{1.145257in}}%
\pgfpathlineto{\pgfqpoint{2.360576in}{1.150208in}}%
\pgfpathlineto{\pgfqpoint{2.353440in}{1.144317in}}%
\pgfpathlineto{\pgfqpoint{2.343298in}{1.151287in}}%
\pgfpathlineto{\pgfqpoint{2.323842in}{1.153426in}}%
\pgfpathlineto{\pgfqpoint{2.317973in}{1.166088in}}%
\pgfpathlineto{\pgfqpoint{2.310538in}{1.171698in}}%
\pgfpathlineto{\pgfqpoint{2.303330in}{1.168107in}}%
\pgfpathlineto{\pgfqpoint{2.291994in}{1.168959in}}%
\pgfusepath{stroke}%
\end{pgfscope}%
\begin{pgfscope}%
\pgfpathrectangle{\pgfqpoint{0.100000in}{0.100000in}}{\pgfqpoint{3.007045in}{1.925000in}}%
\pgfusepath{clip}%
\pgfsetbuttcap%
\pgfsetroundjoin%
\pgfsetlinewidth{0.050187pt}%
\definecolor{currentstroke}{rgb}{1.000000,1.000000,1.000000}%
\pgfsetstrokecolor{currentstroke}%
\pgfsetdash{}{0pt}%
\pgfpathmoveto{\pgfqpoint{2.081137in}{1.010475in}}%
\pgfpathlineto{\pgfqpoint{2.073675in}{1.016537in}}%
\pgfpathlineto{\pgfqpoint{2.067590in}{1.013653in}}%
\pgfpathlineto{\pgfqpoint{2.059812in}{1.028177in}}%
\pgfpathlineto{\pgfqpoint{2.063786in}{1.037307in}}%
\pgfpathlineto{\pgfqpoint{2.058058in}{1.047582in}}%
\pgfpathlineto{\pgfqpoint{2.057751in}{1.055672in}}%
\pgfpathlineto{\pgfqpoint{2.050012in}{1.058555in}}%
\pgfpathlineto{\pgfqpoint{2.046423in}{1.064650in}}%
\pgfpathlineto{\pgfqpoint{2.031291in}{1.072263in}}%
\pgfpathlineto{\pgfqpoint{2.018150in}{1.081642in}}%
\pgfpathlineto{\pgfqpoint{2.012044in}{1.088723in}}%
\pgfpathlineto{\pgfqpoint{2.011918in}{1.097359in}}%
\pgfpathlineto{\pgfqpoint{2.020088in}{1.113940in}}%
\pgfpathlineto{\pgfqpoint{2.022602in}{1.126623in}}%
\pgfpathlineto{\pgfqpoint{2.015945in}{1.133789in}}%
\pgfpathlineto{\pgfqpoint{2.007129in}{1.136489in}}%
\pgfpathlineto{\pgfqpoint{1.996433in}{1.130585in}}%
\pgfpathlineto{\pgfqpoint{1.991906in}{1.140721in}}%
\pgfpathlineto{\pgfqpoint{1.989621in}{1.154462in}}%
\pgfpathlineto{\pgfqpoint{1.973858in}{1.166711in}}%
\pgfpathlineto{\pgfqpoint{1.970708in}{1.172146in}}%
\pgfpathlineto{\pgfqpoint{1.956311in}{1.184451in}}%
\pgfpathlineto{\pgfqpoint{1.951798in}{1.193397in}}%
\pgfpathlineto{\pgfqpoint{1.947730in}{1.211138in}}%
\pgfpathlineto{\pgfqpoint{1.950519in}{1.226895in}}%
\pgfpathlineto{\pgfqpoint{1.954240in}{1.229091in}}%
\pgfpathlineto{\pgfqpoint{1.953571in}{1.242282in}}%
\pgfpathlineto{\pgfqpoint{1.964076in}{1.246200in}}%
\pgfpathlineto{\pgfqpoint{1.967249in}{1.258037in}}%
\pgfpathlineto{\pgfqpoint{1.973292in}{1.266004in}}%
\pgfpathlineto{\pgfqpoint{1.972987in}{1.276140in}}%
\pgfpathlineto{\pgfqpoint{1.965496in}{1.284201in}}%
\pgfpathlineto{\pgfqpoint{1.967275in}{1.295485in}}%
\pgfpathlineto{\pgfqpoint{1.986693in}{1.300392in}}%
\pgfpathlineto{\pgfqpoint{2.001591in}{1.309350in}}%
\pgfpathlineto{\pgfqpoint{2.003154in}{1.320629in}}%
\pgfpathlineto{\pgfqpoint{2.008338in}{1.324180in}}%
\pgfpathlineto{\pgfqpoint{2.010321in}{1.336029in}}%
\pgfpathlineto{\pgfqpoint{2.008727in}{1.343855in}}%
\pgfpathlineto{\pgfqpoint{1.998543in}{1.350356in}}%
\pgfpathlineto{\pgfqpoint{1.994442in}{1.360047in}}%
\pgfpathlineto{\pgfqpoint{1.984391in}{1.369415in}}%
\pgfpathlineto{\pgfqpoint{2.066992in}{1.372935in}}%
\pgfpathlineto{\pgfqpoint{2.122408in}{1.376872in}}%
\pgfpathlineto{\pgfqpoint{2.121395in}{1.365215in}}%
\pgfpathlineto{\pgfqpoint{2.130839in}{1.349136in}}%
\pgfpathlineto{\pgfqpoint{2.134804in}{1.335398in}}%
\pgfpathlineto{\pgfqpoint{2.139527in}{1.327598in}}%
\pgfpathlineto{\pgfqpoint{2.144682in}{1.263329in}}%
\pgfpathlineto{\pgfqpoint{2.151841in}{1.171716in}}%
\pgfpathlineto{\pgfqpoint{2.146318in}{1.157795in}}%
\pgfpathlineto{\pgfqpoint{2.154275in}{1.146352in}}%
\pgfpathlineto{\pgfqpoint{2.156492in}{1.135423in}}%
\pgfpathlineto{\pgfqpoint{2.145520in}{1.112760in}}%
\pgfpathlineto{\pgfqpoint{2.146167in}{1.109630in}}%
\pgfpathlineto{\pgfqpoint{2.134873in}{1.096640in}}%
\pgfpathlineto{\pgfqpoint{2.137970in}{1.092672in}}%
\pgfpathlineto{\pgfqpoint{2.133271in}{1.084575in}}%
\pgfpathlineto{\pgfqpoint{2.134480in}{1.068255in}}%
\pgfpathlineto{\pgfqpoint{2.128773in}{1.058240in}}%
\pgfpathlineto{\pgfqpoint{2.133415in}{1.046407in}}%
\pgfpathlineto{\pgfqpoint{2.113967in}{1.039976in}}%
\pgfpathlineto{\pgfqpoint{2.112175in}{1.032983in}}%
\pgfpathlineto{\pgfqpoint{2.117486in}{1.024119in}}%
\pgfpathlineto{\pgfqpoint{2.115042in}{1.018342in}}%
\pgfpathlineto{\pgfqpoint{2.094189in}{1.025479in}}%
\pgfpathlineto{\pgfqpoint{2.087308in}{1.026197in}}%
\pgfpathlineto{\pgfqpoint{2.078730in}{1.015347in}}%
\pgfpathlineto{\pgfqpoint{2.081137in}{1.010475in}}%
\pgfusepath{stroke}%
\end{pgfscope}%
\begin{pgfscope}%
\pgfpathrectangle{\pgfqpoint{0.100000in}{0.100000in}}{\pgfqpoint{3.007045in}{1.925000in}}%
\pgfusepath{clip}%
\pgfsetbuttcap%
\pgfsetroundjoin%
\pgfsetlinewidth{0.050187pt}%
\definecolor{currentstroke}{rgb}{1.000000,1.000000,1.000000}%
\pgfsetstrokecolor{currentstroke}%
\pgfsetdash{}{0pt}%
\pgfpathmoveto{\pgfqpoint{2.688153in}{1.209860in}}%
\pgfpathlineto{\pgfqpoint{2.682414in}{1.218382in}}%
\pgfpathlineto{\pgfqpoint{2.685652in}{1.223151in}}%
\pgfpathlineto{\pgfqpoint{2.693583in}{1.217795in}}%
\pgfpathlineto{\pgfqpoint{2.688153in}{1.209860in}}%
\pgfusepath{stroke}%
\end{pgfscope}%
\begin{pgfscope}%
\pgfpathrectangle{\pgfqpoint{0.100000in}{0.100000in}}{\pgfqpoint{3.007045in}{1.925000in}}%
\pgfusepath{clip}%
\pgfsetbuttcap%
\pgfsetroundjoin%
\pgfsetlinewidth{0.050187pt}%
\definecolor{currentstroke}{rgb}{1.000000,1.000000,1.000000}%
\pgfsetstrokecolor{currentstroke}%
\pgfsetdash{}{0pt}%
\pgfpathmoveto{\pgfqpoint{2.738266in}{1.283189in}}%
\pgfpathlineto{\pgfqpoint{2.741696in}{1.290435in}}%
\pgfpathlineto{\pgfqpoint{2.755560in}{1.292000in}}%
\pgfpathlineto{\pgfqpoint{2.753338in}{1.285842in}}%
\pgfpathlineto{\pgfqpoint{2.748764in}{1.277976in}}%
\pgfpathlineto{\pgfqpoint{2.751863in}{1.268604in}}%
\pgfpathlineto{\pgfqpoint{2.764095in}{1.257374in}}%
\pgfpathlineto{\pgfqpoint{2.766934in}{1.245549in}}%
\pgfpathlineto{\pgfqpoint{2.781027in}{1.230823in}}%
\pgfpathlineto{\pgfqpoint{2.786556in}{1.231453in}}%
\pgfpathlineto{\pgfqpoint{2.793443in}{1.209347in}}%
\pgfpathlineto{\pgfqpoint{2.792313in}{1.209127in}}%
\pgfpathlineto{\pgfqpoint{2.791057in}{1.208880in}}%
\pgfpathlineto{\pgfqpoint{2.760342in}{1.203000in}}%
\pgfpathlineto{\pgfqpoint{2.743908in}{1.261445in}}%
\pgfpathlineto{\pgfqpoint{2.738266in}{1.283189in}}%
\pgfusepath{stroke}%
\end{pgfscope}%
\begin{pgfscope}%
\pgfpathrectangle{\pgfqpoint{0.100000in}{0.100000in}}{\pgfqpoint{3.007045in}{1.925000in}}%
\pgfusepath{clip}%
\pgfsetbuttcap%
\pgfsetroundjoin%
\pgfsetlinewidth{0.050187pt}%
\definecolor{currentstroke}{rgb}{1.000000,1.000000,1.000000}%
\pgfsetstrokecolor{currentstroke}%
\pgfsetdash{}{0pt}%
\pgfpathmoveto{\pgfqpoint{2.451373in}{1.085086in}}%
\pgfpathlineto{\pgfqpoint{2.441667in}{1.085443in}}%
\pgfpathlineto{\pgfqpoint{2.432723in}{1.091604in}}%
\pgfpathlineto{\pgfqpoint{2.420641in}{1.110328in}}%
\pgfpathlineto{\pgfqpoint{2.410356in}{1.120244in}}%
\pgfpathlineto{\pgfqpoint{2.413044in}{1.127903in}}%
\pgfpathlineto{\pgfqpoint{2.411023in}{1.138034in}}%
\pgfpathlineto{\pgfqpoint{2.423914in}{1.141782in}}%
\pgfpathlineto{\pgfqpoint{2.424649in}{1.151074in}}%
\pgfpathlineto{\pgfqpoint{2.429901in}{1.152811in}}%
\pgfpathlineto{\pgfqpoint{2.426546in}{1.164570in}}%
\pgfpathlineto{\pgfqpoint{2.433183in}{1.181164in}}%
\pgfpathlineto{\pgfqpoint{2.447199in}{1.177168in}}%
\pgfpathlineto{\pgfqpoint{2.446020in}{1.193784in}}%
\pgfpathlineto{\pgfqpoint{2.458605in}{1.211041in}}%
\pgfpathlineto{\pgfqpoint{2.463690in}{1.207732in}}%
\pgfpathlineto{\pgfqpoint{2.470230in}{1.211575in}}%
\pgfpathlineto{\pgfqpoint{2.485764in}{1.229680in}}%
\pgfpathlineto{\pgfqpoint{2.488835in}{1.249980in}}%
\pgfpathlineto{\pgfqpoint{2.489267in}{1.265577in}}%
\pgfpathlineto{\pgfqpoint{2.492536in}{1.276247in}}%
\pgfpathlineto{\pgfqpoint{2.490858in}{1.287608in}}%
\pgfpathlineto{\pgfqpoint{2.486644in}{1.292790in}}%
\pgfpathlineto{\pgfqpoint{2.493170in}{1.298432in}}%
\pgfpathlineto{\pgfqpoint{2.502630in}{1.238834in}}%
\pgfpathlineto{\pgfqpoint{2.554781in}{1.247456in}}%
\pgfpathlineto{\pgfqpoint{2.560186in}{1.213439in}}%
\pgfpathlineto{\pgfqpoint{2.579067in}{1.235892in}}%
\pgfpathlineto{\pgfqpoint{2.583515in}{1.233653in}}%
\pgfpathlineto{\pgfqpoint{2.590060in}{1.246663in}}%
\pgfpathlineto{\pgfqpoint{2.599126in}{1.242584in}}%
\pgfpathlineto{\pgfqpoint{2.607031in}{1.243360in}}%
\pgfpathlineto{\pgfqpoint{2.612246in}{1.252404in}}%
\pgfpathlineto{\pgfqpoint{2.619817in}{1.257435in}}%
\pgfpathlineto{\pgfqpoint{2.630342in}{1.253004in}}%
\pgfpathlineto{\pgfqpoint{2.637267in}{1.254535in}}%
\pgfpathlineto{\pgfqpoint{2.647188in}{1.237365in}}%
\pgfpathlineto{\pgfqpoint{2.644319in}{1.224366in}}%
\pgfpathlineto{\pgfqpoint{2.618406in}{1.239002in}}%
\pgfpathlineto{\pgfqpoint{2.613553in}{1.226987in}}%
\pgfpathlineto{\pgfqpoint{2.615154in}{1.221459in}}%
\pgfpathlineto{\pgfqpoint{2.604125in}{1.202772in}}%
\pgfpathlineto{\pgfqpoint{2.598920in}{1.199055in}}%
\pgfpathlineto{\pgfqpoint{2.596570in}{1.190815in}}%
\pgfpathlineto{\pgfqpoint{2.587696in}{1.191677in}}%
\pgfpathlineto{\pgfqpoint{2.581324in}{1.169185in}}%
\pgfpathlineto{\pgfqpoint{2.577757in}{1.164024in}}%
\pgfpathlineto{\pgfqpoint{2.568590in}{1.165745in}}%
\pgfpathlineto{\pgfqpoint{2.559215in}{1.172839in}}%
\pgfpathlineto{\pgfqpoint{2.558890in}{1.161951in}}%
\pgfpathlineto{\pgfqpoint{2.549814in}{1.143653in}}%
\pgfpathlineto{\pgfqpoint{2.547556in}{1.130626in}}%
\pgfpathlineto{\pgfqpoint{2.540593in}{1.121950in}}%
\pgfpathlineto{\pgfqpoint{2.535313in}{1.108096in}}%
\pgfpathlineto{\pgfqpoint{2.535053in}{1.095270in}}%
\pgfpathlineto{\pgfqpoint{2.526900in}{1.092874in}}%
\pgfpathlineto{\pgfqpoint{2.517650in}{1.085614in}}%
\pgfpathlineto{\pgfqpoint{2.508744in}{1.084231in}}%
\pgfpathlineto{\pgfqpoint{2.506692in}{1.078086in}}%
\pgfpathlineto{\pgfqpoint{2.492341in}{1.071792in}}%
\pgfpathlineto{\pgfqpoint{2.484322in}{1.077155in}}%
\pgfpathlineto{\pgfqpoint{2.475358in}{1.066973in}}%
\pgfpathlineto{\pgfqpoint{2.459950in}{1.069974in}}%
\pgfpathlineto{\pgfqpoint{2.450500in}{1.080661in}}%
\pgfpathlineto{\pgfqpoint{2.451373in}{1.085086in}}%
\pgfusepath{stroke}%
\end{pgfscope}%
\begin{pgfscope}%
\pgfpathrectangle{\pgfqpoint{0.100000in}{0.100000in}}{\pgfqpoint{3.007045in}{1.925000in}}%
\pgfusepath{clip}%
\pgfsetbuttcap%
\pgfsetroundjoin%
\pgfsetlinewidth{0.050187pt}%
\definecolor{currentstroke}{rgb}{1.000000,1.000000,1.000000}%
\pgfsetstrokecolor{currentstroke}%
\pgfsetdash{}{0pt}%
\pgfpathmoveto{\pgfqpoint{2.791057in}{1.208880in}}%
\pgfpathlineto{\pgfqpoint{2.790664in}{1.196870in}}%
\pgfpathlineto{\pgfqpoint{2.786042in}{1.190966in}}%
\pgfpathlineto{\pgfqpoint{2.783092in}{1.177859in}}%
\pgfpathlineto{\pgfqpoint{2.769772in}{1.171817in}}%
\pgfpathlineto{\pgfqpoint{2.758597in}{1.170057in}}%
\pgfpathlineto{\pgfqpoint{2.761853in}{1.178677in}}%
\pgfpathlineto{\pgfqpoint{2.744646in}{1.185900in}}%
\pgfpathlineto{\pgfqpoint{2.730667in}{1.194944in}}%
\pgfpathlineto{\pgfqpoint{2.739321in}{1.208462in}}%
\pgfpathlineto{\pgfqpoint{2.732385in}{1.218971in}}%
\pgfpathlineto{\pgfqpoint{2.733202in}{1.228035in}}%
\pgfpathlineto{\pgfqpoint{2.728187in}{1.230522in}}%
\pgfpathlineto{\pgfqpoint{2.724117in}{1.245248in}}%
\pgfpathlineto{\pgfqpoint{2.733741in}{1.265210in}}%
\pgfpathlineto{\pgfqpoint{2.726510in}{1.268470in}}%
\pgfpathlineto{\pgfqpoint{2.724637in}{1.258622in}}%
\pgfpathlineto{\pgfqpoint{2.714315in}{1.255880in}}%
\pgfpathlineto{\pgfqpoint{2.715401in}{1.223478in}}%
\pgfpathlineto{\pgfqpoint{2.713513in}{1.213050in}}%
\pgfpathlineto{\pgfqpoint{2.718690in}{1.198187in}}%
\pgfpathlineto{\pgfqpoint{2.726654in}{1.191036in}}%
\pgfpathlineto{\pgfqpoint{2.723043in}{1.186546in}}%
\pgfpathlineto{\pgfqpoint{2.731172in}{1.179985in}}%
\pgfpathlineto{\pgfqpoint{2.734111in}{1.169336in}}%
\pgfpathlineto{\pgfqpoint{2.719210in}{1.178155in}}%
\pgfpathlineto{\pgfqpoint{2.709780in}{1.177009in}}%
\pgfpathlineto{\pgfqpoint{2.702458in}{1.186080in}}%
\pgfpathlineto{\pgfqpoint{2.695002in}{1.186964in}}%
\pgfpathlineto{\pgfqpoint{2.684407in}{1.182420in}}%
\pgfpathlineto{\pgfqpoint{2.680297in}{1.188082in}}%
\pgfpathlineto{\pgfqpoint{2.685698in}{1.199967in}}%
\pgfpathlineto{\pgfqpoint{2.688153in}{1.209860in}}%
\pgfpathlineto{\pgfqpoint{2.693583in}{1.217795in}}%
\pgfpathlineto{\pgfqpoint{2.685652in}{1.223151in}}%
\pgfpathlineto{\pgfqpoint{2.682414in}{1.218382in}}%
\pgfpathlineto{\pgfqpoint{2.674491in}{1.223222in}}%
\pgfpathlineto{\pgfqpoint{2.660462in}{1.226448in}}%
\pgfpathlineto{\pgfqpoint{2.661734in}{1.233539in}}%
\pgfpathlineto{\pgfqpoint{2.655375in}{1.237653in}}%
\pgfpathlineto{\pgfqpoint{2.647188in}{1.237365in}}%
\pgfpathlineto{\pgfqpoint{2.637267in}{1.254535in}}%
\pgfpathlineto{\pgfqpoint{2.630342in}{1.253004in}}%
\pgfpathlineto{\pgfqpoint{2.619817in}{1.257435in}}%
\pgfpathlineto{\pgfqpoint{2.612246in}{1.252404in}}%
\pgfpathlineto{\pgfqpoint{2.607031in}{1.243360in}}%
\pgfpathlineto{\pgfqpoint{2.599126in}{1.242584in}}%
\pgfpathlineto{\pgfqpoint{2.590060in}{1.246663in}}%
\pgfpathlineto{\pgfqpoint{2.583515in}{1.233653in}}%
\pgfpathlineto{\pgfqpoint{2.579067in}{1.235892in}}%
\pgfpathlineto{\pgfqpoint{2.560186in}{1.213439in}}%
\pgfpathlineto{\pgfqpoint{2.554781in}{1.247456in}}%
\pgfpathlineto{\pgfqpoint{2.623776in}{1.260244in}}%
\pgfpathlineto{\pgfqpoint{2.654706in}{1.265782in}}%
\pgfpathlineto{\pgfqpoint{2.699704in}{1.274845in}}%
\pgfpathlineto{\pgfqpoint{2.738266in}{1.283189in}}%
\pgfpathlineto{\pgfqpoint{2.743908in}{1.261445in}}%
\pgfpathlineto{\pgfqpoint{2.760342in}{1.203000in}}%
\pgfpathlineto{\pgfqpoint{2.791057in}{1.208880in}}%
\pgfusepath{stroke}%
\end{pgfscope}%
\begin{pgfscope}%
\pgfpathrectangle{\pgfqpoint{0.100000in}{0.100000in}}{\pgfqpoint{3.007045in}{1.925000in}}%
\pgfusepath{clip}%
\pgfsetbuttcap%
\pgfsetroundjoin%
\pgfsetlinewidth{0.050187pt}%
\definecolor{currentstroke}{rgb}{1.000000,1.000000,1.000000}%
\pgfsetstrokecolor{currentstroke}%
\pgfsetdash{}{0pt}%
\pgfpathmoveto{\pgfqpoint{1.402207in}{1.014889in}}%
\pgfpathlineto{\pgfqpoint{1.351869in}{1.019693in}}%
\pgfpathlineto{\pgfqpoint{1.299648in}{1.024321in}}%
\pgfpathlineto{\pgfqpoint{1.239303in}{1.030616in}}%
\pgfpathlineto{\pgfqpoint{1.149648in}{1.041279in}}%
\pgfpathlineto{\pgfqpoint{1.117842in}{1.046179in}}%
\pgfpathlineto{\pgfqpoint{1.035744in}{1.058103in}}%
\pgfpathlineto{\pgfqpoint{1.047732in}{1.133342in}}%
\pgfpathlineto{\pgfqpoint{1.048031in}{1.139428in}}%
\pgfpathlineto{\pgfqpoint{1.059579in}{1.212028in}}%
\pgfpathlineto{\pgfqpoint{1.068205in}{1.267320in}}%
\pgfpathlineto{\pgfqpoint{1.076345in}{1.318594in}}%
\pgfpathlineto{\pgfqpoint{1.131961in}{1.310611in}}%
\pgfpathlineto{\pgfqpoint{1.183815in}{1.303168in}}%
\pgfpathlineto{\pgfqpoint{1.279137in}{1.291403in}}%
\pgfpathlineto{\pgfqpoint{1.322867in}{1.287410in}}%
\pgfpathlineto{\pgfqpoint{1.392146in}{1.280688in}}%
\pgfpathlineto{\pgfqpoint{1.422117in}{1.278255in}}%
\pgfpathlineto{\pgfqpoint{1.416825in}{1.212637in}}%
\pgfpathlineto{\pgfqpoint{1.412045in}{1.149457in}}%
\pgfpathlineto{\pgfqpoint{1.408197in}{1.098013in}}%
\pgfpathlineto{\pgfqpoint{1.402207in}{1.014889in}}%
\pgfusepath{stroke}%
\end{pgfscope}%
\begin{pgfscope}%
\pgfpathrectangle{\pgfqpoint{0.100000in}{0.100000in}}{\pgfqpoint{3.007045in}{1.925000in}}%
\pgfusepath{clip}%
\pgfsetbuttcap%
\pgfsetroundjoin%
\pgfsetlinewidth{0.050187pt}%
\definecolor{currentstroke}{rgb}{1.000000,1.000000,1.000000}%
\pgfsetstrokecolor{currentstroke}%
\pgfsetdash{}{0pt}%
\pgfpathmoveto{\pgfqpoint{2.068136in}{0.977587in}}%
\pgfpathlineto{\pgfqpoint{2.069833in}{0.985194in}}%
\pgfpathlineto{\pgfqpoint{2.078628in}{0.983482in}}%
\pgfpathlineto{\pgfqpoint{2.081137in}{1.010475in}}%
\pgfpathlineto{\pgfqpoint{2.078730in}{1.015347in}}%
\pgfpathlineto{\pgfqpoint{2.087308in}{1.026197in}}%
\pgfpathlineto{\pgfqpoint{2.094189in}{1.025479in}}%
\pgfpathlineto{\pgfqpoint{2.115042in}{1.018342in}}%
\pgfpathlineto{\pgfqpoint{2.117486in}{1.024119in}}%
\pgfpathlineto{\pgfqpoint{2.112175in}{1.032983in}}%
\pgfpathlineto{\pgfqpoint{2.113967in}{1.039976in}}%
\pgfpathlineto{\pgfqpoint{2.133415in}{1.046407in}}%
\pgfpathlineto{\pgfqpoint{2.128773in}{1.058240in}}%
\pgfpathlineto{\pgfqpoint{2.134480in}{1.068255in}}%
\pgfpathlineto{\pgfqpoint{2.145103in}{1.073876in}}%
\pgfpathlineto{\pgfqpoint{2.167393in}{1.079442in}}%
\pgfpathlineto{\pgfqpoint{2.181457in}{1.071040in}}%
\pgfpathlineto{\pgfqpoint{2.185723in}{1.079234in}}%
\pgfpathlineto{\pgfqpoint{2.197401in}{1.084958in}}%
\pgfpathlineto{\pgfqpoint{2.200414in}{1.079910in}}%
\pgfpathlineto{\pgfqpoint{2.212410in}{1.083920in}}%
\pgfpathlineto{\pgfqpoint{2.211653in}{1.090784in}}%
\pgfpathlineto{\pgfqpoint{2.222438in}{1.098640in}}%
\pgfpathlineto{\pgfqpoint{2.228803in}{1.090482in}}%
\pgfpathlineto{\pgfqpoint{2.237177in}{1.089662in}}%
\pgfpathlineto{\pgfqpoint{2.242742in}{1.094978in}}%
\pgfpathlineto{\pgfqpoint{2.242123in}{1.102538in}}%
\pgfpathlineto{\pgfqpoint{2.246853in}{1.110045in}}%
\pgfpathlineto{\pgfqpoint{2.253191in}{1.111671in}}%
\pgfpathlineto{\pgfqpoint{2.255736in}{1.121584in}}%
\pgfpathlineto{\pgfqpoint{2.264955in}{1.130153in}}%
\pgfpathlineto{\pgfqpoint{2.262180in}{1.138682in}}%
\pgfpathlineto{\pgfqpoint{2.271148in}{1.142929in}}%
\pgfpathlineto{\pgfqpoint{2.277140in}{1.140317in}}%
\pgfpathlineto{\pgfqpoint{2.285991in}{1.146952in}}%
\pgfpathlineto{\pgfqpoint{2.293900in}{1.148681in}}%
\pgfpathlineto{\pgfqpoint{2.293928in}{1.155563in}}%
\pgfpathlineto{\pgfqpoint{2.288386in}{1.165112in}}%
\pgfpathlineto{\pgfqpoint{2.291994in}{1.168959in}}%
\pgfpathlineto{\pgfqpoint{2.303330in}{1.168107in}}%
\pgfpathlineto{\pgfqpoint{2.310538in}{1.171698in}}%
\pgfpathlineto{\pgfqpoint{2.317973in}{1.166088in}}%
\pgfpathlineto{\pgfqpoint{2.323842in}{1.153426in}}%
\pgfpathlineto{\pgfqpoint{2.343298in}{1.151287in}}%
\pgfpathlineto{\pgfqpoint{2.353440in}{1.144317in}}%
\pgfpathlineto{\pgfqpoint{2.360576in}{1.150208in}}%
\pgfpathlineto{\pgfqpoint{2.373478in}{1.145257in}}%
\pgfpathlineto{\pgfqpoint{2.380782in}{1.147717in}}%
\pgfpathlineto{\pgfqpoint{2.385887in}{1.154714in}}%
\pgfpathlineto{\pgfqpoint{2.393680in}{1.156052in}}%
\pgfpathlineto{\pgfqpoint{2.395706in}{1.148445in}}%
\pgfpathlineto{\pgfqpoint{2.404325in}{1.145536in}}%
\pgfpathlineto{\pgfqpoint{2.411023in}{1.138034in}}%
\pgfpathlineto{\pgfqpoint{2.413044in}{1.127903in}}%
\pgfpathlineto{\pgfqpoint{2.410356in}{1.120244in}}%
\pgfpathlineto{\pgfqpoint{2.420641in}{1.110328in}}%
\pgfpathlineto{\pgfqpoint{2.432723in}{1.091604in}}%
\pgfpathlineto{\pgfqpoint{2.441667in}{1.085443in}}%
\pgfpathlineto{\pgfqpoint{2.451373in}{1.085086in}}%
\pgfpathlineto{\pgfqpoint{2.433451in}{1.064522in}}%
\pgfpathlineto{\pgfqpoint{2.415810in}{1.052075in}}%
\pgfpathlineto{\pgfqpoint{2.409317in}{1.042191in}}%
\pgfpathlineto{\pgfqpoint{2.409427in}{1.036826in}}%
\pgfpathlineto{\pgfqpoint{2.399872in}{1.032715in}}%
\pgfpathlineto{\pgfqpoint{2.397134in}{1.024978in}}%
\pgfpathlineto{\pgfqpoint{2.377229in}{1.017186in}}%
\pgfpathlineto{\pgfqpoint{2.370173in}{1.012125in}}%
\pgfpathlineto{\pgfqpoint{2.369218in}{1.011045in}}%
\pgfpathlineto{\pgfqpoint{2.311961in}{1.005616in}}%
\pgfpathlineto{\pgfqpoint{2.277384in}{1.002708in}}%
\pgfpathlineto{\pgfqpoint{2.220644in}{0.999497in}}%
\pgfpathlineto{\pgfqpoint{2.149956in}{0.992504in}}%
\pgfpathlineto{\pgfqpoint{2.138278in}{0.994121in}}%
\pgfpathlineto{\pgfqpoint{2.140708in}{0.982200in}}%
\pgfpathlineto{\pgfqpoint{2.068136in}{0.977587in}}%
\pgfusepath{stroke}%
\end{pgfscope}%
\begin{pgfscope}%
\pgfpathrectangle{\pgfqpoint{0.100000in}{0.100000in}}{\pgfqpoint{3.007045in}{1.925000in}}%
\pgfusepath{clip}%
\pgfsetbuttcap%
\pgfsetroundjoin%
\pgfsetlinewidth{0.050187pt}%
\definecolor{currentstroke}{rgb}{1.000000,1.000000,1.000000}%
\pgfsetstrokecolor{currentstroke}%
\pgfsetdash{}{0pt}%
\pgfpathmoveto{\pgfqpoint{1.402207in}{1.014889in}}%
\pgfpathlineto{\pgfqpoint{1.408197in}{1.098013in}}%
\pgfpathlineto{\pgfqpoint{1.412045in}{1.149457in}}%
\pgfpathlineto{\pgfqpoint{1.416825in}{1.212637in}}%
\pgfpathlineto{\pgfqpoint{1.483096in}{1.207983in}}%
\pgfpathlineto{\pgfqpoint{1.567259in}{1.203386in}}%
\pgfpathlineto{\pgfqpoint{1.624502in}{1.201190in}}%
\pgfpathlineto{\pgfqpoint{1.681409in}{1.199446in}}%
\pgfpathlineto{\pgfqpoint{1.732915in}{1.198649in}}%
\pgfpathlineto{\pgfqpoint{1.756737in}{1.198896in}}%
\pgfpathlineto{\pgfqpoint{1.767223in}{1.190339in}}%
\pgfpathlineto{\pgfqpoint{1.775436in}{1.192065in}}%
\pgfpathlineto{\pgfqpoint{1.775693in}{1.184601in}}%
\pgfpathlineto{\pgfqpoint{1.769593in}{1.171690in}}%
\pgfpathlineto{\pgfqpoint{1.770255in}{1.163533in}}%
\pgfpathlineto{\pgfqpoint{1.776095in}{1.158154in}}%
\pgfpathlineto{\pgfqpoint{1.781459in}{1.146947in}}%
\pgfpathlineto{\pgfqpoint{1.792341in}{1.140513in}}%
\pgfpathlineto{\pgfqpoint{1.791978in}{1.098269in}}%
\pgfpathlineto{\pgfqpoint{1.792298in}{1.001125in}}%
\pgfpathlineto{\pgfqpoint{1.719359in}{1.001461in}}%
\pgfpathlineto{\pgfqpoint{1.642551in}{1.002805in}}%
\pgfpathlineto{\pgfqpoint{1.586009in}{1.004800in}}%
\pgfpathlineto{\pgfqpoint{1.538812in}{1.006621in}}%
\pgfpathlineto{\pgfqpoint{1.459299in}{1.011309in}}%
\pgfpathlineto{\pgfqpoint{1.402207in}{1.014889in}}%
\pgfusepath{stroke}%
\end{pgfscope}%
\begin{pgfscope}%
\pgfpathrectangle{\pgfqpoint{0.100000in}{0.100000in}}{\pgfqpoint{3.007045in}{1.925000in}}%
\pgfusepath{clip}%
\pgfsetbuttcap%
\pgfsetroundjoin%
\pgfsetlinewidth{0.050187pt}%
\definecolor{currentstroke}{rgb}{1.000000,1.000000,1.000000}%
\pgfsetstrokecolor{currentstroke}%
\pgfsetdash{}{0pt}%
\pgfpathmoveto{\pgfqpoint{2.475422in}{1.025896in}}%
\pgfpathlineto{\pgfqpoint{2.394637in}{1.014547in}}%
\pgfpathlineto{\pgfqpoint{2.370173in}{1.012125in}}%
\pgfpathlineto{\pgfqpoint{2.377229in}{1.017186in}}%
\pgfpathlineto{\pgfqpoint{2.397134in}{1.024978in}}%
\pgfpathlineto{\pgfqpoint{2.399872in}{1.032715in}}%
\pgfpathlineto{\pgfqpoint{2.409427in}{1.036826in}}%
\pgfpathlineto{\pgfqpoint{2.409317in}{1.042191in}}%
\pgfpathlineto{\pgfqpoint{2.415810in}{1.052075in}}%
\pgfpathlineto{\pgfqpoint{2.433451in}{1.064522in}}%
\pgfpathlineto{\pgfqpoint{2.451373in}{1.085086in}}%
\pgfpathlineto{\pgfqpoint{2.450500in}{1.080661in}}%
\pgfpathlineto{\pgfqpoint{2.459950in}{1.069974in}}%
\pgfpathlineto{\pgfqpoint{2.475358in}{1.066973in}}%
\pgfpathlineto{\pgfqpoint{2.484322in}{1.077155in}}%
\pgfpathlineto{\pgfqpoint{2.492341in}{1.071792in}}%
\pgfpathlineto{\pgfqpoint{2.506692in}{1.078086in}}%
\pgfpathlineto{\pgfqpoint{2.508744in}{1.084231in}}%
\pgfpathlineto{\pgfqpoint{2.517650in}{1.085614in}}%
\pgfpathlineto{\pgfqpoint{2.526900in}{1.092874in}}%
\pgfpathlineto{\pgfqpoint{2.535053in}{1.095270in}}%
\pgfpathlineto{\pgfqpoint{2.535313in}{1.108096in}}%
\pgfpathlineto{\pgfqpoint{2.540593in}{1.121950in}}%
\pgfpathlineto{\pgfqpoint{2.547556in}{1.130626in}}%
\pgfpathlineto{\pgfqpoint{2.549814in}{1.143653in}}%
\pgfpathlineto{\pgfqpoint{2.558890in}{1.161951in}}%
\pgfpathlineto{\pgfqpoint{2.559215in}{1.172839in}}%
\pgfpathlineto{\pgfqpoint{2.568590in}{1.165745in}}%
\pgfpathlineto{\pgfqpoint{2.577757in}{1.164024in}}%
\pgfpathlineto{\pgfqpoint{2.581324in}{1.169185in}}%
\pgfpathlineto{\pgfqpoint{2.587696in}{1.191677in}}%
\pgfpathlineto{\pgfqpoint{2.596570in}{1.190815in}}%
\pgfpathlineto{\pgfqpoint{2.598920in}{1.199055in}}%
\pgfpathlineto{\pgfqpoint{2.604125in}{1.202772in}}%
\pgfpathlineto{\pgfqpoint{2.615154in}{1.221459in}}%
\pgfpathlineto{\pgfqpoint{2.613553in}{1.226987in}}%
\pgfpathlineto{\pgfqpoint{2.618406in}{1.239002in}}%
\pgfpathlineto{\pgfqpoint{2.644319in}{1.224366in}}%
\pgfpathlineto{\pgfqpoint{2.647188in}{1.237365in}}%
\pgfpathlineto{\pgfqpoint{2.655375in}{1.237653in}}%
\pgfpathlineto{\pgfqpoint{2.661734in}{1.233539in}}%
\pgfpathlineto{\pgfqpoint{2.660462in}{1.226448in}}%
\pgfpathlineto{\pgfqpoint{2.674491in}{1.223222in}}%
\pgfpathlineto{\pgfqpoint{2.682414in}{1.218382in}}%
\pgfpathlineto{\pgfqpoint{2.688153in}{1.209860in}}%
\pgfpathlineto{\pgfqpoint{2.685698in}{1.199967in}}%
\pgfpathlineto{\pgfqpoint{2.680740in}{1.199156in}}%
\pgfpathlineto{\pgfqpoint{2.677865in}{1.184228in}}%
\pgfpathlineto{\pgfqpoint{2.684163in}{1.178391in}}%
\pgfpathlineto{\pgfqpoint{2.693030in}{1.183109in}}%
\pgfpathlineto{\pgfqpoint{2.701257in}{1.173146in}}%
\pgfpathlineto{\pgfqpoint{2.719638in}{1.171355in}}%
\pgfpathlineto{\pgfqpoint{2.722805in}{1.165624in}}%
\pgfpathlineto{\pgfqpoint{2.739815in}{1.160050in}}%
\pgfpathlineto{\pgfqpoint{2.737716in}{1.153473in}}%
\pgfpathlineto{\pgfqpoint{2.739929in}{1.134884in}}%
\pgfpathlineto{\pgfqpoint{2.746795in}{1.127865in}}%
\pgfpathlineto{\pgfqpoint{2.736664in}{1.124902in}}%
\pgfpathlineto{\pgfqpoint{2.738007in}{1.116967in}}%
\pgfpathlineto{\pgfqpoint{2.748815in}{1.110238in}}%
\pgfpathlineto{\pgfqpoint{2.749789in}{1.103601in}}%
\pgfpathlineto{\pgfqpoint{2.743687in}{1.098614in}}%
\pgfpathlineto{\pgfqpoint{2.731375in}{1.110429in}}%
\pgfpathlineto{\pgfqpoint{2.730158in}{1.101808in}}%
\pgfpathlineto{\pgfqpoint{2.740465in}{1.097709in}}%
\pgfpathlineto{\pgfqpoint{2.741496in}{1.093469in}}%
\pgfpathlineto{\pgfqpoint{2.755633in}{1.099072in}}%
\pgfpathlineto{\pgfqpoint{2.766466in}{1.100555in}}%
\pgfpathlineto{\pgfqpoint{2.777563in}{1.078165in}}%
\pgfpathlineto{\pgfqpoint{2.776325in}{1.077922in}}%
\pgfpathlineto{\pgfqpoint{2.771311in}{1.076886in}}%
\pgfpathlineto{\pgfqpoint{2.769833in}{1.076577in}}%
\pgfpathlineto{\pgfqpoint{2.768856in}{1.076386in}}%
\pgfpathlineto{\pgfqpoint{2.702785in}{1.062681in}}%
\pgfpathlineto{\pgfqpoint{2.643688in}{1.050585in}}%
\pgfpathlineto{\pgfqpoint{2.561981in}{1.036510in}}%
\pgfpathlineto{\pgfqpoint{2.492587in}{1.027345in}}%
\pgfpathlineto{\pgfqpoint{2.475422in}{1.025896in}}%
\pgfusepath{stroke}%
\end{pgfscope}%
\begin{pgfscope}%
\pgfpathrectangle{\pgfqpoint{0.100000in}{0.100000in}}{\pgfqpoint{3.007045in}{1.925000in}}%
\pgfusepath{clip}%
\pgfsetbuttcap%
\pgfsetroundjoin%
\pgfsetlinewidth{0.050187pt}%
\definecolor{currentstroke}{rgb}{1.000000,1.000000,1.000000}%
\pgfsetstrokecolor{currentstroke}%
\pgfsetdash{}{0pt}%
\pgfpathmoveto{\pgfqpoint{2.769772in}{1.171817in}}%
\pgfpathlineto{\pgfqpoint{2.783092in}{1.177859in}}%
\pgfpathlineto{\pgfqpoint{2.772455in}{1.146743in}}%
\pgfpathlineto{\pgfqpoint{2.765421in}{1.130265in}}%
\pgfpathlineto{\pgfqpoint{2.760353in}{1.139590in}}%
\pgfpathlineto{\pgfqpoint{2.769365in}{1.161915in}}%
\pgfpathlineto{\pgfqpoint{2.769772in}{1.171817in}}%
\pgfusepath{stroke}%
\end{pgfscope}%
\begin{pgfscope}%
\pgfpathrectangle{\pgfqpoint{0.100000in}{0.100000in}}{\pgfqpoint{3.007045in}{1.925000in}}%
\pgfusepath{clip}%
\pgfsetbuttcap%
\pgfsetroundjoin%
\pgfsetlinewidth{0.050187pt}%
\definecolor{currentstroke}{rgb}{1.000000,1.000000,1.000000}%
\pgfsetstrokecolor{currentstroke}%
\pgfsetdash{}{0pt}%
\pgfpathmoveto{\pgfqpoint{2.068136in}{0.977587in}}%
\pgfpathlineto{\pgfqpoint{2.064919in}{0.977122in}}%
\pgfpathlineto{\pgfqpoint{2.061886in}{0.976909in}}%
\pgfpathlineto{\pgfqpoint{2.063184in}{0.967596in}}%
\pgfpathlineto{\pgfqpoint{2.055358in}{0.958208in}}%
\pgfpathlineto{\pgfqpoint{2.053816in}{0.943518in}}%
\pgfpathlineto{\pgfqpoint{2.018831in}{0.940931in}}%
\pgfpathlineto{\pgfqpoint{2.021881in}{0.947828in}}%
\pgfpathlineto{\pgfqpoint{2.034495in}{0.960424in}}%
\pgfpathlineto{\pgfqpoint{2.034844in}{0.967726in}}%
\pgfpathlineto{\pgfqpoint{2.029259in}{0.974639in}}%
\pgfpathlineto{\pgfqpoint{1.977188in}{0.971886in}}%
\pgfpathlineto{\pgfqpoint{1.886150in}{0.968996in}}%
\pgfpathlineto{\pgfqpoint{1.832874in}{0.968025in}}%
\pgfpathlineto{\pgfqpoint{1.792603in}{0.967662in}}%
\pgfpathlineto{\pgfqpoint{1.792298in}{1.001125in}}%
\pgfpathlineto{\pgfqpoint{1.791978in}{1.098269in}}%
\pgfpathlineto{\pgfqpoint{1.792341in}{1.140513in}}%
\pgfpathlineto{\pgfqpoint{1.781459in}{1.146947in}}%
\pgfpathlineto{\pgfqpoint{1.776095in}{1.158154in}}%
\pgfpathlineto{\pgfqpoint{1.770255in}{1.163533in}}%
\pgfpathlineto{\pgfqpoint{1.769593in}{1.171690in}}%
\pgfpathlineto{\pgfqpoint{1.775693in}{1.184601in}}%
\pgfpathlineto{\pgfqpoint{1.775436in}{1.192065in}}%
\pgfpathlineto{\pgfqpoint{1.767223in}{1.190339in}}%
\pgfpathlineto{\pgfqpoint{1.756737in}{1.198896in}}%
\pgfpathlineto{\pgfqpoint{1.748332in}{1.213916in}}%
\pgfpathlineto{\pgfqpoint{1.741281in}{1.220848in}}%
\pgfpathlineto{\pgfqpoint{1.733914in}{1.237878in}}%
\pgfpathlineto{\pgfqpoint{1.810432in}{1.236685in}}%
\pgfpathlineto{\pgfqpoint{1.886493in}{1.239191in}}%
\pgfpathlineto{\pgfqpoint{1.935262in}{1.242004in}}%
\pgfpathlineto{\pgfqpoint{1.950519in}{1.226895in}}%
\pgfpathlineto{\pgfqpoint{1.947730in}{1.211138in}}%
\pgfpathlineto{\pgfqpoint{1.951798in}{1.193397in}}%
\pgfpathlineto{\pgfqpoint{1.956311in}{1.184451in}}%
\pgfpathlineto{\pgfqpoint{1.970708in}{1.172146in}}%
\pgfpathlineto{\pgfqpoint{1.973858in}{1.166711in}}%
\pgfpathlineto{\pgfqpoint{1.989621in}{1.154462in}}%
\pgfpathlineto{\pgfqpoint{1.991906in}{1.140721in}}%
\pgfpathlineto{\pgfqpoint{1.996433in}{1.130585in}}%
\pgfpathlineto{\pgfqpoint{2.007129in}{1.136489in}}%
\pgfpathlineto{\pgfqpoint{2.015945in}{1.133789in}}%
\pgfpathlineto{\pgfqpoint{2.022602in}{1.126623in}}%
\pgfpathlineto{\pgfqpoint{2.020088in}{1.113940in}}%
\pgfpathlineto{\pgfqpoint{2.011918in}{1.097359in}}%
\pgfpathlineto{\pgfqpoint{2.012044in}{1.088723in}}%
\pgfpathlineto{\pgfqpoint{2.018150in}{1.081642in}}%
\pgfpathlineto{\pgfqpoint{2.031291in}{1.072263in}}%
\pgfpathlineto{\pgfqpoint{2.046423in}{1.064650in}}%
\pgfpathlineto{\pgfqpoint{2.050012in}{1.058555in}}%
\pgfpathlineto{\pgfqpoint{2.057751in}{1.055672in}}%
\pgfpathlineto{\pgfqpoint{2.058058in}{1.047582in}}%
\pgfpathlineto{\pgfqpoint{2.063786in}{1.037307in}}%
\pgfpathlineto{\pgfqpoint{2.059812in}{1.028177in}}%
\pgfpathlineto{\pgfqpoint{2.067590in}{1.013653in}}%
\pgfpathlineto{\pgfqpoint{2.073675in}{1.016537in}}%
\pgfpathlineto{\pgfqpoint{2.081137in}{1.010475in}}%
\pgfpathlineto{\pgfqpoint{2.078628in}{0.983482in}}%
\pgfpathlineto{\pgfqpoint{2.069833in}{0.985194in}}%
\pgfpathlineto{\pgfqpoint{2.068136in}{0.977587in}}%
\pgfusepath{stroke}%
\end{pgfscope}%
\begin{pgfscope}%
\pgfpathrectangle{\pgfqpoint{0.100000in}{0.100000in}}{\pgfqpoint{3.007045in}{1.925000in}}%
\pgfusepath{clip}%
\pgfsetbuttcap%
\pgfsetroundjoin%
\pgfsetlinewidth{0.050187pt}%
\definecolor{currentstroke}{rgb}{1.000000,1.000000,1.000000}%
\pgfsetstrokecolor{currentstroke}%
\pgfsetdash{}{0pt}%
\pgfpathmoveto{\pgfqpoint{0.719235in}{0.983407in}}%
\pgfpathlineto{\pgfqpoint{0.725604in}{0.996967in}}%
\pgfpathlineto{\pgfqpoint{0.723861in}{1.017321in}}%
\pgfpathlineto{\pgfqpoint{0.726911in}{1.023110in}}%
\pgfpathlineto{\pgfqpoint{0.727212in}{1.048463in}}%
\pgfpathlineto{\pgfqpoint{0.730103in}{1.055764in}}%
\pgfpathlineto{\pgfqpoint{0.740288in}{1.056900in}}%
\pgfpathlineto{\pgfqpoint{0.749739in}{1.053659in}}%
\pgfpathlineto{\pgfqpoint{0.753855in}{1.044733in}}%
\pgfpathlineto{\pgfqpoint{0.759627in}{1.045075in}}%
\pgfpathlineto{\pgfqpoint{0.766797in}{1.055300in}}%
\pgfpathlineto{\pgfqpoint{0.777160in}{1.105704in}}%
\pgfpathlineto{\pgfqpoint{0.836134in}{1.093528in}}%
\pgfpathlineto{\pgfqpoint{0.870361in}{1.086778in}}%
\pgfpathlineto{\pgfqpoint{0.961251in}{1.070801in}}%
\pgfpathlineto{\pgfqpoint{0.986397in}{1.065751in}}%
\pgfpathlineto{\pgfqpoint{1.035744in}{1.058103in}}%
\pgfpathlineto{\pgfqpoint{1.025626in}{0.992952in}}%
\pgfpathlineto{\pgfqpoint{1.015102in}{0.924993in}}%
\pgfpathlineto{\pgfqpoint{1.002983in}{0.848537in}}%
\pgfpathlineto{\pgfqpoint{0.989353in}{0.760774in}}%
\pgfpathlineto{\pgfqpoint{0.978343in}{0.688693in}}%
\pgfpathlineto{\pgfqpoint{0.899477in}{0.701217in}}%
\pgfpathlineto{\pgfqpoint{0.864826in}{0.707151in}}%
\pgfpathlineto{\pgfqpoint{0.849347in}{0.716418in}}%
\pgfpathlineto{\pgfqpoint{0.748420in}{0.777313in}}%
\pgfpathlineto{\pgfqpoint{0.672680in}{0.823525in}}%
\pgfpathlineto{\pgfqpoint{0.675202in}{0.831715in}}%
\pgfpathlineto{\pgfqpoint{0.681459in}{0.837446in}}%
\pgfpathlineto{\pgfqpoint{0.687967in}{0.836384in}}%
\pgfpathlineto{\pgfqpoint{0.697405in}{0.842403in}}%
\pgfpathlineto{\pgfqpoint{0.698891in}{0.851052in}}%
\pgfpathlineto{\pgfqpoint{0.687380in}{0.861545in}}%
\pgfpathlineto{\pgfqpoint{0.695567in}{0.881688in}}%
\pgfpathlineto{\pgfqpoint{0.703798in}{0.889434in}}%
\pgfpathlineto{\pgfqpoint{0.707704in}{0.898624in}}%
\pgfpathlineto{\pgfqpoint{0.710111in}{0.915475in}}%
\pgfpathlineto{\pgfqpoint{0.717829in}{0.923106in}}%
\pgfpathlineto{\pgfqpoint{0.734029in}{0.930728in}}%
\pgfpathlineto{\pgfqpoint{0.733671in}{0.937444in}}%
\pgfpathlineto{\pgfqpoint{0.724623in}{0.945778in}}%
\pgfpathlineto{\pgfqpoint{0.723391in}{0.962960in}}%
\pgfpathlineto{\pgfqpoint{0.717162in}{0.975520in}}%
\pgfpathlineto{\pgfqpoint{0.719235in}{0.983407in}}%
\pgfusepath{stroke}%
\end{pgfscope}%
\begin{pgfscope}%
\pgfpathrectangle{\pgfqpoint{0.100000in}{0.100000in}}{\pgfqpoint{3.007045in}{1.925000in}}%
\pgfusepath{clip}%
\pgfsetbuttcap%
\pgfsetroundjoin%
\pgfsetlinewidth{0.050187pt}%
\definecolor{currentstroke}{rgb}{1.000000,1.000000,1.000000}%
\pgfsetstrokecolor{currentstroke}%
\pgfsetdash{}{0pt}%
\pgfpathmoveto{\pgfqpoint{1.801141in}{0.779025in}}%
\pgfpathlineto{\pgfqpoint{1.786521in}{0.783547in}}%
\pgfpathlineto{\pgfqpoint{1.767788in}{0.797893in}}%
\pgfpathlineto{\pgfqpoint{1.759467in}{0.800976in}}%
\pgfpathlineto{\pgfqpoint{1.754180in}{0.794780in}}%
\pgfpathlineto{\pgfqpoint{1.747503in}{0.794467in}}%
\pgfpathlineto{\pgfqpoint{1.739055in}{0.799728in}}%
\pgfpathlineto{\pgfqpoint{1.725799in}{0.792988in}}%
\pgfpathlineto{\pgfqpoint{1.721080in}{0.796307in}}%
\pgfpathlineto{\pgfqpoint{1.708036in}{0.788443in}}%
\pgfpathlineto{\pgfqpoint{1.694275in}{0.789894in}}%
\pgfpathlineto{\pgfqpoint{1.683661in}{0.795004in}}%
\pgfpathlineto{\pgfqpoint{1.676227in}{0.793081in}}%
\pgfpathlineto{\pgfqpoint{1.664329in}{0.800308in}}%
\pgfpathlineto{\pgfqpoint{1.657711in}{0.787562in}}%
\pgfpathlineto{\pgfqpoint{1.650942in}{0.798523in}}%
\pgfpathlineto{\pgfqpoint{1.637572in}{0.794271in}}%
\pgfpathlineto{\pgfqpoint{1.625879in}{0.804682in}}%
\pgfpathlineto{\pgfqpoint{1.615661in}{0.796290in}}%
\pgfpathlineto{\pgfqpoint{1.610544in}{0.804002in}}%
\pgfpathlineto{\pgfqpoint{1.603167in}{0.806507in}}%
\pgfpathlineto{\pgfqpoint{1.598675in}{0.813942in}}%
\pgfpathlineto{\pgfqpoint{1.589028in}{0.816058in}}%
\pgfpathlineto{\pgfqpoint{1.583461in}{0.810461in}}%
\pgfpathlineto{\pgfqpoint{1.574003in}{0.817714in}}%
\pgfpathlineto{\pgfqpoint{1.569594in}{0.816058in}}%
\pgfpathlineto{\pgfqpoint{1.553918in}{0.821935in}}%
\pgfpathlineto{\pgfqpoint{1.544100in}{0.822590in}}%
\pgfpathlineto{\pgfqpoint{1.539706in}{0.835071in}}%
\pgfpathlineto{\pgfqpoint{1.531836in}{0.833517in}}%
\pgfpathlineto{\pgfqpoint{1.522821in}{0.836596in}}%
\pgfpathlineto{\pgfqpoint{1.516885in}{0.834817in}}%
\pgfpathlineto{\pgfqpoint{1.504154in}{0.848834in}}%
\pgfpathlineto{\pgfqpoint{1.500607in}{0.847917in}}%
\pgfpathlineto{\pgfqpoint{1.503828in}{0.904762in}}%
\pgfpathlineto{\pgfqpoint{1.507338in}{0.975120in}}%
\pgfpathlineto{\pgfqpoint{1.449726in}{0.978343in}}%
\pgfpathlineto{\pgfqpoint{1.392875in}{0.982641in}}%
\pgfpathlineto{\pgfqpoint{1.348953in}{0.986455in}}%
\pgfpathlineto{\pgfqpoint{1.351869in}{1.019693in}}%
\pgfpathlineto{\pgfqpoint{1.402207in}{1.014889in}}%
\pgfpathlineto{\pgfqpoint{1.459299in}{1.011309in}}%
\pgfpathlineto{\pgfqpoint{1.538812in}{1.006621in}}%
\pgfpathlineto{\pgfqpoint{1.586009in}{1.004800in}}%
\pgfpathlineto{\pgfqpoint{1.642551in}{1.002805in}}%
\pgfpathlineto{\pgfqpoint{1.719359in}{1.001461in}}%
\pgfpathlineto{\pgfqpoint{1.792298in}{1.001125in}}%
\pgfpathlineto{\pgfqpoint{1.792603in}{0.967662in}}%
\pgfpathlineto{\pgfqpoint{1.796697in}{0.942452in}}%
\pgfpathlineto{\pgfqpoint{1.803058in}{0.895889in}}%
\pgfpathlineto{\pgfqpoint{1.801746in}{0.816372in}}%
\pgfpathlineto{\pgfqpoint{1.801141in}{0.779025in}}%
\pgfusepath{stroke}%
\end{pgfscope}%
\begin{pgfscope}%
\pgfpathrectangle{\pgfqpoint{0.100000in}{0.100000in}}{\pgfqpoint{3.007045in}{1.925000in}}%
\pgfusepath{clip}%
\pgfsetbuttcap%
\pgfsetroundjoin%
\pgfsetlinewidth{0.050187pt}%
\definecolor{currentstroke}{rgb}{1.000000,1.000000,1.000000}%
\pgfsetstrokecolor{currentstroke}%
\pgfsetdash{}{0pt}%
\pgfpathmoveto{\pgfqpoint{2.348542in}{0.902516in}}%
\pgfpathlineto{\pgfqpoint{2.348585in}{0.917250in}}%
\pgfpathlineto{\pgfqpoint{2.361388in}{0.922908in}}%
\pgfpathlineto{\pgfqpoint{2.361925in}{0.931952in}}%
\pgfpathlineto{\pgfqpoint{2.373389in}{0.943000in}}%
\pgfpathlineto{\pgfqpoint{2.387718in}{0.945263in}}%
\pgfpathlineto{\pgfqpoint{2.400316in}{0.957439in}}%
\pgfpathlineto{\pgfqpoint{2.413472in}{0.962881in}}%
\pgfpathlineto{\pgfqpoint{2.423346in}{0.979220in}}%
\pgfpathlineto{\pgfqpoint{2.432339in}{0.978033in}}%
\pgfpathlineto{\pgfqpoint{2.451345in}{0.992920in}}%
\pgfpathlineto{\pgfqpoint{2.461381in}{0.993200in}}%
\pgfpathlineto{\pgfqpoint{2.465620in}{1.004543in}}%
\pgfpathlineto{\pgfqpoint{2.473597in}{1.012440in}}%
\pgfpathlineto{\pgfqpoint{2.475422in}{1.025896in}}%
\pgfpathlineto{\pgfqpoint{2.492587in}{1.027345in}}%
\pgfpathlineto{\pgfqpoint{2.561981in}{1.036510in}}%
\pgfpathlineto{\pgfqpoint{2.643688in}{1.050585in}}%
\pgfpathlineto{\pgfqpoint{2.702785in}{1.062681in}}%
\pgfpathlineto{\pgfqpoint{2.768856in}{1.076386in}}%
\pgfpathlineto{\pgfqpoint{2.787736in}{1.050465in}}%
\pgfpathlineto{\pgfqpoint{2.777509in}{1.052120in}}%
\pgfpathlineto{\pgfqpoint{2.764382in}{1.048922in}}%
\pgfpathlineto{\pgfqpoint{2.751446in}{1.035712in}}%
\pgfpathlineto{\pgfqpoint{2.742149in}{1.036662in}}%
\pgfpathlineto{\pgfqpoint{2.740973in}{1.028816in}}%
\pgfpathlineto{\pgfqpoint{2.757777in}{1.035021in}}%
\pgfpathlineto{\pgfqpoint{2.774686in}{1.037579in}}%
\pgfpathlineto{\pgfqpoint{2.780965in}{1.034175in}}%
\pgfpathlineto{\pgfqpoint{2.789414in}{1.038056in}}%
\pgfpathlineto{\pgfqpoint{2.793774in}{1.035336in}}%
\pgfpathlineto{\pgfqpoint{2.797633in}{1.022389in}}%
\pgfpathlineto{\pgfqpoint{2.789602in}{1.018379in}}%
\pgfpathlineto{\pgfqpoint{2.784090in}{1.002591in}}%
\pgfpathlineto{\pgfqpoint{2.778311in}{0.996443in}}%
\pgfpathlineto{\pgfqpoint{2.760595in}{0.997788in}}%
\pgfpathlineto{\pgfqpoint{2.761547in}{1.007059in}}%
\pgfpathlineto{\pgfqpoint{2.751872in}{1.003009in}}%
\pgfpathlineto{\pgfqpoint{2.749767in}{0.995277in}}%
\pgfpathlineto{\pgfqpoint{2.756399in}{0.987268in}}%
\pgfpathlineto{\pgfqpoint{2.758652in}{0.973674in}}%
\pgfpathlineto{\pgfqpoint{2.747800in}{0.965929in}}%
\pgfpathlineto{\pgfqpoint{2.759537in}{0.963205in}}%
\pgfpathlineto{\pgfqpoint{2.764848in}{0.968900in}}%
\pgfpathlineto{\pgfqpoint{2.776591in}{0.969594in}}%
\pgfpathlineto{\pgfqpoint{2.770556in}{0.957295in}}%
\pgfpathlineto{\pgfqpoint{2.763165in}{0.951369in}}%
\pgfpathlineto{\pgfqpoint{2.740807in}{0.945436in}}%
\pgfpathlineto{\pgfqpoint{2.717851in}{0.924616in}}%
\pgfpathlineto{\pgfqpoint{2.703697in}{0.904102in}}%
\pgfpathlineto{\pgfqpoint{2.703613in}{0.895771in}}%
\pgfpathlineto{\pgfqpoint{2.697959in}{0.884275in}}%
\pgfpathlineto{\pgfqpoint{2.668941in}{0.876711in}}%
\pgfpathlineto{\pgfqpoint{2.599931in}{0.926180in}}%
\pgfpathlineto{\pgfqpoint{2.539158in}{0.917251in}}%
\pgfpathlineto{\pgfqpoint{2.538649in}{0.925496in}}%
\pgfpathlineto{\pgfqpoint{2.529411in}{0.934781in}}%
\pgfpathlineto{\pgfqpoint{2.522418in}{0.937054in}}%
\pgfpathlineto{\pgfqpoint{2.456388in}{0.930193in}}%
\pgfpathlineto{\pgfqpoint{2.441211in}{0.925045in}}%
\pgfpathlineto{\pgfqpoint{2.413807in}{0.911413in}}%
\pgfpathlineto{\pgfqpoint{2.390124in}{0.907640in}}%
\pgfpathlineto{\pgfqpoint{2.348542in}{0.902516in}}%
\pgfusepath{stroke}%
\end{pgfscope}%
\begin{pgfscope}%
\pgfpathrectangle{\pgfqpoint{0.100000in}{0.100000in}}{\pgfqpoint{3.007045in}{1.925000in}}%
\pgfusepath{clip}%
\pgfsetbuttcap%
\pgfsetroundjoin%
\pgfsetlinewidth{0.050187pt}%
\definecolor{currentstroke}{rgb}{1.000000,1.000000,1.000000}%
\pgfsetstrokecolor{currentstroke}%
\pgfsetdash{}{0pt}%
\pgfpathmoveto{\pgfqpoint{2.348542in}{0.902516in}}%
\pgfpathlineto{\pgfqpoint{2.279433in}{0.894926in}}%
\pgfpathlineto{\pgfqpoint{2.216213in}{0.889241in}}%
\pgfpathlineto{\pgfqpoint{2.151235in}{0.885039in}}%
\pgfpathlineto{\pgfqpoint{2.140081in}{0.883422in}}%
\pgfpathlineto{\pgfqpoint{2.096280in}{0.880053in}}%
\pgfpathlineto{\pgfqpoint{2.026124in}{0.875959in}}%
\pgfpathlineto{\pgfqpoint{2.038606in}{0.886320in}}%
\pgfpathlineto{\pgfqpoint{2.035758in}{0.897226in}}%
\pgfpathlineto{\pgfqpoint{2.038517in}{0.910446in}}%
\pgfpathlineto{\pgfqpoint{2.042732in}{0.916670in}}%
\pgfpathlineto{\pgfqpoint{2.042574in}{0.925329in}}%
\pgfpathlineto{\pgfqpoint{2.053797in}{0.930775in}}%
\pgfpathlineto{\pgfqpoint{2.053816in}{0.943518in}}%
\pgfpathlineto{\pgfqpoint{2.055358in}{0.958208in}}%
\pgfpathlineto{\pgfqpoint{2.063184in}{0.967596in}}%
\pgfpathlineto{\pgfqpoint{2.061886in}{0.976909in}}%
\pgfpathlineto{\pgfqpoint{2.064919in}{0.977122in}}%
\pgfpathlineto{\pgfqpoint{2.068136in}{0.977587in}}%
\pgfpathlineto{\pgfqpoint{2.140708in}{0.982200in}}%
\pgfpathlineto{\pgfqpoint{2.138278in}{0.994121in}}%
\pgfpathlineto{\pgfqpoint{2.149956in}{0.992504in}}%
\pgfpathlineto{\pgfqpoint{2.220644in}{0.999497in}}%
\pgfpathlineto{\pgfqpoint{2.277384in}{1.002708in}}%
\pgfpathlineto{\pgfqpoint{2.311961in}{1.005616in}}%
\pgfpathlineto{\pgfqpoint{2.369218in}{1.011045in}}%
\pgfpathlineto{\pgfqpoint{2.370173in}{1.012125in}}%
\pgfpathlineto{\pgfqpoint{2.394637in}{1.014547in}}%
\pgfpathlineto{\pgfqpoint{2.475422in}{1.025896in}}%
\pgfpathlineto{\pgfqpoint{2.473597in}{1.012440in}}%
\pgfpathlineto{\pgfqpoint{2.465620in}{1.004543in}}%
\pgfpathlineto{\pgfqpoint{2.461381in}{0.993200in}}%
\pgfpathlineto{\pgfqpoint{2.451345in}{0.992920in}}%
\pgfpathlineto{\pgfqpoint{2.432339in}{0.978033in}}%
\pgfpathlineto{\pgfqpoint{2.423346in}{0.979220in}}%
\pgfpathlineto{\pgfqpoint{2.413472in}{0.962881in}}%
\pgfpathlineto{\pgfqpoint{2.400316in}{0.957439in}}%
\pgfpathlineto{\pgfqpoint{2.387718in}{0.945263in}}%
\pgfpathlineto{\pgfqpoint{2.373389in}{0.943000in}}%
\pgfpathlineto{\pgfqpoint{2.361925in}{0.931952in}}%
\pgfpathlineto{\pgfqpoint{2.361388in}{0.922908in}}%
\pgfpathlineto{\pgfqpoint{2.348585in}{0.917250in}}%
\pgfpathlineto{\pgfqpoint{2.348542in}{0.902516in}}%
\pgfusepath{stroke}%
\end{pgfscope}%
\begin{pgfscope}%
\pgfpathrectangle{\pgfqpoint{0.100000in}{0.100000in}}{\pgfqpoint{3.007045in}{1.925000in}}%
\pgfusepath{clip}%
\pgfsetbuttcap%
\pgfsetroundjoin%
\pgfsetlinewidth{0.050187pt}%
\definecolor{currentstroke}{rgb}{1.000000,1.000000,1.000000}%
\pgfsetstrokecolor{currentstroke}%
\pgfsetdash{}{0pt}%
\pgfpathmoveto{\pgfqpoint{1.122887in}{0.697915in}}%
\pgfpathlineto{\pgfqpoint{1.117958in}{0.705835in}}%
\pgfpathlineto{\pgfqpoint{1.119998in}{0.712658in}}%
\pgfpathlineto{\pgfqpoint{1.219064in}{0.701043in}}%
\pgfpathlineto{\pgfqpoint{1.319441in}{0.691038in}}%
\pgfpathlineto{\pgfqpoint{1.322368in}{0.724970in}}%
\pgfpathlineto{\pgfqpoint{1.337455in}{0.871739in}}%
\pgfpathlineto{\pgfqpoint{1.347370in}{0.986543in}}%
\pgfpathlineto{\pgfqpoint{1.348953in}{0.986455in}}%
\pgfpathlineto{\pgfqpoint{1.449726in}{0.978343in}}%
\pgfpathlineto{\pgfqpoint{1.507338in}{0.975120in}}%
\pgfpathlineto{\pgfqpoint{1.500607in}{0.847917in}}%
\pgfpathlineto{\pgfqpoint{1.504154in}{0.848834in}}%
\pgfpathlineto{\pgfqpoint{1.516885in}{0.834817in}}%
\pgfpathlineto{\pgfqpoint{1.522821in}{0.836596in}}%
\pgfpathlineto{\pgfqpoint{1.531836in}{0.833517in}}%
\pgfpathlineto{\pgfqpoint{1.539706in}{0.835071in}}%
\pgfpathlineto{\pgfqpoint{1.544100in}{0.822590in}}%
\pgfpathlineto{\pgfqpoint{1.553918in}{0.821935in}}%
\pgfpathlineto{\pgfqpoint{1.569594in}{0.816058in}}%
\pgfpathlineto{\pgfqpoint{1.574003in}{0.817714in}}%
\pgfpathlineto{\pgfqpoint{1.583461in}{0.810461in}}%
\pgfpathlineto{\pgfqpoint{1.589028in}{0.816058in}}%
\pgfpathlineto{\pgfqpoint{1.598675in}{0.813942in}}%
\pgfpathlineto{\pgfqpoint{1.603167in}{0.806507in}}%
\pgfpathlineto{\pgfqpoint{1.610544in}{0.804002in}}%
\pgfpathlineto{\pgfqpoint{1.615661in}{0.796290in}}%
\pgfpathlineto{\pgfqpoint{1.625879in}{0.804682in}}%
\pgfpathlineto{\pgfqpoint{1.637572in}{0.794271in}}%
\pgfpathlineto{\pgfqpoint{1.650942in}{0.798523in}}%
\pgfpathlineto{\pgfqpoint{1.657711in}{0.787562in}}%
\pgfpathlineto{\pgfqpoint{1.664329in}{0.800308in}}%
\pgfpathlineto{\pgfqpoint{1.676227in}{0.793081in}}%
\pgfpathlineto{\pgfqpoint{1.683661in}{0.795004in}}%
\pgfpathlineto{\pgfqpoint{1.694275in}{0.789894in}}%
\pgfpathlineto{\pgfqpoint{1.708036in}{0.788443in}}%
\pgfpathlineto{\pgfqpoint{1.721080in}{0.796307in}}%
\pgfpathlineto{\pgfqpoint{1.725799in}{0.792988in}}%
\pgfpathlineto{\pgfqpoint{1.739055in}{0.799728in}}%
\pgfpathlineto{\pgfqpoint{1.747503in}{0.794467in}}%
\pgfpathlineto{\pgfqpoint{1.754180in}{0.794780in}}%
\pgfpathlineto{\pgfqpoint{1.759467in}{0.800976in}}%
\pgfpathlineto{\pgfqpoint{1.767788in}{0.797893in}}%
\pgfpathlineto{\pgfqpoint{1.786521in}{0.783547in}}%
\pgfpathlineto{\pgfqpoint{1.801141in}{0.779025in}}%
\pgfpathlineto{\pgfqpoint{1.807003in}{0.773486in}}%
\pgfpathlineto{\pgfqpoint{1.814342in}{0.776503in}}%
\pgfpathlineto{\pgfqpoint{1.825462in}{0.774193in}}%
\pgfpathlineto{\pgfqpoint{1.826608in}{0.670746in}}%
\pgfpathlineto{\pgfqpoint{1.834331in}{0.664193in}}%
\pgfpathlineto{\pgfqpoint{1.840521in}{0.652119in}}%
\pgfpathlineto{\pgfqpoint{1.838345in}{0.644041in}}%
\pgfpathlineto{\pgfqpoint{1.843206in}{0.640621in}}%
\pgfpathlineto{\pgfqpoint{1.847157in}{0.625724in}}%
\pgfpathlineto{\pgfqpoint{1.854722in}{0.617746in}}%
\pgfpathlineto{\pgfqpoint{1.856418in}{0.600809in}}%
\pgfpathlineto{\pgfqpoint{1.855102in}{0.593637in}}%
\pgfpathlineto{\pgfqpoint{1.844817in}{0.574659in}}%
\pgfpathlineto{\pgfqpoint{1.843635in}{0.561013in}}%
\pgfpathlineto{\pgfqpoint{1.847123in}{0.558165in}}%
\pgfpathlineto{\pgfqpoint{1.847427in}{0.546210in}}%
\pgfpathlineto{\pgfqpoint{1.843881in}{0.538697in}}%
\pgfpathlineto{\pgfqpoint{1.838366in}{0.535867in}}%
\pgfpathlineto{\pgfqpoint{1.832966in}{0.526048in}}%
\pgfpathlineto{\pgfqpoint{1.839848in}{0.516539in}}%
\pgfpathlineto{\pgfqpoint{1.826489in}{0.516351in}}%
\pgfpathlineto{\pgfqpoint{1.790766in}{0.500008in}}%
\pgfpathlineto{\pgfqpoint{1.797587in}{0.509789in}}%
\pgfpathlineto{\pgfqpoint{1.789328in}{0.515062in}}%
\pgfpathlineto{\pgfqpoint{1.787602in}{0.524028in}}%
\pgfpathlineto{\pgfqpoint{1.771483in}{0.508411in}}%
\pgfpathlineto{\pgfqpoint{1.778637in}{0.497737in}}%
\pgfpathlineto{\pgfqpoint{1.768432in}{0.484145in}}%
\pgfpathlineto{\pgfqpoint{1.762940in}{0.484430in}}%
\pgfpathlineto{\pgfqpoint{1.757774in}{0.469628in}}%
\pgfpathlineto{\pgfqpoint{1.741435in}{0.457985in}}%
\pgfpathlineto{\pgfqpoint{1.726171in}{0.453789in}}%
\pgfpathlineto{\pgfqpoint{1.699548in}{0.442873in}}%
\pgfpathlineto{\pgfqpoint{1.696789in}{0.448962in}}%
\pgfpathlineto{\pgfqpoint{1.680205in}{0.443357in}}%
\pgfpathlineto{\pgfqpoint{1.690405in}{0.433814in}}%
\pgfpathlineto{\pgfqpoint{1.674316in}{0.425524in}}%
\pgfpathlineto{\pgfqpoint{1.656962in}{0.413006in}}%
\pgfpathlineto{\pgfqpoint{1.652794in}{0.418825in}}%
\pgfpathlineto{\pgfqpoint{1.638592in}{0.410104in}}%
\pgfpathlineto{\pgfqpoint{1.652512in}{0.406792in}}%
\pgfpathlineto{\pgfqpoint{1.642016in}{0.393080in}}%
\pgfpathlineto{\pgfqpoint{1.636893in}{0.397161in}}%
\pgfpathlineto{\pgfqpoint{1.630012in}{0.390606in}}%
\pgfpathlineto{\pgfqpoint{1.638588in}{0.384881in}}%
\pgfpathlineto{\pgfqpoint{1.633518in}{0.376512in}}%
\pgfpathlineto{\pgfqpoint{1.627286in}{0.356702in}}%
\pgfpathlineto{\pgfqpoint{1.618297in}{0.337622in}}%
\pgfpathlineto{\pgfqpoint{1.622332in}{0.325116in}}%
\pgfpathlineto{\pgfqpoint{1.629192in}{0.295630in}}%
\pgfpathlineto{\pgfqpoint{1.629809in}{0.283697in}}%
\pgfpathlineto{\pgfqpoint{1.640407in}{0.268045in}}%
\pgfpathlineto{\pgfqpoint{1.632235in}{0.268959in}}%
\pgfpathlineto{\pgfqpoint{1.624297in}{0.261063in}}%
\pgfpathlineto{\pgfqpoint{1.615000in}{0.267269in}}%
\pgfpathlineto{\pgfqpoint{1.611663in}{0.273460in}}%
\pgfpathlineto{\pgfqpoint{1.598432in}{0.276340in}}%
\pgfpathlineto{\pgfqpoint{1.578220in}{0.276693in}}%
\pgfpathlineto{\pgfqpoint{1.563329in}{0.288418in}}%
\pgfpathlineto{\pgfqpoint{1.549807in}{0.290376in}}%
\pgfpathlineto{\pgfqpoint{1.541604in}{0.299694in}}%
\pgfpathlineto{\pgfqpoint{1.524424in}{0.303444in}}%
\pgfpathlineto{\pgfqpoint{1.515030in}{0.333403in}}%
\pgfpathlineto{\pgfqpoint{1.505425in}{0.345405in}}%
\pgfpathlineto{\pgfqpoint{1.507057in}{0.356809in}}%
\pgfpathlineto{\pgfqpoint{1.501101in}{0.365143in}}%
\pgfpathlineto{\pgfqpoint{1.504839in}{0.376534in}}%
\pgfpathlineto{\pgfqpoint{1.501755in}{0.384882in}}%
\pgfpathlineto{\pgfqpoint{1.492103in}{0.388664in}}%
\pgfpathlineto{\pgfqpoint{1.483078in}{0.398289in}}%
\pgfpathlineto{\pgfqpoint{1.479790in}{0.411177in}}%
\pgfpathlineto{\pgfqpoint{1.471250in}{0.422889in}}%
\pgfpathlineto{\pgfqpoint{1.459883in}{0.432001in}}%
\pgfpathlineto{\pgfqpoint{1.449644in}{0.458142in}}%
\pgfpathlineto{\pgfqpoint{1.441377in}{0.486731in}}%
\pgfpathlineto{\pgfqpoint{1.434594in}{0.498034in}}%
\pgfpathlineto{\pgfqpoint{1.422834in}{0.507551in}}%
\pgfpathlineto{\pgfqpoint{1.419895in}{0.514459in}}%
\pgfpathlineto{\pgfqpoint{1.408884in}{0.518748in}}%
\pgfpathlineto{\pgfqpoint{1.398016in}{0.537050in}}%
\pgfpathlineto{\pgfqpoint{1.388091in}{0.535645in}}%
\pgfpathlineto{\pgfqpoint{1.378014in}{0.540216in}}%
\pgfpathlineto{\pgfqpoint{1.363728in}{0.539332in}}%
\pgfpathlineto{\pgfqpoint{1.349216in}{0.546876in}}%
\pgfpathlineto{\pgfqpoint{1.345123in}{0.539702in}}%
\pgfpathlineto{\pgfqpoint{1.328163in}{0.539522in}}%
\pgfpathlineto{\pgfqpoint{1.319529in}{0.525922in}}%
\pgfpathlineto{\pgfqpoint{1.312019in}{0.509082in}}%
\pgfpathlineto{\pgfqpoint{1.296050in}{0.491006in}}%
\pgfpathlineto{\pgfqpoint{1.277957in}{0.498939in}}%
\pgfpathlineto{\pgfqpoint{1.275390in}{0.504181in}}%
\pgfpathlineto{\pgfqpoint{1.264395in}{0.508180in}}%
\pgfpathlineto{\pgfqpoint{1.261017in}{0.513650in}}%
\pgfpathlineto{\pgfqpoint{1.246421in}{0.519171in}}%
\pgfpathlineto{\pgfqpoint{1.238256in}{0.530458in}}%
\pgfpathlineto{\pgfqpoint{1.228720in}{0.535904in}}%
\pgfpathlineto{\pgfqpoint{1.220515in}{0.545406in}}%
\pgfpathlineto{\pgfqpoint{1.214124in}{0.561500in}}%
\pgfpathlineto{\pgfqpoint{1.214839in}{0.583492in}}%
\pgfpathlineto{\pgfqpoint{1.207341in}{0.594605in}}%
\pgfpathlineto{\pgfqpoint{1.206478in}{0.606641in}}%
\pgfpathlineto{\pgfqpoint{1.189842in}{0.624662in}}%
\pgfpathlineto{\pgfqpoint{1.180162in}{0.628509in}}%
\pgfpathlineto{\pgfqpoint{1.169890in}{0.645324in}}%
\pgfpathlineto{\pgfqpoint{1.161136in}{0.652005in}}%
\pgfpathlineto{\pgfqpoint{1.150014in}{0.668274in}}%
\pgfpathlineto{\pgfqpoint{1.138634in}{0.675314in}}%
\pgfpathlineto{\pgfqpoint{1.131182in}{0.693355in}}%
\pgfpathlineto{\pgfqpoint{1.122887in}{0.697915in}}%
\pgfpathlineto{\pgfqpoint{1.122887in}{0.697915in}}%
\pgfusepath{stroke}%
\end{pgfscope}%
\begin{pgfscope}%
\pgfpathrectangle{\pgfqpoint{0.100000in}{0.100000in}}{\pgfqpoint{3.007045in}{1.925000in}}%
\pgfusepath{clip}%
\pgfsetbuttcap%
\pgfsetroundjoin%
\pgfsetlinewidth{0.050187pt}%
\definecolor{currentstroke}{rgb}{1.000000,1.000000,1.000000}%
\pgfsetstrokecolor{currentstroke}%
\pgfsetdash{}{0pt}%
\pgfpathmoveto{\pgfqpoint{1.035744in}{1.058103in}}%
\pgfpathlineto{\pgfqpoint{1.117842in}{1.046179in}}%
\pgfpathlineto{\pgfqpoint{1.149648in}{1.041279in}}%
\pgfpathlineto{\pgfqpoint{1.239303in}{1.030616in}}%
\pgfpathlineto{\pgfqpoint{1.299648in}{1.024321in}}%
\pgfpathlineto{\pgfqpoint{1.351869in}{1.019693in}}%
\pgfpathlineto{\pgfqpoint{1.348953in}{0.986455in}}%
\pgfpathlineto{\pgfqpoint{1.347370in}{0.986543in}}%
\pgfpathlineto{\pgfqpoint{1.342587in}{0.929495in}}%
\pgfpathlineto{\pgfqpoint{1.337455in}{0.871739in}}%
\pgfpathlineto{\pgfqpoint{1.331519in}{0.811277in}}%
\pgfpathlineto{\pgfqpoint{1.322368in}{0.724970in}}%
\pgfpathlineto{\pgfqpoint{1.319441in}{0.691038in}}%
\pgfpathlineto{\pgfqpoint{1.265644in}{0.696466in}}%
\pgfpathlineto{\pgfqpoint{1.219064in}{0.701043in}}%
\pgfpathlineto{\pgfqpoint{1.154643in}{0.708360in}}%
\pgfpathlineto{\pgfqpoint{1.119998in}{0.712658in}}%
\pgfpathlineto{\pgfqpoint{1.117958in}{0.705835in}}%
\pgfpathlineto{\pgfqpoint{1.122887in}{0.697915in}}%
\pgfpathlineto{\pgfqpoint{1.081250in}{0.703323in}}%
\pgfpathlineto{\pgfqpoint{1.029878in}{0.710696in}}%
\pgfpathlineto{\pgfqpoint{1.025202in}{0.681640in}}%
\pgfpathlineto{\pgfqpoint{0.978343in}{0.688693in}}%
\pgfpathlineto{\pgfqpoint{0.989353in}{0.760774in}}%
\pgfpathlineto{\pgfqpoint{1.002983in}{0.848537in}}%
\pgfpathlineto{\pgfqpoint{1.015102in}{0.924993in}}%
\pgfpathlineto{\pgfqpoint{1.025626in}{0.992952in}}%
\pgfpathlineto{\pgfqpoint{1.035744in}{1.058103in}}%
\pgfusepath{stroke}%
\end{pgfscope}%
\begin{pgfscope}%
\pgfpathrectangle{\pgfqpoint{0.100000in}{0.100000in}}{\pgfqpoint{3.007045in}{1.925000in}}%
\pgfusepath{clip}%
\pgfsetbuttcap%
\pgfsetroundjoin%
\pgfsetlinewidth{0.050187pt}%
\definecolor{currentstroke}{rgb}{1.000000,1.000000,1.000000}%
\pgfsetstrokecolor{currentstroke}%
\pgfsetdash{}{0pt}%
\pgfpathmoveto{\pgfqpoint{2.341115in}{0.635787in}}%
\pgfpathlineto{\pgfqpoint{2.262793in}{0.627088in}}%
\pgfpathlineto{\pgfqpoint{2.193752in}{0.621728in}}%
\pgfpathlineto{\pgfqpoint{2.192889in}{0.613268in}}%
\pgfpathlineto{\pgfqpoint{2.199230in}{0.605215in}}%
\pgfpathlineto{\pgfqpoint{2.206971in}{0.600484in}}%
\pgfpathlineto{\pgfqpoint{2.206853in}{0.588027in}}%
\pgfpathlineto{\pgfqpoint{2.204803in}{0.579705in}}%
\pgfpathlineto{\pgfqpoint{2.197971in}{0.573707in}}%
\pgfpathlineto{\pgfqpoint{2.188449in}{0.574342in}}%
\pgfpathlineto{\pgfqpoint{2.179446in}{0.581774in}}%
\pgfpathlineto{\pgfqpoint{2.177842in}{0.594998in}}%
\pgfpathlineto{\pgfqpoint{2.171137in}{0.602689in}}%
\pgfpathlineto{\pgfqpoint{2.166573in}{0.575162in}}%
\pgfpathlineto{\pgfqpoint{2.151065in}{0.577778in}}%
\pgfpathlineto{\pgfqpoint{2.145672in}{0.625913in}}%
\pgfpathlineto{\pgfqpoint{2.139832in}{0.676640in}}%
\pgfpathlineto{\pgfqpoint{2.142001in}{0.749828in}}%
\pgfpathlineto{\pgfqpoint{2.143347in}{0.800051in}}%
\pgfpathlineto{\pgfqpoint{2.146210in}{0.876679in}}%
\pgfpathlineto{\pgfqpoint{2.140081in}{0.883422in}}%
\pgfpathlineto{\pgfqpoint{2.151235in}{0.885039in}}%
\pgfpathlineto{\pgfqpoint{2.216213in}{0.889241in}}%
\pgfpathlineto{\pgfqpoint{2.279433in}{0.894926in}}%
\pgfpathlineto{\pgfqpoint{2.296066in}{0.836665in}}%
\pgfpathlineto{\pgfqpoint{2.307410in}{0.793901in}}%
\pgfpathlineto{\pgfqpoint{2.317578in}{0.758099in}}%
\pgfpathlineto{\pgfqpoint{2.323459in}{0.743683in}}%
\pgfpathlineto{\pgfqpoint{2.332718in}{0.730267in}}%
\pgfpathlineto{\pgfqpoint{2.331222in}{0.723440in}}%
\pgfpathlineto{\pgfqpoint{2.337837in}{0.720114in}}%
\pgfpathlineto{\pgfqpoint{2.330020in}{0.709765in}}%
\pgfpathlineto{\pgfqpoint{2.327380in}{0.691335in}}%
\pgfpathlineto{\pgfqpoint{2.335019in}{0.669779in}}%
\pgfpathlineto{\pgfqpoint{2.333958in}{0.648089in}}%
\pgfpathlineto{\pgfqpoint{2.341115in}{0.635787in}}%
\pgfusepath{stroke}%
\end{pgfscope}%
\begin{pgfscope}%
\pgfpathrectangle{\pgfqpoint{0.100000in}{0.100000in}}{\pgfqpoint{3.007045in}{1.925000in}}%
\pgfusepath{clip}%
\pgfsetbuttcap%
\pgfsetroundjoin%
\pgfsetlinewidth{0.050187pt}%
\definecolor{currentstroke}{rgb}{1.000000,1.000000,1.000000}%
\pgfsetstrokecolor{currentstroke}%
\pgfsetdash{}{0pt}%
\pgfpathmoveto{\pgfqpoint{2.151065in}{0.577778in}}%
\pgfpathlineto{\pgfqpoint{2.135151in}{0.573224in}}%
\pgfpathlineto{\pgfqpoint{2.123289in}{0.576148in}}%
\pgfpathlineto{\pgfqpoint{2.101259in}{0.569141in}}%
\pgfpathlineto{\pgfqpoint{2.084645in}{0.560117in}}%
\pgfpathlineto{\pgfqpoint{2.077828in}{0.576424in}}%
\pgfpathlineto{\pgfqpoint{2.070801in}{0.583259in}}%
\pgfpathlineto{\pgfqpoint{2.067411in}{0.591438in}}%
\pgfpathlineto{\pgfqpoint{2.072380in}{0.613576in}}%
\pgfpathlineto{\pgfqpoint{2.025296in}{0.610677in}}%
\pgfpathlineto{\pgfqpoint{1.964244in}{0.608030in}}%
\pgfpathlineto{\pgfqpoint{1.967876in}{0.613540in}}%
\pgfpathlineto{\pgfqpoint{1.963386in}{0.623948in}}%
\pgfpathlineto{\pgfqpoint{1.967860in}{0.626069in}}%
\pgfpathlineto{\pgfqpoint{1.966867in}{0.636068in}}%
\pgfpathlineto{\pgfqpoint{1.974665in}{0.645765in}}%
\pgfpathlineto{\pgfqpoint{1.978931in}{0.664590in}}%
\pgfpathlineto{\pgfqpoint{1.993202in}{0.676993in}}%
\pgfpathlineto{\pgfqpoint{1.987913in}{0.689038in}}%
\pgfpathlineto{\pgfqpoint{1.998033in}{0.695971in}}%
\pgfpathlineto{\pgfqpoint{1.989303in}{0.708523in}}%
\pgfpathlineto{\pgfqpoint{1.982985in}{0.737257in}}%
\pgfpathlineto{\pgfqpoint{1.985381in}{0.742474in}}%
\pgfpathlineto{\pgfqpoint{1.989164in}{0.752482in}}%
\pgfpathlineto{\pgfqpoint{1.985643in}{0.762985in}}%
\pgfpathlineto{\pgfqpoint{1.987365in}{0.767763in}}%
\pgfpathlineto{\pgfqpoint{1.981171in}{0.785815in}}%
\pgfpathlineto{\pgfqpoint{1.999458in}{0.818229in}}%
\pgfpathlineto{\pgfqpoint{2.000308in}{0.826857in}}%
\pgfpathlineto{\pgfqpoint{2.009117in}{0.833130in}}%
\pgfpathlineto{\pgfqpoint{2.018516in}{0.853839in}}%
\pgfpathlineto{\pgfqpoint{2.018070in}{0.862257in}}%
\pgfpathlineto{\pgfqpoint{2.029777in}{0.870862in}}%
\pgfpathlineto{\pgfqpoint{2.026124in}{0.875959in}}%
\pgfpathlineto{\pgfqpoint{2.096280in}{0.880053in}}%
\pgfpathlineto{\pgfqpoint{2.140081in}{0.883422in}}%
\pgfpathlineto{\pgfqpoint{2.146210in}{0.876679in}}%
\pgfpathlineto{\pgfqpoint{2.143347in}{0.800051in}}%
\pgfpathlineto{\pgfqpoint{2.142001in}{0.749828in}}%
\pgfpathlineto{\pgfqpoint{2.139832in}{0.676640in}}%
\pgfpathlineto{\pgfqpoint{2.145672in}{0.625913in}}%
\pgfpathlineto{\pgfqpoint{2.151065in}{0.577778in}}%
\pgfusepath{stroke}%
\end{pgfscope}%
\begin{pgfscope}%
\pgfpathrectangle{\pgfqpoint{0.100000in}{0.100000in}}{\pgfqpoint{3.007045in}{1.925000in}}%
\pgfusepath{clip}%
\pgfsetbuttcap%
\pgfsetroundjoin%
\pgfsetlinewidth{0.050187pt}%
\definecolor{currentstroke}{rgb}{1.000000,1.000000,1.000000}%
\pgfsetstrokecolor{currentstroke}%
\pgfsetdash{}{0pt}%
\pgfpathmoveto{\pgfqpoint{2.279433in}{0.894926in}}%
\pgfpathlineto{\pgfqpoint{2.348542in}{0.902516in}}%
\pgfpathlineto{\pgfqpoint{2.390124in}{0.907640in}}%
\pgfpathlineto{\pgfqpoint{2.413807in}{0.911413in}}%
\pgfpathlineto{\pgfqpoint{2.403997in}{0.895940in}}%
\pgfpathlineto{\pgfqpoint{2.404003in}{0.888645in}}%
\pgfpathlineto{\pgfqpoint{2.414101in}{0.884711in}}%
\pgfpathlineto{\pgfqpoint{2.420968in}{0.878359in}}%
\pgfpathlineto{\pgfqpoint{2.431346in}{0.877574in}}%
\pgfpathlineto{\pgfqpoint{2.441046in}{0.859694in}}%
\pgfpathlineto{\pgfqpoint{2.451568in}{0.847085in}}%
\pgfpathlineto{\pgfqpoint{2.470255in}{0.836488in}}%
\pgfpathlineto{\pgfqpoint{2.475566in}{0.828036in}}%
\pgfpathlineto{\pgfqpoint{2.492230in}{0.818895in}}%
\pgfpathlineto{\pgfqpoint{2.493483in}{0.812266in}}%
\pgfpathlineto{\pgfqpoint{2.507001in}{0.799084in}}%
\pgfpathlineto{\pgfqpoint{2.517356in}{0.795331in}}%
\pgfpathlineto{\pgfqpoint{2.524689in}{0.782884in}}%
\pgfpathlineto{\pgfqpoint{2.527935in}{0.768894in}}%
\pgfpathlineto{\pgfqpoint{2.538690in}{0.763486in}}%
\pgfpathlineto{\pgfqpoint{2.547349in}{0.748448in}}%
\pgfpathlineto{\pgfqpoint{2.550141in}{0.737389in}}%
\pgfpathlineto{\pgfqpoint{2.562829in}{0.732681in}}%
\pgfpathlineto{\pgfqpoint{2.552182in}{0.712295in}}%
\pgfpathlineto{\pgfqpoint{2.548176in}{0.700248in}}%
\pgfpathlineto{\pgfqpoint{2.550790in}{0.694600in}}%
\pgfpathlineto{\pgfqpoint{2.546225in}{0.685215in}}%
\pgfpathlineto{\pgfqpoint{2.544285in}{0.672239in}}%
\pgfpathlineto{\pgfqpoint{2.536167in}{0.669822in}}%
\pgfpathlineto{\pgfqpoint{2.540454in}{0.657942in}}%
\pgfpathlineto{\pgfqpoint{2.540180in}{0.642869in}}%
\pgfpathlineto{\pgfqpoint{2.525751in}{0.643434in}}%
\pgfpathlineto{\pgfqpoint{2.515700in}{0.646184in}}%
\pgfpathlineto{\pgfqpoint{2.511350in}{0.641072in}}%
\pgfpathlineto{\pgfqpoint{2.514596in}{0.628996in}}%
\pgfpathlineto{\pgfqpoint{2.513890in}{0.614955in}}%
\pgfpathlineto{\pgfqpoint{2.507549in}{0.613881in}}%
\pgfpathlineto{\pgfqpoint{2.502419in}{0.626994in}}%
\pgfpathlineto{\pgfqpoint{2.450185in}{0.623513in}}%
\pgfpathlineto{\pgfqpoint{2.384322in}{0.619919in}}%
\pgfpathlineto{\pgfqpoint{2.351105in}{0.617581in}}%
\pgfpathlineto{\pgfqpoint{2.341115in}{0.635787in}}%
\pgfpathlineto{\pgfqpoint{2.333958in}{0.648089in}}%
\pgfpathlineto{\pgfqpoint{2.335019in}{0.669779in}}%
\pgfpathlineto{\pgfqpoint{2.327380in}{0.691335in}}%
\pgfpathlineto{\pgfqpoint{2.330020in}{0.709765in}}%
\pgfpathlineto{\pgfqpoint{2.337837in}{0.720114in}}%
\pgfpathlineto{\pgfqpoint{2.331222in}{0.723440in}}%
\pgfpathlineto{\pgfqpoint{2.332718in}{0.730267in}}%
\pgfpathlineto{\pgfqpoint{2.323459in}{0.743683in}}%
\pgfpathlineto{\pgfqpoint{2.317578in}{0.758099in}}%
\pgfpathlineto{\pgfqpoint{2.307410in}{0.793901in}}%
\pgfpathlineto{\pgfqpoint{2.296066in}{0.836665in}}%
\pgfpathlineto{\pgfqpoint{2.279433in}{0.894926in}}%
\pgfusepath{stroke}%
\end{pgfscope}%
\begin{pgfscope}%
\pgfpathrectangle{\pgfqpoint{0.100000in}{0.100000in}}{\pgfqpoint{3.007045in}{1.925000in}}%
\pgfusepath{clip}%
\pgfsetbuttcap%
\pgfsetroundjoin%
\pgfsetlinewidth{0.050187pt}%
\definecolor{currentstroke}{rgb}{1.000000,1.000000,1.000000}%
\pgfsetstrokecolor{currentstroke}%
\pgfsetdash{}{0pt}%
\pgfpathmoveto{\pgfqpoint{2.413807in}{0.911413in}}%
\pgfpathlineto{\pgfqpoint{2.441211in}{0.925045in}}%
\pgfpathlineto{\pgfqpoint{2.456388in}{0.930193in}}%
\pgfpathlineto{\pgfqpoint{2.522418in}{0.937054in}}%
\pgfpathlineto{\pgfqpoint{2.529411in}{0.934781in}}%
\pgfpathlineto{\pgfqpoint{2.538649in}{0.925496in}}%
\pgfpathlineto{\pgfqpoint{2.539158in}{0.917251in}}%
\pgfpathlineto{\pgfqpoint{2.599931in}{0.926180in}}%
\pgfpathlineto{\pgfqpoint{2.668941in}{0.876711in}}%
\pgfpathlineto{\pgfqpoint{2.655999in}{0.863222in}}%
\pgfpathlineto{\pgfqpoint{2.638241in}{0.831896in}}%
\pgfpathlineto{\pgfqpoint{2.643296in}{0.825164in}}%
\pgfpathlineto{\pgfqpoint{2.633823in}{0.812096in}}%
\pgfpathlineto{\pgfqpoint{2.624400in}{0.810617in}}%
\pgfpathlineto{\pgfqpoint{2.624365in}{0.802765in}}%
\pgfpathlineto{\pgfqpoint{2.617549in}{0.794548in}}%
\pgfpathlineto{\pgfqpoint{2.609066in}{0.792881in}}%
\pgfpathlineto{\pgfqpoint{2.610917in}{0.785581in}}%
\pgfpathlineto{\pgfqpoint{2.606183in}{0.779976in}}%
\pgfpathlineto{\pgfqpoint{2.594860in}{0.775134in}}%
\pgfpathlineto{\pgfqpoint{2.581058in}{0.764120in}}%
\pgfpathlineto{\pgfqpoint{2.583663in}{0.756995in}}%
\pgfpathlineto{\pgfqpoint{2.574977in}{0.752521in}}%
\pgfpathlineto{\pgfqpoint{2.567651in}{0.757225in}}%
\pgfpathlineto{\pgfqpoint{2.562294in}{0.736773in}}%
\pgfpathlineto{\pgfqpoint{2.550141in}{0.737389in}}%
\pgfpathlineto{\pgfqpoint{2.547349in}{0.748448in}}%
\pgfpathlineto{\pgfqpoint{2.538690in}{0.763486in}}%
\pgfpathlineto{\pgfqpoint{2.527935in}{0.768894in}}%
\pgfpathlineto{\pgfqpoint{2.524689in}{0.782884in}}%
\pgfpathlineto{\pgfqpoint{2.517356in}{0.795331in}}%
\pgfpathlineto{\pgfqpoint{2.507001in}{0.799084in}}%
\pgfpathlineto{\pgfqpoint{2.493483in}{0.812266in}}%
\pgfpathlineto{\pgfqpoint{2.492230in}{0.818895in}}%
\pgfpathlineto{\pgfqpoint{2.475566in}{0.828036in}}%
\pgfpathlineto{\pgfqpoint{2.470255in}{0.836488in}}%
\pgfpathlineto{\pgfqpoint{2.451568in}{0.847085in}}%
\pgfpathlineto{\pgfqpoint{2.441046in}{0.859694in}}%
\pgfpathlineto{\pgfqpoint{2.431346in}{0.877574in}}%
\pgfpathlineto{\pgfqpoint{2.420968in}{0.878359in}}%
\pgfpathlineto{\pgfqpoint{2.414101in}{0.884711in}}%
\pgfpathlineto{\pgfqpoint{2.404003in}{0.888645in}}%
\pgfpathlineto{\pgfqpoint{2.403997in}{0.895940in}}%
\pgfpathlineto{\pgfqpoint{2.413807in}{0.911413in}}%
\pgfusepath{stroke}%
\end{pgfscope}%
\begin{pgfscope}%
\pgfpathrectangle{\pgfqpoint{0.100000in}{0.100000in}}{\pgfqpoint{3.007045in}{1.925000in}}%
\pgfusepath{clip}%
\pgfsetbuttcap%
\pgfsetroundjoin%
\pgfsetlinewidth{0.050187pt}%
\definecolor{currentstroke}{rgb}{1.000000,1.000000,1.000000}%
\pgfsetstrokecolor{currentstroke}%
\pgfsetdash{}{0pt}%
\pgfpathmoveto{\pgfqpoint{1.792603in}{0.967662in}}%
\pgfpathlineto{\pgfqpoint{1.832874in}{0.968025in}}%
\pgfpathlineto{\pgfqpoint{1.886150in}{0.968996in}}%
\pgfpathlineto{\pgfqpoint{1.977188in}{0.971886in}}%
\pgfpathlineto{\pgfqpoint{2.029259in}{0.974639in}}%
\pgfpathlineto{\pgfqpoint{2.034844in}{0.967726in}}%
\pgfpathlineto{\pgfqpoint{2.034495in}{0.960424in}}%
\pgfpathlineto{\pgfqpoint{2.021881in}{0.947828in}}%
\pgfpathlineto{\pgfqpoint{2.018831in}{0.940931in}}%
\pgfpathlineto{\pgfqpoint{2.053816in}{0.943518in}}%
\pgfpathlineto{\pgfqpoint{2.053797in}{0.930775in}}%
\pgfpathlineto{\pgfqpoint{2.042574in}{0.925329in}}%
\pgfpathlineto{\pgfqpoint{2.042732in}{0.916670in}}%
\pgfpathlineto{\pgfqpoint{2.038517in}{0.910446in}}%
\pgfpathlineto{\pgfqpoint{2.035758in}{0.897226in}}%
\pgfpathlineto{\pgfqpoint{2.038606in}{0.886320in}}%
\pgfpathlineto{\pgfqpoint{2.026124in}{0.875959in}}%
\pgfpathlineto{\pgfqpoint{2.029777in}{0.870862in}}%
\pgfpathlineto{\pgfqpoint{2.018070in}{0.862257in}}%
\pgfpathlineto{\pgfqpoint{2.018516in}{0.853839in}}%
\pgfpathlineto{\pgfqpoint{2.009117in}{0.833130in}}%
\pgfpathlineto{\pgfqpoint{2.000308in}{0.826857in}}%
\pgfpathlineto{\pgfqpoint{1.999458in}{0.818229in}}%
\pgfpathlineto{\pgfqpoint{1.981171in}{0.785815in}}%
\pgfpathlineto{\pgfqpoint{1.987365in}{0.767763in}}%
\pgfpathlineto{\pgfqpoint{1.985643in}{0.762985in}}%
\pgfpathlineto{\pgfqpoint{1.989164in}{0.752482in}}%
\pgfpathlineto{\pgfqpoint{1.985381in}{0.742474in}}%
\pgfpathlineto{\pgfqpoint{1.935375in}{0.740410in}}%
\pgfpathlineto{\pgfqpoint{1.870456in}{0.739340in}}%
\pgfpathlineto{\pgfqpoint{1.825680in}{0.738936in}}%
\pgfpathlineto{\pgfqpoint{1.825462in}{0.774193in}}%
\pgfpathlineto{\pgfqpoint{1.814342in}{0.776503in}}%
\pgfpathlineto{\pgfqpoint{1.807003in}{0.773486in}}%
\pgfpathlineto{\pgfqpoint{1.801141in}{0.779025in}}%
\pgfpathlineto{\pgfqpoint{1.801746in}{0.816372in}}%
\pgfpathlineto{\pgfqpoint{1.803058in}{0.895889in}}%
\pgfpathlineto{\pgfqpoint{1.796697in}{0.942452in}}%
\pgfpathlineto{\pgfqpoint{1.792603in}{0.967662in}}%
\pgfusepath{stroke}%
\end{pgfscope}%
\begin{pgfscope}%
\pgfpathrectangle{\pgfqpoint{0.100000in}{0.100000in}}{\pgfqpoint{3.007045in}{1.925000in}}%
\pgfusepath{clip}%
\pgfsetbuttcap%
\pgfsetroundjoin%
\pgfsetlinewidth{0.050187pt}%
\definecolor{currentstroke}{rgb}{1.000000,1.000000,1.000000}%
\pgfsetstrokecolor{currentstroke}%
\pgfsetdash{}{0pt}%
\pgfpathmoveto{\pgfqpoint{1.825680in}{0.738936in}}%
\pgfpathlineto{\pgfqpoint{1.870456in}{0.739340in}}%
\pgfpathlineto{\pgfqpoint{1.935375in}{0.740410in}}%
\pgfpathlineto{\pgfqpoint{1.985381in}{0.742474in}}%
\pgfpathlineto{\pgfqpoint{1.982985in}{0.737257in}}%
\pgfpathlineto{\pgfqpoint{1.989303in}{0.708523in}}%
\pgfpathlineto{\pgfqpoint{1.998033in}{0.695971in}}%
\pgfpathlineto{\pgfqpoint{1.987913in}{0.689038in}}%
\pgfpathlineto{\pgfqpoint{1.993202in}{0.676993in}}%
\pgfpathlineto{\pgfqpoint{1.978931in}{0.664590in}}%
\pgfpathlineto{\pgfqpoint{1.974665in}{0.645765in}}%
\pgfpathlineto{\pgfqpoint{1.966867in}{0.636068in}}%
\pgfpathlineto{\pgfqpoint{1.967860in}{0.626069in}}%
\pgfpathlineto{\pgfqpoint{1.963386in}{0.623948in}}%
\pgfpathlineto{\pgfqpoint{1.967876in}{0.613540in}}%
\pgfpathlineto{\pgfqpoint{1.964244in}{0.608030in}}%
\pgfpathlineto{\pgfqpoint{2.025296in}{0.610677in}}%
\pgfpathlineto{\pgfqpoint{2.072380in}{0.613576in}}%
\pgfpathlineto{\pgfqpoint{2.067411in}{0.591438in}}%
\pgfpathlineto{\pgfqpoint{2.070801in}{0.583259in}}%
\pgfpathlineto{\pgfqpoint{2.077828in}{0.576424in}}%
\pgfpathlineto{\pgfqpoint{2.084645in}{0.560117in}}%
\pgfpathlineto{\pgfqpoint{2.063113in}{0.563874in}}%
\pgfpathlineto{\pgfqpoint{2.055170in}{0.570050in}}%
\pgfpathlineto{\pgfqpoint{2.045706in}{0.570344in}}%
\pgfpathlineto{\pgfqpoint{2.035757in}{0.556820in}}%
\pgfpathlineto{\pgfqpoint{2.037740in}{0.550661in}}%
\pgfpathlineto{\pgfqpoint{2.069524in}{0.546935in}}%
\pgfpathlineto{\pgfqpoint{2.077845in}{0.539855in}}%
\pgfpathlineto{\pgfqpoint{2.091386in}{0.536213in}}%
\pgfpathlineto{\pgfqpoint{2.085472in}{0.527832in}}%
\pgfpathlineto{\pgfqpoint{2.075509in}{0.520522in}}%
\pgfpathlineto{\pgfqpoint{2.101808in}{0.504085in}}%
\pgfpathlineto{\pgfqpoint{2.114058in}{0.501516in}}%
\pgfpathlineto{\pgfqpoint{2.119152in}{0.488080in}}%
\pgfpathlineto{\pgfqpoint{2.107618in}{0.485609in}}%
\pgfpathlineto{\pgfqpoint{2.085940in}{0.502507in}}%
\pgfpathlineto{\pgfqpoint{2.077465in}{0.504856in}}%
\pgfpathlineto{\pgfqpoint{2.073363in}{0.511532in}}%
\pgfpathlineto{\pgfqpoint{2.060791in}{0.508792in}}%
\pgfpathlineto{\pgfqpoint{2.056855in}{0.500190in}}%
\pgfpathlineto{\pgfqpoint{2.059348in}{0.490604in}}%
\pgfpathlineto{\pgfqpoint{2.050884in}{0.484948in}}%
\pgfpathlineto{\pgfqpoint{2.043160in}{0.498883in}}%
\pgfpathlineto{\pgfqpoint{2.029631in}{0.494728in}}%
\pgfpathlineto{\pgfqpoint{2.024562in}{0.486414in}}%
\pgfpathlineto{\pgfqpoint{2.014953in}{0.488779in}}%
\pgfpathlineto{\pgfqpoint{2.008810in}{0.499279in}}%
\pgfpathlineto{\pgfqpoint{1.992513in}{0.502746in}}%
\pgfpathlineto{\pgfqpoint{1.989436in}{0.508210in}}%
\pgfpathlineto{\pgfqpoint{1.972456in}{0.517727in}}%
\pgfpathlineto{\pgfqpoint{1.968241in}{0.526049in}}%
\pgfpathlineto{\pgfqpoint{1.954018in}{0.522648in}}%
\pgfpathlineto{\pgfqpoint{1.955825in}{0.530267in}}%
\pgfpathlineto{\pgfqpoint{1.938148in}{0.522451in}}%
\pgfpathlineto{\pgfqpoint{1.942896in}{0.514343in}}%
\pgfpathlineto{\pgfqpoint{1.929245in}{0.509548in}}%
\pgfpathlineto{\pgfqpoint{1.911163in}{0.512183in}}%
\pgfpathlineto{\pgfqpoint{1.874567in}{0.524719in}}%
\pgfpathlineto{\pgfqpoint{1.846328in}{0.522231in}}%
\pgfpathlineto{\pgfqpoint{1.843881in}{0.538697in}}%
\pgfpathlineto{\pgfqpoint{1.847427in}{0.546210in}}%
\pgfpathlineto{\pgfqpoint{1.847123in}{0.558165in}}%
\pgfpathlineto{\pgfqpoint{1.843635in}{0.561013in}}%
\pgfpathlineto{\pgfqpoint{1.844817in}{0.574659in}}%
\pgfpathlineto{\pgfqpoint{1.855102in}{0.593637in}}%
\pgfpathlineto{\pgfqpoint{1.856418in}{0.600809in}}%
\pgfpathlineto{\pgfqpoint{1.854722in}{0.617746in}}%
\pgfpathlineto{\pgfqpoint{1.847157in}{0.625724in}}%
\pgfpathlineto{\pgfqpoint{1.843206in}{0.640621in}}%
\pgfpathlineto{\pgfqpoint{1.838345in}{0.644041in}}%
\pgfpathlineto{\pgfqpoint{1.840521in}{0.652119in}}%
\pgfpathlineto{\pgfqpoint{1.834331in}{0.664193in}}%
\pgfpathlineto{\pgfqpoint{1.826608in}{0.670746in}}%
\pgfpathlineto{\pgfqpoint{1.825680in}{0.738936in}}%
\pgfusepath{stroke}%
\end{pgfscope}%
\begin{pgfscope}%
\pgfpathrectangle{\pgfqpoint{0.100000in}{0.100000in}}{\pgfqpoint{3.007045in}{1.925000in}}%
\pgfusepath{clip}%
\pgfsetbuttcap%
\pgfsetroundjoin%
\pgfsetlinewidth{0.050187pt}%
\definecolor{currentstroke}{rgb}{1.000000,1.000000,1.000000}%
\pgfsetstrokecolor{currentstroke}%
\pgfsetdash{}{0pt}%
\pgfpathmoveto{\pgfqpoint{1.945436in}{0.513722in}}%
\pgfpathlineto{\pgfqpoint{1.951913in}{0.517575in}}%
\pgfpathlineto{\pgfqpoint{1.959776in}{0.513019in}}%
\pgfpathlineto{\pgfqpoint{1.955391in}{0.506751in}}%
\pgfpathlineto{\pgfqpoint{1.945436in}{0.513722in}}%
\pgfusepath{stroke}%
\end{pgfscope}%
\begin{pgfscope}%
\pgfpathrectangle{\pgfqpoint{0.100000in}{0.100000in}}{\pgfqpoint{3.007045in}{1.925000in}}%
\pgfusepath{clip}%
\pgfsetbuttcap%
\pgfsetroundjoin%
\pgfsetlinewidth{0.050187pt}%
\definecolor{currentstroke}{rgb}{1.000000,1.000000,1.000000}%
\pgfsetstrokecolor{currentstroke}%
\pgfsetdash{}{0pt}%
\pgfpathmoveto{\pgfqpoint{2.206853in}{0.588027in}}%
\pgfpathlineto{\pgfqpoint{2.206971in}{0.600484in}}%
\pgfpathlineto{\pgfqpoint{2.199230in}{0.605215in}}%
\pgfpathlineto{\pgfqpoint{2.192889in}{0.613268in}}%
\pgfpathlineto{\pgfqpoint{2.193752in}{0.621728in}}%
\pgfpathlineto{\pgfqpoint{2.262793in}{0.627088in}}%
\pgfpathlineto{\pgfqpoint{2.341115in}{0.635787in}}%
\pgfpathlineto{\pgfqpoint{2.351105in}{0.617581in}}%
\pgfpathlineto{\pgfqpoint{2.384322in}{0.619919in}}%
\pgfpathlineto{\pgfqpoint{2.450185in}{0.623513in}}%
\pgfpathlineto{\pgfqpoint{2.502419in}{0.626994in}}%
\pgfpathlineto{\pgfqpoint{2.507549in}{0.613881in}}%
\pgfpathlineto{\pgfqpoint{2.513890in}{0.614955in}}%
\pgfpathlineto{\pgfqpoint{2.514596in}{0.628996in}}%
\pgfpathlineto{\pgfqpoint{2.511350in}{0.641072in}}%
\pgfpathlineto{\pgfqpoint{2.515700in}{0.646184in}}%
\pgfpathlineto{\pgfqpoint{2.525751in}{0.643434in}}%
\pgfpathlineto{\pgfqpoint{2.540180in}{0.642869in}}%
\pgfpathlineto{\pgfqpoint{2.542399in}{0.632080in}}%
\pgfpathlineto{\pgfqpoint{2.546837in}{0.625915in}}%
\pgfpathlineto{\pgfqpoint{2.550303in}{0.612426in}}%
\pgfpathlineto{\pgfqpoint{2.561042in}{0.591552in}}%
\pgfpathlineto{\pgfqpoint{2.561096in}{0.585906in}}%
\pgfpathlineto{\pgfqpoint{2.576968in}{0.561396in}}%
\pgfpathlineto{\pgfqpoint{2.578503in}{0.556332in}}%
\pgfpathlineto{\pgfqpoint{2.602314in}{0.521715in}}%
\pgfpathlineto{\pgfqpoint{2.598541in}{0.521156in}}%
\pgfpathlineto{\pgfqpoint{2.608579in}{0.496534in}}%
\pgfpathlineto{\pgfqpoint{2.636155in}{0.453729in}}%
\pgfpathlineto{\pgfqpoint{2.652165in}{0.422256in}}%
\pgfpathlineto{\pgfqpoint{2.657516in}{0.416746in}}%
\pgfpathlineto{\pgfqpoint{2.666685in}{0.396991in}}%
\pgfpathlineto{\pgfqpoint{2.669866in}{0.365359in}}%
\pgfpathlineto{\pgfqpoint{2.671132in}{0.341707in}}%
\pgfpathlineto{\pgfqpoint{2.669603in}{0.326563in}}%
\pgfpathlineto{\pgfqpoint{2.664687in}{0.315737in}}%
\pgfpathlineto{\pgfqpoint{2.666971in}{0.301540in}}%
\pgfpathlineto{\pgfqpoint{2.661678in}{0.290303in}}%
\pgfpathlineto{\pgfqpoint{2.645962in}{0.281088in}}%
\pgfpathlineto{\pgfqpoint{2.635740in}{0.281761in}}%
\pgfpathlineto{\pgfqpoint{2.629103in}{0.276944in}}%
\pgfpathlineto{\pgfqpoint{2.627201in}{0.289805in}}%
\pgfpathlineto{\pgfqpoint{2.616214in}{0.293200in}}%
\pgfpathlineto{\pgfqpoint{2.606364in}{0.311145in}}%
\pgfpathlineto{\pgfqpoint{2.605252in}{0.319317in}}%
\pgfpathlineto{\pgfqpoint{2.587624in}{0.324280in}}%
\pgfpathlineto{\pgfqpoint{2.576265in}{0.323205in}}%
\pgfpathlineto{\pgfqpoint{2.569838in}{0.335063in}}%
\pgfpathlineto{\pgfqpoint{2.562480in}{0.356375in}}%
\pgfpathlineto{\pgfqpoint{2.552268in}{0.360695in}}%
\pgfpathlineto{\pgfqpoint{2.546732in}{0.372908in}}%
\pgfpathlineto{\pgfqpoint{2.533193in}{0.380106in}}%
\pgfpathlineto{\pgfqpoint{2.525367in}{0.389110in}}%
\pgfpathlineto{\pgfqpoint{2.511867in}{0.413587in}}%
\pgfpathlineto{\pgfqpoint{2.510279in}{0.430815in}}%
\pgfpathlineto{\pgfqpoint{2.517636in}{0.441838in}}%
\pgfpathlineto{\pgfqpoint{2.516901in}{0.449497in}}%
\pgfpathlineto{\pgfqpoint{2.508399in}{0.450316in}}%
\pgfpathlineto{\pgfqpoint{2.501297in}{0.455644in}}%
\pgfpathlineto{\pgfqpoint{2.497407in}{0.449138in}}%
\pgfpathlineto{\pgfqpoint{2.504224in}{0.443837in}}%
\pgfpathlineto{\pgfqpoint{2.499262in}{0.434305in}}%
\pgfpathlineto{\pgfqpoint{2.491236in}{0.442217in}}%
\pgfpathlineto{\pgfqpoint{2.492167in}{0.464156in}}%
\pgfpathlineto{\pgfqpoint{2.496039in}{0.481956in}}%
\pgfpathlineto{\pgfqpoint{2.494048in}{0.512497in}}%
\pgfpathlineto{\pgfqpoint{2.486048in}{0.519772in}}%
\pgfpathlineto{\pgfqpoint{2.482024in}{0.529107in}}%
\pgfpathlineto{\pgfqpoint{2.468242in}{0.528899in}}%
\pgfpathlineto{\pgfqpoint{2.454618in}{0.544259in}}%
\pgfpathlineto{\pgfqpoint{2.445483in}{0.548863in}}%
\pgfpathlineto{\pgfqpoint{2.442774in}{0.558590in}}%
\pgfpathlineto{\pgfqpoint{2.433840in}{0.562006in}}%
\pgfpathlineto{\pgfqpoint{2.426435in}{0.572746in}}%
\pgfpathlineto{\pgfqpoint{2.406875in}{0.581513in}}%
\pgfpathlineto{\pgfqpoint{2.391668in}{0.581737in}}%
\pgfpathlineto{\pgfqpoint{2.385048in}{0.578370in}}%
\pgfpathlineto{\pgfqpoint{2.386693in}{0.567859in}}%
\pgfpathlineto{\pgfqpoint{2.379790in}{0.568357in}}%
\pgfpathlineto{\pgfqpoint{2.358555in}{0.553661in}}%
\pgfpathlineto{\pgfqpoint{2.333048in}{0.547836in}}%
\pgfpathlineto{\pgfqpoint{2.332629in}{0.555043in}}%
\pgfpathlineto{\pgfqpoint{2.326976in}{0.562087in}}%
\pgfpathlineto{\pgfqpoint{2.311810in}{0.571811in}}%
\pgfpathlineto{\pgfqpoint{2.290041in}{0.581781in}}%
\pgfpathlineto{\pgfqpoint{2.266395in}{0.587059in}}%
\pgfpathlineto{\pgfqpoint{2.261956in}{0.594236in}}%
\pgfpathlineto{\pgfqpoint{2.253431in}{0.588250in}}%
\pgfpathlineto{\pgfqpoint{2.243165in}{0.586930in}}%
\pgfpathlineto{\pgfqpoint{2.207429in}{0.577514in}}%
\pgfpathlineto{\pgfqpoint{2.206853in}{0.588027in}}%
\pgfusepath{stroke}%
\end{pgfscope}%
\begin{pgfscope}%
\pgfpathrectangle{\pgfqpoint{0.100000in}{0.100000in}}{\pgfqpoint{3.007045in}{1.925000in}}%
\pgfusepath{clip}%
\pgfsetbuttcap%
\pgfsetroundjoin%
\pgfsetlinewidth{0.050187pt}%
\definecolor{currentstroke}{rgb}{1.000000,1.000000,1.000000}%
\pgfsetstrokecolor{currentstroke}%
\pgfsetdash{}{0pt}%
\pgfpathmoveto{\pgfqpoint{2.605783in}{0.522207in}}%
\pgfpathlineto{\pgfqpoint{2.615029in}{0.511452in}}%
\pgfpathlineto{\pgfqpoint{2.604349in}{0.510775in}}%
\pgfpathlineto{\pgfqpoint{2.605783in}{0.522207in}}%
\pgfusepath{stroke}%
\end{pgfscope}%
\begin{pgfscope}%
\pgfpathrectangle{\pgfqpoint{0.100000in}{0.100000in}}{\pgfqpoint{3.007045in}{1.925000in}}%
\pgfusepath{clip}%
\pgfsetbuttcap%
\pgfsetroundjoin%
\pgfsetlinewidth{0.050187pt}%
\definecolor{currentstroke}{rgb}{1.000000,1.000000,1.000000}%
\pgfsetstrokecolor{currentstroke}%
\pgfsetdash{}{0pt}%
\pgfpathmoveto{\pgfqpoint{2.061492in}{1.750249in}}%
\pgfpathlineto{\pgfqpoint{2.056476in}{1.740485in}}%
\pgfpathlineto{\pgfqpoint{2.044508in}{1.734787in}}%
\pgfpathlineto{\pgfqpoint{2.039325in}{1.727117in}}%
\pgfpathlineto{\pgfqpoint{2.033253in}{1.732641in}}%
\pgfpathlineto{\pgfqpoint{2.061492in}{1.750249in}}%
\pgfusepath{stroke}%
\end{pgfscope}%
\begin{pgfscope}%
\pgfpathrectangle{\pgfqpoint{0.100000in}{0.100000in}}{\pgfqpoint{3.007045in}{1.925000in}}%
\pgfusepath{clip}%
\pgfsetbuttcap%
\pgfsetroundjoin%
\pgfsetlinewidth{0.050187pt}%
\definecolor{currentstroke}{rgb}{1.000000,1.000000,1.000000}%
\pgfsetstrokecolor{currentstroke}%
\pgfsetdash{}{0pt}%
\pgfpathmoveto{\pgfqpoint{2.065561in}{1.691489in}}%
\pgfpathlineto{\pgfqpoint{2.077807in}{1.702912in}}%
\pgfpathlineto{\pgfqpoint{2.096694in}{1.705904in}}%
\pgfpathlineto{\pgfqpoint{2.091492in}{1.697959in}}%
\pgfpathlineto{\pgfqpoint{2.078543in}{1.686470in}}%
\pgfpathlineto{\pgfqpoint{2.070973in}{1.671723in}}%
\pgfpathlineto{\pgfqpoint{2.067767in}{1.679722in}}%
\pgfpathlineto{\pgfqpoint{2.062075in}{1.680870in}}%
\pgfpathlineto{\pgfqpoint{2.061503in}{1.688091in}}%
\pgfpathlineto{\pgfqpoint{2.065561in}{1.691489in}}%
\pgfusepath{stroke}%
\end{pgfscope}%
\begin{pgfscope}%
\pgfpathrectangle{\pgfqpoint{0.100000in}{0.100000in}}{\pgfqpoint{3.007045in}{1.925000in}}%
\pgfusepath{clip}%
\pgfsetbuttcap%
\pgfsetroundjoin%
\pgfsetlinewidth{0.050187pt}%
\definecolor{currentstroke}{rgb}{1.000000,1.000000,1.000000}%
\pgfsetstrokecolor{currentstroke}%
\pgfsetdash{}{0pt}%
\pgfpathmoveto{\pgfqpoint{2.114334in}{1.552060in}}%
\pgfpathlineto{\pgfqpoint{2.111072in}{1.555679in}}%
\pgfpathlineto{\pgfqpoint{2.114501in}{1.565867in}}%
\pgfpathlineto{\pgfqpoint{2.104324in}{1.566517in}}%
\pgfpathlineto{\pgfqpoint{2.107021in}{1.575301in}}%
\pgfpathlineto{\pgfqpoint{2.105344in}{1.589291in}}%
\pgfpathlineto{\pgfqpoint{2.096212in}{1.594142in}}%
\pgfpathlineto{\pgfqpoint{2.086656in}{1.603958in}}%
\pgfpathlineto{\pgfqpoint{2.071931in}{1.606832in}}%
\pgfpathlineto{\pgfqpoint{2.057600in}{1.606750in}}%
\pgfpathlineto{\pgfqpoint{2.043506in}{1.613716in}}%
\pgfpathlineto{\pgfqpoint{1.996352in}{1.623871in}}%
\pgfpathlineto{\pgfqpoint{1.991199in}{1.634631in}}%
\pgfpathlineto{\pgfqpoint{1.982007in}{1.638302in}}%
\pgfpathlineto{\pgfqpoint{1.999369in}{1.646540in}}%
\pgfpathlineto{\pgfqpoint{2.009177in}{1.656818in}}%
\pgfpathlineto{\pgfqpoint{2.027446in}{1.659605in}}%
\pgfpathlineto{\pgfqpoint{2.038687in}{1.670074in}}%
\pgfpathlineto{\pgfqpoint{2.044585in}{1.670493in}}%
\pgfpathlineto{\pgfqpoint{2.049086in}{1.677958in}}%
\pgfpathlineto{\pgfqpoint{2.060059in}{1.686803in}}%
\pgfpathlineto{\pgfqpoint{2.061015in}{1.680552in}}%
\pgfpathlineto{\pgfqpoint{2.065959in}{1.679261in}}%
\pgfpathlineto{\pgfqpoint{2.069679in}{1.671805in}}%
\pgfpathlineto{\pgfqpoint{2.070349in}{1.659095in}}%
\pgfpathlineto{\pgfqpoint{2.081529in}{1.666691in}}%
\pgfpathlineto{\pgfqpoint{2.094557in}{1.668277in}}%
\pgfpathlineto{\pgfqpoint{2.105677in}{1.664298in}}%
\pgfpathlineto{\pgfqpoint{2.120759in}{1.643609in}}%
\pgfpathlineto{\pgfqpoint{2.137230in}{1.646915in}}%
\pgfpathlineto{\pgfqpoint{2.143923in}{1.641378in}}%
\pgfpathlineto{\pgfqpoint{2.154687in}{1.640884in}}%
\pgfpathlineto{\pgfqpoint{2.161839in}{1.650814in}}%
\pgfpathlineto{\pgfqpoint{2.175410in}{1.659592in}}%
\pgfpathlineto{\pgfqpoint{2.188452in}{1.662342in}}%
\pgfpathlineto{\pgfqpoint{2.204637in}{1.662628in}}%
\pgfpathlineto{\pgfqpoint{2.216463in}{1.669409in}}%
\pgfpathlineto{\pgfqpoint{2.226115in}{1.666285in}}%
\pgfpathlineto{\pgfqpoint{2.228171in}{1.651944in}}%
\pgfpathlineto{\pgfqpoint{2.248877in}{1.649698in}}%
\pgfpathlineto{\pgfqpoint{2.260125in}{1.656411in}}%
\pgfpathlineto{\pgfqpoint{2.267921in}{1.641288in}}%
\pgfpathlineto{\pgfqpoint{2.282798in}{1.623800in}}%
\pgfpathlineto{\pgfqpoint{2.261976in}{1.623901in}}%
\pgfpathlineto{\pgfqpoint{2.255397in}{1.621740in}}%
\pgfpathlineto{\pgfqpoint{2.246379in}{1.624549in}}%
\pgfpathlineto{\pgfqpoint{2.245772in}{1.612443in}}%
\pgfpathlineto{\pgfqpoint{2.229399in}{1.621904in}}%
\pgfpathlineto{\pgfqpoint{2.208341in}{1.624795in}}%
\pgfpathlineto{\pgfqpoint{2.202544in}{1.615581in}}%
\pgfpathlineto{\pgfqpoint{2.174992in}{1.611101in}}%
\pgfpathlineto{\pgfqpoint{2.171822in}{1.603295in}}%
\pgfpathlineto{\pgfqpoint{2.152650in}{1.600953in}}%
\pgfpathlineto{\pgfqpoint{2.146854in}{1.593019in}}%
\pgfpathlineto{\pgfqpoint{2.136713in}{1.590891in}}%
\pgfpathlineto{\pgfqpoint{2.128613in}{1.572076in}}%
\pgfpathlineto{\pgfqpoint{2.118357in}{1.553854in}}%
\pgfpathlineto{\pgfqpoint{2.114334in}{1.552060in}}%
\pgfusepath{stroke}%
\end{pgfscope}%
\begin{pgfscope}%
\pgfpathrectangle{\pgfqpoint{0.100000in}{0.100000in}}{\pgfqpoint{3.007045in}{1.925000in}}%
\pgfusepath{clip}%
\pgfsetbuttcap%
\pgfsetroundjoin%
\pgfsetlinewidth{0.050187pt}%
\definecolor{currentstroke}{rgb}{1.000000,1.000000,1.000000}%
\pgfsetstrokecolor{currentstroke}%
\pgfsetdash{}{0pt}%
\pgfpathmoveto{\pgfqpoint{2.173359in}{1.333287in}}%
\pgfpathlineto{\pgfqpoint{2.183137in}{1.343579in}}%
\pgfpathlineto{\pgfqpoint{2.187591in}{1.358500in}}%
\pgfpathlineto{\pgfqpoint{2.192889in}{1.367149in}}%
\pgfpathlineto{\pgfqpoint{2.196137in}{1.378919in}}%
\pgfpathlineto{\pgfqpoint{2.197147in}{1.402387in}}%
\pgfpathlineto{\pgfqpoint{2.192256in}{1.424878in}}%
\pgfpathlineto{\pgfqpoint{2.176094in}{1.459319in}}%
\pgfpathlineto{\pgfqpoint{2.180471in}{1.470147in}}%
\pgfpathlineto{\pgfqpoint{2.174769in}{1.485115in}}%
\pgfpathlineto{\pgfqpoint{2.184550in}{1.505816in}}%
\pgfpathlineto{\pgfqpoint{2.182956in}{1.528897in}}%
\pgfpathlineto{\pgfqpoint{2.189769in}{1.531802in}}%
\pgfpathlineto{\pgfqpoint{2.190625in}{1.542784in}}%
\pgfpathlineto{\pgfqpoint{2.202718in}{1.549833in}}%
\pgfpathlineto{\pgfqpoint{2.209558in}{1.548695in}}%
\pgfpathlineto{\pgfqpoint{2.211480in}{1.536917in}}%
\pgfpathlineto{\pgfqpoint{2.216819in}{1.536455in}}%
\pgfpathlineto{\pgfqpoint{2.221689in}{1.553444in}}%
\pgfpathlineto{\pgfqpoint{2.220034in}{1.566657in}}%
\pgfpathlineto{\pgfqpoint{2.223201in}{1.574245in}}%
\pgfpathlineto{\pgfqpoint{2.240322in}{1.582083in}}%
\pgfpathlineto{\pgfqpoint{2.232520in}{1.584872in}}%
\pgfpathlineto{\pgfqpoint{2.229979in}{1.591605in}}%
\pgfpathlineto{\pgfqpoint{2.235599in}{1.603430in}}%
\pgfpathlineto{\pgfqpoint{2.246683in}{1.607514in}}%
\pgfpathlineto{\pgfqpoint{2.259554in}{1.600537in}}%
\pgfpathlineto{\pgfqpoint{2.271689in}{1.600450in}}%
\pgfpathlineto{\pgfqpoint{2.277325in}{1.592301in}}%
\pgfpathlineto{\pgfqpoint{2.285823in}{1.592867in}}%
\pgfpathlineto{\pgfqpoint{2.312058in}{1.581538in}}%
\pgfpathlineto{\pgfqpoint{2.317300in}{1.570571in}}%
\pgfpathlineto{\pgfqpoint{2.311575in}{1.566773in}}%
\pgfpathlineto{\pgfqpoint{2.313339in}{1.558555in}}%
\pgfpathlineto{\pgfqpoint{2.318994in}{1.554899in}}%
\pgfpathlineto{\pgfqpoint{2.322146in}{1.544796in}}%
\pgfpathlineto{\pgfqpoint{2.321705in}{1.520181in}}%
\pgfpathlineto{\pgfqpoint{2.314238in}{1.514301in}}%
\pgfpathlineto{\pgfqpoint{2.312562in}{1.501372in}}%
\pgfpathlineto{\pgfqpoint{2.298712in}{1.489423in}}%
\pgfpathlineto{\pgfqpoint{2.299534in}{1.474960in}}%
\pgfpathlineto{\pgfqpoint{2.311673in}{1.469893in}}%
\pgfpathlineto{\pgfqpoint{2.325368in}{1.487939in}}%
\pgfpathlineto{\pgfqpoint{2.326500in}{1.494484in}}%
\pgfpathlineto{\pgfqpoint{2.343579in}{1.505363in}}%
\pgfpathlineto{\pgfqpoint{2.354440in}{1.500324in}}%
\pgfpathlineto{\pgfqpoint{2.361256in}{1.488936in}}%
\pgfpathlineto{\pgfqpoint{2.372251in}{1.449395in}}%
\pgfpathlineto{\pgfqpoint{2.378077in}{1.436886in}}%
\pgfpathlineto{\pgfqpoint{2.376253in}{1.431715in}}%
\pgfpathlineto{\pgfqpoint{2.376430in}{1.414106in}}%
\pgfpathlineto{\pgfqpoint{2.365832in}{1.415796in}}%
\pgfpathlineto{\pgfqpoint{2.359845in}{1.402637in}}%
\pgfpathlineto{\pgfqpoint{2.359015in}{1.393692in}}%
\pgfpathlineto{\pgfqpoint{2.351009in}{1.387949in}}%
\pgfpathlineto{\pgfqpoint{2.349230in}{1.370503in}}%
\pgfpathlineto{\pgfqpoint{2.337587in}{1.348464in}}%
\pgfpathlineto{\pgfqpoint{2.273898in}{1.338979in}}%
\pgfpathlineto{\pgfqpoint{2.273522in}{1.343148in}}%
\pgfpathlineto{\pgfqpoint{2.230918in}{1.338658in}}%
\pgfpathlineto{\pgfqpoint{2.173359in}{1.333287in}}%
\pgfusepath{stroke}%
\end{pgfscope}%
\begin{pgfscope}%
\pgfpathrectangle{\pgfqpoint{0.100000in}{0.100000in}}{\pgfqpoint{3.007045in}{1.925000in}}%
\pgfusepath{clip}%
\pgfsetbuttcap%
\pgfsetroundjoin%
\pgfsetlinewidth{0.050187pt}%
\definecolor{currentstroke}{rgb}{1.000000,1.000000,1.000000}%
\pgfsetstrokecolor{currentstroke}%
\pgfsetdash{}{0pt}%
\pgfpathmoveto{\pgfqpoint{0.000000in}{0.000000in}}%
\pgfusepath{stroke}%
\end{pgfscope}%
\begin{pgfscope}%
\pgfpathrectangle{\pgfqpoint{0.100000in}{0.100000in}}{\pgfqpoint{3.007045in}{1.925000in}}%
\pgfusepath{clip}%
\pgfsetbuttcap%
\pgfsetroundjoin%
\pgfsetlinewidth{0.050187pt}%
\definecolor{currentstroke}{rgb}{1.000000,1.000000,1.000000}%
\pgfsetstrokecolor{currentstroke}%
\pgfsetdash{}{0pt}%
\pgfusepath{stroke}%
\end{pgfscope}%
\begin{pgfscope}%
\pgfpathrectangle{\pgfqpoint{0.100000in}{0.100000in}}{\pgfqpoint{3.007045in}{1.925000in}}%
\pgfusepath{clip}%
\pgfsetbuttcap%
\pgfsetroundjoin%
\pgfsetlinewidth{0.050187pt}%
\definecolor{currentstroke}{rgb}{1.000000,1.000000,1.000000}%
\pgfsetstrokecolor{currentstroke}%
\pgfsetdash{}{0pt}%
\pgfusepath{stroke}%
\end{pgfscope}%
\begin{pgfscope}%
\pgfpathrectangle{\pgfqpoint{0.100000in}{0.100000in}}{\pgfqpoint{3.007045in}{1.925000in}}%
\pgfusepath{clip}%
\pgfsetbuttcap%
\pgfsetroundjoin%
\pgfsetlinewidth{0.050187pt}%
\definecolor{currentstroke}{rgb}{1.000000,1.000000,1.000000}%
\pgfsetstrokecolor{currentstroke}%
\pgfsetdash{}{0pt}%
\pgfusepath{stroke}%
\end{pgfscope}%
\begin{pgfscope}%
\pgfpathrectangle{\pgfqpoint{0.100000in}{0.100000in}}{\pgfqpoint{3.007045in}{1.925000in}}%
\pgfusepath{clip}%
\pgfsetbuttcap%
\pgfsetroundjoin%
\pgfsetlinewidth{0.050187pt}%
\definecolor{currentstroke}{rgb}{1.000000,1.000000,1.000000}%
\pgfsetstrokecolor{currentstroke}%
\pgfsetdash{}{0pt}%
\pgfusepath{stroke}%
\end{pgfscope}%
\begin{pgfscope}%
\pgfpathrectangle{\pgfqpoint{0.100000in}{0.100000in}}{\pgfqpoint{3.007045in}{1.925000in}}%
\pgfusepath{clip}%
\pgfsetbuttcap%
\pgfsetroundjoin%
\pgfsetlinewidth{0.050187pt}%
\definecolor{currentstroke}{rgb}{1.000000,1.000000,1.000000}%
\pgfsetstrokecolor{currentstroke}%
\pgfsetdash{}{0pt}%
\pgfusepath{stroke}%
\end{pgfscope}%
\begin{pgfscope}%
\pgfpathrectangle{\pgfqpoint{0.100000in}{0.100000in}}{\pgfqpoint{3.007045in}{1.925000in}}%
\pgfusepath{clip}%
\pgfsetbuttcap%
\pgfsetroundjoin%
\pgfsetlinewidth{0.050187pt}%
\definecolor{currentstroke}{rgb}{1.000000,1.000000,1.000000}%
\pgfsetstrokecolor{currentstroke}%
\pgfsetdash{}{0pt}%
\pgfusepath{stroke}%
\end{pgfscope}%
\begin{pgfscope}%
\pgfpathrectangle{\pgfqpoint{0.100000in}{0.100000in}}{\pgfqpoint{3.007045in}{1.925000in}}%
\pgfusepath{clip}%
\pgfsetbuttcap%
\pgfsetroundjoin%
\pgfsetlinewidth{0.050187pt}%
\definecolor{currentstroke}{rgb}{1.000000,1.000000,1.000000}%
\pgfsetstrokecolor{currentstroke}%
\pgfsetdash{}{0pt}%
\pgfusepath{stroke}%
\end{pgfscope}%
\begin{pgfscope}%
\pgfpathrectangle{\pgfqpoint{0.100000in}{0.100000in}}{\pgfqpoint{3.007045in}{1.925000in}}%
\pgfusepath{clip}%
\pgfsetbuttcap%
\pgfsetroundjoin%
\pgfsetlinewidth{0.050187pt}%
\definecolor{currentstroke}{rgb}{1.000000,1.000000,1.000000}%
\pgfsetstrokecolor{currentstroke}%
\pgfsetdash{}{0pt}%
\pgfusepath{stroke}%
\end{pgfscope}%
\begin{pgfscope}%
\pgfpathrectangle{\pgfqpoint{0.100000in}{0.100000in}}{\pgfqpoint{3.007045in}{1.925000in}}%
\pgfusepath{clip}%
\pgfsetbuttcap%
\pgfsetroundjoin%
\pgfsetlinewidth{0.050187pt}%
\definecolor{currentstroke}{rgb}{1.000000,1.000000,1.000000}%
\pgfsetstrokecolor{currentstroke}%
\pgfsetdash{}{0pt}%
\pgfusepath{stroke}%
\end{pgfscope}%
\begin{pgfscope}%
\pgfpathrectangle{\pgfqpoint{0.100000in}{0.100000in}}{\pgfqpoint{3.007045in}{1.925000in}}%
\pgfusepath{clip}%
\pgfsetbuttcap%
\pgfsetroundjoin%
\pgfsetlinewidth{0.050187pt}%
\definecolor{currentstroke}{rgb}{1.000000,1.000000,1.000000}%
\pgfsetstrokecolor{currentstroke}%
\pgfsetdash{}{0pt}%
\pgfusepath{stroke}%
\end{pgfscope}%
\begin{pgfscope}%
\pgfpathrectangle{\pgfqpoint{0.100000in}{0.100000in}}{\pgfqpoint{3.007045in}{1.925000in}}%
\pgfusepath{clip}%
\pgfsetbuttcap%
\pgfsetroundjoin%
\pgfsetlinewidth{0.050187pt}%
\definecolor{currentstroke}{rgb}{1.000000,1.000000,1.000000}%
\pgfsetstrokecolor{currentstroke}%
\pgfsetdash{}{0pt}%
\pgfusepath{stroke}%
\end{pgfscope}%
\begin{pgfscope}%
\pgfpathrectangle{\pgfqpoint{0.100000in}{0.100000in}}{\pgfqpoint{3.007045in}{1.925000in}}%
\pgfusepath{clip}%
\pgfsetbuttcap%
\pgfsetroundjoin%
\pgfsetlinewidth{0.050187pt}%
\definecolor{currentstroke}{rgb}{1.000000,1.000000,1.000000}%
\pgfsetstrokecolor{currentstroke}%
\pgfsetdash{}{0pt}%
\pgfusepath{stroke}%
\end{pgfscope}%
\begin{pgfscope}%
\pgfpathrectangle{\pgfqpoint{0.100000in}{0.100000in}}{\pgfqpoint{3.007045in}{1.925000in}}%
\pgfusepath{clip}%
\pgfsetbuttcap%
\pgfsetroundjoin%
\pgfsetlinewidth{0.050187pt}%
\definecolor{currentstroke}{rgb}{1.000000,1.000000,1.000000}%
\pgfsetstrokecolor{currentstroke}%
\pgfsetdash{}{0pt}%
\pgfusepath{stroke}%
\end{pgfscope}%
\begin{pgfscope}%
\pgfpathrectangle{\pgfqpoint{0.100000in}{0.100000in}}{\pgfqpoint{3.007045in}{1.925000in}}%
\pgfusepath{clip}%
\pgfsetbuttcap%
\pgfsetroundjoin%
\pgfsetlinewidth{0.050187pt}%
\definecolor{currentstroke}{rgb}{1.000000,1.000000,1.000000}%
\pgfsetstrokecolor{currentstroke}%
\pgfsetdash{}{0pt}%
\pgfusepath{stroke}%
\end{pgfscope}%
\begin{pgfscope}%
\pgfpathrectangle{\pgfqpoint{0.100000in}{0.100000in}}{\pgfqpoint{3.007045in}{1.925000in}}%
\pgfusepath{clip}%
\pgfsetbuttcap%
\pgfsetroundjoin%
\pgfsetlinewidth{0.050187pt}%
\definecolor{currentstroke}{rgb}{1.000000,1.000000,1.000000}%
\pgfsetstrokecolor{currentstroke}%
\pgfsetdash{}{0pt}%
\pgfusepath{stroke}%
\end{pgfscope}%
\begin{pgfscope}%
\pgfpathrectangle{\pgfqpoint{0.100000in}{0.100000in}}{\pgfqpoint{3.007045in}{1.925000in}}%
\pgfusepath{clip}%
\pgfsetbuttcap%
\pgfsetroundjoin%
\pgfsetlinewidth{0.050187pt}%
\definecolor{currentstroke}{rgb}{1.000000,1.000000,1.000000}%
\pgfsetstrokecolor{currentstroke}%
\pgfsetdash{}{0pt}%
\pgfusepath{stroke}%
\end{pgfscope}%
\begin{pgfscope}%
\pgfpathrectangle{\pgfqpoint{0.100000in}{0.100000in}}{\pgfqpoint{3.007045in}{1.925000in}}%
\pgfusepath{clip}%
\pgfsetbuttcap%
\pgfsetroundjoin%
\pgfsetlinewidth{0.050187pt}%
\definecolor{currentstroke}{rgb}{1.000000,1.000000,1.000000}%
\pgfsetstrokecolor{currentstroke}%
\pgfsetdash{}{0pt}%
\pgfusepath{stroke}%
\end{pgfscope}%
\begin{pgfscope}%
\pgfpathrectangle{\pgfqpoint{0.100000in}{0.100000in}}{\pgfqpoint{3.007045in}{1.925000in}}%
\pgfusepath{clip}%
\pgfsetbuttcap%
\pgfsetroundjoin%
\pgfsetlinewidth{0.050187pt}%
\definecolor{currentstroke}{rgb}{1.000000,1.000000,1.000000}%
\pgfsetstrokecolor{currentstroke}%
\pgfsetdash{}{0pt}%
\pgfusepath{stroke}%
\end{pgfscope}%
\begin{pgfscope}%
\pgfpathrectangle{\pgfqpoint{0.100000in}{0.100000in}}{\pgfqpoint{3.007045in}{1.925000in}}%
\pgfusepath{clip}%
\pgfsetbuttcap%
\pgfsetroundjoin%
\pgfsetlinewidth{0.050187pt}%
\definecolor{currentstroke}{rgb}{1.000000,1.000000,1.000000}%
\pgfsetstrokecolor{currentstroke}%
\pgfsetdash{}{0pt}%
\pgfusepath{stroke}%
\end{pgfscope}%
\begin{pgfscope}%
\pgfpathrectangle{\pgfqpoint{0.100000in}{0.100000in}}{\pgfqpoint{3.007045in}{1.925000in}}%
\pgfusepath{clip}%
\pgfsetbuttcap%
\pgfsetroundjoin%
\pgfsetlinewidth{0.050187pt}%
\definecolor{currentstroke}{rgb}{1.000000,1.000000,1.000000}%
\pgfsetstrokecolor{currentstroke}%
\pgfsetdash{}{0pt}%
\pgfusepath{stroke}%
\end{pgfscope}%
\begin{pgfscope}%
\pgfpathrectangle{\pgfqpoint{0.100000in}{0.100000in}}{\pgfqpoint{3.007045in}{1.925000in}}%
\pgfusepath{clip}%
\pgfsetbuttcap%
\pgfsetroundjoin%
\pgfsetlinewidth{0.050187pt}%
\definecolor{currentstroke}{rgb}{1.000000,1.000000,1.000000}%
\pgfsetstrokecolor{currentstroke}%
\pgfsetdash{}{0pt}%
\pgfusepath{stroke}%
\end{pgfscope}%
\begin{pgfscope}%
\pgfpathrectangle{\pgfqpoint{0.100000in}{0.100000in}}{\pgfqpoint{3.007045in}{1.925000in}}%
\pgfusepath{clip}%
\pgfsetbuttcap%
\pgfsetroundjoin%
\pgfsetlinewidth{0.050187pt}%
\definecolor{currentstroke}{rgb}{1.000000,1.000000,1.000000}%
\pgfsetstrokecolor{currentstroke}%
\pgfsetdash{}{0pt}%
\pgfusepath{stroke}%
\end{pgfscope}%
\begin{pgfscope}%
\pgfpathrectangle{\pgfqpoint{0.100000in}{0.100000in}}{\pgfqpoint{3.007045in}{1.925000in}}%
\pgfusepath{clip}%
\pgfsetbuttcap%
\pgfsetroundjoin%
\pgfsetlinewidth{0.050187pt}%
\definecolor{currentstroke}{rgb}{1.000000,1.000000,1.000000}%
\pgfsetstrokecolor{currentstroke}%
\pgfsetdash{}{0pt}%
\pgfusepath{stroke}%
\end{pgfscope}%
\begin{pgfscope}%
\pgfpathrectangle{\pgfqpoint{0.100000in}{0.100000in}}{\pgfqpoint{3.007045in}{1.925000in}}%
\pgfusepath{clip}%
\pgfsetbuttcap%
\pgfsetroundjoin%
\pgfsetlinewidth{0.050187pt}%
\definecolor{currentstroke}{rgb}{1.000000,1.000000,1.000000}%
\pgfsetstrokecolor{currentstroke}%
\pgfsetdash{}{0pt}%
\pgfusepath{stroke}%
\end{pgfscope}%
\begin{pgfscope}%
\pgfpathrectangle{\pgfqpoint{0.100000in}{0.100000in}}{\pgfqpoint{3.007045in}{1.925000in}}%
\pgfusepath{clip}%
\pgfsetbuttcap%
\pgfsetroundjoin%
\pgfsetlinewidth{0.050187pt}%
\definecolor{currentstroke}{rgb}{1.000000,1.000000,1.000000}%
\pgfsetstrokecolor{currentstroke}%
\pgfsetdash{}{0pt}%
\pgfusepath{stroke}%
\end{pgfscope}%
\begin{pgfscope}%
\pgfpathrectangle{\pgfqpoint{0.100000in}{0.100000in}}{\pgfqpoint{3.007045in}{1.925000in}}%
\pgfusepath{clip}%
\pgfsetbuttcap%
\pgfsetroundjoin%
\pgfsetlinewidth{0.050187pt}%
\definecolor{currentstroke}{rgb}{1.000000,1.000000,1.000000}%
\pgfsetstrokecolor{currentstroke}%
\pgfsetdash{}{0pt}%
\pgfusepath{stroke}%
\end{pgfscope}%
\begin{pgfscope}%
\pgfpathrectangle{\pgfqpoint{0.100000in}{0.100000in}}{\pgfqpoint{3.007045in}{1.925000in}}%
\pgfusepath{clip}%
\pgfsetbuttcap%
\pgfsetroundjoin%
\pgfsetlinewidth{0.050187pt}%
\definecolor{currentstroke}{rgb}{1.000000,1.000000,1.000000}%
\pgfsetstrokecolor{currentstroke}%
\pgfsetdash{}{0pt}%
\pgfusepath{stroke}%
\end{pgfscope}%
\begin{pgfscope}%
\pgfpathrectangle{\pgfqpoint{2.857045in}{0.350000in}}{\pgfqpoint{0.300705in}{0.962500in}}%
\pgfusepath{clip}%
\pgfsetbuttcap%
\pgfsetmiterjoin%
\definecolor{currentfill}{rgb}{0.893887,0.959123,0.874002}%
\pgfsetfillcolor{currentfill}%
\pgfsetlinewidth{0.000000pt}%
\definecolor{currentstroke}{rgb}{0.000000,0.000000,0.000000}%
\pgfsetstrokecolor{currentstroke}%
\pgfsetstrokeopacity{0.000000}%
\pgfsetdash{}{0pt}%
\pgfpathmoveto{\pgfqpoint{3.122141in}{0.350000in}}%
\pgfpathlineto{\pgfqpoint{3.157750in}{0.350000in}}%
\pgfpathlineto{\pgfqpoint{3.157750in}{0.446250in}}%
\pgfpathlineto{\pgfqpoint{3.122141in}{0.446250in}}%
\pgfpathlineto{\pgfqpoint{3.122141in}{0.350000in}}%
\pgfpathclose%
\pgfusepath{fill}%
\end{pgfscope}%
\begin{pgfscope}%
\pgfpathrectangle{\pgfqpoint{2.857045in}{0.350000in}}{\pgfqpoint{0.300705in}{0.962500in}}%
\pgfusepath{clip}%
\pgfsetbuttcap%
\pgfsetmiterjoin%
\definecolor{currentfill}{rgb}{0.893887,0.959123,0.874002}%
\pgfsetfillcolor{currentfill}%
\pgfsetlinewidth{0.000000pt}%
\definecolor{currentstroke}{rgb}{0.000000,0.000000,0.000000}%
\pgfsetstrokecolor{currentstroke}%
\pgfsetstrokeopacity{0.000000}%
\pgfsetdash{}{0pt}%
\pgfpathmoveto{\pgfqpoint{3.122141in}{0.359625in}}%
\pgfpathlineto{\pgfqpoint{3.157750in}{0.359625in}}%
\pgfpathlineto{\pgfqpoint{3.157750in}{0.455875in}}%
\pgfpathlineto{\pgfqpoint{3.122141in}{0.455875in}}%
\pgfpathlineto{\pgfqpoint{3.122141in}{0.359625in}}%
\pgfpathclose%
\pgfusepath{fill}%
\end{pgfscope}%
\begin{pgfscope}%
\pgfpathrectangle{\pgfqpoint{2.857045in}{0.350000in}}{\pgfqpoint{0.300705in}{0.962500in}}%
\pgfusepath{clip}%
\pgfsetbuttcap%
\pgfsetmiterjoin%
\definecolor{currentfill}{rgb}{0.875433,0.951742,0.854318}%
\pgfsetfillcolor{currentfill}%
\pgfsetlinewidth{0.000000pt}%
\definecolor{currentstroke}{rgb}{0.000000,0.000000,0.000000}%
\pgfsetstrokecolor{currentstroke}%
\pgfsetstrokeopacity{0.000000}%
\pgfsetdash{}{0pt}%
\pgfpathmoveto{\pgfqpoint{3.120801in}{0.369250in}}%
\pgfpathlineto{\pgfqpoint{3.157750in}{0.369250in}}%
\pgfpathlineto{\pgfqpoint{3.157750in}{0.465500in}}%
\pgfpathlineto{\pgfqpoint{3.120801in}{0.465500in}}%
\pgfpathlineto{\pgfqpoint{3.120801in}{0.369250in}}%
\pgfpathclose%
\pgfusepath{fill}%
\end{pgfscope}%
\begin{pgfscope}%
\pgfpathrectangle{\pgfqpoint{2.857045in}{0.350000in}}{\pgfqpoint{0.300705in}{0.962500in}}%
\pgfusepath{clip}%
\pgfsetbuttcap%
\pgfsetmiterjoin%
\definecolor{currentfill}{rgb}{0.871742,0.950265,0.850381}%
\pgfsetfillcolor{currentfill}%
\pgfsetlinewidth{0.000000pt}%
\definecolor{currentstroke}{rgb}{0.000000,0.000000,0.000000}%
\pgfsetstrokecolor{currentstroke}%
\pgfsetstrokeopacity{0.000000}%
\pgfsetdash{}{0pt}%
\pgfpathmoveto{\pgfqpoint{3.120341in}{0.378875in}}%
\pgfpathlineto{\pgfqpoint{3.157750in}{0.378875in}}%
\pgfpathlineto{\pgfqpoint{3.157750in}{0.475125in}}%
\pgfpathlineto{\pgfqpoint{3.120341in}{0.475125in}}%
\pgfpathlineto{\pgfqpoint{3.120341in}{0.378875in}}%
\pgfpathclose%
\pgfusepath{fill}%
\end{pgfscope}%
\begin{pgfscope}%
\pgfpathrectangle{\pgfqpoint{2.857045in}{0.350000in}}{\pgfqpoint{0.300705in}{0.962500in}}%
\pgfusepath{clip}%
\pgfsetbuttcap%
\pgfsetmiterjoin%
\definecolor{currentfill}{rgb}{0.871742,0.950265,0.850381}%
\pgfsetfillcolor{currentfill}%
\pgfsetlinewidth{0.000000pt}%
\definecolor{currentstroke}{rgb}{0.000000,0.000000,0.000000}%
\pgfsetstrokecolor{currentstroke}%
\pgfsetstrokeopacity{0.000000}%
\pgfsetdash{}{0pt}%
\pgfpathmoveto{\pgfqpoint{3.120341in}{0.388500in}}%
\pgfpathlineto{\pgfqpoint{3.157750in}{0.388500in}}%
\pgfpathlineto{\pgfqpoint{3.157750in}{0.484750in}}%
\pgfpathlineto{\pgfqpoint{3.120341in}{0.484750in}}%
\pgfpathlineto{\pgfqpoint{3.120341in}{0.388500in}}%
\pgfpathclose%
\pgfusepath{fill}%
\end{pgfscope}%
\begin{pgfscope}%
\pgfpathrectangle{\pgfqpoint{2.857045in}{0.350000in}}{\pgfqpoint{0.300705in}{0.962500in}}%
\pgfusepath{clip}%
\pgfsetbuttcap%
\pgfsetmiterjoin%
\definecolor{currentfill}{rgb}{0.871742,0.950265,0.850381}%
\pgfsetfillcolor{currentfill}%
\pgfsetlinewidth{0.000000pt}%
\definecolor{currentstroke}{rgb}{0.000000,0.000000,0.000000}%
\pgfsetstrokecolor{currentstroke}%
\pgfsetstrokeopacity{0.000000}%
\pgfsetdash{}{0pt}%
\pgfpathmoveto{\pgfqpoint{3.120282in}{0.398125in}}%
\pgfpathlineto{\pgfqpoint{3.157750in}{0.398125in}}%
\pgfpathlineto{\pgfqpoint{3.157750in}{0.494375in}}%
\pgfpathlineto{\pgfqpoint{3.120282in}{0.494375in}}%
\pgfpathlineto{\pgfqpoint{3.120282in}{0.398125in}}%
\pgfpathclose%
\pgfusepath{fill}%
\end{pgfscope}%
\begin{pgfscope}%
\pgfpathrectangle{\pgfqpoint{2.857045in}{0.350000in}}{\pgfqpoint{0.300705in}{0.962500in}}%
\pgfusepath{clip}%
\pgfsetbuttcap%
\pgfsetmiterjoin%
\definecolor{currentfill}{rgb}{0.845905,0.939931,0.822822}%
\pgfsetfillcolor{currentfill}%
\pgfsetlinewidth{0.000000pt}%
\definecolor{currentstroke}{rgb}{0.000000,0.000000,0.000000}%
\pgfsetstrokecolor{currentstroke}%
\pgfsetstrokeopacity{0.000000}%
\pgfsetdash{}{0pt}%
\pgfpathmoveto{\pgfqpoint{3.118183in}{0.407750in}}%
\pgfpathlineto{\pgfqpoint{3.157750in}{0.407750in}}%
\pgfpathlineto{\pgfqpoint{3.157750in}{0.504000in}}%
\pgfpathlineto{\pgfqpoint{3.118183in}{0.504000in}}%
\pgfpathlineto{\pgfqpoint{3.118183in}{0.407750in}}%
\pgfpathclose%
\pgfusepath{fill}%
\end{pgfscope}%
\begin{pgfscope}%
\pgfpathrectangle{\pgfqpoint{2.857045in}{0.350000in}}{\pgfqpoint{0.300705in}{0.962500in}}%
\pgfusepath{clip}%
\pgfsetbuttcap%
\pgfsetmiterjoin%
\definecolor{currentfill}{rgb}{0.838524,0.936978,0.814948}%
\pgfsetfillcolor{currentfill}%
\pgfsetlinewidth{0.000000pt}%
\definecolor{currentstroke}{rgb}{0.000000,0.000000,0.000000}%
\pgfsetstrokecolor{currentstroke}%
\pgfsetstrokeopacity{0.000000}%
\pgfsetdash{}{0pt}%
\pgfpathmoveto{\pgfqpoint{3.117663in}{0.417375in}}%
\pgfpathlineto{\pgfqpoint{3.157750in}{0.417375in}}%
\pgfpathlineto{\pgfqpoint{3.157750in}{0.513625in}}%
\pgfpathlineto{\pgfqpoint{3.117663in}{0.513625in}}%
\pgfpathlineto{\pgfqpoint{3.117663in}{0.417375in}}%
\pgfpathclose%
\pgfusepath{fill}%
\end{pgfscope}%
\begin{pgfscope}%
\pgfpathrectangle{\pgfqpoint{2.857045in}{0.350000in}}{\pgfqpoint{0.300705in}{0.962500in}}%
\pgfusepath{clip}%
\pgfsetbuttcap%
\pgfsetmiterjoin%
\definecolor{currentfill}{rgb}{0.838524,0.936978,0.814948}%
\pgfsetfillcolor{currentfill}%
\pgfsetlinewidth{0.000000pt}%
\definecolor{currentstroke}{rgb}{0.000000,0.000000,0.000000}%
\pgfsetstrokecolor{currentstroke}%
\pgfsetstrokeopacity{0.000000}%
\pgfsetdash{}{0pt}%
\pgfpathmoveto{\pgfqpoint{3.117663in}{0.427000in}}%
\pgfpathlineto{\pgfqpoint{3.157750in}{0.427000in}}%
\pgfpathlineto{\pgfqpoint{3.157750in}{0.523250in}}%
\pgfpathlineto{\pgfqpoint{3.117663in}{0.523250in}}%
\pgfpathlineto{\pgfqpoint{3.117663in}{0.427000in}}%
\pgfpathclose%
\pgfusepath{fill}%
\end{pgfscope}%
\begin{pgfscope}%
\pgfpathrectangle{\pgfqpoint{2.857045in}{0.350000in}}{\pgfqpoint{0.300705in}{0.962500in}}%
\pgfusepath{clip}%
\pgfsetbuttcap%
\pgfsetmiterjoin%
\definecolor{currentfill}{rgb}{0.834833,0.935502,0.811011}%
\pgfsetfillcolor{currentfill}%
\pgfsetlinewidth{0.000000pt}%
\definecolor{currentstroke}{rgb}{0.000000,0.000000,0.000000}%
\pgfsetstrokecolor{currentstroke}%
\pgfsetstrokeopacity{0.000000}%
\pgfsetdash{}{0pt}%
\pgfpathmoveto{\pgfqpoint{3.117227in}{0.436625in}}%
\pgfpathlineto{\pgfqpoint{3.157750in}{0.436625in}}%
\pgfpathlineto{\pgfqpoint{3.157750in}{0.532875in}}%
\pgfpathlineto{\pgfqpoint{3.117227in}{0.532875in}}%
\pgfpathlineto{\pgfqpoint{3.117227in}{0.436625in}}%
\pgfpathclose%
\pgfusepath{fill}%
\end{pgfscope}%
\begin{pgfscope}%
\pgfpathrectangle{\pgfqpoint{2.857045in}{0.350000in}}{\pgfqpoint{0.300705in}{0.962500in}}%
\pgfusepath{clip}%
\pgfsetbuttcap%
\pgfsetmiterjoin%
\definecolor{currentfill}{rgb}{0.820069,0.929596,0.795263}%
\pgfsetfillcolor{currentfill}%
\pgfsetlinewidth{0.000000pt}%
\definecolor{currentstroke}{rgb}{0.000000,0.000000,0.000000}%
\pgfsetstrokecolor{currentstroke}%
\pgfsetstrokeopacity{0.000000}%
\pgfsetdash{}{0pt}%
\pgfpathmoveto{\pgfqpoint{3.115878in}{0.446250in}}%
\pgfpathlineto{\pgfqpoint{3.157750in}{0.446250in}}%
\pgfpathlineto{\pgfqpoint{3.157750in}{0.542500in}}%
\pgfpathlineto{\pgfqpoint{3.115878in}{0.542500in}}%
\pgfpathlineto{\pgfqpoint{3.115878in}{0.446250in}}%
\pgfpathclose%
\pgfusepath{fill}%
\end{pgfscope}%
\begin{pgfscope}%
\pgfpathrectangle{\pgfqpoint{2.857045in}{0.350000in}}{\pgfqpoint{0.300705in}{0.962500in}}%
\pgfusepath{clip}%
\pgfsetbuttcap%
\pgfsetmiterjoin%
\definecolor{currentfill}{rgb}{0.812687,0.926644,0.787389}%
\pgfsetfillcolor{currentfill}%
\pgfsetlinewidth{0.000000pt}%
\definecolor{currentstroke}{rgb}{0.000000,0.000000,0.000000}%
\pgfsetstrokecolor{currentstroke}%
\pgfsetstrokeopacity{0.000000}%
\pgfsetdash{}{0pt}%
\pgfpathmoveto{\pgfqpoint{3.115444in}{0.455875in}}%
\pgfpathlineto{\pgfqpoint{3.157750in}{0.455875in}}%
\pgfpathlineto{\pgfqpoint{3.157750in}{0.552125in}}%
\pgfpathlineto{\pgfqpoint{3.115444in}{0.552125in}}%
\pgfpathlineto{\pgfqpoint{3.115444in}{0.455875in}}%
\pgfpathclose%
\pgfusepath{fill}%
\end{pgfscope}%
\begin{pgfscope}%
\pgfpathrectangle{\pgfqpoint{2.857045in}{0.350000in}}{\pgfqpoint{0.300705in}{0.962500in}}%
\pgfusepath{clip}%
\pgfsetbuttcap%
\pgfsetmiterjoin%
\definecolor{currentfill}{rgb}{0.808997,0.925167,0.783453}%
\pgfsetfillcolor{currentfill}%
\pgfsetlinewidth{0.000000pt}%
\definecolor{currentstroke}{rgb}{0.000000,0.000000,0.000000}%
\pgfsetstrokecolor{currentstroke}%
\pgfsetstrokeopacity{0.000000}%
\pgfsetdash{}{0pt}%
\pgfpathmoveto{\pgfqpoint{3.114889in}{0.465500in}}%
\pgfpathlineto{\pgfqpoint{3.157750in}{0.465500in}}%
\pgfpathlineto{\pgfqpoint{3.157750in}{0.561750in}}%
\pgfpathlineto{\pgfqpoint{3.114889in}{0.561750in}}%
\pgfpathlineto{\pgfqpoint{3.114889in}{0.465500in}}%
\pgfpathclose%
\pgfusepath{fill}%
\end{pgfscope}%
\begin{pgfscope}%
\pgfpathrectangle{\pgfqpoint{2.857045in}{0.350000in}}{\pgfqpoint{0.300705in}{0.962500in}}%
\pgfusepath{clip}%
\pgfsetbuttcap%
\pgfsetmiterjoin%
\definecolor{currentfill}{rgb}{0.797924,0.920738,0.771642}%
\pgfsetfillcolor{currentfill}%
\pgfsetlinewidth{0.000000pt}%
\definecolor{currentstroke}{rgb}{0.000000,0.000000,0.000000}%
\pgfsetstrokecolor{currentstroke}%
\pgfsetstrokeopacity{0.000000}%
\pgfsetdash{}{0pt}%
\pgfpathmoveto{\pgfqpoint{3.114228in}{0.475125in}}%
\pgfpathlineto{\pgfqpoint{3.157750in}{0.475125in}}%
\pgfpathlineto{\pgfqpoint{3.157750in}{0.571375in}}%
\pgfpathlineto{\pgfqpoint{3.114228in}{0.571375in}}%
\pgfpathlineto{\pgfqpoint{3.114228in}{0.475125in}}%
\pgfpathclose%
\pgfusepath{fill}%
\end{pgfscope}%
\begin{pgfscope}%
\pgfpathrectangle{\pgfqpoint{2.857045in}{0.350000in}}{\pgfqpoint{0.300705in}{0.962500in}}%
\pgfusepath{clip}%
\pgfsetbuttcap%
\pgfsetmiterjoin%
\definecolor{currentfill}{rgb}{0.797924,0.920738,0.771642}%
\pgfsetfillcolor{currentfill}%
\pgfsetlinewidth{0.000000pt}%
\definecolor{currentstroke}{rgb}{0.000000,0.000000,0.000000}%
\pgfsetstrokecolor{currentstroke}%
\pgfsetstrokeopacity{0.000000}%
\pgfsetdash{}{0pt}%
\pgfpathmoveto{\pgfqpoint{3.114177in}{0.484750in}}%
\pgfpathlineto{\pgfqpoint{3.157750in}{0.484750in}}%
\pgfpathlineto{\pgfqpoint{3.157750in}{0.581000in}}%
\pgfpathlineto{\pgfqpoint{3.114177in}{0.581000in}}%
\pgfpathlineto{\pgfqpoint{3.114177in}{0.484750in}}%
\pgfpathclose%
\pgfusepath{fill}%
\end{pgfscope}%
\begin{pgfscope}%
\pgfpathrectangle{\pgfqpoint{2.857045in}{0.350000in}}{\pgfqpoint{0.300705in}{0.962500in}}%
\pgfusepath{clip}%
\pgfsetbuttcap%
\pgfsetmiterjoin%
\definecolor{currentfill}{rgb}{0.797924,0.920738,0.771642}%
\pgfsetfillcolor{currentfill}%
\pgfsetlinewidth{0.000000pt}%
\definecolor{currentstroke}{rgb}{0.000000,0.000000,0.000000}%
\pgfsetstrokecolor{currentstroke}%
\pgfsetstrokeopacity{0.000000}%
\pgfsetdash{}{0pt}%
\pgfpathmoveto{\pgfqpoint{3.114177in}{0.494375in}}%
\pgfpathlineto{\pgfqpoint{3.157750in}{0.494375in}}%
\pgfpathlineto{\pgfqpoint{3.157750in}{0.590625in}}%
\pgfpathlineto{\pgfqpoint{3.114177in}{0.590625in}}%
\pgfpathlineto{\pgfqpoint{3.114177in}{0.494375in}}%
\pgfpathclose%
\pgfusepath{fill}%
\end{pgfscope}%
\begin{pgfscope}%
\pgfpathrectangle{\pgfqpoint{2.857045in}{0.350000in}}{\pgfqpoint{0.300705in}{0.962500in}}%
\pgfusepath{clip}%
\pgfsetbuttcap%
\pgfsetmiterjoin%
\definecolor{currentfill}{rgb}{0.797924,0.920738,0.771642}%
\pgfsetfillcolor{currentfill}%
\pgfsetlinewidth{0.000000pt}%
\definecolor{currentstroke}{rgb}{0.000000,0.000000,0.000000}%
\pgfsetstrokecolor{currentstroke}%
\pgfsetstrokeopacity{0.000000}%
\pgfsetdash{}{0pt}%
\pgfpathmoveto{\pgfqpoint{3.113939in}{0.504000in}}%
\pgfpathlineto{\pgfqpoint{3.157750in}{0.504000in}}%
\pgfpathlineto{\pgfqpoint{3.157750in}{0.600250in}}%
\pgfpathlineto{\pgfqpoint{3.113939in}{0.600250in}}%
\pgfpathlineto{\pgfqpoint{3.113939in}{0.504000in}}%
\pgfpathclose%
\pgfusepath{fill}%
\end{pgfscope}%
\begin{pgfscope}%
\pgfpathrectangle{\pgfqpoint{2.857045in}{0.350000in}}{\pgfqpoint{0.300705in}{0.962500in}}%
\pgfusepath{clip}%
\pgfsetbuttcap%
\pgfsetmiterjoin%
\definecolor{currentfill}{rgb}{0.794233,0.919262,0.767705}%
\pgfsetfillcolor{currentfill}%
\pgfsetlinewidth{0.000000pt}%
\definecolor{currentstroke}{rgb}{0.000000,0.000000,0.000000}%
\pgfsetstrokecolor{currentstroke}%
\pgfsetstrokeopacity{0.000000}%
\pgfsetdash{}{0pt}%
\pgfpathmoveto{\pgfqpoint{3.113669in}{0.513625in}}%
\pgfpathlineto{\pgfqpoint{3.157750in}{0.513625in}}%
\pgfpathlineto{\pgfqpoint{3.157750in}{0.609875in}}%
\pgfpathlineto{\pgfqpoint{3.113669in}{0.609875in}}%
\pgfpathlineto{\pgfqpoint{3.113669in}{0.513625in}}%
\pgfpathclose%
\pgfusepath{fill}%
\end{pgfscope}%
\begin{pgfscope}%
\pgfpathrectangle{\pgfqpoint{2.857045in}{0.350000in}}{\pgfqpoint{0.300705in}{0.962500in}}%
\pgfusepath{clip}%
\pgfsetbuttcap%
\pgfsetmiterjoin%
\definecolor{currentfill}{rgb}{0.783160,0.914833,0.755894}%
\pgfsetfillcolor{currentfill}%
\pgfsetlinewidth{0.000000pt}%
\definecolor{currentstroke}{rgb}{0.000000,0.000000,0.000000}%
\pgfsetstrokecolor{currentstroke}%
\pgfsetstrokeopacity{0.000000}%
\pgfsetdash{}{0pt}%
\pgfpathmoveto{\pgfqpoint{3.112932in}{0.523250in}}%
\pgfpathlineto{\pgfqpoint{3.157750in}{0.523250in}}%
\pgfpathlineto{\pgfqpoint{3.157750in}{0.619500in}}%
\pgfpathlineto{\pgfqpoint{3.112932in}{0.619500in}}%
\pgfpathlineto{\pgfqpoint{3.112932in}{0.523250in}}%
\pgfpathclose%
\pgfusepath{fill}%
\end{pgfscope}%
\begin{pgfscope}%
\pgfpathrectangle{\pgfqpoint{2.857045in}{0.350000in}}{\pgfqpoint{0.300705in}{0.962500in}}%
\pgfusepath{clip}%
\pgfsetbuttcap%
\pgfsetmiterjoin%
\definecolor{currentfill}{rgb}{0.769873,0.909296,0.742699}%
\pgfsetfillcolor{currentfill}%
\pgfsetlinewidth{0.000000pt}%
\definecolor{currentstroke}{rgb}{0.000000,0.000000,0.000000}%
\pgfsetstrokecolor{currentstroke}%
\pgfsetstrokeopacity{0.000000}%
\pgfsetdash{}{0pt}%
\pgfpathmoveto{\pgfqpoint{3.111877in}{0.532875in}}%
\pgfpathlineto{\pgfqpoint{3.157750in}{0.532875in}}%
\pgfpathlineto{\pgfqpoint{3.157750in}{0.629125in}}%
\pgfpathlineto{\pgfqpoint{3.111877in}{0.629125in}}%
\pgfpathlineto{\pgfqpoint{3.111877in}{0.532875in}}%
\pgfpathclose%
\pgfusepath{fill}%
\end{pgfscope}%
\begin{pgfscope}%
\pgfpathrectangle{\pgfqpoint{2.857045in}{0.350000in}}{\pgfqpoint{0.300705in}{0.962500in}}%
\pgfusepath{clip}%
\pgfsetbuttcap%
\pgfsetmiterjoin%
\definecolor{currentfill}{rgb}{0.769873,0.909296,0.742699}%
\pgfsetfillcolor{currentfill}%
\pgfsetlinewidth{0.000000pt}%
\definecolor{currentstroke}{rgb}{0.000000,0.000000,0.000000}%
\pgfsetstrokecolor{currentstroke}%
\pgfsetstrokeopacity{0.000000}%
\pgfsetdash{}{0pt}%
\pgfpathmoveto{\pgfqpoint{3.111877in}{0.542500in}}%
\pgfpathlineto{\pgfqpoint{3.157750in}{0.542500in}}%
\pgfpathlineto{\pgfqpoint{3.157750in}{0.638750in}}%
\pgfpathlineto{\pgfqpoint{3.111877in}{0.638750in}}%
\pgfpathlineto{\pgfqpoint{3.111877in}{0.542500in}}%
\pgfpathclose%
\pgfusepath{fill}%
\end{pgfscope}%
\begin{pgfscope}%
\pgfpathrectangle{\pgfqpoint{2.857045in}{0.350000in}}{\pgfqpoint{0.300705in}{0.962500in}}%
\pgfusepath{clip}%
\pgfsetbuttcap%
\pgfsetmiterjoin%
\definecolor{currentfill}{rgb}{0.769873,0.909296,0.742699}%
\pgfsetfillcolor{currentfill}%
\pgfsetlinewidth{0.000000pt}%
\definecolor{currentstroke}{rgb}{0.000000,0.000000,0.000000}%
\pgfsetstrokecolor{currentstroke}%
\pgfsetstrokeopacity{0.000000}%
\pgfsetdash{}{0pt}%
\pgfpathmoveto{\pgfqpoint{3.111877in}{0.552125in}}%
\pgfpathlineto{\pgfqpoint{3.157750in}{0.552125in}}%
\pgfpathlineto{\pgfqpoint{3.157750in}{0.648375in}}%
\pgfpathlineto{\pgfqpoint{3.111877in}{0.648375in}}%
\pgfpathlineto{\pgfqpoint{3.111877in}{0.552125in}}%
\pgfpathclose%
\pgfusepath{fill}%
\end{pgfscope}%
\begin{pgfscope}%
\pgfpathrectangle{\pgfqpoint{2.857045in}{0.350000in}}{\pgfqpoint{0.300705in}{0.962500in}}%
\pgfusepath{clip}%
\pgfsetbuttcap%
\pgfsetmiterjoin%
\definecolor{currentfill}{rgb}{0.769873,0.909296,0.742699}%
\pgfsetfillcolor{currentfill}%
\pgfsetlinewidth{0.000000pt}%
\definecolor{currentstroke}{rgb}{0.000000,0.000000,0.000000}%
\pgfsetstrokecolor{currentstroke}%
\pgfsetstrokeopacity{0.000000}%
\pgfsetdash{}{0pt}%
\pgfpathmoveto{\pgfqpoint{3.111877in}{0.561750in}}%
\pgfpathlineto{\pgfqpoint{3.157750in}{0.561750in}}%
\pgfpathlineto{\pgfqpoint{3.157750in}{0.658000in}}%
\pgfpathlineto{\pgfqpoint{3.111877in}{0.658000in}}%
\pgfpathlineto{\pgfqpoint{3.111877in}{0.561750in}}%
\pgfpathclose%
\pgfusepath{fill}%
\end{pgfscope}%
\begin{pgfscope}%
\pgfpathrectangle{\pgfqpoint{2.857045in}{0.350000in}}{\pgfqpoint{0.300705in}{0.962500in}}%
\pgfusepath{clip}%
\pgfsetbuttcap%
\pgfsetmiterjoin%
\definecolor{currentfill}{rgb}{0.769873,0.909296,0.742699}%
\pgfsetfillcolor{currentfill}%
\pgfsetlinewidth{0.000000pt}%
\definecolor{currentstroke}{rgb}{0.000000,0.000000,0.000000}%
\pgfsetstrokecolor{currentstroke}%
\pgfsetstrokeopacity{0.000000}%
\pgfsetdash{}{0pt}%
\pgfpathmoveto{\pgfqpoint{3.111877in}{0.571375in}}%
\pgfpathlineto{\pgfqpoint{3.157750in}{0.571375in}}%
\pgfpathlineto{\pgfqpoint{3.157750in}{0.667625in}}%
\pgfpathlineto{\pgfqpoint{3.111877in}{0.667625in}}%
\pgfpathlineto{\pgfqpoint{3.111877in}{0.571375in}}%
\pgfpathclose%
\pgfusepath{fill}%
\end{pgfscope}%
\begin{pgfscope}%
\pgfpathrectangle{\pgfqpoint{2.857045in}{0.350000in}}{\pgfqpoint{0.300705in}{0.962500in}}%
\pgfusepath{clip}%
\pgfsetbuttcap%
\pgfsetmiterjoin%
\definecolor{currentfill}{rgb}{0.769873,0.909296,0.742699}%
\pgfsetfillcolor{currentfill}%
\pgfsetlinewidth{0.000000pt}%
\definecolor{currentstroke}{rgb}{0.000000,0.000000,0.000000}%
\pgfsetstrokecolor{currentstroke}%
\pgfsetstrokeopacity{0.000000}%
\pgfsetdash{}{0pt}%
\pgfpathmoveto{\pgfqpoint{3.111877in}{0.581000in}}%
\pgfpathlineto{\pgfqpoint{3.157750in}{0.581000in}}%
\pgfpathlineto{\pgfqpoint{3.157750in}{0.677250in}}%
\pgfpathlineto{\pgfqpoint{3.111877in}{0.677250in}}%
\pgfpathlineto{\pgfqpoint{3.111877in}{0.581000in}}%
\pgfpathclose%
\pgfusepath{fill}%
\end{pgfscope}%
\begin{pgfscope}%
\pgfpathrectangle{\pgfqpoint{2.857045in}{0.350000in}}{\pgfqpoint{0.300705in}{0.962500in}}%
\pgfusepath{clip}%
\pgfsetbuttcap%
\pgfsetmiterjoin%
\definecolor{currentfill}{rgb}{0.760523,0.905359,0.733595}%
\pgfsetfillcolor{currentfill}%
\pgfsetlinewidth{0.000000pt}%
\definecolor{currentstroke}{rgb}{0.000000,0.000000,0.000000}%
\pgfsetstrokecolor{currentstroke}%
\pgfsetstrokeopacity{0.000000}%
\pgfsetdash{}{0pt}%
\pgfpathmoveto{\pgfqpoint{3.111334in}{0.590625in}}%
\pgfpathlineto{\pgfqpoint{3.157750in}{0.590625in}}%
\pgfpathlineto{\pgfqpoint{3.157750in}{0.686875in}}%
\pgfpathlineto{\pgfqpoint{3.111334in}{0.686875in}}%
\pgfpathlineto{\pgfqpoint{3.111334in}{0.590625in}}%
\pgfpathclose%
\pgfusepath{fill}%
\end{pgfscope}%
\begin{pgfscope}%
\pgfpathrectangle{\pgfqpoint{2.857045in}{0.350000in}}{\pgfqpoint{0.300705in}{0.962500in}}%
\pgfusepath{clip}%
\pgfsetbuttcap%
\pgfsetmiterjoin%
\definecolor{currentfill}{rgb}{0.760523,0.905359,0.733595}%
\pgfsetfillcolor{currentfill}%
\pgfsetlinewidth{0.000000pt}%
\definecolor{currentstroke}{rgb}{0.000000,0.000000,0.000000}%
\pgfsetstrokecolor{currentstroke}%
\pgfsetstrokeopacity{0.000000}%
\pgfsetdash{}{0pt}%
\pgfpathmoveto{\pgfqpoint{3.111245in}{0.600250in}}%
\pgfpathlineto{\pgfqpoint{3.157750in}{0.600250in}}%
\pgfpathlineto{\pgfqpoint{3.157750in}{0.696500in}}%
\pgfpathlineto{\pgfqpoint{3.111245in}{0.696500in}}%
\pgfpathlineto{\pgfqpoint{3.111245in}{0.600250in}}%
\pgfpathclose%
\pgfusepath{fill}%
\end{pgfscope}%
\begin{pgfscope}%
\pgfpathrectangle{\pgfqpoint{2.857045in}{0.350000in}}{\pgfqpoint{0.300705in}{0.962500in}}%
\pgfusepath{clip}%
\pgfsetbuttcap%
\pgfsetmiterjoin%
\definecolor{currentfill}{rgb}{0.760523,0.905359,0.733595}%
\pgfsetfillcolor{currentfill}%
\pgfsetlinewidth{0.000000pt}%
\definecolor{currentstroke}{rgb}{0.000000,0.000000,0.000000}%
\pgfsetstrokecolor{currentstroke}%
\pgfsetstrokeopacity{0.000000}%
\pgfsetdash{}{0pt}%
\pgfpathmoveto{\pgfqpoint{3.111155in}{0.609875in}}%
\pgfpathlineto{\pgfqpoint{3.157750in}{0.609875in}}%
\pgfpathlineto{\pgfqpoint{3.157750in}{0.706125in}}%
\pgfpathlineto{\pgfqpoint{3.111155in}{0.706125in}}%
\pgfpathlineto{\pgfqpoint{3.111155in}{0.609875in}}%
\pgfpathclose%
\pgfusepath{fill}%
\end{pgfscope}%
\begin{pgfscope}%
\pgfpathrectangle{\pgfqpoint{2.857045in}{0.350000in}}{\pgfqpoint{0.300705in}{0.962500in}}%
\pgfusepath{clip}%
\pgfsetbuttcap%
\pgfsetmiterjoin%
\definecolor{currentfill}{rgb}{0.755848,0.903391,0.729043}%
\pgfsetfillcolor{currentfill}%
\pgfsetlinewidth{0.000000pt}%
\definecolor{currentstroke}{rgb}{0.000000,0.000000,0.000000}%
\pgfsetstrokecolor{currentstroke}%
\pgfsetstrokeopacity{0.000000}%
\pgfsetdash{}{0pt}%
\pgfpathmoveto{\pgfqpoint{3.110809in}{0.619500in}}%
\pgfpathlineto{\pgfqpoint{3.157750in}{0.619500in}}%
\pgfpathlineto{\pgfqpoint{3.157750in}{0.715750in}}%
\pgfpathlineto{\pgfqpoint{3.110809in}{0.715750in}}%
\pgfpathlineto{\pgfqpoint{3.110809in}{0.619500in}}%
\pgfpathclose%
\pgfusepath{fill}%
\end{pgfscope}%
\begin{pgfscope}%
\pgfpathrectangle{\pgfqpoint{2.857045in}{0.350000in}}{\pgfqpoint{0.300705in}{0.962500in}}%
\pgfusepath{clip}%
\pgfsetbuttcap%
\pgfsetmiterjoin%
\definecolor{currentfill}{rgb}{0.751173,0.901423,0.724491}%
\pgfsetfillcolor{currentfill}%
\pgfsetlinewidth{0.000000pt}%
\definecolor{currentstroke}{rgb}{0.000000,0.000000,0.000000}%
\pgfsetstrokecolor{currentstroke}%
\pgfsetstrokeopacity{0.000000}%
\pgfsetdash{}{0pt}%
\pgfpathmoveto{\pgfqpoint{3.110599in}{0.629125in}}%
\pgfpathlineto{\pgfqpoint{3.157750in}{0.629125in}}%
\pgfpathlineto{\pgfqpoint{3.157750in}{0.725375in}}%
\pgfpathlineto{\pgfqpoint{3.110599in}{0.725375in}}%
\pgfpathlineto{\pgfqpoint{3.110599in}{0.629125in}}%
\pgfpathclose%
\pgfusepath{fill}%
\end{pgfscope}%
\begin{pgfscope}%
\pgfpathrectangle{\pgfqpoint{2.857045in}{0.350000in}}{\pgfqpoint{0.300705in}{0.962500in}}%
\pgfusepath{clip}%
\pgfsetbuttcap%
\pgfsetmiterjoin%
\definecolor{currentfill}{rgb}{0.746498,0.899454,0.719938}%
\pgfsetfillcolor{currentfill}%
\pgfsetlinewidth{0.000000pt}%
\definecolor{currentstroke}{rgb}{0.000000,0.000000,0.000000}%
\pgfsetstrokecolor{currentstroke}%
\pgfsetstrokeopacity{0.000000}%
\pgfsetdash{}{0pt}%
\pgfpathmoveto{\pgfqpoint{3.110190in}{0.638750in}}%
\pgfpathlineto{\pgfqpoint{3.157750in}{0.638750in}}%
\pgfpathlineto{\pgfqpoint{3.157750in}{0.735000in}}%
\pgfpathlineto{\pgfqpoint{3.110190in}{0.735000in}}%
\pgfpathlineto{\pgfqpoint{3.110190in}{0.638750in}}%
\pgfpathclose%
\pgfusepath{fill}%
\end{pgfscope}%
\begin{pgfscope}%
\pgfpathrectangle{\pgfqpoint{2.857045in}{0.350000in}}{\pgfqpoint{0.300705in}{0.962500in}}%
\pgfusepath{clip}%
\pgfsetbuttcap%
\pgfsetmiterjoin%
\definecolor{currentfill}{rgb}{0.741822,0.897486,0.715386}%
\pgfsetfillcolor{currentfill}%
\pgfsetlinewidth{0.000000pt}%
\definecolor{currentstroke}{rgb}{0.000000,0.000000,0.000000}%
\pgfsetstrokecolor{currentstroke}%
\pgfsetstrokeopacity{0.000000}%
\pgfsetdash{}{0pt}%
\pgfpathmoveto{\pgfqpoint{3.110152in}{0.648375in}}%
\pgfpathlineto{\pgfqpoint{3.157750in}{0.648375in}}%
\pgfpathlineto{\pgfqpoint{3.157750in}{0.744625in}}%
\pgfpathlineto{\pgfqpoint{3.110152in}{0.744625in}}%
\pgfpathlineto{\pgfqpoint{3.110152in}{0.648375in}}%
\pgfpathclose%
\pgfusepath{fill}%
\end{pgfscope}%
\begin{pgfscope}%
\pgfpathrectangle{\pgfqpoint{2.857045in}{0.350000in}}{\pgfqpoint{0.300705in}{0.962500in}}%
\pgfusepath{clip}%
\pgfsetbuttcap%
\pgfsetmiterjoin%
\definecolor{currentfill}{rgb}{0.732472,0.893549,0.706282}%
\pgfsetfillcolor{currentfill}%
\pgfsetlinewidth{0.000000pt}%
\definecolor{currentstroke}{rgb}{0.000000,0.000000,0.000000}%
\pgfsetstrokecolor{currentstroke}%
\pgfsetstrokeopacity{0.000000}%
\pgfsetdash{}{0pt}%
\pgfpathmoveto{\pgfqpoint{3.109375in}{0.658000in}}%
\pgfpathlineto{\pgfqpoint{3.157750in}{0.658000in}}%
\pgfpathlineto{\pgfqpoint{3.157750in}{0.754250in}}%
\pgfpathlineto{\pgfqpoint{3.109375in}{0.754250in}}%
\pgfpathlineto{\pgfqpoint{3.109375in}{0.658000in}}%
\pgfpathclose%
\pgfusepath{fill}%
\end{pgfscope}%
\begin{pgfscope}%
\pgfpathrectangle{\pgfqpoint{2.857045in}{0.350000in}}{\pgfqpoint{0.300705in}{0.962500in}}%
\pgfusepath{clip}%
\pgfsetbuttcap%
\pgfsetmiterjoin%
\definecolor{currentfill}{rgb}{0.732472,0.893549,0.706282}%
\pgfsetfillcolor{currentfill}%
\pgfsetlinewidth{0.000000pt}%
\definecolor{currentstroke}{rgb}{0.000000,0.000000,0.000000}%
\pgfsetstrokecolor{currentstroke}%
\pgfsetstrokeopacity{0.000000}%
\pgfsetdash{}{0pt}%
\pgfpathmoveto{\pgfqpoint{3.109264in}{0.667625in}}%
\pgfpathlineto{\pgfqpoint{3.157750in}{0.667625in}}%
\pgfpathlineto{\pgfqpoint{3.157750in}{0.763875in}}%
\pgfpathlineto{\pgfqpoint{3.109264in}{0.763875in}}%
\pgfpathlineto{\pgfqpoint{3.109264in}{0.667625in}}%
\pgfpathclose%
\pgfusepath{fill}%
\end{pgfscope}%
\begin{pgfscope}%
\pgfpathrectangle{\pgfqpoint{2.857045in}{0.350000in}}{\pgfqpoint{0.300705in}{0.962500in}}%
\pgfusepath{clip}%
\pgfsetbuttcap%
\pgfsetmiterjoin%
\definecolor{currentfill}{rgb}{0.732472,0.893549,0.706282}%
\pgfsetfillcolor{currentfill}%
\pgfsetlinewidth{0.000000pt}%
\definecolor{currentstroke}{rgb}{0.000000,0.000000,0.000000}%
\pgfsetstrokecolor{currentstroke}%
\pgfsetstrokeopacity{0.000000}%
\pgfsetdash{}{0pt}%
\pgfpathmoveto{\pgfqpoint{3.109264in}{0.677250in}}%
\pgfpathlineto{\pgfqpoint{3.157750in}{0.677250in}}%
\pgfpathlineto{\pgfqpoint{3.157750in}{0.773500in}}%
\pgfpathlineto{\pgfqpoint{3.109264in}{0.773500in}}%
\pgfpathlineto{\pgfqpoint{3.109264in}{0.677250in}}%
\pgfpathclose%
\pgfusepath{fill}%
\end{pgfscope}%
\begin{pgfscope}%
\pgfpathrectangle{\pgfqpoint{2.857045in}{0.350000in}}{\pgfqpoint{0.300705in}{0.962500in}}%
\pgfusepath{clip}%
\pgfsetbuttcap%
\pgfsetmiterjoin%
\definecolor{currentfill}{rgb}{0.732472,0.893549,0.706282}%
\pgfsetfillcolor{currentfill}%
\pgfsetlinewidth{0.000000pt}%
\definecolor{currentstroke}{rgb}{0.000000,0.000000,0.000000}%
\pgfsetstrokecolor{currentstroke}%
\pgfsetstrokeopacity{0.000000}%
\pgfsetdash{}{0pt}%
\pgfpathmoveto{\pgfqpoint{3.109264in}{0.686875in}}%
\pgfpathlineto{\pgfqpoint{3.157750in}{0.686875in}}%
\pgfpathlineto{\pgfqpoint{3.157750in}{0.783125in}}%
\pgfpathlineto{\pgfqpoint{3.109264in}{0.783125in}}%
\pgfpathlineto{\pgfqpoint{3.109264in}{0.686875in}}%
\pgfpathclose%
\pgfusepath{fill}%
\end{pgfscope}%
\begin{pgfscope}%
\pgfpathrectangle{\pgfqpoint{2.857045in}{0.350000in}}{\pgfqpoint{0.300705in}{0.962500in}}%
\pgfusepath{clip}%
\pgfsetbuttcap%
\pgfsetmiterjoin%
\definecolor{currentfill}{rgb}{0.723122,0.889612,0.697178}%
\pgfsetfillcolor{currentfill}%
\pgfsetlinewidth{0.000000pt}%
\definecolor{currentstroke}{rgb}{0.000000,0.000000,0.000000}%
\pgfsetstrokecolor{currentstroke}%
\pgfsetstrokeopacity{0.000000}%
\pgfsetdash{}{0pt}%
\pgfpathmoveto{\pgfqpoint{3.108883in}{0.696500in}}%
\pgfpathlineto{\pgfqpoint{3.157750in}{0.696500in}}%
\pgfpathlineto{\pgfqpoint{3.157750in}{0.792750in}}%
\pgfpathlineto{\pgfqpoint{3.108883in}{0.792750in}}%
\pgfpathlineto{\pgfqpoint{3.108883in}{0.696500in}}%
\pgfpathclose%
\pgfusepath{fill}%
\end{pgfscope}%
\begin{pgfscope}%
\pgfpathrectangle{\pgfqpoint{2.857045in}{0.350000in}}{\pgfqpoint{0.300705in}{0.962500in}}%
\pgfusepath{clip}%
\pgfsetbuttcap%
\pgfsetmiterjoin%
\definecolor{currentfill}{rgb}{0.723122,0.889612,0.697178}%
\pgfsetfillcolor{currentfill}%
\pgfsetlinewidth{0.000000pt}%
\definecolor{currentstroke}{rgb}{0.000000,0.000000,0.000000}%
\pgfsetstrokecolor{currentstroke}%
\pgfsetstrokeopacity{0.000000}%
\pgfsetdash{}{0pt}%
\pgfpathmoveto{\pgfqpoint{3.108883in}{0.706125in}}%
\pgfpathlineto{\pgfqpoint{3.157750in}{0.706125in}}%
\pgfpathlineto{\pgfqpoint{3.157750in}{0.802375in}}%
\pgfpathlineto{\pgfqpoint{3.108883in}{0.802375in}}%
\pgfpathlineto{\pgfqpoint{3.108883in}{0.706125in}}%
\pgfpathclose%
\pgfusepath{fill}%
\end{pgfscope}%
\begin{pgfscope}%
\pgfpathrectangle{\pgfqpoint{2.857045in}{0.350000in}}{\pgfqpoint{0.300705in}{0.962500in}}%
\pgfusepath{clip}%
\pgfsetbuttcap%
\pgfsetmiterjoin%
\definecolor{currentfill}{rgb}{0.718447,0.887643,0.692626}%
\pgfsetfillcolor{currentfill}%
\pgfsetlinewidth{0.000000pt}%
\definecolor{currentstroke}{rgb}{0.000000,0.000000,0.000000}%
\pgfsetstrokecolor{currentstroke}%
\pgfsetstrokeopacity{0.000000}%
\pgfsetdash{}{0pt}%
\pgfpathmoveto{\pgfqpoint{3.108500in}{0.715750in}}%
\pgfpathlineto{\pgfqpoint{3.157750in}{0.715750in}}%
\pgfpathlineto{\pgfqpoint{3.157750in}{0.812000in}}%
\pgfpathlineto{\pgfqpoint{3.108500in}{0.812000in}}%
\pgfpathlineto{\pgfqpoint{3.108500in}{0.715750in}}%
\pgfpathclose%
\pgfusepath{fill}%
\end{pgfscope}%
\begin{pgfscope}%
\pgfpathrectangle{\pgfqpoint{2.857045in}{0.350000in}}{\pgfqpoint{0.300705in}{0.962500in}}%
\pgfusepath{clip}%
\pgfsetbuttcap%
\pgfsetmiterjoin%
\definecolor{currentfill}{rgb}{0.709097,0.883706,0.683522}%
\pgfsetfillcolor{currentfill}%
\pgfsetlinewidth{0.000000pt}%
\definecolor{currentstroke}{rgb}{0.000000,0.000000,0.000000}%
\pgfsetstrokecolor{currentstroke}%
\pgfsetstrokeopacity{0.000000}%
\pgfsetdash{}{0pt}%
\pgfpathmoveto{\pgfqpoint{3.107910in}{0.725375in}}%
\pgfpathlineto{\pgfqpoint{3.157750in}{0.725375in}}%
\pgfpathlineto{\pgfqpoint{3.157750in}{0.821625in}}%
\pgfpathlineto{\pgfqpoint{3.107910in}{0.821625in}}%
\pgfpathlineto{\pgfqpoint{3.107910in}{0.725375in}}%
\pgfpathclose%
\pgfusepath{fill}%
\end{pgfscope}%
\begin{pgfscope}%
\pgfpathrectangle{\pgfqpoint{2.857045in}{0.350000in}}{\pgfqpoint{0.300705in}{0.962500in}}%
\pgfusepath{clip}%
\pgfsetbuttcap%
\pgfsetmiterjoin%
\definecolor{currentfill}{rgb}{0.709097,0.883706,0.683522}%
\pgfsetfillcolor{currentfill}%
\pgfsetlinewidth{0.000000pt}%
\definecolor{currentstroke}{rgb}{0.000000,0.000000,0.000000}%
\pgfsetstrokecolor{currentstroke}%
\pgfsetstrokeopacity{0.000000}%
\pgfsetdash{}{0pt}%
\pgfpathmoveto{\pgfqpoint{3.107713in}{0.735000in}}%
\pgfpathlineto{\pgfqpoint{3.157750in}{0.735000in}}%
\pgfpathlineto{\pgfqpoint{3.157750in}{0.831250in}}%
\pgfpathlineto{\pgfqpoint{3.107713in}{0.831250in}}%
\pgfpathlineto{\pgfqpoint{3.107713in}{0.735000in}}%
\pgfpathclose%
\pgfusepath{fill}%
\end{pgfscope}%
\begin{pgfscope}%
\pgfpathrectangle{\pgfqpoint{2.857045in}{0.350000in}}{\pgfqpoint{0.300705in}{0.962500in}}%
\pgfusepath{clip}%
\pgfsetbuttcap%
\pgfsetmiterjoin%
\definecolor{currentfill}{rgb}{0.704421,0.881738,0.678970}%
\pgfsetfillcolor{currentfill}%
\pgfsetlinewidth{0.000000pt}%
\definecolor{currentstroke}{rgb}{0.000000,0.000000,0.000000}%
\pgfsetstrokecolor{currentstroke}%
\pgfsetstrokeopacity{0.000000}%
\pgfsetdash{}{0pt}%
\pgfpathmoveto{\pgfqpoint{3.107576in}{0.744625in}}%
\pgfpathlineto{\pgfqpoint{3.157750in}{0.744625in}}%
\pgfpathlineto{\pgfqpoint{3.157750in}{0.840875in}}%
\pgfpathlineto{\pgfqpoint{3.107576in}{0.840875in}}%
\pgfpathlineto{\pgfqpoint{3.107576in}{0.744625in}}%
\pgfpathclose%
\pgfusepath{fill}%
\end{pgfscope}%
\begin{pgfscope}%
\pgfpathrectangle{\pgfqpoint{2.857045in}{0.350000in}}{\pgfqpoint{0.300705in}{0.962500in}}%
\pgfusepath{clip}%
\pgfsetbuttcap%
\pgfsetmiterjoin%
\definecolor{currentfill}{rgb}{0.704421,0.881738,0.678970}%
\pgfsetfillcolor{currentfill}%
\pgfsetlinewidth{0.000000pt}%
\definecolor{currentstroke}{rgb}{0.000000,0.000000,0.000000}%
\pgfsetstrokecolor{currentstroke}%
\pgfsetstrokeopacity{0.000000}%
\pgfsetdash{}{0pt}%
\pgfpathmoveto{\pgfqpoint{3.107575in}{0.754250in}}%
\pgfpathlineto{\pgfqpoint{3.157750in}{0.754250in}}%
\pgfpathlineto{\pgfqpoint{3.157750in}{0.850500in}}%
\pgfpathlineto{\pgfqpoint{3.107575in}{0.850500in}}%
\pgfpathlineto{\pgfqpoint{3.107575in}{0.754250in}}%
\pgfpathclose%
\pgfusepath{fill}%
\end{pgfscope}%
\begin{pgfscope}%
\pgfpathrectangle{\pgfqpoint{2.857045in}{0.350000in}}{\pgfqpoint{0.300705in}{0.962500in}}%
\pgfusepath{clip}%
\pgfsetbuttcap%
\pgfsetmiterjoin%
\definecolor{currentfill}{rgb}{0.704421,0.881738,0.678970}%
\pgfsetfillcolor{currentfill}%
\pgfsetlinewidth{0.000000pt}%
\definecolor{currentstroke}{rgb}{0.000000,0.000000,0.000000}%
\pgfsetstrokecolor{currentstroke}%
\pgfsetstrokeopacity{0.000000}%
\pgfsetdash{}{0pt}%
\pgfpathmoveto{\pgfqpoint{3.107374in}{0.763875in}}%
\pgfpathlineto{\pgfqpoint{3.157750in}{0.763875in}}%
\pgfpathlineto{\pgfqpoint{3.157750in}{0.860125in}}%
\pgfpathlineto{\pgfqpoint{3.107374in}{0.860125in}}%
\pgfpathlineto{\pgfqpoint{3.107374in}{0.763875in}}%
\pgfpathclose%
\pgfusepath{fill}%
\end{pgfscope}%
\begin{pgfscope}%
\pgfpathrectangle{\pgfqpoint{2.857045in}{0.350000in}}{\pgfqpoint{0.300705in}{0.962500in}}%
\pgfusepath{clip}%
\pgfsetbuttcap%
\pgfsetmiterjoin%
\definecolor{currentfill}{rgb}{0.699746,0.879769,0.674418}%
\pgfsetfillcolor{currentfill}%
\pgfsetlinewidth{0.000000pt}%
\definecolor{currentstroke}{rgb}{0.000000,0.000000,0.000000}%
\pgfsetstrokecolor{currentstroke}%
\pgfsetstrokeopacity{0.000000}%
\pgfsetdash{}{0pt}%
\pgfpathmoveto{\pgfqpoint{3.107174in}{0.773500in}}%
\pgfpathlineto{\pgfqpoint{3.157750in}{0.773500in}}%
\pgfpathlineto{\pgfqpoint{3.157750in}{0.869750in}}%
\pgfpathlineto{\pgfqpoint{3.107174in}{0.869750in}}%
\pgfpathlineto{\pgfqpoint{3.107174in}{0.773500in}}%
\pgfpathclose%
\pgfusepath{fill}%
\end{pgfscope}%
\begin{pgfscope}%
\pgfpathrectangle{\pgfqpoint{2.857045in}{0.350000in}}{\pgfqpoint{0.300705in}{0.962500in}}%
\pgfusepath{clip}%
\pgfsetbuttcap%
\pgfsetmiterjoin%
\definecolor{currentfill}{rgb}{0.699746,0.879769,0.674418}%
\pgfsetfillcolor{currentfill}%
\pgfsetlinewidth{0.000000pt}%
\definecolor{currentstroke}{rgb}{0.000000,0.000000,0.000000}%
\pgfsetstrokecolor{currentstroke}%
\pgfsetstrokeopacity{0.000000}%
\pgfsetdash{}{0pt}%
\pgfpathmoveto{\pgfqpoint{3.107171in}{0.783125in}}%
\pgfpathlineto{\pgfqpoint{3.157750in}{0.783125in}}%
\pgfpathlineto{\pgfqpoint{3.157750in}{0.879375in}}%
\pgfpathlineto{\pgfqpoint{3.107171in}{0.879375in}}%
\pgfpathlineto{\pgfqpoint{3.107171in}{0.783125in}}%
\pgfpathclose%
\pgfusepath{fill}%
\end{pgfscope}%
\begin{pgfscope}%
\pgfpathrectangle{\pgfqpoint{2.857045in}{0.350000in}}{\pgfqpoint{0.300705in}{0.962500in}}%
\pgfusepath{clip}%
\pgfsetbuttcap%
\pgfsetmiterjoin%
\definecolor{currentfill}{rgb}{0.695071,0.877801,0.669865}%
\pgfsetfillcolor{currentfill}%
\pgfsetlinewidth{0.000000pt}%
\definecolor{currentstroke}{rgb}{0.000000,0.000000,0.000000}%
\pgfsetstrokecolor{currentstroke}%
\pgfsetstrokeopacity{0.000000}%
\pgfsetdash{}{0pt}%
\pgfpathmoveto{\pgfqpoint{3.106981in}{0.792750in}}%
\pgfpathlineto{\pgfqpoint{3.157750in}{0.792750in}}%
\pgfpathlineto{\pgfqpoint{3.157750in}{0.889000in}}%
\pgfpathlineto{\pgfqpoint{3.106981in}{0.889000in}}%
\pgfpathlineto{\pgfqpoint{3.106981in}{0.792750in}}%
\pgfpathclose%
\pgfusepath{fill}%
\end{pgfscope}%
\begin{pgfscope}%
\pgfpathrectangle{\pgfqpoint{2.857045in}{0.350000in}}{\pgfqpoint{0.300705in}{0.962500in}}%
\pgfusepath{clip}%
\pgfsetbuttcap%
\pgfsetmiterjoin%
\definecolor{currentfill}{rgb}{0.695071,0.877801,0.669865}%
\pgfsetfillcolor{currentfill}%
\pgfsetlinewidth{0.000000pt}%
\definecolor{currentstroke}{rgb}{0.000000,0.000000,0.000000}%
\pgfsetstrokecolor{currentstroke}%
\pgfsetstrokeopacity{0.000000}%
\pgfsetdash{}{0pt}%
\pgfpathmoveto{\pgfqpoint{3.106766in}{0.802375in}}%
\pgfpathlineto{\pgfqpoint{3.157750in}{0.802375in}}%
\pgfpathlineto{\pgfqpoint{3.157750in}{0.898625in}}%
\pgfpathlineto{\pgfqpoint{3.106766in}{0.898625in}}%
\pgfpathlineto{\pgfqpoint{3.106766in}{0.802375in}}%
\pgfpathclose%
\pgfusepath{fill}%
\end{pgfscope}%
\begin{pgfscope}%
\pgfpathrectangle{\pgfqpoint{2.857045in}{0.350000in}}{\pgfqpoint{0.300705in}{0.962500in}}%
\pgfusepath{clip}%
\pgfsetbuttcap%
\pgfsetmiterjoin%
\definecolor{currentfill}{rgb}{0.681046,0.871895,0.656209}%
\pgfsetfillcolor{currentfill}%
\pgfsetlinewidth{0.000000pt}%
\definecolor{currentstroke}{rgb}{0.000000,0.000000,0.000000}%
\pgfsetstrokecolor{currentstroke}%
\pgfsetstrokeopacity{0.000000}%
\pgfsetdash{}{0pt}%
\pgfpathmoveto{\pgfqpoint{3.105989in}{0.812000in}}%
\pgfpathlineto{\pgfqpoint{3.157750in}{0.812000in}}%
\pgfpathlineto{\pgfqpoint{3.157750in}{0.908250in}}%
\pgfpathlineto{\pgfqpoint{3.105989in}{0.908250in}}%
\pgfpathlineto{\pgfqpoint{3.105989in}{0.812000in}}%
\pgfpathclose%
\pgfusepath{fill}%
\end{pgfscope}%
\begin{pgfscope}%
\pgfpathrectangle{\pgfqpoint{2.857045in}{0.350000in}}{\pgfqpoint{0.300705in}{0.962500in}}%
\pgfusepath{clip}%
\pgfsetbuttcap%
\pgfsetmiterjoin%
\definecolor{currentfill}{rgb}{0.681046,0.871895,0.656209}%
\pgfsetfillcolor{currentfill}%
\pgfsetlinewidth{0.000000pt}%
\definecolor{currentstroke}{rgb}{0.000000,0.000000,0.000000}%
\pgfsetstrokecolor{currentstroke}%
\pgfsetstrokeopacity{0.000000}%
\pgfsetdash{}{0pt}%
\pgfpathmoveto{\pgfqpoint{3.105833in}{0.821625in}}%
\pgfpathlineto{\pgfqpoint{3.157750in}{0.821625in}}%
\pgfpathlineto{\pgfqpoint{3.157750in}{0.917875in}}%
\pgfpathlineto{\pgfqpoint{3.105833in}{0.917875in}}%
\pgfpathlineto{\pgfqpoint{3.105833in}{0.821625in}}%
\pgfpathclose%
\pgfusepath{fill}%
\end{pgfscope}%
\begin{pgfscope}%
\pgfpathrectangle{\pgfqpoint{2.857045in}{0.350000in}}{\pgfqpoint{0.300705in}{0.962500in}}%
\pgfusepath{clip}%
\pgfsetbuttcap%
\pgfsetmiterjoin%
\definecolor{currentfill}{rgb}{0.681046,0.871895,0.656209}%
\pgfsetfillcolor{currentfill}%
\pgfsetlinewidth{0.000000pt}%
\definecolor{currentstroke}{rgb}{0.000000,0.000000,0.000000}%
\pgfsetstrokecolor{currentstroke}%
\pgfsetstrokeopacity{0.000000}%
\pgfsetdash{}{0pt}%
\pgfpathmoveto{\pgfqpoint{3.105833in}{0.831250in}}%
\pgfpathlineto{\pgfqpoint{3.157750in}{0.831250in}}%
\pgfpathlineto{\pgfqpoint{3.157750in}{0.927500in}}%
\pgfpathlineto{\pgfqpoint{3.105833in}{0.927500in}}%
\pgfpathlineto{\pgfqpoint{3.105833in}{0.831250in}}%
\pgfpathclose%
\pgfusepath{fill}%
\end{pgfscope}%
\begin{pgfscope}%
\pgfpathrectangle{\pgfqpoint{2.857045in}{0.350000in}}{\pgfqpoint{0.300705in}{0.962500in}}%
\pgfusepath{clip}%
\pgfsetbuttcap%
\pgfsetmiterjoin%
\definecolor{currentfill}{rgb}{0.681046,0.871895,0.656209}%
\pgfsetfillcolor{currentfill}%
\pgfsetlinewidth{0.000000pt}%
\definecolor{currentstroke}{rgb}{0.000000,0.000000,0.000000}%
\pgfsetstrokecolor{currentstroke}%
\pgfsetstrokeopacity{0.000000}%
\pgfsetdash{}{0pt}%
\pgfpathmoveto{\pgfqpoint{3.105833in}{0.840875in}}%
\pgfpathlineto{\pgfqpoint{3.157750in}{0.840875in}}%
\pgfpathlineto{\pgfqpoint{3.157750in}{0.937125in}}%
\pgfpathlineto{\pgfqpoint{3.105833in}{0.937125in}}%
\pgfpathlineto{\pgfqpoint{3.105833in}{0.840875in}}%
\pgfpathclose%
\pgfusepath{fill}%
\end{pgfscope}%
\begin{pgfscope}%
\pgfpathrectangle{\pgfqpoint{2.857045in}{0.350000in}}{\pgfqpoint{0.300705in}{0.962500in}}%
\pgfusepath{clip}%
\pgfsetbuttcap%
\pgfsetmiterjoin%
\definecolor{currentfill}{rgb}{0.676371,0.869927,0.651657}%
\pgfsetfillcolor{currentfill}%
\pgfsetlinewidth{0.000000pt}%
\definecolor{currentstroke}{rgb}{0.000000,0.000000,0.000000}%
\pgfsetstrokecolor{currentstroke}%
\pgfsetstrokeopacity{0.000000}%
\pgfsetdash{}{0pt}%
\pgfpathmoveto{\pgfqpoint{3.105726in}{0.850500in}}%
\pgfpathlineto{\pgfqpoint{3.157750in}{0.850500in}}%
\pgfpathlineto{\pgfqpoint{3.157750in}{0.946750in}}%
\pgfpathlineto{\pgfqpoint{3.105726in}{0.946750in}}%
\pgfpathlineto{\pgfqpoint{3.105726in}{0.850500in}}%
\pgfpathclose%
\pgfusepath{fill}%
\end{pgfscope}%
\begin{pgfscope}%
\pgfpathrectangle{\pgfqpoint{2.857045in}{0.350000in}}{\pgfqpoint{0.300705in}{0.962500in}}%
\pgfusepath{clip}%
\pgfsetbuttcap%
\pgfsetmiterjoin%
\definecolor{currentfill}{rgb}{0.676371,0.869927,0.651657}%
\pgfsetfillcolor{currentfill}%
\pgfsetlinewidth{0.000000pt}%
\definecolor{currentstroke}{rgb}{0.000000,0.000000,0.000000}%
\pgfsetstrokecolor{currentstroke}%
\pgfsetstrokeopacity{0.000000}%
\pgfsetdash{}{0pt}%
\pgfpathmoveto{\pgfqpoint{3.105671in}{0.860125in}}%
\pgfpathlineto{\pgfqpoint{3.157750in}{0.860125in}}%
\pgfpathlineto{\pgfqpoint{3.157750in}{0.956375in}}%
\pgfpathlineto{\pgfqpoint{3.105671in}{0.956375in}}%
\pgfpathlineto{\pgfqpoint{3.105671in}{0.860125in}}%
\pgfpathclose%
\pgfusepath{fill}%
\end{pgfscope}%
\begin{pgfscope}%
\pgfpathrectangle{\pgfqpoint{2.857045in}{0.350000in}}{\pgfqpoint{0.300705in}{0.962500in}}%
\pgfusepath{clip}%
\pgfsetbuttcap%
\pgfsetmiterjoin%
\definecolor{currentfill}{rgb}{0.676371,0.869927,0.651657}%
\pgfsetfillcolor{currentfill}%
\pgfsetlinewidth{0.000000pt}%
\definecolor{currentstroke}{rgb}{0.000000,0.000000,0.000000}%
\pgfsetstrokecolor{currentstroke}%
\pgfsetstrokeopacity{0.000000}%
\pgfsetdash{}{0pt}%
\pgfpathmoveto{\pgfqpoint{3.105478in}{0.869750in}}%
\pgfpathlineto{\pgfqpoint{3.157750in}{0.869750in}}%
\pgfpathlineto{\pgfqpoint{3.157750in}{0.966000in}}%
\pgfpathlineto{\pgfqpoint{3.105478in}{0.966000in}}%
\pgfpathlineto{\pgfqpoint{3.105478in}{0.869750in}}%
\pgfpathclose%
\pgfusepath{fill}%
\end{pgfscope}%
\begin{pgfscope}%
\pgfpathrectangle{\pgfqpoint{2.857045in}{0.350000in}}{\pgfqpoint{0.300705in}{0.962500in}}%
\pgfusepath{clip}%
\pgfsetbuttcap%
\pgfsetmiterjoin%
\definecolor{currentfill}{rgb}{0.671696,0.867958,0.647105}%
\pgfsetfillcolor{currentfill}%
\pgfsetlinewidth{0.000000pt}%
\definecolor{currentstroke}{rgb}{0.000000,0.000000,0.000000}%
\pgfsetstrokecolor{currentstroke}%
\pgfsetstrokeopacity{0.000000}%
\pgfsetdash{}{0pt}%
\pgfpathmoveto{\pgfqpoint{3.105429in}{0.879375in}}%
\pgfpathlineto{\pgfqpoint{3.157750in}{0.879375in}}%
\pgfpathlineto{\pgfqpoint{3.157750in}{0.975625in}}%
\pgfpathlineto{\pgfqpoint{3.105429in}{0.975625in}}%
\pgfpathlineto{\pgfqpoint{3.105429in}{0.879375in}}%
\pgfpathclose%
\pgfusepath{fill}%
\end{pgfscope}%
\begin{pgfscope}%
\pgfpathrectangle{\pgfqpoint{2.857045in}{0.350000in}}{\pgfqpoint{0.300705in}{0.962500in}}%
\pgfusepath{clip}%
\pgfsetbuttcap%
\pgfsetmiterjoin%
\definecolor{currentfill}{rgb}{0.662345,0.864022,0.638001}%
\pgfsetfillcolor{currentfill}%
\pgfsetlinewidth{0.000000pt}%
\definecolor{currentstroke}{rgb}{0.000000,0.000000,0.000000}%
\pgfsetstrokecolor{currentstroke}%
\pgfsetstrokeopacity{0.000000}%
\pgfsetdash{}{0pt}%
\pgfpathmoveto{\pgfqpoint{3.104769in}{0.889000in}}%
\pgfpathlineto{\pgfqpoint{3.157750in}{0.889000in}}%
\pgfpathlineto{\pgfqpoint{3.157750in}{0.985250in}}%
\pgfpathlineto{\pgfqpoint{3.104769in}{0.985250in}}%
\pgfpathlineto{\pgfqpoint{3.104769in}{0.889000in}}%
\pgfpathclose%
\pgfusepath{fill}%
\end{pgfscope}%
\begin{pgfscope}%
\pgfpathrectangle{\pgfqpoint{2.857045in}{0.350000in}}{\pgfqpoint{0.300705in}{0.962500in}}%
\pgfusepath{clip}%
\pgfsetbuttcap%
\pgfsetmiterjoin%
\definecolor{currentfill}{rgb}{0.662345,0.864022,0.638001}%
\pgfsetfillcolor{currentfill}%
\pgfsetlinewidth{0.000000pt}%
\definecolor{currentstroke}{rgb}{0.000000,0.000000,0.000000}%
\pgfsetstrokecolor{currentstroke}%
\pgfsetstrokeopacity{0.000000}%
\pgfsetdash{}{0pt}%
\pgfpathmoveto{\pgfqpoint{3.104769in}{0.898625in}}%
\pgfpathlineto{\pgfqpoint{3.157750in}{0.898625in}}%
\pgfpathlineto{\pgfqpoint{3.157750in}{0.994875in}}%
\pgfpathlineto{\pgfqpoint{3.104769in}{0.994875in}}%
\pgfpathlineto{\pgfqpoint{3.104769in}{0.898625in}}%
\pgfpathclose%
\pgfusepath{fill}%
\end{pgfscope}%
\begin{pgfscope}%
\pgfpathrectangle{\pgfqpoint{2.857045in}{0.350000in}}{\pgfqpoint{0.300705in}{0.962500in}}%
\pgfusepath{clip}%
\pgfsetbuttcap%
\pgfsetmiterjoin%
\definecolor{currentfill}{rgb}{0.662345,0.864022,0.638001}%
\pgfsetfillcolor{currentfill}%
\pgfsetlinewidth{0.000000pt}%
\definecolor{currentstroke}{rgb}{0.000000,0.000000,0.000000}%
\pgfsetstrokecolor{currentstroke}%
\pgfsetstrokeopacity{0.000000}%
\pgfsetdash{}{0pt}%
\pgfpathmoveto{\pgfqpoint{3.104753in}{0.908250in}}%
\pgfpathlineto{\pgfqpoint{3.157750in}{0.908250in}}%
\pgfpathlineto{\pgfqpoint{3.157750in}{1.004500in}}%
\pgfpathlineto{\pgfqpoint{3.104753in}{1.004500in}}%
\pgfpathlineto{\pgfqpoint{3.104753in}{0.908250in}}%
\pgfpathclose%
\pgfusepath{fill}%
\end{pgfscope}%
\begin{pgfscope}%
\pgfpathrectangle{\pgfqpoint{2.857045in}{0.350000in}}{\pgfqpoint{0.300705in}{0.962500in}}%
\pgfusepath{clip}%
\pgfsetbuttcap%
\pgfsetmiterjoin%
\definecolor{currentfill}{rgb}{0.662345,0.864022,0.638001}%
\pgfsetfillcolor{currentfill}%
\pgfsetlinewidth{0.000000pt}%
\definecolor{currentstroke}{rgb}{0.000000,0.000000,0.000000}%
\pgfsetstrokecolor{currentstroke}%
\pgfsetstrokeopacity{0.000000}%
\pgfsetdash{}{0pt}%
\pgfpathmoveto{\pgfqpoint{3.104753in}{0.917875in}}%
\pgfpathlineto{\pgfqpoint{3.157750in}{0.917875in}}%
\pgfpathlineto{\pgfqpoint{3.157750in}{1.014125in}}%
\pgfpathlineto{\pgfqpoint{3.104753in}{1.014125in}}%
\pgfpathlineto{\pgfqpoint{3.104753in}{0.917875in}}%
\pgfpathclose%
\pgfusepath{fill}%
\end{pgfscope}%
\begin{pgfscope}%
\pgfpathrectangle{\pgfqpoint{2.857045in}{0.350000in}}{\pgfqpoint{0.300705in}{0.962500in}}%
\pgfusepath{clip}%
\pgfsetbuttcap%
\pgfsetmiterjoin%
\definecolor{currentfill}{rgb}{0.662345,0.864022,0.638001}%
\pgfsetfillcolor{currentfill}%
\pgfsetlinewidth{0.000000pt}%
\definecolor{currentstroke}{rgb}{0.000000,0.000000,0.000000}%
\pgfsetstrokecolor{currentstroke}%
\pgfsetstrokeopacity{0.000000}%
\pgfsetdash{}{0pt}%
\pgfpathmoveto{\pgfqpoint{3.104715in}{0.927500in}}%
\pgfpathlineto{\pgfqpoint{3.157750in}{0.927500in}}%
\pgfpathlineto{\pgfqpoint{3.157750in}{1.023750in}}%
\pgfpathlineto{\pgfqpoint{3.104715in}{1.023750in}}%
\pgfpathlineto{\pgfqpoint{3.104715in}{0.927500in}}%
\pgfpathclose%
\pgfusepath{fill}%
\end{pgfscope}%
\begin{pgfscope}%
\pgfpathrectangle{\pgfqpoint{2.857045in}{0.350000in}}{\pgfqpoint{0.300705in}{0.962500in}}%
\pgfusepath{clip}%
\pgfsetbuttcap%
\pgfsetmiterjoin%
\definecolor{currentfill}{rgb}{0.662345,0.864022,0.638001}%
\pgfsetfillcolor{currentfill}%
\pgfsetlinewidth{0.000000pt}%
\definecolor{currentstroke}{rgb}{0.000000,0.000000,0.000000}%
\pgfsetstrokecolor{currentstroke}%
\pgfsetstrokeopacity{0.000000}%
\pgfsetdash{}{0pt}%
\pgfpathmoveto{\pgfqpoint{3.104610in}{0.937125in}}%
\pgfpathlineto{\pgfqpoint{3.157750in}{0.937125in}}%
\pgfpathlineto{\pgfqpoint{3.157750in}{1.033375in}}%
\pgfpathlineto{\pgfqpoint{3.104610in}{1.033375in}}%
\pgfpathlineto{\pgfqpoint{3.104610in}{0.937125in}}%
\pgfpathclose%
\pgfusepath{fill}%
\end{pgfscope}%
\begin{pgfscope}%
\pgfpathrectangle{\pgfqpoint{2.857045in}{0.350000in}}{\pgfqpoint{0.300705in}{0.962500in}}%
\pgfusepath{clip}%
\pgfsetbuttcap%
\pgfsetmiterjoin%
\definecolor{currentfill}{rgb}{0.657670,0.862053,0.633449}%
\pgfsetfillcolor{currentfill}%
\pgfsetlinewidth{0.000000pt}%
\definecolor{currentstroke}{rgb}{0.000000,0.000000,0.000000}%
\pgfsetstrokecolor{currentstroke}%
\pgfsetstrokeopacity{0.000000}%
\pgfsetdash{}{0pt}%
\pgfpathmoveto{\pgfqpoint{3.104496in}{0.946750in}}%
\pgfpathlineto{\pgfqpoint{3.157750in}{0.946750in}}%
\pgfpathlineto{\pgfqpoint{3.157750in}{1.043000in}}%
\pgfpathlineto{\pgfqpoint{3.104496in}{1.043000in}}%
\pgfpathlineto{\pgfqpoint{3.104496in}{0.946750in}}%
\pgfpathclose%
\pgfusepath{fill}%
\end{pgfscope}%
\begin{pgfscope}%
\pgfpathrectangle{\pgfqpoint{2.857045in}{0.350000in}}{\pgfqpoint{0.300705in}{0.962500in}}%
\pgfusepath{clip}%
\pgfsetbuttcap%
\pgfsetmiterjoin%
\definecolor{currentfill}{rgb}{0.657670,0.862053,0.633449}%
\pgfsetfillcolor{currentfill}%
\pgfsetlinewidth{0.000000pt}%
\definecolor{currentstroke}{rgb}{0.000000,0.000000,0.000000}%
\pgfsetstrokecolor{currentstroke}%
\pgfsetstrokeopacity{0.000000}%
\pgfsetdash{}{0pt}%
\pgfpathmoveto{\pgfqpoint{3.104266in}{0.956375in}}%
\pgfpathlineto{\pgfqpoint{3.157750in}{0.956375in}}%
\pgfpathlineto{\pgfqpoint{3.157750in}{1.052625in}}%
\pgfpathlineto{\pgfqpoint{3.104266in}{1.052625in}}%
\pgfpathlineto{\pgfqpoint{3.104266in}{0.956375in}}%
\pgfpathclose%
\pgfusepath{fill}%
\end{pgfscope}%
\begin{pgfscope}%
\pgfpathrectangle{\pgfqpoint{2.857045in}{0.350000in}}{\pgfqpoint{0.300705in}{0.962500in}}%
\pgfusepath{clip}%
\pgfsetbuttcap%
\pgfsetmiterjoin%
\definecolor{currentfill}{rgb}{0.657670,0.862053,0.633449}%
\pgfsetfillcolor{currentfill}%
\pgfsetlinewidth{0.000000pt}%
\definecolor{currentstroke}{rgb}{0.000000,0.000000,0.000000}%
\pgfsetstrokecolor{currentstroke}%
\pgfsetstrokeopacity{0.000000}%
\pgfsetdash{}{0pt}%
\pgfpathmoveto{\pgfqpoint{3.104238in}{0.966000in}}%
\pgfpathlineto{\pgfqpoint{3.157750in}{0.966000in}}%
\pgfpathlineto{\pgfqpoint{3.157750in}{1.062250in}}%
\pgfpathlineto{\pgfqpoint{3.104238in}{1.062250in}}%
\pgfpathlineto{\pgfqpoint{3.104238in}{0.966000in}}%
\pgfpathclose%
\pgfusepath{fill}%
\end{pgfscope}%
\begin{pgfscope}%
\pgfpathrectangle{\pgfqpoint{2.857045in}{0.350000in}}{\pgfqpoint{0.300705in}{0.962500in}}%
\pgfusepath{clip}%
\pgfsetbuttcap%
\pgfsetmiterjoin%
\definecolor{currentfill}{rgb}{0.652995,0.860085,0.628897}%
\pgfsetfillcolor{currentfill}%
\pgfsetlinewidth{0.000000pt}%
\definecolor{currentstroke}{rgb}{0.000000,0.000000,0.000000}%
\pgfsetstrokecolor{currentstroke}%
\pgfsetstrokeopacity{0.000000}%
\pgfsetdash{}{0pt}%
\pgfpathmoveto{\pgfqpoint{3.104137in}{0.975625in}}%
\pgfpathlineto{\pgfqpoint{3.157750in}{0.975625in}}%
\pgfpathlineto{\pgfqpoint{3.157750in}{1.071875in}}%
\pgfpathlineto{\pgfqpoint{3.104137in}{1.071875in}}%
\pgfpathlineto{\pgfqpoint{3.104137in}{0.975625in}}%
\pgfpathclose%
\pgfusepath{fill}%
\end{pgfscope}%
\begin{pgfscope}%
\pgfpathrectangle{\pgfqpoint{2.857045in}{0.350000in}}{\pgfqpoint{0.300705in}{0.962500in}}%
\pgfusepath{clip}%
\pgfsetbuttcap%
\pgfsetmiterjoin%
\definecolor{currentfill}{rgb}{0.648320,0.858116,0.624344}%
\pgfsetfillcolor{currentfill}%
\pgfsetlinewidth{0.000000pt}%
\definecolor{currentstroke}{rgb}{0.000000,0.000000,0.000000}%
\pgfsetstrokecolor{currentstroke}%
\pgfsetstrokeopacity{0.000000}%
\pgfsetdash{}{0pt}%
\pgfpathmoveto{\pgfqpoint{3.103800in}{0.985250in}}%
\pgfpathlineto{\pgfqpoint{3.157750in}{0.985250in}}%
\pgfpathlineto{\pgfqpoint{3.157750in}{1.081500in}}%
\pgfpathlineto{\pgfqpoint{3.103800in}{1.081500in}}%
\pgfpathlineto{\pgfqpoint{3.103800in}{0.985250in}}%
\pgfpathclose%
\pgfusepath{fill}%
\end{pgfscope}%
\begin{pgfscope}%
\pgfpathrectangle{\pgfqpoint{2.857045in}{0.350000in}}{\pgfqpoint{0.300705in}{0.962500in}}%
\pgfusepath{clip}%
\pgfsetbuttcap%
\pgfsetmiterjoin%
\definecolor{currentfill}{rgb}{0.643645,0.856148,0.619792}%
\pgfsetfillcolor{currentfill}%
\pgfsetlinewidth{0.000000pt}%
\definecolor{currentstroke}{rgb}{0.000000,0.000000,0.000000}%
\pgfsetstrokecolor{currentstroke}%
\pgfsetstrokeopacity{0.000000}%
\pgfsetdash{}{0pt}%
\pgfpathmoveto{\pgfqpoint{3.103444in}{0.994875in}}%
\pgfpathlineto{\pgfqpoint{3.157750in}{0.994875in}}%
\pgfpathlineto{\pgfqpoint{3.157750in}{1.091125in}}%
\pgfpathlineto{\pgfqpoint{3.103444in}{1.091125in}}%
\pgfpathlineto{\pgfqpoint{3.103444in}{0.994875in}}%
\pgfpathclose%
\pgfusepath{fill}%
\end{pgfscope}%
\begin{pgfscope}%
\pgfpathrectangle{\pgfqpoint{2.857045in}{0.350000in}}{\pgfqpoint{0.300705in}{0.962500in}}%
\pgfusepath{clip}%
\pgfsetbuttcap%
\pgfsetmiterjoin%
\definecolor{currentfill}{rgb}{0.643645,0.856148,0.619792}%
\pgfsetfillcolor{currentfill}%
\pgfsetlinewidth{0.000000pt}%
\definecolor{currentstroke}{rgb}{0.000000,0.000000,0.000000}%
\pgfsetstrokecolor{currentstroke}%
\pgfsetstrokeopacity{0.000000}%
\pgfsetdash{}{0pt}%
\pgfpathmoveto{\pgfqpoint{3.103388in}{1.004500in}}%
\pgfpathlineto{\pgfqpoint{3.157750in}{1.004500in}}%
\pgfpathlineto{\pgfqpoint{3.157750in}{1.100750in}}%
\pgfpathlineto{\pgfqpoint{3.103388in}{1.100750in}}%
\pgfpathlineto{\pgfqpoint{3.103388in}{1.004500in}}%
\pgfpathclose%
\pgfusepath{fill}%
\end{pgfscope}%
\begin{pgfscope}%
\pgfpathrectangle{\pgfqpoint{2.857045in}{0.350000in}}{\pgfqpoint{0.300705in}{0.962500in}}%
\pgfusepath{clip}%
\pgfsetbuttcap%
\pgfsetmiterjoin%
\definecolor{currentfill}{rgb}{0.638970,0.854179,0.615240}%
\pgfsetfillcolor{currentfill}%
\pgfsetlinewidth{0.000000pt}%
\definecolor{currentstroke}{rgb}{0.000000,0.000000,0.000000}%
\pgfsetstrokecolor{currentstroke}%
\pgfsetstrokeopacity{0.000000}%
\pgfsetdash{}{0pt}%
\pgfpathmoveto{\pgfqpoint{3.102963in}{1.014125in}}%
\pgfpathlineto{\pgfqpoint{3.157750in}{1.014125in}}%
\pgfpathlineto{\pgfqpoint{3.157750in}{1.110375in}}%
\pgfpathlineto{\pgfqpoint{3.102963in}{1.110375in}}%
\pgfpathlineto{\pgfqpoint{3.102963in}{1.014125in}}%
\pgfpathclose%
\pgfusepath{fill}%
\end{pgfscope}%
\begin{pgfscope}%
\pgfpathrectangle{\pgfqpoint{2.857045in}{0.350000in}}{\pgfqpoint{0.300705in}{0.962500in}}%
\pgfusepath{clip}%
\pgfsetbuttcap%
\pgfsetmiterjoin%
\definecolor{currentfill}{rgb}{0.634295,0.852211,0.610688}%
\pgfsetfillcolor{currentfill}%
\pgfsetlinewidth{0.000000pt}%
\definecolor{currentstroke}{rgb}{0.000000,0.000000,0.000000}%
\pgfsetstrokecolor{currentstroke}%
\pgfsetstrokeopacity{0.000000}%
\pgfsetdash{}{0pt}%
\pgfpathmoveto{\pgfqpoint{3.102650in}{1.023750in}}%
\pgfpathlineto{\pgfqpoint{3.157750in}{1.023750in}}%
\pgfpathlineto{\pgfqpoint{3.157750in}{1.120000in}}%
\pgfpathlineto{\pgfqpoint{3.102650in}{1.120000in}}%
\pgfpathlineto{\pgfqpoint{3.102650in}{1.023750in}}%
\pgfpathclose%
\pgfusepath{fill}%
\end{pgfscope}%
\begin{pgfscope}%
\pgfpathrectangle{\pgfqpoint{2.857045in}{0.350000in}}{\pgfqpoint{0.300705in}{0.962500in}}%
\pgfusepath{clip}%
\pgfsetbuttcap%
\pgfsetmiterjoin%
\definecolor{currentfill}{rgb}{0.629296,0.850012,0.606136}%
\pgfsetfillcolor{currentfill}%
\pgfsetlinewidth{0.000000pt}%
\definecolor{currentstroke}{rgb}{0.000000,0.000000,0.000000}%
\pgfsetstrokecolor{currentstroke}%
\pgfsetstrokeopacity{0.000000}%
\pgfsetdash{}{0pt}%
\pgfpathmoveto{\pgfqpoint{3.102609in}{1.033375in}}%
\pgfpathlineto{\pgfqpoint{3.157750in}{1.033375in}}%
\pgfpathlineto{\pgfqpoint{3.157750in}{1.129625in}}%
\pgfpathlineto{\pgfqpoint{3.102609in}{1.129625in}}%
\pgfpathlineto{\pgfqpoint{3.102609in}{1.033375in}}%
\pgfpathclose%
\pgfusepath{fill}%
\end{pgfscope}%
\begin{pgfscope}%
\pgfpathrectangle{\pgfqpoint{2.857045in}{0.350000in}}{\pgfqpoint{0.300705in}{0.962500in}}%
\pgfusepath{clip}%
\pgfsetbuttcap%
\pgfsetmiterjoin%
\definecolor{currentfill}{rgb}{0.623760,0.847428,0.601584}%
\pgfsetfillcolor{currentfill}%
\pgfsetlinewidth{0.000000pt}%
\definecolor{currentstroke}{rgb}{0.000000,0.000000,0.000000}%
\pgfsetstrokecolor{currentstroke}%
\pgfsetstrokeopacity{0.000000}%
\pgfsetdash{}{0pt}%
\pgfpathmoveto{\pgfqpoint{3.102119in}{1.043000in}}%
\pgfpathlineto{\pgfqpoint{3.157750in}{1.043000in}}%
\pgfpathlineto{\pgfqpoint{3.157750in}{1.139250in}}%
\pgfpathlineto{\pgfqpoint{3.102119in}{1.139250in}}%
\pgfpathlineto{\pgfqpoint{3.102119in}{1.043000in}}%
\pgfpathclose%
\pgfusepath{fill}%
\end{pgfscope}%
\begin{pgfscope}%
\pgfpathrectangle{\pgfqpoint{2.857045in}{0.350000in}}{\pgfqpoint{0.300705in}{0.962500in}}%
\pgfusepath{clip}%
\pgfsetbuttcap%
\pgfsetmiterjoin%
\definecolor{currentfill}{rgb}{0.618224,0.844844,0.597032}%
\pgfsetfillcolor{currentfill}%
\pgfsetlinewidth{0.000000pt}%
\definecolor{currentstroke}{rgb}{0.000000,0.000000,0.000000}%
\pgfsetstrokecolor{currentstroke}%
\pgfsetstrokeopacity{0.000000}%
\pgfsetdash{}{0pt}%
\pgfpathmoveto{\pgfqpoint{3.101940in}{1.052625in}}%
\pgfpathlineto{\pgfqpoint{3.157750in}{1.052625in}}%
\pgfpathlineto{\pgfqpoint{3.157750in}{1.148875in}}%
\pgfpathlineto{\pgfqpoint{3.101940in}{1.148875in}}%
\pgfpathlineto{\pgfqpoint{3.101940in}{1.052625in}}%
\pgfpathclose%
\pgfusepath{fill}%
\end{pgfscope}%
\begin{pgfscope}%
\pgfpathrectangle{\pgfqpoint{2.857045in}{0.350000in}}{\pgfqpoint{0.300705in}{0.962500in}}%
\pgfusepath{clip}%
\pgfsetbuttcap%
\pgfsetmiterjoin%
\definecolor{currentfill}{rgb}{0.612687,0.842261,0.592480}%
\pgfsetfillcolor{currentfill}%
\pgfsetlinewidth{0.000000pt}%
\definecolor{currentstroke}{rgb}{0.000000,0.000000,0.000000}%
\pgfsetstrokecolor{currentstroke}%
\pgfsetstrokeopacity{0.000000}%
\pgfsetdash{}{0pt}%
\pgfpathmoveto{\pgfqpoint{3.101549in}{1.062250in}}%
\pgfpathlineto{\pgfqpoint{3.157750in}{1.062250in}}%
\pgfpathlineto{\pgfqpoint{3.157750in}{1.158500in}}%
\pgfpathlineto{\pgfqpoint{3.101549in}{1.158500in}}%
\pgfpathlineto{\pgfqpoint{3.101549in}{1.062250in}}%
\pgfpathclose%
\pgfusepath{fill}%
\end{pgfscope}%
\begin{pgfscope}%
\pgfpathrectangle{\pgfqpoint{2.857045in}{0.350000in}}{\pgfqpoint{0.300705in}{0.962500in}}%
\pgfusepath{clip}%
\pgfsetbuttcap%
\pgfsetmiterjoin%
\definecolor{currentfill}{rgb}{0.601615,0.837093,0.583376}%
\pgfsetfillcolor{currentfill}%
\pgfsetlinewidth{0.000000pt}%
\definecolor{currentstroke}{rgb}{0.000000,0.000000,0.000000}%
\pgfsetstrokecolor{currentstroke}%
\pgfsetstrokeopacity{0.000000}%
\pgfsetdash{}{0pt}%
\pgfpathmoveto{\pgfqpoint{3.100977in}{1.071875in}}%
\pgfpathlineto{\pgfqpoint{3.157750in}{1.071875in}}%
\pgfpathlineto{\pgfqpoint{3.157750in}{1.168125in}}%
\pgfpathlineto{\pgfqpoint{3.100977in}{1.168125in}}%
\pgfpathlineto{\pgfqpoint{3.100977in}{1.071875in}}%
\pgfpathclose%
\pgfusepath{fill}%
\end{pgfscope}%
\begin{pgfscope}%
\pgfpathrectangle{\pgfqpoint{2.857045in}{0.350000in}}{\pgfqpoint{0.300705in}{0.962500in}}%
\pgfusepath{clip}%
\pgfsetbuttcap%
\pgfsetmiterjoin%
\definecolor{currentfill}{rgb}{0.590542,0.831926,0.574271}%
\pgfsetfillcolor{currentfill}%
\pgfsetlinewidth{0.000000pt}%
\definecolor{currentstroke}{rgb}{0.000000,0.000000,0.000000}%
\pgfsetstrokecolor{currentstroke}%
\pgfsetstrokeopacity{0.000000}%
\pgfsetdash{}{0pt}%
\pgfpathmoveto{\pgfqpoint{3.100325in}{1.081500in}}%
\pgfpathlineto{\pgfqpoint{3.157750in}{1.081500in}}%
\pgfpathlineto{\pgfqpoint{3.157750in}{1.177750in}}%
\pgfpathlineto{\pgfqpoint{3.100325in}{1.177750in}}%
\pgfpathlineto{\pgfqpoint{3.100325in}{1.081500in}}%
\pgfpathclose%
\pgfusepath{fill}%
\end{pgfscope}%
\begin{pgfscope}%
\pgfpathrectangle{\pgfqpoint{2.857045in}{0.350000in}}{\pgfqpoint{0.300705in}{0.962500in}}%
\pgfusepath{clip}%
\pgfsetbuttcap%
\pgfsetmiterjoin%
\definecolor{currentfill}{rgb}{0.585006,0.829343,0.569719}%
\pgfsetfillcolor{currentfill}%
\pgfsetlinewidth{0.000000pt}%
\definecolor{currentstroke}{rgb}{0.000000,0.000000,0.000000}%
\pgfsetstrokecolor{currentstroke}%
\pgfsetstrokeopacity{0.000000}%
\pgfsetdash{}{0pt}%
\pgfpathmoveto{\pgfqpoint{3.100068in}{1.091125in}}%
\pgfpathlineto{\pgfqpoint{3.157750in}{1.091125in}}%
\pgfpathlineto{\pgfqpoint{3.157750in}{1.187375in}}%
\pgfpathlineto{\pgfqpoint{3.100068in}{1.187375in}}%
\pgfpathlineto{\pgfqpoint{3.100068in}{1.091125in}}%
\pgfpathclose%
\pgfusepath{fill}%
\end{pgfscope}%
\begin{pgfscope}%
\pgfpathrectangle{\pgfqpoint{2.857045in}{0.350000in}}{\pgfqpoint{0.300705in}{0.962500in}}%
\pgfusepath{clip}%
\pgfsetbuttcap%
\pgfsetmiterjoin%
\definecolor{currentfill}{rgb}{0.579469,0.826759,0.565167}%
\pgfsetfillcolor{currentfill}%
\pgfsetlinewidth{0.000000pt}%
\definecolor{currentstroke}{rgb}{0.000000,0.000000,0.000000}%
\pgfsetstrokecolor{currentstroke}%
\pgfsetstrokeopacity{0.000000}%
\pgfsetdash{}{0pt}%
\pgfpathmoveto{\pgfqpoint{3.099540in}{1.100750in}}%
\pgfpathlineto{\pgfqpoint{3.157750in}{1.100750in}}%
\pgfpathlineto{\pgfqpoint{3.157750in}{1.197000in}}%
\pgfpathlineto{\pgfqpoint{3.099540in}{1.197000in}}%
\pgfpathlineto{\pgfqpoint{3.099540in}{1.100750in}}%
\pgfpathclose%
\pgfusepath{fill}%
\end{pgfscope}%
\begin{pgfscope}%
\pgfpathrectangle{\pgfqpoint{2.857045in}{0.350000in}}{\pgfqpoint{0.300705in}{0.962500in}}%
\pgfusepath{clip}%
\pgfsetbuttcap%
\pgfsetmiterjoin%
\definecolor{currentfill}{rgb}{0.573933,0.824175,0.560615}%
\pgfsetfillcolor{currentfill}%
\pgfsetlinewidth{0.000000pt}%
\definecolor{currentstroke}{rgb}{0.000000,0.000000,0.000000}%
\pgfsetstrokecolor{currentstroke}%
\pgfsetstrokeopacity{0.000000}%
\pgfsetdash{}{0pt}%
\pgfpathmoveto{\pgfqpoint{3.099270in}{1.110375in}}%
\pgfpathlineto{\pgfqpoint{3.157750in}{1.110375in}}%
\pgfpathlineto{\pgfqpoint{3.157750in}{1.206625in}}%
\pgfpathlineto{\pgfqpoint{3.099270in}{1.206625in}}%
\pgfpathlineto{\pgfqpoint{3.099270in}{1.110375in}}%
\pgfpathclose%
\pgfusepath{fill}%
\end{pgfscope}%
\begin{pgfscope}%
\pgfpathrectangle{\pgfqpoint{2.857045in}{0.350000in}}{\pgfqpoint{0.300705in}{0.962500in}}%
\pgfusepath{clip}%
\pgfsetbuttcap%
\pgfsetmiterjoin%
\definecolor{currentfill}{rgb}{0.568397,0.821592,0.556063}%
\pgfsetfillcolor{currentfill}%
\pgfsetlinewidth{0.000000pt}%
\definecolor{currentstroke}{rgb}{0.000000,0.000000,0.000000}%
\pgfsetstrokecolor{currentstroke}%
\pgfsetstrokeopacity{0.000000}%
\pgfsetdash{}{0pt}%
\pgfpathmoveto{\pgfqpoint{3.099132in}{1.120000in}}%
\pgfpathlineto{\pgfqpoint{3.157750in}{1.120000in}}%
\pgfpathlineto{\pgfqpoint{3.157750in}{1.216250in}}%
\pgfpathlineto{\pgfqpoint{3.099132in}{1.216250in}}%
\pgfpathlineto{\pgfqpoint{3.099132in}{1.120000in}}%
\pgfpathclose%
\pgfusepath{fill}%
\end{pgfscope}%
\begin{pgfscope}%
\pgfpathrectangle{\pgfqpoint{2.857045in}{0.350000in}}{\pgfqpoint{0.300705in}{0.962500in}}%
\pgfusepath{clip}%
\pgfsetbuttcap%
\pgfsetmiterjoin%
\definecolor{currentfill}{rgb}{0.557324,0.816424,0.546959}%
\pgfsetfillcolor{currentfill}%
\pgfsetlinewidth{0.000000pt}%
\definecolor{currentstroke}{rgb}{0.000000,0.000000,0.000000}%
\pgfsetstrokecolor{currentstroke}%
\pgfsetstrokeopacity{0.000000}%
\pgfsetdash{}{0pt}%
\pgfpathmoveto{\pgfqpoint{3.098413in}{1.129625in}}%
\pgfpathlineto{\pgfqpoint{3.157750in}{1.129625in}}%
\pgfpathlineto{\pgfqpoint{3.157750in}{1.225875in}}%
\pgfpathlineto{\pgfqpoint{3.098413in}{1.225875in}}%
\pgfpathlineto{\pgfqpoint{3.098413in}{1.129625in}}%
\pgfpathclose%
\pgfusepath{fill}%
\end{pgfscope}%
\begin{pgfscope}%
\pgfpathrectangle{\pgfqpoint{2.857045in}{0.350000in}}{\pgfqpoint{0.300705in}{0.962500in}}%
\pgfusepath{clip}%
\pgfsetbuttcap%
\pgfsetmiterjoin%
\definecolor{currentfill}{rgb}{0.557324,0.816424,0.546959}%
\pgfsetfillcolor{currentfill}%
\pgfsetlinewidth{0.000000pt}%
\definecolor{currentstroke}{rgb}{0.000000,0.000000,0.000000}%
\pgfsetstrokecolor{currentstroke}%
\pgfsetstrokeopacity{0.000000}%
\pgfsetdash{}{0pt}%
\pgfpathmoveto{\pgfqpoint{3.098288in}{1.139250in}}%
\pgfpathlineto{\pgfqpoint{3.157750in}{1.139250in}}%
\pgfpathlineto{\pgfqpoint{3.157750in}{1.235500in}}%
\pgfpathlineto{\pgfqpoint{3.098288in}{1.235500in}}%
\pgfpathlineto{\pgfqpoint{3.098288in}{1.139250in}}%
\pgfpathclose%
\pgfusepath{fill}%
\end{pgfscope}%
\begin{pgfscope}%
\pgfpathrectangle{\pgfqpoint{2.857045in}{0.350000in}}{\pgfqpoint{0.300705in}{0.962500in}}%
\pgfusepath{clip}%
\pgfsetbuttcap%
\pgfsetmiterjoin%
\definecolor{currentfill}{rgb}{0.546251,0.811257,0.537855}%
\pgfsetfillcolor{currentfill}%
\pgfsetlinewidth{0.000000pt}%
\definecolor{currentstroke}{rgb}{0.000000,0.000000,0.000000}%
\pgfsetstrokecolor{currentstroke}%
\pgfsetstrokeopacity{0.000000}%
\pgfsetdash{}{0pt}%
\pgfpathmoveto{\pgfqpoint{3.097927in}{1.148875in}}%
\pgfpathlineto{\pgfqpoint{3.157750in}{1.148875in}}%
\pgfpathlineto{\pgfqpoint{3.157750in}{1.245125in}}%
\pgfpathlineto{\pgfqpoint{3.097927in}{1.245125in}}%
\pgfpathlineto{\pgfqpoint{3.097927in}{1.148875in}}%
\pgfpathclose%
\pgfusepath{fill}%
\end{pgfscope}%
\begin{pgfscope}%
\pgfpathrectangle{\pgfqpoint{2.857045in}{0.350000in}}{\pgfqpoint{0.300705in}{0.962500in}}%
\pgfusepath{clip}%
\pgfsetbuttcap%
\pgfsetmiterjoin%
\definecolor{currentfill}{rgb}{0.540715,0.808674,0.533303}%
\pgfsetfillcolor{currentfill}%
\pgfsetlinewidth{0.000000pt}%
\definecolor{currentstroke}{rgb}{0.000000,0.000000,0.000000}%
\pgfsetstrokecolor{currentstroke}%
\pgfsetstrokeopacity{0.000000}%
\pgfsetdash{}{0pt}%
\pgfpathmoveto{\pgfqpoint{3.097593in}{1.158500in}}%
\pgfpathlineto{\pgfqpoint{3.157750in}{1.158500in}}%
\pgfpathlineto{\pgfqpoint{3.157750in}{1.254750in}}%
\pgfpathlineto{\pgfqpoint{3.097593in}{1.254750in}}%
\pgfpathlineto{\pgfqpoint{3.097593in}{1.158500in}}%
\pgfpathclose%
\pgfusepath{fill}%
\end{pgfscope}%
\begin{pgfscope}%
\pgfpathrectangle{\pgfqpoint{2.857045in}{0.350000in}}{\pgfqpoint{0.300705in}{0.962500in}}%
\pgfusepath{clip}%
\pgfsetbuttcap%
\pgfsetmiterjoin%
\definecolor{currentfill}{rgb}{0.535179,0.806090,0.528750}%
\pgfsetfillcolor{currentfill}%
\pgfsetlinewidth{0.000000pt}%
\definecolor{currentstroke}{rgb}{0.000000,0.000000,0.000000}%
\pgfsetstrokecolor{currentstroke}%
\pgfsetstrokeopacity{0.000000}%
\pgfsetdash{}{0pt}%
\pgfpathmoveto{\pgfqpoint{3.097287in}{1.168125in}}%
\pgfpathlineto{\pgfqpoint{3.157750in}{1.168125in}}%
\pgfpathlineto{\pgfqpoint{3.157750in}{1.264375in}}%
\pgfpathlineto{\pgfqpoint{3.097287in}{1.264375in}}%
\pgfpathlineto{\pgfqpoint{3.097287in}{1.168125in}}%
\pgfpathclose%
\pgfusepath{fill}%
\end{pgfscope}%
\begin{pgfscope}%
\pgfpathrectangle{\pgfqpoint{2.857045in}{0.350000in}}{\pgfqpoint{0.300705in}{0.962500in}}%
\pgfusepath{clip}%
\pgfsetbuttcap%
\pgfsetmiterjoin%
\definecolor{currentfill}{rgb}{0.524106,0.800923,0.519646}%
\pgfsetfillcolor{currentfill}%
\pgfsetlinewidth{0.000000pt}%
\definecolor{currentstroke}{rgb}{0.000000,0.000000,0.000000}%
\pgfsetstrokecolor{currentstroke}%
\pgfsetstrokeopacity{0.000000}%
\pgfsetdash{}{0pt}%
\pgfpathmoveto{\pgfqpoint{3.096650in}{1.177750in}}%
\pgfpathlineto{\pgfqpoint{3.157750in}{1.177750in}}%
\pgfpathlineto{\pgfqpoint{3.157750in}{1.274000in}}%
\pgfpathlineto{\pgfqpoint{3.096650in}{1.274000in}}%
\pgfpathlineto{\pgfqpoint{3.096650in}{1.177750in}}%
\pgfpathclose%
\pgfusepath{fill}%
\end{pgfscope}%
\begin{pgfscope}%
\pgfpathrectangle{\pgfqpoint{2.857045in}{0.350000in}}{\pgfqpoint{0.300705in}{0.962500in}}%
\pgfusepath{clip}%
\pgfsetbuttcap%
\pgfsetmiterjoin%
\definecolor{currentfill}{rgb}{0.518570,0.798339,0.515094}%
\pgfsetfillcolor{currentfill}%
\pgfsetlinewidth{0.000000pt}%
\definecolor{currentstroke}{rgb}{0.000000,0.000000,0.000000}%
\pgfsetstrokecolor{currentstroke}%
\pgfsetstrokeopacity{0.000000}%
\pgfsetdash{}{0pt}%
\pgfpathmoveto{\pgfqpoint{3.096154in}{1.187375in}}%
\pgfpathlineto{\pgfqpoint{3.157750in}{1.187375in}}%
\pgfpathlineto{\pgfqpoint{3.157750in}{1.283625in}}%
\pgfpathlineto{\pgfqpoint{3.096154in}{1.283625in}}%
\pgfpathlineto{\pgfqpoint{3.096154in}{1.187375in}}%
\pgfpathclose%
\pgfusepath{fill}%
\end{pgfscope}%
\begin{pgfscope}%
\pgfpathrectangle{\pgfqpoint{2.857045in}{0.350000in}}{\pgfqpoint{0.300705in}{0.962500in}}%
\pgfusepath{clip}%
\pgfsetbuttcap%
\pgfsetmiterjoin%
\definecolor{currentfill}{rgb}{0.501961,0.790588,0.501438}%
\pgfsetfillcolor{currentfill}%
\pgfsetlinewidth{0.000000pt}%
\definecolor{currentstroke}{rgb}{0.000000,0.000000,0.000000}%
\pgfsetstrokecolor{currentstroke}%
\pgfsetstrokeopacity{0.000000}%
\pgfsetdash{}{0pt}%
\pgfpathmoveto{\pgfqpoint{3.095304in}{1.197000in}}%
\pgfpathlineto{\pgfqpoint{3.157750in}{1.197000in}}%
\pgfpathlineto{\pgfqpoint{3.157750in}{1.293250in}}%
\pgfpathlineto{\pgfqpoint{3.095304in}{1.293250in}}%
\pgfpathlineto{\pgfqpoint{3.095304in}{1.197000in}}%
\pgfpathclose%
\pgfusepath{fill}%
\end{pgfscope}%
\begin{pgfscope}%
\pgfpathrectangle{\pgfqpoint{2.857045in}{0.350000in}}{\pgfqpoint{0.300705in}{0.962500in}}%
\pgfusepath{clip}%
\pgfsetbuttcap%
\pgfsetmiterjoin%
\definecolor{currentfill}{rgb}{0.485352,0.782837,0.487782}%
\pgfsetfillcolor{currentfill}%
\pgfsetlinewidth{0.000000pt}%
\definecolor{currentstroke}{rgb}{0.000000,0.000000,0.000000}%
\pgfsetstrokecolor{currentstroke}%
\pgfsetstrokeopacity{0.000000}%
\pgfsetdash{}{0pt}%
\pgfpathmoveto{\pgfqpoint{3.094254in}{1.206625in}}%
\pgfpathlineto{\pgfqpoint{3.157750in}{1.206625in}}%
\pgfpathlineto{\pgfqpoint{3.157750in}{1.302875in}}%
\pgfpathlineto{\pgfqpoint{3.094254in}{1.302875in}}%
\pgfpathlineto{\pgfqpoint{3.094254in}{1.206625in}}%
\pgfpathclose%
\pgfusepath{fill}%
\end{pgfscope}%
\begin{pgfscope}%
\pgfpathrectangle{\pgfqpoint{2.857045in}{0.350000in}}{\pgfqpoint{0.300705in}{0.962500in}}%
\pgfusepath{clip}%
\pgfsetbuttcap%
\pgfsetmiterjoin%
\definecolor{currentfill}{rgb}{0.463206,0.772503,0.469573}%
\pgfsetfillcolor{currentfill}%
\pgfsetlinewidth{0.000000pt}%
\definecolor{currentstroke}{rgb}{0.000000,0.000000,0.000000}%
\pgfsetstrokecolor{currentstroke}%
\pgfsetstrokeopacity{0.000000}%
\pgfsetdash{}{0pt}%
\pgfpathmoveto{\pgfqpoint{3.093230in}{1.216250in}}%
\pgfpathlineto{\pgfqpoint{3.157750in}{1.216250in}}%
\pgfpathlineto{\pgfqpoint{3.157750in}{1.312500in}}%
\pgfpathlineto{\pgfqpoint{3.093230in}{1.312500in}}%
\pgfpathlineto{\pgfqpoint{3.093230in}{1.216250in}}%
\pgfpathclose%
\pgfusepath{fill}%
\end{pgfscope}%
\begin{pgfscope}%
\pgfpathrectangle{\pgfqpoint{2.857045in}{0.350000in}}{\pgfqpoint{0.300705in}{0.962500in}}%
\pgfusepath{clip}%
\pgfsetbuttcap%
\pgfsetmiterjoin%
\definecolor{currentfill}{rgb}{0.457670,0.769919,0.465021}%
\pgfsetfillcolor{currentfill}%
\pgfsetlinewidth{0.000000pt}%
\definecolor{currentstroke}{rgb}{0.000000,0.000000,0.000000}%
\pgfsetstrokecolor{currentstroke}%
\pgfsetstrokeopacity{0.000000}%
\pgfsetdash{}{0pt}%
\pgfpathmoveto{\pgfqpoint{3.092778in}{1.225875in}}%
\pgfpathlineto{\pgfqpoint{3.157750in}{1.225875in}}%
\pgfpathlineto{\pgfqpoint{3.157750in}{1.322125in}}%
\pgfpathlineto{\pgfqpoint{3.092778in}{1.322125in}}%
\pgfpathlineto{\pgfqpoint{3.092778in}{1.225875in}}%
\pgfpathclose%
\pgfusepath{fill}%
\end{pgfscope}%
\begin{pgfscope}%
\pgfpathrectangle{\pgfqpoint{2.857045in}{0.350000in}}{\pgfqpoint{0.300705in}{0.962500in}}%
\pgfusepath{clip}%
\pgfsetbuttcap%
\pgfsetmiterjoin%
\definecolor{currentfill}{rgb}{0.439216,0.760938,0.455056}%
\pgfsetfillcolor{currentfill}%
\pgfsetlinewidth{0.000000pt}%
\definecolor{currentstroke}{rgb}{0.000000,0.000000,0.000000}%
\pgfsetstrokecolor{currentstroke}%
\pgfsetstrokeopacity{0.000000}%
\pgfsetdash{}{0pt}%
\pgfpathmoveto{\pgfqpoint{3.091947in}{1.235500in}}%
\pgfpathlineto{\pgfqpoint{3.157750in}{1.235500in}}%
\pgfpathlineto{\pgfqpoint{3.157750in}{1.331750in}}%
\pgfpathlineto{\pgfqpoint{3.091947in}{1.331750in}}%
\pgfpathlineto{\pgfqpoint{3.091947in}{1.235500in}}%
\pgfpathclose%
\pgfusepath{fill}%
\end{pgfscope}%
\begin{pgfscope}%
\pgfpathrectangle{\pgfqpoint{2.857045in}{0.350000in}}{\pgfqpoint{0.300705in}{0.962500in}}%
\pgfusepath{clip}%
\pgfsetbuttcap%
\pgfsetmiterjoin%
\definecolor{currentfill}{rgb}{0.401569,0.742484,0.436601}%
\pgfsetfillcolor{currentfill}%
\pgfsetlinewidth{0.000000pt}%
\definecolor{currentstroke}{rgb}{0.000000,0.000000,0.000000}%
\pgfsetstrokecolor{currentstroke}%
\pgfsetstrokeopacity{0.000000}%
\pgfsetdash{}{0pt}%
\pgfpathmoveto{\pgfqpoint{3.089899in}{1.245125in}}%
\pgfpathlineto{\pgfqpoint{3.157750in}{1.245125in}}%
\pgfpathlineto{\pgfqpoint{3.157750in}{1.341375in}}%
\pgfpathlineto{\pgfqpoint{3.089899in}{1.341375in}}%
\pgfpathlineto{\pgfqpoint{3.089899in}{1.245125in}}%
\pgfpathclose%
\pgfusepath{fill}%
\end{pgfscope}%
\begin{pgfscope}%
\pgfpathrectangle{\pgfqpoint{2.857045in}{0.350000in}}{\pgfqpoint{0.300705in}{0.962500in}}%
\pgfusepath{clip}%
\pgfsetbuttcap%
\pgfsetmiterjoin%
\definecolor{currentfill}{rgb}{0.363922,0.724029,0.418147}%
\pgfsetfillcolor{currentfill}%
\pgfsetlinewidth{0.000000pt}%
\definecolor{currentstroke}{rgb}{0.000000,0.000000,0.000000}%
\pgfsetstrokecolor{currentstroke}%
\pgfsetstrokeopacity{0.000000}%
\pgfsetdash{}{0pt}%
\pgfpathmoveto{\pgfqpoint{3.088237in}{1.254750in}}%
\pgfpathlineto{\pgfqpoint{3.157750in}{1.254750in}}%
\pgfpathlineto{\pgfqpoint{3.157750in}{1.351000in}}%
\pgfpathlineto{\pgfqpoint{3.088237in}{1.351000in}}%
\pgfpathlineto{\pgfqpoint{3.088237in}{1.254750in}}%
\pgfpathclose%
\pgfusepath{fill}%
\end{pgfscope}%
\begin{pgfscope}%
\pgfpathrectangle{\pgfqpoint{2.857045in}{0.350000in}}{\pgfqpoint{0.300705in}{0.962500in}}%
\pgfusepath{clip}%
\pgfsetbuttcap%
\pgfsetmiterjoin%
\definecolor{currentfill}{rgb}{0.332549,0.708651,0.402768}%
\pgfsetfillcolor{currentfill}%
\pgfsetlinewidth{0.000000pt}%
\definecolor{currentstroke}{rgb}{0.000000,0.000000,0.000000}%
\pgfsetstrokecolor{currentstroke}%
\pgfsetstrokeopacity{0.000000}%
\pgfsetdash{}{0pt}%
\pgfpathmoveto{\pgfqpoint{3.086626in}{1.264375in}}%
\pgfpathlineto{\pgfqpoint{3.157750in}{1.264375in}}%
\pgfpathlineto{\pgfqpoint{3.157750in}{1.360625in}}%
\pgfpathlineto{\pgfqpoint{3.086626in}{1.360625in}}%
\pgfpathlineto{\pgfqpoint{3.086626in}{1.264375in}}%
\pgfpathclose%
\pgfusepath{fill}%
\end{pgfscope}%
\begin{pgfscope}%
\pgfpathrectangle{\pgfqpoint{2.857045in}{0.350000in}}{\pgfqpoint{0.300705in}{0.962500in}}%
\pgfusepath{clip}%
\pgfsetbuttcap%
\pgfsetmiterjoin%
\definecolor{currentfill}{rgb}{0.282353,0.684045,0.378162}%
\pgfsetfillcolor{currentfill}%
\pgfsetlinewidth{0.000000pt}%
\definecolor{currentstroke}{rgb}{0.000000,0.000000,0.000000}%
\pgfsetstrokecolor{currentstroke}%
\pgfsetstrokeopacity{0.000000}%
\pgfsetdash{}{0pt}%
\pgfpathmoveto{\pgfqpoint{3.084131in}{1.274000in}}%
\pgfpathlineto{\pgfqpoint{3.157750in}{1.274000in}}%
\pgfpathlineto{\pgfqpoint{3.157750in}{1.370250in}}%
\pgfpathlineto{\pgfqpoint{3.084131in}{1.370250in}}%
\pgfpathlineto{\pgfqpoint{3.084131in}{1.274000in}}%
\pgfpathclose%
\pgfusepath{fill}%
\end{pgfscope}%
\begin{pgfscope}%
\pgfpathrectangle{\pgfqpoint{2.857045in}{0.350000in}}{\pgfqpoint{0.300705in}{0.962500in}}%
\pgfusepath{clip}%
\pgfsetbuttcap%
\pgfsetmiterjoin%
\definecolor{currentfill}{rgb}{0.234141,0.648443,0.348097}%
\pgfsetfillcolor{currentfill}%
\pgfsetlinewidth{0.000000pt}%
\definecolor{currentstroke}{rgb}{0.000000,0.000000,0.000000}%
\pgfsetstrokecolor{currentstroke}%
\pgfsetstrokeopacity{0.000000}%
\pgfsetdash{}{0pt}%
\pgfpathmoveto{\pgfqpoint{3.080977in}{1.283625in}}%
\pgfpathlineto{\pgfqpoint{3.157750in}{1.283625in}}%
\pgfpathlineto{\pgfqpoint{3.157750in}{1.379875in}}%
\pgfpathlineto{\pgfqpoint{3.080977in}{1.379875in}}%
\pgfpathlineto{\pgfqpoint{3.080977in}{1.283625in}}%
\pgfpathclose%
\pgfusepath{fill}%
\end{pgfscope}%
\begin{pgfscope}%
\pgfpathrectangle{\pgfqpoint{2.857045in}{0.350000in}}{\pgfqpoint{0.300705in}{0.962500in}}%
\pgfusepath{clip}%
\pgfsetbuttcap%
\pgfsetmiterjoin%
\definecolor{currentfill}{rgb}{0.134025,0.542330,0.268281}%
\pgfsetfillcolor{currentfill}%
\pgfsetlinewidth{0.000000pt}%
\definecolor{currentstroke}{rgb}{0.000000,0.000000,0.000000}%
\pgfsetstrokecolor{currentstroke}%
\pgfsetstrokeopacity{0.000000}%
\pgfsetdash{}{0pt}%
\pgfpathmoveto{\pgfqpoint{3.072500in}{1.293250in}}%
\pgfpathlineto{\pgfqpoint{3.157750in}{1.293250in}}%
\pgfpathlineto{\pgfqpoint{3.157750in}{1.389500in}}%
\pgfpathlineto{\pgfqpoint{3.072500in}{1.389500in}}%
\pgfpathlineto{\pgfqpoint{3.072500in}{1.293250in}}%
\pgfpathclose%
\pgfusepath{fill}%
\end{pgfscope}%
\begin{pgfscope}%
\pgfpathrectangle{\pgfqpoint{2.857045in}{0.350000in}}{\pgfqpoint{0.300705in}{0.962500in}}%
\pgfusepath{clip}%
\pgfsetbuttcap%
\pgfsetmiterjoin%
\definecolor{currentfill}{rgb}{0.000000,0.266667,0.105882}%
\pgfsetfillcolor{currentfill}%
\pgfsetlinewidth{0.000000pt}%
\definecolor{currentstroke}{rgb}{0.000000,0.000000,0.000000}%
\pgfsetstrokecolor{currentstroke}%
\pgfsetstrokeopacity{0.000000}%
\pgfsetdash{}{0pt}%
\pgfpathmoveto{\pgfqpoint{3.052541in}{1.302875in}}%
\pgfpathlineto{\pgfqpoint{3.157750in}{1.302875in}}%
\pgfpathlineto{\pgfqpoint{3.157750in}{1.399125in}}%
\pgfpathlineto{\pgfqpoint{3.052541in}{1.399125in}}%
\pgfpathlineto{\pgfqpoint{3.052541in}{1.302875in}}%
\pgfpathclose%
\pgfusepath{fill}%
\end{pgfscope}%
\begin{pgfscope}%
\definecolor{textcolor}{rgb}{0.000000,0.000000,0.000000}%
\pgfsetstrokecolor{textcolor}%
\pgfsetfillcolor{textcolor}%
\pgftext[x=2.977327in,y=1.360625in,left,base]{\color{textcolor}\setmainfont{Lato}\rmfamily\fontsize{7.000000}{8.400000}\selectfont 35.0\%}%
\end{pgfscope}%
\begin{pgfscope}%
\definecolor{textcolor}{rgb}{0.000000,0.000000,0.000000}%
\pgfsetstrokecolor{textcolor}%
\pgfsetfillcolor{textcolor}%
\pgftext[x=2.977327in,y=0.234500in,left,base]{\color{textcolor}\setmainfont{Lato}\rmfamily\fontsize{7.000000}{8.400000}\selectfont 8.3\%}%
\end{pgfscope}%
\end{pgfpicture}%
\makeatother%
\endgroup%


\vspace{-4mm}
\footnotesize{Source: American Community Survey, Dorn} \hspace{61mm} \tbllink{acs_cz_age.csv}
\vspace{3.5mm}

\begin{minipage}{0.76\textwidth}
\subsubsection*{Life expectancy}
\vspace{-0.5mm}

\small \input{text/life_exp.txt}

\begin{minipage}{0.435\textwidth} 
\normalsize \textbf{Life Expectancy at Birth}\\
\footnotesize{\textit{by group, in years}}\\
\noindent \hspace*{-2mm} \begin{tikzpicture}
	\begin{axis}[\bbar{y}{0}, \shdateaxisticks  
		xticklabel={`\short{\year}}, clip=false, 
		yticklabel style={text width=1.1em}, max space between ticks=20,
		height=5.7cm, width=6.0cm, enlarge y limits={0.08}]
	\rbars
	\thickline{cyan}{date}{Total}{data/life_exp.csv}
	\absnode{{1992-01-01}}{77.6}{\footnotesize \color{cyan}Total}
	\thickline{purple}{date}{Men}{data/life_exp.csv}
	\absnode{{2010-07-01}}{76}{\footnotesize \color{purple}Men}
	\thickline{blue!80!black}{date}{Women}{data/life_exp.csv}
	\absnode{{1991-01-01}}{80.6}{\footnotesize \color{blue!80!black}Women}
	\input{text/life_exp_node.txt}
	\end{axis}
\end{tikzpicture}\\
\footnotesize{Source: NCHS} \hfill \tbllink{life_exp.csv}\end{minipage}\hfill
\begin{minipage}{0.5\textwidth}
\small \input{text/life_exp2.txt}
\end{minipage}
\end{minipage}
\newpage
\begin{minipage}{0.76\textwidth}
\subsubsection*{Education}
\small Education is central in many discussions of the future of the US economy. Though very expensive in forgone years of earnings and often also expensive in tuition and other costs, education typically provides individuals with higher earnings. In response to changing job opportunities from globalization and other policy decisions, household spending on education has increased considerably, resulting in a much \textbf{more-educated population}.
\end{minipage}
\vspace{1mm}

\begin{minipage}{0.33\textwidth}
\small \input{text/cps_educ.txt} 
\end{minipage}\hspace{6mm}
\begin{minipage}{0.39\textwidth}
\normalsize \textbf{Highest Level of Education}\\
\footnotesize{\textit{age 25+ population, in millions, \input{text/cps_ltdate.txt}}}\\
\hspace*{-4mm} \begin{tikzpicture}
  	\begin{axis}[\barplotnogrid axis y line=left, \barylab{2.5cm}{1.5ex}
    	width=5.3cm, bar width=3.4ex, height=5.6cm,
    	enlarge y limits={abs=0.35cm}, enlarge x limits=0.05, \bbar{x}{0},
		yticklabels from table={\edtot}{name}, xmin=0, 
		yticklabel style={font=\footnotesize}, nodes near coords, 
		every node near coord/.style={/pgf/number format/fixed,
	    	font=\bfseries\scriptsize, style={white}},
		nodes near coords style={/pgf/number format/.cd, fixed zerofill,
			precision=1, assume math mode}, nodes near coords align={left}]
  	\addplot[fill=red!70!purple, draw=none] 
  		table [y expr=-\coordindex, x index=1] {\edtot};
  	\end{axis}
\end{tikzpicture}\\
\footnotesize{Source: Author's Calculations from CPS} \hfill \tbllink{cps_educ_tot.csv}
\end{minipage}
\vspace{1mm}

\begin{minipage}{0.76\textwidth}
\small \input{text/cps_educ2.txt}

Increased household spending on education may be a response to a weak labor market and lack of worker bargaining power. Behind the increase in education is a large increase in student debt. The burden of this debt is severe for many, as the more-educated workforce is not necessarily receiving the historical wage premium from education. 
\vspace{2mm}

\normalsize \textbf{Education Distribution} \small \hspace{10mm} \cbox{red!70!purple} \input{text/cps_ltdate.txt} \hspace{4mm} \cbox{orange!80} 2000\\
\footnotesize{\textit{share of age 25+ population, percent}}\\
\small \hspace{22mm} Total \hspace{54mm} Employed\\
\hspace*{-4mm} \begin{tikzpicture}
  	\begin{axis}[\barplotnogrid axis y line=left, \barylab{3.0cm}{1.5ex}
    	width=3.9cm, bar width=1.7ex, height=6.9cm,
    	enlarge y limits={abs=0.35cm}, enlarge x limits=0.05, \bbar{x}{0},
		yticklabels from table={\edsh}{name}, xmin=0, xmax=33,
		yticklabel style={font=\footnotesize}, nodes near coords, 
		nodes near coords style={/pgf/number format/.cd, fixed zerofill,
			precision=1, assume math mode}]
  	\addplot[fill=orange!80, draw=none] 
  		table [y expr=-\coordindex, x index=1] {\edsh};
  	\addplot[fill=red!70!purple, draw=none] 
  		table [y expr=-\coordindex, x index=2] {\edsh};
  	\end{axis}
\end{tikzpicture} \hfill
\begin{tikzpicture}
  	\begin{axis}[\barplotnogrid axis y line=left, \barylab{3.0cm}{1.5ex}
    	width=3.9cm, bar width=1.7ex, height=6.9cm,
    	enlarge y limits={abs=0.35cm}, enlarge x limits=0.05, \bbar{x}{0},
		yticklabels from table={\edsh}{name}, xmin=0, xmax=33,
		yticklabel style={font=\footnotesize}, nodes near coords, 
		nodes near coords style={/pgf/number format/.cd, fixed zerofill,
			precision=1, assume math mode}]
  	\addplot[fill=orange!80, draw=none] 
  		table [y expr=-\coordindex, x index=3] {\edsh};
  	\addplot[fill=red!70!purple, draw=none] 
  		table [y expr=-\coordindex, x index=4] {\edsh};
  	\end{axis}
\end{tikzpicture}\\
\footnotesize{Source: Author's Calculations from CPS} \hfill \tbllink{cps_educ.csv}
\end{minipage}
\newpage
\hypertarget{hhinc}{}
\begin{minipage}{0.76\textwidth}
\subsection*{Income, Spending, and Saving}
\small The next subsections cover household and personal income, spending, and saving. This subsection offers an overview, with mean and median per capita measures. 

\input{text/inc_out_save.txt}

The Consumer Expenditure Survey \href{https://www.bls.gov/cex/tables.htm}{provides} data on typical income, spending, and saving, measured as the median rather than the average. Median income, spending, and saving are not affected by the activities of the highest income families, which skew mean or average data. Saving is calculated as after-tax income minus spending, excluding spending on pensions (which count as saving). \input{text/ce_med_inc.txt}
\end{minipage}
\vspace{2mm}

\begin{minipage}{0.43\textwidth}
\normalsize \textbf{Average Income, Spending, and Saving}\\
\footnotesize{\textit{per capita, thousands of \input{text/pi_date.txt}\unskip \ US dollars}}
\vspace{4.8cm}

\hspace{4mm} \begin{tikzpicture}[overlay]
	\begin{axis}[\bbar{y}{0}, \shdateaxisticks ytick={0, 10, 20, 30, 40, 50, 60, 70, 80}, 
		ymin=0, ymax=57,
		yticklabel style={text width=0.9em},
		xticklabel={`\short{\year}}, clip=false, height=6.8cm, width=7.8cm]
	\rbars
	\absnode{{1992-01-01}}{41}{\footnotesize \color{blue!75!black}Income}
	\stdline{blue!75!black}{date}{DPI}{data/inc_out_save.csv}
	\absnode{{2011-01-01}}{41}{\footnotesize \color{orange}Spending}
	\stdline{orange}{date}{PCE}{data/inc_out_save.csv}
	\absnode{{2013-07-01}}{9}{\footnotesize \color{green!80!blue}Saving}
	\stdline{green!80!blue}{date}{PSAVE}{data/inc_out_save.csv}
	\end{axis}
\end{tikzpicture}

\footnotesize{Source: Bureau of Economic Analysis} \hfill \tbllink{inc_out_save.csv}
\end{minipage} \hspace{8mm}
\begin{minipage}{0.43\textwidth}
\normalsize \textbf{Median (Middle Quintile Average)}\\
\footnotesize{\textit{per householder, thousands of \input{text/pi_date.txt}\unskip \ US dollars}}
\vspace{4.8cm}

\hspace{4mm} \begin{tikzpicture}[overlay]
	\begin{axis}[\bbar{y}{0}, \shdateaxisticks ytick={0, 10, 20, 30, 40, 50, 60, 70, 80}, 
		ymin=0, ymax=57,
		yticklabel style={text width=0.9em},
		xticklabel={`\short{\year}}, clip=false, height=6.8cm, width=7.8cm]
	\rbars
	\absnode{{2010-01-01}}{27}{\footnotesize \color{blue!50!cyan}Income}
	\thickline{blue!50!cyan}{date}{INCAFTTX_4}{data/medincexp.csv}
	\absnode{{2010-01-01}}{18}{\footnotesize \color{red!50!orange}Spending}
	\thickline{red!50!orange}{date}{ADJEXP_4}{data/medincexp.csv}
	\absnode{{2012-10-01}}{7}{\footnotesize \color{green!50!blue}Saving}
	\thickline{green!50!blue}{date}{DI_SAVE_4}{data/medincexp.csv}
	\end{axis}
\end{tikzpicture}

\footnotesize{Source: Consumer Expenditure Survey} \hfill \tbllink{medincexp.csv}
\end{minipage}
\vspace{7mm}

\begin{minipage}{0.76\textwidth}
\normalsize \textbf{Average Income, Spending, and Saving}\\
\footnotesize{\textit{per capita, annualized, \input{text/pi_date.txt}\unskip \ US dollars}}\\
\noindent \hspace*{-3mm} \rowcolors{1}{}{black!5} \setlength{\tabcolsep}{3.1pt} \color{black!90}
		{\renewcommand{\arraystretch}{1.65}
		 \begin{tabular}{p{37mm} R{11.4mm} R{11.4mm} R{11.4mm} R{11.4mm} R{11.4mm} R{11.4mm}}
			 \input data/inc_out_save_levels.tex \hline
		\end{tabular}}\vspace{-2mm}
		
\footnotesize{Source: Bureau of Economic Analysis}
\end{minipage}
\newpage
\subsubsection*{Distribution by Income}
\begin{minipage}{0.76\textwidth}
\small Income varies massively by household. While some spending is non-discretionary, spending increases with income. The bottom 40 percent of households, by total money income, have expenses exceeding after-tax income. This includes retirees who are dissaving and low-income families taking on debt to cover expenses. Meanwhile, the top ten percent of households save nearly half of their income.

\input{text/ce_inc_pct.txt}
\end{minipage}
\vspace{1mm}

\begin{minipage}{0.68\textwidth}
\normalsize \textbf{Household Income and Spending, by Income Percentile}\\
\footnotesize{\textit{thousands of \input{text/ce_year.txt}\unskip \ dollars}}\\
\hspace*{-2mm} \begin{tikzpicture}
  	\begin{axis}[ybar=20pt, \bbar{y}{0}, ytick={0, 50, 100, 150, 200, 250}, 
    	width=11.8cm, bar width=20pt, height=5.5cm, ymin=11.5, ymax=260,
    	enlarge x limits={abs=0.8cm}, enlarge y limits=0.05, 
    	symbolic x coords={Bottom 20 Percent, Percentile 20--40, Percentile 40--60, Percentile 60--80, Percentile 80--90, Top 10 Percent},
    	xtick=data, yticklabel style={text width=0.7em}, 
		xticklabels from table={\ceinc}{name}, 
		xticklabel style={black!70, font=\footnotesize, text width=1.8cm, align=center}, 
		clip=false, nodes near coords,  
		nodes near coords style={/pgf/number format/.cd, fixed zerofill,
			precision=1, assume math mode}, legend cell align={left},
		legend style={at={(0.23, 0.97)}, legend columns=1, 
		fill=white, draw=none}]
  	\addplot[ybar, bar shift=-10pt, fill=blue!70!black, draw=none, area legend] 
  		table [y index=1] {\ceinc};
  	\addplot[ybar, bar shift=10pt, fill=blue!30!cyan, draw=none, area legend] 
  		table [y index=2] {\ceinc};
  	\legend{Income, Spending};
  	\end{axis}
\end{tikzpicture}\\
\footnotesize{Income is after taxes; spending does not include spending on pensions} \\
\footnotesize{Source: Consumer Expenditure Survey} \hfill \tbllink{ce_inc.csv}
\end{minipage}
\vspace{2mm}

\subsubsection*{Distribution by Age}
\begin{minipage}{0.76\textwidth}
\small Income and consumer spending both vary by age and peak around age 45 to 54. The oldest age group, households where the reference person is 75 or older, have income below their expenses. Likewise, the youngest group, with the reference person younger than 25, have income around or below expenses.

\input{text/ce_age.txt}
\end{minipage}
\vspace{1mm}

\begin{minipage}{0.68\textwidth}
\normalsize \textbf{Household Income and Spending, by Age of Reference Person}\\
\footnotesize{\textit{thousands of \input{text/ce_year.txt}\unskip \ dollars}}\\
\hspace*{-2mm} \begin{tikzpicture}
  	\begin{axis}[ybar=16pt, \bbar{y}{0}, ytick={0, 20, 40, 60, 80, 100}, 
    	width=11.8cm, bar width=16pt, height=5.5cm, ymin=4.5, ymax=100,
    	enlarge x limits={abs=0.8cm}, enlarge y limits=0.02, 
    	symbolic x coords={Under 25, 25-34, 35-44, 45-54, 55-64, 65-74, 75+},
    	xtick=data, yticklabel style={text width=0.7em}, 
		xticklabels from table={\ceage}{name}, 
		xticklabel style={black!70, font=\small}, clip=false,
		nodes near coords,  
		nodes near coords style={/pgf/number format/.cd, fixed zerofill,
			precision=1, assume math mode}, legend cell align={left},
		legend style={at={(0.21, 1.05)}, legend columns=1, 
		fill=white, draw=none}]
  	\addplot[ybar, bar shift=-8pt, fill=green!80!blue, draw=none, area legend] 
  		table [y index=1] {\ceage};
  	\addplot[ybar, bar shift=8pt, fill=blue!40!magenta, draw=none, area legend] 
  		table [y index=2] {\ceage};
  	\legend{Income, Spending};
  	\end{axis}
\end{tikzpicture}\\
\footnotesize{Income is after taxes; spending does not include spending on pensions} \\
\footnotesize{Source: Consumer Expenditure Survey} \hfill \tbllink{ce_age.csv}
\end{minipage}
\newpage
\subsection*{Income to Persons}
\begin{minipage}{0.3\textwidth}
\small \href{https://www.bea.gov/data/income-saving/personal-income}{Personal income} includes labor income (see\cbox{green!75!black}), measured as compensation of employees, capital income (see\cbox{orange!40!yellow}), measured as the sum of proprietor income, rental income, and dividend and interest income, and welfare income (see\cbox{blue!80!white}), measured primarily as government social benefits less contributions to social insurance. 

\input{text/pi_levels_basic.txt}
\end{minipage} \hspace{6mm}
\begin{minipage}{0.65\textwidth}
\normalsize \textbf{Personal Income}\\
\footnotesize{\textit{real per capita income, thousands of US dollars}}
\vspace{5.2cm}

\hspace{4mm} \begin{tikzpicture}[overlay]
	\begin{axis}[\bbar{y}{0}, \shdateaxisticks ytick={0, 10, 20, 30, 40, 50, 60, 70},
		xmin={1989-08-01},
		xticklabel={`\short{\year}}, clip=false, ymin=3, ymax=66.2,
		yticklabel style={text width=1.1em}, reverse legend, legend cell align=left,
		legend style={at={(0.27, 0.98)}, legend columns=1}, width=11.2cm, height=7.2cm]
	\rbars
	\ctsbar{green!75!black}{date}{LABOR}{data/pi_levels.csv}{0.75pt}
	\ctsbar{orange!40!yellow}{date}{CAPITAL}{data/pi_levels.csv}{0.75pt}
	\ctsbar{blue!80!white}{date}{TRANSFER}{data/pi_levels.csv}{0.75pt}
	\legend{ \ Labor, \ Capital, \ Welfare};
	\end{axis}
\end{tikzpicture}

\footnotesize{Source: Bureau of Economic Analysis; Author} \hfill \tbllink{pi_levels.csv}
\end{minipage}
\vspace{3mm}

\begin{minipage}{0.76\textwidth}
\normalsize \textbf{Personal Income by Source}\\
\footnotesize{\textit{per capita, annualized, \input{text/pi_date.txt}\unskip \ US dollars}}\\
\noindent \hspace*{-3mm} \rowcolors{1}{}{black!5} \setlength{\tabcolsep}{3.1pt} \color{black!90}
		{\renewcommand{\arraystretch}{1.52}
		 \begin{tabular}{p{37mm} R{11.4mm} R{11.4mm} R{11.4mm} R{11.4mm} R{11.4mm} R{11.4mm}}
			 \input data/pi_levels.tex \hline
		\end{tabular}}\vspace{-2mm}
		
\footnotesize{Source: Bureau of Economic Analysis}
\end{minipage}
\newpage
\begin{minipage}{0.76\textwidth}
\normalsize \textbf{Personal Income}\\
\footnotesize{\textit{percentage point contribution to real personal income growth}}
\vspace{6.6cm}

\hspace{4mm} \begin{tikzpicture}[overlay]
	\begin{axis}[\bbar{y}{0}, \dateaxisticks ytick={-30, -20, -10, 0, 10, 20, 30, 40, 50},
		height=8.6cm, ymax=46, reverse legend, legend cell align=left,
		xticklabel={`\short{\year}}, clip=false, yticklabel style={text width=1.4em}, 
		legend style={at={(0.28, 0.99)}, legend columns=1}]
	\rbars
	\sbar{green!75!black}{date}{A033RC}{data/pi.csv}
	\sbar{orange!40!yellow}{date}{CAPITAL}{data/pi.csv}
	\sbar{blue!80!white}{date}{TRANSFER}{data/pi.csv}
	\legend{Labor, Capital$^1$, Welfare$^2$};
	\stdnode{3.4cm}{1.1cm}{\scriptsize $^1$ Includes proprietor, rent, and asset income\\ \scriptsize $^2$ Current transfer receipts minus contributions\\\scriptsize \ \ \ \ to government social insurance}
	\end{axis}
\end{tikzpicture}

\footnotesize{Source: Bureau of Economic Analysis} \hfill \tbllink{pi.csv}
\vspace{3mm}

\small &   2023  Q3 & `23  Q2 & `23  Q1 & `22  Q4 & `22  Q3 &1.0-year&10.0-year&30.0-year\\  \hspace{0.5mm}Personal  Income  (Pre-Tax  Income) &1.36&1.65&0.69&-1.18&-2.25&0.63&2.62&2.69\\  \hspace{-2mm}\cbox{green!55!blue!72!black}Labor &1.23&1.53&0.55&-0.53&0.38&0.70&1.52&1.54\\  \hspace{3mm}  Wages  \&  Salaries &1.08&1.42&0.58&-0.32&0.61&0.69&1.37&1.30\\  \hspace{3mm}  Supplements  to  Wages  \&  Salaries &0.15&0.11&-0.03&-0.21&-0.23&0.01&0.15&0.24\\  \hspace{-2mm}\cbox{yellow!85!orange!88!white}Capital &0.36&0.43&0.51&-0.07&-0.49&0.31&0.71&0.78\\  \hspace{3mm}  Proprietors'  Income &-0.04&-0.08&-0.08&-0.21&-0.46&-0.10&0.07&0.25\\  \hspace{3mm}  Rental  Income &0.23&0.23&0.30&0.10&0.07&0.22&0.11&0.16\\  \hspace{3mm}  Personal  Interest  Income &0.32&0.39&0.51&0.42&0.15&0.41&0.12&0.07\\  \hspace{3mm}  Personal  Dividend  Income &-0.16&-0.11&-0.22&-0.37&-0.26&-0.22&0.40&0.30\\  \hspace{-2mm}\cbox{cyan!90!blue!90!white}Net  Government  Benefits &-0.23&-0.31&-0.38&-0.59&-2.15&-0.38&0.40&0.37\\  \hspace{2mm}  Government  Social  Benefits &-0.11&-0.10&-0.25&-0.50&-1.92&-0.24&0.59&0.55\\  \hspace{3mm}  Social  Security &0.47&0.44&0.36&0.15&0.11&0.36&0.16&0.16\\  \hspace{3mm}  Medicare &-0.05&-0.08&-0.12&-0.09&-0.06&-0.08&0.12&0.14\\  \hspace{3mm}  Medicaid &0.15&0.27&0.18&0.17&0.09&0.19&0.16&0.13\\  \hspace{3mm}  Unemployment  Insurance &0.00&0.00&-0.02&-0.06&-0.99&-0.02&-0.00&0.00\\  \hspace{3mm}  Veterans'  Benefits &-0.02&-0.02&-0.00&0.02&0.03&-0.00&0.04&0.03\\  \hspace{3mm}  Other &-0.65&-0.72&-0.64&-0.69&-1.11&-0.68&0.11&0.10\\  \hspace{2mm}  Less:  Social  Insurance  Contributions &-0.13&-0.20&-0.12&-0.07&-0.22&-0.13&-0.22&-0.19\\  
\end{minipage}

\begin{minipage}{0.81\textwidth}
\normalsize \textbf{Personal Income by Source}\\
\footnotesize{\textit{percentage point contribution to real personal income growth \hspace{22mm} moving averages}}\\
\rowcolors{1}{}{black!5} \setlength{\tabcolsep}{3.1pt} \color{black!90}
	{\renewcommand{\arraystretch}{1.52}
		\begin{tabular}{p{39.5mm} R{8.2mm} R{8.2mm} R{8.2mm} R{8.2mm} R{8.2mm} 
		   R{7.8mm} R{6.7mm} R{6.7mm} }
			 \input data/pi.tex \hline
		\end{tabular}}\vspace{-2mm}
		
\footnotesize{Source: Bureau of Economic Analysis}
\end{minipage}
\newpage
\begin{minipage}{0.76\textwidth}
\small \input{text/rdpigrowth.txt}
\vspace{2mm}

\normalsize \textbf{After-Tax Income Growth}\\
\footnotesize{\textit{annual growth, per capita real disposable personal income, percent}}
\vspace{4.4cm}

\hspace{4mm} \begin{tikzpicture}[overlay]
	\begin{axis}[\bbar{y}{0}, \dateaxisticks ytick={-10, -5, 0, 5, 10, 15, 20, 25}, 
		yticklabel style={text width=1.4em}, ymax=23.2,
		xticklabel={`\short{\year}}, clip=false, height=6.4cm, width=12.1cm]
	\rbars
	\stdline{blue!80!cyan}{date}{rdpi}{data/rdpigrowth.csv}
	\input{text/rdpi_node.txt}
	\end{axis}
\end{tikzpicture}

\footnotesize{Source: Bureau of Economic Analysis}  \hfill \tbllink{rdpigrowth.csv}
\vspace{4mm}

\subsubsection*{Distribution of Personal Income}
\small Labor income, which includes wages and salaries as well as self-employment income, is the vast majority of personal income. \input{text/pearn.txt}
\vspace{2mm}

\normalsize \textbf{Distribution of Personal Income, 2020}\\
\footnotesize{\textit{by percentile of income, thousands of US dollars}}\\
\begin{tikzpicture}
	\begin{axis}[ymin=0, ymax=150, enlarge x limits={0.02},
		width=12.4cm, height=6.4cm, ytick={0, 12.8, 50, 100, 150},
		reverse legend,
		legend cell align={left}, yticklabel style={text width=1.8em}, 
		legend style={legend columns=1, at={(0.25, 0.98)}}]
	\draw [black!60] (axis cs:{50},\pgfkeysvalueof{/pgfplots/ymin}) -- 
		(axis cs:{50}, \pgfkeysvalueof{/pgfplots/ymax});
	\draw [dashed] (axis cs:{79, 150}) -- 
		(axis cs:{100, 150});
	\abar{cyan!50!white}{percentile}{PEARNVAL}{data/pearn_dist.csv}
	\abar{violet!50!purple}{percentile}{Diff}{data/pearn_dist.csv}
	\absnode{40.8}{40}{\scriptsize \color{black!60} median}
	\absnode{79}{148}{\scriptsize \color{black!60} values above \\\scriptsize \color{black!60} \$150,000 \\\scriptsize \color{black!60} not shown}
	\legend{\ Labor Income,\ Total Income};
	\end{axis}
\end{tikzpicture}\\
\footnotesize{Source: Author's Calculations from CPS ASEC} \hfill \tbllink{pearn_dist.csv}
\end{minipage}
\newpage 
\begin{minipage}{0.76\textwidth} 
\subsubsection*{Contributions to Personal Income Growth}
\small Annual data on personal income \href{https://www.census.gov/data/tables/time-series/demo/income-poverty/cps-pinc/pinc-09.html}{detail} the number of people receiving various categories of income and the average payment. As a result, it is possible to match changes in aggregate personal income from changes in payment amounts (see\cbox{blue!60!cyan}) and changes in who is receiving payments (see\cbox{green!90!blue}). 

\input{text/pi_ch_ann.txt}after adjusting for changes in prices. Compared to 2019, many fewer people received earnings from work, while mean earnings fell as many were unemployed for part of the year. These changes were partially offset by unemployment insurance payments, as well as larger property income amounts. 
\vspace{2mm}

\begin{minipage}{0.55\textwidth}
\normalsize \textbf{Sources of Personal Income Growth, 2020}\\
\footnotesize{\textit{percentage point contribution to aggregate growth}}\\
\hspace*{-4mm} \begin{tikzpicture}
  	\begin{axis}[\barplotnogrid axis y line=left, \barylab{2.7cm}{1.5ex}
    	width=5.9cm, bar width=2.4ex, height=5.0cm,
    	enlarge y limits={abs=0.3cm}, enlarge x limits=0.05, \bbar{x}{0},
		yticklabels from table={\pinc}{Name}, xmin=-1.6, 
		yticklabel style={font=\footnotesize}, nodes near coords, 
		nodes near coords style={/pgf/number format/.cd, fixed zerofill,
			precision=1, assume math mode},
		legend style={text=black!70, at={(0.23, 1.2)}, anchor=north, legend columns=-1, 
				fill=none, draw=none,
		        /tikz/every even column/.append style={column sep=0.4cm}}]
  	\addplot[fill=green!90!blue, draw=none] 
  		table [y expr=-\coordindex, x index=1] {\pinc};
  	\addplot[fill=blue!60!cyan, draw=none] 
  		table [y expr=-\coordindex, x index=2] {\pinc};
  	\legend{Number of recipients, Amount per recipient}
  	\end{axis}
\end{tikzpicture}\\
\footnotesize{Source: Census Bureau}
\end{minipage} 
\vspace{5mm}

\subsubsection*{Household Income}
\small Given massive variance in personal income, with half of the population receiving no market income, many people rely on the income of other members of their household for survival. Next, we discuss household income, which measures the combined income of all people in a given housing unit. Like personal income, household income is distributed very unevenly in the US.
\end{minipage}
\vspace{1mm}

\begin{minipage}{0.36\textwidth}
\normalsize \textbf{Real Median Household Income}\\
\footnotesize{\textit{thousands of \input{text/rmhi_ltdt.txt}US dollars}}\\
\hspace*{-3mm} \begin{tikzpicture}
	\begin{axis}[\bbar{y}{1}, \shdateaxisticks ytick={20, 40, 60, 80, 100}, 
		yticklabel style={text width=1.3em}, 
		enlarge y limits={0.04}, 
		xticklabel={`\short{\year}}, width=6.8cm, height=6.4cm]
	\rbars
	\stdline{blue!80!black}{date}{Black_LT}{data/real_med_hh_inc.csv}
	\stdline{blue!80!black}{date}{Black_PR}{data/real_med_hh_inc.csv}
	\stdline{blue!80!black}{date}{Black_PR2}{data/real_med_hh_inc.csv}
	\stdline{green!80!black}{date}{White_LT}{data/real_med_hh_inc.csv}
	\stdline{green!80!black}{date}{White_PR}{data/real_med_hh_inc.csv}
	\stdline{green!80!black}{date}{White_PR2}{data/real_med_hh_inc.csv}
	\stdline{orange}{date}{Asian_LT}{data/real_med_hh_inc.csv}
	\stdline{orange}{date}{Asian_PR}{data/real_med_hh_inc.csv}
	\stdline{orange}{date}{Asian_PR2}{data/real_med_hh_inc.csv}
	\stdline{red!90!black}{date}{Hispanic_LT}{data/real_med_hh_inc.csv}
	\stdline{red!90!black}{date}{Hispanic_PR}{data/real_med_hh_inc.csv}
	\stdline{red!90!black}{date}{Hispanic_PR2}{data/real_med_hh_inc.csv}
	\absnode{{2015-03-01}}{40}{\footnotesize \color{blue!80!black}Black}
	\absnode{{2014-01-01}}{67}{\footnotesize \color{green!80!black}White}
	\absnode{{2011-12-01}}{94}{\footnotesize \color{orange}Asian}
	\absnode{{2010-02-01}}{58.5}{\footnotesize \color{red!90!black}Hispanic}
	\end{axis}
\end{tikzpicture}\\
\footnotesize{Source: Census Bureau} \hfill \tbllink{real_med_hh_inc.csv}
\end{minipage}\hspace{7mm}
\begin{minipage}{0.35\textwidth}
\small \input{text/real_med_hh_inc.txt}

Two values are shown for 2013 and 2017 to mark revisions to the survey design (2013) and the processing of survey data (2017). These data are not perfectly comparable over time. 
\end{minipage}
\newpage 
\begin{minipage}{0.76\textwidth}
\subsection*{Household Spending and Saving}
\hypertarget{hhss}{}
\small The previous subsection examines money coming into households while this section covers money leaving households. \href{https://www.bea.gov/data/consumer-spending/main}{Consumer spending} includes household purchases of goods and services, certain spending for households, such as Medicare and Medicaid, and estimates for services furnished without payment. As an example of a special estimate, homeowners are considered to rent their homes to themselves. 

Over the past 30 years, increased consumer spending represents the majority of economic growth. Consumer spending usually increases when households have more income and falls when households have less income. This effect is visible in both the long-run and during the course of a business cycle, with consumer spending generally falling or slowing during a recession. During the COVID-19 pandemic, some categories of spending were lower as a result of business closures and restrictions.

\href{https://www.bea.gov/data/income-saving/personal-saving-rate}{Personal saving} occurs when households have income in excess of their expenses. Savings are invested, often providing additional income, and are used for future expenses, such as costs incurred during retirement.

\subsubsection*{Spending}
\small Consumer spending is comprised of two broad expenditure types: goods and services. Spending on goods includes durable goods (goods with a useful life of at least three years), such as cars, furniture, or recreational goods, and nondurable goods, such as groceries, clothing, and gasoline. Spending on services includes housing, health care, restaurants and bars, transportation services, financial services, and other services. 
\end{minipage}
\vspace{2mm}

\begin{minipage}{0.37\textwidth}
\normalsize \textbf{Expenditure Types}\\
\footnotesize{\textit{per capita, thousands of \input{text/gdp_ltdt.txt}dollars}}\\
\hspace*{-2mm} \begin{tikzpicture}
	\begin{axis}[\bbar{y}{0}, \shdateaxisticks ytick={0, 10, 20, 30}, ymin=1, 
		yticklabel style={text width=1.0em},
		xticklabel={`\short{\year}}, clip=false, height=5.2cm, width=7.0cm]
	\rbars
	\stdnode{4.3cm}{1.8cm}{\footnotesize \color{red}Goods}
	\stdline{red}{date}{Goods}{data/pce_levels.csv}
	\stdnode{1.05cm}{2.9cm}{\footnotesize \color{orange}Services}
	\stdline{orange}{date}{Services}{data/pce_levels.csv}
	\end{axis}
\end{tikzpicture}\\
\footnotesize{Source: Bureau of Economic Analysis} \hspace{8mm} \tbllink{pce_levels.csv}
\end{minipage} \hspace{8mm}
\begin{minipage}{0.33\textwidth}
\small &2023
Q3&2023
Q2&2022
Q3&2019
Q4&2000
Q1&1989
Q1\\  Consumer  Spending & \$55,676 &55,331&54,751&51,410&38,824&30,333\\  \hspace*{-0.6mm}  {\color{red}\textbf{---}}  Goods &18,575&18,381&18,201&15,910&10,554&7,919\\  \hspace{4mm}  Motor  Vehicles  \&  Parts &2,279&2,289&2,171&2,128&1,781&1,340\\  \hspace{4mm}  Furniture  \&  HH  Equipment &1,427&1,410&1,413&1,232&614&425\\  \hspace{4mm}  Recreational  Durable  Goods &2,056&1,980&1,907&1,266&262&76\\  \hspace{4mm}  Groceries &4,315&4,301&4,356&4,072&3,306&3,284\\  \hspace{4mm}  Clothes  \&  Shoes &1,552&1,531&1,557&1,329&956&710\\  \hspace*{-0.6mm}  {\color{blue!75!white}\textbf{---}}  Services  Excluding  Shelter &27,239&27,146&26,730&26,092&20,040&15,377\\  \hspace{4mm}  Health  Care  Services &8,969&8,922&8,556&8,371&5,514&4,993\\  \hspace{4mm}  Transportation &1,800&1,799&1,814&1,818&1,647&1,144\\  \hspace{4mm}  Recreational &2,153&2,144&2,079&2,137&1,644&1,174\\  \hspace{4mm}  Food  \&  Accommodations &4,107&4,043&4,017&3,840&2,963&2,721\\  \hspace{4mm}  Financial  \&  Insurance &3,982&3,975&3,931&4,021&4,137&2,451\\  \hspace*{-0.6mm}  {\color{green!85!blue}\textbf{---}}  Shelter  Including  Investment &13,089&12,984&13,314&12,815&12,768&10,962\\  \hspace{4mm}  Housing  Services  \&  Utilities   &9,863&9,805&9,820&9,405&8,559&7,439\\  \hspace{4mm}  Residential  Fixed  Investment &3,226&3,179&3,495&3,410&4,208&3,523\\ 
\end{minipage}
\vspace{2mm}

\begin{minipage}{0.37\textwidth}
\normalsize \textbf{Shelter Costs}\\
\footnotesize{\textit{per capita, thousands of \input{text/gdp_ltdt.txt}dollars}}\\
\hspace*{-2mm} \begin{tikzpicture}
	\begin{axis}[\bbar{y}{0}, \shdateaxisticks ytick={0, 2, 4, 6, 8, 10}, ymin=0.3, 
		yticklabel style={text width=1.0em},
		xticklabel={`\short{\year}}, clip=false, height=5.2cm, width=6.8cm]
	\rbars
	\stdnode{3.9cm}{3.0cm}{\footnotesize \color{green!60!blue}Housing and Utilities}
	\stdline{green!60!blue}{date}{Housing}{data/pce_levels.csv}
	\stdnode{1.55cm}{0.65cm}{\footnotesize \color{blue!80!black}Residential Investment}
	\stdline{blue!80!black}{date}{ResInv}{data/pce_levels.csv}
	\end{axis}
\end{tikzpicture}\\
\footnotesize{Source: Bureau of Economic Analysis} \hspace{8mm} \tbllink{pce_levels.csv}
\end{minipage} \hspace{8mm}
\begin{minipage}{0.33\textwidth}
\small \input{text/pce2_levels.txt}
\end{minipage}
\newpage 
\begin{minipage}{0.76\textwidth}
\small The combined categories from the previous two charts cover spending on goods, spending on services other than shelter, and spending on housing, utilities, and residential construction. Including residential construction with other housing costs provides a more broad overview of household sector expenses. 
\end{minipage}
\vspace{2mm}

\begin{minipage}{0.37\textwidth}
\normalsize \textbf{Expenditure Types}\\
\footnotesize{\textit{per capita, thousands of \input{text/gdp_ltdt.txt}dollars}}\\
\hspace*{-2mm} \begin{tikzpicture}
	\begin{axis}[\bbar{y}{0}, \shdateaxisticks ytick={0, 5, 10, 15, 20, 25}, ymin=1, 
		yticklabel style={text width=1.0em},
		xticklabel={`\short{\year}}, clip=false, height=6.4cm, width=6.85cm]
	\rbars
	\absnode{{2010-01-01}}{15.0}{\footnotesize \color{red}Goods}
	\stdline{red}{date}{Goods}{data/pce_levels.csv}
	\absnode{{1998-07-01}}{24}{\footnotesize \color{blue!75!white}Services ex. Shelter}
	\stdline{blue!75!white}{date}{OTHSERV}{data/pce_levels.csv}
	\absnode{{1994-10-01}}{13.8}{\footnotesize \color{green!85!blue}Shelter}
	\stdline{green!85!blue}{date}{SHELTER}{data/pce_levels.csv}
	\end{axis}
\end{tikzpicture}\\
\footnotesize{Source: Bureau of Economic Analysis} \hfill \tbllink{pce_levels.csv}
\end{minipage} \hspace{8mm}
\begin{minipage}{0.33\textwidth}
\small \input{text/pce3_levels.txt}
\end{minipage}

\vspace{3mm}

\begin{minipage}{0.76\textwidth}
\normalsize \textbf{Expenditure Types}\\
\footnotesize{\textit{per capita, annualized, \input{text/gdp_ltdt.txt}dollars}}\\
\hspace{-2mm} \rowcolors{1}{}{black!5} \setlength{\tabcolsep}{3.7pt} \color{black!90}
		{\renewcommand{\arraystretch}{1.55}
		 \begin{tabular}{p{52mm} R{13mm} R{13mm} R{13mm} R{13mm}}
			\input data/pce_levels.tex \hline
		\end{tabular}}\vspace{-2mm}
		
\footnotesize{Source: Bureau of Economic Analysis}
\end{minipage}
\newpage 
\begin{minipage}{0.76\textwidth}
\small Next, we examine how changes in consumer spending on goods (see\cbox{red}), services excluding housing and utilities (see\cbox{blue!75!white}), and shelter (see\cbox{green!85!blue}, calculated as housing and utilities plus residential fixed investment), affect GDP growth. \input{text/pce1.txt}
\vspace{1mm}

\normalsize \textbf{Consumer Spending and Residential Investment}\\
\footnotesize{\textit{percentage point contribution to real GDP growth, one-year moving average}}
\vspace{4.4cm}

\hspace{4mm} \begin{tikzpicture}[overlay]
	\begin{axis}[\bbar{y}{0}, \dateaxisticks 
		xticklabel={`\short{\year}}, clip=false, height=6.4cm, reverse legend,
		legend cell align=left, ymax=9.7, max space between ticks=20,
		yticklabel style={text width=1.5em}, enlarge y limits={0.02},
		legend style={at={(0.25, 0.99)}, legend columns=1}]
	\rbars
	\sbar{red}{date}{DGDSRY}{data/pce_ma.csv}
	\sbar{blue!75!white}{date}{OTHSERV}{data/pce_ma.csv}
	\sbar{green!85!blue}{date}{SHELTER}{data/pce_ma.csv}
	\legend{Goods, Services, Shelter$^1$};
	\stdnode{2.2cm}{0.6cm}{\footnotesize $^1$ Includes residential fixed investment}
	\end{axis}
\end{tikzpicture}

\footnotesize{Source: Bureau of Economic Analysis} \hfill \tbllink{pce.csv}
\vspace{2mm}

\small \input{text/pce2.txt} 
\end{minipage}
\vspace{1mm}

\begin{minipage}{0.80\textwidth}
\normalsize \textbf{Consumer Spending and Residential Investment}\\
\footnotesize{\textit{percentage point contribution to real GDP growth \hspace{36mm} moving averages}}\\
\rowcolors{1}{}{black!5} \setlength{\tabcolsep}{3.0pt} \color{black!90}
	{\renewcommand{\arraystretch}{1.52}
\hspace*{-2mm} \begin{tabular}{p{2.0mm} p{39mm} R{8.2mm} R{8.2mm} R{8.2mm} R{8.2mm} R{8.2mm} 
					R{7.8mm} R{7.2mm} R{7.2mm}}
		\input data/pce.tex \hline
	\end{tabular}}
	\vspace{-2mm}
	
\footnotesize{Source: Bureau of Economic Analysis}
\end{minipage}
\newpage 
\begin{minipage}{0.76\textwidth}
\small \input{text/pcegrowth.txt}
\vspace{1mm}

\normalsize \textbf{Consumer Spending Growth}\\
\footnotesize{\textit{annual growth, per capita real personal consumption expenditures, percent}}
\vspace{4.6cm}

\hspace{4mm} \begin{tikzpicture}[overlay]
	\begin{axis}[\bbar{y}{0}, \dateaxisticks  enlarge y limits={0.03}, ymax=19,
		yticklabel style={text width=1.5em}, max space between ticks=20,
		xticklabel={`\short{\year}}, clip=false, height=6.6cm, width=12.2cm]
	\rbars
	\stdline{green!70!olive!60!black}{date}{rpcepop}{data/pcegrowth.csv}
	\input{text/pcegrowth_node.txt}
	\end{axis}
\end{tikzpicture}

\footnotesize{Source: Bureau of Economic Analysis}  \hfill \tbllink{pcegrowth.csv}
\vspace{4mm}

\small Changes to consumer spending (see {\color{black}\textbf{---}}) are largely the result of changes to income (see\cbox{cyan!28!white}) and changes to the rate at which income is saved (see\cbox{green!72!blue!90!white}). Changes to other outlays (see\cbox{blue!92!violet}) reflect changes in interest payments, fines and fees, and charitable giving. 

\input{text/pcedecomp.txt}
\vspace{1mm} 

\normalsize \textbf{Contributions to Consumer Spending}\\
\footnotesize{\textit{percentage point contribution to real per capita PCE growth, one-year moving average}}
\vspace{6.0cm}

\hspace{4mm} \begin{tikzpicture}[overlay]
	\begin{axis}[\bbar{y}{0}, \dateaxisticks ytick={-20, -10, 0, 10, 20},
		xticklabel={`\short{\year}}, clip=false, height=8.0cm,
		yticklabel style={text width=1.4em}, ymax=19, enlarge y limits={0.03},
		legend style={at={(0.64, 0.96)}, fill=white}]
	\rbars
	\sbar{cyan!28!white}{date}{A067RC}{data/pcedecomp.csv}
	\sbar{green!72!blue!90!white}{date}{SAVING}{data/pcedecomp.csv}
	\sbar{blue!92!violet}{date}{OTHER}{data/pcedecomp.csv}
	\stdline{black}{date}{DPCERC}{data/pcedecomp.csv}
	\legend{Income, Saving, Other Outlays, Total};
	\end{axis}
\end{tikzpicture}

\footnotesize{Source: Bureau of Economic Analysis} \hfill \tbllink{pcedecomp.csv}
\end{minipage}
\newpage
\begin{minipage}{0.76\textwidth}
\subsubsection*{Personal Saving}
\small The after-tax income that people do not spend is \href{https://www.bea.gov/index.php/help/glossary/personal-saving}{considered} \textbf{personal saving}, from an economic accounting perspective. Peoples' savings are invested through the financial system and become the current fixed investment or consumption activities of other groups in the economy. Savers generally receive a return from this investment.

\input{text/psavert.txt}
\vspace{1mm}

\normalsize \textbf{Personal Saving Rate}\\
\footnotesize{\textit{saving as a share of after-tax income, percent}}
\vspace{3.1cm}

\hspace{4mm} \begin{tikzpicture}[overlay]
	\begin{axis}[\bbar{y}{0}, \dateaxisticks ytick={0, 10, 20, 30}, enlarge y limits={0.01},
		yticklabel style={text width=1.0em}, ymin=0, ymax=29.8,
		xticklabel={`\short{\year}}, height=5.1cm, width=12.0cm, clip=false]
	\rbars
	\stdline{red}{date}{A072RC}{data/psavert.csv}
	\input{text/psavert_node.txt}
	\end{axis}
\end{tikzpicture}

\footnotesize{Source: Bureau of Economic Analysis} \hfill \tbllink{psavert.csv}
\vspace{2mm}

\subsubsection*{Distribution of Saving}
\small With such a wide distribution of after-tax income, \textbf{saving rates vary massively between households}. Some households dissave and others save more than two typical incomes. Saving by income quintile is calculated using the \href{https://www.bls.gov/cex/}{Consumer Expenditure Survey} (CE) as after-tax income minus spending (other than spending on pensions). The following chart shows the average saving of each group divided by the US median personal income, thus reporting saving, or dissaving, in terms of a typical annual US income. 
\vspace{1mm}

\begin{minipage}{0.55\textwidth}
\normalsize \textbf{Saving Rate by Income Quintile}\\
\footnotesize{\textit{income quintile average saving, as share of \\ US median personal income, percent}}\\
\hspace*{-3mm} \begin{tikzpicture}
	\begin{axis}[\bbar{y}{0}, \shdateaxisticks ytick={-100, -50, 0, 50, 100, 150, 200, 250, 300}, 
		enlarge y limits={0.02}, ymin=-70,
		xticklabel={`\short{\year}}, yticklabel style={text width=1.5em}, 
		height=7.6cm, width=6.9cm, clip=false]
	\rbars
	\stdline{blue!80!cyan!80!black}{date}{DI_SAVE_RT_6}{data/save_dist.csv}
	\stdline{red!80!black}{date}{DI_SAVE_RT_5}{data/save_dist.csv}
	\stdline{green!80!blue!85!black}{date}{DI_SAVE_RT_4}{data/save_dist.csv}
	\stdline{red!20!orange!95!black}{date}{DI_SAVE_RT_3}{data/save_dist.csv}
	\stdline{violet}{date}{DI_SAVE_RT_2}{data/save_dist.csv}
	\input{text/save_dist_nodes.txt}
	\absnode{{1991-06-01}}{264}{\footnotesize \color{blue!80!cyan!80!black}Top 20 Percent\\\footnotesize \color{blue!80!cyan!80!black}of Households\\\footnotesize \color{blue!80!cyan!80!black}by Income}
	\absnode{{2010-10-01}}{100}{\footnotesize \color{red!80!black}Percentiles\\\footnotesize \color{red!80!black}\ \ \ \ 60--80}
	\absnode{{2012-07-01}}{38}{\footnotesize \color{green!80!blue!85!black}40--60}
	\absnode{{2009-04-01}}{-14}{\footnotesize \color{red!20!orange!95!black}20--40}
	\absnode{{2009-10-01}}{-45}{\footnotesize \color{violet}Bottom 20}
	\end{axis}
\end{tikzpicture}\\
\footnotesize{Source: Author's calculations from CE} \hfill \tbllink{save_dist.csv} 
\end{minipage}\hfill
\begin{minipage}{0.38\textwidth}
\small \input{text/save_dist.txt}
\end{minipage}
\end{minipage}
\newpage
\begin{minipage}{0.76\textwidth}
\subsubsection*{Consumer Sentiment}
\small The University of Michigan conducts a monthly \href{http://www.sca.isr.umich.edu/}{survey} of \textbf{consumer sentiment} (see {\color{violet!60!blue}\textbf{---}}). The survey asks about personal finances, business conditions, and buying conditions. An increase in consumer sentiment means individuals feel more confident about economic conditions and are more willing to make large purchases or take on debt. 

\input{text/umichsoc.txt}

\normalsize \textbf{Consumer Sentiment}\\
\footnotesize{\textit{index, 1966=100}}\\
\hspace*{-2mm} \begin{tikzpicture}
	\begin{axis}[\bbar{y}{0}, \dateaxisticks ytick={60, 80, 100, 120}, 
		yticklabel style={text width=1.6em},
		xticklabel={`\short{\year}}, clip=false, height=4.8cm, width=12.2cm]
	\rbars
	\stdline{violet!60!blue}{date}{ICS_ALL}{data/umichsoc.csv}
	\input{text/soc_node.txt}
	\end{axis}
\end{tikzpicture}\\
\footnotesize{Source: University of Michigan} \hfill \tbllink{umichsoc.csv}
\vspace{2.5mm}

\small \input{text/umichsoc_ice.txt}

\normalsize \textbf{Consumer Sentiment Index Components}\\
\footnotesize{\textit{index, 1966=100}}\\
\hspace*{-2mm} \begin{tikzpicture}
    \begin{groupplot}[group style={group size=1 by 2, 
    x descriptions at=edge bottom, vertical sep=0pt,},]
    \nextgroupplot[\bbar{y}{0}, date coordinates in=x, axis line style={draw=black},
    	xticklabels=\empty, xtick style={draw=none},
		xmax={2022-03-15}, max space between ticks=40,	    
		enlarge y limits={0.16},  
		enlarge x limits={0.01}, ytick={40, 60, 80, 100, 120, 140}, 
		yticklabel style={text width=1.6em},
		clip=false, height=5.2cm, width=12.2cm]
			\rbars
			\stdline{blue!50!cyan}{date}{ICC}{data/umichsoc_ice.csv}
			\input{text/icc_node.txt}
			\node[text width=2.8cm, anchor=west, fill=white] at 
				(axis description cs: -0.01, 0.93){\small{Current Conditions}};
    \nextgroupplot[\bbar{y}{0}, \dateaxisticks ytick={40, 60, 80, 100, 120, 140}, 
		yticklabel style={text width=1.6em}, enlarge y limits={0.16},
		xticklabel={`\short{\year}}, clip=false, height=5.2cm, width=12.2cm]
        	\rbars
			\stdline{green!60!black}{date}{ICE}{data/umichsoc_ice.csv}
			\input{text/ice_node.txt}
			\node[text width=3.2cm, anchor=west, fill=white] at 
				(axis description cs: -0.01, 0.93){\small{Expected Conditions}};
    \end{groupplot}
\end{tikzpicture}\\
\footnotesize{Source: University of Michigan} \hfill \tbllink{umichsoc_ice.csv}
\end{minipage}
\newpage
\begin{minipage}{0.76\textwidth}
\subsection*{Household Balance Sheets}
\hypertarget{hhbs}{}

\subsubsection*{Liabilities}

\small \input{text/hhdebt1.txt} 

\input{text/hhdebt2.txt}
\vspace{1mm}

\normalsize \textbf{Household and Nonprofit Debt}\\
\footnotesize{\textit{by type, as share of disposable personal income}}\\
\hspace*{-2mm} \begin{tikzpicture}
	\begin{axis}[\bbar{y}{0}, \dateaxisticks ytick={0, 50, 100, 150}, ymin=3,
		yticklabel style={text width=1.0em}, width=12.6cm,
		xticklabel={`\short{\year}}, clip=false, height=5.6cm, 
		legend cell align={left}, legend style={legend columns=1, at={(0.32, 1.0)}}]
	\rbars
	\sbar{blue!60!violet}{date}{Mortgages}{data/hhdebt.csv}	
	\sbar{magenta}{date}{Consumer Credit}{data/hhdebt.csv}
	\sbar{orange!80!yellow}{date}{Other}{data/hhdebt.csv}	
	\legend{Home Mortgages, Consumer Credit, Other};
	\end{axis}
\end{tikzpicture}\\
\footnotesize{Source: Federal Reserve, Bureau of Economic Analysis} \hfill \tbllink{hhdebt.csv}
\vspace{3mm}

\small \input{text/hhcdebt3.txt} 
\vspace{1mm}

\normalsize \textbf{Mortgages and Consumer Credit}\\
\footnotesize{\textit{share of disposable personal income, percent}}
\vspace{-0.5mm}

\noindent \normalsize Total Change, \small \input{text/cc_dates.txt} \normalsize \hspace{21mm} Consumer Debt Trends \footnotesize \\
\hspace*{-2mm} \begin{tikzpicture}
	\begin{axis}[\barplotnogrid \barylab{1.9cm}{1.5ex}
		width=4.6cm, height=5.4cm, bar width=2.6ex, y=5.0ex, 
		enlarge y limits={abs=0.35cm}, 
		enlarge x limits=0.3, \bbar{x}{0},
		ytick={0,1,2,3,4,5}, nodes near coords, 
    	yticklabels from table={\ccbar}{name},
		yticklabel style={font=\small, xshift=-4pt},
		every axis plot/.append style={bar shift=0pt, fill},
		nodes near coords style={/pgf/number format/.cd,fixed zerofill,
		precision=1, assume math mode}]
		\addplot [draw=none, fill=gray,discard if not={color}{gray}] 
			table [y expr=\coordindex,x index=2,col sep=comma]{data/ccdebtbar.csv};
		\addplot [draw=none,fill=blue!60!cyan,discard if not={color}{blue!60!cyan}] 
			table [y expr=\coordindex, x index=2,col sep=comma]{data/ccdebtbar.csv};
		\addplot [draw=none,fill=red,discard if not={color}{red}] 
			table [y expr=\coordindex, x index=2,col sep=comma]{data/ccdebtbar.csv};
		\addplot [draw=none,fill=green!80!blue,discard if not={color}{green!80!blue}] 
			table [y expr=\coordindex, x index=2,col sep=comma]{data/ccdebtbar.csv};
	\end{axis}
\end{tikzpicture}
\hfill
\begin{tikzpicture}
	\begin{axis}[\bbar{y}{0}, \shdateaxisticks ytick={2, 4, 6, 8, 10}, 
		clip=false, width=6.5cm, height=5.4cm, xticklabel={`\short{\year}}, 
        minor xtick={}, ymin=2, ymax=10, enlarge y limits=0.01, 
        yticklabel style={text width=1.0em}, 
        enlarge x limits={0.01}]
	\rebars
	\stdline{blue!60!cyan}{date}{Auto Loan}{data/hhcdebt.csv}
	\stdline{green!80!blue}{date}{Student Loan}{data/hhcdebt.csv}
	\stdline{red}{date}{Credit Card}{data/hhcdebt.csv}
	\input{text/cc_nodes.txt}
	\absnode{{2015-12-01}}{7.8}{\scriptsize \color{blue!60!cyan}{auto}}
	\absnode{{2002-10-01}}{5.7}{\scriptsize \color{green!80!blue}{student}}
	\absnode{{2013-01-01}}{5.0}{\scriptsize \color{red}{credit card}}
	\end{axis}
\end{tikzpicture}\\
\footnotesize{Source: Federal Reserve Bank of New York and Bureau of Economic Analysis} \hfill \tbllink{hhcdebt.csv}
\end{minipage}
\newpage
\begin{minipage}{0.76\textwidth}
\small Trends in \textbf{household debt} over the past three years, measured in both the US Financial Accounts and the New York Fed Consumer Credit Panel, show mortgages growing in line with income while consumer credit debt falls relative to income. Minor discrepancies between the two data sources arise because the Financial Accounts include debt of nonprofit institutions and the Consumer Credit Panel does not include people without a social security number.

\input{text/hhcdebt4.txt} 

\input{text/hhcdebt5.txt} 
\vspace{1mm}

\normalsize \textbf{Household Debt Outstanding}\\
\footnotesize{\textit{trillions of US dollars \hspace{46mm} share of disposable personal income}\\
\rowcolors{1}{}{black!5} \setlength{\tabcolsep}{3.1pt} \color{black!90}
	{\renewcommand{\arraystretch}{1.54}
		 \begin{tabular}{p{36mm} R{11.8mm} R{11.8mm}  R{8.0mm} R{7.4mm} R{7.4mm} 
		   R{7.4mm} R{7.4mm}}
			 \input data/hhcdebt.tex \hline
		\end{tabular}
	}}\vspace{-1mm}

\footnotesize{Source: Federal Reserve, Federal Reserve Bank of New York, Bureau of Economic Analysis}
\end{minipage}
\newpage
\begin{minipage}{0.76\textwidth}
\subsubsection*{Consumer Credit}

\small The Federal Reserve also \href{https://www.federalreserve.gov/releases/g19/current/default.htm}{report} consumer credit on a monthly basis. \input{text/cc_dpi.txt}

\input{text/cc_dpi_2.txt}
\vspace{1mm}

\normalsize \textbf{Consumer Credit}\\
\footnotesize{\textit{share of disposable personal income, percent}}\\
\hspace*{-2mm} \begin{tikzpicture}
	\begin{axis}[\bbar{y}{0}, \dateaxisticks ytick={16, 18, 20, 22, 24, 26}, 
		enlarge y limits={0.1}, height=6.0cm, width=12.2cm, 
		yticklabel style={text width=1.0em}, 
		xticklabel={`\short{\year}}, clip=false]
	\rbars
	\thickline{lime!55!green}{date}{Consumer Credit}{data/hhcdebt.csv}
	\thinline{green!60!blue!80!black}{date}{value}{data/cc_dpi_monthly.csv}
	\input{text/cc_dpi_monthly_node.txt}
	\absnode{{1997-01-01}}{18}{\color{lime!55!green}FRBNY}
	\absnode{{2010-07-01}}{25.5}{\color{green!60!blue!80!black}Fed G.19}
	\end{axis}
\end{tikzpicture}\\
\footnotesize{Source: Federal Reserve; FRBNY} \hfill \tbllink{cc_dpi_monthly.csv} \tbllink{hhcdebt.csv}
\end{minipage}
\vspace{5mm}

\subsubsection*{Financial Obligations}
\begin{minipage}{0.29\textwidth}
\small Payments to service debt, along with rent, auto lease payments, homeowner's insurance, and property tax are considered financial obligations. The Federal Reserve \href{https://www.federalreserve.gov/releases/housedebt/default.htm}{report} financial obligations as a share of disposable personal income. The ratio of debt service payments and financial obligations to income gives insight into the current aggregate financial burden facing households. 

\input{text/for.txt} The financial obligations ratio peaked at 18.1 percent in 2007 Q4, during the housing bubble.
\end{minipage}\hspace{8mm}
\begin{minipage}{0.415\textwidth}
\normalsize \textbf{Financial Obligations and}\\
\normalsize \textbf{Debt Service Ratios}\\
\footnotesize{\textit{share of disposable personal income, percent}}\\
\hspace*{-2mm} \begin{tikzpicture}
	\begin{axis}[\shdateaxisticks ytick={6, 8, 10, 12, 14, 16, 18, 20}, 
		enlarge y limits={0.1}, yticklabel style={text width=1.0em}, 
		xticklabel={`\short{\year}}, height=7.5cm, width=7.0cm, clip=false]
	\rbars
	\stdline{blue!80!black}{date}{Financial obligations ratio}{data/for.csv}
	\stdline{red}{date}{Debt service ratio}{data/for.csv}
	\input{text/for_node.txt}
	\stdnode{4.15cm}{4.9cm}{\color{blue!80!black} Financial\\ \color{blue!80!black} Obligations}
	\input{text/dsr_node.txt}
	\stdnode{0.95cm}{1.0cm}{\color{red} Debt\\ \color{red} Service}
	\end{axis}
\end{tikzpicture}\\
\footnotesize{Source: Federal Reserve} \hfill \tbllink{for.csv} 
\end{minipage}
\newpage
\begin{minipage}{0.76\textwidth}
\subsubsection*{Assets}

\small \input{text/hhasset1.txt} 

\input{text/hhasset2.txt} 

\input{text/hhasset3.txt}
\vspace{2mm}

\normalsize \textbf{Selected Household and Nonprofit Assets}\\
\footnotesize{\textit{share of GDP}}\\
\hspace{4mm} \begin{tikzpicture}
	\begin{axis}[\bbar{y}{0}, \dateaxisticks ytick={0, 100, 200, 300, 400, 500},
		 ymin=20, height=5.2cm, reverse legend, 	
		 yticklabel style={text width=1.2em}, 
		 xticklabel={`\short{\year}}, clip=false, legend cell align={left},
		 legend style={legend columns=1, at={(0.41, 1.0)}}]
	\rbars
	\sbar{green!60!teal}{date}{HRE}{data/hhassetsgdp.csv}
	\sbar{cyan!40!white}{date}{DEP}{data/hhassetsgdp.csv}	
	\sbar{blue!55!cyan}{date}{DSL}{data/hhassetsgdp.csv}
	\sbar{blue!65!black}{date}{TEQ}{data/hhassetsgdp.csv}
	\legend{Real Estate, Deposits, Debt Securities \& Loans, Business Equity};
	\end{axis}
\end{tikzpicture}\\
\footnotesize{Source: Federal Reserve, Bureau of Economic Analysis} \hfill \tbllink{hhassetsgdp.csv}
\vspace{3mm}

\small \input{text/hhasset4.txt} 
\vspace{1mm}

\normalsize \textbf{Contributions to Real Growth in Household and Nonprofit Assets}\\
\footnotesize{\textit{contribution to one-year percent change in assets, adjusted by PCE price deflator}}\\
\hspace*{-3mm} \begin{tikzpicture}
	\begin{axis}[\bbar{y}{0}, \dateaxisticks ytick={-15, -10, -5, 0, 5, 10, 15, 20, 25}, 
		height=6.2cm, enlarge y limits={0.02}, 
		xticklabel={`\short{\year}}, yticklabel style={text width=1.4em},
			clip=false, legend cell align={left},
		legend style={legend columns=2, at={(0.51, 1.01)}, fill=white}]
	\rbars
	\sbar{cyan!40!white}{date}{DEP}{data/hh_asset_growth.csv}	
	\sbar{blue!55!cyan}{date}{DSL}{data/hh_asset_growth.csv}
	\sbar{green!60!teal}{date}{HRE}{data/hh_asset_growth.csv}
	\sbar{blue!65!black}{date}{TEQ}{data/hh_asset_growth.csv}
	\legend{Deposits, Debt Securities \& Loans, Real Estate, Business Equity};
	\end{axis}
\end{tikzpicture}\\
\footnotesize{Source: Federal Reserve, Bureau of Economic Analysis} \hfill \tbllink{hh_asset_growth.csv}
\end{minipage}
\newpage
\begin{minipage}{0.77\textwidth}
\normalsize \textbf{Household and Nonprofit Assets}\\
\footnotesize{\textit{various measures: \hspace{25mm} trillions of USD \hspace{4mm} share of GDP \hspace{4mm} real annual growth rate}}\\
\rowcolors{1}{}{black!5} \setlength{\tabcolsep}{3.1pt} \color{black!90}
	{\renewcommand{\arraystretch}{1.55}
		 \hspace*{-2mm} \begin{tabular}{p{2mm} p{46mm} R{11mm} R{9.5mm} R{9.5mm} R{8mm} 
		   R{8mm} R{8mm} }
			 \input data/hhasset.tex \hline
	\end{tabular}}\vspace{-2mm}

\footnotesize{Source: Federal Reserve, Bureau of Economic Analysis}
\end{minipage}
\vspace{7mm}

\begin{minipage}{0.76\textwidth}
\subsubsection*{Return on Assets}
\small Asset prices rising faster than income can be viewed as a decrease in the expected rate of return on total household assets. This can be measured by disposable income as a share of household assets. \input{text/dpishta.txt} 
\vspace{2mm}

\normalsize \textbf{Return on Household Assets}\\
\footnotesize{\textit{disposable personal income as share of household and nonprofit total assets, percent}}\\
\hspace*{-2mm} \begin{tikzpicture}
	\begin{axis}[\bbar{y}{0}, \dateaxisticks ytick={12, 14, 16}, 
		enlarge y limits={0.08}, height=5.4cm, width=12.3cm, 
		yticklabel style={text width=1.0em}, 
		xticklabel={`\short{\year}}, clip=false]
	\rbars
	\stdline{red}{date}{DPIsh}{data/dpish.csv}
	\input{text/dpish_node.txt}
	\end{axis}
\end{tikzpicture}\\
\footnotesize{Source: Federal Reserve, Bureau of Economic Analysis} \hfill \tbllink{dpish.csv}
\end{minipage}
\newpage
\begin{minipage}{0.76\textwidth}
\subsubsection*{Wealth/Net Worth}
\small Wealth \href{https://www.federalreserve.gov/econres/notes/feds-notes/disparities-in-wealth-by-race-and-ethnicity-in-the-2019-survey-of-consumer-finances-20200928.htm}{varies} substantially by race and ethnicity. Net worth, measured as total assets minus total liabilities, captures the overall financial position--or wealth--of families. In 2019, white non-Hispanic families' average net worth was \$980,550, compared to \$142,330 for black non-Hispanic families, and \$165,540 for Hispanic families of any race. Additionally, the wealth of a typical (median) family is much lower than the average wealth of families, the result of a concentration of wealth among the wealthiest families.

\normalsize \textbf{Racial Wealth Gap}\\
\footnotesize{\textit{net worth by race/ethnicity, thousands of US dollars, 2019}}\\
\hspace*{-3mm} \begin{tikzpicture}
\begin{axis}[clip=false,
    symbolic x coords={White, Black, Hispanic, Other},
    xtick=data, ymin=0, width=6.7cm, height=5cm, bar width=20pt,
    yticklabel style={text width=1.5em}, 
    enlarge x limits={0.2}, nodes near coords, 
    nodes near coords style={/pgf/number format/.cd, fixed zerofill,
			precision=1, assume math mode},
	xticklabel style={black, font=\small}]
    \addplot[ybar,fill=blue, draw=blue] coordinates {
        (White,980.6)
        (Black,142.3)
        (Hispanic,165.54)
        (Other,656.6)};
    \stdnode{1.9cm}{3.9cm}{\large Mean};
\end{axis}
\end{tikzpicture} \hfill
\begin{tikzpicture}
\begin{axis}[clip=false,
    symbolic x coords={White, Black, Hispanic, Other},
    xtick=data, ymin=0, width=6.7cm, height=5cm, bar width=20pt,
    yticklabel style={text width=1.5em}, 
    enlarge x limits={0.2}, nodes near coords, 
    nodes near coords style={/pgf/number format/.cd, fixed zerofill,
			precision=1, assume math mode},
	xticklabel style={black, font=\small}]
    \addplot[ybar,fill=blue, draw=blue] coordinates {
        (White,189.1)
        (Black,24.1)
        (Hispanic,36.1)
        (Other,74.5)};
    \stdnode{1.9cm}{3.9cm}{\large Median};
\end{axis}
\end{tikzpicture}\\
\footnotesize{Source: Federal Reserve, Survey of Consumer Finances}
\vspace{8mm}

\small The market value of household assets has risen much faster than the total level of household debt, resulting in a substantial increase in aggregate net worth. \input{text/nw1.txt} 

\input{text/nw2.txt}
\vspace{2mm}

\normalsize \textbf{Net Worth and After-Tax Income Growth}\\
\footnotesize{\textit{one-year percent change, adjusted by PCE price deflator}}
\vspace{4.2cm}

\hspace{4mm} \begin{tikzpicture}[overlay]
	\begin{axis}[\bbar{y}{0}, \dateaxisticks ytick={-10, 0, 10, 20}, height=6.2cm, 
		ymax = 18.2, 
		legend cell align={left}, reverse legend,
		xticklabel={`\short{\year}}, yticklabel style={text width=1.5em}, clip=false, 
		legend style={at={(0.91, 1.21)}, legend columns=1}]
	\rbars 
	\sbaral{cyan!40!white}{date}{NW}{data/rdpi_nw.csv}
	\stdline{blue!50!violet}{date}{DPI}{data/rdpi_nw.csv}
	\legend{ \ Net Worth, \ After-Tax Income};
	\end{axis}
\end{tikzpicture}
\vspace{1mm}

\footnotesize{Source: Federal Reserve, Bureau of Economic Analysis} \hfill \tbllink{rdpi_nw.csv}
\end{minipage}
\newpage
\begin{minipage}{0.35\textwidth}
\small \input{text/dfa_nw_sh.txt}
\end{minipage}\hspace{8mm}
\begin{minipage}{0.35\textwidth}
\normalsize \textbf{Share of US Wealth}\\
\footnotesize{\textit{share of total net worth, by percentile, percent}}\\
\hspace*{-2mm} \begin{tikzpicture}
	\begin{axis}[\bbar{y}{0}, \shdateaxisticks clip=false,
		ymin=0.6, xticklabel={`\short{\year}}, height=6.6cm, width=6.2cm]
	\rbars
	\thickline{red!50!orange}{date}{Top1}{data/dfa_nw_sh.csv}
	\thickline{blue!80!cyan}{date}{Bottom50}{data/dfa_nw_sh.csv}
	\thickline{cyan!80!white}{date}{Next40}{data/dfa_nw_sh.csv}
	\thickline{orange!50!yellow}{date}{Next9}{data/dfa_nw_sh.csv}
	\absnode{{2009-05-01}}{5.5}{\footnotesize \color{blue!80!cyan} Bottom 50\%}
	\absnode{{1993-11-01}}{26.5}{\footnotesize \color{red!50!orange} Top 1\%}
	\absnode{{2012-01-01}}{42}{\footnotesize \color{orange!50!yellow} Next 9\%}
	\absnode{{2012-08-01}}{28}{\footnotesize \color{cyan!80!white} Next 40\%}
	\input{text/dfa_nw_sh_nodes.txt}
	\end{axis}
\end{tikzpicture}\\
\footnotesize{Source: Federal Reserve} \hfill \tbllink{businvsh.csv}
\end{minipage}
\vspace{5mm}

\begin{minipage}{0.76\textwidth}
\subsubsection*{Wealth and Income}
\small Wealth, when measured as assets minus liabilities, can be a source of income, but does not correspond perfectly to income. For example, early-career professionals with student debt may have a negative net worth and a high income. That said, data on family income by wealth percentile show that income tends to increase with wealth. 

Additionally, the before-tax income of the wealthiest ten percent of families (see {\color{blue!80!cyan!80!black}\textbf{---}}) has \href{https://www.federalreserve.gov/econres/scf/dataviz/scf/chart/#series:Before_Tax_Income;demographic:nwcat;population:1,2,3,4,5;units:median}{increased} substantially more than the income of other groups. The top ten percent of families by wealth, percentiles 90 to 100 with a mean wealth of \$5.7 million and a median wealth of \$2.6 million in 2019, \input{text/income_wealthg1.txt} 
\end{minipage}
\vspace{1mm}

\begin{minipage}{0.4\textwidth}
\normalsize \textbf{Before-Tax Income by Wealth Percentile}\\
\footnotesize{\textit{median annual income, thousands of 2019 USD}}\\
\hspace*{-2mm} \begin{tikzpicture}
	\begin{axis}[\bbar{y}{0}, \shdateaxisticks ytick={0, 50, 100, 150, 200, 250}, 
		enlarge y limits={0.1}, ymin=20,
		xticklabel={`\short{\year}}, yticklabel style={text width=1.3em}, 
		height=8.7cm, width=6.8cm, clip=false]
	\rbars
	\stdline{blue!80!cyan!80!black}{date}{90-100}{data/scfwealthmedian.csv}
	\stdline{red!80!black}{date}{75-89.9}{data/scfwealthmedian.csv}
	\stdline{green!80!blue!85!black}{date}{50-74.9}{data/scfwealthmedian.csv}
	\stdline{red!20!orange!95!black}{date}{25-49.9}{data/scfwealthmedian.csv}
	\stdline{violet}{date}{Less than 25}{data/scfwealthmedian.csv}
	\input{text/scfwealthmedian_node1.txt}
	\input{text/scfwealthmedian_node2.txt}
	\input{text/scfwealthmedian_node3.txt}
	\input{text/scfwealthmedian_node4.txt}
	\input{text/scfwealthmedian_node5.txt}
	\stdnode{1.1cm}{4.2cm}{\footnotesize \color{blue!80!cyan!80!black}90-100}
	\stdnode{4.35cm}{2.65cm}{\footnotesize \color{red!80!black}75-89.9}
	\stdnode{4.2cm}{1.75cm}{\footnotesize \color{green!80!blue!85!black}50-74.9}
	\stdnode{1.0cm}{0.7cm}{\footnotesize \color{red!20!orange!95!black}25-49.9}
	\stdnode{4.2cm}{0.15cm}{\footnotesize \color{violet}Less than 25}
	\end{axis}
\end{tikzpicture}\\
\footnotesize{Source: Federal Reserve (SCF)} \hfill \tbllink{scfwealthmedian.csv} 
\end{minipage}\hspace{5.0mm}
\begin{minipage}{0.32\textwidth}
\small \input{text/income_wealthg2.txt}
\end{minipage}
\newpage
\begin{minipage}{0.76\textwidth}
\small The US history of slavery and inter-generational white wealth show up in current data from the Survey of Consumer Finances. The racial wealth gap discussed above cannot be explained by borrowing or access to credit. White families have substantially more financial assets including stocks and are much more likely to receive inheritance. Income for black families is also substantially lower--about half of white family income. Persistent structural inequalities are seen in income data, but are also evident from measures of wealth and assets.
\vspace{1mm}

\normalsize \textbf{Measures of Wealth and Income by Race or Ethnicity}\\
\footnotesize{\textit{by family, mean, thousands of 2019 USD}}

\Large{\color{red!80!white}\textbf{---}} \footnotesize{White, non-Hispanic} \ \ \ \Large{\color{green!90!blue}\textbf{---}} \footnotesize{Black, non-Hispanic} \ \ \ \Large{\color{magenta!50!violet}\textbf{---}} \footnotesize{Hispanic} \ \ \ \Large{\color{blue}\textbf{---}} \footnotesize{Other}
\end{minipage}
\vspace{-1mm}

\begin{minipage}{0.26\textwidth}
\normalsize Net Worth\\
\hspace*{-2mm} \begin{tikzpicture}
	\begin{axis}[\bbar{y}{0}, \shdateaxisticks ytick={0, 250, 500, 750, 1000}, ymin=50,
		yticklabel style={text width=1.5em}, 
		xticklabel={`\short{\year}}, height=4.8cm, width=5.0cm, clip=false]
	\rbars
	\stdline{red!80!white}{date}{White}{data/scfracemean.csv}
	\stdline{blue}{date}{Other}{data/scfracemean.csv}
	\stdline{magenta!50!violet}{date}{Hispanic}{data/scfracemean.csv}
	\stdline{green!90!blue}{date}{Black}{data/scfracemean.csv}
	\end{axis}
\end{tikzpicture}
\end{minipage}\hspace{4mm}
\begin{minipage}{0.26\textwidth}
\normalsize Financial Assets\\
\hspace*{-2mm} \begin{tikzpicture}
	\begin{axis}[\bbar{y}{0}, \shdateaxisticks ytick={0, 100, 200, 300, 400, 500}, ymin=25,
		yticklabel style={text width=1.2em}, 
		xticklabel={`\short{\year}}, height=4.8cm, width=4.95cm, clip=false]
	\rbars
	\stdline{red!80!white}{date}{White}{data/scfracemean0.csv}
	\stdline{blue}{date}{Other}{data/scfracemean0.csv}
	\stdline{magenta!50!violet}{date}{Hispanic}{data/scfracemean0.csv}
	\stdline{green!90!blue}{date}{Black}{data/scfracemean0.csv}
	\end{axis}
\end{tikzpicture}
\end{minipage}\hspace{4mm}
\begin{minipage}{0.26\textwidth}
\normalsize Pre-Tax Income\\
\hspace*{-2mm} \begin{tikzpicture}
	\begin{axis}[\bbar{y}{0}, \shdateaxisticks ytick={0, 25, 50, 75, 100, 125}, ymin=6,
		yticklabel style={text width=1.2em}, 
		xticklabel={`\short{\year}}, height=4.8cm, width=4.95cm, clip=false]
	\rbars
	\stdline{red!80!white}{date}{White}{data/scfracemean1.csv}
	\stdline{blue}{date}{Other}{data/scfracemean1.csv}
	\stdline{magenta!50!violet}{date}{Hispanic}{data/scfracemean1.csv}
	\stdline{green!90!blue}{date}{Black}{data/scfracemean1.csv}
	\end{axis}
\end{tikzpicture}
\end{minipage}\\
\begin{minipage}{0.76\textwidth}
\footnotesize{Source: Federal Reserve} \hfill \tbllink{scfracemean.csv} \tbllink{scfracemean0.csv} \tbllink{scfracemean1.csv}
\end{minipage}
\vspace{4mm}

\begin{minipage}{0.76\textwidth}
\small \input{text/scfracestocks.txt} 
\end{minipage}
\vspace{2mm}

\begin{minipage}{0.36\textwidth}
\normalsize \textbf{Stock Holdings}\\
\footnotesize{\textit{mean value, thousands of 2019 USD}}\\
\hspace*{-2mm} \begin{tikzpicture}
	\begin{axis}[\bbar{y}{0}, \shdateaxisticks ytick={0, 100, 200, 300, 400, 500}, ymin=20,
		yticklabel style={text width=1.2em}, 
		xticklabel={`\short{\year}}, height=5.2cm, width=6.5cm, clip=false]
	\rbars
	\stdline{red!80!white}{date}{White}{data/scfracemean2.csv}
	\stdline{blue}{date}{Other}{data/scfracemean2.csv}
	\stdline{magenta!50!violet}{date}{Hispanic}{data/scfracemean2.csv}
	\stdline{green!90!blue}{date}{Black}{data/scfracemean2.csv}
	\end{axis}
\end{tikzpicture}\\
\footnotesize{Source: Federal Reserve} \hfill \tbllink{scfracemean2.csv} 
\end{minipage}\hspace{7mm}
\begin{minipage}{0.36\textwidth}
\normalsize \textbf{Stock Holdings}\\
\footnotesize{\textit{share of families with holdings, percent}}\\
\hspace*{-2mm} \begin{tikzpicture}
	\begin{axis}[\bbar{y}{0}, \shdateaxisticks ytick={0, 10, 20, 30, 40, 50, 60, 70}, ymin=3,
		yticklabel style={text width=1.0em}, 
		xticklabel={`\short{\year}}, height=5.2cm, width=6.5cm, clip=false]
	\rbars
	\stdline{red!80!white}{date}{White}{data/scfracehave.csv}
	\stdline{blue}{date}{Other}{data/scfracehave.csv}
	\stdline{magenta!50!violet}{date}{Hispanic}{data/scfracehave.csv}
	\stdline{green!90!blue}{date}{Black}{data/scfracehave.csv}
	\end{axis}
\end{tikzpicture}\\
\footnotesize{Source: Federal Reserve} \hfill \tbllink{scfracehave.csv} 
\end{minipage}
\newpage
\begin{minipage}{0.76\textwidth}
\small \textbf{Changes in net worth} come mostly from changes in the value of assets and other holding gains (see\cbox{cyan!40!white}), but also from saving of income, and decisions about how much to spend. Each period, households spend a portion of their income. The saved portion of income, after consumer spending and spending to cover depreciation of households assets, is considered net investment and increases net worth. Since 1989, household net investment has averaged 10 percent of disposable personal income. 

In the chart, income invested at the historical-average rate (see\cbox{magenta!85!white}) is visually separated from investment that is above or below trend (see\cbox{violet}). The separation distinguishes changes in disposable personal income from changes in decisions about how to use that income. Separately, changes in data sources or from natural disasters are identified as other volume changes (see\cbox{cyan!70!blue}).
\vspace{2mm}

\normalsize{\textbf{Net Worth Growth}}\\
\footnotesize{\textit{contribution to one-year percent change in net worth, nominal}}
\vspace{4.8cm}

\hspace{4mm} \begin{tikzpicture}[overlay]
	\begin{axis}[\bbar{y}{0}, \dateaxisticks ytick={-15, -10, -5, 0, 5, 10, 15, 20},
		height=6.8cm, xticklabel={`\short{\year}}, legend cell align={left},
		yticklabel style={text width=1.5em}, clip=false, ymax=18.1, reverse legend,
		legend style={at={(0.55, 0.24)}, fill=white, legend columns=2}]
	\rbars
	\sbar{violet}{date}{INV}{data/nw_gr.csv}
	\sbar{magenta!85!white}{date}{INC}{data/nw_gr.csv}
	\sbar{cyan!70!blue}{date}{OVC}{data/nw_gr.csv}		
	\sbar{cyan!40!white}{date}{RV}{data/nw_gr.csv}
	\legend{ \ Rate of Investment, \ Income, \ Other Volume Changes, \ Holding Gains};
	\end{axis}
\end{tikzpicture}

\footnotesize{Source: Federal Reserve, Bureau of Economic Analysis} \hfill \tbllink{nw_gr.csv}
\vspace{3mm}

\small \input{text/nwcontrib.txt}
\end{minipage}
\newpage
\begin{minipage}{0.76\textwidth}
\subsubsection*{Homeownership}
\small The \textbf{homeownership rate} measures the percent of occupied housing units that are owner-occupied, as opposed to rented. In 2004, near to the peak of the housing bubble, the overall homeownership rate reached 69.2 percent. \input{text/homeown.txt}
\end{minipage}

\begin{minipage}{0.41\textwidth}
\normalsize \textbf{Homeownership Rate}\\
\footnotesize{\textit{share of households occupied by owners, percent}}\\
\hspace*{-2mm} \begin{tikzpicture}
	\begin{axis}[\shdateaxisticks ytick={40, 50, 60, 70}, 
		enlarge y limits={0.03}, ymin=40, xmin={1991-01-01}, 
		yticklabel style={text width=1.0em},
		xticklabel={`\short{\year}}, height=8.6cm, width=7.0cm, clip=false]
	\rbars
	\stdline{violet!60!magenta}{date}{VALUE}{data/homeown.csv}
	\stdline{blue!80!cyan!80!black}{date}{White}{data/homeown2.csv}
	\stdline{red!80!black}{date}{Other}{data/homeown2.csv}
	\stdline{green!80!blue!85!black}{date}{Hispanic}{data/homeown2.csv}
	\stdline{red!20!orange!95!black}{date}{Black}{data/homeown2.csv}
	\input{text/homeown_node.txt}
	\input{text/homeown_white_node.txt}
	\input{text/homeown_other_node.txt}
	\input{text/homeown_hispanic_node.txt}
	\input{text/homeown_black_node.txt}
	\stdnode{4.3cm}{0.4cm}{\footnotesize \color{red!20!orange!95!black}Black}
	\stdnode{4.4cm}{6.6cm}{\footnotesize \color{blue!80!cyan!80!black}White}
	\stdnode{4.7cm}{3.4cm}{\footnotesize \color{red!80!black}Other}
	\stdnode{0.9cm}{5.1cm}{\footnotesize \color{violet!60!magenta}Total}
	\stdnode{4.3cm}{1.92cm}{\footnotesize \color{green!80!blue!85!black}Hispanic}
	\end{axis}
\end{tikzpicture}\\
\footnotesize{Source: Census Bureau} \hspace{26mm} \tbllink{homeown2.csv} 

\end{minipage}\hspace{5.0mm}
\begin{minipage}{0.315\textwidth}
\small Census data also show large differences in homeownership rates by race and ethnicity. \input{text/homeown2.txt}

\input{text/homeown3.txt}
\end{minipage}\hspace{5mm}
\vspace{2mm}

\begin{minipage}{0.76\textwidth}
\small Use caution when interpreting homeownership rates during the COVID-19 pandemic. Individuals who are renters or homeowners are captured by the measure, but when an individual moves in with family and stops being a head of household, they are dropped from the measure. Therefore when renters move in with family the homeownership rate increases. The 2020 spike in homeownership rates reflects renters moving in with family. 

As seen during the collapse of the housing bubble, it is possible for someone to be a homeowner but have no equity in their home, for example if the market price of the home falls below the principal remaining on the mortgage. Trends in owners' equity as a share of the market value of real estate show substantial improvement since the lows following the collapse of the housing bubble, when millions of homes were in foreclosure. 
\end{minipage}
\vspace{1mm}

\begin{minipage}{0.29\textwidth}
\small \input{text/homeeq.txt}
\end{minipage}\hspace{7mm}
\begin{minipage}{0.46\textwidth}
\normalsize \textbf{Owners' Equity Share of Real Estate}\\
\footnotesize{\textit{owners' equity as percent of real estate}}\\
\hspace*{-2mm} \begin{tikzpicture}
	\begin{axis}[\shdateaxisticks ytick={45, 50, 55, 60, 65, 70}, 
		xticklabel={`\short{\year}}, 
		yticklabel style={text width=1.0em}, height=4.4cm, 
		width=7.0cm, clip=false]
	\rbars
	\stdline{blue!80!violet}{date}{VALUE}{data/homeeq.csv}
	\input{text/homeeq_node.txt}
	\end{axis}
\end{tikzpicture}\\
\footnotesize{Source: Federal Reserve} \hspace{30mm} \tbllink{homeeq.csv} 

\end{minipage}
\newpage
\begin{minipage}{0.76\textwidth}
\subsubsection*{Housing Construction}
\small The Census Bureau \href{https://www.census.gov/construction/bps/}{tracks} the issuance of \textbf{new residential building permits}, which offer insight into planned residential construction. \input{text/permits.txt}
\vspace{1mm}

\normalsize{\textbf{Residential Construction Permits}}\\
\footnotesize{\textit{building permits issued, seasonally-adjusted annual rate, in thousands}}\\
\hspace*{-2mm} \begin{tikzpicture}
	\begin{axis}[\bbar{y}{0}, \dateaxisticks ytick={500, 1000, 1500, 2000, 2500}, 
		xticklabel={`\short{\year}}, clip=false, width=11.9cm, 
		yticklabel style={text width=2.2em}]
	\rbars		
	\stdline{blue!70!green}{date}{APERMITS}{data/permits_total.csv}
	\input{text/permits_node.txt}
	\end{axis}
\end{tikzpicture}\\
\footnotesize{Source: Census Bureau} \hfill \tbllink{permits_total.csv}
\vspace{2mm}

\small In addition to data on permits, the Census Bureau also report how many residential construction projects are started and completed. Not all permitted projects are built and completion can be affected by economic conditions. \input{text/completions.txt}
\vspace{1mm}

\normalsize{\textbf{Residential Construction Completions}}\\
\footnotesize{\textit{completed units, seasonally-adjusted annual rate, in thousands}}\\
\hspace*{-2mm} \begin{tikzpicture}
	\begin{axis}[\bbar{y}{0}, \dateaxisticks ytick={500, 1000, 1500, 2000, 2500}, 
		xticklabel={`\short{\year}}, clip=false, width=11.9cm, 
		yticklabel style={text width=2.2em}]
	\rbars	
	\stdline{blue!30!cyan}{date}{ACOMPLETIONS}{data/permits_total.csv}	
	\input{text/completions_node.txt}
	\end{axis}
\end{tikzpicture}\\
\footnotesize{Source: Census Bureau} \hfill \tbllink{permits_total.csv}
\vspace{2mm}

\small \input{text/permits_single.txt}
\vspace{1mm}

\normalsize \textbf{Single-family units}\\
\footnotesize{\textit{permits and completions, seasonally-adjusted annual rate, in thousands}}\\
\hspace*{-2mm} \begin{tikzpicture}
	\begin{axis}[\dateaxisticks ytick={500, 1000, 1500}, 
		yticklabel style={text width=2.1em}, legend style={at={(0.4, 1.0)}},
		xticklabel={`\short{\year}}, width=12.4cm]
	\rbars
	\thinline{blue!40!cyan!80!white}{date}{ACOMPLETIONS}{data/permits_single.csv}
	\stdline{blue!60!green}{date}{APERMITS}{data/permits_single.csv}
	\legend{Completions, Permits}
	\end{axis}
\end{tikzpicture}\\
\footnotesize{Source: Census Bureau} \hfill \tbllink{permits_single.csv}
\end{minipage}
\newpage
\begin{minipage}{0.76\textwidth}
\small \input{text/permits_multi.txt}
\vspace{1mm}

\normalsize \textbf{Multi-family units}\\
\footnotesize{\textit{permits and completions, in thousands}}\\
\hspace*{-2mm} \begin{tikzpicture}
	\begin{axis}[\bbar{y}{0}, \dateaxisticks ytick={0, 200, 400, 600, 800, 1000}, 
		yticklabel style={text width=2.1em}, legend style={at={(0.4, 1.0)}}, ymin=30,
		xticklabel={`\short{\year}}, height=6.2cm, width=12.4cm]
	\rbars
	\thinline{blue!40!cyan!80!white}{date}{ACOMPLETIONS}{data/permits_multi.csv}
	\stdline{blue!60!green}{date}{APERMITS}{data/permits_multi.csv}
	\legend{Completions, Permits}
	\end{axis}
\end{tikzpicture}\\
\footnotesize{Source: Census Bureau} \hfill \tbllink{permits_multi.csv}
\vspace{6mm}

\subsubsection*{Housing Sales}
\small \input{text/nhs.txt}
\vspace{1mm}

\normalsize{\textbf{New Home Sales}}\\
\footnotesize{\textit{new single family homes sold, in thousands, seasonally adjusted annual rate}}\\
\hspace*{-2mm} \begin{tikzpicture}
	\begin{axis}[\dateaxisticks ytick={200, 400, 600, 800, 1000, 1200, 1400}, 
		xticklabel={`\short{\year}}, clip=false, height=5.8cm, width=12.0cm, 
		yticklabel style={text width=2.2em}, enlarge y limits={0.05},]
	\rbars
	\stdline{green!80!blue}{date}{VALUE}{data/nhs.csv}	
	\input{text/nhs_node.txt}
	\end{axis}
\end{tikzpicture}\\
\footnotesize{Source: Census Bureau} \hfill \tbllink{nhs.csv}
\end{minipage}
\newpage
\subsubsection*{Housing Prices}
\begin{minipage}{0.76\textwidth}
\small The Federal Housing Finance Agency (FHFA) \textbf{housing price index} \href{https://www.fhfa.gov/DataTools/Downloads/Pages/House-Price-Index-Datasets.aspx}{measures} changes in the price of the same home. The seasonally-adjusted index \input{text/fhfa.txt}

\normalsize \textbf{House Price Index}\\
\footnotesize{\textit{seasonally adjusted index, one-year percent change and three-month percent change}}\\
\hspace*{-2mm} \begin{tikzpicture}
	\begin{axis}[\bbar{y}{0}, \dateaxisticks ytick={-10, -5, 0, 5, 10, 15, 20, 25}, 
		yticklabel style={text width=1.5em}, clip=false,
		xticklabel={`\short{\year}}, height=5.0cm, width=12.1cm, 
		legend style={at={(0.76, 1.0)}}]
	\rebars
	\stdline{purple}{date}{USA3M}{data/hpi.csv}
	\thickline{blue!35!cyan}{date}{USA}{data/hpi.csv}		
	\input{text/hpi_nodes.txt}
	\legend{Three-month change (3M/3M), One-year change};
	\end{axis}
\end{tikzpicture}\\
\footnotesize{Source: Federal Housing Finance Agency} \hfill \tbllink{hpi.csv}
\vspace{1mm}

\normalsize{\textbf{Housing Price Growth}}\\
\footnotesize{\textit{seasonally adjusted, one-year percent change}}\\
\hspace*{-3mm} \noindent \rowcolors{1}{}{black!5} \setlength{\tabcolsep}{2.1pt} \color{black!90}
		{\renewcommand{\arraystretch}{1.45}
		 \begin{tabular}{p{25mm} R{8mm} R{8mm} R{8mm} R{8mm} R{8mm} R{8mm} R{8mm} R{12.0mm} R{10.8mm}}
			 \input data/hpi.tex \hline
		\end{tabular}}
		
	\vspace{-2mm}
\footnotesize{Source: Federal Housing Finance Agency} \hfill \tbllink{hpi.csv}
\end{minipage}
\vspace{3mm}

\begin{minipage}{0.33\textwidth}
\small The purchase price of housing should move with the rental price. When housing prices exceed the rental equivalent, it may suggest that housing is overvalued. 

During the housing bubble that caused the great recession, housing prices reached more than 40 percent above the rental equivalent. \input{text/hp_rent.txt}
\end{minipage}\hspace{8mm}
\begin{minipage}{0.37\textwidth}
\normalsize \textbf{Housing Price to Rent Ratio}\\
\footnotesize{\textit{index, 1991=1}}\\
\hspace*{-2mm} \begin{tikzpicture}
	\begin{axis}[\shdateaxisticks ytick={0.9, 1.0, 1.1, 1.2, 1.3, 1.4, 1.5}, 
		xticklabel={`\short{\year}}, \bbar{y}{1},
		yticklabel style={text width=1.3em}, height=4.6cm, 
		width=6.8cm, clip=false]
	\rebars
	\stdline{red}{date}{VALUE}{data/hpi_rent_ratio.csv}
	\end{axis}
\end{tikzpicture}\\
\footnotesize{Source: FHFA; BLS} \hfill \tbllink{hpi_rent_ratio.csv} 
\end{minipage}
\newpage
\begin{minipage}{0.76\textwidth}
\subsection*{Poverty}
\hypertarget{hhpov}{}
\small \input{text/poverty_main.txt}

For purposes of program eligibility and economic data, poverty is defined by having income below a certain threshold. The processes for calculating poverty vary, with the Official Poverty Measure (OPM) based on three times a price-adjusted 1963 minimal food budget, and the SPM based on food, shelter, clothing, and utilities costs and additionally capturing program benefits and taxes, along with other adjustments.
\end{minipage} \hspace{6mm}
\begin{minipage}{0.24\textwidth}
\normalsize \textbf{In Poverty, 2020}\\
\footnotesize{\textit{millions of people}}\\
  \hspace*{-4mm} \begin{tikzpicture}
    \begin{axis}[
    	\barplotnogrid axis y line=left, \barylab{1.5cm}{1.5ex} \bbar{x}{0},
		yticklabel style={font=\footnotesize},
        width=4.0cm, 
        height=4.1cm,
        ytick={0,1}, nodes near coords, 
        yticklabels={Disposable Income\\*(SPM), Market Income},
        xmin=0, enlarge x limits={0.02}, enlarge y limits={abs=0.55}, 
		every node near coord/.style={/pgf/number format/fixed,
	    	font=\bfseries\scriptsize, style={white}},
		nodes near coords style={/pgf/number format/.cd, fixed zerofill,
			precision=1, assume math mode},
        every axis plot/.append style={
          xbar, bar width=5.0ex, bar shift=0pt, fill}, 
          nodes near coords align={left}]
      \addplot[green!85!blue] coordinates {(29.8, 0)};
      \addplot[blue] coordinates {(82.5, 1)};
    \end{axis}
  \end{tikzpicture}\\
\footnotesize{Source: CPS ASEC}
\vfill
\end{minipage}
\vspace{1mm}

\begin{minipage}{0.76\textwidth}
\small While some fully-employed people are in poverty, \textbf{the vast majority of poor people are either children, elderly, disabled, caregivers, or students}. \input{text/poverty_groups.txt} If the missing labor income required to keep a person out of poverty is not supplied in the form of capital income or welfare income, the person will be poverty, by definition. As a result, these groups that are disproportionately work-limited in some way have much higher rates of poverty. 
\vspace{1mm}

\normalsize \textbf{Poverty Measures, 2020}\vspace{-1mm}

\ \small \cbox{blue}Market Income \ \hspace{1mm} \cbox{lime!70!green!90!blue}Disposable Income (SPM) \ \hspace{1mm} \cbox{green!60!blue}Disposable Income (OPM)
\end{minipage}\vspace{-1mm}

\begin{minipage}{0.4\textwidth}
\normalsize In Poverty\\
\footnotesize{\textit{share of total, percent}}\\
\hspace*{-2mm} \begin{tikzpicture}
  	\begin{axis}[\barplotnogrid axis y line=left, \barylab{1.9cm}{1.5ex}
    	width=4.6cm, bar width=1.9ex, height=10.4cm, xmin=0,
    	enlarge y limits={abs=0.5cm}, enlarge x limits=0.05, \bbar{x}{0},
		yticklabels from table={\poor}{name}, nodes near coords, 
		yticklabel style={font=\footnotesize},
		nodes near coords style={/pgf/number format/.cd, fixed zerofill,
			precision=1, assume math mode}]
  	\addplot[fill=green!60!blue, draw=none] 
  		table [y expr=-\coordindex, x index=2] {\poor};
  	\addplot[fill=lime!70!green!90!blue, draw=none] 
  		table [y expr=-\coordindex, x index=1] {\poor};
  	\addplot[fill=blue, draw=none] 
  		table [y expr=-\coordindex, x index=3] {\poor};
  	\end{axis}
\end{tikzpicture}
\end{minipage}\hspace{4mm}
\begin{minipage}{0.4\textwidth}
\normalsize Poverty Rate\\
\footnotesize{\textit{share of group below poverty line, percent}}\\
\hspace*{-2mm} \begin{tikzpicture}
  	\begin{axis}[\barplotnogrid axis y line=left, \barylab{1.9cm}{1.5ex}
    	width=4.6cm, bar width=1.9ex, height=10.4cm, xmin=0,
    	enlarge y limits={abs=0.5cm}, enlarge x limits=0.05, \bbar{x}{0},
		yticklabels from table={\pvrt}{name}, nodes near coords, 
		yticklabel style={font=\footnotesize},
		nodes near coords style={/pgf/number format/.cd, fixed zerofill,
			precision=1, assume math mode}]
   	\addplot[fill=green!60!blue, draw=none] 
  		table [y expr=-\coordindex, x index=2] {\pvrt};
  	\addplot[fill=lime!70!green!90!blue, draw=none] 
  		table [y expr=-\coordindex, x index=1] {\pvrt};
  	\addplot[fill=blue, draw=none] 
  		table [y expr=-\coordindex, x index=3] {\pvrt};
  	\end{axis}
\end{tikzpicture}
\end{minipage}\vspace{-1mm}

\footnotesize{Source: CPS ASEC; Author's Replication of Bruenig, see \href{https://github.com/bdecon/US-chartbook/blob/master/notebooks/CPS_ASEC.ipynb}{\faPython}} \hspace{44mm} \tbllink{poor.csv}
\newpage
\begin{minipage}{0.76\textwidth}
\small The share of a group whose combined labor, capital, and welfare income is below the poverty line is the poverty rate for the group. In 2020, students, caregivers, and the disabled had the highest rates of poverty. Those fully-employed have a low rate of poverty. 

By age, market income (see\cbox{blue}) leaves the elderly particularly vulnerable to poverty as they are not as likely to have labor income. After social benefits and taxes (disposable income, see\cbox{green!85!blue}), the elderly have much lower rates of poverty than other age cohorts. Higher survivorship for the wealthy also has the effect of reducing poverty in very old ages. Disposable income still leaves young adults and those just below social security and medicare age (late 50s and early 60s) vulnerable to poverty, relative to other ages.
\vspace{2mm}

\normalsize \textbf{Amount of Poverty by Age} \small \hspace{6mm} \cbox{blue} Market Income \ \hspace{2mm} \cbox{green!85!blue} Disposable Income (SPM)\\
\footnotesize{\textit{in billions of 2020 US dollars}}\\
\hspace*{-2mm} \begin{tikzpicture}
	\begin{axis}[ymin=0, ymax=15, enlarge x limits={0.01},
		width=6.7cm, height=5.9cm, ytick={0, 5, 10, 15}, xtick={0, 10, 20, 30, 40, 50, 60, 70},
		xlabel={Age}]
	\abar{green!85!blue}{age}{disposable}{data/poor_age_latest.csv}
	\abar{blue}{age}{market}{data/poor_age_latest.csv}
	\stdnode{0.5cm}{4.1cm}{\large 2020}
	\end{axis}
\end{tikzpicture}
\hfill
\begin{tikzpicture}
	\begin{axis}[ymin=0, ymax=15, enlarge x limits={0.01},
		width=6.7cm, height=5.9cm, ytick={0, 5, 10, 15}, xtick={0, 10, 20, 30, 40, 50, 60, 70},
		yticklabel style={text width=1.0em}, 
		xlabel={Age}]
	\abar{green!85!blue}{age}{disposable}{data/poor_age_prev.csv}
	\abar{blue}{age}{market}{data/poor_age_prev.csv}
	\stdnode{0.5cm}{4.1cm}{\large 2014}
	\end{axis}
\end{tikzpicture}\\
\footnotesize{Source: CPS ASEC, 2014 adjusted for inflation with CPI-U-RS} \hfill \tbllink{poor_age_latest.csv}
\vspace{6mm}

\small \input{text/poverty_rate_time.txt}
\vspace{1mm}

\normalsize \textbf{Poverty Rates}\\
\footnotesize{\textit{in poverty, share of population, percent}}\\
\hspace*{-2mm} \begin{tikzpicture}
	\begin{axis}[\bbar{y}{0}, \dateaxisticks ytick={0, 2, 4, 6, 8, 10, 12, 14, 16, 18}, 
		xticklabel={`\short{\year}}, clip=false, ymin=8, ymax=18, enlarge y limits={0.01}, 
		height=5.2cm, width=12.7cm,  yticklabel style={text width=1.4em}]
	\rbars
	\thickline{green!80!blue}{date}{OPM_LT}{data/poverty_opm_line.csv}
	\thickline{green!80!blue}{date}{OPM_PR}{data/poverty_opm_line.csv}
	\thickline{green!80!blue}{date}{OPM_PR2}{data/poverty_opm_line.csv}
	\thickline{blue!80!green}{date}{LT}{data/poverty_spm_line.csv}
	\thickline{blue!80!green}{date}{PR}{data/poverty_spm_line.csv}
	\thickline{blue!80!green}{date}{PR2}{data/poverty_spm_line.csv}
	\thickline{blue!80!green}{date}{PR3}{data/poverty_spm_line.csv}
	\thickline{blue!80!green}{date}{PR4}{data/poverty_spm_line.csv}
	\absnode{{2009-10-01}}{17.6}{\footnotesize \color{blue!80!green}\textbf{Supplemental (SPM)}}
	\absnode{{1997-01-01}}{11.4}{\footnotesize \color{green!80!blue}\textbf{Official} \\ \small \color{green!80!blue}\textbf{(OPM)}}
	\input{text/poverty_spm_node.txt}
	\input{text/poverty_opm_node.txt}
	\end{axis}
\end{tikzpicture}\\
\footnotesize{Source: Census Bureau} \hfill \tbllink{poverty_opm_line.csv}
\end{minipage}
\newpage
\begin{minipage}{0.76\textwidth}
\small The Census Bureau \href{https://www.census.gov/library/publications/2021/demo/p60-275.html}{report} the number of people taken out of poverty (according to the supplemental poverty measure) by various government programs, along with how many people are put in poverty by various expenses. In 2020, Social Security payments lift income above the poverty line for 26.5 million people, by far the most effective program for reducing poverty.

Economic impact payments, for example from the CARES act, prevented 11.7 million people from being in poverty. Unemployment benefits removed 5.5 million people from poverty. Refundable tax credits, which often fail to reach the poorest people, would remove 5.3 million people from poverty if participation rates were 100 percent. 

Several elements add to the number of people in poverty, medical expenses are the most significant, and cause the disposable income of five million people to fall below the poverty line. Work expenses additionally put two and a half million people in poverty.

\vspace{5mm}

\normalsize \textbf{Effect of Individual Elements on Poverty Headcount}\\
\footnotesize{\textit{individual element effect on number of people in poverty, millions, 2020}}\\
\vspace{0.5mm}
\hspace*{-12mm}  \begin{tikzpicture}
	\begin{axis}[clip=false,
            xbar stacked, \bbar{x}{0},
            xmin=-31, xmax=7.5, ymin=-0.5, ymax=19.5,
            stack negative=separate, ymajorgrids=false,
            yticklabels from table={\spm}{name},
			y axis line style={opacity=0},   
		    x axis line style={opacity=0}, 
			yticklabel style={black, align=left, anchor=east},            
            ytick=data, tickwidth=0pt,
            legend style={
                at={(0.0, 1.05)},
                anchor=north,legend
                columns=-1, draw=none},
            legend cell align=left, extra x ticks={0},
            extra x tick style={
                grid style={black},
                xticklabel=\empty},
            enlarge y limits={0.02}, enlarge x limits={0.02},
            width=10.4cm, height=11.6cm,
            \barylab{3.8cm}{1.5ex}]
            \addplot[fill=green!75!blue!75!black, draw opacity=0] table 
            	[y expr=\coordindex, x=kids, col sep=comma] {\spm};
            \addplot[fill=cyan!95!blue!70!white, draw opacity=0] table 
            	[y expr=\coordindex, x=adults, col sep=comma] {\spm};
            \addplot[fill=blue!80!purple, draw opacity=0] table 
            	[y expr=\coordindex, x=elderly, col sep=comma] {\spm};
            \legend{Under 18, 18 to 64, 65 or older}
        \end{axis}
        \begin{axis}[xmin=-31, xmax=7.5, ymin=-0.5, 
            ymax=19.5, enlarge y limits={0.02}, enlarge x limits={0.02}, xticklabels=\empty,
            width=10.4cm, height=11.6cm, ymajorgrids=false, yticklabels=\empty,
			y axis line style={opacity=0}, x axis line style={opacity=0}, ]
            \addplot[scatter, no marks, draw=none, nodes near coords*={\Label},
            	nodes near coords align={horizontal},
        		visualization depends on={value \thisrow{label} \as \Label}] table 		
        		[x=xloc,y=yloc,col sep=comma] {data/spmtbl20.csv};
        \end{axis}
    \end{tikzpicture}\\
\footnotesize{Source: Census Bureau Supplemental Poverty Measure} \hfill \tbllink{spmtbl20.csv}
\end{minipage}
\newpage
\begin{minipage}{0.76\textwidth}
\small Poverty can be geographically concentrated. In the United States, some regions have particularly high and persistent poverty rates. In 2020, one third of people in the US have income below \input{text/low_income_cz_cutoff.txt}\unskip. In some regions of the US, more than half of the population has income below this threshold. 

\input{text/low_income_cz.txt}
\end{minipage}
\vspace{2mm}

\normalsize \textbf{Low-Income Share of Commuter Zone, \input{text/acs_ftfy_year.txt}}
\vspace{-2mm}

\footnotesize{\textit{Share of commuting zone householders with income below \input{text/low_income_cz_cutoff.txt}}}

\vspace{-1mm}
\hspace{-15mm} %% Creator: Matplotlib, PGF backend
%%
%% To include the figure in your LaTeX document, write
%%   \input{<filename>.pgf}
%%
%% Make sure the required packages are loaded in your preamble
%%   \usepackage{pgf}
%%
%% Also ensure that all the required font packages are loaded; for instance,
%% the lmodern package is sometimes necessary when using math font.
%%   \usepackage{lmodern}
%%
%% Figures using additional raster images can only be included by \input if
%% they are in the same directory as the main LaTeX file. For loading figures
%% from other directories you can use the `import` package
%%   \usepackage{import}
%%
%% and then include the figures with
%%   \import{<path to file>}{<filename>.pgf}
%%
%% Matplotlib used the following preamble
%%   
%%   \usepackage{fontspec}
%%   \setmainfont{DejaVuSerif.ttf}[Path=\detokenize{/home/brian/miniconda3/lib/python3.8/site-packages/matplotlib/mpl-data/fonts/ttf/}]
%%   \setsansfont{DejaVuSans.ttf}[Path=\detokenize{/home/brian/miniconda3/lib/python3.8/site-packages/matplotlib/mpl-data/fonts/ttf/}]
%%   \setmonofont{DejaVuSansMono.ttf}[Path=\detokenize{/home/brian/miniconda3/lib/python3.8/site-packages/matplotlib/mpl-data/fonts/ttf/}]
%%   \makeatletter\@ifpackageloaded{underscore}{}{\usepackage[strings]{underscore}}\makeatother
%%
\begingroup%
\makeatletter%
\begin{pgfpicture}%
\pgfpathrectangle{\pgfpointorigin}{\pgfqpoint{3.841155in}{2.510000in}}%
\pgfusepath{use as bounding box, clip}%
\begin{pgfscope}%
\pgfsetbuttcap%
\pgfsetmiterjoin%
\pgfsetlinewidth{0.000000pt}%
\definecolor{currentstroke}{rgb}{0.000000,0.000000,0.000000}%
\pgfsetstrokecolor{currentstroke}%
\pgfsetstrokeopacity{0.000000}%
\pgfsetdash{}{0pt}%
\pgfpathmoveto{\pgfqpoint{0.000000in}{0.000000in}}%
\pgfpathlineto{\pgfqpoint{3.841155in}{0.000000in}}%
\pgfpathlineto{\pgfqpoint{3.841155in}{2.510000in}}%
\pgfpathlineto{\pgfqpoint{0.000000in}{2.510000in}}%
\pgfpathlineto{\pgfqpoint{0.000000in}{0.000000in}}%
\pgfpathclose%
\pgfusepath{}%
\end{pgfscope}%
\begin{pgfscope}%
\pgfpathrectangle{\pgfqpoint{0.100000in}{0.100000in}}{\pgfqpoint{3.608454in}{2.310000in}}%
\pgfusepath{clip}%
\pgfsetbuttcap%
\pgfsetroundjoin%
\pgfsetlinewidth{0.010037pt}%
\definecolor{currentstroke}{rgb}{1.000000,1.000000,1.000000}%
\pgfsetstrokecolor{currentstroke}%
\pgfsetdash{}{0pt}%
\pgfusepath{stroke}%
\end{pgfscope}%
\begin{pgfscope}%
\pgfpathrectangle{\pgfqpoint{0.100000in}{0.100000in}}{\pgfqpoint{3.608454in}{2.310000in}}%
\pgfusepath{clip}%
\pgfsetbuttcap%
\pgfsetroundjoin%
\pgfsetlinewidth{0.010037pt}%
\definecolor{currentstroke}{rgb}{1.000000,1.000000,1.000000}%
\pgfsetstrokecolor{currentstroke}%
\pgfsetdash{}{0pt}%
\pgfusepath{stroke}%
\end{pgfscope}%
\begin{pgfscope}%
\pgfpathrectangle{\pgfqpoint{0.100000in}{0.100000in}}{\pgfqpoint{3.608454in}{2.310000in}}%
\pgfusepath{clip}%
\pgfsetbuttcap%
\pgfsetroundjoin%
\pgfsetlinewidth{0.010037pt}%
\definecolor{currentstroke}{rgb}{1.000000,1.000000,1.000000}%
\pgfsetstrokecolor{currentstroke}%
\pgfsetdash{}{0pt}%
\pgfusepath{stroke}%
\end{pgfscope}%
\begin{pgfscope}%
\pgfpathrectangle{\pgfqpoint{0.100000in}{0.100000in}}{\pgfqpoint{3.608454in}{2.310000in}}%
\pgfusepath{clip}%
\pgfsetbuttcap%
\pgfsetroundjoin%
\pgfsetlinewidth{0.010037pt}%
\definecolor{currentstroke}{rgb}{1.000000,1.000000,1.000000}%
\pgfsetstrokecolor{currentstroke}%
\pgfsetdash{}{0pt}%
\pgfusepath{stroke}%
\end{pgfscope}%
\begin{pgfscope}%
\pgfpathrectangle{\pgfqpoint{0.100000in}{0.100000in}}{\pgfqpoint{3.608454in}{2.310000in}}%
\pgfusepath{clip}%
\pgfsetbuttcap%
\pgfsetroundjoin%
\pgfsetlinewidth{0.010037pt}%
\definecolor{currentstroke}{rgb}{1.000000,1.000000,1.000000}%
\pgfsetstrokecolor{currentstroke}%
\pgfsetdash{}{0pt}%
\pgfusepath{stroke}%
\end{pgfscope}%
\begin{pgfscope}%
\pgfpathrectangle{\pgfqpoint{0.100000in}{0.100000in}}{\pgfqpoint{3.608454in}{2.310000in}}%
\pgfusepath{clip}%
\pgfsetbuttcap%
\pgfsetroundjoin%
\pgfsetlinewidth{0.010037pt}%
\definecolor{currentstroke}{rgb}{1.000000,1.000000,1.000000}%
\pgfsetstrokecolor{currentstroke}%
\pgfsetdash{}{0pt}%
\pgfusepath{stroke}%
\end{pgfscope}%
\begin{pgfscope}%
\pgfpathrectangle{\pgfqpoint{0.100000in}{0.100000in}}{\pgfqpoint{3.608454in}{2.310000in}}%
\pgfusepath{clip}%
\pgfsetbuttcap%
\pgfsetroundjoin%
\pgfsetlinewidth{0.010037pt}%
\definecolor{currentstroke}{rgb}{1.000000,1.000000,1.000000}%
\pgfsetstrokecolor{currentstroke}%
\pgfsetdash{}{0pt}%
\pgfpathmoveto{\pgfqpoint{0.957404in}{2.274144in}}%
\pgfpathlineto{\pgfqpoint{0.938691in}{2.199997in}}%
\pgfpathlineto{\pgfqpoint{0.925385in}{2.146869in}}%
\pgfpathlineto{\pgfqpoint{0.909565in}{2.082958in}}%
\pgfpathlineto{\pgfqpoint{0.906458in}{2.067206in}}%
\pgfpathlineto{\pgfqpoint{0.908487in}{2.052757in}}%
\pgfpathlineto{\pgfqpoint{0.905799in}{2.039449in}}%
\pgfpathlineto{\pgfqpoint{0.795364in}{2.068266in}}%
\pgfpathlineto{\pgfqpoint{0.785388in}{2.064919in}}%
\pgfpathlineto{\pgfqpoint{0.776851in}{2.067824in}}%
\pgfpathlineto{\pgfqpoint{0.737324in}{2.068774in}}%
\pgfpathlineto{\pgfqpoint{0.723956in}{2.065005in}}%
\pgfpathlineto{\pgfqpoint{0.710680in}{2.066278in}}%
\pgfpathlineto{\pgfqpoint{0.705053in}{2.072177in}}%
\pgfpathlineto{\pgfqpoint{0.669666in}{2.070891in}}%
\pgfpathlineto{\pgfqpoint{0.664077in}{2.080532in}}%
\pgfpathlineto{\pgfqpoint{0.653468in}{2.085611in}}%
\pgfpathlineto{\pgfqpoint{0.638032in}{2.088706in}}%
\pgfpathlineto{\pgfqpoint{0.611388in}{2.084174in}}%
\pgfpathlineto{\pgfqpoint{0.601553in}{2.088614in}}%
\pgfpathlineto{\pgfqpoint{0.586388in}{2.100423in}}%
\pgfpathlineto{\pgfqpoint{0.590946in}{2.123567in}}%
\pgfpathlineto{\pgfqpoint{0.589348in}{2.135386in}}%
\pgfpathlineto{\pgfqpoint{0.577330in}{2.147934in}}%
\pgfpathlineto{\pgfqpoint{0.569631in}{2.147155in}}%
\pgfpathlineto{\pgfqpoint{0.564109in}{2.159848in}}%
\pgfpathlineto{\pgfqpoint{0.551027in}{2.164978in}}%
\pgfpathlineto{\pgfqpoint{0.541518in}{2.164324in}}%
\pgfpathlineto{\pgfqpoint{0.538442in}{2.177359in}}%
\pgfpathlineto{\pgfqpoint{0.547915in}{2.175972in}}%
\pgfpathlineto{\pgfqpoint{0.547162in}{2.194090in}}%
\pgfpathlineto{\pgfqpoint{0.543033in}{2.204847in}}%
\pgfpathlineto{\pgfqpoint{0.544853in}{2.218617in}}%
\pgfpathlineto{\pgfqpoint{0.549810in}{2.228973in}}%
\pgfpathlineto{\pgfqpoint{0.549530in}{2.248625in}}%
\pgfpathlineto{\pgfqpoint{0.546884in}{2.255755in}}%
\pgfpathlineto{\pgfqpoint{0.551454in}{2.278631in}}%
\pgfpathlineto{\pgfqpoint{0.550135in}{2.293374in}}%
\pgfpathlineto{\pgfqpoint{0.545566in}{2.300443in}}%
\pgfpathlineto{\pgfqpoint{0.546263in}{2.323556in}}%
\pgfpathlineto{\pgfqpoint{0.552909in}{2.340554in}}%
\pgfpathlineto{\pgfqpoint{0.560136in}{2.336394in}}%
\pgfpathlineto{\pgfqpoint{0.584067in}{2.311673in}}%
\pgfpathlineto{\pgfqpoint{0.613066in}{2.298145in}}%
\pgfpathlineto{\pgfqpoint{0.627914in}{2.296544in}}%
\pgfpathlineto{\pgfqpoint{0.641902in}{2.290785in}}%
\pgfpathlineto{\pgfqpoint{0.645873in}{2.271403in}}%
\pgfpathlineto{\pgfqpoint{0.633614in}{2.267732in}}%
\pgfpathlineto{\pgfqpoint{0.625490in}{2.252457in}}%
\pgfpathlineto{\pgfqpoint{0.635028in}{2.253306in}}%
\pgfpathlineto{\pgfqpoint{0.638875in}{2.260162in}}%
\pgfpathlineto{\pgfqpoint{0.652326in}{2.268736in}}%
\pgfpathlineto{\pgfqpoint{0.651631in}{2.256004in}}%
\pgfpathlineto{\pgfqpoint{0.642650in}{2.253951in}}%
\pgfpathlineto{\pgfqpoint{0.644141in}{2.237586in}}%
\pgfpathlineto{\pgfqpoint{0.633204in}{2.221193in}}%
\pgfpathlineto{\pgfqpoint{0.626204in}{2.228988in}}%
\pgfpathlineto{\pgfqpoint{0.605402in}{2.224750in}}%
\pgfpathlineto{\pgfqpoint{0.604308in}{2.215213in}}%
\pgfpathlineto{\pgfqpoint{0.611479in}{2.209403in}}%
\pgfpathlineto{\pgfqpoint{0.622419in}{2.208832in}}%
\pgfpathlineto{\pgfqpoint{0.637546in}{2.221247in}}%
\pgfpathlineto{\pgfqpoint{0.649519in}{2.222278in}}%
\pgfpathlineto{\pgfqpoint{0.649978in}{2.236021in}}%
\pgfpathlineto{\pgfqpoint{0.656149in}{2.256205in}}%
\pgfpathlineto{\pgfqpoint{0.670772in}{2.271217in}}%
\pgfpathlineto{\pgfqpoint{0.665896in}{2.282844in}}%
\pgfpathlineto{\pgfqpoint{0.669227in}{2.295363in}}%
\pgfpathlineto{\pgfqpoint{0.662739in}{2.307839in}}%
\pgfpathlineto{\pgfqpoint{0.671985in}{2.323793in}}%
\pgfpathlineto{\pgfqpoint{0.673370in}{2.333392in}}%
\pgfpathlineto{\pgfqpoint{0.665342in}{2.339700in}}%
\pgfpathlineto{\pgfqpoint{0.666679in}{2.355837in}}%
\pgfpathlineto{\pgfqpoint{0.762858in}{2.326616in}}%
\pgfpathlineto{\pgfqpoint{0.864991in}{2.297975in}}%
\pgfpathlineto{\pgfqpoint{0.957404in}{2.274144in}}%
\pgfusepath{stroke}%
\end{pgfscope}%
\begin{pgfscope}%
\pgfpathrectangle{\pgfqpoint{0.100000in}{0.100000in}}{\pgfqpoint{3.608454in}{2.310000in}}%
\pgfusepath{clip}%
\pgfsetbuttcap%
\pgfsetroundjoin%
\pgfsetlinewidth{0.010037pt}%
\definecolor{currentstroke}{rgb}{1.000000,1.000000,1.000000}%
\pgfsetstrokecolor{currentstroke}%
\pgfsetdash{}{0pt}%
\pgfpathmoveto{\pgfqpoint{0.650857in}{2.300012in}}%
\pgfpathlineto{\pgfqpoint{0.655593in}{2.281754in}}%
\pgfpathlineto{\pgfqpoint{0.663020in}{2.275689in}}%
\pgfpathlineto{\pgfqpoint{0.657094in}{2.267969in}}%
\pgfpathlineto{\pgfqpoint{0.651318in}{2.279281in}}%
\pgfpathlineto{\pgfqpoint{0.650857in}{2.300012in}}%
\pgfusepath{stroke}%
\end{pgfscope}%
\begin{pgfscope}%
\pgfpathrectangle{\pgfqpoint{0.100000in}{0.100000in}}{\pgfqpoint{3.608454in}{2.310000in}}%
\pgfusepath{clip}%
\pgfsetbuttcap%
\pgfsetroundjoin%
\pgfsetlinewidth{0.010037pt}%
\definecolor{currentstroke}{rgb}{1.000000,1.000000,1.000000}%
\pgfsetstrokecolor{currentstroke}%
\pgfsetdash{}{0pt}%
\pgfpathmoveto{\pgfqpoint{1.007001in}{2.262160in}}%
\pgfpathlineto{\pgfqpoint{1.109811in}{2.239088in}}%
\pgfpathlineto{\pgfqpoint{1.206635in}{2.219517in}}%
\pgfpathlineto{\pgfqpoint{1.281130in}{2.205868in}}%
\pgfpathlineto{\pgfqpoint{1.346076in}{2.194957in}}%
\pgfpathlineto{\pgfqpoint{1.411166in}{2.184941in}}%
\pgfpathlineto{\pgfqpoint{1.466597in}{2.177132in}}%
\pgfpathlineto{\pgfqpoint{1.522115in}{2.169972in}}%
\pgfpathlineto{\pgfqpoint{1.577712in}{2.163461in}}%
\pgfpathlineto{\pgfqpoint{1.630106in}{2.157928in}}%
\pgfpathlineto{\pgfqpoint{1.622900in}{2.078078in}}%
\pgfpathlineto{\pgfqpoint{1.612175in}{1.970174in}}%
\pgfpathlineto{\pgfqpoint{1.606550in}{1.914676in}}%
\pgfpathlineto{\pgfqpoint{1.598460in}{1.839883in}}%
\pgfpathlineto{\pgfqpoint{1.541187in}{1.846137in}}%
\pgfpathlineto{\pgfqpoint{1.475618in}{1.853580in}}%
\pgfpathlineto{\pgfqpoint{1.384573in}{1.866040in}}%
\pgfpathlineto{\pgfqpoint{1.343910in}{1.871841in}}%
\pgfpathlineto{\pgfqpoint{1.243728in}{1.887567in}}%
\pgfpathlineto{\pgfqpoint{1.209241in}{1.893876in}}%
\pgfpathlineto{\pgfqpoint{1.202053in}{1.853008in}}%
\pgfpathlineto{\pgfqpoint{1.198123in}{1.855924in}}%
\pgfpathlineto{\pgfqpoint{1.190738in}{1.875584in}}%
\pgfpathlineto{\pgfqpoint{1.181245in}{1.874373in}}%
\pgfpathlineto{\pgfqpoint{1.174451in}{1.864434in}}%
\pgfpathlineto{\pgfqpoint{1.160350in}{1.863105in}}%
\pgfpathlineto{\pgfqpoint{1.157517in}{1.867860in}}%
\pgfpathlineto{\pgfqpoint{1.144309in}{1.867284in}}%
\pgfpathlineto{\pgfqpoint{1.137669in}{1.871889in}}%
\pgfpathlineto{\pgfqpoint{1.128385in}{1.864765in}}%
\pgfpathlineto{\pgfqpoint{1.105833in}{1.871176in}}%
\pgfpathlineto{\pgfqpoint{1.095991in}{1.867733in}}%
\pgfpathlineto{\pgfqpoint{1.090697in}{1.883009in}}%
\pgfpathlineto{\pgfqpoint{1.090449in}{1.897544in}}%
\pgfpathlineto{\pgfqpoint{1.074904in}{1.907997in}}%
\pgfpathlineto{\pgfqpoint{1.077439in}{1.923062in}}%
\pgfpathlineto{\pgfqpoint{1.066370in}{1.947947in}}%
\pgfpathlineto{\pgfqpoint{1.067116in}{1.969101in}}%
\pgfpathlineto{\pgfqpoint{1.057602in}{1.979454in}}%
\pgfpathlineto{\pgfqpoint{1.048676in}{1.970294in}}%
\pgfpathlineto{\pgfqpoint{1.036527in}{1.965340in}}%
\pgfpathlineto{\pgfqpoint{1.027278in}{1.975551in}}%
\pgfpathlineto{\pgfqpoint{1.028782in}{1.991857in}}%
\pgfpathlineto{\pgfqpoint{1.039874in}{1.998281in}}%
\pgfpathlineto{\pgfqpoint{1.036925in}{2.008611in}}%
\pgfpathlineto{\pgfqpoint{1.056205in}{2.058614in}}%
\pgfpathlineto{\pgfqpoint{1.041003in}{2.059867in}}%
\pgfpathlineto{\pgfqpoint{1.039440in}{2.068841in}}%
\pgfpathlineto{\pgfqpoint{1.028366in}{2.076578in}}%
\pgfpathlineto{\pgfqpoint{1.029074in}{2.085229in}}%
\pgfpathlineto{\pgfqpoint{1.023232in}{2.091949in}}%
\pgfpathlineto{\pgfqpoint{1.013723in}{2.116347in}}%
\pgfpathlineto{\pgfqpoint{1.005033in}{2.121323in}}%
\pgfpathlineto{\pgfqpoint{0.998811in}{2.142381in}}%
\pgfpathlineto{\pgfqpoint{1.000279in}{2.156377in}}%
\pgfpathlineto{\pgfqpoint{0.988670in}{2.182155in}}%
\pgfpathlineto{\pgfqpoint{1.007001in}{2.262160in}}%
\pgfusepath{stroke}%
\end{pgfscope}%
\begin{pgfscope}%
\pgfpathrectangle{\pgfqpoint{0.100000in}{0.100000in}}{\pgfqpoint{3.608454in}{2.310000in}}%
\pgfusepath{clip}%
\pgfsetbuttcap%
\pgfsetroundjoin%
\pgfsetlinewidth{0.010037pt}%
\definecolor{currentstroke}{rgb}{1.000000,1.000000,1.000000}%
\pgfsetstrokecolor{currentstroke}%
\pgfsetdash{}{0pt}%
\pgfpathmoveto{\pgfqpoint{3.488916in}{1.850686in}}%
\pgfpathlineto{\pgfqpoint{3.486858in}{1.859435in}}%
\pgfpathlineto{\pgfqpoint{3.475422in}{1.867100in}}%
\pgfpathlineto{\pgfqpoint{3.445201in}{1.965955in}}%
\pgfpathlineto{\pgfqpoint{3.428911in}{2.013660in}}%
\pgfpathlineto{\pgfqpoint{3.442007in}{2.027403in}}%
\pgfpathlineto{\pgfqpoint{3.454737in}{2.061039in}}%
\pgfpathlineto{\pgfqpoint{3.461086in}{2.071788in}}%
\pgfpathlineto{\pgfqpoint{3.456588in}{2.076302in}}%
\pgfpathlineto{\pgfqpoint{3.455082in}{2.106159in}}%
\pgfpathlineto{\pgfqpoint{3.460762in}{2.115312in}}%
\pgfpathlineto{\pgfqpoint{3.458272in}{2.136586in}}%
\pgfpathlineto{\pgfqpoint{3.480942in}{2.206386in}}%
\pgfpathlineto{\pgfqpoint{3.491156in}{2.206767in}}%
\pgfpathlineto{\pgfqpoint{3.495396in}{2.194299in}}%
\pgfpathlineto{\pgfqpoint{3.504473in}{2.190719in}}%
\pgfpathlineto{\pgfqpoint{3.521621in}{2.205384in}}%
\pgfpathlineto{\pgfqpoint{3.535066in}{2.214046in}}%
\pgfpathlineto{\pgfqpoint{3.564692in}{2.198793in}}%
\pgfpathlineto{\pgfqpoint{3.591456in}{2.114100in}}%
\pgfpathlineto{\pgfqpoint{3.596544in}{2.093264in}}%
\pgfpathlineto{\pgfqpoint{3.608267in}{2.090820in}}%
\pgfpathlineto{\pgfqpoint{3.624167in}{2.076668in}}%
\pgfpathlineto{\pgfqpoint{3.623267in}{2.068423in}}%
\pgfpathlineto{\pgfqpoint{3.634098in}{2.058650in}}%
\pgfpathlineto{\pgfqpoint{3.642896in}{2.064257in}}%
\pgfpathlineto{\pgfqpoint{3.661307in}{2.042785in}}%
\pgfpathlineto{\pgfqpoint{3.653078in}{2.025601in}}%
\pgfpathlineto{\pgfqpoint{3.642044in}{2.025261in}}%
\pgfpathlineto{\pgfqpoint{3.633192in}{2.009922in}}%
\pgfpathlineto{\pgfqpoint{3.622468in}{2.007746in}}%
\pgfpathlineto{\pgfqpoint{3.614595in}{1.999578in}}%
\pgfpathlineto{\pgfqpoint{3.591081in}{1.990819in}}%
\pgfpathlineto{\pgfqpoint{3.576849in}{1.976702in}}%
\pgfpathlineto{\pgfqpoint{3.569620in}{1.986835in}}%
\pgfpathlineto{\pgfqpoint{3.563050in}{1.979573in}}%
\pgfpathlineto{\pgfqpoint{3.564849in}{1.950215in}}%
\pgfpathlineto{\pgfqpoint{3.559648in}{1.938586in}}%
\pgfpathlineto{\pgfqpoint{3.548302in}{1.941762in}}%
\pgfpathlineto{\pgfqpoint{3.546471in}{1.929840in}}%
\pgfpathlineto{\pgfqpoint{3.541568in}{1.924938in}}%
\pgfpathlineto{\pgfqpoint{3.531757in}{1.926132in}}%
\pgfpathlineto{\pgfqpoint{3.532375in}{1.914985in}}%
\pgfpathlineto{\pgfqpoint{3.517503in}{1.918151in}}%
\pgfpathlineto{\pgfqpoint{3.509381in}{1.902707in}}%
\pgfpathlineto{\pgfqpoint{3.512448in}{1.894627in}}%
\pgfpathlineto{\pgfqpoint{3.507617in}{1.881187in}}%
\pgfpathlineto{\pgfqpoint{3.500012in}{1.871301in}}%
\pgfpathlineto{\pgfqpoint{3.498098in}{1.850665in}}%
\pgfpathlineto{\pgfqpoint{3.488916in}{1.850686in}}%
\pgfusepath{stroke}%
\end{pgfscope}%
\begin{pgfscope}%
\pgfpathrectangle{\pgfqpoint{0.100000in}{0.100000in}}{\pgfqpoint{3.608454in}{2.310000in}}%
\pgfusepath{clip}%
\pgfsetbuttcap%
\pgfsetroundjoin%
\pgfsetlinewidth{0.010037pt}%
\definecolor{currentstroke}{rgb}{1.000000,1.000000,1.000000}%
\pgfsetstrokecolor{currentstroke}%
\pgfsetdash{}{0pt}%
\pgfpathmoveto{\pgfqpoint{3.595305in}{1.984790in}}%
\pgfpathlineto{\pgfqpoint{3.602019in}{1.991823in}}%
\pgfpathlineto{\pgfqpoint{3.608411in}{1.985213in}}%
\pgfpathlineto{\pgfqpoint{3.596929in}{1.976454in}}%
\pgfpathlineto{\pgfqpoint{3.595305in}{1.984790in}}%
\pgfusepath{stroke}%
\end{pgfscope}%
\begin{pgfscope}%
\pgfpathrectangle{\pgfqpoint{0.100000in}{0.100000in}}{\pgfqpoint{3.608454in}{2.310000in}}%
\pgfusepath{clip}%
\pgfsetbuttcap%
\pgfsetroundjoin%
\pgfsetlinewidth{0.010037pt}%
\definecolor{currentstroke}{rgb}{1.000000,1.000000,1.000000}%
\pgfsetstrokecolor{currentstroke}%
\pgfsetdash{}{0pt}%
\pgfpathmoveto{\pgfqpoint{1.606550in}{1.914676in}}%
\pgfpathlineto{\pgfqpoint{1.612175in}{1.970174in}}%
\pgfpathlineto{\pgfqpoint{1.622900in}{2.078078in}}%
\pgfpathlineto{\pgfqpoint{1.630106in}{2.157928in}}%
\pgfpathlineto{\pgfqpoint{1.689116in}{2.152392in}}%
\pgfpathlineto{\pgfqpoint{1.764609in}{2.146381in}}%
\pgfpathlineto{\pgfqpoint{1.833613in}{2.141936in}}%
\pgfpathlineto{\pgfqpoint{1.896094in}{2.138775in}}%
\pgfpathlineto{\pgfqpoint{1.989337in}{2.135576in}}%
\pgfpathlineto{\pgfqpoint{1.995700in}{2.109297in}}%
\pgfpathlineto{\pgfqpoint{1.993422in}{2.096738in}}%
\pgfpathlineto{\pgfqpoint{1.992673in}{2.070924in}}%
\pgfpathlineto{\pgfqpoint{1.997002in}{2.051591in}}%
\pgfpathlineto{\pgfqpoint{2.006920in}{2.023070in}}%
\pgfpathlineto{\pgfqpoint{2.006897in}{1.987162in}}%
\pgfpathlineto{\pgfqpoint{2.008733in}{1.945448in}}%
\pgfpathlineto{\pgfqpoint{2.011253in}{1.934205in}}%
\pgfpathlineto{\pgfqpoint{2.018625in}{1.921863in}}%
\pgfpathlineto{\pgfqpoint{2.021056in}{1.902632in}}%
\pgfpathlineto{\pgfqpoint{2.020006in}{1.889792in}}%
\pgfpathlineto{\pgfqpoint{1.941838in}{1.891485in}}%
\pgfpathlineto{\pgfqpoint{1.884971in}{1.894368in}}%
\pgfpathlineto{\pgfqpoint{1.801609in}{1.898805in}}%
\pgfpathlineto{\pgfqpoint{1.719401in}{1.904645in}}%
\pgfpathlineto{\pgfqpoint{1.664646in}{1.909062in}}%
\pgfpathlineto{\pgfqpoint{1.606550in}{1.914676in}}%
\pgfusepath{stroke}%
\end{pgfscope}%
\begin{pgfscope}%
\pgfpathrectangle{\pgfqpoint{0.100000in}{0.100000in}}{\pgfqpoint{3.608454in}{2.310000in}}%
\pgfusepath{clip}%
\pgfsetbuttcap%
\pgfsetroundjoin%
\pgfsetlinewidth{0.010037pt}%
\definecolor{currentstroke}{rgb}{1.000000,1.000000,1.000000}%
\pgfsetstrokecolor{currentstroke}%
\pgfsetdash{}{0pt}%
\pgfpathmoveto{\pgfqpoint{1.582921in}{1.682417in}}%
\pgfpathlineto{\pgfqpoint{1.591946in}{1.775447in}}%
\pgfpathlineto{\pgfqpoint{1.598460in}{1.839883in}}%
\pgfpathlineto{\pgfqpoint{1.606550in}{1.914676in}}%
\pgfpathlineto{\pgfqpoint{1.664646in}{1.909062in}}%
\pgfpathlineto{\pgfqpoint{1.719401in}{1.904645in}}%
\pgfpathlineto{\pgfqpoint{1.801609in}{1.898805in}}%
\pgfpathlineto{\pgfqpoint{1.884971in}{1.894368in}}%
\pgfpathlineto{\pgfqpoint{1.941838in}{1.891485in}}%
\pgfpathlineto{\pgfqpoint{2.020006in}{1.889792in}}%
\pgfpathlineto{\pgfqpoint{2.014714in}{1.874348in}}%
\pgfpathlineto{\pgfqpoint{2.004150in}{1.862224in}}%
\pgfpathlineto{\pgfqpoint{2.012240in}{1.848265in}}%
\pgfpathlineto{\pgfqpoint{2.021160in}{1.845282in}}%
\pgfpathlineto{\pgfqpoint{2.025393in}{1.837264in}}%
\pgfpathlineto{\pgfqpoint{2.024421in}{1.778721in}}%
\pgfpathlineto{\pgfqpoint{2.022794in}{1.696330in}}%
\pgfpathlineto{\pgfqpoint{2.016778in}{1.674690in}}%
\pgfpathlineto{\pgfqpoint{2.022159in}{1.662720in}}%
\pgfpathlineto{\pgfqpoint{2.011355in}{1.636997in}}%
\pgfpathlineto{\pgfqpoint{2.022737in}{1.616262in}}%
\pgfpathlineto{\pgfqpoint{2.013067in}{1.617849in}}%
\pgfpathlineto{\pgfqpoint{2.006469in}{1.630750in}}%
\pgfpathlineto{\pgfqpoint{1.982922in}{1.639593in}}%
\pgfpathlineto{\pgfqpoint{1.968073in}{1.647372in}}%
\pgfpathlineto{\pgfqpoint{1.943154in}{1.648019in}}%
\pgfpathlineto{\pgfqpoint{1.934505in}{1.640932in}}%
\pgfpathlineto{\pgfqpoint{1.906300in}{1.654886in}}%
\pgfpathlineto{\pgfqpoint{1.904136in}{1.659299in}}%
\pgfpathlineto{\pgfqpoint{1.805703in}{1.663958in}}%
\pgfpathlineto{\pgfqpoint{1.745900in}{1.667386in}}%
\pgfpathlineto{\pgfqpoint{1.696503in}{1.671254in}}%
\pgfpathlineto{\pgfqpoint{1.614880in}{1.678966in}}%
\pgfpathlineto{\pgfqpoint{1.582921in}{1.682417in}}%
\pgfusepath{stroke}%
\end{pgfscope}%
\begin{pgfscope}%
\pgfpathrectangle{\pgfqpoint{0.100000in}{0.100000in}}{\pgfqpoint{3.608454in}{2.310000in}}%
\pgfusepath{clip}%
\pgfsetbuttcap%
\pgfsetroundjoin%
\pgfsetlinewidth{0.010037pt}%
\definecolor{currentstroke}{rgb}{1.000000,1.000000,1.000000}%
\pgfsetstrokecolor{currentstroke}%
\pgfsetdash{}{0pt}%
\pgfpathmoveto{\pgfqpoint{1.567441in}{1.524892in}}%
\pgfpathlineto{\pgfqpoint{1.514964in}{1.529684in}}%
\pgfpathlineto{\pgfqpoint{1.400578in}{1.543802in}}%
\pgfpathlineto{\pgfqpoint{1.338353in}{1.552733in}}%
\pgfpathlineto{\pgfqpoint{1.271614in}{1.562313in}}%
\pgfpathlineto{\pgfqpoint{1.215397in}{1.571309in}}%
\pgfpathlineto{\pgfqpoint{1.153691in}{1.581874in}}%
\pgfpathlineto{\pgfqpoint{1.167721in}{1.659701in}}%
\pgfpathlineto{\pgfqpoint{1.181932in}{1.739505in}}%
\pgfpathlineto{\pgfqpoint{1.202053in}{1.853008in}}%
\pgfpathlineto{\pgfqpoint{1.209241in}{1.893876in}}%
\pgfpathlineto{\pgfqpoint{1.243728in}{1.887567in}}%
\pgfpathlineto{\pgfqpoint{1.343910in}{1.871841in}}%
\pgfpathlineto{\pgfqpoint{1.384573in}{1.866040in}}%
\pgfpathlineto{\pgfqpoint{1.475618in}{1.853580in}}%
\pgfpathlineto{\pgfqpoint{1.541187in}{1.846137in}}%
\pgfpathlineto{\pgfqpoint{1.598460in}{1.839883in}}%
\pgfpathlineto{\pgfqpoint{1.591946in}{1.775447in}}%
\pgfpathlineto{\pgfqpoint{1.582921in}{1.682417in}}%
\pgfpathlineto{\pgfqpoint{1.575175in}{1.603355in}}%
\pgfpathlineto{\pgfqpoint{1.567441in}{1.524892in}}%
\pgfusepath{stroke}%
\end{pgfscope}%
\begin{pgfscope}%
\pgfpathrectangle{\pgfqpoint{0.100000in}{0.100000in}}{\pgfqpoint{3.608454in}{2.310000in}}%
\pgfusepath{clip}%
\pgfsetbuttcap%
\pgfsetroundjoin%
\pgfsetlinewidth{0.010037pt}%
\definecolor{currentstroke}{rgb}{1.000000,1.000000,1.000000}%
\pgfsetstrokecolor{currentstroke}%
\pgfsetdash{}{0pt}%
\pgfpathmoveto{\pgfqpoint{2.526890in}{1.632246in}}%
\pgfpathlineto{\pgfqpoint{2.460390in}{1.627522in}}%
\pgfpathlineto{\pgfqpoint{2.361269in}{1.623298in}}%
\pgfpathlineto{\pgfqpoint{2.357502in}{1.633311in}}%
\pgfpathlineto{\pgfqpoint{2.335521in}{1.640795in}}%
\pgfpathlineto{\pgfqpoint{2.330673in}{1.654936in}}%
\pgfpathlineto{\pgfqpoint{2.328643in}{1.672430in}}%
\pgfpathlineto{\pgfqpoint{2.333596in}{1.681393in}}%
\pgfpathlineto{\pgfqpoint{2.325773in}{1.689994in}}%
\pgfpathlineto{\pgfqpoint{2.323886in}{1.700253in}}%
\pgfpathlineto{\pgfqpoint{2.321364in}{1.722947in}}%
\pgfpathlineto{\pgfqpoint{2.313867in}{1.735275in}}%
\pgfpathlineto{\pgfqpoint{2.300552in}{1.742194in}}%
\pgfpathlineto{\pgfqpoint{2.286062in}{1.753611in}}%
\pgfpathlineto{\pgfqpoint{2.278572in}{1.767120in}}%
\pgfpathlineto{\pgfqpoint{2.265132in}{1.772560in}}%
\pgfpathlineto{\pgfqpoint{2.257232in}{1.781413in}}%
\pgfpathlineto{\pgfqpoint{2.247661in}{1.782916in}}%
\pgfpathlineto{\pgfqpoint{2.230573in}{1.796062in}}%
\pgfpathlineto{\pgfqpoint{2.233349in}{1.811190in}}%
\pgfpathlineto{\pgfqpoint{2.232814in}{1.839958in}}%
\pgfpathlineto{\pgfqpoint{2.238084in}{1.847876in}}%
\pgfpathlineto{\pgfqpoint{2.233349in}{1.859834in}}%
\pgfpathlineto{\pgfqpoint{2.225008in}{1.862146in}}%
\pgfpathlineto{\pgfqpoint{2.225698in}{1.872648in}}%
\pgfpathlineto{\pgfqpoint{2.236046in}{1.889238in}}%
\pgfpathlineto{\pgfqpoint{2.256539in}{1.902357in}}%
\pgfpathlineto{\pgfqpoint{2.255262in}{1.949012in}}%
\pgfpathlineto{\pgfqpoint{2.265516in}{1.956021in}}%
\pgfpathlineto{\pgfqpoint{2.275220in}{1.951352in}}%
\pgfpathlineto{\pgfqpoint{2.294988in}{1.958187in}}%
\pgfpathlineto{\pgfqpoint{2.332179in}{1.975363in}}%
\pgfpathlineto{\pgfqpoint{2.337024in}{1.970045in}}%
\pgfpathlineto{\pgfqpoint{2.329980in}{1.945955in}}%
\pgfpathlineto{\pgfqpoint{2.340460in}{1.951242in}}%
\pgfpathlineto{\pgfqpoint{2.358409in}{1.945963in}}%
\pgfpathlineto{\pgfqpoint{2.369439in}{1.941557in}}%
\pgfpathlineto{\pgfqpoint{2.375622in}{1.928645in}}%
\pgfpathlineto{\pgfqpoint{2.432207in}{1.916460in}}%
\pgfpathlineto{\pgfqpoint{2.449120in}{1.908100in}}%
\pgfpathlineto{\pgfqpoint{2.466317in}{1.908198in}}%
\pgfpathlineto{\pgfqpoint{2.483987in}{1.904750in}}%
\pgfpathlineto{\pgfqpoint{2.495454in}{1.892970in}}%
\pgfpathlineto{\pgfqpoint{2.506412in}{1.887150in}}%
\pgfpathlineto{\pgfqpoint{2.508426in}{1.870361in}}%
\pgfpathlineto{\pgfqpoint{2.505188in}{1.859820in}}%
\pgfpathlineto{\pgfqpoint{2.517401in}{1.859041in}}%
\pgfpathlineto{\pgfqpoint{2.513287in}{1.846815in}}%
\pgfpathlineto{\pgfqpoint{2.517200in}{1.842472in}}%
\pgfpathlineto{\pgfqpoint{2.521092in}{1.830931in}}%
\pgfpathlineto{\pgfqpoint{2.509190in}{1.824821in}}%
\pgfpathlineto{\pgfqpoint{2.502280in}{1.807797in}}%
\pgfpathlineto{\pgfqpoint{2.500111in}{1.795759in}}%
\pgfpathlineto{\pgfqpoint{2.506741in}{1.793698in}}%
\pgfpathlineto{\pgfqpoint{2.515230in}{1.802725in}}%
\pgfpathlineto{\pgfqpoint{2.522446in}{1.818369in}}%
\pgfpathlineto{\pgfqpoint{2.532211in}{1.823802in}}%
\pgfpathlineto{\pgfqpoint{2.539548in}{1.816735in}}%
\pgfpathlineto{\pgfqpoint{2.532326in}{1.795344in}}%
\pgfpathlineto{\pgfqpoint{2.530060in}{1.778738in}}%
\pgfpathlineto{\pgfqpoint{2.532194in}{1.766801in}}%
\pgfpathlineto{\pgfqpoint{2.525487in}{1.760048in}}%
\pgfpathlineto{\pgfqpoint{2.522137in}{1.743546in}}%
\pgfpathlineto{\pgfqpoint{2.524874in}{1.726204in}}%
\pgfpathlineto{\pgfqpoint{2.516958in}{1.700555in}}%
\pgfpathlineto{\pgfqpoint{2.517125in}{1.687738in}}%
\pgfpathlineto{\pgfqpoint{2.523380in}{1.659966in}}%
\pgfpathlineto{\pgfqpoint{2.527434in}{1.655201in}}%
\pgfpathlineto{\pgfqpoint{2.526890in}{1.632246in}}%
\pgfusepath{stroke}%
\end{pgfscope}%
\begin{pgfscope}%
\pgfpathrectangle{\pgfqpoint{0.100000in}{0.100000in}}{\pgfqpoint{3.608454in}{2.310000in}}%
\pgfusepath{clip}%
\pgfsetbuttcap%
\pgfsetroundjoin%
\pgfsetlinewidth{0.010037pt}%
\definecolor{currentstroke}{rgb}{1.000000,1.000000,1.000000}%
\pgfsetstrokecolor{currentstroke}%
\pgfsetdash{}{0pt}%
\pgfpathmoveto{\pgfqpoint{2.551823in}{1.857339in}}%
\pgfpathlineto{\pgfqpoint{2.553308in}{1.846261in}}%
\pgfpathlineto{\pgfqpoint{2.539709in}{1.817072in}}%
\pgfpathlineto{\pgfqpoint{2.533666in}{1.825527in}}%
\pgfpathlineto{\pgfqpoint{2.551823in}{1.857339in}}%
\pgfusepath{stroke}%
\end{pgfscope}%
\begin{pgfscope}%
\pgfpathrectangle{\pgfqpoint{0.100000in}{0.100000in}}{\pgfqpoint{3.608454in}{2.310000in}}%
\pgfusepath{clip}%
\pgfsetbuttcap%
\pgfsetroundjoin%
\pgfsetlinewidth{0.010037pt}%
\definecolor{currentstroke}{rgb}{1.000000,1.000000,1.000000}%
\pgfsetstrokecolor{currentstroke}%
\pgfsetdash{}{0pt}%
\pgfpathmoveto{\pgfqpoint{0.905799in}{2.039449in}}%
\pgfpathlineto{\pgfqpoint{0.908487in}{2.052757in}}%
\pgfpathlineto{\pgfqpoint{0.906458in}{2.067206in}}%
\pgfpathlineto{\pgfqpoint{0.909565in}{2.082958in}}%
\pgfpathlineto{\pgfqpoint{0.925385in}{2.146869in}}%
\pgfpathlineto{\pgfqpoint{0.938691in}{2.199997in}}%
\pgfpathlineto{\pgfqpoint{0.957404in}{2.274144in}}%
\pgfpathlineto{\pgfqpoint{1.007001in}{2.262160in}}%
\pgfpathlineto{\pgfqpoint{0.988670in}{2.182155in}}%
\pgfpathlineto{\pgfqpoint{1.000279in}{2.156377in}}%
\pgfpathlineto{\pgfqpoint{0.998811in}{2.142381in}}%
\pgfpathlineto{\pgfqpoint{1.005033in}{2.121323in}}%
\pgfpathlineto{\pgfqpoint{1.013723in}{2.116347in}}%
\pgfpathlineto{\pgfqpoint{1.023232in}{2.091949in}}%
\pgfpathlineto{\pgfqpoint{1.029074in}{2.085229in}}%
\pgfpathlineto{\pgfqpoint{1.028366in}{2.076578in}}%
\pgfpathlineto{\pgfqpoint{1.039440in}{2.068841in}}%
\pgfpathlineto{\pgfqpoint{1.041003in}{2.059867in}}%
\pgfpathlineto{\pgfqpoint{1.056205in}{2.058614in}}%
\pgfpathlineto{\pgfqpoint{1.036925in}{2.008611in}}%
\pgfpathlineto{\pgfqpoint{1.039874in}{1.998281in}}%
\pgfpathlineto{\pgfqpoint{1.028782in}{1.991857in}}%
\pgfpathlineto{\pgfqpoint{1.027278in}{1.975551in}}%
\pgfpathlineto{\pgfqpoint{1.036527in}{1.965340in}}%
\pgfpathlineto{\pgfqpoint{1.048676in}{1.970294in}}%
\pgfpathlineto{\pgfqpoint{1.057602in}{1.979454in}}%
\pgfpathlineto{\pgfqpoint{1.067116in}{1.969101in}}%
\pgfpathlineto{\pgfqpoint{1.066370in}{1.947947in}}%
\pgfpathlineto{\pgfqpoint{1.077439in}{1.923062in}}%
\pgfpathlineto{\pgfqpoint{1.074904in}{1.907997in}}%
\pgfpathlineto{\pgfqpoint{1.090449in}{1.897544in}}%
\pgfpathlineto{\pgfqpoint{1.090697in}{1.883009in}}%
\pgfpathlineto{\pgfqpoint{1.095991in}{1.867733in}}%
\pgfpathlineto{\pgfqpoint{1.105833in}{1.871176in}}%
\pgfpathlineto{\pgfqpoint{1.128385in}{1.864765in}}%
\pgfpathlineto{\pgfqpoint{1.137669in}{1.871889in}}%
\pgfpathlineto{\pgfqpoint{1.144309in}{1.867284in}}%
\pgfpathlineto{\pgfqpoint{1.157517in}{1.867860in}}%
\pgfpathlineto{\pgfqpoint{1.160350in}{1.863105in}}%
\pgfpathlineto{\pgfqpoint{1.174451in}{1.864434in}}%
\pgfpathlineto{\pgfqpoint{1.181245in}{1.874373in}}%
\pgfpathlineto{\pgfqpoint{1.190738in}{1.875584in}}%
\pgfpathlineto{\pgfqpoint{1.198123in}{1.855924in}}%
\pgfpathlineto{\pgfqpoint{1.202053in}{1.853008in}}%
\pgfpathlineto{\pgfqpoint{1.181932in}{1.739505in}}%
\pgfpathlineto{\pgfqpoint{1.167721in}{1.659701in}}%
\pgfpathlineto{\pgfqpoint{1.055637in}{1.681339in}}%
\pgfpathlineto{\pgfqpoint{0.995094in}{1.693381in}}%
\pgfpathlineto{\pgfqpoint{0.938464in}{1.705834in}}%
\pgfpathlineto{\pgfqpoint{0.885533in}{1.717818in}}%
\pgfpathlineto{\pgfqpoint{0.824275in}{1.732591in}}%
\pgfpathlineto{\pgfqpoint{0.855872in}{1.862213in}}%
\pgfpathlineto{\pgfqpoint{0.857576in}{1.871686in}}%
\pgfpathlineto{\pgfqpoint{0.871627in}{1.896508in}}%
\pgfpathlineto{\pgfqpoint{0.857078in}{1.911427in}}%
\pgfpathlineto{\pgfqpoint{0.860099in}{1.926077in}}%
\pgfpathlineto{\pgfqpoint{0.865637in}{1.929744in}}%
\pgfpathlineto{\pgfqpoint{0.875529in}{1.944840in}}%
\pgfpathlineto{\pgfqpoint{0.889947in}{1.955284in}}%
\pgfpathlineto{\pgfqpoint{0.890735in}{1.963012in}}%
\pgfpathlineto{\pgfqpoint{0.899437in}{1.970733in}}%
\pgfpathlineto{\pgfqpoint{0.906663in}{1.985185in}}%
\pgfpathlineto{\pgfqpoint{0.921484in}{2.000475in}}%
\pgfpathlineto{\pgfqpoint{0.920432in}{2.015576in}}%
\pgfpathlineto{\pgfqpoint{0.910299in}{2.023965in}}%
\pgfpathlineto{\pgfqpoint{0.905799in}{2.039449in}}%
\pgfusepath{stroke}%
\end{pgfscope}%
\begin{pgfscope}%
\pgfpathrectangle{\pgfqpoint{0.100000in}{0.100000in}}{\pgfqpoint{3.608454in}{2.310000in}}%
\pgfusepath{clip}%
\pgfsetbuttcap%
\pgfsetroundjoin%
\pgfsetlinewidth{0.010037pt}%
\definecolor{currentstroke}{rgb}{1.000000,1.000000,1.000000}%
\pgfsetstrokecolor{currentstroke}%
\pgfsetdash{}{0pt}%
\pgfpathmoveto{\pgfqpoint{3.359367in}{1.786492in}}%
\pgfpathlineto{\pgfqpoint{3.354710in}{1.801193in}}%
\pgfpathlineto{\pgfqpoint{3.346067in}{1.845790in}}%
\pgfpathlineto{\pgfqpoint{3.334894in}{1.859510in}}%
\pgfpathlineto{\pgfqpoint{3.325093in}{1.884184in}}%
\pgfpathlineto{\pgfqpoint{3.328311in}{1.902477in}}%
\pgfpathlineto{\pgfqpoint{3.325730in}{1.915960in}}%
\pgfpathlineto{\pgfqpoint{3.317356in}{1.929202in}}%
\pgfpathlineto{\pgfqpoint{3.317011in}{1.943790in}}%
\pgfpathlineto{\pgfqpoint{3.312149in}{1.959523in}}%
\pgfpathlineto{\pgfqpoint{3.355664in}{1.970232in}}%
\pgfpathlineto{\pgfqpoint{3.412185in}{1.985489in}}%
\pgfpathlineto{\pgfqpoint{3.414441in}{1.976741in}}%
\pgfpathlineto{\pgfqpoint{3.410842in}{1.962813in}}%
\pgfpathlineto{\pgfqpoint{3.419291in}{1.951669in}}%
\pgfpathlineto{\pgfqpoint{3.414804in}{1.937548in}}%
\pgfpathlineto{\pgfqpoint{3.396973in}{1.919797in}}%
\pgfpathlineto{\pgfqpoint{3.401874in}{1.906464in}}%
\pgfpathlineto{\pgfqpoint{3.398242in}{1.878280in}}%
\pgfpathlineto{\pgfqpoint{3.393157in}{1.862389in}}%
\pgfpathlineto{\pgfqpoint{3.399644in}{1.817200in}}%
\pgfpathlineto{\pgfqpoint{3.398770in}{1.801301in}}%
\pgfpathlineto{\pgfqpoint{3.405047in}{1.796199in}}%
\pgfpathlineto{\pgfqpoint{3.359367in}{1.786492in}}%
\pgfusepath{stroke}%
\end{pgfscope}%
\begin{pgfscope}%
\pgfpathrectangle{\pgfqpoint{0.100000in}{0.100000in}}{\pgfqpoint{3.608454in}{2.310000in}}%
\pgfusepath{clip}%
\pgfsetbuttcap%
\pgfsetroundjoin%
\pgfsetlinewidth{0.010037pt}%
\definecolor{currentstroke}{rgb}{1.000000,1.000000,1.000000}%
\pgfsetstrokecolor{currentstroke}%
\pgfsetdash{}{0pt}%
\pgfpathmoveto{\pgfqpoint{2.022794in}{1.696330in}}%
\pgfpathlineto{\pgfqpoint{2.024421in}{1.778721in}}%
\pgfpathlineto{\pgfqpoint{2.025393in}{1.837264in}}%
\pgfpathlineto{\pgfqpoint{2.021160in}{1.845282in}}%
\pgfpathlineto{\pgfqpoint{2.012240in}{1.848265in}}%
\pgfpathlineto{\pgfqpoint{2.004150in}{1.862224in}}%
\pgfpathlineto{\pgfqpoint{2.014714in}{1.874348in}}%
\pgfpathlineto{\pgfqpoint{2.020006in}{1.889792in}}%
\pgfpathlineto{\pgfqpoint{2.021056in}{1.902632in}}%
\pgfpathlineto{\pgfqpoint{2.018625in}{1.921863in}}%
\pgfpathlineto{\pgfqpoint{2.011253in}{1.934205in}}%
\pgfpathlineto{\pgfqpoint{2.008733in}{1.945448in}}%
\pgfpathlineto{\pgfqpoint{2.006897in}{1.987162in}}%
\pgfpathlineto{\pgfqpoint{2.006920in}{2.023070in}}%
\pgfpathlineto{\pgfqpoint{1.997002in}{2.051591in}}%
\pgfpathlineto{\pgfqpoint{1.992673in}{2.070924in}}%
\pgfpathlineto{\pgfqpoint{1.993422in}{2.096738in}}%
\pgfpathlineto{\pgfqpoint{1.995700in}{2.109297in}}%
\pgfpathlineto{\pgfqpoint{1.989337in}{2.135576in}}%
\pgfpathlineto{\pgfqpoint{2.032661in}{2.134709in}}%
\pgfpathlineto{\pgfqpoint{2.098467in}{2.134144in}}%
\pgfpathlineto{\pgfqpoint{2.098827in}{2.164013in}}%
\pgfpathlineto{\pgfqpoint{2.115577in}{2.160724in}}%
\pgfpathlineto{\pgfqpoint{2.123605in}{2.124304in}}%
\pgfpathlineto{\pgfqpoint{2.129524in}{2.111209in}}%
\pgfpathlineto{\pgfqpoint{2.144242in}{2.110823in}}%
\pgfpathlineto{\pgfqpoint{2.147536in}{2.106379in}}%
\pgfpathlineto{\pgfqpoint{2.168066in}{2.104411in}}%
\pgfpathlineto{\pgfqpoint{2.171516in}{2.095382in}}%
\pgfpathlineto{\pgfqpoint{2.185665in}{2.097411in}}%
\pgfpathlineto{\pgfqpoint{2.196654in}{2.105856in}}%
\pgfpathlineto{\pgfqpoint{2.215604in}{2.105540in}}%
\pgfpathlineto{\pgfqpoint{2.227330in}{2.098749in}}%
\pgfpathlineto{\pgfqpoint{2.228678in}{2.092380in}}%
\pgfpathlineto{\pgfqpoint{2.239835in}{2.091046in}}%
\pgfpathlineto{\pgfqpoint{2.247116in}{2.073673in}}%
\pgfpathlineto{\pgfqpoint{2.251814in}{2.084355in}}%
\pgfpathlineto{\pgfqpoint{2.264631in}{2.085012in}}%
\pgfpathlineto{\pgfqpoint{2.267873in}{2.076692in}}%
\pgfpathlineto{\pgfqpoint{2.282295in}{2.072895in}}%
\pgfpathlineto{\pgfqpoint{2.290247in}{2.060912in}}%
\pgfpathlineto{\pgfqpoint{2.307885in}{2.064633in}}%
\pgfpathlineto{\pgfqpoint{2.327294in}{2.079338in}}%
\pgfpathlineto{\pgfqpoint{2.334370in}{2.066375in}}%
\pgfpathlineto{\pgfqpoint{2.366216in}{2.069923in}}%
\pgfpathlineto{\pgfqpoint{2.379829in}{2.061026in}}%
\pgfpathlineto{\pgfqpoint{2.387752in}{2.064202in}}%
\pgfpathlineto{\pgfqpoint{2.394106in}{2.059194in}}%
\pgfpathlineto{\pgfqpoint{2.375257in}{2.047310in}}%
\pgfpathlineto{\pgfqpoint{2.348353in}{2.036722in}}%
\pgfpathlineto{\pgfqpoint{2.321706in}{2.015558in}}%
\pgfpathlineto{\pgfqpoint{2.298605in}{1.987726in}}%
\pgfpathlineto{\pgfqpoint{2.281117in}{1.971272in}}%
\pgfpathlineto{\pgfqpoint{2.265799in}{1.959963in}}%
\pgfpathlineto{\pgfqpoint{2.255262in}{1.949012in}}%
\pgfpathlineto{\pgfqpoint{2.256539in}{1.902357in}}%
\pgfpathlineto{\pgfqpoint{2.236046in}{1.889238in}}%
\pgfpathlineto{\pgfqpoint{2.225698in}{1.872648in}}%
\pgfpathlineto{\pgfqpoint{2.225008in}{1.862146in}}%
\pgfpathlineto{\pgfqpoint{2.233349in}{1.859834in}}%
\pgfpathlineto{\pgfqpoint{2.238084in}{1.847876in}}%
\pgfpathlineto{\pgfqpoint{2.232814in}{1.839958in}}%
\pgfpathlineto{\pgfqpoint{2.233349in}{1.811190in}}%
\pgfpathlineto{\pgfqpoint{2.230573in}{1.796062in}}%
\pgfpathlineto{\pgfqpoint{2.247661in}{1.782916in}}%
\pgfpathlineto{\pgfqpoint{2.257232in}{1.781413in}}%
\pgfpathlineto{\pgfqpoint{2.265132in}{1.772560in}}%
\pgfpathlineto{\pgfqpoint{2.278572in}{1.767120in}}%
\pgfpathlineto{\pgfqpoint{2.286062in}{1.753611in}}%
\pgfpathlineto{\pgfqpoint{2.300552in}{1.742194in}}%
\pgfpathlineto{\pgfqpoint{2.313867in}{1.735275in}}%
\pgfpathlineto{\pgfqpoint{2.321364in}{1.722947in}}%
\pgfpathlineto{\pgfqpoint{2.323886in}{1.700253in}}%
\pgfpathlineto{\pgfqpoint{2.253222in}{1.697687in}}%
\pgfpathlineto{\pgfqpoint{2.192987in}{1.696427in}}%
\pgfpathlineto{\pgfqpoint{2.138106in}{1.695621in}}%
\pgfpathlineto{\pgfqpoint{2.080050in}{1.695709in}}%
\pgfpathlineto{\pgfqpoint{2.022794in}{1.696330in}}%
\pgfusepath{stroke}%
\end{pgfscope}%
\begin{pgfscope}%
\pgfpathrectangle{\pgfqpoint{0.100000in}{0.100000in}}{\pgfqpoint{3.608454in}{2.310000in}}%
\pgfusepath{clip}%
\pgfsetbuttcap%
\pgfsetroundjoin%
\pgfsetlinewidth{0.010037pt}%
\definecolor{currentstroke}{rgb}{1.000000,1.000000,1.000000}%
\pgfsetstrokecolor{currentstroke}%
\pgfsetdash{}{0pt}%
\pgfpathmoveto{\pgfqpoint{0.418363in}{1.850356in}}%
\pgfpathlineto{\pgfqpoint{0.412766in}{1.860654in}}%
\pgfpathlineto{\pgfqpoint{0.416340in}{1.887087in}}%
\pgfpathlineto{\pgfqpoint{0.423237in}{1.900901in}}%
\pgfpathlineto{\pgfqpoint{0.419789in}{1.919587in}}%
\pgfpathlineto{\pgfqpoint{0.426984in}{1.927445in}}%
\pgfpathlineto{\pgfqpoint{0.440110in}{1.948627in}}%
\pgfpathlineto{\pgfqpoint{0.451240in}{1.961413in}}%
\pgfpathlineto{\pgfqpoint{0.467509in}{1.989261in}}%
\pgfpathlineto{\pgfqpoint{0.479986in}{2.019596in}}%
\pgfpathlineto{\pgfqpoint{0.493232in}{2.048048in}}%
\pgfpathlineto{\pgfqpoint{0.495910in}{2.059930in}}%
\pgfpathlineto{\pgfqpoint{0.514183in}{2.093924in}}%
\pgfpathlineto{\pgfqpoint{0.517702in}{2.108854in}}%
\pgfpathlineto{\pgfqpoint{0.525482in}{2.124532in}}%
\pgfpathlineto{\pgfqpoint{0.528263in}{2.143657in}}%
\pgfpathlineto{\pgfqpoint{0.543116in}{2.152982in}}%
\pgfpathlineto{\pgfqpoint{0.560727in}{2.157683in}}%
\pgfpathlineto{\pgfqpoint{0.569631in}{2.147155in}}%
\pgfpathlineto{\pgfqpoint{0.577330in}{2.147934in}}%
\pgfpathlineto{\pgfqpoint{0.589348in}{2.135386in}}%
\pgfpathlineto{\pgfqpoint{0.590946in}{2.123567in}}%
\pgfpathlineto{\pgfqpoint{0.586388in}{2.100423in}}%
\pgfpathlineto{\pgfqpoint{0.601553in}{2.088614in}}%
\pgfpathlineto{\pgfqpoint{0.611388in}{2.084174in}}%
\pgfpathlineto{\pgfqpoint{0.638032in}{2.088706in}}%
\pgfpathlineto{\pgfqpoint{0.653468in}{2.085611in}}%
\pgfpathlineto{\pgfqpoint{0.664077in}{2.080532in}}%
\pgfpathlineto{\pgfqpoint{0.669666in}{2.070891in}}%
\pgfpathlineto{\pgfqpoint{0.705053in}{2.072177in}}%
\pgfpathlineto{\pgfqpoint{0.710680in}{2.066278in}}%
\pgfpathlineto{\pgfqpoint{0.723956in}{2.065005in}}%
\pgfpathlineto{\pgfqpoint{0.737324in}{2.068774in}}%
\pgfpathlineto{\pgfqpoint{0.776851in}{2.067824in}}%
\pgfpathlineto{\pgfqpoint{0.785388in}{2.064919in}}%
\pgfpathlineto{\pgfqpoint{0.795364in}{2.068266in}}%
\pgfpathlineto{\pgfqpoint{0.905799in}{2.039449in}}%
\pgfpathlineto{\pgfqpoint{0.910299in}{2.023965in}}%
\pgfpathlineto{\pgfqpoint{0.920432in}{2.015576in}}%
\pgfpathlineto{\pgfqpoint{0.921484in}{2.000475in}}%
\pgfpathlineto{\pgfqpoint{0.906663in}{1.985185in}}%
\pgfpathlineto{\pgfqpoint{0.899437in}{1.970733in}}%
\pgfpathlineto{\pgfqpoint{0.890735in}{1.963012in}}%
\pgfpathlineto{\pgfqpoint{0.889947in}{1.955284in}}%
\pgfpathlineto{\pgfqpoint{0.875529in}{1.944840in}}%
\pgfpathlineto{\pgfqpoint{0.865637in}{1.929744in}}%
\pgfpathlineto{\pgfqpoint{0.860099in}{1.926077in}}%
\pgfpathlineto{\pgfqpoint{0.857078in}{1.911427in}}%
\pgfpathlineto{\pgfqpoint{0.871627in}{1.896508in}}%
\pgfpathlineto{\pgfqpoint{0.857576in}{1.871686in}}%
\pgfpathlineto{\pgfqpoint{0.855872in}{1.862213in}}%
\pgfpathlineto{\pgfqpoint{0.824275in}{1.732591in}}%
\pgfpathlineto{\pgfqpoint{0.757813in}{1.749621in}}%
\pgfpathlineto{\pgfqpoint{0.693698in}{1.766153in}}%
\pgfpathlineto{\pgfqpoint{0.655019in}{1.776916in}}%
\pgfpathlineto{\pgfqpoint{0.605321in}{1.791073in}}%
\pgfpathlineto{\pgfqpoint{0.526049in}{1.816024in}}%
\pgfpathlineto{\pgfqpoint{0.439875in}{1.842853in}}%
\pgfpathlineto{\pgfqpoint{0.418363in}{1.850356in}}%
\pgfusepath{stroke}%
\end{pgfscope}%
\begin{pgfscope}%
\pgfpathrectangle{\pgfqpoint{0.100000in}{0.100000in}}{\pgfqpoint{3.608454in}{2.310000in}}%
\pgfusepath{clip}%
\pgfsetbuttcap%
\pgfsetroundjoin%
\pgfsetlinewidth{0.010037pt}%
\definecolor{currentstroke}{rgb}{1.000000,1.000000,1.000000}%
\pgfsetstrokecolor{currentstroke}%
\pgfsetdash{}{0pt}%
\pgfpathmoveto{\pgfqpoint{3.405047in}{1.796199in}}%
\pgfpathlineto{\pgfqpoint{3.398770in}{1.801301in}}%
\pgfpathlineto{\pgfqpoint{3.399644in}{1.817200in}}%
\pgfpathlineto{\pgfqpoint{3.393157in}{1.862389in}}%
\pgfpathlineto{\pgfqpoint{3.398242in}{1.878280in}}%
\pgfpathlineto{\pgfqpoint{3.401874in}{1.906464in}}%
\pgfpathlineto{\pgfqpoint{3.396973in}{1.919797in}}%
\pgfpathlineto{\pgfqpoint{3.414804in}{1.937548in}}%
\pgfpathlineto{\pgfqpoint{3.419291in}{1.951669in}}%
\pgfpathlineto{\pgfqpoint{3.410842in}{1.962813in}}%
\pgfpathlineto{\pgfqpoint{3.414441in}{1.976741in}}%
\pgfpathlineto{\pgfqpoint{3.412185in}{1.985489in}}%
\pgfpathlineto{\pgfqpoint{3.414122in}{2.004225in}}%
\pgfpathlineto{\pgfqpoint{3.417744in}{2.010009in}}%
\pgfpathlineto{\pgfqpoint{3.428911in}{2.013660in}}%
\pgfpathlineto{\pgfqpoint{3.445201in}{1.965955in}}%
\pgfpathlineto{\pgfqpoint{3.475422in}{1.867100in}}%
\pgfpathlineto{\pgfqpoint{3.486858in}{1.859435in}}%
\pgfpathlineto{\pgfqpoint{3.488916in}{1.850686in}}%
\pgfpathlineto{\pgfqpoint{3.494952in}{1.847150in}}%
\pgfpathlineto{\pgfqpoint{3.494493in}{1.831282in}}%
\pgfpathlineto{\pgfqpoint{3.488098in}{1.831027in}}%
\pgfpathlineto{\pgfqpoint{3.475140in}{1.821145in}}%
\pgfpathlineto{\pgfqpoint{3.471404in}{1.811231in}}%
\pgfpathlineto{\pgfqpoint{3.436723in}{1.802651in}}%
\pgfpathlineto{\pgfqpoint{3.405047in}{1.796199in}}%
\pgfusepath{stroke}%
\end{pgfscope}%
\begin{pgfscope}%
\pgfpathrectangle{\pgfqpoint{0.100000in}{0.100000in}}{\pgfqpoint{3.608454in}{2.310000in}}%
\pgfusepath{clip}%
\pgfsetbuttcap%
\pgfsetroundjoin%
\pgfsetlinewidth{0.010037pt}%
\definecolor{currentstroke}{rgb}{1.000000,1.000000,1.000000}%
\pgfsetstrokecolor{currentstroke}%
\pgfsetdash{}{0pt}%
\pgfpathmoveto{\pgfqpoint{2.320622in}{1.452274in}}%
\pgfpathlineto{\pgfqpoint{2.302315in}{1.470405in}}%
\pgfpathlineto{\pgfqpoint{2.243792in}{1.467029in}}%
\pgfpathlineto{\pgfqpoint{2.152519in}{1.464021in}}%
\pgfpathlineto{\pgfqpoint{2.060697in}{1.465454in}}%
\pgfpathlineto{\pgfqpoint{2.054252in}{1.476692in}}%
\pgfpathlineto{\pgfqpoint{2.056883in}{1.487728in}}%
\pgfpathlineto{\pgfqpoint{2.055634in}{1.506632in}}%
\pgfpathlineto{\pgfqpoint{2.050781in}{1.524944in}}%
\pgfpathlineto{\pgfqpoint{2.051345in}{1.534613in}}%
\pgfpathlineto{\pgfqpoint{2.041030in}{1.545539in}}%
\pgfpathlineto{\pgfqpoint{2.043274in}{1.560741in}}%
\pgfpathlineto{\pgfqpoint{2.027451in}{1.590769in}}%
\pgfpathlineto{\pgfqpoint{2.022737in}{1.616262in}}%
\pgfpathlineto{\pgfqpoint{2.011355in}{1.636997in}}%
\pgfpathlineto{\pgfqpoint{2.022159in}{1.662720in}}%
\pgfpathlineto{\pgfqpoint{2.016778in}{1.674690in}}%
\pgfpathlineto{\pgfqpoint{2.022794in}{1.696330in}}%
\pgfpathlineto{\pgfqpoint{2.080050in}{1.695709in}}%
\pgfpathlineto{\pgfqpoint{2.138106in}{1.695621in}}%
\pgfpathlineto{\pgfqpoint{2.192987in}{1.696427in}}%
\pgfpathlineto{\pgfqpoint{2.253222in}{1.697687in}}%
\pgfpathlineto{\pgfqpoint{2.323886in}{1.700253in}}%
\pgfpathlineto{\pgfqpoint{2.325773in}{1.689994in}}%
\pgfpathlineto{\pgfqpoint{2.333596in}{1.681393in}}%
\pgfpathlineto{\pgfqpoint{2.328643in}{1.672430in}}%
\pgfpathlineto{\pgfqpoint{2.330673in}{1.654936in}}%
\pgfpathlineto{\pgfqpoint{2.335521in}{1.640795in}}%
\pgfpathlineto{\pgfqpoint{2.357502in}{1.633311in}}%
\pgfpathlineto{\pgfqpoint{2.361269in}{1.623298in}}%
\pgfpathlineto{\pgfqpoint{2.373331in}{1.612057in}}%
\pgfpathlineto{\pgfqpoint{2.378251in}{1.600427in}}%
\pgfpathlineto{\pgfqpoint{2.390473in}{1.592626in}}%
\pgfpathlineto{\pgfqpoint{2.392386in}{1.583234in}}%
\pgfpathlineto{\pgfqpoint{2.390006in}{1.569016in}}%
\pgfpathlineto{\pgfqpoint{2.383785in}{1.564755in}}%
\pgfpathlineto{\pgfqpoint{2.381909in}{1.551220in}}%
\pgfpathlineto{\pgfqpoint{2.364032in}{1.540471in}}%
\pgfpathlineto{\pgfqpoint{2.340730in}{1.534582in}}%
\pgfpathlineto{\pgfqpoint{2.338595in}{1.521042in}}%
\pgfpathlineto{\pgfqpoint{2.347584in}{1.511368in}}%
\pgfpathlineto{\pgfqpoint{2.347951in}{1.499205in}}%
\pgfpathlineto{\pgfqpoint{2.340698in}{1.489644in}}%
\pgfpathlineto{\pgfqpoint{2.336891in}{1.475440in}}%
\pgfpathlineto{\pgfqpoint{2.324285in}{1.470738in}}%
\pgfpathlineto{\pgfqpoint{2.325089in}{1.454910in}}%
\pgfpathlineto{\pgfqpoint{2.320622in}{1.452274in}}%
\pgfusepath{stroke}%
\end{pgfscope}%
\begin{pgfscope}%
\pgfpathrectangle{\pgfqpoint{0.100000in}{0.100000in}}{\pgfqpoint{3.608454in}{2.310000in}}%
\pgfusepath{clip}%
\pgfsetbuttcap%
\pgfsetroundjoin%
\pgfsetlinewidth{0.010037pt}%
\definecolor{currentstroke}{rgb}{1.000000,1.000000,1.000000}%
\pgfsetstrokecolor{currentstroke}%
\pgfsetdash{}{0pt}%
\pgfpathmoveto{\pgfqpoint{3.456549in}{1.750456in}}%
\pgfpathlineto{\pgfqpoint{3.439208in}{1.747720in}}%
\pgfpathlineto{\pgfqpoint{3.359557in}{1.729620in}}%
\pgfpathlineto{\pgfqpoint{3.358165in}{1.731729in}}%
\pgfpathlineto{\pgfqpoint{3.359367in}{1.786492in}}%
\pgfpathlineto{\pgfqpoint{3.405047in}{1.796199in}}%
\pgfpathlineto{\pgfqpoint{3.436723in}{1.802651in}}%
\pgfpathlineto{\pgfqpoint{3.471404in}{1.811231in}}%
\pgfpathlineto{\pgfqpoint{3.475140in}{1.821145in}}%
\pgfpathlineto{\pgfqpoint{3.488098in}{1.831027in}}%
\pgfpathlineto{\pgfqpoint{3.494493in}{1.831282in}}%
\pgfpathlineto{\pgfqpoint{3.502901in}{1.816865in}}%
\pgfpathlineto{\pgfqpoint{3.495271in}{1.795821in}}%
\pgfpathlineto{\pgfqpoint{3.494152in}{1.783481in}}%
\pgfpathlineto{\pgfqpoint{3.509594in}{1.784640in}}%
\pgfpathlineto{\pgfqpoint{3.516597in}{1.778721in}}%
\pgfpathlineto{\pgfqpoint{3.532302in}{1.754525in}}%
\pgfpathlineto{\pgfqpoint{3.540107in}{1.751582in}}%
\pgfpathlineto{\pgfqpoint{3.553182in}{1.752673in}}%
\pgfpathlineto{\pgfqpoint{3.562281in}{1.760896in}}%
\pgfpathlineto{\pgfqpoint{3.568329in}{1.753558in}}%
\pgfpathlineto{\pgfqpoint{3.530295in}{1.733487in}}%
\pgfpathlineto{\pgfqpoint{3.529106in}{1.747889in}}%
\pgfpathlineto{\pgfqpoint{3.511829in}{1.725505in}}%
\pgfpathlineto{\pgfqpoint{3.505754in}{1.721649in}}%
\pgfpathlineto{\pgfqpoint{3.497281in}{1.734565in}}%
\pgfpathlineto{\pgfqpoint{3.494962in}{1.736335in}}%
\pgfpathlineto{\pgfqpoint{3.487074in}{1.740513in}}%
\pgfpathlineto{\pgfqpoint{3.480196in}{1.757459in}}%
\pgfpathlineto{\pgfqpoint{3.456549in}{1.750456in}}%
\pgfusepath{stroke}%
\end{pgfscope}%
\begin{pgfscope}%
\pgfpathrectangle{\pgfqpoint{0.100000in}{0.100000in}}{\pgfqpoint{3.608454in}{2.310000in}}%
\pgfusepath{clip}%
\pgfsetbuttcap%
\pgfsetroundjoin%
\pgfsetlinewidth{0.010037pt}%
\definecolor{currentstroke}{rgb}{1.000000,1.000000,1.000000}%
\pgfsetstrokecolor{currentstroke}%
\pgfsetdash{}{0pt}%
\pgfpathmoveto{\pgfqpoint{1.680190in}{1.435164in}}%
\pgfpathlineto{\pgfqpoint{1.686540in}{1.513906in}}%
\pgfpathlineto{\pgfqpoint{1.650576in}{1.516826in}}%
\pgfpathlineto{\pgfqpoint{1.567441in}{1.524892in}}%
\pgfpathlineto{\pgfqpoint{1.575175in}{1.603355in}}%
\pgfpathlineto{\pgfqpoint{1.582921in}{1.682417in}}%
\pgfpathlineto{\pgfqpoint{1.614880in}{1.678966in}}%
\pgfpathlineto{\pgfqpoint{1.696503in}{1.671254in}}%
\pgfpathlineto{\pgfqpoint{1.745900in}{1.667386in}}%
\pgfpathlineto{\pgfqpoint{1.805703in}{1.663958in}}%
\pgfpathlineto{\pgfqpoint{1.904136in}{1.659299in}}%
\pgfpathlineto{\pgfqpoint{1.906300in}{1.654886in}}%
\pgfpathlineto{\pgfqpoint{1.934505in}{1.640932in}}%
\pgfpathlineto{\pgfqpoint{1.943154in}{1.648019in}}%
\pgfpathlineto{\pgfqpoint{1.968073in}{1.647372in}}%
\pgfpathlineto{\pgfqpoint{1.982922in}{1.639593in}}%
\pgfpathlineto{\pgfqpoint{2.006469in}{1.630750in}}%
\pgfpathlineto{\pgfqpoint{2.013067in}{1.617849in}}%
\pgfpathlineto{\pgfqpoint{2.022737in}{1.616262in}}%
\pgfpathlineto{\pgfqpoint{2.027451in}{1.590769in}}%
\pgfpathlineto{\pgfqpoint{2.043274in}{1.560741in}}%
\pgfpathlineto{\pgfqpoint{2.041030in}{1.545539in}}%
\pgfpathlineto{\pgfqpoint{2.051345in}{1.534613in}}%
\pgfpathlineto{\pgfqpoint{2.050781in}{1.524944in}}%
\pgfpathlineto{\pgfqpoint{2.055634in}{1.506632in}}%
\pgfpathlineto{\pgfqpoint{2.056883in}{1.487728in}}%
\pgfpathlineto{\pgfqpoint{2.054252in}{1.476692in}}%
\pgfpathlineto{\pgfqpoint{2.060697in}{1.465454in}}%
\pgfpathlineto{\pgfqpoint{2.069537in}{1.445017in}}%
\pgfpathlineto{\pgfqpoint{2.077998in}{1.436700in}}%
\pgfpathlineto{\pgfqpoint{2.088084in}{1.418676in}}%
\pgfpathlineto{\pgfqpoint{2.059498in}{1.418379in}}%
\pgfpathlineto{\pgfqpoint{1.997690in}{1.419335in}}%
\pgfpathlineto{\pgfqpoint{1.929402in}{1.421428in}}%
\pgfpathlineto{\pgfqpoint{1.860711in}{1.424063in}}%
\pgfpathlineto{\pgfqpoint{1.759716in}{1.429580in}}%
\pgfpathlineto{\pgfqpoint{1.680190in}{1.435164in}}%
\pgfusepath{stroke}%
\end{pgfscope}%
\begin{pgfscope}%
\pgfpathrectangle{\pgfqpoint{0.100000in}{0.100000in}}{\pgfqpoint{3.608454in}{2.310000in}}%
\pgfusepath{clip}%
\pgfsetbuttcap%
\pgfsetroundjoin%
\pgfsetlinewidth{0.010037pt}%
\definecolor{currentstroke}{rgb}{1.000000,1.000000,1.000000}%
\pgfsetstrokecolor{currentstroke}%
\pgfsetdash{}{0pt}%
\pgfpathmoveto{\pgfqpoint{2.995203in}{1.672680in}}%
\pgfpathlineto{\pgfqpoint{3.026828in}{1.702834in}}%
\pgfpathlineto{\pgfqpoint{3.030759in}{1.713553in}}%
\pgfpathlineto{\pgfqpoint{3.040025in}{1.722729in}}%
\pgfpathlineto{\pgfqpoint{3.033072in}{1.736089in}}%
\pgfpathlineto{\pgfqpoint{3.024354in}{1.743903in}}%
\pgfpathlineto{\pgfqpoint{3.021834in}{1.757751in}}%
\pgfpathlineto{\pgfqpoint{3.054283in}{1.771969in}}%
\pgfpathlineto{\pgfqpoint{3.081149in}{1.776467in}}%
\pgfpathlineto{\pgfqpoint{3.095585in}{1.776768in}}%
\pgfpathlineto{\pgfqpoint{3.106582in}{1.771334in}}%
\pgfpathlineto{\pgfqpoint{3.117302in}{1.776186in}}%
\pgfpathlineto{\pgfqpoint{3.143447in}{1.781623in}}%
\pgfpathlineto{\pgfqpoint{3.152481in}{1.788641in}}%
\pgfpathlineto{\pgfqpoint{3.165877in}{1.804177in}}%
\pgfpathlineto{\pgfqpoint{3.178063in}{1.811042in}}%
\pgfpathlineto{\pgfqpoint{3.178922in}{1.817626in}}%
\pgfpathlineto{\pgfqpoint{3.172527in}{1.832648in}}%
\pgfpathlineto{\pgfqpoint{3.177150in}{1.841488in}}%
\pgfpathlineto{\pgfqpoint{3.170925in}{1.850993in}}%
\pgfpathlineto{\pgfqpoint{3.161391in}{1.851657in}}%
\pgfpathlineto{\pgfqpoint{3.185252in}{1.880369in}}%
\pgfpathlineto{\pgfqpoint{3.188086in}{1.891311in}}%
\pgfpathlineto{\pgfqpoint{3.206865in}{1.919218in}}%
\pgfpathlineto{\pgfqpoint{3.224297in}{1.934319in}}%
\pgfpathlineto{\pgfqpoint{3.236268in}{1.940625in}}%
\pgfpathlineto{\pgfqpoint{3.275437in}{1.949500in}}%
\pgfpathlineto{\pgfqpoint{3.312149in}{1.959523in}}%
\pgfpathlineto{\pgfqpoint{3.317011in}{1.943790in}}%
\pgfpathlineto{\pgfqpoint{3.317356in}{1.929202in}}%
\pgfpathlineto{\pgfqpoint{3.325730in}{1.915960in}}%
\pgfpathlineto{\pgfqpoint{3.328311in}{1.902477in}}%
\pgfpathlineto{\pgfqpoint{3.325093in}{1.884184in}}%
\pgfpathlineto{\pgfqpoint{3.334894in}{1.859510in}}%
\pgfpathlineto{\pgfqpoint{3.346067in}{1.845790in}}%
\pgfpathlineto{\pgfqpoint{3.354710in}{1.801193in}}%
\pgfpathlineto{\pgfqpoint{3.359367in}{1.786492in}}%
\pgfpathlineto{\pgfqpoint{3.358165in}{1.731729in}}%
\pgfpathlineto{\pgfqpoint{3.359557in}{1.729620in}}%
\pgfpathlineto{\pgfqpoint{3.369731in}{1.670722in}}%
\pgfpathlineto{\pgfqpoint{3.375437in}{1.665355in}}%
\pgfpathlineto{\pgfqpoint{3.363164in}{1.653431in}}%
\pgfpathlineto{\pgfqpoint{3.369219in}{1.646587in}}%
\pgfpathlineto{\pgfqpoint{3.363897in}{1.636235in}}%
\pgfpathlineto{\pgfqpoint{3.363942in}{1.631828in}}%
\pgfpathlineto{\pgfqpoint{3.357288in}{1.627852in}}%
\pgfpathlineto{\pgfqpoint{3.354047in}{1.619061in}}%
\pgfpathlineto{\pgfqpoint{3.355075in}{1.643233in}}%
\pgfpathlineto{\pgfqpoint{3.334450in}{1.648555in}}%
\pgfpathlineto{\pgfqpoint{3.302196in}{1.659557in}}%
\pgfpathlineto{\pgfqpoint{3.297634in}{1.664976in}}%
\pgfpathlineto{\pgfqpoint{3.284155in}{1.666275in}}%
\pgfpathlineto{\pgfqpoint{3.276402in}{1.674341in}}%
\pgfpathlineto{\pgfqpoint{3.272219in}{1.690405in}}%
\pgfpathlineto{\pgfqpoint{3.261228in}{1.692404in}}%
\pgfpathlineto{\pgfqpoint{3.253878in}{1.700831in}}%
\pgfpathlineto{\pgfqpoint{3.160304in}{1.681966in}}%
\pgfpathlineto{\pgfqpoint{3.115582in}{1.672717in}}%
\pgfpathlineto{\pgfqpoint{3.047681in}{1.660324in}}%
\pgfpathlineto{\pgfqpoint{2.998785in}{1.652075in}}%
\pgfpathlineto{\pgfqpoint{2.995203in}{1.672680in}}%
\pgfusepath{stroke}%
\end{pgfscope}%
\begin{pgfscope}%
\pgfpathrectangle{\pgfqpoint{0.100000in}{0.100000in}}{\pgfqpoint{3.608454in}{2.310000in}}%
\pgfusepath{clip}%
\pgfsetbuttcap%
\pgfsetroundjoin%
\pgfsetlinewidth{0.010037pt}%
\definecolor{currentstroke}{rgb}{1.000000,1.000000,1.000000}%
\pgfsetstrokecolor{currentstroke}%
\pgfsetdash{}{0pt}%
\pgfpathmoveto{\pgfqpoint{3.370867in}{1.614152in}}%
\pgfpathlineto{\pgfqpoint{3.356395in}{1.609644in}}%
\pgfpathlineto{\pgfqpoint{3.353960in}{1.613788in}}%
\pgfpathlineto{\pgfqpoint{3.358606in}{1.627701in}}%
\pgfpathlineto{\pgfqpoint{3.367208in}{1.629071in}}%
\pgfpathlineto{\pgfqpoint{3.366438in}{1.633451in}}%
\pgfpathlineto{\pgfqpoint{3.374162in}{1.640026in}}%
\pgfpathlineto{\pgfqpoint{3.396494in}{1.645263in}}%
\pgfpathlineto{\pgfqpoint{3.406435in}{1.653192in}}%
\pgfpathlineto{\pgfqpoint{3.428768in}{1.659799in}}%
\pgfpathlineto{\pgfqpoint{3.438954in}{1.657382in}}%
\pgfpathlineto{\pgfqpoint{3.447523in}{1.668034in}}%
\pgfpathlineto{\pgfqpoint{3.460473in}{1.669442in}}%
\pgfpathlineto{\pgfqpoint{3.438379in}{1.648663in}}%
\pgfpathlineto{\pgfqpoint{3.370867in}{1.614152in}}%
\pgfusepath{stroke}%
\end{pgfscope}%
\begin{pgfscope}%
\pgfpathrectangle{\pgfqpoint{0.100000in}{0.100000in}}{\pgfqpoint{3.608454in}{2.310000in}}%
\pgfusepath{clip}%
\pgfsetbuttcap%
\pgfsetroundjoin%
\pgfsetlinewidth{0.010037pt}%
\definecolor{currentstroke}{rgb}{1.000000,1.000000,1.000000}%
\pgfsetstrokecolor{currentstroke}%
\pgfsetdash{}{0pt}%
\pgfpathmoveto{\pgfqpoint{3.045737in}{1.476948in}}%
\pgfpathlineto{\pgfqpoint{2.983156in}{1.466600in}}%
\pgfpathlineto{\pgfqpoint{2.971804in}{1.538119in}}%
\pgfpathlineto{\pgfqpoint{2.954949in}{1.643541in}}%
\pgfpathlineto{\pgfqpoint{2.995203in}{1.672680in}}%
\pgfpathlineto{\pgfqpoint{2.998785in}{1.652075in}}%
\pgfpathlineto{\pgfqpoint{3.047681in}{1.660324in}}%
\pgfpathlineto{\pgfqpoint{3.115582in}{1.672717in}}%
\pgfpathlineto{\pgfqpoint{3.160304in}{1.681966in}}%
\pgfpathlineto{\pgfqpoint{3.253878in}{1.700831in}}%
\pgfpathlineto{\pgfqpoint{3.261228in}{1.692404in}}%
\pgfpathlineto{\pgfqpoint{3.272219in}{1.690405in}}%
\pgfpathlineto{\pgfqpoint{3.276402in}{1.674341in}}%
\pgfpathlineto{\pgfqpoint{3.284155in}{1.666275in}}%
\pgfpathlineto{\pgfqpoint{3.297634in}{1.664976in}}%
\pgfpathlineto{\pgfqpoint{3.302196in}{1.659557in}}%
\pgfpathlineto{\pgfqpoint{3.297564in}{1.655376in}}%
\pgfpathlineto{\pgfqpoint{3.293392in}{1.640582in}}%
\pgfpathlineto{\pgfqpoint{3.283615in}{1.623960in}}%
\pgfpathlineto{\pgfqpoint{3.290178in}{1.616731in}}%
\pgfpathlineto{\pgfqpoint{3.284807in}{1.608986in}}%
\pgfpathlineto{\pgfqpoint{3.287837in}{1.592001in}}%
\pgfpathlineto{\pgfqpoint{3.297508in}{1.583085in}}%
\pgfpathlineto{\pgfqpoint{3.320457in}{1.568600in}}%
\pgfpathlineto{\pgfqpoint{3.301976in}{1.548162in}}%
\pgfpathlineto{\pgfqpoint{3.301717in}{1.540404in}}%
\pgfpathlineto{\pgfqpoint{3.286671in}{1.530400in}}%
\pgfpathlineto{\pgfqpoint{3.270035in}{1.528522in}}%
\pgfpathlineto{\pgfqpoint{3.265920in}{1.519827in}}%
\pgfpathlineto{\pgfqpoint{3.219645in}{1.509814in}}%
\pgfpathlineto{\pgfqpoint{3.165648in}{1.498939in}}%
\pgfpathlineto{\pgfqpoint{3.128532in}{1.492293in}}%
\pgfpathlineto{\pgfqpoint{3.045737in}{1.476948in}}%
\pgfusepath{stroke}%
\end{pgfscope}%
\begin{pgfscope}%
\pgfpathrectangle{\pgfqpoint{0.100000in}{0.100000in}}{\pgfqpoint{3.608454in}{2.310000in}}%
\pgfusepath{clip}%
\pgfsetbuttcap%
\pgfsetroundjoin%
\pgfsetlinewidth{0.010037pt}%
\definecolor{currentstroke}{rgb}{1.000000,1.000000,1.000000}%
\pgfsetstrokecolor{currentstroke}%
\pgfsetdash{}{0pt}%
\pgfpathmoveto{\pgfqpoint{3.359557in}{1.729620in}}%
\pgfpathlineto{\pgfqpoint{3.439208in}{1.747720in}}%
\pgfpathlineto{\pgfqpoint{3.456549in}{1.750456in}}%
\pgfpathlineto{\pgfqpoint{3.468026in}{1.705331in}}%
\pgfpathlineto{\pgfqpoint{3.466209in}{1.697251in}}%
\pgfpathlineto{\pgfqpoint{3.429341in}{1.682977in}}%
\pgfpathlineto{\pgfqpoint{3.407332in}{1.677989in}}%
\pgfpathlineto{\pgfqpoint{3.397980in}{1.666798in}}%
\pgfpathlineto{\pgfqpoint{3.369219in}{1.646587in}}%
\pgfpathlineto{\pgfqpoint{3.363164in}{1.653431in}}%
\pgfpathlineto{\pgfqpoint{3.375437in}{1.665355in}}%
\pgfpathlineto{\pgfqpoint{3.369731in}{1.670722in}}%
\pgfpathlineto{\pgfqpoint{3.359557in}{1.729620in}}%
\pgfusepath{stroke}%
\end{pgfscope}%
\begin{pgfscope}%
\pgfpathrectangle{\pgfqpoint{0.100000in}{0.100000in}}{\pgfqpoint{3.608454in}{2.310000in}}%
\pgfusepath{clip}%
\pgfsetbuttcap%
\pgfsetroundjoin%
\pgfsetlinewidth{0.010037pt}%
\definecolor{currentstroke}{rgb}{1.000000,1.000000,1.000000}%
\pgfsetstrokecolor{currentstroke}%
\pgfsetdash{}{0pt}%
\pgfpathmoveto{\pgfqpoint{3.456549in}{1.750456in}}%
\pgfpathlineto{\pgfqpoint{3.480196in}{1.757459in}}%
\pgfpathlineto{\pgfqpoint{3.487074in}{1.740513in}}%
\pgfpathlineto{\pgfqpoint{3.494962in}{1.736335in}}%
\pgfpathlineto{\pgfqpoint{3.485228in}{1.729197in}}%
\pgfpathlineto{\pgfqpoint{3.486456in}{1.708233in}}%
\pgfpathlineto{\pgfqpoint{3.466209in}{1.697251in}}%
\pgfpathlineto{\pgfqpoint{3.468026in}{1.705331in}}%
\pgfpathlineto{\pgfqpoint{3.456549in}{1.750456in}}%
\pgfusepath{stroke}%
\end{pgfscope}%
\begin{pgfscope}%
\pgfpathrectangle{\pgfqpoint{0.100000in}{0.100000in}}{\pgfqpoint{3.608454in}{2.310000in}}%
\pgfusepath{clip}%
\pgfsetbuttcap%
\pgfsetroundjoin%
\pgfsetlinewidth{0.010037pt}%
\definecolor{currentstroke}{rgb}{1.000000,1.000000,1.000000}%
\pgfsetstrokecolor{currentstroke}%
\pgfsetdash{}{0pt}%
\pgfpathmoveto{\pgfqpoint{3.284005in}{1.523011in}}%
\pgfpathlineto{\pgfqpoint{3.286671in}{1.530400in}}%
\pgfpathlineto{\pgfqpoint{3.301717in}{1.540404in}}%
\pgfpathlineto{\pgfqpoint{3.301976in}{1.548162in}}%
\pgfpathlineto{\pgfqpoint{3.320457in}{1.568600in}}%
\pgfpathlineto{\pgfqpoint{3.297508in}{1.583085in}}%
\pgfpathlineto{\pgfqpoint{3.287837in}{1.592001in}}%
\pgfpathlineto{\pgfqpoint{3.284807in}{1.608986in}}%
\pgfpathlineto{\pgfqpoint{3.290178in}{1.616731in}}%
\pgfpathlineto{\pgfqpoint{3.283615in}{1.623960in}}%
\pgfpathlineto{\pgfqpoint{3.293392in}{1.640582in}}%
\pgfpathlineto{\pgfqpoint{3.297564in}{1.655376in}}%
\pgfpathlineto{\pgfqpoint{3.302196in}{1.659557in}}%
\pgfpathlineto{\pgfqpoint{3.334450in}{1.648555in}}%
\pgfpathlineto{\pgfqpoint{3.355075in}{1.643233in}}%
\pgfpathlineto{\pgfqpoint{3.354047in}{1.619061in}}%
\pgfpathlineto{\pgfqpoint{3.341519in}{1.600731in}}%
\pgfpathlineto{\pgfqpoint{3.351842in}{1.598035in}}%
\pgfpathlineto{\pgfqpoint{3.362561in}{1.590168in}}%
\pgfpathlineto{\pgfqpoint{3.363194in}{1.568654in}}%
\pgfpathlineto{\pgfqpoint{3.359949in}{1.553423in}}%
\pgfpathlineto{\pgfqpoint{3.362108in}{1.540913in}}%
\pgfpathlineto{\pgfqpoint{3.343475in}{1.498697in}}%
\pgfpathlineto{\pgfqpoint{3.336779in}{1.478984in}}%
\pgfpathlineto{\pgfqpoint{3.327429in}{1.488566in}}%
\pgfpathlineto{\pgfqpoint{3.315030in}{1.486933in}}%
\pgfpathlineto{\pgfqpoint{3.284000in}{1.504862in}}%
\pgfpathlineto{\pgfqpoint{3.280825in}{1.514462in}}%
\pgfpathlineto{\pgfqpoint{3.284005in}{1.523011in}}%
\pgfusepath{stroke}%
\end{pgfscope}%
\begin{pgfscope}%
\pgfpathrectangle{\pgfqpoint{0.100000in}{0.100000in}}{\pgfqpoint{3.608454in}{2.310000in}}%
\pgfusepath{clip}%
\pgfsetbuttcap%
\pgfsetroundjoin%
\pgfsetlinewidth{0.010037pt}%
\definecolor{currentstroke}{rgb}{1.000000,1.000000,1.000000}%
\pgfsetstrokecolor{currentstroke}%
\pgfsetdash{}{0pt}%
\pgfpathmoveto{\pgfqpoint{2.541376in}{1.261906in}}%
\pgfpathlineto{\pgfqpoint{2.539926in}{1.281490in}}%
\pgfpathlineto{\pgfqpoint{2.545564in}{1.291207in}}%
\pgfpathlineto{\pgfqpoint{2.541848in}{1.295967in}}%
\pgfpathlineto{\pgfqpoint{2.555400in}{1.311555in}}%
\pgfpathlineto{\pgfqpoint{2.554624in}{1.315312in}}%
\pgfpathlineto{\pgfqpoint{2.567791in}{1.342507in}}%
\pgfpathlineto{\pgfqpoint{2.565129in}{1.355623in}}%
\pgfpathlineto{\pgfqpoint{2.555581in}{1.369354in}}%
\pgfpathlineto{\pgfqpoint{2.562209in}{1.386059in}}%
\pgfpathlineto{\pgfqpoint{2.553618in}{1.495995in}}%
\pgfpathlineto{\pgfqpoint{2.547432in}{1.573118in}}%
\pgfpathlineto{\pgfqpoint{2.555978in}{1.566728in}}%
\pgfpathlineto{\pgfqpoint{2.565512in}{1.566900in}}%
\pgfpathlineto{\pgfqpoint{2.588031in}{1.579945in}}%
\pgfpathlineto{\pgfqpoint{2.657102in}{1.586389in}}%
\pgfpathlineto{\pgfqpoint{2.708226in}{1.591778in}}%
\pgfpathlineto{\pgfqpoint{2.708677in}{1.586775in}}%
\pgfpathlineto{\pgfqpoint{2.720348in}{1.481089in}}%
\pgfpathlineto{\pgfqpoint{2.730393in}{1.382751in}}%
\pgfpathlineto{\pgfqpoint{2.726063in}{1.378135in}}%
\pgfpathlineto{\pgfqpoint{2.732714in}{1.366675in}}%
\pgfpathlineto{\pgfqpoint{2.732680in}{1.358417in}}%
\pgfpathlineto{\pgfqpoint{2.723189in}{1.356342in}}%
\pgfpathlineto{\pgfqpoint{2.712569in}{1.348381in}}%
\pgfpathlineto{\pgfqpoint{2.705378in}{1.351515in}}%
\pgfpathlineto{\pgfqpoint{2.694616in}{1.346419in}}%
\pgfpathlineto{\pgfqpoint{2.697946in}{1.336184in}}%
\pgfpathlineto{\pgfqpoint{2.686884in}{1.325901in}}%
\pgfpathlineto{\pgfqpoint{2.683829in}{1.314006in}}%
\pgfpathlineto{\pgfqpoint{2.676223in}{1.312054in}}%
\pgfpathlineto{\pgfqpoint{2.670547in}{1.303046in}}%
\pgfpathlineto{\pgfqpoint{2.671290in}{1.293974in}}%
\pgfpathlineto{\pgfqpoint{2.664612in}{1.287594in}}%
\pgfpathlineto{\pgfqpoint{2.654564in}{1.288579in}}%
\pgfpathlineto{\pgfqpoint{2.646926in}{1.298367in}}%
\pgfpathlineto{\pgfqpoint{2.633983in}{1.288941in}}%
\pgfpathlineto{\pgfqpoint{2.634892in}{1.280705in}}%
\pgfpathlineto{\pgfqpoint{2.620497in}{1.275892in}}%
\pgfpathlineto{\pgfqpoint{2.616881in}{1.281949in}}%
\pgfpathlineto{\pgfqpoint{2.602868in}{1.275081in}}%
\pgfpathlineto{\pgfqpoint{2.597748in}{1.265248in}}%
\pgfpathlineto{\pgfqpoint{2.580872in}{1.275331in}}%
\pgfpathlineto{\pgfqpoint{2.554124in}{1.268651in}}%
\pgfpathlineto{\pgfqpoint{2.541376in}{1.261906in}}%
\pgfusepath{stroke}%
\end{pgfscope}%
\begin{pgfscope}%
\pgfpathrectangle{\pgfqpoint{0.100000in}{0.100000in}}{\pgfqpoint{3.608454in}{2.310000in}}%
\pgfusepath{clip}%
\pgfsetbuttcap%
\pgfsetroundjoin%
\pgfsetlinewidth{0.010037pt}%
\definecolor{currentstroke}{rgb}{1.000000,1.000000,1.000000}%
\pgfsetstrokecolor{currentstroke}%
\pgfsetdash{}{0pt}%
\pgfpathmoveto{\pgfqpoint{0.655019in}{1.776916in}}%
\pgfpathlineto{\pgfqpoint{0.693698in}{1.766153in}}%
\pgfpathlineto{\pgfqpoint{0.757813in}{1.749621in}}%
\pgfpathlineto{\pgfqpoint{0.824275in}{1.732591in}}%
\pgfpathlineto{\pgfqpoint{0.885533in}{1.717818in}}%
\pgfpathlineto{\pgfqpoint{0.938464in}{1.705834in}}%
\pgfpathlineto{\pgfqpoint{0.995094in}{1.693381in}}%
\pgfpathlineto{\pgfqpoint{0.978732in}{1.616166in}}%
\pgfpathlineto{\pgfqpoint{0.964156in}{1.547595in}}%
\pgfpathlineto{\pgfqpoint{0.940244in}{1.436956in}}%
\pgfpathlineto{\pgfqpoint{0.922291in}{1.353447in}}%
\pgfpathlineto{\pgfqpoint{0.912593in}{1.306845in}}%
\pgfpathlineto{\pgfqpoint{0.900156in}{1.246360in}}%
\pgfpathlineto{\pgfqpoint{0.891552in}{1.234089in}}%
\pgfpathlineto{\pgfqpoint{0.884626in}{1.233679in}}%
\pgfpathlineto{\pgfqpoint{0.879687in}{1.244390in}}%
\pgfpathlineto{\pgfqpoint{0.868345in}{1.248280in}}%
\pgfpathlineto{\pgfqpoint{0.856124in}{1.246917in}}%
\pgfpathlineto{\pgfqpoint{0.852655in}{1.238156in}}%
\pgfpathlineto{\pgfqpoint{0.852293in}{1.207731in}}%
\pgfpathlineto{\pgfqpoint{0.848633in}{1.200785in}}%
\pgfpathlineto{\pgfqpoint{0.850725in}{1.176360in}}%
\pgfpathlineto{\pgfqpoint{0.843082in}{1.160088in}}%
\pgfpathlineto{\pgfqpoint{0.781106in}{1.255521in}}%
\pgfpathlineto{\pgfqpoint{0.720071in}{1.348255in}}%
\pgfpathlineto{\pgfqpoint{0.688300in}{1.396883in}}%
\pgfpathlineto{\pgfqpoint{0.661796in}{1.438856in}}%
\pgfpathlineto{\pgfqpoint{0.628400in}{1.490776in}}%
\pgfpathlineto{\pgfqpoint{0.590596in}{1.549016in}}%
\pgfpathlineto{\pgfqpoint{0.606139in}{1.604297in}}%
\pgfpathlineto{\pgfqpoint{0.637416in}{1.715162in}}%
\pgfpathlineto{\pgfqpoint{0.655019in}{1.776916in}}%
\pgfusepath{stroke}%
\end{pgfscope}%
\begin{pgfscope}%
\pgfpathrectangle{\pgfqpoint{0.100000in}{0.100000in}}{\pgfqpoint{3.608454in}{2.310000in}}%
\pgfusepath{clip}%
\pgfsetbuttcap%
\pgfsetroundjoin%
\pgfsetlinewidth{0.010037pt}%
\definecolor{currentstroke}{rgb}{1.000000,1.000000,1.000000}%
\pgfsetstrokecolor{currentstroke}%
\pgfsetdash{}{0pt}%
\pgfpathmoveto{\pgfqpoint{0.995094in}{1.693381in}}%
\pgfpathlineto{\pgfqpoint{1.055637in}{1.681339in}}%
\pgfpathlineto{\pgfqpoint{1.167721in}{1.659701in}}%
\pgfpathlineto{\pgfqpoint{1.153691in}{1.581874in}}%
\pgfpathlineto{\pgfqpoint{1.215397in}{1.571309in}}%
\pgfpathlineto{\pgfqpoint{1.271614in}{1.562313in}}%
\pgfpathlineto{\pgfqpoint{1.261846in}{1.500784in}}%
\pgfpathlineto{\pgfqpoint{1.251495in}{1.434433in}}%
\pgfpathlineto{\pgfqpoint{1.237637in}{1.347314in}}%
\pgfpathlineto{\pgfqpoint{1.237278in}{1.340011in}}%
\pgfpathlineto{\pgfqpoint{1.222893in}{1.249723in}}%
\pgfpathlineto{\pgfqpoint{1.163676in}{1.258901in}}%
\pgfpathlineto{\pgfqpoint{1.133501in}{1.264961in}}%
\pgfpathlineto{\pgfqpoint{1.024434in}{1.284134in}}%
\pgfpathlineto{\pgfqpoint{0.983361in}{1.292233in}}%
\pgfpathlineto{\pgfqpoint{0.912593in}{1.306845in}}%
\pgfpathlineto{\pgfqpoint{0.922291in}{1.353447in}}%
\pgfpathlineto{\pgfqpoint{0.940244in}{1.436956in}}%
\pgfpathlineto{\pgfqpoint{0.964156in}{1.547595in}}%
\pgfpathlineto{\pgfqpoint{0.978732in}{1.616166in}}%
\pgfpathlineto{\pgfqpoint{0.995094in}{1.693381in}}%
\pgfusepath{stroke}%
\end{pgfscope}%
\begin{pgfscope}%
\pgfpathrectangle{\pgfqpoint{0.100000in}{0.100000in}}{\pgfqpoint{3.608454in}{2.310000in}}%
\pgfusepath{clip}%
\pgfsetbuttcap%
\pgfsetroundjoin%
\pgfsetlinewidth{0.010037pt}%
\definecolor{currentstroke}{rgb}{1.000000,1.000000,1.000000}%
\pgfsetstrokecolor{currentstroke}%
\pgfsetdash{}{0pt}%
\pgfpathmoveto{\pgfqpoint{0.418363in}{1.850356in}}%
\pgfpathlineto{\pgfqpoint{0.439875in}{1.842853in}}%
\pgfpathlineto{\pgfqpoint{0.526049in}{1.816024in}}%
\pgfpathlineto{\pgfqpoint{0.605321in}{1.791073in}}%
\pgfpathlineto{\pgfqpoint{0.655019in}{1.776916in}}%
\pgfpathlineto{\pgfqpoint{0.637416in}{1.715162in}}%
\pgfpathlineto{\pgfqpoint{0.606139in}{1.604297in}}%
\pgfpathlineto{\pgfqpoint{0.590596in}{1.549016in}}%
\pgfpathlineto{\pgfqpoint{0.628400in}{1.490776in}}%
\pgfpathlineto{\pgfqpoint{0.661796in}{1.438856in}}%
\pgfpathlineto{\pgfqpoint{0.688300in}{1.396883in}}%
\pgfpathlineto{\pgfqpoint{0.720071in}{1.348255in}}%
\pgfpathlineto{\pgfqpoint{0.781106in}{1.255521in}}%
\pgfpathlineto{\pgfqpoint{0.843082in}{1.160088in}}%
\pgfpathlineto{\pgfqpoint{0.840594in}{1.150624in}}%
\pgfpathlineto{\pgfqpoint{0.848070in}{1.135552in}}%
\pgfpathlineto{\pgfqpoint{0.849548in}{1.114934in}}%
\pgfpathlineto{\pgfqpoint{0.860405in}{1.104932in}}%
\pgfpathlineto{\pgfqpoint{0.860834in}{1.096874in}}%
\pgfpathlineto{\pgfqpoint{0.841395in}{1.087727in}}%
\pgfpathlineto{\pgfqpoint{0.832133in}{1.078570in}}%
\pgfpathlineto{\pgfqpoint{0.829245in}{1.058349in}}%
\pgfpathlineto{\pgfqpoint{0.824557in}{1.047321in}}%
\pgfpathlineto{\pgfqpoint{0.814680in}{1.038025in}}%
\pgfpathlineto{\pgfqpoint{0.804856in}{1.013854in}}%
\pgfpathlineto{\pgfqpoint{0.818669in}{1.001263in}}%
\pgfpathlineto{\pgfqpoint{0.816886in}{0.990883in}}%
\pgfpathlineto{\pgfqpoint{0.805560in}{0.983661in}}%
\pgfpathlineto{\pgfqpoint{0.797751in}{0.984935in}}%
\pgfpathlineto{\pgfqpoint{0.705575in}{0.997727in}}%
\pgfpathlineto{\pgfqpoint{0.637495in}{1.007385in}}%
\pgfpathlineto{\pgfqpoint{0.640472in}{1.018376in}}%
\pgfpathlineto{\pgfqpoint{0.636040in}{1.036625in}}%
\pgfpathlineto{\pgfqpoint{0.635593in}{1.055032in}}%
\pgfpathlineto{\pgfqpoint{0.632697in}{1.065804in}}%
\pgfpathlineto{\pgfqpoint{0.623775in}{1.081205in}}%
\pgfpathlineto{\pgfqpoint{0.598111in}{1.116738in}}%
\pgfpathlineto{\pgfqpoint{0.589705in}{1.121092in}}%
\pgfpathlineto{\pgfqpoint{0.578875in}{1.121023in}}%
\pgfpathlineto{\pgfqpoint{0.581349in}{1.132233in}}%
\pgfpathlineto{\pgfqpoint{0.576221in}{1.146250in}}%
\pgfpathlineto{\pgfqpoint{0.551032in}{1.153217in}}%
\pgfpathlineto{\pgfqpoint{0.535677in}{1.166118in}}%
\pgfpathlineto{\pgfqpoint{0.534404in}{1.174011in}}%
\pgfpathlineto{\pgfqpoint{0.516724in}{1.193554in}}%
\pgfpathlineto{\pgfqpoint{0.499894in}{1.197294in}}%
\pgfpathlineto{\pgfqpoint{0.484301in}{1.207225in}}%
\pgfpathlineto{\pgfqpoint{0.463793in}{1.210661in}}%
\pgfpathlineto{\pgfqpoint{0.455031in}{1.223907in}}%
\pgfpathlineto{\pgfqpoint{0.463352in}{1.244873in}}%
\pgfpathlineto{\pgfqpoint{0.460807in}{1.249588in}}%
\pgfpathlineto{\pgfqpoint{0.467737in}{1.267066in}}%
\pgfpathlineto{\pgfqpoint{0.455409in}{1.276367in}}%
\pgfpathlineto{\pgfqpoint{0.459410in}{1.293230in}}%
\pgfpathlineto{\pgfqpoint{0.452825in}{1.297541in}}%
\pgfpathlineto{\pgfqpoint{0.447140in}{1.313479in}}%
\pgfpathlineto{\pgfqpoint{0.440280in}{1.318340in}}%
\pgfpathlineto{\pgfqpoint{0.439763in}{1.329853in}}%
\pgfpathlineto{\pgfqpoint{0.434389in}{1.337971in}}%
\pgfpathlineto{\pgfqpoint{0.426289in}{1.365312in}}%
\pgfpathlineto{\pgfqpoint{0.417427in}{1.378407in}}%
\pgfpathlineto{\pgfqpoint{0.419354in}{1.400637in}}%
\pgfpathlineto{\pgfqpoint{0.429775in}{1.402874in}}%
\pgfpathlineto{\pgfqpoint{0.436561in}{1.414932in}}%
\pgfpathlineto{\pgfqpoint{0.432461in}{1.428009in}}%
\pgfpathlineto{\pgfqpoint{0.421377in}{1.430192in}}%
\pgfpathlineto{\pgfqpoint{0.415847in}{1.436309in}}%
\pgfpathlineto{\pgfqpoint{0.406892in}{1.458808in}}%
\pgfpathlineto{\pgfqpoint{0.411078in}{1.466888in}}%
\pgfpathlineto{\pgfqpoint{0.408109in}{1.481933in}}%
\pgfpathlineto{\pgfqpoint{0.414682in}{1.501419in}}%
\pgfpathlineto{\pgfqpoint{0.422383in}{1.494243in}}%
\pgfpathlineto{\pgfqpoint{0.418888in}{1.485771in}}%
\pgfpathlineto{\pgfqpoint{0.432233in}{1.472341in}}%
\pgfpathlineto{\pgfqpoint{0.431381in}{1.492284in}}%
\pgfpathlineto{\pgfqpoint{0.425663in}{1.497632in}}%
\pgfpathlineto{\pgfqpoint{0.432199in}{1.515167in}}%
\pgfpathlineto{\pgfqpoint{0.456921in}{1.512375in}}%
\pgfpathlineto{\pgfqpoint{0.453594in}{1.518895in}}%
\pgfpathlineto{\pgfqpoint{0.437278in}{1.518254in}}%
\pgfpathlineto{\pgfqpoint{0.429516in}{1.528130in}}%
\pgfpathlineto{\pgfqpoint{0.419745in}{1.519360in}}%
\pgfpathlineto{\pgfqpoint{0.414561in}{1.504701in}}%
\pgfpathlineto{\pgfqpoint{0.400724in}{1.524426in}}%
\pgfpathlineto{\pgfqpoint{0.395395in}{1.528017in}}%
\pgfpathlineto{\pgfqpoint{0.397425in}{1.549485in}}%
\pgfpathlineto{\pgfqpoint{0.393207in}{1.562145in}}%
\pgfpathlineto{\pgfqpoint{0.385562in}{1.574039in}}%
\pgfpathlineto{\pgfqpoint{0.369905in}{1.610482in}}%
\pgfpathlineto{\pgfqpoint{0.375024in}{1.618541in}}%
\pgfpathlineto{\pgfqpoint{0.374964in}{1.644010in}}%
\pgfpathlineto{\pgfqpoint{0.383414in}{1.658215in}}%
\pgfpathlineto{\pgfqpoint{0.385352in}{1.680415in}}%
\pgfpathlineto{\pgfqpoint{0.377384in}{1.705851in}}%
\pgfpathlineto{\pgfqpoint{0.366813in}{1.722042in}}%
\pgfpathlineto{\pgfqpoint{0.368698in}{1.736652in}}%
\pgfpathlineto{\pgfqpoint{0.398383in}{1.772011in}}%
\pgfpathlineto{\pgfqpoint{0.399858in}{1.784073in}}%
\pgfpathlineto{\pgfqpoint{0.413234in}{1.807081in}}%
\pgfpathlineto{\pgfqpoint{0.415090in}{1.828890in}}%
\pgfpathlineto{\pgfqpoint{0.410789in}{1.834455in}}%
\pgfpathlineto{\pgfqpoint{0.418363in}{1.850356in}}%
\pgfusepath{stroke}%
\end{pgfscope}%
\begin{pgfscope}%
\pgfpathrectangle{\pgfqpoint{0.100000in}{0.100000in}}{\pgfqpoint{3.608454in}{2.310000in}}%
\pgfusepath{clip}%
\pgfsetbuttcap%
\pgfsetroundjoin%
\pgfsetlinewidth{0.010037pt}%
\definecolor{currentstroke}{rgb}{1.000000,1.000000,1.000000}%
\pgfsetstrokecolor{currentstroke}%
\pgfsetdash{}{0pt}%
\pgfpathmoveto{\pgfqpoint{2.730393in}{1.382751in}}%
\pgfpathlineto{\pgfqpoint{2.720348in}{1.481089in}}%
\pgfpathlineto{\pgfqpoint{2.708677in}{1.586775in}}%
\pgfpathlineto{\pgfqpoint{2.785105in}{1.598157in}}%
\pgfpathlineto{\pgfqpoint{2.805389in}{1.592878in}}%
\pgfpathlineto{\pgfqpoint{2.815117in}{1.587147in}}%
\pgfpathlineto{\pgfqpoint{2.827308in}{1.588736in}}%
\pgfpathlineto{\pgfqpoint{2.843381in}{1.579248in}}%
\pgfpathlineto{\pgfqpoint{2.873320in}{1.593370in}}%
\pgfpathlineto{\pgfqpoint{2.889841in}{1.593841in}}%
\pgfpathlineto{\pgfqpoint{2.909135in}{1.615366in}}%
\pgfpathlineto{\pgfqpoint{2.928765in}{1.628459in}}%
\pgfpathlineto{\pgfqpoint{2.954949in}{1.643541in}}%
\pgfpathlineto{\pgfqpoint{2.971804in}{1.538119in}}%
\pgfpathlineto{\pgfqpoint{2.963973in}{1.531348in}}%
\pgfpathlineto{\pgfqpoint{2.969029in}{1.525129in}}%
\pgfpathlineto{\pgfqpoint{2.971043in}{1.511497in}}%
\pgfpathlineto{\pgfqpoint{2.967120in}{1.498692in}}%
\pgfpathlineto{\pgfqpoint{2.966601in}{1.479976in}}%
\pgfpathlineto{\pgfqpoint{2.962917in}{1.455616in}}%
\pgfpathlineto{\pgfqpoint{2.944276in}{1.433890in}}%
\pgfpathlineto{\pgfqpoint{2.936429in}{1.429278in}}%
\pgfpathlineto{\pgfqpoint{2.930326in}{1.433249in}}%
\pgfpathlineto{\pgfqpoint{2.915224in}{1.412541in}}%
\pgfpathlineto{\pgfqpoint{2.916638in}{1.392602in}}%
\pgfpathlineto{\pgfqpoint{2.899819in}{1.397397in}}%
\pgfpathlineto{\pgfqpoint{2.891855in}{1.377484in}}%
\pgfpathlineto{\pgfqpoint{2.895881in}{1.363373in}}%
\pgfpathlineto{\pgfqpoint{2.889579in}{1.361288in}}%
\pgfpathlineto{\pgfqpoint{2.888697in}{1.350139in}}%
\pgfpathlineto{\pgfqpoint{2.873228in}{1.345641in}}%
\pgfpathlineto{\pgfqpoint{2.865190in}{1.354643in}}%
\pgfpathlineto{\pgfqpoint{2.854847in}{1.358134in}}%
\pgfpathlineto{\pgfqpoint{2.852417in}{1.367262in}}%
\pgfpathlineto{\pgfqpoint{2.843065in}{1.365657in}}%
\pgfpathlineto{\pgfqpoint{2.836938in}{1.357261in}}%
\pgfpathlineto{\pgfqpoint{2.828173in}{1.354308in}}%
\pgfpathlineto{\pgfqpoint{2.812691in}{1.360250in}}%
\pgfpathlineto{\pgfqpoint{2.804128in}{1.353181in}}%
\pgfpathlineto{\pgfqpoint{2.791958in}{1.361544in}}%
\pgfpathlineto{\pgfqpoint{2.768610in}{1.364112in}}%
\pgfpathlineto{\pgfqpoint{2.761568in}{1.379305in}}%
\pgfpathlineto{\pgfqpoint{2.752646in}{1.386037in}}%
\pgfpathlineto{\pgfqpoint{2.743996in}{1.381728in}}%
\pgfpathlineto{\pgfqpoint{2.730393in}{1.382751in}}%
\pgfusepath{stroke}%
\end{pgfscope}%
\begin{pgfscope}%
\pgfpathrectangle{\pgfqpoint{0.100000in}{0.100000in}}{\pgfqpoint{3.608454in}{2.310000in}}%
\pgfusepath{clip}%
\pgfsetbuttcap%
\pgfsetroundjoin%
\pgfsetlinewidth{0.010037pt}%
\definecolor{currentstroke}{rgb}{1.000000,1.000000,1.000000}%
\pgfsetstrokecolor{currentstroke}%
\pgfsetdash{}{0pt}%
\pgfpathmoveto{\pgfqpoint{2.477364in}{1.192570in}}%
\pgfpathlineto{\pgfqpoint{2.468410in}{1.199844in}}%
\pgfpathlineto{\pgfqpoint{2.461108in}{1.196383in}}%
\pgfpathlineto{\pgfqpoint{2.451774in}{1.213813in}}%
\pgfpathlineto{\pgfqpoint{2.456543in}{1.224768in}}%
\pgfpathlineto{\pgfqpoint{2.449670in}{1.237099in}}%
\pgfpathlineto{\pgfqpoint{2.449301in}{1.246806in}}%
\pgfpathlineto{\pgfqpoint{2.440014in}{1.250265in}}%
\pgfpathlineto{\pgfqpoint{2.435708in}{1.257580in}}%
\pgfpathlineto{\pgfqpoint{2.417550in}{1.266715in}}%
\pgfpathlineto{\pgfqpoint{2.401780in}{1.277970in}}%
\pgfpathlineto{\pgfqpoint{2.394453in}{1.286467in}}%
\pgfpathlineto{\pgfqpoint{2.394301in}{1.296831in}}%
\pgfpathlineto{\pgfqpoint{2.404106in}{1.316728in}}%
\pgfpathlineto{\pgfqpoint{2.407122in}{1.331947in}}%
\pgfpathlineto{\pgfqpoint{2.399134in}{1.340547in}}%
\pgfpathlineto{\pgfqpoint{2.388555in}{1.343787in}}%
\pgfpathlineto{\pgfqpoint{2.375719in}{1.336702in}}%
\pgfpathlineto{\pgfqpoint{2.370287in}{1.348866in}}%
\pgfpathlineto{\pgfqpoint{2.367545in}{1.365354in}}%
\pgfpathlineto{\pgfqpoint{2.348630in}{1.380053in}}%
\pgfpathlineto{\pgfqpoint{2.344850in}{1.386576in}}%
\pgfpathlineto{\pgfqpoint{2.327573in}{1.401342in}}%
\pgfpathlineto{\pgfqpoint{2.322158in}{1.412077in}}%
\pgfpathlineto{\pgfqpoint{2.317277in}{1.433365in}}%
\pgfpathlineto{\pgfqpoint{2.320622in}{1.452274in}}%
\pgfpathlineto{\pgfqpoint{2.325089in}{1.454910in}}%
\pgfpathlineto{\pgfqpoint{2.324285in}{1.470738in}}%
\pgfpathlineto{\pgfqpoint{2.336891in}{1.475440in}}%
\pgfpathlineto{\pgfqpoint{2.340698in}{1.489644in}}%
\pgfpathlineto{\pgfqpoint{2.347951in}{1.499205in}}%
\pgfpathlineto{\pgfqpoint{2.347584in}{1.511368in}}%
\pgfpathlineto{\pgfqpoint{2.338595in}{1.521042in}}%
\pgfpathlineto{\pgfqpoint{2.340730in}{1.534582in}}%
\pgfpathlineto{\pgfqpoint{2.364032in}{1.540471in}}%
\pgfpathlineto{\pgfqpoint{2.381909in}{1.551220in}}%
\pgfpathlineto{\pgfqpoint{2.383785in}{1.564755in}}%
\pgfpathlineto{\pgfqpoint{2.390006in}{1.569016in}}%
\pgfpathlineto{\pgfqpoint{2.392386in}{1.583234in}}%
\pgfpathlineto{\pgfqpoint{2.390473in}{1.592626in}}%
\pgfpathlineto{\pgfqpoint{2.378251in}{1.600427in}}%
\pgfpathlineto{\pgfqpoint{2.373331in}{1.612057in}}%
\pgfpathlineto{\pgfqpoint{2.361269in}{1.623298in}}%
\pgfpathlineto{\pgfqpoint{2.460390in}{1.627522in}}%
\pgfpathlineto{\pgfqpoint{2.526890in}{1.632246in}}%
\pgfpathlineto{\pgfqpoint{2.525674in}{1.618258in}}%
\pgfpathlineto{\pgfqpoint{2.537007in}{1.598963in}}%
\pgfpathlineto{\pgfqpoint{2.541765in}{1.582478in}}%
\pgfpathlineto{\pgfqpoint{2.547432in}{1.573118in}}%
\pgfpathlineto{\pgfqpoint{2.553618in}{1.495995in}}%
\pgfpathlineto{\pgfqpoint{2.562209in}{1.386059in}}%
\pgfpathlineto{\pgfqpoint{2.555581in}{1.369354in}}%
\pgfpathlineto{\pgfqpoint{2.565129in}{1.355623in}}%
\pgfpathlineto{\pgfqpoint{2.567791in}{1.342507in}}%
\pgfpathlineto{\pgfqpoint{2.554624in}{1.315312in}}%
\pgfpathlineto{\pgfqpoint{2.555400in}{1.311555in}}%
\pgfpathlineto{\pgfqpoint{2.541848in}{1.295967in}}%
\pgfpathlineto{\pgfqpoint{2.545564in}{1.291207in}}%
\pgfpathlineto{\pgfqpoint{2.539926in}{1.281490in}}%
\pgfpathlineto{\pgfqpoint{2.541376in}{1.261906in}}%
\pgfpathlineto{\pgfqpoint{2.534528in}{1.249888in}}%
\pgfpathlineto{\pgfqpoint{2.540098in}{1.235688in}}%
\pgfpathlineto{\pgfqpoint{2.516761in}{1.227971in}}%
\pgfpathlineto{\pgfqpoint{2.514610in}{1.219580in}}%
\pgfpathlineto{\pgfqpoint{2.520983in}{1.208943in}}%
\pgfpathlineto{\pgfqpoint{2.518051in}{1.202010in}}%
\pgfpathlineto{\pgfqpoint{2.493026in}{1.210574in}}%
\pgfpathlineto{\pgfqpoint{2.484770in}{1.211436in}}%
\pgfpathlineto{\pgfqpoint{2.474476in}{1.198417in}}%
\pgfpathlineto{\pgfqpoint{2.477364in}{1.192570in}}%
\pgfusepath{stroke}%
\end{pgfscope}%
\begin{pgfscope}%
\pgfpathrectangle{\pgfqpoint{0.100000in}{0.100000in}}{\pgfqpoint{3.608454in}{2.310000in}}%
\pgfusepath{clip}%
\pgfsetbuttcap%
\pgfsetroundjoin%
\pgfsetlinewidth{0.010037pt}%
\definecolor{currentstroke}{rgb}{1.000000,1.000000,1.000000}%
\pgfsetstrokecolor{currentstroke}%
\pgfsetdash{}{0pt}%
\pgfpathmoveto{\pgfqpoint{3.205784in}{1.431832in}}%
\pgfpathlineto{\pgfqpoint{3.198897in}{1.442058in}}%
\pgfpathlineto{\pgfqpoint{3.202783in}{1.447781in}}%
\pgfpathlineto{\pgfqpoint{3.212300in}{1.441354in}}%
\pgfpathlineto{\pgfqpoint{3.205784in}{1.431832in}}%
\pgfusepath{stroke}%
\end{pgfscope}%
\begin{pgfscope}%
\pgfpathrectangle{\pgfqpoint{0.100000in}{0.100000in}}{\pgfqpoint{3.608454in}{2.310000in}}%
\pgfusepath{clip}%
\pgfsetbuttcap%
\pgfsetroundjoin%
\pgfsetlinewidth{0.010037pt}%
\definecolor{currentstroke}{rgb}{1.000000,1.000000,1.000000}%
\pgfsetstrokecolor{currentstroke}%
\pgfsetdash{}{0pt}%
\pgfpathmoveto{\pgfqpoint{3.265920in}{1.519827in}}%
\pgfpathlineto{\pgfqpoint{3.270035in}{1.528522in}}%
\pgfpathlineto{\pgfqpoint{3.286671in}{1.530400in}}%
\pgfpathlineto{\pgfqpoint{3.284005in}{1.523011in}}%
\pgfpathlineto{\pgfqpoint{3.278517in}{1.513571in}}%
\pgfpathlineto{\pgfqpoint{3.282236in}{1.502324in}}%
\pgfpathlineto{\pgfqpoint{3.296914in}{1.488849in}}%
\pgfpathlineto{\pgfqpoint{3.300321in}{1.474658in}}%
\pgfpathlineto{\pgfqpoint{3.317232in}{1.456987in}}%
\pgfpathlineto{\pgfqpoint{3.323867in}{1.457744in}}%
\pgfpathlineto{\pgfqpoint{3.332131in}{1.431216in}}%
\pgfpathlineto{\pgfqpoint{3.330776in}{1.430952in}}%
\pgfpathlineto{\pgfqpoint{3.329269in}{1.430656in}}%
\pgfpathlineto{\pgfqpoint{3.292410in}{1.423600in}}%
\pgfpathlineto{\pgfqpoint{3.272690in}{1.493733in}}%
\pgfpathlineto{\pgfqpoint{3.265920in}{1.519827in}}%
\pgfusepath{stroke}%
\end{pgfscope}%
\begin{pgfscope}%
\pgfpathrectangle{\pgfqpoint{0.100000in}{0.100000in}}{\pgfqpoint{3.608454in}{2.310000in}}%
\pgfusepath{clip}%
\pgfsetbuttcap%
\pgfsetroundjoin%
\pgfsetlinewidth{0.010037pt}%
\definecolor{currentstroke}{rgb}{1.000000,1.000000,1.000000}%
\pgfsetstrokecolor{currentstroke}%
\pgfsetdash{}{0pt}%
\pgfpathmoveto{\pgfqpoint{2.921647in}{1.282104in}}%
\pgfpathlineto{\pgfqpoint{2.910000in}{1.282532in}}%
\pgfpathlineto{\pgfqpoint{2.899268in}{1.289925in}}%
\pgfpathlineto{\pgfqpoint{2.884769in}{1.312394in}}%
\pgfpathlineto{\pgfqpoint{2.872427in}{1.324293in}}%
\pgfpathlineto{\pgfqpoint{2.875653in}{1.333484in}}%
\pgfpathlineto{\pgfqpoint{2.873228in}{1.345641in}}%
\pgfpathlineto{\pgfqpoint{2.888697in}{1.350139in}}%
\pgfpathlineto{\pgfqpoint{2.889579in}{1.361288in}}%
\pgfpathlineto{\pgfqpoint{2.895881in}{1.363373in}}%
\pgfpathlineto{\pgfqpoint{2.891855in}{1.377484in}}%
\pgfpathlineto{\pgfqpoint{2.899819in}{1.397397in}}%
\pgfpathlineto{\pgfqpoint{2.916638in}{1.392602in}}%
\pgfpathlineto{\pgfqpoint{2.915224in}{1.412541in}}%
\pgfpathlineto{\pgfqpoint{2.930326in}{1.433249in}}%
\pgfpathlineto{\pgfqpoint{2.936429in}{1.429278in}}%
\pgfpathlineto{\pgfqpoint{2.944276in}{1.433890in}}%
\pgfpathlineto{\pgfqpoint{2.962917in}{1.455616in}}%
\pgfpathlineto{\pgfqpoint{2.966601in}{1.479976in}}%
\pgfpathlineto{\pgfqpoint{2.967120in}{1.498692in}}%
\pgfpathlineto{\pgfqpoint{2.971043in}{1.511497in}}%
\pgfpathlineto{\pgfqpoint{2.969029in}{1.525129in}}%
\pgfpathlineto{\pgfqpoint{2.963973in}{1.531348in}}%
\pgfpathlineto{\pgfqpoint{2.971804in}{1.538119in}}%
\pgfpathlineto{\pgfqpoint{2.983156in}{1.466600in}}%
\pgfpathlineto{\pgfqpoint{3.045737in}{1.476948in}}%
\pgfpathlineto{\pgfqpoint{3.052224in}{1.436127in}}%
\pgfpathlineto{\pgfqpoint{3.074880in}{1.463071in}}%
\pgfpathlineto{\pgfqpoint{3.080218in}{1.460384in}}%
\pgfpathlineto{\pgfqpoint{3.088071in}{1.475996in}}%
\pgfpathlineto{\pgfqpoint{3.098951in}{1.471101in}}%
\pgfpathlineto{\pgfqpoint{3.108437in}{1.472032in}}%
\pgfpathlineto{\pgfqpoint{3.114696in}{1.482885in}}%
\pgfpathlineto{\pgfqpoint{3.123780in}{1.488922in}}%
\pgfpathlineto{\pgfqpoint{3.136410in}{1.483605in}}%
\pgfpathlineto{\pgfqpoint{3.144720in}{1.485442in}}%
\pgfpathlineto{\pgfqpoint{3.156626in}{1.464838in}}%
\pgfpathlineto{\pgfqpoint{3.153183in}{1.449239in}}%
\pgfpathlineto{\pgfqpoint{3.122088in}{1.466802in}}%
\pgfpathlineto{\pgfqpoint{3.116263in}{1.452384in}}%
\pgfpathlineto{\pgfqpoint{3.118185in}{1.445751in}}%
\pgfpathlineto{\pgfqpoint{3.104950in}{1.423326in}}%
\pgfpathlineto{\pgfqpoint{3.098705in}{1.418866in}}%
\pgfpathlineto{\pgfqpoint{3.095884in}{1.408977in}}%
\pgfpathlineto{\pgfqpoint{3.085235in}{1.410012in}}%
\pgfpathlineto{\pgfqpoint{3.077588in}{1.383022in}}%
\pgfpathlineto{\pgfqpoint{3.073308in}{1.376829in}}%
\pgfpathlineto{\pgfqpoint{3.062309in}{1.378894in}}%
\pgfpathlineto{\pgfqpoint{3.051058in}{1.387407in}}%
\pgfpathlineto{\pgfqpoint{3.050668in}{1.374341in}}%
\pgfpathlineto{\pgfqpoint{3.039777in}{1.352383in}}%
\pgfpathlineto{\pgfqpoint{3.037068in}{1.336751in}}%
\pgfpathlineto{\pgfqpoint{3.028711in}{1.326341in}}%
\pgfpathlineto{\pgfqpoint{3.022375in}{1.309715in}}%
\pgfpathlineto{\pgfqpoint{3.022064in}{1.294324in}}%
\pgfpathlineto{\pgfqpoint{3.012280in}{1.291449in}}%
\pgfpathlineto{\pgfqpoint{3.001180in}{1.282737in}}%
\pgfpathlineto{\pgfqpoint{2.990493in}{1.281077in}}%
\pgfpathlineto{\pgfqpoint{2.988030in}{1.273703in}}%
\pgfpathlineto{\pgfqpoint{2.970810in}{1.266151in}}%
\pgfpathlineto{\pgfqpoint{2.961186in}{1.272585in}}%
\pgfpathlineto{\pgfqpoint{2.950430in}{1.260367in}}%
\pgfpathlineto{\pgfqpoint{2.931940in}{1.263968in}}%
\pgfpathlineto{\pgfqpoint{2.920600in}{1.276793in}}%
\pgfpathlineto{\pgfqpoint{2.921647in}{1.282104in}}%
\pgfusepath{stroke}%
\end{pgfscope}%
\begin{pgfscope}%
\pgfpathrectangle{\pgfqpoint{0.100000in}{0.100000in}}{\pgfqpoint{3.608454in}{2.310000in}}%
\pgfusepath{clip}%
\pgfsetbuttcap%
\pgfsetroundjoin%
\pgfsetlinewidth{0.010037pt}%
\definecolor{currentstroke}{rgb}{1.000000,1.000000,1.000000}%
\pgfsetstrokecolor{currentstroke}%
\pgfsetdash{}{0pt}%
\pgfpathmoveto{\pgfqpoint{3.329269in}{1.430656in}}%
\pgfpathlineto{\pgfqpoint{3.328797in}{1.416244in}}%
\pgfpathlineto{\pgfqpoint{3.323250in}{1.409159in}}%
\pgfpathlineto{\pgfqpoint{3.319710in}{1.393431in}}%
\pgfpathlineto{\pgfqpoint{3.303727in}{1.386181in}}%
\pgfpathlineto{\pgfqpoint{3.290316in}{1.384068in}}%
\pgfpathlineto{\pgfqpoint{3.294224in}{1.394412in}}%
\pgfpathlineto{\pgfqpoint{3.273575in}{1.403080in}}%
\pgfpathlineto{\pgfqpoint{3.256800in}{1.413932in}}%
\pgfpathlineto{\pgfqpoint{3.267185in}{1.430155in}}%
\pgfpathlineto{\pgfqpoint{3.258862in}{1.442765in}}%
\pgfpathlineto{\pgfqpoint{3.259842in}{1.453641in}}%
\pgfpathlineto{\pgfqpoint{3.253824in}{1.456627in}}%
\pgfpathlineto{\pgfqpoint{3.248940in}{1.474298in}}%
\pgfpathlineto{\pgfqpoint{3.260489in}{1.498252in}}%
\pgfpathlineto{\pgfqpoint{3.251813in}{1.502164in}}%
\pgfpathlineto{\pgfqpoint{3.249564in}{1.490346in}}%
\pgfpathlineto{\pgfqpoint{3.237178in}{1.487056in}}%
\pgfpathlineto{\pgfqpoint{3.238481in}{1.448173in}}%
\pgfpathlineto{\pgfqpoint{3.236216in}{1.435660in}}%
\pgfpathlineto{\pgfqpoint{3.242428in}{1.417824in}}%
\pgfpathlineto{\pgfqpoint{3.251985in}{1.409243in}}%
\pgfpathlineto{\pgfqpoint{3.247651in}{1.403855in}}%
\pgfpathlineto{\pgfqpoint{3.257407in}{1.395982in}}%
\pgfpathlineto{\pgfqpoint{3.260933in}{1.383203in}}%
\pgfpathlineto{\pgfqpoint{3.243052in}{1.393785in}}%
\pgfpathlineto{\pgfqpoint{3.231736in}{1.392411in}}%
\pgfpathlineto{\pgfqpoint{3.222950in}{1.403297in}}%
\pgfpathlineto{\pgfqpoint{3.214003in}{1.404357in}}%
\pgfpathlineto{\pgfqpoint{3.201288in}{1.398904in}}%
\pgfpathlineto{\pgfqpoint{3.196357in}{1.405699in}}%
\pgfpathlineto{\pgfqpoint{3.202838in}{1.419960in}}%
\pgfpathlineto{\pgfqpoint{3.205784in}{1.431832in}}%
\pgfpathlineto{\pgfqpoint{3.212300in}{1.441354in}}%
\pgfpathlineto{\pgfqpoint{3.202783in}{1.447781in}}%
\pgfpathlineto{\pgfqpoint{3.198897in}{1.442058in}}%
\pgfpathlineto{\pgfqpoint{3.189390in}{1.447867in}}%
\pgfpathlineto{\pgfqpoint{3.172555in}{1.451737in}}%
\pgfpathlineto{\pgfqpoint{3.174081in}{1.460247in}}%
\pgfpathlineto{\pgfqpoint{3.166450in}{1.465183in}}%
\pgfpathlineto{\pgfqpoint{3.156626in}{1.464838in}}%
\pgfpathlineto{\pgfqpoint{3.144720in}{1.485442in}}%
\pgfpathlineto{\pgfqpoint{3.136410in}{1.483605in}}%
\pgfpathlineto{\pgfqpoint{3.123780in}{1.488922in}}%
\pgfpathlineto{\pgfqpoint{3.114696in}{1.482885in}}%
\pgfpathlineto{\pgfqpoint{3.108437in}{1.472032in}}%
\pgfpathlineto{\pgfqpoint{3.098951in}{1.471101in}}%
\pgfpathlineto{\pgfqpoint{3.088071in}{1.475996in}}%
\pgfpathlineto{\pgfqpoint{3.080218in}{1.460384in}}%
\pgfpathlineto{\pgfqpoint{3.074880in}{1.463071in}}%
\pgfpathlineto{\pgfqpoint{3.052224in}{1.436127in}}%
\pgfpathlineto{\pgfqpoint{3.045737in}{1.476948in}}%
\pgfpathlineto{\pgfqpoint{3.128532in}{1.492293in}}%
\pgfpathlineto{\pgfqpoint{3.165648in}{1.498939in}}%
\pgfpathlineto{\pgfqpoint{3.219645in}{1.509814in}}%
\pgfpathlineto{\pgfqpoint{3.265920in}{1.519827in}}%
\pgfpathlineto{\pgfqpoint{3.272690in}{1.493733in}}%
\pgfpathlineto{\pgfqpoint{3.292410in}{1.423600in}}%
\pgfpathlineto{\pgfqpoint{3.329269in}{1.430656in}}%
\pgfusepath{stroke}%
\end{pgfscope}%
\begin{pgfscope}%
\pgfpathrectangle{\pgfqpoint{0.100000in}{0.100000in}}{\pgfqpoint{3.608454in}{2.310000in}}%
\pgfusepath{clip}%
\pgfsetbuttcap%
\pgfsetroundjoin%
\pgfsetlinewidth{0.010037pt}%
\definecolor{currentstroke}{rgb}{1.000000,1.000000,1.000000}%
\pgfsetstrokecolor{currentstroke}%
\pgfsetdash{}{0pt}%
\pgfpathmoveto{\pgfqpoint{1.662648in}{1.197867in}}%
\pgfpathlineto{\pgfqpoint{1.602243in}{1.203632in}}%
\pgfpathlineto{\pgfqpoint{1.539577in}{1.209185in}}%
\pgfpathlineto{\pgfqpoint{1.467164in}{1.216739in}}%
\pgfpathlineto{\pgfqpoint{1.359578in}{1.229535in}}%
\pgfpathlineto{\pgfqpoint{1.321411in}{1.235415in}}%
\pgfpathlineto{\pgfqpoint{1.222893in}{1.249723in}}%
\pgfpathlineto{\pgfqpoint{1.237278in}{1.340011in}}%
\pgfpathlineto{\pgfqpoint{1.237637in}{1.347314in}}%
\pgfpathlineto{\pgfqpoint{1.251495in}{1.434433in}}%
\pgfpathlineto{\pgfqpoint{1.261846in}{1.500784in}}%
\pgfpathlineto{\pgfqpoint{1.271614in}{1.562313in}}%
\pgfpathlineto{\pgfqpoint{1.338353in}{1.552733in}}%
\pgfpathlineto{\pgfqpoint{1.400578in}{1.543802in}}%
\pgfpathlineto{\pgfqpoint{1.514964in}{1.529684in}}%
\pgfpathlineto{\pgfqpoint{1.567441in}{1.524892in}}%
\pgfpathlineto{\pgfqpoint{1.650576in}{1.516826in}}%
\pgfpathlineto{\pgfqpoint{1.686540in}{1.513906in}}%
\pgfpathlineto{\pgfqpoint{1.680190in}{1.435164in}}%
\pgfpathlineto{\pgfqpoint{1.674454in}{1.359348in}}%
\pgfpathlineto{\pgfqpoint{1.669837in}{1.297616in}}%
\pgfpathlineto{\pgfqpoint{1.662648in}{1.197867in}}%
\pgfusepath{stroke}%
\end{pgfscope}%
\begin{pgfscope}%
\pgfpathrectangle{\pgfqpoint{0.100000in}{0.100000in}}{\pgfqpoint{3.608454in}{2.310000in}}%
\pgfusepath{clip}%
\pgfsetbuttcap%
\pgfsetroundjoin%
\pgfsetlinewidth{0.010037pt}%
\definecolor{currentstroke}{rgb}{1.000000,1.000000,1.000000}%
\pgfsetstrokecolor{currentstroke}%
\pgfsetdash{}{0pt}%
\pgfpathmoveto{\pgfqpoint{2.461763in}{1.153105in}}%
\pgfpathlineto{\pgfqpoint{2.463799in}{1.162233in}}%
\pgfpathlineto{\pgfqpoint{2.474354in}{1.160178in}}%
\pgfpathlineto{\pgfqpoint{2.477364in}{1.192570in}}%
\pgfpathlineto{\pgfqpoint{2.474476in}{1.198417in}}%
\pgfpathlineto{\pgfqpoint{2.484770in}{1.211436in}}%
\pgfpathlineto{\pgfqpoint{2.493026in}{1.210574in}}%
\pgfpathlineto{\pgfqpoint{2.518051in}{1.202010in}}%
\pgfpathlineto{\pgfqpoint{2.520983in}{1.208943in}}%
\pgfpathlineto{\pgfqpoint{2.514610in}{1.219580in}}%
\pgfpathlineto{\pgfqpoint{2.516761in}{1.227971in}}%
\pgfpathlineto{\pgfqpoint{2.540098in}{1.235688in}}%
\pgfpathlineto{\pgfqpoint{2.534528in}{1.249888in}}%
\pgfpathlineto{\pgfqpoint{2.541376in}{1.261906in}}%
\pgfpathlineto{\pgfqpoint{2.554124in}{1.268651in}}%
\pgfpathlineto{\pgfqpoint{2.580872in}{1.275331in}}%
\pgfpathlineto{\pgfqpoint{2.597748in}{1.265248in}}%
\pgfpathlineto{\pgfqpoint{2.602868in}{1.275081in}}%
\pgfpathlineto{\pgfqpoint{2.616881in}{1.281949in}}%
\pgfpathlineto{\pgfqpoint{2.620497in}{1.275892in}}%
\pgfpathlineto{\pgfqpoint{2.634892in}{1.280705in}}%
\pgfpathlineto{\pgfqpoint{2.633983in}{1.288941in}}%
\pgfpathlineto{\pgfqpoint{2.646926in}{1.298367in}}%
\pgfpathlineto{\pgfqpoint{2.654564in}{1.288579in}}%
\pgfpathlineto{\pgfqpoint{2.664612in}{1.287594in}}%
\pgfpathlineto{\pgfqpoint{2.671290in}{1.293974in}}%
\pgfpathlineto{\pgfqpoint{2.670547in}{1.303046in}}%
\pgfpathlineto{\pgfqpoint{2.676223in}{1.312054in}}%
\pgfpathlineto{\pgfqpoint{2.683829in}{1.314006in}}%
\pgfpathlineto{\pgfqpoint{2.686884in}{1.325901in}}%
\pgfpathlineto{\pgfqpoint{2.697946in}{1.336184in}}%
\pgfpathlineto{\pgfqpoint{2.694616in}{1.346419in}}%
\pgfpathlineto{\pgfqpoint{2.705378in}{1.351515in}}%
\pgfpathlineto{\pgfqpoint{2.712569in}{1.348381in}}%
\pgfpathlineto{\pgfqpoint{2.723189in}{1.356342in}}%
\pgfpathlineto{\pgfqpoint{2.732680in}{1.358417in}}%
\pgfpathlineto{\pgfqpoint{2.732714in}{1.366675in}}%
\pgfpathlineto{\pgfqpoint{2.726063in}{1.378135in}}%
\pgfpathlineto{\pgfqpoint{2.730393in}{1.382751in}}%
\pgfpathlineto{\pgfqpoint{2.743996in}{1.381728in}}%
\pgfpathlineto{\pgfqpoint{2.752646in}{1.386037in}}%
\pgfpathlineto{\pgfqpoint{2.761568in}{1.379305in}}%
\pgfpathlineto{\pgfqpoint{2.768610in}{1.364112in}}%
\pgfpathlineto{\pgfqpoint{2.791958in}{1.361544in}}%
\pgfpathlineto{\pgfqpoint{2.804128in}{1.353181in}}%
\pgfpathlineto{\pgfqpoint{2.812691in}{1.360250in}}%
\pgfpathlineto{\pgfqpoint{2.828173in}{1.354308in}}%
\pgfpathlineto{\pgfqpoint{2.836938in}{1.357261in}}%
\pgfpathlineto{\pgfqpoint{2.843065in}{1.365657in}}%
\pgfpathlineto{\pgfqpoint{2.852417in}{1.367262in}}%
\pgfpathlineto{\pgfqpoint{2.854847in}{1.358134in}}%
\pgfpathlineto{\pgfqpoint{2.865190in}{1.354643in}}%
\pgfpathlineto{\pgfqpoint{2.873228in}{1.345641in}}%
\pgfpathlineto{\pgfqpoint{2.875653in}{1.333484in}}%
\pgfpathlineto{\pgfqpoint{2.872427in}{1.324293in}}%
\pgfpathlineto{\pgfqpoint{2.884769in}{1.312394in}}%
\pgfpathlineto{\pgfqpoint{2.899268in}{1.289925in}}%
\pgfpathlineto{\pgfqpoint{2.910000in}{1.282532in}}%
\pgfpathlineto{\pgfqpoint{2.921647in}{1.282104in}}%
\pgfpathlineto{\pgfqpoint{2.900141in}{1.257427in}}%
\pgfpathlineto{\pgfqpoint{2.878972in}{1.242490in}}%
\pgfpathlineto{\pgfqpoint{2.871181in}{1.230629in}}%
\pgfpathlineto{\pgfqpoint{2.871313in}{1.224191in}}%
\pgfpathlineto{\pgfqpoint{2.859847in}{1.219258in}}%
\pgfpathlineto{\pgfqpoint{2.856561in}{1.209974in}}%
\pgfpathlineto{\pgfqpoint{2.832674in}{1.200623in}}%
\pgfpathlineto{\pgfqpoint{2.824207in}{1.194550in}}%
\pgfpathlineto{\pgfqpoint{2.823062in}{1.193254in}}%
\pgfpathlineto{\pgfqpoint{2.754353in}{1.186739in}}%
\pgfpathlineto{\pgfqpoint{2.712861in}{1.183250in}}%
\pgfpathlineto{\pgfqpoint{2.644773in}{1.179396in}}%
\pgfpathlineto{\pgfqpoint{2.559948in}{1.171005in}}%
\pgfpathlineto{\pgfqpoint{2.545934in}{1.172945in}}%
\pgfpathlineto{\pgfqpoint{2.548849in}{1.158640in}}%
\pgfpathlineto{\pgfqpoint{2.461763in}{1.153105in}}%
\pgfusepath{stroke}%
\end{pgfscope}%
\begin{pgfscope}%
\pgfpathrectangle{\pgfqpoint{0.100000in}{0.100000in}}{\pgfqpoint{3.608454in}{2.310000in}}%
\pgfusepath{clip}%
\pgfsetbuttcap%
\pgfsetroundjoin%
\pgfsetlinewidth{0.010037pt}%
\definecolor{currentstroke}{rgb}{1.000000,1.000000,1.000000}%
\pgfsetstrokecolor{currentstroke}%
\pgfsetdash{}{0pt}%
\pgfpathmoveto{\pgfqpoint{1.662648in}{1.197867in}}%
\pgfpathlineto{\pgfqpoint{1.669837in}{1.297616in}}%
\pgfpathlineto{\pgfqpoint{1.674454in}{1.359348in}}%
\pgfpathlineto{\pgfqpoint{1.680190in}{1.435164in}}%
\pgfpathlineto{\pgfqpoint{1.759716in}{1.429580in}}%
\pgfpathlineto{\pgfqpoint{1.860711in}{1.424063in}}%
\pgfpathlineto{\pgfqpoint{1.929402in}{1.421428in}}%
\pgfpathlineto{\pgfqpoint{1.997690in}{1.419335in}}%
\pgfpathlineto{\pgfqpoint{2.059498in}{1.418379in}}%
\pgfpathlineto{\pgfqpoint{2.088084in}{1.418676in}}%
\pgfpathlineto{\pgfqpoint{2.100667in}{1.408407in}}%
\pgfpathlineto{\pgfqpoint{2.110523in}{1.410477in}}%
\pgfpathlineto{\pgfqpoint{2.110832in}{1.401521in}}%
\pgfpathlineto{\pgfqpoint{2.103512in}{1.386028in}}%
\pgfpathlineto{\pgfqpoint{2.104306in}{1.376240in}}%
\pgfpathlineto{\pgfqpoint{2.111314in}{1.369785in}}%
\pgfpathlineto{\pgfqpoint{2.117751in}{1.356337in}}%
\pgfpathlineto{\pgfqpoint{2.130810in}{1.348616in}}%
\pgfpathlineto{\pgfqpoint{2.130373in}{1.297922in}}%
\pgfpathlineto{\pgfqpoint{2.130758in}{1.181350in}}%
\pgfpathlineto{\pgfqpoint{2.043230in}{1.181753in}}%
\pgfpathlineto{\pgfqpoint{1.951062in}{1.183366in}}%
\pgfpathlineto{\pgfqpoint{1.883211in}{1.185760in}}%
\pgfpathlineto{\pgfqpoint{1.826575in}{1.187945in}}%
\pgfpathlineto{\pgfqpoint{1.731159in}{1.193571in}}%
\pgfpathlineto{\pgfqpoint{1.662648in}{1.197867in}}%
\pgfusepath{stroke}%
\end{pgfscope}%
\begin{pgfscope}%
\pgfpathrectangle{\pgfqpoint{0.100000in}{0.100000in}}{\pgfqpoint{3.608454in}{2.310000in}}%
\pgfusepath{clip}%
\pgfsetbuttcap%
\pgfsetroundjoin%
\pgfsetlinewidth{0.010037pt}%
\definecolor{currentstroke}{rgb}{1.000000,1.000000,1.000000}%
\pgfsetstrokecolor{currentstroke}%
\pgfsetdash{}{0pt}%
\pgfpathmoveto{\pgfqpoint{2.950506in}{1.211075in}}%
\pgfpathlineto{\pgfqpoint{2.853565in}{1.197457in}}%
\pgfpathlineto{\pgfqpoint{2.824207in}{1.194550in}}%
\pgfpathlineto{\pgfqpoint{2.832674in}{1.200623in}}%
\pgfpathlineto{\pgfqpoint{2.856561in}{1.209974in}}%
\pgfpathlineto{\pgfqpoint{2.859847in}{1.219258in}}%
\pgfpathlineto{\pgfqpoint{2.871313in}{1.224191in}}%
\pgfpathlineto{\pgfqpoint{2.871181in}{1.230629in}}%
\pgfpathlineto{\pgfqpoint{2.878972in}{1.242490in}}%
\pgfpathlineto{\pgfqpoint{2.900141in}{1.257427in}}%
\pgfpathlineto{\pgfqpoint{2.921647in}{1.282104in}}%
\pgfpathlineto{\pgfqpoint{2.920600in}{1.276793in}}%
\pgfpathlineto{\pgfqpoint{2.931940in}{1.263968in}}%
\pgfpathlineto{\pgfqpoint{2.950430in}{1.260367in}}%
\pgfpathlineto{\pgfqpoint{2.961186in}{1.272585in}}%
\pgfpathlineto{\pgfqpoint{2.970810in}{1.266151in}}%
\pgfpathlineto{\pgfqpoint{2.988030in}{1.273703in}}%
\pgfpathlineto{\pgfqpoint{2.990493in}{1.281077in}}%
\pgfpathlineto{\pgfqpoint{3.001180in}{1.282737in}}%
\pgfpathlineto{\pgfqpoint{3.012280in}{1.291449in}}%
\pgfpathlineto{\pgfqpoint{3.022064in}{1.294324in}}%
\pgfpathlineto{\pgfqpoint{3.022375in}{1.309715in}}%
\pgfpathlineto{\pgfqpoint{3.028711in}{1.326341in}}%
\pgfpathlineto{\pgfqpoint{3.037068in}{1.336751in}}%
\pgfpathlineto{\pgfqpoint{3.039777in}{1.352383in}}%
\pgfpathlineto{\pgfqpoint{3.050668in}{1.374341in}}%
\pgfpathlineto{\pgfqpoint{3.051058in}{1.387407in}}%
\pgfpathlineto{\pgfqpoint{3.062309in}{1.378894in}}%
\pgfpathlineto{\pgfqpoint{3.073308in}{1.376829in}}%
\pgfpathlineto{\pgfqpoint{3.077588in}{1.383022in}}%
\pgfpathlineto{\pgfqpoint{3.085235in}{1.410012in}}%
\pgfpathlineto{\pgfqpoint{3.095884in}{1.408977in}}%
\pgfpathlineto{\pgfqpoint{3.098705in}{1.418866in}}%
\pgfpathlineto{\pgfqpoint{3.104950in}{1.423326in}}%
\pgfpathlineto{\pgfqpoint{3.118185in}{1.445751in}}%
\pgfpathlineto{\pgfqpoint{3.116263in}{1.452384in}}%
\pgfpathlineto{\pgfqpoint{3.122088in}{1.466802in}}%
\pgfpathlineto{\pgfqpoint{3.153183in}{1.449239in}}%
\pgfpathlineto{\pgfqpoint{3.156626in}{1.464838in}}%
\pgfpathlineto{\pgfqpoint{3.166450in}{1.465183in}}%
\pgfpathlineto{\pgfqpoint{3.174081in}{1.460247in}}%
\pgfpathlineto{\pgfqpoint{3.172555in}{1.451737in}}%
\pgfpathlineto{\pgfqpoint{3.189390in}{1.447867in}}%
\pgfpathlineto{\pgfqpoint{3.198897in}{1.442058in}}%
\pgfpathlineto{\pgfqpoint{3.205784in}{1.431832in}}%
\pgfpathlineto{\pgfqpoint{3.202838in}{1.419960in}}%
\pgfpathlineto{\pgfqpoint{3.196888in}{1.418988in}}%
\pgfpathlineto{\pgfqpoint{3.193438in}{1.401073in}}%
\pgfpathlineto{\pgfqpoint{3.200996in}{1.394069in}}%
\pgfpathlineto{\pgfqpoint{3.211636in}{1.399731in}}%
\pgfpathlineto{\pgfqpoint{3.221508in}{1.387775in}}%
\pgfpathlineto{\pgfqpoint{3.243566in}{1.385626in}}%
\pgfpathlineto{\pgfqpoint{3.247366in}{1.378749in}}%
\pgfpathlineto{\pgfqpoint{3.267778in}{1.372060in}}%
\pgfpathlineto{\pgfqpoint{3.265259in}{1.364168in}}%
\pgfpathlineto{\pgfqpoint{3.267915in}{1.341861in}}%
\pgfpathlineto{\pgfqpoint{3.276154in}{1.333438in}}%
\pgfpathlineto{\pgfqpoint{3.263997in}{1.329883in}}%
\pgfpathlineto{\pgfqpoint{3.265608in}{1.320360in}}%
\pgfpathlineto{\pgfqpoint{3.278578in}{1.312286in}}%
\pgfpathlineto{\pgfqpoint{3.279747in}{1.304321in}}%
\pgfpathlineto{\pgfqpoint{3.272425in}{1.298336in}}%
\pgfpathlineto{\pgfqpoint{3.257650in}{1.312515in}}%
\pgfpathlineto{\pgfqpoint{3.256190in}{1.302170in}}%
\pgfpathlineto{\pgfqpoint{3.268558in}{1.297251in}}%
\pgfpathlineto{\pgfqpoint{3.269796in}{1.292163in}}%
\pgfpathlineto{\pgfqpoint{3.286760in}{1.298887in}}%
\pgfpathlineto{\pgfqpoint{3.299759in}{1.300665in}}%
\pgfpathlineto{\pgfqpoint{3.313076in}{1.273798in}}%
\pgfpathlineto{\pgfqpoint{3.311590in}{1.273507in}}%
\pgfpathlineto{\pgfqpoint{3.305573in}{1.272263in}}%
\pgfpathlineto{\pgfqpoint{3.303799in}{1.271893in}}%
\pgfpathlineto{\pgfqpoint{3.302627in}{1.271664in}}%
\pgfpathlineto{\pgfqpoint{3.223341in}{1.255217in}}%
\pgfpathlineto{\pgfqpoint{3.152426in}{1.240702in}}%
\pgfpathlineto{\pgfqpoint{3.054377in}{1.223812in}}%
\pgfpathlineto{\pgfqpoint{2.971105in}{1.212813in}}%
\pgfpathlineto{\pgfqpoint{2.950506in}{1.211075in}}%
\pgfusepath{stroke}%
\end{pgfscope}%
\begin{pgfscope}%
\pgfpathrectangle{\pgfqpoint{0.100000in}{0.100000in}}{\pgfqpoint{3.608454in}{2.310000in}}%
\pgfusepath{clip}%
\pgfsetbuttcap%
\pgfsetroundjoin%
\pgfsetlinewidth{0.010037pt}%
\definecolor{currentstroke}{rgb}{1.000000,1.000000,1.000000}%
\pgfsetstrokecolor{currentstroke}%
\pgfsetdash{}{0pt}%
\pgfpathmoveto{\pgfqpoint{3.303727in}{1.386181in}}%
\pgfpathlineto{\pgfqpoint{3.319710in}{1.393431in}}%
\pgfpathlineto{\pgfqpoint{3.306946in}{1.356091in}}%
\pgfpathlineto{\pgfqpoint{3.298505in}{1.336318in}}%
\pgfpathlineto{\pgfqpoint{3.292424in}{1.347508in}}%
\pgfpathlineto{\pgfqpoint{3.303237in}{1.374298in}}%
\pgfpathlineto{\pgfqpoint{3.303727in}{1.386181in}}%
\pgfusepath{stroke}%
\end{pgfscope}%
\begin{pgfscope}%
\pgfpathrectangle{\pgfqpoint{0.100000in}{0.100000in}}{\pgfqpoint{3.608454in}{2.310000in}}%
\pgfusepath{clip}%
\pgfsetbuttcap%
\pgfsetroundjoin%
\pgfsetlinewidth{0.010037pt}%
\definecolor{currentstroke}{rgb}{1.000000,1.000000,1.000000}%
\pgfsetstrokecolor{currentstroke}%
\pgfsetdash{}{0pt}%
\pgfpathmoveto{\pgfqpoint{2.461763in}{1.153105in}}%
\pgfpathlineto{\pgfqpoint{2.457903in}{1.152546in}}%
\pgfpathlineto{\pgfqpoint{2.454263in}{1.152291in}}%
\pgfpathlineto{\pgfqpoint{2.455821in}{1.141115in}}%
\pgfpathlineto{\pgfqpoint{2.446430in}{1.129849in}}%
\pgfpathlineto{\pgfqpoint{2.444579in}{1.112222in}}%
\pgfpathlineto{\pgfqpoint{2.402597in}{1.109117in}}%
\pgfpathlineto{\pgfqpoint{2.406257in}{1.117394in}}%
\pgfpathlineto{\pgfqpoint{2.421394in}{1.132509in}}%
\pgfpathlineto{\pgfqpoint{2.421813in}{1.141271in}}%
\pgfpathlineto{\pgfqpoint{2.415111in}{1.149567in}}%
\pgfpathlineto{\pgfqpoint{2.352626in}{1.146263in}}%
\pgfpathlineto{\pgfqpoint{2.243380in}{1.142795in}}%
\pgfpathlineto{\pgfqpoint{2.179448in}{1.141630in}}%
\pgfpathlineto{\pgfqpoint{2.131124in}{1.141195in}}%
\pgfpathlineto{\pgfqpoint{2.130758in}{1.181350in}}%
\pgfpathlineto{\pgfqpoint{2.130373in}{1.297922in}}%
\pgfpathlineto{\pgfqpoint{2.130810in}{1.348616in}}%
\pgfpathlineto{\pgfqpoint{2.117751in}{1.356337in}}%
\pgfpathlineto{\pgfqpoint{2.111314in}{1.369785in}}%
\pgfpathlineto{\pgfqpoint{2.104306in}{1.376240in}}%
\pgfpathlineto{\pgfqpoint{2.103512in}{1.386028in}}%
\pgfpathlineto{\pgfqpoint{2.110832in}{1.401521in}}%
\pgfpathlineto{\pgfqpoint{2.110523in}{1.410477in}}%
\pgfpathlineto{\pgfqpoint{2.100667in}{1.408407in}}%
\pgfpathlineto{\pgfqpoint{2.088084in}{1.418676in}}%
\pgfpathlineto{\pgfqpoint{2.077998in}{1.436700in}}%
\pgfpathlineto{\pgfqpoint{2.069537in}{1.445017in}}%
\pgfpathlineto{\pgfqpoint{2.060697in}{1.465454in}}%
\pgfpathlineto{\pgfqpoint{2.152519in}{1.464021in}}%
\pgfpathlineto{\pgfqpoint{2.243792in}{1.467029in}}%
\pgfpathlineto{\pgfqpoint{2.302315in}{1.470405in}}%
\pgfpathlineto{\pgfqpoint{2.320622in}{1.452274in}}%
\pgfpathlineto{\pgfqpoint{2.317277in}{1.433365in}}%
\pgfpathlineto{\pgfqpoint{2.322158in}{1.412077in}}%
\pgfpathlineto{\pgfqpoint{2.327573in}{1.401342in}}%
\pgfpathlineto{\pgfqpoint{2.344850in}{1.386576in}}%
\pgfpathlineto{\pgfqpoint{2.348630in}{1.380053in}}%
\pgfpathlineto{\pgfqpoint{2.367545in}{1.365354in}}%
\pgfpathlineto{\pgfqpoint{2.370287in}{1.348866in}}%
\pgfpathlineto{\pgfqpoint{2.375719in}{1.336702in}}%
\pgfpathlineto{\pgfqpoint{2.388555in}{1.343787in}}%
\pgfpathlineto{\pgfqpoint{2.399134in}{1.340547in}}%
\pgfpathlineto{\pgfqpoint{2.407122in}{1.331947in}}%
\pgfpathlineto{\pgfqpoint{2.404106in}{1.316728in}}%
\pgfpathlineto{\pgfqpoint{2.394301in}{1.296831in}}%
\pgfpathlineto{\pgfqpoint{2.394453in}{1.286467in}}%
\pgfpathlineto{\pgfqpoint{2.401780in}{1.277970in}}%
\pgfpathlineto{\pgfqpoint{2.417550in}{1.266715in}}%
\pgfpathlineto{\pgfqpoint{2.435708in}{1.257580in}}%
\pgfpathlineto{\pgfqpoint{2.440014in}{1.250265in}}%
\pgfpathlineto{\pgfqpoint{2.449301in}{1.246806in}}%
\pgfpathlineto{\pgfqpoint{2.449670in}{1.237099in}}%
\pgfpathlineto{\pgfqpoint{2.456543in}{1.224768in}}%
\pgfpathlineto{\pgfqpoint{2.451774in}{1.213813in}}%
\pgfpathlineto{\pgfqpoint{2.461108in}{1.196383in}}%
\pgfpathlineto{\pgfqpoint{2.468410in}{1.199844in}}%
\pgfpathlineto{\pgfqpoint{2.477364in}{1.192570in}}%
\pgfpathlineto{\pgfqpoint{2.474354in}{1.160178in}}%
\pgfpathlineto{\pgfqpoint{2.463799in}{1.162233in}}%
\pgfpathlineto{\pgfqpoint{2.461763in}{1.153105in}}%
\pgfusepath{stroke}%
\end{pgfscope}%
\begin{pgfscope}%
\pgfpathrectangle{\pgfqpoint{0.100000in}{0.100000in}}{\pgfqpoint{3.608454in}{2.310000in}}%
\pgfusepath{clip}%
\pgfsetbuttcap%
\pgfsetroundjoin%
\pgfsetlinewidth{0.010037pt}%
\definecolor{currentstroke}{rgb}{1.000000,1.000000,1.000000}%
\pgfsetstrokecolor{currentstroke}%
\pgfsetdash{}{0pt}%
\pgfpathmoveto{\pgfqpoint{0.843082in}{1.160088in}}%
\pgfpathlineto{\pgfqpoint{0.850725in}{1.176360in}}%
\pgfpathlineto{\pgfqpoint{0.848633in}{1.200785in}}%
\pgfpathlineto{\pgfqpoint{0.852293in}{1.207731in}}%
\pgfpathlineto{\pgfqpoint{0.852655in}{1.238156in}}%
\pgfpathlineto{\pgfqpoint{0.856124in}{1.246917in}}%
\pgfpathlineto{\pgfqpoint{0.868345in}{1.248280in}}%
\pgfpathlineto{\pgfqpoint{0.879687in}{1.244390in}}%
\pgfpathlineto{\pgfqpoint{0.884626in}{1.233679in}}%
\pgfpathlineto{\pgfqpoint{0.891552in}{1.234089in}}%
\pgfpathlineto{\pgfqpoint{0.900156in}{1.246360in}}%
\pgfpathlineto{\pgfqpoint{0.912593in}{1.306845in}}%
\pgfpathlineto{\pgfqpoint{0.983361in}{1.292233in}}%
\pgfpathlineto{\pgfqpoint{1.024434in}{1.284134in}}%
\pgfpathlineto{\pgfqpoint{1.133501in}{1.264961in}}%
\pgfpathlineto{\pgfqpoint{1.163676in}{1.258901in}}%
\pgfpathlineto{\pgfqpoint{1.222893in}{1.249723in}}%
\pgfpathlineto{\pgfqpoint{1.210751in}{1.171542in}}%
\pgfpathlineto{\pgfqpoint{1.198122in}{1.089991in}}%
\pgfpathlineto{\pgfqpoint{1.183580in}{0.998244in}}%
\pgfpathlineto{\pgfqpoint{1.167224in}{0.892929in}}%
\pgfpathlineto{\pgfqpoint{1.154011in}{0.806432in}}%
\pgfpathlineto{\pgfqpoint{1.059373in}{0.821460in}}%
\pgfpathlineto{\pgfqpoint{1.017791in}{0.828581in}}%
\pgfpathlineto{\pgfqpoint{0.999216in}{0.839702in}}%
\pgfpathlineto{\pgfqpoint{0.878105in}{0.912775in}}%
\pgfpathlineto{\pgfqpoint{0.787216in}{0.968230in}}%
\pgfpathlineto{\pgfqpoint{0.790243in}{0.978058in}}%
\pgfpathlineto{\pgfqpoint{0.797751in}{0.984935in}}%
\pgfpathlineto{\pgfqpoint{0.805560in}{0.983661in}}%
\pgfpathlineto{\pgfqpoint{0.816886in}{0.990883in}}%
\pgfpathlineto{\pgfqpoint{0.818669in}{1.001263in}}%
\pgfpathlineto{\pgfqpoint{0.804856in}{1.013854in}}%
\pgfpathlineto{\pgfqpoint{0.814680in}{1.038025in}}%
\pgfpathlineto{\pgfqpoint{0.824557in}{1.047321in}}%
\pgfpathlineto{\pgfqpoint{0.829245in}{1.058349in}}%
\pgfpathlineto{\pgfqpoint{0.832133in}{1.078570in}}%
\pgfpathlineto{\pgfqpoint{0.841395in}{1.087727in}}%
\pgfpathlineto{\pgfqpoint{0.860834in}{1.096874in}}%
\pgfpathlineto{\pgfqpoint{0.860405in}{1.104932in}}%
\pgfpathlineto{\pgfqpoint{0.849548in}{1.114934in}}%
\pgfpathlineto{\pgfqpoint{0.848070in}{1.135552in}}%
\pgfpathlineto{\pgfqpoint{0.840594in}{1.150624in}}%
\pgfpathlineto{\pgfqpoint{0.843082in}{1.160088in}}%
\pgfusepath{stroke}%
\end{pgfscope}%
\begin{pgfscope}%
\pgfpathrectangle{\pgfqpoint{0.100000in}{0.100000in}}{\pgfqpoint{3.608454in}{2.310000in}}%
\pgfusepath{clip}%
\pgfsetbuttcap%
\pgfsetroundjoin%
\pgfsetlinewidth{0.010037pt}%
\definecolor{currentstroke}{rgb}{1.000000,1.000000,1.000000}%
\pgfsetstrokecolor{currentstroke}%
\pgfsetdash{}{0pt}%
\pgfpathmoveto{\pgfqpoint{2.141369in}{0.914830in}}%
\pgfpathlineto{\pgfqpoint{2.123825in}{0.920256in}}%
\pgfpathlineto{\pgfqpoint{2.101346in}{0.937472in}}%
\pgfpathlineto{\pgfqpoint{2.091360in}{0.941172in}}%
\pgfpathlineto{\pgfqpoint{2.085016in}{0.933736in}}%
\pgfpathlineto{\pgfqpoint{2.077003in}{0.933361in}}%
\pgfpathlineto{\pgfqpoint{2.066866in}{0.939673in}}%
\pgfpathlineto{\pgfqpoint{2.050959in}{0.931585in}}%
\pgfpathlineto{\pgfqpoint{2.045296in}{0.935568in}}%
\pgfpathlineto{\pgfqpoint{2.029643in}{0.926132in}}%
\pgfpathlineto{\pgfqpoint{2.013130in}{0.927873in}}%
\pgfpathlineto{\pgfqpoint{2.000393in}{0.934005in}}%
\pgfpathlineto{\pgfqpoint{1.991473in}{0.931697in}}%
\pgfpathlineto{\pgfqpoint{1.977195in}{0.940369in}}%
\pgfpathlineto{\pgfqpoint{1.969253in}{0.925074in}}%
\pgfpathlineto{\pgfqpoint{1.961130in}{0.938228in}}%
\pgfpathlineto{\pgfqpoint{1.945087in}{0.933125in}}%
\pgfpathlineto{\pgfqpoint{1.931055in}{0.945619in}}%
\pgfpathlineto{\pgfqpoint{1.918793in}{0.935548in}}%
\pgfpathlineto{\pgfqpoint{1.912652in}{0.944802in}}%
\pgfpathlineto{\pgfqpoint{1.903801in}{0.947808in}}%
\pgfpathlineto{\pgfqpoint{1.898410in}{0.956730in}}%
\pgfpathlineto{\pgfqpoint{1.886834in}{0.959269in}}%
\pgfpathlineto{\pgfqpoint{1.880154in}{0.952553in}}%
\pgfpathlineto{\pgfqpoint{1.868804in}{0.961257in}}%
\pgfpathlineto{\pgfqpoint{1.863513in}{0.959270in}}%
\pgfpathlineto{\pgfqpoint{1.844702in}{0.966322in}}%
\pgfpathlineto{\pgfqpoint{1.832921in}{0.967109in}}%
\pgfpathlineto{\pgfqpoint{1.827647in}{0.982085in}}%
\pgfpathlineto{\pgfqpoint{1.818203in}{0.980220in}}%
\pgfpathlineto{\pgfqpoint{1.807385in}{0.983916in}}%
\pgfpathlineto{\pgfqpoint{1.800261in}{0.981780in}}%
\pgfpathlineto{\pgfqpoint{1.784985in}{0.998601in}}%
\pgfpathlineto{\pgfqpoint{1.780728in}{0.997501in}}%
\pgfpathlineto{\pgfqpoint{1.784594in}{1.065714in}}%
\pgfpathlineto{\pgfqpoint{1.788806in}{1.150144in}}%
\pgfpathlineto{\pgfqpoint{1.719672in}{1.154011in}}%
\pgfpathlineto{\pgfqpoint{1.651450in}{1.159169in}}%
\pgfpathlineto{\pgfqpoint{1.598743in}{1.163746in}}%
\pgfpathlineto{\pgfqpoint{1.602243in}{1.203632in}}%
\pgfpathlineto{\pgfqpoint{1.662648in}{1.197867in}}%
\pgfpathlineto{\pgfqpoint{1.731159in}{1.193571in}}%
\pgfpathlineto{\pgfqpoint{1.826575in}{1.187945in}}%
\pgfpathlineto{\pgfqpoint{1.883211in}{1.185760in}}%
\pgfpathlineto{\pgfqpoint{1.951062in}{1.183366in}}%
\pgfpathlineto{\pgfqpoint{2.043230in}{1.181753in}}%
\pgfpathlineto{\pgfqpoint{2.130758in}{1.181350in}}%
\pgfpathlineto{\pgfqpoint{2.131124in}{1.141195in}}%
\pgfpathlineto{\pgfqpoint{2.136037in}{1.110943in}}%
\pgfpathlineto{\pgfqpoint{2.143669in}{1.055067in}}%
\pgfpathlineto{\pgfqpoint{2.142095in}{0.959646in}}%
\pgfpathlineto{\pgfqpoint{2.141369in}{0.914830in}}%
\pgfusepath{stroke}%
\end{pgfscope}%
\begin{pgfscope}%
\pgfpathrectangle{\pgfqpoint{0.100000in}{0.100000in}}{\pgfqpoint{3.608454in}{2.310000in}}%
\pgfusepath{clip}%
\pgfsetbuttcap%
\pgfsetroundjoin%
\pgfsetlinewidth{0.010037pt}%
\definecolor{currentstroke}{rgb}{1.000000,1.000000,1.000000}%
\pgfsetstrokecolor{currentstroke}%
\pgfsetdash{}{0pt}%
\pgfpathmoveto{\pgfqpoint{2.798250in}{1.063020in}}%
\pgfpathlineto{\pgfqpoint{2.798302in}{1.080700in}}%
\pgfpathlineto{\pgfqpoint{2.813666in}{1.087489in}}%
\pgfpathlineto{\pgfqpoint{2.814310in}{1.098342in}}%
\pgfpathlineto{\pgfqpoint{2.828067in}{1.111600in}}%
\pgfpathlineto{\pgfqpoint{2.845262in}{1.114316in}}%
\pgfpathlineto{\pgfqpoint{2.860379in}{1.128927in}}%
\pgfpathlineto{\pgfqpoint{2.876166in}{1.135457in}}%
\pgfpathlineto{\pgfqpoint{2.888015in}{1.155063in}}%
\pgfpathlineto{\pgfqpoint{2.898807in}{1.153640in}}%
\pgfpathlineto{\pgfqpoint{2.921614in}{1.171504in}}%
\pgfpathlineto{\pgfqpoint{2.933658in}{1.171840in}}%
\pgfpathlineto{\pgfqpoint{2.938743in}{1.185452in}}%
\pgfpathlineto{\pgfqpoint{2.948316in}{1.194928in}}%
\pgfpathlineto{\pgfqpoint{2.950506in}{1.211075in}}%
\pgfpathlineto{\pgfqpoint{2.971105in}{1.212813in}}%
\pgfpathlineto{\pgfqpoint{3.054377in}{1.223812in}}%
\pgfpathlineto{\pgfqpoint{3.152426in}{1.240702in}}%
\pgfpathlineto{\pgfqpoint{3.223341in}{1.255217in}}%
\pgfpathlineto{\pgfqpoint{3.302627in}{1.271664in}}%
\pgfpathlineto{\pgfqpoint{3.325283in}{1.240558in}}%
\pgfpathlineto{\pgfqpoint{3.313011in}{1.242544in}}%
\pgfpathlineto{\pgfqpoint{3.297258in}{1.238706in}}%
\pgfpathlineto{\pgfqpoint{3.281735in}{1.222854in}}%
\pgfpathlineto{\pgfqpoint{3.270579in}{1.223994in}}%
\pgfpathlineto{\pgfqpoint{3.269168in}{1.214579in}}%
\pgfpathlineto{\pgfqpoint{3.289332in}{1.222025in}}%
\pgfpathlineto{\pgfqpoint{3.309624in}{1.225095in}}%
\pgfpathlineto{\pgfqpoint{3.317158in}{1.221010in}}%
\pgfpathlineto{\pgfqpoint{3.327297in}{1.225667in}}%
\pgfpathlineto{\pgfqpoint{3.332529in}{1.222404in}}%
\pgfpathlineto{\pgfqpoint{3.337160in}{1.206867in}}%
\pgfpathlineto{\pgfqpoint{3.327523in}{1.202055in}}%
\pgfpathlineto{\pgfqpoint{3.320908in}{1.183109in}}%
\pgfpathlineto{\pgfqpoint{3.313973in}{1.175732in}}%
\pgfpathlineto{\pgfqpoint{3.292714in}{1.177346in}}%
\pgfpathlineto{\pgfqpoint{3.293856in}{1.188471in}}%
\pgfpathlineto{\pgfqpoint{3.282246in}{1.183611in}}%
\pgfpathlineto{\pgfqpoint{3.279721in}{1.174333in}}%
\pgfpathlineto{\pgfqpoint{3.287679in}{1.164722in}}%
\pgfpathlineto{\pgfqpoint{3.290383in}{1.148409in}}%
\pgfpathlineto{\pgfqpoint{3.277360in}{1.139115in}}%
\pgfpathlineto{\pgfqpoint{3.291444in}{1.135846in}}%
\pgfpathlineto{\pgfqpoint{3.297817in}{1.142680in}}%
\pgfpathlineto{\pgfqpoint{3.311909in}{1.143512in}}%
\pgfpathlineto{\pgfqpoint{3.304667in}{1.128754in}}%
\pgfpathlineto{\pgfqpoint{3.295798in}{1.121643in}}%
\pgfpathlineto{\pgfqpoint{3.268968in}{1.114523in}}%
\pgfpathlineto{\pgfqpoint{3.241422in}{1.089539in}}%
\pgfpathlineto{\pgfqpoint{3.224436in}{1.064923in}}%
\pgfpathlineto{\pgfqpoint{3.224336in}{1.054925in}}%
\pgfpathlineto{\pgfqpoint{3.217551in}{1.041130in}}%
\pgfpathlineto{\pgfqpoint{3.182729in}{1.032053in}}%
\pgfpathlineto{\pgfqpoint{3.099917in}{1.091416in}}%
\pgfpathlineto{\pgfqpoint{3.026989in}{1.080701in}}%
\pgfpathlineto{\pgfqpoint{3.026378in}{1.090595in}}%
\pgfpathlineto{\pgfqpoint{3.015293in}{1.101737in}}%
\pgfpathlineto{\pgfqpoint{3.006902in}{1.104464in}}%
\pgfpathlineto{\pgfqpoint{2.927665in}{1.096231in}}%
\pgfpathlineto{\pgfqpoint{2.909453in}{1.090054in}}%
\pgfpathlineto{\pgfqpoint{2.876568in}{1.073696in}}%
\pgfpathlineto{\pgfqpoint{2.848149in}{1.069168in}}%
\pgfpathlineto{\pgfqpoint{2.798250in}{1.063020in}}%
\pgfusepath{stroke}%
\end{pgfscope}%
\begin{pgfscope}%
\pgfpathrectangle{\pgfqpoint{0.100000in}{0.100000in}}{\pgfqpoint{3.608454in}{2.310000in}}%
\pgfusepath{clip}%
\pgfsetbuttcap%
\pgfsetroundjoin%
\pgfsetlinewidth{0.010037pt}%
\definecolor{currentstroke}{rgb}{1.000000,1.000000,1.000000}%
\pgfsetstrokecolor{currentstroke}%
\pgfsetdash{}{0pt}%
\pgfpathmoveto{\pgfqpoint{2.798250in}{1.063020in}}%
\pgfpathlineto{\pgfqpoint{2.715320in}{1.053912in}}%
\pgfpathlineto{\pgfqpoint{2.639455in}{1.047089in}}%
\pgfpathlineto{\pgfqpoint{2.561482in}{1.042047in}}%
\pgfpathlineto{\pgfqpoint{2.548098in}{1.040107in}}%
\pgfpathlineto{\pgfqpoint{2.495536in}{1.036064in}}%
\pgfpathlineto{\pgfqpoint{2.411349in}{1.031151in}}%
\pgfpathlineto{\pgfqpoint{2.426327in}{1.043584in}}%
\pgfpathlineto{\pgfqpoint{2.422910in}{1.056671in}}%
\pgfpathlineto{\pgfqpoint{2.426220in}{1.072535in}}%
\pgfpathlineto{\pgfqpoint{2.431278in}{1.080004in}}%
\pgfpathlineto{\pgfqpoint{2.431088in}{1.090395in}}%
\pgfpathlineto{\pgfqpoint{2.444556in}{1.096930in}}%
\pgfpathlineto{\pgfqpoint{2.444579in}{1.112222in}}%
\pgfpathlineto{\pgfqpoint{2.446430in}{1.129849in}}%
\pgfpathlineto{\pgfqpoint{2.455821in}{1.141115in}}%
\pgfpathlineto{\pgfqpoint{2.454263in}{1.152291in}}%
\pgfpathlineto{\pgfqpoint{2.457903in}{1.152546in}}%
\pgfpathlineto{\pgfqpoint{2.461763in}{1.153105in}}%
\pgfpathlineto{\pgfqpoint{2.548849in}{1.158640in}}%
\pgfpathlineto{\pgfqpoint{2.545934in}{1.172945in}}%
\pgfpathlineto{\pgfqpoint{2.559948in}{1.171005in}}%
\pgfpathlineto{\pgfqpoint{2.644773in}{1.179396in}}%
\pgfpathlineto{\pgfqpoint{2.712861in}{1.183250in}}%
\pgfpathlineto{\pgfqpoint{2.754353in}{1.186739in}}%
\pgfpathlineto{\pgfqpoint{2.823062in}{1.193254in}}%
\pgfpathlineto{\pgfqpoint{2.824207in}{1.194550in}}%
\pgfpathlineto{\pgfqpoint{2.853565in}{1.197457in}}%
\pgfpathlineto{\pgfqpoint{2.950506in}{1.211075in}}%
\pgfpathlineto{\pgfqpoint{2.948316in}{1.194928in}}%
\pgfpathlineto{\pgfqpoint{2.938743in}{1.185452in}}%
\pgfpathlineto{\pgfqpoint{2.933658in}{1.171840in}}%
\pgfpathlineto{\pgfqpoint{2.921614in}{1.171504in}}%
\pgfpathlineto{\pgfqpoint{2.898807in}{1.153640in}}%
\pgfpathlineto{\pgfqpoint{2.888015in}{1.155063in}}%
\pgfpathlineto{\pgfqpoint{2.876166in}{1.135457in}}%
\pgfpathlineto{\pgfqpoint{2.860379in}{1.128927in}}%
\pgfpathlineto{\pgfqpoint{2.845262in}{1.114316in}}%
\pgfpathlineto{\pgfqpoint{2.828067in}{1.111600in}}%
\pgfpathlineto{\pgfqpoint{2.814310in}{1.098342in}}%
\pgfpathlineto{\pgfqpoint{2.813666in}{1.087489in}}%
\pgfpathlineto{\pgfqpoint{2.798302in}{1.080700in}}%
\pgfpathlineto{\pgfqpoint{2.798250in}{1.063020in}}%
\pgfusepath{stroke}%
\end{pgfscope}%
\begin{pgfscope}%
\pgfpathrectangle{\pgfqpoint{0.100000in}{0.100000in}}{\pgfqpoint{3.608454in}{2.310000in}}%
\pgfusepath{clip}%
\pgfsetbuttcap%
\pgfsetroundjoin%
\pgfsetlinewidth{0.010037pt}%
\definecolor{currentstroke}{rgb}{1.000000,1.000000,1.000000}%
\pgfsetstrokecolor{currentstroke}%
\pgfsetdash{}{0pt}%
\pgfpathmoveto{\pgfqpoint{1.327465in}{0.817498in}}%
\pgfpathlineto{\pgfqpoint{1.321550in}{0.827003in}}%
\pgfpathlineto{\pgfqpoint{1.323997in}{0.835189in}}%
\pgfpathlineto{\pgfqpoint{1.442877in}{0.821252in}}%
\pgfpathlineto{\pgfqpoint{1.563330in}{0.809245in}}%
\pgfpathlineto{\pgfqpoint{1.566842in}{0.849964in}}%
\pgfpathlineto{\pgfqpoint{1.584946in}{1.026087in}}%
\pgfpathlineto{\pgfqpoint{1.596844in}{1.163852in}}%
\pgfpathlineto{\pgfqpoint{1.598743in}{1.163746in}}%
\pgfpathlineto{\pgfqpoint{1.719672in}{1.154011in}}%
\pgfpathlineto{\pgfqpoint{1.788806in}{1.150144in}}%
\pgfpathlineto{\pgfqpoint{1.780728in}{0.997501in}}%
\pgfpathlineto{\pgfqpoint{1.784985in}{0.998601in}}%
\pgfpathlineto{\pgfqpoint{1.800261in}{0.981780in}}%
\pgfpathlineto{\pgfqpoint{1.807385in}{0.983916in}}%
\pgfpathlineto{\pgfqpoint{1.818203in}{0.980220in}}%
\pgfpathlineto{\pgfqpoint{1.827647in}{0.982085in}}%
\pgfpathlineto{\pgfqpoint{1.832921in}{0.967109in}}%
\pgfpathlineto{\pgfqpoint{1.844702in}{0.966322in}}%
\pgfpathlineto{\pgfqpoint{1.863513in}{0.959270in}}%
\pgfpathlineto{\pgfqpoint{1.868804in}{0.961257in}}%
\pgfpathlineto{\pgfqpoint{1.880154in}{0.952553in}}%
\pgfpathlineto{\pgfqpoint{1.886834in}{0.959269in}}%
\pgfpathlineto{\pgfqpoint{1.898410in}{0.956730in}}%
\pgfpathlineto{\pgfqpoint{1.903801in}{0.947808in}}%
\pgfpathlineto{\pgfqpoint{1.912652in}{0.944802in}}%
\pgfpathlineto{\pgfqpoint{1.918793in}{0.935548in}}%
\pgfpathlineto{\pgfqpoint{1.931055in}{0.945619in}}%
\pgfpathlineto{\pgfqpoint{1.945087in}{0.933125in}}%
\pgfpathlineto{\pgfqpoint{1.961130in}{0.938228in}}%
\pgfpathlineto{\pgfqpoint{1.969253in}{0.925074in}}%
\pgfpathlineto{\pgfqpoint{1.977195in}{0.940369in}}%
\pgfpathlineto{\pgfqpoint{1.991473in}{0.931697in}}%
\pgfpathlineto{\pgfqpoint{2.000393in}{0.934005in}}%
\pgfpathlineto{\pgfqpoint{2.013130in}{0.927873in}}%
\pgfpathlineto{\pgfqpoint{2.029643in}{0.926132in}}%
\pgfpathlineto{\pgfqpoint{2.045296in}{0.935568in}}%
\pgfpathlineto{\pgfqpoint{2.050959in}{0.931585in}}%
\pgfpathlineto{\pgfqpoint{2.066866in}{0.939673in}}%
\pgfpathlineto{\pgfqpoint{2.077003in}{0.933361in}}%
\pgfpathlineto{\pgfqpoint{2.085016in}{0.933736in}}%
\pgfpathlineto{\pgfqpoint{2.091360in}{0.941172in}}%
\pgfpathlineto{\pgfqpoint{2.101346in}{0.937472in}}%
\pgfpathlineto{\pgfqpoint{2.123825in}{0.920256in}}%
\pgfpathlineto{\pgfqpoint{2.141369in}{0.914830in}}%
\pgfpathlineto{\pgfqpoint{2.148404in}{0.908183in}}%
\pgfpathlineto{\pgfqpoint{2.157210in}{0.911803in}}%
\pgfpathlineto{\pgfqpoint{2.170554in}{0.909031in}}%
\pgfpathlineto{\pgfqpoint{2.171929in}{0.784895in}}%
\pgfpathlineto{\pgfqpoint{2.181197in}{0.777032in}}%
\pgfpathlineto{\pgfqpoint{2.188625in}{0.762543in}}%
\pgfpathlineto{\pgfqpoint{2.186014in}{0.752849in}}%
\pgfpathlineto{\pgfqpoint{2.191847in}{0.748745in}}%
\pgfpathlineto{\pgfqpoint{2.196588in}{0.730869in}}%
\pgfpathlineto{\pgfqpoint{2.205667in}{0.721296in}}%
\pgfpathlineto{\pgfqpoint{2.207701in}{0.700970in}}%
\pgfpathlineto{\pgfqpoint{2.206122in}{0.692365in}}%
\pgfpathlineto{\pgfqpoint{2.193781in}{0.669590in}}%
\pgfpathlineto{\pgfqpoint{2.192362in}{0.653216in}}%
\pgfpathlineto{\pgfqpoint{2.196548in}{0.649798in}}%
\pgfpathlineto{\pgfqpoint{2.196912in}{0.635452in}}%
\pgfpathlineto{\pgfqpoint{2.192657in}{0.626436in}}%
\pgfpathlineto{\pgfqpoint{2.186039in}{0.623040in}}%
\pgfpathlineto{\pgfqpoint{2.179559in}{0.611258in}}%
\pgfpathlineto{\pgfqpoint{2.187818in}{0.599847in}}%
\pgfpathlineto{\pgfqpoint{2.171786in}{0.599622in}}%
\pgfpathlineto{\pgfqpoint{2.128920in}{0.580009in}}%
\pgfpathlineto{\pgfqpoint{2.137104in}{0.591747in}}%
\pgfpathlineto{\pgfqpoint{2.127193in}{0.598075in}}%
\pgfpathlineto{\pgfqpoint{2.125123in}{0.608834in}}%
\pgfpathlineto{\pgfqpoint{2.105780in}{0.590093in}}%
\pgfpathlineto{\pgfqpoint{2.114364in}{0.577284in}}%
\pgfpathlineto{\pgfqpoint{2.102118in}{0.560974in}}%
\pgfpathlineto{\pgfqpoint{2.095528in}{0.561316in}}%
\pgfpathlineto{\pgfqpoint{2.089328in}{0.543553in}}%
\pgfpathlineto{\pgfqpoint{2.069722in}{0.529582in}}%
\pgfpathlineto{\pgfqpoint{2.051405in}{0.524547in}}%
\pgfpathlineto{\pgfqpoint{2.019457in}{0.511448in}}%
\pgfpathlineto{\pgfqpoint{2.016147in}{0.518755in}}%
\pgfpathlineto{\pgfqpoint{1.996246in}{0.512029in}}%
\pgfpathlineto{\pgfqpoint{2.008486in}{0.500577in}}%
\pgfpathlineto{\pgfqpoint{1.989179in}{0.490629in}}%
\pgfpathlineto{\pgfqpoint{1.968354in}{0.475608in}}%
\pgfpathlineto{\pgfqpoint{1.963353in}{0.482590in}}%
\pgfpathlineto{\pgfqpoint{1.946310in}{0.472125in}}%
\pgfpathlineto{\pgfqpoint{1.963014in}{0.468151in}}%
\pgfpathlineto{\pgfqpoint{1.950420in}{0.451696in}}%
\pgfpathlineto{\pgfqpoint{1.944272in}{0.456593in}}%
\pgfpathlineto{\pgfqpoint{1.936015in}{0.448727in}}%
\pgfpathlineto{\pgfqpoint{1.946306in}{0.441857in}}%
\pgfpathlineto{\pgfqpoint{1.940221in}{0.431815in}}%
\pgfpathlineto{\pgfqpoint{1.932743in}{0.408043in}}%
\pgfpathlineto{\pgfqpoint{1.921956in}{0.385147in}}%
\pgfpathlineto{\pgfqpoint{1.926799in}{0.370139in}}%
\pgfpathlineto{\pgfqpoint{1.935030in}{0.334756in}}%
\pgfpathlineto{\pgfqpoint{1.935771in}{0.320436in}}%
\pgfpathlineto{\pgfqpoint{1.948489in}{0.301654in}}%
\pgfpathlineto{\pgfqpoint{1.938682in}{0.302751in}}%
\pgfpathlineto{\pgfqpoint{1.929156in}{0.293276in}}%
\pgfpathlineto{\pgfqpoint{1.918000in}{0.300723in}}%
\pgfpathlineto{\pgfqpoint{1.913996in}{0.308152in}}%
\pgfpathlineto{\pgfqpoint{1.898118in}{0.311608in}}%
\pgfpathlineto{\pgfqpoint{1.873865in}{0.312032in}}%
\pgfpathlineto{\pgfqpoint{1.855994in}{0.326101in}}%
\pgfpathlineto{\pgfqpoint{1.839768in}{0.328451in}}%
\pgfpathlineto{\pgfqpoint{1.829925in}{0.339632in}}%
\pgfpathlineto{\pgfqpoint{1.809309in}{0.344132in}}%
\pgfpathlineto{\pgfqpoint{1.798036in}{0.380084in}}%
\pgfpathlineto{\pgfqpoint{1.786510in}{0.394486in}}%
\pgfpathlineto{\pgfqpoint{1.788468in}{0.408170in}}%
\pgfpathlineto{\pgfqpoint{1.781321in}{0.418171in}}%
\pgfpathlineto{\pgfqpoint{1.785806in}{0.431841in}}%
\pgfpathlineto{\pgfqpoint{1.782106in}{0.441858in}}%
\pgfpathlineto{\pgfqpoint{1.770524in}{0.446396in}}%
\pgfpathlineto{\pgfqpoint{1.759694in}{0.457947in}}%
\pgfpathlineto{\pgfqpoint{1.755748in}{0.473412in}}%
\pgfpathlineto{\pgfqpoint{1.745500in}{0.487466in}}%
\pgfpathlineto{\pgfqpoint{1.731859in}{0.498402in}}%
\pgfpathlineto{\pgfqpoint{1.719573in}{0.529770in}}%
\pgfpathlineto{\pgfqpoint{1.709652in}{0.564077in}}%
\pgfpathlineto{\pgfqpoint{1.701513in}{0.577641in}}%
\pgfpathlineto{\pgfqpoint{1.687401in}{0.589061in}}%
\pgfpathlineto{\pgfqpoint{1.683874in}{0.597350in}}%
\pgfpathlineto{\pgfqpoint{1.670661in}{0.602497in}}%
\pgfpathlineto{\pgfqpoint{1.657620in}{0.624460in}}%
\pgfpathlineto{\pgfqpoint{1.645710in}{0.622774in}}%
\pgfpathlineto{\pgfqpoint{1.633616in}{0.628259in}}%
\pgfpathlineto{\pgfqpoint{1.616474in}{0.627199in}}%
\pgfpathlineto{\pgfqpoint{1.599059in}{0.636251in}}%
\pgfpathlineto{\pgfqpoint{1.594148in}{0.627642in}}%
\pgfpathlineto{\pgfqpoint{1.573796in}{0.627427in}}%
\pgfpathlineto{\pgfqpoint{1.563435in}{0.611106in}}%
\pgfpathlineto{\pgfqpoint{1.554423in}{0.590899in}}%
\pgfpathlineto{\pgfqpoint{1.535259in}{0.569207in}}%
\pgfpathlineto{\pgfqpoint{1.513549in}{0.578727in}}%
\pgfpathlineto{\pgfqpoint{1.510468in}{0.585018in}}%
\pgfpathlineto{\pgfqpoint{1.497274in}{0.589816in}}%
\pgfpathlineto{\pgfqpoint{1.493221in}{0.596380in}}%
\pgfpathlineto{\pgfqpoint{1.475705in}{0.603005in}}%
\pgfpathlineto{\pgfqpoint{1.465907in}{0.616550in}}%
\pgfpathlineto{\pgfqpoint{1.454464in}{0.623084in}}%
\pgfpathlineto{\pgfqpoint{1.444618in}{0.634488in}}%
\pgfpathlineto{\pgfqpoint{1.436949in}{0.653800in}}%
\pgfpathlineto{\pgfqpoint{1.437807in}{0.680190in}}%
\pgfpathlineto{\pgfqpoint{1.428809in}{0.693526in}}%
\pgfpathlineto{\pgfqpoint{1.427774in}{0.707969in}}%
\pgfpathlineto{\pgfqpoint{1.407810in}{0.729594in}}%
\pgfpathlineto{\pgfqpoint{1.396194in}{0.734211in}}%
\pgfpathlineto{\pgfqpoint{1.383867in}{0.754389in}}%
\pgfpathlineto{\pgfqpoint{1.373364in}{0.762406in}}%
\pgfpathlineto{\pgfqpoint{1.360017in}{0.781928in}}%
\pgfpathlineto{\pgfqpoint{1.346361in}{0.790377in}}%
\pgfpathlineto{\pgfqpoint{1.337418in}{0.812026in}}%
\pgfpathlineto{\pgfqpoint{1.327465in}{0.817498in}}%
\pgfpathlineto{\pgfqpoint{1.327465in}{0.817498in}}%
\pgfusepath{stroke}%
\end{pgfscope}%
\begin{pgfscope}%
\pgfpathrectangle{\pgfqpoint{0.100000in}{0.100000in}}{\pgfqpoint{3.608454in}{2.310000in}}%
\pgfusepath{clip}%
\pgfsetbuttcap%
\pgfsetroundjoin%
\pgfsetlinewidth{0.010037pt}%
\definecolor{currentstroke}{rgb}{1.000000,1.000000,1.000000}%
\pgfsetstrokecolor{currentstroke}%
\pgfsetdash{}{0pt}%
\pgfpathmoveto{\pgfqpoint{1.222893in}{1.249723in}}%
\pgfpathlineto{\pgfqpoint{1.321411in}{1.235415in}}%
\pgfpathlineto{\pgfqpoint{1.359578in}{1.229535in}}%
\pgfpathlineto{\pgfqpoint{1.467164in}{1.216739in}}%
\pgfpathlineto{\pgfqpoint{1.539577in}{1.209185in}}%
\pgfpathlineto{\pgfqpoint{1.602243in}{1.203632in}}%
\pgfpathlineto{\pgfqpoint{1.598743in}{1.163746in}}%
\pgfpathlineto{\pgfqpoint{1.596844in}{1.163852in}}%
\pgfpathlineto{\pgfqpoint{1.591104in}{1.095394in}}%
\pgfpathlineto{\pgfqpoint{1.584946in}{1.026087in}}%
\pgfpathlineto{\pgfqpoint{1.577823in}{0.953532in}}%
\pgfpathlineto{\pgfqpoint{1.566842in}{0.849964in}}%
\pgfpathlineto{\pgfqpoint{1.563330in}{0.809245in}}%
\pgfpathlineto{\pgfqpoint{1.498773in}{0.815759in}}%
\pgfpathlineto{\pgfqpoint{1.442877in}{0.821252in}}%
\pgfpathlineto{\pgfqpoint{1.365572in}{0.830032in}}%
\pgfpathlineto{\pgfqpoint{1.323997in}{0.835189in}}%
\pgfpathlineto{\pgfqpoint{1.321550in}{0.827003in}}%
\pgfpathlineto{\pgfqpoint{1.327465in}{0.817498in}}%
\pgfpathlineto{\pgfqpoint{1.277499in}{0.823987in}}%
\pgfpathlineto{\pgfqpoint{1.215853in}{0.832835in}}%
\pgfpathlineto{\pgfqpoint{1.210243in}{0.797968in}}%
\pgfpathlineto{\pgfqpoint{1.154011in}{0.806432in}}%
\pgfpathlineto{\pgfqpoint{1.167224in}{0.892929in}}%
\pgfpathlineto{\pgfqpoint{1.183580in}{0.998244in}}%
\pgfpathlineto{\pgfqpoint{1.198122in}{1.089991in}}%
\pgfpathlineto{\pgfqpoint{1.210751in}{1.171542in}}%
\pgfpathlineto{\pgfqpoint{1.222893in}{1.249723in}}%
\pgfusepath{stroke}%
\end{pgfscope}%
\begin{pgfscope}%
\pgfpathrectangle{\pgfqpoint{0.100000in}{0.100000in}}{\pgfqpoint{3.608454in}{2.310000in}}%
\pgfusepath{clip}%
\pgfsetbuttcap%
\pgfsetroundjoin%
\pgfsetlinewidth{0.010037pt}%
\definecolor{currentstroke}{rgb}{1.000000,1.000000,1.000000}%
\pgfsetstrokecolor{currentstroke}%
\pgfsetdash{}{0pt}%
\pgfpathmoveto{\pgfqpoint{2.789338in}{0.742945in}}%
\pgfpathlineto{\pgfqpoint{2.695351in}{0.732506in}}%
\pgfpathlineto{\pgfqpoint{2.612503in}{0.726074in}}%
\pgfpathlineto{\pgfqpoint{2.611467in}{0.715921in}}%
\pgfpathlineto{\pgfqpoint{2.619077in}{0.706258in}}%
\pgfpathlineto{\pgfqpoint{2.628365in}{0.700581in}}%
\pgfpathlineto{\pgfqpoint{2.628224in}{0.685633in}}%
\pgfpathlineto{\pgfqpoint{2.625763in}{0.675646in}}%
\pgfpathlineto{\pgfqpoint{2.617565in}{0.668449in}}%
\pgfpathlineto{\pgfqpoint{2.606139in}{0.669211in}}%
\pgfpathlineto{\pgfqpoint{2.595336in}{0.678128in}}%
\pgfpathlineto{\pgfqpoint{2.593410in}{0.693997in}}%
\pgfpathlineto{\pgfqpoint{2.585364in}{0.703227in}}%
\pgfpathlineto{\pgfqpoint{2.579887in}{0.670195in}}%
\pgfpathlineto{\pgfqpoint{2.561278in}{0.673334in}}%
\pgfpathlineto{\pgfqpoint{2.554807in}{0.731095in}}%
\pgfpathlineto{\pgfqpoint{2.547798in}{0.791968in}}%
\pgfpathlineto{\pgfqpoint{2.550401in}{0.879794in}}%
\pgfpathlineto{\pgfqpoint{2.552017in}{0.940061in}}%
\pgfpathlineto{\pgfqpoint{2.555453in}{1.032015in}}%
\pgfpathlineto{\pgfqpoint{2.548098in}{1.040107in}}%
\pgfpathlineto{\pgfqpoint{2.561482in}{1.042047in}}%
\pgfpathlineto{\pgfqpoint{2.639455in}{1.047089in}}%
\pgfpathlineto{\pgfqpoint{2.715320in}{1.053912in}}%
\pgfpathlineto{\pgfqpoint{2.735280in}{0.983998in}}%
\pgfpathlineto{\pgfqpoint{2.748892in}{0.932681in}}%
\pgfpathlineto{\pgfqpoint{2.761093in}{0.889718in}}%
\pgfpathlineto{\pgfqpoint{2.768151in}{0.872419in}}%
\pgfpathlineto{\pgfqpoint{2.779261in}{0.856321in}}%
\pgfpathlineto{\pgfqpoint{2.777467in}{0.848128in}}%
\pgfpathlineto{\pgfqpoint{2.785404in}{0.844137in}}%
\pgfpathlineto{\pgfqpoint{2.776025in}{0.831718in}}%
\pgfpathlineto{\pgfqpoint{2.772856in}{0.809602in}}%
\pgfpathlineto{\pgfqpoint{2.782023in}{0.783735in}}%
\pgfpathlineto{\pgfqpoint{2.780750in}{0.757706in}}%
\pgfpathlineto{\pgfqpoint{2.789338in}{0.742945in}}%
\pgfusepath{stroke}%
\end{pgfscope}%
\begin{pgfscope}%
\pgfpathrectangle{\pgfqpoint{0.100000in}{0.100000in}}{\pgfqpoint{3.608454in}{2.310000in}}%
\pgfusepath{clip}%
\pgfsetbuttcap%
\pgfsetroundjoin%
\pgfsetlinewidth{0.010037pt}%
\definecolor{currentstroke}{rgb}{1.000000,1.000000,1.000000}%
\pgfsetstrokecolor{currentstroke}%
\pgfsetdash{}{0pt}%
\pgfpathmoveto{\pgfqpoint{2.561278in}{0.673334in}}%
\pgfpathlineto{\pgfqpoint{2.542181in}{0.667869in}}%
\pgfpathlineto{\pgfqpoint{2.527946in}{0.671378in}}%
\pgfpathlineto{\pgfqpoint{2.501511in}{0.662970in}}%
\pgfpathlineto{\pgfqpoint{2.481574in}{0.652140in}}%
\pgfpathlineto{\pgfqpoint{2.473394in}{0.671709in}}%
\pgfpathlineto{\pgfqpoint{2.464961in}{0.679911in}}%
\pgfpathlineto{\pgfqpoint{2.460893in}{0.689726in}}%
\pgfpathlineto{\pgfqpoint{2.466856in}{0.716291in}}%
\pgfpathlineto{\pgfqpoint{2.410356in}{0.712812in}}%
\pgfpathlineto{\pgfqpoint{2.337092in}{0.709636in}}%
\pgfpathlineto{\pgfqpoint{2.341451in}{0.716248in}}%
\pgfpathlineto{\pgfqpoint{2.336063in}{0.728737in}}%
\pgfpathlineto{\pgfqpoint{2.341432in}{0.731282in}}%
\pgfpathlineto{\pgfqpoint{2.340241in}{0.743282in}}%
\pgfpathlineto{\pgfqpoint{2.349598in}{0.754918in}}%
\pgfpathlineto{\pgfqpoint{2.354717in}{0.777508in}}%
\pgfpathlineto{\pgfqpoint{2.371843in}{0.792391in}}%
\pgfpathlineto{\pgfqpoint{2.365495in}{0.806845in}}%
\pgfpathlineto{\pgfqpoint{2.377639in}{0.815165in}}%
\pgfpathlineto{\pgfqpoint{2.367164in}{0.830227in}}%
\pgfpathlineto{\pgfqpoint{2.359582in}{0.864708in}}%
\pgfpathlineto{\pgfqpoint{2.362458in}{0.870969in}}%
\pgfpathlineto{\pgfqpoint{2.366997in}{0.882978in}}%
\pgfpathlineto{\pgfqpoint{2.362772in}{0.895582in}}%
\pgfpathlineto{\pgfqpoint{2.364838in}{0.901316in}}%
\pgfpathlineto{\pgfqpoint{2.357405in}{0.922979in}}%
\pgfpathlineto{\pgfqpoint{2.379349in}{0.961875in}}%
\pgfpathlineto{\pgfqpoint{2.380370in}{0.972228in}}%
\pgfpathlineto{\pgfqpoint{2.390941in}{0.979756in}}%
\pgfpathlineto{\pgfqpoint{2.402220in}{1.004607in}}%
\pgfpathlineto{\pgfqpoint{2.401684in}{1.014709in}}%
\pgfpathlineto{\pgfqpoint{2.415733in}{1.025035in}}%
\pgfpathlineto{\pgfqpoint{2.411349in}{1.031151in}}%
\pgfpathlineto{\pgfqpoint{2.495536in}{1.036064in}}%
\pgfpathlineto{\pgfqpoint{2.548098in}{1.040107in}}%
\pgfpathlineto{\pgfqpoint{2.555453in}{1.032015in}}%
\pgfpathlineto{\pgfqpoint{2.552017in}{0.940061in}}%
\pgfpathlineto{\pgfqpoint{2.550401in}{0.879794in}}%
\pgfpathlineto{\pgfqpoint{2.547798in}{0.791968in}}%
\pgfpathlineto{\pgfqpoint{2.554807in}{0.731095in}}%
\pgfpathlineto{\pgfqpoint{2.561278in}{0.673334in}}%
\pgfusepath{stroke}%
\end{pgfscope}%
\begin{pgfscope}%
\pgfpathrectangle{\pgfqpoint{0.100000in}{0.100000in}}{\pgfqpoint{3.608454in}{2.310000in}}%
\pgfusepath{clip}%
\pgfsetbuttcap%
\pgfsetroundjoin%
\pgfsetlinewidth{0.010037pt}%
\definecolor{currentstroke}{rgb}{1.000000,1.000000,1.000000}%
\pgfsetstrokecolor{currentstroke}%
\pgfsetdash{}{0pt}%
\pgfpathmoveto{\pgfqpoint{2.715320in}{1.053912in}}%
\pgfpathlineto{\pgfqpoint{2.798250in}{1.063020in}}%
\pgfpathlineto{\pgfqpoint{2.848149in}{1.069168in}}%
\pgfpathlineto{\pgfqpoint{2.876568in}{1.073696in}}%
\pgfpathlineto{\pgfqpoint{2.864796in}{1.055128in}}%
\pgfpathlineto{\pgfqpoint{2.864803in}{1.046374in}}%
\pgfpathlineto{\pgfqpoint{2.876921in}{1.041653in}}%
\pgfpathlineto{\pgfqpoint{2.885162in}{1.034031in}}%
\pgfpathlineto{\pgfqpoint{2.897615in}{1.033089in}}%
\pgfpathlineto{\pgfqpoint{2.909255in}{1.011633in}}%
\pgfpathlineto{\pgfqpoint{2.921882in}{0.996503in}}%
\pgfpathlineto{\pgfqpoint{2.944306in}{0.983786in}}%
\pgfpathlineto{\pgfqpoint{2.950679in}{0.973643in}}%
\pgfpathlineto{\pgfqpoint{2.970676in}{0.962674in}}%
\pgfpathlineto{\pgfqpoint{2.972180in}{0.954719in}}%
\pgfpathlineto{\pgfqpoint{2.988402in}{0.938901in}}%
\pgfpathlineto{\pgfqpoint{3.000827in}{0.934397in}}%
\pgfpathlineto{\pgfqpoint{3.009626in}{0.919461in}}%
\pgfpathlineto{\pgfqpoint{3.013522in}{0.902673in}}%
\pgfpathlineto{\pgfqpoint{3.026428in}{0.896184in}}%
\pgfpathlineto{\pgfqpoint{3.036819in}{0.878138in}}%
\pgfpathlineto{\pgfqpoint{3.040169in}{0.864866in}}%
\pgfpathlineto{\pgfqpoint{3.055395in}{0.859218in}}%
\pgfpathlineto{\pgfqpoint{3.042619in}{0.834754in}}%
\pgfpathlineto{\pgfqpoint{3.037811in}{0.820297in}}%
\pgfpathlineto{\pgfqpoint{3.040948in}{0.813520in}}%
\pgfpathlineto{\pgfqpoint{3.035470in}{0.802258in}}%
\pgfpathlineto{\pgfqpoint{3.033142in}{0.786687in}}%
\pgfpathlineto{\pgfqpoint{3.023401in}{0.783787in}}%
\pgfpathlineto{\pgfqpoint{3.028545in}{0.769530in}}%
\pgfpathlineto{\pgfqpoint{3.028216in}{0.751443in}}%
\pgfpathlineto{\pgfqpoint{3.010902in}{0.752121in}}%
\pgfpathlineto{\pgfqpoint{2.998840in}{0.755421in}}%
\pgfpathlineto{\pgfqpoint{2.993620in}{0.749287in}}%
\pgfpathlineto{\pgfqpoint{2.997515in}{0.734795in}}%
\pgfpathlineto{\pgfqpoint{2.996668in}{0.717946in}}%
\pgfpathlineto{\pgfqpoint{2.989059in}{0.716657in}}%
\pgfpathlineto{\pgfqpoint{2.982903in}{0.732393in}}%
\pgfpathlineto{\pgfqpoint{2.920222in}{0.728215in}}%
\pgfpathlineto{\pgfqpoint{2.841186in}{0.723903in}}%
\pgfpathlineto{\pgfqpoint{2.801326in}{0.721097in}}%
\pgfpathlineto{\pgfqpoint{2.789338in}{0.742945in}}%
\pgfpathlineto{\pgfqpoint{2.780750in}{0.757706in}}%
\pgfpathlineto{\pgfqpoint{2.782023in}{0.783735in}}%
\pgfpathlineto{\pgfqpoint{2.772856in}{0.809602in}}%
\pgfpathlineto{\pgfqpoint{2.776025in}{0.831718in}}%
\pgfpathlineto{\pgfqpoint{2.785404in}{0.844137in}}%
\pgfpathlineto{\pgfqpoint{2.777467in}{0.848128in}}%
\pgfpathlineto{\pgfqpoint{2.779261in}{0.856321in}}%
\pgfpathlineto{\pgfqpoint{2.768151in}{0.872419in}}%
\pgfpathlineto{\pgfqpoint{2.761093in}{0.889718in}}%
\pgfpathlineto{\pgfqpoint{2.748892in}{0.932681in}}%
\pgfpathlineto{\pgfqpoint{2.735280in}{0.983998in}}%
\pgfpathlineto{\pgfqpoint{2.715320in}{1.053912in}}%
\pgfusepath{stroke}%
\end{pgfscope}%
\begin{pgfscope}%
\pgfpathrectangle{\pgfqpoint{0.100000in}{0.100000in}}{\pgfqpoint{3.608454in}{2.310000in}}%
\pgfusepath{clip}%
\pgfsetbuttcap%
\pgfsetroundjoin%
\pgfsetlinewidth{0.010037pt}%
\definecolor{currentstroke}{rgb}{1.000000,1.000000,1.000000}%
\pgfsetstrokecolor{currentstroke}%
\pgfsetdash{}{0pt}%
\pgfpathmoveto{\pgfqpoint{2.876568in}{1.073696in}}%
\pgfpathlineto{\pgfqpoint{2.909453in}{1.090054in}}%
\pgfpathlineto{\pgfqpoint{2.927665in}{1.096231in}}%
\pgfpathlineto{\pgfqpoint{3.006902in}{1.104464in}}%
\pgfpathlineto{\pgfqpoint{3.015293in}{1.101737in}}%
\pgfpathlineto{\pgfqpoint{3.026378in}{1.090595in}}%
\pgfpathlineto{\pgfqpoint{3.026989in}{1.080701in}}%
\pgfpathlineto{\pgfqpoint{3.099917in}{1.091416in}}%
\pgfpathlineto{\pgfqpoint{3.182729in}{1.032053in}}%
\pgfpathlineto{\pgfqpoint{3.167199in}{1.015867in}}%
\pgfpathlineto{\pgfqpoint{3.145890in}{0.978275in}}%
\pgfpathlineto{\pgfqpoint{3.151955in}{0.970196in}}%
\pgfpathlineto{\pgfqpoint{3.140588in}{0.954516in}}%
\pgfpathlineto{\pgfqpoint{3.129280in}{0.952741in}}%
\pgfpathlineto{\pgfqpoint{3.129238in}{0.943317in}}%
\pgfpathlineto{\pgfqpoint{3.121059in}{0.933457in}}%
\pgfpathlineto{\pgfqpoint{3.110879in}{0.931458in}}%
\pgfpathlineto{\pgfqpoint{3.113100in}{0.922697in}}%
\pgfpathlineto{\pgfqpoint{3.107420in}{0.915971in}}%
\pgfpathlineto{\pgfqpoint{3.093832in}{0.910161in}}%
\pgfpathlineto{\pgfqpoint{3.077270in}{0.896945in}}%
\pgfpathlineto{\pgfqpoint{3.080396in}{0.888394in}}%
\pgfpathlineto{\pgfqpoint{3.069973in}{0.883025in}}%
\pgfpathlineto{\pgfqpoint{3.061181in}{0.888670in}}%
\pgfpathlineto{\pgfqpoint{3.054753in}{0.864128in}}%
\pgfpathlineto{\pgfqpoint{3.040169in}{0.864866in}}%
\pgfpathlineto{\pgfqpoint{3.036819in}{0.878138in}}%
\pgfpathlineto{\pgfqpoint{3.026428in}{0.896184in}}%
\pgfpathlineto{\pgfqpoint{3.013522in}{0.902673in}}%
\pgfpathlineto{\pgfqpoint{3.009626in}{0.919461in}}%
\pgfpathlineto{\pgfqpoint{3.000827in}{0.934397in}}%
\pgfpathlineto{\pgfqpoint{2.988402in}{0.938901in}}%
\pgfpathlineto{\pgfqpoint{2.972180in}{0.954719in}}%
\pgfpathlineto{\pgfqpoint{2.970676in}{0.962674in}}%
\pgfpathlineto{\pgfqpoint{2.950679in}{0.973643in}}%
\pgfpathlineto{\pgfqpoint{2.944306in}{0.983786in}}%
\pgfpathlineto{\pgfqpoint{2.921882in}{0.996503in}}%
\pgfpathlineto{\pgfqpoint{2.909255in}{1.011633in}}%
\pgfpathlineto{\pgfqpoint{2.897615in}{1.033089in}}%
\pgfpathlineto{\pgfqpoint{2.885162in}{1.034031in}}%
\pgfpathlineto{\pgfqpoint{2.876921in}{1.041653in}}%
\pgfpathlineto{\pgfqpoint{2.864803in}{1.046374in}}%
\pgfpathlineto{\pgfqpoint{2.864796in}{1.055128in}}%
\pgfpathlineto{\pgfqpoint{2.876568in}{1.073696in}}%
\pgfusepath{stroke}%
\end{pgfscope}%
\begin{pgfscope}%
\pgfpathrectangle{\pgfqpoint{0.100000in}{0.100000in}}{\pgfqpoint{3.608454in}{2.310000in}}%
\pgfusepath{clip}%
\pgfsetbuttcap%
\pgfsetroundjoin%
\pgfsetlinewidth{0.010037pt}%
\definecolor{currentstroke}{rgb}{1.000000,1.000000,1.000000}%
\pgfsetstrokecolor{currentstroke}%
\pgfsetdash{}{0pt}%
\pgfpathmoveto{\pgfqpoint{2.131124in}{1.141195in}}%
\pgfpathlineto{\pgfqpoint{2.179448in}{1.141630in}}%
\pgfpathlineto{\pgfqpoint{2.243380in}{1.142795in}}%
\pgfpathlineto{\pgfqpoint{2.352626in}{1.146263in}}%
\pgfpathlineto{\pgfqpoint{2.415111in}{1.149567in}}%
\pgfpathlineto{\pgfqpoint{2.421813in}{1.141271in}}%
\pgfpathlineto{\pgfqpoint{2.421394in}{1.132509in}}%
\pgfpathlineto{\pgfqpoint{2.406257in}{1.117394in}}%
\pgfpathlineto{\pgfqpoint{2.402597in}{1.109117in}}%
\pgfpathlineto{\pgfqpoint{2.444579in}{1.112222in}}%
\pgfpathlineto{\pgfqpoint{2.444556in}{1.096930in}}%
\pgfpathlineto{\pgfqpoint{2.431088in}{1.090395in}}%
\pgfpathlineto{\pgfqpoint{2.431278in}{1.080004in}}%
\pgfpathlineto{\pgfqpoint{2.426220in}{1.072535in}}%
\pgfpathlineto{\pgfqpoint{2.422910in}{1.056671in}}%
\pgfpathlineto{\pgfqpoint{2.426327in}{1.043584in}}%
\pgfpathlineto{\pgfqpoint{2.411349in}{1.031151in}}%
\pgfpathlineto{\pgfqpoint{2.415733in}{1.025035in}}%
\pgfpathlineto{\pgfqpoint{2.401684in}{1.014709in}}%
\pgfpathlineto{\pgfqpoint{2.402220in}{1.004607in}}%
\pgfpathlineto{\pgfqpoint{2.390941in}{0.979756in}}%
\pgfpathlineto{\pgfqpoint{2.380370in}{0.972228in}}%
\pgfpathlineto{\pgfqpoint{2.379349in}{0.961875in}}%
\pgfpathlineto{\pgfqpoint{2.357405in}{0.922979in}}%
\pgfpathlineto{\pgfqpoint{2.364838in}{0.901316in}}%
\pgfpathlineto{\pgfqpoint{2.362772in}{0.895582in}}%
\pgfpathlineto{\pgfqpoint{2.366997in}{0.882978in}}%
\pgfpathlineto{\pgfqpoint{2.362458in}{0.870969in}}%
\pgfpathlineto{\pgfqpoint{2.302450in}{0.868492in}}%
\pgfpathlineto{\pgfqpoint{2.224548in}{0.867208in}}%
\pgfpathlineto{\pgfqpoint{2.170816in}{0.866723in}}%
\pgfpathlineto{\pgfqpoint{2.170554in}{0.909031in}}%
\pgfpathlineto{\pgfqpoint{2.157210in}{0.911803in}}%
\pgfpathlineto{\pgfqpoint{2.148404in}{0.908183in}}%
\pgfpathlineto{\pgfqpoint{2.141369in}{0.914830in}}%
\pgfpathlineto{\pgfqpoint{2.142095in}{0.959646in}}%
\pgfpathlineto{\pgfqpoint{2.143669in}{1.055067in}}%
\pgfpathlineto{\pgfqpoint{2.136037in}{1.110943in}}%
\pgfpathlineto{\pgfqpoint{2.131124in}{1.141195in}}%
\pgfusepath{stroke}%
\end{pgfscope}%
\begin{pgfscope}%
\pgfpathrectangle{\pgfqpoint{0.100000in}{0.100000in}}{\pgfqpoint{3.608454in}{2.310000in}}%
\pgfusepath{clip}%
\pgfsetbuttcap%
\pgfsetroundjoin%
\pgfsetlinewidth{0.010037pt}%
\definecolor{currentstroke}{rgb}{1.000000,1.000000,1.000000}%
\pgfsetstrokecolor{currentstroke}%
\pgfsetdash{}{0pt}%
\pgfpathmoveto{\pgfqpoint{2.170816in}{0.866723in}}%
\pgfpathlineto{\pgfqpoint{2.224548in}{0.867208in}}%
\pgfpathlineto{\pgfqpoint{2.302450in}{0.868492in}}%
\pgfpathlineto{\pgfqpoint{2.362458in}{0.870969in}}%
\pgfpathlineto{\pgfqpoint{2.359582in}{0.864708in}}%
\pgfpathlineto{\pgfqpoint{2.367164in}{0.830227in}}%
\pgfpathlineto{\pgfqpoint{2.377639in}{0.815165in}}%
\pgfpathlineto{\pgfqpoint{2.365495in}{0.806845in}}%
\pgfpathlineto{\pgfqpoint{2.371843in}{0.792391in}}%
\pgfpathlineto{\pgfqpoint{2.354717in}{0.777508in}}%
\pgfpathlineto{\pgfqpoint{2.349598in}{0.754918in}}%
\pgfpathlineto{\pgfqpoint{2.340241in}{0.743282in}}%
\pgfpathlineto{\pgfqpoint{2.341432in}{0.731282in}}%
\pgfpathlineto{\pgfqpoint{2.336063in}{0.728737in}}%
\pgfpathlineto{\pgfqpoint{2.341451in}{0.716248in}}%
\pgfpathlineto{\pgfqpoint{2.337092in}{0.709636in}}%
\pgfpathlineto{\pgfqpoint{2.410356in}{0.712812in}}%
\pgfpathlineto{\pgfqpoint{2.466856in}{0.716291in}}%
\pgfpathlineto{\pgfqpoint{2.460893in}{0.689726in}}%
\pgfpathlineto{\pgfqpoint{2.464961in}{0.679911in}}%
\pgfpathlineto{\pgfqpoint{2.473394in}{0.671709in}}%
\pgfpathlineto{\pgfqpoint{2.481574in}{0.652140in}}%
\pgfpathlineto{\pgfqpoint{2.455735in}{0.656649in}}%
\pgfpathlineto{\pgfqpoint{2.446204in}{0.664060in}}%
\pgfpathlineto{\pgfqpoint{2.434847in}{0.664413in}}%
\pgfpathlineto{\pgfqpoint{2.422908in}{0.648183in}}%
\pgfpathlineto{\pgfqpoint{2.425288in}{0.640793in}}%
\pgfpathlineto{\pgfqpoint{2.463429in}{0.636322in}}%
\pgfpathlineto{\pgfqpoint{2.473414in}{0.627827in}}%
\pgfpathlineto{\pgfqpoint{2.489663in}{0.623456in}}%
\pgfpathlineto{\pgfqpoint{2.482567in}{0.613398in}}%
\pgfpathlineto{\pgfqpoint{2.470611in}{0.604627in}}%
\pgfpathlineto{\pgfqpoint{2.502169in}{0.584902in}}%
\pgfpathlineto{\pgfqpoint{2.516869in}{0.581820in}}%
\pgfpathlineto{\pgfqpoint{2.522982in}{0.565696in}}%
\pgfpathlineto{\pgfqpoint{2.509142in}{0.562730in}}%
\pgfpathlineto{\pgfqpoint{2.483128in}{0.583009in}}%
\pgfpathlineto{\pgfqpoint{2.472958in}{0.585828in}}%
\pgfpathlineto{\pgfqpoint{2.468036in}{0.593839in}}%
\pgfpathlineto{\pgfqpoint{2.452950in}{0.590551in}}%
\pgfpathlineto{\pgfqpoint{2.448226in}{0.580228in}}%
\pgfpathlineto{\pgfqpoint{2.451218in}{0.568725in}}%
\pgfpathlineto{\pgfqpoint{2.441061in}{0.561938in}}%
\pgfpathlineto{\pgfqpoint{2.431792in}{0.578660in}}%
\pgfpathlineto{\pgfqpoint{2.415557in}{0.573674in}}%
\pgfpathlineto{\pgfqpoint{2.409474in}{0.563696in}}%
\pgfpathlineto{\pgfqpoint{2.397944in}{0.566534in}}%
\pgfpathlineto{\pgfqpoint{2.390573in}{0.579135in}}%
\pgfpathlineto{\pgfqpoint{2.371016in}{0.583295in}}%
\pgfpathlineto{\pgfqpoint{2.367323in}{0.589852in}}%
\pgfpathlineto{\pgfqpoint{2.346947in}{0.601273in}}%
\pgfpathlineto{\pgfqpoint{2.341889in}{0.611259in}}%
\pgfpathlineto{\pgfqpoint{2.324822in}{0.607178in}}%
\pgfpathlineto{\pgfqpoint{2.326990in}{0.616320in}}%
\pgfpathlineto{\pgfqpoint{2.305777in}{0.606942in}}%
\pgfpathlineto{\pgfqpoint{2.311475in}{0.597211in}}%
\pgfpathlineto{\pgfqpoint{2.295094in}{0.591458in}}%
\pgfpathlineto{\pgfqpoint{2.273396in}{0.594620in}}%
\pgfpathlineto{\pgfqpoint{2.229480in}{0.609663in}}%
\pgfpathlineto{\pgfqpoint{2.195594in}{0.606677in}}%
\pgfpathlineto{\pgfqpoint{2.192657in}{0.626436in}}%
\pgfpathlineto{\pgfqpoint{2.196912in}{0.635452in}}%
\pgfpathlineto{\pgfqpoint{2.196548in}{0.649798in}}%
\pgfpathlineto{\pgfqpoint{2.192362in}{0.653216in}}%
\pgfpathlineto{\pgfqpoint{2.193781in}{0.669590in}}%
\pgfpathlineto{\pgfqpoint{2.206122in}{0.692365in}}%
\pgfpathlineto{\pgfqpoint{2.207701in}{0.700970in}}%
\pgfpathlineto{\pgfqpoint{2.205667in}{0.721296in}}%
\pgfpathlineto{\pgfqpoint{2.196588in}{0.730869in}}%
\pgfpathlineto{\pgfqpoint{2.191847in}{0.748745in}}%
\pgfpathlineto{\pgfqpoint{2.186014in}{0.752849in}}%
\pgfpathlineto{\pgfqpoint{2.188625in}{0.762543in}}%
\pgfpathlineto{\pgfqpoint{2.181197in}{0.777032in}}%
\pgfpathlineto{\pgfqpoint{2.171929in}{0.784895in}}%
\pgfpathlineto{\pgfqpoint{2.170816in}{0.866723in}}%
\pgfusepath{stroke}%
\end{pgfscope}%
\begin{pgfscope}%
\pgfpathrectangle{\pgfqpoint{0.100000in}{0.100000in}}{\pgfqpoint{3.608454in}{2.310000in}}%
\pgfusepath{clip}%
\pgfsetbuttcap%
\pgfsetroundjoin%
\pgfsetlinewidth{0.010037pt}%
\definecolor{currentstroke}{rgb}{1.000000,1.000000,1.000000}%
\pgfsetstrokecolor{currentstroke}%
\pgfsetdash{}{0pt}%
\pgfpathmoveto{\pgfqpoint{2.314523in}{0.596466in}}%
\pgfpathlineto{\pgfqpoint{2.322296in}{0.601090in}}%
\pgfpathlineto{\pgfqpoint{2.331731in}{0.595623in}}%
\pgfpathlineto{\pgfqpoint{2.326469in}{0.588102in}}%
\pgfpathlineto{\pgfqpoint{2.314523in}{0.596466in}}%
\pgfusepath{stroke}%
\end{pgfscope}%
\begin{pgfscope}%
\pgfpathrectangle{\pgfqpoint{0.100000in}{0.100000in}}{\pgfqpoint{3.608454in}{2.310000in}}%
\pgfusepath{clip}%
\pgfsetbuttcap%
\pgfsetroundjoin%
\pgfsetlinewidth{0.010037pt}%
\definecolor{currentstroke}{rgb}{1.000000,1.000000,1.000000}%
\pgfsetstrokecolor{currentstroke}%
\pgfsetdash{}{0pt}%
\pgfpathmoveto{\pgfqpoint{2.628224in}{0.685633in}}%
\pgfpathlineto{\pgfqpoint{2.628365in}{0.700581in}}%
\pgfpathlineto{\pgfqpoint{2.619077in}{0.706258in}}%
\pgfpathlineto{\pgfqpoint{2.611467in}{0.715921in}}%
\pgfpathlineto{\pgfqpoint{2.612503in}{0.726074in}}%
\pgfpathlineto{\pgfqpoint{2.695351in}{0.732506in}}%
\pgfpathlineto{\pgfqpoint{2.789338in}{0.742945in}}%
\pgfpathlineto{\pgfqpoint{2.801326in}{0.721097in}}%
\pgfpathlineto{\pgfqpoint{2.841186in}{0.723903in}}%
\pgfpathlineto{\pgfqpoint{2.920222in}{0.728215in}}%
\pgfpathlineto{\pgfqpoint{2.982903in}{0.732393in}}%
\pgfpathlineto{\pgfqpoint{2.989059in}{0.716657in}}%
\pgfpathlineto{\pgfqpoint{2.996668in}{0.717946in}}%
\pgfpathlineto{\pgfqpoint{2.997515in}{0.734795in}}%
\pgfpathlineto{\pgfqpoint{2.993620in}{0.749287in}}%
\pgfpathlineto{\pgfqpoint{2.998840in}{0.755421in}}%
\pgfpathlineto{\pgfqpoint{3.010902in}{0.752121in}}%
\pgfpathlineto{\pgfqpoint{3.028216in}{0.751443in}}%
\pgfpathlineto{\pgfqpoint{3.030878in}{0.738496in}}%
\pgfpathlineto{\pgfqpoint{3.036204in}{0.731098in}}%
\pgfpathlineto{\pgfqpoint{3.040364in}{0.714912in}}%
\pgfpathlineto{\pgfqpoint{3.053250in}{0.689862in}}%
\pgfpathlineto{\pgfqpoint{3.053315in}{0.683088in}}%
\pgfpathlineto{\pgfqpoint{3.072361in}{0.653675in}}%
\pgfpathlineto{\pgfqpoint{3.074204in}{0.647599in}}%
\pgfpathlineto{\pgfqpoint{3.102776in}{0.606058in}}%
\pgfpathlineto{\pgfqpoint{3.098249in}{0.605387in}}%
\pgfpathlineto{\pgfqpoint{3.110295in}{0.575840in}}%
\pgfpathlineto{\pgfqpoint{3.143385in}{0.524475in}}%
\pgfpathlineto{\pgfqpoint{3.162598in}{0.486707in}}%
\pgfpathlineto{\pgfqpoint{3.169019in}{0.480095in}}%
\pgfpathlineto{\pgfqpoint{3.180022in}{0.456389in}}%
\pgfpathlineto{\pgfqpoint{3.183839in}{0.418430in}}%
\pgfpathlineto{\pgfqpoint{3.185359in}{0.390048in}}%
\pgfpathlineto{\pgfqpoint{3.183524in}{0.371875in}}%
\pgfpathlineto{\pgfqpoint{3.177624in}{0.358884in}}%
\pgfpathlineto{\pgfqpoint{3.180365in}{0.341848in}}%
\pgfpathlineto{\pgfqpoint{3.174013in}{0.328363in}}%
\pgfpathlineto{\pgfqpoint{3.155154in}{0.317306in}}%
\pgfpathlineto{\pgfqpoint{3.142888in}{0.318113in}}%
\pgfpathlineto{\pgfqpoint{3.134923in}{0.312333in}}%
\pgfpathlineto{\pgfqpoint{3.132641in}{0.327766in}}%
\pgfpathlineto{\pgfqpoint{3.119456in}{0.331840in}}%
\pgfpathlineto{\pgfqpoint{3.107637in}{0.353374in}}%
\pgfpathlineto{\pgfqpoint{3.106303in}{0.363180in}}%
\pgfpathlineto{\pgfqpoint{3.085149in}{0.369136in}}%
\pgfpathlineto{\pgfqpoint{3.071519in}{0.367846in}}%
\pgfpathlineto{\pgfqpoint{3.063806in}{0.382075in}}%
\pgfpathlineto{\pgfqpoint{3.054976in}{0.407650in}}%
\pgfpathlineto{\pgfqpoint{3.042722in}{0.412834in}}%
\pgfpathlineto{\pgfqpoint{3.036079in}{0.427490in}}%
\pgfpathlineto{\pgfqpoint{3.019832in}{0.436127in}}%
\pgfpathlineto{\pgfqpoint{3.010441in}{0.446931in}}%
\pgfpathlineto{\pgfqpoint{2.994241in}{0.476304in}}%
\pgfpathlineto{\pgfqpoint{2.992335in}{0.496978in}}%
\pgfpathlineto{\pgfqpoint{3.001163in}{0.510205in}}%
\pgfpathlineto{\pgfqpoint{3.000281in}{0.519397in}}%
\pgfpathlineto{\pgfqpoint{2.990079in}{0.520379in}}%
\pgfpathlineto{\pgfqpoint{2.981556in}{0.526772in}}%
\pgfpathlineto{\pgfqpoint{2.976889in}{0.518966in}}%
\pgfpathlineto{\pgfqpoint{2.985069in}{0.512605in}}%
\pgfpathlineto{\pgfqpoint{2.979115in}{0.501166in}}%
\pgfpathlineto{\pgfqpoint{2.969483in}{0.510660in}}%
\pgfpathlineto{\pgfqpoint{2.970600in}{0.536987in}}%
\pgfpathlineto{\pgfqpoint{2.975246in}{0.558347in}}%
\pgfpathlineto{\pgfqpoint{2.972858in}{0.594997in}}%
\pgfpathlineto{\pgfqpoint{2.963257in}{0.603726in}}%
\pgfpathlineto{\pgfqpoint{2.958429in}{0.614928in}}%
\pgfpathlineto{\pgfqpoint{2.941890in}{0.614679in}}%
\pgfpathlineto{\pgfqpoint{2.925541in}{0.633111in}}%
\pgfpathlineto{\pgfqpoint{2.914580in}{0.638635in}}%
\pgfpathlineto{\pgfqpoint{2.911329in}{0.650308in}}%
\pgfpathlineto{\pgfqpoint{2.900608in}{0.654407in}}%
\pgfpathlineto{\pgfqpoint{2.891722in}{0.667295in}}%
\pgfpathlineto{\pgfqpoint{2.868250in}{0.677815in}}%
\pgfpathlineto{\pgfqpoint{2.850001in}{0.678084in}}%
\pgfpathlineto{\pgfqpoint{2.842057in}{0.674044in}}%
\pgfpathlineto{\pgfqpoint{2.844031in}{0.661431in}}%
\pgfpathlineto{\pgfqpoint{2.835747in}{0.662028in}}%
\pgfpathlineto{\pgfqpoint{2.810265in}{0.644393in}}%
\pgfpathlineto{\pgfqpoint{2.779657in}{0.637403in}}%
\pgfpathlineto{\pgfqpoint{2.779155in}{0.646052in}}%
\pgfpathlineto{\pgfqpoint{2.772372in}{0.654504in}}%
\pgfpathlineto{\pgfqpoint{2.754172in}{0.666173in}}%
\pgfpathlineto{\pgfqpoint{2.728049in}{0.678137in}}%
\pgfpathlineto{\pgfqpoint{2.699674in}{0.684471in}}%
\pgfpathlineto{\pgfqpoint{2.694348in}{0.693083in}}%
\pgfpathlineto{\pgfqpoint{2.684117in}{0.685900in}}%
\pgfpathlineto{\pgfqpoint{2.671797in}{0.684316in}}%
\pgfpathlineto{\pgfqpoint{2.628914in}{0.673017in}}%
\pgfpathlineto{\pgfqpoint{2.628224in}{0.685633in}}%
\pgfusepath{stroke}%
\end{pgfscope}%
\begin{pgfscope}%
\pgfpathrectangle{\pgfqpoint{0.100000in}{0.100000in}}{\pgfqpoint{3.608454in}{2.310000in}}%
\pgfusepath{clip}%
\pgfsetbuttcap%
\pgfsetroundjoin%
\pgfsetlinewidth{0.010037pt}%
\definecolor{currentstroke}{rgb}{1.000000,1.000000,1.000000}%
\pgfsetstrokecolor{currentstroke}%
\pgfsetdash{}{0pt}%
\pgfpathmoveto{\pgfqpoint{3.106940in}{0.606649in}}%
\pgfpathlineto{\pgfqpoint{3.118035in}{0.593742in}}%
\pgfpathlineto{\pgfqpoint{3.105219in}{0.592930in}}%
\pgfpathlineto{\pgfqpoint{3.106940in}{0.606649in}}%
\pgfusepath{stroke}%
\end{pgfscope}%
\begin{pgfscope}%
\pgfpathrectangle{\pgfqpoint{0.100000in}{0.100000in}}{\pgfqpoint{3.608454in}{2.310000in}}%
\pgfusepath{clip}%
\pgfsetbuttcap%
\pgfsetroundjoin%
\pgfsetlinewidth{0.010037pt}%
\definecolor{currentstroke}{rgb}{1.000000,1.000000,1.000000}%
\pgfsetstrokecolor{currentstroke}%
\pgfsetdash{}{0pt}%
\pgfpathmoveto{\pgfqpoint{2.453791in}{2.080299in}}%
\pgfpathlineto{\pgfqpoint{2.447771in}{2.068582in}}%
\pgfpathlineto{\pgfqpoint{2.433410in}{2.061744in}}%
\pgfpathlineto{\pgfqpoint{2.427189in}{2.052541in}}%
\pgfpathlineto{\pgfqpoint{2.419903in}{2.059169in}}%
\pgfpathlineto{\pgfqpoint{2.453791in}{2.080299in}}%
\pgfusepath{stroke}%
\end{pgfscope}%
\begin{pgfscope}%
\pgfpathrectangle{\pgfqpoint{0.100000in}{0.100000in}}{\pgfqpoint{3.608454in}{2.310000in}}%
\pgfusepath{clip}%
\pgfsetbuttcap%
\pgfsetroundjoin%
\pgfsetlinewidth{0.010037pt}%
\definecolor{currentstroke}{rgb}{1.000000,1.000000,1.000000}%
\pgfsetstrokecolor{currentstroke}%
\pgfsetdash{}{0pt}%
\pgfpathmoveto{\pgfqpoint{2.458673in}{2.009787in}}%
\pgfpathlineto{\pgfqpoint{2.473368in}{2.023495in}}%
\pgfpathlineto{\pgfqpoint{2.496033in}{2.027084in}}%
\pgfpathlineto{\pgfqpoint{2.489790in}{2.017551in}}%
\pgfpathlineto{\pgfqpoint{2.474252in}{2.003764in}}%
\pgfpathlineto{\pgfqpoint{2.465168in}{1.986068in}}%
\pgfpathlineto{\pgfqpoint{2.461321in}{1.995666in}}%
\pgfpathlineto{\pgfqpoint{2.454489in}{1.997044in}}%
\pgfpathlineto{\pgfqpoint{2.453803in}{2.005709in}}%
\pgfpathlineto{\pgfqpoint{2.458673in}{2.009787in}}%
\pgfusepath{stroke}%
\end{pgfscope}%
\begin{pgfscope}%
\pgfpathrectangle{\pgfqpoint{0.100000in}{0.100000in}}{\pgfqpoint{3.608454in}{2.310000in}}%
\pgfusepath{clip}%
\pgfsetbuttcap%
\pgfsetroundjoin%
\pgfsetlinewidth{0.010037pt}%
\definecolor{currentstroke}{rgb}{1.000000,1.000000,1.000000}%
\pgfsetstrokecolor{currentstroke}%
\pgfsetdash{}{0pt}%
\pgfpathmoveto{\pgfqpoint{2.517200in}{1.842472in}}%
\pgfpathlineto{\pgfqpoint{2.513287in}{1.846815in}}%
\pgfpathlineto{\pgfqpoint{2.517401in}{1.859041in}}%
\pgfpathlineto{\pgfqpoint{2.505188in}{1.859820in}}%
\pgfpathlineto{\pgfqpoint{2.508426in}{1.870361in}}%
\pgfpathlineto{\pgfqpoint{2.506412in}{1.887150in}}%
\pgfpathlineto{\pgfqpoint{2.495454in}{1.892970in}}%
\pgfpathlineto{\pgfqpoint{2.483987in}{1.904750in}}%
\pgfpathlineto{\pgfqpoint{2.466317in}{1.908198in}}%
\pgfpathlineto{\pgfqpoint{2.449120in}{1.908100in}}%
\pgfpathlineto{\pgfqpoint{2.432207in}{1.916460in}}%
\pgfpathlineto{\pgfqpoint{2.375622in}{1.928645in}}%
\pgfpathlineto{\pgfqpoint{2.369439in}{1.941557in}}%
\pgfpathlineto{\pgfqpoint{2.358409in}{1.945963in}}%
\pgfpathlineto{\pgfqpoint{2.379243in}{1.955848in}}%
\pgfpathlineto{\pgfqpoint{2.391013in}{1.968181in}}%
\pgfpathlineto{\pgfqpoint{2.412935in}{1.971526in}}%
\pgfpathlineto{\pgfqpoint{2.426425in}{1.984088in}}%
\pgfpathlineto{\pgfqpoint{2.433502in}{1.984591in}}%
\pgfpathlineto{\pgfqpoint{2.438904in}{1.993550in}}%
\pgfpathlineto{\pgfqpoint{2.452070in}{2.004163in}}%
\pgfpathlineto{\pgfqpoint{2.453218in}{1.996662in}}%
\pgfpathlineto{\pgfqpoint{2.459151in}{1.995113in}}%
\pgfpathlineto{\pgfqpoint{2.463615in}{1.986167in}}%
\pgfpathlineto{\pgfqpoint{2.464419in}{1.970914in}}%
\pgfpathlineto{\pgfqpoint{2.477835in}{1.980029in}}%
\pgfpathlineto{\pgfqpoint{2.493468in}{1.981933in}}%
\pgfpathlineto{\pgfqpoint{2.506813in}{1.977158in}}%
\pgfpathlineto{\pgfqpoint{2.524910in}{1.952330in}}%
\pgfpathlineto{\pgfqpoint{2.544676in}{1.956298in}}%
\pgfpathlineto{\pgfqpoint{2.552707in}{1.949654in}}%
\pgfpathlineto{\pgfqpoint{2.565625in}{1.949060in}}%
\pgfpathlineto{\pgfqpoint{2.574207in}{1.960977in}}%
\pgfpathlineto{\pgfqpoint{2.590492in}{1.971511in}}%
\pgfpathlineto{\pgfqpoint{2.606143in}{1.974810in}}%
\pgfpathlineto{\pgfqpoint{2.625564in}{1.975153in}}%
\pgfpathlineto{\pgfqpoint{2.639756in}{1.983291in}}%
\pgfpathlineto{\pgfqpoint{2.651338in}{1.979543in}}%
\pgfpathlineto{\pgfqpoint{2.653805in}{1.962333in}}%
\pgfpathlineto{\pgfqpoint{2.678653in}{1.959638in}}%
\pgfpathlineto{\pgfqpoint{2.692149in}{1.967694in}}%
\pgfpathlineto{\pgfqpoint{2.701505in}{1.949546in}}%
\pgfpathlineto{\pgfqpoint{2.719357in}{1.928560in}}%
\pgfpathlineto{\pgfqpoint{2.694371in}{1.928682in}}%
\pgfpathlineto{\pgfqpoint{2.686477in}{1.926088in}}%
\pgfpathlineto{\pgfqpoint{2.675655in}{1.929459in}}%
\pgfpathlineto{\pgfqpoint{2.674927in}{1.914931in}}%
\pgfpathlineto{\pgfqpoint{2.655278in}{1.926284in}}%
\pgfpathlineto{\pgfqpoint{2.630009in}{1.929754in}}%
\pgfpathlineto{\pgfqpoint{2.623052in}{1.918697in}}%
\pgfpathlineto{\pgfqpoint{2.589991in}{1.913321in}}%
\pgfpathlineto{\pgfqpoint{2.586187in}{1.903954in}}%
\pgfpathlineto{\pgfqpoint{2.563180in}{1.901144in}}%
\pgfpathlineto{\pgfqpoint{2.556225in}{1.891622in}}%
\pgfpathlineto{\pgfqpoint{2.544055in}{1.889069in}}%
\pgfpathlineto{\pgfqpoint{2.534335in}{1.866491in}}%
\pgfpathlineto{\pgfqpoint{2.522029in}{1.844624in}}%
\pgfpathlineto{\pgfqpoint{2.517200in}{1.842472in}}%
\pgfusepath{stroke}%
\end{pgfscope}%
\begin{pgfscope}%
\pgfpathrectangle{\pgfqpoint{0.100000in}{0.100000in}}{\pgfqpoint{3.608454in}{2.310000in}}%
\pgfusepath{clip}%
\pgfsetbuttcap%
\pgfsetroundjoin%
\pgfsetlinewidth{0.010037pt}%
\definecolor{currentstroke}{rgb}{1.000000,1.000000,1.000000}%
\pgfsetstrokecolor{currentstroke}%
\pgfsetdash{}{0pt}%
\pgfpathmoveto{\pgfqpoint{2.588031in}{1.579945in}}%
\pgfpathlineto{\pgfqpoint{2.599764in}{1.592295in}}%
\pgfpathlineto{\pgfqpoint{2.605110in}{1.610201in}}%
\pgfpathlineto{\pgfqpoint{2.611467in}{1.620579in}}%
\pgfpathlineto{\pgfqpoint{2.615365in}{1.634703in}}%
\pgfpathlineto{\pgfqpoint{2.616576in}{1.662865in}}%
\pgfpathlineto{\pgfqpoint{2.610707in}{1.689854in}}%
\pgfpathlineto{\pgfqpoint{2.591312in}{1.731182in}}%
\pgfpathlineto{\pgfqpoint{2.596565in}{1.744177in}}%
\pgfpathlineto{\pgfqpoint{2.589723in}{1.762138in}}%
\pgfpathlineto{\pgfqpoint{2.601460in}{1.786979in}}%
\pgfpathlineto{\pgfqpoint{2.599547in}{1.814676in}}%
\pgfpathlineto{\pgfqpoint{2.607723in}{1.818162in}}%
\pgfpathlineto{\pgfqpoint{2.608750in}{1.831341in}}%
\pgfpathlineto{\pgfqpoint{2.623262in}{1.839799in}}%
\pgfpathlineto{\pgfqpoint{2.631469in}{1.838434in}}%
\pgfpathlineto{\pgfqpoint{2.633776in}{1.824300in}}%
\pgfpathlineto{\pgfqpoint{2.640182in}{1.823746in}}%
\pgfpathlineto{\pgfqpoint{2.646027in}{1.844133in}}%
\pgfpathlineto{\pgfqpoint{2.644040in}{1.859989in}}%
\pgfpathlineto{\pgfqpoint{2.647841in}{1.869093in}}%
\pgfpathlineto{\pgfqpoint{2.668387in}{1.878499in}}%
\pgfpathlineto{\pgfqpoint{2.659024in}{1.881847in}}%
\pgfpathlineto{\pgfqpoint{2.655974in}{1.889927in}}%
\pgfpathlineto{\pgfqpoint{2.662719in}{1.904115in}}%
\pgfpathlineto{\pgfqpoint{2.676020in}{1.909016in}}%
\pgfpathlineto{\pgfqpoint{2.691465in}{1.900645in}}%
\pgfpathlineto{\pgfqpoint{2.706027in}{1.900540in}}%
\pgfpathlineto{\pgfqpoint{2.712790in}{1.890761in}}%
\pgfpathlineto{\pgfqpoint{2.722988in}{1.891440in}}%
\pgfpathlineto{\pgfqpoint{2.754470in}{1.877846in}}%
\pgfpathlineto{\pgfqpoint{2.760759in}{1.864686in}}%
\pgfpathlineto{\pgfqpoint{2.753890in}{1.860128in}}%
\pgfpathlineto{\pgfqpoint{2.756007in}{1.850265in}}%
\pgfpathlineto{\pgfqpoint{2.762793in}{1.845879in}}%
\pgfpathlineto{\pgfqpoint{2.766575in}{1.833755in}}%
\pgfpathlineto{\pgfqpoint{2.766045in}{1.804218in}}%
\pgfpathlineto{\pgfqpoint{2.757086in}{1.797162in}}%
\pgfpathlineto{\pgfqpoint{2.755075in}{1.781646in}}%
\pgfpathlineto{\pgfqpoint{2.738454in}{1.767308in}}%
\pgfpathlineto{\pgfqpoint{2.739441in}{1.749952in}}%
\pgfpathlineto{\pgfqpoint{2.754008in}{1.743872in}}%
\pgfpathlineto{\pgfqpoint{2.770441in}{1.765527in}}%
\pgfpathlineto{\pgfqpoint{2.771800in}{1.773381in}}%
\pgfpathlineto{\pgfqpoint{2.792295in}{1.786435in}}%
\pgfpathlineto{\pgfqpoint{2.805328in}{1.780388in}}%
\pgfpathlineto{\pgfqpoint{2.813507in}{1.766723in}}%
\pgfpathlineto{\pgfqpoint{2.826701in}{1.719275in}}%
\pgfpathlineto{\pgfqpoint{2.833693in}{1.704263in}}%
\pgfpathlineto{\pgfqpoint{2.831504in}{1.698058in}}%
\pgfpathlineto{\pgfqpoint{2.831717in}{1.676927in}}%
\pgfpathlineto{\pgfqpoint{2.818998in}{1.678955in}}%
\pgfpathlineto{\pgfqpoint{2.811814in}{1.663165in}}%
\pgfpathlineto{\pgfqpoint{2.810819in}{1.652430in}}%
\pgfpathlineto{\pgfqpoint{2.801211in}{1.645539in}}%
\pgfpathlineto{\pgfqpoint{2.799076in}{1.624604in}}%
\pgfpathlineto{\pgfqpoint{2.785105in}{1.598157in}}%
\pgfpathlineto{\pgfqpoint{2.708677in}{1.586775in}}%
\pgfpathlineto{\pgfqpoint{2.708226in}{1.591778in}}%
\pgfpathlineto{\pgfqpoint{2.657102in}{1.586389in}}%
\pgfpathlineto{\pgfqpoint{2.588031in}{1.579945in}}%
\pgfusepath{stroke}%
\end{pgfscope}%
\begin{pgfscope}%
\pgfpathrectangle{\pgfqpoint{0.100000in}{0.100000in}}{\pgfqpoint{3.608454in}{2.310000in}}%
\pgfusepath{clip}%
\pgfsetbuttcap%
\pgfsetroundjoin%
\pgfsetlinewidth{0.010037pt}%
\definecolor{currentstroke}{rgb}{1.000000,1.000000,1.000000}%
\pgfsetstrokecolor{currentstroke}%
\pgfsetdash{}{0pt}%
\pgfpathmoveto{\pgfqpoint{0.000000in}{0.000000in}}%
\pgfusepath{stroke}%
\end{pgfscope}%
\begin{pgfscope}%
\pgfpathrectangle{\pgfqpoint{0.100000in}{0.100000in}}{\pgfqpoint{3.608454in}{2.310000in}}%
\pgfusepath{clip}%
\pgfsetbuttcap%
\pgfsetroundjoin%
\pgfsetlinewidth{0.010037pt}%
\definecolor{currentstroke}{rgb}{1.000000,1.000000,1.000000}%
\pgfsetstrokecolor{currentstroke}%
\pgfsetdash{}{0pt}%
\pgfusepath{stroke}%
\end{pgfscope}%
\begin{pgfscope}%
\pgfpathrectangle{\pgfqpoint{0.100000in}{0.100000in}}{\pgfqpoint{3.608454in}{2.310000in}}%
\pgfusepath{clip}%
\pgfsetbuttcap%
\pgfsetroundjoin%
\pgfsetlinewidth{0.010037pt}%
\definecolor{currentstroke}{rgb}{1.000000,1.000000,1.000000}%
\pgfsetstrokecolor{currentstroke}%
\pgfsetdash{}{0pt}%
\pgfusepath{stroke}%
\end{pgfscope}%
\begin{pgfscope}%
\pgfpathrectangle{\pgfqpoint{0.100000in}{0.100000in}}{\pgfqpoint{3.608454in}{2.310000in}}%
\pgfusepath{clip}%
\pgfsetbuttcap%
\pgfsetroundjoin%
\pgfsetlinewidth{0.010037pt}%
\definecolor{currentstroke}{rgb}{1.000000,1.000000,1.000000}%
\pgfsetstrokecolor{currentstroke}%
\pgfsetdash{}{0pt}%
\pgfusepath{stroke}%
\end{pgfscope}%
\begin{pgfscope}%
\pgfpathrectangle{\pgfqpoint{0.100000in}{0.100000in}}{\pgfqpoint{3.608454in}{2.310000in}}%
\pgfusepath{clip}%
\pgfsetbuttcap%
\pgfsetroundjoin%
\pgfsetlinewidth{0.010037pt}%
\definecolor{currentstroke}{rgb}{1.000000,1.000000,1.000000}%
\pgfsetstrokecolor{currentstroke}%
\pgfsetdash{}{0pt}%
\pgfusepath{stroke}%
\end{pgfscope}%
\begin{pgfscope}%
\pgfpathrectangle{\pgfqpoint{0.100000in}{0.100000in}}{\pgfqpoint{3.608454in}{2.310000in}}%
\pgfusepath{clip}%
\pgfsetbuttcap%
\pgfsetroundjoin%
\pgfsetlinewidth{0.010037pt}%
\definecolor{currentstroke}{rgb}{1.000000,1.000000,1.000000}%
\pgfsetstrokecolor{currentstroke}%
\pgfsetdash{}{0pt}%
\pgfusepath{stroke}%
\end{pgfscope}%
\begin{pgfscope}%
\pgfpathrectangle{\pgfqpoint{0.100000in}{0.100000in}}{\pgfqpoint{3.608454in}{2.310000in}}%
\pgfusepath{clip}%
\pgfsetbuttcap%
\pgfsetroundjoin%
\pgfsetlinewidth{0.010037pt}%
\definecolor{currentstroke}{rgb}{1.000000,1.000000,1.000000}%
\pgfsetstrokecolor{currentstroke}%
\pgfsetdash{}{0pt}%
\pgfusepath{stroke}%
\end{pgfscope}%
\begin{pgfscope}%
\pgfpathrectangle{\pgfqpoint{0.100000in}{0.100000in}}{\pgfqpoint{3.608454in}{2.310000in}}%
\pgfusepath{clip}%
\pgfsetbuttcap%
\pgfsetroundjoin%
\pgfsetlinewidth{0.010037pt}%
\definecolor{currentstroke}{rgb}{1.000000,1.000000,1.000000}%
\pgfsetstrokecolor{currentstroke}%
\pgfsetdash{}{0pt}%
\pgfusepath{stroke}%
\end{pgfscope}%
\begin{pgfscope}%
\pgfpathrectangle{\pgfqpoint{0.100000in}{0.100000in}}{\pgfqpoint{3.608454in}{2.310000in}}%
\pgfusepath{clip}%
\pgfsetbuttcap%
\pgfsetroundjoin%
\pgfsetlinewidth{0.010037pt}%
\definecolor{currentstroke}{rgb}{1.000000,1.000000,1.000000}%
\pgfsetstrokecolor{currentstroke}%
\pgfsetdash{}{0pt}%
\pgfusepath{stroke}%
\end{pgfscope}%
\begin{pgfscope}%
\pgfpathrectangle{\pgfqpoint{0.100000in}{0.100000in}}{\pgfqpoint{3.608454in}{2.310000in}}%
\pgfusepath{clip}%
\pgfsetbuttcap%
\pgfsetroundjoin%
\pgfsetlinewidth{0.010037pt}%
\definecolor{currentstroke}{rgb}{1.000000,1.000000,1.000000}%
\pgfsetstrokecolor{currentstroke}%
\pgfsetdash{}{0pt}%
\pgfusepath{stroke}%
\end{pgfscope}%
\begin{pgfscope}%
\pgfpathrectangle{\pgfqpoint{0.100000in}{0.100000in}}{\pgfqpoint{3.608454in}{2.310000in}}%
\pgfusepath{clip}%
\pgfsetbuttcap%
\pgfsetroundjoin%
\pgfsetlinewidth{0.010037pt}%
\definecolor{currentstroke}{rgb}{1.000000,1.000000,1.000000}%
\pgfsetstrokecolor{currentstroke}%
\pgfsetdash{}{0pt}%
\pgfusepath{stroke}%
\end{pgfscope}%
\begin{pgfscope}%
\pgfpathrectangle{\pgfqpoint{0.100000in}{0.100000in}}{\pgfqpoint{3.608454in}{2.310000in}}%
\pgfusepath{clip}%
\pgfsetbuttcap%
\pgfsetroundjoin%
\pgfsetlinewidth{0.010037pt}%
\definecolor{currentstroke}{rgb}{1.000000,1.000000,1.000000}%
\pgfsetstrokecolor{currentstroke}%
\pgfsetdash{}{0pt}%
\pgfusepath{stroke}%
\end{pgfscope}%
\begin{pgfscope}%
\pgfpathrectangle{\pgfqpoint{0.100000in}{0.100000in}}{\pgfqpoint{3.608454in}{2.310000in}}%
\pgfusepath{clip}%
\pgfsetbuttcap%
\pgfsetroundjoin%
\pgfsetlinewidth{0.010037pt}%
\definecolor{currentstroke}{rgb}{1.000000,1.000000,1.000000}%
\pgfsetstrokecolor{currentstroke}%
\pgfsetdash{}{0pt}%
\pgfusepath{stroke}%
\end{pgfscope}%
\begin{pgfscope}%
\pgfpathrectangle{\pgfqpoint{0.100000in}{0.100000in}}{\pgfqpoint{3.608454in}{2.310000in}}%
\pgfusepath{clip}%
\pgfsetbuttcap%
\pgfsetroundjoin%
\pgfsetlinewidth{0.010037pt}%
\definecolor{currentstroke}{rgb}{1.000000,1.000000,1.000000}%
\pgfsetstrokecolor{currentstroke}%
\pgfsetdash{}{0pt}%
\pgfusepath{stroke}%
\end{pgfscope}%
\begin{pgfscope}%
\pgfpathrectangle{\pgfqpoint{0.100000in}{0.100000in}}{\pgfqpoint{3.608454in}{2.310000in}}%
\pgfusepath{clip}%
\pgfsetbuttcap%
\pgfsetroundjoin%
\pgfsetlinewidth{0.010037pt}%
\definecolor{currentstroke}{rgb}{1.000000,1.000000,1.000000}%
\pgfsetstrokecolor{currentstroke}%
\pgfsetdash{}{0pt}%
\pgfusepath{stroke}%
\end{pgfscope}%
\begin{pgfscope}%
\pgfpathrectangle{\pgfqpoint{0.100000in}{0.100000in}}{\pgfqpoint{3.608454in}{2.310000in}}%
\pgfusepath{clip}%
\pgfsetbuttcap%
\pgfsetroundjoin%
\pgfsetlinewidth{0.010037pt}%
\definecolor{currentstroke}{rgb}{1.000000,1.000000,1.000000}%
\pgfsetstrokecolor{currentstroke}%
\pgfsetdash{}{0pt}%
\pgfusepath{stroke}%
\end{pgfscope}%
\begin{pgfscope}%
\pgfpathrectangle{\pgfqpoint{0.100000in}{0.100000in}}{\pgfqpoint{3.608454in}{2.310000in}}%
\pgfusepath{clip}%
\pgfsetbuttcap%
\pgfsetroundjoin%
\pgfsetlinewidth{0.010037pt}%
\definecolor{currentstroke}{rgb}{1.000000,1.000000,1.000000}%
\pgfsetstrokecolor{currentstroke}%
\pgfsetdash{}{0pt}%
\pgfusepath{stroke}%
\end{pgfscope}%
\begin{pgfscope}%
\pgfpathrectangle{\pgfqpoint{0.100000in}{0.100000in}}{\pgfqpoint{3.608454in}{2.310000in}}%
\pgfusepath{clip}%
\pgfsetbuttcap%
\pgfsetroundjoin%
\pgfsetlinewidth{0.010037pt}%
\definecolor{currentstroke}{rgb}{1.000000,1.000000,1.000000}%
\pgfsetstrokecolor{currentstroke}%
\pgfsetdash{}{0pt}%
\pgfusepath{stroke}%
\end{pgfscope}%
\begin{pgfscope}%
\pgfpathrectangle{\pgfqpoint{0.100000in}{0.100000in}}{\pgfqpoint{3.608454in}{2.310000in}}%
\pgfusepath{clip}%
\pgfsetbuttcap%
\pgfsetroundjoin%
\pgfsetlinewidth{0.010037pt}%
\definecolor{currentstroke}{rgb}{1.000000,1.000000,1.000000}%
\pgfsetstrokecolor{currentstroke}%
\pgfsetdash{}{0pt}%
\pgfusepath{stroke}%
\end{pgfscope}%
\begin{pgfscope}%
\pgfpathrectangle{\pgfqpoint{0.100000in}{0.100000in}}{\pgfqpoint{3.608454in}{2.310000in}}%
\pgfusepath{clip}%
\pgfsetbuttcap%
\pgfsetroundjoin%
\pgfsetlinewidth{0.010037pt}%
\definecolor{currentstroke}{rgb}{1.000000,1.000000,1.000000}%
\pgfsetstrokecolor{currentstroke}%
\pgfsetdash{}{0pt}%
\pgfusepath{stroke}%
\end{pgfscope}%
\begin{pgfscope}%
\pgfpathrectangle{\pgfqpoint{0.100000in}{0.100000in}}{\pgfqpoint{3.608454in}{2.310000in}}%
\pgfusepath{clip}%
\pgfsetbuttcap%
\pgfsetroundjoin%
\pgfsetlinewidth{0.010037pt}%
\definecolor{currentstroke}{rgb}{1.000000,1.000000,1.000000}%
\pgfsetstrokecolor{currentstroke}%
\pgfsetdash{}{0pt}%
\pgfusepath{stroke}%
\end{pgfscope}%
\begin{pgfscope}%
\pgfpathrectangle{\pgfqpoint{0.100000in}{0.100000in}}{\pgfqpoint{3.608454in}{2.310000in}}%
\pgfusepath{clip}%
\pgfsetbuttcap%
\pgfsetroundjoin%
\pgfsetlinewidth{0.010037pt}%
\definecolor{currentstroke}{rgb}{1.000000,1.000000,1.000000}%
\pgfsetstrokecolor{currentstroke}%
\pgfsetdash{}{0pt}%
\pgfusepath{stroke}%
\end{pgfscope}%
\begin{pgfscope}%
\pgfpathrectangle{\pgfqpoint{0.100000in}{0.100000in}}{\pgfqpoint{3.608454in}{2.310000in}}%
\pgfusepath{clip}%
\pgfsetbuttcap%
\pgfsetroundjoin%
\pgfsetlinewidth{0.010037pt}%
\definecolor{currentstroke}{rgb}{1.000000,1.000000,1.000000}%
\pgfsetstrokecolor{currentstroke}%
\pgfsetdash{}{0pt}%
\pgfusepath{stroke}%
\end{pgfscope}%
\begin{pgfscope}%
\pgfpathrectangle{\pgfqpoint{0.100000in}{0.100000in}}{\pgfqpoint{3.608454in}{2.310000in}}%
\pgfusepath{clip}%
\pgfsetbuttcap%
\pgfsetroundjoin%
\pgfsetlinewidth{0.010037pt}%
\definecolor{currentstroke}{rgb}{1.000000,1.000000,1.000000}%
\pgfsetstrokecolor{currentstroke}%
\pgfsetdash{}{0pt}%
\pgfusepath{stroke}%
\end{pgfscope}%
\begin{pgfscope}%
\pgfpathrectangle{\pgfqpoint{0.100000in}{0.100000in}}{\pgfqpoint{3.608454in}{2.310000in}}%
\pgfusepath{clip}%
\pgfsetbuttcap%
\pgfsetroundjoin%
\pgfsetlinewidth{0.010037pt}%
\definecolor{currentstroke}{rgb}{1.000000,1.000000,1.000000}%
\pgfsetstrokecolor{currentstroke}%
\pgfsetdash{}{0pt}%
\pgfusepath{stroke}%
\end{pgfscope}%
\begin{pgfscope}%
\pgfpathrectangle{\pgfqpoint{0.100000in}{0.100000in}}{\pgfqpoint{3.608454in}{2.310000in}}%
\pgfusepath{clip}%
\pgfsetbuttcap%
\pgfsetroundjoin%
\pgfsetlinewidth{0.010037pt}%
\definecolor{currentstroke}{rgb}{1.000000,1.000000,1.000000}%
\pgfsetstrokecolor{currentstroke}%
\pgfsetdash{}{0pt}%
\pgfusepath{stroke}%
\end{pgfscope}%
\begin{pgfscope}%
\pgfpathrectangle{\pgfqpoint{0.100000in}{0.100000in}}{\pgfqpoint{3.608454in}{2.310000in}}%
\pgfusepath{clip}%
\pgfsetbuttcap%
\pgfsetroundjoin%
\pgfsetlinewidth{0.010037pt}%
\definecolor{currentstroke}{rgb}{1.000000,1.000000,1.000000}%
\pgfsetstrokecolor{currentstroke}%
\pgfsetdash{}{0pt}%
\pgfusepath{stroke}%
\end{pgfscope}%
\begin{pgfscope}%
\pgfpathrectangle{\pgfqpoint{0.100000in}{0.100000in}}{\pgfqpoint{3.608454in}{2.310000in}}%
\pgfusepath{clip}%
\pgfsetbuttcap%
\pgfsetroundjoin%
\pgfsetlinewidth{0.010037pt}%
\definecolor{currentstroke}{rgb}{1.000000,1.000000,1.000000}%
\pgfsetstrokecolor{currentstroke}%
\pgfsetdash{}{0pt}%
\pgfusepath{stroke}%
\end{pgfscope}%
\begin{pgfscope}%
\pgfpathrectangle{\pgfqpoint{0.100000in}{0.100000in}}{\pgfqpoint{3.608454in}{2.310000in}}%
\pgfusepath{clip}%
\pgfsetbuttcap%
\pgfsetmiterjoin%
\pgfsetlinewidth{0.000000pt}%
\definecolor{currentstroke}{rgb}{0.000000,0.000000,0.000000}%
\pgfsetstrokecolor{currentstroke}%
\pgfsetstrokeopacity{0.000000}%
\pgfsetdash{}{0pt}%
\pgfpathmoveto{\pgfqpoint{0.100000in}{0.100000in}}%
\pgfpathlineto{\pgfqpoint{3.708454in}{0.100000in}}%
\pgfpathlineto{\pgfqpoint{3.708454in}{2.410000in}}%
\pgfpathlineto{\pgfqpoint{0.100000in}{2.410000in}}%
\pgfpathlineto{\pgfqpoint{0.100000in}{0.100000in}}%
\pgfpathclose%
\pgfusepath{}%
\end{pgfscope}%
\begin{pgfscope}%
\pgfpathrectangle{\pgfqpoint{0.100000in}{0.100000in}}{\pgfqpoint{3.608454in}{2.310000in}}%
\pgfusepath{clip}%
\pgfsetbuttcap%
\pgfsetmiterjoin%
\definecolor{currentfill}{rgb}{0.000000,0.364706,0.817647}%
\pgfsetfillcolor{currentfill}%
\pgfsetlinewidth{0.000000pt}%
\definecolor{currentstroke}{rgb}{0.000000,0.000000,0.000000}%
\pgfsetstrokecolor{currentstroke}%
\pgfsetstrokeopacity{0.000000}%
\pgfsetdash{}{0pt}%
\pgfpathmoveto{\pgfqpoint{2.124622in}{1.510289in}}%
\pgfpathlineto{\pgfqpoint{2.097514in}{1.510367in}}%
\pgfpathlineto{\pgfqpoint{2.097619in}{1.537721in}}%
\pgfpathlineto{\pgfqpoint{2.077274in}{1.537859in}}%
\pgfpathlineto{\pgfqpoint{2.073777in}{1.545495in}}%
\pgfpathlineto{\pgfqpoint{2.073952in}{1.566162in}}%
\pgfpathlineto{\pgfqpoint{2.121832in}{1.566001in}}%
\pgfpathlineto{\pgfqpoint{2.121931in}{1.545408in}}%
\pgfpathlineto{\pgfqpoint{2.124501in}{1.545418in}}%
\pgfpathlineto{\pgfqpoint{2.124622in}{1.510289in}}%
\pgfpathclose%
\pgfusepath{fill}%
\end{pgfscope}%
\begin{pgfscope}%
\pgfpathrectangle{\pgfqpoint{0.100000in}{0.100000in}}{\pgfqpoint{3.608454in}{2.310000in}}%
\pgfusepath{clip}%
\pgfsetbuttcap%
\pgfsetmiterjoin%
\definecolor{currentfill}{rgb}{0.000000,0.850980,0.574510}%
\pgfsetfillcolor{currentfill}%
\pgfsetlinewidth{0.000000pt}%
\definecolor{currentstroke}{rgb}{0.000000,0.000000,0.000000}%
\pgfsetstrokecolor{currentstroke}%
\pgfsetstrokeopacity{0.000000}%
\pgfsetdash{}{0pt}%
\pgfpathmoveto{\pgfqpoint{2.724860in}{1.203323in}}%
\pgfpathlineto{\pgfqpoint{2.731198in}{1.201857in}}%
\pgfpathlineto{\pgfqpoint{2.747524in}{1.217198in}}%
\pgfpathlineto{\pgfqpoint{2.742647in}{1.220616in}}%
\pgfpathlineto{\pgfqpoint{2.742337in}{1.226077in}}%
\pgfpathlineto{\pgfqpoint{2.752556in}{1.236924in}}%
\pgfpathlineto{\pgfqpoint{2.761764in}{1.247341in}}%
\pgfpathlineto{\pgfqpoint{2.765354in}{1.245580in}}%
\pgfpathlineto{\pgfqpoint{2.780361in}{1.233265in}}%
\pgfpathlineto{\pgfqpoint{2.778898in}{1.226398in}}%
\pgfpathlineto{\pgfqpoint{2.781636in}{1.214245in}}%
\pgfpathlineto{\pgfqpoint{2.779295in}{1.206052in}}%
\pgfpathlineto{\pgfqpoint{2.783865in}{1.195392in}}%
\pgfpathlineto{\pgfqpoint{2.789452in}{1.189777in}}%
\pgfpathlineto{\pgfqpoint{2.788196in}{1.183891in}}%
\pgfpathlineto{\pgfqpoint{2.783231in}{1.181649in}}%
\pgfpathlineto{\pgfqpoint{2.785832in}{1.164847in}}%
\pgfpathlineto{\pgfqpoint{2.781495in}{1.155039in}}%
\pgfpathlineto{\pgfqpoint{2.771140in}{1.159342in}}%
\pgfpathlineto{\pgfqpoint{2.761563in}{1.168904in}}%
\pgfpathlineto{\pgfqpoint{2.763981in}{1.171120in}}%
\pgfpathlineto{\pgfqpoint{2.758225in}{1.180770in}}%
\pgfpathlineto{\pgfqpoint{2.754147in}{1.186574in}}%
\pgfpathlineto{\pgfqpoint{2.722959in}{1.184982in}}%
\pgfpathlineto{\pgfqpoint{2.724860in}{1.203323in}}%
\pgfpathclose%
\pgfusepath{fill}%
\end{pgfscope}%
\begin{pgfscope}%
\pgfpathrectangle{\pgfqpoint{0.100000in}{0.100000in}}{\pgfqpoint{3.608454in}{2.310000in}}%
\pgfusepath{clip}%
\pgfsetbuttcap%
\pgfsetmiterjoin%
\definecolor{currentfill}{rgb}{0.000000,0.635294,0.682353}%
\pgfsetfillcolor{currentfill}%
\pgfsetlinewidth{0.000000pt}%
\definecolor{currentstroke}{rgb}{0.000000,0.000000,0.000000}%
\pgfsetstrokecolor{currentstroke}%
\pgfsetstrokeopacity{0.000000}%
\pgfsetdash{}{0pt}%
\pgfpathmoveto{\pgfqpoint{1.745610in}{0.979809in}}%
\pgfpathlineto{\pgfqpoint{1.743217in}{0.941925in}}%
\pgfpathlineto{\pgfqpoint{1.708835in}{0.943953in}}%
\pgfpathlineto{\pgfqpoint{1.711338in}{0.981859in}}%
\pgfpathlineto{\pgfqpoint{1.745610in}{0.979809in}}%
\pgfpathclose%
\pgfusepath{fill}%
\end{pgfscope}%
\begin{pgfscope}%
\pgfpathrectangle{\pgfqpoint{0.100000in}{0.100000in}}{\pgfqpoint{3.608454in}{2.310000in}}%
\pgfusepath{clip}%
\pgfsetbuttcap%
\pgfsetmiterjoin%
\definecolor{currentfill}{rgb}{0.000000,0.713725,0.643137}%
\pgfsetfillcolor{currentfill}%
\pgfsetlinewidth{0.000000pt}%
\definecolor{currentstroke}{rgb}{0.000000,0.000000,0.000000}%
\pgfsetstrokecolor{currentstroke}%
\pgfsetstrokeopacity{0.000000}%
\pgfsetdash{}{0pt}%
\pgfpathmoveto{\pgfqpoint{1.670049in}{1.298787in}}%
\pgfpathlineto{\pgfqpoint{1.670021in}{1.298287in}}%
\pgfpathlineto{\pgfqpoint{1.666954in}{1.256981in}}%
\pgfpathlineto{\pgfqpoint{1.666412in}{1.249548in}}%
\pgfpathlineto{\pgfqpoint{1.601877in}{1.254860in}}%
\pgfpathlineto{\pgfqpoint{1.540663in}{1.260697in}}%
\pgfpathlineto{\pgfqpoint{1.541303in}{1.267794in}}%
\pgfpathlineto{\pgfqpoint{1.506183in}{1.284643in}}%
\pgfpathlineto{\pgfqpoint{1.483235in}{1.285604in}}%
\pgfpathlineto{\pgfqpoint{1.481289in}{1.288540in}}%
\pgfpathlineto{\pgfqpoint{1.484320in}{1.315446in}}%
\pgfpathlineto{\pgfqpoint{1.491038in}{1.314728in}}%
\pgfpathlineto{\pgfqpoint{1.493221in}{1.335267in}}%
\pgfpathlineto{\pgfqpoint{1.547834in}{1.329763in}}%
\pgfpathlineto{\pgfqpoint{1.581597in}{1.326016in}}%
\pgfpathlineto{\pgfqpoint{1.579920in}{1.306214in}}%
\pgfpathlineto{\pgfqpoint{1.626877in}{1.302163in}}%
\pgfpathlineto{\pgfqpoint{1.670049in}{1.298787in}}%
\pgfpathclose%
\pgfusepath{fill}%
\end{pgfscope}%
\begin{pgfscope}%
\pgfpathrectangle{\pgfqpoint{0.100000in}{0.100000in}}{\pgfqpoint{3.608454in}{2.310000in}}%
\pgfusepath{clip}%
\pgfsetbuttcap%
\pgfsetmiterjoin%
\definecolor{currentfill}{rgb}{0.000000,0.431373,0.784314}%
\pgfsetfillcolor{currentfill}%
\pgfsetlinewidth{0.000000pt}%
\definecolor{currentstroke}{rgb}{0.000000,0.000000,0.000000}%
\pgfsetstrokecolor{currentstroke}%
\pgfsetstrokeopacity{0.000000}%
\pgfsetdash{}{0pt}%
\pgfpathmoveto{\pgfqpoint{2.103336in}{1.338237in}}%
\pgfpathlineto{\pgfqpoint{2.096838in}{1.338167in}}%
\pgfpathlineto{\pgfqpoint{2.095756in}{1.349849in}}%
\pgfpathlineto{\pgfqpoint{2.095818in}{1.372775in}}%
\pgfpathlineto{\pgfqpoint{2.071958in}{1.372885in}}%
\pgfpathlineto{\pgfqpoint{2.072459in}{1.391367in}}%
\pgfpathlineto{\pgfqpoint{2.086114in}{1.391298in}}%
\pgfpathlineto{\pgfqpoint{2.099783in}{1.391231in}}%
\pgfpathlineto{\pgfqpoint{2.103446in}{1.388933in}}%
\pgfpathlineto{\pgfqpoint{2.100521in}{1.381808in}}%
\pgfpathlineto{\pgfqpoint{2.107556in}{1.374791in}}%
\pgfpathlineto{\pgfqpoint{2.116197in}{1.360877in}}%
\pgfpathlineto{\pgfqpoint{2.112893in}{1.355625in}}%
\pgfpathlineto{\pgfqpoint{2.112374in}{1.338940in}}%
\pgfpathlineto{\pgfqpoint{2.103336in}{1.338237in}}%
\pgfpathclose%
\pgfusepath{fill}%
\end{pgfscope}%
\begin{pgfscope}%
\pgfpathrectangle{\pgfqpoint{0.100000in}{0.100000in}}{\pgfqpoint{3.608454in}{2.310000in}}%
\pgfusepath{clip}%
\pgfsetbuttcap%
\pgfsetmiterjoin%
\definecolor{currentfill}{rgb}{0.000000,0.517647,0.741176}%
\pgfsetfillcolor{currentfill}%
\pgfsetlinewidth{0.000000pt}%
\definecolor{currentstroke}{rgb}{0.000000,0.000000,0.000000}%
\pgfsetstrokecolor{currentstroke}%
\pgfsetstrokeopacity{0.000000}%
\pgfsetdash{}{0pt}%
\pgfpathmoveto{\pgfqpoint{2.760784in}{0.972182in}}%
\pgfpathlineto{\pgfqpoint{2.738788in}{0.969584in}}%
\pgfpathlineto{\pgfqpoint{2.737463in}{0.974424in}}%
\pgfpathlineto{\pgfqpoint{2.728983in}{0.972881in}}%
\pgfpathlineto{\pgfqpoint{2.715165in}{0.972447in}}%
\pgfpathlineto{\pgfqpoint{2.706101in}{0.993510in}}%
\pgfpathlineto{\pgfqpoint{2.710524in}{0.995150in}}%
\pgfpathlineto{\pgfqpoint{2.718615in}{1.004553in}}%
\pgfpathlineto{\pgfqpoint{2.721545in}{1.014250in}}%
\pgfpathlineto{\pgfqpoint{2.725316in}{1.017812in}}%
\pgfpathlineto{\pgfqpoint{2.723904in}{1.022807in}}%
\pgfpathlineto{\pgfqpoint{2.754206in}{1.025936in}}%
\pgfpathlineto{\pgfqpoint{2.753682in}{1.036469in}}%
\pgfpathlineto{\pgfqpoint{2.746037in}{1.035657in}}%
\pgfpathlineto{\pgfqpoint{2.747587in}{1.043751in}}%
\pgfpathlineto{\pgfqpoint{2.752122in}{1.044231in}}%
\pgfpathlineto{\pgfqpoint{2.756131in}{1.058082in}}%
\pgfpathlineto{\pgfqpoint{2.779057in}{1.060719in}}%
\pgfpathlineto{\pgfqpoint{2.780160in}{1.050381in}}%
\pgfpathlineto{\pgfqpoint{2.782890in}{1.048115in}}%
\pgfpathlineto{\pgfqpoint{2.779095in}{1.040030in}}%
\pgfpathlineto{\pgfqpoint{2.780892in}{1.025876in}}%
\pgfpathlineto{\pgfqpoint{2.784817in}{0.988730in}}%
\pgfpathlineto{\pgfqpoint{2.767572in}{0.987144in}}%
\pgfpathlineto{\pgfqpoint{2.765064in}{0.976399in}}%
\pgfpathlineto{\pgfqpoint{2.760784in}{0.972182in}}%
\pgfpathclose%
\pgfusepath{fill}%
\end{pgfscope}%
\begin{pgfscope}%
\pgfpathrectangle{\pgfqpoint{0.100000in}{0.100000in}}{\pgfqpoint{3.608454in}{2.310000in}}%
\pgfusepath{clip}%
\pgfsetbuttcap%
\pgfsetmiterjoin%
\definecolor{currentfill}{rgb}{0.000000,0.431373,0.784314}%
\pgfsetfillcolor{currentfill}%
\pgfsetlinewidth{0.000000pt}%
\definecolor{currentstroke}{rgb}{0.000000,0.000000,0.000000}%
\pgfsetstrokecolor{currentstroke}%
\pgfsetstrokeopacity{0.000000}%
\pgfsetdash{}{0pt}%
\pgfpathmoveto{\pgfqpoint{1.716695in}{1.867335in}}%
\pgfpathlineto{\pgfqpoint{1.714041in}{1.832923in}}%
\pgfpathlineto{\pgfqpoint{1.660454in}{1.837408in}}%
\pgfpathlineto{\pgfqpoint{1.661648in}{1.851140in}}%
\pgfpathlineto{\pgfqpoint{1.667609in}{1.909112in}}%
\pgfpathlineto{\pgfqpoint{1.664670in}{1.909384in}}%
\pgfpathlineto{\pgfqpoint{1.667012in}{1.936010in}}%
\pgfpathlineto{\pgfqpoint{1.670909in}{1.935733in}}%
\pgfpathlineto{\pgfqpoint{1.673140in}{1.963430in}}%
\pgfpathlineto{\pgfqpoint{1.718524in}{1.959750in}}%
\pgfpathlineto{\pgfqpoint{1.721015in}{1.959541in}}%
\pgfpathlineto{\pgfqpoint{1.719106in}{1.931926in}}%
\pgfpathlineto{\pgfqpoint{1.721787in}{1.931717in}}%
\pgfpathlineto{\pgfqpoint{1.720336in}{1.913435in}}%
\pgfpathlineto{\pgfqpoint{1.716695in}{1.867335in}}%
\pgfpathclose%
\pgfusepath{fill}%
\end{pgfscope}%
\begin{pgfscope}%
\pgfpathrectangle{\pgfqpoint{0.100000in}{0.100000in}}{\pgfqpoint{3.608454in}{2.310000in}}%
\pgfusepath{clip}%
\pgfsetbuttcap%
\pgfsetmiterjoin%
\definecolor{currentfill}{rgb}{0.000000,0.647059,0.676471}%
\pgfsetfillcolor{currentfill}%
\pgfsetlinewidth{0.000000pt}%
\definecolor{currentstroke}{rgb}{0.000000,0.000000,0.000000}%
\pgfsetstrokecolor{currentstroke}%
\pgfsetstrokeopacity{0.000000}%
\pgfsetdash{}{0pt}%
\pgfpathmoveto{\pgfqpoint{3.298349in}{1.383037in}}%
\pgfpathlineto{\pgfqpoint{3.304739in}{1.388327in}}%
\pgfpathlineto{\pgfqpoint{3.327417in}{1.395926in}}%
\pgfpathlineto{\pgfqpoint{3.322504in}{1.381252in}}%
\pgfpathlineto{\pgfqpoint{3.318671in}{1.381735in}}%
\pgfpathlineto{\pgfqpoint{3.314476in}{1.374248in}}%
\pgfpathlineto{\pgfqpoint{3.312084in}{1.358499in}}%
\pgfpathlineto{\pgfqpoint{3.313688in}{1.355781in}}%
\pgfpathlineto{\pgfqpoint{3.309070in}{1.336218in}}%
\pgfpathlineto{\pgfqpoint{3.307238in}{1.324528in}}%
\pgfpathlineto{\pgfqpoint{3.302528in}{1.316927in}}%
\pgfpathlineto{\pgfqpoint{3.298046in}{1.315923in}}%
\pgfpathlineto{\pgfqpoint{3.292330in}{1.326062in}}%
\pgfpathlineto{\pgfqpoint{3.292495in}{1.350793in}}%
\pgfpathlineto{\pgfqpoint{3.295985in}{1.361073in}}%
\pgfpathlineto{\pgfqpoint{3.296227in}{1.370099in}}%
\pgfpathlineto{\pgfqpoint{3.301816in}{1.371758in}}%
\pgfpathlineto{\pgfqpoint{3.303055in}{1.378748in}}%
\pgfpathlineto{\pgfqpoint{3.298349in}{1.383037in}}%
\pgfpathclose%
\pgfusepath{fill}%
\end{pgfscope}%
\begin{pgfscope}%
\pgfpathrectangle{\pgfqpoint{0.100000in}{0.100000in}}{\pgfqpoint{3.608454in}{2.310000in}}%
\pgfusepath{clip}%
\pgfsetbuttcap%
\pgfsetmiterjoin%
\definecolor{currentfill}{rgb}{0.000000,0.603922,0.698039}%
\pgfsetfillcolor{currentfill}%
\pgfsetlinewidth{0.000000pt}%
\definecolor{currentstroke}{rgb}{0.000000,0.000000,0.000000}%
\pgfsetstrokecolor{currentstroke}%
\pgfsetstrokeopacity{0.000000}%
\pgfsetdash{}{0pt}%
\pgfpathmoveto{\pgfqpoint{2.488486in}{1.795943in}}%
\pgfpathlineto{\pgfqpoint{2.468096in}{1.794922in}}%
\pgfpathlineto{\pgfqpoint{2.467658in}{1.801883in}}%
\pgfpathlineto{\pgfqpoint{2.432784in}{1.799789in}}%
\pgfpathlineto{\pgfqpoint{2.430551in}{1.834455in}}%
\pgfpathlineto{\pgfqpoint{2.419246in}{1.833798in}}%
\pgfpathlineto{\pgfqpoint{2.417621in}{1.861430in}}%
\pgfpathlineto{\pgfqpoint{2.445567in}{1.863054in}}%
\pgfpathlineto{\pgfqpoint{2.445861in}{1.856178in}}%
\pgfpathlineto{\pgfqpoint{2.459694in}{1.857127in}}%
\pgfpathlineto{\pgfqpoint{2.460529in}{1.843339in}}%
\pgfpathlineto{\pgfqpoint{2.463191in}{1.836540in}}%
\pgfpathlineto{\pgfqpoint{2.472001in}{1.837128in}}%
\pgfpathlineto{\pgfqpoint{2.473187in}{1.816374in}}%
\pgfpathlineto{\pgfqpoint{2.486722in}{1.817142in}}%
\pgfpathlineto{\pgfqpoint{2.488486in}{1.795943in}}%
\pgfpathclose%
\pgfusepath{fill}%
\end{pgfscope}%
\begin{pgfscope}%
\pgfpathrectangle{\pgfqpoint{0.100000in}{0.100000in}}{\pgfqpoint{3.608454in}{2.310000in}}%
\pgfusepath{clip}%
\pgfsetbuttcap%
\pgfsetmiterjoin%
\definecolor{currentfill}{rgb}{0.000000,0.556863,0.721569}%
\pgfsetfillcolor{currentfill}%
\pgfsetlinewidth{0.000000pt}%
\definecolor{currentstroke}{rgb}{0.000000,0.000000,0.000000}%
\pgfsetstrokecolor{currentstroke}%
\pgfsetstrokeopacity{0.000000}%
\pgfsetdash{}{0pt}%
\pgfpathmoveto{\pgfqpoint{3.384285in}{1.735025in}}%
\pgfpathlineto{\pgfqpoint{3.359316in}{1.729735in}}%
\pgfpathlineto{\pgfqpoint{3.358769in}{1.729605in}}%
\pgfpathlineto{\pgfqpoint{3.357447in}{1.732265in}}%
\pgfpathlineto{\pgfqpoint{3.358955in}{1.786375in}}%
\pgfpathlineto{\pgfqpoint{3.356483in}{1.793242in}}%
\pgfpathlineto{\pgfqpoint{3.348859in}{1.830296in}}%
\pgfpathlineto{\pgfqpoint{3.370781in}{1.834556in}}%
\pgfpathlineto{\pgfqpoint{3.374364in}{1.831752in}}%
\pgfpathlineto{\pgfqpoint{3.374319in}{1.820118in}}%
\pgfpathlineto{\pgfqpoint{3.367008in}{1.818791in}}%
\pgfpathlineto{\pgfqpoint{3.369972in}{1.804097in}}%
\pgfpathlineto{\pgfqpoint{3.374916in}{1.805051in}}%
\pgfpathlineto{\pgfqpoint{3.377987in}{1.790511in}}%
\pgfpathlineto{\pgfqpoint{3.373418in}{1.786388in}}%
\pgfpathlineto{\pgfqpoint{3.378669in}{1.782710in}}%
\pgfpathlineto{\pgfqpoint{3.380209in}{1.761732in}}%
\pgfpathlineto{\pgfqpoint{3.384285in}{1.735025in}}%
\pgfpathclose%
\pgfusepath{fill}%
\end{pgfscope}%
\begin{pgfscope}%
\pgfpathrectangle{\pgfqpoint{0.100000in}{0.100000in}}{\pgfqpoint{3.608454in}{2.310000in}}%
\pgfusepath{clip}%
\pgfsetbuttcap%
\pgfsetmiterjoin%
\definecolor{currentfill}{rgb}{0.000000,0.752941,0.623529}%
\pgfsetfillcolor{currentfill}%
\pgfsetlinewidth{0.000000pt}%
\definecolor{currentstroke}{rgb}{0.000000,0.000000,0.000000}%
\pgfsetstrokecolor{currentstroke}%
\pgfsetstrokeopacity{0.000000}%
\pgfsetdash{}{0pt}%
\pgfpathmoveto{\pgfqpoint{2.895571in}{1.311319in}}%
\pgfpathlineto{\pgfqpoint{2.890443in}{1.304155in}}%
\pgfpathlineto{\pgfqpoint{2.884060in}{1.308388in}}%
\pgfpathlineto{\pgfqpoint{2.885378in}{1.313095in}}%
\pgfpathlineto{\pgfqpoint{2.873043in}{1.323790in}}%
\pgfpathlineto{\pgfqpoint{2.874015in}{1.332782in}}%
\pgfpathlineto{\pgfqpoint{2.861780in}{1.330306in}}%
\pgfpathlineto{\pgfqpoint{2.855274in}{1.324556in}}%
\pgfpathlineto{\pgfqpoint{2.857831in}{1.319548in}}%
\pgfpathlineto{\pgfqpoint{2.850698in}{1.310252in}}%
\pgfpathlineto{\pgfqpoint{2.840076in}{1.309133in}}%
\pgfpathlineto{\pgfqpoint{2.834982in}{1.312343in}}%
\pgfpathlineto{\pgfqpoint{2.824999in}{1.309058in}}%
\pgfpathlineto{\pgfqpoint{2.811045in}{1.319515in}}%
\pgfpathlineto{\pgfqpoint{2.818045in}{1.327218in}}%
\pgfpathlineto{\pgfqpoint{2.817804in}{1.333230in}}%
\pgfpathlineto{\pgfqpoint{2.822416in}{1.337584in}}%
\pgfpathlineto{\pgfqpoint{2.826872in}{1.332234in}}%
\pgfpathlineto{\pgfqpoint{2.834430in}{1.334621in}}%
\pgfpathlineto{\pgfqpoint{2.833231in}{1.340587in}}%
\pgfpathlineto{\pgfqpoint{2.835681in}{1.347268in}}%
\pgfpathlineto{\pgfqpoint{2.841606in}{1.355934in}}%
\pgfpathlineto{\pgfqpoint{2.842704in}{1.366486in}}%
\pgfpathlineto{\pgfqpoint{2.850733in}{1.370027in}}%
\pgfpathlineto{\pgfqpoint{2.854867in}{1.358657in}}%
\pgfpathlineto{\pgfqpoint{2.861822in}{1.358808in}}%
\pgfpathlineto{\pgfqpoint{2.864952in}{1.371264in}}%
\pgfpathlineto{\pgfqpoint{2.864652in}{1.379289in}}%
\pgfpathlineto{\pgfqpoint{2.869245in}{1.379552in}}%
\pgfpathlineto{\pgfqpoint{2.869516in}{1.374371in}}%
\pgfpathlineto{\pgfqpoint{2.875638in}{1.374674in}}%
\pgfpathlineto{\pgfqpoint{2.877293in}{1.367722in}}%
\pgfpathlineto{\pgfqpoint{2.884573in}{1.368240in}}%
\pgfpathlineto{\pgfqpoint{2.885129in}{1.361064in}}%
\pgfpathlineto{\pgfqpoint{2.893788in}{1.362802in}}%
\pgfpathlineto{\pgfqpoint{2.905122in}{1.355064in}}%
\pgfpathlineto{\pgfqpoint{2.906738in}{1.347296in}}%
\pgfpathlineto{\pgfqpoint{2.898128in}{1.341949in}}%
\pgfpathlineto{\pgfqpoint{2.894444in}{1.332079in}}%
\pgfpathlineto{\pgfqpoint{2.894629in}{1.327193in}}%
\pgfpathlineto{\pgfqpoint{2.900793in}{1.318863in}}%
\pgfpathlineto{\pgfqpoint{2.895571in}{1.311319in}}%
\pgfpathclose%
\pgfusepath{fill}%
\end{pgfscope}%
\begin{pgfscope}%
\pgfpathrectangle{\pgfqpoint{0.100000in}{0.100000in}}{\pgfqpoint{3.608454in}{2.310000in}}%
\pgfusepath{clip}%
\pgfsetbuttcap%
\pgfsetmiterjoin%
\definecolor{currentfill}{rgb}{0.000000,0.623529,0.688235}%
\pgfsetfillcolor{currentfill}%
\pgfsetlinewidth{0.000000pt}%
\definecolor{currentstroke}{rgb}{0.000000,0.000000,0.000000}%
\pgfsetstrokecolor{currentstroke}%
\pgfsetstrokeopacity{0.000000}%
\pgfsetdash{}{0pt}%
\pgfpathmoveto{\pgfqpoint{1.413364in}{1.067432in}}%
\pgfpathlineto{\pgfqpoint{1.379252in}{1.071453in}}%
\pgfpathlineto{\pgfqpoint{1.385971in}{1.127718in}}%
\pgfpathlineto{\pgfqpoint{1.377042in}{1.134592in}}%
\pgfpathlineto{\pgfqpoint{1.378347in}{1.144995in}}%
\pgfpathlineto{\pgfqpoint{1.388440in}{1.147038in}}%
\pgfpathlineto{\pgfqpoint{1.347890in}{1.152024in}}%
\pgfpathlineto{\pgfqpoint{1.350154in}{1.169097in}}%
\pgfpathlineto{\pgfqpoint{1.303419in}{1.175695in}}%
\pgfpathlineto{\pgfqpoint{1.310197in}{1.221478in}}%
\pgfpathlineto{\pgfqpoint{1.322223in}{1.228643in}}%
\pgfpathlineto{\pgfqpoint{1.324687in}{1.235061in}}%
\pgfpathlineto{\pgfqpoint{1.383780in}{1.226748in}}%
\pgfpathlineto{\pgfqpoint{1.462544in}{1.217413in}}%
\pgfpathlineto{\pgfqpoint{1.462472in}{1.211222in}}%
\pgfpathlineto{\pgfqpoint{1.458207in}{1.195789in}}%
\pgfpathlineto{\pgfqpoint{1.450272in}{1.192914in}}%
\pgfpathlineto{\pgfqpoint{1.450548in}{1.172051in}}%
\pgfpathlineto{\pgfqpoint{1.448410in}{1.160662in}}%
\pgfpathlineto{\pgfqpoint{1.445381in}{1.160260in}}%
\pgfpathlineto{\pgfqpoint{1.442215in}{1.150974in}}%
\pgfpathlineto{\pgfqpoint{1.431076in}{1.139578in}}%
\pgfpathlineto{\pgfqpoint{1.421911in}{1.141159in}}%
\pgfpathlineto{\pgfqpoint{1.413364in}{1.067432in}}%
\pgfpathclose%
\pgfusepath{fill}%
\end{pgfscope}%
\begin{pgfscope}%
\pgfpathrectangle{\pgfqpoint{0.100000in}{0.100000in}}{\pgfqpoint{3.608454in}{2.310000in}}%
\pgfusepath{clip}%
\pgfsetbuttcap%
\pgfsetmiterjoin%
\definecolor{currentfill}{rgb}{0.000000,0.847059,0.576471}%
\pgfsetfillcolor{currentfill}%
\pgfsetlinewidth{0.000000pt}%
\definecolor{currentstroke}{rgb}{0.000000,0.000000,0.000000}%
\pgfsetstrokecolor{currentstroke}%
\pgfsetstrokeopacity{0.000000}%
\pgfsetdash{}{0pt}%
\pgfpathmoveto{\pgfqpoint{2.407413in}{1.214350in}}%
\pgfpathlineto{\pgfqpoint{2.386546in}{1.213497in}}%
\pgfpathlineto{\pgfqpoint{2.386664in}{1.209942in}}%
\pgfpathlineto{\pgfqpoint{2.375052in}{1.209261in}}%
\pgfpathlineto{\pgfqpoint{2.372317in}{1.217256in}}%
\pgfpathlineto{\pgfqpoint{2.371740in}{1.235384in}}%
\pgfpathlineto{\pgfqpoint{2.358074in}{1.234831in}}%
\pgfpathlineto{\pgfqpoint{2.347612in}{1.233131in}}%
\pgfpathlineto{\pgfqpoint{2.347355in}{1.241765in}}%
\pgfpathlineto{\pgfqpoint{2.350544in}{1.245278in}}%
\pgfpathlineto{\pgfqpoint{2.349233in}{1.281942in}}%
\pgfpathlineto{\pgfqpoint{2.368816in}{1.282820in}}%
\pgfpathlineto{\pgfqpoint{2.381801in}{1.267561in}}%
\pgfpathlineto{\pgfqpoint{2.392056in}{1.271141in}}%
\pgfpathlineto{\pgfqpoint{2.397241in}{1.275533in}}%
\pgfpathlineto{\pgfqpoint{2.404880in}{1.275417in}}%
\pgfpathlineto{\pgfqpoint{2.409704in}{1.273645in}}%
\pgfpathlineto{\pgfqpoint{2.418250in}{1.266369in}}%
\pgfpathlineto{\pgfqpoint{2.421525in}{1.260176in}}%
\pgfpathlineto{\pgfqpoint{2.428315in}{1.262231in}}%
\pgfpathlineto{\pgfqpoint{2.434982in}{1.258000in}}%
\pgfpathlineto{\pgfqpoint{2.445368in}{1.248025in}}%
\pgfpathlineto{\pgfqpoint{2.449790in}{1.246442in}}%
\pgfpathlineto{\pgfqpoint{2.452467in}{1.239392in}}%
\pgfpathlineto{\pgfqpoint{2.449905in}{1.236650in}}%
\pgfpathlineto{\pgfqpoint{2.444353in}{1.238826in}}%
\pgfpathlineto{\pgfqpoint{2.410673in}{1.236890in}}%
\pgfpathlineto{\pgfqpoint{2.411944in}{1.214363in}}%
\pgfpathlineto{\pgfqpoint{2.407413in}{1.214350in}}%
\pgfpathclose%
\pgfusepath{fill}%
\end{pgfscope}%
\begin{pgfscope}%
\pgfpathrectangle{\pgfqpoint{0.100000in}{0.100000in}}{\pgfqpoint{3.608454in}{2.310000in}}%
\pgfusepath{clip}%
\pgfsetbuttcap%
\pgfsetmiterjoin%
\definecolor{currentfill}{rgb}{0.000000,0.647059,0.676471}%
\pgfsetfillcolor{currentfill}%
\pgfsetlinewidth{0.000000pt}%
\definecolor{currentstroke}{rgb}{0.000000,0.000000,0.000000}%
\pgfsetstrokecolor{currentstroke}%
\pgfsetstrokeopacity{0.000000}%
\pgfsetdash{}{0pt}%
\pgfpathmoveto{\pgfqpoint{2.751930in}{0.830708in}}%
\pgfpathlineto{\pgfqpoint{2.751810in}{0.820152in}}%
\pgfpathlineto{\pgfqpoint{2.741287in}{0.817640in}}%
\pgfpathlineto{\pgfqpoint{2.737520in}{0.807875in}}%
\pgfpathlineto{\pgfqpoint{2.728545in}{0.806963in}}%
\pgfpathlineto{\pgfqpoint{2.727912in}{0.813855in}}%
\pgfpathlineto{\pgfqpoint{2.721575in}{0.813214in}}%
\pgfpathlineto{\pgfqpoint{2.720311in}{0.819497in}}%
\pgfpathlineto{\pgfqpoint{2.713352in}{0.819097in}}%
\pgfpathlineto{\pgfqpoint{2.711685in}{0.841617in}}%
\pgfpathlineto{\pgfqpoint{2.708963in}{0.847375in}}%
\pgfpathlineto{\pgfqpoint{2.717264in}{0.854842in}}%
\pgfpathlineto{\pgfqpoint{2.723195in}{0.855536in}}%
\pgfpathlineto{\pgfqpoint{2.722392in}{0.862394in}}%
\pgfpathlineto{\pgfqpoint{2.729130in}{0.864274in}}%
\pgfpathlineto{\pgfqpoint{2.728245in}{0.872298in}}%
\pgfpathlineto{\pgfqpoint{2.734886in}{0.875470in}}%
\pgfpathlineto{\pgfqpoint{2.755362in}{0.877808in}}%
\pgfpathlineto{\pgfqpoint{2.765201in}{0.880136in}}%
\pgfpathlineto{\pgfqpoint{2.768066in}{0.872209in}}%
\pgfpathlineto{\pgfqpoint{2.776169in}{0.862453in}}%
\pgfpathlineto{\pgfqpoint{2.772613in}{0.859054in}}%
\pgfpathlineto{\pgfqpoint{2.754330in}{0.856732in}}%
\pgfpathlineto{\pgfqpoint{2.755065in}{0.852187in}}%
\pgfpathlineto{\pgfqpoint{2.748134in}{0.851370in}}%
\pgfpathlineto{\pgfqpoint{2.749632in}{0.837509in}}%
\pgfpathlineto{\pgfqpoint{2.751930in}{0.830708in}}%
\pgfpathclose%
\pgfusepath{fill}%
\end{pgfscope}%
\begin{pgfscope}%
\pgfpathrectangle{\pgfqpoint{0.100000in}{0.100000in}}{\pgfqpoint{3.608454in}{2.310000in}}%
\pgfusepath{clip}%
\pgfsetbuttcap%
\pgfsetmiterjoin%
\definecolor{currentfill}{rgb}{0.000000,0.474510,0.762745}%
\pgfsetfillcolor{currentfill}%
\pgfsetlinewidth{0.000000pt}%
\definecolor{currentstroke}{rgb}{0.000000,0.000000,0.000000}%
\pgfsetstrokecolor{currentstroke}%
\pgfsetstrokeopacity{0.000000}%
\pgfsetdash{}{0pt}%
\pgfpathmoveto{\pgfqpoint{0.975129in}{1.855275in}}%
\pgfpathlineto{\pgfqpoint{0.973897in}{1.850229in}}%
\pgfpathlineto{\pgfqpoint{0.967810in}{1.841171in}}%
\pgfpathlineto{\pgfqpoint{0.968859in}{1.837111in}}%
\pgfpathlineto{\pgfqpoint{0.963091in}{1.830834in}}%
\pgfpathlineto{\pgfqpoint{0.950896in}{1.777540in}}%
\pgfpathlineto{\pgfqpoint{0.953662in}{1.776726in}}%
\pgfpathlineto{\pgfqpoint{0.937739in}{1.705964in}}%
\pgfpathlineto{\pgfqpoint{0.861529in}{1.723655in}}%
\pgfpathlineto{\pgfqpoint{0.824452in}{1.732928in}}%
\pgfpathlineto{\pgfqpoint{0.824016in}{1.733109in}}%
\pgfpathlineto{\pgfqpoint{0.855951in}{1.862354in}}%
\pgfpathlineto{\pgfqpoint{0.860241in}{1.876128in}}%
\pgfpathlineto{\pgfqpoint{0.863735in}{1.876688in}}%
\pgfpathlineto{\pgfqpoint{0.867712in}{1.868584in}}%
\pgfpathlineto{\pgfqpoint{0.875796in}{1.867766in}}%
\pgfpathlineto{\pgfqpoint{0.879074in}{1.881212in}}%
\pgfpathlineto{\pgfqpoint{0.886843in}{1.879251in}}%
\pgfpathlineto{\pgfqpoint{0.890663in}{1.885325in}}%
\pgfpathlineto{\pgfqpoint{0.895063in}{1.884202in}}%
\pgfpathlineto{\pgfqpoint{0.896690in}{1.890839in}}%
\pgfpathlineto{\pgfqpoint{0.901859in}{1.889570in}}%
\pgfpathlineto{\pgfqpoint{0.905607in}{1.902658in}}%
\pgfpathlineto{\pgfqpoint{0.910516in}{1.911417in}}%
\pgfpathlineto{\pgfqpoint{0.919645in}{1.913720in}}%
\pgfpathlineto{\pgfqpoint{0.916455in}{1.900234in}}%
\pgfpathlineto{\pgfqpoint{0.913150in}{1.901010in}}%
\pgfpathlineto{\pgfqpoint{0.909972in}{1.887658in}}%
\pgfpathlineto{\pgfqpoint{0.915455in}{1.886073in}}%
\pgfpathlineto{\pgfqpoint{0.917584in}{1.892808in}}%
\pgfpathlineto{\pgfqpoint{0.949827in}{1.885175in}}%
\pgfpathlineto{\pgfqpoint{0.953430in}{1.884971in}}%
\pgfpathlineto{\pgfqpoint{0.964170in}{1.890384in}}%
\pgfpathlineto{\pgfqpoint{0.969814in}{1.885604in}}%
\pgfpathlineto{\pgfqpoint{0.968833in}{1.878204in}}%
\pgfpathlineto{\pgfqpoint{0.975886in}{1.873015in}}%
\pgfpathlineto{\pgfqpoint{0.972986in}{1.865357in}}%
\pgfpathlineto{\pgfqpoint{0.975129in}{1.855275in}}%
\pgfpathclose%
\pgfusepath{fill}%
\end{pgfscope}%
\begin{pgfscope}%
\pgfpathrectangle{\pgfqpoint{0.100000in}{0.100000in}}{\pgfqpoint{3.608454in}{2.310000in}}%
\pgfusepath{clip}%
\pgfsetbuttcap%
\pgfsetmiterjoin%
\definecolor{currentfill}{rgb}{0.000000,0.607843,0.696078}%
\pgfsetfillcolor{currentfill}%
\pgfsetlinewidth{0.000000pt}%
\definecolor{currentstroke}{rgb}{0.000000,0.000000,0.000000}%
\pgfsetstrokecolor{currentstroke}%
\pgfsetstrokeopacity{0.000000}%
\pgfsetdash{}{0pt}%
\pgfpathmoveto{\pgfqpoint{2.272778in}{0.593877in}}%
\pgfpathlineto{\pgfqpoint{2.264140in}{0.596489in}}%
\pgfpathlineto{\pgfqpoint{2.246043in}{0.604519in}}%
\pgfpathlineto{\pgfqpoint{2.232132in}{0.608488in}}%
\pgfpathlineto{\pgfqpoint{2.221611in}{0.607031in}}%
\pgfpathlineto{\pgfqpoint{2.210990in}{0.607615in}}%
\pgfpathlineto{\pgfqpoint{2.192546in}{0.604264in}}%
\pgfpathlineto{\pgfqpoint{2.187595in}{0.601129in}}%
\pgfpathlineto{\pgfqpoint{2.181395in}{0.610955in}}%
\pgfpathlineto{\pgfqpoint{2.188702in}{0.619056in}}%
\pgfpathlineto{\pgfqpoint{2.190675in}{0.624923in}}%
\pgfpathlineto{\pgfqpoint{2.196138in}{0.630275in}}%
\pgfpathlineto{\pgfqpoint{2.196201in}{0.649140in}}%
\pgfpathlineto{\pgfqpoint{2.191936in}{0.652528in}}%
\pgfpathlineto{\pgfqpoint{2.196505in}{0.661211in}}%
\pgfpathlineto{\pgfqpoint{2.193927in}{0.668423in}}%
\pgfpathlineto{\pgfqpoint{2.197286in}{0.677217in}}%
\pgfpathlineto{\pgfqpoint{2.200790in}{0.680240in}}%
\pgfpathlineto{\pgfqpoint{2.204721in}{0.696052in}}%
\pgfpathlineto{\pgfqpoint{2.207696in}{0.700477in}}%
\pgfpathlineto{\pgfqpoint{2.204984in}{0.705199in}}%
\pgfpathlineto{\pgfqpoint{2.205558in}{0.720837in}}%
\pgfpathlineto{\pgfqpoint{2.206617in}{0.728100in}}%
\pgfpathlineto{\pgfqpoint{2.212945in}{0.728230in}}%
\pgfpathlineto{\pgfqpoint{2.216369in}{0.735263in}}%
\pgfpathlineto{\pgfqpoint{2.226879in}{0.735487in}}%
\pgfpathlineto{\pgfqpoint{2.240752in}{0.735755in}}%
\pgfpathlineto{\pgfqpoint{2.244244in}{0.734424in}}%
\pgfpathlineto{\pgfqpoint{2.251222in}{0.732248in}}%
\pgfpathlineto{\pgfqpoint{2.254758in}{0.727139in}}%
\pgfpathlineto{\pgfqpoint{2.255860in}{0.698248in}}%
\pgfpathlineto{\pgfqpoint{2.271317in}{0.699125in}}%
\pgfpathlineto{\pgfqpoint{2.272007in}{0.674022in}}%
\pgfpathlineto{\pgfqpoint{2.269704in}{0.666105in}}%
\pgfpathlineto{\pgfqpoint{2.271892in}{0.630559in}}%
\pgfpathlineto{\pgfqpoint{2.272778in}{0.593877in}}%
\pgfpathclose%
\pgfusepath{fill}%
\end{pgfscope}%
\begin{pgfscope}%
\pgfpathrectangle{\pgfqpoint{0.100000in}{0.100000in}}{\pgfqpoint{3.608454in}{2.310000in}}%
\pgfusepath{clip}%
\pgfsetbuttcap%
\pgfsetmiterjoin%
\definecolor{currentfill}{rgb}{0.000000,0.850980,0.574510}%
\pgfsetfillcolor{currentfill}%
\pgfsetlinewidth{0.000000pt}%
\definecolor{currentstroke}{rgb}{0.000000,0.000000,0.000000}%
\pgfsetstrokecolor{currentstroke}%
\pgfsetstrokeopacity{0.000000}%
\pgfsetdash{}{0pt}%
\pgfpathmoveto{\pgfqpoint{0.792611in}{1.237224in}}%
\pgfpathlineto{\pgfqpoint{0.786789in}{1.237324in}}%
\pgfpathlineto{\pgfqpoint{0.765878in}{1.242208in}}%
\pgfpathlineto{\pgfqpoint{0.692238in}{1.260226in}}%
\pgfpathlineto{\pgfqpoint{0.668493in}{1.266395in}}%
\pgfpathlineto{\pgfqpoint{0.645007in}{1.271759in}}%
\pgfpathlineto{\pgfqpoint{0.648713in}{1.279430in}}%
\pgfpathlineto{\pgfqpoint{0.647370in}{1.296028in}}%
\pgfpathlineto{\pgfqpoint{0.649389in}{1.301776in}}%
\pgfpathlineto{\pgfqpoint{0.647795in}{1.312826in}}%
\pgfpathlineto{\pgfqpoint{0.650356in}{1.315914in}}%
\pgfpathlineto{\pgfqpoint{0.644042in}{1.329081in}}%
\pgfpathlineto{\pgfqpoint{0.644925in}{1.333356in}}%
\pgfpathlineto{\pgfqpoint{0.641670in}{1.344858in}}%
\pgfpathlineto{\pgfqpoint{0.642344in}{1.357942in}}%
\pgfpathlineto{\pgfqpoint{0.645186in}{1.361551in}}%
\pgfpathlineto{\pgfqpoint{0.643441in}{1.374382in}}%
\pgfpathlineto{\pgfqpoint{0.636091in}{1.384439in}}%
\pgfpathlineto{\pgfqpoint{0.632393in}{1.385724in}}%
\pgfpathlineto{\pgfqpoint{0.632454in}{1.400962in}}%
\pgfpathlineto{\pgfqpoint{0.628697in}{1.402513in}}%
\pgfpathlineto{\pgfqpoint{0.632163in}{1.410294in}}%
\pgfpathlineto{\pgfqpoint{0.627031in}{1.414709in}}%
\pgfpathlineto{\pgfqpoint{0.624814in}{1.421196in}}%
\pgfpathlineto{\pgfqpoint{0.619173in}{1.425619in}}%
\pgfpathlineto{\pgfqpoint{0.618630in}{1.433861in}}%
\pgfpathlineto{\pgfqpoint{0.616083in}{1.438599in}}%
\pgfpathlineto{\pgfqpoint{0.607467in}{1.441382in}}%
\pgfpathlineto{\pgfqpoint{0.612588in}{1.444548in}}%
\pgfpathlineto{\pgfqpoint{0.613111in}{1.453862in}}%
\pgfpathlineto{\pgfqpoint{0.609428in}{1.457861in}}%
\pgfpathlineto{\pgfqpoint{0.610025in}{1.468870in}}%
\pgfpathlineto{\pgfqpoint{0.597769in}{1.478364in}}%
\pgfpathlineto{\pgfqpoint{0.596004in}{1.489421in}}%
\pgfpathlineto{\pgfqpoint{0.597596in}{1.492295in}}%
\pgfpathlineto{\pgfqpoint{0.603382in}{1.497016in}}%
\pgfpathlineto{\pgfqpoint{0.606996in}{1.503863in}}%
\pgfpathlineto{\pgfqpoint{0.604651in}{1.513031in}}%
\pgfpathlineto{\pgfqpoint{0.609163in}{1.520157in}}%
\pgfpathlineto{\pgfqpoint{0.620520in}{1.503015in}}%
\pgfpathlineto{\pgfqpoint{0.628251in}{1.491065in}}%
\pgfpathlineto{\pgfqpoint{0.661400in}{1.439746in}}%
\pgfpathlineto{\pgfqpoint{0.688870in}{1.397366in}}%
\pgfpathlineto{\pgfqpoint{0.720005in}{1.349216in}}%
\pgfpathlineto{\pgfqpoint{0.780428in}{1.255922in}}%
\pgfpathlineto{\pgfqpoint{0.792611in}{1.237224in}}%
\pgfpathclose%
\pgfusepath{fill}%
\end{pgfscope}%
\begin{pgfscope}%
\pgfpathrectangle{\pgfqpoint{0.100000in}{0.100000in}}{\pgfqpoint{3.608454in}{2.310000in}}%
\pgfusepath{clip}%
\pgfsetbuttcap%
\pgfsetmiterjoin%
\definecolor{currentfill}{rgb}{0.000000,0.666667,0.666667}%
\pgfsetfillcolor{currentfill}%
\pgfsetlinewidth{0.000000pt}%
\definecolor{currentstroke}{rgb}{0.000000,0.000000,0.000000}%
\pgfsetstrokecolor{currentstroke}%
\pgfsetstrokeopacity{0.000000}%
\pgfsetdash{}{0pt}%
\pgfpathmoveto{\pgfqpoint{2.910462in}{0.795395in}}%
\pgfpathlineto{\pgfqpoint{2.894796in}{0.794134in}}%
\pgfpathlineto{\pgfqpoint{2.891131in}{0.799434in}}%
\pgfpathlineto{\pgfqpoint{2.885535in}{0.800050in}}%
\pgfpathlineto{\pgfqpoint{2.884962in}{0.810094in}}%
\pgfpathlineto{\pgfqpoint{2.887379in}{0.811341in}}%
\pgfpathlineto{\pgfqpoint{2.884148in}{0.822135in}}%
\pgfpathlineto{\pgfqpoint{2.904439in}{0.824809in}}%
\pgfpathlineto{\pgfqpoint{2.917347in}{0.821092in}}%
\pgfpathlineto{\pgfqpoint{2.918401in}{0.811582in}}%
\pgfpathlineto{\pgfqpoint{2.914876in}{0.807165in}}%
\pgfpathlineto{\pgfqpoint{2.912735in}{0.796174in}}%
\pgfpathlineto{\pgfqpoint{2.910462in}{0.795395in}}%
\pgfpathclose%
\pgfusepath{fill}%
\end{pgfscope}%
\begin{pgfscope}%
\pgfpathrectangle{\pgfqpoint{0.100000in}{0.100000in}}{\pgfqpoint{3.608454in}{2.310000in}}%
\pgfusepath{clip}%
\pgfsetbuttcap%
\pgfsetmiterjoin%
\definecolor{currentfill}{rgb}{0.000000,0.788235,0.605882}%
\pgfsetfillcolor{currentfill}%
\pgfsetlinewidth{0.000000pt}%
\definecolor{currentstroke}{rgb}{0.000000,0.000000,0.000000}%
\pgfsetstrokecolor{currentstroke}%
\pgfsetstrokeopacity{0.000000}%
\pgfsetdash{}{0pt}%
\pgfpathmoveto{\pgfqpoint{0.465159in}{1.900474in}}%
\pgfpathlineto{\pgfqpoint{0.478228in}{1.889504in}}%
\pgfpathlineto{\pgfqpoint{0.483261in}{1.890283in}}%
\pgfpathlineto{\pgfqpoint{0.491499in}{1.886385in}}%
\pgfpathlineto{\pgfqpoint{0.495724in}{1.888283in}}%
\pgfpathlineto{\pgfqpoint{0.503506in}{1.885466in}}%
\pgfpathlineto{\pgfqpoint{0.516046in}{1.883757in}}%
\pgfpathlineto{\pgfqpoint{0.532832in}{1.891746in}}%
\pgfpathlineto{\pgfqpoint{0.538396in}{1.890064in}}%
\pgfpathlineto{\pgfqpoint{0.543027in}{1.893318in}}%
\pgfpathlineto{\pgfqpoint{0.549594in}{1.891292in}}%
\pgfpathlineto{\pgfqpoint{0.526008in}{1.816628in}}%
\pgfpathlineto{\pgfqpoint{0.456983in}{1.837942in}}%
\pgfpathlineto{\pgfqpoint{0.439878in}{1.843077in}}%
\pgfpathlineto{\pgfqpoint{0.441504in}{1.854677in}}%
\pgfpathlineto{\pgfqpoint{0.447509in}{1.860643in}}%
\pgfpathlineto{\pgfqpoint{0.444876in}{1.871177in}}%
\pgfpathlineto{\pgfqpoint{0.437497in}{1.874169in}}%
\pgfpathlineto{\pgfqpoint{0.441915in}{1.883291in}}%
\pgfpathlineto{\pgfqpoint{0.448521in}{1.882309in}}%
\pgfpathlineto{\pgfqpoint{0.455038in}{1.891185in}}%
\pgfpathlineto{\pgfqpoint{0.465159in}{1.900474in}}%
\pgfpathclose%
\pgfusepath{fill}%
\end{pgfscope}%
\begin{pgfscope}%
\pgfpathrectangle{\pgfqpoint{0.100000in}{0.100000in}}{\pgfqpoint{3.608454in}{2.310000in}}%
\pgfusepath{clip}%
\pgfsetbuttcap%
\pgfsetmiterjoin%
\definecolor{currentfill}{rgb}{0.000000,0.662745,0.668627}%
\pgfsetfillcolor{currentfill}%
\pgfsetlinewidth{0.000000pt}%
\definecolor{currentstroke}{rgb}{0.000000,0.000000,0.000000}%
\pgfsetstrokecolor{currentstroke}%
\pgfsetstrokeopacity{0.000000}%
\pgfsetdash{}{0pt}%
\pgfpathmoveto{\pgfqpoint{2.734886in}{0.875470in}}%
\pgfpathlineto{\pgfqpoint{2.728245in}{0.872298in}}%
\pgfpathlineto{\pgfqpoint{2.729130in}{0.864274in}}%
\pgfpathlineto{\pgfqpoint{2.722392in}{0.862394in}}%
\pgfpathlineto{\pgfqpoint{2.723195in}{0.855536in}}%
\pgfpathlineto{\pgfqpoint{2.717264in}{0.854842in}}%
\pgfpathlineto{\pgfqpoint{2.715600in}{0.875575in}}%
\pgfpathlineto{\pgfqpoint{2.682505in}{0.872318in}}%
\pgfpathlineto{\pgfqpoint{2.676747in}{0.876531in}}%
\pgfpathlineto{\pgfqpoint{2.672768in}{0.884091in}}%
\pgfpathlineto{\pgfqpoint{2.670991in}{0.904805in}}%
\pgfpathlineto{\pgfqpoint{2.673491in}{0.910352in}}%
\pgfpathlineto{\pgfqpoint{2.679726in}{0.915262in}}%
\pgfpathlineto{\pgfqpoint{2.677615in}{0.922675in}}%
\pgfpathlineto{\pgfqpoint{2.688862in}{0.935943in}}%
\pgfpathlineto{\pgfqpoint{2.686513in}{0.939240in}}%
\pgfpathlineto{\pgfqpoint{2.690775in}{0.946928in}}%
\pgfpathlineto{\pgfqpoint{2.697028in}{0.947184in}}%
\pgfpathlineto{\pgfqpoint{2.701405in}{0.940587in}}%
\pgfpathlineto{\pgfqpoint{2.714599in}{0.941836in}}%
\pgfpathlineto{\pgfqpoint{2.717216in}{0.935165in}}%
\pgfpathlineto{\pgfqpoint{2.725270in}{0.935776in}}%
\pgfpathlineto{\pgfqpoint{2.727775in}{0.904950in}}%
\pgfpathlineto{\pgfqpoint{2.731773in}{0.905420in}}%
\pgfpathlineto{\pgfqpoint{2.734886in}{0.875470in}}%
\pgfpathclose%
\pgfusepath{fill}%
\end{pgfscope}%
\begin{pgfscope}%
\pgfpathrectangle{\pgfqpoint{0.100000in}{0.100000in}}{\pgfqpoint{3.608454in}{2.310000in}}%
\pgfusepath{clip}%
\pgfsetbuttcap%
\pgfsetmiterjoin%
\definecolor{currentfill}{rgb}{0.000000,0.533333,0.733333}%
\pgfsetfillcolor{currentfill}%
\pgfsetlinewidth{0.000000pt}%
\definecolor{currentstroke}{rgb}{0.000000,0.000000,0.000000}%
\pgfsetstrokecolor{currentstroke}%
\pgfsetstrokeopacity{0.000000}%
\pgfsetdash{}{0pt}%
\pgfpathmoveto{\pgfqpoint{2.460863in}{1.433950in}}%
\pgfpathlineto{\pgfqpoint{2.459428in}{1.452340in}}%
\pgfpathlineto{\pgfqpoint{2.452287in}{1.455302in}}%
\pgfpathlineto{\pgfqpoint{2.450616in}{1.476514in}}%
\pgfpathlineto{\pgfqpoint{2.458726in}{1.477228in}}%
\pgfpathlineto{\pgfqpoint{2.467299in}{1.483195in}}%
\pgfpathlineto{\pgfqpoint{2.470141in}{1.490384in}}%
\pgfpathlineto{\pgfqpoint{2.468185in}{1.518191in}}%
\pgfpathlineto{\pgfqpoint{2.488674in}{1.519766in}}%
\pgfpathlineto{\pgfqpoint{2.508522in}{1.521688in}}%
\pgfpathlineto{\pgfqpoint{2.509486in}{1.512307in}}%
\pgfpathlineto{\pgfqpoint{2.512434in}{1.482660in}}%
\pgfpathlineto{\pgfqpoint{2.498999in}{1.481611in}}%
\pgfpathlineto{\pgfqpoint{2.500893in}{1.455137in}}%
\pgfpathlineto{\pgfqpoint{2.493952in}{1.454636in}}%
\pgfpathlineto{\pgfqpoint{2.484922in}{1.436043in}}%
\pgfpathlineto{\pgfqpoint{2.460863in}{1.433950in}}%
\pgfpathclose%
\pgfusepath{fill}%
\end{pgfscope}%
\begin{pgfscope}%
\pgfpathrectangle{\pgfqpoint{0.100000in}{0.100000in}}{\pgfqpoint{3.608454in}{2.310000in}}%
\pgfusepath{clip}%
\pgfsetbuttcap%
\pgfsetmiterjoin%
\definecolor{currentfill}{rgb}{0.000000,0.643137,0.678431}%
\pgfsetfillcolor{currentfill}%
\pgfsetlinewidth{0.000000pt}%
\definecolor{currentstroke}{rgb}{0.000000,0.000000,0.000000}%
\pgfsetstrokecolor{currentstroke}%
\pgfsetstrokeopacity{0.000000}%
\pgfsetdash{}{0pt}%
\pgfpathmoveto{\pgfqpoint{2.322557in}{0.730252in}}%
\pgfpathlineto{\pgfqpoint{2.333480in}{0.724393in}}%
\pgfpathlineto{\pgfqpoint{2.330790in}{0.720030in}}%
\pgfpathlineto{\pgfqpoint{2.330722in}{0.712813in}}%
\pgfpathlineto{\pgfqpoint{2.325328in}{0.707042in}}%
\pgfpathlineto{\pgfqpoint{2.325022in}{0.696944in}}%
\pgfpathlineto{\pgfqpoint{2.296084in}{0.696005in}}%
\pgfpathlineto{\pgfqpoint{2.292928in}{0.702917in}}%
\pgfpathlineto{\pgfqpoint{2.285347in}{0.707925in}}%
\pgfpathlineto{\pgfqpoint{2.278789in}{0.700272in}}%
\pgfpathlineto{\pgfqpoint{2.271317in}{0.699125in}}%
\pgfpathlineto{\pgfqpoint{2.255860in}{0.698248in}}%
\pgfpathlineto{\pgfqpoint{2.254758in}{0.727139in}}%
\pgfpathlineto{\pgfqpoint{2.251222in}{0.732248in}}%
\pgfpathlineto{\pgfqpoint{2.244244in}{0.734424in}}%
\pgfpathlineto{\pgfqpoint{2.248864in}{0.736251in}}%
\pgfpathlineto{\pgfqpoint{2.261771in}{0.748472in}}%
\pgfpathlineto{\pgfqpoint{2.256019in}{0.754662in}}%
\pgfpathlineto{\pgfqpoint{2.248851in}{0.756615in}}%
\pgfpathlineto{\pgfqpoint{2.244140in}{0.763307in}}%
\pgfpathlineto{\pgfqpoint{2.268095in}{0.763939in}}%
\pgfpathlineto{\pgfqpoint{2.267964in}{0.770915in}}%
\pgfpathlineto{\pgfqpoint{2.285332in}{0.771355in}}%
\pgfpathlineto{\pgfqpoint{2.286412in}{0.757594in}}%
\pgfpathlineto{\pgfqpoint{2.291967in}{0.750029in}}%
\pgfpathlineto{\pgfqpoint{2.299838in}{0.745415in}}%
\pgfpathlineto{\pgfqpoint{2.300756in}{0.740640in}}%
\pgfpathlineto{\pgfqpoint{2.311015in}{0.733955in}}%
\pgfpathlineto{\pgfqpoint{2.322557in}{0.730252in}}%
\pgfpathclose%
\pgfusepath{fill}%
\end{pgfscope}%
\begin{pgfscope}%
\pgfpathrectangle{\pgfqpoint{0.100000in}{0.100000in}}{\pgfqpoint{3.608454in}{2.310000in}}%
\pgfusepath{clip}%
\pgfsetbuttcap%
\pgfsetmiterjoin%
\definecolor{currentfill}{rgb}{0.000000,0.623529,0.688235}%
\pgfsetfillcolor{currentfill}%
\pgfsetlinewidth{0.000000pt}%
\definecolor{currentstroke}{rgb}{0.000000,0.000000,0.000000}%
\pgfsetstrokecolor{currentstroke}%
\pgfsetstrokeopacity{0.000000}%
\pgfsetdash{}{0pt}%
\pgfpathmoveto{\pgfqpoint{2.537345in}{1.268351in}}%
\pgfpathlineto{\pgfqpoint{2.534479in}{1.269453in}}%
\pgfpathlineto{\pgfqpoint{2.519677in}{1.268419in}}%
\pgfpathlineto{\pgfqpoint{2.517928in}{1.295771in}}%
\pgfpathlineto{\pgfqpoint{2.497373in}{1.294428in}}%
\pgfpathlineto{\pgfqpoint{2.495702in}{1.322010in}}%
\pgfpathlineto{\pgfqpoint{2.523132in}{1.323450in}}%
\pgfpathlineto{\pgfqpoint{2.521612in}{1.327886in}}%
\pgfpathlineto{\pgfqpoint{2.521388in}{1.342978in}}%
\pgfpathlineto{\pgfqpoint{2.542895in}{1.344839in}}%
\pgfpathlineto{\pgfqpoint{2.544363in}{1.322774in}}%
\pgfpathlineto{\pgfqpoint{2.542054in}{1.316771in}}%
\pgfpathlineto{\pgfqpoint{2.541523in}{1.307274in}}%
\pgfpathlineto{\pgfqpoint{2.543804in}{1.301922in}}%
\pgfpathlineto{\pgfqpoint{2.541416in}{1.297901in}}%
\pgfpathlineto{\pgfqpoint{2.545025in}{1.287628in}}%
\pgfpathlineto{\pgfqpoint{2.537345in}{1.268351in}}%
\pgfpathclose%
\pgfusepath{fill}%
\end{pgfscope}%
\begin{pgfscope}%
\pgfpathrectangle{\pgfqpoint{0.100000in}{0.100000in}}{\pgfqpoint{3.608454in}{2.310000in}}%
\pgfusepath{clip}%
\pgfsetbuttcap%
\pgfsetmiterjoin%
\definecolor{currentfill}{rgb}{0.000000,0.741176,0.629412}%
\pgfsetfillcolor{currentfill}%
\pgfsetlinewidth{0.000000pt}%
\definecolor{currentstroke}{rgb}{0.000000,0.000000,0.000000}%
\pgfsetstrokecolor{currentstroke}%
\pgfsetstrokeopacity{0.000000}%
\pgfsetdash{}{0pt}%
\pgfpathmoveto{\pgfqpoint{2.920044in}{1.373705in}}%
\pgfpathlineto{\pgfqpoint{2.909141in}{1.370242in}}%
\pgfpathlineto{\pgfqpoint{2.903048in}{1.365422in}}%
\pgfpathlineto{\pgfqpoint{2.905122in}{1.355064in}}%
\pgfpathlineto{\pgfqpoint{2.893788in}{1.362802in}}%
\pgfpathlineto{\pgfqpoint{2.885129in}{1.361064in}}%
\pgfpathlineto{\pgfqpoint{2.884573in}{1.368240in}}%
\pgfpathlineto{\pgfqpoint{2.877293in}{1.367722in}}%
\pgfpathlineto{\pgfqpoint{2.875638in}{1.374674in}}%
\pgfpathlineto{\pgfqpoint{2.869516in}{1.374371in}}%
\pgfpathlineto{\pgfqpoint{2.869245in}{1.379552in}}%
\pgfpathlineto{\pgfqpoint{2.876585in}{1.380575in}}%
\pgfpathlineto{\pgfqpoint{2.875723in}{1.395664in}}%
\pgfpathlineto{\pgfqpoint{2.882640in}{1.396033in}}%
\pgfpathlineto{\pgfqpoint{2.881668in}{1.410091in}}%
\pgfpathlineto{\pgfqpoint{2.883523in}{1.417325in}}%
\pgfpathlineto{\pgfqpoint{2.880748in}{1.424222in}}%
\pgfpathlineto{\pgfqpoint{2.880330in}{1.431219in}}%
\pgfpathlineto{\pgfqpoint{2.887216in}{1.431682in}}%
\pgfpathlineto{\pgfqpoint{2.886708in}{1.438806in}}%
\pgfpathlineto{\pgfqpoint{2.893522in}{1.439324in}}%
\pgfpathlineto{\pgfqpoint{2.894009in}{1.432157in}}%
\pgfpathlineto{\pgfqpoint{2.906878in}{1.433220in}}%
\pgfpathlineto{\pgfqpoint{2.907865in}{1.422811in}}%
\pgfpathlineto{\pgfqpoint{2.916394in}{1.420191in}}%
\pgfpathlineto{\pgfqpoint{2.915344in}{1.412906in}}%
\pgfpathlineto{\pgfqpoint{2.916733in}{1.399937in}}%
\pgfpathlineto{\pgfqpoint{2.910035in}{1.386821in}}%
\pgfpathlineto{\pgfqpoint{2.920044in}{1.373705in}}%
\pgfpathclose%
\pgfusepath{fill}%
\end{pgfscope}%
\begin{pgfscope}%
\pgfpathrectangle{\pgfqpoint{0.100000in}{0.100000in}}{\pgfqpoint{3.608454in}{2.310000in}}%
\pgfusepath{clip}%
\pgfsetbuttcap%
\pgfsetmiterjoin%
\definecolor{currentfill}{rgb}{0.000000,0.443137,0.778431}%
\pgfsetfillcolor{currentfill}%
\pgfsetlinewidth{0.000000pt}%
\definecolor{currentstroke}{rgb}{0.000000,0.000000,0.000000}%
\pgfsetstrokecolor{currentstroke}%
\pgfsetstrokeopacity{0.000000}%
\pgfsetdash{}{0pt}%
\pgfpathmoveto{\pgfqpoint{1.893270in}{1.388628in}}%
\pgfpathlineto{\pgfqpoint{1.928379in}{1.387333in}}%
\pgfpathlineto{\pgfqpoint{1.927666in}{1.366756in}}%
\pgfpathlineto{\pgfqpoint{1.927400in}{1.359864in}}%
\pgfpathlineto{\pgfqpoint{1.893147in}{1.361110in}}%
\pgfpathlineto{\pgfqpoint{1.892889in}{1.354245in}}%
\pgfpathlineto{\pgfqpoint{1.858721in}{1.355685in}}%
\pgfpathlineto{\pgfqpoint{1.860470in}{1.390012in}}%
\pgfpathlineto{\pgfqpoint{1.893270in}{1.388628in}}%
\pgfpathclose%
\pgfusepath{fill}%
\end{pgfscope}%
\begin{pgfscope}%
\pgfpathrectangle{\pgfqpoint{0.100000in}{0.100000in}}{\pgfqpoint{3.608454in}{2.310000in}}%
\pgfusepath{clip}%
\pgfsetbuttcap%
\pgfsetmiterjoin%
\definecolor{currentfill}{rgb}{0.000000,0.341176,0.829412}%
\pgfsetfillcolor{currentfill}%
\pgfsetlinewidth{0.000000pt}%
\definecolor{currentstroke}{rgb}{0.000000,0.000000,0.000000}%
\pgfsetstrokecolor{currentstroke}%
\pgfsetstrokeopacity{0.000000}%
\pgfsetdash{}{0pt}%
\pgfpathmoveto{\pgfqpoint{3.214106in}{1.265774in}}%
\pgfpathlineto{\pgfqpoint{3.210178in}{1.266083in}}%
\pgfpathlineto{\pgfqpoint{3.209973in}{1.277484in}}%
\pgfpathlineto{\pgfqpoint{3.202830in}{1.275174in}}%
\pgfpathlineto{\pgfqpoint{3.197116in}{1.277528in}}%
\pgfpathlineto{\pgfqpoint{3.189244in}{1.282467in}}%
\pgfpathlineto{\pgfqpoint{3.181322in}{1.282092in}}%
\pgfpathlineto{\pgfqpoint{3.178373in}{1.294019in}}%
\pgfpathlineto{\pgfqpoint{3.184110in}{1.299032in}}%
\pgfpathlineto{\pgfqpoint{3.155665in}{1.301994in}}%
\pgfpathlineto{\pgfqpoint{3.154386in}{1.307543in}}%
\pgfpathlineto{\pgfqpoint{3.159486in}{1.315467in}}%
\pgfpathlineto{\pgfqpoint{3.159782in}{1.331466in}}%
\pgfpathlineto{\pgfqpoint{3.153524in}{1.337961in}}%
\pgfpathlineto{\pgfqpoint{3.157131in}{1.351262in}}%
\pgfpathlineto{\pgfqpoint{3.165049in}{1.347998in}}%
\pgfpathlineto{\pgfqpoint{3.169983in}{1.342070in}}%
\pgfpathlineto{\pgfqpoint{3.176043in}{1.340697in}}%
\pgfpathlineto{\pgfqpoint{3.178627in}{1.363629in}}%
\pgfpathlineto{\pgfqpoint{3.195177in}{1.357724in}}%
\pgfpathlineto{\pgfqpoint{3.197288in}{1.352635in}}%
\pgfpathlineto{\pgfqpoint{3.202566in}{1.350706in}}%
\pgfpathlineto{\pgfqpoint{3.210448in}{1.365277in}}%
\pgfpathlineto{\pgfqpoint{3.217018in}{1.367687in}}%
\pgfpathlineto{\pgfqpoint{3.222015in}{1.359018in}}%
\pgfpathlineto{\pgfqpoint{3.228182in}{1.355944in}}%
\pgfpathlineto{\pgfqpoint{3.236077in}{1.357801in}}%
\pgfpathlineto{\pgfqpoint{3.240232in}{1.353321in}}%
\pgfpathlineto{\pgfqpoint{3.248164in}{1.344820in}}%
\pgfpathlineto{\pgfqpoint{3.249638in}{1.333200in}}%
\pgfpathlineto{\pgfqpoint{3.236059in}{1.328225in}}%
\pgfpathlineto{\pgfqpoint{3.240056in}{1.312499in}}%
\pgfpathlineto{\pgfqpoint{3.234147in}{1.313361in}}%
\pgfpathlineto{\pgfqpoint{3.224896in}{1.300661in}}%
\pgfpathlineto{\pgfqpoint{3.236491in}{1.297792in}}%
\pgfpathlineto{\pgfqpoint{3.240057in}{1.290021in}}%
\pgfpathlineto{\pgfqpoint{3.214106in}{1.265774in}}%
\pgfpathclose%
\pgfusepath{fill}%
\end{pgfscope}%
\begin{pgfscope}%
\pgfpathrectangle{\pgfqpoint{0.100000in}{0.100000in}}{\pgfqpoint{3.608454in}{2.310000in}}%
\pgfusepath{clip}%
\pgfsetbuttcap%
\pgfsetmiterjoin%
\definecolor{currentfill}{rgb}{0.000000,0.603922,0.698039}%
\pgfsetfillcolor{currentfill}%
\pgfsetlinewidth{0.000000pt}%
\definecolor{currentstroke}{rgb}{0.000000,0.000000,0.000000}%
\pgfsetstrokecolor{currentstroke}%
\pgfsetstrokeopacity{0.000000}%
\pgfsetdash{}{0pt}%
\pgfpathmoveto{\pgfqpoint{1.997875in}{0.723825in}}%
\pgfpathlineto{\pgfqpoint{1.981851in}{0.714949in}}%
\pgfpathlineto{\pgfqpoint{1.965169in}{0.705757in}}%
\pgfpathlineto{\pgfqpoint{1.951524in}{0.729506in}}%
\pgfpathlineto{\pgfqpoint{1.947012in}{0.726789in}}%
\pgfpathlineto{\pgfqpoint{1.942049in}{0.732981in}}%
\pgfpathlineto{\pgfqpoint{1.924502in}{0.764640in}}%
\pgfpathlineto{\pgfqpoint{1.919219in}{0.761725in}}%
\pgfpathlineto{\pgfqpoint{1.903937in}{0.789730in}}%
\pgfpathlineto{\pgfqpoint{1.913591in}{0.794961in}}%
\pgfpathlineto{\pgfqpoint{1.929646in}{0.804003in}}%
\pgfpathlineto{\pgfqpoint{1.934684in}{0.799256in}}%
\pgfpathlineto{\pgfqpoint{1.940012in}{0.801698in}}%
\pgfpathlineto{\pgfqpoint{1.966416in}{0.807704in}}%
\pgfpathlineto{\pgfqpoint{1.975845in}{0.790506in}}%
\pgfpathlineto{\pgfqpoint{1.978879in}{0.792186in}}%
\pgfpathlineto{\pgfqpoint{1.990434in}{0.771318in}}%
\pgfpathlineto{\pgfqpoint{1.975832in}{0.763164in}}%
\pgfpathlineto{\pgfqpoint{1.997875in}{0.723825in}}%
\pgfpathclose%
\pgfusepath{fill}%
\end{pgfscope}%
\begin{pgfscope}%
\pgfpathrectangle{\pgfqpoint{0.100000in}{0.100000in}}{\pgfqpoint{3.608454in}{2.310000in}}%
\pgfusepath{clip}%
\pgfsetbuttcap%
\pgfsetmiterjoin%
\definecolor{currentfill}{rgb}{0.000000,0.592157,0.703922}%
\pgfsetfillcolor{currentfill}%
\pgfsetlinewidth{0.000000pt}%
\definecolor{currentstroke}{rgb}{0.000000,0.000000,0.000000}%
\pgfsetstrokecolor{currentstroke}%
\pgfsetstrokeopacity{0.000000}%
\pgfsetdash{}{0pt}%
\pgfpathmoveto{\pgfqpoint{2.978628in}{1.415804in}}%
\pgfpathlineto{\pgfqpoint{2.973331in}{1.416066in}}%
\pgfpathlineto{\pgfqpoint{2.965224in}{1.427879in}}%
\pgfpathlineto{\pgfqpoint{2.966480in}{1.429783in}}%
\pgfpathlineto{\pgfqpoint{2.980855in}{1.444545in}}%
\pgfpathlineto{\pgfqpoint{2.987075in}{1.444310in}}%
\pgfpathlineto{\pgfqpoint{2.989302in}{1.450132in}}%
\pgfpathlineto{\pgfqpoint{2.986701in}{1.454490in}}%
\pgfpathlineto{\pgfqpoint{2.992195in}{1.461918in}}%
\pgfpathlineto{\pgfqpoint{2.989320in}{1.467586in}}%
\pgfpathlineto{\pgfqpoint{3.045971in}{1.477046in}}%
\pgfpathlineto{\pgfqpoint{3.052407in}{1.435951in}}%
\pgfpathlineto{\pgfqpoint{3.039311in}{1.439843in}}%
\pgfpathlineto{\pgfqpoint{3.032463in}{1.435406in}}%
\pgfpathlineto{\pgfqpoint{3.025784in}{1.440073in}}%
\pgfpathlineto{\pgfqpoint{3.019335in}{1.434599in}}%
\pgfpathlineto{\pgfqpoint{3.010773in}{1.432733in}}%
\pgfpathlineto{\pgfqpoint{3.008758in}{1.422029in}}%
\pgfpathlineto{\pgfqpoint{3.004458in}{1.420613in}}%
\pgfpathlineto{\pgfqpoint{2.986872in}{1.422883in}}%
\pgfpathlineto{\pgfqpoint{2.978628in}{1.415804in}}%
\pgfpathclose%
\pgfusepath{fill}%
\end{pgfscope}%
\begin{pgfscope}%
\pgfpathrectangle{\pgfqpoint{0.100000in}{0.100000in}}{\pgfqpoint{3.608454in}{2.310000in}}%
\pgfusepath{clip}%
\pgfsetbuttcap%
\pgfsetmiterjoin%
\definecolor{currentfill}{rgb}{0.000000,0.607843,0.696078}%
\pgfsetfillcolor{currentfill}%
\pgfsetlinewidth{0.000000pt}%
\definecolor{currentstroke}{rgb}{0.000000,0.000000,0.000000}%
\pgfsetstrokecolor{currentstroke}%
\pgfsetstrokeopacity{0.000000}%
\pgfsetdash{}{0pt}%
\pgfpathmoveto{\pgfqpoint{3.192630in}{1.620121in}}%
\pgfpathlineto{\pgfqpoint{3.195360in}{1.613257in}}%
\pgfpathlineto{\pgfqpoint{3.203466in}{1.605423in}}%
\pgfpathlineto{\pgfqpoint{3.206320in}{1.596141in}}%
\pgfpathlineto{\pgfqpoint{3.211501in}{1.598059in}}%
\pgfpathlineto{\pgfqpoint{3.213947in}{1.593817in}}%
\pgfpathlineto{\pgfqpoint{3.196959in}{1.580835in}}%
\pgfpathlineto{\pgfqpoint{3.184525in}{1.574369in}}%
\pgfpathlineto{\pgfqpoint{3.176843in}{1.578501in}}%
\pgfpathlineto{\pgfqpoint{3.157897in}{1.576523in}}%
\pgfpathlineto{\pgfqpoint{3.155139in}{1.587716in}}%
\pgfpathlineto{\pgfqpoint{3.165006in}{1.605498in}}%
\pgfpathlineto{\pgfqpoint{3.173979in}{1.610853in}}%
\pgfpathlineto{\pgfqpoint{3.174106in}{1.615801in}}%
\pgfpathlineto{\pgfqpoint{3.183685in}{1.619914in}}%
\pgfpathlineto{\pgfqpoint{3.192630in}{1.620121in}}%
\pgfpathclose%
\pgfusepath{fill}%
\end{pgfscope}%
\begin{pgfscope}%
\pgfpathrectangle{\pgfqpoint{0.100000in}{0.100000in}}{\pgfqpoint{3.608454in}{2.310000in}}%
\pgfusepath{clip}%
\pgfsetbuttcap%
\pgfsetmiterjoin%
\definecolor{currentfill}{rgb}{0.000000,0.784314,0.607843}%
\pgfsetfillcolor{currentfill}%
\pgfsetlinewidth{0.000000pt}%
\definecolor{currentstroke}{rgb}{0.000000,0.000000,0.000000}%
\pgfsetstrokecolor{currentstroke}%
\pgfsetstrokeopacity{0.000000}%
\pgfsetdash{}{0pt}%
\pgfpathmoveto{\pgfqpoint{2.549469in}{0.776821in}}%
\pgfpathlineto{\pgfqpoint{2.547754in}{0.792351in}}%
\pgfpathlineto{\pgfqpoint{2.549490in}{0.847081in}}%
\pgfpathlineto{\pgfqpoint{2.550184in}{0.875112in}}%
\pgfpathlineto{\pgfqpoint{2.550318in}{0.880073in}}%
\pgfpathlineto{\pgfqpoint{2.561497in}{0.881275in}}%
\pgfpathlineto{\pgfqpoint{2.564922in}{0.885225in}}%
\pgfpathlineto{\pgfqpoint{2.570549in}{0.885685in}}%
\pgfpathlineto{\pgfqpoint{2.575754in}{0.891315in}}%
\pgfpathlineto{\pgfqpoint{2.582729in}{0.895488in}}%
\pgfpathlineto{\pgfqpoint{2.583933in}{0.884712in}}%
\pgfpathlineto{\pgfqpoint{2.591982in}{0.885391in}}%
\pgfpathlineto{\pgfqpoint{2.591160in}{0.879866in}}%
\pgfpathlineto{\pgfqpoint{2.586494in}{0.875360in}}%
\pgfpathlineto{\pgfqpoint{2.586111in}{0.864841in}}%
\pgfpathlineto{\pgfqpoint{2.589627in}{0.852306in}}%
\pgfpathlineto{\pgfqpoint{2.588335in}{0.845764in}}%
\pgfpathlineto{\pgfqpoint{2.594215in}{0.842750in}}%
\pgfpathlineto{\pgfqpoint{2.607870in}{0.843949in}}%
\pgfpathlineto{\pgfqpoint{2.609026in}{0.830093in}}%
\pgfpathlineto{\pgfqpoint{2.612796in}{0.827022in}}%
\pgfpathlineto{\pgfqpoint{2.609428in}{0.825937in}}%
\pgfpathlineto{\pgfqpoint{2.610325in}{0.816266in}}%
\pgfpathlineto{\pgfqpoint{2.603471in}{0.815667in}}%
\pgfpathlineto{\pgfqpoint{2.604373in}{0.805606in}}%
\pgfpathlineto{\pgfqpoint{2.601434in}{0.804239in}}%
\pgfpathlineto{\pgfqpoint{2.574126in}{0.801999in}}%
\pgfpathlineto{\pgfqpoint{2.567924in}{0.787803in}}%
\pgfpathlineto{\pgfqpoint{2.573394in}{0.783619in}}%
\pgfpathlineto{\pgfqpoint{2.574889in}{0.778819in}}%
\pgfpathlineto{\pgfqpoint{2.549469in}{0.776821in}}%
\pgfpathclose%
\pgfusepath{fill}%
\end{pgfscope}%
\begin{pgfscope}%
\pgfpathrectangle{\pgfqpoint{0.100000in}{0.100000in}}{\pgfqpoint{3.608454in}{2.310000in}}%
\pgfusepath{clip}%
\pgfsetbuttcap%
\pgfsetmiterjoin%
\definecolor{currentfill}{rgb}{0.000000,0.800000,0.600000}%
\pgfsetfillcolor{currentfill}%
\pgfsetlinewidth{0.000000pt}%
\definecolor{currentstroke}{rgb}{0.000000,0.000000,0.000000}%
\pgfsetstrokecolor{currentstroke}%
\pgfsetstrokeopacity{0.000000}%
\pgfsetdash{}{0pt}%
\pgfpathmoveto{\pgfqpoint{2.361668in}{0.958818in}}%
\pgfpathlineto{\pgfqpoint{2.356734in}{0.953261in}}%
\pgfpathlineto{\pgfqpoint{2.361253in}{0.945877in}}%
\pgfpathlineto{\pgfqpoint{2.349626in}{0.945687in}}%
\pgfpathlineto{\pgfqpoint{2.341075in}{0.950676in}}%
\pgfpathlineto{\pgfqpoint{2.339665in}{0.959110in}}%
\pgfpathlineto{\pgfqpoint{2.335508in}{0.959654in}}%
\pgfpathlineto{\pgfqpoint{2.333604in}{0.966070in}}%
\pgfpathlineto{\pgfqpoint{2.322365in}{0.966382in}}%
\pgfpathlineto{\pgfqpoint{2.322083in}{0.986111in}}%
\pgfpathlineto{\pgfqpoint{2.326278in}{0.991006in}}%
\pgfpathlineto{\pgfqpoint{2.326016in}{0.997543in}}%
\pgfpathlineto{\pgfqpoint{2.321823in}{1.001132in}}%
\pgfpathlineto{\pgfqpoint{2.321574in}{1.020760in}}%
\pgfpathlineto{\pgfqpoint{2.314414in}{1.020654in}}%
\pgfpathlineto{\pgfqpoint{2.314200in}{1.029274in}}%
\pgfpathlineto{\pgfqpoint{2.320051in}{1.032034in}}%
\pgfpathlineto{\pgfqpoint{2.328311in}{1.029359in}}%
\pgfpathlineto{\pgfqpoint{2.328181in}{1.034618in}}%
\pgfpathlineto{\pgfqpoint{2.335892in}{1.034757in}}%
\pgfpathlineto{\pgfqpoint{2.342574in}{1.032044in}}%
\pgfpathlineto{\pgfqpoint{2.347372in}{1.026940in}}%
\pgfpathlineto{\pgfqpoint{2.356606in}{1.028702in}}%
\pgfpathlineto{\pgfqpoint{2.356746in}{1.021522in}}%
\pgfpathlineto{\pgfqpoint{2.360312in}{1.018136in}}%
\pgfpathlineto{\pgfqpoint{2.360588in}{1.007650in}}%
\pgfpathlineto{\pgfqpoint{2.364049in}{1.007736in}}%
\pgfpathlineto{\pgfqpoint{2.364599in}{0.986870in}}%
\pgfpathlineto{\pgfqpoint{2.368151in}{0.986916in}}%
\pgfpathlineto{\pgfqpoint{2.368481in}{0.976354in}}%
\pgfpathlineto{\pgfqpoint{2.364963in}{0.976228in}}%
\pgfpathlineto{\pgfqpoint{2.363796in}{0.962251in}}%
\pgfpathlineto{\pgfqpoint{2.361668in}{0.958818in}}%
\pgfpathclose%
\pgfusepath{fill}%
\end{pgfscope}%
\begin{pgfscope}%
\pgfpathrectangle{\pgfqpoint{0.100000in}{0.100000in}}{\pgfqpoint{3.608454in}{2.310000in}}%
\pgfusepath{clip}%
\pgfsetbuttcap%
\pgfsetmiterjoin%
\definecolor{currentfill}{rgb}{0.000000,0.403922,0.798039}%
\pgfsetfillcolor{currentfill}%
\pgfsetlinewidth{0.000000pt}%
\definecolor{currentstroke}{rgb}{0.000000,0.000000,0.000000}%
\pgfsetstrokecolor{currentstroke}%
\pgfsetstrokeopacity{0.000000}%
\pgfsetdash{}{0pt}%
\pgfpathmoveto{\pgfqpoint{1.670049in}{1.298787in}}%
\pgfpathlineto{\pgfqpoint{1.626877in}{1.302163in}}%
\pgfpathlineto{\pgfqpoint{1.579920in}{1.306214in}}%
\pgfpathlineto{\pgfqpoint{1.581597in}{1.326016in}}%
\pgfpathlineto{\pgfqpoint{1.602094in}{1.324825in}}%
\pgfpathlineto{\pgfqpoint{1.603799in}{1.337764in}}%
\pgfpathlineto{\pgfqpoint{1.606968in}{1.372189in}}%
\pgfpathlineto{\pgfqpoint{1.610603in}{1.406638in}}%
\pgfpathlineto{\pgfqpoint{1.677720in}{1.401672in}}%
\pgfpathlineto{\pgfqpoint{1.677687in}{1.401208in}}%
\pgfpathlineto{\pgfqpoint{1.672634in}{1.332579in}}%
\pgfpathlineto{\pgfqpoint{1.670049in}{1.298787in}}%
\pgfpathclose%
\pgfusepath{fill}%
\end{pgfscope}%
\begin{pgfscope}%
\pgfpathrectangle{\pgfqpoint{0.100000in}{0.100000in}}{\pgfqpoint{3.608454in}{2.310000in}}%
\pgfusepath{clip}%
\pgfsetbuttcap%
\pgfsetmiterjoin%
\definecolor{currentfill}{rgb}{0.000000,0.760784,0.619608}%
\pgfsetfillcolor{currentfill}%
\pgfsetlinewidth{0.000000pt}%
\definecolor{currentstroke}{rgb}{0.000000,0.000000,0.000000}%
\pgfsetstrokecolor{currentstroke}%
\pgfsetstrokeopacity{0.000000}%
\pgfsetdash{}{0pt}%
\pgfpathmoveto{\pgfqpoint{2.396088in}{0.992607in}}%
\pgfpathlineto{\pgfqpoint{2.395789in}{0.983518in}}%
\pgfpathlineto{\pgfqpoint{2.384480in}{0.979541in}}%
\pgfpathlineto{\pgfqpoint{2.384214in}{0.972642in}}%
\pgfpathlineto{\pgfqpoint{2.372289in}{0.959362in}}%
\pgfpathlineto{\pgfqpoint{2.372346in}{0.959262in}}%
\pgfpathlineto{\pgfqpoint{2.361668in}{0.958818in}}%
\pgfpathlineto{\pgfqpoint{2.363796in}{0.962251in}}%
\pgfpathlineto{\pgfqpoint{2.364963in}{0.976228in}}%
\pgfpathlineto{\pgfqpoint{2.368481in}{0.976354in}}%
\pgfpathlineto{\pgfqpoint{2.368151in}{0.986916in}}%
\pgfpathlineto{\pgfqpoint{2.364599in}{0.986870in}}%
\pgfpathlineto{\pgfqpoint{2.364049in}{1.007736in}}%
\pgfpathlineto{\pgfqpoint{2.360588in}{1.007650in}}%
\pgfpathlineto{\pgfqpoint{2.360312in}{1.018136in}}%
\pgfpathlineto{\pgfqpoint{2.356746in}{1.021522in}}%
\pgfpathlineto{\pgfqpoint{2.356606in}{1.028702in}}%
\pgfpathlineto{\pgfqpoint{2.363393in}{1.028869in}}%
\pgfpathlineto{\pgfqpoint{2.362249in}{1.063873in}}%
\pgfpathlineto{\pgfqpoint{2.389852in}{1.065601in}}%
\pgfpathlineto{\pgfqpoint{2.410748in}{1.065840in}}%
\pgfpathlineto{\pgfqpoint{2.420217in}{1.066185in}}%
\pgfpathlineto{\pgfqpoint{2.424647in}{1.062261in}}%
\pgfpathlineto{\pgfqpoint{2.420989in}{1.051604in}}%
\pgfpathlineto{\pgfqpoint{2.425407in}{1.049147in}}%
\pgfpathlineto{\pgfqpoint{2.424139in}{1.041016in}}%
\pgfpathlineto{\pgfqpoint{2.419408in}{1.040872in}}%
\pgfpathlineto{\pgfqpoint{2.417154in}{1.033431in}}%
\pgfpathlineto{\pgfqpoint{2.412041in}{1.029240in}}%
\pgfpathlineto{\pgfqpoint{2.415331in}{1.023934in}}%
\pgfpathlineto{\pgfqpoint{2.404445in}{1.017531in}}%
\pgfpathlineto{\pgfqpoint{2.401872in}{1.004184in}}%
\pgfpathlineto{\pgfqpoint{2.394572in}{0.999746in}}%
\pgfpathlineto{\pgfqpoint{2.396088in}{0.992607in}}%
\pgfpathclose%
\pgfusepath{fill}%
\end{pgfscope}%
\begin{pgfscope}%
\pgfpathrectangle{\pgfqpoint{0.100000in}{0.100000in}}{\pgfqpoint{3.608454in}{2.310000in}}%
\pgfusepath{clip}%
\pgfsetbuttcap%
\pgfsetmiterjoin%
\definecolor{currentfill}{rgb}{0.000000,0.470588,0.764706}%
\pgfsetfillcolor{currentfill}%
\pgfsetlinewidth{0.000000pt}%
\definecolor{currentstroke}{rgb}{0.000000,0.000000,0.000000}%
\pgfsetstrokecolor{currentstroke}%
\pgfsetstrokeopacity{0.000000}%
\pgfsetdash{}{0pt}%
\pgfpathmoveto{\pgfqpoint{2.720376in}{1.354526in}}%
\pgfpathlineto{\pgfqpoint{2.710112in}{1.347801in}}%
\pgfpathlineto{\pgfqpoint{2.708117in}{1.365276in}}%
\pgfpathlineto{\pgfqpoint{2.693342in}{1.363630in}}%
\pgfpathlineto{\pgfqpoint{2.691293in}{1.385879in}}%
\pgfpathlineto{\pgfqpoint{2.684115in}{1.380130in}}%
\pgfpathlineto{\pgfqpoint{2.676781in}{1.379202in}}%
\pgfpathlineto{\pgfqpoint{2.675138in}{1.396451in}}%
\pgfpathlineto{\pgfqpoint{2.677647in}{1.404907in}}%
\pgfpathlineto{\pgfqpoint{2.697801in}{1.407041in}}%
\pgfpathlineto{\pgfqpoint{2.697158in}{1.412713in}}%
\pgfpathlineto{\pgfqpoint{2.726448in}{1.415659in}}%
\pgfpathlineto{\pgfqpoint{2.729836in}{1.382914in}}%
\pgfpathlineto{\pgfqpoint{2.725607in}{1.378219in}}%
\pgfpathlineto{\pgfqpoint{2.729795in}{1.373635in}}%
\pgfpathlineto{\pgfqpoint{2.727957in}{1.367391in}}%
\pgfpathlineto{\pgfqpoint{2.733971in}{1.365444in}}%
\pgfpathlineto{\pgfqpoint{2.728448in}{1.358041in}}%
\pgfpathlineto{\pgfqpoint{2.720376in}{1.354526in}}%
\pgfpathclose%
\pgfusepath{fill}%
\end{pgfscope}%
\begin{pgfscope}%
\pgfpathrectangle{\pgfqpoint{0.100000in}{0.100000in}}{\pgfqpoint{3.608454in}{2.310000in}}%
\pgfusepath{clip}%
\pgfsetbuttcap%
\pgfsetmiterjoin%
\definecolor{currentfill}{rgb}{0.000000,0.670588,0.664706}%
\pgfsetfillcolor{currentfill}%
\pgfsetlinewidth{0.000000pt}%
\definecolor{currentstroke}{rgb}{0.000000,0.000000,0.000000}%
\pgfsetstrokecolor{currentstroke}%
\pgfsetstrokeopacity{0.000000}%
\pgfsetdash{}{0pt}%
\pgfpathmoveto{\pgfqpoint{1.708835in}{0.943953in}}%
\pgfpathlineto{\pgfqpoint{1.743217in}{0.941925in}}%
\pgfpathlineto{\pgfqpoint{1.774163in}{0.940149in}}%
\pgfpathlineto{\pgfqpoint{1.797246in}{0.938094in}}%
\pgfpathlineto{\pgfqpoint{1.810392in}{0.932854in}}%
\pgfpathlineto{\pgfqpoint{1.811548in}{0.930044in}}%
\pgfpathlineto{\pgfqpoint{1.810377in}{0.903512in}}%
\pgfpathlineto{\pgfqpoint{1.808729in}{0.868432in}}%
\pgfpathlineto{\pgfqpoint{1.763926in}{0.871040in}}%
\pgfpathlineto{\pgfqpoint{1.738977in}{0.872740in}}%
\pgfpathlineto{\pgfqpoint{1.741190in}{0.907227in}}%
\pgfpathlineto{\pgfqpoint{1.706689in}{0.909371in}}%
\pgfpathlineto{\pgfqpoint{1.708835in}{0.943953in}}%
\pgfpathclose%
\pgfusepath{fill}%
\end{pgfscope}%
\begin{pgfscope}%
\pgfpathrectangle{\pgfqpoint{0.100000in}{0.100000in}}{\pgfqpoint{3.608454in}{2.310000in}}%
\pgfusepath{clip}%
\pgfsetbuttcap%
\pgfsetmiterjoin%
\definecolor{currentfill}{rgb}{0.000000,0.835294,0.582353}%
\pgfsetfillcolor{currentfill}%
\pgfsetlinewidth{0.000000pt}%
\definecolor{currentstroke}{rgb}{0.000000,0.000000,0.000000}%
\pgfsetstrokecolor{currentstroke}%
\pgfsetstrokeopacity{0.000000}%
\pgfsetdash{}{0pt}%
\pgfpathmoveto{\pgfqpoint{2.336789in}{0.709454in}}%
\pgfpathlineto{\pgfqpoint{2.330722in}{0.712813in}}%
\pgfpathlineto{\pgfqpoint{2.330790in}{0.720030in}}%
\pgfpathlineto{\pgfqpoint{2.333480in}{0.724393in}}%
\pgfpathlineto{\pgfqpoint{2.322557in}{0.730252in}}%
\pgfpathlineto{\pgfqpoint{2.319974in}{0.746106in}}%
\pgfpathlineto{\pgfqpoint{2.325362in}{0.751678in}}%
\pgfpathlineto{\pgfqpoint{2.330853in}{0.768113in}}%
\pgfpathlineto{\pgfqpoint{2.334485in}{0.767523in}}%
\pgfpathlineto{\pgfqpoint{2.338396in}{0.773253in}}%
\pgfpathlineto{\pgfqpoint{2.338322in}{0.779971in}}%
\pgfpathlineto{\pgfqpoint{2.343388in}{0.795089in}}%
\pgfpathlineto{\pgfqpoint{2.342946in}{0.805709in}}%
\pgfpathlineto{\pgfqpoint{2.354522in}{0.806095in}}%
\pgfpathlineto{\pgfqpoint{2.366944in}{0.810346in}}%
\pgfpathlineto{\pgfqpoint{2.365200in}{0.801106in}}%
\pgfpathlineto{\pgfqpoint{2.374090in}{0.798791in}}%
\pgfpathlineto{\pgfqpoint{2.371635in}{0.792020in}}%
\pgfpathlineto{\pgfqpoint{2.364252in}{0.787285in}}%
\pgfpathlineto{\pgfqpoint{2.363437in}{0.782994in}}%
\pgfpathlineto{\pgfqpoint{2.353993in}{0.777417in}}%
\pgfpathlineto{\pgfqpoint{2.357930in}{0.763773in}}%
\pgfpathlineto{\pgfqpoint{2.367747in}{0.759421in}}%
\pgfpathlineto{\pgfqpoint{2.395986in}{0.760766in}}%
\pgfpathlineto{\pgfqpoint{2.403042in}{0.761124in}}%
\pgfpathlineto{\pgfqpoint{2.404037in}{0.740243in}}%
\pgfpathlineto{\pgfqpoint{2.377876in}{0.738907in}}%
\pgfpathlineto{\pgfqpoint{2.371609in}{0.736436in}}%
\pgfpathlineto{\pgfqpoint{2.362838in}{0.740120in}}%
\pgfpathlineto{\pgfqpoint{2.347315in}{0.735413in}}%
\pgfpathlineto{\pgfqpoint{2.346226in}{0.727949in}}%
\pgfpathlineto{\pgfqpoint{2.341299in}{0.728481in}}%
\pgfpathlineto{\pgfqpoint{2.337261in}{0.720197in}}%
\pgfpathlineto{\pgfqpoint{2.341562in}{0.715016in}}%
\pgfpathlineto{\pgfqpoint{2.336789in}{0.709454in}}%
\pgfpathclose%
\pgfusepath{fill}%
\end{pgfscope}%
\begin{pgfscope}%
\pgfpathrectangle{\pgfqpoint{0.100000in}{0.100000in}}{\pgfqpoint{3.608454in}{2.310000in}}%
\pgfusepath{clip}%
\pgfsetbuttcap%
\pgfsetmiterjoin%
\definecolor{currentfill}{rgb}{0.000000,0.533333,0.733333}%
\pgfsetfillcolor{currentfill}%
\pgfsetlinewidth{0.000000pt}%
\definecolor{currentstroke}{rgb}{0.000000,0.000000,0.000000}%
\pgfsetstrokecolor{currentstroke}%
\pgfsetstrokeopacity{0.000000}%
\pgfsetdash{}{0pt}%
\pgfpathmoveto{\pgfqpoint{2.368642in}{1.363523in}}%
\pgfpathlineto{\pgfqpoint{2.340427in}{1.362513in}}%
\pgfpathlineto{\pgfqpoint{2.340562in}{1.355530in}}%
\pgfpathlineto{\pgfqpoint{2.335984in}{1.355415in}}%
\pgfpathlineto{\pgfqpoint{2.326861in}{1.355261in}}%
\pgfpathlineto{\pgfqpoint{2.324527in}{1.369104in}}%
\pgfpathlineto{\pgfqpoint{2.307376in}{1.370072in}}%
\pgfpathlineto{\pgfqpoint{2.306650in}{1.395239in}}%
\pgfpathlineto{\pgfqpoint{2.298632in}{1.395025in}}%
\pgfpathlineto{\pgfqpoint{2.298225in}{1.417924in}}%
\pgfpathlineto{\pgfqpoint{2.291417in}{1.417786in}}%
\pgfpathlineto{\pgfqpoint{2.290718in}{1.442111in}}%
\pgfpathlineto{\pgfqpoint{2.318087in}{1.442426in}}%
\pgfpathlineto{\pgfqpoint{2.317746in}{1.438593in}}%
\pgfpathlineto{\pgfqpoint{2.353565in}{1.439495in}}%
\pgfpathlineto{\pgfqpoint{2.354524in}{1.412008in}}%
\pgfpathlineto{\pgfqpoint{2.375453in}{1.412534in}}%
\pgfpathlineto{\pgfqpoint{2.371479in}{1.401180in}}%
\pgfpathlineto{\pgfqpoint{2.376576in}{1.388834in}}%
\pgfpathlineto{\pgfqpoint{2.374602in}{1.377391in}}%
\pgfpathlineto{\pgfqpoint{2.354879in}{1.376769in}}%
\pgfpathlineto{\pgfqpoint{2.363949in}{1.370060in}}%
\pgfpathlineto{\pgfqpoint{2.368642in}{1.363523in}}%
\pgfpathclose%
\pgfusepath{fill}%
\end{pgfscope}%
\begin{pgfscope}%
\pgfpathrectangle{\pgfqpoint{0.100000in}{0.100000in}}{\pgfqpoint{3.608454in}{2.310000in}}%
\pgfusepath{clip}%
\pgfsetbuttcap%
\pgfsetmiterjoin%
\definecolor{currentfill}{rgb}{0.000000,0.858824,0.570588}%
\pgfsetfillcolor{currentfill}%
\pgfsetlinewidth{0.000000pt}%
\definecolor{currentstroke}{rgb}{0.000000,0.000000,0.000000}%
\pgfsetstrokecolor{currentstroke}%
\pgfsetstrokeopacity{0.000000}%
\pgfsetdash{}{0pt}%
\pgfpathmoveto{\pgfqpoint{1.007808in}{2.262053in}}%
\pgfpathlineto{\pgfqpoint{0.998620in}{2.223097in}}%
\pgfpathlineto{\pgfqpoint{0.960593in}{2.232110in}}%
\pgfpathlineto{\pgfqpoint{0.967184in}{2.259137in}}%
\pgfpathlineto{\pgfqpoint{0.954484in}{2.262189in}}%
\pgfpathlineto{\pgfqpoint{0.957430in}{2.274082in}}%
\pgfpathlineto{\pgfqpoint{1.007808in}{2.262053in}}%
\pgfpathclose%
\pgfusepath{fill}%
\end{pgfscope}%
\begin{pgfscope}%
\pgfpathrectangle{\pgfqpoint{0.100000in}{0.100000in}}{\pgfqpoint{3.608454in}{2.310000in}}%
\pgfusepath{clip}%
\pgfsetbuttcap%
\pgfsetmiterjoin%
\definecolor{currentfill}{rgb}{0.000000,0.537255,0.731373}%
\pgfsetfillcolor{currentfill}%
\pgfsetlinewidth{0.000000pt}%
\definecolor{currentstroke}{rgb}{0.000000,0.000000,0.000000}%
\pgfsetstrokecolor{currentstroke}%
\pgfsetstrokeopacity{0.000000}%
\pgfsetdash{}{0pt}%
\pgfpathmoveto{\pgfqpoint{2.799445in}{1.624219in}}%
\pgfpathlineto{\pgfqpoint{2.792334in}{1.628535in}}%
\pgfpathlineto{\pgfqpoint{2.764552in}{1.624067in}}%
\pgfpathlineto{\pgfqpoint{2.743822in}{1.620714in}}%
\pgfpathlineto{\pgfqpoint{2.739935in}{1.648352in}}%
\pgfpathlineto{\pgfqpoint{2.735608in}{1.675908in}}%
\pgfpathlineto{\pgfqpoint{2.723746in}{1.674452in}}%
\pgfpathlineto{\pgfqpoint{2.720350in}{1.701323in}}%
\pgfpathlineto{\pgfqpoint{2.720255in}{1.702142in}}%
\pgfpathlineto{\pgfqpoint{2.745404in}{1.705512in}}%
\pgfpathlineto{\pgfqpoint{2.744382in}{1.712367in}}%
\pgfpathlineto{\pgfqpoint{2.771368in}{1.715974in}}%
\pgfpathlineto{\pgfqpoint{2.777478in}{1.717721in}}%
\pgfpathlineto{\pgfqpoint{2.776445in}{1.724555in}}%
\pgfpathlineto{\pgfqpoint{2.790730in}{1.726819in}}%
\pgfpathlineto{\pgfqpoint{2.786333in}{1.754076in}}%
\pgfpathlineto{\pgfqpoint{2.815217in}{1.759142in}}%
\pgfpathlineto{\pgfqpoint{2.817633in}{1.749127in}}%
\pgfpathlineto{\pgfqpoint{2.821949in}{1.738559in}}%
\pgfpathlineto{\pgfqpoint{2.825537in}{1.722337in}}%
\pgfpathlineto{\pgfqpoint{2.828453in}{1.713989in}}%
\pgfpathlineto{\pgfqpoint{2.833465in}{1.706759in}}%
\pgfpathlineto{\pgfqpoint{2.831898in}{1.690908in}}%
\pgfpathlineto{\pgfqpoint{2.833245in}{1.687359in}}%
\pgfpathlineto{\pgfqpoint{2.831702in}{1.674798in}}%
\pgfpathlineto{\pgfqpoint{2.820874in}{1.679961in}}%
\pgfpathlineto{\pgfqpoint{2.815658in}{1.676414in}}%
\pgfpathlineto{\pgfqpoint{2.812091in}{1.663621in}}%
\pgfpathlineto{\pgfqpoint{2.813316in}{1.659807in}}%
\pgfpathlineto{\pgfqpoint{2.810815in}{1.651423in}}%
\pgfpathlineto{\pgfqpoint{2.803471in}{1.647833in}}%
\pgfpathlineto{\pgfqpoint{2.800576in}{1.641151in}}%
\pgfpathlineto{\pgfqpoint{2.801791in}{1.629277in}}%
\pgfpathlineto{\pgfqpoint{2.799445in}{1.624219in}}%
\pgfpathclose%
\pgfusepath{fill}%
\end{pgfscope}%
\begin{pgfscope}%
\pgfpathrectangle{\pgfqpoint{0.100000in}{0.100000in}}{\pgfqpoint{3.608454in}{2.310000in}}%
\pgfusepath{clip}%
\pgfsetbuttcap%
\pgfsetmiterjoin%
\definecolor{currentfill}{rgb}{0.000000,0.603922,0.698039}%
\pgfsetfillcolor{currentfill}%
\pgfsetlinewidth{0.000000pt}%
\definecolor{currentstroke}{rgb}{0.000000,0.000000,0.000000}%
\pgfsetstrokecolor{currentstroke}%
\pgfsetstrokeopacity{0.000000}%
\pgfsetdash{}{0pt}%
\pgfpathmoveto{\pgfqpoint{0.763231in}{2.326563in}}%
\pgfpathlineto{\pgfqpoint{0.765166in}{2.300496in}}%
\pgfpathlineto{\pgfqpoint{0.760555in}{2.298682in}}%
\pgfpathlineto{\pgfqpoint{0.752119in}{2.299760in}}%
\pgfpathlineto{\pgfqpoint{0.728782in}{2.306722in}}%
\pgfpathlineto{\pgfqpoint{0.672017in}{2.324335in}}%
\pgfpathlineto{\pgfqpoint{0.672629in}{2.335131in}}%
\pgfpathlineto{\pgfqpoint{0.667385in}{2.333960in}}%
\pgfpathlineto{\pgfqpoint{0.662898in}{2.347488in}}%
\pgfpathlineto{\pgfqpoint{0.667167in}{2.355917in}}%
\pgfpathlineto{\pgfqpoint{0.717719in}{2.339867in}}%
\pgfpathlineto{\pgfqpoint{0.763231in}{2.326563in}}%
\pgfpathclose%
\pgfusepath{fill}%
\end{pgfscope}%
\begin{pgfscope}%
\pgfpathrectangle{\pgfqpoint{0.100000in}{0.100000in}}{\pgfqpoint{3.608454in}{2.310000in}}%
\pgfusepath{clip}%
\pgfsetbuttcap%
\pgfsetmiterjoin%
\definecolor{currentfill}{rgb}{0.000000,0.984314,0.507843}%
\pgfsetfillcolor{currentfill}%
\pgfsetlinewidth{0.000000pt}%
\definecolor{currentstroke}{rgb}{0.000000,0.000000,0.000000}%
\pgfsetstrokecolor{currentstroke}%
\pgfsetstrokeopacity{0.000000}%
\pgfsetdash{}{0pt}%
\pgfpathmoveto{\pgfqpoint{1.223167in}{1.249876in}}%
\pgfpathlineto{\pgfqpoint{1.210980in}{1.171960in}}%
\pgfpathlineto{\pgfqpoint{1.300870in}{1.158512in}}%
\pgfpathlineto{\pgfqpoint{1.321043in}{1.155597in}}%
\pgfpathlineto{\pgfqpoint{1.313512in}{1.101203in}}%
\pgfpathlineto{\pgfqpoint{1.239395in}{1.111911in}}%
\pgfpathlineto{\pgfqpoint{1.235282in}{1.084666in}}%
\pgfpathlineto{\pgfqpoint{1.198276in}{1.090371in}}%
\pgfpathlineto{\pgfqpoint{1.193614in}{1.060581in}}%
\pgfpathlineto{\pgfqpoint{1.183780in}{0.997802in}}%
\pgfpathlineto{\pgfqpoint{1.164171in}{1.000833in}}%
\pgfpathlineto{\pgfqpoint{1.154422in}{0.995185in}}%
\pgfpathlineto{\pgfqpoint{1.150108in}{0.990389in}}%
\pgfpathlineto{\pgfqpoint{1.145330in}{0.990568in}}%
\pgfpathlineto{\pgfqpoint{1.133001in}{0.981587in}}%
\pgfpathlineto{\pgfqpoint{1.125550in}{0.990854in}}%
\pgfpathlineto{\pgfqpoint{1.119039in}{0.992055in}}%
\pgfpathlineto{\pgfqpoint{1.124557in}{1.025178in}}%
\pgfpathlineto{\pgfqpoint{1.075926in}{1.033485in}}%
\pgfpathlineto{\pgfqpoint{1.079528in}{1.054057in}}%
\pgfpathlineto{\pgfqpoint{1.116943in}{1.267777in}}%
\pgfpathlineto{\pgfqpoint{1.163571in}{1.259392in}}%
\pgfpathlineto{\pgfqpoint{1.223167in}{1.249876in}}%
\pgfpathclose%
\pgfusepath{fill}%
\end{pgfscope}%
\begin{pgfscope}%
\pgfpathrectangle{\pgfqpoint{0.100000in}{0.100000in}}{\pgfqpoint{3.608454in}{2.310000in}}%
\pgfusepath{clip}%
\pgfsetbuttcap%
\pgfsetmiterjoin%
\definecolor{currentfill}{rgb}{0.000000,0.541176,0.729412}%
\pgfsetfillcolor{currentfill}%
\pgfsetlinewidth{0.000000pt}%
\definecolor{currentstroke}{rgb}{0.000000,0.000000,0.000000}%
\pgfsetstrokecolor{currentstroke}%
\pgfsetstrokeopacity{0.000000}%
\pgfsetdash{}{0pt}%
\pgfpathmoveto{\pgfqpoint{2.269571in}{1.909501in}}%
\pgfpathlineto{\pgfqpoint{2.297058in}{1.910416in}}%
\pgfpathlineto{\pgfqpoint{2.331620in}{1.911661in}}%
\pgfpathlineto{\pgfqpoint{2.332207in}{1.897841in}}%
\pgfpathlineto{\pgfqpoint{2.345925in}{1.898515in}}%
\pgfpathlineto{\pgfqpoint{2.348110in}{1.850512in}}%
\pgfpathlineto{\pgfqpoint{2.334351in}{1.850010in}}%
\pgfpathlineto{\pgfqpoint{2.334657in}{1.843071in}}%
\pgfpathlineto{\pgfqpoint{2.300225in}{1.841636in}}%
\pgfpathlineto{\pgfqpoint{2.300433in}{1.834842in}}%
\pgfpathlineto{\pgfqpoint{2.266038in}{1.833884in}}%
\pgfpathlineto{\pgfqpoint{2.265057in}{1.868094in}}%
\pgfpathlineto{\pgfqpoint{2.271935in}{1.868329in}}%
\pgfpathlineto{\pgfqpoint{2.270918in}{1.895690in}}%
\pgfpathlineto{\pgfqpoint{2.269571in}{1.909501in}}%
\pgfpathclose%
\pgfusepath{fill}%
\end{pgfscope}%
\begin{pgfscope}%
\pgfpathrectangle{\pgfqpoint{0.100000in}{0.100000in}}{\pgfqpoint{3.608454in}{2.310000in}}%
\pgfusepath{clip}%
\pgfsetbuttcap%
\pgfsetmiterjoin%
\definecolor{currentfill}{rgb}{0.000000,0.674510,0.662745}%
\pgfsetfillcolor{currentfill}%
\pgfsetlinewidth{0.000000pt}%
\definecolor{currentstroke}{rgb}{0.000000,0.000000,0.000000}%
\pgfsetstrokecolor{currentstroke}%
\pgfsetstrokeopacity{0.000000}%
\pgfsetdash{}{0pt}%
\pgfpathmoveto{\pgfqpoint{1.724900in}{0.803670in}}%
\pgfpathlineto{\pgfqpoint{1.690000in}{0.806092in}}%
\pgfpathlineto{\pgfqpoint{1.692995in}{0.840738in}}%
\pgfpathlineto{\pgfqpoint{1.695447in}{0.875598in}}%
\pgfpathlineto{\pgfqpoint{1.704446in}{0.875520in}}%
\pgfpathlineto{\pgfqpoint{1.706689in}{0.909371in}}%
\pgfpathlineto{\pgfqpoint{1.741190in}{0.907227in}}%
\pgfpathlineto{\pgfqpoint{1.738977in}{0.872740in}}%
\pgfpathlineto{\pgfqpoint{1.729883in}{0.873343in}}%
\pgfpathlineto{\pgfqpoint{1.724900in}{0.803670in}}%
\pgfpathclose%
\pgfusepath{fill}%
\end{pgfscope}%
\begin{pgfscope}%
\pgfpathrectangle{\pgfqpoint{0.100000in}{0.100000in}}{\pgfqpoint{3.608454in}{2.310000in}}%
\pgfusepath{clip}%
\pgfsetbuttcap%
\pgfsetmiterjoin%
\definecolor{currentfill}{rgb}{0.000000,0.827451,0.586275}%
\pgfsetfillcolor{currentfill}%
\pgfsetlinewidth{0.000000pt}%
\definecolor{currentstroke}{rgb}{0.000000,0.000000,0.000000}%
\pgfsetstrokecolor{currentstroke}%
\pgfsetstrokeopacity{0.000000}%
\pgfsetdash{}{0pt}%
\pgfpathmoveto{\pgfqpoint{2.474841in}{0.898468in}}%
\pgfpathlineto{\pgfqpoint{2.483949in}{0.899006in}}%
\pgfpathlineto{\pgfqpoint{2.485523in}{0.870923in}}%
\pgfpathlineto{\pgfqpoint{2.458164in}{0.869135in}}%
\pgfpathlineto{\pgfqpoint{2.458381in}{0.865552in}}%
\pgfpathlineto{\pgfqpoint{2.442684in}{0.864247in}}%
\pgfpathlineto{\pgfqpoint{2.417323in}{0.873285in}}%
\pgfpathlineto{\pgfqpoint{2.410750in}{0.879777in}}%
\pgfpathlineto{\pgfqpoint{2.417347in}{0.883236in}}%
\pgfpathlineto{\pgfqpoint{2.416578in}{0.896534in}}%
\pgfpathlineto{\pgfqpoint{2.419871in}{0.899622in}}%
\pgfpathlineto{\pgfqpoint{2.426856in}{0.899501in}}%
\pgfpathlineto{\pgfqpoint{2.455570in}{0.891636in}}%
\pgfpathlineto{\pgfqpoint{2.458419in}{0.897453in}}%
\pgfpathlineto{\pgfqpoint{2.474841in}{0.898468in}}%
\pgfpathclose%
\pgfusepath{fill}%
\end{pgfscope}%
\begin{pgfscope}%
\pgfpathrectangle{\pgfqpoint{0.100000in}{0.100000in}}{\pgfqpoint{3.608454in}{2.310000in}}%
\pgfusepath{clip}%
\pgfsetbuttcap%
\pgfsetmiterjoin%
\definecolor{currentfill}{rgb}{0.000000,0.552941,0.723529}%
\pgfsetfillcolor{currentfill}%
\pgfsetlinewidth{0.000000pt}%
\definecolor{currentstroke}{rgb}{0.000000,0.000000,0.000000}%
\pgfsetstrokecolor{currentstroke}%
\pgfsetstrokeopacity{0.000000}%
\pgfsetdash{}{0pt}%
\pgfpathmoveto{\pgfqpoint{2.348213in}{1.508885in}}%
\pgfpathlineto{\pgfqpoint{2.323109in}{1.507938in}}%
\pgfpathlineto{\pgfqpoint{2.322856in}{1.514824in}}%
\pgfpathlineto{\pgfqpoint{2.302293in}{1.514165in}}%
\pgfpathlineto{\pgfqpoint{2.288611in}{1.513740in}}%
\pgfpathlineto{\pgfqpoint{2.287673in}{1.541193in}}%
\pgfpathlineto{\pgfqpoint{2.294578in}{1.541485in}}%
\pgfpathlineto{\pgfqpoint{2.293482in}{1.569183in}}%
\pgfpathlineto{\pgfqpoint{2.320902in}{1.570069in}}%
\pgfpathlineto{\pgfqpoint{2.320640in}{1.576968in}}%
\pgfpathlineto{\pgfqpoint{2.348123in}{1.578046in}}%
\pgfpathlineto{\pgfqpoint{2.348758in}{1.564236in}}%
\pgfpathlineto{\pgfqpoint{2.349290in}{1.550458in}}%
\pgfpathlineto{\pgfqpoint{2.356154in}{1.550723in}}%
\pgfpathlineto{\pgfqpoint{2.356534in}{1.539366in}}%
\pgfpathlineto{\pgfqpoint{2.341290in}{1.535608in}}%
\pgfpathlineto{\pgfqpoint{2.337899in}{1.521676in}}%
\pgfpathlineto{\pgfqpoint{2.345525in}{1.515337in}}%
\pgfpathlineto{\pgfqpoint{2.348213in}{1.508885in}}%
\pgfpathclose%
\pgfusepath{fill}%
\end{pgfscope}%
\begin{pgfscope}%
\pgfpathrectangle{\pgfqpoint{0.100000in}{0.100000in}}{\pgfqpoint{3.608454in}{2.310000in}}%
\pgfusepath{clip}%
\pgfsetbuttcap%
\pgfsetmiterjoin%
\definecolor{currentfill}{rgb}{0.000000,0.513725,0.743137}%
\pgfsetfillcolor{currentfill}%
\pgfsetlinewidth{0.000000pt}%
\definecolor{currentstroke}{rgb}{0.000000,0.000000,0.000000}%
\pgfsetstrokecolor{currentstroke}%
\pgfsetstrokeopacity{0.000000}%
\pgfsetdash{}{0pt}%
\pgfpathmoveto{\pgfqpoint{2.099037in}{1.876260in}}%
\pgfpathlineto{\pgfqpoint{2.099439in}{1.847533in}}%
\pgfpathlineto{\pgfqpoint{2.044431in}{1.847876in}}%
\pgfpathlineto{\pgfqpoint{2.037488in}{1.848022in}}%
\pgfpathlineto{\pgfqpoint{2.037019in}{1.861775in}}%
\pgfpathlineto{\pgfqpoint{2.037372in}{1.889582in}}%
\pgfpathlineto{\pgfqpoint{2.036803in}{1.896501in}}%
\pgfpathlineto{\pgfqpoint{2.036301in}{1.917281in}}%
\pgfpathlineto{\pgfqpoint{2.036642in}{1.945064in}}%
\pgfpathlineto{\pgfqpoint{2.042594in}{1.944974in}}%
\pgfpathlineto{\pgfqpoint{2.042559in}{1.951936in}}%
\pgfpathlineto{\pgfqpoint{2.097851in}{1.951504in}}%
\pgfpathlineto{\pgfqpoint{2.098775in}{1.916792in}}%
\pgfpathlineto{\pgfqpoint{2.099037in}{1.876260in}}%
\pgfpathclose%
\pgfusepath{fill}%
\end{pgfscope}%
\begin{pgfscope}%
\pgfpathrectangle{\pgfqpoint{0.100000in}{0.100000in}}{\pgfqpoint{3.608454in}{2.310000in}}%
\pgfusepath{clip}%
\pgfsetbuttcap%
\pgfsetmiterjoin%
\definecolor{currentfill}{rgb}{0.000000,0.549020,0.725490}%
\pgfsetfillcolor{currentfill}%
\pgfsetlinewidth{0.000000pt}%
\definecolor{currentstroke}{rgb}{0.000000,0.000000,0.000000}%
\pgfsetstrokecolor{currentstroke}%
\pgfsetstrokeopacity{0.000000}%
\pgfsetdash{}{0pt}%
\pgfpathmoveto{\pgfqpoint{2.817692in}{1.389496in}}%
\pgfpathlineto{\pgfqpoint{2.804293in}{1.384898in}}%
\pgfpathlineto{\pgfqpoint{2.788388in}{1.383110in}}%
\pgfpathlineto{\pgfqpoint{2.786603in}{1.400891in}}%
\pgfpathlineto{\pgfqpoint{2.777953in}{1.400497in}}%
\pgfpathlineto{\pgfqpoint{2.776762in}{1.425313in}}%
\pgfpathlineto{\pgfqpoint{2.795500in}{1.426123in}}%
\pgfpathlineto{\pgfqpoint{2.794880in}{1.439319in}}%
\pgfpathlineto{\pgfqpoint{2.819277in}{1.440727in}}%
\pgfpathlineto{\pgfqpoint{2.820238in}{1.426570in}}%
\pgfpathlineto{\pgfqpoint{2.815226in}{1.414852in}}%
\pgfpathlineto{\pgfqpoint{2.816007in}{1.404773in}}%
\pgfpathlineto{\pgfqpoint{2.818458in}{1.403774in}}%
\pgfpathlineto{\pgfqpoint{2.817692in}{1.389496in}}%
\pgfpathclose%
\pgfusepath{fill}%
\end{pgfscope}%
\begin{pgfscope}%
\pgfpathrectangle{\pgfqpoint{0.100000in}{0.100000in}}{\pgfqpoint{3.608454in}{2.310000in}}%
\pgfusepath{clip}%
\pgfsetbuttcap%
\pgfsetmiterjoin%
\definecolor{currentfill}{rgb}{0.000000,0.670588,0.664706}%
\pgfsetfillcolor{currentfill}%
\pgfsetlinewidth{0.000000pt}%
\definecolor{currentstroke}{rgb}{0.000000,0.000000,0.000000}%
\pgfsetstrokecolor{currentstroke}%
\pgfsetstrokeopacity{0.000000}%
\pgfsetdash{}{0pt}%
\pgfpathmoveto{\pgfqpoint{3.596315in}{2.093414in}}%
\pgfpathlineto{\pgfqpoint{3.619860in}{2.089848in}}%
\pgfpathlineto{\pgfqpoint{3.617148in}{2.084070in}}%
\pgfpathlineto{\pgfqpoint{3.624313in}{2.074071in}}%
\pgfpathlineto{\pgfqpoint{3.623386in}{2.068081in}}%
\pgfpathlineto{\pgfqpoint{3.634273in}{2.058968in}}%
\pgfpathlineto{\pgfqpoint{3.636385in}{2.065102in}}%
\pgfpathlineto{\pgfqpoint{3.643097in}{2.064783in}}%
\pgfpathlineto{\pgfqpoint{3.657645in}{2.048872in}}%
\pgfpathlineto{\pgfqpoint{3.660907in}{2.041358in}}%
\pgfpathlineto{\pgfqpoint{3.656677in}{2.034155in}}%
\pgfpathlineto{\pgfqpoint{3.653836in}{2.025053in}}%
\pgfpathlineto{\pgfqpoint{3.642794in}{2.024053in}}%
\pgfpathlineto{\pgfqpoint{3.643360in}{2.017403in}}%
\pgfpathlineto{\pgfqpoint{3.634988in}{2.017176in}}%
\pgfpathlineto{\pgfqpoint{3.633250in}{2.008704in}}%
\pgfpathlineto{\pgfqpoint{3.624318in}{2.007131in}}%
\pgfpathlineto{\pgfqpoint{3.621576in}{1.995429in}}%
\pgfpathlineto{\pgfqpoint{3.618035in}{1.994535in}}%
\pgfpathlineto{\pgfqpoint{3.610482in}{2.006964in}}%
\pgfpathlineto{\pgfqpoint{3.597567in}{2.032110in}}%
\pgfpathlineto{\pgfqpoint{3.603808in}{2.035391in}}%
\pgfpathlineto{\pgfqpoint{3.592928in}{2.056963in}}%
\pgfpathlineto{\pgfqpoint{3.598893in}{2.059799in}}%
\pgfpathlineto{\pgfqpoint{3.584395in}{2.086358in}}%
\pgfpathlineto{\pgfqpoint{3.596315in}{2.093414in}}%
\pgfpathclose%
\pgfusepath{fill}%
\end{pgfscope}%
\begin{pgfscope}%
\pgfpathrectangle{\pgfqpoint{0.100000in}{0.100000in}}{\pgfqpoint{3.608454in}{2.310000in}}%
\pgfusepath{clip}%
\pgfsetbuttcap%
\pgfsetmiterjoin%
\definecolor{currentfill}{rgb}{0.000000,0.529412,0.735294}%
\pgfsetfillcolor{currentfill}%
\pgfsetlinewidth{0.000000pt}%
\definecolor{currentstroke}{rgb}{0.000000,0.000000,0.000000}%
\pgfsetstrokecolor{currentstroke}%
\pgfsetstrokeopacity{0.000000}%
\pgfsetdash{}{0pt}%
\pgfpathmoveto{\pgfqpoint{2.833971in}{1.586462in}}%
\pgfpathlineto{\pgfqpoint{2.818905in}{1.581252in}}%
\pgfpathlineto{\pgfqpoint{2.796288in}{1.578242in}}%
\pgfpathlineto{\pgfqpoint{2.791356in}{1.581015in}}%
\pgfpathlineto{\pgfqpoint{2.794328in}{1.554792in}}%
\pgfpathlineto{\pgfqpoint{2.767138in}{1.551464in}}%
\pgfpathlineto{\pgfqpoint{2.746637in}{1.548837in}}%
\pgfpathlineto{\pgfqpoint{2.744178in}{1.569410in}}%
\pgfpathlineto{\pgfqpoint{2.737491in}{1.568598in}}%
\pgfpathlineto{\pgfqpoint{2.734410in}{1.575123in}}%
\pgfpathlineto{\pgfqpoint{2.731544in}{1.590059in}}%
\pgfpathlineto{\pgfqpoint{2.768691in}{1.595961in}}%
\pgfpathlineto{\pgfqpoint{2.764552in}{1.624067in}}%
\pgfpathlineto{\pgfqpoint{2.792334in}{1.628535in}}%
\pgfpathlineto{\pgfqpoint{2.799445in}{1.624219in}}%
\pgfpathlineto{\pgfqpoint{2.792747in}{1.616003in}}%
\pgfpathlineto{\pgfqpoint{2.786784in}{1.605434in}}%
\pgfpathlineto{\pgfqpoint{2.786107in}{1.596021in}}%
\pgfpathlineto{\pgfqpoint{2.793257in}{1.596373in}}%
\pgfpathlineto{\pgfqpoint{2.810519in}{1.591373in}}%
\pgfpathlineto{\pgfqpoint{2.817182in}{1.586236in}}%
\pgfpathlineto{\pgfqpoint{2.831148in}{1.589793in}}%
\pgfpathlineto{\pgfqpoint{2.833971in}{1.586462in}}%
\pgfpathclose%
\pgfusepath{fill}%
\end{pgfscope}%
\begin{pgfscope}%
\pgfpathrectangle{\pgfqpoint{0.100000in}{0.100000in}}{\pgfqpoint{3.608454in}{2.310000in}}%
\pgfusepath{clip}%
\pgfsetbuttcap%
\pgfsetmiterjoin%
\definecolor{currentfill}{rgb}{0.000000,0.556863,0.721569}%
\pgfsetfillcolor{currentfill}%
\pgfsetlinewidth{0.000000pt}%
\definecolor{currentstroke}{rgb}{0.000000,0.000000,0.000000}%
\pgfsetstrokecolor{currentstroke}%
\pgfsetstrokeopacity{0.000000}%
\pgfsetdash{}{0pt}%
\pgfpathmoveto{\pgfqpoint{2.380907in}{1.900334in}}%
\pgfpathlineto{\pgfqpoint{2.382793in}{1.866283in}}%
\pgfpathlineto{\pgfqpoint{2.383604in}{1.852652in}}%
\pgfpathlineto{\pgfqpoint{2.348110in}{1.850512in}}%
\pgfpathlineto{\pgfqpoint{2.345925in}{1.898515in}}%
\pgfpathlineto{\pgfqpoint{2.332207in}{1.897841in}}%
\pgfpathlineto{\pgfqpoint{2.331620in}{1.911661in}}%
\pgfpathlineto{\pgfqpoint{2.297058in}{1.910416in}}%
\pgfpathlineto{\pgfqpoint{2.295263in}{1.958166in}}%
\pgfpathlineto{\pgfqpoint{2.305024in}{1.961414in}}%
\pgfpathlineto{\pgfqpoint{2.313626in}{1.967707in}}%
\pgfpathlineto{\pgfqpoint{2.319119in}{1.967177in}}%
\pgfpathlineto{\pgfqpoint{2.326155in}{1.974190in}}%
\pgfpathlineto{\pgfqpoint{2.332310in}{1.976143in}}%
\pgfpathlineto{\pgfqpoint{2.338487in}{1.970865in}}%
\pgfpathlineto{\pgfqpoint{2.331663in}{1.959426in}}%
\pgfpathlineto{\pgfqpoint{2.331044in}{1.945780in}}%
\pgfpathlineto{\pgfqpoint{2.342366in}{1.952806in}}%
\pgfpathlineto{\pgfqpoint{2.350608in}{1.946791in}}%
\pgfpathlineto{\pgfqpoint{2.351813in}{1.919479in}}%
\pgfpathlineto{\pgfqpoint{2.358761in}{1.919792in}}%
\pgfpathlineto{\pgfqpoint{2.359091in}{1.912916in}}%
\pgfpathlineto{\pgfqpoint{2.365878in}{1.913327in}}%
\pgfpathlineto{\pgfqpoint{2.366647in}{1.899528in}}%
\pgfpathlineto{\pgfqpoint{2.380907in}{1.900334in}}%
\pgfpathclose%
\pgfusepath{fill}%
\end{pgfscope}%
\begin{pgfscope}%
\pgfpathrectangle{\pgfqpoint{0.100000in}{0.100000in}}{\pgfqpoint{3.608454in}{2.310000in}}%
\pgfusepath{clip}%
\pgfsetbuttcap%
\pgfsetmiterjoin%
\definecolor{currentfill}{rgb}{0.000000,0.611765,0.694118}%
\pgfsetfillcolor{currentfill}%
\pgfsetlinewidth{0.000000pt}%
\definecolor{currentstroke}{rgb}{0.000000,0.000000,0.000000}%
\pgfsetstrokecolor{currentstroke}%
\pgfsetstrokeopacity{0.000000}%
\pgfsetdash{}{0pt}%
\pgfpathmoveto{\pgfqpoint{2.634617in}{1.461203in}}%
\pgfpathlineto{\pgfqpoint{2.634909in}{1.458446in}}%
\pgfpathlineto{\pgfqpoint{2.607655in}{1.455712in}}%
\pgfpathlineto{\pgfqpoint{2.607399in}{1.458524in}}%
\pgfpathlineto{\pgfqpoint{2.583559in}{1.456426in}}%
\pgfpathlineto{\pgfqpoint{2.585430in}{1.435815in}}%
\pgfpathlineto{\pgfqpoint{2.565557in}{1.434073in}}%
\pgfpathlineto{\pgfqpoint{2.565221in}{1.447932in}}%
\pgfpathlineto{\pgfqpoint{2.557572in}{1.448924in}}%
\pgfpathlineto{\pgfqpoint{2.559183in}{1.428045in}}%
\pgfpathlineto{\pgfqpoint{2.534753in}{1.425828in}}%
\pgfpathlineto{\pgfqpoint{2.532347in}{1.453054in}}%
\pgfpathlineto{\pgfqpoint{2.531157in}{1.473630in}}%
\pgfpathlineto{\pgfqpoint{2.555646in}{1.476018in}}%
\pgfpathlineto{\pgfqpoint{2.554067in}{1.495355in}}%
\pgfpathlineto{\pgfqpoint{2.579594in}{1.497487in}}%
\pgfpathlineto{\pgfqpoint{2.578828in}{1.505444in}}%
\pgfpathlineto{\pgfqpoint{2.585581in}{1.506060in}}%
\pgfpathlineto{\pgfqpoint{2.588380in}{1.512229in}}%
\pgfpathlineto{\pgfqpoint{2.609172in}{1.513996in}}%
\pgfpathlineto{\pgfqpoint{2.610441in}{1.500202in}}%
\pgfpathlineto{\pgfqpoint{2.614139in}{1.497101in}}%
\pgfpathlineto{\pgfqpoint{2.623185in}{1.497972in}}%
\pgfpathlineto{\pgfqpoint{2.625004in}{1.477436in}}%
\pgfpathlineto{\pgfqpoint{2.633399in}{1.473636in}}%
\pgfpathlineto{\pgfqpoint{2.634617in}{1.461203in}}%
\pgfpathclose%
\pgfusepath{fill}%
\end{pgfscope}%
\begin{pgfscope}%
\pgfpathrectangle{\pgfqpoint{0.100000in}{0.100000in}}{\pgfqpoint{3.608454in}{2.310000in}}%
\pgfusepath{clip}%
\pgfsetbuttcap%
\pgfsetmiterjoin%
\definecolor{currentfill}{rgb}{0.000000,0.788235,0.605882}%
\pgfsetfillcolor{currentfill}%
\pgfsetlinewidth{0.000000pt}%
\definecolor{currentstroke}{rgb}{0.000000,0.000000,0.000000}%
\pgfsetstrokecolor{currentstroke}%
\pgfsetstrokeopacity{0.000000}%
\pgfsetdash{}{0pt}%
\pgfpathmoveto{\pgfqpoint{2.202317in}{0.720035in}}%
\pgfpathlineto{\pgfqpoint{2.181123in}{0.718301in}}%
\pgfpathlineto{\pgfqpoint{2.166253in}{0.713450in}}%
\pgfpathlineto{\pgfqpoint{2.152123in}{0.723725in}}%
\pgfpathlineto{\pgfqpoint{2.145077in}{0.730301in}}%
\pgfpathlineto{\pgfqpoint{2.140576in}{0.738107in}}%
\pgfpathlineto{\pgfqpoint{2.125247in}{0.741762in}}%
\pgfpathlineto{\pgfqpoint{2.115910in}{0.747279in}}%
\pgfpathlineto{\pgfqpoint{2.111897in}{0.751748in}}%
\pgfpathlineto{\pgfqpoint{2.108517in}{0.767596in}}%
\pgfpathlineto{\pgfqpoint{2.111050in}{0.772661in}}%
\pgfpathlineto{\pgfqpoint{2.143853in}{0.772670in}}%
\pgfpathlineto{\pgfqpoint{2.139795in}{0.782977in}}%
\pgfpathlineto{\pgfqpoint{2.173280in}{0.783716in}}%
\pgfpathlineto{\pgfqpoint{2.182637in}{0.773027in}}%
\pgfpathlineto{\pgfqpoint{2.186530in}{0.767441in}}%
\pgfpathlineto{\pgfqpoint{2.187852in}{0.760329in}}%
\pgfpathlineto{\pgfqpoint{2.185571in}{0.753524in}}%
\pgfpathlineto{\pgfqpoint{2.191816in}{0.743213in}}%
\pgfpathlineto{\pgfqpoint{2.197023in}{0.737625in}}%
\pgfpathlineto{\pgfqpoint{2.196211in}{0.730685in}}%
\pgfpathlineto{\pgfqpoint{2.202317in}{0.720035in}}%
\pgfpathclose%
\pgfusepath{fill}%
\end{pgfscope}%
\begin{pgfscope}%
\pgfpathrectangle{\pgfqpoint{0.100000in}{0.100000in}}{\pgfqpoint{3.608454in}{2.310000in}}%
\pgfusepath{clip}%
\pgfsetbuttcap%
\pgfsetmiterjoin%
\definecolor{currentfill}{rgb}{0.000000,0.725490,0.637255}%
\pgfsetfillcolor{currentfill}%
\pgfsetlinewidth{0.000000pt}%
\definecolor{currentstroke}{rgb}{0.000000,0.000000,0.000000}%
\pgfsetstrokecolor{currentstroke}%
\pgfsetstrokeopacity{0.000000}%
\pgfsetdash{}{0pt}%
\pgfpathmoveto{\pgfqpoint{3.147047in}{0.962987in}}%
\pgfpathlineto{\pgfqpoint{3.134423in}{0.971285in}}%
\pgfpathlineto{\pgfqpoint{3.115214in}{0.974844in}}%
\pgfpathlineto{\pgfqpoint{3.097264in}{0.986382in}}%
\pgfpathlineto{\pgfqpoint{3.088906in}{0.984668in}}%
\pgfpathlineto{\pgfqpoint{3.094287in}{1.003646in}}%
\pgfpathlineto{\pgfqpoint{3.092298in}{1.007042in}}%
\pgfpathlineto{\pgfqpoint{3.096860in}{1.011889in}}%
\pgfpathlineto{\pgfqpoint{3.091334in}{1.021253in}}%
\pgfpathlineto{\pgfqpoint{3.096046in}{1.025469in}}%
\pgfpathlineto{\pgfqpoint{3.076612in}{1.036110in}}%
\pgfpathlineto{\pgfqpoint{3.078779in}{1.041854in}}%
\pgfpathlineto{\pgfqpoint{3.073425in}{1.048823in}}%
\pgfpathlineto{\pgfqpoint{3.065576in}{1.050744in}}%
\pgfpathlineto{\pgfqpoint{3.079593in}{1.063190in}}%
\pgfpathlineto{\pgfqpoint{3.095167in}{1.065172in}}%
\pgfpathlineto{\pgfqpoint{3.108787in}{1.052476in}}%
\pgfpathlineto{\pgfqpoint{3.116266in}{1.079750in}}%
\pgfpathlineto{\pgfqpoint{3.145061in}{1.058883in}}%
\pgfpathlineto{\pgfqpoint{3.185535in}{1.030022in}}%
\pgfpathlineto{\pgfqpoint{3.175718in}{1.024509in}}%
\pgfpathlineto{\pgfqpoint{3.169582in}{1.018517in}}%
\pgfpathlineto{\pgfqpoint{3.162304in}{1.007975in}}%
\pgfpathlineto{\pgfqpoint{3.155468in}{0.994985in}}%
\pgfpathlineto{\pgfqpoint{3.152816in}{0.985825in}}%
\pgfpathlineto{\pgfqpoint{3.152996in}{0.969508in}}%
\pgfpathlineto{\pgfqpoint{3.147047in}{0.962987in}}%
\pgfpathclose%
\pgfusepath{fill}%
\end{pgfscope}%
\begin{pgfscope}%
\pgfpathrectangle{\pgfqpoint{0.100000in}{0.100000in}}{\pgfqpoint{3.608454in}{2.310000in}}%
\pgfusepath{clip}%
\pgfsetbuttcap%
\pgfsetmiterjoin%
\definecolor{currentfill}{rgb}{0.000000,0.796078,0.601961}%
\pgfsetfillcolor{currentfill}%
\pgfsetlinewidth{0.000000pt}%
\definecolor{currentstroke}{rgb}{0.000000,0.000000,0.000000}%
\pgfsetstrokecolor{currentstroke}%
\pgfsetstrokeopacity{0.000000}%
\pgfsetdash{}{0pt}%
\pgfpathmoveto{\pgfqpoint{2.611522in}{1.917144in}}%
\pgfpathlineto{\pgfqpoint{2.605875in}{1.974408in}}%
\pgfpathlineto{\pgfqpoint{2.625687in}{1.975346in}}%
\pgfpathlineto{\pgfqpoint{2.636425in}{1.981889in}}%
\pgfpathlineto{\pgfqpoint{2.653804in}{1.986050in}}%
\pgfpathlineto{\pgfqpoint{2.651220in}{1.978697in}}%
\pgfpathlineto{\pgfqpoint{2.652749in}{1.962987in}}%
\pgfpathlineto{\pgfqpoint{2.664373in}{1.960858in}}%
\pgfpathlineto{\pgfqpoint{2.672255in}{1.965247in}}%
\pgfpathlineto{\pgfqpoint{2.677033in}{1.960330in}}%
\pgfpathlineto{\pgfqpoint{2.686113in}{1.967978in}}%
\pgfpathlineto{\pgfqpoint{2.702927in}{1.970218in}}%
\pgfpathlineto{\pgfqpoint{2.701836in}{1.963213in}}%
\pgfpathlineto{\pgfqpoint{2.704934in}{1.955834in}}%
\pgfpathlineto{\pgfqpoint{2.703228in}{1.945319in}}%
\pgfpathlineto{\pgfqpoint{2.711142in}{1.941318in}}%
\pgfpathlineto{\pgfqpoint{2.714091in}{1.933704in}}%
\pgfpathlineto{\pgfqpoint{2.724603in}{1.932074in}}%
\pgfpathlineto{\pgfqpoint{2.735969in}{1.940892in}}%
\pgfpathlineto{\pgfqpoint{2.742949in}{1.933510in}}%
\pgfpathlineto{\pgfqpoint{2.741240in}{1.928257in}}%
\pgfpathlineto{\pgfqpoint{2.733601in}{1.928846in}}%
\pgfpathlineto{\pgfqpoint{2.725778in}{1.927125in}}%
\pgfpathlineto{\pgfqpoint{2.720067in}{1.929651in}}%
\pgfpathlineto{\pgfqpoint{2.706856in}{1.928149in}}%
\pgfpathlineto{\pgfqpoint{2.693720in}{1.923132in}}%
\pgfpathlineto{\pgfqpoint{2.684872in}{1.925062in}}%
\pgfpathlineto{\pgfqpoint{2.682294in}{1.929947in}}%
\pgfpathlineto{\pgfqpoint{2.675578in}{1.929116in}}%
\pgfpathlineto{\pgfqpoint{2.673548in}{1.913649in}}%
\pgfpathlineto{\pgfqpoint{2.664389in}{1.919036in}}%
\pgfpathlineto{\pgfqpoint{2.658981in}{1.924851in}}%
\pgfpathlineto{\pgfqpoint{2.650299in}{1.927597in}}%
\pgfpathlineto{\pgfqpoint{2.635790in}{1.929952in}}%
\pgfpathlineto{\pgfqpoint{2.629423in}{1.928847in}}%
\pgfpathlineto{\pgfqpoint{2.622249in}{1.917995in}}%
\pgfpathlineto{\pgfqpoint{2.611522in}{1.917144in}}%
\pgfpathclose%
\pgfusepath{fill}%
\end{pgfscope}%
\begin{pgfscope}%
\pgfpathrectangle{\pgfqpoint{0.100000in}{0.100000in}}{\pgfqpoint{3.608454in}{2.310000in}}%
\pgfusepath{clip}%
\pgfsetbuttcap%
\pgfsetmiterjoin%
\definecolor{currentfill}{rgb}{0.000000,0.662745,0.668627}%
\pgfsetfillcolor{currentfill}%
\pgfsetlinewidth{0.000000pt}%
\definecolor{currentstroke}{rgb}{0.000000,0.000000,0.000000}%
\pgfsetstrokecolor{currentstroke}%
\pgfsetstrokeopacity{0.000000}%
\pgfsetdash{}{0pt}%
\pgfpathmoveto{\pgfqpoint{2.482016in}{1.066685in}}%
\pgfpathlineto{\pgfqpoint{2.463288in}{1.065984in}}%
\pgfpathlineto{\pgfqpoint{2.461856in}{1.079471in}}%
\pgfpathlineto{\pgfqpoint{2.453807in}{1.083497in}}%
\pgfpathlineto{\pgfqpoint{2.441604in}{1.083493in}}%
\pgfpathlineto{\pgfqpoint{2.434616in}{1.075007in}}%
\pgfpathlineto{\pgfqpoint{2.432045in}{1.079916in}}%
\pgfpathlineto{\pgfqpoint{2.431130in}{1.090720in}}%
\pgfpathlineto{\pgfqpoint{2.439328in}{1.092627in}}%
\pgfpathlineto{\pgfqpoint{2.442664in}{1.097005in}}%
\pgfpathlineto{\pgfqpoint{2.442276in}{1.103468in}}%
\pgfpathlineto{\pgfqpoint{2.445220in}{1.109326in}}%
\pgfpathlineto{\pgfqpoint{2.446875in}{1.118895in}}%
\pgfpathlineto{\pgfqpoint{2.452158in}{1.122751in}}%
\pgfpathlineto{\pgfqpoint{2.445412in}{1.130104in}}%
\pgfpathlineto{\pgfqpoint{2.455346in}{1.132898in}}%
\pgfpathlineto{\pgfqpoint{2.454276in}{1.139497in}}%
\pgfpathlineto{\pgfqpoint{2.453898in}{1.152266in}}%
\pgfpathlineto{\pgfqpoint{2.457328in}{1.152433in}}%
\pgfpathlineto{\pgfqpoint{2.461626in}{1.152803in}}%
\pgfpathlineto{\pgfqpoint{2.465696in}{1.145939in}}%
\pgfpathlineto{\pgfqpoint{2.462058in}{1.133851in}}%
\pgfpathlineto{\pgfqpoint{2.458843in}{1.129936in}}%
\pgfpathlineto{\pgfqpoint{2.479806in}{1.130609in}}%
\pgfpathlineto{\pgfqpoint{2.480144in}{1.116463in}}%
\pgfpathlineto{\pgfqpoint{2.469563in}{1.097602in}}%
\pgfpathlineto{\pgfqpoint{2.487860in}{1.090523in}}%
\pgfpathlineto{\pgfqpoint{2.488608in}{1.069834in}}%
\pgfpathlineto{\pgfqpoint{2.482016in}{1.066685in}}%
\pgfpathclose%
\pgfusepath{fill}%
\end{pgfscope}%
\begin{pgfscope}%
\pgfpathrectangle{\pgfqpoint{0.100000in}{0.100000in}}{\pgfqpoint{3.608454in}{2.310000in}}%
\pgfusepath{clip}%
\pgfsetbuttcap%
\pgfsetmiterjoin%
\definecolor{currentfill}{rgb}{0.000000,0.470588,0.764706}%
\pgfsetfillcolor{currentfill}%
\pgfsetlinewidth{0.000000pt}%
\definecolor{currentstroke}{rgb}{0.000000,0.000000,0.000000}%
\pgfsetstrokecolor{currentstroke}%
\pgfsetstrokeopacity{0.000000}%
\pgfsetdash{}{0pt}%
\pgfpathmoveto{\pgfqpoint{2.970938in}{1.132406in}}%
\pgfpathlineto{\pgfqpoint{2.978683in}{1.148660in}}%
\pgfpathlineto{\pgfqpoint{2.973998in}{1.150189in}}%
\pgfpathlineto{\pgfqpoint{2.967129in}{1.148531in}}%
\pgfpathlineto{\pgfqpoint{2.956204in}{1.154739in}}%
\pgfpathlineto{\pgfqpoint{2.949055in}{1.160645in}}%
\pgfpathlineto{\pgfqpoint{2.955589in}{1.174438in}}%
\pgfpathlineto{\pgfqpoint{2.963826in}{1.175254in}}%
\pgfpathlineto{\pgfqpoint{2.972912in}{1.172989in}}%
\pgfpathlineto{\pgfqpoint{2.978982in}{1.167861in}}%
\pgfpathlineto{\pgfqpoint{2.987201in}{1.172347in}}%
\pgfpathlineto{\pgfqpoint{2.993174in}{1.172573in}}%
\pgfpathlineto{\pgfqpoint{3.006498in}{1.176840in}}%
\pgfpathlineto{\pgfqpoint{3.018531in}{1.178366in}}%
\pgfpathlineto{\pgfqpoint{3.020036in}{1.162708in}}%
\pgfpathlineto{\pgfqpoint{3.018191in}{1.148162in}}%
\pgfpathlineto{\pgfqpoint{3.022456in}{1.135384in}}%
\pgfpathlineto{\pgfqpoint{3.008399in}{1.134076in}}%
\pgfpathlineto{\pgfqpoint{3.007763in}{1.136339in}}%
\pgfpathlineto{\pgfqpoint{2.970938in}{1.132406in}}%
\pgfpathclose%
\pgfusepath{fill}%
\end{pgfscope}%
\begin{pgfscope}%
\pgfpathrectangle{\pgfqpoint{0.100000in}{0.100000in}}{\pgfqpoint{3.608454in}{2.310000in}}%
\pgfusepath{clip}%
\pgfsetbuttcap%
\pgfsetmiterjoin%
\definecolor{currentfill}{rgb}{0.000000,0.517647,0.741176}%
\pgfsetfillcolor{currentfill}%
\pgfsetlinewidth{0.000000pt}%
\definecolor{currentstroke}{rgb}{0.000000,0.000000,0.000000}%
\pgfsetstrokecolor{currentstroke}%
\pgfsetstrokeopacity{0.000000}%
\pgfsetdash{}{0pt}%
\pgfpathmoveto{\pgfqpoint{2.801132in}{1.502933in}}%
\pgfpathlineto{\pgfqpoint{2.795078in}{1.502065in}}%
\pgfpathlineto{\pgfqpoint{2.773279in}{1.502104in}}%
\pgfpathlineto{\pgfqpoint{2.767138in}{1.551464in}}%
\pgfpathlineto{\pgfqpoint{2.794328in}{1.554792in}}%
\pgfpathlineto{\pgfqpoint{2.791356in}{1.581015in}}%
\pgfpathlineto{\pgfqpoint{2.796288in}{1.578242in}}%
\pgfpathlineto{\pgfqpoint{2.818905in}{1.581252in}}%
\pgfpathlineto{\pgfqpoint{2.825402in}{1.579904in}}%
\pgfpathlineto{\pgfqpoint{2.831040in}{1.546066in}}%
\pgfpathlineto{\pgfqpoint{2.814322in}{1.543583in}}%
\pgfpathlineto{\pgfqpoint{2.817395in}{1.520810in}}%
\pgfpathlineto{\pgfqpoint{2.799219in}{1.517163in}}%
\pgfpathlineto{\pgfqpoint{2.801132in}{1.502933in}}%
\pgfpathclose%
\pgfusepath{fill}%
\end{pgfscope}%
\begin{pgfscope}%
\pgfpathrectangle{\pgfqpoint{0.100000in}{0.100000in}}{\pgfqpoint{3.608454in}{2.310000in}}%
\pgfusepath{clip}%
\pgfsetbuttcap%
\pgfsetmiterjoin%
\definecolor{currentfill}{rgb}{0.000000,0.709804,0.645098}%
\pgfsetfillcolor{currentfill}%
\pgfsetlinewidth{0.000000pt}%
\definecolor{currentstroke}{rgb}{0.000000,0.000000,0.000000}%
\pgfsetstrokecolor{currentstroke}%
\pgfsetstrokeopacity{0.000000}%
\pgfsetdash{}{0pt}%
\pgfpathmoveto{\pgfqpoint{3.436777in}{1.747450in}}%
\pgfpathlineto{\pgfqpoint{3.384285in}{1.735025in}}%
\pgfpathlineto{\pgfqpoint{3.380209in}{1.761732in}}%
\pgfpathlineto{\pgfqpoint{3.378669in}{1.782710in}}%
\pgfpathlineto{\pgfqpoint{3.373418in}{1.786388in}}%
\pgfpathlineto{\pgfqpoint{3.377987in}{1.790511in}}%
\pgfpathlineto{\pgfqpoint{3.414849in}{1.798333in}}%
\pgfpathlineto{\pgfqpoint{3.419015in}{1.794601in}}%
\pgfpathlineto{\pgfqpoint{3.418010in}{1.787350in}}%
\pgfpathlineto{\pgfqpoint{3.420462in}{1.768872in}}%
\pgfpathlineto{\pgfqpoint{3.427810in}{1.762986in}}%
\pgfpathlineto{\pgfqpoint{3.426466in}{1.757361in}}%
\pgfpathlineto{\pgfqpoint{3.434236in}{1.757543in}}%
\pgfpathlineto{\pgfqpoint{3.436777in}{1.747450in}}%
\pgfpathclose%
\pgfusepath{fill}%
\end{pgfscope}%
\begin{pgfscope}%
\pgfpathrectangle{\pgfqpoint{0.100000in}{0.100000in}}{\pgfqpoint{3.608454in}{2.310000in}}%
\pgfusepath{clip}%
\pgfsetbuttcap%
\pgfsetmiterjoin%
\definecolor{currentfill}{rgb}{0.000000,0.470588,0.764706}%
\pgfsetfillcolor{currentfill}%
\pgfsetlinewidth{0.000000pt}%
\definecolor{currentstroke}{rgb}{0.000000,0.000000,0.000000}%
\pgfsetstrokecolor{currentstroke}%
\pgfsetstrokeopacity{0.000000}%
\pgfsetdash{}{0pt}%
\pgfpathmoveto{\pgfqpoint{1.938292in}{1.256226in}}%
\pgfpathlineto{\pgfqpoint{1.938903in}{1.276826in}}%
\pgfpathlineto{\pgfqpoint{1.925218in}{1.277253in}}%
\pgfpathlineto{\pgfqpoint{1.926190in}{1.311718in}}%
\pgfpathlineto{\pgfqpoint{1.960306in}{1.310682in}}%
\pgfpathlineto{\pgfqpoint{1.987581in}{1.310037in}}%
\pgfpathlineto{\pgfqpoint{1.987451in}{1.303171in}}%
\pgfpathlineto{\pgfqpoint{1.994270in}{1.303029in}}%
\pgfpathlineto{\pgfqpoint{1.992866in}{1.289314in}}%
\pgfpathlineto{\pgfqpoint{1.992256in}{1.268523in}}%
\pgfpathlineto{\pgfqpoint{1.972846in}{1.269115in}}%
\pgfpathlineto{\pgfqpoint{1.972560in}{1.255285in}}%
\pgfpathlineto{\pgfqpoint{1.938292in}{1.256226in}}%
\pgfpathclose%
\pgfusepath{fill}%
\end{pgfscope}%
\begin{pgfscope}%
\pgfpathrectangle{\pgfqpoint{0.100000in}{0.100000in}}{\pgfqpoint{3.608454in}{2.310000in}}%
\pgfusepath{clip}%
\pgfsetbuttcap%
\pgfsetmiterjoin%
\definecolor{currentfill}{rgb}{0.000000,0.431373,0.784314}%
\pgfsetfillcolor{currentfill}%
\pgfsetlinewidth{0.000000pt}%
\definecolor{currentstroke}{rgb}{0.000000,0.000000,0.000000}%
\pgfsetstrokecolor{currentstroke}%
\pgfsetstrokeopacity{0.000000}%
\pgfsetdash{}{0pt}%
\pgfpathmoveto{\pgfqpoint{1.719085in}{1.966639in}}%
\pgfpathlineto{\pgfqpoint{1.718524in}{1.959750in}}%
\pgfpathlineto{\pgfqpoint{1.673140in}{1.963430in}}%
\pgfpathlineto{\pgfqpoint{1.635984in}{1.966799in}}%
\pgfpathlineto{\pgfqpoint{1.638609in}{1.994610in}}%
\pgfpathlineto{\pgfqpoint{1.635493in}{1.994898in}}%
\pgfpathlineto{\pgfqpoint{1.638167in}{2.022688in}}%
\pgfpathlineto{\pgfqpoint{1.668766in}{2.019808in}}%
\pgfpathlineto{\pgfqpoint{1.671188in}{2.047334in}}%
\pgfpathlineto{\pgfqpoint{1.695708in}{2.045165in}}%
\pgfpathlineto{\pgfqpoint{1.700751in}{2.053267in}}%
\pgfpathlineto{\pgfqpoint{1.710047in}{2.050726in}}%
\pgfpathlineto{\pgfqpoint{1.716886in}{2.052052in}}%
\pgfpathlineto{\pgfqpoint{1.715648in}{2.039417in}}%
\pgfpathlineto{\pgfqpoint{1.718555in}{2.035319in}}%
\pgfpathlineto{\pgfqpoint{1.716982in}{2.015681in}}%
\pgfpathlineto{\pgfqpoint{1.720286in}{2.015429in}}%
\pgfpathlineto{\pgfqpoint{1.718098in}{1.987796in}}%
\pgfpathlineto{\pgfqpoint{1.720698in}{1.987602in}}%
\pgfpathlineto{\pgfqpoint{1.719085in}{1.966639in}}%
\pgfpathclose%
\pgfusepath{fill}%
\end{pgfscope}%
\begin{pgfscope}%
\pgfpathrectangle{\pgfqpoint{0.100000in}{0.100000in}}{\pgfqpoint{3.608454in}{2.310000in}}%
\pgfusepath{clip}%
\pgfsetbuttcap%
\pgfsetmiterjoin%
\definecolor{currentfill}{rgb}{0.000000,0.368627,0.815686}%
\pgfsetfillcolor{currentfill}%
\pgfsetlinewidth{0.000000pt}%
\definecolor{currentstroke}{rgb}{0.000000,0.000000,0.000000}%
\pgfsetstrokecolor{currentstroke}%
\pgfsetstrokeopacity{0.000000}%
\pgfsetdash{}{0pt}%
\pgfpathmoveto{\pgfqpoint{0.934361in}{0.385742in}}%
\pgfpathlineto{\pgfqpoint{0.935737in}{0.388869in}}%
\pgfpathlineto{\pgfqpoint{0.937444in}{0.389743in}}%
\pgfpathlineto{\pgfqpoint{0.940486in}{0.388027in}}%
\pgfpathlineto{\pgfqpoint{0.940016in}{0.386614in}}%
\pgfpathlineto{\pgfqpoint{0.942968in}{0.383955in}}%
\pgfpathlineto{\pgfqpoint{0.936317in}{0.384289in}}%
\pgfpathlineto{\pgfqpoint{0.934361in}{0.385742in}}%
\pgfpathclose%
\pgfusepath{fill}%
\end{pgfscope}%
\begin{pgfscope}%
\pgfpathrectangle{\pgfqpoint{0.100000in}{0.100000in}}{\pgfqpoint{3.608454in}{2.310000in}}%
\pgfusepath{clip}%
\pgfsetbuttcap%
\pgfsetmiterjoin%
\definecolor{currentfill}{rgb}{0.000000,0.368627,0.815686}%
\pgfsetfillcolor{currentfill}%
\pgfsetlinewidth{0.000000pt}%
\definecolor{currentstroke}{rgb}{0.000000,0.000000,0.000000}%
\pgfsetstrokecolor{currentstroke}%
\pgfsetstrokeopacity{0.000000}%
\pgfsetdash{}{0pt}%
\pgfpathmoveto{\pgfqpoint{0.911727in}{0.383608in}}%
\pgfpathlineto{\pgfqpoint{0.914559in}{0.385588in}}%
\pgfpathlineto{\pgfqpoint{0.917234in}{0.385734in}}%
\pgfpathlineto{\pgfqpoint{0.927384in}{0.388734in}}%
\pgfpathlineto{\pgfqpoint{0.928405in}{0.390996in}}%
\pgfpathlineto{\pgfqpoint{0.930625in}{0.389633in}}%
\pgfpathlineto{\pgfqpoint{0.930606in}{0.387345in}}%
\pgfpathlineto{\pgfqpoint{0.924754in}{0.385978in}}%
\pgfpathlineto{\pgfqpoint{0.920480in}{0.383050in}}%
\pgfpathlineto{\pgfqpoint{0.918337in}{0.383207in}}%
\pgfpathlineto{\pgfqpoint{0.917432in}{0.381621in}}%
\pgfpathlineto{\pgfqpoint{0.912959in}{0.382564in}}%
\pgfpathlineto{\pgfqpoint{0.911727in}{0.383608in}}%
\pgfpathclose%
\pgfusepath{fill}%
\end{pgfscope}%
\begin{pgfscope}%
\pgfpathrectangle{\pgfqpoint{0.100000in}{0.100000in}}{\pgfqpoint{3.608454in}{2.310000in}}%
\pgfusepath{clip}%
\pgfsetbuttcap%
\pgfsetmiterjoin%
\definecolor{currentfill}{rgb}{0.000000,0.368627,0.815686}%
\pgfsetfillcolor{currentfill}%
\pgfsetlinewidth{0.000000pt}%
\definecolor{currentstroke}{rgb}{0.000000,0.000000,0.000000}%
\pgfsetstrokecolor{currentstroke}%
\pgfsetstrokeopacity{0.000000}%
\pgfsetdash{}{0pt}%
\pgfpathmoveto{\pgfqpoint{0.918694in}{0.390985in}}%
\pgfpathlineto{\pgfqpoint{0.922043in}{0.397743in}}%
\pgfpathlineto{\pgfqpoint{0.922804in}{0.395503in}}%
\pgfpathlineto{\pgfqpoint{0.918694in}{0.390985in}}%
\pgfpathclose%
\pgfusepath{fill}%
\end{pgfscope}%
\begin{pgfscope}%
\pgfpathrectangle{\pgfqpoint{0.100000in}{0.100000in}}{\pgfqpoint{3.608454in}{2.310000in}}%
\pgfusepath{clip}%
\pgfsetbuttcap%
\pgfsetmiterjoin%
\definecolor{currentfill}{rgb}{0.000000,0.368627,0.815686}%
\pgfsetfillcolor{currentfill}%
\pgfsetlinewidth{0.000000pt}%
\definecolor{currentstroke}{rgb}{0.000000,0.000000,0.000000}%
\pgfsetstrokecolor{currentstroke}%
\pgfsetstrokeopacity{0.000000}%
\pgfsetdash{}{0pt}%
\pgfpathmoveto{\pgfqpoint{0.979580in}{0.452265in}}%
\pgfpathlineto{\pgfqpoint{0.983915in}{0.449521in}}%
\pgfpathlineto{\pgfqpoint{0.986635in}{0.453959in}}%
\pgfpathlineto{\pgfqpoint{0.997579in}{0.447110in}}%
\pgfpathlineto{\pgfqpoint{0.996218in}{0.444880in}}%
\pgfpathlineto{\pgfqpoint{1.002410in}{0.441046in}}%
\pgfpathlineto{\pgfqpoint{1.019938in}{0.430914in}}%
\pgfpathlineto{\pgfqpoint{1.017896in}{0.425014in}}%
\pgfpathlineto{\pgfqpoint{1.014239in}{0.423539in}}%
\pgfpathlineto{\pgfqpoint{1.015064in}{0.421049in}}%
\pgfpathlineto{\pgfqpoint{1.012216in}{0.417308in}}%
\pgfpathlineto{\pgfqpoint{1.015043in}{0.415179in}}%
\pgfpathlineto{\pgfqpoint{1.018284in}{0.417154in}}%
\pgfpathlineto{\pgfqpoint{1.018953in}{0.415851in}}%
\pgfpathlineto{\pgfqpoint{1.022698in}{0.413642in}}%
\pgfpathlineto{\pgfqpoint{1.021385in}{0.411488in}}%
\pgfpathlineto{\pgfqpoint{1.023662in}{0.410127in}}%
\pgfpathlineto{\pgfqpoint{1.022393in}{0.408007in}}%
\pgfpathlineto{\pgfqpoint{1.024400in}{0.406832in}}%
\pgfpathlineto{\pgfqpoint{1.019250in}{0.398079in}}%
\pgfpathlineto{\pgfqpoint{1.020524in}{0.396033in}}%
\pgfpathlineto{\pgfqpoint{1.017303in}{0.390307in}}%
\pgfpathlineto{\pgfqpoint{1.027182in}{0.384801in}}%
\pgfpathlineto{\pgfqpoint{1.016055in}{0.364729in}}%
\pgfpathlineto{\pgfqpoint{1.004608in}{0.344079in}}%
\pgfpathlineto{\pgfqpoint{0.997055in}{0.330455in}}%
\pgfpathlineto{\pgfqpoint{0.993111in}{0.336827in}}%
\pgfpathlineto{\pgfqpoint{0.996872in}{0.337073in}}%
\pgfpathlineto{\pgfqpoint{0.997859in}{0.339919in}}%
\pgfpathlineto{\pgfqpoint{0.995627in}{0.341376in}}%
\pgfpathlineto{\pgfqpoint{0.995160in}{0.339365in}}%
\pgfpathlineto{\pgfqpoint{0.990363in}{0.340190in}}%
\pgfpathlineto{\pgfqpoint{0.988940in}{0.342964in}}%
\pgfpathlineto{\pgfqpoint{0.979357in}{0.351803in}}%
\pgfpathlineto{\pgfqpoint{0.975077in}{0.352806in}}%
\pgfpathlineto{\pgfqpoint{0.969268in}{0.355424in}}%
\pgfpathlineto{\pgfqpoint{0.963603in}{0.357154in}}%
\pgfpathlineto{\pgfqpoint{0.962385in}{0.358792in}}%
\pgfpathlineto{\pgfqpoint{0.962290in}{0.362125in}}%
\pgfpathlineto{\pgfqpoint{0.961040in}{0.364752in}}%
\pgfpathlineto{\pgfqpoint{0.957865in}{0.367144in}}%
\pgfpathlineto{\pgfqpoint{0.958163in}{0.370011in}}%
\pgfpathlineto{\pgfqpoint{0.956916in}{0.372017in}}%
\pgfpathlineto{\pgfqpoint{0.959883in}{0.375282in}}%
\pgfpathlineto{\pgfqpoint{0.960378in}{0.377067in}}%
\pgfpathlineto{\pgfqpoint{0.958642in}{0.378412in}}%
\pgfpathlineto{\pgfqpoint{0.956108in}{0.376672in}}%
\pgfpathlineto{\pgfqpoint{0.953407in}{0.377974in}}%
\pgfpathlineto{\pgfqpoint{0.952146in}{0.377061in}}%
\pgfpathlineto{\pgfqpoint{0.949343in}{0.382116in}}%
\pgfpathlineto{\pgfqpoint{0.947228in}{0.383148in}}%
\pgfpathlineto{\pgfqpoint{0.949313in}{0.385149in}}%
\pgfpathlineto{\pgfqpoint{0.949856in}{0.389896in}}%
\pgfpathlineto{\pgfqpoint{0.949598in}{0.393135in}}%
\pgfpathlineto{\pgfqpoint{0.945394in}{0.394108in}}%
\pgfpathlineto{\pgfqpoint{0.947389in}{0.396132in}}%
\pgfpathlineto{\pgfqpoint{0.942948in}{0.398138in}}%
\pgfpathlineto{\pgfqpoint{0.943546in}{0.402478in}}%
\pgfpathlineto{\pgfqpoint{0.941428in}{0.403633in}}%
\pgfpathlineto{\pgfqpoint{0.940901in}{0.405559in}}%
\pgfpathlineto{\pgfqpoint{0.938307in}{0.405202in}}%
\pgfpathlineto{\pgfqpoint{0.937360in}{0.403019in}}%
\pgfpathlineto{\pgfqpoint{0.933986in}{0.407135in}}%
\pgfpathlineto{\pgfqpoint{0.935525in}{0.410297in}}%
\pgfpathlineto{\pgfqpoint{0.930641in}{0.409132in}}%
\pgfpathlineto{\pgfqpoint{0.924871in}{0.408830in}}%
\pgfpathlineto{\pgfqpoint{0.926703in}{0.411773in}}%
\pgfpathlineto{\pgfqpoint{0.928482in}{0.410664in}}%
\pgfpathlineto{\pgfqpoint{0.928700in}{0.414487in}}%
\pgfpathlineto{\pgfqpoint{0.926958in}{0.414550in}}%
\pgfpathlineto{\pgfqpoint{0.922100in}{0.411389in}}%
\pgfpathlineto{\pgfqpoint{0.918534in}{0.408022in}}%
\pgfpathlineto{\pgfqpoint{0.919688in}{0.407473in}}%
\pgfpathlineto{\pgfqpoint{0.921994in}{0.409270in}}%
\pgfpathlineto{\pgfqpoint{0.923912in}{0.407532in}}%
\pgfpathlineto{\pgfqpoint{0.921409in}{0.405514in}}%
\pgfpathlineto{\pgfqpoint{0.919314in}{0.405095in}}%
\pgfpathlineto{\pgfqpoint{0.916088in}{0.406347in}}%
\pgfpathlineto{\pgfqpoint{0.917614in}{0.402235in}}%
\pgfpathlineto{\pgfqpoint{0.919867in}{0.400817in}}%
\pgfpathlineto{\pgfqpoint{0.913772in}{0.398471in}}%
\pgfpathlineto{\pgfqpoint{0.912875in}{0.397284in}}%
\pgfpathlineto{\pgfqpoint{0.916456in}{0.393000in}}%
\pgfpathlineto{\pgfqpoint{0.915325in}{0.388126in}}%
\pgfpathlineto{\pgfqpoint{0.913318in}{0.388120in}}%
\pgfpathlineto{\pgfqpoint{0.911580in}{0.390660in}}%
\pgfpathlineto{\pgfqpoint{0.912461in}{0.393615in}}%
\pgfpathlineto{\pgfqpoint{0.910952in}{0.394996in}}%
\pgfpathlineto{\pgfqpoint{0.908291in}{0.391990in}}%
\pgfpathlineto{\pgfqpoint{0.906140in}{0.393309in}}%
\pgfpathlineto{\pgfqpoint{0.908663in}{0.396558in}}%
\pgfpathlineto{\pgfqpoint{0.913452in}{0.403668in}}%
\pgfpathlineto{\pgfqpoint{0.912214in}{0.404513in}}%
\pgfpathlineto{\pgfqpoint{0.917407in}{0.411973in}}%
\pgfpathlineto{\pgfqpoint{0.916327in}{0.412673in}}%
\pgfpathlineto{\pgfqpoint{0.918151in}{0.415401in}}%
\pgfpathlineto{\pgfqpoint{0.921451in}{0.413351in}}%
\pgfpathlineto{\pgfqpoint{0.926490in}{0.421058in}}%
\pgfpathlineto{\pgfqpoint{0.930988in}{0.427301in}}%
\pgfpathlineto{\pgfqpoint{0.945694in}{0.417578in}}%
\pgfpathlineto{\pgfqpoint{0.946494in}{0.418805in}}%
\pgfpathlineto{\pgfqpoint{0.949551in}{0.416794in}}%
\pgfpathlineto{\pgfqpoint{0.960218in}{0.433861in}}%
\pgfpathlineto{\pgfqpoint{0.961359in}{0.436162in}}%
\pgfpathlineto{\pgfqpoint{0.967777in}{0.431965in}}%
\pgfpathlineto{\pgfqpoint{0.977356in}{0.447349in}}%
\pgfpathlineto{\pgfqpoint{0.979580in}{0.452265in}}%
\pgfpathclose%
\pgfusepath{fill}%
\end{pgfscope}%
\begin{pgfscope}%
\pgfpathrectangle{\pgfqpoint{0.100000in}{0.100000in}}{\pgfqpoint{3.608454in}{2.310000in}}%
\pgfusepath{clip}%
\pgfsetbuttcap%
\pgfsetmiterjoin%
\definecolor{currentfill}{rgb}{0.000000,0.941176,0.529412}%
\pgfsetfillcolor{currentfill}%
\pgfsetlinewidth{0.000000pt}%
\definecolor{currentstroke}{rgb}{0.000000,0.000000,0.000000}%
\pgfsetstrokecolor{currentstroke}%
\pgfsetstrokeopacity{0.000000}%
\pgfsetdash{}{0pt}%
\pgfpathmoveto{\pgfqpoint{2.923822in}{1.280887in}}%
\pgfpathlineto{\pgfqpoint{2.927867in}{1.284414in}}%
\pgfpathlineto{\pgfqpoint{2.929879in}{1.293625in}}%
\pgfpathlineto{\pgfqpoint{2.935674in}{1.297744in}}%
\pgfpathlineto{\pgfqpoint{2.931210in}{1.301107in}}%
\pgfpathlineto{\pgfqpoint{2.946226in}{1.306488in}}%
\pgfpathlineto{\pgfqpoint{2.952046in}{1.304075in}}%
\pgfpathlineto{\pgfqpoint{2.961369in}{1.291611in}}%
\pgfpathlineto{\pgfqpoint{2.967623in}{1.287269in}}%
\pgfpathlineto{\pgfqpoint{2.965505in}{1.286020in}}%
\pgfpathlineto{\pgfqpoint{2.960443in}{1.270838in}}%
\pgfpathlineto{\pgfqpoint{2.950534in}{1.260497in}}%
\pgfpathlineto{\pgfqpoint{2.942904in}{1.258852in}}%
\pgfpathlineto{\pgfqpoint{2.935481in}{1.264222in}}%
\pgfpathlineto{\pgfqpoint{2.926305in}{1.268226in}}%
\pgfpathlineto{\pgfqpoint{2.923526in}{1.275100in}}%
\pgfpathlineto{\pgfqpoint{2.923822in}{1.280887in}}%
\pgfpathclose%
\pgfusepath{fill}%
\end{pgfscope}%
\begin{pgfscope}%
\pgfpathrectangle{\pgfqpoint{0.100000in}{0.100000in}}{\pgfqpoint{3.608454in}{2.310000in}}%
\pgfusepath{clip}%
\pgfsetbuttcap%
\pgfsetmiterjoin%
\definecolor{currentfill}{rgb}{0.000000,0.443137,0.778431}%
\pgfsetfillcolor{currentfill}%
\pgfsetlinewidth{0.000000pt}%
\definecolor{currentstroke}{rgb}{0.000000,0.000000,0.000000}%
\pgfsetstrokecolor{currentstroke}%
\pgfsetstrokeopacity{0.000000}%
\pgfsetdash{}{0pt}%
\pgfpathmoveto{\pgfqpoint{3.019784in}{0.555093in}}%
\pgfpathlineto{\pgfqpoint{3.017489in}{0.571785in}}%
\pgfpathlineto{\pgfqpoint{3.006224in}{0.574274in}}%
\pgfpathlineto{\pgfqpoint{3.001283in}{0.581332in}}%
\pgfpathlineto{\pgfqpoint{3.006411in}{0.592107in}}%
\pgfpathlineto{\pgfqpoint{2.994626in}{0.604123in}}%
\pgfpathlineto{\pgfqpoint{3.039989in}{0.610672in}}%
\pgfpathlineto{\pgfqpoint{3.040352in}{0.617807in}}%
\pgfpathlineto{\pgfqpoint{3.037437in}{0.635958in}}%
\pgfpathlineto{\pgfqpoint{3.040376in}{0.630406in}}%
\pgfpathlineto{\pgfqpoint{3.046213in}{0.629170in}}%
\pgfpathlineto{\pgfqpoint{3.049410in}{0.623089in}}%
\pgfpathlineto{\pgfqpoint{3.060692in}{0.615927in}}%
\pgfpathlineto{\pgfqpoint{3.065010in}{0.604224in}}%
\pgfpathlineto{\pgfqpoint{3.071577in}{0.605034in}}%
\pgfpathlineto{\pgfqpoint{3.078312in}{0.602797in}}%
\pgfpathlineto{\pgfqpoint{3.080990in}{0.605844in}}%
\pgfpathlineto{\pgfqpoint{3.090875in}{0.590177in}}%
\pgfpathlineto{\pgfqpoint{3.099275in}{0.583081in}}%
\pgfpathlineto{\pgfqpoint{3.099105in}{0.576837in}}%
\pgfpathlineto{\pgfqpoint{3.102843in}{0.570453in}}%
\pgfpathlineto{\pgfqpoint{3.110913in}{0.514275in}}%
\pgfpathlineto{\pgfqpoint{3.091930in}{0.511405in}}%
\pgfpathlineto{\pgfqpoint{3.090421in}{0.517236in}}%
\pgfpathlineto{\pgfqpoint{3.084501in}{0.525466in}}%
\pgfpathlineto{\pgfqpoint{3.077818in}{0.527142in}}%
\pgfpathlineto{\pgfqpoint{3.076539in}{0.531713in}}%
\pgfpathlineto{\pgfqpoint{3.070237in}{0.538537in}}%
\pgfpathlineto{\pgfqpoint{3.059031in}{0.547094in}}%
\pgfpathlineto{\pgfqpoint{3.055318in}{0.556018in}}%
\pgfpathlineto{\pgfqpoint{3.048333in}{0.554969in}}%
\pgfpathlineto{\pgfqpoint{3.047294in}{0.561961in}}%
\pgfpathlineto{\pgfqpoint{3.026291in}{0.558732in}}%
\pgfpathlineto{\pgfqpoint{3.019784in}{0.555093in}}%
\pgfpathclose%
\pgfusepath{fill}%
\end{pgfscope}%
\begin{pgfscope}%
\pgfpathrectangle{\pgfqpoint{0.100000in}{0.100000in}}{\pgfqpoint{3.608454in}{2.310000in}}%
\pgfusepath{clip}%
\pgfsetbuttcap%
\pgfsetmiterjoin%
\definecolor{currentfill}{rgb}{0.000000,0.474510,0.762745}%
\pgfsetfillcolor{currentfill}%
\pgfsetlinewidth{0.000000pt}%
\definecolor{currentstroke}{rgb}{0.000000,0.000000,0.000000}%
\pgfsetstrokecolor{currentstroke}%
\pgfsetstrokeopacity{0.000000}%
\pgfsetdash{}{0pt}%
\pgfpathmoveto{\pgfqpoint{1.723873in}{1.194151in}}%
\pgfpathlineto{\pgfqpoint{1.731511in}{1.193664in}}%
\pgfpathlineto{\pgfqpoint{1.728388in}{1.154348in}}%
\pgfpathlineto{\pgfqpoint{1.685892in}{1.157286in}}%
\pgfpathlineto{\pgfqpoint{1.683362in}{1.122232in}}%
\pgfpathlineto{\pgfqpoint{1.648937in}{1.124836in}}%
\pgfpathlineto{\pgfqpoint{1.651744in}{1.159945in}}%
\pgfpathlineto{\pgfqpoint{1.598511in}{1.164402in}}%
\pgfpathlineto{\pgfqpoint{1.602002in}{1.203777in}}%
\pgfpathlineto{\pgfqpoint{1.621102in}{1.201730in}}%
\pgfpathlineto{\pgfqpoint{1.662385in}{1.198210in}}%
\pgfpathlineto{\pgfqpoint{1.664847in}{1.229439in}}%
\pgfpathlineto{\pgfqpoint{1.695230in}{1.227106in}}%
\pgfpathlineto{\pgfqpoint{1.693075in}{1.196095in}}%
\pgfpathlineto{\pgfqpoint{1.723873in}{1.194151in}}%
\pgfpathclose%
\pgfusepath{fill}%
\end{pgfscope}%
\begin{pgfscope}%
\pgfpathrectangle{\pgfqpoint{0.100000in}{0.100000in}}{\pgfqpoint{3.608454in}{2.310000in}}%
\pgfusepath{clip}%
\pgfsetbuttcap%
\pgfsetmiterjoin%
\definecolor{currentfill}{rgb}{0.000000,0.607843,0.696078}%
\pgfsetfillcolor{currentfill}%
\pgfsetlinewidth{0.000000pt}%
\definecolor{currentstroke}{rgb}{0.000000,0.000000,0.000000}%
\pgfsetstrokecolor{currentstroke}%
\pgfsetstrokeopacity{0.000000}%
\pgfsetdash{}{0pt}%
\pgfpathmoveto{\pgfqpoint{2.642473in}{1.232080in}}%
\pgfpathlineto{\pgfqpoint{2.647960in}{1.220595in}}%
\pgfpathlineto{\pgfqpoint{2.639824in}{1.220359in}}%
\pgfpathlineto{\pgfqpoint{2.634641in}{1.217991in}}%
\pgfpathlineto{\pgfqpoint{2.631929in}{1.205576in}}%
\pgfpathlineto{\pgfqpoint{2.620792in}{1.209072in}}%
\pgfpathlineto{\pgfqpoint{2.607679in}{1.208265in}}%
\pgfpathlineto{\pgfqpoint{2.595443in}{1.205574in}}%
\pgfpathlineto{\pgfqpoint{2.589411in}{1.214310in}}%
\pgfpathlineto{\pgfqpoint{2.585378in}{1.221768in}}%
\pgfpathlineto{\pgfqpoint{2.589982in}{1.231634in}}%
\pgfpathlineto{\pgfqpoint{2.586421in}{1.235460in}}%
\pgfpathlineto{\pgfqpoint{2.585128in}{1.245024in}}%
\pgfpathlineto{\pgfqpoint{2.577210in}{1.248460in}}%
\pgfpathlineto{\pgfqpoint{2.585709in}{1.266483in}}%
\pgfpathlineto{\pgfqpoint{2.586412in}{1.272897in}}%
\pgfpathlineto{\pgfqpoint{2.601681in}{1.267043in}}%
\pgfpathlineto{\pgfqpoint{2.602387in}{1.273595in}}%
\pgfpathlineto{\pgfqpoint{2.617290in}{1.282808in}}%
\pgfpathlineto{\pgfqpoint{2.620883in}{1.276667in}}%
\pgfpathlineto{\pgfqpoint{2.627054in}{1.277303in}}%
\pgfpathlineto{\pgfqpoint{2.627023in}{1.272051in}}%
\pgfpathlineto{\pgfqpoint{2.625972in}{1.268011in}}%
\pgfpathlineto{\pgfqpoint{2.630443in}{1.257485in}}%
\pgfpathlineto{\pgfqpoint{2.628815in}{1.247944in}}%
\pgfpathlineto{\pgfqpoint{2.632971in}{1.237093in}}%
\pgfpathlineto{\pgfqpoint{2.641029in}{1.235231in}}%
\pgfpathlineto{\pgfqpoint{2.642473in}{1.232080in}}%
\pgfpathclose%
\pgfusepath{fill}%
\end{pgfscope}%
\begin{pgfscope}%
\pgfpathrectangle{\pgfqpoint{0.100000in}{0.100000in}}{\pgfqpoint{3.608454in}{2.310000in}}%
\pgfusepath{clip}%
\pgfsetbuttcap%
\pgfsetmiterjoin%
\definecolor{currentfill}{rgb}{0.000000,0.415686,0.792157}%
\pgfsetfillcolor{currentfill}%
\pgfsetlinewidth{0.000000pt}%
\definecolor{currentstroke}{rgb}{0.000000,0.000000,0.000000}%
\pgfsetstrokecolor{currentstroke}%
\pgfsetstrokeopacity{0.000000}%
\pgfsetdash{}{0pt}%
\pgfpathmoveto{\pgfqpoint{2.218568in}{1.724511in}}%
\pgfpathlineto{\pgfqpoint{2.238913in}{1.724992in}}%
\pgfpathlineto{\pgfqpoint{2.252724in}{1.724220in}}%
\pgfpathlineto{\pgfqpoint{2.253472in}{1.697776in}}%
\pgfpathlineto{\pgfqpoint{2.220391in}{1.696881in}}%
\pgfpathlineto{\pgfqpoint{2.184501in}{1.696223in}}%
\pgfpathlineto{\pgfqpoint{2.184096in}{1.723863in}}%
\pgfpathlineto{\pgfqpoint{2.218568in}{1.724511in}}%
\pgfpathclose%
\pgfusepath{fill}%
\end{pgfscope}%
\begin{pgfscope}%
\pgfpathrectangle{\pgfqpoint{0.100000in}{0.100000in}}{\pgfqpoint{3.608454in}{2.310000in}}%
\pgfusepath{clip}%
\pgfsetbuttcap%
\pgfsetmiterjoin%
\definecolor{currentfill}{rgb}{0.000000,0.207843,0.896078}%
\pgfsetfillcolor{currentfill}%
\pgfsetlinewidth{0.000000pt}%
\definecolor{currentstroke}{rgb}{0.000000,0.000000,0.000000}%
\pgfsetstrokecolor{currentstroke}%
\pgfsetstrokeopacity{0.000000}%
\pgfsetdash{}{0pt}%
\pgfpathmoveto{\pgfqpoint{1.989587in}{1.638444in}}%
\pgfpathlineto{\pgfqpoint{1.981253in}{1.641675in}}%
\pgfpathlineto{\pgfqpoint{1.981953in}{1.671006in}}%
\pgfpathlineto{\pgfqpoint{1.968146in}{1.671305in}}%
\pgfpathlineto{\pgfqpoint{1.968709in}{1.697472in}}%
\pgfpathlineto{\pgfqpoint{1.956891in}{1.697804in}}%
\pgfpathlineto{\pgfqpoint{1.957610in}{1.725492in}}%
\pgfpathlineto{\pgfqpoint{1.998767in}{1.724416in}}%
\pgfpathlineto{\pgfqpoint{2.023670in}{1.724067in}}%
\pgfpathlineto{\pgfqpoint{2.023239in}{1.696380in}}%
\pgfpathlineto{\pgfqpoint{2.014851in}{1.696524in}}%
\pgfpathlineto{\pgfqpoint{2.015077in}{1.691199in}}%
\pgfpathlineto{\pgfqpoint{2.019001in}{1.688123in}}%
\pgfpathlineto{\pgfqpoint{2.018524in}{1.680571in}}%
\pgfpathlineto{\pgfqpoint{2.021565in}{1.674298in}}%
\pgfpathlineto{\pgfqpoint{2.022652in}{1.663329in}}%
\pgfpathlineto{\pgfqpoint{2.002309in}{1.663696in}}%
\pgfpathlineto{\pgfqpoint{2.001679in}{1.633660in}}%
\pgfpathlineto{\pgfqpoint{1.989587in}{1.638444in}}%
\pgfpathclose%
\pgfusepath{fill}%
\end{pgfscope}%
\begin{pgfscope}%
\pgfpathrectangle{\pgfqpoint{0.100000in}{0.100000in}}{\pgfqpoint{3.608454in}{2.310000in}}%
\pgfusepath{clip}%
\pgfsetbuttcap%
\pgfsetmiterjoin%
\definecolor{currentfill}{rgb}{0.000000,0.490196,0.754902}%
\pgfsetfillcolor{currentfill}%
\pgfsetlinewidth{0.000000pt}%
\definecolor{currentstroke}{rgb}{0.000000,0.000000,0.000000}%
\pgfsetstrokecolor{currentstroke}%
\pgfsetstrokeopacity{0.000000}%
\pgfsetdash{}{0pt}%
\pgfpathmoveto{\pgfqpoint{2.018170in}{1.419481in}}%
\pgfpathlineto{\pgfqpoint{2.031759in}{1.419257in}}%
\pgfpathlineto{\pgfqpoint{2.031304in}{1.384920in}}%
\pgfpathlineto{\pgfqpoint{1.962491in}{1.386340in}}%
\pgfpathlineto{\pgfqpoint{1.963356in}{1.420723in}}%
\pgfpathlineto{\pgfqpoint{2.018170in}{1.419481in}}%
\pgfpathclose%
\pgfusepath{fill}%
\end{pgfscope}%
\begin{pgfscope}%
\pgfpathrectangle{\pgfqpoint{0.100000in}{0.100000in}}{\pgfqpoint{3.608454in}{2.310000in}}%
\pgfusepath{clip}%
\pgfsetbuttcap%
\pgfsetmiterjoin%
\definecolor{currentfill}{rgb}{0.000000,0.592157,0.703922}%
\pgfsetfillcolor{currentfill}%
\pgfsetlinewidth{0.000000pt}%
\definecolor{currentstroke}{rgb}{0.000000,0.000000,0.000000}%
\pgfsetstrokecolor{currentstroke}%
\pgfsetstrokeopacity{0.000000}%
\pgfsetdash{}{0pt}%
\pgfpathmoveto{\pgfqpoint{2.396575in}{1.440884in}}%
\pgfpathlineto{\pgfqpoint{2.381360in}{1.440493in}}%
\pgfpathlineto{\pgfqpoint{2.381085in}{1.447410in}}%
\pgfpathlineto{\pgfqpoint{2.353369in}{1.446719in}}%
\pgfpathlineto{\pgfqpoint{2.352440in}{1.474762in}}%
\pgfpathlineto{\pgfqpoint{2.359331in}{1.474799in}}%
\pgfpathlineto{\pgfqpoint{2.358001in}{1.509016in}}%
\pgfpathlineto{\pgfqpoint{2.378619in}{1.509638in}}%
\pgfpathlineto{\pgfqpoint{2.378376in}{1.516534in}}%
\pgfpathlineto{\pgfqpoint{2.405330in}{1.517808in}}%
\pgfpathlineto{\pgfqpoint{2.406029in}{1.504012in}}%
\pgfpathlineto{\pgfqpoint{2.407342in}{1.476448in}}%
\pgfpathlineto{\pgfqpoint{2.414274in}{1.476742in}}%
\pgfpathlineto{\pgfqpoint{2.414833in}{1.467947in}}%
\pgfpathlineto{\pgfqpoint{2.412968in}{1.462469in}}%
\pgfpathlineto{\pgfqpoint{2.405497in}{1.456494in}}%
\pgfpathlineto{\pgfqpoint{2.401121in}{1.445010in}}%
\pgfpathlineto{\pgfqpoint{2.396575in}{1.440884in}}%
\pgfpathclose%
\pgfusepath{fill}%
\end{pgfscope}%
\begin{pgfscope}%
\pgfpathrectangle{\pgfqpoint{0.100000in}{0.100000in}}{\pgfqpoint{3.608454in}{2.310000in}}%
\pgfusepath{clip}%
\pgfsetbuttcap%
\pgfsetmiterjoin%
\definecolor{currentfill}{rgb}{0.000000,0.572549,0.713725}%
\pgfsetfillcolor{currentfill}%
\pgfsetlinewidth{0.000000pt}%
\definecolor{currentstroke}{rgb}{0.000000,0.000000,0.000000}%
\pgfsetstrokecolor{currentstroke}%
\pgfsetstrokeopacity{0.000000}%
\pgfsetdash{}{0pt}%
\pgfpathmoveto{\pgfqpoint{2.729948in}{1.258824in}}%
\pgfpathlineto{\pgfqpoint{2.730780in}{1.248329in}}%
\pgfpathlineto{\pgfqpoint{2.722412in}{1.251628in}}%
\pgfpathlineto{\pgfqpoint{2.718150in}{1.247622in}}%
\pgfpathlineto{\pgfqpoint{2.711292in}{1.250974in}}%
\pgfpathlineto{\pgfqpoint{2.704025in}{1.249697in}}%
\pgfpathlineto{\pgfqpoint{2.701382in}{1.255633in}}%
\pgfpathlineto{\pgfqpoint{2.694535in}{1.255275in}}%
\pgfpathlineto{\pgfqpoint{2.694291in}{1.266816in}}%
\pgfpathlineto{\pgfqpoint{2.691418in}{1.266749in}}%
\pgfpathlineto{\pgfqpoint{2.684273in}{1.275179in}}%
\pgfpathlineto{\pgfqpoint{2.690035in}{1.277558in}}%
\pgfpathlineto{\pgfqpoint{2.698255in}{1.290856in}}%
\pgfpathlineto{\pgfqpoint{2.704197in}{1.287428in}}%
\pgfpathlineto{\pgfqpoint{2.718323in}{1.291481in}}%
\pgfpathlineto{\pgfqpoint{2.719889in}{1.285770in}}%
\pgfpathlineto{\pgfqpoint{2.727518in}{1.286122in}}%
\pgfpathlineto{\pgfqpoint{2.729699in}{1.283399in}}%
\pgfpathlineto{\pgfqpoint{2.729948in}{1.258824in}}%
\pgfpathclose%
\pgfusepath{fill}%
\end{pgfscope}%
\begin{pgfscope}%
\pgfpathrectangle{\pgfqpoint{0.100000in}{0.100000in}}{\pgfqpoint{3.608454in}{2.310000in}}%
\pgfusepath{clip}%
\pgfsetbuttcap%
\pgfsetmiterjoin%
\definecolor{currentfill}{rgb}{0.000000,0.529412,0.735294}%
\pgfsetfillcolor{currentfill}%
\pgfsetlinewidth{0.000000pt}%
\definecolor{currentstroke}{rgb}{0.000000,0.000000,0.000000}%
\pgfsetstrokecolor{currentstroke}%
\pgfsetstrokeopacity{0.000000}%
\pgfsetdash{}{0pt}%
\pgfpathmoveto{\pgfqpoint{2.868390in}{1.010555in}}%
\pgfpathlineto{\pgfqpoint{2.871596in}{1.014173in}}%
\pgfpathlineto{\pgfqpoint{2.886412in}{1.013136in}}%
\pgfpathlineto{\pgfqpoint{2.885344in}{1.008668in}}%
\pgfpathlineto{\pgfqpoint{2.897463in}{0.997613in}}%
\pgfpathlineto{\pgfqpoint{2.908530in}{0.995600in}}%
\pgfpathlineto{\pgfqpoint{2.902891in}{0.990956in}}%
\pgfpathlineto{\pgfqpoint{2.896953in}{0.978821in}}%
\pgfpathlineto{\pgfqpoint{2.896071in}{0.971539in}}%
\pgfpathlineto{\pgfqpoint{2.889208in}{0.971223in}}%
\pgfpathlineto{\pgfqpoint{2.877953in}{0.974806in}}%
\pgfpathlineto{\pgfqpoint{2.870396in}{0.968762in}}%
\pgfpathlineto{\pgfqpoint{2.864261in}{0.976571in}}%
\pgfpathlineto{\pgfqpoint{2.850509in}{0.983786in}}%
\pgfpathlineto{\pgfqpoint{2.844938in}{0.981568in}}%
\pgfpathlineto{\pgfqpoint{2.837086in}{0.989053in}}%
\pgfpathlineto{\pgfqpoint{2.839244in}{0.999162in}}%
\pgfpathlineto{\pgfqpoint{2.843269in}{0.994686in}}%
\pgfpathlineto{\pgfqpoint{2.856780in}{0.993705in}}%
\pgfpathlineto{\pgfqpoint{2.859121in}{0.988728in}}%
\pgfpathlineto{\pgfqpoint{2.868081in}{0.996615in}}%
\pgfpathlineto{\pgfqpoint{2.869364in}{1.002057in}}%
\pgfpathlineto{\pgfqpoint{2.865578in}{1.008074in}}%
\pgfpathlineto{\pgfqpoint{2.868390in}{1.010555in}}%
\pgfpathclose%
\pgfusepath{fill}%
\end{pgfscope}%
\begin{pgfscope}%
\pgfpathrectangle{\pgfqpoint{0.100000in}{0.100000in}}{\pgfqpoint{3.608454in}{2.310000in}}%
\pgfusepath{clip}%
\pgfsetbuttcap%
\pgfsetmiterjoin%
\definecolor{currentfill}{rgb}{0.000000,0.364706,0.817647}%
\pgfsetfillcolor{currentfill}%
\pgfsetlinewidth{0.000000pt}%
\definecolor{currentstroke}{rgb}{0.000000,0.000000,0.000000}%
\pgfsetstrokecolor{currentstroke}%
\pgfsetstrokeopacity{0.000000}%
\pgfsetdash{}{0pt}%
\pgfpathmoveto{\pgfqpoint{1.858155in}{1.589614in}}%
\pgfpathlineto{\pgfqpoint{1.831612in}{1.590776in}}%
\pgfpathlineto{\pgfqpoint{1.803419in}{1.592288in}}%
\pgfpathlineto{\pgfqpoint{1.804075in}{1.619812in}}%
\pgfpathlineto{\pgfqpoint{1.805638in}{1.647292in}}%
\pgfpathlineto{\pgfqpoint{1.805743in}{1.664546in}}%
\pgfpathlineto{\pgfqpoint{1.860373in}{1.661671in}}%
\pgfpathlineto{\pgfqpoint{1.858155in}{1.589614in}}%
\pgfpathclose%
\pgfusepath{fill}%
\end{pgfscope}%
\begin{pgfscope}%
\pgfpathrectangle{\pgfqpoint{0.100000in}{0.100000in}}{\pgfqpoint{3.608454in}{2.310000in}}%
\pgfusepath{clip}%
\pgfsetbuttcap%
\pgfsetmiterjoin%
\definecolor{currentfill}{rgb}{0.000000,0.325490,0.837255}%
\pgfsetfillcolor{currentfill}%
\pgfsetlinewidth{0.000000pt}%
\definecolor{currentstroke}{rgb}{0.000000,0.000000,0.000000}%
\pgfsetstrokecolor{currentstroke}%
\pgfsetstrokeopacity{0.000000}%
\pgfsetdash{}{0pt}%
\pgfpathmoveto{\pgfqpoint{2.039668in}{1.566696in}}%
\pgfpathlineto{\pgfqpoint{2.043988in}{1.561009in}}%
\pgfpathlineto{\pgfqpoint{2.040513in}{1.552171in}}%
\pgfpathlineto{\pgfqpoint{2.021430in}{1.552509in}}%
\pgfpathlineto{\pgfqpoint{2.021531in}{1.557101in}}%
\pgfpathlineto{\pgfqpoint{1.994270in}{1.557682in}}%
\pgfpathlineto{\pgfqpoint{1.987553in}{1.557852in}}%
\pgfpathlineto{\pgfqpoint{1.988196in}{1.585344in}}%
\pgfpathlineto{\pgfqpoint{1.999669in}{1.585100in}}%
\pgfpathlineto{\pgfqpoint{2.000030in}{1.598859in}}%
\pgfpathlineto{\pgfqpoint{2.005666in}{1.599833in}}%
\pgfpathlineto{\pgfqpoint{2.027350in}{1.599345in}}%
\pgfpathlineto{\pgfqpoint{2.027934in}{1.590280in}}%
\pgfpathlineto{\pgfqpoint{2.036665in}{1.576098in}}%
\pgfpathlineto{\pgfqpoint{2.040288in}{1.575134in}}%
\pgfpathlineto{\pgfqpoint{2.039668in}{1.566696in}}%
\pgfpathclose%
\pgfusepath{fill}%
\end{pgfscope}%
\begin{pgfscope}%
\pgfpathrectangle{\pgfqpoint{0.100000in}{0.100000in}}{\pgfqpoint{3.608454in}{2.310000in}}%
\pgfusepath{clip}%
\pgfsetbuttcap%
\pgfsetmiterjoin%
\definecolor{currentfill}{rgb}{0.000000,0.352941,0.823529}%
\pgfsetfillcolor{currentfill}%
\pgfsetlinewidth{0.000000pt}%
\definecolor{currentstroke}{rgb}{0.000000,0.000000,0.000000}%
\pgfsetstrokecolor{currentstroke}%
\pgfsetstrokeopacity{0.000000}%
\pgfsetdash{}{0pt}%
\pgfpathmoveto{\pgfqpoint{2.258667in}{1.780519in}}%
\pgfpathlineto{\pgfqpoint{2.247155in}{1.782906in}}%
\pgfpathlineto{\pgfqpoint{2.241126in}{1.786250in}}%
\pgfpathlineto{\pgfqpoint{2.231530in}{1.786839in}}%
\pgfpathlineto{\pgfqpoint{2.231724in}{1.779996in}}%
\pgfpathlineto{\pgfqpoint{2.217797in}{1.777469in}}%
\pgfpathlineto{\pgfqpoint{2.217878in}{1.773970in}}%
\pgfpathlineto{\pgfqpoint{2.204133in}{1.773697in}}%
\pgfpathlineto{\pgfqpoint{2.204050in}{1.779422in}}%
\pgfpathlineto{\pgfqpoint{2.190305in}{1.779149in}}%
\pgfpathlineto{\pgfqpoint{2.168433in}{1.778817in}}%
\pgfpathlineto{\pgfqpoint{2.169520in}{1.792650in}}%
\pgfpathlineto{\pgfqpoint{2.162655in}{1.792587in}}%
\pgfpathlineto{\pgfqpoint{2.162319in}{1.813349in}}%
\pgfpathlineto{\pgfqpoint{2.148622in}{1.813278in}}%
\pgfpathlineto{\pgfqpoint{2.148148in}{1.837483in}}%
\pgfpathlineto{\pgfqpoint{2.156256in}{1.841131in}}%
\pgfpathlineto{\pgfqpoint{2.160005in}{1.848588in}}%
\pgfpathlineto{\pgfqpoint{2.176492in}{1.838639in}}%
\pgfpathlineto{\pgfqpoint{2.186149in}{1.838564in}}%
\pgfpathlineto{\pgfqpoint{2.190082in}{1.834950in}}%
\pgfpathlineto{\pgfqpoint{2.189353in}{1.873824in}}%
\pgfpathlineto{\pgfqpoint{2.226726in}{1.874225in}}%
\pgfpathlineto{\pgfqpoint{2.224641in}{1.861874in}}%
\pgfpathlineto{\pgfqpoint{2.230979in}{1.861346in}}%
\pgfpathlineto{\pgfqpoint{2.238078in}{1.851709in}}%
\pgfpathlineto{\pgfqpoint{2.237996in}{1.848119in}}%
\pgfpathlineto{\pgfqpoint{2.232161in}{1.839416in}}%
\pgfpathlineto{\pgfqpoint{2.232360in}{1.832947in}}%
\pgfpathlineto{\pgfqpoint{2.266038in}{1.833884in}}%
\pgfpathlineto{\pgfqpoint{2.267401in}{1.826914in}}%
\pgfpathlineto{\pgfqpoint{2.268903in}{1.780718in}}%
\pgfpathlineto{\pgfqpoint{2.258667in}{1.780519in}}%
\pgfpathclose%
\pgfusepath{fill}%
\end{pgfscope}%
\begin{pgfscope}%
\pgfpathrectangle{\pgfqpoint{0.100000in}{0.100000in}}{\pgfqpoint{3.608454in}{2.310000in}}%
\pgfusepath{clip}%
\pgfsetbuttcap%
\pgfsetmiterjoin%
\definecolor{currentfill}{rgb}{0.000000,0.482353,0.758824}%
\pgfsetfillcolor{currentfill}%
\pgfsetlinewidth{0.000000pt}%
\definecolor{currentstroke}{rgb}{0.000000,0.000000,0.000000}%
\pgfsetstrokecolor{currentstroke}%
\pgfsetstrokeopacity{0.000000}%
\pgfsetdash{}{0pt}%
\pgfpathmoveto{\pgfqpoint{2.424922in}{0.713419in}}%
\pgfpathlineto{\pgfqpoint{2.432451in}{0.691509in}}%
\pgfpathlineto{\pgfqpoint{2.435277in}{0.651886in}}%
\pgfpathlineto{\pgfqpoint{2.431222in}{0.657222in}}%
\pgfpathlineto{\pgfqpoint{2.419484in}{0.654679in}}%
\pgfpathlineto{\pgfqpoint{2.415540in}{0.649843in}}%
\pgfpathlineto{\pgfqpoint{2.407987in}{0.645878in}}%
\pgfpathlineto{\pgfqpoint{2.389280in}{0.640751in}}%
\pgfpathlineto{\pgfqpoint{2.386115in}{0.636858in}}%
\pgfpathlineto{\pgfqpoint{2.383161in}{0.637744in}}%
\pgfpathlineto{\pgfqpoint{2.368222in}{0.632799in}}%
\pgfpathlineto{\pgfqpoint{2.351410in}{0.638299in}}%
\pgfpathlineto{\pgfqpoint{2.350223in}{0.648417in}}%
\pgfpathlineto{\pgfqpoint{2.344148in}{0.649181in}}%
\pgfpathlineto{\pgfqpoint{2.339922in}{0.654462in}}%
\pgfpathlineto{\pgfqpoint{2.338216in}{0.665019in}}%
\pgfpathlineto{\pgfqpoint{2.331431in}{0.670010in}}%
\pgfpathlineto{\pgfqpoint{2.329801in}{0.674430in}}%
\pgfpathlineto{\pgfqpoint{2.331030in}{0.685463in}}%
\pgfpathlineto{\pgfqpoint{2.325022in}{0.696944in}}%
\pgfpathlineto{\pgfqpoint{2.325328in}{0.707042in}}%
\pgfpathlineto{\pgfqpoint{2.330722in}{0.712813in}}%
\pgfpathlineto{\pgfqpoint{2.336789in}{0.709454in}}%
\pgfpathlineto{\pgfqpoint{2.392212in}{0.711750in}}%
\pgfpathlineto{\pgfqpoint{2.424922in}{0.713419in}}%
\pgfpathclose%
\pgfusepath{fill}%
\end{pgfscope}%
\begin{pgfscope}%
\pgfpathrectangle{\pgfqpoint{0.100000in}{0.100000in}}{\pgfqpoint{3.608454in}{2.310000in}}%
\pgfusepath{clip}%
\pgfsetbuttcap%
\pgfsetmiterjoin%
\definecolor{currentfill}{rgb}{0.000000,0.592157,0.703922}%
\pgfsetfillcolor{currentfill}%
\pgfsetlinewidth{0.000000pt}%
\definecolor{currentstroke}{rgb}{0.000000,0.000000,0.000000}%
\pgfsetstrokecolor{currentstroke}%
\pgfsetstrokeopacity{0.000000}%
\pgfsetdash{}{0pt}%
\pgfpathmoveto{\pgfqpoint{1.895416in}{1.185262in}}%
\pgfpathlineto{\pgfqpoint{1.854168in}{1.187060in}}%
\pgfpathlineto{\pgfqpoint{1.855479in}{1.217455in}}%
\pgfpathlineto{\pgfqpoint{1.821457in}{1.218805in}}%
\pgfpathlineto{\pgfqpoint{1.821831in}{1.246792in}}%
\pgfpathlineto{\pgfqpoint{1.855927in}{1.245032in}}%
\pgfpathlineto{\pgfqpoint{1.856262in}{1.252307in}}%
\pgfpathlineto{\pgfqpoint{1.890012in}{1.250882in}}%
\pgfpathlineto{\pgfqpoint{1.890205in}{1.243605in}}%
\pgfpathlineto{\pgfqpoint{1.889149in}{1.216054in}}%
\pgfpathlineto{\pgfqpoint{1.896376in}{1.215785in}}%
\pgfpathlineto{\pgfqpoint{1.895416in}{1.185262in}}%
\pgfpathclose%
\pgfusepath{fill}%
\end{pgfscope}%
\begin{pgfscope}%
\pgfpathrectangle{\pgfqpoint{0.100000in}{0.100000in}}{\pgfqpoint{3.608454in}{2.310000in}}%
\pgfusepath{clip}%
\pgfsetbuttcap%
\pgfsetmiterjoin%
\definecolor{currentfill}{rgb}{0.000000,0.427451,0.786275}%
\pgfsetfillcolor{currentfill}%
\pgfsetlinewidth{0.000000pt}%
\definecolor{currentstroke}{rgb}{0.000000,0.000000,0.000000}%
\pgfsetstrokecolor{currentstroke}%
\pgfsetstrokeopacity{0.000000}%
\pgfsetdash{}{0pt}%
\pgfpathmoveto{\pgfqpoint{2.499443in}{1.713874in}}%
\pgfpathlineto{\pgfqpoint{2.485729in}{1.712888in}}%
\pgfpathlineto{\pgfqpoint{2.485186in}{1.719799in}}%
\pgfpathlineto{\pgfqpoint{2.450492in}{1.717550in}}%
\pgfpathlineto{\pgfqpoint{2.436777in}{1.717456in}}%
\pgfpathlineto{\pgfqpoint{2.436156in}{1.727226in}}%
\pgfpathlineto{\pgfqpoint{2.440017in}{1.736230in}}%
\pgfpathlineto{\pgfqpoint{2.439466in}{1.744650in}}%
\pgfpathlineto{\pgfqpoint{2.414986in}{1.743088in}}%
\pgfpathlineto{\pgfqpoint{2.413744in}{1.763959in}}%
\pgfpathlineto{\pgfqpoint{2.434927in}{1.765090in}}%
\pgfpathlineto{\pgfqpoint{2.432784in}{1.799789in}}%
\pgfpathlineto{\pgfqpoint{2.467658in}{1.801883in}}%
\pgfpathlineto{\pgfqpoint{2.468096in}{1.794922in}}%
\pgfpathlineto{\pgfqpoint{2.488486in}{1.795943in}}%
\pgfpathlineto{\pgfqpoint{2.491553in}{1.796332in}}%
\pgfpathlineto{\pgfqpoint{2.493489in}{1.769026in}}%
\pgfpathlineto{\pgfqpoint{2.501961in}{1.769586in}}%
\pgfpathlineto{\pgfqpoint{2.504166in}{1.741961in}}%
\pgfpathlineto{\pgfqpoint{2.497309in}{1.741424in}}%
\pgfpathlineto{\pgfqpoint{2.499443in}{1.713874in}}%
\pgfpathclose%
\pgfusepath{fill}%
\end{pgfscope}%
\begin{pgfscope}%
\pgfpathrectangle{\pgfqpoint{0.100000in}{0.100000in}}{\pgfqpoint{3.608454in}{2.310000in}}%
\pgfusepath{clip}%
\pgfsetbuttcap%
\pgfsetmiterjoin%
\definecolor{currentfill}{rgb}{0.000000,0.364706,0.817647}%
\pgfsetfillcolor{currentfill}%
\pgfsetlinewidth{0.000000pt}%
\definecolor{currentstroke}{rgb}{0.000000,0.000000,0.000000}%
\pgfsetstrokecolor{currentstroke}%
\pgfsetstrokeopacity{0.000000}%
\pgfsetdash{}{0pt}%
\pgfpathmoveto{\pgfqpoint{3.147047in}{0.962987in}}%
\pgfpathlineto{\pgfqpoint{3.144033in}{0.954523in}}%
\pgfpathlineto{\pgfqpoint{3.134057in}{0.954802in}}%
\pgfpathlineto{\pgfqpoint{3.127866in}{0.947248in}}%
\pgfpathlineto{\pgfqpoint{3.131225in}{0.944243in}}%
\pgfpathlineto{\pgfqpoint{3.123999in}{0.938426in}}%
\pgfpathlineto{\pgfqpoint{3.122700in}{0.934730in}}%
\pgfpathlineto{\pgfqpoint{3.114162in}{0.927752in}}%
\pgfpathlineto{\pgfqpoint{3.113774in}{0.923333in}}%
\pgfpathlineto{\pgfqpoint{3.106625in}{0.916011in}}%
\pgfpathlineto{\pgfqpoint{3.099237in}{0.913335in}}%
\pgfpathlineto{\pgfqpoint{3.086793in}{0.902214in}}%
\pgfpathlineto{\pgfqpoint{3.077208in}{0.901258in}}%
\pgfpathlineto{\pgfqpoint{3.071225in}{0.906814in}}%
\pgfpathlineto{\pgfqpoint{3.066440in}{0.907291in}}%
\pgfpathlineto{\pgfqpoint{3.061084in}{0.913825in}}%
\pgfpathlineto{\pgfqpoint{3.051321in}{0.914811in}}%
\pgfpathlineto{\pgfqpoint{3.044811in}{0.924182in}}%
\pgfpathlineto{\pgfqpoint{3.038279in}{0.928637in}}%
\pgfpathlineto{\pgfqpoint{3.030466in}{0.937480in}}%
\pgfpathlineto{\pgfqpoint{3.047350in}{0.952216in}}%
\pgfpathlineto{\pgfqpoint{3.064743in}{0.967659in}}%
\pgfpathlineto{\pgfqpoint{3.070676in}{0.962442in}}%
\pgfpathlineto{\pgfqpoint{3.081498in}{0.967552in}}%
\pgfpathlineto{\pgfqpoint{3.081713in}{0.979204in}}%
\pgfpathlineto{\pgfqpoint{3.088906in}{0.984668in}}%
\pgfpathlineto{\pgfqpoint{3.097264in}{0.986382in}}%
\pgfpathlineto{\pgfqpoint{3.115214in}{0.974844in}}%
\pgfpathlineto{\pgfqpoint{3.134423in}{0.971285in}}%
\pgfpathlineto{\pgfqpoint{3.147047in}{0.962987in}}%
\pgfpathclose%
\pgfusepath{fill}%
\end{pgfscope}%
\begin{pgfscope}%
\pgfpathrectangle{\pgfqpoint{0.100000in}{0.100000in}}{\pgfqpoint{3.608454in}{2.310000in}}%
\pgfusepath{clip}%
\pgfsetbuttcap%
\pgfsetmiterjoin%
\definecolor{currentfill}{rgb}{0.000000,0.356863,0.821569}%
\pgfsetfillcolor{currentfill}%
\pgfsetlinewidth{0.000000pt}%
\definecolor{currentstroke}{rgb}{0.000000,0.000000,0.000000}%
\pgfsetstrokecolor{currentstroke}%
\pgfsetstrokeopacity{0.000000}%
\pgfsetdash{}{0pt}%
\pgfpathmoveto{\pgfqpoint{3.191049in}{1.170900in}}%
\pgfpathlineto{\pgfqpoint{3.187783in}{1.156766in}}%
\pgfpathlineto{\pgfqpoint{3.188860in}{1.151551in}}%
\pgfpathlineto{\pgfqpoint{3.180850in}{1.147206in}}%
\pgfpathlineto{\pgfqpoint{3.173892in}{1.141190in}}%
\pgfpathlineto{\pgfqpoint{3.168634in}{1.140787in}}%
\pgfpathlineto{\pgfqpoint{3.164716in}{1.144245in}}%
\pgfpathlineto{\pgfqpoint{3.160885in}{1.137946in}}%
\pgfpathlineto{\pgfqpoint{3.149312in}{1.137485in}}%
\pgfpathlineto{\pgfqpoint{3.131128in}{1.128295in}}%
\pgfpathlineto{\pgfqpoint{3.121994in}{1.132796in}}%
\pgfpathlineto{\pgfqpoint{3.110689in}{1.147372in}}%
\pgfpathlineto{\pgfqpoint{3.110510in}{1.150879in}}%
\pgfpathlineto{\pgfqpoint{3.097477in}{1.148400in}}%
\pgfpathlineto{\pgfqpoint{3.093928in}{1.174159in}}%
\pgfpathlineto{\pgfqpoint{3.113209in}{1.177564in}}%
\pgfpathlineto{\pgfqpoint{3.110627in}{1.182204in}}%
\pgfpathlineto{\pgfqpoint{3.106519in}{1.208514in}}%
\pgfpathlineto{\pgfqpoint{3.113032in}{1.209517in}}%
\pgfpathlineto{\pgfqpoint{3.109908in}{1.233039in}}%
\pgfpathlineto{\pgfqpoint{3.152496in}{1.240772in}}%
\pgfpathlineto{\pgfqpoint{3.150361in}{1.233286in}}%
\pgfpathlineto{\pgfqpoint{3.154402in}{1.212187in}}%
\pgfpathlineto{\pgfqpoint{3.160658in}{1.212280in}}%
\pgfpathlineto{\pgfqpoint{3.165876in}{1.221007in}}%
\pgfpathlineto{\pgfqpoint{3.177079in}{1.221508in}}%
\pgfpathlineto{\pgfqpoint{3.185638in}{1.219606in}}%
\pgfpathlineto{\pgfqpoint{3.175634in}{1.186715in}}%
\pgfpathlineto{\pgfqpoint{3.181671in}{1.180528in}}%
\pgfpathlineto{\pgfqpoint{3.186933in}{1.171546in}}%
\pgfpathlineto{\pgfqpoint{3.191049in}{1.170900in}}%
\pgfpathclose%
\pgfusepath{fill}%
\end{pgfscope}%
\begin{pgfscope}%
\pgfpathrectangle{\pgfqpoint{0.100000in}{0.100000in}}{\pgfqpoint{3.608454in}{2.310000in}}%
\pgfusepath{clip}%
\pgfsetbuttcap%
\pgfsetmiterjoin%
\definecolor{currentfill}{rgb}{0.000000,0.505882,0.747059}%
\pgfsetfillcolor{currentfill}%
\pgfsetlinewidth{0.000000pt}%
\definecolor{currentstroke}{rgb}{0.000000,0.000000,0.000000}%
\pgfsetstrokecolor{currentstroke}%
\pgfsetstrokeopacity{0.000000}%
\pgfsetdash{}{0pt}%
\pgfpathmoveto{\pgfqpoint{1.881222in}{1.423490in}}%
\pgfpathlineto{\pgfqpoint{1.853837in}{1.424663in}}%
\pgfpathlineto{\pgfqpoint{1.856208in}{1.476676in}}%
\pgfpathlineto{\pgfqpoint{1.841990in}{1.478161in}}%
\pgfpathlineto{\pgfqpoint{1.842917in}{1.507856in}}%
\pgfpathlineto{\pgfqpoint{1.855925in}{1.507277in}}%
\pgfpathlineto{\pgfqpoint{1.857457in}{1.534695in}}%
\pgfpathlineto{\pgfqpoint{1.884664in}{1.533508in}}%
\pgfpathlineto{\pgfqpoint{1.883282in}{1.506053in}}%
\pgfpathlineto{\pgfqpoint{1.884864in}{1.505981in}}%
\pgfpathlineto{\pgfqpoint{1.883772in}{1.478524in}}%
\pgfpathlineto{\pgfqpoint{1.881222in}{1.423490in}}%
\pgfpathclose%
\pgfusepath{fill}%
\end{pgfscope}%
\begin{pgfscope}%
\pgfpathrectangle{\pgfqpoint{0.100000in}{0.100000in}}{\pgfqpoint{3.608454in}{2.310000in}}%
\pgfusepath{clip}%
\pgfsetbuttcap%
\pgfsetmiterjoin%
\definecolor{currentfill}{rgb}{0.000000,0.423529,0.788235}%
\pgfsetfillcolor{currentfill}%
\pgfsetlinewidth{0.000000pt}%
\definecolor{currentstroke}{rgb}{0.000000,0.000000,0.000000}%
\pgfsetstrokecolor{currentstroke}%
\pgfsetstrokeopacity{0.000000}%
\pgfsetdash{}{0pt}%
\pgfpathmoveto{\pgfqpoint{0.836288in}{2.219545in}}%
\pgfpathlineto{\pgfqpoint{0.829989in}{2.221231in}}%
\pgfpathlineto{\pgfqpoint{0.822836in}{2.217028in}}%
\pgfpathlineto{\pgfqpoint{0.820986in}{2.210320in}}%
\pgfpathlineto{\pgfqpoint{0.814986in}{2.205959in}}%
\pgfpathlineto{\pgfqpoint{0.810303in}{2.197629in}}%
\pgfpathlineto{\pgfqpoint{0.807006in}{2.198546in}}%
\pgfpathlineto{\pgfqpoint{0.798560in}{2.193839in}}%
\pgfpathlineto{\pgfqpoint{0.796683in}{2.187200in}}%
\pgfpathlineto{\pgfqpoint{0.778779in}{2.191862in}}%
\pgfpathlineto{\pgfqpoint{0.774962in}{2.186885in}}%
\pgfpathlineto{\pgfqpoint{0.769524in}{2.184730in}}%
\pgfpathlineto{\pgfqpoint{0.767233in}{2.177075in}}%
\pgfpathlineto{\pgfqpoint{0.763587in}{2.181602in}}%
\pgfpathlineto{\pgfqpoint{0.748201in}{2.185881in}}%
\pgfpathlineto{\pgfqpoint{0.742925in}{2.192347in}}%
\pgfpathlineto{\pgfqpoint{0.735660in}{2.196356in}}%
\pgfpathlineto{\pgfqpoint{0.733242in}{2.201231in}}%
\pgfpathlineto{\pgfqpoint{0.724213in}{2.207197in}}%
\pgfpathlineto{\pgfqpoint{0.722390in}{2.214825in}}%
\pgfpathlineto{\pgfqpoint{0.718104in}{2.222638in}}%
\pgfpathlineto{\pgfqpoint{0.722977in}{2.230297in}}%
\pgfpathlineto{\pgfqpoint{0.721489in}{2.246797in}}%
\pgfpathlineto{\pgfqpoint{0.727318in}{2.257197in}}%
\pgfpathlineto{\pgfqpoint{0.738201in}{2.260046in}}%
\pgfpathlineto{\pgfqpoint{0.741543in}{2.265250in}}%
\pgfpathlineto{\pgfqpoint{0.736966in}{2.278404in}}%
\pgfpathlineto{\pgfqpoint{0.740255in}{2.288121in}}%
\pgfpathlineto{\pgfqpoint{0.752759in}{2.291869in}}%
\pgfpathlineto{\pgfqpoint{0.760304in}{2.288239in}}%
\pgfpathlineto{\pgfqpoint{0.760184in}{2.281456in}}%
\pgfpathlineto{\pgfqpoint{0.765081in}{2.268518in}}%
\pgfpathlineto{\pgfqpoint{0.771087in}{2.257398in}}%
\pgfpathlineto{\pgfqpoint{0.769468in}{2.254538in}}%
\pgfpathlineto{\pgfqpoint{0.780160in}{2.239854in}}%
\pgfpathlineto{\pgfqpoint{0.780931in}{2.234819in}}%
\pgfpathlineto{\pgfqpoint{0.790372in}{2.232053in}}%
\pgfpathlineto{\pgfqpoint{0.791522in}{2.239767in}}%
\pgfpathlineto{\pgfqpoint{0.800018in}{2.241064in}}%
\pgfpathlineto{\pgfqpoint{0.802328in}{2.232125in}}%
\pgfpathlineto{\pgfqpoint{0.815034in}{2.234984in}}%
\pgfpathlineto{\pgfqpoint{0.819597in}{2.231350in}}%
\pgfpathlineto{\pgfqpoint{0.832614in}{2.236217in}}%
\pgfpathlineto{\pgfqpoint{0.836249in}{2.233779in}}%
\pgfpathlineto{\pgfqpoint{0.836288in}{2.219545in}}%
\pgfpathclose%
\pgfusepath{fill}%
\end{pgfscope}%
\begin{pgfscope}%
\pgfpathrectangle{\pgfqpoint{0.100000in}{0.100000in}}{\pgfqpoint{3.608454in}{2.310000in}}%
\pgfusepath{clip}%
\pgfsetbuttcap%
\pgfsetmiterjoin%
\definecolor{currentfill}{rgb}{0.000000,0.631373,0.684314}%
\pgfsetfillcolor{currentfill}%
\pgfsetlinewidth{0.000000pt}%
\definecolor{currentstroke}{rgb}{0.000000,0.000000,0.000000}%
\pgfsetstrokecolor{currentstroke}%
\pgfsetstrokeopacity{0.000000}%
\pgfsetdash{}{0pt}%
\pgfpathmoveto{\pgfqpoint{2.485040in}{0.651337in}}%
\pgfpathlineto{\pgfqpoint{2.478549in}{0.654208in}}%
\pgfpathlineto{\pgfqpoint{2.472822in}{0.672176in}}%
\pgfpathlineto{\pgfqpoint{2.466743in}{0.676753in}}%
\pgfpathlineto{\pgfqpoint{2.460731in}{0.688378in}}%
\pgfpathlineto{\pgfqpoint{2.461305in}{0.698015in}}%
\pgfpathlineto{\pgfqpoint{2.465118in}{0.704833in}}%
\pgfpathlineto{\pgfqpoint{2.467199in}{0.715885in}}%
\pgfpathlineto{\pgfqpoint{2.472257in}{0.716197in}}%
\pgfpathlineto{\pgfqpoint{2.482623in}{0.717728in}}%
\pgfpathlineto{\pgfqpoint{2.493658in}{0.717377in}}%
\pgfpathlineto{\pgfqpoint{2.494101in}{0.710117in}}%
\pgfpathlineto{\pgfqpoint{2.528702in}{0.712404in}}%
\pgfpathlineto{\pgfqpoint{2.525381in}{0.754028in}}%
\pgfpathlineto{\pgfqpoint{2.551986in}{0.756063in}}%
\pgfpathlineto{\pgfqpoint{2.556138in}{0.721372in}}%
\pgfpathlineto{\pgfqpoint{2.561760in}{0.672494in}}%
\pgfpathlineto{\pgfqpoint{2.556857in}{0.667105in}}%
\pgfpathlineto{\pgfqpoint{2.547096in}{0.670673in}}%
\pgfpathlineto{\pgfqpoint{2.541328in}{0.668013in}}%
\pgfpathlineto{\pgfqpoint{2.525544in}{0.670845in}}%
\pgfpathlineto{\pgfqpoint{2.515743in}{0.668541in}}%
\pgfpathlineto{\pgfqpoint{2.497858in}{0.661623in}}%
\pgfpathlineto{\pgfqpoint{2.485040in}{0.651337in}}%
\pgfpathclose%
\pgfusepath{fill}%
\end{pgfscope}%
\begin{pgfscope}%
\pgfpathrectangle{\pgfqpoint{0.100000in}{0.100000in}}{\pgfqpoint{3.608454in}{2.310000in}}%
\pgfusepath{clip}%
\pgfsetbuttcap%
\pgfsetmiterjoin%
\definecolor{currentfill}{rgb}{0.000000,0.690196,0.654902}%
\pgfsetfillcolor{currentfill}%
\pgfsetlinewidth{0.000000pt}%
\definecolor{currentstroke}{rgb}{0.000000,0.000000,0.000000}%
\pgfsetstrokecolor{currentstroke}%
\pgfsetstrokeopacity{0.000000}%
\pgfsetdash{}{0pt}%
\pgfpathmoveto{\pgfqpoint{3.568956in}{1.730471in}}%
\pgfpathlineto{\pgfqpoint{3.576435in}{1.723494in}}%
\pgfpathlineto{\pgfqpoint{3.574246in}{1.719609in}}%
\pgfpathlineto{\pgfqpoint{3.565620in}{1.718239in}}%
\pgfpathlineto{\pgfqpoint{3.570850in}{1.726339in}}%
\pgfpathlineto{\pgfqpoint{3.568956in}{1.730471in}}%
\pgfpathclose%
\pgfusepath{fill}%
\end{pgfscope}%
\begin{pgfscope}%
\pgfpathrectangle{\pgfqpoint{0.100000in}{0.100000in}}{\pgfqpoint{3.608454in}{2.310000in}}%
\pgfusepath{clip}%
\pgfsetbuttcap%
\pgfsetmiterjoin%
\definecolor{currentfill}{rgb}{0.000000,0.525490,0.737255}%
\pgfsetfillcolor{currentfill}%
\pgfsetlinewidth{0.000000pt}%
\definecolor{currentstroke}{rgb}{0.000000,0.000000,0.000000}%
\pgfsetstrokecolor{currentstroke}%
\pgfsetstrokeopacity{0.000000}%
\pgfsetdash{}{0pt}%
\pgfpathmoveto{\pgfqpoint{3.024683in}{1.232048in}}%
\pgfpathlineto{\pgfqpoint{3.011272in}{1.247768in}}%
\pgfpathlineto{\pgfqpoint{3.005275in}{1.242589in}}%
\pgfpathlineto{\pgfqpoint{2.993913in}{1.252683in}}%
\pgfpathlineto{\pgfqpoint{2.990512in}{1.259171in}}%
\pgfpathlineto{\pgfqpoint{2.994849in}{1.262443in}}%
\pgfpathlineto{\pgfqpoint{2.983343in}{1.270885in}}%
\pgfpathlineto{\pgfqpoint{2.992726in}{1.278073in}}%
\pgfpathlineto{\pgfqpoint{2.991086in}{1.284432in}}%
\pgfpathlineto{\pgfqpoint{3.000083in}{1.309688in}}%
\pgfpathlineto{\pgfqpoint{3.013170in}{1.308307in}}%
\pgfpathlineto{\pgfqpoint{3.022678in}{1.310484in}}%
\pgfpathlineto{\pgfqpoint{3.033267in}{1.304579in}}%
\pgfpathlineto{\pgfqpoint{3.045914in}{1.315133in}}%
\pgfpathlineto{\pgfqpoint{3.050832in}{1.323908in}}%
\pgfpathlineto{\pgfqpoint{3.060124in}{1.322440in}}%
\pgfpathlineto{\pgfqpoint{3.060878in}{1.314437in}}%
\pgfpathlineto{\pgfqpoint{3.074021in}{1.306357in}}%
\pgfpathlineto{\pgfqpoint{3.067687in}{1.299266in}}%
\pgfpathlineto{\pgfqpoint{3.064782in}{1.300929in}}%
\pgfpathlineto{\pgfqpoint{3.057612in}{1.293606in}}%
\pgfpathlineto{\pgfqpoint{3.055501in}{1.285981in}}%
\pgfpathlineto{\pgfqpoint{3.057510in}{1.278132in}}%
\pgfpathlineto{\pgfqpoint{3.069985in}{1.273083in}}%
\pgfpathlineto{\pgfqpoint{3.074664in}{1.267099in}}%
\pgfpathlineto{\pgfqpoint{3.074241in}{1.252094in}}%
\pgfpathlineto{\pgfqpoint{3.067580in}{1.245522in}}%
\pgfpathlineto{\pgfqpoint{3.036908in}{1.247132in}}%
\pgfpathlineto{\pgfqpoint{3.032148in}{1.244073in}}%
\pgfpathlineto{\pgfqpoint{3.024683in}{1.232048in}}%
\pgfpathclose%
\pgfusepath{fill}%
\end{pgfscope}%
\begin{pgfscope}%
\pgfpathrectangle{\pgfqpoint{0.100000in}{0.100000in}}{\pgfqpoint{3.608454in}{2.310000in}}%
\pgfusepath{clip}%
\pgfsetbuttcap%
\pgfsetmiterjoin%
\definecolor{currentfill}{rgb}{0.000000,0.560784,0.719608}%
\pgfsetfillcolor{currentfill}%
\pgfsetlinewidth{0.000000pt}%
\definecolor{currentstroke}{rgb}{0.000000,0.000000,0.000000}%
\pgfsetstrokecolor{currentstroke}%
\pgfsetstrokeopacity{0.000000}%
\pgfsetdash{}{0pt}%
\pgfpathmoveto{\pgfqpoint{1.267460in}{1.535366in}}%
\pgfpathlineto{\pgfqpoint{1.261144in}{1.538773in}}%
\pgfpathlineto{\pgfqpoint{1.262249in}{1.546769in}}%
\pgfpathlineto{\pgfqpoint{1.260798in}{1.552822in}}%
\pgfpathlineto{\pgfqpoint{1.249496in}{1.554569in}}%
\pgfpathlineto{\pgfqpoint{1.248664in}{1.549321in}}%
\pgfpathlineto{\pgfqpoint{1.242723in}{1.546157in}}%
\pgfpathlineto{\pgfqpoint{1.241154in}{1.553467in}}%
\pgfpathlineto{\pgfqpoint{1.220332in}{1.551758in}}%
\pgfpathlineto{\pgfqpoint{1.213136in}{1.556950in}}%
\pgfpathlineto{\pgfqpoint{1.215522in}{1.571272in}}%
\pgfpathlineto{\pgfqpoint{1.186062in}{1.576054in}}%
\pgfpathlineto{\pgfqpoint{1.154007in}{1.581932in}}%
\pgfpathlineto{\pgfqpoint{1.157543in}{1.601688in}}%
\pgfpathlineto{\pgfqpoint{1.146992in}{1.596437in}}%
\pgfpathlineto{\pgfqpoint{1.143956in}{1.603549in}}%
\pgfpathlineto{\pgfqpoint{1.144102in}{1.610430in}}%
\pgfpathlineto{\pgfqpoint{1.136553in}{1.614347in}}%
\pgfpathlineto{\pgfqpoint{1.132825in}{1.619941in}}%
\pgfpathlineto{\pgfqpoint{1.138324in}{1.627635in}}%
\pgfpathlineto{\pgfqpoint{1.142480in}{1.636751in}}%
\pgfpathlineto{\pgfqpoint{1.141060in}{1.650733in}}%
\pgfpathlineto{\pgfqpoint{1.139137in}{1.653332in}}%
\pgfpathlineto{\pgfqpoint{1.142493in}{1.660281in}}%
\pgfpathlineto{\pgfqpoint{1.141273in}{1.664762in}}%
\pgfpathlineto{\pgfqpoint{1.167956in}{1.660095in}}%
\pgfpathlineto{\pgfqpoint{1.182268in}{1.739464in}}%
\pgfpathlineto{\pgfqpoint{1.186356in}{1.762535in}}%
\pgfpathlineto{\pgfqpoint{1.199226in}{1.758954in}}%
\pgfpathlineto{\pgfqpoint{1.198346in}{1.753929in}}%
\pgfpathlineto{\pgfqpoint{1.211732in}{1.751597in}}%
\pgfpathlineto{\pgfqpoint{1.212511in}{1.756000in}}%
\pgfpathlineto{\pgfqpoint{1.225832in}{1.753673in}}%
\pgfpathlineto{\pgfqpoint{1.226972in}{1.760454in}}%
\pgfpathlineto{\pgfqpoint{1.243269in}{1.757690in}}%
\pgfpathlineto{\pgfqpoint{1.244733in}{1.764439in}}%
\pgfpathlineto{\pgfqpoint{1.261563in}{1.761526in}}%
\pgfpathlineto{\pgfqpoint{1.263755in}{1.753309in}}%
\pgfpathlineto{\pgfqpoint{1.262997in}{1.742086in}}%
\pgfpathlineto{\pgfqpoint{1.264849in}{1.733953in}}%
\pgfpathlineto{\pgfqpoint{1.270191in}{1.723778in}}%
\pgfpathlineto{\pgfqpoint{1.275945in}{1.717107in}}%
\pgfpathlineto{\pgfqpoint{1.282465in}{1.703069in}}%
\pgfpathlineto{\pgfqpoint{1.290897in}{1.695072in}}%
\pgfpathlineto{\pgfqpoint{1.287915in}{1.674697in}}%
\pgfpathlineto{\pgfqpoint{1.287325in}{1.661183in}}%
\pgfpathlineto{\pgfqpoint{1.347746in}{1.651996in}}%
\pgfpathlineto{\pgfqpoint{1.375467in}{1.648053in}}%
\pgfpathlineto{\pgfqpoint{1.374676in}{1.634274in}}%
\pgfpathlineto{\pgfqpoint{1.369658in}{1.600576in}}%
\pgfpathlineto{\pgfqpoint{1.345097in}{1.604208in}}%
\pgfpathlineto{\pgfqpoint{1.342910in}{1.583654in}}%
\pgfpathlineto{\pgfqpoint{1.338390in}{1.552623in}}%
\pgfpathlineto{\pgfqpoint{1.302483in}{1.557690in}}%
\pgfpathlineto{\pgfqpoint{1.271578in}{1.562500in}}%
\pgfpathlineto{\pgfqpoint{1.267460in}{1.535366in}}%
\pgfpathclose%
\pgfusepath{fill}%
\end{pgfscope}%
\begin{pgfscope}%
\pgfpathrectangle{\pgfqpoint{0.100000in}{0.100000in}}{\pgfqpoint{3.608454in}{2.310000in}}%
\pgfusepath{clip}%
\pgfsetbuttcap%
\pgfsetmiterjoin%
\definecolor{currentfill}{rgb}{0.000000,0.439216,0.780392}%
\pgfsetfillcolor{currentfill}%
\pgfsetlinewidth{0.000000pt}%
\definecolor{currentstroke}{rgb}{0.000000,0.000000,0.000000}%
\pgfsetstrokecolor{currentstroke}%
\pgfsetstrokeopacity{0.000000}%
\pgfsetdash{}{0pt}%
\pgfpathmoveto{\pgfqpoint{2.520499in}{1.715478in}}%
\pgfpathlineto{\pgfqpoint{2.499443in}{1.713874in}}%
\pgfpathlineto{\pgfqpoint{2.497309in}{1.741424in}}%
\pgfpathlineto{\pgfqpoint{2.504166in}{1.741961in}}%
\pgfpathlineto{\pgfqpoint{2.501961in}{1.769586in}}%
\pgfpathlineto{\pgfqpoint{2.510762in}{1.770224in}}%
\pgfpathlineto{\pgfqpoint{2.510243in}{1.777129in}}%
\pgfpathlineto{\pgfqpoint{2.529824in}{1.778725in}}%
\pgfpathlineto{\pgfqpoint{2.532257in}{1.767832in}}%
\pgfpathlineto{\pgfqpoint{2.525321in}{1.760582in}}%
\pgfpathlineto{\pgfqpoint{2.521671in}{1.742517in}}%
\pgfpathlineto{\pgfqpoint{2.524790in}{1.727344in}}%
\pgfpathlineto{\pgfqpoint{2.520499in}{1.715478in}}%
\pgfpathclose%
\pgfusepath{fill}%
\end{pgfscope}%
\begin{pgfscope}%
\pgfpathrectangle{\pgfqpoint{0.100000in}{0.100000in}}{\pgfqpoint{3.608454in}{2.310000in}}%
\pgfusepath{clip}%
\pgfsetbuttcap%
\pgfsetmiterjoin%
\definecolor{currentfill}{rgb}{0.000000,0.600000,0.700000}%
\pgfsetfillcolor{currentfill}%
\pgfsetlinewidth{0.000000pt}%
\definecolor{currentstroke}{rgb}{0.000000,0.000000,0.000000}%
\pgfsetstrokecolor{currentstroke}%
\pgfsetstrokeopacity{0.000000}%
\pgfsetdash{}{0pt}%
\pgfpathmoveto{\pgfqpoint{1.576920in}{0.953485in}}%
\pgfpathlineto{\pgfqpoint{1.580563in}{0.992038in}}%
\pgfpathlineto{\pgfqpoint{1.614395in}{0.989155in}}%
\pgfpathlineto{\pgfqpoint{1.617270in}{1.023437in}}%
\pgfpathlineto{\pgfqpoint{1.651469in}{1.020799in}}%
\pgfpathlineto{\pgfqpoint{1.648845in}{0.986397in}}%
\pgfpathlineto{\pgfqpoint{1.642812in}{0.986873in}}%
\pgfpathlineto{\pgfqpoint{1.640100in}{0.948213in}}%
\pgfpathlineto{\pgfqpoint{1.576920in}{0.953485in}}%
\pgfpathclose%
\pgfusepath{fill}%
\end{pgfscope}%
\begin{pgfscope}%
\pgfpathrectangle{\pgfqpoint{0.100000in}{0.100000in}}{\pgfqpoint{3.608454in}{2.310000in}}%
\pgfusepath{clip}%
\pgfsetbuttcap%
\pgfsetmiterjoin%
\definecolor{currentfill}{rgb}{0.000000,0.792157,0.603922}%
\pgfsetfillcolor{currentfill}%
\pgfsetlinewidth{0.000000pt}%
\definecolor{currentstroke}{rgb}{0.000000,0.000000,0.000000}%
\pgfsetstrokecolor{currentstroke}%
\pgfsetstrokeopacity{0.000000}%
\pgfsetdash{}{0pt}%
\pgfpathmoveto{\pgfqpoint{0.557819in}{2.120432in}}%
\pgfpathlineto{\pgfqpoint{0.538739in}{2.126084in}}%
\pgfpathlineto{\pgfqpoint{0.525797in}{2.131179in}}%
\pgfpathlineto{\pgfqpoint{0.534887in}{2.152356in}}%
\pgfpathlineto{\pgfqpoint{0.534769in}{2.164166in}}%
\pgfpathlineto{\pgfqpoint{0.540718in}{2.157678in}}%
\pgfpathlineto{\pgfqpoint{0.547629in}{2.160026in}}%
\pgfpathlineto{\pgfqpoint{0.555579in}{2.159647in}}%
\pgfpathlineto{\pgfqpoint{0.552341in}{2.166159in}}%
\pgfpathlineto{\pgfqpoint{0.542113in}{2.164053in}}%
\pgfpathlineto{\pgfqpoint{0.534487in}{2.174240in}}%
\pgfpathlineto{\pgfqpoint{0.538478in}{2.184732in}}%
\pgfpathlineto{\pgfqpoint{0.548088in}{2.185551in}}%
\pgfpathlineto{\pgfqpoint{0.547127in}{2.201192in}}%
\pgfpathlineto{\pgfqpoint{0.543078in}{2.205524in}}%
\pgfpathlineto{\pgfqpoint{0.545221in}{2.218928in}}%
\pgfpathlineto{\pgfqpoint{0.549339in}{2.216444in}}%
\pgfpathlineto{\pgfqpoint{0.559409in}{2.219592in}}%
\pgfpathlineto{\pgfqpoint{0.547351in}{2.228134in}}%
\pgfpathlineto{\pgfqpoint{0.549613in}{2.249680in}}%
\pgfpathlineto{\pgfqpoint{0.547461in}{2.260324in}}%
\pgfpathlineto{\pgfqpoint{0.549782in}{2.270356in}}%
\pgfpathlineto{\pgfqpoint{0.565186in}{2.263908in}}%
\pgfpathlineto{\pgfqpoint{0.593017in}{2.254746in}}%
\pgfpathlineto{\pgfqpoint{0.586531in}{2.235034in}}%
\pgfpathlineto{\pgfqpoint{0.583086in}{2.221440in}}%
\pgfpathlineto{\pgfqpoint{0.597993in}{2.216816in}}%
\pgfpathlineto{\pgfqpoint{0.598094in}{2.209351in}}%
\pgfpathlineto{\pgfqpoint{0.593046in}{2.193971in}}%
\pgfpathlineto{\pgfqpoint{0.582091in}{2.197407in}}%
\pgfpathlineto{\pgfqpoint{0.572748in}{2.166216in}}%
\pgfpathlineto{\pgfqpoint{0.631314in}{2.147810in}}%
\pgfpathlineto{\pgfqpoint{0.623310in}{2.122649in}}%
\pgfpathlineto{\pgfqpoint{0.615173in}{2.117075in}}%
\pgfpathlineto{\pgfqpoint{0.609486in}{2.121452in}}%
\pgfpathlineto{\pgfqpoint{0.593066in}{2.119148in}}%
\pgfpathlineto{\pgfqpoint{0.589293in}{2.135786in}}%
\pgfpathlineto{\pgfqpoint{0.578043in}{2.148170in}}%
\pgfpathlineto{\pgfqpoint{0.566644in}{2.148266in}}%
\pgfpathlineto{\pgfqpoint{0.557819in}{2.120432in}}%
\pgfpathclose%
\pgfusepath{fill}%
\end{pgfscope}%
\begin{pgfscope}%
\pgfpathrectangle{\pgfqpoint{0.100000in}{0.100000in}}{\pgfqpoint{3.608454in}{2.310000in}}%
\pgfusepath{clip}%
\pgfsetbuttcap%
\pgfsetmiterjoin%
\definecolor{currentfill}{rgb}{0.000000,0.839216,0.580392}%
\pgfsetfillcolor{currentfill}%
\pgfsetlinewidth{0.000000pt}%
\definecolor{currentstroke}{rgb}{0.000000,0.000000,0.000000}%
\pgfsetstrokecolor{currentstroke}%
\pgfsetstrokeopacity{0.000000}%
\pgfsetdash{}{0pt}%
\pgfpathmoveto{\pgfqpoint{0.954967in}{2.104591in}}%
\pgfpathlineto{\pgfqpoint{0.949294in}{2.080795in}}%
\pgfpathlineto{\pgfqpoint{0.933769in}{2.084665in}}%
\pgfpathlineto{\pgfqpoint{0.926998in}{2.079195in}}%
\pgfpathlineto{\pgfqpoint{0.909838in}{2.083326in}}%
\pgfpathlineto{\pgfqpoint{0.907603in}{2.074322in}}%
\pgfpathlineto{\pgfqpoint{0.899012in}{2.076012in}}%
\pgfpathlineto{\pgfqpoint{0.898511in}{2.086388in}}%
\pgfpathlineto{\pgfqpoint{0.892918in}{2.093767in}}%
\pgfpathlineto{\pgfqpoint{0.890708in}{2.100407in}}%
\pgfpathlineto{\pgfqpoint{0.882536in}{2.100932in}}%
\pgfpathlineto{\pgfqpoint{0.878570in}{2.104452in}}%
\pgfpathlineto{\pgfqpoint{0.868296in}{2.100626in}}%
\pgfpathlineto{\pgfqpoint{0.849927in}{2.099792in}}%
\pgfpathlineto{\pgfqpoint{0.849420in}{2.114482in}}%
\pgfpathlineto{\pgfqpoint{0.851613in}{2.114354in}}%
\pgfpathlineto{\pgfqpoint{0.863100in}{2.115696in}}%
\pgfpathlineto{\pgfqpoint{0.868405in}{2.124498in}}%
\pgfpathlineto{\pgfqpoint{0.875234in}{2.151233in}}%
\pgfpathlineto{\pgfqpoint{0.923573in}{2.138917in}}%
\pgfpathlineto{\pgfqpoint{0.921056in}{2.128678in}}%
\pgfpathlineto{\pgfqpoint{0.933521in}{2.118231in}}%
\pgfpathlineto{\pgfqpoint{0.950036in}{2.114158in}}%
\pgfpathlineto{\pgfqpoint{0.956600in}{2.111418in}}%
\pgfpathlineto{\pgfqpoint{0.954967in}{2.104591in}}%
\pgfpathclose%
\pgfusepath{fill}%
\end{pgfscope}%
\begin{pgfscope}%
\pgfpathrectangle{\pgfqpoint{0.100000in}{0.100000in}}{\pgfqpoint{3.608454in}{2.310000in}}%
\pgfusepath{clip}%
\pgfsetbuttcap%
\pgfsetmiterjoin%
\definecolor{currentfill}{rgb}{0.000000,0.698039,0.650980}%
\pgfsetfillcolor{currentfill}%
\pgfsetlinewidth{0.000000pt}%
\definecolor{currentstroke}{rgb}{0.000000,0.000000,0.000000}%
\pgfsetstrokecolor{currentstroke}%
\pgfsetstrokeopacity{0.000000}%
\pgfsetdash{}{0pt}%
\pgfpathmoveto{\pgfqpoint{0.877772in}{0.912470in}}%
\pgfpathlineto{\pgfqpoint{0.885192in}{0.948795in}}%
\pgfpathlineto{\pgfqpoint{0.934927in}{0.938858in}}%
\pgfpathlineto{\pgfqpoint{1.001421in}{0.926345in}}%
\pgfpathlineto{\pgfqpoint{1.075355in}{0.913752in}}%
\pgfpathlineto{\pgfqpoint{1.074172in}{0.906976in}}%
\pgfpathlineto{\pgfqpoint{1.092774in}{0.903733in}}%
\pgfpathlineto{\pgfqpoint{1.167197in}{0.891666in}}%
\pgfpathlineto{\pgfqpoint{1.153619in}{0.805601in}}%
\pgfpathlineto{\pgfqpoint{1.089707in}{0.815956in}}%
\pgfpathlineto{\pgfqpoint{1.017594in}{0.828333in}}%
\pgfpathlineto{\pgfqpoint{0.999378in}{0.839223in}}%
\pgfpathlineto{\pgfqpoint{0.900066in}{0.898746in}}%
\pgfpathlineto{\pgfqpoint{0.877772in}{0.912470in}}%
\pgfpathclose%
\pgfusepath{fill}%
\end{pgfscope}%
\begin{pgfscope}%
\pgfpathrectangle{\pgfqpoint{0.100000in}{0.100000in}}{\pgfqpoint{3.608454in}{2.310000in}}%
\pgfusepath{clip}%
\pgfsetbuttcap%
\pgfsetmiterjoin%
\definecolor{currentfill}{rgb}{0.000000,0.439216,0.780392}%
\pgfsetfillcolor{currentfill}%
\pgfsetlinewidth{0.000000pt}%
\definecolor{currentstroke}{rgb}{0.000000,0.000000,0.000000}%
\pgfsetstrokecolor{currentstroke}%
\pgfsetstrokeopacity{0.000000}%
\pgfsetdash{}{0pt}%
\pgfpathmoveto{\pgfqpoint{3.153317in}{1.449269in}}%
\pgfpathlineto{\pgfqpoint{3.152017in}{1.444895in}}%
\pgfpathlineto{\pgfqpoint{3.144962in}{1.435381in}}%
\pgfpathlineto{\pgfqpoint{3.141078in}{1.437828in}}%
\pgfpathlineto{\pgfqpoint{3.125843in}{1.434080in}}%
\pgfpathlineto{\pgfqpoint{3.119217in}{1.439782in}}%
\pgfpathlineto{\pgfqpoint{3.117131in}{1.434042in}}%
\pgfpathlineto{\pgfqpoint{3.111250in}{1.435310in}}%
\pgfpathlineto{\pgfqpoint{3.105576in}{1.439244in}}%
\pgfpathlineto{\pgfqpoint{3.087837in}{1.442494in}}%
\pgfpathlineto{\pgfqpoint{3.082421in}{1.444639in}}%
\pgfpathlineto{\pgfqpoint{3.088891in}{1.453283in}}%
\pgfpathlineto{\pgfqpoint{3.091266in}{1.461236in}}%
\pgfpathlineto{\pgfqpoint{3.097628in}{1.471273in}}%
\pgfpathlineto{\pgfqpoint{3.109226in}{1.471916in}}%
\pgfpathlineto{\pgfqpoint{3.117245in}{1.469315in}}%
\pgfpathlineto{\pgfqpoint{3.153317in}{1.449269in}}%
\pgfpathclose%
\pgfusepath{fill}%
\end{pgfscope}%
\begin{pgfscope}%
\pgfpathrectangle{\pgfqpoint{0.100000in}{0.100000in}}{\pgfqpoint{3.608454in}{2.310000in}}%
\pgfusepath{clip}%
\pgfsetbuttcap%
\pgfsetmiterjoin%
\definecolor{currentfill}{rgb}{0.000000,0.501961,0.749020}%
\pgfsetfillcolor{currentfill}%
\pgfsetlinewidth{0.000000pt}%
\definecolor{currentstroke}{rgb}{0.000000,0.000000,0.000000}%
\pgfsetstrokecolor{currentstroke}%
\pgfsetstrokeopacity{0.000000}%
\pgfsetdash{}{0pt}%
\pgfpathmoveto{\pgfqpoint{1.965169in}{0.705757in}}%
\pgfpathlineto{\pgfqpoint{1.952059in}{0.698355in}}%
\pgfpathlineto{\pgfqpoint{1.947915in}{0.687499in}}%
\pgfpathlineto{\pgfqpoint{1.926963in}{0.697491in}}%
\pgfpathlineto{\pgfqpoint{1.913119in}{0.701138in}}%
\pgfpathlineto{\pgfqpoint{1.907744in}{0.711220in}}%
\pgfpathlineto{\pgfqpoint{1.871649in}{0.712052in}}%
\pgfpathlineto{\pgfqpoint{1.866023in}{0.717001in}}%
\pgfpathlineto{\pgfqpoint{1.863869in}{0.728404in}}%
\pgfpathlineto{\pgfqpoint{1.884281in}{0.742465in}}%
\pgfpathlineto{\pgfqpoint{1.919219in}{0.761725in}}%
\pgfpathlineto{\pgfqpoint{1.924502in}{0.764640in}}%
\pgfpathlineto{\pgfqpoint{1.942049in}{0.732981in}}%
\pgfpathlineto{\pgfqpoint{1.947012in}{0.726789in}}%
\pgfpathlineto{\pgfqpoint{1.951524in}{0.729506in}}%
\pgfpathlineto{\pgfqpoint{1.965169in}{0.705757in}}%
\pgfpathclose%
\pgfusepath{fill}%
\end{pgfscope}%
\begin{pgfscope}%
\pgfpathrectangle{\pgfqpoint{0.100000in}{0.100000in}}{\pgfqpoint{3.608454in}{2.310000in}}%
\pgfusepath{clip}%
\pgfsetbuttcap%
\pgfsetmiterjoin%
\definecolor{currentfill}{rgb}{0.000000,0.800000,0.600000}%
\pgfsetfillcolor{currentfill}%
\pgfsetlinewidth{0.000000pt}%
\definecolor{currentstroke}{rgb}{0.000000,0.000000,0.000000}%
\pgfsetstrokecolor{currentstroke}%
\pgfsetstrokeopacity{0.000000}%
\pgfsetdash{}{0pt}%
\pgfpathmoveto{\pgfqpoint{1.805378in}{0.356576in}}%
\pgfpathlineto{\pgfqpoint{1.817184in}{0.364409in}}%
\pgfpathlineto{\pgfqpoint{1.821677in}{0.373163in}}%
\pgfpathlineto{\pgfqpoint{1.859982in}{0.371433in}}%
\pgfpathlineto{\pgfqpoint{1.891484in}{0.370110in}}%
\pgfpathlineto{\pgfqpoint{1.891060in}{0.356772in}}%
\pgfpathlineto{\pgfqpoint{1.899944in}{0.355045in}}%
\pgfpathlineto{\pgfqpoint{1.930477in}{0.354309in}}%
\pgfpathlineto{\pgfqpoint{1.932540in}{0.346180in}}%
\pgfpathlineto{\pgfqpoint{1.930157in}{0.343218in}}%
\pgfpathlineto{\pgfqpoint{1.936961in}{0.334615in}}%
\pgfpathlineto{\pgfqpoint{1.936672in}{0.327595in}}%
\pgfpathlineto{\pgfqpoint{1.940091in}{0.321228in}}%
\pgfpathlineto{\pgfqpoint{1.940482in}{0.313585in}}%
\pgfpathlineto{\pgfqpoint{1.950108in}{0.309249in}}%
\pgfpathlineto{\pgfqpoint{1.950498in}{0.302891in}}%
\pgfpathlineto{\pgfqpoint{1.934207in}{0.298800in}}%
\pgfpathlineto{\pgfqpoint{1.929813in}{0.292862in}}%
\pgfpathlineto{\pgfqpoint{1.923685in}{0.296794in}}%
\pgfpathlineto{\pgfqpoint{1.911823in}{0.308379in}}%
\pgfpathlineto{\pgfqpoint{1.896783in}{0.312095in}}%
\pgfpathlineto{\pgfqpoint{1.883885in}{0.310848in}}%
\pgfpathlineto{\pgfqpoint{1.870440in}{0.313575in}}%
\pgfpathlineto{\pgfqpoint{1.855483in}{0.326486in}}%
\pgfpathlineto{\pgfqpoint{1.840888in}{0.327815in}}%
\pgfpathlineto{\pgfqpoint{1.830852in}{0.339080in}}%
\pgfpathlineto{\pgfqpoint{1.828260in}{0.338364in}}%
\pgfpathlineto{\pgfqpoint{1.809077in}{0.344635in}}%
\pgfpathlineto{\pgfqpoint{1.810652in}{0.348669in}}%
\pgfpathlineto{\pgfqpoint{1.805378in}{0.356576in}}%
\pgfpathclose%
\pgfusepath{fill}%
\end{pgfscope}%
\begin{pgfscope}%
\pgfpathrectangle{\pgfqpoint{0.100000in}{0.100000in}}{\pgfqpoint{3.608454in}{2.310000in}}%
\pgfusepath{clip}%
\pgfsetbuttcap%
\pgfsetmiterjoin%
\definecolor{currentfill}{rgb}{0.000000,0.525490,0.737255}%
\pgfsetfillcolor{currentfill}%
\pgfsetlinewidth{0.000000pt}%
\definecolor{currentstroke}{rgb}{0.000000,0.000000,0.000000}%
\pgfsetstrokecolor{currentstroke}%
\pgfsetstrokeopacity{0.000000}%
\pgfsetdash{}{0pt}%
\pgfpathmoveto{\pgfqpoint{2.495702in}{1.322010in}}%
\pgfpathlineto{\pgfqpoint{2.494687in}{1.339376in}}%
\pgfpathlineto{\pgfqpoint{2.467492in}{1.337391in}}%
\pgfpathlineto{\pgfqpoint{2.467926in}{1.330458in}}%
\pgfpathlineto{\pgfqpoint{2.460771in}{1.330448in}}%
\pgfpathlineto{\pgfqpoint{2.458623in}{1.367966in}}%
\pgfpathlineto{\pgfqpoint{2.465371in}{1.368430in}}%
\pgfpathlineto{\pgfqpoint{2.506265in}{1.370984in}}%
\pgfpathlineto{\pgfqpoint{2.506514in}{1.367545in}}%
\pgfpathlineto{\pgfqpoint{2.550764in}{1.371543in}}%
\pgfpathlineto{\pgfqpoint{2.558058in}{1.370400in}}%
\pgfpathlineto{\pgfqpoint{2.562133in}{1.362791in}}%
\pgfpathlineto{\pgfqpoint{2.562407in}{1.357174in}}%
\pgfpathlineto{\pgfqpoint{2.566176in}{1.350839in}}%
\pgfpathlineto{\pgfqpoint{2.565996in}{1.346868in}}%
\pgfpathlineto{\pgfqpoint{2.542895in}{1.344839in}}%
\pgfpathlineto{\pgfqpoint{2.521388in}{1.342978in}}%
\pgfpathlineto{\pgfqpoint{2.521612in}{1.327886in}}%
\pgfpathlineto{\pgfqpoint{2.523132in}{1.323450in}}%
\pgfpathlineto{\pgfqpoint{2.495702in}{1.322010in}}%
\pgfpathclose%
\pgfusepath{fill}%
\end{pgfscope}%
\begin{pgfscope}%
\pgfpathrectangle{\pgfqpoint{0.100000in}{0.100000in}}{\pgfqpoint{3.608454in}{2.310000in}}%
\pgfusepath{clip}%
\pgfsetbuttcap%
\pgfsetmiterjoin%
\definecolor{currentfill}{rgb}{0.000000,0.568627,0.715686}%
\pgfsetfillcolor{currentfill}%
\pgfsetlinewidth{0.000000pt}%
\definecolor{currentstroke}{rgb}{0.000000,0.000000,0.000000}%
\pgfsetstrokecolor{currentstroke}%
\pgfsetstrokeopacity{0.000000}%
\pgfsetdash{}{0pt}%
\pgfpathmoveto{\pgfqpoint{3.037191in}{1.156638in}}%
\pgfpathlineto{\pgfqpoint{3.024951in}{1.162919in}}%
\pgfpathlineto{\pgfqpoint{3.020036in}{1.162708in}}%
\pgfpathlineto{\pgfqpoint{3.018531in}{1.178366in}}%
\pgfpathlineto{\pgfqpoint{3.006498in}{1.176840in}}%
\pgfpathlineto{\pgfqpoint{3.004020in}{1.198138in}}%
\pgfpathlineto{\pgfqpoint{2.996900in}{1.203053in}}%
\pgfpathlineto{\pgfqpoint{2.995823in}{1.206373in}}%
\pgfpathlineto{\pgfqpoint{2.999075in}{1.216208in}}%
\pgfpathlineto{\pgfqpoint{3.017299in}{1.218774in}}%
\pgfpathlineto{\pgfqpoint{3.037358in}{1.220901in}}%
\pgfpathlineto{\pgfqpoint{3.054182in}{1.223554in}}%
\pgfpathlineto{\pgfqpoint{3.060084in}{1.181854in}}%
\pgfpathlineto{\pgfqpoint{3.049089in}{1.181346in}}%
\pgfpathlineto{\pgfqpoint{3.039491in}{1.174669in}}%
\pgfpathlineto{\pgfqpoint{3.039762in}{1.164964in}}%
\pgfpathlineto{\pgfqpoint{3.035098in}{1.163147in}}%
\pgfpathlineto{\pgfqpoint{3.037191in}{1.156638in}}%
\pgfpathclose%
\pgfusepath{fill}%
\end{pgfscope}%
\begin{pgfscope}%
\pgfpathrectangle{\pgfqpoint{0.100000in}{0.100000in}}{\pgfqpoint{3.608454in}{2.310000in}}%
\pgfusepath{clip}%
\pgfsetbuttcap%
\pgfsetmiterjoin%
\definecolor{currentfill}{rgb}{0.000000,0.607843,0.696078}%
\pgfsetfillcolor{currentfill}%
\pgfsetlinewidth{0.000000pt}%
\definecolor{currentstroke}{rgb}{0.000000,0.000000,0.000000}%
\pgfsetstrokecolor{currentstroke}%
\pgfsetstrokeopacity{0.000000}%
\pgfsetdash{}{0pt}%
\pgfpathmoveto{\pgfqpoint{3.358684in}{1.520578in}}%
\pgfpathlineto{\pgfqpoint{3.353587in}{1.521095in}}%
\pgfpathlineto{\pgfqpoint{3.348171in}{1.542140in}}%
\pgfpathlineto{\pgfqpoint{3.330022in}{1.567556in}}%
\pgfpathlineto{\pgfqpoint{3.327731in}{1.570641in}}%
\pgfpathlineto{\pgfqpoint{3.339300in}{1.581420in}}%
\pgfpathlineto{\pgfqpoint{3.345889in}{1.596696in}}%
\pgfpathlineto{\pgfqpoint{3.350906in}{1.598176in}}%
\pgfpathlineto{\pgfqpoint{3.359788in}{1.596585in}}%
\pgfpathlineto{\pgfqpoint{3.363035in}{1.588297in}}%
\pgfpathlineto{\pgfqpoint{3.363118in}{1.574131in}}%
\pgfpathlineto{\pgfqpoint{3.365563in}{1.546396in}}%
\pgfpathlineto{\pgfqpoint{3.364102in}{1.538156in}}%
\pgfpathlineto{\pgfqpoint{3.358684in}{1.520578in}}%
\pgfpathclose%
\pgfusepath{fill}%
\end{pgfscope}%
\begin{pgfscope}%
\pgfpathrectangle{\pgfqpoint{0.100000in}{0.100000in}}{\pgfqpoint{3.608454in}{2.310000in}}%
\pgfusepath{clip}%
\pgfsetbuttcap%
\pgfsetmiterjoin%
\definecolor{currentfill}{rgb}{0.000000,0.419608,0.790196}%
\pgfsetfillcolor{currentfill}%
\pgfsetlinewidth{0.000000pt}%
\definecolor{currentstroke}{rgb}{0.000000,0.000000,0.000000}%
\pgfsetstrokecolor{currentstroke}%
\pgfsetstrokeopacity{0.000000}%
\pgfsetdash{}{0pt}%
\pgfpathmoveto{\pgfqpoint{1.939662in}{1.829664in}}%
\pgfpathlineto{\pgfqpoint{1.938531in}{1.788233in}}%
\pgfpathlineto{\pgfqpoint{1.897415in}{1.789939in}}%
\pgfpathlineto{\pgfqpoint{1.898285in}{1.810809in}}%
\pgfpathlineto{\pgfqpoint{1.849510in}{1.813058in}}%
\pgfpathlineto{\pgfqpoint{1.850854in}{1.840684in}}%
\pgfpathlineto{\pgfqpoint{1.843275in}{1.841097in}}%
\pgfpathlineto{\pgfqpoint{1.845659in}{1.896366in}}%
\pgfpathlineto{\pgfqpoint{1.885060in}{1.894397in}}%
\pgfpathlineto{\pgfqpoint{1.941869in}{1.891889in}}%
\pgfpathlineto{\pgfqpoint{1.939662in}{1.829664in}}%
\pgfpathclose%
\pgfusepath{fill}%
\end{pgfscope}%
\begin{pgfscope}%
\pgfpathrectangle{\pgfqpoint{0.100000in}{0.100000in}}{\pgfqpoint{3.608454in}{2.310000in}}%
\pgfusepath{clip}%
\pgfsetbuttcap%
\pgfsetmiterjoin%
\definecolor{currentfill}{rgb}{0.000000,0.600000,0.700000}%
\pgfsetfillcolor{currentfill}%
\pgfsetlinewidth{0.000000pt}%
\definecolor{currentstroke}{rgb}{0.000000,0.000000,0.000000}%
\pgfsetstrokecolor{currentstroke}%
\pgfsetstrokeopacity{0.000000}%
\pgfsetdash{}{0pt}%
\pgfpathmoveto{\pgfqpoint{2.173280in}{0.783716in}}%
\pgfpathlineto{\pgfqpoint{2.139795in}{0.782977in}}%
\pgfpathlineto{\pgfqpoint{2.143853in}{0.772670in}}%
\pgfpathlineto{\pgfqpoint{2.111050in}{0.772661in}}%
\pgfpathlineto{\pgfqpoint{2.107805in}{0.772710in}}%
\pgfpathlineto{\pgfqpoint{2.107625in}{0.827720in}}%
\pgfpathlineto{\pgfqpoint{2.096510in}{0.830344in}}%
\pgfpathlineto{\pgfqpoint{2.096641in}{0.856769in}}%
\pgfpathlineto{\pgfqpoint{2.125475in}{0.856951in}}%
\pgfpathlineto{\pgfqpoint{2.126427in}{0.854941in}}%
\pgfpathlineto{\pgfqpoint{2.170684in}{0.855388in}}%
\pgfpathlineto{\pgfqpoint{2.171444in}{0.784825in}}%
\pgfpathlineto{\pgfqpoint{2.173280in}{0.783716in}}%
\pgfpathclose%
\pgfusepath{fill}%
\end{pgfscope}%
\begin{pgfscope}%
\pgfpathrectangle{\pgfqpoint{0.100000in}{0.100000in}}{\pgfqpoint{3.608454in}{2.310000in}}%
\pgfusepath{clip}%
\pgfsetbuttcap%
\pgfsetmiterjoin%
\definecolor{currentfill}{rgb}{0.000000,0.827451,0.586275}%
\pgfsetfillcolor{currentfill}%
\pgfsetlinewidth{0.000000pt}%
\definecolor{currentstroke}{rgb}{0.000000,0.000000,0.000000}%
\pgfsetstrokecolor{currentstroke}%
\pgfsetstrokeopacity{0.000000}%
\pgfsetdash{}{0pt}%
\pgfpathmoveto{\pgfqpoint{0.976297in}{0.641465in}}%
\pgfpathlineto{\pgfqpoint{0.971445in}{0.636313in}}%
\pgfpathlineto{\pgfqpoint{0.973534in}{0.634639in}}%
\pgfpathlineto{\pgfqpoint{0.971825in}{0.632484in}}%
\pgfpathlineto{\pgfqpoint{0.974033in}{0.630789in}}%
\pgfpathlineto{\pgfqpoint{0.972344in}{0.628719in}}%
\pgfpathlineto{\pgfqpoint{0.974510in}{0.627028in}}%
\pgfpathlineto{\pgfqpoint{0.972240in}{0.625296in}}%
\pgfpathlineto{\pgfqpoint{0.967277in}{0.618907in}}%
\pgfpathlineto{\pgfqpoint{0.971558in}{0.615632in}}%
\pgfpathlineto{\pgfqpoint{0.969910in}{0.613404in}}%
\pgfpathlineto{\pgfqpoint{0.971435in}{0.612298in}}%
\pgfpathlineto{\pgfqpoint{0.964303in}{0.603251in}}%
\pgfpathlineto{\pgfqpoint{0.961995in}{0.601456in}}%
\pgfpathlineto{\pgfqpoint{0.955683in}{0.606572in}}%
\pgfpathlineto{\pgfqpoint{0.952323in}{0.602325in}}%
\pgfpathlineto{\pgfqpoint{0.946326in}{0.607097in}}%
\pgfpathlineto{\pgfqpoint{0.940885in}{0.601118in}}%
\pgfpathlineto{\pgfqpoint{0.936932in}{0.604209in}}%
\pgfpathlineto{\pgfqpoint{0.938627in}{0.606373in}}%
\pgfpathlineto{\pgfqpoint{0.936545in}{0.608078in}}%
\pgfpathlineto{\pgfqpoint{0.938173in}{0.610139in}}%
\pgfpathlineto{\pgfqpoint{0.934029in}{0.613492in}}%
\pgfpathlineto{\pgfqpoint{0.930663in}{0.609300in}}%
\pgfpathlineto{\pgfqpoint{0.926478in}{0.612658in}}%
\pgfpathlineto{\pgfqpoint{0.929884in}{0.616921in}}%
\pgfpathlineto{\pgfqpoint{0.921542in}{0.623737in}}%
\pgfpathlineto{\pgfqpoint{0.914790in}{0.615304in}}%
\pgfpathlineto{\pgfqpoint{0.904546in}{0.623960in}}%
\pgfpathlineto{\pgfqpoint{0.901313in}{0.620035in}}%
\pgfpathlineto{\pgfqpoint{0.895296in}{0.625196in}}%
\pgfpathlineto{\pgfqpoint{0.886576in}{0.615036in}}%
\pgfpathlineto{\pgfqpoint{0.884580in}{0.616759in}}%
\pgfpathlineto{\pgfqpoint{0.882807in}{0.614773in}}%
\pgfpathlineto{\pgfqpoint{0.870840in}{0.625356in}}%
\pgfpathlineto{\pgfqpoint{0.853614in}{0.640940in}}%
\pgfpathlineto{\pgfqpoint{0.845457in}{0.648542in}}%
\pgfpathlineto{\pgfqpoint{0.849066in}{0.652382in}}%
\pgfpathlineto{\pgfqpoint{0.847118in}{0.654238in}}%
\pgfpathlineto{\pgfqpoint{0.850812in}{0.658129in}}%
\pgfpathlineto{\pgfqpoint{0.848280in}{0.660517in}}%
\pgfpathlineto{\pgfqpoint{0.855680in}{0.668231in}}%
\pgfpathlineto{\pgfqpoint{0.855048in}{0.668896in}}%
\pgfpathlineto{\pgfqpoint{0.862419in}{0.676643in}}%
\pgfpathlineto{\pgfqpoint{0.863943in}{0.679699in}}%
\pgfpathlineto{\pgfqpoint{0.867167in}{0.677080in}}%
\pgfpathlineto{\pgfqpoint{0.868881in}{0.673318in}}%
\pgfpathlineto{\pgfqpoint{0.866790in}{0.671377in}}%
\pgfpathlineto{\pgfqpoint{0.864000in}{0.670092in}}%
\pgfpathlineto{\pgfqpoint{0.862316in}{0.667806in}}%
\pgfpathlineto{\pgfqpoint{0.859533in}{0.667863in}}%
\pgfpathlineto{\pgfqpoint{0.858589in}{0.662654in}}%
\pgfpathlineto{\pgfqpoint{0.861141in}{0.660797in}}%
\pgfpathlineto{\pgfqpoint{0.865258in}{0.656612in}}%
\pgfpathlineto{\pgfqpoint{0.868725in}{0.654793in}}%
\pgfpathlineto{\pgfqpoint{0.868360in}{0.652816in}}%
\pgfpathlineto{\pgfqpoint{0.870641in}{0.649968in}}%
\pgfpathlineto{\pgfqpoint{0.873832in}{0.649148in}}%
\pgfpathlineto{\pgfqpoint{0.875177in}{0.647013in}}%
\pgfpathlineto{\pgfqpoint{0.873862in}{0.644273in}}%
\pgfpathlineto{\pgfqpoint{0.879648in}{0.646559in}}%
\pgfpathlineto{\pgfqpoint{0.883020in}{0.648596in}}%
\pgfpathlineto{\pgfqpoint{0.884548in}{0.645448in}}%
\pgfpathlineto{\pgfqpoint{0.888156in}{0.644278in}}%
\pgfpathlineto{\pgfqpoint{0.888264in}{0.647509in}}%
\pgfpathlineto{\pgfqpoint{0.885826in}{0.651944in}}%
\pgfpathlineto{\pgfqpoint{0.882189in}{0.653954in}}%
\pgfpathlineto{\pgfqpoint{0.885088in}{0.656299in}}%
\pgfpathlineto{\pgfqpoint{0.885204in}{0.660563in}}%
\pgfpathlineto{\pgfqpoint{0.886202in}{0.663888in}}%
\pgfpathlineto{\pgfqpoint{0.884247in}{0.666557in}}%
\pgfpathlineto{\pgfqpoint{0.886027in}{0.671213in}}%
\pgfpathlineto{\pgfqpoint{0.888440in}{0.671601in}}%
\pgfpathlineto{\pgfqpoint{0.890129in}{0.670612in}}%
\pgfpathlineto{\pgfqpoint{0.889298in}{0.664217in}}%
\pgfpathlineto{\pgfqpoint{0.886556in}{0.660681in}}%
\pgfpathlineto{\pgfqpoint{0.886858in}{0.653674in}}%
\pgfpathlineto{\pgfqpoint{0.891233in}{0.651664in}}%
\pgfpathlineto{\pgfqpoint{0.892089in}{0.649695in}}%
\pgfpathlineto{\pgfqpoint{0.891874in}{0.645714in}}%
\pgfpathlineto{\pgfqpoint{0.897602in}{0.640349in}}%
\pgfpathlineto{\pgfqpoint{0.900250in}{0.642484in}}%
\pgfpathlineto{\pgfqpoint{0.901360in}{0.646766in}}%
\pgfpathlineto{\pgfqpoint{0.898368in}{0.648567in}}%
\pgfpathlineto{\pgfqpoint{0.897419in}{0.652291in}}%
\pgfpathlineto{\pgfqpoint{0.893693in}{0.653344in}}%
\pgfpathlineto{\pgfqpoint{0.891532in}{0.652766in}}%
\pgfpathlineto{\pgfqpoint{0.889610in}{0.655212in}}%
\pgfpathlineto{\pgfqpoint{0.889414in}{0.661725in}}%
\pgfpathlineto{\pgfqpoint{0.891701in}{0.663459in}}%
\pgfpathlineto{\pgfqpoint{0.897189in}{0.666384in}}%
\pgfpathlineto{\pgfqpoint{0.896490in}{0.669062in}}%
\pgfpathlineto{\pgfqpoint{0.890332in}{0.672207in}}%
\pgfpathlineto{\pgfqpoint{0.891668in}{0.673920in}}%
\pgfpathlineto{\pgfqpoint{0.887467in}{0.677256in}}%
\pgfpathlineto{\pgfqpoint{0.884828in}{0.680240in}}%
\pgfpathlineto{\pgfqpoint{0.881627in}{0.686225in}}%
\pgfpathlineto{\pgfqpoint{0.885981in}{0.691916in}}%
\pgfpathlineto{\pgfqpoint{0.888139in}{0.698608in}}%
\pgfpathlineto{\pgfqpoint{0.888448in}{0.701521in}}%
\pgfpathlineto{\pgfqpoint{0.887314in}{0.706916in}}%
\pgfpathlineto{\pgfqpoint{0.887408in}{0.709755in}}%
\pgfpathlineto{\pgfqpoint{0.886741in}{0.720168in}}%
\pgfpathlineto{\pgfqpoint{0.894146in}{0.712972in}}%
\pgfpathlineto{\pgfqpoint{0.898458in}{0.717465in}}%
\pgfpathlineto{\pgfqpoint{0.913662in}{0.703154in}}%
\pgfpathlineto{\pgfqpoint{0.915131in}{0.704745in}}%
\pgfpathlineto{\pgfqpoint{0.921351in}{0.699050in}}%
\pgfpathlineto{\pgfqpoint{0.919895in}{0.697448in}}%
\pgfpathlineto{\pgfqpoint{0.924102in}{0.693649in}}%
\pgfpathlineto{\pgfqpoint{0.925549in}{0.695258in}}%
\pgfpathlineto{\pgfqpoint{0.929228in}{0.691967in}}%
\pgfpathlineto{\pgfqpoint{0.927789in}{0.690351in}}%
\pgfpathlineto{\pgfqpoint{0.936363in}{0.682792in}}%
\pgfpathlineto{\pgfqpoint{0.937775in}{0.684411in}}%
\pgfpathlineto{\pgfqpoint{0.941559in}{0.681130in}}%
\pgfpathlineto{\pgfqpoint{0.940158in}{0.679506in}}%
\pgfpathlineto{\pgfqpoint{0.954477in}{0.667338in}}%
\pgfpathlineto{\pgfqpoint{0.963280in}{0.660093in}}%
\pgfpathlineto{\pgfqpoint{0.961484in}{0.657889in}}%
\pgfpathlineto{\pgfqpoint{0.969832in}{0.651052in}}%
\pgfpathlineto{\pgfqpoint{0.968183in}{0.648986in}}%
\pgfpathlineto{\pgfqpoint{0.976829in}{0.642165in}}%
\pgfpathlineto{\pgfqpoint{0.976297in}{0.641465in}}%
\pgfpathclose%
\pgfusepath{fill}%
\end{pgfscope}%
\begin{pgfscope}%
\pgfpathrectangle{\pgfqpoint{0.100000in}{0.100000in}}{\pgfqpoint{3.608454in}{2.310000in}}%
\pgfusepath{clip}%
\pgfsetbuttcap%
\pgfsetmiterjoin%
\definecolor{currentfill}{rgb}{0.000000,0.796078,0.601961}%
\pgfsetfillcolor{currentfill}%
\pgfsetlinewidth{0.000000pt}%
\definecolor{currentstroke}{rgb}{0.000000,0.000000,0.000000}%
\pgfsetstrokecolor{currentstroke}%
\pgfsetstrokeopacity{0.000000}%
\pgfsetdash{}{0pt}%
\pgfpathmoveto{\pgfqpoint{1.801456in}{1.757724in}}%
\pgfpathlineto{\pgfqpoint{1.800009in}{1.732253in}}%
\pgfpathlineto{\pgfqpoint{1.800957in}{1.721865in}}%
\pgfpathlineto{\pgfqpoint{1.793255in}{1.722620in}}%
\pgfpathlineto{\pgfqpoint{1.787775in}{1.725501in}}%
\pgfpathlineto{\pgfqpoint{1.775409in}{1.724494in}}%
\pgfpathlineto{\pgfqpoint{1.768321in}{1.732235in}}%
\pgfpathlineto{\pgfqpoint{1.760148in}{1.734344in}}%
\pgfpathlineto{\pgfqpoint{1.762841in}{1.760148in}}%
\pgfpathlineto{\pgfqpoint{1.801456in}{1.757724in}}%
\pgfpathclose%
\pgfusepath{fill}%
\end{pgfscope}%
\begin{pgfscope}%
\pgfpathrectangle{\pgfqpoint{0.100000in}{0.100000in}}{\pgfqpoint{3.608454in}{2.310000in}}%
\pgfusepath{clip}%
\pgfsetbuttcap%
\pgfsetmiterjoin%
\definecolor{currentfill}{rgb}{0.000000,0.462745,0.768627}%
\pgfsetfillcolor{currentfill}%
\pgfsetlinewidth{0.000000pt}%
\definecolor{currentstroke}{rgb}{0.000000,0.000000,0.000000}%
\pgfsetstrokecolor{currentstroke}%
\pgfsetstrokeopacity{0.000000}%
\pgfsetdash{}{0pt}%
\pgfpathmoveto{\pgfqpoint{2.387375in}{1.567036in}}%
\pgfpathlineto{\pgfqpoint{2.380372in}{1.562314in}}%
\pgfpathlineto{\pgfqpoint{2.376029in}{1.564684in}}%
\pgfpathlineto{\pgfqpoint{2.348758in}{1.564236in}}%
\pgfpathlineto{\pgfqpoint{2.348123in}{1.578046in}}%
\pgfpathlineto{\pgfqpoint{2.346907in}{1.605656in}}%
\pgfpathlineto{\pgfqpoint{2.360507in}{1.606246in}}%
\pgfpathlineto{\pgfqpoint{2.360207in}{1.613167in}}%
\pgfpathlineto{\pgfqpoint{2.371380in}{1.613652in}}%
\pgfpathlineto{\pgfqpoint{2.374950in}{1.609365in}}%
\pgfpathlineto{\pgfqpoint{2.376710in}{1.601691in}}%
\pgfpathlineto{\pgfqpoint{2.381293in}{1.599262in}}%
\pgfpathlineto{\pgfqpoint{2.390701in}{1.593517in}}%
\pgfpathlineto{\pgfqpoint{2.392472in}{1.584764in}}%
\pgfpathlineto{\pgfqpoint{2.390834in}{1.569129in}}%
\pgfpathlineto{\pgfqpoint{2.387375in}{1.567036in}}%
\pgfpathclose%
\pgfusepath{fill}%
\end{pgfscope}%
\begin{pgfscope}%
\pgfpathrectangle{\pgfqpoint{0.100000in}{0.100000in}}{\pgfqpoint{3.608454in}{2.310000in}}%
\pgfusepath{clip}%
\pgfsetbuttcap%
\pgfsetmiterjoin%
\definecolor{currentfill}{rgb}{0.000000,0.364706,0.817647}%
\pgfsetfillcolor{currentfill}%
\pgfsetlinewidth{0.000000pt}%
\definecolor{currentstroke}{rgb}{0.000000,0.000000,0.000000}%
\pgfsetstrokecolor{currentstroke}%
\pgfsetstrokeopacity{0.000000}%
\pgfsetdash{}{0pt}%
\pgfpathmoveto{\pgfqpoint{1.853837in}{1.424663in}}%
\pgfpathlineto{\pgfqpoint{1.881222in}{1.423490in}}%
\pgfpathlineto{\pgfqpoint{1.894644in}{1.422961in}}%
\pgfpathlineto{\pgfqpoint{1.893270in}{1.388628in}}%
\pgfpathlineto{\pgfqpoint{1.860470in}{1.390012in}}%
\pgfpathlineto{\pgfqpoint{1.824934in}{1.391676in}}%
\pgfpathlineto{\pgfqpoint{1.826801in}{1.425965in}}%
\pgfpathlineto{\pgfqpoint{1.853837in}{1.424663in}}%
\pgfpathclose%
\pgfusepath{fill}%
\end{pgfscope}%
\begin{pgfscope}%
\pgfpathrectangle{\pgfqpoint{0.100000in}{0.100000in}}{\pgfqpoint{3.608454in}{2.310000in}}%
\pgfusepath{clip}%
\pgfsetbuttcap%
\pgfsetmiterjoin%
\definecolor{currentfill}{rgb}{0.000000,0.494118,0.752941}%
\pgfsetfillcolor{currentfill}%
\pgfsetlinewidth{0.000000pt}%
\definecolor{currentstroke}{rgb}{0.000000,0.000000,0.000000}%
\pgfsetstrokecolor{currentstroke}%
\pgfsetstrokeopacity{0.000000}%
\pgfsetdash{}{0pt}%
\pgfpathmoveto{\pgfqpoint{2.374602in}{1.377391in}}%
\pgfpathlineto{\pgfqpoint{2.376576in}{1.388834in}}%
\pgfpathlineto{\pgfqpoint{2.371479in}{1.401180in}}%
\pgfpathlineto{\pgfqpoint{2.375453in}{1.412534in}}%
\pgfpathlineto{\pgfqpoint{2.354524in}{1.412008in}}%
\pgfpathlineto{\pgfqpoint{2.353565in}{1.439495in}}%
\pgfpathlineto{\pgfqpoint{2.353369in}{1.446719in}}%
\pgfpathlineto{\pgfqpoint{2.381085in}{1.447410in}}%
\pgfpathlineto{\pgfqpoint{2.381360in}{1.440493in}}%
\pgfpathlineto{\pgfqpoint{2.396575in}{1.440884in}}%
\pgfpathlineto{\pgfqpoint{2.389932in}{1.435148in}}%
\pgfpathlineto{\pgfqpoint{2.391369in}{1.431061in}}%
\pgfpathlineto{\pgfqpoint{2.397384in}{1.430530in}}%
\pgfpathlineto{\pgfqpoint{2.409257in}{1.435587in}}%
\pgfpathlineto{\pgfqpoint{2.410721in}{1.410092in}}%
\pgfpathlineto{\pgfqpoint{2.415980in}{1.389550in}}%
\pgfpathlineto{\pgfqpoint{2.402134in}{1.388642in}}%
\pgfpathlineto{\pgfqpoint{2.403570in}{1.368279in}}%
\pgfpathlineto{\pgfqpoint{2.393816in}{1.360813in}}%
\pgfpathlineto{\pgfqpoint{2.383166in}{1.360351in}}%
\pgfpathlineto{\pgfqpoint{2.376719in}{1.362091in}}%
\pgfpathlineto{\pgfqpoint{2.374602in}{1.377391in}}%
\pgfpathclose%
\pgfusepath{fill}%
\end{pgfscope}%
\begin{pgfscope}%
\pgfpathrectangle{\pgfqpoint{0.100000in}{0.100000in}}{\pgfqpoint{3.608454in}{2.310000in}}%
\pgfusepath{clip}%
\pgfsetbuttcap%
\pgfsetmiterjoin%
\definecolor{currentfill}{rgb}{0.000000,0.639216,0.680392}%
\pgfsetfillcolor{currentfill}%
\pgfsetlinewidth{0.000000pt}%
\definecolor{currentstroke}{rgb}{0.000000,0.000000,0.000000}%
\pgfsetstrokecolor{currentstroke}%
\pgfsetstrokeopacity{0.000000}%
\pgfsetdash{}{0pt}%
\pgfpathmoveto{\pgfqpoint{3.177265in}{1.821081in}}%
\pgfpathlineto{\pgfqpoint{3.173354in}{1.830395in}}%
\pgfpathlineto{\pgfqpoint{3.177753in}{1.838567in}}%
\pgfpathlineto{\pgfqpoint{3.175101in}{1.843253in}}%
\pgfpathlineto{\pgfqpoint{3.166217in}{1.845402in}}%
\pgfpathlineto{\pgfqpoint{3.162379in}{1.854420in}}%
\pgfpathlineto{\pgfqpoint{3.171115in}{1.863602in}}%
\pgfpathlineto{\pgfqpoint{3.169608in}{1.868162in}}%
\pgfpathlineto{\pgfqpoint{3.181998in}{1.877869in}}%
\pgfpathlineto{\pgfqpoint{3.186045in}{1.883942in}}%
\pgfpathlineto{\pgfqpoint{3.187678in}{1.891082in}}%
\pgfpathlineto{\pgfqpoint{3.198874in}{1.908865in}}%
\pgfpathlineto{\pgfqpoint{3.215748in}{1.929678in}}%
\pgfpathlineto{\pgfqpoint{3.223456in}{1.936959in}}%
\pgfpathlineto{\pgfqpoint{3.230888in}{1.941273in}}%
\pgfpathlineto{\pgfqpoint{3.236756in}{1.940890in}}%
\pgfpathlineto{\pgfqpoint{3.242124in}{1.938669in}}%
\pgfpathlineto{\pgfqpoint{3.262366in}{1.879593in}}%
\pgfpathlineto{\pgfqpoint{3.263105in}{1.873951in}}%
\pgfpathlineto{\pgfqpoint{3.245888in}{1.867665in}}%
\pgfpathlineto{\pgfqpoint{3.234624in}{1.863539in}}%
\pgfpathlineto{\pgfqpoint{3.227847in}{1.865817in}}%
\pgfpathlineto{\pgfqpoint{3.239480in}{1.829258in}}%
\pgfpathlineto{\pgfqpoint{3.219039in}{1.808893in}}%
\pgfpathlineto{\pgfqpoint{3.205664in}{1.810106in}}%
\pgfpathlineto{\pgfqpoint{3.200994in}{1.826823in}}%
\pgfpathlineto{\pgfqpoint{3.186645in}{1.825233in}}%
\pgfpathlineto{\pgfqpoint{3.177265in}{1.821081in}}%
\pgfpathclose%
\pgfusepath{fill}%
\end{pgfscope}%
\begin{pgfscope}%
\pgfpathrectangle{\pgfqpoint{0.100000in}{0.100000in}}{\pgfqpoint{3.608454in}{2.310000in}}%
\pgfusepath{clip}%
\pgfsetbuttcap%
\pgfsetmiterjoin%
\definecolor{currentfill}{rgb}{0.000000,0.647059,0.676471}%
\pgfsetfillcolor{currentfill}%
\pgfsetlinewidth{0.000000pt}%
\definecolor{currentstroke}{rgb}{0.000000,0.000000,0.000000}%
\pgfsetstrokecolor{currentstroke}%
\pgfsetstrokeopacity{0.000000}%
\pgfsetdash{}{0pt}%
\pgfpathmoveto{\pgfqpoint{2.268476in}{1.251773in}}%
\pgfpathlineto{\pgfqpoint{2.260153in}{1.246809in}}%
\pgfpathlineto{\pgfqpoint{2.255244in}{1.254033in}}%
\pgfpathlineto{\pgfqpoint{2.227021in}{1.254178in}}%
\pgfpathlineto{\pgfqpoint{2.227204in}{1.266822in}}%
\pgfpathlineto{\pgfqpoint{2.225942in}{1.276603in}}%
\pgfpathlineto{\pgfqpoint{2.226458in}{1.316665in}}%
\pgfpathlineto{\pgfqpoint{2.240062in}{1.316251in}}%
\pgfpathlineto{\pgfqpoint{2.240098in}{1.309360in}}%
\pgfpathlineto{\pgfqpoint{2.253822in}{1.296576in}}%
\pgfpathlineto{\pgfqpoint{2.267557in}{1.296520in}}%
\pgfpathlineto{\pgfqpoint{2.267441in}{1.289875in}}%
\pgfpathlineto{\pgfqpoint{2.274008in}{1.289693in}}%
\pgfpathlineto{\pgfqpoint{2.280744in}{1.286387in}}%
\pgfpathlineto{\pgfqpoint{2.282058in}{1.264607in}}%
\pgfpathlineto{\pgfqpoint{2.268296in}{1.264580in}}%
\pgfpathlineto{\pgfqpoint{2.268476in}{1.251773in}}%
\pgfpathclose%
\pgfusepath{fill}%
\end{pgfscope}%
\begin{pgfscope}%
\pgfpathrectangle{\pgfqpoint{0.100000in}{0.100000in}}{\pgfqpoint{3.608454in}{2.310000in}}%
\pgfusepath{clip}%
\pgfsetbuttcap%
\pgfsetmiterjoin%
\definecolor{currentfill}{rgb}{0.000000,0.709804,0.645098}%
\pgfsetfillcolor{currentfill}%
\pgfsetlinewidth{0.000000pt}%
\definecolor{currentstroke}{rgb}{0.000000,0.000000,0.000000}%
\pgfsetstrokecolor{currentstroke}%
\pgfsetstrokeopacity{0.000000}%
\pgfsetdash{}{0pt}%
\pgfpathmoveto{\pgfqpoint{2.434616in}{1.075007in}}%
\pgfpathlineto{\pgfqpoint{2.426249in}{1.075405in}}%
\pgfpathlineto{\pgfqpoint{2.427918in}{1.068838in}}%
\pgfpathlineto{\pgfqpoint{2.420217in}{1.066185in}}%
\pgfpathlineto{\pgfqpoint{2.410748in}{1.065840in}}%
\pgfpathlineto{\pgfqpoint{2.408349in}{1.109943in}}%
\pgfpathlineto{\pgfqpoint{2.406026in}{1.117238in}}%
\pgfpathlineto{\pgfqpoint{2.411938in}{1.125058in}}%
\pgfpathlineto{\pgfqpoint{2.420239in}{1.132410in}}%
\pgfpathlineto{\pgfqpoint{2.421101in}{1.141357in}}%
\pgfpathlineto{\pgfqpoint{2.410619in}{1.149798in}}%
\pgfpathlineto{\pgfqpoint{2.414649in}{1.160533in}}%
\pgfpathlineto{\pgfqpoint{2.426666in}{1.161074in}}%
\pgfpathlineto{\pgfqpoint{2.426853in}{1.142086in}}%
\pgfpathlineto{\pgfqpoint{2.440887in}{1.142638in}}%
\pgfpathlineto{\pgfqpoint{2.450987in}{1.144593in}}%
\pgfpathlineto{\pgfqpoint{2.454276in}{1.139497in}}%
\pgfpathlineto{\pgfqpoint{2.455346in}{1.132898in}}%
\pgfpathlineto{\pgfqpoint{2.445412in}{1.130104in}}%
\pgfpathlineto{\pgfqpoint{2.452158in}{1.122751in}}%
\pgfpathlineto{\pgfqpoint{2.446875in}{1.118895in}}%
\pgfpathlineto{\pgfqpoint{2.445220in}{1.109326in}}%
\pgfpathlineto{\pgfqpoint{2.442276in}{1.103468in}}%
\pgfpathlineto{\pgfqpoint{2.442664in}{1.097005in}}%
\pgfpathlineto{\pgfqpoint{2.439328in}{1.092627in}}%
\pgfpathlineto{\pgfqpoint{2.431130in}{1.090720in}}%
\pgfpathlineto{\pgfqpoint{2.432045in}{1.079916in}}%
\pgfpathlineto{\pgfqpoint{2.434616in}{1.075007in}}%
\pgfpathclose%
\pgfusepath{fill}%
\end{pgfscope}%
\begin{pgfscope}%
\pgfpathrectangle{\pgfqpoint{0.100000in}{0.100000in}}{\pgfqpoint{3.608454in}{2.310000in}}%
\pgfusepath{clip}%
\pgfsetbuttcap%
\pgfsetmiterjoin%
\definecolor{currentfill}{rgb}{0.000000,0.552941,0.723529}%
\pgfsetfillcolor{currentfill}%
\pgfsetlinewidth{0.000000pt}%
\definecolor{currentstroke}{rgb}{0.000000,0.000000,0.000000}%
\pgfsetstrokecolor{currentstroke}%
\pgfsetstrokeopacity{0.000000}%
\pgfsetdash{}{0pt}%
\pgfpathmoveto{\pgfqpoint{2.994468in}{0.861760in}}%
\pgfpathlineto{\pgfqpoint{2.991213in}{0.868134in}}%
\pgfpathlineto{\pgfqpoint{2.962520in}{0.870538in}}%
\pgfpathlineto{\pgfqpoint{2.954746in}{0.867311in}}%
\pgfpathlineto{\pgfqpoint{2.950142in}{0.872561in}}%
\pgfpathlineto{\pgfqpoint{2.951258in}{0.878733in}}%
\pgfpathlineto{\pgfqpoint{2.943431in}{0.882150in}}%
\pgfpathlineto{\pgfqpoint{2.932829in}{0.881990in}}%
\pgfpathlineto{\pgfqpoint{2.935640in}{0.887719in}}%
\pgfpathlineto{\pgfqpoint{2.943276in}{0.894232in}}%
\pgfpathlineto{\pgfqpoint{2.940759in}{0.897010in}}%
\pgfpathlineto{\pgfqpoint{2.944266in}{0.903605in}}%
\pgfpathlineto{\pgfqpoint{2.950977in}{0.910809in}}%
\pgfpathlineto{\pgfqpoint{2.959107in}{0.909040in}}%
\pgfpathlineto{\pgfqpoint{2.962570in}{0.910798in}}%
\pgfpathlineto{\pgfqpoint{2.965871in}{0.919067in}}%
\pgfpathlineto{\pgfqpoint{2.980224in}{0.924118in}}%
\pgfpathlineto{\pgfqpoint{2.986683in}{0.921479in}}%
\pgfpathlineto{\pgfqpoint{2.999992in}{0.934399in}}%
\pgfpathlineto{\pgfqpoint{3.003733in}{0.931980in}}%
\pgfpathlineto{\pgfqpoint{3.003885in}{0.926116in}}%
\pgfpathlineto{\pgfqpoint{3.010014in}{0.919433in}}%
\pgfpathlineto{\pgfqpoint{3.011589in}{0.908504in}}%
\pgfpathlineto{\pgfqpoint{3.015378in}{0.900512in}}%
\pgfpathlineto{\pgfqpoint{3.006499in}{0.892231in}}%
\pgfpathlineto{\pgfqpoint{3.008209in}{0.883079in}}%
\pgfpathlineto{\pgfqpoint{3.014555in}{0.877148in}}%
\pgfpathlineto{\pgfqpoint{3.016488in}{0.872188in}}%
\pgfpathlineto{\pgfqpoint{2.994468in}{0.861760in}}%
\pgfpathclose%
\pgfusepath{fill}%
\end{pgfscope}%
\begin{pgfscope}%
\pgfpathrectangle{\pgfqpoint{0.100000in}{0.100000in}}{\pgfqpoint{3.608454in}{2.310000in}}%
\pgfusepath{clip}%
\pgfsetbuttcap%
\pgfsetmiterjoin%
\definecolor{currentfill}{rgb}{0.000000,0.631373,0.684314}%
\pgfsetfillcolor{currentfill}%
\pgfsetlinewidth{0.000000pt}%
\definecolor{currentstroke}{rgb}{0.000000,0.000000,0.000000}%
\pgfsetstrokecolor{currentstroke}%
\pgfsetstrokeopacity{0.000000}%
\pgfsetdash{}{0pt}%
\pgfpathmoveto{\pgfqpoint{1.438374in}{0.821513in}}%
\pgfpathlineto{\pgfqpoint{1.435480in}{0.791748in}}%
\pgfpathlineto{\pgfqpoint{1.429579in}{0.734241in}}%
\pgfpathlineto{\pgfqpoint{1.426837in}{0.715528in}}%
\pgfpathlineto{\pgfqpoint{1.422243in}{0.713258in}}%
\pgfpathlineto{\pgfqpoint{1.413714in}{0.723890in}}%
\pgfpathlineto{\pgfqpoint{1.407758in}{0.729084in}}%
\pgfpathlineto{\pgfqpoint{1.396001in}{0.734660in}}%
\pgfpathlineto{\pgfqpoint{1.395928in}{0.737162in}}%
\pgfpathlineto{\pgfqpoint{1.386189in}{0.746307in}}%
\pgfpathlineto{\pgfqpoint{1.383892in}{0.754060in}}%
\pgfpathlineto{\pgfqpoint{1.373173in}{0.762125in}}%
\pgfpathlineto{\pgfqpoint{1.362767in}{0.779271in}}%
\pgfpathlineto{\pgfqpoint{1.354361in}{0.783205in}}%
\pgfpathlineto{\pgfqpoint{1.346623in}{0.789592in}}%
\pgfpathlineto{\pgfqpoint{1.337607in}{0.811852in}}%
\pgfpathlineto{\pgfqpoint{1.328213in}{0.817202in}}%
\pgfpathlineto{\pgfqpoint{1.276635in}{0.824008in}}%
\pgfpathlineto{\pgfqpoint{1.285282in}{0.888829in}}%
\pgfpathlineto{\pgfqpoint{1.287163in}{0.902597in}}%
\pgfpathlineto{\pgfqpoint{1.314425in}{0.898871in}}%
\pgfpathlineto{\pgfqpoint{1.314868in}{0.902303in}}%
\pgfpathlineto{\pgfqpoint{1.353414in}{0.915777in}}%
\pgfpathlineto{\pgfqpoint{1.350085in}{0.908598in}}%
\pgfpathlineto{\pgfqpoint{1.340564in}{0.833080in}}%
\pgfpathlineto{\pgfqpoint{1.404085in}{0.825375in}}%
\pgfpathlineto{\pgfqpoint{1.438374in}{0.821513in}}%
\pgfpathclose%
\pgfusepath{fill}%
\end{pgfscope}%
\begin{pgfscope}%
\pgfpathrectangle{\pgfqpoint{0.100000in}{0.100000in}}{\pgfqpoint{3.608454in}{2.310000in}}%
\pgfusepath{clip}%
\pgfsetbuttcap%
\pgfsetmiterjoin%
\definecolor{currentfill}{rgb}{0.000000,0.564706,0.717647}%
\pgfsetfillcolor{currentfill}%
\pgfsetlinewidth{0.000000pt}%
\definecolor{currentstroke}{rgb}{0.000000,0.000000,0.000000}%
\pgfsetstrokecolor{currentstroke}%
\pgfsetstrokeopacity{0.000000}%
\pgfsetdash{}{0pt}%
\pgfpathmoveto{\pgfqpoint{2.765201in}{0.880136in}}%
\pgfpathlineto{\pgfqpoint{2.760875in}{0.889582in}}%
\pgfpathlineto{\pgfqpoint{2.782280in}{0.891779in}}%
\pgfpathlineto{\pgfqpoint{2.779766in}{0.919893in}}%
\pgfpathlineto{\pgfqpoint{2.802911in}{0.922322in}}%
\pgfpathlineto{\pgfqpoint{2.802438in}{0.905479in}}%
\pgfpathlineto{\pgfqpoint{2.805495in}{0.895708in}}%
\pgfpathlineto{\pgfqpoint{2.814267in}{0.888921in}}%
\pgfpathlineto{\pgfqpoint{2.821573in}{0.886623in}}%
\pgfpathlineto{\pgfqpoint{2.812799in}{0.870744in}}%
\pgfpathlineto{\pgfqpoint{2.800195in}{0.867107in}}%
\pgfpathlineto{\pgfqpoint{2.799223in}{0.863800in}}%
\pgfpathlineto{\pgfqpoint{2.801485in}{0.843067in}}%
\pgfpathlineto{\pgfqpoint{2.783833in}{0.840916in}}%
\pgfpathlineto{\pgfqpoint{2.775432in}{0.832235in}}%
\pgfpathlineto{\pgfqpoint{2.776392in}{0.826653in}}%
\pgfpathlineto{\pgfqpoint{2.767748in}{0.825602in}}%
\pgfpathlineto{\pgfqpoint{2.762136in}{0.831926in}}%
\pgfpathlineto{\pgfqpoint{2.751930in}{0.830708in}}%
\pgfpathlineto{\pgfqpoint{2.749632in}{0.837509in}}%
\pgfpathlineto{\pgfqpoint{2.748134in}{0.851370in}}%
\pgfpathlineto{\pgfqpoint{2.755065in}{0.852187in}}%
\pgfpathlineto{\pgfqpoint{2.754330in}{0.856732in}}%
\pgfpathlineto{\pgfqpoint{2.772613in}{0.859054in}}%
\pgfpathlineto{\pgfqpoint{2.776169in}{0.862453in}}%
\pgfpathlineto{\pgfqpoint{2.768066in}{0.872209in}}%
\pgfpathlineto{\pgfqpoint{2.765201in}{0.880136in}}%
\pgfpathclose%
\pgfusepath{fill}%
\end{pgfscope}%
\begin{pgfscope}%
\pgfpathrectangle{\pgfqpoint{0.100000in}{0.100000in}}{\pgfqpoint{3.608454in}{2.310000in}}%
\pgfusepath{clip}%
\pgfsetbuttcap%
\pgfsetmiterjoin%
\definecolor{currentfill}{rgb}{0.000000,0.494118,0.752941}%
\pgfsetfillcolor{currentfill}%
\pgfsetlinewidth{0.000000pt}%
\definecolor{currentstroke}{rgb}{0.000000,0.000000,0.000000}%
\pgfsetstrokecolor{currentstroke}%
\pgfsetstrokeopacity{0.000000}%
\pgfsetdash{}{0pt}%
\pgfpathmoveto{\pgfqpoint{2.682505in}{0.872318in}}%
\pgfpathlineto{\pgfqpoint{2.680466in}{0.868375in}}%
\pgfpathlineto{\pgfqpoint{2.660419in}{0.866374in}}%
\pgfpathlineto{\pgfqpoint{2.660793in}{0.862897in}}%
\pgfpathlineto{\pgfqpoint{2.647203in}{0.861864in}}%
\pgfpathlineto{\pgfqpoint{2.640493in}{0.861198in}}%
\pgfpathlineto{\pgfqpoint{2.639182in}{0.874953in}}%
\pgfpathlineto{\pgfqpoint{2.648704in}{0.875737in}}%
\pgfpathlineto{\pgfqpoint{2.647138in}{0.892648in}}%
\pgfpathlineto{\pgfqpoint{2.640666in}{0.895938in}}%
\pgfpathlineto{\pgfqpoint{2.636464in}{0.900943in}}%
\pgfpathlineto{\pgfqpoint{2.635819in}{0.907320in}}%
\pgfpathlineto{\pgfqpoint{2.625197in}{0.912348in}}%
\pgfpathlineto{\pgfqpoint{2.617467in}{0.921208in}}%
\pgfpathlineto{\pgfqpoint{2.618130in}{0.927031in}}%
\pgfpathlineto{\pgfqpoint{2.623213in}{0.935505in}}%
\pgfpathlineto{\pgfqpoint{2.628827in}{0.937113in}}%
\pgfpathlineto{\pgfqpoint{2.631914in}{0.941959in}}%
\pgfpathlineto{\pgfqpoint{2.632392in}{0.948970in}}%
\pgfpathlineto{\pgfqpoint{2.635731in}{0.956087in}}%
\pgfpathlineto{\pgfqpoint{2.627618in}{0.959226in}}%
\pgfpathlineto{\pgfqpoint{2.625447in}{0.963052in}}%
\pgfpathlineto{\pgfqpoint{2.623076in}{0.990093in}}%
\pgfpathlineto{\pgfqpoint{2.622963in}{0.991238in}}%
\pgfpathlineto{\pgfqpoint{2.657488in}{0.993619in}}%
\pgfpathlineto{\pgfqpoint{2.664313in}{0.994103in}}%
\pgfpathlineto{\pgfqpoint{2.677206in}{0.979100in}}%
\pgfpathlineto{\pgfqpoint{2.675715in}{0.974764in}}%
\pgfpathlineto{\pgfqpoint{2.678957in}{0.970072in}}%
\pgfpathlineto{\pgfqpoint{2.684888in}{0.971060in}}%
\pgfpathlineto{\pgfqpoint{2.689264in}{0.964347in}}%
\pgfpathlineto{\pgfqpoint{2.694783in}{0.960330in}}%
\pgfpathlineto{\pgfqpoint{2.696367in}{0.955422in}}%
\pgfpathlineto{\pgfqpoint{2.690775in}{0.946928in}}%
\pgfpathlineto{\pgfqpoint{2.686513in}{0.939240in}}%
\pgfpathlineto{\pgfqpoint{2.688862in}{0.935943in}}%
\pgfpathlineto{\pgfqpoint{2.677615in}{0.922675in}}%
\pgfpathlineto{\pgfqpoint{2.679726in}{0.915262in}}%
\pgfpathlineto{\pgfqpoint{2.673491in}{0.910352in}}%
\pgfpathlineto{\pgfqpoint{2.670991in}{0.904805in}}%
\pgfpathlineto{\pgfqpoint{2.672768in}{0.884091in}}%
\pgfpathlineto{\pgfqpoint{2.676747in}{0.876531in}}%
\pgfpathlineto{\pgfqpoint{2.682505in}{0.872318in}}%
\pgfpathclose%
\pgfusepath{fill}%
\end{pgfscope}%
\begin{pgfscope}%
\pgfpathrectangle{\pgfqpoint{0.100000in}{0.100000in}}{\pgfqpoint{3.608454in}{2.310000in}}%
\pgfusepath{clip}%
\pgfsetbuttcap%
\pgfsetmiterjoin%
\definecolor{currentfill}{rgb}{0.000000,0.458824,0.770588}%
\pgfsetfillcolor{currentfill}%
\pgfsetlinewidth{0.000000pt}%
\definecolor{currentstroke}{rgb}{0.000000,0.000000,0.000000}%
\pgfsetstrokecolor{currentstroke}%
\pgfsetstrokeopacity{0.000000}%
\pgfsetdash{}{0pt}%
\pgfpathmoveto{\pgfqpoint{2.032756in}{2.134709in}}%
\pgfpathlineto{\pgfqpoint{2.032428in}{2.112015in}}%
\pgfpathlineto{\pgfqpoint{2.033124in}{2.098129in}}%
\pgfpathlineto{\pgfqpoint{1.991978in}{2.098902in}}%
\pgfpathlineto{\pgfqpoint{1.993550in}{2.092414in}}%
\pgfpathlineto{\pgfqpoint{1.991966in}{2.086284in}}%
\pgfpathlineto{\pgfqpoint{1.993953in}{2.077757in}}%
\pgfpathlineto{\pgfqpoint{1.992478in}{2.070813in}}%
\pgfpathlineto{\pgfqpoint{1.951799in}{2.072044in}}%
\pgfpathlineto{\pgfqpoint{1.931043in}{2.072760in}}%
\pgfpathlineto{\pgfqpoint{1.931568in}{2.086739in}}%
\pgfpathlineto{\pgfqpoint{1.930695in}{2.100743in}}%
\pgfpathlineto{\pgfqpoint{1.896095in}{2.102164in}}%
\pgfpathlineto{\pgfqpoint{1.895146in}{2.116164in}}%
\pgfpathlineto{\pgfqpoint{1.896146in}{2.138772in}}%
\pgfpathlineto{\pgfqpoint{1.946219in}{2.136867in}}%
\pgfpathlineto{\pgfqpoint{1.989377in}{2.135633in}}%
\pgfpathlineto{\pgfqpoint{2.032756in}{2.134709in}}%
\pgfpathclose%
\pgfusepath{fill}%
\end{pgfscope}%
\begin{pgfscope}%
\pgfpathrectangle{\pgfqpoint{0.100000in}{0.100000in}}{\pgfqpoint{3.608454in}{2.310000in}}%
\pgfusepath{clip}%
\pgfsetbuttcap%
\pgfsetmiterjoin%
\definecolor{currentfill}{rgb}{0.000000,0.592157,0.703922}%
\pgfsetfillcolor{currentfill}%
\pgfsetlinewidth{0.000000pt}%
\definecolor{currentstroke}{rgb}{0.000000,0.000000,0.000000}%
\pgfsetstrokecolor{currentstroke}%
\pgfsetstrokeopacity{0.000000}%
\pgfsetdash{}{0pt}%
\pgfpathmoveto{\pgfqpoint{1.805378in}{0.356576in}}%
\pgfpathlineto{\pgfqpoint{1.803098in}{0.369070in}}%
\pgfpathlineto{\pgfqpoint{1.799243in}{0.378831in}}%
\pgfpathlineto{\pgfqpoint{1.791844in}{0.386485in}}%
\pgfpathlineto{\pgfqpoint{1.787001in}{0.397445in}}%
\pgfpathlineto{\pgfqpoint{1.789197in}{0.406267in}}%
\pgfpathlineto{\pgfqpoint{1.784237in}{0.419600in}}%
\pgfpathlineto{\pgfqpoint{1.786564in}{0.431356in}}%
\pgfpathlineto{\pgfqpoint{1.783350in}{0.432395in}}%
\pgfpathlineto{\pgfqpoint{1.782660in}{0.441230in}}%
\pgfpathlineto{\pgfqpoint{1.770087in}{0.446743in}}%
\pgfpathlineto{\pgfqpoint{1.764022in}{0.455756in}}%
\pgfpathlineto{\pgfqpoint{1.759560in}{0.458124in}}%
\pgfpathlineto{\pgfqpoint{1.758847in}{0.466132in}}%
\pgfpathlineto{\pgfqpoint{1.752364in}{0.474146in}}%
\pgfpathlineto{\pgfqpoint{1.746457in}{0.486543in}}%
\pgfpathlineto{\pgfqpoint{1.737580in}{0.491300in}}%
\pgfpathlineto{\pgfqpoint{1.744553in}{0.490983in}}%
\pgfpathlineto{\pgfqpoint{1.795506in}{0.488863in}}%
\pgfpathlineto{\pgfqpoint{1.795200in}{0.474847in}}%
\pgfpathlineto{\pgfqpoint{1.836840in}{0.475148in}}%
\pgfpathlineto{\pgfqpoint{1.834828in}{0.418520in}}%
\pgfpathlineto{\pgfqpoint{1.852427in}{0.418217in}}%
\pgfpathlineto{\pgfqpoint{1.856313in}{0.408047in}}%
\pgfpathlineto{\pgfqpoint{1.857651in}{0.393438in}}%
\pgfpathlineto{\pgfqpoint{1.861148in}{0.393352in}}%
\pgfpathlineto{\pgfqpoint{1.859982in}{0.371433in}}%
\pgfpathlineto{\pgfqpoint{1.821677in}{0.373163in}}%
\pgfpathlineto{\pgfqpoint{1.817184in}{0.364409in}}%
\pgfpathlineto{\pgfqpoint{1.805378in}{0.356576in}}%
\pgfpathclose%
\pgfusepath{fill}%
\end{pgfscope}%
\begin{pgfscope}%
\pgfpathrectangle{\pgfqpoint{0.100000in}{0.100000in}}{\pgfqpoint{3.608454in}{2.310000in}}%
\pgfusepath{clip}%
\pgfsetbuttcap%
\pgfsetmiterjoin%
\definecolor{currentfill}{rgb}{0.000000,0.596078,0.701961}%
\pgfsetfillcolor{currentfill}%
\pgfsetlinewidth{0.000000pt}%
\definecolor{currentstroke}{rgb}{0.000000,0.000000,0.000000}%
\pgfsetstrokecolor{currentstroke}%
\pgfsetstrokeopacity{0.000000}%
\pgfsetdash{}{0pt}%
\pgfpathmoveto{\pgfqpoint{2.983960in}{0.768316in}}%
\pgfpathlineto{\pgfqpoint{2.977965in}{0.773429in}}%
\pgfpathlineto{\pgfqpoint{2.977554in}{0.783183in}}%
\pgfpathlineto{\pgfqpoint{2.983748in}{0.789742in}}%
\pgfpathlineto{\pgfqpoint{2.990030in}{0.793119in}}%
\pgfpathlineto{\pgfqpoint{2.986023in}{0.797823in}}%
\pgfpathlineto{\pgfqpoint{2.994265in}{0.796799in}}%
\pgfpathlineto{\pgfqpoint{3.002283in}{0.799408in}}%
\pgfpathlineto{\pgfqpoint{3.007825in}{0.797264in}}%
\pgfpathlineto{\pgfqpoint{3.013726in}{0.809687in}}%
\pgfpathlineto{\pgfqpoint{3.009224in}{0.814410in}}%
\pgfpathlineto{\pgfqpoint{3.015577in}{0.818559in}}%
\pgfpathlineto{\pgfqpoint{3.019148in}{0.828947in}}%
\pgfpathlineto{\pgfqpoint{3.023851in}{0.824906in}}%
\pgfpathlineto{\pgfqpoint{3.035156in}{0.827393in}}%
\pgfpathlineto{\pgfqpoint{3.040729in}{0.821509in}}%
\pgfpathlineto{\pgfqpoint{3.042224in}{0.818383in}}%
\pgfpathlineto{\pgfqpoint{3.038091in}{0.806725in}}%
\pgfpathlineto{\pgfqpoint{3.040053in}{0.796620in}}%
\pgfpathlineto{\pgfqpoint{3.031881in}{0.783770in}}%
\pgfpathlineto{\pgfqpoint{3.030760in}{0.775066in}}%
\pgfpathlineto{\pgfqpoint{3.020788in}{0.780741in}}%
\pgfpathlineto{\pgfqpoint{3.005604in}{0.782838in}}%
\pgfpathlineto{\pgfqpoint{2.996804in}{0.773153in}}%
\pgfpathlineto{\pgfqpoint{2.988034in}{0.774133in}}%
\pgfpathlineto{\pgfqpoint{2.983960in}{0.768316in}}%
\pgfpathclose%
\pgfusepath{fill}%
\end{pgfscope}%
\begin{pgfscope}%
\pgfpathrectangle{\pgfqpoint{0.100000in}{0.100000in}}{\pgfqpoint{3.608454in}{2.310000in}}%
\pgfusepath{clip}%
\pgfsetbuttcap%
\pgfsetmiterjoin%
\definecolor{currentfill}{rgb}{0.000000,0.364706,0.817647}%
\pgfsetfillcolor{currentfill}%
\pgfsetlinewidth{0.000000pt}%
\definecolor{currentstroke}{rgb}{0.000000,0.000000,0.000000}%
\pgfsetstrokecolor{currentstroke}%
\pgfsetstrokeopacity{0.000000}%
\pgfsetdash{}{0pt}%
\pgfpathmoveto{\pgfqpoint{1.571835in}{1.569055in}}%
\pgfpathlineto{\pgfqpoint{1.580064in}{1.651542in}}%
\pgfpathlineto{\pgfqpoint{1.583107in}{1.682203in}}%
\pgfpathlineto{\pgfqpoint{1.643808in}{1.676503in}}%
\pgfpathlineto{\pgfqpoint{1.655854in}{1.675393in}}%
\pgfpathlineto{\pgfqpoint{1.655483in}{1.658354in}}%
\pgfpathlineto{\pgfqpoint{1.653148in}{1.631106in}}%
\pgfpathlineto{\pgfqpoint{1.654543in}{1.630993in}}%
\pgfpathlineto{\pgfqpoint{1.652617in}{1.603533in}}%
\pgfpathlineto{\pgfqpoint{1.655822in}{1.596346in}}%
\pgfpathlineto{\pgfqpoint{1.654522in}{1.581965in}}%
\pgfpathlineto{\pgfqpoint{1.656181in}{1.573226in}}%
\pgfpathlineto{\pgfqpoint{1.654457in}{1.551499in}}%
\pgfpathlineto{\pgfqpoint{1.611187in}{1.555209in}}%
\pgfpathlineto{\pgfqpoint{1.610795in}{1.551818in}}%
\pgfpathlineto{\pgfqpoint{1.570505in}{1.555595in}}%
\pgfpathlineto{\pgfqpoint{1.571835in}{1.569055in}}%
\pgfpathclose%
\pgfusepath{fill}%
\end{pgfscope}%
\begin{pgfscope}%
\pgfpathrectangle{\pgfqpoint{0.100000in}{0.100000in}}{\pgfqpoint{3.608454in}{2.310000in}}%
\pgfusepath{clip}%
\pgfsetbuttcap%
\pgfsetmiterjoin%
\definecolor{currentfill}{rgb}{0.000000,0.937255,0.531373}%
\pgfsetfillcolor{currentfill}%
\pgfsetlinewidth{0.000000pt}%
\definecolor{currentstroke}{rgb}{0.000000,0.000000,0.000000}%
\pgfsetstrokecolor{currentstroke}%
\pgfsetstrokeopacity{0.000000}%
\pgfsetdash{}{0pt}%
\pgfpathmoveto{\pgfqpoint{2.393376in}{0.840961in}}%
\pgfpathlineto{\pgfqpoint{2.393171in}{0.837522in}}%
\pgfpathlineto{\pgfqpoint{2.384135in}{0.837054in}}%
\pgfpathlineto{\pgfqpoint{2.384564in}{0.826478in}}%
\pgfpathlineto{\pgfqpoint{2.378401in}{0.826111in}}%
\pgfpathlineto{\pgfqpoint{2.377206in}{0.836777in}}%
\pgfpathlineto{\pgfqpoint{2.371879in}{0.836595in}}%
\pgfpathlineto{\pgfqpoint{2.366196in}{0.837164in}}%
\pgfpathlineto{\pgfqpoint{2.364353in}{0.841406in}}%
\pgfpathlineto{\pgfqpoint{2.370206in}{0.848561in}}%
\pgfpathlineto{\pgfqpoint{2.363126in}{0.851002in}}%
\pgfpathlineto{\pgfqpoint{2.363366in}{0.855813in}}%
\pgfpathlineto{\pgfqpoint{2.368994in}{0.861333in}}%
\pgfpathlineto{\pgfqpoint{2.367614in}{0.868244in}}%
\pgfpathlineto{\pgfqpoint{2.362240in}{0.870230in}}%
\pgfpathlineto{\pgfqpoint{2.366771in}{0.881093in}}%
\pgfpathlineto{\pgfqpoint{2.366452in}{0.887645in}}%
\pgfpathlineto{\pgfqpoint{2.362887in}{0.893699in}}%
\pgfpathlineto{\pgfqpoint{2.362531in}{0.900120in}}%
\pgfpathlineto{\pgfqpoint{2.356517in}{0.914773in}}%
\pgfpathlineto{\pgfqpoint{2.362379in}{0.917175in}}%
\pgfpathlineto{\pgfqpoint{2.356747in}{0.922152in}}%
\pgfpathlineto{\pgfqpoint{2.361535in}{0.931873in}}%
\pgfpathlineto{\pgfqpoint{2.370618in}{0.931411in}}%
\pgfpathlineto{\pgfqpoint{2.365618in}{0.938943in}}%
\pgfpathlineto{\pgfqpoint{2.369367in}{0.943780in}}%
\pgfpathlineto{\pgfqpoint{2.364126in}{0.947506in}}%
\pgfpathlineto{\pgfqpoint{2.376340in}{0.951894in}}%
\pgfpathlineto{\pgfqpoint{2.377959in}{0.956142in}}%
\pgfpathlineto{\pgfqpoint{2.372346in}{0.959262in}}%
\pgfpathlineto{\pgfqpoint{2.372289in}{0.959362in}}%
\pgfpathlineto{\pgfqpoint{2.391745in}{0.960160in}}%
\pgfpathlineto{\pgfqpoint{2.392184in}{0.949754in}}%
\pgfpathlineto{\pgfqpoint{2.405968in}{0.950277in}}%
\pgfpathlineto{\pgfqpoint{2.406598in}{0.936460in}}%
\pgfpathlineto{\pgfqpoint{2.408409in}{0.898240in}}%
\pgfpathlineto{\pgfqpoint{2.401625in}{0.897879in}}%
\pgfpathlineto{\pgfqpoint{2.401868in}{0.893227in}}%
\pgfpathlineto{\pgfqpoint{2.391193in}{0.892674in}}%
\pgfpathlineto{\pgfqpoint{2.395715in}{0.879087in}}%
\pgfpathlineto{\pgfqpoint{2.396470in}{0.865278in}}%
\pgfpathlineto{\pgfqpoint{2.389913in}{0.858059in}}%
\pgfpathlineto{\pgfqpoint{2.393315in}{0.855864in}}%
\pgfpathlineto{\pgfqpoint{2.393376in}{0.840961in}}%
\pgfpathclose%
\pgfusepath{fill}%
\end{pgfscope}%
\begin{pgfscope}%
\pgfpathrectangle{\pgfqpoint{0.100000in}{0.100000in}}{\pgfqpoint{3.608454in}{2.310000in}}%
\pgfusepath{clip}%
\pgfsetbuttcap%
\pgfsetmiterjoin%
\definecolor{currentfill}{rgb}{0.000000,0.733333,0.633333}%
\pgfsetfillcolor{currentfill}%
\pgfsetlinewidth{0.000000pt}%
\definecolor{currentstroke}{rgb}{0.000000,0.000000,0.000000}%
\pgfsetstrokecolor{currentstroke}%
\pgfsetstrokeopacity{0.000000}%
\pgfsetdash{}{0pt}%
\pgfpathmoveto{\pgfqpoint{0.480357in}{2.020001in}}%
\pgfpathlineto{\pgfqpoint{0.492958in}{2.047705in}}%
\pgfpathlineto{\pgfqpoint{0.494929in}{2.056692in}}%
\pgfpathlineto{\pgfqpoint{0.510061in}{2.083946in}}%
\pgfpathlineto{\pgfqpoint{0.514910in}{2.098885in}}%
\pgfpathlineto{\pgfqpoint{0.524396in}{2.121590in}}%
\pgfpathlineto{\pgfqpoint{0.525797in}{2.131179in}}%
\pgfpathlineto{\pgfqpoint{0.538739in}{2.126084in}}%
\pgfpathlineto{\pgfqpoint{0.557819in}{2.120432in}}%
\pgfpathlineto{\pgfqpoint{0.556089in}{2.115055in}}%
\pgfpathlineto{\pgfqpoint{0.551015in}{2.110628in}}%
\pgfpathlineto{\pgfqpoint{0.550017in}{2.103751in}}%
\pgfpathlineto{\pgfqpoint{0.543848in}{2.096029in}}%
\pgfpathlineto{\pgfqpoint{0.538605in}{2.079539in}}%
\pgfpathlineto{\pgfqpoint{0.521499in}{2.085093in}}%
\pgfpathlineto{\pgfqpoint{0.518018in}{2.074545in}}%
\pgfpathlineto{\pgfqpoint{0.521244in}{2.073470in}}%
\pgfpathlineto{\pgfqpoint{0.512881in}{2.048028in}}%
\pgfpathlineto{\pgfqpoint{0.519006in}{2.044536in}}%
\pgfpathlineto{\pgfqpoint{0.512206in}{2.022672in}}%
\pgfpathlineto{\pgfqpoint{0.505570in}{2.024850in}}%
\pgfpathlineto{\pgfqpoint{0.504063in}{2.019289in}}%
\pgfpathlineto{\pgfqpoint{0.498912in}{2.014543in}}%
\pgfpathlineto{\pgfqpoint{0.480357in}{2.020001in}}%
\pgfpathclose%
\pgfusepath{fill}%
\end{pgfscope}%
\begin{pgfscope}%
\pgfpathrectangle{\pgfqpoint{0.100000in}{0.100000in}}{\pgfqpoint{3.608454in}{2.310000in}}%
\pgfusepath{clip}%
\pgfsetbuttcap%
\pgfsetmiterjoin%
\definecolor{currentfill}{rgb}{0.000000,0.827451,0.586275}%
\pgfsetfillcolor{currentfill}%
\pgfsetlinewidth{0.000000pt}%
\definecolor{currentstroke}{rgb}{0.000000,0.000000,0.000000}%
\pgfsetstrokecolor{currentstroke}%
\pgfsetstrokeopacity{0.000000}%
\pgfsetdash{}{0pt}%
\pgfpathmoveto{\pgfqpoint{0.375397in}{0.571894in}}%
\pgfpathlineto{\pgfqpoint{0.374645in}{0.576884in}}%
\pgfpathlineto{\pgfqpoint{0.375793in}{0.577086in}}%
\pgfpathlineto{\pgfqpoint{0.375397in}{0.571894in}}%
\pgfpathclose%
\pgfusepath{fill}%
\end{pgfscope}%
\begin{pgfscope}%
\pgfpathrectangle{\pgfqpoint{0.100000in}{0.100000in}}{\pgfqpoint{3.608454in}{2.310000in}}%
\pgfusepath{clip}%
\pgfsetbuttcap%
\pgfsetmiterjoin%
\definecolor{currentfill}{rgb}{0.000000,0.827451,0.586275}%
\pgfsetfillcolor{currentfill}%
\pgfsetlinewidth{0.000000pt}%
\definecolor{currentstroke}{rgb}{0.000000,0.000000,0.000000}%
\pgfsetstrokecolor{currentstroke}%
\pgfsetstrokeopacity{0.000000}%
\pgfsetdash{}{0pt}%
\pgfpathmoveto{\pgfqpoint{0.409283in}{0.538102in}}%
\pgfpathlineto{\pgfqpoint{0.408841in}{0.542212in}}%
\pgfpathlineto{\pgfqpoint{0.410919in}{0.540658in}}%
\pgfpathlineto{\pgfqpoint{0.412176in}{0.541182in}}%
\pgfpathlineto{\pgfqpoint{0.411443in}{0.546443in}}%
\pgfpathlineto{\pgfqpoint{0.412787in}{0.546879in}}%
\pgfpathlineto{\pgfqpoint{0.415145in}{0.543509in}}%
\pgfpathlineto{\pgfqpoint{0.414661in}{0.542347in}}%
\pgfpathlineto{\pgfqpoint{0.415729in}{0.539263in}}%
\pgfpathlineto{\pgfqpoint{0.410458in}{0.539171in}}%
\pgfpathlineto{\pgfqpoint{0.409283in}{0.538102in}}%
\pgfpathclose%
\pgfusepath{fill}%
\end{pgfscope}%
\begin{pgfscope}%
\pgfpathrectangle{\pgfqpoint{0.100000in}{0.100000in}}{\pgfqpoint{3.608454in}{2.310000in}}%
\pgfusepath{clip}%
\pgfsetbuttcap%
\pgfsetmiterjoin%
\definecolor{currentfill}{rgb}{0.000000,0.827451,0.586275}%
\pgfsetfillcolor{currentfill}%
\pgfsetlinewidth{0.000000pt}%
\definecolor{currentstroke}{rgb}{0.000000,0.000000,0.000000}%
\pgfsetstrokecolor{currentstroke}%
\pgfsetstrokeopacity{0.000000}%
\pgfsetdash{}{0pt}%
\pgfpathmoveto{\pgfqpoint{0.420907in}{0.523072in}}%
\pgfpathlineto{\pgfqpoint{0.420394in}{0.526482in}}%
\pgfpathlineto{\pgfqpoint{0.425789in}{0.526498in}}%
\pgfpathlineto{\pgfqpoint{0.427996in}{0.528757in}}%
\pgfpathlineto{\pgfqpoint{0.429378in}{0.528347in}}%
\pgfpathlineto{\pgfqpoint{0.431337in}{0.526034in}}%
\pgfpathlineto{\pgfqpoint{0.429229in}{0.525606in}}%
\pgfpathlineto{\pgfqpoint{0.428074in}{0.524089in}}%
\pgfpathlineto{\pgfqpoint{0.429993in}{0.521225in}}%
\pgfpathlineto{\pgfqpoint{0.430475in}{0.519141in}}%
\pgfpathlineto{\pgfqpoint{0.428313in}{0.519015in}}%
\pgfpathlineto{\pgfqpoint{0.427037in}{0.520588in}}%
\pgfpathlineto{\pgfqpoint{0.423736in}{0.522918in}}%
\pgfpathlineto{\pgfqpoint{0.420907in}{0.523072in}}%
\pgfpathclose%
\pgfusepath{fill}%
\end{pgfscope}%
\begin{pgfscope}%
\pgfpathrectangle{\pgfqpoint{0.100000in}{0.100000in}}{\pgfqpoint{3.608454in}{2.310000in}}%
\pgfusepath{clip}%
\pgfsetbuttcap%
\pgfsetmiterjoin%
\definecolor{currentfill}{rgb}{0.000000,0.827451,0.586275}%
\pgfsetfillcolor{currentfill}%
\pgfsetlinewidth{0.000000pt}%
\definecolor{currentstroke}{rgb}{0.000000,0.000000,0.000000}%
\pgfsetstrokecolor{currentstroke}%
\pgfsetstrokeopacity{0.000000}%
\pgfsetdash{}{0pt}%
\pgfpathmoveto{\pgfqpoint{0.417557in}{0.530445in}}%
\pgfpathlineto{\pgfqpoint{0.417550in}{0.532867in}}%
\pgfpathlineto{\pgfqpoint{0.412931in}{0.535157in}}%
\pgfpathlineto{\pgfqpoint{0.414886in}{0.536599in}}%
\pgfpathlineto{\pgfqpoint{0.417623in}{0.533713in}}%
\pgfpathlineto{\pgfqpoint{0.421451in}{0.532232in}}%
\pgfpathlineto{\pgfqpoint{0.423992in}{0.533859in}}%
\pgfpathlineto{\pgfqpoint{0.424555in}{0.531246in}}%
\pgfpathlineto{\pgfqpoint{0.421848in}{0.531121in}}%
\pgfpathlineto{\pgfqpoint{0.419783in}{0.529378in}}%
\pgfpathlineto{\pgfqpoint{0.417557in}{0.530445in}}%
\pgfpathclose%
\pgfusepath{fill}%
\end{pgfscope}%
\begin{pgfscope}%
\pgfpathrectangle{\pgfqpoint{0.100000in}{0.100000in}}{\pgfqpoint{3.608454in}{2.310000in}}%
\pgfusepath{clip}%
\pgfsetbuttcap%
\pgfsetmiterjoin%
\definecolor{currentfill}{rgb}{0.000000,0.827451,0.586275}%
\pgfsetfillcolor{currentfill}%
\pgfsetlinewidth{0.000000pt}%
\definecolor{currentstroke}{rgb}{0.000000,0.000000,0.000000}%
\pgfsetstrokecolor{currentstroke}%
\pgfsetstrokeopacity{0.000000}%
\pgfsetdash{}{0pt}%
\pgfpathmoveto{\pgfqpoint{0.389946in}{0.575682in}}%
\pgfpathlineto{\pgfqpoint{0.389432in}{0.578330in}}%
\pgfpathlineto{\pgfqpoint{0.392641in}{0.578848in}}%
\pgfpathlineto{\pgfqpoint{0.393455in}{0.577648in}}%
\pgfpathlineto{\pgfqpoint{0.392694in}{0.575332in}}%
\pgfpathlineto{\pgfqpoint{0.389946in}{0.575682in}}%
\pgfpathclose%
\pgfusepath{fill}%
\end{pgfscope}%
\begin{pgfscope}%
\pgfpathrectangle{\pgfqpoint{0.100000in}{0.100000in}}{\pgfqpoint{3.608454in}{2.310000in}}%
\pgfusepath{clip}%
\pgfsetbuttcap%
\pgfsetmiterjoin%
\definecolor{currentfill}{rgb}{0.000000,0.827451,0.586275}%
\pgfsetfillcolor{currentfill}%
\pgfsetlinewidth{0.000000pt}%
\definecolor{currentstroke}{rgb}{0.000000,0.000000,0.000000}%
\pgfsetstrokecolor{currentstroke}%
\pgfsetstrokeopacity{0.000000}%
\pgfsetdash{}{0pt}%
\pgfpathmoveto{\pgfqpoint{0.432376in}{0.516958in}}%
\pgfpathlineto{\pgfqpoint{0.431679in}{0.519437in}}%
\pgfpathlineto{\pgfqpoint{0.436217in}{0.517285in}}%
\pgfpathlineto{\pgfqpoint{0.434890in}{0.515977in}}%
\pgfpathlineto{\pgfqpoint{0.432376in}{0.516958in}}%
\pgfpathclose%
\pgfusepath{fill}%
\end{pgfscope}%
\begin{pgfscope}%
\pgfpathrectangle{\pgfqpoint{0.100000in}{0.100000in}}{\pgfqpoint{3.608454in}{2.310000in}}%
\pgfusepath{clip}%
\pgfsetbuttcap%
\pgfsetmiterjoin%
\definecolor{currentfill}{rgb}{0.000000,0.827451,0.586275}%
\pgfsetfillcolor{currentfill}%
\pgfsetlinewidth{0.000000pt}%
\definecolor{currentstroke}{rgb}{0.000000,0.000000,0.000000}%
\pgfsetstrokecolor{currentstroke}%
\pgfsetstrokeopacity{0.000000}%
\pgfsetdash{}{0pt}%
\pgfpathmoveto{\pgfqpoint{0.378939in}{0.591628in}}%
\pgfpathlineto{\pgfqpoint{0.379145in}{0.593528in}}%
\pgfpathlineto{\pgfqpoint{0.381451in}{0.592746in}}%
\pgfpathlineto{\pgfqpoint{0.380887in}{0.590969in}}%
\pgfpathlineto{\pgfqpoint{0.378939in}{0.591628in}}%
\pgfpathclose%
\pgfusepath{fill}%
\end{pgfscope}%
\begin{pgfscope}%
\pgfpathrectangle{\pgfqpoint{0.100000in}{0.100000in}}{\pgfqpoint{3.608454in}{2.310000in}}%
\pgfusepath{clip}%
\pgfsetbuttcap%
\pgfsetmiterjoin%
\definecolor{currentfill}{rgb}{0.000000,0.827451,0.586275}%
\pgfsetfillcolor{currentfill}%
\pgfsetlinewidth{0.000000pt}%
\definecolor{currentstroke}{rgb}{0.000000,0.000000,0.000000}%
\pgfsetstrokecolor{currentstroke}%
\pgfsetstrokeopacity{0.000000}%
\pgfsetdash{}{0pt}%
\pgfpathmoveto{\pgfqpoint{0.436549in}{0.519300in}}%
\pgfpathlineto{\pgfqpoint{0.436374in}{0.521884in}}%
\pgfpathlineto{\pgfqpoint{0.438554in}{0.523113in}}%
\pgfpathlineto{\pgfqpoint{0.439378in}{0.520715in}}%
\pgfpathlineto{\pgfqpoint{0.436549in}{0.519300in}}%
\pgfpathclose%
\pgfusepath{fill}%
\end{pgfscope}%
\begin{pgfscope}%
\pgfpathrectangle{\pgfqpoint{0.100000in}{0.100000in}}{\pgfqpoint{3.608454in}{2.310000in}}%
\pgfusepath{clip}%
\pgfsetbuttcap%
\pgfsetmiterjoin%
\definecolor{currentfill}{rgb}{0.000000,0.827451,0.586275}%
\pgfsetfillcolor{currentfill}%
\pgfsetlinewidth{0.000000pt}%
\definecolor{currentstroke}{rgb}{0.000000,0.000000,0.000000}%
\pgfsetstrokecolor{currentstroke}%
\pgfsetstrokeopacity{0.000000}%
\pgfsetdash{}{0pt}%
\pgfpathmoveto{\pgfqpoint{0.365932in}{0.607246in}}%
\pgfpathlineto{\pgfqpoint{0.366004in}{0.609622in}}%
\pgfpathlineto{\pgfqpoint{0.369338in}{0.609095in}}%
\pgfpathlineto{\pgfqpoint{0.371550in}{0.607762in}}%
\pgfpathlineto{\pgfqpoint{0.374186in}{0.608460in}}%
\pgfpathlineto{\pgfqpoint{0.375594in}{0.606756in}}%
\pgfpathlineto{\pgfqpoint{0.371710in}{0.606454in}}%
\pgfpathlineto{\pgfqpoint{0.370752in}{0.604856in}}%
\pgfpathlineto{\pgfqpoint{0.368700in}{0.607913in}}%
\pgfpathlineto{\pgfqpoint{0.365932in}{0.607246in}}%
\pgfpathclose%
\pgfusepath{fill}%
\end{pgfscope}%
\begin{pgfscope}%
\pgfpathrectangle{\pgfqpoint{0.100000in}{0.100000in}}{\pgfqpoint{3.608454in}{2.310000in}}%
\pgfusepath{clip}%
\pgfsetbuttcap%
\pgfsetmiterjoin%
\definecolor{currentfill}{rgb}{0.000000,0.827451,0.586275}%
\pgfsetfillcolor{currentfill}%
\pgfsetlinewidth{0.000000pt}%
\definecolor{currentstroke}{rgb}{0.000000,0.000000,0.000000}%
\pgfsetstrokecolor{currentstroke}%
\pgfsetstrokeopacity{0.000000}%
\pgfsetdash{}{0pt}%
\pgfpathmoveto{\pgfqpoint{0.480883in}{0.478483in}}%
\pgfpathlineto{\pgfqpoint{0.480302in}{0.480201in}}%
\pgfpathlineto{\pgfqpoint{0.484465in}{0.480081in}}%
\pgfpathlineto{\pgfqpoint{0.485718in}{0.477885in}}%
\pgfpathlineto{\pgfqpoint{0.484820in}{0.476951in}}%
\pgfpathlineto{\pgfqpoint{0.480883in}{0.478483in}}%
\pgfpathclose%
\pgfusepath{fill}%
\end{pgfscope}%
\begin{pgfscope}%
\pgfpathrectangle{\pgfqpoint{0.100000in}{0.100000in}}{\pgfqpoint{3.608454in}{2.310000in}}%
\pgfusepath{clip}%
\pgfsetbuttcap%
\pgfsetmiterjoin%
\definecolor{currentfill}{rgb}{0.000000,0.827451,0.586275}%
\pgfsetfillcolor{currentfill}%
\pgfsetlinewidth{0.000000pt}%
\definecolor{currentstroke}{rgb}{0.000000,0.000000,0.000000}%
\pgfsetstrokecolor{currentstroke}%
\pgfsetstrokeopacity{0.000000}%
\pgfsetdash{}{0pt}%
\pgfpathmoveto{\pgfqpoint{0.344293in}{0.666210in}}%
\pgfpathlineto{\pgfqpoint{0.343816in}{0.669244in}}%
\pgfpathlineto{\pgfqpoint{0.344625in}{0.670409in}}%
\pgfpathlineto{\pgfqpoint{0.346843in}{0.669814in}}%
\pgfpathlineto{\pgfqpoint{0.346503in}{0.668280in}}%
\pgfpathlineto{\pgfqpoint{0.344293in}{0.666210in}}%
\pgfpathclose%
\pgfusepath{fill}%
\end{pgfscope}%
\begin{pgfscope}%
\pgfpathrectangle{\pgfqpoint{0.100000in}{0.100000in}}{\pgfqpoint{3.608454in}{2.310000in}}%
\pgfusepath{clip}%
\pgfsetbuttcap%
\pgfsetmiterjoin%
\definecolor{currentfill}{rgb}{0.000000,0.827451,0.586275}%
\pgfsetfillcolor{currentfill}%
\pgfsetlinewidth{0.000000pt}%
\definecolor{currentstroke}{rgb}{0.000000,0.000000,0.000000}%
\pgfsetstrokecolor{currentstroke}%
\pgfsetstrokeopacity{0.000000}%
\pgfsetdash{}{0pt}%
\pgfpathmoveto{\pgfqpoint{0.506806in}{0.461794in}}%
\pgfpathlineto{\pgfqpoint{0.507677in}{0.463806in}}%
\pgfpathlineto{\pgfqpoint{0.511473in}{0.462938in}}%
\pgfpathlineto{\pgfqpoint{0.510522in}{0.460924in}}%
\pgfpathlineto{\pgfqpoint{0.506806in}{0.461794in}}%
\pgfpathclose%
\pgfusepath{fill}%
\end{pgfscope}%
\begin{pgfscope}%
\pgfpathrectangle{\pgfqpoint{0.100000in}{0.100000in}}{\pgfqpoint{3.608454in}{2.310000in}}%
\pgfusepath{clip}%
\pgfsetbuttcap%
\pgfsetmiterjoin%
\definecolor{currentfill}{rgb}{0.000000,0.827451,0.586275}%
\pgfsetfillcolor{currentfill}%
\pgfsetlinewidth{0.000000pt}%
\definecolor{currentstroke}{rgb}{0.000000,0.000000,0.000000}%
\pgfsetstrokecolor{currentstroke}%
\pgfsetstrokeopacity{0.000000}%
\pgfsetdash{}{0pt}%
\pgfpathmoveto{\pgfqpoint{0.349460in}{0.680635in}}%
\pgfpathlineto{\pgfqpoint{0.349236in}{0.682931in}}%
\pgfpathlineto{\pgfqpoint{0.345635in}{0.686541in}}%
\pgfpathlineto{\pgfqpoint{0.348354in}{0.687196in}}%
\pgfpathlineto{\pgfqpoint{0.346567in}{0.689336in}}%
\pgfpathlineto{\pgfqpoint{0.347637in}{0.690563in}}%
\pgfpathlineto{\pgfqpoint{0.350022in}{0.690962in}}%
\pgfpathlineto{\pgfqpoint{0.351654in}{0.687198in}}%
\pgfpathlineto{\pgfqpoint{0.353748in}{0.683821in}}%
\pgfpathlineto{\pgfqpoint{0.353994in}{0.680863in}}%
\pgfpathlineto{\pgfqpoint{0.352204in}{0.679863in}}%
\pgfpathlineto{\pgfqpoint{0.349460in}{0.680635in}}%
\pgfpathclose%
\pgfusepath{fill}%
\end{pgfscope}%
\begin{pgfscope}%
\pgfpathrectangle{\pgfqpoint{0.100000in}{0.100000in}}{\pgfqpoint{3.608454in}{2.310000in}}%
\pgfusepath{clip}%
\pgfsetbuttcap%
\pgfsetmiterjoin%
\definecolor{currentfill}{rgb}{0.000000,0.827451,0.586275}%
\pgfsetfillcolor{currentfill}%
\pgfsetlinewidth{0.000000pt}%
\definecolor{currentstroke}{rgb}{0.000000,0.000000,0.000000}%
\pgfsetstrokecolor{currentstroke}%
\pgfsetstrokeopacity{0.000000}%
\pgfsetdash{}{0pt}%
\pgfpathmoveto{\pgfqpoint{0.556547in}{0.436939in}}%
\pgfpathlineto{\pgfqpoint{0.555985in}{0.439788in}}%
\pgfpathlineto{\pgfqpoint{0.557578in}{0.440221in}}%
\pgfpathlineto{\pgfqpoint{0.561322in}{0.439031in}}%
\pgfpathlineto{\pgfqpoint{0.563123in}{0.437521in}}%
\pgfpathlineto{\pgfqpoint{0.565195in}{0.437153in}}%
\pgfpathlineto{\pgfqpoint{0.566888in}{0.435455in}}%
\pgfpathlineto{\pgfqpoint{0.569188in}{0.437256in}}%
\pgfpathlineto{\pgfqpoint{0.572877in}{0.438751in}}%
\pgfpathlineto{\pgfqpoint{0.575932in}{0.435745in}}%
\pgfpathlineto{\pgfqpoint{0.574169in}{0.442233in}}%
\pgfpathlineto{\pgfqpoint{0.576542in}{0.443329in}}%
\pgfpathlineto{\pgfqpoint{0.579175in}{0.442765in}}%
\pgfpathlineto{\pgfqpoint{0.583426in}{0.440455in}}%
\pgfpathlineto{\pgfqpoint{0.582448in}{0.438627in}}%
\pgfpathlineto{\pgfqpoint{0.584561in}{0.436213in}}%
\pgfpathlineto{\pgfqpoint{0.586754in}{0.436810in}}%
\pgfpathlineto{\pgfqpoint{0.586008in}{0.433801in}}%
\pgfpathlineto{\pgfqpoint{0.582411in}{0.433473in}}%
\pgfpathlineto{\pgfqpoint{0.580948in}{0.429661in}}%
\pgfpathlineto{\pgfqpoint{0.574828in}{0.431855in}}%
\pgfpathlineto{\pgfqpoint{0.572584in}{0.432073in}}%
\pgfpathlineto{\pgfqpoint{0.571112in}{0.430136in}}%
\pgfpathlineto{\pgfqpoint{0.569377in}{0.433000in}}%
\pgfpathlineto{\pgfqpoint{0.567679in}{0.433841in}}%
\pgfpathlineto{\pgfqpoint{0.562725in}{0.434707in}}%
\pgfpathlineto{\pgfqpoint{0.560495in}{0.435926in}}%
\pgfpathlineto{\pgfqpoint{0.556547in}{0.436939in}}%
\pgfpathclose%
\pgfusepath{fill}%
\end{pgfscope}%
\begin{pgfscope}%
\pgfpathrectangle{\pgfqpoint{0.100000in}{0.100000in}}{\pgfqpoint{3.608454in}{2.310000in}}%
\pgfusepath{clip}%
\pgfsetbuttcap%
\pgfsetmiterjoin%
\definecolor{currentfill}{rgb}{0.000000,0.827451,0.586275}%
\pgfsetfillcolor{currentfill}%
\pgfsetlinewidth{0.000000pt}%
\definecolor{currentstroke}{rgb}{0.000000,0.000000,0.000000}%
\pgfsetstrokecolor{currentstroke}%
\pgfsetstrokeopacity{0.000000}%
\pgfsetdash{}{0pt}%
\pgfpathmoveto{\pgfqpoint{0.523468in}{0.452981in}}%
\pgfpathlineto{\pgfqpoint{0.523015in}{0.455703in}}%
\pgfpathlineto{\pgfqpoint{0.525225in}{0.455976in}}%
\pgfpathlineto{\pgfqpoint{0.525280in}{0.453766in}}%
\pgfpathlineto{\pgfqpoint{0.523468in}{0.452981in}}%
\pgfpathclose%
\pgfusepath{fill}%
\end{pgfscope}%
\begin{pgfscope}%
\pgfpathrectangle{\pgfqpoint{0.100000in}{0.100000in}}{\pgfqpoint{3.608454in}{2.310000in}}%
\pgfusepath{clip}%
\pgfsetbuttcap%
\pgfsetmiterjoin%
\definecolor{currentfill}{rgb}{0.000000,0.827451,0.586275}%
\pgfsetfillcolor{currentfill}%
\pgfsetlinewidth{0.000000pt}%
\definecolor{currentstroke}{rgb}{0.000000,0.000000,0.000000}%
\pgfsetstrokecolor{currentstroke}%
\pgfsetstrokeopacity{0.000000}%
\pgfsetdash{}{0pt}%
\pgfpathmoveto{\pgfqpoint{0.532273in}{0.445921in}}%
\pgfpathlineto{\pgfqpoint{0.532818in}{0.446891in}}%
\pgfpathlineto{\pgfqpoint{0.538427in}{0.447049in}}%
\pgfpathlineto{\pgfqpoint{0.541753in}{0.449221in}}%
\pgfpathlineto{\pgfqpoint{0.546018in}{0.448961in}}%
\pgfpathlineto{\pgfqpoint{0.548669in}{0.445595in}}%
\pgfpathlineto{\pgfqpoint{0.551094in}{0.448774in}}%
\pgfpathlineto{\pgfqpoint{0.553730in}{0.449651in}}%
\pgfpathlineto{\pgfqpoint{0.556658in}{0.449030in}}%
\pgfpathlineto{\pgfqpoint{0.558881in}{0.447475in}}%
\pgfpathlineto{\pgfqpoint{0.559652in}{0.444922in}}%
\pgfpathlineto{\pgfqpoint{0.559201in}{0.443279in}}%
\pgfpathlineto{\pgfqpoint{0.556808in}{0.441953in}}%
\pgfpathlineto{\pgfqpoint{0.548779in}{0.444295in}}%
\pgfpathlineto{\pgfqpoint{0.543555in}{0.442792in}}%
\pgfpathlineto{\pgfqpoint{0.541563in}{0.443699in}}%
\pgfpathlineto{\pgfqpoint{0.532273in}{0.445921in}}%
\pgfpathclose%
\pgfusepath{fill}%
\end{pgfscope}%
\begin{pgfscope}%
\pgfpathrectangle{\pgfqpoint{0.100000in}{0.100000in}}{\pgfqpoint{3.608454in}{2.310000in}}%
\pgfusepath{clip}%
\pgfsetbuttcap%
\pgfsetmiterjoin%
\definecolor{currentfill}{rgb}{0.000000,0.827451,0.586275}%
\pgfsetfillcolor{currentfill}%
\pgfsetlinewidth{0.000000pt}%
\definecolor{currentstroke}{rgb}{0.000000,0.000000,0.000000}%
\pgfsetstrokecolor{currentstroke}%
\pgfsetstrokeopacity{0.000000}%
\pgfsetdash{}{0pt}%
\pgfpathmoveto{\pgfqpoint{0.467157in}{0.485560in}}%
\pgfpathlineto{\pgfqpoint{0.465100in}{0.487760in}}%
\pgfpathlineto{\pgfqpoint{0.463877in}{0.490340in}}%
\pgfpathlineto{\pgfqpoint{0.461910in}{0.492281in}}%
\pgfpathlineto{\pgfqpoint{0.462381in}{0.494960in}}%
\pgfpathlineto{\pgfqpoint{0.466531in}{0.490940in}}%
\pgfpathlineto{\pgfqpoint{0.468975in}{0.485923in}}%
\pgfpathlineto{\pgfqpoint{0.467157in}{0.485560in}}%
\pgfpathclose%
\pgfusepath{fill}%
\end{pgfscope}%
\begin{pgfscope}%
\pgfpathrectangle{\pgfqpoint{0.100000in}{0.100000in}}{\pgfqpoint{3.608454in}{2.310000in}}%
\pgfusepath{clip}%
\pgfsetbuttcap%
\pgfsetmiterjoin%
\definecolor{currentfill}{rgb}{0.000000,0.827451,0.586275}%
\pgfsetfillcolor{currentfill}%
\pgfsetlinewidth{0.000000pt}%
\definecolor{currentstroke}{rgb}{0.000000,0.000000,0.000000}%
\pgfsetstrokecolor{currentstroke}%
\pgfsetstrokeopacity{0.000000}%
\pgfsetdash{}{0pt}%
\pgfpathmoveto{\pgfqpoint{0.456040in}{0.499043in}}%
\pgfpathlineto{\pgfqpoint{0.454673in}{0.501491in}}%
\pgfpathlineto{\pgfqpoint{0.451914in}{0.502371in}}%
\pgfpathlineto{\pgfqpoint{0.451045in}{0.506296in}}%
\pgfpathlineto{\pgfqpoint{0.454103in}{0.504564in}}%
\pgfpathlineto{\pgfqpoint{0.455369in}{0.502279in}}%
\pgfpathlineto{\pgfqpoint{0.457486in}{0.502766in}}%
\pgfpathlineto{\pgfqpoint{0.461584in}{0.502365in}}%
\pgfpathlineto{\pgfqpoint{0.463798in}{0.503703in}}%
\pgfpathlineto{\pgfqpoint{0.465884in}{0.502834in}}%
\pgfpathlineto{\pgfqpoint{0.466975in}{0.500406in}}%
\pgfpathlineto{\pgfqpoint{0.466593in}{0.498917in}}%
\pgfpathlineto{\pgfqpoint{0.464307in}{0.497731in}}%
\pgfpathlineto{\pgfqpoint{0.462347in}{0.499412in}}%
\pgfpathlineto{\pgfqpoint{0.461307in}{0.496536in}}%
\pgfpathlineto{\pgfqpoint{0.459853in}{0.497056in}}%
\pgfpathlineto{\pgfqpoint{0.457626in}{0.499895in}}%
\pgfpathlineto{\pgfqpoint{0.456040in}{0.499043in}}%
\pgfpathclose%
\pgfusepath{fill}%
\end{pgfscope}%
\begin{pgfscope}%
\pgfpathrectangle{\pgfqpoint{0.100000in}{0.100000in}}{\pgfqpoint{3.608454in}{2.310000in}}%
\pgfusepath{clip}%
\pgfsetbuttcap%
\pgfsetmiterjoin%
\definecolor{currentfill}{rgb}{0.000000,0.776471,0.611765}%
\pgfsetfillcolor{currentfill}%
\pgfsetlinewidth{0.000000pt}%
\definecolor{currentstroke}{rgb}{0.000000,0.000000,0.000000}%
\pgfsetstrokecolor{currentstroke}%
\pgfsetstrokeopacity{0.000000}%
\pgfsetdash{}{0pt}%
\pgfpathmoveto{\pgfqpoint{2.231585in}{0.901788in}}%
\pgfpathlineto{\pgfqpoint{2.224654in}{0.900658in}}%
\pgfpathlineto{\pgfqpoint{2.207240in}{0.900881in}}%
\pgfpathlineto{\pgfqpoint{2.207261in}{0.903173in}}%
\pgfpathlineto{\pgfqpoint{2.191374in}{0.903396in}}%
\pgfpathlineto{\pgfqpoint{2.183884in}{0.913246in}}%
\pgfpathlineto{\pgfqpoint{2.175110in}{0.918121in}}%
\pgfpathlineto{\pgfqpoint{2.175503in}{0.924555in}}%
\pgfpathlineto{\pgfqpoint{2.174128in}{0.929756in}}%
\pgfpathlineto{\pgfqpoint{2.169177in}{0.935220in}}%
\pgfpathlineto{\pgfqpoint{2.166804in}{0.946907in}}%
\pgfpathlineto{\pgfqpoint{2.165791in}{0.959401in}}%
\pgfpathlineto{\pgfqpoint{2.156410in}{0.959401in}}%
\pgfpathlineto{\pgfqpoint{2.155587in}{0.972321in}}%
\pgfpathlineto{\pgfqpoint{2.176495in}{0.972103in}}%
\pgfpathlineto{\pgfqpoint{2.176519in}{0.997121in}}%
\pgfpathlineto{\pgfqpoint{2.180404in}{1.000901in}}%
\pgfpathlineto{\pgfqpoint{2.190786in}{1.003624in}}%
\pgfpathlineto{\pgfqpoint{2.211366in}{1.003801in}}%
\pgfpathlineto{\pgfqpoint{2.218257in}{1.006289in}}%
\pgfpathlineto{\pgfqpoint{2.238992in}{1.006643in}}%
\pgfpathlineto{\pgfqpoint{2.238537in}{0.995926in}}%
\pgfpathlineto{\pgfqpoint{2.245351in}{0.995915in}}%
\pgfpathlineto{\pgfqpoint{2.249833in}{0.992475in}}%
\pgfpathlineto{\pgfqpoint{2.250959in}{0.985619in}}%
\pgfpathlineto{\pgfqpoint{2.258094in}{0.983370in}}%
\pgfpathlineto{\pgfqpoint{2.259281in}{0.978738in}}%
\pgfpathlineto{\pgfqpoint{2.259184in}{0.957741in}}%
\pgfpathlineto{\pgfqpoint{2.245528in}{0.957904in}}%
\pgfpathlineto{\pgfqpoint{2.245189in}{0.934157in}}%
\pgfpathlineto{\pgfqpoint{2.241965in}{0.929323in}}%
\pgfpathlineto{\pgfqpoint{2.231849in}{0.927551in}}%
\pgfpathlineto{\pgfqpoint{2.231585in}{0.901788in}}%
\pgfpathclose%
\pgfusepath{fill}%
\end{pgfscope}%
\begin{pgfscope}%
\pgfpathrectangle{\pgfqpoint{0.100000in}{0.100000in}}{\pgfqpoint{3.608454in}{2.310000in}}%
\pgfusepath{clip}%
\pgfsetbuttcap%
\pgfsetmiterjoin%
\definecolor{currentfill}{rgb}{0.000000,0.721569,0.639216}%
\pgfsetfillcolor{currentfill}%
\pgfsetlinewidth{0.000000pt}%
\definecolor{currentstroke}{rgb}{0.000000,0.000000,0.000000}%
\pgfsetstrokecolor{currentstroke}%
\pgfsetstrokeopacity{0.000000}%
\pgfsetdash{}{0pt}%
\pgfpathmoveto{\pgfqpoint{0.695564in}{1.998853in}}%
\pgfpathlineto{\pgfqpoint{0.656754in}{2.010285in}}%
\pgfpathlineto{\pgfqpoint{0.620648in}{2.021281in}}%
\pgfpathlineto{\pgfqpoint{0.623427in}{2.025636in}}%
\pgfpathlineto{\pgfqpoint{0.622394in}{2.035924in}}%
\pgfpathlineto{\pgfqpoint{0.630035in}{2.038493in}}%
\pgfpathlineto{\pgfqpoint{0.629691in}{2.048489in}}%
\pgfpathlineto{\pgfqpoint{0.634662in}{2.052316in}}%
\pgfpathlineto{\pgfqpoint{0.636428in}{2.062774in}}%
\pgfpathlineto{\pgfqpoint{0.628826in}{2.076535in}}%
\pgfpathlineto{\pgfqpoint{0.630988in}{2.086576in}}%
\pgfpathlineto{\pgfqpoint{0.634943in}{2.089043in}}%
\pgfpathlineto{\pgfqpoint{0.641984in}{2.086967in}}%
\pgfpathlineto{\pgfqpoint{0.653479in}{2.086201in}}%
\pgfpathlineto{\pgfqpoint{0.650732in}{2.091686in}}%
\pgfpathlineto{\pgfqpoint{0.656547in}{2.111660in}}%
\pgfpathlineto{\pgfqpoint{0.661241in}{2.110265in}}%
\pgfpathlineto{\pgfqpoint{0.719253in}{2.092888in}}%
\pgfpathlineto{\pgfqpoint{0.749167in}{2.084439in}}%
\pgfpathlineto{\pgfqpoint{0.744663in}{2.068773in}}%
\pgfpathlineto{\pgfqpoint{0.737141in}{2.068948in}}%
\pgfpathlineto{\pgfqpoint{0.723978in}{2.065521in}}%
\pgfpathlineto{\pgfqpoint{0.713071in}{2.066429in}}%
\pgfpathlineto{\pgfqpoint{0.707552in}{2.064386in}}%
\pgfpathlineto{\pgfqpoint{0.710428in}{2.058581in}}%
\pgfpathlineto{\pgfqpoint{0.710674in}{2.050338in}}%
\pgfpathlineto{\pgfqpoint{0.702993in}{2.049946in}}%
\pgfpathlineto{\pgfqpoint{0.695191in}{2.034615in}}%
\pgfpathlineto{\pgfqpoint{0.697834in}{2.026472in}}%
\pgfpathlineto{\pgfqpoint{0.692055in}{2.010580in}}%
\pgfpathlineto{\pgfqpoint{0.692032in}{2.003466in}}%
\pgfpathlineto{\pgfqpoint{0.695564in}{1.998853in}}%
\pgfpathclose%
\pgfusepath{fill}%
\end{pgfscope}%
\begin{pgfscope}%
\pgfpathrectangle{\pgfqpoint{0.100000in}{0.100000in}}{\pgfqpoint{3.608454in}{2.310000in}}%
\pgfusepath{clip}%
\pgfsetbuttcap%
\pgfsetmiterjoin%
\definecolor{currentfill}{rgb}{0.000000,0.505882,0.747059}%
\pgfsetfillcolor{currentfill}%
\pgfsetlinewidth{0.000000pt}%
\definecolor{currentstroke}{rgb}{0.000000,0.000000,0.000000}%
\pgfsetstrokecolor{currentstroke}%
\pgfsetstrokeopacity{0.000000}%
\pgfsetdash{}{0pt}%
\pgfpathmoveto{\pgfqpoint{1.748391in}{1.048651in}}%
\pgfpathlineto{\pgfqpoint{1.713023in}{1.050837in}}%
\pgfpathlineto{\pgfqpoint{1.678367in}{1.053341in}}%
\pgfpathlineto{\pgfqpoint{1.680896in}{1.087866in}}%
\pgfpathlineto{\pgfqpoint{1.683362in}{1.122232in}}%
\pgfpathlineto{\pgfqpoint{1.752524in}{1.117682in}}%
\pgfpathlineto{\pgfqpoint{1.787003in}{1.115616in}}%
\pgfpathlineto{\pgfqpoint{1.786242in}{1.101806in}}%
\pgfpathlineto{\pgfqpoint{1.785102in}{1.081107in}}%
\pgfpathlineto{\pgfqpoint{1.750416in}{1.083136in}}%
\pgfpathlineto{\pgfqpoint{1.748391in}{1.048651in}}%
\pgfpathclose%
\pgfusepath{fill}%
\end{pgfscope}%
\begin{pgfscope}%
\pgfpathrectangle{\pgfqpoint{0.100000in}{0.100000in}}{\pgfqpoint{3.608454in}{2.310000in}}%
\pgfusepath{clip}%
\pgfsetbuttcap%
\pgfsetmiterjoin%
\definecolor{currentfill}{rgb}{0.000000,0.654902,0.672549}%
\pgfsetfillcolor{currentfill}%
\pgfsetlinewidth{0.000000pt}%
\definecolor{currentstroke}{rgb}{0.000000,0.000000,0.000000}%
\pgfsetstrokecolor{currentstroke}%
\pgfsetstrokeopacity{0.000000}%
\pgfsetdash{}{0pt}%
\pgfpathmoveto{\pgfqpoint{1.884281in}{0.742465in}}%
\pgfpathlineto{\pgfqpoint{1.863869in}{0.728404in}}%
\pgfpathlineto{\pgfqpoint{1.866023in}{0.717001in}}%
\pgfpathlineto{\pgfqpoint{1.871649in}{0.712052in}}%
\pgfpathlineto{\pgfqpoint{1.870905in}{0.703446in}}%
\pgfpathlineto{\pgfqpoint{1.835408in}{0.704889in}}%
\pgfpathlineto{\pgfqpoint{1.826743in}{0.706842in}}%
\pgfpathlineto{\pgfqpoint{1.828725in}{0.748261in}}%
\pgfpathlineto{\pgfqpoint{1.818920in}{0.750454in}}%
\pgfpathlineto{\pgfqpoint{1.813328in}{0.744891in}}%
\pgfpathlineto{\pgfqpoint{1.804967in}{0.750514in}}%
\pgfpathlineto{\pgfqpoint{1.792433in}{0.750538in}}%
\pgfpathlineto{\pgfqpoint{1.786260in}{0.758457in}}%
\pgfpathlineto{\pgfqpoint{1.788954in}{0.799743in}}%
\pgfpathlineto{\pgfqpoint{1.842140in}{0.796889in}}%
\pgfpathlineto{\pgfqpoint{1.858178in}{0.766124in}}%
\pgfpathlineto{\pgfqpoint{1.870147in}{0.767494in}}%
\pgfpathlineto{\pgfqpoint{1.884281in}{0.742465in}}%
\pgfpathclose%
\pgfusepath{fill}%
\end{pgfscope}%
\begin{pgfscope}%
\pgfpathrectangle{\pgfqpoint{0.100000in}{0.100000in}}{\pgfqpoint{3.608454in}{2.310000in}}%
\pgfusepath{clip}%
\pgfsetbuttcap%
\pgfsetmiterjoin%
\definecolor{currentfill}{rgb}{0.000000,0.905882,0.547059}%
\pgfsetfillcolor{currentfill}%
\pgfsetlinewidth{0.000000pt}%
\definecolor{currentstroke}{rgb}{0.000000,0.000000,0.000000}%
\pgfsetstrokecolor{currentstroke}%
\pgfsetstrokeopacity{0.000000}%
\pgfsetdash{}{0pt}%
\pgfpathmoveto{\pgfqpoint{2.910647in}{0.649877in}}%
\pgfpathlineto{\pgfqpoint{2.900528in}{0.654212in}}%
\pgfpathlineto{\pgfqpoint{2.893044in}{0.665397in}}%
\pgfpathlineto{\pgfqpoint{2.882864in}{0.671272in}}%
\pgfpathlineto{\pgfqpoint{2.871699in}{0.674920in}}%
\pgfpathlineto{\pgfqpoint{2.866769in}{0.678489in}}%
\pgfpathlineto{\pgfqpoint{2.870453in}{0.691569in}}%
\pgfpathlineto{\pgfqpoint{2.878191in}{0.698489in}}%
\pgfpathlineto{\pgfqpoint{2.876684in}{0.701966in}}%
\pgfpathlineto{\pgfqpoint{2.881876in}{0.708207in}}%
\pgfpathlineto{\pgfqpoint{2.880133in}{0.713968in}}%
\pgfpathlineto{\pgfqpoint{2.888448in}{0.720934in}}%
\pgfpathlineto{\pgfqpoint{2.887245in}{0.726366in}}%
\pgfpathlineto{\pgfqpoint{2.957900in}{0.730883in}}%
\pgfpathlineto{\pgfqpoint{2.966459in}{0.731456in}}%
\pgfpathlineto{\pgfqpoint{2.971461in}{0.696012in}}%
\pgfpathlineto{\pgfqpoint{2.965631in}{0.682186in}}%
\pgfpathlineto{\pgfqpoint{2.965236in}{0.672986in}}%
\pgfpathlineto{\pgfqpoint{2.962277in}{0.669805in}}%
\pgfpathlineto{\pgfqpoint{2.955684in}{0.671251in}}%
\pgfpathlineto{\pgfqpoint{2.951329in}{0.676818in}}%
\pgfpathlineto{\pgfqpoint{2.946140in}{0.674377in}}%
\pgfpathlineto{\pgfqpoint{2.944970in}{0.667447in}}%
\pgfpathlineto{\pgfqpoint{2.915694in}{0.663204in}}%
\pgfpathlineto{\pgfqpoint{2.913541in}{0.650279in}}%
\pgfpathlineto{\pgfqpoint{2.910647in}{0.649877in}}%
\pgfpathclose%
\pgfusepath{fill}%
\end{pgfscope}%
\begin{pgfscope}%
\pgfpathrectangle{\pgfqpoint{0.100000in}{0.100000in}}{\pgfqpoint{3.608454in}{2.310000in}}%
\pgfusepath{clip}%
\pgfsetbuttcap%
\pgfsetmiterjoin%
\definecolor{currentfill}{rgb}{0.000000,0.686275,0.656863}%
\pgfsetfillcolor{currentfill}%
\pgfsetlinewidth{0.000000pt}%
\definecolor{currentstroke}{rgb}{0.000000,0.000000,0.000000}%
\pgfsetstrokecolor{currentstroke}%
\pgfsetstrokeopacity{0.000000}%
\pgfsetdash{}{0pt}%
\pgfpathmoveto{\pgfqpoint{2.543041in}{1.178054in}}%
\pgfpathlineto{\pgfqpoint{2.549236in}{1.162164in}}%
\pgfpathlineto{\pgfqpoint{2.548186in}{1.158629in}}%
\pgfpathlineto{\pgfqpoint{2.520509in}{1.156888in}}%
\pgfpathlineto{\pgfqpoint{2.519321in}{1.176572in}}%
\pgfpathlineto{\pgfqpoint{2.543041in}{1.178054in}}%
\pgfpathclose%
\pgfusepath{fill}%
\end{pgfscope}%
\begin{pgfscope}%
\pgfpathrectangle{\pgfqpoint{0.100000in}{0.100000in}}{\pgfqpoint{3.608454in}{2.310000in}}%
\pgfusepath{clip}%
\pgfsetbuttcap%
\pgfsetmiterjoin%
\definecolor{currentfill}{rgb}{0.000000,0.729412,0.635294}%
\pgfsetfillcolor{currentfill}%
\pgfsetlinewidth{0.000000pt}%
\definecolor{currentstroke}{rgb}{0.000000,0.000000,0.000000}%
\pgfsetstrokecolor{currentstroke}%
\pgfsetstrokeopacity{0.000000}%
\pgfsetdash{}{0pt}%
\pgfpathmoveto{\pgfqpoint{2.213676in}{1.928846in}}%
\pgfpathlineto{\pgfqpoint{2.214175in}{1.908040in}}%
\pgfpathlineto{\pgfqpoint{2.193329in}{1.907307in}}%
\pgfpathlineto{\pgfqpoint{2.193274in}{1.914649in}}%
\pgfpathlineto{\pgfqpoint{2.172320in}{1.914091in}}%
\pgfpathlineto{\pgfqpoint{2.172440in}{1.907196in}}%
\pgfpathlineto{\pgfqpoint{2.141306in}{1.906925in}}%
\pgfpathlineto{\pgfqpoint{2.142555in}{1.917955in}}%
\pgfpathlineto{\pgfqpoint{2.138352in}{1.920870in}}%
\pgfpathlineto{\pgfqpoint{2.132410in}{1.918482in}}%
\pgfpathlineto{\pgfqpoint{2.118937in}{1.925718in}}%
\pgfpathlineto{\pgfqpoint{2.118427in}{1.958403in}}%
\pgfpathlineto{\pgfqpoint{2.125367in}{1.958378in}}%
\pgfpathlineto{\pgfqpoint{2.124628in}{2.006820in}}%
\pgfpathlineto{\pgfqpoint{2.138354in}{2.007024in}}%
\pgfpathlineto{\pgfqpoint{2.138015in}{2.041726in}}%
\pgfpathlineto{\pgfqpoint{2.172503in}{2.042117in}}%
\pgfpathlineto{\pgfqpoint{2.172471in}{2.046322in}}%
\pgfpathlineto{\pgfqpoint{2.209745in}{2.046401in}}%
\pgfpathlineto{\pgfqpoint{2.211106in}{2.032691in}}%
\pgfpathlineto{\pgfqpoint{2.211215in}{2.006566in}}%
\pgfpathlineto{\pgfqpoint{2.212580in}{1.977269in}}%
\pgfpathlineto{\pgfqpoint{2.212805in}{1.945663in}}%
\pgfpathlineto{\pgfqpoint{2.213676in}{1.928846in}}%
\pgfpathclose%
\pgfusepath{fill}%
\end{pgfscope}%
\begin{pgfscope}%
\pgfpathrectangle{\pgfqpoint{0.100000in}{0.100000in}}{\pgfqpoint{3.608454in}{2.310000in}}%
\pgfusepath{clip}%
\pgfsetbuttcap%
\pgfsetmiterjoin%
\definecolor{currentfill}{rgb}{0.000000,0.521569,0.739216}%
\pgfsetfillcolor{currentfill}%
\pgfsetlinewidth{0.000000pt}%
\definecolor{currentstroke}{rgb}{0.000000,0.000000,0.000000}%
\pgfsetstrokecolor{currentstroke}%
\pgfsetstrokeopacity{0.000000}%
\pgfsetdash{}{0pt}%
\pgfpathmoveto{\pgfqpoint{1.598608in}{1.839877in}}%
\pgfpathlineto{\pgfqpoint{1.543825in}{1.845887in}}%
\pgfpathlineto{\pgfqpoint{1.546917in}{1.873677in}}%
\pgfpathlineto{\pgfqpoint{1.549754in}{1.873361in}}%
\pgfpathlineto{\pgfqpoint{1.553831in}{1.907730in}}%
\pgfpathlineto{\pgfqpoint{1.556192in}{1.907466in}}%
\pgfpathlineto{\pgfqpoint{1.559222in}{1.935142in}}%
\pgfpathlineto{\pgfqpoint{1.562162in}{1.934829in}}%
\pgfpathlineto{\pgfqpoint{1.582466in}{1.932625in}}%
\pgfpathlineto{\pgfqpoint{1.581763in}{1.926153in}}%
\pgfpathlineto{\pgfqpoint{1.588691in}{1.925458in}}%
\pgfpathlineto{\pgfqpoint{1.587963in}{1.918755in}}%
\pgfpathlineto{\pgfqpoint{1.602057in}{1.917205in}}%
\pgfpathlineto{\pgfqpoint{1.606786in}{1.914838in}}%
\pgfpathlineto{\pgfqpoint{1.664670in}{1.909384in}}%
\pgfpathlineto{\pgfqpoint{1.667609in}{1.909112in}}%
\pgfpathlineto{\pgfqpoint{1.661648in}{1.851140in}}%
\pgfpathlineto{\pgfqpoint{1.601301in}{1.856823in}}%
\pgfpathlineto{\pgfqpoint{1.598608in}{1.839877in}}%
\pgfpathclose%
\pgfusepath{fill}%
\end{pgfscope}%
\begin{pgfscope}%
\pgfpathrectangle{\pgfqpoint{0.100000in}{0.100000in}}{\pgfqpoint{3.608454in}{2.310000in}}%
\pgfusepath{clip}%
\pgfsetbuttcap%
\pgfsetmiterjoin%
\definecolor{currentfill}{rgb}{0.000000,0.768627,0.615686}%
\pgfsetfillcolor{currentfill}%
\pgfsetlinewidth{0.000000pt}%
\definecolor{currentstroke}{rgb}{0.000000,0.000000,0.000000}%
\pgfsetstrokecolor{currentstroke}%
\pgfsetstrokeopacity{0.000000}%
\pgfsetdash{}{0pt}%
\pgfpathmoveto{\pgfqpoint{2.618822in}{1.745974in}}%
\pgfpathlineto{\pgfqpoint{2.596450in}{1.744135in}}%
\pgfpathlineto{\pgfqpoint{2.593518in}{1.755812in}}%
\pgfpathlineto{\pgfqpoint{2.590087in}{1.761390in}}%
\pgfpathlineto{\pgfqpoint{2.594695in}{1.768652in}}%
\pgfpathlineto{\pgfqpoint{2.601730in}{1.786502in}}%
\pgfpathlineto{\pgfqpoint{2.602311in}{1.800403in}}%
\pgfpathlineto{\pgfqpoint{2.625838in}{1.802295in}}%
\pgfpathlineto{\pgfqpoint{2.628418in}{1.774736in}}%
\pgfpathlineto{\pgfqpoint{2.615872in}{1.773777in}}%
\pgfpathlineto{\pgfqpoint{2.618822in}{1.745974in}}%
\pgfpathclose%
\pgfusepath{fill}%
\end{pgfscope}%
\begin{pgfscope}%
\pgfpathrectangle{\pgfqpoint{0.100000in}{0.100000in}}{\pgfqpoint{3.608454in}{2.310000in}}%
\pgfusepath{clip}%
\pgfsetbuttcap%
\pgfsetmiterjoin%
\definecolor{currentfill}{rgb}{0.000000,0.682353,0.658824}%
\pgfsetfillcolor{currentfill}%
\pgfsetlinewidth{0.000000pt}%
\definecolor{currentstroke}{rgb}{0.000000,0.000000,0.000000}%
\pgfsetstrokecolor{currentstroke}%
\pgfsetstrokeopacity{0.000000}%
\pgfsetdash{}{0pt}%
\pgfpathmoveto{\pgfqpoint{1.194163in}{1.451751in}}%
\pgfpathlineto{\pgfqpoint{1.192855in}{1.443180in}}%
\pgfpathlineto{\pgfqpoint{1.188609in}{1.438112in}}%
\pgfpathlineto{\pgfqpoint{1.189342in}{1.433929in}}%
\pgfpathlineto{\pgfqpoint{1.182094in}{1.423685in}}%
\pgfpathlineto{\pgfqpoint{1.177987in}{1.406555in}}%
\pgfpathlineto{\pgfqpoint{1.181137in}{1.398548in}}%
\pgfpathlineto{\pgfqpoint{1.179215in}{1.395609in}}%
\pgfpathlineto{\pgfqpoint{1.183363in}{1.379999in}}%
\pgfpathlineto{\pgfqpoint{1.181575in}{1.376785in}}%
\pgfpathlineto{\pgfqpoint{1.141855in}{1.383558in}}%
\pgfpathlineto{\pgfqpoint{1.103543in}{1.390370in}}%
\pgfpathlineto{\pgfqpoint{1.103689in}{1.391163in}}%
\pgfpathlineto{\pgfqpoint{1.111444in}{1.431696in}}%
\pgfpathlineto{\pgfqpoint{1.077995in}{1.437925in}}%
\pgfpathlineto{\pgfqpoint{1.068484in}{1.440713in}}%
\pgfpathlineto{\pgfqpoint{1.072686in}{1.463362in}}%
\pgfpathlineto{\pgfqpoint{1.078936in}{1.465774in}}%
\pgfpathlineto{\pgfqpoint{1.089487in}{1.463907in}}%
\pgfpathlineto{\pgfqpoint{1.092637in}{1.469533in}}%
\pgfpathlineto{\pgfqpoint{1.095705in}{1.487818in}}%
\pgfpathlineto{\pgfqpoint{1.104984in}{1.493265in}}%
\pgfpathlineto{\pgfqpoint{1.101981in}{1.496172in}}%
\pgfpathlineto{\pgfqpoint{1.148850in}{1.487724in}}%
\pgfpathlineto{\pgfqpoint{1.169072in}{1.483661in}}%
\pgfpathlineto{\pgfqpoint{1.206990in}{1.477188in}}%
\pgfpathlineto{\pgfqpoint{1.199560in}{1.472867in}}%
\pgfpathlineto{\pgfqpoint{1.199187in}{1.467135in}}%
\pgfpathlineto{\pgfqpoint{1.194698in}{1.460074in}}%
\pgfpathlineto{\pgfqpoint{1.194163in}{1.451751in}}%
\pgfpathclose%
\pgfusepath{fill}%
\end{pgfscope}%
\begin{pgfscope}%
\pgfpathrectangle{\pgfqpoint{0.100000in}{0.100000in}}{\pgfqpoint{3.608454in}{2.310000in}}%
\pgfusepath{clip}%
\pgfsetbuttcap%
\pgfsetmiterjoin%
\definecolor{currentfill}{rgb}{0.000000,0.945098,0.527451}%
\pgfsetfillcolor{currentfill}%
\pgfsetlinewidth{0.000000pt}%
\definecolor{currentstroke}{rgb}{0.000000,0.000000,0.000000}%
\pgfsetstrokecolor{currentstroke}%
\pgfsetstrokeopacity{0.000000}%
\pgfsetdash{}{0pt}%
\pgfpathmoveto{\pgfqpoint{0.456983in}{1.837942in}}%
\pgfpathlineto{\pgfqpoint{0.452336in}{1.831556in}}%
\pgfpathlineto{\pgfqpoint{0.442382in}{1.828434in}}%
\pgfpathlineto{\pgfqpoint{0.441986in}{1.818955in}}%
\pgfpathlineto{\pgfqpoint{0.436698in}{1.812439in}}%
\pgfpathlineto{\pgfqpoint{0.438276in}{1.805265in}}%
\pgfpathlineto{\pgfqpoint{0.437233in}{1.797086in}}%
\pgfpathlineto{\pgfqpoint{0.433774in}{1.794008in}}%
\pgfpathlineto{\pgfqpoint{0.427599in}{1.795918in}}%
\pgfpathlineto{\pgfqpoint{0.429636in}{1.802194in}}%
\pgfpathlineto{\pgfqpoint{0.413061in}{1.807713in}}%
\pgfpathlineto{\pgfqpoint{0.414756in}{1.827792in}}%
\pgfpathlineto{\pgfqpoint{0.410295in}{1.835155in}}%
\pgfpathlineto{\pgfqpoint{0.417067in}{1.845633in}}%
\pgfpathlineto{\pgfqpoint{0.418133in}{1.851097in}}%
\pgfpathlineto{\pgfqpoint{0.412925in}{1.860819in}}%
\pgfpathlineto{\pgfqpoint{0.414420in}{1.867082in}}%
\pgfpathlineto{\pgfqpoint{0.414023in}{1.879041in}}%
\pgfpathlineto{\pgfqpoint{0.418855in}{1.891440in}}%
\pgfpathlineto{\pgfqpoint{0.422736in}{1.897083in}}%
\pgfpathlineto{\pgfqpoint{0.423493in}{1.904021in}}%
\pgfpathlineto{\pgfqpoint{0.420036in}{1.911157in}}%
\pgfpathlineto{\pgfqpoint{0.420507in}{1.920108in}}%
\pgfpathlineto{\pgfqpoint{0.427334in}{1.927274in}}%
\pgfpathlineto{\pgfqpoint{0.439586in}{1.923181in}}%
\pgfpathlineto{\pgfqpoint{0.443285in}{1.914029in}}%
\pgfpathlineto{\pgfqpoint{0.438998in}{1.899672in}}%
\pgfpathlineto{\pgfqpoint{0.447212in}{1.898187in}}%
\pgfpathlineto{\pgfqpoint{0.456477in}{1.902564in}}%
\pgfpathlineto{\pgfqpoint{0.465159in}{1.900474in}}%
\pgfpathlineto{\pgfqpoint{0.455038in}{1.891185in}}%
\pgfpathlineto{\pgfqpoint{0.448521in}{1.882309in}}%
\pgfpathlineto{\pgfqpoint{0.441915in}{1.883291in}}%
\pgfpathlineto{\pgfqpoint{0.437497in}{1.874169in}}%
\pgfpathlineto{\pgfqpoint{0.444876in}{1.871177in}}%
\pgfpathlineto{\pgfqpoint{0.447509in}{1.860643in}}%
\pgfpathlineto{\pgfqpoint{0.441504in}{1.854677in}}%
\pgfpathlineto{\pgfqpoint{0.439878in}{1.843077in}}%
\pgfpathlineto{\pgfqpoint{0.456983in}{1.837942in}}%
\pgfpathclose%
\pgfusepath{fill}%
\end{pgfscope}%
\begin{pgfscope}%
\pgfpathrectangle{\pgfqpoint{0.100000in}{0.100000in}}{\pgfqpoint{3.608454in}{2.310000in}}%
\pgfusepath{clip}%
\pgfsetbuttcap%
\pgfsetmiterjoin%
\definecolor{currentfill}{rgb}{0.000000,0.639216,0.680392}%
\pgfsetfillcolor{currentfill}%
\pgfsetlinewidth{0.000000pt}%
\definecolor{currentstroke}{rgb}{0.000000,0.000000,0.000000}%
\pgfsetstrokecolor{currentstroke}%
\pgfsetstrokeopacity{0.000000}%
\pgfsetdash{}{0pt}%
\pgfpathmoveto{\pgfqpoint{2.720165in}{1.477649in}}%
\pgfpathlineto{\pgfqpoint{2.722450in}{1.453627in}}%
\pgfpathlineto{\pgfqpoint{2.698943in}{1.451000in}}%
\pgfpathlineto{\pgfqpoint{2.699598in}{1.433892in}}%
\pgfpathlineto{\pgfqpoint{2.694747in}{1.433296in}}%
\pgfpathlineto{\pgfqpoint{2.676915in}{1.431338in}}%
\pgfpathlineto{\pgfqpoint{2.676851in}{1.443992in}}%
\pgfpathlineto{\pgfqpoint{2.659619in}{1.442049in}}%
\pgfpathlineto{\pgfqpoint{2.657467in}{1.463709in}}%
\pgfpathlineto{\pgfqpoint{2.653860in}{1.497742in}}%
\pgfpathlineto{\pgfqpoint{2.667350in}{1.499209in}}%
\pgfpathlineto{\pgfqpoint{2.678668in}{1.500441in}}%
\pgfpathlineto{\pgfqpoint{2.679497in}{1.493619in}}%
\pgfpathlineto{\pgfqpoint{2.694173in}{1.495218in}}%
\pgfpathlineto{\pgfqpoint{2.695336in}{1.474638in}}%
\pgfpathlineto{\pgfqpoint{2.720165in}{1.477649in}}%
\pgfpathclose%
\pgfusepath{fill}%
\end{pgfscope}%
\begin{pgfscope}%
\pgfpathrectangle{\pgfqpoint{0.100000in}{0.100000in}}{\pgfqpoint{3.608454in}{2.310000in}}%
\pgfusepath{clip}%
\pgfsetbuttcap%
\pgfsetmiterjoin%
\definecolor{currentfill}{rgb}{0.000000,0.635294,0.682353}%
\pgfsetfillcolor{currentfill}%
\pgfsetlinewidth{0.000000pt}%
\definecolor{currentstroke}{rgb}{0.000000,0.000000,0.000000}%
\pgfsetstrokecolor{currentstroke}%
\pgfsetstrokeopacity{0.000000}%
\pgfsetdash{}{0pt}%
\pgfpathmoveto{\pgfqpoint{2.809988in}{1.252552in}}%
\pgfpathlineto{\pgfqpoint{2.801049in}{1.243631in}}%
\pgfpathlineto{\pgfqpoint{2.797121in}{1.243378in}}%
\pgfpathlineto{\pgfqpoint{2.788710in}{1.248796in}}%
\pgfpathlineto{\pgfqpoint{2.781851in}{1.240501in}}%
\pgfpathlineto{\pgfqpoint{2.780361in}{1.233265in}}%
\pgfpathlineto{\pgfqpoint{2.765354in}{1.245580in}}%
\pgfpathlineto{\pgfqpoint{2.767417in}{1.258344in}}%
\pgfpathlineto{\pgfqpoint{2.773095in}{1.263245in}}%
\pgfpathlineto{\pgfqpoint{2.759930in}{1.280063in}}%
\pgfpathlineto{\pgfqpoint{2.774318in}{1.293523in}}%
\pgfpathlineto{\pgfqpoint{2.791570in}{1.288910in}}%
\pgfpathlineto{\pgfqpoint{2.801059in}{1.283583in}}%
\pgfpathlineto{\pgfqpoint{2.810307in}{1.281879in}}%
\pgfpathlineto{\pgfqpoint{2.804559in}{1.276722in}}%
\pgfpathlineto{\pgfqpoint{2.800586in}{1.266861in}}%
\pgfpathlineto{\pgfqpoint{2.802573in}{1.261875in}}%
\pgfpathlineto{\pgfqpoint{2.809988in}{1.252552in}}%
\pgfpathclose%
\pgfusepath{fill}%
\end{pgfscope}%
\begin{pgfscope}%
\pgfpathrectangle{\pgfqpoint{0.100000in}{0.100000in}}{\pgfqpoint{3.608454in}{2.310000in}}%
\pgfusepath{clip}%
\pgfsetbuttcap%
\pgfsetmiterjoin%
\definecolor{currentfill}{rgb}{0.000000,0.964706,0.517647}%
\pgfsetfillcolor{currentfill}%
\pgfsetlinewidth{0.000000pt}%
\definecolor{currentstroke}{rgb}{0.000000,0.000000,0.000000}%
\pgfsetstrokecolor{currentstroke}%
\pgfsetstrokeopacity{0.000000}%
\pgfsetdash{}{0pt}%
\pgfpathmoveto{\pgfqpoint{2.182637in}{0.773027in}}%
\pgfpathlineto{\pgfqpoint{2.212262in}{0.773539in}}%
\pgfpathlineto{\pgfqpoint{2.212402in}{0.763123in}}%
\pgfpathlineto{\pgfqpoint{2.219367in}{0.763229in}}%
\pgfpathlineto{\pgfqpoint{2.219726in}{0.749195in}}%
\pgfpathlineto{\pgfqpoint{2.226572in}{0.749343in}}%
\pgfpathlineto{\pgfqpoint{2.226879in}{0.735487in}}%
\pgfpathlineto{\pgfqpoint{2.216369in}{0.735263in}}%
\pgfpathlineto{\pgfqpoint{2.212945in}{0.728230in}}%
\pgfpathlineto{\pgfqpoint{2.206617in}{0.728100in}}%
\pgfpathlineto{\pgfqpoint{2.205558in}{0.720837in}}%
\pgfpathlineto{\pgfqpoint{2.202317in}{0.720035in}}%
\pgfpathlineto{\pgfqpoint{2.196211in}{0.730685in}}%
\pgfpathlineto{\pgfqpoint{2.197023in}{0.737625in}}%
\pgfpathlineto{\pgfqpoint{2.191816in}{0.743213in}}%
\pgfpathlineto{\pgfqpoint{2.185571in}{0.753524in}}%
\pgfpathlineto{\pgfqpoint{2.187852in}{0.760329in}}%
\pgfpathlineto{\pgfqpoint{2.186530in}{0.767441in}}%
\pgfpathlineto{\pgfqpoint{2.182637in}{0.773027in}}%
\pgfpathclose%
\pgfusepath{fill}%
\end{pgfscope}%
\begin{pgfscope}%
\pgfpathrectangle{\pgfqpoint{0.100000in}{0.100000in}}{\pgfqpoint{3.608454in}{2.310000in}}%
\pgfusepath{clip}%
\pgfsetbuttcap%
\pgfsetmiterjoin%
\definecolor{currentfill}{rgb}{0.000000,0.486275,0.756863}%
\pgfsetfillcolor{currentfill}%
\pgfsetlinewidth{0.000000pt}%
\definecolor{currentstroke}{rgb}{0.000000,0.000000,0.000000}%
\pgfsetstrokecolor{currentstroke}%
\pgfsetstrokeopacity{0.000000}%
\pgfsetdash{}{0pt}%
\pgfpathmoveto{\pgfqpoint{1.732544in}{1.486727in}}%
\pgfpathlineto{\pgfqpoint{1.684512in}{1.490228in}}%
\pgfpathlineto{\pgfqpoint{1.684833in}{1.494331in}}%
\pgfpathlineto{\pgfqpoint{1.649013in}{1.497206in}}%
\pgfpathlineto{\pgfqpoint{1.650561in}{1.517162in}}%
\pgfpathlineto{\pgfqpoint{1.653438in}{1.520561in}}%
\pgfpathlineto{\pgfqpoint{1.654575in}{1.534236in}}%
\pgfpathlineto{\pgfqpoint{1.687474in}{1.531563in}}%
\pgfpathlineto{\pgfqpoint{1.688556in}{1.545253in}}%
\pgfpathlineto{\pgfqpoint{1.692731in}{1.544894in}}%
\pgfpathlineto{\pgfqpoint{1.726904in}{1.542427in}}%
\pgfpathlineto{\pgfqpoint{1.735045in}{1.541820in}}%
\pgfpathlineto{\pgfqpoint{1.733109in}{1.514401in}}%
\pgfpathlineto{\pgfqpoint{1.734320in}{1.514317in}}%
\pgfpathlineto{\pgfqpoint{1.732544in}{1.486727in}}%
\pgfpathclose%
\pgfusepath{fill}%
\end{pgfscope}%
\begin{pgfscope}%
\pgfpathrectangle{\pgfqpoint{0.100000in}{0.100000in}}{\pgfqpoint{3.608454in}{2.310000in}}%
\pgfusepath{clip}%
\pgfsetbuttcap%
\pgfsetmiterjoin%
\definecolor{currentfill}{rgb}{0.000000,0.596078,0.701961}%
\pgfsetfillcolor{currentfill}%
\pgfsetlinewidth{0.000000pt}%
\definecolor{currentstroke}{rgb}{0.000000,0.000000,0.000000}%
\pgfsetstrokecolor{currentstroke}%
\pgfsetstrokeopacity{0.000000}%
\pgfsetdash{}{0pt}%
\pgfpathmoveto{\pgfqpoint{2.610086in}{1.175304in}}%
\pgfpathlineto{\pgfqpoint{2.607679in}{1.208265in}}%
\pgfpathlineto{\pgfqpoint{2.620792in}{1.209072in}}%
\pgfpathlineto{\pgfqpoint{2.631929in}{1.205576in}}%
\pgfpathlineto{\pgfqpoint{2.634641in}{1.217991in}}%
\pgfpathlineto{\pgfqpoint{2.639824in}{1.220359in}}%
\pgfpathlineto{\pgfqpoint{2.647960in}{1.220595in}}%
\pgfpathlineto{\pgfqpoint{2.642473in}{1.232080in}}%
\pgfpathlineto{\pgfqpoint{2.661850in}{1.235103in}}%
\pgfpathlineto{\pgfqpoint{2.669571in}{1.226318in}}%
\pgfpathlineto{\pgfqpoint{2.669252in}{1.213390in}}%
\pgfpathlineto{\pgfqpoint{2.664358in}{1.203424in}}%
\pgfpathlineto{\pgfqpoint{2.677443in}{1.189404in}}%
\pgfpathlineto{\pgfqpoint{2.678723in}{1.180579in}}%
\pgfpathlineto{\pgfqpoint{2.656122in}{1.179965in}}%
\pgfpathlineto{\pgfqpoint{2.639785in}{1.178722in}}%
\pgfpathlineto{\pgfqpoint{2.610086in}{1.175304in}}%
\pgfpathclose%
\pgfusepath{fill}%
\end{pgfscope}%
\begin{pgfscope}%
\pgfpathrectangle{\pgfqpoint{0.100000in}{0.100000in}}{\pgfqpoint{3.608454in}{2.310000in}}%
\pgfusepath{clip}%
\pgfsetbuttcap%
\pgfsetmiterjoin%
\definecolor{currentfill}{rgb}{0.000000,0.564706,0.717647}%
\pgfsetfillcolor{currentfill}%
\pgfsetlinewidth{0.000000pt}%
\definecolor{currentstroke}{rgb}{0.000000,0.000000,0.000000}%
\pgfsetstrokecolor{currentstroke}%
\pgfsetstrokeopacity{0.000000}%
\pgfsetdash{}{0pt}%
\pgfpathmoveto{\pgfqpoint{0.975129in}{1.855275in}}%
\pgfpathlineto{\pgfqpoint{0.981526in}{1.858299in}}%
\pgfpathlineto{\pgfqpoint{0.983732in}{1.852523in}}%
\pgfpathlineto{\pgfqpoint{0.991130in}{1.850884in}}%
\pgfpathlineto{\pgfqpoint{0.999848in}{1.846301in}}%
\pgfpathlineto{\pgfqpoint{1.012788in}{1.842181in}}%
\pgfpathlineto{\pgfqpoint{1.013709in}{1.836495in}}%
\pgfpathlineto{\pgfqpoint{1.027024in}{1.825238in}}%
\pgfpathlineto{\pgfqpoint{1.034613in}{1.812423in}}%
\pgfpathlineto{\pgfqpoint{1.038897in}{1.813232in}}%
\pgfpathlineto{\pgfqpoint{1.041894in}{1.801489in}}%
\pgfpathlineto{\pgfqpoint{1.040506in}{1.794892in}}%
\pgfpathlineto{\pgfqpoint{1.055937in}{1.791680in}}%
\pgfpathlineto{\pgfqpoint{1.054623in}{1.785268in}}%
\pgfpathlineto{\pgfqpoint{1.074676in}{1.781213in}}%
\pgfpathlineto{\pgfqpoint{1.071945in}{1.767715in}}%
\pgfpathlineto{\pgfqpoint{1.058573in}{1.770425in}}%
\pgfpathlineto{\pgfqpoint{1.053439in}{1.743275in}}%
\pgfpathlineto{\pgfqpoint{1.055372in}{1.735866in}}%
\pgfpathlineto{\pgfqpoint{1.051274in}{1.732636in}}%
\pgfpathlineto{\pgfqpoint{1.044474in}{1.739072in}}%
\pgfpathlineto{\pgfqpoint{1.038549in}{1.738796in}}%
\pgfpathlineto{\pgfqpoint{1.041428in}{1.752757in}}%
\pgfpathlineto{\pgfqpoint{1.044792in}{1.752067in}}%
\pgfpathlineto{\pgfqpoint{1.050343in}{1.779280in}}%
\pgfpathlineto{\pgfqpoint{1.033321in}{1.782826in}}%
\pgfpathlineto{\pgfqpoint{1.027689in}{1.755650in}}%
\pgfpathlineto{\pgfqpoint{1.023421in}{1.749624in}}%
\pgfpathlineto{\pgfqpoint{1.013857in}{1.751700in}}%
\pgfpathlineto{\pgfqpoint{1.010156in}{1.733884in}}%
\pgfpathlineto{\pgfqpoint{1.002092in}{1.733915in}}%
\pgfpathlineto{\pgfqpoint{1.000366in}{1.726400in}}%
\pgfpathlineto{\pgfqpoint{0.988011in}{1.729060in}}%
\pgfpathlineto{\pgfqpoint{0.981161in}{1.696275in}}%
\pgfpathlineto{\pgfqpoint{0.949739in}{1.703791in}}%
\pgfpathlineto{\pgfqpoint{0.937739in}{1.705964in}}%
\pgfpathlineto{\pgfqpoint{0.953662in}{1.776726in}}%
\pgfpathlineto{\pgfqpoint{0.950896in}{1.777540in}}%
\pgfpathlineto{\pgfqpoint{0.963091in}{1.830834in}}%
\pgfpathlineto{\pgfqpoint{0.968859in}{1.837111in}}%
\pgfpathlineto{\pgfqpoint{0.967810in}{1.841171in}}%
\pgfpathlineto{\pgfqpoint{0.973897in}{1.850229in}}%
\pgfpathlineto{\pgfqpoint{0.975129in}{1.855275in}}%
\pgfpathclose%
\pgfusepath{fill}%
\end{pgfscope}%
\begin{pgfscope}%
\pgfpathrectangle{\pgfqpoint{0.100000in}{0.100000in}}{\pgfqpoint{3.608454in}{2.310000in}}%
\pgfusepath{clip}%
\pgfsetbuttcap%
\pgfsetmiterjoin%
\definecolor{currentfill}{rgb}{0.000000,0.909804,0.545098}%
\pgfsetfillcolor{currentfill}%
\pgfsetlinewidth{0.000000pt}%
\definecolor{currentstroke}{rgb}{0.000000,0.000000,0.000000}%
\pgfsetstrokecolor{currentstroke}%
\pgfsetstrokeopacity{0.000000}%
\pgfsetdash{}{0pt}%
\pgfpathmoveto{\pgfqpoint{0.494828in}{1.757961in}}%
\pgfpathlineto{\pgfqpoint{0.489189in}{1.745779in}}%
\pgfpathlineto{\pgfqpoint{0.485004in}{1.744684in}}%
\pgfpathlineto{\pgfqpoint{0.475334in}{1.730099in}}%
\pgfpathlineto{\pgfqpoint{0.468741in}{1.723271in}}%
\pgfpathlineto{\pgfqpoint{0.469494in}{1.715995in}}%
\pgfpathlineto{\pgfqpoint{0.458930in}{1.713085in}}%
\pgfpathlineto{\pgfqpoint{0.454296in}{1.709584in}}%
\pgfpathlineto{\pgfqpoint{0.448092in}{1.709303in}}%
\pgfpathlineto{\pgfqpoint{0.442506in}{1.705022in}}%
\pgfpathlineto{\pgfqpoint{0.441132in}{1.699873in}}%
\pgfpathlineto{\pgfqpoint{0.444780in}{1.696306in}}%
\pgfpathlineto{\pgfqpoint{0.442069in}{1.688570in}}%
\pgfpathlineto{\pgfqpoint{0.441276in}{1.674970in}}%
\pgfpathlineto{\pgfqpoint{0.406038in}{1.686194in}}%
\pgfpathlineto{\pgfqpoint{0.406652in}{1.688065in}}%
\pgfpathlineto{\pgfqpoint{0.379033in}{1.697172in}}%
\pgfpathlineto{\pgfqpoint{0.376641in}{1.705935in}}%
\pgfpathlineto{\pgfqpoint{0.372917in}{1.709765in}}%
\pgfpathlineto{\pgfqpoint{0.365968in}{1.722668in}}%
\pgfpathlineto{\pgfqpoint{0.368189in}{1.726700in}}%
\pgfpathlineto{\pgfqpoint{0.367902in}{1.737177in}}%
\pgfpathlineto{\pgfqpoint{0.378459in}{1.750564in}}%
\pgfpathlineto{\pgfqpoint{0.394964in}{1.767791in}}%
\pgfpathlineto{\pgfqpoint{0.399441in}{1.775849in}}%
\pgfpathlineto{\pgfqpoint{0.397857in}{1.779799in}}%
\pgfpathlineto{\pgfqpoint{0.406645in}{1.794088in}}%
\pgfpathlineto{\pgfqpoint{0.413061in}{1.807713in}}%
\pgfpathlineto{\pgfqpoint{0.429636in}{1.802194in}}%
\pgfpathlineto{\pgfqpoint{0.427599in}{1.795918in}}%
\pgfpathlineto{\pgfqpoint{0.433774in}{1.794008in}}%
\pgfpathlineto{\pgfqpoint{0.442827in}{1.791103in}}%
\pgfpathlineto{\pgfqpoint{0.443568in}{1.772460in}}%
\pgfpathlineto{\pgfqpoint{0.449403in}{1.767395in}}%
\pgfpathlineto{\pgfqpoint{0.449915in}{1.763500in}}%
\pgfpathlineto{\pgfqpoint{0.457583in}{1.760831in}}%
\pgfpathlineto{\pgfqpoint{0.459959in}{1.754157in}}%
\pgfpathlineto{\pgfqpoint{0.468505in}{1.753515in}}%
\pgfpathlineto{\pgfqpoint{0.466061in}{1.760966in}}%
\pgfpathlineto{\pgfqpoint{0.469737in}{1.765482in}}%
\pgfpathlineto{\pgfqpoint{0.477526in}{1.765034in}}%
\pgfpathlineto{\pgfqpoint{0.494062in}{1.772814in}}%
\pgfpathlineto{\pgfqpoint{0.498294in}{1.770162in}}%
\pgfpathlineto{\pgfqpoint{0.494828in}{1.757961in}}%
\pgfpathclose%
\pgfusepath{fill}%
\end{pgfscope}%
\begin{pgfscope}%
\pgfpathrectangle{\pgfqpoint{0.100000in}{0.100000in}}{\pgfqpoint{3.608454in}{2.310000in}}%
\pgfusepath{clip}%
\pgfsetbuttcap%
\pgfsetmiterjoin%
\definecolor{currentfill}{rgb}{0.000000,0.603922,0.698039}%
\pgfsetfillcolor{currentfill}%
\pgfsetlinewidth{0.000000pt}%
\definecolor{currentstroke}{rgb}{0.000000,0.000000,0.000000}%
\pgfsetstrokecolor{currentstroke}%
\pgfsetstrokeopacity{0.000000}%
\pgfsetdash{}{0pt}%
\pgfpathmoveto{\pgfqpoint{2.246959in}{1.366834in}}%
\pgfpathlineto{\pgfqpoint{2.241843in}{1.368445in}}%
\pgfpathlineto{\pgfqpoint{2.238456in}{1.359055in}}%
\pgfpathlineto{\pgfqpoint{2.233381in}{1.359171in}}%
\pgfpathlineto{\pgfqpoint{2.230926in}{1.365957in}}%
\pgfpathlineto{\pgfqpoint{2.224430in}{1.367560in}}%
\pgfpathlineto{\pgfqpoint{2.221539in}{1.373866in}}%
\pgfpathlineto{\pgfqpoint{2.215503in}{1.374729in}}%
\pgfpathlineto{\pgfqpoint{2.211917in}{1.379066in}}%
\pgfpathlineto{\pgfqpoint{2.212943in}{1.384548in}}%
\pgfpathlineto{\pgfqpoint{2.211089in}{1.394791in}}%
\pgfpathlineto{\pgfqpoint{2.206226in}{1.396171in}}%
\pgfpathlineto{\pgfqpoint{2.206018in}{1.417015in}}%
\pgfpathlineto{\pgfqpoint{2.205934in}{1.422201in}}%
\pgfpathlineto{\pgfqpoint{2.236527in}{1.423164in}}%
\pgfpathlineto{\pgfqpoint{2.237055in}{1.423181in}}%
\pgfpathlineto{\pgfqpoint{2.237052in}{1.396510in}}%
\pgfpathlineto{\pgfqpoint{2.247305in}{1.396610in}}%
\pgfpathlineto{\pgfqpoint{2.246959in}{1.366834in}}%
\pgfpathclose%
\pgfusepath{fill}%
\end{pgfscope}%
\begin{pgfscope}%
\pgfpathrectangle{\pgfqpoint{0.100000in}{0.100000in}}{\pgfqpoint{3.608454in}{2.310000in}}%
\pgfusepath{clip}%
\pgfsetbuttcap%
\pgfsetmiterjoin%
\definecolor{currentfill}{rgb}{0.000000,0.572549,0.713725}%
\pgfsetfillcolor{currentfill}%
\pgfsetlinewidth{0.000000pt}%
\definecolor{currentstroke}{rgb}{0.000000,0.000000,0.000000}%
\pgfsetstrokecolor{currentstroke}%
\pgfsetstrokeopacity{0.000000}%
\pgfsetdash{}{0pt}%
\pgfpathmoveto{\pgfqpoint{1.857172in}{0.996308in}}%
\pgfpathlineto{\pgfqpoint{1.856890in}{0.989419in}}%
\pgfpathlineto{\pgfqpoint{1.867617in}{0.988984in}}%
\pgfpathlineto{\pgfqpoint{1.867064in}{0.975207in}}%
\pgfpathlineto{\pgfqpoint{1.870481in}{0.975087in}}%
\pgfpathlineto{\pgfqpoint{1.869882in}{0.961093in}}%
\pgfpathlineto{\pgfqpoint{1.864563in}{0.959406in}}%
\pgfpathlineto{\pgfqpoint{1.852516in}{0.961694in}}%
\pgfpathlineto{\pgfqpoint{1.847621in}{0.966431in}}%
\pgfpathlineto{\pgfqpoint{1.832097in}{0.967259in}}%
\pgfpathlineto{\pgfqpoint{1.831143in}{0.977031in}}%
\pgfpathlineto{\pgfqpoint{1.831022in}{0.989497in}}%
\pgfpathlineto{\pgfqpoint{1.835171in}{0.996785in}}%
\pgfpathlineto{\pgfqpoint{1.857172in}{0.996308in}}%
\pgfpathclose%
\pgfusepath{fill}%
\end{pgfscope}%
\begin{pgfscope}%
\pgfpathrectangle{\pgfqpoint{0.100000in}{0.100000in}}{\pgfqpoint{3.608454in}{2.310000in}}%
\pgfusepath{clip}%
\pgfsetbuttcap%
\pgfsetmiterjoin%
\definecolor{currentfill}{rgb}{0.000000,0.533333,0.733333}%
\pgfsetfillcolor{currentfill}%
\pgfsetlinewidth{0.000000pt}%
\definecolor{currentstroke}{rgb}{0.000000,0.000000,0.000000}%
\pgfsetstrokecolor{currentstroke}%
\pgfsetstrokeopacity{0.000000}%
\pgfsetdash{}{0pt}%
\pgfpathmoveto{\pgfqpoint{2.046244in}{1.696054in}}%
\pgfpathlineto{\pgfqpoint{2.023239in}{1.696380in}}%
\pgfpathlineto{\pgfqpoint{2.023670in}{1.724067in}}%
\pgfpathlineto{\pgfqpoint{2.024621in}{1.779133in}}%
\pgfpathlineto{\pgfqpoint{2.024747in}{1.786093in}}%
\pgfpathlineto{\pgfqpoint{2.072659in}{1.785449in}}%
\pgfpathlineto{\pgfqpoint{2.073124in}{1.778481in}}%
\pgfpathlineto{\pgfqpoint{2.073025in}{1.750974in}}%
\pgfpathlineto{\pgfqpoint{2.046268in}{1.751283in}}%
\pgfpathlineto{\pgfqpoint{2.046612in}{1.723715in}}%
\pgfpathlineto{\pgfqpoint{2.046244in}{1.696054in}}%
\pgfpathclose%
\pgfusepath{fill}%
\end{pgfscope}%
\begin{pgfscope}%
\pgfpathrectangle{\pgfqpoint{0.100000in}{0.100000in}}{\pgfqpoint{3.608454in}{2.310000in}}%
\pgfusepath{clip}%
\pgfsetbuttcap%
\pgfsetmiterjoin%
\definecolor{currentfill}{rgb}{0.000000,0.490196,0.754902}%
\pgfsetfillcolor{currentfill}%
\pgfsetlinewidth{0.000000pt}%
\definecolor{currentstroke}{rgb}{0.000000,0.000000,0.000000}%
\pgfsetstrokecolor{currentstroke}%
\pgfsetstrokeopacity{0.000000}%
\pgfsetdash{}{0pt}%
\pgfpathmoveto{\pgfqpoint{1.461309in}{1.650537in}}%
\pgfpathlineto{\pgfqpoint{1.457886in}{1.623314in}}%
\pgfpathlineto{\pgfqpoint{1.451617in}{1.568947in}}%
\pgfpathlineto{\pgfqpoint{1.436631in}{1.570812in}}%
\pgfpathlineto{\pgfqpoint{1.432841in}{1.539719in}}%
\pgfpathlineto{\pgfqpoint{1.401111in}{1.544070in}}%
\pgfpathlineto{\pgfqpoint{1.338390in}{1.552623in}}%
\pgfpathlineto{\pgfqpoint{1.342910in}{1.583654in}}%
\pgfpathlineto{\pgfqpoint{1.345097in}{1.604208in}}%
\pgfpathlineto{\pgfqpoint{1.369658in}{1.600576in}}%
\pgfpathlineto{\pgfqpoint{1.374676in}{1.634274in}}%
\pgfpathlineto{\pgfqpoint{1.375467in}{1.648053in}}%
\pgfpathlineto{\pgfqpoint{1.377332in}{1.661590in}}%
\pgfpathlineto{\pgfqpoint{1.429302in}{1.654336in}}%
\pgfpathlineto{\pgfqpoint{1.461309in}{1.650537in}}%
\pgfpathclose%
\pgfusepath{fill}%
\end{pgfscope}%
\begin{pgfscope}%
\pgfpathrectangle{\pgfqpoint{0.100000in}{0.100000in}}{\pgfqpoint{3.608454in}{2.310000in}}%
\pgfusepath{clip}%
\pgfsetbuttcap%
\pgfsetmiterjoin%
\definecolor{currentfill}{rgb}{0.000000,0.301961,0.849020}%
\pgfsetfillcolor{currentfill}%
\pgfsetlinewidth{0.000000pt}%
\definecolor{currentstroke}{rgb}{0.000000,0.000000,0.000000}%
\pgfsetstrokecolor{currentstroke}%
\pgfsetstrokeopacity{0.000000}%
\pgfsetdash{}{0pt}%
\pgfpathmoveto{\pgfqpoint{1.963356in}{1.420723in}}%
\pgfpathlineto{\pgfqpoint{1.964117in}{1.448246in}}%
\pgfpathlineto{\pgfqpoint{1.965564in}{1.503301in}}%
\pgfpathlineto{\pgfqpoint{1.992970in}{1.502588in}}%
\pgfpathlineto{\pgfqpoint{2.019500in}{1.502021in}}%
\pgfpathlineto{\pgfqpoint{2.019155in}{1.481378in}}%
\pgfpathlineto{\pgfqpoint{2.018170in}{1.419481in}}%
\pgfpathlineto{\pgfqpoint{1.963356in}{1.420723in}}%
\pgfpathclose%
\pgfusepath{fill}%
\end{pgfscope}%
\begin{pgfscope}%
\pgfpathrectangle{\pgfqpoint{0.100000in}{0.100000in}}{\pgfqpoint{3.608454in}{2.310000in}}%
\pgfusepath{clip}%
\pgfsetbuttcap%
\pgfsetmiterjoin%
\definecolor{currentfill}{rgb}{0.000000,0.458824,0.770588}%
\pgfsetfillcolor{currentfill}%
\pgfsetlinewidth{0.000000pt}%
\definecolor{currentstroke}{rgb}{0.000000,0.000000,0.000000}%
\pgfsetstrokecolor{currentstroke}%
\pgfsetstrokeopacity{0.000000}%
\pgfsetdash{}{0pt}%
\pgfpathmoveto{\pgfqpoint{1.960057in}{1.995378in}}%
\pgfpathlineto{\pgfqpoint{1.946235in}{1.995849in}}%
\pgfpathlineto{\pgfqpoint{1.945081in}{2.002729in}}%
\pgfpathlineto{\pgfqpoint{1.946085in}{2.030374in}}%
\pgfpathlineto{\pgfqpoint{1.951584in}{2.030204in}}%
\pgfpathlineto{\pgfqpoint{1.952493in}{2.057946in}}%
\pgfpathlineto{\pgfqpoint{1.951799in}{2.072044in}}%
\pgfpathlineto{\pgfqpoint{1.992478in}{2.070813in}}%
\pgfpathlineto{\pgfqpoint{1.992147in}{2.069188in}}%
\pgfpathlineto{\pgfqpoint{2.026639in}{2.068568in}}%
\pgfpathlineto{\pgfqpoint{2.027218in}{2.042403in}}%
\pgfpathlineto{\pgfqpoint{2.034152in}{2.042276in}}%
\pgfpathlineto{\pgfqpoint{2.034103in}{2.035335in}}%
\pgfpathlineto{\pgfqpoint{2.061781in}{2.035022in}}%
\pgfpathlineto{\pgfqpoint{2.061858in}{2.042025in}}%
\pgfpathlineto{\pgfqpoint{2.068716in}{2.041994in}}%
\pgfpathlineto{\pgfqpoint{2.068749in}{2.048976in}}%
\pgfpathlineto{\pgfqpoint{2.075577in}{2.048780in}}%
\pgfpathlineto{\pgfqpoint{2.075622in}{2.027945in}}%
\pgfpathlineto{\pgfqpoint{2.076864in}{2.014087in}}%
\pgfpathlineto{\pgfqpoint{2.049145in}{2.014243in}}%
\pgfpathlineto{\pgfqpoint{2.006954in}{2.014902in}}%
\pgfpathlineto{\pgfqpoint{2.007393in}{1.994001in}}%
\pgfpathlineto{\pgfqpoint{1.960057in}{1.995378in}}%
\pgfpathclose%
\pgfusepath{fill}%
\end{pgfscope}%
\begin{pgfscope}%
\pgfpathrectangle{\pgfqpoint{0.100000in}{0.100000in}}{\pgfqpoint{3.608454in}{2.310000in}}%
\pgfusepath{clip}%
\pgfsetbuttcap%
\pgfsetmiterjoin%
\definecolor{currentfill}{rgb}{0.000000,0.454902,0.772549}%
\pgfsetfillcolor{currentfill}%
\pgfsetlinewidth{0.000000pt}%
\definecolor{currentstroke}{rgb}{0.000000,0.000000,0.000000}%
\pgfsetstrokecolor{currentstroke}%
\pgfsetstrokeopacity{0.000000}%
\pgfsetdash{}{0pt}%
\pgfpathmoveto{\pgfqpoint{2.100603in}{1.750815in}}%
\pgfpathlineto{\pgfqpoint{2.073025in}{1.750974in}}%
\pgfpathlineto{\pgfqpoint{2.073124in}{1.778481in}}%
\pgfpathlineto{\pgfqpoint{2.086396in}{1.778361in}}%
\pgfpathlineto{\pgfqpoint{2.086458in}{1.790823in}}%
\pgfpathlineto{\pgfqpoint{2.079554in}{1.795195in}}%
\pgfpathlineto{\pgfqpoint{2.079580in}{1.806133in}}%
\pgfpathlineto{\pgfqpoint{2.120405in}{1.806180in}}%
\pgfpathlineto{\pgfqpoint{2.135067in}{1.806256in}}%
\pgfpathlineto{\pgfqpoint{2.135124in}{1.792351in}}%
\pgfpathlineto{\pgfqpoint{2.127691in}{1.792284in}}%
\pgfpathlineto{\pgfqpoint{2.128088in}{1.771585in}}%
\pgfpathlineto{\pgfqpoint{2.118639in}{1.771585in}}%
\pgfpathlineto{\pgfqpoint{2.114392in}{1.774876in}}%
\pgfpathlineto{\pgfqpoint{2.114393in}{1.757738in}}%
\pgfpathlineto{\pgfqpoint{2.100687in}{1.757698in}}%
\pgfpathlineto{\pgfqpoint{2.100603in}{1.750815in}}%
\pgfpathclose%
\pgfusepath{fill}%
\end{pgfscope}%
\begin{pgfscope}%
\pgfpathrectangle{\pgfqpoint{0.100000in}{0.100000in}}{\pgfqpoint{3.608454in}{2.310000in}}%
\pgfusepath{clip}%
\pgfsetbuttcap%
\pgfsetmiterjoin%
\definecolor{currentfill}{rgb}{0.000000,0.796078,0.601961}%
\pgfsetfillcolor{currentfill}%
\pgfsetlinewidth{0.000000pt}%
\definecolor{currentstroke}{rgb}{0.000000,0.000000,0.000000}%
\pgfsetstrokecolor{currentstroke}%
\pgfsetstrokeopacity{0.000000}%
\pgfsetdash{}{0pt}%
\pgfpathmoveto{\pgfqpoint{2.318616in}{2.073687in}}%
\pgfpathlineto{\pgfqpoint{2.328248in}{2.078608in}}%
\pgfpathlineto{\pgfqpoint{2.333925in}{2.066565in}}%
\pgfpathlineto{\pgfqpoint{2.354068in}{2.067548in}}%
\pgfpathlineto{\pgfqpoint{2.366718in}{2.069902in}}%
\pgfpathlineto{\pgfqpoint{2.372839in}{2.068040in}}%
\pgfpathlineto{\pgfqpoint{2.374681in}{2.063610in}}%
\pgfpathlineto{\pgfqpoint{2.381315in}{2.060550in}}%
\pgfpathlineto{\pgfqpoint{2.386736in}{2.063805in}}%
\pgfpathlineto{\pgfqpoint{2.394855in}{2.063381in}}%
\pgfpathlineto{\pgfqpoint{2.394022in}{2.058782in}}%
\pgfpathlineto{\pgfqpoint{2.376788in}{2.048043in}}%
\pgfpathlineto{\pgfqpoint{2.365376in}{2.042875in}}%
\pgfpathlineto{\pgfqpoint{2.346958in}{2.036101in}}%
\pgfpathlineto{\pgfqpoint{2.336678in}{2.029365in}}%
\pgfpathlineto{\pgfqpoint{2.321610in}{2.015873in}}%
\pgfpathlineto{\pgfqpoint{2.320258in}{2.049681in}}%
\pgfpathlineto{\pgfqpoint{2.318616in}{2.073687in}}%
\pgfpathclose%
\pgfusepath{fill}%
\end{pgfscope}%
\begin{pgfscope}%
\pgfpathrectangle{\pgfqpoint{0.100000in}{0.100000in}}{\pgfqpoint{3.608454in}{2.310000in}}%
\pgfusepath{clip}%
\pgfsetbuttcap%
\pgfsetmiterjoin%
\definecolor{currentfill}{rgb}{0.000000,0.737255,0.631373}%
\pgfsetfillcolor{currentfill}%
\pgfsetlinewidth{0.000000pt}%
\definecolor{currentstroke}{rgb}{0.000000,0.000000,0.000000}%
\pgfsetstrokecolor{currentstroke}%
\pgfsetstrokeopacity{0.000000}%
\pgfsetdash{}{0pt}%
\pgfpathmoveto{\pgfqpoint{2.842704in}{1.366486in}}%
\pgfpathlineto{\pgfqpoint{2.836077in}{1.357238in}}%
\pgfpathlineto{\pgfqpoint{2.829317in}{1.356171in}}%
\pgfpathlineto{\pgfqpoint{2.825080in}{1.387318in}}%
\pgfpathlineto{\pgfqpoint{2.817692in}{1.389496in}}%
\pgfpathlineto{\pgfqpoint{2.818458in}{1.403774in}}%
\pgfpathlineto{\pgfqpoint{2.816007in}{1.404773in}}%
\pgfpathlineto{\pgfqpoint{2.815226in}{1.414852in}}%
\pgfpathlineto{\pgfqpoint{2.820238in}{1.426570in}}%
\pgfpathlineto{\pgfqpoint{2.837395in}{1.428067in}}%
\pgfpathlineto{\pgfqpoint{2.852553in}{1.427030in}}%
\pgfpathlineto{\pgfqpoint{2.853130in}{1.419128in}}%
\pgfpathlineto{\pgfqpoint{2.867100in}{1.421144in}}%
\pgfpathlineto{\pgfqpoint{2.867003in}{1.423229in}}%
\pgfpathlineto{\pgfqpoint{2.880748in}{1.424222in}}%
\pgfpathlineto{\pgfqpoint{2.883523in}{1.417325in}}%
\pgfpathlineto{\pgfqpoint{2.881668in}{1.410091in}}%
\pgfpathlineto{\pgfqpoint{2.882640in}{1.396033in}}%
\pgfpathlineto{\pgfqpoint{2.875723in}{1.395664in}}%
\pgfpathlineto{\pgfqpoint{2.876585in}{1.380575in}}%
\pgfpathlineto{\pgfqpoint{2.869245in}{1.379552in}}%
\pgfpathlineto{\pgfqpoint{2.864652in}{1.379289in}}%
\pgfpathlineto{\pgfqpoint{2.864952in}{1.371264in}}%
\pgfpathlineto{\pgfqpoint{2.861822in}{1.358808in}}%
\pgfpathlineto{\pgfqpoint{2.854867in}{1.358657in}}%
\pgfpathlineto{\pgfqpoint{2.850733in}{1.370027in}}%
\pgfpathlineto{\pgfqpoint{2.842704in}{1.366486in}}%
\pgfpathclose%
\pgfusepath{fill}%
\end{pgfscope}%
\begin{pgfscope}%
\pgfpathrectangle{\pgfqpoint{0.100000in}{0.100000in}}{\pgfqpoint{3.608454in}{2.310000in}}%
\pgfusepath{clip}%
\pgfsetbuttcap%
\pgfsetmiterjoin%
\definecolor{currentfill}{rgb}{0.000000,0.537255,0.731373}%
\pgfsetfillcolor{currentfill}%
\pgfsetlinewidth{0.000000pt}%
\definecolor{currentstroke}{rgb}{0.000000,0.000000,0.000000}%
\pgfsetstrokecolor{currentstroke}%
\pgfsetstrokeopacity{0.000000}%
\pgfsetdash{}{0pt}%
\pgfpathmoveto{\pgfqpoint{2.402794in}{1.625171in}}%
\pgfpathlineto{\pgfqpoint{2.404561in}{1.600729in}}%
\pgfpathlineto{\pgfqpoint{2.381293in}{1.599262in}}%
\pgfpathlineto{\pgfqpoint{2.376710in}{1.601691in}}%
\pgfpathlineto{\pgfqpoint{2.374950in}{1.609365in}}%
\pgfpathlineto{\pgfqpoint{2.371380in}{1.613652in}}%
\pgfpathlineto{\pgfqpoint{2.360207in}{1.613167in}}%
\pgfpathlineto{\pgfqpoint{2.360507in}{1.606246in}}%
\pgfpathlineto{\pgfqpoint{2.346907in}{1.605656in}}%
\pgfpathlineto{\pgfqpoint{2.306014in}{1.604048in}}%
\pgfpathlineto{\pgfqpoint{2.304826in}{1.617701in}}%
\pgfpathlineto{\pgfqpoint{2.303199in}{1.666191in}}%
\pgfpathlineto{\pgfqpoint{2.301639in}{1.699356in}}%
\pgfpathlineto{\pgfqpoint{2.324231in}{1.700249in}}%
\pgfpathlineto{\pgfqpoint{2.325187in}{1.694132in}}%
\pgfpathlineto{\pgfqpoint{2.325644in}{1.688292in}}%
\pgfpathlineto{\pgfqpoint{2.334254in}{1.681222in}}%
\pgfpathlineto{\pgfqpoint{2.327890in}{1.671376in}}%
\pgfpathlineto{\pgfqpoint{2.330505in}{1.653247in}}%
\pgfpathlineto{\pgfqpoint{2.333060in}{1.651644in}}%
\pgfpathlineto{\pgfqpoint{2.336274in}{1.640299in}}%
\pgfpathlineto{\pgfqpoint{2.342986in}{1.636283in}}%
\pgfpathlineto{\pgfqpoint{2.356020in}{1.633465in}}%
\pgfpathlineto{\pgfqpoint{2.361136in}{1.623207in}}%
\pgfpathlineto{\pgfqpoint{2.402794in}{1.625171in}}%
\pgfpathclose%
\pgfusepath{fill}%
\end{pgfscope}%
\begin{pgfscope}%
\pgfpathrectangle{\pgfqpoint{0.100000in}{0.100000in}}{\pgfqpoint{3.608454in}{2.310000in}}%
\pgfusepath{clip}%
\pgfsetbuttcap%
\pgfsetmiterjoin%
\definecolor{currentfill}{rgb}{0.000000,0.576471,0.711765}%
\pgfsetfillcolor{currentfill}%
\pgfsetlinewidth{0.000000pt}%
\definecolor{currentstroke}{rgb}{0.000000,0.000000,0.000000}%
\pgfsetstrokecolor{currentstroke}%
\pgfsetstrokeopacity{0.000000}%
\pgfsetdash{}{0pt}%
\pgfpathmoveto{\pgfqpoint{3.105607in}{1.274029in}}%
\pgfpathlineto{\pgfqpoint{3.099433in}{1.272390in}}%
\pgfpathlineto{\pgfqpoint{3.095309in}{1.277259in}}%
\pgfpathlineto{\pgfqpoint{3.089238in}{1.276542in}}%
\pgfpathlineto{\pgfqpoint{3.081510in}{1.270060in}}%
\pgfpathlineto{\pgfqpoint{3.074664in}{1.267099in}}%
\pgfpathlineto{\pgfqpoint{3.069985in}{1.273083in}}%
\pgfpathlineto{\pgfqpoint{3.057510in}{1.278132in}}%
\pgfpathlineto{\pgfqpoint{3.055501in}{1.285981in}}%
\pgfpathlineto{\pgfqpoint{3.057612in}{1.293606in}}%
\pgfpathlineto{\pgfqpoint{3.064782in}{1.300929in}}%
\pgfpathlineto{\pgfqpoint{3.067687in}{1.299266in}}%
\pgfpathlineto{\pgfqpoint{3.074021in}{1.306357in}}%
\pgfpathlineto{\pgfqpoint{3.075683in}{1.311970in}}%
\pgfpathlineto{\pgfqpoint{3.081757in}{1.317949in}}%
\pgfpathlineto{\pgfqpoint{3.084921in}{1.329255in}}%
\pgfpathlineto{\pgfqpoint{3.092431in}{1.332135in}}%
\pgfpathlineto{\pgfqpoint{3.098792in}{1.325084in}}%
\pgfpathlineto{\pgfqpoint{3.103381in}{1.324700in}}%
\pgfpathlineto{\pgfqpoint{3.110914in}{1.313734in}}%
\pgfpathlineto{\pgfqpoint{3.115324in}{1.315206in}}%
\pgfpathlineto{\pgfqpoint{3.125802in}{1.307188in}}%
\pgfpathlineto{\pgfqpoint{3.131753in}{1.305762in}}%
\pgfpathlineto{\pgfqpoint{3.128475in}{1.293170in}}%
\pgfpathlineto{\pgfqpoint{3.120247in}{1.288198in}}%
\pgfpathlineto{\pgfqpoint{3.117815in}{1.273052in}}%
\pgfpathlineto{\pgfqpoint{3.105607in}{1.274029in}}%
\pgfpathclose%
\pgfusepath{fill}%
\end{pgfscope}%
\begin{pgfscope}%
\pgfpathrectangle{\pgfqpoint{0.100000in}{0.100000in}}{\pgfqpoint{3.608454in}{2.310000in}}%
\pgfusepath{clip}%
\pgfsetbuttcap%
\pgfsetmiterjoin%
\definecolor{currentfill}{rgb}{0.000000,0.909804,0.545098}%
\pgfsetfillcolor{currentfill}%
\pgfsetlinewidth{0.000000pt}%
\definecolor{currentstroke}{rgb}{0.000000,0.000000,0.000000}%
\pgfsetstrokecolor{currentstroke}%
\pgfsetstrokeopacity{0.000000}%
\pgfsetdash{}{0pt}%
\pgfpathmoveto{\pgfqpoint{0.616244in}{1.922404in}}%
\pgfpathlineto{\pgfqpoint{0.668725in}{1.906595in}}%
\pgfpathlineto{\pgfqpoint{0.695319in}{1.898980in}}%
\pgfpathlineto{\pgfqpoint{0.686047in}{1.866076in}}%
\pgfpathlineto{\pgfqpoint{0.684062in}{1.866632in}}%
\pgfpathlineto{\pgfqpoint{0.674134in}{1.833816in}}%
\pgfpathlineto{\pgfqpoint{0.706614in}{1.825006in}}%
\pgfpathlineto{\pgfqpoint{0.691075in}{1.767312in}}%
\pgfpathlineto{\pgfqpoint{0.654872in}{1.777406in}}%
\pgfpathlineto{\pgfqpoint{0.605093in}{1.791654in}}%
\pgfpathlineto{\pgfqpoint{0.621562in}{1.848713in}}%
\pgfpathlineto{\pgfqpoint{0.595602in}{1.856614in}}%
\pgfpathlineto{\pgfqpoint{0.609494in}{1.902912in}}%
\pgfpathlineto{\pgfqpoint{0.616244in}{1.922404in}}%
\pgfpathclose%
\pgfusepath{fill}%
\end{pgfscope}%
\begin{pgfscope}%
\pgfpathrectangle{\pgfqpoint{0.100000in}{0.100000in}}{\pgfqpoint{3.608454in}{2.310000in}}%
\pgfusepath{clip}%
\pgfsetbuttcap%
\pgfsetmiterjoin%
\definecolor{currentfill}{rgb}{0.000000,0.403922,0.798039}%
\pgfsetfillcolor{currentfill}%
\pgfsetlinewidth{0.000000pt}%
\definecolor{currentstroke}{rgb}{0.000000,0.000000,0.000000}%
\pgfsetstrokecolor{currentstroke}%
\pgfsetstrokeopacity{0.000000}%
\pgfsetdash{}{0pt}%
\pgfpathmoveto{\pgfqpoint{2.139356in}{1.083433in}}%
\pgfpathlineto{\pgfqpoint{2.134726in}{1.115276in}}%
\pgfpathlineto{\pgfqpoint{2.106071in}{1.115212in}}%
\pgfpathlineto{\pgfqpoint{2.106820in}{1.155478in}}%
\pgfpathlineto{\pgfqpoint{2.131026in}{1.155307in}}%
\pgfpathlineto{\pgfqpoint{2.130979in}{1.163114in}}%
\pgfpathlineto{\pgfqpoint{2.165811in}{1.161887in}}%
\pgfpathlineto{\pgfqpoint{2.165488in}{1.142173in}}%
\pgfpathlineto{\pgfqpoint{2.178864in}{1.142325in}}%
\pgfpathlineto{\pgfqpoint{2.178906in}{1.127188in}}%
\pgfpathlineto{\pgfqpoint{2.192391in}{1.127027in}}%
\pgfpathlineto{\pgfqpoint{2.194911in}{1.119004in}}%
\pgfpathlineto{\pgfqpoint{2.199417in}{1.113130in}}%
\pgfpathlineto{\pgfqpoint{2.206212in}{1.109649in}}%
\pgfpathlineto{\pgfqpoint{2.206038in}{1.100680in}}%
\pgfpathlineto{\pgfqpoint{2.201478in}{1.093804in}}%
\pgfpathlineto{\pgfqpoint{2.201915in}{1.084462in}}%
\pgfpathlineto{\pgfqpoint{2.176762in}{1.083808in}}%
\pgfpathlineto{\pgfqpoint{2.163020in}{1.082921in}}%
\pgfpathlineto{\pgfqpoint{2.139356in}{1.083433in}}%
\pgfpathclose%
\pgfusepath{fill}%
\end{pgfscope}%
\begin{pgfscope}%
\pgfpathrectangle{\pgfqpoint{0.100000in}{0.100000in}}{\pgfqpoint{3.608454in}{2.310000in}}%
\pgfusepath{clip}%
\pgfsetbuttcap%
\pgfsetmiterjoin%
\definecolor{currentfill}{rgb}{0.000000,0.533333,0.733333}%
\pgfsetfillcolor{currentfill}%
\pgfsetlinewidth{0.000000pt}%
\definecolor{currentstroke}{rgb}{0.000000,0.000000,0.000000}%
\pgfsetstrokecolor{currentstroke}%
\pgfsetstrokeopacity{0.000000}%
\pgfsetdash{}{0pt}%
\pgfpathmoveto{\pgfqpoint{3.081713in}{0.979204in}}%
\pgfpathlineto{\pgfqpoint{3.081498in}{0.967552in}}%
\pgfpathlineto{\pgfqpoint{3.070676in}{0.962442in}}%
\pgfpathlineto{\pgfqpoint{3.064743in}{0.967659in}}%
\pgfpathlineto{\pgfqpoint{3.047350in}{0.952216in}}%
\pgfpathlineto{\pgfqpoint{3.034445in}{0.961587in}}%
\pgfpathlineto{\pgfqpoint{3.018836in}{0.966601in}}%
\pgfpathlineto{\pgfqpoint{3.010939in}{0.970937in}}%
\pgfpathlineto{\pgfqpoint{3.005840in}{0.971004in}}%
\pgfpathlineto{\pgfqpoint{3.016045in}{0.985596in}}%
\pgfpathlineto{\pgfqpoint{2.999338in}{0.990301in}}%
\pgfpathlineto{\pgfqpoint{2.988314in}{0.999433in}}%
\pgfpathlineto{\pgfqpoint{2.983807in}{0.993720in}}%
\pgfpathlineto{\pgfqpoint{2.971145in}{0.996030in}}%
\pgfpathlineto{\pgfqpoint{2.966134in}{1.004041in}}%
\pgfpathlineto{\pgfqpoint{2.958909in}{1.001900in}}%
\pgfpathlineto{\pgfqpoint{2.963991in}{1.017466in}}%
\pgfpathlineto{\pgfqpoint{2.961263in}{1.022919in}}%
\pgfpathlineto{\pgfqpoint{2.962380in}{1.029555in}}%
\pgfpathlineto{\pgfqpoint{2.968957in}{1.035717in}}%
\pgfpathlineto{\pgfqpoint{2.975999in}{1.050234in}}%
\pgfpathlineto{\pgfqpoint{2.982517in}{1.046654in}}%
\pgfpathlineto{\pgfqpoint{2.990720in}{1.049995in}}%
\pgfpathlineto{\pgfqpoint{2.989806in}{1.055375in}}%
\pgfpathlineto{\pgfqpoint{3.025339in}{1.058515in}}%
\pgfpathlineto{\pgfqpoint{3.026275in}{1.051810in}}%
\pgfpathlineto{\pgfqpoint{3.034331in}{1.046602in}}%
\pgfpathlineto{\pgfqpoint{3.030270in}{1.030329in}}%
\pgfpathlineto{\pgfqpoint{3.041428in}{1.022979in}}%
\pgfpathlineto{\pgfqpoint{3.046504in}{1.026091in}}%
\pgfpathlineto{\pgfqpoint{3.049502in}{1.021849in}}%
\pgfpathlineto{\pgfqpoint{3.049719in}{1.006843in}}%
\pgfpathlineto{\pgfqpoint{3.052596in}{1.003504in}}%
\pgfpathlineto{\pgfqpoint{3.051594in}{0.997455in}}%
\pgfpathlineto{\pgfqpoint{3.059264in}{0.991086in}}%
\pgfpathlineto{\pgfqpoint{3.065412in}{0.981312in}}%
\pgfpathlineto{\pgfqpoint{3.073391in}{0.976759in}}%
\pgfpathlineto{\pgfqpoint{3.081713in}{0.979204in}}%
\pgfpathclose%
\pgfusepath{fill}%
\end{pgfscope}%
\begin{pgfscope}%
\pgfpathrectangle{\pgfqpoint{0.100000in}{0.100000in}}{\pgfqpoint{3.608454in}{2.310000in}}%
\pgfusepath{clip}%
\pgfsetbuttcap%
\pgfsetmiterjoin%
\definecolor{currentfill}{rgb}{0.000000,0.811765,0.594118}%
\pgfsetfillcolor{currentfill}%
\pgfsetlinewidth{0.000000pt}%
\definecolor{currentstroke}{rgb}{0.000000,0.000000,0.000000}%
\pgfsetstrokecolor{currentstroke}%
\pgfsetstrokeopacity{0.000000}%
\pgfsetdash{}{0pt}%
\pgfpathmoveto{\pgfqpoint{2.876403in}{1.073700in}}%
\pgfpathlineto{\pgfqpoint{2.876200in}{1.069013in}}%
\pgfpathlineto{\pgfqpoint{2.869918in}{1.063176in}}%
\pgfpathlineto{\pgfqpoint{2.864683in}{1.055271in}}%
\pgfpathlineto{\pgfqpoint{2.863588in}{1.049271in}}%
\pgfpathlineto{\pgfqpoint{2.850394in}{1.050302in}}%
\pgfpathlineto{\pgfqpoint{2.847354in}{1.055558in}}%
\pgfpathlineto{\pgfqpoint{2.843160in}{1.055221in}}%
\pgfpathlineto{\pgfqpoint{2.842470in}{1.059416in}}%
\pgfpathlineto{\pgfqpoint{2.850115in}{1.070383in}}%
\pgfpathlineto{\pgfqpoint{2.840206in}{1.081206in}}%
\pgfpathlineto{\pgfqpoint{2.834354in}{1.080705in}}%
\pgfpathlineto{\pgfqpoint{2.835727in}{1.088275in}}%
\pgfpathlineto{\pgfqpoint{2.837326in}{1.091183in}}%
\pgfpathlineto{\pgfqpoint{2.851025in}{1.094550in}}%
\pgfpathlineto{\pgfqpoint{2.858178in}{1.097707in}}%
\pgfpathlineto{\pgfqpoint{2.867388in}{1.085037in}}%
\pgfpathlineto{\pgfqpoint{2.872545in}{1.080983in}}%
\pgfpathlineto{\pgfqpoint{2.876403in}{1.073700in}}%
\pgfpathclose%
\pgfusepath{fill}%
\end{pgfscope}%
\begin{pgfscope}%
\pgfpathrectangle{\pgfqpoint{0.100000in}{0.100000in}}{\pgfqpoint{3.608454in}{2.310000in}}%
\pgfusepath{clip}%
\pgfsetbuttcap%
\pgfsetmiterjoin%
\definecolor{currentfill}{rgb}{0.000000,0.768627,0.615686}%
\pgfsetfillcolor{currentfill}%
\pgfsetlinewidth{0.000000pt}%
\definecolor{currentstroke}{rgb}{0.000000,0.000000,0.000000}%
\pgfsetstrokecolor{currentstroke}%
\pgfsetstrokeopacity{0.000000}%
\pgfsetdash{}{0pt}%
\pgfpathmoveto{\pgfqpoint{2.288166in}{1.125922in}}%
\pgfpathlineto{\pgfqpoint{2.286012in}{1.115699in}}%
\pgfpathlineto{\pgfqpoint{2.286066in}{1.109865in}}%
\pgfpathlineto{\pgfqpoint{2.279199in}{1.109799in}}%
\pgfpathlineto{\pgfqpoint{2.279258in}{1.102907in}}%
\pgfpathlineto{\pgfqpoint{2.272396in}{1.102853in}}%
\pgfpathlineto{\pgfqpoint{2.272381in}{1.109733in}}%
\pgfpathlineto{\pgfqpoint{2.245103in}{1.109604in}}%
\pgfpathlineto{\pgfqpoint{2.241638in}{1.113006in}}%
\pgfpathlineto{\pgfqpoint{2.241311in}{1.142821in}}%
\pgfpathlineto{\pgfqpoint{2.243255in}{1.143432in}}%
\pgfpathlineto{\pgfqpoint{2.248459in}{1.143536in}}%
\pgfpathlineto{\pgfqpoint{2.248336in}{1.167927in}}%
\pgfpathlineto{\pgfqpoint{2.289673in}{1.168159in}}%
\pgfpathlineto{\pgfqpoint{2.289909in}{1.144788in}}%
\pgfpathlineto{\pgfqpoint{2.288004in}{1.144708in}}%
\pgfpathlineto{\pgfqpoint{2.288166in}{1.125922in}}%
\pgfpathclose%
\pgfusepath{fill}%
\end{pgfscope}%
\begin{pgfscope}%
\pgfpathrectangle{\pgfqpoint{0.100000in}{0.100000in}}{\pgfqpoint{3.608454in}{2.310000in}}%
\pgfusepath{clip}%
\pgfsetbuttcap%
\pgfsetmiterjoin%
\definecolor{currentfill}{rgb}{0.000000,0.847059,0.576471}%
\pgfsetfillcolor{currentfill}%
\pgfsetlinewidth{0.000000pt}%
\definecolor{currentstroke}{rgb}{0.000000,0.000000,0.000000}%
\pgfsetstrokecolor{currentstroke}%
\pgfsetstrokeopacity{0.000000}%
\pgfsetdash{}{0pt}%
\pgfpathmoveto{\pgfqpoint{3.197116in}{1.277528in}}%
\pgfpathlineto{\pgfqpoint{3.199968in}{1.263498in}}%
\pgfpathlineto{\pgfqpoint{3.195545in}{1.249102in}}%
\pgfpathlineto{\pgfqpoint{3.187275in}{1.247460in}}%
\pgfpathlineto{\pgfqpoint{3.152496in}{1.240772in}}%
\pgfpathlineto{\pgfqpoint{3.109908in}{1.233039in}}%
\pgfpathlineto{\pgfqpoint{3.104858in}{1.232143in}}%
\pgfpathlineto{\pgfqpoint{3.105607in}{1.274029in}}%
\pgfpathlineto{\pgfqpoint{3.117815in}{1.273052in}}%
\pgfpathlineto{\pgfqpoint{3.120247in}{1.288198in}}%
\pgfpathlineto{\pgfqpoint{3.128475in}{1.293170in}}%
\pgfpathlineto{\pgfqpoint{3.131753in}{1.305762in}}%
\pgfpathlineto{\pgfqpoint{3.140579in}{1.302693in}}%
\pgfpathlineto{\pgfqpoint{3.149648in}{1.328164in}}%
\pgfpathlineto{\pgfqpoint{3.149422in}{1.332537in}}%
\pgfpathlineto{\pgfqpoint{3.153524in}{1.337961in}}%
\pgfpathlineto{\pgfqpoint{3.159782in}{1.331466in}}%
\pgfpathlineto{\pgfqpoint{3.159486in}{1.315467in}}%
\pgfpathlineto{\pgfqpoint{3.154386in}{1.307543in}}%
\pgfpathlineto{\pgfqpoint{3.155665in}{1.301994in}}%
\pgfpathlineto{\pgfqpoint{3.184110in}{1.299032in}}%
\pgfpathlineto{\pgfqpoint{3.178373in}{1.294019in}}%
\pgfpathlineto{\pgfqpoint{3.181322in}{1.282092in}}%
\pgfpathlineto{\pgfqpoint{3.189244in}{1.282467in}}%
\pgfpathlineto{\pgfqpoint{3.197116in}{1.277528in}}%
\pgfpathclose%
\pgfusepath{fill}%
\end{pgfscope}%
\begin{pgfscope}%
\pgfpathrectangle{\pgfqpoint{0.100000in}{0.100000in}}{\pgfqpoint{3.608454in}{2.310000in}}%
\pgfusepath{clip}%
\pgfsetbuttcap%
\pgfsetmiterjoin%
\definecolor{currentfill}{rgb}{0.000000,0.713725,0.643137}%
\pgfsetfillcolor{currentfill}%
\pgfsetlinewidth{0.000000pt}%
\definecolor{currentstroke}{rgb}{0.000000,0.000000,0.000000}%
\pgfsetstrokecolor{currentstroke}%
\pgfsetstrokeopacity{0.000000}%
\pgfsetdash{}{0pt}%
\pgfpathmoveto{\pgfqpoint{1.034613in}{1.812423in}}%
\pgfpathlineto{\pgfqpoint{1.027024in}{1.825238in}}%
\pgfpathlineto{\pgfqpoint{1.013709in}{1.836495in}}%
\pgfpathlineto{\pgfqpoint{1.012788in}{1.842181in}}%
\pgfpathlineto{\pgfqpoint{0.999848in}{1.846301in}}%
\pgfpathlineto{\pgfqpoint{0.991130in}{1.850884in}}%
\pgfpathlineto{\pgfqpoint{0.983732in}{1.852523in}}%
\pgfpathlineto{\pgfqpoint{0.981526in}{1.858299in}}%
\pgfpathlineto{\pgfqpoint{0.975129in}{1.855275in}}%
\pgfpathlineto{\pgfqpoint{0.972986in}{1.865357in}}%
\pgfpathlineto{\pgfqpoint{0.975886in}{1.873015in}}%
\pgfpathlineto{\pgfqpoint{0.968833in}{1.878204in}}%
\pgfpathlineto{\pgfqpoint{0.969814in}{1.885604in}}%
\pgfpathlineto{\pgfqpoint{0.964170in}{1.890384in}}%
\pgfpathlineto{\pgfqpoint{0.969006in}{1.895451in}}%
\pgfpathlineto{\pgfqpoint{0.970494in}{1.901193in}}%
\pgfpathlineto{\pgfqpoint{0.967552in}{1.905305in}}%
\pgfpathlineto{\pgfqpoint{0.969772in}{1.911307in}}%
\pgfpathlineto{\pgfqpoint{0.976747in}{1.912867in}}%
\pgfpathlineto{\pgfqpoint{0.981235in}{1.920433in}}%
\pgfpathlineto{\pgfqpoint{0.989615in}{1.915702in}}%
\pgfpathlineto{\pgfqpoint{1.004445in}{1.925280in}}%
\pgfpathlineto{\pgfqpoint{1.009235in}{1.946422in}}%
\pgfpathlineto{\pgfqpoint{1.012082in}{1.949489in}}%
\pgfpathlineto{\pgfqpoint{1.018689in}{1.956057in}}%
\pgfpathlineto{\pgfqpoint{1.011182in}{1.972459in}}%
\pgfpathlineto{\pgfqpoint{1.018436in}{1.970411in}}%
\pgfpathlineto{\pgfqpoint{1.025284in}{1.975956in}}%
\pgfpathlineto{\pgfqpoint{1.028023in}{1.976205in}}%
\pgfpathlineto{\pgfqpoint{1.034877in}{1.968576in}}%
\pgfpathlineto{\pgfqpoint{1.048671in}{1.970878in}}%
\pgfpathlineto{\pgfqpoint{1.053520in}{1.975559in}}%
\pgfpathlineto{\pgfqpoint{1.061696in}{1.979359in}}%
\pgfpathlineto{\pgfqpoint{1.067306in}{1.970565in}}%
\pgfpathlineto{\pgfqpoint{1.064456in}{1.964544in}}%
\pgfpathlineto{\pgfqpoint{1.067866in}{1.953218in}}%
\pgfpathlineto{\pgfqpoint{1.066629in}{1.948325in}}%
\pgfpathlineto{\pgfqpoint{1.073148in}{1.929599in}}%
\pgfpathlineto{\pgfqpoint{1.077934in}{1.924256in}}%
\pgfpathlineto{\pgfqpoint{1.074674in}{1.909170in}}%
\pgfpathlineto{\pgfqpoint{1.079391in}{1.904532in}}%
\pgfpathlineto{\pgfqpoint{1.088102in}{1.901286in}}%
\pgfpathlineto{\pgfqpoint{1.091986in}{1.894613in}}%
\pgfpathlineto{\pgfqpoint{1.094071in}{1.877735in}}%
\pgfpathlineto{\pgfqpoint{1.093092in}{1.871853in}}%
\pgfpathlineto{\pgfqpoint{1.100477in}{1.862882in}}%
\pgfpathlineto{\pgfqpoint{1.102494in}{1.863989in}}%
\pgfpathlineto{\pgfqpoint{1.100225in}{1.852591in}}%
\pgfpathlineto{\pgfqpoint{1.066665in}{1.859594in}}%
\pgfpathlineto{\pgfqpoint{1.063909in}{1.846064in}}%
\pgfpathlineto{\pgfqpoint{1.061051in}{1.840480in}}%
\pgfpathlineto{\pgfqpoint{1.057180in}{1.823892in}}%
\pgfpathlineto{\pgfqpoint{1.049281in}{1.819620in}}%
\pgfpathlineto{\pgfqpoint{1.037674in}{1.816193in}}%
\pgfpathlineto{\pgfqpoint{1.034613in}{1.812423in}}%
\pgfpathclose%
\pgfusepath{fill}%
\end{pgfscope}%
\begin{pgfscope}%
\pgfpathrectangle{\pgfqpoint{0.100000in}{0.100000in}}{\pgfqpoint{3.608454in}{2.310000in}}%
\pgfusepath{clip}%
\pgfsetbuttcap%
\pgfsetmiterjoin%
\definecolor{currentfill}{rgb}{0.000000,0.666667,0.666667}%
\pgfsetfillcolor{currentfill}%
\pgfsetlinewidth{0.000000pt}%
\definecolor{currentstroke}{rgb}{0.000000,0.000000,0.000000}%
\pgfsetstrokecolor{currentstroke}%
\pgfsetstrokeopacity{0.000000}%
\pgfsetdash{}{0pt}%
\pgfpathmoveto{\pgfqpoint{1.183026in}{1.348649in}}%
\pgfpathlineto{\pgfqpoint{1.177278in}{1.346155in}}%
\pgfpathlineto{\pgfqpoint{1.170873in}{1.338075in}}%
\pgfpathlineto{\pgfqpoint{1.168304in}{1.339366in}}%
\pgfpathlineto{\pgfqpoint{1.159137in}{1.332405in}}%
\pgfpathlineto{\pgfqpoint{1.150177in}{1.332380in}}%
\pgfpathlineto{\pgfqpoint{1.141396in}{1.315621in}}%
\pgfpathlineto{\pgfqpoint{1.134341in}{1.312779in}}%
\pgfpathlineto{\pgfqpoint{1.130720in}{1.308550in}}%
\pgfpathlineto{\pgfqpoint{1.104844in}{1.313028in}}%
\pgfpathlineto{\pgfqpoint{1.024916in}{1.327610in}}%
\pgfpathlineto{\pgfqpoint{1.005205in}{1.331999in}}%
\pgfpathlineto{\pgfqpoint{1.008850in}{1.352444in}}%
\pgfpathlineto{\pgfqpoint{1.015614in}{1.351012in}}%
\pgfpathlineto{\pgfqpoint{1.016803in}{1.357032in}}%
\pgfpathlineto{\pgfqpoint{1.023679in}{1.356359in}}%
\pgfpathlineto{\pgfqpoint{1.026917in}{1.376417in}}%
\pgfpathlineto{\pgfqpoint{1.032898in}{1.379983in}}%
\pgfpathlineto{\pgfqpoint{1.038444in}{1.389127in}}%
\pgfpathlineto{\pgfqpoint{1.034771in}{1.396805in}}%
\pgfpathlineto{\pgfqpoint{1.030622in}{1.400197in}}%
\pgfpathlineto{\pgfqpoint{1.031123in}{1.409823in}}%
\pgfpathlineto{\pgfqpoint{1.036855in}{1.417471in}}%
\pgfpathlineto{\pgfqpoint{1.045136in}{1.415471in}}%
\pgfpathlineto{\pgfqpoint{1.051426in}{1.418402in}}%
\pgfpathlineto{\pgfqpoint{1.053099in}{1.429145in}}%
\pgfpathlineto{\pgfqpoint{1.060054in}{1.435184in}}%
\pgfpathlineto{\pgfqpoint{1.064107in}{1.434490in}}%
\pgfpathlineto{\pgfqpoint{1.068484in}{1.440713in}}%
\pgfpathlineto{\pgfqpoint{1.077995in}{1.437925in}}%
\pgfpathlineto{\pgfqpoint{1.111444in}{1.431696in}}%
\pgfpathlineto{\pgfqpoint{1.103689in}{1.391163in}}%
\pgfpathlineto{\pgfqpoint{1.076288in}{1.396187in}}%
\pgfpathlineto{\pgfqpoint{1.069417in}{1.390629in}}%
\pgfpathlineto{\pgfqpoint{1.065800in}{1.369302in}}%
\pgfpathlineto{\pgfqpoint{1.136475in}{1.356525in}}%
\pgfpathlineto{\pgfqpoint{1.183026in}{1.348649in}}%
\pgfpathclose%
\pgfusepath{fill}%
\end{pgfscope}%
\begin{pgfscope}%
\pgfpathrectangle{\pgfqpoint{0.100000in}{0.100000in}}{\pgfqpoint{3.608454in}{2.310000in}}%
\pgfusepath{clip}%
\pgfsetbuttcap%
\pgfsetmiterjoin%
\definecolor{currentfill}{rgb}{0.000000,0.631373,0.684314}%
\pgfsetfillcolor{currentfill}%
\pgfsetlinewidth{0.000000pt}%
\definecolor{currentstroke}{rgb}{0.000000,0.000000,0.000000}%
\pgfsetstrokecolor{currentstroke}%
\pgfsetstrokeopacity{0.000000}%
\pgfsetdash{}{0pt}%
\pgfpathmoveto{\pgfqpoint{2.102150in}{1.181491in}}%
\pgfpathlineto{\pgfqpoint{2.081042in}{1.181525in}}%
\pgfpathlineto{\pgfqpoint{2.079502in}{1.177066in}}%
\pgfpathlineto{\pgfqpoint{2.079386in}{1.149757in}}%
\pgfpathlineto{\pgfqpoint{2.055192in}{1.149919in}}%
\pgfpathlineto{\pgfqpoint{2.055083in}{1.136146in}}%
\pgfpathlineto{\pgfqpoint{2.043055in}{1.136287in}}%
\pgfpathlineto{\pgfqpoint{2.043578in}{1.181785in}}%
\pgfpathlineto{\pgfqpoint{2.045886in}{1.181780in}}%
\pgfpathlineto{\pgfqpoint{2.046716in}{1.239913in}}%
\pgfpathlineto{\pgfqpoint{2.101288in}{1.239463in}}%
\pgfpathlineto{\pgfqpoint{2.102082in}{1.211883in}}%
\pgfpathlineto{\pgfqpoint{2.102150in}{1.181491in}}%
\pgfpathclose%
\pgfusepath{fill}%
\end{pgfscope}%
\begin{pgfscope}%
\pgfpathrectangle{\pgfqpoint{0.100000in}{0.100000in}}{\pgfqpoint{3.608454in}{2.310000in}}%
\pgfusepath{clip}%
\pgfsetbuttcap%
\pgfsetmiterjoin%
\definecolor{currentfill}{rgb}{0.000000,0.654902,0.672549}%
\pgfsetfillcolor{currentfill}%
\pgfsetlinewidth{0.000000pt}%
\definecolor{currentstroke}{rgb}{0.000000,0.000000,0.000000}%
\pgfsetstrokecolor{currentstroke}%
\pgfsetstrokeopacity{0.000000}%
\pgfsetdash{}{0pt}%
\pgfpathmoveto{\pgfqpoint{1.411450in}{2.059486in}}%
\pgfpathlineto{\pgfqpoint{1.406013in}{2.044122in}}%
\pgfpathlineto{\pgfqpoint{1.405800in}{2.033757in}}%
\pgfpathlineto{\pgfqpoint{1.407470in}{2.024546in}}%
\pgfpathlineto{\pgfqpoint{1.403304in}{2.014885in}}%
\pgfpathlineto{\pgfqpoint{1.405583in}{2.011993in}}%
\pgfpathlineto{\pgfqpoint{1.407997in}{2.003940in}}%
\pgfpathlineto{\pgfqpoint{1.364475in}{2.009800in}}%
\pgfpathlineto{\pgfqpoint{1.323552in}{2.016288in}}%
\pgfpathlineto{\pgfqpoint{1.322819in}{2.011773in}}%
\pgfpathlineto{\pgfqpoint{1.303565in}{2.014897in}}%
\pgfpathlineto{\pgfqpoint{1.303017in}{2.028995in}}%
\pgfpathlineto{\pgfqpoint{1.305651in}{2.030921in}}%
\pgfpathlineto{\pgfqpoint{1.309306in}{2.053731in}}%
\pgfpathlineto{\pgfqpoint{1.303684in}{2.058169in}}%
\pgfpathlineto{\pgfqpoint{1.296843in}{2.059295in}}%
\pgfpathlineto{\pgfqpoint{1.290585in}{2.063837in}}%
\pgfpathlineto{\pgfqpoint{1.293100in}{2.076361in}}%
\pgfpathlineto{\pgfqpoint{1.309109in}{2.082641in}}%
\pgfpathlineto{\pgfqpoint{1.312471in}{2.096193in}}%
\pgfpathlineto{\pgfqpoint{1.328054in}{2.095490in}}%
\pgfpathlineto{\pgfqpoint{1.332093in}{2.093598in}}%
\pgfpathlineto{\pgfqpoint{1.342674in}{2.097531in}}%
\pgfpathlineto{\pgfqpoint{1.350200in}{2.095416in}}%
\pgfpathlineto{\pgfqpoint{1.358594in}{2.095823in}}%
\pgfpathlineto{\pgfqpoint{1.362402in}{2.089710in}}%
\pgfpathlineto{\pgfqpoint{1.367934in}{2.081540in}}%
\pgfpathlineto{\pgfqpoint{1.377859in}{2.077250in}}%
\pgfpathlineto{\pgfqpoint{1.385004in}{2.076421in}}%
\pgfpathlineto{\pgfqpoint{1.389242in}{2.072875in}}%
\pgfpathlineto{\pgfqpoint{1.404841in}{2.072375in}}%
\pgfpathlineto{\pgfqpoint{1.411723in}{2.068893in}}%
\pgfpathlineto{\pgfqpoint{1.411450in}{2.059486in}}%
\pgfpathclose%
\pgfusepath{fill}%
\end{pgfscope}%
\begin{pgfscope}%
\pgfpathrectangle{\pgfqpoint{0.100000in}{0.100000in}}{\pgfqpoint{3.608454in}{2.310000in}}%
\pgfusepath{clip}%
\pgfsetbuttcap%
\pgfsetmiterjoin%
\definecolor{currentfill}{rgb}{0.000000,0.521569,0.739216}%
\pgfsetfillcolor{currentfill}%
\pgfsetlinewidth{0.000000pt}%
\definecolor{currentstroke}{rgb}{0.000000,0.000000,0.000000}%
\pgfsetstrokecolor{currentstroke}%
\pgfsetstrokeopacity{0.000000}%
\pgfsetdash{}{0pt}%
\pgfpathmoveto{\pgfqpoint{2.002513in}{1.074759in}}%
\pgfpathlineto{\pgfqpoint{2.002912in}{1.098705in}}%
\pgfpathlineto{\pgfqpoint{2.003198in}{1.115990in}}%
\pgfpathlineto{\pgfqpoint{1.990571in}{1.116149in}}%
\pgfpathlineto{\pgfqpoint{1.990717in}{1.123080in}}%
\pgfpathlineto{\pgfqpoint{1.983944in}{1.123221in}}%
\pgfpathlineto{\pgfqpoint{1.984100in}{1.130082in}}%
\pgfpathlineto{\pgfqpoint{1.977272in}{1.130240in}}%
\pgfpathlineto{\pgfqpoint{1.977572in}{1.143978in}}%
\pgfpathlineto{\pgfqpoint{1.985971in}{1.147335in}}%
\pgfpathlineto{\pgfqpoint{1.984832in}{1.151230in}}%
\pgfpathlineto{\pgfqpoint{1.950456in}{1.151489in}}%
\pgfpathlineto{\pgfqpoint{1.951284in}{1.183534in}}%
\pgfpathlineto{\pgfqpoint{1.971148in}{1.183063in}}%
\pgfpathlineto{\pgfqpoint{2.043578in}{1.181785in}}%
\pgfpathlineto{\pgfqpoint{2.043055in}{1.136287in}}%
\pgfpathlineto{\pgfqpoint{2.055083in}{1.136146in}}%
\pgfpathlineto{\pgfqpoint{2.055192in}{1.149919in}}%
\pgfpathlineto{\pgfqpoint{2.079386in}{1.149757in}}%
\pgfpathlineto{\pgfqpoint{2.079502in}{1.177066in}}%
\pgfpathlineto{\pgfqpoint{2.081042in}{1.181525in}}%
\pgfpathlineto{\pgfqpoint{2.102150in}{1.181491in}}%
\pgfpathlineto{\pgfqpoint{2.106837in}{1.177041in}}%
\pgfpathlineto{\pgfqpoint{2.106820in}{1.155478in}}%
\pgfpathlineto{\pgfqpoint{2.106071in}{1.115212in}}%
\pgfpathlineto{\pgfqpoint{2.099224in}{1.115204in}}%
\pgfpathlineto{\pgfqpoint{2.099232in}{1.108302in}}%
\pgfpathlineto{\pgfqpoint{2.093490in}{1.108334in}}%
\pgfpathlineto{\pgfqpoint{2.088937in}{1.099591in}}%
\pgfpathlineto{\pgfqpoint{2.089737in}{1.087648in}}%
\pgfpathlineto{\pgfqpoint{2.076132in}{1.086986in}}%
\pgfpathlineto{\pgfqpoint{2.068981in}{1.083880in}}%
\pgfpathlineto{\pgfqpoint{2.065063in}{1.091216in}}%
\pgfpathlineto{\pgfqpoint{2.057637in}{1.091246in}}%
\pgfpathlineto{\pgfqpoint{2.030288in}{1.091590in}}%
\pgfpathlineto{\pgfqpoint{2.030056in}{1.074337in}}%
\pgfpathlineto{\pgfqpoint{2.002513in}{1.074759in}}%
\pgfpathclose%
\pgfusepath{fill}%
\end{pgfscope}%
\begin{pgfscope}%
\pgfpathrectangle{\pgfqpoint{0.100000in}{0.100000in}}{\pgfqpoint{3.608454in}{2.310000in}}%
\pgfusepath{clip}%
\pgfsetbuttcap%
\pgfsetmiterjoin%
\definecolor{currentfill}{rgb}{0.000000,0.796078,0.601961}%
\pgfsetfillcolor{currentfill}%
\pgfsetlinewidth{0.000000pt}%
\definecolor{currentstroke}{rgb}{0.000000,0.000000,0.000000}%
\pgfsetstrokecolor{currentstroke}%
\pgfsetstrokeopacity{0.000000}%
\pgfsetdash{}{0pt}%
\pgfpathmoveto{\pgfqpoint{1.841646in}{1.759396in}}%
\pgfpathlineto{\pgfqpoint{1.843095in}{1.749824in}}%
\pgfpathlineto{\pgfqpoint{1.835653in}{1.753170in}}%
\pgfpathlineto{\pgfqpoint{1.830770in}{1.752552in}}%
\pgfpathlineto{\pgfqpoint{1.825951in}{1.758455in}}%
\pgfpathlineto{\pgfqpoint{1.801576in}{1.759791in}}%
\pgfpathlineto{\pgfqpoint{1.801456in}{1.757724in}}%
\pgfpathlineto{\pgfqpoint{1.762841in}{1.760148in}}%
\pgfpathlineto{\pgfqpoint{1.756032in}{1.760574in}}%
\pgfpathlineto{\pgfqpoint{1.757629in}{1.774393in}}%
\pgfpathlineto{\pgfqpoint{1.759531in}{1.801865in}}%
\pgfpathlineto{\pgfqpoint{1.760814in}{1.806325in}}%
\pgfpathlineto{\pgfqpoint{1.767130in}{1.804067in}}%
\pgfpathlineto{\pgfqpoint{1.778292in}{1.808049in}}%
\pgfpathlineto{\pgfqpoint{1.781507in}{1.805854in}}%
\pgfpathlineto{\pgfqpoint{1.787641in}{1.810142in}}%
\pgfpathlineto{\pgfqpoint{1.795074in}{1.805632in}}%
\pgfpathlineto{\pgfqpoint{1.803220in}{1.813204in}}%
\pgfpathlineto{\pgfqpoint{1.802764in}{1.815650in}}%
\pgfpathlineto{\pgfqpoint{1.843805in}{1.813366in}}%
\pgfpathlineto{\pgfqpoint{1.842265in}{1.785638in}}%
\pgfpathlineto{\pgfqpoint{1.841646in}{1.759396in}}%
\pgfpathclose%
\pgfusepath{fill}%
\end{pgfscope}%
\begin{pgfscope}%
\pgfpathrectangle{\pgfqpoint{0.100000in}{0.100000in}}{\pgfqpoint{3.608454in}{2.310000in}}%
\pgfusepath{clip}%
\pgfsetbuttcap%
\pgfsetmiterjoin%
\definecolor{currentfill}{rgb}{0.000000,0.850980,0.574510}%
\pgfsetfillcolor{currentfill}%
\pgfsetlinewidth{0.000000pt}%
\definecolor{currentstroke}{rgb}{0.000000,0.000000,0.000000}%
\pgfsetstrokecolor{currentstroke}%
\pgfsetstrokeopacity{0.000000}%
\pgfsetdash{}{0pt}%
\pgfpathmoveto{\pgfqpoint{2.404037in}{0.740243in}}%
\pgfpathlineto{\pgfqpoint{2.403042in}{0.761124in}}%
\pgfpathlineto{\pgfqpoint{2.395986in}{0.760766in}}%
\pgfpathlineto{\pgfqpoint{2.395693in}{0.767722in}}%
\pgfpathlineto{\pgfqpoint{2.426622in}{0.769362in}}%
\pgfpathlineto{\pgfqpoint{2.436793in}{0.774096in}}%
\pgfpathlineto{\pgfqpoint{2.447021in}{0.775387in}}%
\pgfpathlineto{\pgfqpoint{2.447975in}{0.756350in}}%
\pgfpathlineto{\pgfqpoint{2.449720in}{0.745908in}}%
\pgfpathlineto{\pgfqpoint{2.444160in}{0.745936in}}%
\pgfpathlineto{\pgfqpoint{2.439716in}{0.742196in}}%
\pgfpathlineto{\pgfqpoint{2.404037in}{0.740243in}}%
\pgfpathclose%
\pgfusepath{fill}%
\end{pgfscope}%
\begin{pgfscope}%
\pgfpathrectangle{\pgfqpoint{0.100000in}{0.100000in}}{\pgfqpoint{3.608454in}{2.310000in}}%
\pgfusepath{clip}%
\pgfsetbuttcap%
\pgfsetmiterjoin%
\definecolor{currentfill}{rgb}{0.000000,0.427451,0.786275}%
\pgfsetfillcolor{currentfill}%
\pgfsetlinewidth{0.000000pt}%
\definecolor{currentstroke}{rgb}{0.000000,0.000000,0.000000}%
\pgfsetstrokecolor{currentstroke}%
\pgfsetstrokeopacity{0.000000}%
\pgfsetdash{}{0pt}%
\pgfpathmoveto{\pgfqpoint{1.644176in}{0.728867in}}%
\pgfpathlineto{\pgfqpoint{1.601306in}{0.732857in}}%
\pgfpathlineto{\pgfqpoint{1.599669in}{0.742990in}}%
\pgfpathlineto{\pgfqpoint{1.584520in}{0.753390in}}%
\pgfpathlineto{\pgfqpoint{1.578119in}{0.751288in}}%
\pgfpathlineto{\pgfqpoint{1.580562in}{0.779720in}}%
\pgfpathlineto{\pgfqpoint{1.578420in}{0.779909in}}%
\pgfpathlineto{\pgfqpoint{1.581375in}{0.814353in}}%
\pgfpathlineto{\pgfqpoint{1.563559in}{0.816012in}}%
\pgfpathlineto{\pgfqpoint{1.566607in}{0.850506in}}%
\pgfpathlineto{\pgfqpoint{1.624223in}{0.845743in}}%
\pgfpathlineto{\pgfqpoint{1.658603in}{0.843265in}}%
\pgfpathlineto{\pgfqpoint{1.655619in}{0.808558in}}%
\pgfpathlineto{\pgfqpoint{1.650167in}{0.808902in}}%
\pgfpathlineto{\pgfqpoint{1.644176in}{0.728867in}}%
\pgfpathclose%
\pgfusepath{fill}%
\end{pgfscope}%
\begin{pgfscope}%
\pgfpathrectangle{\pgfqpoint{0.100000in}{0.100000in}}{\pgfqpoint{3.608454in}{2.310000in}}%
\pgfusepath{clip}%
\pgfsetbuttcap%
\pgfsetmiterjoin%
\definecolor{currentfill}{rgb}{0.000000,0.521569,0.739216}%
\pgfsetfillcolor{currentfill}%
\pgfsetlinewidth{0.000000pt}%
\definecolor{currentstroke}{rgb}{0.000000,0.000000,0.000000}%
\pgfsetstrokecolor{currentstroke}%
\pgfsetstrokeopacity{0.000000}%
\pgfsetdash{}{0pt}%
\pgfpathmoveto{\pgfqpoint{2.433309in}{1.321900in}}%
\pgfpathlineto{\pgfqpoint{2.404058in}{1.320682in}}%
\pgfpathlineto{\pgfqpoint{2.402286in}{1.326271in}}%
\pgfpathlineto{\pgfqpoint{2.407445in}{1.335820in}}%
\pgfpathlineto{\pgfqpoint{2.398503in}{1.340979in}}%
\pgfpathlineto{\pgfqpoint{2.386232in}{1.344093in}}%
\pgfpathlineto{\pgfqpoint{2.380869in}{1.336490in}}%
\pgfpathlineto{\pgfqpoint{2.373739in}{1.339722in}}%
\pgfpathlineto{\pgfqpoint{2.369779in}{1.350114in}}%
\pgfpathlineto{\pgfqpoint{2.371543in}{1.352918in}}%
\pgfpathlineto{\pgfqpoint{2.368642in}{1.363523in}}%
\pgfpathlineto{\pgfqpoint{2.363949in}{1.370060in}}%
\pgfpathlineto{\pgfqpoint{2.354879in}{1.376769in}}%
\pgfpathlineto{\pgfqpoint{2.374602in}{1.377391in}}%
\pgfpathlineto{\pgfqpoint{2.376719in}{1.362091in}}%
\pgfpathlineto{\pgfqpoint{2.383166in}{1.360351in}}%
\pgfpathlineto{\pgfqpoint{2.393816in}{1.360813in}}%
\pgfpathlineto{\pgfqpoint{2.403570in}{1.368279in}}%
\pgfpathlineto{\pgfqpoint{2.402134in}{1.388642in}}%
\pgfpathlineto{\pgfqpoint{2.415980in}{1.389550in}}%
\pgfpathlineto{\pgfqpoint{2.439852in}{1.391129in}}%
\pgfpathlineto{\pgfqpoint{2.440862in}{1.377267in}}%
\pgfpathlineto{\pgfqpoint{2.464725in}{1.378777in}}%
\pgfpathlineto{\pgfqpoint{2.465371in}{1.368430in}}%
\pgfpathlineto{\pgfqpoint{2.458623in}{1.367966in}}%
\pgfpathlineto{\pgfqpoint{2.460771in}{1.330448in}}%
\pgfpathlineto{\pgfqpoint{2.439649in}{1.329250in}}%
\pgfpathlineto{\pgfqpoint{2.440188in}{1.322366in}}%
\pgfpathlineto{\pgfqpoint{2.433309in}{1.321900in}}%
\pgfpathclose%
\pgfusepath{fill}%
\end{pgfscope}%
\begin{pgfscope}%
\pgfpathrectangle{\pgfqpoint{0.100000in}{0.100000in}}{\pgfqpoint{3.608454in}{2.310000in}}%
\pgfusepath{clip}%
\pgfsetbuttcap%
\pgfsetmiterjoin%
\definecolor{currentfill}{rgb}{0.000000,0.882353,0.558824}%
\pgfsetfillcolor{currentfill}%
\pgfsetlinewidth{0.000000pt}%
\definecolor{currentstroke}{rgb}{0.000000,0.000000,0.000000}%
\pgfsetstrokecolor{currentstroke}%
\pgfsetstrokeopacity{0.000000}%
\pgfsetdash{}{0pt}%
\pgfpathmoveto{\pgfqpoint{2.342635in}{0.869578in}}%
\pgfpathlineto{\pgfqpoint{2.302086in}{0.868312in}}%
\pgfpathlineto{\pgfqpoint{2.299404in}{0.870120in}}%
\pgfpathlineto{\pgfqpoint{2.297141in}{0.880461in}}%
\pgfpathlineto{\pgfqpoint{2.297645in}{0.883786in}}%
\pgfpathlineto{\pgfqpoint{2.307139in}{0.891576in}}%
\pgfpathlineto{\pgfqpoint{2.306319in}{0.899456in}}%
\pgfpathlineto{\pgfqpoint{2.341651in}{0.899996in}}%
\pgfpathlineto{\pgfqpoint{2.342635in}{0.869578in}}%
\pgfpathclose%
\pgfusepath{fill}%
\end{pgfscope}%
\begin{pgfscope}%
\pgfpathrectangle{\pgfqpoint{0.100000in}{0.100000in}}{\pgfqpoint{3.608454in}{2.310000in}}%
\pgfusepath{clip}%
\pgfsetbuttcap%
\pgfsetmiterjoin%
\definecolor{currentfill}{rgb}{0.000000,0.403922,0.798039}%
\pgfsetfillcolor{currentfill}%
\pgfsetlinewidth{0.000000pt}%
\definecolor{currentstroke}{rgb}{0.000000,0.000000,0.000000}%
\pgfsetstrokecolor{currentstroke}%
\pgfsetstrokeopacity{0.000000}%
\pgfsetdash{}{0pt}%
\pgfpathmoveto{\pgfqpoint{1.726904in}{1.542427in}}%
\pgfpathlineto{\pgfqpoint{1.692731in}{1.544894in}}%
\pgfpathlineto{\pgfqpoint{1.694791in}{1.572316in}}%
\pgfpathlineto{\pgfqpoint{1.689939in}{1.572676in}}%
\pgfpathlineto{\pgfqpoint{1.691656in}{1.593744in}}%
\pgfpathlineto{\pgfqpoint{1.693738in}{1.600532in}}%
\pgfpathlineto{\pgfqpoint{1.695984in}{1.627874in}}%
\pgfpathlineto{\pgfqpoint{1.694322in}{1.628015in}}%
\pgfpathlineto{\pgfqpoint{1.696414in}{1.655020in}}%
\pgfpathlineto{\pgfqpoint{1.696825in}{1.671982in}}%
\pgfpathlineto{\pgfqpoint{1.746207in}{1.668241in}}%
\pgfpathlineto{\pgfqpoint{1.805743in}{1.664546in}}%
\pgfpathlineto{\pgfqpoint{1.805638in}{1.647292in}}%
\pgfpathlineto{\pgfqpoint{1.804075in}{1.619812in}}%
\pgfpathlineto{\pgfqpoint{1.803419in}{1.592288in}}%
\pgfpathlineto{\pgfqpoint{1.763649in}{1.594885in}}%
\pgfpathlineto{\pgfqpoint{1.762096in}{1.567407in}}%
\pgfpathlineto{\pgfqpoint{1.728889in}{1.569871in}}%
\pgfpathlineto{\pgfqpoint{1.726904in}{1.542427in}}%
\pgfpathclose%
\pgfusepath{fill}%
\end{pgfscope}%
\begin{pgfscope}%
\pgfpathrectangle{\pgfqpoint{0.100000in}{0.100000in}}{\pgfqpoint{3.608454in}{2.310000in}}%
\pgfusepath{clip}%
\pgfsetbuttcap%
\pgfsetmiterjoin%
\definecolor{currentfill}{rgb}{0.000000,0.639216,0.680392}%
\pgfsetfillcolor{currentfill}%
\pgfsetlinewidth{0.000000pt}%
\definecolor{currentstroke}{rgb}{0.000000,0.000000,0.000000}%
\pgfsetstrokecolor{currentstroke}%
\pgfsetstrokeopacity{0.000000}%
\pgfsetdash{}{0pt}%
\pgfpathmoveto{\pgfqpoint{2.731544in}{1.590059in}}%
\pgfpathlineto{\pgfqpoint{2.707766in}{1.586556in}}%
\pgfpathlineto{\pgfqpoint{2.706079in}{1.591455in}}%
\pgfpathlineto{\pgfqpoint{2.703247in}{1.616030in}}%
\pgfpathlineto{\pgfqpoint{2.710082in}{1.616628in}}%
\pgfpathlineto{\pgfqpoint{2.706426in}{1.644227in}}%
\pgfpathlineto{\pgfqpoint{2.739935in}{1.648352in}}%
\pgfpathlineto{\pgfqpoint{2.743822in}{1.620714in}}%
\pgfpathlineto{\pgfqpoint{2.764552in}{1.624067in}}%
\pgfpathlineto{\pgfqpoint{2.768691in}{1.595961in}}%
\pgfpathlineto{\pgfqpoint{2.731544in}{1.590059in}}%
\pgfpathclose%
\pgfusepath{fill}%
\end{pgfscope}%
\begin{pgfscope}%
\pgfpathrectangle{\pgfqpoint{0.100000in}{0.100000in}}{\pgfqpoint{3.608454in}{2.310000in}}%
\pgfusepath{clip}%
\pgfsetbuttcap%
\pgfsetmiterjoin%
\definecolor{currentfill}{rgb}{0.000000,0.470588,0.764706}%
\pgfsetfillcolor{currentfill}%
\pgfsetlinewidth{0.000000pt}%
\definecolor{currentstroke}{rgb}{0.000000,0.000000,0.000000}%
\pgfsetstrokecolor{currentstroke}%
\pgfsetstrokeopacity{0.000000}%
\pgfsetdash{}{0pt}%
\pgfpathmoveto{\pgfqpoint{3.269282in}{1.579037in}}%
\pgfpathlineto{\pgfqpoint{3.244201in}{1.592250in}}%
\pgfpathlineto{\pgfqpoint{3.251074in}{1.598426in}}%
\pgfpathlineto{\pgfqpoint{3.234122in}{1.609162in}}%
\pgfpathlineto{\pgfqpoint{3.248038in}{1.619832in}}%
\pgfpathlineto{\pgfqpoint{3.244631in}{1.624617in}}%
\pgfpathlineto{\pgfqpoint{3.248335in}{1.629894in}}%
\pgfpathlineto{\pgfqpoint{3.252558in}{1.629161in}}%
\pgfpathlineto{\pgfqpoint{3.262436in}{1.621007in}}%
\pgfpathlineto{\pgfqpoint{3.258629in}{1.617098in}}%
\pgfpathlineto{\pgfqpoint{3.266342in}{1.608149in}}%
\pgfpathlineto{\pgfqpoint{3.276398in}{1.614345in}}%
\pgfpathlineto{\pgfqpoint{3.284382in}{1.624557in}}%
\pgfpathlineto{\pgfqpoint{3.289903in}{1.618119in}}%
\pgfpathlineto{\pgfqpoint{3.288120in}{1.611032in}}%
\pgfpathlineto{\pgfqpoint{3.284662in}{1.609183in}}%
\pgfpathlineto{\pgfqpoint{3.286080in}{1.595797in}}%
\pgfpathlineto{\pgfqpoint{3.275573in}{1.583652in}}%
\pgfpathlineto{\pgfqpoint{3.269282in}{1.579037in}}%
\pgfpathclose%
\pgfusepath{fill}%
\end{pgfscope}%
\begin{pgfscope}%
\pgfpathrectangle{\pgfqpoint{0.100000in}{0.100000in}}{\pgfqpoint{3.608454in}{2.310000in}}%
\pgfusepath{clip}%
\pgfsetbuttcap%
\pgfsetmiterjoin%
\definecolor{currentfill}{rgb}{0.000000,0.556863,0.721569}%
\pgfsetfillcolor{currentfill}%
\pgfsetlinewidth{0.000000pt}%
\definecolor{currentstroke}{rgb}{0.000000,0.000000,0.000000}%
\pgfsetstrokecolor{currentstroke}%
\pgfsetstrokeopacity{0.000000}%
\pgfsetdash{}{0pt}%
\pgfpathmoveto{\pgfqpoint{1.826743in}{0.706842in}}%
\pgfpathlineto{\pgfqpoint{1.835408in}{0.704889in}}%
\pgfpathlineto{\pgfqpoint{1.833962in}{0.671147in}}%
\pgfpathlineto{\pgfqpoint{1.810600in}{0.672318in}}%
\pgfpathlineto{\pgfqpoint{1.798187in}{0.672904in}}%
\pgfpathlineto{\pgfqpoint{1.799907in}{0.708077in}}%
\pgfpathlineto{\pgfqpoint{1.826743in}{0.706842in}}%
\pgfpathclose%
\pgfusepath{fill}%
\end{pgfscope}%
\begin{pgfscope}%
\pgfpathrectangle{\pgfqpoint{0.100000in}{0.100000in}}{\pgfqpoint{3.608454in}{2.310000in}}%
\pgfusepath{clip}%
\pgfsetbuttcap%
\pgfsetmiterjoin%
\definecolor{currentfill}{rgb}{0.000000,0.627451,0.686275}%
\pgfsetfillcolor{currentfill}%
\pgfsetlinewidth{0.000000pt}%
\definecolor{currentstroke}{rgb}{0.000000,0.000000,0.000000}%
\pgfsetstrokecolor{currentstroke}%
\pgfsetstrokeopacity{0.000000}%
\pgfsetdash{}{0pt}%
\pgfpathmoveto{\pgfqpoint{2.236527in}{1.423164in}}%
\pgfpathlineto{\pgfqpoint{2.205934in}{1.422201in}}%
\pgfpathlineto{\pgfqpoint{2.204405in}{1.465435in}}%
\pgfpathlineto{\pgfqpoint{2.244011in}{1.467014in}}%
\pgfpathlineto{\pgfqpoint{2.246988in}{1.453136in}}%
\pgfpathlineto{\pgfqpoint{2.246319in}{1.447619in}}%
\pgfpathlineto{\pgfqpoint{2.235992in}{1.447300in}}%
\pgfpathlineto{\pgfqpoint{2.236527in}{1.423164in}}%
\pgfpathclose%
\pgfusepath{fill}%
\end{pgfscope}%
\begin{pgfscope}%
\pgfpathrectangle{\pgfqpoint{0.100000in}{0.100000in}}{\pgfqpoint{3.608454in}{2.310000in}}%
\pgfusepath{clip}%
\pgfsetbuttcap%
\pgfsetmiterjoin%
\definecolor{currentfill}{rgb}{0.000000,0.513725,0.743137}%
\pgfsetfillcolor{currentfill}%
\pgfsetlinewidth{0.000000pt}%
\definecolor{currentstroke}{rgb}{0.000000,0.000000,0.000000}%
\pgfsetstrokecolor{currentstroke}%
\pgfsetstrokeopacity{0.000000}%
\pgfsetdash{}{0pt}%
\pgfpathmoveto{\pgfqpoint{3.115324in}{1.315206in}}%
\pgfpathlineto{\pgfqpoint{3.110914in}{1.313734in}}%
\pgfpathlineto{\pgfqpoint{3.103381in}{1.324700in}}%
\pgfpathlineto{\pgfqpoint{3.098792in}{1.325084in}}%
\pgfpathlineto{\pgfqpoint{3.092431in}{1.332135in}}%
\pgfpathlineto{\pgfqpoint{3.084921in}{1.329255in}}%
\pgfpathlineto{\pgfqpoint{3.081757in}{1.317949in}}%
\pgfpathlineto{\pgfqpoint{3.075683in}{1.311970in}}%
\pgfpathlineto{\pgfqpoint{3.074021in}{1.306357in}}%
\pgfpathlineto{\pgfqpoint{3.060878in}{1.314437in}}%
\pgfpathlineto{\pgfqpoint{3.060124in}{1.322440in}}%
\pgfpathlineto{\pgfqpoint{3.050832in}{1.323908in}}%
\pgfpathlineto{\pgfqpoint{3.045914in}{1.315133in}}%
\pgfpathlineto{\pgfqpoint{3.033267in}{1.304579in}}%
\pgfpathlineto{\pgfqpoint{3.022678in}{1.310484in}}%
\pgfpathlineto{\pgfqpoint{3.025954in}{1.320474in}}%
\pgfpathlineto{\pgfqpoint{3.037059in}{1.337144in}}%
\pgfpathlineto{\pgfqpoint{3.038592in}{1.342913in}}%
\pgfpathlineto{\pgfqpoint{3.039437in}{1.352438in}}%
\pgfpathlineto{\pgfqpoint{3.046369in}{1.360627in}}%
\pgfpathlineto{\pgfqpoint{3.044609in}{1.363268in}}%
\pgfpathlineto{\pgfqpoint{3.050299in}{1.374365in}}%
\pgfpathlineto{\pgfqpoint{3.050650in}{1.387261in}}%
\pgfpathlineto{\pgfqpoint{3.058026in}{1.385268in}}%
\pgfpathlineto{\pgfqpoint{3.062954in}{1.378601in}}%
\pgfpathlineto{\pgfqpoint{3.073594in}{1.376808in}}%
\pgfpathlineto{\pgfqpoint{3.078058in}{1.382882in}}%
\pgfpathlineto{\pgfqpoint{3.110807in}{1.366936in}}%
\pgfpathlineto{\pgfqpoint{3.110708in}{1.357752in}}%
\pgfpathlineto{\pgfqpoint{3.107539in}{1.353570in}}%
\pgfpathlineto{\pgfqpoint{3.120128in}{1.335691in}}%
\pgfpathlineto{\pgfqpoint{3.121341in}{1.327517in}}%
\pgfpathlineto{\pgfqpoint{3.116356in}{1.322582in}}%
\pgfpathlineto{\pgfqpoint{3.115324in}{1.315206in}}%
\pgfpathclose%
\pgfusepath{fill}%
\end{pgfscope}%
\begin{pgfscope}%
\pgfpathrectangle{\pgfqpoint{0.100000in}{0.100000in}}{\pgfqpoint{3.608454in}{2.310000in}}%
\pgfusepath{clip}%
\pgfsetbuttcap%
\pgfsetmiterjoin%
\definecolor{currentfill}{rgb}{0.000000,0.556863,0.721569}%
\pgfsetfillcolor{currentfill}%
\pgfsetlinewidth{0.000000pt}%
\definecolor{currentstroke}{rgb}{0.000000,0.000000,0.000000}%
\pgfsetstrokecolor{currentstroke}%
\pgfsetstrokeopacity{0.000000}%
\pgfsetdash{}{0pt}%
\pgfpathmoveto{\pgfqpoint{1.810600in}{0.672318in}}%
\pgfpathlineto{\pgfqpoint{1.809932in}{0.655294in}}%
\pgfpathlineto{\pgfqpoint{1.778772in}{0.657165in}}%
\pgfpathlineto{\pgfqpoint{1.753800in}{0.658495in}}%
\pgfpathlineto{\pgfqpoint{1.755704in}{0.692007in}}%
\pgfpathlineto{\pgfqpoint{1.757470in}{0.722100in}}%
\pgfpathlineto{\pgfqpoint{1.792400in}{0.720174in}}%
\pgfpathlineto{\pgfqpoint{1.791780in}{0.708494in}}%
\pgfpathlineto{\pgfqpoint{1.799907in}{0.708077in}}%
\pgfpathlineto{\pgfqpoint{1.798187in}{0.672904in}}%
\pgfpathlineto{\pgfqpoint{1.810600in}{0.672318in}}%
\pgfpathclose%
\pgfusepath{fill}%
\end{pgfscope}%
\begin{pgfscope}%
\pgfpathrectangle{\pgfqpoint{0.100000in}{0.100000in}}{\pgfqpoint{3.608454in}{2.310000in}}%
\pgfusepath{clip}%
\pgfsetbuttcap%
\pgfsetmiterjoin%
\definecolor{currentfill}{rgb}{0.000000,0.521569,0.739216}%
\pgfsetfillcolor{currentfill}%
\pgfsetlinewidth{0.000000pt}%
\definecolor{currentstroke}{rgb}{0.000000,0.000000,0.000000}%
\pgfsetstrokecolor{currentstroke}%
\pgfsetstrokeopacity{0.000000}%
\pgfsetdash{}{0pt}%
\pgfpathmoveto{\pgfqpoint{3.185535in}{1.030022in}}%
\pgfpathlineto{\pgfqpoint{3.145061in}{1.058883in}}%
\pgfpathlineto{\pgfqpoint{3.150261in}{1.071349in}}%
\pgfpathlineto{\pgfqpoint{3.156843in}{1.077749in}}%
\pgfpathlineto{\pgfqpoint{3.156648in}{1.092471in}}%
\pgfpathlineto{\pgfqpoint{3.148549in}{1.102631in}}%
\pgfpathlineto{\pgfqpoint{3.154787in}{1.104979in}}%
\pgfpathlineto{\pgfqpoint{3.174249in}{1.109020in}}%
\pgfpathlineto{\pgfqpoint{3.182406in}{1.101945in}}%
\pgfpathlineto{\pgfqpoint{3.194034in}{1.089009in}}%
\pgfpathlineto{\pgfqpoint{3.200611in}{1.103096in}}%
\pgfpathlineto{\pgfqpoint{3.228319in}{1.108263in}}%
\pgfpathlineto{\pgfqpoint{3.237815in}{1.092298in}}%
\pgfpathlineto{\pgfqpoint{3.243003in}{1.088398in}}%
\pgfpathlineto{\pgfqpoint{3.231067in}{1.072331in}}%
\pgfpathlineto{\pgfqpoint{3.226781in}{1.061809in}}%
\pgfpathlineto{\pgfqpoint{3.223101in}{1.039243in}}%
\pgfpathlineto{\pgfqpoint{3.209280in}{1.039463in}}%
\pgfpathlineto{\pgfqpoint{3.196751in}{1.036307in}}%
\pgfpathlineto{\pgfqpoint{3.185535in}{1.030022in}}%
\pgfpathclose%
\pgfusepath{fill}%
\end{pgfscope}%
\begin{pgfscope}%
\pgfpathrectangle{\pgfqpoint{0.100000in}{0.100000in}}{\pgfqpoint{3.608454in}{2.310000in}}%
\pgfusepath{clip}%
\pgfsetbuttcap%
\pgfsetmiterjoin%
\definecolor{currentfill}{rgb}{0.000000,0.490196,0.754902}%
\pgfsetfillcolor{currentfill}%
\pgfsetlinewidth{0.000000pt}%
\definecolor{currentstroke}{rgb}{0.000000,0.000000,0.000000}%
\pgfsetstrokecolor{currentstroke}%
\pgfsetstrokeopacity{0.000000}%
\pgfsetdash{}{0pt}%
\pgfpathmoveto{\pgfqpoint{0.628251in}{1.491065in}}%
\pgfpathlineto{\pgfqpoint{0.620520in}{1.503015in}}%
\pgfpathlineto{\pgfqpoint{0.623308in}{1.518183in}}%
\pgfpathlineto{\pgfqpoint{0.620041in}{1.519423in}}%
\pgfpathlineto{\pgfqpoint{0.623821in}{1.536529in}}%
\pgfpathlineto{\pgfqpoint{0.631100in}{1.537480in}}%
\pgfpathlineto{\pgfqpoint{0.632981in}{1.544578in}}%
\pgfpathlineto{\pgfqpoint{0.606969in}{1.551881in}}%
\pgfpathlineto{\pgfqpoint{0.607050in}{1.554247in}}%
\pgfpathlineto{\pgfqpoint{0.592721in}{1.558163in}}%
\pgfpathlineto{\pgfqpoint{0.599878in}{1.583444in}}%
\pgfpathlineto{\pgfqpoint{0.605938in}{1.604509in}}%
\pgfpathlineto{\pgfqpoint{0.622947in}{1.663573in}}%
\pgfpathlineto{\pgfqpoint{0.637413in}{1.715698in}}%
\pgfpathlineto{\pgfqpoint{0.654872in}{1.777406in}}%
\pgfpathlineto{\pgfqpoint{0.691075in}{1.767312in}}%
\pgfpathlineto{\pgfqpoint{0.693113in}{1.766738in}}%
\pgfpathlineto{\pgfqpoint{0.681032in}{1.722500in}}%
\pgfpathlineto{\pgfqpoint{0.671164in}{1.687956in}}%
\pgfpathlineto{\pgfqpoint{0.672323in}{1.687633in}}%
\pgfpathlineto{\pgfqpoint{0.662884in}{1.653976in}}%
\pgfpathlineto{\pgfqpoint{0.661398in}{1.654341in}}%
\pgfpathlineto{\pgfqpoint{0.650489in}{1.614771in}}%
\pgfpathlineto{\pgfqpoint{0.693366in}{1.603165in}}%
\pgfpathlineto{\pgfqpoint{0.755783in}{1.587192in}}%
\pgfpathlineto{\pgfqpoint{0.758589in}{1.580048in}}%
\pgfpathlineto{\pgfqpoint{0.753615in}{1.567143in}}%
\pgfpathlineto{\pgfqpoint{0.753176in}{1.557006in}}%
\pgfpathlineto{\pgfqpoint{0.749642in}{1.550822in}}%
\pgfpathlineto{\pgfqpoint{0.738497in}{1.547647in}}%
\pgfpathlineto{\pgfqpoint{0.730097in}{1.540859in}}%
\pgfpathlineto{\pgfqpoint{0.730363in}{1.531286in}}%
\pgfpathlineto{\pgfqpoint{0.725481in}{1.527601in}}%
\pgfpathlineto{\pgfqpoint{0.724276in}{1.521184in}}%
\pgfpathlineto{\pgfqpoint{0.718561in}{1.521037in}}%
\pgfpathlineto{\pgfqpoint{0.655392in}{1.537674in}}%
\pgfpathlineto{\pgfqpoint{0.647659in}{1.529088in}}%
\pgfpathlineto{\pgfqpoint{0.645777in}{1.522062in}}%
\pgfpathlineto{\pgfqpoint{0.652474in}{1.520325in}}%
\pgfpathlineto{\pgfqpoint{0.643201in}{1.487011in}}%
\pgfpathlineto{\pgfqpoint{0.628251in}{1.491065in}}%
\pgfpathclose%
\pgfusepath{fill}%
\end{pgfscope}%
\begin{pgfscope}%
\pgfpathrectangle{\pgfqpoint{0.100000in}{0.100000in}}{\pgfqpoint{3.608454in}{2.310000in}}%
\pgfusepath{clip}%
\pgfsetbuttcap%
\pgfsetmiterjoin%
\definecolor{currentfill}{rgb}{0.000000,0.470588,0.764706}%
\pgfsetfillcolor{currentfill}%
\pgfsetlinewidth{0.000000pt}%
\definecolor{currentstroke}{rgb}{0.000000,0.000000,0.000000}%
\pgfsetstrokecolor{currentstroke}%
\pgfsetstrokeopacity{0.000000}%
\pgfsetdash{}{0pt}%
\pgfpathmoveto{\pgfqpoint{1.889333in}{1.734748in}}%
\pgfpathlineto{\pgfqpoint{1.882514in}{1.735032in}}%
\pgfpathlineto{\pgfqpoint{1.883458in}{1.755759in}}%
\pgfpathlineto{\pgfqpoint{1.896283in}{1.755231in}}%
\pgfpathlineto{\pgfqpoint{1.897415in}{1.789939in}}%
\pgfpathlineto{\pgfqpoint{1.938531in}{1.788233in}}%
\pgfpathlineto{\pgfqpoint{1.945390in}{1.787969in}}%
\pgfpathlineto{\pgfqpoint{1.945343in}{1.781068in}}%
\pgfpathlineto{\pgfqpoint{1.944449in}{1.753368in}}%
\pgfpathlineto{\pgfqpoint{1.917260in}{1.754415in}}%
\pgfpathlineto{\pgfqpoint{1.916523in}{1.733891in}}%
\pgfpathlineto{\pgfqpoint{1.889333in}{1.734748in}}%
\pgfpathclose%
\pgfusepath{fill}%
\end{pgfscope}%
\begin{pgfscope}%
\pgfpathrectangle{\pgfqpoint{0.100000in}{0.100000in}}{\pgfqpoint{3.608454in}{2.310000in}}%
\pgfusepath{clip}%
\pgfsetbuttcap%
\pgfsetmiterjoin%
\definecolor{currentfill}{rgb}{0.000000,0.647059,0.676471}%
\pgfsetfillcolor{currentfill}%
\pgfsetlinewidth{0.000000pt}%
\definecolor{currentstroke}{rgb}{0.000000,0.000000,0.000000}%
\pgfsetstrokecolor{currentstroke}%
\pgfsetstrokeopacity{0.000000}%
\pgfsetdash{}{0pt}%
\pgfpathmoveto{\pgfqpoint{3.262571in}{1.349107in}}%
\pgfpathlineto{\pgfqpoint{3.255578in}{1.348500in}}%
\pgfpathlineto{\pgfqpoint{3.250135in}{1.356031in}}%
\pgfpathlineto{\pgfqpoint{3.244523in}{1.358215in}}%
\pgfpathlineto{\pgfqpoint{3.240232in}{1.353321in}}%
\pgfpathlineto{\pgfqpoint{3.236077in}{1.357801in}}%
\pgfpathlineto{\pgfqpoint{3.228182in}{1.355944in}}%
\pgfpathlineto{\pgfqpoint{3.222015in}{1.359018in}}%
\pgfpathlineto{\pgfqpoint{3.217018in}{1.367687in}}%
\pgfpathlineto{\pgfqpoint{3.209188in}{1.375295in}}%
\pgfpathlineto{\pgfqpoint{3.214965in}{1.384823in}}%
\pgfpathlineto{\pgfqpoint{3.216571in}{1.392855in}}%
\pgfpathlineto{\pgfqpoint{3.225674in}{1.386154in}}%
\pgfpathlineto{\pgfqpoint{3.242132in}{1.387601in}}%
\pgfpathlineto{\pgfqpoint{3.253966in}{1.377009in}}%
\pgfpathlineto{\pgfqpoint{3.269135in}{1.372212in}}%
\pgfpathlineto{\pgfqpoint{3.266819in}{1.363023in}}%
\pgfpathlineto{\pgfqpoint{3.268336in}{1.356399in}}%
\pgfpathlineto{\pgfqpoint{3.266028in}{1.348977in}}%
\pgfpathlineto{\pgfqpoint{3.262571in}{1.349107in}}%
\pgfpathclose%
\pgfusepath{fill}%
\end{pgfscope}%
\begin{pgfscope}%
\pgfpathrectangle{\pgfqpoint{0.100000in}{0.100000in}}{\pgfqpoint{3.608454in}{2.310000in}}%
\pgfusepath{clip}%
\pgfsetbuttcap%
\pgfsetmiterjoin%
\definecolor{currentfill}{rgb}{0.000000,0.498039,0.750980}%
\pgfsetfillcolor{currentfill}%
\pgfsetlinewidth{0.000000pt}%
\definecolor{currentstroke}{rgb}{0.000000,0.000000,0.000000}%
\pgfsetstrokecolor{currentstroke}%
\pgfsetstrokeopacity{0.000000}%
\pgfsetdash{}{0pt}%
\pgfpathmoveto{\pgfqpoint{1.938292in}{1.256226in}}%
\pgfpathlineto{\pgfqpoint{1.938084in}{1.242176in}}%
\pgfpathlineto{\pgfqpoint{1.890205in}{1.243605in}}%
\pgfpathlineto{\pgfqpoint{1.890012in}{1.250882in}}%
\pgfpathlineto{\pgfqpoint{1.891797in}{1.319820in}}%
\pgfpathlineto{\pgfqpoint{1.892401in}{1.333545in}}%
\pgfpathlineto{\pgfqpoint{1.926584in}{1.332368in}}%
\pgfpathlineto{\pgfqpoint{1.926190in}{1.311718in}}%
\pgfpathlineto{\pgfqpoint{1.925218in}{1.277253in}}%
\pgfpathlineto{\pgfqpoint{1.938903in}{1.276826in}}%
\pgfpathlineto{\pgfqpoint{1.938292in}{1.256226in}}%
\pgfpathclose%
\pgfusepath{fill}%
\end{pgfscope}%
\begin{pgfscope}%
\pgfpathrectangle{\pgfqpoint{0.100000in}{0.100000in}}{\pgfqpoint{3.608454in}{2.310000in}}%
\pgfusepath{clip}%
\pgfsetbuttcap%
\pgfsetmiterjoin%
\definecolor{currentfill}{rgb}{0.000000,0.631373,0.684314}%
\pgfsetfillcolor{currentfill}%
\pgfsetlinewidth{0.000000pt}%
\definecolor{currentstroke}{rgb}{0.000000,0.000000,0.000000}%
\pgfsetstrokecolor{currentstroke}%
\pgfsetstrokeopacity{0.000000}%
\pgfsetdash{}{0pt}%
\pgfpathmoveto{\pgfqpoint{2.856403in}{0.756814in}}%
\pgfpathlineto{\pgfqpoint{2.831112in}{0.753917in}}%
\pgfpathlineto{\pgfqpoint{2.819817in}{0.752621in}}%
\pgfpathlineto{\pgfqpoint{2.818573in}{0.766684in}}%
\pgfpathlineto{\pgfqpoint{2.811462in}{0.766109in}}%
\pgfpathlineto{\pgfqpoint{2.810152in}{0.780002in}}%
\pgfpathlineto{\pgfqpoint{2.820126in}{0.780667in}}%
\pgfpathlineto{\pgfqpoint{2.825145in}{0.783605in}}%
\pgfpathlineto{\pgfqpoint{2.821443in}{0.790502in}}%
\pgfpathlineto{\pgfqpoint{2.821071in}{0.796344in}}%
\pgfpathlineto{\pgfqpoint{2.814786in}{0.800405in}}%
\pgfpathlineto{\pgfqpoint{2.809408in}{0.807037in}}%
\pgfpathlineto{\pgfqpoint{2.808283in}{0.818788in}}%
\pgfpathlineto{\pgfqpoint{2.853827in}{0.823500in}}%
\pgfpathlineto{\pgfqpoint{2.854299in}{0.818491in}}%
\pgfpathlineto{\pgfqpoint{2.862925in}{0.815951in}}%
\pgfpathlineto{\pgfqpoint{2.865088in}{0.801722in}}%
\pgfpathlineto{\pgfqpoint{2.870385in}{0.802308in}}%
\pgfpathlineto{\pgfqpoint{2.875554in}{0.798633in}}%
\pgfpathlineto{\pgfqpoint{2.877599in}{0.779820in}}%
\pgfpathlineto{\pgfqpoint{2.854226in}{0.777256in}}%
\pgfpathlineto{\pgfqpoint{2.856403in}{0.756814in}}%
\pgfpathclose%
\pgfusepath{fill}%
\end{pgfscope}%
\begin{pgfscope}%
\pgfpathrectangle{\pgfqpoint{0.100000in}{0.100000in}}{\pgfqpoint{3.608454in}{2.310000in}}%
\pgfusepath{clip}%
\pgfsetbuttcap%
\pgfsetmiterjoin%
\definecolor{currentfill}{rgb}{0.000000,0.619608,0.690196}%
\pgfsetfillcolor{currentfill}%
\pgfsetlinewidth{0.000000pt}%
\definecolor{currentstroke}{rgb}{0.000000,0.000000,0.000000}%
\pgfsetstrokecolor{currentstroke}%
\pgfsetstrokeopacity{0.000000}%
\pgfsetdash{}{0pt}%
\pgfpathmoveto{\pgfqpoint{2.631561in}{1.887489in}}%
\pgfpathlineto{\pgfqpoint{2.627523in}{1.889173in}}%
\pgfpathlineto{\pgfqpoint{2.629695in}{1.902348in}}%
\pgfpathlineto{\pgfqpoint{2.633726in}{1.900071in}}%
\pgfpathlineto{\pgfqpoint{2.635138in}{1.890834in}}%
\pgfpathlineto{\pgfqpoint{2.631561in}{1.887489in}}%
\pgfpathclose%
\pgfusepath{fill}%
\end{pgfscope}%
\begin{pgfscope}%
\pgfpathrectangle{\pgfqpoint{0.100000in}{0.100000in}}{\pgfqpoint{3.608454in}{2.310000in}}%
\pgfusepath{clip}%
\pgfsetbuttcap%
\pgfsetmiterjoin%
\definecolor{currentfill}{rgb}{0.000000,0.619608,0.690196}%
\pgfsetfillcolor{currentfill}%
\pgfsetlinewidth{0.000000pt}%
\definecolor{currentstroke}{rgb}{0.000000,0.000000,0.000000}%
\pgfsetstrokecolor{currentstroke}%
\pgfsetstrokeopacity{0.000000}%
\pgfsetdash{}{0pt}%
\pgfpathmoveto{\pgfqpoint{2.706092in}{1.900170in}}%
\pgfpathlineto{\pgfqpoint{2.703507in}{1.899857in}}%
\pgfpathlineto{\pgfqpoint{2.707810in}{1.866142in}}%
\pgfpathlineto{\pgfqpoint{2.680662in}{1.863242in}}%
\pgfpathlineto{\pgfqpoint{2.681411in}{1.856342in}}%
\pgfpathlineto{\pgfqpoint{2.674595in}{1.855705in}}%
\pgfpathlineto{\pgfqpoint{2.654113in}{1.853564in}}%
\pgfpathlineto{\pgfqpoint{2.653419in}{1.860520in}}%
\pgfpathlineto{\pgfqpoint{2.644203in}{1.859587in}}%
\pgfpathlineto{\pgfqpoint{2.648170in}{1.869172in}}%
\pgfpathlineto{\pgfqpoint{2.656698in}{1.874267in}}%
\pgfpathlineto{\pgfqpoint{2.661603in}{1.879931in}}%
\pgfpathlineto{\pgfqpoint{2.658078in}{1.883394in}}%
\pgfpathlineto{\pgfqpoint{2.656143in}{1.891002in}}%
\pgfpathlineto{\pgfqpoint{2.664544in}{1.902846in}}%
\pgfpathlineto{\pgfqpoint{2.671993in}{1.905904in}}%
\pgfpathlineto{\pgfqpoint{2.673199in}{1.909457in}}%
\pgfpathlineto{\pgfqpoint{2.691843in}{1.900527in}}%
\pgfpathlineto{\pgfqpoint{2.698822in}{1.902391in}}%
\pgfpathlineto{\pgfqpoint{2.706092in}{1.900170in}}%
\pgfpathclose%
\pgfusepath{fill}%
\end{pgfscope}%
\begin{pgfscope}%
\pgfpathrectangle{\pgfqpoint{0.100000in}{0.100000in}}{\pgfqpoint{3.608454in}{2.310000in}}%
\pgfusepath{clip}%
\pgfsetbuttcap%
\pgfsetmiterjoin%
\definecolor{currentfill}{rgb}{0.000000,0.670588,0.664706}%
\pgfsetfillcolor{currentfill}%
\pgfsetlinewidth{0.000000pt}%
\definecolor{currentstroke}{rgb}{0.000000,0.000000,0.000000}%
\pgfsetstrokecolor{currentstroke}%
\pgfsetstrokeopacity{0.000000}%
\pgfsetdash{}{0pt}%
\pgfpathmoveto{\pgfqpoint{2.878161in}{0.725802in}}%
\pgfpathlineto{\pgfqpoint{2.887245in}{0.726366in}}%
\pgfpathlineto{\pgfqpoint{2.888448in}{0.720934in}}%
\pgfpathlineto{\pgfqpoint{2.880133in}{0.713968in}}%
\pgfpathlineto{\pgfqpoint{2.881876in}{0.708207in}}%
\pgfpathlineto{\pgfqpoint{2.876684in}{0.701966in}}%
\pgfpathlineto{\pgfqpoint{2.878191in}{0.698489in}}%
\pgfpathlineto{\pgfqpoint{2.870453in}{0.691569in}}%
\pgfpathlineto{\pgfqpoint{2.866769in}{0.678489in}}%
\pgfpathlineto{\pgfqpoint{2.854948in}{0.675589in}}%
\pgfpathlineto{\pgfqpoint{2.847384in}{0.676389in}}%
\pgfpathlineto{\pgfqpoint{2.841677in}{0.670279in}}%
\pgfpathlineto{\pgfqpoint{2.844033in}{0.665474in}}%
\pgfpathlineto{\pgfqpoint{2.829477in}{0.667841in}}%
\pgfpathlineto{\pgfqpoint{2.797575in}{0.664470in}}%
\pgfpathlineto{\pgfqpoint{2.788514in}{0.665867in}}%
\pgfpathlineto{\pgfqpoint{2.787376in}{0.676778in}}%
\pgfpathlineto{\pgfqpoint{2.769691in}{0.676534in}}%
\pgfpathlineto{\pgfqpoint{2.766984in}{0.705619in}}%
\pgfpathlineto{\pgfqpoint{2.760006in}{0.704933in}}%
\pgfpathlineto{\pgfqpoint{2.760979in}{0.694551in}}%
\pgfpathlineto{\pgfqpoint{2.735873in}{0.692265in}}%
\pgfpathlineto{\pgfqpoint{2.726751in}{0.687235in}}%
\pgfpathlineto{\pgfqpoint{2.735056in}{0.696150in}}%
\pgfpathlineto{\pgfqpoint{2.730084in}{0.707397in}}%
\pgfpathlineto{\pgfqpoint{2.733506in}{0.712730in}}%
\pgfpathlineto{\pgfqpoint{2.721186in}{0.711980in}}%
\pgfpathlineto{\pgfqpoint{2.719091in}{0.734938in}}%
\pgfpathlineto{\pgfqpoint{2.789288in}{0.742878in}}%
\pgfpathlineto{\pgfqpoint{2.794346in}{0.734429in}}%
\pgfpathlineto{\pgfqpoint{2.795289in}{0.729104in}}%
\pgfpathlineto{\pgfqpoint{2.801217in}{0.720980in}}%
\pgfpathlineto{\pgfqpoint{2.834431in}{0.723057in}}%
\pgfpathlineto{\pgfqpoint{2.878161in}{0.725802in}}%
\pgfpathclose%
\pgfusepath{fill}%
\end{pgfscope}%
\begin{pgfscope}%
\pgfpathrectangle{\pgfqpoint{0.100000in}{0.100000in}}{\pgfqpoint{3.608454in}{2.310000in}}%
\pgfusepath{clip}%
\pgfsetbuttcap%
\pgfsetmiterjoin%
\definecolor{currentfill}{rgb}{0.000000,0.431373,0.784314}%
\pgfsetfillcolor{currentfill}%
\pgfsetlinewidth{0.000000pt}%
\definecolor{currentstroke}{rgb}{0.000000,0.000000,0.000000}%
\pgfsetstrokecolor{currentstroke}%
\pgfsetstrokeopacity{0.000000}%
\pgfsetdash{}{0pt}%
\pgfpathmoveto{\pgfqpoint{1.718524in}{1.959750in}}%
\pgfpathlineto{\pgfqpoint{1.719085in}{1.966639in}}%
\pgfpathlineto{\pgfqpoint{1.739705in}{1.965094in}}%
\pgfpathlineto{\pgfqpoint{1.739184in}{1.958172in}}%
\pgfpathlineto{\pgfqpoint{1.762088in}{1.956478in}}%
\pgfpathlineto{\pgfqpoint{1.760656in}{1.935885in}}%
\pgfpathlineto{\pgfqpoint{1.774399in}{1.934946in}}%
\pgfpathlineto{\pgfqpoint{1.775193in}{1.927930in}}%
\pgfpathlineto{\pgfqpoint{1.759368in}{1.916843in}}%
\pgfpathlineto{\pgfqpoint{1.756437in}{1.912852in}}%
\pgfpathlineto{\pgfqpoint{1.744713in}{1.908214in}}%
\pgfpathlineto{\pgfqpoint{1.730862in}{1.910799in}}%
\pgfpathlineto{\pgfqpoint{1.726123in}{1.914246in}}%
\pgfpathlineto{\pgfqpoint{1.720336in}{1.913435in}}%
\pgfpathlineto{\pgfqpoint{1.721787in}{1.931717in}}%
\pgfpathlineto{\pgfqpoint{1.719106in}{1.931926in}}%
\pgfpathlineto{\pgfqpoint{1.721015in}{1.959541in}}%
\pgfpathlineto{\pgfqpoint{1.718524in}{1.959750in}}%
\pgfpathclose%
\pgfusepath{fill}%
\end{pgfscope}%
\begin{pgfscope}%
\pgfpathrectangle{\pgfqpoint{0.100000in}{0.100000in}}{\pgfqpoint{3.608454in}{2.310000in}}%
\pgfusepath{clip}%
\pgfsetbuttcap%
\pgfsetmiterjoin%
\definecolor{currentfill}{rgb}{0.000000,0.364706,0.817647}%
\pgfsetfillcolor{currentfill}%
\pgfsetlinewidth{0.000000pt}%
\definecolor{currentstroke}{rgb}{0.000000,0.000000,0.000000}%
\pgfsetstrokecolor{currentstroke}%
\pgfsetstrokeopacity{0.000000}%
\pgfsetdash{}{0pt}%
\pgfpathmoveto{\pgfqpoint{1.670021in}{1.298287in}}%
\pgfpathlineto{\pgfqpoint{1.670049in}{1.298787in}}%
\pgfpathlineto{\pgfqpoint{1.672634in}{1.332579in}}%
\pgfpathlineto{\pgfqpoint{1.707133in}{1.330203in}}%
\pgfpathlineto{\pgfqpoint{1.756162in}{1.326912in}}%
\pgfpathlineto{\pgfqpoint{1.754213in}{1.292521in}}%
\pgfpathlineto{\pgfqpoint{1.699539in}{1.296124in}}%
\pgfpathlineto{\pgfqpoint{1.670021in}{1.298287in}}%
\pgfpathclose%
\pgfusepath{fill}%
\end{pgfscope}%
\begin{pgfscope}%
\pgfpathrectangle{\pgfqpoint{0.100000in}{0.100000in}}{\pgfqpoint{3.608454in}{2.310000in}}%
\pgfusepath{clip}%
\pgfsetbuttcap%
\pgfsetmiterjoin%
\definecolor{currentfill}{rgb}{0.000000,0.454902,0.772549}%
\pgfsetfillcolor{currentfill}%
\pgfsetlinewidth{0.000000pt}%
\definecolor{currentstroke}{rgb}{0.000000,0.000000,0.000000}%
\pgfsetstrokecolor{currentstroke}%
\pgfsetstrokeopacity{0.000000}%
\pgfsetdash{}{0pt}%
\pgfpathmoveto{\pgfqpoint{2.084392in}{1.676317in}}%
\pgfpathlineto{\pgfqpoint{2.084271in}{1.648958in}}%
\pgfpathlineto{\pgfqpoint{2.056914in}{1.649093in}}%
\pgfpathlineto{\pgfqpoint{2.017541in}{1.649592in}}%
\pgfpathlineto{\pgfqpoint{2.019889in}{1.653420in}}%
\pgfpathlineto{\pgfqpoint{2.019320in}{1.660716in}}%
\pgfpathlineto{\pgfqpoint{2.022652in}{1.663329in}}%
\pgfpathlineto{\pgfqpoint{2.021565in}{1.674298in}}%
\pgfpathlineto{\pgfqpoint{2.018524in}{1.680571in}}%
\pgfpathlineto{\pgfqpoint{2.019001in}{1.688123in}}%
\pgfpathlineto{\pgfqpoint{2.015077in}{1.691199in}}%
\pgfpathlineto{\pgfqpoint{2.014851in}{1.696524in}}%
\pgfpathlineto{\pgfqpoint{2.023239in}{1.696380in}}%
\pgfpathlineto{\pgfqpoint{2.046244in}{1.696054in}}%
\pgfpathlineto{\pgfqpoint{2.057273in}{1.695912in}}%
\pgfpathlineto{\pgfqpoint{2.057034in}{1.676704in}}%
\pgfpathlineto{\pgfqpoint{2.084392in}{1.676317in}}%
\pgfpathclose%
\pgfusepath{fill}%
\end{pgfscope}%
\begin{pgfscope}%
\pgfpathrectangle{\pgfqpoint{0.100000in}{0.100000in}}{\pgfqpoint{3.608454in}{2.310000in}}%
\pgfusepath{clip}%
\pgfsetbuttcap%
\pgfsetmiterjoin%
\definecolor{currentfill}{rgb}{0.000000,0.439216,0.780392}%
\pgfsetfillcolor{currentfill}%
\pgfsetlinewidth{0.000000pt}%
\definecolor{currentstroke}{rgb}{0.000000,0.000000,0.000000}%
\pgfsetstrokecolor{currentstroke}%
\pgfsetstrokeopacity{0.000000}%
\pgfsetdash{}{0pt}%
\pgfpathmoveto{\pgfqpoint{3.420582in}{1.833115in}}%
\pgfpathlineto{\pgfqpoint{3.422865in}{1.818650in}}%
\pgfpathlineto{\pgfqpoint{3.426046in}{1.819395in}}%
\pgfpathlineto{\pgfqpoint{3.430894in}{1.808006in}}%
\pgfpathlineto{\pgfqpoint{3.434992in}{1.802755in}}%
\pgfpathlineto{\pgfqpoint{3.414849in}{1.798333in}}%
\pgfpathlineto{\pgfqpoint{3.377987in}{1.790511in}}%
\pgfpathlineto{\pgfqpoint{3.374916in}{1.805051in}}%
\pgfpathlineto{\pgfqpoint{3.369972in}{1.804097in}}%
\pgfpathlineto{\pgfqpoint{3.367008in}{1.818791in}}%
\pgfpathlineto{\pgfqpoint{3.374319in}{1.820118in}}%
\pgfpathlineto{\pgfqpoint{3.374364in}{1.831752in}}%
\pgfpathlineto{\pgfqpoint{3.382485in}{1.831082in}}%
\pgfpathlineto{\pgfqpoint{3.396386in}{1.835470in}}%
\pgfpathlineto{\pgfqpoint{3.401069in}{1.830482in}}%
\pgfpathlineto{\pgfqpoint{3.412904in}{1.835488in}}%
\pgfpathlineto{\pgfqpoint{3.420582in}{1.833115in}}%
\pgfpathclose%
\pgfusepath{fill}%
\end{pgfscope}%
\begin{pgfscope}%
\pgfpathrectangle{\pgfqpoint{0.100000in}{0.100000in}}{\pgfqpoint{3.608454in}{2.310000in}}%
\pgfusepath{clip}%
\pgfsetbuttcap%
\pgfsetmiterjoin%
\definecolor{currentfill}{rgb}{0.000000,0.403922,0.798039}%
\pgfsetfillcolor{currentfill}%
\pgfsetlinewidth{0.000000pt}%
\definecolor{currentstroke}{rgb}{0.000000,0.000000,0.000000}%
\pgfsetstrokecolor{currentstroke}%
\pgfsetstrokeopacity{0.000000}%
\pgfsetdash{}{0pt}%
\pgfpathmoveto{\pgfqpoint{1.357331in}{0.522521in}}%
\pgfpathlineto{\pgfqpoint{1.350294in}{0.534978in}}%
\pgfpathlineto{\pgfqpoint{1.341731in}{0.539953in}}%
\pgfpathlineto{\pgfqpoint{1.343780in}{0.546075in}}%
\pgfpathlineto{\pgfqpoint{1.344206in}{0.557855in}}%
\pgfpathlineto{\pgfqpoint{1.348048in}{0.565231in}}%
\pgfpathlineto{\pgfqpoint{1.356655in}{0.558900in}}%
\pgfpathlineto{\pgfqpoint{1.368635in}{0.560958in}}%
\pgfpathlineto{\pgfqpoint{1.371871in}{0.559296in}}%
\pgfpathlineto{\pgfqpoint{1.370602in}{0.544014in}}%
\pgfpathlineto{\pgfqpoint{1.366190in}{0.537636in}}%
\pgfpathlineto{\pgfqpoint{1.369907in}{0.525625in}}%
\pgfpathlineto{\pgfqpoint{1.369026in}{0.517805in}}%
\pgfpathlineto{\pgfqpoint{1.357331in}{0.522521in}}%
\pgfpathclose%
\pgfusepath{fill}%
\end{pgfscope}%
\begin{pgfscope}%
\pgfpathrectangle{\pgfqpoint{0.100000in}{0.100000in}}{\pgfqpoint{3.608454in}{2.310000in}}%
\pgfusepath{clip}%
\pgfsetbuttcap%
\pgfsetmiterjoin%
\definecolor{currentfill}{rgb}{0.000000,0.521569,0.739216}%
\pgfsetfillcolor{currentfill}%
\pgfsetlinewidth{0.000000pt}%
\definecolor{currentstroke}{rgb}{0.000000,0.000000,0.000000}%
\pgfsetstrokecolor{currentstroke}%
\pgfsetstrokeopacity{0.000000}%
\pgfsetdash{}{0pt}%
\pgfpathmoveto{\pgfqpoint{0.575043in}{1.085831in}}%
\pgfpathlineto{\pgfqpoint{0.566660in}{1.090402in}}%
\pgfpathlineto{\pgfqpoint{0.566840in}{1.100543in}}%
\pgfpathlineto{\pgfqpoint{0.574437in}{1.095234in}}%
\pgfpathlineto{\pgfqpoint{0.575043in}{1.085831in}}%
\pgfpathclose%
\pgfusepath{fill}%
\end{pgfscope}%
\begin{pgfscope}%
\pgfpathrectangle{\pgfqpoint{0.100000in}{0.100000in}}{\pgfqpoint{3.608454in}{2.310000in}}%
\pgfusepath{clip}%
\pgfsetbuttcap%
\pgfsetmiterjoin%
\definecolor{currentfill}{rgb}{0.000000,0.521569,0.739216}%
\pgfsetfillcolor{currentfill}%
\pgfsetlinewidth{0.000000pt}%
\definecolor{currentstroke}{rgb}{0.000000,0.000000,0.000000}%
\pgfsetstrokecolor{currentstroke}%
\pgfsetstrokeopacity{0.000000}%
\pgfsetdash{}{0pt}%
\pgfpathmoveto{\pgfqpoint{0.623749in}{1.080645in}}%
\pgfpathlineto{\pgfqpoint{0.614676in}{1.095570in}}%
\pgfpathlineto{\pgfqpoint{0.606778in}{1.102859in}}%
\pgfpathlineto{\pgfqpoint{0.598881in}{1.114997in}}%
\pgfpathlineto{\pgfqpoint{0.593720in}{1.119151in}}%
\pgfpathlineto{\pgfqpoint{0.585255in}{1.116816in}}%
\pgfpathlineto{\pgfqpoint{0.578599in}{1.121274in}}%
\pgfpathlineto{\pgfqpoint{0.581839in}{1.128446in}}%
\pgfpathlineto{\pgfqpoint{0.579418in}{1.141321in}}%
\pgfpathlineto{\pgfqpoint{0.576385in}{1.146196in}}%
\pgfpathlineto{\pgfqpoint{0.563289in}{1.149271in}}%
\pgfpathlineto{\pgfqpoint{0.558328in}{1.148636in}}%
\pgfpathlineto{\pgfqpoint{0.535794in}{1.165997in}}%
\pgfpathlineto{\pgfqpoint{0.534238in}{1.176440in}}%
\pgfpathlineto{\pgfqpoint{0.524045in}{1.187970in}}%
\pgfpathlineto{\pgfqpoint{0.527921in}{1.193369in}}%
\pgfpathlineto{\pgfqpoint{0.537333in}{1.227732in}}%
\pgfpathlineto{\pgfqpoint{0.547247in}{1.223234in}}%
\pgfpathlineto{\pgfqpoint{0.547967in}{1.217656in}}%
\pgfpathlineto{\pgfqpoint{0.598059in}{1.204495in}}%
\pgfpathlineto{\pgfqpoint{0.649520in}{1.191532in}}%
\pgfpathlineto{\pgfqpoint{0.668493in}{1.266395in}}%
\pgfpathlineto{\pgfqpoint{0.692238in}{1.260226in}}%
\pgfpathlineto{\pgfqpoint{0.765878in}{1.242208in}}%
\pgfpathlineto{\pgfqpoint{0.786789in}{1.237324in}}%
\pgfpathlineto{\pgfqpoint{0.792611in}{1.237224in}}%
\pgfpathlineto{\pgfqpoint{0.842429in}{1.160488in}}%
\pgfpathlineto{\pgfqpoint{0.840274in}{1.151795in}}%
\pgfpathlineto{\pgfqpoint{0.843543in}{1.141162in}}%
\pgfpathlineto{\pgfqpoint{0.847819in}{1.135796in}}%
\pgfpathlineto{\pgfqpoint{0.848041in}{1.126325in}}%
\pgfpathlineto{\pgfqpoint{0.851945in}{1.113704in}}%
\pgfpathlineto{\pgfqpoint{0.862231in}{1.099514in}}%
\pgfpathlineto{\pgfqpoint{0.868397in}{1.094701in}}%
\pgfpathlineto{\pgfqpoint{0.895182in}{1.086583in}}%
\pgfpathlineto{\pgfqpoint{0.900636in}{1.091652in}}%
\pgfpathlineto{\pgfqpoint{0.914070in}{1.089859in}}%
\pgfpathlineto{\pgfqpoint{0.908989in}{1.065097in}}%
\pgfpathlineto{\pgfqpoint{0.899087in}{1.016739in}}%
\pgfpathlineto{\pgfqpoint{0.858565in}{1.025184in}}%
\pgfpathlineto{\pgfqpoint{0.859997in}{1.031950in}}%
\pgfpathlineto{\pgfqpoint{0.839911in}{1.036150in}}%
\pgfpathlineto{\pgfqpoint{0.832599in}{1.002476in}}%
\pgfpathlineto{\pgfqpoint{0.807337in}{1.008296in}}%
\pgfpathlineto{\pgfqpoint{0.805414in}{1.014538in}}%
\pgfpathlineto{\pgfqpoint{0.809954in}{1.027225in}}%
\pgfpathlineto{\pgfqpoint{0.809266in}{1.037990in}}%
\pgfpathlineto{\pgfqpoint{0.787832in}{1.044715in}}%
\pgfpathlineto{\pgfqpoint{0.721788in}{1.059750in}}%
\pgfpathlineto{\pgfqpoint{0.660882in}{1.074556in}}%
\pgfpathlineto{\pgfqpoint{0.647432in}{1.078308in}}%
\pgfpathlineto{\pgfqpoint{0.640257in}{1.084919in}}%
\pgfpathlineto{\pgfqpoint{0.631612in}{1.088268in}}%
\pgfpathlineto{\pgfqpoint{0.626202in}{1.085453in}}%
\pgfpathlineto{\pgfqpoint{0.623749in}{1.080645in}}%
\pgfpathclose%
\pgfusepath{fill}%
\end{pgfscope}%
\begin{pgfscope}%
\pgfpathrectangle{\pgfqpoint{0.100000in}{0.100000in}}{\pgfqpoint{3.608454in}{2.310000in}}%
\pgfusepath{clip}%
\pgfsetbuttcap%
\pgfsetmiterjoin%
\definecolor{currentfill}{rgb}{0.000000,0.764706,0.617647}%
\pgfsetfillcolor{currentfill}%
\pgfsetlinewidth{0.000000pt}%
\definecolor{currentstroke}{rgb}{0.000000,0.000000,0.000000}%
\pgfsetstrokecolor{currentstroke}%
\pgfsetstrokeopacity{0.000000}%
\pgfsetdash{}{0pt}%
\pgfpathmoveto{\pgfqpoint{2.669844in}{0.772749in}}%
\pgfpathlineto{\pgfqpoint{2.683521in}{0.774139in}}%
\pgfpathlineto{\pgfqpoint{2.682532in}{0.784446in}}%
\pgfpathlineto{\pgfqpoint{2.685966in}{0.784799in}}%
\pgfpathlineto{\pgfqpoint{2.683683in}{0.809309in}}%
\pgfpathlineto{\pgfqpoint{2.685856in}{0.816401in}}%
\pgfpathlineto{\pgfqpoint{2.692745in}{0.817134in}}%
\pgfpathlineto{\pgfqpoint{2.693497in}{0.810278in}}%
\pgfpathlineto{\pgfqpoint{2.714111in}{0.812491in}}%
\pgfpathlineto{\pgfqpoint{2.713352in}{0.819097in}}%
\pgfpathlineto{\pgfqpoint{2.720311in}{0.819497in}}%
\pgfpathlineto{\pgfqpoint{2.721575in}{0.813214in}}%
\pgfpathlineto{\pgfqpoint{2.727912in}{0.813855in}}%
\pgfpathlineto{\pgfqpoint{2.728545in}{0.806963in}}%
\pgfpathlineto{\pgfqpoint{2.737520in}{0.807875in}}%
\pgfpathlineto{\pgfqpoint{2.737974in}{0.800305in}}%
\pgfpathlineto{\pgfqpoint{2.733549in}{0.786465in}}%
\pgfpathlineto{\pgfqpoint{2.706700in}{0.783758in}}%
\pgfpathlineto{\pgfqpoint{2.704075in}{0.776471in}}%
\pgfpathlineto{\pgfqpoint{2.708723in}{0.733977in}}%
\pgfpathlineto{\pgfqpoint{2.674593in}{0.730790in}}%
\pgfpathlineto{\pgfqpoint{2.673739in}{0.731186in}}%
\pgfpathlineto{\pgfqpoint{2.669844in}{0.772749in}}%
\pgfpathclose%
\pgfusepath{fill}%
\end{pgfscope}%
\begin{pgfscope}%
\pgfpathrectangle{\pgfqpoint{0.100000in}{0.100000in}}{\pgfqpoint{3.608454in}{2.310000in}}%
\pgfusepath{clip}%
\pgfsetbuttcap%
\pgfsetmiterjoin%
\definecolor{currentfill}{rgb}{0.000000,0.431373,0.784314}%
\pgfsetfillcolor{currentfill}%
\pgfsetlinewidth{0.000000pt}%
\definecolor{currentstroke}{rgb}{0.000000,0.000000,0.000000}%
\pgfsetstrokecolor{currentstroke}%
\pgfsetstrokeopacity{0.000000}%
\pgfsetdash{}{0pt}%
\pgfpathmoveto{\pgfqpoint{2.420217in}{1.066185in}}%
\pgfpathlineto{\pgfqpoint{2.427918in}{1.068838in}}%
\pgfpathlineto{\pgfqpoint{2.426249in}{1.075405in}}%
\pgfpathlineto{\pgfqpoint{2.434616in}{1.075007in}}%
\pgfpathlineto{\pgfqpoint{2.441604in}{1.083493in}}%
\pgfpathlineto{\pgfqpoint{2.453807in}{1.083497in}}%
\pgfpathlineto{\pgfqpoint{2.461856in}{1.079471in}}%
\pgfpathlineto{\pgfqpoint{2.463288in}{1.065984in}}%
\pgfpathlineto{\pgfqpoint{2.482016in}{1.066685in}}%
\pgfpathlineto{\pgfqpoint{2.483127in}{1.034794in}}%
\pgfpathlineto{\pgfqpoint{2.473129in}{1.034113in}}%
\pgfpathlineto{\pgfqpoint{2.473863in}{1.022295in}}%
\pgfpathlineto{\pgfqpoint{2.477335in}{1.022510in}}%
\pgfpathlineto{\pgfqpoint{2.478692in}{1.001827in}}%
\pgfpathlineto{\pgfqpoint{2.482530in}{0.995135in}}%
\pgfpathlineto{\pgfqpoint{2.467215in}{0.998833in}}%
\pgfpathlineto{\pgfqpoint{2.420148in}{0.996176in}}%
\pgfpathlineto{\pgfqpoint{2.413440in}{0.992269in}}%
\pgfpathlineto{\pgfqpoint{2.413949in}{0.985540in}}%
\pgfpathlineto{\pgfqpoint{2.407574in}{0.985155in}}%
\pgfpathlineto{\pgfqpoint{2.396088in}{0.992607in}}%
\pgfpathlineto{\pgfqpoint{2.394572in}{0.999746in}}%
\pgfpathlineto{\pgfqpoint{2.401872in}{1.004184in}}%
\pgfpathlineto{\pgfqpoint{2.404445in}{1.017531in}}%
\pgfpathlineto{\pgfqpoint{2.415331in}{1.023934in}}%
\pgfpathlineto{\pgfqpoint{2.412041in}{1.029240in}}%
\pgfpathlineto{\pgfqpoint{2.417154in}{1.033431in}}%
\pgfpathlineto{\pgfqpoint{2.419408in}{1.040872in}}%
\pgfpathlineto{\pgfqpoint{2.424139in}{1.041016in}}%
\pgfpathlineto{\pgfqpoint{2.425407in}{1.049147in}}%
\pgfpathlineto{\pgfqpoint{2.420989in}{1.051604in}}%
\pgfpathlineto{\pgfqpoint{2.424647in}{1.062261in}}%
\pgfpathlineto{\pgfqpoint{2.420217in}{1.066185in}}%
\pgfpathclose%
\pgfusepath{fill}%
\end{pgfscope}%
\begin{pgfscope}%
\pgfpathrectangle{\pgfqpoint{0.100000in}{0.100000in}}{\pgfqpoint{3.608454in}{2.310000in}}%
\pgfusepath{clip}%
\pgfsetbuttcap%
\pgfsetmiterjoin%
\definecolor{currentfill}{rgb}{0.000000,0.937255,0.531373}%
\pgfsetfillcolor{currentfill}%
\pgfsetlinewidth{0.000000pt}%
\definecolor{currentstroke}{rgb}{0.000000,0.000000,0.000000}%
\pgfsetstrokecolor{currentstroke}%
\pgfsetstrokeopacity{0.000000}%
\pgfsetdash{}{0pt}%
\pgfpathmoveto{\pgfqpoint{2.826981in}{1.247494in}}%
\pgfpathlineto{\pgfqpoint{2.816644in}{1.252449in}}%
\pgfpathlineto{\pgfqpoint{2.809988in}{1.252552in}}%
\pgfpathlineto{\pgfqpoint{2.802573in}{1.261875in}}%
\pgfpathlineto{\pgfqpoint{2.800586in}{1.266861in}}%
\pgfpathlineto{\pgfqpoint{2.804559in}{1.276722in}}%
\pgfpathlineto{\pgfqpoint{2.810307in}{1.281879in}}%
\pgfpathlineto{\pgfqpoint{2.814875in}{1.291146in}}%
\pgfpathlineto{\pgfqpoint{2.820760in}{1.291911in}}%
\pgfpathlineto{\pgfqpoint{2.825277in}{1.301789in}}%
\pgfpathlineto{\pgfqpoint{2.824999in}{1.309058in}}%
\pgfpathlineto{\pgfqpoint{2.834982in}{1.312343in}}%
\pgfpathlineto{\pgfqpoint{2.840076in}{1.309133in}}%
\pgfpathlineto{\pgfqpoint{2.850698in}{1.310252in}}%
\pgfpathlineto{\pgfqpoint{2.855369in}{1.302588in}}%
\pgfpathlineto{\pgfqpoint{2.847809in}{1.300549in}}%
\pgfpathlineto{\pgfqpoint{2.837862in}{1.289352in}}%
\pgfpathlineto{\pgfqpoint{2.839446in}{1.282704in}}%
\pgfpathlineto{\pgfqpoint{2.849787in}{1.279563in}}%
\pgfpathlineto{\pgfqpoint{2.860524in}{1.271364in}}%
\pgfpathlineto{\pgfqpoint{2.851560in}{1.271717in}}%
\pgfpathlineto{\pgfqpoint{2.850587in}{1.262241in}}%
\pgfpathlineto{\pgfqpoint{2.841576in}{1.260538in}}%
\pgfpathlineto{\pgfqpoint{2.833788in}{1.254755in}}%
\pgfpathlineto{\pgfqpoint{2.830889in}{1.256941in}}%
\pgfpathlineto{\pgfqpoint{2.825042in}{1.252479in}}%
\pgfpathlineto{\pgfqpoint{2.826981in}{1.247494in}}%
\pgfpathclose%
\pgfusepath{fill}%
\end{pgfscope}%
\begin{pgfscope}%
\pgfpathrectangle{\pgfqpoint{0.100000in}{0.100000in}}{\pgfqpoint{3.608454in}{2.310000in}}%
\pgfusepath{clip}%
\pgfsetbuttcap%
\pgfsetmiterjoin%
\definecolor{currentfill}{rgb}{0.000000,0.678431,0.660784}%
\pgfsetfillcolor{currentfill}%
\pgfsetlinewidth{0.000000pt}%
\definecolor{currentstroke}{rgb}{0.000000,0.000000,0.000000}%
\pgfsetstrokecolor{currentstroke}%
\pgfsetstrokeopacity{0.000000}%
\pgfsetdash{}{0pt}%
\pgfpathmoveto{\pgfqpoint{3.127726in}{0.403873in}}%
\pgfpathlineto{\pgfqpoint{3.133900in}{0.367620in}}%
\pgfpathlineto{\pgfqpoint{3.099940in}{0.362168in}}%
\pgfpathlineto{\pgfqpoint{3.093963in}{0.366222in}}%
\pgfpathlineto{\pgfqpoint{3.076512in}{0.366235in}}%
\pgfpathlineto{\pgfqpoint{3.068767in}{0.369764in}}%
\pgfpathlineto{\pgfqpoint{3.063694in}{0.379993in}}%
\pgfpathlineto{\pgfqpoint{3.058788in}{0.396588in}}%
\pgfpathlineto{\pgfqpoint{3.055285in}{0.402609in}}%
\pgfpathlineto{\pgfqpoint{3.046296in}{0.409678in}}%
\pgfpathlineto{\pgfqpoint{3.037249in}{0.409747in}}%
\pgfpathlineto{\pgfqpoint{3.032422in}{0.421403in}}%
\pgfpathlineto{\pgfqpoint{3.036040in}{0.423473in}}%
\pgfpathlineto{\pgfqpoint{3.037010in}{0.432074in}}%
\pgfpathlineto{\pgfqpoint{3.072400in}{0.437158in}}%
\pgfpathlineto{\pgfqpoint{3.075618in}{0.416692in}}%
\pgfpathlineto{\pgfqpoint{3.096468in}{0.420138in}}%
\pgfpathlineto{\pgfqpoint{3.099914in}{0.399053in}}%
\pgfpathlineto{\pgfqpoint{3.127726in}{0.403873in}}%
\pgfpathclose%
\pgfusepath{fill}%
\end{pgfscope}%
\begin{pgfscope}%
\pgfpathrectangle{\pgfqpoint{0.100000in}{0.100000in}}{\pgfqpoint{3.608454in}{2.310000in}}%
\pgfusepath{clip}%
\pgfsetbuttcap%
\pgfsetmiterjoin%
\definecolor{currentfill}{rgb}{0.000000,0.427451,0.786275}%
\pgfsetfillcolor{currentfill}%
\pgfsetlinewidth{0.000000pt}%
\definecolor{currentstroke}{rgb}{0.000000,0.000000,0.000000}%
\pgfsetstrokecolor{currentstroke}%
\pgfsetstrokeopacity{0.000000}%
\pgfsetdash{}{0pt}%
\pgfpathmoveto{\pgfqpoint{0.508309in}{1.308909in}}%
\pgfpathlineto{\pgfqpoint{0.468241in}{1.320678in}}%
\pgfpathlineto{\pgfqpoint{0.438474in}{1.329831in}}%
\pgfpathlineto{\pgfqpoint{0.433387in}{1.338858in}}%
\pgfpathlineto{\pgfqpoint{0.433303in}{1.348727in}}%
\pgfpathlineto{\pgfqpoint{0.430147in}{1.350732in}}%
\pgfpathlineto{\pgfqpoint{0.425066in}{1.366592in}}%
\pgfpathlineto{\pgfqpoint{0.420021in}{1.371609in}}%
\pgfpathlineto{\pgfqpoint{0.416793in}{1.379069in}}%
\pgfpathlineto{\pgfqpoint{0.418273in}{1.400974in}}%
\pgfpathlineto{\pgfqpoint{0.425399in}{1.400805in}}%
\pgfpathlineto{\pgfqpoint{0.429729in}{1.404759in}}%
\pgfpathlineto{\pgfqpoint{0.434644in}{1.414437in}}%
\pgfpathlineto{\pgfqpoint{0.433004in}{1.425553in}}%
\pgfpathlineto{\pgfqpoint{0.418942in}{1.431643in}}%
\pgfpathlineto{\pgfqpoint{0.413490in}{1.439050in}}%
\pgfpathlineto{\pgfqpoint{0.411162in}{1.446533in}}%
\pgfpathlineto{\pgfqpoint{0.411566in}{1.453040in}}%
\pgfpathlineto{\pgfqpoint{0.422155in}{1.452127in}}%
\pgfpathlineto{\pgfqpoint{0.422616in}{1.464232in}}%
\pgfpathlineto{\pgfqpoint{0.424905in}{1.469119in}}%
\pgfpathlineto{\pgfqpoint{0.432540in}{1.470642in}}%
\pgfpathlineto{\pgfqpoint{0.441374in}{1.466010in}}%
\pgfpathlineto{\pgfqpoint{0.445670in}{1.467188in}}%
\pgfpathlineto{\pgfqpoint{0.469193in}{1.459876in}}%
\pgfpathlineto{\pgfqpoint{0.470704in}{1.451769in}}%
\pgfpathlineto{\pgfqpoint{0.465818in}{1.444533in}}%
\pgfpathlineto{\pgfqpoint{0.465842in}{1.434190in}}%
\pgfpathlineto{\pgfqpoint{0.476039in}{1.431016in}}%
\pgfpathlineto{\pgfqpoint{0.475697in}{1.423778in}}%
\pgfpathlineto{\pgfqpoint{0.471004in}{1.413448in}}%
\pgfpathlineto{\pgfqpoint{0.474622in}{1.404992in}}%
\pgfpathlineto{\pgfqpoint{0.486008in}{1.393700in}}%
\pgfpathlineto{\pgfqpoint{0.499634in}{1.368954in}}%
\pgfpathlineto{\pgfqpoint{0.496557in}{1.356685in}}%
\pgfpathlineto{\pgfqpoint{0.492171in}{1.356794in}}%
\pgfpathlineto{\pgfqpoint{0.488481in}{1.345190in}}%
\pgfpathlineto{\pgfqpoint{0.493971in}{1.332372in}}%
\pgfpathlineto{\pgfqpoint{0.498712in}{1.326448in}}%
\pgfpathlineto{\pgfqpoint{0.503087in}{1.324821in}}%
\pgfpathlineto{\pgfqpoint{0.508309in}{1.308909in}}%
\pgfpathclose%
\pgfusepath{fill}%
\end{pgfscope}%
\begin{pgfscope}%
\pgfpathrectangle{\pgfqpoint{0.100000in}{0.100000in}}{\pgfqpoint{3.608454in}{2.310000in}}%
\pgfusepath{clip}%
\pgfsetbuttcap%
\pgfsetmiterjoin%
\definecolor{currentfill}{rgb}{0.000000,0.364706,0.817647}%
\pgfsetfillcolor{currentfill}%
\pgfsetlinewidth{0.000000pt}%
\definecolor{currentstroke}{rgb}{0.000000,0.000000,0.000000}%
\pgfsetstrokecolor{currentstroke}%
\pgfsetstrokeopacity{0.000000}%
\pgfsetdash{}{0pt}%
\pgfpathmoveto{\pgfqpoint{1.082084in}{0.210879in}}%
\pgfpathlineto{\pgfqpoint{1.080487in}{0.211762in}}%
\pgfpathlineto{\pgfqpoint{1.079012in}{0.214678in}}%
\pgfpathlineto{\pgfqpoint{1.073207in}{0.214262in}}%
\pgfpathlineto{\pgfqpoint{1.071517in}{0.212096in}}%
\pgfpathlineto{\pgfqpoint{1.069651in}{0.211374in}}%
\pgfpathlineto{\pgfqpoint{1.067915in}{0.213771in}}%
\pgfpathlineto{\pgfqpoint{1.064056in}{0.213557in}}%
\pgfpathlineto{\pgfqpoint{1.066029in}{0.216398in}}%
\pgfpathlineto{\pgfqpoint{1.068443in}{0.216507in}}%
\pgfpathlineto{\pgfqpoint{1.067670in}{0.218823in}}%
\pgfpathlineto{\pgfqpoint{1.067684in}{0.221839in}}%
\pgfpathlineto{\pgfqpoint{1.068559in}{0.223929in}}%
\pgfpathlineto{\pgfqpoint{1.067133in}{0.226051in}}%
\pgfpathlineto{\pgfqpoint{1.068011in}{0.227410in}}%
\pgfpathlineto{\pgfqpoint{1.070899in}{0.228330in}}%
\pgfpathlineto{\pgfqpoint{1.071501in}{0.231767in}}%
\pgfpathlineto{\pgfqpoint{1.068523in}{0.231093in}}%
\pgfpathlineto{\pgfqpoint{1.068667in}{0.235893in}}%
\pgfpathlineto{\pgfqpoint{1.067104in}{0.238920in}}%
\pgfpathlineto{\pgfqpoint{1.067789in}{0.242382in}}%
\pgfpathlineto{\pgfqpoint{1.067367in}{0.244154in}}%
\pgfpathlineto{\pgfqpoint{1.065528in}{0.244746in}}%
\pgfpathlineto{\pgfqpoint{1.061828in}{0.247605in}}%
\pgfpathlineto{\pgfqpoint{1.061314in}{0.250040in}}%
\pgfpathlineto{\pgfqpoint{1.062351in}{0.251212in}}%
\pgfpathlineto{\pgfqpoint{1.061213in}{0.253041in}}%
\pgfpathlineto{\pgfqpoint{1.062006in}{0.255168in}}%
\pgfpathlineto{\pgfqpoint{1.061710in}{0.260818in}}%
\pgfpathlineto{\pgfqpoint{1.062710in}{0.261351in}}%
\pgfpathlineto{\pgfqpoint{1.063933in}{0.265257in}}%
\pgfpathlineto{\pgfqpoint{1.061849in}{0.265034in}}%
\pgfpathlineto{\pgfqpoint{1.062080in}{0.267893in}}%
\pgfpathlineto{\pgfqpoint{1.064059in}{0.276502in}}%
\pgfpathlineto{\pgfqpoint{1.064827in}{0.282112in}}%
\pgfpathlineto{\pgfqpoint{1.062541in}{0.281407in}}%
\pgfpathlineto{\pgfqpoint{1.062850in}{0.278917in}}%
\pgfpathlineto{\pgfqpoint{1.061375in}{0.278552in}}%
\pgfpathlineto{\pgfqpoint{1.061437in}{0.275880in}}%
\pgfpathlineto{\pgfqpoint{1.059564in}{0.269832in}}%
\pgfpathlineto{\pgfqpoint{1.059156in}{0.263608in}}%
\pgfpathlineto{\pgfqpoint{1.057760in}{0.259141in}}%
\pgfpathlineto{\pgfqpoint{1.056698in}{0.253549in}}%
\pgfpathlineto{\pgfqpoint{1.054233in}{0.250993in}}%
\pgfpathlineto{\pgfqpoint{1.052026in}{0.254747in}}%
\pgfpathlineto{\pgfqpoint{1.052338in}{0.257691in}}%
\pgfpathlineto{\pgfqpoint{1.050359in}{0.258935in}}%
\pgfpathlineto{\pgfqpoint{1.045539in}{0.260432in}}%
\pgfpathlineto{\pgfqpoint{1.048292in}{0.263786in}}%
\pgfpathlineto{\pgfqpoint{1.048887in}{0.269049in}}%
\pgfpathlineto{\pgfqpoint{1.048759in}{0.271883in}}%
\pgfpathlineto{\pgfqpoint{1.045964in}{0.271941in}}%
\pgfpathlineto{\pgfqpoint{1.044789in}{0.274978in}}%
\pgfpathlineto{\pgfqpoint{1.043552in}{0.275786in}}%
\pgfpathlineto{\pgfqpoint{1.044076in}{0.279452in}}%
\pgfpathlineto{\pgfqpoint{1.040938in}{0.280527in}}%
\pgfpathlineto{\pgfqpoint{1.038980in}{0.283005in}}%
\pgfpathlineto{\pgfqpoint{1.037883in}{0.281325in}}%
\pgfpathlineto{\pgfqpoint{1.040576in}{0.278511in}}%
\pgfpathlineto{\pgfqpoint{1.040732in}{0.274755in}}%
\pgfpathlineto{\pgfqpoint{1.043542in}{0.270946in}}%
\pgfpathlineto{\pgfqpoint{1.040999in}{0.270087in}}%
\pgfpathlineto{\pgfqpoint{1.042205in}{0.268402in}}%
\pgfpathlineto{\pgfqpoint{1.044236in}{0.268878in}}%
\pgfpathlineto{\pgfqpoint{1.043891in}{0.261248in}}%
\pgfpathlineto{\pgfqpoint{1.039661in}{0.261025in}}%
\pgfpathlineto{\pgfqpoint{1.040400in}{0.263502in}}%
\pgfpathlineto{\pgfqpoint{1.037395in}{0.262751in}}%
\pgfpathlineto{\pgfqpoint{1.033326in}{0.260670in}}%
\pgfpathlineto{\pgfqpoint{1.033334in}{0.264446in}}%
\pgfpathlineto{\pgfqpoint{1.032173in}{0.267309in}}%
\pgfpathlineto{\pgfqpoint{1.029605in}{0.268061in}}%
\pgfpathlineto{\pgfqpoint{1.025973in}{0.276463in}}%
\pgfpathlineto{\pgfqpoint{1.026377in}{0.278219in}}%
\pgfpathlineto{\pgfqpoint{1.024142in}{0.282697in}}%
\pgfpathlineto{\pgfqpoint{1.024419in}{0.288147in}}%
\pgfpathlineto{\pgfqpoint{1.022440in}{0.293955in}}%
\pgfpathlineto{\pgfqpoint{1.018989in}{0.298401in}}%
\pgfpathlineto{\pgfqpoint{1.017453in}{0.301716in}}%
\pgfpathlineto{\pgfqpoint{1.012827in}{0.307991in}}%
\pgfpathlineto{\pgfqpoint{1.008318in}{0.316535in}}%
\pgfpathlineto{\pgfqpoint{1.011237in}{0.316069in}}%
\pgfpathlineto{\pgfqpoint{1.015372in}{0.318061in}}%
\pgfpathlineto{\pgfqpoint{1.015809in}{0.323540in}}%
\pgfpathlineto{\pgfqpoint{1.016829in}{0.325554in}}%
\pgfpathlineto{\pgfqpoint{1.011575in}{0.323679in}}%
\pgfpathlineto{\pgfqpoint{1.009826in}{0.326603in}}%
\pgfpathlineto{\pgfqpoint{1.006161in}{0.325696in}}%
\pgfpathlineto{\pgfqpoint{1.005156in}{0.323648in}}%
\pgfpathlineto{\pgfqpoint{1.000051in}{0.326986in}}%
\pgfpathlineto{\pgfqpoint{0.997055in}{0.330455in}}%
\pgfpathlineto{\pgfqpoint{1.004608in}{0.344079in}}%
\pgfpathlineto{\pgfqpoint{1.009161in}{0.338796in}}%
\pgfpathlineto{\pgfqpoint{1.012055in}{0.339391in}}%
\pgfpathlineto{\pgfqpoint{1.015589in}{0.334105in}}%
\pgfpathlineto{\pgfqpoint{1.021713in}{0.335942in}}%
\pgfpathlineto{\pgfqpoint{1.027923in}{0.333090in}}%
\pgfpathlineto{\pgfqpoint{1.029124in}{0.331408in}}%
\pgfpathlineto{\pgfqpoint{1.024345in}{0.326303in}}%
\pgfpathlineto{\pgfqpoint{1.025043in}{0.322901in}}%
\pgfpathlineto{\pgfqpoint{1.028234in}{0.318423in}}%
\pgfpathlineto{\pgfqpoint{1.027542in}{0.314145in}}%
\pgfpathlineto{\pgfqpoint{1.033899in}{0.293981in}}%
\pgfpathlineto{\pgfqpoint{1.032438in}{0.287802in}}%
\pgfpathlineto{\pgfqpoint{1.030648in}{0.284485in}}%
\pgfpathlineto{\pgfqpoint{1.031736in}{0.284026in}}%
\pgfpathlineto{\pgfqpoint{1.035421in}{0.285214in}}%
\pgfpathlineto{\pgfqpoint{1.041713in}{0.285984in}}%
\pgfpathlineto{\pgfqpoint{1.046557in}{0.284922in}}%
\pgfpathlineto{\pgfqpoint{1.049608in}{0.287280in}}%
\pgfpathlineto{\pgfqpoint{1.051192in}{0.290700in}}%
\pgfpathlineto{\pgfqpoint{1.054086in}{0.290900in}}%
\pgfpathlineto{\pgfqpoint{1.058160in}{0.294663in}}%
\pgfpathlineto{\pgfqpoint{1.061301in}{0.293799in}}%
\pgfpathlineto{\pgfqpoint{1.068256in}{0.294295in}}%
\pgfpathlineto{\pgfqpoint{1.069607in}{0.293456in}}%
\pgfpathlineto{\pgfqpoint{1.071897in}{0.287157in}}%
\pgfpathlineto{\pgfqpoint{1.071746in}{0.284296in}}%
\pgfpathlineto{\pgfqpoint{1.069635in}{0.282181in}}%
\pgfpathlineto{\pgfqpoint{1.069594in}{0.277534in}}%
\pgfpathlineto{\pgfqpoint{1.072600in}{0.274944in}}%
\pgfpathlineto{\pgfqpoint{1.072637in}{0.271740in}}%
\pgfpathlineto{\pgfqpoint{1.074084in}{0.269420in}}%
\pgfpathlineto{\pgfqpoint{1.073539in}{0.266568in}}%
\pgfpathlineto{\pgfqpoint{1.073899in}{0.262067in}}%
\pgfpathlineto{\pgfqpoint{1.077857in}{0.256052in}}%
\pgfpathlineto{\pgfqpoint{1.078297in}{0.250259in}}%
\pgfpathlineto{\pgfqpoint{1.080003in}{0.246871in}}%
\pgfpathlineto{\pgfqpoint{1.078915in}{0.244757in}}%
\pgfpathlineto{\pgfqpoint{1.079981in}{0.235727in}}%
\pgfpathlineto{\pgfqpoint{1.079690in}{0.232569in}}%
\pgfpathlineto{\pgfqpoint{1.080751in}{0.226649in}}%
\pgfpathlineto{\pgfqpoint{1.081182in}{0.216121in}}%
\pgfpathlineto{\pgfqpoint{1.082084in}{0.210879in}}%
\pgfpathclose%
\pgfusepath{fill}%
\end{pgfscope}%
\begin{pgfscope}%
\pgfpathrectangle{\pgfqpoint{0.100000in}{0.100000in}}{\pgfqpoint{3.608454in}{2.310000in}}%
\pgfusepath{clip}%
\pgfsetbuttcap%
\pgfsetmiterjoin%
\definecolor{currentfill}{rgb}{0.000000,0.364706,0.817647}%
\pgfsetfillcolor{currentfill}%
\pgfsetlinewidth{0.000000pt}%
\definecolor{currentstroke}{rgb}{0.000000,0.000000,0.000000}%
\pgfsetstrokecolor{currentstroke}%
\pgfsetstrokeopacity{0.000000}%
\pgfsetdash{}{0pt}%
\pgfpathmoveto{\pgfqpoint{1.046772in}{0.216481in}}%
\pgfpathlineto{\pgfqpoint{1.049473in}{0.221987in}}%
\pgfpathlineto{\pgfqpoint{1.051491in}{0.223563in}}%
\pgfpathlineto{\pgfqpoint{1.053235in}{0.226047in}}%
\pgfpathlineto{\pgfqpoint{1.053669in}{0.234273in}}%
\pgfpathlineto{\pgfqpoint{1.055631in}{0.239433in}}%
\pgfpathlineto{\pgfqpoint{1.056457in}{0.243986in}}%
\pgfpathlineto{\pgfqpoint{1.058002in}{0.245736in}}%
\pgfpathlineto{\pgfqpoint{1.057331in}{0.248039in}}%
\pgfpathlineto{\pgfqpoint{1.057875in}{0.253294in}}%
\pgfpathlineto{\pgfqpoint{1.059869in}{0.251052in}}%
\pgfpathlineto{\pgfqpoint{1.059422in}{0.246810in}}%
\pgfpathlineto{\pgfqpoint{1.061246in}{0.245939in}}%
\pgfpathlineto{\pgfqpoint{1.065594in}{0.242418in}}%
\pgfpathlineto{\pgfqpoint{1.066004in}{0.236870in}}%
\pgfpathlineto{\pgfqpoint{1.065636in}{0.230146in}}%
\pgfpathlineto{\pgfqpoint{1.064005in}{0.231270in}}%
\pgfpathlineto{\pgfqpoint{1.064341in}{0.239679in}}%
\pgfpathlineto{\pgfqpoint{1.062886in}{0.239330in}}%
\pgfpathlineto{\pgfqpoint{1.060924in}{0.235715in}}%
\pgfpathlineto{\pgfqpoint{1.062789in}{0.235055in}}%
\pgfpathlineto{\pgfqpoint{1.062708in}{0.232793in}}%
\pgfpathlineto{\pgfqpoint{1.061477in}{0.229384in}}%
\pgfpathlineto{\pgfqpoint{1.062971in}{0.228484in}}%
\pgfpathlineto{\pgfqpoint{1.062440in}{0.223965in}}%
\pgfpathlineto{\pgfqpoint{1.059310in}{0.224988in}}%
\pgfpathlineto{\pgfqpoint{1.061475in}{0.221010in}}%
\pgfpathlineto{\pgfqpoint{1.059297in}{0.219771in}}%
\pgfpathlineto{\pgfqpoint{1.058337in}{0.221103in}}%
\pgfpathlineto{\pgfqpoint{1.056457in}{0.219297in}}%
\pgfpathlineto{\pgfqpoint{1.053559in}{0.218893in}}%
\pgfpathlineto{\pgfqpoint{1.050357in}{0.216700in}}%
\pgfpathlineto{\pgfqpoint{1.046772in}{0.216481in}}%
\pgfpathclose%
\pgfusepath{fill}%
\end{pgfscope}%
\begin{pgfscope}%
\pgfpathrectangle{\pgfqpoint{0.100000in}{0.100000in}}{\pgfqpoint{3.608454in}{2.310000in}}%
\pgfusepath{clip}%
\pgfsetbuttcap%
\pgfsetmiterjoin%
\definecolor{currentfill}{rgb}{0.000000,0.364706,0.817647}%
\pgfsetfillcolor{currentfill}%
\pgfsetlinewidth{0.000000pt}%
\definecolor{currentstroke}{rgb}{0.000000,0.000000,0.000000}%
\pgfsetstrokecolor{currentstroke}%
\pgfsetstrokeopacity{0.000000}%
\pgfsetdash{}{0pt}%
\pgfpathmoveto{\pgfqpoint{1.031898in}{0.249479in}}%
\pgfpathlineto{\pgfqpoint{1.031917in}{0.252366in}}%
\pgfpathlineto{\pgfqpoint{1.033185in}{0.256544in}}%
\pgfpathlineto{\pgfqpoint{1.037420in}{0.256450in}}%
\pgfpathlineto{\pgfqpoint{1.039053in}{0.255636in}}%
\pgfpathlineto{\pgfqpoint{1.040317in}{0.257461in}}%
\pgfpathlineto{\pgfqpoint{1.041757in}{0.255742in}}%
\pgfpathlineto{\pgfqpoint{1.045494in}{0.257061in}}%
\pgfpathlineto{\pgfqpoint{1.048001in}{0.253354in}}%
\pgfpathlineto{\pgfqpoint{1.048579in}{0.249693in}}%
\pgfpathlineto{\pgfqpoint{1.052991in}{0.247043in}}%
\pgfpathlineto{\pgfqpoint{1.054422in}{0.245099in}}%
\pgfpathlineto{\pgfqpoint{1.053698in}{0.242447in}}%
\pgfpathlineto{\pgfqpoint{1.052033in}{0.241302in}}%
\pgfpathlineto{\pgfqpoint{1.050384in}{0.238572in}}%
\pgfpathlineto{\pgfqpoint{1.047084in}{0.239857in}}%
\pgfpathlineto{\pgfqpoint{1.044554in}{0.244818in}}%
\pgfpathlineto{\pgfqpoint{1.042604in}{0.247304in}}%
\pgfpathlineto{\pgfqpoint{1.041099in}{0.250574in}}%
\pgfpathlineto{\pgfqpoint{1.037351in}{0.246961in}}%
\pgfpathlineto{\pgfqpoint{1.031898in}{0.249479in}}%
\pgfpathclose%
\pgfusepath{fill}%
\end{pgfscope}%
\begin{pgfscope}%
\pgfpathrectangle{\pgfqpoint{0.100000in}{0.100000in}}{\pgfqpoint{3.608454in}{2.310000in}}%
\pgfusepath{clip}%
\pgfsetbuttcap%
\pgfsetmiterjoin%
\definecolor{currentfill}{rgb}{0.000000,0.482353,0.758824}%
\pgfsetfillcolor{currentfill}%
\pgfsetlinewidth{0.000000pt}%
\definecolor{currentstroke}{rgb}{0.000000,0.000000,0.000000}%
\pgfsetstrokecolor{currentstroke}%
\pgfsetstrokeopacity{0.000000}%
\pgfsetdash{}{0pt}%
\pgfpathmoveto{\pgfqpoint{3.344220in}{1.667169in}}%
\pgfpathlineto{\pgfqpoint{3.344000in}{1.680413in}}%
\pgfpathlineto{\pgfqpoint{3.366550in}{1.688953in}}%
\pgfpathlineto{\pgfqpoint{3.369634in}{1.670838in}}%
\pgfpathlineto{\pgfqpoint{3.375111in}{1.665432in}}%
\pgfpathlineto{\pgfqpoint{3.362988in}{1.653409in}}%
\pgfpathlineto{\pgfqpoint{3.369189in}{1.645562in}}%
\pgfpathlineto{\pgfqpoint{3.372252in}{1.639364in}}%
\pgfpathlineto{\pgfqpoint{3.379684in}{1.642975in}}%
\pgfpathlineto{\pgfqpoint{3.395641in}{1.645504in}}%
\pgfpathlineto{\pgfqpoint{3.404650in}{1.652435in}}%
\pgfpathlineto{\pgfqpoint{3.420950in}{1.656635in}}%
\pgfpathlineto{\pgfqpoint{3.428750in}{1.659772in}}%
\pgfpathlineto{\pgfqpoint{3.439665in}{1.671536in}}%
\pgfpathlineto{\pgfqpoint{3.441705in}{1.675928in}}%
\pgfpathlineto{\pgfqpoint{3.448912in}{1.669829in}}%
\pgfpathlineto{\pgfqpoint{3.454605in}{1.672223in}}%
\pgfpathlineto{\pgfqpoint{3.458853in}{1.668402in}}%
\pgfpathlineto{\pgfqpoint{3.465859in}{1.676647in}}%
\pgfpathlineto{\pgfqpoint{3.471136in}{1.676378in}}%
\pgfpathlineto{\pgfqpoint{3.438928in}{1.650105in}}%
\pgfpathlineto{\pgfqpoint{3.418976in}{1.636559in}}%
\pgfpathlineto{\pgfqpoint{3.411351in}{1.630261in}}%
\pgfpathlineto{\pgfqpoint{3.399067in}{1.623101in}}%
\pgfpathlineto{\pgfqpoint{3.394098in}{1.622800in}}%
\pgfpathlineto{\pgfqpoint{3.382218in}{1.615957in}}%
\pgfpathlineto{\pgfqpoint{3.369738in}{1.613583in}}%
\pgfpathlineto{\pgfqpoint{3.361694in}{1.608438in}}%
\pgfpathlineto{\pgfqpoint{3.353119in}{1.610609in}}%
\pgfpathlineto{\pgfqpoint{3.350697in}{1.605499in}}%
\pgfpathlineto{\pgfqpoint{3.343167in}{1.599400in}}%
\pgfpathlineto{\pgfqpoint{3.344569in}{1.612129in}}%
\pgfpathlineto{\pgfqpoint{3.352075in}{1.614317in}}%
\pgfpathlineto{\pgfqpoint{3.355371in}{1.641928in}}%
\pgfpathlineto{\pgfqpoint{3.353043in}{1.654103in}}%
\pgfpathlineto{\pgfqpoint{3.346226in}{1.661401in}}%
\pgfpathlineto{\pgfqpoint{3.344220in}{1.667169in}}%
\pgfpathclose%
\pgfusepath{fill}%
\end{pgfscope}%
\begin{pgfscope}%
\pgfpathrectangle{\pgfqpoint{0.100000in}{0.100000in}}{\pgfqpoint{3.608454in}{2.310000in}}%
\pgfusepath{clip}%
\pgfsetbuttcap%
\pgfsetmiterjoin%
\definecolor{currentfill}{rgb}{0.000000,0.478431,0.760784}%
\pgfsetfillcolor{currentfill}%
\pgfsetlinewidth{0.000000pt}%
\definecolor{currentstroke}{rgb}{0.000000,0.000000,0.000000}%
\pgfsetstrokecolor{currentstroke}%
\pgfsetstrokeopacity{0.000000}%
\pgfsetdash{}{0pt}%
\pgfpathmoveto{\pgfqpoint{1.157736in}{1.964285in}}%
\pgfpathlineto{\pgfqpoint{1.146770in}{1.973621in}}%
\pgfpathlineto{\pgfqpoint{1.147664in}{1.994687in}}%
\pgfpathlineto{\pgfqpoint{1.143228in}{2.002013in}}%
\pgfpathlineto{\pgfqpoint{1.145294in}{2.011278in}}%
\pgfpathlineto{\pgfqpoint{1.148015in}{2.013748in}}%
\pgfpathlineto{\pgfqpoint{1.158552in}{2.015787in}}%
\pgfpathlineto{\pgfqpoint{1.162054in}{2.020418in}}%
\pgfpathlineto{\pgfqpoint{1.164980in}{2.032458in}}%
\pgfpathlineto{\pgfqpoint{1.160126in}{2.041022in}}%
\pgfpathlineto{\pgfqpoint{1.153421in}{2.042231in}}%
\pgfpathlineto{\pgfqpoint{1.155484in}{2.052872in}}%
\pgfpathlineto{\pgfqpoint{1.141533in}{2.055406in}}%
\pgfpathlineto{\pgfqpoint{1.146831in}{2.082471in}}%
\pgfpathlineto{\pgfqpoint{1.132903in}{2.085429in}}%
\pgfpathlineto{\pgfqpoint{1.137742in}{2.109618in}}%
\pgfpathlineto{\pgfqpoint{1.134932in}{2.119015in}}%
\pgfpathlineto{\pgfqpoint{1.135799in}{2.128987in}}%
\pgfpathlineto{\pgfqpoint{1.140741in}{2.130368in}}%
\pgfpathlineto{\pgfqpoint{1.143835in}{2.140807in}}%
\pgfpathlineto{\pgfqpoint{1.149026in}{2.145360in}}%
\pgfpathlineto{\pgfqpoint{1.150853in}{2.134585in}}%
\pgfpathlineto{\pgfqpoint{1.148672in}{2.124295in}}%
\pgfpathlineto{\pgfqpoint{1.152365in}{2.116451in}}%
\pgfpathlineto{\pgfqpoint{1.165476in}{2.115944in}}%
\pgfpathlineto{\pgfqpoint{1.173281in}{2.111946in}}%
\pgfpathlineto{\pgfqpoint{1.173331in}{2.107985in}}%
\pgfpathlineto{\pgfqpoint{1.179403in}{2.102869in}}%
\pgfpathlineto{\pgfqpoint{1.191810in}{2.101317in}}%
\pgfpathlineto{\pgfqpoint{1.187119in}{2.075916in}}%
\pgfpathlineto{\pgfqpoint{1.192015in}{2.070842in}}%
\pgfpathlineto{\pgfqpoint{1.199816in}{2.068342in}}%
\pgfpathlineto{\pgfqpoint{1.196622in}{2.051472in}}%
\pgfpathlineto{\pgfqpoint{1.203686in}{2.050112in}}%
\pgfpathlineto{\pgfqpoint{1.203339in}{2.044520in}}%
\pgfpathlineto{\pgfqpoint{1.210034in}{2.040419in}}%
\pgfpathlineto{\pgfqpoint{1.213019in}{2.026776in}}%
\pgfpathlineto{\pgfqpoint{1.216917in}{2.019801in}}%
\pgfpathlineto{\pgfqpoint{1.219896in}{2.007840in}}%
\pgfpathlineto{\pgfqpoint{1.224493in}{2.008821in}}%
\pgfpathlineto{\pgfqpoint{1.228711in}{2.003790in}}%
\pgfpathlineto{\pgfqpoint{1.224264in}{1.998009in}}%
\pgfpathlineto{\pgfqpoint{1.225529in}{1.987742in}}%
\pgfpathlineto{\pgfqpoint{1.210052in}{1.990143in}}%
\pgfpathlineto{\pgfqpoint{1.208958in}{1.985065in}}%
\pgfpathlineto{\pgfqpoint{1.203995in}{1.980934in}}%
\pgfpathlineto{\pgfqpoint{1.199884in}{1.972842in}}%
\pgfpathlineto{\pgfqpoint{1.192609in}{1.969852in}}%
\pgfpathlineto{\pgfqpoint{1.185357in}{1.963304in}}%
\pgfpathlineto{\pgfqpoint{1.176413in}{1.966592in}}%
\pgfpathlineto{\pgfqpoint{1.173618in}{1.970061in}}%
\pgfpathlineto{\pgfqpoint{1.165857in}{1.971182in}}%
\pgfpathlineto{\pgfqpoint{1.157736in}{1.964285in}}%
\pgfpathclose%
\pgfusepath{fill}%
\end{pgfscope}%
\begin{pgfscope}%
\pgfpathrectangle{\pgfqpoint{0.100000in}{0.100000in}}{\pgfqpoint{3.608454in}{2.310000in}}%
\pgfusepath{clip}%
\pgfsetbuttcap%
\pgfsetmiterjoin%
\definecolor{currentfill}{rgb}{0.000000,0.372549,0.813725}%
\pgfsetfillcolor{currentfill}%
\pgfsetlinewidth{0.000000pt}%
\definecolor{currentstroke}{rgb}{0.000000,0.000000,0.000000}%
\pgfsetstrokecolor{currentstroke}%
\pgfsetstrokeopacity{0.000000}%
\pgfsetdash{}{0pt}%
\pgfpathmoveto{\pgfqpoint{1.964117in}{1.448246in}}%
\pgfpathlineto{\pgfqpoint{1.936670in}{1.449037in}}%
\pgfpathlineto{\pgfqpoint{1.937428in}{1.476563in}}%
\pgfpathlineto{\pgfqpoint{1.938437in}{1.514106in}}%
\pgfpathlineto{\pgfqpoint{1.942299in}{1.518835in}}%
\pgfpathlineto{\pgfqpoint{1.952480in}{1.526297in}}%
\pgfpathlineto{\pgfqpoint{1.966281in}{1.530825in}}%
\pgfpathlineto{\pgfqpoint{1.965564in}{1.503301in}}%
\pgfpathlineto{\pgfqpoint{1.964117in}{1.448246in}}%
\pgfpathclose%
\pgfusepath{fill}%
\end{pgfscope}%
\begin{pgfscope}%
\pgfpathrectangle{\pgfqpoint{0.100000in}{0.100000in}}{\pgfqpoint{3.608454in}{2.310000in}}%
\pgfusepath{clip}%
\pgfsetbuttcap%
\pgfsetmiterjoin%
\definecolor{currentfill}{rgb}{0.000000,0.678431,0.660784}%
\pgfsetfillcolor{currentfill}%
\pgfsetlinewidth{0.000000pt}%
\definecolor{currentstroke}{rgb}{0.000000,0.000000,0.000000}%
\pgfsetstrokecolor{currentstroke}%
\pgfsetstrokeopacity{0.000000}%
\pgfsetdash{}{0pt}%
\pgfpathmoveto{\pgfqpoint{0.737189in}{2.001359in}}%
\pgfpathlineto{\pgfqpoint{0.734618in}{1.987950in}}%
\pgfpathlineto{\pgfqpoint{0.723444in}{1.948508in}}%
\pgfpathlineto{\pgfqpoint{0.710132in}{1.952205in}}%
\pgfpathlineto{\pgfqpoint{0.711901in}{1.958542in}}%
\pgfpathlineto{\pgfqpoint{0.706253in}{1.964177in}}%
\pgfpathlineto{\pgfqpoint{0.686387in}{1.969750in}}%
\pgfpathlineto{\pgfqpoint{0.693514in}{1.994463in}}%
\pgfpathlineto{\pgfqpoint{0.695564in}{1.998853in}}%
\pgfpathlineto{\pgfqpoint{0.692032in}{2.003466in}}%
\pgfpathlineto{\pgfqpoint{0.692055in}{2.010580in}}%
\pgfpathlineto{\pgfqpoint{0.697834in}{2.026472in}}%
\pgfpathlineto{\pgfqpoint{0.695191in}{2.034615in}}%
\pgfpathlineto{\pgfqpoint{0.702993in}{2.049946in}}%
\pgfpathlineto{\pgfqpoint{0.710674in}{2.050338in}}%
\pgfpathlineto{\pgfqpoint{0.710428in}{2.058581in}}%
\pgfpathlineto{\pgfqpoint{0.707552in}{2.064386in}}%
\pgfpathlineto{\pgfqpoint{0.713071in}{2.066429in}}%
\pgfpathlineto{\pgfqpoint{0.723978in}{2.065521in}}%
\pgfpathlineto{\pgfqpoint{0.737141in}{2.068948in}}%
\pgfpathlineto{\pgfqpoint{0.724757in}{2.026561in}}%
\pgfpathlineto{\pgfqpoint{0.731442in}{2.024554in}}%
\pgfpathlineto{\pgfqpoint{0.729534in}{2.017963in}}%
\pgfpathlineto{\pgfqpoint{0.736122in}{2.015946in}}%
\pgfpathlineto{\pgfqpoint{0.730749in}{2.003174in}}%
\pgfpathlineto{\pgfqpoint{0.737189in}{2.001359in}}%
\pgfpathclose%
\pgfusepath{fill}%
\end{pgfscope}%
\begin{pgfscope}%
\pgfpathrectangle{\pgfqpoint{0.100000in}{0.100000in}}{\pgfqpoint{3.608454in}{2.310000in}}%
\pgfusepath{clip}%
\pgfsetbuttcap%
\pgfsetmiterjoin%
\definecolor{currentfill}{rgb}{0.000000,0.521569,0.739216}%
\pgfsetfillcolor{currentfill}%
\pgfsetlinewidth{0.000000pt}%
\definecolor{currentstroke}{rgb}{0.000000,0.000000,0.000000}%
\pgfsetstrokecolor{currentstroke}%
\pgfsetstrokeopacity{0.000000}%
\pgfsetdash{}{0pt}%
\pgfpathmoveto{\pgfqpoint{1.792481in}{1.427797in}}%
\pgfpathlineto{\pgfqpoint{1.724173in}{1.432221in}}%
\pgfpathlineto{\pgfqpoint{1.726124in}{1.459646in}}%
\pgfpathlineto{\pgfqpoint{1.682362in}{1.462771in}}%
\pgfpathlineto{\pgfqpoint{1.684512in}{1.490228in}}%
\pgfpathlineto{\pgfqpoint{1.732544in}{1.486727in}}%
\pgfpathlineto{\pgfqpoint{1.760676in}{1.485050in}}%
\pgfpathlineto{\pgfqpoint{1.758946in}{1.457374in}}%
\pgfpathlineto{\pgfqpoint{1.793800in}{1.455314in}}%
\pgfpathlineto{\pgfqpoint{1.792481in}{1.427797in}}%
\pgfpathclose%
\pgfusepath{fill}%
\end{pgfscope}%
\begin{pgfscope}%
\pgfpathrectangle{\pgfqpoint{0.100000in}{0.100000in}}{\pgfqpoint{3.608454in}{2.310000in}}%
\pgfusepath{clip}%
\pgfsetbuttcap%
\pgfsetmiterjoin%
\definecolor{currentfill}{rgb}{0.000000,0.505882,0.747059}%
\pgfsetfillcolor{currentfill}%
\pgfsetlinewidth{0.000000pt}%
\definecolor{currentstroke}{rgb}{0.000000,0.000000,0.000000}%
\pgfsetstrokecolor{currentstroke}%
\pgfsetstrokeopacity{0.000000}%
\pgfsetdash{}{0pt}%
\pgfpathmoveto{\pgfqpoint{1.223167in}{1.249876in}}%
\pgfpathlineto{\pgfqpoint{1.163571in}{1.259392in}}%
\pgfpathlineto{\pgfqpoint{1.116943in}{1.267777in}}%
\pgfpathlineto{\pgfqpoint{1.078577in}{1.274542in}}%
\pgfpathlineto{\pgfqpoint{1.087638in}{1.275829in}}%
\pgfpathlineto{\pgfqpoint{1.091741in}{1.280236in}}%
\pgfpathlineto{\pgfqpoint{1.099188in}{1.278642in}}%
\pgfpathlineto{\pgfqpoint{1.109948in}{1.282813in}}%
\pgfpathlineto{\pgfqpoint{1.113627in}{1.291412in}}%
\pgfpathlineto{\pgfqpoint{1.122123in}{1.294263in}}%
\pgfpathlineto{\pgfqpoint{1.130720in}{1.308550in}}%
\pgfpathlineto{\pgfqpoint{1.134341in}{1.312779in}}%
\pgfpathlineto{\pgfqpoint{1.141396in}{1.315621in}}%
\pgfpathlineto{\pgfqpoint{1.150177in}{1.332380in}}%
\pgfpathlineto{\pgfqpoint{1.159137in}{1.332405in}}%
\pgfpathlineto{\pgfqpoint{1.168304in}{1.339366in}}%
\pgfpathlineto{\pgfqpoint{1.170873in}{1.338075in}}%
\pgfpathlineto{\pgfqpoint{1.177278in}{1.346155in}}%
\pgfpathlineto{\pgfqpoint{1.183026in}{1.348649in}}%
\pgfpathlineto{\pgfqpoint{1.186021in}{1.352257in}}%
\pgfpathlineto{\pgfqpoint{1.183516in}{1.359664in}}%
\pgfpathlineto{\pgfqpoint{1.181575in}{1.376785in}}%
\pgfpathlineto{\pgfqpoint{1.240533in}{1.367269in}}%
\pgfpathlineto{\pgfqpoint{1.237792in}{1.349736in}}%
\pgfpathlineto{\pgfqpoint{1.237566in}{1.340907in}}%
\pgfpathlineto{\pgfqpoint{1.234133in}{1.318758in}}%
\pgfpathlineto{\pgfqpoint{1.248158in}{1.317741in}}%
\pgfpathlineto{\pgfqpoint{1.282717in}{1.312481in}}%
\pgfpathlineto{\pgfqpoint{1.284736in}{1.306313in}}%
\pgfpathlineto{\pgfqpoint{1.298775in}{1.307121in}}%
\pgfpathlineto{\pgfqpoint{1.305818in}{1.299625in}}%
\pgfpathlineto{\pgfqpoint{1.303897in}{1.295339in}}%
\pgfpathlineto{\pgfqpoint{1.293840in}{1.286121in}}%
\pgfpathlineto{\pgfqpoint{1.291651in}{1.276850in}}%
\pgfpathlineto{\pgfqpoint{1.286167in}{1.271003in}}%
\pgfpathlineto{\pgfqpoint{1.280183in}{1.269745in}}%
\pgfpathlineto{\pgfqpoint{1.272888in}{1.260292in}}%
\pgfpathlineto{\pgfqpoint{1.268205in}{1.247503in}}%
\pgfpathlineto{\pgfqpoint{1.264762in}{1.243586in}}%
\pgfpathlineto{\pgfqpoint{1.223167in}{1.249876in}}%
\pgfpathclose%
\pgfusepath{fill}%
\end{pgfscope}%
\begin{pgfscope}%
\pgfpathrectangle{\pgfqpoint{0.100000in}{0.100000in}}{\pgfqpoint{3.608454in}{2.310000in}}%
\pgfusepath{clip}%
\pgfsetbuttcap%
\pgfsetmiterjoin%
\definecolor{currentfill}{rgb}{0.000000,0.423529,0.788235}%
\pgfsetfillcolor{currentfill}%
\pgfsetlinewidth{0.000000pt}%
\definecolor{currentstroke}{rgb}{0.000000,0.000000,0.000000}%
\pgfsetstrokecolor{currentstroke}%
\pgfsetstrokeopacity{0.000000}%
\pgfsetdash{}{0pt}%
\pgfpathmoveto{\pgfqpoint{2.669947in}{1.049451in}}%
\pgfpathlineto{\pgfqpoint{2.668245in}{1.039940in}}%
\pgfpathlineto{\pgfqpoint{2.670176in}{1.036633in}}%
\pgfpathlineto{\pgfqpoint{2.669558in}{1.020872in}}%
\pgfpathlineto{\pgfqpoint{2.673359in}{1.013797in}}%
\pgfpathlineto{\pgfqpoint{2.666467in}{1.008380in}}%
\pgfpathlineto{\pgfqpoint{2.659190in}{1.009555in}}%
\pgfpathlineto{\pgfqpoint{2.657488in}{0.993619in}}%
\pgfpathlineto{\pgfqpoint{2.622963in}{0.991238in}}%
\pgfpathlineto{\pgfqpoint{2.623076in}{0.990093in}}%
\pgfpathlineto{\pgfqpoint{2.595636in}{0.988190in}}%
\pgfpathlineto{\pgfqpoint{2.593933in}{1.008935in}}%
\pgfpathlineto{\pgfqpoint{2.599703in}{1.020694in}}%
\pgfpathlineto{\pgfqpoint{2.599104in}{1.027892in}}%
\pgfpathlineto{\pgfqpoint{2.607222in}{1.025513in}}%
\pgfpathlineto{\pgfqpoint{2.612717in}{1.030317in}}%
\pgfpathlineto{\pgfqpoint{2.611745in}{1.044790in}}%
\pgfpathlineto{\pgfqpoint{2.636003in}{1.046348in}}%
\pgfpathlineto{\pgfqpoint{2.634590in}{1.067850in}}%
\pgfpathlineto{\pgfqpoint{2.639911in}{1.067474in}}%
\pgfpathlineto{\pgfqpoint{2.648663in}{1.074175in}}%
\pgfpathlineto{\pgfqpoint{2.648588in}{1.077186in}}%
\pgfpathlineto{\pgfqpoint{2.662888in}{1.068219in}}%
\pgfpathlineto{\pgfqpoint{2.665061in}{1.060409in}}%
\pgfpathlineto{\pgfqpoint{2.668473in}{1.060141in}}%
\pgfpathlineto{\pgfqpoint{2.669947in}{1.049451in}}%
\pgfpathclose%
\pgfusepath{fill}%
\end{pgfscope}%
\begin{pgfscope}%
\pgfpathrectangle{\pgfqpoint{0.100000in}{0.100000in}}{\pgfqpoint{3.608454in}{2.310000in}}%
\pgfusepath{clip}%
\pgfsetbuttcap%
\pgfsetmiterjoin%
\definecolor{currentfill}{rgb}{0.000000,0.349020,0.825490}%
\pgfsetfillcolor{currentfill}%
\pgfsetlinewidth{0.000000pt}%
\definecolor{currentstroke}{rgb}{0.000000,0.000000,0.000000}%
\pgfsetstrokecolor{currentstroke}%
\pgfsetstrokeopacity{0.000000}%
\pgfsetdash{}{0pt}%
\pgfpathmoveto{\pgfqpoint{1.255196in}{1.457818in}}%
\pgfpathlineto{\pgfqpoint{1.267460in}{1.535366in}}%
\pgfpathlineto{\pgfqpoint{1.271578in}{1.562500in}}%
\pgfpathlineto{\pgfqpoint{1.302483in}{1.557690in}}%
\pgfpathlineto{\pgfqpoint{1.338390in}{1.552623in}}%
\pgfpathlineto{\pgfqpoint{1.401111in}{1.544070in}}%
\pgfpathlineto{\pgfqpoint{1.400660in}{1.537925in}}%
\pgfpathlineto{\pgfqpoint{1.411585in}{1.530462in}}%
\pgfpathlineto{\pgfqpoint{1.412111in}{1.525836in}}%
\pgfpathlineto{\pgfqpoint{1.405998in}{1.509570in}}%
\pgfpathlineto{\pgfqpoint{1.407647in}{1.498770in}}%
\pgfpathlineto{\pgfqpoint{1.407785in}{1.490488in}}%
\pgfpathlineto{\pgfqpoint{1.404073in}{1.463909in}}%
\pgfpathlineto{\pgfqpoint{1.403922in}{1.457766in}}%
\pgfpathlineto{\pgfqpoint{1.379506in}{1.460495in}}%
\pgfpathlineto{\pgfqpoint{1.380971in}{1.474066in}}%
\pgfpathlineto{\pgfqpoint{1.364204in}{1.476264in}}%
\pgfpathlineto{\pgfqpoint{1.362220in}{1.462510in}}%
\pgfpathlineto{\pgfqpoint{1.355468in}{1.463655in}}%
\pgfpathlineto{\pgfqpoint{1.354499in}{1.456659in}}%
\pgfpathlineto{\pgfqpoint{1.324000in}{1.460674in}}%
\pgfpathlineto{\pgfqpoint{1.322533in}{1.450461in}}%
\pgfpathlineto{\pgfqpoint{1.285446in}{1.455840in}}%
\pgfpathlineto{\pgfqpoint{1.284941in}{1.452441in}}%
\pgfpathlineto{\pgfqpoint{1.255196in}{1.457818in}}%
\pgfpathclose%
\pgfusepath{fill}%
\end{pgfscope}%
\begin{pgfscope}%
\pgfpathrectangle{\pgfqpoint{0.100000in}{0.100000in}}{\pgfqpoint{3.608454in}{2.310000in}}%
\pgfusepath{clip}%
\pgfsetbuttcap%
\pgfsetmiterjoin%
\definecolor{currentfill}{rgb}{0.000000,0.717647,0.641176}%
\pgfsetfillcolor{currentfill}%
\pgfsetlinewidth{0.000000pt}%
\definecolor{currentstroke}{rgb}{0.000000,0.000000,0.000000}%
\pgfsetstrokecolor{currentstroke}%
\pgfsetstrokeopacity{0.000000}%
\pgfsetdash{}{0pt}%
\pgfpathmoveto{\pgfqpoint{1.709808in}{0.564368in}}%
\pgfpathlineto{\pgfqpoint{1.748728in}{0.562219in}}%
\pgfpathlineto{\pgfqpoint{1.744553in}{0.490983in}}%
\pgfpathlineto{\pgfqpoint{1.737580in}{0.491300in}}%
\pgfpathlineto{\pgfqpoint{1.732274in}{0.498247in}}%
\pgfpathlineto{\pgfqpoint{1.729908in}{0.510537in}}%
\pgfpathlineto{\pgfqpoint{1.725203in}{0.520544in}}%
\pgfpathlineto{\pgfqpoint{1.726134in}{0.523344in}}%
\pgfpathlineto{\pgfqpoint{1.719531in}{0.529734in}}%
\pgfpathlineto{\pgfqpoint{1.716987in}{0.543109in}}%
\pgfpathlineto{\pgfqpoint{1.710227in}{0.553019in}}%
\pgfpathlineto{\pgfqpoint{1.709808in}{0.564368in}}%
\pgfpathclose%
\pgfusepath{fill}%
\end{pgfscope}%
\begin{pgfscope}%
\pgfpathrectangle{\pgfqpoint{0.100000in}{0.100000in}}{\pgfqpoint{3.608454in}{2.310000in}}%
\pgfusepath{clip}%
\pgfsetbuttcap%
\pgfsetmiterjoin%
\definecolor{currentfill}{rgb}{0.000000,0.631373,0.684314}%
\pgfsetfillcolor{currentfill}%
\pgfsetlinewidth{0.000000pt}%
\definecolor{currentstroke}{rgb}{0.000000,0.000000,0.000000}%
\pgfsetstrokecolor{currentstroke}%
\pgfsetstrokeopacity{0.000000}%
\pgfsetdash{}{0pt}%
\pgfpathmoveto{\pgfqpoint{2.111050in}{0.772661in}}%
\pgfpathlineto{\pgfqpoint{2.108517in}{0.767596in}}%
\pgfpathlineto{\pgfqpoint{2.111897in}{0.751748in}}%
\pgfpathlineto{\pgfqpoint{2.115910in}{0.747279in}}%
\pgfpathlineto{\pgfqpoint{2.106573in}{0.739232in}}%
\pgfpathlineto{\pgfqpoint{2.099586in}{0.743769in}}%
\pgfpathlineto{\pgfqpoint{2.096621in}{0.750671in}}%
\pgfpathlineto{\pgfqpoint{2.088232in}{0.752553in}}%
\pgfpathlineto{\pgfqpoint{2.062519in}{0.748614in}}%
\pgfpathlineto{\pgfqpoint{2.056507in}{0.745639in}}%
\pgfpathlineto{\pgfqpoint{2.058523in}{0.753655in}}%
\pgfpathlineto{\pgfqpoint{2.048151in}{0.760702in}}%
\pgfpathlineto{\pgfqpoint{2.047446in}{0.765794in}}%
\pgfpathlineto{\pgfqpoint{2.042310in}{0.767446in}}%
\pgfpathlineto{\pgfqpoint{2.034850in}{0.788420in}}%
\pgfpathlineto{\pgfqpoint{2.023652in}{0.805506in}}%
\pgfpathlineto{\pgfqpoint{2.013777in}{0.811827in}}%
\pgfpathlineto{\pgfqpoint{2.034479in}{0.813847in}}%
\pgfpathlineto{\pgfqpoint{2.034871in}{0.852102in}}%
\pgfpathlineto{\pgfqpoint{2.044379in}{0.851895in}}%
\pgfpathlineto{\pgfqpoint{2.050385in}{0.846426in}}%
\pgfpathlineto{\pgfqpoint{2.053670in}{0.846942in}}%
\pgfpathlineto{\pgfqpoint{2.064298in}{0.842428in}}%
\pgfpathlineto{\pgfqpoint{2.062420in}{0.861532in}}%
\pgfpathlineto{\pgfqpoint{2.086189in}{0.861574in}}%
\pgfpathlineto{\pgfqpoint{2.095246in}{0.861494in}}%
\pgfpathlineto{\pgfqpoint{2.096641in}{0.856769in}}%
\pgfpathlineto{\pgfqpoint{2.096510in}{0.830344in}}%
\pgfpathlineto{\pgfqpoint{2.107625in}{0.827720in}}%
\pgfpathlineto{\pgfqpoint{2.107805in}{0.772710in}}%
\pgfpathlineto{\pgfqpoint{2.111050in}{0.772661in}}%
\pgfpathclose%
\pgfusepath{fill}%
\end{pgfscope}%
\begin{pgfscope}%
\pgfpathrectangle{\pgfqpoint{0.100000in}{0.100000in}}{\pgfqpoint{3.608454in}{2.310000in}}%
\pgfusepath{clip}%
\pgfsetbuttcap%
\pgfsetmiterjoin%
\definecolor{currentfill}{rgb}{0.000000,0.600000,0.700000}%
\pgfsetfillcolor{currentfill}%
\pgfsetlinewidth{0.000000pt}%
\definecolor{currentstroke}{rgb}{0.000000,0.000000,0.000000}%
\pgfsetstrokecolor{currentstroke}%
\pgfsetstrokeopacity{0.000000}%
\pgfsetdash{}{0pt}%
\pgfpathmoveto{\pgfqpoint{0.469229in}{1.459834in}}%
\pgfpathlineto{\pgfqpoint{0.465525in}{1.465978in}}%
\pgfpathlineto{\pgfqpoint{0.471844in}{1.486991in}}%
\pgfpathlineto{\pgfqpoint{0.471692in}{1.491720in}}%
\pgfpathlineto{\pgfqpoint{0.476612in}{1.508153in}}%
\pgfpathlineto{\pgfqpoint{0.471851in}{1.509726in}}%
\pgfpathlineto{\pgfqpoint{0.471900in}{1.514999in}}%
\pgfpathlineto{\pgfqpoint{0.477138in}{1.516663in}}%
\pgfpathlineto{\pgfqpoint{0.480918in}{1.524945in}}%
\pgfpathlineto{\pgfqpoint{0.474910in}{1.526813in}}%
\pgfpathlineto{\pgfqpoint{0.479765in}{1.542960in}}%
\pgfpathlineto{\pgfqpoint{0.465266in}{1.547913in}}%
\pgfpathlineto{\pgfqpoint{0.460110in}{1.545792in}}%
\pgfpathlineto{\pgfqpoint{0.455284in}{1.549289in}}%
\pgfpathlineto{\pgfqpoint{0.454748in}{1.561199in}}%
\pgfpathlineto{\pgfqpoint{0.452551in}{1.565246in}}%
\pgfpathlineto{\pgfqpoint{0.452005in}{1.577321in}}%
\pgfpathlineto{\pgfqpoint{0.446281in}{1.581110in}}%
\pgfpathlineto{\pgfqpoint{0.450911in}{1.584637in}}%
\pgfpathlineto{\pgfqpoint{0.480562in}{1.575547in}}%
\pgfpathlineto{\pgfqpoint{0.485256in}{1.568714in}}%
\pgfpathlineto{\pgfqpoint{0.484414in}{1.564406in}}%
\pgfpathlineto{\pgfqpoint{0.489967in}{1.557121in}}%
\pgfpathlineto{\pgfqpoint{0.498193in}{1.555922in}}%
\pgfpathlineto{\pgfqpoint{0.502169in}{1.569148in}}%
\pgfpathlineto{\pgfqpoint{0.507112in}{1.573994in}}%
\pgfpathlineto{\pgfqpoint{0.514461in}{1.575916in}}%
\pgfpathlineto{\pgfqpoint{0.520214in}{1.588877in}}%
\pgfpathlineto{\pgfqpoint{0.532464in}{1.597564in}}%
\pgfpathlineto{\pgfqpoint{0.538914in}{1.596597in}}%
\pgfpathlineto{\pgfqpoint{0.556871in}{1.597127in}}%
\pgfpathlineto{\pgfqpoint{0.563417in}{1.600581in}}%
\pgfpathlineto{\pgfqpoint{0.567874in}{1.599010in}}%
\pgfpathlineto{\pgfqpoint{0.570419in}{1.591928in}}%
\pgfpathlineto{\pgfqpoint{0.599878in}{1.583444in}}%
\pgfpathlineto{\pgfqpoint{0.592721in}{1.558163in}}%
\pgfpathlineto{\pgfqpoint{0.607050in}{1.554247in}}%
\pgfpathlineto{\pgfqpoint{0.606969in}{1.551881in}}%
\pgfpathlineto{\pgfqpoint{0.632981in}{1.544578in}}%
\pgfpathlineto{\pgfqpoint{0.631100in}{1.537480in}}%
\pgfpathlineto{\pgfqpoint{0.623821in}{1.536529in}}%
\pgfpathlineto{\pgfqpoint{0.620041in}{1.519423in}}%
\pgfpathlineto{\pgfqpoint{0.623308in}{1.518183in}}%
\pgfpathlineto{\pgfqpoint{0.620520in}{1.503015in}}%
\pgfpathlineto{\pgfqpoint{0.609163in}{1.520157in}}%
\pgfpathlineto{\pgfqpoint{0.604651in}{1.513031in}}%
\pgfpathlineto{\pgfqpoint{0.606996in}{1.503863in}}%
\pgfpathlineto{\pgfqpoint{0.603382in}{1.497016in}}%
\pgfpathlineto{\pgfqpoint{0.597596in}{1.492295in}}%
\pgfpathlineto{\pgfqpoint{0.592836in}{1.501027in}}%
\pgfpathlineto{\pgfqpoint{0.583557in}{1.498623in}}%
\pgfpathlineto{\pgfqpoint{0.577147in}{1.506768in}}%
\pgfpathlineto{\pgfqpoint{0.566537in}{1.504878in}}%
\pgfpathlineto{\pgfqpoint{0.550096in}{1.491227in}}%
\pgfpathlineto{\pgfqpoint{0.537164in}{1.478720in}}%
\pgfpathlineto{\pgfqpoint{0.526181in}{1.471930in}}%
\pgfpathlineto{\pgfqpoint{0.515295in}{1.495292in}}%
\pgfpathlineto{\pgfqpoint{0.508101in}{1.469498in}}%
\pgfpathlineto{\pgfqpoint{0.504262in}{1.472446in}}%
\pgfpathlineto{\pgfqpoint{0.489989in}{1.471204in}}%
\pgfpathlineto{\pgfqpoint{0.469229in}{1.459834in}}%
\pgfpathclose%
\pgfusepath{fill}%
\end{pgfscope}%
\begin{pgfscope}%
\pgfpathrectangle{\pgfqpoint{0.100000in}{0.100000in}}{\pgfqpoint{3.608454in}{2.310000in}}%
\pgfusepath{clip}%
\pgfsetbuttcap%
\pgfsetmiterjoin%
\definecolor{currentfill}{rgb}{0.000000,0.478431,0.760784}%
\pgfsetfillcolor{currentfill}%
\pgfsetlinewidth{0.000000pt}%
\definecolor{currentstroke}{rgb}{0.000000,0.000000,0.000000}%
\pgfsetstrokecolor{currentstroke}%
\pgfsetstrokeopacity{0.000000}%
\pgfsetdash{}{0pt}%
\pgfpathmoveto{\pgfqpoint{1.886398in}{1.921692in}}%
\pgfpathlineto{\pgfqpoint{1.884519in}{1.921788in}}%
\pgfpathlineto{\pgfqpoint{1.885796in}{1.949443in}}%
\pgfpathlineto{\pgfqpoint{1.863197in}{1.950599in}}%
\pgfpathlineto{\pgfqpoint{1.864737in}{1.978421in}}%
\pgfpathlineto{\pgfqpoint{1.862851in}{1.978498in}}%
\pgfpathlineto{\pgfqpoint{1.864232in}{2.006135in}}%
\pgfpathlineto{\pgfqpoint{1.875835in}{2.005576in}}%
\pgfpathlineto{\pgfqpoint{1.917334in}{2.003750in}}%
\pgfpathlineto{\pgfqpoint{1.918791in}{1.996815in}}%
\pgfpathlineto{\pgfqpoint{1.946235in}{1.995849in}}%
\pgfpathlineto{\pgfqpoint{1.960057in}{1.995378in}}%
\pgfpathlineto{\pgfqpoint{1.960762in}{1.974477in}}%
\pgfpathlineto{\pgfqpoint{1.959926in}{1.946668in}}%
\pgfpathlineto{\pgfqpoint{1.940671in}{1.947331in}}%
\pgfpathlineto{\pgfqpoint{1.939707in}{1.919594in}}%
\pgfpathlineto{\pgfqpoint{1.886398in}{1.921692in}}%
\pgfpathclose%
\pgfusepath{fill}%
\end{pgfscope}%
\begin{pgfscope}%
\pgfpathrectangle{\pgfqpoint{0.100000in}{0.100000in}}{\pgfqpoint{3.608454in}{2.310000in}}%
\pgfusepath{clip}%
\pgfsetbuttcap%
\pgfsetmiterjoin%
\definecolor{currentfill}{rgb}{0.000000,0.529412,0.735294}%
\pgfsetfillcolor{currentfill}%
\pgfsetlinewidth{0.000000pt}%
\definecolor{currentstroke}{rgb}{0.000000,0.000000,0.000000}%
\pgfsetstrokecolor{currentstroke}%
\pgfsetstrokeopacity{0.000000}%
\pgfsetdash{}{0pt}%
\pgfpathmoveto{\pgfqpoint{3.420582in}{1.833115in}}%
\pgfpathlineto{\pgfqpoint{3.412904in}{1.835488in}}%
\pgfpathlineto{\pgfqpoint{3.401069in}{1.830482in}}%
\pgfpathlineto{\pgfqpoint{3.396386in}{1.835470in}}%
\pgfpathlineto{\pgfqpoint{3.382485in}{1.831082in}}%
\pgfpathlineto{\pgfqpoint{3.374364in}{1.831752in}}%
\pgfpathlineto{\pgfqpoint{3.370781in}{1.834556in}}%
\pgfpathlineto{\pgfqpoint{3.375120in}{1.840353in}}%
\pgfpathlineto{\pgfqpoint{3.375426in}{1.849437in}}%
\pgfpathlineto{\pgfqpoint{3.371356in}{1.850568in}}%
\pgfpathlineto{\pgfqpoint{3.372767in}{1.866271in}}%
\pgfpathlineto{\pgfqpoint{3.365489in}{1.867233in}}%
\pgfpathlineto{\pgfqpoint{3.365877in}{1.874467in}}%
\pgfpathlineto{\pgfqpoint{3.354162in}{1.876806in}}%
\pgfpathlineto{\pgfqpoint{3.358819in}{1.880708in}}%
\pgfpathlineto{\pgfqpoint{3.363839in}{1.892385in}}%
\pgfpathlineto{\pgfqpoint{3.353338in}{1.893246in}}%
\pgfpathlineto{\pgfqpoint{3.350605in}{1.900061in}}%
\pgfpathlineto{\pgfqpoint{3.349620in}{1.909475in}}%
\pgfpathlineto{\pgfqpoint{3.352226in}{1.924022in}}%
\pgfpathlineto{\pgfqpoint{3.361540in}{1.921976in}}%
\pgfpathlineto{\pgfqpoint{3.363137in}{1.929015in}}%
\pgfpathlineto{\pgfqpoint{3.370017in}{1.926361in}}%
\pgfpathlineto{\pgfqpoint{3.371539in}{1.933332in}}%
\pgfpathlineto{\pgfqpoint{3.384688in}{1.930010in}}%
\pgfpathlineto{\pgfqpoint{3.381069in}{1.911013in}}%
\pgfpathlineto{\pgfqpoint{3.384973in}{1.910295in}}%
\pgfpathlineto{\pgfqpoint{3.398006in}{1.914180in}}%
\pgfpathlineto{\pgfqpoint{3.396256in}{1.920610in}}%
\pgfpathlineto{\pgfqpoint{3.399587in}{1.926596in}}%
\pgfpathlineto{\pgfqpoint{3.405777in}{1.928235in}}%
\pgfpathlineto{\pgfqpoint{3.410255in}{1.935145in}}%
\pgfpathlineto{\pgfqpoint{3.415486in}{1.930416in}}%
\pgfpathlineto{\pgfqpoint{3.422059in}{1.929884in}}%
\pgfpathlineto{\pgfqpoint{3.423697in}{1.926204in}}%
\pgfpathlineto{\pgfqpoint{3.431085in}{1.927713in}}%
\pgfpathlineto{\pgfqpoint{3.433355in}{1.925260in}}%
\pgfpathlineto{\pgfqpoint{3.438972in}{1.915003in}}%
\pgfpathlineto{\pgfqpoint{3.442370in}{1.903310in}}%
\pgfpathlineto{\pgfqpoint{3.431626in}{1.898622in}}%
\pgfpathlineto{\pgfqpoint{3.437941in}{1.884217in}}%
\pgfpathlineto{\pgfqpoint{3.431173in}{1.881530in}}%
\pgfpathlineto{\pgfqpoint{3.432335in}{1.876511in}}%
\pgfpathlineto{\pgfqpoint{3.424616in}{1.868740in}}%
\pgfpathlineto{\pgfqpoint{3.421501in}{1.872886in}}%
\pgfpathlineto{\pgfqpoint{3.415970in}{1.857152in}}%
\pgfpathlineto{\pgfqpoint{3.413847in}{1.847236in}}%
\pgfpathlineto{\pgfqpoint{3.421042in}{1.838069in}}%
\pgfpathlineto{\pgfqpoint{3.420582in}{1.833115in}}%
\pgfpathclose%
\pgfusepath{fill}%
\end{pgfscope}%
\begin{pgfscope}%
\pgfpathrectangle{\pgfqpoint{0.100000in}{0.100000in}}{\pgfqpoint{3.608454in}{2.310000in}}%
\pgfusepath{clip}%
\pgfsetbuttcap%
\pgfsetmiterjoin%
\definecolor{currentfill}{rgb}{0.000000,0.666667,0.666667}%
\pgfsetfillcolor{currentfill}%
\pgfsetlinewidth{0.000000pt}%
\definecolor{currentstroke}{rgb}{0.000000,0.000000,0.000000}%
\pgfsetstrokecolor{currentstroke}%
\pgfsetstrokeopacity{0.000000}%
\pgfsetdash{}{0pt}%
\pgfpathmoveto{\pgfqpoint{3.043474in}{1.583306in}}%
\pgfpathlineto{\pgfqpoint{3.040035in}{1.579416in}}%
\pgfpathlineto{\pgfqpoint{3.028102in}{1.574866in}}%
\pgfpathlineto{\pgfqpoint{3.020349in}{1.574020in}}%
\pgfpathlineto{\pgfqpoint{3.013646in}{1.587974in}}%
\pgfpathlineto{\pgfqpoint{2.995768in}{1.584878in}}%
\pgfpathlineto{\pgfqpoint{2.991632in}{1.609704in}}%
\pgfpathlineto{\pgfqpoint{2.984410in}{1.607734in}}%
\pgfpathlineto{\pgfqpoint{2.961131in}{1.605506in}}%
\pgfpathlineto{\pgfqpoint{2.955120in}{1.642813in}}%
\pgfpathlineto{\pgfqpoint{2.973125in}{1.655287in}}%
\pgfpathlineto{\pgfqpoint{2.989337in}{1.668857in}}%
\pgfpathlineto{\pgfqpoint{2.999470in}{1.675879in}}%
\pgfpathlineto{\pgfqpoint{3.011435in}{1.687328in}}%
\pgfpathlineto{\pgfqpoint{3.021181in}{1.698369in}}%
\pgfpathlineto{\pgfqpoint{3.027192in}{1.702502in}}%
\pgfpathlineto{\pgfqpoint{3.032000in}{1.700788in}}%
\pgfpathlineto{\pgfqpoint{3.039413in}{1.658799in}}%
\pgfpathlineto{\pgfqpoint{3.047674in}{1.660178in}}%
\pgfpathlineto{\pgfqpoint{3.050043in}{1.647962in}}%
\pgfpathlineto{\pgfqpoint{3.048278in}{1.646809in}}%
\pgfpathlineto{\pgfqpoint{3.050721in}{1.630639in}}%
\pgfpathlineto{\pgfqpoint{3.051910in}{1.617021in}}%
\pgfpathlineto{\pgfqpoint{3.045476in}{1.613591in}}%
\pgfpathlineto{\pgfqpoint{3.046524in}{1.607153in}}%
\pgfpathlineto{\pgfqpoint{3.039983in}{1.605278in}}%
\pgfpathlineto{\pgfqpoint{3.043474in}{1.583306in}}%
\pgfpathclose%
\pgfusepath{fill}%
\end{pgfscope}%
\begin{pgfscope}%
\pgfpathrectangle{\pgfqpoint{0.100000in}{0.100000in}}{\pgfqpoint{3.608454in}{2.310000in}}%
\pgfusepath{clip}%
\pgfsetbuttcap%
\pgfsetmiterjoin%
\definecolor{currentfill}{rgb}{0.000000,0.333333,0.833333}%
\pgfsetfillcolor{currentfill}%
\pgfsetlinewidth{0.000000pt}%
\definecolor{currentstroke}{rgb}{0.000000,0.000000,0.000000}%
\pgfsetstrokecolor{currentstroke}%
\pgfsetstrokeopacity{0.000000}%
\pgfsetdash{}{0pt}%
\pgfpathmoveto{\pgfqpoint{1.992970in}{1.502588in}}%
\pgfpathlineto{\pgfqpoint{1.965564in}{1.503301in}}%
\pgfpathlineto{\pgfqpoint{1.966281in}{1.530825in}}%
\pgfpathlineto{\pgfqpoint{1.952480in}{1.526297in}}%
\pgfpathlineto{\pgfqpoint{1.952621in}{1.531206in}}%
\pgfpathlineto{\pgfqpoint{1.946376in}{1.531379in}}%
\pgfpathlineto{\pgfqpoint{1.946070in}{1.524600in}}%
\pgfpathlineto{\pgfqpoint{1.925386in}{1.523978in}}%
\pgfpathlineto{\pgfqpoint{1.922421in}{1.521780in}}%
\pgfpathlineto{\pgfqpoint{1.912100in}{1.522170in}}%
\pgfpathlineto{\pgfqpoint{1.911513in}{1.532437in}}%
\pgfpathlineto{\pgfqpoint{1.912258in}{1.559919in}}%
\pgfpathlineto{\pgfqpoint{1.912766in}{1.573654in}}%
\pgfpathlineto{\pgfqpoint{1.939982in}{1.572847in}}%
\pgfpathlineto{\pgfqpoint{1.939698in}{1.559064in}}%
\pgfpathlineto{\pgfqpoint{1.987553in}{1.557852in}}%
\pgfpathlineto{\pgfqpoint{1.994270in}{1.557682in}}%
\pgfpathlineto{\pgfqpoint{1.992970in}{1.502588in}}%
\pgfpathclose%
\pgfusepath{fill}%
\end{pgfscope}%
\begin{pgfscope}%
\pgfpathrectangle{\pgfqpoint{0.100000in}{0.100000in}}{\pgfqpoint{3.608454in}{2.310000in}}%
\pgfusepath{clip}%
\pgfsetbuttcap%
\pgfsetmiterjoin%
\definecolor{currentfill}{rgb}{0.000000,0.462745,0.768627}%
\pgfsetfillcolor{currentfill}%
\pgfsetlinewidth{0.000000pt}%
\definecolor{currentstroke}{rgb}{0.000000,0.000000,0.000000}%
\pgfsetstrokecolor{currentstroke}%
\pgfsetstrokeopacity{0.000000}%
\pgfsetdash{}{0pt}%
\pgfpathmoveto{\pgfqpoint{2.115120in}{1.723283in}}%
\pgfpathlineto{\pgfqpoint{2.114776in}{1.743905in}}%
\pgfpathlineto{\pgfqpoint{2.101083in}{1.743915in}}%
\pgfpathlineto{\pgfqpoint{2.100603in}{1.750815in}}%
\pgfpathlineto{\pgfqpoint{2.100687in}{1.757698in}}%
\pgfpathlineto{\pgfqpoint{2.114393in}{1.757738in}}%
\pgfpathlineto{\pgfqpoint{2.114392in}{1.774876in}}%
\pgfpathlineto{\pgfqpoint{2.118639in}{1.771585in}}%
\pgfpathlineto{\pgfqpoint{2.128088in}{1.771585in}}%
\pgfpathlineto{\pgfqpoint{2.167464in}{1.771934in}}%
\pgfpathlineto{\pgfqpoint{2.168433in}{1.778817in}}%
\pgfpathlineto{\pgfqpoint{2.190305in}{1.779149in}}%
\pgfpathlineto{\pgfqpoint{2.190701in}{1.751566in}}%
\pgfpathlineto{\pgfqpoint{2.176891in}{1.751349in}}%
\pgfpathlineto{\pgfqpoint{2.177278in}{1.723751in}}%
\pgfpathlineto{\pgfqpoint{2.115120in}{1.723283in}}%
\pgfpathclose%
\pgfusepath{fill}%
\end{pgfscope}%
\begin{pgfscope}%
\pgfpathrectangle{\pgfqpoint{0.100000in}{0.100000in}}{\pgfqpoint{3.608454in}{2.310000in}}%
\pgfusepath{clip}%
\pgfsetbuttcap%
\pgfsetmiterjoin%
\definecolor{currentfill}{rgb}{0.000000,0.529412,0.735294}%
\pgfsetfillcolor{currentfill}%
\pgfsetlinewidth{0.000000pt}%
\definecolor{currentstroke}{rgb}{0.000000,0.000000,0.000000}%
\pgfsetstrokecolor{currentstroke}%
\pgfsetstrokeopacity{0.000000}%
\pgfsetdash{}{0pt}%
\pgfpathmoveto{\pgfqpoint{2.248598in}{1.467244in}}%
\pgfpathlineto{\pgfqpoint{2.247888in}{1.491592in}}%
\pgfpathlineto{\pgfqpoint{2.220496in}{1.490915in}}%
\pgfpathlineto{\pgfqpoint{2.219965in}{1.511654in}}%
\pgfpathlineto{\pgfqpoint{2.206358in}{1.511376in}}%
\pgfpathlineto{\pgfqpoint{2.205803in}{1.538835in}}%
\pgfpathlineto{\pgfqpoint{2.266869in}{1.540448in}}%
\pgfpathlineto{\pgfqpoint{2.287673in}{1.541193in}}%
\pgfpathlineto{\pgfqpoint{2.288611in}{1.513740in}}%
\pgfpathlineto{\pgfqpoint{2.302293in}{1.514165in}}%
\pgfpathlineto{\pgfqpoint{2.303069in}{1.470734in}}%
\pgfpathlineto{\pgfqpoint{2.290253in}{1.469670in}}%
\pgfpathlineto{\pgfqpoint{2.248598in}{1.467244in}}%
\pgfpathclose%
\pgfusepath{fill}%
\end{pgfscope}%
\begin{pgfscope}%
\pgfpathrectangle{\pgfqpoint{0.100000in}{0.100000in}}{\pgfqpoint{3.608454in}{2.310000in}}%
\pgfusepath{clip}%
\pgfsetbuttcap%
\pgfsetmiterjoin%
\definecolor{currentfill}{rgb}{0.000000,0.627451,0.686275}%
\pgfsetfillcolor{currentfill}%
\pgfsetlinewidth{0.000000pt}%
\definecolor{currentstroke}{rgb}{0.000000,0.000000,0.000000}%
\pgfsetstrokecolor{currentstroke}%
\pgfsetstrokeopacity{0.000000}%
\pgfsetdash{}{0pt}%
\pgfpathmoveto{\pgfqpoint{1.141273in}{1.664762in}}%
\pgfpathlineto{\pgfqpoint{1.142493in}{1.660281in}}%
\pgfpathlineto{\pgfqpoint{1.139137in}{1.653332in}}%
\pgfpathlineto{\pgfqpoint{1.141060in}{1.650733in}}%
\pgfpathlineto{\pgfqpoint{1.142480in}{1.636751in}}%
\pgfpathlineto{\pgfqpoint{1.138324in}{1.627635in}}%
\pgfpathlineto{\pgfqpoint{1.132825in}{1.619941in}}%
\pgfpathlineto{\pgfqpoint{1.121923in}{1.621467in}}%
\pgfpathlineto{\pgfqpoint{1.117387in}{1.618454in}}%
\pgfpathlineto{\pgfqpoint{1.110470in}{1.623700in}}%
\pgfpathlineto{\pgfqpoint{1.101371in}{1.618838in}}%
\pgfpathlineto{\pgfqpoint{1.089153in}{1.621193in}}%
\pgfpathlineto{\pgfqpoint{1.070321in}{1.603891in}}%
\pgfpathlineto{\pgfqpoint{1.051249in}{1.601475in}}%
\pgfpathlineto{\pgfqpoint{0.978586in}{1.616384in}}%
\pgfpathlineto{\pgfqpoint{0.994964in}{1.693281in}}%
\pgfpathlineto{\pgfqpoint{1.004416in}{1.690817in}}%
\pgfpathlineto{\pgfqpoint{1.054992in}{1.681239in}}%
\pgfpathlineto{\pgfqpoint{1.060120in}{1.706792in}}%
\pgfpathlineto{\pgfqpoint{1.080051in}{1.702943in}}%
\pgfpathlineto{\pgfqpoint{1.081382in}{1.709710in}}%
\pgfpathlineto{\pgfqpoint{1.090342in}{1.707956in}}%
\pgfpathlineto{\pgfqpoint{1.091640in}{1.714726in}}%
\pgfpathlineto{\pgfqpoint{1.098147in}{1.713450in}}%
\pgfpathlineto{\pgfqpoint{1.101946in}{1.704218in}}%
\pgfpathlineto{\pgfqpoint{1.109727in}{1.693683in}}%
\pgfpathlineto{\pgfqpoint{1.117899in}{1.692151in}}%
\pgfpathlineto{\pgfqpoint{1.122487in}{1.688955in}}%
\pgfpathlineto{\pgfqpoint{1.126068in}{1.701146in}}%
\pgfpathlineto{\pgfqpoint{1.141976in}{1.698141in}}%
\pgfpathlineto{\pgfqpoint{1.143732in}{1.692681in}}%
\pgfpathlineto{\pgfqpoint{1.141935in}{1.686685in}}%
\pgfpathlineto{\pgfqpoint{1.137407in}{1.682303in}}%
\pgfpathlineto{\pgfqpoint{1.137582in}{1.673009in}}%
\pgfpathlineto{\pgfqpoint{1.141528in}{1.670743in}}%
\pgfpathlineto{\pgfqpoint{1.141273in}{1.664762in}}%
\pgfpathclose%
\pgfusepath{fill}%
\end{pgfscope}%
\begin{pgfscope}%
\pgfpathrectangle{\pgfqpoint{0.100000in}{0.100000in}}{\pgfqpoint{3.608454in}{2.310000in}}%
\pgfusepath{clip}%
\pgfsetbuttcap%
\pgfsetmiterjoin%
\definecolor{currentfill}{rgb}{0.000000,0.713725,0.643137}%
\pgfsetfillcolor{currentfill}%
\pgfsetlinewidth{0.000000pt}%
\definecolor{currentstroke}{rgb}{0.000000,0.000000,0.000000}%
\pgfsetstrokecolor{currentstroke}%
\pgfsetstrokeopacity{0.000000}%
\pgfsetdash{}{0pt}%
\pgfpathmoveto{\pgfqpoint{2.811045in}{1.319515in}}%
\pgfpathlineto{\pgfqpoint{2.804061in}{1.326162in}}%
\pgfpathlineto{\pgfqpoint{2.796848in}{1.326453in}}%
\pgfpathlineto{\pgfqpoint{2.787097in}{1.337635in}}%
\pgfpathlineto{\pgfqpoint{2.789928in}{1.341177in}}%
\pgfpathlineto{\pgfqpoint{2.786720in}{1.347970in}}%
\pgfpathlineto{\pgfqpoint{2.779558in}{1.348364in}}%
\pgfpathlineto{\pgfqpoint{2.772188in}{1.346181in}}%
\pgfpathlineto{\pgfqpoint{2.768060in}{1.364802in}}%
\pgfpathlineto{\pgfqpoint{2.777359in}{1.362093in}}%
\pgfpathlineto{\pgfqpoint{2.784131in}{1.363939in}}%
\pgfpathlineto{\pgfqpoint{2.790830in}{1.362794in}}%
\pgfpathlineto{\pgfqpoint{2.798224in}{1.355170in}}%
\pgfpathlineto{\pgfqpoint{2.804813in}{1.353816in}}%
\pgfpathlineto{\pgfqpoint{2.807405in}{1.358642in}}%
\pgfpathlineto{\pgfqpoint{2.812960in}{1.360831in}}%
\pgfpathlineto{\pgfqpoint{2.829317in}{1.356171in}}%
\pgfpathlineto{\pgfqpoint{2.836077in}{1.357238in}}%
\pgfpathlineto{\pgfqpoint{2.842704in}{1.366486in}}%
\pgfpathlineto{\pgfqpoint{2.841606in}{1.355934in}}%
\pgfpathlineto{\pgfqpoint{2.835681in}{1.347268in}}%
\pgfpathlineto{\pgfqpoint{2.833231in}{1.340587in}}%
\pgfpathlineto{\pgfqpoint{2.834430in}{1.334621in}}%
\pgfpathlineto{\pgfqpoint{2.826872in}{1.332234in}}%
\pgfpathlineto{\pgfqpoint{2.822416in}{1.337584in}}%
\pgfpathlineto{\pgfqpoint{2.817804in}{1.333230in}}%
\pgfpathlineto{\pgfqpoint{2.818045in}{1.327218in}}%
\pgfpathlineto{\pgfqpoint{2.811045in}{1.319515in}}%
\pgfpathclose%
\pgfusepath{fill}%
\end{pgfscope}%
\begin{pgfscope}%
\pgfpathrectangle{\pgfqpoint{0.100000in}{0.100000in}}{\pgfqpoint{3.608454in}{2.310000in}}%
\pgfusepath{clip}%
\pgfsetbuttcap%
\pgfsetmiterjoin%
\definecolor{currentfill}{rgb}{0.000000,0.560784,0.719608}%
\pgfsetfillcolor{currentfill}%
\pgfsetlinewidth{0.000000pt}%
\definecolor{currentstroke}{rgb}{0.000000,0.000000,0.000000}%
\pgfsetstrokecolor{currentstroke}%
\pgfsetstrokeopacity{0.000000}%
\pgfsetdash{}{0pt}%
\pgfpathmoveto{\pgfqpoint{2.839212in}{0.870674in}}%
\pgfpathlineto{\pgfqpoint{2.834772in}{0.868856in}}%
\pgfpathlineto{\pgfqpoint{2.830726in}{0.859944in}}%
\pgfpathlineto{\pgfqpoint{2.822592in}{0.856703in}}%
\pgfpathlineto{\pgfqpoint{2.816271in}{0.861716in}}%
\pgfpathlineto{\pgfqpoint{2.812799in}{0.870744in}}%
\pgfpathlineto{\pgfqpoint{2.821573in}{0.886623in}}%
\pgfpathlineto{\pgfqpoint{2.827685in}{0.882746in}}%
\pgfpathlineto{\pgfqpoint{2.830174in}{0.891948in}}%
\pgfpathlineto{\pgfqpoint{2.830655in}{0.902493in}}%
\pgfpathlineto{\pgfqpoint{2.835752in}{0.904367in}}%
\pgfpathlineto{\pgfqpoint{2.833490in}{0.924483in}}%
\pgfpathlineto{\pgfqpoint{2.844937in}{0.925826in}}%
\pgfpathlineto{\pgfqpoint{2.848183in}{0.924494in}}%
\pgfpathlineto{\pgfqpoint{2.849065in}{0.920723in}}%
\pgfpathlineto{\pgfqpoint{2.866524in}{0.926111in}}%
\pgfpathlineto{\pgfqpoint{2.874063in}{0.928161in}}%
\pgfpathlineto{\pgfqpoint{2.881402in}{0.908328in}}%
\pgfpathlineto{\pgfqpoint{2.878409in}{0.905708in}}%
\pgfpathlineto{\pgfqpoint{2.893588in}{0.882495in}}%
\pgfpathlineto{\pgfqpoint{2.901077in}{0.870542in}}%
\pgfpathlineto{\pgfqpoint{2.898491in}{0.872568in}}%
\pgfpathlineto{\pgfqpoint{2.890867in}{0.862703in}}%
\pgfpathlineto{\pgfqpoint{2.887187in}{0.854824in}}%
\pgfpathlineto{\pgfqpoint{2.893436in}{0.849951in}}%
\pgfpathlineto{\pgfqpoint{2.891100in}{0.843787in}}%
\pgfpathlineto{\pgfqpoint{2.872842in}{0.842457in}}%
\pgfpathlineto{\pgfqpoint{2.871284in}{0.855909in}}%
\pgfpathlineto{\pgfqpoint{2.855107in}{0.854082in}}%
\pgfpathlineto{\pgfqpoint{2.854444in}{0.865881in}}%
\pgfpathlineto{\pgfqpoint{2.851231in}{0.869522in}}%
\pgfpathlineto{\pgfqpoint{2.839212in}{0.870674in}}%
\pgfpathclose%
\pgfusepath{fill}%
\end{pgfscope}%
\begin{pgfscope}%
\pgfpathrectangle{\pgfqpoint{0.100000in}{0.100000in}}{\pgfqpoint{3.608454in}{2.310000in}}%
\pgfusepath{clip}%
\pgfsetbuttcap%
\pgfsetmiterjoin%
\definecolor{currentfill}{rgb}{0.000000,0.419608,0.790196}%
\pgfsetfillcolor{currentfill}%
\pgfsetlinewidth{0.000000pt}%
\definecolor{currentstroke}{rgb}{0.000000,0.000000,0.000000}%
\pgfsetstrokecolor{currentstroke}%
\pgfsetstrokeopacity{0.000000}%
\pgfsetdash{}{0pt}%
\pgfpathmoveto{\pgfqpoint{1.981983in}{1.828333in}}%
\pgfpathlineto{\pgfqpoint{1.939662in}{1.829664in}}%
\pgfpathlineto{\pgfqpoint{1.941869in}{1.891889in}}%
\pgfpathlineto{\pgfqpoint{1.983409in}{1.890641in}}%
\pgfpathlineto{\pgfqpoint{2.020229in}{1.889881in}}%
\pgfpathlineto{\pgfqpoint{2.018973in}{1.880728in}}%
\pgfpathlineto{\pgfqpoint{2.015031in}{1.874970in}}%
\pgfpathlineto{\pgfqpoint{2.004519in}{1.867267in}}%
\pgfpathlineto{\pgfqpoint{2.003389in}{1.863983in}}%
\pgfpathlineto{\pgfqpoint{2.012969in}{1.848244in}}%
\pgfpathlineto{\pgfqpoint{2.021780in}{1.845315in}}%
\pgfpathlineto{\pgfqpoint{2.024583in}{1.841364in}}%
\pgfpathlineto{\pgfqpoint{1.995378in}{1.841947in}}%
\pgfpathlineto{\pgfqpoint{1.994509in}{1.839579in}}%
\pgfpathlineto{\pgfqpoint{1.982265in}{1.839933in}}%
\pgfpathlineto{\pgfqpoint{1.981983in}{1.828333in}}%
\pgfpathclose%
\pgfusepath{fill}%
\end{pgfscope}%
\begin{pgfscope}%
\pgfpathrectangle{\pgfqpoint{0.100000in}{0.100000in}}{\pgfqpoint{3.608454in}{2.310000in}}%
\pgfusepath{clip}%
\pgfsetbuttcap%
\pgfsetmiterjoin%
\definecolor{currentfill}{rgb}{0.000000,0.682353,0.658824}%
\pgfsetfillcolor{currentfill}%
\pgfsetlinewidth{0.000000pt}%
\definecolor{currentstroke}{rgb}{0.000000,0.000000,0.000000}%
\pgfsetstrokecolor{currentstroke}%
\pgfsetstrokeopacity{0.000000}%
\pgfsetdash{}{0pt}%
\pgfpathmoveto{\pgfqpoint{2.971951in}{1.538265in}}%
\pgfpathlineto{\pgfqpoint{2.963825in}{1.532458in}}%
\pgfpathlineto{\pgfqpoint{2.952109in}{1.531961in}}%
\pgfpathlineto{\pgfqpoint{2.948242in}{1.534882in}}%
\pgfpathlineto{\pgfqpoint{2.947280in}{1.541412in}}%
\pgfpathlineto{\pgfqpoint{2.937196in}{1.539921in}}%
\pgfpathlineto{\pgfqpoint{2.934085in}{1.560268in}}%
\pgfpathlineto{\pgfqpoint{2.939221in}{1.561031in}}%
\pgfpathlineto{\pgfqpoint{2.932782in}{1.601168in}}%
\pgfpathlineto{\pgfqpoint{2.961131in}{1.605506in}}%
\pgfpathlineto{\pgfqpoint{2.984410in}{1.607734in}}%
\pgfpathlineto{\pgfqpoint{2.991632in}{1.609704in}}%
\pgfpathlineto{\pgfqpoint{2.995768in}{1.584878in}}%
\pgfpathlineto{\pgfqpoint{2.988249in}{1.569893in}}%
\pgfpathlineto{\pgfqpoint{2.990516in}{1.558619in}}%
\pgfpathlineto{\pgfqpoint{2.969282in}{1.554858in}}%
\pgfpathlineto{\pgfqpoint{2.971951in}{1.538265in}}%
\pgfpathclose%
\pgfusepath{fill}%
\end{pgfscope}%
\begin{pgfscope}%
\pgfpathrectangle{\pgfqpoint{0.100000in}{0.100000in}}{\pgfqpoint{3.608454in}{2.310000in}}%
\pgfusepath{clip}%
\pgfsetbuttcap%
\pgfsetmiterjoin%
\definecolor{currentfill}{rgb}{0.000000,0.443137,0.778431}%
\pgfsetfillcolor{currentfill}%
\pgfsetlinewidth{0.000000pt}%
\definecolor{currentstroke}{rgb}{0.000000,0.000000,0.000000}%
\pgfsetstrokecolor{currentstroke}%
\pgfsetstrokeopacity{0.000000}%
\pgfsetdash{}{0pt}%
\pgfpathmoveto{\pgfqpoint{1.643808in}{1.676503in}}%
\pgfpathlineto{\pgfqpoint{1.583107in}{1.682203in}}%
\pgfpathlineto{\pgfqpoint{1.586973in}{1.721839in}}%
\pgfpathlineto{\pgfqpoint{1.592333in}{1.775294in}}%
\pgfpathlineto{\pgfqpoint{1.598608in}{1.839877in}}%
\pgfpathlineto{\pgfqpoint{1.601301in}{1.856823in}}%
\pgfpathlineto{\pgfqpoint{1.661648in}{1.851140in}}%
\pgfpathlineto{\pgfqpoint{1.660454in}{1.837408in}}%
\pgfpathlineto{\pgfqpoint{1.714041in}{1.832923in}}%
\pgfpathlineto{\pgfqpoint{1.710799in}{1.791190in}}%
\pgfpathlineto{\pgfqpoint{1.707249in}{1.750308in}}%
\pgfpathlineto{\pgfqpoint{1.704828in}{1.727690in}}%
\pgfpathlineto{\pgfqpoint{1.697918in}{1.727667in}}%
\pgfpathlineto{\pgfqpoint{1.695616in}{1.726789in}}%
\pgfpathlineto{\pgfqpoint{1.659487in}{1.729827in}}%
\pgfpathlineto{\pgfqpoint{1.654175in}{1.729234in}}%
\pgfpathlineto{\pgfqpoint{1.648016in}{1.724287in}}%
\pgfpathlineto{\pgfqpoint{1.643808in}{1.676503in}}%
\pgfpathclose%
\pgfusepath{fill}%
\end{pgfscope}%
\begin{pgfscope}%
\pgfpathrectangle{\pgfqpoint{0.100000in}{0.100000in}}{\pgfqpoint{3.608454in}{2.310000in}}%
\pgfusepath{clip}%
\pgfsetbuttcap%
\pgfsetmiterjoin%
\definecolor{currentfill}{rgb}{0.000000,0.768627,0.615686}%
\pgfsetfillcolor{currentfill}%
\pgfsetlinewidth{0.000000pt}%
\definecolor{currentstroke}{rgb}{0.000000,0.000000,0.000000}%
\pgfsetstrokecolor{currentstroke}%
\pgfsetstrokeopacity{0.000000}%
\pgfsetdash{}{0pt}%
\pgfpathmoveto{\pgfqpoint{2.867721in}{1.199764in}}%
\pgfpathlineto{\pgfqpoint{2.859737in}{1.192156in}}%
\pgfpathlineto{\pgfqpoint{2.856311in}{1.184739in}}%
\pgfpathlineto{\pgfqpoint{2.851139in}{1.181704in}}%
\pgfpathlineto{\pgfqpoint{2.846368in}{1.183090in}}%
\pgfpathlineto{\pgfqpoint{2.836596in}{1.177507in}}%
\pgfpathlineto{\pgfqpoint{2.827195in}{1.174638in}}%
\pgfpathlineto{\pgfqpoint{2.814063in}{1.180206in}}%
\pgfpathlineto{\pgfqpoint{2.810250in}{1.179188in}}%
\pgfpathlineto{\pgfqpoint{2.804581in}{1.191394in}}%
\pgfpathlineto{\pgfqpoint{2.806029in}{1.196000in}}%
\pgfpathlineto{\pgfqpoint{2.812724in}{1.202234in}}%
\pgfpathlineto{\pgfqpoint{2.815210in}{1.209500in}}%
\pgfpathlineto{\pgfqpoint{2.831285in}{1.222519in}}%
\pgfpathlineto{\pgfqpoint{2.833882in}{1.215723in}}%
\pgfpathlineto{\pgfqpoint{2.833845in}{1.206501in}}%
\pgfpathlineto{\pgfqpoint{2.837048in}{1.201444in}}%
\pgfpathlineto{\pgfqpoint{2.828062in}{1.198088in}}%
\pgfpathlineto{\pgfqpoint{2.827558in}{1.194906in}}%
\pgfpathlineto{\pgfqpoint{2.867721in}{1.199764in}}%
\pgfpathclose%
\pgfusepath{fill}%
\end{pgfscope}%
\begin{pgfscope}%
\pgfpathrectangle{\pgfqpoint{0.100000in}{0.100000in}}{\pgfqpoint{3.608454in}{2.310000in}}%
\pgfusepath{clip}%
\pgfsetbuttcap%
\pgfsetmiterjoin%
\definecolor{currentfill}{rgb}{0.000000,0.784314,0.607843}%
\pgfsetfillcolor{currentfill}%
\pgfsetlinewidth{0.000000pt}%
\definecolor{currentstroke}{rgb}{0.000000,0.000000,0.000000}%
\pgfsetstrokecolor{currentstroke}%
\pgfsetstrokeopacity{0.000000}%
\pgfsetdash{}{0pt}%
\pgfpathmoveto{\pgfqpoint{2.843160in}{1.055221in}}%
\pgfpathlineto{\pgfqpoint{2.833248in}{1.051039in}}%
\pgfpathlineto{\pgfqpoint{2.830948in}{1.045651in}}%
\pgfpathlineto{\pgfqpoint{2.825426in}{1.046467in}}%
\pgfpathlineto{\pgfqpoint{2.820073in}{1.037913in}}%
\pgfpathlineto{\pgfqpoint{2.812171in}{1.037511in}}%
\pgfpathlineto{\pgfqpoint{2.814885in}{1.050035in}}%
\pgfpathlineto{\pgfqpoint{2.808303in}{1.061863in}}%
\pgfpathlineto{\pgfqpoint{2.798375in}{1.062974in}}%
\pgfpathlineto{\pgfqpoint{2.798205in}{1.081874in}}%
\pgfpathlineto{\pgfqpoint{2.802072in}{1.085804in}}%
\pgfpathlineto{\pgfqpoint{2.807766in}{1.084602in}}%
\pgfpathlineto{\pgfqpoint{2.814347in}{1.089123in}}%
\pgfpathlineto{\pgfqpoint{2.819581in}{1.083755in}}%
\pgfpathlineto{\pgfqpoint{2.827437in}{1.087830in}}%
\pgfpathlineto{\pgfqpoint{2.835727in}{1.088275in}}%
\pgfpathlineto{\pgfqpoint{2.834354in}{1.080705in}}%
\pgfpathlineto{\pgfqpoint{2.840206in}{1.081206in}}%
\pgfpathlineto{\pgfqpoint{2.850115in}{1.070383in}}%
\pgfpathlineto{\pgfqpoint{2.842470in}{1.059416in}}%
\pgfpathlineto{\pgfqpoint{2.843160in}{1.055221in}}%
\pgfpathclose%
\pgfusepath{fill}%
\end{pgfscope}%
\begin{pgfscope}%
\pgfpathrectangle{\pgfqpoint{0.100000in}{0.100000in}}{\pgfqpoint{3.608454in}{2.310000in}}%
\pgfusepath{clip}%
\pgfsetbuttcap%
\pgfsetmiterjoin%
\definecolor{currentfill}{rgb}{0.000000,0.478431,0.760784}%
\pgfsetfillcolor{currentfill}%
\pgfsetlinewidth{0.000000pt}%
\definecolor{currentstroke}{rgb}{0.000000,0.000000,0.000000}%
\pgfsetstrokecolor{currentstroke}%
\pgfsetstrokeopacity{0.000000}%
\pgfsetdash{}{0pt}%
\pgfpathmoveto{\pgfqpoint{2.872479in}{1.551436in}}%
\pgfpathlineto{\pgfqpoint{2.869199in}{1.556678in}}%
\pgfpathlineto{\pgfqpoint{2.868477in}{1.562448in}}%
\pgfpathlineto{\pgfqpoint{2.873516in}{1.568184in}}%
\pgfpathlineto{\pgfqpoint{2.879290in}{1.569013in}}%
\pgfpathlineto{\pgfqpoint{2.878552in}{1.574944in}}%
\pgfpathlineto{\pgfqpoint{2.884117in}{1.575727in}}%
\pgfpathlineto{\pgfqpoint{2.883315in}{1.581642in}}%
\pgfpathlineto{\pgfqpoint{2.877765in}{1.580902in}}%
\pgfpathlineto{\pgfqpoint{2.876219in}{1.593055in}}%
\pgfpathlineto{\pgfqpoint{2.889848in}{1.593630in}}%
\pgfpathlineto{\pgfqpoint{2.897159in}{1.600476in}}%
\pgfpathlineto{\pgfqpoint{2.905051in}{1.610697in}}%
\pgfpathlineto{\pgfqpoint{2.928077in}{1.628592in}}%
\pgfpathlineto{\pgfqpoint{2.934104in}{1.630691in}}%
\pgfpathlineto{\pgfqpoint{2.955120in}{1.642813in}}%
\pgfpathlineto{\pgfqpoint{2.961131in}{1.605506in}}%
\pgfpathlineto{\pgfqpoint{2.932782in}{1.601168in}}%
\pgfpathlineto{\pgfqpoint{2.939221in}{1.561031in}}%
\pgfpathlineto{\pgfqpoint{2.934085in}{1.560268in}}%
\pgfpathlineto{\pgfqpoint{2.914636in}{1.557327in}}%
\pgfpathlineto{\pgfqpoint{2.915375in}{1.550896in}}%
\pgfpathlineto{\pgfqpoint{2.901826in}{1.549489in}}%
\pgfpathlineto{\pgfqpoint{2.901003in}{1.555251in}}%
\pgfpathlineto{\pgfqpoint{2.872479in}{1.551436in}}%
\pgfpathclose%
\pgfusepath{fill}%
\end{pgfscope}%
\begin{pgfscope}%
\pgfpathrectangle{\pgfqpoint{0.100000in}{0.100000in}}{\pgfqpoint{3.608454in}{2.310000in}}%
\pgfusepath{clip}%
\pgfsetbuttcap%
\pgfsetmiterjoin%
\definecolor{currentfill}{rgb}{0.000000,0.639216,0.680392}%
\pgfsetfillcolor{currentfill}%
\pgfsetlinewidth{0.000000pt}%
\definecolor{currentstroke}{rgb}{0.000000,0.000000,0.000000}%
\pgfsetstrokecolor{currentstroke}%
\pgfsetstrokeopacity{0.000000}%
\pgfsetdash{}{0pt}%
\pgfpathmoveto{\pgfqpoint{0.737189in}{2.001359in}}%
\pgfpathlineto{\pgfqpoint{0.730749in}{2.003174in}}%
\pgfpathlineto{\pgfqpoint{0.736122in}{2.015946in}}%
\pgfpathlineto{\pgfqpoint{0.729534in}{2.017963in}}%
\pgfpathlineto{\pgfqpoint{0.731442in}{2.024554in}}%
\pgfpathlineto{\pgfqpoint{0.724757in}{2.026561in}}%
\pgfpathlineto{\pgfqpoint{0.737141in}{2.068948in}}%
\pgfpathlineto{\pgfqpoint{0.744663in}{2.068773in}}%
\pgfpathlineto{\pgfqpoint{0.749167in}{2.084439in}}%
\pgfpathlineto{\pgfqpoint{0.761418in}{2.129669in}}%
\pgfpathlineto{\pgfqpoint{0.774845in}{2.127276in}}%
\pgfpathlineto{\pgfqpoint{0.782427in}{2.131997in}}%
\pgfpathlineto{\pgfqpoint{0.784697in}{2.127385in}}%
\pgfpathlineto{\pgfqpoint{0.790335in}{2.130635in}}%
\pgfpathlineto{\pgfqpoint{0.823108in}{2.121596in}}%
\pgfpathlineto{\pgfqpoint{0.851613in}{2.114354in}}%
\pgfpathlineto{\pgfqpoint{0.849420in}{2.114482in}}%
\pgfpathlineto{\pgfqpoint{0.849927in}{2.099792in}}%
\pgfpathlineto{\pgfqpoint{0.868296in}{2.100626in}}%
\pgfpathlineto{\pgfqpoint{0.864637in}{2.088901in}}%
\pgfpathlineto{\pgfqpoint{0.871344in}{2.087217in}}%
\pgfpathlineto{\pgfqpoint{0.872020in}{2.076043in}}%
\pgfpathlineto{\pgfqpoint{0.875412in}{2.075250in}}%
\pgfpathlineto{\pgfqpoint{0.869214in}{2.048910in}}%
\pgfpathlineto{\pgfqpoint{0.849176in}{2.054195in}}%
\pgfpathlineto{\pgfqpoint{0.846645in}{2.043323in}}%
\pgfpathlineto{\pgfqpoint{0.841907in}{2.040973in}}%
\pgfpathlineto{\pgfqpoint{0.838162in}{2.031358in}}%
\pgfpathlineto{\pgfqpoint{0.835457in}{2.032072in}}%
\pgfpathlineto{\pgfqpoint{0.831103in}{2.015303in}}%
\pgfpathlineto{\pgfqpoint{0.825859in}{2.013085in}}%
\pgfpathlineto{\pgfqpoint{0.813450in}{2.016488in}}%
\pgfpathlineto{\pgfqpoint{0.811964in}{2.010809in}}%
\pgfpathlineto{\pgfqpoint{0.797207in}{2.013901in}}%
\pgfpathlineto{\pgfqpoint{0.796435in}{2.001902in}}%
\pgfpathlineto{\pgfqpoint{0.802344in}{2.000341in}}%
\pgfpathlineto{\pgfqpoint{0.799750in}{1.984499in}}%
\pgfpathlineto{\pgfqpoint{0.737189in}{2.001359in}}%
\pgfpathclose%
\pgfusepath{fill}%
\end{pgfscope}%
\begin{pgfscope}%
\pgfpathrectangle{\pgfqpoint{0.100000in}{0.100000in}}{\pgfqpoint{3.608454in}{2.310000in}}%
\pgfusepath{clip}%
\pgfsetbuttcap%
\pgfsetmiterjoin%
\definecolor{currentfill}{rgb}{0.000000,0.670588,0.664706}%
\pgfsetfillcolor{currentfill}%
\pgfsetlinewidth{0.000000pt}%
\definecolor{currentstroke}{rgb}{0.000000,0.000000,0.000000}%
\pgfsetstrokecolor{currentstroke}%
\pgfsetstrokeopacity{0.000000}%
\pgfsetdash{}{0pt}%
\pgfpathmoveto{\pgfqpoint{1.601877in}{1.254860in}}%
\pgfpathlineto{\pgfqpoint{1.596724in}{1.204258in}}%
\pgfpathlineto{\pgfqpoint{1.538731in}{1.209408in}}%
\pgfpathlineto{\pgfqpoint{1.513551in}{1.211748in}}%
\pgfpathlineto{\pgfqpoint{1.466660in}{1.216957in}}%
\pgfpathlineto{\pgfqpoint{1.469341in}{1.240362in}}%
\pgfpathlineto{\pgfqpoint{1.471264in}{1.248705in}}%
\pgfpathlineto{\pgfqpoint{1.469792in}{1.261228in}}%
\pgfpathlineto{\pgfqpoint{1.463748in}{1.269672in}}%
\pgfpathlineto{\pgfqpoint{1.450742in}{1.268060in}}%
\pgfpathlineto{\pgfqpoint{1.457227in}{1.286192in}}%
\pgfpathlineto{\pgfqpoint{1.454881in}{1.289984in}}%
\pgfpathlineto{\pgfqpoint{1.458524in}{1.289130in}}%
\pgfpathlineto{\pgfqpoint{1.470019in}{1.293572in}}%
\pgfpathlineto{\pgfqpoint{1.474843in}{1.297491in}}%
\pgfpathlineto{\pgfqpoint{1.481289in}{1.288540in}}%
\pgfpathlineto{\pgfqpoint{1.483235in}{1.285604in}}%
\pgfpathlineto{\pgfqpoint{1.506183in}{1.284643in}}%
\pgfpathlineto{\pgfqpoint{1.541303in}{1.267794in}}%
\pgfpathlineto{\pgfqpoint{1.540663in}{1.260697in}}%
\pgfpathlineto{\pgfqpoint{1.601877in}{1.254860in}}%
\pgfpathclose%
\pgfusepath{fill}%
\end{pgfscope}%
\begin{pgfscope}%
\pgfpathrectangle{\pgfqpoint{0.100000in}{0.100000in}}{\pgfqpoint{3.608454in}{2.310000in}}%
\pgfusepath{clip}%
\pgfsetbuttcap%
\pgfsetmiterjoin%
\definecolor{currentfill}{rgb}{0.000000,0.380392,0.809804}%
\pgfsetfillcolor{currentfill}%
\pgfsetlinewidth{0.000000pt}%
\definecolor{currentstroke}{rgb}{0.000000,0.000000,0.000000}%
\pgfsetstrokecolor{currentstroke}%
\pgfsetstrokeopacity{0.000000}%
\pgfsetdash{}{0pt}%
\pgfpathmoveto{\pgfqpoint{2.347355in}{1.241765in}}%
\pgfpathlineto{\pgfqpoint{2.337514in}{1.241571in}}%
\pgfpathlineto{\pgfqpoint{2.337381in}{1.248460in}}%
\pgfpathlineto{\pgfqpoint{2.323611in}{1.248060in}}%
\pgfpathlineto{\pgfqpoint{2.322748in}{1.276780in}}%
\pgfpathlineto{\pgfqpoint{2.315298in}{1.276599in}}%
\pgfpathlineto{\pgfqpoint{2.314915in}{1.287249in}}%
\pgfpathlineto{\pgfqpoint{2.313973in}{1.320063in}}%
\pgfpathlineto{\pgfqpoint{2.319587in}{1.318096in}}%
\pgfpathlineto{\pgfqpoint{2.327621in}{1.321058in}}%
\pgfpathlineto{\pgfqpoint{2.327371in}{1.331768in}}%
\pgfpathlineto{\pgfqpoint{2.336652in}{1.332026in}}%
\pgfpathlineto{\pgfqpoint{2.335984in}{1.355415in}}%
\pgfpathlineto{\pgfqpoint{2.340562in}{1.355530in}}%
\pgfpathlineto{\pgfqpoint{2.340427in}{1.362513in}}%
\pgfpathlineto{\pgfqpoint{2.368642in}{1.363523in}}%
\pgfpathlineto{\pgfqpoint{2.371543in}{1.352918in}}%
\pgfpathlineto{\pgfqpoint{2.369779in}{1.350114in}}%
\pgfpathlineto{\pgfqpoint{2.373739in}{1.339722in}}%
\pgfpathlineto{\pgfqpoint{2.380869in}{1.336490in}}%
\pgfpathlineto{\pgfqpoint{2.386232in}{1.344093in}}%
\pgfpathlineto{\pgfqpoint{2.398503in}{1.340979in}}%
\pgfpathlineto{\pgfqpoint{2.407445in}{1.335820in}}%
\pgfpathlineto{\pgfqpoint{2.402286in}{1.326271in}}%
\pgfpathlineto{\pgfqpoint{2.404058in}{1.320682in}}%
\pgfpathlineto{\pgfqpoint{2.433309in}{1.321900in}}%
\pgfpathlineto{\pgfqpoint{2.435549in}{1.287521in}}%
\pgfpathlineto{\pgfqpoint{2.414877in}{1.286654in}}%
\pgfpathlineto{\pgfqpoint{2.415279in}{1.279749in}}%
\pgfpathlineto{\pgfqpoint{2.404880in}{1.275417in}}%
\pgfpathlineto{\pgfqpoint{2.397241in}{1.275533in}}%
\pgfpathlineto{\pgfqpoint{2.392056in}{1.271141in}}%
\pgfpathlineto{\pgfqpoint{2.381801in}{1.267561in}}%
\pgfpathlineto{\pgfqpoint{2.368816in}{1.282820in}}%
\pgfpathlineto{\pgfqpoint{2.349233in}{1.281942in}}%
\pgfpathlineto{\pgfqpoint{2.350544in}{1.245278in}}%
\pgfpathlineto{\pgfqpoint{2.347355in}{1.241765in}}%
\pgfpathclose%
\pgfusepath{fill}%
\end{pgfscope}%
\begin{pgfscope}%
\pgfpathrectangle{\pgfqpoint{0.100000in}{0.100000in}}{\pgfqpoint{3.608454in}{2.310000in}}%
\pgfusepath{clip}%
\pgfsetbuttcap%
\pgfsetmiterjoin%
\definecolor{currentfill}{rgb}{0.000000,0.301961,0.849020}%
\pgfsetfillcolor{currentfill}%
\pgfsetlinewidth{0.000000pt}%
\definecolor{currentstroke}{rgb}{0.000000,0.000000,0.000000}%
\pgfsetstrokecolor{currentstroke}%
\pgfsetstrokeopacity{0.000000}%
\pgfsetdash{}{0pt}%
\pgfpathmoveto{\pgfqpoint{2.130827in}{1.263573in}}%
\pgfpathlineto{\pgfqpoint{2.101955in}{1.263584in}}%
\pgfpathlineto{\pgfqpoint{2.102752in}{1.281057in}}%
\pgfpathlineto{\pgfqpoint{2.102737in}{1.315533in}}%
\pgfpathlineto{\pgfqpoint{2.103312in}{1.318980in}}%
\pgfpathlineto{\pgfqpoint{2.103336in}{1.338237in}}%
\pgfpathlineto{\pgfqpoint{2.112374in}{1.338940in}}%
\pgfpathlineto{\pgfqpoint{2.112893in}{1.355625in}}%
\pgfpathlineto{\pgfqpoint{2.116197in}{1.360877in}}%
\pgfpathlineto{\pgfqpoint{2.107556in}{1.374791in}}%
\pgfpathlineto{\pgfqpoint{2.100521in}{1.381808in}}%
\pgfpathlineto{\pgfqpoint{2.131100in}{1.381764in}}%
\pgfpathlineto{\pgfqpoint{2.131152in}{1.375734in}}%
\pgfpathlineto{\pgfqpoint{2.154965in}{1.375795in}}%
\pgfpathlineto{\pgfqpoint{2.155044in}{1.381527in}}%
\pgfpathlineto{\pgfqpoint{2.182439in}{1.381621in}}%
\pgfpathlineto{\pgfqpoint{2.182827in}{1.356530in}}%
\pgfpathlineto{\pgfqpoint{2.188988in}{1.359699in}}%
\pgfpathlineto{\pgfqpoint{2.198627in}{1.357636in}}%
\pgfpathlineto{\pgfqpoint{2.199219in}{1.335898in}}%
\pgfpathlineto{\pgfqpoint{2.198808in}{1.305319in}}%
\pgfpathlineto{\pgfqpoint{2.164617in}{1.305750in}}%
\pgfpathlineto{\pgfqpoint{2.164228in}{1.277853in}}%
\pgfpathlineto{\pgfqpoint{2.165424in}{1.263791in}}%
\pgfpathlineto{\pgfqpoint{2.154102in}{1.264813in}}%
\pgfpathlineto{\pgfqpoint{2.130815in}{1.265402in}}%
\pgfpathlineto{\pgfqpoint{2.130827in}{1.263573in}}%
\pgfpathclose%
\pgfusepath{fill}%
\end{pgfscope}%
\begin{pgfscope}%
\pgfpathrectangle{\pgfqpoint{0.100000in}{0.100000in}}{\pgfqpoint{3.608454in}{2.310000in}}%
\pgfusepath{clip}%
\pgfsetbuttcap%
\pgfsetmiterjoin%
\definecolor{currentfill}{rgb}{0.000000,0.611765,0.694118}%
\pgfsetfillcolor{currentfill}%
\pgfsetlinewidth{0.000000pt}%
\definecolor{currentstroke}{rgb}{0.000000,0.000000,0.000000}%
\pgfsetstrokecolor{currentstroke}%
\pgfsetstrokeopacity{0.000000}%
\pgfsetdash{}{0pt}%
\pgfpathmoveto{\pgfqpoint{1.522840in}{2.169802in}}%
\pgfpathlineto{\pgfqpoint{1.578022in}{2.163356in}}%
\pgfpathlineto{\pgfqpoint{1.578243in}{2.156172in}}%
\pgfpathlineto{\pgfqpoint{1.575926in}{2.135412in}}%
\pgfpathlineto{\pgfqpoint{1.578597in}{2.128096in}}%
\pgfpathlineto{\pgfqpoint{1.534783in}{2.133134in}}%
\pgfpathlineto{\pgfqpoint{1.523357in}{2.134506in}}%
\pgfpathlineto{\pgfqpoint{1.525884in}{2.155265in}}%
\pgfpathlineto{\pgfqpoint{1.519036in}{2.156103in}}%
\pgfpathlineto{\pgfqpoint{1.522840in}{2.169802in}}%
\pgfpathclose%
\pgfusepath{fill}%
\end{pgfscope}%
\begin{pgfscope}%
\pgfpathrectangle{\pgfqpoint{0.100000in}{0.100000in}}{\pgfqpoint{3.608454in}{2.310000in}}%
\pgfusepath{clip}%
\pgfsetbuttcap%
\pgfsetmiterjoin%
\definecolor{currentfill}{rgb}{0.000000,0.419608,0.790196}%
\pgfsetfillcolor{currentfill}%
\pgfsetlinewidth{0.000000pt}%
\definecolor{currentstroke}{rgb}{0.000000,0.000000,0.000000}%
\pgfsetstrokecolor{currentstroke}%
\pgfsetstrokeopacity{0.000000}%
\pgfsetdash{}{0pt}%
\pgfpathmoveto{\pgfqpoint{3.062330in}{0.671592in}}%
\pgfpathlineto{\pgfqpoint{3.055227in}{0.666862in}}%
\pgfpathlineto{\pgfqpoint{3.041523in}{0.664546in}}%
\pgfpathlineto{\pgfqpoint{3.039796in}{0.675377in}}%
\pgfpathlineto{\pgfqpoint{3.035034in}{0.681162in}}%
\pgfpathlineto{\pgfqpoint{3.019158in}{0.678567in}}%
\pgfpathlineto{\pgfqpoint{3.011416in}{0.670241in}}%
\pgfpathlineto{\pgfqpoint{3.004229in}{0.666876in}}%
\pgfpathlineto{\pgfqpoint{2.999393in}{0.700496in}}%
\pgfpathlineto{\pgfqpoint{2.971461in}{0.696012in}}%
\pgfpathlineto{\pgfqpoint{2.966459in}{0.731456in}}%
\pgfpathlineto{\pgfqpoint{2.983300in}{0.732552in}}%
\pgfpathlineto{\pgfqpoint{2.985291in}{0.750228in}}%
\pgfpathlineto{\pgfqpoint{2.966799in}{0.748504in}}%
\pgfpathlineto{\pgfqpoint{2.964698in}{0.765826in}}%
\pgfpathlineto{\pgfqpoint{2.983960in}{0.768316in}}%
\pgfpathlineto{\pgfqpoint{2.988034in}{0.774133in}}%
\pgfpathlineto{\pgfqpoint{2.996804in}{0.773153in}}%
\pgfpathlineto{\pgfqpoint{3.005604in}{0.782838in}}%
\pgfpathlineto{\pgfqpoint{3.020788in}{0.780741in}}%
\pgfpathlineto{\pgfqpoint{3.030760in}{0.775066in}}%
\pgfpathlineto{\pgfqpoint{3.034315in}{0.767379in}}%
\pgfpathlineto{\pgfqpoint{3.032529in}{0.754216in}}%
\pgfpathlineto{\pgfqpoint{3.035355in}{0.750593in}}%
\pgfpathlineto{\pgfqpoint{3.035557in}{0.741809in}}%
\pgfpathlineto{\pgfqpoint{3.041274in}{0.727152in}}%
\pgfpathlineto{\pgfqpoint{3.042489in}{0.719793in}}%
\pgfpathlineto{\pgfqpoint{3.052515in}{0.693343in}}%
\pgfpathlineto{\pgfqpoint{3.062330in}{0.671592in}}%
\pgfpathclose%
\pgfusepath{fill}%
\end{pgfscope}%
\begin{pgfscope}%
\pgfpathrectangle{\pgfqpoint{0.100000in}{0.100000in}}{\pgfqpoint{3.608454in}{2.310000in}}%
\pgfusepath{clip}%
\pgfsetbuttcap%
\pgfsetmiterjoin%
\definecolor{currentfill}{rgb}{0.000000,0.721569,0.639216}%
\pgfsetfillcolor{currentfill}%
\pgfsetlinewidth{0.000000pt}%
\definecolor{currentstroke}{rgb}{0.000000,0.000000,0.000000}%
\pgfsetstrokecolor{currentstroke}%
\pgfsetstrokeopacity{0.000000}%
\pgfsetdash{}{0pt}%
\pgfpathmoveto{\pgfqpoint{1.590719in}{1.104616in}}%
\pgfpathlineto{\pgfqpoint{1.589893in}{1.095406in}}%
\pgfpathlineto{\pgfqpoint{1.585141in}{1.042675in}}%
\pgfpathlineto{\pgfqpoint{1.569604in}{1.044093in}}%
\pgfpathlineto{\pgfqpoint{1.568889in}{1.037039in}}%
\pgfpathlineto{\pgfqpoint{1.562067in}{1.037697in}}%
\pgfpathlineto{\pgfqpoint{1.561377in}{1.030839in}}%
\pgfpathlineto{\pgfqpoint{1.554565in}{1.031506in}}%
\pgfpathlineto{\pgfqpoint{1.553889in}{1.024630in}}%
\pgfpathlineto{\pgfqpoint{1.540290in}{1.025994in}}%
\pgfpathlineto{\pgfqpoint{1.539569in}{1.019066in}}%
\pgfpathlineto{\pgfqpoint{1.523829in}{1.020634in}}%
\pgfpathlineto{\pgfqpoint{1.512139in}{1.021816in}}%
\pgfpathlineto{\pgfqpoint{1.516667in}{1.064131in}}%
\pgfpathlineto{\pgfqpoint{1.534671in}{1.070306in}}%
\pgfpathlineto{\pgfqpoint{1.548858in}{1.068831in}}%
\pgfpathlineto{\pgfqpoint{1.550211in}{1.083114in}}%
\pgfpathlineto{\pgfqpoint{1.544598in}{1.083669in}}%
\pgfpathlineto{\pgfqpoint{1.531357in}{1.115154in}}%
\pgfpathlineto{\pgfqpoint{1.528359in}{1.113844in}}%
\pgfpathlineto{\pgfqpoint{1.506250in}{1.116084in}}%
\pgfpathlineto{\pgfqpoint{1.508110in}{1.125109in}}%
\pgfpathlineto{\pgfqpoint{1.508502in}{1.136752in}}%
\pgfpathlineto{\pgfqpoint{1.505428in}{1.150972in}}%
\pgfpathlineto{\pgfqpoint{1.532570in}{1.148149in}}%
\pgfpathlineto{\pgfqpoint{1.532225in}{1.144723in}}%
\pgfpathlineto{\pgfqpoint{1.545891in}{1.143339in}}%
\pgfpathlineto{\pgfqpoint{1.545193in}{1.136393in}}%
\pgfpathlineto{\pgfqpoint{1.572537in}{1.133851in}}%
\pgfpathlineto{\pgfqpoint{1.569376in}{1.106608in}}%
\pgfpathlineto{\pgfqpoint{1.590719in}{1.104616in}}%
\pgfpathclose%
\pgfusepath{fill}%
\end{pgfscope}%
\begin{pgfscope}%
\pgfpathrectangle{\pgfqpoint{0.100000in}{0.100000in}}{\pgfqpoint{3.608454in}{2.310000in}}%
\pgfusepath{clip}%
\pgfsetbuttcap%
\pgfsetmiterjoin%
\definecolor{currentfill}{rgb}{0.000000,0.631373,0.684314}%
\pgfsetfillcolor{currentfill}%
\pgfsetlinewidth{0.000000pt}%
\definecolor{currentstroke}{rgb}{0.000000,0.000000,0.000000}%
\pgfsetstrokecolor{currentstroke}%
\pgfsetstrokeopacity{0.000000}%
\pgfsetdash{}{0pt}%
\pgfpathmoveto{\pgfqpoint{3.130816in}{1.638979in}}%
\pgfpathlineto{\pgfqpoint{3.123343in}{1.674373in}}%
\pgfpathlineto{\pgfqpoint{3.115261in}{1.672766in}}%
\pgfpathlineto{\pgfqpoint{3.109751in}{1.709791in}}%
\pgfpathlineto{\pgfqpoint{3.111728in}{1.718970in}}%
\pgfpathlineto{\pgfqpoint{3.141312in}{1.724493in}}%
\pgfpathlineto{\pgfqpoint{3.142173in}{1.718464in}}%
\pgfpathlineto{\pgfqpoint{3.156855in}{1.718354in}}%
\pgfpathlineto{\pgfqpoint{3.155225in}{1.724423in}}%
\pgfpathlineto{\pgfqpoint{3.173622in}{1.728849in}}%
\pgfpathlineto{\pgfqpoint{3.167852in}{1.733612in}}%
\pgfpathlineto{\pgfqpoint{3.190729in}{1.738354in}}%
\pgfpathlineto{\pgfqpoint{3.194930in}{1.721757in}}%
\pgfpathlineto{\pgfqpoint{3.196916in}{1.713203in}}%
\pgfpathlineto{\pgfqpoint{3.187096in}{1.712911in}}%
\pgfpathlineto{\pgfqpoint{3.180627in}{1.708614in}}%
\pgfpathlineto{\pgfqpoint{3.183986in}{1.686549in}}%
\pgfpathlineto{\pgfqpoint{3.162699in}{1.682336in}}%
\pgfpathlineto{\pgfqpoint{3.172073in}{1.651526in}}%
\pgfpathlineto{\pgfqpoint{3.167806in}{1.647018in}}%
\pgfpathlineto{\pgfqpoint{3.130816in}{1.638979in}}%
\pgfpathclose%
\pgfusepath{fill}%
\end{pgfscope}%
\begin{pgfscope}%
\pgfpathrectangle{\pgfqpoint{0.100000in}{0.100000in}}{\pgfqpoint{3.608454in}{2.310000in}}%
\pgfusepath{clip}%
\pgfsetbuttcap%
\pgfsetmiterjoin%
\definecolor{currentfill}{rgb}{0.000000,0.796078,0.601961}%
\pgfsetfillcolor{currentfill}%
\pgfsetlinewidth{0.000000pt}%
\definecolor{currentstroke}{rgb}{0.000000,0.000000,0.000000}%
\pgfsetstrokecolor{currentstroke}%
\pgfsetstrokeopacity{0.000000}%
\pgfsetdash{}{0pt}%
\pgfpathmoveto{\pgfqpoint{0.800486in}{1.923543in}}%
\pgfpathlineto{\pgfqpoint{0.805150in}{1.923524in}}%
\pgfpathlineto{\pgfqpoint{0.814693in}{1.932896in}}%
\pgfpathlineto{\pgfqpoint{0.818538in}{1.934210in}}%
\pgfpathlineto{\pgfqpoint{0.839459in}{1.928946in}}%
\pgfpathlineto{\pgfqpoint{0.844130in}{1.923115in}}%
\pgfpathlineto{\pgfqpoint{0.842438in}{1.916478in}}%
\pgfpathlineto{\pgfqpoint{0.857097in}{1.912795in}}%
\pgfpathlineto{\pgfqpoint{0.860164in}{1.926417in}}%
\pgfpathlineto{\pgfqpoint{0.866297in}{1.932367in}}%
\pgfpathlineto{\pgfqpoint{0.875640in}{1.944745in}}%
\pgfpathlineto{\pgfqpoint{0.884981in}{1.950005in}}%
\pgfpathlineto{\pgfqpoint{0.899991in}{1.946190in}}%
\pgfpathlineto{\pgfqpoint{0.896730in}{1.932620in}}%
\pgfpathlineto{\pgfqpoint{0.901098in}{1.927625in}}%
\pgfpathlineto{\pgfqpoint{0.898851in}{1.918021in}}%
\pgfpathlineto{\pgfqpoint{0.904574in}{1.912817in}}%
\pgfpathlineto{\pgfqpoint{0.910516in}{1.911417in}}%
\pgfpathlineto{\pgfqpoint{0.905607in}{1.902658in}}%
\pgfpathlineto{\pgfqpoint{0.901859in}{1.889570in}}%
\pgfpathlineto{\pgfqpoint{0.896690in}{1.890839in}}%
\pgfpathlineto{\pgfqpoint{0.895063in}{1.884202in}}%
\pgfpathlineto{\pgfqpoint{0.890663in}{1.885325in}}%
\pgfpathlineto{\pgfqpoint{0.886843in}{1.879251in}}%
\pgfpathlineto{\pgfqpoint{0.879074in}{1.881212in}}%
\pgfpathlineto{\pgfqpoint{0.875796in}{1.867766in}}%
\pgfpathlineto{\pgfqpoint{0.867712in}{1.868584in}}%
\pgfpathlineto{\pgfqpoint{0.863735in}{1.876688in}}%
\pgfpathlineto{\pgfqpoint{0.860241in}{1.876128in}}%
\pgfpathlineto{\pgfqpoint{0.855951in}{1.862354in}}%
\pgfpathlineto{\pgfqpoint{0.824016in}{1.733109in}}%
\pgfpathlineto{\pgfqpoint{0.757203in}{1.749818in}}%
\pgfpathlineto{\pgfqpoint{0.761786in}{1.771448in}}%
\pgfpathlineto{\pgfqpoint{0.774453in}{1.820402in}}%
\pgfpathlineto{\pgfqpoint{0.773761in}{1.820585in}}%
\pgfpathlineto{\pgfqpoint{0.790924in}{1.886917in}}%
\pgfpathlineto{\pgfqpoint{0.796408in}{1.906884in}}%
\pgfpathlineto{\pgfqpoint{0.800486in}{1.923543in}}%
\pgfpathclose%
\pgfusepath{fill}%
\end{pgfscope}%
\begin{pgfscope}%
\pgfpathrectangle{\pgfqpoint{0.100000in}{0.100000in}}{\pgfqpoint{3.608454in}{2.310000in}}%
\pgfusepath{clip}%
\pgfsetbuttcap%
\pgfsetmiterjoin%
\definecolor{currentfill}{rgb}{0.000000,0.635294,0.682353}%
\pgfsetfillcolor{currentfill}%
\pgfsetlinewidth{0.000000pt}%
\definecolor{currentstroke}{rgb}{0.000000,0.000000,0.000000}%
\pgfsetstrokecolor{currentstroke}%
\pgfsetstrokeopacity{0.000000}%
\pgfsetdash{}{0pt}%
\pgfpathmoveto{\pgfqpoint{2.862925in}{0.815951in}}%
\pgfpathlineto{\pgfqpoint{2.854299in}{0.818491in}}%
\pgfpathlineto{\pgfqpoint{2.853827in}{0.823500in}}%
\pgfpathlineto{\pgfqpoint{2.847520in}{0.837484in}}%
\pgfpathlineto{\pgfqpoint{2.844094in}{0.841666in}}%
\pgfpathlineto{\pgfqpoint{2.847916in}{0.849972in}}%
\pgfpathlineto{\pgfqpoint{2.855107in}{0.854082in}}%
\pgfpathlineto{\pgfqpoint{2.871284in}{0.855909in}}%
\pgfpathlineto{\pgfqpoint{2.872842in}{0.842457in}}%
\pgfpathlineto{\pgfqpoint{2.891100in}{0.843787in}}%
\pgfpathlineto{\pgfqpoint{2.894217in}{0.842012in}}%
\pgfpathlineto{\pgfqpoint{2.895939in}{0.833978in}}%
\pgfpathlineto{\pgfqpoint{2.900223in}{0.831517in}}%
\pgfpathlineto{\pgfqpoint{2.904439in}{0.824809in}}%
\pgfpathlineto{\pgfqpoint{2.884148in}{0.822135in}}%
\pgfpathlineto{\pgfqpoint{2.875215in}{0.821548in}}%
\pgfpathlineto{\pgfqpoint{2.875686in}{0.817596in}}%
\pgfpathlineto{\pgfqpoint{2.862925in}{0.815951in}}%
\pgfpathclose%
\pgfusepath{fill}%
\end{pgfscope}%
\begin{pgfscope}%
\pgfpathrectangle{\pgfqpoint{0.100000in}{0.100000in}}{\pgfqpoint{3.608454in}{2.310000in}}%
\pgfusepath{clip}%
\pgfsetbuttcap%
\pgfsetmiterjoin%
\definecolor{currentfill}{rgb}{0.000000,0.517647,0.741176}%
\pgfsetfillcolor{currentfill}%
\pgfsetlinewidth{0.000000pt}%
\definecolor{currentstroke}{rgb}{0.000000,0.000000,0.000000}%
\pgfsetstrokecolor{currentstroke}%
\pgfsetstrokeopacity{0.000000}%
\pgfsetdash{}{0pt}%
\pgfpathmoveto{\pgfqpoint{1.793434in}{1.427744in}}%
\pgfpathlineto{\pgfqpoint{1.792481in}{1.427797in}}%
\pgfpathlineto{\pgfqpoint{1.793800in}{1.455314in}}%
\pgfpathlineto{\pgfqpoint{1.758946in}{1.457374in}}%
\pgfpathlineto{\pgfqpoint{1.760676in}{1.485050in}}%
\pgfpathlineto{\pgfqpoint{1.793857in}{1.483052in}}%
\pgfpathlineto{\pgfqpoint{1.795374in}{1.510393in}}%
\pgfpathlineto{\pgfqpoint{1.842917in}{1.507856in}}%
\pgfpathlineto{\pgfqpoint{1.841990in}{1.478161in}}%
\pgfpathlineto{\pgfqpoint{1.856208in}{1.476676in}}%
\pgfpathlineto{\pgfqpoint{1.853837in}{1.424663in}}%
\pgfpathlineto{\pgfqpoint{1.826801in}{1.425965in}}%
\pgfpathlineto{\pgfqpoint{1.793434in}{1.427744in}}%
\pgfpathclose%
\pgfusepath{fill}%
\end{pgfscope}%
\begin{pgfscope}%
\pgfpathrectangle{\pgfqpoint{0.100000in}{0.100000in}}{\pgfqpoint{3.608454in}{2.310000in}}%
\pgfusepath{clip}%
\pgfsetbuttcap%
\pgfsetmiterjoin%
\definecolor{currentfill}{rgb}{0.000000,0.341176,0.829412}%
\pgfsetfillcolor{currentfill}%
\pgfsetlinewidth{0.000000pt}%
\definecolor{currentstroke}{rgb}{0.000000,0.000000,0.000000}%
\pgfsetstrokecolor{currentstroke}%
\pgfsetstrokeopacity{0.000000}%
\pgfsetdash{}{0pt}%
\pgfpathmoveto{\pgfqpoint{1.538610in}{1.726627in}}%
\pgfpathlineto{\pgfqpoint{1.586973in}{1.721839in}}%
\pgfpathlineto{\pgfqpoint{1.583107in}{1.682203in}}%
\pgfpathlineto{\pgfqpoint{1.580064in}{1.651542in}}%
\pgfpathlineto{\pgfqpoint{1.531372in}{1.656425in}}%
\pgfpathlineto{\pgfqpoint{1.532653in}{1.671867in}}%
\pgfpathlineto{\pgfqpoint{1.538610in}{1.726627in}}%
\pgfpathclose%
\pgfusepath{fill}%
\end{pgfscope}%
\begin{pgfscope}%
\pgfpathrectangle{\pgfqpoint{0.100000in}{0.100000in}}{\pgfqpoint{3.608454in}{2.310000in}}%
\pgfusepath{clip}%
\pgfsetbuttcap%
\pgfsetmiterjoin%
\definecolor{currentfill}{rgb}{0.000000,0.392157,0.803922}%
\pgfsetfillcolor{currentfill}%
\pgfsetlinewidth{0.000000pt}%
\definecolor{currentstroke}{rgb}{0.000000,0.000000,0.000000}%
\pgfsetstrokecolor{currentstroke}%
\pgfsetstrokeopacity{0.000000}%
\pgfsetdash{}{0pt}%
\pgfpathmoveto{\pgfqpoint{0.429033in}{1.527264in}}%
\pgfpathlineto{\pgfqpoint{0.432587in}{1.529018in}}%
\pgfpathlineto{\pgfqpoint{0.433169in}{1.538861in}}%
\pgfpathlineto{\pgfqpoint{0.431374in}{1.550880in}}%
\pgfpathlineto{\pgfqpoint{0.425226in}{1.565689in}}%
\pgfpathlineto{\pgfqpoint{0.427977in}{1.570534in}}%
\pgfpathlineto{\pgfqpoint{0.438495in}{1.570367in}}%
\pgfpathlineto{\pgfqpoint{0.446440in}{1.577149in}}%
\pgfpathlineto{\pgfqpoint{0.446281in}{1.581110in}}%
\pgfpathlineto{\pgfqpoint{0.452005in}{1.577321in}}%
\pgfpathlineto{\pgfqpoint{0.452551in}{1.565246in}}%
\pgfpathlineto{\pgfqpoint{0.454748in}{1.561199in}}%
\pgfpathlineto{\pgfqpoint{0.455284in}{1.549289in}}%
\pgfpathlineto{\pgfqpoint{0.460110in}{1.545792in}}%
\pgfpathlineto{\pgfqpoint{0.465266in}{1.547913in}}%
\pgfpathlineto{\pgfqpoint{0.479765in}{1.542960in}}%
\pgfpathlineto{\pgfqpoint{0.474910in}{1.526813in}}%
\pgfpathlineto{\pgfqpoint{0.480918in}{1.524945in}}%
\pgfpathlineto{\pgfqpoint{0.477138in}{1.516663in}}%
\pgfpathlineto{\pgfqpoint{0.471900in}{1.514999in}}%
\pgfpathlineto{\pgfqpoint{0.471851in}{1.509726in}}%
\pgfpathlineto{\pgfqpoint{0.476612in}{1.508153in}}%
\pgfpathlineto{\pgfqpoint{0.471692in}{1.491720in}}%
\pgfpathlineto{\pgfqpoint{0.471844in}{1.486991in}}%
\pgfpathlineto{\pgfqpoint{0.465525in}{1.465978in}}%
\pgfpathlineto{\pgfqpoint{0.469229in}{1.459834in}}%
\pgfpathlineto{\pgfqpoint{0.469193in}{1.459876in}}%
\pgfpathlineto{\pgfqpoint{0.445670in}{1.467188in}}%
\pgfpathlineto{\pgfqpoint{0.441374in}{1.466010in}}%
\pgfpathlineto{\pgfqpoint{0.432540in}{1.470642in}}%
\pgfpathlineto{\pgfqpoint{0.424905in}{1.469119in}}%
\pgfpathlineto{\pgfqpoint{0.422616in}{1.464232in}}%
\pgfpathlineto{\pgfqpoint{0.422155in}{1.452127in}}%
\pgfpathlineto{\pgfqpoint{0.411566in}{1.453040in}}%
\pgfpathlineto{\pgfqpoint{0.411162in}{1.446533in}}%
\pgfpathlineto{\pgfqpoint{0.406795in}{1.454089in}}%
\pgfpathlineto{\pgfqpoint{0.406757in}{1.459003in}}%
\pgfpathlineto{\pgfqpoint{0.410394in}{1.466958in}}%
\pgfpathlineto{\pgfqpoint{0.410387in}{1.477682in}}%
\pgfpathlineto{\pgfqpoint{0.407424in}{1.482308in}}%
\pgfpathlineto{\pgfqpoint{0.411743in}{1.490790in}}%
\pgfpathlineto{\pgfqpoint{0.413684in}{1.501555in}}%
\pgfpathlineto{\pgfqpoint{0.420869in}{1.501805in}}%
\pgfpathlineto{\pgfqpoint{0.418472in}{1.484385in}}%
\pgfpathlineto{\pgfqpoint{0.423806in}{1.481158in}}%
\pgfpathlineto{\pgfqpoint{0.428604in}{1.473723in}}%
\pgfpathlineto{\pgfqpoint{0.431932in}{1.475097in}}%
\pgfpathlineto{\pgfqpoint{0.432089in}{1.485171in}}%
\pgfpathlineto{\pgfqpoint{0.424993in}{1.499408in}}%
\pgfpathlineto{\pgfqpoint{0.428477in}{1.506604in}}%
\pgfpathlineto{\pgfqpoint{0.424773in}{1.512850in}}%
\pgfpathlineto{\pgfqpoint{0.435086in}{1.517401in}}%
\pgfpathlineto{\pgfqpoint{0.433615in}{1.522888in}}%
\pgfpathlineto{\pgfqpoint{0.429033in}{1.527264in}}%
\pgfpathclose%
\pgfusepath{fill}%
\end{pgfscope}%
\begin{pgfscope}%
\pgfpathrectangle{\pgfqpoint{0.100000in}{0.100000in}}{\pgfqpoint{3.608454in}{2.310000in}}%
\pgfusepath{clip}%
\pgfsetbuttcap%
\pgfsetmiterjoin%
\definecolor{currentfill}{rgb}{0.000000,0.392157,0.803922}%
\pgfsetfillcolor{currentfill}%
\pgfsetlinewidth{0.000000pt}%
\definecolor{currentstroke}{rgb}{0.000000,0.000000,0.000000}%
\pgfsetstrokecolor{currentstroke}%
\pgfsetstrokeopacity{0.000000}%
\pgfsetdash{}{0pt}%
\pgfpathmoveto{\pgfqpoint{0.422980in}{1.525827in}}%
\pgfpathlineto{\pgfqpoint{0.419309in}{1.509781in}}%
\pgfpathlineto{\pgfqpoint{0.414223in}{1.505038in}}%
\pgfpathlineto{\pgfqpoint{0.408463in}{1.513232in}}%
\pgfpathlineto{\pgfqpoint{0.405158in}{1.513449in}}%
\pgfpathlineto{\pgfqpoint{0.399143in}{1.525491in}}%
\pgfpathlineto{\pgfqpoint{0.393060in}{1.528517in}}%
\pgfpathlineto{\pgfqpoint{0.396639in}{1.537475in}}%
\pgfpathlineto{\pgfqpoint{0.396978in}{1.549508in}}%
\pgfpathlineto{\pgfqpoint{0.403577in}{1.549156in}}%
\pgfpathlineto{\pgfqpoint{0.410415in}{1.537899in}}%
\pgfpathlineto{\pgfqpoint{0.419978in}{1.533247in}}%
\pgfpathlineto{\pgfqpoint{0.422980in}{1.525827in}}%
\pgfpathclose%
\pgfusepath{fill}%
\end{pgfscope}%
\begin{pgfscope}%
\pgfpathrectangle{\pgfqpoint{0.100000in}{0.100000in}}{\pgfqpoint{3.608454in}{2.310000in}}%
\pgfusepath{clip}%
\pgfsetbuttcap%
\pgfsetmiterjoin%
\definecolor{currentfill}{rgb}{0.000000,0.592157,0.703922}%
\pgfsetfillcolor{currentfill}%
\pgfsetlinewidth{0.000000pt}%
\definecolor{currentstroke}{rgb}{0.000000,0.000000,0.000000}%
\pgfsetstrokecolor{currentstroke}%
\pgfsetstrokeopacity{0.000000}%
\pgfsetdash{}{0pt}%
\pgfpathmoveto{\pgfqpoint{2.410650in}{1.552521in}}%
\pgfpathlineto{\pgfqpoint{2.391551in}{1.551544in}}%
\pgfpathlineto{\pgfqpoint{2.387375in}{1.567036in}}%
\pgfpathlineto{\pgfqpoint{2.390834in}{1.569129in}}%
\pgfpathlineto{\pgfqpoint{2.392472in}{1.584764in}}%
\pgfpathlineto{\pgfqpoint{2.390701in}{1.593517in}}%
\pgfpathlineto{\pgfqpoint{2.381293in}{1.599262in}}%
\pgfpathlineto{\pgfqpoint{2.404561in}{1.600729in}}%
\pgfpathlineto{\pgfqpoint{2.402794in}{1.625171in}}%
\pgfpathlineto{\pgfqpoint{2.435480in}{1.626662in}}%
\pgfpathlineto{\pgfqpoint{2.473885in}{1.628654in}}%
\pgfpathlineto{\pgfqpoint{2.475849in}{1.601806in}}%
\pgfpathlineto{\pgfqpoint{2.482705in}{1.602290in}}%
\pgfpathlineto{\pgfqpoint{2.482413in}{1.595355in}}%
\pgfpathlineto{\pgfqpoint{2.484797in}{1.561007in}}%
\pgfpathlineto{\pgfqpoint{2.451488in}{1.558550in}}%
\pgfpathlineto{\pgfqpoint{2.451750in}{1.555104in}}%
\pgfpathlineto{\pgfqpoint{2.410650in}{1.552521in}}%
\pgfpathclose%
\pgfusepath{fill}%
\end{pgfscope}%
\begin{pgfscope}%
\pgfpathrectangle{\pgfqpoint{0.100000in}{0.100000in}}{\pgfqpoint{3.608454in}{2.310000in}}%
\pgfusepath{clip}%
\pgfsetbuttcap%
\pgfsetmiterjoin%
\definecolor{currentfill}{rgb}{0.000000,0.462745,0.768627}%
\pgfsetfillcolor{currentfill}%
\pgfsetlinewidth{0.000000pt}%
\definecolor{currentstroke}{rgb}{0.000000,0.000000,0.000000}%
\pgfsetstrokecolor{currentstroke}%
\pgfsetstrokeopacity{0.000000}%
\pgfsetdash{}{0pt}%
\pgfpathmoveto{\pgfqpoint{2.583434in}{1.352723in}}%
\pgfpathlineto{\pgfqpoint{2.590232in}{1.353266in}}%
\pgfpathlineto{\pgfqpoint{2.590013in}{1.348185in}}%
\pgfpathlineto{\pgfqpoint{2.584141in}{1.345785in}}%
\pgfpathlineto{\pgfqpoint{2.582510in}{1.335565in}}%
\pgfpathlineto{\pgfqpoint{2.585800in}{1.324135in}}%
\pgfpathlineto{\pgfqpoint{2.572272in}{1.322088in}}%
\pgfpathlineto{\pgfqpoint{2.565371in}{1.317560in}}%
\pgfpathlineto{\pgfqpoint{2.560427in}{1.323864in}}%
\pgfpathlineto{\pgfqpoint{2.544363in}{1.322774in}}%
\pgfpathlineto{\pgfqpoint{2.542895in}{1.344839in}}%
\pgfpathlineto{\pgfqpoint{2.565996in}{1.346868in}}%
\pgfpathlineto{\pgfqpoint{2.566176in}{1.350839in}}%
\pgfpathlineto{\pgfqpoint{2.583434in}{1.352723in}}%
\pgfpathclose%
\pgfusepath{fill}%
\end{pgfscope}%
\begin{pgfscope}%
\pgfpathrectangle{\pgfqpoint{0.100000in}{0.100000in}}{\pgfqpoint{3.608454in}{2.310000in}}%
\pgfusepath{clip}%
\pgfsetbuttcap%
\pgfsetmiterjoin%
\definecolor{currentfill}{rgb}{0.000000,0.701961,0.649020}%
\pgfsetfillcolor{currentfill}%
\pgfsetlinewidth{0.000000pt}%
\definecolor{currentstroke}{rgb}{0.000000,0.000000,0.000000}%
\pgfsetstrokecolor{currentstroke}%
\pgfsetstrokeopacity{0.000000}%
\pgfsetdash{}{0pt}%
\pgfpathmoveto{\pgfqpoint{2.062941in}{0.991185in}}%
\pgfpathlineto{\pgfqpoint{2.062908in}{0.984282in}}%
\pgfpathlineto{\pgfqpoint{2.055870in}{0.984320in}}%
\pgfpathlineto{\pgfqpoint{2.055634in}{0.956575in}}%
\pgfpathlineto{\pgfqpoint{2.041706in}{0.956720in}}%
\pgfpathlineto{\pgfqpoint{2.014398in}{0.957120in}}%
\pgfpathlineto{\pgfqpoint{2.014740in}{0.977835in}}%
\pgfpathlineto{\pgfqpoint{2.007778in}{0.978034in}}%
\pgfpathlineto{\pgfqpoint{2.007887in}{0.984833in}}%
\pgfpathlineto{\pgfqpoint{1.987413in}{0.985257in}}%
\pgfpathlineto{\pgfqpoint{1.987552in}{0.992168in}}%
\pgfpathlineto{\pgfqpoint{1.980692in}{0.992304in}}%
\pgfpathlineto{\pgfqpoint{1.981120in}{1.013018in}}%
\pgfpathlineto{\pgfqpoint{1.984031in}{1.018827in}}%
\pgfpathlineto{\pgfqpoint{1.988856in}{1.020093in}}%
\pgfpathlineto{\pgfqpoint{1.994760in}{1.012926in}}%
\pgfpathlineto{\pgfqpoint{1.999760in}{1.016432in}}%
\pgfpathlineto{\pgfqpoint{2.009922in}{1.016922in}}%
\pgfpathlineto{\pgfqpoint{2.010212in}{1.033190in}}%
\pgfpathlineto{\pgfqpoint{2.013392in}{1.033129in}}%
\pgfpathlineto{\pgfqpoint{2.013608in}{1.046935in}}%
\pgfpathlineto{\pgfqpoint{2.043347in}{1.046522in}}%
\pgfpathlineto{\pgfqpoint{2.043077in}{1.035565in}}%
\pgfpathlineto{\pgfqpoint{2.036171in}{1.027587in}}%
\pgfpathlineto{\pgfqpoint{2.035603in}{0.998326in}}%
\pgfpathlineto{\pgfqpoint{2.049296in}{0.998169in}}%
\pgfpathlineto{\pgfqpoint{2.049224in}{0.991276in}}%
\pgfpathlineto{\pgfqpoint{2.062941in}{0.991185in}}%
\pgfpathclose%
\pgfusepath{fill}%
\end{pgfscope}%
\begin{pgfscope}%
\pgfpathrectangle{\pgfqpoint{0.100000in}{0.100000in}}{\pgfqpoint{3.608454in}{2.310000in}}%
\pgfusepath{clip}%
\pgfsetbuttcap%
\pgfsetmiterjoin%
\definecolor{currentfill}{rgb}{0.000000,0.580392,0.709804}%
\pgfsetfillcolor{currentfill}%
\pgfsetlinewidth{0.000000pt}%
\definecolor{currentstroke}{rgb}{0.000000,0.000000,0.000000}%
\pgfsetstrokecolor{currentstroke}%
\pgfsetstrokeopacity{0.000000}%
\pgfsetdash{}{0pt}%
\pgfpathmoveto{\pgfqpoint{3.152496in}{1.240772in}}%
\pgfpathlineto{\pgfqpoint{3.187275in}{1.247460in}}%
\pgfpathlineto{\pgfqpoint{3.188912in}{1.234693in}}%
\pgfpathlineto{\pgfqpoint{3.185638in}{1.219606in}}%
\pgfpathlineto{\pgfqpoint{3.177079in}{1.221508in}}%
\pgfpathlineto{\pgfqpoint{3.165876in}{1.221007in}}%
\pgfpathlineto{\pgfqpoint{3.160658in}{1.212280in}}%
\pgfpathlineto{\pgfqpoint{3.154402in}{1.212187in}}%
\pgfpathlineto{\pgfqpoint{3.150361in}{1.233286in}}%
\pgfpathlineto{\pgfqpoint{3.152496in}{1.240772in}}%
\pgfpathclose%
\pgfusepath{fill}%
\end{pgfscope}%
\begin{pgfscope}%
\pgfpathrectangle{\pgfqpoint{0.100000in}{0.100000in}}{\pgfqpoint{3.608454in}{2.310000in}}%
\pgfusepath{clip}%
\pgfsetbuttcap%
\pgfsetmiterjoin%
\definecolor{currentfill}{rgb}{0.000000,0.792157,0.603922}%
\pgfsetfillcolor{currentfill}%
\pgfsetlinewidth{0.000000pt}%
\definecolor{currentstroke}{rgb}{0.000000,0.000000,0.000000}%
\pgfsetstrokecolor{currentstroke}%
\pgfsetstrokeopacity{0.000000}%
\pgfsetdash{}{0pt}%
\pgfpathmoveto{\pgfqpoint{1.183780in}{0.997802in}}%
\pgfpathlineto{\pgfqpoint{1.176800in}{0.953138in}}%
\pgfpathlineto{\pgfqpoint{1.167197in}{0.891666in}}%
\pgfpathlineto{\pgfqpoint{1.092774in}{0.903733in}}%
\pgfpathlineto{\pgfqpoint{1.074172in}{0.906976in}}%
\pgfpathlineto{\pgfqpoint{1.075355in}{0.913752in}}%
\pgfpathlineto{\pgfqpoint{1.084657in}{0.967062in}}%
\pgfpathlineto{\pgfqpoint{1.087666in}{0.981611in}}%
\pgfpathlineto{\pgfqpoint{1.096233in}{0.984306in}}%
\pgfpathlineto{\pgfqpoint{1.106555in}{0.985188in}}%
\pgfpathlineto{\pgfqpoint{1.117590in}{0.983356in}}%
\pgfpathlineto{\pgfqpoint{1.119039in}{0.992055in}}%
\pgfpathlineto{\pgfqpoint{1.125550in}{0.990854in}}%
\pgfpathlineto{\pgfqpoint{1.133001in}{0.981587in}}%
\pgfpathlineto{\pgfqpoint{1.145330in}{0.990568in}}%
\pgfpathlineto{\pgfqpoint{1.150108in}{0.990389in}}%
\pgfpathlineto{\pgfqpoint{1.154422in}{0.995185in}}%
\pgfpathlineto{\pgfqpoint{1.164171in}{1.000833in}}%
\pgfpathlineto{\pgfqpoint{1.183780in}{0.997802in}}%
\pgfpathclose%
\pgfusepath{fill}%
\end{pgfscope}%
\begin{pgfscope}%
\pgfpathrectangle{\pgfqpoint{0.100000in}{0.100000in}}{\pgfqpoint{3.608454in}{2.310000in}}%
\pgfusepath{clip}%
\pgfsetbuttcap%
\pgfsetmiterjoin%
\definecolor{currentfill}{rgb}{0.000000,0.760784,0.619608}%
\pgfsetfillcolor{currentfill}%
\pgfsetlinewidth{0.000000pt}%
\definecolor{currentstroke}{rgb}{0.000000,0.000000,0.000000}%
\pgfsetstrokecolor{currentstroke}%
\pgfsetstrokeopacity{0.000000}%
\pgfsetdash{}{0pt}%
\pgfpathmoveto{\pgfqpoint{2.056507in}{0.745639in}}%
\pgfpathlineto{\pgfqpoint{2.061056in}{0.729018in}}%
\pgfpathlineto{\pgfqpoint{2.057335in}{0.727032in}}%
\pgfpathlineto{\pgfqpoint{2.053852in}{0.716384in}}%
\pgfpathlineto{\pgfqpoint{2.054567in}{0.712981in}}%
\pgfpathlineto{\pgfqpoint{2.040106in}{0.712927in}}%
\pgfpathlineto{\pgfqpoint{2.021875in}{0.703710in}}%
\pgfpathlineto{\pgfqpoint{2.016298in}{0.728776in}}%
\pgfpathlineto{\pgfqpoint{2.016972in}{0.734385in}}%
\pgfpathlineto{\pgfqpoint{1.997875in}{0.723825in}}%
\pgfpathlineto{\pgfqpoint{1.975832in}{0.763164in}}%
\pgfpathlineto{\pgfqpoint{1.990434in}{0.771318in}}%
\pgfpathlineto{\pgfqpoint{1.978879in}{0.792186in}}%
\pgfpathlineto{\pgfqpoint{2.013777in}{0.811827in}}%
\pgfpathlineto{\pgfqpoint{2.023652in}{0.805506in}}%
\pgfpathlineto{\pgfqpoint{2.034850in}{0.788420in}}%
\pgfpathlineto{\pgfqpoint{2.042310in}{0.767446in}}%
\pgfpathlineto{\pgfqpoint{2.047446in}{0.765794in}}%
\pgfpathlineto{\pgfqpoint{2.048151in}{0.760702in}}%
\pgfpathlineto{\pgfqpoint{2.058523in}{0.753655in}}%
\pgfpathlineto{\pgfqpoint{2.056507in}{0.745639in}}%
\pgfpathclose%
\pgfusepath{fill}%
\end{pgfscope}%
\begin{pgfscope}%
\pgfpathrectangle{\pgfqpoint{0.100000in}{0.100000in}}{\pgfqpoint{3.608454in}{2.310000in}}%
\pgfusepath{clip}%
\pgfsetbuttcap%
\pgfsetmiterjoin%
\definecolor{currentfill}{rgb}{0.000000,0.466667,0.766667}%
\pgfsetfillcolor{currentfill}%
\pgfsetlinewidth{0.000000pt}%
\definecolor{currentstroke}{rgb}{0.000000,0.000000,0.000000}%
\pgfsetstrokecolor{currentstroke}%
\pgfsetstrokeopacity{0.000000}%
\pgfsetdash{}{0pt}%
\pgfpathmoveto{\pgfqpoint{1.951799in}{2.072044in}}%
\pgfpathlineto{\pgfqpoint{1.952493in}{2.057946in}}%
\pgfpathlineto{\pgfqpoint{1.951584in}{2.030204in}}%
\pgfpathlineto{\pgfqpoint{1.946085in}{2.030374in}}%
\pgfpathlineto{\pgfqpoint{1.918355in}{2.031373in}}%
\pgfpathlineto{\pgfqpoint{1.917529in}{2.045335in}}%
\pgfpathlineto{\pgfqpoint{1.876114in}{2.047161in}}%
\pgfpathlineto{\pgfqpoint{1.848523in}{2.048601in}}%
\pgfpathlineto{\pgfqpoint{1.847326in}{2.062601in}}%
\pgfpathlineto{\pgfqpoint{1.848899in}{2.090580in}}%
\pgfpathlineto{\pgfqpoint{1.867703in}{2.089585in}}%
\pgfpathlineto{\pgfqpoint{1.895459in}{2.088252in}}%
\pgfpathlineto{\pgfqpoint{1.896095in}{2.102164in}}%
\pgfpathlineto{\pgfqpoint{1.930695in}{2.100743in}}%
\pgfpathlineto{\pgfqpoint{1.931568in}{2.086739in}}%
\pgfpathlineto{\pgfqpoint{1.931043in}{2.072760in}}%
\pgfpathlineto{\pgfqpoint{1.951799in}{2.072044in}}%
\pgfpathclose%
\pgfusepath{fill}%
\end{pgfscope}%
\begin{pgfscope}%
\pgfpathrectangle{\pgfqpoint{0.100000in}{0.100000in}}{\pgfqpoint{3.608454in}{2.310000in}}%
\pgfusepath{clip}%
\pgfsetbuttcap%
\pgfsetmiterjoin%
\definecolor{currentfill}{rgb}{0.000000,0.654902,0.672549}%
\pgfsetfillcolor{currentfill}%
\pgfsetlinewidth{0.000000pt}%
\definecolor{currentstroke}{rgb}{0.000000,0.000000,0.000000}%
\pgfsetstrokecolor{currentstroke}%
\pgfsetstrokeopacity{0.000000}%
\pgfsetdash{}{0pt}%
\pgfpathmoveto{\pgfqpoint{2.452854in}{0.589739in}}%
\pgfpathlineto{\pgfqpoint{2.452510in}{0.580153in}}%
\pgfpathlineto{\pgfqpoint{2.447694in}{0.578303in}}%
\pgfpathlineto{\pgfqpoint{2.452088in}{0.571325in}}%
\pgfpathlineto{\pgfqpoint{2.450488in}{0.567507in}}%
\pgfpathlineto{\pgfqpoint{2.441407in}{0.560816in}}%
\pgfpathlineto{\pgfqpoint{2.437817in}{0.565645in}}%
\pgfpathlineto{\pgfqpoint{2.435406in}{0.574170in}}%
\pgfpathlineto{\pgfqpoint{2.431293in}{0.576903in}}%
\pgfpathlineto{\pgfqpoint{2.428352in}{0.570405in}}%
\pgfpathlineto{\pgfqpoint{2.421368in}{0.576035in}}%
\pgfpathlineto{\pgfqpoint{2.411945in}{0.568422in}}%
\pgfpathlineto{\pgfqpoint{2.411612in}{0.562816in}}%
\pgfpathlineto{\pgfqpoint{2.405732in}{0.562002in}}%
\pgfpathlineto{\pgfqpoint{2.397767in}{0.557468in}}%
\pgfpathlineto{\pgfqpoint{2.389463in}{0.565906in}}%
\pgfpathlineto{\pgfqpoint{2.380049in}{0.566008in}}%
\pgfpathlineto{\pgfqpoint{2.367175in}{0.570371in}}%
\pgfpathlineto{\pgfqpoint{2.363150in}{0.573969in}}%
\pgfpathlineto{\pgfqpoint{2.369537in}{0.580614in}}%
\pgfpathlineto{\pgfqpoint{2.366976in}{0.588543in}}%
\pgfpathlineto{\pgfqpoint{2.347924in}{0.595251in}}%
\pgfpathlineto{\pgfqpoint{2.341614in}{0.599740in}}%
\pgfpathlineto{\pgfqpoint{2.341637in}{0.608452in}}%
\pgfpathlineto{\pgfqpoint{2.330757in}{0.608626in}}%
\pgfpathlineto{\pgfqpoint{2.323648in}{0.606472in}}%
\pgfpathlineto{\pgfqpoint{2.335052in}{0.617153in}}%
\pgfpathlineto{\pgfqpoint{2.340753in}{0.625307in}}%
\pgfpathlineto{\pgfqpoint{2.351303in}{0.626566in}}%
\pgfpathlineto{\pgfqpoint{2.358501in}{0.617112in}}%
\pgfpathlineto{\pgfqpoint{2.357771in}{0.613019in}}%
\pgfpathlineto{\pgfqpoint{2.372197in}{0.609200in}}%
\pgfpathlineto{\pgfqpoint{2.378486in}{0.611623in}}%
\pgfpathlineto{\pgfqpoint{2.376744in}{0.617237in}}%
\pgfpathlineto{\pgfqpoint{2.371581in}{0.618635in}}%
\pgfpathlineto{\pgfqpoint{2.365882in}{0.630690in}}%
\pgfpathlineto{\pgfqpoint{2.368222in}{0.632799in}}%
\pgfpathlineto{\pgfqpoint{2.383161in}{0.637744in}}%
\pgfpathlineto{\pgfqpoint{2.386115in}{0.636858in}}%
\pgfpathlineto{\pgfqpoint{2.386695in}{0.629800in}}%
\pgfpathlineto{\pgfqpoint{2.392099in}{0.624223in}}%
\pgfpathlineto{\pgfqpoint{2.416070in}{0.624294in}}%
\pgfpathlineto{\pgfqpoint{2.420890in}{0.618181in}}%
\pgfpathlineto{\pgfqpoint{2.428180in}{0.614445in}}%
\pgfpathlineto{\pgfqpoint{2.430825in}{0.609081in}}%
\pgfpathlineto{\pgfqpoint{2.440322in}{0.609463in}}%
\pgfpathlineto{\pgfqpoint{2.443757in}{0.607123in}}%
\pgfpathlineto{\pgfqpoint{2.442006in}{0.599263in}}%
\pgfpathlineto{\pgfqpoint{2.445754in}{0.592557in}}%
\pgfpathlineto{\pgfqpoint{2.452854in}{0.589739in}}%
\pgfpathclose%
\pgfusepath{fill}%
\end{pgfscope}%
\begin{pgfscope}%
\pgfpathrectangle{\pgfqpoint{0.100000in}{0.100000in}}{\pgfqpoint{3.608454in}{2.310000in}}%
\pgfusepath{clip}%
\pgfsetbuttcap%
\pgfsetmiterjoin%
\definecolor{currentfill}{rgb}{0.000000,0.568627,0.715686}%
\pgfsetfillcolor{currentfill}%
\pgfsetlinewidth{0.000000pt}%
\definecolor{currentstroke}{rgb}{0.000000,0.000000,0.000000}%
\pgfsetstrokecolor{currentstroke}%
\pgfsetstrokeopacity{0.000000}%
\pgfsetdash{}{0pt}%
\pgfpathmoveto{\pgfqpoint{1.782538in}{1.034546in}}%
\pgfpathlineto{\pgfqpoint{1.785102in}{1.081107in}}%
\pgfpathlineto{\pgfqpoint{1.786242in}{1.101806in}}%
\pgfpathlineto{\pgfqpoint{1.792285in}{1.107896in}}%
\pgfpathlineto{\pgfqpoint{1.798623in}{1.104938in}}%
\pgfpathlineto{\pgfqpoint{1.804299in}{1.099161in}}%
\pgfpathlineto{\pgfqpoint{1.814486in}{1.099133in}}%
\pgfpathlineto{\pgfqpoint{1.818722in}{1.107444in}}%
\pgfpathlineto{\pgfqpoint{1.826418in}{1.110191in}}%
\pgfpathlineto{\pgfqpoint{1.825972in}{1.094304in}}%
\pgfpathlineto{\pgfqpoint{1.873684in}{1.092215in}}%
\pgfpathlineto{\pgfqpoint{1.873435in}{1.071561in}}%
\pgfpathlineto{\pgfqpoint{1.872137in}{1.035604in}}%
\pgfpathlineto{\pgfqpoint{1.863914in}{1.037452in}}%
\pgfpathlineto{\pgfqpoint{1.796992in}{1.040622in}}%
\pgfpathlineto{\pgfqpoint{1.796623in}{1.033765in}}%
\pgfpathlineto{\pgfqpoint{1.782538in}{1.034546in}}%
\pgfpathclose%
\pgfusepath{fill}%
\end{pgfscope}%
\begin{pgfscope}%
\pgfpathrectangle{\pgfqpoint{0.100000in}{0.100000in}}{\pgfqpoint{3.608454in}{2.310000in}}%
\pgfusepath{clip}%
\pgfsetbuttcap%
\pgfsetmiterjoin%
\definecolor{currentfill}{rgb}{0.000000,0.682353,0.658824}%
\pgfsetfillcolor{currentfill}%
\pgfsetlinewidth{0.000000pt}%
\definecolor{currentstroke}{rgb}{0.000000,0.000000,0.000000}%
\pgfsetstrokecolor{currentstroke}%
\pgfsetstrokeopacity{0.000000}%
\pgfsetdash{}{0pt}%
\pgfpathmoveto{\pgfqpoint{3.008758in}{1.422029in}}%
\pgfpathlineto{\pgfqpoint{3.021096in}{1.413057in}}%
\pgfpathlineto{\pgfqpoint{3.018667in}{1.406375in}}%
\pgfpathlineto{\pgfqpoint{3.019598in}{1.394420in}}%
\pgfpathlineto{\pgfqpoint{3.012891in}{1.393149in}}%
\pgfpathlineto{\pgfqpoint{3.010884in}{1.388904in}}%
\pgfpathlineto{\pgfqpoint{3.001117in}{1.390072in}}%
\pgfpathlineto{\pgfqpoint{2.997346in}{1.397787in}}%
\pgfpathlineto{\pgfqpoint{2.989078in}{1.401477in}}%
\pgfpathlineto{\pgfqpoint{2.980404in}{1.409102in}}%
\pgfpathlineto{\pgfqpoint{2.978628in}{1.415804in}}%
\pgfpathlineto{\pgfqpoint{2.986872in}{1.422883in}}%
\pgfpathlineto{\pgfqpoint{3.004458in}{1.420613in}}%
\pgfpathlineto{\pgfqpoint{3.008758in}{1.422029in}}%
\pgfpathclose%
\pgfusepath{fill}%
\end{pgfscope}%
\begin{pgfscope}%
\pgfpathrectangle{\pgfqpoint{0.100000in}{0.100000in}}{\pgfqpoint{3.608454in}{2.310000in}}%
\pgfusepath{clip}%
\pgfsetbuttcap%
\pgfsetmiterjoin%
\definecolor{currentfill}{rgb}{0.000000,0.478431,0.760784}%
\pgfsetfillcolor{currentfill}%
\pgfsetlinewidth{0.000000pt}%
\definecolor{currentstroke}{rgb}{0.000000,0.000000,0.000000}%
\pgfsetstrokecolor{currentstroke}%
\pgfsetstrokeopacity{0.000000}%
\pgfsetdash{}{0pt}%
\pgfpathmoveto{\pgfqpoint{2.814875in}{1.291146in}}%
\pgfpathlineto{\pgfqpoint{2.810307in}{1.281879in}}%
\pgfpathlineto{\pgfqpoint{2.801059in}{1.283583in}}%
\pgfpathlineto{\pgfqpoint{2.791570in}{1.288910in}}%
\pgfpathlineto{\pgfqpoint{2.774318in}{1.293523in}}%
\pgfpathlineto{\pgfqpoint{2.759930in}{1.280063in}}%
\pgfpathlineto{\pgfqpoint{2.757404in}{1.276520in}}%
\pgfpathlineto{\pgfqpoint{2.740147in}{1.292411in}}%
\pgfpathlineto{\pgfqpoint{2.733018in}{1.291688in}}%
\pgfpathlineto{\pgfqpoint{2.727518in}{1.286122in}}%
\pgfpathlineto{\pgfqpoint{2.719889in}{1.285770in}}%
\pgfpathlineto{\pgfqpoint{2.718323in}{1.291481in}}%
\pgfpathlineto{\pgfqpoint{2.724571in}{1.300390in}}%
\pgfpathlineto{\pgfqpoint{2.726400in}{1.307715in}}%
\pgfpathlineto{\pgfqpoint{2.726160in}{1.316456in}}%
\pgfpathlineto{\pgfqpoint{2.730992in}{1.324932in}}%
\pgfpathlineto{\pgfqpoint{2.724859in}{1.333241in}}%
\pgfpathlineto{\pgfqpoint{2.722897in}{1.338929in}}%
\pgfpathlineto{\pgfqpoint{2.718732in}{1.341235in}}%
\pgfpathlineto{\pgfqpoint{2.726746in}{1.347306in}}%
\pgfpathlineto{\pgfqpoint{2.735344in}{1.352880in}}%
\pgfpathlineto{\pgfqpoint{2.736943in}{1.344983in}}%
\pgfpathlineto{\pgfqpoint{2.750154in}{1.334910in}}%
\pgfpathlineto{\pgfqpoint{2.760419in}{1.343088in}}%
\pgfpathlineto{\pgfqpoint{2.772188in}{1.346181in}}%
\pgfpathlineto{\pgfqpoint{2.779558in}{1.348364in}}%
\pgfpathlineto{\pgfqpoint{2.786720in}{1.347970in}}%
\pgfpathlineto{\pgfqpoint{2.789928in}{1.341177in}}%
\pgfpathlineto{\pgfqpoint{2.787097in}{1.337635in}}%
\pgfpathlineto{\pgfqpoint{2.796848in}{1.326453in}}%
\pgfpathlineto{\pgfqpoint{2.784277in}{1.310449in}}%
\pgfpathlineto{\pgfqpoint{2.793467in}{1.297141in}}%
\pgfpathlineto{\pgfqpoint{2.806402in}{1.296003in}}%
\pgfpathlineto{\pgfqpoint{2.814875in}{1.291146in}}%
\pgfpathclose%
\pgfusepath{fill}%
\end{pgfscope}%
\begin{pgfscope}%
\pgfpathrectangle{\pgfqpoint{0.100000in}{0.100000in}}{\pgfqpoint{3.608454in}{2.310000in}}%
\pgfusepath{clip}%
\pgfsetbuttcap%
\pgfsetmiterjoin%
\definecolor{currentfill}{rgb}{0.000000,0.619608,0.690196}%
\pgfsetfillcolor{currentfill}%
\pgfsetlinewidth{0.000000pt}%
\definecolor{currentstroke}{rgb}{0.000000,0.000000,0.000000}%
\pgfsetstrokecolor{currentstroke}%
\pgfsetstrokeopacity{0.000000}%
\pgfsetdash{}{0pt}%
\pgfpathmoveto{\pgfqpoint{2.657467in}{1.463709in}}%
\pgfpathlineto{\pgfqpoint{2.634617in}{1.461203in}}%
\pgfpathlineto{\pgfqpoint{2.633399in}{1.473636in}}%
\pgfpathlineto{\pgfqpoint{2.625004in}{1.477436in}}%
\pgfpathlineto{\pgfqpoint{2.623185in}{1.497972in}}%
\pgfpathlineto{\pgfqpoint{2.614139in}{1.497101in}}%
\pgfpathlineto{\pgfqpoint{2.610441in}{1.500202in}}%
\pgfpathlineto{\pgfqpoint{2.609172in}{1.513996in}}%
\pgfpathlineto{\pgfqpoint{2.588380in}{1.512229in}}%
\pgfpathlineto{\pgfqpoint{2.586580in}{1.532701in}}%
\pgfpathlineto{\pgfqpoint{2.614051in}{1.535158in}}%
\pgfpathlineto{\pgfqpoint{2.615905in}{1.514539in}}%
\pgfpathlineto{\pgfqpoint{2.633720in}{1.516230in}}%
\pgfpathlineto{\pgfqpoint{2.633101in}{1.523054in}}%
\pgfpathlineto{\pgfqpoint{2.646259in}{1.524584in}}%
\pgfpathlineto{\pgfqpoint{2.649395in}{1.497238in}}%
\pgfpathlineto{\pgfqpoint{2.653860in}{1.497742in}}%
\pgfpathlineto{\pgfqpoint{2.657467in}{1.463709in}}%
\pgfpathclose%
\pgfusepath{fill}%
\end{pgfscope}%
\begin{pgfscope}%
\pgfpathrectangle{\pgfqpoint{0.100000in}{0.100000in}}{\pgfqpoint{3.608454in}{2.310000in}}%
\pgfusepath{clip}%
\pgfsetbuttcap%
\pgfsetmiterjoin%
\definecolor{currentfill}{rgb}{0.000000,0.713725,0.643137}%
\pgfsetfillcolor{currentfill}%
\pgfsetlinewidth{0.000000pt}%
\definecolor{currentstroke}{rgb}{0.000000,0.000000,0.000000}%
\pgfsetstrokecolor{currentstroke}%
\pgfsetstrokeopacity{0.000000}%
\pgfsetdash{}{0pt}%
\pgfpathmoveto{\pgfqpoint{2.356606in}{1.028702in}}%
\pgfpathlineto{\pgfqpoint{2.347372in}{1.026940in}}%
\pgfpathlineto{\pgfqpoint{2.342574in}{1.032044in}}%
\pgfpathlineto{\pgfqpoint{2.335892in}{1.034757in}}%
\pgfpathlineto{\pgfqpoint{2.328181in}{1.034618in}}%
\pgfpathlineto{\pgfqpoint{2.328311in}{1.029359in}}%
\pgfpathlineto{\pgfqpoint{2.320051in}{1.032034in}}%
\pgfpathlineto{\pgfqpoint{2.314200in}{1.029274in}}%
\pgfpathlineto{\pgfqpoint{2.309489in}{1.028555in}}%
\pgfpathlineto{\pgfqpoint{2.300760in}{1.031780in}}%
\pgfpathlineto{\pgfqpoint{2.293540in}{1.031438in}}%
\pgfpathlineto{\pgfqpoint{2.293330in}{1.054914in}}%
\pgfpathlineto{\pgfqpoint{2.310094in}{1.055119in}}%
\pgfpathlineto{\pgfqpoint{2.313444in}{1.062948in}}%
\pgfpathlineto{\pgfqpoint{2.313347in}{1.069167in}}%
\pgfpathlineto{\pgfqpoint{2.340700in}{1.069582in}}%
\pgfpathlineto{\pgfqpoint{2.340391in}{1.078576in}}%
\pgfpathlineto{\pgfqpoint{2.343663in}{1.088662in}}%
\pgfpathlineto{\pgfqpoint{2.350507in}{1.098801in}}%
\pgfpathlineto{\pgfqpoint{2.361164in}{1.098864in}}%
\pgfpathlineto{\pgfqpoint{2.362249in}{1.063873in}}%
\pgfpathlineto{\pgfqpoint{2.363393in}{1.028869in}}%
\pgfpathlineto{\pgfqpoint{2.356606in}{1.028702in}}%
\pgfpathclose%
\pgfusepath{fill}%
\end{pgfscope}%
\begin{pgfscope}%
\pgfpathrectangle{\pgfqpoint{0.100000in}{0.100000in}}{\pgfqpoint{3.608454in}{2.310000in}}%
\pgfusepath{clip}%
\pgfsetbuttcap%
\pgfsetmiterjoin%
\definecolor{currentfill}{rgb}{0.000000,0.337255,0.831373}%
\pgfsetfillcolor{currentfill}%
\pgfsetlinewidth{0.000000pt}%
\definecolor{currentstroke}{rgb}{0.000000,0.000000,0.000000}%
\pgfsetstrokecolor{currentstroke}%
\pgfsetstrokeopacity{0.000000}%
\pgfsetdash{}{0pt}%
\pgfpathmoveto{\pgfqpoint{1.844714in}{0.901800in}}%
\pgfpathlineto{\pgfqpoint{1.846276in}{0.936427in}}%
\pgfpathlineto{\pgfqpoint{1.847621in}{0.966431in}}%
\pgfpathlineto{\pgfqpoint{1.852516in}{0.961694in}}%
\pgfpathlineto{\pgfqpoint{1.864563in}{0.959406in}}%
\pgfpathlineto{\pgfqpoint{1.869882in}{0.961093in}}%
\pgfpathlineto{\pgfqpoint{1.877716in}{0.953284in}}%
\pgfpathlineto{\pgfqpoint{1.883512in}{0.956032in}}%
\pgfpathlineto{\pgfqpoint{1.885921in}{0.960469in}}%
\pgfpathlineto{\pgfqpoint{1.898581in}{0.956614in}}%
\pgfpathlineto{\pgfqpoint{1.906317in}{0.946847in}}%
\pgfpathlineto{\pgfqpoint{1.913020in}{0.946294in}}%
\pgfpathlineto{\pgfqpoint{1.910664in}{0.938402in}}%
\pgfpathlineto{\pgfqpoint{1.909495in}{0.904858in}}%
\pgfpathlineto{\pgfqpoint{1.880200in}{0.905864in}}%
\pgfpathlineto{\pgfqpoint{1.880039in}{0.900259in}}%
\pgfpathlineto{\pgfqpoint{1.844714in}{0.901800in}}%
\pgfpathclose%
\pgfusepath{fill}%
\end{pgfscope}%
\begin{pgfscope}%
\pgfpathrectangle{\pgfqpoint{0.100000in}{0.100000in}}{\pgfqpoint{3.608454in}{2.310000in}}%
\pgfusepath{clip}%
\pgfsetbuttcap%
\pgfsetmiterjoin%
\definecolor{currentfill}{rgb}{0.000000,0.627451,0.686275}%
\pgfsetfillcolor{currentfill}%
\pgfsetlinewidth{0.000000pt}%
\definecolor{currentstroke}{rgb}{0.000000,0.000000,0.000000}%
\pgfsetstrokecolor{currentstroke}%
\pgfsetstrokeopacity{0.000000}%
\pgfsetdash{}{0pt}%
\pgfpathmoveto{\pgfqpoint{2.227204in}{1.266822in}}%
\pgfpathlineto{\pgfqpoint{2.199885in}{1.267142in}}%
\pgfpathlineto{\pgfqpoint{2.199587in}{1.254527in}}%
\pgfpathlineto{\pgfqpoint{2.196058in}{1.254577in}}%
\pgfpathlineto{\pgfqpoint{2.192435in}{1.247763in}}%
\pgfpathlineto{\pgfqpoint{2.180991in}{1.248016in}}%
\pgfpathlineto{\pgfqpoint{2.181129in}{1.252610in}}%
\pgfpathlineto{\pgfqpoint{2.165075in}{1.253073in}}%
\pgfpathlineto{\pgfqpoint{2.165424in}{1.263791in}}%
\pgfpathlineto{\pgfqpoint{2.164228in}{1.277853in}}%
\pgfpathlineto{\pgfqpoint{2.164617in}{1.305750in}}%
\pgfpathlineto{\pgfqpoint{2.198808in}{1.305319in}}%
\pgfpathlineto{\pgfqpoint{2.199219in}{1.335898in}}%
\pgfpathlineto{\pgfqpoint{2.226654in}{1.335261in}}%
\pgfpathlineto{\pgfqpoint{2.226458in}{1.316665in}}%
\pgfpathlineto{\pgfqpoint{2.225942in}{1.276603in}}%
\pgfpathlineto{\pgfqpoint{2.227204in}{1.266822in}}%
\pgfpathclose%
\pgfusepath{fill}%
\end{pgfscope}%
\begin{pgfscope}%
\pgfpathrectangle{\pgfqpoint{0.100000in}{0.100000in}}{\pgfqpoint{3.608454in}{2.310000in}}%
\pgfusepath{clip}%
\pgfsetbuttcap%
\pgfsetmiterjoin%
\definecolor{currentfill}{rgb}{0.000000,0.674510,0.662745}%
\pgfsetfillcolor{currentfill}%
\pgfsetlinewidth{0.000000pt}%
\definecolor{currentstroke}{rgb}{0.000000,0.000000,0.000000}%
\pgfsetstrokecolor{currentstroke}%
\pgfsetstrokeopacity{0.000000}%
\pgfsetdash{}{0pt}%
\pgfpathmoveto{\pgfqpoint{0.499716in}{1.160624in}}%
\pgfpathlineto{\pgfqpoint{0.490524in}{1.164969in}}%
\pgfpathlineto{\pgfqpoint{0.493536in}{1.171597in}}%
\pgfpathlineto{\pgfqpoint{0.503853in}{1.164499in}}%
\pgfpathlineto{\pgfqpoint{0.499716in}{1.160624in}}%
\pgfpathclose%
\pgfusepath{fill}%
\end{pgfscope}%
\begin{pgfscope}%
\pgfpathrectangle{\pgfqpoint{0.100000in}{0.100000in}}{\pgfqpoint{3.608454in}{2.310000in}}%
\pgfusepath{clip}%
\pgfsetbuttcap%
\pgfsetmiterjoin%
\definecolor{currentfill}{rgb}{0.000000,0.674510,0.662745}%
\pgfsetfillcolor{currentfill}%
\pgfsetlinewidth{0.000000pt}%
\definecolor{currentstroke}{rgb}{0.000000,0.000000,0.000000}%
\pgfsetstrokecolor{currentstroke}%
\pgfsetstrokeopacity{0.000000}%
\pgfsetdash{}{0pt}%
\pgfpathmoveto{\pgfqpoint{0.473802in}{1.162641in}}%
\pgfpathlineto{\pgfqpoint{0.468569in}{1.171802in}}%
\pgfpathlineto{\pgfqpoint{0.474440in}{1.173296in}}%
\pgfpathlineto{\pgfqpoint{0.483849in}{1.167231in}}%
\pgfpathlineto{\pgfqpoint{0.482039in}{1.164194in}}%
\pgfpathlineto{\pgfqpoint{0.473802in}{1.162641in}}%
\pgfpathclose%
\pgfusepath{fill}%
\end{pgfscope}%
\begin{pgfscope}%
\pgfpathrectangle{\pgfqpoint{0.100000in}{0.100000in}}{\pgfqpoint{3.608454in}{2.310000in}}%
\pgfusepath{clip}%
\pgfsetbuttcap%
\pgfsetmiterjoin%
\definecolor{currentfill}{rgb}{0.000000,0.674510,0.662745}%
\pgfsetfillcolor{currentfill}%
\pgfsetlinewidth{0.000000pt}%
\definecolor{currentstroke}{rgb}{0.000000,0.000000,0.000000}%
\pgfsetstrokecolor{currentstroke}%
\pgfsetstrokeopacity{0.000000}%
\pgfsetdash{}{0pt}%
\pgfpathmoveto{\pgfqpoint{0.537333in}{1.227732in}}%
\pgfpathlineto{\pgfqpoint{0.527921in}{1.193369in}}%
\pgfpathlineto{\pgfqpoint{0.524045in}{1.187970in}}%
\pgfpathlineto{\pgfqpoint{0.516406in}{1.193887in}}%
\pgfpathlineto{\pgfqpoint{0.508386in}{1.194136in}}%
\pgfpathlineto{\pgfqpoint{0.499604in}{1.197528in}}%
\pgfpathlineto{\pgfqpoint{0.493082in}{1.203453in}}%
\pgfpathlineto{\pgfqpoint{0.479238in}{1.208664in}}%
\pgfpathlineto{\pgfqpoint{0.464027in}{1.210445in}}%
\pgfpathlineto{\pgfqpoint{0.462181in}{1.217763in}}%
\pgfpathlineto{\pgfqpoint{0.454827in}{1.224336in}}%
\pgfpathlineto{\pgfqpoint{0.460433in}{1.233714in}}%
\pgfpathlineto{\pgfqpoint{0.459304in}{1.237786in}}%
\pgfpathlineto{\pgfqpoint{0.463136in}{1.244048in}}%
\pgfpathlineto{\pgfqpoint{0.460490in}{1.249495in}}%
\pgfpathlineto{\pgfqpoint{0.464659in}{1.256580in}}%
\pgfpathlineto{\pgfqpoint{0.467179in}{1.267774in}}%
\pgfpathlineto{\pgfqpoint{0.460247in}{1.270875in}}%
\pgfpathlineto{\pgfqpoint{0.454219in}{1.280640in}}%
\pgfpathlineto{\pgfqpoint{0.458963in}{1.289133in}}%
\pgfpathlineto{\pgfqpoint{0.457966in}{1.295353in}}%
\pgfpathlineto{\pgfqpoint{0.452130in}{1.297956in}}%
\pgfpathlineto{\pgfqpoint{0.446221in}{1.313638in}}%
\pgfpathlineto{\pgfqpoint{0.439451in}{1.318713in}}%
\pgfpathlineto{\pgfqpoint{0.438474in}{1.329831in}}%
\pgfpathlineto{\pgfqpoint{0.468241in}{1.320678in}}%
\pgfpathlineto{\pgfqpoint{0.508309in}{1.308909in}}%
\pgfpathlineto{\pgfqpoint{0.509532in}{1.308555in}}%
\pgfpathlineto{\pgfqpoint{0.505758in}{1.295236in}}%
\pgfpathlineto{\pgfqpoint{0.512448in}{1.293350in}}%
\pgfpathlineto{\pgfqpoint{0.510504in}{1.286661in}}%
\pgfpathlineto{\pgfqpoint{0.514797in}{1.280686in}}%
\pgfpathlineto{\pgfqpoint{0.521448in}{1.276381in}}%
\pgfpathlineto{\pgfqpoint{0.519557in}{1.269679in}}%
\pgfpathlineto{\pgfqpoint{0.523955in}{1.268417in}}%
\pgfpathlineto{\pgfqpoint{0.522085in}{1.261761in}}%
\pgfpathlineto{\pgfqpoint{0.530958in}{1.259199in}}%
\pgfpathlineto{\pgfqpoint{0.529070in}{1.252510in}}%
\pgfpathlineto{\pgfqpoint{0.536255in}{1.250941in}}%
\pgfpathlineto{\pgfqpoint{0.533846in}{1.244020in}}%
\pgfpathlineto{\pgfqpoint{0.539132in}{1.241671in}}%
\pgfpathlineto{\pgfqpoint{0.537333in}{1.227732in}}%
\pgfpathclose%
\pgfusepath{fill}%
\end{pgfscope}%
\begin{pgfscope}%
\pgfpathrectangle{\pgfqpoint{0.100000in}{0.100000in}}{\pgfqpoint{3.608454in}{2.310000in}}%
\pgfusepath{clip}%
\pgfsetbuttcap%
\pgfsetmiterjoin%
\definecolor{currentfill}{rgb}{0.000000,0.415686,0.792157}%
\pgfsetfillcolor{currentfill}%
\pgfsetlinewidth{0.000000pt}%
\definecolor{currentstroke}{rgb}{0.000000,0.000000,0.000000}%
\pgfsetstrokecolor{currentstroke}%
\pgfsetstrokeopacity{0.000000}%
\pgfsetdash{}{0pt}%
\pgfpathmoveto{\pgfqpoint{3.370781in}{1.834556in}}%
\pgfpathlineto{\pgfqpoint{3.348859in}{1.830296in}}%
\pgfpathlineto{\pgfqpoint{3.344700in}{1.848799in}}%
\pgfpathlineto{\pgfqpoint{3.333464in}{1.857292in}}%
\pgfpathlineto{\pgfqpoint{3.332649in}{1.866641in}}%
\pgfpathlineto{\pgfqpoint{3.324985in}{1.884098in}}%
\pgfpathlineto{\pgfqpoint{3.324873in}{1.895982in}}%
\pgfpathlineto{\pgfqpoint{3.327709in}{1.903478in}}%
\pgfpathlineto{\pgfqpoint{3.324626in}{1.910739in}}%
\pgfpathlineto{\pgfqpoint{3.323703in}{1.920977in}}%
\pgfpathlineto{\pgfqpoint{3.316899in}{1.928969in}}%
\pgfpathlineto{\pgfqpoint{3.316446in}{1.944043in}}%
\pgfpathlineto{\pgfqpoint{3.313379in}{1.946070in}}%
\pgfpathlineto{\pgfqpoint{3.314088in}{1.952894in}}%
\pgfpathlineto{\pgfqpoint{3.312126in}{1.959898in}}%
\pgfpathlineto{\pgfqpoint{3.351137in}{1.969833in}}%
\pgfpathlineto{\pgfqpoint{3.355196in}{1.970349in}}%
\pgfpathlineto{\pgfqpoint{3.359777in}{1.957182in}}%
\pgfpathlineto{\pgfqpoint{3.358298in}{1.952554in}}%
\pgfpathlineto{\pgfqpoint{3.368238in}{1.950258in}}%
\pgfpathlineto{\pgfqpoint{3.366415in}{1.942507in}}%
\pgfpathlineto{\pgfqpoint{3.373495in}{1.940927in}}%
\pgfpathlineto{\pgfqpoint{3.371539in}{1.933332in}}%
\pgfpathlineto{\pgfqpoint{3.370017in}{1.926361in}}%
\pgfpathlineto{\pgfqpoint{3.363137in}{1.929015in}}%
\pgfpathlineto{\pgfqpoint{3.361540in}{1.921976in}}%
\pgfpathlineto{\pgfqpoint{3.352226in}{1.924022in}}%
\pgfpathlineto{\pgfqpoint{3.349620in}{1.909475in}}%
\pgfpathlineto{\pgfqpoint{3.350605in}{1.900061in}}%
\pgfpathlineto{\pgfqpoint{3.353338in}{1.893246in}}%
\pgfpathlineto{\pgfqpoint{3.363839in}{1.892385in}}%
\pgfpathlineto{\pgfqpoint{3.358819in}{1.880708in}}%
\pgfpathlineto{\pgfqpoint{3.354162in}{1.876806in}}%
\pgfpathlineto{\pgfqpoint{3.365877in}{1.874467in}}%
\pgfpathlineto{\pgfqpoint{3.365489in}{1.867233in}}%
\pgfpathlineto{\pgfqpoint{3.372767in}{1.866271in}}%
\pgfpathlineto{\pgfqpoint{3.371356in}{1.850568in}}%
\pgfpathlineto{\pgfqpoint{3.375426in}{1.849437in}}%
\pgfpathlineto{\pgfqpoint{3.375120in}{1.840353in}}%
\pgfpathlineto{\pgfqpoint{3.370781in}{1.834556in}}%
\pgfpathclose%
\pgfusepath{fill}%
\end{pgfscope}%
\begin{pgfscope}%
\pgfpathrectangle{\pgfqpoint{0.100000in}{0.100000in}}{\pgfqpoint{3.608454in}{2.310000in}}%
\pgfusepath{clip}%
\pgfsetbuttcap%
\pgfsetmiterjoin%
\definecolor{currentfill}{rgb}{0.000000,0.745098,0.627451}%
\pgfsetfillcolor{currentfill}%
\pgfsetlinewidth{0.000000pt}%
\definecolor{currentstroke}{rgb}{0.000000,0.000000,0.000000}%
\pgfsetstrokecolor{currentstroke}%
\pgfsetstrokeopacity{0.000000}%
\pgfsetdash{}{0pt}%
\pgfpathmoveto{\pgfqpoint{2.342946in}{0.805709in}}%
\pgfpathlineto{\pgfqpoint{2.340198in}{0.812791in}}%
\pgfpathlineto{\pgfqpoint{2.344014in}{0.814843in}}%
\pgfpathlineto{\pgfqpoint{2.342569in}{0.828917in}}%
\pgfpathlineto{\pgfqpoint{2.344269in}{0.832393in}}%
\pgfpathlineto{\pgfqpoint{2.365973in}{0.833290in}}%
\pgfpathlineto{\pgfqpoint{2.371879in}{0.836595in}}%
\pgfpathlineto{\pgfqpoint{2.377206in}{0.836777in}}%
\pgfpathlineto{\pgfqpoint{2.378401in}{0.826111in}}%
\pgfpathlineto{\pgfqpoint{2.384564in}{0.826478in}}%
\pgfpathlineto{\pgfqpoint{2.384135in}{0.837054in}}%
\pgfpathlineto{\pgfqpoint{2.393171in}{0.837522in}}%
\pgfpathlineto{\pgfqpoint{2.393376in}{0.840961in}}%
\pgfpathlineto{\pgfqpoint{2.404935in}{0.832654in}}%
\pgfpathlineto{\pgfqpoint{2.401669in}{0.830402in}}%
\pgfpathlineto{\pgfqpoint{2.399936in}{0.821863in}}%
\pgfpathlineto{\pgfqpoint{2.395425in}{0.819431in}}%
\pgfpathlineto{\pgfqpoint{2.394278in}{0.809641in}}%
\pgfpathlineto{\pgfqpoint{2.396163in}{0.790004in}}%
\pgfpathlineto{\pgfqpoint{2.395693in}{0.767722in}}%
\pgfpathlineto{\pgfqpoint{2.395986in}{0.760766in}}%
\pgfpathlineto{\pgfqpoint{2.367747in}{0.759421in}}%
\pgfpathlineto{\pgfqpoint{2.357930in}{0.763773in}}%
\pgfpathlineto{\pgfqpoint{2.353993in}{0.777417in}}%
\pgfpathlineto{\pgfqpoint{2.363437in}{0.782994in}}%
\pgfpathlineto{\pgfqpoint{2.364252in}{0.787285in}}%
\pgfpathlineto{\pgfqpoint{2.371635in}{0.792020in}}%
\pgfpathlineto{\pgfqpoint{2.374090in}{0.798791in}}%
\pgfpathlineto{\pgfqpoint{2.365200in}{0.801106in}}%
\pgfpathlineto{\pgfqpoint{2.366944in}{0.810346in}}%
\pgfpathlineto{\pgfqpoint{2.354522in}{0.806095in}}%
\pgfpathlineto{\pgfqpoint{2.342946in}{0.805709in}}%
\pgfpathclose%
\pgfusepath{fill}%
\end{pgfscope}%
\begin{pgfscope}%
\pgfpathrectangle{\pgfqpoint{0.100000in}{0.100000in}}{\pgfqpoint{3.608454in}{2.310000in}}%
\pgfusepath{clip}%
\pgfsetbuttcap%
\pgfsetmiterjoin%
\definecolor{currentfill}{rgb}{0.000000,0.760784,0.619608}%
\pgfsetfillcolor{currentfill}%
\pgfsetlinewidth{0.000000pt}%
\definecolor{currentstroke}{rgb}{0.000000,0.000000,0.000000}%
\pgfsetstrokecolor{currentstroke}%
\pgfsetstrokeopacity{0.000000}%
\pgfsetdash{}{0pt}%
\pgfpathmoveto{\pgfqpoint{0.509532in}{1.308555in}}%
\pgfpathlineto{\pgfqpoint{0.508309in}{1.308909in}}%
\pgfpathlineto{\pgfqpoint{0.503087in}{1.324821in}}%
\pgfpathlineto{\pgfqpoint{0.498712in}{1.326448in}}%
\pgfpathlineto{\pgfqpoint{0.493971in}{1.332372in}}%
\pgfpathlineto{\pgfqpoint{0.488481in}{1.345190in}}%
\pgfpathlineto{\pgfqpoint{0.492171in}{1.356794in}}%
\pgfpathlineto{\pgfqpoint{0.496557in}{1.356685in}}%
\pgfpathlineto{\pgfqpoint{0.499634in}{1.368954in}}%
\pgfpathlineto{\pgfqpoint{0.486008in}{1.393700in}}%
\pgfpathlineto{\pgfqpoint{0.520707in}{1.412942in}}%
\pgfpathlineto{\pgfqpoint{0.533257in}{1.413969in}}%
\pgfpathlineto{\pgfqpoint{0.548331in}{1.412190in}}%
\pgfpathlineto{\pgfqpoint{0.570915in}{1.425063in}}%
\pgfpathlineto{\pgfqpoint{0.577622in}{1.423252in}}%
\pgfpathlineto{\pgfqpoint{0.583298in}{1.427984in}}%
\pgfpathlineto{\pgfqpoint{0.584724in}{1.432975in}}%
\pgfpathlineto{\pgfqpoint{0.605864in}{1.444998in}}%
\pgfpathlineto{\pgfqpoint{0.607467in}{1.441382in}}%
\pgfpathlineto{\pgfqpoint{0.616083in}{1.438599in}}%
\pgfpathlineto{\pgfqpoint{0.618630in}{1.433861in}}%
\pgfpathlineto{\pgfqpoint{0.619173in}{1.425619in}}%
\pgfpathlineto{\pgfqpoint{0.624814in}{1.421196in}}%
\pgfpathlineto{\pgfqpoint{0.627031in}{1.414709in}}%
\pgfpathlineto{\pgfqpoint{0.632163in}{1.410294in}}%
\pgfpathlineto{\pgfqpoint{0.628697in}{1.402513in}}%
\pgfpathlineto{\pgfqpoint{0.632454in}{1.400962in}}%
\pgfpathlineto{\pgfqpoint{0.632393in}{1.385724in}}%
\pgfpathlineto{\pgfqpoint{0.636091in}{1.384439in}}%
\pgfpathlineto{\pgfqpoint{0.643441in}{1.374382in}}%
\pgfpathlineto{\pgfqpoint{0.645186in}{1.361551in}}%
\pgfpathlineto{\pgfqpoint{0.642344in}{1.357942in}}%
\pgfpathlineto{\pgfqpoint{0.641670in}{1.344858in}}%
\pgfpathlineto{\pgfqpoint{0.644925in}{1.333356in}}%
\pgfpathlineto{\pgfqpoint{0.644042in}{1.329081in}}%
\pgfpathlineto{\pgfqpoint{0.650356in}{1.315914in}}%
\pgfpathlineto{\pgfqpoint{0.647795in}{1.312826in}}%
\pgfpathlineto{\pgfqpoint{0.649389in}{1.301776in}}%
\pgfpathlineto{\pgfqpoint{0.647370in}{1.296028in}}%
\pgfpathlineto{\pgfqpoint{0.648713in}{1.279430in}}%
\pgfpathlineto{\pgfqpoint{0.645007in}{1.271759in}}%
\pgfpathlineto{\pgfqpoint{0.606581in}{1.281871in}}%
\pgfpathlineto{\pgfqpoint{0.550101in}{1.297200in}}%
\pgfpathlineto{\pgfqpoint{0.509532in}{1.308555in}}%
\pgfpathclose%
\pgfusepath{fill}%
\end{pgfscope}%
\begin{pgfscope}%
\pgfpathrectangle{\pgfqpoint{0.100000in}{0.100000in}}{\pgfqpoint{3.608454in}{2.310000in}}%
\pgfusepath{clip}%
\pgfsetbuttcap%
\pgfsetmiterjoin%
\definecolor{currentfill}{rgb}{0.000000,0.564706,0.717647}%
\pgfsetfillcolor{currentfill}%
\pgfsetlinewidth{0.000000pt}%
\definecolor{currentstroke}{rgb}{0.000000,0.000000,0.000000}%
\pgfsetstrokecolor{currentstroke}%
\pgfsetstrokeopacity{0.000000}%
\pgfsetdash{}{0pt}%
\pgfpathmoveto{\pgfqpoint{1.007808in}{2.262053in}}%
\pgfpathlineto{\pgfqpoint{1.050960in}{2.251947in}}%
\pgfpathlineto{\pgfqpoint{1.075629in}{2.246526in}}%
\pgfpathlineto{\pgfqpoint{1.074794in}{2.233716in}}%
\pgfpathlineto{\pgfqpoint{1.075971in}{2.228263in}}%
\pgfpathlineto{\pgfqpoint{1.072055in}{2.221371in}}%
\pgfpathlineto{\pgfqpoint{1.061432in}{2.221573in}}%
\pgfpathlineto{\pgfqpoint{1.061972in}{2.214474in}}%
\pgfpathlineto{\pgfqpoint{1.056030in}{2.187273in}}%
\pgfpathlineto{\pgfqpoint{1.047061in}{2.189256in}}%
\pgfpathlineto{\pgfqpoint{1.043899in}{2.172864in}}%
\pgfpathlineto{\pgfqpoint{1.036396in}{2.174561in}}%
\pgfpathlineto{\pgfqpoint{1.033549in}{2.168208in}}%
\pgfpathlineto{\pgfqpoint{1.027896in}{2.166208in}}%
\pgfpathlineto{\pgfqpoint{1.015079in}{2.170487in}}%
\pgfpathlineto{\pgfqpoint{1.015735in}{2.176278in}}%
\pgfpathlineto{\pgfqpoint{1.012094in}{2.185135in}}%
\pgfpathlineto{\pgfqpoint{1.011635in}{2.199142in}}%
\pgfpathlineto{\pgfqpoint{0.994462in}{2.195826in}}%
\pgfpathlineto{\pgfqpoint{0.993310in}{2.200722in}}%
\pgfpathlineto{\pgfqpoint{0.998620in}{2.223097in}}%
\pgfpathlineto{\pgfqpoint{1.007808in}{2.262053in}}%
\pgfpathclose%
\pgfusepath{fill}%
\end{pgfscope}%
\begin{pgfscope}%
\pgfpathrectangle{\pgfqpoint{0.100000in}{0.100000in}}{\pgfqpoint{3.608454in}{2.310000in}}%
\pgfusepath{clip}%
\pgfsetbuttcap%
\pgfsetmiterjoin%
\definecolor{currentfill}{rgb}{0.000000,0.631373,0.684314}%
\pgfsetfillcolor{currentfill}%
\pgfsetlinewidth{0.000000pt}%
\definecolor{currentstroke}{rgb}{0.000000,0.000000,0.000000}%
\pgfsetstrokecolor{currentstroke}%
\pgfsetstrokeopacity{0.000000}%
\pgfsetdash{}{0pt}%
\pgfpathmoveto{\pgfqpoint{2.515690in}{0.998385in}}%
\pgfpathlineto{\pgfqpoint{2.515503in}{1.018069in}}%
\pgfpathlineto{\pgfqpoint{2.512709in}{1.025980in}}%
\pgfpathlineto{\pgfqpoint{2.508135in}{1.025655in}}%
\pgfpathlineto{\pgfqpoint{2.507414in}{1.036455in}}%
\pgfpathlineto{\pgfqpoint{2.509787in}{1.036620in}}%
\pgfpathlineto{\pgfqpoint{2.508734in}{1.056586in}}%
\pgfpathlineto{\pgfqpoint{2.512736in}{1.056859in}}%
\pgfpathlineto{\pgfqpoint{2.517201in}{1.066138in}}%
\pgfpathlineto{\pgfqpoint{2.521534in}{1.068158in}}%
\pgfpathlineto{\pgfqpoint{2.535012in}{1.069064in}}%
\pgfpathlineto{\pgfqpoint{2.534914in}{1.072050in}}%
\pgfpathlineto{\pgfqpoint{2.550818in}{1.070538in}}%
\pgfpathlineto{\pgfqpoint{2.556697in}{1.071568in}}%
\pgfpathlineto{\pgfqpoint{2.560066in}{1.064151in}}%
\pgfpathlineto{\pgfqpoint{2.561622in}{1.041248in}}%
\pgfpathlineto{\pgfqpoint{2.547499in}{1.040339in}}%
\pgfpathlineto{\pgfqpoint{2.554995in}{1.031701in}}%
\pgfpathlineto{\pgfqpoint{2.553769in}{0.997504in}}%
\pgfpathlineto{\pgfqpoint{2.528607in}{0.995823in}}%
\pgfpathlineto{\pgfqpoint{2.528379in}{0.999297in}}%
\pgfpathlineto{\pgfqpoint{2.515690in}{0.998385in}}%
\pgfpathclose%
\pgfusepath{fill}%
\end{pgfscope}%
\begin{pgfscope}%
\pgfpathrectangle{\pgfqpoint{0.100000in}{0.100000in}}{\pgfqpoint{3.608454in}{2.310000in}}%
\pgfusepath{clip}%
\pgfsetbuttcap%
\pgfsetmiterjoin%
\definecolor{currentfill}{rgb}{0.000000,0.992157,0.503922}%
\pgfsetfillcolor{currentfill}%
\pgfsetlinewidth{0.000000pt}%
\definecolor{currentstroke}{rgb}{0.000000,0.000000,0.000000}%
\pgfsetstrokecolor{currentstroke}%
\pgfsetstrokeopacity{0.000000}%
\pgfsetdash{}{0pt}%
\pgfpathmoveto{\pgfqpoint{0.647635in}{2.272305in}}%
\pgfpathlineto{\pgfqpoint{0.648220in}{2.267247in}}%
\pgfpathlineto{\pgfqpoint{0.639739in}{2.262486in}}%
\pgfpathlineto{\pgfqpoint{0.635231in}{2.255252in}}%
\pgfpathlineto{\pgfqpoint{0.623636in}{2.252362in}}%
\pgfpathlineto{\pgfqpoint{0.595222in}{2.261545in}}%
\pgfpathlineto{\pgfqpoint{0.593017in}{2.254746in}}%
\pgfpathlineto{\pgfqpoint{0.565186in}{2.263908in}}%
\pgfpathlineto{\pgfqpoint{0.549782in}{2.270356in}}%
\pgfpathlineto{\pgfqpoint{0.551434in}{2.287191in}}%
\pgfpathlineto{\pgfqpoint{0.549994in}{2.294555in}}%
\pgfpathlineto{\pgfqpoint{0.544785in}{2.302357in}}%
\pgfpathlineto{\pgfqpoint{0.544644in}{2.308702in}}%
\pgfpathlineto{\pgfqpoint{0.548618in}{2.325631in}}%
\pgfpathlineto{\pgfqpoint{0.557234in}{2.339937in}}%
\pgfpathlineto{\pgfqpoint{0.566992in}{2.328262in}}%
\pgfpathlineto{\pgfqpoint{0.573972in}{2.323699in}}%
\pgfpathlineto{\pgfqpoint{0.582903in}{2.313221in}}%
\pgfpathlineto{\pgfqpoint{0.591705in}{2.308913in}}%
\pgfpathlineto{\pgfqpoint{0.598948in}{2.307790in}}%
\pgfpathlineto{\pgfqpoint{0.604150in}{2.302970in}}%
\pgfpathlineto{\pgfqpoint{0.620717in}{2.296209in}}%
\pgfpathlineto{\pgfqpoint{0.627218in}{2.298959in}}%
\pgfpathlineto{\pgfqpoint{0.630726in}{2.290276in}}%
\pgfpathlineto{\pgfqpoint{0.646614in}{2.290179in}}%
\pgfpathlineto{\pgfqpoint{0.648837in}{2.285967in}}%
\pgfpathlineto{\pgfqpoint{0.646729in}{2.278658in}}%
\pgfpathlineto{\pgfqpoint{0.647635in}{2.272305in}}%
\pgfpathclose%
\pgfusepath{fill}%
\end{pgfscope}%
\begin{pgfscope}%
\pgfpathrectangle{\pgfqpoint{0.100000in}{0.100000in}}{\pgfqpoint{3.608454in}{2.310000in}}%
\pgfusepath{clip}%
\pgfsetbuttcap%
\pgfsetmiterjoin%
\definecolor{currentfill}{rgb}{0.000000,0.639216,0.680392}%
\pgfsetfillcolor{currentfill}%
\pgfsetlinewidth{0.000000pt}%
\definecolor{currentstroke}{rgb}{0.000000,0.000000,0.000000}%
\pgfsetstrokecolor{currentstroke}%
\pgfsetstrokeopacity{0.000000}%
\pgfsetdash{}{0pt}%
\pgfpathmoveto{\pgfqpoint{2.726071in}{1.419220in}}%
\pgfpathlineto{\pgfqpoint{2.726448in}{1.415659in}}%
\pgfpathlineto{\pgfqpoint{2.697158in}{1.412713in}}%
\pgfpathlineto{\pgfqpoint{2.694747in}{1.433296in}}%
\pgfpathlineto{\pgfqpoint{2.699598in}{1.433892in}}%
\pgfpathlineto{\pgfqpoint{2.698943in}{1.451000in}}%
\pgfpathlineto{\pgfqpoint{2.722450in}{1.453627in}}%
\pgfpathlineto{\pgfqpoint{2.726071in}{1.419220in}}%
\pgfpathclose%
\pgfusepath{fill}%
\end{pgfscope}%
\begin{pgfscope}%
\pgfpathrectangle{\pgfqpoint{0.100000in}{0.100000in}}{\pgfqpoint{3.608454in}{2.310000in}}%
\pgfusepath{clip}%
\pgfsetbuttcap%
\pgfsetmiterjoin%
\definecolor{currentfill}{rgb}{0.000000,0.529412,0.735294}%
\pgfsetfillcolor{currentfill}%
\pgfsetlinewidth{0.000000pt}%
\definecolor{currentstroke}{rgb}{0.000000,0.000000,0.000000}%
\pgfsetstrokecolor{currentstroke}%
\pgfsetstrokeopacity{0.000000}%
\pgfsetdash{}{0pt}%
\pgfpathmoveto{\pgfqpoint{2.266038in}{1.833884in}}%
\pgfpathlineto{\pgfqpoint{2.232360in}{1.832947in}}%
\pgfpathlineto{\pgfqpoint{2.232161in}{1.839416in}}%
\pgfpathlineto{\pgfqpoint{2.237996in}{1.848119in}}%
\pgfpathlineto{\pgfqpoint{2.238078in}{1.851709in}}%
\pgfpathlineto{\pgfqpoint{2.230979in}{1.861346in}}%
\pgfpathlineto{\pgfqpoint{2.224641in}{1.861874in}}%
\pgfpathlineto{\pgfqpoint{2.226726in}{1.874225in}}%
\pgfpathlineto{\pgfqpoint{2.229731in}{1.876962in}}%
\pgfpathlineto{\pgfqpoint{2.233501in}{1.887214in}}%
\pgfpathlineto{\pgfqpoint{2.237467in}{1.890615in}}%
\pgfpathlineto{\pgfqpoint{2.253310in}{1.897735in}}%
\pgfpathlineto{\pgfqpoint{2.256290in}{1.902446in}}%
\pgfpathlineto{\pgfqpoint{2.256111in}{1.909051in}}%
\pgfpathlineto{\pgfqpoint{2.269571in}{1.909501in}}%
\pgfpathlineto{\pgfqpoint{2.270918in}{1.895690in}}%
\pgfpathlineto{\pgfqpoint{2.271935in}{1.868329in}}%
\pgfpathlineto{\pgfqpoint{2.265057in}{1.868094in}}%
\pgfpathlineto{\pgfqpoint{2.266038in}{1.833884in}}%
\pgfpathclose%
\pgfusepath{fill}%
\end{pgfscope}%
\begin{pgfscope}%
\pgfpathrectangle{\pgfqpoint{0.100000in}{0.100000in}}{\pgfqpoint{3.608454in}{2.310000in}}%
\pgfusepath{clip}%
\pgfsetbuttcap%
\pgfsetmiterjoin%
\definecolor{currentfill}{rgb}{0.000000,0.560784,0.719608}%
\pgfsetfillcolor{currentfill}%
\pgfsetlinewidth{0.000000pt}%
\definecolor{currentstroke}{rgb}{0.000000,0.000000,0.000000}%
\pgfsetstrokecolor{currentstroke}%
\pgfsetstrokeopacity{0.000000}%
\pgfsetdash{}{0pt}%
\pgfpathmoveto{\pgfqpoint{2.442684in}{0.864247in}}%
\pgfpathlineto{\pgfqpoint{2.438751in}{0.854506in}}%
\pgfpathlineto{\pgfqpoint{2.428148in}{0.846306in}}%
\pgfpathlineto{\pgfqpoint{2.412628in}{0.839279in}}%
\pgfpathlineto{\pgfqpoint{2.404935in}{0.832654in}}%
\pgfpathlineto{\pgfqpoint{2.393376in}{0.840961in}}%
\pgfpathlineto{\pgfqpoint{2.393315in}{0.855864in}}%
\pgfpathlineto{\pgfqpoint{2.389913in}{0.858059in}}%
\pgfpathlineto{\pgfqpoint{2.396470in}{0.865278in}}%
\pgfpathlineto{\pgfqpoint{2.395715in}{0.879087in}}%
\pgfpathlineto{\pgfqpoint{2.391193in}{0.892674in}}%
\pgfpathlineto{\pgfqpoint{2.401868in}{0.893227in}}%
\pgfpathlineto{\pgfqpoint{2.401625in}{0.897879in}}%
\pgfpathlineto{\pgfqpoint{2.408409in}{0.898240in}}%
\pgfpathlineto{\pgfqpoint{2.408616in}{0.893304in}}%
\pgfpathlineto{\pgfqpoint{2.416578in}{0.896534in}}%
\pgfpathlineto{\pgfqpoint{2.417347in}{0.883236in}}%
\pgfpathlineto{\pgfqpoint{2.410750in}{0.879777in}}%
\pgfpathlineto{\pgfqpoint{2.417323in}{0.873285in}}%
\pgfpathlineto{\pgfqpoint{2.442684in}{0.864247in}}%
\pgfpathclose%
\pgfusepath{fill}%
\end{pgfscope}%
\begin{pgfscope}%
\pgfpathrectangle{\pgfqpoint{0.100000in}{0.100000in}}{\pgfqpoint{3.608454in}{2.310000in}}%
\pgfusepath{clip}%
\pgfsetbuttcap%
\pgfsetmiterjoin%
\definecolor{currentfill}{rgb}{0.000000,0.741176,0.629412}%
\pgfsetfillcolor{currentfill}%
\pgfsetlinewidth{0.000000pt}%
\definecolor{currentstroke}{rgb}{0.000000,0.000000,0.000000}%
\pgfsetstrokecolor{currentstroke}%
\pgfsetstrokeopacity{0.000000}%
\pgfsetdash{}{0pt}%
\pgfpathmoveto{\pgfqpoint{2.172440in}{1.907196in}}%
\pgfpathlineto{\pgfqpoint{2.172623in}{1.893374in}}%
\pgfpathlineto{\pgfqpoint{2.174916in}{1.893366in}}%
\pgfpathlineto{\pgfqpoint{2.175363in}{1.880599in}}%
\pgfpathlineto{\pgfqpoint{2.142690in}{1.880249in}}%
\pgfpathlineto{\pgfqpoint{2.147223in}{1.876558in}}%
\pgfpathlineto{\pgfqpoint{2.099037in}{1.876260in}}%
\pgfpathlineto{\pgfqpoint{2.098775in}{1.916792in}}%
\pgfpathlineto{\pgfqpoint{2.097851in}{1.951504in}}%
\pgfpathlineto{\pgfqpoint{2.097902in}{1.958459in}}%
\pgfpathlineto{\pgfqpoint{2.118427in}{1.958403in}}%
\pgfpathlineto{\pgfqpoint{2.118937in}{1.925718in}}%
\pgfpathlineto{\pgfqpoint{2.132410in}{1.918482in}}%
\pgfpathlineto{\pgfqpoint{2.138352in}{1.920870in}}%
\pgfpathlineto{\pgfqpoint{2.142555in}{1.917955in}}%
\pgfpathlineto{\pgfqpoint{2.141306in}{1.906925in}}%
\pgfpathlineto{\pgfqpoint{2.172440in}{1.907196in}}%
\pgfpathclose%
\pgfusepath{fill}%
\end{pgfscope}%
\begin{pgfscope}%
\pgfpathrectangle{\pgfqpoint{0.100000in}{0.100000in}}{\pgfqpoint{3.608454in}{2.310000in}}%
\pgfusepath{clip}%
\pgfsetbuttcap%
\pgfsetmiterjoin%
\definecolor{currentfill}{rgb}{0.000000,0.745098,0.627451}%
\pgfsetfillcolor{currentfill}%
\pgfsetlinewidth{0.000000pt}%
\definecolor{currentstroke}{rgb}{0.000000,0.000000,0.000000}%
\pgfsetstrokecolor{currentstroke}%
\pgfsetstrokeopacity{0.000000}%
\pgfsetdash{}{0pt}%
\pgfpathmoveto{\pgfqpoint{2.231585in}{0.901788in}}%
\pgfpathlineto{\pgfqpoint{2.231529in}{0.896018in}}%
\pgfpathlineto{\pgfqpoint{2.240853in}{0.895980in}}%
\pgfpathlineto{\pgfqpoint{2.240807in}{0.867381in}}%
\pgfpathlineto{\pgfqpoint{2.186423in}{0.866575in}}%
\pgfpathlineto{\pgfqpoint{2.182676in}{0.872245in}}%
\pgfpathlineto{\pgfqpoint{2.183288in}{0.876547in}}%
\pgfpathlineto{\pgfqpoint{2.189397in}{0.878197in}}%
\pgfpathlineto{\pgfqpoint{2.195345in}{0.891554in}}%
\pgfpathlineto{\pgfqpoint{2.190887in}{0.898874in}}%
\pgfpathlineto{\pgfqpoint{2.191374in}{0.903396in}}%
\pgfpathlineto{\pgfqpoint{2.207261in}{0.903173in}}%
\pgfpathlineto{\pgfqpoint{2.207240in}{0.900881in}}%
\pgfpathlineto{\pgfqpoint{2.224654in}{0.900658in}}%
\pgfpathlineto{\pgfqpoint{2.231585in}{0.901788in}}%
\pgfpathclose%
\pgfusepath{fill}%
\end{pgfscope}%
\begin{pgfscope}%
\pgfpathrectangle{\pgfqpoint{0.100000in}{0.100000in}}{\pgfqpoint{3.608454in}{2.310000in}}%
\pgfusepath{clip}%
\pgfsetbuttcap%
\pgfsetmiterjoin%
\definecolor{currentfill}{rgb}{0.000000,0.282353,0.858824}%
\pgfsetfillcolor{currentfill}%
\pgfsetlinewidth{0.000000pt}%
\definecolor{currentstroke}{rgb}{0.000000,0.000000,0.000000}%
\pgfsetstrokecolor{currentstroke}%
\pgfsetstrokeopacity{0.000000}%
\pgfsetdash{}{0pt}%
\pgfpathmoveto{\pgfqpoint{2.454210in}{1.655405in}}%
\pgfpathlineto{\pgfqpoint{2.433615in}{1.653903in}}%
\pgfpathlineto{\pgfqpoint{2.406360in}{1.653264in}}%
\pgfpathlineto{\pgfqpoint{2.406543in}{1.649815in}}%
\pgfpathlineto{\pgfqpoint{2.372395in}{1.647911in}}%
\pgfpathlineto{\pgfqpoint{2.370685in}{1.678611in}}%
\pgfpathlineto{\pgfqpoint{2.363976in}{1.679112in}}%
\pgfpathlineto{\pgfqpoint{2.357123in}{1.675633in}}%
\pgfpathlineto{\pgfqpoint{2.356072in}{1.695517in}}%
\pgfpathlineto{\pgfqpoint{2.355403in}{1.705801in}}%
\pgfpathlineto{\pgfqpoint{2.376100in}{1.706906in}}%
\pgfpathlineto{\pgfqpoint{2.375661in}{1.713794in}}%
\pgfpathlineto{\pgfqpoint{2.405834in}{1.715444in}}%
\pgfpathlineto{\pgfqpoint{2.416473in}{1.716189in}}%
\pgfpathlineto{\pgfqpoint{2.414986in}{1.743088in}}%
\pgfpathlineto{\pgfqpoint{2.439466in}{1.744650in}}%
\pgfpathlineto{\pgfqpoint{2.440017in}{1.736230in}}%
\pgfpathlineto{\pgfqpoint{2.436156in}{1.727226in}}%
\pgfpathlineto{\pgfqpoint{2.436777in}{1.717456in}}%
\pgfpathlineto{\pgfqpoint{2.450492in}{1.717550in}}%
\pgfpathlineto{\pgfqpoint{2.454210in}{1.655405in}}%
\pgfpathclose%
\pgfusepath{fill}%
\end{pgfscope}%
\begin{pgfscope}%
\pgfpathrectangle{\pgfqpoint{0.100000in}{0.100000in}}{\pgfqpoint{3.608454in}{2.310000in}}%
\pgfusepath{clip}%
\pgfsetbuttcap%
\pgfsetmiterjoin%
\definecolor{currentfill}{rgb}{0.000000,0.388235,0.805882}%
\pgfsetfillcolor{currentfill}%
\pgfsetlinewidth{0.000000pt}%
\definecolor{currentstroke}{rgb}{0.000000,0.000000,0.000000}%
\pgfsetstrokecolor{currentstroke}%
\pgfsetstrokeopacity{0.000000}%
\pgfsetdash{}{0pt}%
\pgfpathmoveto{\pgfqpoint{1.278870in}{1.882470in}}%
\pgfpathlineto{\pgfqpoint{1.317490in}{1.876575in}}%
\pgfpathlineto{\pgfqpoint{1.318547in}{1.875904in}}%
\pgfpathlineto{\pgfqpoint{1.364865in}{1.868873in}}%
\pgfpathlineto{\pgfqpoint{1.383661in}{1.866297in}}%
\pgfpathlineto{\pgfqpoint{1.384830in}{1.855882in}}%
\pgfpathlineto{\pgfqpoint{1.389944in}{1.843550in}}%
\pgfpathlineto{\pgfqpoint{1.396824in}{1.840362in}}%
\pgfpathlineto{\pgfqpoint{1.406562in}{1.831802in}}%
\pgfpathlineto{\pgfqpoint{1.412375in}{1.821775in}}%
\pgfpathlineto{\pgfqpoint{1.417230in}{1.817131in}}%
\pgfpathlineto{\pgfqpoint{1.419171in}{1.807060in}}%
\pgfpathlineto{\pgfqpoint{1.417395in}{1.794618in}}%
\pgfpathlineto{\pgfqpoint{1.338318in}{1.806097in}}%
\pgfpathlineto{\pgfqpoint{1.337257in}{1.799111in}}%
\pgfpathlineto{\pgfqpoint{1.323648in}{1.801191in}}%
\pgfpathlineto{\pgfqpoint{1.322614in}{1.794241in}}%
\pgfpathlineto{\pgfqpoint{1.315604in}{1.795301in}}%
\pgfpathlineto{\pgfqpoint{1.314743in}{1.788523in}}%
\pgfpathlineto{\pgfqpoint{1.304569in}{1.790081in}}%
\pgfpathlineto{\pgfqpoint{1.303529in}{1.783157in}}%
\pgfpathlineto{\pgfqpoint{1.291011in}{1.784895in}}%
\pgfpathlineto{\pgfqpoint{1.285167in}{1.795237in}}%
\pgfpathlineto{\pgfqpoint{1.279919in}{1.798565in}}%
\pgfpathlineto{\pgfqpoint{1.273821in}{1.796453in}}%
\pgfpathlineto{\pgfqpoint{1.271394in}{1.792453in}}%
\pgfpathlineto{\pgfqpoint{1.264050in}{1.788274in}}%
\pgfpathlineto{\pgfqpoint{1.260477in}{1.800714in}}%
\pgfpathlineto{\pgfqpoint{1.254530in}{1.801740in}}%
\pgfpathlineto{\pgfqpoint{1.251746in}{1.806946in}}%
\pgfpathlineto{\pgfqpoint{1.253366in}{1.816700in}}%
\pgfpathlineto{\pgfqpoint{1.250644in}{1.821729in}}%
\pgfpathlineto{\pgfqpoint{1.249318in}{1.830814in}}%
\pgfpathlineto{\pgfqpoint{1.243670in}{1.842268in}}%
\pgfpathlineto{\pgfqpoint{1.245880in}{1.848777in}}%
\pgfpathlineto{\pgfqpoint{1.241305in}{1.854999in}}%
\pgfpathlineto{\pgfqpoint{1.224995in}{1.857794in}}%
\pgfpathlineto{\pgfqpoint{1.226134in}{1.864328in}}%
\pgfpathlineto{\pgfqpoint{1.204580in}{1.868127in}}%
\pgfpathlineto{\pgfqpoint{1.209266in}{1.894331in}}%
\pgfpathlineto{\pgfqpoint{1.228507in}{1.890208in}}%
\pgfpathlineto{\pgfqpoint{1.278870in}{1.882470in}}%
\pgfpathclose%
\pgfusepath{fill}%
\end{pgfscope}%
\begin{pgfscope}%
\pgfpathrectangle{\pgfqpoint{0.100000in}{0.100000in}}{\pgfqpoint{3.608454in}{2.310000in}}%
\pgfusepath{clip}%
\pgfsetbuttcap%
\pgfsetmiterjoin%
\definecolor{currentfill}{rgb}{0.000000,0.470588,0.764706}%
\pgfsetfillcolor{currentfill}%
\pgfsetlinewidth{0.000000pt}%
\definecolor{currentstroke}{rgb}{0.000000,0.000000,0.000000}%
\pgfsetstrokecolor{currentstroke}%
\pgfsetstrokeopacity{0.000000}%
\pgfsetdash{}{0pt}%
\pgfpathmoveto{\pgfqpoint{2.551959in}{0.940228in}}%
\pgfpathlineto{\pgfqpoint{2.546945in}{0.939900in}}%
\pgfpathlineto{\pgfqpoint{2.544297in}{0.934956in}}%
\pgfpathlineto{\pgfqpoint{2.535053in}{0.931443in}}%
\pgfpathlineto{\pgfqpoint{2.521327in}{0.932876in}}%
\pgfpathlineto{\pgfqpoint{2.520583in}{0.943229in}}%
\pgfpathlineto{\pgfqpoint{2.506684in}{0.942391in}}%
\pgfpathlineto{\pgfqpoint{2.506964in}{0.937791in}}%
\pgfpathlineto{\pgfqpoint{2.489985in}{0.934122in}}%
\pgfpathlineto{\pgfqpoint{2.469213in}{0.932802in}}%
\pgfpathlineto{\pgfqpoint{2.468504in}{0.944341in}}%
\pgfpathlineto{\pgfqpoint{2.466935in}{0.967664in}}%
\pgfpathlineto{\pgfqpoint{2.484251in}{0.968666in}}%
\pgfpathlineto{\pgfqpoint{2.483100in}{0.985946in}}%
\pgfpathlineto{\pgfqpoint{2.503832in}{0.987282in}}%
\pgfpathlineto{\pgfqpoint{2.512839in}{0.989031in}}%
\pgfpathlineto{\pgfqpoint{2.512313in}{0.997036in}}%
\pgfpathlineto{\pgfqpoint{2.515690in}{0.998385in}}%
\pgfpathlineto{\pgfqpoint{2.528379in}{0.999297in}}%
\pgfpathlineto{\pgfqpoint{2.528607in}{0.995823in}}%
\pgfpathlineto{\pgfqpoint{2.553769in}{0.997504in}}%
\pgfpathlineto{\pgfqpoint{2.553484in}{0.986177in}}%
\pgfpathlineto{\pgfqpoint{2.551959in}{0.940228in}}%
\pgfpathclose%
\pgfusepath{fill}%
\end{pgfscope}%
\begin{pgfscope}%
\pgfpathrectangle{\pgfqpoint{0.100000in}{0.100000in}}{\pgfqpoint{3.608454in}{2.310000in}}%
\pgfusepath{clip}%
\pgfsetbuttcap%
\pgfsetmiterjoin%
\definecolor{currentfill}{rgb}{0.000000,0.670588,0.664706}%
\pgfsetfillcolor{currentfill}%
\pgfsetlinewidth{0.000000pt}%
\definecolor{currentstroke}{rgb}{0.000000,0.000000,0.000000}%
\pgfsetstrokecolor{currentstroke}%
\pgfsetstrokeopacity{0.000000}%
\pgfsetdash{}{0pt}%
\pgfpathmoveto{\pgfqpoint{3.458788in}{2.127000in}}%
\pgfpathlineto{\pgfqpoint{3.457596in}{2.136892in}}%
\pgfpathlineto{\pgfqpoint{3.481080in}{2.206952in}}%
\pgfpathlineto{\pgfqpoint{3.491936in}{2.206386in}}%
\pgfpathlineto{\pgfqpoint{3.494641in}{2.194072in}}%
\pgfpathlineto{\pgfqpoint{3.504213in}{2.190328in}}%
\pgfpathlineto{\pgfqpoint{3.518463in}{2.203609in}}%
\pgfpathlineto{\pgfqpoint{3.528567in}{2.206604in}}%
\pgfpathlineto{\pgfqpoint{3.527922in}{2.211919in}}%
\pgfpathlineto{\pgfqpoint{3.535054in}{2.214135in}}%
\pgfpathlineto{\pgfqpoint{3.552913in}{2.206672in}}%
\pgfpathlineto{\pgfqpoint{3.558787in}{2.200716in}}%
\pgfpathlineto{\pgfqpoint{3.564560in}{2.199223in}}%
\pgfpathlineto{\pgfqpoint{3.575635in}{2.163717in}}%
\pgfpathlineto{\pgfqpoint{3.591481in}{2.113768in}}%
\pgfpathlineto{\pgfqpoint{3.591685in}{2.108859in}}%
\pgfpathlineto{\pgfqpoint{3.596315in}{2.093414in}}%
\pgfpathlineto{\pgfqpoint{3.584395in}{2.086358in}}%
\pgfpathlineto{\pgfqpoint{3.566650in}{2.075414in}}%
\pgfpathlineto{\pgfqpoint{3.565274in}{2.075679in}}%
\pgfpathlineto{\pgfqpoint{3.557203in}{2.103454in}}%
\pgfpathlineto{\pgfqpoint{3.546699in}{2.136666in}}%
\pgfpathlineto{\pgfqpoint{3.526210in}{2.131697in}}%
\pgfpathlineto{\pgfqpoint{3.522068in}{2.145170in}}%
\pgfpathlineto{\pgfqpoint{3.474712in}{2.131510in}}%
\pgfpathlineto{\pgfqpoint{3.458788in}{2.127000in}}%
\pgfpathclose%
\pgfusepath{fill}%
\end{pgfscope}%
\begin{pgfscope}%
\pgfpathrectangle{\pgfqpoint{0.100000in}{0.100000in}}{\pgfqpoint{3.608454in}{2.310000in}}%
\pgfusepath{clip}%
\pgfsetbuttcap%
\pgfsetmiterjoin%
\definecolor{currentfill}{rgb}{0.000000,0.482353,0.758824}%
\pgfsetfillcolor{currentfill}%
\pgfsetlinewidth{0.000000pt}%
\definecolor{currentstroke}{rgb}{0.000000,0.000000,0.000000}%
\pgfsetstrokecolor{currentstroke}%
\pgfsetstrokeopacity{0.000000}%
\pgfsetdash{}{0pt}%
\pgfpathmoveto{\pgfqpoint{2.714189in}{1.531011in}}%
\pgfpathlineto{\pgfqpoint{2.707766in}{1.586556in}}%
\pgfpathlineto{\pgfqpoint{2.731544in}{1.590059in}}%
\pgfpathlineto{\pgfqpoint{2.734410in}{1.575123in}}%
\pgfpathlineto{\pgfqpoint{2.737491in}{1.568598in}}%
\pgfpathlineto{\pgfqpoint{2.744178in}{1.569410in}}%
\pgfpathlineto{\pgfqpoint{2.746637in}{1.548837in}}%
\pgfpathlineto{\pgfqpoint{2.739953in}{1.548013in}}%
\pgfpathlineto{\pgfqpoint{2.741540in}{1.534284in}}%
\pgfpathlineto{\pgfqpoint{2.714189in}{1.531011in}}%
\pgfpathclose%
\pgfusepath{fill}%
\end{pgfscope}%
\begin{pgfscope}%
\pgfpathrectangle{\pgfqpoint{0.100000in}{0.100000in}}{\pgfqpoint{3.608454in}{2.310000in}}%
\pgfusepath{clip}%
\pgfsetbuttcap%
\pgfsetmiterjoin%
\definecolor{currentfill}{rgb}{0.000000,0.650980,0.674510}%
\pgfsetfillcolor{currentfill}%
\pgfsetlinewidth{0.000000pt}%
\definecolor{currentstroke}{rgb}{0.000000,0.000000,0.000000}%
\pgfsetstrokecolor{currentstroke}%
\pgfsetstrokeopacity{0.000000}%
\pgfsetdash{}{0pt}%
\pgfpathmoveto{\pgfqpoint{1.303419in}{1.175695in}}%
\pgfpathlineto{\pgfqpoint{1.300870in}{1.158512in}}%
\pgfpathlineto{\pgfqpoint{1.210980in}{1.171960in}}%
\pgfpathlineto{\pgfqpoint{1.223167in}{1.249876in}}%
\pgfpathlineto{\pgfqpoint{1.264762in}{1.243586in}}%
\pgfpathlineto{\pgfqpoint{1.268205in}{1.247503in}}%
\pgfpathlineto{\pgfqpoint{1.272888in}{1.260292in}}%
\pgfpathlineto{\pgfqpoint{1.280183in}{1.269745in}}%
\pgfpathlineto{\pgfqpoint{1.286167in}{1.271003in}}%
\pgfpathlineto{\pgfqpoint{1.291651in}{1.276850in}}%
\pgfpathlineto{\pgfqpoint{1.293840in}{1.286121in}}%
\pgfpathlineto{\pgfqpoint{1.303897in}{1.295339in}}%
\pgfpathlineto{\pgfqpoint{1.305818in}{1.299625in}}%
\pgfpathlineto{\pgfqpoint{1.314836in}{1.308720in}}%
\pgfpathlineto{\pgfqpoint{1.325950in}{1.311919in}}%
\pgfpathlineto{\pgfqpoint{1.327961in}{1.309954in}}%
\pgfpathlineto{\pgfqpoint{1.329166in}{1.295757in}}%
\pgfpathlineto{\pgfqpoint{1.325433in}{1.268700in}}%
\pgfpathlineto{\pgfqpoint{1.347428in}{1.265725in}}%
\pgfpathlineto{\pgfqpoint{1.347103in}{1.263333in}}%
\pgfpathlineto{\pgfqpoint{1.375196in}{1.259964in}}%
\pgfpathlineto{\pgfqpoint{1.373495in}{1.246788in}}%
\pgfpathlineto{\pgfqpoint{1.383780in}{1.226748in}}%
\pgfpathlineto{\pgfqpoint{1.324687in}{1.235061in}}%
\pgfpathlineto{\pgfqpoint{1.322223in}{1.228643in}}%
\pgfpathlineto{\pgfqpoint{1.310197in}{1.221478in}}%
\pgfpathlineto{\pgfqpoint{1.303419in}{1.175695in}}%
\pgfpathclose%
\pgfusepath{fill}%
\end{pgfscope}%
\begin{pgfscope}%
\pgfpathrectangle{\pgfqpoint{0.100000in}{0.100000in}}{\pgfqpoint{3.608454in}{2.310000in}}%
\pgfusepath{clip}%
\pgfsetbuttcap%
\pgfsetmiterjoin%
\definecolor{currentfill}{rgb}{0.000000,0.760784,0.619608}%
\pgfsetfillcolor{currentfill}%
\pgfsetlinewidth{0.000000pt}%
\definecolor{currentstroke}{rgb}{0.000000,0.000000,0.000000}%
\pgfsetstrokecolor{currentstroke}%
\pgfsetstrokeopacity{0.000000}%
\pgfsetdash{}{0pt}%
\pgfpathmoveto{\pgfqpoint{2.420115in}{2.049767in}}%
\pgfpathlineto{\pgfqpoint{2.414517in}{2.052884in}}%
\pgfpathlineto{\pgfqpoint{2.417513in}{2.057828in}}%
\pgfpathlineto{\pgfqpoint{2.438907in}{2.070311in}}%
\pgfpathlineto{\pgfqpoint{2.446529in}{2.077612in}}%
\pgfpathlineto{\pgfqpoint{2.455726in}{2.081379in}}%
\pgfpathlineto{\pgfqpoint{2.450327in}{2.070558in}}%
\pgfpathlineto{\pgfqpoint{2.444429in}{2.065660in}}%
\pgfpathlineto{\pgfqpoint{2.430864in}{2.058882in}}%
\pgfpathlineto{\pgfqpoint{2.433030in}{2.055826in}}%
\pgfpathlineto{\pgfqpoint{2.420115in}{2.049767in}}%
\pgfpathclose%
\pgfusepath{fill}%
\end{pgfscope}%
\begin{pgfscope}%
\pgfpathrectangle{\pgfqpoint{0.100000in}{0.100000in}}{\pgfqpoint{3.608454in}{2.310000in}}%
\pgfusepath{clip}%
\pgfsetbuttcap%
\pgfsetmiterjoin%
\definecolor{currentfill}{rgb}{0.000000,0.760784,0.619608}%
\pgfsetfillcolor{currentfill}%
\pgfsetlinewidth{0.000000pt}%
\definecolor{currentstroke}{rgb}{0.000000,0.000000,0.000000}%
\pgfsetstrokecolor{currentstroke}%
\pgfsetstrokeopacity{0.000000}%
\pgfsetdash{}{0pt}%
\pgfpathmoveto{\pgfqpoint{2.485656in}{1.981514in}}%
\pgfpathlineto{\pgfqpoint{2.486521in}{1.969842in}}%
\pgfpathlineto{\pgfqpoint{2.489953in}{1.963218in}}%
\pgfpathlineto{\pgfqpoint{2.483177in}{1.962621in}}%
\pgfpathlineto{\pgfqpoint{2.484672in}{1.942048in}}%
\pgfpathlineto{\pgfqpoint{2.436761in}{1.938649in}}%
\pgfpathlineto{\pgfqpoint{2.435480in}{1.959378in}}%
\pgfpathlineto{\pgfqpoint{2.442299in}{1.959772in}}%
\pgfpathlineto{\pgfqpoint{2.438144in}{1.966329in}}%
\pgfpathlineto{\pgfqpoint{2.436662in}{1.987686in}}%
\pgfpathlineto{\pgfqpoint{2.438695in}{1.993269in}}%
\pgfpathlineto{\pgfqpoint{2.450631in}{2.004021in}}%
\pgfpathlineto{\pgfqpoint{2.454724in}{2.005602in}}%
\pgfpathlineto{\pgfqpoint{2.462516in}{2.016619in}}%
\pgfpathlineto{\pgfqpoint{2.472955in}{2.023672in}}%
\pgfpathlineto{\pgfqpoint{2.487300in}{2.027022in}}%
\pgfpathlineto{\pgfqpoint{2.495202in}{2.026998in}}%
\pgfpathlineto{\pgfqpoint{2.500592in}{2.021625in}}%
\pgfpathlineto{\pgfqpoint{2.489998in}{2.020382in}}%
\pgfpathlineto{\pgfqpoint{2.488577in}{2.015517in}}%
\pgfpathlineto{\pgfqpoint{2.482568in}{2.011992in}}%
\pgfpathlineto{\pgfqpoint{2.473948in}{2.003663in}}%
\pgfpathlineto{\pgfqpoint{2.473922in}{1.999171in}}%
\pgfpathlineto{\pgfqpoint{2.463650in}{1.984401in}}%
\pgfpathlineto{\pgfqpoint{2.462316in}{1.973885in}}%
\pgfpathlineto{\pgfqpoint{2.465262in}{1.970695in}}%
\pgfpathlineto{\pgfqpoint{2.474637in}{1.982064in}}%
\pgfpathlineto{\pgfqpoint{2.479239in}{1.985245in}}%
\pgfpathlineto{\pgfqpoint{2.485656in}{1.981514in}}%
\pgfpathclose%
\pgfusepath{fill}%
\end{pgfscope}%
\begin{pgfscope}%
\pgfpathrectangle{\pgfqpoint{0.100000in}{0.100000in}}{\pgfqpoint{3.608454in}{2.310000in}}%
\pgfusepath{clip}%
\pgfsetbuttcap%
\pgfsetmiterjoin%
\definecolor{currentfill}{rgb}{0.000000,0.372549,0.813725}%
\pgfsetfillcolor{currentfill}%
\pgfsetlinewidth{0.000000pt}%
\definecolor{currentstroke}{rgb}{0.000000,0.000000,0.000000}%
\pgfsetstrokecolor{currentstroke}%
\pgfsetstrokeopacity{0.000000}%
\pgfsetdash{}{0pt}%
\pgfpathmoveto{\pgfqpoint{1.481289in}{1.288540in}}%
\pgfpathlineto{\pgfqpoint{1.474843in}{1.297491in}}%
\pgfpathlineto{\pgfqpoint{1.470019in}{1.293572in}}%
\pgfpathlineto{\pgfqpoint{1.458524in}{1.289130in}}%
\pgfpathlineto{\pgfqpoint{1.454881in}{1.289984in}}%
\pgfpathlineto{\pgfqpoint{1.448675in}{1.296375in}}%
\pgfpathlineto{\pgfqpoint{1.444613in}{1.311073in}}%
\pgfpathlineto{\pgfqpoint{1.426840in}{1.337144in}}%
\pgfpathlineto{\pgfqpoint{1.435665in}{1.345243in}}%
\pgfpathlineto{\pgfqpoint{1.436154in}{1.349875in}}%
\pgfpathlineto{\pgfqpoint{1.431712in}{1.356218in}}%
\pgfpathlineto{\pgfqpoint{1.470992in}{1.351885in}}%
\pgfpathlineto{\pgfqpoint{1.474789in}{1.385857in}}%
\pgfpathlineto{\pgfqpoint{1.552744in}{1.377402in}}%
\pgfpathlineto{\pgfqpoint{1.550467in}{1.357020in}}%
\pgfpathlineto{\pgfqpoint{1.547834in}{1.329763in}}%
\pgfpathlineto{\pgfqpoint{1.493221in}{1.335267in}}%
\pgfpathlineto{\pgfqpoint{1.491038in}{1.314728in}}%
\pgfpathlineto{\pgfqpoint{1.484320in}{1.315446in}}%
\pgfpathlineto{\pgfqpoint{1.481289in}{1.288540in}}%
\pgfpathclose%
\pgfusepath{fill}%
\end{pgfscope}%
\begin{pgfscope}%
\pgfpathrectangle{\pgfqpoint{0.100000in}{0.100000in}}{\pgfqpoint{3.608454in}{2.310000in}}%
\pgfusepath{clip}%
\pgfsetbuttcap%
\pgfsetmiterjoin%
\definecolor{currentfill}{rgb}{0.000000,0.717647,0.641176}%
\pgfsetfillcolor{currentfill}%
\pgfsetlinewidth{0.000000pt}%
\definecolor{currentstroke}{rgb}{0.000000,0.000000,0.000000}%
\pgfsetstrokecolor{currentstroke}%
\pgfsetstrokeopacity{0.000000}%
\pgfsetdash{}{0pt}%
\pgfpathmoveto{\pgfqpoint{2.977955in}{1.500863in}}%
\pgfpathlineto{\pgfqpoint{2.956251in}{1.497406in}}%
\pgfpathlineto{\pgfqpoint{2.954062in}{1.518144in}}%
\pgfpathlineto{\pgfqpoint{2.949522in}{1.517575in}}%
\pgfpathlineto{\pgfqpoint{2.949140in}{1.528048in}}%
\pgfpathlineto{\pgfqpoint{2.952109in}{1.531961in}}%
\pgfpathlineto{\pgfqpoint{2.963825in}{1.532458in}}%
\pgfpathlineto{\pgfqpoint{2.971951in}{1.538265in}}%
\pgfpathlineto{\pgfqpoint{2.977955in}{1.500863in}}%
\pgfpathclose%
\pgfusepath{fill}%
\end{pgfscope}%
\begin{pgfscope}%
\pgfpathrectangle{\pgfqpoint{0.100000in}{0.100000in}}{\pgfqpoint{3.608454in}{2.310000in}}%
\pgfusepath{clip}%
\pgfsetbuttcap%
\pgfsetmiterjoin%
\definecolor{currentfill}{rgb}{0.000000,0.760784,0.619608}%
\pgfsetfillcolor{currentfill}%
\pgfsetlinewidth{0.000000pt}%
\definecolor{currentstroke}{rgb}{0.000000,0.000000,0.000000}%
\pgfsetstrokecolor{currentstroke}%
\pgfsetstrokeopacity{0.000000}%
\pgfsetdash{}{0pt}%
\pgfpathmoveto{\pgfqpoint{2.573691in}{1.108364in}}%
\pgfpathlineto{\pgfqpoint{2.573068in}{1.094430in}}%
\pgfpathlineto{\pgfqpoint{2.579650in}{1.089158in}}%
\pgfpathlineto{\pgfqpoint{2.575892in}{1.080309in}}%
\pgfpathlineto{\pgfqpoint{2.559562in}{1.077181in}}%
\pgfpathlineto{\pgfqpoint{2.556697in}{1.071568in}}%
\pgfpathlineto{\pgfqpoint{2.550818in}{1.070538in}}%
\pgfpathlineto{\pgfqpoint{2.534914in}{1.072050in}}%
\pgfpathlineto{\pgfqpoint{2.531082in}{1.076803in}}%
\pgfpathlineto{\pgfqpoint{2.525076in}{1.078728in}}%
\pgfpathlineto{\pgfqpoint{2.517692in}{1.084251in}}%
\pgfpathlineto{\pgfqpoint{2.517006in}{1.100169in}}%
\pgfpathlineto{\pgfqpoint{2.519200in}{1.103140in}}%
\pgfpathlineto{\pgfqpoint{2.544235in}{1.104329in}}%
\pgfpathlineto{\pgfqpoint{2.552190in}{1.107568in}}%
\pgfpathlineto{\pgfqpoint{2.553934in}{1.104609in}}%
\pgfpathlineto{\pgfqpoint{2.561312in}{1.107350in}}%
\pgfpathlineto{\pgfqpoint{2.573691in}{1.108364in}}%
\pgfpathclose%
\pgfusepath{fill}%
\end{pgfscope}%
\begin{pgfscope}%
\pgfpathrectangle{\pgfqpoint{0.100000in}{0.100000in}}{\pgfqpoint{3.608454in}{2.310000in}}%
\pgfusepath{clip}%
\pgfsetbuttcap%
\pgfsetmiterjoin%
\definecolor{currentfill}{rgb}{0.000000,0.662745,0.668627}%
\pgfsetfillcolor{currentfill}%
\pgfsetlinewidth{0.000000pt}%
\definecolor{currentstroke}{rgb}{0.000000,0.000000,0.000000}%
\pgfsetstrokecolor{currentstroke}%
\pgfsetstrokeopacity{0.000000}%
\pgfsetdash{}{0pt}%
\pgfpathmoveto{\pgfqpoint{1.518670in}{0.813225in}}%
\pgfpathlineto{\pgfqpoint{1.498395in}{0.815195in}}%
\pgfpathlineto{\pgfqpoint{1.443108in}{0.820998in}}%
\pgfpathlineto{\pgfqpoint{1.447334in}{0.862167in}}%
\pgfpathlineto{\pgfqpoint{1.413830in}{0.865768in}}%
\pgfpathlineto{\pgfqpoint{1.417927in}{0.900770in}}%
\pgfpathlineto{\pgfqpoint{1.420339in}{0.900498in}}%
\pgfpathlineto{\pgfqpoint{1.421822in}{0.913943in}}%
\pgfpathlineto{\pgfqpoint{1.449063in}{0.911411in}}%
\pgfpathlineto{\pgfqpoint{1.452184in}{0.931773in}}%
\pgfpathlineto{\pgfqpoint{1.458090in}{0.986273in}}%
\pgfpathlineto{\pgfqpoint{1.465068in}{0.985473in}}%
\pgfpathlineto{\pgfqpoint{1.464338in}{0.978551in}}%
\pgfpathlineto{\pgfqpoint{1.505497in}{0.973876in}}%
\pgfpathlineto{\pgfqpoint{1.506226in}{0.980812in}}%
\pgfpathlineto{\pgfqpoint{1.519945in}{0.979398in}}%
\pgfpathlineto{\pgfqpoint{1.526759in}{0.978642in}}%
\pgfpathlineto{\pgfqpoint{1.524717in}{0.957918in}}%
\pgfpathlineto{\pgfqpoint{1.532929in}{0.957126in}}%
\pgfpathlineto{\pgfqpoint{1.531592in}{0.944243in}}%
\pgfpathlineto{\pgfqpoint{1.545317in}{0.943021in}}%
\pgfpathlineto{\pgfqpoint{1.544636in}{0.936153in}}%
\pgfpathlineto{\pgfqpoint{1.530771in}{0.937426in}}%
\pgfpathlineto{\pgfqpoint{1.529406in}{0.923693in}}%
\pgfpathlineto{\pgfqpoint{1.526429in}{0.923948in}}%
\pgfpathlineto{\pgfqpoint{1.523143in}{0.889921in}}%
\pgfpathlineto{\pgfqpoint{1.519940in}{0.890229in}}%
\pgfpathlineto{\pgfqpoint{1.516573in}{0.855187in}}%
\pgfpathlineto{\pgfqpoint{1.522650in}{0.854609in}}%
\pgfpathlineto{\pgfqpoint{1.518670in}{0.813225in}}%
\pgfpathclose%
\pgfusepath{fill}%
\end{pgfscope}%
\begin{pgfscope}%
\pgfpathrectangle{\pgfqpoint{0.100000in}{0.100000in}}{\pgfqpoint{3.608454in}{2.310000in}}%
\pgfusepath{clip}%
\pgfsetbuttcap%
\pgfsetmiterjoin%
\definecolor{currentfill}{rgb}{0.000000,0.654902,0.672549}%
\pgfsetfillcolor{currentfill}%
\pgfsetlinewidth{0.000000pt}%
\definecolor{currentstroke}{rgb}{0.000000,0.000000,0.000000}%
\pgfsetstrokecolor{currentstroke}%
\pgfsetstrokeopacity{0.000000}%
\pgfsetdash{}{0pt}%
\pgfpathmoveto{\pgfqpoint{1.411936in}{2.184794in}}%
\pgfpathlineto{\pgfqpoint{1.407741in}{2.164527in}}%
\pgfpathlineto{\pgfqpoint{1.404613in}{2.164996in}}%
\pgfpathlineto{\pgfqpoint{1.402592in}{2.151198in}}%
\pgfpathlineto{\pgfqpoint{1.399456in}{2.139818in}}%
\pgfpathlineto{\pgfqpoint{1.396285in}{2.142117in}}%
\pgfpathlineto{\pgfqpoint{1.392437in}{2.124428in}}%
\pgfpathlineto{\pgfqpoint{1.389566in}{2.105191in}}%
\pgfpathlineto{\pgfqpoint{1.380515in}{2.107723in}}%
\pgfpathlineto{\pgfqpoint{1.378739in}{2.102305in}}%
\pgfpathlineto{\pgfqpoint{1.364847in}{2.104580in}}%
\pgfpathlineto{\pgfqpoint{1.362402in}{2.089710in}}%
\pgfpathlineto{\pgfqpoint{1.358594in}{2.095823in}}%
\pgfpathlineto{\pgfqpoint{1.350200in}{2.095416in}}%
\pgfpathlineto{\pgfqpoint{1.342674in}{2.097531in}}%
\pgfpathlineto{\pgfqpoint{1.332093in}{2.093598in}}%
\pgfpathlineto{\pgfqpoint{1.328054in}{2.095490in}}%
\pgfpathlineto{\pgfqpoint{1.333325in}{2.126759in}}%
\pgfpathlineto{\pgfqpoint{1.323124in}{2.128216in}}%
\pgfpathlineto{\pgfqpoint{1.324278in}{2.135234in}}%
\pgfpathlineto{\pgfqpoint{1.318309in}{2.136230in}}%
\pgfpathlineto{\pgfqpoint{1.319123in}{2.142944in}}%
\pgfpathlineto{\pgfqpoint{1.278167in}{2.149917in}}%
\pgfpathlineto{\pgfqpoint{1.276965in}{2.143064in}}%
\pgfpathlineto{\pgfqpoint{1.266372in}{2.144919in}}%
\pgfpathlineto{\pgfqpoint{1.265169in}{2.138097in}}%
\pgfpathlineto{\pgfqpoint{1.234540in}{2.143503in}}%
\pgfpathlineto{\pgfqpoint{1.235802in}{2.150383in}}%
\pgfpathlineto{\pgfqpoint{1.242847in}{2.149099in}}%
\pgfpathlineto{\pgfqpoint{1.254169in}{2.210455in}}%
\pgfpathlineto{\pgfqpoint{1.308831in}{2.201031in}}%
\pgfpathlineto{\pgfqpoint{1.372332in}{2.190734in}}%
\pgfpathlineto{\pgfqpoint{1.411936in}{2.184794in}}%
\pgfpathclose%
\pgfusepath{fill}%
\end{pgfscope}%
\begin{pgfscope}%
\pgfpathrectangle{\pgfqpoint{0.100000in}{0.100000in}}{\pgfqpoint{3.608454in}{2.310000in}}%
\pgfusepath{clip}%
\pgfsetbuttcap%
\pgfsetmiterjoin%
\definecolor{currentfill}{rgb}{0.000000,0.541176,0.729412}%
\pgfsetfillcolor{currentfill}%
\pgfsetlinewidth{0.000000pt}%
\definecolor{currentstroke}{rgb}{0.000000,0.000000,0.000000}%
\pgfsetstrokecolor{currentstroke}%
\pgfsetstrokeopacity{0.000000}%
\pgfsetdash{}{0pt}%
\pgfpathmoveto{\pgfqpoint{3.286885in}{1.169498in}}%
\pgfpathlineto{\pgfqpoint{3.287098in}{1.161529in}}%
\pgfpathlineto{\pgfqpoint{3.273900in}{1.157290in}}%
\pgfpathlineto{\pgfqpoint{3.265343in}{1.162433in}}%
\pgfpathlineto{\pgfqpoint{3.248882in}{1.168666in}}%
\pgfpathlineto{\pgfqpoint{3.248243in}{1.176631in}}%
\pgfpathlineto{\pgfqpoint{3.250804in}{1.179929in}}%
\pgfpathlineto{\pgfqpoint{3.246587in}{1.185703in}}%
\pgfpathlineto{\pgfqpoint{3.244554in}{1.193790in}}%
\pgfpathlineto{\pgfqpoint{3.232624in}{1.197760in}}%
\pgfpathlineto{\pgfqpoint{3.231820in}{1.204552in}}%
\pgfpathlineto{\pgfqpoint{3.226514in}{1.211542in}}%
\pgfpathlineto{\pgfqpoint{3.230833in}{1.216563in}}%
\pgfpathlineto{\pgfqpoint{3.238791in}{1.214139in}}%
\pgfpathlineto{\pgfqpoint{3.239989in}{1.209565in}}%
\pgfpathlineto{\pgfqpoint{3.247206in}{1.209157in}}%
\pgfpathlineto{\pgfqpoint{3.253779in}{1.205127in}}%
\pgfpathlineto{\pgfqpoint{3.258356in}{1.209245in}}%
\pgfpathlineto{\pgfqpoint{3.262155in}{1.203038in}}%
\pgfpathlineto{\pgfqpoint{3.265874in}{1.209619in}}%
\pgfpathlineto{\pgfqpoint{3.269662in}{1.209194in}}%
\pgfpathlineto{\pgfqpoint{3.272014in}{1.215831in}}%
\pgfpathlineto{\pgfqpoint{3.282173in}{1.217740in}}%
\pgfpathlineto{\pgfqpoint{3.289627in}{1.222451in}}%
\pgfpathlineto{\pgfqpoint{3.293663in}{1.220232in}}%
\pgfpathlineto{\pgfqpoint{3.301128in}{1.225302in}}%
\pgfpathlineto{\pgfqpoint{3.310778in}{1.227788in}}%
\pgfpathlineto{\pgfqpoint{3.318505in}{1.221809in}}%
\pgfpathlineto{\pgfqpoint{3.321566in}{1.228593in}}%
\pgfpathlineto{\pgfqpoint{3.327494in}{1.228385in}}%
\pgfpathlineto{\pgfqpoint{3.332202in}{1.223602in}}%
\pgfpathlineto{\pgfqpoint{3.337579in}{1.209759in}}%
\pgfpathlineto{\pgfqpoint{3.336122in}{1.199785in}}%
\pgfpathlineto{\pgfqpoint{3.328969in}{1.197798in}}%
\pgfpathlineto{\pgfqpoint{3.323739in}{1.188050in}}%
\pgfpathlineto{\pgfqpoint{3.323024in}{1.183127in}}%
\pgfpathlineto{\pgfqpoint{3.315978in}{1.174823in}}%
\pgfpathlineto{\pgfqpoint{3.307317in}{1.174597in}}%
\pgfpathlineto{\pgfqpoint{3.297293in}{1.177393in}}%
\pgfpathlineto{\pgfqpoint{3.295040in}{1.173975in}}%
\pgfpathlineto{\pgfqpoint{3.285428in}{1.174379in}}%
\pgfpathlineto{\pgfqpoint{3.286885in}{1.169498in}}%
\pgfpathclose%
\pgfusepath{fill}%
\end{pgfscope}%
\begin{pgfscope}%
\pgfpathrectangle{\pgfqpoint{0.100000in}{0.100000in}}{\pgfqpoint{3.608454in}{2.310000in}}%
\pgfusepath{clip}%
\pgfsetbuttcap%
\pgfsetmiterjoin%
\definecolor{currentfill}{rgb}{0.000000,0.541176,0.729412}%
\pgfsetfillcolor{currentfill}%
\pgfsetlinewidth{0.000000pt}%
\definecolor{currentstroke}{rgb}{0.000000,0.000000,0.000000}%
\pgfsetstrokecolor{currentstroke}%
\pgfsetstrokeopacity{0.000000}%
\pgfsetdash{}{0pt}%
\pgfpathmoveto{\pgfqpoint{3.324933in}{1.250084in}}%
\pgfpathlineto{\pgfqpoint{3.338120in}{1.231991in}}%
\pgfpathlineto{\pgfqpoint{3.331617in}{1.232924in}}%
\pgfpathlineto{\pgfqpoint{3.324107in}{1.249673in}}%
\pgfpathlineto{\pgfqpoint{3.324933in}{1.250084in}}%
\pgfpathclose%
\pgfusepath{fill}%
\end{pgfscope}%
\begin{pgfscope}%
\pgfpathrectangle{\pgfqpoint{0.100000in}{0.100000in}}{\pgfqpoint{3.608454in}{2.310000in}}%
\pgfusepath{clip}%
\pgfsetbuttcap%
\pgfsetmiterjoin%
\definecolor{currentfill}{rgb}{0.000000,0.505882,0.747059}%
\pgfsetfillcolor{currentfill}%
\pgfsetlinewidth{0.000000pt}%
\definecolor{currentstroke}{rgb}{0.000000,0.000000,0.000000}%
\pgfsetstrokecolor{currentstroke}%
\pgfsetstrokeopacity{0.000000}%
\pgfsetdash{}{0pt}%
\pgfpathmoveto{\pgfqpoint{1.923573in}{0.599973in}}%
\pgfpathlineto{\pgfqpoint{1.925990in}{0.598363in}}%
\pgfpathlineto{\pgfqpoint{1.945778in}{0.610158in}}%
\pgfpathlineto{\pgfqpoint{1.957611in}{0.597327in}}%
\pgfpathlineto{\pgfqpoint{1.953286in}{0.593627in}}%
\pgfpathlineto{\pgfqpoint{1.950391in}{0.577980in}}%
\pgfpathlineto{\pgfqpoint{1.923720in}{0.556660in}}%
\pgfpathlineto{\pgfqpoint{1.915913in}{0.565873in}}%
\pgfpathlineto{\pgfqpoint{1.907128in}{0.576661in}}%
\pgfpathlineto{\pgfqpoint{1.923573in}{0.599973in}}%
\pgfpathclose%
\pgfusepath{fill}%
\end{pgfscope}%
\begin{pgfscope}%
\pgfpathrectangle{\pgfqpoint{0.100000in}{0.100000in}}{\pgfqpoint{3.608454in}{2.310000in}}%
\pgfusepath{clip}%
\pgfsetbuttcap%
\pgfsetmiterjoin%
\definecolor{currentfill}{rgb}{0.000000,0.623529,0.688235}%
\pgfsetfillcolor{currentfill}%
\pgfsetlinewidth{0.000000pt}%
\definecolor{currentstroke}{rgb}{0.000000,0.000000,0.000000}%
\pgfsetstrokecolor{currentstroke}%
\pgfsetstrokeopacity{0.000000}%
\pgfsetdash{}{0pt}%
\pgfpathmoveto{\pgfqpoint{1.888636in}{1.700186in}}%
\pgfpathlineto{\pgfqpoint{1.956891in}{1.697804in}}%
\pgfpathlineto{\pgfqpoint{1.968709in}{1.697472in}}%
\pgfpathlineto{\pgfqpoint{1.968146in}{1.671305in}}%
\pgfpathlineto{\pgfqpoint{1.928813in}{1.672403in}}%
\pgfpathlineto{\pgfqpoint{1.926808in}{1.664949in}}%
\pgfpathlineto{\pgfqpoint{1.927556in}{1.657034in}}%
\pgfpathlineto{\pgfqpoint{1.925767in}{1.649254in}}%
\pgfpathlineto{\pgfqpoint{1.922640in}{1.646360in}}%
\pgfpathlineto{\pgfqpoint{1.914725in}{1.650247in}}%
\pgfpathlineto{\pgfqpoint{1.905777in}{1.655731in}}%
\pgfpathlineto{\pgfqpoint{1.904126in}{1.659838in}}%
\pgfpathlineto{\pgfqpoint{1.902600in}{1.663282in}}%
\pgfpathlineto{\pgfqpoint{1.893266in}{1.666452in}}%
\pgfpathlineto{\pgfqpoint{1.888697in}{1.672141in}}%
\pgfpathlineto{\pgfqpoint{1.884096in}{1.673023in}}%
\pgfpathlineto{\pgfqpoint{1.881583in}{1.680508in}}%
\pgfpathlineto{\pgfqpoint{1.872226in}{1.686715in}}%
\pgfpathlineto{\pgfqpoint{1.867137in}{1.695946in}}%
\pgfpathlineto{\pgfqpoint{1.862167in}{1.696367in}}%
\pgfpathlineto{\pgfqpoint{1.859723in}{1.701494in}}%
\pgfpathlineto{\pgfqpoint{1.888636in}{1.700186in}}%
\pgfpathclose%
\pgfusepath{fill}%
\end{pgfscope}%
\begin{pgfscope}%
\pgfpathrectangle{\pgfqpoint{0.100000in}{0.100000in}}{\pgfqpoint{3.608454in}{2.310000in}}%
\pgfusepath{clip}%
\pgfsetbuttcap%
\pgfsetmiterjoin%
\definecolor{currentfill}{rgb}{0.000000,0.486275,0.756863}%
\pgfsetfillcolor{currentfill}%
\pgfsetlinewidth{0.000000pt}%
\definecolor{currentstroke}{rgb}{0.000000,0.000000,0.000000}%
\pgfsetstrokecolor{currentstroke}%
\pgfsetstrokeopacity{0.000000}%
\pgfsetdash{}{0pt}%
\pgfpathmoveto{\pgfqpoint{2.688201in}{1.338715in}}%
\pgfpathlineto{\pgfqpoint{2.674475in}{1.337189in}}%
\pgfpathlineto{\pgfqpoint{2.669050in}{1.334303in}}%
\pgfpathlineto{\pgfqpoint{2.667619in}{1.346789in}}%
\pgfpathlineto{\pgfqpoint{2.662944in}{1.348688in}}%
\pgfpathlineto{\pgfqpoint{2.643662in}{1.346766in}}%
\pgfpathlineto{\pgfqpoint{2.641627in}{1.364614in}}%
\pgfpathlineto{\pgfqpoint{2.638885in}{1.368939in}}%
\pgfpathlineto{\pgfqpoint{2.653405in}{1.370465in}}%
\pgfpathlineto{\pgfqpoint{2.650805in}{1.393553in}}%
\pgfpathlineto{\pgfqpoint{2.675138in}{1.396451in}}%
\pgfpathlineto{\pgfqpoint{2.676781in}{1.379202in}}%
\pgfpathlineto{\pgfqpoint{2.684115in}{1.380130in}}%
\pgfpathlineto{\pgfqpoint{2.691293in}{1.385879in}}%
\pgfpathlineto{\pgfqpoint{2.693342in}{1.363630in}}%
\pgfpathlineto{\pgfqpoint{2.683930in}{1.355595in}}%
\pgfpathlineto{\pgfqpoint{2.679482in}{1.355109in}}%
\pgfpathlineto{\pgfqpoint{2.680182in}{1.348259in}}%
\pgfpathlineto{\pgfqpoint{2.687551in}{1.344405in}}%
\pgfpathlineto{\pgfqpoint{2.688201in}{1.338715in}}%
\pgfpathclose%
\pgfusepath{fill}%
\end{pgfscope}%
\begin{pgfscope}%
\pgfpathrectangle{\pgfqpoint{0.100000in}{0.100000in}}{\pgfqpoint{3.608454in}{2.310000in}}%
\pgfusepath{clip}%
\pgfsetbuttcap%
\pgfsetmiterjoin%
\definecolor{currentfill}{rgb}{0.000000,0.596078,0.701961}%
\pgfsetfillcolor{currentfill}%
\pgfsetlinewidth{0.000000pt}%
\definecolor{currentstroke}{rgb}{0.000000,0.000000,0.000000}%
\pgfsetstrokecolor{currentstroke}%
\pgfsetstrokeopacity{0.000000}%
\pgfsetdash{}{0pt}%
\pgfpathmoveto{\pgfqpoint{2.043578in}{1.181785in}}%
\pgfpathlineto{\pgfqpoint{1.971148in}{1.183063in}}%
\pgfpathlineto{\pgfqpoint{1.971690in}{1.220737in}}%
\pgfpathlineto{\pgfqpoint{2.011063in}{1.220009in}}%
\pgfpathlineto{\pgfqpoint{2.012052in}{1.268224in}}%
\pgfpathlineto{\pgfqpoint{2.022297in}{1.268017in}}%
\pgfpathlineto{\pgfqpoint{2.022450in}{1.274884in}}%
\pgfpathlineto{\pgfqpoint{2.047190in}{1.274443in}}%
\pgfpathlineto{\pgfqpoint{2.047067in}{1.264108in}}%
\pgfpathlineto{\pgfqpoint{2.046716in}{1.239913in}}%
\pgfpathlineto{\pgfqpoint{2.045886in}{1.181780in}}%
\pgfpathlineto{\pgfqpoint{2.043578in}{1.181785in}}%
\pgfpathclose%
\pgfusepath{fill}%
\end{pgfscope}%
\begin{pgfscope}%
\pgfpathrectangle{\pgfqpoint{0.100000in}{0.100000in}}{\pgfqpoint{3.608454in}{2.310000in}}%
\pgfusepath{clip}%
\pgfsetbuttcap%
\pgfsetmiterjoin%
\definecolor{currentfill}{rgb}{0.000000,0.901961,0.549020}%
\pgfsetfillcolor{currentfill}%
\pgfsetlinewidth{0.000000pt}%
\definecolor{currentstroke}{rgb}{0.000000,0.000000,0.000000}%
\pgfsetstrokecolor{currentstroke}%
\pgfsetstrokeopacity{0.000000}%
\pgfsetdash{}{0pt}%
\pgfpathmoveto{\pgfqpoint{2.867721in}{1.199764in}}%
\pgfpathlineto{\pgfqpoint{2.827558in}{1.194906in}}%
\pgfpathlineto{\pgfqpoint{2.828062in}{1.198088in}}%
\pgfpathlineto{\pgfqpoint{2.837048in}{1.201444in}}%
\pgfpathlineto{\pgfqpoint{2.856686in}{1.210206in}}%
\pgfpathlineto{\pgfqpoint{2.859365in}{1.219154in}}%
\pgfpathlineto{\pgfqpoint{2.871264in}{1.223809in}}%
\pgfpathlineto{\pgfqpoint{2.870950in}{1.230065in}}%
\pgfpathlineto{\pgfqpoint{2.878218in}{1.237030in}}%
\pgfpathlineto{\pgfqpoint{2.878670in}{1.243251in}}%
\pgfpathlineto{\pgfqpoint{2.887645in}{1.250541in}}%
\pgfpathlineto{\pgfqpoint{2.902185in}{1.260526in}}%
\pgfpathlineto{\pgfqpoint{2.908022in}{1.254600in}}%
\pgfpathlineto{\pgfqpoint{2.916003in}{1.242873in}}%
\pgfpathlineto{\pgfqpoint{2.904928in}{1.234962in}}%
\pgfpathlineto{\pgfqpoint{2.907489in}{1.230091in}}%
\pgfpathlineto{\pgfqpoint{2.901040in}{1.226866in}}%
\pgfpathlineto{\pgfqpoint{2.886736in}{1.225061in}}%
\pgfpathlineto{\pgfqpoint{2.880082in}{1.220503in}}%
\pgfpathlineto{\pgfqpoint{2.876442in}{1.212071in}}%
\pgfpathlineto{\pgfqpoint{2.869228in}{1.206518in}}%
\pgfpathlineto{\pgfqpoint{2.867721in}{1.199764in}}%
\pgfpathclose%
\pgfusepath{fill}%
\end{pgfscope}%
\begin{pgfscope}%
\pgfpathrectangle{\pgfqpoint{0.100000in}{0.100000in}}{\pgfqpoint{3.608454in}{2.310000in}}%
\pgfusepath{clip}%
\pgfsetbuttcap%
\pgfsetmiterjoin%
\definecolor{currentfill}{rgb}{0.000000,0.505882,0.747059}%
\pgfsetfillcolor{currentfill}%
\pgfsetlinewidth{0.000000pt}%
\definecolor{currentstroke}{rgb}{0.000000,0.000000,0.000000}%
\pgfsetstrokecolor{currentstroke}%
\pgfsetstrokeopacity{0.000000}%
\pgfsetdash{}{0pt}%
\pgfpathmoveto{\pgfqpoint{1.683362in}{1.122232in}}%
\pgfpathlineto{\pgfqpoint{1.680896in}{1.087866in}}%
\pgfpathlineto{\pgfqpoint{1.646284in}{1.091077in}}%
\pgfpathlineto{\pgfqpoint{1.589893in}{1.095406in}}%
\pgfpathlineto{\pgfqpoint{1.590719in}{1.104616in}}%
\pgfpathlineto{\pgfqpoint{1.596009in}{1.164627in}}%
\pgfpathlineto{\pgfqpoint{1.598511in}{1.164402in}}%
\pgfpathlineto{\pgfqpoint{1.651744in}{1.159945in}}%
\pgfpathlineto{\pgfqpoint{1.648937in}{1.124836in}}%
\pgfpathlineto{\pgfqpoint{1.683362in}{1.122232in}}%
\pgfpathclose%
\pgfusepath{fill}%
\end{pgfscope}%
\begin{pgfscope}%
\pgfpathrectangle{\pgfqpoint{0.100000in}{0.100000in}}{\pgfqpoint{3.608454in}{2.310000in}}%
\pgfusepath{clip}%
\pgfsetbuttcap%
\pgfsetmiterjoin%
\definecolor{currentfill}{rgb}{0.000000,0.447059,0.776471}%
\pgfsetfillcolor{currentfill}%
\pgfsetlinewidth{0.000000pt}%
\definecolor{currentstroke}{rgb}{0.000000,0.000000,0.000000}%
\pgfsetstrokecolor{currentstroke}%
\pgfsetstrokeopacity{0.000000}%
\pgfsetdash{}{0pt}%
\pgfpathmoveto{\pgfqpoint{1.907060in}{0.482191in}}%
\pgfpathlineto{\pgfqpoint{1.900788in}{0.477356in}}%
\pgfpathlineto{\pgfqpoint{1.902126in}{0.472699in}}%
\pgfpathlineto{\pgfqpoint{1.836840in}{0.475148in}}%
\pgfpathlineto{\pgfqpoint{1.839000in}{0.522528in}}%
\pgfpathlineto{\pgfqpoint{1.840202in}{0.558059in}}%
\pgfpathlineto{\pgfqpoint{1.797569in}{0.560061in}}%
\pgfpathlineto{\pgfqpoint{1.799798in}{0.602976in}}%
\pgfpathlineto{\pgfqpoint{1.786496in}{0.603609in}}%
\pgfpathlineto{\pgfqpoint{1.787654in}{0.626021in}}%
\pgfpathlineto{\pgfqpoint{1.817099in}{0.623588in}}%
\pgfpathlineto{\pgfqpoint{1.834695in}{0.613722in}}%
\pgfpathlineto{\pgfqpoint{1.835751in}{0.642247in}}%
\pgfpathlineto{\pgfqpoint{1.858713in}{0.641352in}}%
\pgfpathlineto{\pgfqpoint{1.870125in}{0.624788in}}%
\pgfpathlineto{\pgfqpoint{1.878483in}{0.632521in}}%
\pgfpathlineto{\pgfqpoint{1.896461in}{0.616726in}}%
\pgfpathlineto{\pgfqpoint{1.898375in}{0.608966in}}%
\pgfpathlineto{\pgfqpoint{1.907278in}{0.617155in}}%
\pgfpathlineto{\pgfqpoint{1.918807in}{0.601623in}}%
\pgfpathlineto{\pgfqpoint{1.923573in}{0.599973in}}%
\pgfpathlineto{\pgfqpoint{1.907128in}{0.576661in}}%
\pgfpathlineto{\pgfqpoint{1.915913in}{0.565873in}}%
\pgfpathlineto{\pgfqpoint{1.882610in}{0.539717in}}%
\pgfpathlineto{\pgfqpoint{1.895145in}{0.523860in}}%
\pgfpathlineto{\pgfqpoint{1.889121in}{0.521871in}}%
\pgfpathlineto{\pgfqpoint{1.907719in}{0.482706in}}%
\pgfpathlineto{\pgfqpoint{1.907060in}{0.482191in}}%
\pgfpathclose%
\pgfusepath{fill}%
\end{pgfscope}%
\begin{pgfscope}%
\pgfpathrectangle{\pgfqpoint{0.100000in}{0.100000in}}{\pgfqpoint{3.608454in}{2.310000in}}%
\pgfusepath{clip}%
\pgfsetbuttcap%
\pgfsetmiterjoin%
\definecolor{currentfill}{rgb}{0.000000,0.572549,0.713725}%
\pgfsetfillcolor{currentfill}%
\pgfsetlinewidth{0.000000pt}%
\definecolor{currentstroke}{rgb}{0.000000,0.000000,0.000000}%
\pgfsetstrokecolor{currentstroke}%
\pgfsetstrokeopacity{0.000000}%
\pgfsetdash{}{0pt}%
\pgfpathmoveto{\pgfqpoint{1.831143in}{0.977031in}}%
\pgfpathlineto{\pgfqpoint{1.820478in}{0.987070in}}%
\pgfpathlineto{\pgfqpoint{1.818337in}{0.980492in}}%
\pgfpathlineto{\pgfqpoint{1.811536in}{0.984126in}}%
\pgfpathlineto{\pgfqpoint{1.805846in}{0.981264in}}%
\pgfpathlineto{\pgfqpoint{1.799655in}{0.981871in}}%
\pgfpathlineto{\pgfqpoint{1.785322in}{0.998334in}}%
\pgfpathlineto{\pgfqpoint{1.780677in}{0.997355in}}%
\pgfpathlineto{\pgfqpoint{1.781300in}{1.012079in}}%
\pgfpathlineto{\pgfqpoint{1.782538in}{1.034546in}}%
\pgfpathlineto{\pgfqpoint{1.796623in}{1.033765in}}%
\pgfpathlineto{\pgfqpoint{1.796992in}{1.040622in}}%
\pgfpathlineto{\pgfqpoint{1.863914in}{1.037452in}}%
\pgfpathlineto{\pgfqpoint{1.872137in}{1.035604in}}%
\pgfpathlineto{\pgfqpoint{1.871452in}{1.016408in}}%
\pgfpathlineto{\pgfqpoint{1.858037in}{1.016994in}}%
\pgfpathlineto{\pgfqpoint{1.857172in}{0.996308in}}%
\pgfpathlineto{\pgfqpoint{1.835171in}{0.996785in}}%
\pgfpathlineto{\pgfqpoint{1.831022in}{0.989497in}}%
\pgfpathlineto{\pgfqpoint{1.831143in}{0.977031in}}%
\pgfpathclose%
\pgfusepath{fill}%
\end{pgfscope}%
\begin{pgfscope}%
\pgfpathrectangle{\pgfqpoint{0.100000in}{0.100000in}}{\pgfqpoint{3.608454in}{2.310000in}}%
\pgfusepath{clip}%
\pgfsetbuttcap%
\pgfsetmiterjoin%
\definecolor{currentfill}{rgb}{0.000000,0.349020,0.825490}%
\pgfsetfillcolor{currentfill}%
\pgfsetlinewidth{0.000000pt}%
\definecolor{currentstroke}{rgb}{0.000000,0.000000,0.000000}%
\pgfsetstrokecolor{currentstroke}%
\pgfsetstrokeopacity{0.000000}%
\pgfsetdash{}{0pt}%
\pgfpathmoveto{\pgfqpoint{2.023407in}{1.319659in}}%
\pgfpathlineto{\pgfqpoint{2.023077in}{1.302462in}}%
\pgfpathlineto{\pgfqpoint{1.994270in}{1.303029in}}%
\pgfpathlineto{\pgfqpoint{1.987451in}{1.303171in}}%
\pgfpathlineto{\pgfqpoint{1.987581in}{1.310037in}}%
\pgfpathlineto{\pgfqpoint{1.960306in}{1.310682in}}%
\pgfpathlineto{\pgfqpoint{1.961886in}{1.365738in}}%
\pgfpathlineto{\pgfqpoint{1.962491in}{1.386340in}}%
\pgfpathlineto{\pgfqpoint{2.031304in}{1.384920in}}%
\pgfpathlineto{\pgfqpoint{2.043672in}{1.384754in}}%
\pgfpathlineto{\pgfqpoint{2.043381in}{1.357123in}}%
\pgfpathlineto{\pgfqpoint{2.034708in}{1.355023in}}%
\pgfpathlineto{\pgfqpoint{2.030596in}{1.356610in}}%
\pgfpathlineto{\pgfqpoint{2.021694in}{1.353951in}}%
\pgfpathlineto{\pgfqpoint{2.021423in}{1.343723in}}%
\pgfpathlineto{\pgfqpoint{2.014646in}{1.343860in}}%
\pgfpathlineto{\pgfqpoint{2.014334in}{1.326679in}}%
\pgfpathlineto{\pgfqpoint{2.021173in}{1.326536in}}%
\pgfpathlineto{\pgfqpoint{2.023407in}{1.319659in}}%
\pgfpathclose%
\pgfusepath{fill}%
\end{pgfscope}%
\begin{pgfscope}%
\pgfpathrectangle{\pgfqpoint{0.100000in}{0.100000in}}{\pgfqpoint{3.608454in}{2.310000in}}%
\pgfusepath{clip}%
\pgfsetbuttcap%
\pgfsetmiterjoin%
\definecolor{currentfill}{rgb}{0.000000,0.482353,0.758824}%
\pgfsetfillcolor{currentfill}%
\pgfsetlinewidth{0.000000pt}%
\definecolor{currentstroke}{rgb}{0.000000,0.000000,0.000000}%
\pgfsetstrokecolor{currentstroke}%
\pgfsetstrokeopacity{0.000000}%
\pgfsetdash{}{0pt}%
\pgfpathmoveto{\pgfqpoint{2.986399in}{1.102377in}}%
\pgfpathlineto{\pgfqpoint{2.947537in}{1.098430in}}%
\pgfpathlineto{\pgfqpoint{2.946773in}{1.103937in}}%
\pgfpathlineto{\pgfqpoint{2.932868in}{1.114078in}}%
\pgfpathlineto{\pgfqpoint{2.926537in}{1.111891in}}%
\pgfpathlineto{\pgfqpoint{2.924776in}{1.115579in}}%
\pgfpathlineto{\pgfqpoint{2.926944in}{1.122067in}}%
\pgfpathlineto{\pgfqpoint{2.930944in}{1.123301in}}%
\pgfpathlineto{\pgfqpoint{2.943450in}{1.124830in}}%
\pgfpathlineto{\pgfqpoint{2.951598in}{1.127420in}}%
\pgfpathlineto{\pgfqpoint{2.956399in}{1.133737in}}%
\pgfpathlineto{\pgfqpoint{2.963707in}{1.130755in}}%
\pgfpathlineto{\pgfqpoint{2.970938in}{1.132406in}}%
\pgfpathlineto{\pgfqpoint{3.007763in}{1.136339in}}%
\pgfpathlineto{\pgfqpoint{3.008399in}{1.134076in}}%
\pgfpathlineto{\pgfqpoint{3.009870in}{1.122844in}}%
\pgfpathlineto{\pgfqpoint{3.007642in}{1.113989in}}%
\pgfpathlineto{\pgfqpoint{3.008536in}{1.104446in}}%
\pgfpathlineto{\pgfqpoint{2.986399in}{1.102377in}}%
\pgfpathclose%
\pgfusepath{fill}%
\end{pgfscope}%
\begin{pgfscope}%
\pgfpathrectangle{\pgfqpoint{0.100000in}{0.100000in}}{\pgfqpoint{3.608454in}{2.310000in}}%
\pgfusepath{clip}%
\pgfsetbuttcap%
\pgfsetmiterjoin%
\definecolor{currentfill}{rgb}{0.000000,0.631373,0.684314}%
\pgfsetfillcolor{currentfill}%
\pgfsetlinewidth{0.000000pt}%
\definecolor{currentstroke}{rgb}{0.000000,0.000000,0.000000}%
\pgfsetstrokecolor{currentstroke}%
\pgfsetstrokeopacity{0.000000}%
\pgfsetdash{}{0pt}%
\pgfpathmoveto{\pgfqpoint{2.922696in}{0.996327in}}%
\pgfpathlineto{\pgfqpoint{2.916400in}{0.998276in}}%
\pgfpathlineto{\pgfqpoint{2.908530in}{0.995600in}}%
\pgfpathlineto{\pgfqpoint{2.897463in}{0.997613in}}%
\pgfpathlineto{\pgfqpoint{2.885344in}{1.008668in}}%
\pgfpathlineto{\pgfqpoint{2.886412in}{1.013136in}}%
\pgfpathlineto{\pgfqpoint{2.871596in}{1.014173in}}%
\pgfpathlineto{\pgfqpoint{2.868390in}{1.010555in}}%
\pgfpathlineto{\pgfqpoint{2.864861in}{1.018162in}}%
\pgfpathlineto{\pgfqpoint{2.864516in}{1.026478in}}%
\pgfpathlineto{\pgfqpoint{2.859003in}{1.029918in}}%
\pgfpathlineto{\pgfqpoint{2.863172in}{1.045964in}}%
\pgfpathlineto{\pgfqpoint{2.864794in}{1.047108in}}%
\pgfpathlineto{\pgfqpoint{2.872404in}{1.042068in}}%
\pgfpathlineto{\pgfqpoint{2.876324in}{1.042276in}}%
\pgfpathlineto{\pgfqpoint{2.884093in}{1.034847in}}%
\pgfpathlineto{\pgfqpoint{2.888625in}{1.033073in}}%
\pgfpathlineto{\pgfqpoint{2.895066in}{1.035008in}}%
\pgfpathlineto{\pgfqpoint{2.907921in}{1.017944in}}%
\pgfpathlineto{\pgfqpoint{2.910642in}{1.010268in}}%
\pgfpathlineto{\pgfqpoint{2.919954in}{1.001789in}}%
\pgfpathlineto{\pgfqpoint{2.922696in}{0.996327in}}%
\pgfpathclose%
\pgfusepath{fill}%
\end{pgfscope}%
\begin{pgfscope}%
\pgfpathrectangle{\pgfqpoint{0.100000in}{0.100000in}}{\pgfqpoint{3.608454in}{2.310000in}}%
\pgfusepath{clip}%
\pgfsetbuttcap%
\pgfsetmiterjoin%
\definecolor{currentfill}{rgb}{0.000000,0.647059,0.676471}%
\pgfsetfillcolor{currentfill}%
\pgfsetlinewidth{0.000000pt}%
\definecolor{currentstroke}{rgb}{0.000000,0.000000,0.000000}%
\pgfsetstrokecolor{currentstroke}%
\pgfsetstrokeopacity{0.000000}%
\pgfsetdash{}{0pt}%
\pgfpathmoveto{\pgfqpoint{0.960301in}{1.531851in}}%
\pgfpathlineto{\pgfqpoint{0.954290in}{1.503771in}}%
\pgfpathlineto{\pgfqpoint{0.939931in}{1.436894in}}%
\pgfpathlineto{\pgfqpoint{0.882599in}{1.449402in}}%
\pgfpathlineto{\pgfqpoint{0.836825in}{1.499170in}}%
\pgfpathlineto{\pgfqpoint{0.795453in}{1.508995in}}%
\pgfpathlineto{\pgfqpoint{0.752002in}{1.519867in}}%
\pgfpathlineto{\pgfqpoint{0.724276in}{1.521184in}}%
\pgfpathlineto{\pgfqpoint{0.725481in}{1.527601in}}%
\pgfpathlineto{\pgfqpoint{0.730363in}{1.531286in}}%
\pgfpathlineto{\pgfqpoint{0.730097in}{1.540859in}}%
\pgfpathlineto{\pgfqpoint{0.738497in}{1.547647in}}%
\pgfpathlineto{\pgfqpoint{0.749642in}{1.550822in}}%
\pgfpathlineto{\pgfqpoint{0.753176in}{1.557006in}}%
\pgfpathlineto{\pgfqpoint{0.753615in}{1.567143in}}%
\pgfpathlineto{\pgfqpoint{0.758589in}{1.580048in}}%
\pgfpathlineto{\pgfqpoint{0.755783in}{1.587192in}}%
\pgfpathlineto{\pgfqpoint{0.785394in}{1.632033in}}%
\pgfpathlineto{\pgfqpoint{0.798688in}{1.628778in}}%
\pgfpathlineto{\pgfqpoint{0.805480in}{1.656197in}}%
\pgfpathlineto{\pgfqpoint{0.824452in}{1.732928in}}%
\pgfpathlineto{\pgfqpoint{0.861529in}{1.723655in}}%
\pgfpathlineto{\pgfqpoint{0.937739in}{1.705964in}}%
\pgfpathlineto{\pgfqpoint{0.949739in}{1.703791in}}%
\pgfpathlineto{\pgfqpoint{0.981161in}{1.696275in}}%
\pgfpathlineto{\pgfqpoint{0.994964in}{1.693281in}}%
\pgfpathlineto{\pgfqpoint{0.978586in}{1.616384in}}%
\pgfpathlineto{\pgfqpoint{0.960301in}{1.531851in}}%
\pgfpathclose%
\pgfusepath{fill}%
\end{pgfscope}%
\begin{pgfscope}%
\pgfpathrectangle{\pgfqpoint{0.100000in}{0.100000in}}{\pgfqpoint{3.608454in}{2.310000in}}%
\pgfusepath{clip}%
\pgfsetbuttcap%
\pgfsetmiterjoin%
\definecolor{currentfill}{rgb}{0.000000,0.796078,0.601961}%
\pgfsetfillcolor{currentfill}%
\pgfsetlinewidth{0.000000pt}%
\definecolor{currentstroke}{rgb}{0.000000,0.000000,0.000000}%
\pgfsetstrokecolor{currentstroke}%
\pgfsetstrokeopacity{0.000000}%
\pgfsetdash{}{0pt}%
\pgfpathmoveto{\pgfqpoint{1.859723in}{1.701494in}}%
\pgfpathlineto{\pgfqpoint{1.862167in}{1.696367in}}%
\pgfpathlineto{\pgfqpoint{1.867137in}{1.695946in}}%
\pgfpathlineto{\pgfqpoint{1.872226in}{1.686715in}}%
\pgfpathlineto{\pgfqpoint{1.881583in}{1.680508in}}%
\pgfpathlineto{\pgfqpoint{1.884096in}{1.673023in}}%
\pgfpathlineto{\pgfqpoint{1.888697in}{1.672141in}}%
\pgfpathlineto{\pgfqpoint{1.893266in}{1.666452in}}%
\pgfpathlineto{\pgfqpoint{1.902600in}{1.663282in}}%
\pgfpathlineto{\pgfqpoint{1.904126in}{1.659838in}}%
\pgfpathlineto{\pgfqpoint{1.860373in}{1.661671in}}%
\pgfpathlineto{\pgfqpoint{1.805743in}{1.664546in}}%
\pgfpathlineto{\pgfqpoint{1.805445in}{1.674881in}}%
\pgfpathlineto{\pgfqpoint{1.806996in}{1.702523in}}%
\pgfpathlineto{\pgfqpoint{1.807130in}{1.721260in}}%
\pgfpathlineto{\pgfqpoint{1.817373in}{1.718961in}}%
\pgfpathlineto{\pgfqpoint{1.829398in}{1.718027in}}%
\pgfpathlineto{\pgfqpoint{1.842006in}{1.721715in}}%
\pgfpathlineto{\pgfqpoint{1.841021in}{1.702423in}}%
\pgfpathlineto{\pgfqpoint{1.859723in}{1.701494in}}%
\pgfpathclose%
\pgfusepath{fill}%
\end{pgfscope}%
\begin{pgfscope}%
\pgfpathrectangle{\pgfqpoint{0.100000in}{0.100000in}}{\pgfqpoint{3.608454in}{2.310000in}}%
\pgfusepath{clip}%
\pgfsetbuttcap%
\pgfsetmiterjoin%
\definecolor{currentfill}{rgb}{0.000000,0.709804,0.645098}%
\pgfsetfillcolor{currentfill}%
\pgfsetlinewidth{0.000000pt}%
\definecolor{currentstroke}{rgb}{0.000000,0.000000,0.000000}%
\pgfsetstrokecolor{currentstroke}%
\pgfsetstrokeopacity{0.000000}%
\pgfsetdash{}{0pt}%
\pgfpathmoveto{\pgfqpoint{2.811462in}{0.766109in}}%
\pgfpathlineto{\pgfqpoint{2.792796in}{0.763762in}}%
\pgfpathlineto{\pgfqpoint{2.794089in}{0.749184in}}%
\pgfpathlineto{\pgfqpoint{2.786852in}{0.748618in}}%
\pgfpathlineto{\pgfqpoint{2.781206in}{0.755220in}}%
\pgfpathlineto{\pgfqpoint{2.779513in}{0.764362in}}%
\pgfpathlineto{\pgfqpoint{2.781911in}{0.785540in}}%
\pgfpathlineto{\pgfqpoint{2.780161in}{0.791452in}}%
\pgfpathlineto{\pgfqpoint{2.774910in}{0.796997in}}%
\pgfpathlineto{\pgfqpoint{2.755229in}{0.795820in}}%
\pgfpathlineto{\pgfqpoint{2.755934in}{0.788907in}}%
\pgfpathlineto{\pgfqpoint{2.733549in}{0.786465in}}%
\pgfpathlineto{\pgfqpoint{2.737974in}{0.800305in}}%
\pgfpathlineto{\pgfqpoint{2.737520in}{0.807875in}}%
\pgfpathlineto{\pgfqpoint{2.741287in}{0.817640in}}%
\pgfpathlineto{\pgfqpoint{2.751810in}{0.820152in}}%
\pgfpathlineto{\pgfqpoint{2.751930in}{0.830708in}}%
\pgfpathlineto{\pgfqpoint{2.762136in}{0.831926in}}%
\pgfpathlineto{\pgfqpoint{2.767748in}{0.825602in}}%
\pgfpathlineto{\pgfqpoint{2.776392in}{0.826653in}}%
\pgfpathlineto{\pgfqpoint{2.778092in}{0.819839in}}%
\pgfpathlineto{\pgfqpoint{2.786918in}{0.816730in}}%
\pgfpathlineto{\pgfqpoint{2.808283in}{0.818788in}}%
\pgfpathlineto{\pgfqpoint{2.809408in}{0.807037in}}%
\pgfpathlineto{\pgfqpoint{2.814786in}{0.800405in}}%
\pgfpathlineto{\pgfqpoint{2.821071in}{0.796344in}}%
\pgfpathlineto{\pgfqpoint{2.821443in}{0.790502in}}%
\pgfpathlineto{\pgfqpoint{2.825145in}{0.783605in}}%
\pgfpathlineto{\pgfqpoint{2.820126in}{0.780667in}}%
\pgfpathlineto{\pgfqpoint{2.810152in}{0.780002in}}%
\pgfpathlineto{\pgfqpoint{2.811462in}{0.766109in}}%
\pgfpathclose%
\pgfusepath{fill}%
\end{pgfscope}%
\begin{pgfscope}%
\pgfpathrectangle{\pgfqpoint{0.100000in}{0.100000in}}{\pgfqpoint{3.608454in}{2.310000in}}%
\pgfusepath{clip}%
\pgfsetbuttcap%
\pgfsetmiterjoin%
\definecolor{currentfill}{rgb}{0.000000,0.635294,0.682353}%
\pgfsetfillcolor{currentfill}%
\pgfsetlinewidth{0.000000pt}%
\definecolor{currentstroke}{rgb}{0.000000,0.000000,0.000000}%
\pgfsetstrokecolor{currentstroke}%
\pgfsetstrokeopacity{0.000000}%
\pgfsetdash{}{0pt}%
\pgfpathmoveto{\pgfqpoint{1.088551in}{1.993499in}}%
\pgfpathlineto{\pgfqpoint{1.086269in}{1.988945in}}%
\pgfpathlineto{\pgfqpoint{1.076963in}{1.988148in}}%
\pgfpathlineto{\pgfqpoint{1.069612in}{1.982716in}}%
\pgfpathlineto{\pgfqpoint{1.061696in}{1.979359in}}%
\pgfpathlineto{\pgfqpoint{1.053520in}{1.975559in}}%
\pgfpathlineto{\pgfqpoint{1.048671in}{1.970878in}}%
\pgfpathlineto{\pgfqpoint{1.034877in}{1.968576in}}%
\pgfpathlineto{\pgfqpoint{1.028023in}{1.976205in}}%
\pgfpathlineto{\pgfqpoint{1.025284in}{1.975956in}}%
\pgfpathlineto{\pgfqpoint{1.030719in}{1.983835in}}%
\pgfpathlineto{\pgfqpoint{1.028886in}{1.992775in}}%
\pgfpathlineto{\pgfqpoint{1.033186in}{1.997675in}}%
\pgfpathlineto{\pgfqpoint{1.040540in}{1.999733in}}%
\pgfpathlineto{\pgfqpoint{1.037550in}{2.013584in}}%
\pgfpathlineto{\pgfqpoint{1.042048in}{2.022019in}}%
\pgfpathlineto{\pgfqpoint{1.042825in}{2.030433in}}%
\pgfpathlineto{\pgfqpoint{1.049410in}{2.040224in}}%
\pgfpathlineto{\pgfqpoint{1.056865in}{2.058071in}}%
\pgfpathlineto{\pgfqpoint{1.055012in}{2.060043in}}%
\pgfpathlineto{\pgfqpoint{1.041035in}{2.060695in}}%
\pgfpathlineto{\pgfqpoint{1.041857in}{2.066086in}}%
\pgfpathlineto{\pgfqpoint{1.027942in}{2.078990in}}%
\pgfpathlineto{\pgfqpoint{1.027701in}{2.087646in}}%
\pgfpathlineto{\pgfqpoint{1.023030in}{2.093793in}}%
\pgfpathlineto{\pgfqpoint{1.014258in}{2.117180in}}%
\pgfpathlineto{\pgfqpoint{1.003956in}{2.123021in}}%
\pgfpathlineto{\pgfqpoint{1.003713in}{2.127225in}}%
\pgfpathlineto{\pgfqpoint{0.994972in}{2.136779in}}%
\pgfpathlineto{\pgfqpoint{1.001953in}{2.138275in}}%
\pgfpathlineto{\pgfqpoint{1.011451in}{2.134988in}}%
\pgfpathlineto{\pgfqpoint{1.020317in}{2.134665in}}%
\pgfpathlineto{\pgfqpoint{1.034696in}{2.123940in}}%
\pgfpathlineto{\pgfqpoint{1.033798in}{2.117353in}}%
\pgfpathlineto{\pgfqpoint{1.039741in}{2.112438in}}%
\pgfpathlineto{\pgfqpoint{1.048078in}{2.111068in}}%
\pgfpathlineto{\pgfqpoint{1.060904in}{2.103004in}}%
\pgfpathlineto{\pgfqpoint{1.068378in}{2.094075in}}%
\pgfpathlineto{\pgfqpoint{1.072235in}{2.094680in}}%
\pgfpathlineto{\pgfqpoint{1.085713in}{2.091817in}}%
\pgfpathlineto{\pgfqpoint{1.092364in}{2.093954in}}%
\pgfpathlineto{\pgfqpoint{1.091464in}{2.102885in}}%
\pgfpathlineto{\pgfqpoint{1.093234in}{2.110956in}}%
\pgfpathlineto{\pgfqpoint{1.090372in}{2.118413in}}%
\pgfpathlineto{\pgfqpoint{1.093543in}{2.127613in}}%
\pgfpathlineto{\pgfqpoint{1.118004in}{2.122768in}}%
\pgfpathlineto{\pgfqpoint{1.111233in}{2.089809in}}%
\pgfpathlineto{\pgfqpoint{1.119922in}{2.088067in}}%
\pgfpathlineto{\pgfqpoint{1.114370in}{2.060903in}}%
\pgfpathlineto{\pgfqpoint{1.108312in}{2.062127in}}%
\pgfpathlineto{\pgfqpoint{1.103510in}{2.056028in}}%
\pgfpathlineto{\pgfqpoint{1.092981in}{2.055856in}}%
\pgfpathlineto{\pgfqpoint{1.092070in}{2.051353in}}%
\pgfpathlineto{\pgfqpoint{1.085312in}{2.052758in}}%
\pgfpathlineto{\pgfqpoint{1.078479in}{2.042039in}}%
\pgfpathlineto{\pgfqpoint{1.077963in}{2.034819in}}%
\pgfpathlineto{\pgfqpoint{1.081054in}{2.029697in}}%
\pgfpathlineto{\pgfqpoint{1.081296in}{2.021948in}}%
\pgfpathlineto{\pgfqpoint{1.076933in}{2.013289in}}%
\pgfpathlineto{\pgfqpoint{1.076031in}{2.007092in}}%
\pgfpathlineto{\pgfqpoint{1.080846in}{2.002972in}}%
\pgfpathlineto{\pgfqpoint{1.081829in}{1.997042in}}%
\pgfpathlineto{\pgfqpoint{1.088551in}{1.993499in}}%
\pgfpathclose%
\pgfusepath{fill}%
\end{pgfscope}%
\begin{pgfscope}%
\pgfpathrectangle{\pgfqpoint{0.100000in}{0.100000in}}{\pgfqpoint{3.608454in}{2.310000in}}%
\pgfusepath{clip}%
\pgfsetbuttcap%
\pgfsetmiterjoin%
\definecolor{currentfill}{rgb}{0.000000,0.701961,0.649020}%
\pgfsetfillcolor{currentfill}%
\pgfsetlinewidth{0.000000pt}%
\definecolor{currentstroke}{rgb}{0.000000,0.000000,0.000000}%
\pgfsetstrokecolor{currentstroke}%
\pgfsetstrokeopacity{0.000000}%
\pgfsetdash{}{0pt}%
\pgfpathmoveto{\pgfqpoint{2.193329in}{1.907307in}}%
\pgfpathlineto{\pgfqpoint{2.214175in}{1.908040in}}%
\pgfpathlineto{\pgfqpoint{2.213676in}{1.928846in}}%
\pgfpathlineto{\pgfqpoint{2.255541in}{1.929754in}}%
\pgfpathlineto{\pgfqpoint{2.256111in}{1.909051in}}%
\pgfpathlineto{\pgfqpoint{2.256290in}{1.902446in}}%
\pgfpathlineto{\pgfqpoint{2.253310in}{1.897735in}}%
\pgfpathlineto{\pgfqpoint{2.237467in}{1.890615in}}%
\pgfpathlineto{\pgfqpoint{2.233501in}{1.887214in}}%
\pgfpathlineto{\pgfqpoint{2.229731in}{1.876962in}}%
\pgfpathlineto{\pgfqpoint{2.226726in}{1.874225in}}%
\pgfpathlineto{\pgfqpoint{2.189353in}{1.873824in}}%
\pgfpathlineto{\pgfqpoint{2.188794in}{1.893513in}}%
\pgfpathlineto{\pgfqpoint{2.193529in}{1.893561in}}%
\pgfpathlineto{\pgfqpoint{2.193329in}{1.907307in}}%
\pgfpathclose%
\pgfusepath{fill}%
\end{pgfscope}%
\begin{pgfscope}%
\pgfpathrectangle{\pgfqpoint{0.100000in}{0.100000in}}{\pgfqpoint{3.608454in}{2.310000in}}%
\pgfusepath{clip}%
\pgfsetbuttcap%
\pgfsetmiterjoin%
\definecolor{currentfill}{rgb}{0.000000,0.592157,0.703922}%
\pgfsetfillcolor{currentfill}%
\pgfsetlinewidth{0.000000pt}%
\definecolor{currentstroke}{rgb}{0.000000,0.000000,0.000000}%
\pgfsetstrokecolor{currentstroke}%
\pgfsetstrokeopacity{0.000000}%
\pgfsetdash{}{0pt}%
\pgfpathmoveto{\pgfqpoint{1.854168in}{1.187060in}}%
\pgfpathlineto{\pgfqpoint{1.825437in}{1.188388in}}%
\pgfpathlineto{\pgfqpoint{1.820080in}{1.188664in}}%
\pgfpathlineto{\pgfqpoint{1.821457in}{1.218805in}}%
\pgfpathlineto{\pgfqpoint{1.855479in}{1.217455in}}%
\pgfpathlineto{\pgfqpoint{1.854168in}{1.187060in}}%
\pgfpathclose%
\pgfusepath{fill}%
\end{pgfscope}%
\begin{pgfscope}%
\pgfpathrectangle{\pgfqpoint{0.100000in}{0.100000in}}{\pgfqpoint{3.608454in}{2.310000in}}%
\pgfusepath{clip}%
\pgfsetbuttcap%
\pgfsetmiterjoin%
\definecolor{currentfill}{rgb}{0.000000,0.701961,0.649020}%
\pgfsetfillcolor{currentfill}%
\pgfsetlinewidth{0.000000pt}%
\definecolor{currentstroke}{rgb}{0.000000,0.000000,0.000000}%
\pgfsetstrokecolor{currentstroke}%
\pgfsetstrokeopacity{0.000000}%
\pgfsetdash{}{0pt}%
\pgfpathmoveto{\pgfqpoint{2.589411in}{1.214310in}}%
\pgfpathlineto{\pgfqpoint{2.577947in}{1.209252in}}%
\pgfpathlineto{\pgfqpoint{2.567654in}{1.211922in}}%
\pgfpathlineto{\pgfqpoint{2.566571in}{1.201683in}}%
\pgfpathlineto{\pgfqpoint{2.563459in}{1.197714in}}%
\pgfpathlineto{\pgfqpoint{2.555635in}{1.195943in}}%
\pgfpathlineto{\pgfqpoint{2.539380in}{1.187318in}}%
\pgfpathlineto{\pgfqpoint{2.533377in}{1.204248in}}%
\pgfpathlineto{\pgfqpoint{2.535791in}{1.209919in}}%
\pgfpathlineto{\pgfqpoint{2.533129in}{1.219044in}}%
\pgfpathlineto{\pgfqpoint{2.523585in}{1.228784in}}%
\pgfpathlineto{\pgfqpoint{2.528115in}{1.232902in}}%
\pgfpathlineto{\pgfqpoint{2.540147in}{1.235271in}}%
\pgfpathlineto{\pgfqpoint{2.534558in}{1.250105in}}%
\pgfpathlineto{\pgfqpoint{2.541890in}{1.261426in}}%
\pgfpathlineto{\pgfqpoint{2.549388in}{1.262674in}}%
\pgfpathlineto{\pgfqpoint{2.546912in}{1.269077in}}%
\pgfpathlineto{\pgfqpoint{2.553706in}{1.268477in}}%
\pgfpathlineto{\pgfqpoint{2.562848in}{1.271335in}}%
\pgfpathlineto{\pgfqpoint{2.567447in}{1.266096in}}%
\pgfpathlineto{\pgfqpoint{2.570793in}{1.273710in}}%
\pgfpathlineto{\pgfqpoint{2.578231in}{1.275886in}}%
\pgfpathlineto{\pgfqpoint{2.586412in}{1.272897in}}%
\pgfpathlineto{\pgfqpoint{2.585709in}{1.266483in}}%
\pgfpathlineto{\pgfqpoint{2.577210in}{1.248460in}}%
\pgfpathlineto{\pgfqpoint{2.585128in}{1.245024in}}%
\pgfpathlineto{\pgfqpoint{2.586421in}{1.235460in}}%
\pgfpathlineto{\pgfqpoint{2.589982in}{1.231634in}}%
\pgfpathlineto{\pgfqpoint{2.585378in}{1.221768in}}%
\pgfpathlineto{\pgfqpoint{2.589411in}{1.214310in}}%
\pgfpathclose%
\pgfusepath{fill}%
\end{pgfscope}%
\begin{pgfscope}%
\pgfpathrectangle{\pgfqpoint{0.100000in}{0.100000in}}{\pgfqpoint{3.608454in}{2.310000in}}%
\pgfusepath{clip}%
\pgfsetbuttcap%
\pgfsetmiterjoin%
\definecolor{currentfill}{rgb}{0.000000,0.588235,0.705882}%
\pgfsetfillcolor{currentfill}%
\pgfsetlinewidth{0.000000pt}%
\definecolor{currentstroke}{rgb}{0.000000,0.000000,0.000000}%
\pgfsetstrokecolor{currentstroke}%
\pgfsetstrokeopacity{0.000000}%
\pgfsetdash{}{0pt}%
\pgfpathmoveto{\pgfqpoint{1.638167in}{2.022688in}}%
\pgfpathlineto{\pgfqpoint{1.635493in}{1.994898in}}%
\pgfpathlineto{\pgfqpoint{1.638609in}{1.994610in}}%
\pgfpathlineto{\pgfqpoint{1.635984in}{1.966799in}}%
\pgfpathlineto{\pgfqpoint{1.625532in}{1.967784in}}%
\pgfpathlineto{\pgfqpoint{1.624849in}{1.960734in}}%
\pgfpathlineto{\pgfqpoint{1.611520in}{1.962062in}}%
\pgfpathlineto{\pgfqpoint{1.612308in}{1.970040in}}%
\pgfpathlineto{\pgfqpoint{1.595495in}{1.971747in}}%
\pgfpathlineto{\pgfqpoint{1.592373in}{1.975555in}}%
\pgfpathlineto{\pgfqpoint{1.582076in}{1.976639in}}%
\pgfpathlineto{\pgfqpoint{1.583733in}{1.990580in}}%
\pgfpathlineto{\pgfqpoint{1.560751in}{1.993122in}}%
\pgfpathlineto{\pgfqpoint{1.558985in}{1.997946in}}%
\pgfpathlineto{\pgfqpoint{1.552079in}{1.998723in}}%
\pgfpathlineto{\pgfqpoint{1.552609in}{2.003322in}}%
\pgfpathlineto{\pgfqpoint{1.545674in}{2.004107in}}%
\pgfpathlineto{\pgfqpoint{1.547528in}{2.020313in}}%
\pgfpathlineto{\pgfqpoint{1.543137in}{2.020826in}}%
\pgfpathlineto{\pgfqpoint{1.546296in}{2.048425in}}%
\pgfpathlineto{\pgfqpoint{1.548607in}{2.048174in}}%
\pgfpathlineto{\pgfqpoint{1.551007in}{2.068911in}}%
\pgfpathlineto{\pgfqpoint{1.557865in}{2.068157in}}%
\pgfpathlineto{\pgfqpoint{1.557068in}{2.061233in}}%
\pgfpathlineto{\pgfqpoint{1.570807in}{2.059681in}}%
\pgfpathlineto{\pgfqpoint{1.570016in}{2.052759in}}%
\pgfpathlineto{\pgfqpoint{1.590639in}{2.050503in}}%
\pgfpathlineto{\pgfqpoint{1.589896in}{2.043588in}}%
\pgfpathlineto{\pgfqpoint{1.594243in}{2.043152in}}%
\pgfpathlineto{\pgfqpoint{1.592763in}{2.029281in}}%
\pgfpathlineto{\pgfqpoint{1.613315in}{2.027119in}}%
\pgfpathlineto{\pgfqpoint{1.617764in}{2.024744in}}%
\pgfpathlineto{\pgfqpoint{1.638167in}{2.022688in}}%
\pgfpathclose%
\pgfusepath{fill}%
\end{pgfscope}%
\begin{pgfscope}%
\pgfpathrectangle{\pgfqpoint{0.100000in}{0.100000in}}{\pgfqpoint{3.608454in}{2.310000in}}%
\pgfusepath{clip}%
\pgfsetbuttcap%
\pgfsetmiterjoin%
\definecolor{currentfill}{rgb}{0.000000,0.647059,0.676471}%
\pgfsetfillcolor{currentfill}%
\pgfsetlinewidth{0.000000pt}%
\definecolor{currentstroke}{rgb}{0.000000,0.000000,0.000000}%
\pgfsetstrokecolor{currentstroke}%
\pgfsetstrokeopacity{0.000000}%
\pgfsetdash{}{0pt}%
\pgfpathmoveto{\pgfqpoint{2.534753in}{1.425828in}}%
\pgfpathlineto{\pgfqpoint{2.534185in}{1.410410in}}%
\pgfpathlineto{\pgfqpoint{2.528514in}{1.407327in}}%
\pgfpathlineto{\pgfqpoint{2.503724in}{1.405402in}}%
\pgfpathlineto{\pgfqpoint{2.506265in}{1.370984in}}%
\pgfpathlineto{\pgfqpoint{2.465371in}{1.368430in}}%
\pgfpathlineto{\pgfqpoint{2.464725in}{1.378777in}}%
\pgfpathlineto{\pgfqpoint{2.471713in}{1.379003in}}%
\pgfpathlineto{\pgfqpoint{2.470095in}{1.403294in}}%
\pgfpathlineto{\pgfqpoint{2.463205in}{1.402934in}}%
\pgfpathlineto{\pgfqpoint{2.462297in}{1.414401in}}%
\pgfpathlineto{\pgfqpoint{2.457647in}{1.415090in}}%
\pgfpathlineto{\pgfqpoint{2.457108in}{1.423280in}}%
\pgfpathlineto{\pgfqpoint{2.461611in}{1.423638in}}%
\pgfpathlineto{\pgfqpoint{2.460863in}{1.433950in}}%
\pgfpathlineto{\pgfqpoint{2.484922in}{1.436043in}}%
\pgfpathlineto{\pgfqpoint{2.493952in}{1.454636in}}%
\pgfpathlineto{\pgfqpoint{2.500893in}{1.455137in}}%
\pgfpathlineto{\pgfqpoint{2.498999in}{1.481611in}}%
\pgfpathlineto{\pgfqpoint{2.512434in}{1.482660in}}%
\pgfpathlineto{\pgfqpoint{2.509486in}{1.512307in}}%
\pgfpathlineto{\pgfqpoint{2.516336in}{1.513050in}}%
\pgfpathlineto{\pgfqpoint{2.520205in}{1.472936in}}%
\pgfpathlineto{\pgfqpoint{2.531157in}{1.473630in}}%
\pgfpathlineto{\pgfqpoint{2.532347in}{1.453054in}}%
\pgfpathlineto{\pgfqpoint{2.534753in}{1.425828in}}%
\pgfpathclose%
\pgfusepath{fill}%
\end{pgfscope}%
\begin{pgfscope}%
\pgfpathrectangle{\pgfqpoint{0.100000in}{0.100000in}}{\pgfqpoint{3.608454in}{2.310000in}}%
\pgfusepath{clip}%
\pgfsetbuttcap%
\pgfsetmiterjoin%
\definecolor{currentfill}{rgb}{0.000000,0.478431,0.760784}%
\pgfsetfillcolor{currentfill}%
\pgfsetlinewidth{0.000000pt}%
\definecolor{currentstroke}{rgb}{0.000000,0.000000,0.000000}%
\pgfsetstrokecolor{currentstroke}%
\pgfsetstrokeopacity{0.000000}%
\pgfsetdash{}{0pt}%
\pgfpathmoveto{\pgfqpoint{1.638167in}{2.022688in}}%
\pgfpathlineto{\pgfqpoint{1.617764in}{2.024744in}}%
\pgfpathlineto{\pgfqpoint{1.613315in}{2.027119in}}%
\pgfpathlineto{\pgfqpoint{1.592763in}{2.029281in}}%
\pgfpathlineto{\pgfqpoint{1.594243in}{2.043152in}}%
\pgfpathlineto{\pgfqpoint{1.589896in}{2.043588in}}%
\pgfpathlineto{\pgfqpoint{1.590639in}{2.050503in}}%
\pgfpathlineto{\pgfqpoint{1.570016in}{2.052759in}}%
\pgfpathlineto{\pgfqpoint{1.570807in}{2.059681in}}%
\pgfpathlineto{\pgfqpoint{1.557068in}{2.061233in}}%
\pgfpathlineto{\pgfqpoint{1.557865in}{2.068157in}}%
\pgfpathlineto{\pgfqpoint{1.560678in}{2.074073in}}%
\pgfpathlineto{\pgfqpoint{1.562458in}{2.089780in}}%
\pgfpathlineto{\pgfqpoint{1.569948in}{2.088334in}}%
\pgfpathlineto{\pgfqpoint{1.582871in}{2.091842in}}%
\pgfpathlineto{\pgfqpoint{1.590181in}{2.089744in}}%
\pgfpathlineto{\pgfqpoint{1.601867in}{2.089959in}}%
\pgfpathlineto{\pgfqpoint{1.607083in}{2.083595in}}%
\pgfpathlineto{\pgfqpoint{1.619223in}{2.081466in}}%
\pgfpathlineto{\pgfqpoint{1.623124in}{2.077735in}}%
\pgfpathlineto{\pgfqpoint{1.632154in}{2.074810in}}%
\pgfpathlineto{\pgfqpoint{1.643336in}{2.085862in}}%
\pgfpathlineto{\pgfqpoint{1.647283in}{2.084906in}}%
\pgfpathlineto{\pgfqpoint{1.648734in}{2.078432in}}%
\pgfpathlineto{\pgfqpoint{1.653010in}{2.075697in}}%
\pgfpathlineto{\pgfqpoint{1.664466in}{2.078762in}}%
\pgfpathlineto{\pgfqpoint{1.669040in}{2.084553in}}%
\pgfpathlineto{\pgfqpoint{1.686894in}{2.083116in}}%
\pgfpathlineto{\pgfqpoint{1.690580in}{2.101766in}}%
\pgfpathlineto{\pgfqpoint{1.687537in}{2.102006in}}%
\pgfpathlineto{\pgfqpoint{1.688717in}{2.115878in}}%
\pgfpathlineto{\pgfqpoint{1.723190in}{2.112972in}}%
\pgfpathlineto{\pgfqpoint{1.737002in}{2.111989in}}%
\pgfpathlineto{\pgfqpoint{1.735922in}{2.098000in}}%
\pgfpathlineto{\pgfqpoint{1.738640in}{2.097811in}}%
\pgfpathlineto{\pgfqpoint{1.736501in}{2.069931in}}%
\pgfpathlineto{\pgfqpoint{1.739211in}{2.069716in}}%
\pgfpathlineto{\pgfqpoint{1.738152in}{2.055718in}}%
\pgfpathlineto{\pgfqpoint{1.710634in}{2.057876in}}%
\pgfpathlineto{\pgfqpoint{1.710047in}{2.050726in}}%
\pgfpathlineto{\pgfqpoint{1.700751in}{2.053267in}}%
\pgfpathlineto{\pgfqpoint{1.695708in}{2.045165in}}%
\pgfpathlineto{\pgfqpoint{1.671188in}{2.047334in}}%
\pgfpathlineto{\pgfqpoint{1.668766in}{2.019808in}}%
\pgfpathlineto{\pgfqpoint{1.638167in}{2.022688in}}%
\pgfpathclose%
\pgfusepath{fill}%
\end{pgfscope}%
\begin{pgfscope}%
\pgfpathrectangle{\pgfqpoint{0.100000in}{0.100000in}}{\pgfqpoint{3.608454in}{2.310000in}}%
\pgfusepath{clip}%
\pgfsetbuttcap%
\pgfsetmiterjoin%
\definecolor{currentfill}{rgb}{0.000000,0.690196,0.654902}%
\pgfsetfillcolor{currentfill}%
\pgfsetlinewidth{0.000000pt}%
\definecolor{currentstroke}{rgb}{0.000000,0.000000,0.000000}%
\pgfsetstrokecolor{currentstroke}%
\pgfsetstrokeopacity{0.000000}%
\pgfsetdash{}{0pt}%
\pgfpathmoveto{\pgfqpoint{3.535136in}{1.728569in}}%
\pgfpathlineto{\pgfqpoint{3.546887in}{1.721036in}}%
\pgfpathlineto{\pgfqpoint{3.532181in}{1.716409in}}%
\pgfpathlineto{\pgfqpoint{3.530762in}{1.723556in}}%
\pgfpathlineto{\pgfqpoint{3.535136in}{1.728569in}}%
\pgfpathclose%
\pgfusepath{fill}%
\end{pgfscope}%
\begin{pgfscope}%
\pgfpathrectangle{\pgfqpoint{0.100000in}{0.100000in}}{\pgfqpoint{3.608454in}{2.310000in}}%
\pgfusepath{clip}%
\pgfsetbuttcap%
\pgfsetmiterjoin%
\definecolor{currentfill}{rgb}{0.000000,0.494118,0.752941}%
\pgfsetfillcolor{currentfill}%
\pgfsetlinewidth{0.000000pt}%
\definecolor{currentstroke}{rgb}{0.000000,0.000000,0.000000}%
\pgfsetstrokecolor{currentstroke}%
\pgfsetstrokeopacity{0.000000}%
\pgfsetdash{}{0pt}%
\pgfpathmoveto{\pgfqpoint{3.153524in}{1.337961in}}%
\pgfpathlineto{\pgfqpoint{3.149422in}{1.332537in}}%
\pgfpathlineto{\pgfqpoint{3.149648in}{1.328164in}}%
\pgfpathlineto{\pgfqpoint{3.140579in}{1.302693in}}%
\pgfpathlineto{\pgfqpoint{3.131753in}{1.305762in}}%
\pgfpathlineto{\pgfqpoint{3.125802in}{1.307188in}}%
\pgfpathlineto{\pgfqpoint{3.115324in}{1.315206in}}%
\pgfpathlineto{\pgfqpoint{3.116356in}{1.322582in}}%
\pgfpathlineto{\pgfqpoint{3.121341in}{1.327517in}}%
\pgfpathlineto{\pgfqpoint{3.120128in}{1.335691in}}%
\pgfpathlineto{\pgfqpoint{3.107539in}{1.353570in}}%
\pgfpathlineto{\pgfqpoint{3.110708in}{1.357752in}}%
\pgfpathlineto{\pgfqpoint{3.110807in}{1.366936in}}%
\pgfpathlineto{\pgfqpoint{3.115799in}{1.372179in}}%
\pgfpathlineto{\pgfqpoint{3.123623in}{1.385333in}}%
\pgfpathlineto{\pgfqpoint{3.125203in}{1.394430in}}%
\pgfpathlineto{\pgfqpoint{3.129858in}{1.404228in}}%
\pgfpathlineto{\pgfqpoint{3.137714in}{1.398047in}}%
\pgfpathlineto{\pgfqpoint{3.141594in}{1.398020in}}%
\pgfpathlineto{\pgfqpoint{3.153282in}{1.414127in}}%
\pgfpathlineto{\pgfqpoint{3.158541in}{1.406177in}}%
\pgfpathlineto{\pgfqpoint{3.169602in}{1.394541in}}%
\pgfpathlineto{\pgfqpoint{3.175771in}{1.395430in}}%
\pgfpathlineto{\pgfqpoint{3.161085in}{1.369292in}}%
\pgfpathlineto{\pgfqpoint{3.168135in}{1.369114in}}%
\pgfpathlineto{\pgfqpoint{3.178627in}{1.363629in}}%
\pgfpathlineto{\pgfqpoint{3.176043in}{1.340697in}}%
\pgfpathlineto{\pgfqpoint{3.169983in}{1.342070in}}%
\pgfpathlineto{\pgfqpoint{3.165049in}{1.347998in}}%
\pgfpathlineto{\pgfqpoint{3.157131in}{1.351262in}}%
\pgfpathlineto{\pgfqpoint{3.153524in}{1.337961in}}%
\pgfpathclose%
\pgfusepath{fill}%
\end{pgfscope}%
\begin{pgfscope}%
\pgfpathrectangle{\pgfqpoint{0.100000in}{0.100000in}}{\pgfqpoint{3.608454in}{2.310000in}}%
\pgfusepath{clip}%
\pgfsetbuttcap%
\pgfsetmiterjoin%
\definecolor{currentfill}{rgb}{0.000000,0.427451,0.786275}%
\pgfsetfillcolor{currentfill}%
\pgfsetlinewidth{0.000000pt}%
\definecolor{currentstroke}{rgb}{0.000000,0.000000,0.000000}%
\pgfsetstrokecolor{currentstroke}%
\pgfsetstrokeopacity{0.000000}%
\pgfsetdash{}{0pt}%
\pgfpathmoveto{\pgfqpoint{1.713023in}{1.050837in}}%
\pgfpathlineto{\pgfqpoint{1.710432in}{1.016551in}}%
\pgfpathlineto{\pgfqpoint{1.685686in}{1.018201in}}%
\pgfpathlineto{\pgfqpoint{1.651469in}{1.020799in}}%
\pgfpathlineto{\pgfqpoint{1.617270in}{1.023437in}}%
\pgfpathlineto{\pgfqpoint{1.583688in}{1.026354in}}%
\pgfpathlineto{\pgfqpoint{1.585141in}{1.042675in}}%
\pgfpathlineto{\pgfqpoint{1.589893in}{1.095406in}}%
\pgfpathlineto{\pgfqpoint{1.646284in}{1.091077in}}%
\pgfpathlineto{\pgfqpoint{1.680896in}{1.087866in}}%
\pgfpathlineto{\pgfqpoint{1.678367in}{1.053341in}}%
\pgfpathlineto{\pgfqpoint{1.713023in}{1.050837in}}%
\pgfpathclose%
\pgfusepath{fill}%
\end{pgfscope}%
\begin{pgfscope}%
\pgfpathrectangle{\pgfqpoint{0.100000in}{0.100000in}}{\pgfqpoint{3.608454in}{2.310000in}}%
\pgfusepath{clip}%
\pgfsetbuttcap%
\pgfsetmiterjoin%
\definecolor{currentfill}{rgb}{0.000000,0.470588,0.764706}%
\pgfsetfillcolor{currentfill}%
\pgfsetlinewidth{0.000000pt}%
\definecolor{currentstroke}{rgb}{0.000000,0.000000,0.000000}%
\pgfsetstrokecolor{currentstroke}%
\pgfsetstrokeopacity{0.000000}%
\pgfsetdash{}{0pt}%
\pgfpathmoveto{\pgfqpoint{2.281822in}{0.951221in}}%
\pgfpathlineto{\pgfqpoint{2.271324in}{0.957636in}}%
\pgfpathlineto{\pgfqpoint{2.259184in}{0.957741in}}%
\pgfpathlineto{\pgfqpoint{2.259281in}{0.978738in}}%
\pgfpathlineto{\pgfqpoint{2.258094in}{0.983370in}}%
\pgfpathlineto{\pgfqpoint{2.250959in}{0.985619in}}%
\pgfpathlineto{\pgfqpoint{2.249833in}{0.992475in}}%
\pgfpathlineto{\pgfqpoint{2.245351in}{0.995915in}}%
\pgfpathlineto{\pgfqpoint{2.238537in}{0.995926in}}%
\pgfpathlineto{\pgfqpoint{2.238992in}{1.006643in}}%
\pgfpathlineto{\pgfqpoint{2.218257in}{1.006289in}}%
\pgfpathlineto{\pgfqpoint{2.216940in}{1.014382in}}%
\pgfpathlineto{\pgfqpoint{2.226237in}{1.022520in}}%
\pgfpathlineto{\pgfqpoint{2.241035in}{1.037964in}}%
\pgfpathlineto{\pgfqpoint{2.245679in}{1.038542in}}%
\pgfpathlineto{\pgfqpoint{2.245419in}{1.061527in}}%
\pgfpathlineto{\pgfqpoint{2.269359in}{1.061552in}}%
\pgfpathlineto{\pgfqpoint{2.269431in}{1.054652in}}%
\pgfpathlineto{\pgfqpoint{2.293330in}{1.054914in}}%
\pgfpathlineto{\pgfqpoint{2.293540in}{1.031438in}}%
\pgfpathlineto{\pgfqpoint{2.300760in}{1.031780in}}%
\pgfpathlineto{\pgfqpoint{2.309489in}{1.028555in}}%
\pgfpathlineto{\pgfqpoint{2.314200in}{1.029274in}}%
\pgfpathlineto{\pgfqpoint{2.314414in}{1.020654in}}%
\pgfpathlineto{\pgfqpoint{2.321574in}{1.020760in}}%
\pgfpathlineto{\pgfqpoint{2.321823in}{1.001132in}}%
\pgfpathlineto{\pgfqpoint{2.326016in}{0.997543in}}%
\pgfpathlineto{\pgfqpoint{2.326278in}{0.991006in}}%
\pgfpathlineto{\pgfqpoint{2.322083in}{0.986111in}}%
\pgfpathlineto{\pgfqpoint{2.289234in}{0.985811in}}%
\pgfpathlineto{\pgfqpoint{2.289533in}{0.964998in}}%
\pgfpathlineto{\pgfqpoint{2.288597in}{0.951635in}}%
\pgfpathlineto{\pgfqpoint{2.281822in}{0.951221in}}%
\pgfpathclose%
\pgfusepath{fill}%
\end{pgfscope}%
\begin{pgfscope}%
\pgfpathrectangle{\pgfqpoint{0.100000in}{0.100000in}}{\pgfqpoint{3.608454in}{2.310000in}}%
\pgfusepath{clip}%
\pgfsetbuttcap%
\pgfsetmiterjoin%
\definecolor{currentfill}{rgb}{0.000000,0.611765,0.694118}%
\pgfsetfillcolor{currentfill}%
\pgfsetlinewidth{0.000000pt}%
\definecolor{currentstroke}{rgb}{0.000000,0.000000,0.000000}%
\pgfsetstrokecolor{currentstroke}%
\pgfsetstrokeopacity{0.000000}%
\pgfsetdash{}{0pt}%
\pgfpathmoveto{\pgfqpoint{1.553831in}{1.907730in}}%
\pgfpathlineto{\pgfqpoint{1.549754in}{1.873361in}}%
\pgfpathlineto{\pgfqpoint{1.546917in}{1.873677in}}%
\pgfpathlineto{\pgfqpoint{1.543825in}{1.845887in}}%
\pgfpathlineto{\pgfqpoint{1.494987in}{1.851589in}}%
\pgfpathlineto{\pgfqpoint{1.488701in}{1.851762in}}%
\pgfpathlineto{\pgfqpoint{1.475384in}{1.853420in}}%
\pgfpathlineto{\pgfqpoint{1.477206in}{1.868092in}}%
\pgfpathlineto{\pgfqpoint{1.449171in}{1.871656in}}%
\pgfpathlineto{\pgfqpoint{1.453065in}{1.884941in}}%
\pgfpathlineto{\pgfqpoint{1.456350in}{1.910291in}}%
\pgfpathlineto{\pgfqpoint{1.446358in}{1.912489in}}%
\pgfpathlineto{\pgfqpoint{1.446859in}{1.927356in}}%
\pgfpathlineto{\pgfqpoint{1.434747in}{1.935937in}}%
\pgfpathlineto{\pgfqpoint{1.421111in}{1.937794in}}%
\pgfpathlineto{\pgfqpoint{1.422053in}{1.944628in}}%
\pgfpathlineto{\pgfqpoint{1.417308in}{1.945271in}}%
\pgfpathlineto{\pgfqpoint{1.421242in}{1.955525in}}%
\pgfpathlineto{\pgfqpoint{1.414681in}{1.966931in}}%
\pgfpathlineto{\pgfqpoint{1.410750in}{1.980307in}}%
\pgfpathlineto{\pgfqpoint{1.406757in}{1.983167in}}%
\pgfpathlineto{\pgfqpoint{1.409028in}{1.997695in}}%
\pgfpathlineto{\pgfqpoint{1.407997in}{2.003940in}}%
\pgfpathlineto{\pgfqpoint{1.405583in}{2.011993in}}%
\pgfpathlineto{\pgfqpoint{1.469044in}{2.003996in}}%
\pgfpathlineto{\pgfqpoint{1.468884in}{2.002862in}}%
\pgfpathlineto{\pgfqpoint{1.503208in}{1.998589in}}%
\pgfpathlineto{\pgfqpoint{1.501201in}{1.997658in}}%
\pgfpathlineto{\pgfqpoint{1.497843in}{1.970127in}}%
\pgfpathlineto{\pgfqpoint{1.495943in}{1.970351in}}%
\pgfpathlineto{\pgfqpoint{1.492673in}{1.942949in}}%
\pgfpathlineto{\pgfqpoint{1.490580in}{1.943193in}}%
\pgfpathlineto{\pgfqpoint{1.487149in}{1.915660in}}%
\pgfpathlineto{\pgfqpoint{1.553831in}{1.907730in}}%
\pgfpathclose%
\pgfusepath{fill}%
\end{pgfscope}%
\begin{pgfscope}%
\pgfpathrectangle{\pgfqpoint{0.100000in}{0.100000in}}{\pgfqpoint{3.608454in}{2.310000in}}%
\pgfusepath{clip}%
\pgfsetbuttcap%
\pgfsetmiterjoin%
\definecolor{currentfill}{rgb}{0.000000,0.776471,0.611765}%
\pgfsetfillcolor{currentfill}%
\pgfsetlinewidth{0.000000pt}%
\definecolor{currentstroke}{rgb}{0.000000,0.000000,0.000000}%
\pgfsetstrokecolor{currentstroke}%
\pgfsetstrokeopacity{0.000000}%
\pgfsetdash{}{0pt}%
\pgfpathmoveto{\pgfqpoint{2.978722in}{0.804662in}}%
\pgfpathlineto{\pgfqpoint{2.970476in}{0.810318in}}%
\pgfpathlineto{\pgfqpoint{2.953500in}{0.817959in}}%
\pgfpathlineto{\pgfqpoint{2.929182in}{0.813904in}}%
\pgfpathlineto{\pgfqpoint{2.926243in}{0.822499in}}%
\pgfpathlineto{\pgfqpoint{2.917347in}{0.821092in}}%
\pgfpathlineto{\pgfqpoint{2.904439in}{0.824809in}}%
\pgfpathlineto{\pgfqpoint{2.900223in}{0.831517in}}%
\pgfpathlineto{\pgfqpoint{2.895939in}{0.833978in}}%
\pgfpathlineto{\pgfqpoint{2.894217in}{0.842012in}}%
\pgfpathlineto{\pgfqpoint{2.891100in}{0.843787in}}%
\pgfpathlineto{\pgfqpoint{2.893436in}{0.849951in}}%
\pgfpathlineto{\pgfqpoint{2.887187in}{0.854824in}}%
\pgfpathlineto{\pgfqpoint{2.890867in}{0.862703in}}%
\pgfpathlineto{\pgfqpoint{2.898491in}{0.872568in}}%
\pgfpathlineto{\pgfqpoint{2.901077in}{0.870542in}}%
\pgfpathlineto{\pgfqpoint{2.913824in}{0.850086in}}%
\pgfpathlineto{\pgfqpoint{2.922187in}{0.854988in}}%
\pgfpathlineto{\pgfqpoint{2.930056in}{0.865279in}}%
\pgfpathlineto{\pgfqpoint{2.928863in}{0.867621in}}%
\pgfpathlineto{\pgfqpoint{2.932829in}{0.881990in}}%
\pgfpathlineto{\pgfqpoint{2.943431in}{0.882150in}}%
\pgfpathlineto{\pgfqpoint{2.951258in}{0.878733in}}%
\pgfpathlineto{\pgfqpoint{2.950142in}{0.872561in}}%
\pgfpathlineto{\pgfqpoint{2.954746in}{0.867311in}}%
\pgfpathlineto{\pgfqpoint{2.962520in}{0.870538in}}%
\pgfpathlineto{\pgfqpoint{2.967645in}{0.858971in}}%
\pgfpathlineto{\pgfqpoint{2.967579in}{0.838638in}}%
\pgfpathlineto{\pgfqpoint{2.975754in}{0.837437in}}%
\pgfpathlineto{\pgfqpoint{2.979065in}{0.834250in}}%
\pgfpathlineto{\pgfqpoint{2.975294in}{0.828495in}}%
\pgfpathlineto{\pgfqpoint{2.978722in}{0.804662in}}%
\pgfpathclose%
\pgfusepath{fill}%
\end{pgfscope}%
\begin{pgfscope}%
\pgfpathrectangle{\pgfqpoint{0.100000in}{0.100000in}}{\pgfqpoint{3.608454in}{2.310000in}}%
\pgfusepath{clip}%
\pgfsetbuttcap%
\pgfsetmiterjoin%
\definecolor{currentfill}{rgb}{0.000000,0.458824,0.770588}%
\pgfsetfillcolor{currentfill}%
\pgfsetlinewidth{0.000000pt}%
\definecolor{currentstroke}{rgb}{0.000000,0.000000,0.000000}%
\pgfsetstrokecolor{currentstroke}%
\pgfsetstrokeopacity{0.000000}%
\pgfsetdash{}{0pt}%
\pgfpathmoveto{\pgfqpoint{3.269493in}{1.294748in}}%
\pgfpathlineto{\pgfqpoint{3.266535in}{1.293082in}}%
\pgfpathlineto{\pgfqpoint{3.248223in}{1.267423in}}%
\pgfpathlineto{\pgfqpoint{3.248646in}{1.259422in}}%
\pgfpathlineto{\pgfqpoint{3.221614in}{1.254572in}}%
\pgfpathlineto{\pgfqpoint{3.214106in}{1.265774in}}%
\pgfpathlineto{\pgfqpoint{3.240057in}{1.290021in}}%
\pgfpathlineto{\pgfqpoint{3.236491in}{1.297792in}}%
\pgfpathlineto{\pgfqpoint{3.224896in}{1.300661in}}%
\pgfpathlineto{\pgfqpoint{3.234147in}{1.313361in}}%
\pgfpathlineto{\pgfqpoint{3.240056in}{1.312499in}}%
\pgfpathlineto{\pgfqpoint{3.236059in}{1.328225in}}%
\pgfpathlineto{\pgfqpoint{3.249638in}{1.333200in}}%
\pgfpathlineto{\pgfqpoint{3.248164in}{1.344820in}}%
\pgfpathlineto{\pgfqpoint{3.240232in}{1.353321in}}%
\pgfpathlineto{\pgfqpoint{3.244523in}{1.358215in}}%
\pgfpathlineto{\pgfqpoint{3.250135in}{1.356031in}}%
\pgfpathlineto{\pgfqpoint{3.255578in}{1.348500in}}%
\pgfpathlineto{\pgfqpoint{3.262571in}{1.349107in}}%
\pgfpathlineto{\pgfqpoint{3.267940in}{1.341269in}}%
\pgfpathlineto{\pgfqpoint{3.273671in}{1.340685in}}%
\pgfpathlineto{\pgfqpoint{3.276753in}{1.332527in}}%
\pgfpathlineto{\pgfqpoint{3.273954in}{1.327806in}}%
\pgfpathlineto{\pgfqpoint{3.268757in}{1.331100in}}%
\pgfpathlineto{\pgfqpoint{3.270994in}{1.316465in}}%
\pgfpathlineto{\pgfqpoint{3.280061in}{1.309564in}}%
\pgfpathlineto{\pgfqpoint{3.278895in}{1.302449in}}%
\pgfpathlineto{\pgfqpoint{3.274672in}{1.301154in}}%
\pgfpathlineto{\pgfqpoint{3.269493in}{1.294748in}}%
\pgfpathclose%
\pgfusepath{fill}%
\end{pgfscope}%
\begin{pgfscope}%
\pgfpathrectangle{\pgfqpoint{0.100000in}{0.100000in}}{\pgfqpoint{3.608454in}{2.310000in}}%
\pgfusepath{clip}%
\pgfsetbuttcap%
\pgfsetmiterjoin%
\definecolor{currentfill}{rgb}{0.000000,0.447059,0.776471}%
\pgfsetfillcolor{currentfill}%
\pgfsetlinewidth{0.000000pt}%
\definecolor{currentstroke}{rgb}{0.000000,0.000000,0.000000}%
\pgfsetstrokecolor{currentstroke}%
\pgfsetstrokeopacity{0.000000}%
\pgfsetdash{}{0pt}%
\pgfpathmoveto{\pgfqpoint{2.844033in}{0.665474in}}%
\pgfpathlineto{\pgfqpoint{2.844132in}{0.660367in}}%
\pgfpathlineto{\pgfqpoint{2.835832in}{0.661752in}}%
\pgfpathlineto{\pgfqpoint{2.827136in}{0.657037in}}%
\pgfpathlineto{\pgfqpoint{2.809853in}{0.644061in}}%
\pgfpathlineto{\pgfqpoint{2.799185in}{0.640400in}}%
\pgfpathlineto{\pgfqpoint{2.793287in}{0.640173in}}%
\pgfpathlineto{\pgfqpoint{2.783205in}{0.636006in}}%
\pgfpathlineto{\pgfqpoint{2.779514in}{0.637469in}}%
\pgfpathlineto{\pgfqpoint{2.779023in}{0.645884in}}%
\pgfpathlineto{\pgfqpoint{2.769443in}{0.656196in}}%
\pgfpathlineto{\pgfqpoint{2.764564in}{0.656808in}}%
\pgfpathlineto{\pgfqpoint{2.757459in}{0.662818in}}%
\pgfpathlineto{\pgfqpoint{2.742727in}{0.670681in}}%
\pgfpathlineto{\pgfqpoint{2.727447in}{0.677658in}}%
\pgfpathlineto{\pgfqpoint{2.726751in}{0.687235in}}%
\pgfpathlineto{\pgfqpoint{2.735873in}{0.692265in}}%
\pgfpathlineto{\pgfqpoint{2.760979in}{0.694551in}}%
\pgfpathlineto{\pgfqpoint{2.760006in}{0.704933in}}%
\pgfpathlineto{\pgfqpoint{2.766984in}{0.705619in}}%
\pgfpathlineto{\pgfqpoint{2.769691in}{0.676534in}}%
\pgfpathlineto{\pgfqpoint{2.787376in}{0.676778in}}%
\pgfpathlineto{\pgfqpoint{2.788514in}{0.665867in}}%
\pgfpathlineto{\pgfqpoint{2.797575in}{0.664470in}}%
\pgfpathlineto{\pgfqpoint{2.829477in}{0.667841in}}%
\pgfpathlineto{\pgfqpoint{2.844033in}{0.665474in}}%
\pgfpathclose%
\pgfusepath{fill}%
\end{pgfscope}%
\begin{pgfscope}%
\pgfpathrectangle{\pgfqpoint{0.100000in}{0.100000in}}{\pgfqpoint{3.608454in}{2.310000in}}%
\pgfusepath{clip}%
\pgfsetbuttcap%
\pgfsetmiterjoin%
\definecolor{currentfill}{rgb}{0.000000,0.737255,0.631373}%
\pgfsetfillcolor{currentfill}%
\pgfsetlinewidth{0.000000pt}%
\definecolor{currentstroke}{rgb}{0.000000,0.000000,0.000000}%
\pgfsetstrokecolor{currentstroke}%
\pgfsetstrokeopacity{0.000000}%
\pgfsetdash{}{0pt}%
\pgfpathmoveto{\pgfqpoint{3.074241in}{1.252094in}}%
\pgfpathlineto{\pgfqpoint{3.073773in}{1.226802in}}%
\pgfpathlineto{\pgfqpoint{3.054182in}{1.223554in}}%
\pgfpathlineto{\pgfqpoint{3.037358in}{1.220901in}}%
\pgfpathlineto{\pgfqpoint{3.017299in}{1.218774in}}%
\pgfpathlineto{\pgfqpoint{3.015913in}{1.223529in}}%
\pgfpathlineto{\pgfqpoint{3.025017in}{1.229246in}}%
\pgfpathlineto{\pgfqpoint{3.024683in}{1.232048in}}%
\pgfpathlineto{\pgfqpoint{3.032148in}{1.244073in}}%
\pgfpathlineto{\pgfqpoint{3.036908in}{1.247132in}}%
\pgfpathlineto{\pgfqpoint{3.067580in}{1.245522in}}%
\pgfpathlineto{\pgfqpoint{3.074241in}{1.252094in}}%
\pgfpathclose%
\pgfusepath{fill}%
\end{pgfscope}%
\begin{pgfscope}%
\pgfpathrectangle{\pgfqpoint{0.100000in}{0.100000in}}{\pgfqpoint{3.608454in}{2.310000in}}%
\pgfusepath{clip}%
\pgfsetbuttcap%
\pgfsetmiterjoin%
\definecolor{currentfill}{rgb}{0.000000,0.721569,0.639216}%
\pgfsetfillcolor{currentfill}%
\pgfsetlinewidth{0.000000pt}%
\definecolor{currentstroke}{rgb}{0.000000,0.000000,0.000000}%
\pgfsetstrokecolor{currentstroke}%
\pgfsetstrokeopacity{0.000000}%
\pgfsetdash{}{0pt}%
\pgfpathmoveto{\pgfqpoint{1.523829in}{1.020634in}}%
\pgfpathlineto{\pgfqpoint{1.519945in}{0.979398in}}%
\pgfpathlineto{\pgfqpoint{1.506226in}{0.980812in}}%
\pgfpathlineto{\pgfqpoint{1.505497in}{0.973876in}}%
\pgfpathlineto{\pgfqpoint{1.464338in}{0.978551in}}%
\pgfpathlineto{\pgfqpoint{1.465068in}{0.985473in}}%
\pgfpathlineto{\pgfqpoint{1.458090in}{0.986273in}}%
\pgfpathlineto{\pgfqpoint{1.460404in}{1.006662in}}%
\pgfpathlineto{\pgfqpoint{1.432998in}{1.009699in}}%
\pgfpathlineto{\pgfqpoint{1.435358in}{1.030206in}}%
\pgfpathlineto{\pgfqpoint{1.436822in}{1.030070in}}%
\pgfpathlineto{\pgfqpoint{1.440664in}{1.064305in}}%
\pgfpathlineto{\pgfqpoint{1.442195in}{1.078025in}}%
\pgfpathlineto{\pgfqpoint{1.517207in}{1.069928in}}%
\pgfpathlineto{\pgfqpoint{1.516667in}{1.064131in}}%
\pgfpathlineto{\pgfqpoint{1.512139in}{1.021816in}}%
\pgfpathlineto{\pgfqpoint{1.523829in}{1.020634in}}%
\pgfpathclose%
\pgfusepath{fill}%
\end{pgfscope}%
\begin{pgfscope}%
\pgfpathrectangle{\pgfqpoint{0.100000in}{0.100000in}}{\pgfqpoint{3.608454in}{2.310000in}}%
\pgfusepath{clip}%
\pgfsetbuttcap%
\pgfsetmiterjoin%
\definecolor{currentfill}{rgb}{0.000000,0.352941,0.823529}%
\pgfsetfillcolor{currentfill}%
\pgfsetlinewidth{0.000000pt}%
\definecolor{currentstroke}{rgb}{0.000000,0.000000,0.000000}%
\pgfsetstrokecolor{currentstroke}%
\pgfsetstrokeopacity{0.000000}%
\pgfsetdash{}{0pt}%
\pgfpathmoveto{\pgfqpoint{3.530808in}{1.730215in}}%
\pgfpathlineto{\pgfqpoint{3.529206in}{1.730530in}}%
\pgfpathlineto{\pgfqpoint{3.531215in}{1.735382in}}%
\pgfpathlineto{\pgfqpoint{3.526862in}{1.745686in}}%
\pgfpathlineto{\pgfqpoint{3.518613in}{1.736187in}}%
\pgfpathlineto{\pgfqpoint{3.513375in}{1.745585in}}%
\pgfpathlineto{\pgfqpoint{3.499874in}{1.758602in}}%
\pgfpathlineto{\pgfqpoint{3.495434in}{1.768247in}}%
\pgfpathlineto{\pgfqpoint{3.481572in}{1.755490in}}%
\pgfpathlineto{\pgfqpoint{3.479909in}{1.757791in}}%
\pgfpathlineto{\pgfqpoint{3.456367in}{1.750710in}}%
\pgfpathlineto{\pgfqpoint{3.436777in}{1.747450in}}%
\pgfpathlineto{\pgfqpoint{3.434236in}{1.757543in}}%
\pgfpathlineto{\pgfqpoint{3.426466in}{1.757361in}}%
\pgfpathlineto{\pgfqpoint{3.427810in}{1.762986in}}%
\pgfpathlineto{\pgfqpoint{3.420462in}{1.768872in}}%
\pgfpathlineto{\pgfqpoint{3.418010in}{1.787350in}}%
\pgfpathlineto{\pgfqpoint{3.419015in}{1.794601in}}%
\pgfpathlineto{\pgfqpoint{3.414849in}{1.798333in}}%
\pgfpathlineto{\pgfqpoint{3.434992in}{1.802755in}}%
\pgfpathlineto{\pgfqpoint{3.471016in}{1.810988in}}%
\pgfpathlineto{\pgfqpoint{3.477550in}{1.822939in}}%
\pgfpathlineto{\pgfqpoint{3.487553in}{1.830826in}}%
\pgfpathlineto{\pgfqpoint{3.494171in}{1.831598in}}%
\pgfpathlineto{\pgfqpoint{3.501460in}{1.818237in}}%
\pgfpathlineto{\pgfqpoint{3.511382in}{1.816812in}}%
\pgfpathlineto{\pgfqpoint{3.507584in}{1.811929in}}%
\pgfpathlineto{\pgfqpoint{3.498548in}{1.805841in}}%
\pgfpathlineto{\pgfqpoint{3.500415in}{1.803471in}}%
\pgfpathlineto{\pgfqpoint{3.494052in}{1.793429in}}%
\pgfpathlineto{\pgfqpoint{3.496118in}{1.786298in}}%
\pgfpathlineto{\pgfqpoint{3.502227in}{1.787338in}}%
\pgfpathlineto{\pgfqpoint{3.513856in}{1.782235in}}%
\pgfpathlineto{\pgfqpoint{3.520466in}{1.774625in}}%
\pgfpathlineto{\pgfqpoint{3.519312in}{1.765836in}}%
\pgfpathlineto{\pgfqpoint{3.529828in}{1.763255in}}%
\pgfpathlineto{\pgfqpoint{3.531850in}{1.755199in}}%
\pgfpathlineto{\pgfqpoint{3.540059in}{1.751718in}}%
\pgfpathlineto{\pgfqpoint{3.547764in}{1.752941in}}%
\pgfpathlineto{\pgfqpoint{3.566780in}{1.763298in}}%
\pgfpathlineto{\pgfqpoint{3.569480in}{1.754165in}}%
\pgfpathlineto{\pgfqpoint{3.542220in}{1.740804in}}%
\pgfpathlineto{\pgfqpoint{3.540912in}{1.736515in}}%
\pgfpathlineto{\pgfqpoint{3.530808in}{1.730215in}}%
\pgfpathclose%
\pgfusepath{fill}%
\end{pgfscope}%
\begin{pgfscope}%
\pgfpathrectangle{\pgfqpoint{0.100000in}{0.100000in}}{\pgfqpoint{3.608454in}{2.310000in}}%
\pgfusepath{clip}%
\pgfsetbuttcap%
\pgfsetmiterjoin%
\definecolor{currentfill}{rgb}{0.000000,0.827451,0.586275}%
\pgfsetfillcolor{currentfill}%
\pgfsetlinewidth{0.000000pt}%
\definecolor{currentstroke}{rgb}{0.000000,0.000000,0.000000}%
\pgfsetstrokecolor{currentstroke}%
\pgfsetstrokeopacity{0.000000}%
\pgfsetdash{}{0pt}%
\pgfpathmoveto{\pgfqpoint{0.737385in}{0.650268in}}%
\pgfpathlineto{\pgfqpoint{0.738818in}{0.653776in}}%
\pgfpathlineto{\pgfqpoint{0.738884in}{0.657312in}}%
\pgfpathlineto{\pgfqpoint{0.737658in}{0.662474in}}%
\pgfpathlineto{\pgfqpoint{0.738418in}{0.665625in}}%
\pgfpathlineto{\pgfqpoint{0.736741in}{0.671963in}}%
\pgfpathlineto{\pgfqpoint{0.734805in}{0.674605in}}%
\pgfpathlineto{\pgfqpoint{0.729012in}{0.675764in}}%
\pgfpathlineto{\pgfqpoint{0.728310in}{0.678102in}}%
\pgfpathlineto{\pgfqpoint{0.728866in}{0.682342in}}%
\pgfpathlineto{\pgfqpoint{0.731069in}{0.684627in}}%
\pgfpathlineto{\pgfqpoint{0.733927in}{0.685526in}}%
\pgfpathlineto{\pgfqpoint{0.736427in}{0.687461in}}%
\pgfpathlineto{\pgfqpoint{0.737267in}{0.686485in}}%
\pgfpathlineto{\pgfqpoint{0.736221in}{0.683768in}}%
\pgfpathlineto{\pgfqpoint{0.738962in}{0.676360in}}%
\pgfpathlineto{\pgfqpoint{0.740767in}{0.675472in}}%
\pgfpathlineto{\pgfqpoint{0.746631in}{0.673961in}}%
\pgfpathlineto{\pgfqpoint{0.747962in}{0.671116in}}%
\pgfpathlineto{\pgfqpoint{0.747857in}{0.668032in}}%
\pgfpathlineto{\pgfqpoint{0.745828in}{0.665738in}}%
\pgfpathlineto{\pgfqpoint{0.747273in}{0.661310in}}%
\pgfpathlineto{\pgfqpoint{0.747476in}{0.658092in}}%
\pgfpathlineto{\pgfqpoint{0.749575in}{0.653902in}}%
\pgfpathlineto{\pgfqpoint{0.752259in}{0.652120in}}%
\pgfpathlineto{\pgfqpoint{0.753791in}{0.649042in}}%
\pgfpathlineto{\pgfqpoint{0.752833in}{0.647643in}}%
\pgfpathlineto{\pgfqpoint{0.749601in}{0.646731in}}%
\pgfpathlineto{\pgfqpoint{0.748012in}{0.649373in}}%
\pgfpathlineto{\pgfqpoint{0.745554in}{0.651542in}}%
\pgfpathlineto{\pgfqpoint{0.742391in}{0.652099in}}%
\pgfpathlineto{\pgfqpoint{0.737385in}{0.650268in}}%
\pgfpathclose%
\pgfusepath{fill}%
\end{pgfscope}%
\begin{pgfscope}%
\pgfpathrectangle{\pgfqpoint{0.100000in}{0.100000in}}{\pgfqpoint{3.608454in}{2.310000in}}%
\pgfusepath{clip}%
\pgfsetbuttcap%
\pgfsetmiterjoin%
\definecolor{currentfill}{rgb}{0.000000,0.827451,0.586275}%
\pgfsetfillcolor{currentfill}%
\pgfsetlinewidth{0.000000pt}%
\definecolor{currentstroke}{rgb}{0.000000,0.000000,0.000000}%
\pgfsetstrokecolor{currentstroke}%
\pgfsetstrokeopacity{0.000000}%
\pgfsetdash{}{0pt}%
\pgfpathmoveto{\pgfqpoint{0.863943in}{0.679699in}}%
\pgfpathlineto{\pgfqpoint{0.862419in}{0.676643in}}%
\pgfpathlineto{\pgfqpoint{0.855048in}{0.668896in}}%
\pgfpathlineto{\pgfqpoint{0.855680in}{0.668231in}}%
\pgfpathlineto{\pgfqpoint{0.848280in}{0.660517in}}%
\pgfpathlineto{\pgfqpoint{0.850812in}{0.658129in}}%
\pgfpathlineto{\pgfqpoint{0.847118in}{0.654238in}}%
\pgfpathlineto{\pgfqpoint{0.849066in}{0.652382in}}%
\pgfpathlineto{\pgfqpoint{0.845457in}{0.648542in}}%
\pgfpathlineto{\pgfqpoint{0.853614in}{0.640940in}}%
\pgfpathlineto{\pgfqpoint{0.870840in}{0.625356in}}%
\pgfpathlineto{\pgfqpoint{0.882807in}{0.614773in}}%
\pgfpathlineto{\pgfqpoint{0.884580in}{0.616759in}}%
\pgfpathlineto{\pgfqpoint{0.886576in}{0.615036in}}%
\pgfpathlineto{\pgfqpoint{0.888579in}{0.613334in}}%
\pgfpathlineto{\pgfqpoint{0.883388in}{0.607238in}}%
\pgfpathlineto{\pgfqpoint{0.881527in}{0.608877in}}%
\pgfpathlineto{\pgfqpoint{0.874796in}{0.601023in}}%
\pgfpathlineto{\pgfqpoint{0.872825in}{0.602697in}}%
\pgfpathlineto{\pgfqpoint{0.869420in}{0.598690in}}%
\pgfpathlineto{\pgfqpoint{0.867382in}{0.600383in}}%
\pgfpathlineto{\pgfqpoint{0.853750in}{0.584249in}}%
\pgfpathlineto{\pgfqpoint{0.856024in}{0.582334in}}%
\pgfpathlineto{\pgfqpoint{0.850691in}{0.576622in}}%
\pgfpathlineto{\pgfqpoint{0.848929in}{0.578173in}}%
\pgfpathlineto{\pgfqpoint{0.847165in}{0.576109in}}%
\pgfpathlineto{\pgfqpoint{0.845387in}{0.577604in}}%
\pgfpathlineto{\pgfqpoint{0.843712in}{0.575563in}}%
\pgfpathlineto{\pgfqpoint{0.839724in}{0.579022in}}%
\pgfpathlineto{\pgfqpoint{0.838193in}{0.577309in}}%
\pgfpathlineto{\pgfqpoint{0.836205in}{0.579038in}}%
\pgfpathlineto{\pgfqpoint{0.831146in}{0.572881in}}%
\pgfpathlineto{\pgfqpoint{0.829339in}{0.574413in}}%
\pgfpathlineto{\pgfqpoint{0.824219in}{0.568707in}}%
\pgfpathlineto{\pgfqpoint{0.812726in}{0.579011in}}%
\pgfpathlineto{\pgfqpoint{0.814280in}{0.580959in}}%
\pgfpathlineto{\pgfqpoint{0.812542in}{0.582561in}}%
\pgfpathlineto{\pgfqpoint{0.814282in}{0.584612in}}%
\pgfpathlineto{\pgfqpoint{0.812396in}{0.586335in}}%
\pgfpathlineto{\pgfqpoint{0.814106in}{0.588275in}}%
\pgfpathlineto{\pgfqpoint{0.812082in}{0.590252in}}%
\pgfpathlineto{\pgfqpoint{0.818724in}{0.592259in}}%
\pgfpathlineto{\pgfqpoint{0.818875in}{0.594792in}}%
\pgfpathlineto{\pgfqpoint{0.816853in}{0.596475in}}%
\pgfpathlineto{\pgfqpoint{0.818540in}{0.598214in}}%
\pgfpathlineto{\pgfqpoint{0.821571in}{0.596416in}}%
\pgfpathlineto{\pgfqpoint{0.820998in}{0.592054in}}%
\pgfpathlineto{\pgfqpoint{0.823073in}{0.586861in}}%
\pgfpathlineto{\pgfqpoint{0.829242in}{0.582849in}}%
\pgfpathlineto{\pgfqpoint{0.831926in}{0.582649in}}%
\pgfpathlineto{\pgfqpoint{0.836286in}{0.583592in}}%
\pgfpathlineto{\pgfqpoint{0.839848in}{0.584937in}}%
\pgfpathlineto{\pgfqpoint{0.841321in}{0.586303in}}%
\pgfpathlineto{\pgfqpoint{0.844401in}{0.593367in}}%
\pgfpathlineto{\pgfqpoint{0.847330in}{0.597336in}}%
\pgfpathlineto{\pgfqpoint{0.848129in}{0.603625in}}%
\pgfpathlineto{\pgfqpoint{0.845939in}{0.606018in}}%
\pgfpathlineto{\pgfqpoint{0.848681in}{0.607604in}}%
\pgfpathlineto{\pgfqpoint{0.851564in}{0.603948in}}%
\pgfpathlineto{\pgfqpoint{0.856618in}{0.604295in}}%
\pgfpathlineto{\pgfqpoint{0.858912in}{0.606588in}}%
\pgfpathlineto{\pgfqpoint{0.860360in}{0.609378in}}%
\pgfpathlineto{\pgfqpoint{0.859005in}{0.614759in}}%
\pgfpathlineto{\pgfqpoint{0.854769in}{0.613300in}}%
\pgfpathlineto{\pgfqpoint{0.852074in}{0.613498in}}%
\pgfpathlineto{\pgfqpoint{0.851004in}{0.616189in}}%
\pgfpathlineto{\pgfqpoint{0.847338in}{0.616442in}}%
\pgfpathlineto{\pgfqpoint{0.843639in}{0.617453in}}%
\pgfpathlineto{\pgfqpoint{0.840644in}{0.617353in}}%
\pgfpathlineto{\pgfqpoint{0.831603in}{0.615210in}}%
\pgfpathlineto{\pgfqpoint{0.833848in}{0.619214in}}%
\pgfpathlineto{\pgfqpoint{0.834393in}{0.625065in}}%
\pgfpathlineto{\pgfqpoint{0.831455in}{0.625440in}}%
\pgfpathlineto{\pgfqpoint{0.831880in}{0.622625in}}%
\pgfpathlineto{\pgfqpoint{0.829678in}{0.621808in}}%
\pgfpathlineto{\pgfqpoint{0.829726in}{0.624782in}}%
\pgfpathlineto{\pgfqpoint{0.828624in}{0.627896in}}%
\pgfpathlineto{\pgfqpoint{0.822394in}{0.633936in}}%
\pgfpathlineto{\pgfqpoint{0.815333in}{0.636154in}}%
\pgfpathlineto{\pgfqpoint{0.813484in}{0.637133in}}%
\pgfpathlineto{\pgfqpoint{0.805786in}{0.649884in}}%
\pgfpathlineto{\pgfqpoint{0.804377in}{0.652802in}}%
\pgfpathlineto{\pgfqpoint{0.804337in}{0.655328in}}%
\pgfpathlineto{\pgfqpoint{0.806949in}{0.658877in}}%
\pgfpathlineto{\pgfqpoint{0.809295in}{0.660793in}}%
\pgfpathlineto{\pgfqpoint{0.809700in}{0.668848in}}%
\pgfpathlineto{\pgfqpoint{0.813648in}{0.667058in}}%
\pgfpathlineto{\pgfqpoint{0.818571in}{0.669230in}}%
\pgfpathlineto{\pgfqpoint{0.813780in}{0.677479in}}%
\pgfpathlineto{\pgfqpoint{0.811754in}{0.680104in}}%
\pgfpathlineto{\pgfqpoint{0.810600in}{0.686441in}}%
\pgfpathlineto{\pgfqpoint{0.809362in}{0.690533in}}%
\pgfpathlineto{\pgfqpoint{0.811858in}{0.693024in}}%
\pgfpathlineto{\pgfqpoint{0.822834in}{0.689897in}}%
\pgfpathlineto{\pgfqpoint{0.823528in}{0.687901in}}%
\pgfpathlineto{\pgfqpoint{0.827552in}{0.687211in}}%
\pgfpathlineto{\pgfqpoint{0.827406in}{0.689342in}}%
\pgfpathlineto{\pgfqpoint{0.838119in}{0.687842in}}%
\pgfpathlineto{\pgfqpoint{0.840746in}{0.682200in}}%
\pgfpathlineto{\pgfqpoint{0.842608in}{0.680338in}}%
\pgfpathlineto{\pgfqpoint{0.844777in}{0.680608in}}%
\pgfpathlineto{\pgfqpoint{0.843442in}{0.683454in}}%
\pgfpathlineto{\pgfqpoint{0.846822in}{0.686101in}}%
\pgfpathlineto{\pgfqpoint{0.855456in}{0.683423in}}%
\pgfpathlineto{\pgfqpoint{0.860934in}{0.681377in}}%
\pgfpathlineto{\pgfqpoint{0.863943in}{0.679699in}}%
\pgfpathclose%
\pgfusepath{fill}%
\end{pgfscope}%
\begin{pgfscope}%
\pgfpathrectangle{\pgfqpoint{0.100000in}{0.100000in}}{\pgfqpoint{3.608454in}{2.310000in}}%
\pgfusepath{clip}%
\pgfsetbuttcap%
\pgfsetmiterjoin%
\definecolor{currentfill}{rgb}{0.000000,0.796078,0.601961}%
\pgfsetfillcolor{currentfill}%
\pgfsetlinewidth{0.000000pt}%
\definecolor{currentstroke}{rgb}{0.000000,0.000000,0.000000}%
\pgfsetstrokecolor{currentstroke}%
\pgfsetstrokeopacity{0.000000}%
\pgfsetdash{}{0pt}%
\pgfpathmoveto{\pgfqpoint{1.805743in}{1.664546in}}%
\pgfpathlineto{\pgfqpoint{1.746207in}{1.668241in}}%
\pgfpathlineto{\pgfqpoint{1.748309in}{1.699187in}}%
\pgfpathlineto{\pgfqpoint{1.750543in}{1.731001in}}%
\pgfpathlineto{\pgfqpoint{1.760148in}{1.734344in}}%
\pgfpathlineto{\pgfqpoint{1.768321in}{1.732235in}}%
\pgfpathlineto{\pgfqpoint{1.775409in}{1.724494in}}%
\pgfpathlineto{\pgfqpoint{1.787775in}{1.725501in}}%
\pgfpathlineto{\pgfqpoint{1.793255in}{1.722620in}}%
\pgfpathlineto{\pgfqpoint{1.800957in}{1.721865in}}%
\pgfpathlineto{\pgfqpoint{1.807130in}{1.721260in}}%
\pgfpathlineto{\pgfqpoint{1.806996in}{1.702523in}}%
\pgfpathlineto{\pgfqpoint{1.805445in}{1.674881in}}%
\pgfpathlineto{\pgfqpoint{1.805743in}{1.664546in}}%
\pgfpathclose%
\pgfusepath{fill}%
\end{pgfscope}%
\begin{pgfscope}%
\pgfpathrectangle{\pgfqpoint{0.100000in}{0.100000in}}{\pgfqpoint{3.608454in}{2.310000in}}%
\pgfusepath{clip}%
\pgfsetbuttcap%
\pgfsetmiterjoin%
\definecolor{currentfill}{rgb}{0.000000,0.607843,0.696078}%
\pgfsetfillcolor{currentfill}%
\pgfsetlinewidth{0.000000pt}%
\definecolor{currentstroke}{rgb}{0.000000,0.000000,0.000000}%
\pgfsetstrokecolor{currentstroke}%
\pgfsetstrokeopacity{0.000000}%
\pgfsetdash{}{0pt}%
\pgfpathmoveto{\pgfqpoint{0.895815in}{1.966858in}}%
\pgfpathlineto{\pgfqpoint{0.891257in}{1.963356in}}%
\pgfpathlineto{\pgfqpoint{0.890375in}{1.956866in}}%
\pgfpathlineto{\pgfqpoint{0.884981in}{1.950005in}}%
\pgfpathlineto{\pgfqpoint{0.875640in}{1.944745in}}%
\pgfpathlineto{\pgfqpoint{0.866297in}{1.932367in}}%
\pgfpathlineto{\pgfqpoint{0.860164in}{1.926417in}}%
\pgfpathlineto{\pgfqpoint{0.857097in}{1.912795in}}%
\pgfpathlineto{\pgfqpoint{0.842438in}{1.916478in}}%
\pgfpathlineto{\pgfqpoint{0.844130in}{1.923115in}}%
\pgfpathlineto{\pgfqpoint{0.839459in}{1.928946in}}%
\pgfpathlineto{\pgfqpoint{0.818538in}{1.934210in}}%
\pgfpathlineto{\pgfqpoint{0.814693in}{1.932896in}}%
\pgfpathlineto{\pgfqpoint{0.805150in}{1.923524in}}%
\pgfpathlineto{\pgfqpoint{0.800486in}{1.923543in}}%
\pgfpathlineto{\pgfqpoint{0.785890in}{1.927298in}}%
\pgfpathlineto{\pgfqpoint{0.792882in}{1.935165in}}%
\pgfpathlineto{\pgfqpoint{0.793934in}{1.941045in}}%
\pgfpathlineto{\pgfqpoint{0.803077in}{1.950037in}}%
\pgfpathlineto{\pgfqpoint{0.801936in}{1.953999in}}%
\pgfpathlineto{\pgfqpoint{0.793832in}{1.958919in}}%
\pgfpathlineto{\pgfqpoint{0.793910in}{1.961931in}}%
\pgfpathlineto{\pgfqpoint{0.807519in}{1.962306in}}%
\pgfpathlineto{\pgfqpoint{0.807404in}{1.969358in}}%
\pgfpathlineto{\pgfqpoint{0.813336in}{1.971025in}}%
\pgfpathlineto{\pgfqpoint{0.814622in}{1.975723in}}%
\pgfpathlineto{\pgfqpoint{0.808158in}{1.981911in}}%
\pgfpathlineto{\pgfqpoint{0.799750in}{1.984499in}}%
\pgfpathlineto{\pgfqpoint{0.802344in}{2.000341in}}%
\pgfpathlineto{\pgfqpoint{0.796435in}{2.001902in}}%
\pgfpathlineto{\pgfqpoint{0.797207in}{2.013901in}}%
\pgfpathlineto{\pgfqpoint{0.811964in}{2.010809in}}%
\pgfpathlineto{\pgfqpoint{0.813450in}{2.016488in}}%
\pgfpathlineto{\pgfqpoint{0.825859in}{2.013085in}}%
\pgfpathlineto{\pgfqpoint{0.831103in}{2.015303in}}%
\pgfpathlineto{\pgfqpoint{0.835457in}{2.032072in}}%
\pgfpathlineto{\pgfqpoint{0.838162in}{2.031358in}}%
\pgfpathlineto{\pgfqpoint{0.841907in}{2.040973in}}%
\pgfpathlineto{\pgfqpoint{0.846645in}{2.043323in}}%
\pgfpathlineto{\pgfqpoint{0.858767in}{2.040242in}}%
\pgfpathlineto{\pgfqpoint{0.853215in}{2.027577in}}%
\pgfpathlineto{\pgfqpoint{0.855049in}{2.021169in}}%
\pgfpathlineto{\pgfqpoint{0.853023in}{2.013264in}}%
\pgfpathlineto{\pgfqpoint{0.854053in}{1.998430in}}%
\pgfpathlineto{\pgfqpoint{0.859301in}{1.991402in}}%
\pgfpathlineto{\pgfqpoint{0.859744in}{1.982888in}}%
\pgfpathlineto{\pgfqpoint{0.871160in}{1.980064in}}%
\pgfpathlineto{\pgfqpoint{0.869552in}{1.973496in}}%
\pgfpathlineto{\pgfqpoint{0.895815in}{1.966858in}}%
\pgfpathclose%
\pgfusepath{fill}%
\end{pgfscope}%
\begin{pgfscope}%
\pgfpathrectangle{\pgfqpoint{0.100000in}{0.100000in}}{\pgfqpoint{3.608454in}{2.310000in}}%
\pgfusepath{clip}%
\pgfsetbuttcap%
\pgfsetmiterjoin%
\definecolor{currentfill}{rgb}{0.000000,0.364706,0.817647}%
\pgfsetfillcolor{currentfill}%
\pgfsetlinewidth{0.000000pt}%
\definecolor{currentstroke}{rgb}{0.000000,0.000000,0.000000}%
\pgfsetstrokecolor{currentstroke}%
\pgfsetstrokeopacity{0.000000}%
\pgfsetdash{}{0pt}%
\pgfpathmoveto{\pgfqpoint{1.791350in}{1.359031in}}%
\pgfpathlineto{\pgfqpoint{1.789033in}{1.324734in}}%
\pgfpathlineto{\pgfqpoint{1.756162in}{1.326912in}}%
\pgfpathlineto{\pgfqpoint{1.707133in}{1.330203in}}%
\pgfpathlineto{\pgfqpoint{1.709976in}{1.364447in}}%
\pgfpathlineto{\pgfqpoint{1.715265in}{1.364101in}}%
\pgfpathlineto{\pgfqpoint{1.717833in}{1.398242in}}%
\pgfpathlineto{\pgfqpoint{1.791324in}{1.393448in}}%
\pgfpathlineto{\pgfqpoint{1.792462in}{1.393414in}}%
\pgfpathlineto{\pgfqpoint{1.790342in}{1.359067in}}%
\pgfpathlineto{\pgfqpoint{1.791350in}{1.359031in}}%
\pgfpathclose%
\pgfusepath{fill}%
\end{pgfscope}%
\begin{pgfscope}%
\pgfpathrectangle{\pgfqpoint{0.100000in}{0.100000in}}{\pgfqpoint{3.608454in}{2.310000in}}%
\pgfusepath{clip}%
\pgfsetbuttcap%
\pgfsetmiterjoin%
\definecolor{currentfill}{rgb}{0.000000,0.486275,0.756863}%
\pgfsetfillcolor{currentfill}%
\pgfsetlinewidth{0.000000pt}%
\definecolor{currentstroke}{rgb}{0.000000,0.000000,0.000000}%
\pgfsetstrokecolor{currentstroke}%
\pgfsetstrokeopacity{0.000000}%
\pgfsetdash{}{0pt}%
\pgfpathmoveto{\pgfqpoint{2.464725in}{1.378777in}}%
\pgfpathlineto{\pgfqpoint{2.440862in}{1.377267in}}%
\pgfpathlineto{\pgfqpoint{2.439852in}{1.391129in}}%
\pgfpathlineto{\pgfqpoint{2.415980in}{1.389550in}}%
\pgfpathlineto{\pgfqpoint{2.410721in}{1.410092in}}%
\pgfpathlineto{\pgfqpoint{2.409257in}{1.435587in}}%
\pgfpathlineto{\pgfqpoint{2.410568in}{1.438329in}}%
\pgfpathlineto{\pgfqpoint{2.418888in}{1.437396in}}%
\pgfpathlineto{\pgfqpoint{2.427687in}{1.440218in}}%
\pgfpathlineto{\pgfqpoint{2.432920in}{1.438056in}}%
\pgfpathlineto{\pgfqpoint{2.431914in}{1.453643in}}%
\pgfpathlineto{\pgfqpoint{2.452287in}{1.455302in}}%
\pgfpathlineto{\pgfqpoint{2.459428in}{1.452340in}}%
\pgfpathlineto{\pgfqpoint{2.460863in}{1.433950in}}%
\pgfpathlineto{\pgfqpoint{2.461611in}{1.423638in}}%
\pgfpathlineto{\pgfqpoint{2.457108in}{1.423280in}}%
\pgfpathlineto{\pgfqpoint{2.457647in}{1.415090in}}%
\pgfpathlineto{\pgfqpoint{2.462297in}{1.414401in}}%
\pgfpathlineto{\pgfqpoint{2.463205in}{1.402934in}}%
\pgfpathlineto{\pgfqpoint{2.470095in}{1.403294in}}%
\pgfpathlineto{\pgfqpoint{2.471713in}{1.379003in}}%
\pgfpathlineto{\pgfqpoint{2.464725in}{1.378777in}}%
\pgfpathclose%
\pgfusepath{fill}%
\end{pgfscope}%
\begin{pgfscope}%
\pgfpathrectangle{\pgfqpoint{0.100000in}{0.100000in}}{\pgfqpoint{3.608454in}{2.310000in}}%
\pgfusepath{clip}%
\pgfsetbuttcap%
\pgfsetmiterjoin%
\definecolor{currentfill}{rgb}{0.000000,0.556863,0.721569}%
\pgfsetfillcolor{currentfill}%
\pgfsetlinewidth{0.000000pt}%
\definecolor{currentstroke}{rgb}{0.000000,0.000000,0.000000}%
\pgfsetstrokecolor{currentstroke}%
\pgfsetstrokeopacity{0.000000}%
\pgfsetdash{}{0pt}%
\pgfpathmoveto{\pgfqpoint{3.344220in}{1.667169in}}%
\pgfpathlineto{\pgfqpoint{3.332839in}{1.649737in}}%
\pgfpathlineto{\pgfqpoint{3.302353in}{1.660141in}}%
\pgfpathlineto{\pgfqpoint{3.298905in}{1.664653in}}%
\pgfpathlineto{\pgfqpoint{3.292838in}{1.664005in}}%
\pgfpathlineto{\pgfqpoint{3.282783in}{1.667998in}}%
\pgfpathlineto{\pgfqpoint{3.277583in}{1.676067in}}%
\pgfpathlineto{\pgfqpoint{3.274806in}{1.685909in}}%
\pgfpathlineto{\pgfqpoint{3.267733in}{1.692380in}}%
\pgfpathlineto{\pgfqpoint{3.302028in}{1.725960in}}%
\pgfpathlineto{\pgfqpoint{3.311248in}{1.723631in}}%
\pgfpathlineto{\pgfqpoint{3.324878in}{1.725364in}}%
\pgfpathlineto{\pgfqpoint{3.325396in}{1.731486in}}%
\pgfpathlineto{\pgfqpoint{3.333233in}{1.729959in}}%
\pgfpathlineto{\pgfqpoint{3.335015in}{1.726252in}}%
\pgfpathlineto{\pgfqpoint{3.347348in}{1.723269in}}%
\pgfpathlineto{\pgfqpoint{3.358373in}{1.723673in}}%
\pgfpathlineto{\pgfqpoint{3.358769in}{1.729605in}}%
\pgfpathlineto{\pgfqpoint{3.359316in}{1.729735in}}%
\pgfpathlineto{\pgfqpoint{3.366550in}{1.688953in}}%
\pgfpathlineto{\pgfqpoint{3.344000in}{1.680413in}}%
\pgfpathlineto{\pgfqpoint{3.344220in}{1.667169in}}%
\pgfpathclose%
\pgfusepath{fill}%
\end{pgfscope}%
\begin{pgfscope}%
\pgfpathrectangle{\pgfqpoint{0.100000in}{0.100000in}}{\pgfqpoint{3.608454in}{2.310000in}}%
\pgfusepath{clip}%
\pgfsetbuttcap%
\pgfsetmiterjoin%
\definecolor{currentfill}{rgb}{0.000000,0.674510,0.662745}%
\pgfsetfillcolor{currentfill}%
\pgfsetlinewidth{0.000000pt}%
\definecolor{currentstroke}{rgb}{0.000000,0.000000,0.000000}%
\pgfsetstrokecolor{currentstroke}%
\pgfsetstrokeopacity{0.000000}%
\pgfsetdash{}{0pt}%
\pgfpathmoveto{\pgfqpoint{1.443108in}{0.820998in}}%
\pgfpathlineto{\pgfqpoint{1.438374in}{0.821513in}}%
\pgfpathlineto{\pgfqpoint{1.404085in}{0.825375in}}%
\pgfpathlineto{\pgfqpoint{1.340564in}{0.833080in}}%
\pgfpathlineto{\pgfqpoint{1.350085in}{0.908598in}}%
\pgfpathlineto{\pgfqpoint{1.353414in}{0.915777in}}%
\pgfpathlineto{\pgfqpoint{1.356403in}{0.942467in}}%
\pgfpathlineto{\pgfqpoint{1.355496in}{0.949729in}}%
\pgfpathlineto{\pgfqpoint{1.357228in}{0.963939in}}%
\pgfpathlineto{\pgfqpoint{1.378188in}{0.960531in}}%
\pgfpathlineto{\pgfqpoint{1.379878in}{0.974364in}}%
\pgfpathlineto{\pgfqpoint{1.388260in}{0.973319in}}%
\pgfpathlineto{\pgfqpoint{1.392442in}{1.007538in}}%
\pgfpathlineto{\pgfqpoint{1.432205in}{1.002872in}}%
\pgfpathlineto{\pgfqpoint{1.432998in}{1.009699in}}%
\pgfpathlineto{\pgfqpoint{1.460404in}{1.006662in}}%
\pgfpathlineto{\pgfqpoint{1.458090in}{0.986273in}}%
\pgfpathlineto{\pgfqpoint{1.452184in}{0.931773in}}%
\pgfpathlineto{\pgfqpoint{1.449063in}{0.911411in}}%
\pgfpathlineto{\pgfqpoint{1.421822in}{0.913943in}}%
\pgfpathlineto{\pgfqpoint{1.420339in}{0.900498in}}%
\pgfpathlineto{\pgfqpoint{1.417927in}{0.900770in}}%
\pgfpathlineto{\pgfqpoint{1.413830in}{0.865768in}}%
\pgfpathlineto{\pgfqpoint{1.447334in}{0.862167in}}%
\pgfpathlineto{\pgfqpoint{1.443108in}{0.820998in}}%
\pgfpathclose%
\pgfusepath{fill}%
\end{pgfscope}%
\begin{pgfscope}%
\pgfpathrectangle{\pgfqpoint{0.100000in}{0.100000in}}{\pgfqpoint{3.608454in}{2.310000in}}%
\pgfusepath{clip}%
\pgfsetbuttcap%
\pgfsetmiterjoin%
\definecolor{currentfill}{rgb}{0.000000,0.627451,0.686275}%
\pgfsetfillcolor{currentfill}%
\pgfsetlinewidth{0.000000pt}%
\definecolor{currentstroke}{rgb}{0.000000,0.000000,0.000000}%
\pgfsetstrokecolor{currentstroke}%
\pgfsetstrokeopacity{0.000000}%
\pgfsetdash{}{0pt}%
\pgfpathmoveto{\pgfqpoint{2.853827in}{0.823500in}}%
\pgfpathlineto{\pgfqpoint{2.808283in}{0.818788in}}%
\pgfpathlineto{\pgfqpoint{2.786918in}{0.816730in}}%
\pgfpathlineto{\pgfqpoint{2.778092in}{0.819839in}}%
\pgfpathlineto{\pgfqpoint{2.776392in}{0.826653in}}%
\pgfpathlineto{\pgfqpoint{2.775432in}{0.832235in}}%
\pgfpathlineto{\pgfqpoint{2.783833in}{0.840916in}}%
\pgfpathlineto{\pgfqpoint{2.801485in}{0.843067in}}%
\pgfpathlineto{\pgfqpoint{2.799223in}{0.863800in}}%
\pgfpathlineto{\pgfqpoint{2.800195in}{0.867107in}}%
\pgfpathlineto{\pgfqpoint{2.812799in}{0.870744in}}%
\pgfpathlineto{\pgfqpoint{2.816271in}{0.861716in}}%
\pgfpathlineto{\pgfqpoint{2.822592in}{0.856703in}}%
\pgfpathlineto{\pgfqpoint{2.830726in}{0.859944in}}%
\pgfpathlineto{\pgfqpoint{2.834772in}{0.868856in}}%
\pgfpathlineto{\pgfqpoint{2.839212in}{0.870674in}}%
\pgfpathlineto{\pgfqpoint{2.851231in}{0.869522in}}%
\pgfpathlineto{\pgfqpoint{2.854444in}{0.865881in}}%
\pgfpathlineto{\pgfqpoint{2.855107in}{0.854082in}}%
\pgfpathlineto{\pgfqpoint{2.847916in}{0.849972in}}%
\pgfpathlineto{\pgfqpoint{2.844094in}{0.841666in}}%
\pgfpathlineto{\pgfqpoint{2.847520in}{0.837484in}}%
\pgfpathlineto{\pgfqpoint{2.853827in}{0.823500in}}%
\pgfpathclose%
\pgfusepath{fill}%
\end{pgfscope}%
\begin{pgfscope}%
\pgfpathrectangle{\pgfqpoint{0.100000in}{0.100000in}}{\pgfqpoint{3.608454in}{2.310000in}}%
\pgfusepath{clip}%
\pgfsetbuttcap%
\pgfsetmiterjoin%
\definecolor{currentfill}{rgb}{0.000000,0.843137,0.578431}%
\pgfsetfillcolor{currentfill}%
\pgfsetlinewidth{0.000000pt}%
\definecolor{currentstroke}{rgb}{0.000000,0.000000,0.000000}%
\pgfsetstrokecolor{currentstroke}%
\pgfsetstrokeopacity{0.000000}%
\pgfsetdash{}{0pt}%
\pgfpathmoveto{\pgfqpoint{2.756041in}{1.789197in}}%
\pgfpathlineto{\pgfqpoint{2.738013in}{1.786827in}}%
\pgfpathlineto{\pgfqpoint{2.729298in}{1.855252in}}%
\pgfpathlineto{\pgfqpoint{2.728097in}{1.868854in}}%
\pgfpathlineto{\pgfqpoint{2.707810in}{1.866142in}}%
\pgfpathlineto{\pgfqpoint{2.703507in}{1.899857in}}%
\pgfpathlineto{\pgfqpoint{2.706092in}{1.900170in}}%
\pgfpathlineto{\pgfqpoint{2.712457in}{1.891199in}}%
\pgfpathlineto{\pgfqpoint{2.723164in}{1.891127in}}%
\pgfpathlineto{\pgfqpoint{2.731458in}{1.886319in}}%
\pgfpathlineto{\pgfqpoint{2.749951in}{1.882041in}}%
\pgfpathlineto{\pgfqpoint{2.753811in}{1.874987in}}%
\pgfpathlineto{\pgfqpoint{2.760125in}{1.868000in}}%
\pgfpathlineto{\pgfqpoint{2.764215in}{1.858878in}}%
\pgfpathlineto{\pgfqpoint{2.754329in}{1.860433in}}%
\pgfpathlineto{\pgfqpoint{2.755769in}{1.850748in}}%
\pgfpathlineto{\pgfqpoint{2.762952in}{1.847681in}}%
\pgfpathlineto{\pgfqpoint{2.767020in}{1.834435in}}%
\pgfpathlineto{\pgfqpoint{2.765648in}{1.826069in}}%
\pgfpathlineto{\pgfqpoint{2.767351in}{1.804576in}}%
\pgfpathlineto{\pgfqpoint{2.762782in}{1.798736in}}%
\pgfpathlineto{\pgfqpoint{2.757883in}{1.797924in}}%
\pgfpathlineto{\pgfqpoint{2.756041in}{1.789197in}}%
\pgfpathclose%
\pgfusepath{fill}%
\end{pgfscope}%
\begin{pgfscope}%
\pgfpathrectangle{\pgfqpoint{0.100000in}{0.100000in}}{\pgfqpoint{3.608454in}{2.310000in}}%
\pgfusepath{clip}%
\pgfsetbuttcap%
\pgfsetmiterjoin%
\definecolor{currentfill}{rgb}{0.000000,0.513725,0.743137}%
\pgfsetfillcolor{currentfill}%
\pgfsetlinewidth{0.000000pt}%
\definecolor{currentstroke}{rgb}{0.000000,0.000000,0.000000}%
\pgfsetstrokecolor{currentstroke}%
\pgfsetstrokeopacity{0.000000}%
\pgfsetdash{}{0pt}%
\pgfpathmoveto{\pgfqpoint{2.773279in}{1.502104in}}%
\pgfpathlineto{\pgfqpoint{2.766536in}{1.500999in}}%
\pgfpathlineto{\pgfqpoint{2.767029in}{1.496892in}}%
\pgfpathlineto{\pgfqpoint{2.746418in}{1.494296in}}%
\pgfpathlineto{\pgfqpoint{2.747392in}{1.486269in}}%
\pgfpathlineto{\pgfqpoint{2.741875in}{1.483652in}}%
\pgfpathlineto{\pgfqpoint{2.719775in}{1.481004in}}%
\pgfpathlineto{\pgfqpoint{2.714189in}{1.531011in}}%
\pgfpathlineto{\pgfqpoint{2.741540in}{1.534284in}}%
\pgfpathlineto{\pgfqpoint{2.739953in}{1.548013in}}%
\pgfpathlineto{\pgfqpoint{2.746637in}{1.548837in}}%
\pgfpathlineto{\pgfqpoint{2.767138in}{1.551464in}}%
\pgfpathlineto{\pgfqpoint{2.773279in}{1.502104in}}%
\pgfpathclose%
\pgfusepath{fill}%
\end{pgfscope}%
\begin{pgfscope}%
\pgfpathrectangle{\pgfqpoint{0.100000in}{0.100000in}}{\pgfqpoint{3.608454in}{2.310000in}}%
\pgfusepath{clip}%
\pgfsetbuttcap%
\pgfsetmiterjoin%
\definecolor{currentfill}{rgb}{0.000000,0.768627,0.615686}%
\pgfsetfillcolor{currentfill}%
\pgfsetlinewidth{0.000000pt}%
\definecolor{currentstroke}{rgb}{0.000000,0.000000,0.000000}%
\pgfsetstrokecolor{currentstroke}%
\pgfsetstrokeopacity{0.000000}%
\pgfsetdash{}{0pt}%
\pgfpathmoveto{\pgfqpoint{3.130816in}{1.638979in}}%
\pgfpathlineto{\pgfqpoint{3.131860in}{1.634049in}}%
\pgfpathlineto{\pgfqpoint{3.109153in}{1.629395in}}%
\pgfpathlineto{\pgfqpoint{3.105576in}{1.628770in}}%
\pgfpathlineto{\pgfqpoint{3.094591in}{1.638194in}}%
\pgfpathlineto{\pgfqpoint{3.078140in}{1.636089in}}%
\pgfpathlineto{\pgfqpoint{3.050721in}{1.630639in}}%
\pgfpathlineto{\pgfqpoint{3.048278in}{1.646809in}}%
\pgfpathlineto{\pgfqpoint{3.050043in}{1.647962in}}%
\pgfpathlineto{\pgfqpoint{3.047674in}{1.660178in}}%
\pgfpathlineto{\pgfqpoint{3.039413in}{1.658799in}}%
\pgfpathlineto{\pgfqpoint{3.032000in}{1.700788in}}%
\pgfpathlineto{\pgfqpoint{3.036017in}{1.700877in}}%
\pgfpathlineto{\pgfqpoint{3.042918in}{1.694827in}}%
\pgfpathlineto{\pgfqpoint{3.049340in}{1.697252in}}%
\pgfpathlineto{\pgfqpoint{3.064906in}{1.706470in}}%
\pgfpathlineto{\pgfqpoint{3.066495in}{1.705513in}}%
\pgfpathlineto{\pgfqpoint{3.102306in}{1.712116in}}%
\pgfpathlineto{\pgfqpoint{3.109751in}{1.709791in}}%
\pgfpathlineto{\pgfqpoint{3.115261in}{1.672766in}}%
\pgfpathlineto{\pgfqpoint{3.123343in}{1.674373in}}%
\pgfpathlineto{\pgfqpoint{3.130816in}{1.638979in}}%
\pgfpathclose%
\pgfusepath{fill}%
\end{pgfscope}%
\begin{pgfscope}%
\pgfpathrectangle{\pgfqpoint{0.100000in}{0.100000in}}{\pgfqpoint{3.608454in}{2.310000in}}%
\pgfusepath{clip}%
\pgfsetbuttcap%
\pgfsetmiterjoin%
\definecolor{currentfill}{rgb}{0.000000,0.552941,0.723529}%
\pgfsetfillcolor{currentfill}%
\pgfsetlinewidth{0.000000pt}%
\definecolor{currentstroke}{rgb}{0.000000,0.000000,0.000000}%
\pgfsetstrokecolor{currentstroke}%
\pgfsetstrokeopacity{0.000000}%
\pgfsetdash{}{0pt}%
\pgfpathmoveto{\pgfqpoint{2.682759in}{1.669371in}}%
\pgfpathlineto{\pgfqpoint{2.696454in}{1.670883in}}%
\pgfpathlineto{\pgfqpoint{2.693396in}{1.698299in}}%
\pgfpathlineto{\pgfqpoint{2.720350in}{1.701323in}}%
\pgfpathlineto{\pgfqpoint{2.723746in}{1.674452in}}%
\pgfpathlineto{\pgfqpoint{2.735608in}{1.675908in}}%
\pgfpathlineto{\pgfqpoint{2.739935in}{1.648352in}}%
\pgfpathlineto{\pgfqpoint{2.706426in}{1.644227in}}%
\pgfpathlineto{\pgfqpoint{2.685918in}{1.641892in}}%
\pgfpathlineto{\pgfqpoint{2.682759in}{1.669371in}}%
\pgfpathclose%
\pgfusepath{fill}%
\end{pgfscope}%
\begin{pgfscope}%
\pgfpathrectangle{\pgfqpoint{0.100000in}{0.100000in}}{\pgfqpoint{3.608454in}{2.310000in}}%
\pgfusepath{clip}%
\pgfsetbuttcap%
\pgfsetmiterjoin%
\definecolor{currentfill}{rgb}{0.000000,0.447059,0.776471}%
\pgfsetfillcolor{currentfill}%
\pgfsetlinewidth{0.000000pt}%
\definecolor{currentstroke}{rgb}{0.000000,0.000000,0.000000}%
\pgfsetstrokecolor{currentstroke}%
\pgfsetstrokeopacity{0.000000}%
\pgfsetdash{}{0pt}%
\pgfpathmoveto{\pgfqpoint{1.992478in}{2.070813in}}%
\pgfpathlineto{\pgfqpoint{1.993953in}{2.077757in}}%
\pgfpathlineto{\pgfqpoint{1.991966in}{2.086284in}}%
\pgfpathlineto{\pgfqpoint{1.993550in}{2.092414in}}%
\pgfpathlineto{\pgfqpoint{1.991978in}{2.098902in}}%
\pgfpathlineto{\pgfqpoint{2.033124in}{2.098129in}}%
\pgfpathlineto{\pgfqpoint{2.074816in}{2.097239in}}%
\pgfpathlineto{\pgfqpoint{2.075068in}{2.055715in}}%
\pgfpathlineto{\pgfqpoint{2.075577in}{2.048780in}}%
\pgfpathlineto{\pgfqpoint{2.068749in}{2.048976in}}%
\pgfpathlineto{\pgfqpoint{2.068716in}{2.041994in}}%
\pgfpathlineto{\pgfqpoint{2.061858in}{2.042025in}}%
\pgfpathlineto{\pgfqpoint{2.061781in}{2.035022in}}%
\pgfpathlineto{\pgfqpoint{2.034103in}{2.035335in}}%
\pgfpathlineto{\pgfqpoint{2.034152in}{2.042276in}}%
\pgfpathlineto{\pgfqpoint{2.027218in}{2.042403in}}%
\pgfpathlineto{\pgfqpoint{2.026639in}{2.068568in}}%
\pgfpathlineto{\pgfqpoint{1.992147in}{2.069188in}}%
\pgfpathlineto{\pgfqpoint{1.992478in}{2.070813in}}%
\pgfpathclose%
\pgfusepath{fill}%
\end{pgfscope}%
\begin{pgfscope}%
\pgfpathrectangle{\pgfqpoint{0.100000in}{0.100000in}}{\pgfqpoint{3.608454in}{2.310000in}}%
\pgfusepath{clip}%
\pgfsetbuttcap%
\pgfsetmiterjoin%
\definecolor{currentfill}{rgb}{0.000000,0.439216,0.780392}%
\pgfsetfillcolor{currentfill}%
\pgfsetlinewidth{0.000000pt}%
\definecolor{currentstroke}{rgb}{0.000000,0.000000,0.000000}%
\pgfsetstrokecolor{currentstroke}%
\pgfsetstrokeopacity{0.000000}%
\pgfsetdash{}{0pt}%
\pgfpathmoveto{\pgfqpoint{3.181076in}{1.501996in}}%
\pgfpathlineto{\pgfqpoint{3.166603in}{1.499150in}}%
\pgfpathlineto{\pgfqpoint{3.162472in}{1.516430in}}%
\pgfpathlineto{\pgfqpoint{3.146592in}{1.539583in}}%
\pgfpathlineto{\pgfqpoint{3.143924in}{1.538541in}}%
\pgfpathlineto{\pgfqpoint{3.138889in}{1.546016in}}%
\pgfpathlineto{\pgfqpoint{3.133310in}{1.544667in}}%
\pgfpathlineto{\pgfqpoint{3.130234in}{1.548464in}}%
\pgfpathlineto{\pgfqpoint{3.133732in}{1.555574in}}%
\pgfpathlineto{\pgfqpoint{3.131483in}{1.559276in}}%
\pgfpathlineto{\pgfqpoint{3.139700in}{1.576466in}}%
\pgfpathlineto{\pgfqpoint{3.155139in}{1.587716in}}%
\pgfpathlineto{\pgfqpoint{3.157897in}{1.576523in}}%
\pgfpathlineto{\pgfqpoint{3.176843in}{1.578501in}}%
\pgfpathlineto{\pgfqpoint{3.184525in}{1.574369in}}%
\pgfpathlineto{\pgfqpoint{3.196959in}{1.580835in}}%
\pgfpathlineto{\pgfqpoint{3.208359in}{1.574869in}}%
\pgfpathlineto{\pgfqpoint{3.201240in}{1.566916in}}%
\pgfpathlineto{\pgfqpoint{3.212062in}{1.546662in}}%
\pgfpathlineto{\pgfqpoint{3.204878in}{1.540940in}}%
\pgfpathlineto{\pgfqpoint{3.208552in}{1.535223in}}%
\pgfpathlineto{\pgfqpoint{3.216418in}{1.535851in}}%
\pgfpathlineto{\pgfqpoint{3.221777in}{1.528537in}}%
\pgfpathlineto{\pgfqpoint{3.227209in}{1.526852in}}%
\pgfpathlineto{\pgfqpoint{3.239245in}{1.513965in}}%
\pgfpathlineto{\pgfqpoint{3.181076in}{1.501996in}}%
\pgfpathclose%
\pgfusepath{fill}%
\end{pgfscope}%
\begin{pgfscope}%
\pgfpathrectangle{\pgfqpoint{0.100000in}{0.100000in}}{\pgfqpoint{3.608454in}{2.310000in}}%
\pgfusepath{clip}%
\pgfsetbuttcap%
\pgfsetmiterjoin%
\definecolor{currentfill}{rgb}{0.000000,0.278431,0.860784}%
\pgfsetfillcolor{currentfill}%
\pgfsetlinewidth{0.000000pt}%
\definecolor{currentstroke}{rgb}{0.000000,0.000000,0.000000}%
\pgfsetstrokecolor{currentstroke}%
\pgfsetstrokeopacity{0.000000}%
\pgfsetdash{}{0pt}%
\pgfpathmoveto{\pgfqpoint{1.654457in}{1.551499in}}%
\pgfpathlineto{\pgfqpoint{1.655781in}{1.547900in}}%
\pgfpathlineto{\pgfqpoint{1.654575in}{1.534236in}}%
\pgfpathlineto{\pgfqpoint{1.653438in}{1.520561in}}%
\pgfpathlineto{\pgfqpoint{1.650561in}{1.517162in}}%
\pgfpathlineto{\pgfqpoint{1.595804in}{1.522015in}}%
\pgfpathlineto{\pgfqpoint{1.533391in}{1.528263in}}%
\pgfpathlineto{\pgfqpoint{1.494704in}{1.532257in}}%
\pgfpathlineto{\pgfqpoint{1.500517in}{1.584019in}}%
\pgfpathlineto{\pgfqpoint{1.537205in}{1.579689in}}%
\pgfpathlineto{\pgfqpoint{1.536486in}{1.572765in}}%
\pgfpathlineto{\pgfqpoint{1.571835in}{1.569055in}}%
\pgfpathlineto{\pgfqpoint{1.570505in}{1.555595in}}%
\pgfpathlineto{\pgfqpoint{1.610795in}{1.551818in}}%
\pgfpathlineto{\pgfqpoint{1.611187in}{1.555209in}}%
\pgfpathlineto{\pgfqpoint{1.654457in}{1.551499in}}%
\pgfpathclose%
\pgfusepath{fill}%
\end{pgfscope}%
\begin{pgfscope}%
\pgfpathrectangle{\pgfqpoint{0.100000in}{0.100000in}}{\pgfqpoint{3.608454in}{2.310000in}}%
\pgfusepath{clip}%
\pgfsetbuttcap%
\pgfsetmiterjoin%
\definecolor{currentfill}{rgb}{0.000000,0.713725,0.643137}%
\pgfsetfillcolor{currentfill}%
\pgfsetlinewidth{0.000000pt}%
\definecolor{currentstroke}{rgb}{0.000000,0.000000,0.000000}%
\pgfsetstrokecolor{currentstroke}%
\pgfsetstrokeopacity{0.000000}%
\pgfsetdash{}{0pt}%
\pgfpathmoveto{\pgfqpoint{1.421683in}{1.386315in}}%
\pgfpathlineto{\pgfqpoint{1.421188in}{1.383353in}}%
\pgfpathlineto{\pgfqpoint{1.430521in}{1.376213in}}%
\pgfpathlineto{\pgfqpoint{1.436137in}{1.369904in}}%
\pgfpathlineto{\pgfqpoint{1.436653in}{1.364650in}}%
\pgfpathlineto{\pgfqpoint{1.431712in}{1.356218in}}%
\pgfpathlineto{\pgfqpoint{1.436154in}{1.349875in}}%
\pgfpathlineto{\pgfqpoint{1.435665in}{1.345243in}}%
\pgfpathlineto{\pgfqpoint{1.426840in}{1.337144in}}%
\pgfpathlineto{\pgfqpoint{1.422616in}{1.335846in}}%
\pgfpathlineto{\pgfqpoint{1.412097in}{1.337063in}}%
\pgfpathlineto{\pgfqpoint{1.406993in}{1.344082in}}%
\pgfpathlineto{\pgfqpoint{1.404621in}{1.362130in}}%
\pgfpathlineto{\pgfqpoint{1.405672in}{1.369838in}}%
\pgfpathlineto{\pgfqpoint{1.412208in}{1.374642in}}%
\pgfpathlineto{\pgfqpoint{1.403546in}{1.376980in}}%
\pgfpathlineto{\pgfqpoint{1.396241in}{1.384895in}}%
\pgfpathlineto{\pgfqpoint{1.398138in}{1.389426in}}%
\pgfpathlineto{\pgfqpoint{1.421683in}{1.386315in}}%
\pgfpathclose%
\pgfusepath{fill}%
\end{pgfscope}%
\begin{pgfscope}%
\pgfpathrectangle{\pgfqpoint{0.100000in}{0.100000in}}{\pgfqpoint{3.608454in}{2.310000in}}%
\pgfusepath{clip}%
\pgfsetbuttcap%
\pgfsetmiterjoin%
\definecolor{currentfill}{rgb}{0.000000,0.572549,0.713725}%
\pgfsetfillcolor{currentfill}%
\pgfsetlinewidth{0.000000pt}%
\definecolor{currentstroke}{rgb}{0.000000,0.000000,0.000000}%
\pgfsetstrokecolor{currentstroke}%
\pgfsetstrokeopacity{0.000000}%
\pgfsetdash{}{0pt}%
\pgfpathmoveto{\pgfqpoint{2.517006in}{1.100169in}}%
\pgfpathlineto{\pgfqpoint{2.510594in}{1.099857in}}%
\pgfpathlineto{\pgfqpoint{2.510006in}{1.121373in}}%
\pgfpathlineto{\pgfqpoint{2.512080in}{1.125849in}}%
\pgfpathlineto{\pgfqpoint{2.519951in}{1.126253in}}%
\pgfpathlineto{\pgfqpoint{2.518781in}{1.156793in}}%
\pgfpathlineto{\pgfqpoint{2.520509in}{1.156888in}}%
\pgfpathlineto{\pgfqpoint{2.548186in}{1.158629in}}%
\pgfpathlineto{\pgfqpoint{2.548297in}{1.153932in}}%
\pgfpathlineto{\pgfqpoint{2.555391in}{1.144733in}}%
\pgfpathlineto{\pgfqpoint{2.556313in}{1.139100in}}%
\pgfpathlineto{\pgfqpoint{2.570262in}{1.139592in}}%
\pgfpathlineto{\pgfqpoint{2.579740in}{1.137359in}}%
\pgfpathlineto{\pgfqpoint{2.583376in}{1.141723in}}%
\pgfpathlineto{\pgfqpoint{2.583459in}{1.148544in}}%
\pgfpathlineto{\pgfqpoint{2.597965in}{1.148745in}}%
\pgfpathlineto{\pgfqpoint{2.606315in}{1.148332in}}%
\pgfpathlineto{\pgfqpoint{2.604989in}{1.140525in}}%
\pgfpathlineto{\pgfqpoint{2.606773in}{1.136211in}}%
\pgfpathlineto{\pgfqpoint{2.605651in}{1.120589in}}%
\pgfpathlineto{\pgfqpoint{2.584331in}{1.121514in}}%
\pgfpathlineto{\pgfqpoint{2.573891in}{1.113688in}}%
\pgfpathlineto{\pgfqpoint{2.573691in}{1.108364in}}%
\pgfpathlineto{\pgfqpoint{2.561312in}{1.107350in}}%
\pgfpathlineto{\pgfqpoint{2.553934in}{1.104609in}}%
\pgfpathlineto{\pgfqpoint{2.552190in}{1.107568in}}%
\pgfpathlineto{\pgfqpoint{2.544235in}{1.104329in}}%
\pgfpathlineto{\pgfqpoint{2.519200in}{1.103140in}}%
\pgfpathlineto{\pgfqpoint{2.517006in}{1.100169in}}%
\pgfpathclose%
\pgfusepath{fill}%
\end{pgfscope}%
\begin{pgfscope}%
\pgfpathrectangle{\pgfqpoint{0.100000in}{0.100000in}}{\pgfqpoint{3.608454in}{2.310000in}}%
\pgfusepath{clip}%
\pgfsetbuttcap%
\pgfsetmiterjoin%
\definecolor{currentfill}{rgb}{0.000000,0.360784,0.819608}%
\pgfsetfillcolor{currentfill}%
\pgfsetlinewidth{0.000000pt}%
\definecolor{currentstroke}{rgb}{0.000000,0.000000,0.000000}%
\pgfsetstrokecolor{currentstroke}%
\pgfsetstrokeopacity{0.000000}%
\pgfsetdash{}{0pt}%
\pgfpathmoveto{\pgfqpoint{1.027182in}{0.384801in}}%
\pgfpathlineto{\pgfqpoint{1.017303in}{0.390307in}}%
\pgfpathlineto{\pgfqpoint{1.020524in}{0.396033in}}%
\pgfpathlineto{\pgfqpoint{1.019250in}{0.398079in}}%
\pgfpathlineto{\pgfqpoint{1.024400in}{0.406832in}}%
\pgfpathlineto{\pgfqpoint{1.022393in}{0.408007in}}%
\pgfpathlineto{\pgfqpoint{1.023662in}{0.410127in}}%
\pgfpathlineto{\pgfqpoint{1.021385in}{0.411488in}}%
\pgfpathlineto{\pgfqpoint{1.022698in}{0.413642in}}%
\pgfpathlineto{\pgfqpoint{1.018953in}{0.415851in}}%
\pgfpathlineto{\pgfqpoint{1.018284in}{0.417154in}}%
\pgfpathlineto{\pgfqpoint{1.015043in}{0.415179in}}%
\pgfpathlineto{\pgfqpoint{1.012216in}{0.417308in}}%
\pgfpathlineto{\pgfqpoint{1.015064in}{0.421049in}}%
\pgfpathlineto{\pgfqpoint{1.014239in}{0.423539in}}%
\pgfpathlineto{\pgfqpoint{1.017896in}{0.425014in}}%
\pgfpathlineto{\pgfqpoint{1.019938in}{0.430914in}}%
\pgfpathlineto{\pgfqpoint{1.002410in}{0.441046in}}%
\pgfpathlineto{\pgfqpoint{0.996218in}{0.444880in}}%
\pgfpathlineto{\pgfqpoint{0.997579in}{0.447110in}}%
\pgfpathlineto{\pgfqpoint{0.986635in}{0.453959in}}%
\pgfpathlineto{\pgfqpoint{0.983915in}{0.449521in}}%
\pgfpathlineto{\pgfqpoint{0.979580in}{0.452265in}}%
\pgfpathlineto{\pgfqpoint{0.986452in}{0.463301in}}%
\pgfpathlineto{\pgfqpoint{0.972860in}{0.471797in}}%
\pgfpathlineto{\pgfqpoint{0.968975in}{0.474416in}}%
\pgfpathlineto{\pgfqpoint{0.966684in}{0.470967in}}%
\pgfpathlineto{\pgfqpoint{0.950166in}{0.482212in}}%
\pgfpathlineto{\pgfqpoint{0.940198in}{0.483801in}}%
\pgfpathlineto{\pgfqpoint{0.915352in}{0.487908in}}%
\pgfpathlineto{\pgfqpoint{0.914054in}{0.486250in}}%
\pgfpathlineto{\pgfqpoint{0.902205in}{0.495355in}}%
\pgfpathlineto{\pgfqpoint{0.894650in}{0.485392in}}%
\pgfpathlineto{\pgfqpoint{0.892390in}{0.483966in}}%
\pgfpathlineto{\pgfqpoint{0.888822in}{0.486756in}}%
\pgfpathlineto{\pgfqpoint{0.887256in}{0.484755in}}%
\pgfpathlineto{\pgfqpoint{0.883274in}{0.487912in}}%
\pgfpathlineto{\pgfqpoint{0.880946in}{0.484915in}}%
\pgfpathlineto{\pgfqpoint{0.878680in}{0.486403in}}%
\pgfpathlineto{\pgfqpoint{0.867623in}{0.494883in}}%
\pgfpathlineto{\pgfqpoint{0.854046in}{0.505533in}}%
\pgfpathlineto{\pgfqpoint{0.845087in}{0.513012in}}%
\pgfpathlineto{\pgfqpoint{0.846928in}{0.515267in}}%
\pgfpathlineto{\pgfqpoint{0.831654in}{0.528086in}}%
\pgfpathlineto{\pgfqpoint{0.830011in}{0.526141in}}%
\pgfpathlineto{\pgfqpoint{0.821946in}{0.532976in}}%
\pgfpathlineto{\pgfqpoint{0.819842in}{0.530526in}}%
\pgfpathlineto{\pgfqpoint{0.812784in}{0.536599in}}%
\pgfpathlineto{\pgfqpoint{0.802699in}{0.545595in}}%
\pgfpathlineto{\pgfqpoint{0.807682in}{0.551257in}}%
\pgfpathlineto{\pgfqpoint{0.806268in}{0.552525in}}%
\pgfpathlineto{\pgfqpoint{0.813107in}{0.560166in}}%
\pgfpathlineto{\pgfqpoint{0.813592in}{0.559787in}}%
\pgfpathlineto{\pgfqpoint{0.820507in}{0.567425in}}%
\pgfpathlineto{\pgfqpoint{0.823025in}{0.569214in}}%
\pgfpathlineto{\pgfqpoint{0.824219in}{0.568707in}}%
\pgfpathlineto{\pgfqpoint{0.829339in}{0.574413in}}%
\pgfpathlineto{\pgfqpoint{0.831146in}{0.572881in}}%
\pgfpathlineto{\pgfqpoint{0.836205in}{0.579038in}}%
\pgfpathlineto{\pgfqpoint{0.838193in}{0.577309in}}%
\pgfpathlineto{\pgfqpoint{0.839724in}{0.579022in}}%
\pgfpathlineto{\pgfqpoint{0.843712in}{0.575563in}}%
\pgfpathlineto{\pgfqpoint{0.845387in}{0.577604in}}%
\pgfpathlineto{\pgfqpoint{0.847165in}{0.576109in}}%
\pgfpathlineto{\pgfqpoint{0.848929in}{0.578173in}}%
\pgfpathlineto{\pgfqpoint{0.850691in}{0.576622in}}%
\pgfpathlineto{\pgfqpoint{0.856024in}{0.582334in}}%
\pgfpathlineto{\pgfqpoint{0.853750in}{0.584249in}}%
\pgfpathlineto{\pgfqpoint{0.867382in}{0.600383in}}%
\pgfpathlineto{\pgfqpoint{0.869420in}{0.598690in}}%
\pgfpathlineto{\pgfqpoint{0.872825in}{0.602697in}}%
\pgfpathlineto{\pgfqpoint{0.874796in}{0.601023in}}%
\pgfpathlineto{\pgfqpoint{0.881527in}{0.608877in}}%
\pgfpathlineto{\pgfqpoint{0.883388in}{0.607238in}}%
\pgfpathlineto{\pgfqpoint{0.888579in}{0.613334in}}%
\pgfpathlineto{\pgfqpoint{0.886576in}{0.615036in}}%
\pgfpathlineto{\pgfqpoint{0.895296in}{0.625196in}}%
\pgfpathlineto{\pgfqpoint{0.901313in}{0.620035in}}%
\pgfpathlineto{\pgfqpoint{0.904546in}{0.623960in}}%
\pgfpathlineto{\pgfqpoint{0.914790in}{0.615304in}}%
\pgfpathlineto{\pgfqpoint{0.921542in}{0.623737in}}%
\pgfpathlineto{\pgfqpoint{0.929884in}{0.616921in}}%
\pgfpathlineto{\pgfqpoint{0.926478in}{0.612658in}}%
\pgfpathlineto{\pgfqpoint{0.930663in}{0.609300in}}%
\pgfpathlineto{\pgfqpoint{0.934029in}{0.613492in}}%
\pgfpathlineto{\pgfqpoint{0.938173in}{0.610139in}}%
\pgfpathlineto{\pgfqpoint{0.936545in}{0.608078in}}%
\pgfpathlineto{\pgfqpoint{0.938627in}{0.606373in}}%
\pgfpathlineto{\pgfqpoint{0.936932in}{0.604209in}}%
\pgfpathlineto{\pgfqpoint{0.940885in}{0.601118in}}%
\pgfpathlineto{\pgfqpoint{0.946326in}{0.607097in}}%
\pgfpathlineto{\pgfqpoint{0.952323in}{0.602325in}}%
\pgfpathlineto{\pgfqpoint{0.955683in}{0.606572in}}%
\pgfpathlineto{\pgfqpoint{0.961995in}{0.601456in}}%
\pgfpathlineto{\pgfqpoint{0.964303in}{0.603251in}}%
\pgfpathlineto{\pgfqpoint{0.971435in}{0.612298in}}%
\pgfpathlineto{\pgfqpoint{0.969910in}{0.613404in}}%
\pgfpathlineto{\pgfqpoint{0.971558in}{0.615632in}}%
\pgfpathlineto{\pgfqpoint{0.967277in}{0.618907in}}%
\pgfpathlineto{\pgfqpoint{0.972240in}{0.625296in}}%
\pgfpathlineto{\pgfqpoint{0.974510in}{0.627028in}}%
\pgfpathlineto{\pgfqpoint{0.972344in}{0.628719in}}%
\pgfpathlineto{\pgfqpoint{0.974033in}{0.630789in}}%
\pgfpathlineto{\pgfqpoint{0.971825in}{0.632484in}}%
\pgfpathlineto{\pgfqpoint{0.973534in}{0.634639in}}%
\pgfpathlineto{\pgfqpoint{0.971445in}{0.636313in}}%
\pgfpathlineto{\pgfqpoint{0.976297in}{0.641465in}}%
\pgfpathlineto{\pgfqpoint{0.997892in}{0.624739in}}%
\pgfpathlineto{\pgfqpoint{1.026608in}{0.603847in}}%
\pgfpathlineto{\pgfqpoint{1.052163in}{0.586365in}}%
\pgfpathlineto{\pgfqpoint{1.056440in}{0.584227in}}%
\pgfpathlineto{\pgfqpoint{1.066205in}{0.578011in}}%
\pgfpathlineto{\pgfqpoint{1.074192in}{0.589904in}}%
\pgfpathlineto{\pgfqpoint{1.083089in}{0.584700in}}%
\pgfpathlineto{\pgfqpoint{1.110388in}{0.567904in}}%
\pgfpathlineto{\pgfqpoint{1.124077in}{0.559590in}}%
\pgfpathlineto{\pgfqpoint{1.114317in}{0.541984in}}%
\pgfpathlineto{\pgfqpoint{1.101033in}{0.518021in}}%
\pgfpathlineto{\pgfqpoint{1.087919in}{0.494364in}}%
\pgfpathlineto{\pgfqpoint{1.075939in}{0.472754in}}%
\pgfpathlineto{\pgfqpoint{1.063110in}{0.449610in}}%
\pgfpathlineto{\pgfqpoint{1.050408in}{0.426699in}}%
\pgfpathlineto{\pgfqpoint{1.037825in}{0.403999in}}%
\pgfpathlineto{\pgfqpoint{1.027182in}{0.384801in}}%
\pgfpathclose%
\pgfusepath{fill}%
\end{pgfscope}%
\begin{pgfscope}%
\pgfpathrectangle{\pgfqpoint{0.100000in}{0.100000in}}{\pgfqpoint{3.608454in}{2.310000in}}%
\pgfusepath{clip}%
\pgfsetbuttcap%
\pgfsetmiterjoin%
\definecolor{currentfill}{rgb}{0.000000,0.580392,0.709804}%
\pgfsetfillcolor{currentfill}%
\pgfsetlinewidth{0.000000pt}%
\definecolor{currentstroke}{rgb}{0.000000,0.000000,0.000000}%
\pgfsetstrokecolor{currentstroke}%
\pgfsetstrokeopacity{0.000000}%
\pgfsetdash{}{0pt}%
\pgfpathmoveto{\pgfqpoint{2.289673in}{1.168159in}}%
\pgfpathlineto{\pgfqpoint{2.248336in}{1.167927in}}%
\pgfpathlineto{\pgfqpoint{2.211773in}{1.168036in}}%
\pgfpathlineto{\pgfqpoint{2.211810in}{1.181895in}}%
\pgfpathlineto{\pgfqpoint{2.194584in}{1.181869in}}%
\pgfpathlineto{\pgfqpoint{2.194554in}{1.204465in}}%
\pgfpathlineto{\pgfqpoint{2.193249in}{1.204475in}}%
\pgfpathlineto{\pgfqpoint{2.193719in}{1.240848in}}%
\pgfpathlineto{\pgfqpoint{2.192435in}{1.247763in}}%
\pgfpathlineto{\pgfqpoint{2.196058in}{1.254577in}}%
\pgfpathlineto{\pgfqpoint{2.199587in}{1.254527in}}%
\pgfpathlineto{\pgfqpoint{2.199885in}{1.267142in}}%
\pgfpathlineto{\pgfqpoint{2.227204in}{1.266822in}}%
\pgfpathlineto{\pgfqpoint{2.227021in}{1.254178in}}%
\pgfpathlineto{\pgfqpoint{2.255244in}{1.254033in}}%
\pgfpathlineto{\pgfqpoint{2.260153in}{1.246809in}}%
\pgfpathlineto{\pgfqpoint{2.268476in}{1.251773in}}%
\pgfpathlineto{\pgfqpoint{2.268526in}{1.240287in}}%
\pgfpathlineto{\pgfqpoint{2.278985in}{1.235414in}}%
\pgfpathlineto{\pgfqpoint{2.279053in}{1.231943in}}%
\pgfpathlineto{\pgfqpoint{2.280237in}{1.188871in}}%
\pgfpathlineto{\pgfqpoint{2.290504in}{1.188921in}}%
\pgfpathlineto{\pgfqpoint{2.289673in}{1.168159in}}%
\pgfpathclose%
\pgfusepath{fill}%
\end{pgfscope}%
\begin{pgfscope}%
\pgfpathrectangle{\pgfqpoint{0.100000in}{0.100000in}}{\pgfqpoint{3.608454in}{2.310000in}}%
\pgfusepath{clip}%
\pgfsetbuttcap%
\pgfsetmiterjoin%
\definecolor{currentfill}{rgb}{0.000000,0.800000,0.600000}%
\pgfsetfillcolor{currentfill}%
\pgfsetlinewidth{0.000000pt}%
\definecolor{currentstroke}{rgb}{0.000000,0.000000,0.000000}%
\pgfsetstrokecolor{currentstroke}%
\pgfsetstrokeopacity{0.000000}%
\pgfsetdash{}{0pt}%
\pgfpathmoveto{\pgfqpoint{2.175503in}{0.924555in}}%
\pgfpathlineto{\pgfqpoint{2.170065in}{0.926960in}}%
\pgfpathlineto{\pgfqpoint{2.157026in}{0.928424in}}%
\pgfpathlineto{\pgfqpoint{2.147313in}{0.939341in}}%
\pgfpathlineto{\pgfqpoint{2.141202in}{0.939334in}}%
\pgfpathlineto{\pgfqpoint{2.140773in}{0.915277in}}%
\pgfpathlineto{\pgfqpoint{2.123559in}{0.920377in}}%
\pgfpathlineto{\pgfqpoint{2.115386in}{0.923776in}}%
\pgfpathlineto{\pgfqpoint{2.109858in}{0.929360in}}%
\pgfpathlineto{\pgfqpoint{2.092356in}{0.940891in}}%
\pgfpathlineto{\pgfqpoint{2.087358in}{0.933839in}}%
\pgfpathlineto{\pgfqpoint{2.079921in}{0.933374in}}%
\pgfpathlineto{\pgfqpoint{2.070424in}{0.935180in}}%
\pgfpathlineto{\pgfqpoint{2.068217in}{0.938848in}}%
\pgfpathlineto{\pgfqpoint{2.056583in}{0.934051in}}%
\pgfpathlineto{\pgfqpoint{2.051407in}{0.934194in}}%
\pgfpathlineto{\pgfqpoint{2.045055in}{0.940136in}}%
\pgfpathlineto{\pgfqpoint{2.045185in}{0.954206in}}%
\pgfpathlineto{\pgfqpoint{2.041706in}{0.956720in}}%
\pgfpathlineto{\pgfqpoint{2.055634in}{0.956575in}}%
\pgfpathlineto{\pgfqpoint{2.055870in}{0.984320in}}%
\pgfpathlineto{\pgfqpoint{2.062908in}{0.984282in}}%
\pgfpathlineto{\pgfqpoint{2.062941in}{0.991185in}}%
\pgfpathlineto{\pgfqpoint{2.073230in}{0.991141in}}%
\pgfpathlineto{\pgfqpoint{2.073255in}{0.998037in}}%
\pgfpathlineto{\pgfqpoint{2.110850in}{0.997922in}}%
\pgfpathlineto{\pgfqpoint{2.110869in}{0.984129in}}%
\pgfpathlineto{\pgfqpoint{2.141998in}{0.984277in}}%
\pgfpathlineto{\pgfqpoint{2.142324in}{1.001839in}}%
\pgfpathlineto{\pgfqpoint{2.151260in}{0.999398in}}%
\pgfpathlineto{\pgfqpoint{2.170781in}{0.999186in}}%
\pgfpathlineto{\pgfqpoint{2.176519in}{0.997121in}}%
\pgfpathlineto{\pgfqpoint{2.176495in}{0.972103in}}%
\pgfpathlineto{\pgfqpoint{2.155587in}{0.972321in}}%
\pgfpathlineto{\pgfqpoint{2.156410in}{0.959401in}}%
\pgfpathlineto{\pgfqpoint{2.165791in}{0.959401in}}%
\pgfpathlineto{\pgfqpoint{2.166804in}{0.946907in}}%
\pgfpathlineto{\pgfqpoint{2.169177in}{0.935220in}}%
\pgfpathlineto{\pgfqpoint{2.174128in}{0.929756in}}%
\pgfpathlineto{\pgfqpoint{2.175503in}{0.924555in}}%
\pgfpathclose%
\pgfusepath{fill}%
\end{pgfscope}%
\begin{pgfscope}%
\pgfpathrectangle{\pgfqpoint{0.100000in}{0.100000in}}{\pgfqpoint{3.608454in}{2.310000in}}%
\pgfusepath{clip}%
\pgfsetbuttcap%
\pgfsetmiterjoin%
\definecolor{currentfill}{rgb}{0.000000,0.282353,0.858824}%
\pgfsetfillcolor{currentfill}%
\pgfsetlinewidth{0.000000pt}%
\definecolor{currentstroke}{rgb}{0.000000,0.000000,0.000000}%
\pgfsetstrokecolor{currentstroke}%
\pgfsetstrokeopacity{0.000000}%
\pgfsetdash{}{0pt}%
\pgfpathmoveto{\pgfqpoint{2.138316in}{1.510307in}}%
\pgfpathlineto{\pgfqpoint{2.124622in}{1.510289in}}%
\pgfpathlineto{\pgfqpoint{2.124501in}{1.545418in}}%
\pgfpathlineto{\pgfqpoint{2.121931in}{1.545408in}}%
\pgfpathlineto{\pgfqpoint{2.121832in}{1.566001in}}%
\pgfpathlineto{\pgfqpoint{2.156052in}{1.566272in}}%
\pgfpathlineto{\pgfqpoint{2.155774in}{1.593709in}}%
\pgfpathlineto{\pgfqpoint{2.196902in}{1.594193in}}%
\pgfpathlineto{\pgfqpoint{2.210467in}{1.594411in}}%
\pgfpathlineto{\pgfqpoint{2.210987in}{1.567017in}}%
\pgfpathlineto{\pgfqpoint{2.204148in}{1.566920in}}%
\pgfpathlineto{\pgfqpoint{2.204489in}{1.546166in}}%
\pgfpathlineto{\pgfqpoint{2.205803in}{1.538835in}}%
\pgfpathlineto{\pgfqpoint{2.206358in}{1.511376in}}%
\pgfpathlineto{\pgfqpoint{2.192689in}{1.511198in}}%
\pgfpathlineto{\pgfqpoint{2.172375in}{1.510519in}}%
\pgfpathlineto{\pgfqpoint{2.138316in}{1.510307in}}%
\pgfpathclose%
\pgfusepath{fill}%
\end{pgfscope}%
\begin{pgfscope}%
\pgfpathrectangle{\pgfqpoint{0.100000in}{0.100000in}}{\pgfqpoint{3.608454in}{2.310000in}}%
\pgfusepath{clip}%
\pgfsetbuttcap%
\pgfsetmiterjoin%
\definecolor{currentfill}{rgb}{0.000000,0.505882,0.747059}%
\pgfsetfillcolor{currentfill}%
\pgfsetlinewidth{0.000000pt}%
\definecolor{currentstroke}{rgb}{0.000000,0.000000,0.000000}%
\pgfsetstrokecolor{currentstroke}%
\pgfsetstrokeopacity{0.000000}%
\pgfsetdash{}{0pt}%
\pgfpathmoveto{\pgfqpoint{1.782538in}{1.034546in}}%
\pgfpathlineto{\pgfqpoint{1.781300in}{1.012079in}}%
\pgfpathlineto{\pgfqpoint{1.754289in}{1.013713in}}%
\pgfpathlineto{\pgfqpoint{1.746176in}{1.014219in}}%
\pgfpathlineto{\pgfqpoint{1.748391in}{1.048651in}}%
\pgfpathlineto{\pgfqpoint{1.750416in}{1.083136in}}%
\pgfpathlineto{\pgfqpoint{1.785102in}{1.081107in}}%
\pgfpathlineto{\pgfqpoint{1.782538in}{1.034546in}}%
\pgfpathclose%
\pgfusepath{fill}%
\end{pgfscope}%
\begin{pgfscope}%
\pgfpathrectangle{\pgfqpoint{0.100000in}{0.100000in}}{\pgfqpoint{3.608454in}{2.310000in}}%
\pgfusepath{clip}%
\pgfsetbuttcap%
\pgfsetmiterjoin%
\definecolor{currentfill}{rgb}{0.000000,0.486275,0.756863}%
\pgfsetfillcolor{currentfill}%
\pgfsetlinewidth{0.000000pt}%
\definecolor{currentstroke}{rgb}{0.000000,0.000000,0.000000}%
\pgfsetstrokecolor{currentstroke}%
\pgfsetstrokeopacity{0.000000}%
\pgfsetdash{}{0pt}%
\pgfpathmoveto{\pgfqpoint{1.893386in}{1.070793in}}%
\pgfpathlineto{\pgfqpoint{1.873435in}{1.071561in}}%
\pgfpathlineto{\pgfqpoint{1.873684in}{1.092215in}}%
\pgfpathlineto{\pgfqpoint{1.825972in}{1.094304in}}%
\pgfpathlineto{\pgfqpoint{1.826418in}{1.110191in}}%
\pgfpathlineto{\pgfqpoint{1.826923in}{1.122162in}}%
\pgfpathlineto{\pgfqpoint{1.854121in}{1.120697in}}%
\pgfpathlineto{\pgfqpoint{1.855366in}{1.155288in}}%
\pgfpathlineto{\pgfqpoint{1.844616in}{1.170342in}}%
\pgfpathlineto{\pgfqpoint{1.837450in}{1.170858in}}%
\pgfpathlineto{\pgfqpoint{1.830118in}{1.178718in}}%
\pgfpathlineto{\pgfqpoint{1.825437in}{1.188388in}}%
\pgfpathlineto{\pgfqpoint{1.854168in}{1.187060in}}%
\pgfpathlineto{\pgfqpoint{1.895416in}{1.185262in}}%
\pgfpathlineto{\pgfqpoint{1.951284in}{1.183534in}}%
\pgfpathlineto{\pgfqpoint{1.950456in}{1.151489in}}%
\pgfpathlineto{\pgfqpoint{1.949595in}{1.117555in}}%
\pgfpathlineto{\pgfqpoint{1.901714in}{1.119083in}}%
\pgfpathlineto{\pgfqpoint{1.900701in}{1.084287in}}%
\pgfpathlineto{\pgfqpoint{1.893880in}{1.084537in}}%
\pgfpathlineto{\pgfqpoint{1.893386in}{1.070793in}}%
\pgfpathclose%
\pgfusepath{fill}%
\end{pgfscope}%
\begin{pgfscope}%
\pgfpathrectangle{\pgfqpoint{0.100000in}{0.100000in}}{\pgfqpoint{3.608454in}{2.310000in}}%
\pgfusepath{clip}%
\pgfsetbuttcap%
\pgfsetmiterjoin%
\definecolor{currentfill}{rgb}{0.000000,0.678431,0.660784}%
\pgfsetfillcolor{currentfill}%
\pgfsetlinewidth{0.000000pt}%
\definecolor{currentstroke}{rgb}{0.000000,0.000000,0.000000}%
\pgfsetstrokecolor{currentstroke}%
\pgfsetstrokeopacity{0.000000}%
\pgfsetdash{}{0pt}%
\pgfpathmoveto{\pgfqpoint{1.068484in}{1.440713in}}%
\pgfpathlineto{\pgfqpoint{1.064107in}{1.434490in}}%
\pgfpathlineto{\pgfqpoint{1.060054in}{1.435184in}}%
\pgfpathlineto{\pgfqpoint{1.053099in}{1.429145in}}%
\pgfpathlineto{\pgfqpoint{1.051426in}{1.418402in}}%
\pgfpathlineto{\pgfqpoint{1.045136in}{1.415471in}}%
\pgfpathlineto{\pgfqpoint{1.036855in}{1.417471in}}%
\pgfpathlineto{\pgfqpoint{1.031123in}{1.409823in}}%
\pgfpathlineto{\pgfqpoint{1.000345in}{1.416064in}}%
\pgfpathlineto{\pgfqpoint{0.938194in}{1.428823in}}%
\pgfpathlineto{\pgfqpoint{0.939931in}{1.436894in}}%
\pgfpathlineto{\pgfqpoint{0.954290in}{1.503771in}}%
\pgfpathlineto{\pgfqpoint{0.960301in}{1.531851in}}%
\pgfpathlineto{\pgfqpoint{1.034454in}{1.516384in}}%
\pgfpathlineto{\pgfqpoint{1.061605in}{1.511149in}}%
\pgfpathlineto{\pgfqpoint{1.072847in}{1.517633in}}%
\pgfpathlineto{\pgfqpoint{1.077321in}{1.529704in}}%
\pgfpathlineto{\pgfqpoint{1.075863in}{1.537379in}}%
\pgfpathlineto{\pgfqpoint{1.077577in}{1.552756in}}%
\pgfpathlineto{\pgfqpoint{1.087943in}{1.551215in}}%
\pgfpathlineto{\pgfqpoint{1.090523in}{1.546657in}}%
\pgfpathlineto{\pgfqpoint{1.105133in}{1.552914in}}%
\pgfpathlineto{\pgfqpoint{1.116941in}{1.552445in}}%
\pgfpathlineto{\pgfqpoint{1.112401in}{1.545424in}}%
\pgfpathlineto{\pgfqpoint{1.120187in}{1.531417in}}%
\pgfpathlineto{\pgfqpoint{1.129198in}{1.527981in}}%
\pgfpathlineto{\pgfqpoint{1.131819in}{1.519099in}}%
\pgfpathlineto{\pgfqpoint{1.128853in}{1.510727in}}%
\pgfpathlineto{\pgfqpoint{1.137161in}{1.499935in}}%
\pgfpathlineto{\pgfqpoint{1.136590in}{1.496867in}}%
\pgfpathlineto{\pgfqpoint{1.148005in}{1.494799in}}%
\pgfpathlineto{\pgfqpoint{1.148850in}{1.487724in}}%
\pgfpathlineto{\pgfqpoint{1.101981in}{1.496172in}}%
\pgfpathlineto{\pgfqpoint{1.104984in}{1.493265in}}%
\pgfpathlineto{\pgfqpoint{1.095705in}{1.487818in}}%
\pgfpathlineto{\pgfqpoint{1.092637in}{1.469533in}}%
\pgfpathlineto{\pgfqpoint{1.089487in}{1.463907in}}%
\pgfpathlineto{\pgfqpoint{1.078936in}{1.465774in}}%
\pgfpathlineto{\pgfqpoint{1.072686in}{1.463362in}}%
\pgfpathlineto{\pgfqpoint{1.068484in}{1.440713in}}%
\pgfpathclose%
\pgfusepath{fill}%
\end{pgfscope}%
\begin{pgfscope}%
\pgfpathrectangle{\pgfqpoint{0.100000in}{0.100000in}}{\pgfqpoint{3.608454in}{2.310000in}}%
\pgfusepath{clip}%
\pgfsetbuttcap%
\pgfsetmiterjoin%
\definecolor{currentfill}{rgb}{0.000000,0.560784,0.719608}%
\pgfsetfillcolor{currentfill}%
\pgfsetlinewidth{0.000000pt}%
\definecolor{currentstroke}{rgb}{0.000000,0.000000,0.000000}%
\pgfsetstrokecolor{currentstroke}%
\pgfsetstrokeopacity{0.000000}%
\pgfsetdash{}{0pt}%
\pgfpathmoveto{\pgfqpoint{2.831040in}{1.546066in}}%
\pgfpathlineto{\pgfqpoint{2.825402in}{1.579904in}}%
\pgfpathlineto{\pgfqpoint{2.818905in}{1.581252in}}%
\pgfpathlineto{\pgfqpoint{2.833971in}{1.586462in}}%
\pgfpathlineto{\pgfqpoint{2.840470in}{1.580694in}}%
\pgfpathlineto{\pgfqpoint{2.847896in}{1.579132in}}%
\pgfpathlineto{\pgfqpoint{2.860187in}{1.585033in}}%
\pgfpathlineto{\pgfqpoint{2.873522in}{1.593422in}}%
\pgfpathlineto{\pgfqpoint{2.876219in}{1.593055in}}%
\pgfpathlineto{\pgfqpoint{2.877765in}{1.580902in}}%
\pgfpathlineto{\pgfqpoint{2.883315in}{1.581642in}}%
\pgfpathlineto{\pgfqpoint{2.884117in}{1.575727in}}%
\pgfpathlineto{\pgfqpoint{2.878552in}{1.574944in}}%
\pgfpathlineto{\pgfqpoint{2.879290in}{1.569013in}}%
\pgfpathlineto{\pgfqpoint{2.873516in}{1.568184in}}%
\pgfpathlineto{\pgfqpoint{2.868477in}{1.562448in}}%
\pgfpathlineto{\pgfqpoint{2.869199in}{1.556678in}}%
\pgfpathlineto{\pgfqpoint{2.853507in}{1.554639in}}%
\pgfpathlineto{\pgfqpoint{2.854535in}{1.548999in}}%
\pgfpathlineto{\pgfqpoint{2.831040in}{1.546066in}}%
\pgfpathclose%
\pgfusepath{fill}%
\end{pgfscope}%
\begin{pgfscope}%
\pgfpathrectangle{\pgfqpoint{0.100000in}{0.100000in}}{\pgfqpoint{3.608454in}{2.310000in}}%
\pgfusepath{clip}%
\pgfsetbuttcap%
\pgfsetmiterjoin%
\definecolor{currentfill}{rgb}{0.000000,0.611765,0.694118}%
\pgfsetfillcolor{currentfill}%
\pgfsetlinewidth{0.000000pt}%
\definecolor{currentstroke}{rgb}{0.000000,0.000000,0.000000}%
\pgfsetstrokecolor{currentstroke}%
\pgfsetstrokeopacity{0.000000}%
\pgfsetdash{}{0pt}%
\pgfpathmoveto{\pgfqpoint{1.611520in}{1.962062in}}%
\pgfpathlineto{\pgfqpoint{1.606786in}{1.914838in}}%
\pgfpathlineto{\pgfqpoint{1.602057in}{1.917205in}}%
\pgfpathlineto{\pgfqpoint{1.587963in}{1.918755in}}%
\pgfpathlineto{\pgfqpoint{1.588691in}{1.925458in}}%
\pgfpathlineto{\pgfqpoint{1.581763in}{1.926153in}}%
\pgfpathlineto{\pgfqpoint{1.582466in}{1.932625in}}%
\pgfpathlineto{\pgfqpoint{1.562162in}{1.934829in}}%
\pgfpathlineto{\pgfqpoint{1.564057in}{1.962319in}}%
\pgfpathlineto{\pgfqpoint{1.573489in}{1.961301in}}%
\pgfpathlineto{\pgfqpoint{1.574985in}{1.975097in}}%
\pgfpathlineto{\pgfqpoint{1.582076in}{1.976639in}}%
\pgfpathlineto{\pgfqpoint{1.592373in}{1.975555in}}%
\pgfpathlineto{\pgfqpoint{1.595495in}{1.971747in}}%
\pgfpathlineto{\pgfqpoint{1.612308in}{1.970040in}}%
\pgfpathlineto{\pgfqpoint{1.611520in}{1.962062in}}%
\pgfpathclose%
\pgfusepath{fill}%
\end{pgfscope}%
\begin{pgfscope}%
\pgfpathrectangle{\pgfqpoint{0.100000in}{0.100000in}}{\pgfqpoint{3.608454in}{2.310000in}}%
\pgfusepath{clip}%
\pgfsetbuttcap%
\pgfsetmiterjoin%
\definecolor{currentfill}{rgb}{0.000000,0.560784,0.719608}%
\pgfsetfillcolor{currentfill}%
\pgfsetlinewidth{0.000000pt}%
\definecolor{currentstroke}{rgb}{0.000000,0.000000,0.000000}%
\pgfsetstrokecolor{currentstroke}%
\pgfsetstrokeopacity{0.000000}%
\pgfsetdash{}{0pt}%
\pgfpathmoveto{\pgfqpoint{3.298349in}{1.383037in}}%
\pgfpathlineto{\pgfqpoint{3.291222in}{1.379099in}}%
\pgfpathlineto{\pgfqpoint{3.288344in}{1.383217in}}%
\pgfpathlineto{\pgfqpoint{3.291045in}{1.389057in}}%
\pgfpathlineto{\pgfqpoint{3.289366in}{1.396494in}}%
\pgfpathlineto{\pgfqpoint{3.281967in}{1.394943in}}%
\pgfpathlineto{\pgfqpoint{3.282347in}{1.403895in}}%
\pgfpathlineto{\pgfqpoint{3.281272in}{1.408524in}}%
\pgfpathlineto{\pgfqpoint{3.284538in}{1.423341in}}%
\pgfpathlineto{\pgfqpoint{3.290599in}{1.431089in}}%
\pgfpathlineto{\pgfqpoint{3.274847in}{1.487194in}}%
\pgfpathlineto{\pgfqpoint{3.281434in}{1.488221in}}%
\pgfpathlineto{\pgfqpoint{3.288460in}{1.495799in}}%
\pgfpathlineto{\pgfqpoint{3.293773in}{1.492637in}}%
\pgfpathlineto{\pgfqpoint{3.298424in}{1.484780in}}%
\pgfpathlineto{\pgfqpoint{3.300159in}{1.474127in}}%
\pgfpathlineto{\pgfqpoint{3.304555in}{1.471386in}}%
\pgfpathlineto{\pgfqpoint{3.308496in}{1.463629in}}%
\pgfpathlineto{\pgfqpoint{3.317538in}{1.456649in}}%
\pgfpathlineto{\pgfqpoint{3.324360in}{1.455596in}}%
\pgfpathlineto{\pgfqpoint{3.331882in}{1.431286in}}%
\pgfpathlineto{\pgfqpoint{3.329200in}{1.401923in}}%
\pgfpathlineto{\pgfqpoint{3.327417in}{1.395926in}}%
\pgfpathlineto{\pgfqpoint{3.304739in}{1.388327in}}%
\pgfpathlineto{\pgfqpoint{3.298349in}{1.383037in}}%
\pgfpathclose%
\pgfusepath{fill}%
\end{pgfscope}%
\begin{pgfscope}%
\pgfpathrectangle{\pgfqpoint{0.100000in}{0.100000in}}{\pgfqpoint{3.608454in}{2.310000in}}%
\pgfusepath{clip}%
\pgfsetbuttcap%
\pgfsetmiterjoin%
\definecolor{currentfill}{rgb}{0.000000,0.443137,0.778431}%
\pgfsetfillcolor{currentfill}%
\pgfsetlinewidth{0.000000pt}%
\definecolor{currentstroke}{rgb}{0.000000,0.000000,0.000000}%
\pgfsetstrokecolor{currentstroke}%
\pgfsetstrokeopacity{0.000000}%
\pgfsetdash{}{0pt}%
\pgfpathmoveto{\pgfqpoint{2.115090in}{1.593467in}}%
\pgfpathlineto{\pgfqpoint{2.067449in}{1.593735in}}%
\pgfpathlineto{\pgfqpoint{2.067679in}{1.614586in}}%
\pgfpathlineto{\pgfqpoint{2.063540in}{1.621479in}}%
\pgfpathlineto{\pgfqpoint{2.056624in}{1.621513in}}%
\pgfpathlineto{\pgfqpoint{2.056914in}{1.649093in}}%
\pgfpathlineto{\pgfqpoint{2.084271in}{1.648958in}}%
\pgfpathlineto{\pgfqpoint{2.111776in}{1.648912in}}%
\pgfpathlineto{\pgfqpoint{2.111772in}{1.614418in}}%
\pgfpathlineto{\pgfqpoint{2.115152in}{1.614399in}}%
\pgfpathlineto{\pgfqpoint{2.115090in}{1.593467in}}%
\pgfpathclose%
\pgfusepath{fill}%
\end{pgfscope}%
\begin{pgfscope}%
\pgfpathrectangle{\pgfqpoint{0.100000in}{0.100000in}}{\pgfqpoint{3.608454in}{2.310000in}}%
\pgfusepath{clip}%
\pgfsetbuttcap%
\pgfsetmiterjoin%
\definecolor{currentfill}{rgb}{0.000000,0.698039,0.650980}%
\pgfsetfillcolor{currentfill}%
\pgfsetlinewidth{0.000000pt}%
\definecolor{currentstroke}{rgb}{0.000000,0.000000,0.000000}%
\pgfsetstrokecolor{currentstroke}%
\pgfsetstrokeopacity{0.000000}%
\pgfsetdash{}{0pt}%
\pgfpathmoveto{\pgfqpoint{1.194163in}{1.451751in}}%
\pgfpathlineto{\pgfqpoint{1.249416in}{1.442841in}}%
\pgfpathlineto{\pgfqpoint{1.253203in}{1.445105in}}%
\pgfpathlineto{\pgfqpoint{1.251645in}{1.434868in}}%
\pgfpathlineto{\pgfqpoint{1.240533in}{1.367269in}}%
\pgfpathlineto{\pgfqpoint{1.181575in}{1.376785in}}%
\pgfpathlineto{\pgfqpoint{1.183363in}{1.379999in}}%
\pgfpathlineto{\pgfqpoint{1.179215in}{1.395609in}}%
\pgfpathlineto{\pgfqpoint{1.181137in}{1.398548in}}%
\pgfpathlineto{\pgfqpoint{1.177987in}{1.406555in}}%
\pgfpathlineto{\pgfqpoint{1.182094in}{1.423685in}}%
\pgfpathlineto{\pgfqpoint{1.189342in}{1.433929in}}%
\pgfpathlineto{\pgfqpoint{1.188609in}{1.438112in}}%
\pgfpathlineto{\pgfqpoint{1.192855in}{1.443180in}}%
\pgfpathlineto{\pgfqpoint{1.194163in}{1.451751in}}%
\pgfpathclose%
\pgfusepath{fill}%
\end{pgfscope}%
\begin{pgfscope}%
\pgfpathrectangle{\pgfqpoint{0.100000in}{0.100000in}}{\pgfqpoint{3.608454in}{2.310000in}}%
\pgfusepath{clip}%
\pgfsetbuttcap%
\pgfsetmiterjoin%
\definecolor{currentfill}{rgb}{0.000000,0.564706,0.717647}%
\pgfsetfillcolor{currentfill}%
\pgfsetlinewidth{0.000000pt}%
\definecolor{currentstroke}{rgb}{0.000000,0.000000,0.000000}%
\pgfsetstrokecolor{currentstroke}%
\pgfsetstrokeopacity{0.000000}%
\pgfsetdash{}{0pt}%
\pgfpathmoveto{\pgfqpoint{2.780892in}{1.025876in}}%
\pgfpathlineto{\pgfqpoint{2.779095in}{1.040030in}}%
\pgfpathlineto{\pgfqpoint{2.782890in}{1.048115in}}%
\pgfpathlineto{\pgfqpoint{2.780160in}{1.050381in}}%
\pgfpathlineto{\pgfqpoint{2.779057in}{1.060719in}}%
\pgfpathlineto{\pgfqpoint{2.798375in}{1.062974in}}%
\pgfpathlineto{\pgfqpoint{2.808303in}{1.061863in}}%
\pgfpathlineto{\pgfqpoint{2.814885in}{1.050035in}}%
\pgfpathlineto{\pgfqpoint{2.812171in}{1.037511in}}%
\pgfpathlineto{\pgfqpoint{2.810594in}{1.033632in}}%
\pgfpathlineto{\pgfqpoint{2.801686in}{1.031063in}}%
\pgfpathlineto{\pgfqpoint{2.799254in}{1.027968in}}%
\pgfpathlineto{\pgfqpoint{2.780892in}{1.025876in}}%
\pgfpathclose%
\pgfusepath{fill}%
\end{pgfscope}%
\begin{pgfscope}%
\pgfpathrectangle{\pgfqpoint{0.100000in}{0.100000in}}{\pgfqpoint{3.608454in}{2.310000in}}%
\pgfusepath{clip}%
\pgfsetbuttcap%
\pgfsetmiterjoin%
\definecolor{currentfill}{rgb}{0.000000,0.917647,0.541176}%
\pgfsetfillcolor{currentfill}%
\pgfsetlinewidth{0.000000pt}%
\definecolor{currentstroke}{rgb}{0.000000,0.000000,0.000000}%
\pgfsetstrokecolor{currentstroke}%
\pgfsetstrokeopacity{0.000000}%
\pgfsetdash{}{0pt}%
\pgfpathmoveto{\pgfqpoint{0.427334in}{1.927274in}}%
\pgfpathlineto{\pgfqpoint{0.440684in}{1.949326in}}%
\pgfpathlineto{\pgfqpoint{0.456163in}{1.968351in}}%
\pgfpathlineto{\pgfqpoint{0.467729in}{1.989714in}}%
\pgfpathlineto{\pgfqpoint{0.479271in}{1.985850in}}%
\pgfpathlineto{\pgfqpoint{0.487714in}{1.989945in}}%
\pgfpathlineto{\pgfqpoint{0.498412in}{1.984367in}}%
\pgfpathlineto{\pgfqpoint{0.501091in}{1.975934in}}%
\pgfpathlineto{\pgfqpoint{0.509770in}{1.969016in}}%
\pgfpathlineto{\pgfqpoint{0.521206in}{1.965301in}}%
\pgfpathlineto{\pgfqpoint{0.517042in}{1.952171in}}%
\pgfpathlineto{\pgfqpoint{0.517068in}{1.946672in}}%
\pgfpathlineto{\pgfqpoint{0.537209in}{1.940659in}}%
\pgfpathlineto{\pgfqpoint{0.534705in}{1.932520in}}%
\pgfpathlineto{\pgfqpoint{0.568235in}{1.922352in}}%
\pgfpathlineto{\pgfqpoint{0.572489in}{1.914545in}}%
\pgfpathlineto{\pgfqpoint{0.572519in}{1.906519in}}%
\pgfpathlineto{\pgfqpoint{0.565189in}{1.899442in}}%
\pgfpathlineto{\pgfqpoint{0.562002in}{1.893971in}}%
\pgfpathlineto{\pgfqpoint{0.551234in}{1.896670in}}%
\pgfpathlineto{\pgfqpoint{0.549594in}{1.891292in}}%
\pgfpathlineto{\pgfqpoint{0.543027in}{1.893318in}}%
\pgfpathlineto{\pgfqpoint{0.538396in}{1.890064in}}%
\pgfpathlineto{\pgfqpoint{0.532832in}{1.891746in}}%
\pgfpathlineto{\pgfqpoint{0.516046in}{1.883757in}}%
\pgfpathlineto{\pgfqpoint{0.503506in}{1.885466in}}%
\pgfpathlineto{\pgfqpoint{0.495724in}{1.888283in}}%
\pgfpathlineto{\pgfqpoint{0.491499in}{1.886385in}}%
\pgfpathlineto{\pgfqpoint{0.483261in}{1.890283in}}%
\pgfpathlineto{\pgfqpoint{0.478228in}{1.889504in}}%
\pgfpathlineto{\pgfqpoint{0.465159in}{1.900474in}}%
\pgfpathlineto{\pgfqpoint{0.456477in}{1.902564in}}%
\pgfpathlineto{\pgfqpoint{0.447212in}{1.898187in}}%
\pgfpathlineto{\pgfqpoint{0.438998in}{1.899672in}}%
\pgfpathlineto{\pgfqpoint{0.443285in}{1.914029in}}%
\pgfpathlineto{\pgfqpoint{0.439586in}{1.923181in}}%
\pgfpathlineto{\pgfqpoint{0.427334in}{1.927274in}}%
\pgfpathclose%
\pgfusepath{fill}%
\end{pgfscope}%
\begin{pgfscope}%
\pgfpathrectangle{\pgfqpoint{0.100000in}{0.100000in}}{\pgfqpoint{3.608454in}{2.310000in}}%
\pgfusepath{clip}%
\pgfsetbuttcap%
\pgfsetmiterjoin%
\definecolor{currentfill}{rgb}{0.000000,0.666667,0.666667}%
\pgfsetfillcolor{currentfill}%
\pgfsetlinewidth{0.000000pt}%
\definecolor{currentstroke}{rgb}{0.000000,0.000000,0.000000}%
\pgfsetstrokecolor{currentstroke}%
\pgfsetstrokeopacity{0.000000}%
\pgfsetdash{}{0pt}%
\pgfpathmoveto{\pgfqpoint{3.106445in}{0.607053in}}%
\pgfpathlineto{\pgfqpoint{3.090036in}{0.604447in}}%
\pgfpathlineto{\pgfqpoint{3.090875in}{0.590177in}}%
\pgfpathlineto{\pgfqpoint{3.080990in}{0.605844in}}%
\pgfpathlineto{\pgfqpoint{3.078312in}{0.602797in}}%
\pgfpathlineto{\pgfqpoint{3.071577in}{0.605034in}}%
\pgfpathlineto{\pgfqpoint{3.065010in}{0.604224in}}%
\pgfpathlineto{\pgfqpoint{3.060692in}{0.615927in}}%
\pgfpathlineto{\pgfqpoint{3.049410in}{0.623089in}}%
\pgfpathlineto{\pgfqpoint{3.046213in}{0.629170in}}%
\pgfpathlineto{\pgfqpoint{3.040376in}{0.630406in}}%
\pgfpathlineto{\pgfqpoint{3.037437in}{0.635958in}}%
\pgfpathlineto{\pgfqpoint{3.029457in}{0.642400in}}%
\pgfpathlineto{\pgfqpoint{3.025702in}{0.651277in}}%
\pgfpathlineto{\pgfqpoint{3.020704in}{0.653264in}}%
\pgfpathlineto{\pgfqpoint{3.006579in}{0.647225in}}%
\pgfpathlineto{\pgfqpoint{3.004229in}{0.666876in}}%
\pgfpathlineto{\pgfqpoint{3.011416in}{0.670241in}}%
\pgfpathlineto{\pgfqpoint{3.019158in}{0.678567in}}%
\pgfpathlineto{\pgfqpoint{3.035034in}{0.681162in}}%
\pgfpathlineto{\pgfqpoint{3.039796in}{0.675377in}}%
\pgfpathlineto{\pgfqpoint{3.041523in}{0.664546in}}%
\pgfpathlineto{\pgfqpoint{3.055227in}{0.666862in}}%
\pgfpathlineto{\pgfqpoint{3.062330in}{0.671592in}}%
\pgfpathlineto{\pgfqpoint{3.078208in}{0.644633in}}%
\pgfpathlineto{\pgfqpoint{3.086273in}{0.632061in}}%
\pgfpathlineto{\pgfqpoint{3.106445in}{0.607053in}}%
\pgfpathclose%
\pgfusepath{fill}%
\end{pgfscope}%
\begin{pgfscope}%
\pgfpathrectangle{\pgfqpoint{0.100000in}{0.100000in}}{\pgfqpoint{3.608454in}{2.310000in}}%
\pgfusepath{clip}%
\pgfsetbuttcap%
\pgfsetmiterjoin%
\definecolor{currentfill}{rgb}{0.000000,0.549020,0.725490}%
\pgfsetfillcolor{currentfill}%
\pgfsetlinewidth{0.000000pt}%
\definecolor{currentstroke}{rgb}{0.000000,0.000000,0.000000}%
\pgfsetstrokecolor{currentstroke}%
\pgfsetstrokeopacity{0.000000}%
\pgfsetdash{}{0pt}%
\pgfpathmoveto{\pgfqpoint{0.914070in}{1.089859in}}%
\pgfpathlineto{\pgfqpoint{0.900636in}{1.091652in}}%
\pgfpathlineto{\pgfqpoint{0.895182in}{1.086583in}}%
\pgfpathlineto{\pgfqpoint{0.868397in}{1.094701in}}%
\pgfpathlineto{\pgfqpoint{0.862231in}{1.099514in}}%
\pgfpathlineto{\pgfqpoint{0.851945in}{1.113704in}}%
\pgfpathlineto{\pgfqpoint{0.848041in}{1.126325in}}%
\pgfpathlineto{\pgfqpoint{0.847819in}{1.135796in}}%
\pgfpathlineto{\pgfqpoint{0.843543in}{1.141162in}}%
\pgfpathlineto{\pgfqpoint{0.840274in}{1.151795in}}%
\pgfpathlineto{\pgfqpoint{0.842429in}{1.160488in}}%
\pgfpathlineto{\pgfqpoint{0.792611in}{1.237224in}}%
\pgfpathlineto{\pgfqpoint{0.780428in}{1.255922in}}%
\pgfpathlineto{\pgfqpoint{0.720005in}{1.349216in}}%
\pgfpathlineto{\pgfqpoint{0.688870in}{1.397366in}}%
\pgfpathlineto{\pgfqpoint{0.661400in}{1.439746in}}%
\pgfpathlineto{\pgfqpoint{0.665974in}{1.438337in}}%
\pgfpathlineto{\pgfqpoint{0.717194in}{1.472445in}}%
\pgfpathlineto{\pgfqpoint{0.695807in}{1.514308in}}%
\pgfpathlineto{\pgfqpoint{0.697419in}{1.520493in}}%
\pgfpathlineto{\pgfqpoint{0.718561in}{1.521037in}}%
\pgfpathlineto{\pgfqpoint{0.724276in}{1.521184in}}%
\pgfpathlineto{\pgfqpoint{0.752002in}{1.519867in}}%
\pgfpathlineto{\pgfqpoint{0.795453in}{1.508995in}}%
\pgfpathlineto{\pgfqpoint{0.836825in}{1.499170in}}%
\pgfpathlineto{\pgfqpoint{0.882599in}{1.449402in}}%
\pgfpathlineto{\pgfqpoint{0.939931in}{1.436894in}}%
\pgfpathlineto{\pgfqpoint{0.938194in}{1.428823in}}%
\pgfpathlineto{\pgfqpoint{0.925741in}{1.369767in}}%
\pgfpathlineto{\pgfqpoint{0.912295in}{1.307222in}}%
\pgfpathlineto{\pgfqpoint{0.983512in}{1.292540in}}%
\pgfpathlineto{\pgfqpoint{1.005762in}{1.288169in}}%
\pgfpathlineto{\pgfqpoint{0.994707in}{1.264046in}}%
\pgfpathlineto{\pgfqpoint{0.991172in}{1.251855in}}%
\pgfpathlineto{\pgfqpoint{0.991491in}{1.243095in}}%
\pgfpathlineto{\pgfqpoint{0.986930in}{1.238736in}}%
\pgfpathlineto{\pgfqpoint{0.978736in}{1.235692in}}%
\pgfpathlineto{\pgfqpoint{0.961661in}{1.233095in}}%
\pgfpathlineto{\pgfqpoint{0.951426in}{1.230336in}}%
\pgfpathlineto{\pgfqpoint{0.949023in}{1.226105in}}%
\pgfpathlineto{\pgfqpoint{0.943617in}{1.227923in}}%
\pgfpathlineto{\pgfqpoint{0.940117in}{1.223821in}}%
\pgfpathlineto{\pgfqpoint{0.941162in}{1.217180in}}%
\pgfpathlineto{\pgfqpoint{0.937616in}{1.205156in}}%
\pgfpathlineto{\pgfqpoint{0.933245in}{1.183825in}}%
\pgfpathlineto{\pgfqpoint{0.914070in}{1.089859in}}%
\pgfpathclose%
\pgfusepath{fill}%
\end{pgfscope}%
\begin{pgfscope}%
\pgfpathrectangle{\pgfqpoint{0.100000in}{0.100000in}}{\pgfqpoint{3.608454in}{2.310000in}}%
\pgfusepath{clip}%
\pgfsetbuttcap%
\pgfsetmiterjoin%
\definecolor{currentfill}{rgb}{0.000000,0.552941,0.723529}%
\pgfsetfillcolor{currentfill}%
\pgfsetlinewidth{0.000000pt}%
\definecolor{currentstroke}{rgb}{0.000000,0.000000,0.000000}%
\pgfsetstrokecolor{currentstroke}%
\pgfsetstrokeopacity{0.000000}%
\pgfsetdash{}{0pt}%
\pgfpathmoveto{\pgfqpoint{2.956251in}{1.497406in}}%
\pgfpathlineto{\pgfqpoint{2.935644in}{1.495065in}}%
\pgfpathlineto{\pgfqpoint{2.928425in}{1.497499in}}%
\pgfpathlineto{\pgfqpoint{2.908656in}{1.495198in}}%
\pgfpathlineto{\pgfqpoint{2.909284in}{1.489578in}}%
\pgfpathlineto{\pgfqpoint{2.878220in}{1.486351in}}%
\pgfpathlineto{\pgfqpoint{2.876912in}{1.491928in}}%
\pgfpathlineto{\pgfqpoint{2.874225in}{1.518225in}}%
\pgfpathlineto{\pgfqpoint{2.871819in}{1.517489in}}%
\pgfpathlineto{\pgfqpoint{2.870685in}{1.525264in}}%
\pgfpathlineto{\pgfqpoint{2.876286in}{1.526111in}}%
\pgfpathlineto{\pgfqpoint{2.872479in}{1.551436in}}%
\pgfpathlineto{\pgfqpoint{2.901003in}{1.555251in}}%
\pgfpathlineto{\pgfqpoint{2.901826in}{1.549489in}}%
\pgfpathlineto{\pgfqpoint{2.915375in}{1.550896in}}%
\pgfpathlineto{\pgfqpoint{2.914636in}{1.557327in}}%
\pgfpathlineto{\pgfqpoint{2.934085in}{1.560268in}}%
\pgfpathlineto{\pgfqpoint{2.937196in}{1.539921in}}%
\pgfpathlineto{\pgfqpoint{2.947280in}{1.541412in}}%
\pgfpathlineto{\pgfqpoint{2.948242in}{1.534882in}}%
\pgfpathlineto{\pgfqpoint{2.952109in}{1.531961in}}%
\pgfpathlineto{\pgfqpoint{2.949140in}{1.528048in}}%
\pgfpathlineto{\pgfqpoint{2.949522in}{1.517575in}}%
\pgfpathlineto{\pgfqpoint{2.954062in}{1.518144in}}%
\pgfpathlineto{\pgfqpoint{2.956251in}{1.497406in}}%
\pgfpathclose%
\pgfusepath{fill}%
\end{pgfscope}%
\begin{pgfscope}%
\pgfpathrectangle{\pgfqpoint{0.100000in}{0.100000in}}{\pgfqpoint{3.608454in}{2.310000in}}%
\pgfusepath{clip}%
\pgfsetbuttcap%
\pgfsetmiterjoin%
\definecolor{currentfill}{rgb}{0.000000,0.658824,0.670588}%
\pgfsetfillcolor{currentfill}%
\pgfsetlinewidth{0.000000pt}%
\definecolor{currentstroke}{rgb}{0.000000,0.000000,0.000000}%
\pgfsetstrokecolor{currentstroke}%
\pgfsetstrokeopacity{0.000000}%
\pgfsetdash{}{0pt}%
\pgfpathmoveto{\pgfqpoint{2.996900in}{1.203053in}}%
\pgfpathlineto{\pgfqpoint{3.004020in}{1.198138in}}%
\pgfpathlineto{\pgfqpoint{3.006498in}{1.176840in}}%
\pgfpathlineto{\pgfqpoint{2.993174in}{1.172573in}}%
\pgfpathlineto{\pgfqpoint{2.987201in}{1.172347in}}%
\pgfpathlineto{\pgfqpoint{2.978982in}{1.167861in}}%
\pgfpathlineto{\pgfqpoint{2.972912in}{1.172989in}}%
\pgfpathlineto{\pgfqpoint{2.963826in}{1.175254in}}%
\pgfpathlineto{\pgfqpoint{2.968126in}{1.182900in}}%
\pgfpathlineto{\pgfqpoint{2.955269in}{1.192553in}}%
\pgfpathlineto{\pgfqpoint{2.948361in}{1.195220in}}%
\pgfpathlineto{\pgfqpoint{2.950303in}{1.201290in}}%
\pgfpathlineto{\pgfqpoint{2.950033in}{1.210864in}}%
\pgfpathlineto{\pgfqpoint{2.970559in}{1.212860in}}%
\pgfpathlineto{\pgfqpoint{2.970773in}{1.210739in}}%
\pgfpathlineto{\pgfqpoint{2.980945in}{1.198576in}}%
\pgfpathlineto{\pgfqpoint{2.987922in}{1.204284in}}%
\pgfpathlineto{\pgfqpoint{2.996900in}{1.203053in}}%
\pgfpathclose%
\pgfusepath{fill}%
\end{pgfscope}%
\begin{pgfscope}%
\pgfpathrectangle{\pgfqpoint{0.100000in}{0.100000in}}{\pgfqpoint{3.608454in}{2.310000in}}%
\pgfusepath{clip}%
\pgfsetbuttcap%
\pgfsetmiterjoin%
\definecolor{currentfill}{rgb}{0.000000,0.364706,0.817647}%
\pgfsetfillcolor{currentfill}%
\pgfsetlinewidth{0.000000pt}%
\definecolor{currentstroke}{rgb}{0.000000,0.000000,0.000000}%
\pgfsetstrokecolor{currentstroke}%
\pgfsetstrokeopacity{0.000000}%
\pgfsetdash{}{0pt}%
\pgfpathmoveto{\pgfqpoint{1.655822in}{1.596346in}}%
\pgfpathlineto{\pgfqpoint{1.691656in}{1.593744in}}%
\pgfpathlineto{\pgfqpoint{1.689939in}{1.572676in}}%
\pgfpathlineto{\pgfqpoint{1.694791in}{1.572316in}}%
\pgfpathlineto{\pgfqpoint{1.692731in}{1.544894in}}%
\pgfpathlineto{\pgfqpoint{1.688556in}{1.545253in}}%
\pgfpathlineto{\pgfqpoint{1.687474in}{1.531563in}}%
\pgfpathlineto{\pgfqpoint{1.654575in}{1.534236in}}%
\pgfpathlineto{\pgfqpoint{1.655781in}{1.547900in}}%
\pgfpathlineto{\pgfqpoint{1.654457in}{1.551499in}}%
\pgfpathlineto{\pgfqpoint{1.656181in}{1.573226in}}%
\pgfpathlineto{\pgfqpoint{1.654522in}{1.581965in}}%
\pgfpathlineto{\pgfqpoint{1.655822in}{1.596346in}}%
\pgfpathclose%
\pgfusepath{fill}%
\end{pgfscope}%
\begin{pgfscope}%
\pgfpathrectangle{\pgfqpoint{0.100000in}{0.100000in}}{\pgfqpoint{3.608454in}{2.310000in}}%
\pgfusepath{clip}%
\pgfsetbuttcap%
\pgfsetmiterjoin%
\definecolor{currentfill}{rgb}{0.000000,0.615686,0.692157}%
\pgfsetfillcolor{currentfill}%
\pgfsetlinewidth{0.000000pt}%
\definecolor{currentstroke}{rgb}{0.000000,0.000000,0.000000}%
\pgfsetstrokecolor{currentstroke}%
\pgfsetstrokeopacity{0.000000}%
\pgfsetdash{}{0pt}%
\pgfpathmoveto{\pgfqpoint{2.722959in}{1.184982in}}%
\pgfpathlineto{\pgfqpoint{2.722324in}{1.177007in}}%
\pgfpathlineto{\pgfqpoint{2.713741in}{1.174022in}}%
\pgfpathlineto{\pgfqpoint{2.710970in}{1.165941in}}%
\pgfpathlineto{\pgfqpoint{2.705300in}{1.170782in}}%
\pgfpathlineto{\pgfqpoint{2.696790in}{1.173951in}}%
\pgfpathlineto{\pgfqpoint{2.690027in}{1.171438in}}%
\pgfpathlineto{\pgfqpoint{2.690217in}{1.180528in}}%
\pgfpathlineto{\pgfqpoint{2.678723in}{1.180579in}}%
\pgfpathlineto{\pgfqpoint{2.677443in}{1.189404in}}%
\pgfpathlineto{\pgfqpoint{2.664358in}{1.203424in}}%
\pgfpathlineto{\pgfqpoint{2.669252in}{1.213390in}}%
\pgfpathlineto{\pgfqpoint{2.669571in}{1.226318in}}%
\pgfpathlineto{\pgfqpoint{2.661850in}{1.235103in}}%
\pgfpathlineto{\pgfqpoint{2.665870in}{1.236725in}}%
\pgfpathlineto{\pgfqpoint{2.667849in}{1.244767in}}%
\pgfpathlineto{\pgfqpoint{2.692446in}{1.245040in}}%
\pgfpathlineto{\pgfqpoint{2.690908in}{1.235308in}}%
\pgfpathlineto{\pgfqpoint{2.692552in}{1.225974in}}%
\pgfpathlineto{\pgfqpoint{2.696385in}{1.221833in}}%
\pgfpathlineto{\pgfqpoint{2.703168in}{1.221287in}}%
\pgfpathlineto{\pgfqpoint{2.705489in}{1.213403in}}%
\pgfpathlineto{\pgfqpoint{2.709289in}{1.208321in}}%
\pgfpathlineto{\pgfqpoint{2.723224in}{1.208766in}}%
\pgfpathlineto{\pgfqpoint{2.724860in}{1.203323in}}%
\pgfpathlineto{\pgfqpoint{2.722959in}{1.184982in}}%
\pgfpathclose%
\pgfusepath{fill}%
\end{pgfscope}%
\begin{pgfscope}%
\pgfpathrectangle{\pgfqpoint{0.100000in}{0.100000in}}{\pgfqpoint{3.608454in}{2.310000in}}%
\pgfusepath{clip}%
\pgfsetbuttcap%
\pgfsetmiterjoin%
\definecolor{currentfill}{rgb}{0.000000,0.505882,0.747059}%
\pgfsetfillcolor{currentfill}%
\pgfsetlinewidth{0.000000pt}%
\definecolor{currentstroke}{rgb}{0.000000,0.000000,0.000000}%
\pgfsetstrokecolor{currentstroke}%
\pgfsetstrokeopacity{0.000000}%
\pgfsetdash{}{0pt}%
\pgfpathmoveto{\pgfqpoint{3.206484in}{1.172807in}}%
\pgfpathlineto{\pgfqpoint{3.191049in}{1.170900in}}%
\pgfpathlineto{\pgfqpoint{3.186933in}{1.171546in}}%
\pgfpathlineto{\pgfqpoint{3.181671in}{1.180528in}}%
\pgfpathlineto{\pgfqpoint{3.175634in}{1.186715in}}%
\pgfpathlineto{\pgfqpoint{3.185638in}{1.219606in}}%
\pgfpathlineto{\pgfqpoint{3.195538in}{1.217166in}}%
\pgfpathlineto{\pgfqpoint{3.206943in}{1.219779in}}%
\pgfpathlineto{\pgfqpoint{3.216999in}{1.216602in}}%
\pgfpathlineto{\pgfqpoint{3.218220in}{1.212547in}}%
\pgfpathlineto{\pgfqpoint{3.226514in}{1.211542in}}%
\pgfpathlineto{\pgfqpoint{3.231820in}{1.204552in}}%
\pgfpathlineto{\pgfqpoint{3.232624in}{1.197760in}}%
\pgfpathlineto{\pgfqpoint{3.229923in}{1.198333in}}%
\pgfpathlineto{\pgfqpoint{3.206484in}{1.172807in}}%
\pgfpathclose%
\pgfusepath{fill}%
\end{pgfscope}%
\begin{pgfscope}%
\pgfpathrectangle{\pgfqpoint{0.100000in}{0.100000in}}{\pgfqpoint{3.608454in}{2.310000in}}%
\pgfusepath{clip}%
\pgfsetbuttcap%
\pgfsetmiterjoin%
\definecolor{currentfill}{rgb}{0.000000,0.568627,0.715686}%
\pgfsetfillcolor{currentfill}%
\pgfsetlinewidth{0.000000pt}%
\definecolor{currentstroke}{rgb}{0.000000,0.000000,0.000000}%
\pgfsetstrokecolor{currentstroke}%
\pgfsetstrokeopacity{0.000000}%
\pgfsetdash{}{0pt}%
\pgfpathmoveto{\pgfqpoint{1.051274in}{1.732636in}}%
\pgfpathlineto{\pgfqpoint{1.054206in}{1.729137in}}%
\pgfpathlineto{\pgfqpoint{1.064066in}{1.727041in}}%
\pgfpathlineto{\pgfqpoint{1.060120in}{1.706792in}}%
\pgfpathlineto{\pgfqpoint{1.054992in}{1.681239in}}%
\pgfpathlineto{\pgfqpoint{1.004416in}{1.690817in}}%
\pgfpathlineto{\pgfqpoint{0.994964in}{1.693281in}}%
\pgfpathlineto{\pgfqpoint{0.981161in}{1.696275in}}%
\pgfpathlineto{\pgfqpoint{0.988011in}{1.729060in}}%
\pgfpathlineto{\pgfqpoint{1.000366in}{1.726400in}}%
\pgfpathlineto{\pgfqpoint{1.002092in}{1.733915in}}%
\pgfpathlineto{\pgfqpoint{1.010156in}{1.733884in}}%
\pgfpathlineto{\pgfqpoint{1.013857in}{1.751700in}}%
\pgfpathlineto{\pgfqpoint{1.023421in}{1.749624in}}%
\pgfpathlineto{\pgfqpoint{1.027689in}{1.755650in}}%
\pgfpathlineto{\pgfqpoint{1.033321in}{1.782826in}}%
\pgfpathlineto{\pgfqpoint{1.050343in}{1.779280in}}%
\pgfpathlineto{\pgfqpoint{1.044792in}{1.752067in}}%
\pgfpathlineto{\pgfqpoint{1.041428in}{1.752757in}}%
\pgfpathlineto{\pgfqpoint{1.038549in}{1.738796in}}%
\pgfpathlineto{\pgfqpoint{1.044474in}{1.739072in}}%
\pgfpathlineto{\pgfqpoint{1.051274in}{1.732636in}}%
\pgfpathclose%
\pgfusepath{fill}%
\end{pgfscope}%
\begin{pgfscope}%
\pgfpathrectangle{\pgfqpoint{0.100000in}{0.100000in}}{\pgfqpoint{3.608454in}{2.310000in}}%
\pgfusepath{clip}%
\pgfsetbuttcap%
\pgfsetmiterjoin%
\definecolor{currentfill}{rgb}{0.000000,0.443137,0.778431}%
\pgfsetfillcolor{currentfill}%
\pgfsetlinewidth{0.000000pt}%
\definecolor{currentstroke}{rgb}{0.000000,0.000000,0.000000}%
\pgfsetstrokecolor{currentstroke}%
\pgfsetstrokeopacity{0.000000}%
\pgfsetdash{}{0pt}%
\pgfpathmoveto{\pgfqpoint{2.723904in}{1.022807in}}%
\pgfpathlineto{\pgfqpoint{2.715608in}{1.053498in}}%
\pgfpathlineto{\pgfqpoint{2.698781in}{1.052062in}}%
\pgfpathlineto{\pgfqpoint{2.696968in}{1.070359in}}%
\pgfpathlineto{\pgfqpoint{2.693282in}{1.075443in}}%
\pgfpathlineto{\pgfqpoint{2.694900in}{1.080379in}}%
\pgfpathlineto{\pgfqpoint{2.693772in}{1.094220in}}%
\pgfpathlineto{\pgfqpoint{2.697680in}{1.093767in}}%
\pgfpathlineto{\pgfqpoint{2.714180in}{1.097018in}}%
\pgfpathlineto{\pgfqpoint{2.722462in}{1.100629in}}%
\pgfpathlineto{\pgfqpoint{2.724406in}{1.094850in}}%
\pgfpathlineto{\pgfqpoint{2.736951in}{1.085202in}}%
\pgfpathlineto{\pgfqpoint{2.739787in}{1.093594in}}%
\pgfpathlineto{\pgfqpoint{2.749929in}{1.091001in}}%
\pgfpathlineto{\pgfqpoint{2.755774in}{1.081517in}}%
\pgfpathlineto{\pgfqpoint{2.751742in}{1.068739in}}%
\pgfpathlineto{\pgfqpoint{2.756131in}{1.058082in}}%
\pgfpathlineto{\pgfqpoint{2.752122in}{1.044231in}}%
\pgfpathlineto{\pgfqpoint{2.747587in}{1.043751in}}%
\pgfpathlineto{\pgfqpoint{2.746037in}{1.035657in}}%
\pgfpathlineto{\pgfqpoint{2.753682in}{1.036469in}}%
\pgfpathlineto{\pgfqpoint{2.754206in}{1.025936in}}%
\pgfpathlineto{\pgfqpoint{2.723904in}{1.022807in}}%
\pgfpathclose%
\pgfusepath{fill}%
\end{pgfscope}%
\begin{pgfscope}%
\pgfpathrectangle{\pgfqpoint{0.100000in}{0.100000in}}{\pgfqpoint{3.608454in}{2.310000in}}%
\pgfusepath{clip}%
\pgfsetbuttcap%
\pgfsetmiterjoin%
\definecolor{currentfill}{rgb}{0.000000,0.466667,0.766667}%
\pgfsetfillcolor{currentfill}%
\pgfsetlinewidth{0.000000pt}%
\definecolor{currentstroke}{rgb}{0.000000,0.000000,0.000000}%
\pgfsetstrokecolor{currentstroke}%
\pgfsetstrokeopacity{0.000000}%
\pgfsetdash{}{0pt}%
\pgfpathmoveto{\pgfqpoint{2.667350in}{1.499209in}}%
\pgfpathlineto{\pgfqpoint{2.653860in}{1.497742in}}%
\pgfpathlineto{\pgfqpoint{2.649395in}{1.497238in}}%
\pgfpathlineto{\pgfqpoint{2.646259in}{1.524584in}}%
\pgfpathlineto{\pgfqpoint{2.633101in}{1.523054in}}%
\pgfpathlineto{\pgfqpoint{2.633720in}{1.516230in}}%
\pgfpathlineto{\pgfqpoint{2.615905in}{1.514539in}}%
\pgfpathlineto{\pgfqpoint{2.614051in}{1.535158in}}%
\pgfpathlineto{\pgfqpoint{2.638514in}{1.537637in}}%
\pgfpathlineto{\pgfqpoint{2.636183in}{1.558287in}}%
\pgfpathlineto{\pgfqpoint{2.660092in}{1.560783in}}%
\pgfpathlineto{\pgfqpoint{2.661344in}{1.549614in}}%
\pgfpathlineto{\pgfqpoint{2.660236in}{1.540245in}}%
\pgfpathlineto{\pgfqpoint{2.661863in}{1.526375in}}%
\pgfpathlineto{\pgfqpoint{2.664197in}{1.526649in}}%
\pgfpathlineto{\pgfqpoint{2.667350in}{1.499209in}}%
\pgfpathclose%
\pgfusepath{fill}%
\end{pgfscope}%
\begin{pgfscope}%
\pgfpathrectangle{\pgfqpoint{0.100000in}{0.100000in}}{\pgfqpoint{3.608454in}{2.310000in}}%
\pgfusepath{clip}%
\pgfsetbuttcap%
\pgfsetmiterjoin%
\definecolor{currentfill}{rgb}{0.000000,0.690196,0.654902}%
\pgfsetfillcolor{currentfill}%
\pgfsetlinewidth{0.000000pt}%
\definecolor{currentstroke}{rgb}{0.000000,0.000000,0.000000}%
\pgfsetstrokecolor{currentstroke}%
\pgfsetstrokeopacity{0.000000}%
\pgfsetdash{}{0pt}%
\pgfpathmoveto{\pgfqpoint{0.429033in}{1.527264in}}%
\pgfpathlineto{\pgfqpoint{0.422980in}{1.525827in}}%
\pgfpathlineto{\pgfqpoint{0.419978in}{1.533247in}}%
\pgfpathlineto{\pgfqpoint{0.410415in}{1.537899in}}%
\pgfpathlineto{\pgfqpoint{0.403577in}{1.549156in}}%
\pgfpathlineto{\pgfqpoint{0.396978in}{1.549508in}}%
\pgfpathlineto{\pgfqpoint{0.393248in}{1.563598in}}%
\pgfpathlineto{\pgfqpoint{0.387572in}{1.570213in}}%
\pgfpathlineto{\pgfqpoint{0.382759in}{1.579907in}}%
\pgfpathlineto{\pgfqpoint{0.380186in}{1.589872in}}%
\pgfpathlineto{\pgfqpoint{0.372415in}{1.602800in}}%
\pgfpathlineto{\pgfqpoint{0.369555in}{1.609925in}}%
\pgfpathlineto{\pgfqpoint{0.374700in}{1.618945in}}%
\pgfpathlineto{\pgfqpoint{0.373785in}{1.631630in}}%
\pgfpathlineto{\pgfqpoint{0.374411in}{1.644351in}}%
\pgfpathlineto{\pgfqpoint{0.377409in}{1.651582in}}%
\pgfpathlineto{\pgfqpoint{0.382897in}{1.658786in}}%
\pgfpathlineto{\pgfqpoint{0.384563in}{1.669377in}}%
\pgfpathlineto{\pgfqpoint{0.384726in}{1.681179in}}%
\pgfpathlineto{\pgfqpoint{0.379033in}{1.697172in}}%
\pgfpathlineto{\pgfqpoint{0.406652in}{1.688065in}}%
\pgfpathlineto{\pgfqpoint{0.406038in}{1.686194in}}%
\pgfpathlineto{\pgfqpoint{0.441276in}{1.674970in}}%
\pgfpathlineto{\pgfqpoint{0.436761in}{1.661477in}}%
\pgfpathlineto{\pgfqpoint{0.437233in}{1.653925in}}%
\pgfpathlineto{\pgfqpoint{0.434583in}{1.644086in}}%
\pgfpathlineto{\pgfqpoint{0.443292in}{1.641381in}}%
\pgfpathlineto{\pgfqpoint{0.440384in}{1.631531in}}%
\pgfpathlineto{\pgfqpoint{0.435002in}{1.619433in}}%
\pgfpathlineto{\pgfqpoint{0.440654in}{1.612562in}}%
\pgfpathlineto{\pgfqpoint{0.447344in}{1.609339in}}%
\pgfpathlineto{\pgfqpoint{0.445039in}{1.597163in}}%
\pgfpathlineto{\pgfqpoint{0.448894in}{1.592474in}}%
\pgfpathlineto{\pgfqpoint{0.450911in}{1.584637in}}%
\pgfpathlineto{\pgfqpoint{0.446281in}{1.581110in}}%
\pgfpathlineto{\pgfqpoint{0.446440in}{1.577149in}}%
\pgfpathlineto{\pgfqpoint{0.438495in}{1.570367in}}%
\pgfpathlineto{\pgfqpoint{0.427977in}{1.570534in}}%
\pgfpathlineto{\pgfqpoint{0.425226in}{1.565689in}}%
\pgfpathlineto{\pgfqpoint{0.431374in}{1.550880in}}%
\pgfpathlineto{\pgfqpoint{0.433169in}{1.538861in}}%
\pgfpathlineto{\pgfqpoint{0.432587in}{1.529018in}}%
\pgfpathlineto{\pgfqpoint{0.429033in}{1.527264in}}%
\pgfpathclose%
\pgfusepath{fill}%
\end{pgfscope}%
\begin{pgfscope}%
\pgfpathrectangle{\pgfqpoint{0.100000in}{0.100000in}}{\pgfqpoint{3.608454in}{2.310000in}}%
\pgfusepath{clip}%
\pgfsetbuttcap%
\pgfsetmiterjoin%
\definecolor{currentfill}{rgb}{0.000000,0.686275,0.656863}%
\pgfsetfillcolor{currentfill}%
\pgfsetlinewidth{0.000000pt}%
\definecolor{currentstroke}{rgb}{0.000000,0.000000,0.000000}%
\pgfsetstrokecolor{currentstroke}%
\pgfsetstrokeopacity{0.000000}%
\pgfsetdash{}{0pt}%
\pgfpathmoveto{\pgfqpoint{2.929879in}{1.293625in}}%
\pgfpathlineto{\pgfqpoint{2.916184in}{1.292250in}}%
\pgfpathlineto{\pgfqpoint{2.909448in}{1.297631in}}%
\pgfpathlineto{\pgfqpoint{2.908038in}{1.306264in}}%
\pgfpathlineto{\pgfqpoint{2.902369in}{1.314765in}}%
\pgfpathlineto{\pgfqpoint{2.895571in}{1.311319in}}%
\pgfpathlineto{\pgfqpoint{2.900793in}{1.318863in}}%
\pgfpathlineto{\pgfqpoint{2.894629in}{1.327193in}}%
\pgfpathlineto{\pgfqpoint{2.894444in}{1.332079in}}%
\pgfpathlineto{\pgfqpoint{2.898128in}{1.341949in}}%
\pgfpathlineto{\pgfqpoint{2.906738in}{1.347296in}}%
\pgfpathlineto{\pgfqpoint{2.905122in}{1.355064in}}%
\pgfpathlineto{\pgfqpoint{2.903048in}{1.365422in}}%
\pgfpathlineto{\pgfqpoint{2.909141in}{1.370242in}}%
\pgfpathlineto{\pgfqpoint{2.920044in}{1.373705in}}%
\pgfpathlineto{\pgfqpoint{2.930069in}{1.365085in}}%
\pgfpathlineto{\pgfqpoint{2.936140in}{1.370645in}}%
\pgfpathlineto{\pgfqpoint{2.940346in}{1.366252in}}%
\pgfpathlineto{\pgfqpoint{2.957191in}{1.367025in}}%
\pgfpathlineto{\pgfqpoint{2.965422in}{1.379505in}}%
\pgfpathlineto{\pgfqpoint{2.972924in}{1.373592in}}%
\pgfpathlineto{\pgfqpoint{2.976752in}{1.368166in}}%
\pgfpathlineto{\pgfqpoint{2.975847in}{1.360884in}}%
\pgfpathlineto{\pgfqpoint{2.958041in}{1.346022in}}%
\pgfpathlineto{\pgfqpoint{2.953019in}{1.337104in}}%
\pgfpathlineto{\pgfqpoint{2.952368in}{1.321633in}}%
\pgfpathlineto{\pgfqpoint{2.944352in}{1.321971in}}%
\pgfpathlineto{\pgfqpoint{2.941096in}{1.316634in}}%
\pgfpathlineto{\pgfqpoint{2.946226in}{1.306488in}}%
\pgfpathlineto{\pgfqpoint{2.931210in}{1.301107in}}%
\pgfpathlineto{\pgfqpoint{2.935674in}{1.297744in}}%
\pgfpathlineto{\pgfqpoint{2.929879in}{1.293625in}}%
\pgfpathclose%
\pgfusepath{fill}%
\end{pgfscope}%
\begin{pgfscope}%
\pgfpathrectangle{\pgfqpoint{0.100000in}{0.100000in}}{\pgfqpoint{3.608454in}{2.310000in}}%
\pgfusepath{clip}%
\pgfsetbuttcap%
\pgfsetmiterjoin%
\definecolor{currentfill}{rgb}{0.000000,0.490196,0.754902}%
\pgfsetfillcolor{currentfill}%
\pgfsetlinewidth{0.000000pt}%
\definecolor{currentstroke}{rgb}{0.000000,0.000000,0.000000}%
\pgfsetstrokecolor{currentstroke}%
\pgfsetstrokeopacity{0.000000}%
\pgfsetdash{}{0pt}%
\pgfpathmoveto{\pgfqpoint{1.538610in}{1.726627in}}%
\pgfpathlineto{\pgfqpoint{1.532653in}{1.671867in}}%
\pgfpathlineto{\pgfqpoint{1.531372in}{1.656425in}}%
\pgfpathlineto{\pgfqpoint{1.508517in}{1.658724in}}%
\pgfpathlineto{\pgfqpoint{1.507108in}{1.644984in}}%
\pgfpathlineto{\pgfqpoint{1.504277in}{1.645283in}}%
\pgfpathlineto{\pgfqpoint{1.500514in}{1.634480in}}%
\pgfpathlineto{\pgfqpoint{1.486783in}{1.637214in}}%
\pgfpathlineto{\pgfqpoint{1.487410in}{1.642958in}}%
\pgfpathlineto{\pgfqpoint{1.491303in}{1.646739in}}%
\pgfpathlineto{\pgfqpoint{1.461309in}{1.650537in}}%
\pgfpathlineto{\pgfqpoint{1.429302in}{1.654336in}}%
\pgfpathlineto{\pgfqpoint{1.377332in}{1.661590in}}%
\pgfpathlineto{\pgfqpoint{1.376226in}{1.661744in}}%
\pgfpathlineto{\pgfqpoint{1.379922in}{1.688990in}}%
\pgfpathlineto{\pgfqpoint{1.382345in}{1.688646in}}%
\pgfpathlineto{\pgfqpoint{1.386059in}{1.715853in}}%
\pgfpathlineto{\pgfqpoint{1.388966in}{1.743023in}}%
\pgfpathlineto{\pgfqpoint{1.388262in}{1.745418in}}%
\pgfpathlineto{\pgfqpoint{1.412417in}{1.742036in}}%
\pgfpathlineto{\pgfqpoint{1.441739in}{1.737485in}}%
\pgfpathlineto{\pgfqpoint{1.474731in}{1.733639in}}%
\pgfpathlineto{\pgfqpoint{1.538610in}{1.726627in}}%
\pgfpathclose%
\pgfusepath{fill}%
\end{pgfscope}%
\begin{pgfscope}%
\pgfpathrectangle{\pgfqpoint{0.100000in}{0.100000in}}{\pgfqpoint{3.608454in}{2.310000in}}%
\pgfusepath{clip}%
\pgfsetbuttcap%
\pgfsetmiterjoin%
\definecolor{currentfill}{rgb}{0.000000,0.372549,0.813725}%
\pgfsetfillcolor{currentfill}%
\pgfsetlinewidth{0.000000pt}%
\definecolor{currentstroke}{rgb}{0.000000,0.000000,0.000000}%
\pgfsetstrokecolor{currentstroke}%
\pgfsetstrokeopacity{0.000000}%
\pgfsetdash{}{0pt}%
\pgfpathmoveto{\pgfqpoint{2.088125in}{1.418722in}}%
\pgfpathlineto{\pgfqpoint{2.086211in}{1.418729in}}%
\pgfpathlineto{\pgfqpoint{2.031759in}{1.419257in}}%
\pgfpathlineto{\pgfqpoint{2.018170in}{1.419481in}}%
\pgfpathlineto{\pgfqpoint{2.019155in}{1.481378in}}%
\pgfpathlineto{\pgfqpoint{2.056854in}{1.480886in}}%
\pgfpathlineto{\pgfqpoint{2.053550in}{1.477162in}}%
\pgfpathlineto{\pgfqpoint{2.059082in}{1.471267in}}%
\pgfpathlineto{\pgfqpoint{2.060821in}{1.465149in}}%
\pgfpathlineto{\pgfqpoint{2.069544in}{1.443619in}}%
\pgfpathlineto{\pgfqpoint{2.077207in}{1.438452in}}%
\pgfpathlineto{\pgfqpoint{2.077853in}{1.433443in}}%
\pgfpathlineto{\pgfqpoint{2.082351in}{1.428944in}}%
\pgfpathlineto{\pgfqpoint{2.081282in}{1.423410in}}%
\pgfpathlineto{\pgfqpoint{2.088125in}{1.418722in}}%
\pgfpathclose%
\pgfusepath{fill}%
\end{pgfscope}%
\begin{pgfscope}%
\pgfpathrectangle{\pgfqpoint{0.100000in}{0.100000in}}{\pgfqpoint{3.608454in}{2.310000in}}%
\pgfusepath{clip}%
\pgfsetbuttcap%
\pgfsetmiterjoin%
\definecolor{currentfill}{rgb}{0.000000,0.619608,0.690196}%
\pgfsetfillcolor{currentfill}%
\pgfsetlinewidth{0.000000pt}%
\definecolor{currentstroke}{rgb}{0.000000,0.000000,0.000000}%
\pgfsetstrokecolor{currentstroke}%
\pgfsetstrokeopacity{0.000000}%
\pgfsetdash{}{0pt}%
\pgfpathmoveto{\pgfqpoint{1.624223in}{0.845743in}}%
\pgfpathlineto{\pgfqpoint{1.566607in}{0.850506in}}%
\pgfpathlineto{\pgfqpoint{1.563559in}{0.816012in}}%
\pgfpathlineto{\pgfqpoint{1.562944in}{0.809149in}}%
\pgfpathlineto{\pgfqpoint{1.518670in}{0.813225in}}%
\pgfpathlineto{\pgfqpoint{1.522650in}{0.854609in}}%
\pgfpathlineto{\pgfqpoint{1.516573in}{0.855187in}}%
\pgfpathlineto{\pgfqpoint{1.519940in}{0.890229in}}%
\pgfpathlineto{\pgfqpoint{1.523143in}{0.889921in}}%
\pgfpathlineto{\pgfqpoint{1.526429in}{0.923948in}}%
\pgfpathlineto{\pgfqpoint{1.529406in}{0.923693in}}%
\pgfpathlineto{\pgfqpoint{1.530771in}{0.937426in}}%
\pgfpathlineto{\pgfqpoint{1.544636in}{0.936153in}}%
\pgfpathlineto{\pgfqpoint{1.574789in}{0.933425in}}%
\pgfpathlineto{\pgfqpoint{1.573241in}{0.919040in}}%
\pgfpathlineto{\pgfqpoint{1.638056in}{0.913724in}}%
\pgfpathlineto{\pgfqpoint{1.635380in}{0.879674in}}%
\pgfpathlineto{\pgfqpoint{1.626580in}{0.880308in}}%
\pgfpathlineto{\pgfqpoint{1.624223in}{0.845743in}}%
\pgfpathclose%
\pgfusepath{fill}%
\end{pgfscope}%
\begin{pgfscope}%
\pgfpathrectangle{\pgfqpoint{0.100000in}{0.100000in}}{\pgfqpoint{3.608454in}{2.310000in}}%
\pgfusepath{clip}%
\pgfsetbuttcap%
\pgfsetmiterjoin%
\definecolor{currentfill}{rgb}{0.000000,0.388235,0.805882}%
\pgfsetfillcolor{currentfill}%
\pgfsetlinewidth{0.000000pt}%
\definecolor{currentstroke}{rgb}{0.000000,0.000000,0.000000}%
\pgfsetstrokecolor{currentstroke}%
\pgfsetstrokeopacity{0.000000}%
\pgfsetdash{}{0pt}%
\pgfpathmoveto{\pgfqpoint{2.547881in}{1.571936in}}%
\pgfpathlineto{\pgfqpoint{2.550384in}{1.539587in}}%
\pgfpathlineto{\pgfqpoint{2.521564in}{1.536820in}}%
\pgfpathlineto{\pgfqpoint{2.522216in}{1.529995in}}%
\pgfpathlineto{\pgfqpoint{2.508442in}{1.528613in}}%
\pgfpathlineto{\pgfqpoint{2.508522in}{1.521688in}}%
\pgfpathlineto{\pgfqpoint{2.488674in}{1.519766in}}%
\pgfpathlineto{\pgfqpoint{2.484797in}{1.561007in}}%
\pgfpathlineto{\pgfqpoint{2.482413in}{1.595355in}}%
\pgfpathlineto{\pgfqpoint{2.482705in}{1.602290in}}%
\pgfpathlineto{\pgfqpoint{2.475849in}{1.601806in}}%
\pgfpathlineto{\pgfqpoint{2.473885in}{1.628654in}}%
\pgfpathlineto{\pgfqpoint{2.526558in}{1.632431in}}%
\pgfpathlineto{\pgfqpoint{2.525980in}{1.616699in}}%
\pgfpathlineto{\pgfqpoint{2.528499in}{1.610045in}}%
\pgfpathlineto{\pgfqpoint{2.536221in}{1.600158in}}%
\pgfpathlineto{\pgfqpoint{2.541972in}{1.582322in}}%
\pgfpathlineto{\pgfqpoint{2.547881in}{1.571936in}}%
\pgfpathclose%
\pgfusepath{fill}%
\end{pgfscope}%
\begin{pgfscope}%
\pgfpathrectangle{\pgfqpoint{0.100000in}{0.100000in}}{\pgfqpoint{3.608454in}{2.310000in}}%
\pgfusepath{clip}%
\pgfsetbuttcap%
\pgfsetmiterjoin%
\definecolor{currentfill}{rgb}{0.000000,0.490196,0.754902}%
\pgfsetfillcolor{currentfill}%
\pgfsetlinewidth{0.000000pt}%
\definecolor{currentstroke}{rgb}{0.000000,0.000000,0.000000}%
\pgfsetstrokecolor{currentstroke}%
\pgfsetstrokeopacity{0.000000}%
\pgfsetdash{}{0pt}%
\pgfpathmoveto{\pgfqpoint{1.962491in}{1.386340in}}%
\pgfpathlineto{\pgfqpoint{1.961886in}{1.365738in}}%
\pgfpathlineto{\pgfqpoint{1.927666in}{1.366756in}}%
\pgfpathlineto{\pgfqpoint{1.928379in}{1.387333in}}%
\pgfpathlineto{\pgfqpoint{1.893270in}{1.388628in}}%
\pgfpathlineto{\pgfqpoint{1.894644in}{1.422961in}}%
\pgfpathlineto{\pgfqpoint{1.908588in}{1.422459in}}%
\pgfpathlineto{\pgfqpoint{1.963356in}{1.420723in}}%
\pgfpathlineto{\pgfqpoint{1.962491in}{1.386340in}}%
\pgfpathclose%
\pgfusepath{fill}%
\end{pgfscope}%
\begin{pgfscope}%
\pgfpathrectangle{\pgfqpoint{0.100000in}{0.100000in}}{\pgfqpoint{3.608454in}{2.310000in}}%
\pgfusepath{clip}%
\pgfsetbuttcap%
\pgfsetmiterjoin%
\definecolor{currentfill}{rgb}{0.000000,0.698039,0.650980}%
\pgfsetfillcolor{currentfill}%
\pgfsetlinewidth{0.000000pt}%
\definecolor{currentstroke}{rgb}{0.000000,0.000000,0.000000}%
\pgfsetstrokecolor{currentstroke}%
\pgfsetstrokeopacity{0.000000}%
\pgfsetdash{}{0pt}%
\pgfpathmoveto{\pgfqpoint{1.691656in}{1.593744in}}%
\pgfpathlineto{\pgfqpoint{1.655822in}{1.596346in}}%
\pgfpathlineto{\pgfqpoint{1.652617in}{1.603533in}}%
\pgfpathlineto{\pgfqpoint{1.654543in}{1.630993in}}%
\pgfpathlineto{\pgfqpoint{1.653148in}{1.631106in}}%
\pgfpathlineto{\pgfqpoint{1.655483in}{1.658354in}}%
\pgfpathlineto{\pgfqpoint{1.655854in}{1.675393in}}%
\pgfpathlineto{\pgfqpoint{1.643808in}{1.676503in}}%
\pgfpathlineto{\pgfqpoint{1.648016in}{1.724287in}}%
\pgfpathlineto{\pgfqpoint{1.654175in}{1.729234in}}%
\pgfpathlineto{\pgfqpoint{1.659487in}{1.729827in}}%
\pgfpathlineto{\pgfqpoint{1.695616in}{1.726789in}}%
\pgfpathlineto{\pgfqpoint{1.697918in}{1.727667in}}%
\pgfpathlineto{\pgfqpoint{1.696530in}{1.709950in}}%
\pgfpathlineto{\pgfqpoint{1.697612in}{1.702965in}}%
\pgfpathlineto{\pgfqpoint{1.748309in}{1.699187in}}%
\pgfpathlineto{\pgfqpoint{1.746207in}{1.668241in}}%
\pgfpathlineto{\pgfqpoint{1.696825in}{1.671982in}}%
\pgfpathlineto{\pgfqpoint{1.696414in}{1.655020in}}%
\pgfpathlineto{\pgfqpoint{1.694322in}{1.628015in}}%
\pgfpathlineto{\pgfqpoint{1.695984in}{1.627874in}}%
\pgfpathlineto{\pgfqpoint{1.693738in}{1.600532in}}%
\pgfpathlineto{\pgfqpoint{1.691656in}{1.593744in}}%
\pgfpathclose%
\pgfusepath{fill}%
\end{pgfscope}%
\begin{pgfscope}%
\pgfpathrectangle{\pgfqpoint{0.100000in}{0.100000in}}{\pgfqpoint{3.608454in}{2.310000in}}%
\pgfusepath{clip}%
\pgfsetbuttcap%
\pgfsetmiterjoin%
\definecolor{currentfill}{rgb}{0.000000,0.721569,0.639216}%
\pgfsetfillcolor{currentfill}%
\pgfsetlinewidth{0.000000pt}%
\definecolor{currentstroke}{rgb}{0.000000,0.000000,0.000000}%
\pgfsetstrokecolor{currentstroke}%
\pgfsetstrokeopacity{0.000000}%
\pgfsetdash{}{0pt}%
\pgfpathmoveto{\pgfqpoint{2.136953in}{2.110228in}}%
\pgfpathlineto{\pgfqpoint{2.145780in}{2.110086in}}%
\pgfpathlineto{\pgfqpoint{2.146826in}{2.106493in}}%
\pgfpathlineto{\pgfqpoint{2.168582in}{2.104438in}}%
\pgfpathlineto{\pgfqpoint{2.170811in}{2.095732in}}%
\pgfpathlineto{\pgfqpoint{2.188133in}{2.098370in}}%
\pgfpathlineto{\pgfqpoint{2.188169in}{2.101766in}}%
\pgfpathlineto{\pgfqpoint{2.201785in}{2.106387in}}%
\pgfpathlineto{\pgfqpoint{2.208113in}{2.105322in}}%
\pgfpathlineto{\pgfqpoint{2.208627in}{2.060466in}}%
\pgfpathlineto{\pgfqpoint{2.209745in}{2.046401in}}%
\pgfpathlineto{\pgfqpoint{2.172471in}{2.046322in}}%
\pgfpathlineto{\pgfqpoint{2.172503in}{2.042117in}}%
\pgfpathlineto{\pgfqpoint{2.138015in}{2.041726in}}%
\pgfpathlineto{\pgfqpoint{2.137206in}{2.083483in}}%
\pgfpathlineto{\pgfqpoint{2.136953in}{2.110228in}}%
\pgfpathclose%
\pgfusepath{fill}%
\end{pgfscope}%
\begin{pgfscope}%
\pgfpathrectangle{\pgfqpoint{0.100000in}{0.100000in}}{\pgfqpoint{3.608454in}{2.310000in}}%
\pgfusepath{clip}%
\pgfsetbuttcap%
\pgfsetmiterjoin%
\definecolor{currentfill}{rgb}{0.000000,0.545098,0.727451}%
\pgfsetfillcolor{currentfill}%
\pgfsetlinewidth{0.000000pt}%
\definecolor{currentstroke}{rgb}{0.000000,0.000000,0.000000}%
\pgfsetstrokecolor{currentstroke}%
\pgfsetstrokeopacity{0.000000}%
\pgfsetdash{}{0pt}%
\pgfpathmoveto{\pgfqpoint{2.554067in}{1.495355in}}%
\pgfpathlineto{\pgfqpoint{2.555646in}{1.476018in}}%
\pgfpathlineto{\pgfqpoint{2.531157in}{1.473630in}}%
\pgfpathlineto{\pgfqpoint{2.520205in}{1.472936in}}%
\pgfpathlineto{\pgfqpoint{2.516336in}{1.513050in}}%
\pgfpathlineto{\pgfqpoint{2.509486in}{1.512307in}}%
\pgfpathlineto{\pgfqpoint{2.508522in}{1.521688in}}%
\pgfpathlineto{\pgfqpoint{2.508442in}{1.528613in}}%
\pgfpathlineto{\pgfqpoint{2.522216in}{1.529995in}}%
\pgfpathlineto{\pgfqpoint{2.521564in}{1.536820in}}%
\pgfpathlineto{\pgfqpoint{2.550384in}{1.539587in}}%
\pgfpathlineto{\pgfqpoint{2.554067in}{1.495355in}}%
\pgfpathclose%
\pgfusepath{fill}%
\end{pgfscope}%
\begin{pgfscope}%
\pgfpathrectangle{\pgfqpoint{0.100000in}{0.100000in}}{\pgfqpoint{3.608454in}{2.310000in}}%
\pgfusepath{clip}%
\pgfsetbuttcap%
\pgfsetmiterjoin%
\definecolor{currentfill}{rgb}{0.000000,0.537255,0.731373}%
\pgfsetfillcolor{currentfill}%
\pgfsetlinewidth{0.000000pt}%
\definecolor{currentstroke}{rgb}{0.000000,0.000000,0.000000}%
\pgfsetstrokecolor{currentstroke}%
\pgfsetstrokeopacity{0.000000}%
\pgfsetdash{}{0pt}%
\pgfpathmoveto{\pgfqpoint{2.452854in}{0.589739in}}%
\pgfpathlineto{\pgfqpoint{2.445754in}{0.592557in}}%
\pgfpathlineto{\pgfqpoint{2.442006in}{0.599263in}}%
\pgfpathlineto{\pgfqpoint{2.443757in}{0.607123in}}%
\pgfpathlineto{\pgfqpoint{2.440322in}{0.609463in}}%
\pgfpathlineto{\pgfqpoint{2.430825in}{0.609081in}}%
\pgfpathlineto{\pgfqpoint{2.428180in}{0.614445in}}%
\pgfpathlineto{\pgfqpoint{2.420890in}{0.618181in}}%
\pgfpathlineto{\pgfqpoint{2.416070in}{0.624294in}}%
\pgfpathlineto{\pgfqpoint{2.392099in}{0.624223in}}%
\pgfpathlineto{\pgfqpoint{2.386695in}{0.629800in}}%
\pgfpathlineto{\pgfqpoint{2.386115in}{0.636858in}}%
\pgfpathlineto{\pgfqpoint{2.389280in}{0.640751in}}%
\pgfpathlineto{\pgfqpoint{2.407987in}{0.645878in}}%
\pgfpathlineto{\pgfqpoint{2.415540in}{0.649843in}}%
\pgfpathlineto{\pgfqpoint{2.419484in}{0.654679in}}%
\pgfpathlineto{\pgfqpoint{2.431222in}{0.657222in}}%
\pgfpathlineto{\pgfqpoint{2.435277in}{0.651886in}}%
\pgfpathlineto{\pgfqpoint{2.432451in}{0.691509in}}%
\pgfpathlineto{\pgfqpoint{2.424922in}{0.713419in}}%
\pgfpathlineto{\pgfqpoint{2.467199in}{0.715885in}}%
\pgfpathlineto{\pgfqpoint{2.465118in}{0.704833in}}%
\pgfpathlineto{\pgfqpoint{2.461305in}{0.698015in}}%
\pgfpathlineto{\pgfqpoint{2.460731in}{0.688378in}}%
\pgfpathlineto{\pgfqpoint{2.466743in}{0.676753in}}%
\pgfpathlineto{\pgfqpoint{2.472822in}{0.672176in}}%
\pgfpathlineto{\pgfqpoint{2.478549in}{0.654208in}}%
\pgfpathlineto{\pgfqpoint{2.485040in}{0.651337in}}%
\pgfpathlineto{\pgfqpoint{2.478638in}{0.648879in}}%
\pgfpathlineto{\pgfqpoint{2.472203in}{0.638511in}}%
\pgfpathlineto{\pgfqpoint{2.465852in}{0.639118in}}%
\pgfpathlineto{\pgfqpoint{2.464257in}{0.631545in}}%
\pgfpathlineto{\pgfqpoint{2.470761in}{0.631547in}}%
\pgfpathlineto{\pgfqpoint{2.473516in}{0.626309in}}%
\pgfpathlineto{\pgfqpoint{2.481391in}{0.627004in}}%
\pgfpathlineto{\pgfqpoint{2.482013in}{0.636124in}}%
\pgfpathlineto{\pgfqpoint{2.487826in}{0.640181in}}%
\pgfpathlineto{\pgfqpoint{2.500761in}{0.631538in}}%
\pgfpathlineto{\pgfqpoint{2.493847in}{0.616542in}}%
\pgfpathlineto{\pgfqpoint{2.478334in}{0.616245in}}%
\pgfpathlineto{\pgfqpoint{2.480862in}{0.608985in}}%
\pgfpathlineto{\pgfqpoint{2.477552in}{0.599501in}}%
\pgfpathlineto{\pgfqpoint{2.485322in}{0.596196in}}%
\pgfpathlineto{\pgfqpoint{2.492517in}{0.589207in}}%
\pgfpathlineto{\pgfqpoint{2.511782in}{0.586087in}}%
\pgfpathlineto{\pgfqpoint{2.520554in}{0.574448in}}%
\pgfpathlineto{\pgfqpoint{2.526501in}{0.573415in}}%
\pgfpathlineto{\pgfqpoint{2.522992in}{0.564279in}}%
\pgfpathlineto{\pgfqpoint{2.516518in}{0.560200in}}%
\pgfpathlineto{\pgfqpoint{2.511813in}{0.562316in}}%
\pgfpathlineto{\pgfqpoint{2.503209in}{0.561121in}}%
\pgfpathlineto{\pgfqpoint{2.491659in}{0.574580in}}%
\pgfpathlineto{\pgfqpoint{2.475551in}{0.580137in}}%
\pgfpathlineto{\pgfqpoint{2.466715in}{0.580608in}}%
\pgfpathlineto{\pgfqpoint{2.468264in}{0.585589in}}%
\pgfpathlineto{\pgfqpoint{2.466744in}{0.592994in}}%
\pgfpathlineto{\pgfqpoint{2.452854in}{0.589739in}}%
\pgfpathclose%
\pgfusepath{fill}%
\end{pgfscope}%
\begin{pgfscope}%
\pgfpathrectangle{\pgfqpoint{0.100000in}{0.100000in}}{\pgfqpoint{3.608454in}{2.310000in}}%
\pgfusepath{clip}%
\pgfsetbuttcap%
\pgfsetmiterjoin%
\definecolor{currentfill}{rgb}{0.000000,0.607843,0.696078}%
\pgfsetfillcolor{currentfill}%
\pgfsetlinewidth{0.000000pt}%
\definecolor{currentstroke}{rgb}{0.000000,0.000000,0.000000}%
\pgfsetstrokecolor{currentstroke}%
\pgfsetstrokeopacity{0.000000}%
\pgfsetdash{}{0pt}%
\pgfpathmoveto{\pgfqpoint{2.306014in}{1.604048in}}%
\pgfpathlineto{\pgfqpoint{2.264937in}{1.602761in}}%
\pgfpathlineto{\pgfqpoint{2.237469in}{1.601952in}}%
\pgfpathlineto{\pgfqpoint{2.237674in}{1.595106in}}%
\pgfpathlineto{\pgfqpoint{2.223956in}{1.594717in}}%
\pgfpathlineto{\pgfqpoint{2.223427in}{1.615413in}}%
\pgfpathlineto{\pgfqpoint{2.222007in}{1.615392in}}%
\pgfpathlineto{\pgfqpoint{2.221382in}{1.649990in}}%
\pgfpathlineto{\pgfqpoint{2.276089in}{1.651404in}}%
\pgfpathlineto{\pgfqpoint{2.275677in}{1.665325in}}%
\pgfpathlineto{\pgfqpoint{2.303199in}{1.666191in}}%
\pgfpathlineto{\pgfqpoint{2.304826in}{1.617701in}}%
\pgfpathlineto{\pgfqpoint{2.306014in}{1.604048in}}%
\pgfpathclose%
\pgfusepath{fill}%
\end{pgfscope}%
\begin{pgfscope}%
\pgfpathrectangle{\pgfqpoint{0.100000in}{0.100000in}}{\pgfqpoint{3.608454in}{2.310000in}}%
\pgfusepath{clip}%
\pgfsetbuttcap%
\pgfsetmiterjoin%
\definecolor{currentfill}{rgb}{0.000000,0.333333,0.833333}%
\pgfsetfillcolor{currentfill}%
\pgfsetlinewidth{0.000000pt}%
\definecolor{currentstroke}{rgb}{0.000000,0.000000,0.000000}%
\pgfsetstrokecolor{currentstroke}%
\pgfsetstrokeopacity{0.000000}%
\pgfsetdash{}{0pt}%
\pgfpathmoveto{\pgfqpoint{2.866524in}{0.926111in}}%
\pgfpathlineto{\pgfqpoint{2.849065in}{0.920723in}}%
\pgfpathlineto{\pgfqpoint{2.848183in}{0.924494in}}%
\pgfpathlineto{\pgfqpoint{2.848096in}{0.929927in}}%
\pgfpathlineto{\pgfqpoint{2.843454in}{0.939687in}}%
\pgfpathlineto{\pgfqpoint{2.838874in}{0.945035in}}%
\pgfpathlineto{\pgfqpoint{2.830443in}{0.932055in}}%
\pgfpathlineto{\pgfqpoint{2.818507in}{0.934685in}}%
\pgfpathlineto{\pgfqpoint{2.809234in}{0.933588in}}%
\pgfpathlineto{\pgfqpoint{2.807171in}{0.925663in}}%
\pgfpathlineto{\pgfqpoint{2.802892in}{0.925244in}}%
\pgfpathlineto{\pgfqpoint{2.796122in}{0.932528in}}%
\pgfpathlineto{\pgfqpoint{2.791999in}{0.944257in}}%
\pgfpathlineto{\pgfqpoint{2.778830in}{0.942946in}}%
\pgfpathlineto{\pgfqpoint{2.773147in}{0.947109in}}%
\pgfpathlineto{\pgfqpoint{2.771603in}{0.963496in}}%
\pgfpathlineto{\pgfqpoint{2.762398in}{0.964970in}}%
\pgfpathlineto{\pgfqpoint{2.760784in}{0.972182in}}%
\pgfpathlineto{\pgfqpoint{2.765064in}{0.976399in}}%
\pgfpathlineto{\pgfqpoint{2.767572in}{0.987144in}}%
\pgfpathlineto{\pgfqpoint{2.784817in}{0.988730in}}%
\pgfpathlineto{\pgfqpoint{2.780892in}{1.025876in}}%
\pgfpathlineto{\pgfqpoint{2.799254in}{1.027968in}}%
\pgfpathlineto{\pgfqpoint{2.801686in}{1.031063in}}%
\pgfpathlineto{\pgfqpoint{2.810594in}{1.033632in}}%
\pgfpathlineto{\pgfqpoint{2.810860in}{1.028741in}}%
\pgfpathlineto{\pgfqpoint{2.818142in}{1.023149in}}%
\pgfpathlineto{\pgfqpoint{2.825899in}{1.020792in}}%
\pgfpathlineto{\pgfqpoint{2.830646in}{1.010243in}}%
\pgfpathlineto{\pgfqpoint{2.827384in}{1.003120in}}%
\pgfpathlineto{\pgfqpoint{2.836248in}{0.996454in}}%
\pgfpathlineto{\pgfqpoint{2.839244in}{0.999162in}}%
\pgfpathlineto{\pgfqpoint{2.837086in}{0.989053in}}%
\pgfpathlineto{\pgfqpoint{2.844938in}{0.981568in}}%
\pgfpathlineto{\pgfqpoint{2.850509in}{0.983786in}}%
\pgfpathlineto{\pgfqpoint{2.864261in}{0.976571in}}%
\pgfpathlineto{\pgfqpoint{2.870396in}{0.968762in}}%
\pgfpathlineto{\pgfqpoint{2.880855in}{0.952903in}}%
\pgfpathlineto{\pgfqpoint{2.864668in}{0.946906in}}%
\pgfpathlineto{\pgfqpoint{2.866524in}{0.926111in}}%
\pgfpathclose%
\pgfusepath{fill}%
\end{pgfscope}%
\begin{pgfscope}%
\pgfpathrectangle{\pgfqpoint{0.100000in}{0.100000in}}{\pgfqpoint{3.608454in}{2.310000in}}%
\pgfusepath{clip}%
\pgfsetbuttcap%
\pgfsetmiterjoin%
\definecolor{currentfill}{rgb}{0.000000,0.803922,0.598039}%
\pgfsetfillcolor{currentfill}%
\pgfsetlinewidth{0.000000pt}%
\definecolor{currentstroke}{rgb}{0.000000,0.000000,0.000000}%
\pgfsetstrokecolor{currentstroke}%
\pgfsetstrokeopacity{0.000000}%
\pgfsetdash{}{0pt}%
\pgfpathmoveto{\pgfqpoint{2.882548in}{1.076650in}}%
\pgfpathlineto{\pgfqpoint{2.876403in}{1.073700in}}%
\pgfpathlineto{\pgfqpoint{2.872545in}{1.080983in}}%
\pgfpathlineto{\pgfqpoint{2.867388in}{1.085037in}}%
\pgfpathlineto{\pgfqpoint{2.858178in}{1.097707in}}%
\pgfpathlineto{\pgfqpoint{2.851025in}{1.094550in}}%
\pgfpathlineto{\pgfqpoint{2.837326in}{1.091183in}}%
\pgfpathlineto{\pgfqpoint{2.835727in}{1.088275in}}%
\pgfpathlineto{\pgfqpoint{2.827437in}{1.087830in}}%
\pgfpathlineto{\pgfqpoint{2.819581in}{1.083755in}}%
\pgfpathlineto{\pgfqpoint{2.814347in}{1.089123in}}%
\pgfpathlineto{\pgfqpoint{2.813980in}{1.097663in}}%
\pgfpathlineto{\pgfqpoint{2.817068in}{1.103142in}}%
\pgfpathlineto{\pgfqpoint{2.828227in}{1.112365in}}%
\pgfpathlineto{\pgfqpoint{2.835111in}{1.113784in}}%
\pgfpathlineto{\pgfqpoint{2.845692in}{1.114629in}}%
\pgfpathlineto{\pgfqpoint{2.859932in}{1.127064in}}%
\pgfpathlineto{\pgfqpoint{2.864797in}{1.125728in}}%
\pgfpathlineto{\pgfqpoint{2.866841in}{1.114054in}}%
\pgfpathlineto{\pgfqpoint{2.884882in}{1.098977in}}%
\pgfpathlineto{\pgfqpoint{2.885380in}{1.093119in}}%
\pgfpathlineto{\pgfqpoint{2.879644in}{1.078954in}}%
\pgfpathlineto{\pgfqpoint{2.882548in}{1.076650in}}%
\pgfpathclose%
\pgfusepath{fill}%
\end{pgfscope}%
\begin{pgfscope}%
\pgfpathrectangle{\pgfqpoint{0.100000in}{0.100000in}}{\pgfqpoint{3.608454in}{2.310000in}}%
\pgfusepath{clip}%
\pgfsetbuttcap%
\pgfsetmiterjoin%
\definecolor{currentfill}{rgb}{0.000000,0.419608,0.790196}%
\pgfsetfillcolor{currentfill}%
\pgfsetlinewidth{0.000000pt}%
\definecolor{currentstroke}{rgb}{0.000000,0.000000,0.000000}%
\pgfsetstrokecolor{currentstroke}%
\pgfsetstrokeopacity{0.000000}%
\pgfsetdash{}{0pt}%
\pgfpathmoveto{\pgfqpoint{2.433615in}{1.653903in}}%
\pgfpathlineto{\pgfqpoint{2.435480in}{1.626662in}}%
\pgfpathlineto{\pgfqpoint{2.402794in}{1.625171in}}%
\pgfpathlineto{\pgfqpoint{2.361136in}{1.623207in}}%
\pgfpathlineto{\pgfqpoint{2.356020in}{1.633465in}}%
\pgfpathlineto{\pgfqpoint{2.342986in}{1.636283in}}%
\pgfpathlineto{\pgfqpoint{2.336274in}{1.640299in}}%
\pgfpathlineto{\pgfqpoint{2.333060in}{1.651644in}}%
\pgfpathlineto{\pgfqpoint{2.330505in}{1.653247in}}%
\pgfpathlineto{\pgfqpoint{2.327890in}{1.671376in}}%
\pgfpathlineto{\pgfqpoint{2.334254in}{1.681222in}}%
\pgfpathlineto{\pgfqpoint{2.325644in}{1.688292in}}%
\pgfpathlineto{\pgfqpoint{2.325187in}{1.694132in}}%
\pgfpathlineto{\pgfqpoint{2.356072in}{1.695517in}}%
\pgfpathlineto{\pgfqpoint{2.357123in}{1.675633in}}%
\pgfpathlineto{\pgfqpoint{2.363976in}{1.679112in}}%
\pgfpathlineto{\pgfqpoint{2.370685in}{1.678611in}}%
\pgfpathlineto{\pgfqpoint{2.372395in}{1.647911in}}%
\pgfpathlineto{\pgfqpoint{2.406543in}{1.649815in}}%
\pgfpathlineto{\pgfqpoint{2.406360in}{1.653264in}}%
\pgfpathlineto{\pgfqpoint{2.433615in}{1.653903in}}%
\pgfpathclose%
\pgfusepath{fill}%
\end{pgfscope}%
\begin{pgfscope}%
\pgfpathrectangle{\pgfqpoint{0.100000in}{0.100000in}}{\pgfqpoint{3.608454in}{2.310000in}}%
\pgfusepath{clip}%
\pgfsetbuttcap%
\pgfsetmiterjoin%
\definecolor{currentfill}{rgb}{0.000000,0.807843,0.596078}%
\pgfsetfillcolor{currentfill}%
\pgfsetlinewidth{0.000000pt}%
\definecolor{currentstroke}{rgb}{0.000000,0.000000,0.000000}%
\pgfsetstrokecolor{currentstroke}%
\pgfsetstrokeopacity{0.000000}%
\pgfsetdash{}{0pt}%
\pgfpathmoveto{\pgfqpoint{2.288166in}{1.125922in}}%
\pgfpathlineto{\pgfqpoint{2.288004in}{1.144708in}}%
\pgfpathlineto{\pgfqpoint{2.289909in}{1.144788in}}%
\pgfpathlineto{\pgfqpoint{2.289673in}{1.168159in}}%
\pgfpathlineto{\pgfqpoint{2.290504in}{1.188921in}}%
\pgfpathlineto{\pgfqpoint{2.280237in}{1.188871in}}%
\pgfpathlineto{\pgfqpoint{2.279053in}{1.231943in}}%
\pgfpathlineto{\pgfqpoint{2.310037in}{1.232515in}}%
\pgfpathlineto{\pgfqpoint{2.310515in}{1.218707in}}%
\pgfpathlineto{\pgfqpoint{2.346916in}{1.219387in}}%
\pgfpathlineto{\pgfqpoint{2.347208in}{1.211310in}}%
\pgfpathlineto{\pgfqpoint{2.350740in}{1.202506in}}%
\pgfpathlineto{\pgfqpoint{2.357801in}{1.198383in}}%
\pgfpathlineto{\pgfqpoint{2.357912in}{1.194600in}}%
\pgfpathlineto{\pgfqpoint{2.345378in}{1.193313in}}%
\pgfpathlineto{\pgfqpoint{2.345750in}{1.172705in}}%
\pgfpathlineto{\pgfqpoint{2.352645in}{1.172837in}}%
\pgfpathlineto{\pgfqpoint{2.353073in}{1.147045in}}%
\pgfpathlineto{\pgfqpoint{2.335244in}{1.146265in}}%
\pgfpathlineto{\pgfqpoint{2.332518in}{1.146200in}}%
\pgfpathlineto{\pgfqpoint{2.332736in}{1.133366in}}%
\pgfpathlineto{\pgfqpoint{2.325633in}{1.133246in}}%
\pgfpathlineto{\pgfqpoint{2.325882in}{1.126330in}}%
\pgfpathlineto{\pgfqpoint{2.288166in}{1.125922in}}%
\pgfpathclose%
\pgfusepath{fill}%
\end{pgfscope}%
\begin{pgfscope}%
\pgfpathrectangle{\pgfqpoint{0.100000in}{0.100000in}}{\pgfqpoint{3.608454in}{2.310000in}}%
\pgfusepath{clip}%
\pgfsetbuttcap%
\pgfsetmiterjoin%
\definecolor{currentfill}{rgb}{0.000000,0.509804,0.745098}%
\pgfsetfillcolor{currentfill}%
\pgfsetlinewidth{0.000000pt}%
\definecolor{currentstroke}{rgb}{0.000000,0.000000,0.000000}%
\pgfsetstrokecolor{currentstroke}%
\pgfsetstrokeopacity{0.000000}%
\pgfsetdash{}{0pt}%
\pgfpathmoveto{\pgfqpoint{1.202252in}{1.853038in}}%
\pgfpathlineto{\pgfqpoint{1.204580in}{1.868127in}}%
\pgfpathlineto{\pgfqpoint{1.226134in}{1.864328in}}%
\pgfpathlineto{\pgfqpoint{1.224995in}{1.857794in}}%
\pgfpathlineto{\pgfqpoint{1.241305in}{1.854999in}}%
\pgfpathlineto{\pgfqpoint{1.245880in}{1.848777in}}%
\pgfpathlineto{\pgfqpoint{1.243670in}{1.842268in}}%
\pgfpathlineto{\pgfqpoint{1.249318in}{1.830814in}}%
\pgfpathlineto{\pgfqpoint{1.250644in}{1.821729in}}%
\pgfpathlineto{\pgfqpoint{1.253366in}{1.816700in}}%
\pgfpathlineto{\pgfqpoint{1.251746in}{1.806946in}}%
\pgfpathlineto{\pgfqpoint{1.244733in}{1.764439in}}%
\pgfpathlineto{\pgfqpoint{1.243269in}{1.757690in}}%
\pgfpathlineto{\pgfqpoint{1.226972in}{1.760454in}}%
\pgfpathlineto{\pgfqpoint{1.225832in}{1.753673in}}%
\pgfpathlineto{\pgfqpoint{1.212511in}{1.756000in}}%
\pgfpathlineto{\pgfqpoint{1.211732in}{1.751597in}}%
\pgfpathlineto{\pgfqpoint{1.198346in}{1.753929in}}%
\pgfpathlineto{\pgfqpoint{1.199226in}{1.758954in}}%
\pgfpathlineto{\pgfqpoint{1.186356in}{1.762535in}}%
\pgfpathlineto{\pgfqpoint{1.188888in}{1.777013in}}%
\pgfpathlineto{\pgfqpoint{1.180737in}{1.781578in}}%
\pgfpathlineto{\pgfqpoint{1.179359in}{1.788436in}}%
\pgfpathlineto{\pgfqpoint{1.170551in}{1.790061in}}%
\pgfpathlineto{\pgfqpoint{1.174897in}{1.813546in}}%
\pgfpathlineto{\pgfqpoint{1.186131in}{1.812373in}}%
\pgfpathlineto{\pgfqpoint{1.195518in}{1.814680in}}%
\pgfpathlineto{\pgfqpoint{1.202252in}{1.853038in}}%
\pgfpathclose%
\pgfusepath{fill}%
\end{pgfscope}%
\begin{pgfscope}%
\pgfpathrectangle{\pgfqpoint{0.100000in}{0.100000in}}{\pgfqpoint{3.608454in}{2.310000in}}%
\pgfusepath{clip}%
\pgfsetbuttcap%
\pgfsetmiterjoin%
\definecolor{currentfill}{rgb}{0.000000,0.572549,0.713725}%
\pgfsetfillcolor{currentfill}%
\pgfsetlinewidth{0.000000pt}%
\definecolor{currentstroke}{rgb}{0.000000,0.000000,0.000000}%
\pgfsetstrokecolor{currentstroke}%
\pgfsetstrokeopacity{0.000000}%
\pgfsetdash{}{0pt}%
\pgfpathmoveto{\pgfqpoint{1.873743in}{0.830250in}}%
\pgfpathlineto{\pgfqpoint{1.873122in}{0.813238in}}%
\pgfpathlineto{\pgfqpoint{1.842140in}{0.796889in}}%
\pgfpathlineto{\pgfqpoint{1.788954in}{0.799743in}}%
\pgfpathlineto{\pgfqpoint{1.759459in}{0.801389in}}%
\pgfpathlineto{\pgfqpoint{1.763926in}{0.871040in}}%
\pgfpathlineto{\pgfqpoint{1.808729in}{0.868432in}}%
\pgfpathlineto{\pgfqpoint{1.868391in}{0.865576in}}%
\pgfpathlineto{\pgfqpoint{1.866975in}{0.830682in}}%
\pgfpathlineto{\pgfqpoint{1.873743in}{0.830250in}}%
\pgfpathclose%
\pgfusepath{fill}%
\end{pgfscope}%
\begin{pgfscope}%
\pgfpathrectangle{\pgfqpoint{0.100000in}{0.100000in}}{\pgfqpoint{3.608454in}{2.310000in}}%
\pgfusepath{clip}%
\pgfsetbuttcap%
\pgfsetmiterjoin%
\definecolor{currentfill}{rgb}{0.000000,0.666667,0.666667}%
\pgfsetfillcolor{currentfill}%
\pgfsetlinewidth{0.000000pt}%
\definecolor{currentstroke}{rgb}{0.000000,0.000000,0.000000}%
\pgfsetstrokecolor{currentstroke}%
\pgfsetstrokeopacity{0.000000}%
\pgfsetdash{}{0pt}%
\pgfpathmoveto{\pgfqpoint{1.103689in}{1.391163in}}%
\pgfpathlineto{\pgfqpoint{1.103543in}{1.390370in}}%
\pgfpathlineto{\pgfqpoint{1.141855in}{1.383558in}}%
\pgfpathlineto{\pgfqpoint{1.181575in}{1.376785in}}%
\pgfpathlineto{\pgfqpoint{1.183516in}{1.359664in}}%
\pgfpathlineto{\pgfqpoint{1.186021in}{1.352257in}}%
\pgfpathlineto{\pgfqpoint{1.183026in}{1.348649in}}%
\pgfpathlineto{\pgfqpoint{1.136475in}{1.356525in}}%
\pgfpathlineto{\pgfqpoint{1.065800in}{1.369302in}}%
\pgfpathlineto{\pgfqpoint{1.069417in}{1.390629in}}%
\pgfpathlineto{\pgfqpoint{1.076288in}{1.396187in}}%
\pgfpathlineto{\pgfqpoint{1.103689in}{1.391163in}}%
\pgfpathclose%
\pgfusepath{fill}%
\end{pgfscope}%
\begin{pgfscope}%
\pgfpathrectangle{\pgfqpoint{0.100000in}{0.100000in}}{\pgfqpoint{3.608454in}{2.310000in}}%
\pgfusepath{clip}%
\pgfsetbuttcap%
\pgfsetmiterjoin%
\definecolor{currentfill}{rgb}{0.000000,0.643137,0.678431}%
\pgfsetfillcolor{currentfill}%
\pgfsetlinewidth{0.000000pt}%
\definecolor{currentstroke}{rgb}{0.000000,0.000000,0.000000}%
\pgfsetstrokecolor{currentstroke}%
\pgfsetstrokeopacity{0.000000}%
\pgfsetdash{}{0pt}%
\pgfpathmoveto{\pgfqpoint{2.123559in}{0.920377in}}%
\pgfpathlineto{\pgfqpoint{2.140773in}{0.915277in}}%
\pgfpathlineto{\pgfqpoint{2.141202in}{0.939334in}}%
\pgfpathlineto{\pgfqpoint{2.147313in}{0.939341in}}%
\pgfpathlineto{\pgfqpoint{2.157026in}{0.928424in}}%
\pgfpathlineto{\pgfqpoint{2.170065in}{0.926960in}}%
\pgfpathlineto{\pgfqpoint{2.175503in}{0.924555in}}%
\pgfpathlineto{\pgfqpoint{2.175110in}{0.918121in}}%
\pgfpathlineto{\pgfqpoint{2.183884in}{0.913246in}}%
\pgfpathlineto{\pgfqpoint{2.191374in}{0.903396in}}%
\pgfpathlineto{\pgfqpoint{2.190887in}{0.898874in}}%
\pgfpathlineto{\pgfqpoint{2.195345in}{0.891554in}}%
\pgfpathlineto{\pgfqpoint{2.189397in}{0.878197in}}%
\pgfpathlineto{\pgfqpoint{2.183288in}{0.876547in}}%
\pgfpathlineto{\pgfqpoint{2.182676in}{0.872245in}}%
\pgfpathlineto{\pgfqpoint{2.186423in}{0.866575in}}%
\pgfpathlineto{\pgfqpoint{2.170562in}{0.866376in}}%
\pgfpathlineto{\pgfqpoint{2.170684in}{0.855388in}}%
\pgfpathlineto{\pgfqpoint{2.126427in}{0.854941in}}%
\pgfpathlineto{\pgfqpoint{2.125475in}{0.856951in}}%
\pgfpathlineto{\pgfqpoint{2.096641in}{0.856769in}}%
\pgfpathlineto{\pgfqpoint{2.095246in}{0.861494in}}%
\pgfpathlineto{\pgfqpoint{2.086189in}{0.861574in}}%
\pgfpathlineto{\pgfqpoint{2.086328in}{0.894484in}}%
\pgfpathlineto{\pgfqpoint{2.100603in}{0.896104in}}%
\pgfpathlineto{\pgfqpoint{2.115003in}{0.895858in}}%
\pgfpathlineto{\pgfqpoint{2.123587in}{0.890672in}}%
\pgfpathlineto{\pgfqpoint{2.123559in}{0.920377in}}%
\pgfpathclose%
\pgfusepath{fill}%
\end{pgfscope}%
\begin{pgfscope}%
\pgfpathrectangle{\pgfqpoint{0.100000in}{0.100000in}}{\pgfqpoint{3.608454in}{2.310000in}}%
\pgfusepath{clip}%
\pgfsetbuttcap%
\pgfsetmiterjoin%
\definecolor{currentfill}{rgb}{0.000000,0.890196,0.554902}%
\pgfsetfillcolor{currentfill}%
\pgfsetlinewidth{0.000000pt}%
\definecolor{currentstroke}{rgb}{0.000000,0.000000,0.000000}%
\pgfsetstrokecolor{currentstroke}%
\pgfsetstrokeopacity{0.000000}%
\pgfsetdash{}{0pt}%
\pgfpathmoveto{\pgfqpoint{3.001283in}{0.581332in}}%
\pgfpathlineto{\pgfqpoint{2.990145in}{0.581978in}}%
\pgfpathlineto{\pgfqpoint{2.974786in}{0.579836in}}%
\pgfpathlineto{\pgfqpoint{2.972832in}{0.594908in}}%
\pgfpathlineto{\pgfqpoint{2.963258in}{0.602568in}}%
\pgfpathlineto{\pgfqpoint{2.965965in}{0.605834in}}%
\pgfpathlineto{\pgfqpoint{2.973159in}{0.605106in}}%
\pgfpathlineto{\pgfqpoint{2.978125in}{0.608670in}}%
\pgfpathlineto{\pgfqpoint{2.976237in}{0.622163in}}%
\pgfpathlineto{\pgfqpoint{2.985432in}{0.623511in}}%
\pgfpathlineto{\pgfqpoint{2.982209in}{0.644864in}}%
\pgfpathlineto{\pgfqpoint{2.995671in}{0.646679in}}%
\pgfpathlineto{\pgfqpoint{2.996330in}{0.642369in}}%
\pgfpathlineto{\pgfqpoint{3.006579in}{0.647225in}}%
\pgfpathlineto{\pgfqpoint{3.020704in}{0.653264in}}%
\pgfpathlineto{\pgfqpoint{3.025702in}{0.651277in}}%
\pgfpathlineto{\pgfqpoint{3.029457in}{0.642400in}}%
\pgfpathlineto{\pgfqpoint{3.037437in}{0.635958in}}%
\pgfpathlineto{\pgfqpoint{3.040352in}{0.617807in}}%
\pgfpathlineto{\pgfqpoint{3.039989in}{0.610672in}}%
\pgfpathlineto{\pgfqpoint{2.994626in}{0.604123in}}%
\pgfpathlineto{\pgfqpoint{3.006411in}{0.592107in}}%
\pgfpathlineto{\pgfqpoint{3.001283in}{0.581332in}}%
\pgfpathclose%
\pgfusepath{fill}%
\end{pgfscope}%
\begin{pgfscope}%
\pgfpathrectangle{\pgfqpoint{0.100000in}{0.100000in}}{\pgfqpoint{3.608454in}{2.310000in}}%
\pgfusepath{clip}%
\pgfsetbuttcap%
\pgfsetmiterjoin%
\definecolor{currentfill}{rgb}{0.000000,0.615686,0.692157}%
\pgfsetfillcolor{currentfill}%
\pgfsetlinewidth{0.000000pt}%
\definecolor{currentstroke}{rgb}{0.000000,0.000000,0.000000}%
\pgfsetstrokecolor{currentstroke}%
\pgfsetstrokeopacity{0.000000}%
\pgfsetdash{}{0pt}%
\pgfpathmoveto{\pgfqpoint{2.002513in}{1.074759in}}%
\pgfpathlineto{\pgfqpoint{2.030056in}{1.074337in}}%
\pgfpathlineto{\pgfqpoint{2.030288in}{1.091590in}}%
\pgfpathlineto{\pgfqpoint{2.057637in}{1.091246in}}%
\pgfpathlineto{\pgfqpoint{2.057547in}{1.080899in}}%
\pgfpathlineto{\pgfqpoint{2.060962in}{1.080885in}}%
\pgfpathlineto{\pgfqpoint{2.060863in}{1.067101in}}%
\pgfpathlineto{\pgfqpoint{2.053817in}{1.067155in}}%
\pgfpathlineto{\pgfqpoint{2.053684in}{1.060206in}}%
\pgfpathlineto{\pgfqpoint{2.050250in}{1.053326in}}%
\pgfpathlineto{\pgfqpoint{2.043425in}{1.053406in}}%
\pgfpathlineto{\pgfqpoint{2.043347in}{1.046522in}}%
\pgfpathlineto{\pgfqpoint{2.013608in}{1.046935in}}%
\pgfpathlineto{\pgfqpoint{2.010159in}{1.057965in}}%
\pgfpathlineto{\pgfqpoint{2.001980in}{1.055914in}}%
\pgfpathlineto{\pgfqpoint{2.002513in}{1.074759in}}%
\pgfpathclose%
\pgfusepath{fill}%
\end{pgfscope}%
\begin{pgfscope}%
\pgfpathrectangle{\pgfqpoint{0.100000in}{0.100000in}}{\pgfqpoint{3.608454in}{2.310000in}}%
\pgfusepath{clip}%
\pgfsetbuttcap%
\pgfsetmiterjoin%
\definecolor{currentfill}{rgb}{0.000000,0.721569,0.639216}%
\pgfsetfillcolor{currentfill}%
\pgfsetlinewidth{0.000000pt}%
\definecolor{currentstroke}{rgb}{0.000000,0.000000,0.000000}%
\pgfsetstrokecolor{currentstroke}%
\pgfsetstrokeopacity{0.000000}%
\pgfsetdash{}{0pt}%
\pgfpathmoveto{\pgfqpoint{2.949055in}{1.160645in}}%
\pgfpathlineto{\pgfqpoint{2.942760in}{1.162773in}}%
\pgfpathlineto{\pgfqpoint{2.938556in}{1.155078in}}%
\pgfpathlineto{\pgfqpoint{2.934882in}{1.156995in}}%
\pgfpathlineto{\pgfqpoint{2.930144in}{1.169183in}}%
\pgfpathlineto{\pgfqpoint{2.933147in}{1.171253in}}%
\pgfpathlineto{\pgfqpoint{2.936507in}{1.182790in}}%
\pgfpathlineto{\pgfqpoint{2.929794in}{1.191773in}}%
\pgfpathlineto{\pgfqpoint{2.933739in}{1.199536in}}%
\pgfpathlineto{\pgfqpoint{2.934516in}{1.206979in}}%
\pgfpathlineto{\pgfqpoint{2.940357in}{1.211530in}}%
\pgfpathlineto{\pgfqpoint{2.951682in}{1.213011in}}%
\pgfpathlineto{\pgfqpoint{2.950033in}{1.210864in}}%
\pgfpathlineto{\pgfqpoint{2.950303in}{1.201290in}}%
\pgfpathlineto{\pgfqpoint{2.948361in}{1.195220in}}%
\pgfpathlineto{\pgfqpoint{2.955269in}{1.192553in}}%
\pgfpathlineto{\pgfqpoint{2.968126in}{1.182900in}}%
\pgfpathlineto{\pgfqpoint{2.963826in}{1.175254in}}%
\pgfpathlineto{\pgfqpoint{2.955589in}{1.174438in}}%
\pgfpathlineto{\pgfqpoint{2.949055in}{1.160645in}}%
\pgfpathclose%
\pgfusepath{fill}%
\end{pgfscope}%
\begin{pgfscope}%
\pgfpathrectangle{\pgfqpoint{0.100000in}{0.100000in}}{\pgfqpoint{3.608454in}{2.310000in}}%
\pgfusepath{clip}%
\pgfsetbuttcap%
\pgfsetmiterjoin%
\definecolor{currentfill}{rgb}{0.000000,0.419608,0.790196}%
\pgfsetfillcolor{currentfill}%
\pgfsetlinewidth{0.000000pt}%
\definecolor{currentstroke}{rgb}{0.000000,0.000000,0.000000}%
\pgfsetstrokecolor{currentstroke}%
\pgfsetstrokeopacity{0.000000}%
\pgfsetdash{}{0pt}%
\pgfpathmoveto{\pgfqpoint{1.812489in}{1.842876in}}%
\pgfpathlineto{\pgfqpoint{1.843275in}{1.841097in}}%
\pgfpathlineto{\pgfqpoint{1.850854in}{1.840684in}}%
\pgfpathlineto{\pgfqpoint{1.849510in}{1.813058in}}%
\pgfpathlineto{\pgfqpoint{1.843805in}{1.813366in}}%
\pgfpathlineto{\pgfqpoint{1.802764in}{1.815650in}}%
\pgfpathlineto{\pgfqpoint{1.802162in}{1.823198in}}%
\pgfpathlineto{\pgfqpoint{1.809782in}{1.824532in}}%
\pgfpathlineto{\pgfqpoint{1.808240in}{1.831561in}}%
\pgfpathlineto{\pgfqpoint{1.812489in}{1.842876in}}%
\pgfpathclose%
\pgfusepath{fill}%
\end{pgfscope}%
\begin{pgfscope}%
\pgfpathrectangle{\pgfqpoint{0.100000in}{0.100000in}}{\pgfqpoint{3.608454in}{2.310000in}}%
\pgfusepath{clip}%
\pgfsetbuttcap%
\pgfsetmiterjoin%
\definecolor{currentfill}{rgb}{0.000000,0.666667,0.666667}%
\pgfsetfillcolor{currentfill}%
\pgfsetlinewidth{0.000000pt}%
\definecolor{currentstroke}{rgb}{0.000000,0.000000,0.000000}%
\pgfsetstrokecolor{currentstroke}%
\pgfsetstrokeopacity{0.000000}%
\pgfsetdash{}{0pt}%
\pgfpathmoveto{\pgfqpoint{1.498395in}{0.815195in}}%
\pgfpathlineto{\pgfqpoint{1.518670in}{0.813225in}}%
\pgfpathlineto{\pgfqpoint{1.562944in}{0.809149in}}%
\pgfpathlineto{\pgfqpoint{1.563559in}{0.816012in}}%
\pgfpathlineto{\pgfqpoint{1.581375in}{0.814353in}}%
\pgfpathlineto{\pgfqpoint{1.578420in}{0.779909in}}%
\pgfpathlineto{\pgfqpoint{1.580562in}{0.779720in}}%
\pgfpathlineto{\pgfqpoint{1.578119in}{0.751288in}}%
\pgfpathlineto{\pgfqpoint{1.574315in}{0.749565in}}%
\pgfpathlineto{\pgfqpoint{1.567213in}{0.756734in}}%
\pgfpathlineto{\pgfqpoint{1.562114in}{0.758888in}}%
\pgfpathlineto{\pgfqpoint{1.518651in}{0.714404in}}%
\pgfpathlineto{\pgfqpoint{1.486032in}{0.744751in}}%
\pgfpathlineto{\pgfqpoint{1.498395in}{0.815195in}}%
\pgfpathclose%
\pgfusepath{fill}%
\end{pgfscope}%
\begin{pgfscope}%
\pgfpathrectangle{\pgfqpoint{0.100000in}{0.100000in}}{\pgfqpoint{3.608454in}{2.310000in}}%
\pgfusepath{clip}%
\pgfsetbuttcap%
\pgfsetmiterjoin%
\definecolor{currentfill}{rgb}{0.000000,0.345098,0.827451}%
\pgfsetfillcolor{currentfill}%
\pgfsetlinewidth{0.000000pt}%
\definecolor{currentstroke}{rgb}{0.000000,0.000000,0.000000}%
\pgfsetstrokecolor{currentstroke}%
\pgfsetstrokeopacity{0.000000}%
\pgfsetdash{}{0pt}%
\pgfpathmoveto{\pgfqpoint{3.271268in}{1.222696in}}%
\pgfpathlineto{\pgfqpoint{3.265839in}{1.231923in}}%
\pgfpathlineto{\pgfqpoint{3.266234in}{1.243000in}}%
\pgfpathlineto{\pgfqpoint{3.259748in}{1.247063in}}%
\pgfpathlineto{\pgfqpoint{3.249894in}{1.248153in}}%
\pgfpathlineto{\pgfqpoint{3.248646in}{1.259422in}}%
\pgfpathlineto{\pgfqpoint{3.248223in}{1.267423in}}%
\pgfpathlineto{\pgfqpoint{3.266535in}{1.293082in}}%
\pgfpathlineto{\pgfqpoint{3.269493in}{1.294748in}}%
\pgfpathlineto{\pgfqpoint{3.275926in}{1.293375in}}%
\pgfpathlineto{\pgfqpoint{3.281080in}{1.300393in}}%
\pgfpathlineto{\pgfqpoint{3.293773in}{1.298157in}}%
\pgfpathlineto{\pgfqpoint{3.299684in}{1.300538in}}%
\pgfpathlineto{\pgfqpoint{3.305220in}{1.288703in}}%
\pgfpathlineto{\pgfqpoint{3.312824in}{1.275983in}}%
\pgfpathlineto{\pgfqpoint{3.324933in}{1.250084in}}%
\pgfpathlineto{\pgfqpoint{3.324107in}{1.249673in}}%
\pgfpathlineto{\pgfqpoint{3.316533in}{1.263940in}}%
\pgfpathlineto{\pgfqpoint{3.313407in}{1.258882in}}%
\pgfpathlineto{\pgfqpoint{3.326492in}{1.238123in}}%
\pgfpathlineto{\pgfqpoint{3.319080in}{1.242583in}}%
\pgfpathlineto{\pgfqpoint{3.311911in}{1.242756in}}%
\pgfpathlineto{\pgfqpoint{3.299704in}{1.233838in}}%
\pgfpathlineto{\pgfqpoint{3.289367in}{1.230321in}}%
\pgfpathlineto{\pgfqpoint{3.283250in}{1.223621in}}%
\pgfpathlineto{\pgfqpoint{3.271268in}{1.222696in}}%
\pgfpathclose%
\pgfusepath{fill}%
\end{pgfscope}%
\begin{pgfscope}%
\pgfpathrectangle{\pgfqpoint{0.100000in}{0.100000in}}{\pgfqpoint{3.608454in}{2.310000in}}%
\pgfusepath{clip}%
\pgfsetbuttcap%
\pgfsetmiterjoin%
\definecolor{currentfill}{rgb}{0.000000,0.776471,0.611765}%
\pgfsetfillcolor{currentfill}%
\pgfsetlinewidth{0.000000pt}%
\definecolor{currentstroke}{rgb}{0.000000,0.000000,0.000000}%
\pgfsetstrokecolor{currentstroke}%
\pgfsetstrokeopacity{0.000000}%
\pgfsetdash{}{0pt}%
\pgfpathmoveto{\pgfqpoint{2.002912in}{1.098705in}}%
\pgfpathlineto{\pgfqpoint{1.969632in}{1.099419in}}%
\pgfpathlineto{\pgfqpoint{1.956052in}{1.103169in}}%
\pgfpathlineto{\pgfqpoint{1.956419in}{1.117360in}}%
\pgfpathlineto{\pgfqpoint{1.949595in}{1.117555in}}%
\pgfpathlineto{\pgfqpoint{1.950456in}{1.151489in}}%
\pgfpathlineto{\pgfqpoint{1.984832in}{1.151230in}}%
\pgfpathlineto{\pgfqpoint{1.985971in}{1.147335in}}%
\pgfpathlineto{\pgfqpoint{1.977572in}{1.143978in}}%
\pgfpathlineto{\pgfqpoint{1.977272in}{1.130240in}}%
\pgfpathlineto{\pgfqpoint{1.984100in}{1.130082in}}%
\pgfpathlineto{\pgfqpoint{1.983944in}{1.123221in}}%
\pgfpathlineto{\pgfqpoint{1.990717in}{1.123080in}}%
\pgfpathlineto{\pgfqpoint{1.990571in}{1.116149in}}%
\pgfpathlineto{\pgfqpoint{2.003198in}{1.115990in}}%
\pgfpathlineto{\pgfqpoint{2.002912in}{1.098705in}}%
\pgfpathclose%
\pgfusepath{fill}%
\end{pgfscope}%
\begin{pgfscope}%
\pgfpathrectangle{\pgfqpoint{0.100000in}{0.100000in}}{\pgfqpoint{3.608454in}{2.310000in}}%
\pgfusepath{clip}%
\pgfsetbuttcap%
\pgfsetmiterjoin%
\definecolor{currentfill}{rgb}{0.000000,0.439216,0.780392}%
\pgfsetfillcolor{currentfill}%
\pgfsetlinewidth{0.000000pt}%
\definecolor{currentstroke}{rgb}{0.000000,0.000000,0.000000}%
\pgfsetstrokecolor{currentstroke}%
\pgfsetstrokeopacity{0.000000}%
\pgfsetdash{}{0pt}%
\pgfpathmoveto{\pgfqpoint{1.857457in}{1.534695in}}%
\pgfpathlineto{\pgfqpoint{1.855925in}{1.507277in}}%
\pgfpathlineto{\pgfqpoint{1.842917in}{1.507856in}}%
\pgfpathlineto{\pgfqpoint{1.795374in}{1.510393in}}%
\pgfpathlineto{\pgfqpoint{1.796934in}{1.565224in}}%
\pgfpathlineto{\pgfqpoint{1.762096in}{1.567407in}}%
\pgfpathlineto{\pgfqpoint{1.763649in}{1.594885in}}%
\pgfpathlineto{\pgfqpoint{1.803419in}{1.592288in}}%
\pgfpathlineto{\pgfqpoint{1.831612in}{1.590776in}}%
\pgfpathlineto{\pgfqpoint{1.830208in}{1.563450in}}%
\pgfpathlineto{\pgfqpoint{1.858208in}{1.562101in}}%
\pgfpathlineto{\pgfqpoint{1.857457in}{1.534695in}}%
\pgfpathclose%
\pgfusepath{fill}%
\end{pgfscope}%
\begin{pgfscope}%
\pgfpathrectangle{\pgfqpoint{0.100000in}{0.100000in}}{\pgfqpoint{3.608454in}{2.310000in}}%
\pgfusepath{clip}%
\pgfsetbuttcap%
\pgfsetmiterjoin%
\definecolor{currentfill}{rgb}{0.000000,0.576471,0.711765}%
\pgfsetfillcolor{currentfill}%
\pgfsetlinewidth{0.000000pt}%
\definecolor{currentstroke}{rgb}{0.000000,0.000000,0.000000}%
\pgfsetstrokecolor{currentstroke}%
\pgfsetstrokeopacity{0.000000}%
\pgfsetdash{}{0pt}%
\pgfpathmoveto{\pgfqpoint{2.410748in}{1.065840in}}%
\pgfpathlineto{\pgfqpoint{2.389852in}{1.065601in}}%
\pgfpathlineto{\pgfqpoint{2.362249in}{1.063873in}}%
\pgfpathlineto{\pgfqpoint{2.361164in}{1.098864in}}%
\pgfpathlineto{\pgfqpoint{2.350507in}{1.098801in}}%
\pgfpathlineto{\pgfqpoint{2.340339in}{1.098431in}}%
\pgfpathlineto{\pgfqpoint{2.339731in}{1.125367in}}%
\pgfpathlineto{\pgfqpoint{2.345455in}{1.126590in}}%
\pgfpathlineto{\pgfqpoint{2.343960in}{1.136912in}}%
\pgfpathlineto{\pgfqpoint{2.335244in}{1.146265in}}%
\pgfpathlineto{\pgfqpoint{2.353073in}{1.147045in}}%
\pgfpathlineto{\pgfqpoint{2.410619in}{1.149798in}}%
\pgfpathlineto{\pgfqpoint{2.421101in}{1.141357in}}%
\pgfpathlineto{\pgfqpoint{2.420239in}{1.132410in}}%
\pgfpathlineto{\pgfqpoint{2.411938in}{1.125058in}}%
\pgfpathlineto{\pgfqpoint{2.406026in}{1.117238in}}%
\pgfpathlineto{\pgfqpoint{2.408349in}{1.109943in}}%
\pgfpathlineto{\pgfqpoint{2.410748in}{1.065840in}}%
\pgfpathclose%
\pgfusepath{fill}%
\end{pgfscope}%
\begin{pgfscope}%
\pgfpathrectangle{\pgfqpoint{0.100000in}{0.100000in}}{\pgfqpoint{3.608454in}{2.310000in}}%
\pgfusepath{clip}%
\pgfsetbuttcap%
\pgfsetmiterjoin%
\definecolor{currentfill}{rgb}{0.000000,0.650980,0.674510}%
\pgfsetfillcolor{currentfill}%
\pgfsetlinewidth{0.000000pt}%
\definecolor{currentstroke}{rgb}{0.000000,0.000000,0.000000}%
\pgfsetstrokecolor{currentstroke}%
\pgfsetstrokeopacity{0.000000}%
\pgfsetdash{}{0pt}%
\pgfpathmoveto{\pgfqpoint{1.383780in}{1.226748in}}%
\pgfpathlineto{\pgfqpoint{1.373495in}{1.246788in}}%
\pgfpathlineto{\pgfqpoint{1.375196in}{1.259964in}}%
\pgfpathlineto{\pgfqpoint{1.347103in}{1.263333in}}%
\pgfpathlineto{\pgfqpoint{1.347428in}{1.265725in}}%
\pgfpathlineto{\pgfqpoint{1.349065in}{1.285654in}}%
\pgfpathlineto{\pgfqpoint{1.352269in}{1.306231in}}%
\pgfpathlineto{\pgfqpoint{1.360950in}{1.306443in}}%
\pgfpathlineto{\pgfqpoint{1.365837in}{1.343240in}}%
\pgfpathlineto{\pgfqpoint{1.412097in}{1.337063in}}%
\pgfpathlineto{\pgfqpoint{1.422616in}{1.335846in}}%
\pgfpathlineto{\pgfqpoint{1.426840in}{1.337144in}}%
\pgfpathlineto{\pgfqpoint{1.444613in}{1.311073in}}%
\pgfpathlineto{\pgfqpoint{1.448675in}{1.296375in}}%
\pgfpathlineto{\pgfqpoint{1.454881in}{1.289984in}}%
\pgfpathlineto{\pgfqpoint{1.457227in}{1.286192in}}%
\pgfpathlineto{\pgfqpoint{1.450742in}{1.268060in}}%
\pgfpathlineto{\pgfqpoint{1.463748in}{1.269672in}}%
\pgfpathlineto{\pgfqpoint{1.469792in}{1.261228in}}%
\pgfpathlineto{\pgfqpoint{1.471264in}{1.248705in}}%
\pgfpathlineto{\pgfqpoint{1.469341in}{1.240362in}}%
\pgfpathlineto{\pgfqpoint{1.466660in}{1.216957in}}%
\pgfpathlineto{\pgfqpoint{1.462544in}{1.217413in}}%
\pgfpathlineto{\pgfqpoint{1.383780in}{1.226748in}}%
\pgfpathclose%
\pgfusepath{fill}%
\end{pgfscope}%
\begin{pgfscope}%
\pgfpathrectangle{\pgfqpoint{0.100000in}{0.100000in}}{\pgfqpoint{3.608454in}{2.310000in}}%
\pgfusepath{clip}%
\pgfsetbuttcap%
\pgfsetmiterjoin%
\definecolor{currentfill}{rgb}{0.000000,0.427451,0.786275}%
\pgfsetfillcolor{currentfill}%
\pgfsetlinewidth{0.000000pt}%
\definecolor{currentstroke}{rgb}{0.000000,0.000000,0.000000}%
\pgfsetstrokecolor{currentstroke}%
\pgfsetstrokeopacity{0.000000}%
\pgfsetdash{}{0pt}%
\pgfpathmoveto{\pgfqpoint{2.204405in}{1.465435in}}%
\pgfpathlineto{\pgfqpoint{2.193447in}{1.465250in}}%
\pgfpathlineto{\pgfqpoint{2.192689in}{1.511198in}}%
\pgfpathlineto{\pgfqpoint{2.206358in}{1.511376in}}%
\pgfpathlineto{\pgfqpoint{2.219965in}{1.511654in}}%
\pgfpathlineto{\pgfqpoint{2.220496in}{1.490915in}}%
\pgfpathlineto{\pgfqpoint{2.247888in}{1.491592in}}%
\pgfpathlineto{\pgfqpoint{2.248598in}{1.467244in}}%
\pgfpathlineto{\pgfqpoint{2.244011in}{1.467014in}}%
\pgfpathlineto{\pgfqpoint{2.204405in}{1.465435in}}%
\pgfpathclose%
\pgfusepath{fill}%
\end{pgfscope}%
\begin{pgfscope}%
\pgfpathrectangle{\pgfqpoint{0.100000in}{0.100000in}}{\pgfqpoint{3.608454in}{2.310000in}}%
\pgfusepath{clip}%
\pgfsetbuttcap%
\pgfsetmiterjoin%
\definecolor{currentfill}{rgb}{0.000000,0.180392,0.909804}%
\pgfsetfillcolor{currentfill}%
\pgfsetlinewidth{0.000000pt}%
\definecolor{currentstroke}{rgb}{0.000000,0.000000,0.000000}%
\pgfsetstrokecolor{currentstroke}%
\pgfsetstrokeopacity{0.000000}%
\pgfsetdash{}{0pt}%
\pgfpathmoveto{\pgfqpoint{1.945778in}{0.610158in}}%
\pgfpathlineto{\pgfqpoint{1.925990in}{0.598363in}}%
\pgfpathlineto{\pgfqpoint{1.923573in}{0.599973in}}%
\pgfpathlineto{\pgfqpoint{1.918807in}{0.601623in}}%
\pgfpathlineto{\pgfqpoint{1.907278in}{0.617155in}}%
\pgfpathlineto{\pgfqpoint{1.898375in}{0.608966in}}%
\pgfpathlineto{\pgfqpoint{1.896461in}{0.616726in}}%
\pgfpathlineto{\pgfqpoint{1.878483in}{0.632521in}}%
\pgfpathlineto{\pgfqpoint{1.870125in}{0.624788in}}%
\pgfpathlineto{\pgfqpoint{1.858713in}{0.641352in}}%
\pgfpathlineto{\pgfqpoint{1.859593in}{0.670192in}}%
\pgfpathlineto{\pgfqpoint{1.871994in}{0.669892in}}%
\pgfpathlineto{\pgfqpoint{1.891457in}{0.663121in}}%
\pgfpathlineto{\pgfqpoint{1.892292in}{0.669010in}}%
\pgfpathlineto{\pgfqpoint{1.903598in}{0.691454in}}%
\pgfpathlineto{\pgfqpoint{1.913119in}{0.701138in}}%
\pgfpathlineto{\pgfqpoint{1.926963in}{0.697491in}}%
\pgfpathlineto{\pgfqpoint{1.947915in}{0.687499in}}%
\pgfpathlineto{\pgfqpoint{1.952059in}{0.698355in}}%
\pgfpathlineto{\pgfqpoint{1.965169in}{0.705757in}}%
\pgfpathlineto{\pgfqpoint{1.981851in}{0.714949in}}%
\pgfpathlineto{\pgfqpoint{1.984720in}{0.711352in}}%
\pgfpathlineto{\pgfqpoint{1.986509in}{0.700655in}}%
\pgfpathlineto{\pgfqpoint{1.993815in}{0.692311in}}%
\pgfpathlineto{\pgfqpoint{1.995500in}{0.684736in}}%
\pgfpathlineto{\pgfqpoint{1.971709in}{0.671338in}}%
\pgfpathlineto{\pgfqpoint{1.975419in}{0.668492in}}%
\pgfpathlineto{\pgfqpoint{1.978079in}{0.661346in}}%
\pgfpathlineto{\pgfqpoint{1.986455in}{0.652084in}}%
\pgfpathlineto{\pgfqpoint{1.993561in}{0.650091in}}%
\pgfpathlineto{\pgfqpoint{1.987116in}{0.645931in}}%
\pgfpathlineto{\pgfqpoint{1.982721in}{0.639405in}}%
\pgfpathlineto{\pgfqpoint{1.979102in}{0.635917in}}%
\pgfpathlineto{\pgfqpoint{1.966633in}{0.631014in}}%
\pgfpathlineto{\pgfqpoint{1.945778in}{0.610158in}}%
\pgfpathclose%
\pgfusepath{fill}%
\end{pgfscope}%
\begin{pgfscope}%
\pgfpathrectangle{\pgfqpoint{0.100000in}{0.100000in}}{\pgfqpoint{3.608454in}{2.310000in}}%
\pgfusepath{clip}%
\pgfsetbuttcap%
\pgfsetmiterjoin%
\definecolor{currentfill}{rgb}{0.000000,0.690196,0.654902}%
\pgfsetfillcolor{currentfill}%
\pgfsetlinewidth{0.000000pt}%
\definecolor{currentstroke}{rgb}{0.000000,0.000000,0.000000}%
\pgfsetstrokecolor{currentstroke}%
\pgfsetstrokeopacity{0.000000}%
\pgfsetdash{}{0pt}%
\pgfpathmoveto{\pgfqpoint{2.288408in}{0.781863in}}%
\pgfpathlineto{\pgfqpoint{2.287535in}{0.813168in}}%
\pgfpathlineto{\pgfqpoint{2.280375in}{0.819890in}}%
\pgfpathlineto{\pgfqpoint{2.279806in}{0.840743in}}%
\pgfpathlineto{\pgfqpoint{2.272892in}{0.840599in}}%
\pgfpathlineto{\pgfqpoint{2.267654in}{0.847225in}}%
\pgfpathlineto{\pgfqpoint{2.258858in}{0.847301in}}%
\pgfpathlineto{\pgfqpoint{2.258381in}{0.867551in}}%
\pgfpathlineto{\pgfqpoint{2.302086in}{0.868312in}}%
\pgfpathlineto{\pgfqpoint{2.342635in}{0.869578in}}%
\pgfpathlineto{\pgfqpoint{2.344274in}{0.869656in}}%
\pgfpathlineto{\pgfqpoint{2.340706in}{0.860346in}}%
\pgfpathlineto{\pgfqpoint{2.337645in}{0.859095in}}%
\pgfpathlineto{\pgfqpoint{2.331640in}{0.845019in}}%
\pgfpathlineto{\pgfqpoint{2.334864in}{0.835567in}}%
\pgfpathlineto{\pgfqpoint{2.344959in}{0.835935in}}%
\pgfpathlineto{\pgfqpoint{2.344269in}{0.832393in}}%
\pgfpathlineto{\pgfqpoint{2.342569in}{0.828917in}}%
\pgfpathlineto{\pgfqpoint{2.344014in}{0.814843in}}%
\pgfpathlineto{\pgfqpoint{2.340198in}{0.812791in}}%
\pgfpathlineto{\pgfqpoint{2.342946in}{0.805709in}}%
\pgfpathlineto{\pgfqpoint{2.343388in}{0.795089in}}%
\pgfpathlineto{\pgfqpoint{2.338322in}{0.779971in}}%
\pgfpathlineto{\pgfqpoint{2.336629in}{0.786857in}}%
\pgfpathlineto{\pgfqpoint{2.324523in}{0.778629in}}%
\pgfpathlineto{\pgfqpoint{2.319400in}{0.785321in}}%
\pgfpathlineto{\pgfqpoint{2.315987in}{0.782756in}}%
\pgfpathlineto{\pgfqpoint{2.288408in}{0.781863in}}%
\pgfpathclose%
\pgfusepath{fill}%
\end{pgfscope}%
\begin{pgfscope}%
\pgfpathrectangle{\pgfqpoint{0.100000in}{0.100000in}}{\pgfqpoint{3.608454in}{2.310000in}}%
\pgfusepath{clip}%
\pgfsetbuttcap%
\pgfsetmiterjoin%
\definecolor{currentfill}{rgb}{0.000000,0.776471,0.611765}%
\pgfsetfillcolor{currentfill}%
\pgfsetlinewidth{0.000000pt}%
\definecolor{currentstroke}{rgb}{0.000000,0.000000,0.000000}%
\pgfsetstrokecolor{currentstroke}%
\pgfsetstrokeopacity{0.000000}%
\pgfsetdash{}{0pt}%
\pgfpathmoveto{\pgfqpoint{2.611522in}{1.917144in}}%
\pgfpathlineto{\pgfqpoint{2.589883in}{1.913322in}}%
\pgfpathlineto{\pgfqpoint{2.586900in}{1.909857in}}%
\pgfpathlineto{\pgfqpoint{2.586216in}{1.900993in}}%
\pgfpathlineto{\pgfqpoint{2.576778in}{1.896148in}}%
\pgfpathlineto{\pgfqpoint{2.572058in}{1.888198in}}%
\pgfpathlineto{\pgfqpoint{2.566665in}{1.889771in}}%
\pgfpathlineto{\pgfqpoint{2.572918in}{1.898549in}}%
\pgfpathlineto{\pgfqpoint{2.575837in}{1.905644in}}%
\pgfpathlineto{\pgfqpoint{2.562752in}{1.899838in}}%
\pgfpathlineto{\pgfqpoint{2.559587in}{1.892271in}}%
\pgfpathlineto{\pgfqpoint{2.552985in}{1.887697in}}%
\pgfpathlineto{\pgfqpoint{2.547425in}{1.892494in}}%
\pgfpathlineto{\pgfqpoint{2.541575in}{1.885933in}}%
\pgfpathlineto{\pgfqpoint{2.537283in}{1.876850in}}%
\pgfpathlineto{\pgfqpoint{2.533757in}{1.876549in}}%
\pgfpathlineto{\pgfqpoint{2.531516in}{1.904166in}}%
\pgfpathlineto{\pgfqpoint{2.528620in}{1.910846in}}%
\pgfpathlineto{\pgfqpoint{2.514867in}{1.909757in}}%
\pgfpathlineto{\pgfqpoint{2.513243in}{1.930478in}}%
\pgfpathlineto{\pgfqpoint{2.485703in}{1.928298in}}%
\pgfpathlineto{\pgfqpoint{2.484672in}{1.942048in}}%
\pgfpathlineto{\pgfqpoint{2.483177in}{1.962621in}}%
\pgfpathlineto{\pgfqpoint{2.489953in}{1.963218in}}%
\pgfpathlineto{\pgfqpoint{2.486521in}{1.969842in}}%
\pgfpathlineto{\pgfqpoint{2.485656in}{1.981514in}}%
\pgfpathlineto{\pgfqpoint{2.493445in}{1.982175in}}%
\pgfpathlineto{\pgfqpoint{2.508763in}{1.975427in}}%
\pgfpathlineto{\pgfqpoint{2.511603in}{1.969882in}}%
\pgfpathlineto{\pgfqpoint{2.526869in}{1.951632in}}%
\pgfpathlineto{\pgfqpoint{2.538957in}{1.952508in}}%
\pgfpathlineto{\pgfqpoint{2.544000in}{1.955889in}}%
\pgfpathlineto{\pgfqpoint{2.552498in}{1.948973in}}%
\pgfpathlineto{\pgfqpoint{2.564052in}{1.959757in}}%
\pgfpathlineto{\pgfqpoint{2.568367in}{1.953169in}}%
\pgfpathlineto{\pgfqpoint{2.574148in}{1.960469in}}%
\pgfpathlineto{\pgfqpoint{2.592197in}{1.970111in}}%
\pgfpathlineto{\pgfqpoint{2.605875in}{1.974408in}}%
\pgfpathlineto{\pgfqpoint{2.611522in}{1.917144in}}%
\pgfpathclose%
\pgfusepath{fill}%
\end{pgfscope}%
\begin{pgfscope}%
\pgfpathrectangle{\pgfqpoint{0.100000in}{0.100000in}}{\pgfqpoint{3.608454in}{2.310000in}}%
\pgfusepath{clip}%
\pgfsetbuttcap%
\pgfsetmiterjoin%
\definecolor{currentfill}{rgb}{0.000000,0.537255,0.731373}%
\pgfsetfillcolor{currentfill}%
\pgfsetlinewidth{0.000000pt}%
\definecolor{currentstroke}{rgb}{0.000000,0.000000,0.000000}%
\pgfsetstrokecolor{currentstroke}%
\pgfsetstrokeopacity{0.000000}%
\pgfsetdash{}{0pt}%
\pgfpathmoveto{\pgfqpoint{3.284382in}{1.624557in}}%
\pgfpathlineto{\pgfqpoint{3.276398in}{1.614345in}}%
\pgfpathlineto{\pgfqpoint{3.266342in}{1.608149in}}%
\pgfpathlineto{\pgfqpoint{3.258629in}{1.617098in}}%
\pgfpathlineto{\pgfqpoint{3.262436in}{1.621007in}}%
\pgfpathlineto{\pgfqpoint{3.252558in}{1.629161in}}%
\pgfpathlineto{\pgfqpoint{3.248335in}{1.629894in}}%
\pgfpathlineto{\pgfqpoint{3.244631in}{1.624617in}}%
\pgfpathlineto{\pgfqpoint{3.248038in}{1.619832in}}%
\pgfpathlineto{\pgfqpoint{3.234122in}{1.609162in}}%
\pgfpathlineto{\pgfqpoint{3.221218in}{1.609411in}}%
\pgfpathlineto{\pgfqpoint{3.217800in}{1.603350in}}%
\pgfpathlineto{\pgfqpoint{3.217800in}{1.596744in}}%
\pgfpathlineto{\pgfqpoint{3.213947in}{1.593817in}}%
\pgfpathlineto{\pgfqpoint{3.211501in}{1.598059in}}%
\pgfpathlineto{\pgfqpoint{3.206320in}{1.596141in}}%
\pgfpathlineto{\pgfqpoint{3.203466in}{1.605423in}}%
\pgfpathlineto{\pgfqpoint{3.195360in}{1.613257in}}%
\pgfpathlineto{\pgfqpoint{3.192630in}{1.620121in}}%
\pgfpathlineto{\pgfqpoint{3.195407in}{1.620855in}}%
\pgfpathlineto{\pgfqpoint{3.203745in}{1.634729in}}%
\pgfpathlineto{\pgfqpoint{3.209405in}{1.635945in}}%
\pgfpathlineto{\pgfqpoint{3.210875in}{1.652502in}}%
\pgfpathlineto{\pgfqpoint{3.210379in}{1.663470in}}%
\pgfpathlineto{\pgfqpoint{3.215154in}{1.664914in}}%
\pgfpathlineto{\pgfqpoint{3.253079in}{1.672195in}}%
\pgfpathlineto{\pgfqpoint{3.245833in}{1.699582in}}%
\pgfpathlineto{\pgfqpoint{3.252945in}{1.701154in}}%
\pgfpathlineto{\pgfqpoint{3.258600in}{1.697493in}}%
\pgfpathlineto{\pgfqpoint{3.260640in}{1.692420in}}%
\pgfpathlineto{\pgfqpoint{3.267733in}{1.692380in}}%
\pgfpathlineto{\pgfqpoint{3.274806in}{1.685909in}}%
\pgfpathlineto{\pgfqpoint{3.277583in}{1.676067in}}%
\pgfpathlineto{\pgfqpoint{3.282783in}{1.667998in}}%
\pgfpathlineto{\pgfqpoint{3.292838in}{1.664005in}}%
\pgfpathlineto{\pgfqpoint{3.298905in}{1.664653in}}%
\pgfpathlineto{\pgfqpoint{3.302353in}{1.660141in}}%
\pgfpathlineto{\pgfqpoint{3.295750in}{1.652950in}}%
\pgfpathlineto{\pgfqpoint{3.294577in}{1.644062in}}%
\pgfpathlineto{\pgfqpoint{3.284382in}{1.624557in}}%
\pgfpathclose%
\pgfusepath{fill}%
\end{pgfscope}%
\begin{pgfscope}%
\pgfpathrectangle{\pgfqpoint{0.100000in}{0.100000in}}{\pgfqpoint{3.608454in}{2.310000in}}%
\pgfusepath{clip}%
\pgfsetbuttcap%
\pgfsetmiterjoin%
\definecolor{currentfill}{rgb}{0.000000,0.364706,0.817647}%
\pgfsetfillcolor{currentfill}%
\pgfsetlinewidth{0.000000pt}%
\definecolor{currentstroke}{rgb}{0.000000,0.000000,0.000000}%
\pgfsetstrokecolor{currentstroke}%
\pgfsetstrokeopacity{0.000000}%
\pgfsetdash{}{0pt}%
\pgfpathmoveto{\pgfqpoint{1.858721in}{1.355685in}}%
\pgfpathlineto{\pgfqpoint{1.892889in}{1.354245in}}%
\pgfpathlineto{\pgfqpoint{1.892401in}{1.333545in}}%
\pgfpathlineto{\pgfqpoint{1.891797in}{1.319820in}}%
\pgfpathlineto{\pgfqpoint{1.824072in}{1.322793in}}%
\pgfpathlineto{\pgfqpoint{1.789033in}{1.324734in}}%
\pgfpathlineto{\pgfqpoint{1.791350in}{1.359031in}}%
\pgfpathlineto{\pgfqpoint{1.858721in}{1.355685in}}%
\pgfpathclose%
\pgfusepath{fill}%
\end{pgfscope}%
\begin{pgfscope}%
\pgfpathrectangle{\pgfqpoint{0.100000in}{0.100000in}}{\pgfqpoint{3.608454in}{2.310000in}}%
\pgfusepath{clip}%
\pgfsetbuttcap%
\pgfsetmiterjoin%
\definecolor{currentfill}{rgb}{0.000000,0.827451,0.586275}%
\pgfsetfillcolor{currentfill}%
\pgfsetlinewidth{0.000000pt}%
\definecolor{currentstroke}{rgb}{0.000000,0.000000,0.000000}%
\pgfsetstrokecolor{currentstroke}%
\pgfsetstrokeopacity{0.000000}%
\pgfsetdash{}{0pt}%
\pgfpathmoveto{\pgfqpoint{0.592949in}{0.431644in}}%
\pgfpathlineto{\pgfqpoint{0.591752in}{0.435712in}}%
\pgfpathlineto{\pgfqpoint{0.594116in}{0.436502in}}%
\pgfpathlineto{\pgfqpoint{0.596254in}{0.435454in}}%
\pgfpathlineto{\pgfqpoint{0.596813in}{0.431481in}}%
\pgfpathlineto{\pgfqpoint{0.595244in}{0.430874in}}%
\pgfpathlineto{\pgfqpoint{0.592949in}{0.431644in}}%
\pgfpathclose%
\pgfusepath{fill}%
\end{pgfscope}%
\begin{pgfscope}%
\pgfpathrectangle{\pgfqpoint{0.100000in}{0.100000in}}{\pgfqpoint{3.608454in}{2.310000in}}%
\pgfusepath{clip}%
\pgfsetbuttcap%
\pgfsetmiterjoin%
\definecolor{currentfill}{rgb}{0.000000,0.827451,0.586275}%
\pgfsetfillcolor{currentfill}%
\pgfsetlinewidth{0.000000pt}%
\definecolor{currentstroke}{rgb}{0.000000,0.000000,0.000000}%
\pgfsetstrokecolor{currentstroke}%
\pgfsetstrokeopacity{0.000000}%
\pgfsetdash{}{0pt}%
\pgfpathmoveto{\pgfqpoint{0.598016in}{0.429456in}}%
\pgfpathlineto{\pgfqpoint{0.599382in}{0.432202in}}%
\pgfpathlineto{\pgfqpoint{0.600924in}{0.431678in}}%
\pgfpathlineto{\pgfqpoint{0.601971in}{0.429670in}}%
\pgfpathlineto{\pgfqpoint{0.598016in}{0.429456in}}%
\pgfpathclose%
\pgfusepath{fill}%
\end{pgfscope}%
\begin{pgfscope}%
\pgfpathrectangle{\pgfqpoint{0.100000in}{0.100000in}}{\pgfqpoint{3.608454in}{2.310000in}}%
\pgfusepath{clip}%
\pgfsetbuttcap%
\pgfsetmiterjoin%
\definecolor{currentfill}{rgb}{0.000000,0.827451,0.586275}%
\pgfsetfillcolor{currentfill}%
\pgfsetlinewidth{0.000000pt}%
\definecolor{currentstroke}{rgb}{0.000000,0.000000,0.000000}%
\pgfsetstrokecolor{currentstroke}%
\pgfsetstrokeopacity{0.000000}%
\pgfsetdash{}{0pt}%
\pgfpathmoveto{\pgfqpoint{0.679037in}{0.397130in}}%
\pgfpathlineto{\pgfqpoint{0.677306in}{0.394866in}}%
\pgfpathlineto{\pgfqpoint{0.676295in}{0.398443in}}%
\pgfpathlineto{\pgfqpoint{0.673518in}{0.397756in}}%
\pgfpathlineto{\pgfqpoint{0.676763in}{0.402425in}}%
\pgfpathlineto{\pgfqpoint{0.680361in}{0.402575in}}%
\pgfpathlineto{\pgfqpoint{0.679428in}{0.400744in}}%
\pgfpathlineto{\pgfqpoint{0.680747in}{0.399061in}}%
\pgfpathlineto{\pgfqpoint{0.683725in}{0.397976in}}%
\pgfpathlineto{\pgfqpoint{0.681665in}{0.395162in}}%
\pgfpathlineto{\pgfqpoint{0.680358in}{0.397560in}}%
\pgfpathlineto{\pgfqpoint{0.679037in}{0.397130in}}%
\pgfpathclose%
\pgfusepath{fill}%
\end{pgfscope}%
\begin{pgfscope}%
\pgfpathrectangle{\pgfqpoint{0.100000in}{0.100000in}}{\pgfqpoint{3.608454in}{2.310000in}}%
\pgfusepath{clip}%
\pgfsetbuttcap%
\pgfsetmiterjoin%
\definecolor{currentfill}{rgb}{0.000000,0.827451,0.586275}%
\pgfsetfillcolor{currentfill}%
\pgfsetlinewidth{0.000000pt}%
\definecolor{currentstroke}{rgb}{0.000000,0.000000,0.000000}%
\pgfsetstrokecolor{currentstroke}%
\pgfsetstrokeopacity{0.000000}%
\pgfsetdash{}{0pt}%
\pgfpathmoveto{\pgfqpoint{0.650134in}{0.407855in}}%
\pgfpathlineto{\pgfqpoint{0.650363in}{0.410897in}}%
\pgfpathlineto{\pgfqpoint{0.653136in}{0.411103in}}%
\pgfpathlineto{\pgfqpoint{0.652771in}{0.408615in}}%
\pgfpathlineto{\pgfqpoint{0.650134in}{0.407855in}}%
\pgfpathclose%
\pgfusepath{fill}%
\end{pgfscope}%
\begin{pgfscope}%
\pgfpathrectangle{\pgfqpoint{0.100000in}{0.100000in}}{\pgfqpoint{3.608454in}{2.310000in}}%
\pgfusepath{clip}%
\pgfsetbuttcap%
\pgfsetmiterjoin%
\definecolor{currentfill}{rgb}{0.000000,0.827451,0.586275}%
\pgfsetfillcolor{currentfill}%
\pgfsetlinewidth{0.000000pt}%
\definecolor{currentstroke}{rgb}{0.000000,0.000000,0.000000}%
\pgfsetstrokecolor{currentstroke}%
\pgfsetstrokeopacity{0.000000}%
\pgfsetdash{}{0pt}%
\pgfpathmoveto{\pgfqpoint{0.676688in}{0.386431in}}%
\pgfpathlineto{\pgfqpoint{0.683893in}{0.390538in}}%
\pgfpathlineto{\pgfqpoint{0.685712in}{0.388279in}}%
\pgfpathlineto{\pgfqpoint{0.683374in}{0.387317in}}%
\pgfpathlineto{\pgfqpoint{0.676688in}{0.386431in}}%
\pgfpathclose%
\pgfusepath{fill}%
\end{pgfscope}%
\begin{pgfscope}%
\pgfpathrectangle{\pgfqpoint{0.100000in}{0.100000in}}{\pgfqpoint{3.608454in}{2.310000in}}%
\pgfusepath{clip}%
\pgfsetbuttcap%
\pgfsetmiterjoin%
\definecolor{currentfill}{rgb}{0.000000,0.827451,0.586275}%
\pgfsetfillcolor{currentfill}%
\pgfsetlinewidth{0.000000pt}%
\definecolor{currentstroke}{rgb}{0.000000,0.000000,0.000000}%
\pgfsetstrokecolor{currentstroke}%
\pgfsetstrokeopacity{0.000000}%
\pgfsetdash{}{0pt}%
\pgfpathmoveto{\pgfqpoint{0.698256in}{0.395576in}}%
\pgfpathlineto{\pgfqpoint{0.696998in}{0.397305in}}%
\pgfpathlineto{\pgfqpoint{0.701060in}{0.399308in}}%
\pgfpathlineto{\pgfqpoint{0.699966in}{0.402043in}}%
\pgfpathlineto{\pgfqpoint{0.697786in}{0.402037in}}%
\pgfpathlineto{\pgfqpoint{0.696446in}{0.403146in}}%
\pgfpathlineto{\pgfqpoint{0.691042in}{0.403748in}}%
\pgfpathlineto{\pgfqpoint{0.686539in}{0.403836in}}%
\pgfpathlineto{\pgfqpoint{0.684692in}{0.405317in}}%
\pgfpathlineto{\pgfqpoint{0.682874in}{0.404197in}}%
\pgfpathlineto{\pgfqpoint{0.680377in}{0.404275in}}%
\pgfpathlineto{\pgfqpoint{0.680706in}{0.406672in}}%
\pgfpathlineto{\pgfqpoint{0.675294in}{0.406801in}}%
\pgfpathlineto{\pgfqpoint{0.672429in}{0.407601in}}%
\pgfpathlineto{\pgfqpoint{0.669990in}{0.409978in}}%
\pgfpathlineto{\pgfqpoint{0.672102in}{0.411831in}}%
\pgfpathlineto{\pgfqpoint{0.676413in}{0.414087in}}%
\pgfpathlineto{\pgfqpoint{0.674276in}{0.416221in}}%
\pgfpathlineto{\pgfqpoint{0.672028in}{0.416339in}}%
\pgfpathlineto{\pgfqpoint{0.667446in}{0.412698in}}%
\pgfpathlineto{\pgfqpoint{0.664420in}{0.412115in}}%
\pgfpathlineto{\pgfqpoint{0.661755in}{0.412674in}}%
\pgfpathlineto{\pgfqpoint{0.660202in}{0.409906in}}%
\pgfpathlineto{\pgfqpoint{0.657982in}{0.409975in}}%
\pgfpathlineto{\pgfqpoint{0.652958in}{0.413217in}}%
\pgfpathlineto{\pgfqpoint{0.652756in}{0.415126in}}%
\pgfpathlineto{\pgfqpoint{0.654783in}{0.415315in}}%
\pgfpathlineto{\pgfqpoint{0.655313in}{0.417608in}}%
\pgfpathlineto{\pgfqpoint{0.654939in}{0.420817in}}%
\pgfpathlineto{\pgfqpoint{0.650913in}{0.416098in}}%
\pgfpathlineto{\pgfqpoint{0.650688in}{0.414349in}}%
\pgfpathlineto{\pgfqpoint{0.647513in}{0.414740in}}%
\pgfpathlineto{\pgfqpoint{0.645402in}{0.416502in}}%
\pgfpathlineto{\pgfqpoint{0.646738in}{0.419922in}}%
\pgfpathlineto{\pgfqpoint{0.645503in}{0.423049in}}%
\pgfpathlineto{\pgfqpoint{0.644228in}{0.422685in}}%
\pgfpathlineto{\pgfqpoint{0.644008in}{0.418053in}}%
\pgfpathlineto{\pgfqpoint{0.636879in}{0.418407in}}%
\pgfpathlineto{\pgfqpoint{0.638459in}{0.416408in}}%
\pgfpathlineto{\pgfqpoint{0.636161in}{0.415808in}}%
\pgfpathlineto{\pgfqpoint{0.635340in}{0.417833in}}%
\pgfpathlineto{\pgfqpoint{0.633246in}{0.416909in}}%
\pgfpathlineto{\pgfqpoint{0.629966in}{0.419033in}}%
\pgfpathlineto{\pgfqpoint{0.625891in}{0.422864in}}%
\pgfpathlineto{\pgfqpoint{0.621156in}{0.424136in}}%
\pgfpathlineto{\pgfqpoint{0.617033in}{0.423185in}}%
\pgfpathlineto{\pgfqpoint{0.614359in}{0.424662in}}%
\pgfpathlineto{\pgfqpoint{0.613000in}{0.429251in}}%
\pgfpathlineto{\pgfqpoint{0.614183in}{0.431739in}}%
\pgfpathlineto{\pgfqpoint{0.619670in}{0.431167in}}%
\pgfpathlineto{\pgfqpoint{0.626353in}{0.433684in}}%
\pgfpathlineto{\pgfqpoint{0.627719in}{0.431363in}}%
\pgfpathlineto{\pgfqpoint{0.630050in}{0.430192in}}%
\pgfpathlineto{\pgfqpoint{0.631614in}{0.430608in}}%
\pgfpathlineto{\pgfqpoint{0.636883in}{0.428698in}}%
\pgfpathlineto{\pgfqpoint{0.640000in}{0.425903in}}%
\pgfpathlineto{\pgfqpoint{0.638907in}{0.421869in}}%
\pgfpathlineto{\pgfqpoint{0.640826in}{0.421267in}}%
\pgfpathlineto{\pgfqpoint{0.643623in}{0.424929in}}%
\pgfpathlineto{\pgfqpoint{0.647952in}{0.423982in}}%
\pgfpathlineto{\pgfqpoint{0.649233in}{0.421491in}}%
\pgfpathlineto{\pgfqpoint{0.651625in}{0.421666in}}%
\pgfpathlineto{\pgfqpoint{0.667655in}{0.424721in}}%
\pgfpathlineto{\pgfqpoint{0.673676in}{0.424199in}}%
\pgfpathlineto{\pgfqpoint{0.676435in}{0.424381in}}%
\pgfpathlineto{\pgfqpoint{0.682895in}{0.421450in}}%
\pgfpathlineto{\pgfqpoint{0.686743in}{0.417731in}}%
\pgfpathlineto{\pgfqpoint{0.685490in}{0.416186in}}%
\pgfpathlineto{\pgfqpoint{0.685449in}{0.411697in}}%
\pgfpathlineto{\pgfqpoint{0.687707in}{0.412299in}}%
\pgfpathlineto{\pgfqpoint{0.691151in}{0.410326in}}%
\pgfpathlineto{\pgfqpoint{0.690257in}{0.408099in}}%
\pgfpathlineto{\pgfqpoint{0.691712in}{0.406391in}}%
\pgfpathlineto{\pgfqpoint{0.693589in}{0.407652in}}%
\pgfpathlineto{\pgfqpoint{0.692837in}{0.410463in}}%
\pgfpathlineto{\pgfqpoint{0.691256in}{0.412047in}}%
\pgfpathlineto{\pgfqpoint{0.692624in}{0.413404in}}%
\pgfpathlineto{\pgfqpoint{0.696548in}{0.415075in}}%
\pgfpathlineto{\pgfqpoint{0.699114in}{0.416794in}}%
\pgfpathlineto{\pgfqpoint{0.705119in}{0.417836in}}%
\pgfpathlineto{\pgfqpoint{0.708125in}{0.417265in}}%
\pgfpathlineto{\pgfqpoint{0.711215in}{0.417647in}}%
\pgfpathlineto{\pgfqpoint{0.714242in}{0.417166in}}%
\pgfpathlineto{\pgfqpoint{0.721904in}{0.414573in}}%
\pgfpathlineto{\pgfqpoint{0.721574in}{0.415801in}}%
\pgfpathlineto{\pgfqpoint{0.725765in}{0.415354in}}%
\pgfpathlineto{\pgfqpoint{0.727666in}{0.413737in}}%
\pgfpathlineto{\pgfqpoint{0.725147in}{0.414065in}}%
\pgfpathlineto{\pgfqpoint{0.723593in}{0.412237in}}%
\pgfpathlineto{\pgfqpoint{0.720711in}{0.414781in}}%
\pgfpathlineto{\pgfqpoint{0.719039in}{0.412825in}}%
\pgfpathlineto{\pgfqpoint{0.715284in}{0.416046in}}%
\pgfpathlineto{\pgfqpoint{0.713680in}{0.414200in}}%
\pgfpathlineto{\pgfqpoint{0.710015in}{0.417271in}}%
\pgfpathlineto{\pgfqpoint{0.706826in}{0.413694in}}%
\pgfpathlineto{\pgfqpoint{0.707923in}{0.412710in}}%
\pgfpathlineto{\pgfqpoint{0.701711in}{0.405421in}}%
\pgfpathlineto{\pgfqpoint{0.699223in}{0.404064in}}%
\pgfpathlineto{\pgfqpoint{0.702855in}{0.400929in}}%
\pgfpathlineto{\pgfqpoint{0.698256in}{0.395576in}}%
\pgfpathclose%
\pgfusepath{fill}%
\end{pgfscope}%
\begin{pgfscope}%
\pgfpathrectangle{\pgfqpoint{0.100000in}{0.100000in}}{\pgfqpoint{3.608454in}{2.310000in}}%
\pgfusepath{clip}%
\pgfsetbuttcap%
\pgfsetmiterjoin%
\definecolor{currentfill}{rgb}{0.000000,0.478431,0.760784}%
\pgfsetfillcolor{currentfill}%
\pgfsetlinewidth{0.000000pt}%
\definecolor{currentstroke}{rgb}{0.000000,0.000000,0.000000}%
\pgfsetstrokecolor{currentstroke}%
\pgfsetstrokeopacity{0.000000}%
\pgfsetdash{}{0pt}%
\pgfpathmoveto{\pgfqpoint{1.886398in}{1.921692in}}%
\pgfpathlineto{\pgfqpoint{1.885060in}{1.894397in}}%
\pgfpathlineto{\pgfqpoint{1.845659in}{1.896366in}}%
\pgfpathlineto{\pgfqpoint{1.802493in}{1.899008in}}%
\pgfpathlineto{\pgfqpoint{1.796957in}{1.908802in}}%
\pgfpathlineto{\pgfqpoint{1.796911in}{1.916921in}}%
\pgfpathlineto{\pgfqpoint{1.800710in}{1.922898in}}%
\pgfpathlineto{\pgfqpoint{1.801655in}{1.932573in}}%
\pgfpathlineto{\pgfqpoint{1.799830in}{1.937828in}}%
\pgfpathlineto{\pgfqpoint{1.802847in}{1.946100in}}%
\pgfpathlineto{\pgfqpoint{1.801850in}{1.950975in}}%
\pgfpathlineto{\pgfqpoint{1.796923in}{1.954510in}}%
\pgfpathlineto{\pgfqpoint{1.828664in}{1.952547in}}%
\pgfpathlineto{\pgfqpoint{1.863197in}{1.950599in}}%
\pgfpathlineto{\pgfqpoint{1.885796in}{1.949443in}}%
\pgfpathlineto{\pgfqpoint{1.884519in}{1.921788in}}%
\pgfpathlineto{\pgfqpoint{1.886398in}{1.921692in}}%
\pgfpathclose%
\pgfusepath{fill}%
\end{pgfscope}%
\begin{pgfscope}%
\pgfpathrectangle{\pgfqpoint{0.100000in}{0.100000in}}{\pgfqpoint{3.608454in}{2.310000in}}%
\pgfusepath{clip}%
\pgfsetbuttcap%
\pgfsetmiterjoin%
\definecolor{currentfill}{rgb}{0.000000,0.356863,0.821569}%
\pgfsetfillcolor{currentfill}%
\pgfsetlinewidth{0.000000pt}%
\definecolor{currentstroke}{rgb}{0.000000,0.000000,0.000000}%
\pgfsetstrokecolor{currentstroke}%
\pgfsetstrokeopacity{0.000000}%
\pgfsetdash{}{0pt}%
\pgfpathmoveto{\pgfqpoint{1.790983in}{1.190359in}}%
\pgfpathlineto{\pgfqpoint{1.788727in}{1.150707in}}%
\pgfpathlineto{\pgfqpoint{1.728388in}{1.154348in}}%
\pgfpathlineto{\pgfqpoint{1.731511in}{1.193664in}}%
\pgfpathlineto{\pgfqpoint{1.723873in}{1.194151in}}%
\pgfpathlineto{\pgfqpoint{1.725856in}{1.224921in}}%
\pgfpathlineto{\pgfqpoint{1.751890in}{1.223219in}}%
\pgfpathlineto{\pgfqpoint{1.752303in}{1.230115in}}%
\pgfpathlineto{\pgfqpoint{1.786488in}{1.228084in}}%
\pgfpathlineto{\pgfqpoint{1.787082in}{1.221084in}}%
\pgfpathlineto{\pgfqpoint{1.785504in}{1.190712in}}%
\pgfpathlineto{\pgfqpoint{1.790983in}{1.190359in}}%
\pgfpathclose%
\pgfusepath{fill}%
\end{pgfscope}%
\begin{pgfscope}%
\pgfpathrectangle{\pgfqpoint{0.100000in}{0.100000in}}{\pgfqpoint{3.608454in}{2.310000in}}%
\pgfusepath{clip}%
\pgfsetbuttcap%
\pgfsetmiterjoin%
\definecolor{currentfill}{rgb}{0.000000,0.525490,0.737255}%
\pgfsetfillcolor{currentfill}%
\pgfsetlinewidth{0.000000pt}%
\definecolor{currentstroke}{rgb}{0.000000,0.000000,0.000000}%
\pgfsetstrokecolor{currentstroke}%
\pgfsetstrokeopacity{0.000000}%
\pgfsetdash{}{0pt}%
\pgfpathmoveto{\pgfqpoint{2.084392in}{1.676317in}}%
\pgfpathlineto{\pgfqpoint{2.057034in}{1.676704in}}%
\pgfpathlineto{\pgfqpoint{2.057273in}{1.695912in}}%
\pgfpathlineto{\pgfqpoint{2.046244in}{1.696054in}}%
\pgfpathlineto{\pgfqpoint{2.046612in}{1.723715in}}%
\pgfpathlineto{\pgfqpoint{2.046268in}{1.751283in}}%
\pgfpathlineto{\pgfqpoint{2.073025in}{1.750974in}}%
\pgfpathlineto{\pgfqpoint{2.100603in}{1.750815in}}%
\pgfpathlineto{\pgfqpoint{2.101083in}{1.743915in}}%
\pgfpathlineto{\pgfqpoint{2.114776in}{1.743905in}}%
\pgfpathlineto{\pgfqpoint{2.115120in}{1.723283in}}%
\pgfpathlineto{\pgfqpoint{2.115164in}{1.695734in}}%
\pgfpathlineto{\pgfqpoint{2.111687in}{1.695719in}}%
\pgfpathlineto{\pgfqpoint{2.084490in}{1.695746in}}%
\pgfpathlineto{\pgfqpoint{2.084392in}{1.676317in}}%
\pgfpathclose%
\pgfusepath{fill}%
\end{pgfscope}%
\begin{pgfscope}%
\pgfpathrectangle{\pgfqpoint{0.100000in}{0.100000in}}{\pgfqpoint{3.608454in}{2.310000in}}%
\pgfusepath{clip}%
\pgfsetbuttcap%
\pgfsetmiterjoin%
\definecolor{currentfill}{rgb}{0.000000,0.427451,0.786275}%
\pgfsetfillcolor{currentfill}%
\pgfsetlinewidth{0.000000pt}%
\definecolor{currentstroke}{rgb}{0.000000,0.000000,0.000000}%
\pgfsetstrokecolor{currentstroke}%
\pgfsetstrokeopacity{0.000000}%
\pgfsetdash{}{0pt}%
\pgfpathmoveto{\pgfqpoint{2.138550in}{1.463912in}}%
\pgfpathlineto{\pgfqpoint{2.128892in}{1.463893in}}%
\pgfpathlineto{\pgfqpoint{2.060821in}{1.465149in}}%
\pgfpathlineto{\pgfqpoint{2.059082in}{1.471267in}}%
\pgfpathlineto{\pgfqpoint{2.053550in}{1.477162in}}%
\pgfpathlineto{\pgfqpoint{2.056854in}{1.480886in}}%
\pgfpathlineto{\pgfqpoint{2.058174in}{1.490199in}}%
\pgfpathlineto{\pgfqpoint{2.083784in}{1.490019in}}%
\pgfpathlineto{\pgfqpoint{2.083971in}{1.510379in}}%
\pgfpathlineto{\pgfqpoint{2.097514in}{1.510367in}}%
\pgfpathlineto{\pgfqpoint{2.124622in}{1.510289in}}%
\pgfpathlineto{\pgfqpoint{2.138316in}{1.510307in}}%
\pgfpathlineto{\pgfqpoint{2.138550in}{1.463912in}}%
\pgfpathclose%
\pgfusepath{fill}%
\end{pgfscope}%
\begin{pgfscope}%
\pgfpathrectangle{\pgfqpoint{0.100000in}{0.100000in}}{\pgfqpoint{3.608454in}{2.310000in}}%
\pgfusepath{clip}%
\pgfsetbuttcap%
\pgfsetmiterjoin%
\definecolor{currentfill}{rgb}{0.000000,0.556863,0.721569}%
\pgfsetfillcolor{currentfill}%
\pgfsetlinewidth{0.000000pt}%
\definecolor{currentstroke}{rgb}{0.000000,0.000000,0.000000}%
\pgfsetstrokecolor{currentstroke}%
\pgfsetstrokeopacity{0.000000}%
\pgfsetdash{}{0pt}%
\pgfpathmoveto{\pgfqpoint{1.941489in}{0.354044in}}%
\pgfpathlineto{\pgfqpoint{1.939142in}{0.354101in}}%
\pgfpathlineto{\pgfqpoint{1.934663in}{0.372389in}}%
\pgfpathlineto{\pgfqpoint{1.935228in}{0.399177in}}%
\pgfpathlineto{\pgfqpoint{1.938012in}{0.416549in}}%
\pgfpathlineto{\pgfqpoint{1.956579in}{0.451354in}}%
\pgfpathlineto{\pgfqpoint{1.961353in}{0.455090in}}%
\pgfpathlineto{\pgfqpoint{1.966578in}{0.467186in}}%
\pgfpathlineto{\pgfqpoint{1.973872in}{0.476940in}}%
\pgfpathlineto{\pgfqpoint{1.976992in}{0.475338in}}%
\pgfpathlineto{\pgfqpoint{1.975277in}{0.471174in}}%
\pgfpathlineto{\pgfqpoint{1.962738in}{0.456890in}}%
\pgfpathlineto{\pgfqpoint{1.952264in}{0.440704in}}%
\pgfpathlineto{\pgfqpoint{1.945455in}{0.427440in}}%
\pgfpathlineto{\pgfqpoint{1.940167in}{0.414451in}}%
\pgfpathlineto{\pgfqpoint{1.937055in}{0.402500in}}%
\pgfpathlineto{\pgfqpoint{1.935766in}{0.383962in}}%
\pgfpathlineto{\pgfqpoint{1.938010in}{0.367087in}}%
\pgfpathlineto{\pgfqpoint{1.941489in}{0.354044in}}%
\pgfpathclose%
\pgfusepath{fill}%
\end{pgfscope}%
\begin{pgfscope}%
\pgfpathrectangle{\pgfqpoint{0.100000in}{0.100000in}}{\pgfqpoint{3.608454in}{2.310000in}}%
\pgfusepath{clip}%
\pgfsetbuttcap%
\pgfsetmiterjoin%
\definecolor{currentfill}{rgb}{0.000000,0.556863,0.721569}%
\pgfsetfillcolor{currentfill}%
\pgfsetlinewidth{0.000000pt}%
\definecolor{currentstroke}{rgb}{0.000000,0.000000,0.000000}%
\pgfsetstrokecolor{currentstroke}%
\pgfsetstrokeopacity{0.000000}%
\pgfsetdash{}{0pt}%
\pgfpathmoveto{\pgfqpoint{1.930477in}{0.354309in}}%
\pgfpathlineto{\pgfqpoint{1.899944in}{0.355045in}}%
\pgfpathlineto{\pgfqpoint{1.891060in}{0.356772in}}%
\pgfpathlineto{\pgfqpoint{1.891484in}{0.370110in}}%
\pgfpathlineto{\pgfqpoint{1.859982in}{0.371433in}}%
\pgfpathlineto{\pgfqpoint{1.861148in}{0.393352in}}%
\pgfpathlineto{\pgfqpoint{1.857651in}{0.393438in}}%
\pgfpathlineto{\pgfqpoint{1.856313in}{0.408047in}}%
\pgfpathlineto{\pgfqpoint{1.852427in}{0.418217in}}%
\pgfpathlineto{\pgfqpoint{1.834828in}{0.418520in}}%
\pgfpathlineto{\pgfqpoint{1.836840in}{0.475148in}}%
\pgfpathlineto{\pgfqpoint{1.902126in}{0.472699in}}%
\pgfpathlineto{\pgfqpoint{1.900788in}{0.477356in}}%
\pgfpathlineto{\pgfqpoint{1.907060in}{0.482191in}}%
\pgfpathlineto{\pgfqpoint{1.924683in}{0.477880in}}%
\pgfpathlineto{\pgfqpoint{1.938836in}{0.498314in}}%
\pgfpathlineto{\pgfqpoint{1.954355in}{0.511185in}}%
\pgfpathlineto{\pgfqpoint{1.962232in}{0.510712in}}%
\pgfpathlineto{\pgfqpoint{1.966594in}{0.506597in}}%
\pgfpathlineto{\pgfqpoint{1.973994in}{0.506774in}}%
\pgfpathlineto{\pgfqpoint{1.981043in}{0.498344in}}%
\pgfpathlineto{\pgfqpoint{1.979461in}{0.488391in}}%
\pgfpathlineto{\pgfqpoint{1.978605in}{0.483866in}}%
\pgfpathlineto{\pgfqpoint{1.969758in}{0.476284in}}%
\pgfpathlineto{\pgfqpoint{1.965121in}{0.476210in}}%
\pgfpathlineto{\pgfqpoint{1.961299in}{0.468390in}}%
\pgfpathlineto{\pgfqpoint{1.951003in}{0.452565in}}%
\pgfpathlineto{\pgfqpoint{1.945104in}{0.457357in}}%
\pgfpathlineto{\pgfqpoint{1.938160in}{0.454348in}}%
\pgfpathlineto{\pgfqpoint{1.938434in}{0.445675in}}%
\pgfpathlineto{\pgfqpoint{1.946055in}{0.442493in}}%
\pgfpathlineto{\pgfqpoint{1.941015in}{0.432395in}}%
\pgfpathlineto{\pgfqpoint{1.934658in}{0.413330in}}%
\pgfpathlineto{\pgfqpoint{1.931292in}{0.396174in}}%
\pgfpathlineto{\pgfqpoint{1.923085in}{0.385830in}}%
\pgfpathlineto{\pgfqpoint{1.922580in}{0.373704in}}%
\pgfpathlineto{\pgfqpoint{1.928475in}{0.367207in}}%
\pgfpathlineto{\pgfqpoint{1.930477in}{0.354309in}}%
\pgfpathclose%
\pgfusepath{fill}%
\end{pgfscope}%
\begin{pgfscope}%
\pgfpathrectangle{\pgfqpoint{0.100000in}{0.100000in}}{\pgfqpoint{3.608454in}{2.310000in}}%
\pgfusepath{clip}%
\pgfsetbuttcap%
\pgfsetmiterjoin%
\definecolor{currentfill}{rgb}{0.000000,0.470588,0.764706}%
\pgfsetfillcolor{currentfill}%
\pgfsetlinewidth{0.000000pt}%
\definecolor{currentstroke}{rgb}{0.000000,0.000000,0.000000}%
\pgfsetstrokecolor{currentstroke}%
\pgfsetstrokeopacity{0.000000}%
\pgfsetdash{}{0pt}%
\pgfpathmoveto{\pgfqpoint{2.023670in}{1.724067in}}%
\pgfpathlineto{\pgfqpoint{1.998767in}{1.724416in}}%
\pgfpathlineto{\pgfqpoint{1.999302in}{1.751993in}}%
\pgfpathlineto{\pgfqpoint{1.944686in}{1.753364in}}%
\pgfpathlineto{\pgfqpoint{1.944449in}{1.753368in}}%
\pgfpathlineto{\pgfqpoint{1.945343in}{1.781068in}}%
\pgfpathlineto{\pgfqpoint{1.986361in}{1.779846in}}%
\pgfpathlineto{\pgfqpoint{2.024621in}{1.779133in}}%
\pgfpathlineto{\pgfqpoint{2.023670in}{1.724067in}}%
\pgfpathclose%
\pgfusepath{fill}%
\end{pgfscope}%
\begin{pgfscope}%
\pgfpathrectangle{\pgfqpoint{0.100000in}{0.100000in}}{\pgfqpoint{3.608454in}{2.310000in}}%
\pgfusepath{clip}%
\pgfsetbuttcap%
\pgfsetmiterjoin%
\definecolor{currentfill}{rgb}{0.000000,0.862745,0.568627}%
\pgfsetfillcolor{currentfill}%
\pgfsetlinewidth{0.000000pt}%
\definecolor{currentstroke}{rgb}{0.000000,0.000000,0.000000}%
\pgfsetstrokecolor{currentstroke}%
\pgfsetstrokeopacity{0.000000}%
\pgfsetdash{}{0pt}%
\pgfpathmoveto{\pgfqpoint{2.313347in}{1.069167in}}%
\pgfpathlineto{\pgfqpoint{2.313444in}{1.062948in}}%
\pgfpathlineto{\pgfqpoint{2.310094in}{1.055119in}}%
\pgfpathlineto{\pgfqpoint{2.293330in}{1.054914in}}%
\pgfpathlineto{\pgfqpoint{2.269431in}{1.054652in}}%
\pgfpathlineto{\pgfqpoint{2.269359in}{1.061552in}}%
\pgfpathlineto{\pgfqpoint{2.245419in}{1.061527in}}%
\pgfpathlineto{\pgfqpoint{2.247654in}{1.068396in}}%
\pgfpathlineto{\pgfqpoint{2.247558in}{1.082128in}}%
\pgfpathlineto{\pgfqpoint{2.259040in}{1.082097in}}%
\pgfpathlineto{\pgfqpoint{2.258979in}{1.087834in}}%
\pgfpathlineto{\pgfqpoint{2.272664in}{1.088030in}}%
\pgfpathlineto{\pgfqpoint{2.272396in}{1.102853in}}%
\pgfpathlineto{\pgfqpoint{2.279258in}{1.102907in}}%
\pgfpathlineto{\pgfqpoint{2.279199in}{1.109799in}}%
\pgfpathlineto{\pgfqpoint{2.286066in}{1.109865in}}%
\pgfpathlineto{\pgfqpoint{2.286012in}{1.115699in}}%
\pgfpathlineto{\pgfqpoint{2.295109in}{1.108191in}}%
\pgfpathlineto{\pgfqpoint{2.293303in}{1.102454in}}%
\pgfpathlineto{\pgfqpoint{2.302022in}{1.100303in}}%
\pgfpathlineto{\pgfqpoint{2.308868in}{1.093507in}}%
\pgfpathlineto{\pgfqpoint{2.307635in}{1.090812in}}%
\pgfpathlineto{\pgfqpoint{2.313117in}{1.083232in}}%
\pgfpathlineto{\pgfqpoint{2.313347in}{1.069167in}}%
\pgfpathclose%
\pgfusepath{fill}%
\end{pgfscope}%
\begin{pgfscope}%
\pgfpathrectangle{\pgfqpoint{0.100000in}{0.100000in}}{\pgfqpoint{3.608454in}{2.310000in}}%
\pgfusepath{clip}%
\pgfsetbuttcap%
\pgfsetmiterjoin%
\definecolor{currentfill}{rgb}{0.000000,0.568627,0.715686}%
\pgfsetfillcolor{currentfill}%
\pgfsetlinewidth{0.000000pt}%
\definecolor{currentstroke}{rgb}{0.000000,0.000000,0.000000}%
\pgfsetstrokecolor{currentstroke}%
\pgfsetstrokeopacity{0.000000}%
\pgfsetdash{}{0pt}%
\pgfpathmoveto{\pgfqpoint{2.958909in}{1.001900in}}%
\pgfpathlineto{\pgfqpoint{2.949644in}{0.997974in}}%
\pgfpathlineto{\pgfqpoint{2.953478in}{0.990526in}}%
\pgfpathlineto{\pgfqpoint{2.950161in}{0.981612in}}%
\pgfpathlineto{\pgfqpoint{2.956878in}{0.975335in}}%
\pgfpathlineto{\pgfqpoint{2.954883in}{0.971559in}}%
\pgfpathlineto{\pgfqpoint{2.948243in}{0.978295in}}%
\pgfpathlineto{\pgfqpoint{2.941645in}{0.986529in}}%
\pgfpathlineto{\pgfqpoint{2.922696in}{0.996327in}}%
\pgfpathlineto{\pgfqpoint{2.919954in}{1.001789in}}%
\pgfpathlineto{\pgfqpoint{2.910642in}{1.010268in}}%
\pgfpathlineto{\pgfqpoint{2.907921in}{1.017944in}}%
\pgfpathlineto{\pgfqpoint{2.895066in}{1.035008in}}%
\pgfpathlineto{\pgfqpoint{2.888625in}{1.033073in}}%
\pgfpathlineto{\pgfqpoint{2.884093in}{1.034847in}}%
\pgfpathlineto{\pgfqpoint{2.876324in}{1.042276in}}%
\pgfpathlineto{\pgfqpoint{2.872404in}{1.042068in}}%
\pgfpathlineto{\pgfqpoint{2.864794in}{1.047108in}}%
\pgfpathlineto{\pgfqpoint{2.863588in}{1.049271in}}%
\pgfpathlineto{\pgfqpoint{2.864683in}{1.055271in}}%
\pgfpathlineto{\pgfqpoint{2.869918in}{1.063176in}}%
\pgfpathlineto{\pgfqpoint{2.876200in}{1.069013in}}%
\pgfpathlineto{\pgfqpoint{2.876403in}{1.073700in}}%
\pgfpathlineto{\pgfqpoint{2.882548in}{1.076650in}}%
\pgfpathlineto{\pgfqpoint{2.903153in}{1.086612in}}%
\pgfpathlineto{\pgfqpoint{2.922867in}{1.095204in}}%
\pgfpathlineto{\pgfqpoint{2.931616in}{1.096854in}}%
\pgfpathlineto{\pgfqpoint{2.934795in}{1.069673in}}%
\pgfpathlineto{\pgfqpoint{2.951248in}{1.056924in}}%
\pgfpathlineto{\pgfqpoint{2.957922in}{1.053177in}}%
\pgfpathlineto{\pgfqpoint{2.975999in}{1.050234in}}%
\pgfpathlineto{\pgfqpoint{2.968957in}{1.035717in}}%
\pgfpathlineto{\pgfqpoint{2.962380in}{1.029555in}}%
\pgfpathlineto{\pgfqpoint{2.961263in}{1.022919in}}%
\pgfpathlineto{\pgfqpoint{2.963991in}{1.017466in}}%
\pgfpathlineto{\pgfqpoint{2.958909in}{1.001900in}}%
\pgfpathclose%
\pgfusepath{fill}%
\end{pgfscope}%
\begin{pgfscope}%
\pgfpathrectangle{\pgfqpoint{0.100000in}{0.100000in}}{\pgfqpoint{3.608454in}{2.310000in}}%
\pgfusepath{clip}%
\pgfsetbuttcap%
\pgfsetmiterjoin%
\definecolor{currentfill}{rgb}{0.000000,0.407843,0.796078}%
\pgfsetfillcolor{currentfill}%
\pgfsetlinewidth{0.000000pt}%
\definecolor{currentstroke}{rgb}{0.000000,0.000000,0.000000}%
\pgfsetstrokecolor{currentstroke}%
\pgfsetstrokeopacity{0.000000}%
\pgfsetdash{}{0pt}%
\pgfpathmoveto{\pgfqpoint{2.684273in}{1.275179in}}%
\pgfpathlineto{\pgfqpoint{2.678952in}{1.278333in}}%
\pgfpathlineto{\pgfqpoint{2.675593in}{1.285078in}}%
\pgfpathlineto{\pgfqpoint{2.669652in}{1.289191in}}%
\pgfpathlineto{\pgfqpoint{2.655765in}{1.288068in}}%
\pgfpathlineto{\pgfqpoint{2.649917in}{1.290900in}}%
\pgfpathlineto{\pgfqpoint{2.648711in}{1.297118in}}%
\pgfpathlineto{\pgfqpoint{2.643610in}{1.301513in}}%
\pgfpathlineto{\pgfqpoint{2.640611in}{1.294225in}}%
\pgfpathlineto{\pgfqpoint{2.636960in}{1.295013in}}%
\pgfpathlineto{\pgfqpoint{2.636434in}{1.301852in}}%
\pgfpathlineto{\pgfqpoint{2.629550in}{1.301264in}}%
\pgfpathlineto{\pgfqpoint{2.629213in}{1.305901in}}%
\pgfpathlineto{\pgfqpoint{2.622478in}{1.305021in}}%
\pgfpathlineto{\pgfqpoint{2.621408in}{1.315375in}}%
\pgfpathlineto{\pgfqpoint{2.644392in}{1.317477in}}%
\pgfpathlineto{\pgfqpoint{2.641760in}{1.344149in}}%
\pgfpathlineto{\pgfqpoint{2.643662in}{1.346766in}}%
\pgfpathlineto{\pgfqpoint{2.662944in}{1.348688in}}%
\pgfpathlineto{\pgfqpoint{2.667619in}{1.346789in}}%
\pgfpathlineto{\pgfqpoint{2.669050in}{1.334303in}}%
\pgfpathlineto{\pgfqpoint{2.674475in}{1.337189in}}%
\pgfpathlineto{\pgfqpoint{2.688201in}{1.338715in}}%
\pgfpathlineto{\pgfqpoint{2.697081in}{1.338065in}}%
\pgfpathlineto{\pgfqpoint{2.697292in}{1.333124in}}%
\pgfpathlineto{\pgfqpoint{2.707002in}{1.331992in}}%
\pgfpathlineto{\pgfqpoint{2.713052in}{1.339727in}}%
\pgfpathlineto{\pgfqpoint{2.718732in}{1.341235in}}%
\pgfpathlineto{\pgfqpoint{2.722897in}{1.338929in}}%
\pgfpathlineto{\pgfqpoint{2.724859in}{1.333241in}}%
\pgfpathlineto{\pgfqpoint{2.730992in}{1.324932in}}%
\pgfpathlineto{\pgfqpoint{2.726160in}{1.316456in}}%
\pgfpathlineto{\pgfqpoint{2.726400in}{1.307715in}}%
\pgfpathlineto{\pgfqpoint{2.724571in}{1.300390in}}%
\pgfpathlineto{\pgfqpoint{2.718323in}{1.291481in}}%
\pgfpathlineto{\pgfqpoint{2.704197in}{1.287428in}}%
\pgfpathlineto{\pgfqpoint{2.698255in}{1.290856in}}%
\pgfpathlineto{\pgfqpoint{2.690035in}{1.277558in}}%
\pgfpathlineto{\pgfqpoint{2.684273in}{1.275179in}}%
\pgfpathclose%
\pgfusepath{fill}%
\end{pgfscope}%
\begin{pgfscope}%
\pgfpathrectangle{\pgfqpoint{0.100000in}{0.100000in}}{\pgfqpoint{3.608454in}{2.310000in}}%
\pgfusepath{clip}%
\pgfsetbuttcap%
\pgfsetmiterjoin%
\definecolor{currentfill}{rgb}{0.000000,0.368627,0.815686}%
\pgfsetfillcolor{currentfill}%
\pgfsetlinewidth{0.000000pt}%
\definecolor{currentstroke}{rgb}{0.000000,0.000000,0.000000}%
\pgfsetstrokecolor{currentstroke}%
\pgfsetstrokeopacity{0.000000}%
\pgfsetdash{}{0pt}%
\pgfpathmoveto{\pgfqpoint{1.912766in}{1.573654in}}%
\pgfpathlineto{\pgfqpoint{1.912258in}{1.559919in}}%
\pgfpathlineto{\pgfqpoint{1.884919in}{1.560911in}}%
\pgfpathlineto{\pgfqpoint{1.885929in}{1.588436in}}%
\pgfpathlineto{\pgfqpoint{1.858155in}{1.589614in}}%
\pgfpathlineto{\pgfqpoint{1.860373in}{1.661671in}}%
\pgfpathlineto{\pgfqpoint{1.904126in}{1.659838in}}%
\pgfpathlineto{\pgfqpoint{1.905777in}{1.655731in}}%
\pgfpathlineto{\pgfqpoint{1.914725in}{1.650247in}}%
\pgfpathlineto{\pgfqpoint{1.914002in}{1.614954in}}%
\pgfpathlineto{\pgfqpoint{1.912766in}{1.573654in}}%
\pgfpathclose%
\pgfusepath{fill}%
\end{pgfscope}%
\begin{pgfscope}%
\pgfpathrectangle{\pgfqpoint{0.100000in}{0.100000in}}{\pgfqpoint{3.608454in}{2.310000in}}%
\pgfusepath{clip}%
\pgfsetbuttcap%
\pgfsetmiterjoin%
\definecolor{currentfill}{rgb}{0.000000,0.458824,0.770588}%
\pgfsetfillcolor{currentfill}%
\pgfsetlinewidth{0.000000pt}%
\definecolor{currentstroke}{rgb}{0.000000,0.000000,0.000000}%
\pgfsetstrokecolor{currentstroke}%
\pgfsetstrokeopacity{0.000000}%
\pgfsetdash{}{0pt}%
\pgfpathmoveto{\pgfqpoint{2.719775in}{1.481004in}}%
\pgfpathlineto{\pgfqpoint{2.720165in}{1.477649in}}%
\pgfpathlineto{\pgfqpoint{2.695336in}{1.474638in}}%
\pgfpathlineto{\pgfqpoint{2.694173in}{1.495218in}}%
\pgfpathlineto{\pgfqpoint{2.679497in}{1.493619in}}%
\pgfpathlineto{\pgfqpoint{2.678668in}{1.500441in}}%
\pgfpathlineto{\pgfqpoint{2.667350in}{1.499209in}}%
\pgfpathlineto{\pgfqpoint{2.664197in}{1.526649in}}%
\pgfpathlineto{\pgfqpoint{2.661863in}{1.526375in}}%
\pgfpathlineto{\pgfqpoint{2.660236in}{1.540245in}}%
\pgfpathlineto{\pgfqpoint{2.661344in}{1.549614in}}%
\pgfpathlineto{\pgfqpoint{2.668362in}{1.548067in}}%
\pgfpathlineto{\pgfqpoint{2.688778in}{1.550091in}}%
\pgfpathlineto{\pgfqpoint{2.684290in}{1.589030in}}%
\pgfpathlineto{\pgfqpoint{2.706079in}{1.591455in}}%
\pgfpathlineto{\pgfqpoint{2.707766in}{1.586556in}}%
\pgfpathlineto{\pgfqpoint{2.714189in}{1.531011in}}%
\pgfpathlineto{\pgfqpoint{2.719775in}{1.481004in}}%
\pgfpathclose%
\pgfusepath{fill}%
\end{pgfscope}%
\begin{pgfscope}%
\pgfpathrectangle{\pgfqpoint{0.100000in}{0.100000in}}{\pgfqpoint{3.608454in}{2.310000in}}%
\pgfusepath{clip}%
\pgfsetbuttcap%
\pgfsetmiterjoin%
\definecolor{currentfill}{rgb}{0.000000,1.000000,0.500000}%
\pgfsetfillcolor{currentfill}%
\pgfsetlinewidth{0.000000pt}%
\definecolor{currentstroke}{rgb}{0.000000,0.000000,0.000000}%
\pgfsetstrokecolor{currentstroke}%
\pgfsetstrokeopacity{0.000000}%
\pgfsetdash{}{0pt}%
\pgfpathmoveto{\pgfqpoint{2.837048in}{1.201444in}}%
\pgfpathlineto{\pgfqpoint{2.833845in}{1.206501in}}%
\pgfpathlineto{\pgfqpoint{2.833882in}{1.215723in}}%
\pgfpathlineto{\pgfqpoint{2.831285in}{1.222519in}}%
\pgfpathlineto{\pgfqpoint{2.826345in}{1.240166in}}%
\pgfpathlineto{\pgfqpoint{2.826981in}{1.247494in}}%
\pgfpathlineto{\pgfqpoint{2.825042in}{1.252479in}}%
\pgfpathlineto{\pgfqpoint{2.830889in}{1.256941in}}%
\pgfpathlineto{\pgfqpoint{2.833788in}{1.254755in}}%
\pgfpathlineto{\pgfqpoint{2.841576in}{1.260538in}}%
\pgfpathlineto{\pgfqpoint{2.850587in}{1.262241in}}%
\pgfpathlineto{\pgfqpoint{2.851560in}{1.271717in}}%
\pgfpathlineto{\pgfqpoint{2.860524in}{1.271364in}}%
\pgfpathlineto{\pgfqpoint{2.872542in}{1.266952in}}%
\pgfpathlineto{\pgfqpoint{2.874724in}{1.257922in}}%
\pgfpathlineto{\pgfqpoint{2.880671in}{1.252069in}}%
\pgfpathlineto{\pgfqpoint{2.887645in}{1.250541in}}%
\pgfpathlineto{\pgfqpoint{2.878670in}{1.243251in}}%
\pgfpathlineto{\pgfqpoint{2.878218in}{1.237030in}}%
\pgfpathlineto{\pgfqpoint{2.870950in}{1.230065in}}%
\pgfpathlineto{\pgfqpoint{2.871264in}{1.223809in}}%
\pgfpathlineto{\pgfqpoint{2.859365in}{1.219154in}}%
\pgfpathlineto{\pgfqpoint{2.856686in}{1.210206in}}%
\pgfpathlineto{\pgfqpoint{2.837048in}{1.201444in}}%
\pgfpathclose%
\pgfusepath{fill}%
\end{pgfscope}%
\begin{pgfscope}%
\pgfpathrectangle{\pgfqpoint{0.100000in}{0.100000in}}{\pgfqpoint{3.608454in}{2.310000in}}%
\pgfusepath{clip}%
\pgfsetbuttcap%
\pgfsetmiterjoin%
\definecolor{currentfill}{rgb}{0.000000,0.564706,0.717647}%
\pgfsetfillcolor{currentfill}%
\pgfsetlinewidth{0.000000pt}%
\definecolor{currentstroke}{rgb}{0.000000,0.000000,0.000000}%
\pgfsetstrokecolor{currentstroke}%
\pgfsetstrokeopacity{0.000000}%
\pgfsetdash{}{0pt}%
\pgfpathmoveto{\pgfqpoint{2.989806in}{1.055375in}}%
\pgfpathlineto{\pgfqpoint{2.990720in}{1.049995in}}%
\pgfpathlineto{\pgfqpoint{2.982517in}{1.046654in}}%
\pgfpathlineto{\pgfqpoint{2.975999in}{1.050234in}}%
\pgfpathlineto{\pgfqpoint{2.957922in}{1.053177in}}%
\pgfpathlineto{\pgfqpoint{2.951248in}{1.056924in}}%
\pgfpathlineto{\pgfqpoint{2.934795in}{1.069673in}}%
\pgfpathlineto{\pgfqpoint{2.931616in}{1.096854in}}%
\pgfpathlineto{\pgfqpoint{2.922867in}{1.095204in}}%
\pgfpathlineto{\pgfqpoint{2.922309in}{1.102781in}}%
\pgfpathlineto{\pgfqpoint{2.926537in}{1.111891in}}%
\pgfpathlineto{\pgfqpoint{2.932868in}{1.114078in}}%
\pgfpathlineto{\pgfqpoint{2.946773in}{1.103937in}}%
\pgfpathlineto{\pgfqpoint{2.947537in}{1.098430in}}%
\pgfpathlineto{\pgfqpoint{2.986399in}{1.102377in}}%
\pgfpathlineto{\pgfqpoint{2.985051in}{1.092680in}}%
\pgfpathlineto{\pgfqpoint{2.980220in}{1.091029in}}%
\pgfpathlineto{\pgfqpoint{2.984590in}{1.076632in}}%
\pgfpathlineto{\pgfqpoint{2.986559in}{1.063157in}}%
\pgfpathlineto{\pgfqpoint{2.989806in}{1.055375in}}%
\pgfpathclose%
\pgfusepath{fill}%
\end{pgfscope}%
\begin{pgfscope}%
\pgfpathrectangle{\pgfqpoint{0.100000in}{0.100000in}}{\pgfqpoint{3.608454in}{2.310000in}}%
\pgfusepath{clip}%
\pgfsetbuttcap%
\pgfsetmiterjoin%
\definecolor{currentfill}{rgb}{0.000000,0.827451,0.586275}%
\pgfsetfillcolor{currentfill}%
\pgfsetlinewidth{0.000000pt}%
\definecolor{currentstroke}{rgb}{0.000000,0.000000,0.000000}%
\pgfsetstrokecolor{currentstroke}%
\pgfsetstrokeopacity{0.000000}%
\pgfsetdash{}{0pt}%
\pgfpathmoveto{\pgfqpoint{1.037345in}{0.183729in}}%
\pgfpathlineto{\pgfqpoint{1.037725in}{0.186685in}}%
\pgfpathlineto{\pgfqpoint{1.039756in}{0.184678in}}%
\pgfpathlineto{\pgfqpoint{1.037345in}{0.183729in}}%
\pgfpathclose%
\pgfusepath{fill}%
\end{pgfscope}%
\begin{pgfscope}%
\pgfpathrectangle{\pgfqpoint{0.100000in}{0.100000in}}{\pgfqpoint{3.608454in}{2.310000in}}%
\pgfusepath{clip}%
\pgfsetbuttcap%
\pgfsetmiterjoin%
\definecolor{currentfill}{rgb}{0.000000,0.827451,0.586275}%
\pgfsetfillcolor{currentfill}%
\pgfsetlinewidth{0.000000pt}%
\definecolor{currentstroke}{rgb}{0.000000,0.000000,0.000000}%
\pgfsetstrokecolor{currentstroke}%
\pgfsetstrokeopacity{0.000000}%
\pgfsetdash{}{0pt}%
\pgfpathmoveto{\pgfqpoint{1.035979in}{0.195213in}}%
\pgfpathlineto{\pgfqpoint{1.035168in}{0.197283in}}%
\pgfpathlineto{\pgfqpoint{1.034662in}{0.206903in}}%
\pgfpathlineto{\pgfqpoint{1.035313in}{0.207991in}}%
\pgfpathlineto{\pgfqpoint{1.034954in}{0.212110in}}%
\pgfpathlineto{\pgfqpoint{1.035841in}{0.213389in}}%
\pgfpathlineto{\pgfqpoint{1.034185in}{0.215042in}}%
\pgfpathlineto{\pgfqpoint{1.034010in}{0.217336in}}%
\pgfpathlineto{\pgfqpoint{1.037398in}{0.218597in}}%
\pgfpathlineto{\pgfqpoint{1.038802in}{0.221050in}}%
\pgfpathlineto{\pgfqpoint{1.037009in}{0.222070in}}%
\pgfpathlineto{\pgfqpoint{1.036796in}{0.223966in}}%
\pgfpathlineto{\pgfqpoint{1.038521in}{0.226447in}}%
\pgfpathlineto{\pgfqpoint{1.036445in}{0.231584in}}%
\pgfpathlineto{\pgfqpoint{1.041639in}{0.235717in}}%
\pgfpathlineto{\pgfqpoint{1.043821in}{0.231654in}}%
\pgfpathlineto{\pgfqpoint{1.048322in}{0.228489in}}%
\pgfpathlineto{\pgfqpoint{1.045161in}{0.227298in}}%
\pgfpathlineto{\pgfqpoint{1.046187in}{0.224668in}}%
\pgfpathlineto{\pgfqpoint{1.045919in}{0.221757in}}%
\pgfpathlineto{\pgfqpoint{1.044595in}{0.217845in}}%
\pgfpathlineto{\pgfqpoint{1.044114in}{0.214196in}}%
\pgfpathlineto{\pgfqpoint{1.042540in}{0.210865in}}%
\pgfpathlineto{\pgfqpoint{1.043020in}{0.209624in}}%
\pgfpathlineto{\pgfqpoint{1.041842in}{0.206198in}}%
\pgfpathlineto{\pgfqpoint{1.039406in}{0.205292in}}%
\pgfpathlineto{\pgfqpoint{1.040077in}{0.203620in}}%
\pgfpathlineto{\pgfqpoint{1.037635in}{0.200811in}}%
\pgfpathlineto{\pgfqpoint{1.035979in}{0.195213in}}%
\pgfpathclose%
\pgfusepath{fill}%
\end{pgfscope}%
\begin{pgfscope}%
\pgfpathrectangle{\pgfqpoint{0.100000in}{0.100000in}}{\pgfqpoint{3.608454in}{2.310000in}}%
\pgfusepath{clip}%
\pgfsetbuttcap%
\pgfsetmiterjoin%
\definecolor{currentfill}{rgb}{0.000000,0.827451,0.586275}%
\pgfsetfillcolor{currentfill}%
\pgfsetlinewidth{0.000000pt}%
\definecolor{currentstroke}{rgb}{0.000000,0.000000,0.000000}%
\pgfsetstrokecolor{currentstroke}%
\pgfsetstrokeopacity{0.000000}%
\pgfsetdash{}{0pt}%
\pgfpathmoveto{\pgfqpoint{1.029215in}{0.224020in}}%
\pgfpathlineto{\pgfqpoint{1.032249in}{0.226756in}}%
\pgfpathlineto{\pgfqpoint{1.032356in}{0.230382in}}%
\pgfpathlineto{\pgfqpoint{1.033329in}{0.232191in}}%
\pgfpathlineto{\pgfqpoint{1.035308in}{0.231235in}}%
\pgfpathlineto{\pgfqpoint{1.034682in}{0.226435in}}%
\pgfpathlineto{\pgfqpoint{1.033872in}{0.224494in}}%
\pgfpathlineto{\pgfqpoint{1.031932in}{0.223434in}}%
\pgfpathlineto{\pgfqpoint{1.029215in}{0.224020in}}%
\pgfpathclose%
\pgfusepath{fill}%
\end{pgfscope}%
\begin{pgfscope}%
\pgfpathrectangle{\pgfqpoint{0.100000in}{0.100000in}}{\pgfqpoint{3.608454in}{2.310000in}}%
\pgfusepath{clip}%
\pgfsetbuttcap%
\pgfsetmiterjoin%
\definecolor{currentfill}{rgb}{0.000000,0.827451,0.586275}%
\pgfsetfillcolor{currentfill}%
\pgfsetlinewidth{0.000000pt}%
\definecolor{currentstroke}{rgb}{0.000000,0.000000,0.000000}%
\pgfsetstrokecolor{currentstroke}%
\pgfsetstrokeopacity{0.000000}%
\pgfsetdash{}{0pt}%
\pgfpathmoveto{\pgfqpoint{1.031898in}{0.249479in}}%
\pgfpathlineto{\pgfqpoint{1.037351in}{0.246961in}}%
\pgfpathlineto{\pgfqpoint{1.041099in}{0.250574in}}%
\pgfpathlineto{\pgfqpoint{1.042604in}{0.247304in}}%
\pgfpathlineto{\pgfqpoint{1.044554in}{0.244818in}}%
\pgfpathlineto{\pgfqpoint{1.047084in}{0.239857in}}%
\pgfpathlineto{\pgfqpoint{1.050384in}{0.238572in}}%
\pgfpathlineto{\pgfqpoint{1.051184in}{0.237535in}}%
\pgfpathlineto{\pgfqpoint{1.049176in}{0.229939in}}%
\pgfpathlineto{\pgfqpoint{1.046352in}{0.230420in}}%
\pgfpathlineto{\pgfqpoint{1.040103in}{0.241756in}}%
\pgfpathlineto{\pgfqpoint{1.039585in}{0.239633in}}%
\pgfpathlineto{\pgfqpoint{1.040523in}{0.237630in}}%
\pgfpathlineto{\pgfqpoint{1.039041in}{0.234732in}}%
\pgfpathlineto{\pgfqpoint{1.036002in}{0.232709in}}%
\pgfpathlineto{\pgfqpoint{1.034369in}{0.234138in}}%
\pgfpathlineto{\pgfqpoint{1.033903in}{0.237752in}}%
\pgfpathlineto{\pgfqpoint{1.034739in}{0.241664in}}%
\pgfpathlineto{\pgfqpoint{1.031898in}{0.249479in}}%
\pgfpathclose%
\pgfusepath{fill}%
\end{pgfscope}%
\begin{pgfscope}%
\pgfpathrectangle{\pgfqpoint{0.100000in}{0.100000in}}{\pgfqpoint{3.608454in}{2.310000in}}%
\pgfusepath{clip}%
\pgfsetbuttcap%
\pgfsetmiterjoin%
\definecolor{currentfill}{rgb}{0.000000,0.827451,0.586275}%
\pgfsetfillcolor{currentfill}%
\pgfsetlinewidth{0.000000pt}%
\definecolor{currentstroke}{rgb}{0.000000,0.000000,0.000000}%
\pgfsetstrokecolor{currentstroke}%
\pgfsetstrokeopacity{0.000000}%
\pgfsetdash{}{0pt}%
\pgfpathmoveto{\pgfqpoint{1.062273in}{0.185108in}}%
\pgfpathlineto{\pgfqpoint{1.060831in}{0.189361in}}%
\pgfpathlineto{\pgfqpoint{1.064048in}{0.191263in}}%
\pgfpathlineto{\pgfqpoint{1.066461in}{0.190126in}}%
\pgfpathlineto{\pgfqpoint{1.067412in}{0.187983in}}%
\pgfpathlineto{\pgfqpoint{1.065144in}{0.184845in}}%
\pgfpathlineto{\pgfqpoint{1.062273in}{0.185108in}}%
\pgfpathclose%
\pgfusepath{fill}%
\end{pgfscope}%
\begin{pgfscope}%
\pgfpathrectangle{\pgfqpoint{0.100000in}{0.100000in}}{\pgfqpoint{3.608454in}{2.310000in}}%
\pgfusepath{clip}%
\pgfsetbuttcap%
\pgfsetmiterjoin%
\definecolor{currentfill}{rgb}{0.000000,0.827451,0.586275}%
\pgfsetfillcolor{currentfill}%
\pgfsetlinewidth{0.000000pt}%
\definecolor{currentstroke}{rgb}{0.000000,0.000000,0.000000}%
\pgfsetstrokecolor{currentstroke}%
\pgfsetstrokeopacity{0.000000}%
\pgfsetdash{}{0pt}%
\pgfpathmoveto{\pgfqpoint{1.053060in}{0.199716in}}%
\pgfpathlineto{\pgfqpoint{1.049731in}{0.199089in}}%
\pgfpathlineto{\pgfqpoint{1.049501in}{0.194791in}}%
\pgfpathlineto{\pgfqpoint{1.047203in}{0.194053in}}%
\pgfpathlineto{\pgfqpoint{1.047508in}{0.191183in}}%
\pgfpathlineto{\pgfqpoint{1.045251in}{0.188644in}}%
\pgfpathlineto{\pgfqpoint{1.044390in}{0.191995in}}%
\pgfpathlineto{\pgfqpoint{1.042486in}{0.190633in}}%
\pgfpathlineto{\pgfqpoint{1.041396in}{0.191487in}}%
\pgfpathlineto{\pgfqpoint{1.042996in}{0.194874in}}%
\pgfpathlineto{\pgfqpoint{1.045532in}{0.196881in}}%
\pgfpathlineto{\pgfqpoint{1.047284in}{0.196813in}}%
\pgfpathlineto{\pgfqpoint{1.048561in}{0.199863in}}%
\pgfpathlineto{\pgfqpoint{1.047065in}{0.200403in}}%
\pgfpathlineto{\pgfqpoint{1.046014in}{0.202505in}}%
\pgfpathlineto{\pgfqpoint{1.046696in}{0.204982in}}%
\pgfpathlineto{\pgfqpoint{1.047059in}{0.209727in}}%
\pgfpathlineto{\pgfqpoint{1.050063in}{0.211128in}}%
\pgfpathlineto{\pgfqpoint{1.054012in}{0.204553in}}%
\pgfpathlineto{\pgfqpoint{1.054925in}{0.201896in}}%
\pgfpathlineto{\pgfqpoint{1.053060in}{0.199716in}}%
\pgfpathclose%
\pgfusepath{fill}%
\end{pgfscope}%
\begin{pgfscope}%
\pgfpathrectangle{\pgfqpoint{0.100000in}{0.100000in}}{\pgfqpoint{3.608454in}{2.310000in}}%
\pgfusepath{clip}%
\pgfsetbuttcap%
\pgfsetmiterjoin%
\definecolor{currentfill}{rgb}{0.000000,0.827451,0.586275}%
\pgfsetfillcolor{currentfill}%
\pgfsetlinewidth{0.000000pt}%
\definecolor{currentstroke}{rgb}{0.000000,0.000000,0.000000}%
\pgfsetstrokecolor{currentstroke}%
\pgfsetstrokeopacity{0.000000}%
\pgfsetdash{}{0pt}%
\pgfpathmoveto{\pgfqpoint{1.064056in}{0.213557in}}%
\pgfpathlineto{\pgfqpoint{1.067915in}{0.213771in}}%
\pgfpathlineto{\pgfqpoint{1.069651in}{0.211374in}}%
\pgfpathlineto{\pgfqpoint{1.071517in}{0.212096in}}%
\pgfpathlineto{\pgfqpoint{1.073207in}{0.214262in}}%
\pgfpathlineto{\pgfqpoint{1.079012in}{0.214678in}}%
\pgfpathlineto{\pgfqpoint{1.080487in}{0.211762in}}%
\pgfpathlineto{\pgfqpoint{1.082084in}{0.210879in}}%
\pgfpathlineto{\pgfqpoint{1.082183in}{0.207755in}}%
\pgfpathlineto{\pgfqpoint{1.079181in}{0.204297in}}%
\pgfpathlineto{\pgfqpoint{1.082604in}{0.201601in}}%
\pgfpathlineto{\pgfqpoint{1.080183in}{0.196814in}}%
\pgfpathlineto{\pgfqpoint{1.082403in}{0.194305in}}%
\pgfpathlineto{\pgfqpoint{1.081247in}{0.188005in}}%
\pgfpathlineto{\pgfqpoint{1.083332in}{0.187330in}}%
\pgfpathlineto{\pgfqpoint{1.085981in}{0.184345in}}%
\pgfpathlineto{\pgfqpoint{1.090049in}{0.178164in}}%
\pgfpathlineto{\pgfqpoint{1.089922in}{0.175542in}}%
\pgfpathlineto{\pgfqpoint{1.087857in}{0.173091in}}%
\pgfpathlineto{\pgfqpoint{1.085139in}{0.174080in}}%
\pgfpathlineto{\pgfqpoint{1.081805in}{0.173502in}}%
\pgfpathlineto{\pgfqpoint{1.082038in}{0.171601in}}%
\pgfpathlineto{\pgfqpoint{1.079838in}{0.170598in}}%
\pgfpathlineto{\pgfqpoint{1.078315in}{0.173884in}}%
\pgfpathlineto{\pgfqpoint{1.075481in}{0.173524in}}%
\pgfpathlineto{\pgfqpoint{1.073238in}{0.172315in}}%
\pgfpathlineto{\pgfqpoint{1.072944in}{0.170978in}}%
\pgfpathlineto{\pgfqpoint{1.070260in}{0.170374in}}%
\pgfpathlineto{\pgfqpoint{1.070834in}{0.168040in}}%
\pgfpathlineto{\pgfqpoint{1.069320in}{0.165633in}}%
\pgfpathlineto{\pgfqpoint{1.067405in}{0.165345in}}%
\pgfpathlineto{\pgfqpoint{1.065482in}{0.169372in}}%
\pgfpathlineto{\pgfqpoint{1.069168in}{0.169743in}}%
\pgfpathlineto{\pgfqpoint{1.070313in}{0.171536in}}%
\pgfpathlineto{\pgfqpoint{1.072838in}{0.172914in}}%
\pgfpathlineto{\pgfqpoint{1.074001in}{0.176696in}}%
\pgfpathlineto{\pgfqpoint{1.075966in}{0.177890in}}%
\pgfpathlineto{\pgfqpoint{1.073883in}{0.179967in}}%
\pgfpathlineto{\pgfqpoint{1.073154in}{0.177269in}}%
\pgfpathlineto{\pgfqpoint{1.069781in}{0.176867in}}%
\pgfpathlineto{\pgfqpoint{1.068604in}{0.173680in}}%
\pgfpathlineto{\pgfqpoint{1.065898in}{0.174566in}}%
\pgfpathlineto{\pgfqpoint{1.065624in}{0.179594in}}%
\pgfpathlineto{\pgfqpoint{1.063296in}{0.179530in}}%
\pgfpathlineto{\pgfqpoint{1.064425in}{0.183198in}}%
\pgfpathlineto{\pgfqpoint{1.069941in}{0.185345in}}%
\pgfpathlineto{\pgfqpoint{1.069443in}{0.182567in}}%
\pgfpathlineto{\pgfqpoint{1.070640in}{0.180586in}}%
\pgfpathlineto{\pgfqpoint{1.071934in}{0.187975in}}%
\pgfpathlineto{\pgfqpoint{1.073144in}{0.190094in}}%
\pgfpathlineto{\pgfqpoint{1.076606in}{0.193050in}}%
\pgfpathlineto{\pgfqpoint{1.072321in}{0.195142in}}%
\pgfpathlineto{\pgfqpoint{1.072458in}{0.197025in}}%
\pgfpathlineto{\pgfqpoint{1.070874in}{0.199642in}}%
\pgfpathlineto{\pgfqpoint{1.070121in}{0.204047in}}%
\pgfpathlineto{\pgfqpoint{1.071795in}{0.204349in}}%
\pgfpathlineto{\pgfqpoint{1.071791in}{0.206742in}}%
\pgfpathlineto{\pgfqpoint{1.069524in}{0.207403in}}%
\pgfpathlineto{\pgfqpoint{1.067853in}{0.209464in}}%
\pgfpathlineto{\pgfqpoint{1.067534in}{0.211631in}}%
\pgfpathlineto{\pgfqpoint{1.065578in}{0.211187in}}%
\pgfpathlineto{\pgfqpoint{1.064056in}{0.213557in}}%
\pgfpathclose%
\pgfusepath{fill}%
\end{pgfscope}%
\begin{pgfscope}%
\pgfpathrectangle{\pgfqpoint{0.100000in}{0.100000in}}{\pgfqpoint{3.608454in}{2.310000in}}%
\pgfusepath{clip}%
\pgfsetbuttcap%
\pgfsetmiterjoin%
\definecolor{currentfill}{rgb}{0.000000,0.827451,0.586275}%
\pgfsetfillcolor{currentfill}%
\pgfsetlinewidth{0.000000pt}%
\definecolor{currentstroke}{rgb}{0.000000,0.000000,0.000000}%
\pgfsetstrokecolor{currentstroke}%
\pgfsetstrokeopacity{0.000000}%
\pgfsetdash{}{0pt}%
\pgfpathmoveto{\pgfqpoint{1.053511in}{0.195656in}}%
\pgfpathlineto{\pgfqpoint{1.056239in}{0.203313in}}%
\pgfpathlineto{\pgfqpoint{1.055808in}{0.208025in}}%
\pgfpathlineto{\pgfqpoint{1.056404in}{0.208822in}}%
\pgfpathlineto{\pgfqpoint{1.055449in}{0.213304in}}%
\pgfpathlineto{\pgfqpoint{1.057580in}{0.213988in}}%
\pgfpathlineto{\pgfqpoint{1.064166in}{0.208247in}}%
\pgfpathlineto{\pgfqpoint{1.067947in}{0.205632in}}%
\pgfpathlineto{\pgfqpoint{1.068462in}{0.202960in}}%
\pgfpathlineto{\pgfqpoint{1.067374in}{0.200817in}}%
\pgfpathlineto{\pgfqpoint{1.069430in}{0.198670in}}%
\pgfpathlineto{\pgfqpoint{1.071014in}{0.192172in}}%
\pgfpathlineto{\pgfqpoint{1.066528in}{0.191514in}}%
\pgfpathlineto{\pgfqpoint{1.064349in}{0.192706in}}%
\pgfpathlineto{\pgfqpoint{1.066503in}{0.196582in}}%
\pgfpathlineto{\pgfqpoint{1.064154in}{0.196268in}}%
\pgfpathlineto{\pgfqpoint{1.062590in}{0.194222in}}%
\pgfpathlineto{\pgfqpoint{1.059615in}{0.193121in}}%
\pgfpathlineto{\pgfqpoint{1.057190in}{0.195536in}}%
\pgfpathlineto{\pgfqpoint{1.055925in}{0.194675in}}%
\pgfpathlineto{\pgfqpoint{1.053511in}{0.195656in}}%
\pgfpathclose%
\pgfusepath{fill}%
\end{pgfscope}%
\begin{pgfscope}%
\pgfpathrectangle{\pgfqpoint{0.100000in}{0.100000in}}{\pgfqpoint{3.608454in}{2.310000in}}%
\pgfusepath{clip}%
\pgfsetbuttcap%
\pgfsetmiterjoin%
\definecolor{currentfill}{rgb}{0.000000,0.694118,0.652941}%
\pgfsetfillcolor{currentfill}%
\pgfsetlinewidth{0.000000pt}%
\definecolor{currentstroke}{rgb}{0.000000,0.000000,0.000000}%
\pgfsetstrokecolor{currentstroke}%
\pgfsetstrokeopacity{0.000000}%
\pgfsetdash{}{0pt}%
\pgfpathmoveto{\pgfqpoint{1.157736in}{1.964285in}}%
\pgfpathlineto{\pgfqpoint{1.140024in}{1.967674in}}%
\pgfpathlineto{\pgfqpoint{1.137125in}{1.963079in}}%
\pgfpathlineto{\pgfqpoint{1.128763in}{1.959930in}}%
\pgfpathlineto{\pgfqpoint{1.126712in}{1.970229in}}%
\pgfpathlineto{\pgfqpoint{1.113810in}{1.978599in}}%
\pgfpathlineto{\pgfqpoint{1.110227in}{1.984421in}}%
\pgfpathlineto{\pgfqpoint{1.103252in}{1.984857in}}%
\pgfpathlineto{\pgfqpoint{1.092311in}{1.978823in}}%
\pgfpathlineto{\pgfqpoint{1.088551in}{1.993499in}}%
\pgfpathlineto{\pgfqpoint{1.081829in}{1.997042in}}%
\pgfpathlineto{\pgfqpoint{1.080846in}{2.002972in}}%
\pgfpathlineto{\pgfqpoint{1.076031in}{2.007092in}}%
\pgfpathlineto{\pgfqpoint{1.076933in}{2.013289in}}%
\pgfpathlineto{\pgfqpoint{1.081296in}{2.021948in}}%
\pgfpathlineto{\pgfqpoint{1.081054in}{2.029697in}}%
\pgfpathlineto{\pgfqpoint{1.077963in}{2.034819in}}%
\pgfpathlineto{\pgfqpoint{1.078479in}{2.042039in}}%
\pgfpathlineto{\pgfqpoint{1.085312in}{2.052758in}}%
\pgfpathlineto{\pgfqpoint{1.092070in}{2.051353in}}%
\pgfpathlineto{\pgfqpoint{1.092981in}{2.055856in}}%
\pgfpathlineto{\pgfqpoint{1.103510in}{2.056028in}}%
\pgfpathlineto{\pgfqpoint{1.108312in}{2.062127in}}%
\pgfpathlineto{\pgfqpoint{1.114370in}{2.060903in}}%
\pgfpathlineto{\pgfqpoint{1.119922in}{2.088067in}}%
\pgfpathlineto{\pgfqpoint{1.111233in}{2.089809in}}%
\pgfpathlineto{\pgfqpoint{1.118004in}{2.122768in}}%
\pgfpathlineto{\pgfqpoint{1.134932in}{2.119015in}}%
\pgfpathlineto{\pgfqpoint{1.137742in}{2.109618in}}%
\pgfpathlineto{\pgfqpoint{1.132903in}{2.085429in}}%
\pgfpathlineto{\pgfqpoint{1.146831in}{2.082471in}}%
\pgfpathlineto{\pgfqpoint{1.141533in}{2.055406in}}%
\pgfpathlineto{\pgfqpoint{1.155484in}{2.052872in}}%
\pgfpathlineto{\pgfqpoint{1.153421in}{2.042231in}}%
\pgfpathlineto{\pgfqpoint{1.160126in}{2.041022in}}%
\pgfpathlineto{\pgfqpoint{1.164980in}{2.032458in}}%
\pgfpathlineto{\pgfqpoint{1.162054in}{2.020418in}}%
\pgfpathlineto{\pgfqpoint{1.158552in}{2.015787in}}%
\pgfpathlineto{\pgfqpoint{1.148015in}{2.013748in}}%
\pgfpathlineto{\pgfqpoint{1.145294in}{2.011278in}}%
\pgfpathlineto{\pgfqpoint{1.143228in}{2.002013in}}%
\pgfpathlineto{\pgfqpoint{1.147664in}{1.994687in}}%
\pgfpathlineto{\pgfqpoint{1.146770in}{1.973621in}}%
\pgfpathlineto{\pgfqpoint{1.157736in}{1.964285in}}%
\pgfpathclose%
\pgfusepath{fill}%
\end{pgfscope}%
\begin{pgfscope}%
\pgfpathrectangle{\pgfqpoint{0.100000in}{0.100000in}}{\pgfqpoint{3.608454in}{2.310000in}}%
\pgfusepath{clip}%
\pgfsetbuttcap%
\pgfsetmiterjoin%
\definecolor{currentfill}{rgb}{0.000000,0.435294,0.782353}%
\pgfsetfillcolor{currentfill}%
\pgfsetlinewidth{0.000000pt}%
\definecolor{currentstroke}{rgb}{0.000000,0.000000,0.000000}%
\pgfsetstrokecolor{currentstroke}%
\pgfsetstrokeopacity{0.000000}%
\pgfsetdash{}{0pt}%
\pgfpathmoveto{\pgfqpoint{1.937428in}{1.476563in}}%
\pgfpathlineto{\pgfqpoint{1.883772in}{1.478524in}}%
\pgfpathlineto{\pgfqpoint{1.884864in}{1.505981in}}%
\pgfpathlineto{\pgfqpoint{1.883282in}{1.506053in}}%
\pgfpathlineto{\pgfqpoint{1.884664in}{1.533508in}}%
\pgfpathlineto{\pgfqpoint{1.911513in}{1.532437in}}%
\pgfpathlineto{\pgfqpoint{1.912100in}{1.522170in}}%
\pgfpathlineto{\pgfqpoint{1.922421in}{1.521780in}}%
\pgfpathlineto{\pgfqpoint{1.925386in}{1.523978in}}%
\pgfpathlineto{\pgfqpoint{1.946070in}{1.524600in}}%
\pgfpathlineto{\pgfqpoint{1.946376in}{1.531379in}}%
\pgfpathlineto{\pgfqpoint{1.952621in}{1.531206in}}%
\pgfpathlineto{\pgfqpoint{1.952480in}{1.526297in}}%
\pgfpathlineto{\pgfqpoint{1.942299in}{1.518835in}}%
\pgfpathlineto{\pgfqpoint{1.938437in}{1.514106in}}%
\pgfpathlineto{\pgfqpoint{1.937428in}{1.476563in}}%
\pgfpathclose%
\pgfusepath{fill}%
\end{pgfscope}%
\begin{pgfscope}%
\pgfpathrectangle{\pgfqpoint{0.100000in}{0.100000in}}{\pgfqpoint{3.608454in}{2.310000in}}%
\pgfusepath{clip}%
\pgfsetbuttcap%
\pgfsetmiterjoin%
\definecolor{currentfill}{rgb}{0.000000,0.678431,0.660784}%
\pgfsetfillcolor{currentfill}%
\pgfsetlinewidth{0.000000pt}%
\definecolor{currentstroke}{rgb}{0.000000,0.000000,0.000000}%
\pgfsetstrokecolor{currentstroke}%
\pgfsetstrokeopacity{0.000000}%
\pgfsetdash{}{0pt}%
\pgfpathmoveto{\pgfqpoint{0.799750in}{1.984499in}}%
\pgfpathlineto{\pgfqpoint{0.808158in}{1.981911in}}%
\pgfpathlineto{\pgfqpoint{0.814622in}{1.975723in}}%
\pgfpathlineto{\pgfqpoint{0.813336in}{1.971025in}}%
\pgfpathlineto{\pgfqpoint{0.807404in}{1.969358in}}%
\pgfpathlineto{\pgfqpoint{0.807519in}{1.962306in}}%
\pgfpathlineto{\pgfqpoint{0.793910in}{1.961931in}}%
\pgfpathlineto{\pgfqpoint{0.793832in}{1.958919in}}%
\pgfpathlineto{\pgfqpoint{0.801936in}{1.953999in}}%
\pgfpathlineto{\pgfqpoint{0.803077in}{1.950037in}}%
\pgfpathlineto{\pgfqpoint{0.793934in}{1.941045in}}%
\pgfpathlineto{\pgfqpoint{0.792882in}{1.935165in}}%
\pgfpathlineto{\pgfqpoint{0.785890in}{1.927298in}}%
\pgfpathlineto{\pgfqpoint{0.800486in}{1.923543in}}%
\pgfpathlineto{\pgfqpoint{0.796408in}{1.906884in}}%
\pgfpathlineto{\pgfqpoint{0.764102in}{1.916069in}}%
\pgfpathlineto{\pgfqpoint{0.762277in}{1.909354in}}%
\pgfpathlineto{\pgfqpoint{0.715935in}{1.921940in}}%
\pgfpathlineto{\pgfqpoint{0.723444in}{1.948508in}}%
\pgfpathlineto{\pgfqpoint{0.734618in}{1.987950in}}%
\pgfpathlineto{\pgfqpoint{0.737189in}{2.001359in}}%
\pgfpathlineto{\pgfqpoint{0.799750in}{1.984499in}}%
\pgfpathclose%
\pgfusepath{fill}%
\end{pgfscope}%
\begin{pgfscope}%
\pgfpathrectangle{\pgfqpoint{0.100000in}{0.100000in}}{\pgfqpoint{3.608454in}{2.310000in}}%
\pgfusepath{clip}%
\pgfsetbuttcap%
\pgfsetmiterjoin%
\definecolor{currentfill}{rgb}{0.000000,0.698039,0.650980}%
\pgfsetfillcolor{currentfill}%
\pgfsetlinewidth{0.000000pt}%
\definecolor{currentstroke}{rgb}{0.000000,0.000000,0.000000}%
\pgfsetstrokecolor{currentstroke}%
\pgfsetstrokeopacity{0.000000}%
\pgfsetdash{}{0pt}%
\pgfpathmoveto{\pgfqpoint{1.267460in}{1.535366in}}%
\pgfpathlineto{\pgfqpoint{1.255196in}{1.457818in}}%
\pgfpathlineto{\pgfqpoint{1.253203in}{1.445105in}}%
\pgfpathlineto{\pgfqpoint{1.249416in}{1.442841in}}%
\pgfpathlineto{\pgfqpoint{1.194163in}{1.451751in}}%
\pgfpathlineto{\pgfqpoint{1.194698in}{1.460074in}}%
\pgfpathlineto{\pgfqpoint{1.199187in}{1.467135in}}%
\pgfpathlineto{\pgfqpoint{1.199560in}{1.472867in}}%
\pgfpathlineto{\pgfqpoint{1.206990in}{1.477188in}}%
\pgfpathlineto{\pgfqpoint{1.169072in}{1.483661in}}%
\pgfpathlineto{\pgfqpoint{1.148850in}{1.487724in}}%
\pgfpathlineto{\pgfqpoint{1.148005in}{1.494799in}}%
\pgfpathlineto{\pgfqpoint{1.158113in}{1.555519in}}%
\pgfpathlineto{\pgfqpoint{1.167943in}{1.559368in}}%
\pgfpathlineto{\pgfqpoint{1.173397in}{1.557849in}}%
\pgfpathlineto{\pgfqpoint{1.190708in}{1.558686in}}%
\pgfpathlineto{\pgfqpoint{1.198089in}{1.561222in}}%
\pgfpathlineto{\pgfqpoint{1.203263in}{1.558508in}}%
\pgfpathlineto{\pgfqpoint{1.213136in}{1.556950in}}%
\pgfpathlineto{\pgfqpoint{1.220332in}{1.551758in}}%
\pgfpathlineto{\pgfqpoint{1.241154in}{1.553467in}}%
\pgfpathlineto{\pgfqpoint{1.242723in}{1.546157in}}%
\pgfpathlineto{\pgfqpoint{1.248664in}{1.549321in}}%
\pgfpathlineto{\pgfqpoint{1.249496in}{1.554569in}}%
\pgfpathlineto{\pgfqpoint{1.260798in}{1.552822in}}%
\pgfpathlineto{\pgfqpoint{1.262249in}{1.546769in}}%
\pgfpathlineto{\pgfqpoint{1.261144in}{1.538773in}}%
\pgfpathlineto{\pgfqpoint{1.267460in}{1.535366in}}%
\pgfpathclose%
\pgfusepath{fill}%
\end{pgfscope}%
\begin{pgfscope}%
\pgfpathrectangle{\pgfqpoint{0.100000in}{0.100000in}}{\pgfqpoint{3.608454in}{2.310000in}}%
\pgfusepath{clip}%
\pgfsetbuttcap%
\pgfsetmiterjoin%
\definecolor{currentfill}{rgb}{0.000000,0.631373,0.684314}%
\pgfsetfillcolor{currentfill}%
\pgfsetlinewidth{0.000000pt}%
\definecolor{currentstroke}{rgb}{0.000000,0.000000,0.000000}%
\pgfsetstrokecolor{currentstroke}%
\pgfsetstrokeopacity{0.000000}%
\pgfsetdash{}{0pt}%
\pgfpathmoveto{\pgfqpoint{2.106820in}{1.155478in}}%
\pgfpathlineto{\pgfqpoint{2.106837in}{1.177041in}}%
\pgfpathlineto{\pgfqpoint{2.102150in}{1.181491in}}%
\pgfpathlineto{\pgfqpoint{2.102082in}{1.211883in}}%
\pgfpathlineto{\pgfqpoint{2.101288in}{1.239463in}}%
\pgfpathlineto{\pgfqpoint{2.101955in}{1.263584in}}%
\pgfpathlineto{\pgfqpoint{2.130827in}{1.263573in}}%
\pgfpathlineto{\pgfqpoint{2.130830in}{1.210394in}}%
\pgfpathlineto{\pgfqpoint{2.164491in}{1.209439in}}%
\pgfpathlineto{\pgfqpoint{2.166408in}{1.204778in}}%
\pgfpathlineto{\pgfqpoint{2.165811in}{1.161887in}}%
\pgfpathlineto{\pgfqpoint{2.130979in}{1.163114in}}%
\pgfpathlineto{\pgfqpoint{2.131026in}{1.155307in}}%
\pgfpathlineto{\pgfqpoint{2.106820in}{1.155478in}}%
\pgfpathclose%
\pgfusepath{fill}%
\end{pgfscope}%
\begin{pgfscope}%
\pgfpathrectangle{\pgfqpoint{0.100000in}{0.100000in}}{\pgfqpoint{3.608454in}{2.310000in}}%
\pgfusepath{clip}%
\pgfsetbuttcap%
\pgfsetmiterjoin%
\definecolor{currentfill}{rgb}{0.000000,0.333333,0.833333}%
\pgfsetfillcolor{currentfill}%
\pgfsetlinewidth{0.000000pt}%
\definecolor{currentstroke}{rgb}{0.000000,0.000000,0.000000}%
\pgfsetstrokecolor{currentstroke}%
\pgfsetstrokeopacity{0.000000}%
\pgfsetdash{}{0pt}%
\pgfpathmoveto{\pgfqpoint{3.080944in}{0.255920in}}%
\pgfpathlineto{\pgfqpoint{3.081905in}{0.258847in}}%
\pgfpathlineto{\pgfqpoint{3.093100in}{0.264126in}}%
\pgfpathlineto{\pgfqpoint{3.097421in}{0.270578in}}%
\pgfpathlineto{\pgfqpoint{3.104568in}{0.275364in}}%
\pgfpathlineto{\pgfqpoint{3.114895in}{0.274075in}}%
\pgfpathlineto{\pgfqpoint{3.119568in}{0.271124in}}%
\pgfpathlineto{\pgfqpoint{3.110703in}{0.266525in}}%
\pgfpathlineto{\pgfqpoint{3.103147in}{0.268644in}}%
\pgfpathlineto{\pgfqpoint{3.101531in}{0.264974in}}%
\pgfpathlineto{\pgfqpoint{3.089987in}{0.258378in}}%
\pgfpathlineto{\pgfqpoint{3.080944in}{0.255920in}}%
\pgfpathclose%
\pgfusepath{fill}%
\end{pgfscope}%
\begin{pgfscope}%
\pgfpathrectangle{\pgfqpoint{0.100000in}{0.100000in}}{\pgfqpoint{3.608454in}{2.310000in}}%
\pgfusepath{clip}%
\pgfsetbuttcap%
\pgfsetmiterjoin%
\definecolor{currentfill}{rgb}{0.000000,0.333333,0.833333}%
\pgfsetfillcolor{currentfill}%
\pgfsetlinewidth{0.000000pt}%
\definecolor{currentstroke}{rgb}{0.000000,0.000000,0.000000}%
\pgfsetstrokecolor{currentstroke}%
\pgfsetstrokeopacity{0.000000}%
\pgfsetdash{}{0pt}%
\pgfpathmoveto{\pgfqpoint{3.099940in}{0.362168in}}%
\pgfpathlineto{\pgfqpoint{3.133900in}{0.367620in}}%
\pgfpathlineto{\pgfqpoint{3.127726in}{0.403873in}}%
\pgfpathlineto{\pgfqpoint{3.126693in}{0.409803in}}%
\pgfpathlineto{\pgfqpoint{3.173551in}{0.417385in}}%
\pgfpathlineto{\pgfqpoint{3.184616in}{0.418072in}}%
\pgfpathlineto{\pgfqpoint{3.185207in}{0.398995in}}%
\pgfpathlineto{\pgfqpoint{3.187948in}{0.377138in}}%
\pgfpathlineto{\pgfqpoint{3.187162in}{0.370691in}}%
\pgfpathlineto{\pgfqpoint{3.180144in}{0.367508in}}%
\pgfpathlineto{\pgfqpoint{3.176678in}{0.346943in}}%
\pgfpathlineto{\pgfqpoint{3.178542in}{0.336148in}}%
\pgfpathlineto{\pgfqpoint{3.168359in}{0.323529in}}%
\pgfpathlineto{\pgfqpoint{3.163072in}{0.325621in}}%
\pgfpathlineto{\pgfqpoint{3.160642in}{0.320071in}}%
\pgfpathlineto{\pgfqpoint{3.154594in}{0.316205in}}%
\pgfpathlineto{\pgfqpoint{3.143564in}{0.317485in}}%
\pgfpathlineto{\pgfqpoint{3.140330in}{0.313854in}}%
\pgfpathlineto{\pgfqpoint{3.126999in}{0.309857in}}%
\pgfpathlineto{\pgfqpoint{3.119526in}{0.317390in}}%
\pgfpathlineto{\pgfqpoint{3.119713in}{0.326025in}}%
\pgfpathlineto{\pgfqpoint{3.116825in}{0.336890in}}%
\pgfpathlineto{\pgfqpoint{3.113204in}{0.339888in}}%
\pgfpathlineto{\pgfqpoint{3.109115in}{0.349472in}}%
\pgfpathlineto{\pgfqpoint{3.101613in}{0.355282in}}%
\pgfpathlineto{\pgfqpoint{3.099940in}{0.362168in}}%
\pgfpathclose%
\pgfusepath{fill}%
\end{pgfscope}%
\begin{pgfscope}%
\pgfpathrectangle{\pgfqpoint{0.100000in}{0.100000in}}{\pgfqpoint{3.608454in}{2.310000in}}%
\pgfusepath{clip}%
\pgfsetbuttcap%
\pgfsetmiterjoin%
\definecolor{currentfill}{rgb}{0.000000,0.698039,0.650980}%
\pgfsetfillcolor{currentfill}%
\pgfsetlinewidth{0.000000pt}%
\definecolor{currentstroke}{rgb}{0.000000,0.000000,0.000000}%
\pgfsetstrokecolor{currentstroke}%
\pgfsetstrokeopacity{0.000000}%
\pgfsetdash{}{0pt}%
\pgfpathmoveto{\pgfqpoint{2.547754in}{0.792351in}}%
\pgfpathlineto{\pgfqpoint{2.536679in}{0.788374in}}%
\pgfpathlineto{\pgfqpoint{2.518619in}{0.784958in}}%
\pgfpathlineto{\pgfqpoint{2.516272in}{0.816527in}}%
\pgfpathlineto{\pgfqpoint{2.488805in}{0.814690in}}%
\pgfpathlineto{\pgfqpoint{2.485523in}{0.870923in}}%
\pgfpathlineto{\pgfqpoint{2.519263in}{0.872679in}}%
\pgfpathlineto{\pgfqpoint{2.550184in}{0.875112in}}%
\pgfpathlineto{\pgfqpoint{2.549490in}{0.847081in}}%
\pgfpathlineto{\pgfqpoint{2.547754in}{0.792351in}}%
\pgfpathclose%
\pgfusepath{fill}%
\end{pgfscope}%
\begin{pgfscope}%
\pgfpathrectangle{\pgfqpoint{0.100000in}{0.100000in}}{\pgfqpoint{3.608454in}{2.310000in}}%
\pgfusepath{clip}%
\pgfsetbuttcap%
\pgfsetmiterjoin%
\definecolor{currentfill}{rgb}{0.000000,0.490196,0.754902}%
\pgfsetfillcolor{currentfill}%
\pgfsetlinewidth{0.000000pt}%
\definecolor{currentstroke}{rgb}{0.000000,0.000000,0.000000}%
\pgfsetstrokecolor{currentstroke}%
\pgfsetstrokeopacity{0.000000}%
\pgfsetdash{}{0pt}%
\pgfpathmoveto{\pgfqpoint{1.244902in}{0.674701in}}%
\pgfpathlineto{\pgfqpoint{1.249588in}{0.682742in}}%
\pgfpathlineto{\pgfqpoint{1.263659in}{0.679913in}}%
\pgfpathlineto{\pgfqpoint{1.258636in}{0.675464in}}%
\pgfpathlineto{\pgfqpoint{1.252077in}{0.677773in}}%
\pgfpathlineto{\pgfqpoint{1.244902in}{0.674701in}}%
\pgfpathclose%
\pgfusepath{fill}%
\end{pgfscope}%
\begin{pgfscope}%
\pgfpathrectangle{\pgfqpoint{0.100000in}{0.100000in}}{\pgfqpoint{3.608454in}{2.310000in}}%
\pgfusepath{clip}%
\pgfsetbuttcap%
\pgfsetmiterjoin%
\definecolor{currentfill}{rgb}{0.000000,0.490196,0.754902}%
\pgfsetfillcolor{currentfill}%
\pgfsetlinewidth{0.000000pt}%
\definecolor{currentstroke}{rgb}{0.000000,0.000000,0.000000}%
\pgfsetstrokecolor{currentstroke}%
\pgfsetstrokeopacity{0.000000}%
\pgfsetdash{}{0pt}%
\pgfpathmoveto{\pgfqpoint{1.293790in}{0.642427in}}%
\pgfpathlineto{\pgfqpoint{1.285285in}{0.654449in}}%
\pgfpathlineto{\pgfqpoint{1.281587in}{0.664883in}}%
\pgfpathlineto{\pgfqpoint{1.282478in}{0.669549in}}%
\pgfpathlineto{\pgfqpoint{1.292812in}{0.673886in}}%
\pgfpathlineto{\pgfqpoint{1.304918in}{0.671461in}}%
\pgfpathlineto{\pgfqpoint{1.308962in}{0.666408in}}%
\pgfpathlineto{\pgfqpoint{1.315946in}{0.662668in}}%
\pgfpathlineto{\pgfqpoint{1.318323in}{0.655591in}}%
\pgfpathlineto{\pgfqpoint{1.315484in}{0.650135in}}%
\pgfpathlineto{\pgfqpoint{1.309751in}{0.648208in}}%
\pgfpathlineto{\pgfqpoint{1.305345in}{0.642514in}}%
\pgfpathlineto{\pgfqpoint{1.293790in}{0.642427in}}%
\pgfpathclose%
\pgfusepath{fill}%
\end{pgfscope}%
\begin{pgfscope}%
\pgfpathrectangle{\pgfqpoint{0.100000in}{0.100000in}}{\pgfqpoint{3.608454in}{2.310000in}}%
\pgfusepath{clip}%
\pgfsetbuttcap%
\pgfsetmiterjoin%
\definecolor{currentfill}{rgb}{0.000000,0.478431,0.760784}%
\pgfsetfillcolor{currentfill}%
\pgfsetlinewidth{0.000000pt}%
\definecolor{currentstroke}{rgb}{0.000000,0.000000,0.000000}%
\pgfsetstrokecolor{currentstroke}%
\pgfsetstrokeopacity{0.000000}%
\pgfsetdash{}{0pt}%
\pgfpathmoveto{\pgfqpoint{1.946085in}{2.030374in}}%
\pgfpathlineto{\pgfqpoint{1.945081in}{2.002729in}}%
\pgfpathlineto{\pgfqpoint{1.946235in}{1.995849in}}%
\pgfpathlineto{\pgfqpoint{1.918791in}{1.996815in}}%
\pgfpathlineto{\pgfqpoint{1.917334in}{2.003750in}}%
\pgfpathlineto{\pgfqpoint{1.918355in}{2.031373in}}%
\pgfpathlineto{\pgfqpoint{1.946085in}{2.030374in}}%
\pgfpathclose%
\pgfusepath{fill}%
\end{pgfscope}%
\begin{pgfscope}%
\pgfpathrectangle{\pgfqpoint{0.100000in}{0.100000in}}{\pgfqpoint{3.608454in}{2.310000in}}%
\pgfusepath{clip}%
\pgfsetbuttcap%
\pgfsetmiterjoin%
\definecolor{currentfill}{rgb}{0.000000,0.439216,0.780392}%
\pgfsetfillcolor{currentfill}%
\pgfsetlinewidth{0.000000pt}%
\definecolor{currentstroke}{rgb}{0.000000,0.000000,0.000000}%
\pgfsetstrokecolor{currentstroke}%
\pgfsetstrokeopacity{0.000000}%
\pgfsetdash{}{0pt}%
\pgfpathmoveto{\pgfqpoint{1.474731in}{1.733639in}}%
\pgfpathlineto{\pgfqpoint{1.441739in}{1.737485in}}%
\pgfpathlineto{\pgfqpoint{1.412417in}{1.742036in}}%
\pgfpathlineto{\pgfqpoint{1.417395in}{1.794618in}}%
\pgfpathlineto{\pgfqpoint{1.419171in}{1.807060in}}%
\pgfpathlineto{\pgfqpoint{1.417230in}{1.817131in}}%
\pgfpathlineto{\pgfqpoint{1.412375in}{1.821775in}}%
\pgfpathlineto{\pgfqpoint{1.406562in}{1.831802in}}%
\pgfpathlineto{\pgfqpoint{1.396824in}{1.840362in}}%
\pgfpathlineto{\pgfqpoint{1.389944in}{1.843550in}}%
\pgfpathlineto{\pgfqpoint{1.384830in}{1.855882in}}%
\pgfpathlineto{\pgfqpoint{1.383661in}{1.866297in}}%
\pgfpathlineto{\pgfqpoint{1.364865in}{1.868873in}}%
\pgfpathlineto{\pgfqpoint{1.372850in}{1.883266in}}%
\pgfpathlineto{\pgfqpoint{1.375793in}{1.885029in}}%
\pgfpathlineto{\pgfqpoint{1.342451in}{1.889782in}}%
\pgfpathlineto{\pgfqpoint{1.345983in}{1.906551in}}%
\pgfpathlineto{\pgfqpoint{1.348402in}{1.908679in}}%
\pgfpathlineto{\pgfqpoint{1.366491in}{1.905688in}}%
\pgfpathlineto{\pgfqpoint{1.380767in}{1.908242in}}%
\pgfpathlineto{\pgfqpoint{1.383830in}{1.929135in}}%
\pgfpathlineto{\pgfqpoint{1.386336in}{1.938080in}}%
\pgfpathlineto{\pgfqpoint{1.395427in}{1.936746in}}%
\pgfpathlineto{\pgfqpoint{1.400977in}{1.942930in}}%
\pgfpathlineto{\pgfqpoint{1.407750in}{1.941975in}}%
\pgfpathlineto{\pgfqpoint{1.408401in}{1.946543in}}%
\pgfpathlineto{\pgfqpoint{1.417308in}{1.945271in}}%
\pgfpathlineto{\pgfqpoint{1.422053in}{1.944628in}}%
\pgfpathlineto{\pgfqpoint{1.421111in}{1.937794in}}%
\pgfpathlineto{\pgfqpoint{1.434747in}{1.935937in}}%
\pgfpathlineto{\pgfqpoint{1.446859in}{1.927356in}}%
\pgfpathlineto{\pgfqpoint{1.446358in}{1.912489in}}%
\pgfpathlineto{\pgfqpoint{1.456350in}{1.910291in}}%
\pgfpathlineto{\pgfqpoint{1.453065in}{1.884941in}}%
\pgfpathlineto{\pgfqpoint{1.449171in}{1.871656in}}%
\pgfpathlineto{\pgfqpoint{1.477206in}{1.868092in}}%
\pgfpathlineto{\pgfqpoint{1.475384in}{1.853420in}}%
\pgfpathlineto{\pgfqpoint{1.488701in}{1.851762in}}%
\pgfpathlineto{\pgfqpoint{1.486605in}{1.827974in}}%
\pgfpathlineto{\pgfqpoint{1.481253in}{1.786760in}}%
\pgfpathlineto{\pgfqpoint{1.474731in}{1.733639in}}%
\pgfpathclose%
\pgfusepath{fill}%
\end{pgfscope}%
\begin{pgfscope}%
\pgfpathrectangle{\pgfqpoint{0.100000in}{0.100000in}}{\pgfqpoint{3.608454in}{2.310000in}}%
\pgfusepath{clip}%
\pgfsetbuttcap%
\pgfsetmiterjoin%
\definecolor{currentfill}{rgb}{0.000000,0.603922,0.698039}%
\pgfsetfillcolor{currentfill}%
\pgfsetlinewidth{0.000000pt}%
\definecolor{currentstroke}{rgb}{0.000000,0.000000,0.000000}%
\pgfsetstrokecolor{currentstroke}%
\pgfsetstrokeopacity{0.000000}%
\pgfsetdash{}{0pt}%
\pgfpathmoveto{\pgfqpoint{2.184096in}{1.723863in}}%
\pgfpathlineto{\pgfqpoint{2.177278in}{1.723751in}}%
\pgfpathlineto{\pgfqpoint{2.176891in}{1.751349in}}%
\pgfpathlineto{\pgfqpoint{2.190701in}{1.751566in}}%
\pgfpathlineto{\pgfqpoint{2.190305in}{1.779149in}}%
\pgfpathlineto{\pgfqpoint{2.204050in}{1.779422in}}%
\pgfpathlineto{\pgfqpoint{2.204133in}{1.773697in}}%
\pgfpathlineto{\pgfqpoint{2.217878in}{1.773970in}}%
\pgfpathlineto{\pgfqpoint{2.218568in}{1.724511in}}%
\pgfpathlineto{\pgfqpoint{2.184096in}{1.723863in}}%
\pgfpathclose%
\pgfusepath{fill}%
\end{pgfscope}%
\begin{pgfscope}%
\pgfpathrectangle{\pgfqpoint{0.100000in}{0.100000in}}{\pgfqpoint{3.608454in}{2.310000in}}%
\pgfusepath{clip}%
\pgfsetbuttcap%
\pgfsetmiterjoin%
\definecolor{currentfill}{rgb}{0.000000,0.364706,0.817647}%
\pgfsetfillcolor{currentfill}%
\pgfsetlinewidth{0.000000pt}%
\definecolor{currentstroke}{rgb}{0.000000,0.000000,0.000000}%
\pgfsetstrokecolor{currentstroke}%
\pgfsetstrokeopacity{0.000000}%
\pgfsetdash{}{0pt}%
\pgfpathmoveto{\pgfqpoint{1.789033in}{1.324734in}}%
\pgfpathlineto{\pgfqpoint{1.824072in}{1.322793in}}%
\pgfpathlineto{\pgfqpoint{1.822395in}{1.288502in}}%
\pgfpathlineto{\pgfqpoint{1.782549in}{1.290657in}}%
\pgfpathlineto{\pgfqpoint{1.754213in}{1.292521in}}%
\pgfpathlineto{\pgfqpoint{1.756162in}{1.326912in}}%
\pgfpathlineto{\pgfqpoint{1.789033in}{1.324734in}}%
\pgfpathclose%
\pgfusepath{fill}%
\end{pgfscope}%
\begin{pgfscope}%
\pgfpathrectangle{\pgfqpoint{0.100000in}{0.100000in}}{\pgfqpoint{3.608454in}{2.310000in}}%
\pgfusepath{clip}%
\pgfsetbuttcap%
\pgfsetmiterjoin%
\definecolor{currentfill}{rgb}{0.000000,0.619608,0.690196}%
\pgfsetfillcolor{currentfill}%
\pgfsetlinewidth{0.000000pt}%
\definecolor{currentstroke}{rgb}{0.000000,0.000000,0.000000}%
\pgfsetstrokecolor{currentstroke}%
\pgfsetstrokeopacity{0.000000}%
\pgfsetdash{}{0pt}%
\pgfpathmoveto{\pgfqpoint{1.720336in}{1.913435in}}%
\pgfpathlineto{\pgfqpoint{1.726123in}{1.914246in}}%
\pgfpathlineto{\pgfqpoint{1.730862in}{1.910799in}}%
\pgfpathlineto{\pgfqpoint{1.744713in}{1.908214in}}%
\pgfpathlineto{\pgfqpoint{1.756437in}{1.912852in}}%
\pgfpathlineto{\pgfqpoint{1.759368in}{1.916843in}}%
\pgfpathlineto{\pgfqpoint{1.775193in}{1.927930in}}%
\pgfpathlineto{\pgfqpoint{1.783790in}{1.937380in}}%
\pgfpathlineto{\pgfqpoint{1.787591in}{1.935130in}}%
\pgfpathlineto{\pgfqpoint{1.799830in}{1.937828in}}%
\pgfpathlineto{\pgfqpoint{1.801655in}{1.932573in}}%
\pgfpathlineto{\pgfqpoint{1.800710in}{1.922898in}}%
\pgfpathlineto{\pgfqpoint{1.796911in}{1.916921in}}%
\pgfpathlineto{\pgfqpoint{1.796957in}{1.908802in}}%
\pgfpathlineto{\pgfqpoint{1.802493in}{1.899008in}}%
\pgfpathlineto{\pgfqpoint{1.808345in}{1.892715in}}%
\pgfpathlineto{\pgfqpoint{1.811895in}{1.879653in}}%
\pgfpathlineto{\pgfqpoint{1.807298in}{1.876074in}}%
\pgfpathlineto{\pgfqpoint{1.801983in}{1.865242in}}%
\pgfpathlineto{\pgfqpoint{1.808999in}{1.861084in}}%
\pgfpathlineto{\pgfqpoint{1.773598in}{1.863245in}}%
\pgfpathlineto{\pgfqpoint{1.716695in}{1.867335in}}%
\pgfpathlineto{\pgfqpoint{1.720336in}{1.913435in}}%
\pgfpathclose%
\pgfusepath{fill}%
\end{pgfscope}%
\begin{pgfscope}%
\pgfpathrectangle{\pgfqpoint{0.100000in}{0.100000in}}{\pgfqpoint{3.608454in}{2.310000in}}%
\pgfusepath{clip}%
\pgfsetbuttcap%
\pgfsetmiterjoin%
\definecolor{currentfill}{rgb}{0.000000,0.407843,0.796078}%
\pgfsetfillcolor{currentfill}%
\pgfsetlinewidth{0.000000pt}%
\definecolor{currentstroke}{rgb}{0.000000,0.000000,0.000000}%
\pgfsetstrokecolor{currentstroke}%
\pgfsetstrokeopacity{0.000000}%
\pgfsetdash{}{0pt}%
\pgfpathmoveto{\pgfqpoint{1.347428in}{1.265725in}}%
\pgfpathlineto{\pgfqpoint{1.325433in}{1.268700in}}%
\pgfpathlineto{\pgfqpoint{1.329166in}{1.295757in}}%
\pgfpathlineto{\pgfqpoint{1.327961in}{1.309954in}}%
\pgfpathlineto{\pgfqpoint{1.325950in}{1.311919in}}%
\pgfpathlineto{\pgfqpoint{1.328108in}{1.320759in}}%
\pgfpathlineto{\pgfqpoint{1.328146in}{1.332828in}}%
\pgfpathlineto{\pgfqpoint{1.325556in}{1.338863in}}%
\pgfpathlineto{\pgfqpoint{1.333852in}{1.337711in}}%
\pgfpathlineto{\pgfqpoint{1.343717in}{1.409436in}}%
\pgfpathlineto{\pgfqpoint{1.350593in}{1.411552in}}%
\pgfpathlineto{\pgfqpoint{1.352191in}{1.404691in}}%
\pgfpathlineto{\pgfqpoint{1.355918in}{1.399849in}}%
\pgfpathlineto{\pgfqpoint{1.369057in}{1.398085in}}%
\pgfpathlineto{\pgfqpoint{1.377365in}{1.386918in}}%
\pgfpathlineto{\pgfqpoint{1.383841in}{1.385038in}}%
\pgfpathlineto{\pgfqpoint{1.387549in}{1.389567in}}%
\pgfpathlineto{\pgfqpoint{1.393282in}{1.388976in}}%
\pgfpathlineto{\pgfqpoint{1.396241in}{1.384895in}}%
\pgfpathlineto{\pgfqpoint{1.403546in}{1.376980in}}%
\pgfpathlineto{\pgfqpoint{1.412208in}{1.374642in}}%
\pgfpathlineto{\pgfqpoint{1.405672in}{1.369838in}}%
\pgfpathlineto{\pgfqpoint{1.404621in}{1.362130in}}%
\pgfpathlineto{\pgfqpoint{1.406993in}{1.344082in}}%
\pgfpathlineto{\pgfqpoint{1.412097in}{1.337063in}}%
\pgfpathlineto{\pgfqpoint{1.365837in}{1.343240in}}%
\pgfpathlineto{\pgfqpoint{1.360950in}{1.306443in}}%
\pgfpathlineto{\pgfqpoint{1.352269in}{1.306231in}}%
\pgfpathlineto{\pgfqpoint{1.349065in}{1.285654in}}%
\pgfpathlineto{\pgfqpoint{1.347428in}{1.265725in}}%
\pgfpathclose%
\pgfusepath{fill}%
\end{pgfscope}%
\begin{pgfscope}%
\pgfpathrectangle{\pgfqpoint{0.100000in}{0.100000in}}{\pgfqpoint{3.608454in}{2.310000in}}%
\pgfusepath{clip}%
\pgfsetbuttcap%
\pgfsetmiterjoin%
\definecolor{currentfill}{rgb}{0.000000,0.796078,0.601961}%
\pgfsetfillcolor{currentfill}%
\pgfsetlinewidth{0.000000pt}%
\definecolor{currentstroke}{rgb}{0.000000,0.000000,0.000000}%
\pgfsetstrokecolor{currentstroke}%
\pgfsetstrokeopacity{0.000000}%
\pgfsetdash{}{0pt}%
\pgfpathmoveto{\pgfqpoint{1.716695in}{1.867335in}}%
\pgfpathlineto{\pgfqpoint{1.773598in}{1.863245in}}%
\pgfpathlineto{\pgfqpoint{1.808999in}{1.861084in}}%
\pgfpathlineto{\pgfqpoint{1.812344in}{1.853586in}}%
\pgfpathlineto{\pgfqpoint{1.809666in}{1.848703in}}%
\pgfpathlineto{\pgfqpoint{1.812489in}{1.842876in}}%
\pgfpathlineto{\pgfqpoint{1.808240in}{1.831561in}}%
\pgfpathlineto{\pgfqpoint{1.809782in}{1.824532in}}%
\pgfpathlineto{\pgfqpoint{1.802162in}{1.823198in}}%
\pgfpathlineto{\pgfqpoint{1.802764in}{1.815650in}}%
\pgfpathlineto{\pgfqpoint{1.803220in}{1.813204in}}%
\pgfpathlineto{\pgfqpoint{1.795074in}{1.805632in}}%
\pgfpathlineto{\pgfqpoint{1.787641in}{1.810142in}}%
\pgfpathlineto{\pgfqpoint{1.781507in}{1.805854in}}%
\pgfpathlineto{\pgfqpoint{1.778292in}{1.808049in}}%
\pgfpathlineto{\pgfqpoint{1.767130in}{1.804067in}}%
\pgfpathlineto{\pgfqpoint{1.760814in}{1.806325in}}%
\pgfpathlineto{\pgfqpoint{1.754993in}{1.802225in}}%
\pgfpathlineto{\pgfqpoint{1.747956in}{1.802589in}}%
\pgfpathlineto{\pgfqpoint{1.737625in}{1.793677in}}%
\pgfpathlineto{\pgfqpoint{1.729800in}{1.795494in}}%
\pgfpathlineto{\pgfqpoint{1.723276in}{1.792199in}}%
\pgfpathlineto{\pgfqpoint{1.710799in}{1.791190in}}%
\pgfpathlineto{\pgfqpoint{1.714041in}{1.832923in}}%
\pgfpathlineto{\pgfqpoint{1.716695in}{1.867335in}}%
\pgfpathclose%
\pgfusepath{fill}%
\end{pgfscope}%
\begin{pgfscope}%
\pgfpathrectangle{\pgfqpoint{0.100000in}{0.100000in}}{\pgfqpoint{3.608454in}{2.310000in}}%
\pgfusepath{clip}%
\pgfsetbuttcap%
\pgfsetmiterjoin%
\definecolor{currentfill}{rgb}{0.000000,0.490196,0.754902}%
\pgfsetfillcolor{currentfill}%
\pgfsetlinewidth{0.000000pt}%
\definecolor{currentstroke}{rgb}{0.000000,0.000000,0.000000}%
\pgfsetstrokecolor{currentstroke}%
\pgfsetstrokeopacity{0.000000}%
\pgfsetdash{}{0pt}%
\pgfpathmoveto{\pgfqpoint{2.835111in}{1.113784in}}%
\pgfpathlineto{\pgfqpoint{2.828227in}{1.112365in}}%
\pgfpathlineto{\pgfqpoint{2.817068in}{1.103142in}}%
\pgfpathlineto{\pgfqpoint{2.805797in}{1.111156in}}%
\pgfpathlineto{\pgfqpoint{2.790002in}{1.116268in}}%
\pgfpathlineto{\pgfqpoint{2.775615in}{1.111731in}}%
\pgfpathlineto{\pgfqpoint{2.772428in}{1.117636in}}%
\pgfpathlineto{\pgfqpoint{2.765621in}{1.120136in}}%
\pgfpathlineto{\pgfqpoint{2.761095in}{1.125193in}}%
\pgfpathlineto{\pgfqpoint{2.766985in}{1.132686in}}%
\pgfpathlineto{\pgfqpoint{2.763659in}{1.139195in}}%
\pgfpathlineto{\pgfqpoint{2.750357in}{1.150558in}}%
\pgfpathlineto{\pgfqpoint{2.751113in}{1.160516in}}%
\pgfpathlineto{\pgfqpoint{2.761563in}{1.168904in}}%
\pgfpathlineto{\pgfqpoint{2.771140in}{1.159342in}}%
\pgfpathlineto{\pgfqpoint{2.781495in}{1.155039in}}%
\pgfpathlineto{\pgfqpoint{2.785832in}{1.164847in}}%
\pgfpathlineto{\pgfqpoint{2.783231in}{1.181649in}}%
\pgfpathlineto{\pgfqpoint{2.788196in}{1.183891in}}%
\pgfpathlineto{\pgfqpoint{2.789452in}{1.189777in}}%
\pgfpathlineto{\pgfqpoint{2.804581in}{1.191394in}}%
\pgfpathlineto{\pgfqpoint{2.810250in}{1.179188in}}%
\pgfpathlineto{\pgfqpoint{2.814063in}{1.180206in}}%
\pgfpathlineto{\pgfqpoint{2.827195in}{1.174638in}}%
\pgfpathlineto{\pgfqpoint{2.824792in}{1.161761in}}%
\pgfpathlineto{\pgfqpoint{2.831674in}{1.145493in}}%
\pgfpathlineto{\pgfqpoint{2.823619in}{1.137722in}}%
\pgfpathlineto{\pgfqpoint{2.829784in}{1.127369in}}%
\pgfpathlineto{\pgfqpoint{2.833923in}{1.123330in}}%
\pgfpathlineto{\pgfqpoint{2.835111in}{1.113784in}}%
\pgfpathclose%
\pgfusepath{fill}%
\end{pgfscope}%
\begin{pgfscope}%
\pgfpathrectangle{\pgfqpoint{0.100000in}{0.100000in}}{\pgfqpoint{3.608454in}{2.310000in}}%
\pgfusepath{clip}%
\pgfsetbuttcap%
\pgfsetmiterjoin%
\definecolor{currentfill}{rgb}{0.000000,0.364706,0.817647}%
\pgfsetfillcolor{currentfill}%
\pgfsetlinewidth{0.000000pt}%
\definecolor{currentstroke}{rgb}{0.000000,0.000000,0.000000}%
\pgfsetstrokecolor{currentstroke}%
\pgfsetstrokeopacity{0.000000}%
\pgfsetdash{}{0pt}%
\pgfpathmoveto{\pgfqpoint{1.793434in}{1.427744in}}%
\pgfpathlineto{\pgfqpoint{1.791324in}{1.393448in}}%
\pgfpathlineto{\pgfqpoint{1.717833in}{1.398242in}}%
\pgfpathlineto{\pgfqpoint{1.716323in}{1.398352in}}%
\pgfpathlineto{\pgfqpoint{1.718916in}{1.432564in}}%
\pgfpathlineto{\pgfqpoint{1.724173in}{1.432221in}}%
\pgfpathlineto{\pgfqpoint{1.792481in}{1.427797in}}%
\pgfpathlineto{\pgfqpoint{1.793434in}{1.427744in}}%
\pgfpathclose%
\pgfusepath{fill}%
\end{pgfscope}%
\begin{pgfscope}%
\pgfpathrectangle{\pgfqpoint{0.100000in}{0.100000in}}{\pgfqpoint{3.608454in}{2.310000in}}%
\pgfusepath{clip}%
\pgfsetbuttcap%
\pgfsetmiterjoin%
\definecolor{currentfill}{rgb}{0.000000,0.607843,0.696078}%
\pgfsetfillcolor{currentfill}%
\pgfsetlinewidth{0.000000pt}%
\definecolor{currentstroke}{rgb}{0.000000,0.000000,0.000000}%
\pgfsetstrokecolor{currentstroke}%
\pgfsetstrokeopacity{0.000000}%
\pgfsetdash{}{0pt}%
\pgfpathmoveto{\pgfqpoint{3.083124in}{1.089029in}}%
\pgfpathlineto{\pgfqpoint{3.099339in}{1.091621in}}%
\pgfpathlineto{\pgfqpoint{3.116266in}{1.079750in}}%
\pgfpathlineto{\pgfqpoint{3.108787in}{1.052476in}}%
\pgfpathlineto{\pgfqpoint{3.095167in}{1.065172in}}%
\pgfpathlineto{\pgfqpoint{3.079593in}{1.063190in}}%
\pgfpathlineto{\pgfqpoint{3.065576in}{1.050744in}}%
\pgfpathlineto{\pgfqpoint{3.060908in}{1.061022in}}%
\pgfpathlineto{\pgfqpoint{3.054685in}{1.068940in}}%
\pgfpathlineto{\pgfqpoint{3.042269in}{1.083215in}}%
\pgfpathlineto{\pgfqpoint{3.083124in}{1.089029in}}%
\pgfpathclose%
\pgfusepath{fill}%
\end{pgfscope}%
\begin{pgfscope}%
\pgfpathrectangle{\pgfqpoint{0.100000in}{0.100000in}}{\pgfqpoint{3.608454in}{2.310000in}}%
\pgfusepath{clip}%
\pgfsetbuttcap%
\pgfsetmiterjoin%
\definecolor{currentfill}{rgb}{0.000000,0.713725,0.643137}%
\pgfsetfillcolor{currentfill}%
\pgfsetlinewidth{0.000000pt}%
\definecolor{currentstroke}{rgb}{0.000000,0.000000,0.000000}%
\pgfsetstrokecolor{currentstroke}%
\pgfsetstrokeopacity{0.000000}%
\pgfsetdash{}{0pt}%
\pgfpathmoveto{\pgfqpoint{1.748728in}{0.562219in}}%
\pgfpathlineto{\pgfqpoint{1.709808in}{0.564368in}}%
\pgfpathlineto{\pgfqpoint{1.702681in}{0.571973in}}%
\pgfpathlineto{\pgfqpoint{1.701376in}{0.578127in}}%
\pgfpathlineto{\pgfqpoint{1.687298in}{0.588703in}}%
\pgfpathlineto{\pgfqpoint{1.682525in}{0.597073in}}%
\pgfpathlineto{\pgfqpoint{1.678141in}{0.597930in}}%
\pgfpathlineto{\pgfqpoint{1.672342in}{0.605442in}}%
\pgfpathlineto{\pgfqpoint{1.667527in}{0.607537in}}%
\pgfpathlineto{\pgfqpoint{1.661236in}{0.620787in}}%
\pgfpathlineto{\pgfqpoint{1.653250in}{0.625462in}}%
\pgfpathlineto{\pgfqpoint{1.648358in}{0.623948in}}%
\pgfpathlineto{\pgfqpoint{1.623671in}{0.627789in}}%
\pgfpathlineto{\pgfqpoint{1.617429in}{0.627185in}}%
\pgfpathlineto{\pgfqpoint{1.599353in}{0.636098in}}%
\pgfpathlineto{\pgfqpoint{1.583492in}{0.651351in}}%
\pgfpathlineto{\pgfqpoint{1.585018in}{0.669731in}}%
\pgfpathlineto{\pgfqpoint{1.600425in}{0.668531in}}%
\pgfpathlineto{\pgfqpoint{1.602461in}{0.693594in}}%
\pgfpathlineto{\pgfqpoint{1.616475in}{0.692376in}}%
\pgfpathlineto{\pgfqpoint{1.616885in}{0.697027in}}%
\pgfpathlineto{\pgfqpoint{1.644365in}{0.694777in}}%
\pgfpathlineto{\pgfqpoint{1.648797in}{0.693179in}}%
\pgfpathlineto{\pgfqpoint{1.649534in}{0.685985in}}%
\pgfpathlineto{\pgfqpoint{1.642706in}{0.676280in}}%
\pgfpathlineto{\pgfqpoint{1.647908in}{0.670831in}}%
\pgfpathlineto{\pgfqpoint{1.640690in}{0.665716in}}%
\pgfpathlineto{\pgfqpoint{1.695622in}{0.661844in}}%
\pgfpathlineto{\pgfqpoint{1.713553in}{0.660732in}}%
\pgfpathlineto{\pgfqpoint{1.710253in}{0.607672in}}%
\pgfpathlineto{\pgfqpoint{1.751098in}{0.605189in}}%
\pgfpathlineto{\pgfqpoint{1.748728in}{0.562219in}}%
\pgfpathclose%
\pgfusepath{fill}%
\end{pgfscope}%
\begin{pgfscope}%
\pgfpathrectangle{\pgfqpoint{0.100000in}{0.100000in}}{\pgfqpoint{3.608454in}{2.310000in}}%
\pgfusepath{clip}%
\pgfsetbuttcap%
\pgfsetmiterjoin%
\definecolor{currentfill}{rgb}{0.000000,0.654902,0.672549}%
\pgfsetfillcolor{currentfill}%
\pgfsetlinewidth{0.000000pt}%
\definecolor{currentstroke}{rgb}{0.000000,0.000000,0.000000}%
\pgfsetstrokecolor{currentstroke}%
\pgfsetstrokeopacity{0.000000}%
\pgfsetdash{}{0pt}%
\pgfpathmoveto{\pgfqpoint{3.051009in}{1.477958in}}%
\pgfpathlineto{\pgfqpoint{3.051273in}{1.485944in}}%
\pgfpathlineto{\pgfqpoint{3.048216in}{1.491838in}}%
\pgfpathlineto{\pgfqpoint{3.054176in}{1.507940in}}%
\pgfpathlineto{\pgfqpoint{3.061504in}{1.519627in}}%
\pgfpathlineto{\pgfqpoint{3.068856in}{1.543568in}}%
\pgfpathlineto{\pgfqpoint{3.072318in}{1.562246in}}%
\pgfpathlineto{\pgfqpoint{3.098915in}{1.567130in}}%
\pgfpathlineto{\pgfqpoint{3.107213in}{1.564628in}}%
\pgfpathlineto{\pgfqpoint{3.111267in}{1.560885in}}%
\pgfpathlineto{\pgfqpoint{3.110368in}{1.555879in}}%
\pgfpathlineto{\pgfqpoint{3.115439in}{1.550679in}}%
\pgfpathlineto{\pgfqpoint{3.110584in}{1.534960in}}%
\pgfpathlineto{\pgfqpoint{3.112902in}{1.529498in}}%
\pgfpathlineto{\pgfqpoint{3.119792in}{1.526024in}}%
\pgfpathlineto{\pgfqpoint{3.116662in}{1.519085in}}%
\pgfpathlineto{\pgfqpoint{3.117244in}{1.511717in}}%
\pgfpathlineto{\pgfqpoint{3.114952in}{1.497798in}}%
\pgfpathlineto{\pgfqpoint{3.111561in}{1.488897in}}%
\pgfpathlineto{\pgfqpoint{3.051009in}{1.477958in}}%
\pgfpathclose%
\pgfusepath{fill}%
\end{pgfscope}%
\begin{pgfscope}%
\pgfpathrectangle{\pgfqpoint{0.100000in}{0.100000in}}{\pgfqpoint{3.608454in}{2.310000in}}%
\pgfusepath{clip}%
\pgfsetbuttcap%
\pgfsetmiterjoin%
\definecolor{currentfill}{rgb}{0.000000,0.400000,0.800000}%
\pgfsetfillcolor{currentfill}%
\pgfsetlinewidth{0.000000pt}%
\definecolor{currentstroke}{rgb}{0.000000,0.000000,0.000000}%
\pgfsetstrokecolor{currentstroke}%
\pgfsetstrokeopacity{0.000000}%
\pgfsetdash{}{0pt}%
\pgfpathmoveto{\pgfqpoint{2.335984in}{1.355415in}}%
\pgfpathlineto{\pgfqpoint{2.336652in}{1.332026in}}%
\pgfpathlineto{\pgfqpoint{2.327371in}{1.331768in}}%
\pgfpathlineto{\pgfqpoint{2.327621in}{1.321058in}}%
\pgfpathlineto{\pgfqpoint{2.319587in}{1.318096in}}%
\pgfpathlineto{\pgfqpoint{2.313973in}{1.320063in}}%
\pgfpathlineto{\pgfqpoint{2.313325in}{1.348137in}}%
\pgfpathlineto{\pgfqpoint{2.284099in}{1.347536in}}%
\pgfpathlineto{\pgfqpoint{2.284008in}{1.361432in}}%
\pgfpathlineto{\pgfqpoint{2.271150in}{1.361570in}}%
\pgfpathlineto{\pgfqpoint{2.271164in}{1.369568in}}%
\pgfpathlineto{\pgfqpoint{2.307376in}{1.370072in}}%
\pgfpathlineto{\pgfqpoint{2.324527in}{1.369104in}}%
\pgfpathlineto{\pgfqpoint{2.326861in}{1.355261in}}%
\pgfpathlineto{\pgfqpoint{2.335984in}{1.355415in}}%
\pgfpathclose%
\pgfusepath{fill}%
\end{pgfscope}%
\begin{pgfscope}%
\pgfpathrectangle{\pgfqpoint{0.100000in}{0.100000in}}{\pgfqpoint{3.608454in}{2.310000in}}%
\pgfusepath{clip}%
\pgfsetbuttcap%
\pgfsetmiterjoin%
\definecolor{currentfill}{rgb}{0.000000,0.537255,0.731373}%
\pgfsetfillcolor{currentfill}%
\pgfsetlinewidth{0.000000pt}%
\definecolor{currentstroke}{rgb}{0.000000,0.000000,0.000000}%
\pgfsetstrokecolor{currentstroke}%
\pgfsetstrokeopacity{0.000000}%
\pgfsetdash{}{0pt}%
\pgfpathmoveto{\pgfqpoint{2.182439in}{1.381621in}}%
\pgfpathlineto{\pgfqpoint{2.155044in}{1.381527in}}%
\pgfpathlineto{\pgfqpoint{2.154965in}{1.375795in}}%
\pgfpathlineto{\pgfqpoint{2.131152in}{1.375734in}}%
\pgfpathlineto{\pgfqpoint{2.131100in}{1.381764in}}%
\pgfpathlineto{\pgfqpoint{2.100521in}{1.381808in}}%
\pgfpathlineto{\pgfqpoint{2.103446in}{1.388933in}}%
\pgfpathlineto{\pgfqpoint{2.099783in}{1.391231in}}%
\pgfpathlineto{\pgfqpoint{2.086114in}{1.391298in}}%
\pgfpathlineto{\pgfqpoint{2.086211in}{1.418729in}}%
\pgfpathlineto{\pgfqpoint{2.088125in}{1.418722in}}%
\pgfpathlineto{\pgfqpoint{2.094548in}{1.411120in}}%
\pgfpathlineto{\pgfqpoint{2.101635in}{1.407776in}}%
\pgfpathlineto{\pgfqpoint{2.107194in}{1.411974in}}%
\pgfpathlineto{\pgfqpoint{2.104194in}{1.428993in}}%
\pgfpathlineto{\pgfqpoint{2.130612in}{1.428716in}}%
\pgfpathlineto{\pgfqpoint{2.130676in}{1.421843in}}%
\pgfpathlineto{\pgfqpoint{2.154074in}{1.421652in}}%
\pgfpathlineto{\pgfqpoint{2.154147in}{1.429674in}}%
\pgfpathlineto{\pgfqpoint{2.181491in}{1.429705in}}%
\pgfpathlineto{\pgfqpoint{2.181819in}{1.416001in}}%
\pgfpathlineto{\pgfqpoint{2.182439in}{1.381621in}}%
\pgfpathclose%
\pgfusepath{fill}%
\end{pgfscope}%
\begin{pgfscope}%
\pgfpathrectangle{\pgfqpoint{0.100000in}{0.100000in}}{\pgfqpoint{3.608454in}{2.310000in}}%
\pgfusepath{clip}%
\pgfsetbuttcap%
\pgfsetmiterjoin%
\definecolor{currentfill}{rgb}{0.000000,0.521569,0.739216}%
\pgfsetfillcolor{currentfill}%
\pgfsetlinewidth{0.000000pt}%
\definecolor{currentstroke}{rgb}{0.000000,0.000000,0.000000}%
\pgfsetstrokecolor{currentstroke}%
\pgfsetstrokeopacity{0.000000}%
\pgfsetdash{}{0pt}%
\pgfpathmoveto{\pgfqpoint{2.461626in}{1.152803in}}%
\pgfpathlineto{\pgfqpoint{2.457328in}{1.152433in}}%
\pgfpathlineto{\pgfqpoint{2.453898in}{1.152266in}}%
\pgfpathlineto{\pgfqpoint{2.454276in}{1.139497in}}%
\pgfpathlineto{\pgfqpoint{2.450987in}{1.144593in}}%
\pgfpathlineto{\pgfqpoint{2.440887in}{1.142638in}}%
\pgfpathlineto{\pgfqpoint{2.426853in}{1.142086in}}%
\pgfpathlineto{\pgfqpoint{2.426666in}{1.161074in}}%
\pgfpathlineto{\pgfqpoint{2.414649in}{1.160533in}}%
\pgfpathlineto{\pgfqpoint{2.413654in}{1.167791in}}%
\pgfpathlineto{\pgfqpoint{2.406457in}{1.183183in}}%
\pgfpathlineto{\pgfqpoint{2.415125in}{1.196618in}}%
\pgfpathlineto{\pgfqpoint{2.408011in}{1.196285in}}%
\pgfpathlineto{\pgfqpoint{2.407413in}{1.214350in}}%
\pgfpathlineto{\pgfqpoint{2.411944in}{1.214363in}}%
\pgfpathlineto{\pgfqpoint{2.410673in}{1.236890in}}%
\pgfpathlineto{\pgfqpoint{2.444353in}{1.238826in}}%
\pgfpathlineto{\pgfqpoint{2.449905in}{1.236650in}}%
\pgfpathlineto{\pgfqpoint{2.457038in}{1.223290in}}%
\pgfpathlineto{\pgfqpoint{2.451445in}{1.214650in}}%
\pgfpathlineto{\pgfqpoint{2.460932in}{1.196917in}}%
\pgfpathlineto{\pgfqpoint{2.467123in}{1.192232in}}%
\pgfpathlineto{\pgfqpoint{2.472760in}{1.194489in}}%
\pgfpathlineto{\pgfqpoint{2.477226in}{1.192084in}}%
\pgfpathlineto{\pgfqpoint{2.479456in}{1.190579in}}%
\pgfpathlineto{\pgfqpoint{2.477580in}{1.181419in}}%
\pgfpathlineto{\pgfqpoint{2.478809in}{1.176580in}}%
\pgfpathlineto{\pgfqpoint{2.474357in}{1.170803in}}%
\pgfpathlineto{\pgfqpoint{2.477166in}{1.167044in}}%
\pgfpathlineto{\pgfqpoint{2.473332in}{1.159111in}}%
\pgfpathlineto{\pgfqpoint{2.464384in}{1.163029in}}%
\pgfpathlineto{\pgfqpoint{2.461626in}{1.152803in}}%
\pgfpathclose%
\pgfusepath{fill}%
\end{pgfscope}%
\begin{pgfscope}%
\pgfpathrectangle{\pgfqpoint{0.100000in}{0.100000in}}{\pgfqpoint{3.608454in}{2.310000in}}%
\pgfusepath{clip}%
\pgfsetbuttcap%
\pgfsetmiterjoin%
\definecolor{currentfill}{rgb}{0.000000,0.686275,0.656863}%
\pgfsetfillcolor{currentfill}%
\pgfsetlinewidth{0.000000pt}%
\definecolor{currentstroke}{rgb}{0.000000,0.000000,0.000000}%
\pgfsetstrokecolor{currentstroke}%
\pgfsetstrokeopacity{0.000000}%
\pgfsetdash{}{0pt}%
\pgfpathmoveto{\pgfqpoint{2.547754in}{0.792351in}}%
\pgfpathlineto{\pgfqpoint{2.549469in}{0.776821in}}%
\pgfpathlineto{\pgfqpoint{2.551986in}{0.756063in}}%
\pgfpathlineto{\pgfqpoint{2.525381in}{0.754028in}}%
\pgfpathlineto{\pgfqpoint{2.487539in}{0.751603in}}%
\pgfpathlineto{\pgfqpoint{2.485619in}{0.780455in}}%
\pgfpathlineto{\pgfqpoint{2.468672in}{0.778139in}}%
\pgfpathlineto{\pgfqpoint{2.467492in}{0.796662in}}%
\pgfpathlineto{\pgfqpoint{2.462239in}{0.799019in}}%
\pgfpathlineto{\pgfqpoint{2.461411in}{0.812905in}}%
\pgfpathlineto{\pgfqpoint{2.488805in}{0.814690in}}%
\pgfpathlineto{\pgfqpoint{2.516272in}{0.816527in}}%
\pgfpathlineto{\pgfqpoint{2.518619in}{0.784958in}}%
\pgfpathlineto{\pgfqpoint{2.536679in}{0.788374in}}%
\pgfpathlineto{\pgfqpoint{2.547754in}{0.792351in}}%
\pgfpathclose%
\pgfusepath{fill}%
\end{pgfscope}%
\begin{pgfscope}%
\pgfpathrectangle{\pgfqpoint{0.100000in}{0.100000in}}{\pgfqpoint{3.608454in}{2.310000in}}%
\pgfusepath{clip}%
\pgfsetbuttcap%
\pgfsetmiterjoin%
\definecolor{currentfill}{rgb}{0.000000,0.611765,0.694118}%
\pgfsetfillcolor{currentfill}%
\pgfsetlinewidth{0.000000pt}%
\definecolor{currentstroke}{rgb}{0.000000,0.000000,0.000000}%
\pgfsetstrokecolor{currentstroke}%
\pgfsetstrokeopacity{0.000000}%
\pgfsetdash{}{0pt}%
\pgfpathmoveto{\pgfqpoint{1.557865in}{2.068157in}}%
\pgfpathlineto{\pgfqpoint{1.551007in}{2.068911in}}%
\pgfpathlineto{\pgfqpoint{1.548607in}{2.048174in}}%
\pgfpathlineto{\pgfqpoint{1.546296in}{2.048425in}}%
\pgfpathlineto{\pgfqpoint{1.543137in}{2.020826in}}%
\pgfpathlineto{\pgfqpoint{1.520131in}{2.023482in}}%
\pgfpathlineto{\pgfqpoint{1.519298in}{2.016515in}}%
\pgfpathlineto{\pgfqpoint{1.512427in}{2.017326in}}%
\pgfpathlineto{\pgfqpoint{1.513276in}{2.024315in}}%
\pgfpathlineto{\pgfqpoint{1.506437in}{2.025137in}}%
\pgfpathlineto{\pgfqpoint{1.501940in}{2.025689in}}%
\pgfpathlineto{\pgfqpoint{1.503629in}{2.039468in}}%
\pgfpathlineto{\pgfqpoint{1.496777in}{2.040294in}}%
\pgfpathlineto{\pgfqpoint{1.500521in}{2.053891in}}%
\pgfpathlineto{\pgfqpoint{1.503848in}{2.080824in}}%
\pgfpathlineto{\pgfqpoint{1.498672in}{2.081480in}}%
\pgfpathlineto{\pgfqpoint{1.496354in}{2.088951in}}%
\pgfpathlineto{\pgfqpoint{1.498726in}{2.097310in}}%
\pgfpathlineto{\pgfqpoint{1.507101in}{2.093171in}}%
\pgfpathlineto{\pgfqpoint{1.527435in}{2.089391in}}%
\pgfpathlineto{\pgfqpoint{1.529442in}{2.105984in}}%
\pgfpathlineto{\pgfqpoint{1.531517in}{2.105747in}}%
\pgfpathlineto{\pgfqpoint{1.534783in}{2.133134in}}%
\pgfpathlineto{\pgfqpoint{1.578597in}{2.128096in}}%
\pgfpathlineto{\pgfqpoint{1.589978in}{2.126848in}}%
\pgfpathlineto{\pgfqpoint{1.589220in}{2.119949in}}%
\pgfpathlineto{\pgfqpoint{1.596092in}{2.119196in}}%
\pgfpathlineto{\pgfqpoint{1.595340in}{2.112255in}}%
\pgfpathlineto{\pgfqpoint{1.626106in}{2.109098in}}%
\pgfpathlineto{\pgfqpoint{1.623124in}{2.077735in}}%
\pgfpathlineto{\pgfqpoint{1.619223in}{2.081466in}}%
\pgfpathlineto{\pgfqpoint{1.607083in}{2.083595in}}%
\pgfpathlineto{\pgfqpoint{1.601867in}{2.089959in}}%
\pgfpathlineto{\pgfqpoint{1.590181in}{2.089744in}}%
\pgfpathlineto{\pgfqpoint{1.582871in}{2.091842in}}%
\pgfpathlineto{\pgfqpoint{1.569948in}{2.088334in}}%
\pgfpathlineto{\pgfqpoint{1.562458in}{2.089780in}}%
\pgfpathlineto{\pgfqpoint{1.560678in}{2.074073in}}%
\pgfpathlineto{\pgfqpoint{1.557865in}{2.068157in}}%
\pgfpathclose%
\pgfusepath{fill}%
\end{pgfscope}%
\begin{pgfscope}%
\pgfpathrectangle{\pgfqpoint{0.100000in}{0.100000in}}{\pgfqpoint{3.608454in}{2.310000in}}%
\pgfusepath{clip}%
\pgfsetbuttcap%
\pgfsetmiterjoin%
\definecolor{currentfill}{rgb}{0.000000,0.854902,0.572549}%
\pgfsetfillcolor{currentfill}%
\pgfsetlinewidth{0.000000pt}%
\definecolor{currentstroke}{rgb}{0.000000,0.000000,0.000000}%
\pgfsetstrokecolor{currentstroke}%
\pgfsetstrokeopacity{0.000000}%
\pgfsetdash{}{0pt}%
\pgfpathmoveto{\pgfqpoint{1.795506in}{0.488863in}}%
\pgfpathlineto{\pgfqpoint{1.744553in}{0.490983in}}%
\pgfpathlineto{\pgfqpoint{1.748728in}{0.562219in}}%
\pgfpathlineto{\pgfqpoint{1.797569in}{0.560061in}}%
\pgfpathlineto{\pgfqpoint{1.796154in}{0.523908in}}%
\pgfpathlineto{\pgfqpoint{1.797091in}{0.523880in}}%
\pgfpathlineto{\pgfqpoint{1.795506in}{0.488863in}}%
\pgfpathclose%
\pgfusepath{fill}%
\end{pgfscope}%
\begin{pgfscope}%
\pgfpathrectangle{\pgfqpoint{0.100000in}{0.100000in}}{\pgfqpoint{3.608454in}{2.310000in}}%
\pgfusepath{clip}%
\pgfsetbuttcap%
\pgfsetmiterjoin%
\definecolor{currentfill}{rgb}{0.000000,0.666667,0.666667}%
\pgfsetfillcolor{currentfill}%
\pgfsetlinewidth{0.000000pt}%
\definecolor{currentstroke}{rgb}{0.000000,0.000000,0.000000}%
\pgfsetstrokecolor{currentstroke}%
\pgfsetstrokeopacity{0.000000}%
\pgfsetdash{}{0pt}%
\pgfpathmoveto{\pgfqpoint{1.443108in}{0.820998in}}%
\pgfpathlineto{\pgfqpoint{1.498395in}{0.815195in}}%
\pgfpathlineto{\pgfqpoint{1.486032in}{0.744751in}}%
\pgfpathlineto{\pgfqpoint{1.426837in}{0.715528in}}%
\pgfpathlineto{\pgfqpoint{1.429579in}{0.734241in}}%
\pgfpathlineto{\pgfqpoint{1.435480in}{0.791748in}}%
\pgfpathlineto{\pgfqpoint{1.438374in}{0.821513in}}%
\pgfpathlineto{\pgfqpoint{1.443108in}{0.820998in}}%
\pgfpathclose%
\pgfusepath{fill}%
\end{pgfscope}%
\begin{pgfscope}%
\pgfpathrectangle{\pgfqpoint{0.100000in}{0.100000in}}{\pgfqpoint{3.608454in}{2.310000in}}%
\pgfusepath{clip}%
\pgfsetbuttcap%
\pgfsetmiterjoin%
\definecolor{currentfill}{rgb}{0.000000,0.721569,0.639216}%
\pgfsetfillcolor{currentfill}%
\pgfsetlinewidth{0.000000pt}%
\definecolor{currentstroke}{rgb}{0.000000,0.000000,0.000000}%
\pgfsetstrokecolor{currentstroke}%
\pgfsetstrokeopacity{0.000000}%
\pgfsetdash{}{0pt}%
\pgfpathmoveto{\pgfqpoint{1.596724in}{1.204258in}}%
\pgfpathlineto{\pgfqpoint{1.602002in}{1.203777in}}%
\pgfpathlineto{\pgfqpoint{1.598511in}{1.164402in}}%
\pgfpathlineto{\pgfqpoint{1.596009in}{1.164627in}}%
\pgfpathlineto{\pgfqpoint{1.590719in}{1.104616in}}%
\pgfpathlineto{\pgfqpoint{1.569376in}{1.106608in}}%
\pgfpathlineto{\pgfqpoint{1.572537in}{1.133851in}}%
\pgfpathlineto{\pgfqpoint{1.545193in}{1.136393in}}%
\pgfpathlineto{\pgfqpoint{1.545891in}{1.143339in}}%
\pgfpathlineto{\pgfqpoint{1.532225in}{1.144723in}}%
\pgfpathlineto{\pgfqpoint{1.532570in}{1.148149in}}%
\pgfpathlineto{\pgfqpoint{1.538731in}{1.209408in}}%
\pgfpathlineto{\pgfqpoint{1.596724in}{1.204258in}}%
\pgfpathclose%
\pgfusepath{fill}%
\end{pgfscope}%
\begin{pgfscope}%
\pgfpathrectangle{\pgfqpoint{0.100000in}{0.100000in}}{\pgfqpoint{3.608454in}{2.310000in}}%
\pgfusepath{clip}%
\pgfsetbuttcap%
\pgfsetmiterjoin%
\definecolor{currentfill}{rgb}{0.000000,0.827451,0.586275}%
\pgfsetfillcolor{currentfill}%
\pgfsetlinewidth{0.000000pt}%
\definecolor{currentstroke}{rgb}{0.000000,0.000000,0.000000}%
\pgfsetstrokecolor{currentstroke}%
\pgfsetstrokeopacity{0.000000}%
\pgfsetdash{}{0pt}%
\pgfpathmoveto{\pgfqpoint{0.655369in}{0.626110in}}%
\pgfpathlineto{\pgfqpoint{0.652280in}{0.629547in}}%
\pgfpathlineto{\pgfqpoint{0.652480in}{0.635784in}}%
\pgfpathlineto{\pgfqpoint{0.651564in}{0.638676in}}%
\pgfpathlineto{\pgfqpoint{0.654993in}{0.639490in}}%
\pgfpathlineto{\pgfqpoint{0.653508in}{0.635458in}}%
\pgfpathlineto{\pgfqpoint{0.653829in}{0.632282in}}%
\pgfpathlineto{\pgfqpoint{0.656074in}{0.628020in}}%
\pgfpathlineto{\pgfqpoint{0.655369in}{0.626110in}}%
\pgfpathclose%
\pgfusepath{fill}%
\end{pgfscope}%
\begin{pgfscope}%
\pgfpathrectangle{\pgfqpoint{0.100000in}{0.100000in}}{\pgfqpoint{3.608454in}{2.310000in}}%
\pgfusepath{clip}%
\pgfsetbuttcap%
\pgfsetmiterjoin%
\definecolor{currentfill}{rgb}{0.000000,0.827451,0.586275}%
\pgfsetfillcolor{currentfill}%
\pgfsetlinewidth{0.000000pt}%
\definecolor{currentstroke}{rgb}{0.000000,0.000000,0.000000}%
\pgfsetstrokecolor{currentstroke}%
\pgfsetstrokeopacity{0.000000}%
\pgfsetdash{}{0pt}%
\pgfpathmoveto{\pgfqpoint{0.705347in}{0.551735in}}%
\pgfpathlineto{\pgfqpoint{0.705017in}{0.553533in}}%
\pgfpathlineto{\pgfqpoint{0.702425in}{0.555599in}}%
\pgfpathlineto{\pgfqpoint{0.701422in}{0.558633in}}%
\pgfpathlineto{\pgfqpoint{0.701257in}{0.561453in}}%
\pgfpathlineto{\pgfqpoint{0.699278in}{0.565095in}}%
\pgfpathlineto{\pgfqpoint{0.699632in}{0.569687in}}%
\pgfpathlineto{\pgfqpoint{0.700852in}{0.572122in}}%
\pgfpathlineto{\pgfqpoint{0.702514in}{0.571731in}}%
\pgfpathlineto{\pgfqpoint{0.706248in}{0.566407in}}%
\pgfpathlineto{\pgfqpoint{0.710913in}{0.566989in}}%
\pgfpathlineto{\pgfqpoint{0.713889in}{0.566223in}}%
\pgfpathlineto{\pgfqpoint{0.714374in}{0.564683in}}%
\pgfpathlineto{\pgfqpoint{0.716943in}{0.563500in}}%
\pgfpathlineto{\pgfqpoint{0.716755in}{0.561123in}}%
\pgfpathlineto{\pgfqpoint{0.719375in}{0.558723in}}%
\pgfpathlineto{\pgfqpoint{0.719272in}{0.555594in}}%
\pgfpathlineto{\pgfqpoint{0.717154in}{0.553156in}}%
\pgfpathlineto{\pgfqpoint{0.715054in}{0.551897in}}%
\pgfpathlineto{\pgfqpoint{0.714314in}{0.546900in}}%
\pgfpathlineto{\pgfqpoint{0.711935in}{0.547634in}}%
\pgfpathlineto{\pgfqpoint{0.708937in}{0.550219in}}%
\pgfpathlineto{\pgfqpoint{0.705347in}{0.551735in}}%
\pgfpathclose%
\pgfusepath{fill}%
\end{pgfscope}%
\begin{pgfscope}%
\pgfpathrectangle{\pgfqpoint{0.100000in}{0.100000in}}{\pgfqpoint{3.608454in}{2.310000in}}%
\pgfusepath{clip}%
\pgfsetbuttcap%
\pgfsetmiterjoin%
\definecolor{currentfill}{rgb}{0.000000,0.827451,0.586275}%
\pgfsetfillcolor{currentfill}%
\pgfsetlinewidth{0.000000pt}%
\definecolor{currentstroke}{rgb}{0.000000,0.000000,0.000000}%
\pgfsetstrokecolor{currentstroke}%
\pgfsetstrokeopacity{0.000000}%
\pgfsetdash{}{0pt}%
\pgfpathmoveto{\pgfqpoint{0.894650in}{0.485392in}}%
\pgfpathlineto{\pgfqpoint{0.878747in}{0.464023in}}%
\pgfpathlineto{\pgfqpoint{0.878374in}{0.466226in}}%
\pgfpathlineto{\pgfqpoint{0.874246in}{0.469508in}}%
\pgfpathlineto{\pgfqpoint{0.870359in}{0.464352in}}%
\pgfpathlineto{\pgfqpoint{0.869217in}{0.461424in}}%
\pgfpathlineto{\pgfqpoint{0.864761in}{0.455531in}}%
\pgfpathlineto{\pgfqpoint{0.853555in}{0.463966in}}%
\pgfpathlineto{\pgfqpoint{0.836300in}{0.477446in}}%
\pgfpathlineto{\pgfqpoint{0.822897in}{0.488128in}}%
\pgfpathlineto{\pgfqpoint{0.824629in}{0.490280in}}%
\pgfpathlineto{\pgfqpoint{0.820776in}{0.493434in}}%
\pgfpathlineto{\pgfqpoint{0.819027in}{0.491316in}}%
\pgfpathlineto{\pgfqpoint{0.816987in}{0.492818in}}%
\pgfpathlineto{\pgfqpoint{0.815390in}{0.490823in}}%
\pgfpathlineto{\pgfqpoint{0.813395in}{0.492489in}}%
\pgfpathlineto{\pgfqpoint{0.814957in}{0.494457in}}%
\pgfpathlineto{\pgfqpoint{0.804074in}{0.503387in}}%
\pgfpathlineto{\pgfqpoint{0.802432in}{0.501473in}}%
\pgfpathlineto{\pgfqpoint{0.800941in}{0.502662in}}%
\pgfpathlineto{\pgfqpoint{0.799211in}{0.500636in}}%
\pgfpathlineto{\pgfqpoint{0.797356in}{0.502139in}}%
\pgfpathlineto{\pgfqpoint{0.795190in}{0.500666in}}%
\pgfpathlineto{\pgfqpoint{0.791902in}{0.496847in}}%
\pgfpathlineto{\pgfqpoint{0.789763in}{0.498705in}}%
\pgfpathlineto{\pgfqpoint{0.788112in}{0.496794in}}%
\pgfpathlineto{\pgfqpoint{0.786206in}{0.498438in}}%
\pgfpathlineto{\pgfqpoint{0.782945in}{0.494696in}}%
\pgfpathlineto{\pgfqpoint{0.781042in}{0.496290in}}%
\pgfpathlineto{\pgfqpoint{0.778697in}{0.494994in}}%
\pgfpathlineto{\pgfqpoint{0.775519in}{0.491429in}}%
\pgfpathlineto{\pgfqpoint{0.773696in}{0.492989in}}%
\pgfpathlineto{\pgfqpoint{0.770402in}{0.489224in}}%
\pgfpathlineto{\pgfqpoint{0.767845in}{0.491365in}}%
\pgfpathlineto{\pgfqpoint{0.764542in}{0.487613in}}%
\pgfpathlineto{\pgfqpoint{0.762636in}{0.489281in}}%
\pgfpathlineto{\pgfqpoint{0.759307in}{0.485494in}}%
\pgfpathlineto{\pgfqpoint{0.756995in}{0.484116in}}%
\pgfpathlineto{\pgfqpoint{0.755169in}{0.485784in}}%
\pgfpathlineto{\pgfqpoint{0.752873in}{0.484455in}}%
\pgfpathlineto{\pgfqpoint{0.749547in}{0.480704in}}%
\pgfpathlineto{\pgfqpoint{0.747682in}{0.482304in}}%
\pgfpathlineto{\pgfqpoint{0.746173in}{0.480590in}}%
\pgfpathlineto{\pgfqpoint{0.744357in}{0.482168in}}%
\pgfpathlineto{\pgfqpoint{0.740722in}{0.478712in}}%
\pgfpathlineto{\pgfqpoint{0.740538in}{0.478732in}}%
\pgfpathlineto{\pgfqpoint{0.737273in}{0.481661in}}%
\pgfpathlineto{\pgfqpoint{0.735963in}{0.479492in}}%
\pgfpathlineto{\pgfqpoint{0.734892in}{0.478188in}}%
\pgfpathlineto{\pgfqpoint{0.727742in}{0.479830in}}%
\pgfpathlineto{\pgfqpoint{0.728503in}{0.483239in}}%
\pgfpathlineto{\pgfqpoint{0.731793in}{0.484480in}}%
\pgfpathlineto{\pgfqpoint{0.735386in}{0.488777in}}%
\pgfpathlineto{\pgfqpoint{0.736702in}{0.491209in}}%
\pgfpathlineto{\pgfqpoint{0.736725in}{0.494624in}}%
\pgfpathlineto{\pgfqpoint{0.737998in}{0.497489in}}%
\pgfpathlineto{\pgfqpoint{0.739339in}{0.496648in}}%
\pgfpathlineto{\pgfqpoint{0.741092in}{0.498876in}}%
\pgfpathlineto{\pgfqpoint{0.745616in}{0.498996in}}%
\pgfpathlineto{\pgfqpoint{0.746961in}{0.501426in}}%
\pgfpathlineto{\pgfqpoint{0.748923in}{0.507053in}}%
\pgfpathlineto{\pgfqpoint{0.749487in}{0.509973in}}%
\pgfpathlineto{\pgfqpoint{0.750888in}{0.512260in}}%
\pgfpathlineto{\pgfqpoint{0.751044in}{0.515415in}}%
\pgfpathlineto{\pgfqpoint{0.754037in}{0.517140in}}%
\pgfpathlineto{\pgfqpoint{0.753941in}{0.521042in}}%
\pgfpathlineto{\pgfqpoint{0.752010in}{0.522096in}}%
\pgfpathlineto{\pgfqpoint{0.748553in}{0.518499in}}%
\pgfpathlineto{\pgfqpoint{0.745838in}{0.518264in}}%
\pgfpathlineto{\pgfqpoint{0.744687in}{0.519881in}}%
\pgfpathlineto{\pgfqpoint{0.738995in}{0.520517in}}%
\pgfpathlineto{\pgfqpoint{0.733381in}{0.523327in}}%
\pgfpathlineto{\pgfqpoint{0.730861in}{0.525258in}}%
\pgfpathlineto{\pgfqpoint{0.728015in}{0.528863in}}%
\pgfpathlineto{\pgfqpoint{0.727312in}{0.530894in}}%
\pgfpathlineto{\pgfqpoint{0.728494in}{0.533400in}}%
\pgfpathlineto{\pgfqpoint{0.730236in}{0.534212in}}%
\pgfpathlineto{\pgfqpoint{0.729733in}{0.537152in}}%
\pgfpathlineto{\pgfqpoint{0.730007in}{0.540442in}}%
\pgfpathlineto{\pgfqpoint{0.731612in}{0.540941in}}%
\pgfpathlineto{\pgfqpoint{0.730213in}{0.545263in}}%
\pgfpathlineto{\pgfqpoint{0.731563in}{0.546123in}}%
\pgfpathlineto{\pgfqpoint{0.728919in}{0.548253in}}%
\pgfpathlineto{\pgfqpoint{0.728522in}{0.551246in}}%
\pgfpathlineto{\pgfqpoint{0.729262in}{0.553763in}}%
\pgfpathlineto{\pgfqpoint{0.731148in}{0.553340in}}%
\pgfpathlineto{\pgfqpoint{0.731806in}{0.555149in}}%
\pgfpathlineto{\pgfqpoint{0.729059in}{0.555932in}}%
\pgfpathlineto{\pgfqpoint{0.728087in}{0.558977in}}%
\pgfpathlineto{\pgfqpoint{0.731970in}{0.558050in}}%
\pgfpathlineto{\pgfqpoint{0.736704in}{0.560032in}}%
\pgfpathlineto{\pgfqpoint{0.739886in}{0.560095in}}%
\pgfpathlineto{\pgfqpoint{0.738328in}{0.563808in}}%
\pgfpathlineto{\pgfqpoint{0.739124in}{0.565657in}}%
\pgfpathlineto{\pgfqpoint{0.741979in}{0.566643in}}%
\pgfpathlineto{\pgfqpoint{0.742241in}{0.570962in}}%
\pgfpathlineto{\pgfqpoint{0.739760in}{0.569380in}}%
\pgfpathlineto{\pgfqpoint{0.738485in}{0.573179in}}%
\pgfpathlineto{\pgfqpoint{0.741686in}{0.576289in}}%
\pgfpathlineto{\pgfqpoint{0.739656in}{0.579992in}}%
\pgfpathlineto{\pgfqpoint{0.741135in}{0.581696in}}%
\pgfpathlineto{\pgfqpoint{0.743178in}{0.581393in}}%
\pgfpathlineto{\pgfqpoint{0.744280in}{0.583186in}}%
\pgfpathlineto{\pgfqpoint{0.743630in}{0.585155in}}%
\pgfpathlineto{\pgfqpoint{0.740303in}{0.585479in}}%
\pgfpathlineto{\pgfqpoint{0.741843in}{0.588299in}}%
\pgfpathlineto{\pgfqpoint{0.746374in}{0.589014in}}%
\pgfpathlineto{\pgfqpoint{0.747420in}{0.592192in}}%
\pgfpathlineto{\pgfqpoint{0.750741in}{0.589193in}}%
\pgfpathlineto{\pgfqpoint{0.753431in}{0.589913in}}%
\pgfpathlineto{\pgfqpoint{0.754707in}{0.592801in}}%
\pgfpathlineto{\pgfqpoint{0.757034in}{0.594176in}}%
\pgfpathlineto{\pgfqpoint{0.767817in}{0.596834in}}%
\pgfpathlineto{\pgfqpoint{0.772792in}{0.597266in}}%
\pgfpathlineto{\pgfqpoint{0.774461in}{0.595375in}}%
\pgfpathlineto{\pgfqpoint{0.777267in}{0.596507in}}%
\pgfpathlineto{\pgfqpoint{0.778048in}{0.599219in}}%
\pgfpathlineto{\pgfqpoint{0.783684in}{0.603781in}}%
\pgfpathlineto{\pgfqpoint{0.789079in}{0.605814in}}%
\pgfpathlineto{\pgfqpoint{0.794866in}{0.605498in}}%
\pgfpathlineto{\pgfqpoint{0.798355in}{0.602512in}}%
\pgfpathlineto{\pgfqpoint{0.799353in}{0.597971in}}%
\pgfpathlineto{\pgfqpoint{0.799751in}{0.593446in}}%
\pgfpathlineto{\pgfqpoint{0.801076in}{0.590791in}}%
\pgfpathlineto{\pgfqpoint{0.803852in}{0.589451in}}%
\pgfpathlineto{\pgfqpoint{0.808679in}{0.590732in}}%
\pgfpathlineto{\pgfqpoint{0.812082in}{0.590252in}}%
\pgfpathlineto{\pgfqpoint{0.814106in}{0.588275in}}%
\pgfpathlineto{\pgfqpoint{0.812396in}{0.586335in}}%
\pgfpathlineto{\pgfqpoint{0.814282in}{0.584612in}}%
\pgfpathlineto{\pgfqpoint{0.812542in}{0.582561in}}%
\pgfpathlineto{\pgfqpoint{0.814280in}{0.580959in}}%
\pgfpathlineto{\pgfqpoint{0.812726in}{0.579011in}}%
\pgfpathlineto{\pgfqpoint{0.824219in}{0.568707in}}%
\pgfpathlineto{\pgfqpoint{0.823025in}{0.569214in}}%
\pgfpathlineto{\pgfqpoint{0.820507in}{0.567425in}}%
\pgfpathlineto{\pgfqpoint{0.813592in}{0.559787in}}%
\pgfpathlineto{\pgfqpoint{0.813107in}{0.560166in}}%
\pgfpathlineto{\pgfqpoint{0.806268in}{0.552525in}}%
\pgfpathlineto{\pgfqpoint{0.807682in}{0.551257in}}%
\pgfpathlineto{\pgfqpoint{0.802699in}{0.545595in}}%
\pgfpathlineto{\pgfqpoint{0.812784in}{0.536599in}}%
\pgfpathlineto{\pgfqpoint{0.819842in}{0.530526in}}%
\pgfpathlineto{\pgfqpoint{0.821946in}{0.532976in}}%
\pgfpathlineto{\pgfqpoint{0.830011in}{0.526141in}}%
\pgfpathlineto{\pgfqpoint{0.831654in}{0.528086in}}%
\pgfpathlineto{\pgfqpoint{0.846928in}{0.515267in}}%
\pgfpathlineto{\pgfqpoint{0.845087in}{0.513012in}}%
\pgfpathlineto{\pgfqpoint{0.854046in}{0.505533in}}%
\pgfpathlineto{\pgfqpoint{0.867623in}{0.494883in}}%
\pgfpathlineto{\pgfqpoint{0.878680in}{0.486403in}}%
\pgfpathlineto{\pgfqpoint{0.880946in}{0.484915in}}%
\pgfpathlineto{\pgfqpoint{0.883274in}{0.487912in}}%
\pgfpathlineto{\pgfqpoint{0.887256in}{0.484755in}}%
\pgfpathlineto{\pgfqpoint{0.888822in}{0.486756in}}%
\pgfpathlineto{\pgfqpoint{0.892390in}{0.483966in}}%
\pgfpathlineto{\pgfqpoint{0.894650in}{0.485392in}}%
\pgfpathclose%
\pgfusepath{fill}%
\end{pgfscope}%
\begin{pgfscope}%
\pgfpathrectangle{\pgfqpoint{0.100000in}{0.100000in}}{\pgfqpoint{3.608454in}{2.310000in}}%
\pgfusepath{clip}%
\pgfsetbuttcap%
\pgfsetmiterjoin%
\definecolor{currentfill}{rgb}{0.000000,0.713725,0.643137}%
\pgfsetfillcolor{currentfill}%
\pgfsetlinewidth{0.000000pt}%
\definecolor{currentstroke}{rgb}{0.000000,0.000000,0.000000}%
\pgfsetstrokecolor{currentstroke}%
\pgfsetstrokeopacity{0.000000}%
\pgfsetdash{}{0pt}%
\pgfpathmoveto{\pgfqpoint{1.012082in}{1.949489in}}%
\pgfpathlineto{\pgfqpoint{1.009235in}{1.946422in}}%
\pgfpathlineto{\pgfqpoint{1.004445in}{1.925280in}}%
\pgfpathlineto{\pgfqpoint{0.989615in}{1.915702in}}%
\pgfpathlineto{\pgfqpoint{0.981235in}{1.920433in}}%
\pgfpathlineto{\pgfqpoint{0.976747in}{1.912867in}}%
\pgfpathlineto{\pgfqpoint{0.969772in}{1.911307in}}%
\pgfpathlineto{\pgfqpoint{0.967552in}{1.905305in}}%
\pgfpathlineto{\pgfqpoint{0.970494in}{1.901193in}}%
\pgfpathlineto{\pgfqpoint{0.969006in}{1.895451in}}%
\pgfpathlineto{\pgfqpoint{0.964170in}{1.890384in}}%
\pgfpathlineto{\pgfqpoint{0.953430in}{1.884971in}}%
\pgfpathlineto{\pgfqpoint{0.949827in}{1.885175in}}%
\pgfpathlineto{\pgfqpoint{0.917584in}{1.892808in}}%
\pgfpathlineto{\pgfqpoint{0.915455in}{1.886073in}}%
\pgfpathlineto{\pgfqpoint{0.909972in}{1.887658in}}%
\pgfpathlineto{\pgfqpoint{0.913150in}{1.901010in}}%
\pgfpathlineto{\pgfqpoint{0.916455in}{1.900234in}}%
\pgfpathlineto{\pgfqpoint{0.919645in}{1.913720in}}%
\pgfpathlineto{\pgfqpoint{0.910516in}{1.911417in}}%
\pgfpathlineto{\pgfqpoint{0.904574in}{1.912817in}}%
\pgfpathlineto{\pgfqpoint{0.898851in}{1.918021in}}%
\pgfpathlineto{\pgfqpoint{0.901098in}{1.927625in}}%
\pgfpathlineto{\pgfqpoint{0.896730in}{1.932620in}}%
\pgfpathlineto{\pgfqpoint{0.899991in}{1.946190in}}%
\pgfpathlineto{\pgfqpoint{0.884981in}{1.950005in}}%
\pgfpathlineto{\pgfqpoint{0.890375in}{1.956866in}}%
\pgfpathlineto{\pgfqpoint{0.891257in}{1.963356in}}%
\pgfpathlineto{\pgfqpoint{0.895815in}{1.966858in}}%
\pgfpathlineto{\pgfqpoint{0.899963in}{1.971006in}}%
\pgfpathlineto{\pgfqpoint{0.904646in}{1.980156in}}%
\pgfpathlineto{\pgfqpoint{0.923373in}{1.975654in}}%
\pgfpathlineto{\pgfqpoint{0.923969in}{1.964164in}}%
\pgfpathlineto{\pgfqpoint{0.931228in}{1.960671in}}%
\pgfpathlineto{\pgfqpoint{0.942731in}{1.965121in}}%
\pgfpathlineto{\pgfqpoint{1.012082in}{1.949489in}}%
\pgfpathclose%
\pgfusepath{fill}%
\end{pgfscope}%
\begin{pgfscope}%
\pgfpathrectangle{\pgfqpoint{0.100000in}{0.100000in}}{\pgfqpoint{3.608454in}{2.310000in}}%
\pgfusepath{clip}%
\pgfsetbuttcap%
\pgfsetmiterjoin%
\definecolor{currentfill}{rgb}{0.000000,0.709804,0.645098}%
\pgfsetfillcolor{currentfill}%
\pgfsetlinewidth{0.000000pt}%
\definecolor{currentstroke}{rgb}{0.000000,0.000000,0.000000}%
\pgfsetstrokecolor{currentstroke}%
\pgfsetstrokeopacity{0.000000}%
\pgfsetdash{}{0pt}%
\pgfpathmoveto{\pgfqpoint{2.477226in}{1.192084in}}%
\pgfpathlineto{\pgfqpoint{2.472760in}{1.194489in}}%
\pgfpathlineto{\pgfqpoint{2.467123in}{1.192232in}}%
\pgfpathlineto{\pgfqpoint{2.460932in}{1.196917in}}%
\pgfpathlineto{\pgfqpoint{2.451445in}{1.214650in}}%
\pgfpathlineto{\pgfqpoint{2.457038in}{1.223290in}}%
\pgfpathlineto{\pgfqpoint{2.449905in}{1.236650in}}%
\pgfpathlineto{\pgfqpoint{2.452467in}{1.239392in}}%
\pgfpathlineto{\pgfqpoint{2.449790in}{1.246442in}}%
\pgfpathlineto{\pgfqpoint{2.445368in}{1.248025in}}%
\pgfpathlineto{\pgfqpoint{2.434982in}{1.258000in}}%
\pgfpathlineto{\pgfqpoint{2.428315in}{1.262231in}}%
\pgfpathlineto{\pgfqpoint{2.421525in}{1.260176in}}%
\pgfpathlineto{\pgfqpoint{2.418250in}{1.266369in}}%
\pgfpathlineto{\pgfqpoint{2.409704in}{1.273645in}}%
\pgfpathlineto{\pgfqpoint{2.404880in}{1.275417in}}%
\pgfpathlineto{\pgfqpoint{2.415279in}{1.279749in}}%
\pgfpathlineto{\pgfqpoint{2.414877in}{1.286654in}}%
\pgfpathlineto{\pgfqpoint{2.435549in}{1.287521in}}%
\pgfpathlineto{\pgfqpoint{2.470034in}{1.289131in}}%
\pgfpathlineto{\pgfqpoint{2.470338in}{1.282183in}}%
\pgfpathlineto{\pgfqpoint{2.497937in}{1.284068in}}%
\pgfpathlineto{\pgfqpoint{2.499015in}{1.266823in}}%
\pgfpathlineto{\pgfqpoint{2.500587in}{1.242556in}}%
\pgfpathlineto{\pgfqpoint{2.501888in}{1.221865in}}%
\pgfpathlineto{\pgfqpoint{2.489480in}{1.220900in}}%
\pgfpathlineto{\pgfqpoint{2.488838in}{1.212212in}}%
\pgfpathlineto{\pgfqpoint{2.482512in}{1.210574in}}%
\pgfpathlineto{\pgfqpoint{2.473800in}{1.196498in}}%
\pgfpathlineto{\pgfqpoint{2.477226in}{1.192084in}}%
\pgfpathclose%
\pgfusepath{fill}%
\end{pgfscope}%
\begin{pgfscope}%
\pgfpathrectangle{\pgfqpoint{0.100000in}{0.100000in}}{\pgfqpoint{3.608454in}{2.310000in}}%
\pgfusepath{clip}%
\pgfsetbuttcap%
\pgfsetmiterjoin%
\definecolor{currentfill}{rgb}{0.000000,0.454902,0.772549}%
\pgfsetfillcolor{currentfill}%
\pgfsetlinewidth{0.000000pt}%
\definecolor{currentstroke}{rgb}{0.000000,0.000000,0.000000}%
\pgfsetstrokecolor{currentstroke}%
\pgfsetstrokeopacity{0.000000}%
\pgfsetdash{}{0pt}%
\pgfpathmoveto{\pgfqpoint{3.245459in}{1.515286in}}%
\pgfpathlineto{\pgfqpoint{3.239646in}{1.514041in}}%
\pgfpathlineto{\pgfqpoint{3.239245in}{1.513965in}}%
\pgfpathlineto{\pgfqpoint{3.227209in}{1.526852in}}%
\pgfpathlineto{\pgfqpoint{3.221777in}{1.528537in}}%
\pgfpathlineto{\pgfqpoint{3.216418in}{1.535851in}}%
\pgfpathlineto{\pgfqpoint{3.208552in}{1.535223in}}%
\pgfpathlineto{\pgfqpoint{3.204878in}{1.540940in}}%
\pgfpathlineto{\pgfqpoint{3.212062in}{1.546662in}}%
\pgfpathlineto{\pgfqpoint{3.201240in}{1.566916in}}%
\pgfpathlineto{\pgfqpoint{3.208359in}{1.574869in}}%
\pgfpathlineto{\pgfqpoint{3.196959in}{1.580835in}}%
\pgfpathlineto{\pgfqpoint{3.213947in}{1.593817in}}%
\pgfpathlineto{\pgfqpoint{3.217800in}{1.596744in}}%
\pgfpathlineto{\pgfqpoint{3.217800in}{1.603350in}}%
\pgfpathlineto{\pgfqpoint{3.221218in}{1.609411in}}%
\pgfpathlineto{\pgfqpoint{3.234122in}{1.609162in}}%
\pgfpathlineto{\pgfqpoint{3.251074in}{1.598426in}}%
\pgfpathlineto{\pgfqpoint{3.244201in}{1.592250in}}%
\pgfpathlineto{\pgfqpoint{3.269282in}{1.579037in}}%
\pgfpathlineto{\pgfqpoint{3.262873in}{1.561037in}}%
\pgfpathlineto{\pgfqpoint{3.250514in}{1.547606in}}%
\pgfpathlineto{\pgfqpoint{3.252158in}{1.542409in}}%
\pgfpathlineto{\pgfqpoint{3.250059in}{1.534993in}}%
\pgfpathlineto{\pgfqpoint{3.250934in}{1.527034in}}%
\pgfpathlineto{\pgfqpoint{3.248921in}{1.519116in}}%
\pgfpathlineto{\pgfqpoint{3.245459in}{1.515286in}}%
\pgfpathclose%
\pgfusepath{fill}%
\end{pgfscope}%
\begin{pgfscope}%
\pgfpathrectangle{\pgfqpoint{0.100000in}{0.100000in}}{\pgfqpoint{3.608454in}{2.310000in}}%
\pgfusepath{clip}%
\pgfsetbuttcap%
\pgfsetmiterjoin%
\definecolor{currentfill}{rgb}{0.000000,0.592157,0.703922}%
\pgfsetfillcolor{currentfill}%
\pgfsetlinewidth{0.000000pt}%
\definecolor{currentstroke}{rgb}{0.000000,0.000000,0.000000}%
\pgfsetstrokecolor{currentstroke}%
\pgfsetstrokeopacity{0.000000}%
\pgfsetdash{}{0pt}%
\pgfpathmoveto{\pgfqpoint{2.887245in}{0.726366in}}%
\pgfpathlineto{\pgfqpoint{2.878161in}{0.725802in}}%
\pgfpathlineto{\pgfqpoint{2.874929in}{0.755916in}}%
\pgfpathlineto{\pgfqpoint{2.856733in}{0.753994in}}%
\pgfpathlineto{\pgfqpoint{2.856403in}{0.756814in}}%
\pgfpathlineto{\pgfqpoint{2.854226in}{0.777256in}}%
\pgfpathlineto{\pgfqpoint{2.877599in}{0.779820in}}%
\pgfpathlineto{\pgfqpoint{2.875554in}{0.798633in}}%
\pgfpathlineto{\pgfqpoint{2.870385in}{0.802308in}}%
\pgfpathlineto{\pgfqpoint{2.865088in}{0.801722in}}%
\pgfpathlineto{\pgfqpoint{2.862925in}{0.815951in}}%
\pgfpathlineto{\pgfqpoint{2.875686in}{0.817596in}}%
\pgfpathlineto{\pgfqpoint{2.875215in}{0.821548in}}%
\pgfpathlineto{\pgfqpoint{2.884148in}{0.822135in}}%
\pgfpathlineto{\pgfqpoint{2.887379in}{0.811341in}}%
\pgfpathlineto{\pgfqpoint{2.884962in}{0.810094in}}%
\pgfpathlineto{\pgfqpoint{2.885535in}{0.800050in}}%
\pgfpathlineto{\pgfqpoint{2.891131in}{0.799434in}}%
\pgfpathlineto{\pgfqpoint{2.894796in}{0.794134in}}%
\pgfpathlineto{\pgfqpoint{2.910462in}{0.795395in}}%
\pgfpathlineto{\pgfqpoint{2.911798in}{0.789691in}}%
\pgfpathlineto{\pgfqpoint{2.920094in}{0.780812in}}%
\pgfpathlineto{\pgfqpoint{2.920160in}{0.773477in}}%
\pgfpathlineto{\pgfqpoint{2.925268in}{0.774204in}}%
\pgfpathlineto{\pgfqpoint{2.928578in}{0.749275in}}%
\pgfpathlineto{\pgfqpoint{2.937427in}{0.747697in}}%
\pgfpathlineto{\pgfqpoint{2.943923in}{0.740325in}}%
\pgfpathlineto{\pgfqpoint{2.955222in}{0.740323in}}%
\pgfpathlineto{\pgfqpoint{2.957900in}{0.730883in}}%
\pgfpathlineto{\pgfqpoint{2.887245in}{0.726366in}}%
\pgfpathclose%
\pgfusepath{fill}%
\end{pgfscope}%
\begin{pgfscope}%
\pgfpathrectangle{\pgfqpoint{0.100000in}{0.100000in}}{\pgfqpoint{3.608454in}{2.310000in}}%
\pgfusepath{clip}%
\pgfsetbuttcap%
\pgfsetmiterjoin%
\definecolor{currentfill}{rgb}{0.000000,0.827451,0.586275}%
\pgfsetfillcolor{currentfill}%
\pgfsetlinewidth{0.000000pt}%
\definecolor{currentstroke}{rgb}{0.000000,0.000000,0.000000}%
\pgfsetstrokecolor{currentstroke}%
\pgfsetstrokeopacity{0.000000}%
\pgfsetdash{}{0pt}%
\pgfpathmoveto{\pgfqpoint{2.416578in}{0.896534in}}%
\pgfpathlineto{\pgfqpoint{2.408616in}{0.893304in}}%
\pgfpathlineto{\pgfqpoint{2.408409in}{0.898240in}}%
\pgfpathlineto{\pgfqpoint{2.406598in}{0.936460in}}%
\pgfpathlineto{\pgfqpoint{2.424326in}{0.937210in}}%
\pgfpathlineto{\pgfqpoint{2.416509in}{0.931498in}}%
\pgfpathlineto{\pgfqpoint{2.427713in}{0.930421in}}%
\pgfpathlineto{\pgfqpoint{2.427361in}{0.937371in}}%
\pgfpathlineto{\pgfqpoint{2.441137in}{0.939280in}}%
\pgfpathlineto{\pgfqpoint{2.440822in}{0.945153in}}%
\pgfpathlineto{\pgfqpoint{2.451281in}{0.943313in}}%
\pgfpathlineto{\pgfqpoint{2.468504in}{0.944341in}}%
\pgfpathlineto{\pgfqpoint{2.469213in}{0.932802in}}%
\pgfpathlineto{\pgfqpoint{2.470515in}{0.912039in}}%
\pgfpathlineto{\pgfqpoint{2.474111in}{0.910147in}}%
\pgfpathlineto{\pgfqpoint{2.474841in}{0.898468in}}%
\pgfpathlineto{\pgfqpoint{2.458419in}{0.897453in}}%
\pgfpathlineto{\pgfqpoint{2.455570in}{0.891636in}}%
\pgfpathlineto{\pgfqpoint{2.426856in}{0.899501in}}%
\pgfpathlineto{\pgfqpoint{2.419871in}{0.899622in}}%
\pgfpathlineto{\pgfqpoint{2.416578in}{0.896534in}}%
\pgfpathclose%
\pgfusepath{fill}%
\end{pgfscope}%
\begin{pgfscope}%
\pgfpathrectangle{\pgfqpoint{0.100000in}{0.100000in}}{\pgfqpoint{3.608454in}{2.310000in}}%
\pgfusepath{clip}%
\pgfsetbuttcap%
\pgfsetmiterjoin%
\definecolor{currentfill}{rgb}{0.000000,0.333333,0.833333}%
\pgfsetfillcolor{currentfill}%
\pgfsetlinewidth{0.000000pt}%
\definecolor{currentstroke}{rgb}{0.000000,0.000000,0.000000}%
\pgfsetstrokecolor{currentstroke}%
\pgfsetstrokeopacity{0.000000}%
\pgfsetdash{}{0pt}%
\pgfpathmoveto{\pgfqpoint{1.913591in}{0.794961in}}%
\pgfpathlineto{\pgfqpoint{1.908521in}{0.806762in}}%
\pgfpathlineto{\pgfqpoint{1.900998in}{0.829156in}}%
\pgfpathlineto{\pgfqpoint{1.873743in}{0.830250in}}%
\pgfpathlineto{\pgfqpoint{1.866975in}{0.830682in}}%
\pgfpathlineto{\pgfqpoint{1.868391in}{0.865576in}}%
\pgfpathlineto{\pgfqpoint{1.878275in}{0.865031in}}%
\pgfpathlineto{\pgfqpoint{1.878486in}{0.869447in}}%
\pgfpathlineto{\pgfqpoint{1.912098in}{0.867745in}}%
\pgfpathlineto{\pgfqpoint{1.913431in}{0.902088in}}%
\pgfpathlineto{\pgfqpoint{1.909495in}{0.904858in}}%
\pgfpathlineto{\pgfqpoint{1.910664in}{0.938402in}}%
\pgfpathlineto{\pgfqpoint{1.917943in}{0.934949in}}%
\pgfpathlineto{\pgfqpoint{1.931087in}{0.945835in}}%
\pgfpathlineto{\pgfqpoint{1.938143in}{0.938336in}}%
\pgfpathlineto{\pgfqpoint{1.943227in}{0.939444in}}%
\pgfpathlineto{\pgfqpoint{1.941990in}{0.901235in}}%
\pgfpathlineto{\pgfqpoint{1.948881in}{0.900975in}}%
\pgfpathlineto{\pgfqpoint{1.946927in}{0.865843in}}%
\pgfpathlineto{\pgfqpoint{1.971315in}{0.864806in}}%
\pgfpathlineto{\pgfqpoint{1.970138in}{0.830158in}}%
\pgfpathlineto{\pgfqpoint{1.966891in}{0.830295in}}%
\pgfpathlineto{\pgfqpoint{1.966416in}{0.807704in}}%
\pgfpathlineto{\pgfqpoint{1.940012in}{0.801698in}}%
\pgfpathlineto{\pgfqpoint{1.934684in}{0.799256in}}%
\pgfpathlineto{\pgfqpoint{1.929646in}{0.804003in}}%
\pgfpathlineto{\pgfqpoint{1.913591in}{0.794961in}}%
\pgfpathclose%
\pgfusepath{fill}%
\end{pgfscope}%
\begin{pgfscope}%
\pgfpathrectangle{\pgfqpoint{0.100000in}{0.100000in}}{\pgfqpoint{3.608454in}{2.310000in}}%
\pgfusepath{clip}%
\pgfsetbuttcap%
\pgfsetmiterjoin%
\definecolor{currentfill}{rgb}{0.000000,0.454902,0.772549}%
\pgfsetfillcolor{currentfill}%
\pgfsetlinewidth{0.000000pt}%
\definecolor{currentstroke}{rgb}{0.000000,0.000000,0.000000}%
\pgfsetstrokecolor{currentstroke}%
\pgfsetstrokeopacity{0.000000}%
\pgfsetdash{}{0pt}%
\pgfpathmoveto{\pgfqpoint{2.698781in}{1.052062in}}%
\pgfpathlineto{\pgfqpoint{2.669947in}{1.049451in}}%
\pgfpathlineto{\pgfqpoint{2.668473in}{1.060141in}}%
\pgfpathlineto{\pgfqpoint{2.665061in}{1.060409in}}%
\pgfpathlineto{\pgfqpoint{2.662888in}{1.068219in}}%
\pgfpathlineto{\pgfqpoint{2.648588in}{1.077186in}}%
\pgfpathlineto{\pgfqpoint{2.644297in}{1.082383in}}%
\pgfpathlineto{\pgfqpoint{2.643683in}{1.102257in}}%
\pgfpathlineto{\pgfqpoint{2.651600in}{1.103191in}}%
\pgfpathlineto{\pgfqpoint{2.654621in}{1.100124in}}%
\pgfpathlineto{\pgfqpoint{2.668887in}{1.104876in}}%
\pgfpathlineto{\pgfqpoint{2.681536in}{1.102539in}}%
\pgfpathlineto{\pgfqpoint{2.685775in}{1.104241in}}%
\pgfpathlineto{\pgfqpoint{2.688788in}{1.097666in}}%
\pgfpathlineto{\pgfqpoint{2.693772in}{1.094220in}}%
\pgfpathlineto{\pgfqpoint{2.694900in}{1.080379in}}%
\pgfpathlineto{\pgfqpoint{2.693282in}{1.075443in}}%
\pgfpathlineto{\pgfqpoint{2.696968in}{1.070359in}}%
\pgfpathlineto{\pgfqpoint{2.698781in}{1.052062in}}%
\pgfpathclose%
\pgfusepath{fill}%
\end{pgfscope}%
\begin{pgfscope}%
\pgfpathrectangle{\pgfqpoint{0.100000in}{0.100000in}}{\pgfqpoint{3.608454in}{2.310000in}}%
\pgfusepath{clip}%
\pgfsetbuttcap%
\pgfsetmiterjoin%
\definecolor{currentfill}{rgb}{0.000000,0.647059,0.676471}%
\pgfsetfillcolor{currentfill}%
\pgfsetlinewidth{0.000000pt}%
\definecolor{currentstroke}{rgb}{0.000000,0.000000,0.000000}%
\pgfsetstrokecolor{currentstroke}%
\pgfsetstrokeopacity{0.000000}%
\pgfsetdash{}{0pt}%
\pgfpathmoveto{\pgfqpoint{2.550184in}{0.875112in}}%
\pgfpathlineto{\pgfqpoint{2.519263in}{0.872679in}}%
\pgfpathlineto{\pgfqpoint{2.517297in}{0.901251in}}%
\pgfpathlineto{\pgfqpoint{2.526832in}{0.901949in}}%
\pgfpathlineto{\pgfqpoint{2.525476in}{0.919309in}}%
\pgfpathlineto{\pgfqpoint{2.534527in}{0.922423in}}%
\pgfpathlineto{\pgfqpoint{2.537884in}{0.927921in}}%
\pgfpathlineto{\pgfqpoint{2.535053in}{0.931443in}}%
\pgfpathlineto{\pgfqpoint{2.544297in}{0.934956in}}%
\pgfpathlineto{\pgfqpoint{2.546945in}{0.939900in}}%
\pgfpathlineto{\pgfqpoint{2.551959in}{0.940228in}}%
\pgfpathlineto{\pgfqpoint{2.551506in}{0.923399in}}%
\pgfpathlineto{\pgfqpoint{2.580198in}{0.924865in}}%
\pgfpathlineto{\pgfqpoint{2.582729in}{0.895488in}}%
\pgfpathlineto{\pgfqpoint{2.575754in}{0.891315in}}%
\pgfpathlineto{\pgfqpoint{2.570549in}{0.885685in}}%
\pgfpathlineto{\pgfqpoint{2.564922in}{0.885225in}}%
\pgfpathlineto{\pgfqpoint{2.561497in}{0.881275in}}%
\pgfpathlineto{\pgfqpoint{2.550318in}{0.880073in}}%
\pgfpathlineto{\pgfqpoint{2.550184in}{0.875112in}}%
\pgfpathclose%
\pgfusepath{fill}%
\end{pgfscope}%
\begin{pgfscope}%
\pgfpathrectangle{\pgfqpoint{0.100000in}{0.100000in}}{\pgfqpoint{3.608454in}{2.310000in}}%
\pgfusepath{clip}%
\pgfsetbuttcap%
\pgfsetmiterjoin%
\definecolor{currentfill}{rgb}{0.000000,0.725490,0.637255}%
\pgfsetfillcolor{currentfill}%
\pgfsetlinewidth{0.000000pt}%
\definecolor{currentstroke}{rgb}{0.000000,0.000000,0.000000}%
\pgfsetstrokecolor{currentstroke}%
\pgfsetstrokeopacity{0.000000}%
\pgfsetdash{}{0pt}%
\pgfpathmoveto{\pgfqpoint{3.215154in}{1.664914in}}%
\pgfpathlineto{\pgfqpoint{3.207689in}{1.691369in}}%
\pgfpathlineto{\pgfqpoint{3.183986in}{1.686549in}}%
\pgfpathlineto{\pgfqpoint{3.180627in}{1.708614in}}%
\pgfpathlineto{\pgfqpoint{3.187096in}{1.712911in}}%
\pgfpathlineto{\pgfqpoint{3.196916in}{1.713203in}}%
\pgfpathlineto{\pgfqpoint{3.194930in}{1.721757in}}%
\pgfpathlineto{\pgfqpoint{3.217063in}{1.727096in}}%
\pgfpathlineto{\pgfqpoint{3.220803in}{1.715267in}}%
\pgfpathlineto{\pgfqpoint{3.232758in}{1.716945in}}%
\pgfpathlineto{\pgfqpoint{3.233713in}{1.712856in}}%
\pgfpathlineto{\pgfqpoint{3.246255in}{1.715536in}}%
\pgfpathlineto{\pgfqpoint{3.248659in}{1.703631in}}%
\pgfpathlineto{\pgfqpoint{3.252945in}{1.701154in}}%
\pgfpathlineto{\pgfqpoint{3.245833in}{1.699582in}}%
\pgfpathlineto{\pgfqpoint{3.253079in}{1.672195in}}%
\pgfpathlineto{\pgfqpoint{3.215154in}{1.664914in}}%
\pgfpathclose%
\pgfusepath{fill}%
\end{pgfscope}%
\begin{pgfscope}%
\pgfpathrectangle{\pgfqpoint{0.100000in}{0.100000in}}{\pgfqpoint{3.608454in}{2.310000in}}%
\pgfusepath{clip}%
\pgfsetbuttcap%
\pgfsetmiterjoin%
\definecolor{currentfill}{rgb}{0.000000,0.521569,0.739216}%
\pgfsetfillcolor{currentfill}%
\pgfsetlinewidth{0.000000pt}%
\definecolor{currentstroke}{rgb}{0.000000,0.000000,0.000000}%
\pgfsetstrokecolor{currentstroke}%
\pgfsetstrokeopacity{0.000000}%
\pgfsetdash{}{0pt}%
\pgfpathmoveto{\pgfqpoint{2.020851in}{0.893088in}}%
\pgfpathlineto{\pgfqpoint{2.014996in}{0.896820in}}%
\pgfpathlineto{\pgfqpoint{1.984737in}{0.897974in}}%
\pgfpathlineto{\pgfqpoint{1.977949in}{0.898974in}}%
\pgfpathlineto{\pgfqpoint{1.978804in}{0.941263in}}%
\pgfpathlineto{\pgfqpoint{1.982724in}{0.934735in}}%
\pgfpathlineto{\pgfqpoint{1.993908in}{0.931626in}}%
\pgfpathlineto{\pgfqpoint{2.000025in}{0.945069in}}%
\pgfpathlineto{\pgfqpoint{2.014398in}{0.957120in}}%
\pgfpathlineto{\pgfqpoint{2.041706in}{0.956720in}}%
\pgfpathlineto{\pgfqpoint{2.045185in}{0.954206in}}%
\pgfpathlineto{\pgfqpoint{2.045055in}{0.940136in}}%
\pgfpathlineto{\pgfqpoint{2.051407in}{0.934194in}}%
\pgfpathlineto{\pgfqpoint{2.056583in}{0.934051in}}%
\pgfpathlineto{\pgfqpoint{2.050404in}{0.930484in}}%
\pgfpathlineto{\pgfqpoint{2.049865in}{0.897339in}}%
\pgfpathlineto{\pgfqpoint{2.020851in}{0.893088in}}%
\pgfpathclose%
\pgfusepath{fill}%
\end{pgfscope}%
\begin{pgfscope}%
\pgfpathrectangle{\pgfqpoint{0.100000in}{0.100000in}}{\pgfqpoint{3.608454in}{2.310000in}}%
\pgfusepath{clip}%
\pgfsetbuttcap%
\pgfsetmiterjoin%
\definecolor{currentfill}{rgb}{0.000000,0.564706,0.717647}%
\pgfsetfillcolor{currentfill}%
\pgfsetlinewidth{0.000000pt}%
\definecolor{currentstroke}{rgb}{0.000000,0.000000,0.000000}%
\pgfsetstrokecolor{currentstroke}%
\pgfsetstrokeopacity{0.000000}%
\pgfsetdash{}{0pt}%
\pgfpathmoveto{\pgfqpoint{2.561760in}{0.672494in}}%
\pgfpathlineto{\pgfqpoint{2.556138in}{0.721372in}}%
\pgfpathlineto{\pgfqpoint{2.551986in}{0.756063in}}%
\pgfpathlineto{\pgfqpoint{2.549469in}{0.776821in}}%
\pgfpathlineto{\pgfqpoint{2.574889in}{0.778819in}}%
\pgfpathlineto{\pgfqpoint{2.573394in}{0.783619in}}%
\pgfpathlineto{\pgfqpoint{2.567924in}{0.787803in}}%
\pgfpathlineto{\pgfqpoint{2.574126in}{0.801999in}}%
\pgfpathlineto{\pgfqpoint{2.601434in}{0.804239in}}%
\pgfpathlineto{\pgfqpoint{2.604373in}{0.805606in}}%
\pgfpathlineto{\pgfqpoint{2.603471in}{0.815667in}}%
\pgfpathlineto{\pgfqpoint{2.610325in}{0.816266in}}%
\pgfpathlineto{\pgfqpoint{2.609428in}{0.825937in}}%
\pgfpathlineto{\pgfqpoint{2.612796in}{0.827022in}}%
\pgfpathlineto{\pgfqpoint{2.622706in}{0.822514in}}%
\pgfpathlineto{\pgfqpoint{2.634007in}{0.811449in}}%
\pgfpathlineto{\pgfqpoint{2.652230in}{0.813085in}}%
\pgfpathlineto{\pgfqpoint{2.653773in}{0.795845in}}%
\pgfpathlineto{\pgfqpoint{2.613728in}{0.792291in}}%
\pgfpathlineto{\pgfqpoint{2.613565in}{0.781771in}}%
\pgfpathlineto{\pgfqpoint{2.610143in}{0.781433in}}%
\pgfpathlineto{\pgfqpoint{2.611445in}{0.760754in}}%
\pgfpathlineto{\pgfqpoint{2.605232in}{0.759125in}}%
\pgfpathlineto{\pgfqpoint{2.601460in}{0.754867in}}%
\pgfpathlineto{\pgfqpoint{2.602701in}{0.749237in}}%
\pgfpathlineto{\pgfqpoint{2.609757in}{0.745194in}}%
\pgfpathlineto{\pgfqpoint{2.612513in}{0.725646in}}%
\pgfpathlineto{\pgfqpoint{2.610908in}{0.714979in}}%
\pgfpathlineto{\pgfqpoint{2.618895in}{0.705619in}}%
\pgfpathlineto{\pgfqpoint{2.627748in}{0.701087in}}%
\pgfpathlineto{\pgfqpoint{2.628964in}{0.696389in}}%
\pgfpathlineto{\pgfqpoint{2.626017in}{0.689071in}}%
\pgfpathlineto{\pgfqpoint{2.632066in}{0.682154in}}%
\pgfpathlineto{\pgfqpoint{2.627322in}{0.678308in}}%
\pgfpathlineto{\pgfqpoint{2.622714in}{0.669017in}}%
\pgfpathlineto{\pgfqpoint{2.613254in}{0.665766in}}%
\pgfpathlineto{\pgfqpoint{2.606414in}{0.667773in}}%
\pgfpathlineto{\pgfqpoint{2.595006in}{0.677521in}}%
\pgfpathlineto{\pgfqpoint{2.592621in}{0.682803in}}%
\pgfpathlineto{\pgfqpoint{2.594483in}{0.687365in}}%
\pgfpathlineto{\pgfqpoint{2.592065in}{0.696291in}}%
\pgfpathlineto{\pgfqpoint{2.588127in}{0.698733in}}%
\pgfpathlineto{\pgfqpoint{2.583318in}{0.695147in}}%
\pgfpathlineto{\pgfqpoint{2.581260in}{0.684211in}}%
\pgfpathlineto{\pgfqpoint{2.581606in}{0.673454in}}%
\pgfpathlineto{\pgfqpoint{2.580066in}{0.669312in}}%
\pgfpathlineto{\pgfqpoint{2.564924in}{0.674357in}}%
\pgfpathlineto{\pgfqpoint{2.561760in}{0.672494in}}%
\pgfpathclose%
\pgfusepath{fill}%
\end{pgfscope}%
\begin{pgfscope}%
\pgfpathrectangle{\pgfqpoint{0.100000in}{0.100000in}}{\pgfqpoint{3.608454in}{2.310000in}}%
\pgfusepath{clip}%
\pgfsetbuttcap%
\pgfsetmiterjoin%
\definecolor{currentfill}{rgb}{0.000000,0.545098,0.727451}%
\pgfsetfillcolor{currentfill}%
\pgfsetlinewidth{0.000000pt}%
\definecolor{currentstroke}{rgb}{0.000000,0.000000,0.000000}%
\pgfsetstrokecolor{currentstroke}%
\pgfsetstrokeopacity{0.000000}%
\pgfsetdash{}{0pt}%
\pgfpathmoveto{\pgfqpoint{1.440664in}{1.064305in}}%
\pgfpathlineto{\pgfqpoint{1.436822in}{1.030070in}}%
\pgfpathlineto{\pgfqpoint{1.435358in}{1.030206in}}%
\pgfpathlineto{\pgfqpoint{1.432998in}{1.009699in}}%
\pgfpathlineto{\pgfqpoint{1.432205in}{1.002872in}}%
\pgfpathlineto{\pgfqpoint{1.392442in}{1.007538in}}%
\pgfpathlineto{\pgfqpoint{1.360437in}{1.011466in}}%
\pgfpathlineto{\pgfqpoint{1.362154in}{1.025644in}}%
\pgfpathlineto{\pgfqpoint{1.340304in}{1.035758in}}%
\pgfpathlineto{\pgfqpoint{1.333401in}{1.040515in}}%
\pgfpathlineto{\pgfqpoint{1.278910in}{1.047826in}}%
\pgfpathlineto{\pgfqpoint{1.236876in}{1.053946in}}%
\pgfpathlineto{\pgfqpoint{1.193614in}{1.060581in}}%
\pgfpathlineto{\pgfqpoint{1.198276in}{1.090371in}}%
\pgfpathlineto{\pgfqpoint{1.235282in}{1.084666in}}%
\pgfpathlineto{\pgfqpoint{1.239395in}{1.111911in}}%
\pgfpathlineto{\pgfqpoint{1.313512in}{1.101203in}}%
\pgfpathlineto{\pgfqpoint{1.321043in}{1.155597in}}%
\pgfpathlineto{\pgfqpoint{1.300870in}{1.158512in}}%
\pgfpathlineto{\pgfqpoint{1.303419in}{1.175695in}}%
\pgfpathlineto{\pgfqpoint{1.350154in}{1.169097in}}%
\pgfpathlineto{\pgfqpoint{1.347890in}{1.152024in}}%
\pgfpathlineto{\pgfqpoint{1.388440in}{1.147038in}}%
\pgfpathlineto{\pgfqpoint{1.378347in}{1.144995in}}%
\pgfpathlineto{\pgfqpoint{1.377042in}{1.134592in}}%
\pgfpathlineto{\pgfqpoint{1.385971in}{1.127718in}}%
\pgfpathlineto{\pgfqpoint{1.379252in}{1.071453in}}%
\pgfpathlineto{\pgfqpoint{1.413364in}{1.067432in}}%
\pgfpathlineto{\pgfqpoint{1.440664in}{1.064305in}}%
\pgfpathclose%
\pgfusepath{fill}%
\end{pgfscope}%
\begin{pgfscope}%
\pgfpathrectangle{\pgfqpoint{0.100000in}{0.100000in}}{\pgfqpoint{3.608454in}{2.310000in}}%
\pgfusepath{clip}%
\pgfsetbuttcap%
\pgfsetmiterjoin%
\definecolor{currentfill}{rgb}{0.000000,0.627451,0.686275}%
\pgfsetfillcolor{currentfill}%
\pgfsetlinewidth{0.000000pt}%
\definecolor{currentstroke}{rgb}{0.000000,0.000000,0.000000}%
\pgfsetstrokecolor{currentstroke}%
\pgfsetstrokeopacity{0.000000}%
\pgfsetdash{}{0pt}%
\pgfpathmoveto{\pgfqpoint{1.993561in}{0.650091in}}%
\pgfpathlineto{\pgfqpoint{1.986455in}{0.652084in}}%
\pgfpathlineto{\pgfqpoint{1.978079in}{0.661346in}}%
\pgfpathlineto{\pgfqpoint{1.975419in}{0.668492in}}%
\pgfpathlineto{\pgfqpoint{1.971709in}{0.671338in}}%
\pgfpathlineto{\pgfqpoint{1.995500in}{0.684736in}}%
\pgfpathlineto{\pgfqpoint{1.993815in}{0.692311in}}%
\pgfpathlineto{\pgfqpoint{1.986509in}{0.700655in}}%
\pgfpathlineto{\pgfqpoint{1.984720in}{0.711352in}}%
\pgfpathlineto{\pgfqpoint{1.981851in}{0.714949in}}%
\pgfpathlineto{\pgfqpoint{1.997875in}{0.723825in}}%
\pgfpathlineto{\pgfqpoint{2.016972in}{0.734385in}}%
\pgfpathlineto{\pgfqpoint{2.016298in}{0.728776in}}%
\pgfpathlineto{\pgfqpoint{2.021875in}{0.703710in}}%
\pgfpathlineto{\pgfqpoint{2.026699in}{0.691630in}}%
\pgfpathlineto{\pgfqpoint{2.047653in}{0.694679in}}%
\pgfpathlineto{\pgfqpoint{2.049731in}{0.675985in}}%
\pgfpathlineto{\pgfqpoint{2.051242in}{0.645234in}}%
\pgfpathlineto{\pgfqpoint{2.031567in}{0.644083in}}%
\pgfpathlineto{\pgfqpoint{2.031314in}{0.647827in}}%
\pgfpathlineto{\pgfqpoint{2.025563in}{0.657506in}}%
\pgfpathlineto{\pgfqpoint{2.016277in}{0.656724in}}%
\pgfpathlineto{\pgfqpoint{2.008653in}{0.653539in}}%
\pgfpathlineto{\pgfqpoint{1.993561in}{0.650091in}}%
\pgfpathclose%
\pgfusepath{fill}%
\end{pgfscope}%
\begin{pgfscope}%
\pgfpathrectangle{\pgfqpoint{0.100000in}{0.100000in}}{\pgfqpoint{3.608454in}{2.310000in}}%
\pgfusepath{clip}%
\pgfsetbuttcap%
\pgfsetmiterjoin%
\definecolor{currentfill}{rgb}{0.000000,0.850980,0.574510}%
\pgfsetfillcolor{currentfill}%
\pgfsetlinewidth{0.000000pt}%
\definecolor{currentstroke}{rgb}{0.000000,0.000000,0.000000}%
\pgfsetstrokecolor{currentstroke}%
\pgfsetstrokeopacity{0.000000}%
\pgfsetdash{}{0pt}%
\pgfpathmoveto{\pgfqpoint{2.991086in}{1.284432in}}%
\pgfpathlineto{\pgfqpoint{2.974604in}{1.294585in}}%
\pgfpathlineto{\pgfqpoint{2.968540in}{1.292624in}}%
\pgfpathlineto{\pgfqpoint{2.967623in}{1.287269in}}%
\pgfpathlineto{\pgfqpoint{2.961369in}{1.291611in}}%
\pgfpathlineto{\pgfqpoint{2.952046in}{1.304075in}}%
\pgfpathlineto{\pgfqpoint{2.946226in}{1.306488in}}%
\pgfpathlineto{\pgfqpoint{2.941096in}{1.316634in}}%
\pgfpathlineto{\pgfqpoint{2.944352in}{1.321971in}}%
\pgfpathlineto{\pgfqpoint{2.952368in}{1.321633in}}%
\pgfpathlineto{\pgfqpoint{2.953019in}{1.337104in}}%
\pgfpathlineto{\pgfqpoint{2.958041in}{1.346022in}}%
\pgfpathlineto{\pgfqpoint{2.968495in}{1.342247in}}%
\pgfpathlineto{\pgfqpoint{2.973652in}{1.345179in}}%
\pgfpathlineto{\pgfqpoint{2.981470in}{1.336849in}}%
\pgfpathlineto{\pgfqpoint{2.992200in}{1.336598in}}%
\pgfpathlineto{\pgfqpoint{3.006750in}{1.353994in}}%
\pgfpathlineto{\pgfqpoint{3.011961in}{1.351516in}}%
\pgfpathlineto{\pgfqpoint{3.013175in}{1.342782in}}%
\pgfpathlineto{\pgfqpoint{3.020093in}{1.338488in}}%
\pgfpathlineto{\pgfqpoint{3.026105in}{1.338632in}}%
\pgfpathlineto{\pgfqpoint{3.038592in}{1.342913in}}%
\pgfpathlineto{\pgfqpoint{3.037059in}{1.337144in}}%
\pgfpathlineto{\pgfqpoint{3.025954in}{1.320474in}}%
\pgfpathlineto{\pgfqpoint{3.022678in}{1.310484in}}%
\pgfpathlineto{\pgfqpoint{3.013170in}{1.308307in}}%
\pgfpathlineto{\pgfqpoint{3.000083in}{1.309688in}}%
\pgfpathlineto{\pgfqpoint{2.991086in}{1.284432in}}%
\pgfpathclose%
\pgfusepath{fill}%
\end{pgfscope}%
\begin{pgfscope}%
\pgfpathrectangle{\pgfqpoint{0.100000in}{0.100000in}}{\pgfqpoint{3.608454in}{2.310000in}}%
\pgfusepath{clip}%
\pgfsetbuttcap%
\pgfsetmiterjoin%
\definecolor{currentfill}{rgb}{0.000000,0.592157,0.703922}%
\pgfsetfillcolor{currentfill}%
\pgfsetlinewidth{0.000000pt}%
\definecolor{currentstroke}{rgb}{0.000000,0.000000,0.000000}%
\pgfsetstrokecolor{currentstroke}%
\pgfsetstrokeopacity{0.000000}%
\pgfsetdash{}{0pt}%
\pgfpathmoveto{\pgfqpoint{3.302028in}{1.725960in}}%
\pgfpathlineto{\pgfqpoint{3.267733in}{1.692380in}}%
\pgfpathlineto{\pgfqpoint{3.260640in}{1.692420in}}%
\pgfpathlineto{\pgfqpoint{3.258600in}{1.697493in}}%
\pgfpathlineto{\pgfqpoint{3.252945in}{1.701154in}}%
\pgfpathlineto{\pgfqpoint{3.248659in}{1.703631in}}%
\pgfpathlineto{\pgfqpoint{3.246255in}{1.715536in}}%
\pgfpathlineto{\pgfqpoint{3.233713in}{1.712856in}}%
\pgfpathlineto{\pgfqpoint{3.232758in}{1.716945in}}%
\pgfpathlineto{\pgfqpoint{3.220803in}{1.715267in}}%
\pgfpathlineto{\pgfqpoint{3.217063in}{1.727096in}}%
\pgfpathlineto{\pgfqpoint{3.210494in}{1.750645in}}%
\pgfpathlineto{\pgfqpoint{3.243936in}{1.759541in}}%
\pgfpathlineto{\pgfqpoint{3.246838in}{1.769301in}}%
\pgfpathlineto{\pgfqpoint{3.252161in}{1.774686in}}%
\pgfpathlineto{\pgfqpoint{3.264610in}{1.770699in}}%
\pgfpathlineto{\pgfqpoint{3.264917in}{1.776723in}}%
\pgfpathlineto{\pgfqpoint{3.272069in}{1.775491in}}%
\pgfpathlineto{\pgfqpoint{3.279173in}{1.774369in}}%
\pgfpathlineto{\pgfqpoint{3.279470in}{1.768032in}}%
\pgfpathlineto{\pgfqpoint{3.283745in}{1.758930in}}%
\pgfpathlineto{\pgfqpoint{3.281055in}{1.749476in}}%
\pgfpathlineto{\pgfqpoint{3.287975in}{1.743465in}}%
\pgfpathlineto{\pgfqpoint{3.300176in}{1.740144in}}%
\pgfpathlineto{\pgfqpoint{3.296564in}{1.727321in}}%
\pgfpathlineto{\pgfqpoint{3.302028in}{1.725960in}}%
\pgfpathclose%
\pgfusepath{fill}%
\end{pgfscope}%
\begin{pgfscope}%
\pgfpathrectangle{\pgfqpoint{0.100000in}{0.100000in}}{\pgfqpoint{3.608454in}{2.310000in}}%
\pgfusepath{clip}%
\pgfsetbuttcap%
\pgfsetmiterjoin%
\definecolor{currentfill}{rgb}{0.000000,0.478431,0.760784}%
\pgfsetfillcolor{currentfill}%
\pgfsetlinewidth{0.000000pt}%
\definecolor{currentstroke}{rgb}{0.000000,0.000000,0.000000}%
\pgfsetstrokecolor{currentstroke}%
\pgfsetstrokeopacity{0.000000}%
\pgfsetdash{}{0pt}%
\pgfpathmoveto{\pgfqpoint{1.917334in}{2.003750in}}%
\pgfpathlineto{\pgfqpoint{1.875835in}{2.005576in}}%
\pgfpathlineto{\pgfqpoint{1.877148in}{2.033168in}}%
\pgfpathlineto{\pgfqpoint{1.876114in}{2.047161in}}%
\pgfpathlineto{\pgfqpoint{1.917529in}{2.045335in}}%
\pgfpathlineto{\pgfqpoint{1.918355in}{2.031373in}}%
\pgfpathlineto{\pgfqpoint{1.917334in}{2.003750in}}%
\pgfpathclose%
\pgfusepath{fill}%
\end{pgfscope}%
\begin{pgfscope}%
\pgfpathrectangle{\pgfqpoint{0.100000in}{0.100000in}}{\pgfqpoint{3.608454in}{2.310000in}}%
\pgfusepath{clip}%
\pgfsetbuttcap%
\pgfsetmiterjoin%
\definecolor{currentfill}{rgb}{0.000000,0.576471,0.711765}%
\pgfsetfillcolor{currentfill}%
\pgfsetlinewidth{0.000000pt}%
\definecolor{currentstroke}{rgb}{0.000000,0.000000,0.000000}%
\pgfsetstrokecolor{currentstroke}%
\pgfsetstrokeopacity{0.000000}%
\pgfsetdash{}{0pt}%
\pgfpathmoveto{\pgfqpoint{1.933314in}{1.014281in}}%
\pgfpathlineto{\pgfqpoint{1.932881in}{1.000512in}}%
\pgfpathlineto{\pgfqpoint{1.905475in}{1.001362in}}%
\pgfpathlineto{\pgfqpoint{1.905912in}{1.015133in}}%
\pgfpathlineto{\pgfqpoint{1.871452in}{1.016408in}}%
\pgfpathlineto{\pgfqpoint{1.872137in}{1.035604in}}%
\pgfpathlineto{\pgfqpoint{1.873435in}{1.071561in}}%
\pgfpathlineto{\pgfqpoint{1.893386in}{1.070793in}}%
\pgfpathlineto{\pgfqpoint{1.893326in}{1.056979in}}%
\pgfpathlineto{\pgfqpoint{1.911760in}{1.056365in}}%
\pgfpathlineto{\pgfqpoint{1.923203in}{1.052514in}}%
\pgfpathlineto{\pgfqpoint{1.934200in}{1.052365in}}%
\pgfpathlineto{\pgfqpoint{1.933314in}{1.014281in}}%
\pgfpathclose%
\pgfusepath{fill}%
\end{pgfscope}%
\begin{pgfscope}%
\pgfpathrectangle{\pgfqpoint{0.100000in}{0.100000in}}{\pgfqpoint{3.608454in}{2.310000in}}%
\pgfusepath{clip}%
\pgfsetbuttcap%
\pgfsetmiterjoin%
\definecolor{currentfill}{rgb}{0.000000,0.458824,0.770588}%
\pgfsetfillcolor{currentfill}%
\pgfsetlinewidth{0.000000pt}%
\definecolor{currentstroke}{rgb}{0.000000,0.000000,0.000000}%
\pgfsetstrokecolor{currentstroke}%
\pgfsetstrokeopacity{0.000000}%
\pgfsetdash{}{0pt}%
\pgfpathmoveto{\pgfqpoint{2.165979in}{1.695985in}}%
\pgfpathlineto{\pgfqpoint{2.166463in}{1.649078in}}%
\pgfpathlineto{\pgfqpoint{2.111776in}{1.648912in}}%
\pgfpathlineto{\pgfqpoint{2.084271in}{1.648958in}}%
\pgfpathlineto{\pgfqpoint{2.084392in}{1.676317in}}%
\pgfpathlineto{\pgfqpoint{2.084490in}{1.695746in}}%
\pgfpathlineto{\pgfqpoint{2.111687in}{1.695719in}}%
\pgfpathlineto{\pgfqpoint{2.111765in}{1.676267in}}%
\pgfpathlineto{\pgfqpoint{2.138939in}{1.676359in}}%
\pgfpathlineto{\pgfqpoint{2.138833in}{1.695825in}}%
\pgfpathlineto{\pgfqpoint{2.165979in}{1.695985in}}%
\pgfpathclose%
\pgfusepath{fill}%
\end{pgfscope}%
\begin{pgfscope}%
\pgfpathrectangle{\pgfqpoint{0.100000in}{0.100000in}}{\pgfqpoint{3.608454in}{2.310000in}}%
\pgfusepath{clip}%
\pgfsetbuttcap%
\pgfsetmiterjoin%
\definecolor{currentfill}{rgb}{0.000000,0.713725,0.643137}%
\pgfsetfillcolor{currentfill}%
\pgfsetlinewidth{0.000000pt}%
\definecolor{currentstroke}{rgb}{0.000000,0.000000,0.000000}%
\pgfsetstrokecolor{currentstroke}%
\pgfsetstrokeopacity{0.000000}%
\pgfsetdash{}{0pt}%
\pgfpathmoveto{\pgfqpoint{2.824999in}{1.309058in}}%
\pgfpathlineto{\pgfqpoint{2.825277in}{1.301789in}}%
\pgfpathlineto{\pgfqpoint{2.820760in}{1.291911in}}%
\pgfpathlineto{\pgfqpoint{2.814875in}{1.291146in}}%
\pgfpathlineto{\pgfqpoint{2.806402in}{1.296003in}}%
\pgfpathlineto{\pgfqpoint{2.793467in}{1.297141in}}%
\pgfpathlineto{\pgfqpoint{2.784277in}{1.310449in}}%
\pgfpathlineto{\pgfqpoint{2.796848in}{1.326453in}}%
\pgfpathlineto{\pgfqpoint{2.804061in}{1.326162in}}%
\pgfpathlineto{\pgfqpoint{2.811045in}{1.319515in}}%
\pgfpathlineto{\pgfqpoint{2.824999in}{1.309058in}}%
\pgfpathclose%
\pgfusepath{fill}%
\end{pgfscope}%
\begin{pgfscope}%
\pgfpathrectangle{\pgfqpoint{0.100000in}{0.100000in}}{\pgfqpoint{3.608454in}{2.310000in}}%
\pgfusepath{clip}%
\pgfsetbuttcap%
\pgfsetmiterjoin%
\definecolor{currentfill}{rgb}{0.000000,0.607843,0.696078}%
\pgfsetfillcolor{currentfill}%
\pgfsetlinewidth{0.000000pt}%
\definecolor{currentstroke}{rgb}{0.000000,0.000000,0.000000}%
\pgfsetstrokecolor{currentstroke}%
\pgfsetstrokeopacity{0.000000}%
\pgfsetdash{}{0pt}%
\pgfpathmoveto{\pgfqpoint{2.586580in}{1.532701in}}%
\pgfpathlineto{\pgfqpoint{2.588380in}{1.512229in}}%
\pgfpathlineto{\pgfqpoint{2.585581in}{1.506060in}}%
\pgfpathlineto{\pgfqpoint{2.578828in}{1.505444in}}%
\pgfpathlineto{\pgfqpoint{2.579594in}{1.497487in}}%
\pgfpathlineto{\pgfqpoint{2.554067in}{1.495355in}}%
\pgfpathlineto{\pgfqpoint{2.550384in}{1.539587in}}%
\pgfpathlineto{\pgfqpoint{2.547881in}{1.571936in}}%
\pgfpathlineto{\pgfqpoint{2.554797in}{1.567328in}}%
\pgfpathlineto{\pgfqpoint{2.566678in}{1.566887in}}%
\pgfpathlineto{\pgfqpoint{2.582676in}{1.575047in}}%
\pgfpathlineto{\pgfqpoint{2.586580in}{1.532701in}}%
\pgfpathclose%
\pgfusepath{fill}%
\end{pgfscope}%
\begin{pgfscope}%
\pgfpathrectangle{\pgfqpoint{0.100000in}{0.100000in}}{\pgfqpoint{3.608454in}{2.310000in}}%
\pgfusepath{clip}%
\pgfsetbuttcap%
\pgfsetmiterjoin%
\definecolor{currentfill}{rgb}{0.000000,0.478431,0.760784}%
\pgfsetfillcolor{currentfill}%
\pgfsetlinewidth{0.000000pt}%
\definecolor{currentstroke}{rgb}{0.000000,0.000000,0.000000}%
\pgfsetstrokecolor{currentstroke}%
\pgfsetstrokeopacity{0.000000}%
\pgfsetdash{}{0pt}%
\pgfpathmoveto{\pgfqpoint{1.844714in}{0.901800in}}%
\pgfpathlineto{\pgfqpoint{1.880039in}{0.900259in}}%
\pgfpathlineto{\pgfqpoint{1.880200in}{0.905864in}}%
\pgfpathlineto{\pgfqpoint{1.909495in}{0.904858in}}%
\pgfpathlineto{\pgfqpoint{1.913431in}{0.902088in}}%
\pgfpathlineto{\pgfqpoint{1.912098in}{0.867745in}}%
\pgfpathlineto{\pgfqpoint{1.878486in}{0.869447in}}%
\pgfpathlineto{\pgfqpoint{1.878275in}{0.865031in}}%
\pgfpathlineto{\pgfqpoint{1.868391in}{0.865576in}}%
\pgfpathlineto{\pgfqpoint{1.808729in}{0.868432in}}%
\pgfpathlineto{\pgfqpoint{1.810377in}{0.903512in}}%
\pgfpathlineto{\pgfqpoint{1.844714in}{0.901800in}}%
\pgfpathclose%
\pgfusepath{fill}%
\end{pgfscope}%
\begin{pgfscope}%
\pgfpathrectangle{\pgfqpoint{0.100000in}{0.100000in}}{\pgfqpoint{3.608454in}{2.310000in}}%
\pgfusepath{clip}%
\pgfsetbuttcap%
\pgfsetmiterjoin%
\definecolor{currentfill}{rgb}{0.000000,0.572549,0.713725}%
\pgfsetfillcolor{currentfill}%
\pgfsetlinewidth{0.000000pt}%
\definecolor{currentstroke}{rgb}{0.000000,0.000000,0.000000}%
\pgfsetstrokecolor{currentstroke}%
\pgfsetstrokeopacity{0.000000}%
\pgfsetdash{}{0pt}%
\pgfpathmoveto{\pgfqpoint{2.307376in}{1.370072in}}%
\pgfpathlineto{\pgfqpoint{2.271164in}{1.369568in}}%
\pgfpathlineto{\pgfqpoint{2.271150in}{1.361570in}}%
\pgfpathlineto{\pgfqpoint{2.264061in}{1.361534in}}%
\pgfpathlineto{\pgfqpoint{2.246959in}{1.366834in}}%
\pgfpathlineto{\pgfqpoint{2.247305in}{1.396610in}}%
\pgfpathlineto{\pgfqpoint{2.237052in}{1.396510in}}%
\pgfpathlineto{\pgfqpoint{2.237055in}{1.423181in}}%
\pgfpathlineto{\pgfqpoint{2.267438in}{1.423982in}}%
\pgfpathlineto{\pgfqpoint{2.267659in}{1.417346in}}%
\pgfpathlineto{\pgfqpoint{2.291417in}{1.417786in}}%
\pgfpathlineto{\pgfqpoint{2.298225in}{1.417924in}}%
\pgfpathlineto{\pgfqpoint{2.298632in}{1.395025in}}%
\pgfpathlineto{\pgfqpoint{2.306650in}{1.395239in}}%
\pgfpathlineto{\pgfqpoint{2.307376in}{1.370072in}}%
\pgfpathclose%
\pgfusepath{fill}%
\end{pgfscope}%
\begin{pgfscope}%
\pgfpathrectangle{\pgfqpoint{0.100000in}{0.100000in}}{\pgfqpoint{3.608454in}{2.310000in}}%
\pgfusepath{clip}%
\pgfsetbuttcap%
\pgfsetmiterjoin%
\definecolor{currentfill}{rgb}{0.000000,0.666667,0.666667}%
\pgfsetfillcolor{currentfill}%
\pgfsetlinewidth{0.000000pt}%
\definecolor{currentstroke}{rgb}{0.000000,0.000000,0.000000}%
\pgfsetstrokecolor{currentstroke}%
\pgfsetstrokeopacity{0.000000}%
\pgfsetdash{}{0pt}%
\pgfpathmoveto{\pgfqpoint{2.043077in}{1.035565in}}%
\pgfpathlineto{\pgfqpoint{2.043347in}{1.046522in}}%
\pgfpathlineto{\pgfqpoint{2.043425in}{1.053406in}}%
\pgfpathlineto{\pgfqpoint{2.050250in}{1.053326in}}%
\pgfpathlineto{\pgfqpoint{2.053684in}{1.060206in}}%
\pgfpathlineto{\pgfqpoint{2.053817in}{1.067155in}}%
\pgfpathlineto{\pgfqpoint{2.060863in}{1.067101in}}%
\pgfpathlineto{\pgfqpoint{2.060962in}{1.080885in}}%
\pgfpathlineto{\pgfqpoint{2.057547in}{1.080899in}}%
\pgfpathlineto{\pgfqpoint{2.057637in}{1.091246in}}%
\pgfpathlineto{\pgfqpoint{2.065063in}{1.091216in}}%
\pgfpathlineto{\pgfqpoint{2.068981in}{1.083880in}}%
\pgfpathlineto{\pgfqpoint{2.076132in}{1.086986in}}%
\pgfpathlineto{\pgfqpoint{2.089737in}{1.087648in}}%
\pgfpathlineto{\pgfqpoint{2.088937in}{1.099591in}}%
\pgfpathlineto{\pgfqpoint{2.093490in}{1.108334in}}%
\pgfpathlineto{\pgfqpoint{2.099232in}{1.108302in}}%
\pgfpathlineto{\pgfqpoint{2.099224in}{1.115204in}}%
\pgfpathlineto{\pgfqpoint{2.106071in}{1.115212in}}%
\pgfpathlineto{\pgfqpoint{2.134726in}{1.115276in}}%
\pgfpathlineto{\pgfqpoint{2.139356in}{1.083433in}}%
\pgfpathlineto{\pgfqpoint{2.140727in}{1.073891in}}%
\pgfpathlineto{\pgfqpoint{2.098590in}{1.073824in}}%
\pgfpathlineto{\pgfqpoint{2.098258in}{1.064900in}}%
\pgfpathlineto{\pgfqpoint{2.101602in}{1.057429in}}%
\pgfpathlineto{\pgfqpoint{2.096410in}{1.052482in}}%
\pgfpathlineto{\pgfqpoint{2.091976in}{1.044393in}}%
\pgfpathlineto{\pgfqpoint{2.081865in}{1.049099in}}%
\pgfpathlineto{\pgfqpoint{2.069896in}{1.042778in}}%
\pgfpathlineto{\pgfqpoint{2.058740in}{1.038661in}}%
\pgfpathlineto{\pgfqpoint{2.050477in}{1.038266in}}%
\pgfpathlineto{\pgfqpoint{2.043077in}{1.035565in}}%
\pgfpathclose%
\pgfusepath{fill}%
\end{pgfscope}%
\begin{pgfscope}%
\pgfpathrectangle{\pgfqpoint{0.100000in}{0.100000in}}{\pgfqpoint{3.608454in}{2.310000in}}%
\pgfusepath{clip}%
\pgfsetbuttcap%
\pgfsetmiterjoin%
\definecolor{currentfill}{rgb}{0.000000,0.627451,0.686275}%
\pgfsetfillcolor{currentfill}%
\pgfsetlinewidth{0.000000pt}%
\definecolor{currentstroke}{rgb}{0.000000,0.000000,0.000000}%
\pgfsetstrokecolor{currentstroke}%
\pgfsetstrokeopacity{0.000000}%
\pgfsetdash{}{0pt}%
\pgfpathmoveto{\pgfqpoint{2.088125in}{1.418722in}}%
\pgfpathlineto{\pgfqpoint{2.081282in}{1.423410in}}%
\pgfpathlineto{\pgfqpoint{2.082351in}{1.428944in}}%
\pgfpathlineto{\pgfqpoint{2.077853in}{1.433443in}}%
\pgfpathlineto{\pgfqpoint{2.077207in}{1.438452in}}%
\pgfpathlineto{\pgfqpoint{2.069544in}{1.443619in}}%
\pgfpathlineto{\pgfqpoint{2.060821in}{1.465149in}}%
\pgfpathlineto{\pgfqpoint{2.128892in}{1.463893in}}%
\pgfpathlineto{\pgfqpoint{2.130971in}{1.449316in}}%
\pgfpathlineto{\pgfqpoint{2.130612in}{1.428716in}}%
\pgfpathlineto{\pgfqpoint{2.104194in}{1.428993in}}%
\pgfpathlineto{\pgfqpoint{2.107194in}{1.411974in}}%
\pgfpathlineto{\pgfqpoint{2.101635in}{1.407776in}}%
\pgfpathlineto{\pgfqpoint{2.094548in}{1.411120in}}%
\pgfpathlineto{\pgfqpoint{2.088125in}{1.418722in}}%
\pgfpathclose%
\pgfusepath{fill}%
\end{pgfscope}%
\begin{pgfscope}%
\pgfpathrectangle{\pgfqpoint{0.100000in}{0.100000in}}{\pgfqpoint{3.608454in}{2.310000in}}%
\pgfusepath{clip}%
\pgfsetbuttcap%
\pgfsetmiterjoin%
\definecolor{currentfill}{rgb}{0.000000,0.733333,0.633333}%
\pgfsetfillcolor{currentfill}%
\pgfsetlinewidth{0.000000pt}%
\definecolor{currentstroke}{rgb}{0.000000,0.000000,0.000000}%
\pgfsetstrokecolor{currentstroke}%
\pgfsetstrokeopacity{0.000000}%
\pgfsetdash{}{0pt}%
\pgfpathmoveto{\pgfqpoint{2.722462in}{1.100629in}}%
\pgfpathlineto{\pgfqpoint{2.714180in}{1.097018in}}%
\pgfpathlineto{\pgfqpoint{2.697680in}{1.093767in}}%
\pgfpathlineto{\pgfqpoint{2.693772in}{1.094220in}}%
\pgfpathlineto{\pgfqpoint{2.688788in}{1.097666in}}%
\pgfpathlineto{\pgfqpoint{2.685775in}{1.104241in}}%
\pgfpathlineto{\pgfqpoint{2.685202in}{1.110108in}}%
\pgfpathlineto{\pgfqpoint{2.689025in}{1.121466in}}%
\pgfpathlineto{\pgfqpoint{2.681516in}{1.127762in}}%
\pgfpathlineto{\pgfqpoint{2.678100in}{1.135237in}}%
\pgfpathlineto{\pgfqpoint{2.683825in}{1.139915in}}%
\pgfpathlineto{\pgfqpoint{2.690045in}{1.139163in}}%
\pgfpathlineto{\pgfqpoint{2.693376in}{1.142517in}}%
\pgfpathlineto{\pgfqpoint{2.695803in}{1.138527in}}%
\pgfpathlineto{\pgfqpoint{2.701451in}{1.138405in}}%
\pgfpathlineto{\pgfqpoint{2.703971in}{1.134522in}}%
\pgfpathlineto{\pgfqpoint{2.712831in}{1.140583in}}%
\pgfpathlineto{\pgfqpoint{2.722754in}{1.138754in}}%
\pgfpathlineto{\pgfqpoint{2.729200in}{1.134130in}}%
\pgfpathlineto{\pgfqpoint{2.732727in}{1.126386in}}%
\pgfpathlineto{\pgfqpoint{2.731314in}{1.113832in}}%
\pgfpathlineto{\pgfqpoint{2.722462in}{1.100629in}}%
\pgfpathclose%
\pgfusepath{fill}%
\end{pgfscope}%
\begin{pgfscope}%
\pgfpathrectangle{\pgfqpoint{0.100000in}{0.100000in}}{\pgfqpoint{3.608454in}{2.310000in}}%
\pgfusepath{clip}%
\pgfsetbuttcap%
\pgfsetmiterjoin%
\definecolor{currentfill}{rgb}{0.000000,0.627451,0.686275}%
\pgfsetfillcolor{currentfill}%
\pgfsetlinewidth{0.000000pt}%
\definecolor{currentstroke}{rgb}{0.000000,0.000000,0.000000}%
\pgfsetstrokecolor{currentstroke}%
\pgfsetstrokeopacity{0.000000}%
\pgfsetdash{}{0pt}%
\pgfpathmoveto{\pgfqpoint{1.411936in}{2.184794in}}%
\pgfpathlineto{\pgfqpoint{1.473904in}{2.176145in}}%
\pgfpathlineto{\pgfqpoint{1.522840in}{2.169802in}}%
\pgfpathlineto{\pgfqpoint{1.519036in}{2.156103in}}%
\pgfpathlineto{\pgfqpoint{1.525884in}{2.155265in}}%
\pgfpathlineto{\pgfqpoint{1.523357in}{2.134506in}}%
\pgfpathlineto{\pgfqpoint{1.534783in}{2.133134in}}%
\pgfpathlineto{\pgfqpoint{1.531517in}{2.105747in}}%
\pgfpathlineto{\pgfqpoint{1.529442in}{2.105984in}}%
\pgfpathlineto{\pgfqpoint{1.527435in}{2.089391in}}%
\pgfpathlineto{\pgfqpoint{1.507101in}{2.093171in}}%
\pgfpathlineto{\pgfqpoint{1.498726in}{2.097310in}}%
\pgfpathlineto{\pgfqpoint{1.496354in}{2.088951in}}%
\pgfpathlineto{\pgfqpoint{1.492326in}{2.087153in}}%
\pgfpathlineto{\pgfqpoint{1.487031in}{2.078600in}}%
\pgfpathlineto{\pgfqpoint{1.476331in}{2.073148in}}%
\pgfpathlineto{\pgfqpoint{1.470183in}{2.074395in}}%
\pgfpathlineto{\pgfqpoint{1.468292in}{2.069917in}}%
\pgfpathlineto{\pgfqpoint{1.449108in}{2.071253in}}%
\pgfpathlineto{\pgfqpoint{1.440516in}{2.074870in}}%
\pgfpathlineto{\pgfqpoint{1.437875in}{2.069615in}}%
\pgfpathlineto{\pgfqpoint{1.428142in}{2.073167in}}%
\pgfpathlineto{\pgfqpoint{1.422183in}{2.067656in}}%
\pgfpathlineto{\pgfqpoint{1.414702in}{2.064290in}}%
\pgfpathlineto{\pgfqpoint{1.411450in}{2.059486in}}%
\pgfpathlineto{\pgfqpoint{1.411723in}{2.068893in}}%
\pgfpathlineto{\pgfqpoint{1.404841in}{2.072375in}}%
\pgfpathlineto{\pgfqpoint{1.389242in}{2.072875in}}%
\pgfpathlineto{\pgfqpoint{1.385004in}{2.076421in}}%
\pgfpathlineto{\pgfqpoint{1.377859in}{2.077250in}}%
\pgfpathlineto{\pgfqpoint{1.367934in}{2.081540in}}%
\pgfpathlineto{\pgfqpoint{1.362402in}{2.089710in}}%
\pgfpathlineto{\pgfqpoint{1.364847in}{2.104580in}}%
\pgfpathlineto{\pgfqpoint{1.378739in}{2.102305in}}%
\pgfpathlineto{\pgfqpoint{1.380515in}{2.107723in}}%
\pgfpathlineto{\pgfqpoint{1.389566in}{2.105191in}}%
\pgfpathlineto{\pgfqpoint{1.392437in}{2.124428in}}%
\pgfpathlineto{\pgfqpoint{1.396285in}{2.142117in}}%
\pgfpathlineto{\pgfqpoint{1.399456in}{2.139818in}}%
\pgfpathlineto{\pgfqpoint{1.402592in}{2.151198in}}%
\pgfpathlineto{\pgfqpoint{1.404613in}{2.164996in}}%
\pgfpathlineto{\pgfqpoint{1.407741in}{2.164527in}}%
\pgfpathlineto{\pgfqpoint{1.411936in}{2.184794in}}%
\pgfpathclose%
\pgfusepath{fill}%
\end{pgfscope}%
\begin{pgfscope}%
\pgfpathrectangle{\pgfqpoint{0.100000in}{0.100000in}}{\pgfqpoint{3.608454in}{2.310000in}}%
\pgfusepath{clip}%
\pgfsetbuttcap%
\pgfsetmiterjoin%
\definecolor{currentfill}{rgb}{0.000000,0.572549,0.713725}%
\pgfsetfillcolor{currentfill}%
\pgfsetlinewidth{0.000000pt}%
\definecolor{currentstroke}{rgb}{0.000000,0.000000,0.000000}%
\pgfsetstrokecolor{currentstroke}%
\pgfsetstrokeopacity{0.000000}%
\pgfsetdash{}{0pt}%
\pgfpathmoveto{\pgfqpoint{3.150978in}{0.513727in}}%
\pgfpathlineto{\pgfqpoint{3.118731in}{0.508688in}}%
\pgfpathlineto{\pgfqpoint{3.117607in}{0.515401in}}%
\pgfpathlineto{\pgfqpoint{3.110913in}{0.514275in}}%
\pgfpathlineto{\pgfqpoint{3.102843in}{0.570453in}}%
\pgfpathlineto{\pgfqpoint{3.099105in}{0.576837in}}%
\pgfpathlineto{\pgfqpoint{3.099275in}{0.583081in}}%
\pgfpathlineto{\pgfqpoint{3.090875in}{0.590177in}}%
\pgfpathlineto{\pgfqpoint{3.090036in}{0.604447in}}%
\pgfpathlineto{\pgfqpoint{3.106445in}{0.607053in}}%
\pgfpathlineto{\pgfqpoint{3.119882in}{0.592598in}}%
\pgfpathlineto{\pgfqpoint{3.125061in}{0.582905in}}%
\pgfpathlineto{\pgfqpoint{3.120768in}{0.574633in}}%
\pgfpathlineto{\pgfqpoint{3.121728in}{0.567411in}}%
\pgfpathlineto{\pgfqpoint{3.127445in}{0.552410in}}%
\pgfpathlineto{\pgfqpoint{3.144148in}{0.527548in}}%
\pgfpathlineto{\pgfqpoint{3.150978in}{0.513727in}}%
\pgfpathclose%
\pgfusepath{fill}%
\end{pgfscope}%
\begin{pgfscope}%
\pgfpathrectangle{\pgfqpoint{0.100000in}{0.100000in}}{\pgfqpoint{3.608454in}{2.310000in}}%
\pgfusepath{clip}%
\pgfsetbuttcap%
\pgfsetmiterjoin%
\definecolor{currentfill}{rgb}{0.000000,0.674510,0.662745}%
\pgfsetfillcolor{currentfill}%
\pgfsetlinewidth{0.000000pt}%
\definecolor{currentstroke}{rgb}{0.000000,0.000000,0.000000}%
\pgfsetstrokecolor{currentstroke}%
\pgfsetstrokeopacity{0.000000}%
\pgfsetdash{}{0pt}%
\pgfpathmoveto{\pgfqpoint{2.218257in}{1.006289in}}%
\pgfpathlineto{\pgfqpoint{2.211366in}{1.003801in}}%
\pgfpathlineto{\pgfqpoint{2.190786in}{1.003624in}}%
\pgfpathlineto{\pgfqpoint{2.190828in}{1.025374in}}%
\pgfpathlineto{\pgfqpoint{2.166850in}{1.025625in}}%
\pgfpathlineto{\pgfqpoint{2.162454in}{1.027888in}}%
\pgfpathlineto{\pgfqpoint{2.162548in}{1.034836in}}%
\pgfpathlineto{\pgfqpoint{2.166153in}{1.040540in}}%
\pgfpathlineto{\pgfqpoint{2.166477in}{1.058801in}}%
\pgfpathlineto{\pgfqpoint{2.169375in}{1.063715in}}%
\pgfpathlineto{\pgfqpoint{2.166287in}{1.069101in}}%
\pgfpathlineto{\pgfqpoint{2.166395in}{1.075969in}}%
\pgfpathlineto{\pgfqpoint{2.172137in}{1.075918in}}%
\pgfpathlineto{\pgfqpoint{2.176762in}{1.083808in}}%
\pgfpathlineto{\pgfqpoint{2.201915in}{1.084462in}}%
\pgfpathlineto{\pgfqpoint{2.201928in}{1.082158in}}%
\pgfpathlineto{\pgfqpoint{2.247558in}{1.082128in}}%
\pgfpathlineto{\pgfqpoint{2.247654in}{1.068396in}}%
\pgfpathlineto{\pgfqpoint{2.245419in}{1.061527in}}%
\pgfpathlineto{\pgfqpoint{2.245679in}{1.038542in}}%
\pgfpathlineto{\pgfqpoint{2.241035in}{1.037964in}}%
\pgfpathlineto{\pgfqpoint{2.226237in}{1.022520in}}%
\pgfpathlineto{\pgfqpoint{2.216940in}{1.014382in}}%
\pgfpathlineto{\pgfqpoint{2.218257in}{1.006289in}}%
\pgfpathclose%
\pgfusepath{fill}%
\end{pgfscope}%
\begin{pgfscope}%
\pgfpathrectangle{\pgfqpoint{0.100000in}{0.100000in}}{\pgfqpoint{3.608454in}{2.310000in}}%
\pgfusepath{clip}%
\pgfsetbuttcap%
\pgfsetmiterjoin%
\definecolor{currentfill}{rgb}{0.000000,0.411765,0.794118}%
\pgfsetfillcolor{currentfill}%
\pgfsetlinewidth{0.000000pt}%
\definecolor{currentstroke}{rgb}{0.000000,0.000000,0.000000}%
\pgfsetstrokecolor{currentstroke}%
\pgfsetstrokeopacity{0.000000}%
\pgfsetdash{}{0pt}%
\pgfpathmoveto{\pgfqpoint{1.075355in}{0.913752in}}%
\pgfpathlineto{\pgfqpoint{1.001421in}{0.926345in}}%
\pgfpathlineto{\pgfqpoint{0.934927in}{0.938858in}}%
\pgfpathlineto{\pgfqpoint{0.885192in}{0.948795in}}%
\pgfpathlineto{\pgfqpoint{0.899087in}{1.016739in}}%
\pgfpathlineto{\pgfqpoint{0.908989in}{1.065097in}}%
\pgfpathlineto{\pgfqpoint{0.947034in}{1.057499in}}%
\pgfpathlineto{\pgfqpoint{0.975479in}{1.042443in}}%
\pgfpathlineto{\pgfqpoint{0.985165in}{1.053983in}}%
\pgfpathlineto{\pgfqpoint{1.012835in}{1.044864in}}%
\pgfpathlineto{\pgfqpoint{1.027746in}{1.042105in}}%
\pgfpathlineto{\pgfqpoint{1.031005in}{1.053711in}}%
\pgfpathlineto{\pgfqpoint{1.015879in}{1.056517in}}%
\pgfpathlineto{\pgfqpoint{1.022683in}{1.074223in}}%
\pgfpathlineto{\pgfqpoint{1.031414in}{1.079449in}}%
\pgfpathlineto{\pgfqpoint{1.039948in}{1.073968in}}%
\pgfpathlineto{\pgfqpoint{1.049066in}{1.073939in}}%
\pgfpathlineto{\pgfqpoint{1.055533in}{1.068101in}}%
\pgfpathlineto{\pgfqpoint{1.063896in}{1.064728in}}%
\pgfpathlineto{\pgfqpoint{1.073305in}{1.055337in}}%
\pgfpathlineto{\pgfqpoint{1.079528in}{1.054057in}}%
\pgfpathlineto{\pgfqpoint{1.075926in}{1.033485in}}%
\pgfpathlineto{\pgfqpoint{1.124557in}{1.025178in}}%
\pgfpathlineto{\pgfqpoint{1.119039in}{0.992055in}}%
\pgfpathlineto{\pgfqpoint{1.117590in}{0.983356in}}%
\pgfpathlineto{\pgfqpoint{1.106555in}{0.985188in}}%
\pgfpathlineto{\pgfqpoint{1.096233in}{0.984306in}}%
\pgfpathlineto{\pgfqpoint{1.087666in}{0.981611in}}%
\pgfpathlineto{\pgfqpoint{1.084657in}{0.967062in}}%
\pgfpathlineto{\pgfqpoint{1.075355in}{0.913752in}}%
\pgfpathclose%
\pgfusepath{fill}%
\end{pgfscope}%
\begin{pgfscope}%
\pgfpathrectangle{\pgfqpoint{0.100000in}{0.100000in}}{\pgfqpoint{3.608454in}{2.310000in}}%
\pgfusepath{clip}%
\pgfsetbuttcap%
\pgfsetmiterjoin%
\definecolor{currentfill}{rgb}{0.000000,0.447059,0.776471}%
\pgfsetfillcolor{currentfill}%
\pgfsetlinewidth{0.000000pt}%
\definecolor{currentstroke}{rgb}{0.000000,0.000000,0.000000}%
\pgfsetstrokecolor{currentstroke}%
\pgfsetstrokeopacity{0.000000}%
\pgfsetdash{}{0pt}%
\pgfpathmoveto{\pgfqpoint{2.032756in}{2.134709in}}%
\pgfpathlineto{\pgfqpoint{2.098692in}{2.134056in}}%
\pgfpathlineto{\pgfqpoint{2.098766in}{2.165029in}}%
\pgfpathlineto{\pgfqpoint{2.103732in}{2.162515in}}%
\pgfpathlineto{\pgfqpoint{2.109017in}{2.163883in}}%
\pgfpathlineto{\pgfqpoint{2.115955in}{2.157789in}}%
\pgfpathlineto{\pgfqpoint{2.123516in}{2.124861in}}%
\pgfpathlineto{\pgfqpoint{2.123134in}{2.116351in}}%
\pgfpathlineto{\pgfqpoint{2.128642in}{2.111522in}}%
\pgfpathlineto{\pgfqpoint{2.136953in}{2.110228in}}%
\pgfpathlineto{\pgfqpoint{2.137206in}{2.083483in}}%
\pgfpathlineto{\pgfqpoint{2.095560in}{2.083291in}}%
\pgfpathlineto{\pgfqpoint{2.095537in}{2.097237in}}%
\pgfpathlineto{\pgfqpoint{2.074816in}{2.097239in}}%
\pgfpathlineto{\pgfqpoint{2.033124in}{2.098129in}}%
\pgfpathlineto{\pgfqpoint{2.032428in}{2.112015in}}%
\pgfpathlineto{\pgfqpoint{2.032756in}{2.134709in}}%
\pgfpathclose%
\pgfusepath{fill}%
\end{pgfscope}%
\begin{pgfscope}%
\pgfpathrectangle{\pgfqpoint{0.100000in}{0.100000in}}{\pgfqpoint{3.608454in}{2.310000in}}%
\pgfusepath{clip}%
\pgfsetbuttcap%
\pgfsetmiterjoin%
\definecolor{currentfill}{rgb}{0.000000,0.576471,0.711765}%
\pgfsetfillcolor{currentfill}%
\pgfsetlinewidth{0.000000pt}%
\definecolor{currentstroke}{rgb}{0.000000,0.000000,0.000000}%
\pgfsetstrokecolor{currentstroke}%
\pgfsetstrokeopacity{0.000000}%
\pgfsetdash{}{0pt}%
\pgfpathmoveto{\pgfqpoint{1.001953in}{2.138275in}}%
\pgfpathlineto{\pgfqpoint{0.996502in}{2.145181in}}%
\pgfpathlineto{\pgfqpoint{1.000929in}{2.147949in}}%
\pgfpathlineto{\pgfqpoint{1.000904in}{2.156205in}}%
\pgfpathlineto{\pgfqpoint{0.996594in}{2.161889in}}%
\pgfpathlineto{\pgfqpoint{0.988985in}{2.182170in}}%
\pgfpathlineto{\pgfqpoint{0.993310in}{2.200722in}}%
\pgfpathlineto{\pgfqpoint{0.994462in}{2.195826in}}%
\pgfpathlineto{\pgfqpoint{1.011635in}{2.199142in}}%
\pgfpathlineto{\pgfqpoint{1.012094in}{2.185135in}}%
\pgfpathlineto{\pgfqpoint{1.015735in}{2.176278in}}%
\pgfpathlineto{\pgfqpoint{1.015079in}{2.170487in}}%
\pgfpathlineto{\pgfqpoint{1.027896in}{2.166208in}}%
\pgfpathlineto{\pgfqpoint{1.033549in}{2.168208in}}%
\pgfpathlineto{\pgfqpoint{1.036396in}{2.174561in}}%
\pgfpathlineto{\pgfqpoint{1.043899in}{2.172864in}}%
\pgfpathlineto{\pgfqpoint{1.047061in}{2.189256in}}%
\pgfpathlineto{\pgfqpoint{1.056030in}{2.187273in}}%
\pgfpathlineto{\pgfqpoint{1.061972in}{2.214474in}}%
\pgfpathlineto{\pgfqpoint{1.061432in}{2.221573in}}%
\pgfpathlineto{\pgfqpoint{1.072055in}{2.221371in}}%
\pgfpathlineto{\pgfqpoint{1.075971in}{2.228263in}}%
\pgfpathlineto{\pgfqpoint{1.074794in}{2.233716in}}%
\pgfpathlineto{\pgfqpoint{1.075629in}{2.246526in}}%
\pgfpathlineto{\pgfqpoint{1.109601in}{2.239048in}}%
\pgfpathlineto{\pgfqpoint{1.109470in}{2.229400in}}%
\pgfpathlineto{\pgfqpoint{1.112359in}{2.223590in}}%
\pgfpathlineto{\pgfqpoint{1.118945in}{2.225334in}}%
\pgfpathlineto{\pgfqpoint{1.123862in}{2.213141in}}%
\pgfpathlineto{\pgfqpoint{1.119541in}{2.205769in}}%
\pgfpathlineto{\pgfqpoint{1.133270in}{2.197089in}}%
\pgfpathlineto{\pgfqpoint{1.131142in}{2.190334in}}%
\pgfpathlineto{\pgfqpoint{1.136982in}{2.181768in}}%
\pgfpathlineto{\pgfqpoint{1.134122in}{2.180102in}}%
\pgfpathlineto{\pgfqpoint{1.140465in}{2.171053in}}%
\pgfpathlineto{\pgfqpoint{1.139517in}{2.164977in}}%
\pgfpathlineto{\pgfqpoint{1.150149in}{2.159708in}}%
\pgfpathlineto{\pgfqpoint{1.154266in}{2.150464in}}%
\pgfpathlineto{\pgfqpoint{1.149026in}{2.145360in}}%
\pgfpathlineto{\pgfqpoint{1.143835in}{2.140807in}}%
\pgfpathlineto{\pgfqpoint{1.140741in}{2.130368in}}%
\pgfpathlineto{\pgfqpoint{1.135799in}{2.128987in}}%
\pgfpathlineto{\pgfqpoint{1.134932in}{2.119015in}}%
\pgfpathlineto{\pgfqpoint{1.118004in}{2.122768in}}%
\pgfpathlineto{\pgfqpoint{1.093543in}{2.127613in}}%
\pgfpathlineto{\pgfqpoint{1.090372in}{2.118413in}}%
\pgfpathlineto{\pgfqpoint{1.093234in}{2.110956in}}%
\pgfpathlineto{\pgfqpoint{1.091464in}{2.102885in}}%
\pgfpathlineto{\pgfqpoint{1.092364in}{2.093954in}}%
\pgfpathlineto{\pgfqpoint{1.085713in}{2.091817in}}%
\pgfpathlineto{\pgfqpoint{1.072235in}{2.094680in}}%
\pgfpathlineto{\pgfqpoint{1.068378in}{2.094075in}}%
\pgfpathlineto{\pgfqpoint{1.060904in}{2.103004in}}%
\pgfpathlineto{\pgfqpoint{1.048078in}{2.111068in}}%
\pgfpathlineto{\pgfqpoint{1.039741in}{2.112438in}}%
\pgfpathlineto{\pgfqpoint{1.033798in}{2.117353in}}%
\pgfpathlineto{\pgfqpoint{1.034696in}{2.123940in}}%
\pgfpathlineto{\pgfqpoint{1.020317in}{2.134665in}}%
\pgfpathlineto{\pgfqpoint{1.011451in}{2.134988in}}%
\pgfpathlineto{\pgfqpoint{1.001953in}{2.138275in}}%
\pgfpathclose%
\pgfusepath{fill}%
\end{pgfscope}%
\begin{pgfscope}%
\pgfpathrectangle{\pgfqpoint{0.100000in}{0.100000in}}{\pgfqpoint{3.608454in}{2.310000in}}%
\pgfusepath{clip}%
\pgfsetbuttcap%
\pgfsetmiterjoin%
\definecolor{currentfill}{rgb}{0.000000,0.666667,0.666667}%
\pgfsetfillcolor{currentfill}%
\pgfsetlinewidth{0.000000pt}%
\definecolor{currentstroke}{rgb}{0.000000,0.000000,0.000000}%
\pgfsetstrokecolor{currentstroke}%
\pgfsetstrokeopacity{0.000000}%
\pgfsetdash{}{0pt}%
\pgfpathmoveto{\pgfqpoint{1.583492in}{0.651351in}}%
\pgfpathlineto{\pgfqpoint{1.599353in}{0.636098in}}%
\pgfpathlineto{\pgfqpoint{1.593544in}{0.627739in}}%
\pgfpathlineto{\pgfqpoint{1.573827in}{0.627230in}}%
\pgfpathlineto{\pgfqpoint{1.567953in}{0.615750in}}%
\pgfpathlineto{\pgfqpoint{1.563163in}{0.610468in}}%
\pgfpathlineto{\pgfqpoint{1.559813in}{0.597646in}}%
\pgfpathlineto{\pgfqpoint{1.546004in}{0.580048in}}%
\pgfpathlineto{\pgfqpoint{1.539619in}{0.575381in}}%
\pgfpathlineto{\pgfqpoint{1.538032in}{0.569527in}}%
\pgfpathlineto{\pgfqpoint{1.529354in}{0.569946in}}%
\pgfpathlineto{\pgfqpoint{1.514568in}{0.578283in}}%
\pgfpathlineto{\pgfqpoint{1.508590in}{0.586063in}}%
\pgfpathlineto{\pgfqpoint{1.497563in}{0.588987in}}%
\pgfpathlineto{\pgfqpoint{1.493508in}{0.596165in}}%
\pgfpathlineto{\pgfqpoint{1.480379in}{0.599940in}}%
\pgfpathlineto{\pgfqpoint{1.469478in}{0.607802in}}%
\pgfpathlineto{\pgfqpoint{1.463905in}{0.617758in}}%
\pgfpathlineto{\pgfqpoint{1.457236in}{0.620260in}}%
\pgfpathlineto{\pgfqpoint{1.445703in}{0.631256in}}%
\pgfpathlineto{\pgfqpoint{1.443059in}{0.642585in}}%
\pgfpathlineto{\pgfqpoint{1.437063in}{0.653883in}}%
\pgfpathlineto{\pgfqpoint{1.436330in}{0.666295in}}%
\pgfpathlineto{\pgfqpoint{1.437587in}{0.679889in}}%
\pgfpathlineto{\pgfqpoint{1.428449in}{0.693343in}}%
\pgfpathlineto{\pgfqpoint{1.428817in}{0.701695in}}%
\pgfpathlineto{\pgfqpoint{1.425660in}{0.710940in}}%
\pgfpathlineto{\pgfqpoint{1.422243in}{0.713258in}}%
\pgfpathlineto{\pgfqpoint{1.426837in}{0.715528in}}%
\pgfpathlineto{\pgfqpoint{1.486032in}{0.744751in}}%
\pgfpathlineto{\pgfqpoint{1.518651in}{0.714404in}}%
\pgfpathlineto{\pgfqpoint{1.583492in}{0.651351in}}%
\pgfpathclose%
\pgfusepath{fill}%
\end{pgfscope}%
\begin{pgfscope}%
\pgfpathrectangle{\pgfqpoint{0.100000in}{0.100000in}}{\pgfqpoint{3.608454in}{2.310000in}}%
\pgfusepath{clip}%
\pgfsetbuttcap%
\pgfsetmiterjoin%
\definecolor{currentfill}{rgb}{0.000000,0.639216,0.680392}%
\pgfsetfillcolor{currentfill}%
\pgfsetlinewidth{0.000000pt}%
\definecolor{currentstroke}{rgb}{0.000000,0.000000,0.000000}%
\pgfsetstrokecolor{currentstroke}%
\pgfsetstrokeopacity{0.000000}%
\pgfsetdash{}{0pt}%
\pgfpathmoveto{\pgfqpoint{2.715165in}{0.972447in}}%
\pgfpathlineto{\pgfqpoint{2.705960in}{0.970065in}}%
\pgfpathlineto{\pgfqpoint{2.694783in}{0.960330in}}%
\pgfpathlineto{\pgfqpoint{2.689264in}{0.964347in}}%
\pgfpathlineto{\pgfqpoint{2.684888in}{0.971060in}}%
\pgfpathlineto{\pgfqpoint{2.678957in}{0.970072in}}%
\pgfpathlineto{\pgfqpoint{2.675715in}{0.974764in}}%
\pgfpathlineto{\pgfqpoint{2.677206in}{0.979100in}}%
\pgfpathlineto{\pgfqpoint{2.664313in}{0.994103in}}%
\pgfpathlineto{\pgfqpoint{2.657488in}{0.993619in}}%
\pgfpathlineto{\pgfqpoint{2.659190in}{1.009555in}}%
\pgfpathlineto{\pgfqpoint{2.666467in}{1.008380in}}%
\pgfpathlineto{\pgfqpoint{2.673359in}{1.013797in}}%
\pgfpathlineto{\pgfqpoint{2.669558in}{1.020872in}}%
\pgfpathlineto{\pgfqpoint{2.670176in}{1.036633in}}%
\pgfpathlineto{\pgfqpoint{2.668245in}{1.039940in}}%
\pgfpathlineto{\pgfqpoint{2.669947in}{1.049451in}}%
\pgfpathlineto{\pgfqpoint{2.698781in}{1.052062in}}%
\pgfpathlineto{\pgfqpoint{2.715608in}{1.053498in}}%
\pgfpathlineto{\pgfqpoint{2.723904in}{1.022807in}}%
\pgfpathlineto{\pgfqpoint{2.725316in}{1.017812in}}%
\pgfpathlineto{\pgfqpoint{2.721545in}{1.014250in}}%
\pgfpathlineto{\pgfqpoint{2.718615in}{1.004553in}}%
\pgfpathlineto{\pgfqpoint{2.710524in}{0.995150in}}%
\pgfpathlineto{\pgfqpoint{2.706101in}{0.993510in}}%
\pgfpathlineto{\pgfqpoint{2.715165in}{0.972447in}}%
\pgfpathclose%
\pgfusepath{fill}%
\end{pgfscope}%
\begin{pgfscope}%
\pgfpathrectangle{\pgfqpoint{0.100000in}{0.100000in}}{\pgfqpoint{3.608454in}{2.310000in}}%
\pgfusepath{clip}%
\pgfsetbuttcap%
\pgfsetmiterjoin%
\definecolor{currentfill}{rgb}{0.000000,0.945098,0.527451}%
\pgfsetfillcolor{currentfill}%
\pgfsetlinewidth{0.000000pt}%
\definecolor{currentstroke}{rgb}{0.000000,0.000000,0.000000}%
\pgfsetstrokecolor{currentstroke}%
\pgfsetstrokeopacity{0.000000}%
\pgfsetdash{}{0pt}%
\pgfpathmoveto{\pgfqpoint{0.605938in}{1.604509in}}%
\pgfpathlineto{\pgfqpoint{0.600479in}{1.609616in}}%
\pgfpathlineto{\pgfqpoint{0.604334in}{1.622797in}}%
\pgfpathlineto{\pgfqpoint{0.600541in}{1.629953in}}%
\pgfpathlineto{\pgfqpoint{0.601624in}{1.635606in}}%
\pgfpathlineto{\pgfqpoint{0.594599in}{1.640016in}}%
\pgfpathlineto{\pgfqpoint{0.584658in}{1.656876in}}%
\pgfpathlineto{\pgfqpoint{0.574470in}{1.662350in}}%
\pgfpathlineto{\pgfqpoint{0.566606in}{1.658522in}}%
\pgfpathlineto{\pgfqpoint{0.558505in}{1.659196in}}%
\pgfpathlineto{\pgfqpoint{0.555953in}{1.662869in}}%
\pgfpathlineto{\pgfqpoint{0.560209in}{1.677284in}}%
\pgfpathlineto{\pgfqpoint{0.544816in}{1.681752in}}%
\pgfpathlineto{\pgfqpoint{0.555668in}{1.716515in}}%
\pgfpathlineto{\pgfqpoint{0.561272in}{1.737785in}}%
\pgfpathlineto{\pgfqpoint{0.494828in}{1.757961in}}%
\pgfpathlineto{\pgfqpoint{0.498294in}{1.770162in}}%
\pgfpathlineto{\pgfqpoint{0.494062in}{1.772814in}}%
\pgfpathlineto{\pgfqpoint{0.477526in}{1.765034in}}%
\pgfpathlineto{\pgfqpoint{0.469737in}{1.765482in}}%
\pgfpathlineto{\pgfqpoint{0.466061in}{1.760966in}}%
\pgfpathlineto{\pgfqpoint{0.468505in}{1.753515in}}%
\pgfpathlineto{\pgfqpoint{0.459959in}{1.754157in}}%
\pgfpathlineto{\pgfqpoint{0.457583in}{1.760831in}}%
\pgfpathlineto{\pgfqpoint{0.449915in}{1.763500in}}%
\pgfpathlineto{\pgfqpoint{0.449403in}{1.767395in}}%
\pgfpathlineto{\pgfqpoint{0.443568in}{1.772460in}}%
\pgfpathlineto{\pgfqpoint{0.442827in}{1.791103in}}%
\pgfpathlineto{\pgfqpoint{0.433774in}{1.794008in}}%
\pgfpathlineto{\pgfqpoint{0.437233in}{1.797086in}}%
\pgfpathlineto{\pgfqpoint{0.438276in}{1.805265in}}%
\pgfpathlineto{\pgfqpoint{0.436698in}{1.812439in}}%
\pgfpathlineto{\pgfqpoint{0.441986in}{1.818955in}}%
\pgfpathlineto{\pgfqpoint{0.442382in}{1.828434in}}%
\pgfpathlineto{\pgfqpoint{0.452336in}{1.831556in}}%
\pgfpathlineto{\pgfqpoint{0.456983in}{1.837942in}}%
\pgfpathlineto{\pgfqpoint{0.526008in}{1.816628in}}%
\pgfpathlineto{\pgfqpoint{0.549594in}{1.891292in}}%
\pgfpathlineto{\pgfqpoint{0.551234in}{1.896670in}}%
\pgfpathlineto{\pgfqpoint{0.562002in}{1.893971in}}%
\pgfpathlineto{\pgfqpoint{0.565189in}{1.899442in}}%
\pgfpathlineto{\pgfqpoint{0.572519in}{1.906519in}}%
\pgfpathlineto{\pgfqpoint{0.572489in}{1.914545in}}%
\pgfpathlineto{\pgfqpoint{0.568235in}{1.922352in}}%
\pgfpathlineto{\pgfqpoint{0.571017in}{1.931177in}}%
\pgfpathlineto{\pgfqpoint{0.579411in}{1.933417in}}%
\pgfpathlineto{\pgfqpoint{0.616244in}{1.922404in}}%
\pgfpathlineto{\pgfqpoint{0.609494in}{1.902912in}}%
\pgfpathlineto{\pgfqpoint{0.595602in}{1.856614in}}%
\pgfpathlineto{\pgfqpoint{0.621562in}{1.848713in}}%
\pgfpathlineto{\pgfqpoint{0.605093in}{1.791654in}}%
\pgfpathlineto{\pgfqpoint{0.654872in}{1.777406in}}%
\pgfpathlineto{\pgfqpoint{0.637413in}{1.715698in}}%
\pgfpathlineto{\pgfqpoint{0.622947in}{1.663573in}}%
\pgfpathlineto{\pgfqpoint{0.605938in}{1.604509in}}%
\pgfpathclose%
\pgfusepath{fill}%
\end{pgfscope}%
\begin{pgfscope}%
\pgfpathrectangle{\pgfqpoint{0.100000in}{0.100000in}}{\pgfqpoint{3.608454in}{2.310000in}}%
\pgfusepath{clip}%
\pgfsetbuttcap%
\pgfsetmiterjoin%
\definecolor{currentfill}{rgb}{0.000000,0.580392,0.709804}%
\pgfsetfillcolor{currentfill}%
\pgfsetlinewidth{0.000000pt}%
\definecolor{currentstroke}{rgb}{0.000000,0.000000,0.000000}%
\pgfsetstrokecolor{currentstroke}%
\pgfsetstrokeopacity{0.000000}%
\pgfsetdash{}{0pt}%
\pgfpathmoveto{\pgfqpoint{2.636183in}{1.558287in}}%
\pgfpathlineto{\pgfqpoint{2.638514in}{1.537637in}}%
\pgfpathlineto{\pgfqpoint{2.614051in}{1.535158in}}%
\pgfpathlineto{\pgfqpoint{2.586580in}{1.532701in}}%
\pgfpathlineto{\pgfqpoint{2.582676in}{1.575047in}}%
\pgfpathlineto{\pgfqpoint{2.596356in}{1.586648in}}%
\pgfpathlineto{\pgfqpoint{2.601592in}{1.595258in}}%
\pgfpathlineto{\pgfqpoint{2.605772in}{1.609448in}}%
\pgfpathlineto{\pgfqpoint{2.612010in}{1.620134in}}%
\pgfpathlineto{\pgfqpoint{2.620303in}{1.620947in}}%
\pgfpathlineto{\pgfqpoint{2.621651in}{1.607409in}}%
\pgfpathlineto{\pgfqpoint{2.648568in}{1.609914in}}%
\pgfpathlineto{\pgfqpoint{2.650897in}{1.588647in}}%
\pgfpathlineto{\pgfqpoint{2.649386in}{1.585311in}}%
\pgfpathlineto{\pgfqpoint{2.633441in}{1.583836in}}%
\pgfpathlineto{\pgfqpoint{2.636183in}{1.558287in}}%
\pgfpathclose%
\pgfusepath{fill}%
\end{pgfscope}%
\begin{pgfscope}%
\pgfpathrectangle{\pgfqpoint{0.100000in}{0.100000in}}{\pgfqpoint{3.608454in}{2.310000in}}%
\pgfusepath{clip}%
\pgfsetbuttcap%
\pgfsetmiterjoin%
\definecolor{currentfill}{rgb}{0.000000,0.525490,0.737255}%
\pgfsetfillcolor{currentfill}%
\pgfsetlinewidth{0.000000pt}%
\definecolor{currentstroke}{rgb}{0.000000,0.000000,0.000000}%
\pgfsetstrokecolor{currentstroke}%
\pgfsetstrokeopacity{0.000000}%
\pgfsetdash{}{0pt}%
\pgfpathmoveto{\pgfqpoint{2.395693in}{0.767722in}}%
\pgfpathlineto{\pgfqpoint{2.396163in}{0.790004in}}%
\pgfpathlineto{\pgfqpoint{2.394278in}{0.809641in}}%
\pgfpathlineto{\pgfqpoint{2.395425in}{0.819431in}}%
\pgfpathlineto{\pgfqpoint{2.399936in}{0.821863in}}%
\pgfpathlineto{\pgfqpoint{2.401669in}{0.830402in}}%
\pgfpathlineto{\pgfqpoint{2.404935in}{0.832654in}}%
\pgfpathlineto{\pgfqpoint{2.412628in}{0.839279in}}%
\pgfpathlineto{\pgfqpoint{2.428148in}{0.846306in}}%
\pgfpathlineto{\pgfqpoint{2.438751in}{0.854506in}}%
\pgfpathlineto{\pgfqpoint{2.442684in}{0.864247in}}%
\pgfpathlineto{\pgfqpoint{2.458381in}{0.865552in}}%
\pgfpathlineto{\pgfqpoint{2.458164in}{0.869135in}}%
\pgfpathlineto{\pgfqpoint{2.485523in}{0.870923in}}%
\pgfpathlineto{\pgfqpoint{2.488805in}{0.814690in}}%
\pgfpathlineto{\pgfqpoint{2.461411in}{0.812905in}}%
\pgfpathlineto{\pgfqpoint{2.462239in}{0.799019in}}%
\pgfpathlineto{\pgfqpoint{2.467492in}{0.796662in}}%
\pgfpathlineto{\pgfqpoint{2.468672in}{0.778139in}}%
\pgfpathlineto{\pgfqpoint{2.447021in}{0.775387in}}%
\pgfpathlineto{\pgfqpoint{2.436793in}{0.774096in}}%
\pgfpathlineto{\pgfqpoint{2.426622in}{0.769362in}}%
\pgfpathlineto{\pgfqpoint{2.395693in}{0.767722in}}%
\pgfpathclose%
\pgfusepath{fill}%
\end{pgfscope}%
\begin{pgfscope}%
\pgfpathrectangle{\pgfqpoint{0.100000in}{0.100000in}}{\pgfqpoint{3.608454in}{2.310000in}}%
\pgfusepath{clip}%
\pgfsetbuttcap%
\pgfsetmiterjoin%
\definecolor{currentfill}{rgb}{0.000000,0.666667,0.666667}%
\pgfsetfillcolor{currentfill}%
\pgfsetlinewidth{0.000000pt}%
\definecolor{currentstroke}{rgb}{0.000000,0.000000,0.000000}%
\pgfsetstrokecolor{currentstroke}%
\pgfsetstrokeopacity{0.000000}%
\pgfsetdash{}{0pt}%
\pgfpathmoveto{\pgfqpoint{2.244140in}{0.763307in}}%
\pgfpathlineto{\pgfqpoint{2.248851in}{0.756615in}}%
\pgfpathlineto{\pgfqpoint{2.256019in}{0.754662in}}%
\pgfpathlineto{\pgfqpoint{2.261771in}{0.748472in}}%
\pgfpathlineto{\pgfqpoint{2.248864in}{0.736251in}}%
\pgfpathlineto{\pgfqpoint{2.244244in}{0.734424in}}%
\pgfpathlineto{\pgfqpoint{2.240752in}{0.735755in}}%
\pgfpathlineto{\pgfqpoint{2.226879in}{0.735487in}}%
\pgfpathlineto{\pgfqpoint{2.226572in}{0.749343in}}%
\pgfpathlineto{\pgfqpoint{2.219726in}{0.749195in}}%
\pgfpathlineto{\pgfqpoint{2.219367in}{0.763229in}}%
\pgfpathlineto{\pgfqpoint{2.212402in}{0.763123in}}%
\pgfpathlineto{\pgfqpoint{2.212262in}{0.773539in}}%
\pgfpathlineto{\pgfqpoint{2.182637in}{0.773027in}}%
\pgfpathlineto{\pgfqpoint{2.173280in}{0.783716in}}%
\pgfpathlineto{\pgfqpoint{2.171444in}{0.784825in}}%
\pgfpathlineto{\pgfqpoint{2.170684in}{0.855388in}}%
\pgfpathlineto{\pgfqpoint{2.170562in}{0.866376in}}%
\pgfpathlineto{\pgfqpoint{2.186423in}{0.866575in}}%
\pgfpathlineto{\pgfqpoint{2.240807in}{0.867381in}}%
\pgfpathlineto{\pgfqpoint{2.258381in}{0.867551in}}%
\pgfpathlineto{\pgfqpoint{2.258858in}{0.847301in}}%
\pgfpathlineto{\pgfqpoint{2.251967in}{0.847074in}}%
\pgfpathlineto{\pgfqpoint{2.248708in}{0.834367in}}%
\pgfpathlineto{\pgfqpoint{2.249143in}{0.822767in}}%
\pgfpathlineto{\pgfqpoint{2.255994in}{0.822878in}}%
\pgfpathlineto{\pgfqpoint{2.256664in}{0.805503in}}%
\pgfpathlineto{\pgfqpoint{2.254067in}{0.798428in}}%
\pgfpathlineto{\pgfqpoint{2.245714in}{0.798315in}}%
\pgfpathlineto{\pgfqpoint{2.249736in}{0.781810in}}%
\pgfpathlineto{\pgfqpoint{2.246031in}{0.774913in}}%
\pgfpathlineto{\pgfqpoint{2.244140in}{0.763307in}}%
\pgfpathclose%
\pgfusepath{fill}%
\end{pgfscope}%
\begin{pgfscope}%
\pgfpathrectangle{\pgfqpoint{0.100000in}{0.100000in}}{\pgfqpoint{3.608454in}{2.310000in}}%
\pgfusepath{clip}%
\pgfsetbuttcap%
\pgfsetmiterjoin%
\definecolor{currentfill}{rgb}{0.000000,0.541176,0.729412}%
\pgfsetfillcolor{currentfill}%
\pgfsetlinewidth{0.000000pt}%
\definecolor{currentstroke}{rgb}{0.000000,0.000000,0.000000}%
\pgfsetstrokecolor{currentstroke}%
\pgfsetstrokeopacity{0.000000}%
\pgfsetdash{}{0pt}%
\pgfpathmoveto{\pgfqpoint{2.859932in}{1.127064in}}%
\pgfpathlineto{\pgfqpoint{2.845692in}{1.114629in}}%
\pgfpathlineto{\pgfqpoint{2.835111in}{1.113784in}}%
\pgfpathlineto{\pgfqpoint{2.833923in}{1.123330in}}%
\pgfpathlineto{\pgfqpoint{2.829784in}{1.127369in}}%
\pgfpathlineto{\pgfqpoint{2.823619in}{1.137722in}}%
\pgfpathlineto{\pgfqpoint{2.831674in}{1.145493in}}%
\pgfpathlineto{\pgfqpoint{2.824792in}{1.161761in}}%
\pgfpathlineto{\pgfqpoint{2.827195in}{1.174638in}}%
\pgfpathlineto{\pgfqpoint{2.836596in}{1.177507in}}%
\pgfpathlineto{\pgfqpoint{2.846368in}{1.183090in}}%
\pgfpathlineto{\pgfqpoint{2.851139in}{1.181704in}}%
\pgfpathlineto{\pgfqpoint{2.853813in}{1.173625in}}%
\pgfpathlineto{\pgfqpoint{2.859396in}{1.177849in}}%
\pgfpathlineto{\pgfqpoint{2.865237in}{1.171385in}}%
\pgfpathlineto{\pgfqpoint{2.860386in}{1.164003in}}%
\pgfpathlineto{\pgfqpoint{2.868331in}{1.159628in}}%
\pgfpathlineto{\pgfqpoint{2.879988in}{1.149617in}}%
\pgfpathlineto{\pgfqpoint{2.877093in}{1.136816in}}%
\pgfpathlineto{\pgfqpoint{2.865280in}{1.133278in}}%
\pgfpathlineto{\pgfqpoint{2.859932in}{1.127064in}}%
\pgfpathclose%
\pgfusepath{fill}%
\end{pgfscope}%
\begin{pgfscope}%
\pgfpathrectangle{\pgfqpoint{0.100000in}{0.100000in}}{\pgfqpoint{3.608454in}{2.310000in}}%
\pgfusepath{clip}%
\pgfsetbuttcap%
\pgfsetmiterjoin%
\definecolor{currentfill}{rgb}{0.000000,0.439216,0.780392}%
\pgfsetfillcolor{currentfill}%
\pgfsetlinewidth{0.000000pt}%
\definecolor{currentstroke}{rgb}{0.000000,0.000000,0.000000}%
\pgfsetstrokecolor{currentstroke}%
\pgfsetstrokeopacity{0.000000}%
\pgfsetdash{}{0pt}%
\pgfpathmoveto{\pgfqpoint{1.291011in}{1.784895in}}%
\pgfpathlineto{\pgfqpoint{1.303529in}{1.783157in}}%
\pgfpathlineto{\pgfqpoint{1.304569in}{1.790081in}}%
\pgfpathlineto{\pgfqpoint{1.314743in}{1.788523in}}%
\pgfpathlineto{\pgfqpoint{1.315604in}{1.795301in}}%
\pgfpathlineto{\pgfqpoint{1.322614in}{1.794241in}}%
\pgfpathlineto{\pgfqpoint{1.323648in}{1.801191in}}%
\pgfpathlineto{\pgfqpoint{1.337257in}{1.799111in}}%
\pgfpathlineto{\pgfqpoint{1.338318in}{1.806097in}}%
\pgfpathlineto{\pgfqpoint{1.417395in}{1.794618in}}%
\pgfpathlineto{\pgfqpoint{1.412417in}{1.742036in}}%
\pgfpathlineto{\pgfqpoint{1.388262in}{1.745418in}}%
\pgfpathlineto{\pgfqpoint{1.384511in}{1.743611in}}%
\pgfpathlineto{\pgfqpoint{1.352534in}{1.748041in}}%
\pgfpathlineto{\pgfqpoint{1.342072in}{1.748707in}}%
\pgfpathlineto{\pgfqpoint{1.322608in}{1.757434in}}%
\pgfpathlineto{\pgfqpoint{1.323149in}{1.760841in}}%
\pgfpathlineto{\pgfqpoint{1.310078in}{1.766325in}}%
\pgfpathlineto{\pgfqpoint{1.296615in}{1.768471in}}%
\pgfpathlineto{\pgfqpoint{1.297605in}{1.774698in}}%
\pgfpathlineto{\pgfqpoint{1.291011in}{1.784895in}}%
\pgfpathclose%
\pgfusepath{fill}%
\end{pgfscope}%
\begin{pgfscope}%
\pgfpathrectangle{\pgfqpoint{0.100000in}{0.100000in}}{\pgfqpoint{3.608454in}{2.310000in}}%
\pgfusepath{clip}%
\pgfsetbuttcap%
\pgfsetmiterjoin%
\definecolor{currentfill}{rgb}{0.000000,0.411765,0.794118}%
\pgfsetfillcolor{currentfill}%
\pgfsetlinewidth{0.000000pt}%
\definecolor{currentstroke}{rgb}{0.000000,0.000000,0.000000}%
\pgfsetstrokecolor{currentstroke}%
\pgfsetstrokeopacity{0.000000}%
\pgfsetdash{}{0pt}%
\pgfpathmoveto{\pgfqpoint{2.522980in}{1.660007in}}%
\pgfpathlineto{\pgfqpoint{2.467930in}{1.655975in}}%
\pgfpathlineto{\pgfqpoint{2.454210in}{1.655405in}}%
\pgfpathlineto{\pgfqpoint{2.450492in}{1.717550in}}%
\pgfpathlineto{\pgfqpoint{2.485186in}{1.719799in}}%
\pgfpathlineto{\pgfqpoint{2.485729in}{1.712888in}}%
\pgfpathlineto{\pgfqpoint{2.499443in}{1.713874in}}%
\pgfpathlineto{\pgfqpoint{2.520499in}{1.715478in}}%
\pgfpathlineto{\pgfqpoint{2.520327in}{1.709190in}}%
\pgfpathlineto{\pgfqpoint{2.517313in}{1.703277in}}%
\pgfpathlineto{\pgfqpoint{2.515606in}{1.691180in}}%
\pgfpathlineto{\pgfqpoint{2.522980in}{1.660007in}}%
\pgfpathclose%
\pgfusepath{fill}%
\end{pgfscope}%
\begin{pgfscope}%
\pgfpathrectangle{\pgfqpoint{0.100000in}{0.100000in}}{\pgfqpoint{3.608454in}{2.310000in}}%
\pgfusepath{clip}%
\pgfsetbuttcap%
\pgfsetmiterjoin%
\definecolor{currentfill}{rgb}{0.000000,0.419608,0.790196}%
\pgfsetfillcolor{currentfill}%
\pgfsetlinewidth{0.000000pt}%
\definecolor{currentstroke}{rgb}{0.000000,0.000000,0.000000}%
\pgfsetstrokecolor{currentstroke}%
\pgfsetstrokeopacity{0.000000}%
\pgfsetdash{}{0pt}%
\pgfpathmoveto{\pgfqpoint{0.721788in}{1.059750in}}%
\pgfpathlineto{\pgfqpoint{0.715639in}{1.032543in}}%
\pgfpathlineto{\pgfqpoint{0.714200in}{1.032869in}}%
\pgfpathlineto{\pgfqpoint{0.705665in}{0.997629in}}%
\pgfpathlineto{\pgfqpoint{0.637733in}{1.007252in}}%
\pgfpathlineto{\pgfqpoint{0.638583in}{1.013548in}}%
\pgfpathlineto{\pgfqpoint{0.632350in}{1.019382in}}%
\pgfpathlineto{\pgfqpoint{0.635822in}{1.034911in}}%
\pgfpathlineto{\pgfqpoint{0.636748in}{1.045662in}}%
\pgfpathlineto{\pgfqpoint{0.634791in}{1.058910in}}%
\pgfpathlineto{\pgfqpoint{0.628784in}{1.074472in}}%
\pgfpathlineto{\pgfqpoint{0.623749in}{1.080645in}}%
\pgfpathlineto{\pgfqpoint{0.626202in}{1.085453in}}%
\pgfpathlineto{\pgfqpoint{0.631612in}{1.088268in}}%
\pgfpathlineto{\pgfqpoint{0.640257in}{1.084919in}}%
\pgfpathlineto{\pgfqpoint{0.647432in}{1.078308in}}%
\pgfpathlineto{\pgfqpoint{0.660882in}{1.074556in}}%
\pgfpathlineto{\pgfqpoint{0.721788in}{1.059750in}}%
\pgfpathclose%
\pgfusepath{fill}%
\end{pgfscope}%
\begin{pgfscope}%
\pgfpathrectangle{\pgfqpoint{0.100000in}{0.100000in}}{\pgfqpoint{3.608454in}{2.310000in}}%
\pgfusepath{clip}%
\pgfsetbuttcap%
\pgfsetmiterjoin%
\definecolor{currentfill}{rgb}{0.000000,0.713725,0.643137}%
\pgfsetfillcolor{currentfill}%
\pgfsetlinewidth{0.000000pt}%
\definecolor{currentstroke}{rgb}{0.000000,0.000000,0.000000}%
\pgfsetstrokecolor{currentstroke}%
\pgfsetstrokeopacity{0.000000}%
\pgfsetdash{}{0pt}%
\pgfpathmoveto{\pgfqpoint{2.380907in}{1.900334in}}%
\pgfpathlineto{\pgfqpoint{2.366647in}{1.899528in}}%
\pgfpathlineto{\pgfqpoint{2.365878in}{1.913327in}}%
\pgfpathlineto{\pgfqpoint{2.359091in}{1.912916in}}%
\pgfpathlineto{\pgfqpoint{2.358761in}{1.919792in}}%
\pgfpathlineto{\pgfqpoint{2.351813in}{1.919479in}}%
\pgfpathlineto{\pgfqpoint{2.350608in}{1.946791in}}%
\pgfpathlineto{\pgfqpoint{2.356731in}{1.945339in}}%
\pgfpathlineto{\pgfqpoint{2.361222in}{1.948386in}}%
\pgfpathlineto{\pgfqpoint{2.378630in}{1.955337in}}%
\pgfpathlineto{\pgfqpoint{2.391147in}{1.967641in}}%
\pgfpathlineto{\pgfqpoint{2.411560in}{1.970914in}}%
\pgfpathlineto{\pgfqpoint{2.420196in}{1.976038in}}%
\pgfpathlineto{\pgfqpoint{2.425368in}{1.983108in}}%
\pgfpathlineto{\pgfqpoint{2.433892in}{1.984602in}}%
\pgfpathlineto{\pgfqpoint{2.436662in}{1.987686in}}%
\pgfpathlineto{\pgfqpoint{2.438144in}{1.966329in}}%
\pgfpathlineto{\pgfqpoint{2.442299in}{1.959772in}}%
\pgfpathlineto{\pgfqpoint{2.435480in}{1.959378in}}%
\pgfpathlineto{\pgfqpoint{2.436761in}{1.938649in}}%
\pgfpathlineto{\pgfqpoint{2.438472in}{1.913015in}}%
\pgfpathlineto{\pgfqpoint{2.432701in}{1.915926in}}%
\pgfpathlineto{\pgfqpoint{2.385839in}{1.925963in}}%
\pgfpathlineto{\pgfqpoint{2.387268in}{1.900687in}}%
\pgfpathlineto{\pgfqpoint{2.380907in}{1.900334in}}%
\pgfpathclose%
\pgfusepath{fill}%
\end{pgfscope}%
\begin{pgfscope}%
\pgfpathrectangle{\pgfqpoint{0.100000in}{0.100000in}}{\pgfqpoint{3.608454in}{2.310000in}}%
\pgfusepath{clip}%
\pgfsetbuttcap%
\pgfsetmiterjoin%
\definecolor{currentfill}{rgb}{0.000000,0.403922,0.798039}%
\pgfsetfillcolor{currentfill}%
\pgfsetlinewidth{0.000000pt}%
\definecolor{currentstroke}{rgb}{0.000000,0.000000,0.000000}%
\pgfsetstrokecolor{currentstroke}%
\pgfsetstrokeopacity{0.000000}%
\pgfsetdash{}{0pt}%
\pgfpathmoveto{\pgfqpoint{1.577074in}{1.409804in}}%
\pgfpathlineto{\pgfqpoint{1.610603in}{1.406638in}}%
\pgfpathlineto{\pgfqpoint{1.606968in}{1.372189in}}%
\pgfpathlineto{\pgfqpoint{1.603799in}{1.337764in}}%
\pgfpathlineto{\pgfqpoint{1.602094in}{1.324825in}}%
\pgfpathlineto{\pgfqpoint{1.581597in}{1.326016in}}%
\pgfpathlineto{\pgfqpoint{1.547834in}{1.329763in}}%
\pgfpathlineto{\pgfqpoint{1.550467in}{1.357020in}}%
\pgfpathlineto{\pgfqpoint{1.571014in}{1.354845in}}%
\pgfpathlineto{\pgfqpoint{1.577074in}{1.409804in}}%
\pgfpathclose%
\pgfusepath{fill}%
\end{pgfscope}%
\begin{pgfscope}%
\pgfpathrectangle{\pgfqpoint{0.100000in}{0.100000in}}{\pgfqpoint{3.608454in}{2.310000in}}%
\pgfusepath{clip}%
\pgfsetbuttcap%
\pgfsetmiterjoin%
\definecolor{currentfill}{rgb}{0.000000,0.537255,0.731373}%
\pgfsetfillcolor{currentfill}%
\pgfsetlinewidth{0.000000pt}%
\definecolor{currentstroke}{rgb}{0.000000,0.000000,0.000000}%
\pgfsetstrokecolor{currentstroke}%
\pgfsetstrokeopacity{0.000000}%
\pgfsetdash{}{0pt}%
\pgfpathmoveto{\pgfqpoint{2.101955in}{1.263584in}}%
\pgfpathlineto{\pgfqpoint{2.101288in}{1.239463in}}%
\pgfpathlineto{\pgfqpoint{2.046716in}{1.239913in}}%
\pgfpathlineto{\pgfqpoint{2.047067in}{1.264108in}}%
\pgfpathlineto{\pgfqpoint{2.074465in}{1.263717in}}%
\pgfpathlineto{\pgfqpoint{2.075326in}{1.295024in}}%
\pgfpathlineto{\pgfqpoint{2.075903in}{1.319066in}}%
\pgfpathlineto{\pgfqpoint{2.103312in}{1.318980in}}%
\pgfpathlineto{\pgfqpoint{2.102737in}{1.315533in}}%
\pgfpathlineto{\pgfqpoint{2.102752in}{1.281057in}}%
\pgfpathlineto{\pgfqpoint{2.101955in}{1.263584in}}%
\pgfpathclose%
\pgfusepath{fill}%
\end{pgfscope}%
\begin{pgfscope}%
\pgfpathrectangle{\pgfqpoint{0.100000in}{0.100000in}}{\pgfqpoint{3.608454in}{2.310000in}}%
\pgfusepath{clip}%
\pgfsetbuttcap%
\pgfsetmiterjoin%
\definecolor{currentfill}{rgb}{0.000000,0.517647,0.741176}%
\pgfsetfillcolor{currentfill}%
\pgfsetlinewidth{0.000000pt}%
\definecolor{currentstroke}{rgb}{0.000000,0.000000,0.000000}%
\pgfsetstrokecolor{currentstroke}%
\pgfsetstrokeopacity{0.000000}%
\pgfsetdash{}{0pt}%
\pgfpathmoveto{\pgfqpoint{2.356072in}{1.695517in}}%
\pgfpathlineto{\pgfqpoint{2.325187in}{1.694132in}}%
\pgfpathlineto{\pgfqpoint{2.324231in}{1.700249in}}%
\pgfpathlineto{\pgfqpoint{2.301639in}{1.699356in}}%
\pgfpathlineto{\pgfqpoint{2.294773in}{1.699116in}}%
\pgfpathlineto{\pgfqpoint{2.293781in}{1.726557in}}%
\pgfpathlineto{\pgfqpoint{2.319291in}{1.727547in}}%
\pgfpathlineto{\pgfqpoint{2.310801in}{1.738104in}}%
\pgfpathlineto{\pgfqpoint{2.316678in}{1.738350in}}%
\pgfpathlineto{\pgfqpoint{2.317255in}{1.744761in}}%
\pgfpathlineto{\pgfqpoint{2.336281in}{1.746036in}}%
\pgfpathlineto{\pgfqpoint{2.339782in}{1.753146in}}%
\pgfpathlineto{\pgfqpoint{2.373549in}{1.754537in}}%
\pgfpathlineto{\pgfqpoint{2.373196in}{1.761948in}}%
\pgfpathlineto{\pgfqpoint{2.396439in}{1.763257in}}%
\pgfpathlineto{\pgfqpoint{2.389968in}{1.750300in}}%
\pgfpathlineto{\pgfqpoint{2.394954in}{1.739243in}}%
\pgfpathlineto{\pgfqpoint{2.395001in}{1.732218in}}%
\pgfpathlineto{\pgfqpoint{2.400634in}{1.726438in}}%
\pgfpathlineto{\pgfqpoint{2.405834in}{1.715444in}}%
\pgfpathlineto{\pgfqpoint{2.375661in}{1.713794in}}%
\pgfpathlineto{\pgfqpoint{2.376100in}{1.706906in}}%
\pgfpathlineto{\pgfqpoint{2.355403in}{1.705801in}}%
\pgfpathlineto{\pgfqpoint{2.356072in}{1.695517in}}%
\pgfpathclose%
\pgfusepath{fill}%
\end{pgfscope}%
\begin{pgfscope}%
\pgfpathrectangle{\pgfqpoint{0.100000in}{0.100000in}}{\pgfqpoint{3.608454in}{2.310000in}}%
\pgfusepath{clip}%
\pgfsetbuttcap%
\pgfsetmiterjoin%
\definecolor{currentfill}{rgb}{0.000000,0.603922,0.698039}%
\pgfsetfillcolor{currentfill}%
\pgfsetlinewidth{0.000000pt}%
\definecolor{currentstroke}{rgb}{0.000000,0.000000,0.000000}%
\pgfsetstrokecolor{currentstroke}%
\pgfsetstrokeopacity{0.000000}%
\pgfsetdash{}{0pt}%
\pgfpathmoveto{\pgfqpoint{1.234540in}{2.143503in}}%
\pgfpathlineto{\pgfqpoint{1.232482in}{2.132070in}}%
\pgfpathlineto{\pgfqpoint{1.202148in}{2.137522in}}%
\pgfpathlineto{\pgfqpoint{1.200330in}{2.146187in}}%
\pgfpathlineto{\pgfqpoint{1.194140in}{2.150958in}}%
\pgfpathlineto{\pgfqpoint{1.150149in}{2.159708in}}%
\pgfpathlineto{\pgfqpoint{1.139517in}{2.164977in}}%
\pgfpathlineto{\pgfqpoint{1.140465in}{2.171053in}}%
\pgfpathlineto{\pgfqpoint{1.134122in}{2.180102in}}%
\pgfpathlineto{\pgfqpoint{1.136982in}{2.181768in}}%
\pgfpathlineto{\pgfqpoint{1.131142in}{2.190334in}}%
\pgfpathlineto{\pgfqpoint{1.133270in}{2.197089in}}%
\pgfpathlineto{\pgfqpoint{1.119541in}{2.205769in}}%
\pgfpathlineto{\pgfqpoint{1.123862in}{2.213141in}}%
\pgfpathlineto{\pgfqpoint{1.118945in}{2.225334in}}%
\pgfpathlineto{\pgfqpoint{1.112359in}{2.223590in}}%
\pgfpathlineto{\pgfqpoint{1.109470in}{2.229400in}}%
\pgfpathlineto{\pgfqpoint{1.109601in}{2.239048in}}%
\pgfpathlineto{\pgfqpoint{1.157261in}{2.229106in}}%
\pgfpathlineto{\pgfqpoint{1.193179in}{2.222014in}}%
\pgfpathlineto{\pgfqpoint{1.254169in}{2.210455in}}%
\pgfpathlineto{\pgfqpoint{1.242847in}{2.149099in}}%
\pgfpathlineto{\pgfqpoint{1.235802in}{2.150383in}}%
\pgfpathlineto{\pgfqpoint{1.234540in}{2.143503in}}%
\pgfpathclose%
\pgfusepath{fill}%
\end{pgfscope}%
\begin{pgfscope}%
\pgfpathrectangle{\pgfqpoint{0.100000in}{0.100000in}}{\pgfqpoint{3.608454in}{2.310000in}}%
\pgfusepath{clip}%
\pgfsetbuttcap%
\pgfsetmiterjoin%
\definecolor{currentfill}{rgb}{0.000000,0.380392,0.809804}%
\pgfsetfillcolor{currentfill}%
\pgfsetlinewidth{0.000000pt}%
\definecolor{currentstroke}{rgb}{0.000000,0.000000,0.000000}%
\pgfsetstrokecolor{currentstroke}%
\pgfsetstrokeopacity{0.000000}%
\pgfsetdash{}{0pt}%
\pgfpathmoveto{\pgfqpoint{1.807069in}{2.050968in}}%
\pgfpathlineto{\pgfqpoint{1.738152in}{2.055718in}}%
\pgfpathlineto{\pgfqpoint{1.739211in}{2.069716in}}%
\pgfpathlineto{\pgfqpoint{1.736501in}{2.069931in}}%
\pgfpathlineto{\pgfqpoint{1.738640in}{2.097811in}}%
\pgfpathlineto{\pgfqpoint{1.735922in}{2.098000in}}%
\pgfpathlineto{\pgfqpoint{1.737002in}{2.111989in}}%
\pgfpathlineto{\pgfqpoint{1.723190in}{2.112972in}}%
\pgfpathlineto{\pgfqpoint{1.724331in}{2.126897in}}%
\pgfpathlineto{\pgfqpoint{1.728606in}{2.126565in}}%
\pgfpathlineto{\pgfqpoint{1.729180in}{2.133512in}}%
\pgfpathlineto{\pgfqpoint{1.736050in}{2.132965in}}%
\pgfpathlineto{\pgfqpoint{1.737269in}{2.148337in}}%
\pgfpathlineto{\pgfqpoint{1.790039in}{2.144561in}}%
\pgfpathlineto{\pgfqpoint{1.833904in}{2.141855in}}%
\pgfpathlineto{\pgfqpoint{1.832606in}{2.119399in}}%
\pgfpathlineto{\pgfqpoint{1.833772in}{2.105376in}}%
\pgfpathlineto{\pgfqpoint{1.826839in}{2.105738in}}%
\pgfpathlineto{\pgfqpoint{1.828130in}{2.091731in}}%
\pgfpathlineto{\pgfqpoint{1.826539in}{2.063794in}}%
\pgfpathlineto{\pgfqpoint{1.827828in}{2.049744in}}%
\pgfpathlineto{\pgfqpoint{1.807069in}{2.050968in}}%
\pgfpathclose%
\pgfusepath{fill}%
\end{pgfscope}%
\begin{pgfscope}%
\pgfpathrectangle{\pgfqpoint{0.100000in}{0.100000in}}{\pgfqpoint{3.608454in}{2.310000in}}%
\pgfusepath{clip}%
\pgfsetbuttcap%
\pgfsetmiterjoin%
\definecolor{currentfill}{rgb}{0.000000,0.639216,0.680392}%
\pgfsetfillcolor{currentfill}%
\pgfsetlinewidth{0.000000pt}%
\definecolor{currentstroke}{rgb}{0.000000,0.000000,0.000000}%
\pgfsetstrokecolor{currentstroke}%
\pgfsetstrokeopacity{0.000000}%
\pgfsetdash{}{0pt}%
\pgfpathmoveto{\pgfqpoint{1.889333in}{1.734748in}}%
\pgfpathlineto{\pgfqpoint{1.888636in}{1.700186in}}%
\pgfpathlineto{\pgfqpoint{1.859723in}{1.701494in}}%
\pgfpathlineto{\pgfqpoint{1.841021in}{1.702423in}}%
\pgfpathlineto{\pgfqpoint{1.842006in}{1.721715in}}%
\pgfpathlineto{\pgfqpoint{1.829398in}{1.718027in}}%
\pgfpathlineto{\pgfqpoint{1.817373in}{1.718961in}}%
\pgfpathlineto{\pgfqpoint{1.807130in}{1.721260in}}%
\pgfpathlineto{\pgfqpoint{1.800957in}{1.721865in}}%
\pgfpathlineto{\pgfqpoint{1.800009in}{1.732253in}}%
\pgfpathlineto{\pgfqpoint{1.801456in}{1.757724in}}%
\pgfpathlineto{\pgfqpoint{1.801576in}{1.759791in}}%
\pgfpathlineto{\pgfqpoint{1.825951in}{1.758455in}}%
\pgfpathlineto{\pgfqpoint{1.830770in}{1.752552in}}%
\pgfpathlineto{\pgfqpoint{1.835653in}{1.753170in}}%
\pgfpathlineto{\pgfqpoint{1.843095in}{1.749824in}}%
\pgfpathlineto{\pgfqpoint{1.841646in}{1.759396in}}%
\pgfpathlineto{\pgfqpoint{1.846482in}{1.757163in}}%
\pgfpathlineto{\pgfqpoint{1.883458in}{1.755759in}}%
\pgfpathlineto{\pgfqpoint{1.882514in}{1.735032in}}%
\pgfpathlineto{\pgfqpoint{1.889333in}{1.734748in}}%
\pgfpathclose%
\pgfusepath{fill}%
\end{pgfscope}%
\begin{pgfscope}%
\pgfpathrectangle{\pgfqpoint{0.100000in}{0.100000in}}{\pgfqpoint{3.608454in}{2.310000in}}%
\pgfusepath{clip}%
\pgfsetbuttcap%
\pgfsetmiterjoin%
\definecolor{currentfill}{rgb}{0.000000,0.925490,0.537255}%
\pgfsetfillcolor{currentfill}%
\pgfsetlinewidth{0.000000pt}%
\definecolor{currentstroke}{rgb}{0.000000,0.000000,0.000000}%
\pgfsetstrokecolor{currentstroke}%
\pgfsetstrokeopacity{0.000000}%
\pgfsetdash{}{0pt}%
\pgfpathmoveto{\pgfqpoint{2.895571in}{1.311319in}}%
\pgfpathlineto{\pgfqpoint{2.902369in}{1.314765in}}%
\pgfpathlineto{\pgfqpoint{2.908038in}{1.306264in}}%
\pgfpathlineto{\pgfqpoint{2.909448in}{1.297631in}}%
\pgfpathlineto{\pgfqpoint{2.916184in}{1.292250in}}%
\pgfpathlineto{\pgfqpoint{2.929879in}{1.293625in}}%
\pgfpathlineto{\pgfqpoint{2.927867in}{1.284414in}}%
\pgfpathlineto{\pgfqpoint{2.923822in}{1.280887in}}%
\pgfpathlineto{\pgfqpoint{2.921023in}{1.282534in}}%
\pgfpathlineto{\pgfqpoint{2.902185in}{1.260526in}}%
\pgfpathlineto{\pgfqpoint{2.887645in}{1.250541in}}%
\pgfpathlineto{\pgfqpoint{2.880671in}{1.252069in}}%
\pgfpathlineto{\pgfqpoint{2.874724in}{1.257922in}}%
\pgfpathlineto{\pgfqpoint{2.872542in}{1.266952in}}%
\pgfpathlineto{\pgfqpoint{2.860524in}{1.271364in}}%
\pgfpathlineto{\pgfqpoint{2.849787in}{1.279563in}}%
\pgfpathlineto{\pgfqpoint{2.839446in}{1.282704in}}%
\pgfpathlineto{\pgfqpoint{2.837862in}{1.289352in}}%
\pgfpathlineto{\pgfqpoint{2.847809in}{1.300549in}}%
\pgfpathlineto{\pgfqpoint{2.855369in}{1.302588in}}%
\pgfpathlineto{\pgfqpoint{2.850698in}{1.310252in}}%
\pgfpathlineto{\pgfqpoint{2.857831in}{1.319548in}}%
\pgfpathlineto{\pgfqpoint{2.855274in}{1.324556in}}%
\pgfpathlineto{\pgfqpoint{2.861780in}{1.330306in}}%
\pgfpathlineto{\pgfqpoint{2.874015in}{1.332782in}}%
\pgfpathlineto{\pgfqpoint{2.873043in}{1.323790in}}%
\pgfpathlineto{\pgfqpoint{2.885378in}{1.313095in}}%
\pgfpathlineto{\pgfqpoint{2.884060in}{1.308388in}}%
\pgfpathlineto{\pgfqpoint{2.890443in}{1.304155in}}%
\pgfpathlineto{\pgfqpoint{2.895571in}{1.311319in}}%
\pgfpathclose%
\pgfusepath{fill}%
\end{pgfscope}%
\begin{pgfscope}%
\pgfpathrectangle{\pgfqpoint{0.100000in}{0.100000in}}{\pgfqpoint{3.608454in}{2.310000in}}%
\pgfusepath{clip}%
\pgfsetbuttcap%
\pgfsetmiterjoin%
\definecolor{currentfill}{rgb}{0.000000,0.462745,0.768627}%
\pgfsetfillcolor{currentfill}%
\pgfsetlinewidth{0.000000pt}%
\definecolor{currentstroke}{rgb}{0.000000,0.000000,0.000000}%
\pgfsetstrokecolor{currentstroke}%
\pgfsetstrokeopacity{0.000000}%
\pgfsetdash{}{0pt}%
\pgfpathmoveto{\pgfqpoint{2.193901in}{1.649520in}}%
\pgfpathlineto{\pgfqpoint{2.166463in}{1.649078in}}%
\pgfpathlineto{\pgfqpoint{2.165979in}{1.695985in}}%
\pgfpathlineto{\pgfqpoint{2.184501in}{1.696223in}}%
\pgfpathlineto{\pgfqpoint{2.220391in}{1.696881in}}%
\pgfpathlineto{\pgfqpoint{2.221382in}{1.649990in}}%
\pgfpathlineto{\pgfqpoint{2.193901in}{1.649520in}}%
\pgfpathclose%
\pgfusepath{fill}%
\end{pgfscope}%
\begin{pgfscope}%
\pgfpathrectangle{\pgfqpoint{0.100000in}{0.100000in}}{\pgfqpoint{3.608454in}{2.310000in}}%
\pgfusepath{clip}%
\pgfsetbuttcap%
\pgfsetmiterjoin%
\definecolor{currentfill}{rgb}{0.000000,0.768627,0.615686}%
\pgfsetfillcolor{currentfill}%
\pgfsetlinewidth{0.000000pt}%
\definecolor{currentstroke}{rgb}{0.000000,0.000000,0.000000}%
\pgfsetstrokecolor{currentstroke}%
\pgfsetstrokeopacity{0.000000}%
\pgfsetdash{}{0pt}%
\pgfpathmoveto{\pgfqpoint{2.936140in}{1.370645in}}%
\pgfpathlineto{\pgfqpoint{2.930069in}{1.365085in}}%
\pgfpathlineto{\pgfqpoint{2.920044in}{1.373705in}}%
\pgfpathlineto{\pgfqpoint{2.910035in}{1.386821in}}%
\pgfpathlineto{\pgfqpoint{2.916733in}{1.399937in}}%
\pgfpathlineto{\pgfqpoint{2.915344in}{1.412906in}}%
\pgfpathlineto{\pgfqpoint{2.916394in}{1.420191in}}%
\pgfpathlineto{\pgfqpoint{2.907865in}{1.422811in}}%
\pgfpathlineto{\pgfqpoint{2.906878in}{1.433220in}}%
\pgfpathlineto{\pgfqpoint{2.907685in}{1.436860in}}%
\pgfpathlineto{\pgfqpoint{2.914756in}{1.436835in}}%
\pgfpathlineto{\pgfqpoint{2.914676in}{1.445590in}}%
\pgfpathlineto{\pgfqpoint{2.927738in}{1.446958in}}%
\pgfpathlineto{\pgfqpoint{2.938939in}{1.450447in}}%
\pgfpathlineto{\pgfqpoint{2.941441in}{1.448688in}}%
\pgfpathlineto{\pgfqpoint{2.954171in}{1.450091in}}%
\pgfpathlineto{\pgfqpoint{2.951761in}{1.443677in}}%
\pgfpathlineto{\pgfqpoint{2.962415in}{1.436530in}}%
\pgfpathlineto{\pgfqpoint{2.966480in}{1.429783in}}%
\pgfpathlineto{\pgfqpoint{2.965224in}{1.427879in}}%
\pgfpathlineto{\pgfqpoint{2.973331in}{1.416066in}}%
\pgfpathlineto{\pgfqpoint{2.967704in}{1.409901in}}%
\pgfpathlineto{\pgfqpoint{2.961168in}{1.406181in}}%
\pgfpathlineto{\pgfqpoint{2.954499in}{1.407993in}}%
\pgfpathlineto{\pgfqpoint{2.943680in}{1.394391in}}%
\pgfpathlineto{\pgfqpoint{2.936799in}{1.397031in}}%
\pgfpathlineto{\pgfqpoint{2.932362in}{1.392633in}}%
\pgfpathlineto{\pgfqpoint{2.933951in}{1.375085in}}%
\pgfpathlineto{\pgfqpoint{2.936140in}{1.370645in}}%
\pgfpathclose%
\pgfusepath{fill}%
\end{pgfscope}%
\begin{pgfscope}%
\pgfpathrectangle{\pgfqpoint{0.100000in}{0.100000in}}{\pgfqpoint{3.608454in}{2.310000in}}%
\pgfusepath{clip}%
\pgfsetbuttcap%
\pgfsetmiterjoin%
\definecolor{currentfill}{rgb}{0.000000,0.431373,0.784314}%
\pgfsetfillcolor{currentfill}%
\pgfsetlinewidth{0.000000pt}%
\definecolor{currentstroke}{rgb}{0.000000,0.000000,0.000000}%
\pgfsetstrokecolor{currentstroke}%
\pgfsetstrokeopacity{0.000000}%
\pgfsetdash{}{0pt}%
\pgfpathmoveto{\pgfqpoint{2.258667in}{1.780519in}}%
\pgfpathlineto{\pgfqpoint{2.268903in}{1.780718in}}%
\pgfpathlineto{\pgfqpoint{2.267401in}{1.826914in}}%
\pgfpathlineto{\pgfqpoint{2.266038in}{1.833884in}}%
\pgfpathlineto{\pgfqpoint{2.300433in}{1.834842in}}%
\pgfpathlineto{\pgfqpoint{2.300225in}{1.841636in}}%
\pgfpathlineto{\pgfqpoint{2.334657in}{1.843071in}}%
\pgfpathlineto{\pgfqpoint{2.337567in}{1.780978in}}%
\pgfpathlineto{\pgfqpoint{2.344426in}{1.781264in}}%
\pgfpathlineto{\pgfqpoint{2.344731in}{1.774369in}}%
\pgfpathlineto{\pgfqpoint{2.372249in}{1.775882in}}%
\pgfpathlineto{\pgfqpoint{2.373196in}{1.761948in}}%
\pgfpathlineto{\pgfqpoint{2.373549in}{1.754537in}}%
\pgfpathlineto{\pgfqpoint{2.339782in}{1.753146in}}%
\pgfpathlineto{\pgfqpoint{2.336281in}{1.746036in}}%
\pgfpathlineto{\pgfqpoint{2.317255in}{1.744761in}}%
\pgfpathlineto{\pgfqpoint{2.316678in}{1.738350in}}%
\pgfpathlineto{\pgfqpoint{2.310801in}{1.738104in}}%
\pgfpathlineto{\pgfqpoint{2.319291in}{1.727547in}}%
\pgfpathlineto{\pgfqpoint{2.293781in}{1.726557in}}%
\pgfpathlineto{\pgfqpoint{2.273836in}{1.725888in}}%
\pgfpathlineto{\pgfqpoint{2.272950in}{1.753323in}}%
\pgfpathlineto{\pgfqpoint{2.285525in}{1.753739in}}%
\pgfpathlineto{\pgfqpoint{2.281318in}{1.764757in}}%
\pgfpathlineto{\pgfqpoint{2.273713in}{1.770000in}}%
\pgfpathlineto{\pgfqpoint{2.264630in}{1.772555in}}%
\pgfpathlineto{\pgfqpoint{2.258667in}{1.780519in}}%
\pgfpathclose%
\pgfusepath{fill}%
\end{pgfscope}%
\begin{pgfscope}%
\pgfpathrectangle{\pgfqpoint{0.100000in}{0.100000in}}{\pgfqpoint{3.608454in}{2.310000in}}%
\pgfusepath{clip}%
\pgfsetbuttcap%
\pgfsetmiterjoin%
\definecolor{currentfill}{rgb}{0.000000,0.678431,0.660784}%
\pgfsetfillcolor{currentfill}%
\pgfsetlinewidth{0.000000pt}%
\definecolor{currentstroke}{rgb}{0.000000,0.000000,0.000000}%
\pgfsetstrokecolor{currentstroke}%
\pgfsetstrokeopacity{0.000000}%
\pgfsetdash{}{0pt}%
\pgfpathmoveto{\pgfqpoint{1.938836in}{0.498314in}}%
\pgfpathlineto{\pgfqpoint{1.924683in}{0.477880in}}%
\pgfpathlineto{\pgfqpoint{1.907060in}{0.482191in}}%
\pgfpathlineto{\pgfqpoint{1.907719in}{0.482706in}}%
\pgfpathlineto{\pgfqpoint{1.889121in}{0.521871in}}%
\pgfpathlineto{\pgfqpoint{1.895145in}{0.523860in}}%
\pgfpathlineto{\pgfqpoint{1.882610in}{0.539717in}}%
\pgfpathlineto{\pgfqpoint{1.915913in}{0.565873in}}%
\pgfpathlineto{\pgfqpoint{1.923720in}{0.556660in}}%
\pgfpathlineto{\pgfqpoint{1.913527in}{0.548723in}}%
\pgfpathlineto{\pgfqpoint{1.925742in}{0.532805in}}%
\pgfpathlineto{\pgfqpoint{1.911040in}{0.521565in}}%
\pgfpathlineto{\pgfqpoint{1.916523in}{0.511256in}}%
\pgfpathlineto{\pgfqpoint{1.926187in}{0.508771in}}%
\pgfpathlineto{\pgfqpoint{1.926522in}{0.503678in}}%
\pgfpathlineto{\pgfqpoint{1.938836in}{0.498314in}}%
\pgfpathclose%
\pgfusepath{fill}%
\end{pgfscope}%
\begin{pgfscope}%
\pgfpathrectangle{\pgfqpoint{0.100000in}{0.100000in}}{\pgfqpoint{3.608454in}{2.310000in}}%
\pgfusepath{clip}%
\pgfsetbuttcap%
\pgfsetmiterjoin%
\definecolor{currentfill}{rgb}{0.000000,0.596078,0.701961}%
\pgfsetfillcolor{currentfill}%
\pgfsetlinewidth{0.000000pt}%
\definecolor{currentstroke}{rgb}{0.000000,0.000000,0.000000}%
\pgfsetstrokecolor{currentstroke}%
\pgfsetstrokeopacity{0.000000}%
\pgfsetdash{}{0pt}%
\pgfpathmoveto{\pgfqpoint{2.999992in}{0.934399in}}%
\pgfpathlineto{\pgfqpoint{2.986683in}{0.921479in}}%
\pgfpathlineto{\pgfqpoint{2.980224in}{0.924118in}}%
\pgfpathlineto{\pgfqpoint{2.965871in}{0.919067in}}%
\pgfpathlineto{\pgfqpoint{2.962570in}{0.910798in}}%
\pgfpathlineto{\pgfqpoint{2.959107in}{0.909040in}}%
\pgfpathlineto{\pgfqpoint{2.950977in}{0.910809in}}%
\pgfpathlineto{\pgfqpoint{2.944266in}{0.903605in}}%
\pgfpathlineto{\pgfqpoint{2.937410in}{0.908377in}}%
\pgfpathlineto{\pgfqpoint{2.937493in}{0.915146in}}%
\pgfpathlineto{\pgfqpoint{2.933286in}{0.922296in}}%
\pgfpathlineto{\pgfqpoint{2.925307in}{0.930408in}}%
\pgfpathlineto{\pgfqpoint{2.919634in}{0.933284in}}%
\pgfpathlineto{\pgfqpoint{2.918368in}{0.938239in}}%
\pgfpathlineto{\pgfqpoint{2.909363in}{0.953452in}}%
\pgfpathlineto{\pgfqpoint{2.910547in}{0.958327in}}%
\pgfpathlineto{\pgfqpoint{2.915431in}{0.959683in}}%
\pgfpathlineto{\pgfqpoint{2.918991in}{0.967320in}}%
\pgfpathlineto{\pgfqpoint{2.922909in}{0.970556in}}%
\pgfpathlineto{\pgfqpoint{2.932420in}{0.972207in}}%
\pgfpathlineto{\pgfqpoint{2.942345in}{0.978613in}}%
\pgfpathlineto{\pgfqpoint{2.948243in}{0.978295in}}%
\pgfpathlineto{\pgfqpoint{2.954883in}{0.971559in}}%
\pgfpathlineto{\pgfqpoint{2.956878in}{0.975335in}}%
\pgfpathlineto{\pgfqpoint{2.950161in}{0.981612in}}%
\pgfpathlineto{\pgfqpoint{2.953478in}{0.990526in}}%
\pgfpathlineto{\pgfqpoint{2.949644in}{0.997974in}}%
\pgfpathlineto{\pgfqpoint{2.958909in}{1.001900in}}%
\pgfpathlineto{\pgfqpoint{2.966134in}{1.004041in}}%
\pgfpathlineto{\pgfqpoint{2.971145in}{0.996030in}}%
\pgfpathlineto{\pgfqpoint{2.983807in}{0.993720in}}%
\pgfpathlineto{\pgfqpoint{2.988314in}{0.999433in}}%
\pgfpathlineto{\pgfqpoint{2.999338in}{0.990301in}}%
\pgfpathlineto{\pgfqpoint{3.016045in}{0.985596in}}%
\pgfpathlineto{\pgfqpoint{3.005840in}{0.971004in}}%
\pgfpathlineto{\pgfqpoint{2.984051in}{0.944289in}}%
\pgfpathlineto{\pgfqpoint{2.988444in}{0.938374in}}%
\pgfpathlineto{\pgfqpoint{2.999992in}{0.934399in}}%
\pgfpathclose%
\pgfusepath{fill}%
\end{pgfscope}%
\begin{pgfscope}%
\pgfpathrectangle{\pgfqpoint{0.100000in}{0.100000in}}{\pgfqpoint{3.608454in}{2.310000in}}%
\pgfusepath{clip}%
\pgfsetbuttcap%
\pgfsetmiterjoin%
\definecolor{currentfill}{rgb}{0.000000,0.466667,0.766667}%
\pgfsetfillcolor{currentfill}%
\pgfsetlinewidth{0.000000pt}%
\definecolor{currentstroke}{rgb}{0.000000,0.000000,0.000000}%
\pgfsetstrokecolor{currentstroke}%
\pgfsetstrokeopacity{0.000000}%
\pgfsetdash{}{0pt}%
\pgfpathmoveto{\pgfqpoint{3.243003in}{1.088398in}}%
\pgfpathlineto{\pgfqpoint{3.237815in}{1.092298in}}%
\pgfpathlineto{\pgfqpoint{3.228319in}{1.108263in}}%
\pgfpathlineto{\pgfqpoint{3.227150in}{1.124441in}}%
\pgfpathlineto{\pgfqpoint{3.219619in}{1.130949in}}%
\pgfpathlineto{\pgfqpoint{3.216293in}{1.140248in}}%
\pgfpathlineto{\pgfqpoint{3.211683in}{1.141958in}}%
\pgfpathlineto{\pgfqpoint{3.206484in}{1.172807in}}%
\pgfpathlineto{\pgfqpoint{3.229923in}{1.198333in}}%
\pgfpathlineto{\pgfqpoint{3.232624in}{1.197760in}}%
\pgfpathlineto{\pgfqpoint{3.244554in}{1.193790in}}%
\pgfpathlineto{\pgfqpoint{3.246587in}{1.185703in}}%
\pgfpathlineto{\pgfqpoint{3.250804in}{1.179929in}}%
\pgfpathlineto{\pgfqpoint{3.248243in}{1.176631in}}%
\pgfpathlineto{\pgfqpoint{3.248882in}{1.168666in}}%
\pgfpathlineto{\pgfqpoint{3.265343in}{1.162433in}}%
\pgfpathlineto{\pgfqpoint{3.273900in}{1.157290in}}%
\pgfpathlineto{\pgfqpoint{3.287098in}{1.161529in}}%
\pgfpathlineto{\pgfqpoint{3.286885in}{1.169498in}}%
\pgfpathlineto{\pgfqpoint{3.295207in}{1.169584in}}%
\pgfpathlineto{\pgfqpoint{3.297149in}{1.166013in}}%
\pgfpathlineto{\pgfqpoint{3.293259in}{1.151590in}}%
\pgfpathlineto{\pgfqpoint{3.288603in}{1.146209in}}%
\pgfpathlineto{\pgfqpoint{3.292278in}{1.142409in}}%
\pgfpathlineto{\pgfqpoint{3.303080in}{1.146473in}}%
\pgfpathlineto{\pgfqpoint{3.308404in}{1.145005in}}%
\pgfpathlineto{\pgfqpoint{3.312700in}{1.148475in}}%
\pgfpathlineto{\pgfqpoint{3.320379in}{1.148245in}}%
\pgfpathlineto{\pgfqpoint{3.305824in}{1.120574in}}%
\pgfpathlineto{\pgfqpoint{3.300774in}{1.117204in}}%
\pgfpathlineto{\pgfqpoint{3.293113in}{1.118696in}}%
\pgfpathlineto{\pgfqpoint{3.285161in}{1.117485in}}%
\pgfpathlineto{\pgfqpoint{3.263408in}{1.107685in}}%
\pgfpathlineto{\pgfqpoint{3.243003in}{1.088398in}}%
\pgfpathclose%
\pgfusepath{fill}%
\end{pgfscope}%
\begin{pgfscope}%
\pgfpathrectangle{\pgfqpoint{0.100000in}{0.100000in}}{\pgfqpoint{3.608454in}{2.310000in}}%
\pgfusepath{clip}%
\pgfsetbuttcap%
\pgfsetmiterjoin%
\definecolor{currentfill}{rgb}{0.000000,0.509804,0.745098}%
\pgfsetfillcolor{currentfill}%
\pgfsetlinewidth{0.000000pt}%
\definecolor{currentstroke}{rgb}{0.000000,0.000000,0.000000}%
\pgfsetstrokecolor{currentstroke}%
\pgfsetstrokeopacity{0.000000}%
\pgfsetdash{}{0pt}%
\pgfpathmoveto{\pgfqpoint{1.757470in}{0.722100in}}%
\pgfpathlineto{\pgfqpoint{1.755704in}{0.692007in}}%
\pgfpathlineto{\pgfqpoint{1.697817in}{0.695195in}}%
\pgfpathlineto{\pgfqpoint{1.699678in}{0.725186in}}%
\pgfpathlineto{\pgfqpoint{1.678348in}{0.726376in}}%
\pgfpathlineto{\pgfqpoint{1.681284in}{0.762092in}}%
\pgfpathlineto{\pgfqpoint{1.684601in}{0.806456in}}%
\pgfpathlineto{\pgfqpoint{1.690000in}{0.806092in}}%
\pgfpathlineto{\pgfqpoint{1.724900in}{0.803670in}}%
\pgfpathlineto{\pgfqpoint{1.753852in}{0.801689in}}%
\pgfpathlineto{\pgfqpoint{1.751590in}{0.761847in}}%
\pgfpathlineto{\pgfqpoint{1.759948in}{0.761279in}}%
\pgfpathlineto{\pgfqpoint{1.757470in}{0.722100in}}%
\pgfpathclose%
\pgfusepath{fill}%
\end{pgfscope}%
\begin{pgfscope}%
\pgfpathrectangle{\pgfqpoint{0.100000in}{0.100000in}}{\pgfqpoint{3.608454in}{2.310000in}}%
\pgfusepath{clip}%
\pgfsetbuttcap%
\pgfsetmiterjoin%
\definecolor{currentfill}{rgb}{0.000000,0.505882,0.747059}%
\pgfsetfillcolor{currentfill}%
\pgfsetlinewidth{0.000000pt}%
\definecolor{currentstroke}{rgb}{0.000000,0.000000,0.000000}%
\pgfsetstrokecolor{currentstroke}%
\pgfsetstrokeopacity{0.000000}%
\pgfsetdash{}{0pt}%
\pgfpathmoveto{\pgfqpoint{2.839244in}{0.999162in}}%
\pgfpathlineto{\pgfqpoint{2.836248in}{0.996454in}}%
\pgfpathlineto{\pgfqpoint{2.827384in}{1.003120in}}%
\pgfpathlineto{\pgfqpoint{2.830646in}{1.010243in}}%
\pgfpathlineto{\pgfqpoint{2.825899in}{1.020792in}}%
\pgfpathlineto{\pgfqpoint{2.818142in}{1.023149in}}%
\pgfpathlineto{\pgfqpoint{2.810860in}{1.028741in}}%
\pgfpathlineto{\pgfqpoint{2.810594in}{1.033632in}}%
\pgfpathlineto{\pgfqpoint{2.812171in}{1.037511in}}%
\pgfpathlineto{\pgfqpoint{2.820073in}{1.037913in}}%
\pgfpathlineto{\pgfqpoint{2.825426in}{1.046467in}}%
\pgfpathlineto{\pgfqpoint{2.830948in}{1.045651in}}%
\pgfpathlineto{\pgfqpoint{2.833248in}{1.051039in}}%
\pgfpathlineto{\pgfqpoint{2.843160in}{1.055221in}}%
\pgfpathlineto{\pgfqpoint{2.847354in}{1.055558in}}%
\pgfpathlineto{\pgfqpoint{2.850394in}{1.050302in}}%
\pgfpathlineto{\pgfqpoint{2.863588in}{1.049271in}}%
\pgfpathlineto{\pgfqpoint{2.864794in}{1.047108in}}%
\pgfpathlineto{\pgfqpoint{2.863172in}{1.045964in}}%
\pgfpathlineto{\pgfqpoint{2.859003in}{1.029918in}}%
\pgfpathlineto{\pgfqpoint{2.864516in}{1.026478in}}%
\pgfpathlineto{\pgfqpoint{2.864861in}{1.018162in}}%
\pgfpathlineto{\pgfqpoint{2.868390in}{1.010555in}}%
\pgfpathlineto{\pgfqpoint{2.865578in}{1.008074in}}%
\pgfpathlineto{\pgfqpoint{2.869364in}{1.002057in}}%
\pgfpathlineto{\pgfqpoint{2.868081in}{0.996615in}}%
\pgfpathlineto{\pgfqpoint{2.859121in}{0.988728in}}%
\pgfpathlineto{\pgfqpoint{2.856780in}{0.993705in}}%
\pgfpathlineto{\pgfqpoint{2.843269in}{0.994686in}}%
\pgfpathlineto{\pgfqpoint{2.839244in}{0.999162in}}%
\pgfpathclose%
\pgfusepath{fill}%
\end{pgfscope}%
\begin{pgfscope}%
\pgfpathrectangle{\pgfqpoint{0.100000in}{0.100000in}}{\pgfqpoint{3.608454in}{2.310000in}}%
\pgfusepath{clip}%
\pgfsetbuttcap%
\pgfsetmiterjoin%
\definecolor{currentfill}{rgb}{0.000000,0.545098,0.727451}%
\pgfsetfillcolor{currentfill}%
\pgfsetlinewidth{0.000000pt}%
\definecolor{currentstroke}{rgb}{0.000000,0.000000,0.000000}%
\pgfsetstrokecolor{currentstroke}%
\pgfsetstrokeopacity{0.000000}%
\pgfsetdash{}{0pt}%
\pgfpathmoveto{\pgfqpoint{2.075326in}{1.295024in}}%
\pgfpathlineto{\pgfqpoint{2.047953in}{1.295247in}}%
\pgfpathlineto{\pgfqpoint{2.048478in}{1.319317in}}%
\pgfpathlineto{\pgfqpoint{2.023407in}{1.319659in}}%
\pgfpathlineto{\pgfqpoint{2.021173in}{1.326536in}}%
\pgfpathlineto{\pgfqpoint{2.014334in}{1.326679in}}%
\pgfpathlineto{\pgfqpoint{2.014646in}{1.343860in}}%
\pgfpathlineto{\pgfqpoint{2.021423in}{1.343723in}}%
\pgfpathlineto{\pgfqpoint{2.021694in}{1.353951in}}%
\pgfpathlineto{\pgfqpoint{2.030596in}{1.356610in}}%
\pgfpathlineto{\pgfqpoint{2.034708in}{1.355023in}}%
\pgfpathlineto{\pgfqpoint{2.043381in}{1.357123in}}%
\pgfpathlineto{\pgfqpoint{2.043672in}{1.384754in}}%
\pgfpathlineto{\pgfqpoint{2.058697in}{1.384598in}}%
\pgfpathlineto{\pgfqpoint{2.058778in}{1.391470in}}%
\pgfpathlineto{\pgfqpoint{2.072459in}{1.391367in}}%
\pgfpathlineto{\pgfqpoint{2.071958in}{1.372885in}}%
\pgfpathlineto{\pgfqpoint{2.095818in}{1.372775in}}%
\pgfpathlineto{\pgfqpoint{2.095756in}{1.349849in}}%
\pgfpathlineto{\pgfqpoint{2.096838in}{1.338167in}}%
\pgfpathlineto{\pgfqpoint{2.103336in}{1.338237in}}%
\pgfpathlineto{\pgfqpoint{2.103312in}{1.318980in}}%
\pgfpathlineto{\pgfqpoint{2.075903in}{1.319066in}}%
\pgfpathlineto{\pgfqpoint{2.075326in}{1.295024in}}%
\pgfpathclose%
\pgfusepath{fill}%
\end{pgfscope}%
\begin{pgfscope}%
\pgfpathrectangle{\pgfqpoint{0.100000in}{0.100000in}}{\pgfqpoint{3.608454in}{2.310000in}}%
\pgfusepath{clip}%
\pgfsetbuttcap%
\pgfsetmiterjoin%
\definecolor{currentfill}{rgb}{0.000000,0.478431,0.760784}%
\pgfsetfillcolor{currentfill}%
\pgfsetlinewidth{0.000000pt}%
\definecolor{currentstroke}{rgb}{0.000000,0.000000,0.000000}%
\pgfsetstrokecolor{currentstroke}%
\pgfsetstrokeopacity{0.000000}%
\pgfsetdash{}{0pt}%
\pgfpathmoveto{\pgfqpoint{2.024967in}{0.596552in}}%
\pgfpathlineto{\pgfqpoint{2.015636in}{0.587127in}}%
\pgfpathlineto{\pgfqpoint{2.012654in}{0.578657in}}%
\pgfpathlineto{\pgfqpoint{1.992065in}{0.566132in}}%
\pgfpathlineto{\pgfqpoint{1.997763in}{0.573057in}}%
\pgfpathlineto{\pgfqpoint{1.981908in}{0.592727in}}%
\pgfpathlineto{\pgfqpoint{1.976342in}{0.597283in}}%
\pgfpathlineto{\pgfqpoint{1.957611in}{0.597327in}}%
\pgfpathlineto{\pgfqpoint{1.945778in}{0.610158in}}%
\pgfpathlineto{\pgfqpoint{1.966633in}{0.631014in}}%
\pgfpathlineto{\pgfqpoint{1.979102in}{0.635917in}}%
\pgfpathlineto{\pgfqpoint{1.982721in}{0.639405in}}%
\pgfpathlineto{\pgfqpoint{1.990391in}{0.638498in}}%
\pgfpathlineto{\pgfqpoint{2.000273in}{0.618823in}}%
\pgfpathlineto{\pgfqpoint{2.008648in}{0.612091in}}%
\pgfpathlineto{\pgfqpoint{2.014382in}{0.611205in}}%
\pgfpathlineto{\pgfqpoint{2.019290in}{0.599307in}}%
\pgfpathlineto{\pgfqpoint{2.024967in}{0.596552in}}%
\pgfpathclose%
\pgfusepath{fill}%
\end{pgfscope}%
\begin{pgfscope}%
\pgfpathrectangle{\pgfqpoint{0.100000in}{0.100000in}}{\pgfqpoint{3.608454in}{2.310000in}}%
\pgfusepath{clip}%
\pgfsetbuttcap%
\pgfsetmiterjoin%
\definecolor{currentfill}{rgb}{0.000000,0.647059,0.676471}%
\pgfsetfillcolor{currentfill}%
\pgfsetlinewidth{0.000000pt}%
\definecolor{currentstroke}{rgb}{0.000000,0.000000,0.000000}%
\pgfsetstrokecolor{currentstroke}%
\pgfsetstrokeopacity{0.000000}%
\pgfsetdash{}{0pt}%
\pgfpathmoveto{\pgfqpoint{0.755783in}{1.587192in}}%
\pgfpathlineto{\pgfqpoint{0.693366in}{1.603165in}}%
\pgfpathlineto{\pgfqpoint{0.650489in}{1.614771in}}%
\pgfpathlineto{\pgfqpoint{0.661398in}{1.654341in}}%
\pgfpathlineto{\pgfqpoint{0.662884in}{1.653976in}}%
\pgfpathlineto{\pgfqpoint{0.672323in}{1.687633in}}%
\pgfpathlineto{\pgfqpoint{0.671164in}{1.687956in}}%
\pgfpathlineto{\pgfqpoint{0.681032in}{1.722500in}}%
\pgfpathlineto{\pgfqpoint{0.693113in}{1.766738in}}%
\pgfpathlineto{\pgfqpoint{0.728707in}{1.756911in}}%
\pgfpathlineto{\pgfqpoint{0.757203in}{1.749818in}}%
\pgfpathlineto{\pgfqpoint{0.824016in}{1.733109in}}%
\pgfpathlineto{\pgfqpoint{0.824452in}{1.732928in}}%
\pgfpathlineto{\pgfqpoint{0.805480in}{1.656197in}}%
\pgfpathlineto{\pgfqpoint{0.798688in}{1.628778in}}%
\pgfpathlineto{\pgfqpoint{0.785394in}{1.632033in}}%
\pgfpathlineto{\pgfqpoint{0.755783in}{1.587192in}}%
\pgfpathclose%
\pgfusepath{fill}%
\end{pgfscope}%
\begin{pgfscope}%
\pgfpathrectangle{\pgfqpoint{0.100000in}{0.100000in}}{\pgfqpoint{3.608454in}{2.310000in}}%
\pgfusepath{clip}%
\pgfsetbuttcap%
\pgfsetmiterjoin%
\definecolor{currentfill}{rgb}{0.000000,0.329412,0.835294}%
\pgfsetfillcolor{currentfill}%
\pgfsetlinewidth{0.000000pt}%
\definecolor{currentstroke}{rgb}{0.000000,0.000000,0.000000}%
\pgfsetstrokecolor{currentstroke}%
\pgfsetstrokeopacity{0.000000}%
\pgfsetdash{}{0pt}%
\pgfpathmoveto{\pgfqpoint{1.989587in}{1.638444in}}%
\pgfpathlineto{\pgfqpoint{1.988757in}{1.605977in}}%
\pgfpathlineto{\pgfqpoint{1.968280in}{1.606524in}}%
\pgfpathlineto{\pgfqpoint{1.968471in}{1.613417in}}%
\pgfpathlineto{\pgfqpoint{1.914002in}{1.614954in}}%
\pgfpathlineto{\pgfqpoint{1.914725in}{1.650247in}}%
\pgfpathlineto{\pgfqpoint{1.922640in}{1.646360in}}%
\pgfpathlineto{\pgfqpoint{1.925767in}{1.649254in}}%
\pgfpathlineto{\pgfqpoint{1.927556in}{1.657034in}}%
\pgfpathlineto{\pgfqpoint{1.926808in}{1.664949in}}%
\pgfpathlineto{\pgfqpoint{1.928813in}{1.672403in}}%
\pgfpathlineto{\pgfqpoint{1.968146in}{1.671305in}}%
\pgfpathlineto{\pgfqpoint{1.981953in}{1.671006in}}%
\pgfpathlineto{\pgfqpoint{1.981253in}{1.641675in}}%
\pgfpathlineto{\pgfqpoint{1.989587in}{1.638444in}}%
\pgfpathclose%
\pgfusepath{fill}%
\end{pgfscope}%
\begin{pgfscope}%
\pgfpathrectangle{\pgfqpoint{0.100000in}{0.100000in}}{\pgfqpoint{3.608454in}{2.310000in}}%
\pgfusepath{clip}%
\pgfsetbuttcap%
\pgfsetmiterjoin%
\definecolor{currentfill}{rgb}{0.000000,0.545098,0.727451}%
\pgfsetfillcolor{currentfill}%
\pgfsetlinewidth{0.000000pt}%
\definecolor{currentstroke}{rgb}{0.000000,0.000000,0.000000}%
\pgfsetstrokecolor{currentstroke}%
\pgfsetstrokeopacity{0.000000}%
\pgfsetdash{}{0pt}%
\pgfpathmoveto{\pgfqpoint{3.191049in}{1.170900in}}%
\pgfpathlineto{\pgfqpoint{3.206484in}{1.172807in}}%
\pgfpathlineto{\pgfqpoint{3.211683in}{1.141958in}}%
\pgfpathlineto{\pgfqpoint{3.216293in}{1.140248in}}%
\pgfpathlineto{\pgfqpoint{3.219619in}{1.130949in}}%
\pgfpathlineto{\pgfqpoint{3.227150in}{1.124441in}}%
\pgfpathlineto{\pgfqpoint{3.228319in}{1.108263in}}%
\pgfpathlineto{\pgfqpoint{3.200611in}{1.103096in}}%
\pgfpathlineto{\pgfqpoint{3.194034in}{1.089009in}}%
\pgfpathlineto{\pgfqpoint{3.182406in}{1.101945in}}%
\pgfpathlineto{\pgfqpoint{3.174249in}{1.109020in}}%
\pgfpathlineto{\pgfqpoint{3.162834in}{1.117655in}}%
\pgfpathlineto{\pgfqpoint{3.160679in}{1.126465in}}%
\pgfpathlineto{\pgfqpoint{3.160885in}{1.137946in}}%
\pgfpathlineto{\pgfqpoint{3.164716in}{1.144245in}}%
\pgfpathlineto{\pgfqpoint{3.168634in}{1.140787in}}%
\pgfpathlineto{\pgfqpoint{3.173892in}{1.141190in}}%
\pgfpathlineto{\pgfqpoint{3.180850in}{1.147206in}}%
\pgfpathlineto{\pgfqpoint{3.188860in}{1.151551in}}%
\pgfpathlineto{\pgfqpoint{3.187783in}{1.156766in}}%
\pgfpathlineto{\pgfqpoint{3.191049in}{1.170900in}}%
\pgfpathclose%
\pgfusepath{fill}%
\end{pgfscope}%
\begin{pgfscope}%
\pgfpathrectangle{\pgfqpoint{0.100000in}{0.100000in}}{\pgfqpoint{3.608454in}{2.310000in}}%
\pgfusepath{clip}%
\pgfsetbuttcap%
\pgfsetmiterjoin%
\definecolor{currentfill}{rgb}{0.000000,0.576471,0.711765}%
\pgfsetfillcolor{currentfill}%
\pgfsetlinewidth{0.000000pt}%
\definecolor{currentstroke}{rgb}{0.000000,0.000000,0.000000}%
\pgfsetstrokecolor{currentstroke}%
\pgfsetstrokeopacity{0.000000}%
\pgfsetdash{}{0pt}%
\pgfpathmoveto{\pgfqpoint{2.786852in}{0.748618in}}%
\pgfpathlineto{\pgfqpoint{2.789288in}{0.742878in}}%
\pgfpathlineto{\pgfqpoint{2.719091in}{0.734938in}}%
\pgfpathlineto{\pgfqpoint{2.708723in}{0.733977in}}%
\pgfpathlineto{\pgfqpoint{2.704075in}{0.776471in}}%
\pgfpathlineto{\pgfqpoint{2.706700in}{0.783758in}}%
\pgfpathlineto{\pgfqpoint{2.733549in}{0.786465in}}%
\pgfpathlineto{\pgfqpoint{2.755934in}{0.788907in}}%
\pgfpathlineto{\pgfqpoint{2.755229in}{0.795820in}}%
\pgfpathlineto{\pgfqpoint{2.774910in}{0.796997in}}%
\pgfpathlineto{\pgfqpoint{2.780161in}{0.791452in}}%
\pgfpathlineto{\pgfqpoint{2.781911in}{0.785540in}}%
\pgfpathlineto{\pgfqpoint{2.779513in}{0.764362in}}%
\pgfpathlineto{\pgfqpoint{2.781206in}{0.755220in}}%
\pgfpathlineto{\pgfqpoint{2.786852in}{0.748618in}}%
\pgfpathclose%
\pgfusepath{fill}%
\end{pgfscope}%
\begin{pgfscope}%
\pgfpathrectangle{\pgfqpoint{0.100000in}{0.100000in}}{\pgfqpoint{3.608454in}{2.310000in}}%
\pgfusepath{clip}%
\pgfsetbuttcap%
\pgfsetmiterjoin%
\definecolor{currentfill}{rgb}{0.000000,0.654902,0.672549}%
\pgfsetfillcolor{currentfill}%
\pgfsetlinewidth{0.000000pt}%
\definecolor{currentstroke}{rgb}{0.000000,0.000000,0.000000}%
\pgfsetstrokecolor{currentstroke}%
\pgfsetstrokeopacity{0.000000}%
\pgfsetdash{}{0pt}%
\pgfpathmoveto{\pgfqpoint{0.875234in}{2.151233in}}%
\pgfpathlineto{\pgfqpoint{0.821826in}{2.165426in}}%
\pgfpathlineto{\pgfqpoint{0.836363in}{2.217995in}}%
\pgfpathlineto{\pgfqpoint{0.843741in}{2.216808in}}%
\pgfpathlineto{\pgfqpoint{0.857008in}{2.211536in}}%
\pgfpathlineto{\pgfqpoint{0.862221in}{2.202501in}}%
\pgfpathlineto{\pgfqpoint{0.873890in}{2.207982in}}%
\pgfpathlineto{\pgfqpoint{0.878250in}{2.201957in}}%
\pgfpathlineto{\pgfqpoint{0.879920in}{2.194564in}}%
\pgfpathlineto{\pgfqpoint{0.891121in}{2.195618in}}%
\pgfpathlineto{\pgfqpoint{0.894472in}{2.191910in}}%
\pgfpathlineto{\pgfqpoint{0.902377in}{2.193011in}}%
\pgfpathlineto{\pgfqpoint{0.908021in}{2.186968in}}%
\pgfpathlineto{\pgfqpoint{0.913019in}{2.206570in}}%
\pgfpathlineto{\pgfqpoint{0.916424in}{2.212894in}}%
\pgfpathlineto{\pgfqpoint{0.926800in}{2.253370in}}%
\pgfpathlineto{\pgfqpoint{0.920139in}{2.255027in}}%
\pgfpathlineto{\pgfqpoint{0.925222in}{2.260906in}}%
\pgfpathlineto{\pgfqpoint{0.927075in}{2.268237in}}%
\pgfpathlineto{\pgfqpoint{0.937150in}{2.279229in}}%
\pgfpathlineto{\pgfqpoint{0.957430in}{2.274082in}}%
\pgfpathlineto{\pgfqpoint{0.954484in}{2.262189in}}%
\pgfpathlineto{\pgfqpoint{0.967184in}{2.259137in}}%
\pgfpathlineto{\pgfqpoint{0.960593in}{2.232110in}}%
\pgfpathlineto{\pgfqpoint{0.998620in}{2.223097in}}%
\pgfpathlineto{\pgfqpoint{0.993310in}{2.200722in}}%
\pgfpathlineto{\pgfqpoint{0.988985in}{2.182170in}}%
\pgfpathlineto{\pgfqpoint{0.996594in}{2.161889in}}%
\pgfpathlineto{\pgfqpoint{1.000904in}{2.156205in}}%
\pgfpathlineto{\pgfqpoint{1.000929in}{2.147949in}}%
\pgfpathlineto{\pgfqpoint{0.996502in}{2.145181in}}%
\pgfpathlineto{\pgfqpoint{1.001953in}{2.138275in}}%
\pgfpathlineto{\pgfqpoint{0.994972in}{2.136779in}}%
\pgfpathlineto{\pgfqpoint{1.003713in}{2.127225in}}%
\pgfpathlineto{\pgfqpoint{1.003956in}{2.123021in}}%
\pgfpathlineto{\pgfqpoint{1.014258in}{2.117180in}}%
\pgfpathlineto{\pgfqpoint{1.023030in}{2.093793in}}%
\pgfpathlineto{\pgfqpoint{1.027701in}{2.087646in}}%
\pgfpathlineto{\pgfqpoint{0.993104in}{2.095521in}}%
\pgfpathlineto{\pgfqpoint{0.954967in}{2.104591in}}%
\pgfpathlineto{\pgfqpoint{0.956600in}{2.111418in}}%
\pgfpathlineto{\pgfqpoint{0.950036in}{2.114158in}}%
\pgfpathlineto{\pgfqpoint{0.933521in}{2.118231in}}%
\pgfpathlineto{\pgfqpoint{0.921056in}{2.128678in}}%
\pgfpathlineto{\pgfqpoint{0.923573in}{2.138917in}}%
\pgfpathlineto{\pgfqpoint{0.875234in}{2.151233in}}%
\pgfpathclose%
\pgfusepath{fill}%
\end{pgfscope}%
\begin{pgfscope}%
\pgfpathrectangle{\pgfqpoint{0.100000in}{0.100000in}}{\pgfqpoint{3.608454in}{2.310000in}}%
\pgfusepath{clip}%
\pgfsetbuttcap%
\pgfsetmiterjoin%
\definecolor{currentfill}{rgb}{0.000000,0.682353,0.658824}%
\pgfsetfillcolor{currentfill}%
\pgfsetlinewidth{0.000000pt}%
\definecolor{currentstroke}{rgb}{0.000000,0.000000,0.000000}%
\pgfsetstrokecolor{currentstroke}%
\pgfsetstrokeopacity{0.000000}%
\pgfsetdash{}{0pt}%
\pgfpathmoveto{\pgfqpoint{2.187595in}{0.601129in}}%
\pgfpathlineto{\pgfqpoint{2.169818in}{0.598417in}}%
\pgfpathlineto{\pgfqpoint{2.151781in}{0.590398in}}%
\pgfpathlineto{\pgfqpoint{2.151337in}{0.616492in}}%
\pgfpathlineto{\pgfqpoint{2.145490in}{0.616619in}}%
\pgfpathlineto{\pgfqpoint{2.145181in}{0.634496in}}%
\pgfpathlineto{\pgfqpoint{2.134719in}{0.634369in}}%
\pgfpathlineto{\pgfqpoint{2.125194in}{0.664549in}}%
\pgfpathlineto{\pgfqpoint{2.109149in}{0.671435in}}%
\pgfpathlineto{\pgfqpoint{2.105170in}{0.678118in}}%
\pgfpathlineto{\pgfqpoint{2.098891in}{0.678184in}}%
\pgfpathlineto{\pgfqpoint{2.097748in}{0.687767in}}%
\pgfpathlineto{\pgfqpoint{2.093068in}{0.691255in}}%
\pgfpathlineto{\pgfqpoint{2.110954in}{0.708601in}}%
\pgfpathlineto{\pgfqpoint{2.117525in}{0.716961in}}%
\pgfpathlineto{\pgfqpoint{2.123775in}{0.713712in}}%
\pgfpathlineto{\pgfqpoint{2.132588in}{0.712274in}}%
\pgfpathlineto{\pgfqpoint{2.143868in}{0.708025in}}%
\pgfpathlineto{\pgfqpoint{2.166253in}{0.713450in}}%
\pgfpathlineto{\pgfqpoint{2.181123in}{0.718301in}}%
\pgfpathlineto{\pgfqpoint{2.202317in}{0.720035in}}%
\pgfpathlineto{\pgfqpoint{2.205558in}{0.720837in}}%
\pgfpathlineto{\pgfqpoint{2.204984in}{0.705199in}}%
\pgfpathlineto{\pgfqpoint{2.207696in}{0.700477in}}%
\pgfpathlineto{\pgfqpoint{2.204721in}{0.696052in}}%
\pgfpathlineto{\pgfqpoint{2.200790in}{0.680240in}}%
\pgfpathlineto{\pgfqpoint{2.197286in}{0.677217in}}%
\pgfpathlineto{\pgfqpoint{2.193927in}{0.668423in}}%
\pgfpathlineto{\pgfqpoint{2.196505in}{0.661211in}}%
\pgfpathlineto{\pgfqpoint{2.191936in}{0.652528in}}%
\pgfpathlineto{\pgfqpoint{2.196201in}{0.649140in}}%
\pgfpathlineto{\pgfqpoint{2.196138in}{0.630275in}}%
\pgfpathlineto{\pgfqpoint{2.190675in}{0.624923in}}%
\pgfpathlineto{\pgfqpoint{2.188702in}{0.619056in}}%
\pgfpathlineto{\pgfqpoint{2.181395in}{0.610955in}}%
\pgfpathlineto{\pgfqpoint{2.187595in}{0.601129in}}%
\pgfpathclose%
\pgfusepath{fill}%
\end{pgfscope}%
\begin{pgfscope}%
\pgfpathrectangle{\pgfqpoint{0.100000in}{0.100000in}}{\pgfqpoint{3.608454in}{2.310000in}}%
\pgfusepath{clip}%
\pgfsetbuttcap%
\pgfsetmiterjoin%
\definecolor{currentfill}{rgb}{0.000000,0.301961,0.849020}%
\pgfsetfillcolor{currentfill}%
\pgfsetlinewidth{0.000000pt}%
\definecolor{currentstroke}{rgb}{0.000000,0.000000,0.000000}%
\pgfsetstrokecolor{currentstroke}%
\pgfsetstrokeopacity{0.000000}%
\pgfsetdash{}{0pt}%
\pgfpathmoveto{\pgfqpoint{3.217018in}{1.367687in}}%
\pgfpathlineto{\pgfqpoint{3.210448in}{1.365277in}}%
\pgfpathlineto{\pgfqpoint{3.202566in}{1.350706in}}%
\pgfpathlineto{\pgfqpoint{3.197288in}{1.352635in}}%
\pgfpathlineto{\pgfqpoint{3.195177in}{1.357724in}}%
\pgfpathlineto{\pgfqpoint{3.178627in}{1.363629in}}%
\pgfpathlineto{\pgfqpoint{3.168135in}{1.369114in}}%
\pgfpathlineto{\pgfqpoint{3.161085in}{1.369292in}}%
\pgfpathlineto{\pgfqpoint{3.175771in}{1.395430in}}%
\pgfpathlineto{\pgfqpoint{3.175115in}{1.399614in}}%
\pgfpathlineto{\pgfqpoint{3.182634in}{1.411460in}}%
\pgfpathlineto{\pgfqpoint{3.194258in}{1.405953in}}%
\pgfpathlineto{\pgfqpoint{3.196379in}{1.396687in}}%
\pgfpathlineto{\pgfqpoint{3.199440in}{1.393885in}}%
\pgfpathlineto{\pgfqpoint{3.209863in}{1.398797in}}%
\pgfpathlineto{\pgfqpoint{3.214679in}{1.396873in}}%
\pgfpathlineto{\pgfqpoint{3.216571in}{1.392855in}}%
\pgfpathlineto{\pgfqpoint{3.214965in}{1.384823in}}%
\pgfpathlineto{\pgfqpoint{3.209188in}{1.375295in}}%
\pgfpathlineto{\pgfqpoint{3.217018in}{1.367687in}}%
\pgfpathclose%
\pgfusepath{fill}%
\end{pgfscope}%
\begin{pgfscope}%
\pgfpathrectangle{\pgfqpoint{0.100000in}{0.100000in}}{\pgfqpoint{3.608454in}{2.310000in}}%
\pgfusepath{clip}%
\pgfsetbuttcap%
\pgfsetmiterjoin%
\definecolor{currentfill}{rgb}{0.000000,0.333333,0.833333}%
\pgfsetfillcolor{currentfill}%
\pgfsetlinewidth{0.000000pt}%
\definecolor{currentstroke}{rgb}{0.000000,0.000000,0.000000}%
\pgfsetstrokecolor{currentstroke}%
\pgfsetstrokeopacity{0.000000}%
\pgfsetdash{}{0pt}%
\pgfpathmoveto{\pgfqpoint{1.821457in}{1.218805in}}%
\pgfpathlineto{\pgfqpoint{1.820080in}{1.188664in}}%
\pgfpathlineto{\pgfqpoint{1.790983in}{1.190359in}}%
\pgfpathlineto{\pgfqpoint{1.785504in}{1.190712in}}%
\pgfpathlineto{\pgfqpoint{1.787082in}{1.221084in}}%
\pgfpathlineto{\pgfqpoint{1.786488in}{1.228084in}}%
\pgfpathlineto{\pgfqpoint{1.752303in}{1.230115in}}%
\pgfpathlineto{\pgfqpoint{1.754126in}{1.258033in}}%
\pgfpathlineto{\pgfqpoint{1.754193in}{1.271800in}}%
\pgfpathlineto{\pgfqpoint{1.781396in}{1.270043in}}%
\pgfpathlineto{\pgfqpoint{1.782549in}{1.290657in}}%
\pgfpathlineto{\pgfqpoint{1.822395in}{1.288502in}}%
\pgfpathlineto{\pgfqpoint{1.823241in}{1.288458in}}%
\pgfpathlineto{\pgfqpoint{1.821521in}{1.253978in}}%
\pgfpathlineto{\pgfqpoint{1.821831in}{1.246792in}}%
\pgfpathlineto{\pgfqpoint{1.821457in}{1.218805in}}%
\pgfpathclose%
\pgfusepath{fill}%
\end{pgfscope}%
\begin{pgfscope}%
\pgfpathrectangle{\pgfqpoint{0.100000in}{0.100000in}}{\pgfqpoint{3.608454in}{2.310000in}}%
\pgfusepath{clip}%
\pgfsetbuttcap%
\pgfsetmiterjoin%
\definecolor{currentfill}{rgb}{0.000000,0.674510,0.662745}%
\pgfsetfillcolor{currentfill}%
\pgfsetlinewidth{0.000000pt}%
\definecolor{currentstroke}{rgb}{0.000000,0.000000,0.000000}%
\pgfsetstrokecolor{currentstroke}%
\pgfsetstrokeopacity{0.000000}%
\pgfsetdash{}{0pt}%
\pgfpathmoveto{\pgfqpoint{1.759459in}{0.801389in}}%
\pgfpathlineto{\pgfqpoint{1.753852in}{0.801689in}}%
\pgfpathlineto{\pgfqpoint{1.724900in}{0.803670in}}%
\pgfpathlineto{\pgfqpoint{1.729883in}{0.873343in}}%
\pgfpathlineto{\pgfqpoint{1.738977in}{0.872740in}}%
\pgfpathlineto{\pgfqpoint{1.763926in}{0.871040in}}%
\pgfpathlineto{\pgfqpoint{1.759459in}{0.801389in}}%
\pgfpathclose%
\pgfusepath{fill}%
\end{pgfscope}%
\begin{pgfscope}%
\pgfpathrectangle{\pgfqpoint{0.100000in}{0.100000in}}{\pgfqpoint{3.608454in}{2.310000in}}%
\pgfusepath{clip}%
\pgfsetbuttcap%
\pgfsetmiterjoin%
\definecolor{currentfill}{rgb}{0.000000,0.333333,0.833333}%
\pgfsetfillcolor{currentfill}%
\pgfsetlinewidth{0.000000pt}%
\definecolor{currentstroke}{rgb}{0.000000,0.000000,0.000000}%
\pgfsetstrokecolor{currentstroke}%
\pgfsetstrokeopacity{0.000000}%
\pgfsetdash{}{0pt}%
\pgfpathmoveto{\pgfqpoint{1.752303in}{1.230115in}}%
\pgfpathlineto{\pgfqpoint{1.751890in}{1.223219in}}%
\pgfpathlineto{\pgfqpoint{1.725856in}{1.224921in}}%
\pgfpathlineto{\pgfqpoint{1.724493in}{1.225005in}}%
\pgfpathlineto{\pgfqpoint{1.726337in}{1.252520in}}%
\pgfpathlineto{\pgfqpoint{1.666954in}{1.256981in}}%
\pgfpathlineto{\pgfqpoint{1.670021in}{1.298287in}}%
\pgfpathlineto{\pgfqpoint{1.699539in}{1.296124in}}%
\pgfpathlineto{\pgfqpoint{1.754213in}{1.292521in}}%
\pgfpathlineto{\pgfqpoint{1.782549in}{1.290657in}}%
\pgfpathlineto{\pgfqpoint{1.781396in}{1.270043in}}%
\pgfpathlineto{\pgfqpoint{1.754193in}{1.271800in}}%
\pgfpathlineto{\pgfqpoint{1.754126in}{1.258033in}}%
\pgfpathlineto{\pgfqpoint{1.752303in}{1.230115in}}%
\pgfpathclose%
\pgfusepath{fill}%
\end{pgfscope}%
\begin{pgfscope}%
\pgfpathrectangle{\pgfqpoint{0.100000in}{0.100000in}}{\pgfqpoint{3.608454in}{2.310000in}}%
\pgfusepath{clip}%
\pgfsetbuttcap%
\pgfsetmiterjoin%
\definecolor{currentfill}{rgb}{0.000000,0.521569,0.739216}%
\pgfsetfillcolor{currentfill}%
\pgfsetlinewidth{0.000000pt}%
\definecolor{currentstroke}{rgb}{0.000000,0.000000,0.000000}%
\pgfsetstrokecolor{currentstroke}%
\pgfsetstrokeopacity{0.000000}%
\pgfsetdash{}{0pt}%
\pgfpathmoveto{\pgfqpoint{2.280744in}{1.286387in}}%
\pgfpathlineto{\pgfqpoint{2.314915in}{1.287249in}}%
\pgfpathlineto{\pgfqpoint{2.315298in}{1.276599in}}%
\pgfpathlineto{\pgfqpoint{2.322748in}{1.276780in}}%
\pgfpathlineto{\pgfqpoint{2.323611in}{1.248060in}}%
\pgfpathlineto{\pgfqpoint{2.337381in}{1.248460in}}%
\pgfpathlineto{\pgfqpoint{2.337514in}{1.241571in}}%
\pgfpathlineto{\pgfqpoint{2.347355in}{1.241765in}}%
\pgfpathlineto{\pgfqpoint{2.347612in}{1.233131in}}%
\pgfpathlineto{\pgfqpoint{2.358074in}{1.234831in}}%
\pgfpathlineto{\pgfqpoint{2.371740in}{1.235384in}}%
\pgfpathlineto{\pgfqpoint{2.372317in}{1.217256in}}%
\pgfpathlineto{\pgfqpoint{2.375052in}{1.209261in}}%
\pgfpathlineto{\pgfqpoint{2.375334in}{1.200990in}}%
\pgfpathlineto{\pgfqpoint{2.373090in}{1.191695in}}%
\pgfpathlineto{\pgfqpoint{2.361419in}{1.191422in}}%
\pgfpathlineto{\pgfqpoint{2.357912in}{1.194600in}}%
\pgfpathlineto{\pgfqpoint{2.357801in}{1.198383in}}%
\pgfpathlineto{\pgfqpoint{2.350740in}{1.202506in}}%
\pgfpathlineto{\pgfqpoint{2.347208in}{1.211310in}}%
\pgfpathlineto{\pgfqpoint{2.346916in}{1.219387in}}%
\pgfpathlineto{\pgfqpoint{2.310515in}{1.218707in}}%
\pgfpathlineto{\pgfqpoint{2.310037in}{1.232515in}}%
\pgfpathlineto{\pgfqpoint{2.279053in}{1.231943in}}%
\pgfpathlineto{\pgfqpoint{2.278985in}{1.235414in}}%
\pgfpathlineto{\pgfqpoint{2.268526in}{1.240287in}}%
\pgfpathlineto{\pgfqpoint{2.268476in}{1.251773in}}%
\pgfpathlineto{\pgfqpoint{2.268296in}{1.264580in}}%
\pgfpathlineto{\pgfqpoint{2.282058in}{1.264607in}}%
\pgfpathlineto{\pgfqpoint{2.280744in}{1.286387in}}%
\pgfpathclose%
\pgfusepath{fill}%
\end{pgfscope}%
\begin{pgfscope}%
\pgfpathrectangle{\pgfqpoint{0.100000in}{0.100000in}}{\pgfqpoint{3.608454in}{2.310000in}}%
\pgfusepath{clip}%
\pgfsetbuttcap%
\pgfsetmiterjoin%
\definecolor{currentfill}{rgb}{0.000000,0.682353,0.658824}%
\pgfsetfillcolor{currentfill}%
\pgfsetlinewidth{0.000000pt}%
\definecolor{currentstroke}{rgb}{0.000000,0.000000,0.000000}%
\pgfsetstrokecolor{currentstroke}%
\pgfsetstrokeopacity{0.000000}%
\pgfsetdash{}{0pt}%
\pgfpathmoveto{\pgfqpoint{2.831112in}{0.753917in}}%
\pgfpathlineto{\pgfqpoint{2.834431in}{0.723057in}}%
\pgfpathlineto{\pgfqpoint{2.801217in}{0.720980in}}%
\pgfpathlineto{\pgfqpoint{2.795289in}{0.729104in}}%
\pgfpathlineto{\pgfqpoint{2.794346in}{0.734429in}}%
\pgfpathlineto{\pgfqpoint{2.789288in}{0.742878in}}%
\pgfpathlineto{\pgfqpoint{2.786852in}{0.748618in}}%
\pgfpathlineto{\pgfqpoint{2.794089in}{0.749184in}}%
\pgfpathlineto{\pgfqpoint{2.792796in}{0.763762in}}%
\pgfpathlineto{\pgfqpoint{2.811462in}{0.766109in}}%
\pgfpathlineto{\pgfqpoint{2.818573in}{0.766684in}}%
\pgfpathlineto{\pgfqpoint{2.819817in}{0.752621in}}%
\pgfpathlineto{\pgfqpoint{2.831112in}{0.753917in}}%
\pgfpathclose%
\pgfusepath{fill}%
\end{pgfscope}%
\begin{pgfscope}%
\pgfpathrectangle{\pgfqpoint{0.100000in}{0.100000in}}{\pgfqpoint{3.608454in}{2.310000in}}%
\pgfusepath{clip}%
\pgfsetbuttcap%
\pgfsetmiterjoin%
\definecolor{currentfill}{rgb}{0.000000,0.431373,0.784314}%
\pgfsetfillcolor{currentfill}%
\pgfsetlinewidth{0.000000pt}%
\definecolor{currentstroke}{rgb}{0.000000,0.000000,0.000000}%
\pgfsetstrokecolor{currentstroke}%
\pgfsetstrokeopacity{0.000000}%
\pgfsetdash{}{0pt}%
\pgfpathmoveto{\pgfqpoint{2.024967in}{0.596552in}}%
\pgfpathlineto{\pgfqpoint{2.019290in}{0.599307in}}%
\pgfpathlineto{\pgfqpoint{2.014382in}{0.611205in}}%
\pgfpathlineto{\pgfqpoint{2.008648in}{0.612091in}}%
\pgfpathlineto{\pgfqpoint{2.000273in}{0.618823in}}%
\pgfpathlineto{\pgfqpoint{1.990391in}{0.638498in}}%
\pgfpathlineto{\pgfqpoint{1.982721in}{0.639405in}}%
\pgfpathlineto{\pgfqpoint{1.987116in}{0.645931in}}%
\pgfpathlineto{\pgfqpoint{1.993561in}{0.650091in}}%
\pgfpathlineto{\pgfqpoint{2.008653in}{0.653539in}}%
\pgfpathlineto{\pgfqpoint{2.016277in}{0.656724in}}%
\pgfpathlineto{\pgfqpoint{2.025563in}{0.657506in}}%
\pgfpathlineto{\pgfqpoint{2.031314in}{0.647827in}}%
\pgfpathlineto{\pgfqpoint{2.031567in}{0.644083in}}%
\pgfpathlineto{\pgfqpoint{2.051242in}{0.645234in}}%
\pgfpathlineto{\pgfqpoint{2.049731in}{0.675985in}}%
\pgfpathlineto{\pgfqpoint{2.065563in}{0.666170in}}%
\pgfpathlineto{\pgfqpoint{2.081886in}{0.666080in}}%
\pgfpathlineto{\pgfqpoint{2.084351in}{0.694073in}}%
\pgfpathlineto{\pgfqpoint{2.089584in}{0.697805in}}%
\pgfpathlineto{\pgfqpoint{2.093068in}{0.691255in}}%
\pgfpathlineto{\pgfqpoint{2.097748in}{0.687767in}}%
\pgfpathlineto{\pgfqpoint{2.098891in}{0.678184in}}%
\pgfpathlineto{\pgfqpoint{2.105170in}{0.678118in}}%
\pgfpathlineto{\pgfqpoint{2.109149in}{0.671435in}}%
\pgfpathlineto{\pgfqpoint{2.125194in}{0.664549in}}%
\pgfpathlineto{\pgfqpoint{2.134719in}{0.634369in}}%
\pgfpathlineto{\pgfqpoint{2.145181in}{0.634496in}}%
\pgfpathlineto{\pgfqpoint{2.145490in}{0.616619in}}%
\pgfpathlineto{\pgfqpoint{2.151337in}{0.616492in}}%
\pgfpathlineto{\pgfqpoint{2.151781in}{0.590398in}}%
\pgfpathlineto{\pgfqpoint{2.133973in}{0.582254in}}%
\pgfpathlineto{\pgfqpoint{2.137484in}{0.591301in}}%
\pgfpathlineto{\pgfqpoint{2.124852in}{0.587349in}}%
\pgfpathlineto{\pgfqpoint{2.128361in}{0.601241in}}%
\pgfpathlineto{\pgfqpoint{2.125042in}{0.608247in}}%
\pgfpathlineto{\pgfqpoint{2.119618in}{0.605904in}}%
\pgfpathlineto{\pgfqpoint{2.115649in}{0.599011in}}%
\pgfpathlineto{\pgfqpoint{2.109756in}{0.601205in}}%
\pgfpathlineto{\pgfqpoint{2.105684in}{0.596360in}}%
\pgfpathlineto{\pgfqpoint{2.105660in}{0.588738in}}%
\pgfpathlineto{\pgfqpoint{2.113115in}{0.585047in}}%
\pgfpathlineto{\pgfqpoint{2.114156in}{0.572287in}}%
\pgfpathlineto{\pgfqpoint{2.125821in}{0.572042in}}%
\pgfpathlineto{\pgfqpoint{2.098432in}{0.551471in}}%
\pgfpathlineto{\pgfqpoint{2.102863in}{0.561208in}}%
\pgfpathlineto{\pgfqpoint{2.090532in}{0.582574in}}%
\pgfpathlineto{\pgfqpoint{2.091703in}{0.589805in}}%
\pgfpathlineto{\pgfqpoint{2.088573in}{0.592795in}}%
\pgfpathlineto{\pgfqpoint{2.077251in}{0.591805in}}%
\pgfpathlineto{\pgfqpoint{2.074520in}{0.580523in}}%
\pgfpathlineto{\pgfqpoint{2.068450in}{0.580491in}}%
\pgfpathlineto{\pgfqpoint{2.066381in}{0.572643in}}%
\pgfpathlineto{\pgfqpoint{2.061080in}{0.568664in}}%
\pgfpathlineto{\pgfqpoint{2.054185in}{0.571463in}}%
\pgfpathlineto{\pgfqpoint{2.047518in}{0.566563in}}%
\pgfpathlineto{\pgfqpoint{2.040450in}{0.572451in}}%
\pgfpathlineto{\pgfqpoint{2.037581in}{0.580386in}}%
\pgfpathlineto{\pgfqpoint{2.032880in}{0.583631in}}%
\pgfpathlineto{\pgfqpoint{2.032870in}{0.593014in}}%
\pgfpathlineto{\pgfqpoint{2.024967in}{0.596552in}}%
\pgfpathclose%
\pgfusepath{fill}%
\end{pgfscope}%
\begin{pgfscope}%
\pgfpathrectangle{\pgfqpoint{0.100000in}{0.100000in}}{\pgfqpoint{3.608454in}{2.310000in}}%
\pgfusepath{clip}%
\pgfsetbuttcap%
\pgfsetmiterjoin%
\definecolor{currentfill}{rgb}{0.000000,0.439216,0.780392}%
\pgfsetfillcolor{currentfill}%
\pgfsetlinewidth{0.000000pt}%
\definecolor{currentstroke}{rgb}{0.000000,0.000000,0.000000}%
\pgfsetstrokecolor{currentstroke}%
\pgfsetstrokeopacity{0.000000}%
\pgfsetdash{}{0pt}%
\pgfpathmoveto{\pgfqpoint{1.938143in}{0.938336in}}%
\pgfpathlineto{\pgfqpoint{1.931087in}{0.945835in}}%
\pgfpathlineto{\pgfqpoint{1.917943in}{0.934949in}}%
\pgfpathlineto{\pgfqpoint{1.910664in}{0.938402in}}%
\pgfpathlineto{\pgfqpoint{1.913020in}{0.946294in}}%
\pgfpathlineto{\pgfqpoint{1.906317in}{0.946847in}}%
\pgfpathlineto{\pgfqpoint{1.898581in}{0.956614in}}%
\pgfpathlineto{\pgfqpoint{1.885921in}{0.960469in}}%
\pgfpathlineto{\pgfqpoint{1.883512in}{0.956032in}}%
\pgfpathlineto{\pgfqpoint{1.877716in}{0.953284in}}%
\pgfpathlineto{\pgfqpoint{1.869882in}{0.961093in}}%
\pgfpathlineto{\pgfqpoint{1.870481in}{0.975087in}}%
\pgfpathlineto{\pgfqpoint{1.867064in}{0.975207in}}%
\pgfpathlineto{\pgfqpoint{1.867617in}{0.988984in}}%
\pgfpathlineto{\pgfqpoint{1.856890in}{0.989419in}}%
\pgfpathlineto{\pgfqpoint{1.857172in}{0.996308in}}%
\pgfpathlineto{\pgfqpoint{1.858037in}{1.016994in}}%
\pgfpathlineto{\pgfqpoint{1.871452in}{1.016408in}}%
\pgfpathlineto{\pgfqpoint{1.905912in}{1.015133in}}%
\pgfpathlineto{\pgfqpoint{1.905475in}{1.001362in}}%
\pgfpathlineto{\pgfqpoint{1.932881in}{1.000512in}}%
\pgfpathlineto{\pgfqpoint{1.939766in}{1.000299in}}%
\pgfpathlineto{\pgfqpoint{1.938143in}{0.938336in}}%
\pgfpathclose%
\pgfusepath{fill}%
\end{pgfscope}%
\begin{pgfscope}%
\pgfpathrectangle{\pgfqpoint{0.100000in}{0.100000in}}{\pgfqpoint{3.608454in}{2.310000in}}%
\pgfusepath{clip}%
\pgfsetbuttcap%
\pgfsetmiterjoin%
\definecolor{currentfill}{rgb}{0.000000,0.498039,0.750980}%
\pgfsetfillcolor{currentfill}%
\pgfsetlinewidth{0.000000pt}%
\definecolor{currentstroke}{rgb}{0.000000,0.000000,0.000000}%
\pgfsetstrokecolor{currentstroke}%
\pgfsetstrokeopacity{0.000000}%
\pgfsetdash{}{0pt}%
\pgfpathmoveto{\pgfqpoint{3.358769in}{1.729605in}}%
\pgfpathlineto{\pgfqpoint{3.358373in}{1.723673in}}%
\pgfpathlineto{\pgfqpoint{3.347348in}{1.723269in}}%
\pgfpathlineto{\pgfqpoint{3.335015in}{1.726252in}}%
\pgfpathlineto{\pgfqpoint{3.333233in}{1.729959in}}%
\pgfpathlineto{\pgfqpoint{3.325396in}{1.731486in}}%
\pgfpathlineto{\pgfqpoint{3.324878in}{1.725364in}}%
\pgfpathlineto{\pgfqpoint{3.311248in}{1.723631in}}%
\pgfpathlineto{\pgfqpoint{3.302028in}{1.725960in}}%
\pgfpathlineto{\pgfqpoint{3.296564in}{1.727321in}}%
\pgfpathlineto{\pgfqpoint{3.300176in}{1.740144in}}%
\pgfpathlineto{\pgfqpoint{3.287975in}{1.743465in}}%
\pgfpathlineto{\pgfqpoint{3.281055in}{1.749476in}}%
\pgfpathlineto{\pgfqpoint{3.283745in}{1.758930in}}%
\pgfpathlineto{\pgfqpoint{3.279470in}{1.768032in}}%
\pgfpathlineto{\pgfqpoint{3.279173in}{1.774369in}}%
\pgfpathlineto{\pgfqpoint{3.291128in}{1.772508in}}%
\pgfpathlineto{\pgfqpoint{3.300308in}{1.775508in}}%
\pgfpathlineto{\pgfqpoint{3.309890in}{1.786988in}}%
\pgfpathlineto{\pgfqpoint{3.287566in}{1.849431in}}%
\pgfpathlineto{\pgfqpoint{3.296075in}{1.852707in}}%
\pgfpathlineto{\pgfqpoint{3.332649in}{1.866641in}}%
\pgfpathlineto{\pgfqpoint{3.333464in}{1.857292in}}%
\pgfpathlineto{\pgfqpoint{3.344700in}{1.848799in}}%
\pgfpathlineto{\pgfqpoint{3.348859in}{1.830296in}}%
\pgfpathlineto{\pgfqpoint{3.356483in}{1.793242in}}%
\pgfpathlineto{\pgfqpoint{3.358955in}{1.786375in}}%
\pgfpathlineto{\pgfqpoint{3.357447in}{1.732265in}}%
\pgfpathlineto{\pgfqpoint{3.358769in}{1.729605in}}%
\pgfpathclose%
\pgfusepath{fill}%
\end{pgfscope}%
\begin{pgfscope}%
\pgfpathrectangle{\pgfqpoint{0.100000in}{0.100000in}}{\pgfqpoint{3.608454in}{2.310000in}}%
\pgfusepath{clip}%
\pgfsetbuttcap%
\pgfsetmiterjoin%
\definecolor{currentfill}{rgb}{0.000000,0.466667,0.766667}%
\pgfsetfillcolor{currentfill}%
\pgfsetlinewidth{0.000000pt}%
\definecolor{currentstroke}{rgb}{0.000000,0.000000,0.000000}%
\pgfsetstrokecolor{currentstroke}%
\pgfsetstrokeopacity{0.000000}%
\pgfsetdash{}{0pt}%
\pgfpathmoveto{\pgfqpoint{1.364865in}{1.868873in}}%
\pgfpathlineto{\pgfqpoint{1.318547in}{1.875904in}}%
\pgfpathlineto{\pgfqpoint{1.317490in}{1.876575in}}%
\pgfpathlineto{\pgfqpoint{1.278870in}{1.882470in}}%
\pgfpathlineto{\pgfqpoint{1.281086in}{1.895779in}}%
\pgfpathlineto{\pgfqpoint{1.266400in}{1.898228in}}%
\pgfpathlineto{\pgfqpoint{1.268744in}{1.912078in}}%
\pgfpathlineto{\pgfqpoint{1.275970in}{1.910859in}}%
\pgfpathlineto{\pgfqpoint{1.278245in}{1.924515in}}%
\pgfpathlineto{\pgfqpoint{1.292921in}{1.929061in}}%
\pgfpathlineto{\pgfqpoint{1.299700in}{1.927978in}}%
\pgfpathlineto{\pgfqpoint{1.301786in}{1.941660in}}%
\pgfpathlineto{\pgfqpoint{1.307165in}{1.954995in}}%
\pgfpathlineto{\pgfqpoint{1.311739in}{1.954347in}}%
\pgfpathlineto{\pgfqpoint{1.313126in}{1.960922in}}%
\pgfpathlineto{\pgfqpoint{1.302894in}{1.962642in}}%
\pgfpathlineto{\pgfqpoint{1.301322in}{1.969905in}}%
\pgfpathlineto{\pgfqpoint{1.302428in}{1.976762in}}%
\pgfpathlineto{\pgfqpoint{1.316114in}{1.974564in}}%
\pgfpathlineto{\pgfqpoint{1.322819in}{2.011773in}}%
\pgfpathlineto{\pgfqpoint{1.323552in}{2.016288in}}%
\pgfpathlineto{\pgfqpoint{1.364475in}{2.009800in}}%
\pgfpathlineto{\pgfqpoint{1.407997in}{2.003940in}}%
\pgfpathlineto{\pgfqpoint{1.409028in}{1.997695in}}%
\pgfpathlineto{\pgfqpoint{1.406757in}{1.983167in}}%
\pgfpathlineto{\pgfqpoint{1.410750in}{1.980307in}}%
\pgfpathlineto{\pgfqpoint{1.414681in}{1.966931in}}%
\pgfpathlineto{\pgfqpoint{1.421242in}{1.955525in}}%
\pgfpathlineto{\pgfqpoint{1.417308in}{1.945271in}}%
\pgfpathlineto{\pgfqpoint{1.408401in}{1.946543in}}%
\pgfpathlineto{\pgfqpoint{1.407750in}{1.941975in}}%
\pgfpathlineto{\pgfqpoint{1.400977in}{1.942930in}}%
\pgfpathlineto{\pgfqpoint{1.395427in}{1.936746in}}%
\pgfpathlineto{\pgfqpoint{1.386336in}{1.938080in}}%
\pgfpathlineto{\pgfqpoint{1.383830in}{1.929135in}}%
\pgfpathlineto{\pgfqpoint{1.380767in}{1.908242in}}%
\pgfpathlineto{\pgfqpoint{1.366491in}{1.905688in}}%
\pgfpathlineto{\pgfqpoint{1.348402in}{1.908679in}}%
\pgfpathlineto{\pgfqpoint{1.345983in}{1.906551in}}%
\pgfpathlineto{\pgfqpoint{1.342451in}{1.889782in}}%
\pgfpathlineto{\pgfqpoint{1.375793in}{1.885029in}}%
\pgfpathlineto{\pgfqpoint{1.372850in}{1.883266in}}%
\pgfpathlineto{\pgfqpoint{1.364865in}{1.868873in}}%
\pgfpathclose%
\pgfusepath{fill}%
\end{pgfscope}%
\begin{pgfscope}%
\pgfpathrectangle{\pgfqpoint{0.100000in}{0.100000in}}{\pgfqpoint{3.608454in}{2.310000in}}%
\pgfusepath{clip}%
\pgfsetbuttcap%
\pgfsetmiterjoin%
\definecolor{currentfill}{rgb}{0.000000,0.647059,0.676471}%
\pgfsetfillcolor{currentfill}%
\pgfsetlinewidth{0.000000pt}%
\definecolor{currentstroke}{rgb}{0.000000,0.000000,0.000000}%
\pgfsetstrokecolor{currentstroke}%
\pgfsetstrokeopacity{0.000000}%
\pgfsetdash{}{0pt}%
\pgfpathmoveto{\pgfqpoint{2.893588in}{0.882495in}}%
\pgfpathlineto{\pgfqpoint{2.878409in}{0.905708in}}%
\pgfpathlineto{\pgfqpoint{2.881402in}{0.908328in}}%
\pgfpathlineto{\pgfqpoint{2.874063in}{0.928161in}}%
\pgfpathlineto{\pgfqpoint{2.866524in}{0.926111in}}%
\pgfpathlineto{\pgfqpoint{2.864668in}{0.946906in}}%
\pgfpathlineto{\pgfqpoint{2.880855in}{0.952903in}}%
\pgfpathlineto{\pgfqpoint{2.889760in}{0.943778in}}%
\pgfpathlineto{\pgfqpoint{2.898524in}{0.954072in}}%
\pgfpathlineto{\pgfqpoint{2.909363in}{0.953452in}}%
\pgfpathlineto{\pgfqpoint{2.918368in}{0.938239in}}%
\pgfpathlineto{\pgfqpoint{2.911845in}{0.934034in}}%
\pgfpathlineto{\pgfqpoint{2.910167in}{0.929143in}}%
\pgfpathlineto{\pgfqpoint{2.905128in}{0.926113in}}%
\pgfpathlineto{\pgfqpoint{2.899562in}{0.918887in}}%
\pgfpathlineto{\pgfqpoint{2.899980in}{0.912387in}}%
\pgfpathlineto{\pgfqpoint{2.903959in}{0.905713in}}%
\pgfpathlineto{\pgfqpoint{2.909397in}{0.902008in}}%
\pgfpathlineto{\pgfqpoint{2.910191in}{0.894634in}}%
\pgfpathlineto{\pgfqpoint{2.893588in}{0.882495in}}%
\pgfpathclose%
\pgfusepath{fill}%
\end{pgfscope}%
\begin{pgfscope}%
\pgfpathrectangle{\pgfqpoint{0.100000in}{0.100000in}}{\pgfqpoint{3.608454in}{2.310000in}}%
\pgfusepath{clip}%
\pgfsetbuttcap%
\pgfsetmiterjoin%
\definecolor{currentfill}{rgb}{0.000000,0.611765,0.694118}%
\pgfsetfillcolor{currentfill}%
\pgfsetlinewidth{0.000000pt}%
\definecolor{currentstroke}{rgb}{0.000000,0.000000,0.000000}%
\pgfsetstrokecolor{currentstroke}%
\pgfsetstrokeopacity{0.000000}%
\pgfsetdash{}{0pt}%
\pgfpathmoveto{\pgfqpoint{1.496354in}{2.088951in}}%
\pgfpathlineto{\pgfqpoint{1.498672in}{2.081480in}}%
\pgfpathlineto{\pgfqpoint{1.503848in}{2.080824in}}%
\pgfpathlineto{\pgfqpoint{1.500521in}{2.053891in}}%
\pgfpathlineto{\pgfqpoint{1.496777in}{2.040294in}}%
\pgfpathlineto{\pgfqpoint{1.503629in}{2.039468in}}%
\pgfpathlineto{\pgfqpoint{1.501940in}{2.025689in}}%
\pgfpathlineto{\pgfqpoint{1.506437in}{2.025137in}}%
\pgfpathlineto{\pgfqpoint{1.503208in}{1.998589in}}%
\pgfpathlineto{\pgfqpoint{1.468884in}{2.002862in}}%
\pgfpathlineto{\pgfqpoint{1.469044in}{2.003996in}}%
\pgfpathlineto{\pgfqpoint{1.405583in}{2.011993in}}%
\pgfpathlineto{\pgfqpoint{1.403304in}{2.014885in}}%
\pgfpathlineto{\pgfqpoint{1.407470in}{2.024546in}}%
\pgfpathlineto{\pgfqpoint{1.405800in}{2.033757in}}%
\pgfpathlineto{\pgfqpoint{1.406013in}{2.044122in}}%
\pgfpathlineto{\pgfqpoint{1.411450in}{2.059486in}}%
\pgfpathlineto{\pgfqpoint{1.414702in}{2.064290in}}%
\pgfpathlineto{\pgfqpoint{1.422183in}{2.067656in}}%
\pgfpathlineto{\pgfqpoint{1.428142in}{2.073167in}}%
\pgfpathlineto{\pgfqpoint{1.437875in}{2.069615in}}%
\pgfpathlineto{\pgfqpoint{1.440516in}{2.074870in}}%
\pgfpathlineto{\pgfqpoint{1.449108in}{2.071253in}}%
\pgfpathlineto{\pgfqpoint{1.468292in}{2.069917in}}%
\pgfpathlineto{\pgfqpoint{1.470183in}{2.074395in}}%
\pgfpathlineto{\pgfqpoint{1.476331in}{2.073148in}}%
\pgfpathlineto{\pgfqpoint{1.487031in}{2.078600in}}%
\pgfpathlineto{\pgfqpoint{1.492326in}{2.087153in}}%
\pgfpathlineto{\pgfqpoint{1.496354in}{2.088951in}}%
\pgfpathclose%
\pgfusepath{fill}%
\end{pgfscope}%
\begin{pgfscope}%
\pgfpathrectangle{\pgfqpoint{0.100000in}{0.100000in}}{\pgfqpoint{3.608454in}{2.310000in}}%
\pgfusepath{clip}%
\pgfsetbuttcap%
\pgfsetmiterjoin%
\definecolor{currentfill}{rgb}{0.000000,0.423529,0.788235}%
\pgfsetfillcolor{currentfill}%
\pgfsetlinewidth{0.000000pt}%
\definecolor{currentstroke}{rgb}{0.000000,0.000000,0.000000}%
\pgfsetstrokecolor{currentstroke}%
\pgfsetstrokeopacity{0.000000}%
\pgfsetdash{}{0pt}%
\pgfpathmoveto{\pgfqpoint{3.153317in}{1.449269in}}%
\pgfpathlineto{\pgfqpoint{3.117245in}{1.469315in}}%
\pgfpathlineto{\pgfqpoint{3.109226in}{1.471916in}}%
\pgfpathlineto{\pgfqpoint{3.110506in}{1.479451in}}%
\pgfpathlineto{\pgfqpoint{3.116008in}{1.487243in}}%
\pgfpathlineto{\pgfqpoint{3.111561in}{1.488897in}}%
\pgfpathlineto{\pgfqpoint{3.114952in}{1.497798in}}%
\pgfpathlineto{\pgfqpoint{3.117244in}{1.511717in}}%
\pgfpathlineto{\pgfqpoint{3.116662in}{1.519085in}}%
\pgfpathlineto{\pgfqpoint{3.119792in}{1.526024in}}%
\pgfpathlineto{\pgfqpoint{3.133137in}{1.523448in}}%
\pgfpathlineto{\pgfqpoint{3.137378in}{1.521010in}}%
\pgfpathlineto{\pgfqpoint{3.143924in}{1.538541in}}%
\pgfpathlineto{\pgfqpoint{3.146592in}{1.539583in}}%
\pgfpathlineto{\pgfqpoint{3.162472in}{1.516430in}}%
\pgfpathlineto{\pgfqpoint{3.166603in}{1.499150in}}%
\pgfpathlineto{\pgfqpoint{3.161516in}{1.490087in}}%
\pgfpathlineto{\pgfqpoint{3.159520in}{1.465911in}}%
\pgfpathlineto{\pgfqpoint{3.156244in}{1.464323in}}%
\pgfpathlineto{\pgfqpoint{3.153317in}{1.449269in}}%
\pgfpathclose%
\pgfusepath{fill}%
\end{pgfscope}%
\begin{pgfscope}%
\pgfpathrectangle{\pgfqpoint{0.100000in}{0.100000in}}{\pgfqpoint{3.608454in}{2.310000in}}%
\pgfusepath{clip}%
\pgfsetbuttcap%
\pgfsetmiterjoin%
\definecolor{currentfill}{rgb}{0.000000,0.827451,0.586275}%
\pgfsetfillcolor{currentfill}%
\pgfsetlinewidth{0.000000pt}%
\definecolor{currentstroke}{rgb}{0.000000,0.000000,0.000000}%
\pgfsetstrokecolor{currentstroke}%
\pgfsetstrokeopacity{0.000000}%
\pgfsetdash{}{0pt}%
\pgfpathmoveto{\pgfqpoint{0.735331in}{0.472883in}}%
\pgfpathlineto{\pgfqpoint{0.738160in}{0.475971in}}%
\pgfpathlineto{\pgfqpoint{0.743525in}{0.474782in}}%
\pgfpathlineto{\pgfqpoint{0.740586in}{0.472887in}}%
\pgfpathlineto{\pgfqpoint{0.735331in}{0.472883in}}%
\pgfpathclose%
\pgfusepath{fill}%
\end{pgfscope}%
\begin{pgfscope}%
\pgfpathrectangle{\pgfqpoint{0.100000in}{0.100000in}}{\pgfqpoint{3.608454in}{2.310000in}}%
\pgfusepath{clip}%
\pgfsetbuttcap%
\pgfsetmiterjoin%
\definecolor{currentfill}{rgb}{0.000000,0.827451,0.586275}%
\pgfsetfillcolor{currentfill}%
\pgfsetlinewidth{0.000000pt}%
\definecolor{currentstroke}{rgb}{0.000000,0.000000,0.000000}%
\pgfsetstrokecolor{currentstroke}%
\pgfsetstrokeopacity{0.000000}%
\pgfsetdash{}{0pt}%
\pgfpathmoveto{\pgfqpoint{0.864761in}{0.455531in}}%
\pgfpathlineto{\pgfqpoint{0.863736in}{0.453184in}}%
\pgfpathlineto{\pgfqpoint{0.857604in}{0.444984in}}%
\pgfpathlineto{\pgfqpoint{0.856061in}{0.446120in}}%
\pgfpathlineto{\pgfqpoint{0.849915in}{0.438009in}}%
\pgfpathlineto{\pgfqpoint{0.848446in}{0.439078in}}%
\pgfpathlineto{\pgfqpoint{0.842589in}{0.431260in}}%
\pgfpathlineto{\pgfqpoint{0.841100in}{0.432388in}}%
\pgfpathlineto{\pgfqpoint{0.839547in}{0.430349in}}%
\pgfpathlineto{\pgfqpoint{0.837531in}{0.431857in}}%
\pgfpathlineto{\pgfqpoint{0.834469in}{0.427785in}}%
\pgfpathlineto{\pgfqpoint{0.832475in}{0.429280in}}%
\pgfpathlineto{\pgfqpoint{0.830966in}{0.427222in}}%
\pgfpathlineto{\pgfqpoint{0.827494in}{0.429779in}}%
\pgfpathlineto{\pgfqpoint{0.824498in}{0.425742in}}%
\pgfpathlineto{\pgfqpoint{0.820512in}{0.428750in}}%
\pgfpathlineto{\pgfqpoint{0.817435in}{0.424709in}}%
\pgfpathlineto{\pgfqpoint{0.818042in}{0.424203in}}%
\pgfpathlineto{\pgfqpoint{0.812034in}{0.416427in}}%
\pgfpathlineto{\pgfqpoint{0.811149in}{0.413923in}}%
\pgfpathlineto{\pgfqpoint{0.815079in}{0.410909in}}%
\pgfpathlineto{\pgfqpoint{0.811859in}{0.406633in}}%
\pgfpathlineto{\pgfqpoint{0.809601in}{0.406866in}}%
\pgfpathlineto{\pgfqpoint{0.808196in}{0.405860in}}%
\pgfpathlineto{\pgfqpoint{0.804050in}{0.407001in}}%
\pgfpathlineto{\pgfqpoint{0.800501in}{0.409249in}}%
\pgfpathlineto{\pgfqpoint{0.798949in}{0.408471in}}%
\pgfpathlineto{\pgfqpoint{0.797005in}{0.409992in}}%
\pgfpathlineto{\pgfqpoint{0.794932in}{0.410045in}}%
\pgfpathlineto{\pgfqpoint{0.789736in}{0.403394in}}%
\pgfpathlineto{\pgfqpoint{0.786031in}{0.403802in}}%
\pgfpathlineto{\pgfqpoint{0.784195in}{0.405252in}}%
\pgfpathlineto{\pgfqpoint{0.781062in}{0.403097in}}%
\pgfpathlineto{\pgfqpoint{0.777277in}{0.402089in}}%
\pgfpathlineto{\pgfqpoint{0.777113in}{0.402156in}}%
\pgfpathlineto{\pgfqpoint{0.777652in}{0.403284in}}%
\pgfpathlineto{\pgfqpoint{0.773253in}{0.404050in}}%
\pgfpathlineto{\pgfqpoint{0.771380in}{0.405564in}}%
\pgfpathlineto{\pgfqpoint{0.768215in}{0.404903in}}%
\pgfpathlineto{\pgfqpoint{0.766993in}{0.403276in}}%
\pgfpathlineto{\pgfqpoint{0.765003in}{0.403792in}}%
\pgfpathlineto{\pgfqpoint{0.761612in}{0.402864in}}%
\pgfpathlineto{\pgfqpoint{0.760327in}{0.401259in}}%
\pgfpathlineto{\pgfqpoint{0.764308in}{0.397983in}}%
\pgfpathlineto{\pgfqpoint{0.759811in}{0.395476in}}%
\pgfpathlineto{\pgfqpoint{0.758622in}{0.397717in}}%
\pgfpathlineto{\pgfqpoint{0.754992in}{0.396142in}}%
\pgfpathlineto{\pgfqpoint{0.753722in}{0.397568in}}%
\pgfpathlineto{\pgfqpoint{0.743096in}{0.399381in}}%
\pgfpathlineto{\pgfqpoint{0.741766in}{0.396518in}}%
\pgfpathlineto{\pgfqpoint{0.739332in}{0.398261in}}%
\pgfpathlineto{\pgfqpoint{0.739251in}{0.400641in}}%
\pgfpathlineto{\pgfqpoint{0.736510in}{0.401202in}}%
\pgfpathlineto{\pgfqpoint{0.736143in}{0.398969in}}%
\pgfpathlineto{\pgfqpoint{0.734429in}{0.397349in}}%
\pgfpathlineto{\pgfqpoint{0.730331in}{0.398641in}}%
\pgfpathlineto{\pgfqpoint{0.727510in}{0.401229in}}%
\pgfpathlineto{\pgfqpoint{0.723876in}{0.399715in}}%
\pgfpathlineto{\pgfqpoint{0.726580in}{0.395915in}}%
\pgfpathlineto{\pgfqpoint{0.720792in}{0.394591in}}%
\pgfpathlineto{\pgfqpoint{0.718356in}{0.391937in}}%
\pgfpathlineto{\pgfqpoint{0.717498in}{0.394606in}}%
\pgfpathlineto{\pgfqpoint{0.714729in}{0.394241in}}%
\pgfpathlineto{\pgfqpoint{0.712644in}{0.395735in}}%
\pgfpathlineto{\pgfqpoint{0.710848in}{0.395775in}}%
\pgfpathlineto{\pgfqpoint{0.705221in}{0.398449in}}%
\pgfpathlineto{\pgfqpoint{0.702788in}{0.397277in}}%
\pgfpathlineto{\pgfqpoint{0.700932in}{0.397327in}}%
\pgfpathlineto{\pgfqpoint{0.698256in}{0.395576in}}%
\pgfpathlineto{\pgfqpoint{0.702855in}{0.400929in}}%
\pgfpathlineto{\pgfqpoint{0.699223in}{0.404064in}}%
\pgfpathlineto{\pgfqpoint{0.701711in}{0.405421in}}%
\pgfpathlineto{\pgfqpoint{0.707923in}{0.412710in}}%
\pgfpathlineto{\pgfqpoint{0.706826in}{0.413694in}}%
\pgfpathlineto{\pgfqpoint{0.710015in}{0.417271in}}%
\pgfpathlineto{\pgfqpoint{0.713680in}{0.414200in}}%
\pgfpathlineto{\pgfqpoint{0.715284in}{0.416046in}}%
\pgfpathlineto{\pgfqpoint{0.719039in}{0.412825in}}%
\pgfpathlineto{\pgfqpoint{0.720711in}{0.414781in}}%
\pgfpathlineto{\pgfqpoint{0.723593in}{0.412237in}}%
\pgfpathlineto{\pgfqpoint{0.725147in}{0.414065in}}%
\pgfpathlineto{\pgfqpoint{0.727666in}{0.413737in}}%
\pgfpathlineto{\pgfqpoint{0.730412in}{0.411143in}}%
\pgfpathlineto{\pgfqpoint{0.733440in}{0.415672in}}%
\pgfpathlineto{\pgfqpoint{0.739202in}{0.417909in}}%
\pgfpathlineto{\pgfqpoint{0.745161in}{0.418898in}}%
\pgfpathlineto{\pgfqpoint{0.748345in}{0.417706in}}%
\pgfpathlineto{\pgfqpoint{0.751670in}{0.419346in}}%
\pgfpathlineto{\pgfqpoint{0.756304in}{0.418993in}}%
\pgfpathlineto{\pgfqpoint{0.757747in}{0.421251in}}%
\pgfpathlineto{\pgfqpoint{0.768447in}{0.430721in}}%
\pgfpathlineto{\pgfqpoint{0.770758in}{0.430833in}}%
\pgfpathlineto{\pgfqpoint{0.770614in}{0.432762in}}%
\pgfpathlineto{\pgfqpoint{0.774098in}{0.436410in}}%
\pgfpathlineto{\pgfqpoint{0.779568in}{0.437653in}}%
\pgfpathlineto{\pgfqpoint{0.790893in}{0.428785in}}%
\pgfpathlineto{\pgfqpoint{0.796011in}{0.435184in}}%
\pgfpathlineto{\pgfqpoint{0.787216in}{0.442053in}}%
\pgfpathlineto{\pgfqpoint{0.772048in}{0.445375in}}%
\pgfpathlineto{\pgfqpoint{0.769184in}{0.446551in}}%
\pgfpathlineto{\pgfqpoint{0.768571in}{0.448855in}}%
\pgfpathlineto{\pgfqpoint{0.769448in}{0.452003in}}%
\pgfpathlineto{\pgfqpoint{0.768194in}{0.453396in}}%
\pgfpathlineto{\pgfqpoint{0.771817in}{0.456869in}}%
\pgfpathlineto{\pgfqpoint{0.770582in}{0.457460in}}%
\pgfpathlineto{\pgfqpoint{0.763393in}{0.455439in}}%
\pgfpathlineto{\pgfqpoint{0.761052in}{0.449862in}}%
\pgfpathlineto{\pgfqpoint{0.759380in}{0.447981in}}%
\pgfpathlineto{\pgfqpoint{0.757127in}{0.448341in}}%
\pgfpathlineto{\pgfqpoint{0.755962in}{0.451093in}}%
\pgfpathlineto{\pgfqpoint{0.756950in}{0.451960in}}%
\pgfpathlineto{\pgfqpoint{0.756995in}{0.455328in}}%
\pgfpathlineto{\pgfqpoint{0.758416in}{0.463072in}}%
\pgfpathlineto{\pgfqpoint{0.757580in}{0.464630in}}%
\pgfpathlineto{\pgfqpoint{0.759553in}{0.467203in}}%
\pgfpathlineto{\pgfqpoint{0.757448in}{0.468248in}}%
\pgfpathlineto{\pgfqpoint{0.756266in}{0.467141in}}%
\pgfpathlineto{\pgfqpoint{0.752660in}{0.466199in}}%
\pgfpathlineto{\pgfqpoint{0.752967in}{0.468537in}}%
\pgfpathlineto{\pgfqpoint{0.751538in}{0.473475in}}%
\pgfpathlineto{\pgfqpoint{0.753567in}{0.476503in}}%
\pgfpathlineto{\pgfqpoint{0.750953in}{0.477091in}}%
\pgfpathlineto{\pgfqpoint{0.745672in}{0.477306in}}%
\pgfpathlineto{\pgfqpoint{0.740722in}{0.478712in}}%
\pgfpathlineto{\pgfqpoint{0.744357in}{0.482168in}}%
\pgfpathlineto{\pgfqpoint{0.746173in}{0.480590in}}%
\pgfpathlineto{\pgfqpoint{0.747682in}{0.482304in}}%
\pgfpathlineto{\pgfqpoint{0.749547in}{0.480704in}}%
\pgfpathlineto{\pgfqpoint{0.752873in}{0.484455in}}%
\pgfpathlineto{\pgfqpoint{0.755169in}{0.485784in}}%
\pgfpathlineto{\pgfqpoint{0.756995in}{0.484116in}}%
\pgfpathlineto{\pgfqpoint{0.759307in}{0.485494in}}%
\pgfpathlineto{\pgfqpoint{0.762636in}{0.489281in}}%
\pgfpathlineto{\pgfqpoint{0.764542in}{0.487613in}}%
\pgfpathlineto{\pgfqpoint{0.767845in}{0.491365in}}%
\pgfpathlineto{\pgfqpoint{0.770402in}{0.489224in}}%
\pgfpathlineto{\pgfqpoint{0.773696in}{0.492989in}}%
\pgfpathlineto{\pgfqpoint{0.775519in}{0.491429in}}%
\pgfpathlineto{\pgfqpoint{0.778697in}{0.494994in}}%
\pgfpathlineto{\pgfqpoint{0.781042in}{0.496290in}}%
\pgfpathlineto{\pgfqpoint{0.782945in}{0.494696in}}%
\pgfpathlineto{\pgfqpoint{0.786206in}{0.498438in}}%
\pgfpathlineto{\pgfqpoint{0.788112in}{0.496794in}}%
\pgfpathlineto{\pgfqpoint{0.789763in}{0.498705in}}%
\pgfpathlineto{\pgfqpoint{0.791902in}{0.496847in}}%
\pgfpathlineto{\pgfqpoint{0.795190in}{0.500666in}}%
\pgfpathlineto{\pgfqpoint{0.797356in}{0.502139in}}%
\pgfpathlineto{\pgfqpoint{0.799211in}{0.500636in}}%
\pgfpathlineto{\pgfqpoint{0.800941in}{0.502662in}}%
\pgfpathlineto{\pgfqpoint{0.802432in}{0.501473in}}%
\pgfpathlineto{\pgfqpoint{0.804074in}{0.503387in}}%
\pgfpathlineto{\pgfqpoint{0.814957in}{0.494457in}}%
\pgfpathlineto{\pgfqpoint{0.813395in}{0.492489in}}%
\pgfpathlineto{\pgfqpoint{0.815390in}{0.490823in}}%
\pgfpathlineto{\pgfqpoint{0.816987in}{0.492818in}}%
\pgfpathlineto{\pgfqpoint{0.819027in}{0.491316in}}%
\pgfpathlineto{\pgfqpoint{0.820776in}{0.493434in}}%
\pgfpathlineto{\pgfqpoint{0.824629in}{0.490280in}}%
\pgfpathlineto{\pgfqpoint{0.822897in}{0.488128in}}%
\pgfpathlineto{\pgfqpoint{0.836300in}{0.477446in}}%
\pgfpathlineto{\pgfqpoint{0.853555in}{0.463966in}}%
\pgfpathlineto{\pgfqpoint{0.864761in}{0.455531in}}%
\pgfpathclose%
\pgfusepath{fill}%
\end{pgfscope}%
\begin{pgfscope}%
\pgfpathrectangle{\pgfqpoint{0.100000in}{0.100000in}}{\pgfqpoint{3.608454in}{2.310000in}}%
\pgfusepath{clip}%
\pgfsetbuttcap%
\pgfsetmiterjoin%
\definecolor{currentfill}{rgb}{0.000000,0.827451,0.586275}%
\pgfsetfillcolor{currentfill}%
\pgfsetlinewidth{0.000000pt}%
\definecolor{currentstroke}{rgb}{0.000000,0.000000,0.000000}%
\pgfsetstrokecolor{currentstroke}%
\pgfsetstrokeopacity{0.000000}%
\pgfsetdash{}{0pt}%
\pgfpathmoveto{\pgfqpoint{0.735963in}{0.479492in}}%
\pgfpathlineto{\pgfqpoint{0.737273in}{0.481661in}}%
\pgfpathlineto{\pgfqpoint{0.740538in}{0.478732in}}%
\pgfpathlineto{\pgfqpoint{0.735963in}{0.479492in}}%
\pgfpathclose%
\pgfusepath{fill}%
\end{pgfscope}%
\begin{pgfscope}%
\pgfpathrectangle{\pgfqpoint{0.100000in}{0.100000in}}{\pgfqpoint{3.608454in}{2.310000in}}%
\pgfusepath{clip}%
\pgfsetbuttcap%
\pgfsetmiterjoin%
\definecolor{currentfill}{rgb}{0.000000,0.694118,0.652941}%
\pgfsetfillcolor{currentfill}%
\pgfsetlinewidth{0.000000pt}%
\definecolor{currentstroke}{rgb}{0.000000,0.000000,0.000000}%
\pgfsetstrokecolor{currentstroke}%
\pgfsetstrokeopacity{0.000000}%
\pgfsetdash{}{0pt}%
\pgfpathmoveto{\pgfqpoint{1.128763in}{1.959930in}}%
\pgfpathlineto{\pgfqpoint{1.126427in}{1.948422in}}%
\pgfpathlineto{\pgfqpoint{1.131774in}{1.943515in}}%
\pgfpathlineto{\pgfqpoint{1.135886in}{1.943764in}}%
\pgfpathlineto{\pgfqpoint{1.137082in}{1.935988in}}%
\pgfpathlineto{\pgfqpoint{1.133161in}{1.915579in}}%
\pgfpathlineto{\pgfqpoint{1.139942in}{1.914263in}}%
\pgfpathlineto{\pgfqpoint{1.138589in}{1.907407in}}%
\pgfpathlineto{\pgfqpoint{1.146489in}{1.905873in}}%
\pgfpathlineto{\pgfqpoint{1.145219in}{1.892062in}}%
\pgfpathlineto{\pgfqpoint{1.150185in}{1.887641in}}%
\pgfpathlineto{\pgfqpoint{1.165886in}{1.884470in}}%
\pgfpathlineto{\pgfqpoint{1.164838in}{1.878825in}}%
\pgfpathlineto{\pgfqpoint{1.181983in}{1.875682in}}%
\pgfpathlineto{\pgfqpoint{1.177604in}{1.866287in}}%
\pgfpathlineto{\pgfqpoint{1.166980in}{1.866417in}}%
\pgfpathlineto{\pgfqpoint{1.159714in}{1.863660in}}%
\pgfpathlineto{\pgfqpoint{1.157878in}{1.868529in}}%
\pgfpathlineto{\pgfqpoint{1.144630in}{1.867633in}}%
\pgfpathlineto{\pgfqpoint{1.134588in}{1.873151in}}%
\pgfpathlineto{\pgfqpoint{1.130114in}{1.870874in}}%
\pgfpathlineto{\pgfqpoint{1.127333in}{1.864844in}}%
\pgfpathlineto{\pgfqpoint{1.123023in}{1.868301in}}%
\pgfpathlineto{\pgfqpoint{1.105971in}{1.872044in}}%
\pgfpathlineto{\pgfqpoint{1.102494in}{1.863989in}}%
\pgfpathlineto{\pgfqpoint{1.100477in}{1.862882in}}%
\pgfpathlineto{\pgfqpoint{1.093092in}{1.871853in}}%
\pgfpathlineto{\pgfqpoint{1.094071in}{1.877735in}}%
\pgfpathlineto{\pgfqpoint{1.091986in}{1.894613in}}%
\pgfpathlineto{\pgfqpoint{1.088102in}{1.901286in}}%
\pgfpathlineto{\pgfqpoint{1.079391in}{1.904532in}}%
\pgfpathlineto{\pgfqpoint{1.074674in}{1.909170in}}%
\pgfpathlineto{\pgfqpoint{1.077934in}{1.924256in}}%
\pgfpathlineto{\pgfqpoint{1.073148in}{1.929599in}}%
\pgfpathlineto{\pgfqpoint{1.066629in}{1.948325in}}%
\pgfpathlineto{\pgfqpoint{1.067866in}{1.953218in}}%
\pgfpathlineto{\pgfqpoint{1.064456in}{1.964544in}}%
\pgfpathlineto{\pgfqpoint{1.067306in}{1.970565in}}%
\pgfpathlineto{\pgfqpoint{1.061696in}{1.979359in}}%
\pgfpathlineto{\pgfqpoint{1.069612in}{1.982716in}}%
\pgfpathlineto{\pgfqpoint{1.076963in}{1.988148in}}%
\pgfpathlineto{\pgfqpoint{1.086269in}{1.988945in}}%
\pgfpathlineto{\pgfqpoint{1.088551in}{1.993499in}}%
\pgfpathlineto{\pgfqpoint{1.092311in}{1.978823in}}%
\pgfpathlineto{\pgfqpoint{1.103252in}{1.984857in}}%
\pgfpathlineto{\pgfqpoint{1.110227in}{1.984421in}}%
\pgfpathlineto{\pgfqpoint{1.113810in}{1.978599in}}%
\pgfpathlineto{\pgfqpoint{1.126712in}{1.970229in}}%
\pgfpathlineto{\pgfqpoint{1.128763in}{1.959930in}}%
\pgfpathclose%
\pgfusepath{fill}%
\end{pgfscope}%
\begin{pgfscope}%
\pgfpathrectangle{\pgfqpoint{0.100000in}{0.100000in}}{\pgfqpoint{3.608454in}{2.310000in}}%
\pgfusepath{clip}%
\pgfsetbuttcap%
\pgfsetmiterjoin%
\definecolor{currentfill}{rgb}{0.000000,0.607843,0.696078}%
\pgfsetfillcolor{currentfill}%
\pgfsetlinewidth{0.000000pt}%
\definecolor{currentstroke}{rgb}{0.000000,0.000000,0.000000}%
\pgfsetstrokecolor{currentstroke}%
\pgfsetstrokeopacity{0.000000}%
\pgfsetdash{}{0pt}%
\pgfpathmoveto{\pgfqpoint{0.869214in}{2.048910in}}%
\pgfpathlineto{\pgfqpoint{0.906027in}{2.039422in}}%
\pgfpathlineto{\pgfqpoint{0.909977in}{2.024514in}}%
\pgfpathlineto{\pgfqpoint{0.914485in}{2.022985in}}%
\pgfpathlineto{\pgfqpoint{0.921587in}{2.014227in}}%
\pgfpathlineto{\pgfqpoint{0.923281in}{2.004122in}}%
\pgfpathlineto{\pgfqpoint{0.916109in}{1.996717in}}%
\pgfpathlineto{\pgfqpoint{0.904646in}{1.980156in}}%
\pgfpathlineto{\pgfqpoint{0.899963in}{1.971006in}}%
\pgfpathlineto{\pgfqpoint{0.895815in}{1.966858in}}%
\pgfpathlineto{\pgfqpoint{0.869552in}{1.973496in}}%
\pgfpathlineto{\pgfqpoint{0.871160in}{1.980064in}}%
\pgfpathlineto{\pgfqpoint{0.859744in}{1.982888in}}%
\pgfpathlineto{\pgfqpoint{0.859301in}{1.991402in}}%
\pgfpathlineto{\pgfqpoint{0.854053in}{1.998430in}}%
\pgfpathlineto{\pgfqpoint{0.853023in}{2.013264in}}%
\pgfpathlineto{\pgfqpoint{0.855049in}{2.021169in}}%
\pgfpathlineto{\pgfqpoint{0.853215in}{2.027577in}}%
\pgfpathlineto{\pgfqpoint{0.858767in}{2.040242in}}%
\pgfpathlineto{\pgfqpoint{0.846645in}{2.043323in}}%
\pgfpathlineto{\pgfqpoint{0.849176in}{2.054195in}}%
\pgfpathlineto{\pgfqpoint{0.869214in}{2.048910in}}%
\pgfpathclose%
\pgfusepath{fill}%
\end{pgfscope}%
\begin{pgfscope}%
\pgfpathrectangle{\pgfqpoint{0.100000in}{0.100000in}}{\pgfqpoint{3.608454in}{2.310000in}}%
\pgfusepath{clip}%
\pgfsetbuttcap%
\pgfsetmiterjoin%
\definecolor{currentfill}{rgb}{0.000000,0.364706,0.817647}%
\pgfsetfillcolor{currentfill}%
\pgfsetlinewidth{0.000000pt}%
\definecolor{currentstroke}{rgb}{0.000000,0.000000,0.000000}%
\pgfsetstrokecolor{currentstroke}%
\pgfsetstrokeopacity{0.000000}%
\pgfsetdash{}{0pt}%
\pgfpathmoveto{\pgfqpoint{1.707133in}{1.330203in}}%
\pgfpathlineto{\pgfqpoint{1.672634in}{1.332579in}}%
\pgfpathlineto{\pgfqpoint{1.677687in}{1.401208in}}%
\pgfpathlineto{\pgfqpoint{1.716323in}{1.398352in}}%
\pgfpathlineto{\pgfqpoint{1.717833in}{1.398242in}}%
\pgfpathlineto{\pgfqpoint{1.715265in}{1.364101in}}%
\pgfpathlineto{\pgfqpoint{1.709976in}{1.364447in}}%
\pgfpathlineto{\pgfqpoint{1.707133in}{1.330203in}}%
\pgfpathclose%
\pgfusepath{fill}%
\end{pgfscope}%
\begin{pgfscope}%
\pgfpathrectangle{\pgfqpoint{0.100000in}{0.100000in}}{\pgfqpoint{3.608454in}{2.310000in}}%
\pgfusepath{clip}%
\pgfsetbuttcap%
\pgfsetmiterjoin%
\definecolor{currentfill}{rgb}{0.000000,0.243137,0.878431}%
\pgfsetfillcolor{currentfill}%
\pgfsetlinewidth{0.000000pt}%
\definecolor{currentstroke}{rgb}{0.000000,0.000000,0.000000}%
\pgfsetstrokecolor{currentstroke}%
\pgfsetstrokeopacity{0.000000}%
\pgfsetdash{}{0pt}%
\pgfpathmoveto{\pgfqpoint{2.034871in}{0.852102in}}%
\pgfpathlineto{\pgfqpoint{2.034479in}{0.813847in}}%
\pgfpathlineto{\pgfqpoint{2.013777in}{0.811827in}}%
\pgfpathlineto{\pgfqpoint{1.978879in}{0.792186in}}%
\pgfpathlineto{\pgfqpoint{1.975845in}{0.790506in}}%
\pgfpathlineto{\pgfqpoint{1.966416in}{0.807704in}}%
\pgfpathlineto{\pgfqpoint{1.966891in}{0.830295in}}%
\pgfpathlineto{\pgfqpoint{1.970138in}{0.830158in}}%
\pgfpathlineto{\pgfqpoint{1.971315in}{0.864806in}}%
\pgfpathlineto{\pgfqpoint{1.946927in}{0.865843in}}%
\pgfpathlineto{\pgfqpoint{1.948881in}{0.900975in}}%
\pgfpathlineto{\pgfqpoint{1.941990in}{0.901235in}}%
\pgfpathlineto{\pgfqpoint{1.943227in}{0.939444in}}%
\pgfpathlineto{\pgfqpoint{1.945662in}{0.932105in}}%
\pgfpathlineto{\pgfqpoint{1.950349in}{0.931638in}}%
\pgfpathlineto{\pgfqpoint{1.958950in}{0.938059in}}%
\pgfpathlineto{\pgfqpoint{1.963046in}{0.937524in}}%
\pgfpathlineto{\pgfqpoint{1.961356in}{0.929934in}}%
\pgfpathlineto{\pgfqpoint{1.964587in}{0.923699in}}%
\pgfpathlineto{\pgfqpoint{1.968361in}{0.923919in}}%
\pgfpathlineto{\pgfqpoint{1.978804in}{0.941263in}}%
\pgfpathlineto{\pgfqpoint{1.977949in}{0.898974in}}%
\pgfpathlineto{\pgfqpoint{1.984737in}{0.897974in}}%
\pgfpathlineto{\pgfqpoint{2.014996in}{0.896820in}}%
\pgfpathlineto{\pgfqpoint{2.020851in}{0.893088in}}%
\pgfpathlineto{\pgfqpoint{2.020147in}{0.852551in}}%
\pgfpathlineto{\pgfqpoint{2.034871in}{0.852102in}}%
\pgfpathclose%
\pgfusepath{fill}%
\end{pgfscope}%
\begin{pgfscope}%
\pgfpathrectangle{\pgfqpoint{0.100000in}{0.100000in}}{\pgfqpoint{3.608454in}{2.310000in}}%
\pgfusepath{clip}%
\pgfsetbuttcap%
\pgfsetmiterjoin%
\definecolor{currentfill}{rgb}{0.000000,0.545098,0.727451}%
\pgfsetfillcolor{currentfill}%
\pgfsetlinewidth{0.000000pt}%
\definecolor{currentstroke}{rgb}{0.000000,0.000000,0.000000}%
\pgfsetstrokecolor{currentstroke}%
\pgfsetstrokeopacity{0.000000}%
\pgfsetdash{}{0pt}%
\pgfpathmoveto{\pgfqpoint{1.978804in}{0.941263in}}%
\pgfpathlineto{\pgfqpoint{1.968361in}{0.923919in}}%
\pgfpathlineto{\pgfqpoint{1.964587in}{0.923699in}}%
\pgfpathlineto{\pgfqpoint{1.961356in}{0.929934in}}%
\pgfpathlineto{\pgfqpoint{1.963046in}{0.937524in}}%
\pgfpathlineto{\pgfqpoint{1.958950in}{0.938059in}}%
\pgfpathlineto{\pgfqpoint{1.950349in}{0.931638in}}%
\pgfpathlineto{\pgfqpoint{1.945662in}{0.932105in}}%
\pgfpathlineto{\pgfqpoint{1.943227in}{0.939444in}}%
\pgfpathlineto{\pgfqpoint{1.938143in}{0.938336in}}%
\pgfpathlineto{\pgfqpoint{1.939766in}{1.000299in}}%
\pgfpathlineto{\pgfqpoint{1.932881in}{1.000512in}}%
\pgfpathlineto{\pgfqpoint{1.933314in}{1.014281in}}%
\pgfpathlineto{\pgfqpoint{1.981120in}{1.013018in}}%
\pgfpathlineto{\pgfqpoint{1.980692in}{0.992304in}}%
\pgfpathlineto{\pgfqpoint{1.987552in}{0.992168in}}%
\pgfpathlineto{\pgfqpoint{1.987413in}{0.985257in}}%
\pgfpathlineto{\pgfqpoint{2.007887in}{0.984833in}}%
\pgfpathlineto{\pgfqpoint{2.007778in}{0.978034in}}%
\pgfpathlineto{\pgfqpoint{2.014740in}{0.977835in}}%
\pgfpathlineto{\pgfqpoint{2.014398in}{0.957120in}}%
\pgfpathlineto{\pgfqpoint{2.000025in}{0.945069in}}%
\pgfpathlineto{\pgfqpoint{1.993908in}{0.931626in}}%
\pgfpathlineto{\pgfqpoint{1.982724in}{0.934735in}}%
\pgfpathlineto{\pgfqpoint{1.978804in}{0.941263in}}%
\pgfpathclose%
\pgfusepath{fill}%
\end{pgfscope}%
\begin{pgfscope}%
\pgfpathrectangle{\pgfqpoint{0.100000in}{0.100000in}}{\pgfqpoint{3.608454in}{2.310000in}}%
\pgfusepath{clip}%
\pgfsetbuttcap%
\pgfsetmiterjoin%
\definecolor{currentfill}{rgb}{0.000000,0.674510,0.662745}%
\pgfsetfillcolor{currentfill}%
\pgfsetlinewidth{0.000000pt}%
\definecolor{currentstroke}{rgb}{0.000000,0.000000,0.000000}%
\pgfsetstrokecolor{currentstroke}%
\pgfsetstrokeopacity{0.000000}%
\pgfsetdash{}{0pt}%
\pgfpathmoveto{\pgfqpoint{1.355496in}{0.949729in}}%
\pgfpathlineto{\pgfqpoint{1.356403in}{0.942467in}}%
\pgfpathlineto{\pgfqpoint{1.353414in}{0.915777in}}%
\pgfpathlineto{\pgfqpoint{1.314868in}{0.902303in}}%
\pgfpathlineto{\pgfqpoint{1.314425in}{0.898871in}}%
\pgfpathlineto{\pgfqpoint{1.287163in}{0.902597in}}%
\pgfpathlineto{\pgfqpoint{1.285282in}{0.888829in}}%
\pgfpathlineto{\pgfqpoint{1.264784in}{0.891644in}}%
\pgfpathlineto{\pgfqpoint{1.256509in}{0.894758in}}%
\pgfpathlineto{\pgfqpoint{1.259177in}{0.906354in}}%
\pgfpathlineto{\pgfqpoint{1.256598in}{0.912844in}}%
\pgfpathlineto{\pgfqpoint{1.257767in}{0.919573in}}%
\pgfpathlineto{\pgfqpoint{1.253634in}{0.922827in}}%
\pgfpathlineto{\pgfqpoint{1.255093in}{0.930635in}}%
\pgfpathlineto{\pgfqpoint{1.253754in}{0.941144in}}%
\pgfpathlineto{\pgfqpoint{1.245517in}{0.942241in}}%
\pgfpathlineto{\pgfqpoint{1.248648in}{0.963976in}}%
\pgfpathlineto{\pgfqpoint{1.252648in}{0.963109in}}%
\pgfpathlineto{\pgfqpoint{1.318032in}{0.954222in}}%
\pgfpathlineto{\pgfqpoint{1.355496in}{0.949729in}}%
\pgfpathclose%
\pgfusepath{fill}%
\end{pgfscope}%
\begin{pgfscope}%
\pgfpathrectangle{\pgfqpoint{0.100000in}{0.100000in}}{\pgfqpoint{3.608454in}{2.310000in}}%
\pgfusepath{clip}%
\pgfsetbuttcap%
\pgfsetmiterjoin%
\definecolor{currentfill}{rgb}{0.000000,0.329412,0.835294}%
\pgfsetfillcolor{currentfill}%
\pgfsetlinewidth{0.000000pt}%
\definecolor{currentstroke}{rgb}{0.000000,0.000000,0.000000}%
\pgfsetstrokecolor{currentstroke}%
\pgfsetstrokeopacity{0.000000}%
\pgfsetdash{}{0pt}%
\pgfpathmoveto{\pgfqpoint{2.042559in}{1.951936in}}%
\pgfpathlineto{\pgfqpoint{2.042594in}{1.944974in}}%
\pgfpathlineto{\pgfqpoint{2.036642in}{1.945064in}}%
\pgfpathlineto{\pgfqpoint{2.036301in}{1.917281in}}%
\pgfpathlineto{\pgfqpoint{2.036803in}{1.896501in}}%
\pgfpathlineto{\pgfqpoint{2.019610in}{1.896790in}}%
\pgfpathlineto{\pgfqpoint{2.020229in}{1.889881in}}%
\pgfpathlineto{\pgfqpoint{1.983409in}{1.890641in}}%
\pgfpathlineto{\pgfqpoint{1.981775in}{1.895615in}}%
\pgfpathlineto{\pgfqpoint{1.981211in}{1.918353in}}%
\pgfpathlineto{\pgfqpoint{1.981936in}{1.946008in}}%
\pgfpathlineto{\pgfqpoint{1.959926in}{1.946668in}}%
\pgfpathlineto{\pgfqpoint{1.960762in}{1.974477in}}%
\pgfpathlineto{\pgfqpoint{1.960057in}{1.995378in}}%
\pgfpathlineto{\pgfqpoint{2.007393in}{1.994001in}}%
\pgfpathlineto{\pgfqpoint{2.006954in}{2.014902in}}%
\pgfpathlineto{\pgfqpoint{2.049145in}{2.014243in}}%
\pgfpathlineto{\pgfqpoint{2.048844in}{1.986478in}}%
\pgfpathlineto{\pgfqpoint{2.041913in}{1.986517in}}%
\pgfpathlineto{\pgfqpoint{2.042858in}{1.972645in}}%
\pgfpathlineto{\pgfqpoint{2.042559in}{1.951936in}}%
\pgfpathclose%
\pgfusepath{fill}%
\end{pgfscope}%
\begin{pgfscope}%
\pgfpathrectangle{\pgfqpoint{0.100000in}{0.100000in}}{\pgfqpoint{3.608454in}{2.310000in}}%
\pgfusepath{clip}%
\pgfsetbuttcap%
\pgfsetmiterjoin%
\definecolor{currentfill}{rgb}{0.000000,0.854902,0.572549}%
\pgfsetfillcolor{currentfill}%
\pgfsetlinewidth{0.000000pt}%
\definecolor{currentstroke}{rgb}{0.000000,0.000000,0.000000}%
\pgfsetstrokecolor{currentstroke}%
\pgfsetstrokeopacity{0.000000}%
\pgfsetdash{}{0pt}%
\pgfpathmoveto{\pgfqpoint{2.738013in}{1.786827in}}%
\pgfpathlineto{\pgfqpoint{2.697115in}{1.781860in}}%
\pgfpathlineto{\pgfqpoint{2.683319in}{1.780357in}}%
\pgfpathlineto{\pgfqpoint{2.680280in}{1.807943in}}%
\pgfpathlineto{\pgfqpoint{2.677389in}{1.835395in}}%
\pgfpathlineto{\pgfqpoint{2.674595in}{1.855705in}}%
\pgfpathlineto{\pgfqpoint{2.681411in}{1.856342in}}%
\pgfpathlineto{\pgfqpoint{2.680662in}{1.863242in}}%
\pgfpathlineto{\pgfqpoint{2.707810in}{1.866142in}}%
\pgfpathlineto{\pgfqpoint{2.728097in}{1.868854in}}%
\pgfpathlineto{\pgfqpoint{2.729298in}{1.855252in}}%
\pgfpathlineto{\pgfqpoint{2.738013in}{1.786827in}}%
\pgfpathclose%
\pgfusepath{fill}%
\end{pgfscope}%
\begin{pgfscope}%
\pgfpathrectangle{\pgfqpoint{0.100000in}{0.100000in}}{\pgfqpoint{3.608454in}{2.310000in}}%
\pgfusepath{clip}%
\pgfsetbuttcap%
\pgfsetmiterjoin%
\definecolor{currentfill}{rgb}{0.000000,0.815686,0.592157}%
\pgfsetfillcolor{currentfill}%
\pgfsetlinewidth{0.000000pt}%
\definecolor{currentstroke}{rgb}{0.000000,0.000000,0.000000}%
\pgfsetstrokecolor{currentstroke}%
\pgfsetstrokeopacity{0.000000}%
\pgfsetdash{}{0pt}%
\pgfpathmoveto{\pgfqpoint{3.069493in}{0.458265in}}%
\pgfpathlineto{\pgfqpoint{3.072400in}{0.437158in}}%
\pgfpathlineto{\pgfqpoint{3.037010in}{0.432074in}}%
\pgfpathlineto{\pgfqpoint{3.036461in}{0.440065in}}%
\pgfpathlineto{\pgfqpoint{3.030688in}{0.445529in}}%
\pgfpathlineto{\pgfqpoint{3.027406in}{0.441545in}}%
\pgfpathlineto{\pgfqpoint{3.030649in}{0.432623in}}%
\pgfpathlineto{\pgfqpoint{3.024312in}{0.431402in}}%
\pgfpathlineto{\pgfqpoint{3.021799in}{0.431432in}}%
\pgfpathlineto{\pgfqpoint{3.006353in}{0.451398in}}%
\pgfpathlineto{\pgfqpoint{2.995081in}{0.469136in}}%
\pgfpathlineto{\pgfqpoint{2.984705in}{0.479666in}}%
\pgfpathlineto{\pgfqpoint{2.986857in}{0.487228in}}%
\pgfpathlineto{\pgfqpoint{2.992126in}{0.496993in}}%
\pgfpathlineto{\pgfqpoint{3.027503in}{0.502091in}}%
\pgfpathlineto{\pgfqpoint{3.030929in}{0.477492in}}%
\pgfpathlineto{\pgfqpoint{3.065776in}{0.482786in}}%
\pgfpathlineto{\pgfqpoint{3.069493in}{0.458265in}}%
\pgfpathclose%
\pgfusepath{fill}%
\end{pgfscope}%
\begin{pgfscope}%
\pgfpathrectangle{\pgfqpoint{0.100000in}{0.100000in}}{\pgfqpoint{3.608454in}{2.310000in}}%
\pgfusepath{clip}%
\pgfsetbuttcap%
\pgfsetmiterjoin%
\definecolor{currentfill}{rgb}{0.000000,0.662745,0.668627}%
\pgfsetfillcolor{currentfill}%
\pgfsetlinewidth{0.000000pt}%
\definecolor{currentstroke}{rgb}{0.000000,0.000000,0.000000}%
\pgfsetstrokecolor{currentstroke}%
\pgfsetstrokeopacity{0.000000}%
\pgfsetdash{}{0pt}%
\pgfpathmoveto{\pgfqpoint{3.046524in}{1.607153in}}%
\pgfpathlineto{\pgfqpoint{3.045476in}{1.613591in}}%
\pgfpathlineto{\pgfqpoint{3.051910in}{1.617021in}}%
\pgfpathlineto{\pgfqpoint{3.050721in}{1.630639in}}%
\pgfpathlineto{\pgfqpoint{3.078140in}{1.636089in}}%
\pgfpathlineto{\pgfqpoint{3.094591in}{1.638194in}}%
\pgfpathlineto{\pgfqpoint{3.105576in}{1.628770in}}%
\pgfpathlineto{\pgfqpoint{3.109153in}{1.629395in}}%
\pgfpathlineto{\pgfqpoint{3.110757in}{1.621055in}}%
\pgfpathlineto{\pgfqpoint{3.106858in}{1.608202in}}%
\pgfpathlineto{\pgfqpoint{3.068655in}{1.604850in}}%
\pgfpathlineto{\pgfqpoint{3.054170in}{1.610224in}}%
\pgfpathlineto{\pgfqpoint{3.046524in}{1.607153in}}%
\pgfpathclose%
\pgfusepath{fill}%
\end{pgfscope}%
\begin{pgfscope}%
\pgfpathrectangle{\pgfqpoint{0.100000in}{0.100000in}}{\pgfqpoint{3.608454in}{2.310000in}}%
\pgfusepath{clip}%
\pgfsetbuttcap%
\pgfsetmiterjoin%
\definecolor{currentfill}{rgb}{0.000000,0.603922,0.698039}%
\pgfsetfillcolor{currentfill}%
\pgfsetlinewidth{0.000000pt}%
\definecolor{currentstroke}{rgb}{0.000000,0.000000,0.000000}%
\pgfsetstrokecolor{currentstroke}%
\pgfsetstrokeopacity{0.000000}%
\pgfsetdash{}{0pt}%
\pgfpathmoveto{\pgfqpoint{2.583434in}{1.352723in}}%
\pgfpathlineto{\pgfqpoint{2.566176in}{1.350839in}}%
\pgfpathlineto{\pgfqpoint{2.562407in}{1.357174in}}%
\pgfpathlineto{\pgfqpoint{2.562133in}{1.362791in}}%
\pgfpathlineto{\pgfqpoint{2.558058in}{1.370400in}}%
\pgfpathlineto{\pgfqpoint{2.560407in}{1.373999in}}%
\pgfpathlineto{\pgfqpoint{2.558464in}{1.382618in}}%
\pgfpathlineto{\pgfqpoint{2.562752in}{1.385889in}}%
\pgfpathlineto{\pgfqpoint{2.559183in}{1.428045in}}%
\pgfpathlineto{\pgfqpoint{2.557572in}{1.448924in}}%
\pgfpathlineto{\pgfqpoint{2.565221in}{1.447932in}}%
\pgfpathlineto{\pgfqpoint{2.565557in}{1.434073in}}%
\pgfpathlineto{\pgfqpoint{2.585430in}{1.435815in}}%
\pgfpathlineto{\pgfqpoint{2.583559in}{1.456426in}}%
\pgfpathlineto{\pgfqpoint{2.607399in}{1.458524in}}%
\pgfpathlineto{\pgfqpoint{2.607655in}{1.455712in}}%
\pgfpathlineto{\pgfqpoint{2.612117in}{1.412594in}}%
\pgfpathlineto{\pgfqpoint{2.614715in}{1.406613in}}%
\pgfpathlineto{\pgfqpoint{2.613308in}{1.399991in}}%
\pgfpathlineto{\pgfqpoint{2.597905in}{1.398852in}}%
\pgfpathlineto{\pgfqpoint{2.598672in}{1.388501in}}%
\pgfpathlineto{\pgfqpoint{2.591897in}{1.387949in}}%
\pgfpathlineto{\pgfqpoint{2.593047in}{1.374201in}}%
\pgfpathlineto{\pgfqpoint{2.581657in}{1.373507in}}%
\pgfpathlineto{\pgfqpoint{2.583434in}{1.352723in}}%
\pgfpathclose%
\pgfusepath{fill}%
\end{pgfscope}%
\begin{pgfscope}%
\pgfpathrectangle{\pgfqpoint{0.100000in}{0.100000in}}{\pgfqpoint{3.608454in}{2.310000in}}%
\pgfusepath{clip}%
\pgfsetbuttcap%
\pgfsetmiterjoin%
\definecolor{currentfill}{rgb}{0.000000,0.341176,0.829412}%
\pgfsetfillcolor{currentfill}%
\pgfsetlinewidth{0.000000pt}%
\definecolor{currentstroke}{rgb}{0.000000,0.000000,0.000000}%
\pgfsetstrokecolor{currentstroke}%
\pgfsetstrokeopacity{0.000000}%
\pgfsetdash{}{0pt}%
\pgfpathmoveto{\pgfqpoint{1.580064in}{1.651542in}}%
\pgfpathlineto{\pgfqpoint{1.571835in}{1.569055in}}%
\pgfpathlineto{\pgfqpoint{1.536486in}{1.572765in}}%
\pgfpathlineto{\pgfqpoint{1.537205in}{1.579689in}}%
\pgfpathlineto{\pgfqpoint{1.500517in}{1.584019in}}%
\pgfpathlineto{\pgfqpoint{1.503839in}{1.609878in}}%
\pgfpathlineto{\pgfqpoint{1.507108in}{1.644984in}}%
\pgfpathlineto{\pgfqpoint{1.508517in}{1.658724in}}%
\pgfpathlineto{\pgfqpoint{1.531372in}{1.656425in}}%
\pgfpathlineto{\pgfqpoint{1.580064in}{1.651542in}}%
\pgfpathclose%
\pgfusepath{fill}%
\end{pgfscope}%
\begin{pgfscope}%
\pgfpathrectangle{\pgfqpoint{0.100000in}{0.100000in}}{\pgfqpoint{3.608454in}{2.310000in}}%
\pgfusepath{clip}%
\pgfsetbuttcap%
\pgfsetmiterjoin%
\definecolor{currentfill}{rgb}{0.000000,0.615686,0.692157}%
\pgfsetfillcolor{currentfill}%
\pgfsetlinewidth{0.000000pt}%
\definecolor{currentstroke}{rgb}{0.000000,0.000000,0.000000}%
\pgfsetstrokecolor{currentstroke}%
\pgfsetstrokeopacity{0.000000}%
\pgfsetdash{}{0pt}%
\pgfpathmoveto{\pgfqpoint{1.651469in}{1.020799in}}%
\pgfpathlineto{\pgfqpoint{1.685686in}{1.018201in}}%
\pgfpathlineto{\pgfqpoint{1.683230in}{0.983785in}}%
\pgfpathlineto{\pgfqpoint{1.711338in}{0.981859in}}%
\pgfpathlineto{\pgfqpoint{1.708835in}{0.943953in}}%
\pgfpathlineto{\pgfqpoint{1.640100in}{0.948213in}}%
\pgfpathlineto{\pgfqpoint{1.642812in}{0.986873in}}%
\pgfpathlineto{\pgfqpoint{1.648845in}{0.986397in}}%
\pgfpathlineto{\pgfqpoint{1.651469in}{1.020799in}}%
\pgfpathclose%
\pgfusepath{fill}%
\end{pgfscope}%
\begin{pgfscope}%
\pgfpathrectangle{\pgfqpoint{0.100000in}{0.100000in}}{\pgfqpoint{3.608454in}{2.310000in}}%
\pgfusepath{clip}%
\pgfsetbuttcap%
\pgfsetmiterjoin%
\definecolor{currentfill}{rgb}{0.000000,0.156863,0.921569}%
\pgfsetfillcolor{currentfill}%
\pgfsetlinewidth{0.000000pt}%
\definecolor{currentstroke}{rgb}{0.000000,0.000000,0.000000}%
\pgfsetstrokecolor{currentstroke}%
\pgfsetstrokeopacity{0.000000}%
\pgfsetdash{}{0pt}%
\pgfpathmoveto{\pgfqpoint{1.396241in}{1.384895in}}%
\pgfpathlineto{\pgfqpoint{1.393282in}{1.388976in}}%
\pgfpathlineto{\pgfqpoint{1.387549in}{1.389567in}}%
\pgfpathlineto{\pgfqpoint{1.383841in}{1.385038in}}%
\pgfpathlineto{\pgfqpoint{1.377365in}{1.386918in}}%
\pgfpathlineto{\pgfqpoint{1.369057in}{1.398085in}}%
\pgfpathlineto{\pgfqpoint{1.355918in}{1.399849in}}%
\pgfpathlineto{\pgfqpoint{1.352191in}{1.404691in}}%
\pgfpathlineto{\pgfqpoint{1.350593in}{1.411552in}}%
\pgfpathlineto{\pgfqpoint{1.346932in}{1.416687in}}%
\pgfpathlineto{\pgfqpoint{1.349566in}{1.420460in}}%
\pgfpathlineto{\pgfqpoint{1.290367in}{1.428918in}}%
\pgfpathlineto{\pgfqpoint{1.251645in}{1.434868in}}%
\pgfpathlineto{\pgfqpoint{1.253203in}{1.445105in}}%
\pgfpathlineto{\pgfqpoint{1.255196in}{1.457818in}}%
\pgfpathlineto{\pgfqpoint{1.284941in}{1.452441in}}%
\pgfpathlineto{\pgfqpoint{1.285446in}{1.455840in}}%
\pgfpathlineto{\pgfqpoint{1.322533in}{1.450461in}}%
\pgfpathlineto{\pgfqpoint{1.324000in}{1.460674in}}%
\pgfpathlineto{\pgfqpoint{1.354499in}{1.456659in}}%
\pgfpathlineto{\pgfqpoint{1.355468in}{1.463655in}}%
\pgfpathlineto{\pgfqpoint{1.362220in}{1.462510in}}%
\pgfpathlineto{\pgfqpoint{1.364204in}{1.476264in}}%
\pgfpathlineto{\pgfqpoint{1.380971in}{1.474066in}}%
\pgfpathlineto{\pgfqpoint{1.379506in}{1.460495in}}%
\pgfpathlineto{\pgfqpoint{1.403922in}{1.457766in}}%
\pgfpathlineto{\pgfqpoint{1.426636in}{1.453930in}}%
\pgfpathlineto{\pgfqpoint{1.435412in}{1.444400in}}%
\pgfpathlineto{\pgfqpoint{1.437752in}{1.434786in}}%
\pgfpathlineto{\pgfqpoint{1.444005in}{1.434780in}}%
\pgfpathlineto{\pgfqpoint{1.451897in}{1.428337in}}%
\pgfpathlineto{\pgfqpoint{1.439656in}{1.412644in}}%
\pgfpathlineto{\pgfqpoint{1.424568in}{1.406130in}}%
\pgfpathlineto{\pgfqpoint{1.423585in}{1.392666in}}%
\pgfpathlineto{\pgfqpoint{1.421683in}{1.386315in}}%
\pgfpathlineto{\pgfqpoint{1.398138in}{1.389426in}}%
\pgfpathlineto{\pgfqpoint{1.396241in}{1.384895in}}%
\pgfpathclose%
\pgfusepath{fill}%
\end{pgfscope}%
\begin{pgfscope}%
\pgfpathrectangle{\pgfqpoint{0.100000in}{0.100000in}}{\pgfqpoint{3.608454in}{2.310000in}}%
\pgfusepath{clip}%
\pgfsetbuttcap%
\pgfsetmiterjoin%
\definecolor{currentfill}{rgb}{0.000000,0.603922,0.698039}%
\pgfsetfillcolor{currentfill}%
\pgfsetlinewidth{0.000000pt}%
\definecolor{currentstroke}{rgb}{0.000000,0.000000,0.000000}%
\pgfsetstrokecolor{currentstroke}%
\pgfsetstrokeopacity{0.000000}%
\pgfsetdash{}{0pt}%
\pgfpathmoveto{\pgfqpoint{3.091930in}{0.511405in}}%
\pgfpathlineto{\pgfqpoint{3.088462in}{0.501164in}}%
\pgfpathlineto{\pgfqpoint{3.097174in}{0.491133in}}%
\pgfpathlineto{\pgfqpoint{3.100740in}{0.491314in}}%
\pgfpathlineto{\pgfqpoint{3.109319in}{0.479362in}}%
\pgfpathlineto{\pgfqpoint{3.095343in}{0.476565in}}%
\pgfpathlineto{\pgfqpoint{3.096465in}{0.469557in}}%
\pgfpathlineto{\pgfqpoint{3.089414in}{0.468427in}}%
\pgfpathlineto{\pgfqpoint{3.090529in}{0.461360in}}%
\pgfpathlineto{\pgfqpoint{3.069493in}{0.458265in}}%
\pgfpathlineto{\pgfqpoint{3.065776in}{0.482786in}}%
\pgfpathlineto{\pgfqpoint{3.030929in}{0.477492in}}%
\pgfpathlineto{\pgfqpoint{3.027503in}{0.502091in}}%
\pgfpathlineto{\pgfqpoint{3.021367in}{0.543921in}}%
\pgfpathlineto{\pgfqpoint{3.016866in}{0.550384in}}%
\pgfpathlineto{\pgfqpoint{3.019784in}{0.555093in}}%
\pgfpathlineto{\pgfqpoint{3.026291in}{0.558732in}}%
\pgfpathlineto{\pgfqpoint{3.047294in}{0.561961in}}%
\pgfpathlineto{\pgfqpoint{3.048333in}{0.554969in}}%
\pgfpathlineto{\pgfqpoint{3.055318in}{0.556018in}}%
\pgfpathlineto{\pgfqpoint{3.059031in}{0.547094in}}%
\pgfpathlineto{\pgfqpoint{3.070237in}{0.538537in}}%
\pgfpathlineto{\pgfqpoint{3.076539in}{0.531713in}}%
\pgfpathlineto{\pgfqpoint{3.077818in}{0.527142in}}%
\pgfpathlineto{\pgfqpoint{3.084501in}{0.525466in}}%
\pgfpathlineto{\pgfqpoint{3.090421in}{0.517236in}}%
\pgfpathlineto{\pgfqpoint{3.091930in}{0.511405in}}%
\pgfpathclose%
\pgfusepath{fill}%
\end{pgfscope}%
\begin{pgfscope}%
\pgfpathrectangle{\pgfqpoint{0.100000in}{0.100000in}}{\pgfqpoint{3.608454in}{2.310000in}}%
\pgfusepath{clip}%
\pgfsetbuttcap%
\pgfsetmiterjoin%
\definecolor{currentfill}{rgb}{0.000000,0.835294,0.582353}%
\pgfsetfillcolor{currentfill}%
\pgfsetlinewidth{0.000000pt}%
\definecolor{currentstroke}{rgb}{0.000000,0.000000,0.000000}%
\pgfsetstrokecolor{currentstroke}%
\pgfsetstrokeopacity{0.000000}%
\pgfsetdash{}{0pt}%
\pgfpathmoveto{\pgfqpoint{0.509532in}{1.308555in}}%
\pgfpathlineto{\pgfqpoint{0.550101in}{1.297200in}}%
\pgfpathlineto{\pgfqpoint{0.606581in}{1.281871in}}%
\pgfpathlineto{\pgfqpoint{0.645007in}{1.271759in}}%
\pgfpathlineto{\pgfqpoint{0.668493in}{1.266395in}}%
\pgfpathlineto{\pgfqpoint{0.649520in}{1.191532in}}%
\pgfpathlineto{\pgfqpoint{0.598059in}{1.204495in}}%
\pgfpathlineto{\pgfqpoint{0.547967in}{1.217656in}}%
\pgfpathlineto{\pgfqpoint{0.547247in}{1.223234in}}%
\pgfpathlineto{\pgfqpoint{0.537333in}{1.227732in}}%
\pgfpathlineto{\pgfqpoint{0.539132in}{1.241671in}}%
\pgfpathlineto{\pgfqpoint{0.533846in}{1.244020in}}%
\pgfpathlineto{\pgfqpoint{0.536255in}{1.250941in}}%
\pgfpathlineto{\pgfqpoint{0.529070in}{1.252510in}}%
\pgfpathlineto{\pgfqpoint{0.530958in}{1.259199in}}%
\pgfpathlineto{\pgfqpoint{0.522085in}{1.261761in}}%
\pgfpathlineto{\pgfqpoint{0.523955in}{1.268417in}}%
\pgfpathlineto{\pgfqpoint{0.519557in}{1.269679in}}%
\pgfpathlineto{\pgfqpoint{0.521448in}{1.276381in}}%
\pgfpathlineto{\pgfqpoint{0.514797in}{1.280686in}}%
\pgfpathlineto{\pgfqpoint{0.510504in}{1.286661in}}%
\pgfpathlineto{\pgfqpoint{0.512448in}{1.293350in}}%
\pgfpathlineto{\pgfqpoint{0.505758in}{1.295236in}}%
\pgfpathlineto{\pgfqpoint{0.509532in}{1.308555in}}%
\pgfpathclose%
\pgfusepath{fill}%
\end{pgfscope}%
\begin{pgfscope}%
\pgfpathrectangle{\pgfqpoint{0.100000in}{0.100000in}}{\pgfqpoint{3.608454in}{2.310000in}}%
\pgfusepath{clip}%
\pgfsetbuttcap%
\pgfsetmiterjoin%
\definecolor{currentfill}{rgb}{0.000000,0.427451,0.786275}%
\pgfsetfillcolor{currentfill}%
\pgfsetlinewidth{0.000000pt}%
\definecolor{currentstroke}{rgb}{0.000000,0.000000,0.000000}%
\pgfsetstrokecolor{currentstroke}%
\pgfsetstrokeopacity{0.000000}%
\pgfsetdash{}{0pt}%
\pgfpathmoveto{\pgfqpoint{2.193447in}{1.465250in}}%
\pgfpathlineto{\pgfqpoint{2.180394in}{1.464846in}}%
\pgfpathlineto{\pgfqpoint{2.138550in}{1.463912in}}%
\pgfpathlineto{\pgfqpoint{2.138316in}{1.510307in}}%
\pgfpathlineto{\pgfqpoint{2.172375in}{1.510519in}}%
\pgfpathlineto{\pgfqpoint{2.192689in}{1.511198in}}%
\pgfpathlineto{\pgfqpoint{2.193447in}{1.465250in}}%
\pgfpathclose%
\pgfusepath{fill}%
\end{pgfscope}%
\begin{pgfscope}%
\pgfpathrectangle{\pgfqpoint{0.100000in}{0.100000in}}{\pgfqpoint{3.608454in}{2.310000in}}%
\pgfusepath{clip}%
\pgfsetbuttcap%
\pgfsetmiterjoin%
\definecolor{currentfill}{rgb}{0.000000,0.709804,0.645098}%
\pgfsetfillcolor{currentfill}%
\pgfsetlinewidth{0.000000pt}%
\definecolor{currentstroke}{rgb}{0.000000,0.000000,0.000000}%
\pgfsetstrokecolor{currentstroke}%
\pgfsetstrokeopacity{0.000000}%
\pgfsetdash{}{0pt}%
\pgfpathmoveto{\pgfqpoint{2.724860in}{1.203323in}}%
\pgfpathlineto{\pgfqpoint{2.723224in}{1.208766in}}%
\pgfpathlineto{\pgfqpoint{2.709289in}{1.208321in}}%
\pgfpathlineto{\pgfqpoint{2.705489in}{1.213403in}}%
\pgfpathlineto{\pgfqpoint{2.703168in}{1.221287in}}%
\pgfpathlineto{\pgfqpoint{2.696385in}{1.221833in}}%
\pgfpathlineto{\pgfqpoint{2.692552in}{1.225974in}}%
\pgfpathlineto{\pgfqpoint{2.690908in}{1.235308in}}%
\pgfpathlineto{\pgfqpoint{2.692446in}{1.245040in}}%
\pgfpathlineto{\pgfqpoint{2.693473in}{1.249077in}}%
\pgfpathlineto{\pgfqpoint{2.704025in}{1.249697in}}%
\pgfpathlineto{\pgfqpoint{2.711292in}{1.250974in}}%
\pgfpathlineto{\pgfqpoint{2.718150in}{1.247622in}}%
\pgfpathlineto{\pgfqpoint{2.722412in}{1.251628in}}%
\pgfpathlineto{\pgfqpoint{2.730780in}{1.248329in}}%
\pgfpathlineto{\pgfqpoint{2.729948in}{1.258824in}}%
\pgfpathlineto{\pgfqpoint{2.741472in}{1.258776in}}%
\pgfpathlineto{\pgfqpoint{2.742755in}{1.250848in}}%
\pgfpathlineto{\pgfqpoint{2.751350in}{1.246500in}}%
\pgfpathlineto{\pgfqpoint{2.752556in}{1.236924in}}%
\pgfpathlineto{\pgfqpoint{2.742337in}{1.226077in}}%
\pgfpathlineto{\pgfqpoint{2.742647in}{1.220616in}}%
\pgfpathlineto{\pgfqpoint{2.747524in}{1.217198in}}%
\pgfpathlineto{\pgfqpoint{2.731198in}{1.201857in}}%
\pgfpathlineto{\pgfqpoint{2.724860in}{1.203323in}}%
\pgfpathclose%
\pgfusepath{fill}%
\end{pgfscope}%
\begin{pgfscope}%
\pgfpathrectangle{\pgfqpoint{0.100000in}{0.100000in}}{\pgfqpoint{3.608454in}{2.310000in}}%
\pgfusepath{clip}%
\pgfsetbuttcap%
\pgfsetmiterjoin%
\definecolor{currentfill}{rgb}{0.000000,0.505882,0.747059}%
\pgfsetfillcolor{currentfill}%
\pgfsetlinewidth{0.000000pt}%
\definecolor{currentstroke}{rgb}{0.000000,0.000000,0.000000}%
\pgfsetstrokecolor{currentstroke}%
\pgfsetstrokeopacity{0.000000}%
\pgfsetdash{}{0pt}%
\pgfpathmoveto{\pgfqpoint{1.711338in}{0.981859in}}%
\pgfpathlineto{\pgfqpoint{1.683230in}{0.983785in}}%
\pgfpathlineto{\pgfqpoint{1.685686in}{1.018201in}}%
\pgfpathlineto{\pgfqpoint{1.710432in}{1.016551in}}%
\pgfpathlineto{\pgfqpoint{1.713023in}{1.050837in}}%
\pgfpathlineto{\pgfqpoint{1.748391in}{1.048651in}}%
\pgfpathlineto{\pgfqpoint{1.746176in}{1.014219in}}%
\pgfpathlineto{\pgfqpoint{1.754289in}{1.013713in}}%
\pgfpathlineto{\pgfqpoint{1.752117in}{0.979371in}}%
\pgfpathlineto{\pgfqpoint{1.745610in}{0.979809in}}%
\pgfpathlineto{\pgfqpoint{1.711338in}{0.981859in}}%
\pgfpathclose%
\pgfusepath{fill}%
\end{pgfscope}%
\begin{pgfscope}%
\pgfpathrectangle{\pgfqpoint{0.100000in}{0.100000in}}{\pgfqpoint{3.608454in}{2.310000in}}%
\pgfusepath{clip}%
\pgfsetbuttcap%
\pgfsetmiterjoin%
\definecolor{currentfill}{rgb}{0.000000,0.650980,0.674510}%
\pgfsetfillcolor{currentfill}%
\pgfsetlinewidth{0.000000pt}%
\definecolor{currentstroke}{rgb}{0.000000,0.000000,0.000000}%
\pgfsetstrokecolor{currentstroke}%
\pgfsetstrokeopacity{0.000000}%
\pgfsetdash{}{0pt}%
\pgfpathmoveto{\pgfqpoint{1.666412in}{1.249548in}}%
\pgfpathlineto{\pgfqpoint{1.664847in}{1.229439in}}%
\pgfpathlineto{\pgfqpoint{1.662385in}{1.198210in}}%
\pgfpathlineto{\pgfqpoint{1.621102in}{1.201730in}}%
\pgfpathlineto{\pgfqpoint{1.602002in}{1.203777in}}%
\pgfpathlineto{\pgfqpoint{1.596724in}{1.204258in}}%
\pgfpathlineto{\pgfqpoint{1.601877in}{1.254860in}}%
\pgfpathlineto{\pgfqpoint{1.666412in}{1.249548in}}%
\pgfpathclose%
\pgfusepath{fill}%
\end{pgfscope}%
\begin{pgfscope}%
\pgfpathrectangle{\pgfqpoint{0.100000in}{0.100000in}}{\pgfqpoint{3.608454in}{2.310000in}}%
\pgfusepath{clip}%
\pgfsetbuttcap%
\pgfsetmiterjoin%
\definecolor{currentfill}{rgb}{0.000000,0.309804,0.845098}%
\pgfsetfillcolor{currentfill}%
\pgfsetlinewidth{0.000000pt}%
\definecolor{currentstroke}{rgb}{0.000000,0.000000,0.000000}%
\pgfsetstrokecolor{currentstroke}%
\pgfsetstrokeopacity{0.000000}%
\pgfsetdash{}{0pt}%
\pgfpathmoveto{\pgfqpoint{2.685775in}{1.104241in}}%
\pgfpathlineto{\pgfqpoint{2.681536in}{1.102539in}}%
\pgfpathlineto{\pgfqpoint{2.668887in}{1.104876in}}%
\pgfpathlineto{\pgfqpoint{2.654621in}{1.100124in}}%
\pgfpathlineto{\pgfqpoint{2.651600in}{1.103191in}}%
\pgfpathlineto{\pgfqpoint{2.643683in}{1.102257in}}%
\pgfpathlineto{\pgfqpoint{2.632392in}{1.103973in}}%
\pgfpathlineto{\pgfqpoint{2.605693in}{1.111968in}}%
\pgfpathlineto{\pgfqpoint{2.605651in}{1.120589in}}%
\pgfpathlineto{\pgfqpoint{2.606773in}{1.136211in}}%
\pgfpathlineto{\pgfqpoint{2.604989in}{1.140525in}}%
\pgfpathlineto{\pgfqpoint{2.606315in}{1.148332in}}%
\pgfpathlineto{\pgfqpoint{2.597965in}{1.148745in}}%
\pgfpathlineto{\pgfqpoint{2.599908in}{1.153357in}}%
\pgfpathlineto{\pgfqpoint{2.607032in}{1.160645in}}%
\pgfpathlineto{\pgfqpoint{2.606663in}{1.174945in}}%
\pgfpathlineto{\pgfqpoint{2.610086in}{1.175304in}}%
\pgfpathlineto{\pgfqpoint{2.639785in}{1.178722in}}%
\pgfpathlineto{\pgfqpoint{2.656122in}{1.179965in}}%
\pgfpathlineto{\pgfqpoint{2.678723in}{1.180579in}}%
\pgfpathlineto{\pgfqpoint{2.690217in}{1.180528in}}%
\pgfpathlineto{\pgfqpoint{2.690027in}{1.171438in}}%
\pgfpathlineto{\pgfqpoint{2.689646in}{1.154594in}}%
\pgfpathlineto{\pgfqpoint{2.693285in}{1.148452in}}%
\pgfpathlineto{\pgfqpoint{2.693376in}{1.142517in}}%
\pgfpathlineto{\pgfqpoint{2.690045in}{1.139163in}}%
\pgfpathlineto{\pgfqpoint{2.683825in}{1.139915in}}%
\pgfpathlineto{\pgfqpoint{2.678100in}{1.135237in}}%
\pgfpathlineto{\pgfqpoint{2.681516in}{1.127762in}}%
\pgfpathlineto{\pgfqpoint{2.689025in}{1.121466in}}%
\pgfpathlineto{\pgfqpoint{2.685202in}{1.110108in}}%
\pgfpathlineto{\pgfqpoint{2.685775in}{1.104241in}}%
\pgfpathclose%
\pgfusepath{fill}%
\end{pgfscope}%
\begin{pgfscope}%
\pgfpathrectangle{\pgfqpoint{0.100000in}{0.100000in}}{\pgfqpoint{3.608454in}{2.310000in}}%
\pgfusepath{clip}%
\pgfsetbuttcap%
\pgfsetmiterjoin%
\definecolor{currentfill}{rgb}{0.000000,0.752941,0.623529}%
\pgfsetfillcolor{currentfill}%
\pgfsetlinewidth{0.000000pt}%
\definecolor{currentstroke}{rgb}{0.000000,0.000000,0.000000}%
\pgfsetstrokecolor{currentstroke}%
\pgfsetstrokeopacity{0.000000}%
\pgfsetdash{}{0pt}%
\pgfpathmoveto{\pgfqpoint{0.486008in}{1.393700in}}%
\pgfpathlineto{\pgfqpoint{0.474622in}{1.404992in}}%
\pgfpathlineto{\pgfqpoint{0.471004in}{1.413448in}}%
\pgfpathlineto{\pgfqpoint{0.475697in}{1.423778in}}%
\pgfpathlineto{\pgfqpoint{0.476039in}{1.431016in}}%
\pgfpathlineto{\pgfqpoint{0.465842in}{1.434190in}}%
\pgfpathlineto{\pgfqpoint{0.465818in}{1.444533in}}%
\pgfpathlineto{\pgfqpoint{0.470704in}{1.451769in}}%
\pgfpathlineto{\pgfqpoint{0.469193in}{1.459876in}}%
\pgfpathlineto{\pgfqpoint{0.469229in}{1.459834in}}%
\pgfpathlineto{\pgfqpoint{0.489989in}{1.471204in}}%
\pgfpathlineto{\pgfqpoint{0.504262in}{1.472446in}}%
\pgfpathlineto{\pgfqpoint{0.508101in}{1.469498in}}%
\pgfpathlineto{\pgfqpoint{0.515295in}{1.495292in}}%
\pgfpathlineto{\pgfqpoint{0.526181in}{1.471930in}}%
\pgfpathlineto{\pgfqpoint{0.537164in}{1.478720in}}%
\pgfpathlineto{\pgfqpoint{0.550096in}{1.491227in}}%
\pgfpathlineto{\pgfqpoint{0.566537in}{1.504878in}}%
\pgfpathlineto{\pgfqpoint{0.577147in}{1.506768in}}%
\pgfpathlineto{\pgfqpoint{0.583557in}{1.498623in}}%
\pgfpathlineto{\pgfqpoint{0.592836in}{1.501027in}}%
\pgfpathlineto{\pgfqpoint{0.597596in}{1.492295in}}%
\pgfpathlineto{\pgfqpoint{0.596004in}{1.489421in}}%
\pgfpathlineto{\pgfqpoint{0.597769in}{1.478364in}}%
\pgfpathlineto{\pgfqpoint{0.610025in}{1.468870in}}%
\pgfpathlineto{\pgfqpoint{0.609428in}{1.457861in}}%
\pgfpathlineto{\pgfqpoint{0.613111in}{1.453862in}}%
\pgfpathlineto{\pgfqpoint{0.612588in}{1.444548in}}%
\pgfpathlineto{\pgfqpoint{0.607467in}{1.441382in}}%
\pgfpathlineto{\pgfqpoint{0.605864in}{1.444998in}}%
\pgfpathlineto{\pgfqpoint{0.584724in}{1.432975in}}%
\pgfpathlineto{\pgfqpoint{0.583298in}{1.427984in}}%
\pgfpathlineto{\pgfqpoint{0.577622in}{1.423252in}}%
\pgfpathlineto{\pgfqpoint{0.570915in}{1.425063in}}%
\pgfpathlineto{\pgfqpoint{0.548331in}{1.412190in}}%
\pgfpathlineto{\pgfqpoint{0.533257in}{1.413969in}}%
\pgfpathlineto{\pgfqpoint{0.520707in}{1.412942in}}%
\pgfpathlineto{\pgfqpoint{0.486008in}{1.393700in}}%
\pgfpathclose%
\pgfusepath{fill}%
\end{pgfscope}%
\begin{pgfscope}%
\pgfpathrectangle{\pgfqpoint{0.100000in}{0.100000in}}{\pgfqpoint{3.608454in}{2.310000in}}%
\pgfusepath{clip}%
\pgfsetbuttcap%
\pgfsetmiterjoin%
\definecolor{currentfill}{rgb}{0.000000,0.431373,0.784314}%
\pgfsetfillcolor{currentfill}%
\pgfsetlinewidth{0.000000pt}%
\definecolor{currentstroke}{rgb}{0.000000,0.000000,0.000000}%
\pgfsetstrokecolor{currentstroke}%
\pgfsetstrokeopacity{0.000000}%
\pgfsetdash{}{0pt}%
\pgfpathmoveto{\pgfqpoint{1.664670in}{1.909384in}}%
\pgfpathlineto{\pgfqpoint{1.606786in}{1.914838in}}%
\pgfpathlineto{\pgfqpoint{1.611520in}{1.962062in}}%
\pgfpathlineto{\pgfqpoint{1.624849in}{1.960734in}}%
\pgfpathlineto{\pgfqpoint{1.625532in}{1.967784in}}%
\pgfpathlineto{\pgfqpoint{1.635984in}{1.966799in}}%
\pgfpathlineto{\pgfqpoint{1.673140in}{1.963430in}}%
\pgfpathlineto{\pgfqpoint{1.670909in}{1.935733in}}%
\pgfpathlineto{\pgfqpoint{1.667012in}{1.936010in}}%
\pgfpathlineto{\pgfqpoint{1.664670in}{1.909384in}}%
\pgfpathclose%
\pgfusepath{fill}%
\end{pgfscope}%
\begin{pgfscope}%
\pgfpathrectangle{\pgfqpoint{0.100000in}{0.100000in}}{\pgfqpoint{3.608454in}{2.310000in}}%
\pgfusepath{clip}%
\pgfsetbuttcap%
\pgfsetmiterjoin%
\definecolor{currentfill}{rgb}{0.000000,0.521569,0.739216}%
\pgfsetfillcolor{currentfill}%
\pgfsetlinewidth{0.000000pt}%
\definecolor{currentstroke}{rgb}{0.000000,0.000000,0.000000}%
\pgfsetstrokecolor{currentstroke}%
\pgfsetstrokeopacity{0.000000}%
\pgfsetdash{}{0pt}%
\pgfpathmoveto{\pgfqpoint{3.282347in}{1.403895in}}%
\pgfpathlineto{\pgfqpoint{3.278306in}{1.404567in}}%
\pgfpathlineto{\pgfqpoint{3.267145in}{1.400703in}}%
\pgfpathlineto{\pgfqpoint{3.253828in}{1.416317in}}%
\pgfpathlineto{\pgfqpoint{3.255954in}{1.419237in}}%
\pgfpathlineto{\pgfqpoint{3.254537in}{1.428747in}}%
\pgfpathlineto{\pgfqpoint{3.261997in}{1.431811in}}%
\pgfpathlineto{\pgfqpoint{3.248879in}{1.439534in}}%
\pgfpathlineto{\pgfqpoint{3.254321in}{1.452877in}}%
\pgfpathlineto{\pgfqpoint{3.250356in}{1.457926in}}%
\pgfpathlineto{\pgfqpoint{3.250047in}{1.464723in}}%
\pgfpathlineto{\pgfqpoint{3.246140in}{1.469571in}}%
\pgfpathlineto{\pgfqpoint{3.249034in}{1.483689in}}%
\pgfpathlineto{\pgfqpoint{3.256327in}{1.489994in}}%
\pgfpathlineto{\pgfqpoint{3.266239in}{1.490550in}}%
\pgfpathlineto{\pgfqpoint{3.273108in}{1.493337in}}%
\pgfpathlineto{\pgfqpoint{3.274847in}{1.487194in}}%
\pgfpathlineto{\pgfqpoint{3.290599in}{1.431089in}}%
\pgfpathlineto{\pgfqpoint{3.284538in}{1.423341in}}%
\pgfpathlineto{\pgfqpoint{3.281272in}{1.408524in}}%
\pgfpathlineto{\pgfqpoint{3.282347in}{1.403895in}}%
\pgfpathclose%
\pgfusepath{fill}%
\end{pgfscope}%
\begin{pgfscope}%
\pgfpathrectangle{\pgfqpoint{0.100000in}{0.100000in}}{\pgfqpoint{3.608454in}{2.310000in}}%
\pgfusepath{clip}%
\pgfsetbuttcap%
\pgfsetmiterjoin%
\definecolor{currentfill}{rgb}{0.000000,0.443137,0.778431}%
\pgfsetfillcolor{currentfill}%
\pgfsetlinewidth{0.000000pt}%
\definecolor{currentstroke}{rgb}{0.000000,0.000000,0.000000}%
\pgfsetstrokecolor{currentstroke}%
\pgfsetstrokeopacity{0.000000}%
\pgfsetdash{}{0pt}%
\pgfpathmoveto{\pgfqpoint{2.529824in}{1.778725in}}%
\pgfpathlineto{\pgfqpoint{2.510243in}{1.777129in}}%
\pgfpathlineto{\pgfqpoint{2.510762in}{1.770224in}}%
\pgfpathlineto{\pgfqpoint{2.501961in}{1.769586in}}%
\pgfpathlineto{\pgfqpoint{2.493489in}{1.769026in}}%
\pgfpathlineto{\pgfqpoint{2.491553in}{1.796332in}}%
\pgfpathlineto{\pgfqpoint{2.488486in}{1.795943in}}%
\pgfpathlineto{\pgfqpoint{2.486722in}{1.817142in}}%
\pgfpathlineto{\pgfqpoint{2.473187in}{1.816374in}}%
\pgfpathlineto{\pgfqpoint{2.472001in}{1.837128in}}%
\pgfpathlineto{\pgfqpoint{2.463191in}{1.836540in}}%
\pgfpathlineto{\pgfqpoint{2.460529in}{1.843339in}}%
\pgfpathlineto{\pgfqpoint{2.459694in}{1.857127in}}%
\pgfpathlineto{\pgfqpoint{2.473620in}{1.857983in}}%
\pgfpathlineto{\pgfqpoint{2.480448in}{1.858288in}}%
\pgfpathlineto{\pgfqpoint{2.481279in}{1.844556in}}%
\pgfpathlineto{\pgfqpoint{2.488144in}{1.844796in}}%
\pgfpathlineto{\pgfqpoint{2.492447in}{1.838114in}}%
\pgfpathlineto{\pgfqpoint{2.492843in}{1.831235in}}%
\pgfpathlineto{\pgfqpoint{2.513293in}{1.828191in}}%
\pgfpathlineto{\pgfqpoint{2.502443in}{1.807863in}}%
\pgfpathlineto{\pgfqpoint{2.500063in}{1.795356in}}%
\pgfpathlineto{\pgfqpoint{2.505725in}{1.792829in}}%
\pgfpathlineto{\pgfqpoint{2.509822in}{1.800026in}}%
\pgfpathlineto{\pgfqpoint{2.514631in}{1.802211in}}%
\pgfpathlineto{\pgfqpoint{2.522578in}{1.818649in}}%
\pgfpathlineto{\pgfqpoint{2.529869in}{1.821187in}}%
\pgfpathlineto{\pgfqpoint{2.533595in}{1.825201in}}%
\pgfpathlineto{\pgfqpoint{2.540562in}{1.840166in}}%
\pgfpathlineto{\pgfqpoint{2.541258in}{1.846577in}}%
\pgfpathlineto{\pgfqpoint{2.552435in}{1.851077in}}%
\pgfpathlineto{\pgfqpoint{2.552379in}{1.842511in}}%
\pgfpathlineto{\pgfqpoint{2.545403in}{1.830966in}}%
\pgfpathlineto{\pgfqpoint{2.545230in}{1.823543in}}%
\pgfpathlineto{\pgfqpoint{2.541717in}{1.820773in}}%
\pgfpathlineto{\pgfqpoint{2.533894in}{1.801143in}}%
\pgfpathlineto{\pgfqpoint{2.529824in}{1.778725in}}%
\pgfpathclose%
\pgfusepath{fill}%
\end{pgfscope}%
\begin{pgfscope}%
\pgfpathrectangle{\pgfqpoint{0.100000in}{0.100000in}}{\pgfqpoint{3.608454in}{2.310000in}}%
\pgfusepath{clip}%
\pgfsetbuttcap%
\pgfsetmiterjoin%
\definecolor{currentfill}{rgb}{0.000000,0.654902,0.672549}%
\pgfsetfillcolor{currentfill}%
\pgfsetlinewidth{0.000000pt}%
\definecolor{currentstroke}{rgb}{0.000000,0.000000,0.000000}%
\pgfsetstrokecolor{currentstroke}%
\pgfsetstrokeopacity{0.000000}%
\pgfsetdash{}{0pt}%
\pgfpathmoveto{\pgfqpoint{3.143924in}{1.538541in}}%
\pgfpathlineto{\pgfqpoint{3.137378in}{1.521010in}}%
\pgfpathlineto{\pgfqpoint{3.133137in}{1.523448in}}%
\pgfpathlineto{\pgfqpoint{3.119792in}{1.526024in}}%
\pgfpathlineto{\pgfqpoint{3.112902in}{1.529498in}}%
\pgfpathlineto{\pgfqpoint{3.110584in}{1.534960in}}%
\pgfpathlineto{\pgfqpoint{3.115439in}{1.550679in}}%
\pgfpathlineto{\pgfqpoint{3.110368in}{1.555879in}}%
\pgfpathlineto{\pgfqpoint{3.111267in}{1.560885in}}%
\pgfpathlineto{\pgfqpoint{3.107213in}{1.564628in}}%
\pgfpathlineto{\pgfqpoint{3.098915in}{1.567130in}}%
\pgfpathlineto{\pgfqpoint{3.072318in}{1.562246in}}%
\pgfpathlineto{\pgfqpoint{3.069571in}{1.576329in}}%
\pgfpathlineto{\pgfqpoint{3.045363in}{1.572457in}}%
\pgfpathlineto{\pgfqpoint{3.043474in}{1.583306in}}%
\pgfpathlineto{\pgfqpoint{3.039983in}{1.605278in}}%
\pgfpathlineto{\pgfqpoint{3.046524in}{1.607153in}}%
\pgfpathlineto{\pgfqpoint{3.054170in}{1.610224in}}%
\pgfpathlineto{\pgfqpoint{3.068655in}{1.604850in}}%
\pgfpathlineto{\pgfqpoint{3.106858in}{1.608202in}}%
\pgfpathlineto{\pgfqpoint{3.111013in}{1.603891in}}%
\pgfpathlineto{\pgfqpoint{3.125203in}{1.608488in}}%
\pgfpathlineto{\pgfqpoint{3.128786in}{1.603472in}}%
\pgfpathlineto{\pgfqpoint{3.133554in}{1.604605in}}%
\pgfpathlineto{\pgfqpoint{3.144862in}{1.595009in}}%
\pgfpathlineto{\pgfqpoint{3.154199in}{1.598846in}}%
\pgfpathlineto{\pgfqpoint{3.165006in}{1.605498in}}%
\pgfpathlineto{\pgfqpoint{3.155139in}{1.587716in}}%
\pgfpathlineto{\pgfqpoint{3.139700in}{1.576466in}}%
\pgfpathlineto{\pgfqpoint{3.131483in}{1.559276in}}%
\pgfpathlineto{\pgfqpoint{3.133732in}{1.555574in}}%
\pgfpathlineto{\pgfqpoint{3.130234in}{1.548464in}}%
\pgfpathlineto{\pgfqpoint{3.133310in}{1.544667in}}%
\pgfpathlineto{\pgfqpoint{3.138889in}{1.546016in}}%
\pgfpathlineto{\pgfqpoint{3.143924in}{1.538541in}}%
\pgfpathclose%
\pgfusepath{fill}%
\end{pgfscope}%
\begin{pgfscope}%
\pgfpathrectangle{\pgfqpoint{0.100000in}{0.100000in}}{\pgfqpoint{3.608454in}{2.310000in}}%
\pgfusepath{clip}%
\pgfsetbuttcap%
\pgfsetmiterjoin%
\definecolor{currentfill}{rgb}{0.000000,0.345098,0.827451}%
\pgfsetfillcolor{currentfill}%
\pgfsetlinewidth{0.000000pt}%
\definecolor{currentstroke}{rgb}{0.000000,0.000000,0.000000}%
\pgfsetstrokecolor{currentstroke}%
\pgfsetstrokeopacity{0.000000}%
\pgfsetdash{}{0pt}%
\pgfpathmoveto{\pgfqpoint{2.196902in}{1.594193in}}%
\pgfpathlineto{\pgfqpoint{2.155774in}{1.593709in}}%
\pgfpathlineto{\pgfqpoint{2.115090in}{1.593467in}}%
\pgfpathlineto{\pgfqpoint{2.115152in}{1.614399in}}%
\pgfpathlineto{\pgfqpoint{2.111772in}{1.614418in}}%
\pgfpathlineto{\pgfqpoint{2.111776in}{1.648912in}}%
\pgfpathlineto{\pgfqpoint{2.166463in}{1.649078in}}%
\pgfpathlineto{\pgfqpoint{2.193901in}{1.649520in}}%
\pgfpathlineto{\pgfqpoint{2.194446in}{1.614873in}}%
\pgfpathlineto{\pgfqpoint{2.196575in}{1.614898in}}%
\pgfpathlineto{\pgfqpoint{2.196902in}{1.594193in}}%
\pgfpathclose%
\pgfusepath{fill}%
\end{pgfscope}%
\begin{pgfscope}%
\pgfpathrectangle{\pgfqpoint{0.100000in}{0.100000in}}{\pgfqpoint{3.608454in}{2.310000in}}%
\pgfusepath{clip}%
\pgfsetbuttcap%
\pgfsetmiterjoin%
\definecolor{currentfill}{rgb}{0.000000,0.447059,0.776471}%
\pgfsetfillcolor{currentfill}%
\pgfsetlinewidth{0.000000pt}%
\definecolor{currentstroke}{rgb}{0.000000,0.000000,0.000000}%
\pgfsetstrokecolor{currentstroke}%
\pgfsetstrokeopacity{0.000000}%
\pgfsetdash{}{0pt}%
\pgfpathmoveto{\pgfqpoint{3.466647in}{1.697217in}}%
\pgfpathlineto{\pgfqpoint{3.461826in}{1.697323in}}%
\pgfpathlineto{\pgfqpoint{3.452537in}{1.691767in}}%
\pgfpathlineto{\pgfqpoint{3.447030in}{1.691452in}}%
\pgfpathlineto{\pgfqpoint{3.437379in}{1.684748in}}%
\pgfpathlineto{\pgfqpoint{3.433364in}{1.685208in}}%
\pgfpathlineto{\pgfqpoint{3.406113in}{1.677314in}}%
\pgfpathlineto{\pgfqpoint{3.369189in}{1.645562in}}%
\pgfpathlineto{\pgfqpoint{3.362988in}{1.653409in}}%
\pgfpathlineto{\pgfqpoint{3.375111in}{1.665432in}}%
\pgfpathlineto{\pgfqpoint{3.369634in}{1.670838in}}%
\pgfpathlineto{\pgfqpoint{3.366550in}{1.688953in}}%
\pgfpathlineto{\pgfqpoint{3.359316in}{1.729735in}}%
\pgfpathlineto{\pgfqpoint{3.384285in}{1.735025in}}%
\pgfpathlineto{\pgfqpoint{3.436777in}{1.747450in}}%
\pgfpathlineto{\pgfqpoint{3.456367in}{1.750710in}}%
\pgfpathlineto{\pgfqpoint{3.464103in}{1.723936in}}%
\pgfpathlineto{\pgfqpoint{3.468281in}{1.705449in}}%
\pgfpathlineto{\pgfqpoint{3.466647in}{1.697217in}}%
\pgfpathclose%
\pgfusepath{fill}%
\end{pgfscope}%
\begin{pgfscope}%
\pgfpathrectangle{\pgfqpoint{0.100000in}{0.100000in}}{\pgfqpoint{3.608454in}{2.310000in}}%
\pgfusepath{clip}%
\pgfsetbuttcap%
\pgfsetmiterjoin%
\definecolor{currentfill}{rgb}{0.000000,0.580392,0.709804}%
\pgfsetfillcolor{currentfill}%
\pgfsetlinewidth{0.000000pt}%
\definecolor{currentstroke}{rgb}{0.000000,0.000000,0.000000}%
\pgfsetstrokecolor{currentstroke}%
\pgfsetstrokeopacity{0.000000}%
\pgfsetdash{}{0pt}%
\pgfpathmoveto{\pgfqpoint{2.353369in}{1.446719in}}%
\pgfpathlineto{\pgfqpoint{2.353565in}{1.439495in}}%
\pgfpathlineto{\pgfqpoint{2.317746in}{1.438593in}}%
\pgfpathlineto{\pgfqpoint{2.318087in}{1.442426in}}%
\pgfpathlineto{\pgfqpoint{2.290718in}{1.442111in}}%
\pgfpathlineto{\pgfqpoint{2.290253in}{1.469670in}}%
\pgfpathlineto{\pgfqpoint{2.303069in}{1.470734in}}%
\pgfpathlineto{\pgfqpoint{2.302293in}{1.514165in}}%
\pgfpathlineto{\pgfqpoint{2.322856in}{1.514824in}}%
\pgfpathlineto{\pgfqpoint{2.323109in}{1.507938in}}%
\pgfpathlineto{\pgfqpoint{2.348213in}{1.508885in}}%
\pgfpathlineto{\pgfqpoint{2.358001in}{1.509016in}}%
\pgfpathlineto{\pgfqpoint{2.359331in}{1.474799in}}%
\pgfpathlineto{\pgfqpoint{2.352440in}{1.474762in}}%
\pgfpathlineto{\pgfqpoint{2.353369in}{1.446719in}}%
\pgfpathclose%
\pgfusepath{fill}%
\end{pgfscope}%
\begin{pgfscope}%
\pgfpathrectangle{\pgfqpoint{0.100000in}{0.100000in}}{\pgfqpoint{3.608454in}{2.310000in}}%
\pgfusepath{clip}%
\pgfsetbuttcap%
\pgfsetmiterjoin%
\definecolor{currentfill}{rgb}{0.000000,0.635294,0.682353}%
\pgfsetfillcolor{currentfill}%
\pgfsetlinewidth{0.000000pt}%
\definecolor{currentstroke}{rgb}{0.000000,0.000000,0.000000}%
\pgfsetstrokecolor{currentstroke}%
\pgfsetstrokeopacity{0.000000}%
\pgfsetdash{}{0pt}%
\pgfpathmoveto{\pgfqpoint{2.798375in}{1.062974in}}%
\pgfpathlineto{\pgfqpoint{2.779057in}{1.060719in}}%
\pgfpathlineto{\pgfqpoint{2.756131in}{1.058082in}}%
\pgfpathlineto{\pgfqpoint{2.751742in}{1.068739in}}%
\pgfpathlineto{\pgfqpoint{2.755774in}{1.081517in}}%
\pgfpathlineto{\pgfqpoint{2.749929in}{1.091001in}}%
\pgfpathlineto{\pgfqpoint{2.755630in}{1.093871in}}%
\pgfpathlineto{\pgfqpoint{2.763588in}{1.107918in}}%
\pgfpathlineto{\pgfqpoint{2.762622in}{1.113861in}}%
\pgfpathlineto{\pgfqpoint{2.765621in}{1.120136in}}%
\pgfpathlineto{\pgfqpoint{2.772428in}{1.117636in}}%
\pgfpathlineto{\pgfqpoint{2.775615in}{1.111731in}}%
\pgfpathlineto{\pgfqpoint{2.790002in}{1.116268in}}%
\pgfpathlineto{\pgfqpoint{2.805797in}{1.111156in}}%
\pgfpathlineto{\pgfqpoint{2.817068in}{1.103142in}}%
\pgfpathlineto{\pgfqpoint{2.813980in}{1.097663in}}%
\pgfpathlineto{\pgfqpoint{2.814347in}{1.089123in}}%
\pgfpathlineto{\pgfqpoint{2.807766in}{1.084602in}}%
\pgfpathlineto{\pgfqpoint{2.802072in}{1.085804in}}%
\pgfpathlineto{\pgfqpoint{2.798205in}{1.081874in}}%
\pgfpathlineto{\pgfqpoint{2.798375in}{1.062974in}}%
\pgfpathclose%
\pgfusepath{fill}%
\end{pgfscope}%
\begin{pgfscope}%
\pgfpathrectangle{\pgfqpoint{0.100000in}{0.100000in}}{\pgfqpoint{3.608454in}{2.310000in}}%
\pgfusepath{clip}%
\pgfsetbuttcap%
\pgfsetmiterjoin%
\definecolor{currentfill}{rgb}{0.000000,0.682353,0.658824}%
\pgfsetfillcolor{currentfill}%
\pgfsetlinewidth{0.000000pt}%
\definecolor{currentstroke}{rgb}{0.000000,0.000000,0.000000}%
\pgfsetstrokecolor{currentstroke}%
\pgfsetstrokeopacity{0.000000}%
\pgfsetdash{}{0pt}%
\pgfpathmoveto{\pgfqpoint{1.516667in}{1.064131in}}%
\pgfpathlineto{\pgfqpoint{1.517207in}{1.069928in}}%
\pgfpathlineto{\pgfqpoint{1.442195in}{1.078025in}}%
\pgfpathlineto{\pgfqpoint{1.440664in}{1.064305in}}%
\pgfpathlineto{\pgfqpoint{1.413364in}{1.067432in}}%
\pgfpathlineto{\pgfqpoint{1.421911in}{1.141159in}}%
\pgfpathlineto{\pgfqpoint{1.431076in}{1.139578in}}%
\pgfpathlineto{\pgfqpoint{1.442215in}{1.150974in}}%
\pgfpathlineto{\pgfqpoint{1.445381in}{1.160260in}}%
\pgfpathlineto{\pgfqpoint{1.448410in}{1.160662in}}%
\pgfpathlineto{\pgfqpoint{1.450548in}{1.172051in}}%
\pgfpathlineto{\pgfqpoint{1.450272in}{1.192914in}}%
\pgfpathlineto{\pgfqpoint{1.458207in}{1.195789in}}%
\pgfpathlineto{\pgfqpoint{1.462472in}{1.211222in}}%
\pgfpathlineto{\pgfqpoint{1.462544in}{1.217413in}}%
\pgfpathlineto{\pgfqpoint{1.466660in}{1.216957in}}%
\pgfpathlineto{\pgfqpoint{1.513551in}{1.211748in}}%
\pgfpathlineto{\pgfqpoint{1.538731in}{1.209408in}}%
\pgfpathlineto{\pgfqpoint{1.532570in}{1.148149in}}%
\pgfpathlineto{\pgfqpoint{1.505428in}{1.150972in}}%
\pgfpathlineto{\pgfqpoint{1.508502in}{1.136752in}}%
\pgfpathlineto{\pgfqpoint{1.508110in}{1.125109in}}%
\pgfpathlineto{\pgfqpoint{1.506250in}{1.116084in}}%
\pgfpathlineto{\pgfqpoint{1.528359in}{1.113844in}}%
\pgfpathlineto{\pgfqpoint{1.531357in}{1.115154in}}%
\pgfpathlineto{\pgfqpoint{1.544598in}{1.083669in}}%
\pgfpathlineto{\pgfqpoint{1.550211in}{1.083114in}}%
\pgfpathlineto{\pgfqpoint{1.548858in}{1.068831in}}%
\pgfpathlineto{\pgfqpoint{1.534671in}{1.070306in}}%
\pgfpathlineto{\pgfqpoint{1.516667in}{1.064131in}}%
\pgfpathclose%
\pgfusepath{fill}%
\end{pgfscope}%
\begin{pgfscope}%
\pgfpathrectangle{\pgfqpoint{0.100000in}{0.100000in}}{\pgfqpoint{3.608454in}{2.310000in}}%
\pgfusepath{clip}%
\pgfsetbuttcap%
\pgfsetmiterjoin%
\definecolor{currentfill}{rgb}{0.000000,0.411765,0.794118}%
\pgfsetfillcolor{currentfill}%
\pgfsetlinewidth{0.000000pt}%
\definecolor{currentstroke}{rgb}{0.000000,0.000000,0.000000}%
\pgfsetstrokecolor{currentstroke}%
\pgfsetstrokeopacity{0.000000}%
\pgfsetdash{}{0pt}%
\pgfpathmoveto{\pgfqpoint{2.266869in}{1.540448in}}%
\pgfpathlineto{\pgfqpoint{2.205803in}{1.538835in}}%
\pgfpathlineto{\pgfqpoint{2.204489in}{1.546166in}}%
\pgfpathlineto{\pgfqpoint{2.204148in}{1.566920in}}%
\pgfpathlineto{\pgfqpoint{2.210987in}{1.567017in}}%
\pgfpathlineto{\pgfqpoint{2.210467in}{1.594411in}}%
\pgfpathlineto{\pgfqpoint{2.223956in}{1.594717in}}%
\pgfpathlineto{\pgfqpoint{2.237674in}{1.595106in}}%
\pgfpathlineto{\pgfqpoint{2.237469in}{1.601952in}}%
\pgfpathlineto{\pgfqpoint{2.264937in}{1.602761in}}%
\pgfpathlineto{\pgfqpoint{2.266869in}{1.540448in}}%
\pgfpathclose%
\pgfusepath{fill}%
\end{pgfscope}%
\begin{pgfscope}%
\pgfpathrectangle{\pgfqpoint{0.100000in}{0.100000in}}{\pgfqpoint{3.608454in}{2.310000in}}%
\pgfusepath{clip}%
\pgfsetbuttcap%
\pgfsetmiterjoin%
\definecolor{currentfill}{rgb}{0.000000,0.647059,0.676471}%
\pgfsetfillcolor{currentfill}%
\pgfsetlinewidth{0.000000pt}%
\definecolor{currentstroke}{rgb}{0.000000,0.000000,0.000000}%
\pgfsetstrokecolor{currentstroke}%
\pgfsetstrokeopacity{0.000000}%
\pgfsetdash{}{0pt}%
\pgfpathmoveto{\pgfqpoint{0.718561in}{1.521037in}}%
\pgfpathlineto{\pgfqpoint{0.697419in}{1.520493in}}%
\pgfpathlineto{\pgfqpoint{0.695807in}{1.514308in}}%
\pgfpathlineto{\pgfqpoint{0.717194in}{1.472445in}}%
\pgfpathlineto{\pgfqpoint{0.665974in}{1.438337in}}%
\pgfpathlineto{\pgfqpoint{0.661400in}{1.439746in}}%
\pgfpathlineto{\pgfqpoint{0.628251in}{1.491065in}}%
\pgfpathlineto{\pgfqpoint{0.643201in}{1.487011in}}%
\pgfpathlineto{\pgfqpoint{0.652474in}{1.520325in}}%
\pgfpathlineto{\pgfqpoint{0.645777in}{1.522062in}}%
\pgfpathlineto{\pgfqpoint{0.647659in}{1.529088in}}%
\pgfpathlineto{\pgfqpoint{0.655392in}{1.537674in}}%
\pgfpathlineto{\pgfqpoint{0.718561in}{1.521037in}}%
\pgfpathclose%
\pgfusepath{fill}%
\end{pgfscope}%
\begin{pgfscope}%
\pgfpathrectangle{\pgfqpoint{0.100000in}{0.100000in}}{\pgfqpoint{3.608454in}{2.310000in}}%
\pgfusepath{clip}%
\pgfsetbuttcap%
\pgfsetmiterjoin%
\definecolor{currentfill}{rgb}{0.000000,0.737255,0.631373}%
\pgfsetfillcolor{currentfill}%
\pgfsetlinewidth{0.000000pt}%
\definecolor{currentstroke}{rgb}{0.000000,0.000000,0.000000}%
\pgfsetstrokecolor{currentstroke}%
\pgfsetstrokeopacity{0.000000}%
\pgfsetdash{}{0pt}%
\pgfpathmoveto{\pgfqpoint{2.996900in}{1.203053in}}%
\pgfpathlineto{\pgfqpoint{2.987922in}{1.204284in}}%
\pgfpathlineto{\pgfqpoint{2.980945in}{1.198576in}}%
\pgfpathlineto{\pgfqpoint{2.970773in}{1.210739in}}%
\pgfpathlineto{\pgfqpoint{2.970559in}{1.212860in}}%
\pgfpathlineto{\pgfqpoint{2.950033in}{1.210864in}}%
\pgfpathlineto{\pgfqpoint{2.951682in}{1.213011in}}%
\pgfpathlineto{\pgfqpoint{2.957411in}{1.217461in}}%
\pgfpathlineto{\pgfqpoint{2.958209in}{1.221651in}}%
\pgfpathlineto{\pgfqpoint{2.974075in}{1.228469in}}%
\pgfpathlineto{\pgfqpoint{2.984220in}{1.230496in}}%
\pgfpathlineto{\pgfqpoint{2.986436in}{1.233543in}}%
\pgfpathlineto{\pgfqpoint{3.005275in}{1.242589in}}%
\pgfpathlineto{\pgfqpoint{3.011272in}{1.247768in}}%
\pgfpathlineto{\pgfqpoint{3.024683in}{1.232048in}}%
\pgfpathlineto{\pgfqpoint{3.025017in}{1.229246in}}%
\pgfpathlineto{\pgfqpoint{3.015913in}{1.223529in}}%
\pgfpathlineto{\pgfqpoint{3.017299in}{1.218774in}}%
\pgfpathlineto{\pgfqpoint{2.999075in}{1.216208in}}%
\pgfpathlineto{\pgfqpoint{2.995823in}{1.206373in}}%
\pgfpathlineto{\pgfqpoint{2.996900in}{1.203053in}}%
\pgfpathclose%
\pgfusepath{fill}%
\end{pgfscope}%
\begin{pgfscope}%
\pgfpathrectangle{\pgfqpoint{0.100000in}{0.100000in}}{\pgfqpoint{3.608454in}{2.310000in}}%
\pgfusepath{clip}%
\pgfsetbuttcap%
\pgfsetmiterjoin%
\definecolor{currentfill}{rgb}{0.000000,0.678431,0.660784}%
\pgfsetfillcolor{currentfill}%
\pgfsetlinewidth{0.000000pt}%
\definecolor{currentstroke}{rgb}{0.000000,0.000000,0.000000}%
\pgfsetstrokecolor{currentstroke}%
\pgfsetstrokeopacity{0.000000}%
\pgfsetdash{}{0pt}%
\pgfpathmoveto{\pgfqpoint{2.869199in}{1.556678in}}%
\pgfpathlineto{\pgfqpoint{2.872479in}{1.551436in}}%
\pgfpathlineto{\pgfqpoint{2.876286in}{1.526111in}}%
\pgfpathlineto{\pgfqpoint{2.870685in}{1.525264in}}%
\pgfpathlineto{\pgfqpoint{2.871819in}{1.517489in}}%
\pgfpathlineto{\pgfqpoint{2.862804in}{1.514836in}}%
\pgfpathlineto{\pgfqpoint{2.846858in}{1.512561in}}%
\pgfpathlineto{\pgfqpoint{2.848845in}{1.496480in}}%
\pgfpathlineto{\pgfqpoint{2.831782in}{1.495222in}}%
\pgfpathlineto{\pgfqpoint{2.831500in}{1.499705in}}%
\pgfpathlineto{\pgfqpoint{2.811880in}{1.499432in}}%
\pgfpathlineto{\pgfqpoint{2.811270in}{1.504385in}}%
\pgfpathlineto{\pgfqpoint{2.801132in}{1.502933in}}%
\pgfpathlineto{\pgfqpoint{2.799219in}{1.517163in}}%
\pgfpathlineto{\pgfqpoint{2.817395in}{1.520810in}}%
\pgfpathlineto{\pgfqpoint{2.814322in}{1.543583in}}%
\pgfpathlineto{\pgfqpoint{2.831040in}{1.546066in}}%
\pgfpathlineto{\pgfqpoint{2.854535in}{1.548999in}}%
\pgfpathlineto{\pgfqpoint{2.853507in}{1.554639in}}%
\pgfpathlineto{\pgfqpoint{2.869199in}{1.556678in}}%
\pgfpathclose%
\pgfusepath{fill}%
\end{pgfscope}%
\begin{pgfscope}%
\pgfpathrectangle{\pgfqpoint{0.100000in}{0.100000in}}{\pgfqpoint{3.608454in}{2.310000in}}%
\pgfusepath{clip}%
\pgfsetbuttcap%
\pgfsetmiterjoin%
\definecolor{currentfill}{rgb}{0.000000,0.878431,0.560784}%
\pgfsetfillcolor{currentfill}%
\pgfsetlinewidth{0.000000pt}%
\definecolor{currentstroke}{rgb}{0.000000,0.000000,0.000000}%
\pgfsetstrokecolor{currentstroke}%
\pgfsetstrokeopacity{0.000000}%
\pgfsetdash{}{0pt}%
\pgfpathmoveto{\pgfqpoint{2.722462in}{1.100629in}}%
\pgfpathlineto{\pgfqpoint{2.731314in}{1.113832in}}%
\pgfpathlineto{\pgfqpoint{2.732727in}{1.126386in}}%
\pgfpathlineto{\pgfqpoint{2.729200in}{1.134130in}}%
\pgfpathlineto{\pgfqpoint{2.729930in}{1.144288in}}%
\pgfpathlineto{\pgfqpoint{2.735267in}{1.144371in}}%
\pgfpathlineto{\pgfqpoint{2.738900in}{1.152461in}}%
\pgfpathlineto{\pgfqpoint{2.735445in}{1.160714in}}%
\pgfpathlineto{\pgfqpoint{2.734679in}{1.168795in}}%
\pgfpathlineto{\pgfqpoint{2.738211in}{1.181250in}}%
\pgfpathlineto{\pgfqpoint{2.744380in}{1.183721in}}%
\pgfpathlineto{\pgfqpoint{2.758225in}{1.180770in}}%
\pgfpathlineto{\pgfqpoint{2.763981in}{1.171120in}}%
\pgfpathlineto{\pgfqpoint{2.761563in}{1.168904in}}%
\pgfpathlineto{\pgfqpoint{2.751113in}{1.160516in}}%
\pgfpathlineto{\pgfqpoint{2.750357in}{1.150558in}}%
\pgfpathlineto{\pgfqpoint{2.763659in}{1.139195in}}%
\pgfpathlineto{\pgfqpoint{2.766985in}{1.132686in}}%
\pgfpathlineto{\pgfqpoint{2.761095in}{1.125193in}}%
\pgfpathlineto{\pgfqpoint{2.765621in}{1.120136in}}%
\pgfpathlineto{\pgfqpoint{2.762622in}{1.113861in}}%
\pgfpathlineto{\pgfqpoint{2.763588in}{1.107918in}}%
\pgfpathlineto{\pgfqpoint{2.755630in}{1.093871in}}%
\pgfpathlineto{\pgfqpoint{2.749929in}{1.091001in}}%
\pgfpathlineto{\pgfqpoint{2.739787in}{1.093594in}}%
\pgfpathlineto{\pgfqpoint{2.736951in}{1.085202in}}%
\pgfpathlineto{\pgfqpoint{2.724406in}{1.094850in}}%
\pgfpathlineto{\pgfqpoint{2.722462in}{1.100629in}}%
\pgfpathclose%
\pgfusepath{fill}%
\end{pgfscope}%
\begin{pgfscope}%
\pgfpathrectangle{\pgfqpoint{0.100000in}{0.100000in}}{\pgfqpoint{3.608454in}{2.310000in}}%
\pgfusepath{clip}%
\pgfsetbuttcap%
\pgfsetmiterjoin%
\definecolor{currentfill}{rgb}{0.000000,0.690196,0.654902}%
\pgfsetfillcolor{currentfill}%
\pgfsetlinewidth{0.000000pt}%
\definecolor{currentstroke}{rgb}{0.000000,0.000000,0.000000}%
\pgfsetstrokecolor{currentstroke}%
\pgfsetstrokeopacity{0.000000}%
\pgfsetdash{}{0pt}%
\pgfpathmoveto{\pgfqpoint{3.271268in}{1.222696in}}%
\pgfpathlineto{\pgfqpoint{3.272014in}{1.215831in}}%
\pgfpathlineto{\pgfqpoint{3.269662in}{1.209194in}}%
\pgfpathlineto{\pgfqpoint{3.265874in}{1.209619in}}%
\pgfpathlineto{\pgfqpoint{3.262155in}{1.203038in}}%
\pgfpathlineto{\pgfqpoint{3.258356in}{1.209245in}}%
\pgfpathlineto{\pgfqpoint{3.253779in}{1.205127in}}%
\pgfpathlineto{\pgfqpoint{3.247206in}{1.209157in}}%
\pgfpathlineto{\pgfqpoint{3.239989in}{1.209565in}}%
\pgfpathlineto{\pgfqpoint{3.238791in}{1.214139in}}%
\pgfpathlineto{\pgfqpoint{3.230833in}{1.216563in}}%
\pgfpathlineto{\pgfqpoint{3.226514in}{1.211542in}}%
\pgfpathlineto{\pgfqpoint{3.218220in}{1.212547in}}%
\pgfpathlineto{\pgfqpoint{3.216999in}{1.216602in}}%
\pgfpathlineto{\pgfqpoint{3.206943in}{1.219779in}}%
\pgfpathlineto{\pgfqpoint{3.195538in}{1.217166in}}%
\pgfpathlineto{\pgfqpoint{3.185638in}{1.219606in}}%
\pgfpathlineto{\pgfqpoint{3.188912in}{1.234693in}}%
\pgfpathlineto{\pgfqpoint{3.187275in}{1.247460in}}%
\pgfpathlineto{\pgfqpoint{3.195545in}{1.249102in}}%
\pgfpathlineto{\pgfqpoint{3.199968in}{1.263498in}}%
\pgfpathlineto{\pgfqpoint{3.197116in}{1.277528in}}%
\pgfpathlineto{\pgfqpoint{3.202830in}{1.275174in}}%
\pgfpathlineto{\pgfqpoint{3.209973in}{1.277484in}}%
\pgfpathlineto{\pgfqpoint{3.210178in}{1.266083in}}%
\pgfpathlineto{\pgfqpoint{3.214106in}{1.265774in}}%
\pgfpathlineto{\pgfqpoint{3.221614in}{1.254572in}}%
\pgfpathlineto{\pgfqpoint{3.248646in}{1.259422in}}%
\pgfpathlineto{\pgfqpoint{3.249894in}{1.248153in}}%
\pgfpathlineto{\pgfqpoint{3.259748in}{1.247063in}}%
\pgfpathlineto{\pgfqpoint{3.266234in}{1.243000in}}%
\pgfpathlineto{\pgfqpoint{3.265839in}{1.231923in}}%
\pgfpathlineto{\pgfqpoint{3.271268in}{1.222696in}}%
\pgfpathclose%
\pgfusepath{fill}%
\end{pgfscope}%
\begin{pgfscope}%
\pgfpathrectangle{\pgfqpoint{0.100000in}{0.100000in}}{\pgfqpoint{3.608454in}{2.310000in}}%
\pgfusepath{clip}%
\pgfsetbuttcap%
\pgfsetmiterjoin%
\definecolor{currentfill}{rgb}{0.000000,0.486275,0.756863}%
\pgfsetfillcolor{currentfill}%
\pgfsetlinewidth{0.000000pt}%
\definecolor{currentstroke}{rgb}{0.000000,0.000000,0.000000}%
\pgfsetstrokecolor{currentstroke}%
\pgfsetstrokeopacity{0.000000}%
\pgfsetdash{}{0pt}%
\pgfpathmoveto{\pgfqpoint{1.826418in}{1.110191in}}%
\pgfpathlineto{\pgfqpoint{1.818722in}{1.107444in}}%
\pgfpathlineto{\pgfqpoint{1.814486in}{1.099133in}}%
\pgfpathlineto{\pgfqpoint{1.804299in}{1.099161in}}%
\pgfpathlineto{\pgfqpoint{1.798623in}{1.104938in}}%
\pgfpathlineto{\pgfqpoint{1.792285in}{1.107896in}}%
\pgfpathlineto{\pgfqpoint{1.786242in}{1.101806in}}%
\pgfpathlineto{\pgfqpoint{1.787003in}{1.115616in}}%
\pgfpathlineto{\pgfqpoint{1.788727in}{1.150707in}}%
\pgfpathlineto{\pgfqpoint{1.790983in}{1.190359in}}%
\pgfpathlineto{\pgfqpoint{1.820080in}{1.188664in}}%
\pgfpathlineto{\pgfqpoint{1.825437in}{1.188388in}}%
\pgfpathlineto{\pgfqpoint{1.830118in}{1.178718in}}%
\pgfpathlineto{\pgfqpoint{1.837450in}{1.170858in}}%
\pgfpathlineto{\pgfqpoint{1.844616in}{1.170342in}}%
\pgfpathlineto{\pgfqpoint{1.855366in}{1.155288in}}%
\pgfpathlineto{\pgfqpoint{1.854121in}{1.120697in}}%
\pgfpathlineto{\pgfqpoint{1.826923in}{1.122162in}}%
\pgfpathlineto{\pgfqpoint{1.826418in}{1.110191in}}%
\pgfpathclose%
\pgfusepath{fill}%
\end{pgfscope}%
\begin{pgfscope}%
\pgfpathrectangle{\pgfqpoint{0.100000in}{0.100000in}}{\pgfqpoint{3.608454in}{2.310000in}}%
\pgfusepath{clip}%
\pgfsetbuttcap%
\pgfsetmiterjoin%
\definecolor{currentfill}{rgb}{0.000000,0.313725,0.843137}%
\pgfsetfillcolor{currentfill}%
\pgfsetlinewidth{0.000000pt}%
\definecolor{currentstroke}{rgb}{0.000000,0.000000,0.000000}%
\pgfsetstrokecolor{currentstroke}%
\pgfsetstrokeopacity{0.000000}%
\pgfsetdash{}{0pt}%
\pgfpathmoveto{\pgfqpoint{3.083124in}{1.089029in}}%
\pgfpathlineto{\pgfqpoint{3.042269in}{1.083215in}}%
\pgfpathlineto{\pgfqpoint{3.054685in}{1.068940in}}%
\pgfpathlineto{\pgfqpoint{3.046213in}{1.063185in}}%
\pgfpathlineto{\pgfqpoint{3.047179in}{1.057968in}}%
\pgfpathlineto{\pgfqpoint{3.040979in}{1.056871in}}%
\pgfpathlineto{\pgfqpoint{3.034032in}{1.059624in}}%
\pgfpathlineto{\pgfqpoint{3.026275in}{1.051810in}}%
\pgfpathlineto{\pgfqpoint{3.025339in}{1.058515in}}%
\pgfpathlineto{\pgfqpoint{2.989806in}{1.055375in}}%
\pgfpathlineto{\pgfqpoint{2.986559in}{1.063157in}}%
\pgfpathlineto{\pgfqpoint{2.984590in}{1.076632in}}%
\pgfpathlineto{\pgfqpoint{2.980220in}{1.091029in}}%
\pgfpathlineto{\pgfqpoint{2.985051in}{1.092680in}}%
\pgfpathlineto{\pgfqpoint{2.986399in}{1.102377in}}%
\pgfpathlineto{\pgfqpoint{3.008536in}{1.104446in}}%
\pgfpathlineto{\pgfqpoint{3.007642in}{1.113989in}}%
\pgfpathlineto{\pgfqpoint{3.009870in}{1.122844in}}%
\pgfpathlineto{\pgfqpoint{3.008399in}{1.134076in}}%
\pgfpathlineto{\pgfqpoint{3.022456in}{1.135384in}}%
\pgfpathlineto{\pgfqpoint{3.018191in}{1.148162in}}%
\pgfpathlineto{\pgfqpoint{3.020036in}{1.162708in}}%
\pgfpathlineto{\pgfqpoint{3.024951in}{1.162919in}}%
\pgfpathlineto{\pgfqpoint{3.037191in}{1.156638in}}%
\pgfpathlineto{\pgfqpoint{3.046178in}{1.152956in}}%
\pgfpathlineto{\pgfqpoint{3.055178in}{1.146538in}}%
\pgfpathlineto{\pgfqpoint{3.057775in}{1.140860in}}%
\pgfpathlineto{\pgfqpoint{3.062964in}{1.138236in}}%
\pgfpathlineto{\pgfqpoint{3.067859in}{1.131866in}}%
\pgfpathlineto{\pgfqpoint{3.066320in}{1.122502in}}%
\pgfpathlineto{\pgfqpoint{3.070229in}{1.116768in}}%
\pgfpathlineto{\pgfqpoint{3.070169in}{1.111400in}}%
\pgfpathlineto{\pgfqpoint{3.078110in}{1.113287in}}%
\pgfpathlineto{\pgfqpoint{3.083344in}{1.107498in}}%
\pgfpathlineto{\pgfqpoint{3.086909in}{1.097540in}}%
\pgfpathlineto{\pgfqpoint{3.083124in}{1.089029in}}%
\pgfpathclose%
\pgfusepath{fill}%
\end{pgfscope}%
\begin{pgfscope}%
\pgfpathrectangle{\pgfqpoint{0.100000in}{0.100000in}}{\pgfqpoint{3.608454in}{2.310000in}}%
\pgfusepath{clip}%
\pgfsetbuttcap%
\pgfsetmiterjoin%
\definecolor{currentfill}{rgb}{0.000000,0.439216,0.780392}%
\pgfsetfillcolor{currentfill}%
\pgfsetlinewidth{0.000000pt}%
\definecolor{currentstroke}{rgb}{0.000000,0.000000,0.000000}%
\pgfsetstrokecolor{currentstroke}%
\pgfsetstrokeopacity{0.000000}%
\pgfsetdash{}{0pt}%
\pgfpathmoveto{\pgfqpoint{2.034871in}{0.852102in}}%
\pgfpathlineto{\pgfqpoint{2.020147in}{0.852551in}}%
\pgfpathlineto{\pgfqpoint{2.020851in}{0.893088in}}%
\pgfpathlineto{\pgfqpoint{2.049865in}{0.897339in}}%
\pgfpathlineto{\pgfqpoint{2.050404in}{0.930484in}}%
\pgfpathlineto{\pgfqpoint{2.056583in}{0.934051in}}%
\pgfpathlineto{\pgfqpoint{2.068217in}{0.938848in}}%
\pgfpathlineto{\pgfqpoint{2.070424in}{0.935180in}}%
\pgfpathlineto{\pgfqpoint{2.079921in}{0.933374in}}%
\pgfpathlineto{\pgfqpoint{2.087358in}{0.933839in}}%
\pgfpathlineto{\pgfqpoint{2.092356in}{0.940891in}}%
\pgfpathlineto{\pgfqpoint{2.109858in}{0.929360in}}%
\pgfpathlineto{\pgfqpoint{2.115386in}{0.923776in}}%
\pgfpathlineto{\pgfqpoint{2.123559in}{0.920377in}}%
\pgfpathlineto{\pgfqpoint{2.123587in}{0.890672in}}%
\pgfpathlineto{\pgfqpoint{2.115003in}{0.895858in}}%
\pgfpathlineto{\pgfqpoint{2.100603in}{0.896104in}}%
\pgfpathlineto{\pgfqpoint{2.086328in}{0.894484in}}%
\pgfpathlineto{\pgfqpoint{2.086189in}{0.861574in}}%
\pgfpathlineto{\pgfqpoint{2.062420in}{0.861532in}}%
\pgfpathlineto{\pgfqpoint{2.064298in}{0.842428in}}%
\pgfpathlineto{\pgfqpoint{2.053670in}{0.846942in}}%
\pgfpathlineto{\pgfqpoint{2.050385in}{0.846426in}}%
\pgfpathlineto{\pgfqpoint{2.044379in}{0.851895in}}%
\pgfpathlineto{\pgfqpoint{2.034871in}{0.852102in}}%
\pgfpathclose%
\pgfusepath{fill}%
\end{pgfscope}%
\begin{pgfscope}%
\pgfpathrectangle{\pgfqpoint{0.100000in}{0.100000in}}{\pgfqpoint{3.608454in}{2.310000in}}%
\pgfusepath{clip}%
\pgfsetbuttcap%
\pgfsetmiterjoin%
\definecolor{currentfill}{rgb}{0.000000,0.956863,0.521569}%
\pgfsetfillcolor{currentfill}%
\pgfsetlinewidth{0.000000pt}%
\definecolor{currentstroke}{rgb}{0.000000,0.000000,0.000000}%
\pgfsetstrokecolor{currentstroke}%
\pgfsetstrokeopacity{0.000000}%
\pgfsetdash{}{0pt}%
\pgfpathmoveto{\pgfqpoint{2.961168in}{1.406181in}}%
\pgfpathlineto{\pgfqpoint{2.965023in}{1.393727in}}%
\pgfpathlineto{\pgfqpoint{2.962379in}{1.389071in}}%
\pgfpathlineto{\pgfqpoint{2.967663in}{1.383858in}}%
\pgfpathlineto{\pgfqpoint{2.965422in}{1.379505in}}%
\pgfpathlineto{\pgfqpoint{2.957191in}{1.367025in}}%
\pgfpathlineto{\pgfqpoint{2.940346in}{1.366252in}}%
\pgfpathlineto{\pgfqpoint{2.936140in}{1.370645in}}%
\pgfpathlineto{\pgfqpoint{2.933951in}{1.375085in}}%
\pgfpathlineto{\pgfqpoint{2.932362in}{1.392633in}}%
\pgfpathlineto{\pgfqpoint{2.936799in}{1.397031in}}%
\pgfpathlineto{\pgfqpoint{2.943680in}{1.394391in}}%
\pgfpathlineto{\pgfqpoint{2.954499in}{1.407993in}}%
\pgfpathlineto{\pgfqpoint{2.961168in}{1.406181in}}%
\pgfpathclose%
\pgfusepath{fill}%
\end{pgfscope}%
\begin{pgfscope}%
\pgfpathrectangle{\pgfqpoint{0.100000in}{0.100000in}}{\pgfqpoint{3.608454in}{2.310000in}}%
\pgfusepath{clip}%
\pgfsetbuttcap%
\pgfsetmiterjoin%
\definecolor{currentfill}{rgb}{0.000000,0.427451,0.786275}%
\pgfsetfillcolor{currentfill}%
\pgfsetlinewidth{0.000000pt}%
\definecolor{currentstroke}{rgb}{0.000000,0.000000,0.000000}%
\pgfsetstrokecolor{currentstroke}%
\pgfsetstrokeopacity{0.000000}%
\pgfsetdash{}{0pt}%
\pgfpathmoveto{\pgfqpoint{1.623124in}{2.077735in}}%
\pgfpathlineto{\pgfqpoint{1.626106in}{2.109098in}}%
\pgfpathlineto{\pgfqpoint{1.630877in}{2.157836in}}%
\pgfpathlineto{\pgfqpoint{1.689092in}{2.152335in}}%
\pgfpathlineto{\pgfqpoint{1.737269in}{2.148337in}}%
\pgfpathlineto{\pgfqpoint{1.736050in}{2.132965in}}%
\pgfpathlineto{\pgfqpoint{1.729180in}{2.133512in}}%
\pgfpathlineto{\pgfqpoint{1.728606in}{2.126565in}}%
\pgfpathlineto{\pgfqpoint{1.724331in}{2.126897in}}%
\pgfpathlineto{\pgfqpoint{1.723190in}{2.112972in}}%
\pgfpathlineto{\pgfqpoint{1.688717in}{2.115878in}}%
\pgfpathlineto{\pgfqpoint{1.687537in}{2.102006in}}%
\pgfpathlineto{\pgfqpoint{1.690580in}{2.101766in}}%
\pgfpathlineto{\pgfqpoint{1.686894in}{2.083116in}}%
\pgfpathlineto{\pgfqpoint{1.669040in}{2.084553in}}%
\pgfpathlineto{\pgfqpoint{1.664466in}{2.078762in}}%
\pgfpathlineto{\pgfqpoint{1.653010in}{2.075697in}}%
\pgfpathlineto{\pgfqpoint{1.648734in}{2.078432in}}%
\pgfpathlineto{\pgfqpoint{1.647283in}{2.084906in}}%
\pgfpathlineto{\pgfqpoint{1.643336in}{2.085862in}}%
\pgfpathlineto{\pgfqpoint{1.632154in}{2.074810in}}%
\pgfpathlineto{\pgfqpoint{1.623124in}{2.077735in}}%
\pgfpathclose%
\pgfusepath{fill}%
\end{pgfscope}%
\begin{pgfscope}%
\pgfpathrectangle{\pgfqpoint{0.100000in}{0.100000in}}{\pgfqpoint{3.608454in}{2.310000in}}%
\pgfusepath{clip}%
\pgfsetbuttcap%
\pgfsetmiterjoin%
\definecolor{currentfill}{rgb}{0.000000,0.505882,0.747059}%
\pgfsetfillcolor{currentfill}%
\pgfsetlinewidth{0.000000pt}%
\definecolor{currentstroke}{rgb}{0.000000,0.000000,0.000000}%
\pgfsetstrokecolor{currentstroke}%
\pgfsetstrokeopacity{0.000000}%
\pgfsetdash{}{0pt}%
\pgfpathmoveto{\pgfqpoint{1.576920in}{0.953485in}}%
\pgfpathlineto{\pgfqpoint{1.640100in}{0.948213in}}%
\pgfpathlineto{\pgfqpoint{1.708835in}{0.943953in}}%
\pgfpathlineto{\pgfqpoint{1.706689in}{0.909371in}}%
\pgfpathlineto{\pgfqpoint{1.704446in}{0.875520in}}%
\pgfpathlineto{\pgfqpoint{1.695447in}{0.875598in}}%
\pgfpathlineto{\pgfqpoint{1.635380in}{0.879674in}}%
\pgfpathlineto{\pgfqpoint{1.638056in}{0.913724in}}%
\pgfpathlineto{\pgfqpoint{1.573241in}{0.919040in}}%
\pgfpathlineto{\pgfqpoint{1.574789in}{0.933425in}}%
\pgfpathlineto{\pgfqpoint{1.576920in}{0.953485in}}%
\pgfpathclose%
\pgfusepath{fill}%
\end{pgfscope}%
\begin{pgfscope}%
\pgfpathrectangle{\pgfqpoint{0.100000in}{0.100000in}}{\pgfqpoint{3.608454in}{2.310000in}}%
\pgfusepath{clip}%
\pgfsetbuttcap%
\pgfsetmiterjoin%
\definecolor{currentfill}{rgb}{0.000000,0.584314,0.707843}%
\pgfsetfillcolor{currentfill}%
\pgfsetlinewidth{0.000000pt}%
\definecolor{currentstroke}{rgb}{0.000000,0.000000,0.000000}%
\pgfsetstrokecolor{currentstroke}%
\pgfsetstrokeopacity{0.000000}%
\pgfsetdash{}{0pt}%
\pgfpathmoveto{\pgfqpoint{3.145061in}{1.058883in}}%
\pgfpathlineto{\pgfqpoint{3.116266in}{1.079750in}}%
\pgfpathlineto{\pgfqpoint{3.099339in}{1.091621in}}%
\pgfpathlineto{\pgfqpoint{3.083124in}{1.089029in}}%
\pgfpathlineto{\pgfqpoint{3.086909in}{1.097540in}}%
\pgfpathlineto{\pgfqpoint{3.083344in}{1.107498in}}%
\pgfpathlineto{\pgfqpoint{3.078110in}{1.113287in}}%
\pgfpathlineto{\pgfqpoint{3.070169in}{1.111400in}}%
\pgfpathlineto{\pgfqpoint{3.070229in}{1.116768in}}%
\pgfpathlineto{\pgfqpoint{3.066320in}{1.122502in}}%
\pgfpathlineto{\pgfqpoint{3.067859in}{1.131866in}}%
\pgfpathlineto{\pgfqpoint{3.062964in}{1.138236in}}%
\pgfpathlineto{\pgfqpoint{3.057775in}{1.140860in}}%
\pgfpathlineto{\pgfqpoint{3.097477in}{1.148400in}}%
\pgfpathlineto{\pgfqpoint{3.110510in}{1.150879in}}%
\pgfpathlineto{\pgfqpoint{3.110689in}{1.147372in}}%
\pgfpathlineto{\pgfqpoint{3.121994in}{1.132796in}}%
\pgfpathlineto{\pgfqpoint{3.131128in}{1.128295in}}%
\pgfpathlineto{\pgfqpoint{3.149312in}{1.137485in}}%
\pgfpathlineto{\pgfqpoint{3.160885in}{1.137946in}}%
\pgfpathlineto{\pgfqpoint{3.160679in}{1.126465in}}%
\pgfpathlineto{\pgfqpoint{3.162834in}{1.117655in}}%
\pgfpathlineto{\pgfqpoint{3.174249in}{1.109020in}}%
\pgfpathlineto{\pgfqpoint{3.154787in}{1.104979in}}%
\pgfpathlineto{\pgfqpoint{3.148549in}{1.102631in}}%
\pgfpathlineto{\pgfqpoint{3.156648in}{1.092471in}}%
\pgfpathlineto{\pgfqpoint{3.156843in}{1.077749in}}%
\pgfpathlineto{\pgfqpoint{3.150261in}{1.071349in}}%
\pgfpathlineto{\pgfqpoint{3.145061in}{1.058883in}}%
\pgfpathclose%
\pgfusepath{fill}%
\end{pgfscope}%
\begin{pgfscope}%
\pgfpathrectangle{\pgfqpoint{0.100000in}{0.100000in}}{\pgfqpoint{3.608454in}{2.310000in}}%
\pgfusepath{clip}%
\pgfsetbuttcap%
\pgfsetmiterjoin%
\definecolor{currentfill}{rgb}{0.000000,0.435294,0.782353}%
\pgfsetfillcolor{currentfill}%
\pgfsetlinewidth{0.000000pt}%
\definecolor{currentstroke}{rgb}{0.000000,0.000000,0.000000}%
\pgfsetstrokecolor{currentstroke}%
\pgfsetstrokeopacity{0.000000}%
\pgfsetdash{}{0pt}%
\pgfpathmoveto{\pgfqpoint{1.964117in}{1.448246in}}%
\pgfpathlineto{\pgfqpoint{1.963356in}{1.420723in}}%
\pgfpathlineto{\pgfqpoint{1.908588in}{1.422459in}}%
\pgfpathlineto{\pgfqpoint{1.909610in}{1.449945in}}%
\pgfpathlineto{\pgfqpoint{1.936670in}{1.449037in}}%
\pgfpathlineto{\pgfqpoint{1.964117in}{1.448246in}}%
\pgfpathclose%
\pgfusepath{fill}%
\end{pgfscope}%
\begin{pgfscope}%
\pgfpathrectangle{\pgfqpoint{0.100000in}{0.100000in}}{\pgfqpoint{3.608454in}{2.310000in}}%
\pgfusepath{clip}%
\pgfsetbuttcap%
\pgfsetmiterjoin%
\definecolor{currentfill}{rgb}{0.000000,0.717647,0.641176}%
\pgfsetfillcolor{currentfill}%
\pgfsetlinewidth{0.000000pt}%
\definecolor{currentstroke}{rgb}{0.000000,0.000000,0.000000}%
\pgfsetstrokecolor{currentstroke}%
\pgfsetstrokeopacity{0.000000}%
\pgfsetdash{}{0pt}%
\pgfpathmoveto{\pgfqpoint{2.510006in}{1.121373in}}%
\pgfpathlineto{\pgfqpoint{2.498683in}{1.126546in}}%
\pgfpathlineto{\pgfqpoint{2.492139in}{1.131174in}}%
\pgfpathlineto{\pgfqpoint{2.479806in}{1.130609in}}%
\pgfpathlineto{\pgfqpoint{2.458843in}{1.129936in}}%
\pgfpathlineto{\pgfqpoint{2.462058in}{1.133851in}}%
\pgfpathlineto{\pgfqpoint{2.465696in}{1.145939in}}%
\pgfpathlineto{\pgfqpoint{2.461626in}{1.152803in}}%
\pgfpathlineto{\pgfqpoint{2.464384in}{1.163029in}}%
\pgfpathlineto{\pgfqpoint{2.473332in}{1.159111in}}%
\pgfpathlineto{\pgfqpoint{2.477166in}{1.167044in}}%
\pgfpathlineto{\pgfqpoint{2.474357in}{1.170803in}}%
\pgfpathlineto{\pgfqpoint{2.478809in}{1.176580in}}%
\pgfpathlineto{\pgfqpoint{2.498496in}{1.176925in}}%
\pgfpathlineto{\pgfqpoint{2.497579in}{1.190610in}}%
\pgfpathlineto{\pgfqpoint{2.518313in}{1.191718in}}%
\pgfpathlineto{\pgfqpoint{2.519321in}{1.176572in}}%
\pgfpathlineto{\pgfqpoint{2.520509in}{1.156888in}}%
\pgfpathlineto{\pgfqpoint{2.518781in}{1.156793in}}%
\pgfpathlineto{\pgfqpoint{2.519951in}{1.126253in}}%
\pgfpathlineto{\pgfqpoint{2.512080in}{1.125849in}}%
\pgfpathlineto{\pgfqpoint{2.510006in}{1.121373in}}%
\pgfpathclose%
\pgfusepath{fill}%
\end{pgfscope}%
\begin{pgfscope}%
\pgfpathrectangle{\pgfqpoint{0.100000in}{0.100000in}}{\pgfqpoint{3.608454in}{2.310000in}}%
\pgfusepath{clip}%
\pgfsetbuttcap%
\pgfsetmiterjoin%
\definecolor{currentfill}{rgb}{0.000000,0.474510,0.762745}%
\pgfsetfillcolor{currentfill}%
\pgfsetlinewidth{0.000000pt}%
\definecolor{currentstroke}{rgb}{0.000000,0.000000,0.000000}%
\pgfsetstrokecolor{currentstroke}%
\pgfsetstrokeopacity{0.000000}%
\pgfsetdash{}{0pt}%
\pgfpathmoveto{\pgfqpoint{1.401111in}{1.544070in}}%
\pgfpathlineto{\pgfqpoint{1.432841in}{1.539719in}}%
\pgfpathlineto{\pgfqpoint{1.436631in}{1.570812in}}%
\pgfpathlineto{\pgfqpoint{1.451617in}{1.568947in}}%
\pgfpathlineto{\pgfqpoint{1.457886in}{1.623314in}}%
\pgfpathlineto{\pgfqpoint{1.461309in}{1.650537in}}%
\pgfpathlineto{\pgfqpoint{1.491303in}{1.646739in}}%
\pgfpathlineto{\pgfqpoint{1.487410in}{1.642958in}}%
\pgfpathlineto{\pgfqpoint{1.486783in}{1.637214in}}%
\pgfpathlineto{\pgfqpoint{1.500514in}{1.634480in}}%
\pgfpathlineto{\pgfqpoint{1.504277in}{1.645283in}}%
\pgfpathlineto{\pgfqpoint{1.507108in}{1.644984in}}%
\pgfpathlineto{\pgfqpoint{1.503839in}{1.609878in}}%
\pgfpathlineto{\pgfqpoint{1.500517in}{1.584019in}}%
\pgfpathlineto{\pgfqpoint{1.494704in}{1.532257in}}%
\pgfpathlineto{\pgfqpoint{1.440561in}{1.538635in}}%
\pgfpathlineto{\pgfqpoint{1.440255in}{1.533604in}}%
\pgfpathlineto{\pgfqpoint{1.446700in}{1.523501in}}%
\pgfpathlineto{\pgfqpoint{1.455629in}{1.496066in}}%
\pgfpathlineto{\pgfqpoint{1.447426in}{1.486216in}}%
\pgfpathlineto{\pgfqpoint{1.444777in}{1.488383in}}%
\pgfpathlineto{\pgfqpoint{1.436560in}{1.486098in}}%
\pgfpathlineto{\pgfqpoint{1.426803in}{1.487168in}}%
\pgfpathlineto{\pgfqpoint{1.410993in}{1.493320in}}%
\pgfpathlineto{\pgfqpoint{1.407647in}{1.498770in}}%
\pgfpathlineto{\pgfqpoint{1.405998in}{1.509570in}}%
\pgfpathlineto{\pgfqpoint{1.412111in}{1.525836in}}%
\pgfpathlineto{\pgfqpoint{1.411585in}{1.530462in}}%
\pgfpathlineto{\pgfqpoint{1.400660in}{1.537925in}}%
\pgfpathlineto{\pgfqpoint{1.401111in}{1.544070in}}%
\pgfpathclose%
\pgfusepath{fill}%
\end{pgfscope}%
\begin{pgfscope}%
\pgfpathrectangle{\pgfqpoint{0.100000in}{0.100000in}}{\pgfqpoint{3.608454in}{2.310000in}}%
\pgfusepath{clip}%
\pgfsetbuttcap%
\pgfsetmiterjoin%
\definecolor{currentfill}{rgb}{0.000000,0.749020,0.625490}%
\pgfsetfillcolor{currentfill}%
\pgfsetlinewidth{0.000000pt}%
\definecolor{currentstroke}{rgb}{0.000000,0.000000,0.000000}%
\pgfsetstrokecolor{currentstroke}%
\pgfsetstrokeopacity{0.000000}%
\pgfsetdash{}{0pt}%
\pgfpathmoveto{\pgfqpoint{2.815217in}{1.759142in}}%
\pgfpathlineto{\pgfqpoint{2.786333in}{1.754076in}}%
\pgfpathlineto{\pgfqpoint{2.790730in}{1.726819in}}%
\pgfpathlineto{\pgfqpoint{2.776445in}{1.724555in}}%
\pgfpathlineto{\pgfqpoint{2.777478in}{1.717721in}}%
\pgfpathlineto{\pgfqpoint{2.771368in}{1.715974in}}%
\pgfpathlineto{\pgfqpoint{2.744382in}{1.712367in}}%
\pgfpathlineto{\pgfqpoint{2.745404in}{1.705512in}}%
\pgfpathlineto{\pgfqpoint{2.720255in}{1.702142in}}%
\pgfpathlineto{\pgfqpoint{2.717010in}{1.728719in}}%
\pgfpathlineto{\pgfqpoint{2.703428in}{1.727130in}}%
\pgfpathlineto{\pgfqpoint{2.697115in}{1.781860in}}%
\pgfpathlineto{\pgfqpoint{2.738013in}{1.786827in}}%
\pgfpathlineto{\pgfqpoint{2.756041in}{1.789197in}}%
\pgfpathlineto{\pgfqpoint{2.756426in}{1.784651in}}%
\pgfpathlineto{\pgfqpoint{2.746228in}{1.773703in}}%
\pgfpathlineto{\pgfqpoint{2.742788in}{1.773732in}}%
\pgfpathlineto{\pgfqpoint{2.738774in}{1.765761in}}%
\pgfpathlineto{\pgfqpoint{2.738244in}{1.753504in}}%
\pgfpathlineto{\pgfqpoint{2.741189in}{1.748238in}}%
\pgfpathlineto{\pgfqpoint{2.756060in}{1.743595in}}%
\pgfpathlineto{\pgfqpoint{2.775084in}{1.776539in}}%
\pgfpathlineto{\pgfqpoint{2.786967in}{1.781353in}}%
\pgfpathlineto{\pgfqpoint{2.792933in}{1.786497in}}%
\pgfpathlineto{\pgfqpoint{2.804037in}{1.782140in}}%
\pgfpathlineto{\pgfqpoint{2.812223in}{1.770016in}}%
\pgfpathlineto{\pgfqpoint{2.815217in}{1.759142in}}%
\pgfpathclose%
\pgfusepath{fill}%
\end{pgfscope}%
\begin{pgfscope}%
\pgfpathrectangle{\pgfqpoint{0.100000in}{0.100000in}}{\pgfqpoint{3.608454in}{2.310000in}}%
\pgfusepath{clip}%
\pgfsetbuttcap%
\pgfsetmiterjoin%
\definecolor{currentfill}{rgb}{0.000000,0.701961,0.649020}%
\pgfsetfillcolor{currentfill}%
\pgfsetlinewidth{0.000000pt}%
\definecolor{currentstroke}{rgb}{0.000000,0.000000,0.000000}%
\pgfsetstrokecolor{currentstroke}%
\pgfsetstrokeopacity{0.000000}%
\pgfsetdash{}{0pt}%
\pgfpathmoveto{\pgfqpoint{3.050650in}{1.387261in}}%
\pgfpathlineto{\pgfqpoint{3.050299in}{1.374365in}}%
\pgfpathlineto{\pgfqpoint{3.044609in}{1.363268in}}%
\pgfpathlineto{\pgfqpoint{3.046369in}{1.360627in}}%
\pgfpathlineto{\pgfqpoint{3.039437in}{1.352438in}}%
\pgfpathlineto{\pgfqpoint{3.038592in}{1.342913in}}%
\pgfpathlineto{\pgfqpoint{3.026105in}{1.338632in}}%
\pgfpathlineto{\pgfqpoint{3.020093in}{1.338488in}}%
\pgfpathlineto{\pgfqpoint{3.013175in}{1.342782in}}%
\pgfpathlineto{\pgfqpoint{3.011961in}{1.351516in}}%
\pgfpathlineto{\pgfqpoint{3.012394in}{1.360390in}}%
\pgfpathlineto{\pgfqpoint{3.016902in}{1.365346in}}%
\pgfpathlineto{\pgfqpoint{3.018910in}{1.376644in}}%
\pgfpathlineto{\pgfqpoint{3.010884in}{1.388904in}}%
\pgfpathlineto{\pgfqpoint{3.012891in}{1.393149in}}%
\pgfpathlineto{\pgfqpoint{3.019598in}{1.394420in}}%
\pgfpathlineto{\pgfqpoint{3.018667in}{1.406375in}}%
\pgfpathlineto{\pgfqpoint{3.021096in}{1.413057in}}%
\pgfpathlineto{\pgfqpoint{3.008758in}{1.422029in}}%
\pgfpathlineto{\pgfqpoint{3.010773in}{1.432733in}}%
\pgfpathlineto{\pgfqpoint{3.019335in}{1.434599in}}%
\pgfpathlineto{\pgfqpoint{3.025784in}{1.440073in}}%
\pgfpathlineto{\pgfqpoint{3.032463in}{1.435406in}}%
\pgfpathlineto{\pgfqpoint{3.039311in}{1.439843in}}%
\pgfpathlineto{\pgfqpoint{3.052407in}{1.435951in}}%
\pgfpathlineto{\pgfqpoint{3.060760in}{1.438938in}}%
\pgfpathlineto{\pgfqpoint{3.063986in}{1.437446in}}%
\pgfpathlineto{\pgfqpoint{3.062165in}{1.429526in}}%
\pgfpathlineto{\pgfqpoint{3.065468in}{1.428454in}}%
\pgfpathlineto{\pgfqpoint{3.063909in}{1.418899in}}%
\pgfpathlineto{\pgfqpoint{3.055324in}{1.410700in}}%
\pgfpathlineto{\pgfqpoint{3.056652in}{1.404504in}}%
\pgfpathlineto{\pgfqpoint{3.050650in}{1.387261in}}%
\pgfpathclose%
\pgfusepath{fill}%
\end{pgfscope}%
\begin{pgfscope}%
\pgfpathrectangle{\pgfqpoint{0.100000in}{0.100000in}}{\pgfqpoint{3.608454in}{2.310000in}}%
\pgfusepath{clip}%
\pgfsetbuttcap%
\pgfsetmiterjoin%
\definecolor{currentfill}{rgb}{0.000000,0.600000,0.700000}%
\pgfsetfillcolor{currentfill}%
\pgfsetlinewidth{0.000000pt}%
\definecolor{currentstroke}{rgb}{0.000000,0.000000,0.000000}%
\pgfsetstrokecolor{currentstroke}%
\pgfsetstrokeopacity{0.000000}%
\pgfsetdash{}{0pt}%
\pgfpathmoveto{\pgfqpoint{3.217063in}{1.727096in}}%
\pgfpathlineto{\pgfqpoint{3.194930in}{1.721757in}}%
\pgfpathlineto{\pgfqpoint{3.190729in}{1.738354in}}%
\pgfpathlineto{\pgfqpoint{3.167852in}{1.733612in}}%
\pgfpathlineto{\pgfqpoint{3.162874in}{1.739825in}}%
\pgfpathlineto{\pgfqpoint{3.158361in}{1.759713in}}%
\pgfpathlineto{\pgfqpoint{3.158574in}{1.768645in}}%
\pgfpathlineto{\pgfqpoint{3.153298in}{1.788862in}}%
\pgfpathlineto{\pgfqpoint{3.157336in}{1.795352in}}%
\pgfpathlineto{\pgfqpoint{3.167511in}{1.806116in}}%
\pgfpathlineto{\pgfqpoint{3.177989in}{1.809166in}}%
\pgfpathlineto{\pgfqpoint{3.177265in}{1.821081in}}%
\pgfpathlineto{\pgfqpoint{3.186645in}{1.825233in}}%
\pgfpathlineto{\pgfqpoint{3.200994in}{1.826823in}}%
\pgfpathlineto{\pgfqpoint{3.205664in}{1.810106in}}%
\pgfpathlineto{\pgfqpoint{3.219039in}{1.808893in}}%
\pgfpathlineto{\pgfqpoint{3.239480in}{1.829258in}}%
\pgfpathlineto{\pgfqpoint{3.227847in}{1.865817in}}%
\pgfpathlineto{\pgfqpoint{3.234624in}{1.863539in}}%
\pgfpathlineto{\pgfqpoint{3.245888in}{1.867665in}}%
\pgfpathlineto{\pgfqpoint{3.260468in}{1.823521in}}%
\pgfpathlineto{\pgfqpoint{3.257858in}{1.810998in}}%
\pgfpathlineto{\pgfqpoint{3.267521in}{1.808820in}}%
\pgfpathlineto{\pgfqpoint{3.270444in}{1.800501in}}%
\pgfpathlineto{\pgfqpoint{3.267628in}{1.792368in}}%
\pgfpathlineto{\pgfqpoint{3.271017in}{1.786225in}}%
\pgfpathlineto{\pgfqpoint{3.272069in}{1.775491in}}%
\pgfpathlineto{\pgfqpoint{3.264917in}{1.776723in}}%
\pgfpathlineto{\pgfqpoint{3.264610in}{1.770699in}}%
\pgfpathlineto{\pgfqpoint{3.252161in}{1.774686in}}%
\pgfpathlineto{\pgfqpoint{3.246838in}{1.769301in}}%
\pgfpathlineto{\pgfqpoint{3.243936in}{1.759541in}}%
\pgfpathlineto{\pgfqpoint{3.210494in}{1.750645in}}%
\pgfpathlineto{\pgfqpoint{3.217063in}{1.727096in}}%
\pgfpathclose%
\pgfusepath{fill}%
\end{pgfscope}%
\begin{pgfscope}%
\pgfpathrectangle{\pgfqpoint{0.100000in}{0.100000in}}{\pgfqpoint{3.608454in}{2.310000in}}%
\pgfusepath{clip}%
\pgfsetbuttcap%
\pgfsetmiterjoin%
\definecolor{currentfill}{rgb}{0.000000,0.443137,0.778431}%
\pgfsetfillcolor{currentfill}%
\pgfsetlinewidth{0.000000pt}%
\definecolor{currentstroke}{rgb}{0.000000,0.000000,0.000000}%
\pgfsetstrokecolor{currentstroke}%
\pgfsetstrokeopacity{0.000000}%
\pgfsetdash{}{0pt}%
\pgfpathmoveto{\pgfqpoint{2.115090in}{1.593467in}}%
\pgfpathlineto{\pgfqpoint{2.155774in}{1.593709in}}%
\pgfpathlineto{\pgfqpoint{2.156052in}{1.566272in}}%
\pgfpathlineto{\pgfqpoint{2.121832in}{1.566001in}}%
\pgfpathlineto{\pgfqpoint{2.073952in}{1.566162in}}%
\pgfpathlineto{\pgfqpoint{2.067148in}{1.566188in}}%
\pgfpathlineto{\pgfqpoint{2.067449in}{1.593735in}}%
\pgfpathlineto{\pgfqpoint{2.115090in}{1.593467in}}%
\pgfpathclose%
\pgfusepath{fill}%
\end{pgfscope}%
\begin{pgfscope}%
\pgfpathrectangle{\pgfqpoint{0.100000in}{0.100000in}}{\pgfqpoint{3.608454in}{2.310000in}}%
\pgfusepath{clip}%
\pgfsetbuttcap%
\pgfsetmiterjoin%
\definecolor{currentfill}{rgb}{0.000000,0.603922,0.698039}%
\pgfsetfillcolor{currentfill}%
\pgfsetlinewidth{0.000000pt}%
\definecolor{currentstroke}{rgb}{0.000000,0.000000,0.000000}%
\pgfsetstrokecolor{currentstroke}%
\pgfsetstrokeopacity{0.000000}%
\pgfsetdash{}{0pt}%
\pgfpathmoveto{\pgfqpoint{2.647203in}{0.861864in}}%
\pgfpathlineto{\pgfqpoint{2.648743in}{0.852393in}}%
\pgfpathlineto{\pgfqpoint{2.656320in}{0.836856in}}%
\pgfpathlineto{\pgfqpoint{2.650875in}{0.827115in}}%
\pgfpathlineto{\pgfqpoint{2.652230in}{0.813085in}}%
\pgfpathlineto{\pgfqpoint{2.634007in}{0.811449in}}%
\pgfpathlineto{\pgfqpoint{2.622706in}{0.822514in}}%
\pgfpathlineto{\pgfqpoint{2.612796in}{0.827022in}}%
\pgfpathlineto{\pgfqpoint{2.609026in}{0.830093in}}%
\pgfpathlineto{\pgfqpoint{2.607870in}{0.843949in}}%
\pgfpathlineto{\pgfqpoint{2.594215in}{0.842750in}}%
\pgfpathlineto{\pgfqpoint{2.588335in}{0.845764in}}%
\pgfpathlineto{\pgfqpoint{2.589627in}{0.852306in}}%
\pgfpathlineto{\pgfqpoint{2.586111in}{0.864841in}}%
\pgfpathlineto{\pgfqpoint{2.586494in}{0.875360in}}%
\pgfpathlineto{\pgfqpoint{2.591160in}{0.879866in}}%
\pgfpathlineto{\pgfqpoint{2.591982in}{0.885391in}}%
\pgfpathlineto{\pgfqpoint{2.583933in}{0.884712in}}%
\pgfpathlineto{\pgfqpoint{2.582729in}{0.895488in}}%
\pgfpathlineto{\pgfqpoint{2.580198in}{0.924865in}}%
\pgfpathlineto{\pgfqpoint{2.591702in}{0.925524in}}%
\pgfpathlineto{\pgfqpoint{2.593439in}{0.932696in}}%
\pgfpathlineto{\pgfqpoint{2.610666in}{0.933472in}}%
\pgfpathlineto{\pgfqpoint{2.618130in}{0.927031in}}%
\pgfpathlineto{\pgfqpoint{2.617467in}{0.921208in}}%
\pgfpathlineto{\pgfqpoint{2.625197in}{0.912348in}}%
\pgfpathlineto{\pgfqpoint{2.635819in}{0.907320in}}%
\pgfpathlineto{\pgfqpoint{2.636464in}{0.900943in}}%
\pgfpathlineto{\pgfqpoint{2.640666in}{0.895938in}}%
\pgfpathlineto{\pgfqpoint{2.647138in}{0.892648in}}%
\pgfpathlineto{\pgfqpoint{2.648704in}{0.875737in}}%
\pgfpathlineto{\pgfqpoint{2.639182in}{0.874953in}}%
\pgfpathlineto{\pgfqpoint{2.640493in}{0.861198in}}%
\pgfpathlineto{\pgfqpoint{2.647203in}{0.861864in}}%
\pgfpathclose%
\pgfusepath{fill}%
\end{pgfscope}%
\begin{pgfscope}%
\pgfpathrectangle{\pgfqpoint{0.100000in}{0.100000in}}{\pgfqpoint{3.608454in}{2.310000in}}%
\pgfusepath{clip}%
\pgfsetbuttcap%
\pgfsetmiterjoin%
\definecolor{currentfill}{rgb}{0.000000,0.321569,0.839216}%
\pgfsetfillcolor{currentfill}%
\pgfsetlinewidth{0.000000pt}%
\definecolor{currentstroke}{rgb}{0.000000,0.000000,0.000000}%
\pgfsetstrokecolor{currentstroke}%
\pgfsetstrokeopacity{0.000000}%
\pgfsetdash{}{0pt}%
\pgfpathmoveto{\pgfqpoint{2.223956in}{1.594717in}}%
\pgfpathlineto{\pgfqpoint{2.210467in}{1.594411in}}%
\pgfpathlineto{\pgfqpoint{2.196902in}{1.594193in}}%
\pgfpathlineto{\pgfqpoint{2.196575in}{1.614898in}}%
\pgfpathlineto{\pgfqpoint{2.194446in}{1.614873in}}%
\pgfpathlineto{\pgfqpoint{2.193901in}{1.649520in}}%
\pgfpathlineto{\pgfqpoint{2.221382in}{1.649990in}}%
\pgfpathlineto{\pgfqpoint{2.222007in}{1.615392in}}%
\pgfpathlineto{\pgfqpoint{2.223427in}{1.615413in}}%
\pgfpathlineto{\pgfqpoint{2.223956in}{1.594717in}}%
\pgfpathclose%
\pgfusepath{fill}%
\end{pgfscope}%
\begin{pgfscope}%
\pgfpathrectangle{\pgfqpoint{0.100000in}{0.100000in}}{\pgfqpoint{3.608454in}{2.310000in}}%
\pgfusepath{clip}%
\pgfsetbuttcap%
\pgfsetmiterjoin%
\definecolor{currentfill}{rgb}{0.000000,0.694118,0.652941}%
\pgfsetfillcolor{currentfill}%
\pgfsetlinewidth{0.000000pt}%
\definecolor{currentstroke}{rgb}{0.000000,0.000000,0.000000}%
\pgfsetstrokecolor{currentstroke}%
\pgfsetstrokeopacity{0.000000}%
\pgfsetdash{}{0pt}%
\pgfpathmoveto{\pgfqpoint{2.720255in}{1.702142in}}%
\pgfpathlineto{\pgfqpoint{2.720350in}{1.701323in}}%
\pgfpathlineto{\pgfqpoint{2.693396in}{1.698299in}}%
\pgfpathlineto{\pgfqpoint{2.689796in}{1.725591in}}%
\pgfpathlineto{\pgfqpoint{2.676133in}{1.724089in}}%
\pgfpathlineto{\pgfqpoint{2.669984in}{1.779116in}}%
\pgfpathlineto{\pgfqpoint{2.683319in}{1.780357in}}%
\pgfpathlineto{\pgfqpoint{2.697115in}{1.781860in}}%
\pgfpathlineto{\pgfqpoint{2.703428in}{1.727130in}}%
\pgfpathlineto{\pgfqpoint{2.717010in}{1.728719in}}%
\pgfpathlineto{\pgfqpoint{2.720255in}{1.702142in}}%
\pgfpathclose%
\pgfusepath{fill}%
\end{pgfscope}%
\begin{pgfscope}%
\pgfpathrectangle{\pgfqpoint{0.100000in}{0.100000in}}{\pgfqpoint{3.608454in}{2.310000in}}%
\pgfusepath{clip}%
\pgfsetbuttcap%
\pgfsetmiterjoin%
\definecolor{currentfill}{rgb}{0.000000,0.635294,0.682353}%
\pgfsetfillcolor{currentfill}%
\pgfsetlinewidth{0.000000pt}%
\definecolor{currentstroke}{rgb}{0.000000,0.000000,0.000000}%
\pgfsetstrokecolor{currentstroke}%
\pgfsetstrokeopacity{0.000000}%
\pgfsetdash{}{0pt}%
\pgfpathmoveto{\pgfqpoint{2.882548in}{1.076650in}}%
\pgfpathlineto{\pgfqpoint{2.879644in}{1.078954in}}%
\pgfpathlineto{\pgfqpoint{2.885380in}{1.093119in}}%
\pgfpathlineto{\pgfqpoint{2.884882in}{1.098977in}}%
\pgfpathlineto{\pgfqpoint{2.866841in}{1.114054in}}%
\pgfpathlineto{\pgfqpoint{2.864797in}{1.125728in}}%
\pgfpathlineto{\pgfqpoint{2.859932in}{1.127064in}}%
\pgfpathlineto{\pgfqpoint{2.865280in}{1.133278in}}%
\pgfpathlineto{\pgfqpoint{2.877093in}{1.136816in}}%
\pgfpathlineto{\pgfqpoint{2.879988in}{1.149617in}}%
\pgfpathlineto{\pgfqpoint{2.895600in}{1.161365in}}%
\pgfpathlineto{\pgfqpoint{2.898012in}{1.153842in}}%
\pgfpathlineto{\pgfqpoint{2.904596in}{1.155570in}}%
\pgfpathlineto{\pgfqpoint{2.905552in}{1.151887in}}%
\pgfpathlineto{\pgfqpoint{2.916363in}{1.143909in}}%
\pgfpathlineto{\pgfqpoint{2.925369in}{1.125057in}}%
\pgfpathlineto{\pgfqpoint{2.930944in}{1.123301in}}%
\pgfpathlineto{\pgfqpoint{2.926944in}{1.122067in}}%
\pgfpathlineto{\pgfqpoint{2.924776in}{1.115579in}}%
\pgfpathlineto{\pgfqpoint{2.926537in}{1.111891in}}%
\pgfpathlineto{\pgfqpoint{2.922309in}{1.102781in}}%
\pgfpathlineto{\pgfqpoint{2.922867in}{1.095204in}}%
\pgfpathlineto{\pgfqpoint{2.903153in}{1.086612in}}%
\pgfpathlineto{\pgfqpoint{2.882548in}{1.076650in}}%
\pgfpathclose%
\pgfusepath{fill}%
\end{pgfscope}%
\begin{pgfscope}%
\pgfpathrectangle{\pgfqpoint{0.100000in}{0.100000in}}{\pgfqpoint{3.608454in}{2.310000in}}%
\pgfusepath{clip}%
\pgfsetbuttcap%
\pgfsetmiterjoin%
\definecolor{currentfill}{rgb}{0.000000,0.764706,0.617647}%
\pgfsetfillcolor{currentfill}%
\pgfsetlinewidth{0.000000pt}%
\definecolor{currentstroke}{rgb}{0.000000,0.000000,0.000000}%
\pgfsetstrokecolor{currentstroke}%
\pgfsetstrokeopacity{0.000000}%
\pgfsetdash{}{0pt}%
\pgfpathmoveto{\pgfqpoint{2.272396in}{1.102853in}}%
\pgfpathlineto{\pgfqpoint{2.272664in}{1.088030in}}%
\pgfpathlineto{\pgfqpoint{2.258979in}{1.087834in}}%
\pgfpathlineto{\pgfqpoint{2.259040in}{1.082097in}}%
\pgfpathlineto{\pgfqpoint{2.247558in}{1.082128in}}%
\pgfpathlineto{\pgfqpoint{2.201928in}{1.082158in}}%
\pgfpathlineto{\pgfqpoint{2.201915in}{1.084462in}}%
\pgfpathlineto{\pgfqpoint{2.201478in}{1.093804in}}%
\pgfpathlineto{\pgfqpoint{2.206038in}{1.100680in}}%
\pgfpathlineto{\pgfqpoint{2.206212in}{1.109649in}}%
\pgfpathlineto{\pgfqpoint{2.199417in}{1.113130in}}%
\pgfpathlineto{\pgfqpoint{2.194911in}{1.119004in}}%
\pgfpathlineto{\pgfqpoint{2.192391in}{1.127027in}}%
\pgfpathlineto{\pgfqpoint{2.178906in}{1.127188in}}%
\pgfpathlineto{\pgfqpoint{2.178864in}{1.142325in}}%
\pgfpathlineto{\pgfqpoint{2.243255in}{1.143432in}}%
\pgfpathlineto{\pgfqpoint{2.241311in}{1.142821in}}%
\pgfpathlineto{\pgfqpoint{2.241638in}{1.113006in}}%
\pgfpathlineto{\pgfqpoint{2.245103in}{1.109604in}}%
\pgfpathlineto{\pgfqpoint{2.272381in}{1.109733in}}%
\pgfpathlineto{\pgfqpoint{2.272396in}{1.102853in}}%
\pgfpathclose%
\pgfusepath{fill}%
\end{pgfscope}%
\begin{pgfscope}%
\pgfpathrectangle{\pgfqpoint{0.100000in}{0.100000in}}{\pgfqpoint{3.608454in}{2.310000in}}%
\pgfusepath{clip}%
\pgfsetbuttcap%
\pgfsetmiterjoin%
\definecolor{currentfill}{rgb}{0.000000,0.447059,0.776471}%
\pgfsetfillcolor{currentfill}%
\pgfsetlinewidth{0.000000pt}%
\definecolor{currentstroke}{rgb}{0.000000,0.000000,0.000000}%
\pgfsetstrokecolor{currentstroke}%
\pgfsetstrokeopacity{0.000000}%
\pgfsetdash{}{0pt}%
\pgfpathmoveto{\pgfqpoint{2.526558in}{1.632431in}}%
\pgfpathlineto{\pgfqpoint{2.473885in}{1.628654in}}%
\pgfpathlineto{\pgfqpoint{2.435480in}{1.626662in}}%
\pgfpathlineto{\pgfqpoint{2.433615in}{1.653903in}}%
\pgfpathlineto{\pgfqpoint{2.454210in}{1.655405in}}%
\pgfpathlineto{\pgfqpoint{2.467930in}{1.655975in}}%
\pgfpathlineto{\pgfqpoint{2.522980in}{1.660007in}}%
\pgfpathlineto{\pgfqpoint{2.527343in}{1.655339in}}%
\pgfpathlineto{\pgfqpoint{2.524832in}{1.642175in}}%
\pgfpathlineto{\pgfqpoint{2.526558in}{1.632431in}}%
\pgfpathclose%
\pgfusepath{fill}%
\end{pgfscope}%
\begin{pgfscope}%
\pgfpathrectangle{\pgfqpoint{0.100000in}{0.100000in}}{\pgfqpoint{3.608454in}{2.310000in}}%
\pgfusepath{clip}%
\pgfsetbuttcap%
\pgfsetmiterjoin%
\definecolor{currentfill}{rgb}{0.000000,0.501961,0.749020}%
\pgfsetfillcolor{currentfill}%
\pgfsetlinewidth{0.000000pt}%
\definecolor{currentstroke}{rgb}{0.000000,0.000000,0.000000}%
\pgfsetstrokecolor{currentstroke}%
\pgfsetstrokeopacity{0.000000}%
\pgfsetdash{}{0pt}%
\pgfpathmoveto{\pgfqpoint{1.774163in}{0.940149in}}%
\pgfpathlineto{\pgfqpoint{1.743217in}{0.941925in}}%
\pgfpathlineto{\pgfqpoint{1.745610in}{0.979809in}}%
\pgfpathlineto{\pgfqpoint{1.752117in}{0.979371in}}%
\pgfpathlineto{\pgfqpoint{1.754289in}{1.013713in}}%
\pgfpathlineto{\pgfqpoint{1.781300in}{1.012079in}}%
\pgfpathlineto{\pgfqpoint{1.780677in}{0.997355in}}%
\pgfpathlineto{\pgfqpoint{1.779205in}{0.970733in}}%
\pgfpathlineto{\pgfqpoint{1.775993in}{0.971342in}}%
\pgfpathlineto{\pgfqpoint{1.774163in}{0.940149in}}%
\pgfpathclose%
\pgfusepath{fill}%
\end{pgfscope}%
\begin{pgfscope}%
\pgfpathrectangle{\pgfqpoint{0.100000in}{0.100000in}}{\pgfqpoint{3.608454in}{2.310000in}}%
\pgfusepath{clip}%
\pgfsetbuttcap%
\pgfsetmiterjoin%
\definecolor{currentfill}{rgb}{0.000000,0.811765,0.594118}%
\pgfsetfillcolor{currentfill}%
\pgfsetlinewidth{0.000000pt}%
\definecolor{currentstroke}{rgb}{0.000000,0.000000,0.000000}%
\pgfsetstrokecolor{currentstroke}%
\pgfsetstrokeopacity{0.000000}%
\pgfsetdash{}{0pt}%
\pgfpathmoveto{\pgfqpoint{2.338322in}{0.779971in}}%
\pgfpathlineto{\pgfqpoint{2.338396in}{0.773253in}}%
\pgfpathlineto{\pgfqpoint{2.334485in}{0.767523in}}%
\pgfpathlineto{\pgfqpoint{2.330853in}{0.768113in}}%
\pgfpathlineto{\pgfqpoint{2.325362in}{0.751678in}}%
\pgfpathlineto{\pgfqpoint{2.319974in}{0.746106in}}%
\pgfpathlineto{\pgfqpoint{2.322557in}{0.730252in}}%
\pgfpathlineto{\pgfqpoint{2.311015in}{0.733955in}}%
\pgfpathlineto{\pgfqpoint{2.300756in}{0.740640in}}%
\pgfpathlineto{\pgfqpoint{2.299838in}{0.745415in}}%
\pgfpathlineto{\pgfqpoint{2.291967in}{0.750029in}}%
\pgfpathlineto{\pgfqpoint{2.286412in}{0.757594in}}%
\pgfpathlineto{\pgfqpoint{2.285332in}{0.771355in}}%
\pgfpathlineto{\pgfqpoint{2.288408in}{0.781863in}}%
\pgfpathlineto{\pgfqpoint{2.315987in}{0.782756in}}%
\pgfpathlineto{\pgfqpoint{2.319400in}{0.785321in}}%
\pgfpathlineto{\pgfqpoint{2.324523in}{0.778629in}}%
\pgfpathlineto{\pgfqpoint{2.336629in}{0.786857in}}%
\pgfpathlineto{\pgfqpoint{2.338322in}{0.779971in}}%
\pgfpathclose%
\pgfusepath{fill}%
\end{pgfscope}%
\begin{pgfscope}%
\pgfpathrectangle{\pgfqpoint{0.100000in}{0.100000in}}{\pgfqpoint{3.608454in}{2.310000in}}%
\pgfusepath{clip}%
\pgfsetbuttcap%
\pgfsetmiterjoin%
\definecolor{currentfill}{rgb}{0.000000,0.654902,0.672549}%
\pgfsetfillcolor{currentfill}%
\pgfsetlinewidth{0.000000pt}%
\definecolor{currentstroke}{rgb}{0.000000,0.000000,0.000000}%
\pgfsetstrokecolor{currentstroke}%
\pgfsetstrokeopacity{0.000000}%
\pgfsetdash{}{0pt}%
\pgfpathmoveto{\pgfqpoint{2.618130in}{0.927031in}}%
\pgfpathlineto{\pgfqpoint{2.610666in}{0.933472in}}%
\pgfpathlineto{\pgfqpoint{2.593439in}{0.932696in}}%
\pgfpathlineto{\pgfqpoint{2.591702in}{0.925524in}}%
\pgfpathlineto{\pgfqpoint{2.580198in}{0.924865in}}%
\pgfpathlineto{\pgfqpoint{2.551506in}{0.923399in}}%
\pgfpathlineto{\pgfqpoint{2.551959in}{0.940228in}}%
\pgfpathlineto{\pgfqpoint{2.553484in}{0.986177in}}%
\pgfpathlineto{\pgfqpoint{2.595636in}{0.988190in}}%
\pgfpathlineto{\pgfqpoint{2.623076in}{0.990093in}}%
\pgfpathlineto{\pgfqpoint{2.625447in}{0.963052in}}%
\pgfpathlineto{\pgfqpoint{2.627618in}{0.959226in}}%
\pgfpathlineto{\pgfqpoint{2.635731in}{0.956087in}}%
\pgfpathlineto{\pgfqpoint{2.632392in}{0.948970in}}%
\pgfpathlineto{\pgfqpoint{2.631914in}{0.941959in}}%
\pgfpathlineto{\pgfqpoint{2.628827in}{0.937113in}}%
\pgfpathlineto{\pgfqpoint{2.623213in}{0.935505in}}%
\pgfpathlineto{\pgfqpoint{2.618130in}{0.927031in}}%
\pgfpathclose%
\pgfusepath{fill}%
\end{pgfscope}%
\begin{pgfscope}%
\pgfpathrectangle{\pgfqpoint{0.100000in}{0.100000in}}{\pgfqpoint{3.608454in}{2.310000in}}%
\pgfusepath{clip}%
\pgfsetbuttcap%
\pgfsetmiterjoin%
\definecolor{currentfill}{rgb}{0.000000,0.525490,0.737255}%
\pgfsetfillcolor{currentfill}%
\pgfsetlinewidth{0.000000pt}%
\definecolor{currentstroke}{rgb}{0.000000,0.000000,0.000000}%
\pgfsetstrokecolor{currentstroke}%
\pgfsetstrokeopacity{0.000000}%
\pgfsetdash{}{0pt}%
\pgfpathmoveto{\pgfqpoint{2.949055in}{1.160645in}}%
\pgfpathlineto{\pgfqpoint{2.956204in}{1.154739in}}%
\pgfpathlineto{\pgfqpoint{2.967129in}{1.148531in}}%
\pgfpathlineto{\pgfqpoint{2.973998in}{1.150189in}}%
\pgfpathlineto{\pgfqpoint{2.978683in}{1.148660in}}%
\pgfpathlineto{\pgfqpoint{2.970938in}{1.132406in}}%
\pgfpathlineto{\pgfqpoint{2.963707in}{1.130755in}}%
\pgfpathlineto{\pgfqpoint{2.956399in}{1.133737in}}%
\pgfpathlineto{\pgfqpoint{2.951598in}{1.127420in}}%
\pgfpathlineto{\pgfqpoint{2.943450in}{1.124830in}}%
\pgfpathlineto{\pgfqpoint{2.930944in}{1.123301in}}%
\pgfpathlineto{\pgfqpoint{2.925369in}{1.125057in}}%
\pgfpathlineto{\pgfqpoint{2.916363in}{1.143909in}}%
\pgfpathlineto{\pgfqpoint{2.905552in}{1.151887in}}%
\pgfpathlineto{\pgfqpoint{2.904596in}{1.155570in}}%
\pgfpathlineto{\pgfqpoint{2.909797in}{1.165077in}}%
\pgfpathlineto{\pgfqpoint{2.920582in}{1.172131in}}%
\pgfpathlineto{\pgfqpoint{2.930144in}{1.169183in}}%
\pgfpathlineto{\pgfqpoint{2.934882in}{1.156995in}}%
\pgfpathlineto{\pgfqpoint{2.938556in}{1.155078in}}%
\pgfpathlineto{\pgfqpoint{2.942760in}{1.162773in}}%
\pgfpathlineto{\pgfqpoint{2.949055in}{1.160645in}}%
\pgfpathclose%
\pgfusepath{fill}%
\end{pgfscope}%
\begin{pgfscope}%
\pgfpathrectangle{\pgfqpoint{0.100000in}{0.100000in}}{\pgfqpoint{3.608454in}{2.310000in}}%
\pgfusepath{clip}%
\pgfsetbuttcap%
\pgfsetmiterjoin%
\definecolor{currentfill}{rgb}{0.000000,0.603922,0.698039}%
\pgfsetfillcolor{currentfill}%
\pgfsetlinewidth{0.000000pt}%
\definecolor{currentstroke}{rgb}{0.000000,0.000000,0.000000}%
\pgfsetstrokecolor{currentstroke}%
\pgfsetstrokeopacity{0.000000}%
\pgfsetdash{}{0pt}%
\pgfpathmoveto{\pgfqpoint{0.650466in}{2.312218in}}%
\pgfpathlineto{\pgfqpoint{0.640068in}{2.318805in}}%
\pgfpathlineto{\pgfqpoint{0.635270in}{2.326403in}}%
\pgfpathlineto{\pgfqpoint{0.636669in}{2.333571in}}%
\pgfpathlineto{\pgfqpoint{0.643160in}{2.329044in}}%
\pgfpathlineto{\pgfqpoint{0.650928in}{2.336446in}}%
\pgfpathlineto{\pgfqpoint{0.659748in}{2.329464in}}%
\pgfpathlineto{\pgfqpoint{0.650466in}{2.312218in}}%
\pgfpathclose%
\pgfusepath{fill}%
\end{pgfscope}%
\begin{pgfscope}%
\pgfpathrectangle{\pgfqpoint{0.100000in}{0.100000in}}{\pgfqpoint{3.608454in}{2.310000in}}%
\pgfusepath{clip}%
\pgfsetbuttcap%
\pgfsetmiterjoin%
\definecolor{currentfill}{rgb}{0.000000,0.556863,0.721569}%
\pgfsetfillcolor{currentfill}%
\pgfsetlinewidth{0.000000pt}%
\definecolor{currentstroke}{rgb}{0.000000,0.000000,0.000000}%
\pgfsetstrokecolor{currentstroke}%
\pgfsetstrokeopacity{0.000000}%
\pgfsetdash{}{0pt}%
\pgfpathmoveto{\pgfqpoint{1.644176in}{0.728867in}}%
\pgfpathlineto{\pgfqpoint{1.678348in}{0.726376in}}%
\pgfpathlineto{\pgfqpoint{1.699678in}{0.725186in}}%
\pgfpathlineto{\pgfqpoint{1.697817in}{0.695195in}}%
\pgfpathlineto{\pgfqpoint{1.695622in}{0.661844in}}%
\pgfpathlineto{\pgfqpoint{1.640690in}{0.665716in}}%
\pgfpathlineto{\pgfqpoint{1.647908in}{0.670831in}}%
\pgfpathlineto{\pgfqpoint{1.642706in}{0.676280in}}%
\pgfpathlineto{\pgfqpoint{1.649534in}{0.685985in}}%
\pgfpathlineto{\pgfqpoint{1.648797in}{0.693179in}}%
\pgfpathlineto{\pgfqpoint{1.644365in}{0.694777in}}%
\pgfpathlineto{\pgfqpoint{1.616885in}{0.697027in}}%
\pgfpathlineto{\pgfqpoint{1.616475in}{0.692376in}}%
\pgfpathlineto{\pgfqpoint{1.602461in}{0.693594in}}%
\pgfpathlineto{\pgfqpoint{1.600425in}{0.668531in}}%
\pgfpathlineto{\pgfqpoint{1.585018in}{0.669731in}}%
\pgfpathlineto{\pgfqpoint{1.583492in}{0.651351in}}%
\pgfpathlineto{\pgfqpoint{1.518651in}{0.714404in}}%
\pgfpathlineto{\pgfqpoint{1.562114in}{0.758888in}}%
\pgfpathlineto{\pgfqpoint{1.567213in}{0.756734in}}%
\pgfpathlineto{\pgfqpoint{1.574315in}{0.749565in}}%
\pgfpathlineto{\pgfqpoint{1.578119in}{0.751288in}}%
\pgfpathlineto{\pgfqpoint{1.584520in}{0.753390in}}%
\pgfpathlineto{\pgfqpoint{1.599669in}{0.742990in}}%
\pgfpathlineto{\pgfqpoint{1.601306in}{0.732857in}}%
\pgfpathlineto{\pgfqpoint{1.644176in}{0.728867in}}%
\pgfpathclose%
\pgfusepath{fill}%
\end{pgfscope}%
\begin{pgfscope}%
\pgfpathrectangle{\pgfqpoint{0.100000in}{0.100000in}}{\pgfqpoint{3.608454in}{2.310000in}}%
\pgfusepath{clip}%
\pgfsetbuttcap%
\pgfsetmiterjoin%
\definecolor{currentfill}{rgb}{0.000000,0.745098,0.627451}%
\pgfsetfillcolor{currentfill}%
\pgfsetlinewidth{0.000000pt}%
\definecolor{currentstroke}{rgb}{0.000000,0.000000,0.000000}%
\pgfsetstrokecolor{currentstroke}%
\pgfsetstrokeopacity{0.000000}%
\pgfsetdash{}{0pt}%
\pgfpathmoveto{\pgfqpoint{2.297141in}{0.880461in}}%
\pgfpathlineto{\pgfqpoint{2.299404in}{0.870120in}}%
\pgfpathlineto{\pgfqpoint{2.302086in}{0.868312in}}%
\pgfpathlineto{\pgfqpoint{2.258381in}{0.867551in}}%
\pgfpathlineto{\pgfqpoint{2.240807in}{0.867381in}}%
\pgfpathlineto{\pgfqpoint{2.240853in}{0.895980in}}%
\pgfpathlineto{\pgfqpoint{2.231529in}{0.896018in}}%
\pgfpathlineto{\pgfqpoint{2.231585in}{0.901788in}}%
\pgfpathlineto{\pgfqpoint{2.231849in}{0.927551in}}%
\pgfpathlineto{\pgfqpoint{2.241965in}{0.929323in}}%
\pgfpathlineto{\pgfqpoint{2.245189in}{0.934157in}}%
\pgfpathlineto{\pgfqpoint{2.245528in}{0.957904in}}%
\pgfpathlineto{\pgfqpoint{2.259184in}{0.957741in}}%
\pgfpathlineto{\pgfqpoint{2.271324in}{0.957636in}}%
\pgfpathlineto{\pgfqpoint{2.281822in}{0.951221in}}%
\pgfpathlineto{\pgfqpoint{2.272954in}{0.950762in}}%
\pgfpathlineto{\pgfqpoint{2.274161in}{0.944083in}}%
\pgfpathlineto{\pgfqpoint{2.282183in}{0.933447in}}%
\pgfpathlineto{\pgfqpoint{2.284812in}{0.906848in}}%
\pgfpathlineto{\pgfqpoint{2.280766in}{0.896956in}}%
\pgfpathlineto{\pgfqpoint{2.282169in}{0.889435in}}%
\pgfpathlineto{\pgfqpoint{2.286293in}{0.889498in}}%
\pgfpathlineto{\pgfqpoint{2.297141in}{0.880461in}}%
\pgfpathclose%
\pgfusepath{fill}%
\end{pgfscope}%
\begin{pgfscope}%
\pgfpathrectangle{\pgfqpoint{0.100000in}{0.100000in}}{\pgfqpoint{3.608454in}{2.310000in}}%
\pgfusepath{clip}%
\pgfsetbuttcap%
\pgfsetmiterjoin%
\definecolor{currentfill}{rgb}{0.000000,0.521569,0.739216}%
\pgfsetfillcolor{currentfill}%
\pgfsetlinewidth{0.000000pt}%
\definecolor{currentstroke}{rgb}{0.000000,0.000000,0.000000}%
\pgfsetstrokecolor{currentstroke}%
\pgfsetstrokeopacity{0.000000}%
\pgfsetdash{}{0pt}%
\pgfpathmoveto{\pgfqpoint{1.795374in}{1.510393in}}%
\pgfpathlineto{\pgfqpoint{1.793857in}{1.483052in}}%
\pgfpathlineto{\pgfqpoint{1.760676in}{1.485050in}}%
\pgfpathlineto{\pgfqpoint{1.732544in}{1.486727in}}%
\pgfpathlineto{\pgfqpoint{1.734320in}{1.514317in}}%
\pgfpathlineto{\pgfqpoint{1.733109in}{1.514401in}}%
\pgfpathlineto{\pgfqpoint{1.735045in}{1.541820in}}%
\pgfpathlineto{\pgfqpoint{1.726904in}{1.542427in}}%
\pgfpathlineto{\pgfqpoint{1.728889in}{1.569871in}}%
\pgfpathlineto{\pgfqpoint{1.762096in}{1.567407in}}%
\pgfpathlineto{\pgfqpoint{1.796934in}{1.565224in}}%
\pgfpathlineto{\pgfqpoint{1.795374in}{1.510393in}}%
\pgfpathclose%
\pgfusepath{fill}%
\end{pgfscope}%
\begin{pgfscope}%
\pgfpathrectangle{\pgfqpoint{0.100000in}{0.100000in}}{\pgfqpoint{3.608454in}{2.310000in}}%
\pgfusepath{clip}%
\pgfsetbuttcap%
\pgfsetmiterjoin%
\definecolor{currentfill}{rgb}{0.000000,0.513725,0.743137}%
\pgfsetfillcolor{currentfill}%
\pgfsetlinewidth{0.000000pt}%
\definecolor{currentstroke}{rgb}{0.000000,0.000000,0.000000}%
\pgfsetstrokecolor{currentstroke}%
\pgfsetstrokeopacity{0.000000}%
\pgfsetdash{}{0pt}%
\pgfpathmoveto{\pgfqpoint{2.794880in}{1.439319in}}%
\pgfpathlineto{\pgfqpoint{2.795500in}{1.426123in}}%
\pgfpathlineto{\pgfqpoint{2.776762in}{1.425313in}}%
\pgfpathlineto{\pgfqpoint{2.757924in}{1.424641in}}%
\pgfpathlineto{\pgfqpoint{2.746401in}{1.421668in}}%
\pgfpathlineto{\pgfqpoint{2.726071in}{1.419220in}}%
\pgfpathlineto{\pgfqpoint{2.722450in}{1.453627in}}%
\pgfpathlineto{\pgfqpoint{2.720165in}{1.477649in}}%
\pgfpathlineto{\pgfqpoint{2.719775in}{1.481004in}}%
\pgfpathlineto{\pgfqpoint{2.741875in}{1.483652in}}%
\pgfpathlineto{\pgfqpoint{2.747392in}{1.486269in}}%
\pgfpathlineto{\pgfqpoint{2.746418in}{1.494296in}}%
\pgfpathlineto{\pgfqpoint{2.767029in}{1.496892in}}%
\pgfpathlineto{\pgfqpoint{2.766536in}{1.500999in}}%
\pgfpathlineto{\pgfqpoint{2.773279in}{1.502104in}}%
\pgfpathlineto{\pgfqpoint{2.795078in}{1.502065in}}%
\pgfpathlineto{\pgfqpoint{2.795940in}{1.480285in}}%
\pgfpathlineto{\pgfqpoint{2.799407in}{1.480427in}}%
\pgfpathlineto{\pgfqpoint{2.800713in}{1.463456in}}%
\pgfpathlineto{\pgfqpoint{2.798306in}{1.443885in}}%
\pgfpathlineto{\pgfqpoint{2.794880in}{1.439319in}}%
\pgfpathclose%
\pgfusepath{fill}%
\end{pgfscope}%
\begin{pgfscope}%
\pgfpathrectangle{\pgfqpoint{0.100000in}{0.100000in}}{\pgfqpoint{3.608454in}{2.310000in}}%
\pgfusepath{clip}%
\pgfsetbuttcap%
\pgfsetmiterjoin%
\definecolor{currentfill}{rgb}{0.000000,0.584314,0.707843}%
\pgfsetfillcolor{currentfill}%
\pgfsetlinewidth{0.000000pt}%
\definecolor{currentstroke}{rgb}{0.000000,0.000000,0.000000}%
\pgfsetstrokecolor{currentstroke}%
\pgfsetstrokeopacity{0.000000}%
\pgfsetdash{}{0pt}%
\pgfpathmoveto{\pgfqpoint{2.290718in}{1.442111in}}%
\pgfpathlineto{\pgfqpoint{2.291417in}{1.417786in}}%
\pgfpathlineto{\pgfqpoint{2.267659in}{1.417346in}}%
\pgfpathlineto{\pgfqpoint{2.267438in}{1.423982in}}%
\pgfpathlineto{\pgfqpoint{2.237055in}{1.423181in}}%
\pgfpathlineto{\pgfqpoint{2.236527in}{1.423164in}}%
\pgfpathlineto{\pgfqpoint{2.235992in}{1.447300in}}%
\pgfpathlineto{\pgfqpoint{2.246319in}{1.447619in}}%
\pgfpathlineto{\pgfqpoint{2.246988in}{1.453136in}}%
\pgfpathlineto{\pgfqpoint{2.244011in}{1.467014in}}%
\pgfpathlineto{\pgfqpoint{2.248598in}{1.467244in}}%
\pgfpathlineto{\pgfqpoint{2.290253in}{1.469670in}}%
\pgfpathlineto{\pgfqpoint{2.290718in}{1.442111in}}%
\pgfpathclose%
\pgfusepath{fill}%
\end{pgfscope}%
\begin{pgfscope}%
\pgfpathrectangle{\pgfqpoint{0.100000in}{0.100000in}}{\pgfqpoint{3.608454in}{2.310000in}}%
\pgfusepath{clip}%
\pgfsetbuttcap%
\pgfsetmiterjoin%
\definecolor{currentfill}{rgb}{0.000000,0.709804,0.645098}%
\pgfsetfillcolor{currentfill}%
\pgfsetlinewidth{0.000000pt}%
\definecolor{currentstroke}{rgb}{0.000000,0.000000,0.000000}%
\pgfsetstrokecolor{currentstroke}%
\pgfsetstrokeopacity{0.000000}%
\pgfsetdash{}{0pt}%
\pgfpathmoveto{\pgfqpoint{2.190786in}{1.003624in}}%
\pgfpathlineto{\pgfqpoint{2.180404in}{1.000901in}}%
\pgfpathlineto{\pgfqpoint{2.176519in}{0.997121in}}%
\pgfpathlineto{\pgfqpoint{2.170781in}{0.999186in}}%
\pgfpathlineto{\pgfqpoint{2.151260in}{0.999398in}}%
\pgfpathlineto{\pgfqpoint{2.142324in}{1.001839in}}%
\pgfpathlineto{\pgfqpoint{2.141998in}{0.984277in}}%
\pgfpathlineto{\pgfqpoint{2.110869in}{0.984129in}}%
\pgfpathlineto{\pgfqpoint{2.110850in}{0.997922in}}%
\pgfpathlineto{\pgfqpoint{2.073255in}{0.998037in}}%
\pgfpathlineto{\pgfqpoint{2.073230in}{0.991141in}}%
\pgfpathlineto{\pgfqpoint{2.062941in}{0.991185in}}%
\pgfpathlineto{\pgfqpoint{2.049224in}{0.991276in}}%
\pgfpathlineto{\pgfqpoint{2.049296in}{0.998169in}}%
\pgfpathlineto{\pgfqpoint{2.035603in}{0.998326in}}%
\pgfpathlineto{\pgfqpoint{2.036171in}{1.027587in}}%
\pgfpathlineto{\pgfqpoint{2.043077in}{1.035565in}}%
\pgfpathlineto{\pgfqpoint{2.050477in}{1.038266in}}%
\pgfpathlineto{\pgfqpoint{2.058740in}{1.038661in}}%
\pgfpathlineto{\pgfqpoint{2.069896in}{1.042778in}}%
\pgfpathlineto{\pgfqpoint{2.081865in}{1.049099in}}%
\pgfpathlineto{\pgfqpoint{2.091976in}{1.044393in}}%
\pgfpathlineto{\pgfqpoint{2.096410in}{1.052482in}}%
\pgfpathlineto{\pgfqpoint{2.101602in}{1.057429in}}%
\pgfpathlineto{\pgfqpoint{2.098258in}{1.064900in}}%
\pgfpathlineto{\pgfqpoint{2.098590in}{1.073824in}}%
\pgfpathlineto{\pgfqpoint{2.140727in}{1.073891in}}%
\pgfpathlineto{\pgfqpoint{2.139356in}{1.083433in}}%
\pgfpathlineto{\pgfqpoint{2.163020in}{1.082921in}}%
\pgfpathlineto{\pgfqpoint{2.176762in}{1.083808in}}%
\pgfpathlineto{\pgfqpoint{2.172137in}{1.075918in}}%
\pgfpathlineto{\pgfqpoint{2.166395in}{1.075969in}}%
\pgfpathlineto{\pgfqpoint{2.166287in}{1.069101in}}%
\pgfpathlineto{\pgfqpoint{2.169375in}{1.063715in}}%
\pgfpathlineto{\pgfqpoint{2.166477in}{1.058801in}}%
\pgfpathlineto{\pgfqpoint{2.166153in}{1.040540in}}%
\pgfpathlineto{\pgfqpoint{2.162548in}{1.034836in}}%
\pgfpathlineto{\pgfqpoint{2.162454in}{1.027888in}}%
\pgfpathlineto{\pgfqpoint{2.166850in}{1.025625in}}%
\pgfpathlineto{\pgfqpoint{2.190828in}{1.025374in}}%
\pgfpathlineto{\pgfqpoint{2.190786in}{1.003624in}}%
\pgfpathclose%
\pgfusepath{fill}%
\end{pgfscope}%
\begin{pgfscope}%
\pgfpathrectangle{\pgfqpoint{0.100000in}{0.100000in}}{\pgfqpoint{3.608454in}{2.310000in}}%
\pgfusepath{clip}%
\pgfsetbuttcap%
\pgfsetmiterjoin%
\definecolor{currentfill}{rgb}{0.000000,0.623529,0.688235}%
\pgfsetfillcolor{currentfill}%
\pgfsetlinewidth{0.000000pt}%
\definecolor{currentstroke}{rgb}{0.000000,0.000000,0.000000}%
\pgfsetstrokecolor{currentstroke}%
\pgfsetstrokeopacity{0.000000}%
\pgfsetdash{}{0pt}%
\pgfpathmoveto{\pgfqpoint{3.573009in}{1.957078in}}%
\pgfpathlineto{\pgfqpoint{3.580481in}{1.956181in}}%
\pgfpathlineto{\pgfqpoint{3.577419in}{1.950172in}}%
\pgfpathlineto{\pgfqpoint{3.573009in}{1.957078in}}%
\pgfpathclose%
\pgfusepath{fill}%
\end{pgfscope}%
\begin{pgfscope}%
\pgfpathrectangle{\pgfqpoint{0.100000in}{0.100000in}}{\pgfqpoint{3.608454in}{2.310000in}}%
\pgfusepath{clip}%
\pgfsetbuttcap%
\pgfsetmiterjoin%
\definecolor{currentfill}{rgb}{0.000000,0.623529,0.688235}%
\pgfsetfillcolor{currentfill}%
\pgfsetlinewidth{0.000000pt}%
\definecolor{currentstroke}{rgb}{0.000000,0.000000,0.000000}%
\pgfsetstrokecolor{currentstroke}%
\pgfsetstrokeopacity{0.000000}%
\pgfsetdash{}{0pt}%
\pgfpathmoveto{\pgfqpoint{3.550714in}{1.937007in}}%
\pgfpathlineto{\pgfqpoint{3.553924in}{1.946267in}}%
\pgfpathlineto{\pgfqpoint{3.547311in}{1.956692in}}%
\pgfpathlineto{\pgfqpoint{3.543555in}{1.955713in}}%
\pgfpathlineto{\pgfqpoint{3.538805in}{1.962876in}}%
\pgfpathlineto{\pgfqpoint{3.535140in}{1.964201in}}%
\pgfpathlineto{\pgfqpoint{3.535820in}{1.974394in}}%
\pgfpathlineto{\pgfqpoint{3.534775in}{1.988683in}}%
\pgfpathlineto{\pgfqpoint{3.532167in}{1.996604in}}%
\pgfpathlineto{\pgfqpoint{3.539818in}{1.996897in}}%
\pgfpathlineto{\pgfqpoint{3.527215in}{2.022326in}}%
\pgfpathlineto{\pgfqpoint{3.514278in}{2.013412in}}%
\pgfpathlineto{\pgfqpoint{3.494071in}{2.051630in}}%
\pgfpathlineto{\pgfqpoint{3.496496in}{2.059377in}}%
\pgfpathlineto{\pgfqpoint{3.490714in}{2.062918in}}%
\pgfpathlineto{\pgfqpoint{3.491275in}{2.076574in}}%
\pgfpathlineto{\pgfqpoint{3.487328in}{2.085547in}}%
\pgfpathlineto{\pgfqpoint{3.478067in}{2.117502in}}%
\pgfpathlineto{\pgfqpoint{3.474712in}{2.131510in}}%
\pgfpathlineto{\pgfqpoint{3.522068in}{2.145170in}}%
\pgfpathlineto{\pgfqpoint{3.526210in}{2.131697in}}%
\pgfpathlineto{\pgfqpoint{3.546699in}{2.136666in}}%
\pgfpathlineto{\pgfqpoint{3.557203in}{2.103454in}}%
\pgfpathlineto{\pgfqpoint{3.565274in}{2.075679in}}%
\pgfpathlineto{\pgfqpoint{3.566650in}{2.075414in}}%
\pgfpathlineto{\pgfqpoint{3.584395in}{2.086358in}}%
\pgfpathlineto{\pgfqpoint{3.598893in}{2.059799in}}%
\pgfpathlineto{\pgfqpoint{3.592928in}{2.056963in}}%
\pgfpathlineto{\pgfqpoint{3.603808in}{2.035391in}}%
\pgfpathlineto{\pgfqpoint{3.597567in}{2.032110in}}%
\pgfpathlineto{\pgfqpoint{3.610482in}{2.006964in}}%
\pgfpathlineto{\pgfqpoint{3.618035in}{1.994535in}}%
\pgfpathlineto{\pgfqpoint{3.613315in}{1.987169in}}%
\pgfpathlineto{\pgfqpoint{3.607145in}{1.996865in}}%
\pgfpathlineto{\pgfqpoint{3.599803in}{1.992874in}}%
\pgfpathlineto{\pgfqpoint{3.608045in}{1.984589in}}%
\pgfpathlineto{\pgfqpoint{3.601528in}{1.973738in}}%
\pgfpathlineto{\pgfqpoint{3.594431in}{1.979179in}}%
\pgfpathlineto{\pgfqpoint{3.594740in}{1.984518in}}%
\pgfpathlineto{\pgfqpoint{3.589736in}{1.989331in}}%
\pgfpathlineto{\pgfqpoint{3.587038in}{1.983926in}}%
\pgfpathlineto{\pgfqpoint{3.586028in}{1.973437in}}%
\pgfpathlineto{\pgfqpoint{3.576059in}{1.975363in}}%
\pgfpathlineto{\pgfqpoint{3.564314in}{1.981978in}}%
\pgfpathlineto{\pgfqpoint{3.561572in}{1.978778in}}%
\pgfpathlineto{\pgfqpoint{3.565543in}{1.970193in}}%
\pgfpathlineto{\pgfqpoint{3.562497in}{1.956796in}}%
\pgfpathlineto{\pgfqpoint{3.565408in}{1.948473in}}%
\pgfpathlineto{\pgfqpoint{3.560888in}{1.939136in}}%
\pgfpathlineto{\pgfqpoint{3.550714in}{1.937007in}}%
\pgfpathclose%
\pgfusepath{fill}%
\end{pgfscope}%
\begin{pgfscope}%
\pgfpathrectangle{\pgfqpoint{0.100000in}{0.100000in}}{\pgfqpoint{3.608454in}{2.310000in}}%
\pgfusepath{clip}%
\pgfsetbuttcap%
\pgfsetmiterjoin%
\definecolor{currentfill}{rgb}{0.000000,0.564706,0.717647}%
\pgfsetfillcolor{currentfill}%
\pgfsetlinewidth{0.000000pt}%
\definecolor{currentstroke}{rgb}{0.000000,0.000000,0.000000}%
\pgfsetstrokecolor{currentstroke}%
\pgfsetstrokeopacity{0.000000}%
\pgfsetdash{}{0pt}%
\pgfpathmoveto{\pgfqpoint{1.859593in}{0.670192in}}%
\pgfpathlineto{\pgfqpoint{1.833962in}{0.671147in}}%
\pgfpathlineto{\pgfqpoint{1.835408in}{0.704889in}}%
\pgfpathlineto{\pgfqpoint{1.870905in}{0.703446in}}%
\pgfpathlineto{\pgfqpoint{1.871649in}{0.712052in}}%
\pgfpathlineto{\pgfqpoint{1.907744in}{0.711220in}}%
\pgfpathlineto{\pgfqpoint{1.913119in}{0.701138in}}%
\pgfpathlineto{\pgfqpoint{1.903598in}{0.691454in}}%
\pgfpathlineto{\pgfqpoint{1.892292in}{0.669010in}}%
\pgfpathlineto{\pgfqpoint{1.891457in}{0.663121in}}%
\pgfpathlineto{\pgfqpoint{1.871994in}{0.669892in}}%
\pgfpathlineto{\pgfqpoint{1.859593in}{0.670192in}}%
\pgfpathclose%
\pgfusepath{fill}%
\end{pgfscope}%
\begin{pgfscope}%
\pgfpathrectangle{\pgfqpoint{0.100000in}{0.100000in}}{\pgfqpoint{3.608454in}{2.310000in}}%
\pgfusepath{clip}%
\pgfsetbuttcap%
\pgfsetmiterjoin%
\definecolor{currentfill}{rgb}{0.000000,0.803922,0.598039}%
\pgfsetfillcolor{currentfill}%
\pgfsetlinewidth{0.000000pt}%
\definecolor{currentstroke}{rgb}{0.000000,0.000000,0.000000}%
\pgfsetstrokecolor{currentstroke}%
\pgfsetstrokeopacity{0.000000}%
\pgfsetdash{}{0pt}%
\pgfpathmoveto{\pgfqpoint{2.362240in}{0.870230in}}%
\pgfpathlineto{\pgfqpoint{2.367614in}{0.868244in}}%
\pgfpathlineto{\pgfqpoint{2.368994in}{0.861333in}}%
\pgfpathlineto{\pgfqpoint{2.363366in}{0.855813in}}%
\pgfpathlineto{\pgfqpoint{2.363126in}{0.851002in}}%
\pgfpathlineto{\pgfqpoint{2.370206in}{0.848561in}}%
\pgfpathlineto{\pgfqpoint{2.364353in}{0.841406in}}%
\pgfpathlineto{\pgfqpoint{2.366196in}{0.837164in}}%
\pgfpathlineto{\pgfqpoint{2.371879in}{0.836595in}}%
\pgfpathlineto{\pgfqpoint{2.365973in}{0.833290in}}%
\pgfpathlineto{\pgfqpoint{2.344269in}{0.832393in}}%
\pgfpathlineto{\pgfqpoint{2.344959in}{0.835935in}}%
\pgfpathlineto{\pgfqpoint{2.334864in}{0.835567in}}%
\pgfpathlineto{\pgfqpoint{2.331640in}{0.845019in}}%
\pgfpathlineto{\pgfqpoint{2.337645in}{0.859095in}}%
\pgfpathlineto{\pgfqpoint{2.340706in}{0.860346in}}%
\pgfpathlineto{\pgfqpoint{2.344274in}{0.869656in}}%
\pgfpathlineto{\pgfqpoint{2.362240in}{0.870230in}}%
\pgfpathclose%
\pgfusepath{fill}%
\end{pgfscope}%
\begin{pgfscope}%
\pgfpathrectangle{\pgfqpoint{0.100000in}{0.100000in}}{\pgfqpoint{3.608454in}{2.310000in}}%
\pgfusepath{clip}%
\pgfsetbuttcap%
\pgfsetmiterjoin%
\definecolor{currentfill}{rgb}{0.000000,0.749020,0.625490}%
\pgfsetfillcolor{currentfill}%
\pgfsetlinewidth{0.000000pt}%
\definecolor{currentstroke}{rgb}{0.000000,0.000000,0.000000}%
\pgfsetstrokecolor{currentstroke}%
\pgfsetstrokeopacity{0.000000}%
\pgfsetdash{}{0pt}%
\pgfpathmoveto{\pgfqpoint{2.093068in}{0.691255in}}%
\pgfpathlineto{\pgfqpoint{2.089584in}{0.697805in}}%
\pgfpathlineto{\pgfqpoint{2.084351in}{0.694073in}}%
\pgfpathlineto{\pgfqpoint{2.081886in}{0.666080in}}%
\pgfpathlineto{\pgfqpoint{2.065563in}{0.666170in}}%
\pgfpathlineto{\pgfqpoint{2.049731in}{0.675985in}}%
\pgfpathlineto{\pgfqpoint{2.047653in}{0.694679in}}%
\pgfpathlineto{\pgfqpoint{2.026699in}{0.691630in}}%
\pgfpathlineto{\pgfqpoint{2.021875in}{0.703710in}}%
\pgfpathlineto{\pgfqpoint{2.040106in}{0.712927in}}%
\pgfpathlineto{\pgfqpoint{2.054567in}{0.712981in}}%
\pgfpathlineto{\pgfqpoint{2.053852in}{0.716384in}}%
\pgfpathlineto{\pgfqpoint{2.057335in}{0.727032in}}%
\pgfpathlineto{\pgfqpoint{2.061056in}{0.729018in}}%
\pgfpathlineto{\pgfqpoint{2.056507in}{0.745639in}}%
\pgfpathlineto{\pgfqpoint{2.062519in}{0.748614in}}%
\pgfpathlineto{\pgfqpoint{2.088232in}{0.752553in}}%
\pgfpathlineto{\pgfqpoint{2.096621in}{0.750671in}}%
\pgfpathlineto{\pgfqpoint{2.099586in}{0.743769in}}%
\pgfpathlineto{\pgfqpoint{2.106573in}{0.739232in}}%
\pgfpathlineto{\pgfqpoint{2.115910in}{0.747279in}}%
\pgfpathlineto{\pgfqpoint{2.125247in}{0.741762in}}%
\pgfpathlineto{\pgfqpoint{2.140576in}{0.738107in}}%
\pgfpathlineto{\pgfqpoint{2.145077in}{0.730301in}}%
\pgfpathlineto{\pgfqpoint{2.152123in}{0.723725in}}%
\pgfpathlineto{\pgfqpoint{2.166253in}{0.713450in}}%
\pgfpathlineto{\pgfqpoint{2.143868in}{0.708025in}}%
\pgfpathlineto{\pgfqpoint{2.132588in}{0.712274in}}%
\pgfpathlineto{\pgfqpoint{2.123775in}{0.713712in}}%
\pgfpathlineto{\pgfqpoint{2.117525in}{0.716961in}}%
\pgfpathlineto{\pgfqpoint{2.110954in}{0.708601in}}%
\pgfpathlineto{\pgfqpoint{2.093068in}{0.691255in}}%
\pgfpathclose%
\pgfusepath{fill}%
\end{pgfscope}%
\begin{pgfscope}%
\pgfpathrectangle{\pgfqpoint{0.100000in}{0.100000in}}{\pgfqpoint{3.608454in}{2.310000in}}%
\pgfusepath{clip}%
\pgfsetbuttcap%
\pgfsetmiterjoin%
\definecolor{currentfill}{rgb}{0.000000,0.654902,0.672549}%
\pgfsetfillcolor{currentfill}%
\pgfsetlinewidth{0.000000pt}%
\definecolor{currentstroke}{rgb}{0.000000,0.000000,0.000000}%
\pgfsetstrokecolor{currentstroke}%
\pgfsetstrokeopacity{0.000000}%
\pgfsetdash{}{0pt}%
\pgfpathmoveto{\pgfqpoint{2.904596in}{1.155570in}}%
\pgfpathlineto{\pgfqpoint{2.898012in}{1.153842in}}%
\pgfpathlineto{\pgfqpoint{2.895600in}{1.161365in}}%
\pgfpathlineto{\pgfqpoint{2.879988in}{1.149617in}}%
\pgfpathlineto{\pgfqpoint{2.868331in}{1.159628in}}%
\pgfpathlineto{\pgfqpoint{2.860386in}{1.164003in}}%
\pgfpathlineto{\pgfqpoint{2.865237in}{1.171385in}}%
\pgfpathlineto{\pgfqpoint{2.859396in}{1.177849in}}%
\pgfpathlineto{\pgfqpoint{2.853813in}{1.173625in}}%
\pgfpathlineto{\pgfqpoint{2.851139in}{1.181704in}}%
\pgfpathlineto{\pgfqpoint{2.856311in}{1.184739in}}%
\pgfpathlineto{\pgfqpoint{2.859737in}{1.192156in}}%
\pgfpathlineto{\pgfqpoint{2.867721in}{1.199764in}}%
\pgfpathlineto{\pgfqpoint{2.869228in}{1.206518in}}%
\pgfpathlineto{\pgfqpoint{2.876442in}{1.212071in}}%
\pgfpathlineto{\pgfqpoint{2.880082in}{1.220503in}}%
\pgfpathlineto{\pgfqpoint{2.886736in}{1.225061in}}%
\pgfpathlineto{\pgfqpoint{2.901040in}{1.226866in}}%
\pgfpathlineto{\pgfqpoint{2.907547in}{1.214437in}}%
\pgfpathlineto{\pgfqpoint{2.920870in}{1.223465in}}%
\pgfpathlineto{\pgfqpoint{2.921871in}{1.228104in}}%
\pgfpathlineto{\pgfqpoint{2.933591in}{1.232567in}}%
\pgfpathlineto{\pgfqpoint{2.937030in}{1.237689in}}%
\pgfpathlineto{\pgfqpoint{2.943559in}{1.240996in}}%
\pgfpathlineto{\pgfqpoint{2.960170in}{1.246122in}}%
\pgfpathlineto{\pgfqpoint{2.968379in}{1.238517in}}%
\pgfpathlineto{\pgfqpoint{2.974075in}{1.228469in}}%
\pgfpathlineto{\pgfqpoint{2.958209in}{1.221651in}}%
\pgfpathlineto{\pgfqpoint{2.957411in}{1.217461in}}%
\pgfpathlineto{\pgfqpoint{2.951682in}{1.213011in}}%
\pgfpathlineto{\pgfqpoint{2.940357in}{1.211530in}}%
\pgfpathlineto{\pgfqpoint{2.934516in}{1.206979in}}%
\pgfpathlineto{\pgfqpoint{2.933739in}{1.199536in}}%
\pgfpathlineto{\pgfqpoint{2.929794in}{1.191773in}}%
\pgfpathlineto{\pgfqpoint{2.936507in}{1.182790in}}%
\pgfpathlineto{\pgfqpoint{2.933147in}{1.171253in}}%
\pgfpathlineto{\pgfqpoint{2.930144in}{1.169183in}}%
\pgfpathlineto{\pgfqpoint{2.920582in}{1.172131in}}%
\pgfpathlineto{\pgfqpoint{2.909797in}{1.165077in}}%
\pgfpathlineto{\pgfqpoint{2.904596in}{1.155570in}}%
\pgfpathclose%
\pgfusepath{fill}%
\end{pgfscope}%
\begin{pgfscope}%
\pgfpathrectangle{\pgfqpoint{0.100000in}{0.100000in}}{\pgfqpoint{3.608454in}{2.310000in}}%
\pgfusepath{clip}%
\pgfsetbuttcap%
\pgfsetmiterjoin%
\definecolor{currentfill}{rgb}{0.000000,0.701961,0.649020}%
\pgfsetfillcolor{currentfill}%
\pgfsetlinewidth{0.000000pt}%
\definecolor{currentstroke}{rgb}{0.000000,0.000000,0.000000}%
\pgfsetstrokecolor{currentstroke}%
\pgfsetstrokeopacity{0.000000}%
\pgfsetdash{}{0pt}%
\pgfpathmoveto{\pgfqpoint{2.537283in}{1.876850in}}%
\pgfpathlineto{\pgfqpoint{2.528827in}{1.854971in}}%
\pgfpathlineto{\pgfqpoint{2.521970in}{1.844929in}}%
\pgfpathlineto{\pgfqpoint{2.520382in}{1.829615in}}%
\pgfpathlineto{\pgfqpoint{2.513293in}{1.828191in}}%
\pgfpathlineto{\pgfqpoint{2.492843in}{1.831235in}}%
\pgfpathlineto{\pgfqpoint{2.492447in}{1.838114in}}%
\pgfpathlineto{\pgfqpoint{2.488144in}{1.844796in}}%
\pgfpathlineto{\pgfqpoint{2.481279in}{1.844556in}}%
\pgfpathlineto{\pgfqpoint{2.480448in}{1.858288in}}%
\pgfpathlineto{\pgfqpoint{2.473620in}{1.857983in}}%
\pgfpathlineto{\pgfqpoint{2.471787in}{1.885405in}}%
\pgfpathlineto{\pgfqpoint{2.457894in}{1.884450in}}%
\pgfpathlineto{\pgfqpoint{2.456551in}{1.904950in}}%
\pgfpathlineto{\pgfqpoint{2.438472in}{1.913015in}}%
\pgfpathlineto{\pgfqpoint{2.436761in}{1.938649in}}%
\pgfpathlineto{\pgfqpoint{2.484672in}{1.942048in}}%
\pgfpathlineto{\pgfqpoint{2.485703in}{1.928298in}}%
\pgfpathlineto{\pgfqpoint{2.513243in}{1.930478in}}%
\pgfpathlineto{\pgfqpoint{2.514867in}{1.909757in}}%
\pgfpathlineto{\pgfqpoint{2.528620in}{1.910846in}}%
\pgfpathlineto{\pgfqpoint{2.531516in}{1.904166in}}%
\pgfpathlineto{\pgfqpoint{2.533757in}{1.876549in}}%
\pgfpathlineto{\pgfqpoint{2.537283in}{1.876850in}}%
\pgfpathclose%
\pgfusepath{fill}%
\end{pgfscope}%
\begin{pgfscope}%
\pgfpathrectangle{\pgfqpoint{0.100000in}{0.100000in}}{\pgfqpoint{3.608454in}{2.310000in}}%
\pgfusepath{clip}%
\pgfsetbuttcap%
\pgfsetmiterjoin%
\definecolor{currentfill}{rgb}{0.000000,0.545098,0.727451}%
\pgfsetfillcolor{currentfill}%
\pgfsetlinewidth{0.000000pt}%
\definecolor{currentstroke}{rgb}{0.000000,0.000000,0.000000}%
\pgfsetstrokecolor{currentstroke}%
\pgfsetstrokeopacity{0.000000}%
\pgfsetdash{}{0pt}%
\pgfpathmoveto{\pgfqpoint{2.507414in}{1.036455in}}%
\pgfpathlineto{\pgfqpoint{2.483127in}{1.034794in}}%
\pgfpathlineto{\pgfqpoint{2.482016in}{1.066685in}}%
\pgfpathlineto{\pgfqpoint{2.488608in}{1.069834in}}%
\pgfpathlineto{\pgfqpoint{2.487860in}{1.090523in}}%
\pgfpathlineto{\pgfqpoint{2.469563in}{1.097602in}}%
\pgfpathlineto{\pgfqpoint{2.480144in}{1.116463in}}%
\pgfpathlineto{\pgfqpoint{2.479806in}{1.130609in}}%
\pgfpathlineto{\pgfqpoint{2.492139in}{1.131174in}}%
\pgfpathlineto{\pgfqpoint{2.498683in}{1.126546in}}%
\pgfpathlineto{\pgfqpoint{2.510006in}{1.121373in}}%
\pgfpathlineto{\pgfqpoint{2.510594in}{1.099857in}}%
\pgfpathlineto{\pgfqpoint{2.517006in}{1.100169in}}%
\pgfpathlineto{\pgfqpoint{2.517692in}{1.084251in}}%
\pgfpathlineto{\pgfqpoint{2.525076in}{1.078728in}}%
\pgfpathlineto{\pgfqpoint{2.531082in}{1.076803in}}%
\pgfpathlineto{\pgfqpoint{2.534914in}{1.072050in}}%
\pgfpathlineto{\pgfqpoint{2.535012in}{1.069064in}}%
\pgfpathlineto{\pgfqpoint{2.521534in}{1.068158in}}%
\pgfpathlineto{\pgfqpoint{2.517201in}{1.066138in}}%
\pgfpathlineto{\pgfqpoint{2.512736in}{1.056859in}}%
\pgfpathlineto{\pgfqpoint{2.508734in}{1.056586in}}%
\pgfpathlineto{\pgfqpoint{2.509787in}{1.036620in}}%
\pgfpathlineto{\pgfqpoint{2.507414in}{1.036455in}}%
\pgfpathclose%
\pgfusepath{fill}%
\end{pgfscope}%
\begin{pgfscope}%
\pgfpathrectangle{\pgfqpoint{0.100000in}{0.100000in}}{\pgfqpoint{3.608454in}{2.310000in}}%
\pgfusepath{clip}%
\pgfsetbuttcap%
\pgfsetmiterjoin%
\definecolor{currentfill}{rgb}{0.000000,0.490196,0.754902}%
\pgfsetfillcolor{currentfill}%
\pgfsetlinewidth{0.000000pt}%
\definecolor{currentstroke}{rgb}{0.000000,0.000000,0.000000}%
\pgfsetstrokecolor{currentstroke}%
\pgfsetstrokeopacity{0.000000}%
\pgfsetdash{}{0pt}%
\pgfpathmoveto{\pgfqpoint{1.960306in}{1.310682in}}%
\pgfpathlineto{\pgfqpoint{1.926190in}{1.311718in}}%
\pgfpathlineto{\pgfqpoint{1.926584in}{1.332368in}}%
\pgfpathlineto{\pgfqpoint{1.892401in}{1.333545in}}%
\pgfpathlineto{\pgfqpoint{1.892889in}{1.354245in}}%
\pgfpathlineto{\pgfqpoint{1.893147in}{1.361110in}}%
\pgfpathlineto{\pgfqpoint{1.927400in}{1.359864in}}%
\pgfpathlineto{\pgfqpoint{1.927666in}{1.366756in}}%
\pgfpathlineto{\pgfqpoint{1.961886in}{1.365738in}}%
\pgfpathlineto{\pgfqpoint{1.960306in}{1.310682in}}%
\pgfpathclose%
\pgfusepath{fill}%
\end{pgfscope}%
\begin{pgfscope}%
\pgfpathrectangle{\pgfqpoint{0.100000in}{0.100000in}}{\pgfqpoint{3.608454in}{2.310000in}}%
\pgfusepath{clip}%
\pgfsetbuttcap%
\pgfsetmiterjoin%
\definecolor{currentfill}{rgb}{0.000000,0.494118,0.752941}%
\pgfsetfillcolor{currentfill}%
\pgfsetlinewidth{0.000000pt}%
\definecolor{currentstroke}{rgb}{0.000000,0.000000,0.000000}%
\pgfsetstrokecolor{currentstroke}%
\pgfsetstrokeopacity{0.000000}%
\pgfsetdash{}{0pt}%
\pgfpathmoveto{\pgfqpoint{2.548186in}{1.158629in}}%
\pgfpathlineto{\pgfqpoint{2.549236in}{1.162164in}}%
\pgfpathlineto{\pgfqpoint{2.543041in}{1.178054in}}%
\pgfpathlineto{\pgfqpoint{2.539380in}{1.187318in}}%
\pgfpathlineto{\pgfqpoint{2.555635in}{1.195943in}}%
\pgfpathlineto{\pgfqpoint{2.563459in}{1.197714in}}%
\pgfpathlineto{\pgfqpoint{2.566571in}{1.201683in}}%
\pgfpathlineto{\pgfqpoint{2.567654in}{1.211922in}}%
\pgfpathlineto{\pgfqpoint{2.577947in}{1.209252in}}%
\pgfpathlineto{\pgfqpoint{2.589411in}{1.214310in}}%
\pgfpathlineto{\pgfqpoint{2.595443in}{1.205574in}}%
\pgfpathlineto{\pgfqpoint{2.607679in}{1.208265in}}%
\pgfpathlineto{\pgfqpoint{2.610086in}{1.175304in}}%
\pgfpathlineto{\pgfqpoint{2.606663in}{1.174945in}}%
\pgfpathlineto{\pgfqpoint{2.607032in}{1.160645in}}%
\pgfpathlineto{\pgfqpoint{2.599908in}{1.153357in}}%
\pgfpathlineto{\pgfqpoint{2.597965in}{1.148745in}}%
\pgfpathlineto{\pgfqpoint{2.583459in}{1.148544in}}%
\pgfpathlineto{\pgfqpoint{2.583376in}{1.141723in}}%
\pgfpathlineto{\pgfqpoint{2.579740in}{1.137359in}}%
\pgfpathlineto{\pgfqpoint{2.570262in}{1.139592in}}%
\pgfpathlineto{\pgfqpoint{2.556313in}{1.139100in}}%
\pgfpathlineto{\pgfqpoint{2.555391in}{1.144733in}}%
\pgfpathlineto{\pgfqpoint{2.548297in}{1.153932in}}%
\pgfpathlineto{\pgfqpoint{2.548186in}{1.158629in}}%
\pgfpathclose%
\pgfusepath{fill}%
\end{pgfscope}%
\begin{pgfscope}%
\pgfpathrectangle{\pgfqpoint{0.100000in}{0.100000in}}{\pgfqpoint{3.608454in}{2.310000in}}%
\pgfusepath{clip}%
\pgfsetbuttcap%
\pgfsetmiterjoin%
\definecolor{currentfill}{rgb}{0.000000,0.317647,0.841176}%
\pgfsetfillcolor{currentfill}%
\pgfsetlinewidth{0.000000pt}%
\definecolor{currentstroke}{rgb}{0.000000,0.000000,0.000000}%
\pgfsetstrokecolor{currentstroke}%
\pgfsetstrokeopacity{0.000000}%
\pgfsetdash{}{0pt}%
\pgfpathmoveto{\pgfqpoint{3.245459in}{1.515286in}}%
\pgfpathlineto{\pgfqpoint{3.266966in}{1.519876in}}%
\pgfpathlineto{\pgfqpoint{3.271013in}{1.528198in}}%
\pgfpathlineto{\pgfqpoint{3.275612in}{1.531070in}}%
\pgfpathlineto{\pgfqpoint{3.284318in}{1.531602in}}%
\pgfpathlineto{\pgfqpoint{3.284745in}{1.526913in}}%
\pgfpathlineto{\pgfqpoint{3.281349in}{1.515956in}}%
\pgfpathlineto{\pgfqpoint{3.285117in}{1.511332in}}%
\pgfpathlineto{\pgfqpoint{3.285323in}{1.505755in}}%
\pgfpathlineto{\pgfqpoint{3.282148in}{1.502290in}}%
\pgfpathlineto{\pgfqpoint{3.288460in}{1.495799in}}%
\pgfpathlineto{\pgfqpoint{3.281434in}{1.488221in}}%
\pgfpathlineto{\pgfqpoint{3.274847in}{1.487194in}}%
\pgfpathlineto{\pgfqpoint{3.273108in}{1.493337in}}%
\pgfpathlineto{\pgfqpoint{3.266239in}{1.490550in}}%
\pgfpathlineto{\pgfqpoint{3.256327in}{1.489994in}}%
\pgfpathlineto{\pgfqpoint{3.258440in}{1.494561in}}%
\pgfpathlineto{\pgfqpoint{3.257846in}{1.502321in}}%
\pgfpathlineto{\pgfqpoint{3.251965in}{1.502268in}}%
\pgfpathlineto{\pgfqpoint{3.239646in}{1.514041in}}%
\pgfpathlineto{\pgfqpoint{3.245459in}{1.515286in}}%
\pgfpathclose%
\pgfusepath{fill}%
\end{pgfscope}%
\begin{pgfscope}%
\pgfpathrectangle{\pgfqpoint{0.100000in}{0.100000in}}{\pgfqpoint{3.608454in}{2.310000in}}%
\pgfusepath{clip}%
\pgfsetbuttcap%
\pgfsetmiterjoin%
\definecolor{currentfill}{rgb}{0.000000,0.364706,0.817647}%
\pgfsetfillcolor{currentfill}%
\pgfsetlinewidth{0.000000pt}%
\definecolor{currentstroke}{rgb}{0.000000,0.000000,0.000000}%
\pgfsetstrokecolor{currentstroke}%
\pgfsetstrokeopacity{0.000000}%
\pgfsetdash{}{0pt}%
\pgfpathmoveto{\pgfqpoint{1.791350in}{1.359031in}}%
\pgfpathlineto{\pgfqpoint{1.790342in}{1.359067in}}%
\pgfpathlineto{\pgfqpoint{1.792462in}{1.393414in}}%
\pgfpathlineto{\pgfqpoint{1.824934in}{1.391676in}}%
\pgfpathlineto{\pgfqpoint{1.860470in}{1.390012in}}%
\pgfpathlineto{\pgfqpoint{1.858721in}{1.355685in}}%
\pgfpathlineto{\pgfqpoint{1.791350in}{1.359031in}}%
\pgfpathclose%
\pgfusepath{fill}%
\end{pgfscope}%
\begin{pgfscope}%
\pgfpathrectangle{\pgfqpoint{0.100000in}{0.100000in}}{\pgfqpoint{3.608454in}{2.310000in}}%
\pgfusepath{clip}%
\pgfsetbuttcap%
\pgfsetmiterjoin%
\definecolor{currentfill}{rgb}{0.000000,0.478431,0.760784}%
\pgfsetfillcolor{currentfill}%
\pgfsetlinewidth{0.000000pt}%
\definecolor{currentstroke}{rgb}{0.000000,0.000000,0.000000}%
\pgfsetstrokecolor{currentstroke}%
\pgfsetstrokeopacity{0.000000}%
\pgfsetdash{}{0pt}%
\pgfpathmoveto{\pgfqpoint{1.895416in}{1.185262in}}%
\pgfpathlineto{\pgfqpoint{1.896376in}{1.215785in}}%
\pgfpathlineto{\pgfqpoint{1.889149in}{1.216054in}}%
\pgfpathlineto{\pgfqpoint{1.890205in}{1.243605in}}%
\pgfpathlineto{\pgfqpoint{1.938084in}{1.242176in}}%
\pgfpathlineto{\pgfqpoint{1.938292in}{1.256226in}}%
\pgfpathlineto{\pgfqpoint{1.972560in}{1.255285in}}%
\pgfpathlineto{\pgfqpoint{1.972846in}{1.269115in}}%
\pgfpathlineto{\pgfqpoint{1.992256in}{1.268523in}}%
\pgfpathlineto{\pgfqpoint{2.012052in}{1.268224in}}%
\pgfpathlineto{\pgfqpoint{2.011063in}{1.220009in}}%
\pgfpathlineto{\pgfqpoint{1.971690in}{1.220737in}}%
\pgfpathlineto{\pgfqpoint{1.971148in}{1.183063in}}%
\pgfpathlineto{\pgfqpoint{1.951284in}{1.183534in}}%
\pgfpathlineto{\pgfqpoint{1.895416in}{1.185262in}}%
\pgfpathclose%
\pgfusepath{fill}%
\end{pgfscope}%
\begin{pgfscope}%
\pgfpathrectangle{\pgfqpoint{0.100000in}{0.100000in}}{\pgfqpoint{3.608454in}{2.310000in}}%
\pgfusepath{clip}%
\pgfsetbuttcap%
\pgfsetmiterjoin%
\definecolor{currentfill}{rgb}{0.000000,0.470588,0.764706}%
\pgfsetfillcolor{currentfill}%
\pgfsetlinewidth{0.000000pt}%
\definecolor{currentstroke}{rgb}{0.000000,0.000000,0.000000}%
\pgfsetstrokecolor{currentstroke}%
\pgfsetstrokeopacity{0.000000}%
\pgfsetdash{}{0pt}%
\pgfpathmoveto{\pgfqpoint{1.811548in}{0.930044in}}%
\pgfpathlineto{\pgfqpoint{1.810392in}{0.932854in}}%
\pgfpathlineto{\pgfqpoint{1.797246in}{0.938094in}}%
\pgfpathlineto{\pgfqpoint{1.774163in}{0.940149in}}%
\pgfpathlineto{\pgfqpoint{1.775993in}{0.971342in}}%
\pgfpathlineto{\pgfqpoint{1.779205in}{0.970733in}}%
\pgfpathlineto{\pgfqpoint{1.780677in}{0.997355in}}%
\pgfpathlineto{\pgfqpoint{1.785322in}{0.998334in}}%
\pgfpathlineto{\pgfqpoint{1.799655in}{0.981871in}}%
\pgfpathlineto{\pgfqpoint{1.805846in}{0.981264in}}%
\pgfpathlineto{\pgfqpoint{1.811536in}{0.984126in}}%
\pgfpathlineto{\pgfqpoint{1.818337in}{0.980492in}}%
\pgfpathlineto{\pgfqpoint{1.820478in}{0.987070in}}%
\pgfpathlineto{\pgfqpoint{1.831143in}{0.977031in}}%
\pgfpathlineto{\pgfqpoint{1.832097in}{0.967259in}}%
\pgfpathlineto{\pgfqpoint{1.847621in}{0.966431in}}%
\pgfpathlineto{\pgfqpoint{1.846276in}{0.936427in}}%
\pgfpathlineto{\pgfqpoint{1.811868in}{0.938019in}}%
\pgfpathlineto{\pgfqpoint{1.811548in}{0.930044in}}%
\pgfpathclose%
\pgfusepath{fill}%
\end{pgfscope}%
\begin{pgfscope}%
\pgfpathrectangle{\pgfqpoint{0.100000in}{0.100000in}}{\pgfqpoint{3.608454in}{2.310000in}}%
\pgfusepath{clip}%
\pgfsetbuttcap%
\pgfsetmiterjoin%
\definecolor{currentfill}{rgb}{0.000000,0.537255,0.731373}%
\pgfsetfillcolor{currentfill}%
\pgfsetlinewidth{0.000000pt}%
\definecolor{currentstroke}{rgb}{0.000000,0.000000,0.000000}%
\pgfsetstrokecolor{currentstroke}%
\pgfsetstrokeopacity{0.000000}%
\pgfsetdash{}{0pt}%
\pgfpathmoveto{\pgfqpoint{0.906140in}{0.393309in}}%
\pgfpathlineto{\pgfqpoint{0.902140in}{0.395283in}}%
\pgfpathlineto{\pgfqpoint{0.900276in}{0.397748in}}%
\pgfpathlineto{\pgfqpoint{0.897685in}{0.398394in}}%
\pgfpathlineto{\pgfqpoint{0.898240in}{0.401261in}}%
\pgfpathlineto{\pgfqpoint{0.897317in}{0.402234in}}%
\pgfpathlineto{\pgfqpoint{0.894866in}{0.400086in}}%
\pgfpathlineto{\pgfqpoint{0.892669in}{0.399270in}}%
\pgfpathlineto{\pgfqpoint{0.891166in}{0.400988in}}%
\pgfpathlineto{\pgfqpoint{0.889388in}{0.400332in}}%
\pgfpathlineto{\pgfqpoint{0.888087in}{0.396647in}}%
\pgfpathlineto{\pgfqpoint{0.885872in}{0.401723in}}%
\pgfpathlineto{\pgfqpoint{0.885272in}{0.397710in}}%
\pgfpathlineto{\pgfqpoint{0.881126in}{0.397511in}}%
\pgfpathlineto{\pgfqpoint{0.880525in}{0.396050in}}%
\pgfpathlineto{\pgfqpoint{0.876957in}{0.396163in}}%
\pgfpathlineto{\pgfqpoint{0.872853in}{0.397496in}}%
\pgfpathlineto{\pgfqpoint{0.866722in}{0.396548in}}%
\pgfpathlineto{\pgfqpoint{0.864048in}{0.397682in}}%
\pgfpathlineto{\pgfqpoint{0.863015in}{0.395943in}}%
\pgfpathlineto{\pgfqpoint{0.860626in}{0.397438in}}%
\pgfpathlineto{\pgfqpoint{0.860500in}{0.399959in}}%
\pgfpathlineto{\pgfqpoint{0.857986in}{0.398830in}}%
\pgfpathlineto{\pgfqpoint{0.854460in}{0.399934in}}%
\pgfpathlineto{\pgfqpoint{0.854188in}{0.405776in}}%
\pgfpathlineto{\pgfqpoint{0.857163in}{0.407525in}}%
\pgfpathlineto{\pgfqpoint{0.861136in}{0.406969in}}%
\pgfpathlineto{\pgfqpoint{0.862984in}{0.405050in}}%
\pgfpathlineto{\pgfqpoint{0.865125in}{0.405991in}}%
\pgfpathlineto{\pgfqpoint{0.867877in}{0.404847in}}%
\pgfpathlineto{\pgfqpoint{0.868825in}{0.406515in}}%
\pgfpathlineto{\pgfqpoint{0.874638in}{0.407994in}}%
\pgfpathlineto{\pgfqpoint{0.872011in}{0.409155in}}%
\pgfpathlineto{\pgfqpoint{0.868887in}{0.408859in}}%
\pgfpathlineto{\pgfqpoint{0.865823in}{0.407518in}}%
\pgfpathlineto{\pgfqpoint{0.863498in}{0.411733in}}%
\pgfpathlineto{\pgfqpoint{0.863048in}{0.414812in}}%
\pgfpathlineto{\pgfqpoint{0.869799in}{0.420224in}}%
\pgfpathlineto{\pgfqpoint{0.874638in}{0.421627in}}%
\pgfpathlineto{\pgfqpoint{0.880313in}{0.424977in}}%
\pgfpathlineto{\pgfqpoint{0.883440in}{0.427931in}}%
\pgfpathlineto{\pgfqpoint{0.884480in}{0.433448in}}%
\pgfpathlineto{\pgfqpoint{0.887231in}{0.433726in}}%
\pgfpathlineto{\pgfqpoint{0.889978in}{0.432755in}}%
\pgfpathlineto{\pgfqpoint{0.896244in}{0.433381in}}%
\pgfpathlineto{\pgfqpoint{0.902329in}{0.432923in}}%
\pgfpathlineto{\pgfqpoint{0.902058in}{0.429768in}}%
\pgfpathlineto{\pgfqpoint{0.903873in}{0.427224in}}%
\pgfpathlineto{\pgfqpoint{0.909287in}{0.427185in}}%
\pgfpathlineto{\pgfqpoint{0.907671in}{0.433445in}}%
\pgfpathlineto{\pgfqpoint{0.909967in}{0.434891in}}%
\pgfpathlineto{\pgfqpoint{0.905039in}{0.438552in}}%
\pgfpathlineto{\pgfqpoint{0.903341in}{0.440916in}}%
\pgfpathlineto{\pgfqpoint{0.897892in}{0.441660in}}%
\pgfpathlineto{\pgfqpoint{0.894137in}{0.440185in}}%
\pgfpathlineto{\pgfqpoint{0.888440in}{0.441543in}}%
\pgfpathlineto{\pgfqpoint{0.882939in}{0.440631in}}%
\pgfpathlineto{\pgfqpoint{0.881028in}{0.436394in}}%
\pgfpathlineto{\pgfqpoint{0.879917in}{0.438011in}}%
\pgfpathlineto{\pgfqpoint{0.876955in}{0.437975in}}%
\pgfpathlineto{\pgfqpoint{0.871270in}{0.436691in}}%
\pgfpathlineto{\pgfqpoint{0.868136in}{0.433288in}}%
\pgfpathlineto{\pgfqpoint{0.865686in}{0.432241in}}%
\pgfpathlineto{\pgfqpoint{0.862763in}{0.432076in}}%
\pgfpathlineto{\pgfqpoint{0.860841in}{0.433543in}}%
\pgfpathlineto{\pgfqpoint{0.860190in}{0.428478in}}%
\pgfpathlineto{\pgfqpoint{0.855724in}{0.426050in}}%
\pgfpathlineto{\pgfqpoint{0.853611in}{0.426386in}}%
\pgfpathlineto{\pgfqpoint{0.851971in}{0.427933in}}%
\pgfpathlineto{\pgfqpoint{0.849085in}{0.424656in}}%
\pgfpathlineto{\pgfqpoint{0.846580in}{0.424147in}}%
\pgfpathlineto{\pgfqpoint{0.840273in}{0.426897in}}%
\pgfpathlineto{\pgfqpoint{0.838838in}{0.425763in}}%
\pgfpathlineto{\pgfqpoint{0.836699in}{0.425925in}}%
\pgfpathlineto{\pgfqpoint{0.835510in}{0.423824in}}%
\pgfpathlineto{\pgfqpoint{0.831134in}{0.425525in}}%
\pgfpathlineto{\pgfqpoint{0.826995in}{0.422232in}}%
\pgfpathlineto{\pgfqpoint{0.824051in}{0.421527in}}%
\pgfpathlineto{\pgfqpoint{0.825544in}{0.418315in}}%
\pgfpathlineto{\pgfqpoint{0.828849in}{0.415327in}}%
\pgfpathlineto{\pgfqpoint{0.831233in}{0.409187in}}%
\pgfpathlineto{\pgfqpoint{0.831242in}{0.406233in}}%
\pgfpathlineto{\pgfqpoint{0.828822in}{0.408037in}}%
\pgfpathlineto{\pgfqpoint{0.826771in}{0.405229in}}%
\pgfpathlineto{\pgfqpoint{0.822419in}{0.408486in}}%
\pgfpathlineto{\pgfqpoint{0.820929in}{0.406477in}}%
\pgfpathlineto{\pgfqpoint{0.815079in}{0.410909in}}%
\pgfpathlineto{\pgfqpoint{0.811149in}{0.413923in}}%
\pgfpathlineto{\pgfqpoint{0.812034in}{0.416427in}}%
\pgfpathlineto{\pgfqpoint{0.818042in}{0.424203in}}%
\pgfpathlineto{\pgfqpoint{0.817435in}{0.424709in}}%
\pgfpathlineto{\pgfqpoint{0.820512in}{0.428750in}}%
\pgfpathlineto{\pgfqpoint{0.824498in}{0.425742in}}%
\pgfpathlineto{\pgfqpoint{0.827494in}{0.429779in}}%
\pgfpathlineto{\pgfqpoint{0.830966in}{0.427222in}}%
\pgfpathlineto{\pgfqpoint{0.832475in}{0.429280in}}%
\pgfpathlineto{\pgfqpoint{0.834469in}{0.427785in}}%
\pgfpathlineto{\pgfqpoint{0.837531in}{0.431857in}}%
\pgfpathlineto{\pgfqpoint{0.839547in}{0.430349in}}%
\pgfpathlineto{\pgfqpoint{0.841100in}{0.432388in}}%
\pgfpathlineto{\pgfqpoint{0.842589in}{0.431260in}}%
\pgfpathlineto{\pgfqpoint{0.848446in}{0.439078in}}%
\pgfpathlineto{\pgfqpoint{0.849915in}{0.438009in}}%
\pgfpathlineto{\pgfqpoint{0.856061in}{0.446120in}}%
\pgfpathlineto{\pgfqpoint{0.857604in}{0.444984in}}%
\pgfpathlineto{\pgfqpoint{0.863736in}{0.453184in}}%
\pgfpathlineto{\pgfqpoint{0.864761in}{0.455531in}}%
\pgfpathlineto{\pgfqpoint{0.869217in}{0.461424in}}%
\pgfpathlineto{\pgfqpoint{0.870359in}{0.464352in}}%
\pgfpathlineto{\pgfqpoint{0.874246in}{0.469508in}}%
\pgfpathlineto{\pgfqpoint{0.878374in}{0.466226in}}%
\pgfpathlineto{\pgfqpoint{0.878747in}{0.464023in}}%
\pgfpathlineto{\pgfqpoint{0.894650in}{0.485392in}}%
\pgfpathlineto{\pgfqpoint{0.902205in}{0.495355in}}%
\pgfpathlineto{\pgfqpoint{0.914054in}{0.486250in}}%
\pgfpathlineto{\pgfqpoint{0.915352in}{0.487908in}}%
\pgfpathlineto{\pgfqpoint{0.940198in}{0.483801in}}%
\pgfpathlineto{\pgfqpoint{0.950166in}{0.482212in}}%
\pgfpathlineto{\pgfqpoint{0.966684in}{0.470967in}}%
\pgfpathlineto{\pgfqpoint{0.968975in}{0.474416in}}%
\pgfpathlineto{\pgfqpoint{0.972860in}{0.471797in}}%
\pgfpathlineto{\pgfqpoint{0.986452in}{0.463301in}}%
\pgfpathlineto{\pgfqpoint{0.979580in}{0.452265in}}%
\pgfpathlineto{\pgfqpoint{0.977356in}{0.447349in}}%
\pgfpathlineto{\pgfqpoint{0.967777in}{0.431965in}}%
\pgfpathlineto{\pgfqpoint{0.961359in}{0.436162in}}%
\pgfpathlineto{\pgfqpoint{0.960218in}{0.433861in}}%
\pgfpathlineto{\pgfqpoint{0.949551in}{0.416794in}}%
\pgfpathlineto{\pgfqpoint{0.946494in}{0.418805in}}%
\pgfpathlineto{\pgfqpoint{0.945694in}{0.417578in}}%
\pgfpathlineto{\pgfqpoint{0.930988in}{0.427301in}}%
\pgfpathlineto{\pgfqpoint{0.926490in}{0.421058in}}%
\pgfpathlineto{\pgfqpoint{0.921451in}{0.413351in}}%
\pgfpathlineto{\pgfqpoint{0.918151in}{0.415401in}}%
\pgfpathlineto{\pgfqpoint{0.916327in}{0.412673in}}%
\pgfpathlineto{\pgfqpoint{0.917407in}{0.411973in}}%
\pgfpathlineto{\pgfqpoint{0.912214in}{0.404513in}}%
\pgfpathlineto{\pgfqpoint{0.913452in}{0.403668in}}%
\pgfpathlineto{\pgfqpoint{0.908663in}{0.396558in}}%
\pgfpathlineto{\pgfqpoint{0.906140in}{0.393309in}}%
\pgfpathclose%
\pgfusepath{fill}%
\end{pgfscope}%
\begin{pgfscope}%
\pgfpathrectangle{\pgfqpoint{0.100000in}{0.100000in}}{\pgfqpoint{3.608454in}{2.310000in}}%
\pgfusepath{clip}%
\pgfsetbuttcap%
\pgfsetmiterjoin%
\definecolor{currentfill}{rgb}{0.000000,0.874510,0.562745}%
\pgfsetfillcolor{currentfill}%
\pgfsetlinewidth{0.000000pt}%
\definecolor{currentstroke}{rgb}{0.000000,0.000000,0.000000}%
\pgfsetstrokecolor{currentstroke}%
\pgfsetstrokeopacity{0.000000}%
\pgfsetdash{}{0pt}%
\pgfpathmoveto{\pgfqpoint{2.901040in}{1.226866in}}%
\pgfpathlineto{\pgfqpoint{2.907489in}{1.230091in}}%
\pgfpathlineto{\pgfqpoint{2.904928in}{1.234962in}}%
\pgfpathlineto{\pgfqpoint{2.916003in}{1.242873in}}%
\pgfpathlineto{\pgfqpoint{2.908022in}{1.254600in}}%
\pgfpathlineto{\pgfqpoint{2.902185in}{1.260526in}}%
\pgfpathlineto{\pgfqpoint{2.921023in}{1.282534in}}%
\pgfpathlineto{\pgfqpoint{2.923822in}{1.280887in}}%
\pgfpathlineto{\pgfqpoint{2.923526in}{1.275100in}}%
\pgfpathlineto{\pgfqpoint{2.926305in}{1.268226in}}%
\pgfpathlineto{\pgfqpoint{2.935481in}{1.264222in}}%
\pgfpathlineto{\pgfqpoint{2.942904in}{1.258852in}}%
\pgfpathlineto{\pgfqpoint{2.950534in}{1.260497in}}%
\pgfpathlineto{\pgfqpoint{2.960443in}{1.270838in}}%
\pgfpathlineto{\pgfqpoint{2.965505in}{1.286020in}}%
\pgfpathlineto{\pgfqpoint{2.967623in}{1.287269in}}%
\pgfpathlineto{\pgfqpoint{2.968540in}{1.292624in}}%
\pgfpathlineto{\pgfqpoint{2.974604in}{1.294585in}}%
\pgfpathlineto{\pgfqpoint{2.991086in}{1.284432in}}%
\pgfpathlineto{\pgfqpoint{2.992726in}{1.278073in}}%
\pgfpathlineto{\pgfqpoint{2.983343in}{1.270885in}}%
\pgfpathlineto{\pgfqpoint{2.994849in}{1.262443in}}%
\pgfpathlineto{\pgfqpoint{2.990512in}{1.259171in}}%
\pgfpathlineto{\pgfqpoint{2.993913in}{1.252683in}}%
\pgfpathlineto{\pgfqpoint{3.005275in}{1.242589in}}%
\pgfpathlineto{\pgfqpoint{2.986436in}{1.233543in}}%
\pgfpathlineto{\pgfqpoint{2.984220in}{1.230496in}}%
\pgfpathlineto{\pgfqpoint{2.974075in}{1.228469in}}%
\pgfpathlineto{\pgfqpoint{2.968379in}{1.238517in}}%
\pgfpathlineto{\pgfqpoint{2.960170in}{1.246122in}}%
\pgfpathlineto{\pgfqpoint{2.943559in}{1.240996in}}%
\pgfpathlineto{\pgfqpoint{2.937030in}{1.237689in}}%
\pgfpathlineto{\pgfqpoint{2.933591in}{1.232567in}}%
\pgfpathlineto{\pgfqpoint{2.921871in}{1.228104in}}%
\pgfpathlineto{\pgfqpoint{2.920870in}{1.223465in}}%
\pgfpathlineto{\pgfqpoint{2.907547in}{1.214437in}}%
\pgfpathlineto{\pgfqpoint{2.901040in}{1.226866in}}%
\pgfpathclose%
\pgfusepath{fill}%
\end{pgfscope}%
\begin{pgfscope}%
\pgfpathrectangle{\pgfqpoint{0.100000in}{0.100000in}}{\pgfqpoint{3.608454in}{2.310000in}}%
\pgfusepath{clip}%
\pgfsetbuttcap%
\pgfsetmiterjoin%
\definecolor{currentfill}{rgb}{0.000000,0.478431,0.760784}%
\pgfsetfillcolor{currentfill}%
\pgfsetlinewidth{0.000000pt}%
\definecolor{currentstroke}{rgb}{0.000000,0.000000,0.000000}%
\pgfsetstrokecolor{currentstroke}%
\pgfsetstrokeopacity{0.000000}%
\pgfsetdash{}{0pt}%
\pgfpathmoveto{\pgfqpoint{1.864232in}{2.006135in}}%
\pgfpathlineto{\pgfqpoint{1.862851in}{1.978498in}}%
\pgfpathlineto{\pgfqpoint{1.864737in}{1.978421in}}%
\pgfpathlineto{\pgfqpoint{1.863197in}{1.950599in}}%
\pgfpathlineto{\pgfqpoint{1.828664in}{1.952547in}}%
\pgfpathlineto{\pgfqpoint{1.830555in}{1.980267in}}%
\pgfpathlineto{\pgfqpoint{1.828413in}{1.980392in}}%
\pgfpathlineto{\pgfqpoint{1.830021in}{2.007934in}}%
\pgfpathlineto{\pgfqpoint{1.864232in}{2.006135in}}%
\pgfpathclose%
\pgfusepath{fill}%
\end{pgfscope}%
\begin{pgfscope}%
\pgfpathrectangle{\pgfqpoint{0.100000in}{0.100000in}}{\pgfqpoint{3.608454in}{2.310000in}}%
\pgfusepath{clip}%
\pgfsetbuttcap%
\pgfsetmiterjoin%
\definecolor{currentfill}{rgb}{0.000000,0.505882,0.747059}%
\pgfsetfillcolor{currentfill}%
\pgfsetlinewidth{0.000000pt}%
\definecolor{currentstroke}{rgb}{0.000000,0.000000,0.000000}%
\pgfsetstrokecolor{currentstroke}%
\pgfsetstrokeopacity{0.000000}%
\pgfsetdash{}{0pt}%
\pgfpathmoveto{\pgfqpoint{1.728388in}{1.154348in}}%
\pgfpathlineto{\pgfqpoint{1.788727in}{1.150707in}}%
\pgfpathlineto{\pgfqpoint{1.787003in}{1.115616in}}%
\pgfpathlineto{\pgfqpoint{1.752524in}{1.117682in}}%
\pgfpathlineto{\pgfqpoint{1.683362in}{1.122232in}}%
\pgfpathlineto{\pgfqpoint{1.685892in}{1.157286in}}%
\pgfpathlineto{\pgfqpoint{1.728388in}{1.154348in}}%
\pgfpathclose%
\pgfusepath{fill}%
\end{pgfscope}%
\begin{pgfscope}%
\pgfpathrectangle{\pgfqpoint{0.100000in}{0.100000in}}{\pgfqpoint{3.608454in}{2.310000in}}%
\pgfusepath{clip}%
\pgfsetbuttcap%
\pgfsetmiterjoin%
\definecolor{currentfill}{rgb}{0.000000,0.482353,0.758824}%
\pgfsetfillcolor{currentfill}%
\pgfsetlinewidth{0.000000pt}%
\definecolor{currentstroke}{rgb}{0.000000,0.000000,0.000000}%
\pgfsetstrokecolor{currentstroke}%
\pgfsetstrokeopacity{0.000000}%
\pgfsetdash{}{0pt}%
\pgfpathmoveto{\pgfqpoint{2.615996in}{1.634560in}}%
\pgfpathlineto{\pgfqpoint{2.617722in}{1.657792in}}%
\pgfpathlineto{\pgfqpoint{2.615667in}{1.673618in}}%
\pgfpathlineto{\pgfqpoint{2.611600in}{1.688088in}}%
\pgfpathlineto{\pgfqpoint{2.600135in}{1.708748in}}%
\pgfpathlineto{\pgfqpoint{2.592486in}{1.725771in}}%
\pgfpathlineto{\pgfqpoint{2.591900in}{1.732408in}}%
\pgfpathlineto{\pgfqpoint{2.596450in}{1.744135in}}%
\pgfpathlineto{\pgfqpoint{2.618822in}{1.745974in}}%
\pgfpathlineto{\pgfqpoint{2.645951in}{1.748717in}}%
\pgfpathlineto{\pgfqpoint{2.648810in}{1.721346in}}%
\pgfpathlineto{\pgfqpoint{2.676133in}{1.724089in}}%
\pgfpathlineto{\pgfqpoint{2.689796in}{1.725591in}}%
\pgfpathlineto{\pgfqpoint{2.693396in}{1.698299in}}%
\pgfpathlineto{\pgfqpoint{2.696454in}{1.670883in}}%
\pgfpathlineto{\pgfqpoint{2.682759in}{1.669371in}}%
\pgfpathlineto{\pgfqpoint{2.655531in}{1.666251in}}%
\pgfpathlineto{\pgfqpoint{2.658513in}{1.638934in}}%
\pgfpathlineto{\pgfqpoint{2.615996in}{1.634560in}}%
\pgfpathclose%
\pgfusepath{fill}%
\end{pgfscope}%
\begin{pgfscope}%
\pgfpathrectangle{\pgfqpoint{0.100000in}{0.100000in}}{\pgfqpoint{3.608454in}{2.310000in}}%
\pgfusepath{clip}%
\pgfsetbuttcap%
\pgfsetmiterjoin%
\definecolor{currentfill}{rgb}{0.000000,0.372549,0.813725}%
\pgfsetfillcolor{currentfill}%
\pgfsetlinewidth{0.000000pt}%
\definecolor{currentstroke}{rgb}{0.000000,0.000000,0.000000}%
\pgfsetstrokecolor{currentstroke}%
\pgfsetstrokeopacity{0.000000}%
\pgfsetdash{}{0pt}%
\pgfpathmoveto{\pgfqpoint{1.148005in}{1.494799in}}%
\pgfpathlineto{\pgfqpoint{1.136590in}{1.496867in}}%
\pgfpathlineto{\pgfqpoint{1.137161in}{1.499935in}}%
\pgfpathlineto{\pgfqpoint{1.128853in}{1.510727in}}%
\pgfpathlineto{\pgfqpoint{1.131819in}{1.519099in}}%
\pgfpathlineto{\pgfqpoint{1.129198in}{1.527981in}}%
\pgfpathlineto{\pgfqpoint{1.120187in}{1.531417in}}%
\pgfpathlineto{\pgfqpoint{1.112401in}{1.545424in}}%
\pgfpathlineto{\pgfqpoint{1.116941in}{1.552445in}}%
\pgfpathlineto{\pgfqpoint{1.105133in}{1.552914in}}%
\pgfpathlineto{\pgfqpoint{1.090523in}{1.546657in}}%
\pgfpathlineto{\pgfqpoint{1.087943in}{1.551215in}}%
\pgfpathlineto{\pgfqpoint{1.077577in}{1.552756in}}%
\pgfpathlineto{\pgfqpoint{1.075863in}{1.537379in}}%
\pgfpathlineto{\pgfqpoint{1.077321in}{1.529704in}}%
\pgfpathlineto{\pgfqpoint{1.072847in}{1.517633in}}%
\pgfpathlineto{\pgfqpoint{1.061605in}{1.511149in}}%
\pgfpathlineto{\pgfqpoint{1.034454in}{1.516384in}}%
\pgfpathlineto{\pgfqpoint{0.960301in}{1.531851in}}%
\pgfpathlineto{\pgfqpoint{0.978586in}{1.616384in}}%
\pgfpathlineto{\pgfqpoint{1.051249in}{1.601475in}}%
\pgfpathlineto{\pgfqpoint{1.070321in}{1.603891in}}%
\pgfpathlineto{\pgfqpoint{1.089153in}{1.621193in}}%
\pgfpathlineto{\pgfqpoint{1.101371in}{1.618838in}}%
\pgfpathlineto{\pgfqpoint{1.110470in}{1.623700in}}%
\pgfpathlineto{\pgfqpoint{1.117387in}{1.618454in}}%
\pgfpathlineto{\pgfqpoint{1.121923in}{1.621467in}}%
\pgfpathlineto{\pgfqpoint{1.132825in}{1.619941in}}%
\pgfpathlineto{\pgfqpoint{1.136553in}{1.614347in}}%
\pgfpathlineto{\pgfqpoint{1.144102in}{1.610430in}}%
\pgfpathlineto{\pgfqpoint{1.143956in}{1.603549in}}%
\pgfpathlineto{\pgfqpoint{1.146992in}{1.596437in}}%
\pgfpathlineto{\pgfqpoint{1.157543in}{1.601688in}}%
\pgfpathlineto{\pgfqpoint{1.154007in}{1.581932in}}%
\pgfpathlineto{\pgfqpoint{1.186062in}{1.576054in}}%
\pgfpathlineto{\pgfqpoint{1.215522in}{1.571272in}}%
\pgfpathlineto{\pgfqpoint{1.213136in}{1.556950in}}%
\pgfpathlineto{\pgfqpoint{1.203263in}{1.558508in}}%
\pgfpathlineto{\pgfqpoint{1.198089in}{1.561222in}}%
\pgfpathlineto{\pgfqpoint{1.190708in}{1.558686in}}%
\pgfpathlineto{\pgfqpoint{1.173397in}{1.557849in}}%
\pgfpathlineto{\pgfqpoint{1.167943in}{1.559368in}}%
\pgfpathlineto{\pgfqpoint{1.158113in}{1.555519in}}%
\pgfpathlineto{\pgfqpoint{1.148005in}{1.494799in}}%
\pgfpathclose%
\pgfusepath{fill}%
\end{pgfscope}%
\begin{pgfscope}%
\pgfpathrectangle{\pgfqpoint{0.100000in}{0.100000in}}{\pgfqpoint{3.608454in}{2.310000in}}%
\pgfusepath{clip}%
\pgfsetbuttcap%
\pgfsetmiterjoin%
\definecolor{currentfill}{rgb}{0.000000,0.286275,0.856863}%
\pgfsetfillcolor{currentfill}%
\pgfsetlinewidth{0.000000pt}%
\definecolor{currentstroke}{rgb}{0.000000,0.000000,0.000000}%
\pgfsetstrokecolor{currentstroke}%
\pgfsetstrokeopacity{0.000000}%
\pgfsetdash{}{0pt}%
\pgfpathmoveto{\pgfqpoint{2.056854in}{1.480886in}}%
\pgfpathlineto{\pgfqpoint{2.019155in}{1.481378in}}%
\pgfpathlineto{\pgfqpoint{2.019500in}{1.502021in}}%
\pgfpathlineto{\pgfqpoint{1.992970in}{1.502588in}}%
\pgfpathlineto{\pgfqpoint{1.994270in}{1.557682in}}%
\pgfpathlineto{\pgfqpoint{2.021531in}{1.557101in}}%
\pgfpathlineto{\pgfqpoint{2.021430in}{1.552509in}}%
\pgfpathlineto{\pgfqpoint{2.040513in}{1.552171in}}%
\pgfpathlineto{\pgfqpoint{2.043988in}{1.561009in}}%
\pgfpathlineto{\pgfqpoint{2.039668in}{1.566696in}}%
\pgfpathlineto{\pgfqpoint{2.067148in}{1.566188in}}%
\pgfpathlineto{\pgfqpoint{2.073952in}{1.566162in}}%
\pgfpathlineto{\pgfqpoint{2.073777in}{1.545495in}}%
\pgfpathlineto{\pgfqpoint{2.077274in}{1.537859in}}%
\pgfpathlineto{\pgfqpoint{2.097619in}{1.537721in}}%
\pgfpathlineto{\pgfqpoint{2.097514in}{1.510367in}}%
\pgfpathlineto{\pgfqpoint{2.083971in}{1.510379in}}%
\pgfpathlineto{\pgfqpoint{2.083784in}{1.490019in}}%
\pgfpathlineto{\pgfqpoint{2.058174in}{1.490199in}}%
\pgfpathlineto{\pgfqpoint{2.056854in}{1.480886in}}%
\pgfpathclose%
\pgfusepath{fill}%
\end{pgfscope}%
\begin{pgfscope}%
\pgfpathrectangle{\pgfqpoint{0.100000in}{0.100000in}}{\pgfqpoint{3.608454in}{2.310000in}}%
\pgfusepath{clip}%
\pgfsetbuttcap%
\pgfsetmiterjoin%
\definecolor{currentfill}{rgb}{0.000000,0.490196,0.754902}%
\pgfsetfillcolor{currentfill}%
\pgfsetlinewidth{0.000000pt}%
\definecolor{currentstroke}{rgb}{0.000000,0.000000,0.000000}%
\pgfsetstrokecolor{currentstroke}%
\pgfsetstrokeopacity{0.000000}%
\pgfsetdash{}{0pt}%
\pgfpathmoveto{\pgfqpoint{2.293781in}{1.726557in}}%
\pgfpathlineto{\pgfqpoint{2.294773in}{1.699116in}}%
\pgfpathlineto{\pgfqpoint{2.253472in}{1.697776in}}%
\pgfpathlineto{\pgfqpoint{2.252724in}{1.724220in}}%
\pgfpathlineto{\pgfqpoint{2.238913in}{1.724992in}}%
\pgfpathlineto{\pgfqpoint{2.218568in}{1.724511in}}%
\pgfpathlineto{\pgfqpoint{2.217878in}{1.773970in}}%
\pgfpathlineto{\pgfqpoint{2.217797in}{1.777469in}}%
\pgfpathlineto{\pgfqpoint{2.231724in}{1.779996in}}%
\pgfpathlineto{\pgfqpoint{2.231530in}{1.786839in}}%
\pgfpathlineto{\pgfqpoint{2.241126in}{1.786250in}}%
\pgfpathlineto{\pgfqpoint{2.247155in}{1.782906in}}%
\pgfpathlineto{\pgfqpoint{2.258667in}{1.780519in}}%
\pgfpathlineto{\pgfqpoint{2.264630in}{1.772555in}}%
\pgfpathlineto{\pgfqpoint{2.273713in}{1.770000in}}%
\pgfpathlineto{\pgfqpoint{2.281318in}{1.764757in}}%
\pgfpathlineto{\pgfqpoint{2.285525in}{1.753739in}}%
\pgfpathlineto{\pgfqpoint{2.272950in}{1.753323in}}%
\pgfpathlineto{\pgfqpoint{2.273836in}{1.725888in}}%
\pgfpathlineto{\pgfqpoint{2.293781in}{1.726557in}}%
\pgfpathclose%
\pgfusepath{fill}%
\end{pgfscope}%
\begin{pgfscope}%
\pgfpathrectangle{\pgfqpoint{0.100000in}{0.100000in}}{\pgfqpoint{3.608454in}{2.310000in}}%
\pgfusepath{clip}%
\pgfsetbuttcap%
\pgfsetmiterjoin%
\definecolor{currentfill}{rgb}{0.000000,0.470588,0.764706}%
\pgfsetfillcolor{currentfill}%
\pgfsetlinewidth{0.000000pt}%
\definecolor{currentstroke}{rgb}{0.000000,0.000000,0.000000}%
\pgfsetstrokecolor{currentstroke}%
\pgfsetstrokeopacity{0.000000}%
\pgfsetdash{}{0pt}%
\pgfpathmoveto{\pgfqpoint{1.944686in}{1.753364in}}%
\pgfpathlineto{\pgfqpoint{1.999302in}{1.751993in}}%
\pgfpathlineto{\pgfqpoint{1.998767in}{1.724416in}}%
\pgfpathlineto{\pgfqpoint{1.957610in}{1.725492in}}%
\pgfpathlineto{\pgfqpoint{1.943727in}{1.725985in}}%
\pgfpathlineto{\pgfqpoint{1.944686in}{1.753364in}}%
\pgfpathclose%
\pgfusepath{fill}%
\end{pgfscope}%
\begin{pgfscope}%
\pgfpathrectangle{\pgfqpoint{0.100000in}{0.100000in}}{\pgfqpoint{3.608454in}{2.310000in}}%
\pgfusepath{clip}%
\pgfsetbuttcap%
\pgfsetmiterjoin%
\definecolor{currentfill}{rgb}{0.000000,0.501961,0.749020}%
\pgfsetfillcolor{currentfill}%
\pgfsetlinewidth{0.000000pt}%
\definecolor{currentstroke}{rgb}{0.000000,0.000000,0.000000}%
\pgfsetstrokecolor{currentstroke}%
\pgfsetstrokeopacity{0.000000}%
\pgfsetdash{}{0pt}%
\pgfpathmoveto{\pgfqpoint{2.086211in}{1.418729in}}%
\pgfpathlineto{\pgfqpoint{2.086114in}{1.391298in}}%
\pgfpathlineto{\pgfqpoint{2.072459in}{1.391367in}}%
\pgfpathlineto{\pgfqpoint{2.058778in}{1.391470in}}%
\pgfpathlineto{\pgfqpoint{2.058697in}{1.384598in}}%
\pgfpathlineto{\pgfqpoint{2.043672in}{1.384754in}}%
\pgfpathlineto{\pgfqpoint{2.031304in}{1.384920in}}%
\pgfpathlineto{\pgfqpoint{2.031759in}{1.419257in}}%
\pgfpathlineto{\pgfqpoint{2.086211in}{1.418729in}}%
\pgfpathclose%
\pgfusepath{fill}%
\end{pgfscope}%
\begin{pgfscope}%
\pgfpathrectangle{\pgfqpoint{0.100000in}{0.100000in}}{\pgfqpoint{3.608454in}{2.310000in}}%
\pgfusepath{clip}%
\pgfsetbuttcap%
\pgfsetmiterjoin%
\definecolor{currentfill}{rgb}{0.000000,0.458824,0.770588}%
\pgfsetfillcolor{currentfill}%
\pgfsetlinewidth{0.000000pt}%
\definecolor{currentstroke}{rgb}{0.000000,0.000000,0.000000}%
\pgfsetstrokecolor{currentstroke}%
\pgfsetstrokeopacity{0.000000}%
\pgfsetdash{}{0pt}%
\pgfpathmoveto{\pgfqpoint{2.414986in}{1.743088in}}%
\pgfpathlineto{\pgfqpoint{2.416473in}{1.716189in}}%
\pgfpathlineto{\pgfqpoint{2.405834in}{1.715444in}}%
\pgfpathlineto{\pgfqpoint{2.400634in}{1.726438in}}%
\pgfpathlineto{\pgfqpoint{2.395001in}{1.732218in}}%
\pgfpathlineto{\pgfqpoint{2.394954in}{1.739243in}}%
\pgfpathlineto{\pgfqpoint{2.389968in}{1.750300in}}%
\pgfpathlineto{\pgfqpoint{2.396439in}{1.763257in}}%
\pgfpathlineto{\pgfqpoint{2.373196in}{1.761948in}}%
\pgfpathlineto{\pgfqpoint{2.372249in}{1.775882in}}%
\pgfpathlineto{\pgfqpoint{2.344731in}{1.774369in}}%
\pgfpathlineto{\pgfqpoint{2.344426in}{1.781264in}}%
\pgfpathlineto{\pgfqpoint{2.337567in}{1.780978in}}%
\pgfpathlineto{\pgfqpoint{2.334657in}{1.843071in}}%
\pgfpathlineto{\pgfqpoint{2.334351in}{1.850010in}}%
\pgfpathlineto{\pgfqpoint{2.348110in}{1.850512in}}%
\pgfpathlineto{\pgfqpoint{2.383604in}{1.852652in}}%
\pgfpathlineto{\pgfqpoint{2.382793in}{1.866283in}}%
\pgfpathlineto{\pgfqpoint{2.417000in}{1.868413in}}%
\pgfpathlineto{\pgfqpoint{2.417621in}{1.861430in}}%
\pgfpathlineto{\pgfqpoint{2.419246in}{1.833798in}}%
\pgfpathlineto{\pgfqpoint{2.430551in}{1.834455in}}%
\pgfpathlineto{\pgfqpoint{2.432784in}{1.799789in}}%
\pgfpathlineto{\pgfqpoint{2.434927in}{1.765090in}}%
\pgfpathlineto{\pgfqpoint{2.413744in}{1.763959in}}%
\pgfpathlineto{\pgfqpoint{2.414986in}{1.743088in}}%
\pgfpathclose%
\pgfusepath{fill}%
\end{pgfscope}%
\begin{pgfscope}%
\pgfpathrectangle{\pgfqpoint{0.100000in}{0.100000in}}{\pgfqpoint{3.608454in}{2.310000in}}%
\pgfusepath{clip}%
\pgfsetbuttcap%
\pgfsetmiterjoin%
\definecolor{currentfill}{rgb}{0.000000,0.894118,0.552941}%
\pgfsetfillcolor{currentfill}%
\pgfsetlinewidth{0.000000pt}%
\definecolor{currentstroke}{rgb}{0.000000,0.000000,0.000000}%
\pgfsetstrokecolor{currentstroke}%
\pgfsetstrokeopacity{0.000000}%
\pgfsetdash{}{0pt}%
\pgfpathmoveto{\pgfqpoint{0.531886in}{1.661402in}}%
\pgfpathlineto{\pgfqpoint{0.522050in}{1.660114in}}%
\pgfpathlineto{\pgfqpoint{0.515901in}{1.652222in}}%
\pgfpathlineto{\pgfqpoint{0.512624in}{1.653371in}}%
\pgfpathlineto{\pgfqpoint{0.504486in}{1.647574in}}%
\pgfpathlineto{\pgfqpoint{0.490498in}{1.651799in}}%
\pgfpathlineto{\pgfqpoint{0.488413in}{1.645315in}}%
\pgfpathlineto{\pgfqpoint{0.436761in}{1.661477in}}%
\pgfpathlineto{\pgfqpoint{0.441276in}{1.674970in}}%
\pgfpathlineto{\pgfqpoint{0.442069in}{1.688570in}}%
\pgfpathlineto{\pgfqpoint{0.444780in}{1.696306in}}%
\pgfpathlineto{\pgfqpoint{0.441132in}{1.699873in}}%
\pgfpathlineto{\pgfqpoint{0.442506in}{1.705022in}}%
\pgfpathlineto{\pgfqpoint{0.448092in}{1.709303in}}%
\pgfpathlineto{\pgfqpoint{0.454296in}{1.709584in}}%
\pgfpathlineto{\pgfqpoint{0.458930in}{1.713085in}}%
\pgfpathlineto{\pgfqpoint{0.469494in}{1.715995in}}%
\pgfpathlineto{\pgfqpoint{0.468741in}{1.723271in}}%
\pgfpathlineto{\pgfqpoint{0.475334in}{1.730099in}}%
\pgfpathlineto{\pgfqpoint{0.485004in}{1.744684in}}%
\pgfpathlineto{\pgfqpoint{0.489189in}{1.745779in}}%
\pgfpathlineto{\pgfqpoint{0.494828in}{1.757961in}}%
\pgfpathlineto{\pgfqpoint{0.561272in}{1.737785in}}%
\pgfpathlineto{\pgfqpoint{0.555668in}{1.716515in}}%
\pgfpathlineto{\pgfqpoint{0.544816in}{1.681752in}}%
\pgfpathlineto{\pgfqpoint{0.535033in}{1.684697in}}%
\pgfpathlineto{\pgfqpoint{0.534540in}{1.676936in}}%
\pgfpathlineto{\pgfqpoint{0.540383in}{1.672372in}}%
\pgfpathlineto{\pgfqpoint{0.537339in}{1.664898in}}%
\pgfpathlineto{\pgfqpoint{0.531886in}{1.661402in}}%
\pgfpathclose%
\pgfusepath{fill}%
\end{pgfscope}%
\begin{pgfscope}%
\pgfpathrectangle{\pgfqpoint{0.100000in}{0.100000in}}{\pgfqpoint{3.608454in}{2.310000in}}%
\pgfusepath{clip}%
\pgfsetbuttcap%
\pgfsetmiterjoin%
\definecolor{currentfill}{rgb}{0.000000,0.415686,0.792157}%
\pgfsetfillcolor{currentfill}%
\pgfsetlinewidth{0.000000pt}%
\definecolor{currentstroke}{rgb}{0.000000,0.000000,0.000000}%
\pgfsetstrokecolor{currentstroke}%
\pgfsetstrokeopacity{0.000000}%
\pgfsetdash{}{0pt}%
\pgfpathmoveto{\pgfqpoint{2.688201in}{1.338715in}}%
\pgfpathlineto{\pgfqpoint{2.687551in}{1.344405in}}%
\pgfpathlineto{\pgfqpoint{2.680182in}{1.348259in}}%
\pgfpathlineto{\pgfqpoint{2.679482in}{1.355109in}}%
\pgfpathlineto{\pgfqpoint{2.683930in}{1.355595in}}%
\pgfpathlineto{\pgfqpoint{2.693342in}{1.363630in}}%
\pgfpathlineto{\pgfqpoint{2.708117in}{1.365276in}}%
\pgfpathlineto{\pgfqpoint{2.710112in}{1.347801in}}%
\pgfpathlineto{\pgfqpoint{2.720376in}{1.354526in}}%
\pgfpathlineto{\pgfqpoint{2.726746in}{1.347306in}}%
\pgfpathlineto{\pgfqpoint{2.718732in}{1.341235in}}%
\pgfpathlineto{\pgfqpoint{2.713052in}{1.339727in}}%
\pgfpathlineto{\pgfqpoint{2.707002in}{1.331992in}}%
\pgfpathlineto{\pgfqpoint{2.697292in}{1.333124in}}%
\pgfpathlineto{\pgfqpoint{2.697081in}{1.338065in}}%
\pgfpathlineto{\pgfqpoint{2.688201in}{1.338715in}}%
\pgfpathclose%
\pgfusepath{fill}%
\end{pgfscope}%
\begin{pgfscope}%
\pgfpathrectangle{\pgfqpoint{0.100000in}{0.100000in}}{\pgfqpoint{3.608454in}{2.310000in}}%
\pgfusepath{clip}%
\pgfsetbuttcap%
\pgfsetmiterjoin%
\definecolor{currentfill}{rgb}{0.000000,0.435294,0.782353}%
\pgfsetfillcolor{currentfill}%
\pgfsetlinewidth{0.000000pt}%
\definecolor{currentstroke}{rgb}{0.000000,0.000000,0.000000}%
\pgfsetstrokecolor{currentstroke}%
\pgfsetstrokeopacity{0.000000}%
\pgfsetdash{}{0pt}%
\pgfpathmoveto{\pgfqpoint{1.912258in}{1.559919in}}%
\pgfpathlineto{\pgfqpoint{1.911513in}{1.532437in}}%
\pgfpathlineto{\pgfqpoint{1.884664in}{1.533508in}}%
\pgfpathlineto{\pgfqpoint{1.857457in}{1.534695in}}%
\pgfpathlineto{\pgfqpoint{1.858208in}{1.562101in}}%
\pgfpathlineto{\pgfqpoint{1.830208in}{1.563450in}}%
\pgfpathlineto{\pgfqpoint{1.831612in}{1.590776in}}%
\pgfpathlineto{\pgfqpoint{1.858155in}{1.589614in}}%
\pgfpathlineto{\pgfqpoint{1.885929in}{1.588436in}}%
\pgfpathlineto{\pgfqpoint{1.884919in}{1.560911in}}%
\pgfpathlineto{\pgfqpoint{1.912258in}{1.559919in}}%
\pgfpathclose%
\pgfusepath{fill}%
\end{pgfscope}%
\begin{pgfscope}%
\pgfpathrectangle{\pgfqpoint{0.100000in}{0.100000in}}{\pgfqpoint{3.608454in}{2.310000in}}%
\pgfusepath{clip}%
\pgfsetbuttcap%
\pgfsetmiterjoin%
\definecolor{currentfill}{rgb}{0.000000,0.313725,0.843137}%
\pgfsetfillcolor{currentfill}%
\pgfsetlinewidth{0.000000pt}%
\definecolor{currentstroke}{rgb}{0.000000,0.000000,0.000000}%
\pgfsetstrokecolor{currentstroke}%
\pgfsetstrokeopacity{0.000000}%
\pgfsetdash{}{0pt}%
\pgfpathmoveto{\pgfqpoint{1.968280in}{1.606524in}}%
\pgfpathlineto{\pgfqpoint{1.988757in}{1.605977in}}%
\pgfpathlineto{\pgfqpoint{1.989587in}{1.638444in}}%
\pgfpathlineto{\pgfqpoint{2.001679in}{1.633660in}}%
\pgfpathlineto{\pgfqpoint{2.002309in}{1.663696in}}%
\pgfpathlineto{\pgfqpoint{2.022652in}{1.663329in}}%
\pgfpathlineto{\pgfqpoint{2.019320in}{1.660716in}}%
\pgfpathlineto{\pgfqpoint{2.019889in}{1.653420in}}%
\pgfpathlineto{\pgfqpoint{2.017541in}{1.649592in}}%
\pgfpathlineto{\pgfqpoint{2.056914in}{1.649093in}}%
\pgfpathlineto{\pgfqpoint{2.056624in}{1.621513in}}%
\pgfpathlineto{\pgfqpoint{2.063540in}{1.621479in}}%
\pgfpathlineto{\pgfqpoint{2.067679in}{1.614586in}}%
\pgfpathlineto{\pgfqpoint{2.067449in}{1.593735in}}%
\pgfpathlineto{\pgfqpoint{2.067148in}{1.566188in}}%
\pgfpathlineto{\pgfqpoint{2.039668in}{1.566696in}}%
\pgfpathlineto{\pgfqpoint{2.040288in}{1.575134in}}%
\pgfpathlineto{\pgfqpoint{2.036665in}{1.576098in}}%
\pgfpathlineto{\pgfqpoint{2.027934in}{1.590280in}}%
\pgfpathlineto{\pgfqpoint{2.027350in}{1.599345in}}%
\pgfpathlineto{\pgfqpoint{2.005666in}{1.599833in}}%
\pgfpathlineto{\pgfqpoint{2.000030in}{1.598859in}}%
\pgfpathlineto{\pgfqpoint{1.999669in}{1.585100in}}%
\pgfpathlineto{\pgfqpoint{1.988196in}{1.585344in}}%
\pgfpathlineto{\pgfqpoint{1.967706in}{1.585865in}}%
\pgfpathlineto{\pgfqpoint{1.968280in}{1.606524in}}%
\pgfpathclose%
\pgfusepath{fill}%
\end{pgfscope}%
\begin{pgfscope}%
\pgfpathrectangle{\pgfqpoint{0.100000in}{0.100000in}}{\pgfqpoint{3.608454in}{2.310000in}}%
\pgfusepath{clip}%
\pgfsetbuttcap%
\pgfsetmiterjoin%
\definecolor{currentfill}{rgb}{0.000000,0.650980,0.674510}%
\pgfsetfillcolor{currentfill}%
\pgfsetlinewidth{0.000000pt}%
\definecolor{currentstroke}{rgb}{0.000000,0.000000,0.000000}%
\pgfsetstrokecolor{currentstroke}%
\pgfsetstrokeopacity{0.000000}%
\pgfsetdash{}{0pt}%
\pgfpathmoveto{\pgfqpoint{3.051321in}{0.914811in}}%
\pgfpathlineto{\pgfqpoint{3.047304in}{0.907519in}}%
\pgfpathlineto{\pgfqpoint{3.038277in}{0.916697in}}%
\pgfpathlineto{\pgfqpoint{3.022585in}{0.898037in}}%
\pgfpathlineto{\pgfqpoint{3.015378in}{0.900512in}}%
\pgfpathlineto{\pgfqpoint{3.011589in}{0.908504in}}%
\pgfpathlineto{\pgfqpoint{3.010014in}{0.919433in}}%
\pgfpathlineto{\pgfqpoint{3.003885in}{0.926116in}}%
\pgfpathlineto{\pgfqpoint{3.003733in}{0.931980in}}%
\pgfpathlineto{\pgfqpoint{2.999992in}{0.934399in}}%
\pgfpathlineto{\pgfqpoint{2.988444in}{0.938374in}}%
\pgfpathlineto{\pgfqpoint{2.984051in}{0.944289in}}%
\pgfpathlineto{\pgfqpoint{3.005840in}{0.971004in}}%
\pgfpathlineto{\pgfqpoint{3.010939in}{0.970937in}}%
\pgfpathlineto{\pgfqpoint{3.018836in}{0.966601in}}%
\pgfpathlineto{\pgfqpoint{3.034445in}{0.961587in}}%
\pgfpathlineto{\pgfqpoint{3.047350in}{0.952216in}}%
\pgfpathlineto{\pgfqpoint{3.030466in}{0.937480in}}%
\pgfpathlineto{\pgfqpoint{3.038279in}{0.928637in}}%
\pgfpathlineto{\pgfqpoint{3.044811in}{0.924182in}}%
\pgfpathlineto{\pgfqpoint{3.051321in}{0.914811in}}%
\pgfpathclose%
\pgfusepath{fill}%
\end{pgfscope}%
\begin{pgfscope}%
\pgfpathrectangle{\pgfqpoint{0.100000in}{0.100000in}}{\pgfqpoint{3.608454in}{2.310000in}}%
\pgfusepath{clip}%
\pgfsetbuttcap%
\pgfsetmiterjoin%
\definecolor{currentfill}{rgb}{0.000000,0.478431,0.760784}%
\pgfsetfillcolor{currentfill}%
\pgfsetlinewidth{0.000000pt}%
\definecolor{currentstroke}{rgb}{0.000000,0.000000,0.000000}%
\pgfsetstrokecolor{currentstroke}%
\pgfsetstrokeopacity{0.000000}%
\pgfsetdash{}{0pt}%
\pgfpathmoveto{\pgfqpoint{1.046091in}{0.176022in}}%
\pgfpathlineto{\pgfqpoint{1.046239in}{0.179505in}}%
\pgfpathlineto{\pgfqpoint{1.048650in}{0.177612in}}%
\pgfpathlineto{\pgfqpoint{1.046091in}{0.176022in}}%
\pgfpathclose%
\pgfusepath{fill}%
\end{pgfscope}%
\begin{pgfscope}%
\pgfpathrectangle{\pgfqpoint{0.100000in}{0.100000in}}{\pgfqpoint{3.608454in}{2.310000in}}%
\pgfusepath{clip}%
\pgfsetbuttcap%
\pgfsetmiterjoin%
\definecolor{currentfill}{rgb}{0.000000,0.478431,0.760784}%
\pgfsetfillcolor{currentfill}%
\pgfsetlinewidth{0.000000pt}%
\definecolor{currentstroke}{rgb}{0.000000,0.000000,0.000000}%
\pgfsetstrokecolor{currentstroke}%
\pgfsetstrokeopacity{0.000000}%
\pgfsetdash{}{0pt}%
\pgfpathmoveto{\pgfqpoint{1.042447in}{0.162554in}}%
\pgfpathlineto{\pgfqpoint{1.043313in}{0.165192in}}%
\pgfpathlineto{\pgfqpoint{1.044966in}{0.165456in}}%
\pgfpathlineto{\pgfqpoint{1.044338in}{0.161816in}}%
\pgfpathlineto{\pgfqpoint{1.042447in}{0.162554in}}%
\pgfpathclose%
\pgfusepath{fill}%
\end{pgfscope}%
\begin{pgfscope}%
\pgfpathrectangle{\pgfqpoint{0.100000in}{0.100000in}}{\pgfqpoint{3.608454in}{2.310000in}}%
\pgfusepath{clip}%
\pgfsetbuttcap%
\pgfsetmiterjoin%
\definecolor{currentfill}{rgb}{0.000000,0.478431,0.760784}%
\pgfsetfillcolor{currentfill}%
\pgfsetlinewidth{0.000000pt}%
\definecolor{currentstroke}{rgb}{0.000000,0.000000,0.000000}%
\pgfsetstrokecolor{currentstroke}%
\pgfsetstrokeopacity{0.000000}%
\pgfsetdash{}{0pt}%
\pgfpathmoveto{\pgfqpoint{1.045943in}{0.145022in}}%
\pgfpathlineto{\pgfqpoint{1.044467in}{0.150857in}}%
\pgfpathlineto{\pgfqpoint{1.045363in}{0.152527in}}%
\pgfpathlineto{\pgfqpoint{1.043425in}{0.154046in}}%
\pgfpathlineto{\pgfqpoint{1.045163in}{0.157939in}}%
\pgfpathlineto{\pgfqpoint{1.045050in}{0.160374in}}%
\pgfpathlineto{\pgfqpoint{1.047193in}{0.161279in}}%
\pgfpathlineto{\pgfqpoint{1.047667in}{0.157760in}}%
\pgfpathlineto{\pgfqpoint{1.045744in}{0.156113in}}%
\pgfpathlineto{\pgfqpoint{1.046983in}{0.154946in}}%
\pgfpathlineto{\pgfqpoint{1.047011in}{0.149220in}}%
\pgfpathlineto{\pgfqpoint{1.049976in}{0.148483in}}%
\pgfpathlineto{\pgfqpoint{1.045943in}{0.145022in}}%
\pgfpathclose%
\pgfusepath{fill}%
\end{pgfscope}%
\begin{pgfscope}%
\pgfpathrectangle{\pgfqpoint{0.100000in}{0.100000in}}{\pgfqpoint{3.608454in}{2.310000in}}%
\pgfusepath{clip}%
\pgfsetbuttcap%
\pgfsetmiterjoin%
\definecolor{currentfill}{rgb}{0.000000,0.478431,0.760784}%
\pgfsetfillcolor{currentfill}%
\pgfsetlinewidth{0.000000pt}%
\definecolor{currentstroke}{rgb}{0.000000,0.000000,0.000000}%
\pgfsetstrokecolor{currentstroke}%
\pgfsetstrokeopacity{0.000000}%
\pgfsetdash{}{0pt}%
\pgfpathmoveto{\pgfqpoint{1.071810in}{0.150553in}}%
\pgfpathlineto{\pgfqpoint{1.072025in}{0.147586in}}%
\pgfpathlineto{\pgfqpoint{1.070746in}{0.144824in}}%
\pgfpathlineto{\pgfqpoint{1.069328in}{0.144797in}}%
\pgfpathlineto{\pgfqpoint{1.068043in}{0.147923in}}%
\pgfpathlineto{\pgfqpoint{1.069455in}{0.148655in}}%
\pgfpathlineto{\pgfqpoint{1.069622in}{0.151729in}}%
\pgfpathlineto{\pgfqpoint{1.071810in}{0.150553in}}%
\pgfpathclose%
\pgfusepath{fill}%
\end{pgfscope}%
\begin{pgfscope}%
\pgfpathrectangle{\pgfqpoint{0.100000in}{0.100000in}}{\pgfqpoint{3.608454in}{2.310000in}}%
\pgfusepath{clip}%
\pgfsetbuttcap%
\pgfsetmiterjoin%
\definecolor{currentfill}{rgb}{0.000000,0.478431,0.760784}%
\pgfsetfillcolor{currentfill}%
\pgfsetlinewidth{0.000000pt}%
\definecolor{currentstroke}{rgb}{0.000000,0.000000,0.000000}%
\pgfsetstrokecolor{currentstroke}%
\pgfsetstrokeopacity{0.000000}%
\pgfsetdash{}{0pt}%
\pgfpathmoveto{\pgfqpoint{1.065482in}{0.169372in}}%
\pgfpathlineto{\pgfqpoint{1.067405in}{0.165345in}}%
\pgfpathlineto{\pgfqpoint{1.069320in}{0.165633in}}%
\pgfpathlineto{\pgfqpoint{1.070834in}{0.168040in}}%
\pgfpathlineto{\pgfqpoint{1.070260in}{0.170374in}}%
\pgfpathlineto{\pgfqpoint{1.072944in}{0.170978in}}%
\pgfpathlineto{\pgfqpoint{1.073238in}{0.172315in}}%
\pgfpathlineto{\pgfqpoint{1.075481in}{0.173524in}}%
\pgfpathlineto{\pgfqpoint{1.078315in}{0.173884in}}%
\pgfpathlineto{\pgfqpoint{1.079838in}{0.170598in}}%
\pgfpathlineto{\pgfqpoint{1.082038in}{0.171601in}}%
\pgfpathlineto{\pgfqpoint{1.081805in}{0.173502in}}%
\pgfpathlineto{\pgfqpoint{1.085139in}{0.174080in}}%
\pgfpathlineto{\pgfqpoint{1.087857in}{0.173091in}}%
\pgfpathlineto{\pgfqpoint{1.089922in}{0.175542in}}%
\pgfpathlineto{\pgfqpoint{1.090049in}{0.178164in}}%
\pgfpathlineto{\pgfqpoint{1.093019in}{0.176002in}}%
\pgfpathlineto{\pgfqpoint{1.096013in}{0.170850in}}%
\pgfpathlineto{\pgfqpoint{1.096478in}{0.167420in}}%
\pgfpathlineto{\pgfqpoint{1.098614in}{0.165154in}}%
\pgfpathlineto{\pgfqpoint{1.100805in}{0.164982in}}%
\pgfpathlineto{\pgfqpoint{1.101142in}{0.162059in}}%
\pgfpathlineto{\pgfqpoint{1.100030in}{0.159422in}}%
\pgfpathlineto{\pgfqpoint{1.096111in}{0.156382in}}%
\pgfpathlineto{\pgfqpoint{1.095842in}{0.153748in}}%
\pgfpathlineto{\pgfqpoint{1.094580in}{0.151215in}}%
\pgfpathlineto{\pgfqpoint{1.094242in}{0.146717in}}%
\pgfpathlineto{\pgfqpoint{1.093131in}{0.142662in}}%
\pgfpathlineto{\pgfqpoint{1.090632in}{0.141147in}}%
\pgfpathlineto{\pgfqpoint{1.084308in}{0.135553in}}%
\pgfpathlineto{\pgfqpoint{1.081972in}{0.136078in}}%
\pgfpathlineto{\pgfqpoint{1.079389in}{0.134752in}}%
\pgfpathlineto{\pgfqpoint{1.077089in}{0.136527in}}%
\pgfpathlineto{\pgfqpoint{1.074765in}{0.136279in}}%
\pgfpathlineto{\pgfqpoint{1.076502in}{0.144158in}}%
\pgfpathlineto{\pgfqpoint{1.076368in}{0.146457in}}%
\pgfpathlineto{\pgfqpoint{1.079906in}{0.148920in}}%
\pgfpathlineto{\pgfqpoint{1.081825in}{0.148421in}}%
\pgfpathlineto{\pgfqpoint{1.081856in}{0.151528in}}%
\pgfpathlineto{\pgfqpoint{1.083827in}{0.155130in}}%
\pgfpathlineto{\pgfqpoint{1.085699in}{0.162221in}}%
\pgfpathlineto{\pgfqpoint{1.083172in}{0.161819in}}%
\pgfpathlineto{\pgfqpoint{1.083697in}{0.159212in}}%
\pgfpathlineto{\pgfqpoint{1.082643in}{0.155646in}}%
\pgfpathlineto{\pgfqpoint{1.078109in}{0.149553in}}%
\pgfpathlineto{\pgfqpoint{1.074518in}{0.148168in}}%
\pgfpathlineto{\pgfqpoint{1.073929in}{0.150391in}}%
\pgfpathlineto{\pgfqpoint{1.070564in}{0.153361in}}%
\pgfpathlineto{\pgfqpoint{1.067603in}{0.152672in}}%
\pgfpathlineto{\pgfqpoint{1.065718in}{0.151421in}}%
\pgfpathlineto{\pgfqpoint{1.067029in}{0.156508in}}%
\pgfpathlineto{\pgfqpoint{1.068970in}{0.159001in}}%
\pgfpathlineto{\pgfqpoint{1.070918in}{0.159489in}}%
\pgfpathlineto{\pgfqpoint{1.072200in}{0.162116in}}%
\pgfpathlineto{\pgfqpoint{1.074427in}{0.163792in}}%
\pgfpathlineto{\pgfqpoint{1.074481in}{0.165957in}}%
\pgfpathlineto{\pgfqpoint{1.077573in}{0.166662in}}%
\pgfpathlineto{\pgfqpoint{1.077538in}{0.168796in}}%
\pgfpathlineto{\pgfqpoint{1.074777in}{0.170290in}}%
\pgfpathlineto{\pgfqpoint{1.071926in}{0.169175in}}%
\pgfpathlineto{\pgfqpoint{1.073508in}{0.167354in}}%
\pgfpathlineto{\pgfqpoint{1.070989in}{0.164095in}}%
\pgfpathlineto{\pgfqpoint{1.067088in}{0.160807in}}%
\pgfpathlineto{\pgfqpoint{1.065060in}{0.165053in}}%
\pgfpathlineto{\pgfqpoint{1.065482in}{0.169372in}}%
\pgfpathclose%
\pgfusepath{fill}%
\end{pgfscope}%
\begin{pgfscope}%
\pgfpathrectangle{\pgfqpoint{0.100000in}{0.100000in}}{\pgfqpoint{3.608454in}{2.310000in}}%
\pgfusepath{clip}%
\pgfsetbuttcap%
\pgfsetmiterjoin%
\definecolor{currentfill}{rgb}{0.000000,0.478431,0.760784}%
\pgfsetfillcolor{currentfill}%
\pgfsetlinewidth{0.000000pt}%
\definecolor{currentstroke}{rgb}{0.000000,0.000000,0.000000}%
\pgfsetstrokecolor{currentstroke}%
\pgfsetstrokeopacity{0.000000}%
\pgfsetdash{}{0pt}%
\pgfpathmoveto{\pgfqpoint{1.041943in}{0.169830in}}%
\pgfpathlineto{\pgfqpoint{1.039521in}{0.170855in}}%
\pgfpathlineto{\pgfqpoint{1.042951in}{0.172853in}}%
\pgfpathlineto{\pgfqpoint{1.043959in}{0.170752in}}%
\pgfpathlineto{\pgfqpoint{1.047379in}{0.171633in}}%
\pgfpathlineto{\pgfqpoint{1.048310in}{0.169620in}}%
\pgfpathlineto{\pgfqpoint{1.045621in}{0.169580in}}%
\pgfpathlineto{\pgfqpoint{1.044188in}{0.167269in}}%
\pgfpathlineto{\pgfqpoint{1.041932in}{0.166429in}}%
\pgfpathlineto{\pgfqpoint{1.040306in}{0.167304in}}%
\pgfpathlineto{\pgfqpoint{1.041943in}{0.169830in}}%
\pgfpathclose%
\pgfusepath{fill}%
\end{pgfscope}%
\begin{pgfscope}%
\pgfpathrectangle{\pgfqpoint{0.100000in}{0.100000in}}{\pgfqpoint{3.608454in}{2.310000in}}%
\pgfusepath{clip}%
\pgfsetbuttcap%
\pgfsetmiterjoin%
\definecolor{currentfill}{rgb}{0.000000,0.478431,0.760784}%
\pgfsetfillcolor{currentfill}%
\pgfsetlinewidth{0.000000pt}%
\definecolor{currentstroke}{rgb}{0.000000,0.000000,0.000000}%
\pgfsetstrokecolor{currentstroke}%
\pgfsetstrokeopacity{0.000000}%
\pgfsetdash{}{0pt}%
\pgfpathmoveto{\pgfqpoint{1.051425in}{0.157660in}}%
\pgfpathlineto{\pgfqpoint{1.051517in}{0.160010in}}%
\pgfpathlineto{\pgfqpoint{1.046502in}{0.163773in}}%
\pgfpathlineto{\pgfqpoint{1.050660in}{0.165847in}}%
\pgfpathlineto{\pgfqpoint{1.049847in}{0.167856in}}%
\pgfpathlineto{\pgfqpoint{1.050846in}{0.170035in}}%
\pgfpathlineto{\pgfqpoint{1.048993in}{0.170836in}}%
\pgfpathlineto{\pgfqpoint{1.047832in}{0.173344in}}%
\pgfpathlineto{\pgfqpoint{1.051430in}{0.175827in}}%
\pgfpathlineto{\pgfqpoint{1.051845in}{0.178802in}}%
\pgfpathlineto{\pgfqpoint{1.053869in}{0.179753in}}%
\pgfpathlineto{\pgfqpoint{1.053809in}{0.182723in}}%
\pgfpathlineto{\pgfqpoint{1.051267in}{0.183712in}}%
\pgfpathlineto{\pgfqpoint{1.045462in}{0.182905in}}%
\pgfpathlineto{\pgfqpoint{1.047279in}{0.185118in}}%
\pgfpathlineto{\pgfqpoint{1.049511in}{0.186205in}}%
\pgfpathlineto{\pgfqpoint{1.052044in}{0.186328in}}%
\pgfpathlineto{\pgfqpoint{1.051609in}{0.190204in}}%
\pgfpathlineto{\pgfqpoint{1.052862in}{0.192676in}}%
\pgfpathlineto{\pgfqpoint{1.054788in}{0.191850in}}%
\pgfpathlineto{\pgfqpoint{1.059113in}{0.188805in}}%
\pgfpathlineto{\pgfqpoint{1.059157in}{0.184625in}}%
\pgfpathlineto{\pgfqpoint{1.057111in}{0.183014in}}%
\pgfpathlineto{\pgfqpoint{1.057987in}{0.179356in}}%
\pgfpathlineto{\pgfqpoint{1.059495in}{0.179632in}}%
\pgfpathlineto{\pgfqpoint{1.061360in}{0.177220in}}%
\pgfpathlineto{\pgfqpoint{1.062648in}{0.172710in}}%
\pgfpathlineto{\pgfqpoint{1.062360in}{0.170277in}}%
\pgfpathlineto{\pgfqpoint{1.060929in}{0.168990in}}%
\pgfpathlineto{\pgfqpoint{1.062884in}{0.167409in}}%
\pgfpathlineto{\pgfqpoint{1.062703in}{0.164914in}}%
\pgfpathlineto{\pgfqpoint{1.060871in}{0.164522in}}%
\pgfpathlineto{\pgfqpoint{1.059752in}{0.166916in}}%
\pgfpathlineto{\pgfqpoint{1.058308in}{0.164637in}}%
\pgfpathlineto{\pgfqpoint{1.061166in}{0.162445in}}%
\pgfpathlineto{\pgfqpoint{1.061172in}{0.160447in}}%
\pgfpathlineto{\pgfqpoint{1.062512in}{0.157360in}}%
\pgfpathlineto{\pgfqpoint{1.060934in}{0.153696in}}%
\pgfpathlineto{\pgfqpoint{1.064018in}{0.154722in}}%
\pgfpathlineto{\pgfqpoint{1.063461in}{0.152526in}}%
\pgfpathlineto{\pgfqpoint{1.061554in}{0.151190in}}%
\pgfpathlineto{\pgfqpoint{1.060026in}{0.151947in}}%
\pgfpathlineto{\pgfqpoint{1.058945in}{0.148455in}}%
\pgfpathlineto{\pgfqpoint{1.061641in}{0.148773in}}%
\pgfpathlineto{\pgfqpoint{1.059952in}{0.145277in}}%
\pgfpathlineto{\pgfqpoint{1.059401in}{0.142596in}}%
\pgfpathlineto{\pgfqpoint{1.057983in}{0.140914in}}%
\pgfpathlineto{\pgfqpoint{1.055057in}{0.141929in}}%
\pgfpathlineto{\pgfqpoint{1.054548in}{0.147234in}}%
\pgfpathlineto{\pgfqpoint{1.052684in}{0.150955in}}%
\pgfpathlineto{\pgfqpoint{1.052499in}{0.153302in}}%
\pgfpathlineto{\pgfqpoint{1.054358in}{0.157385in}}%
\pgfpathlineto{\pgfqpoint{1.051425in}{0.157660in}}%
\pgfpathclose%
\pgfusepath{fill}%
\end{pgfscope}%
\begin{pgfscope}%
\pgfpathrectangle{\pgfqpoint{0.100000in}{0.100000in}}{\pgfqpoint{3.608454in}{2.310000in}}%
\pgfusepath{clip}%
\pgfsetbuttcap%
\pgfsetmiterjoin%
\definecolor{currentfill}{rgb}{0.000000,0.682353,0.658824}%
\pgfsetfillcolor{currentfill}%
\pgfsetlinewidth{0.000000pt}%
\definecolor{currentstroke}{rgb}{0.000000,0.000000,0.000000}%
\pgfsetstrokecolor{currentstroke}%
\pgfsetstrokeopacity{0.000000}%
\pgfsetdash{}{0pt}%
\pgfpathmoveto{\pgfqpoint{0.875234in}{2.151233in}}%
\pgfpathlineto{\pgfqpoint{0.868405in}{2.124498in}}%
\pgfpathlineto{\pgfqpoint{0.863100in}{2.115696in}}%
\pgfpathlineto{\pgfqpoint{0.851613in}{2.114354in}}%
\pgfpathlineto{\pgfqpoint{0.823108in}{2.121596in}}%
\pgfpathlineto{\pgfqpoint{0.790335in}{2.130635in}}%
\pgfpathlineto{\pgfqpoint{0.784697in}{2.127385in}}%
\pgfpathlineto{\pgfqpoint{0.782427in}{2.131997in}}%
\pgfpathlineto{\pgfqpoint{0.774845in}{2.127276in}}%
\pgfpathlineto{\pgfqpoint{0.761418in}{2.129669in}}%
\pgfpathlineto{\pgfqpoint{0.758684in}{2.134138in}}%
\pgfpathlineto{\pgfqpoint{0.762232in}{2.143437in}}%
\pgfpathlineto{\pgfqpoint{0.763393in}{2.155197in}}%
\pgfpathlineto{\pgfqpoint{0.762473in}{2.166902in}}%
\pgfpathlineto{\pgfqpoint{0.767233in}{2.177075in}}%
\pgfpathlineto{\pgfqpoint{0.769524in}{2.184730in}}%
\pgfpathlineto{\pgfqpoint{0.774962in}{2.186885in}}%
\pgfpathlineto{\pgfqpoint{0.778779in}{2.191862in}}%
\pgfpathlineto{\pgfqpoint{0.796683in}{2.187200in}}%
\pgfpathlineto{\pgfqpoint{0.798560in}{2.193839in}}%
\pgfpathlineto{\pgfqpoint{0.807006in}{2.198546in}}%
\pgfpathlineto{\pgfqpoint{0.810303in}{2.197629in}}%
\pgfpathlineto{\pgfqpoint{0.814986in}{2.205959in}}%
\pgfpathlineto{\pgfqpoint{0.820986in}{2.210320in}}%
\pgfpathlineto{\pgfqpoint{0.822836in}{2.217028in}}%
\pgfpathlineto{\pgfqpoint{0.829989in}{2.221231in}}%
\pgfpathlineto{\pgfqpoint{0.836288in}{2.219545in}}%
\pgfpathlineto{\pgfqpoint{0.836363in}{2.217995in}}%
\pgfpathlineto{\pgfqpoint{0.821826in}{2.165426in}}%
\pgfpathlineto{\pgfqpoint{0.875234in}{2.151233in}}%
\pgfpathclose%
\pgfusepath{fill}%
\end{pgfscope}%
\begin{pgfscope}%
\pgfpathrectangle{\pgfqpoint{0.100000in}{0.100000in}}{\pgfqpoint{3.608454in}{2.310000in}}%
\pgfusepath{clip}%
\pgfsetbuttcap%
\pgfsetmiterjoin%
\definecolor{currentfill}{rgb}{0.000000,0.709804,0.645098}%
\pgfsetfillcolor{currentfill}%
\pgfsetlinewidth{0.000000pt}%
\definecolor{currentstroke}{rgb}{0.000000,0.000000,0.000000}%
\pgfsetstrokecolor{currentstroke}%
\pgfsetstrokeopacity{0.000000}%
\pgfsetdash{}{0pt}%
\pgfpathmoveto{\pgfqpoint{2.966480in}{1.429783in}}%
\pgfpathlineto{\pgfqpoint{2.962415in}{1.436530in}}%
\pgfpathlineto{\pgfqpoint{2.951761in}{1.443677in}}%
\pgfpathlineto{\pgfqpoint{2.954171in}{1.450091in}}%
\pgfpathlineto{\pgfqpoint{2.941441in}{1.448688in}}%
\pgfpathlineto{\pgfqpoint{2.938939in}{1.450447in}}%
\pgfpathlineto{\pgfqpoint{2.935564in}{1.458054in}}%
\pgfpathlineto{\pgfqpoint{2.933975in}{1.470774in}}%
\pgfpathlineto{\pgfqpoint{2.938569in}{1.471342in}}%
\pgfpathlineto{\pgfqpoint{2.935644in}{1.495065in}}%
\pgfpathlineto{\pgfqpoint{2.956251in}{1.497406in}}%
\pgfpathlineto{\pgfqpoint{2.977955in}{1.500863in}}%
\pgfpathlineto{\pgfqpoint{2.983438in}{1.466654in}}%
\pgfpathlineto{\pgfqpoint{2.989320in}{1.467586in}}%
\pgfpathlineto{\pgfqpoint{2.992195in}{1.461918in}}%
\pgfpathlineto{\pgfqpoint{2.986701in}{1.454490in}}%
\pgfpathlineto{\pgfqpoint{2.989302in}{1.450132in}}%
\pgfpathlineto{\pgfqpoint{2.987075in}{1.444310in}}%
\pgfpathlineto{\pgfqpoint{2.980855in}{1.444545in}}%
\pgfpathlineto{\pgfqpoint{2.966480in}{1.429783in}}%
\pgfpathclose%
\pgfusepath{fill}%
\end{pgfscope}%
\begin{pgfscope}%
\pgfpathrectangle{\pgfqpoint{0.100000in}{0.100000in}}{\pgfqpoint{3.608454in}{2.310000in}}%
\pgfusepath{clip}%
\pgfsetbuttcap%
\pgfsetmiterjoin%
\definecolor{currentfill}{rgb}{0.000000,0.411765,0.794118}%
\pgfsetfillcolor{currentfill}%
\pgfsetlinewidth{0.000000pt}%
\definecolor{currentstroke}{rgb}{0.000000,0.000000,0.000000}%
\pgfsetstrokecolor{currentstroke}%
\pgfsetstrokeopacity{0.000000}%
\pgfsetdash{}{0pt}%
\pgfpathmoveto{\pgfqpoint{3.466647in}{1.697217in}}%
\pgfpathlineto{\pgfqpoint{3.468281in}{1.705449in}}%
\pgfpathlineto{\pgfqpoint{3.464103in}{1.723936in}}%
\pgfpathlineto{\pgfqpoint{3.456367in}{1.750710in}}%
\pgfpathlineto{\pgfqpoint{3.479909in}{1.757791in}}%
\pgfpathlineto{\pgfqpoint{3.481572in}{1.755490in}}%
\pgfpathlineto{\pgfqpoint{3.495434in}{1.768247in}}%
\pgfpathlineto{\pgfqpoint{3.499874in}{1.758602in}}%
\pgfpathlineto{\pgfqpoint{3.513375in}{1.745585in}}%
\pgfpathlineto{\pgfqpoint{3.518613in}{1.736187in}}%
\pgfpathlineto{\pgfqpoint{3.513987in}{1.733317in}}%
\pgfpathlineto{\pgfqpoint{3.514812in}{1.727935in}}%
\pgfpathlineto{\pgfqpoint{3.510151in}{1.723018in}}%
\pgfpathlineto{\pgfqpoint{3.494576in}{1.717460in}}%
\pgfpathlineto{\pgfqpoint{3.494096in}{1.728843in}}%
\pgfpathlineto{\pgfqpoint{3.490331in}{1.736004in}}%
\pgfpathlineto{\pgfqpoint{3.485585in}{1.730146in}}%
\pgfpathlineto{\pgfqpoint{3.485574in}{1.722330in}}%
\pgfpathlineto{\pgfqpoint{3.488922in}{1.715480in}}%
\pgfpathlineto{\pgfqpoint{3.487219in}{1.706707in}}%
\pgfpathlineto{\pgfqpoint{3.482998in}{1.705766in}}%
\pgfpathlineto{\pgfqpoint{3.466647in}{1.697217in}}%
\pgfpathclose%
\pgfusepath{fill}%
\end{pgfscope}%
\begin{pgfscope}%
\pgfpathrectangle{\pgfqpoint{0.100000in}{0.100000in}}{\pgfqpoint{3.608454in}{2.310000in}}%
\pgfusepath{clip}%
\pgfsetbuttcap%
\pgfsetmiterjoin%
\definecolor{currentfill}{rgb}{0.000000,0.458824,0.770588}%
\pgfsetfillcolor{currentfill}%
\pgfsetlinewidth{0.000000pt}%
\definecolor{currentstroke}{rgb}{0.000000,0.000000,0.000000}%
\pgfsetstrokecolor{currentstroke}%
\pgfsetstrokeopacity{0.000000}%
\pgfsetdash{}{0pt}%
\pgfpathmoveto{\pgfqpoint{2.148148in}{1.837483in}}%
\pgfpathlineto{\pgfqpoint{2.141258in}{1.840834in}}%
\pgfpathlineto{\pgfqpoint{2.120021in}{1.840682in}}%
\pgfpathlineto{\pgfqpoint{2.120016in}{1.847577in}}%
\pgfpathlineto{\pgfqpoint{2.099439in}{1.847533in}}%
\pgfpathlineto{\pgfqpoint{2.099037in}{1.876260in}}%
\pgfpathlineto{\pgfqpoint{2.147223in}{1.876558in}}%
\pgfpathlineto{\pgfqpoint{2.142690in}{1.880249in}}%
\pgfpathlineto{\pgfqpoint{2.175363in}{1.880599in}}%
\pgfpathlineto{\pgfqpoint{2.174916in}{1.893366in}}%
\pgfpathlineto{\pgfqpoint{2.172623in}{1.893374in}}%
\pgfpathlineto{\pgfqpoint{2.172440in}{1.907196in}}%
\pgfpathlineto{\pgfqpoint{2.172320in}{1.914091in}}%
\pgfpathlineto{\pgfqpoint{2.193274in}{1.914649in}}%
\pgfpathlineto{\pgfqpoint{2.193329in}{1.907307in}}%
\pgfpathlineto{\pgfqpoint{2.193529in}{1.893561in}}%
\pgfpathlineto{\pgfqpoint{2.188794in}{1.893513in}}%
\pgfpathlineto{\pgfqpoint{2.189353in}{1.873824in}}%
\pgfpathlineto{\pgfqpoint{2.190082in}{1.834950in}}%
\pgfpathlineto{\pgfqpoint{2.186149in}{1.838564in}}%
\pgfpathlineto{\pgfqpoint{2.176492in}{1.838639in}}%
\pgfpathlineto{\pgfqpoint{2.160005in}{1.848588in}}%
\pgfpathlineto{\pgfqpoint{2.156256in}{1.841131in}}%
\pgfpathlineto{\pgfqpoint{2.148148in}{1.837483in}}%
\pgfpathclose%
\pgfusepath{fill}%
\end{pgfscope}%
\begin{pgfscope}%
\pgfpathrectangle{\pgfqpoint{0.100000in}{0.100000in}}{\pgfqpoint{3.608454in}{2.310000in}}%
\pgfusepath{clip}%
\pgfsetbuttcap%
\pgfsetmiterjoin%
\definecolor{currentfill}{rgb}{0.000000,0.678431,0.660784}%
\pgfsetfillcolor{currentfill}%
\pgfsetlinewidth{0.000000pt}%
\definecolor{currentstroke}{rgb}{0.000000,0.000000,0.000000}%
\pgfsetstrokecolor{currentstroke}%
\pgfsetstrokeopacity{0.000000}%
\pgfsetdash{}{0pt}%
\pgfpathmoveto{\pgfqpoint{1.141273in}{1.664762in}}%
\pgfpathlineto{\pgfqpoint{1.141528in}{1.670743in}}%
\pgfpathlineto{\pgfqpoint{1.137582in}{1.673009in}}%
\pgfpathlineto{\pgfqpoint{1.137407in}{1.682303in}}%
\pgfpathlineto{\pgfqpoint{1.141935in}{1.686685in}}%
\pgfpathlineto{\pgfqpoint{1.143732in}{1.692681in}}%
\pgfpathlineto{\pgfqpoint{1.141976in}{1.698141in}}%
\pgfpathlineto{\pgfqpoint{1.126068in}{1.701146in}}%
\pgfpathlineto{\pgfqpoint{1.126636in}{1.717223in}}%
\pgfpathlineto{\pgfqpoint{1.125111in}{1.725638in}}%
\pgfpathlineto{\pgfqpoint{1.117055in}{1.725647in}}%
\pgfpathlineto{\pgfqpoint{1.116145in}{1.733195in}}%
\pgfpathlineto{\pgfqpoint{1.117477in}{1.741580in}}%
\pgfpathlineto{\pgfqpoint{1.124259in}{1.750336in}}%
\pgfpathlineto{\pgfqpoint{1.182268in}{1.739464in}}%
\pgfpathlineto{\pgfqpoint{1.167956in}{1.660095in}}%
\pgfpathlineto{\pgfqpoint{1.141273in}{1.664762in}}%
\pgfpathclose%
\pgfusepath{fill}%
\end{pgfscope}%
\begin{pgfscope}%
\pgfpathrectangle{\pgfqpoint{0.100000in}{0.100000in}}{\pgfqpoint{3.608454in}{2.310000in}}%
\pgfusepath{clip}%
\pgfsetbuttcap%
\pgfsetmiterjoin%
\definecolor{currentfill}{rgb}{0.000000,0.713725,0.643137}%
\pgfsetfillcolor{currentfill}%
\pgfsetlinewidth{0.000000pt}%
\definecolor{currentstroke}{rgb}{0.000000,0.000000,0.000000}%
\pgfsetstrokecolor{currentstroke}%
\pgfsetstrokeopacity{0.000000}%
\pgfsetdash{}{0pt}%
\pgfpathmoveto{\pgfqpoint{2.193249in}{1.204475in}}%
\pgfpathlineto{\pgfqpoint{2.166408in}{1.204778in}}%
\pgfpathlineto{\pgfqpoint{2.164491in}{1.209439in}}%
\pgfpathlineto{\pgfqpoint{2.130830in}{1.210394in}}%
\pgfpathlineto{\pgfqpoint{2.130827in}{1.263573in}}%
\pgfpathlineto{\pgfqpoint{2.130815in}{1.265402in}}%
\pgfpathlineto{\pgfqpoint{2.154102in}{1.264813in}}%
\pgfpathlineto{\pgfqpoint{2.165424in}{1.263791in}}%
\pgfpathlineto{\pgfqpoint{2.165075in}{1.253073in}}%
\pgfpathlineto{\pgfqpoint{2.181129in}{1.252610in}}%
\pgfpathlineto{\pgfqpoint{2.180991in}{1.248016in}}%
\pgfpathlineto{\pgfqpoint{2.192435in}{1.247763in}}%
\pgfpathlineto{\pgfqpoint{2.193719in}{1.240848in}}%
\pgfpathlineto{\pgfqpoint{2.193249in}{1.204475in}}%
\pgfpathclose%
\pgfusepath{fill}%
\end{pgfscope}%
\begin{pgfscope}%
\pgfpathrectangle{\pgfqpoint{0.100000in}{0.100000in}}{\pgfqpoint{3.608454in}{2.310000in}}%
\pgfusepath{clip}%
\pgfsetbuttcap%
\pgfsetmiterjoin%
\definecolor{currentfill}{rgb}{0.000000,0.552941,0.723529}%
\pgfsetfillcolor{currentfill}%
\pgfsetlinewidth{0.000000pt}%
\definecolor{currentstroke}{rgb}{0.000000,0.000000,0.000000}%
\pgfsetstrokecolor{currentstroke}%
\pgfsetstrokeopacity{0.000000}%
\pgfsetdash{}{0pt}%
\pgfpathmoveto{\pgfqpoint{3.097477in}{1.148400in}}%
\pgfpathlineto{\pgfqpoint{3.057775in}{1.140860in}}%
\pgfpathlineto{\pgfqpoint{3.055178in}{1.146538in}}%
\pgfpathlineto{\pgfqpoint{3.046178in}{1.152956in}}%
\pgfpathlineto{\pgfqpoint{3.037191in}{1.156638in}}%
\pgfpathlineto{\pgfqpoint{3.035098in}{1.163147in}}%
\pgfpathlineto{\pgfqpoint{3.039762in}{1.164964in}}%
\pgfpathlineto{\pgfqpoint{3.039491in}{1.174669in}}%
\pgfpathlineto{\pgfqpoint{3.049089in}{1.181346in}}%
\pgfpathlineto{\pgfqpoint{3.060084in}{1.181854in}}%
\pgfpathlineto{\pgfqpoint{3.054182in}{1.223554in}}%
\pgfpathlineto{\pgfqpoint{3.073773in}{1.226802in}}%
\pgfpathlineto{\pgfqpoint{3.074241in}{1.252094in}}%
\pgfpathlineto{\pgfqpoint{3.074664in}{1.267099in}}%
\pgfpathlineto{\pgfqpoint{3.081510in}{1.270060in}}%
\pgfpathlineto{\pgfqpoint{3.089238in}{1.276542in}}%
\pgfpathlineto{\pgfqpoint{3.095309in}{1.277259in}}%
\pgfpathlineto{\pgfqpoint{3.099433in}{1.272390in}}%
\pgfpathlineto{\pgfqpoint{3.105607in}{1.274029in}}%
\pgfpathlineto{\pgfqpoint{3.104858in}{1.232143in}}%
\pgfpathlineto{\pgfqpoint{3.109908in}{1.233039in}}%
\pgfpathlineto{\pgfqpoint{3.113032in}{1.209517in}}%
\pgfpathlineto{\pgfqpoint{3.106519in}{1.208514in}}%
\pgfpathlineto{\pgfqpoint{3.110627in}{1.182204in}}%
\pgfpathlineto{\pgfqpoint{3.113209in}{1.177564in}}%
\pgfpathlineto{\pgfqpoint{3.093928in}{1.174159in}}%
\pgfpathlineto{\pgfqpoint{3.097477in}{1.148400in}}%
\pgfpathclose%
\pgfusepath{fill}%
\end{pgfscope}%
\begin{pgfscope}%
\pgfpathrectangle{\pgfqpoint{0.100000in}{0.100000in}}{\pgfqpoint{3.608454in}{2.310000in}}%
\pgfusepath{clip}%
\pgfsetbuttcap%
\pgfsetmiterjoin%
\definecolor{currentfill}{rgb}{0.000000,0.490196,0.754902}%
\pgfsetfillcolor{currentfill}%
\pgfsetlinewidth{0.000000pt}%
\definecolor{currentstroke}{rgb}{0.000000,0.000000,0.000000}%
\pgfsetstrokecolor{currentstroke}%
\pgfsetstrokeopacity{0.000000}%
\pgfsetdash{}{0pt}%
\pgfpathmoveto{\pgfqpoint{1.380303in}{0.447788in}}%
\pgfpathlineto{\pgfqpoint{1.383896in}{0.457637in}}%
\pgfpathlineto{\pgfqpoint{1.384080in}{0.463783in}}%
\pgfpathlineto{\pgfqpoint{1.393078in}{0.457489in}}%
\pgfpathlineto{\pgfqpoint{1.394896in}{0.449171in}}%
\pgfpathlineto{\pgfqpoint{1.389591in}{0.443421in}}%
\pgfpathlineto{\pgfqpoint{1.380303in}{0.447788in}}%
\pgfpathclose%
\pgfusepath{fill}%
\end{pgfscope}%
\begin{pgfscope}%
\pgfpathrectangle{\pgfqpoint{0.100000in}{0.100000in}}{\pgfqpoint{3.608454in}{2.310000in}}%
\pgfusepath{clip}%
\pgfsetbuttcap%
\pgfsetmiterjoin%
\definecolor{currentfill}{rgb}{0.000000,0.490196,0.754902}%
\pgfsetfillcolor{currentfill}%
\pgfsetlinewidth{0.000000pt}%
\definecolor{currentstroke}{rgb}{0.000000,0.000000,0.000000}%
\pgfsetstrokecolor{currentstroke}%
\pgfsetstrokeopacity{0.000000}%
\pgfsetdash{}{0pt}%
\pgfpathmoveto{\pgfqpoint{1.385807in}{0.418653in}}%
\pgfpathlineto{\pgfqpoint{1.386519in}{0.422119in}}%
\pgfpathlineto{\pgfqpoint{1.397550in}{0.420104in}}%
\pgfpathlineto{\pgfqpoint{1.394836in}{0.415075in}}%
\pgfpathlineto{\pgfqpoint{1.385807in}{0.418653in}}%
\pgfpathclose%
\pgfusepath{fill}%
\end{pgfscope}%
\begin{pgfscope}%
\pgfpathrectangle{\pgfqpoint{0.100000in}{0.100000in}}{\pgfqpoint{3.608454in}{2.310000in}}%
\pgfusepath{clip}%
\pgfsetbuttcap%
\pgfsetmiterjoin%
\definecolor{currentfill}{rgb}{0.000000,0.490196,0.754902}%
\pgfsetfillcolor{currentfill}%
\pgfsetlinewidth{0.000000pt}%
\definecolor{currentstroke}{rgb}{0.000000,0.000000,0.000000}%
\pgfsetstrokecolor{currentstroke}%
\pgfsetstrokeopacity{0.000000}%
\pgfsetdash{}{0pt}%
\pgfpathmoveto{\pgfqpoint{1.407961in}{0.408746in}}%
\pgfpathlineto{\pgfqpoint{1.405253in}{0.414109in}}%
\pgfpathlineto{\pgfqpoint{1.413113in}{0.426464in}}%
\pgfpathlineto{\pgfqpoint{1.404816in}{0.436091in}}%
\pgfpathlineto{\pgfqpoint{1.404807in}{0.443981in}}%
\pgfpathlineto{\pgfqpoint{1.412629in}{0.450646in}}%
\pgfpathlineto{\pgfqpoint{1.418487in}{0.448805in}}%
\pgfpathlineto{\pgfqpoint{1.419893in}{0.442848in}}%
\pgfpathlineto{\pgfqpoint{1.418178in}{0.434960in}}%
\pgfpathlineto{\pgfqpoint{1.432543in}{0.428456in}}%
\pgfpathlineto{\pgfqpoint{1.436516in}{0.412431in}}%
\pgfpathlineto{\pgfqpoint{1.441689in}{0.404958in}}%
\pgfpathlineto{\pgfqpoint{1.438189in}{0.398895in}}%
\pgfpathlineto{\pgfqpoint{1.431764in}{0.398037in}}%
\pgfpathlineto{\pgfqpoint{1.421058in}{0.403867in}}%
\pgfpathlineto{\pgfqpoint{1.407961in}{0.408746in}}%
\pgfpathclose%
\pgfusepath{fill}%
\end{pgfscope}%
\begin{pgfscope}%
\pgfpathrectangle{\pgfqpoint{0.100000in}{0.100000in}}{\pgfqpoint{3.608454in}{2.310000in}}%
\pgfusepath{clip}%
\pgfsetbuttcap%
\pgfsetmiterjoin%
\definecolor{currentfill}{rgb}{0.000000,0.490196,0.754902}%
\pgfsetfillcolor{currentfill}%
\pgfsetlinewidth{0.000000pt}%
\definecolor{currentstroke}{rgb}{0.000000,0.000000,0.000000}%
\pgfsetstrokecolor{currentstroke}%
\pgfsetstrokeopacity{0.000000}%
\pgfsetdash{}{0pt}%
\pgfpathmoveto{\pgfqpoint{1.406047in}{0.473337in}}%
\pgfpathlineto{\pgfqpoint{1.416814in}{0.465266in}}%
\pgfpathlineto{\pgfqpoint{1.412882in}{0.461231in}}%
\pgfpathlineto{\pgfqpoint{1.404816in}{0.462303in}}%
\pgfpathlineto{\pgfqpoint{1.398464in}{0.468448in}}%
\pgfpathlineto{\pgfqpoint{1.393702in}{0.476372in}}%
\pgfpathlineto{\pgfqpoint{1.378874in}{0.487793in}}%
\pgfpathlineto{\pgfqpoint{1.389247in}{0.492428in}}%
\pgfpathlineto{\pgfqpoint{1.406047in}{0.473337in}}%
\pgfpathclose%
\pgfusepath{fill}%
\end{pgfscope}%
\begin{pgfscope}%
\pgfpathrectangle{\pgfqpoint{0.100000in}{0.100000in}}{\pgfqpoint{3.608454in}{2.310000in}}%
\pgfusepath{clip}%
\pgfsetbuttcap%
\pgfsetmiterjoin%
\definecolor{currentfill}{rgb}{0.000000,0.694118,0.652941}%
\pgfsetfillcolor{currentfill}%
\pgfsetlinewidth{0.000000pt}%
\definecolor{currentstroke}{rgb}{0.000000,0.000000,0.000000}%
\pgfsetstrokecolor{currentstroke}%
\pgfsetstrokeopacity{0.000000}%
\pgfsetdash{}{0pt}%
\pgfpathmoveto{\pgfqpoint{1.576920in}{0.953485in}}%
\pgfpathlineto{\pgfqpoint{1.574789in}{0.933425in}}%
\pgfpathlineto{\pgfqpoint{1.544636in}{0.936153in}}%
\pgfpathlineto{\pgfqpoint{1.545317in}{0.943021in}}%
\pgfpathlineto{\pgfqpoint{1.531592in}{0.944243in}}%
\pgfpathlineto{\pgfqpoint{1.532929in}{0.957126in}}%
\pgfpathlineto{\pgfqpoint{1.524717in}{0.957918in}}%
\pgfpathlineto{\pgfqpoint{1.526759in}{0.978642in}}%
\pgfpathlineto{\pgfqpoint{1.519945in}{0.979398in}}%
\pgfpathlineto{\pgfqpoint{1.523829in}{1.020634in}}%
\pgfpathlineto{\pgfqpoint{1.539569in}{1.019066in}}%
\pgfpathlineto{\pgfqpoint{1.540290in}{1.025994in}}%
\pgfpathlineto{\pgfqpoint{1.553889in}{1.024630in}}%
\pgfpathlineto{\pgfqpoint{1.554565in}{1.031506in}}%
\pgfpathlineto{\pgfqpoint{1.561377in}{1.030839in}}%
\pgfpathlineto{\pgfqpoint{1.562067in}{1.037697in}}%
\pgfpathlineto{\pgfqpoint{1.568889in}{1.037039in}}%
\pgfpathlineto{\pgfqpoint{1.569604in}{1.044093in}}%
\pgfpathlineto{\pgfqpoint{1.585141in}{1.042675in}}%
\pgfpathlineto{\pgfqpoint{1.583688in}{1.026354in}}%
\pgfpathlineto{\pgfqpoint{1.617270in}{1.023437in}}%
\pgfpathlineto{\pgfqpoint{1.614395in}{0.989155in}}%
\pgfpathlineto{\pgfqpoint{1.580563in}{0.992038in}}%
\pgfpathlineto{\pgfqpoint{1.576920in}{0.953485in}}%
\pgfpathclose%
\pgfusepath{fill}%
\end{pgfscope}%
\begin{pgfscope}%
\pgfpathrectangle{\pgfqpoint{0.100000in}{0.100000in}}{\pgfqpoint{3.608454in}{2.310000in}}%
\pgfusepath{clip}%
\pgfsetbuttcap%
\pgfsetmiterjoin%
\definecolor{currentfill}{rgb}{0.000000,0.611765,0.694118}%
\pgfsetfillcolor{currentfill}%
\pgfsetlinewidth{0.000000pt}%
\definecolor{currentstroke}{rgb}{0.000000,0.000000,0.000000}%
\pgfsetstrokecolor{currentstroke}%
\pgfsetstrokeopacity{0.000000}%
\pgfsetdash{}{0pt}%
\pgfpathmoveto{\pgfqpoint{3.088906in}{0.984668in}}%
\pgfpathlineto{\pgfqpoint{3.081713in}{0.979204in}}%
\pgfpathlineto{\pgfqpoint{3.073391in}{0.976759in}}%
\pgfpathlineto{\pgfqpoint{3.065412in}{0.981312in}}%
\pgfpathlineto{\pgfqpoint{3.059264in}{0.991086in}}%
\pgfpathlineto{\pgfqpoint{3.051594in}{0.997455in}}%
\pgfpathlineto{\pgfqpoint{3.052596in}{1.003504in}}%
\pgfpathlineto{\pgfqpoint{3.049719in}{1.006843in}}%
\pgfpathlineto{\pgfqpoint{3.049502in}{1.021849in}}%
\pgfpathlineto{\pgfqpoint{3.046504in}{1.026091in}}%
\pgfpathlineto{\pgfqpoint{3.041428in}{1.022979in}}%
\pgfpathlineto{\pgfqpoint{3.030270in}{1.030329in}}%
\pgfpathlineto{\pgfqpoint{3.034331in}{1.046602in}}%
\pgfpathlineto{\pgfqpoint{3.026275in}{1.051810in}}%
\pgfpathlineto{\pgfqpoint{3.034032in}{1.059624in}}%
\pgfpathlineto{\pgfqpoint{3.040979in}{1.056871in}}%
\pgfpathlineto{\pgfqpoint{3.047179in}{1.057968in}}%
\pgfpathlineto{\pgfqpoint{3.046213in}{1.063185in}}%
\pgfpathlineto{\pgfqpoint{3.054685in}{1.068940in}}%
\pgfpathlineto{\pgfqpoint{3.060908in}{1.061022in}}%
\pgfpathlineto{\pgfqpoint{3.065576in}{1.050744in}}%
\pgfpathlineto{\pgfqpoint{3.073425in}{1.048823in}}%
\pgfpathlineto{\pgfqpoint{3.078779in}{1.041854in}}%
\pgfpathlineto{\pgfqpoint{3.076612in}{1.036110in}}%
\pgfpathlineto{\pgfqpoint{3.096046in}{1.025469in}}%
\pgfpathlineto{\pgfqpoint{3.091334in}{1.021253in}}%
\pgfpathlineto{\pgfqpoint{3.096860in}{1.011889in}}%
\pgfpathlineto{\pgfqpoint{3.092298in}{1.007042in}}%
\pgfpathlineto{\pgfqpoint{3.094287in}{1.003646in}}%
\pgfpathlineto{\pgfqpoint{3.088906in}{0.984668in}}%
\pgfpathclose%
\pgfusepath{fill}%
\end{pgfscope}%
\begin{pgfscope}%
\pgfpathrectangle{\pgfqpoint{0.100000in}{0.100000in}}{\pgfqpoint{3.608454in}{2.310000in}}%
\pgfusepath{clip}%
\pgfsetbuttcap%
\pgfsetmiterjoin%
\definecolor{currentfill}{rgb}{0.000000,0.611765,0.694118}%
\pgfsetfillcolor{currentfill}%
\pgfsetlinewidth{0.000000pt}%
\definecolor{currentstroke}{rgb}{0.000000,0.000000,0.000000}%
\pgfsetstrokecolor{currentstroke}%
\pgfsetstrokeopacity{0.000000}%
\pgfsetdash{}{0pt}%
\pgfpathmoveto{\pgfqpoint{1.562162in}{1.934829in}}%
\pgfpathlineto{\pgfqpoint{1.559222in}{1.935142in}}%
\pgfpathlineto{\pgfqpoint{1.556192in}{1.907466in}}%
\pgfpathlineto{\pgfqpoint{1.553831in}{1.907730in}}%
\pgfpathlineto{\pgfqpoint{1.487149in}{1.915660in}}%
\pgfpathlineto{\pgfqpoint{1.490580in}{1.943193in}}%
\pgfpathlineto{\pgfqpoint{1.492673in}{1.942949in}}%
\pgfpathlineto{\pgfqpoint{1.495943in}{1.970351in}}%
\pgfpathlineto{\pgfqpoint{1.497843in}{1.970127in}}%
\pgfpathlineto{\pgfqpoint{1.501201in}{1.997658in}}%
\pgfpathlineto{\pgfqpoint{1.503208in}{1.998589in}}%
\pgfpathlineto{\pgfqpoint{1.506437in}{2.025137in}}%
\pgfpathlineto{\pgfqpoint{1.513276in}{2.024315in}}%
\pgfpathlineto{\pgfqpoint{1.512427in}{2.017326in}}%
\pgfpathlineto{\pgfqpoint{1.519298in}{2.016515in}}%
\pgfpathlineto{\pgfqpoint{1.520131in}{2.023482in}}%
\pgfpathlineto{\pgfqpoint{1.543137in}{2.020826in}}%
\pgfpathlineto{\pgfqpoint{1.547528in}{2.020313in}}%
\pgfpathlineto{\pgfqpoint{1.545674in}{2.004107in}}%
\pgfpathlineto{\pgfqpoint{1.552609in}{2.003322in}}%
\pgfpathlineto{\pgfqpoint{1.552079in}{1.998723in}}%
\pgfpathlineto{\pgfqpoint{1.558985in}{1.997946in}}%
\pgfpathlineto{\pgfqpoint{1.560751in}{1.993122in}}%
\pgfpathlineto{\pgfqpoint{1.583733in}{1.990580in}}%
\pgfpathlineto{\pgfqpoint{1.582076in}{1.976639in}}%
\pgfpathlineto{\pgfqpoint{1.574985in}{1.975097in}}%
\pgfpathlineto{\pgfqpoint{1.573489in}{1.961301in}}%
\pgfpathlineto{\pgfqpoint{1.564057in}{1.962319in}}%
\pgfpathlineto{\pgfqpoint{1.562162in}{1.934829in}}%
\pgfpathclose%
\pgfusepath{fill}%
\end{pgfscope}%
\begin{pgfscope}%
\pgfpathrectangle{\pgfqpoint{0.100000in}{0.100000in}}{\pgfqpoint{3.608454in}{2.310000in}}%
\pgfusepath{clip}%
\pgfsetbuttcap%
\pgfsetmiterjoin%
\definecolor{currentfill}{rgb}{0.000000,0.431373,0.784314}%
\pgfsetfillcolor{currentfill}%
\pgfsetlinewidth{0.000000pt}%
\definecolor{currentstroke}{rgb}{0.000000,0.000000,0.000000}%
\pgfsetstrokecolor{currentstroke}%
\pgfsetstrokeopacity{0.000000}%
\pgfsetdash{}{0pt}%
\pgfpathmoveto{\pgfqpoint{2.013559in}{0.517174in}}%
\pgfpathlineto{\pgfqpoint{2.014894in}{0.542939in}}%
\pgfpathlineto{\pgfqpoint{1.992065in}{0.566132in}}%
\pgfpathlineto{\pgfqpoint{2.012654in}{0.578657in}}%
\pgfpathlineto{\pgfqpoint{2.015636in}{0.587127in}}%
\pgfpathlineto{\pgfqpoint{2.024967in}{0.596552in}}%
\pgfpathlineto{\pgfqpoint{2.032870in}{0.593014in}}%
\pgfpathlineto{\pgfqpoint{2.032880in}{0.583631in}}%
\pgfpathlineto{\pgfqpoint{2.037581in}{0.580386in}}%
\pgfpathlineto{\pgfqpoint{2.040450in}{0.572451in}}%
\pgfpathlineto{\pgfqpoint{2.047518in}{0.566563in}}%
\pgfpathlineto{\pgfqpoint{2.054185in}{0.571463in}}%
\pgfpathlineto{\pgfqpoint{2.061080in}{0.568664in}}%
\pgfpathlineto{\pgfqpoint{2.066381in}{0.572643in}}%
\pgfpathlineto{\pgfqpoint{2.068450in}{0.580491in}}%
\pgfpathlineto{\pgfqpoint{2.074520in}{0.580523in}}%
\pgfpathlineto{\pgfqpoint{2.077251in}{0.591805in}}%
\pgfpathlineto{\pgfqpoint{2.088573in}{0.592795in}}%
\pgfpathlineto{\pgfqpoint{2.091703in}{0.589805in}}%
\pgfpathlineto{\pgfqpoint{2.090532in}{0.582574in}}%
\pgfpathlineto{\pgfqpoint{2.102863in}{0.561208in}}%
\pgfpathlineto{\pgfqpoint{2.098432in}{0.551471in}}%
\pgfpathlineto{\pgfqpoint{2.079822in}{0.534485in}}%
\pgfpathlineto{\pgfqpoint{2.074339in}{0.533228in}}%
\pgfpathlineto{\pgfqpoint{2.049544in}{0.518585in}}%
\pgfpathlineto{\pgfqpoint{2.036951in}{0.514000in}}%
\pgfpathlineto{\pgfqpoint{2.033736in}{0.517714in}}%
\pgfpathlineto{\pgfqpoint{2.020863in}{0.511930in}}%
\pgfpathlineto{\pgfqpoint{2.021021in}{0.521517in}}%
\pgfpathlineto{\pgfqpoint{2.013559in}{0.517174in}}%
\pgfpathclose%
\pgfusepath{fill}%
\end{pgfscope}%
\begin{pgfscope}%
\pgfpathrectangle{\pgfqpoint{0.100000in}{0.100000in}}{\pgfqpoint{3.608454in}{2.310000in}}%
\pgfusepath{clip}%
\pgfsetbuttcap%
\pgfsetmiterjoin%
\definecolor{currentfill}{rgb}{0.000000,0.686275,0.656863}%
\pgfsetfillcolor{currentfill}%
\pgfsetlinewidth{0.000000pt}%
\definecolor{currentstroke}{rgb}{0.000000,0.000000,0.000000}%
\pgfsetstrokecolor{currentstroke}%
\pgfsetstrokeopacity{0.000000}%
\pgfsetdash{}{0pt}%
\pgfpathmoveto{\pgfqpoint{3.111250in}{1.435310in}}%
\pgfpathlineto{\pgfqpoint{3.103999in}{1.422556in}}%
\pgfpathlineto{\pgfqpoint{3.099019in}{1.419473in}}%
\pgfpathlineto{\pgfqpoint{3.095711in}{1.408904in}}%
\pgfpathlineto{\pgfqpoint{3.086938in}{1.414331in}}%
\pgfpathlineto{\pgfqpoint{3.084328in}{1.406748in}}%
\pgfpathlineto{\pgfqpoint{3.063909in}{1.418899in}}%
\pgfpathlineto{\pgfqpoint{3.065468in}{1.428454in}}%
\pgfpathlineto{\pgfqpoint{3.062165in}{1.429526in}}%
\pgfpathlineto{\pgfqpoint{3.063986in}{1.437446in}}%
\pgfpathlineto{\pgfqpoint{3.060760in}{1.438938in}}%
\pgfpathlineto{\pgfqpoint{3.052407in}{1.435951in}}%
\pgfpathlineto{\pgfqpoint{3.045971in}{1.477046in}}%
\pgfpathlineto{\pgfqpoint{3.051009in}{1.477958in}}%
\pgfpathlineto{\pgfqpoint{3.111561in}{1.488897in}}%
\pgfpathlineto{\pgfqpoint{3.116008in}{1.487243in}}%
\pgfpathlineto{\pgfqpoint{3.110506in}{1.479451in}}%
\pgfpathlineto{\pgfqpoint{3.109226in}{1.471916in}}%
\pgfpathlineto{\pgfqpoint{3.097628in}{1.471273in}}%
\pgfpathlineto{\pgfqpoint{3.091266in}{1.461236in}}%
\pgfpathlineto{\pgfqpoint{3.088891in}{1.453283in}}%
\pgfpathlineto{\pgfqpoint{3.082421in}{1.444639in}}%
\pgfpathlineto{\pgfqpoint{3.087837in}{1.442494in}}%
\pgfpathlineto{\pgfqpoint{3.105576in}{1.439244in}}%
\pgfpathlineto{\pgfqpoint{3.111250in}{1.435310in}}%
\pgfpathclose%
\pgfusepath{fill}%
\end{pgfscope}%
\begin{pgfscope}%
\pgfpathrectangle{\pgfqpoint{0.100000in}{0.100000in}}{\pgfqpoint{3.608454in}{2.310000in}}%
\pgfusepath{clip}%
\pgfsetbuttcap%
\pgfsetmiterjoin%
\definecolor{currentfill}{rgb}{0.000000,0.611765,0.694118}%
\pgfsetfillcolor{currentfill}%
\pgfsetlinewidth{0.000000pt}%
\definecolor{currentstroke}{rgb}{0.000000,0.000000,0.000000}%
\pgfsetstrokecolor{currentstroke}%
\pgfsetstrokeopacity{0.000000}%
\pgfsetdash{}{0pt}%
\pgfpathmoveto{\pgfqpoint{3.279173in}{1.774369in}}%
\pgfpathlineto{\pgfqpoint{3.272069in}{1.775491in}}%
\pgfpathlineto{\pgfqpoint{3.271017in}{1.786225in}}%
\pgfpathlineto{\pgfqpoint{3.267628in}{1.792368in}}%
\pgfpathlineto{\pgfqpoint{3.270444in}{1.800501in}}%
\pgfpathlineto{\pgfqpoint{3.267521in}{1.808820in}}%
\pgfpathlineto{\pgfqpoint{3.257858in}{1.810998in}}%
\pgfpathlineto{\pgfqpoint{3.260468in}{1.823521in}}%
\pgfpathlineto{\pgfqpoint{3.245888in}{1.867665in}}%
\pgfpathlineto{\pgfqpoint{3.263105in}{1.873951in}}%
\pgfpathlineto{\pgfqpoint{3.276788in}{1.878840in}}%
\pgfpathlineto{\pgfqpoint{3.280930in}{1.867522in}}%
\pgfpathlineto{\pgfqpoint{3.277231in}{1.863034in}}%
\pgfpathlineto{\pgfqpoint{3.286142in}{1.855787in}}%
\pgfpathlineto{\pgfqpoint{3.295690in}{1.856839in}}%
\pgfpathlineto{\pgfqpoint{3.296075in}{1.852707in}}%
\pgfpathlineto{\pgfqpoint{3.287566in}{1.849431in}}%
\pgfpathlineto{\pgfqpoint{3.309890in}{1.786988in}}%
\pgfpathlineto{\pgfqpoint{3.300308in}{1.775508in}}%
\pgfpathlineto{\pgfqpoint{3.291128in}{1.772508in}}%
\pgfpathlineto{\pgfqpoint{3.279173in}{1.774369in}}%
\pgfpathclose%
\pgfusepath{fill}%
\end{pgfscope}%
\begin{pgfscope}%
\pgfpathrectangle{\pgfqpoint{0.100000in}{0.100000in}}{\pgfqpoint{3.608454in}{2.310000in}}%
\pgfusepath{clip}%
\pgfsetbuttcap%
\pgfsetmiterjoin%
\definecolor{currentfill}{rgb}{0.000000,0.517647,0.741176}%
\pgfsetfillcolor{currentfill}%
\pgfsetlinewidth{0.000000pt}%
\definecolor{currentstroke}{rgb}{0.000000,0.000000,0.000000}%
\pgfsetstrokecolor{currentstroke}%
\pgfsetstrokeopacity{0.000000}%
\pgfsetdash{}{0pt}%
\pgfpathmoveto{\pgfqpoint{2.558058in}{1.370400in}}%
\pgfpathlineto{\pgfqpoint{2.550764in}{1.371543in}}%
\pgfpathlineto{\pgfqpoint{2.506514in}{1.367545in}}%
\pgfpathlineto{\pgfqpoint{2.506265in}{1.370984in}}%
\pgfpathlineto{\pgfqpoint{2.503724in}{1.405402in}}%
\pgfpathlineto{\pgfqpoint{2.528514in}{1.407327in}}%
\pgfpathlineto{\pgfqpoint{2.534185in}{1.410410in}}%
\pgfpathlineto{\pgfqpoint{2.534753in}{1.425828in}}%
\pgfpathlineto{\pgfqpoint{2.559183in}{1.428045in}}%
\pgfpathlineto{\pgfqpoint{2.562752in}{1.385889in}}%
\pgfpathlineto{\pgfqpoint{2.558464in}{1.382618in}}%
\pgfpathlineto{\pgfqpoint{2.560407in}{1.373999in}}%
\pgfpathlineto{\pgfqpoint{2.558058in}{1.370400in}}%
\pgfpathclose%
\pgfusepath{fill}%
\end{pgfscope}%
\begin{pgfscope}%
\pgfpathrectangle{\pgfqpoint{0.100000in}{0.100000in}}{\pgfqpoint{3.608454in}{2.310000in}}%
\pgfusepath{clip}%
\pgfsetbuttcap%
\pgfsetmiterjoin%
\definecolor{currentfill}{rgb}{0.000000,0.419608,0.790196}%
\pgfsetfillcolor{currentfill}%
\pgfsetlinewidth{0.000000pt}%
\definecolor{currentstroke}{rgb}{0.000000,0.000000,0.000000}%
\pgfsetstrokecolor{currentstroke}%
\pgfsetstrokeopacity{0.000000}%
\pgfsetdash{}{0pt}%
\pgfpathmoveto{\pgfqpoint{1.843275in}{1.841097in}}%
\pgfpathlineto{\pgfqpoint{1.812489in}{1.842876in}}%
\pgfpathlineto{\pgfqpoint{1.809666in}{1.848703in}}%
\pgfpathlineto{\pgfqpoint{1.812344in}{1.853586in}}%
\pgfpathlineto{\pgfqpoint{1.808999in}{1.861084in}}%
\pgfpathlineto{\pgfqpoint{1.801983in}{1.865242in}}%
\pgfpathlineto{\pgfqpoint{1.807298in}{1.876074in}}%
\pgfpathlineto{\pgfqpoint{1.811895in}{1.879653in}}%
\pgfpathlineto{\pgfqpoint{1.808345in}{1.892715in}}%
\pgfpathlineto{\pgfqpoint{1.802493in}{1.899008in}}%
\pgfpathlineto{\pgfqpoint{1.845659in}{1.896366in}}%
\pgfpathlineto{\pgfqpoint{1.843275in}{1.841097in}}%
\pgfpathclose%
\pgfusepath{fill}%
\end{pgfscope}%
\begin{pgfscope}%
\pgfpathrectangle{\pgfqpoint{0.100000in}{0.100000in}}{\pgfqpoint{3.608454in}{2.310000in}}%
\pgfusepath{clip}%
\pgfsetbuttcap%
\pgfsetmiterjoin%
\definecolor{currentfill}{rgb}{0.000000,0.827451,0.586275}%
\pgfsetfillcolor{currentfill}%
\pgfsetlinewidth{0.000000pt}%
\definecolor{currentstroke}{rgb}{0.000000,0.000000,0.000000}%
\pgfsetstrokecolor{currentstroke}%
\pgfsetstrokeopacity{0.000000}%
\pgfsetdash{}{0pt}%
\pgfpathmoveto{\pgfqpoint{0.787216in}{0.442053in}}%
\pgfpathlineto{\pgfqpoint{0.796011in}{0.435184in}}%
\pgfpathlineto{\pgfqpoint{0.790893in}{0.428785in}}%
\pgfpathlineto{\pgfqpoint{0.779568in}{0.437653in}}%
\pgfpathlineto{\pgfqpoint{0.783683in}{0.438305in}}%
\pgfpathlineto{\pgfqpoint{0.787048in}{0.440340in}}%
\pgfpathlineto{\pgfqpoint{0.787216in}{0.442053in}}%
\pgfpathclose%
\pgfusepath{fill}%
\end{pgfscope}%
\begin{pgfscope}%
\pgfpathrectangle{\pgfqpoint{0.100000in}{0.100000in}}{\pgfqpoint{3.608454in}{2.310000in}}%
\pgfusepath{clip}%
\pgfsetbuttcap%
\pgfsetmiterjoin%
\definecolor{currentfill}{rgb}{0.000000,0.796078,0.601961}%
\pgfsetfillcolor{currentfill}%
\pgfsetlinewidth{0.000000pt}%
\definecolor{currentstroke}{rgb}{0.000000,0.000000,0.000000}%
\pgfsetstrokecolor{currentstroke}%
\pgfsetstrokeopacity{0.000000}%
\pgfsetdash{}{0pt}%
\pgfpathmoveto{\pgfqpoint{1.760148in}{1.734344in}}%
\pgfpathlineto{\pgfqpoint{1.750543in}{1.731001in}}%
\pgfpathlineto{\pgfqpoint{1.748309in}{1.699187in}}%
\pgfpathlineto{\pgfqpoint{1.697612in}{1.702965in}}%
\pgfpathlineto{\pgfqpoint{1.696530in}{1.709950in}}%
\pgfpathlineto{\pgfqpoint{1.697918in}{1.727667in}}%
\pgfpathlineto{\pgfqpoint{1.704828in}{1.727690in}}%
\pgfpathlineto{\pgfqpoint{1.707249in}{1.750308in}}%
\pgfpathlineto{\pgfqpoint{1.710799in}{1.791190in}}%
\pgfpathlineto{\pgfqpoint{1.723276in}{1.792199in}}%
\pgfpathlineto{\pgfqpoint{1.729800in}{1.795494in}}%
\pgfpathlineto{\pgfqpoint{1.737625in}{1.793677in}}%
\pgfpathlineto{\pgfqpoint{1.747956in}{1.802589in}}%
\pgfpathlineto{\pgfqpoint{1.754993in}{1.802225in}}%
\pgfpathlineto{\pgfqpoint{1.760814in}{1.806325in}}%
\pgfpathlineto{\pgfqpoint{1.759531in}{1.801865in}}%
\pgfpathlineto{\pgfqpoint{1.757629in}{1.774393in}}%
\pgfpathlineto{\pgfqpoint{1.756032in}{1.760574in}}%
\pgfpathlineto{\pgfqpoint{1.762841in}{1.760148in}}%
\pgfpathlineto{\pgfqpoint{1.760148in}{1.734344in}}%
\pgfpathclose%
\pgfusepath{fill}%
\end{pgfscope}%
\begin{pgfscope}%
\pgfpathrectangle{\pgfqpoint{0.100000in}{0.100000in}}{\pgfqpoint{3.608454in}{2.310000in}}%
\pgfusepath{clip}%
\pgfsetbuttcap%
\pgfsetmiterjoin%
\definecolor{currentfill}{rgb}{0.000000,0.541176,0.729412}%
\pgfsetfillcolor{currentfill}%
\pgfsetlinewidth{0.000000pt}%
\definecolor{currentstroke}{rgb}{0.000000,0.000000,0.000000}%
\pgfsetstrokecolor{currentstroke}%
\pgfsetstrokeopacity{0.000000}%
\pgfsetdash{}{0pt}%
\pgfpathmoveto{\pgfqpoint{2.406029in}{1.504012in}}%
\pgfpathlineto{\pgfqpoint{2.405330in}{1.517808in}}%
\pgfpathlineto{\pgfqpoint{2.378376in}{1.516534in}}%
\pgfpathlineto{\pgfqpoint{2.378619in}{1.509638in}}%
\pgfpathlineto{\pgfqpoint{2.358001in}{1.509016in}}%
\pgfpathlineto{\pgfqpoint{2.348213in}{1.508885in}}%
\pgfpathlineto{\pgfqpoint{2.345525in}{1.515337in}}%
\pgfpathlineto{\pgfqpoint{2.337899in}{1.521676in}}%
\pgfpathlineto{\pgfqpoint{2.341290in}{1.535608in}}%
\pgfpathlineto{\pgfqpoint{2.356534in}{1.539366in}}%
\pgfpathlineto{\pgfqpoint{2.356154in}{1.550723in}}%
\pgfpathlineto{\pgfqpoint{2.349290in}{1.550458in}}%
\pgfpathlineto{\pgfqpoint{2.348758in}{1.564236in}}%
\pgfpathlineto{\pgfqpoint{2.376029in}{1.564684in}}%
\pgfpathlineto{\pgfqpoint{2.380372in}{1.562314in}}%
\pgfpathlineto{\pgfqpoint{2.387375in}{1.567036in}}%
\pgfpathlineto{\pgfqpoint{2.391551in}{1.551544in}}%
\pgfpathlineto{\pgfqpoint{2.410650in}{1.552521in}}%
\pgfpathlineto{\pgfqpoint{2.412483in}{1.524944in}}%
\pgfpathlineto{\pgfqpoint{2.425489in}{1.525655in}}%
\pgfpathlineto{\pgfqpoint{2.426710in}{1.505121in}}%
\pgfpathlineto{\pgfqpoint{2.406029in}{1.504012in}}%
\pgfpathclose%
\pgfusepath{fill}%
\end{pgfscope}%
\begin{pgfscope}%
\pgfpathrectangle{\pgfqpoint{0.100000in}{0.100000in}}{\pgfqpoint{3.608454in}{2.310000in}}%
\pgfusepath{clip}%
\pgfsetbuttcap%
\pgfsetmiterjoin%
\definecolor{currentfill}{rgb}{0.000000,0.494118,0.752941}%
\pgfsetfillcolor{currentfill}%
\pgfsetlinewidth{0.000000pt}%
\definecolor{currentstroke}{rgb}{0.000000,0.000000,0.000000}%
\pgfsetstrokecolor{currentstroke}%
\pgfsetstrokeopacity{0.000000}%
\pgfsetdash{}{0pt}%
\pgfpathmoveto{\pgfqpoint{2.708723in}{0.733977in}}%
\pgfpathlineto{\pgfqpoint{2.719091in}{0.734938in}}%
\pgfpathlineto{\pgfqpoint{2.721186in}{0.711980in}}%
\pgfpathlineto{\pgfqpoint{2.733506in}{0.712730in}}%
\pgfpathlineto{\pgfqpoint{2.730084in}{0.707397in}}%
\pgfpathlineto{\pgfqpoint{2.735056in}{0.696150in}}%
\pgfpathlineto{\pgfqpoint{2.726751in}{0.687235in}}%
\pgfpathlineto{\pgfqpoint{2.727447in}{0.677658in}}%
\pgfpathlineto{\pgfqpoint{2.710794in}{0.682159in}}%
\pgfpathlineto{\pgfqpoint{2.699063in}{0.683673in}}%
\pgfpathlineto{\pgfqpoint{2.682757in}{0.683563in}}%
\pgfpathlineto{\pgfqpoint{2.666620in}{0.680661in}}%
\pgfpathlineto{\pgfqpoint{2.645613in}{0.674473in}}%
\pgfpathlineto{\pgfqpoint{2.635436in}{0.672947in}}%
\pgfpathlineto{\pgfqpoint{2.622714in}{0.669017in}}%
\pgfpathlineto{\pgfqpoint{2.627322in}{0.678308in}}%
\pgfpathlineto{\pgfqpoint{2.632066in}{0.682154in}}%
\pgfpathlineto{\pgfqpoint{2.626017in}{0.689071in}}%
\pgfpathlineto{\pgfqpoint{2.628964in}{0.696389in}}%
\pgfpathlineto{\pgfqpoint{2.627748in}{0.701087in}}%
\pgfpathlineto{\pgfqpoint{2.618895in}{0.705619in}}%
\pgfpathlineto{\pgfqpoint{2.610908in}{0.714979in}}%
\pgfpathlineto{\pgfqpoint{2.612513in}{0.725646in}}%
\pgfpathlineto{\pgfqpoint{2.674593in}{0.730790in}}%
\pgfpathlineto{\pgfqpoint{2.708723in}{0.733977in}}%
\pgfpathclose%
\pgfusepath{fill}%
\end{pgfscope}%
\begin{pgfscope}%
\pgfpathrectangle{\pgfqpoint{0.100000in}{0.100000in}}{\pgfqpoint{3.608454in}{2.310000in}}%
\pgfusepath{clip}%
\pgfsetbuttcap%
\pgfsetmiterjoin%
\definecolor{currentfill}{rgb}{0.000000,0.364706,0.817647}%
\pgfsetfillcolor{currentfill}%
\pgfsetlinewidth{0.000000pt}%
\definecolor{currentstroke}{rgb}{0.000000,0.000000,0.000000}%
\pgfsetstrokecolor{currentstroke}%
\pgfsetstrokeopacity{0.000000}%
\pgfsetdash{}{0pt}%
\pgfpathmoveto{\pgfqpoint{1.826801in}{1.425965in}}%
\pgfpathlineto{\pgfqpoint{1.824934in}{1.391676in}}%
\pgfpathlineto{\pgfqpoint{1.792462in}{1.393414in}}%
\pgfpathlineto{\pgfqpoint{1.791324in}{1.393448in}}%
\pgfpathlineto{\pgfqpoint{1.793434in}{1.427744in}}%
\pgfpathlineto{\pgfqpoint{1.826801in}{1.425965in}}%
\pgfpathclose%
\pgfusepath{fill}%
\end{pgfscope}%
\begin{pgfscope}%
\pgfpathrectangle{\pgfqpoint{0.100000in}{0.100000in}}{\pgfqpoint{3.608454in}{2.310000in}}%
\pgfusepath{clip}%
\pgfsetbuttcap%
\pgfsetmiterjoin%
\definecolor{currentfill}{rgb}{0.000000,0.815686,0.592157}%
\pgfsetfillcolor{currentfill}%
\pgfsetlinewidth{0.000000pt}%
\definecolor{currentstroke}{rgb}{0.000000,0.000000,0.000000}%
\pgfsetstrokecolor{currentstroke}%
\pgfsetstrokeopacity{0.000000}%
\pgfsetdash{}{0pt}%
\pgfpathmoveto{\pgfqpoint{2.350507in}{1.098801in}}%
\pgfpathlineto{\pgfqpoint{2.343663in}{1.088662in}}%
\pgfpathlineto{\pgfqpoint{2.340391in}{1.078576in}}%
\pgfpathlineto{\pgfqpoint{2.340700in}{1.069582in}}%
\pgfpathlineto{\pgfqpoint{2.313347in}{1.069167in}}%
\pgfpathlineto{\pgfqpoint{2.313117in}{1.083232in}}%
\pgfpathlineto{\pgfqpoint{2.307635in}{1.090812in}}%
\pgfpathlineto{\pgfqpoint{2.308868in}{1.093507in}}%
\pgfpathlineto{\pgfqpoint{2.302022in}{1.100303in}}%
\pgfpathlineto{\pgfqpoint{2.293303in}{1.102454in}}%
\pgfpathlineto{\pgfqpoint{2.295109in}{1.108191in}}%
\pgfpathlineto{\pgfqpoint{2.286012in}{1.115699in}}%
\pgfpathlineto{\pgfqpoint{2.288166in}{1.125922in}}%
\pgfpathlineto{\pgfqpoint{2.325882in}{1.126330in}}%
\pgfpathlineto{\pgfqpoint{2.325633in}{1.133246in}}%
\pgfpathlineto{\pgfqpoint{2.332736in}{1.133366in}}%
\pgfpathlineto{\pgfqpoint{2.332518in}{1.146200in}}%
\pgfpathlineto{\pgfqpoint{2.335244in}{1.146265in}}%
\pgfpathlineto{\pgfqpoint{2.343960in}{1.136912in}}%
\pgfpathlineto{\pgfqpoint{2.345455in}{1.126590in}}%
\pgfpathlineto{\pgfqpoint{2.339731in}{1.125367in}}%
\pgfpathlineto{\pgfqpoint{2.340339in}{1.098431in}}%
\pgfpathlineto{\pgfqpoint{2.350507in}{1.098801in}}%
\pgfpathclose%
\pgfusepath{fill}%
\end{pgfscope}%
\begin{pgfscope}%
\pgfpathrectangle{\pgfqpoint{0.100000in}{0.100000in}}{\pgfqpoint{3.608454in}{2.310000in}}%
\pgfusepath{clip}%
\pgfsetbuttcap%
\pgfsetmiterjoin%
\definecolor{currentfill}{rgb}{0.000000,0.411765,0.794118}%
\pgfsetfillcolor{currentfill}%
\pgfsetlinewidth{0.000000pt}%
\definecolor{currentstroke}{rgb}{0.000000,0.000000,0.000000}%
\pgfsetstrokecolor{currentstroke}%
\pgfsetstrokeopacity{0.000000}%
\pgfsetdash{}{0pt}%
\pgfpathmoveto{\pgfqpoint{2.313973in}{1.320063in}}%
\pgfpathlineto{\pgfqpoint{2.314915in}{1.287249in}}%
\pgfpathlineto{\pgfqpoint{2.280744in}{1.286387in}}%
\pgfpathlineto{\pgfqpoint{2.274008in}{1.289693in}}%
\pgfpathlineto{\pgfqpoint{2.267441in}{1.289875in}}%
\pgfpathlineto{\pgfqpoint{2.267557in}{1.296520in}}%
\pgfpathlineto{\pgfqpoint{2.253822in}{1.296576in}}%
\pgfpathlineto{\pgfqpoint{2.240098in}{1.309360in}}%
\pgfpathlineto{\pgfqpoint{2.240062in}{1.316251in}}%
\pgfpathlineto{\pgfqpoint{2.226458in}{1.316665in}}%
\pgfpathlineto{\pgfqpoint{2.226654in}{1.335261in}}%
\pgfpathlineto{\pgfqpoint{2.226701in}{1.338754in}}%
\pgfpathlineto{\pgfqpoint{2.233532in}{1.346189in}}%
\pgfpathlineto{\pgfqpoint{2.232957in}{1.349258in}}%
\pgfpathlineto{\pgfqpoint{2.238456in}{1.359055in}}%
\pgfpathlineto{\pgfqpoint{2.241843in}{1.368445in}}%
\pgfpathlineto{\pgfqpoint{2.246959in}{1.366834in}}%
\pgfpathlineto{\pgfqpoint{2.264061in}{1.361534in}}%
\pgfpathlineto{\pgfqpoint{2.271150in}{1.361570in}}%
\pgfpathlineto{\pgfqpoint{2.284008in}{1.361432in}}%
\pgfpathlineto{\pgfqpoint{2.284099in}{1.347536in}}%
\pgfpathlineto{\pgfqpoint{2.313325in}{1.348137in}}%
\pgfpathlineto{\pgfqpoint{2.313973in}{1.320063in}}%
\pgfpathclose%
\pgfusepath{fill}%
\end{pgfscope}%
\begin{pgfscope}%
\pgfpathrectangle{\pgfqpoint{0.100000in}{0.100000in}}{\pgfqpoint{3.608454in}{2.310000in}}%
\pgfusepath{clip}%
\pgfsetbuttcap%
\pgfsetmiterjoin%
\definecolor{currentfill}{rgb}{0.000000,0.521569,0.739216}%
\pgfsetfillcolor{currentfill}%
\pgfsetlinewidth{0.000000pt}%
\definecolor{currentstroke}{rgb}{0.000000,0.000000,0.000000}%
\pgfsetstrokecolor{currentstroke}%
\pgfsetstrokeopacity{0.000000}%
\pgfsetdash{}{0pt}%
\pgfpathmoveto{\pgfqpoint{2.986023in}{0.797823in}}%
\pgfpathlineto{\pgfqpoint{2.983849in}{0.804589in}}%
\pgfpathlineto{\pgfqpoint{2.978722in}{0.804662in}}%
\pgfpathlineto{\pgfqpoint{2.975294in}{0.828495in}}%
\pgfpathlineto{\pgfqpoint{2.979065in}{0.834250in}}%
\pgfpathlineto{\pgfqpoint{2.975754in}{0.837437in}}%
\pgfpathlineto{\pgfqpoint{2.967579in}{0.838638in}}%
\pgfpathlineto{\pgfqpoint{2.967645in}{0.858971in}}%
\pgfpathlineto{\pgfqpoint{2.962520in}{0.870538in}}%
\pgfpathlineto{\pgfqpoint{2.991213in}{0.868134in}}%
\pgfpathlineto{\pgfqpoint{2.994468in}{0.861760in}}%
\pgfpathlineto{\pgfqpoint{2.996708in}{0.858258in}}%
\pgfpathlineto{\pgfqpoint{3.005118in}{0.855085in}}%
\pgfpathlineto{\pgfqpoint{3.011213in}{0.848731in}}%
\pgfpathlineto{\pgfqpoint{3.020884in}{0.848437in}}%
\pgfpathlineto{\pgfqpoint{3.025040in}{0.841844in}}%
\pgfpathlineto{\pgfqpoint{3.033775in}{0.835716in}}%
\pgfpathlineto{\pgfqpoint{3.039659in}{0.833647in}}%
\pgfpathlineto{\pgfqpoint{3.044197in}{0.827018in}}%
\pgfpathlineto{\pgfqpoint{3.040729in}{0.821509in}}%
\pgfpathlineto{\pgfqpoint{3.035156in}{0.827393in}}%
\pgfpathlineto{\pgfqpoint{3.023851in}{0.824906in}}%
\pgfpathlineto{\pgfqpoint{3.019148in}{0.828947in}}%
\pgfpathlineto{\pgfqpoint{3.015577in}{0.818559in}}%
\pgfpathlineto{\pgfqpoint{3.009224in}{0.814410in}}%
\pgfpathlineto{\pgfqpoint{3.013726in}{0.809687in}}%
\pgfpathlineto{\pgfqpoint{3.007825in}{0.797264in}}%
\pgfpathlineto{\pgfqpoint{3.002283in}{0.799408in}}%
\pgfpathlineto{\pgfqpoint{2.994265in}{0.796799in}}%
\pgfpathlineto{\pgfqpoint{2.986023in}{0.797823in}}%
\pgfpathclose%
\pgfusepath{fill}%
\end{pgfscope}%
\begin{pgfscope}%
\pgfpathrectangle{\pgfqpoint{0.100000in}{0.100000in}}{\pgfqpoint{3.608454in}{2.310000in}}%
\pgfusepath{clip}%
\pgfsetbuttcap%
\pgfsetmiterjoin%
\definecolor{currentfill}{rgb}{0.000000,0.419608,0.790196}%
\pgfsetfillcolor{currentfill}%
\pgfsetlinewidth{0.000000pt}%
\definecolor{currentstroke}{rgb}{0.000000,0.000000,0.000000}%
\pgfsetstrokecolor{currentstroke}%
\pgfsetstrokeopacity{0.000000}%
\pgfsetdash{}{0pt}%
\pgfpathmoveto{\pgfqpoint{2.025168in}{1.813601in}}%
\pgfpathlineto{\pgfqpoint{2.024747in}{1.786093in}}%
\pgfpathlineto{\pgfqpoint{2.024621in}{1.779133in}}%
\pgfpathlineto{\pgfqpoint{1.986361in}{1.779846in}}%
\pgfpathlineto{\pgfqpoint{1.945343in}{1.781068in}}%
\pgfpathlineto{\pgfqpoint{1.945390in}{1.787969in}}%
\pgfpathlineto{\pgfqpoint{1.938531in}{1.788233in}}%
\pgfpathlineto{\pgfqpoint{1.939662in}{1.829664in}}%
\pgfpathlineto{\pgfqpoint{1.981983in}{1.828333in}}%
\pgfpathlineto{\pgfqpoint{2.001203in}{1.827773in}}%
\pgfpathlineto{\pgfqpoint{2.000983in}{1.813993in}}%
\pgfpathlineto{\pgfqpoint{2.025168in}{1.813601in}}%
\pgfpathclose%
\pgfusepath{fill}%
\end{pgfscope}%
\begin{pgfscope}%
\pgfpathrectangle{\pgfqpoint{0.100000in}{0.100000in}}{\pgfqpoint{3.608454in}{2.310000in}}%
\pgfusepath{clip}%
\pgfsetbuttcap%
\pgfsetmiterjoin%
\definecolor{currentfill}{rgb}{0.000000,0.678431,0.660784}%
\pgfsetfillcolor{currentfill}%
\pgfsetlinewidth{0.000000pt}%
\definecolor{currentstroke}{rgb}{0.000000,0.000000,0.000000}%
\pgfsetstrokecolor{currentstroke}%
\pgfsetstrokeopacity{0.000000}%
\pgfsetdash{}{0pt}%
\pgfpathmoveto{\pgfqpoint{1.919219in}{0.761725in}}%
\pgfpathlineto{\pgfqpoint{1.884281in}{0.742465in}}%
\pgfpathlineto{\pgfqpoint{1.870147in}{0.767494in}}%
\pgfpathlineto{\pgfqpoint{1.858178in}{0.766124in}}%
\pgfpathlineto{\pgfqpoint{1.842140in}{0.796889in}}%
\pgfpathlineto{\pgfqpoint{1.873122in}{0.813238in}}%
\pgfpathlineto{\pgfqpoint{1.873743in}{0.830250in}}%
\pgfpathlineto{\pgfqpoint{1.900998in}{0.829156in}}%
\pgfpathlineto{\pgfqpoint{1.908521in}{0.806762in}}%
\pgfpathlineto{\pgfqpoint{1.913591in}{0.794961in}}%
\pgfpathlineto{\pgfqpoint{1.903937in}{0.789730in}}%
\pgfpathlineto{\pgfqpoint{1.919219in}{0.761725in}}%
\pgfpathclose%
\pgfusepath{fill}%
\end{pgfscope}%
\begin{pgfscope}%
\pgfpathrectangle{\pgfqpoint{0.100000in}{0.100000in}}{\pgfqpoint{3.608454in}{2.310000in}}%
\pgfusepath{clip}%
\pgfsetbuttcap%
\pgfsetmiterjoin%
\definecolor{currentfill}{rgb}{0.000000,0.580392,0.709804}%
\pgfsetfillcolor{currentfill}%
\pgfsetlinewidth{0.000000pt}%
\definecolor{currentstroke}{rgb}{0.000000,0.000000,0.000000}%
\pgfsetstrokecolor{currentstroke}%
\pgfsetstrokeopacity{0.000000}%
\pgfsetdash{}{0pt}%
\pgfpathmoveto{\pgfqpoint{3.111250in}{1.435310in}}%
\pgfpathlineto{\pgfqpoint{3.117131in}{1.434042in}}%
\pgfpathlineto{\pgfqpoint{3.119217in}{1.439782in}}%
\pgfpathlineto{\pgfqpoint{3.125843in}{1.434080in}}%
\pgfpathlineto{\pgfqpoint{3.127584in}{1.429231in}}%
\pgfpathlineto{\pgfqpoint{3.124006in}{1.419453in}}%
\pgfpathlineto{\pgfqpoint{3.131195in}{1.415085in}}%
\pgfpathlineto{\pgfqpoint{3.129858in}{1.404228in}}%
\pgfpathlineto{\pgfqpoint{3.125203in}{1.394430in}}%
\pgfpathlineto{\pgfqpoint{3.123623in}{1.385333in}}%
\pgfpathlineto{\pgfqpoint{3.115799in}{1.372179in}}%
\pgfpathlineto{\pgfqpoint{3.110807in}{1.366936in}}%
\pgfpathlineto{\pgfqpoint{3.078058in}{1.382882in}}%
\pgfpathlineto{\pgfqpoint{3.073594in}{1.376808in}}%
\pgfpathlineto{\pgfqpoint{3.062954in}{1.378601in}}%
\pgfpathlineto{\pgfqpoint{3.058026in}{1.385268in}}%
\pgfpathlineto{\pgfqpoint{3.050650in}{1.387261in}}%
\pgfpathlineto{\pgfqpoint{3.056652in}{1.404504in}}%
\pgfpathlineto{\pgfqpoint{3.055324in}{1.410700in}}%
\pgfpathlineto{\pgfqpoint{3.063909in}{1.418899in}}%
\pgfpathlineto{\pgfqpoint{3.084328in}{1.406748in}}%
\pgfpathlineto{\pgfqpoint{3.086938in}{1.414331in}}%
\pgfpathlineto{\pgfqpoint{3.095711in}{1.408904in}}%
\pgfpathlineto{\pgfqpoint{3.099019in}{1.419473in}}%
\pgfpathlineto{\pgfqpoint{3.103999in}{1.422556in}}%
\pgfpathlineto{\pgfqpoint{3.111250in}{1.435310in}}%
\pgfpathclose%
\pgfusepath{fill}%
\end{pgfscope}%
\begin{pgfscope}%
\pgfpathrectangle{\pgfqpoint{0.100000in}{0.100000in}}{\pgfqpoint{3.608454in}{2.310000in}}%
\pgfusepath{clip}%
\pgfsetbuttcap%
\pgfsetmiterjoin%
\definecolor{currentfill}{rgb}{0.000000,0.741176,0.629412}%
\pgfsetfillcolor{currentfill}%
\pgfsetlinewidth{0.000000pt}%
\definecolor{currentstroke}{rgb}{0.000000,0.000000,0.000000}%
\pgfsetstrokecolor{currentstroke}%
\pgfsetstrokeopacity{0.000000}%
\pgfsetdash{}{0pt}%
\pgfpathmoveto{\pgfqpoint{2.414649in}{1.160533in}}%
\pgfpathlineto{\pgfqpoint{2.410619in}{1.149798in}}%
\pgfpathlineto{\pgfqpoint{2.353073in}{1.147045in}}%
\pgfpathlineto{\pgfqpoint{2.352645in}{1.172837in}}%
\pgfpathlineto{\pgfqpoint{2.345750in}{1.172705in}}%
\pgfpathlineto{\pgfqpoint{2.345378in}{1.193313in}}%
\pgfpathlineto{\pgfqpoint{2.357912in}{1.194600in}}%
\pgfpathlineto{\pgfqpoint{2.361419in}{1.191422in}}%
\pgfpathlineto{\pgfqpoint{2.373090in}{1.191695in}}%
\pgfpathlineto{\pgfqpoint{2.375334in}{1.200990in}}%
\pgfpathlineto{\pgfqpoint{2.375052in}{1.209261in}}%
\pgfpathlineto{\pgfqpoint{2.386664in}{1.209942in}}%
\pgfpathlineto{\pgfqpoint{2.386546in}{1.213497in}}%
\pgfpathlineto{\pgfqpoint{2.407413in}{1.214350in}}%
\pgfpathlineto{\pgfqpoint{2.408011in}{1.196285in}}%
\pgfpathlineto{\pgfqpoint{2.415125in}{1.196618in}}%
\pgfpathlineto{\pgfqpoint{2.406457in}{1.183183in}}%
\pgfpathlineto{\pgfqpoint{2.413654in}{1.167791in}}%
\pgfpathlineto{\pgfqpoint{2.414649in}{1.160533in}}%
\pgfpathclose%
\pgfusepath{fill}%
\end{pgfscope}%
\begin{pgfscope}%
\pgfpathrectangle{\pgfqpoint{0.100000in}{0.100000in}}{\pgfqpoint{3.608454in}{2.310000in}}%
\pgfusepath{clip}%
\pgfsetbuttcap%
\pgfsetmiterjoin%
\definecolor{currentfill}{rgb}{0.000000,0.611765,0.694118}%
\pgfsetfillcolor{currentfill}%
\pgfsetlinewidth{0.000000pt}%
\definecolor{currentstroke}{rgb}{0.000000,0.000000,0.000000}%
\pgfsetstrokecolor{currentstroke}%
\pgfsetstrokeopacity{0.000000}%
\pgfsetdash{}{0pt}%
\pgfpathmoveto{\pgfqpoint{1.578022in}{2.163356in}}%
\pgfpathlineto{\pgfqpoint{1.630877in}{2.157836in}}%
\pgfpathlineto{\pgfqpoint{1.626106in}{2.109098in}}%
\pgfpathlineto{\pgfqpoint{1.595340in}{2.112255in}}%
\pgfpathlineto{\pgfqpoint{1.596092in}{2.119196in}}%
\pgfpathlineto{\pgfqpoint{1.589220in}{2.119949in}}%
\pgfpathlineto{\pgfqpoint{1.589978in}{2.126848in}}%
\pgfpathlineto{\pgfqpoint{1.578597in}{2.128096in}}%
\pgfpathlineto{\pgfqpoint{1.575926in}{2.135412in}}%
\pgfpathlineto{\pgfqpoint{1.578243in}{2.156172in}}%
\pgfpathlineto{\pgfqpoint{1.578022in}{2.163356in}}%
\pgfpathclose%
\pgfusepath{fill}%
\end{pgfscope}%
\begin{pgfscope}%
\pgfpathrectangle{\pgfqpoint{0.100000in}{0.100000in}}{\pgfqpoint{3.608454in}{2.310000in}}%
\pgfusepath{clip}%
\pgfsetbuttcap%
\pgfsetmiterjoin%
\definecolor{currentfill}{rgb}{0.000000,0.584314,0.707843}%
\pgfsetfillcolor{currentfill}%
\pgfsetlinewidth{0.000000pt}%
\definecolor{currentstroke}{rgb}{0.000000,0.000000,0.000000}%
\pgfsetstrokecolor{currentstroke}%
\pgfsetstrokeopacity{0.000000}%
\pgfsetdash{}{0pt}%
\pgfpathmoveto{\pgfqpoint{2.118427in}{1.958403in}}%
\pgfpathlineto{\pgfqpoint{2.097902in}{1.958459in}}%
\pgfpathlineto{\pgfqpoint{2.097851in}{1.951504in}}%
\pgfpathlineto{\pgfqpoint{2.042559in}{1.951936in}}%
\pgfpathlineto{\pgfqpoint{2.042858in}{1.972645in}}%
\pgfpathlineto{\pgfqpoint{2.041913in}{1.986517in}}%
\pgfpathlineto{\pgfqpoint{2.048844in}{1.986478in}}%
\pgfpathlineto{\pgfqpoint{2.049145in}{2.014243in}}%
\pgfpathlineto{\pgfqpoint{2.076864in}{2.014087in}}%
\pgfpathlineto{\pgfqpoint{2.075622in}{2.027945in}}%
\pgfpathlineto{\pgfqpoint{2.075577in}{2.048780in}}%
\pgfpathlineto{\pgfqpoint{2.075068in}{2.055715in}}%
\pgfpathlineto{\pgfqpoint{2.074816in}{2.097239in}}%
\pgfpathlineto{\pgfqpoint{2.095537in}{2.097237in}}%
\pgfpathlineto{\pgfqpoint{2.095560in}{2.083291in}}%
\pgfpathlineto{\pgfqpoint{2.137206in}{2.083483in}}%
\pgfpathlineto{\pgfqpoint{2.138015in}{2.041726in}}%
\pgfpathlineto{\pgfqpoint{2.138354in}{2.007024in}}%
\pgfpathlineto{\pgfqpoint{2.124628in}{2.006820in}}%
\pgfpathlineto{\pgfqpoint{2.125367in}{1.958378in}}%
\pgfpathlineto{\pgfqpoint{2.118427in}{1.958403in}}%
\pgfpathclose%
\pgfusepath{fill}%
\end{pgfscope}%
\begin{pgfscope}%
\pgfpathrectangle{\pgfqpoint{0.100000in}{0.100000in}}{\pgfqpoint{3.608454in}{2.310000in}}%
\pgfusepath{clip}%
\pgfsetbuttcap%
\pgfsetmiterjoin%
\definecolor{currentfill}{rgb}{0.000000,0.666667,0.666667}%
\pgfsetfillcolor{currentfill}%
\pgfsetlinewidth{0.000000pt}%
\definecolor{currentstroke}{rgb}{0.000000,0.000000,0.000000}%
\pgfsetstrokecolor{currentstroke}%
\pgfsetstrokeopacity{0.000000}%
\pgfsetdash{}{0pt}%
\pgfpathmoveto{\pgfqpoint{2.272778in}{0.593877in}}%
\pgfpathlineto{\pgfqpoint{2.271892in}{0.630559in}}%
\pgfpathlineto{\pgfqpoint{2.269704in}{0.666105in}}%
\pgfpathlineto{\pgfqpoint{2.272007in}{0.674022in}}%
\pgfpathlineto{\pgfqpoint{2.271317in}{0.699125in}}%
\pgfpathlineto{\pgfqpoint{2.278789in}{0.700272in}}%
\pgfpathlineto{\pgfqpoint{2.285347in}{0.707925in}}%
\pgfpathlineto{\pgfqpoint{2.292928in}{0.702917in}}%
\pgfpathlineto{\pgfqpoint{2.296084in}{0.696005in}}%
\pgfpathlineto{\pgfqpoint{2.325022in}{0.696944in}}%
\pgfpathlineto{\pgfqpoint{2.331030in}{0.685463in}}%
\pgfpathlineto{\pgfqpoint{2.329801in}{0.674430in}}%
\pgfpathlineto{\pgfqpoint{2.331431in}{0.670010in}}%
\pgfpathlineto{\pgfqpoint{2.338216in}{0.665019in}}%
\pgfpathlineto{\pgfqpoint{2.339922in}{0.654462in}}%
\pgfpathlineto{\pgfqpoint{2.344148in}{0.649181in}}%
\pgfpathlineto{\pgfqpoint{2.350223in}{0.648417in}}%
\pgfpathlineto{\pgfqpoint{2.351410in}{0.638299in}}%
\pgfpathlineto{\pgfqpoint{2.368222in}{0.632799in}}%
\pgfpathlineto{\pgfqpoint{2.365882in}{0.630690in}}%
\pgfpathlineto{\pgfqpoint{2.371581in}{0.618635in}}%
\pgfpathlineto{\pgfqpoint{2.376744in}{0.617237in}}%
\pgfpathlineto{\pgfqpoint{2.378486in}{0.611623in}}%
\pgfpathlineto{\pgfqpoint{2.372197in}{0.609200in}}%
\pgfpathlineto{\pgfqpoint{2.357771in}{0.613019in}}%
\pgfpathlineto{\pgfqpoint{2.358501in}{0.617112in}}%
\pgfpathlineto{\pgfqpoint{2.351303in}{0.626566in}}%
\pgfpathlineto{\pgfqpoint{2.340753in}{0.625307in}}%
\pgfpathlineto{\pgfqpoint{2.335052in}{0.617153in}}%
\pgfpathlineto{\pgfqpoint{2.323648in}{0.606472in}}%
\pgfpathlineto{\pgfqpoint{2.323245in}{0.615819in}}%
\pgfpathlineto{\pgfqpoint{2.312586in}{0.611195in}}%
\pgfpathlineto{\pgfqpoint{2.305299in}{0.605725in}}%
\pgfpathlineto{\pgfqpoint{2.307552in}{0.598624in}}%
\pgfpathlineto{\pgfqpoint{2.315848in}{0.597697in}}%
\pgfpathlineto{\pgfqpoint{2.326725in}{0.599490in}}%
\pgfpathlineto{\pgfqpoint{2.336035in}{0.593898in}}%
\pgfpathlineto{\pgfqpoint{2.328591in}{0.587197in}}%
\pgfpathlineto{\pgfqpoint{2.311674in}{0.595112in}}%
\pgfpathlineto{\pgfqpoint{2.304943in}{0.595012in}}%
\pgfpathlineto{\pgfqpoint{2.294524in}{0.590708in}}%
\pgfpathlineto{\pgfqpoint{2.272778in}{0.593877in}}%
\pgfpathclose%
\pgfusepath{fill}%
\end{pgfscope}%
\begin{pgfscope}%
\pgfpathrectangle{\pgfqpoint{0.100000in}{0.100000in}}{\pgfqpoint{3.608454in}{2.310000in}}%
\pgfusepath{clip}%
\pgfsetbuttcap%
\pgfsetmiterjoin%
\definecolor{currentfill}{rgb}{0.000000,0.682353,0.658824}%
\pgfsetfillcolor{currentfill}%
\pgfsetlinewidth{0.000000pt}%
\definecolor{currentstroke}{rgb}{0.000000,0.000000,0.000000}%
\pgfsetstrokecolor{currentstroke}%
\pgfsetstrokeopacity{0.000000}%
\pgfsetdash{}{0pt}%
\pgfpathmoveto{\pgfqpoint{1.836840in}{0.475148in}}%
\pgfpathlineto{\pgfqpoint{1.795200in}{0.474847in}}%
\pgfpathlineto{\pgfqpoint{1.795506in}{0.488863in}}%
\pgfpathlineto{\pgfqpoint{1.797091in}{0.523880in}}%
\pgfpathlineto{\pgfqpoint{1.796154in}{0.523908in}}%
\pgfpathlineto{\pgfqpoint{1.797569in}{0.560061in}}%
\pgfpathlineto{\pgfqpoint{1.840202in}{0.558059in}}%
\pgfpathlineto{\pgfqpoint{1.839000in}{0.522528in}}%
\pgfpathlineto{\pgfqpoint{1.836840in}{0.475148in}}%
\pgfpathclose%
\pgfusepath{fill}%
\end{pgfscope}%
\begin{pgfscope}%
\pgfpathrectangle{\pgfqpoint{0.100000in}{0.100000in}}{\pgfqpoint{3.608454in}{2.310000in}}%
\pgfusepath{clip}%
\pgfsetbuttcap%
\pgfsetmiterjoin%
\definecolor{currentfill}{rgb}{0.000000,0.588235,0.705882}%
\pgfsetfillcolor{currentfill}%
\pgfsetlinewidth{0.000000pt}%
\definecolor{currentstroke}{rgb}{0.000000,0.000000,0.000000}%
\pgfsetstrokecolor{currentstroke}%
\pgfsetstrokeopacity{0.000000}%
\pgfsetdash{}{0pt}%
\pgfpathmoveto{\pgfqpoint{2.966459in}{0.731456in}}%
\pgfpathlineto{\pgfqpoint{2.957900in}{0.730883in}}%
\pgfpathlineto{\pgfqpoint{2.955222in}{0.740323in}}%
\pgfpathlineto{\pgfqpoint{2.943923in}{0.740325in}}%
\pgfpathlineto{\pgfqpoint{2.937427in}{0.747697in}}%
\pgfpathlineto{\pgfqpoint{2.928578in}{0.749275in}}%
\pgfpathlineto{\pgfqpoint{2.925268in}{0.774204in}}%
\pgfpathlineto{\pgfqpoint{2.920160in}{0.773477in}}%
\pgfpathlineto{\pgfqpoint{2.920094in}{0.780812in}}%
\pgfpathlineto{\pgfqpoint{2.911798in}{0.789691in}}%
\pgfpathlineto{\pgfqpoint{2.910462in}{0.795395in}}%
\pgfpathlineto{\pgfqpoint{2.912735in}{0.796174in}}%
\pgfpathlineto{\pgfqpoint{2.914876in}{0.807165in}}%
\pgfpathlineto{\pgfqpoint{2.918401in}{0.811582in}}%
\pgfpathlineto{\pgfqpoint{2.917347in}{0.821092in}}%
\pgfpathlineto{\pgfqpoint{2.926243in}{0.822499in}}%
\pgfpathlineto{\pgfqpoint{2.929182in}{0.813904in}}%
\pgfpathlineto{\pgfqpoint{2.953500in}{0.817959in}}%
\pgfpathlineto{\pgfqpoint{2.970476in}{0.810318in}}%
\pgfpathlineto{\pgfqpoint{2.978722in}{0.804662in}}%
\pgfpathlineto{\pgfqpoint{2.983849in}{0.804589in}}%
\pgfpathlineto{\pgfqpoint{2.986023in}{0.797823in}}%
\pgfpathlineto{\pgfqpoint{2.990030in}{0.793119in}}%
\pgfpathlineto{\pgfqpoint{2.983748in}{0.789742in}}%
\pgfpathlineto{\pgfqpoint{2.977554in}{0.783183in}}%
\pgfpathlineto{\pgfqpoint{2.977965in}{0.773429in}}%
\pgfpathlineto{\pgfqpoint{2.983960in}{0.768316in}}%
\pgfpathlineto{\pgfqpoint{2.964698in}{0.765826in}}%
\pgfpathlineto{\pgfqpoint{2.966799in}{0.748504in}}%
\pgfpathlineto{\pgfqpoint{2.985291in}{0.750228in}}%
\pgfpathlineto{\pgfqpoint{2.983300in}{0.732552in}}%
\pgfpathlineto{\pgfqpoint{2.966459in}{0.731456in}}%
\pgfpathclose%
\pgfusepath{fill}%
\end{pgfscope}%
\begin{pgfscope}%
\pgfpathrectangle{\pgfqpoint{0.100000in}{0.100000in}}{\pgfqpoint{3.608454in}{2.310000in}}%
\pgfusepath{clip}%
\pgfsetbuttcap%
\pgfsetmiterjoin%
\definecolor{currentfill}{rgb}{0.000000,0.670588,0.664706}%
\pgfsetfillcolor{currentfill}%
\pgfsetlinewidth{0.000000pt}%
\definecolor{currentstroke}{rgb}{0.000000,0.000000,0.000000}%
\pgfsetstrokecolor{currentstroke}%
\pgfsetstrokeopacity{0.000000}%
\pgfsetdash{}{0pt}%
\pgfpathmoveto{\pgfqpoint{0.579411in}{1.933417in}}%
\pgfpathlineto{\pgfqpoint{0.571017in}{1.931177in}}%
\pgfpathlineto{\pgfqpoint{0.568235in}{1.922352in}}%
\pgfpathlineto{\pgfqpoint{0.534705in}{1.932520in}}%
\pgfpathlineto{\pgfqpoint{0.537209in}{1.940659in}}%
\pgfpathlineto{\pgfqpoint{0.517068in}{1.946672in}}%
\pgfpathlineto{\pgfqpoint{0.517042in}{1.952171in}}%
\pgfpathlineto{\pgfqpoint{0.521206in}{1.965301in}}%
\pgfpathlineto{\pgfqpoint{0.509770in}{1.969016in}}%
\pgfpathlineto{\pgfqpoint{0.501091in}{1.975934in}}%
\pgfpathlineto{\pgfqpoint{0.498412in}{1.984367in}}%
\pgfpathlineto{\pgfqpoint{0.487714in}{1.989945in}}%
\pgfpathlineto{\pgfqpoint{0.479271in}{1.985850in}}%
\pgfpathlineto{\pgfqpoint{0.467729in}{1.989714in}}%
\pgfpathlineto{\pgfqpoint{0.475658in}{2.007052in}}%
\pgfpathlineto{\pgfqpoint{0.480357in}{2.020001in}}%
\pgfpathlineto{\pgfqpoint{0.498912in}{2.014543in}}%
\pgfpathlineto{\pgfqpoint{0.504063in}{2.019289in}}%
\pgfpathlineto{\pgfqpoint{0.505570in}{2.024850in}}%
\pgfpathlineto{\pgfqpoint{0.512206in}{2.022672in}}%
\pgfpathlineto{\pgfqpoint{0.519006in}{2.044536in}}%
\pgfpathlineto{\pgfqpoint{0.512881in}{2.048028in}}%
\pgfpathlineto{\pgfqpoint{0.521244in}{2.073470in}}%
\pgfpathlineto{\pgfqpoint{0.518018in}{2.074545in}}%
\pgfpathlineto{\pgfqpoint{0.521499in}{2.085093in}}%
\pgfpathlineto{\pgfqpoint{0.538605in}{2.079539in}}%
\pgfpathlineto{\pgfqpoint{0.543848in}{2.096029in}}%
\pgfpathlineto{\pgfqpoint{0.561348in}{2.090351in}}%
\pgfpathlineto{\pgfqpoint{0.566154in}{2.086439in}}%
\pgfpathlineto{\pgfqpoint{0.568024in}{2.080965in}}%
\pgfpathlineto{\pgfqpoint{0.576179in}{2.072349in}}%
\pgfpathlineto{\pgfqpoint{0.575705in}{2.064800in}}%
\pgfpathlineto{\pgfqpoint{0.572541in}{2.061484in}}%
\pgfpathlineto{\pgfqpoint{0.575268in}{2.054730in}}%
\pgfpathlineto{\pgfqpoint{0.580214in}{2.049876in}}%
\pgfpathlineto{\pgfqpoint{0.581481in}{2.042623in}}%
\pgfpathlineto{\pgfqpoint{0.587828in}{2.036379in}}%
\pgfpathlineto{\pgfqpoint{0.623427in}{2.025636in}}%
\pgfpathlineto{\pgfqpoint{0.620648in}{2.021281in}}%
\pgfpathlineto{\pgfqpoint{0.616042in}{2.014692in}}%
\pgfpathlineto{\pgfqpoint{0.613506in}{2.004223in}}%
\pgfpathlineto{\pgfqpoint{0.608267in}{1.997300in}}%
\pgfpathlineto{\pgfqpoint{0.605387in}{1.987819in}}%
\pgfpathlineto{\pgfqpoint{0.605729in}{1.977984in}}%
\pgfpathlineto{\pgfqpoint{0.603007in}{1.965607in}}%
\pgfpathlineto{\pgfqpoint{0.595256in}{1.957298in}}%
\pgfpathlineto{\pgfqpoint{0.586510in}{1.951287in}}%
\pgfpathlineto{\pgfqpoint{0.585151in}{1.943983in}}%
\pgfpathlineto{\pgfqpoint{0.579411in}{1.933417in}}%
\pgfpathclose%
\pgfusepath{fill}%
\end{pgfscope}%
\begin{pgfscope}%
\pgfpathrectangle{\pgfqpoint{0.100000in}{0.100000in}}{\pgfqpoint{3.608454in}{2.310000in}}%
\pgfusepath{clip}%
\pgfsetbuttcap%
\pgfsetmiterjoin%
\definecolor{currentfill}{rgb}{0.000000,0.545098,0.727451}%
\pgfsetfillcolor{currentfill}%
\pgfsetlinewidth{0.000000pt}%
\definecolor{currentstroke}{rgb}{0.000000,0.000000,0.000000}%
\pgfsetstrokecolor{currentstroke}%
\pgfsetstrokeopacity{0.000000}%
\pgfsetdash{}{0pt}%
\pgfpathmoveto{\pgfqpoint{2.765201in}{0.880136in}}%
\pgfpathlineto{\pgfqpoint{2.755362in}{0.877808in}}%
\pgfpathlineto{\pgfqpoint{2.734886in}{0.875470in}}%
\pgfpathlineto{\pgfqpoint{2.731773in}{0.905420in}}%
\pgfpathlineto{\pgfqpoint{2.727775in}{0.904950in}}%
\pgfpathlineto{\pgfqpoint{2.725270in}{0.935776in}}%
\pgfpathlineto{\pgfqpoint{2.717216in}{0.935165in}}%
\pgfpathlineto{\pgfqpoint{2.714599in}{0.941836in}}%
\pgfpathlineto{\pgfqpoint{2.701405in}{0.940587in}}%
\pgfpathlineto{\pgfqpoint{2.697028in}{0.947184in}}%
\pgfpathlineto{\pgfqpoint{2.690775in}{0.946928in}}%
\pgfpathlineto{\pgfqpoint{2.696367in}{0.955422in}}%
\pgfpathlineto{\pgfqpoint{2.694783in}{0.960330in}}%
\pgfpathlineto{\pgfqpoint{2.705960in}{0.970065in}}%
\pgfpathlineto{\pgfqpoint{2.715165in}{0.972447in}}%
\pgfpathlineto{\pgfqpoint{2.728983in}{0.972881in}}%
\pgfpathlineto{\pgfqpoint{2.737463in}{0.974424in}}%
\pgfpathlineto{\pgfqpoint{2.738788in}{0.969584in}}%
\pgfpathlineto{\pgfqpoint{2.760784in}{0.972182in}}%
\pgfpathlineto{\pgfqpoint{2.762398in}{0.964970in}}%
\pgfpathlineto{\pgfqpoint{2.771603in}{0.963496in}}%
\pgfpathlineto{\pgfqpoint{2.773147in}{0.947109in}}%
\pgfpathlineto{\pgfqpoint{2.778830in}{0.942946in}}%
\pgfpathlineto{\pgfqpoint{2.791999in}{0.944257in}}%
\pgfpathlineto{\pgfqpoint{2.796122in}{0.932528in}}%
\pgfpathlineto{\pgfqpoint{2.802892in}{0.925244in}}%
\pgfpathlineto{\pgfqpoint{2.802911in}{0.922322in}}%
\pgfpathlineto{\pgfqpoint{2.779766in}{0.919893in}}%
\pgfpathlineto{\pgfqpoint{2.782280in}{0.891779in}}%
\pgfpathlineto{\pgfqpoint{2.760875in}{0.889582in}}%
\pgfpathlineto{\pgfqpoint{2.765201in}{0.880136in}}%
\pgfpathclose%
\pgfusepath{fill}%
\end{pgfscope}%
\begin{pgfscope}%
\pgfpathrectangle{\pgfqpoint{0.100000in}{0.100000in}}{\pgfqpoint{3.608454in}{2.310000in}}%
\pgfusepath{clip}%
\pgfsetbuttcap%
\pgfsetmiterjoin%
\definecolor{currentfill}{rgb}{0.000000,0.556863,0.721569}%
\pgfsetfillcolor{currentfill}%
\pgfsetlinewidth{0.000000pt}%
\definecolor{currentstroke}{rgb}{0.000000,0.000000,0.000000}%
\pgfsetstrokecolor{currentstroke}%
\pgfsetstrokeopacity{0.000000}%
\pgfsetdash{}{0pt}%
\pgfpathmoveto{\pgfqpoint{1.305818in}{1.299625in}}%
\pgfpathlineto{\pgfqpoint{1.298775in}{1.307121in}}%
\pgfpathlineto{\pgfqpoint{1.284736in}{1.306313in}}%
\pgfpathlineto{\pgfqpoint{1.282717in}{1.312481in}}%
\pgfpathlineto{\pgfqpoint{1.248158in}{1.317741in}}%
\pgfpathlineto{\pgfqpoint{1.234133in}{1.318758in}}%
\pgfpathlineto{\pgfqpoint{1.237566in}{1.340907in}}%
\pgfpathlineto{\pgfqpoint{1.237792in}{1.349736in}}%
\pgfpathlineto{\pgfqpoint{1.240533in}{1.367269in}}%
\pgfpathlineto{\pgfqpoint{1.251645in}{1.434868in}}%
\pgfpathlineto{\pgfqpoint{1.290367in}{1.428918in}}%
\pgfpathlineto{\pgfqpoint{1.349566in}{1.420460in}}%
\pgfpathlineto{\pgfqpoint{1.346932in}{1.416687in}}%
\pgfpathlineto{\pgfqpoint{1.350593in}{1.411552in}}%
\pgfpathlineto{\pgfqpoint{1.343717in}{1.409436in}}%
\pgfpathlineto{\pgfqpoint{1.333852in}{1.337711in}}%
\pgfpathlineto{\pgfqpoint{1.325556in}{1.338863in}}%
\pgfpathlineto{\pgfqpoint{1.328146in}{1.332828in}}%
\pgfpathlineto{\pgfqpoint{1.328108in}{1.320759in}}%
\pgfpathlineto{\pgfqpoint{1.325950in}{1.311919in}}%
\pgfpathlineto{\pgfqpoint{1.314836in}{1.308720in}}%
\pgfpathlineto{\pgfqpoint{1.305818in}{1.299625in}}%
\pgfpathclose%
\pgfusepath{fill}%
\end{pgfscope}%
\begin{pgfscope}%
\pgfpathrectangle{\pgfqpoint{0.100000in}{0.100000in}}{\pgfqpoint{3.608454in}{2.310000in}}%
\pgfusepath{clip}%
\pgfsetbuttcap%
\pgfsetmiterjoin%
\definecolor{currentfill}{rgb}{0.000000,0.827451,0.586275}%
\pgfsetfillcolor{currentfill}%
\pgfsetlinewidth{0.000000pt}%
\definecolor{currentstroke}{rgb}{0.000000,0.000000,0.000000}%
\pgfsetstrokecolor{currentstroke}%
\pgfsetstrokeopacity{0.000000}%
\pgfsetdash{}{0pt}%
\pgfpathmoveto{\pgfqpoint{1.124077in}{0.559590in}}%
\pgfpathlineto{\pgfqpoint{1.110388in}{0.567904in}}%
\pgfpathlineto{\pgfqpoint{1.083089in}{0.584700in}}%
\pgfpathlineto{\pgfqpoint{1.074192in}{0.589904in}}%
\pgfpathlineto{\pgfqpoint{1.066205in}{0.578011in}}%
\pgfpathlineto{\pgfqpoint{1.056440in}{0.584227in}}%
\pgfpathlineto{\pgfqpoint{1.052163in}{0.586365in}}%
\pgfpathlineto{\pgfqpoint{1.026608in}{0.603847in}}%
\pgfpathlineto{\pgfqpoint{0.997892in}{0.624739in}}%
\pgfpathlineto{\pgfqpoint{0.976297in}{0.641465in}}%
\pgfpathlineto{\pgfqpoint{0.976829in}{0.642165in}}%
\pgfpathlineto{\pgfqpoint{0.968183in}{0.648986in}}%
\pgfpathlineto{\pgfqpoint{0.969832in}{0.651052in}}%
\pgfpathlineto{\pgfqpoint{0.961484in}{0.657889in}}%
\pgfpathlineto{\pgfqpoint{0.963280in}{0.660093in}}%
\pgfpathlineto{\pgfqpoint{0.954477in}{0.667338in}}%
\pgfpathlineto{\pgfqpoint{0.940158in}{0.679506in}}%
\pgfpathlineto{\pgfqpoint{0.941559in}{0.681130in}}%
\pgfpathlineto{\pgfqpoint{0.937775in}{0.684411in}}%
\pgfpathlineto{\pgfqpoint{0.936363in}{0.682792in}}%
\pgfpathlineto{\pgfqpoint{0.927789in}{0.690351in}}%
\pgfpathlineto{\pgfqpoint{0.929228in}{0.691967in}}%
\pgfpathlineto{\pgfqpoint{0.925549in}{0.695258in}}%
\pgfpathlineto{\pgfqpoint{0.924102in}{0.693649in}}%
\pgfpathlineto{\pgfqpoint{0.919895in}{0.697448in}}%
\pgfpathlineto{\pgfqpoint{0.921351in}{0.699050in}}%
\pgfpathlineto{\pgfqpoint{0.915131in}{0.704745in}}%
\pgfpathlineto{\pgfqpoint{0.913662in}{0.703154in}}%
\pgfpathlineto{\pgfqpoint{0.898458in}{0.717465in}}%
\pgfpathlineto{\pgfqpoint{0.894146in}{0.712972in}}%
\pgfpathlineto{\pgfqpoint{0.886741in}{0.720168in}}%
\pgfpathlineto{\pgfqpoint{0.883880in}{0.727259in}}%
\pgfpathlineto{\pgfqpoint{0.884750in}{0.730564in}}%
\pgfpathlineto{\pgfqpoint{0.884173in}{0.733970in}}%
\pgfpathlineto{\pgfqpoint{0.881741in}{0.738770in}}%
\pgfpathlineto{\pgfqpoint{0.884973in}{0.737504in}}%
\pgfpathlineto{\pgfqpoint{0.888394in}{0.737302in}}%
\pgfpathlineto{\pgfqpoint{0.891146in}{0.738897in}}%
\pgfpathlineto{\pgfqpoint{0.897248in}{0.744755in}}%
\pgfpathlineto{\pgfqpoint{0.899391in}{0.744053in}}%
\pgfpathlineto{\pgfqpoint{0.904078in}{0.738998in}}%
\pgfpathlineto{\pgfqpoint{0.914603in}{0.732006in}}%
\pgfpathlineto{\pgfqpoint{0.917793in}{0.730737in}}%
\pgfpathlineto{\pgfqpoint{0.924697in}{0.730323in}}%
\pgfpathlineto{\pgfqpoint{0.934212in}{0.732247in}}%
\pgfpathlineto{\pgfqpoint{0.938989in}{0.735849in}}%
\pgfpathlineto{\pgfqpoint{0.940161in}{0.737749in}}%
\pgfpathlineto{\pgfqpoint{0.943533in}{0.740578in}}%
\pgfpathlineto{\pgfqpoint{0.945793in}{0.741496in}}%
\pgfpathlineto{\pgfqpoint{0.952899in}{0.741432in}}%
\pgfpathlineto{\pgfqpoint{0.955066in}{0.743565in}}%
\pgfpathlineto{\pgfqpoint{0.959761in}{0.743939in}}%
\pgfpathlineto{\pgfqpoint{0.965368in}{0.745161in}}%
\pgfpathlineto{\pgfqpoint{0.968573in}{0.740621in}}%
\pgfpathlineto{\pgfqpoint{0.974554in}{0.739226in}}%
\pgfpathlineto{\pgfqpoint{0.985159in}{0.739890in}}%
\pgfpathlineto{\pgfqpoint{0.995988in}{0.741765in}}%
\pgfpathlineto{\pgfqpoint{0.999534in}{0.740148in}}%
\pgfpathlineto{\pgfqpoint{0.999964in}{0.737676in}}%
\pgfpathlineto{\pgfqpoint{1.004192in}{0.734276in}}%
\pgfpathlineto{\pgfqpoint{1.008232in}{0.732003in}}%
\pgfpathlineto{\pgfqpoint{1.011631in}{0.730934in}}%
\pgfpathlineto{\pgfqpoint{1.018774in}{0.731477in}}%
\pgfpathlineto{\pgfqpoint{1.028230in}{0.734980in}}%
\pgfpathlineto{\pgfqpoint{1.031476in}{0.735474in}}%
\pgfpathlineto{\pgfqpoint{1.033074in}{0.730543in}}%
\pgfpathlineto{\pgfqpoint{1.032211in}{0.727412in}}%
\pgfpathlineto{\pgfqpoint{1.034678in}{0.727466in}}%
\pgfpathlineto{\pgfqpoint{1.036005in}{0.725068in}}%
\pgfpathlineto{\pgfqpoint{1.036019in}{0.722231in}}%
\pgfpathlineto{\pgfqpoint{1.030729in}{0.720182in}}%
\pgfpathlineto{\pgfqpoint{1.026664in}{0.721652in}}%
\pgfpathlineto{\pgfqpoint{1.027301in}{0.717005in}}%
\pgfpathlineto{\pgfqpoint{1.031426in}{0.715395in}}%
\pgfpathlineto{\pgfqpoint{1.033829in}{0.717336in}}%
\pgfpathlineto{\pgfqpoint{1.039293in}{0.717228in}}%
\pgfpathlineto{\pgfqpoint{1.042163in}{0.713874in}}%
\pgfpathlineto{\pgfqpoint{1.038932in}{0.710194in}}%
\pgfpathlineto{\pgfqpoint{1.038327in}{0.708052in}}%
\pgfpathlineto{\pgfqpoint{1.040735in}{0.706277in}}%
\pgfpathlineto{\pgfqpoint{1.041143in}{0.703816in}}%
\pgfpathlineto{\pgfqpoint{1.044548in}{0.704736in}}%
\pgfpathlineto{\pgfqpoint{1.049310in}{0.702353in}}%
\pgfpathlineto{\pgfqpoint{1.052657in}{0.701365in}}%
\pgfpathlineto{\pgfqpoint{1.054785in}{0.696074in}}%
\pgfpathlineto{\pgfqpoint{1.056904in}{0.696328in}}%
\pgfpathlineto{\pgfqpoint{1.058734in}{0.693326in}}%
\pgfpathlineto{\pgfqpoint{1.052615in}{0.689071in}}%
\pgfpathlineto{\pgfqpoint{1.050791in}{0.686844in}}%
\pgfpathlineto{\pgfqpoint{1.057270in}{0.682397in}}%
\pgfpathlineto{\pgfqpoint{1.056506in}{0.680783in}}%
\pgfpathlineto{\pgfqpoint{1.053335in}{0.680298in}}%
\pgfpathlineto{\pgfqpoint{1.058156in}{0.676641in}}%
\pgfpathlineto{\pgfqpoint{1.058655in}{0.673622in}}%
\pgfpathlineto{\pgfqpoint{1.061740in}{0.674956in}}%
\pgfpathlineto{\pgfqpoint{1.066675in}{0.673514in}}%
\pgfpathlineto{\pgfqpoint{1.067308in}{0.671656in}}%
\pgfpathlineto{\pgfqpoint{1.064920in}{0.669776in}}%
\pgfpathlineto{\pgfqpoint{1.069974in}{0.668359in}}%
\pgfpathlineto{\pgfqpoint{1.073388in}{0.668942in}}%
\pgfpathlineto{\pgfqpoint{1.077468in}{0.664660in}}%
\pgfpathlineto{\pgfqpoint{1.078685in}{0.664612in}}%
\pgfpathlineto{\pgfqpoint{1.079652in}{0.659570in}}%
\pgfpathlineto{\pgfqpoint{1.082792in}{0.657912in}}%
\pgfpathlineto{\pgfqpoint{1.082915in}{0.655259in}}%
\pgfpathlineto{\pgfqpoint{1.085209in}{0.653941in}}%
\pgfpathlineto{\pgfqpoint{1.085707in}{0.649109in}}%
\pgfpathlineto{\pgfqpoint{1.092374in}{0.641077in}}%
\pgfpathlineto{\pgfqpoint{1.095064in}{0.640992in}}%
\pgfpathlineto{\pgfqpoint{1.098844in}{0.638529in}}%
\pgfpathlineto{\pgfqpoint{1.101206in}{0.635275in}}%
\pgfpathlineto{\pgfqpoint{1.104111in}{0.633554in}}%
\pgfpathlineto{\pgfqpoint{1.106568in}{0.626099in}}%
\pgfpathlineto{\pgfqpoint{1.108579in}{0.623361in}}%
\pgfpathlineto{\pgfqpoint{1.113936in}{0.621999in}}%
\pgfpathlineto{\pgfqpoint{1.120984in}{0.619591in}}%
\pgfpathlineto{\pgfqpoint{1.123831in}{0.619357in}}%
\pgfpathlineto{\pgfqpoint{1.127068in}{0.617346in}}%
\pgfpathlineto{\pgfqpoint{1.130555in}{0.612352in}}%
\pgfpathlineto{\pgfqpoint{1.132716in}{0.606694in}}%
\pgfpathlineto{\pgfqpoint{1.132681in}{0.605047in}}%
\pgfpathlineto{\pgfqpoint{1.134831in}{0.600830in}}%
\pgfpathlineto{\pgfqpoint{1.137136in}{0.597981in}}%
\pgfpathlineto{\pgfqpoint{1.137981in}{0.595190in}}%
\pgfpathlineto{\pgfqpoint{1.141010in}{0.591617in}}%
\pgfpathlineto{\pgfqpoint{1.139745in}{0.587853in}}%
\pgfpathlineto{\pgfqpoint{1.124077in}{0.559590in}}%
\pgfpathclose%
\pgfusepath{fill}%
\end{pgfscope}%
\begin{pgfscope}%
\pgfpathrectangle{\pgfqpoint{0.100000in}{0.100000in}}{\pgfqpoint{3.608454in}{2.310000in}}%
\pgfusepath{clip}%
\pgfsetbuttcap%
\pgfsetmiterjoin%
\definecolor{currentfill}{rgb}{0.000000,0.776471,0.611765}%
\pgfsetfillcolor{currentfill}%
\pgfsetlinewidth{0.000000pt}%
\definecolor{currentstroke}{rgb}{0.000000,0.000000,0.000000}%
\pgfsetstrokecolor{currentstroke}%
\pgfsetstrokeopacity{0.000000}%
\pgfsetdash{}{0pt}%
\pgfpathmoveto{\pgfqpoint{2.523585in}{1.228784in}}%
\pgfpathlineto{\pgfqpoint{2.519923in}{1.230014in}}%
\pgfpathlineto{\pgfqpoint{2.519105in}{1.243921in}}%
\pgfpathlineto{\pgfqpoint{2.500587in}{1.242556in}}%
\pgfpathlineto{\pgfqpoint{2.499015in}{1.266823in}}%
\pgfpathlineto{\pgfqpoint{2.519677in}{1.268419in}}%
\pgfpathlineto{\pgfqpoint{2.534479in}{1.269453in}}%
\pgfpathlineto{\pgfqpoint{2.537345in}{1.268351in}}%
\pgfpathlineto{\pgfqpoint{2.541890in}{1.261426in}}%
\pgfpathlineto{\pgfqpoint{2.534558in}{1.250105in}}%
\pgfpathlineto{\pgfqpoint{2.540147in}{1.235271in}}%
\pgfpathlineto{\pgfqpoint{2.528115in}{1.232902in}}%
\pgfpathlineto{\pgfqpoint{2.523585in}{1.228784in}}%
\pgfpathclose%
\pgfusepath{fill}%
\end{pgfscope}%
\begin{pgfscope}%
\pgfpathrectangle{\pgfqpoint{0.100000in}{0.100000in}}{\pgfqpoint{3.608454in}{2.310000in}}%
\pgfusepath{clip}%
\pgfsetbuttcap%
\pgfsetmiterjoin%
\definecolor{currentfill}{rgb}{0.000000,0.643137,0.678431}%
\pgfsetfillcolor{currentfill}%
\pgfsetlinewidth{0.000000pt}%
\definecolor{currentstroke}{rgb}{0.000000,0.000000,0.000000}%
\pgfsetstrokecolor{currentstroke}%
\pgfsetstrokeopacity{0.000000}%
\pgfsetdash{}{0pt}%
\pgfpathmoveto{\pgfqpoint{2.468504in}{0.944341in}}%
\pgfpathlineto{\pgfqpoint{2.451281in}{0.943313in}}%
\pgfpathlineto{\pgfqpoint{2.440822in}{0.945153in}}%
\pgfpathlineto{\pgfqpoint{2.441137in}{0.939280in}}%
\pgfpathlineto{\pgfqpoint{2.427361in}{0.937371in}}%
\pgfpathlineto{\pgfqpoint{2.427713in}{0.930421in}}%
\pgfpathlineto{\pgfqpoint{2.416509in}{0.931498in}}%
\pgfpathlineto{\pgfqpoint{2.424326in}{0.937210in}}%
\pgfpathlineto{\pgfqpoint{2.406598in}{0.936460in}}%
\pgfpathlineto{\pgfqpoint{2.405968in}{0.950277in}}%
\pgfpathlineto{\pgfqpoint{2.392184in}{0.949754in}}%
\pgfpathlineto{\pgfqpoint{2.391745in}{0.960160in}}%
\pgfpathlineto{\pgfqpoint{2.372289in}{0.959362in}}%
\pgfpathlineto{\pgfqpoint{2.384214in}{0.972642in}}%
\pgfpathlineto{\pgfqpoint{2.384480in}{0.979541in}}%
\pgfpathlineto{\pgfqpoint{2.395789in}{0.983518in}}%
\pgfpathlineto{\pgfqpoint{2.396088in}{0.992607in}}%
\pgfpathlineto{\pgfqpoint{2.407574in}{0.985155in}}%
\pgfpathlineto{\pgfqpoint{2.413949in}{0.985540in}}%
\pgfpathlineto{\pgfqpoint{2.413440in}{0.992269in}}%
\pgfpathlineto{\pgfqpoint{2.420148in}{0.996176in}}%
\pgfpathlineto{\pgfqpoint{2.467215in}{0.998833in}}%
\pgfpathlineto{\pgfqpoint{2.482530in}{0.995135in}}%
\pgfpathlineto{\pgfqpoint{2.483100in}{0.985946in}}%
\pgfpathlineto{\pgfqpoint{2.484251in}{0.968666in}}%
\pgfpathlineto{\pgfqpoint{2.466935in}{0.967664in}}%
\pgfpathlineto{\pgfqpoint{2.468504in}{0.944341in}}%
\pgfpathclose%
\pgfusepath{fill}%
\end{pgfscope}%
\begin{pgfscope}%
\pgfpathrectangle{\pgfqpoint{0.100000in}{0.100000in}}{\pgfqpoint{3.608454in}{2.310000in}}%
\pgfusepath{clip}%
\pgfsetbuttcap%
\pgfsetmiterjoin%
\definecolor{currentfill}{rgb}{0.000000,0.427451,0.786275}%
\pgfsetfillcolor{currentfill}%
\pgfsetlinewidth{0.000000pt}%
\definecolor{currentstroke}{rgb}{0.000000,0.000000,0.000000}%
\pgfsetstrokecolor{currentstroke}%
\pgfsetstrokeopacity{0.000000}%
\pgfsetdash{}{0pt}%
\pgfpathmoveto{\pgfqpoint{3.506058in}{1.889401in}}%
\pgfpathlineto{\pgfqpoint{3.496870in}{1.892724in}}%
\pgfpathlineto{\pgfqpoint{3.498018in}{1.895957in}}%
\pgfpathlineto{\pgfqpoint{3.491360in}{1.900167in}}%
\pgfpathlineto{\pgfqpoint{3.485597in}{1.897568in}}%
\pgfpathlineto{\pgfqpoint{3.484176in}{1.904080in}}%
\pgfpathlineto{\pgfqpoint{3.474707in}{1.903468in}}%
\pgfpathlineto{\pgfqpoint{3.465497in}{1.899405in}}%
\pgfpathlineto{\pgfqpoint{3.454146in}{1.936802in}}%
\pgfpathlineto{\pgfqpoint{3.442334in}{1.974192in}}%
\pgfpathlineto{\pgfqpoint{3.428979in}{2.013989in}}%
\pgfpathlineto{\pgfqpoint{3.435281in}{2.018401in}}%
\pgfpathlineto{\pgfqpoint{3.440288in}{2.012007in}}%
\pgfpathlineto{\pgfqpoint{3.443602in}{2.017361in}}%
\pgfpathlineto{\pgfqpoint{3.442691in}{2.028080in}}%
\pgfpathlineto{\pgfqpoint{3.447268in}{2.025988in}}%
\pgfpathlineto{\pgfqpoint{3.451169in}{2.029587in}}%
\pgfpathlineto{\pgfqpoint{3.444152in}{2.035398in}}%
\pgfpathlineto{\pgfqpoint{3.446349in}{2.043593in}}%
\pgfpathlineto{\pgfqpoint{3.457409in}{2.056674in}}%
\pgfpathlineto{\pgfqpoint{3.455828in}{2.064013in}}%
\pgfpathlineto{\pgfqpoint{3.460024in}{2.069258in}}%
\pgfpathlineto{\pgfqpoint{3.460838in}{2.075413in}}%
\pgfpathlineto{\pgfqpoint{3.455091in}{2.080073in}}%
\pgfpathlineto{\pgfqpoint{3.456695in}{2.090917in}}%
\pgfpathlineto{\pgfqpoint{3.452786in}{2.093614in}}%
\pgfpathlineto{\pgfqpoint{3.454249in}{2.105442in}}%
\pgfpathlineto{\pgfqpoint{3.460446in}{2.114346in}}%
\pgfpathlineto{\pgfqpoint{3.458788in}{2.127000in}}%
\pgfpathlineto{\pgfqpoint{3.474712in}{2.131510in}}%
\pgfpathlineto{\pgfqpoint{3.478067in}{2.117502in}}%
\pgfpathlineto{\pgfqpoint{3.487328in}{2.085547in}}%
\pgfpathlineto{\pgfqpoint{3.491275in}{2.076574in}}%
\pgfpathlineto{\pgfqpoint{3.490714in}{2.062918in}}%
\pgfpathlineto{\pgfqpoint{3.496496in}{2.059377in}}%
\pgfpathlineto{\pgfqpoint{3.494071in}{2.051630in}}%
\pgfpathlineto{\pgfqpoint{3.514278in}{2.013412in}}%
\pgfpathlineto{\pgfqpoint{3.527215in}{2.022326in}}%
\pgfpathlineto{\pgfqpoint{3.539818in}{1.996897in}}%
\pgfpathlineto{\pgfqpoint{3.532167in}{1.996604in}}%
\pgfpathlineto{\pgfqpoint{3.534775in}{1.988683in}}%
\pgfpathlineto{\pgfqpoint{3.535820in}{1.974394in}}%
\pgfpathlineto{\pgfqpoint{3.535140in}{1.964201in}}%
\pgfpathlineto{\pgfqpoint{3.538805in}{1.962876in}}%
\pgfpathlineto{\pgfqpoint{3.543555in}{1.955713in}}%
\pgfpathlineto{\pgfqpoint{3.547311in}{1.956692in}}%
\pgfpathlineto{\pgfqpoint{3.553924in}{1.946267in}}%
\pgfpathlineto{\pgfqpoint{3.550714in}{1.937007in}}%
\pgfpathlineto{\pgfqpoint{3.547033in}{1.936458in}}%
\pgfpathlineto{\pgfqpoint{3.546482in}{1.926493in}}%
\pgfpathlineto{\pgfqpoint{3.535703in}{1.920334in}}%
\pgfpathlineto{\pgfqpoint{3.534745in}{1.914657in}}%
\pgfpathlineto{\pgfqpoint{3.529959in}{1.911075in}}%
\pgfpathlineto{\pgfqpoint{3.516703in}{1.917002in}}%
\pgfpathlineto{\pgfqpoint{3.510197in}{1.910231in}}%
\pgfpathlineto{\pgfqpoint{3.508494in}{1.902315in}}%
\pgfpathlineto{\pgfqpoint{3.514068in}{1.894022in}}%
\pgfpathlineto{\pgfqpoint{3.506058in}{1.889401in}}%
\pgfpathclose%
\pgfusepath{fill}%
\end{pgfscope}%
\begin{pgfscope}%
\pgfpathrectangle{\pgfqpoint{0.100000in}{0.100000in}}{\pgfqpoint{3.608454in}{2.310000in}}%
\pgfusepath{clip}%
\pgfsetbuttcap%
\pgfsetmiterjoin%
\definecolor{currentfill}{rgb}{0.000000,0.615686,0.692157}%
\pgfsetfillcolor{currentfill}%
\pgfsetlinewidth{0.000000pt}%
\definecolor{currentstroke}{rgb}{0.000000,0.000000,0.000000}%
\pgfsetstrokecolor{currentstroke}%
\pgfsetstrokeopacity{0.000000}%
\pgfsetdash{}{0pt}%
\pgfpathmoveto{\pgfqpoint{2.519263in}{0.872679in}}%
\pgfpathlineto{\pgfqpoint{2.485523in}{0.870923in}}%
\pgfpathlineto{\pgfqpoint{2.483949in}{0.899006in}}%
\pgfpathlineto{\pgfqpoint{2.474841in}{0.898468in}}%
\pgfpathlineto{\pgfqpoint{2.474111in}{0.910147in}}%
\pgfpathlineto{\pgfqpoint{2.470515in}{0.912039in}}%
\pgfpathlineto{\pgfqpoint{2.469213in}{0.932802in}}%
\pgfpathlineto{\pgfqpoint{2.489985in}{0.934122in}}%
\pgfpathlineto{\pgfqpoint{2.506964in}{0.937791in}}%
\pgfpathlineto{\pgfqpoint{2.506684in}{0.942391in}}%
\pgfpathlineto{\pgfqpoint{2.520583in}{0.943229in}}%
\pgfpathlineto{\pgfqpoint{2.521327in}{0.932876in}}%
\pgfpathlineto{\pgfqpoint{2.535053in}{0.931443in}}%
\pgfpathlineto{\pgfqpoint{2.537884in}{0.927921in}}%
\pgfpathlineto{\pgfqpoint{2.534527in}{0.922423in}}%
\pgfpathlineto{\pgfqpoint{2.525476in}{0.919309in}}%
\pgfpathlineto{\pgfqpoint{2.526832in}{0.901949in}}%
\pgfpathlineto{\pgfqpoint{2.517297in}{0.901251in}}%
\pgfpathlineto{\pgfqpoint{2.519263in}{0.872679in}}%
\pgfpathclose%
\pgfusepath{fill}%
\end{pgfscope}%
\begin{pgfscope}%
\pgfpathrectangle{\pgfqpoint{0.100000in}{0.100000in}}{\pgfqpoint{3.608454in}{2.310000in}}%
\pgfusepath{clip}%
\pgfsetbuttcap%
\pgfsetmiterjoin%
\definecolor{currentfill}{rgb}{0.000000,0.403922,0.798039}%
\pgfsetfillcolor{currentfill}%
\pgfsetlinewidth{0.000000pt}%
\definecolor{currentstroke}{rgb}{0.000000,0.000000,0.000000}%
\pgfsetstrokecolor{currentstroke}%
\pgfsetstrokeopacity{0.000000}%
\pgfsetdash{}{0pt}%
\pgfpathmoveto{\pgfqpoint{0.623427in}{2.025636in}}%
\pgfpathlineto{\pgfqpoint{0.587828in}{2.036379in}}%
\pgfpathlineto{\pgfqpoint{0.581481in}{2.042623in}}%
\pgfpathlineto{\pgfqpoint{0.580214in}{2.049876in}}%
\pgfpathlineto{\pgfqpoint{0.575268in}{2.054730in}}%
\pgfpathlineto{\pgfqpoint{0.572541in}{2.061484in}}%
\pgfpathlineto{\pgfqpoint{0.575705in}{2.064800in}}%
\pgfpathlineto{\pgfqpoint{0.576179in}{2.072349in}}%
\pgfpathlineto{\pgfqpoint{0.568024in}{2.080965in}}%
\pgfpathlineto{\pgfqpoint{0.566154in}{2.086439in}}%
\pgfpathlineto{\pgfqpoint{0.561348in}{2.090351in}}%
\pgfpathlineto{\pgfqpoint{0.543848in}{2.096029in}}%
\pgfpathlineto{\pgfqpoint{0.550017in}{2.103751in}}%
\pgfpathlineto{\pgfqpoint{0.551015in}{2.110628in}}%
\pgfpathlineto{\pgfqpoint{0.556089in}{2.115055in}}%
\pgfpathlineto{\pgfqpoint{0.557819in}{2.120432in}}%
\pgfpathlineto{\pgfqpoint{0.566644in}{2.148266in}}%
\pgfpathlineto{\pgfqpoint{0.578043in}{2.148170in}}%
\pgfpathlineto{\pgfqpoint{0.589293in}{2.135786in}}%
\pgfpathlineto{\pgfqpoint{0.593066in}{2.119148in}}%
\pgfpathlineto{\pgfqpoint{0.609486in}{2.121452in}}%
\pgfpathlineto{\pgfqpoint{0.615173in}{2.117075in}}%
\pgfpathlineto{\pgfqpoint{0.623310in}{2.122649in}}%
\pgfpathlineto{\pgfqpoint{0.631314in}{2.147810in}}%
\pgfpathlineto{\pgfqpoint{0.669147in}{2.136525in}}%
\pgfpathlineto{\pgfqpoint{0.661241in}{2.110265in}}%
\pgfpathlineto{\pgfqpoint{0.656547in}{2.111660in}}%
\pgfpathlineto{\pgfqpoint{0.650732in}{2.091686in}}%
\pgfpathlineto{\pgfqpoint{0.653479in}{2.086201in}}%
\pgfpathlineto{\pgfqpoint{0.641984in}{2.086967in}}%
\pgfpathlineto{\pgfqpoint{0.634943in}{2.089043in}}%
\pgfpathlineto{\pgfqpoint{0.630988in}{2.086576in}}%
\pgfpathlineto{\pgfqpoint{0.628826in}{2.076535in}}%
\pgfpathlineto{\pgfqpoint{0.636428in}{2.062774in}}%
\pgfpathlineto{\pgfqpoint{0.634662in}{2.052316in}}%
\pgfpathlineto{\pgfqpoint{0.629691in}{2.048489in}}%
\pgfpathlineto{\pgfqpoint{0.630035in}{2.038493in}}%
\pgfpathlineto{\pgfqpoint{0.622394in}{2.035924in}}%
\pgfpathlineto{\pgfqpoint{0.623427in}{2.025636in}}%
\pgfpathclose%
\pgfusepath{fill}%
\end{pgfscope}%
\begin{pgfscope}%
\pgfpathrectangle{\pgfqpoint{0.100000in}{0.100000in}}{\pgfqpoint{3.608454in}{2.310000in}}%
\pgfusepath{clip}%
\pgfsetbuttcap%
\pgfsetmiterjoin%
\definecolor{currentfill}{rgb}{0.000000,0.466667,0.766667}%
\pgfsetfillcolor{currentfill}%
\pgfsetlinewidth{0.000000pt}%
\definecolor{currentstroke}{rgb}{0.000000,0.000000,0.000000}%
\pgfsetstrokecolor{currentstroke}%
\pgfsetstrokeopacity{0.000000}%
\pgfsetdash{}{0pt}%
\pgfpathmoveto{\pgfqpoint{3.027503in}{0.502091in}}%
\pgfpathlineto{\pgfqpoint{2.992126in}{0.496993in}}%
\pgfpathlineto{\pgfqpoint{2.998004in}{0.510943in}}%
\pgfpathlineto{\pgfqpoint{2.992287in}{0.512337in}}%
\pgfpathlineto{\pgfqpoint{2.985816in}{0.507630in}}%
\pgfpathlineto{\pgfqpoint{2.984712in}{0.500436in}}%
\pgfpathlineto{\pgfqpoint{2.980383in}{0.498839in}}%
\pgfpathlineto{\pgfqpoint{2.968704in}{0.512418in}}%
\pgfpathlineto{\pgfqpoint{2.971482in}{0.527436in}}%
\pgfpathlineto{\pgfqpoint{2.969072in}{0.537771in}}%
\pgfpathlineto{\pgfqpoint{2.971652in}{0.541641in}}%
\pgfpathlineto{\pgfqpoint{2.975315in}{0.570540in}}%
\pgfpathlineto{\pgfqpoint{2.974786in}{0.579836in}}%
\pgfpathlineto{\pgfqpoint{2.990145in}{0.581978in}}%
\pgfpathlineto{\pgfqpoint{3.001283in}{0.581332in}}%
\pgfpathlineto{\pgfqpoint{3.006224in}{0.574274in}}%
\pgfpathlineto{\pgfqpoint{3.017489in}{0.571785in}}%
\pgfpathlineto{\pgfqpoint{3.019784in}{0.555093in}}%
\pgfpathlineto{\pgfqpoint{3.016866in}{0.550384in}}%
\pgfpathlineto{\pgfqpoint{3.021367in}{0.543921in}}%
\pgfpathlineto{\pgfqpoint{3.027503in}{0.502091in}}%
\pgfpathclose%
\pgfusepath{fill}%
\end{pgfscope}%
\begin{pgfscope}%
\pgfpathrectangle{\pgfqpoint{0.100000in}{0.100000in}}{\pgfqpoint{3.608454in}{2.310000in}}%
\pgfusepath{clip}%
\pgfsetbuttcap%
\pgfsetmiterjoin%
\definecolor{currentfill}{rgb}{0.000000,0.290196,0.854902}%
\pgfsetfillcolor{currentfill}%
\pgfsetlinewidth{0.000000pt}%
\definecolor{currentstroke}{rgb}{0.000000,0.000000,0.000000}%
\pgfsetstrokecolor{currentstroke}%
\pgfsetstrokeopacity{0.000000}%
\pgfsetdash{}{0pt}%
\pgfpathmoveto{\pgfqpoint{3.420582in}{1.833115in}}%
\pgfpathlineto{\pgfqpoint{3.421042in}{1.838069in}}%
\pgfpathlineto{\pgfqpoint{3.413847in}{1.847236in}}%
\pgfpathlineto{\pgfqpoint{3.415970in}{1.857152in}}%
\pgfpathlineto{\pgfqpoint{3.421501in}{1.872886in}}%
\pgfpathlineto{\pgfqpoint{3.424616in}{1.868740in}}%
\pgfpathlineto{\pgfqpoint{3.432335in}{1.876511in}}%
\pgfpathlineto{\pgfqpoint{3.431173in}{1.881530in}}%
\pgfpathlineto{\pgfqpoint{3.437941in}{1.884217in}}%
\pgfpathlineto{\pgfqpoint{3.431626in}{1.898622in}}%
\pgfpathlineto{\pgfqpoint{3.442370in}{1.903310in}}%
\pgfpathlineto{\pgfqpoint{3.438972in}{1.915003in}}%
\pgfpathlineto{\pgfqpoint{3.433355in}{1.925260in}}%
\pgfpathlineto{\pgfqpoint{3.440922in}{1.921126in}}%
\pgfpathlineto{\pgfqpoint{3.441678in}{1.930119in}}%
\pgfpathlineto{\pgfqpoint{3.453336in}{1.932723in}}%
\pgfpathlineto{\pgfqpoint{3.454146in}{1.936802in}}%
\pgfpathlineto{\pgfqpoint{3.465497in}{1.899405in}}%
\pgfpathlineto{\pgfqpoint{3.474707in}{1.903468in}}%
\pgfpathlineto{\pgfqpoint{3.484176in}{1.904080in}}%
\pgfpathlineto{\pgfqpoint{3.485597in}{1.897568in}}%
\pgfpathlineto{\pgfqpoint{3.491360in}{1.900167in}}%
\pgfpathlineto{\pgfqpoint{3.498018in}{1.895957in}}%
\pgfpathlineto{\pgfqpoint{3.496870in}{1.892724in}}%
\pgfpathlineto{\pgfqpoint{3.506058in}{1.889401in}}%
\pgfpathlineto{\pgfqpoint{3.507318in}{1.881477in}}%
\pgfpathlineto{\pgfqpoint{3.505212in}{1.874640in}}%
\pgfpathlineto{\pgfqpoint{3.499936in}{1.871265in}}%
\pgfpathlineto{\pgfqpoint{3.499231in}{1.853431in}}%
\pgfpathlineto{\pgfqpoint{3.493789in}{1.834815in}}%
\pgfpathlineto{\pgfqpoint{3.494171in}{1.831598in}}%
\pgfpathlineto{\pgfqpoint{3.487553in}{1.830826in}}%
\pgfpathlineto{\pgfqpoint{3.477550in}{1.822939in}}%
\pgfpathlineto{\pgfqpoint{3.471016in}{1.810988in}}%
\pgfpathlineto{\pgfqpoint{3.434992in}{1.802755in}}%
\pgfpathlineto{\pgfqpoint{3.430894in}{1.808006in}}%
\pgfpathlineto{\pgfqpoint{3.426046in}{1.819395in}}%
\pgfpathlineto{\pgfqpoint{3.422865in}{1.818650in}}%
\pgfpathlineto{\pgfqpoint{3.420582in}{1.833115in}}%
\pgfpathclose%
\pgfusepath{fill}%
\end{pgfscope}%
\begin{pgfscope}%
\pgfpathrectangle{\pgfqpoint{0.100000in}{0.100000in}}{\pgfqpoint{3.608454in}{2.310000in}}%
\pgfusepath{clip}%
\pgfsetbuttcap%
\pgfsetmiterjoin%
\definecolor{currentfill}{rgb}{0.000000,0.698039,0.650980}%
\pgfsetfillcolor{currentfill}%
\pgfsetlinewidth{0.000000pt}%
\definecolor{currentstroke}{rgb}{0.000000,0.000000,0.000000}%
\pgfsetstrokecolor{currentstroke}%
\pgfsetstrokeopacity{0.000000}%
\pgfsetdash{}{0pt}%
\pgfpathmoveto{\pgfqpoint{2.478809in}{1.176580in}}%
\pgfpathlineto{\pgfqpoint{2.477580in}{1.181419in}}%
\pgfpathlineto{\pgfqpoint{2.479456in}{1.190579in}}%
\pgfpathlineto{\pgfqpoint{2.477226in}{1.192084in}}%
\pgfpathlineto{\pgfqpoint{2.473800in}{1.196498in}}%
\pgfpathlineto{\pgfqpoint{2.482512in}{1.210574in}}%
\pgfpathlineto{\pgfqpoint{2.488838in}{1.212212in}}%
\pgfpathlineto{\pgfqpoint{2.489480in}{1.220900in}}%
\pgfpathlineto{\pgfqpoint{2.501888in}{1.221865in}}%
\pgfpathlineto{\pgfqpoint{2.500587in}{1.242556in}}%
\pgfpathlineto{\pgfqpoint{2.519105in}{1.243921in}}%
\pgfpathlineto{\pgfqpoint{2.519923in}{1.230014in}}%
\pgfpathlineto{\pgfqpoint{2.523585in}{1.228784in}}%
\pgfpathlineto{\pgfqpoint{2.533129in}{1.219044in}}%
\pgfpathlineto{\pgfqpoint{2.535791in}{1.209919in}}%
\pgfpathlineto{\pgfqpoint{2.533377in}{1.204248in}}%
\pgfpathlineto{\pgfqpoint{2.539380in}{1.187318in}}%
\pgfpathlineto{\pgfqpoint{2.543041in}{1.178054in}}%
\pgfpathlineto{\pgfqpoint{2.519321in}{1.176572in}}%
\pgfpathlineto{\pgfqpoint{2.518313in}{1.191718in}}%
\pgfpathlineto{\pgfqpoint{2.497579in}{1.190610in}}%
\pgfpathlineto{\pgfqpoint{2.498496in}{1.176925in}}%
\pgfpathlineto{\pgfqpoint{2.478809in}{1.176580in}}%
\pgfpathclose%
\pgfusepath{fill}%
\end{pgfscope}%
\begin{pgfscope}%
\pgfpathrectangle{\pgfqpoint{0.100000in}{0.100000in}}{\pgfqpoint{3.608454in}{2.310000in}}%
\pgfusepath{clip}%
\pgfsetbuttcap%
\pgfsetmiterjoin%
\definecolor{currentfill}{rgb}{0.000000,0.505882,0.747059}%
\pgfsetfillcolor{currentfill}%
\pgfsetlinewidth{0.000000pt}%
\definecolor{currentstroke}{rgb}{0.000000,0.000000,0.000000}%
\pgfsetstrokecolor{currentstroke}%
\pgfsetstrokeopacity{0.000000}%
\pgfsetdash{}{0pt}%
\pgfpathmoveto{\pgfqpoint{3.165006in}{1.605498in}}%
\pgfpathlineto{\pgfqpoint{3.154199in}{1.598846in}}%
\pgfpathlineto{\pgfqpoint{3.144862in}{1.595009in}}%
\pgfpathlineto{\pgfqpoint{3.133554in}{1.604605in}}%
\pgfpathlineto{\pgfqpoint{3.128786in}{1.603472in}}%
\pgfpathlineto{\pgfqpoint{3.125203in}{1.608488in}}%
\pgfpathlineto{\pgfqpoint{3.111013in}{1.603891in}}%
\pgfpathlineto{\pgfqpoint{3.106858in}{1.608202in}}%
\pgfpathlineto{\pgfqpoint{3.110757in}{1.621055in}}%
\pgfpathlineto{\pgfqpoint{3.109153in}{1.629395in}}%
\pgfpathlineto{\pgfqpoint{3.131860in}{1.634049in}}%
\pgfpathlineto{\pgfqpoint{3.130816in}{1.638979in}}%
\pgfpathlineto{\pgfqpoint{3.167806in}{1.647018in}}%
\pgfpathlineto{\pgfqpoint{3.172073in}{1.651526in}}%
\pgfpathlineto{\pgfqpoint{3.162699in}{1.682336in}}%
\pgfpathlineto{\pgfqpoint{3.183986in}{1.686549in}}%
\pgfpathlineto{\pgfqpoint{3.207689in}{1.691369in}}%
\pgfpathlineto{\pgfqpoint{3.215154in}{1.664914in}}%
\pgfpathlineto{\pgfqpoint{3.210379in}{1.663470in}}%
\pgfpathlineto{\pgfqpoint{3.210875in}{1.652502in}}%
\pgfpathlineto{\pgfqpoint{3.209405in}{1.635945in}}%
\pgfpathlineto{\pgfqpoint{3.203745in}{1.634729in}}%
\pgfpathlineto{\pgfqpoint{3.195407in}{1.620855in}}%
\pgfpathlineto{\pgfqpoint{3.192630in}{1.620121in}}%
\pgfpathlineto{\pgfqpoint{3.183685in}{1.619914in}}%
\pgfpathlineto{\pgfqpoint{3.174106in}{1.615801in}}%
\pgfpathlineto{\pgfqpoint{3.173979in}{1.610853in}}%
\pgfpathlineto{\pgfqpoint{3.165006in}{1.605498in}}%
\pgfpathclose%
\pgfusepath{fill}%
\end{pgfscope}%
\begin{pgfscope}%
\pgfpathrectangle{\pgfqpoint{0.100000in}{0.100000in}}{\pgfqpoint{3.608454in}{2.310000in}}%
\pgfusepath{clip}%
\pgfsetbuttcap%
\pgfsetmiterjoin%
\definecolor{currentfill}{rgb}{0.000000,0.823529,0.588235}%
\pgfsetfillcolor{currentfill}%
\pgfsetlinewidth{0.000000pt}%
\definecolor{currentstroke}{rgb}{0.000000,0.000000,0.000000}%
\pgfsetstrokecolor{currentstroke}%
\pgfsetstrokeopacity{0.000000}%
\pgfsetdash{}{0pt}%
\pgfpathmoveto{\pgfqpoint{2.674593in}{0.730790in}}%
\pgfpathlineto{\pgfqpoint{2.612513in}{0.725646in}}%
\pgfpathlineto{\pgfqpoint{2.609757in}{0.745194in}}%
\pgfpathlineto{\pgfqpoint{2.602701in}{0.749237in}}%
\pgfpathlineto{\pgfqpoint{2.601460in}{0.754867in}}%
\pgfpathlineto{\pgfqpoint{2.605232in}{0.759125in}}%
\pgfpathlineto{\pgfqpoint{2.611445in}{0.760754in}}%
\pgfpathlineto{\pgfqpoint{2.610143in}{0.781433in}}%
\pgfpathlineto{\pgfqpoint{2.613565in}{0.781771in}}%
\pgfpathlineto{\pgfqpoint{2.613728in}{0.792291in}}%
\pgfpathlineto{\pgfqpoint{2.653773in}{0.795845in}}%
\pgfpathlineto{\pgfqpoint{2.655184in}{0.780136in}}%
\pgfpathlineto{\pgfqpoint{2.660506in}{0.772013in}}%
\pgfpathlineto{\pgfqpoint{2.669844in}{0.772749in}}%
\pgfpathlineto{\pgfqpoint{2.673739in}{0.731186in}}%
\pgfpathlineto{\pgfqpoint{2.674593in}{0.730790in}}%
\pgfpathclose%
\pgfusepath{fill}%
\end{pgfscope}%
\begin{pgfscope}%
\pgfpathrectangle{\pgfqpoint{0.100000in}{0.100000in}}{\pgfqpoint{3.608454in}{2.310000in}}%
\pgfusepath{clip}%
\pgfsetbuttcap%
\pgfsetmiterjoin%
\definecolor{currentfill}{rgb}{0.000000,0.674510,0.662745}%
\pgfsetfillcolor{currentfill}%
\pgfsetlinewidth{0.000000pt}%
\definecolor{currentstroke}{rgb}{0.000000,0.000000,0.000000}%
\pgfsetstrokecolor{currentstroke}%
\pgfsetstrokeopacity{0.000000}%
\pgfsetdash{}{0pt}%
\pgfpathmoveto{\pgfqpoint{1.245517in}{0.942241in}}%
\pgfpathlineto{\pgfqpoint{1.176800in}{0.953138in}}%
\pgfpathlineto{\pgfqpoint{1.183780in}{0.997802in}}%
\pgfpathlineto{\pgfqpoint{1.193614in}{1.060581in}}%
\pgfpathlineto{\pgfqpoint{1.236876in}{1.053946in}}%
\pgfpathlineto{\pgfqpoint{1.278910in}{1.047826in}}%
\pgfpathlineto{\pgfqpoint{1.333401in}{1.040515in}}%
\pgfpathlineto{\pgfqpoint{1.340304in}{1.035758in}}%
\pgfpathlineto{\pgfqpoint{1.362154in}{1.025644in}}%
\pgfpathlineto{\pgfqpoint{1.360437in}{1.011466in}}%
\pgfpathlineto{\pgfqpoint{1.392442in}{1.007538in}}%
\pgfpathlineto{\pgfqpoint{1.388260in}{0.973319in}}%
\pgfpathlineto{\pgfqpoint{1.379878in}{0.974364in}}%
\pgfpathlineto{\pgfqpoint{1.378188in}{0.960531in}}%
\pgfpathlineto{\pgfqpoint{1.357228in}{0.963939in}}%
\pgfpathlineto{\pgfqpoint{1.355496in}{0.949729in}}%
\pgfpathlineto{\pgfqpoint{1.318032in}{0.954222in}}%
\pgfpathlineto{\pgfqpoint{1.252648in}{0.963109in}}%
\pgfpathlineto{\pgfqpoint{1.248648in}{0.963976in}}%
\pgfpathlineto{\pgfqpoint{1.245517in}{0.942241in}}%
\pgfpathclose%
\pgfusepath{fill}%
\end{pgfscope}%
\begin{pgfscope}%
\pgfpathrectangle{\pgfqpoint{0.100000in}{0.100000in}}{\pgfqpoint{3.608454in}{2.310000in}}%
\pgfusepath{clip}%
\pgfsetbuttcap%
\pgfsetmiterjoin%
\definecolor{currentfill}{rgb}{0.000000,0.670588,0.664706}%
\pgfsetfillcolor{currentfill}%
\pgfsetlinewidth{0.000000pt}%
\definecolor{currentstroke}{rgb}{0.000000,0.000000,0.000000}%
\pgfsetstrokecolor{currentstroke}%
\pgfsetstrokeopacity{0.000000}%
\pgfsetdash{}{0pt}%
\pgfpathmoveto{\pgfqpoint{2.676133in}{1.724089in}}%
\pgfpathlineto{\pgfqpoint{2.648810in}{1.721346in}}%
\pgfpathlineto{\pgfqpoint{2.645951in}{1.748717in}}%
\pgfpathlineto{\pgfqpoint{2.618822in}{1.745974in}}%
\pgfpathlineto{\pgfqpoint{2.615872in}{1.773777in}}%
\pgfpathlineto{\pgfqpoint{2.628418in}{1.774736in}}%
\pgfpathlineto{\pgfqpoint{2.625838in}{1.802295in}}%
\pgfpathlineto{\pgfqpoint{2.680280in}{1.807943in}}%
\pgfpathlineto{\pgfqpoint{2.683319in}{1.780357in}}%
\pgfpathlineto{\pgfqpoint{2.669984in}{1.779116in}}%
\pgfpathlineto{\pgfqpoint{2.676133in}{1.724089in}}%
\pgfpathclose%
\pgfusepath{fill}%
\end{pgfscope}%
\begin{pgfscope}%
\pgfpathrectangle{\pgfqpoint{0.100000in}{0.100000in}}{\pgfqpoint{3.608454in}{2.310000in}}%
\pgfusepath{clip}%
\pgfsetbuttcap%
\pgfsetmiterjoin%
\definecolor{currentfill}{rgb}{0.000000,0.411765,0.794118}%
\pgfsetfillcolor{currentfill}%
\pgfsetlinewidth{0.000000pt}%
\definecolor{currentstroke}{rgb}{0.000000,0.000000,0.000000}%
\pgfsetstrokecolor{currentstroke}%
\pgfsetstrokeopacity{0.000000}%
\pgfsetdash{}{0pt}%
\pgfpathmoveto{\pgfqpoint{2.575892in}{1.080309in}}%
\pgfpathlineto{\pgfqpoint{2.579650in}{1.089158in}}%
\pgfpathlineto{\pgfqpoint{2.573068in}{1.094430in}}%
\pgfpathlineto{\pgfqpoint{2.573691in}{1.108364in}}%
\pgfpathlineto{\pgfqpoint{2.573891in}{1.113688in}}%
\pgfpathlineto{\pgfqpoint{2.584331in}{1.121514in}}%
\pgfpathlineto{\pgfqpoint{2.605651in}{1.120589in}}%
\pgfpathlineto{\pgfqpoint{2.605693in}{1.111968in}}%
\pgfpathlineto{\pgfqpoint{2.632392in}{1.103973in}}%
\pgfpathlineto{\pgfqpoint{2.643683in}{1.102257in}}%
\pgfpathlineto{\pgfqpoint{2.644297in}{1.082383in}}%
\pgfpathlineto{\pgfqpoint{2.648588in}{1.077186in}}%
\pgfpathlineto{\pgfqpoint{2.648663in}{1.074175in}}%
\pgfpathlineto{\pgfqpoint{2.639911in}{1.067474in}}%
\pgfpathlineto{\pgfqpoint{2.634590in}{1.067850in}}%
\pgfpathlineto{\pgfqpoint{2.636003in}{1.046348in}}%
\pgfpathlineto{\pgfqpoint{2.611745in}{1.044790in}}%
\pgfpathlineto{\pgfqpoint{2.610505in}{1.048649in}}%
\pgfpathlineto{\pgfqpoint{2.608880in}{1.074941in}}%
\pgfpathlineto{\pgfqpoint{2.607711in}{1.078660in}}%
\pgfpathlineto{\pgfqpoint{2.601253in}{1.080261in}}%
\pgfpathlineto{\pgfqpoint{2.585581in}{1.074349in}}%
\pgfpathlineto{\pgfqpoint{2.580566in}{1.079844in}}%
\pgfpathlineto{\pgfqpoint{2.575892in}{1.080309in}}%
\pgfpathclose%
\pgfusepath{fill}%
\end{pgfscope}%
\begin{pgfscope}%
\pgfpathrectangle{\pgfqpoint{0.100000in}{0.100000in}}{\pgfqpoint{3.608454in}{2.310000in}}%
\pgfusepath{clip}%
\pgfsetbuttcap%
\pgfsetmiterjoin%
\definecolor{currentfill}{rgb}{0.000000,0.949020,0.525490}%
\pgfsetfillcolor{currentfill}%
\pgfsetlinewidth{0.000000pt}%
\definecolor{currentstroke}{rgb}{0.000000,0.000000,0.000000}%
\pgfsetstrokecolor{currentstroke}%
\pgfsetstrokeopacity{0.000000}%
\pgfsetdash{}{0pt}%
\pgfpathmoveto{\pgfqpoint{3.010884in}{1.388904in}}%
\pgfpathlineto{\pgfqpoint{3.018910in}{1.376644in}}%
\pgfpathlineto{\pgfqpoint{3.016902in}{1.365346in}}%
\pgfpathlineto{\pgfqpoint{3.012394in}{1.360390in}}%
\pgfpathlineto{\pgfqpoint{3.011961in}{1.351516in}}%
\pgfpathlineto{\pgfqpoint{3.006750in}{1.353994in}}%
\pgfpathlineto{\pgfqpoint{2.992200in}{1.336598in}}%
\pgfpathlineto{\pgfqpoint{2.981470in}{1.336849in}}%
\pgfpathlineto{\pgfqpoint{2.973652in}{1.345179in}}%
\pgfpathlineto{\pgfqpoint{2.968495in}{1.342247in}}%
\pgfpathlineto{\pgfqpoint{2.958041in}{1.346022in}}%
\pgfpathlineto{\pgfqpoint{2.975847in}{1.360884in}}%
\pgfpathlineto{\pgfqpoint{2.976752in}{1.368166in}}%
\pgfpathlineto{\pgfqpoint{2.972924in}{1.373592in}}%
\pgfpathlineto{\pgfqpoint{2.965422in}{1.379505in}}%
\pgfpathlineto{\pgfqpoint{2.967663in}{1.383858in}}%
\pgfpathlineto{\pgfqpoint{2.962379in}{1.389071in}}%
\pgfpathlineto{\pgfqpoint{2.965023in}{1.393727in}}%
\pgfpathlineto{\pgfqpoint{2.961168in}{1.406181in}}%
\pgfpathlineto{\pgfqpoint{2.967704in}{1.409901in}}%
\pgfpathlineto{\pgfqpoint{2.973331in}{1.416066in}}%
\pgfpathlineto{\pgfqpoint{2.978628in}{1.415804in}}%
\pgfpathlineto{\pgfqpoint{2.980404in}{1.409102in}}%
\pgfpathlineto{\pgfqpoint{2.989078in}{1.401477in}}%
\pgfpathlineto{\pgfqpoint{2.997346in}{1.397787in}}%
\pgfpathlineto{\pgfqpoint{3.001117in}{1.390072in}}%
\pgfpathlineto{\pgfqpoint{3.010884in}{1.388904in}}%
\pgfpathclose%
\pgfusepath{fill}%
\end{pgfscope}%
\begin{pgfscope}%
\pgfpathrectangle{\pgfqpoint{0.100000in}{0.100000in}}{\pgfqpoint{3.608454in}{2.310000in}}%
\pgfusepath{clip}%
\pgfsetbuttcap%
\pgfsetmiterjoin%
\definecolor{currentfill}{rgb}{0.000000,0.784314,0.607843}%
\pgfsetfillcolor{currentfill}%
\pgfsetlinewidth{0.000000pt}%
\definecolor{currentstroke}{rgb}{0.000000,0.000000,0.000000}%
\pgfsetstrokecolor{currentstroke}%
\pgfsetstrokeopacity{0.000000}%
\pgfsetdash{}{0pt}%
\pgfpathmoveto{\pgfqpoint{0.983512in}{1.292540in}}%
\pgfpathlineto{\pgfqpoint{0.912295in}{1.307222in}}%
\pgfpathlineto{\pgfqpoint{0.925741in}{1.369767in}}%
\pgfpathlineto{\pgfqpoint{0.938194in}{1.428823in}}%
\pgfpathlineto{\pgfqpoint{1.000345in}{1.416064in}}%
\pgfpathlineto{\pgfqpoint{1.031123in}{1.409823in}}%
\pgfpathlineto{\pgfqpoint{1.030622in}{1.400197in}}%
\pgfpathlineto{\pgfqpoint{1.034771in}{1.396805in}}%
\pgfpathlineto{\pgfqpoint{1.038444in}{1.389127in}}%
\pgfpathlineto{\pgfqpoint{1.032898in}{1.379983in}}%
\pgfpathlineto{\pgfqpoint{1.026917in}{1.376417in}}%
\pgfpathlineto{\pgfqpoint{1.023679in}{1.356359in}}%
\pgfpathlineto{\pgfqpoint{1.016803in}{1.357032in}}%
\pgfpathlineto{\pgfqpoint{1.015614in}{1.351012in}}%
\pgfpathlineto{\pgfqpoint{1.008850in}{1.352444in}}%
\pgfpathlineto{\pgfqpoint{1.005205in}{1.331999in}}%
\pgfpathlineto{\pgfqpoint{0.991804in}{1.334673in}}%
\pgfpathlineto{\pgfqpoint{0.983512in}{1.292540in}}%
\pgfpathclose%
\pgfusepath{fill}%
\end{pgfscope}%
\begin{pgfscope}%
\pgfpathrectangle{\pgfqpoint{0.100000in}{0.100000in}}{\pgfqpoint{3.608454in}{2.310000in}}%
\pgfusepath{clip}%
\pgfsetbuttcap%
\pgfsetmiterjoin%
\definecolor{currentfill}{rgb}{0.000000,0.490196,0.754902}%
\pgfsetfillcolor{currentfill}%
\pgfsetlinewidth{0.000000pt}%
\definecolor{currentstroke}{rgb}{0.000000,0.000000,0.000000}%
\pgfsetstrokecolor{currentstroke}%
\pgfsetstrokeopacity{0.000000}%
\pgfsetdash{}{0pt}%
\pgfpathmoveto{\pgfqpoint{2.684273in}{1.275179in}}%
\pgfpathlineto{\pgfqpoint{2.691418in}{1.266749in}}%
\pgfpathlineto{\pgfqpoint{2.694291in}{1.266816in}}%
\pgfpathlineto{\pgfqpoint{2.694535in}{1.255275in}}%
\pgfpathlineto{\pgfqpoint{2.701382in}{1.255633in}}%
\pgfpathlineto{\pgfqpoint{2.704025in}{1.249697in}}%
\pgfpathlineto{\pgfqpoint{2.693473in}{1.249077in}}%
\pgfpathlineto{\pgfqpoint{2.692446in}{1.245040in}}%
\pgfpathlineto{\pgfqpoint{2.667849in}{1.244767in}}%
\pgfpathlineto{\pgfqpoint{2.665870in}{1.236725in}}%
\pgfpathlineto{\pgfqpoint{2.661850in}{1.235103in}}%
\pgfpathlineto{\pgfqpoint{2.642473in}{1.232080in}}%
\pgfpathlineto{\pgfqpoint{2.641029in}{1.235231in}}%
\pgfpathlineto{\pgfqpoint{2.632971in}{1.237093in}}%
\pgfpathlineto{\pgfqpoint{2.628815in}{1.247944in}}%
\pgfpathlineto{\pgfqpoint{2.630443in}{1.257485in}}%
\pgfpathlineto{\pgfqpoint{2.625972in}{1.268011in}}%
\pgfpathlineto{\pgfqpoint{2.627023in}{1.272051in}}%
\pgfpathlineto{\pgfqpoint{2.630595in}{1.278625in}}%
\pgfpathlineto{\pgfqpoint{2.635552in}{1.279697in}}%
\pgfpathlineto{\pgfqpoint{2.634127in}{1.288599in}}%
\pgfpathlineto{\pgfqpoint{2.636960in}{1.295013in}}%
\pgfpathlineto{\pgfqpoint{2.640611in}{1.294225in}}%
\pgfpathlineto{\pgfqpoint{2.643610in}{1.301513in}}%
\pgfpathlineto{\pgfqpoint{2.648711in}{1.297118in}}%
\pgfpathlineto{\pgfqpoint{2.649917in}{1.290900in}}%
\pgfpathlineto{\pgfqpoint{2.655765in}{1.288068in}}%
\pgfpathlineto{\pgfqpoint{2.669652in}{1.289191in}}%
\pgfpathlineto{\pgfqpoint{2.675593in}{1.285078in}}%
\pgfpathlineto{\pgfqpoint{2.678952in}{1.278333in}}%
\pgfpathlineto{\pgfqpoint{2.684273in}{1.275179in}}%
\pgfpathclose%
\pgfusepath{fill}%
\end{pgfscope}%
\begin{pgfscope}%
\pgfpathrectangle{\pgfqpoint{0.100000in}{0.100000in}}{\pgfqpoint{3.608454in}{2.310000in}}%
\pgfusepath{clip}%
\pgfsetbuttcap%
\pgfsetmiterjoin%
\definecolor{currentfill}{rgb}{0.000000,0.478431,0.760784}%
\pgfsetfillcolor{currentfill}%
\pgfsetlinewidth{0.000000pt}%
\definecolor{currentstroke}{rgb}{0.000000,0.000000,0.000000}%
\pgfsetstrokecolor{currentstroke}%
\pgfsetstrokeopacity{0.000000}%
\pgfsetdash{}{0pt}%
\pgfpathmoveto{\pgfqpoint{1.896146in}{2.138772in}}%
\pgfpathlineto{\pgfqpoint{1.895146in}{2.116164in}}%
\pgfpathlineto{\pgfqpoint{1.896095in}{2.102164in}}%
\pgfpathlineto{\pgfqpoint{1.895459in}{2.088252in}}%
\pgfpathlineto{\pgfqpoint{1.867703in}{2.089585in}}%
\pgfpathlineto{\pgfqpoint{1.868375in}{2.103508in}}%
\pgfpathlineto{\pgfqpoint{1.833772in}{2.105376in}}%
\pgfpathlineto{\pgfqpoint{1.832606in}{2.119399in}}%
\pgfpathlineto{\pgfqpoint{1.833904in}{2.141855in}}%
\pgfpathlineto{\pgfqpoint{1.896146in}{2.138772in}}%
\pgfpathclose%
\pgfusepath{fill}%
\end{pgfscope}%
\begin{pgfscope}%
\pgfpathrectangle{\pgfqpoint{0.100000in}{0.100000in}}{\pgfqpoint{3.608454in}{2.310000in}}%
\pgfusepath{clip}%
\pgfsetbuttcap%
\pgfsetmiterjoin%
\definecolor{currentfill}{rgb}{0.000000,0.490196,0.754902}%
\pgfsetfillcolor{currentfill}%
\pgfsetlinewidth{0.000000pt}%
\definecolor{currentstroke}{rgb}{0.000000,0.000000,0.000000}%
\pgfsetstrokecolor{currentstroke}%
\pgfsetstrokeopacity{0.000000}%
\pgfsetdash{}{0pt}%
\pgfpathmoveto{\pgfqpoint{2.013608in}{1.046935in}}%
\pgfpathlineto{\pgfqpoint{2.013392in}{1.033129in}}%
\pgfpathlineto{\pgfqpoint{2.010212in}{1.033190in}}%
\pgfpathlineto{\pgfqpoint{2.009922in}{1.016922in}}%
\pgfpathlineto{\pgfqpoint{1.999760in}{1.016432in}}%
\pgfpathlineto{\pgfqpoint{1.994760in}{1.012926in}}%
\pgfpathlineto{\pgfqpoint{1.988856in}{1.020093in}}%
\pgfpathlineto{\pgfqpoint{1.984031in}{1.018827in}}%
\pgfpathlineto{\pgfqpoint{1.981120in}{1.013018in}}%
\pgfpathlineto{\pgfqpoint{1.933314in}{1.014281in}}%
\pgfpathlineto{\pgfqpoint{1.934200in}{1.052365in}}%
\pgfpathlineto{\pgfqpoint{1.923203in}{1.052514in}}%
\pgfpathlineto{\pgfqpoint{1.911760in}{1.056365in}}%
\pgfpathlineto{\pgfqpoint{1.893326in}{1.056979in}}%
\pgfpathlineto{\pgfqpoint{1.893386in}{1.070793in}}%
\pgfpathlineto{\pgfqpoint{1.893880in}{1.084537in}}%
\pgfpathlineto{\pgfqpoint{1.900701in}{1.084287in}}%
\pgfpathlineto{\pgfqpoint{1.901714in}{1.119083in}}%
\pgfpathlineto{\pgfqpoint{1.949595in}{1.117555in}}%
\pgfpathlineto{\pgfqpoint{1.956419in}{1.117360in}}%
\pgfpathlineto{\pgfqpoint{1.956052in}{1.103169in}}%
\pgfpathlineto{\pgfqpoint{1.969632in}{1.099419in}}%
\pgfpathlineto{\pgfqpoint{2.002912in}{1.098705in}}%
\pgfpathlineto{\pgfqpoint{2.002513in}{1.074759in}}%
\pgfpathlineto{\pgfqpoint{2.001980in}{1.055914in}}%
\pgfpathlineto{\pgfqpoint{2.010159in}{1.057965in}}%
\pgfpathlineto{\pgfqpoint{2.013608in}{1.046935in}}%
\pgfpathclose%
\pgfusepath{fill}%
\end{pgfscope}%
\begin{pgfscope}%
\pgfpathrectangle{\pgfqpoint{0.100000in}{0.100000in}}{\pgfqpoint{3.608454in}{2.310000in}}%
\pgfusepath{clip}%
\pgfsetbuttcap%
\pgfsetmiterjoin%
\definecolor{currentfill}{rgb}{0.000000,0.482353,0.758824}%
\pgfsetfillcolor{currentfill}%
\pgfsetlinewidth{0.000000pt}%
\definecolor{currentstroke}{rgb}{0.000000,0.000000,0.000000}%
\pgfsetstrokecolor{currentstroke}%
\pgfsetstrokeopacity{0.000000}%
\pgfsetdash{}{0pt}%
\pgfpathmoveto{\pgfqpoint{2.111687in}{1.695719in}}%
\pgfpathlineto{\pgfqpoint{2.115164in}{1.695734in}}%
\pgfpathlineto{\pgfqpoint{2.115120in}{1.723283in}}%
\pgfpathlineto{\pgfqpoint{2.177278in}{1.723751in}}%
\pgfpathlineto{\pgfqpoint{2.184096in}{1.723863in}}%
\pgfpathlineto{\pgfqpoint{2.184501in}{1.696223in}}%
\pgfpathlineto{\pgfqpoint{2.165979in}{1.695985in}}%
\pgfpathlineto{\pgfqpoint{2.138833in}{1.695825in}}%
\pgfpathlineto{\pgfqpoint{2.138939in}{1.676359in}}%
\pgfpathlineto{\pgfqpoint{2.111765in}{1.676267in}}%
\pgfpathlineto{\pgfqpoint{2.111687in}{1.695719in}}%
\pgfpathclose%
\pgfusepath{fill}%
\end{pgfscope}%
\begin{pgfscope}%
\pgfpathrectangle{\pgfqpoint{0.100000in}{0.100000in}}{\pgfqpoint{3.608454in}{2.310000in}}%
\pgfusepath{clip}%
\pgfsetbuttcap%
\pgfsetmiterjoin%
\definecolor{currentfill}{rgb}{0.000000,0.823529,0.588235}%
\pgfsetfillcolor{currentfill}%
\pgfsetlinewidth{0.000000pt}%
\definecolor{currentstroke}{rgb}{0.000000,0.000000,0.000000}%
\pgfsetstrokecolor{currentstroke}%
\pgfsetstrokeopacity{0.000000}%
\pgfsetdash{}{0pt}%
\pgfpathmoveto{\pgfqpoint{2.804581in}{1.191394in}}%
\pgfpathlineto{\pgfqpoint{2.789452in}{1.189777in}}%
\pgfpathlineto{\pgfqpoint{2.783865in}{1.195392in}}%
\pgfpathlineto{\pgfqpoint{2.779295in}{1.206052in}}%
\pgfpathlineto{\pgfqpoint{2.781636in}{1.214245in}}%
\pgfpathlineto{\pgfqpoint{2.778898in}{1.226398in}}%
\pgfpathlineto{\pgfqpoint{2.780361in}{1.233265in}}%
\pgfpathlineto{\pgfqpoint{2.781851in}{1.240501in}}%
\pgfpathlineto{\pgfqpoint{2.788710in}{1.248796in}}%
\pgfpathlineto{\pgfqpoint{2.797121in}{1.243378in}}%
\pgfpathlineto{\pgfqpoint{2.801049in}{1.243631in}}%
\pgfpathlineto{\pgfqpoint{2.809988in}{1.252552in}}%
\pgfpathlineto{\pgfqpoint{2.816644in}{1.252449in}}%
\pgfpathlineto{\pgfqpoint{2.826981in}{1.247494in}}%
\pgfpathlineto{\pgfqpoint{2.826345in}{1.240166in}}%
\pgfpathlineto{\pgfqpoint{2.831285in}{1.222519in}}%
\pgfpathlineto{\pgfqpoint{2.815210in}{1.209500in}}%
\pgfpathlineto{\pgfqpoint{2.812724in}{1.202234in}}%
\pgfpathlineto{\pgfqpoint{2.806029in}{1.196000in}}%
\pgfpathlineto{\pgfqpoint{2.804581in}{1.191394in}}%
\pgfpathclose%
\pgfusepath{fill}%
\end{pgfscope}%
\begin{pgfscope}%
\pgfpathrectangle{\pgfqpoint{0.100000in}{0.100000in}}{\pgfqpoint{3.608454in}{2.310000in}}%
\pgfusepath{clip}%
\pgfsetbuttcap%
\pgfsetmiterjoin%
\definecolor{currentfill}{rgb}{0.000000,0.615686,0.692157}%
\pgfsetfillcolor{currentfill}%
\pgfsetlinewidth{0.000000pt}%
\definecolor{currentstroke}{rgb}{0.000000,0.000000,0.000000}%
\pgfsetstrokecolor{currentstroke}%
\pgfsetstrokeopacity{0.000000}%
\pgfsetdash{}{0pt}%
\pgfpathmoveto{\pgfqpoint{2.765354in}{1.245580in}}%
\pgfpathlineto{\pgfqpoint{2.761764in}{1.247341in}}%
\pgfpathlineto{\pgfqpoint{2.752556in}{1.236924in}}%
\pgfpathlineto{\pgfqpoint{2.751350in}{1.246500in}}%
\pgfpathlineto{\pgfqpoint{2.742755in}{1.250848in}}%
\pgfpathlineto{\pgfqpoint{2.741472in}{1.258776in}}%
\pgfpathlineto{\pgfqpoint{2.729948in}{1.258824in}}%
\pgfpathlineto{\pgfqpoint{2.729699in}{1.283399in}}%
\pgfpathlineto{\pgfqpoint{2.727518in}{1.286122in}}%
\pgfpathlineto{\pgfqpoint{2.733018in}{1.291688in}}%
\pgfpathlineto{\pgfqpoint{2.740147in}{1.292411in}}%
\pgfpathlineto{\pgfqpoint{2.757404in}{1.276520in}}%
\pgfpathlineto{\pgfqpoint{2.759930in}{1.280063in}}%
\pgfpathlineto{\pgfqpoint{2.773095in}{1.263245in}}%
\pgfpathlineto{\pgfqpoint{2.767417in}{1.258344in}}%
\pgfpathlineto{\pgfqpoint{2.765354in}{1.245580in}}%
\pgfpathclose%
\pgfusepath{fill}%
\end{pgfscope}%
\begin{pgfscope}%
\pgfpathrectangle{\pgfqpoint{0.100000in}{0.100000in}}{\pgfqpoint{3.608454in}{2.310000in}}%
\pgfusepath{clip}%
\pgfsetbuttcap%
\pgfsetmiterjoin%
\definecolor{currentfill}{rgb}{0.000000,0.403922,0.798039}%
\pgfsetfillcolor{currentfill}%
\pgfsetlinewidth{0.000000pt}%
\definecolor{currentstroke}{rgb}{0.000000,0.000000,0.000000}%
\pgfsetstrokecolor{currentstroke}%
\pgfsetstrokeopacity{0.000000}%
\pgfsetdash{}{0pt}%
\pgfpathmoveto{\pgfqpoint{1.682362in}{1.462771in}}%
\pgfpathlineto{\pgfqpoint{1.677720in}{1.401672in}}%
\pgfpathlineto{\pgfqpoint{1.610603in}{1.406638in}}%
\pgfpathlineto{\pgfqpoint{1.577074in}{1.409804in}}%
\pgfpathlineto{\pgfqpoint{1.580447in}{1.444028in}}%
\pgfpathlineto{\pgfqpoint{1.594536in}{1.442692in}}%
\pgfpathlineto{\pgfqpoint{1.598652in}{1.483744in}}%
\pgfpathlineto{\pgfqpoint{1.591803in}{1.484401in}}%
\pgfpathlineto{\pgfqpoint{1.595804in}{1.522015in}}%
\pgfpathlineto{\pgfqpoint{1.650561in}{1.517162in}}%
\pgfpathlineto{\pgfqpoint{1.649013in}{1.497206in}}%
\pgfpathlineto{\pgfqpoint{1.684833in}{1.494331in}}%
\pgfpathlineto{\pgfqpoint{1.684512in}{1.490228in}}%
\pgfpathlineto{\pgfqpoint{1.682362in}{1.462771in}}%
\pgfpathclose%
\pgfusepath{fill}%
\end{pgfscope}%
\begin{pgfscope}%
\pgfpathrectangle{\pgfqpoint{0.100000in}{0.100000in}}{\pgfqpoint{3.608454in}{2.310000in}}%
\pgfusepath{clip}%
\pgfsetbuttcap%
\pgfsetmiterjoin%
\definecolor{currentfill}{rgb}{0.000000,0.470588,0.764706}%
\pgfsetfillcolor{currentfill}%
\pgfsetlinewidth{0.000000pt}%
\definecolor{currentstroke}{rgb}{0.000000,0.000000,0.000000}%
\pgfsetstrokecolor{currentstroke}%
\pgfsetstrokeopacity{0.000000}%
\pgfsetdash{}{0pt}%
\pgfpathmoveto{\pgfqpoint{1.846276in}{0.936427in}}%
\pgfpathlineto{\pgfqpoint{1.844714in}{0.901800in}}%
\pgfpathlineto{\pgfqpoint{1.810377in}{0.903512in}}%
\pgfpathlineto{\pgfqpoint{1.811548in}{0.930044in}}%
\pgfpathlineto{\pgfqpoint{1.811868in}{0.938019in}}%
\pgfpathlineto{\pgfqpoint{1.846276in}{0.936427in}}%
\pgfpathclose%
\pgfusepath{fill}%
\end{pgfscope}%
\begin{pgfscope}%
\pgfpathrectangle{\pgfqpoint{0.100000in}{0.100000in}}{\pgfqpoint{3.608454in}{2.310000in}}%
\pgfusepath{clip}%
\pgfsetbuttcap%
\pgfsetmiterjoin%
\definecolor{currentfill}{rgb}{0.000000,0.588235,0.705882}%
\pgfsetfillcolor{currentfill}%
\pgfsetlinewidth{0.000000pt}%
\definecolor{currentstroke}{rgb}{0.000000,0.000000,0.000000}%
\pgfsetstrokecolor{currentstroke}%
\pgfsetstrokeopacity{0.000000}%
\pgfsetdash{}{0pt}%
\pgfpathmoveto{\pgfqpoint{2.238456in}{1.359055in}}%
\pgfpathlineto{\pgfqpoint{2.232957in}{1.349258in}}%
\pgfpathlineto{\pgfqpoint{2.233532in}{1.346189in}}%
\pgfpathlineto{\pgfqpoint{2.226701in}{1.338754in}}%
\pgfpathlineto{\pgfqpoint{2.226654in}{1.335261in}}%
\pgfpathlineto{\pgfqpoint{2.199219in}{1.335898in}}%
\pgfpathlineto{\pgfqpoint{2.198627in}{1.357636in}}%
\pgfpathlineto{\pgfqpoint{2.188988in}{1.359699in}}%
\pgfpathlineto{\pgfqpoint{2.182827in}{1.356530in}}%
\pgfpathlineto{\pgfqpoint{2.182439in}{1.381621in}}%
\pgfpathlineto{\pgfqpoint{2.181819in}{1.416001in}}%
\pgfpathlineto{\pgfqpoint{2.206018in}{1.417015in}}%
\pgfpathlineto{\pgfqpoint{2.206226in}{1.396171in}}%
\pgfpathlineto{\pgfqpoint{2.211089in}{1.394791in}}%
\pgfpathlineto{\pgfqpoint{2.212943in}{1.384548in}}%
\pgfpathlineto{\pgfqpoint{2.211917in}{1.379066in}}%
\pgfpathlineto{\pgfqpoint{2.215503in}{1.374729in}}%
\pgfpathlineto{\pgfqpoint{2.221539in}{1.373866in}}%
\pgfpathlineto{\pgfqpoint{2.224430in}{1.367560in}}%
\pgfpathlineto{\pgfqpoint{2.230926in}{1.365957in}}%
\pgfpathlineto{\pgfqpoint{2.233381in}{1.359171in}}%
\pgfpathlineto{\pgfqpoint{2.238456in}{1.359055in}}%
\pgfpathclose%
\pgfusepath{fill}%
\end{pgfscope}%
\begin{pgfscope}%
\pgfpathrectangle{\pgfqpoint{0.100000in}{0.100000in}}{\pgfqpoint{3.608454in}{2.310000in}}%
\pgfusepath{clip}%
\pgfsetbuttcap%
\pgfsetmiterjoin%
\definecolor{currentfill}{rgb}{0.000000,0.356863,0.821569}%
\pgfsetfillcolor{currentfill}%
\pgfsetlinewidth{0.000000pt}%
\definecolor{currentstroke}{rgb}{0.000000,0.000000,0.000000}%
\pgfsetstrokecolor{currentstroke}%
\pgfsetstrokeopacity{0.000000}%
\pgfsetdash{}{0pt}%
\pgfpathmoveto{\pgfqpoint{2.886708in}{1.438806in}}%
\pgfpathlineto{\pgfqpoint{2.887216in}{1.431682in}}%
\pgfpathlineto{\pgfqpoint{2.880330in}{1.431219in}}%
\pgfpathlineto{\pgfqpoint{2.880748in}{1.424222in}}%
\pgfpathlineto{\pgfqpoint{2.867003in}{1.423229in}}%
\pgfpathlineto{\pgfqpoint{2.867100in}{1.421144in}}%
\pgfpathlineto{\pgfqpoint{2.853130in}{1.419128in}}%
\pgfpathlineto{\pgfqpoint{2.852553in}{1.427030in}}%
\pgfpathlineto{\pgfqpoint{2.837395in}{1.428067in}}%
\pgfpathlineto{\pgfqpoint{2.820238in}{1.426570in}}%
\pgfpathlineto{\pgfqpoint{2.819277in}{1.440727in}}%
\pgfpathlineto{\pgfqpoint{2.794880in}{1.439319in}}%
\pgfpathlineto{\pgfqpoint{2.798306in}{1.443885in}}%
\pgfpathlineto{\pgfqpoint{2.800713in}{1.463456in}}%
\pgfpathlineto{\pgfqpoint{2.799407in}{1.480427in}}%
\pgfpathlineto{\pgfqpoint{2.795940in}{1.480285in}}%
\pgfpathlineto{\pgfqpoint{2.795078in}{1.502065in}}%
\pgfpathlineto{\pgfqpoint{2.801132in}{1.502933in}}%
\pgfpathlineto{\pgfqpoint{2.811270in}{1.504385in}}%
\pgfpathlineto{\pgfqpoint{2.811880in}{1.499432in}}%
\pgfpathlineto{\pgfqpoint{2.831500in}{1.499705in}}%
\pgfpathlineto{\pgfqpoint{2.831782in}{1.495222in}}%
\pgfpathlineto{\pgfqpoint{2.848845in}{1.496480in}}%
\pgfpathlineto{\pgfqpoint{2.846858in}{1.512561in}}%
\pgfpathlineto{\pgfqpoint{2.862804in}{1.514836in}}%
\pgfpathlineto{\pgfqpoint{2.871819in}{1.517489in}}%
\pgfpathlineto{\pgfqpoint{2.874225in}{1.518225in}}%
\pgfpathlineto{\pgfqpoint{2.876912in}{1.491928in}}%
\pgfpathlineto{\pgfqpoint{2.878220in}{1.486351in}}%
\pgfpathlineto{\pgfqpoint{2.879950in}{1.469292in}}%
\pgfpathlineto{\pgfqpoint{2.878244in}{1.466104in}}%
\pgfpathlineto{\pgfqpoint{2.864318in}{1.465519in}}%
\pgfpathlineto{\pgfqpoint{2.864747in}{1.458144in}}%
\pgfpathlineto{\pgfqpoint{2.869385in}{1.458446in}}%
\pgfpathlineto{\pgfqpoint{2.872185in}{1.451622in}}%
\pgfpathlineto{\pgfqpoint{2.872947in}{1.440094in}}%
\pgfpathlineto{\pgfqpoint{2.886708in}{1.438806in}}%
\pgfpathclose%
\pgfusepath{fill}%
\end{pgfscope}%
\begin{pgfscope}%
\pgfpathrectangle{\pgfqpoint{0.100000in}{0.100000in}}{\pgfqpoint{3.608454in}{2.310000in}}%
\pgfusepath{clip}%
\pgfsetbuttcap%
\pgfsetmiterjoin%
\definecolor{currentfill}{rgb}{0.000000,0.627451,0.686275}%
\pgfsetfillcolor{currentfill}%
\pgfsetlinewidth{0.000000pt}%
\definecolor{currentstroke}{rgb}{0.000000,0.000000,0.000000}%
\pgfsetstrokecolor{currentstroke}%
\pgfsetstrokeopacity{0.000000}%
\pgfsetdash{}{0pt}%
\pgfpathmoveto{\pgfqpoint{2.130612in}{1.428716in}}%
\pgfpathlineto{\pgfqpoint{2.130971in}{1.449316in}}%
\pgfpathlineto{\pgfqpoint{2.128892in}{1.463893in}}%
\pgfpathlineto{\pgfqpoint{2.138550in}{1.463912in}}%
\pgfpathlineto{\pgfqpoint{2.180394in}{1.464846in}}%
\pgfpathlineto{\pgfqpoint{2.181491in}{1.429705in}}%
\pgfpathlineto{\pgfqpoint{2.154147in}{1.429674in}}%
\pgfpathlineto{\pgfqpoint{2.154074in}{1.421652in}}%
\pgfpathlineto{\pgfqpoint{2.130676in}{1.421843in}}%
\pgfpathlineto{\pgfqpoint{2.130612in}{1.428716in}}%
\pgfpathclose%
\pgfusepath{fill}%
\end{pgfscope}%
\begin{pgfscope}%
\pgfpathrectangle{\pgfqpoint{0.100000in}{0.100000in}}{\pgfqpoint{3.608454in}{2.310000in}}%
\pgfusepath{clip}%
\pgfsetbuttcap%
\pgfsetmiterjoin%
\definecolor{currentfill}{rgb}{0.000000,0.427451,0.786275}%
\pgfsetfillcolor{currentfill}%
\pgfsetlinewidth{0.000000pt}%
\definecolor{currentstroke}{rgb}{0.000000,0.000000,0.000000}%
\pgfsetstrokecolor{currentstroke}%
\pgfsetstrokeopacity{0.000000}%
\pgfsetdash{}{0pt}%
\pgfpathmoveto{\pgfqpoint{2.073124in}{1.778481in}}%
\pgfpathlineto{\pgfqpoint{2.072659in}{1.785449in}}%
\pgfpathlineto{\pgfqpoint{2.024747in}{1.786093in}}%
\pgfpathlineto{\pgfqpoint{2.025168in}{1.813601in}}%
\pgfpathlineto{\pgfqpoint{2.025490in}{1.836778in}}%
\pgfpathlineto{\pgfqpoint{2.034935in}{1.834889in}}%
\pgfpathlineto{\pgfqpoint{2.041094in}{1.830482in}}%
\pgfpathlineto{\pgfqpoint{2.044201in}{1.834191in}}%
\pgfpathlineto{\pgfqpoint{2.044431in}{1.847876in}}%
\pgfpathlineto{\pgfqpoint{2.099439in}{1.847533in}}%
\pgfpathlineto{\pgfqpoint{2.120016in}{1.847577in}}%
\pgfpathlineto{\pgfqpoint{2.120021in}{1.840682in}}%
\pgfpathlineto{\pgfqpoint{2.120405in}{1.806180in}}%
\pgfpathlineto{\pgfqpoint{2.079580in}{1.806133in}}%
\pgfpathlineto{\pgfqpoint{2.079554in}{1.795195in}}%
\pgfpathlineto{\pgfqpoint{2.086458in}{1.790823in}}%
\pgfpathlineto{\pgfqpoint{2.086396in}{1.778361in}}%
\pgfpathlineto{\pgfqpoint{2.073124in}{1.778481in}}%
\pgfpathclose%
\pgfusepath{fill}%
\end{pgfscope}%
\begin{pgfscope}%
\pgfpathrectangle{\pgfqpoint{0.100000in}{0.100000in}}{\pgfqpoint{3.608454in}{2.310000in}}%
\pgfusepath{clip}%
\pgfsetbuttcap%
\pgfsetmiterjoin%
\definecolor{currentfill}{rgb}{0.000000,0.705882,0.647059}%
\pgfsetfillcolor{currentfill}%
\pgfsetlinewidth{0.000000pt}%
\definecolor{currentstroke}{rgb}{0.000000,0.000000,0.000000}%
\pgfsetstrokecolor{currentstroke}%
\pgfsetstrokeopacity{0.000000}%
\pgfsetdash{}{0pt}%
\pgfpathmoveto{\pgfqpoint{2.248336in}{1.167927in}}%
\pgfpathlineto{\pgfqpoint{2.248459in}{1.143536in}}%
\pgfpathlineto{\pgfqpoint{2.243255in}{1.143432in}}%
\pgfpathlineto{\pgfqpoint{2.178864in}{1.142325in}}%
\pgfpathlineto{\pgfqpoint{2.165488in}{1.142173in}}%
\pgfpathlineto{\pgfqpoint{2.165811in}{1.161887in}}%
\pgfpathlineto{\pgfqpoint{2.166408in}{1.204778in}}%
\pgfpathlineto{\pgfqpoint{2.193249in}{1.204475in}}%
\pgfpathlineto{\pgfqpoint{2.194554in}{1.204465in}}%
\pgfpathlineto{\pgfqpoint{2.194584in}{1.181869in}}%
\pgfpathlineto{\pgfqpoint{2.211810in}{1.181895in}}%
\pgfpathlineto{\pgfqpoint{2.211773in}{1.168036in}}%
\pgfpathlineto{\pgfqpoint{2.248336in}{1.167927in}}%
\pgfpathclose%
\pgfusepath{fill}%
\end{pgfscope}%
\begin{pgfscope}%
\pgfpathrectangle{\pgfqpoint{0.100000in}{0.100000in}}{\pgfqpoint{3.608454in}{2.310000in}}%
\pgfusepath{clip}%
\pgfsetbuttcap%
\pgfsetmiterjoin%
\definecolor{currentfill}{rgb}{0.000000,0.517647,0.741176}%
\pgfsetfillcolor{currentfill}%
\pgfsetlinewidth{0.000000pt}%
\definecolor{currentstroke}{rgb}{0.000000,0.000000,0.000000}%
\pgfsetstrokecolor{currentstroke}%
\pgfsetstrokeopacity{0.000000}%
\pgfsetdash{}{0pt}%
\pgfpathmoveto{\pgfqpoint{2.722959in}{1.184982in}}%
\pgfpathlineto{\pgfqpoint{2.754147in}{1.186574in}}%
\pgfpathlineto{\pgfqpoint{2.758225in}{1.180770in}}%
\pgfpathlineto{\pgfqpoint{2.744380in}{1.183721in}}%
\pgfpathlineto{\pgfqpoint{2.738211in}{1.181250in}}%
\pgfpathlineto{\pgfqpoint{2.734679in}{1.168795in}}%
\pgfpathlineto{\pgfqpoint{2.735445in}{1.160714in}}%
\pgfpathlineto{\pgfqpoint{2.738900in}{1.152461in}}%
\pgfpathlineto{\pgfqpoint{2.735267in}{1.144371in}}%
\pgfpathlineto{\pgfqpoint{2.729930in}{1.144288in}}%
\pgfpathlineto{\pgfqpoint{2.729200in}{1.134130in}}%
\pgfpathlineto{\pgfqpoint{2.722754in}{1.138754in}}%
\pgfpathlineto{\pgfqpoint{2.712831in}{1.140583in}}%
\pgfpathlineto{\pgfqpoint{2.703971in}{1.134522in}}%
\pgfpathlineto{\pgfqpoint{2.701451in}{1.138405in}}%
\pgfpathlineto{\pgfqpoint{2.695803in}{1.138527in}}%
\pgfpathlineto{\pgfqpoint{2.693376in}{1.142517in}}%
\pgfpathlineto{\pgfqpoint{2.693285in}{1.148452in}}%
\pgfpathlineto{\pgfqpoint{2.689646in}{1.154594in}}%
\pgfpathlineto{\pgfqpoint{2.690027in}{1.171438in}}%
\pgfpathlineto{\pgfqpoint{2.696790in}{1.173951in}}%
\pgfpathlineto{\pgfqpoint{2.705300in}{1.170782in}}%
\pgfpathlineto{\pgfqpoint{2.710970in}{1.165941in}}%
\pgfpathlineto{\pgfqpoint{2.713741in}{1.174022in}}%
\pgfpathlineto{\pgfqpoint{2.722324in}{1.177007in}}%
\pgfpathlineto{\pgfqpoint{2.722959in}{1.184982in}}%
\pgfpathclose%
\pgfusepath{fill}%
\end{pgfscope}%
\begin{pgfscope}%
\pgfpathrectangle{\pgfqpoint{0.100000in}{0.100000in}}{\pgfqpoint{3.608454in}{2.310000in}}%
\pgfusepath{clip}%
\pgfsetbuttcap%
\pgfsetmiterjoin%
\definecolor{currentfill}{rgb}{0.000000,0.341176,0.829412}%
\pgfsetfillcolor{currentfill}%
\pgfsetlinewidth{0.000000pt}%
\definecolor{currentstroke}{rgb}{0.000000,0.000000,0.000000}%
\pgfsetstrokecolor{currentstroke}%
\pgfsetstrokeopacity{0.000000}%
\pgfsetdash{}{0pt}%
\pgfpathmoveto{\pgfqpoint{1.586973in}{1.721839in}}%
\pgfpathlineto{\pgfqpoint{1.538610in}{1.726627in}}%
\pgfpathlineto{\pgfqpoint{1.474731in}{1.733639in}}%
\pgfpathlineto{\pgfqpoint{1.481253in}{1.786760in}}%
\pgfpathlineto{\pgfqpoint{1.486605in}{1.827974in}}%
\pgfpathlineto{\pgfqpoint{1.488701in}{1.851762in}}%
\pgfpathlineto{\pgfqpoint{1.494987in}{1.851589in}}%
\pgfpathlineto{\pgfqpoint{1.543825in}{1.845887in}}%
\pgfpathlineto{\pgfqpoint{1.598608in}{1.839877in}}%
\pgfpathlineto{\pgfqpoint{1.592333in}{1.775294in}}%
\pgfpathlineto{\pgfqpoint{1.586973in}{1.721839in}}%
\pgfpathclose%
\pgfusepath{fill}%
\end{pgfscope}%
\begin{pgfscope}%
\pgfpathrectangle{\pgfqpoint{0.100000in}{0.100000in}}{\pgfqpoint{3.608454in}{2.310000in}}%
\pgfusepath{clip}%
\pgfsetbuttcap%
\pgfsetmiterjoin%
\definecolor{currentfill}{rgb}{0.000000,0.776471,0.611765}%
\pgfsetfillcolor{currentfill}%
\pgfsetlinewidth{0.000000pt}%
\definecolor{currentstroke}{rgb}{0.000000,0.000000,0.000000}%
\pgfsetstrokecolor{currentstroke}%
\pgfsetstrokeopacity{0.000000}%
\pgfsetdash{}{0pt}%
\pgfpathmoveto{\pgfqpoint{1.365564in}{0.266179in}}%
\pgfpathlineto{\pgfqpoint{1.364279in}{0.272148in}}%
\pgfpathlineto{\pgfqpoint{1.358083in}{0.284462in}}%
\pgfpathlineto{\pgfqpoint{1.360076in}{0.291758in}}%
\pgfpathlineto{\pgfqpoint{1.373325in}{0.304088in}}%
\pgfpathlineto{\pgfqpoint{1.377160in}{0.317108in}}%
\pgfpathlineto{\pgfqpoint{1.381805in}{0.326247in}}%
\pgfpathlineto{\pgfqpoint{1.380932in}{0.332508in}}%
\pgfpathlineto{\pgfqpoint{1.385909in}{0.340951in}}%
\pgfpathlineto{\pgfqpoint{1.410377in}{0.343837in}}%
\pgfpathlineto{\pgfqpoint{1.412863in}{0.346234in}}%
\pgfpathlineto{\pgfqpoint{1.413879in}{0.356269in}}%
\pgfpathlineto{\pgfqpoint{1.419360in}{0.363772in}}%
\pgfpathlineto{\pgfqpoint{1.424674in}{0.362812in}}%
\pgfpathlineto{\pgfqpoint{1.428368in}{0.354664in}}%
\pgfpathlineto{\pgfqpoint{1.431876in}{0.341640in}}%
\pgfpathlineto{\pgfqpoint{1.439179in}{0.333553in}}%
\pgfpathlineto{\pgfqpoint{1.446187in}{0.319277in}}%
\pgfpathlineto{\pgfqpoint{1.449694in}{0.299992in}}%
\pgfpathlineto{\pgfqpoint{1.442998in}{0.291249in}}%
\pgfpathlineto{\pgfqpoint{1.448609in}{0.287822in}}%
\pgfpathlineto{\pgfqpoint{1.445177in}{0.280721in}}%
\pgfpathlineto{\pgfqpoint{1.446258in}{0.271967in}}%
\pgfpathlineto{\pgfqpoint{1.449844in}{0.264785in}}%
\pgfpathlineto{\pgfqpoint{1.446708in}{0.261925in}}%
\pgfpathlineto{\pgfqpoint{1.430233in}{0.260876in}}%
\pgfpathlineto{\pgfqpoint{1.416961in}{0.263446in}}%
\pgfpathlineto{\pgfqpoint{1.406310in}{0.270716in}}%
\pgfpathlineto{\pgfqpoint{1.400271in}{0.270027in}}%
\pgfpathlineto{\pgfqpoint{1.387382in}{0.272320in}}%
\pgfpathlineto{\pgfqpoint{1.365564in}{0.266179in}}%
\pgfpathclose%
\pgfusepath{fill}%
\end{pgfscope}%
\begin{pgfscope}%
\pgfpathrectangle{\pgfqpoint{0.100000in}{0.100000in}}{\pgfqpoint{3.608454in}{2.310000in}}%
\pgfusepath{clip}%
\pgfsetbuttcap%
\pgfsetmiterjoin%
\definecolor{currentfill}{rgb}{0.000000,0.368627,0.815686}%
\pgfsetfillcolor{currentfill}%
\pgfsetlinewidth{0.000000pt}%
\definecolor{currentstroke}{rgb}{0.000000,0.000000,0.000000}%
\pgfsetstrokecolor{currentstroke}%
\pgfsetstrokeopacity{0.000000}%
\pgfsetdash{}{0pt}%
\pgfpathmoveto{\pgfqpoint{2.168433in}{1.778817in}}%
\pgfpathlineto{\pgfqpoint{2.167464in}{1.771934in}}%
\pgfpathlineto{\pgfqpoint{2.128088in}{1.771585in}}%
\pgfpathlineto{\pgfqpoint{2.127691in}{1.792284in}}%
\pgfpathlineto{\pgfqpoint{2.135124in}{1.792351in}}%
\pgfpathlineto{\pgfqpoint{2.135067in}{1.806256in}}%
\pgfpathlineto{\pgfqpoint{2.120405in}{1.806180in}}%
\pgfpathlineto{\pgfqpoint{2.120021in}{1.840682in}}%
\pgfpathlineto{\pgfqpoint{2.141258in}{1.840834in}}%
\pgfpathlineto{\pgfqpoint{2.148148in}{1.837483in}}%
\pgfpathlineto{\pgfqpoint{2.148622in}{1.813278in}}%
\pgfpathlineto{\pgfqpoint{2.162319in}{1.813349in}}%
\pgfpathlineto{\pgfqpoint{2.162655in}{1.792587in}}%
\pgfpathlineto{\pgfqpoint{2.169520in}{1.792650in}}%
\pgfpathlineto{\pgfqpoint{2.168433in}{1.778817in}}%
\pgfpathclose%
\pgfusepath{fill}%
\end{pgfscope}%
\begin{pgfscope}%
\pgfpathrectangle{\pgfqpoint{0.100000in}{0.100000in}}{\pgfqpoint{3.608454in}{2.310000in}}%
\pgfusepath{clip}%
\pgfsetbuttcap%
\pgfsetmiterjoin%
\definecolor{currentfill}{rgb}{0.000000,0.454902,0.772549}%
\pgfsetfillcolor{currentfill}%
\pgfsetlinewidth{0.000000pt}%
\definecolor{currentstroke}{rgb}{0.000000,0.000000,0.000000}%
\pgfsetstrokecolor{currentstroke}%
\pgfsetstrokeopacity{0.000000}%
\pgfsetdash{}{0pt}%
\pgfpathmoveto{\pgfqpoint{2.636960in}{1.295013in}}%
\pgfpathlineto{\pgfqpoint{2.634127in}{1.288599in}}%
\pgfpathlineto{\pgfqpoint{2.635552in}{1.279697in}}%
\pgfpathlineto{\pgfqpoint{2.630595in}{1.278625in}}%
\pgfpathlineto{\pgfqpoint{2.627023in}{1.272051in}}%
\pgfpathlineto{\pgfqpoint{2.627054in}{1.277303in}}%
\pgfpathlineto{\pgfqpoint{2.620883in}{1.276667in}}%
\pgfpathlineto{\pgfqpoint{2.617290in}{1.282808in}}%
\pgfpathlineto{\pgfqpoint{2.602387in}{1.273595in}}%
\pgfpathlineto{\pgfqpoint{2.601681in}{1.267043in}}%
\pgfpathlineto{\pgfqpoint{2.586412in}{1.272897in}}%
\pgfpathlineto{\pgfqpoint{2.578231in}{1.275886in}}%
\pgfpathlineto{\pgfqpoint{2.570793in}{1.273710in}}%
\pgfpathlineto{\pgfqpoint{2.567447in}{1.266096in}}%
\pgfpathlineto{\pgfqpoint{2.562848in}{1.271335in}}%
\pgfpathlineto{\pgfqpoint{2.553706in}{1.268477in}}%
\pgfpathlineto{\pgfqpoint{2.546912in}{1.269077in}}%
\pgfpathlineto{\pgfqpoint{2.549388in}{1.262674in}}%
\pgfpathlineto{\pgfqpoint{2.541890in}{1.261426in}}%
\pgfpathlineto{\pgfqpoint{2.537345in}{1.268351in}}%
\pgfpathlineto{\pgfqpoint{2.545025in}{1.287628in}}%
\pgfpathlineto{\pgfqpoint{2.541416in}{1.297901in}}%
\pgfpathlineto{\pgfqpoint{2.543804in}{1.301922in}}%
\pgfpathlineto{\pgfqpoint{2.541523in}{1.307274in}}%
\pgfpathlineto{\pgfqpoint{2.542054in}{1.316771in}}%
\pgfpathlineto{\pgfqpoint{2.544363in}{1.322774in}}%
\pgfpathlineto{\pgfqpoint{2.560427in}{1.323864in}}%
\pgfpathlineto{\pgfqpoint{2.565371in}{1.317560in}}%
\pgfpathlineto{\pgfqpoint{2.572272in}{1.322088in}}%
\pgfpathlineto{\pgfqpoint{2.585800in}{1.324135in}}%
\pgfpathlineto{\pgfqpoint{2.594128in}{1.324280in}}%
\pgfpathlineto{\pgfqpoint{2.608516in}{1.322156in}}%
\pgfpathlineto{\pgfqpoint{2.610839in}{1.324868in}}%
\pgfpathlineto{\pgfqpoint{2.620401in}{1.325720in}}%
\pgfpathlineto{\pgfqpoint{2.621408in}{1.315375in}}%
\pgfpathlineto{\pgfqpoint{2.622478in}{1.305021in}}%
\pgfpathlineto{\pgfqpoint{2.629213in}{1.305901in}}%
\pgfpathlineto{\pgfqpoint{2.629550in}{1.301264in}}%
\pgfpathlineto{\pgfqpoint{2.636434in}{1.301852in}}%
\pgfpathlineto{\pgfqpoint{2.636960in}{1.295013in}}%
\pgfpathclose%
\pgfusepath{fill}%
\end{pgfscope}%
\begin{pgfscope}%
\pgfpathrectangle{\pgfqpoint{0.100000in}{0.100000in}}{\pgfqpoint{3.608454in}{2.310000in}}%
\pgfusepath{clip}%
\pgfsetbuttcap%
\pgfsetmiterjoin%
\definecolor{currentfill}{rgb}{0.000000,0.721569,0.639216}%
\pgfsetfillcolor{currentfill}%
\pgfsetlinewidth{0.000000pt}%
\definecolor{currentstroke}{rgb}{0.000000,0.000000,0.000000}%
\pgfsetstrokecolor{currentstroke}%
\pgfsetstrokeopacity{0.000000}%
\pgfsetdash{}{0pt}%
\pgfpathmoveto{\pgfqpoint{2.295263in}{1.958166in}}%
\pgfpathlineto{\pgfqpoint{2.297058in}{1.910416in}}%
\pgfpathlineto{\pgfqpoint{2.269571in}{1.909501in}}%
\pgfpathlineto{\pgfqpoint{2.256111in}{1.909051in}}%
\pgfpathlineto{\pgfqpoint{2.255541in}{1.929754in}}%
\pgfpathlineto{\pgfqpoint{2.213676in}{1.928846in}}%
\pgfpathlineto{\pgfqpoint{2.212805in}{1.945663in}}%
\pgfpathlineto{\pgfqpoint{2.212580in}{1.977269in}}%
\pgfpathlineto{\pgfqpoint{2.211215in}{2.006566in}}%
\pgfpathlineto{\pgfqpoint{2.211106in}{2.032691in}}%
\pgfpathlineto{\pgfqpoint{2.209745in}{2.046401in}}%
\pgfpathlineto{\pgfqpoint{2.208627in}{2.060466in}}%
\pgfpathlineto{\pgfqpoint{2.208113in}{2.105322in}}%
\pgfpathlineto{\pgfqpoint{2.213604in}{2.105132in}}%
\pgfpathlineto{\pgfqpoint{2.227390in}{2.098684in}}%
\pgfpathlineto{\pgfqpoint{2.232324in}{2.099097in}}%
\pgfpathlineto{\pgfqpoint{2.231425in}{2.090560in}}%
\pgfpathlineto{\pgfqpoint{2.239307in}{2.091665in}}%
\pgfpathlineto{\pgfqpoint{2.247149in}{2.073655in}}%
\pgfpathlineto{\pgfqpoint{2.251836in}{2.075702in}}%
\pgfpathlineto{\pgfqpoint{2.251160in}{2.083595in}}%
\pgfpathlineto{\pgfqpoint{2.263561in}{2.085282in}}%
\pgfpathlineto{\pgfqpoint{2.266399in}{2.077856in}}%
\pgfpathlineto{\pgfqpoint{2.274138in}{2.073437in}}%
\pgfpathlineto{\pgfqpoint{2.282131in}{2.073030in}}%
\pgfpathlineto{\pgfqpoint{2.282543in}{2.066327in}}%
\pgfpathlineto{\pgfqpoint{2.297364in}{2.061677in}}%
\pgfpathlineto{\pgfqpoint{2.307312in}{2.064841in}}%
\pgfpathlineto{\pgfqpoint{2.318616in}{2.073687in}}%
\pgfpathlineto{\pgfqpoint{2.320258in}{2.049681in}}%
\pgfpathlineto{\pgfqpoint{2.321610in}{2.015873in}}%
\pgfpathlineto{\pgfqpoint{2.309782in}{2.000701in}}%
\pgfpathlineto{\pgfqpoint{2.298141in}{1.987730in}}%
\pgfpathlineto{\pgfqpoint{2.293110in}{1.984555in}}%
\pgfpathlineto{\pgfqpoint{2.282209in}{1.972690in}}%
\pgfpathlineto{\pgfqpoint{2.265774in}{1.959901in}}%
\pgfpathlineto{\pgfqpoint{2.269412in}{1.953632in}}%
\pgfpathlineto{\pgfqpoint{2.274053in}{1.951302in}}%
\pgfpathlineto{\pgfqpoint{2.282372in}{1.952765in}}%
\pgfpathlineto{\pgfqpoint{2.295263in}{1.958166in}}%
\pgfpathclose%
\pgfusepath{fill}%
\end{pgfscope}%
\begin{pgfscope}%
\pgfpathrectangle{\pgfqpoint{0.100000in}{0.100000in}}{\pgfqpoint{3.608454in}{2.310000in}}%
\pgfusepath{clip}%
\pgfsetbuttcap%
\pgfsetmiterjoin%
\definecolor{currentfill}{rgb}{0.000000,0.372549,0.813725}%
\pgfsetfillcolor{currentfill}%
\pgfsetlinewidth{0.000000pt}%
\definecolor{currentstroke}{rgb}{0.000000,0.000000,0.000000}%
\pgfsetstrokecolor{currentstroke}%
\pgfsetstrokeopacity{0.000000}%
\pgfsetdash{}{0pt}%
\pgfpathmoveto{\pgfqpoint{2.287673in}{1.541193in}}%
\pgfpathlineto{\pgfqpoint{2.266869in}{1.540448in}}%
\pgfpathlineto{\pgfqpoint{2.264937in}{1.602761in}}%
\pgfpathlineto{\pgfqpoint{2.306014in}{1.604048in}}%
\pgfpathlineto{\pgfqpoint{2.346907in}{1.605656in}}%
\pgfpathlineto{\pgfqpoint{2.348123in}{1.578046in}}%
\pgfpathlineto{\pgfqpoint{2.320640in}{1.576968in}}%
\pgfpathlineto{\pgfqpoint{2.320902in}{1.570069in}}%
\pgfpathlineto{\pgfqpoint{2.293482in}{1.569183in}}%
\pgfpathlineto{\pgfqpoint{2.294578in}{1.541485in}}%
\pgfpathlineto{\pgfqpoint{2.287673in}{1.541193in}}%
\pgfpathclose%
\pgfusepath{fill}%
\end{pgfscope}%
\begin{pgfscope}%
\pgfpathrectangle{\pgfqpoint{0.100000in}{0.100000in}}{\pgfqpoint{3.608454in}{2.310000in}}%
\pgfusepath{clip}%
\pgfsetbuttcap%
\pgfsetmiterjoin%
\definecolor{currentfill}{rgb}{0.000000,0.549020,0.725490}%
\pgfsetfillcolor{currentfill}%
\pgfsetlinewidth{0.000000pt}%
\definecolor{currentstroke}{rgb}{0.000000,0.000000,0.000000}%
\pgfsetstrokecolor{currentstroke}%
\pgfsetstrokeopacity{0.000000}%
\pgfsetdash{}{0pt}%
\pgfpathmoveto{\pgfqpoint{3.039659in}{0.833647in}}%
\pgfpathlineto{\pgfqpoint{3.033775in}{0.835716in}}%
\pgfpathlineto{\pgfqpoint{3.025040in}{0.841844in}}%
\pgfpathlineto{\pgfqpoint{3.020884in}{0.848437in}}%
\pgfpathlineto{\pgfqpoint{3.011213in}{0.848731in}}%
\pgfpathlineto{\pgfqpoint{3.005118in}{0.855085in}}%
\pgfpathlineto{\pgfqpoint{2.996708in}{0.858258in}}%
\pgfpathlineto{\pgfqpoint{2.994468in}{0.861760in}}%
\pgfpathlineto{\pgfqpoint{3.016488in}{0.872188in}}%
\pgfpathlineto{\pgfqpoint{3.014555in}{0.877148in}}%
\pgfpathlineto{\pgfqpoint{3.008209in}{0.883079in}}%
\pgfpathlineto{\pgfqpoint{3.006499in}{0.892231in}}%
\pgfpathlineto{\pgfqpoint{3.015378in}{0.900512in}}%
\pgfpathlineto{\pgfqpoint{3.022585in}{0.898037in}}%
\pgfpathlineto{\pgfqpoint{3.038277in}{0.916697in}}%
\pgfpathlineto{\pgfqpoint{3.047304in}{0.907519in}}%
\pgfpathlineto{\pgfqpoint{3.051321in}{0.914811in}}%
\pgfpathlineto{\pgfqpoint{3.061084in}{0.913825in}}%
\pgfpathlineto{\pgfqpoint{3.066440in}{0.907291in}}%
\pgfpathlineto{\pgfqpoint{3.071225in}{0.906814in}}%
\pgfpathlineto{\pgfqpoint{3.077208in}{0.901258in}}%
\pgfpathlineto{\pgfqpoint{3.081530in}{0.894395in}}%
\pgfpathlineto{\pgfqpoint{3.080811in}{0.888841in}}%
\pgfpathlineto{\pgfqpoint{3.069735in}{0.881587in}}%
\pgfpathlineto{\pgfqpoint{3.064847in}{0.872766in}}%
\pgfpathlineto{\pgfqpoint{3.055250in}{0.864816in}}%
\pgfpathlineto{\pgfqpoint{3.057975in}{0.856666in}}%
\pgfpathlineto{\pgfqpoint{3.043134in}{0.834472in}}%
\pgfpathlineto{\pgfqpoint{3.039659in}{0.833647in}}%
\pgfpathclose%
\pgfusepath{fill}%
\end{pgfscope}%
\begin{pgfscope}%
\pgfpathrectangle{\pgfqpoint{0.100000in}{0.100000in}}{\pgfqpoint{3.608454in}{2.310000in}}%
\pgfusepath{clip}%
\pgfsetbuttcap%
\pgfsetmiterjoin%
\definecolor{currentfill}{rgb}{0.000000,0.909804,0.545098}%
\pgfsetfillcolor{currentfill}%
\pgfsetlinewidth{0.000000pt}%
\definecolor{currentstroke}{rgb}{0.000000,0.000000,0.000000}%
\pgfsetstrokecolor{currentstroke}%
\pgfsetstrokeopacity{0.000000}%
\pgfsetdash{}{0pt}%
\pgfpathmoveto{\pgfqpoint{0.715935in}{1.921940in}}%
\pgfpathlineto{\pgfqpoint{0.762277in}{1.909354in}}%
\pgfpathlineto{\pgfqpoint{0.764102in}{1.916069in}}%
\pgfpathlineto{\pgfqpoint{0.796408in}{1.906884in}}%
\pgfpathlineto{\pgfqpoint{0.790924in}{1.886917in}}%
\pgfpathlineto{\pgfqpoint{0.773761in}{1.820585in}}%
\pgfpathlineto{\pgfqpoint{0.774453in}{1.820402in}}%
\pgfpathlineto{\pgfqpoint{0.761786in}{1.771448in}}%
\pgfpathlineto{\pgfqpoint{0.757203in}{1.749818in}}%
\pgfpathlineto{\pgfqpoint{0.728707in}{1.756911in}}%
\pgfpathlineto{\pgfqpoint{0.693113in}{1.766738in}}%
\pgfpathlineto{\pgfqpoint{0.691075in}{1.767312in}}%
\pgfpathlineto{\pgfqpoint{0.706614in}{1.825006in}}%
\pgfpathlineto{\pgfqpoint{0.674134in}{1.833816in}}%
\pgfpathlineto{\pgfqpoint{0.684062in}{1.866632in}}%
\pgfpathlineto{\pgfqpoint{0.686047in}{1.866076in}}%
\pgfpathlineto{\pgfqpoint{0.695319in}{1.898980in}}%
\pgfpathlineto{\pgfqpoint{0.697108in}{1.905730in}}%
\pgfpathlineto{\pgfqpoint{0.703763in}{1.903845in}}%
\pgfpathlineto{\pgfqpoint{0.709441in}{1.923753in}}%
\pgfpathlineto{\pgfqpoint{0.715935in}{1.921940in}}%
\pgfpathclose%
\pgfusepath{fill}%
\end{pgfscope}%
\begin{pgfscope}%
\pgfpathrectangle{\pgfqpoint{0.100000in}{0.100000in}}{\pgfqpoint{3.608454in}{2.310000in}}%
\pgfusepath{clip}%
\pgfsetbuttcap%
\pgfsetmiterjoin%
\definecolor{currentfill}{rgb}{0.000000,0.568627,0.715686}%
\pgfsetfillcolor{currentfill}%
\pgfsetlinewidth{0.000000pt}%
\definecolor{currentstroke}{rgb}{0.000000,0.000000,0.000000}%
\pgfsetstrokecolor{currentstroke}%
\pgfsetstrokeopacity{0.000000}%
\pgfsetdash{}{0pt}%
\pgfpathmoveto{\pgfqpoint{3.127726in}{0.403873in}}%
\pgfpathlineto{\pgfqpoint{3.099914in}{0.399053in}}%
\pgfpathlineto{\pgfqpoint{3.096468in}{0.420138in}}%
\pgfpathlineto{\pgfqpoint{3.075618in}{0.416692in}}%
\pgfpathlineto{\pgfqpoint{3.072400in}{0.437158in}}%
\pgfpathlineto{\pgfqpoint{3.069493in}{0.458265in}}%
\pgfpathlineto{\pgfqpoint{3.090529in}{0.461360in}}%
\pgfpathlineto{\pgfqpoint{3.089414in}{0.468427in}}%
\pgfpathlineto{\pgfqpoint{3.096465in}{0.469557in}}%
\pgfpathlineto{\pgfqpoint{3.095343in}{0.476565in}}%
\pgfpathlineto{\pgfqpoint{3.109319in}{0.479362in}}%
\pgfpathlineto{\pgfqpoint{3.100740in}{0.491314in}}%
\pgfpathlineto{\pgfqpoint{3.097174in}{0.491133in}}%
\pgfpathlineto{\pgfqpoint{3.088462in}{0.501164in}}%
\pgfpathlineto{\pgfqpoint{3.091930in}{0.511405in}}%
\pgfpathlineto{\pgfqpoint{3.110913in}{0.514275in}}%
\pgfpathlineto{\pgfqpoint{3.117607in}{0.515401in}}%
\pgfpathlineto{\pgfqpoint{3.118731in}{0.508688in}}%
\pgfpathlineto{\pgfqpoint{3.150978in}{0.513727in}}%
\pgfpathlineto{\pgfqpoint{3.156513in}{0.502751in}}%
\pgfpathlineto{\pgfqpoint{3.167005in}{0.486501in}}%
\pgfpathlineto{\pgfqpoint{3.173443in}{0.474576in}}%
\pgfpathlineto{\pgfqpoint{3.181412in}{0.456522in}}%
\pgfpathlineto{\pgfqpoint{3.183733in}{0.440357in}}%
\pgfpathlineto{\pgfqpoint{3.184616in}{0.418072in}}%
\pgfpathlineto{\pgfqpoint{3.173551in}{0.417385in}}%
\pgfpathlineto{\pgfqpoint{3.126693in}{0.409803in}}%
\pgfpathlineto{\pgfqpoint{3.127726in}{0.403873in}}%
\pgfpathclose%
\pgfusepath{fill}%
\end{pgfscope}%
\begin{pgfscope}%
\pgfpathrectangle{\pgfqpoint{0.100000in}{0.100000in}}{\pgfqpoint{3.608454in}{2.310000in}}%
\pgfusepath{clip}%
\pgfsetbuttcap%
\pgfsetmiterjoin%
\definecolor{currentfill}{rgb}{0.000000,0.521569,0.739216}%
\pgfsetfillcolor{currentfill}%
\pgfsetlinewidth{0.000000pt}%
\definecolor{currentstroke}{rgb}{0.000000,0.000000,0.000000}%
\pgfsetstrokecolor{currentstroke}%
\pgfsetstrokeopacity{0.000000}%
\pgfsetdash{}{0pt}%
\pgfpathmoveto{\pgfqpoint{2.621408in}{1.315375in}}%
\pgfpathlineto{\pgfqpoint{2.620401in}{1.325720in}}%
\pgfpathlineto{\pgfqpoint{2.610839in}{1.324868in}}%
\pgfpathlineto{\pgfqpoint{2.608516in}{1.322156in}}%
\pgfpathlineto{\pgfqpoint{2.594128in}{1.324280in}}%
\pgfpathlineto{\pgfqpoint{2.585800in}{1.324135in}}%
\pgfpathlineto{\pgfqpoint{2.582510in}{1.335565in}}%
\pgfpathlineto{\pgfqpoint{2.584141in}{1.345785in}}%
\pgfpathlineto{\pgfqpoint{2.590013in}{1.348185in}}%
\pgfpathlineto{\pgfqpoint{2.590232in}{1.353266in}}%
\pgfpathlineto{\pgfqpoint{2.583434in}{1.352723in}}%
\pgfpathlineto{\pgfqpoint{2.581657in}{1.373507in}}%
\pgfpathlineto{\pgfqpoint{2.593047in}{1.374201in}}%
\pgfpathlineto{\pgfqpoint{2.591897in}{1.387949in}}%
\pgfpathlineto{\pgfqpoint{2.598672in}{1.388501in}}%
\pgfpathlineto{\pgfqpoint{2.597905in}{1.398852in}}%
\pgfpathlineto{\pgfqpoint{2.613308in}{1.399991in}}%
\pgfpathlineto{\pgfqpoint{2.616632in}{1.400311in}}%
\pgfpathlineto{\pgfqpoint{2.617652in}{1.390119in}}%
\pgfpathlineto{\pgfqpoint{2.632842in}{1.391462in}}%
\pgfpathlineto{\pgfqpoint{2.635715in}{1.368602in}}%
\pgfpathlineto{\pgfqpoint{2.638885in}{1.368939in}}%
\pgfpathlineto{\pgfqpoint{2.641627in}{1.364614in}}%
\pgfpathlineto{\pgfqpoint{2.643662in}{1.346766in}}%
\pgfpathlineto{\pgfqpoint{2.641760in}{1.344149in}}%
\pgfpathlineto{\pgfqpoint{2.644392in}{1.317477in}}%
\pgfpathlineto{\pgfqpoint{2.621408in}{1.315375in}}%
\pgfpathclose%
\pgfusepath{fill}%
\end{pgfscope}%
\begin{pgfscope}%
\pgfpathrectangle{\pgfqpoint{0.100000in}{0.100000in}}{\pgfqpoint{3.608454in}{2.310000in}}%
\pgfusepath{clip}%
\pgfsetbuttcap%
\pgfsetmiterjoin%
\definecolor{currentfill}{rgb}{0.000000,0.905882,0.547059}%
\pgfsetfillcolor{currentfill}%
\pgfsetlinewidth{0.000000pt}%
\definecolor{currentstroke}{rgb}{0.000000,0.000000,0.000000}%
\pgfsetstrokecolor{currentstroke}%
\pgfsetstrokeopacity{0.000000}%
\pgfsetdash{}{0pt}%
\pgfpathmoveto{\pgfqpoint{2.258858in}{0.847301in}}%
\pgfpathlineto{\pgfqpoint{2.267654in}{0.847225in}}%
\pgfpathlineto{\pgfqpoint{2.272892in}{0.840599in}}%
\pgfpathlineto{\pgfqpoint{2.279806in}{0.840743in}}%
\pgfpathlineto{\pgfqpoint{2.280375in}{0.819890in}}%
\pgfpathlineto{\pgfqpoint{2.287535in}{0.813168in}}%
\pgfpathlineto{\pgfqpoint{2.288408in}{0.781863in}}%
\pgfpathlineto{\pgfqpoint{2.285332in}{0.771355in}}%
\pgfpathlineto{\pgfqpoint{2.267964in}{0.770915in}}%
\pgfpathlineto{\pgfqpoint{2.268095in}{0.763939in}}%
\pgfpathlineto{\pgfqpoint{2.244140in}{0.763307in}}%
\pgfpathlineto{\pgfqpoint{2.246031in}{0.774913in}}%
\pgfpathlineto{\pgfqpoint{2.249736in}{0.781810in}}%
\pgfpathlineto{\pgfqpoint{2.245714in}{0.798315in}}%
\pgfpathlineto{\pgfqpoint{2.254067in}{0.798428in}}%
\pgfpathlineto{\pgfqpoint{2.256664in}{0.805503in}}%
\pgfpathlineto{\pgfqpoint{2.255994in}{0.822878in}}%
\pgfpathlineto{\pgfqpoint{2.249143in}{0.822767in}}%
\pgfpathlineto{\pgfqpoint{2.248708in}{0.834367in}}%
\pgfpathlineto{\pgfqpoint{2.251967in}{0.847074in}}%
\pgfpathlineto{\pgfqpoint{2.258858in}{0.847301in}}%
\pgfpathclose%
\pgfusepath{fill}%
\end{pgfscope}%
\begin{pgfscope}%
\pgfpathrectangle{\pgfqpoint{0.100000in}{0.100000in}}{\pgfqpoint{3.608454in}{2.310000in}}%
\pgfusepath{clip}%
\pgfsetbuttcap%
\pgfsetmiterjoin%
\definecolor{currentfill}{rgb}{0.000000,0.603922,0.698039}%
\pgfsetfillcolor{currentfill}%
\pgfsetlinewidth{0.000000pt}%
\definecolor{currentstroke}{rgb}{0.000000,0.000000,0.000000}%
\pgfsetstrokecolor{currentstroke}%
\pgfsetstrokeopacity{0.000000}%
\pgfsetdash{}{0pt}%
\pgfpathmoveto{\pgfqpoint{1.291011in}{1.784895in}}%
\pgfpathlineto{\pgfqpoint{1.297605in}{1.774698in}}%
\pgfpathlineto{\pgfqpoint{1.296615in}{1.768471in}}%
\pgfpathlineto{\pgfqpoint{1.310078in}{1.766325in}}%
\pgfpathlineto{\pgfqpoint{1.323149in}{1.760841in}}%
\pgfpathlineto{\pgfqpoint{1.322608in}{1.757434in}}%
\pgfpathlineto{\pgfqpoint{1.342072in}{1.748707in}}%
\pgfpathlineto{\pgfqpoint{1.352534in}{1.748041in}}%
\pgfpathlineto{\pgfqpoint{1.384511in}{1.743611in}}%
\pgfpathlineto{\pgfqpoint{1.388262in}{1.745418in}}%
\pgfpathlineto{\pgfqpoint{1.388966in}{1.743023in}}%
\pgfpathlineto{\pgfqpoint{1.386059in}{1.715853in}}%
\pgfpathlineto{\pgfqpoint{1.382345in}{1.688646in}}%
\pgfpathlineto{\pgfqpoint{1.379922in}{1.688990in}}%
\pgfpathlineto{\pgfqpoint{1.376226in}{1.661744in}}%
\pgfpathlineto{\pgfqpoint{1.377332in}{1.661590in}}%
\pgfpathlineto{\pgfqpoint{1.375467in}{1.648053in}}%
\pgfpathlineto{\pgfqpoint{1.347746in}{1.651996in}}%
\pgfpathlineto{\pgfqpoint{1.287325in}{1.661183in}}%
\pgfpathlineto{\pgfqpoint{1.287915in}{1.674697in}}%
\pgfpathlineto{\pgfqpoint{1.290897in}{1.695072in}}%
\pgfpathlineto{\pgfqpoint{1.282465in}{1.703069in}}%
\pgfpathlineto{\pgfqpoint{1.275945in}{1.717107in}}%
\pgfpathlineto{\pgfqpoint{1.270191in}{1.723778in}}%
\pgfpathlineto{\pgfqpoint{1.264849in}{1.733953in}}%
\pgfpathlineto{\pgfqpoint{1.262997in}{1.742086in}}%
\pgfpathlineto{\pgfqpoint{1.263755in}{1.753309in}}%
\pgfpathlineto{\pgfqpoint{1.261563in}{1.761526in}}%
\pgfpathlineto{\pgfqpoint{1.244733in}{1.764439in}}%
\pgfpathlineto{\pgfqpoint{1.251746in}{1.806946in}}%
\pgfpathlineto{\pgfqpoint{1.254530in}{1.801740in}}%
\pgfpathlineto{\pgfqpoint{1.260477in}{1.800714in}}%
\pgfpathlineto{\pgfqpoint{1.264050in}{1.788274in}}%
\pgfpathlineto{\pgfqpoint{1.271394in}{1.792453in}}%
\pgfpathlineto{\pgfqpoint{1.273821in}{1.796453in}}%
\pgfpathlineto{\pgfqpoint{1.279919in}{1.798565in}}%
\pgfpathlineto{\pgfqpoint{1.285167in}{1.795237in}}%
\pgfpathlineto{\pgfqpoint{1.291011in}{1.784895in}}%
\pgfpathclose%
\pgfusepath{fill}%
\end{pgfscope}%
\begin{pgfscope}%
\pgfpathrectangle{\pgfqpoint{0.100000in}{0.100000in}}{\pgfqpoint{3.608454in}{2.310000in}}%
\pgfusepath{clip}%
\pgfsetbuttcap%
\pgfsetmiterjoin%
\definecolor{currentfill}{rgb}{0.000000,0.588235,0.705882}%
\pgfsetfillcolor{currentfill}%
\pgfsetlinewidth{0.000000pt}%
\definecolor{currentstroke}{rgb}{0.000000,0.000000,0.000000}%
\pgfsetstrokecolor{currentstroke}%
\pgfsetstrokeopacity{0.000000}%
\pgfsetdash{}{0pt}%
\pgfpathmoveto{\pgfqpoint{2.417621in}{1.861430in}}%
\pgfpathlineto{\pgfqpoint{2.417000in}{1.868413in}}%
\pgfpathlineto{\pgfqpoint{2.382793in}{1.866283in}}%
\pgfpathlineto{\pgfqpoint{2.380907in}{1.900334in}}%
\pgfpathlineto{\pgfqpoint{2.387268in}{1.900687in}}%
\pgfpathlineto{\pgfqpoint{2.385839in}{1.925963in}}%
\pgfpathlineto{\pgfqpoint{2.432701in}{1.915926in}}%
\pgfpathlineto{\pgfqpoint{2.438472in}{1.913015in}}%
\pgfpathlineto{\pgfqpoint{2.456551in}{1.904950in}}%
\pgfpathlineto{\pgfqpoint{2.457894in}{1.884450in}}%
\pgfpathlineto{\pgfqpoint{2.471787in}{1.885405in}}%
\pgfpathlineto{\pgfqpoint{2.473620in}{1.857983in}}%
\pgfpathlineto{\pgfqpoint{2.459694in}{1.857127in}}%
\pgfpathlineto{\pgfqpoint{2.445861in}{1.856178in}}%
\pgfpathlineto{\pgfqpoint{2.445567in}{1.863054in}}%
\pgfpathlineto{\pgfqpoint{2.417621in}{1.861430in}}%
\pgfpathclose%
\pgfusepath{fill}%
\end{pgfscope}%
\begin{pgfscope}%
\pgfpathrectangle{\pgfqpoint{0.100000in}{0.100000in}}{\pgfqpoint{3.608454in}{2.310000in}}%
\pgfusepath{clip}%
\pgfsetbuttcap%
\pgfsetmiterjoin%
\definecolor{currentfill}{rgb}{0.000000,0.701961,0.649020}%
\pgfsetfillcolor{currentfill}%
\pgfsetlinewidth{0.000000pt}%
\definecolor{currentstroke}{rgb}{0.000000,0.000000,0.000000}%
\pgfsetstrokecolor{currentstroke}%
\pgfsetstrokeopacity{0.000000}%
\pgfsetdash{}{0pt}%
\pgfpathmoveto{\pgfqpoint{1.126068in}{1.701146in}}%
\pgfpathlineto{\pgfqpoint{1.122487in}{1.688955in}}%
\pgfpathlineto{\pgfqpoint{1.117899in}{1.692151in}}%
\pgfpathlineto{\pgfqpoint{1.109727in}{1.693683in}}%
\pgfpathlineto{\pgfqpoint{1.101946in}{1.704218in}}%
\pgfpathlineto{\pgfqpoint{1.098147in}{1.713450in}}%
\pgfpathlineto{\pgfqpoint{1.091640in}{1.714726in}}%
\pgfpathlineto{\pgfqpoint{1.090342in}{1.707956in}}%
\pgfpathlineto{\pgfqpoint{1.081382in}{1.709710in}}%
\pgfpathlineto{\pgfqpoint{1.080051in}{1.702943in}}%
\pgfpathlineto{\pgfqpoint{1.060120in}{1.706792in}}%
\pgfpathlineto{\pgfqpoint{1.064066in}{1.727041in}}%
\pgfpathlineto{\pgfqpoint{1.054206in}{1.729137in}}%
\pgfpathlineto{\pgfqpoint{1.051274in}{1.732636in}}%
\pgfpathlineto{\pgfqpoint{1.055372in}{1.735866in}}%
\pgfpathlineto{\pgfqpoint{1.053439in}{1.743275in}}%
\pgfpathlineto{\pgfqpoint{1.058573in}{1.770425in}}%
\pgfpathlineto{\pgfqpoint{1.071945in}{1.767715in}}%
\pgfpathlineto{\pgfqpoint{1.074676in}{1.781213in}}%
\pgfpathlineto{\pgfqpoint{1.054623in}{1.785268in}}%
\pgfpathlineto{\pgfqpoint{1.055937in}{1.791680in}}%
\pgfpathlineto{\pgfqpoint{1.040506in}{1.794892in}}%
\pgfpathlineto{\pgfqpoint{1.041894in}{1.801489in}}%
\pgfpathlineto{\pgfqpoint{1.038897in}{1.813232in}}%
\pgfpathlineto{\pgfqpoint{1.034613in}{1.812423in}}%
\pgfpathlineto{\pgfqpoint{1.037674in}{1.816193in}}%
\pgfpathlineto{\pgfqpoint{1.049281in}{1.819620in}}%
\pgfpathlineto{\pgfqpoint{1.057180in}{1.823892in}}%
\pgfpathlineto{\pgfqpoint{1.061051in}{1.840480in}}%
\pgfpathlineto{\pgfqpoint{1.063909in}{1.846064in}}%
\pgfpathlineto{\pgfqpoint{1.066665in}{1.859594in}}%
\pgfpathlineto{\pgfqpoint{1.100225in}{1.852591in}}%
\pgfpathlineto{\pgfqpoint{1.102494in}{1.863989in}}%
\pgfpathlineto{\pgfqpoint{1.105971in}{1.872044in}}%
\pgfpathlineto{\pgfqpoint{1.123023in}{1.868301in}}%
\pgfpathlineto{\pgfqpoint{1.127333in}{1.864844in}}%
\pgfpathlineto{\pgfqpoint{1.130114in}{1.870874in}}%
\pgfpathlineto{\pgfqpoint{1.134588in}{1.873151in}}%
\pgfpathlineto{\pgfqpoint{1.144630in}{1.867633in}}%
\pgfpathlineto{\pgfqpoint{1.157878in}{1.868529in}}%
\pgfpathlineto{\pgfqpoint{1.159714in}{1.863660in}}%
\pgfpathlineto{\pgfqpoint{1.166980in}{1.866417in}}%
\pgfpathlineto{\pgfqpoint{1.177604in}{1.866287in}}%
\pgfpathlineto{\pgfqpoint{1.181983in}{1.875682in}}%
\pgfpathlineto{\pgfqpoint{1.187877in}{1.878039in}}%
\pgfpathlineto{\pgfqpoint{1.191695in}{1.873351in}}%
\pgfpathlineto{\pgfqpoint{1.198018in}{1.855881in}}%
\pgfpathlineto{\pgfqpoint{1.202252in}{1.853038in}}%
\pgfpathlineto{\pgfqpoint{1.195518in}{1.814680in}}%
\pgfpathlineto{\pgfqpoint{1.186131in}{1.812373in}}%
\pgfpathlineto{\pgfqpoint{1.174897in}{1.813546in}}%
\pgfpathlineto{\pgfqpoint{1.170551in}{1.790061in}}%
\pgfpathlineto{\pgfqpoint{1.179359in}{1.788436in}}%
\pgfpathlineto{\pgfqpoint{1.180737in}{1.781578in}}%
\pgfpathlineto{\pgfqpoint{1.188888in}{1.777013in}}%
\pgfpathlineto{\pgfqpoint{1.186356in}{1.762535in}}%
\pgfpathlineto{\pgfqpoint{1.182268in}{1.739464in}}%
\pgfpathlineto{\pgfqpoint{1.124259in}{1.750336in}}%
\pgfpathlineto{\pgfqpoint{1.117477in}{1.741580in}}%
\pgfpathlineto{\pgfqpoint{1.116145in}{1.733195in}}%
\pgfpathlineto{\pgfqpoint{1.117055in}{1.725647in}}%
\pgfpathlineto{\pgfqpoint{1.125111in}{1.725638in}}%
\pgfpathlineto{\pgfqpoint{1.126636in}{1.717223in}}%
\pgfpathlineto{\pgfqpoint{1.126068in}{1.701146in}}%
\pgfpathclose%
\pgfusepath{fill}%
\end{pgfscope}%
\begin{pgfscope}%
\pgfpathrectangle{\pgfqpoint{0.100000in}{0.100000in}}{\pgfqpoint{3.608454in}{2.310000in}}%
\pgfusepath{clip}%
\pgfsetbuttcap%
\pgfsetmiterjoin%
\definecolor{currentfill}{rgb}{0.000000,0.701961,0.649020}%
\pgfsetfillcolor{currentfill}%
\pgfsetlinewidth{0.000000pt}%
\definecolor{currentstroke}{rgb}{0.000000,0.000000,0.000000}%
\pgfsetstrokecolor{currentstroke}%
\pgfsetstrokeopacity{0.000000}%
\pgfsetdash{}{0pt}%
\pgfpathmoveto{\pgfqpoint{0.718104in}{2.222638in}}%
\pgfpathlineto{\pgfqpoint{0.722390in}{2.214825in}}%
\pgfpathlineto{\pgfqpoint{0.724213in}{2.207197in}}%
\pgfpathlineto{\pgfqpoint{0.733242in}{2.201231in}}%
\pgfpathlineto{\pgfqpoint{0.735660in}{2.196356in}}%
\pgfpathlineto{\pgfqpoint{0.742925in}{2.192347in}}%
\pgfpathlineto{\pgfqpoint{0.748201in}{2.185881in}}%
\pgfpathlineto{\pgfqpoint{0.763587in}{2.181602in}}%
\pgfpathlineto{\pgfqpoint{0.767233in}{2.177075in}}%
\pgfpathlineto{\pgfqpoint{0.762473in}{2.166902in}}%
\pgfpathlineto{\pgfqpoint{0.763393in}{2.155197in}}%
\pgfpathlineto{\pgfqpoint{0.762232in}{2.143437in}}%
\pgfpathlineto{\pgfqpoint{0.758684in}{2.134138in}}%
\pgfpathlineto{\pgfqpoint{0.761418in}{2.129669in}}%
\pgfpathlineto{\pgfqpoint{0.749167in}{2.084439in}}%
\pgfpathlineto{\pgfqpoint{0.719253in}{2.092888in}}%
\pgfpathlineto{\pgfqpoint{0.661241in}{2.110265in}}%
\pgfpathlineto{\pgfqpoint{0.669147in}{2.136525in}}%
\pgfpathlineto{\pgfqpoint{0.675514in}{2.134601in}}%
\pgfpathlineto{\pgfqpoint{0.675711in}{2.144078in}}%
\pgfpathlineto{\pgfqpoint{0.682016in}{2.153295in}}%
\pgfpathlineto{\pgfqpoint{0.682663in}{2.166183in}}%
\pgfpathlineto{\pgfqpoint{0.680298in}{2.173569in}}%
\pgfpathlineto{\pgfqpoint{0.697139in}{2.190297in}}%
\pgfpathlineto{\pgfqpoint{0.699138in}{2.201963in}}%
\pgfpathlineto{\pgfqpoint{0.694658in}{2.204301in}}%
\pgfpathlineto{\pgfqpoint{0.694418in}{2.209381in}}%
\pgfpathlineto{\pgfqpoint{0.701083in}{2.215781in}}%
\pgfpathlineto{\pgfqpoint{0.712161in}{2.222097in}}%
\pgfpathlineto{\pgfqpoint{0.718104in}{2.222638in}}%
\pgfpathclose%
\pgfusepath{fill}%
\end{pgfscope}%
\begin{pgfscope}%
\pgfpathrectangle{\pgfqpoint{0.100000in}{0.100000in}}{\pgfqpoint{3.608454in}{2.310000in}}%
\pgfusepath{clip}%
\pgfsetbuttcap%
\pgfsetmiterjoin%
\definecolor{currentfill}{rgb}{0.000000,0.776471,0.611765}%
\pgfsetfillcolor{currentfill}%
\pgfsetlinewidth{0.000000pt}%
\definecolor{currentstroke}{rgb}{0.000000,0.000000,0.000000}%
\pgfsetstrokecolor{currentstroke}%
\pgfsetstrokeopacity{0.000000}%
\pgfsetdash{}{0pt}%
\pgfpathmoveto{\pgfqpoint{0.538914in}{1.596597in}}%
\pgfpathlineto{\pgfqpoint{0.532464in}{1.597564in}}%
\pgfpathlineto{\pgfqpoint{0.520214in}{1.588877in}}%
\pgfpathlineto{\pgfqpoint{0.514461in}{1.575916in}}%
\pgfpathlineto{\pgfqpoint{0.507112in}{1.573994in}}%
\pgfpathlineto{\pgfqpoint{0.502169in}{1.569148in}}%
\pgfpathlineto{\pgfqpoint{0.498193in}{1.555922in}}%
\pgfpathlineto{\pgfqpoint{0.489967in}{1.557121in}}%
\pgfpathlineto{\pgfqpoint{0.484414in}{1.564406in}}%
\pgfpathlineto{\pgfqpoint{0.485256in}{1.568714in}}%
\pgfpathlineto{\pgfqpoint{0.480562in}{1.575547in}}%
\pgfpathlineto{\pgfqpoint{0.450911in}{1.584637in}}%
\pgfpathlineto{\pgfqpoint{0.448894in}{1.592474in}}%
\pgfpathlineto{\pgfqpoint{0.445039in}{1.597163in}}%
\pgfpathlineto{\pgfqpoint{0.447344in}{1.609339in}}%
\pgfpathlineto{\pgfqpoint{0.440654in}{1.612562in}}%
\pgfpathlineto{\pgfqpoint{0.435002in}{1.619433in}}%
\pgfpathlineto{\pgfqpoint{0.440384in}{1.631531in}}%
\pgfpathlineto{\pgfqpoint{0.443292in}{1.641381in}}%
\pgfpathlineto{\pgfqpoint{0.434583in}{1.644086in}}%
\pgfpathlineto{\pgfqpoint{0.437233in}{1.653925in}}%
\pgfpathlineto{\pgfqpoint{0.436761in}{1.661477in}}%
\pgfpathlineto{\pgfqpoint{0.488413in}{1.645315in}}%
\pgfpathlineto{\pgfqpoint{0.490498in}{1.651799in}}%
\pgfpathlineto{\pgfqpoint{0.504486in}{1.647574in}}%
\pgfpathlineto{\pgfqpoint{0.512624in}{1.653371in}}%
\pgfpathlineto{\pgfqpoint{0.515901in}{1.652222in}}%
\pgfpathlineto{\pgfqpoint{0.522050in}{1.660114in}}%
\pgfpathlineto{\pgfqpoint{0.531886in}{1.661402in}}%
\pgfpathlineto{\pgfqpoint{0.534394in}{1.657082in}}%
\pgfpathlineto{\pgfqpoint{0.529796in}{1.650793in}}%
\pgfpathlineto{\pgfqpoint{0.526964in}{1.639730in}}%
\pgfpathlineto{\pgfqpoint{0.529553in}{1.637642in}}%
\pgfpathlineto{\pgfqpoint{0.540950in}{1.606159in}}%
\pgfpathlineto{\pgfqpoint{0.538914in}{1.596597in}}%
\pgfpathclose%
\pgfusepath{fill}%
\end{pgfscope}%
\begin{pgfscope}%
\pgfpathrectangle{\pgfqpoint{0.100000in}{0.100000in}}{\pgfqpoint{3.608454in}{2.310000in}}%
\pgfusepath{clip}%
\pgfsetbuttcap%
\pgfsetmiterjoin%
\definecolor{currentfill}{rgb}{0.000000,0.596078,0.701961}%
\pgfsetfillcolor{currentfill}%
\pgfsetlinewidth{0.000000pt}%
\definecolor{currentstroke}{rgb}{0.000000,0.000000,0.000000}%
\pgfsetstrokecolor{currentstroke}%
\pgfsetstrokeopacity{0.000000}%
\pgfsetdash{}{0pt}%
\pgfpathmoveto{\pgfqpoint{2.075326in}{1.295024in}}%
\pgfpathlineto{\pgfqpoint{2.074465in}{1.263717in}}%
\pgfpathlineto{\pgfqpoint{2.047067in}{1.264108in}}%
\pgfpathlineto{\pgfqpoint{2.047190in}{1.274443in}}%
\pgfpathlineto{\pgfqpoint{2.022450in}{1.274884in}}%
\pgfpathlineto{\pgfqpoint{2.022297in}{1.268017in}}%
\pgfpathlineto{\pgfqpoint{2.012052in}{1.268224in}}%
\pgfpathlineto{\pgfqpoint{1.992256in}{1.268523in}}%
\pgfpathlineto{\pgfqpoint{1.992866in}{1.289314in}}%
\pgfpathlineto{\pgfqpoint{1.994270in}{1.303029in}}%
\pgfpathlineto{\pgfqpoint{2.023077in}{1.302462in}}%
\pgfpathlineto{\pgfqpoint{2.023407in}{1.319659in}}%
\pgfpathlineto{\pgfqpoint{2.048478in}{1.319317in}}%
\pgfpathlineto{\pgfqpoint{2.047953in}{1.295247in}}%
\pgfpathlineto{\pgfqpoint{2.075326in}{1.295024in}}%
\pgfpathclose%
\pgfusepath{fill}%
\end{pgfscope}%
\begin{pgfscope}%
\pgfpathrectangle{\pgfqpoint{0.100000in}{0.100000in}}{\pgfqpoint{3.608454in}{2.310000in}}%
\pgfusepath{clip}%
\pgfsetbuttcap%
\pgfsetmiterjoin%
\definecolor{currentfill}{rgb}{0.000000,0.462745,0.768627}%
\pgfsetfillcolor{currentfill}%
\pgfsetlinewidth{0.000000pt}%
\definecolor{currentstroke}{rgb}{0.000000,0.000000,0.000000}%
\pgfsetstrokecolor{currentstroke}%
\pgfsetstrokeopacity{0.000000}%
\pgfsetdash{}{0pt}%
\pgfpathmoveto{\pgfqpoint{1.303565in}{2.014897in}}%
\pgfpathlineto{\pgfqpoint{1.299591in}{2.015508in}}%
\pgfpathlineto{\pgfqpoint{1.292244in}{2.021526in}}%
\pgfpathlineto{\pgfqpoint{1.264971in}{2.020151in}}%
\pgfpathlineto{\pgfqpoint{1.259993in}{2.025470in}}%
\pgfpathlineto{\pgfqpoint{1.251605in}{2.040078in}}%
\pgfpathlineto{\pgfqpoint{1.244581in}{2.048177in}}%
\pgfpathlineto{\pgfqpoint{1.231866in}{2.052930in}}%
\pgfpathlineto{\pgfqpoint{1.211322in}{2.055765in}}%
\pgfpathlineto{\pgfqpoint{1.210055in}{2.048940in}}%
\pgfpathlineto{\pgfqpoint{1.203686in}{2.050112in}}%
\pgfpathlineto{\pgfqpoint{1.196622in}{2.051472in}}%
\pgfpathlineto{\pgfqpoint{1.199816in}{2.068342in}}%
\pgfpathlineto{\pgfqpoint{1.192015in}{2.070842in}}%
\pgfpathlineto{\pgfqpoint{1.187119in}{2.075916in}}%
\pgfpathlineto{\pgfqpoint{1.191810in}{2.101317in}}%
\pgfpathlineto{\pgfqpoint{1.179403in}{2.102869in}}%
\pgfpathlineto{\pgfqpoint{1.173331in}{2.107985in}}%
\pgfpathlineto{\pgfqpoint{1.173281in}{2.111946in}}%
\pgfpathlineto{\pgfqpoint{1.165476in}{2.115944in}}%
\pgfpathlineto{\pgfqpoint{1.152365in}{2.116451in}}%
\pgfpathlineto{\pgfqpoint{1.148672in}{2.124295in}}%
\pgfpathlineto{\pgfqpoint{1.150853in}{2.134585in}}%
\pgfpathlineto{\pgfqpoint{1.149026in}{2.145360in}}%
\pgfpathlineto{\pgfqpoint{1.154266in}{2.150464in}}%
\pgfpathlineto{\pgfqpoint{1.150149in}{2.159708in}}%
\pgfpathlineto{\pgfqpoint{1.194140in}{2.150958in}}%
\pgfpathlineto{\pgfqpoint{1.200330in}{2.146187in}}%
\pgfpathlineto{\pgfqpoint{1.202148in}{2.137522in}}%
\pgfpathlineto{\pgfqpoint{1.232482in}{2.132070in}}%
\pgfpathlineto{\pgfqpoint{1.234540in}{2.143503in}}%
\pgfpathlineto{\pgfqpoint{1.265169in}{2.138097in}}%
\pgfpathlineto{\pgfqpoint{1.266372in}{2.144919in}}%
\pgfpathlineto{\pgfqpoint{1.276965in}{2.143064in}}%
\pgfpathlineto{\pgfqpoint{1.278167in}{2.149917in}}%
\pgfpathlineto{\pgfqpoint{1.319123in}{2.142944in}}%
\pgfpathlineto{\pgfqpoint{1.318309in}{2.136230in}}%
\pgfpathlineto{\pgfqpoint{1.324278in}{2.135234in}}%
\pgfpathlineto{\pgfqpoint{1.323124in}{2.128216in}}%
\pgfpathlineto{\pgfqpoint{1.333325in}{2.126759in}}%
\pgfpathlineto{\pgfqpoint{1.328054in}{2.095490in}}%
\pgfpathlineto{\pgfqpoint{1.312471in}{2.096193in}}%
\pgfpathlineto{\pgfqpoint{1.309109in}{2.082641in}}%
\pgfpathlineto{\pgfqpoint{1.293100in}{2.076361in}}%
\pgfpathlineto{\pgfqpoint{1.290585in}{2.063837in}}%
\pgfpathlineto{\pgfqpoint{1.296843in}{2.059295in}}%
\pgfpathlineto{\pgfqpoint{1.303684in}{2.058169in}}%
\pgfpathlineto{\pgfqpoint{1.309306in}{2.053731in}}%
\pgfpathlineto{\pgfqpoint{1.305651in}{2.030921in}}%
\pgfpathlineto{\pgfqpoint{1.303017in}{2.028995in}}%
\pgfpathlineto{\pgfqpoint{1.303565in}{2.014897in}}%
\pgfpathclose%
\pgfusepath{fill}%
\end{pgfscope}%
\begin{pgfscope}%
\pgfpathrectangle{\pgfqpoint{0.100000in}{0.100000in}}{\pgfqpoint{3.608454in}{2.310000in}}%
\pgfusepath{clip}%
\pgfsetbuttcap%
\pgfsetmiterjoin%
\definecolor{currentfill}{rgb}{0.000000,0.462745,0.768627}%
\pgfsetfillcolor{currentfill}%
\pgfsetlinewidth{0.000000pt}%
\definecolor{currentstroke}{rgb}{0.000000,0.000000,0.000000}%
\pgfsetstrokecolor{currentstroke}%
\pgfsetstrokeopacity{0.000000}%
\pgfsetdash{}{0pt}%
\pgfpathmoveto{\pgfqpoint{1.807069in}{2.050968in}}%
\pgfpathlineto{\pgfqpoint{1.827828in}{2.049744in}}%
\pgfpathlineto{\pgfqpoint{1.826539in}{2.063794in}}%
\pgfpathlineto{\pgfqpoint{1.828130in}{2.091731in}}%
\pgfpathlineto{\pgfqpoint{1.826839in}{2.105738in}}%
\pgfpathlineto{\pgfqpoint{1.833772in}{2.105376in}}%
\pgfpathlineto{\pgfqpoint{1.868375in}{2.103508in}}%
\pgfpathlineto{\pgfqpoint{1.867703in}{2.089585in}}%
\pgfpathlineto{\pgfqpoint{1.848899in}{2.090580in}}%
\pgfpathlineto{\pgfqpoint{1.847326in}{2.062601in}}%
\pgfpathlineto{\pgfqpoint{1.848523in}{2.048601in}}%
\pgfpathlineto{\pgfqpoint{1.876114in}{2.047161in}}%
\pgfpathlineto{\pgfqpoint{1.877148in}{2.033168in}}%
\pgfpathlineto{\pgfqpoint{1.875835in}{2.005576in}}%
\pgfpathlineto{\pgfqpoint{1.864232in}{2.006135in}}%
\pgfpathlineto{\pgfqpoint{1.830021in}{2.007934in}}%
\pgfpathlineto{\pgfqpoint{1.799746in}{2.009759in}}%
\pgfpathlineto{\pgfqpoint{1.801529in}{2.037408in}}%
\pgfpathlineto{\pgfqpoint{1.806219in}{2.037109in}}%
\pgfpathlineto{\pgfqpoint{1.807069in}{2.050968in}}%
\pgfpathclose%
\pgfusepath{fill}%
\end{pgfscope}%
\begin{pgfscope}%
\pgfpathrectangle{\pgfqpoint{0.100000in}{0.100000in}}{\pgfqpoint{3.608454in}{2.310000in}}%
\pgfusepath{clip}%
\pgfsetbuttcap%
\pgfsetmiterjoin%
\definecolor{currentfill}{rgb}{0.000000,0.580392,0.709804}%
\pgfsetfillcolor{currentfill}%
\pgfsetlinewidth{0.000000pt}%
\definecolor{currentstroke}{rgb}{0.000000,0.000000,0.000000}%
\pgfsetstrokecolor{currentstroke}%
\pgfsetstrokeopacity{0.000000}%
\pgfsetdash{}{0pt}%
\pgfpathmoveto{\pgfqpoint{2.409257in}{1.435587in}}%
\pgfpathlineto{\pgfqpoint{2.397384in}{1.430530in}}%
\pgfpathlineto{\pgfqpoint{2.391369in}{1.431061in}}%
\pgfpathlineto{\pgfqpoint{2.389932in}{1.435148in}}%
\pgfpathlineto{\pgfqpoint{2.396575in}{1.440884in}}%
\pgfpathlineto{\pgfqpoint{2.401121in}{1.445010in}}%
\pgfpathlineto{\pgfqpoint{2.405497in}{1.456494in}}%
\pgfpathlineto{\pgfqpoint{2.412968in}{1.462469in}}%
\pgfpathlineto{\pgfqpoint{2.414833in}{1.467947in}}%
\pgfpathlineto{\pgfqpoint{2.414274in}{1.476742in}}%
\pgfpathlineto{\pgfqpoint{2.407342in}{1.476448in}}%
\pgfpathlineto{\pgfqpoint{2.406029in}{1.504012in}}%
\pgfpathlineto{\pgfqpoint{2.426710in}{1.505121in}}%
\pgfpathlineto{\pgfqpoint{2.425489in}{1.525655in}}%
\pgfpathlineto{\pgfqpoint{2.412483in}{1.524944in}}%
\pgfpathlineto{\pgfqpoint{2.410650in}{1.552521in}}%
\pgfpathlineto{\pgfqpoint{2.451750in}{1.555104in}}%
\pgfpathlineto{\pgfqpoint{2.451488in}{1.558550in}}%
\pgfpathlineto{\pgfqpoint{2.484797in}{1.561007in}}%
\pgfpathlineto{\pgfqpoint{2.488674in}{1.519766in}}%
\pgfpathlineto{\pgfqpoint{2.468185in}{1.518191in}}%
\pgfpathlineto{\pgfqpoint{2.470141in}{1.490384in}}%
\pgfpathlineto{\pgfqpoint{2.467299in}{1.483195in}}%
\pgfpathlineto{\pgfqpoint{2.458726in}{1.477228in}}%
\pgfpathlineto{\pgfqpoint{2.450616in}{1.476514in}}%
\pgfpathlineto{\pgfqpoint{2.452287in}{1.455302in}}%
\pgfpathlineto{\pgfqpoint{2.431914in}{1.453643in}}%
\pgfpathlineto{\pgfqpoint{2.432920in}{1.438056in}}%
\pgfpathlineto{\pgfqpoint{2.427687in}{1.440218in}}%
\pgfpathlineto{\pgfqpoint{2.418888in}{1.437396in}}%
\pgfpathlineto{\pgfqpoint{2.410568in}{1.438329in}}%
\pgfpathlineto{\pgfqpoint{2.409257in}{1.435587in}}%
\pgfpathclose%
\pgfusepath{fill}%
\end{pgfscope}%
\begin{pgfscope}%
\pgfpathrectangle{\pgfqpoint{0.100000in}{0.100000in}}{\pgfqpoint{3.608454in}{2.310000in}}%
\pgfusepath{clip}%
\pgfsetbuttcap%
\pgfsetmiterjoin%
\definecolor{currentfill}{rgb}{0.000000,0.560784,0.719608}%
\pgfsetfillcolor{currentfill}%
\pgfsetlinewidth{0.000000pt}%
\definecolor{currentstroke}{rgb}{0.000000,0.000000,0.000000}%
\pgfsetstrokecolor{currentstroke}%
\pgfsetstrokeopacity{0.000000}%
\pgfsetdash{}{0pt}%
\pgfpathmoveto{\pgfqpoint{1.843805in}{1.813366in}}%
\pgfpathlineto{\pgfqpoint{1.849510in}{1.813058in}}%
\pgfpathlineto{\pgfqpoint{1.898285in}{1.810809in}}%
\pgfpathlineto{\pgfqpoint{1.897415in}{1.789939in}}%
\pgfpathlineto{\pgfqpoint{1.896283in}{1.755231in}}%
\pgfpathlineto{\pgfqpoint{1.883458in}{1.755759in}}%
\pgfpathlineto{\pgfqpoint{1.846482in}{1.757163in}}%
\pgfpathlineto{\pgfqpoint{1.841646in}{1.759396in}}%
\pgfpathlineto{\pgfqpoint{1.842265in}{1.785638in}}%
\pgfpathlineto{\pgfqpoint{1.843805in}{1.813366in}}%
\pgfpathclose%
\pgfusepath{fill}%
\end{pgfscope}%
\begin{pgfscope}%
\pgfpathrectangle{\pgfqpoint{0.100000in}{0.100000in}}{\pgfqpoint{3.608454in}{2.310000in}}%
\pgfusepath{clip}%
\pgfsetbuttcap%
\pgfsetmiterjoin%
\definecolor{currentfill}{rgb}{0.000000,0.796078,0.601961}%
\pgfsetfillcolor{currentfill}%
\pgfsetlinewidth{0.000000pt}%
\definecolor{currentstroke}{rgb}{0.000000,0.000000,0.000000}%
\pgfsetstrokecolor{currentstroke}%
\pgfsetstrokeopacity{0.000000}%
\pgfsetdash{}{0pt}%
\pgfpathmoveto{\pgfqpoint{2.336789in}{0.709454in}}%
\pgfpathlineto{\pgfqpoint{2.341562in}{0.715016in}}%
\pgfpathlineto{\pgfqpoint{2.337261in}{0.720197in}}%
\pgfpathlineto{\pgfqpoint{2.341299in}{0.728481in}}%
\pgfpathlineto{\pgfqpoint{2.346226in}{0.727949in}}%
\pgfpathlineto{\pgfqpoint{2.347315in}{0.735413in}}%
\pgfpathlineto{\pgfqpoint{2.362838in}{0.740120in}}%
\pgfpathlineto{\pgfqpoint{2.371609in}{0.736436in}}%
\pgfpathlineto{\pgfqpoint{2.377876in}{0.738907in}}%
\pgfpathlineto{\pgfqpoint{2.404037in}{0.740243in}}%
\pgfpathlineto{\pgfqpoint{2.439716in}{0.742196in}}%
\pgfpathlineto{\pgfqpoint{2.444160in}{0.745936in}}%
\pgfpathlineto{\pgfqpoint{2.449720in}{0.745908in}}%
\pgfpathlineto{\pgfqpoint{2.458936in}{0.746389in}}%
\pgfpathlineto{\pgfqpoint{2.458660in}{0.749893in}}%
\pgfpathlineto{\pgfqpoint{2.470222in}{0.750540in}}%
\pgfpathlineto{\pgfqpoint{2.472257in}{0.716197in}}%
\pgfpathlineto{\pgfqpoint{2.467199in}{0.715885in}}%
\pgfpathlineto{\pgfqpoint{2.424922in}{0.713419in}}%
\pgfpathlineto{\pgfqpoint{2.392212in}{0.711750in}}%
\pgfpathlineto{\pgfqpoint{2.336789in}{0.709454in}}%
\pgfpathclose%
\pgfusepath{fill}%
\end{pgfscope}%
\begin{pgfscope}%
\pgfpathrectangle{\pgfqpoint{0.100000in}{0.100000in}}{\pgfqpoint{3.608454in}{2.310000in}}%
\pgfusepath{clip}%
\pgfsetbuttcap%
\pgfsetmiterjoin%
\definecolor{currentfill}{rgb}{0.000000,0.439216,0.780392}%
\pgfsetfillcolor{currentfill}%
\pgfsetlinewidth{0.000000pt}%
\definecolor{currentstroke}{rgb}{0.000000,0.000000,0.000000}%
\pgfsetstrokecolor{currentstroke}%
\pgfsetstrokeopacity{0.000000}%
\pgfsetdash{}{0pt}%
\pgfpathmoveto{\pgfqpoint{2.221382in}{1.649990in}}%
\pgfpathlineto{\pgfqpoint{2.220391in}{1.696881in}}%
\pgfpathlineto{\pgfqpoint{2.253472in}{1.697776in}}%
\pgfpathlineto{\pgfqpoint{2.294773in}{1.699116in}}%
\pgfpathlineto{\pgfqpoint{2.301639in}{1.699356in}}%
\pgfpathlineto{\pgfqpoint{2.303199in}{1.666191in}}%
\pgfpathlineto{\pgfqpoint{2.275677in}{1.665325in}}%
\pgfpathlineto{\pgfqpoint{2.276089in}{1.651404in}}%
\pgfpathlineto{\pgfqpoint{2.221382in}{1.649990in}}%
\pgfpathclose%
\pgfusepath{fill}%
\end{pgfscope}%
\begin{pgfscope}%
\pgfpathrectangle{\pgfqpoint{0.100000in}{0.100000in}}{\pgfqpoint{3.608454in}{2.310000in}}%
\pgfusepath{clip}%
\pgfsetbuttcap%
\pgfsetmiterjoin%
\definecolor{currentfill}{rgb}{0.000000,0.529412,0.735294}%
\pgfsetfillcolor{currentfill}%
\pgfsetlinewidth{0.000000pt}%
\definecolor{currentstroke}{rgb}{0.000000,0.000000,0.000000}%
\pgfsetstrokecolor{currentstroke}%
\pgfsetstrokeopacity{0.000000}%
\pgfsetdash{}{0pt}%
\pgfpathmoveto{\pgfqpoint{3.109751in}{1.709791in}}%
\pgfpathlineto{\pgfqpoint{3.102306in}{1.712116in}}%
\pgfpathlineto{\pgfqpoint{3.066495in}{1.705513in}}%
\pgfpathlineto{\pgfqpoint{3.064906in}{1.706470in}}%
\pgfpathlineto{\pgfqpoint{3.049340in}{1.697252in}}%
\pgfpathlineto{\pgfqpoint{3.042918in}{1.694827in}}%
\pgfpathlineto{\pgfqpoint{3.036017in}{1.700877in}}%
\pgfpathlineto{\pgfqpoint{3.032000in}{1.700788in}}%
\pgfpathlineto{\pgfqpoint{3.027192in}{1.702502in}}%
\pgfpathlineto{\pgfqpoint{3.030633in}{1.712923in}}%
\pgfpathlineto{\pgfqpoint{3.037383in}{1.717726in}}%
\pgfpathlineto{\pgfqpoint{3.040436in}{1.722640in}}%
\pgfpathlineto{\pgfqpoint{3.033554in}{1.734692in}}%
\pgfpathlineto{\pgfqpoint{3.028625in}{1.736179in}}%
\pgfpathlineto{\pgfqpoint{3.027760in}{1.742504in}}%
\pgfpathlineto{\pgfqpoint{3.022366in}{1.756709in}}%
\pgfpathlineto{\pgfqpoint{3.033985in}{1.763896in}}%
\pgfpathlineto{\pgfqpoint{3.044744in}{1.769031in}}%
\pgfpathlineto{\pgfqpoint{3.058321in}{1.773021in}}%
\pgfpathlineto{\pgfqpoint{3.081941in}{1.776783in}}%
\pgfpathlineto{\pgfqpoint{3.094629in}{1.776945in}}%
\pgfpathlineto{\pgfqpoint{3.106459in}{1.771415in}}%
\pgfpathlineto{\pgfqpoint{3.117744in}{1.776426in}}%
\pgfpathlineto{\pgfqpoint{3.141833in}{1.780817in}}%
\pgfpathlineto{\pgfqpoint{3.153298in}{1.788862in}}%
\pgfpathlineto{\pgfqpoint{3.158574in}{1.768645in}}%
\pgfpathlineto{\pgfqpoint{3.158361in}{1.759713in}}%
\pgfpathlineto{\pgfqpoint{3.162874in}{1.739825in}}%
\pgfpathlineto{\pgfqpoint{3.167852in}{1.733612in}}%
\pgfpathlineto{\pgfqpoint{3.173622in}{1.728849in}}%
\pgfpathlineto{\pgfqpoint{3.155225in}{1.724423in}}%
\pgfpathlineto{\pgfqpoint{3.156855in}{1.718354in}}%
\pgfpathlineto{\pgfqpoint{3.142173in}{1.718464in}}%
\pgfpathlineto{\pgfqpoint{3.141312in}{1.724493in}}%
\pgfpathlineto{\pgfqpoint{3.111728in}{1.718970in}}%
\pgfpathlineto{\pgfqpoint{3.109751in}{1.709791in}}%
\pgfpathclose%
\pgfusepath{fill}%
\end{pgfscope}%
\begin{pgfscope}%
\pgfpathrectangle{\pgfqpoint{0.100000in}{0.100000in}}{\pgfqpoint{3.608454in}{2.310000in}}%
\pgfusepath{clip}%
\pgfsetbuttcap%
\pgfsetmiterjoin%
\definecolor{currentfill}{rgb}{0.000000,0.466667,0.766667}%
\pgfsetfillcolor{currentfill}%
\pgfsetlinewidth{0.000000pt}%
\definecolor{currentstroke}{rgb}{0.000000,0.000000,0.000000}%
\pgfsetstrokecolor{currentstroke}%
\pgfsetstrokeopacity{0.000000}%
\pgfsetdash{}{0pt}%
\pgfpathmoveto{\pgfqpoint{2.989320in}{1.467586in}}%
\pgfpathlineto{\pgfqpoint{2.983438in}{1.466654in}}%
\pgfpathlineto{\pgfqpoint{2.977955in}{1.500863in}}%
\pgfpathlineto{\pgfqpoint{2.971951in}{1.538265in}}%
\pgfpathlineto{\pgfqpoint{2.969282in}{1.554858in}}%
\pgfpathlineto{\pgfqpoint{2.990516in}{1.558619in}}%
\pgfpathlineto{\pgfqpoint{2.988249in}{1.569893in}}%
\pgfpathlineto{\pgfqpoint{2.995768in}{1.584878in}}%
\pgfpathlineto{\pgfqpoint{3.013646in}{1.587974in}}%
\pgfpathlineto{\pgfqpoint{3.020349in}{1.574020in}}%
\pgfpathlineto{\pgfqpoint{3.028102in}{1.574866in}}%
\pgfpathlineto{\pgfqpoint{3.040035in}{1.579416in}}%
\pgfpathlineto{\pgfqpoint{3.043474in}{1.583306in}}%
\pgfpathlineto{\pgfqpoint{3.045363in}{1.572457in}}%
\pgfpathlineto{\pgfqpoint{3.069571in}{1.576329in}}%
\pgfpathlineto{\pgfqpoint{3.072318in}{1.562246in}}%
\pgfpathlineto{\pgfqpoint{3.068856in}{1.543568in}}%
\pgfpathlineto{\pgfqpoint{3.061504in}{1.519627in}}%
\pgfpathlineto{\pgfqpoint{3.054176in}{1.507940in}}%
\pgfpathlineto{\pgfqpoint{3.048216in}{1.491838in}}%
\pgfpathlineto{\pgfqpoint{3.051273in}{1.485944in}}%
\pgfpathlineto{\pgfqpoint{3.051009in}{1.477958in}}%
\pgfpathlineto{\pgfqpoint{3.045971in}{1.477046in}}%
\pgfpathlineto{\pgfqpoint{2.989320in}{1.467586in}}%
\pgfpathclose%
\pgfusepath{fill}%
\end{pgfscope}%
\begin{pgfscope}%
\pgfpathrectangle{\pgfqpoint{0.100000in}{0.100000in}}{\pgfqpoint{3.608454in}{2.310000in}}%
\pgfusepath{clip}%
\pgfsetbuttcap%
\pgfsetmiterjoin%
\definecolor{currentfill}{rgb}{0.000000,0.596078,0.701961}%
\pgfsetfillcolor{currentfill}%
\pgfsetlinewidth{0.000000pt}%
\definecolor{currentstroke}{rgb}{0.000000,0.000000,0.000000}%
\pgfsetstrokecolor{currentstroke}%
\pgfsetstrokeopacity{0.000000}%
\pgfsetdash{}{0pt}%
\pgfpathmoveto{\pgfqpoint{1.787654in}{0.626021in}}%
\pgfpathlineto{\pgfqpoint{1.781527in}{0.626427in}}%
\pgfpathlineto{\pgfqpoint{1.782309in}{0.639635in}}%
\pgfpathlineto{\pgfqpoint{1.777624in}{0.639888in}}%
\pgfpathlineto{\pgfqpoint{1.778772in}{0.657165in}}%
\pgfpathlineto{\pgfqpoint{1.809932in}{0.655294in}}%
\pgfpathlineto{\pgfqpoint{1.810600in}{0.672318in}}%
\pgfpathlineto{\pgfqpoint{1.833962in}{0.671147in}}%
\pgfpathlineto{\pgfqpoint{1.859593in}{0.670192in}}%
\pgfpathlineto{\pgfqpoint{1.858713in}{0.641352in}}%
\pgfpathlineto{\pgfqpoint{1.835751in}{0.642247in}}%
\pgfpathlineto{\pgfqpoint{1.834695in}{0.613722in}}%
\pgfpathlineto{\pgfqpoint{1.817099in}{0.623588in}}%
\pgfpathlineto{\pgfqpoint{1.787654in}{0.626021in}}%
\pgfpathclose%
\pgfusepath{fill}%
\end{pgfscope}%
\begin{pgfscope}%
\pgfpathrectangle{\pgfqpoint{0.100000in}{0.100000in}}{\pgfqpoint{3.608454in}{2.310000in}}%
\pgfusepath{clip}%
\pgfsetbuttcap%
\pgfsetmiterjoin%
\definecolor{currentfill}{rgb}{0.000000,0.619608,0.690196}%
\pgfsetfillcolor{currentfill}%
\pgfsetlinewidth{0.000000pt}%
\definecolor{currentstroke}{rgb}{0.000000,0.000000,0.000000}%
\pgfsetstrokecolor{currentstroke}%
\pgfsetstrokeopacity{0.000000}%
\pgfsetdash{}{0pt}%
\pgfpathmoveto{\pgfqpoint{2.870396in}{0.968762in}}%
\pgfpathlineto{\pgfqpoint{2.877953in}{0.974806in}}%
\pgfpathlineto{\pgfqpoint{2.889208in}{0.971223in}}%
\pgfpathlineto{\pgfqpoint{2.896071in}{0.971539in}}%
\pgfpathlineto{\pgfqpoint{2.896953in}{0.978821in}}%
\pgfpathlineto{\pgfqpoint{2.902891in}{0.990956in}}%
\pgfpathlineto{\pgfqpoint{2.908530in}{0.995600in}}%
\pgfpathlineto{\pgfqpoint{2.916400in}{0.998276in}}%
\pgfpathlineto{\pgfqpoint{2.922696in}{0.996327in}}%
\pgfpathlineto{\pgfqpoint{2.941645in}{0.986529in}}%
\pgfpathlineto{\pgfqpoint{2.948243in}{0.978295in}}%
\pgfpathlineto{\pgfqpoint{2.942345in}{0.978613in}}%
\pgfpathlineto{\pgfqpoint{2.932420in}{0.972207in}}%
\pgfpathlineto{\pgfqpoint{2.922909in}{0.970556in}}%
\pgfpathlineto{\pgfqpoint{2.918991in}{0.967320in}}%
\pgfpathlineto{\pgfqpoint{2.915431in}{0.959683in}}%
\pgfpathlineto{\pgfqpoint{2.910547in}{0.958327in}}%
\pgfpathlineto{\pgfqpoint{2.909363in}{0.953452in}}%
\pgfpathlineto{\pgfqpoint{2.898524in}{0.954072in}}%
\pgfpathlineto{\pgfqpoint{2.889760in}{0.943778in}}%
\pgfpathlineto{\pgfqpoint{2.880855in}{0.952903in}}%
\pgfpathlineto{\pgfqpoint{2.870396in}{0.968762in}}%
\pgfpathclose%
\pgfusepath{fill}%
\end{pgfscope}%
\begin{pgfscope}%
\pgfpathrectangle{\pgfqpoint{0.100000in}{0.100000in}}{\pgfqpoint{3.608454in}{2.310000in}}%
\pgfusepath{clip}%
\pgfsetbuttcap%
\pgfsetmiterjoin%
\definecolor{currentfill}{rgb}{0.000000,0.627451,0.686275}%
\pgfsetfillcolor{currentfill}%
\pgfsetlinewidth{0.000000pt}%
\definecolor{currentstroke}{rgb}{0.000000,0.000000,0.000000}%
\pgfsetstrokecolor{currentstroke}%
\pgfsetstrokeopacity{0.000000}%
\pgfsetdash{}{0pt}%
\pgfpathmoveto{\pgfqpoint{2.205934in}{1.422201in}}%
\pgfpathlineto{\pgfqpoint{2.206018in}{1.417015in}}%
\pgfpathlineto{\pgfqpoint{2.181819in}{1.416001in}}%
\pgfpathlineto{\pgfqpoint{2.181491in}{1.429705in}}%
\pgfpathlineto{\pgfqpoint{2.180394in}{1.464846in}}%
\pgfpathlineto{\pgfqpoint{2.193447in}{1.465250in}}%
\pgfpathlineto{\pgfqpoint{2.204405in}{1.465435in}}%
\pgfpathlineto{\pgfqpoint{2.205934in}{1.422201in}}%
\pgfpathclose%
\pgfusepath{fill}%
\end{pgfscope}%
\begin{pgfscope}%
\pgfpathrectangle{\pgfqpoint{0.100000in}{0.100000in}}{\pgfqpoint{3.608454in}{2.310000in}}%
\pgfusepath{clip}%
\pgfsetbuttcap%
\pgfsetmiterjoin%
\definecolor{currentfill}{rgb}{0.000000,0.615686,0.692157}%
\pgfsetfillcolor{currentfill}%
\pgfsetlinewidth{0.000000pt}%
\definecolor{currentstroke}{rgb}{0.000000,0.000000,0.000000}%
\pgfsetstrokecolor{currentstroke}%
\pgfsetstrokeopacity{0.000000}%
\pgfsetdash{}{0pt}%
\pgfpathmoveto{\pgfqpoint{1.826743in}{0.706842in}}%
\pgfpathlineto{\pgfqpoint{1.799907in}{0.708077in}}%
\pgfpathlineto{\pgfqpoint{1.791780in}{0.708494in}}%
\pgfpathlineto{\pgfqpoint{1.792400in}{0.720174in}}%
\pgfpathlineto{\pgfqpoint{1.757470in}{0.722100in}}%
\pgfpathlineto{\pgfqpoint{1.759948in}{0.761279in}}%
\pgfpathlineto{\pgfqpoint{1.751590in}{0.761847in}}%
\pgfpathlineto{\pgfqpoint{1.753852in}{0.801689in}}%
\pgfpathlineto{\pgfqpoint{1.759459in}{0.801389in}}%
\pgfpathlineto{\pgfqpoint{1.788954in}{0.799743in}}%
\pgfpathlineto{\pgfqpoint{1.786260in}{0.758457in}}%
\pgfpathlineto{\pgfqpoint{1.792433in}{0.750538in}}%
\pgfpathlineto{\pgfqpoint{1.804967in}{0.750514in}}%
\pgfpathlineto{\pgfqpoint{1.813328in}{0.744891in}}%
\pgfpathlineto{\pgfqpoint{1.818920in}{0.750454in}}%
\pgfpathlineto{\pgfqpoint{1.828725in}{0.748261in}}%
\pgfpathlineto{\pgfqpoint{1.826743in}{0.706842in}}%
\pgfpathclose%
\pgfusepath{fill}%
\end{pgfscope}%
\begin{pgfscope}%
\pgfpathrectangle{\pgfqpoint{0.100000in}{0.100000in}}{\pgfqpoint{3.608454in}{2.310000in}}%
\pgfusepath{clip}%
\pgfsetbuttcap%
\pgfsetmiterjoin%
\definecolor{currentfill}{rgb}{0.000000,0.639216,0.680392}%
\pgfsetfillcolor{currentfill}%
\pgfsetlinewidth{0.000000pt}%
\definecolor{currentstroke}{rgb}{0.000000,0.000000,0.000000}%
\pgfsetstrokecolor{currentstroke}%
\pgfsetstrokeopacity{0.000000}%
\pgfsetdash{}{0pt}%
\pgfpathmoveto{\pgfqpoint{2.472257in}{0.716197in}}%
\pgfpathlineto{\pgfqpoint{2.470222in}{0.750540in}}%
\pgfpathlineto{\pgfqpoint{2.458660in}{0.749893in}}%
\pgfpathlineto{\pgfqpoint{2.458936in}{0.746389in}}%
\pgfpathlineto{\pgfqpoint{2.449720in}{0.745908in}}%
\pgfpathlineto{\pgfqpoint{2.447975in}{0.756350in}}%
\pgfpathlineto{\pgfqpoint{2.447021in}{0.775387in}}%
\pgfpathlineto{\pgfqpoint{2.468672in}{0.778139in}}%
\pgfpathlineto{\pgfqpoint{2.485619in}{0.780455in}}%
\pgfpathlineto{\pgfqpoint{2.487539in}{0.751603in}}%
\pgfpathlineto{\pgfqpoint{2.525381in}{0.754028in}}%
\pgfpathlineto{\pgfqpoint{2.528702in}{0.712404in}}%
\pgfpathlineto{\pgfqpoint{2.494101in}{0.710117in}}%
\pgfpathlineto{\pgfqpoint{2.493658in}{0.717377in}}%
\pgfpathlineto{\pgfqpoint{2.482623in}{0.717728in}}%
\pgfpathlineto{\pgfqpoint{2.472257in}{0.716197in}}%
\pgfpathclose%
\pgfusepath{fill}%
\end{pgfscope}%
\begin{pgfscope}%
\pgfpathrectangle{\pgfqpoint{0.100000in}{0.100000in}}{\pgfqpoint{3.608454in}{2.310000in}}%
\pgfusepath{clip}%
\pgfsetbuttcap%
\pgfsetmiterjoin%
\definecolor{currentfill}{rgb}{0.000000,0.196078,0.901961}%
\pgfsetfillcolor{currentfill}%
\pgfsetlinewidth{0.000000pt}%
\definecolor{currentstroke}{rgb}{0.000000,0.000000,0.000000}%
\pgfsetstrokecolor{currentstroke}%
\pgfsetstrokeopacity{0.000000}%
\pgfsetdash{}{0pt}%
\pgfpathmoveto{\pgfqpoint{3.129858in}{1.404228in}}%
\pgfpathlineto{\pgfqpoint{3.131195in}{1.415085in}}%
\pgfpathlineto{\pgfqpoint{3.124006in}{1.419453in}}%
\pgfpathlineto{\pgfqpoint{3.127584in}{1.429231in}}%
\pgfpathlineto{\pgfqpoint{3.125843in}{1.434080in}}%
\pgfpathlineto{\pgfqpoint{3.141078in}{1.437828in}}%
\pgfpathlineto{\pgfqpoint{3.144962in}{1.435381in}}%
\pgfpathlineto{\pgfqpoint{3.152017in}{1.444895in}}%
\pgfpathlineto{\pgfqpoint{3.153317in}{1.449269in}}%
\pgfpathlineto{\pgfqpoint{3.156244in}{1.464323in}}%
\pgfpathlineto{\pgfqpoint{3.159520in}{1.465911in}}%
\pgfpathlineto{\pgfqpoint{3.161516in}{1.490087in}}%
\pgfpathlineto{\pgfqpoint{3.166603in}{1.499150in}}%
\pgfpathlineto{\pgfqpoint{3.181076in}{1.501996in}}%
\pgfpathlineto{\pgfqpoint{3.177742in}{1.492668in}}%
\pgfpathlineto{\pgfqpoint{3.186509in}{1.489666in}}%
\pgfpathlineto{\pgfqpoint{3.191082in}{1.485805in}}%
\pgfpathlineto{\pgfqpoint{3.188742in}{1.472751in}}%
\pgfpathlineto{\pgfqpoint{3.192873in}{1.468222in}}%
\pgfpathlineto{\pgfqpoint{3.201406in}{1.464696in}}%
\pgfpathlineto{\pgfqpoint{3.204352in}{1.460668in}}%
\pgfpathlineto{\pgfqpoint{3.223628in}{1.451194in}}%
\pgfpathlineto{\pgfqpoint{3.226611in}{1.445349in}}%
\pgfpathlineto{\pgfqpoint{3.225469in}{1.437928in}}%
\pgfpathlineto{\pgfqpoint{3.227384in}{1.433116in}}%
\pgfpathlineto{\pgfqpoint{3.238143in}{1.432326in}}%
\pgfpathlineto{\pgfqpoint{3.241599in}{1.419010in}}%
\pgfpathlineto{\pgfqpoint{3.251419in}{1.409012in}}%
\pgfpathlineto{\pgfqpoint{3.254865in}{1.395654in}}%
\pgfpathlineto{\pgfqpoint{3.260051in}{1.390426in}}%
\pgfpathlineto{\pgfqpoint{3.258704in}{1.383890in}}%
\pgfpathlineto{\pgfqpoint{3.252557in}{1.388608in}}%
\pgfpathlineto{\pgfqpoint{3.237222in}{1.393476in}}%
\pgfpathlineto{\pgfqpoint{3.230597in}{1.391657in}}%
\pgfpathlineto{\pgfqpoint{3.214679in}{1.396873in}}%
\pgfpathlineto{\pgfqpoint{3.209863in}{1.398797in}}%
\pgfpathlineto{\pgfqpoint{3.199440in}{1.393885in}}%
\pgfpathlineto{\pgfqpoint{3.196379in}{1.396687in}}%
\pgfpathlineto{\pgfqpoint{3.194258in}{1.405953in}}%
\pgfpathlineto{\pgfqpoint{3.182634in}{1.411460in}}%
\pgfpathlineto{\pgfqpoint{3.175115in}{1.399614in}}%
\pgfpathlineto{\pgfqpoint{3.175771in}{1.395430in}}%
\pgfpathlineto{\pgfqpoint{3.169602in}{1.394541in}}%
\pgfpathlineto{\pgfqpoint{3.158541in}{1.406177in}}%
\pgfpathlineto{\pgfqpoint{3.153282in}{1.414127in}}%
\pgfpathlineto{\pgfqpoint{3.141594in}{1.398020in}}%
\pgfpathlineto{\pgfqpoint{3.137714in}{1.398047in}}%
\pgfpathlineto{\pgfqpoint{3.129858in}{1.404228in}}%
\pgfpathclose%
\pgfusepath{fill}%
\end{pgfscope}%
\begin{pgfscope}%
\pgfpathrectangle{\pgfqpoint{0.100000in}{0.100000in}}{\pgfqpoint{3.608454in}{2.310000in}}%
\pgfusepath{clip}%
\pgfsetbuttcap%
\pgfsetmiterjoin%
\definecolor{currentfill}{rgb}{0.000000,0.952941,0.523529}%
\pgfsetfillcolor{currentfill}%
\pgfsetlinewidth{0.000000pt}%
\definecolor{currentstroke}{rgb}{0.000000,0.000000,0.000000}%
\pgfsetstrokecolor{currentstroke}%
\pgfsetstrokeopacity{0.000000}%
\pgfsetdash{}{0pt}%
\pgfpathmoveto{\pgfqpoint{0.605938in}{1.604509in}}%
\pgfpathlineto{\pgfqpoint{0.599878in}{1.583444in}}%
\pgfpathlineto{\pgfqpoint{0.570419in}{1.591928in}}%
\pgfpathlineto{\pgfqpoint{0.567874in}{1.599010in}}%
\pgfpathlineto{\pgfqpoint{0.563417in}{1.600581in}}%
\pgfpathlineto{\pgfqpoint{0.556871in}{1.597127in}}%
\pgfpathlineto{\pgfqpoint{0.538914in}{1.596597in}}%
\pgfpathlineto{\pgfqpoint{0.540950in}{1.606159in}}%
\pgfpathlineto{\pgfqpoint{0.529553in}{1.637642in}}%
\pgfpathlineto{\pgfqpoint{0.526964in}{1.639730in}}%
\pgfpathlineto{\pgfqpoint{0.529796in}{1.650793in}}%
\pgfpathlineto{\pgfqpoint{0.534394in}{1.657082in}}%
\pgfpathlineto{\pgfqpoint{0.531886in}{1.661402in}}%
\pgfpathlineto{\pgfqpoint{0.537339in}{1.664898in}}%
\pgfpathlineto{\pgfqpoint{0.540383in}{1.672372in}}%
\pgfpathlineto{\pgfqpoint{0.534540in}{1.676936in}}%
\pgfpathlineto{\pgfqpoint{0.535033in}{1.684697in}}%
\pgfpathlineto{\pgfqpoint{0.544816in}{1.681752in}}%
\pgfpathlineto{\pgfqpoint{0.560209in}{1.677284in}}%
\pgfpathlineto{\pgfqpoint{0.555953in}{1.662869in}}%
\pgfpathlineto{\pgfqpoint{0.558505in}{1.659196in}}%
\pgfpathlineto{\pgfqpoint{0.566606in}{1.658522in}}%
\pgfpathlineto{\pgfqpoint{0.574470in}{1.662350in}}%
\pgfpathlineto{\pgfqpoint{0.584658in}{1.656876in}}%
\pgfpathlineto{\pgfqpoint{0.594599in}{1.640016in}}%
\pgfpathlineto{\pgfqpoint{0.601624in}{1.635606in}}%
\pgfpathlineto{\pgfqpoint{0.600541in}{1.629953in}}%
\pgfpathlineto{\pgfqpoint{0.604334in}{1.622797in}}%
\pgfpathlineto{\pgfqpoint{0.600479in}{1.609616in}}%
\pgfpathlineto{\pgfqpoint{0.605938in}{1.604509in}}%
\pgfpathclose%
\pgfusepath{fill}%
\end{pgfscope}%
\begin{pgfscope}%
\pgfpathrectangle{\pgfqpoint{0.100000in}{0.100000in}}{\pgfqpoint{3.608454in}{2.310000in}}%
\pgfusepath{clip}%
\pgfsetbuttcap%
\pgfsetmiterjoin%
\definecolor{currentfill}{rgb}{0.000000,0.470588,0.764706}%
\pgfsetfillcolor{currentfill}%
\pgfsetlinewidth{0.000000pt}%
\definecolor{currentstroke}{rgb}{0.000000,0.000000,0.000000}%
\pgfsetstrokecolor{currentstroke}%
\pgfsetstrokeopacity{0.000000}%
\pgfsetdash{}{0pt}%
\pgfpathmoveto{\pgfqpoint{1.888636in}{1.700186in}}%
\pgfpathlineto{\pgfqpoint{1.889333in}{1.734748in}}%
\pgfpathlineto{\pgfqpoint{1.916523in}{1.733891in}}%
\pgfpathlineto{\pgfqpoint{1.917260in}{1.754415in}}%
\pgfpathlineto{\pgfqpoint{1.944449in}{1.753368in}}%
\pgfpathlineto{\pgfqpoint{1.944686in}{1.753364in}}%
\pgfpathlineto{\pgfqpoint{1.943727in}{1.725985in}}%
\pgfpathlineto{\pgfqpoint{1.957610in}{1.725492in}}%
\pgfpathlineto{\pgfqpoint{1.956891in}{1.697804in}}%
\pgfpathlineto{\pgfqpoint{1.888636in}{1.700186in}}%
\pgfpathclose%
\pgfusepath{fill}%
\end{pgfscope}%
\begin{pgfscope}%
\pgfpathrectangle{\pgfqpoint{0.100000in}{0.100000in}}{\pgfqpoint{3.608454in}{2.310000in}}%
\pgfusepath{clip}%
\pgfsetbuttcap%
\pgfsetmiterjoin%
\definecolor{currentfill}{rgb}{0.000000,0.909804,0.545098}%
\pgfsetfillcolor{currentfill}%
\pgfsetlinewidth{0.000000pt}%
\definecolor{currentstroke}{rgb}{0.000000,0.000000,0.000000}%
\pgfsetstrokecolor{currentstroke}%
\pgfsetstrokeopacity{0.000000}%
\pgfsetdash{}{0pt}%
\pgfpathmoveto{\pgfqpoint{0.937150in}{2.279229in}}%
\pgfpathlineto{\pgfqpoint{0.927075in}{2.268237in}}%
\pgfpathlineto{\pgfqpoint{0.925222in}{2.260906in}}%
\pgfpathlineto{\pgfqpoint{0.920139in}{2.255027in}}%
\pgfpathlineto{\pgfqpoint{0.926800in}{2.253370in}}%
\pgfpathlineto{\pgfqpoint{0.916424in}{2.212894in}}%
\pgfpathlineto{\pgfqpoint{0.913019in}{2.206570in}}%
\pgfpathlineto{\pgfqpoint{0.908021in}{2.186968in}}%
\pgfpathlineto{\pgfqpoint{0.902377in}{2.193011in}}%
\pgfpathlineto{\pgfqpoint{0.894472in}{2.191910in}}%
\pgfpathlineto{\pgfqpoint{0.891121in}{2.195618in}}%
\pgfpathlineto{\pgfqpoint{0.879920in}{2.194564in}}%
\pgfpathlineto{\pgfqpoint{0.878250in}{2.201957in}}%
\pgfpathlineto{\pgfqpoint{0.873890in}{2.207982in}}%
\pgfpathlineto{\pgfqpoint{0.862221in}{2.202501in}}%
\pgfpathlineto{\pgfqpoint{0.857008in}{2.211536in}}%
\pgfpathlineto{\pgfqpoint{0.843741in}{2.216808in}}%
\pgfpathlineto{\pgfqpoint{0.836363in}{2.217995in}}%
\pgfpathlineto{\pgfqpoint{0.836288in}{2.219545in}}%
\pgfpathlineto{\pgfqpoint{0.836249in}{2.233779in}}%
\pgfpathlineto{\pgfqpoint{0.832614in}{2.236217in}}%
\pgfpathlineto{\pgfqpoint{0.819597in}{2.231350in}}%
\pgfpathlineto{\pgfqpoint{0.815034in}{2.234984in}}%
\pgfpathlineto{\pgfqpoint{0.802328in}{2.232125in}}%
\pgfpathlineto{\pgfqpoint{0.800018in}{2.241064in}}%
\pgfpathlineto{\pgfqpoint{0.791522in}{2.239767in}}%
\pgfpathlineto{\pgfqpoint{0.790372in}{2.232053in}}%
\pgfpathlineto{\pgfqpoint{0.780931in}{2.234819in}}%
\pgfpathlineto{\pgfqpoint{0.780160in}{2.239854in}}%
\pgfpathlineto{\pgfqpoint{0.769468in}{2.254538in}}%
\pgfpathlineto{\pgfqpoint{0.771087in}{2.257398in}}%
\pgfpathlineto{\pgfqpoint{0.765081in}{2.268518in}}%
\pgfpathlineto{\pgfqpoint{0.760184in}{2.281456in}}%
\pgfpathlineto{\pgfqpoint{0.760304in}{2.288239in}}%
\pgfpathlineto{\pgfqpoint{0.757956in}{2.296770in}}%
\pgfpathlineto{\pgfqpoint{0.760555in}{2.298682in}}%
\pgfpathlineto{\pgfqpoint{0.765166in}{2.300496in}}%
\pgfpathlineto{\pgfqpoint{0.763231in}{2.326563in}}%
\pgfpathlineto{\pgfqpoint{0.801462in}{2.315507in}}%
\pgfpathlineto{\pgfqpoint{0.847920in}{2.302634in}}%
\pgfpathlineto{\pgfqpoint{0.882493in}{2.293352in}}%
\pgfpathlineto{\pgfqpoint{0.937150in}{2.279229in}}%
\pgfpathclose%
\pgfusepath{fill}%
\end{pgfscope}%
\begin{pgfscope}%
\pgfpathrectangle{\pgfqpoint{0.100000in}{0.100000in}}{\pgfqpoint{3.608454in}{2.310000in}}%
\pgfusepath{clip}%
\pgfsetbuttcap%
\pgfsetmiterjoin%
\definecolor{currentfill}{rgb}{0.000000,0.580392,0.709804}%
\pgfsetfillcolor{currentfill}%
\pgfsetlinewidth{0.000000pt}%
\definecolor{currentstroke}{rgb}{0.000000,0.000000,0.000000}%
\pgfsetstrokecolor{currentstroke}%
\pgfsetstrokeopacity{0.000000}%
\pgfsetdash{}{0pt}%
\pgfpathmoveto{\pgfqpoint{2.821573in}{0.886623in}}%
\pgfpathlineto{\pgfqpoint{2.814267in}{0.888921in}}%
\pgfpathlineto{\pgfqpoint{2.805495in}{0.895708in}}%
\pgfpathlineto{\pgfqpoint{2.802438in}{0.905479in}}%
\pgfpathlineto{\pgfqpoint{2.802911in}{0.922322in}}%
\pgfpathlineto{\pgfqpoint{2.802892in}{0.925244in}}%
\pgfpathlineto{\pgfqpoint{2.807171in}{0.925663in}}%
\pgfpathlineto{\pgfqpoint{2.809234in}{0.933588in}}%
\pgfpathlineto{\pgfqpoint{2.818507in}{0.934685in}}%
\pgfpathlineto{\pgfqpoint{2.830443in}{0.932055in}}%
\pgfpathlineto{\pgfqpoint{2.838874in}{0.945035in}}%
\pgfpathlineto{\pgfqpoint{2.843454in}{0.939687in}}%
\pgfpathlineto{\pgfqpoint{2.848096in}{0.929927in}}%
\pgfpathlineto{\pgfqpoint{2.848183in}{0.924494in}}%
\pgfpathlineto{\pgfqpoint{2.844937in}{0.925826in}}%
\pgfpathlineto{\pgfqpoint{2.833490in}{0.924483in}}%
\pgfpathlineto{\pgfqpoint{2.835752in}{0.904367in}}%
\pgfpathlineto{\pgfqpoint{2.830655in}{0.902493in}}%
\pgfpathlineto{\pgfqpoint{2.830174in}{0.891948in}}%
\pgfpathlineto{\pgfqpoint{2.827685in}{0.882746in}}%
\pgfpathlineto{\pgfqpoint{2.821573in}{0.886623in}}%
\pgfpathclose%
\pgfusepath{fill}%
\end{pgfscope}%
\begin{pgfscope}%
\pgfpathrectangle{\pgfqpoint{0.100000in}{0.100000in}}{\pgfqpoint{3.608454in}{2.310000in}}%
\pgfusepath{clip}%
\pgfsetbuttcap%
\pgfsetmiterjoin%
\definecolor{currentfill}{rgb}{0.000000,0.286275,0.856863}%
\pgfsetfillcolor{currentfill}%
\pgfsetlinewidth{0.000000pt}%
\definecolor{currentstroke}{rgb}{0.000000,0.000000,0.000000}%
\pgfsetstrokecolor{currentstroke}%
\pgfsetstrokeopacity{0.000000}%
\pgfsetdash{}{0pt}%
\pgfpathmoveto{\pgfqpoint{3.238143in}{1.432326in}}%
\pgfpathlineto{\pgfqpoint{3.227384in}{1.433116in}}%
\pgfpathlineto{\pgfqpoint{3.225469in}{1.437928in}}%
\pgfpathlineto{\pgfqpoint{3.226611in}{1.445349in}}%
\pgfpathlineto{\pgfqpoint{3.223628in}{1.451194in}}%
\pgfpathlineto{\pgfqpoint{3.204352in}{1.460668in}}%
\pgfpathlineto{\pgfqpoint{3.201406in}{1.464696in}}%
\pgfpathlineto{\pgfqpoint{3.192873in}{1.468222in}}%
\pgfpathlineto{\pgfqpoint{3.188742in}{1.472751in}}%
\pgfpathlineto{\pgfqpoint{3.191082in}{1.485805in}}%
\pgfpathlineto{\pgfqpoint{3.186509in}{1.489666in}}%
\pgfpathlineto{\pgfqpoint{3.177742in}{1.492668in}}%
\pgfpathlineto{\pgfqpoint{3.181076in}{1.501996in}}%
\pgfpathlineto{\pgfqpoint{3.239245in}{1.513965in}}%
\pgfpathlineto{\pgfqpoint{3.239646in}{1.514041in}}%
\pgfpathlineto{\pgfqpoint{3.251965in}{1.502268in}}%
\pgfpathlineto{\pgfqpoint{3.252119in}{1.493455in}}%
\pgfpathlineto{\pgfqpoint{3.248435in}{1.488117in}}%
\pgfpathlineto{\pgfqpoint{3.238301in}{1.480374in}}%
\pgfpathlineto{\pgfqpoint{3.235613in}{1.470733in}}%
\pgfpathlineto{\pgfqpoint{3.238603in}{1.459477in}}%
\pgfpathlineto{\pgfqpoint{3.239193in}{1.448194in}}%
\pgfpathlineto{\pgfqpoint{3.235959in}{1.434274in}}%
\pgfpathlineto{\pgfqpoint{3.238143in}{1.432326in}}%
\pgfpathclose%
\pgfusepath{fill}%
\end{pgfscope}%
\begin{pgfscope}%
\pgfpathrectangle{\pgfqpoint{0.100000in}{0.100000in}}{\pgfqpoint{3.608454in}{2.310000in}}%
\pgfusepath{clip}%
\pgfsetbuttcap%
\pgfsetmiterjoin%
\definecolor{currentfill}{rgb}{0.000000,0.592157,0.703922}%
\pgfsetfillcolor{currentfill}%
\pgfsetlinewidth{0.000000pt}%
\definecolor{currentstroke}{rgb}{0.000000,0.000000,0.000000}%
\pgfsetstrokecolor{currentstroke}%
\pgfsetstrokeopacity{0.000000}%
\pgfsetdash{}{0pt}%
\pgfpathmoveto{\pgfqpoint{1.821831in}{1.246792in}}%
\pgfpathlineto{\pgfqpoint{1.821521in}{1.253978in}}%
\pgfpathlineto{\pgfqpoint{1.823241in}{1.288458in}}%
\pgfpathlineto{\pgfqpoint{1.822395in}{1.288502in}}%
\pgfpathlineto{\pgfqpoint{1.824072in}{1.322793in}}%
\pgfpathlineto{\pgfqpoint{1.891797in}{1.319820in}}%
\pgfpathlineto{\pgfqpoint{1.890012in}{1.250882in}}%
\pgfpathlineto{\pgfqpoint{1.856262in}{1.252307in}}%
\pgfpathlineto{\pgfqpoint{1.855927in}{1.245032in}}%
\pgfpathlineto{\pgfqpoint{1.821831in}{1.246792in}}%
\pgfpathclose%
\pgfusepath{fill}%
\end{pgfscope}%
\begin{pgfscope}%
\pgfpathrectangle{\pgfqpoint{0.100000in}{0.100000in}}{\pgfqpoint{3.608454in}{2.310000in}}%
\pgfusepath{clip}%
\pgfsetbuttcap%
\pgfsetmiterjoin%
\definecolor{currentfill}{rgb}{0.000000,0.678431,0.660784}%
\pgfsetfillcolor{currentfill}%
\pgfsetlinewidth{0.000000pt}%
\definecolor{currentstroke}{rgb}{0.000000,0.000000,0.000000}%
\pgfsetstrokecolor{currentstroke}%
\pgfsetstrokeopacity{0.000000}%
\pgfsetdash{}{0pt}%
\pgfpathmoveto{\pgfqpoint{2.625838in}{1.802295in}}%
\pgfpathlineto{\pgfqpoint{2.602311in}{1.800403in}}%
\pgfpathlineto{\pgfqpoint{2.600324in}{1.814817in}}%
\pgfpathlineto{\pgfqpoint{2.608715in}{1.818757in}}%
\pgfpathlineto{\pgfqpoint{2.608533in}{1.830495in}}%
\pgfpathlineto{\pgfqpoint{2.613100in}{1.831951in}}%
\pgfpathlineto{\pgfqpoint{2.615491in}{1.837204in}}%
\pgfpathlineto{\pgfqpoint{2.623020in}{1.837013in}}%
\pgfpathlineto{\pgfqpoint{2.625532in}{1.845348in}}%
\pgfpathlineto{\pgfqpoint{2.632241in}{1.851902in}}%
\pgfpathlineto{\pgfqpoint{2.634383in}{1.847792in}}%
\pgfpathlineto{\pgfqpoint{2.634627in}{1.836092in}}%
\pgfpathlineto{\pgfqpoint{2.632415in}{1.832536in}}%
\pgfpathlineto{\pgfqpoint{2.633665in}{1.824033in}}%
\pgfpathlineto{\pgfqpoint{2.640260in}{1.822509in}}%
\pgfpathlineto{\pgfqpoint{2.645608in}{1.835898in}}%
\pgfpathlineto{\pgfqpoint{2.646411in}{1.852152in}}%
\pgfpathlineto{\pgfqpoint{2.644203in}{1.859587in}}%
\pgfpathlineto{\pgfqpoint{2.653419in}{1.860520in}}%
\pgfpathlineto{\pgfqpoint{2.654113in}{1.853564in}}%
\pgfpathlineto{\pgfqpoint{2.674595in}{1.855705in}}%
\pgfpathlineto{\pgfqpoint{2.677389in}{1.835395in}}%
\pgfpathlineto{\pgfqpoint{2.680280in}{1.807943in}}%
\pgfpathlineto{\pgfqpoint{2.625838in}{1.802295in}}%
\pgfpathclose%
\pgfusepath{fill}%
\end{pgfscope}%
\begin{pgfscope}%
\pgfpathrectangle{\pgfqpoint{0.100000in}{0.100000in}}{\pgfqpoint{3.608454in}{2.310000in}}%
\pgfusepath{clip}%
\pgfsetbuttcap%
\pgfsetmiterjoin%
\definecolor{currentfill}{rgb}{0.000000,0.305882,0.847059}%
\pgfsetfillcolor{currentfill}%
\pgfsetlinewidth{0.000000pt}%
\definecolor{currentstroke}{rgb}{0.000000,0.000000,0.000000}%
\pgfsetstrokecolor{currentstroke}%
\pgfsetstrokeopacity{0.000000}%
\pgfsetdash{}{0pt}%
\pgfpathmoveto{\pgfqpoint{1.830021in}{2.007934in}}%
\pgfpathlineto{\pgfqpoint{1.828413in}{1.980392in}}%
\pgfpathlineto{\pgfqpoint{1.830555in}{1.980267in}}%
\pgfpathlineto{\pgfqpoint{1.828664in}{1.952547in}}%
\pgfpathlineto{\pgfqpoint{1.796923in}{1.954510in}}%
\pgfpathlineto{\pgfqpoint{1.801850in}{1.950975in}}%
\pgfpathlineto{\pgfqpoint{1.802847in}{1.946100in}}%
\pgfpathlineto{\pgfqpoint{1.799830in}{1.937828in}}%
\pgfpathlineto{\pgfqpoint{1.787591in}{1.935130in}}%
\pgfpathlineto{\pgfqpoint{1.783790in}{1.937380in}}%
\pgfpathlineto{\pgfqpoint{1.775193in}{1.927930in}}%
\pgfpathlineto{\pgfqpoint{1.774399in}{1.934946in}}%
\pgfpathlineto{\pgfqpoint{1.760656in}{1.935885in}}%
\pgfpathlineto{\pgfqpoint{1.762088in}{1.956478in}}%
\pgfpathlineto{\pgfqpoint{1.739184in}{1.958172in}}%
\pgfpathlineto{\pgfqpoint{1.739705in}{1.965094in}}%
\pgfpathlineto{\pgfqpoint{1.719085in}{1.966639in}}%
\pgfpathlineto{\pgfqpoint{1.720698in}{1.987602in}}%
\pgfpathlineto{\pgfqpoint{1.718098in}{1.987796in}}%
\pgfpathlineto{\pgfqpoint{1.720286in}{2.015429in}}%
\pgfpathlineto{\pgfqpoint{1.716982in}{2.015681in}}%
\pgfpathlineto{\pgfqpoint{1.718555in}{2.035319in}}%
\pgfpathlineto{\pgfqpoint{1.715648in}{2.039417in}}%
\pgfpathlineto{\pgfqpoint{1.716886in}{2.052052in}}%
\pgfpathlineto{\pgfqpoint{1.710047in}{2.050726in}}%
\pgfpathlineto{\pgfqpoint{1.710634in}{2.057876in}}%
\pgfpathlineto{\pgfqpoint{1.738152in}{2.055718in}}%
\pgfpathlineto{\pgfqpoint{1.807069in}{2.050968in}}%
\pgfpathlineto{\pgfqpoint{1.806219in}{2.037109in}}%
\pgfpathlineto{\pgfqpoint{1.801529in}{2.037408in}}%
\pgfpathlineto{\pgfqpoint{1.799746in}{2.009759in}}%
\pgfpathlineto{\pgfqpoint{1.830021in}{2.007934in}}%
\pgfpathclose%
\pgfusepath{fill}%
\end{pgfscope}%
\begin{pgfscope}%
\pgfpathrectangle{\pgfqpoint{0.100000in}{0.100000in}}{\pgfqpoint{3.608454in}{2.310000in}}%
\pgfusepath{clip}%
\pgfsetbuttcap%
\pgfsetmiterjoin%
\definecolor{currentfill}{rgb}{0.000000,0.423529,0.788235}%
\pgfsetfillcolor{currentfill}%
\pgfsetlinewidth{0.000000pt}%
\definecolor{currentstroke}{rgb}{0.000000,0.000000,0.000000}%
\pgfsetstrokecolor{currentstroke}%
\pgfsetstrokeopacity{0.000000}%
\pgfsetdash{}{0pt}%
\pgfpathmoveto{\pgfqpoint{3.285323in}{1.505755in}}%
\pgfpathlineto{\pgfqpoint{3.285117in}{1.511332in}}%
\pgfpathlineto{\pgfqpoint{3.281349in}{1.515956in}}%
\pgfpathlineto{\pgfqpoint{3.284745in}{1.526913in}}%
\pgfpathlineto{\pgfqpoint{3.284318in}{1.531602in}}%
\pgfpathlineto{\pgfqpoint{3.275612in}{1.531070in}}%
\pgfpathlineto{\pgfqpoint{3.271013in}{1.528198in}}%
\pgfpathlineto{\pgfqpoint{3.266966in}{1.519876in}}%
\pgfpathlineto{\pgfqpoint{3.245459in}{1.515286in}}%
\pgfpathlineto{\pgfqpoint{3.248921in}{1.519116in}}%
\pgfpathlineto{\pgfqpoint{3.250934in}{1.527034in}}%
\pgfpathlineto{\pgfqpoint{3.250059in}{1.534993in}}%
\pgfpathlineto{\pgfqpoint{3.252158in}{1.542409in}}%
\pgfpathlineto{\pgfqpoint{3.250514in}{1.547606in}}%
\pgfpathlineto{\pgfqpoint{3.262873in}{1.561037in}}%
\pgfpathlineto{\pgfqpoint{3.269282in}{1.579037in}}%
\pgfpathlineto{\pgfqpoint{3.275573in}{1.583652in}}%
\pgfpathlineto{\pgfqpoint{3.286080in}{1.595797in}}%
\pgfpathlineto{\pgfqpoint{3.294732in}{1.592344in}}%
\pgfpathlineto{\pgfqpoint{3.297442in}{1.582891in}}%
\pgfpathlineto{\pgfqpoint{3.303159in}{1.582589in}}%
\pgfpathlineto{\pgfqpoint{3.313073in}{1.572820in}}%
\pgfpathlineto{\pgfqpoint{3.330022in}{1.567556in}}%
\pgfpathlineto{\pgfqpoint{3.348171in}{1.542140in}}%
\pgfpathlineto{\pgfqpoint{3.353587in}{1.521095in}}%
\pgfpathlineto{\pgfqpoint{3.358684in}{1.520578in}}%
\pgfpathlineto{\pgfqpoint{3.354142in}{1.510014in}}%
\pgfpathlineto{\pgfqpoint{3.344357in}{1.498736in}}%
\pgfpathlineto{\pgfqpoint{3.337885in}{1.476407in}}%
\pgfpathlineto{\pgfqpoint{3.334527in}{1.471502in}}%
\pgfpathlineto{\pgfqpoint{3.328098in}{1.470069in}}%
\pgfpathlineto{\pgfqpoint{3.329144in}{1.488418in}}%
\pgfpathlineto{\pgfqpoint{3.310575in}{1.490636in}}%
\pgfpathlineto{\pgfqpoint{3.301327in}{1.493916in}}%
\pgfpathlineto{\pgfqpoint{3.288516in}{1.501805in}}%
\pgfpathlineto{\pgfqpoint{3.285323in}{1.505755in}}%
\pgfpathclose%
\pgfusepath{fill}%
\end{pgfscope}%
\begin{pgfscope}%
\pgfpathrectangle{\pgfqpoint{0.100000in}{0.100000in}}{\pgfqpoint{3.608454in}{2.310000in}}%
\pgfusepath{clip}%
\pgfsetbuttcap%
\pgfsetmiterjoin%
\definecolor{currentfill}{rgb}{0.000000,0.333333,0.833333}%
\pgfsetfillcolor{currentfill}%
\pgfsetlinewidth{0.000000pt}%
\definecolor{currentstroke}{rgb}{0.000000,0.000000,0.000000}%
\pgfsetstrokecolor{currentstroke}%
\pgfsetstrokeopacity{0.000000}%
\pgfsetdash{}{0pt}%
\pgfpathmoveto{\pgfqpoint{1.725856in}{1.224921in}}%
\pgfpathlineto{\pgfqpoint{1.723873in}{1.194151in}}%
\pgfpathlineto{\pgfqpoint{1.693075in}{1.196095in}}%
\pgfpathlineto{\pgfqpoint{1.695230in}{1.227106in}}%
\pgfpathlineto{\pgfqpoint{1.664847in}{1.229439in}}%
\pgfpathlineto{\pgfqpoint{1.666412in}{1.249548in}}%
\pgfpathlineto{\pgfqpoint{1.666954in}{1.256981in}}%
\pgfpathlineto{\pgfqpoint{1.726337in}{1.252520in}}%
\pgfpathlineto{\pgfqpoint{1.724493in}{1.225005in}}%
\pgfpathlineto{\pgfqpoint{1.725856in}{1.224921in}}%
\pgfpathclose%
\pgfusepath{fill}%
\end{pgfscope}%
\begin{pgfscope}%
\pgfpathrectangle{\pgfqpoint{0.100000in}{0.100000in}}{\pgfqpoint{3.608454in}{2.310000in}}%
\pgfusepath{clip}%
\pgfsetbuttcap%
\pgfsetmiterjoin%
\definecolor{currentfill}{rgb}{0.000000,0.596078,0.701961}%
\pgfsetfillcolor{currentfill}%
\pgfsetlinewidth{0.000000pt}%
\definecolor{currentstroke}{rgb}{0.000000,0.000000,0.000000}%
\pgfsetstrokecolor{currentstroke}%
\pgfsetstrokeopacity{0.000000}%
\pgfsetdash{}{0pt}%
\pgfpathmoveto{\pgfqpoint{2.856403in}{0.756814in}}%
\pgfpathlineto{\pgfqpoint{2.856733in}{0.753994in}}%
\pgfpathlineto{\pgfqpoint{2.874929in}{0.755916in}}%
\pgfpathlineto{\pgfqpoint{2.878161in}{0.725802in}}%
\pgfpathlineto{\pgfqpoint{2.834431in}{0.723057in}}%
\pgfpathlineto{\pgfqpoint{2.831112in}{0.753917in}}%
\pgfpathlineto{\pgfqpoint{2.856403in}{0.756814in}}%
\pgfpathclose%
\pgfusepath{fill}%
\end{pgfscope}%
\begin{pgfscope}%
\pgfpathrectangle{\pgfqpoint{0.100000in}{0.100000in}}{\pgfqpoint{3.608454in}{2.310000in}}%
\pgfusepath{clip}%
\pgfsetbuttcap%
\pgfsetmiterjoin%
\definecolor{currentfill}{rgb}{0.000000,0.552941,0.723529}%
\pgfsetfillcolor{currentfill}%
\pgfsetlinewidth{0.000000pt}%
\definecolor{currentstroke}{rgb}{0.000000,0.000000,0.000000}%
\pgfsetstrokecolor{currentstroke}%
\pgfsetstrokeopacity{0.000000}%
\pgfsetdash{}{0pt}%
\pgfpathmoveto{\pgfqpoint{2.519677in}{1.268419in}}%
\pgfpathlineto{\pgfqpoint{2.499015in}{1.266823in}}%
\pgfpathlineto{\pgfqpoint{2.497937in}{1.284068in}}%
\pgfpathlineto{\pgfqpoint{2.470338in}{1.282183in}}%
\pgfpathlineto{\pgfqpoint{2.470034in}{1.289131in}}%
\pgfpathlineto{\pgfqpoint{2.435549in}{1.287521in}}%
\pgfpathlineto{\pgfqpoint{2.433309in}{1.321900in}}%
\pgfpathlineto{\pgfqpoint{2.440188in}{1.322366in}}%
\pgfpathlineto{\pgfqpoint{2.439649in}{1.329250in}}%
\pgfpathlineto{\pgfqpoint{2.460771in}{1.330448in}}%
\pgfpathlineto{\pgfqpoint{2.467926in}{1.330458in}}%
\pgfpathlineto{\pgfqpoint{2.467492in}{1.337391in}}%
\pgfpathlineto{\pgfqpoint{2.494687in}{1.339376in}}%
\pgfpathlineto{\pgfqpoint{2.495702in}{1.322010in}}%
\pgfpathlineto{\pgfqpoint{2.497373in}{1.294428in}}%
\pgfpathlineto{\pgfqpoint{2.517928in}{1.295771in}}%
\pgfpathlineto{\pgfqpoint{2.519677in}{1.268419in}}%
\pgfpathclose%
\pgfusepath{fill}%
\end{pgfscope}%
\begin{pgfscope}%
\pgfpathrectangle{\pgfqpoint{0.100000in}{0.100000in}}{\pgfqpoint{3.608454in}{2.310000in}}%
\pgfusepath{clip}%
\pgfsetbuttcap%
\pgfsetmiterjoin%
\definecolor{currentfill}{rgb}{0.000000,0.647059,0.676471}%
\pgfsetfillcolor{currentfill}%
\pgfsetlinewidth{0.000000pt}%
\definecolor{currentstroke}{rgb}{0.000000,0.000000,0.000000}%
\pgfsetstrokecolor{currentstroke}%
\pgfsetstrokeopacity{0.000000}%
\pgfsetdash{}{0pt}%
\pgfpathmoveto{\pgfqpoint{2.669844in}{0.772749in}}%
\pgfpathlineto{\pgfqpoint{2.660506in}{0.772013in}}%
\pgfpathlineto{\pgfqpoint{2.655184in}{0.780136in}}%
\pgfpathlineto{\pgfqpoint{2.653773in}{0.795845in}}%
\pgfpathlineto{\pgfqpoint{2.652230in}{0.813085in}}%
\pgfpathlineto{\pgfqpoint{2.650875in}{0.827115in}}%
\pgfpathlineto{\pgfqpoint{2.656320in}{0.836856in}}%
\pgfpathlineto{\pgfqpoint{2.648743in}{0.852393in}}%
\pgfpathlineto{\pgfqpoint{2.647203in}{0.861864in}}%
\pgfpathlineto{\pgfqpoint{2.660793in}{0.862897in}}%
\pgfpathlineto{\pgfqpoint{2.660419in}{0.866374in}}%
\pgfpathlineto{\pgfqpoint{2.680466in}{0.868375in}}%
\pgfpathlineto{\pgfqpoint{2.682505in}{0.872318in}}%
\pgfpathlineto{\pgfqpoint{2.715600in}{0.875575in}}%
\pgfpathlineto{\pgfqpoint{2.717264in}{0.854842in}}%
\pgfpathlineto{\pgfqpoint{2.708963in}{0.847375in}}%
\pgfpathlineto{\pgfqpoint{2.711685in}{0.841617in}}%
\pgfpathlineto{\pgfqpoint{2.713352in}{0.819097in}}%
\pgfpathlineto{\pgfqpoint{2.714111in}{0.812491in}}%
\pgfpathlineto{\pgfqpoint{2.693497in}{0.810278in}}%
\pgfpathlineto{\pgfqpoint{2.692745in}{0.817134in}}%
\pgfpathlineto{\pgfqpoint{2.685856in}{0.816401in}}%
\pgfpathlineto{\pgfqpoint{2.683683in}{0.809309in}}%
\pgfpathlineto{\pgfqpoint{2.685966in}{0.784799in}}%
\pgfpathlineto{\pgfqpoint{2.682532in}{0.784446in}}%
\pgfpathlineto{\pgfqpoint{2.683521in}{0.774139in}}%
\pgfpathlineto{\pgfqpoint{2.669844in}{0.772749in}}%
\pgfpathclose%
\pgfusepath{fill}%
\end{pgfscope}%
\begin{pgfscope}%
\pgfpathrectangle{\pgfqpoint{0.100000in}{0.100000in}}{\pgfqpoint{3.608454in}{2.310000in}}%
\pgfusepath{clip}%
\pgfsetbuttcap%
\pgfsetmiterjoin%
\definecolor{currentfill}{rgb}{0.000000,0.368627,0.815686}%
\pgfsetfillcolor{currentfill}%
\pgfsetlinewidth{0.000000pt}%
\definecolor{currentstroke}{rgb}{0.000000,0.000000,0.000000}%
\pgfsetstrokecolor{currentstroke}%
\pgfsetstrokeopacity{0.000000}%
\pgfsetdash{}{0pt}%
\pgfpathmoveto{\pgfqpoint{0.658484in}{2.308460in}}%
\pgfpathlineto{\pgfqpoint{0.663035in}{2.299847in}}%
\pgfpathlineto{\pgfqpoint{0.657090in}{2.299503in}}%
\pgfpathlineto{\pgfqpoint{0.654017in}{2.294863in}}%
\pgfpathlineto{\pgfqpoint{0.658915in}{2.282360in}}%
\pgfpathlineto{\pgfqpoint{0.663385in}{2.275790in}}%
\pgfpathlineto{\pgfqpoint{0.662941in}{2.269541in}}%
\pgfpathlineto{\pgfqpoint{0.657685in}{2.267521in}}%
\pgfpathlineto{\pgfqpoint{0.656759in}{2.274209in}}%
\pgfpathlineto{\pgfqpoint{0.651475in}{2.279265in}}%
\pgfpathlineto{\pgfqpoint{0.653699in}{2.284940in}}%
\pgfpathlineto{\pgfqpoint{0.647875in}{2.296843in}}%
\pgfpathlineto{\pgfqpoint{0.658484in}{2.308460in}}%
\pgfpathclose%
\pgfusepath{fill}%
\end{pgfscope}%
\begin{pgfscope}%
\pgfpathrectangle{\pgfqpoint{0.100000in}{0.100000in}}{\pgfqpoint{3.608454in}{2.310000in}}%
\pgfusepath{clip}%
\pgfsetbuttcap%
\pgfsetmiterjoin%
\definecolor{currentfill}{rgb}{0.000000,0.368627,0.815686}%
\pgfsetfillcolor{currentfill}%
\pgfsetlinewidth{0.000000pt}%
\definecolor{currentstroke}{rgb}{0.000000,0.000000,0.000000}%
\pgfsetstrokecolor{currentstroke}%
\pgfsetstrokeopacity{0.000000}%
\pgfsetdash{}{0pt}%
\pgfpathmoveto{\pgfqpoint{0.669147in}{2.136525in}}%
\pgfpathlineto{\pgfqpoint{0.631314in}{2.147810in}}%
\pgfpathlineto{\pgfqpoint{0.572748in}{2.166216in}}%
\pgfpathlineto{\pgfqpoint{0.582091in}{2.197407in}}%
\pgfpathlineto{\pgfqpoint{0.593046in}{2.193971in}}%
\pgfpathlineto{\pgfqpoint{0.598094in}{2.209351in}}%
\pgfpathlineto{\pgfqpoint{0.597993in}{2.216816in}}%
\pgfpathlineto{\pgfqpoint{0.583086in}{2.221440in}}%
\pgfpathlineto{\pgfqpoint{0.586531in}{2.235034in}}%
\pgfpathlineto{\pgfqpoint{0.593017in}{2.254746in}}%
\pgfpathlineto{\pgfqpoint{0.595222in}{2.261545in}}%
\pgfpathlineto{\pgfqpoint{0.623636in}{2.252362in}}%
\pgfpathlineto{\pgfqpoint{0.635231in}{2.255252in}}%
\pgfpathlineto{\pgfqpoint{0.639739in}{2.262486in}}%
\pgfpathlineto{\pgfqpoint{0.648220in}{2.267247in}}%
\pgfpathlineto{\pgfqpoint{0.647635in}{2.272305in}}%
\pgfpathlineto{\pgfqpoint{0.652787in}{2.268601in}}%
\pgfpathlineto{\pgfqpoint{0.651158in}{2.248072in}}%
\pgfpathlineto{\pgfqpoint{0.656337in}{2.252428in}}%
\pgfpathlineto{\pgfqpoint{0.657168in}{2.258990in}}%
\pgfpathlineto{\pgfqpoint{0.664956in}{2.268322in}}%
\pgfpathlineto{\pgfqpoint{0.670631in}{2.272725in}}%
\pgfpathlineto{\pgfqpoint{0.666622in}{2.280032in}}%
\pgfpathlineto{\pgfqpoint{0.667712in}{2.286805in}}%
\pgfpathlineto{\pgfqpoint{0.660569in}{2.294722in}}%
\pgfpathlineto{\pgfqpoint{0.669465in}{2.296315in}}%
\pgfpathlineto{\pgfqpoint{0.662370in}{2.309255in}}%
\pgfpathlineto{\pgfqpoint{0.674084in}{2.320271in}}%
\pgfpathlineto{\pgfqpoint{0.672017in}{2.324335in}}%
\pgfpathlineto{\pgfqpoint{0.728782in}{2.306722in}}%
\pgfpathlineto{\pgfqpoint{0.752119in}{2.299760in}}%
\pgfpathlineto{\pgfqpoint{0.760555in}{2.298682in}}%
\pgfpathlineto{\pgfqpoint{0.757956in}{2.296770in}}%
\pgfpathlineto{\pgfqpoint{0.760304in}{2.288239in}}%
\pgfpathlineto{\pgfqpoint{0.752759in}{2.291869in}}%
\pgfpathlineto{\pgfqpoint{0.740255in}{2.288121in}}%
\pgfpathlineto{\pgfqpoint{0.736966in}{2.278404in}}%
\pgfpathlineto{\pgfqpoint{0.741543in}{2.265250in}}%
\pgfpathlineto{\pgfqpoint{0.738201in}{2.260046in}}%
\pgfpathlineto{\pgfqpoint{0.727318in}{2.257197in}}%
\pgfpathlineto{\pgfqpoint{0.721489in}{2.246797in}}%
\pgfpathlineto{\pgfqpoint{0.722977in}{2.230297in}}%
\pgfpathlineto{\pgfqpoint{0.718104in}{2.222638in}}%
\pgfpathlineto{\pgfqpoint{0.712161in}{2.222097in}}%
\pgfpathlineto{\pgfqpoint{0.701083in}{2.215781in}}%
\pgfpathlineto{\pgfqpoint{0.694418in}{2.209381in}}%
\pgfpathlineto{\pgfqpoint{0.694658in}{2.204301in}}%
\pgfpathlineto{\pgfqpoint{0.699138in}{2.201963in}}%
\pgfpathlineto{\pgfqpoint{0.697139in}{2.190297in}}%
\pgfpathlineto{\pgfqpoint{0.680298in}{2.173569in}}%
\pgfpathlineto{\pgfqpoint{0.682663in}{2.166183in}}%
\pgfpathlineto{\pgfqpoint{0.682016in}{2.153295in}}%
\pgfpathlineto{\pgfqpoint{0.675711in}{2.144078in}}%
\pgfpathlineto{\pgfqpoint{0.675514in}{2.134601in}}%
\pgfpathlineto{\pgfqpoint{0.669147in}{2.136525in}}%
\pgfpathclose%
\pgfusepath{fill}%
\end{pgfscope}%
\begin{pgfscope}%
\pgfpathrectangle{\pgfqpoint{0.100000in}{0.100000in}}{\pgfqpoint{3.608454in}{2.310000in}}%
\pgfusepath{clip}%
\pgfsetbuttcap%
\pgfsetmiterjoin%
\definecolor{currentfill}{rgb}{0.000000,0.541176,0.729412}%
\pgfsetfillcolor{currentfill}%
\pgfsetlinewidth{0.000000pt}%
\definecolor{currentstroke}{rgb}{0.000000,0.000000,0.000000}%
\pgfsetstrokecolor{currentstroke}%
\pgfsetstrokeopacity{0.000000}%
\pgfsetdash{}{0pt}%
\pgfpathmoveto{\pgfqpoint{2.615996in}{1.634560in}}%
\pgfpathlineto{\pgfqpoint{2.658513in}{1.638934in}}%
\pgfpathlineto{\pgfqpoint{2.655531in}{1.666251in}}%
\pgfpathlineto{\pgfqpoint{2.682759in}{1.669371in}}%
\pgfpathlineto{\pgfqpoint{2.685918in}{1.641892in}}%
\pgfpathlineto{\pgfqpoint{2.706426in}{1.644227in}}%
\pgfpathlineto{\pgfqpoint{2.710082in}{1.616628in}}%
\pgfpathlineto{\pgfqpoint{2.703247in}{1.616030in}}%
\pgfpathlineto{\pgfqpoint{2.648568in}{1.609914in}}%
\pgfpathlineto{\pgfqpoint{2.621651in}{1.607409in}}%
\pgfpathlineto{\pgfqpoint{2.620303in}{1.620947in}}%
\pgfpathlineto{\pgfqpoint{2.612010in}{1.620134in}}%
\pgfpathlineto{\pgfqpoint{2.615996in}{1.634560in}}%
\pgfpathclose%
\pgfusepath{fill}%
\end{pgfscope}%
\begin{pgfscope}%
\pgfpathrectangle{\pgfqpoint{0.100000in}{0.100000in}}{\pgfqpoint{3.608454in}{2.310000in}}%
\pgfusepath{clip}%
\pgfsetbuttcap%
\pgfsetmiterjoin%
\definecolor{currentfill}{rgb}{0.000000,0.654902,0.672549}%
\pgfsetfillcolor{currentfill}%
\pgfsetlinewidth{0.000000pt}%
\definecolor{currentstroke}{rgb}{0.000000,0.000000,0.000000}%
\pgfsetstrokecolor{currentstroke}%
\pgfsetstrokeopacity{0.000000}%
\pgfsetdash{}{0pt}%
\pgfpathmoveto{\pgfqpoint{1.303565in}{2.014897in}}%
\pgfpathlineto{\pgfqpoint{1.322819in}{2.011773in}}%
\pgfpathlineto{\pgfqpoint{1.316114in}{1.974564in}}%
\pgfpathlineto{\pgfqpoint{1.302428in}{1.976762in}}%
\pgfpathlineto{\pgfqpoint{1.268315in}{1.982523in}}%
\pgfpathlineto{\pgfqpoint{1.261011in}{1.981452in}}%
\pgfpathlineto{\pgfqpoint{1.225529in}{1.987742in}}%
\pgfpathlineto{\pgfqpoint{1.224264in}{1.998009in}}%
\pgfpathlineto{\pgfqpoint{1.228711in}{2.003790in}}%
\pgfpathlineto{\pgfqpoint{1.224493in}{2.008821in}}%
\pgfpathlineto{\pgfqpoint{1.219896in}{2.007840in}}%
\pgfpathlineto{\pgfqpoint{1.216917in}{2.019801in}}%
\pgfpathlineto{\pgfqpoint{1.213019in}{2.026776in}}%
\pgfpathlineto{\pgfqpoint{1.210034in}{2.040419in}}%
\pgfpathlineto{\pgfqpoint{1.203339in}{2.044520in}}%
\pgfpathlineto{\pgfqpoint{1.203686in}{2.050112in}}%
\pgfpathlineto{\pgfqpoint{1.210055in}{2.048940in}}%
\pgfpathlineto{\pgfqpoint{1.211322in}{2.055765in}}%
\pgfpathlineto{\pgfqpoint{1.231866in}{2.052930in}}%
\pgfpathlineto{\pgfqpoint{1.244581in}{2.048177in}}%
\pgfpathlineto{\pgfqpoint{1.251605in}{2.040078in}}%
\pgfpathlineto{\pgfqpoint{1.259993in}{2.025470in}}%
\pgfpathlineto{\pgfqpoint{1.264971in}{2.020151in}}%
\pgfpathlineto{\pgfqpoint{1.292244in}{2.021526in}}%
\pgfpathlineto{\pgfqpoint{1.299591in}{2.015508in}}%
\pgfpathlineto{\pgfqpoint{1.303565in}{2.014897in}}%
\pgfpathclose%
\pgfusepath{fill}%
\end{pgfscope}%
\begin{pgfscope}%
\pgfpathrectangle{\pgfqpoint{0.100000in}{0.100000in}}{\pgfqpoint{3.608454in}{2.310000in}}%
\pgfusepath{clip}%
\pgfsetbuttcap%
\pgfsetmiterjoin%
\definecolor{currentfill}{rgb}{0.000000,0.380392,0.809804}%
\pgfsetfillcolor{currentfill}%
\pgfsetlinewidth{0.000000pt}%
\definecolor{currentstroke}{rgb}{0.000000,0.000000,0.000000}%
\pgfsetstrokecolor{currentstroke}%
\pgfsetstrokeopacity{0.000000}%
\pgfsetdash{}{0pt}%
\pgfpathmoveto{\pgfqpoint{3.343167in}{1.599400in}}%
\pgfpathlineto{\pgfqpoint{3.345889in}{1.596696in}}%
\pgfpathlineto{\pgfqpoint{3.339300in}{1.581420in}}%
\pgfpathlineto{\pgfqpoint{3.327731in}{1.570641in}}%
\pgfpathlineto{\pgfqpoint{3.330022in}{1.567556in}}%
\pgfpathlineto{\pgfqpoint{3.313073in}{1.572820in}}%
\pgfpathlineto{\pgfqpoint{3.303159in}{1.582589in}}%
\pgfpathlineto{\pgfqpoint{3.297442in}{1.582891in}}%
\pgfpathlineto{\pgfqpoint{3.294732in}{1.592344in}}%
\pgfpathlineto{\pgfqpoint{3.286080in}{1.595797in}}%
\pgfpathlineto{\pgfqpoint{3.284662in}{1.609183in}}%
\pgfpathlineto{\pgfqpoint{3.288120in}{1.611032in}}%
\pgfpathlineto{\pgfqpoint{3.289903in}{1.618119in}}%
\pgfpathlineto{\pgfqpoint{3.284382in}{1.624557in}}%
\pgfpathlineto{\pgfqpoint{3.294577in}{1.644062in}}%
\pgfpathlineto{\pgfqpoint{3.295750in}{1.652950in}}%
\pgfpathlineto{\pgfqpoint{3.302353in}{1.660141in}}%
\pgfpathlineto{\pgfqpoint{3.332839in}{1.649737in}}%
\pgfpathlineto{\pgfqpoint{3.344220in}{1.667169in}}%
\pgfpathlineto{\pgfqpoint{3.346226in}{1.661401in}}%
\pgfpathlineto{\pgfqpoint{3.353043in}{1.654103in}}%
\pgfpathlineto{\pgfqpoint{3.355371in}{1.641928in}}%
\pgfpathlineto{\pgfqpoint{3.352075in}{1.614317in}}%
\pgfpathlineto{\pgfqpoint{3.344569in}{1.612129in}}%
\pgfpathlineto{\pgfqpoint{3.343167in}{1.599400in}}%
\pgfpathclose%
\pgfusepath{fill}%
\end{pgfscope}%
\begin{pgfscope}%
\pgfpathrectangle{\pgfqpoint{0.100000in}{0.100000in}}{\pgfqpoint{3.608454in}{2.310000in}}%
\pgfusepath{clip}%
\pgfsetbuttcap%
\pgfsetmiterjoin%
\definecolor{currentfill}{rgb}{0.000000,0.670588,0.664706}%
\pgfsetfillcolor{currentfill}%
\pgfsetlinewidth{0.000000pt}%
\definecolor{currentstroke}{rgb}{0.000000,0.000000,0.000000}%
\pgfsetstrokecolor{currentstroke}%
\pgfsetstrokeopacity{0.000000}%
\pgfsetdash{}{0pt}%
\pgfpathmoveto{\pgfqpoint{2.893588in}{0.882495in}}%
\pgfpathlineto{\pgfqpoint{2.910191in}{0.894634in}}%
\pgfpathlineto{\pgfqpoint{2.909397in}{0.902008in}}%
\pgfpathlineto{\pgfqpoint{2.903959in}{0.905713in}}%
\pgfpathlineto{\pgfqpoint{2.899980in}{0.912387in}}%
\pgfpathlineto{\pgfqpoint{2.899562in}{0.918887in}}%
\pgfpathlineto{\pgfqpoint{2.905128in}{0.926113in}}%
\pgfpathlineto{\pgfqpoint{2.910167in}{0.929143in}}%
\pgfpathlineto{\pgfqpoint{2.911845in}{0.934034in}}%
\pgfpathlineto{\pgfqpoint{2.918368in}{0.938239in}}%
\pgfpathlineto{\pgfqpoint{2.919634in}{0.933284in}}%
\pgfpathlineto{\pgfqpoint{2.925307in}{0.930408in}}%
\pgfpathlineto{\pgfqpoint{2.933286in}{0.922296in}}%
\pgfpathlineto{\pgfqpoint{2.937493in}{0.915146in}}%
\pgfpathlineto{\pgfqpoint{2.937410in}{0.908377in}}%
\pgfpathlineto{\pgfqpoint{2.944266in}{0.903605in}}%
\pgfpathlineto{\pgfqpoint{2.940759in}{0.897010in}}%
\pgfpathlineto{\pgfqpoint{2.943276in}{0.894232in}}%
\pgfpathlineto{\pgfqpoint{2.935640in}{0.887719in}}%
\pgfpathlineto{\pgfqpoint{2.932829in}{0.881990in}}%
\pgfpathlineto{\pgfqpoint{2.928863in}{0.867621in}}%
\pgfpathlineto{\pgfqpoint{2.930056in}{0.865279in}}%
\pgfpathlineto{\pgfqpoint{2.922187in}{0.854988in}}%
\pgfpathlineto{\pgfqpoint{2.913824in}{0.850086in}}%
\pgfpathlineto{\pgfqpoint{2.901077in}{0.870542in}}%
\pgfpathlineto{\pgfqpoint{2.893588in}{0.882495in}}%
\pgfpathclose%
\pgfusepath{fill}%
\end{pgfscope}%
\begin{pgfscope}%
\pgfpathrectangle{\pgfqpoint{0.100000in}{0.100000in}}{\pgfqpoint{3.608454in}{2.310000in}}%
\pgfusepath{clip}%
\pgfsetbuttcap%
\pgfsetmiterjoin%
\definecolor{currentfill}{rgb}{0.000000,0.509804,0.745098}%
\pgfsetfillcolor{currentfill}%
\pgfsetlinewidth{0.000000pt}%
\definecolor{currentstroke}{rgb}{0.000000,0.000000,0.000000}%
\pgfsetstrokecolor{currentstroke}%
\pgfsetstrokeopacity{0.000000}%
\pgfsetdash{}{0pt}%
\pgfpathmoveto{\pgfqpoint{2.515690in}{0.998385in}}%
\pgfpathlineto{\pgfqpoint{2.512313in}{0.997036in}}%
\pgfpathlineto{\pgfqpoint{2.512839in}{0.989031in}}%
\pgfpathlineto{\pgfqpoint{2.503832in}{0.987282in}}%
\pgfpathlineto{\pgfqpoint{2.483100in}{0.985946in}}%
\pgfpathlineto{\pgfqpoint{2.482530in}{0.995135in}}%
\pgfpathlineto{\pgfqpoint{2.478692in}{1.001827in}}%
\pgfpathlineto{\pgfqpoint{2.477335in}{1.022510in}}%
\pgfpathlineto{\pgfqpoint{2.473863in}{1.022295in}}%
\pgfpathlineto{\pgfqpoint{2.473129in}{1.034113in}}%
\pgfpathlineto{\pgfqpoint{2.483127in}{1.034794in}}%
\pgfpathlineto{\pgfqpoint{2.507414in}{1.036455in}}%
\pgfpathlineto{\pgfqpoint{2.508135in}{1.025655in}}%
\pgfpathlineto{\pgfqpoint{2.512709in}{1.025980in}}%
\pgfpathlineto{\pgfqpoint{2.515503in}{1.018069in}}%
\pgfpathlineto{\pgfqpoint{2.515690in}{0.998385in}}%
\pgfpathclose%
\pgfusepath{fill}%
\end{pgfscope}%
\begin{pgfscope}%
\pgfpathrectangle{\pgfqpoint{0.100000in}{0.100000in}}{\pgfqpoint{3.608454in}{2.310000in}}%
\pgfusepath{clip}%
\pgfsetbuttcap%
\pgfsetmiterjoin%
\definecolor{currentfill}{rgb}{0.000000,0.564706,0.717647}%
\pgfsetfillcolor{currentfill}%
\pgfsetlinewidth{0.000000pt}%
\definecolor{currentstroke}{rgb}{0.000000,0.000000,0.000000}%
\pgfsetstrokecolor{currentstroke}%
\pgfsetstrokeopacity{0.000000}%
\pgfsetdash{}{0pt}%
\pgfpathmoveto{\pgfqpoint{2.595636in}{0.988190in}}%
\pgfpathlineto{\pgfqpoint{2.553484in}{0.986177in}}%
\pgfpathlineto{\pgfqpoint{2.553769in}{0.997504in}}%
\pgfpathlineto{\pgfqpoint{2.554995in}{1.031701in}}%
\pgfpathlineto{\pgfqpoint{2.547499in}{1.040339in}}%
\pgfpathlineto{\pgfqpoint{2.561622in}{1.041248in}}%
\pgfpathlineto{\pgfqpoint{2.560066in}{1.064151in}}%
\pgfpathlineto{\pgfqpoint{2.556697in}{1.071568in}}%
\pgfpathlineto{\pgfqpoint{2.559562in}{1.077181in}}%
\pgfpathlineto{\pgfqpoint{2.575892in}{1.080309in}}%
\pgfpathlineto{\pgfqpoint{2.580566in}{1.079844in}}%
\pgfpathlineto{\pgfqpoint{2.585581in}{1.074349in}}%
\pgfpathlineto{\pgfqpoint{2.601253in}{1.080261in}}%
\pgfpathlineto{\pgfqpoint{2.607711in}{1.078660in}}%
\pgfpathlineto{\pgfqpoint{2.608880in}{1.074941in}}%
\pgfpathlineto{\pgfqpoint{2.610505in}{1.048649in}}%
\pgfpathlineto{\pgfqpoint{2.611745in}{1.044790in}}%
\pgfpathlineto{\pgfqpoint{2.612717in}{1.030317in}}%
\pgfpathlineto{\pgfqpoint{2.607222in}{1.025513in}}%
\pgfpathlineto{\pgfqpoint{2.599104in}{1.027892in}}%
\pgfpathlineto{\pgfqpoint{2.599703in}{1.020694in}}%
\pgfpathlineto{\pgfqpoint{2.593933in}{1.008935in}}%
\pgfpathlineto{\pgfqpoint{2.595636in}{0.988190in}}%
\pgfpathclose%
\pgfusepath{fill}%
\end{pgfscope}%
\begin{pgfscope}%
\pgfpathrectangle{\pgfqpoint{0.100000in}{0.100000in}}{\pgfqpoint{3.608454in}{2.310000in}}%
\pgfusepath{clip}%
\pgfsetbuttcap%
\pgfsetmiterjoin%
\definecolor{currentfill}{rgb}{0.000000,0.549020,0.725490}%
\pgfsetfillcolor{currentfill}%
\pgfsetlinewidth{0.000000pt}%
\definecolor{currentstroke}{rgb}{0.000000,0.000000,0.000000}%
\pgfsetstrokecolor{currentstroke}%
\pgfsetstrokeopacity{0.000000}%
\pgfsetdash{}{0pt}%
\pgfpathmoveto{\pgfqpoint{1.975277in}{0.471174in}}%
\pgfpathlineto{\pgfqpoint{1.976992in}{0.475338in}}%
\pgfpathlineto{\pgfqpoint{1.978118in}{0.480166in}}%
\pgfpathlineto{\pgfqpoint{1.984185in}{0.481426in}}%
\pgfpathlineto{\pgfqpoint{1.994367in}{0.489945in}}%
\pgfpathlineto{\pgfqpoint{1.996428in}{0.487179in}}%
\pgfpathlineto{\pgfqpoint{1.984727in}{0.479179in}}%
\pgfpathlineto{\pgfqpoint{1.975277in}{0.471174in}}%
\pgfpathclose%
\pgfusepath{fill}%
\end{pgfscope}%
\begin{pgfscope}%
\pgfpathrectangle{\pgfqpoint{0.100000in}{0.100000in}}{\pgfqpoint{3.608454in}{2.310000in}}%
\pgfusepath{clip}%
\pgfsetbuttcap%
\pgfsetmiterjoin%
\definecolor{currentfill}{rgb}{0.000000,0.549020,0.725490}%
\pgfsetfillcolor{currentfill}%
\pgfsetlinewidth{0.000000pt}%
\definecolor{currentstroke}{rgb}{0.000000,0.000000,0.000000}%
\pgfsetstrokecolor{currentstroke}%
\pgfsetstrokeopacity{0.000000}%
\pgfsetdash{}{0pt}%
\pgfpathmoveto{\pgfqpoint{2.013559in}{0.517174in}}%
\pgfpathlineto{\pgfqpoint{2.002833in}{0.511803in}}%
\pgfpathlineto{\pgfqpoint{1.996542in}{0.517508in}}%
\pgfpathlineto{\pgfqpoint{1.995156in}{0.524126in}}%
\pgfpathlineto{\pgfqpoint{1.989574in}{0.521754in}}%
\pgfpathlineto{\pgfqpoint{1.994165in}{0.512797in}}%
\pgfpathlineto{\pgfqpoint{2.008150in}{0.500988in}}%
\pgfpathlineto{\pgfqpoint{1.995428in}{0.494588in}}%
\pgfpathlineto{\pgfqpoint{1.992059in}{0.491469in}}%
\pgfpathlineto{\pgfqpoint{1.981043in}{0.498344in}}%
\pgfpathlineto{\pgfqpoint{1.973994in}{0.506774in}}%
\pgfpathlineto{\pgfqpoint{1.966594in}{0.506597in}}%
\pgfpathlineto{\pgfqpoint{1.962232in}{0.510712in}}%
\pgfpathlineto{\pgfqpoint{1.954355in}{0.511185in}}%
\pgfpathlineto{\pgfqpoint{1.938836in}{0.498314in}}%
\pgfpathlineto{\pgfqpoint{1.926522in}{0.503678in}}%
\pgfpathlineto{\pgfqpoint{1.926187in}{0.508771in}}%
\pgfpathlineto{\pgfqpoint{1.916523in}{0.511256in}}%
\pgfpathlineto{\pgfqpoint{1.911040in}{0.521565in}}%
\pgfpathlineto{\pgfqpoint{1.925742in}{0.532805in}}%
\pgfpathlineto{\pgfqpoint{1.913527in}{0.548723in}}%
\pgfpathlineto{\pgfqpoint{1.923720in}{0.556660in}}%
\pgfpathlineto{\pgfqpoint{1.950391in}{0.577980in}}%
\pgfpathlineto{\pgfqpoint{1.953286in}{0.593627in}}%
\pgfpathlineto{\pgfqpoint{1.957611in}{0.597327in}}%
\pgfpathlineto{\pgfqpoint{1.976342in}{0.597283in}}%
\pgfpathlineto{\pgfqpoint{1.981908in}{0.592727in}}%
\pgfpathlineto{\pgfqpoint{1.997763in}{0.573057in}}%
\pgfpathlineto{\pgfqpoint{1.992065in}{0.566132in}}%
\pgfpathlineto{\pgfqpoint{2.014894in}{0.542939in}}%
\pgfpathlineto{\pgfqpoint{2.013559in}{0.517174in}}%
\pgfpathclose%
\pgfusepath{fill}%
\end{pgfscope}%
\begin{pgfscope}%
\pgfpathrectangle{\pgfqpoint{0.100000in}{0.100000in}}{\pgfqpoint{3.608454in}{2.310000in}}%
\pgfusepath{clip}%
\pgfsetbuttcap%
\pgfsetmiterjoin%
\definecolor{currentfill}{rgb}{0.000000,0.882353,0.558824}%
\pgfsetfillcolor{currentfill}%
\pgfsetlinewidth{0.000000pt}%
\definecolor{currentstroke}{rgb}{0.000000,0.000000,0.000000}%
\pgfsetstrokecolor{currentstroke}%
\pgfsetstrokeopacity{0.000000}%
\pgfsetdash{}{0pt}%
\pgfpathmoveto{\pgfqpoint{2.362240in}{0.870230in}}%
\pgfpathlineto{\pgfqpoint{2.344274in}{0.869656in}}%
\pgfpathlineto{\pgfqpoint{2.342635in}{0.869578in}}%
\pgfpathlineto{\pgfqpoint{2.341651in}{0.899996in}}%
\pgfpathlineto{\pgfqpoint{2.306319in}{0.899456in}}%
\pgfpathlineto{\pgfqpoint{2.307139in}{0.891576in}}%
\pgfpathlineto{\pgfqpoint{2.297645in}{0.883786in}}%
\pgfpathlineto{\pgfqpoint{2.297141in}{0.880461in}}%
\pgfpathlineto{\pgfqpoint{2.286293in}{0.889498in}}%
\pgfpathlineto{\pgfqpoint{2.282169in}{0.889435in}}%
\pgfpathlineto{\pgfqpoint{2.280766in}{0.896956in}}%
\pgfpathlineto{\pgfqpoint{2.284812in}{0.906848in}}%
\pgfpathlineto{\pgfqpoint{2.282183in}{0.933447in}}%
\pgfpathlineto{\pgfqpoint{2.274161in}{0.944083in}}%
\pgfpathlineto{\pgfqpoint{2.272954in}{0.950762in}}%
\pgfpathlineto{\pgfqpoint{2.281822in}{0.951221in}}%
\pgfpathlineto{\pgfqpoint{2.288597in}{0.951635in}}%
\pgfpathlineto{\pgfqpoint{2.289533in}{0.964998in}}%
\pgfpathlineto{\pgfqpoint{2.289234in}{0.985811in}}%
\pgfpathlineto{\pgfqpoint{2.322083in}{0.986111in}}%
\pgfpathlineto{\pgfqpoint{2.322365in}{0.966382in}}%
\pgfpathlineto{\pgfqpoint{2.333604in}{0.966070in}}%
\pgfpathlineto{\pgfqpoint{2.335508in}{0.959654in}}%
\pgfpathlineto{\pgfqpoint{2.339665in}{0.959110in}}%
\pgfpathlineto{\pgfqpoint{2.341075in}{0.950676in}}%
\pgfpathlineto{\pgfqpoint{2.349626in}{0.945687in}}%
\pgfpathlineto{\pgfqpoint{2.361253in}{0.945877in}}%
\pgfpathlineto{\pgfqpoint{2.356734in}{0.953261in}}%
\pgfpathlineto{\pgfqpoint{2.361668in}{0.958818in}}%
\pgfpathlineto{\pgfqpoint{2.372346in}{0.959262in}}%
\pgfpathlineto{\pgfqpoint{2.377959in}{0.956142in}}%
\pgfpathlineto{\pgfqpoint{2.376340in}{0.951894in}}%
\pgfpathlineto{\pgfqpoint{2.364126in}{0.947506in}}%
\pgfpathlineto{\pgfqpoint{2.369367in}{0.943780in}}%
\pgfpathlineto{\pgfqpoint{2.365618in}{0.938943in}}%
\pgfpathlineto{\pgfqpoint{2.370618in}{0.931411in}}%
\pgfpathlineto{\pgfqpoint{2.361535in}{0.931873in}}%
\pgfpathlineto{\pgfqpoint{2.356747in}{0.922152in}}%
\pgfpathlineto{\pgfqpoint{2.362379in}{0.917175in}}%
\pgfpathlineto{\pgfqpoint{2.356517in}{0.914773in}}%
\pgfpathlineto{\pgfqpoint{2.362531in}{0.900120in}}%
\pgfpathlineto{\pgfqpoint{2.362887in}{0.893699in}}%
\pgfpathlineto{\pgfqpoint{2.366452in}{0.887645in}}%
\pgfpathlineto{\pgfqpoint{2.366771in}{0.881093in}}%
\pgfpathlineto{\pgfqpoint{2.362240in}{0.870230in}}%
\pgfpathclose%
\pgfusepath{fill}%
\end{pgfscope}%
\begin{pgfscope}%
\pgfpathrectangle{\pgfqpoint{0.100000in}{0.100000in}}{\pgfqpoint{3.608454in}{2.310000in}}%
\pgfusepath{clip}%
\pgfsetbuttcap%
\pgfsetmiterjoin%
\definecolor{currentfill}{rgb}{0.000000,0.686275,0.656863}%
\pgfsetfillcolor{currentfill}%
\pgfsetlinewidth{0.000000pt}%
\definecolor{currentstroke}{rgb}{0.000000,0.000000,0.000000}%
\pgfsetstrokecolor{currentstroke}%
\pgfsetstrokeopacity{0.000000}%
\pgfsetdash{}{0pt}%
\pgfpathmoveto{\pgfqpoint{3.332649in}{1.866641in}}%
\pgfpathlineto{\pgfqpoint{3.296075in}{1.852707in}}%
\pgfpathlineto{\pgfqpoint{3.295690in}{1.856839in}}%
\pgfpathlineto{\pgfqpoint{3.286142in}{1.855787in}}%
\pgfpathlineto{\pgfqpoint{3.277231in}{1.863034in}}%
\pgfpathlineto{\pgfqpoint{3.280930in}{1.867522in}}%
\pgfpathlineto{\pgfqpoint{3.276788in}{1.878840in}}%
\pgfpathlineto{\pgfqpoint{3.263105in}{1.873951in}}%
\pgfpathlineto{\pgfqpoint{3.262366in}{1.879593in}}%
\pgfpathlineto{\pgfqpoint{3.242124in}{1.938669in}}%
\pgfpathlineto{\pgfqpoint{3.236756in}{1.940890in}}%
\pgfpathlineto{\pgfqpoint{3.268366in}{1.947882in}}%
\pgfpathlineto{\pgfqpoint{3.312126in}{1.959898in}}%
\pgfpathlineto{\pgfqpoint{3.314088in}{1.952894in}}%
\pgfpathlineto{\pgfqpoint{3.313379in}{1.946070in}}%
\pgfpathlineto{\pgfqpoint{3.316446in}{1.944043in}}%
\pgfpathlineto{\pgfqpoint{3.316899in}{1.928969in}}%
\pgfpathlineto{\pgfqpoint{3.323703in}{1.920977in}}%
\pgfpathlineto{\pgfqpoint{3.324626in}{1.910739in}}%
\pgfpathlineto{\pgfqpoint{3.327709in}{1.903478in}}%
\pgfpathlineto{\pgfqpoint{3.324873in}{1.895982in}}%
\pgfpathlineto{\pgfqpoint{3.324985in}{1.884098in}}%
\pgfpathlineto{\pgfqpoint{3.332649in}{1.866641in}}%
\pgfpathclose%
\pgfusepath{fill}%
\end{pgfscope}%
\begin{pgfscope}%
\pgfpathrectangle{\pgfqpoint{0.100000in}{0.100000in}}{\pgfqpoint{3.608454in}{2.310000in}}%
\pgfusepath{clip}%
\pgfsetbuttcap%
\pgfsetmiterjoin%
\definecolor{currentfill}{rgb}{0.000000,0.556863,0.721569}%
\pgfsetfillcolor{currentfill}%
\pgfsetlinewidth{0.000000pt}%
\definecolor{currentstroke}{rgb}{0.000000,0.000000,0.000000}%
\pgfsetstrokecolor{currentstroke}%
\pgfsetstrokeopacity{0.000000}%
\pgfsetdash{}{0pt}%
\pgfpathmoveto{\pgfqpoint{1.678348in}{0.726376in}}%
\pgfpathlineto{\pgfqpoint{1.644176in}{0.728867in}}%
\pgfpathlineto{\pgfqpoint{1.650167in}{0.808902in}}%
\pgfpathlineto{\pgfqpoint{1.655619in}{0.808558in}}%
\pgfpathlineto{\pgfqpoint{1.658603in}{0.843265in}}%
\pgfpathlineto{\pgfqpoint{1.624223in}{0.845743in}}%
\pgfpathlineto{\pgfqpoint{1.626580in}{0.880308in}}%
\pgfpathlineto{\pgfqpoint{1.635380in}{0.879674in}}%
\pgfpathlineto{\pgfqpoint{1.695447in}{0.875598in}}%
\pgfpathlineto{\pgfqpoint{1.692995in}{0.840738in}}%
\pgfpathlineto{\pgfqpoint{1.690000in}{0.806092in}}%
\pgfpathlineto{\pgfqpoint{1.684601in}{0.806456in}}%
\pgfpathlineto{\pgfqpoint{1.681284in}{0.762092in}}%
\pgfpathlineto{\pgfqpoint{1.678348in}{0.726376in}}%
\pgfpathclose%
\pgfusepath{fill}%
\end{pgfscope}%
\begin{pgfscope}%
\pgfpathrectangle{\pgfqpoint{0.100000in}{0.100000in}}{\pgfqpoint{3.608454in}{2.310000in}}%
\pgfusepath{clip}%
\pgfsetbuttcap%
\pgfsetmiterjoin%
\definecolor{currentfill}{rgb}{0.000000,0.827451,0.586275}%
\pgfsetfillcolor{currentfill}%
\pgfsetlinewidth{0.000000pt}%
\definecolor{currentstroke}{rgb}{0.000000,0.000000,0.000000}%
\pgfsetstrokecolor{currentstroke}%
\pgfsetstrokeopacity{0.000000}%
\pgfsetdash{}{0pt}%
\pgfpathmoveto{\pgfqpoint{0.779307in}{0.362404in}}%
\pgfpathlineto{\pgfqpoint{0.779048in}{0.364878in}}%
\pgfpathlineto{\pgfqpoint{0.781778in}{0.364420in}}%
\pgfpathlineto{\pgfqpoint{0.784239in}{0.359816in}}%
\pgfpathlineto{\pgfqpoint{0.779307in}{0.362404in}}%
\pgfpathclose%
\pgfusepath{fill}%
\end{pgfscope}%
\begin{pgfscope}%
\pgfpathrectangle{\pgfqpoint{0.100000in}{0.100000in}}{\pgfqpoint{3.608454in}{2.310000in}}%
\pgfusepath{clip}%
\pgfsetbuttcap%
\pgfsetmiterjoin%
\definecolor{currentfill}{rgb}{0.000000,0.827451,0.586275}%
\pgfsetfillcolor{currentfill}%
\pgfsetlinewidth{0.000000pt}%
\definecolor{currentstroke}{rgb}{0.000000,0.000000,0.000000}%
\pgfsetstrokecolor{currentstroke}%
\pgfsetstrokeopacity{0.000000}%
\pgfsetdash{}{0pt}%
\pgfpathmoveto{\pgfqpoint{0.771381in}{0.365345in}}%
\pgfpathlineto{\pgfqpoint{0.773191in}{0.367159in}}%
\pgfpathlineto{\pgfqpoint{0.777403in}{0.367463in}}%
\pgfpathlineto{\pgfqpoint{0.776562in}{0.365684in}}%
\pgfpathlineto{\pgfqpoint{0.771381in}{0.365345in}}%
\pgfpathclose%
\pgfusepath{fill}%
\end{pgfscope}%
\begin{pgfscope}%
\pgfpathrectangle{\pgfqpoint{0.100000in}{0.100000in}}{\pgfqpoint{3.608454in}{2.310000in}}%
\pgfusepath{clip}%
\pgfsetbuttcap%
\pgfsetmiterjoin%
\definecolor{currentfill}{rgb}{0.000000,0.827451,0.586275}%
\pgfsetfillcolor{currentfill}%
\pgfsetlinewidth{0.000000pt}%
\definecolor{currentstroke}{rgb}{0.000000,0.000000,0.000000}%
\pgfsetstrokecolor{currentstroke}%
\pgfsetstrokeopacity{0.000000}%
\pgfsetdash{}{0pt}%
\pgfpathmoveto{\pgfqpoint{0.764308in}{0.397983in}}%
\pgfpathlineto{\pgfqpoint{0.760327in}{0.401259in}}%
\pgfpathlineto{\pgfqpoint{0.761612in}{0.402864in}}%
\pgfpathlineto{\pgfqpoint{0.765003in}{0.403792in}}%
\pgfpathlineto{\pgfqpoint{0.766993in}{0.403276in}}%
\pgfpathlineto{\pgfqpoint{0.768215in}{0.404903in}}%
\pgfpathlineto{\pgfqpoint{0.771380in}{0.405564in}}%
\pgfpathlineto{\pgfqpoint{0.773253in}{0.404050in}}%
\pgfpathlineto{\pgfqpoint{0.777652in}{0.403284in}}%
\pgfpathlineto{\pgfqpoint{0.777113in}{0.402156in}}%
\pgfpathlineto{\pgfqpoint{0.772733in}{0.402967in}}%
\pgfpathlineto{\pgfqpoint{0.767839in}{0.402629in}}%
\pgfpathlineto{\pgfqpoint{0.767479in}{0.399727in}}%
\pgfpathlineto{\pgfqpoint{0.764308in}{0.397983in}}%
\pgfpathclose%
\pgfusepath{fill}%
\end{pgfscope}%
\begin{pgfscope}%
\pgfpathrectangle{\pgfqpoint{0.100000in}{0.100000in}}{\pgfqpoint{3.608454in}{2.310000in}}%
\pgfusepath{clip}%
\pgfsetbuttcap%
\pgfsetmiterjoin%
\definecolor{currentfill}{rgb}{0.000000,0.827451,0.586275}%
\pgfsetfillcolor{currentfill}%
\pgfsetlinewidth{0.000000pt}%
\definecolor{currentstroke}{rgb}{0.000000,0.000000,0.000000}%
\pgfsetstrokecolor{currentstroke}%
\pgfsetstrokeopacity{0.000000}%
\pgfsetdash{}{0pt}%
\pgfpathmoveto{\pgfqpoint{0.833159in}{0.390083in}}%
\pgfpathlineto{\pgfqpoint{0.831710in}{0.387820in}}%
\pgfpathlineto{\pgfqpoint{0.834237in}{0.387444in}}%
\pgfpathlineto{\pgfqpoint{0.834627in}{0.385754in}}%
\pgfpathlineto{\pgfqpoint{0.836745in}{0.385201in}}%
\pgfpathlineto{\pgfqpoint{0.837269in}{0.381931in}}%
\pgfpathlineto{\pgfqpoint{0.835377in}{0.380107in}}%
\pgfpathlineto{\pgfqpoint{0.832974in}{0.379753in}}%
\pgfpathlineto{\pgfqpoint{0.827040in}{0.383181in}}%
\pgfpathlineto{\pgfqpoint{0.822060in}{0.382223in}}%
\pgfpathlineto{\pgfqpoint{0.818304in}{0.385966in}}%
\pgfpathlineto{\pgfqpoint{0.817822in}{0.388703in}}%
\pgfpathlineto{\pgfqpoint{0.819792in}{0.390087in}}%
\pgfpathlineto{\pgfqpoint{0.823330in}{0.389436in}}%
\pgfpathlineto{\pgfqpoint{0.823863in}{0.391532in}}%
\pgfpathlineto{\pgfqpoint{0.826549in}{0.387455in}}%
\pgfpathlineto{\pgfqpoint{0.828702in}{0.391954in}}%
\pgfpathlineto{\pgfqpoint{0.832601in}{0.391664in}}%
\pgfpathlineto{\pgfqpoint{0.834521in}{0.394379in}}%
\pgfpathlineto{\pgfqpoint{0.835969in}{0.392848in}}%
\pgfpathlineto{\pgfqpoint{0.836007in}{0.391059in}}%
\pgfpathlineto{\pgfqpoint{0.833159in}{0.390083in}}%
\pgfpathclose%
\pgfusepath{fill}%
\end{pgfscope}%
\begin{pgfscope}%
\pgfpathrectangle{\pgfqpoint{0.100000in}{0.100000in}}{\pgfqpoint{3.608454in}{2.310000in}}%
\pgfusepath{clip}%
\pgfsetbuttcap%
\pgfsetmiterjoin%
\definecolor{currentfill}{rgb}{0.000000,0.827451,0.586275}%
\pgfsetfillcolor{currentfill}%
\pgfsetlinewidth{0.000000pt}%
\definecolor{currentstroke}{rgb}{0.000000,0.000000,0.000000}%
\pgfsetstrokecolor{currentstroke}%
\pgfsetstrokeopacity{0.000000}%
\pgfsetdash{}{0pt}%
\pgfpathmoveto{\pgfqpoint{0.824220in}{0.377345in}}%
\pgfpathlineto{\pgfqpoint{0.824880in}{0.374202in}}%
\pgfpathlineto{\pgfqpoint{0.822730in}{0.374227in}}%
\pgfpathlineto{\pgfqpoint{0.819725in}{0.370152in}}%
\pgfpathlineto{\pgfqpoint{0.823688in}{0.368053in}}%
\pgfpathlineto{\pgfqpoint{0.818723in}{0.365758in}}%
\pgfpathlineto{\pgfqpoint{0.816629in}{0.366618in}}%
\pgfpathlineto{\pgfqpoint{0.815348in}{0.369220in}}%
\pgfpathlineto{\pgfqpoint{0.813787in}{0.368393in}}%
\pgfpathlineto{\pgfqpoint{0.813447in}{0.366368in}}%
\pgfpathlineto{\pgfqpoint{0.811430in}{0.365767in}}%
\pgfpathlineto{\pgfqpoint{0.803663in}{0.368137in}}%
\pgfpathlineto{\pgfqpoint{0.804036in}{0.365808in}}%
\pgfpathlineto{\pgfqpoint{0.801470in}{0.365130in}}%
\pgfpathlineto{\pgfqpoint{0.797979in}{0.368165in}}%
\pgfpathlineto{\pgfqpoint{0.796078in}{0.368564in}}%
\pgfpathlineto{\pgfqpoint{0.791325in}{0.365921in}}%
\pgfpathlineto{\pgfqpoint{0.786527in}{0.365586in}}%
\pgfpathlineto{\pgfqpoint{0.787213in}{0.367927in}}%
\pgfpathlineto{\pgfqpoint{0.792491in}{0.369127in}}%
\pgfpathlineto{\pgfqpoint{0.792763in}{0.372545in}}%
\pgfpathlineto{\pgfqpoint{0.789745in}{0.372900in}}%
\pgfpathlineto{\pgfqpoint{0.786003in}{0.370933in}}%
\pgfpathlineto{\pgfqpoint{0.784989in}{0.376973in}}%
\pgfpathlineto{\pgfqpoint{0.787640in}{0.381629in}}%
\pgfpathlineto{\pgfqpoint{0.786375in}{0.384997in}}%
\pgfpathlineto{\pgfqpoint{0.792271in}{0.388607in}}%
\pgfpathlineto{\pgfqpoint{0.796911in}{0.388654in}}%
\pgfpathlineto{\pgfqpoint{0.799734in}{0.387836in}}%
\pgfpathlineto{\pgfqpoint{0.801573in}{0.385770in}}%
\pgfpathlineto{\pgfqpoint{0.799429in}{0.380434in}}%
\pgfpathlineto{\pgfqpoint{0.803009in}{0.382399in}}%
\pgfpathlineto{\pgfqpoint{0.804618in}{0.387969in}}%
\pgfpathlineto{\pgfqpoint{0.807031in}{0.389222in}}%
\pgfpathlineto{\pgfqpoint{0.809285in}{0.388827in}}%
\pgfpathlineto{\pgfqpoint{0.809591in}{0.384951in}}%
\pgfpathlineto{\pgfqpoint{0.811916in}{0.384419in}}%
\pgfpathlineto{\pgfqpoint{0.812085in}{0.387811in}}%
\pgfpathlineto{\pgfqpoint{0.816863in}{0.385495in}}%
\pgfpathlineto{\pgfqpoint{0.818200in}{0.383011in}}%
\pgfpathlineto{\pgfqpoint{0.820527in}{0.380666in}}%
\pgfpathlineto{\pgfqpoint{0.818276in}{0.378662in}}%
\pgfpathlineto{\pgfqpoint{0.824441in}{0.378263in}}%
\pgfpathlineto{\pgfqpoint{0.824220in}{0.377345in}}%
\pgfpathclose%
\pgfusepath{fill}%
\end{pgfscope}%
\begin{pgfscope}%
\pgfpathrectangle{\pgfqpoint{0.100000in}{0.100000in}}{\pgfqpoint{3.608454in}{2.310000in}}%
\pgfusepath{clip}%
\pgfsetbuttcap%
\pgfsetmiterjoin%
\definecolor{currentfill}{rgb}{0.000000,0.827451,0.586275}%
\pgfsetfillcolor{currentfill}%
\pgfsetlinewidth{0.000000pt}%
\definecolor{currentstroke}{rgb}{0.000000,0.000000,0.000000}%
\pgfsetstrokecolor{currentstroke}%
\pgfsetstrokeopacity{0.000000}%
\pgfsetdash{}{0pt}%
\pgfpathmoveto{\pgfqpoint{0.831242in}{0.406233in}}%
\pgfpathlineto{\pgfqpoint{0.829444in}{0.406598in}}%
\pgfpathlineto{\pgfqpoint{0.826973in}{0.404781in}}%
\pgfpathlineto{\pgfqpoint{0.821584in}{0.404692in}}%
\pgfpathlineto{\pgfqpoint{0.820126in}{0.405792in}}%
\pgfpathlineto{\pgfqpoint{0.817124in}{0.404505in}}%
\pgfpathlineto{\pgfqpoint{0.814345in}{0.404191in}}%
\pgfpathlineto{\pgfqpoint{0.812992in}{0.402013in}}%
\pgfpathlineto{\pgfqpoint{0.811166in}{0.403256in}}%
\pgfpathlineto{\pgfqpoint{0.810441in}{0.400723in}}%
\pgfpathlineto{\pgfqpoint{0.803324in}{0.399430in}}%
\pgfpathlineto{\pgfqpoint{0.800398in}{0.399952in}}%
\pgfpathlineto{\pgfqpoint{0.797438in}{0.402092in}}%
\pgfpathlineto{\pgfqpoint{0.794014in}{0.402817in}}%
\pgfpathlineto{\pgfqpoint{0.793204in}{0.401308in}}%
\pgfpathlineto{\pgfqpoint{0.788959in}{0.401438in}}%
\pgfpathlineto{\pgfqpoint{0.786334in}{0.400255in}}%
\pgfpathlineto{\pgfqpoint{0.785696in}{0.403237in}}%
\pgfpathlineto{\pgfqpoint{0.783361in}{0.401415in}}%
\pgfpathlineto{\pgfqpoint{0.780693in}{0.402239in}}%
\pgfpathlineto{\pgfqpoint{0.779576in}{0.400441in}}%
\pgfpathlineto{\pgfqpoint{0.777277in}{0.402089in}}%
\pgfpathlineto{\pgfqpoint{0.781062in}{0.403097in}}%
\pgfpathlineto{\pgfqpoint{0.784195in}{0.405252in}}%
\pgfpathlineto{\pgfqpoint{0.786031in}{0.403802in}}%
\pgfpathlineto{\pgfqpoint{0.789736in}{0.403394in}}%
\pgfpathlineto{\pgfqpoint{0.794932in}{0.410045in}}%
\pgfpathlineto{\pgfqpoint{0.797005in}{0.409992in}}%
\pgfpathlineto{\pgfqpoint{0.798949in}{0.408471in}}%
\pgfpathlineto{\pgfqpoint{0.800501in}{0.409249in}}%
\pgfpathlineto{\pgfqpoint{0.804050in}{0.407001in}}%
\pgfpathlineto{\pgfqpoint{0.808196in}{0.405860in}}%
\pgfpathlineto{\pgfqpoint{0.809601in}{0.406866in}}%
\pgfpathlineto{\pgfqpoint{0.811859in}{0.406633in}}%
\pgfpathlineto{\pgfqpoint{0.815079in}{0.410909in}}%
\pgfpathlineto{\pgfqpoint{0.820929in}{0.406477in}}%
\pgfpathlineto{\pgfqpoint{0.822419in}{0.408486in}}%
\pgfpathlineto{\pgfqpoint{0.826771in}{0.405229in}}%
\pgfpathlineto{\pgfqpoint{0.828822in}{0.408037in}}%
\pgfpathlineto{\pgfqpoint{0.831242in}{0.406233in}}%
\pgfpathclose%
\pgfusepath{fill}%
\end{pgfscope}%
\begin{pgfscope}%
\pgfpathrectangle{\pgfqpoint{0.100000in}{0.100000in}}{\pgfqpoint{3.608454in}{2.310000in}}%
\pgfusepath{clip}%
\pgfsetbuttcap%
\pgfsetmiterjoin%
\definecolor{currentfill}{rgb}{0.000000,0.454902,0.772549}%
\pgfsetfillcolor{currentfill}%
\pgfsetlinewidth{0.000000pt}%
\definecolor{currentstroke}{rgb}{0.000000,0.000000,0.000000}%
\pgfsetstrokecolor{currentstroke}%
\pgfsetstrokeopacity{0.000000}%
\pgfsetdash{}{0pt}%
\pgfpathmoveto{\pgfqpoint{1.204580in}{1.868127in}}%
\pgfpathlineto{\pgfqpoint{1.202252in}{1.853038in}}%
\pgfpathlineto{\pgfqpoint{1.198018in}{1.855881in}}%
\pgfpathlineto{\pgfqpoint{1.191695in}{1.873351in}}%
\pgfpathlineto{\pgfqpoint{1.187877in}{1.878039in}}%
\pgfpathlineto{\pgfqpoint{1.181983in}{1.875682in}}%
\pgfpathlineto{\pgfqpoint{1.164838in}{1.878825in}}%
\pgfpathlineto{\pgfqpoint{1.165886in}{1.884470in}}%
\pgfpathlineto{\pgfqpoint{1.150185in}{1.887641in}}%
\pgfpathlineto{\pgfqpoint{1.145219in}{1.892062in}}%
\pgfpathlineto{\pgfqpoint{1.146489in}{1.905873in}}%
\pgfpathlineto{\pgfqpoint{1.138589in}{1.907407in}}%
\pgfpathlineto{\pgfqpoint{1.139942in}{1.914263in}}%
\pgfpathlineto{\pgfqpoint{1.133161in}{1.915579in}}%
\pgfpathlineto{\pgfqpoint{1.137082in}{1.935988in}}%
\pgfpathlineto{\pgfqpoint{1.135886in}{1.943764in}}%
\pgfpathlineto{\pgfqpoint{1.131774in}{1.943515in}}%
\pgfpathlineto{\pgfqpoint{1.126427in}{1.948422in}}%
\pgfpathlineto{\pgfqpoint{1.128763in}{1.959930in}}%
\pgfpathlineto{\pgfqpoint{1.137125in}{1.963079in}}%
\pgfpathlineto{\pgfqpoint{1.140024in}{1.967674in}}%
\pgfpathlineto{\pgfqpoint{1.157736in}{1.964285in}}%
\pgfpathlineto{\pgfqpoint{1.165857in}{1.971182in}}%
\pgfpathlineto{\pgfqpoint{1.173618in}{1.970061in}}%
\pgfpathlineto{\pgfqpoint{1.176413in}{1.966592in}}%
\pgfpathlineto{\pgfqpoint{1.185357in}{1.963304in}}%
\pgfpathlineto{\pgfqpoint{1.192609in}{1.969852in}}%
\pgfpathlineto{\pgfqpoint{1.199884in}{1.972842in}}%
\pgfpathlineto{\pgfqpoint{1.203995in}{1.980934in}}%
\pgfpathlineto{\pgfqpoint{1.208958in}{1.985065in}}%
\pgfpathlineto{\pgfqpoint{1.210052in}{1.990143in}}%
\pgfpathlineto{\pgfqpoint{1.225529in}{1.987742in}}%
\pgfpathlineto{\pgfqpoint{1.261011in}{1.981452in}}%
\pgfpathlineto{\pgfqpoint{1.268315in}{1.982523in}}%
\pgfpathlineto{\pgfqpoint{1.302428in}{1.976762in}}%
\pgfpathlineto{\pgfqpoint{1.301322in}{1.969905in}}%
\pgfpathlineto{\pgfqpoint{1.302894in}{1.962642in}}%
\pgfpathlineto{\pgfqpoint{1.313126in}{1.960922in}}%
\pgfpathlineto{\pgfqpoint{1.311739in}{1.954347in}}%
\pgfpathlineto{\pgfqpoint{1.307165in}{1.954995in}}%
\pgfpathlineto{\pgfqpoint{1.301786in}{1.941660in}}%
\pgfpathlineto{\pgfqpoint{1.299700in}{1.927978in}}%
\pgfpathlineto{\pgfqpoint{1.292921in}{1.929061in}}%
\pgfpathlineto{\pgfqpoint{1.278245in}{1.924515in}}%
\pgfpathlineto{\pgfqpoint{1.275970in}{1.910859in}}%
\pgfpathlineto{\pgfqpoint{1.268744in}{1.912078in}}%
\pgfpathlineto{\pgfqpoint{1.266400in}{1.898228in}}%
\pgfpathlineto{\pgfqpoint{1.281086in}{1.895779in}}%
\pgfpathlineto{\pgfqpoint{1.278870in}{1.882470in}}%
\pgfpathlineto{\pgfqpoint{1.228507in}{1.890208in}}%
\pgfpathlineto{\pgfqpoint{1.209266in}{1.894331in}}%
\pgfpathlineto{\pgfqpoint{1.204580in}{1.868127in}}%
\pgfpathclose%
\pgfusepath{fill}%
\end{pgfscope}%
\begin{pgfscope}%
\pgfpathrectangle{\pgfqpoint{0.100000in}{0.100000in}}{\pgfqpoint{3.608454in}{2.310000in}}%
\pgfusepath{clip}%
\pgfsetbuttcap%
\pgfsetmiterjoin%
\definecolor{currentfill}{rgb}{0.000000,0.509804,0.745098}%
\pgfsetfillcolor{currentfill}%
\pgfsetlinewidth{0.000000pt}%
\definecolor{currentstroke}{rgb}{0.000000,0.000000,0.000000}%
\pgfsetstrokecolor{currentstroke}%
\pgfsetstrokeopacity{0.000000}%
\pgfsetdash{}{0pt}%
\pgfpathmoveto{\pgfqpoint{1.908588in}{1.422459in}}%
\pgfpathlineto{\pgfqpoint{1.894644in}{1.422961in}}%
\pgfpathlineto{\pgfqpoint{1.881222in}{1.423490in}}%
\pgfpathlineto{\pgfqpoint{1.883772in}{1.478524in}}%
\pgfpathlineto{\pgfqpoint{1.937428in}{1.476563in}}%
\pgfpathlineto{\pgfqpoint{1.936670in}{1.449037in}}%
\pgfpathlineto{\pgfqpoint{1.909610in}{1.449945in}}%
\pgfpathlineto{\pgfqpoint{1.908588in}{1.422459in}}%
\pgfpathclose%
\pgfusepath{fill}%
\end{pgfscope}%
\begin{pgfscope}%
\pgfpathrectangle{\pgfqpoint{0.100000in}{0.100000in}}{\pgfqpoint{3.608454in}{2.310000in}}%
\pgfusepath{clip}%
\pgfsetbuttcap%
\pgfsetmiterjoin%
\definecolor{currentfill}{rgb}{0.000000,0.576471,0.711765}%
\pgfsetfillcolor{currentfill}%
\pgfsetlinewidth{0.000000pt}%
\definecolor{currentstroke}{rgb}{0.000000,0.000000,0.000000}%
\pgfsetstrokecolor{currentstroke}%
\pgfsetstrokeopacity{0.000000}%
\pgfsetdash{}{0pt}%
\pgfpathmoveto{\pgfqpoint{0.695319in}{1.898980in}}%
\pgfpathlineto{\pgfqpoint{0.668725in}{1.906595in}}%
\pgfpathlineto{\pgfqpoint{0.616244in}{1.922404in}}%
\pgfpathlineto{\pgfqpoint{0.579411in}{1.933417in}}%
\pgfpathlineto{\pgfqpoint{0.585151in}{1.943983in}}%
\pgfpathlineto{\pgfqpoint{0.586510in}{1.951287in}}%
\pgfpathlineto{\pgfqpoint{0.595256in}{1.957298in}}%
\pgfpathlineto{\pgfqpoint{0.603007in}{1.965607in}}%
\pgfpathlineto{\pgfqpoint{0.605729in}{1.977984in}}%
\pgfpathlineto{\pgfqpoint{0.605387in}{1.987819in}}%
\pgfpathlineto{\pgfqpoint{0.608267in}{1.997300in}}%
\pgfpathlineto{\pgfqpoint{0.613506in}{2.004223in}}%
\pgfpathlineto{\pgfqpoint{0.616042in}{2.014692in}}%
\pgfpathlineto{\pgfqpoint{0.620648in}{2.021281in}}%
\pgfpathlineto{\pgfqpoint{0.656754in}{2.010285in}}%
\pgfpathlineto{\pgfqpoint{0.695564in}{1.998853in}}%
\pgfpathlineto{\pgfqpoint{0.693514in}{1.994463in}}%
\pgfpathlineto{\pgfqpoint{0.686387in}{1.969750in}}%
\pgfpathlineto{\pgfqpoint{0.706253in}{1.964177in}}%
\pgfpathlineto{\pgfqpoint{0.711901in}{1.958542in}}%
\pgfpathlineto{\pgfqpoint{0.710132in}{1.952205in}}%
\pgfpathlineto{\pgfqpoint{0.723444in}{1.948508in}}%
\pgfpathlineto{\pgfqpoint{0.715935in}{1.921940in}}%
\pgfpathlineto{\pgfqpoint{0.709441in}{1.923753in}}%
\pgfpathlineto{\pgfqpoint{0.703763in}{1.903845in}}%
\pgfpathlineto{\pgfqpoint{0.697108in}{1.905730in}}%
\pgfpathlineto{\pgfqpoint{0.695319in}{1.898980in}}%
\pgfpathclose%
\pgfusepath{fill}%
\end{pgfscope}%
\begin{pgfscope}%
\pgfpathrectangle{\pgfqpoint{0.100000in}{0.100000in}}{\pgfqpoint{3.608454in}{2.310000in}}%
\pgfusepath{clip}%
\pgfsetbuttcap%
\pgfsetmiterjoin%
\definecolor{currentfill}{rgb}{0.000000,0.423529,0.788235}%
\pgfsetfillcolor{currentfill}%
\pgfsetlinewidth{0.000000pt}%
\definecolor{currentstroke}{rgb}{0.000000,0.000000,0.000000}%
\pgfsetstrokecolor{currentstroke}%
\pgfsetstrokeopacity{0.000000}%
\pgfsetdash{}{0pt}%
\pgfpathmoveto{\pgfqpoint{1.716323in}{1.398352in}}%
\pgfpathlineto{\pgfqpoint{1.677687in}{1.401208in}}%
\pgfpathlineto{\pgfqpoint{1.677720in}{1.401672in}}%
\pgfpathlineto{\pgfqpoint{1.682362in}{1.462771in}}%
\pgfpathlineto{\pgfqpoint{1.726124in}{1.459646in}}%
\pgfpathlineto{\pgfqpoint{1.724173in}{1.432221in}}%
\pgfpathlineto{\pgfqpoint{1.718916in}{1.432564in}}%
\pgfpathlineto{\pgfqpoint{1.716323in}{1.398352in}}%
\pgfpathclose%
\pgfusepath{fill}%
\end{pgfscope}%
\begin{pgfscope}%
\pgfpathrectangle{\pgfqpoint{0.100000in}{0.100000in}}{\pgfqpoint{3.608454in}{2.310000in}}%
\pgfusepath{clip}%
\pgfsetbuttcap%
\pgfsetmiterjoin%
\definecolor{currentfill}{rgb}{0.000000,0.815686,0.592157}%
\pgfsetfillcolor{currentfill}%
\pgfsetlinewidth{0.000000pt}%
\definecolor{currentstroke}{rgb}{0.000000,0.000000,0.000000}%
\pgfsetstrokecolor{currentstroke}%
\pgfsetstrokeopacity{0.000000}%
\pgfsetdash{}{0pt}%
\pgfpathmoveto{\pgfqpoint{0.908989in}{1.065097in}}%
\pgfpathlineto{\pgfqpoint{0.914070in}{1.089859in}}%
\pgfpathlineto{\pgfqpoint{0.933245in}{1.183825in}}%
\pgfpathlineto{\pgfqpoint{0.937616in}{1.205156in}}%
\pgfpathlineto{\pgfqpoint{0.941162in}{1.217180in}}%
\pgfpathlineto{\pgfqpoint{0.940117in}{1.223821in}}%
\pgfpathlineto{\pgfqpoint{0.943617in}{1.227923in}}%
\pgfpathlineto{\pgfqpoint{0.949023in}{1.226105in}}%
\pgfpathlineto{\pgfqpoint{0.951426in}{1.230336in}}%
\pgfpathlineto{\pgfqpoint{0.961661in}{1.233095in}}%
\pgfpathlineto{\pgfqpoint{0.978736in}{1.235692in}}%
\pgfpathlineto{\pgfqpoint{0.986930in}{1.238736in}}%
\pgfpathlineto{\pgfqpoint{0.991491in}{1.243095in}}%
\pgfpathlineto{\pgfqpoint{0.991172in}{1.251855in}}%
\pgfpathlineto{\pgfqpoint{0.994707in}{1.264046in}}%
\pgfpathlineto{\pgfqpoint{1.005762in}{1.288169in}}%
\pgfpathlineto{\pgfqpoint{0.983512in}{1.292540in}}%
\pgfpathlineto{\pgfqpoint{0.991804in}{1.334673in}}%
\pgfpathlineto{\pgfqpoint{1.005205in}{1.331999in}}%
\pgfpathlineto{\pgfqpoint{1.024916in}{1.327610in}}%
\pgfpathlineto{\pgfqpoint{1.104844in}{1.313028in}}%
\pgfpathlineto{\pgfqpoint{1.130720in}{1.308550in}}%
\pgfpathlineto{\pgfqpoint{1.122123in}{1.294263in}}%
\pgfpathlineto{\pgfqpoint{1.113627in}{1.291412in}}%
\pgfpathlineto{\pgfqpoint{1.109948in}{1.282813in}}%
\pgfpathlineto{\pgfqpoint{1.099188in}{1.278642in}}%
\pgfpathlineto{\pgfqpoint{1.091741in}{1.280236in}}%
\pgfpathlineto{\pgfqpoint{1.087638in}{1.275829in}}%
\pgfpathlineto{\pgfqpoint{1.078577in}{1.274542in}}%
\pgfpathlineto{\pgfqpoint{1.116943in}{1.267777in}}%
\pgfpathlineto{\pgfqpoint{1.079528in}{1.054057in}}%
\pgfpathlineto{\pgfqpoint{1.073305in}{1.055337in}}%
\pgfpathlineto{\pgfqpoint{1.063896in}{1.064728in}}%
\pgfpathlineto{\pgfqpoint{1.055533in}{1.068101in}}%
\pgfpathlineto{\pgfqpoint{1.049066in}{1.073939in}}%
\pgfpathlineto{\pgfqpoint{1.039948in}{1.073968in}}%
\pgfpathlineto{\pgfqpoint{1.031414in}{1.079449in}}%
\pgfpathlineto{\pgfqpoint{1.022683in}{1.074223in}}%
\pgfpathlineto{\pgfqpoint{1.015879in}{1.056517in}}%
\pgfpathlineto{\pgfqpoint{1.031005in}{1.053711in}}%
\pgfpathlineto{\pgfqpoint{1.027746in}{1.042105in}}%
\pgfpathlineto{\pgfqpoint{1.012835in}{1.044864in}}%
\pgfpathlineto{\pgfqpoint{0.985165in}{1.053983in}}%
\pgfpathlineto{\pgfqpoint{0.975479in}{1.042443in}}%
\pgfpathlineto{\pgfqpoint{0.947034in}{1.057499in}}%
\pgfpathlineto{\pgfqpoint{0.908989in}{1.065097in}}%
\pgfpathclose%
\pgfusepath{fill}%
\end{pgfscope}%
\begin{pgfscope}%
\pgfpathrectangle{\pgfqpoint{0.100000in}{0.100000in}}{\pgfqpoint{3.608454in}{2.310000in}}%
\pgfusepath{clip}%
\pgfsetbuttcap%
\pgfsetmiterjoin%
\definecolor{currentfill}{rgb}{0.000000,0.219608,0.890196}%
\pgfsetfillcolor{currentfill}%
\pgfsetlinewidth{0.000000pt}%
\definecolor{currentstroke}{rgb}{0.000000,0.000000,0.000000}%
\pgfsetstrokecolor{currentstroke}%
\pgfsetstrokeopacity{0.000000}%
\pgfsetdash{}{0pt}%
\pgfpathmoveto{\pgfqpoint{1.577074in}{1.409804in}}%
\pgfpathlineto{\pgfqpoint{1.571014in}{1.354845in}}%
\pgfpathlineto{\pgfqpoint{1.550467in}{1.357020in}}%
\pgfpathlineto{\pgfqpoint{1.552744in}{1.377402in}}%
\pgfpathlineto{\pgfqpoint{1.474789in}{1.385857in}}%
\pgfpathlineto{\pgfqpoint{1.470992in}{1.351885in}}%
\pgfpathlineto{\pgfqpoint{1.431712in}{1.356218in}}%
\pgfpathlineto{\pgfqpoint{1.436653in}{1.364650in}}%
\pgfpathlineto{\pgfqpoint{1.436137in}{1.369904in}}%
\pgfpathlineto{\pgfqpoint{1.430521in}{1.376213in}}%
\pgfpathlineto{\pgfqpoint{1.421188in}{1.383353in}}%
\pgfpathlineto{\pgfqpoint{1.421683in}{1.386315in}}%
\pgfpathlineto{\pgfqpoint{1.423585in}{1.392666in}}%
\pgfpathlineto{\pgfqpoint{1.424568in}{1.406130in}}%
\pgfpathlineto{\pgfqpoint{1.439656in}{1.412644in}}%
\pgfpathlineto{\pgfqpoint{1.451897in}{1.428337in}}%
\pgfpathlineto{\pgfqpoint{1.444005in}{1.434780in}}%
\pgfpathlineto{\pgfqpoint{1.450957in}{1.440841in}}%
\pgfpathlineto{\pgfqpoint{1.459571in}{1.445102in}}%
\pgfpathlineto{\pgfqpoint{1.461324in}{1.456249in}}%
\pgfpathlineto{\pgfqpoint{1.464403in}{1.459400in}}%
\pgfpathlineto{\pgfqpoint{1.465579in}{1.476906in}}%
\pgfpathlineto{\pgfqpoint{1.501446in}{1.472932in}}%
\pgfpathlineto{\pgfqpoint{1.499141in}{1.452348in}}%
\pgfpathlineto{\pgfqpoint{1.553648in}{1.446628in}}%
\pgfpathlineto{\pgfqpoint{1.580447in}{1.444028in}}%
\pgfpathlineto{\pgfqpoint{1.577074in}{1.409804in}}%
\pgfpathclose%
\pgfusepath{fill}%
\end{pgfscope}%
\begin{pgfscope}%
\pgfpathrectangle{\pgfqpoint{0.100000in}{0.100000in}}{\pgfqpoint{3.608454in}{2.310000in}}%
\pgfusepath{clip}%
\pgfsetbuttcap%
\pgfsetmiterjoin%
\definecolor{currentfill}{rgb}{0.000000,0.701961,0.649020}%
\pgfsetfillcolor{currentfill}%
\pgfsetlinewidth{0.000000pt}%
\definecolor{currentstroke}{rgb}{0.000000,0.000000,0.000000}%
\pgfsetstrokecolor{currentstroke}%
\pgfsetstrokeopacity{0.000000}%
\pgfsetdash{}{0pt}%
\pgfpathmoveto{\pgfqpoint{1.797569in}{0.560061in}}%
\pgfpathlineto{\pgfqpoint{1.748728in}{0.562219in}}%
\pgfpathlineto{\pgfqpoint{1.751098in}{0.605189in}}%
\pgfpathlineto{\pgfqpoint{1.710253in}{0.607672in}}%
\pgfpathlineto{\pgfqpoint{1.713553in}{0.660732in}}%
\pgfpathlineto{\pgfqpoint{1.695622in}{0.661844in}}%
\pgfpathlineto{\pgfqpoint{1.697817in}{0.695195in}}%
\pgfpathlineto{\pgfqpoint{1.755704in}{0.692007in}}%
\pgfpathlineto{\pgfqpoint{1.753800in}{0.658495in}}%
\pgfpathlineto{\pgfqpoint{1.778772in}{0.657165in}}%
\pgfpathlineto{\pgfqpoint{1.777624in}{0.639888in}}%
\pgfpathlineto{\pgfqpoint{1.782309in}{0.639635in}}%
\pgfpathlineto{\pgfqpoint{1.781527in}{0.626427in}}%
\pgfpathlineto{\pgfqpoint{1.787654in}{0.626021in}}%
\pgfpathlineto{\pgfqpoint{1.786496in}{0.603609in}}%
\pgfpathlineto{\pgfqpoint{1.799798in}{0.602976in}}%
\pgfpathlineto{\pgfqpoint{1.797569in}{0.560061in}}%
\pgfpathclose%
\pgfusepath{fill}%
\end{pgfscope}%
\begin{pgfscope}%
\pgfpathrectangle{\pgfqpoint{0.100000in}{0.100000in}}{\pgfqpoint{3.608454in}{2.310000in}}%
\pgfusepath{clip}%
\pgfsetbuttcap%
\pgfsetmiterjoin%
\definecolor{currentfill}{rgb}{0.000000,0.674510,0.662745}%
\pgfsetfillcolor{currentfill}%
\pgfsetlinewidth{0.000000pt}%
\definecolor{currentstroke}{rgb}{0.000000,0.000000,0.000000}%
\pgfsetstrokecolor{currentstroke}%
\pgfsetstrokeopacity{0.000000}%
\pgfsetdash{}{0pt}%
\pgfpathmoveto{\pgfqpoint{1.167197in}{0.891666in}}%
\pgfpathlineto{\pgfqpoint{1.176800in}{0.953138in}}%
\pgfpathlineto{\pgfqpoint{1.245517in}{0.942241in}}%
\pgfpathlineto{\pgfqpoint{1.253754in}{0.941144in}}%
\pgfpathlineto{\pgfqpoint{1.255093in}{0.930635in}}%
\pgfpathlineto{\pgfqpoint{1.253634in}{0.922827in}}%
\pgfpathlineto{\pgfqpoint{1.257767in}{0.919573in}}%
\pgfpathlineto{\pgfqpoint{1.256598in}{0.912844in}}%
\pgfpathlineto{\pgfqpoint{1.259177in}{0.906354in}}%
\pgfpathlineto{\pgfqpoint{1.256509in}{0.894758in}}%
\pgfpathlineto{\pgfqpoint{1.264784in}{0.891644in}}%
\pgfpathlineto{\pgfqpoint{1.285282in}{0.888829in}}%
\pgfpathlineto{\pgfqpoint{1.276635in}{0.824008in}}%
\pgfpathlineto{\pgfqpoint{1.215509in}{0.832612in}}%
\pgfpathlineto{\pgfqpoint{1.210301in}{0.797108in}}%
\pgfpathlineto{\pgfqpoint{1.153619in}{0.805601in}}%
\pgfpathlineto{\pgfqpoint{1.167197in}{0.891666in}}%
\pgfpathclose%
\pgfusepath{fill}%
\end{pgfscope}%
\begin{pgfscope}%
\pgfpathrectangle{\pgfqpoint{0.100000in}{0.100000in}}{\pgfqpoint{3.608454in}{2.310000in}}%
\pgfusepath{clip}%
\pgfsetbuttcap%
\pgfsetmiterjoin%
\definecolor{currentfill}{rgb}{0.000000,0.478431,0.760784}%
\pgfsetfillcolor{currentfill}%
\pgfsetlinewidth{0.000000pt}%
\definecolor{currentstroke}{rgb}{0.000000,0.000000,0.000000}%
\pgfsetstrokecolor{currentstroke}%
\pgfsetstrokeopacity{0.000000}%
\pgfsetdash{}{0pt}%
\pgfpathmoveto{\pgfqpoint{1.983409in}{1.890641in}}%
\pgfpathlineto{\pgfqpoint{1.941869in}{1.891889in}}%
\pgfpathlineto{\pgfqpoint{1.885060in}{1.894397in}}%
\pgfpathlineto{\pgfqpoint{1.886398in}{1.921692in}}%
\pgfpathlineto{\pgfqpoint{1.939707in}{1.919594in}}%
\pgfpathlineto{\pgfqpoint{1.940671in}{1.947331in}}%
\pgfpathlineto{\pgfqpoint{1.959926in}{1.946668in}}%
\pgfpathlineto{\pgfqpoint{1.981936in}{1.946008in}}%
\pgfpathlineto{\pgfqpoint{1.981211in}{1.918353in}}%
\pgfpathlineto{\pgfqpoint{1.981775in}{1.895615in}}%
\pgfpathlineto{\pgfqpoint{1.983409in}{1.890641in}}%
\pgfpathclose%
\pgfusepath{fill}%
\end{pgfscope}%
\begin{pgfscope}%
\pgfpathrectangle{\pgfqpoint{0.100000in}{0.100000in}}{\pgfqpoint{3.608454in}{2.310000in}}%
\pgfusepath{clip}%
\pgfsetbuttcap%
\pgfsetmiterjoin%
\definecolor{currentfill}{rgb}{0.000000,0.556863,0.721569}%
\pgfsetfillcolor{currentfill}%
\pgfsetlinewidth{0.000000pt}%
\definecolor{currentstroke}{rgb}{0.000000,0.000000,0.000000}%
\pgfsetstrokecolor{currentstroke}%
\pgfsetstrokeopacity{0.000000}%
\pgfsetdash{}{0pt}%
\pgfpathmoveto{\pgfqpoint{3.371539in}{1.933332in}}%
\pgfpathlineto{\pgfqpoint{3.373495in}{1.940927in}}%
\pgfpathlineto{\pgfqpoint{3.366415in}{1.942507in}}%
\pgfpathlineto{\pgfqpoint{3.368238in}{1.950258in}}%
\pgfpathlineto{\pgfqpoint{3.358298in}{1.952554in}}%
\pgfpathlineto{\pgfqpoint{3.359777in}{1.957182in}}%
\pgfpathlineto{\pgfqpoint{3.355196in}{1.970349in}}%
\pgfpathlineto{\pgfqpoint{3.368542in}{1.973373in}}%
\pgfpathlineto{\pgfqpoint{3.412317in}{1.985521in}}%
\pgfpathlineto{\pgfqpoint{3.415090in}{2.005571in}}%
\pgfpathlineto{\pgfqpoint{3.427643in}{2.008700in}}%
\pgfpathlineto{\pgfqpoint{3.428979in}{2.013989in}}%
\pgfpathlineto{\pgfqpoint{3.442334in}{1.974192in}}%
\pgfpathlineto{\pgfqpoint{3.454146in}{1.936802in}}%
\pgfpathlineto{\pgfqpoint{3.453336in}{1.932723in}}%
\pgfpathlineto{\pgfqpoint{3.441678in}{1.930119in}}%
\pgfpathlineto{\pgfqpoint{3.440922in}{1.921126in}}%
\pgfpathlineto{\pgfqpoint{3.433355in}{1.925260in}}%
\pgfpathlineto{\pgfqpoint{3.431085in}{1.927713in}}%
\pgfpathlineto{\pgfqpoint{3.423697in}{1.926204in}}%
\pgfpathlineto{\pgfqpoint{3.422059in}{1.929884in}}%
\pgfpathlineto{\pgfqpoint{3.415486in}{1.930416in}}%
\pgfpathlineto{\pgfqpoint{3.410255in}{1.935145in}}%
\pgfpathlineto{\pgfqpoint{3.405777in}{1.928235in}}%
\pgfpathlineto{\pgfqpoint{3.399587in}{1.926596in}}%
\pgfpathlineto{\pgfqpoint{3.396256in}{1.920610in}}%
\pgfpathlineto{\pgfqpoint{3.398006in}{1.914180in}}%
\pgfpathlineto{\pgfqpoint{3.384973in}{1.910295in}}%
\pgfpathlineto{\pgfqpoint{3.381069in}{1.911013in}}%
\pgfpathlineto{\pgfqpoint{3.384688in}{1.930010in}}%
\pgfpathlineto{\pgfqpoint{3.371539in}{1.933332in}}%
\pgfpathclose%
\pgfusepath{fill}%
\end{pgfscope}%
\begin{pgfscope}%
\pgfpathrectangle{\pgfqpoint{0.100000in}{0.100000in}}{\pgfqpoint{3.608454in}{2.310000in}}%
\pgfusepath{clip}%
\pgfsetbuttcap%
\pgfsetmiterjoin%
\definecolor{currentfill}{rgb}{0.000000,0.325490,0.837255}%
\pgfsetfillcolor{currentfill}%
\pgfsetlinewidth{0.000000pt}%
\definecolor{currentstroke}{rgb}{0.000000,0.000000,0.000000}%
\pgfsetstrokecolor{currentstroke}%
\pgfsetstrokeopacity{0.000000}%
\pgfsetdash{}{0pt}%
\pgfpathmoveto{\pgfqpoint{1.912766in}{1.573654in}}%
\pgfpathlineto{\pgfqpoint{1.914002in}{1.614954in}}%
\pgfpathlineto{\pgfqpoint{1.968471in}{1.613417in}}%
\pgfpathlineto{\pgfqpoint{1.968280in}{1.606524in}}%
\pgfpathlineto{\pgfqpoint{1.967706in}{1.585865in}}%
\pgfpathlineto{\pgfqpoint{1.988196in}{1.585344in}}%
\pgfpathlineto{\pgfqpoint{1.987553in}{1.557852in}}%
\pgfpathlineto{\pgfqpoint{1.939698in}{1.559064in}}%
\pgfpathlineto{\pgfqpoint{1.939982in}{1.572847in}}%
\pgfpathlineto{\pgfqpoint{1.912766in}{1.573654in}}%
\pgfpathclose%
\pgfusepath{fill}%
\end{pgfscope}%
\begin{pgfscope}%
\pgfpathrectangle{\pgfqpoint{0.100000in}{0.100000in}}{\pgfqpoint{3.608454in}{2.310000in}}%
\pgfusepath{clip}%
\pgfsetbuttcap%
\pgfsetmiterjoin%
\definecolor{currentfill}{rgb}{0.000000,0.450980,0.774510}%
\pgfsetfillcolor{currentfill}%
\pgfsetlinewidth{0.000000pt}%
\definecolor{currentstroke}{rgb}{0.000000,0.000000,0.000000}%
\pgfsetstrokecolor{currentstroke}%
\pgfsetstrokeopacity{0.000000}%
\pgfsetdash{}{0pt}%
\pgfpathmoveto{\pgfqpoint{1.580447in}{1.444028in}}%
\pgfpathlineto{\pgfqpoint{1.553648in}{1.446628in}}%
\pgfpathlineto{\pgfqpoint{1.499141in}{1.452348in}}%
\pgfpathlineto{\pgfqpoint{1.501446in}{1.472932in}}%
\pgfpathlineto{\pgfqpoint{1.465579in}{1.476906in}}%
\pgfpathlineto{\pgfqpoint{1.459655in}{1.485514in}}%
\pgfpathlineto{\pgfqpoint{1.455629in}{1.496066in}}%
\pgfpathlineto{\pgfqpoint{1.446700in}{1.523501in}}%
\pgfpathlineto{\pgfqpoint{1.440255in}{1.533604in}}%
\pgfpathlineto{\pgfqpoint{1.440561in}{1.538635in}}%
\pgfpathlineto{\pgfqpoint{1.494704in}{1.532257in}}%
\pgfpathlineto{\pgfqpoint{1.533391in}{1.528263in}}%
\pgfpathlineto{\pgfqpoint{1.595804in}{1.522015in}}%
\pgfpathlineto{\pgfqpoint{1.591803in}{1.484401in}}%
\pgfpathlineto{\pgfqpoint{1.598652in}{1.483744in}}%
\pgfpathlineto{\pgfqpoint{1.594536in}{1.442692in}}%
\pgfpathlineto{\pgfqpoint{1.580447in}{1.444028in}}%
\pgfpathclose%
\pgfusepath{fill}%
\end{pgfscope}%
\begin{pgfscope}%
\pgfpathrectangle{\pgfqpoint{0.100000in}{0.100000in}}{\pgfqpoint{3.608454in}{2.310000in}}%
\pgfusepath{clip}%
\pgfsetbuttcap%
\pgfsetmiterjoin%
\definecolor{currentfill}{rgb}{0.000000,0.741176,0.629412}%
\pgfsetfillcolor{currentfill}%
\pgfsetlinewidth{0.000000pt}%
\definecolor{currentstroke}{rgb}{0.000000,0.000000,0.000000}%
\pgfsetstrokecolor{currentstroke}%
\pgfsetstrokeopacity{0.000000}%
\pgfsetdash{}{0pt}%
\pgfpathmoveto{\pgfqpoint{0.869214in}{2.048910in}}%
\pgfpathlineto{\pgfqpoint{0.875412in}{2.075250in}}%
\pgfpathlineto{\pgfqpoint{0.872020in}{2.076043in}}%
\pgfpathlineto{\pgfqpoint{0.871344in}{2.087217in}}%
\pgfpathlineto{\pgfqpoint{0.864637in}{2.088901in}}%
\pgfpathlineto{\pgfqpoint{0.868296in}{2.100626in}}%
\pgfpathlineto{\pgfqpoint{0.878570in}{2.104452in}}%
\pgfpathlineto{\pgfqpoint{0.882536in}{2.100932in}}%
\pgfpathlineto{\pgfqpoint{0.890708in}{2.100407in}}%
\pgfpathlineto{\pgfqpoint{0.892918in}{2.093767in}}%
\pgfpathlineto{\pgfqpoint{0.898511in}{2.086388in}}%
\pgfpathlineto{\pgfqpoint{0.899012in}{2.076012in}}%
\pgfpathlineto{\pgfqpoint{0.907603in}{2.074322in}}%
\pgfpathlineto{\pgfqpoint{0.909838in}{2.083326in}}%
\pgfpathlineto{\pgfqpoint{0.926998in}{2.079195in}}%
\pgfpathlineto{\pgfqpoint{0.933769in}{2.084665in}}%
\pgfpathlineto{\pgfqpoint{0.949294in}{2.080795in}}%
\pgfpathlineto{\pgfqpoint{0.954967in}{2.104591in}}%
\pgfpathlineto{\pgfqpoint{0.993104in}{2.095521in}}%
\pgfpathlineto{\pgfqpoint{1.027701in}{2.087646in}}%
\pgfpathlineto{\pgfqpoint{1.027942in}{2.078990in}}%
\pgfpathlineto{\pgfqpoint{1.041857in}{2.066086in}}%
\pgfpathlineto{\pgfqpoint{1.041035in}{2.060695in}}%
\pgfpathlineto{\pgfqpoint{1.055012in}{2.060043in}}%
\pgfpathlineto{\pgfqpoint{1.056865in}{2.058071in}}%
\pgfpathlineto{\pgfqpoint{1.049410in}{2.040224in}}%
\pgfpathlineto{\pgfqpoint{1.042825in}{2.030433in}}%
\pgfpathlineto{\pgfqpoint{1.042048in}{2.022019in}}%
\pgfpathlineto{\pgfqpoint{1.037550in}{2.013584in}}%
\pgfpathlineto{\pgfqpoint{1.040540in}{1.999733in}}%
\pgfpathlineto{\pgfqpoint{1.033186in}{1.997675in}}%
\pgfpathlineto{\pgfqpoint{1.028886in}{1.992775in}}%
\pgfpathlineto{\pgfqpoint{1.030719in}{1.983835in}}%
\pgfpathlineto{\pgfqpoint{1.025284in}{1.975956in}}%
\pgfpathlineto{\pgfqpoint{1.018436in}{1.970411in}}%
\pgfpathlineto{\pgfqpoint{1.011182in}{1.972459in}}%
\pgfpathlineto{\pgfqpoint{1.018689in}{1.956057in}}%
\pgfpathlineto{\pgfqpoint{1.012082in}{1.949489in}}%
\pgfpathlineto{\pgfqpoint{0.942731in}{1.965121in}}%
\pgfpathlineto{\pgfqpoint{0.931228in}{1.960671in}}%
\pgfpathlineto{\pgfqpoint{0.923969in}{1.964164in}}%
\pgfpathlineto{\pgfqpoint{0.923373in}{1.975654in}}%
\pgfpathlineto{\pgfqpoint{0.904646in}{1.980156in}}%
\pgfpathlineto{\pgfqpoint{0.916109in}{1.996717in}}%
\pgfpathlineto{\pgfqpoint{0.923281in}{2.004122in}}%
\pgfpathlineto{\pgfqpoint{0.921587in}{2.014227in}}%
\pgfpathlineto{\pgfqpoint{0.914485in}{2.022985in}}%
\pgfpathlineto{\pgfqpoint{0.909977in}{2.024514in}}%
\pgfpathlineto{\pgfqpoint{0.906027in}{2.039422in}}%
\pgfpathlineto{\pgfqpoint{0.869214in}{2.048910in}}%
\pgfpathclose%
\pgfusepath{fill}%
\end{pgfscope}%
\begin{pgfscope}%
\pgfpathrectangle{\pgfqpoint{0.100000in}{0.100000in}}{\pgfqpoint{3.608454in}{2.310000in}}%
\pgfusepath{clip}%
\pgfsetbuttcap%
\pgfsetmiterjoin%
\definecolor{currentfill}{rgb}{0.000000,0.698039,0.650980}%
\pgfsetfillcolor{currentfill}%
\pgfsetlinewidth{0.000000pt}%
\definecolor{currentstroke}{rgb}{0.000000,0.000000,0.000000}%
\pgfsetstrokecolor{currentstroke}%
\pgfsetstrokeopacity{0.000000}%
\pgfsetdash{}{0pt}%
\pgfpathmoveto{\pgfqpoint{2.971461in}{0.696012in}}%
\pgfpathlineto{\pgfqpoint{2.999393in}{0.700496in}}%
\pgfpathlineto{\pgfqpoint{3.004229in}{0.666876in}}%
\pgfpathlineto{\pgfqpoint{3.006579in}{0.647225in}}%
\pgfpathlineto{\pgfqpoint{2.996330in}{0.642369in}}%
\pgfpathlineto{\pgfqpoint{2.995671in}{0.646679in}}%
\pgfpathlineto{\pgfqpoint{2.982209in}{0.644864in}}%
\pgfpathlineto{\pgfqpoint{2.985432in}{0.623511in}}%
\pgfpathlineto{\pgfqpoint{2.976237in}{0.622163in}}%
\pgfpathlineto{\pgfqpoint{2.978125in}{0.608670in}}%
\pgfpathlineto{\pgfqpoint{2.973159in}{0.605106in}}%
\pgfpathlineto{\pgfqpoint{2.965965in}{0.605834in}}%
\pgfpathlineto{\pgfqpoint{2.963258in}{0.602568in}}%
\pgfpathlineto{\pgfqpoint{2.962394in}{0.607479in}}%
\pgfpathlineto{\pgfqpoint{2.957118in}{0.615111in}}%
\pgfpathlineto{\pgfqpoint{2.942926in}{0.613795in}}%
\pgfpathlineto{\pgfqpoint{2.937858in}{0.621000in}}%
\pgfpathlineto{\pgfqpoint{2.930039in}{0.626757in}}%
\pgfpathlineto{\pgfqpoint{2.926843in}{0.632121in}}%
\pgfpathlineto{\pgfqpoint{2.921124in}{0.633342in}}%
\pgfpathlineto{\pgfqpoint{2.913615in}{0.638189in}}%
\pgfpathlineto{\pgfqpoint{2.910647in}{0.649877in}}%
\pgfpathlineto{\pgfqpoint{2.913541in}{0.650279in}}%
\pgfpathlineto{\pgfqpoint{2.915694in}{0.663204in}}%
\pgfpathlineto{\pgfqpoint{2.944970in}{0.667447in}}%
\pgfpathlineto{\pgfqpoint{2.946140in}{0.674377in}}%
\pgfpathlineto{\pgfqpoint{2.951329in}{0.676818in}}%
\pgfpathlineto{\pgfqpoint{2.955684in}{0.671251in}}%
\pgfpathlineto{\pgfqpoint{2.962277in}{0.669805in}}%
\pgfpathlineto{\pgfqpoint{2.965236in}{0.672986in}}%
\pgfpathlineto{\pgfqpoint{2.965631in}{0.682186in}}%
\pgfpathlineto{\pgfqpoint{2.971461in}{0.696012in}}%
\pgfpathclose%
\pgfusepath{fill}%
\end{pgfscope}%
\begin{pgfscope}%
\pgfpathrectangle{\pgfqpoint{0.100000in}{0.100000in}}{\pgfqpoint{3.608454in}{2.310000in}}%
\pgfusepath{clip}%
\pgfsetbuttcap%
\pgfsetmiterjoin%
\definecolor{currentfill}{rgb}{0.000000,0.427451,0.786275}%
\pgfsetfillcolor{currentfill}%
\pgfsetlinewidth{0.000000pt}%
\definecolor{currentstroke}{rgb}{0.000000,0.000000,0.000000}%
\pgfsetstrokecolor{currentstroke}%
\pgfsetstrokeopacity{0.000000}%
\pgfsetdash{}{0pt}%
\pgfpathmoveto{\pgfqpoint{2.817692in}{1.389496in}}%
\pgfpathlineto{\pgfqpoint{2.825080in}{1.387318in}}%
\pgfpathlineto{\pgfqpoint{2.829317in}{1.356171in}}%
\pgfpathlineto{\pgfqpoint{2.812960in}{1.360831in}}%
\pgfpathlineto{\pgfqpoint{2.807405in}{1.358642in}}%
\pgfpathlineto{\pgfqpoint{2.804813in}{1.353816in}}%
\pgfpathlineto{\pgfqpoint{2.798224in}{1.355170in}}%
\pgfpathlineto{\pgfqpoint{2.790830in}{1.362794in}}%
\pgfpathlineto{\pgfqpoint{2.784131in}{1.363939in}}%
\pgfpathlineto{\pgfqpoint{2.777359in}{1.362093in}}%
\pgfpathlineto{\pgfqpoint{2.768060in}{1.364802in}}%
\pgfpathlineto{\pgfqpoint{2.772188in}{1.346181in}}%
\pgfpathlineto{\pgfqpoint{2.760419in}{1.343088in}}%
\pgfpathlineto{\pgfqpoint{2.750154in}{1.334910in}}%
\pgfpathlineto{\pgfqpoint{2.736943in}{1.344983in}}%
\pgfpathlineto{\pgfqpoint{2.735344in}{1.352880in}}%
\pgfpathlineto{\pgfqpoint{2.726746in}{1.347306in}}%
\pgfpathlineto{\pgfqpoint{2.720376in}{1.354526in}}%
\pgfpathlineto{\pgfqpoint{2.728448in}{1.358041in}}%
\pgfpathlineto{\pgfqpoint{2.733971in}{1.365444in}}%
\pgfpathlineto{\pgfqpoint{2.727957in}{1.367391in}}%
\pgfpathlineto{\pgfqpoint{2.729795in}{1.373635in}}%
\pgfpathlineto{\pgfqpoint{2.725607in}{1.378219in}}%
\pgfpathlineto{\pgfqpoint{2.729836in}{1.382914in}}%
\pgfpathlineto{\pgfqpoint{2.726448in}{1.415659in}}%
\pgfpathlineto{\pgfqpoint{2.726071in}{1.419220in}}%
\pgfpathlineto{\pgfqpoint{2.746401in}{1.421668in}}%
\pgfpathlineto{\pgfqpoint{2.757924in}{1.424641in}}%
\pgfpathlineto{\pgfqpoint{2.776762in}{1.425313in}}%
\pgfpathlineto{\pgfqpoint{2.777953in}{1.400497in}}%
\pgfpathlineto{\pgfqpoint{2.786603in}{1.400891in}}%
\pgfpathlineto{\pgfqpoint{2.788388in}{1.383110in}}%
\pgfpathlineto{\pgfqpoint{2.804293in}{1.384898in}}%
\pgfpathlineto{\pgfqpoint{2.817692in}{1.389496in}}%
\pgfpathclose%
\pgfusepath{fill}%
\end{pgfscope}%
\begin{pgfscope}%
\pgfpathrectangle{\pgfqpoint{0.100000in}{0.100000in}}{\pgfqpoint{3.608454in}{2.310000in}}%
\pgfusepath{clip}%
\pgfsetbuttcap%
\pgfsetmiterjoin%
\definecolor{currentfill}{rgb}{0.000000,0.611765,0.694118}%
\pgfsetfillcolor{currentfill}%
\pgfsetlinewidth{0.000000pt}%
\definecolor{currentstroke}{rgb}{0.000000,0.000000,0.000000}%
\pgfsetstrokecolor{currentstroke}%
\pgfsetstrokeopacity{0.000000}%
\pgfsetdash{}{0pt}%
\pgfpathmoveto{\pgfqpoint{2.906878in}{1.433220in}}%
\pgfpathlineto{\pgfqpoint{2.894009in}{1.432157in}}%
\pgfpathlineto{\pgfqpoint{2.893522in}{1.439324in}}%
\pgfpathlineto{\pgfqpoint{2.886708in}{1.438806in}}%
\pgfpathlineto{\pgfqpoint{2.872947in}{1.440094in}}%
\pgfpathlineto{\pgfqpoint{2.872185in}{1.451622in}}%
\pgfpathlineto{\pgfqpoint{2.869385in}{1.458446in}}%
\pgfpathlineto{\pgfqpoint{2.864747in}{1.458144in}}%
\pgfpathlineto{\pgfqpoint{2.864318in}{1.465519in}}%
\pgfpathlineto{\pgfqpoint{2.878244in}{1.466104in}}%
\pgfpathlineto{\pgfqpoint{2.879950in}{1.469292in}}%
\pgfpathlineto{\pgfqpoint{2.878220in}{1.486351in}}%
\pgfpathlineto{\pgfqpoint{2.909284in}{1.489578in}}%
\pgfpathlineto{\pgfqpoint{2.908656in}{1.495198in}}%
\pgfpathlineto{\pgfqpoint{2.928425in}{1.497499in}}%
\pgfpathlineto{\pgfqpoint{2.935644in}{1.495065in}}%
\pgfpathlineto{\pgfqpoint{2.938569in}{1.471342in}}%
\pgfpathlineto{\pgfqpoint{2.933975in}{1.470774in}}%
\pgfpathlineto{\pgfqpoint{2.935564in}{1.458054in}}%
\pgfpathlineto{\pgfqpoint{2.938939in}{1.450447in}}%
\pgfpathlineto{\pgfqpoint{2.927738in}{1.446958in}}%
\pgfpathlineto{\pgfqpoint{2.914676in}{1.445590in}}%
\pgfpathlineto{\pgfqpoint{2.914756in}{1.436835in}}%
\pgfpathlineto{\pgfqpoint{2.907685in}{1.436860in}}%
\pgfpathlineto{\pgfqpoint{2.906878in}{1.433220in}}%
\pgfpathclose%
\pgfusepath{fill}%
\end{pgfscope}%
\begin{pgfscope}%
\pgfpathrectangle{\pgfqpoint{0.100000in}{0.100000in}}{\pgfqpoint{3.608454in}{2.310000in}}%
\pgfusepath{clip}%
\pgfsetbuttcap%
\pgfsetmiterjoin%
\definecolor{currentfill}{rgb}{0.000000,0.458824,0.770588}%
\pgfsetfillcolor{currentfill}%
\pgfsetlinewidth{0.000000pt}%
\definecolor{currentstroke}{rgb}{0.000000,0.000000,0.000000}%
\pgfsetstrokecolor{currentstroke}%
\pgfsetstrokeopacity{0.000000}%
\pgfsetdash{}{0pt}%
\pgfpathmoveto{\pgfqpoint{2.025168in}{1.813601in}}%
\pgfpathlineto{\pgfqpoint{2.000983in}{1.813993in}}%
\pgfpathlineto{\pgfqpoint{2.001203in}{1.827773in}}%
\pgfpathlineto{\pgfqpoint{1.981983in}{1.828333in}}%
\pgfpathlineto{\pgfqpoint{1.982265in}{1.839933in}}%
\pgfpathlineto{\pgfqpoint{1.994509in}{1.839579in}}%
\pgfpathlineto{\pgfqpoint{1.995378in}{1.841947in}}%
\pgfpathlineto{\pgfqpoint{2.024583in}{1.841364in}}%
\pgfpathlineto{\pgfqpoint{2.021780in}{1.845315in}}%
\pgfpathlineto{\pgfqpoint{2.012969in}{1.848244in}}%
\pgfpathlineto{\pgfqpoint{2.003389in}{1.863983in}}%
\pgfpathlineto{\pgfqpoint{2.004519in}{1.867267in}}%
\pgfpathlineto{\pgfqpoint{2.015031in}{1.874970in}}%
\pgfpathlineto{\pgfqpoint{2.018973in}{1.880728in}}%
\pgfpathlineto{\pgfqpoint{2.020229in}{1.889881in}}%
\pgfpathlineto{\pgfqpoint{2.019610in}{1.896790in}}%
\pgfpathlineto{\pgfqpoint{2.036803in}{1.896501in}}%
\pgfpathlineto{\pgfqpoint{2.037372in}{1.889582in}}%
\pgfpathlineto{\pgfqpoint{2.037019in}{1.861775in}}%
\pgfpathlineto{\pgfqpoint{2.037488in}{1.848022in}}%
\pgfpathlineto{\pgfqpoint{2.044431in}{1.847876in}}%
\pgfpathlineto{\pgfqpoint{2.044201in}{1.834191in}}%
\pgfpathlineto{\pgfqpoint{2.041094in}{1.830482in}}%
\pgfpathlineto{\pgfqpoint{2.034935in}{1.834889in}}%
\pgfpathlineto{\pgfqpoint{2.025490in}{1.836778in}}%
\pgfpathlineto{\pgfqpoint{2.025168in}{1.813601in}}%
\pgfpathclose%
\pgfusepath{fill}%
\end{pgfscope}%
\begin{pgfscope}%
\pgfpathrectangle{\pgfqpoint{0.100000in}{0.100000in}}{\pgfqpoint{3.608454in}{2.310000in}}%
\pgfusepath{clip}%
\pgfsetbuttcap%
\pgfsetmiterjoin%
\definecolor{currentfill}{rgb}{0.000000,0.647059,0.676471}%
\pgfsetfillcolor{currentfill}%
\pgfsetlinewidth{0.000000pt}%
\definecolor{currentstroke}{rgb}{0.000000,0.000000,0.000000}%
\pgfsetstrokecolor{currentstroke}%
\pgfsetstrokeopacity{0.000000}%
\pgfsetdash{}{0pt}%
\pgfpathmoveto{\pgfqpoint{2.661344in}{1.549614in}}%
\pgfpathlineto{\pgfqpoint{2.660092in}{1.560783in}}%
\pgfpathlineto{\pgfqpoint{2.636183in}{1.558287in}}%
\pgfpathlineto{\pgfqpoint{2.633441in}{1.583836in}}%
\pgfpathlineto{\pgfqpoint{2.649386in}{1.585311in}}%
\pgfpathlineto{\pgfqpoint{2.650897in}{1.588647in}}%
\pgfpathlineto{\pgfqpoint{2.648568in}{1.609914in}}%
\pgfpathlineto{\pgfqpoint{2.703247in}{1.616030in}}%
\pgfpathlineto{\pgfqpoint{2.706079in}{1.591455in}}%
\pgfpathlineto{\pgfqpoint{2.684290in}{1.589030in}}%
\pgfpathlineto{\pgfqpoint{2.688778in}{1.550091in}}%
\pgfpathlineto{\pgfqpoint{2.668362in}{1.548067in}}%
\pgfpathlineto{\pgfqpoint{2.661344in}{1.549614in}}%
\pgfpathclose%
\pgfusepath{fill}%
\end{pgfscope}%
\begin{pgfscope}%
\pgfpathrectangle{\pgfqpoint{0.100000in}{0.100000in}}{\pgfqpoint{3.608454in}{2.310000in}}%
\pgfusepath{clip}%
\pgfsetbuttcap%
\pgfsetmiterjoin%
\definecolor{currentfill}{rgb}{0.000000,0.000000,1.000000}%
\pgfsetfillcolor{currentfill}%
\pgfsetlinewidth{0.000000pt}%
\definecolor{currentstroke}{rgb}{0.000000,0.000000,0.000000}%
\pgfsetstrokecolor{currentstroke}%
\pgfsetstrokeopacity{0.000000}%
\pgfsetdash{}{0pt}%
\pgfpathmoveto{\pgfqpoint{1.455629in}{1.496066in}}%
\pgfpathlineto{\pgfqpoint{1.459655in}{1.485514in}}%
\pgfpathlineto{\pgfqpoint{1.465579in}{1.476906in}}%
\pgfpathlineto{\pgfqpoint{1.464403in}{1.459400in}}%
\pgfpathlineto{\pgfqpoint{1.461324in}{1.456249in}}%
\pgfpathlineto{\pgfqpoint{1.459571in}{1.445102in}}%
\pgfpathlineto{\pgfqpoint{1.450957in}{1.440841in}}%
\pgfpathlineto{\pgfqpoint{1.444005in}{1.434780in}}%
\pgfpathlineto{\pgfqpoint{1.437752in}{1.434786in}}%
\pgfpathlineto{\pgfqpoint{1.435412in}{1.444400in}}%
\pgfpathlineto{\pgfqpoint{1.426636in}{1.453930in}}%
\pgfpathlineto{\pgfqpoint{1.403922in}{1.457766in}}%
\pgfpathlineto{\pgfqpoint{1.404073in}{1.463909in}}%
\pgfpathlineto{\pgfqpoint{1.407785in}{1.490488in}}%
\pgfpathlineto{\pgfqpoint{1.407647in}{1.498770in}}%
\pgfpathlineto{\pgfqpoint{1.410993in}{1.493320in}}%
\pgfpathlineto{\pgfqpoint{1.426803in}{1.487168in}}%
\pgfpathlineto{\pgfqpoint{1.436560in}{1.486098in}}%
\pgfpathlineto{\pgfqpoint{1.444777in}{1.488383in}}%
\pgfpathlineto{\pgfqpoint{1.447426in}{1.486216in}}%
\pgfpathlineto{\pgfqpoint{1.455629in}{1.496066in}}%
\pgfpathclose%
\pgfusepath{fill}%
\end{pgfscope}%
\begin{pgfscope}%
\pgfpathrectangle{\pgfqpoint{0.100000in}{0.100000in}}{\pgfqpoint{3.608454in}{2.310000in}}%
\pgfusepath{clip}%
\pgfsetbuttcap%
\pgfsetmiterjoin%
\definecolor{currentfill}{rgb}{0.000000,0.384314,0.807843}%
\pgfsetfillcolor{currentfill}%
\pgfsetlinewidth{0.000000pt}%
\definecolor{currentstroke}{rgb}{0.000000,0.000000,0.000000}%
\pgfsetstrokecolor{currentstroke}%
\pgfsetstrokeopacity{0.000000}%
\pgfsetdash{}{0pt}%
\pgfpathmoveto{\pgfqpoint{2.638885in}{1.368939in}}%
\pgfpathlineto{\pgfqpoint{2.635715in}{1.368602in}}%
\pgfpathlineto{\pgfqpoint{2.632842in}{1.391462in}}%
\pgfpathlineto{\pgfqpoint{2.617652in}{1.390119in}}%
\pgfpathlineto{\pgfqpoint{2.616632in}{1.400311in}}%
\pgfpathlineto{\pgfqpoint{2.613308in}{1.399991in}}%
\pgfpathlineto{\pgfqpoint{2.614715in}{1.406613in}}%
\pgfpathlineto{\pgfqpoint{2.612117in}{1.412594in}}%
\pgfpathlineto{\pgfqpoint{2.607655in}{1.455712in}}%
\pgfpathlineto{\pgfqpoint{2.634909in}{1.458446in}}%
\pgfpathlineto{\pgfqpoint{2.634617in}{1.461203in}}%
\pgfpathlineto{\pgfqpoint{2.657467in}{1.463709in}}%
\pgfpathlineto{\pgfqpoint{2.659619in}{1.442049in}}%
\pgfpathlineto{\pgfqpoint{2.676851in}{1.443992in}}%
\pgfpathlineto{\pgfqpoint{2.676915in}{1.431338in}}%
\pgfpathlineto{\pgfqpoint{2.694747in}{1.433296in}}%
\pgfpathlineto{\pgfqpoint{2.697158in}{1.412713in}}%
\pgfpathlineto{\pgfqpoint{2.697801in}{1.407041in}}%
\pgfpathlineto{\pgfqpoint{2.677647in}{1.404907in}}%
\pgfpathlineto{\pgfqpoint{2.675138in}{1.396451in}}%
\pgfpathlineto{\pgfqpoint{2.650805in}{1.393553in}}%
\pgfpathlineto{\pgfqpoint{2.653405in}{1.370465in}}%
\pgfpathlineto{\pgfqpoint{2.638885in}{1.368939in}}%
\pgfpathclose%
\pgfusepath{fill}%
\end{pgfscope}%
\begin{pgfscope}%
\pgfpathrectangle{\pgfqpoint{0.100000in}{0.100000in}}{\pgfqpoint{3.608454in}{2.310000in}}%
\pgfusepath{clip}%
\pgfsetbuttcap%
\pgfsetmiterjoin%
\definecolor{currentfill}{rgb}{0.000000,0.827451,0.586275}%
\pgfsetfillcolor{currentfill}%
\pgfsetlinewidth{0.000000pt}%
\definecolor{currentstroke}{rgb}{0.000000,0.000000,0.000000}%
\pgfsetstrokecolor{currentstroke}%
\pgfsetstrokeopacity{0.000000}%
\pgfsetdash{}{0pt}%
\pgfpathmoveto{\pgfqpoint{0.899087in}{1.016739in}}%
\pgfpathlineto{\pgfqpoint{0.885192in}{0.948795in}}%
\pgfpathlineto{\pgfqpoint{0.877772in}{0.912470in}}%
\pgfpathlineto{\pgfqpoint{0.787923in}{0.968647in}}%
\pgfpathlineto{\pgfqpoint{0.790528in}{0.978335in}}%
\pgfpathlineto{\pgfqpoint{0.797923in}{0.984723in}}%
\pgfpathlineto{\pgfqpoint{0.705665in}{0.997629in}}%
\pgfpathlineto{\pgfqpoint{0.714200in}{1.032869in}}%
\pgfpathlineto{\pgfqpoint{0.715639in}{1.032543in}}%
\pgfpathlineto{\pgfqpoint{0.721788in}{1.059750in}}%
\pgfpathlineto{\pgfqpoint{0.787832in}{1.044715in}}%
\pgfpathlineto{\pgfqpoint{0.809266in}{1.037990in}}%
\pgfpathlineto{\pgfqpoint{0.809954in}{1.027225in}}%
\pgfpathlineto{\pgfqpoint{0.805414in}{1.014538in}}%
\pgfpathlineto{\pgfqpoint{0.807337in}{1.008296in}}%
\pgfpathlineto{\pgfqpoint{0.832599in}{1.002476in}}%
\pgfpathlineto{\pgfqpoint{0.839911in}{1.036150in}}%
\pgfpathlineto{\pgfqpoint{0.859997in}{1.031950in}}%
\pgfpathlineto{\pgfqpoint{0.858565in}{1.025184in}}%
\pgfpathlineto{\pgfqpoint{0.899087in}{1.016739in}}%
\pgfpathclose%
\pgfusepath{fill}%
\end{pgfscope}%
\begin{pgfscope}%
\pgfpathrectangle{\pgfqpoint{3.319565in}{0.488889in}}{\pgfqpoint{0.360845in}{0.924000in}}%
\pgfusepath{clip}%
\pgfsetbuttcap%
\pgfsetmiterjoin%
\definecolor{currentfill}{rgb}{0.000000,0.854902,0.572549}%
\pgfsetfillcolor{currentfill}%
\pgfsetlinewidth{0.000000pt}%
\definecolor{currentstroke}{rgb}{0.000000,0.000000,0.000000}%
\pgfsetstrokecolor{currentstroke}%
\pgfsetstrokeopacity{0.000000}%
\pgfsetdash{}{0pt}%
\pgfpathmoveto{\pgfqpoint{3.581475in}{0.488889in}}%
\pgfpathlineto{\pgfqpoint{3.680411in}{0.488889in}}%
\pgfpathlineto{\pgfqpoint{3.680411in}{0.581289in}}%
\pgfpathlineto{\pgfqpoint{3.581475in}{0.581289in}}%
\pgfpathlineto{\pgfqpoint{3.581475in}{0.488889in}}%
\pgfpathclose%
\pgfusepath{fill}%
\end{pgfscope}%
\begin{pgfscope}%
\pgfpathrectangle{\pgfqpoint{3.319565in}{0.488889in}}{\pgfqpoint{0.360845in}{0.924000in}}%
\pgfusepath{clip}%
\pgfsetbuttcap%
\pgfsetmiterjoin%
\definecolor{currentfill}{rgb}{0.000000,0.815686,0.592157}%
\pgfsetfillcolor{currentfill}%
\pgfsetlinewidth{0.000000pt}%
\definecolor{currentstroke}{rgb}{0.000000,0.000000,0.000000}%
\pgfsetstrokecolor{currentstroke}%
\pgfsetstrokeopacity{0.000000}%
\pgfsetdash{}{0pt}%
\pgfpathmoveto{\pgfqpoint{3.578539in}{0.498129in}}%
\pgfpathlineto{\pgfqpoint{3.680411in}{0.498129in}}%
\pgfpathlineto{\pgfqpoint{3.680411in}{0.590529in}}%
\pgfpathlineto{\pgfqpoint{3.578539in}{0.590529in}}%
\pgfpathlineto{\pgfqpoint{3.578539in}{0.498129in}}%
\pgfpathclose%
\pgfusepath{fill}%
\end{pgfscope}%
\begin{pgfscope}%
\pgfpathrectangle{\pgfqpoint{3.319565in}{0.488889in}}{\pgfqpoint{0.360845in}{0.924000in}}%
\pgfusepath{clip}%
\pgfsetbuttcap%
\pgfsetmiterjoin%
\definecolor{currentfill}{rgb}{0.000000,0.788235,0.605882}%
\pgfsetfillcolor{currentfill}%
\pgfsetlinewidth{0.000000pt}%
\definecolor{currentstroke}{rgb}{0.000000,0.000000,0.000000}%
\pgfsetstrokecolor{currentstroke}%
\pgfsetstrokeopacity{0.000000}%
\pgfsetdash{}{0pt}%
\pgfpathmoveto{\pgfqpoint{3.576250in}{0.507369in}}%
\pgfpathlineto{\pgfqpoint{3.680411in}{0.507369in}}%
\pgfpathlineto{\pgfqpoint{3.680411in}{0.599769in}}%
\pgfpathlineto{\pgfqpoint{3.576250in}{0.599769in}}%
\pgfpathlineto{\pgfqpoint{3.576250in}{0.507369in}}%
\pgfpathclose%
\pgfusepath{fill}%
\end{pgfscope}%
\begin{pgfscope}%
\pgfpathrectangle{\pgfqpoint{3.319565in}{0.488889in}}{\pgfqpoint{0.360845in}{0.924000in}}%
\pgfusepath{clip}%
\pgfsetbuttcap%
\pgfsetmiterjoin%
\definecolor{currentfill}{rgb}{0.000000,0.760784,0.619608}%
\pgfsetfillcolor{currentfill}%
\pgfsetlinewidth{0.000000pt}%
\definecolor{currentstroke}{rgb}{0.000000,0.000000,0.000000}%
\pgfsetstrokecolor{currentstroke}%
\pgfsetstrokeopacity{0.000000}%
\pgfsetdash{}{0pt}%
\pgfpathmoveto{\pgfqpoint{3.574081in}{0.516609in}}%
\pgfpathlineto{\pgfqpoint{3.680411in}{0.516609in}}%
\pgfpathlineto{\pgfqpoint{3.680411in}{0.609009in}}%
\pgfpathlineto{\pgfqpoint{3.574081in}{0.609009in}}%
\pgfpathlineto{\pgfqpoint{3.574081in}{0.516609in}}%
\pgfpathclose%
\pgfusepath{fill}%
\end{pgfscope}%
\begin{pgfscope}%
\pgfpathrectangle{\pgfqpoint{3.319565in}{0.488889in}}{\pgfqpoint{0.360845in}{0.924000in}}%
\pgfusepath{clip}%
\pgfsetbuttcap%
\pgfsetmiterjoin%
\definecolor{currentfill}{rgb}{0.000000,0.749020,0.625490}%
\pgfsetfillcolor{currentfill}%
\pgfsetlinewidth{0.000000pt}%
\definecolor{currentstroke}{rgb}{0.000000,0.000000,0.000000}%
\pgfsetstrokecolor{currentstroke}%
\pgfsetstrokeopacity{0.000000}%
\pgfsetdash{}{0pt}%
\pgfpathmoveto{\pgfqpoint{3.573102in}{0.525849in}}%
\pgfpathlineto{\pgfqpoint{3.680411in}{0.525849in}}%
\pgfpathlineto{\pgfqpoint{3.680411in}{0.618249in}}%
\pgfpathlineto{\pgfqpoint{3.573102in}{0.618249in}}%
\pgfpathlineto{\pgfqpoint{3.573102in}{0.525849in}}%
\pgfpathclose%
\pgfusepath{fill}%
\end{pgfscope}%
\begin{pgfscope}%
\pgfpathrectangle{\pgfqpoint{3.319565in}{0.488889in}}{\pgfqpoint{0.360845in}{0.924000in}}%
\pgfusepath{clip}%
\pgfsetbuttcap%
\pgfsetmiterjoin%
\definecolor{currentfill}{rgb}{0.000000,0.713725,0.643137}%
\pgfsetfillcolor{currentfill}%
\pgfsetlinewidth{0.000000pt}%
\definecolor{currentstroke}{rgb}{0.000000,0.000000,0.000000}%
\pgfsetstrokecolor{currentstroke}%
\pgfsetstrokeopacity{0.000000}%
\pgfsetdash{}{0pt}%
\pgfpathmoveto{\pgfqpoint{3.570208in}{0.535089in}}%
\pgfpathlineto{\pgfqpoint{3.680411in}{0.535089in}}%
\pgfpathlineto{\pgfqpoint{3.680411in}{0.627489in}}%
\pgfpathlineto{\pgfqpoint{3.570208in}{0.627489in}}%
\pgfpathlineto{\pgfqpoint{3.570208in}{0.535089in}}%
\pgfpathclose%
\pgfusepath{fill}%
\end{pgfscope}%
\begin{pgfscope}%
\pgfpathrectangle{\pgfqpoint{3.319565in}{0.488889in}}{\pgfqpoint{0.360845in}{0.924000in}}%
\pgfusepath{clip}%
\pgfsetbuttcap%
\pgfsetmiterjoin%
\definecolor{currentfill}{rgb}{0.000000,0.701961,0.649020}%
\pgfsetfillcolor{currentfill}%
\pgfsetlinewidth{0.000000pt}%
\definecolor{currentstroke}{rgb}{0.000000,0.000000,0.000000}%
\pgfsetstrokecolor{currentstroke}%
\pgfsetstrokeopacity{0.000000}%
\pgfsetdash{}{0pt}%
\pgfpathmoveto{\pgfqpoint{3.569139in}{0.544329in}}%
\pgfpathlineto{\pgfqpoint{3.680411in}{0.544329in}}%
\pgfpathlineto{\pgfqpoint{3.680411in}{0.636729in}}%
\pgfpathlineto{\pgfqpoint{3.569139in}{0.636729in}}%
\pgfpathlineto{\pgfqpoint{3.569139in}{0.544329in}}%
\pgfpathclose%
\pgfusepath{fill}%
\end{pgfscope}%
\begin{pgfscope}%
\pgfpathrectangle{\pgfqpoint{3.319565in}{0.488889in}}{\pgfqpoint{0.360845in}{0.924000in}}%
\pgfusepath{clip}%
\pgfsetbuttcap%
\pgfsetmiterjoin%
\definecolor{currentfill}{rgb}{0.000000,0.690196,0.654902}%
\pgfsetfillcolor{currentfill}%
\pgfsetlinewidth{0.000000pt}%
\definecolor{currentstroke}{rgb}{0.000000,0.000000,0.000000}%
\pgfsetstrokecolor{currentstroke}%
\pgfsetstrokeopacity{0.000000}%
\pgfsetdash{}{0pt}%
\pgfpathmoveto{\pgfqpoint{3.568218in}{0.553569in}}%
\pgfpathlineto{\pgfqpoint{3.680411in}{0.553569in}}%
\pgfpathlineto{\pgfqpoint{3.680411in}{0.645969in}}%
\pgfpathlineto{\pgfqpoint{3.568218in}{0.645969in}}%
\pgfpathlineto{\pgfqpoint{3.568218in}{0.553569in}}%
\pgfpathclose%
\pgfusepath{fill}%
\end{pgfscope}%
\begin{pgfscope}%
\pgfpathrectangle{\pgfqpoint{3.319565in}{0.488889in}}{\pgfqpoint{0.360845in}{0.924000in}}%
\pgfusepath{clip}%
\pgfsetbuttcap%
\pgfsetmiterjoin%
\definecolor{currentfill}{rgb}{0.000000,0.678431,0.660784}%
\pgfsetfillcolor{currentfill}%
\pgfsetlinewidth{0.000000pt}%
\definecolor{currentstroke}{rgb}{0.000000,0.000000,0.000000}%
\pgfsetstrokecolor{currentstroke}%
\pgfsetstrokeopacity{0.000000}%
\pgfsetdash{}{0pt}%
\pgfpathmoveto{\pgfqpoint{3.567329in}{0.562809in}}%
\pgfpathlineto{\pgfqpoint{3.680411in}{0.562809in}}%
\pgfpathlineto{\pgfqpoint{3.680411in}{0.655209in}}%
\pgfpathlineto{\pgfqpoint{3.567329in}{0.655209in}}%
\pgfpathlineto{\pgfqpoint{3.567329in}{0.562809in}}%
\pgfpathclose%
\pgfusepath{fill}%
\end{pgfscope}%
\begin{pgfscope}%
\pgfpathrectangle{\pgfqpoint{3.319565in}{0.488889in}}{\pgfqpoint{0.360845in}{0.924000in}}%
\pgfusepath{clip}%
\pgfsetbuttcap%
\pgfsetmiterjoin%
\definecolor{currentfill}{rgb}{0.000000,0.670588,0.664706}%
\pgfsetfillcolor{currentfill}%
\pgfsetlinewidth{0.000000pt}%
\definecolor{currentstroke}{rgb}{0.000000,0.000000,0.000000}%
\pgfsetstrokecolor{currentstroke}%
\pgfsetstrokeopacity{0.000000}%
\pgfsetdash{}{0pt}%
\pgfpathmoveto{\pgfqpoint{3.566822in}{0.572049in}}%
\pgfpathlineto{\pgfqpoint{3.680411in}{0.572049in}}%
\pgfpathlineto{\pgfqpoint{3.680411in}{0.664449in}}%
\pgfpathlineto{\pgfqpoint{3.566822in}{0.664449in}}%
\pgfpathlineto{\pgfqpoint{3.566822in}{0.572049in}}%
\pgfpathclose%
\pgfusepath{fill}%
\end{pgfscope}%
\begin{pgfscope}%
\pgfpathrectangle{\pgfqpoint{3.319565in}{0.488889in}}{\pgfqpoint{0.360845in}{0.924000in}}%
\pgfusepath{clip}%
\pgfsetbuttcap%
\pgfsetmiterjoin%
\definecolor{currentfill}{rgb}{0.000000,0.666667,0.666667}%
\pgfsetfillcolor{currentfill}%
\pgfsetlinewidth{0.000000pt}%
\definecolor{currentstroke}{rgb}{0.000000,0.000000,0.000000}%
\pgfsetstrokecolor{currentstroke}%
\pgfsetstrokeopacity{0.000000}%
\pgfsetdash{}{0pt}%
\pgfpathmoveto{\pgfqpoint{3.566410in}{0.581289in}}%
\pgfpathlineto{\pgfqpoint{3.680411in}{0.581289in}}%
\pgfpathlineto{\pgfqpoint{3.680411in}{0.673689in}}%
\pgfpathlineto{\pgfqpoint{3.566410in}{0.673689in}}%
\pgfpathlineto{\pgfqpoint{3.566410in}{0.581289in}}%
\pgfpathclose%
\pgfusepath{fill}%
\end{pgfscope}%
\begin{pgfscope}%
\pgfpathrectangle{\pgfqpoint{3.319565in}{0.488889in}}{\pgfqpoint{0.360845in}{0.924000in}}%
\pgfusepath{clip}%
\pgfsetbuttcap%
\pgfsetmiterjoin%
\definecolor{currentfill}{rgb}{0.000000,0.654902,0.672549}%
\pgfsetfillcolor{currentfill}%
\pgfsetlinewidth{0.000000pt}%
\definecolor{currentstroke}{rgb}{0.000000,0.000000,0.000000}%
\pgfsetstrokecolor{currentstroke}%
\pgfsetstrokeopacity{0.000000}%
\pgfsetdash{}{0pt}%
\pgfpathmoveto{\pgfqpoint{3.565309in}{0.590529in}}%
\pgfpathlineto{\pgfqpoint{3.680411in}{0.590529in}}%
\pgfpathlineto{\pgfqpoint{3.680411in}{0.682929in}}%
\pgfpathlineto{\pgfqpoint{3.565309in}{0.682929in}}%
\pgfpathlineto{\pgfqpoint{3.565309in}{0.590529in}}%
\pgfpathclose%
\pgfusepath{fill}%
\end{pgfscope}%
\begin{pgfscope}%
\pgfpathrectangle{\pgfqpoint{3.319565in}{0.488889in}}{\pgfqpoint{0.360845in}{0.924000in}}%
\pgfusepath{clip}%
\pgfsetbuttcap%
\pgfsetmiterjoin%
\definecolor{currentfill}{rgb}{0.000000,0.643137,0.678431}%
\pgfsetfillcolor{currentfill}%
\pgfsetlinewidth{0.000000pt}%
\definecolor{currentstroke}{rgb}{0.000000,0.000000,0.000000}%
\pgfsetstrokecolor{currentstroke}%
\pgfsetstrokeopacity{0.000000}%
\pgfsetdash{}{0pt}%
\pgfpathmoveto{\pgfqpoint{3.564406in}{0.599769in}}%
\pgfpathlineto{\pgfqpoint{3.680411in}{0.599769in}}%
\pgfpathlineto{\pgfqpoint{3.680411in}{0.692169in}}%
\pgfpathlineto{\pgfqpoint{3.564406in}{0.692169in}}%
\pgfpathlineto{\pgfqpoint{3.564406in}{0.599769in}}%
\pgfpathclose%
\pgfusepath{fill}%
\end{pgfscope}%
\begin{pgfscope}%
\pgfpathrectangle{\pgfqpoint{3.319565in}{0.488889in}}{\pgfqpoint{0.360845in}{0.924000in}}%
\pgfusepath{clip}%
\pgfsetbuttcap%
\pgfsetmiterjoin%
\definecolor{currentfill}{rgb}{0.000000,0.631373,0.684314}%
\pgfsetfillcolor{currentfill}%
\pgfsetlinewidth{0.000000pt}%
\definecolor{currentstroke}{rgb}{0.000000,0.000000,0.000000}%
\pgfsetstrokecolor{currentstroke}%
\pgfsetstrokeopacity{0.000000}%
\pgfsetdash{}{0pt}%
\pgfpathmoveto{\pgfqpoint{3.563660in}{0.609009in}}%
\pgfpathlineto{\pgfqpoint{3.680411in}{0.609009in}}%
\pgfpathlineto{\pgfqpoint{3.680411in}{0.701409in}}%
\pgfpathlineto{\pgfqpoint{3.563660in}{0.701409in}}%
\pgfpathlineto{\pgfqpoint{3.563660in}{0.609009in}}%
\pgfpathclose%
\pgfusepath{fill}%
\end{pgfscope}%
\begin{pgfscope}%
\pgfpathrectangle{\pgfqpoint{3.319565in}{0.488889in}}{\pgfqpoint{0.360845in}{0.924000in}}%
\pgfusepath{clip}%
\pgfsetbuttcap%
\pgfsetmiterjoin%
\definecolor{currentfill}{rgb}{0.000000,0.627451,0.686275}%
\pgfsetfillcolor{currentfill}%
\pgfsetlinewidth{0.000000pt}%
\definecolor{currentstroke}{rgb}{0.000000,0.000000,0.000000}%
\pgfsetstrokecolor{currentstroke}%
\pgfsetstrokeopacity{0.000000}%
\pgfsetdash{}{0pt}%
\pgfpathmoveto{\pgfqpoint{3.563249in}{0.618249in}}%
\pgfpathlineto{\pgfqpoint{3.680411in}{0.618249in}}%
\pgfpathlineto{\pgfqpoint{3.680411in}{0.710649in}}%
\pgfpathlineto{\pgfqpoint{3.563249in}{0.710649in}}%
\pgfpathlineto{\pgfqpoint{3.563249in}{0.618249in}}%
\pgfpathclose%
\pgfusepath{fill}%
\end{pgfscope}%
\begin{pgfscope}%
\pgfpathrectangle{\pgfqpoint{3.319565in}{0.488889in}}{\pgfqpoint{0.360845in}{0.924000in}}%
\pgfusepath{clip}%
\pgfsetbuttcap%
\pgfsetmiterjoin%
\definecolor{currentfill}{rgb}{0.000000,0.607843,0.696078}%
\pgfsetfillcolor{currentfill}%
\pgfsetlinewidth{0.000000pt}%
\definecolor{currentstroke}{rgb}{0.000000,0.000000,0.000000}%
\pgfsetstrokecolor{currentstroke}%
\pgfsetstrokeopacity{0.000000}%
\pgfsetdash{}{0pt}%
\pgfpathmoveto{\pgfqpoint{3.561733in}{0.627489in}}%
\pgfpathlineto{\pgfqpoint{3.680411in}{0.627489in}}%
\pgfpathlineto{\pgfqpoint{3.680411in}{0.719889in}}%
\pgfpathlineto{\pgfqpoint{3.561733in}{0.719889in}}%
\pgfpathlineto{\pgfqpoint{3.561733in}{0.627489in}}%
\pgfpathclose%
\pgfusepath{fill}%
\end{pgfscope}%
\begin{pgfscope}%
\pgfpathrectangle{\pgfqpoint{3.319565in}{0.488889in}}{\pgfqpoint{0.360845in}{0.924000in}}%
\pgfusepath{clip}%
\pgfsetbuttcap%
\pgfsetmiterjoin%
\definecolor{currentfill}{rgb}{0.000000,0.603922,0.698039}%
\pgfsetfillcolor{currentfill}%
\pgfsetlinewidth{0.000000pt}%
\definecolor{currentstroke}{rgb}{0.000000,0.000000,0.000000}%
\pgfsetstrokecolor{currentstroke}%
\pgfsetstrokeopacity{0.000000}%
\pgfsetdash{}{0pt}%
\pgfpathmoveto{\pgfqpoint{3.561363in}{0.636729in}}%
\pgfpathlineto{\pgfqpoint{3.680411in}{0.636729in}}%
\pgfpathlineto{\pgfqpoint{3.680411in}{0.729129in}}%
\pgfpathlineto{\pgfqpoint{3.561363in}{0.729129in}}%
\pgfpathlineto{\pgfqpoint{3.561363in}{0.636729in}}%
\pgfpathclose%
\pgfusepath{fill}%
\end{pgfscope}%
\begin{pgfscope}%
\pgfpathrectangle{\pgfqpoint{3.319565in}{0.488889in}}{\pgfqpoint{0.360845in}{0.924000in}}%
\pgfusepath{clip}%
\pgfsetbuttcap%
\pgfsetmiterjoin%
\definecolor{currentfill}{rgb}{0.000000,0.600000,0.700000}%
\pgfsetfillcolor{currentfill}%
\pgfsetlinewidth{0.000000pt}%
\definecolor{currentstroke}{rgb}{0.000000,0.000000,0.000000}%
\pgfsetstrokecolor{currentstroke}%
\pgfsetstrokeopacity{0.000000}%
\pgfsetdash{}{0pt}%
\pgfpathmoveto{\pgfqpoint{3.560871in}{0.645969in}}%
\pgfpathlineto{\pgfqpoint{3.680411in}{0.645969in}}%
\pgfpathlineto{\pgfqpoint{3.680411in}{0.738369in}}%
\pgfpathlineto{\pgfqpoint{3.560871in}{0.738369in}}%
\pgfpathlineto{\pgfqpoint{3.560871in}{0.645969in}}%
\pgfpathclose%
\pgfusepath{fill}%
\end{pgfscope}%
\begin{pgfscope}%
\pgfpathrectangle{\pgfqpoint{3.319565in}{0.488889in}}{\pgfqpoint{0.360845in}{0.924000in}}%
\pgfusepath{clip}%
\pgfsetbuttcap%
\pgfsetmiterjoin%
\definecolor{currentfill}{rgb}{0.000000,0.596078,0.701961}%
\pgfsetfillcolor{currentfill}%
\pgfsetlinewidth{0.000000pt}%
\definecolor{currentstroke}{rgb}{0.000000,0.000000,0.000000}%
\pgfsetstrokecolor{currentstroke}%
\pgfsetstrokeopacity{0.000000}%
\pgfsetdash{}{0pt}%
\pgfpathmoveto{\pgfqpoint{3.560811in}{0.655209in}}%
\pgfpathlineto{\pgfqpoint{3.680411in}{0.655209in}}%
\pgfpathlineto{\pgfqpoint{3.680411in}{0.747609in}}%
\pgfpathlineto{\pgfqpoint{3.560811in}{0.747609in}}%
\pgfpathlineto{\pgfqpoint{3.560811in}{0.655209in}}%
\pgfpathclose%
\pgfusepath{fill}%
\end{pgfscope}%
\begin{pgfscope}%
\pgfpathrectangle{\pgfqpoint{3.319565in}{0.488889in}}{\pgfqpoint{0.360845in}{0.924000in}}%
\pgfusepath{clip}%
\pgfsetbuttcap%
\pgfsetmiterjoin%
\definecolor{currentfill}{rgb}{0.000000,0.584314,0.707843}%
\pgfsetfillcolor{currentfill}%
\pgfsetlinewidth{0.000000pt}%
\definecolor{currentstroke}{rgb}{0.000000,0.000000,0.000000}%
\pgfsetstrokecolor{currentstroke}%
\pgfsetstrokeopacity{0.000000}%
\pgfsetdash{}{0pt}%
\pgfpathmoveto{\pgfqpoint{3.559810in}{0.664449in}}%
\pgfpathlineto{\pgfqpoint{3.680411in}{0.664449in}}%
\pgfpathlineto{\pgfqpoint{3.680411in}{0.756849in}}%
\pgfpathlineto{\pgfqpoint{3.559810in}{0.756849in}}%
\pgfpathlineto{\pgfqpoint{3.559810in}{0.664449in}}%
\pgfpathclose%
\pgfusepath{fill}%
\end{pgfscope}%
\begin{pgfscope}%
\pgfpathrectangle{\pgfqpoint{3.319565in}{0.488889in}}{\pgfqpoint{0.360845in}{0.924000in}}%
\pgfusepath{clip}%
\pgfsetbuttcap%
\pgfsetmiterjoin%
\definecolor{currentfill}{rgb}{0.000000,0.576471,0.711765}%
\pgfsetfillcolor{currentfill}%
\pgfsetlinewidth{0.000000pt}%
\definecolor{currentstroke}{rgb}{0.000000,0.000000,0.000000}%
\pgfsetstrokecolor{currentstroke}%
\pgfsetstrokeopacity{0.000000}%
\pgfsetdash{}{0pt}%
\pgfpathmoveto{\pgfqpoint{3.559144in}{0.673689in}}%
\pgfpathlineto{\pgfqpoint{3.680411in}{0.673689in}}%
\pgfpathlineto{\pgfqpoint{3.680411in}{0.766089in}}%
\pgfpathlineto{\pgfqpoint{3.559144in}{0.766089in}}%
\pgfpathlineto{\pgfqpoint{3.559144in}{0.673689in}}%
\pgfpathclose%
\pgfusepath{fill}%
\end{pgfscope}%
\begin{pgfscope}%
\pgfpathrectangle{\pgfqpoint{3.319565in}{0.488889in}}{\pgfqpoint{0.360845in}{0.924000in}}%
\pgfusepath{clip}%
\pgfsetbuttcap%
\pgfsetmiterjoin%
\definecolor{currentfill}{rgb}{0.000000,0.568627,0.715686}%
\pgfsetfillcolor{currentfill}%
\pgfsetlinewidth{0.000000pt}%
\definecolor{currentstroke}{rgb}{0.000000,0.000000,0.000000}%
\pgfsetstrokecolor{currentstroke}%
\pgfsetstrokeopacity{0.000000}%
\pgfsetdash{}{0pt}%
\pgfpathmoveto{\pgfqpoint{3.558529in}{0.682929in}}%
\pgfpathlineto{\pgfqpoint{3.680411in}{0.682929in}}%
\pgfpathlineto{\pgfqpoint{3.680411in}{0.775329in}}%
\pgfpathlineto{\pgfqpoint{3.558529in}{0.775329in}}%
\pgfpathlineto{\pgfqpoint{3.558529in}{0.682929in}}%
\pgfpathclose%
\pgfusepath{fill}%
\end{pgfscope}%
\begin{pgfscope}%
\pgfpathrectangle{\pgfqpoint{3.319565in}{0.488889in}}{\pgfqpoint{0.360845in}{0.924000in}}%
\pgfusepath{clip}%
\pgfsetbuttcap%
\pgfsetmiterjoin%
\definecolor{currentfill}{rgb}{0.000000,0.564706,0.717647}%
\pgfsetfillcolor{currentfill}%
\pgfsetlinewidth{0.000000pt}%
\definecolor{currentstroke}{rgb}{0.000000,0.000000,0.000000}%
\pgfsetstrokecolor{currentstroke}%
\pgfsetstrokeopacity{0.000000}%
\pgfsetdash{}{0pt}%
\pgfpathmoveto{\pgfqpoint{3.558230in}{0.692169in}}%
\pgfpathlineto{\pgfqpoint{3.680411in}{0.692169in}}%
\pgfpathlineto{\pgfqpoint{3.680411in}{0.784569in}}%
\pgfpathlineto{\pgfqpoint{3.558230in}{0.784569in}}%
\pgfpathlineto{\pgfqpoint{3.558230in}{0.692169in}}%
\pgfpathclose%
\pgfusepath{fill}%
\end{pgfscope}%
\begin{pgfscope}%
\pgfpathrectangle{\pgfqpoint{3.319565in}{0.488889in}}{\pgfqpoint{0.360845in}{0.924000in}}%
\pgfusepath{clip}%
\pgfsetbuttcap%
\pgfsetmiterjoin%
\definecolor{currentfill}{rgb}{0.000000,0.560784,0.719608}%
\pgfsetfillcolor{currentfill}%
\pgfsetlinewidth{0.000000pt}%
\definecolor{currentstroke}{rgb}{0.000000,0.000000,0.000000}%
\pgfsetstrokecolor{currentstroke}%
\pgfsetstrokeopacity{0.000000}%
\pgfsetdash{}{0pt}%
\pgfpathmoveto{\pgfqpoint{3.557690in}{0.701409in}}%
\pgfpathlineto{\pgfqpoint{3.680411in}{0.701409in}}%
\pgfpathlineto{\pgfqpoint{3.680411in}{0.793809in}}%
\pgfpathlineto{\pgfqpoint{3.557690in}{0.793809in}}%
\pgfpathlineto{\pgfqpoint{3.557690in}{0.701409in}}%
\pgfpathclose%
\pgfusepath{fill}%
\end{pgfscope}%
\begin{pgfscope}%
\pgfpathrectangle{\pgfqpoint{3.319565in}{0.488889in}}{\pgfqpoint{0.360845in}{0.924000in}}%
\pgfusepath{clip}%
\pgfsetbuttcap%
\pgfsetmiterjoin%
\definecolor{currentfill}{rgb}{0.000000,0.552941,0.723529}%
\pgfsetfillcolor{currentfill}%
\pgfsetlinewidth{0.000000pt}%
\definecolor{currentstroke}{rgb}{0.000000,0.000000,0.000000}%
\pgfsetstrokecolor{currentstroke}%
\pgfsetstrokeopacity{0.000000}%
\pgfsetdash{}{0pt}%
\pgfpathmoveto{\pgfqpoint{3.557112in}{0.710649in}}%
\pgfpathlineto{\pgfqpoint{3.680411in}{0.710649in}}%
\pgfpathlineto{\pgfqpoint{3.680411in}{0.803049in}}%
\pgfpathlineto{\pgfqpoint{3.557112in}{0.803049in}}%
\pgfpathlineto{\pgfqpoint{3.557112in}{0.710649in}}%
\pgfpathclose%
\pgfusepath{fill}%
\end{pgfscope}%
\begin{pgfscope}%
\pgfpathrectangle{\pgfqpoint{3.319565in}{0.488889in}}{\pgfqpoint{0.360845in}{0.924000in}}%
\pgfusepath{clip}%
\pgfsetbuttcap%
\pgfsetmiterjoin%
\definecolor{currentfill}{rgb}{0.000000,0.549020,0.725490}%
\pgfsetfillcolor{currentfill}%
\pgfsetlinewidth{0.000000pt}%
\definecolor{currentstroke}{rgb}{0.000000,0.000000,0.000000}%
\pgfsetstrokecolor{currentstroke}%
\pgfsetstrokeopacity{0.000000}%
\pgfsetdash{}{0pt}%
\pgfpathmoveto{\pgfqpoint{3.556888in}{0.719889in}}%
\pgfpathlineto{\pgfqpoint{3.680411in}{0.719889in}}%
\pgfpathlineto{\pgfqpoint{3.680411in}{0.812289in}}%
\pgfpathlineto{\pgfqpoint{3.556888in}{0.812289in}}%
\pgfpathlineto{\pgfqpoint{3.556888in}{0.719889in}}%
\pgfpathclose%
\pgfusepath{fill}%
\end{pgfscope}%
\begin{pgfscope}%
\pgfpathrectangle{\pgfqpoint{3.319565in}{0.488889in}}{\pgfqpoint{0.360845in}{0.924000in}}%
\pgfusepath{clip}%
\pgfsetbuttcap%
\pgfsetmiterjoin%
\definecolor{currentfill}{rgb}{0.000000,0.545098,0.727451}%
\pgfsetfillcolor{currentfill}%
\pgfsetlinewidth{0.000000pt}%
\definecolor{currentstroke}{rgb}{0.000000,0.000000,0.000000}%
\pgfsetstrokecolor{currentstroke}%
\pgfsetstrokeopacity{0.000000}%
\pgfsetdash{}{0pt}%
\pgfpathmoveto{\pgfqpoint{3.556562in}{0.729129in}}%
\pgfpathlineto{\pgfqpoint{3.680411in}{0.729129in}}%
\pgfpathlineto{\pgfqpoint{3.680411in}{0.821529in}}%
\pgfpathlineto{\pgfqpoint{3.556562in}{0.821529in}}%
\pgfpathlineto{\pgfqpoint{3.556562in}{0.729129in}}%
\pgfpathclose%
\pgfusepath{fill}%
\end{pgfscope}%
\begin{pgfscope}%
\pgfpathrectangle{\pgfqpoint{3.319565in}{0.488889in}}{\pgfqpoint{0.360845in}{0.924000in}}%
\pgfusepath{clip}%
\pgfsetbuttcap%
\pgfsetmiterjoin%
\definecolor{currentfill}{rgb}{0.000000,0.537255,0.731373}%
\pgfsetfillcolor{currentfill}%
\pgfsetlinewidth{0.000000pt}%
\definecolor{currentstroke}{rgb}{0.000000,0.000000,0.000000}%
\pgfsetstrokecolor{currentstroke}%
\pgfsetstrokeopacity{0.000000}%
\pgfsetdash{}{0pt}%
\pgfpathmoveto{\pgfqpoint{3.556029in}{0.738369in}}%
\pgfpathlineto{\pgfqpoint{3.680411in}{0.738369in}}%
\pgfpathlineto{\pgfqpoint{3.680411in}{0.830769in}}%
\pgfpathlineto{\pgfqpoint{3.556029in}{0.830769in}}%
\pgfpathlineto{\pgfqpoint{3.556029in}{0.738369in}}%
\pgfpathclose%
\pgfusepath{fill}%
\end{pgfscope}%
\begin{pgfscope}%
\pgfpathrectangle{\pgfqpoint{3.319565in}{0.488889in}}{\pgfqpoint{0.360845in}{0.924000in}}%
\pgfusepath{clip}%
\pgfsetbuttcap%
\pgfsetmiterjoin%
\definecolor{currentfill}{rgb}{0.000000,0.537255,0.731373}%
\pgfsetfillcolor{currentfill}%
\pgfsetlinewidth{0.000000pt}%
\definecolor{currentstroke}{rgb}{0.000000,0.000000,0.000000}%
\pgfsetstrokecolor{currentstroke}%
\pgfsetstrokeopacity{0.000000}%
\pgfsetdash{}{0pt}%
\pgfpathmoveto{\pgfqpoint{3.555811in}{0.747609in}}%
\pgfpathlineto{\pgfqpoint{3.680411in}{0.747609in}}%
\pgfpathlineto{\pgfqpoint{3.680411in}{0.840009in}}%
\pgfpathlineto{\pgfqpoint{3.555811in}{0.840009in}}%
\pgfpathlineto{\pgfqpoint{3.555811in}{0.747609in}}%
\pgfpathclose%
\pgfusepath{fill}%
\end{pgfscope}%
\begin{pgfscope}%
\pgfpathrectangle{\pgfqpoint{3.319565in}{0.488889in}}{\pgfqpoint{0.360845in}{0.924000in}}%
\pgfusepath{clip}%
\pgfsetbuttcap%
\pgfsetmiterjoin%
\definecolor{currentfill}{rgb}{0.000000,0.537255,0.731373}%
\pgfsetfillcolor{currentfill}%
\pgfsetlinewidth{0.000000pt}%
\definecolor{currentstroke}{rgb}{0.000000,0.000000,0.000000}%
\pgfsetstrokecolor{currentstroke}%
\pgfsetstrokeopacity{0.000000}%
\pgfsetdash{}{0pt}%
\pgfpathmoveto{\pgfqpoint{3.555811in}{0.756849in}}%
\pgfpathlineto{\pgfqpoint{3.680411in}{0.756849in}}%
\pgfpathlineto{\pgfqpoint{3.680411in}{0.849249in}}%
\pgfpathlineto{\pgfqpoint{3.555811in}{0.849249in}}%
\pgfpathlineto{\pgfqpoint{3.555811in}{0.756849in}}%
\pgfpathclose%
\pgfusepath{fill}%
\end{pgfscope}%
\begin{pgfscope}%
\pgfpathrectangle{\pgfqpoint{3.319565in}{0.488889in}}{\pgfqpoint{0.360845in}{0.924000in}}%
\pgfusepath{clip}%
\pgfsetbuttcap%
\pgfsetmiterjoin%
\definecolor{currentfill}{rgb}{0.000000,0.529412,0.735294}%
\pgfsetfillcolor{currentfill}%
\pgfsetlinewidth{0.000000pt}%
\definecolor{currentstroke}{rgb}{0.000000,0.000000,0.000000}%
\pgfsetstrokecolor{currentstroke}%
\pgfsetstrokeopacity{0.000000}%
\pgfsetdash{}{0pt}%
\pgfpathmoveto{\pgfqpoint{3.555410in}{0.766089in}}%
\pgfpathlineto{\pgfqpoint{3.680411in}{0.766089in}}%
\pgfpathlineto{\pgfqpoint{3.680411in}{0.858489in}}%
\pgfpathlineto{\pgfqpoint{3.555410in}{0.858489in}}%
\pgfpathlineto{\pgfqpoint{3.555410in}{0.766089in}}%
\pgfpathclose%
\pgfusepath{fill}%
\end{pgfscope}%
\begin{pgfscope}%
\pgfpathrectangle{\pgfqpoint{3.319565in}{0.488889in}}{\pgfqpoint{0.360845in}{0.924000in}}%
\pgfusepath{clip}%
\pgfsetbuttcap%
\pgfsetmiterjoin%
\definecolor{currentfill}{rgb}{0.000000,0.525490,0.737255}%
\pgfsetfillcolor{currentfill}%
\pgfsetlinewidth{0.000000pt}%
\definecolor{currentstroke}{rgb}{0.000000,0.000000,0.000000}%
\pgfsetstrokecolor{currentstroke}%
\pgfsetstrokeopacity{0.000000}%
\pgfsetdash{}{0pt}%
\pgfpathmoveto{\pgfqpoint{3.555056in}{0.775329in}}%
\pgfpathlineto{\pgfqpoint{3.680411in}{0.775329in}}%
\pgfpathlineto{\pgfqpoint{3.680411in}{0.867729in}}%
\pgfpathlineto{\pgfqpoint{3.555056in}{0.867729in}}%
\pgfpathlineto{\pgfqpoint{3.555056in}{0.775329in}}%
\pgfpathclose%
\pgfusepath{fill}%
\end{pgfscope}%
\begin{pgfscope}%
\pgfpathrectangle{\pgfqpoint{3.319565in}{0.488889in}}{\pgfqpoint{0.360845in}{0.924000in}}%
\pgfusepath{clip}%
\pgfsetbuttcap%
\pgfsetmiterjoin%
\definecolor{currentfill}{rgb}{0.000000,0.521569,0.739216}%
\pgfsetfillcolor{currentfill}%
\pgfsetlinewidth{0.000000pt}%
\definecolor{currentstroke}{rgb}{0.000000,0.000000,0.000000}%
\pgfsetstrokecolor{currentstroke}%
\pgfsetstrokeopacity{0.000000}%
\pgfsetdash{}{0pt}%
\pgfpathmoveto{\pgfqpoint{3.554599in}{0.784569in}}%
\pgfpathlineto{\pgfqpoint{3.680411in}{0.784569in}}%
\pgfpathlineto{\pgfqpoint{3.680411in}{0.876969in}}%
\pgfpathlineto{\pgfqpoint{3.554599in}{0.876969in}}%
\pgfpathlineto{\pgfqpoint{3.554599in}{0.784569in}}%
\pgfpathclose%
\pgfusepath{fill}%
\end{pgfscope}%
\begin{pgfscope}%
\pgfpathrectangle{\pgfqpoint{3.319565in}{0.488889in}}{\pgfqpoint{0.360845in}{0.924000in}}%
\pgfusepath{clip}%
\pgfsetbuttcap%
\pgfsetmiterjoin%
\definecolor{currentfill}{rgb}{0.000000,0.521569,0.739216}%
\pgfsetfillcolor{currentfill}%
\pgfsetlinewidth{0.000000pt}%
\definecolor{currentstroke}{rgb}{0.000000,0.000000,0.000000}%
\pgfsetstrokecolor{currentstroke}%
\pgfsetstrokeopacity{0.000000}%
\pgfsetdash{}{0pt}%
\pgfpathmoveto{\pgfqpoint{3.554552in}{0.793809in}}%
\pgfpathlineto{\pgfqpoint{3.680411in}{0.793809in}}%
\pgfpathlineto{\pgfqpoint{3.680411in}{0.886209in}}%
\pgfpathlineto{\pgfqpoint{3.554552in}{0.886209in}}%
\pgfpathlineto{\pgfqpoint{3.554552in}{0.793809in}}%
\pgfpathclose%
\pgfusepath{fill}%
\end{pgfscope}%
\begin{pgfscope}%
\pgfpathrectangle{\pgfqpoint{3.319565in}{0.488889in}}{\pgfqpoint{0.360845in}{0.924000in}}%
\pgfusepath{clip}%
\pgfsetbuttcap%
\pgfsetmiterjoin%
\definecolor{currentfill}{rgb}{0.000000,0.521569,0.739216}%
\pgfsetfillcolor{currentfill}%
\pgfsetlinewidth{0.000000pt}%
\definecolor{currentstroke}{rgb}{0.000000,0.000000,0.000000}%
\pgfsetstrokecolor{currentstroke}%
\pgfsetstrokeopacity{0.000000}%
\pgfsetdash{}{0pt}%
\pgfpathmoveto{\pgfqpoint{3.554552in}{0.803049in}}%
\pgfpathlineto{\pgfqpoint{3.680411in}{0.803049in}}%
\pgfpathlineto{\pgfqpoint{3.680411in}{0.895449in}}%
\pgfpathlineto{\pgfqpoint{3.554552in}{0.895449in}}%
\pgfpathlineto{\pgfqpoint{3.554552in}{0.803049in}}%
\pgfpathclose%
\pgfusepath{fill}%
\end{pgfscope}%
\begin{pgfscope}%
\pgfpathrectangle{\pgfqpoint{3.319565in}{0.488889in}}{\pgfqpoint{0.360845in}{0.924000in}}%
\pgfusepath{clip}%
\pgfsetbuttcap%
\pgfsetmiterjoin%
\definecolor{currentfill}{rgb}{0.000000,0.521569,0.739216}%
\pgfsetfillcolor{currentfill}%
\pgfsetlinewidth{0.000000pt}%
\definecolor{currentstroke}{rgb}{0.000000,0.000000,0.000000}%
\pgfsetstrokecolor{currentstroke}%
\pgfsetstrokeopacity{0.000000}%
\pgfsetdash{}{0pt}%
\pgfpathmoveto{\pgfqpoint{3.554552in}{0.812289in}}%
\pgfpathlineto{\pgfqpoint{3.680411in}{0.812289in}}%
\pgfpathlineto{\pgfqpoint{3.680411in}{0.904689in}}%
\pgfpathlineto{\pgfqpoint{3.554552in}{0.904689in}}%
\pgfpathlineto{\pgfqpoint{3.554552in}{0.812289in}}%
\pgfpathclose%
\pgfusepath{fill}%
\end{pgfscope}%
\begin{pgfscope}%
\pgfpathrectangle{\pgfqpoint{3.319565in}{0.488889in}}{\pgfqpoint{0.360845in}{0.924000in}}%
\pgfusepath{clip}%
\pgfsetbuttcap%
\pgfsetmiterjoin%
\definecolor{currentfill}{rgb}{0.000000,0.521569,0.739216}%
\pgfsetfillcolor{currentfill}%
\pgfsetlinewidth{0.000000pt}%
\definecolor{currentstroke}{rgb}{0.000000,0.000000,0.000000}%
\pgfsetstrokecolor{currentstroke}%
\pgfsetstrokeopacity{0.000000}%
\pgfsetdash{}{0pt}%
\pgfpathmoveto{\pgfqpoint{3.554552in}{0.821529in}}%
\pgfpathlineto{\pgfqpoint{3.680411in}{0.821529in}}%
\pgfpathlineto{\pgfqpoint{3.680411in}{0.913929in}}%
\pgfpathlineto{\pgfqpoint{3.554552in}{0.913929in}}%
\pgfpathlineto{\pgfqpoint{3.554552in}{0.821529in}}%
\pgfpathclose%
\pgfusepath{fill}%
\end{pgfscope}%
\begin{pgfscope}%
\pgfpathrectangle{\pgfqpoint{3.319565in}{0.488889in}}{\pgfqpoint{0.360845in}{0.924000in}}%
\pgfusepath{clip}%
\pgfsetbuttcap%
\pgfsetmiterjoin%
\definecolor{currentfill}{rgb}{0.000000,0.521569,0.739216}%
\pgfsetfillcolor{currentfill}%
\pgfsetlinewidth{0.000000pt}%
\definecolor{currentstroke}{rgb}{0.000000,0.000000,0.000000}%
\pgfsetstrokecolor{currentstroke}%
\pgfsetstrokeopacity{0.000000}%
\pgfsetdash{}{0pt}%
\pgfpathmoveto{\pgfqpoint{3.554552in}{0.830769in}}%
\pgfpathlineto{\pgfqpoint{3.680411in}{0.830769in}}%
\pgfpathlineto{\pgfqpoint{3.680411in}{0.923169in}}%
\pgfpathlineto{\pgfqpoint{3.554552in}{0.923169in}}%
\pgfpathlineto{\pgfqpoint{3.554552in}{0.830769in}}%
\pgfpathclose%
\pgfusepath{fill}%
\end{pgfscope}%
\begin{pgfscope}%
\pgfpathrectangle{\pgfqpoint{3.319565in}{0.488889in}}{\pgfqpoint{0.360845in}{0.924000in}}%
\pgfusepath{clip}%
\pgfsetbuttcap%
\pgfsetmiterjoin%
\definecolor{currentfill}{rgb}{0.000000,0.517647,0.741176}%
\pgfsetfillcolor{currentfill}%
\pgfsetlinewidth{0.000000pt}%
\definecolor{currentstroke}{rgb}{0.000000,0.000000,0.000000}%
\pgfsetstrokecolor{currentstroke}%
\pgfsetstrokeopacity{0.000000}%
\pgfsetdash{}{0pt}%
\pgfpathmoveto{\pgfqpoint{3.554430in}{0.840009in}}%
\pgfpathlineto{\pgfqpoint{3.680411in}{0.840009in}}%
\pgfpathlineto{\pgfqpoint{3.680411in}{0.932409in}}%
\pgfpathlineto{\pgfqpoint{3.554430in}{0.932409in}}%
\pgfpathlineto{\pgfqpoint{3.554430in}{0.840009in}}%
\pgfpathclose%
\pgfusepath{fill}%
\end{pgfscope}%
\begin{pgfscope}%
\pgfpathrectangle{\pgfqpoint{3.319565in}{0.488889in}}{\pgfqpoint{0.360845in}{0.924000in}}%
\pgfusepath{clip}%
\pgfsetbuttcap%
\pgfsetmiterjoin%
\definecolor{currentfill}{rgb}{0.000000,0.505882,0.747059}%
\pgfsetfillcolor{currentfill}%
\pgfsetlinewidth{0.000000pt}%
\definecolor{currentstroke}{rgb}{0.000000,0.000000,0.000000}%
\pgfsetstrokecolor{currentstroke}%
\pgfsetstrokeopacity{0.000000}%
\pgfsetdash{}{0pt}%
\pgfpathmoveto{\pgfqpoint{3.553331in}{0.849249in}}%
\pgfpathlineto{\pgfqpoint{3.680411in}{0.849249in}}%
\pgfpathlineto{\pgfqpoint{3.680411in}{0.941649in}}%
\pgfpathlineto{\pgfqpoint{3.553331in}{0.941649in}}%
\pgfpathlineto{\pgfqpoint{3.553331in}{0.849249in}}%
\pgfpathclose%
\pgfusepath{fill}%
\end{pgfscope}%
\begin{pgfscope}%
\pgfpathrectangle{\pgfqpoint{3.319565in}{0.488889in}}{\pgfqpoint{0.360845in}{0.924000in}}%
\pgfusepath{clip}%
\pgfsetbuttcap%
\pgfsetmiterjoin%
\definecolor{currentfill}{rgb}{0.000000,0.494118,0.752941}%
\pgfsetfillcolor{currentfill}%
\pgfsetlinewidth{0.000000pt}%
\definecolor{currentstroke}{rgb}{0.000000,0.000000,0.000000}%
\pgfsetstrokecolor{currentstroke}%
\pgfsetstrokeopacity{0.000000}%
\pgfsetdash{}{0pt}%
\pgfpathmoveto{\pgfqpoint{3.552507in}{0.858489in}}%
\pgfpathlineto{\pgfqpoint{3.680411in}{0.858489in}}%
\pgfpathlineto{\pgfqpoint{3.680411in}{0.950889in}}%
\pgfpathlineto{\pgfqpoint{3.552507in}{0.950889in}}%
\pgfpathlineto{\pgfqpoint{3.552507in}{0.858489in}}%
\pgfpathclose%
\pgfusepath{fill}%
\end{pgfscope}%
\begin{pgfscope}%
\pgfpathrectangle{\pgfqpoint{3.319565in}{0.488889in}}{\pgfqpoint{0.360845in}{0.924000in}}%
\pgfusepath{clip}%
\pgfsetbuttcap%
\pgfsetmiterjoin%
\definecolor{currentfill}{rgb}{0.000000,0.490196,0.754902}%
\pgfsetfillcolor{currentfill}%
\pgfsetlinewidth{0.000000pt}%
\definecolor{currentstroke}{rgb}{0.000000,0.000000,0.000000}%
\pgfsetstrokecolor{currentstroke}%
\pgfsetstrokeopacity{0.000000}%
\pgfsetdash{}{0pt}%
\pgfpathmoveto{\pgfqpoint{3.552233in}{0.867729in}}%
\pgfpathlineto{\pgfqpoint{3.680411in}{0.867729in}}%
\pgfpathlineto{\pgfqpoint{3.680411in}{0.960129in}}%
\pgfpathlineto{\pgfqpoint{3.552233in}{0.960129in}}%
\pgfpathlineto{\pgfqpoint{3.552233in}{0.867729in}}%
\pgfpathclose%
\pgfusepath{fill}%
\end{pgfscope}%
\begin{pgfscope}%
\pgfpathrectangle{\pgfqpoint{3.319565in}{0.488889in}}{\pgfqpoint{0.360845in}{0.924000in}}%
\pgfusepath{clip}%
\pgfsetbuttcap%
\pgfsetmiterjoin%
\definecolor{currentfill}{rgb}{0.000000,0.482353,0.758824}%
\pgfsetfillcolor{currentfill}%
\pgfsetlinewidth{0.000000pt}%
\definecolor{currentstroke}{rgb}{0.000000,0.000000,0.000000}%
\pgfsetstrokecolor{currentstroke}%
\pgfsetstrokeopacity{0.000000}%
\pgfsetdash{}{0pt}%
\pgfpathmoveto{\pgfqpoint{3.551600in}{0.876969in}}%
\pgfpathlineto{\pgfqpoint{3.680411in}{0.876969in}}%
\pgfpathlineto{\pgfqpoint{3.680411in}{0.969369in}}%
\pgfpathlineto{\pgfqpoint{3.551600in}{0.969369in}}%
\pgfpathlineto{\pgfqpoint{3.551600in}{0.876969in}}%
\pgfpathclose%
\pgfusepath{fill}%
\end{pgfscope}%
\begin{pgfscope}%
\pgfpathrectangle{\pgfqpoint{3.319565in}{0.488889in}}{\pgfqpoint{0.360845in}{0.924000in}}%
\pgfusepath{clip}%
\pgfsetbuttcap%
\pgfsetmiterjoin%
\definecolor{currentfill}{rgb}{0.000000,0.482353,0.758824}%
\pgfsetfillcolor{currentfill}%
\pgfsetlinewidth{0.000000pt}%
\definecolor{currentstroke}{rgb}{0.000000,0.000000,0.000000}%
\pgfsetstrokecolor{currentstroke}%
\pgfsetstrokeopacity{0.000000}%
\pgfsetdash{}{0pt}%
\pgfpathmoveto{\pgfqpoint{3.551392in}{0.886209in}}%
\pgfpathlineto{\pgfqpoint{3.680411in}{0.886209in}}%
\pgfpathlineto{\pgfqpoint{3.680411in}{0.978609in}}%
\pgfpathlineto{\pgfqpoint{3.551392in}{0.978609in}}%
\pgfpathlineto{\pgfqpoint{3.551392in}{0.886209in}}%
\pgfpathclose%
\pgfusepath{fill}%
\end{pgfscope}%
\begin{pgfscope}%
\pgfpathrectangle{\pgfqpoint{3.319565in}{0.488889in}}{\pgfqpoint{0.360845in}{0.924000in}}%
\pgfusepath{clip}%
\pgfsetbuttcap%
\pgfsetmiterjoin%
\definecolor{currentfill}{rgb}{0.000000,0.482353,0.758824}%
\pgfsetfillcolor{currentfill}%
\pgfsetlinewidth{0.000000pt}%
\definecolor{currentstroke}{rgb}{0.000000,0.000000,0.000000}%
\pgfsetstrokecolor{currentstroke}%
\pgfsetstrokeopacity{0.000000}%
\pgfsetdash{}{0pt}%
\pgfpathmoveto{\pgfqpoint{3.551392in}{0.895449in}}%
\pgfpathlineto{\pgfqpoint{3.680411in}{0.895449in}}%
\pgfpathlineto{\pgfqpoint{3.680411in}{0.987849in}}%
\pgfpathlineto{\pgfqpoint{3.551392in}{0.987849in}}%
\pgfpathlineto{\pgfqpoint{3.551392in}{0.895449in}}%
\pgfpathclose%
\pgfusepath{fill}%
\end{pgfscope}%
\begin{pgfscope}%
\pgfpathrectangle{\pgfqpoint{3.319565in}{0.488889in}}{\pgfqpoint{0.360845in}{0.924000in}}%
\pgfusepath{clip}%
\pgfsetbuttcap%
\pgfsetmiterjoin%
\definecolor{currentfill}{rgb}{0.000000,0.482353,0.758824}%
\pgfsetfillcolor{currentfill}%
\pgfsetlinewidth{0.000000pt}%
\definecolor{currentstroke}{rgb}{0.000000,0.000000,0.000000}%
\pgfsetstrokecolor{currentstroke}%
\pgfsetstrokeopacity{0.000000}%
\pgfsetdash{}{0pt}%
\pgfpathmoveto{\pgfqpoint{3.551392in}{0.904689in}}%
\pgfpathlineto{\pgfqpoint{3.680411in}{0.904689in}}%
\pgfpathlineto{\pgfqpoint{3.680411in}{0.997089in}}%
\pgfpathlineto{\pgfqpoint{3.551392in}{0.997089in}}%
\pgfpathlineto{\pgfqpoint{3.551392in}{0.904689in}}%
\pgfpathclose%
\pgfusepath{fill}%
\end{pgfscope}%
\begin{pgfscope}%
\pgfpathrectangle{\pgfqpoint{3.319565in}{0.488889in}}{\pgfqpoint{0.360845in}{0.924000in}}%
\pgfusepath{clip}%
\pgfsetbuttcap%
\pgfsetmiterjoin%
\definecolor{currentfill}{rgb}{0.000000,0.482353,0.758824}%
\pgfsetfillcolor{currentfill}%
\pgfsetlinewidth{0.000000pt}%
\definecolor{currentstroke}{rgb}{0.000000,0.000000,0.000000}%
\pgfsetstrokecolor{currentstroke}%
\pgfsetstrokeopacity{0.000000}%
\pgfsetdash{}{0pt}%
\pgfpathmoveto{\pgfqpoint{3.551392in}{0.913929in}}%
\pgfpathlineto{\pgfqpoint{3.680411in}{0.913929in}}%
\pgfpathlineto{\pgfqpoint{3.680411in}{1.006329in}}%
\pgfpathlineto{\pgfqpoint{3.551392in}{1.006329in}}%
\pgfpathlineto{\pgfqpoint{3.551392in}{0.913929in}}%
\pgfpathclose%
\pgfusepath{fill}%
\end{pgfscope}%
\begin{pgfscope}%
\pgfpathrectangle{\pgfqpoint{3.319565in}{0.488889in}}{\pgfqpoint{0.360845in}{0.924000in}}%
\pgfusepath{clip}%
\pgfsetbuttcap%
\pgfsetmiterjoin%
\definecolor{currentfill}{rgb}{0.000000,0.478431,0.760784}%
\pgfsetfillcolor{currentfill}%
\pgfsetlinewidth{0.000000pt}%
\definecolor{currentstroke}{rgb}{0.000000,0.000000,0.000000}%
\pgfsetstrokecolor{currentstroke}%
\pgfsetstrokeopacity{0.000000}%
\pgfsetdash{}{0pt}%
\pgfpathmoveto{\pgfqpoint{3.551102in}{0.923169in}}%
\pgfpathlineto{\pgfqpoint{3.680411in}{0.923169in}}%
\pgfpathlineto{\pgfqpoint{3.680411in}{1.015569in}}%
\pgfpathlineto{\pgfqpoint{3.551102in}{1.015569in}}%
\pgfpathlineto{\pgfqpoint{3.551102in}{0.923169in}}%
\pgfpathclose%
\pgfusepath{fill}%
\end{pgfscope}%
\begin{pgfscope}%
\pgfpathrectangle{\pgfqpoint{3.319565in}{0.488889in}}{\pgfqpoint{0.360845in}{0.924000in}}%
\pgfusepath{clip}%
\pgfsetbuttcap%
\pgfsetmiterjoin%
\definecolor{currentfill}{rgb}{0.000000,0.474510,0.762745}%
\pgfsetfillcolor{currentfill}%
\pgfsetlinewidth{0.000000pt}%
\definecolor{currentstroke}{rgb}{0.000000,0.000000,0.000000}%
\pgfsetstrokecolor{currentstroke}%
\pgfsetstrokeopacity{0.000000}%
\pgfsetdash{}{0pt}%
\pgfpathmoveto{\pgfqpoint{3.550994in}{0.932409in}}%
\pgfpathlineto{\pgfqpoint{3.680411in}{0.932409in}}%
\pgfpathlineto{\pgfqpoint{3.680411in}{1.024809in}}%
\pgfpathlineto{\pgfqpoint{3.550994in}{1.024809in}}%
\pgfpathlineto{\pgfqpoint{3.550994in}{0.932409in}}%
\pgfpathclose%
\pgfusepath{fill}%
\end{pgfscope}%
\begin{pgfscope}%
\pgfpathrectangle{\pgfqpoint{3.319565in}{0.488889in}}{\pgfqpoint{0.360845in}{0.924000in}}%
\pgfusepath{clip}%
\pgfsetbuttcap%
\pgfsetmiterjoin%
\definecolor{currentfill}{rgb}{0.000000,0.466667,0.766667}%
\pgfsetfillcolor{currentfill}%
\pgfsetlinewidth{0.000000pt}%
\definecolor{currentstroke}{rgb}{0.000000,0.000000,0.000000}%
\pgfsetstrokecolor{currentstroke}%
\pgfsetstrokeopacity{0.000000}%
\pgfsetdash{}{0pt}%
\pgfpathmoveto{\pgfqpoint{3.550302in}{0.941649in}}%
\pgfpathlineto{\pgfqpoint{3.680411in}{0.941649in}}%
\pgfpathlineto{\pgfqpoint{3.680411in}{1.034049in}}%
\pgfpathlineto{\pgfqpoint{3.550302in}{1.034049in}}%
\pgfpathlineto{\pgfqpoint{3.550302in}{0.941649in}}%
\pgfpathclose%
\pgfusepath{fill}%
\end{pgfscope}%
\begin{pgfscope}%
\pgfpathrectangle{\pgfqpoint{3.319565in}{0.488889in}}{\pgfqpoint{0.360845in}{0.924000in}}%
\pgfusepath{clip}%
\pgfsetbuttcap%
\pgfsetmiterjoin%
\definecolor{currentfill}{rgb}{0.000000,0.466667,0.766667}%
\pgfsetfillcolor{currentfill}%
\pgfsetlinewidth{0.000000pt}%
\definecolor{currentstroke}{rgb}{0.000000,0.000000,0.000000}%
\pgfsetstrokecolor{currentstroke}%
\pgfsetstrokeopacity{0.000000}%
\pgfsetdash{}{0pt}%
\pgfpathmoveto{\pgfqpoint{3.550168in}{0.950889in}}%
\pgfpathlineto{\pgfqpoint{3.680411in}{0.950889in}}%
\pgfpathlineto{\pgfqpoint{3.680411in}{1.043289in}}%
\pgfpathlineto{\pgfqpoint{3.550168in}{1.043289in}}%
\pgfpathlineto{\pgfqpoint{3.550168in}{0.950889in}}%
\pgfpathclose%
\pgfusepath{fill}%
\end{pgfscope}%
\begin{pgfscope}%
\pgfpathrectangle{\pgfqpoint{3.319565in}{0.488889in}}{\pgfqpoint{0.360845in}{0.924000in}}%
\pgfusepath{clip}%
\pgfsetbuttcap%
\pgfsetmiterjoin%
\definecolor{currentfill}{rgb}{0.000000,0.462745,0.768627}%
\pgfsetfillcolor{currentfill}%
\pgfsetlinewidth{0.000000pt}%
\definecolor{currentstroke}{rgb}{0.000000,0.000000,0.000000}%
\pgfsetstrokecolor{currentstroke}%
\pgfsetstrokeopacity{0.000000}%
\pgfsetdash{}{0pt}%
\pgfpathmoveto{\pgfqpoint{3.549899in}{0.960129in}}%
\pgfpathlineto{\pgfqpoint{3.680411in}{0.960129in}}%
\pgfpathlineto{\pgfqpoint{3.680411in}{1.052529in}}%
\pgfpathlineto{\pgfqpoint{3.549899in}{1.052529in}}%
\pgfpathlineto{\pgfqpoint{3.549899in}{0.960129in}}%
\pgfpathclose%
\pgfusepath{fill}%
\end{pgfscope}%
\begin{pgfscope}%
\pgfpathrectangle{\pgfqpoint{3.319565in}{0.488889in}}{\pgfqpoint{0.360845in}{0.924000in}}%
\pgfusepath{clip}%
\pgfsetbuttcap%
\pgfsetmiterjoin%
\definecolor{currentfill}{rgb}{0.000000,0.454902,0.772549}%
\pgfsetfillcolor{currentfill}%
\pgfsetlinewidth{0.000000pt}%
\definecolor{currentstroke}{rgb}{0.000000,0.000000,0.000000}%
\pgfsetstrokecolor{currentstroke}%
\pgfsetstrokeopacity{0.000000}%
\pgfsetdash{}{0pt}%
\pgfpathmoveto{\pgfqpoint{3.549135in}{0.969369in}}%
\pgfpathlineto{\pgfqpoint{3.680411in}{0.969369in}}%
\pgfpathlineto{\pgfqpoint{3.680411in}{1.061769in}}%
\pgfpathlineto{\pgfqpoint{3.549135in}{1.061769in}}%
\pgfpathlineto{\pgfqpoint{3.549135in}{0.969369in}}%
\pgfpathclose%
\pgfusepath{fill}%
\end{pgfscope}%
\begin{pgfscope}%
\pgfpathrectangle{\pgfqpoint{3.319565in}{0.488889in}}{\pgfqpoint{0.360845in}{0.924000in}}%
\pgfusepath{clip}%
\pgfsetbuttcap%
\pgfsetmiterjoin%
\definecolor{currentfill}{rgb}{0.000000,0.447059,0.776471}%
\pgfsetfillcolor{currentfill}%
\pgfsetlinewidth{0.000000pt}%
\definecolor{currentstroke}{rgb}{0.000000,0.000000,0.000000}%
\pgfsetstrokecolor{currentstroke}%
\pgfsetstrokeopacity{0.000000}%
\pgfsetdash{}{0pt}%
\pgfpathmoveto{\pgfqpoint{3.548594in}{0.978609in}}%
\pgfpathlineto{\pgfqpoint{3.680411in}{0.978609in}}%
\pgfpathlineto{\pgfqpoint{3.680411in}{1.071009in}}%
\pgfpathlineto{\pgfqpoint{3.548594in}{1.071009in}}%
\pgfpathlineto{\pgfqpoint{3.548594in}{0.978609in}}%
\pgfpathclose%
\pgfusepath{fill}%
\end{pgfscope}%
\begin{pgfscope}%
\pgfpathrectangle{\pgfqpoint{3.319565in}{0.488889in}}{\pgfqpoint{0.360845in}{0.924000in}}%
\pgfusepath{clip}%
\pgfsetbuttcap%
\pgfsetmiterjoin%
\definecolor{currentfill}{rgb}{0.000000,0.447059,0.776471}%
\pgfsetfillcolor{currentfill}%
\pgfsetlinewidth{0.000000pt}%
\definecolor{currentstroke}{rgb}{0.000000,0.000000,0.000000}%
\pgfsetstrokecolor{currentstroke}%
\pgfsetstrokeopacity{0.000000}%
\pgfsetdash{}{0pt}%
\pgfpathmoveto{\pgfqpoint{3.548591in}{0.987849in}}%
\pgfpathlineto{\pgfqpoint{3.680411in}{0.987849in}}%
\pgfpathlineto{\pgfqpoint{3.680411in}{1.080249in}}%
\pgfpathlineto{\pgfqpoint{3.548591in}{1.080249in}}%
\pgfpathlineto{\pgfqpoint{3.548591in}{0.987849in}}%
\pgfpathclose%
\pgfusepath{fill}%
\end{pgfscope}%
\begin{pgfscope}%
\pgfpathrectangle{\pgfqpoint{3.319565in}{0.488889in}}{\pgfqpoint{0.360845in}{0.924000in}}%
\pgfusepath{clip}%
\pgfsetbuttcap%
\pgfsetmiterjoin%
\definecolor{currentfill}{rgb}{0.000000,0.443137,0.778431}%
\pgfsetfillcolor{currentfill}%
\pgfsetlinewidth{0.000000pt}%
\definecolor{currentstroke}{rgb}{0.000000,0.000000,0.000000}%
\pgfsetstrokecolor{currentstroke}%
\pgfsetstrokeopacity{0.000000}%
\pgfsetdash{}{0pt}%
\pgfpathmoveto{\pgfqpoint{3.548193in}{0.997089in}}%
\pgfpathlineto{\pgfqpoint{3.680411in}{0.997089in}}%
\pgfpathlineto{\pgfqpoint{3.680411in}{1.089489in}}%
\pgfpathlineto{\pgfqpoint{3.548193in}{1.089489in}}%
\pgfpathlineto{\pgfqpoint{3.548193in}{0.997089in}}%
\pgfpathclose%
\pgfusepath{fill}%
\end{pgfscope}%
\begin{pgfscope}%
\pgfpathrectangle{\pgfqpoint{3.319565in}{0.488889in}}{\pgfqpoint{0.360845in}{0.924000in}}%
\pgfusepath{clip}%
\pgfsetbuttcap%
\pgfsetmiterjoin%
\definecolor{currentfill}{rgb}{0.000000,0.439216,0.780392}%
\pgfsetfillcolor{currentfill}%
\pgfsetlinewidth{0.000000pt}%
\definecolor{currentstroke}{rgb}{0.000000,0.000000,0.000000}%
\pgfsetstrokecolor{currentstroke}%
\pgfsetstrokeopacity{0.000000}%
\pgfsetdash{}{0pt}%
\pgfpathmoveto{\pgfqpoint{3.548087in}{1.006329in}}%
\pgfpathlineto{\pgfqpoint{3.680411in}{1.006329in}}%
\pgfpathlineto{\pgfqpoint{3.680411in}{1.098729in}}%
\pgfpathlineto{\pgfqpoint{3.548087in}{1.098729in}}%
\pgfpathlineto{\pgfqpoint{3.548087in}{1.006329in}}%
\pgfpathclose%
\pgfusepath{fill}%
\end{pgfscope}%
\begin{pgfscope}%
\pgfpathrectangle{\pgfqpoint{3.319565in}{0.488889in}}{\pgfqpoint{0.360845in}{0.924000in}}%
\pgfusepath{clip}%
\pgfsetbuttcap%
\pgfsetmiterjoin%
\definecolor{currentfill}{rgb}{0.000000,0.431373,0.784314}%
\pgfsetfillcolor{currentfill}%
\pgfsetlinewidth{0.000000pt}%
\definecolor{currentstroke}{rgb}{0.000000,0.000000,0.000000}%
\pgfsetstrokecolor{currentstroke}%
\pgfsetstrokeopacity{0.000000}%
\pgfsetdash{}{0pt}%
\pgfpathmoveto{\pgfqpoint{3.547398in}{1.015569in}}%
\pgfpathlineto{\pgfqpoint{3.680411in}{1.015569in}}%
\pgfpathlineto{\pgfqpoint{3.680411in}{1.107969in}}%
\pgfpathlineto{\pgfqpoint{3.547398in}{1.107969in}}%
\pgfpathlineto{\pgfqpoint{3.547398in}{1.015569in}}%
\pgfpathclose%
\pgfusepath{fill}%
\end{pgfscope}%
\begin{pgfscope}%
\pgfpathrectangle{\pgfqpoint{3.319565in}{0.488889in}}{\pgfqpoint{0.360845in}{0.924000in}}%
\pgfusepath{clip}%
\pgfsetbuttcap%
\pgfsetmiterjoin%
\definecolor{currentfill}{rgb}{0.000000,0.431373,0.784314}%
\pgfsetfillcolor{currentfill}%
\pgfsetlinewidth{0.000000pt}%
\definecolor{currentstroke}{rgb}{0.000000,0.000000,0.000000}%
\pgfsetstrokecolor{currentstroke}%
\pgfsetstrokeopacity{0.000000}%
\pgfsetdash{}{0pt}%
\pgfpathmoveto{\pgfqpoint{3.547398in}{1.024809in}}%
\pgfpathlineto{\pgfqpoint{3.680411in}{1.024809in}}%
\pgfpathlineto{\pgfqpoint{3.680411in}{1.117209in}}%
\pgfpathlineto{\pgfqpoint{3.547398in}{1.117209in}}%
\pgfpathlineto{\pgfqpoint{3.547398in}{1.024809in}}%
\pgfpathclose%
\pgfusepath{fill}%
\end{pgfscope}%
\begin{pgfscope}%
\pgfpathrectangle{\pgfqpoint{3.319565in}{0.488889in}}{\pgfqpoint{0.360845in}{0.924000in}}%
\pgfusepath{clip}%
\pgfsetbuttcap%
\pgfsetmiterjoin%
\definecolor{currentfill}{rgb}{0.000000,0.431373,0.784314}%
\pgfsetfillcolor{currentfill}%
\pgfsetlinewidth{0.000000pt}%
\definecolor{currentstroke}{rgb}{0.000000,0.000000,0.000000}%
\pgfsetstrokecolor{currentstroke}%
\pgfsetstrokeopacity{0.000000}%
\pgfsetdash{}{0pt}%
\pgfpathmoveto{\pgfqpoint{3.547398in}{1.034049in}}%
\pgfpathlineto{\pgfqpoint{3.680411in}{1.034049in}}%
\pgfpathlineto{\pgfqpoint{3.680411in}{1.126449in}}%
\pgfpathlineto{\pgfqpoint{3.547398in}{1.126449in}}%
\pgfpathlineto{\pgfqpoint{3.547398in}{1.034049in}}%
\pgfpathclose%
\pgfusepath{fill}%
\end{pgfscope}%
\begin{pgfscope}%
\pgfpathrectangle{\pgfqpoint{3.319565in}{0.488889in}}{\pgfqpoint{0.360845in}{0.924000in}}%
\pgfusepath{clip}%
\pgfsetbuttcap%
\pgfsetmiterjoin%
\definecolor{currentfill}{rgb}{0.000000,0.427451,0.786275}%
\pgfsetfillcolor{currentfill}%
\pgfsetlinewidth{0.000000pt}%
\definecolor{currentstroke}{rgb}{0.000000,0.000000,0.000000}%
\pgfsetstrokecolor{currentstroke}%
\pgfsetstrokeopacity{0.000000}%
\pgfsetdash{}{0pt}%
\pgfpathmoveto{\pgfqpoint{3.547013in}{1.043289in}}%
\pgfpathlineto{\pgfqpoint{3.680411in}{1.043289in}}%
\pgfpathlineto{\pgfqpoint{3.680411in}{1.135689in}}%
\pgfpathlineto{\pgfqpoint{3.547013in}{1.135689in}}%
\pgfpathlineto{\pgfqpoint{3.547013in}{1.043289in}}%
\pgfpathclose%
\pgfusepath{fill}%
\end{pgfscope}%
\begin{pgfscope}%
\pgfpathrectangle{\pgfqpoint{3.319565in}{0.488889in}}{\pgfqpoint{0.360845in}{0.924000in}}%
\pgfusepath{clip}%
\pgfsetbuttcap%
\pgfsetmiterjoin%
\definecolor{currentfill}{rgb}{0.000000,0.427451,0.786275}%
\pgfsetfillcolor{currentfill}%
\pgfsetlinewidth{0.000000pt}%
\definecolor{currentstroke}{rgb}{0.000000,0.000000,0.000000}%
\pgfsetstrokecolor{currentstroke}%
\pgfsetstrokeopacity{0.000000}%
\pgfsetdash{}{0pt}%
\pgfpathmoveto{\pgfqpoint{3.546966in}{1.052529in}}%
\pgfpathlineto{\pgfqpoint{3.680411in}{1.052529in}}%
\pgfpathlineto{\pgfqpoint{3.680411in}{1.144929in}}%
\pgfpathlineto{\pgfqpoint{3.546966in}{1.144929in}}%
\pgfpathlineto{\pgfqpoint{3.546966in}{1.052529in}}%
\pgfpathclose%
\pgfusepath{fill}%
\end{pgfscope}%
\begin{pgfscope}%
\pgfpathrectangle{\pgfqpoint{3.319565in}{0.488889in}}{\pgfqpoint{0.360845in}{0.924000in}}%
\pgfusepath{clip}%
\pgfsetbuttcap%
\pgfsetmiterjoin%
\definecolor{currentfill}{rgb}{0.000000,0.423529,0.788235}%
\pgfsetfillcolor{currentfill}%
\pgfsetlinewidth{0.000000pt}%
\definecolor{currentstroke}{rgb}{0.000000,0.000000,0.000000}%
\pgfsetstrokecolor{currentstroke}%
\pgfsetstrokeopacity{0.000000}%
\pgfsetdash{}{0pt}%
\pgfpathmoveto{\pgfqpoint{3.546797in}{1.061769in}}%
\pgfpathlineto{\pgfqpoint{3.680411in}{1.061769in}}%
\pgfpathlineto{\pgfqpoint{3.680411in}{1.154169in}}%
\pgfpathlineto{\pgfqpoint{3.546797in}{1.154169in}}%
\pgfpathlineto{\pgfqpoint{3.546797in}{1.061769in}}%
\pgfpathclose%
\pgfusepath{fill}%
\end{pgfscope}%
\begin{pgfscope}%
\pgfpathrectangle{\pgfqpoint{3.319565in}{0.488889in}}{\pgfqpoint{0.360845in}{0.924000in}}%
\pgfusepath{clip}%
\pgfsetbuttcap%
\pgfsetmiterjoin%
\definecolor{currentfill}{rgb}{0.000000,0.423529,0.788235}%
\pgfsetfillcolor{currentfill}%
\pgfsetlinewidth{0.000000pt}%
\definecolor{currentstroke}{rgb}{0.000000,0.000000,0.000000}%
\pgfsetstrokecolor{currentstroke}%
\pgfsetstrokeopacity{0.000000}%
\pgfsetdash{}{0pt}%
\pgfpathmoveto{\pgfqpoint{3.546797in}{1.071009in}}%
\pgfpathlineto{\pgfqpoint{3.680411in}{1.071009in}}%
\pgfpathlineto{\pgfqpoint{3.680411in}{1.163409in}}%
\pgfpathlineto{\pgfqpoint{3.546797in}{1.163409in}}%
\pgfpathlineto{\pgfqpoint{3.546797in}{1.071009in}}%
\pgfpathclose%
\pgfusepath{fill}%
\end{pgfscope}%
\begin{pgfscope}%
\pgfpathrectangle{\pgfqpoint{3.319565in}{0.488889in}}{\pgfqpoint{0.360845in}{0.924000in}}%
\pgfusepath{clip}%
\pgfsetbuttcap%
\pgfsetmiterjoin%
\definecolor{currentfill}{rgb}{0.000000,0.419608,0.790196}%
\pgfsetfillcolor{currentfill}%
\pgfsetlinewidth{0.000000pt}%
\definecolor{currentstroke}{rgb}{0.000000,0.000000,0.000000}%
\pgfsetstrokecolor{currentstroke}%
\pgfsetstrokeopacity{0.000000}%
\pgfsetdash{}{0pt}%
\pgfpathmoveto{\pgfqpoint{3.546460in}{1.080249in}}%
\pgfpathlineto{\pgfqpoint{3.680411in}{1.080249in}}%
\pgfpathlineto{\pgfqpoint{3.680411in}{1.172649in}}%
\pgfpathlineto{\pgfqpoint{3.546460in}{1.172649in}}%
\pgfpathlineto{\pgfqpoint{3.546460in}{1.080249in}}%
\pgfpathclose%
\pgfusepath{fill}%
\end{pgfscope}%
\begin{pgfscope}%
\pgfpathrectangle{\pgfqpoint{3.319565in}{0.488889in}}{\pgfqpoint{0.360845in}{0.924000in}}%
\pgfusepath{clip}%
\pgfsetbuttcap%
\pgfsetmiterjoin%
\definecolor{currentfill}{rgb}{0.000000,0.419608,0.790196}%
\pgfsetfillcolor{currentfill}%
\pgfsetlinewidth{0.000000pt}%
\definecolor{currentstroke}{rgb}{0.000000,0.000000,0.000000}%
\pgfsetstrokecolor{currentstroke}%
\pgfsetstrokeopacity{0.000000}%
\pgfsetdash{}{0pt}%
\pgfpathmoveto{\pgfqpoint{3.546402in}{1.089489in}}%
\pgfpathlineto{\pgfqpoint{3.680411in}{1.089489in}}%
\pgfpathlineto{\pgfqpoint{3.680411in}{1.181889in}}%
\pgfpathlineto{\pgfqpoint{3.546402in}{1.181889in}}%
\pgfpathlineto{\pgfqpoint{3.546402in}{1.089489in}}%
\pgfpathclose%
\pgfusepath{fill}%
\end{pgfscope}%
\begin{pgfscope}%
\pgfpathrectangle{\pgfqpoint{3.319565in}{0.488889in}}{\pgfqpoint{0.360845in}{0.924000in}}%
\pgfusepath{clip}%
\pgfsetbuttcap%
\pgfsetmiterjoin%
\definecolor{currentfill}{rgb}{0.000000,0.411765,0.794118}%
\pgfsetfillcolor{currentfill}%
\pgfsetlinewidth{0.000000pt}%
\definecolor{currentstroke}{rgb}{0.000000,0.000000,0.000000}%
\pgfsetstrokecolor{currentstroke}%
\pgfsetstrokeopacity{0.000000}%
\pgfsetdash{}{0pt}%
\pgfpathmoveto{\pgfqpoint{3.545929in}{1.098729in}}%
\pgfpathlineto{\pgfqpoint{3.680411in}{1.098729in}}%
\pgfpathlineto{\pgfqpoint{3.680411in}{1.191129in}}%
\pgfpathlineto{\pgfqpoint{3.545929in}{1.191129in}}%
\pgfpathlineto{\pgfqpoint{3.545929in}{1.098729in}}%
\pgfpathclose%
\pgfusepath{fill}%
\end{pgfscope}%
\begin{pgfscope}%
\pgfpathrectangle{\pgfqpoint{3.319565in}{0.488889in}}{\pgfqpoint{0.360845in}{0.924000in}}%
\pgfusepath{clip}%
\pgfsetbuttcap%
\pgfsetmiterjoin%
\definecolor{currentfill}{rgb}{0.000000,0.411765,0.794118}%
\pgfsetfillcolor{currentfill}%
\pgfsetlinewidth{0.000000pt}%
\definecolor{currentstroke}{rgb}{0.000000,0.000000,0.000000}%
\pgfsetstrokecolor{currentstroke}%
\pgfsetstrokeopacity{0.000000}%
\pgfsetdash{}{0pt}%
\pgfpathmoveto{\pgfqpoint{3.545690in}{1.107969in}}%
\pgfpathlineto{\pgfqpoint{3.680411in}{1.107969in}}%
\pgfpathlineto{\pgfqpoint{3.680411in}{1.200369in}}%
\pgfpathlineto{\pgfqpoint{3.545690in}{1.200369in}}%
\pgfpathlineto{\pgfqpoint{3.545690in}{1.107969in}}%
\pgfpathclose%
\pgfusepath{fill}%
\end{pgfscope}%
\begin{pgfscope}%
\pgfpathrectangle{\pgfqpoint{3.319565in}{0.488889in}}{\pgfqpoint{0.360845in}{0.924000in}}%
\pgfusepath{clip}%
\pgfsetbuttcap%
\pgfsetmiterjoin%
\definecolor{currentfill}{rgb}{0.000000,0.411765,0.794118}%
\pgfsetfillcolor{currentfill}%
\pgfsetlinewidth{0.000000pt}%
\definecolor{currentstroke}{rgb}{0.000000,0.000000,0.000000}%
\pgfsetstrokecolor{currentstroke}%
\pgfsetstrokeopacity{0.000000}%
\pgfsetdash{}{0pt}%
\pgfpathmoveto{\pgfqpoint{3.545690in}{1.117209in}}%
\pgfpathlineto{\pgfqpoint{3.680411in}{1.117209in}}%
\pgfpathlineto{\pgfqpoint{3.680411in}{1.209609in}}%
\pgfpathlineto{\pgfqpoint{3.545690in}{1.209609in}}%
\pgfpathlineto{\pgfqpoint{3.545690in}{1.117209in}}%
\pgfpathclose%
\pgfusepath{fill}%
\end{pgfscope}%
\begin{pgfscope}%
\pgfpathrectangle{\pgfqpoint{3.319565in}{0.488889in}}{\pgfqpoint{0.360845in}{0.924000in}}%
\pgfusepath{clip}%
\pgfsetbuttcap%
\pgfsetmiterjoin%
\definecolor{currentfill}{rgb}{0.000000,0.403922,0.798039}%
\pgfsetfillcolor{currentfill}%
\pgfsetlinewidth{0.000000pt}%
\definecolor{currentstroke}{rgb}{0.000000,0.000000,0.000000}%
\pgfsetstrokecolor{currentstroke}%
\pgfsetstrokeopacity{0.000000}%
\pgfsetdash{}{0pt}%
\pgfpathmoveto{\pgfqpoint{3.545305in}{1.126449in}}%
\pgfpathlineto{\pgfqpoint{3.680411in}{1.126449in}}%
\pgfpathlineto{\pgfqpoint{3.680411in}{1.218849in}}%
\pgfpathlineto{\pgfqpoint{3.545305in}{1.218849in}}%
\pgfpathlineto{\pgfqpoint{3.545305in}{1.126449in}}%
\pgfpathclose%
\pgfusepath{fill}%
\end{pgfscope}%
\begin{pgfscope}%
\pgfpathrectangle{\pgfqpoint{3.319565in}{0.488889in}}{\pgfqpoint{0.360845in}{0.924000in}}%
\pgfusepath{clip}%
\pgfsetbuttcap%
\pgfsetmiterjoin%
\definecolor{currentfill}{rgb}{0.000000,0.392157,0.803922}%
\pgfsetfillcolor{currentfill}%
\pgfsetlinewidth{0.000000pt}%
\definecolor{currentstroke}{rgb}{0.000000,0.000000,0.000000}%
\pgfsetstrokecolor{currentstroke}%
\pgfsetstrokeopacity{0.000000}%
\pgfsetdash{}{0pt}%
\pgfpathmoveto{\pgfqpoint{3.544072in}{1.135689in}}%
\pgfpathlineto{\pgfqpoint{3.680411in}{1.135689in}}%
\pgfpathlineto{\pgfqpoint{3.680411in}{1.228089in}}%
\pgfpathlineto{\pgfqpoint{3.544072in}{1.228089in}}%
\pgfpathlineto{\pgfqpoint{3.544072in}{1.135689in}}%
\pgfpathclose%
\pgfusepath{fill}%
\end{pgfscope}%
\begin{pgfscope}%
\pgfpathrectangle{\pgfqpoint{3.319565in}{0.488889in}}{\pgfqpoint{0.360845in}{0.924000in}}%
\pgfusepath{clip}%
\pgfsetbuttcap%
\pgfsetmiterjoin%
\definecolor{currentfill}{rgb}{0.000000,0.392157,0.803922}%
\pgfsetfillcolor{currentfill}%
\pgfsetlinewidth{0.000000pt}%
\definecolor{currentstroke}{rgb}{0.000000,0.000000,0.000000}%
\pgfsetstrokecolor{currentstroke}%
\pgfsetstrokeopacity{0.000000}%
\pgfsetdash{}{0pt}%
\pgfpathmoveto{\pgfqpoint{3.544072in}{1.144929in}}%
\pgfpathlineto{\pgfqpoint{3.680411in}{1.144929in}}%
\pgfpathlineto{\pgfqpoint{3.680411in}{1.237329in}}%
\pgfpathlineto{\pgfqpoint{3.544072in}{1.237329in}}%
\pgfpathlineto{\pgfqpoint{3.544072in}{1.144929in}}%
\pgfpathclose%
\pgfusepath{fill}%
\end{pgfscope}%
\begin{pgfscope}%
\pgfpathrectangle{\pgfqpoint{3.319565in}{0.488889in}}{\pgfqpoint{0.360845in}{0.924000in}}%
\pgfusepath{clip}%
\pgfsetbuttcap%
\pgfsetmiterjoin%
\definecolor{currentfill}{rgb}{0.000000,0.388235,0.805882}%
\pgfsetfillcolor{currentfill}%
\pgfsetlinewidth{0.000000pt}%
\definecolor{currentstroke}{rgb}{0.000000,0.000000,0.000000}%
\pgfsetstrokecolor{currentstroke}%
\pgfsetstrokeopacity{0.000000}%
\pgfsetdash{}{0pt}%
\pgfpathmoveto{\pgfqpoint{3.543791in}{1.154169in}}%
\pgfpathlineto{\pgfqpoint{3.680411in}{1.154169in}}%
\pgfpathlineto{\pgfqpoint{3.680411in}{1.246569in}}%
\pgfpathlineto{\pgfqpoint{3.543791in}{1.246569in}}%
\pgfpathlineto{\pgfqpoint{3.543791in}{1.154169in}}%
\pgfpathclose%
\pgfusepath{fill}%
\end{pgfscope}%
\begin{pgfscope}%
\pgfpathrectangle{\pgfqpoint{3.319565in}{0.488889in}}{\pgfqpoint{0.360845in}{0.924000in}}%
\pgfusepath{clip}%
\pgfsetbuttcap%
\pgfsetmiterjoin%
\definecolor{currentfill}{rgb}{0.000000,0.388235,0.805882}%
\pgfsetfillcolor{currentfill}%
\pgfsetlinewidth{0.000000pt}%
\definecolor{currentstroke}{rgb}{0.000000,0.000000,0.000000}%
\pgfsetstrokecolor{currentstroke}%
\pgfsetstrokeopacity{0.000000}%
\pgfsetdash{}{0pt}%
\pgfpathmoveto{\pgfqpoint{3.543791in}{1.163409in}}%
\pgfpathlineto{\pgfqpoint{3.680411in}{1.163409in}}%
\pgfpathlineto{\pgfqpoint{3.680411in}{1.255809in}}%
\pgfpathlineto{\pgfqpoint{3.543791in}{1.255809in}}%
\pgfpathlineto{\pgfqpoint{3.543791in}{1.163409in}}%
\pgfpathclose%
\pgfusepath{fill}%
\end{pgfscope}%
\begin{pgfscope}%
\pgfpathrectangle{\pgfqpoint{3.319565in}{0.488889in}}{\pgfqpoint{0.360845in}{0.924000in}}%
\pgfusepath{clip}%
\pgfsetbuttcap%
\pgfsetmiterjoin%
\definecolor{currentfill}{rgb}{0.000000,0.388235,0.805882}%
\pgfsetfillcolor{currentfill}%
\pgfsetlinewidth{0.000000pt}%
\definecolor{currentstroke}{rgb}{0.000000,0.000000,0.000000}%
\pgfsetstrokecolor{currentstroke}%
\pgfsetstrokeopacity{0.000000}%
\pgfsetdash{}{0pt}%
\pgfpathmoveto{\pgfqpoint{3.543791in}{1.172649in}}%
\pgfpathlineto{\pgfqpoint{3.680411in}{1.172649in}}%
\pgfpathlineto{\pgfqpoint{3.680411in}{1.265049in}}%
\pgfpathlineto{\pgfqpoint{3.543791in}{1.265049in}}%
\pgfpathlineto{\pgfqpoint{3.543791in}{1.172649in}}%
\pgfpathclose%
\pgfusepath{fill}%
\end{pgfscope}%
\begin{pgfscope}%
\pgfpathrectangle{\pgfqpoint{3.319565in}{0.488889in}}{\pgfqpoint{0.360845in}{0.924000in}}%
\pgfusepath{clip}%
\pgfsetbuttcap%
\pgfsetmiterjoin%
\definecolor{currentfill}{rgb}{0.000000,0.380392,0.809804}%
\pgfsetfillcolor{currentfill}%
\pgfsetlinewidth{0.000000pt}%
\definecolor{currentstroke}{rgb}{0.000000,0.000000,0.000000}%
\pgfsetstrokecolor{currentstroke}%
\pgfsetstrokeopacity{0.000000}%
\pgfsetdash{}{0pt}%
\pgfpathmoveto{\pgfqpoint{3.543213in}{1.181889in}}%
\pgfpathlineto{\pgfqpoint{3.680411in}{1.181889in}}%
\pgfpathlineto{\pgfqpoint{3.680411in}{1.274289in}}%
\pgfpathlineto{\pgfqpoint{3.543213in}{1.274289in}}%
\pgfpathlineto{\pgfqpoint{3.543213in}{1.181889in}}%
\pgfpathclose%
\pgfusepath{fill}%
\end{pgfscope}%
\begin{pgfscope}%
\pgfpathrectangle{\pgfqpoint{3.319565in}{0.488889in}}{\pgfqpoint{0.360845in}{0.924000in}}%
\pgfusepath{clip}%
\pgfsetbuttcap%
\pgfsetmiterjoin%
\definecolor{currentfill}{rgb}{0.000000,0.380392,0.809804}%
\pgfsetfillcolor{currentfill}%
\pgfsetlinewidth{0.000000pt}%
\definecolor{currentstroke}{rgb}{0.000000,0.000000,0.000000}%
\pgfsetstrokecolor{currentstroke}%
\pgfsetstrokeopacity{0.000000}%
\pgfsetdash{}{0pt}%
\pgfpathmoveto{\pgfqpoint{3.543213in}{1.191129in}}%
\pgfpathlineto{\pgfqpoint{3.680411in}{1.191129in}}%
\pgfpathlineto{\pgfqpoint{3.680411in}{1.283529in}}%
\pgfpathlineto{\pgfqpoint{3.543213in}{1.283529in}}%
\pgfpathlineto{\pgfqpoint{3.543213in}{1.191129in}}%
\pgfpathclose%
\pgfusepath{fill}%
\end{pgfscope}%
\begin{pgfscope}%
\pgfpathrectangle{\pgfqpoint{3.319565in}{0.488889in}}{\pgfqpoint{0.360845in}{0.924000in}}%
\pgfusepath{clip}%
\pgfsetbuttcap%
\pgfsetmiterjoin%
\definecolor{currentfill}{rgb}{0.000000,0.380392,0.809804}%
\pgfsetfillcolor{currentfill}%
\pgfsetlinewidth{0.000000pt}%
\definecolor{currentstroke}{rgb}{0.000000,0.000000,0.000000}%
\pgfsetstrokecolor{currentstroke}%
\pgfsetstrokeopacity{0.000000}%
\pgfsetdash{}{0pt}%
\pgfpathmoveto{\pgfqpoint{3.543104in}{1.200369in}}%
\pgfpathlineto{\pgfqpoint{3.680411in}{1.200369in}}%
\pgfpathlineto{\pgfqpoint{3.680411in}{1.292769in}}%
\pgfpathlineto{\pgfqpoint{3.543104in}{1.292769in}}%
\pgfpathlineto{\pgfqpoint{3.543104in}{1.200369in}}%
\pgfpathclose%
\pgfusepath{fill}%
\end{pgfscope}%
\begin{pgfscope}%
\pgfpathrectangle{\pgfqpoint{3.319565in}{0.488889in}}{\pgfqpoint{0.360845in}{0.924000in}}%
\pgfusepath{clip}%
\pgfsetbuttcap%
\pgfsetmiterjoin%
\definecolor{currentfill}{rgb}{0.000000,0.372549,0.813725}%
\pgfsetfillcolor{currentfill}%
\pgfsetlinewidth{0.000000pt}%
\definecolor{currentstroke}{rgb}{0.000000,0.000000,0.000000}%
\pgfsetstrokecolor{currentstroke}%
\pgfsetstrokeopacity{0.000000}%
\pgfsetdash{}{0pt}%
\pgfpathmoveto{\pgfqpoint{3.542527in}{1.209609in}}%
\pgfpathlineto{\pgfqpoint{3.680411in}{1.209609in}}%
\pgfpathlineto{\pgfqpoint{3.680411in}{1.302009in}}%
\pgfpathlineto{\pgfqpoint{3.542527in}{1.302009in}}%
\pgfpathlineto{\pgfqpoint{3.542527in}{1.209609in}}%
\pgfpathclose%
\pgfusepath{fill}%
\end{pgfscope}%
\begin{pgfscope}%
\pgfpathrectangle{\pgfqpoint{3.319565in}{0.488889in}}{\pgfqpoint{0.360845in}{0.924000in}}%
\pgfusepath{clip}%
\pgfsetbuttcap%
\pgfsetmiterjoin%
\definecolor{currentfill}{rgb}{0.000000,0.368627,0.815686}%
\pgfsetfillcolor{currentfill}%
\pgfsetlinewidth{0.000000pt}%
\definecolor{currentstroke}{rgb}{0.000000,0.000000,0.000000}%
\pgfsetstrokecolor{currentstroke}%
\pgfsetstrokeopacity{0.000000}%
\pgfsetdash{}{0pt}%
\pgfpathmoveto{\pgfqpoint{3.542427in}{1.218849in}}%
\pgfpathlineto{\pgfqpoint{3.680411in}{1.218849in}}%
\pgfpathlineto{\pgfqpoint{3.680411in}{1.311249in}}%
\pgfpathlineto{\pgfqpoint{3.542427in}{1.311249in}}%
\pgfpathlineto{\pgfqpoint{3.542427in}{1.218849in}}%
\pgfpathclose%
\pgfusepath{fill}%
\end{pgfscope}%
\begin{pgfscope}%
\pgfpathrectangle{\pgfqpoint{3.319565in}{0.488889in}}{\pgfqpoint{0.360845in}{0.924000in}}%
\pgfusepath{clip}%
\pgfsetbuttcap%
\pgfsetmiterjoin%
\definecolor{currentfill}{rgb}{0.000000,0.364706,0.817647}%
\pgfsetfillcolor{currentfill}%
\pgfsetlinewidth{0.000000pt}%
\definecolor{currentstroke}{rgb}{0.000000,0.000000,0.000000}%
\pgfsetstrokecolor{currentstroke}%
\pgfsetstrokeopacity{0.000000}%
\pgfsetdash{}{0pt}%
\pgfpathmoveto{\pgfqpoint{3.541911in}{1.228089in}}%
\pgfpathlineto{\pgfqpoint{3.680411in}{1.228089in}}%
\pgfpathlineto{\pgfqpoint{3.680411in}{1.320489in}}%
\pgfpathlineto{\pgfqpoint{3.541911in}{1.320489in}}%
\pgfpathlineto{\pgfqpoint{3.541911in}{1.228089in}}%
\pgfpathclose%
\pgfusepath{fill}%
\end{pgfscope}%
\begin{pgfscope}%
\pgfpathrectangle{\pgfqpoint{3.319565in}{0.488889in}}{\pgfqpoint{0.360845in}{0.924000in}}%
\pgfusepath{clip}%
\pgfsetbuttcap%
\pgfsetmiterjoin%
\definecolor{currentfill}{rgb}{0.000000,0.356863,0.821569}%
\pgfsetfillcolor{currentfill}%
\pgfsetlinewidth{0.000000pt}%
\definecolor{currentstroke}{rgb}{0.000000,0.000000,0.000000}%
\pgfsetstrokecolor{currentstroke}%
\pgfsetstrokeopacity{0.000000}%
\pgfsetdash{}{0pt}%
\pgfpathmoveto{\pgfqpoint{3.541325in}{1.237329in}}%
\pgfpathlineto{\pgfqpoint{3.680411in}{1.237329in}}%
\pgfpathlineto{\pgfqpoint{3.680411in}{1.329729in}}%
\pgfpathlineto{\pgfqpoint{3.541325in}{1.329729in}}%
\pgfpathlineto{\pgfqpoint{3.541325in}{1.237329in}}%
\pgfpathclose%
\pgfusepath{fill}%
\end{pgfscope}%
\begin{pgfscope}%
\pgfpathrectangle{\pgfqpoint{3.319565in}{0.488889in}}{\pgfqpoint{0.360845in}{0.924000in}}%
\pgfusepath{clip}%
\pgfsetbuttcap%
\pgfsetmiterjoin%
\definecolor{currentfill}{rgb}{0.000000,0.352941,0.823529}%
\pgfsetfillcolor{currentfill}%
\pgfsetlinewidth{0.000000pt}%
\definecolor{currentstroke}{rgb}{0.000000,0.000000,0.000000}%
\pgfsetstrokecolor{currentstroke}%
\pgfsetstrokeopacity{0.000000}%
\pgfsetdash{}{0pt}%
\pgfpathmoveto{\pgfqpoint{3.540997in}{1.246569in}}%
\pgfpathlineto{\pgfqpoint{3.680411in}{1.246569in}}%
\pgfpathlineto{\pgfqpoint{3.680411in}{1.338969in}}%
\pgfpathlineto{\pgfqpoint{3.540997in}{1.338969in}}%
\pgfpathlineto{\pgfqpoint{3.540997in}{1.246569in}}%
\pgfpathclose%
\pgfusepath{fill}%
\end{pgfscope}%
\begin{pgfscope}%
\pgfpathrectangle{\pgfqpoint{3.319565in}{0.488889in}}{\pgfqpoint{0.360845in}{0.924000in}}%
\pgfusepath{clip}%
\pgfsetbuttcap%
\pgfsetmiterjoin%
\definecolor{currentfill}{rgb}{0.000000,0.352941,0.823529}%
\pgfsetfillcolor{currentfill}%
\pgfsetlinewidth{0.000000pt}%
\definecolor{currentstroke}{rgb}{0.000000,0.000000,0.000000}%
\pgfsetstrokecolor{currentstroke}%
\pgfsetstrokeopacity{0.000000}%
\pgfsetdash{}{0pt}%
\pgfpathmoveto{\pgfqpoint{3.540986in}{1.255809in}}%
\pgfpathlineto{\pgfqpoint{3.680411in}{1.255809in}}%
\pgfpathlineto{\pgfqpoint{3.680411in}{1.348209in}}%
\pgfpathlineto{\pgfqpoint{3.540986in}{1.348209in}}%
\pgfpathlineto{\pgfqpoint{3.540986in}{1.255809in}}%
\pgfpathclose%
\pgfusepath{fill}%
\end{pgfscope}%
\begin{pgfscope}%
\pgfpathrectangle{\pgfqpoint{3.319565in}{0.488889in}}{\pgfqpoint{0.360845in}{0.924000in}}%
\pgfusepath{clip}%
\pgfsetbuttcap%
\pgfsetmiterjoin%
\definecolor{currentfill}{rgb}{0.000000,0.352941,0.823529}%
\pgfsetfillcolor{currentfill}%
\pgfsetlinewidth{0.000000pt}%
\definecolor{currentstroke}{rgb}{0.000000,0.000000,0.000000}%
\pgfsetstrokecolor{currentstroke}%
\pgfsetstrokeopacity{0.000000}%
\pgfsetdash{}{0pt}%
\pgfpathmoveto{\pgfqpoint{3.540986in}{1.265049in}}%
\pgfpathlineto{\pgfqpoint{3.680411in}{1.265049in}}%
\pgfpathlineto{\pgfqpoint{3.680411in}{1.357449in}}%
\pgfpathlineto{\pgfqpoint{3.540986in}{1.357449in}}%
\pgfpathlineto{\pgfqpoint{3.540986in}{1.265049in}}%
\pgfpathclose%
\pgfusepath{fill}%
\end{pgfscope}%
\begin{pgfscope}%
\pgfpathrectangle{\pgfqpoint{3.319565in}{0.488889in}}{\pgfqpoint{0.360845in}{0.924000in}}%
\pgfusepath{clip}%
\pgfsetbuttcap%
\pgfsetmiterjoin%
\definecolor{currentfill}{rgb}{0.000000,0.341176,0.829412}%
\pgfsetfillcolor{currentfill}%
\pgfsetlinewidth{0.000000pt}%
\definecolor{currentstroke}{rgb}{0.000000,0.000000,0.000000}%
\pgfsetstrokecolor{currentstroke}%
\pgfsetstrokeopacity{0.000000}%
\pgfsetdash{}{0pt}%
\pgfpathmoveto{\pgfqpoint{3.540080in}{1.274289in}}%
\pgfpathlineto{\pgfqpoint{3.680411in}{1.274289in}}%
\pgfpathlineto{\pgfqpoint{3.680411in}{1.366689in}}%
\pgfpathlineto{\pgfqpoint{3.540080in}{1.366689in}}%
\pgfpathlineto{\pgfqpoint{3.540080in}{1.274289in}}%
\pgfpathclose%
\pgfusepath{fill}%
\end{pgfscope}%
\begin{pgfscope}%
\pgfpathrectangle{\pgfqpoint{3.319565in}{0.488889in}}{\pgfqpoint{0.360845in}{0.924000in}}%
\pgfusepath{clip}%
\pgfsetbuttcap%
\pgfsetmiterjoin%
\definecolor{currentfill}{rgb}{0.000000,0.333333,0.833333}%
\pgfsetfillcolor{currentfill}%
\pgfsetlinewidth{0.000000pt}%
\definecolor{currentstroke}{rgb}{0.000000,0.000000,0.000000}%
\pgfsetstrokecolor{currentstroke}%
\pgfsetstrokeopacity{0.000000}%
\pgfsetdash{}{0pt}%
\pgfpathmoveto{\pgfqpoint{3.539586in}{1.283529in}}%
\pgfpathlineto{\pgfqpoint{3.680411in}{1.283529in}}%
\pgfpathlineto{\pgfqpoint{3.680411in}{1.375929in}}%
\pgfpathlineto{\pgfqpoint{3.539586in}{1.375929in}}%
\pgfpathlineto{\pgfqpoint{3.539586in}{1.283529in}}%
\pgfpathclose%
\pgfusepath{fill}%
\end{pgfscope}%
\begin{pgfscope}%
\pgfpathrectangle{\pgfqpoint{3.319565in}{0.488889in}}{\pgfqpoint{0.360845in}{0.924000in}}%
\pgfusepath{clip}%
\pgfsetbuttcap%
\pgfsetmiterjoin%
\definecolor{currentfill}{rgb}{0.000000,0.333333,0.833333}%
\pgfsetfillcolor{currentfill}%
\pgfsetlinewidth{0.000000pt}%
\definecolor{currentstroke}{rgb}{0.000000,0.000000,0.000000}%
\pgfsetstrokecolor{currentstroke}%
\pgfsetstrokeopacity{0.000000}%
\pgfsetdash{}{0pt}%
\pgfpathmoveto{\pgfqpoint{3.539576in}{1.292769in}}%
\pgfpathlineto{\pgfqpoint{3.680411in}{1.292769in}}%
\pgfpathlineto{\pgfqpoint{3.680411in}{1.385169in}}%
\pgfpathlineto{\pgfqpoint{3.539576in}{1.385169in}}%
\pgfpathlineto{\pgfqpoint{3.539576in}{1.292769in}}%
\pgfpathclose%
\pgfusepath{fill}%
\end{pgfscope}%
\begin{pgfscope}%
\pgfpathrectangle{\pgfqpoint{3.319565in}{0.488889in}}{\pgfqpoint{0.360845in}{0.924000in}}%
\pgfusepath{clip}%
\pgfsetbuttcap%
\pgfsetmiterjoin%
\definecolor{currentfill}{rgb}{0.000000,0.333333,0.833333}%
\pgfsetfillcolor{currentfill}%
\pgfsetlinewidth{0.000000pt}%
\definecolor{currentstroke}{rgb}{0.000000,0.000000,0.000000}%
\pgfsetstrokecolor{currentstroke}%
\pgfsetstrokeopacity{0.000000}%
\pgfsetdash{}{0pt}%
\pgfpathmoveto{\pgfqpoint{3.539297in}{1.302009in}}%
\pgfpathlineto{\pgfqpoint{3.680411in}{1.302009in}}%
\pgfpathlineto{\pgfqpoint{3.680411in}{1.394409in}}%
\pgfpathlineto{\pgfqpoint{3.539297in}{1.394409in}}%
\pgfpathlineto{\pgfqpoint{3.539297in}{1.302009in}}%
\pgfpathclose%
\pgfusepath{fill}%
\end{pgfscope}%
\begin{pgfscope}%
\pgfpathrectangle{\pgfqpoint{3.319565in}{0.488889in}}{\pgfqpoint{0.360845in}{0.924000in}}%
\pgfusepath{clip}%
\pgfsetbuttcap%
\pgfsetmiterjoin%
\definecolor{currentfill}{rgb}{0.000000,0.333333,0.833333}%
\pgfsetfillcolor{currentfill}%
\pgfsetlinewidth{0.000000pt}%
\definecolor{currentstroke}{rgb}{0.000000,0.000000,0.000000}%
\pgfsetstrokecolor{currentstroke}%
\pgfsetstrokeopacity{0.000000}%
\pgfsetdash{}{0pt}%
\pgfpathmoveto{\pgfqpoint{3.539297in}{1.311249in}}%
\pgfpathlineto{\pgfqpoint{3.680411in}{1.311249in}}%
\pgfpathlineto{\pgfqpoint{3.680411in}{1.403649in}}%
\pgfpathlineto{\pgfqpoint{3.539297in}{1.403649in}}%
\pgfpathlineto{\pgfqpoint{3.539297in}{1.311249in}}%
\pgfpathclose%
\pgfusepath{fill}%
\end{pgfscope}%
\begin{pgfscope}%
\pgfpathrectangle{\pgfqpoint{3.319565in}{0.488889in}}{\pgfqpoint{0.360845in}{0.924000in}}%
\pgfusepath{clip}%
\pgfsetbuttcap%
\pgfsetmiterjoin%
\definecolor{currentfill}{rgb}{0.000000,0.313725,0.843137}%
\pgfsetfillcolor{currentfill}%
\pgfsetlinewidth{0.000000pt}%
\definecolor{currentstroke}{rgb}{0.000000,0.000000,0.000000}%
\pgfsetstrokecolor{currentstroke}%
\pgfsetstrokeopacity{0.000000}%
\pgfsetdash{}{0pt}%
\pgfpathmoveto{\pgfqpoint{3.537936in}{1.320489in}}%
\pgfpathlineto{\pgfqpoint{3.680411in}{1.320489in}}%
\pgfpathlineto{\pgfqpoint{3.680411in}{1.412889in}}%
\pgfpathlineto{\pgfqpoint{3.537936in}{1.412889in}}%
\pgfpathlineto{\pgfqpoint{3.537936in}{1.320489in}}%
\pgfpathclose%
\pgfusepath{fill}%
\end{pgfscope}%
\begin{pgfscope}%
\pgfpathrectangle{\pgfqpoint{3.319565in}{0.488889in}}{\pgfqpoint{0.360845in}{0.924000in}}%
\pgfusepath{clip}%
\pgfsetbuttcap%
\pgfsetmiterjoin%
\definecolor{currentfill}{rgb}{0.000000,0.301961,0.849020}%
\pgfsetfillcolor{currentfill}%
\pgfsetlinewidth{0.000000pt}%
\definecolor{currentstroke}{rgb}{0.000000,0.000000,0.000000}%
\pgfsetstrokecolor{currentstroke}%
\pgfsetstrokeopacity{0.000000}%
\pgfsetdash{}{0pt}%
\pgfpathmoveto{\pgfqpoint{3.536971in}{1.329729in}}%
\pgfpathlineto{\pgfqpoint{3.680411in}{1.329729in}}%
\pgfpathlineto{\pgfqpoint{3.680411in}{1.422129in}}%
\pgfpathlineto{\pgfqpoint{3.536971in}{1.422129in}}%
\pgfpathlineto{\pgfqpoint{3.536971in}{1.329729in}}%
\pgfpathclose%
\pgfusepath{fill}%
\end{pgfscope}%
\begin{pgfscope}%
\pgfpathrectangle{\pgfqpoint{3.319565in}{0.488889in}}{\pgfqpoint{0.360845in}{0.924000in}}%
\pgfusepath{clip}%
\pgfsetbuttcap%
\pgfsetmiterjoin%
\definecolor{currentfill}{rgb}{0.000000,0.290196,0.854902}%
\pgfsetfillcolor{currentfill}%
\pgfsetlinewidth{0.000000pt}%
\definecolor{currentstroke}{rgb}{0.000000,0.000000,0.000000}%
\pgfsetstrokecolor{currentstroke}%
\pgfsetstrokeopacity{0.000000}%
\pgfsetdash{}{0pt}%
\pgfpathmoveto{\pgfqpoint{3.535827in}{1.338969in}}%
\pgfpathlineto{\pgfqpoint{3.680411in}{1.338969in}}%
\pgfpathlineto{\pgfqpoint{3.680411in}{1.431369in}}%
\pgfpathlineto{\pgfqpoint{3.535827in}{1.431369in}}%
\pgfpathlineto{\pgfqpoint{3.535827in}{1.338969in}}%
\pgfpathclose%
\pgfusepath{fill}%
\end{pgfscope}%
\begin{pgfscope}%
\pgfpathrectangle{\pgfqpoint{3.319565in}{0.488889in}}{\pgfqpoint{0.360845in}{0.924000in}}%
\pgfusepath{clip}%
\pgfsetbuttcap%
\pgfsetmiterjoin%
\definecolor{currentfill}{rgb}{0.000000,0.286275,0.856863}%
\pgfsetfillcolor{currentfill}%
\pgfsetlinewidth{0.000000pt}%
\definecolor{currentstroke}{rgb}{0.000000,0.000000,0.000000}%
\pgfsetstrokecolor{currentstroke}%
\pgfsetstrokeopacity{0.000000}%
\pgfsetdash{}{0pt}%
\pgfpathmoveto{\pgfqpoint{3.535510in}{1.348209in}}%
\pgfpathlineto{\pgfqpoint{3.680411in}{1.348209in}}%
\pgfpathlineto{\pgfqpoint{3.680411in}{1.440609in}}%
\pgfpathlineto{\pgfqpoint{3.535510in}{1.440609in}}%
\pgfpathlineto{\pgfqpoint{3.535510in}{1.348209in}}%
\pgfpathclose%
\pgfusepath{fill}%
\end{pgfscope}%
\begin{pgfscope}%
\pgfpathrectangle{\pgfqpoint{3.319565in}{0.488889in}}{\pgfqpoint{0.360845in}{0.924000in}}%
\pgfusepath{clip}%
\pgfsetbuttcap%
\pgfsetmiterjoin%
\definecolor{currentfill}{rgb}{0.000000,0.243137,0.878431}%
\pgfsetfillcolor{currentfill}%
\pgfsetlinewidth{0.000000pt}%
\definecolor{currentstroke}{rgb}{0.000000,0.000000,0.000000}%
\pgfsetstrokecolor{currentstroke}%
\pgfsetstrokeopacity{0.000000}%
\pgfsetdash{}{0pt}%
\pgfpathmoveto{\pgfqpoint{3.532156in}{1.357449in}}%
\pgfpathlineto{\pgfqpoint{3.680411in}{1.357449in}}%
\pgfpathlineto{\pgfqpoint{3.680411in}{1.449849in}}%
\pgfpathlineto{\pgfqpoint{3.532156in}{1.449849in}}%
\pgfpathlineto{\pgfqpoint{3.532156in}{1.357449in}}%
\pgfpathclose%
\pgfusepath{fill}%
\end{pgfscope}%
\begin{pgfscope}%
\pgfpathrectangle{\pgfqpoint{3.319565in}{0.488889in}}{\pgfqpoint{0.360845in}{0.924000in}}%
\pgfusepath{clip}%
\pgfsetbuttcap%
\pgfsetmiterjoin%
\definecolor{currentfill}{rgb}{0.000000,0.243137,0.878431}%
\pgfsetfillcolor{currentfill}%
\pgfsetlinewidth{0.000000pt}%
\definecolor{currentstroke}{rgb}{0.000000,0.000000,0.000000}%
\pgfsetstrokecolor{currentstroke}%
\pgfsetstrokeopacity{0.000000}%
\pgfsetdash{}{0pt}%
\pgfpathmoveto{\pgfqpoint{3.532156in}{1.366689in}}%
\pgfpathlineto{\pgfqpoint{3.680411in}{1.366689in}}%
\pgfpathlineto{\pgfqpoint{3.680411in}{1.459089in}}%
\pgfpathlineto{\pgfqpoint{3.532156in}{1.459089in}}%
\pgfpathlineto{\pgfqpoint{3.532156in}{1.366689in}}%
\pgfpathclose%
\pgfusepath{fill}%
\end{pgfscope}%
\begin{pgfscope}%
\pgfpathrectangle{\pgfqpoint{3.319565in}{0.488889in}}{\pgfqpoint{0.360845in}{0.924000in}}%
\pgfusepath{clip}%
\pgfsetbuttcap%
\pgfsetmiterjoin%
\definecolor{currentfill}{rgb}{0.000000,0.219608,0.890196}%
\pgfsetfillcolor{currentfill}%
\pgfsetlinewidth{0.000000pt}%
\definecolor{currentstroke}{rgb}{0.000000,0.000000,0.000000}%
\pgfsetstrokecolor{currentstroke}%
\pgfsetstrokeopacity{0.000000}%
\pgfsetdash{}{0pt}%
\pgfpathmoveto{\pgfqpoint{3.530237in}{1.375929in}}%
\pgfpathlineto{\pgfqpoint{3.680411in}{1.375929in}}%
\pgfpathlineto{\pgfqpoint{3.680411in}{1.468329in}}%
\pgfpathlineto{\pgfqpoint{3.530237in}{1.468329in}}%
\pgfpathlineto{\pgfqpoint{3.530237in}{1.375929in}}%
\pgfpathclose%
\pgfusepath{fill}%
\end{pgfscope}%
\begin{pgfscope}%
\pgfpathrectangle{\pgfqpoint{3.319565in}{0.488889in}}{\pgfqpoint{0.360845in}{0.924000in}}%
\pgfusepath{clip}%
\pgfsetbuttcap%
\pgfsetmiterjoin%
\definecolor{currentfill}{rgb}{0.000000,0.196078,0.901961}%
\pgfsetfillcolor{currentfill}%
\pgfsetlinewidth{0.000000pt}%
\definecolor{currentstroke}{rgb}{0.000000,0.000000,0.000000}%
\pgfsetstrokecolor{currentstroke}%
\pgfsetstrokeopacity{0.000000}%
\pgfsetdash{}{0pt}%
\pgfpathmoveto{\pgfqpoint{3.528298in}{1.385169in}}%
\pgfpathlineto{\pgfqpoint{3.680411in}{1.385169in}}%
\pgfpathlineto{\pgfqpoint{3.680411in}{1.477569in}}%
\pgfpathlineto{\pgfqpoint{3.528298in}{1.477569in}}%
\pgfpathlineto{\pgfqpoint{3.528298in}{1.385169in}}%
\pgfpathclose%
\pgfusepath{fill}%
\end{pgfscope}%
\begin{pgfscope}%
\pgfpathrectangle{\pgfqpoint{3.319565in}{0.488889in}}{\pgfqpoint{0.360845in}{0.924000in}}%
\pgfusepath{clip}%
\pgfsetbuttcap%
\pgfsetmiterjoin%
\definecolor{currentfill}{rgb}{0.000000,0.196078,0.901961}%
\pgfsetfillcolor{currentfill}%
\pgfsetlinewidth{0.000000pt}%
\definecolor{currentstroke}{rgb}{0.000000,0.000000,0.000000}%
\pgfsetstrokecolor{currentstroke}%
\pgfsetstrokeopacity{0.000000}%
\pgfsetdash{}{0pt}%
\pgfpathmoveto{\pgfqpoint{3.528298in}{1.394409in}}%
\pgfpathlineto{\pgfqpoint{3.680411in}{1.394409in}}%
\pgfpathlineto{\pgfqpoint{3.680411in}{1.486809in}}%
\pgfpathlineto{\pgfqpoint{3.528298in}{1.486809in}}%
\pgfpathlineto{\pgfqpoint{3.528298in}{1.394409in}}%
\pgfpathclose%
\pgfusepath{fill}%
\end{pgfscope}%
\begin{pgfscope}%
\pgfpathrectangle{\pgfqpoint{3.319565in}{0.488889in}}{\pgfqpoint{0.360845in}{0.924000in}}%
\pgfusepath{clip}%
\pgfsetbuttcap%
\pgfsetmiterjoin%
\definecolor{currentfill}{rgb}{0.000000,0.000000,1.000000}%
\pgfsetfillcolor{currentfill}%
\pgfsetlinewidth{0.000000pt}%
\definecolor{currentstroke}{rgb}{0.000000,0.000000,0.000000}%
\pgfsetstrokecolor{currentstroke}%
\pgfsetstrokeopacity{0.000000}%
\pgfsetdash{}{0pt}%
\pgfpathmoveto{\pgfqpoint{3.512377in}{1.403649in}}%
\pgfpathlineto{\pgfqpoint{3.680411in}{1.403649in}}%
\pgfpathlineto{\pgfqpoint{3.680411in}{1.496049in}}%
\pgfpathlineto{\pgfqpoint{3.512377in}{1.496049in}}%
\pgfpathlineto{\pgfqpoint{3.512377in}{1.403649in}}%
\pgfpathclose%
\pgfusepath{fill}%
\end{pgfscope}%
\begin{pgfscope}%
\definecolor{textcolor}{rgb}{0.000000,0.000000,0.000000}%
\pgfsetstrokecolor{textcolor}%
\pgfsetfillcolor{textcolor}%
\pgftext[x=3.471121in,y=1.468329in,left,base]{\color{textcolor}\setmainfont{Lato}\rmfamily\fontsize{7.000000}{8.400000}\selectfont 46.6\%}%
\end{pgfscope}%
\begin{pgfscope}%
\definecolor{textcolor}{rgb}{0.000000,0.000000,0.000000}%
\pgfsetstrokecolor{textcolor}%
\pgfsetfillcolor{textcolor}%
\pgftext[x=3.471121in,y=0.359529in,left,base]{\color{textcolor}\setmainfont{Lato}\rmfamily\fontsize{7.000000}{8.400000}\selectfont 24.1\%}%
\end{pgfscope}%
\end{pgfpicture}%
\makeatother%
\endgroup%


\vspace{-5mm}
\begin{minipage}{0.76\textwidth}
\footnotesize{Source: American Community Survey, Author's Calculations} \hfill \tbllink{low_income_cz.csv} 
\end{minipage}

\newpage
\begin{minipage}{0.76\textwidth}
\section*{Businesses}
\hypertarget{bus}{}
\small The factories, offices, and equipment that workers use to produce goods and services are all important to the economy. This section looks at the business sector, with data covering business investment, retail sales, industrial production, corporate profits, and the financial activities of private businesses. 

\subsection*{Investment}
\hypertarget{busin}{}
\small Production by private businesses requires capital goods, such as buildings, equipment, and software. When businesses purchase such items, defined as having a useful life of more than one year, it is considered \href{https://www.bea.gov/help/glossary/nonresidential-fixed-investment}{fixed investment}, or investment in fixed assets. From an accounting perspective, these transactions are considered an exchange of assets (cash in exchange for capital goods) rather than an expense.

Over time, these capital goods wear down, a process called consumption of fixed capital or depreciation. Depreciation is the expense related to capital goods from an accounting perspective. Businesses must decide whether to replace or add to the existing stock of capital goods, and their new purchases of capital goods and inventory investment are considered \href{https://www.bea.gov/help/glossary/gross-private-domestic-investment}{gross investment}. Accordingly, net investment is gross investment minus depreciation, and measures whether the stock of capital goods is expanding.

Net investment is important for many reasons. In the short run, the production and installation of capital goods adds directly to GDP and increases the level of economic activity. In the long run, \textbf{investments in fixed assets make workers more productive}, as they allow businesses to produce more goods and services with the same hours of work. 
\end{minipage}
\vspace{1mm}

\begin{minipage}{0.37\textwidth}
\normalsize \textbf{Investment and Depreciation}\\
\footnotesize{\textit{share of GDP, percent}}\\
\hspace*{-2mm} \begin{tikzpicture}
	\begin{axis}[\bbar{y}{0}, \shdateaxisticks ytick={0, 5, 10, 15}, 
		height=8.6cm, width=6.8cm, xmin={1989-12-01}, clip=false,
		yticklabel style={text width=1.0em},
		xticklabel={`\short{\year}}]
	\rbars
	\thickline{cyan!80!white}{date}{A008RC}{data/businv_main.csv}
	\stdline{blue!60!violet}{date}{W987RC}{data/businv_main.csv}
	\stdline{magenta}{date}{W276RC}{data/businv_main.csv}
	\stdline{green!80!blue}{date}{W790RC}{data/businv_main.csv}
	\stdnode{3.1cm}{6.5cm}{\footnotesize \color{blue!60!violet}Gross Investment}
	\stdnode{1.5cm}{5.2cm}{\footnotesize \color{cyan!80!white}\ \ \ Fixed\\ \footnotesize \color{cyan!80!white} Investment}
	\stdnode{1.2cm}{3.95cm}{\footnotesize \color{magenta}Depreciation}
	\stdnode{4.3cm}{2.6cm}{\footnotesize \color{green!80!blue}Net \\ \footnotesize \color{green!80!blue}Investment}
	\end{axis}
\end{tikzpicture}\\
\footnotesize{Source: Bureau of Economic Analysis} \hfill \tbllink{businv_main.csv}
\end{minipage} \hspace{4.5mm}
\begin{minipage}{0.35\textwidth}
\small \input{text/businv_levels.txt}
\end{minipage}
\vspace{1mm}

\begin{minipage}{0.76\textwidth}
\small Note that gross investment includes fixed investment and inventory investment, or the \href{https://www.bea.gov/help/glossary/change-private-inventories-cipi}{change in private inventories}. Changes to private inventories capture the difference between sales and production. Reduced production of new inventory explains much of the overall reduction in gross investment during the COVID-19 pandemic. 
\end{minipage}
\newpage
\begin{minipage}{0.3\textwidth}
\small \input{text/businv_sh.txt}
\end{minipage} \hspace{6mm}
\begin{minipage}{0.41\textwidth}
\normalsize \textbf{Business Fixed Investment by Type}\\
\footnotesize{\textit{share of GDP, percent}}\\
\hspace*{-2mm} \begin{tikzpicture}
	\begin{axis}[\bbar{y}{0}, \shdateaxisticks 
		ytick={0, 2, 4, 6}, xmin={1989-12-01}, clip=false,
		ymin=0.4, xticklabel={`\short{\year}}, height=5.0cm, width=7.5cm]
	\rbars
	\stdline{yellow!50!orange}{date}{B009RC}{data/businvsh.csv}
	\stdline{cyan!60!white}{date}{Y033RC}{data/businvsh.csv}
	\stdline{violet}{date}{Y001RC}{data/businvsh.csv}
	\stdnode{4.95cm}{1.1cm}{\footnotesize \color{yellow!50!orange} Structures}
	\stdnode{4.7cm}{2.9cm}{\footnotesize \color{cyan!60!white} Equipment}
	\stdnode{1.35cm}{2.1cm}{\footnotesize \color{violet} Intellectual \\ \footnotesize \color{violet} Property}
	\end{axis}
\end{tikzpicture}\\
\footnotesize{Source: Bureau of Economic Analysis} \hfill \tbllink{businvsh.csv}
\end{minipage}

\begin{minipage}{0.76\textwidth}
\subsubsection*{Contribution to Growth}
\vspace{-0.5mm}
\small &   2023  Q3 & `23  Q2 & `23  Q1 & `22  Q3 & `21  Q3 &1-year&10-year&30-year\\ Total&0.21&0.98&0.76&0.62&-0.15&0.55&0.54&0.60\\  \hspace{-2mm}\cbox{yellow!50!orange}Structures &0.33&0.46&0.77&-0.03&-0.12&0.43&0.05&0.03\\  \hspace{-2mm}\cbox{cyan!60!white}Equipment &-0.22&0.38&-0.21&0.28&-0.40&-0.08&0.16&0.31\\  \hspace{4mm}  Information  Processing &-0.12&-0.11&-0.02&0.09&-0.09&-0.16&0.10&0.20\\  \hspace{6mm}  Computers  \&  Peripherals &-0.11&0.02&-0.05&0.09&-0.03&-0.09&0.02&0.10\\  \hspace{4mm}  Industrial  Equipment &-0.06&-0.06&0.04&-0.07&0.08&-0.01&0.02&0.03\\  \hspace{4mm}  Transportation  Equipment &-0.02&0.54&-0.14&0.30&-0.41&0.13&0.02&0.05\\  \hspace{-2mm}\cbox{violet}Intellectual  Property  Products &0.10&0.15&0.20&0.37&0.37&0.19&0.34&0.25\\  \hspace{4mm}  Software &0.17&0.13&0.16&0.28&0.18&0.19&0.21&0.15\\  \hspace{4mm}  Research  \&  Development &-0.06&0.00&0.04&0.05&0.14&0.00&0.12&0.09\\  
\vspace{1mm}

\normalsize \textbf{Business Fixed Investment}\\
\footnotesize{\textit{percentage point contribution to real GDP growth}}\\
\hspace*{-2mm} \begin{tikzpicture}
	\begin{axis}[\bbar{y}{0}, \dateaxisticks ytick={-4, -2, 0, 2, 4},
		xticklabel={`\short{\year}}, clip=false, height=5.5cm,
		legend style={at={(0.72, 1.0)}, fill=white}]
	\rbars
	\sbar{yellow!50!orange}{date}{A009RY}{data/businv.csv}
	\sbar{cyan!60!white}{date}{Y033RY}{data/businv.csv}
	\sbar{violet}{date}{Y001RY}{data/businv.csv}
	\legend{Structures, Equipment, Intellectual Property Products};
	\end{axis}
\end{tikzpicture}\\
\footnotesize{Source: Bureau of Economic Analysis} \hfill \tbllink{businv.csv}
\vspace{2mm}

\normalsize \textbf{Business Fixed Investment}\\
\footnotesize{\textit{percentage point contribution to real GDP growth \hspace{36mm} moving averages}}\\ 
\rowcolors{1}{}{black!5} \setlength{\tabcolsep}{3.1pt} \color{black!90}
		{\renewcommand{\arraystretch}{1.55}
		 \begin{tabular}{p{42mm} R{6.7mm} R{6.7mm} R{6.7mm} R{6.7mm} R{6.7mm} 
		   R{7.8mm} R{6.7mm} R{6.7mm} }
			 \input data/businv.tex \hline
		\end{tabular}}\vspace{-2mm}
		
\footnotesize{Source: Bureau of Economic Analysis}
\end{minipage}
\newpage
\begin{minipage}{0.32\textwidth}
\small Productive business investments also show up as \textbf{new orders for core capital goods}. The category excludes the more-volatile aircraft orders as well as defense-related orders, and is derived from a Census Bureau \href{https://www.census.gov/manufacturing/m3/index.html}{survey} of shipments, inventories, and orders.  

\input{text/dgno.txt}
\end{minipage} \hspace{6mm}
\begin{minipage}{0.39\textwidth}
\normalsize \textbf{New Orders for Core Capital Goods}\\
\footnotesize{\textit{nondefense capital goods ex-aircraft, share of GDP}}\\
\hspace*{-2mm} \begin{tikzpicture}
	\begin{axis}[\bbar{y}{0}, \shdateaxisticks ytick={4, 6, 8}, 
		width=6.8cm, height=5.8cm, clip=false, 
		xmin={1993-05-01},
		xticklabel={`\short{\year}}]
	\rebars
	\thickline{purple!50!violet}{date}{value}{data/dgno.csv}
	\input{text/dgno_node.txt}
	\end{axis}
\end{tikzpicture}\\
\footnotesize{Source: Census Bureau} \hfill \tbllink{dgno.csv}
\end{minipage}
\vspace{10mm}

\begin{minipage}{0.76\textwidth}
\subsection*{Inventories}
\small Changes in private inventories are often grouped with investment in national accounts. For purposes of flow measures such as GDP, inventory changes capture a situation where goods were produced but not sold and therefore are not included in consumption or investment. This situation can allow a drop in spending that was not predicted to be obscured in GDP growth figures. However, BEA \href{https://apps.bea.gov/iTable/index_UD.cfm}{report} an inflation-adjusted ratio of inventories to sales in manufacturing and trade businesses (see {\color{blue!60!green}\textbf{---}}). 

When thinking about the longer-term trends in these data, it's important to note that sales of businesses include services while inventories include only goods. Over the 30 years prior to COVID-19, sales shifted towards services, which reduced the inventories to sales ratio, all else equal. Following COVID-19, sales shifted back towards goods, increasing the inventories to sales ratio, all else equal, and affecting comparability of data over this period. 
\end{minipage}

\begin{minipage}{0.4\textwidth}
\normalsize \textbf{Inventories to Sales Ratio}\\
\footnotesize{\textit{total business}}\\
\hspace*{-2mm} \begin{tikzpicture}
	\begin{axis}[\shdateaxisticks ytick={1.2, 1.3, 1.4, 1.5, 1.6, 1.7, 1.8}, 
		yticklabel style={text width=1.2em}, legend style={at={(0.6, 0.98)}},
		xticklabel={`\short{\year}}, width=7.3cm, ymin=1.2, enlarge y limits={0.04},
		height=6.2cm]
	\rbars
	\thinline{blue!40!cyan!80!white}{date}{VALUE}{data/isratio.csv}
	\stdline{blue!60!green}{date}{Value}{data/real_isratio.csv}
	\legend{Nominal, Real}
	\end{axis}
\end{tikzpicture}\\
\footnotesize{Source: FRED, BEA} \hfill \tbllink{isratio.csv}
\end{minipage} \hspace{7mm}
\begin{minipage}{0.31\textwidth}
\small Census \href{https://www.census.gov/mtis/www/data/pdf/mtis_current.pdf}{report} the nominal \textbf{ratio of inventories to sales} for the total business sector (see {\color{blue!40!cyan!80!white}\textbf{---}}). \input{text/isratio.txt}
\end{minipage}
\newpage 
\hypertarget{buspr}{}
\begin{minipage}{0.76\textwidth}
\subsection*{Corporate Profits}
\small The national accounts include detailed information on \textbf{aggregate \href{https://www.bea.gov/data/income-saving/corporate-profits}{corporate profits}}. \input{text/cprof.txt} 
\vspace{1mm}

\normalsize \textbf{Destination of Corporate Profits}\\
\footnotesize{\textit{share of net national income}}\\
\hspace*{-2mm} \begin{tikzpicture}
	\begin{axis}[\bbar{y}{0}, \dateaxisticks ytick={0, 5, 10, 15}, ymin=0.4, 
		yticklabel style={text width=1.0em}, 
		height=5.6cm,
		xticklabel={`\short{\year}}, clip=false, legend style={at={(0.68, 1.12)}}]
	\rbars
	\sbar{blue!70!purple}{date}{B056RC}{data/cprof.csv}
	\sbar{cyan!50!white}{date}{A127RC}{data/cprof.csv}
	\sbar{red!80!orange}{date}{A054RC}{data/cprof.csv}	
	\legend{Net Dividends, Retained Earnings, Corporate Taxes};
	\end{axis}
\end{tikzpicture}\\
\footnotesize{Source: Bureau of Economic Analysis} \hfill \tbllink{cprof.csv}
\vspace{4mm}

\small \textbf{Aggregate corporate savings} (corporate profits less dividends and corporate profit tax) are the result of net investment and non-business saving. Investment (see\cbox{green!50!teal}) is a source of aggregate profit because it is revenue for one party but not an expense for the other. Non-business saving, which includes household (see\cbox{blue!80!cyan}), government (see\cbox{cyan!80!white}), and rest of world saving (see\cbox{yellow!70!orange}), necessarily reduces aggregate corporate profits because it is money that did not return to businesses as revenue. 
\vspace{1mm}

\normalsize \textbf{Sources of Corporate Saving}\\
\footnotesize{\textit{contribution to corporate saving, as share of gross national income}}
\vspace{5.0cm}

\hspace{4mm} \begin{tikzpicture}[overlay]
	\begin{axis}[\bbar{y}{0}, \dateaxisticks 
		yticklabel style={text width=1.5em}, height=7.0cm, ymax=17.6,
		xticklabel={`\short{\year}}, clip=false, legend style={at={(0.9, 0.99)}}]
	\rbars
	\sbar{cyan!80!white}{date}{Gov Saving}{data/cprof2.csv}
	\sbar{yellow!70!orange}{date}{ROW Saving}{data/cprof2.csv}
	\sbar{green!50!teal!90!white}{date}{Investment}{data/cprof2.csv}
	\sbar{blue!80!green}{date}{HH Saving}{data/cprof2.csv}	
	\legend{Government Saving, ROW Saving, Investment, Household Saving};
	\end{axis}
\end{tikzpicture}

\footnotesize{Source: Bureau of Economic Analysis} \hfill \tbllink{cprof2.csv}
\end{minipage}
\newpage
\begin{minipage}{0.76\textwidth}
\subsection*{Business Debt}
\small \input{text/busdebtgdp.txt} 
\vspace{2mm}

\normalsize \textbf{Nonfinancial Business Debt}\\
\footnotesize{\textit{by type, as share of GDP}}\\
\hspace*{-2mm} \begin{tikzpicture}
	\begin{axis}[\bbar{y}{0}, \dateaxisticks ytick={0, 25, 50, 75, 100}, ymin=2, 
		yticklabel style={text width=1.0em}, height=5.5cm,
		xticklabel={`\short{\year}}, clip=false, legend style={at={(0.93, 1.13)}}]
	\rbars
	\sbar{violet}{date}{Bank Loans and Mortgages}{data/busdebtgdp2.csv}	
	\sbar{magenta!70!purple}{date}{Debt Securities}{data/busdebtgdp2.csv}	
	\sbar{blue}{date}{Nonbank Loans}{data/busdebtgdp2.csv}	
	\legend{Mortgages and Bank Loans, Corporate Debt Securities, Nonbank Loans};
	\end{axis}
\end{tikzpicture}\\
\footnotesize{Source: Federal Reserve, Bureau of Economic Analysis} \hfill \tbllink{busdebtgdp2.csv}
\vspace{4mm}

\small The debt of the domestic financial sector includes agency and government-sponsored enterprise (GSE) backed securities (see\cbox{red!60!purple}), corporate and foreign bonds, loans (see\cbox{orange!70!yellow}), and open market paper. The long-term increase in financial sector debt reflects the emergence and growth of various asset-backed securities. In addition to home mortgage-backed securities, the domestic financial sector issues debt securities based on commercial mortgages, auto loans, credit cards, student debt, and even restaurant revenue. 

\input{text/findebtgdp.txt} 
\vspace{2mm}

\normalsize \textbf{Financial Sector Debt}\\
\footnotesize{\textit{by type, as share of GDP}}\\
\hspace*{-2mm} \begin{tikzpicture}
	\begin{axis}[\bbar{y}{0}, \dateaxisticks ytick={0, 50, 100, 150}, ymin=2, 
		yticklabel style={text width=1.6em}, height=5.5cm,
		xticklabel={`\short{\year}}, clip=false, legend style={at={(0.87, 1.13)}}]
	\rbars
	\sbar{orange!70!yellow}{date}{Financial Loans}{data/busdebtgdp2.csv}	
	\sbar{red!60!purple}{date}{Agency MBS}{data/busdebtgdp2.csv}
	\sbar{blue!40!violet}{date}{Other}{data/busdebtgdp2.csv}
	\legend{Loans, Agency- and GSE-backed Securities, Other Debt Securities};
	\end{axis}
\end{tikzpicture}\\
\footnotesize{Source: Federal Reserve, Bureau of Economic Analysis} \hfill \tbllink{busdebtgdp2.csv}
\end{minipage}
\newpage
\subsection*{Industrial Production}
\hypertarget{busip}{}
\begin{minipage}{0.6\textwidth}
\small & Dec  `23 & Nov  `23 & Oct  `23 & Dec  `22 &   Dec  `23 &   Nov  `23 &   Oct  `23 &   Dec  `22 \\  Total  Index &1.0&-0.6&-0.9&0.6&1.0&-0.6&-0.9&0.6\\  \hspace{2mm}Manufacturing &0.9&-0.7&-1.5&-1.0&1.2&-1.0&-1.9&-1.3\\  \hspace{-3mm}  \cbox{blue!62!black}  Durable  Manufacturing &0.4&0.1&-0.7&0.3&1.0&0.2&-2.0&0.7\\  \hspace{5mm}Motor  Vehicles  \&  Parts &0.3&0.1&-0.3&0.3&4.8&2.7&-6.4&5.7\\  \hspace{-3mm}  \cbox{blue!20!cyan!80!white}  Nondurable  Manufacturing &0.6&-0.7&-0.7&-1.2&1.7&-2.0&-1.8&-3.4\\  \hspace{-3mm}  \cbox{orange!30!yellow}  Mining &0.6&0.2&0.2&0.6&4.3&1.3&1.6&3.8\\  \hspace{-3mm}  \cbox{green!70!blue}  Utilities &-0.5&-0.1&0.4&1.0&-4.9&-0.8&3.2&9.2\\  \hspace{-3mm}  \cbox{violet!70!purple!90!black}  Consumer  Goods &-0.3&-0.5&-0.6&0.2&-1.0&-1.8&-2.1&0.9\\    \hspace{4mm}Consumer  Durables &0.1&-0.0&-0.3&-0.1&1.6&-0.4&-5.5&-1.7\\    \hspace{6mm}Automotive  Products &0.1&0.1&-0.2&0.0&4.5&2.0&-7.6&0.9\\    \hspace{4mm}Consumer  Nondurables &-0.4&-0.5&-0.2&0.3&-1.7&-2.2&-1.1&1.6\\    \hspace{6mm}Foods  \&  Tobacco &-0.2&-0.3&-0.3&-0.1&-1.7&-3.6&-3.4&-0.9\\    \hspace{6mm}Chemical  Products &-0.0&0.0&0.1&0.1&-0.3&0.7&2.4&3.1\\    \hspace{6mm}Consumer  Energy  Products &-0.1&-0.1&0.0&0.3&-2.4&-2.0&0.0&5.0\\  \hspace{-3mm}  \cbox{magenta}  Business  Equipment  \&  Supplies &0.1&-0.4&-0.5&0.3&0.2&-1.4&-1.9&1.0\\  \hspace{4mm}Equipment &0.1&-0.0&-0.2&0.5&1.1&-0.2&-1.9&4.5\\    \hspace{6mm}Industrial  Equipment &-0.1&-0.1&-0.1&0.1&-1.8&-3.3&-4.7&2.0\\    \hspace{4mm}Nonindustrial  Supplies &-0.1&-0.4&-0.3&-0.2&-0.4&-2.2&-1.9&-1.5\\    \hspace{6mm}Construction  Supplies &-0.1&-0.1&-0.1&-0.2&-1.0&-2.7&-2.6&-3.4\\    \hspace{6mm}Business  Supplies &-0.0&-0.2&-0.2&-0.1&-0.1&-2.0&-1.6&-0.5\\    \hspace{3mm}Materials &1.2&0.3&0.2&0.1&2.7&0.6&0.4&0.2\\    \hspace{6mm}Consumer  Parts &0.1&0.1&0.0&0.1&3.3&3.8&0.5&3.8\\    \hspace{6mm}Equipment  Parts &0.1&0.1&0.0&-0.0&3.1&1.9&0.2&-1.0\\    \hspace{6mm}Chemical  Materials &0.5&0.1&0.1&-0.6&9.1&1.8&1.2&-10.8\\  \hspace{-3mm}  \cbox{blue!70!violet}  Energy  Materials &0.6&0.4&0.6&1.1&3.2&2.2&3.1&5.4\\  
\end{minipage}\hspace{7mm}
\begin{minipage}{0.37\textwidth}
\normalsize \textbf{Industrial Production}\\
\footnotesize{\textit{index, 2012=100}}\\
\noindent \hspace*{-2mm} \begin{tikzpicture}
	\begin{axis}[\bbar{y}{0}, \shdateaxisticks ytick={60, 80, 100}, 
	enlarge y limits={0.05}, enlarge x limits={0}, legend cell align={left},
	yticklabel style={text width=1.0em}, xticklabel={`\short{\year}}, 
	clip=false, height=5.4cm, width=6.8cm,
	legend style={fill=white, legend columns=1, at={(1.02, 0.33)}}]
	\rbars
	\thinline{red}{date}{Manufacturing}{data/indpro.csv}
	\thinline{blue!90!black}{date}{Total index}{data/indpro.csv}
	\legend{Manufacturing, Total Index};
\end{axis}
\end{tikzpicture}\\
\footnotesize{Source: Federal Reserve} \hfill \tbllink{indpro.csv}
\end{minipage}
\vspace{1mm} 

\begin{minipage}{0.76\textwidth}
\normalsize \textbf{Industrial Production Growth}\\
\footnotesize{\textit{One-year growth \hspace{34mm} contribution to total \hspace{12mm} rate, percent}\\
\noindent \rowcolors{1}{}{black!5} \setlength{\tabcolsep}{3.1pt} \color{black!90}
		{\renewcommand{\arraystretch}{1.54}
		 \begin{tabular}{p{2mm} p{44mm} R{7.3mm} R{7.3mm} R{7.3mm} R{7.3mm} | R{8mm} R{7.3mm} 
		   R{7.3mm} R{7.3mm} }
			 \input data/indpro.tex \hline
		\end{tabular}}\vspace{-2mm}
		
\footnotesize{Source: Federal Reserve}}
\end{minipage}
\newpage
\begin{minipage}{0.76\textwidth}
\vspace{3mm}
\small Market group data show the lack of growth in the production of consumer goods, equipment, and nonindustrial supplies over the past decade.
\vspace{1mm}

\normalsize \textbf{Industrial Production Growth, Market Group}\\
\footnotesize{\textit{percentage point contribution to one-year growth}}\\\
\hspace*{-2mm} \begin{tikzpicture}
	\begin{axis}[\bbar{y}{0}, \dateaxisticks ytick={-15, -10, -5, 0, 5, 10, 15},
		xticklabel={`\short{\year}}, clip=false, height=5.6cm,
		yticklabel style={text width=1.5em}, legend cell align={left},
		legend style={legend columns=1, at={(0.53, 0.4)}}]
	\rbars
	\sbar{violet!60!black}{date}{Consumer goods}{data/indprogr.csv}
	\sbar{magenta}{date}{ENS}{data/indprogr.csv}
	\sbar{orange!70!yellow}{date}{Materials}{data/indprogr.csv}
	\legend{Consumer Goods, Equipment \& Nonindustrial Supplies, Materials};
	\end{axis}
\end{tikzpicture}\\
\footnotesize{Source: Federal Reserve} \hfill \tbllink{indprogr.csv}
\vspace{3mm}

\small Industry group data show a long-term change in the composition of new industrial activity, towards mining and away from manufacturing.
\vspace{1mm}

\normalsize \textbf{Industrial Production Growth, Industry Group}\\
\footnotesize{\textit{percentage point contribution to one-year growth}}\\
\hspace*{-2mm} \begin{tikzpicture}
	\begin{axis}[\bbar{y}{0}, \dateaxisticks ytick={-15, -10, -5, 0, 5, 10, 15},
		xticklabel={`\short{\year}}, clip=false, height=5.6cm,
		yticklabel style={text width=1.5em}, legend cell align={left},
		legend style={legend columns=1, at={(0.43, 0.44)}}]
	\rbars
	\sbar{blue!60!black}{date}{Durable manufacturing}{data/indprogr.csv}
	\sbar{blue!20!cyan!80!white}{date}{Nondurable manufacturing}{data/indprogr.csv}
	\sbar{orange!20!yellow}{date}{Mining}{data/indprogr.csv}
	\sbar{green!80!blue}{date}{Electric and gas utilities}{data/indprogr.csv}
	\legend{Durable Manufacturing, Nondurable Manufacturing, Mining, Utilities};
	\end{axis}
\end{tikzpicture}\\
\footnotesize{Source: Federal Reserve} \hfill \tbllink{indprogr.csv}
\vspace{3mm}

\small Production slowed across industry and market groups during the pandemic. Monthly data are shown below.
\vspace{1mm}

\normalsize \textbf{Recent Data in Detail}
\vspace{1mm}

\normalsize \hspace{14mm} Market Group \hspace{38mm} Industry Group \footnotesize \\
\noindent \hspace*{-2mm} \begin{tikzpicture} 
	\begin{axis}[\bbar{y}{0}, \ltdateaxisticks max space between ticks=20,
		clip=false, width=6.5cm, height=5.5cm, yticklabel style={text width=1.28em}, 
		yticklabel style={text width=1.5em}, xticklabel={`\short{\year}}, ymin=-18, ymax=18,
		minor xtick={}, enlarge y limits=0.02,  enlarge x limits={0.01}] 
	\rbar
	\ctsbar{violet!60!black}{date}{Consumer goods}{data/indprogr_rec.csv}{1.6pt} 
	\ctsbar{magenta}{date}{ENS}{data/indprogr_rec.csv}{1.6pt} 
	\ctsbar{orange!70!yellow}{date}{Materials}{data/indprogr_rec.csv}{1.6pt} 
	\end{axis}
\end{tikzpicture} \hfill 
\begin{tikzpicture} \begin{axis}[\ltdateaxisticks 
	clip=false, width=6.5cm, max space between ticks=20,	 
	height=5.5cm, yticklabel style={text width=1.28em}, 
	yticklabel style={text width=1.5em}, xticklabel={`\short{\year}}, ymin=-18, ymax=18,
	minor xtick={}, enlarge y limits=0.02,  enlarge x limits={0.01}] 
	\rbar
	\ctsbar{blue!60!black}{date}{Durable manufacturing}{data/indprogr_rec.csv}{1.6pt} 
	\ctsbar{blue!20!cyan!80!white}{date}{Nondurable manufacturing}{data/indprogr_rec.csv}{1.6pt} 
	\ctsbar{orange!20!yellow}{date}{Mining}{data/indprogr_rec.csv}{1.6pt} 
	\ctsbar{green!80!blue}{date}{Electric and gas utilities}{data/indprogr_rec.csv}{1.6pt}
	\end{axis} 
\end{tikzpicture}\\
\footnotesize{Source: Federal Reserve} \hfill \tbllink{indprogr_rec.csv}
\end{minipage}
\newpage
\begin{minipage}{0.32\textwidth}
\small \input{text/ip_comp.txt}
\end{minipage} \hspace{6mm}
\begin{minipage}{0.44\textwidth}
\normalsize \textbf{Industrial Production and Capacity}\\
\footnotesize{\textit{As of \input{text/ip_ind_ldate.txt}\unskip, percent change since February 2020}}\\
\hspace*{-4mm} \begin{tikzpicture}
  	\begin{axis}[\barplotnogrid axis y line=left, \barylab{3.3cm}{1.5ex}
    	width=5.6cm, bar width=1.8ex, height=9.5cm,
    	enlarge y limits={abs=0.32cm}, enlarge x limits=0.22, \bbar{x}{0},
		yticklabels from table={\ip}{name}, 
		yticklabel style={font=\scriptsize}, nodes near coords, 
		nodes near coords style={/pgf/number format/.cd, fixed zerofill,
			precision=1, assume math mode},
		legend style={text=black!70, at={(0,1.05)}, anchor=north, legend columns=-1, 
				fill=none, draw=none,
		        /tikz/every even column/.append style={column sep=0.4cm}}]
  	\addplot[fill=cyan!90!blue, draw=none] 
  		table [y expr=-\coordindex, x index=2] {\ip};
  	\addplot[fill=green!60!lime, draw=none] 
  		table [y expr=-\coordindex, x index=1] {\ip};
 	\legend{Capacity, Production}
  	\end{axis}
\end{tikzpicture}\\
\footnotesize{Source: Federal Reserve} \hfill \tbllink{ip_comp.csv}
\end{minipage}
\vspace{2mm}

\begin{minipage}{0.76\textwidth}
\small The Federal Reserve's monthly industrial production \href{https://www.federalreserve.gov/releases/g17/}{report} also measures the economy's total industrial capacity. The extent to which the economy is using its industrial capacity is called \textbf{capacity utilization}, and calculated as industrial production as a share of total industrial capacity. Long-term, capacity utilization has fallen as many US factories and industrial production facilities closed.

\input{text/tcu.txt}
\vspace{1mm}

\normalsize \textbf{Capacity Utilization}\\
\footnotesize{\textit{industrial production as a share of total capacity, percent}}\\
\hspace*{-2mm} \begin{tikzpicture}
	\begin{axis}[\bbar{y}{0}, \dateaxisticks width=12.2cm, 
		yticklabel style={text width=1.2em}, max space between ticks=30,	
		height=6.6cm, clip=false, xticklabel={`\short{\year}}]
	\rbars
	\stdline{blue!80!black}{date}{Total index}{data/tcu.csv}
	\thinline{blue!40!cyan}{date}{Manufacturing}{data/tcu.csv}
	\input{text/tcu_nodes.txt}
	\absnode{{2011-01-01}}{73}{\scriptsize \color{blue!40!cyan}\textbf{Manufacturing}}
	\absnode{{2009-08-01}}{79.5}{\scriptsize \color{blue!80!black}\textbf{Total index}}
	\end{axis}
\end{tikzpicture}\\
\footnotesize{Source: Census Bureau} \hfill \tbllink{tcu.csv}
\end{minipage}
\newpage
\subsubsection*{Energy Production and Use}
\begin{minipage}{0.76\textwidth}
\small The Energy Information Administration \href{https://www.eia.gov/dnav/pet/pet_crd_crpdn_adc_mbblpd_m.htm}{report} a large increase in US \textbf{crude oil production}, from around five million barrels per day in 2007 to nearly 13 million barrels per day at the end of 2019. Much of the increase comes from Texas, New Mexico, North Dakota, and Colorado. \input{text/oil_prod.txt}
\vspace{1mm}

\normalsize \textbf{Crude Oil Production}\\
\footnotesize{\textit{millions of barrels per day}}\\
\hspace*{-2mm} \begin{tikzpicture}
	\begin{axis}[\dateaxisticks ymin=0.6, enlarge x limits={0.01},
		width=13.0cm, height=6.8cm, ytick={0, 5, 10}, xticklabel={`\short{\year}},
		legend style={at={(0.82, 1)}, fill=white}]
	\rbars
	\abar{teal!90!green!96!black}{date}{ND_NM_CO}{data/oil_prod.csv}
	\abar{green!21!yellow!71!white}{date}{TX}{data/oil_prod.csv}
	\abar{blue!46!cyan!88!white}{date}{OTH}{data/oil_prod.csv}
	\abar{blue!47!violet}{date}{GM}{data/oil_prod.csv}
	\legend{ND + NM + CO, Texas, All Other, Gulf of Mexico};
	\end{axis}
\end{tikzpicture}\\
\footnotesize{Source: Energy Information Administration} \hfill \tbllink{oil_prod.csv}
\vspace{8mm}

\small \input{text/elec_prod.txt}
\end{minipage}
\vspace{2mm}

\begin{minipage}{0.6\textwidth}
\normalsize \textbf{Electricity Generation, Population, and Production}\\
\footnotesize{\textit{cumulative percent change since 2011}}\\
\hspace*{-3mm} \begin{tikzpicture}
	\begin{axis}[\bbar{y}{0}, \tydateaxisticks width=9.7cm, height=7.9cm,
		xticklabel={`\short{\year}}, clip=false, 
		yticklabel style={text width=1.6em}]
	\rbar
	\thickline{green!60!blue}{date}{gdp}{data/elec.csv}
	\thickline{blue!90!black}{date}{pop}{data/elec.csv}
	\thickline{red}{date}{elec}{data/elec.csv}
	\thickline{orange}{date}{elecpop}{data/elec.csv}
	\thickline{cyan}{date}{elecgdp}{data/elec.csv}
	\input{text/elec_nodes.txt}
	\absnode{{2016-01-01}}{16}{\footnotesize \color{green!60!blue}\textbf{GDP}}
	\absnode{{2017-10-01}}{10}{\footnotesize \color{blue!90!black}\textbf{Population}}
	\absnode{{2020-07-01}}{4}{\footnotesize \color{red}\textbf{Electricity}}
	\absnode{{2018-01-01}}{-6}{\footnotesize \color{orange}\textbf{Electricity}\\ \footnotesize \color{orange}\textbf{per capita}}
	\absnode{{2016-02-01}}{-16}{\footnotesize \color{cyan}\textbf{Electricity / GDP}}	
	\end{axis}
\end{tikzpicture}\\
\footnotesize{Source: EIA, BEA} \hfill \tbllink{elec.csv}
\end{minipage}
\newpage
\begin{minipage}{0.76\textwidth}
\small \input{text/elec_prod2.txt}
\end{minipage}
\vspace{2mm}

\normalsize \textbf{Electricity Generation by Source \hspace{31mm} Renewable Sources}\\
\footnotesize{\textit{billion kilowatt hours, 12-month moving sum}}\\
\hspace*{-4mm} \begin{tikzpicture}
	\begin{axis}[\dateaxisticks ymin=118, enlarge x limits={0.01},
		width=8.5cm, height=6.6cm, ytick={0, 1000, 2000, 3000, 4000, 5000},
		xticklabel={`\short{\year}}, yticklabel style={text width=2.2em},
		legend cell align={left}, legend style={legend columns=3, at={(1, 1.15)}}]
	\recbars
	\abar{brown}{date}{Coal}{data/elec_prod.csv}
	\abar{blue!40!cyan!60!white}{date}{Natural Gas}{data/elec_prod.csv}
	\abar{red!80!blue!70!white}{date}{Nuclear}{data/elec_prod.csv}
	\abar{green!75!blue}{date}{Renewables}{data/elec_prod.csv}
	\abar{orange!90!white}{date}{Petroleum and other}{data/elec_prod.csv}
	\legend{Coal, Natural Gas, Nuclear, Renewables, Petroleum \& Oth.};
	\end{axis}
\end{tikzpicture} \hfill 
\begin{tikzpicture}
	\begin{axis}[\dateaxisticks ymin=24, enlarge x limits={0.01},
		width=8.4cm, height=6.6cm, 
		ytick={0, 100, 200, 300, 400, 500, 600, 700, 800, 1000, 1100},
		xticklabel={`\short{\year}}, yticklabel style={text width=1.6em},
		legend cell align={left}, legend style={legend columns=3, at={(0.98, 1.15)}}]
	\recbars
	\abar{cyan}{date}{Hydroelectric}{data/elec_prod.csv}
	\abar{brown!92!black}{date}{Biomass}{data/elec_prod.csv}
	\abar{orange}{date}{Geothermal}{data/elec_prod.csv}
	\abar{green!65!blue}{date}{Wind}{data/elec_prod.csv}vv
	\abar{yellow!70!orange}{date}{Solar}{data/elec_prod.csv}
	\legend{Hydroelectric, Biomass, Geothermal, Wind, Solar};
	\end{axis}
\end{tikzpicture}\\
\footnotesize{Source: Energy Information Administration} \hfill \tbllink{elec_prod.csv}
\vspace{1mm}

\begin{minipage}{0.76\textwidth}
\small \input{text/elec_prod3.txt}
\vspace{3mm}

\small The Energy Information Administration \href{https://www.eia.gov/electricity/monthly/}{report} the \textbf{retail sales of electricity} to each major sector. These data show a decrease in electricity sales to commercial and industrial sectors during the pandemic, offset in part by increased electricity sales to the residential sector. 
\end{minipage}

\begin{minipage}{0.415\textwidth}
\normalsize \textbf{Electricity Retail Sales by Sector}\\
\footnotesize{\textit{billion kilowatt hours, 12-month moving sum}}\\
\hspace*{-2mm} \begin{tikzpicture}
	\begin{axis}[\dateaxisticks ytick={800, 1000, 1200, 1400, 1600}, width=6.5cm, 
		height=5.3cm, clip=false,
		xticklabel={`\short{\year}}, yticklabel style={text width=2.2em}]
	\recbars
	\thickline{green!90!black}{date}{Residential}{data/elec_sales.csv}
	\thickline{blue!80!cyan}{date}{Commercial}{data/elec_sales.csv}
	\thickline{black!35!white}{date}{Industrial}{data/elec_sales.csv}
	\input{text/elec_ret_nodes.txt}
	\absnode{{2011-01-01}}{1325}{\footnotesize \color{blue!80!cyan}\textbf{Commercial}}
	\absnode{{2012-01-01}}{1500}{\footnotesize \color{green!90!black}\textbf{Residential}}
	\absnode{{2014-01-01}}{975}{\footnotesize \color{black!35!white}\textbf{Industrial}}
	\end{axis}
\end{tikzpicture}\\
\footnotesize{Source: Energy Information Administration} \hfill \tbllink{elec_sales.csv}
\end{minipage} \hspace{5mm}
\begin{minipage}{0.3\textwidth}
\small \input{text/elec_ret.txt}
\end{minipage}
\newpage
\hypertarget{busrs}{}
\begin{minipage}{0.76\textwidth}
\subsection*{Retail Sales}
\small The Census Bureau \href{https://www.census.gov/retail/index.html}{report} the monthly sales of retail businesses, restaurants, and bars. These retail trade figures can be a useful economic indicator. Retail trade includes brick and mortar stores as well as e-commerce and other nonstore sales to the general public. 

\input{text/rs_sh.txt} During the first two months of the US COVID-19 pandemic, retail sales were a smaller portion of overall economic activity, as many businesses were closed. Since the initial reopening, retail sales have comprised a larger share of economic activity, in part as other activities, like transportation, have recovered less. 

Retail and food service sales, \textbf{adjusted for population growth and inflation}, provide additional context on economic developments. \input{text/rs_pc.txt}
\vspace{0.5mm}

\normalsize \textbf{Retail and Food Services Sales}\\
\footnotesize{\textit{share of aggregate measure, percent \hspace{17mm} real*, per capita, monthly, \input{text/rs_dt.txt}\unskip \ dollars}}\\
\hspace*{-2mm} \begin{tikzpicture}
	\begin{axis}[\bbar{y}{0}, \shdateaxisticks 
		enlarge y limits={0.12}, yticklabel style={text width=1.2em}, 
		xticklabel={`\short{\year}}, height=4.8cm, width=5.7cm, clip=false]
	\rebars
	\stdline{blue}{date}{GDP}{data/rs_sh.csv}
	\stdline{cyan}{date}{CE}{data/rs_sh.csv}
	\stdline{magenta!80!purple}{date}{DPI}{data/rs_sh.csv}
	\input{text/rs_sh_nodes.txt}
	\absnode{{1995-01-01}}{36.5}{\footnotesize \color{blue}\textbf{GDP}}
	\absnode{{2012-01-01}}{37.5}{\footnotesize \color{magenta!80!purple}\textbf{DPI}}
	\absnode{{2012-01-01}}{48}{\footnotesize \color{cyan}\textbf{PCE}}
	\end{axis}
\end{tikzpicture} \hfill \begin{tikzpicture}
	\begin{axis}[\bbar{y}{0}, \shdateaxisticks ymin=180,
		enlarge y limits={0.12}, yticklabel style={text width=2.2em}, 
		max space between ticks=20,	 
		xticklabel={`\short{\year}}, height=4.8cm, width=5.7cm, clip=false]
	\rebars
	\stdline{violet!60!blue}{date}{RealPC_Tot}{data/rs_sh.csv}
	\thickline{green!80!olive}{date}{RealPC_ExAutoGas}{data/rs_sh.csv}
	\input{text/rs_pc_nodes.txt}
	\absnode{{1993-01-01}}{800}{\scriptsize \color{green!80!olive}\textbf{Ex. Auto}\\\scriptsize \color{green!80!olive}\textbf{\& Gas}}
	\absnode{{2010-01-01}}{1840}{\scriptsize \color{violet!60!blue}\textbf{Total}}
	\absnode{{1992-01-01}}{300}{\scriptsize \color{black!70}*Adjusted by PCE price index}
	\end{axis}
\end{tikzpicture}\\
\footnotesize{Source: Bureau of Labor Statistics} \hfill \tbllink{rs_sh.csv}
\vspace{3mm}

\small \input{text/rs_gr.txt}
\vspace{0.5mm}

\normalsize \textbf{Retail and Food Services Sales Growth}\\
\footnotesize{\textit{annual growth, percent, 3-month moving average}}
\vspace{28mm}

\hspace{4mm} \begin{tikzpicture}[overlay]
	\begin{axis}[\bbar{y}{0}, \dateaxisticks ytick={-20, -10, 0, 10, 20, 30, 40}, 
		xticklabel={`\short{\year}}, yticklabel style={text width=1.4em}, ymax=28,
		width=12.1cm, height=4.8cm, clip=false]
	\rebars
	\stdline{blue!70!black}{date}{454}{data/marts.csv}
	\thickline{green!60!black}{date}{44X72}{data/marts.csv}
	\input{text/rs_nodes.txt}
	\absnode{{2017-12-01}}{-2}{\footnotesize \color{green!60!black}\textbf{Total}}
	\absnode{{2012-10-01}}{16}{\footnotesize \color{blue!70!black}\textbf{Nonstore}}
	\end{axis}
\end{tikzpicture}

\footnotesize{Source: Census Bureau} \hfill \tbllink{marts.csv}
\end{minipage}
\newpage
\begin{minipage}{0.76\textwidth}
\small Since 1992, the share of after-tax income spent at different \textbf{kinds of businesses} has diverged wildly. In large part, this is due to the growth of e-commerce, with online sales replacing brick and mortar sales. However, there have also been shifts in other consumer preferences and relative prices. 

\input{text/rs_kb.txt} 

Some sales categories were boosted by the housing bubble during the 2000s and its associated wealth effects, then fell sharply following the collapse of the bubble. Building and garden equipment, furniture and home furnishings stores, and motor vehicle sales all claimed a larger share of income during the 1990s and 2000s than during the 2010s. Meanwhile, food service and drinking places and health and personal care stores received a relatively stable share of income from 2000 until the COVID-19 pandemic, which hit restaurants and bars particularly hard. 

Lastly, some categories are more affected by changes in relative prices. Sales at gasoline stations, for example, move with gasoline prices. Likewise, an increase in building material prices during the pandemic partially boosted the share of income spent at building and garden equipment stores. 
\end{minipage}
\vspace{2mm}

\begin{minipage}{0.8\textwidth}
\normalsize \textbf{Retail Sales by Kind of Business}\\
\footnotesize{\textit{share of disposable personal income, \input{text/rs_dates.txt}}}
\vspace{-1mm}

\hspace*{-2mm} \begin{tikzpicture}
\input{text/rs_ind.txt}
\end{tikzpicture}\\
\footnotesize{Source: Census Bureau, Bureau of Economic Analysis} \hfill \tbllink{rs_ind.csv}
\end{minipage}
\newpage
\begin{minipage}{0.76\textwidth}
\section*{Government}
\hypertarget{gov}{}
\small Public institutions are collectively referred to as the \textit{public-sector} or the \textit{government}. In the United States, the government has the authority to spend, tax, and create money, as well as to regulate economic and financial activity. The government also enforces and determines the ownership of property. These activities are all extremely important to production and distribution in the economy.

This chartbook section covers various government statistics, including contributions to current economic activity, receipts and expenditures, assets and liabilities, and government jobs.
\subsection*{Current Economic Activity}
\small The government has several roles in society and there are multiple ways to interpret its contribution to current economic activity. The contribution to overall economic activity from the government sector can be measured as: 1) the gross output of the sector minus intermediate inputs used in production (value added); 2) the government income payments to people and on behalf of people and taxes and social insurance contributions; or 3) the sum of government expenditures on final goods and services, including investment. 

The value added in production by the government sector is composed primarily of the compensation of government employees. These employees provide all types of services: education, health care, transportation, utilities, sanitation, etc. Government value added also includes a residual term called \href{https://www.bea.gov/help/glossary/gross-operating-surplus}{gross operating surplus}, much of which reflects money spent replacing and maintaining government fixed assets. 
\end{minipage}
\vspace{2mm}

\begin{minipage}{0.38\textwidth}
\normalsize \textbf{Value Added in Domestic Production}\\
\footnotesize{\textit{per capita, thousands of \input{text/gdpva_ltdt.txt}dollars}}\\
\hspace*{-2mm} \begin{tikzpicture}
	\begin{axis}[\bbar{y}{0}, \shdateaxisticks ytick={0, 2, 4, 6},
		ymin=0.3, 
		yticklabel style={text width=1.0em}, 
		xticklabel={`\short{\year}}, clip=false, height=5.6cm, width=7.0cm]
	\rebars
	\stdnode{1.7cm}{1.8cm}{\footnotesize \color{blue!70!white}Federal}
	\stdline{blue!70!white}{date}{Federal}{data/gdpva_pc.csv}
	\stdnode{3.9cm}{3.3cm}{\footnotesize \color{red}State and Local}
	\stdline{red}{date}{State and local}{data/gdpva_pc.csv}
	\end{axis}
\end{tikzpicture}\\
\footnotesize{Source: Bureau of Economic Analysis} \hfill \tbllink{gdpva_pc.csv}
\end{minipage} \hspace{6mm}
\begin{minipage}{0.33\textwidth}
\small \input{text/govva.txt}
\end{minipage}
\newpage

\begin{minipage}{0.76\textwidth}
\small Payments from the government to people include both the wages and salaries of government workers and transfer payments, also referred to as government social benefits or welfare. Welfare is critically important to society, yet the word is stigmatized in the US. It's worth noting that a huge portion of US welfare payments are not cash transfers to people but payments to businesses on behalf of people. For example, the US overpays for prescription drugs and medical services on behalf of people. Those over-payments benefit shareholders and doctors. The US also structures certain welfare payments as complex phased-in tax credits, paid the following year, which have very high private administrative costs (around ten percent goes to tax preparation services) and exclude the poorest people.

Over the past thirty years, welfare payments, including those made to businesses, have mostly kept pace with consumer spending, while tax collection lagged behind income growth. Increased welfare payments on a per capita basis reflect price protections for the medical industry, as well as a larger share of the population receiving social security benefits. During the COVID-19 pandemic, consumer spending on services collapsed as businesses were closed, while government transfer payments to people hit record highs. As a result, estimates show seven million people taken out of poverty.
\end{minipage}
\vspace{1mm}

\begin{minipage}{0.38\textwidth}
\normalsize \textbf{Personal Income and Outlays}\\
\footnotesize{\textit{per capita, thousands of \input{text/gdp_ltdt.txt}dollars}}\\
\hspace*{-2mm} \begin{tikzpicture}
	\begin{axis}[\bbar{y}{0}, \shdateaxisticks ytick={0, 5, 10, 15, 20},
		ymin=0.8, 
		yticklabel style={text width=1.0em}, 
		xticklabel={`\short{\year}}, clip=false, height=5.6cm, width=7.0cm]
	\rbars
	\stdnode{4.4cm}{3.1cm}{\footnotesize \color{violet}Welfare}
	\stdline{violet}{date}{Welfare}{data/govpi.csv}
	\stdnode{1.3cm}{3.1cm}{\footnotesize \color{cyan!80!blue}Tax \& Social\\ \footnotesize \color{cyan!80!blue} Insurance\\ \footnotesize \color{cyan!80!blue} Contributions}
	\stdline{cyan!80!blue}{date}{TaxSI}{data/govpi.csv}
	\stdnode{1.2cm}{0.65cm}{\footnotesize \color{orange}Wages and \\ \footnotesize \color{orange}Salaries}
	\stdline{orange}{date}{WandS}{data/govpi.csv}
	\end{axis}
\end{tikzpicture}\\
\footnotesize{Source: Bureau of Economic Analysis} \hfill \tbllink{govpi.csv}
\end{minipage}\hspace{5.5mm}
\begin{minipage}{0.34\textwidth}
\small \input{text/govpi.txt}
\end{minipage}
\vspace{3mm}

\begin{minipage}{0.76\textwidth}
\small Another approach to calculating the government sector effect on current economic activity is to add up spending on final goods and services. Government consumption and investment tends to be more stable than consumer spending or private investment, and thus tends to rise as a share of economic activity during recessions. This category does not include government transfer payments, which mostly become consumer spending.
\end{minipage}
\vspace{1mm}

\begin{minipage}{0.38\textwidth}
\normalsize \textbf{Consumption and Investment}\\
\footnotesize{\textit{share of GDP, percent}}\\
\hspace*{-2mm} \begin{tikzpicture}
	\begin{axis}[\bbar{y}{0}, \shdateaxisticks ytick={0, 5, 10, 15, 20},
		ymin=0.6, 
		yticklabel style={text width=1.0em}, 
		xticklabel={`\short{\year}}, clip=false, height=5.8cm, width=7.0cm]
	\rbars
	\stdnode{1.3cm}{0.6cm}{\footnotesize \color{green!85!black}Federal\\ \footnotesize \color{green!85!black} Non-Defense}
	\stdline{green!85!black}{date}{A825RC}{data/govci.csv}
	\stdnode{4.3cm}{3.1cm}{\footnotesize \color{purple!70!magenta}State \& Local}
	\stdline{purple!70!magenta}{date}{A829RC}{data/govci.csv}
	\stdnode{1.2cm}{2.0cm}{\footnotesize \color{blue!60!black}Defense}
	\stdline{blue!60!black}{date}{A824RC}{data/govci.csv}
	\end{axis}
\end{tikzpicture}\\
\footnotesize{Source: Bureau of Economic Analysis} \hfill \tbllink{govci.csv}
\end{minipage}\hspace{5.5mm}
\begin{minipage}{0.34\textwidth}
\small \input{text/govci.txt}
\end{minipage}
\newpage
\begin{minipage}{0.76\textwidth}
\subsubsection*{Government Consumption and Investment}
\small &   2023  Q3 & `23  Q2 & `23  Q1 & `22  Q3 & `21  Q3 &1-year&10-year&30-year\\  Consolidated  Government  Total &0.99&0.57&0.82&0.49&-0.26&0.82&0.30&0.25\\  \hspace{2mm}Federal  Total &0.45&0.07&0.33&0.07&-0.48&0.36&0.11&0.10\\  \cbox{blue!60!black}National  Defense &0.30&0.08&0.07&-0.01&-0.12&0.18&0.02&0.04\\  \hspace{6mm}Consumption  Expenditures &0.23&0.02&0.16&0.05&-0.15&0.12&0.01&0.02\\  \hspace{6mm}Gross  Investment &0.07&0.06&-0.09&-0.06&0.03&0.06&0.02&0.02\\  \cbox{green!85!black}Federal  Non-Defense &0.15&-0.01&0.26&0.08&-0.36&0.18&0.08&0.07\\  \hspace{6mm}Consumption  Expenditures &0.13&-0.07&0.17&0.05&-0.41&0.12&0.05&0.04\\  \hspace{6mm}Gross  Investment &0.02&0.05&0.08&0.03&0.05&0.06&0.04&0.02\\  \hspace{-2mm}\cbox{purple!70!magenta}State  \&  Local  Total &0.53&0.50&0.49&0.41&0.22&0.46&0.19&0.15\\  \hspace{4mm}Consumption  Expenditures &0.23&0.16&0.29&0.20&0.49&0.20&0.15&0.12\\  \hspace{4mm}Gross  Investment &0.30&0.34&0.21&0.22&-0.27&0.26&0.04&0.03\\   
\vspace{2mm}

\normalsize \textbf{Government Consumption and Investment}\\
\footnotesize{\textit{percentage point contribution to real GDP growth, one-year moving average}}\\
\hspace*{-2mm} \begin{tikzpicture}
	\begin{axis}[\bbar{y}{0}, \dateaxisticks 
		xticklabel={`\short{\year}}, clip=false, height=5.2cm, width=12.7cm,
		yticklabel style={text width=1.6em},
		legend style={at={(0.67, 1.09)}, fill=white}]
	\rbars
	\sbar{blue!60!black}{date}{A824RY}{data/gov.csv}
	\sbar{green!85!black}{date}{A825RY}{data/gov.csv}
	\sbar{purple!70!magenta}{date}{A829RY}{data/gov.csv}
	\legend{Defense, Federal Non-Defense, State and Local};
	\end{axis}
\end{tikzpicture}\\
\footnotesize{Source: Bureau of Economic Analysis} \hfill \tbllink{gov.csv}
\vspace{2mm}

\normalsize \textbf{Government Consumption and Investment}\\
\footnotesize{\textit{percentage point contribution to real GDP growth \hspace{35mm} moving averages}}\\ 
\rowcolors{1}{}{black!5} \setlength{\tabcolsep}{3.2pt} \color{black!90}
	{\renewcommand{\arraystretch}{1.7}
		\begin{tabular}{p{42mm} R{6.7mm} R{6.7mm} R{6.7mm} R{6.7mm} R{6.7mm} 
		   R{7.8mm} R{6.7mm} R{6.7mm} }
			 \input data/gov.tex \hline
		\end{tabular}}\vspace{-2mm}
		
\footnotesize{Source: Bureau of Economic Analysis}
\end{minipage}
\newpage
\begin{minipage}{0.76\textwidth}
\small Government current expenditures include consumption and investment as well as transfers such as government social benefits to persons. Government spending provides services and income to people. Government current receipts come primarily from taxes. When government expenditures exceed receipts, it is referred to as a \textit{government deficit}, and corresponds to a private sector surplus. A large government deficit, relative to GDP, means the government is increasing current household income and corporate profits.

\input{text/fedgdp.txt} 

\input{text/slggdp.txt} 
\vspace{1mm}

\normalsize \textbf{Receipts and Expenditures as Share of GDP}

\normalsize \hspace{20mm} Federal \hspace{44mm} State and Local \footnotesize \\
\vspace{1mm}
\hspace*{-2mm} \begin{tikzpicture}
	\begin{axis}[\bbar{y}{0}, \shdateaxisticks 
		clip=false, width=6.65cm, height=5.0cm, xticklabel={`\short{\year}}, 
		yticklabel style={text width=1.0em}, legend style={at={(0.95, 1.13)}}]
	\rbars
	\stdline{blue!90!cyan}{date}{W005RC}{data/fedgdp.csv}
	\stdline{red}{date}{W013RC}{data/fedgdp.csv}
	\legend{Receipts, Expenditures};
	\end{axis}
\end{tikzpicture}
\hfill
\begin{tikzpicture}
	\begin{axis}[\bbar{y}{0}, \shdateaxisticks 
		clip=false, width=6.65cm, height=5.0cm, xticklabel={`\short{\year}}, 
		yticklabel style={text width=1.0em}]
	\rbars
	\stdline{blue!90!cyan}{date}{W023RC}{data/slggdp.csv}
	\stdline{red}{date}{W024RC}{data/slggdp.csv}
	\end{axis}
\end{tikzpicture}\\
\footnotesize{Source: Bureau of Economic Analysis} \hfill \tbllink{fedgdp.csv} \ \tbllink{slggdp.csv}
\vspace{5mm}

\small \input{text/tmb_rec.txt}
\end{minipage}  
\vspace{1mm}

\begin{minipage}{0.36\textwidth}
\normalsize \textbf{Federal Government Receipts}\\
\footnotesize{\textit{trillions of USD, one-year moving sum}}\\
\hspace*{-2mm} \begin{tikzpicture}
	\begin{axis}[\bbar{y}{0}, \ltdateaxisticks xmin={2016-01-01},
		clip=false, xticklabel={`\short{\year}}, height=6.2cm, width=6.4cm]
	\rbar
	\stdline{blue!70!black}{date}{Individual Income Taxes}{data/tmb_rec.csv}
	\input{text/tmb_rec_iit.txt}
	\stdline{blue!40!cyan}{date}{Other}{data/tmb_rec.csv}
	\input{text/tmb_rec_oth.txt}
	\stdnode{2.3cm}{1.1cm}{\footnotesize \color{blue!70!black}Individual\\\color{blue!70!black}Income\\\color{blue!70!black}Taxes}
	\stdnode{1.6cm}{2.1cm}{\footnotesize \color{blue!40!cyan}Other}
	\end{axis}
\end{tikzpicture}\\
\footnotesize{Source: Treasury} \hfill \tbllink{tmb_rec.csv}
\end{minipage}
\newpage
\begin{minipage}{0.76\textwidth}
\subsubsection*{Composition of Federal Government Spending}
\small Over the long-term, there have been important shifts in the \textbf{composition of federal spending}. The ways federal spending varies from these long-term trends, in the short-term, are also important. 

Over the long-term, Office of Management and Budget (OMB) data \href{https://www.whitehouse.gov/omb/historical-tables/}{show} \input{text/gov_exp_comp1.txt}
\vspace{1mm}

\normalsize \textbf{Composition of Federal Government Outlays}\\
\footnotesize{\textit{share of total, percent}}
\vspace{1mm}

\begin{minipage}{0.58\textwidth}
\small {Since 1990, by Budget Enforcement Act Type}\\
\vspace{5.0cm}

\hspace{4mm} \begin{tikzpicture}[overlay]
	\begin{axis}[\shdateaxisticks, enlarge y limits=0.02, enlarge x limits={0.02}, 
		yticklabel style={text width=1.5em}, ymin=0, ymax=100, xmin={1990-01-01}, 
		xmax={2020-10-01}, clip=false,  ytick={0, 50, 100}, 
		width=7.4cm, height=7.2cm, xticklabel={\year}, ]
	\rbars
	\abar{green!80!blue!68!black}{date}{National Defense}{data/gov_exp_omb.csv}
	\abar{green!78!black}{date}{Non- defense}{data/gov_exp_omb.csv}
	\abar{blue!65!black}{date}{MP Medicare}{data/gov_exp_omb.csv}	
	\abar{cyan!50!blue!95!white}{date}{MP Social Security}{data/gov_exp_omb.csv}
	\abar{cyan!95!white}{date}{Social Safety Net}{data/gov_exp_omb.csv}	
	\abar{cyan!40!white}{date}{Other}{data/gov_exp_omb.csv}	
	\abar{orange!80!yellow}{date}{Net Interest}{data/gov_exp_omb.csv}	
	\absnode{{1990-01-01}}{{98}}{\footnotesize \color{white}\textbf{Interest}}
	\absnode{{1989-05-01}}{{86.5}}{\footnotesize \color{black!90!white}\textbf{Other}}
	\absnode{{1991-01-01}}{{78}}{\footnotesize \color{white}\textbf{Social Safety Net}}
	\absnode{{1990-01-01}}{{62.5}}{\footnotesize \color{white}\textbf{Social Security}}
	\absnode{{1995-01-01}}{{44.5}}{\footnotesize \color{white}\textbf{Medicare}}
	\absnode{{1993-01-01}}{{30.5}}{\footnotesize \color{white}\textbf{Non-Defense}}
	\absnode{{1990-01-01}}{{14}}{\footnotesize \color{white}\textbf{National Defense}}
  	\node[rotate=90, anchor=west, shift={(0cm,0.5cm)}] at (axis cs:{1991-04-01},4){\scriptsize \color{green!80!blue!70!black} Discretionary};
	\input{text/gov_exp_nodes.txt}
	\end{axis}
\end{tikzpicture}

\footnotesize{Source: Office of Management and Budget} \hfill \tbllink{gov_exp_omb.csv}
\end{minipage} \hfill \begin{minipage}{0.335\textwidth}
\small {Recent Monthly, by Type*}\\
\vspace{5.0cm}

\begin{tikzpicture}[overlay]
	\begin{axis}[\ltdateaxisticks, enlarge y limits=0.02, enlarge x limits={0.01}, 
		yticklabel style={text width=1.5em}, ymin=0, ymax=100, ytick=false,
  		ymajorticks=false, clip=false,
		width=4.7cm, height=7.2cm, xticklabel={`\short{\year}}, ]
	\rbar
	\abar{green!80!blue!68!black}{date}{National Defense}{data/gov_exp_mts.csv}
	\abar{violet!50!black}{date}{Health}{data/gov_exp_mts.csv}
	\abar{blue!65!black}{date}{Medicare}{data/gov_exp_mts.csv}	
	\abar{cyan!50!blue!95!white}{date}{Social Security}{data/gov_exp_mts.csv}
	\abar{brown}{date}{Income Security}{data/gov_exp_mts.csv}	
	\abar{yellow!50!brown!55!white}{date}{Other}{data/gov_exp_mts.csv}	
	\abar{orange!80!yellow}{date}{Net Interest}{data/gov_exp_mts.csv}	
	\absnode{{2017-01-01}}{{100.9}}{\footnotesize \color{white}\textbf{Interest}}
	\absnode{{2017-01-01}}{{91}}{\footnotesize \color{black!90!white}\textbf{Other*}}
	\absnode{{2017-01-01}}{{76}}{\footnotesize \color{white}\textbf{Income Security}}
	\absnode{{2017-01-01}}{{58}}{\footnotesize \color{white}\textbf{Social Security}}
	\absnode{{2017-01-01}}{{40}}{\footnotesize \color{white}\textbf{Medicare}}
	\absnode{{2017-01-01}}{{26}}{\footnotesize \color{white}\textbf{Health}}
	\absnode{{2017-01-01}}{{11}}{\footnotesize \color{white}\textbf{National Defense}}
	\input{text/gov_exp_nodes2.txt}
	\end{axis}
\end{tikzpicture}

\footnotesize{Source: Treasury} \hfill \tbllink{gov_exp_mts.csv}
\end{minipage}

\footnotesize{* The two charts use different data sources with different categories and therefore do not match.}
\vspace{3mm}

\small The Treasury Bureau of Fiscal Service \href{https://www.fiscal.treasury.gov/reports-statements/mts/}{report} federal outlays by type on a monthly basis (see right chart above). The categories used in the Treasury monthly report are not the same as those used in the OMB data, so the two charts above should not be compared. The higher-frequency Treasury data, however, are helpful for showing short-term changes, and \textbf{recent changes in the composition of federal spending}. 

\input{text/gov_exp_comp2.txt}

\end{minipage}
\newpage
\begin{minipage}{0.76\textwidth}
\small \input{text/gov_exp_comp3.txt}

Within the ``other'' category, the biggest changes during the pandemic came from business and housing subsidies (commerce and housing credit) and transfers to state and local governments (general government). The category is described in the following table. 
\vspace{2mm}

\normalsize \textbf{Composition of Federal Government Outlays}\\
\footnotesize{\textit{share of total, percent}\\
\rowcolors{1}{}{black!5} \setlength{\tabcolsep}{1.2pt} \color{black!90}
{\renewcommand{\arraystretch}{1.55}
\hspace*{-2mm} \begin{tabular}{p{58mm} R{9.0mm} R{9.0mm} R{9.0mm} R{9.0mm} R{9.0mm} R{9.0mm}}
		\input data/gov_exp_comp.tex \hline
	\end{tabular}
}}
		
\vspace{-2mm}
\footnotesize{Source: Treasury Bureau of Fiscal Service}
\end{minipage}
\newpage
\begin{minipage}{0.76\textwidth}
\subsubsection*{Government Jobs}
\small Government workers provide public services to the population. As examples, federal government jobs include the Postal Service, state government jobs include teachers and social workers, and local governments employ firefighters and utilities workers. Additionally, government employment is traditionally a relatively-stable source of aggregate household income. Government jobs are also disproportionately likely to provide health insurance and retirement benefits.  
\end{minipage}
\vspace{1mm}

\begin{minipage}{0.38\textwidth}
\normalsize \textbf{Government Employment}\\
\footnotesize{\textit{share of age 16+ population, percent}}\\
\hspace*{-2mm} \begin{tikzpicture}
	\begin{axis}[\shdateaxisticks \bbar{y}{0}, clip=false, 
		xticklabel={`\short{\year}}, height=11.4cm, width=6.4cm, 
		yticklabel style={text width=1.0em}, ymin=0,
        enlarge y limits={0.04, upper}]
	\rbars
	\stdline{blue!50!cyan}{date}{govjobs}{data/govjobs.csv}
	\stdline{green!80!blue}{date}{fedjobs}{data/govjobs.csv}
	\stdline{orange}{date}{stjobs}{data/govjobs.csv}
	\stdline{red}{date}{locjobs}{data/govjobs.csv}
	\input{text/govjobs_node.txt}
	\input{text/fedjobs_node.txt}
	\input{text/stjobs_node.txt}
	\input{text/locjobs_node.txt}
	\absnode{{1992-01-01}}{9.7}{\footnotesize \color{blue!50!cyan}Total}
	\absnode{{1993-10-01}}{6.5}{\footnotesize \color{red}Local}
	\absnode{{1994-01-01}}{2.8}{\footnotesize \color{orange}State}
	\absnode{{2012-01-01}}{1.1}{\footnotesize \color{green!80!blue}Federal}
	\end{axis}
\end{tikzpicture}\\
\footnotesize{Source: Bureau of Labor Statistics} \hfill \tbllink{govjobs.csv}
\end{minipage}\hspace{7mm}
\begin{minipage}{0.33\textwidth}
\small \input{text/govjobs.txt}
\end{minipage}
\newpage
\begin{minipage}{0.76\textwidth}

\subsubsection*{Government Net Investment}

\small \input{text/govnetinv.txt}

\end{minipage}
\vspace{1mm}

\begin{minipage}{0.47\textwidth}

\normalsize \textbf{Government Net Investment}\\
\footnotesize{\textit{government investment less depreciation, share of GDP}}\\
\hspace*{-2mm} \begin{tikzpicture}
	\begin{axis}[\bbar{y}{0}, \shdateaxisticks ytick={0, 0.5, 1.0, 1.5, 2.0}, width=7.8cm, 
		height=5.4cm, clip=false,
		xticklabel={`\short{\year}}, yticklabel style={text width=1.5em}]
	\rbars
	\thickline{red}{date}{Value}{data/govnetinv.csv}
	\input{text/govnetinv_node.txt}
	\end{axis}
\end{tikzpicture}\\
\footnotesize{Source: Bureau of Economic Analysis} \hfill \tbllink{govnetinv.csv}
\end{minipage}
\vspace{7mm}

\begin{minipage}{0.76\textwidth}

\subsection*{Government Balance Sheets}

\subsubsection*{Government Sector Net Worth}
\hypertarget{govbs}{}

\small \input{text/govshnw.txt}

\vspace{2mm}

\normalsize \textbf{Government Share of US Wealth}\\
\footnotesize{\textit{government share of US net worth, percent}}\\
\hspace*{-2mm} \begin{tikzpicture}
	\begin{axis}[\bbar{y}{0}, \dateaxisticks ytick={10, 5, 0, -5, -10, -15}, width=12.1cm, 
		height=5.4cm, clip=false,
		xticklabel={`\short{\year}}, yticklabel style={text width=1.5em}]
	\rbars
	\thickline{green!80!blue}{date}{Gov NW}{data/govshnw.csv}
	\absnode{{2011-07-01}}{-4.5}{\footnotesize \color{green!80!blue}General Government}
	\stdline{orange!80!white}{date}{FG NW}{data/govshnw.csv}
	\absnode{{2013-01-01}}{-11.5}{\footnotesize \color{orange!80!white}Federal}
	\stdline{blue!90!black}{date}{SLG NW}{data/govshnw.csv}
	\absnode{{1992-01-01}}{6.5}{\footnotesize \color{blue!90!black}State \& Local}
	\input{text/govshnw_node.txt}
	\input{text/fgshnw_node.txt}
	\input{text/slgshnw_node.txt}
	\end{axis}
\end{tikzpicture}\\
\footnotesize{Source: Federal Reserve} \hfill \tbllink{govshnw.csv}

\end{minipage}
\newpage
\begin{minipage}{0.76\textwidth}
\subsubsection*{Liabilities}
\small \textbf{Treasury bonds are held by a mixture of investors}, including private domestic investors, overseas investors, the Federal Reserve, and government agencies and trusts. \input{text/pubdebt.txt}
\vspace{2mm}

\normalsize \textbf{Federal Debt by Holder}\\
\footnotesize{\textit{percent of GDP}}\\
\hspace*{-2mm} \begin{tikzpicture}
	\begin{axis}[\bbar{y}{0}, \dateaxisticks ytick={0, 25, 50, 75, 100, 125, 150},
		xticklabel={`\short{\year}}, clip=false, ymin=4, height=6.0cm,
		legend cell align={left}, yticklabel style={text width=1.5em},
		legend style={legend columns=2, at={(0.52, 1.0)}, fill=white}]
	\rbars
	\sbar{orange!70!white}{date}{FDHBFIN}{data/pubdebt.csv}
	\sbar{green!60!black}{date}{PD}{data/pubdebt.csv}
	\sbar{cyan!50!white}{date}{IG}{data/pubdebt.csv}
	\sbar{blue}{date}{FDHBFRBN}{data/pubdebt.csv}
	\legend{Foreign, Private Domestic, Intragovernmental, Federal Reserve};
	\end{axis}
\end{tikzpicture}\\
\footnotesize{Source: Treasury} \hfill \tbllink{pubdebt.csv}
\vspace{2mm}

\small The ratio of public debt to GDP increased during the COVID-19 response, while the typical interest income from holding public debt fell because of lower interest rates. Treasuries and other government debt securities provide a safe asset for the balance sheets domestic households and businesses, and for foreign investors.
\end{minipage}
\subsubsection*{Interest Expense}
\begin{minipage}{0.33\textwidth}
\small \input{text/fedintexp.txt}
\end{minipage} \hspace{6mm} \begin{minipage}{0.38\textwidth}
\normalsize \textbf{Federal Interest Outlays}\\
\footnotesize{\textit{share of GDP, percent}}\\
\hspace*{-2mm} \begin{tikzpicture}
	\begin{axis}[\shdateaxisticks ytick={1, 2, 3}, clip=false, ymin=1.1,
		xmin={1990-03-01},
		xticklabel={`\short{\year}}, height=4.8cm, width=6.7cm]
	\rbars
	\thickline{magenta}{date}{value}{data/fedintexp.csv}
	\input{text/fedintexp_node.txt}
	\end{axis}
\end{tikzpicture}\\
\footnotesize{Source: FRED; OMB} \hfill \tbllink{fedintexp.csv}
\end{minipage} 
\vspace{1mm}

\begin{minipage}{0.76\textwidth}
\small While debt levels are much lower for the consolidated state and local government sectors, interest rates on state and local government debt are typically higher. As a result, interest paid to investors can claim a larger share of government expenses at the state and local level.
\end{minipage}
\newpage
\begin{minipage}{0.76\textwidth}
\subsubsection*{Debt Sustainability}
\small Changes in the ratio of federal government debt to GDP can be \href{https://www.imf.org/en/Publications/TNM/Issues/2016/12/31/A-Practical-Guide-to-Public-Debt-Dynamics-Fiscal-Sustainability-and-Cyclical-Adjustment-of-23498}{decomposed} to understand how various economic forces affect the trajectory of the debt relative to our ability to service it. Specifically, \textbf{debt sustainability} is affected not only by borrowing but also by changes in real interest rates and economic growth. 

In a mechanical way, government debt is the result of the accumulation of past deficits. When the government spends more than it takes in through taxes, it borrows the difference, which adds to the debt. Importantly, some government spending is interest payments on the debt. The \textit{primary balance} measures the gap between spending excluding interest payments and revenue. Interest payments are a product of the interest rate and the existing debt. Higher real interest rates mean larger interest payments which increase deficits and, in turn, increase debt.

Federal debt is often divided by GDP as a way to capture the ability to repay the debt. The basic idea is that \textbf{a growing economy gradually erodes the burden of its debt}. As the economy grows, it is better able to produce the resources needed to repay its debt. Finally, there are often discrepancies between when borrowing occurs and when spending occurs, and the account balances at the Treasury vary over time. For example, the Treasury Secretary made more cash available to cover any potential short-term needs during the peak of the COVID-19 pandemic. Stock-flow adjustments correct for the difference between the change in liabilities (the stock) and the current federal deficit or surplus (the flow).

\input{text/debt_dynamics.txt}
\vspace{2mm}

\normalsize \textbf{Federal Government Debt Dynamics}\\
\footnotesize{\textit{contribution to change in debt to GDP ratio, percentage points}}
\vspace{5.6cm}

\hspace*{4mm} \begin{tikzpicture}[overlay]
	\begin{axis}[\bbar{y}{0}, \dateaxisticks 
		yticklabel style={text width=1.4em}, ymax=18.2,
		%ytick={-4, -2, 0, 2, 4, 6, 8},
		xticklabel={`\short{\year}}, clip=false, height=7.6cm, width=12.9cm,
		legend cell align=left, transpose legend,
		legend style={at={(0.66, 0.99)}, fill=white, legend columns=2}]
	\rbars
	\ctsbar{orange!60!yellow!86!white}{date}{e_p}{data/debt_dynamics.csv}{8.0pt}
	\ctsbar{violet!80!white}{date}{e_gdp}{data/debt_dynamics.csv}{8.0pt}
	\ctsbar{green!80!yellow!68!black}{date}{e_r}{data/debt_dynamics.csv}{8.0pt}
	\ctsbar{blue!90!cyan!78!white}{date}{sfa}{data/debt_dynamics.csv}{8.0pt}
	\addplot[only marks, draw=black, fill=white] 
		table [x=date, y=total, col sep=comma] {data/debt_dynamics.csv};
	\legend{Primary Balance, Economic Growth, Real Interest Rate, Stock-Flow Adjustment, Total};
	\end{axis}
\end{tikzpicture}

\footnotesize{Source: Author's Calculations, Federal Reserve, Bureau of Economic Analysis} \hfill \tbllink{debt_dynamics.csv} 

\end{minipage}
\newpage
\begin{minipage}{0.76\textwidth}
\subsubsection*{Assets}
\small US government assets include financial assets but are mostly comprised of the non-financial assets of state and local governments (SLG), such as buildings and equipment. Land is not included in US measures of government assets. 

\input{text/govassets.txt}
\vspace{2mm}

\normalsize \textbf{Government Assets}\\
\footnotesize{\textit{does not include land, percent of GDP}}\\
\hspace*{-2mm} \begin{tikzpicture}
	\begin{axis}[\bbar{y}{0}, \dateaxisticks ytick={0, 25, 50, 75, 100}, width=12.8cm,
		xticklabel={`\short{\year}}, clip=false, ymin=4, ymax=118, height=6.0cm,
		legend cell align={left}, yticklabel style={text width=1.5em}, 
		legend style={legend columns=2, at={(0.6, 1.04)}}]
	\rbars
	\sbar{cyan!36!white}{date}{SLG NFA}{data/govassets.csv}
	\sbar{blue!75!cyan}{date}{SLG FA}{data/govassets.csv}
	\sbar{violet!60!white}{date}{FG NFA}{data/govassets.csv}
	\sbar{blue!28!violet!90!black}{date}{FG FA}{data/govassets.csv}
	\legend{SLG Nonfinancial, SLG Financial, Federal Nonfinancial, Federal Financial};
	\end{axis}
\end{tikzpicture}\\
\footnotesize{Source: Federal Reserve} \hfill \tbllink{govassets.csv}
\end{minipage}
\newpage
\hypertarget{ext}{}
\begin{minipage}{0.76\textwidth}
\section*{International Transactions}
\small Transactions between the US and the rest of the world are recorded in the balance of payments as either current account transactions (which measure income) or capital and financial account transactions (which measure change in ownership of assets). This section discusses both types of transactions with a focus on the difference between payments from residents and payments to residents, referred to as the balance of payments. The section also covers trade in more detail and discusses trends in exchanges rates. 

\subsection*{Balance of Payments}
\small The \textbf{\href{https://www.bea.gov/help/glossary/balance-current-account-national-income-and-product-accounts}{current account balance}} captures international payments for goods and services as well as international transfers of money and net income received from ownership of foreign assets. It is the result of a set of payments to US residents, called current receipts, which are largely for exports of goods and services and returns on foreign assets, and a set of payments from US residents to the rest of the world, called current payments. Current payments are for goods and services imports, as well as transfers and payments for foreign ownership of US-based assets.

The current account balance can be decomposed into the balance on trade in goods (see\cbox{yellow!60!lime!85!white}), the balance on trade in services (see\cbox{cyan!50!white}), the balance on primary income (such as wages or income from assets [see\cbox{blue!80!cyan}]), and secondary income (such as remittances and taxes [see\cbox{green!80!blue}]). 

&   2023  Q3 & `23  Q2 & `23  Q1 & `22  Q4 & `22  Q3 & `22  Q2 &3-year&10-year\\  Current  Account  Balance &-2.97&-3.09&-3.25&-3.34&-3.60&-3.96&-3.55&-2.56\\  Current  Receipts &17.09&17.01&17.31&17.31&17.53&17.46&16.87&17.62\\  \hspace{1mm}Exports &10.98&10.94&11.43&11.41&11.87&11.92&11.19&11.83\\  \hspace{3mm}Goods &7.34&7.27&7.80&7.75&8.23&8.28&7.64&7.87\\  \hspace{5mm}Durable &4.10&4.07&4.16&4.10&4.19&4.14&4.11&4.68\\  \hspace{5mm}Non-Durable &3.24&3.21&3.63&3.66&4.04&4.14&3.53&3.19\\  \hspace{3mm}Services &3.64&3.67&3.63&3.66&3.63&3.64&3.55&3.96\\  \hspace{1mm}Income  Receipts &5.43&5.37&5.19&5.08&4.95&4.81&4.94&5.02\\  \hspace{1mm}Transfer  Receipts &0.68&0.70&0.69&0.82&0.72&0.73&0.74&0.78\\  Current  payments &20.06&20.11&20.56&20.65&21.13&21.42&20.41&20.18\\  \hspace{1mm}Imports &13.80&13.92&14.51&14.73&15.30&15.93&14.69&14.83\\  \hspace{3mm}Goods &11.26&11.30&11.78&11.98&12.52&13.19&12.11&12.15\\  \hspace{5mm}Durable &7.30&7.40&7.51&7.52&7.89&8.22&7.75&7.81\\  \hspace{5mm}Non-Durable &3.96&3.90&4.27&4.45&4.63&4.97&4.36&4.34\\  \hspace{3mm}Services &2.54&2.62&2.73&2.76&2.78&2.74&2.58&2.68\\  \hspace{1mm}Income  Payments &4.84&4.73&4.59&4.37&4.21&4.02&4.23&3.89\\  \hspace{1mm}Transfer  Payments &1.42&1.46&1.46&1.55&1.62&1.47&1.49&1.46\\  
\vspace{2mm}

\normalsize \textbf{Current Account Balance}\\
\footnotesize{\textit{balance on individual current account component, as percent of GDP}}\\
\hspace*{-2mm} \begin{tikzpicture}
	\begin{axis}[\bbar{y}{0}, \dateaxisticks ytick={-8, -6, -4, -2, 0, 2, 4}, 
		yticklabel style={text width=1.0em}, ymax=4, ymin=-8.1,
		height=5.2cm, width=12.4cm, enlarge y limits={0.01}, 
		xticklabel={`\short{\year}}, clip=false, 
		legend style={at={(0.71, 1.13)}}]
	\rbars
	\ctsbar{yellow!60!lime!85!white}{date}{GOODS}{data/cab.csv}{2.35pt}
	\ctsbar{cyan!50!white}{date}{SERVICES}{data/cab.csv}{2.35pt}
	\ctsbar{blue!80!cyan}{date}{INCOME}{data/cab.csv}{2.35pt}
	\ctsbar{green!80!blue}{date}{TRANSFERS}{data/cab.csv}{2.35pt}
	\stdline{black}{date}{A124RC}{data/cab.csv}
	\input{text/cab_node.txt}
	\legend{Goods, Services, Income, Transfers, Total};
	\end{axis}
\end{tikzpicture}\\
\footnotesize{Source: Bureau of Economic Analysis} \hfill \tbllink{cab.csv}
\vspace{2mm}

\small US current payments exceed current receipts and the US runs a persistent current account deficit. As a result, the extra flow of income (including in the form of goods and services) is balanced by an outflow of assets. Economic theory suggests that investment flows towards countries with lower labor costs and less capital per worker, as they have higher marginal productivity from additional capital. However, in the case of the US the opposite is happening, as net investment is flowing from less-developed countries with lower wages into the US. In other words, the US is borrowing money from less-developed countries to balance out its trade deficit, as shown in the financial accounts. 
\end{minipage}
\newpage
\begin{minipage}{0.76\textwidth}
\normalsize \textbf{Components of Current Account}\\
\footnotesize{\textit{share of GDP, percent \hspace{70mm} moving averages}}\\ 
\rowcolors{1}{}{black!5} \setlength{\tabcolsep}{3.1pt} \color{black!90}
	{\renewcommand{\arraystretch}{1.48}
		\begin{tabular}{p{32mm} R{7.6mm} R{7.6mm} R{7.6mm} R{7.6mm} R{7.6mm} 
		   R{7.6mm} R{9.1mm} R{8.0mm} }
			 \input data/cab.tex \hline
		\end{tabular}}\vspace{-2mm}
		
\footnotesize{Source: Bureau of Economic Analysis}
\vspace{4mm}

\small Financial account transactions include the net domestic acquisition of foreign assets and the net domestic incurrence of foreign liabilities. The US \textbf{financial account balance} (see\cbox{yellow!80!orange}) is the net lending or borrowing of the combined domestic sectors with the rest of the world. The timing of payments lead to a statistical discrepancy (see\cbox{red}), but the financial and capital account balance and current account balance otherwise sum to zero. 

\input{text/fab.txt}
\vspace{1mm}

\normalsize \textbf{Financial Account Balance}\\
\footnotesize{\textit{percent of GDP}}\\
\hspace*{-2mm} \begin{tikzpicture}
	\begin{axis}[\bbar{y}{0}, \dateaxisticks ytick={0, 5},
		xticklabel={`\short{\year}}, clip=false, height=5.8cm,
		legend style={at={(0.95, 1.13)}}]
	\rbars
	\sbar{yellow!80!orange}{date}{FAB}{data/fab.csv}
	\sbar{violet}{date}{FinDeriv}{data/fab.csv}
	\sbar{red}{date}{StatDisc}{data/fab.csv}
	\stdline{black}{date}{TOT}{data/fab.csv}
	\legend{Balance excl. Derivatives, Derivatives, Statistical Discrepancy, Total};
	\end{axis}
\end{tikzpicture}\\
\footnotesize{Source: Bureau of Economic Analysis} \hfill \tbllink{fab.csv}
\end{minipage}
\newpage
\hypertarget{extt}{}
\begin{minipage}{0.76\textwidth}
\subsection*{Trade}
\small Goods and services trade between the US and the rest of the world is \href{https://www.census.gov/foreign-trade/Press-Release/current\_press\_release/index.html}{reported} each month by the Census Bureau. US purchases of foreign goods and services are imports and foreign purchases of US goods and services are exports. Trade in goods includes consumer goods, industrial equipment, and agricultural products. Trade in services includes travel and tourism, business services, and charges for the use of intellectual property, among other services. 
\end{minipage}
\vspace{1mm}

\begin{minipage}{0.39\textwidth}
\normalsize \textbf{US Imports and Exports, Monthly}\\
\footnotesize{\textit{billions of USD, seasonally adjusted}}\\
\hspace*{-2mm} \begin{tikzpicture}
	\begin{axis}[\bbar{y}{0}, \ltdateaxisticks ytick={100, 150, 200, 250, 300}, 
		enlarge y limits={0.05}, clip=false, 
		yticklabel style={text width=1.3em},
		xticklabel={`\short{\year}}, height=6.6cm, width=6.6cm]
	\rbar
	\stdline{green!80!blue}{date}{EXP}{data/tradelt.csv}
	\input{text/explt.txt}
	\stdline{blue!80!violet}{date}{IMP}{data/tradelt.csv}
	\input{text/implt.txt}
	\absnode{{2017-01-01}}{190}{\footnotesize \color{green!80!blue}Exports}
	\absnode{{2016-03-01}}{265}{\footnotesize \color{blue!80!violet}Imports}
	\end{axis}
\end{tikzpicture}
\vspace{0.1mm}

\normalsize \textbf{Trade Balance, Monthly}\\
\footnotesize{\textit{billions of USD, seasonally adjusted}}\\
\hspace*{-3mm} \begin{tikzpicture}
	\begin{axis}[\bbar{y}{0}, \ltdateaxisticks 
		enlarge y limits={0.06}, clip=false, 
		yticklabel style={text width=1.6em},
		xticklabel={`\short{\year}}, height=5.6cm, width=6.6cm]
	\rbar
	\stdline{red}{date}{BAL}{data/tradelt.csv}
	\input{text/ballt.txt}
	\end{axis}
\end{tikzpicture}\\
\footnotesize{Source: Census Bureau} \hfill \tbllink{tradelt.csv}
\end{minipage} \hspace{6mm}
\begin{minipage}{0.32\textwidth}
\small \input{text/tradeltlevels.txt}
\end{minipage}
\newpage
%\begin{minipage}{0.76\textwidth}
%\small While the US is less dependent on exports than most major economies, trade is far from inconsequential to domestic economic activity and quality of life. Additionally, over the past 40 years, a major shift in trade policy, the reduction of tariffs, lowered prices substantially but closed a large portion of the US factories with very little compensation for workers.
%\end{minipage}
%\vspace{1mm}
%
%\begin{minipage}{0.33\textwidth}
%\small \input{text/tariff.txt}
%\end{minipage} \hspace{7mm} \begin{minipage}{0.37\textwidth}
%\normalsize \textbf{Average Import Tariff Rate}\\
%\footnotesize{\textit{weighted average applied rate, percent}}\\
%\hspace*{-2mm} \begin{tikzpicture}
%	\begin{axis}[\bbar{y}{0}, \shdateaxisticks ytick={0, 2, 4, 6, 8, 10, 12, 14}, clip=false, ymin=0.15,
%		xticklabel={`\short{\year}}, height=5.5cm, width=6.5cm]
%	\rbars
%	\thickline{violet}{date}{tariff}{data/tariff.csv}
%	\input{text/tariff_node.txt}
%	\end{axis}
%\end{tikzpicture}\\
%\footnotesize{Source: World Bank} \hfill \tbllink{tariff.csv}
%\end{minipage} 
%\vspace{7mm}

\begin{minipage}{0.76\textwidth}
\small & `23  Q3 & `23  Q2 & `22  Q3 &2016& 2012  --13 & 2005  --06 & 1998  --99 & 1989  --93 \\  Exports  of  Goods  \&  Services &10.98&10.94&11.87&11.89&13.60&10.31&10.41&9.42\\  Exports  of  Goods &7.34&7.27&8.23&7.70&9.34&7.30&7.52&6.84\\  \hspace{2mm}Foods,  Feeds,  \&  Beverages &0.56&0.58&0.69&0.69&0.81&0.46&0.50&0.60\\  \hspace{2mm}Industrial  Supplies  \&  Materials &2.55&2.52&3.27&2.06&2.94&1.92&1.55&1.65\\  \hspace{4mm}Petroleum  \&  Products &1.05&0.98&1.33&0.53&0.89&0.28&0.11&0.12\\  \hspace{2mm}Capital  Goods,  Except  Automotive &2.20&2.19&2.25&2.77&3.21&2.84&3.27&2.61\\  \hspace{2mm}Automotive  Vehicles,  \&  Parts &0.70&0.66&0.63&0.80&0.90&0.77&0.79&0.67\\  \hspace{2mm}Consumer  Goods,  Ex.  Food  \&  Auto &0.96&0.94&0.97&1.03&1.11&0.91&0.86&0.74\\  \hspace{4mm}Durable  Goods &0.43&0.43&0.45&0.55&0.61&0.49&0.44&0.39\\  \hspace{4mm}Nondurable  Goods &0.53&0.51&0.52&0.47&0.50&0.41&0.42&0.35\\  Exports  of  Services &3.64&3.67&3.63&4.19&4.26&3.01&2.90&2.58\\  \hspace{2mm}Transport &0.37&0.37&0.36&0.43&0.54&0.46&0.49&0.59\\  \hspace{2mm}Travel &0.66&0.63&0.56&1.03&0.98&0.71&0.93&0.90\\  \hspace{2mm}Intellectual  Property  Charges &0.43&0.46&0.48&0.60&0.67&0.50&0.40&0.29\\  \hspace{2mm}Other  Business  Services &2.04&2.08&2.11&2.00&1.92&1.19&0.92&0.60\\  Imports  of  Goods  \&  Services &13.80&13.92&15.30&14.56&16.71&15.99&12.65&10.38\\  Imports  of  Goods &11.26&11.30&12.52&11.80&13.85&13.48&10.59&8.45\\  \hspace{2mm}Foods,  Feeds,  \&  Beverages &0.73&0.72&0.80&0.70&0.69&0.54&0.46&0.43\\  \hspace{2mm}Industrial  Supplies  \&  Materials &2.34&2.41&3.13&2.32&4.24&4.24&2.22&2.16\\  \hspace{4mm}Petroleum  \&  Products &0.94&0.87&1.25&0.85&2.49&2.15&0.65&0.87\\  \hspace{2mm}Capital  Goods,  Except  Automotive &3.11&3.16&3.39&3.16&3.35&3.00&3.03&2.04\\  \hspace{2mm}Automotive  Vehicles,  \&  Parts &1.72&1.68&1.57&1.87&1.84&1.84&1.74&1.46\\  \hspace{2mm}Consumer  Goods,  Ex.  Food  \&  Auto &2.78&2.80&3.14&3.11&3.17&3.20&2.47&1.83\\  \hspace{4mm}Durable  Goods &1.37&1.40&1.62&1.63&1.70&1.75&1.29&0.97\\  \hspace{4mm}Nondurable  Goods &1.41&1.40&1.52&1.48&1.47&1.46&1.18&0.86\\  Imports  of  Services &2.54&2.62&2.78&2.77&2.86&2.51&2.06&1.93\\  \hspace{2mm}Transport &0.48&0.51&0.64&0.49&0.59&0.60&0.54&0.55\\  \hspace{2mm}Travel &0.52&0.53&0.48&0.58&0.55&0.57&0.63&0.61\\  \hspace{2mm}Intellectual  Property  Charges &0.16&0.16&0.19&0.22&0.21&0.18&0.13&0.06\\  \hspace{2mm}Other  Business  Services &1.27&1.30&1.35&1.32&1.32&0.91&0.57&0.38\\ 
\vspace{1mm}

\normalsize \textbf{Imports and Exports, Nonpetroleum}\\
\footnotesize{\textit{includes goods and services, but excludes petroleum products, share of GDP}}\\
\hspace*{-2mm} \begin{tikzpicture}
	\begin{axis}[\bbar{y}{0}, \dateaxisticks ytick={8, 10, 12, 14}, enlarge y limits={0.01},
		xticklabel={`\short{\year}}, clip=false, height=5.4cm, width=12.4cm, ymin=8, ymax=15,
		yticklabel style={text width=1.0em},
		legend style={legend columns=1, at={(0.28, 0.84)}}]
	\rbars
	\thickline{green!60!teal!80!black}{date}{EX}{data/eximgdp.csv}
	\thickline{blue!90!cyan}{date}{IM}{data/eximgdp.csv}
	\input{text/np_im_node.txt}
	\input{text/np_ex_node.txt}
	\legend{Exports, Imports};
	\end{axis}
\end{tikzpicture}\\
\footnotesize{Source: Bureau of Economic Analysis} \hfill \tbllink{eximgdp.csv}

\subsubsection*{Contribution to Overall Growth}
\small The \textbf{trade balance} (exports of goods\hspace*{-0.5mm}\cbox{green!70!white} and services\hspace*{-0.5mm}\cbox{green!50!black} minus imports of goods\hspace*{-0.5mm}\cbox{cyan!70!white} and services\hspace*{-0.5mm}\cbox{blue!70!black}) acts as an adjustment to consumption and investment when calculating domestic production using the expenditure approach. A country with a positive trade balance, or trade surplus, produces more exports than its residents purchase in imports, therefore its trade balance is added to domestic purchases to calculate domestic production. The US runs a persistent trade deficit, which is subtracted from spending data to estimate domestic production. 

\input{text/trade.txt} 
\vspace{1mm}

\normalsize \textbf{International Trade}\\
\footnotesize{\textit{percentage point contribution to real GDP growth}}\\
\hspace*{-2mm} \begin{tikzpicture}
	\begin{axis}[\bbar{y}{0}, \dateaxisticks ytick={-10, -5, 0, 5, 10},
		yticklabel style={text width=1.4em},
		height=8.0cm,
		xticklabel={`\short{\year}}, clip=false, 
		legend style={at={(0.91, 1.01)}, fill=white}]
	\rbars
	\sbar{green!70!white}{date}{A253RY}{data/nx.csv}
	\sbar{green!50!black}{date}{A646RY}{data/nx.csv}
	\sbar{cyan!70!white}{date}{A255RY}{data/nx.csv}
	\sbar{blue!70!black}{date}{A656RY}{data/nx.csv}	
	\legend{Goods Exports, Services Exports, Goods Imports, Services Imports};
	\end{axis}
\end{tikzpicture}\\
\footnotesize{Source: Bureau of Economic Analysis} \hfill \tbllink{nx.csv}
\end{minipage}
\newpage
\begin{minipage}{0.76\textwidth}
\small Changes to the trade balance can come from many sources, such as changes in domestic or foreign preferences and income, changes in exchange rates, and changes in trade policy. The following table shows major categories of trade as a share of gross domestic product at various points over the past 30 years. 
\vspace{1mm}

\normalsize \textbf{Exports and Imports by Type}\\
\footnotesize{\textit{percentage point share of GDP \hspace{40mm} period averages}\\ 
\rowcolors{1}{}{black!5} \setlength{\tabcolsep}{3.1pt} \color{black!90}
	{\renewcommand{\arraystretch}{1.55}
		\begin{tabular}{p{42mm} R{7.1mm} R{7.1mm} R{7.1mm} R{7.1mm} R{7.1mm} 
		   R{7.1mm} R{7.1mm} R{7.1mm} }
			 \input data/exim.tex \hline
		\end{tabular}}\vspace{-2mm}
		
\footnotesize{Source: Bureau of Economic Analysis}}
\end{minipage}
\newpage
\subsubsection*{Import Penetration}
\begin{minipage}{0.76\textwidth}
\small Goods can be produced domestically or imported or some combination of the two. The import share of the total US demand for goods, measured as US produced goods and imported goods less exported goods, is also referred to as \textbf{import penetration}. This measure has risen considerably over the past thirty years. The majority of the long-term increase has been concentrated in consumer goods, while the decrease since 2011 has come primarily from petroleum and related products. 

\input{text/goodsimpsh.txt} 

\input{text/goodsimpsh1.txt}

\vspace{3mm}

\normalsize \textbf{Import Share of Goods}\\
\footnotesize{\textit{import share of domestic-produced and imported goods less exports, by category}}\\
\hspace*{-2mm} \begin{tikzpicture}
	\begin{axis}[\bbar{y}{0}, \dateaxisticks ytick={0, 10, 20, 30, 40}, ymin=1,
		yticklabel style={text width=1.0em}, height=5.4cm, 
		xticklabel={`\short{\year}}, clip=false, legend style={at={(0.79, 1.12)}}]
	\rbars
	\sbar{cyan!40!white}{date}{Consumer}{data/goodsimpsh.csv}	
	\sbar{blue!50!cyan}{date}{Capital}{data/goodsimpsh.csv}	
	\sbar{purple}{date}{B648RC}{data/goodsimpsh.csv}	
	\legend{Consumer, Capital/Supplies/Other, Petroleum and Products};
	\end{axis}
\end{tikzpicture}\\
\footnotesize{Source: Bureau of Economic Analysis} \hfill \tbllink{goodsimpsh.csv}
\end{minipage}
\newpage
\begin{minipage}{0.76\textwidth}
\small The US Census Bureau \href{https://www.census.gov/foreign-trade/data/index.html}{report} monthly data on US trade in goods, including by partner country. \input{text/trade_partner_sh.txt} 
\vspace{2mm}

\normalsize \textbf{US Trade in Goods} \hspace{12mm} \input{text/trade_date.txt} \hspace{22mm} \input{text/trade_date2.txt}\\
\footnotesize{\textit{census basis, millions of USD,}}\\
\footnotesize{\textit{not seasonally adjusted}}\\
\rowcolors{1}{}{black!5} \setlength{\tabcolsep}{3.0pt} \color{black!90}
	{\renewcommand{\arraystretch}{1.54}
		\begin{tabular}{p{30mm} R{12mm} R{12mm} R{12mm} R{15mm} R{12mm} R{12mm}}
			 \input data/trade_partner.tex \hline
		\end{tabular}}
	\vspace{-2mm}		
	
\footnotesize{Source: Census Bureau}
\end{minipage}
\newpage
\begin{minipage}{0.76\textwidth}
\subsection*{International Investment Position}
\vspace{-0.5mm}
\small The US \textbf{net international investment position} (IIP) measures the difference between residents' foreign assets and liabilities. The Bureau of Economic Analysis \href{https://www.bea.gov/data/intl-trade-investment/international-investment-position}{report} US IIP data on a quarterly basis beginning in 2006, while prior data are annual. 

\input{text/niip.txt}
\vspace{1mm}

\normalsize \textbf{International Investment}\\
\footnotesize{\textit{share of GDP, percent}}\\
\hspace*{-2mm} \begin{tikzpicture}
	\begin{axis}[\bbar{y}{0}, \dateaxisticks ytick={50, 100, 150, 200}, 
		enlarge y limits={0.1},
		xticklabel={`\short{\year}}, clip=false, height=5.2cm, 
		yticklabel style={text width=1.5em}, width=12.9cm,
		legend style={legend columns=1, at={(0.28, 0.92)}, fill=white}, 
		fill between/on layer={main}]
	\rbars
	\addplot[name path=F, ultra thick, no markers, color=red] 
		table [x=date, y=FinLiabs, col sep=comma] {data/iip.csv};
	\addplot[name path=G, ultra thick, no markers, color=blue!90!violet] 
		table [x=date, y=FinAssets, col sep=comma] {data/iip.csv};	
	\addplot[color=cyan!25!white]
		fill between[of=F and G];
	\legend{Liabilities, Assets, Net IIP};
	\end{axis}
\end{tikzpicture}\\
\footnotesize{Source: Bureau of Economic Analysis} \hfill \tbllink{iip.csv}
\newpage
\subsection*{Capital Flows}
\vspace{-0.5mm}
\small The \href{https://home.treasury.gov/data/treasury-international-capital-tic-system}{purchases and sales} of US bonds by the rest of the world give insight into overall \textbf{capital flows} and appetite for different types of debt. During the 2000s, other countries were accumulating US corporate and government bonds, as the US was borrowing from the rest of the world and running a very wide trade deficit. During the pandemic, the rest of the world has been a net purchaser of US government agency bonds but a net seller of treasuries. 
\vspace{1mm}

\normalsize \textbf{Long-term Bond Flows}\\
\footnotesize{\textit{purchases by foreigners minus sales by foreigners}}\\\footnotesize{\textit{trailing 12-month sum, billions of \input{text/tic_date.txt}US dollars}}\\
\hspace*{-3mm} \begin{tikzpicture}
	\begin{axis}[\bbar{y}{0}, \dateaxisticks max space between ticks=20,	
		xticklabel={`\short{\year}}, clip=false, yticklabel style={text width=2.2em},
		legend cell align={left}, reverse legend, width=12.6cm, height=5.8cm,
		legend style={legend columns=1, at={(0.29, 1.0)}, fill=white}]
	\rbars
	\ctsbar{black!25!white}{date}{Treasury_Bonds}{data/tic_bond.csv}{0.8pt}
	\ctsbar{orange!30!yellow!90!white}{date}{Agency_Bonds}{data/tic_bond.csv}{0.8pt}
	\ctsbar{blue!70!white}{date}{Corporate_Bonds}{data/tic_bond.csv}{0.8pt}
	\legend{US Treasury Bonds, US Agency Bonds, US Corporate Bonds};
	\end{axis}
\end{tikzpicture}\\
\footnotesize{Source: Treasury International Capital, Setser} \hfill \tbllink{tic_bond.csv}
\end{minipage}
\newpage
\hypertarget{extfx}{}
\begin{minipage}{0.76\textwidth}
\subsection*{Exchange Rates}
\small Changes in the strength or weakness of the US dollar (USD) can affect trade and financial flows. The dollar is said to be relatively strong when more units of foreign currency, for example Japanese yen (JPY), British pounds (GBP), euros (EUR), or Canadian dollars (CAD), are required to buy one USD.

\input{text/selcurr_basic.txt}
\vspace{1mm}

\normalsize \textbf{Selected Exchange Rates}\\
\footnotesize{\textit{units of foreign currency required to purchase one US dollar}}\\
\hspace*{-2mm} \begin{tikzpicture}
	\begin{axis}[\bbar{y}{1}, \dateaxisticks ytick={0.5, 1, 1.5}, 
		yticklabel style={text width=1.2em}, 
		xticklabel={`\short{\year}}, clip=false, height=4.8cm, width=12.0cm]
	\rbars
	\stdline{blue!90!cyan}{date}{RXI_N.B.UK}{data/fx1.csv}
	\stdline{green!80!blue}{date}{RXI_N.B.CA}{data/fx1.csv}
	\stdline{red}{date}{RXI_N.B.JA}{data/fx1.csv}
	\stdline{cyan!90!white}{date}{RXI_N.B.EU}{data/fx1.csv}
	\absnode{{2016-06-01}}{0.7}{\footnotesize \color{blue!90!cyan}GBP}
	\absnode{{2003-01-01}}{1.62}{\footnotesize \color{green!80!blue}CAD}
	\absnode{{1996-12-01}}{0.9}{\footnotesize \color{cyan!90!white}EUR}
	\absnode{{1990-12-01}}{0.94}{\footnotesize \color{red}JPY (100)}
	\input{text/fx_nodes.txt}
	\end{axis}
\end{tikzpicture}\\
\footnotesize{Source: Federal Reserve} \hfill \tbllink{fx1.csv}
\vspace{4mm} 

\small The Federal Reserve \textbf{trade-weighted dollar} indices \href{https://www.federalreserve.gov/releases/h10/summary/default.htm}{track} weighted-average foreign exchange rates based on 26 currencies that are important to US trade. The \href{https://www.federalreserve.gov/releases/h10/weights/default.htm}{weight} of each currency in the index is based on the bilateral trade share of total trade in goods and services. These US dollar indices can simplify analysis of the overall role of foreign exchange rates on US trade. 
\vspace{0.5mm}

\begin{minipage}{0.46\textwidth}
\normalsize \textbf{Dollar Indices}\\
\footnotesize{\textit{trade-weighted foreign exchange rate,}}\\
\footnotesize{\textit{index, January 2006=100}}\\
\hspace*{-2mm} \begin{tikzpicture}
\begin{groupplot}[group style={group size=1 by 2, 
    x descriptions at=edge bottom, vertical sep=12pt,},]
	\nextgroupplot[\bbar{y}{100}, \dateaxisticks ytick={90, 100, 110, 120}, 
		yticklabel style={text width=1.5em}, enlarge y limits={0.1},
		xticklabel={`\short{\year}}, width=5.8cm, height=4.5cm,
		clip=false, xticklabels=\empty, xtick style={draw=none},]
	\recbars
	\stdline{blue!70!black}{date}{JRXWTFB_N.B}{data/fx_idx.csv}
	\absnode{{2016-06-01}}{108}{\footnotesize \color{blue!70!black}Broad}
	\input{text/twd_node.txt}
\nextgroupplot[\bbar{y}{100}, \dateaxisticks ytick={80, 100, 120, 140}, 
		yticklabel style={text width=1.5em}, enlarge y limits={0.1},
		xticklabel={`\short{\year}}, width=5.8cm, height=4.5cm,
		clip=false]
	\recbars
	\stdline{violet!90!blue}{date}{JRXWTFN_N.B}{data/fx_idx.csv}
	\stdline{green!75!yellow!92!black}{date}{JRXWTFO_N.B}{data/fx_idx.csv}
	\absnode{{2014-06-01}}{134}{\footnotesize \color{green!75!yellow!90!black}Emerging}
	\absnode{{2017-01-01}}{102}{\footnotesize \color{violet!90!blue}Advanced}
	\input{text/twd_nodes.txt}
\end{groupplot}
\end{tikzpicture}\\
\footnotesize{Source: Federal Reserve} \hfill \tbllink{fx_idx.csv}
\end{minipage}\hfill
\begin{minipage}{0.45\textwidth}
\small The \textbf{broad dollar index} (see {\color{blue!60!black}\textbf{---}}) summarizes foreign exchange rates between the US and trading partners by weighting foreign currencies in the index by the total amount of goods and services trade with the relevant countries.

\input{text/twdbasic.txt}

\input{text/twd_adv_em.txt}
\end{minipage}
\end{minipage}
\newpage
\begin{minipage}{0.76\textwidth}
\small Summarizing exchange rates using a trade-weighted index can be complicated by shifts in relative consumer prices between the US and trading partners. For example, the exchange rate between the US dollar and the Japanese yen has been relatively stable in nominal terms since the 1990s, but there has been less inflation in consumer prices in Japan than in the US. As a result, 100 yen buys more consumer goods in Japan than one dollar will buy in the US. 

By incorporating the consumer price index in the US and in trading partners, again weighted by the amount of trade with each partner, it is possible to compare a basket of consumer goods across currencies. The measure is still affected by differences in the consumer basket between countries.  
\vspace{0.5mm}

\begin{minipage}{0.58\textwidth}
\normalsize \textbf{USD Real Effective Exchange Rate}\\
\footnotesize{\textit{trade-weighted foreign exchange rate,}}\\
\footnotesize{\textit{index, January 2010=100}}\\
\hspace*{-2mm} \begin{tikzpicture}
\begin{axis}[\bbar{y}{100}, \shdateaxisticks %ytick={90, 100, 110, 120}, 
		yticklabel style={text width=1.5em}, enlarge y limits={0.1},
		xticklabel={`\short{\year}}, width=7.5cm, height=3.9cm,
		clip=false]
	\rebars
	\stdline{red}{date}{VALUE}{data/reer.csv}
	%\absnode{{2016-06-01}}{108}{\footnotesize \color{blue!70!black}Broad}
	\input{text/reer_node.txt}
\end{axis}
\end{tikzpicture}\\
\footnotesize{Source: BIS} \hfill \tbllink{reer.csv}
\end{minipage}\hfill
\begin{minipage}{0.33\textwidth}
\small \input{text/reer.txt}
\end{minipage}
\vspace{4mm}

\normalsize \textbf{Selected Exchange Rates}\\
\footnotesize{\textit{units of foreign currency required to buy one US dollar}}\\
\rowcolors{1}{}{black!5} \setlength{\tabcolsep}{1.0pt} \color{black!90}
	{\renewcommand{\arraystretch}{1.52}
		\begin{tabular}{p{12mm} R{14.5mm} R{14.5mm} R{14.5mm} R{14.5mm} R{14.0mm} 
		 				 R{14.0mm} R{14.0mm}}
			 \input data/fx_table.tex \hline
		\end{tabular}}\vspace{-2mm}
		
\footnotesize{Source: Federal Reserve}
\end{minipage}
\newpage
\hypertarget{lab}{}
\begin{minipage}{0.34\textwidth}
\section*{Labor Markets}
\vspace{-2mm}
\small Labor is the primary source of income for US households and is essential to the production of goods and services. The portion of labor that is provided by a household member to others outside of the household or to other households is considered \textit{employment}. \input{text/cps_lfs0.txt}

\input{text/cps_lfs.txt} Note that these values are not seasonally adjusted and include children, while BLS published values refer to those 16 or older.
\end{minipage} \hspace{6mm}
\begin{minipage}{0.57\textwidth}
\normalsize \textbf{Labor Force Status of Population}\\
\footnotesize{\textit{millions of people}}\\
\hspace*{-3mm} \begin{tikzpicture}
	\begin{axis}[\bbar{y}{0}, \dateaxisticks ytick={0, 100, 200, 300}, 
		height=7.4cm, ymin=10, width=10.0cm,
		clip=false, enlarge x limits={0.01}, 
		yticklabel style={text width=1.5em},
		xtick={{1989-01-01}, {1995-01-01}, {2000-01-01}, 
				{2005-01-01}, {2010-01-01}, {2015-01-01}, {2020-01-01}},
        minor xtick={},
        xticklabels={`89, `95, `00, `05, `10, `15, `20},
        legend cell align={left},
		legend style={fill=white, legend columns=2, at={(0.54, 1.09)}}]
	\rbars
	\ctsbar{green!50!teal!85!blue}{date}{Employed}{data/cps_lfs.csv}{0.7pt}
	\ctsbar{lime!80!yellow!90!white}{date}{Unemployed}{data/cps_lfs.csv}{0.7pt}
	\ctsbar{blue!80!cyan}{date}{NILF}{data/cps_lfs.csv}{0.7pt}
	\ctsbar{cyan!80!white}{date}{Children}{data/cps_lfs.csv}{0.7pt}
	\legend{Employed, Unemployed, Non-Participant, Under 15};
	\end{axis}
\end{tikzpicture}\\
\footnotesize{Source: Author's Calculations from CPS} \hfill \tbllink{cps_lfs.csv}
\end{minipage}
\vspace{1mm}

\begin{minipage}{0.76\textwidth}
\small When a member of a household is not employed but looked for a job during the past four weeks or is on temporary layoff, they are considered \href{https://www.bls.gov/cps/cps_htgm.htm}{\textbf{unemployed}}. \input{text/cps_lfs2.txt}

People who are not employed and not unemployed are considered to be outside of the labor force. \input{text/cps_lfs3.txt} 
\vspace{1mm}

\normalsize \textbf{Labor Force Status Changes}\\
\footnotesize{\textit{percentage point contribution to one-year growth of age 15+ population}}\\
\hspace*{-2mm} \begin{tikzpicture}
	\begin{axis}[\bbar{y}{0}, \dateaxisticks ytick={-8, -6, -4, -2, 0, 2, 4, 6, 8},
		xticklabel={`\short{\year}}, clip=false, height=7.4cm,
		yticklabel style={text width=1.0em},
		legend cell align={left}, yticklabel style={text width=1.3em},
		legend style={legend columns=1, at={(0.29, 0.92)}}]
	\rbars
	\ctsbar{green!50!teal!85!blue}{date}{Employed}{data/cps_lfs2.csv}{2.8pt}
	\ctsbar{lime!80!yellow!90!white}{date}{Unemployed}{data/cps_lfs2.csv}{2.8pt}
	\ctsbar{blue!80!cyan}{date}{NILF}{data/cps_lfs2.csv}{2.8pt}
	\legend{Employed, Unemployed, Non-Participant};
	\end{axis}
\end{tikzpicture}\\
\footnotesize{Source: Author's Calculations from CPS} \hfill \tbllink{cps_lfs2.csv}
\end{minipage}
\newpage
\begin{minipage}{0.76\textwidth}
\small The labor force status of the US population varies by age, sex, and over time. Because very few people receive capital income, the share of the population with labor income is particularly important to overall levels of economic activity. 
\vspace{2mm}

\normalsize \textbf{Labor Force Status}\\
\footnotesize{\textit{\input{text/cps_lfs_date.txt}\unskip, thousands of people, not seasonally adjusted}}\\
\noindent \rowcolors{1}{}{black!5} \setlength{\tabcolsep}{3.0pt} \color{black!90}
		{\renewcommand{\arraystretch}{1.52}
		 \begin{tabular}{p{27.8mm} R{11.2mm} R{10.4mm} R{10.4mm} R{10.2mm} 
		 				 R{10.6mm} R{10.4mm} R{10.4mm}}
			 \input data/lfs_table1.tex \hline
		\end{tabular}}
\vspace{-2mm}	
	
\footnotesize{Source: Author's Calculations from CPS}
\vspace{4mm}

\small Additionally, changes over time in labor force status are particularly important to understanding both secular and cyclical trends in the economy. The following table presents the net two-year change in labor force status, in number of people.
\vspace{2mm}

\normalsize \textbf{Labor Force Changes}\\
\footnotesize{\textit{Change from \input{text/cps_lfs_date2pr.txt}\unskip, thousands of people}}\\
\noindent \rowcolors{1}{}{black!5} \setlength{\tabcolsep}{3.0pt} \color{black!90}
		{\renewcommand{\arraystretch}{1.52}
		 \begin{tabular}{p{27.8mm} R{11.2mm} R{10.4mm} R{10.4mm} R{10.2mm} 
		 				 R{10.6mm} R{10.4mm} R{10.4mm}}
			 \input data/lfs_table2.tex \hline
		\end{tabular}}
		\vspace{-2mm}
		
\footnotesize{Source: Author's Calculations from CPS}
\end{minipage}
\newpage
\begin{minipage}{0.76\textwidth}
\subsubsection*{Labor Force Status and Age}
\small There is a clear relationship between age and employment. Children are not permitted to work and many young people attend school full time. During ages 25 to 54, around 80 percent of the population is employed. The remaining 20 percent include caregivers and those unable to work due to disability or illness. Retirement becomes more likely as workers reach their 60s and 70s; less than 10 percent of people continue to work into their 80s. 
\vspace{1mm}

\normalsize \textbf{Labor Force Status, by Age}\\
\footnotesize{\textit{share of same-age population, percent, \input{text/cps_lfs_date.txt}}}\\
\hspace*{-2mm} \begin{tikzpicture}
	\begin{axis}[ymin=0, ymax=100, xmin=0, xmax=80, 
		yticklabel style={text width=1.5em},
		width=12.6cm, height=11.4cm, xlabel={Age}]
	\abar{blue!80!black!90!white}{AGE}{Employed}{data/lfs_age.csv}
	\abar{orange!70!yellow!70!white}{AGE}{Disabled/Ill}{data/lfs_age.csv}
	\abar{violet!80!purple!60!white}{AGE}{Retired}{data/lfs_age.csv}
	\abar{green!60!yellow!80!gray!80!white}{AGE}{Family}{data/lfs_age.csv}
	\abar{green!60!blue!80!black}{AGE}{Unemployed/Discouraged/Other}{data/lfs_age.csv}
	\abar{cyan!50!white}{AGE}{Child/Student}{data/lfs_age.csv}
	\stdnode{{40}}{{50}}{\large \color{white}\textbf{Employed}}
	\stdnode{{10}}{{85}}{\large \color{black}\textbf{Child/}\\ 
		\large \color{black}\textbf{Student}}
	\stdnode{{42}}{{99}}{\small \color{white}\textbf{Unemployed/Other}}
	\stdnode{{70}}{{85}}{\large \color{white}\textbf{Retired}}
	\stdnode{{40}}{{92}}{\small \color{black}\textbf{Carer}}
	\stdnode{{47}}{{85}}{\small \color{black}\textbf{Disabled/Ill}}
	\end{axis}
\end{tikzpicture}
\vspace{-1mm}

\footnotesize{\textit{change since \input{text/cps_lfs_datepc.txt} in share of age group, percentage points}}\\
\hspace*{-2mm} \begin{tikzpicture}
	\begin{axis}[\bbar{y}{0}, height=6.8cm, width=12.6cm, xmin=0, xmax=80, 
		enlarge y limits={0.05}, xlabel={Age}, axis line style={draw=none},
		yticklabel style={text width=1.5em},
		legend style={legend columns=1, at={(0.28, 0.99)}, reverse legend,
		legend cell align={left}}]
	\ctsbar{blue!80!black!90!white}{AGE}{Employed}{data/lfs_age_ch.csv}{4.0pt}
	\ctsbar{cyan!50!white}{AGE}{Student}{data/lfs_age_ch.csv}{4.0pt}
	\ctsbar{orange!70!yellow!70!white}{AGE}{Disabled/Ill}{data/lfs_age_ch.csv}{4.0pt}
	\ctsbar{violet!80!purple!60!white}{AGE}{Retired}{data/lfs_age_ch.csv}{4.0pt}
	\ctsbar{green!60!yellow!80!gray!80!white}{AGE}{Family}{data/lfs_age_ch.csv}{4.0pt}
	\ctsbar{green!60!blue!80!black}{AGE}{Unemployed/Discouraged/Other}{data/lfs_age_ch.csv}{4.0pt}
	\legend{Employed, Student, Disabled/Ill, Retired, Carer, Unemployed/Other};
	\end{axis}
\end{tikzpicture}\\
\footnotesize{Source: Author's Calculations from CPS, Bruenig} \hfill \tbllink{lfs_age.csv} \tbllink{lfs_age_ch.csv}
\end{minipage}
\newpage
\subsection*{Gross Labor Force Status Changes}
\begin{minipage}{0.76\textwidth}
\small The changes in labor force status described above are net changes, rather than gross changes. That is, the one-year change in unemployment is the result of some people staying unemployed, some becoming unemployed, and some no longer being unemployed, as well as other flows like deaths and people becoming working age. This section looks at six categories of gross changes, specifically the share of people with a given labor force status who move to a different status in the following month. 
\end{minipage}
\vspace{2mm}

\normalsize \textbf{Gross Labor Force Flows}\\
\footnotesize{\textit{share of initial labor force status, percent}}

\begin{minipage}{0.345\textwidth}
\small Employed to Unemployed\\
\hspace*{-2mm} \begin{tikzpicture}
	\begin{axis}[\bbar{y}{0}, \lfdateaxisticks ytick={0, 5, 10}, 
		xticklabel={\year}, 
		legend style={fill=white, legend columns=1, at={(0.45, 0.99)},
		legend cell align={left}}]
	\rbars
	\stdline{blue!40!cyan}{date}{MenEUE}{data/grosslf.csv}
	\stdline{green!55!black}{date}{WomenEUE}{data/grosslf.csv}
	\legend{Men, Women};
	\end{axis}
\end{tikzpicture}
\end{minipage} \hspace{8mm}
\begin{minipage}{0.345\textwidth}
\small Employed to Out of Labor Force\\
\hspace*{-2mm} \begin{tikzpicture}
	\begin{axis}[\bbar{y}{0}, \lfdateaxisticks xticklabel={\year}, 
	ytick={0, 2, 4, 6}]
	\rbars
	\stdline{blue!40!cyan}{date}{MenENE}{data/grosslf.csv}
	\stdline{green!55!black}{date}{WomenENE}{data/grosslf.csv}
	\end{axis}
\end{tikzpicture}
\end{minipage}
\vspace{1mm}

\begin{minipage}{0.345\textwidth}
\small Unemployed to Employed\\
\hspace*{-2mm} \begin{tikzpicture}
	\begin{axis}[\bbar{y}{0}, \lfdateaxisticks xticklabel={\year}, 
	ytick={0, 10, 20, 30, 40}]
	\rbars
	\stdline{blue!40!cyan}{date}{MenUEU}{data/grosslf.csv}
	\stdline{green!55!black}{date}{WomenUEU}{data/grosslf.csv}
	\end{axis}
\end{tikzpicture}
\end{minipage} \hspace{8mm}
\begin{minipage}{0.345\textwidth}
\small Unemployed to Out of Labor Force\\
\hspace*{-2mm} \begin{tikzpicture}
	\begin{axis}[\bbar{y}{0}, \lfdateaxisticks xticklabel={\year}, 
	ytick={0, 10, 20, 30, 40}]
	\rbars
	\stdline{blue!40!cyan}{date}{MenUNU}{data/grosslf.csv}
	\stdline{green!55!black}{date}{WomenUNU}{data/grosslf.csv}
	\end{axis}
\end{tikzpicture}
\end{minipage}
\vspace{1mm}

\begin{minipage}{0.345\textwidth}
\small Out of Labor Force to Employed\\
\hspace*{-2mm} \begin{tikzpicture}
	\begin{axis}[\bbar{y}{0}, \lfdateaxisticks xticklabel={\year}, 
	ytick={0, 2, 4, 6}]
	\rbars
	\stdline{blue!40!cyan}{date}{MenNEN}{data/grosslf.csv}
	\stdline{green!55!black}{date}{WomenNEN}{data/grosslf.csv}
	\end{axis}
\end{tikzpicture}
\end{minipage} \hspace{8mm}
\begin{minipage}{0.345\textwidth}
\small Out of Labor Force to Unemployed\\
\hspace*{-2mm} \begin{tikzpicture}
	\begin{axis}[\bbar{y}{0}, \lfdateaxisticks xticklabel={\year}, 
	ytick={0, 1, 2, 3, 4, 5}]
	\rbars
	\stdline{blue!40!cyan}{date}{MenNUN}{data/grosslf.csv}
	\stdline{green!55!black}{date}{WomenNUN}{data/grosslf.csv}
	\end{axis}
\end{tikzpicture}
\end{minipage}\\
\footnotesize{Source: Bureau of Labor Statistics} \hspace{72mm} \tbllink{grosslf.csv}
\newpage
\begin{minipage}{0.43\textwidth}
\normalsize \textbf{Labor Share of Income}\\
\footnotesize{\textit{compensation of employees as share of}}\\
\footnotesize{\textit{net domestic income, percent}}\\
\hspace*{-2mm} \begin{tikzpicture}
	\begin{axis}[\shdateaxisticks ytick={60, 62, 64, 66, 68}, 
		enlarge y limits={0.08}, xmin={1989-12-30},
		yticklabel style={text width=0.8em},
		xticklabel={`\short{\year}}, height=6.8cm, width=7.6cm, clip=false]
	\rbars
	\stdline{blue!60!cyan}{date}{Share}{data/laborshare.csv}
	\input{text/laborshare_node.txt}
	\end{axis}
\end{tikzpicture}\\
\footnotesize{Source: Bureau of Economic Analysis} \hfill \tbllink{laborshare.csv}
\end{minipage}\hspace{7mm}
\begin{minipage}{0.27\textwidth}
\small The labor share (see {\color{blue!60!cyan}\textbf{---}}) measures the portion of net income paid to labor. While the laborer share of the population has fallen, cyclical patterns suggest worker bargaining power also affects labor's share.  

\input{text/laborshare.txt}
\end{minipage}
\vspace{3mm}

\begin{minipage}{0.76\textwidth}
\subsection*{Gross Labor Income}
\small Wages (the unit price of labor) tend not to be cut in response to a short-term decrease in demand; businesses typically instead employ fewer workers and/or cut hours. As a result, wage data give only a partial picture of the labor income received by households.

\textbf{Gross labor income} (compensation of employees in the national accounts), which captures both the amount of employment (see\cbox{teal!80!blue!85!white}) and the rate of compensation (see\cbox{green!80!lime!90!white}), \input{text/gli.txt} 
\vspace{1mm}

\normalsize \textbf{Gross Labor Income Growth}\\
\footnotesize{\textit{percentage point contribution to gross labor income growth}}\\
\hspace*{-2mm} \begin{tikzpicture}
	\begin{axis}[\bbar{y}{0}, \dateaxisticks ytick={-50, -40, -30, -20, -10, 0, 10, 20, 30, 40},
		height=8.6cm,
		xticklabel={`\short{\year}}, clip=false,  yticklabel style={text width=1.4em},
		legend style={at={(0.51, 1.0)}, fill=white}]
	\rbars
	\sbar{green!80!lime!90!white}{date}{wage}{data/gli.csv}
	\sbar{teal!80!blue!85!white}{date}{work}{data/gli.csv}
	\legend{Rate of Compensation, Hours of Work};
	\end{axis}
\end{tikzpicture}\\
\footnotesize{Source: Author's Calculations} \hfill \tbllink{gli.csv}

\end{minipage}
\newpage
\begin{minipage}{0.76\textwidth}
\subsection*{Employment}
\hypertarget{labe}{}
\vspace{-0.5mm}
\small The employed share of those ages 25 to 54 is an important measure of labor market utilization. In a tight labor market, the age group is employed at a very high rate. \input{text/epop_text.txt} 

\normalsize \textbf{Age 25--54 Employment Rate}\\
\footnotesize{\textit{employed share of age 25--54 population}}\\
\hspace*{-2mm} \begin{tikzpicture}
	\begin{axis}[\bbar{y}{0}, \dateaxisticks ytick={70, 75, 80}, 
		yticklabel style={text width=1.0em}, 
		width=12.2cm, enlarge y limits={0.04},
		xticklabel={`\short{\year}}, clip=false, height=4.6cm]
	\rbars
	\stdline{blue!90!cyan}{date}{PA_EPOP}{data/epop.csv}
	\input{text/epop_node.txt}
	\end{axis}
\end{tikzpicture}\\
\footnotesize{Source: Bureau of Labor Statistics} \hfill \tbllink{epop.csv}
\vspace{1.5mm}

\small The monthly jobs report describes employment at a given point in time, by asking about activities during a specific week of the previous month. To instead examine activities over a period of time, annual data on weeks worked per year and hours worked per week can be combined to identify the \textit{fully-employed}, or \textit{full-time, full-year workers}, who usually work 35 hours per week or more for 50 weeks per year or more. The Census Bureau \href{https://www.census.gov/data/tables/time-series/demo/income-poverty/cps-pinc/pinc-01.html#par_textimage_14}{report} \input{text/asec_ftfy_tot.txt}

Employment rates vary dramatically by location. \input{text/acs_ftfy_text.txt}The top ten and bottom ten commuter zones by fully-employed rate are listed below.
\end{minipage}
\vspace{1mm}

\begin{minipage}{0.52\textwidth}
\normalsize \textbf{Commuter Zone Fully-Employed Rate}\\
\footnotesize{\textit{full-time, full-year worker share of population, \input{text/acs_ftfy_year.txt}}}\\
\vspace{1mm}

\hspace{-11mm} %% Creator: Matplotlib, PGF backend
%%
%% To include the figure in your LaTeX document, write
%%   \input{<filename>.pgf}
%%
%% Make sure the required packages are loaded in your preamble
%%   \usepackage{pgf}
%%
%% Also ensure that all the required font packages are loaded; for instance,
%% the lmodern package is sometimes necessary when using math font.
%%   \usepackage{lmodern}
%%
%% Figures using additional raster images can only be included by \input if
%% they are in the same directory as the main LaTeX file. For loading figures
%% from other directories you can use the `import` package
%%   \usepackage{import}
%%
%% and then include the figures with
%%   \import{<path to file>}{<filename>.pgf}
%%
%% Matplotlib used the following preamble
%%   
%%   \usepackage{fontspec}
%%   \setmainfont{DejaVuSerif.ttf}[Path=\detokenize{/home/brian/miniconda3/lib/python3.8/site-packages/matplotlib/mpl-data/fonts/ttf/}]
%%   \setsansfont{DejaVuSans.ttf}[Path=\detokenize{/home/brian/miniconda3/lib/python3.8/site-packages/matplotlib/mpl-data/fonts/ttf/}]
%%   \setmonofont{DejaVuSansMono.ttf}[Path=\detokenize{/home/brian/miniconda3/lib/python3.8/site-packages/matplotlib/mpl-data/fonts/ttf/}]
%%   \makeatletter\@ifpackageloaded{underscore}{}{\usepackage[strings]{underscore}}\makeatother
%%
\begingroup%
\makeatletter%
\begin{pgfpicture}%
\pgfpathrectangle{\pgfpointorigin}{\pgfqpoint{3.841155in}{2.510000in}}%
\pgfusepath{use as bounding box, clip}%
\begin{pgfscope}%
\pgfsetbuttcap%
\pgfsetmiterjoin%
\pgfsetlinewidth{0.000000pt}%
\definecolor{currentstroke}{rgb}{0.000000,0.000000,0.000000}%
\pgfsetstrokecolor{currentstroke}%
\pgfsetstrokeopacity{0.000000}%
\pgfsetdash{}{0pt}%
\pgfpathmoveto{\pgfqpoint{0.000000in}{0.000000in}}%
\pgfpathlineto{\pgfqpoint{3.841155in}{0.000000in}}%
\pgfpathlineto{\pgfqpoint{3.841155in}{2.510000in}}%
\pgfpathlineto{\pgfqpoint{0.000000in}{2.510000in}}%
\pgfpathlineto{\pgfqpoint{0.000000in}{0.000000in}}%
\pgfpathclose%
\pgfusepath{}%
\end{pgfscope}%
\begin{pgfscope}%
\pgfpathrectangle{\pgfqpoint{0.100000in}{0.100000in}}{\pgfqpoint{3.608454in}{2.310000in}}%
\pgfusepath{clip}%
\pgfsetbuttcap%
\pgfsetroundjoin%
\pgfsetlinewidth{0.010037pt}%
\definecolor{currentstroke}{rgb}{1.000000,1.000000,1.000000}%
\pgfsetstrokecolor{currentstroke}%
\pgfsetdash{}{0pt}%
\pgfusepath{stroke}%
\end{pgfscope}%
\begin{pgfscope}%
\pgfpathrectangle{\pgfqpoint{0.100000in}{0.100000in}}{\pgfqpoint{3.608454in}{2.310000in}}%
\pgfusepath{clip}%
\pgfsetbuttcap%
\pgfsetroundjoin%
\pgfsetlinewidth{0.010037pt}%
\definecolor{currentstroke}{rgb}{1.000000,1.000000,1.000000}%
\pgfsetstrokecolor{currentstroke}%
\pgfsetdash{}{0pt}%
\pgfusepath{stroke}%
\end{pgfscope}%
\begin{pgfscope}%
\pgfpathrectangle{\pgfqpoint{0.100000in}{0.100000in}}{\pgfqpoint{3.608454in}{2.310000in}}%
\pgfusepath{clip}%
\pgfsetbuttcap%
\pgfsetroundjoin%
\pgfsetlinewidth{0.010037pt}%
\definecolor{currentstroke}{rgb}{1.000000,1.000000,1.000000}%
\pgfsetstrokecolor{currentstroke}%
\pgfsetdash{}{0pt}%
\pgfusepath{stroke}%
\end{pgfscope}%
\begin{pgfscope}%
\pgfpathrectangle{\pgfqpoint{0.100000in}{0.100000in}}{\pgfqpoint{3.608454in}{2.310000in}}%
\pgfusepath{clip}%
\pgfsetbuttcap%
\pgfsetroundjoin%
\pgfsetlinewidth{0.010037pt}%
\definecolor{currentstroke}{rgb}{1.000000,1.000000,1.000000}%
\pgfsetstrokecolor{currentstroke}%
\pgfsetdash{}{0pt}%
\pgfusepath{stroke}%
\end{pgfscope}%
\begin{pgfscope}%
\pgfpathrectangle{\pgfqpoint{0.100000in}{0.100000in}}{\pgfqpoint{3.608454in}{2.310000in}}%
\pgfusepath{clip}%
\pgfsetbuttcap%
\pgfsetroundjoin%
\pgfsetlinewidth{0.010037pt}%
\definecolor{currentstroke}{rgb}{1.000000,1.000000,1.000000}%
\pgfsetstrokecolor{currentstroke}%
\pgfsetdash{}{0pt}%
\pgfusepath{stroke}%
\end{pgfscope}%
\begin{pgfscope}%
\pgfpathrectangle{\pgfqpoint{0.100000in}{0.100000in}}{\pgfqpoint{3.608454in}{2.310000in}}%
\pgfusepath{clip}%
\pgfsetbuttcap%
\pgfsetroundjoin%
\pgfsetlinewidth{0.010037pt}%
\definecolor{currentstroke}{rgb}{1.000000,1.000000,1.000000}%
\pgfsetstrokecolor{currentstroke}%
\pgfsetdash{}{0pt}%
\pgfusepath{stroke}%
\end{pgfscope}%
\begin{pgfscope}%
\pgfpathrectangle{\pgfqpoint{0.100000in}{0.100000in}}{\pgfqpoint{3.608454in}{2.310000in}}%
\pgfusepath{clip}%
\pgfsetbuttcap%
\pgfsetroundjoin%
\pgfsetlinewidth{0.010037pt}%
\definecolor{currentstroke}{rgb}{1.000000,1.000000,1.000000}%
\pgfsetstrokecolor{currentstroke}%
\pgfsetdash{}{0pt}%
\pgfpathmoveto{\pgfqpoint{0.957404in}{2.274144in}}%
\pgfpathlineto{\pgfqpoint{0.938691in}{2.199997in}}%
\pgfpathlineto{\pgfqpoint{0.925385in}{2.146869in}}%
\pgfpathlineto{\pgfqpoint{0.909565in}{2.082958in}}%
\pgfpathlineto{\pgfqpoint{0.906458in}{2.067206in}}%
\pgfpathlineto{\pgfqpoint{0.908487in}{2.052757in}}%
\pgfpathlineto{\pgfqpoint{0.905799in}{2.039449in}}%
\pgfpathlineto{\pgfqpoint{0.795364in}{2.068266in}}%
\pgfpathlineto{\pgfqpoint{0.785388in}{2.064919in}}%
\pgfpathlineto{\pgfqpoint{0.776851in}{2.067824in}}%
\pgfpathlineto{\pgfqpoint{0.737324in}{2.068774in}}%
\pgfpathlineto{\pgfqpoint{0.723956in}{2.065005in}}%
\pgfpathlineto{\pgfqpoint{0.710680in}{2.066278in}}%
\pgfpathlineto{\pgfqpoint{0.705053in}{2.072177in}}%
\pgfpathlineto{\pgfqpoint{0.669666in}{2.070891in}}%
\pgfpathlineto{\pgfqpoint{0.664077in}{2.080532in}}%
\pgfpathlineto{\pgfqpoint{0.653468in}{2.085611in}}%
\pgfpathlineto{\pgfqpoint{0.638032in}{2.088706in}}%
\pgfpathlineto{\pgfqpoint{0.611388in}{2.084174in}}%
\pgfpathlineto{\pgfqpoint{0.601553in}{2.088614in}}%
\pgfpathlineto{\pgfqpoint{0.586388in}{2.100423in}}%
\pgfpathlineto{\pgfqpoint{0.590946in}{2.123567in}}%
\pgfpathlineto{\pgfqpoint{0.589348in}{2.135386in}}%
\pgfpathlineto{\pgfqpoint{0.577330in}{2.147934in}}%
\pgfpathlineto{\pgfqpoint{0.569631in}{2.147155in}}%
\pgfpathlineto{\pgfqpoint{0.564109in}{2.159848in}}%
\pgfpathlineto{\pgfqpoint{0.551027in}{2.164978in}}%
\pgfpathlineto{\pgfqpoint{0.541518in}{2.164324in}}%
\pgfpathlineto{\pgfqpoint{0.538442in}{2.177359in}}%
\pgfpathlineto{\pgfqpoint{0.547915in}{2.175972in}}%
\pgfpathlineto{\pgfqpoint{0.547162in}{2.194090in}}%
\pgfpathlineto{\pgfqpoint{0.543033in}{2.204847in}}%
\pgfpathlineto{\pgfqpoint{0.544853in}{2.218617in}}%
\pgfpathlineto{\pgfqpoint{0.549810in}{2.228973in}}%
\pgfpathlineto{\pgfqpoint{0.549530in}{2.248625in}}%
\pgfpathlineto{\pgfqpoint{0.546884in}{2.255755in}}%
\pgfpathlineto{\pgfqpoint{0.551454in}{2.278631in}}%
\pgfpathlineto{\pgfqpoint{0.550135in}{2.293374in}}%
\pgfpathlineto{\pgfqpoint{0.545566in}{2.300443in}}%
\pgfpathlineto{\pgfqpoint{0.546263in}{2.323556in}}%
\pgfpathlineto{\pgfqpoint{0.552909in}{2.340554in}}%
\pgfpathlineto{\pgfqpoint{0.560136in}{2.336394in}}%
\pgfpathlineto{\pgfqpoint{0.584067in}{2.311673in}}%
\pgfpathlineto{\pgfqpoint{0.613066in}{2.298145in}}%
\pgfpathlineto{\pgfqpoint{0.627914in}{2.296544in}}%
\pgfpathlineto{\pgfqpoint{0.641902in}{2.290785in}}%
\pgfpathlineto{\pgfqpoint{0.645873in}{2.271403in}}%
\pgfpathlineto{\pgfqpoint{0.633614in}{2.267732in}}%
\pgfpathlineto{\pgfqpoint{0.625490in}{2.252457in}}%
\pgfpathlineto{\pgfqpoint{0.635028in}{2.253306in}}%
\pgfpathlineto{\pgfqpoint{0.638875in}{2.260162in}}%
\pgfpathlineto{\pgfqpoint{0.652326in}{2.268736in}}%
\pgfpathlineto{\pgfqpoint{0.651631in}{2.256004in}}%
\pgfpathlineto{\pgfqpoint{0.642650in}{2.253951in}}%
\pgfpathlineto{\pgfqpoint{0.644141in}{2.237586in}}%
\pgfpathlineto{\pgfqpoint{0.633204in}{2.221193in}}%
\pgfpathlineto{\pgfqpoint{0.626204in}{2.228988in}}%
\pgfpathlineto{\pgfqpoint{0.605402in}{2.224750in}}%
\pgfpathlineto{\pgfqpoint{0.604308in}{2.215213in}}%
\pgfpathlineto{\pgfqpoint{0.611479in}{2.209403in}}%
\pgfpathlineto{\pgfqpoint{0.622419in}{2.208832in}}%
\pgfpathlineto{\pgfqpoint{0.637546in}{2.221247in}}%
\pgfpathlineto{\pgfqpoint{0.649519in}{2.222278in}}%
\pgfpathlineto{\pgfqpoint{0.649978in}{2.236021in}}%
\pgfpathlineto{\pgfqpoint{0.656149in}{2.256205in}}%
\pgfpathlineto{\pgfqpoint{0.670772in}{2.271217in}}%
\pgfpathlineto{\pgfqpoint{0.665896in}{2.282844in}}%
\pgfpathlineto{\pgfqpoint{0.669227in}{2.295363in}}%
\pgfpathlineto{\pgfqpoint{0.662739in}{2.307839in}}%
\pgfpathlineto{\pgfqpoint{0.671985in}{2.323793in}}%
\pgfpathlineto{\pgfqpoint{0.673370in}{2.333392in}}%
\pgfpathlineto{\pgfqpoint{0.665342in}{2.339700in}}%
\pgfpathlineto{\pgfqpoint{0.666679in}{2.355837in}}%
\pgfpathlineto{\pgfqpoint{0.762858in}{2.326616in}}%
\pgfpathlineto{\pgfqpoint{0.864991in}{2.297975in}}%
\pgfpathlineto{\pgfqpoint{0.957404in}{2.274144in}}%
\pgfusepath{stroke}%
\end{pgfscope}%
\begin{pgfscope}%
\pgfpathrectangle{\pgfqpoint{0.100000in}{0.100000in}}{\pgfqpoint{3.608454in}{2.310000in}}%
\pgfusepath{clip}%
\pgfsetbuttcap%
\pgfsetroundjoin%
\pgfsetlinewidth{0.010037pt}%
\definecolor{currentstroke}{rgb}{1.000000,1.000000,1.000000}%
\pgfsetstrokecolor{currentstroke}%
\pgfsetdash{}{0pt}%
\pgfpathmoveto{\pgfqpoint{0.650857in}{2.300012in}}%
\pgfpathlineto{\pgfqpoint{0.655593in}{2.281754in}}%
\pgfpathlineto{\pgfqpoint{0.663020in}{2.275689in}}%
\pgfpathlineto{\pgfqpoint{0.657094in}{2.267969in}}%
\pgfpathlineto{\pgfqpoint{0.651318in}{2.279281in}}%
\pgfpathlineto{\pgfqpoint{0.650857in}{2.300012in}}%
\pgfusepath{stroke}%
\end{pgfscope}%
\begin{pgfscope}%
\pgfpathrectangle{\pgfqpoint{0.100000in}{0.100000in}}{\pgfqpoint{3.608454in}{2.310000in}}%
\pgfusepath{clip}%
\pgfsetbuttcap%
\pgfsetroundjoin%
\pgfsetlinewidth{0.010037pt}%
\definecolor{currentstroke}{rgb}{1.000000,1.000000,1.000000}%
\pgfsetstrokecolor{currentstroke}%
\pgfsetdash{}{0pt}%
\pgfpathmoveto{\pgfqpoint{1.007001in}{2.262160in}}%
\pgfpathlineto{\pgfqpoint{1.109811in}{2.239088in}}%
\pgfpathlineto{\pgfqpoint{1.206635in}{2.219517in}}%
\pgfpathlineto{\pgfqpoint{1.281130in}{2.205868in}}%
\pgfpathlineto{\pgfqpoint{1.346076in}{2.194957in}}%
\pgfpathlineto{\pgfqpoint{1.411166in}{2.184941in}}%
\pgfpathlineto{\pgfqpoint{1.466597in}{2.177132in}}%
\pgfpathlineto{\pgfqpoint{1.522115in}{2.169972in}}%
\pgfpathlineto{\pgfqpoint{1.577712in}{2.163461in}}%
\pgfpathlineto{\pgfqpoint{1.630106in}{2.157928in}}%
\pgfpathlineto{\pgfqpoint{1.622900in}{2.078078in}}%
\pgfpathlineto{\pgfqpoint{1.612175in}{1.970174in}}%
\pgfpathlineto{\pgfqpoint{1.606550in}{1.914676in}}%
\pgfpathlineto{\pgfqpoint{1.598460in}{1.839883in}}%
\pgfpathlineto{\pgfqpoint{1.541187in}{1.846137in}}%
\pgfpathlineto{\pgfqpoint{1.475618in}{1.853580in}}%
\pgfpathlineto{\pgfqpoint{1.384573in}{1.866040in}}%
\pgfpathlineto{\pgfqpoint{1.343910in}{1.871841in}}%
\pgfpathlineto{\pgfqpoint{1.243728in}{1.887567in}}%
\pgfpathlineto{\pgfqpoint{1.209241in}{1.893876in}}%
\pgfpathlineto{\pgfqpoint{1.202053in}{1.853008in}}%
\pgfpathlineto{\pgfqpoint{1.198123in}{1.855924in}}%
\pgfpathlineto{\pgfqpoint{1.190738in}{1.875584in}}%
\pgfpathlineto{\pgfqpoint{1.181245in}{1.874373in}}%
\pgfpathlineto{\pgfqpoint{1.174451in}{1.864434in}}%
\pgfpathlineto{\pgfqpoint{1.160350in}{1.863105in}}%
\pgfpathlineto{\pgfqpoint{1.157517in}{1.867860in}}%
\pgfpathlineto{\pgfqpoint{1.144309in}{1.867284in}}%
\pgfpathlineto{\pgfqpoint{1.137669in}{1.871889in}}%
\pgfpathlineto{\pgfqpoint{1.128385in}{1.864765in}}%
\pgfpathlineto{\pgfqpoint{1.105833in}{1.871176in}}%
\pgfpathlineto{\pgfqpoint{1.095991in}{1.867733in}}%
\pgfpathlineto{\pgfqpoint{1.090697in}{1.883009in}}%
\pgfpathlineto{\pgfqpoint{1.090449in}{1.897544in}}%
\pgfpathlineto{\pgfqpoint{1.074904in}{1.907997in}}%
\pgfpathlineto{\pgfqpoint{1.077439in}{1.923062in}}%
\pgfpathlineto{\pgfqpoint{1.066370in}{1.947947in}}%
\pgfpathlineto{\pgfqpoint{1.067116in}{1.969101in}}%
\pgfpathlineto{\pgfqpoint{1.057602in}{1.979454in}}%
\pgfpathlineto{\pgfqpoint{1.048676in}{1.970294in}}%
\pgfpathlineto{\pgfqpoint{1.036527in}{1.965340in}}%
\pgfpathlineto{\pgfqpoint{1.027278in}{1.975551in}}%
\pgfpathlineto{\pgfqpoint{1.028782in}{1.991857in}}%
\pgfpathlineto{\pgfqpoint{1.039874in}{1.998281in}}%
\pgfpathlineto{\pgfqpoint{1.036925in}{2.008611in}}%
\pgfpathlineto{\pgfqpoint{1.056205in}{2.058614in}}%
\pgfpathlineto{\pgfqpoint{1.041003in}{2.059867in}}%
\pgfpathlineto{\pgfqpoint{1.039440in}{2.068841in}}%
\pgfpathlineto{\pgfqpoint{1.028366in}{2.076578in}}%
\pgfpathlineto{\pgfqpoint{1.029074in}{2.085229in}}%
\pgfpathlineto{\pgfqpoint{1.023232in}{2.091949in}}%
\pgfpathlineto{\pgfqpoint{1.013723in}{2.116347in}}%
\pgfpathlineto{\pgfqpoint{1.005033in}{2.121323in}}%
\pgfpathlineto{\pgfqpoint{0.998811in}{2.142381in}}%
\pgfpathlineto{\pgfqpoint{1.000279in}{2.156377in}}%
\pgfpathlineto{\pgfqpoint{0.988670in}{2.182155in}}%
\pgfpathlineto{\pgfqpoint{1.007001in}{2.262160in}}%
\pgfusepath{stroke}%
\end{pgfscope}%
\begin{pgfscope}%
\pgfpathrectangle{\pgfqpoint{0.100000in}{0.100000in}}{\pgfqpoint{3.608454in}{2.310000in}}%
\pgfusepath{clip}%
\pgfsetbuttcap%
\pgfsetroundjoin%
\pgfsetlinewidth{0.010037pt}%
\definecolor{currentstroke}{rgb}{1.000000,1.000000,1.000000}%
\pgfsetstrokecolor{currentstroke}%
\pgfsetdash{}{0pt}%
\pgfpathmoveto{\pgfqpoint{3.488916in}{1.850686in}}%
\pgfpathlineto{\pgfqpoint{3.486858in}{1.859435in}}%
\pgfpathlineto{\pgfqpoint{3.475422in}{1.867100in}}%
\pgfpathlineto{\pgfqpoint{3.445201in}{1.965955in}}%
\pgfpathlineto{\pgfqpoint{3.428911in}{2.013660in}}%
\pgfpathlineto{\pgfqpoint{3.442007in}{2.027403in}}%
\pgfpathlineto{\pgfqpoint{3.454737in}{2.061039in}}%
\pgfpathlineto{\pgfqpoint{3.461086in}{2.071788in}}%
\pgfpathlineto{\pgfqpoint{3.456588in}{2.076302in}}%
\pgfpathlineto{\pgfqpoint{3.455082in}{2.106159in}}%
\pgfpathlineto{\pgfqpoint{3.460762in}{2.115312in}}%
\pgfpathlineto{\pgfqpoint{3.458272in}{2.136586in}}%
\pgfpathlineto{\pgfqpoint{3.480942in}{2.206386in}}%
\pgfpathlineto{\pgfqpoint{3.491156in}{2.206767in}}%
\pgfpathlineto{\pgfqpoint{3.495396in}{2.194299in}}%
\pgfpathlineto{\pgfqpoint{3.504473in}{2.190719in}}%
\pgfpathlineto{\pgfqpoint{3.521621in}{2.205384in}}%
\pgfpathlineto{\pgfqpoint{3.535066in}{2.214046in}}%
\pgfpathlineto{\pgfqpoint{3.564692in}{2.198793in}}%
\pgfpathlineto{\pgfqpoint{3.591456in}{2.114100in}}%
\pgfpathlineto{\pgfqpoint{3.596544in}{2.093264in}}%
\pgfpathlineto{\pgfqpoint{3.608267in}{2.090820in}}%
\pgfpathlineto{\pgfqpoint{3.624167in}{2.076668in}}%
\pgfpathlineto{\pgfqpoint{3.623267in}{2.068423in}}%
\pgfpathlineto{\pgfqpoint{3.634098in}{2.058650in}}%
\pgfpathlineto{\pgfqpoint{3.642896in}{2.064257in}}%
\pgfpathlineto{\pgfqpoint{3.661307in}{2.042785in}}%
\pgfpathlineto{\pgfqpoint{3.653078in}{2.025601in}}%
\pgfpathlineto{\pgfqpoint{3.642044in}{2.025261in}}%
\pgfpathlineto{\pgfqpoint{3.633192in}{2.009922in}}%
\pgfpathlineto{\pgfqpoint{3.622468in}{2.007746in}}%
\pgfpathlineto{\pgfqpoint{3.614595in}{1.999578in}}%
\pgfpathlineto{\pgfqpoint{3.591081in}{1.990819in}}%
\pgfpathlineto{\pgfqpoint{3.576849in}{1.976702in}}%
\pgfpathlineto{\pgfqpoint{3.569620in}{1.986835in}}%
\pgfpathlineto{\pgfqpoint{3.563050in}{1.979573in}}%
\pgfpathlineto{\pgfqpoint{3.564849in}{1.950215in}}%
\pgfpathlineto{\pgfqpoint{3.559648in}{1.938586in}}%
\pgfpathlineto{\pgfqpoint{3.548302in}{1.941762in}}%
\pgfpathlineto{\pgfqpoint{3.546471in}{1.929840in}}%
\pgfpathlineto{\pgfqpoint{3.541568in}{1.924938in}}%
\pgfpathlineto{\pgfqpoint{3.531757in}{1.926132in}}%
\pgfpathlineto{\pgfqpoint{3.532375in}{1.914985in}}%
\pgfpathlineto{\pgfqpoint{3.517503in}{1.918151in}}%
\pgfpathlineto{\pgfqpoint{3.509381in}{1.902707in}}%
\pgfpathlineto{\pgfqpoint{3.512448in}{1.894627in}}%
\pgfpathlineto{\pgfqpoint{3.507617in}{1.881187in}}%
\pgfpathlineto{\pgfqpoint{3.500012in}{1.871301in}}%
\pgfpathlineto{\pgfqpoint{3.498098in}{1.850665in}}%
\pgfpathlineto{\pgfqpoint{3.488916in}{1.850686in}}%
\pgfusepath{stroke}%
\end{pgfscope}%
\begin{pgfscope}%
\pgfpathrectangle{\pgfqpoint{0.100000in}{0.100000in}}{\pgfqpoint{3.608454in}{2.310000in}}%
\pgfusepath{clip}%
\pgfsetbuttcap%
\pgfsetroundjoin%
\pgfsetlinewidth{0.010037pt}%
\definecolor{currentstroke}{rgb}{1.000000,1.000000,1.000000}%
\pgfsetstrokecolor{currentstroke}%
\pgfsetdash{}{0pt}%
\pgfpathmoveto{\pgfqpoint{3.595305in}{1.984790in}}%
\pgfpathlineto{\pgfqpoint{3.602019in}{1.991823in}}%
\pgfpathlineto{\pgfqpoint{3.608411in}{1.985213in}}%
\pgfpathlineto{\pgfqpoint{3.596929in}{1.976454in}}%
\pgfpathlineto{\pgfqpoint{3.595305in}{1.984790in}}%
\pgfusepath{stroke}%
\end{pgfscope}%
\begin{pgfscope}%
\pgfpathrectangle{\pgfqpoint{0.100000in}{0.100000in}}{\pgfqpoint{3.608454in}{2.310000in}}%
\pgfusepath{clip}%
\pgfsetbuttcap%
\pgfsetroundjoin%
\pgfsetlinewidth{0.010037pt}%
\definecolor{currentstroke}{rgb}{1.000000,1.000000,1.000000}%
\pgfsetstrokecolor{currentstroke}%
\pgfsetdash{}{0pt}%
\pgfpathmoveto{\pgfqpoint{1.606550in}{1.914676in}}%
\pgfpathlineto{\pgfqpoint{1.612175in}{1.970174in}}%
\pgfpathlineto{\pgfqpoint{1.622900in}{2.078078in}}%
\pgfpathlineto{\pgfqpoint{1.630106in}{2.157928in}}%
\pgfpathlineto{\pgfqpoint{1.689116in}{2.152392in}}%
\pgfpathlineto{\pgfqpoint{1.764609in}{2.146381in}}%
\pgfpathlineto{\pgfqpoint{1.833613in}{2.141936in}}%
\pgfpathlineto{\pgfqpoint{1.896094in}{2.138775in}}%
\pgfpathlineto{\pgfqpoint{1.989337in}{2.135576in}}%
\pgfpathlineto{\pgfqpoint{1.995700in}{2.109297in}}%
\pgfpathlineto{\pgfqpoint{1.993422in}{2.096738in}}%
\pgfpathlineto{\pgfqpoint{1.992673in}{2.070924in}}%
\pgfpathlineto{\pgfqpoint{1.997002in}{2.051591in}}%
\pgfpathlineto{\pgfqpoint{2.006920in}{2.023070in}}%
\pgfpathlineto{\pgfqpoint{2.006897in}{1.987162in}}%
\pgfpathlineto{\pgfqpoint{2.008733in}{1.945448in}}%
\pgfpathlineto{\pgfqpoint{2.011253in}{1.934205in}}%
\pgfpathlineto{\pgfqpoint{2.018625in}{1.921863in}}%
\pgfpathlineto{\pgfqpoint{2.021056in}{1.902632in}}%
\pgfpathlineto{\pgfqpoint{2.020006in}{1.889792in}}%
\pgfpathlineto{\pgfqpoint{1.941838in}{1.891485in}}%
\pgfpathlineto{\pgfqpoint{1.884971in}{1.894368in}}%
\pgfpathlineto{\pgfqpoint{1.801609in}{1.898805in}}%
\pgfpathlineto{\pgfqpoint{1.719401in}{1.904645in}}%
\pgfpathlineto{\pgfqpoint{1.664646in}{1.909062in}}%
\pgfpathlineto{\pgfqpoint{1.606550in}{1.914676in}}%
\pgfusepath{stroke}%
\end{pgfscope}%
\begin{pgfscope}%
\pgfpathrectangle{\pgfqpoint{0.100000in}{0.100000in}}{\pgfqpoint{3.608454in}{2.310000in}}%
\pgfusepath{clip}%
\pgfsetbuttcap%
\pgfsetroundjoin%
\pgfsetlinewidth{0.010037pt}%
\definecolor{currentstroke}{rgb}{1.000000,1.000000,1.000000}%
\pgfsetstrokecolor{currentstroke}%
\pgfsetdash{}{0pt}%
\pgfpathmoveto{\pgfqpoint{1.582921in}{1.682417in}}%
\pgfpathlineto{\pgfqpoint{1.591946in}{1.775447in}}%
\pgfpathlineto{\pgfqpoint{1.598460in}{1.839883in}}%
\pgfpathlineto{\pgfqpoint{1.606550in}{1.914676in}}%
\pgfpathlineto{\pgfqpoint{1.664646in}{1.909062in}}%
\pgfpathlineto{\pgfqpoint{1.719401in}{1.904645in}}%
\pgfpathlineto{\pgfqpoint{1.801609in}{1.898805in}}%
\pgfpathlineto{\pgfqpoint{1.884971in}{1.894368in}}%
\pgfpathlineto{\pgfqpoint{1.941838in}{1.891485in}}%
\pgfpathlineto{\pgfqpoint{2.020006in}{1.889792in}}%
\pgfpathlineto{\pgfqpoint{2.014714in}{1.874348in}}%
\pgfpathlineto{\pgfqpoint{2.004150in}{1.862224in}}%
\pgfpathlineto{\pgfqpoint{2.012240in}{1.848265in}}%
\pgfpathlineto{\pgfqpoint{2.021160in}{1.845282in}}%
\pgfpathlineto{\pgfqpoint{2.025393in}{1.837264in}}%
\pgfpathlineto{\pgfqpoint{2.024421in}{1.778721in}}%
\pgfpathlineto{\pgfqpoint{2.022794in}{1.696330in}}%
\pgfpathlineto{\pgfqpoint{2.016778in}{1.674690in}}%
\pgfpathlineto{\pgfqpoint{2.022159in}{1.662720in}}%
\pgfpathlineto{\pgfqpoint{2.011355in}{1.636997in}}%
\pgfpathlineto{\pgfqpoint{2.022737in}{1.616262in}}%
\pgfpathlineto{\pgfqpoint{2.013067in}{1.617849in}}%
\pgfpathlineto{\pgfqpoint{2.006469in}{1.630750in}}%
\pgfpathlineto{\pgfqpoint{1.982922in}{1.639593in}}%
\pgfpathlineto{\pgfqpoint{1.968073in}{1.647372in}}%
\pgfpathlineto{\pgfqpoint{1.943154in}{1.648019in}}%
\pgfpathlineto{\pgfqpoint{1.934505in}{1.640932in}}%
\pgfpathlineto{\pgfqpoint{1.906300in}{1.654886in}}%
\pgfpathlineto{\pgfqpoint{1.904136in}{1.659299in}}%
\pgfpathlineto{\pgfqpoint{1.805703in}{1.663958in}}%
\pgfpathlineto{\pgfqpoint{1.745900in}{1.667386in}}%
\pgfpathlineto{\pgfqpoint{1.696503in}{1.671254in}}%
\pgfpathlineto{\pgfqpoint{1.614880in}{1.678966in}}%
\pgfpathlineto{\pgfqpoint{1.582921in}{1.682417in}}%
\pgfusepath{stroke}%
\end{pgfscope}%
\begin{pgfscope}%
\pgfpathrectangle{\pgfqpoint{0.100000in}{0.100000in}}{\pgfqpoint{3.608454in}{2.310000in}}%
\pgfusepath{clip}%
\pgfsetbuttcap%
\pgfsetroundjoin%
\pgfsetlinewidth{0.010037pt}%
\definecolor{currentstroke}{rgb}{1.000000,1.000000,1.000000}%
\pgfsetstrokecolor{currentstroke}%
\pgfsetdash{}{0pt}%
\pgfpathmoveto{\pgfqpoint{1.567441in}{1.524892in}}%
\pgfpathlineto{\pgfqpoint{1.514964in}{1.529684in}}%
\pgfpathlineto{\pgfqpoint{1.400578in}{1.543802in}}%
\pgfpathlineto{\pgfqpoint{1.338353in}{1.552733in}}%
\pgfpathlineto{\pgfqpoint{1.271614in}{1.562313in}}%
\pgfpathlineto{\pgfqpoint{1.215397in}{1.571309in}}%
\pgfpathlineto{\pgfqpoint{1.153691in}{1.581874in}}%
\pgfpathlineto{\pgfqpoint{1.167721in}{1.659701in}}%
\pgfpathlineto{\pgfqpoint{1.181932in}{1.739505in}}%
\pgfpathlineto{\pgfqpoint{1.202053in}{1.853008in}}%
\pgfpathlineto{\pgfqpoint{1.209241in}{1.893876in}}%
\pgfpathlineto{\pgfqpoint{1.243728in}{1.887567in}}%
\pgfpathlineto{\pgfqpoint{1.343910in}{1.871841in}}%
\pgfpathlineto{\pgfqpoint{1.384573in}{1.866040in}}%
\pgfpathlineto{\pgfqpoint{1.475618in}{1.853580in}}%
\pgfpathlineto{\pgfqpoint{1.541187in}{1.846137in}}%
\pgfpathlineto{\pgfqpoint{1.598460in}{1.839883in}}%
\pgfpathlineto{\pgfqpoint{1.591946in}{1.775447in}}%
\pgfpathlineto{\pgfqpoint{1.582921in}{1.682417in}}%
\pgfpathlineto{\pgfqpoint{1.575175in}{1.603355in}}%
\pgfpathlineto{\pgfqpoint{1.567441in}{1.524892in}}%
\pgfusepath{stroke}%
\end{pgfscope}%
\begin{pgfscope}%
\pgfpathrectangle{\pgfqpoint{0.100000in}{0.100000in}}{\pgfqpoint{3.608454in}{2.310000in}}%
\pgfusepath{clip}%
\pgfsetbuttcap%
\pgfsetroundjoin%
\pgfsetlinewidth{0.010037pt}%
\definecolor{currentstroke}{rgb}{1.000000,1.000000,1.000000}%
\pgfsetstrokecolor{currentstroke}%
\pgfsetdash{}{0pt}%
\pgfpathmoveto{\pgfqpoint{2.526890in}{1.632246in}}%
\pgfpathlineto{\pgfqpoint{2.460390in}{1.627522in}}%
\pgfpathlineto{\pgfqpoint{2.361269in}{1.623298in}}%
\pgfpathlineto{\pgfqpoint{2.357502in}{1.633311in}}%
\pgfpathlineto{\pgfqpoint{2.335521in}{1.640795in}}%
\pgfpathlineto{\pgfqpoint{2.330673in}{1.654936in}}%
\pgfpathlineto{\pgfqpoint{2.328643in}{1.672430in}}%
\pgfpathlineto{\pgfqpoint{2.333596in}{1.681393in}}%
\pgfpathlineto{\pgfqpoint{2.325773in}{1.689994in}}%
\pgfpathlineto{\pgfqpoint{2.323886in}{1.700253in}}%
\pgfpathlineto{\pgfqpoint{2.321364in}{1.722947in}}%
\pgfpathlineto{\pgfqpoint{2.313867in}{1.735275in}}%
\pgfpathlineto{\pgfqpoint{2.300552in}{1.742194in}}%
\pgfpathlineto{\pgfqpoint{2.286062in}{1.753611in}}%
\pgfpathlineto{\pgfqpoint{2.278572in}{1.767120in}}%
\pgfpathlineto{\pgfqpoint{2.265132in}{1.772560in}}%
\pgfpathlineto{\pgfqpoint{2.257232in}{1.781413in}}%
\pgfpathlineto{\pgfqpoint{2.247661in}{1.782916in}}%
\pgfpathlineto{\pgfqpoint{2.230573in}{1.796062in}}%
\pgfpathlineto{\pgfqpoint{2.233349in}{1.811190in}}%
\pgfpathlineto{\pgfqpoint{2.232814in}{1.839958in}}%
\pgfpathlineto{\pgfqpoint{2.238084in}{1.847876in}}%
\pgfpathlineto{\pgfqpoint{2.233349in}{1.859834in}}%
\pgfpathlineto{\pgfqpoint{2.225008in}{1.862146in}}%
\pgfpathlineto{\pgfqpoint{2.225698in}{1.872648in}}%
\pgfpathlineto{\pgfqpoint{2.236046in}{1.889238in}}%
\pgfpathlineto{\pgfqpoint{2.256539in}{1.902357in}}%
\pgfpathlineto{\pgfqpoint{2.255262in}{1.949012in}}%
\pgfpathlineto{\pgfqpoint{2.265516in}{1.956021in}}%
\pgfpathlineto{\pgfqpoint{2.275220in}{1.951352in}}%
\pgfpathlineto{\pgfqpoint{2.294988in}{1.958187in}}%
\pgfpathlineto{\pgfqpoint{2.332179in}{1.975363in}}%
\pgfpathlineto{\pgfqpoint{2.337024in}{1.970045in}}%
\pgfpathlineto{\pgfqpoint{2.329980in}{1.945955in}}%
\pgfpathlineto{\pgfqpoint{2.340460in}{1.951242in}}%
\pgfpathlineto{\pgfqpoint{2.358409in}{1.945963in}}%
\pgfpathlineto{\pgfqpoint{2.369439in}{1.941557in}}%
\pgfpathlineto{\pgfqpoint{2.375622in}{1.928645in}}%
\pgfpathlineto{\pgfqpoint{2.432207in}{1.916460in}}%
\pgfpathlineto{\pgfqpoint{2.449120in}{1.908100in}}%
\pgfpathlineto{\pgfqpoint{2.466317in}{1.908198in}}%
\pgfpathlineto{\pgfqpoint{2.483987in}{1.904750in}}%
\pgfpathlineto{\pgfqpoint{2.495454in}{1.892970in}}%
\pgfpathlineto{\pgfqpoint{2.506412in}{1.887150in}}%
\pgfpathlineto{\pgfqpoint{2.508426in}{1.870361in}}%
\pgfpathlineto{\pgfqpoint{2.505188in}{1.859820in}}%
\pgfpathlineto{\pgfqpoint{2.517401in}{1.859041in}}%
\pgfpathlineto{\pgfqpoint{2.513287in}{1.846815in}}%
\pgfpathlineto{\pgfqpoint{2.517200in}{1.842472in}}%
\pgfpathlineto{\pgfqpoint{2.521092in}{1.830931in}}%
\pgfpathlineto{\pgfqpoint{2.509190in}{1.824821in}}%
\pgfpathlineto{\pgfqpoint{2.502280in}{1.807797in}}%
\pgfpathlineto{\pgfqpoint{2.500111in}{1.795759in}}%
\pgfpathlineto{\pgfqpoint{2.506741in}{1.793698in}}%
\pgfpathlineto{\pgfqpoint{2.515230in}{1.802725in}}%
\pgfpathlineto{\pgfqpoint{2.522446in}{1.818369in}}%
\pgfpathlineto{\pgfqpoint{2.532211in}{1.823802in}}%
\pgfpathlineto{\pgfqpoint{2.539548in}{1.816735in}}%
\pgfpathlineto{\pgfqpoint{2.532326in}{1.795344in}}%
\pgfpathlineto{\pgfqpoint{2.530060in}{1.778738in}}%
\pgfpathlineto{\pgfqpoint{2.532194in}{1.766801in}}%
\pgfpathlineto{\pgfqpoint{2.525487in}{1.760048in}}%
\pgfpathlineto{\pgfqpoint{2.522137in}{1.743546in}}%
\pgfpathlineto{\pgfqpoint{2.524874in}{1.726204in}}%
\pgfpathlineto{\pgfqpoint{2.516958in}{1.700555in}}%
\pgfpathlineto{\pgfqpoint{2.517125in}{1.687738in}}%
\pgfpathlineto{\pgfqpoint{2.523380in}{1.659966in}}%
\pgfpathlineto{\pgfqpoint{2.527434in}{1.655201in}}%
\pgfpathlineto{\pgfqpoint{2.526890in}{1.632246in}}%
\pgfusepath{stroke}%
\end{pgfscope}%
\begin{pgfscope}%
\pgfpathrectangle{\pgfqpoint{0.100000in}{0.100000in}}{\pgfqpoint{3.608454in}{2.310000in}}%
\pgfusepath{clip}%
\pgfsetbuttcap%
\pgfsetroundjoin%
\pgfsetlinewidth{0.010037pt}%
\definecolor{currentstroke}{rgb}{1.000000,1.000000,1.000000}%
\pgfsetstrokecolor{currentstroke}%
\pgfsetdash{}{0pt}%
\pgfpathmoveto{\pgfqpoint{2.551823in}{1.857339in}}%
\pgfpathlineto{\pgfqpoint{2.553308in}{1.846261in}}%
\pgfpathlineto{\pgfqpoint{2.539709in}{1.817072in}}%
\pgfpathlineto{\pgfqpoint{2.533666in}{1.825527in}}%
\pgfpathlineto{\pgfqpoint{2.551823in}{1.857339in}}%
\pgfusepath{stroke}%
\end{pgfscope}%
\begin{pgfscope}%
\pgfpathrectangle{\pgfqpoint{0.100000in}{0.100000in}}{\pgfqpoint{3.608454in}{2.310000in}}%
\pgfusepath{clip}%
\pgfsetbuttcap%
\pgfsetroundjoin%
\pgfsetlinewidth{0.010037pt}%
\definecolor{currentstroke}{rgb}{1.000000,1.000000,1.000000}%
\pgfsetstrokecolor{currentstroke}%
\pgfsetdash{}{0pt}%
\pgfpathmoveto{\pgfqpoint{0.905799in}{2.039449in}}%
\pgfpathlineto{\pgfqpoint{0.908487in}{2.052757in}}%
\pgfpathlineto{\pgfqpoint{0.906458in}{2.067206in}}%
\pgfpathlineto{\pgfqpoint{0.909565in}{2.082958in}}%
\pgfpathlineto{\pgfqpoint{0.925385in}{2.146869in}}%
\pgfpathlineto{\pgfqpoint{0.938691in}{2.199997in}}%
\pgfpathlineto{\pgfqpoint{0.957404in}{2.274144in}}%
\pgfpathlineto{\pgfqpoint{1.007001in}{2.262160in}}%
\pgfpathlineto{\pgfqpoint{0.988670in}{2.182155in}}%
\pgfpathlineto{\pgfqpoint{1.000279in}{2.156377in}}%
\pgfpathlineto{\pgfqpoint{0.998811in}{2.142381in}}%
\pgfpathlineto{\pgfqpoint{1.005033in}{2.121323in}}%
\pgfpathlineto{\pgfqpoint{1.013723in}{2.116347in}}%
\pgfpathlineto{\pgfqpoint{1.023232in}{2.091949in}}%
\pgfpathlineto{\pgfqpoint{1.029074in}{2.085229in}}%
\pgfpathlineto{\pgfqpoint{1.028366in}{2.076578in}}%
\pgfpathlineto{\pgfqpoint{1.039440in}{2.068841in}}%
\pgfpathlineto{\pgfqpoint{1.041003in}{2.059867in}}%
\pgfpathlineto{\pgfqpoint{1.056205in}{2.058614in}}%
\pgfpathlineto{\pgfqpoint{1.036925in}{2.008611in}}%
\pgfpathlineto{\pgfqpoint{1.039874in}{1.998281in}}%
\pgfpathlineto{\pgfqpoint{1.028782in}{1.991857in}}%
\pgfpathlineto{\pgfqpoint{1.027278in}{1.975551in}}%
\pgfpathlineto{\pgfqpoint{1.036527in}{1.965340in}}%
\pgfpathlineto{\pgfqpoint{1.048676in}{1.970294in}}%
\pgfpathlineto{\pgfqpoint{1.057602in}{1.979454in}}%
\pgfpathlineto{\pgfqpoint{1.067116in}{1.969101in}}%
\pgfpathlineto{\pgfqpoint{1.066370in}{1.947947in}}%
\pgfpathlineto{\pgfqpoint{1.077439in}{1.923062in}}%
\pgfpathlineto{\pgfqpoint{1.074904in}{1.907997in}}%
\pgfpathlineto{\pgfqpoint{1.090449in}{1.897544in}}%
\pgfpathlineto{\pgfqpoint{1.090697in}{1.883009in}}%
\pgfpathlineto{\pgfqpoint{1.095991in}{1.867733in}}%
\pgfpathlineto{\pgfqpoint{1.105833in}{1.871176in}}%
\pgfpathlineto{\pgfqpoint{1.128385in}{1.864765in}}%
\pgfpathlineto{\pgfqpoint{1.137669in}{1.871889in}}%
\pgfpathlineto{\pgfqpoint{1.144309in}{1.867284in}}%
\pgfpathlineto{\pgfqpoint{1.157517in}{1.867860in}}%
\pgfpathlineto{\pgfqpoint{1.160350in}{1.863105in}}%
\pgfpathlineto{\pgfqpoint{1.174451in}{1.864434in}}%
\pgfpathlineto{\pgfqpoint{1.181245in}{1.874373in}}%
\pgfpathlineto{\pgfqpoint{1.190738in}{1.875584in}}%
\pgfpathlineto{\pgfqpoint{1.198123in}{1.855924in}}%
\pgfpathlineto{\pgfqpoint{1.202053in}{1.853008in}}%
\pgfpathlineto{\pgfqpoint{1.181932in}{1.739505in}}%
\pgfpathlineto{\pgfqpoint{1.167721in}{1.659701in}}%
\pgfpathlineto{\pgfqpoint{1.055637in}{1.681339in}}%
\pgfpathlineto{\pgfqpoint{0.995094in}{1.693381in}}%
\pgfpathlineto{\pgfqpoint{0.938464in}{1.705834in}}%
\pgfpathlineto{\pgfqpoint{0.885533in}{1.717818in}}%
\pgfpathlineto{\pgfqpoint{0.824275in}{1.732591in}}%
\pgfpathlineto{\pgfqpoint{0.855872in}{1.862213in}}%
\pgfpathlineto{\pgfqpoint{0.857576in}{1.871686in}}%
\pgfpathlineto{\pgfqpoint{0.871627in}{1.896508in}}%
\pgfpathlineto{\pgfqpoint{0.857078in}{1.911427in}}%
\pgfpathlineto{\pgfqpoint{0.860099in}{1.926077in}}%
\pgfpathlineto{\pgfqpoint{0.865637in}{1.929744in}}%
\pgfpathlineto{\pgfqpoint{0.875529in}{1.944840in}}%
\pgfpathlineto{\pgfqpoint{0.889947in}{1.955284in}}%
\pgfpathlineto{\pgfqpoint{0.890735in}{1.963012in}}%
\pgfpathlineto{\pgfqpoint{0.899437in}{1.970733in}}%
\pgfpathlineto{\pgfqpoint{0.906663in}{1.985185in}}%
\pgfpathlineto{\pgfqpoint{0.921484in}{2.000475in}}%
\pgfpathlineto{\pgfqpoint{0.920432in}{2.015576in}}%
\pgfpathlineto{\pgfqpoint{0.910299in}{2.023965in}}%
\pgfpathlineto{\pgfqpoint{0.905799in}{2.039449in}}%
\pgfusepath{stroke}%
\end{pgfscope}%
\begin{pgfscope}%
\pgfpathrectangle{\pgfqpoint{0.100000in}{0.100000in}}{\pgfqpoint{3.608454in}{2.310000in}}%
\pgfusepath{clip}%
\pgfsetbuttcap%
\pgfsetroundjoin%
\pgfsetlinewidth{0.010037pt}%
\definecolor{currentstroke}{rgb}{1.000000,1.000000,1.000000}%
\pgfsetstrokecolor{currentstroke}%
\pgfsetdash{}{0pt}%
\pgfpathmoveto{\pgfqpoint{3.359367in}{1.786492in}}%
\pgfpathlineto{\pgfqpoint{3.354710in}{1.801193in}}%
\pgfpathlineto{\pgfqpoint{3.346067in}{1.845790in}}%
\pgfpathlineto{\pgfqpoint{3.334894in}{1.859510in}}%
\pgfpathlineto{\pgfqpoint{3.325093in}{1.884184in}}%
\pgfpathlineto{\pgfqpoint{3.328311in}{1.902477in}}%
\pgfpathlineto{\pgfqpoint{3.325730in}{1.915960in}}%
\pgfpathlineto{\pgfqpoint{3.317356in}{1.929202in}}%
\pgfpathlineto{\pgfqpoint{3.317011in}{1.943790in}}%
\pgfpathlineto{\pgfqpoint{3.312149in}{1.959523in}}%
\pgfpathlineto{\pgfqpoint{3.355664in}{1.970232in}}%
\pgfpathlineto{\pgfqpoint{3.412185in}{1.985489in}}%
\pgfpathlineto{\pgfqpoint{3.414441in}{1.976741in}}%
\pgfpathlineto{\pgfqpoint{3.410842in}{1.962813in}}%
\pgfpathlineto{\pgfqpoint{3.419291in}{1.951669in}}%
\pgfpathlineto{\pgfqpoint{3.414804in}{1.937548in}}%
\pgfpathlineto{\pgfqpoint{3.396973in}{1.919797in}}%
\pgfpathlineto{\pgfqpoint{3.401874in}{1.906464in}}%
\pgfpathlineto{\pgfqpoint{3.398242in}{1.878280in}}%
\pgfpathlineto{\pgfqpoint{3.393157in}{1.862389in}}%
\pgfpathlineto{\pgfqpoint{3.399644in}{1.817200in}}%
\pgfpathlineto{\pgfqpoint{3.398770in}{1.801301in}}%
\pgfpathlineto{\pgfqpoint{3.405047in}{1.796199in}}%
\pgfpathlineto{\pgfqpoint{3.359367in}{1.786492in}}%
\pgfusepath{stroke}%
\end{pgfscope}%
\begin{pgfscope}%
\pgfpathrectangle{\pgfqpoint{0.100000in}{0.100000in}}{\pgfqpoint{3.608454in}{2.310000in}}%
\pgfusepath{clip}%
\pgfsetbuttcap%
\pgfsetroundjoin%
\pgfsetlinewidth{0.010037pt}%
\definecolor{currentstroke}{rgb}{1.000000,1.000000,1.000000}%
\pgfsetstrokecolor{currentstroke}%
\pgfsetdash{}{0pt}%
\pgfpathmoveto{\pgfqpoint{2.022794in}{1.696330in}}%
\pgfpathlineto{\pgfqpoint{2.024421in}{1.778721in}}%
\pgfpathlineto{\pgfqpoint{2.025393in}{1.837264in}}%
\pgfpathlineto{\pgfqpoint{2.021160in}{1.845282in}}%
\pgfpathlineto{\pgfqpoint{2.012240in}{1.848265in}}%
\pgfpathlineto{\pgfqpoint{2.004150in}{1.862224in}}%
\pgfpathlineto{\pgfqpoint{2.014714in}{1.874348in}}%
\pgfpathlineto{\pgfqpoint{2.020006in}{1.889792in}}%
\pgfpathlineto{\pgfqpoint{2.021056in}{1.902632in}}%
\pgfpathlineto{\pgfqpoint{2.018625in}{1.921863in}}%
\pgfpathlineto{\pgfqpoint{2.011253in}{1.934205in}}%
\pgfpathlineto{\pgfqpoint{2.008733in}{1.945448in}}%
\pgfpathlineto{\pgfqpoint{2.006897in}{1.987162in}}%
\pgfpathlineto{\pgfqpoint{2.006920in}{2.023070in}}%
\pgfpathlineto{\pgfqpoint{1.997002in}{2.051591in}}%
\pgfpathlineto{\pgfqpoint{1.992673in}{2.070924in}}%
\pgfpathlineto{\pgfqpoint{1.993422in}{2.096738in}}%
\pgfpathlineto{\pgfqpoint{1.995700in}{2.109297in}}%
\pgfpathlineto{\pgfqpoint{1.989337in}{2.135576in}}%
\pgfpathlineto{\pgfqpoint{2.032661in}{2.134709in}}%
\pgfpathlineto{\pgfqpoint{2.098467in}{2.134144in}}%
\pgfpathlineto{\pgfqpoint{2.098827in}{2.164013in}}%
\pgfpathlineto{\pgfqpoint{2.115577in}{2.160724in}}%
\pgfpathlineto{\pgfqpoint{2.123605in}{2.124304in}}%
\pgfpathlineto{\pgfqpoint{2.129524in}{2.111209in}}%
\pgfpathlineto{\pgfqpoint{2.144242in}{2.110823in}}%
\pgfpathlineto{\pgfqpoint{2.147536in}{2.106379in}}%
\pgfpathlineto{\pgfqpoint{2.168066in}{2.104411in}}%
\pgfpathlineto{\pgfqpoint{2.171516in}{2.095382in}}%
\pgfpathlineto{\pgfqpoint{2.185665in}{2.097411in}}%
\pgfpathlineto{\pgfqpoint{2.196654in}{2.105856in}}%
\pgfpathlineto{\pgfqpoint{2.215604in}{2.105540in}}%
\pgfpathlineto{\pgfqpoint{2.227330in}{2.098749in}}%
\pgfpathlineto{\pgfqpoint{2.228678in}{2.092380in}}%
\pgfpathlineto{\pgfqpoint{2.239835in}{2.091046in}}%
\pgfpathlineto{\pgfqpoint{2.247116in}{2.073673in}}%
\pgfpathlineto{\pgfqpoint{2.251814in}{2.084355in}}%
\pgfpathlineto{\pgfqpoint{2.264631in}{2.085012in}}%
\pgfpathlineto{\pgfqpoint{2.267873in}{2.076692in}}%
\pgfpathlineto{\pgfqpoint{2.282295in}{2.072895in}}%
\pgfpathlineto{\pgfqpoint{2.290247in}{2.060912in}}%
\pgfpathlineto{\pgfqpoint{2.307885in}{2.064633in}}%
\pgfpathlineto{\pgfqpoint{2.327294in}{2.079338in}}%
\pgfpathlineto{\pgfqpoint{2.334370in}{2.066375in}}%
\pgfpathlineto{\pgfqpoint{2.366216in}{2.069923in}}%
\pgfpathlineto{\pgfqpoint{2.379829in}{2.061026in}}%
\pgfpathlineto{\pgfqpoint{2.387752in}{2.064202in}}%
\pgfpathlineto{\pgfqpoint{2.394106in}{2.059194in}}%
\pgfpathlineto{\pgfqpoint{2.375257in}{2.047310in}}%
\pgfpathlineto{\pgfqpoint{2.348353in}{2.036722in}}%
\pgfpathlineto{\pgfqpoint{2.321706in}{2.015558in}}%
\pgfpathlineto{\pgfqpoint{2.298605in}{1.987726in}}%
\pgfpathlineto{\pgfqpoint{2.281117in}{1.971272in}}%
\pgfpathlineto{\pgfqpoint{2.265799in}{1.959963in}}%
\pgfpathlineto{\pgfqpoint{2.255262in}{1.949012in}}%
\pgfpathlineto{\pgfqpoint{2.256539in}{1.902357in}}%
\pgfpathlineto{\pgfqpoint{2.236046in}{1.889238in}}%
\pgfpathlineto{\pgfqpoint{2.225698in}{1.872648in}}%
\pgfpathlineto{\pgfqpoint{2.225008in}{1.862146in}}%
\pgfpathlineto{\pgfqpoint{2.233349in}{1.859834in}}%
\pgfpathlineto{\pgfqpoint{2.238084in}{1.847876in}}%
\pgfpathlineto{\pgfqpoint{2.232814in}{1.839958in}}%
\pgfpathlineto{\pgfqpoint{2.233349in}{1.811190in}}%
\pgfpathlineto{\pgfqpoint{2.230573in}{1.796062in}}%
\pgfpathlineto{\pgfqpoint{2.247661in}{1.782916in}}%
\pgfpathlineto{\pgfqpoint{2.257232in}{1.781413in}}%
\pgfpathlineto{\pgfqpoint{2.265132in}{1.772560in}}%
\pgfpathlineto{\pgfqpoint{2.278572in}{1.767120in}}%
\pgfpathlineto{\pgfqpoint{2.286062in}{1.753611in}}%
\pgfpathlineto{\pgfqpoint{2.300552in}{1.742194in}}%
\pgfpathlineto{\pgfqpoint{2.313867in}{1.735275in}}%
\pgfpathlineto{\pgfqpoint{2.321364in}{1.722947in}}%
\pgfpathlineto{\pgfqpoint{2.323886in}{1.700253in}}%
\pgfpathlineto{\pgfqpoint{2.253222in}{1.697687in}}%
\pgfpathlineto{\pgfqpoint{2.192987in}{1.696427in}}%
\pgfpathlineto{\pgfqpoint{2.138106in}{1.695621in}}%
\pgfpathlineto{\pgfqpoint{2.080050in}{1.695709in}}%
\pgfpathlineto{\pgfqpoint{2.022794in}{1.696330in}}%
\pgfusepath{stroke}%
\end{pgfscope}%
\begin{pgfscope}%
\pgfpathrectangle{\pgfqpoint{0.100000in}{0.100000in}}{\pgfqpoint{3.608454in}{2.310000in}}%
\pgfusepath{clip}%
\pgfsetbuttcap%
\pgfsetroundjoin%
\pgfsetlinewidth{0.010037pt}%
\definecolor{currentstroke}{rgb}{1.000000,1.000000,1.000000}%
\pgfsetstrokecolor{currentstroke}%
\pgfsetdash{}{0pt}%
\pgfpathmoveto{\pgfqpoint{0.418363in}{1.850356in}}%
\pgfpathlineto{\pgfqpoint{0.412766in}{1.860654in}}%
\pgfpathlineto{\pgfqpoint{0.416340in}{1.887087in}}%
\pgfpathlineto{\pgfqpoint{0.423237in}{1.900901in}}%
\pgfpathlineto{\pgfqpoint{0.419789in}{1.919587in}}%
\pgfpathlineto{\pgfqpoint{0.426984in}{1.927445in}}%
\pgfpathlineto{\pgfqpoint{0.440110in}{1.948627in}}%
\pgfpathlineto{\pgfqpoint{0.451240in}{1.961413in}}%
\pgfpathlineto{\pgfqpoint{0.467509in}{1.989261in}}%
\pgfpathlineto{\pgfqpoint{0.479986in}{2.019596in}}%
\pgfpathlineto{\pgfqpoint{0.493232in}{2.048048in}}%
\pgfpathlineto{\pgfqpoint{0.495910in}{2.059930in}}%
\pgfpathlineto{\pgfqpoint{0.514183in}{2.093924in}}%
\pgfpathlineto{\pgfqpoint{0.517702in}{2.108854in}}%
\pgfpathlineto{\pgfqpoint{0.525482in}{2.124532in}}%
\pgfpathlineto{\pgfqpoint{0.528263in}{2.143657in}}%
\pgfpathlineto{\pgfqpoint{0.543116in}{2.152982in}}%
\pgfpathlineto{\pgfqpoint{0.560727in}{2.157683in}}%
\pgfpathlineto{\pgfqpoint{0.569631in}{2.147155in}}%
\pgfpathlineto{\pgfqpoint{0.577330in}{2.147934in}}%
\pgfpathlineto{\pgfqpoint{0.589348in}{2.135386in}}%
\pgfpathlineto{\pgfqpoint{0.590946in}{2.123567in}}%
\pgfpathlineto{\pgfqpoint{0.586388in}{2.100423in}}%
\pgfpathlineto{\pgfqpoint{0.601553in}{2.088614in}}%
\pgfpathlineto{\pgfqpoint{0.611388in}{2.084174in}}%
\pgfpathlineto{\pgfqpoint{0.638032in}{2.088706in}}%
\pgfpathlineto{\pgfqpoint{0.653468in}{2.085611in}}%
\pgfpathlineto{\pgfqpoint{0.664077in}{2.080532in}}%
\pgfpathlineto{\pgfqpoint{0.669666in}{2.070891in}}%
\pgfpathlineto{\pgfqpoint{0.705053in}{2.072177in}}%
\pgfpathlineto{\pgfqpoint{0.710680in}{2.066278in}}%
\pgfpathlineto{\pgfqpoint{0.723956in}{2.065005in}}%
\pgfpathlineto{\pgfqpoint{0.737324in}{2.068774in}}%
\pgfpathlineto{\pgfqpoint{0.776851in}{2.067824in}}%
\pgfpathlineto{\pgfqpoint{0.785388in}{2.064919in}}%
\pgfpathlineto{\pgfqpoint{0.795364in}{2.068266in}}%
\pgfpathlineto{\pgfqpoint{0.905799in}{2.039449in}}%
\pgfpathlineto{\pgfqpoint{0.910299in}{2.023965in}}%
\pgfpathlineto{\pgfqpoint{0.920432in}{2.015576in}}%
\pgfpathlineto{\pgfqpoint{0.921484in}{2.000475in}}%
\pgfpathlineto{\pgfqpoint{0.906663in}{1.985185in}}%
\pgfpathlineto{\pgfqpoint{0.899437in}{1.970733in}}%
\pgfpathlineto{\pgfqpoint{0.890735in}{1.963012in}}%
\pgfpathlineto{\pgfqpoint{0.889947in}{1.955284in}}%
\pgfpathlineto{\pgfqpoint{0.875529in}{1.944840in}}%
\pgfpathlineto{\pgfqpoint{0.865637in}{1.929744in}}%
\pgfpathlineto{\pgfqpoint{0.860099in}{1.926077in}}%
\pgfpathlineto{\pgfqpoint{0.857078in}{1.911427in}}%
\pgfpathlineto{\pgfqpoint{0.871627in}{1.896508in}}%
\pgfpathlineto{\pgfqpoint{0.857576in}{1.871686in}}%
\pgfpathlineto{\pgfqpoint{0.855872in}{1.862213in}}%
\pgfpathlineto{\pgfqpoint{0.824275in}{1.732591in}}%
\pgfpathlineto{\pgfqpoint{0.757813in}{1.749621in}}%
\pgfpathlineto{\pgfqpoint{0.693698in}{1.766153in}}%
\pgfpathlineto{\pgfqpoint{0.655019in}{1.776916in}}%
\pgfpathlineto{\pgfqpoint{0.605321in}{1.791073in}}%
\pgfpathlineto{\pgfqpoint{0.526049in}{1.816024in}}%
\pgfpathlineto{\pgfqpoint{0.439875in}{1.842853in}}%
\pgfpathlineto{\pgfqpoint{0.418363in}{1.850356in}}%
\pgfusepath{stroke}%
\end{pgfscope}%
\begin{pgfscope}%
\pgfpathrectangle{\pgfqpoint{0.100000in}{0.100000in}}{\pgfqpoint{3.608454in}{2.310000in}}%
\pgfusepath{clip}%
\pgfsetbuttcap%
\pgfsetroundjoin%
\pgfsetlinewidth{0.010037pt}%
\definecolor{currentstroke}{rgb}{1.000000,1.000000,1.000000}%
\pgfsetstrokecolor{currentstroke}%
\pgfsetdash{}{0pt}%
\pgfpathmoveto{\pgfqpoint{3.405047in}{1.796199in}}%
\pgfpathlineto{\pgfqpoint{3.398770in}{1.801301in}}%
\pgfpathlineto{\pgfqpoint{3.399644in}{1.817200in}}%
\pgfpathlineto{\pgfqpoint{3.393157in}{1.862389in}}%
\pgfpathlineto{\pgfqpoint{3.398242in}{1.878280in}}%
\pgfpathlineto{\pgfqpoint{3.401874in}{1.906464in}}%
\pgfpathlineto{\pgfqpoint{3.396973in}{1.919797in}}%
\pgfpathlineto{\pgfqpoint{3.414804in}{1.937548in}}%
\pgfpathlineto{\pgfqpoint{3.419291in}{1.951669in}}%
\pgfpathlineto{\pgfqpoint{3.410842in}{1.962813in}}%
\pgfpathlineto{\pgfqpoint{3.414441in}{1.976741in}}%
\pgfpathlineto{\pgfqpoint{3.412185in}{1.985489in}}%
\pgfpathlineto{\pgfqpoint{3.414122in}{2.004225in}}%
\pgfpathlineto{\pgfqpoint{3.417744in}{2.010009in}}%
\pgfpathlineto{\pgfqpoint{3.428911in}{2.013660in}}%
\pgfpathlineto{\pgfqpoint{3.445201in}{1.965955in}}%
\pgfpathlineto{\pgfqpoint{3.475422in}{1.867100in}}%
\pgfpathlineto{\pgfqpoint{3.486858in}{1.859435in}}%
\pgfpathlineto{\pgfqpoint{3.488916in}{1.850686in}}%
\pgfpathlineto{\pgfqpoint{3.494952in}{1.847150in}}%
\pgfpathlineto{\pgfqpoint{3.494493in}{1.831282in}}%
\pgfpathlineto{\pgfqpoint{3.488098in}{1.831027in}}%
\pgfpathlineto{\pgfqpoint{3.475140in}{1.821145in}}%
\pgfpathlineto{\pgfqpoint{3.471404in}{1.811231in}}%
\pgfpathlineto{\pgfqpoint{3.436723in}{1.802651in}}%
\pgfpathlineto{\pgfqpoint{3.405047in}{1.796199in}}%
\pgfusepath{stroke}%
\end{pgfscope}%
\begin{pgfscope}%
\pgfpathrectangle{\pgfqpoint{0.100000in}{0.100000in}}{\pgfqpoint{3.608454in}{2.310000in}}%
\pgfusepath{clip}%
\pgfsetbuttcap%
\pgfsetroundjoin%
\pgfsetlinewidth{0.010037pt}%
\definecolor{currentstroke}{rgb}{1.000000,1.000000,1.000000}%
\pgfsetstrokecolor{currentstroke}%
\pgfsetdash{}{0pt}%
\pgfpathmoveto{\pgfqpoint{2.320622in}{1.452274in}}%
\pgfpathlineto{\pgfqpoint{2.302315in}{1.470405in}}%
\pgfpathlineto{\pgfqpoint{2.243792in}{1.467029in}}%
\pgfpathlineto{\pgfqpoint{2.152519in}{1.464021in}}%
\pgfpathlineto{\pgfqpoint{2.060697in}{1.465454in}}%
\pgfpathlineto{\pgfqpoint{2.054252in}{1.476692in}}%
\pgfpathlineto{\pgfqpoint{2.056883in}{1.487728in}}%
\pgfpathlineto{\pgfqpoint{2.055634in}{1.506632in}}%
\pgfpathlineto{\pgfqpoint{2.050781in}{1.524944in}}%
\pgfpathlineto{\pgfqpoint{2.051345in}{1.534613in}}%
\pgfpathlineto{\pgfqpoint{2.041030in}{1.545539in}}%
\pgfpathlineto{\pgfqpoint{2.043274in}{1.560741in}}%
\pgfpathlineto{\pgfqpoint{2.027451in}{1.590769in}}%
\pgfpathlineto{\pgfqpoint{2.022737in}{1.616262in}}%
\pgfpathlineto{\pgfqpoint{2.011355in}{1.636997in}}%
\pgfpathlineto{\pgfqpoint{2.022159in}{1.662720in}}%
\pgfpathlineto{\pgfqpoint{2.016778in}{1.674690in}}%
\pgfpathlineto{\pgfqpoint{2.022794in}{1.696330in}}%
\pgfpathlineto{\pgfqpoint{2.080050in}{1.695709in}}%
\pgfpathlineto{\pgfqpoint{2.138106in}{1.695621in}}%
\pgfpathlineto{\pgfqpoint{2.192987in}{1.696427in}}%
\pgfpathlineto{\pgfqpoint{2.253222in}{1.697687in}}%
\pgfpathlineto{\pgfqpoint{2.323886in}{1.700253in}}%
\pgfpathlineto{\pgfqpoint{2.325773in}{1.689994in}}%
\pgfpathlineto{\pgfqpoint{2.333596in}{1.681393in}}%
\pgfpathlineto{\pgfqpoint{2.328643in}{1.672430in}}%
\pgfpathlineto{\pgfqpoint{2.330673in}{1.654936in}}%
\pgfpathlineto{\pgfqpoint{2.335521in}{1.640795in}}%
\pgfpathlineto{\pgfqpoint{2.357502in}{1.633311in}}%
\pgfpathlineto{\pgfqpoint{2.361269in}{1.623298in}}%
\pgfpathlineto{\pgfqpoint{2.373331in}{1.612057in}}%
\pgfpathlineto{\pgfqpoint{2.378251in}{1.600427in}}%
\pgfpathlineto{\pgfqpoint{2.390473in}{1.592626in}}%
\pgfpathlineto{\pgfqpoint{2.392386in}{1.583234in}}%
\pgfpathlineto{\pgfqpoint{2.390006in}{1.569016in}}%
\pgfpathlineto{\pgfqpoint{2.383785in}{1.564755in}}%
\pgfpathlineto{\pgfqpoint{2.381909in}{1.551220in}}%
\pgfpathlineto{\pgfqpoint{2.364032in}{1.540471in}}%
\pgfpathlineto{\pgfqpoint{2.340730in}{1.534582in}}%
\pgfpathlineto{\pgfqpoint{2.338595in}{1.521042in}}%
\pgfpathlineto{\pgfqpoint{2.347584in}{1.511368in}}%
\pgfpathlineto{\pgfqpoint{2.347951in}{1.499205in}}%
\pgfpathlineto{\pgfqpoint{2.340698in}{1.489644in}}%
\pgfpathlineto{\pgfqpoint{2.336891in}{1.475440in}}%
\pgfpathlineto{\pgfqpoint{2.324285in}{1.470738in}}%
\pgfpathlineto{\pgfqpoint{2.325089in}{1.454910in}}%
\pgfpathlineto{\pgfqpoint{2.320622in}{1.452274in}}%
\pgfusepath{stroke}%
\end{pgfscope}%
\begin{pgfscope}%
\pgfpathrectangle{\pgfqpoint{0.100000in}{0.100000in}}{\pgfqpoint{3.608454in}{2.310000in}}%
\pgfusepath{clip}%
\pgfsetbuttcap%
\pgfsetroundjoin%
\pgfsetlinewidth{0.010037pt}%
\definecolor{currentstroke}{rgb}{1.000000,1.000000,1.000000}%
\pgfsetstrokecolor{currentstroke}%
\pgfsetdash{}{0pt}%
\pgfpathmoveto{\pgfqpoint{3.456549in}{1.750456in}}%
\pgfpathlineto{\pgfqpoint{3.439208in}{1.747720in}}%
\pgfpathlineto{\pgfqpoint{3.359557in}{1.729620in}}%
\pgfpathlineto{\pgfqpoint{3.358165in}{1.731729in}}%
\pgfpathlineto{\pgfqpoint{3.359367in}{1.786492in}}%
\pgfpathlineto{\pgfqpoint{3.405047in}{1.796199in}}%
\pgfpathlineto{\pgfqpoint{3.436723in}{1.802651in}}%
\pgfpathlineto{\pgfqpoint{3.471404in}{1.811231in}}%
\pgfpathlineto{\pgfqpoint{3.475140in}{1.821145in}}%
\pgfpathlineto{\pgfqpoint{3.488098in}{1.831027in}}%
\pgfpathlineto{\pgfqpoint{3.494493in}{1.831282in}}%
\pgfpathlineto{\pgfqpoint{3.502901in}{1.816865in}}%
\pgfpathlineto{\pgfqpoint{3.495271in}{1.795821in}}%
\pgfpathlineto{\pgfqpoint{3.494152in}{1.783481in}}%
\pgfpathlineto{\pgfqpoint{3.509594in}{1.784640in}}%
\pgfpathlineto{\pgfqpoint{3.516597in}{1.778721in}}%
\pgfpathlineto{\pgfqpoint{3.532302in}{1.754525in}}%
\pgfpathlineto{\pgfqpoint{3.540107in}{1.751582in}}%
\pgfpathlineto{\pgfqpoint{3.553182in}{1.752673in}}%
\pgfpathlineto{\pgfqpoint{3.562281in}{1.760896in}}%
\pgfpathlineto{\pgfqpoint{3.568329in}{1.753558in}}%
\pgfpathlineto{\pgfqpoint{3.530295in}{1.733487in}}%
\pgfpathlineto{\pgfqpoint{3.529106in}{1.747889in}}%
\pgfpathlineto{\pgfqpoint{3.511829in}{1.725505in}}%
\pgfpathlineto{\pgfqpoint{3.505754in}{1.721649in}}%
\pgfpathlineto{\pgfqpoint{3.497281in}{1.734565in}}%
\pgfpathlineto{\pgfqpoint{3.494962in}{1.736335in}}%
\pgfpathlineto{\pgfqpoint{3.487074in}{1.740513in}}%
\pgfpathlineto{\pgfqpoint{3.480196in}{1.757459in}}%
\pgfpathlineto{\pgfqpoint{3.456549in}{1.750456in}}%
\pgfusepath{stroke}%
\end{pgfscope}%
\begin{pgfscope}%
\pgfpathrectangle{\pgfqpoint{0.100000in}{0.100000in}}{\pgfqpoint{3.608454in}{2.310000in}}%
\pgfusepath{clip}%
\pgfsetbuttcap%
\pgfsetroundjoin%
\pgfsetlinewidth{0.010037pt}%
\definecolor{currentstroke}{rgb}{1.000000,1.000000,1.000000}%
\pgfsetstrokecolor{currentstroke}%
\pgfsetdash{}{0pt}%
\pgfpathmoveto{\pgfqpoint{1.680190in}{1.435164in}}%
\pgfpathlineto{\pgfqpoint{1.686540in}{1.513906in}}%
\pgfpathlineto{\pgfqpoint{1.650576in}{1.516826in}}%
\pgfpathlineto{\pgfqpoint{1.567441in}{1.524892in}}%
\pgfpathlineto{\pgfqpoint{1.575175in}{1.603355in}}%
\pgfpathlineto{\pgfqpoint{1.582921in}{1.682417in}}%
\pgfpathlineto{\pgfqpoint{1.614880in}{1.678966in}}%
\pgfpathlineto{\pgfqpoint{1.696503in}{1.671254in}}%
\pgfpathlineto{\pgfqpoint{1.745900in}{1.667386in}}%
\pgfpathlineto{\pgfqpoint{1.805703in}{1.663958in}}%
\pgfpathlineto{\pgfqpoint{1.904136in}{1.659299in}}%
\pgfpathlineto{\pgfqpoint{1.906300in}{1.654886in}}%
\pgfpathlineto{\pgfqpoint{1.934505in}{1.640932in}}%
\pgfpathlineto{\pgfqpoint{1.943154in}{1.648019in}}%
\pgfpathlineto{\pgfqpoint{1.968073in}{1.647372in}}%
\pgfpathlineto{\pgfqpoint{1.982922in}{1.639593in}}%
\pgfpathlineto{\pgfqpoint{2.006469in}{1.630750in}}%
\pgfpathlineto{\pgfqpoint{2.013067in}{1.617849in}}%
\pgfpathlineto{\pgfqpoint{2.022737in}{1.616262in}}%
\pgfpathlineto{\pgfqpoint{2.027451in}{1.590769in}}%
\pgfpathlineto{\pgfqpoint{2.043274in}{1.560741in}}%
\pgfpathlineto{\pgfqpoint{2.041030in}{1.545539in}}%
\pgfpathlineto{\pgfqpoint{2.051345in}{1.534613in}}%
\pgfpathlineto{\pgfqpoint{2.050781in}{1.524944in}}%
\pgfpathlineto{\pgfqpoint{2.055634in}{1.506632in}}%
\pgfpathlineto{\pgfqpoint{2.056883in}{1.487728in}}%
\pgfpathlineto{\pgfqpoint{2.054252in}{1.476692in}}%
\pgfpathlineto{\pgfqpoint{2.060697in}{1.465454in}}%
\pgfpathlineto{\pgfqpoint{2.069537in}{1.445017in}}%
\pgfpathlineto{\pgfqpoint{2.077998in}{1.436700in}}%
\pgfpathlineto{\pgfqpoint{2.088084in}{1.418676in}}%
\pgfpathlineto{\pgfqpoint{2.059498in}{1.418379in}}%
\pgfpathlineto{\pgfqpoint{1.997690in}{1.419335in}}%
\pgfpathlineto{\pgfqpoint{1.929402in}{1.421428in}}%
\pgfpathlineto{\pgfqpoint{1.860711in}{1.424063in}}%
\pgfpathlineto{\pgfqpoint{1.759716in}{1.429580in}}%
\pgfpathlineto{\pgfqpoint{1.680190in}{1.435164in}}%
\pgfusepath{stroke}%
\end{pgfscope}%
\begin{pgfscope}%
\pgfpathrectangle{\pgfqpoint{0.100000in}{0.100000in}}{\pgfqpoint{3.608454in}{2.310000in}}%
\pgfusepath{clip}%
\pgfsetbuttcap%
\pgfsetroundjoin%
\pgfsetlinewidth{0.010037pt}%
\definecolor{currentstroke}{rgb}{1.000000,1.000000,1.000000}%
\pgfsetstrokecolor{currentstroke}%
\pgfsetdash{}{0pt}%
\pgfpathmoveto{\pgfqpoint{2.995203in}{1.672680in}}%
\pgfpathlineto{\pgfqpoint{3.026828in}{1.702834in}}%
\pgfpathlineto{\pgfqpoint{3.030759in}{1.713553in}}%
\pgfpathlineto{\pgfqpoint{3.040025in}{1.722729in}}%
\pgfpathlineto{\pgfqpoint{3.033072in}{1.736089in}}%
\pgfpathlineto{\pgfqpoint{3.024354in}{1.743903in}}%
\pgfpathlineto{\pgfqpoint{3.021834in}{1.757751in}}%
\pgfpathlineto{\pgfqpoint{3.054283in}{1.771969in}}%
\pgfpathlineto{\pgfqpoint{3.081149in}{1.776467in}}%
\pgfpathlineto{\pgfqpoint{3.095585in}{1.776768in}}%
\pgfpathlineto{\pgfqpoint{3.106582in}{1.771334in}}%
\pgfpathlineto{\pgfqpoint{3.117302in}{1.776186in}}%
\pgfpathlineto{\pgfqpoint{3.143447in}{1.781623in}}%
\pgfpathlineto{\pgfqpoint{3.152481in}{1.788641in}}%
\pgfpathlineto{\pgfqpoint{3.165877in}{1.804177in}}%
\pgfpathlineto{\pgfqpoint{3.178063in}{1.811042in}}%
\pgfpathlineto{\pgfqpoint{3.178922in}{1.817626in}}%
\pgfpathlineto{\pgfqpoint{3.172527in}{1.832648in}}%
\pgfpathlineto{\pgfqpoint{3.177150in}{1.841488in}}%
\pgfpathlineto{\pgfqpoint{3.170925in}{1.850993in}}%
\pgfpathlineto{\pgfqpoint{3.161391in}{1.851657in}}%
\pgfpathlineto{\pgfqpoint{3.185252in}{1.880369in}}%
\pgfpathlineto{\pgfqpoint{3.188086in}{1.891311in}}%
\pgfpathlineto{\pgfqpoint{3.206865in}{1.919218in}}%
\pgfpathlineto{\pgfqpoint{3.224297in}{1.934319in}}%
\pgfpathlineto{\pgfqpoint{3.236268in}{1.940625in}}%
\pgfpathlineto{\pgfqpoint{3.275437in}{1.949500in}}%
\pgfpathlineto{\pgfqpoint{3.312149in}{1.959523in}}%
\pgfpathlineto{\pgfqpoint{3.317011in}{1.943790in}}%
\pgfpathlineto{\pgfqpoint{3.317356in}{1.929202in}}%
\pgfpathlineto{\pgfqpoint{3.325730in}{1.915960in}}%
\pgfpathlineto{\pgfqpoint{3.328311in}{1.902477in}}%
\pgfpathlineto{\pgfqpoint{3.325093in}{1.884184in}}%
\pgfpathlineto{\pgfqpoint{3.334894in}{1.859510in}}%
\pgfpathlineto{\pgfqpoint{3.346067in}{1.845790in}}%
\pgfpathlineto{\pgfqpoint{3.354710in}{1.801193in}}%
\pgfpathlineto{\pgfqpoint{3.359367in}{1.786492in}}%
\pgfpathlineto{\pgfqpoint{3.358165in}{1.731729in}}%
\pgfpathlineto{\pgfqpoint{3.359557in}{1.729620in}}%
\pgfpathlineto{\pgfqpoint{3.369731in}{1.670722in}}%
\pgfpathlineto{\pgfqpoint{3.375437in}{1.665355in}}%
\pgfpathlineto{\pgfqpoint{3.363164in}{1.653431in}}%
\pgfpathlineto{\pgfqpoint{3.369219in}{1.646587in}}%
\pgfpathlineto{\pgfqpoint{3.363897in}{1.636235in}}%
\pgfpathlineto{\pgfqpoint{3.363942in}{1.631828in}}%
\pgfpathlineto{\pgfqpoint{3.357288in}{1.627852in}}%
\pgfpathlineto{\pgfqpoint{3.354047in}{1.619061in}}%
\pgfpathlineto{\pgfqpoint{3.355075in}{1.643233in}}%
\pgfpathlineto{\pgfqpoint{3.334450in}{1.648555in}}%
\pgfpathlineto{\pgfqpoint{3.302196in}{1.659557in}}%
\pgfpathlineto{\pgfqpoint{3.297634in}{1.664976in}}%
\pgfpathlineto{\pgfqpoint{3.284155in}{1.666275in}}%
\pgfpathlineto{\pgfqpoint{3.276402in}{1.674341in}}%
\pgfpathlineto{\pgfqpoint{3.272219in}{1.690405in}}%
\pgfpathlineto{\pgfqpoint{3.261228in}{1.692404in}}%
\pgfpathlineto{\pgfqpoint{3.253878in}{1.700831in}}%
\pgfpathlineto{\pgfqpoint{3.160304in}{1.681966in}}%
\pgfpathlineto{\pgfqpoint{3.115582in}{1.672717in}}%
\pgfpathlineto{\pgfqpoint{3.047681in}{1.660324in}}%
\pgfpathlineto{\pgfqpoint{2.998785in}{1.652075in}}%
\pgfpathlineto{\pgfqpoint{2.995203in}{1.672680in}}%
\pgfusepath{stroke}%
\end{pgfscope}%
\begin{pgfscope}%
\pgfpathrectangle{\pgfqpoint{0.100000in}{0.100000in}}{\pgfqpoint{3.608454in}{2.310000in}}%
\pgfusepath{clip}%
\pgfsetbuttcap%
\pgfsetroundjoin%
\pgfsetlinewidth{0.010037pt}%
\definecolor{currentstroke}{rgb}{1.000000,1.000000,1.000000}%
\pgfsetstrokecolor{currentstroke}%
\pgfsetdash{}{0pt}%
\pgfpathmoveto{\pgfqpoint{3.370867in}{1.614152in}}%
\pgfpathlineto{\pgfqpoint{3.356395in}{1.609644in}}%
\pgfpathlineto{\pgfqpoint{3.353960in}{1.613788in}}%
\pgfpathlineto{\pgfqpoint{3.358606in}{1.627701in}}%
\pgfpathlineto{\pgfqpoint{3.367208in}{1.629071in}}%
\pgfpathlineto{\pgfqpoint{3.366438in}{1.633451in}}%
\pgfpathlineto{\pgfqpoint{3.374162in}{1.640026in}}%
\pgfpathlineto{\pgfqpoint{3.396494in}{1.645263in}}%
\pgfpathlineto{\pgfqpoint{3.406435in}{1.653192in}}%
\pgfpathlineto{\pgfqpoint{3.428768in}{1.659799in}}%
\pgfpathlineto{\pgfqpoint{3.438954in}{1.657382in}}%
\pgfpathlineto{\pgfqpoint{3.447523in}{1.668034in}}%
\pgfpathlineto{\pgfqpoint{3.460473in}{1.669442in}}%
\pgfpathlineto{\pgfqpoint{3.438379in}{1.648663in}}%
\pgfpathlineto{\pgfqpoint{3.370867in}{1.614152in}}%
\pgfusepath{stroke}%
\end{pgfscope}%
\begin{pgfscope}%
\pgfpathrectangle{\pgfqpoint{0.100000in}{0.100000in}}{\pgfqpoint{3.608454in}{2.310000in}}%
\pgfusepath{clip}%
\pgfsetbuttcap%
\pgfsetroundjoin%
\pgfsetlinewidth{0.010037pt}%
\definecolor{currentstroke}{rgb}{1.000000,1.000000,1.000000}%
\pgfsetstrokecolor{currentstroke}%
\pgfsetdash{}{0pt}%
\pgfpathmoveto{\pgfqpoint{3.045737in}{1.476948in}}%
\pgfpathlineto{\pgfqpoint{2.983156in}{1.466600in}}%
\pgfpathlineto{\pgfqpoint{2.971804in}{1.538119in}}%
\pgfpathlineto{\pgfqpoint{2.954949in}{1.643541in}}%
\pgfpathlineto{\pgfqpoint{2.995203in}{1.672680in}}%
\pgfpathlineto{\pgfqpoint{2.998785in}{1.652075in}}%
\pgfpathlineto{\pgfqpoint{3.047681in}{1.660324in}}%
\pgfpathlineto{\pgfqpoint{3.115582in}{1.672717in}}%
\pgfpathlineto{\pgfqpoint{3.160304in}{1.681966in}}%
\pgfpathlineto{\pgfqpoint{3.253878in}{1.700831in}}%
\pgfpathlineto{\pgfqpoint{3.261228in}{1.692404in}}%
\pgfpathlineto{\pgfqpoint{3.272219in}{1.690405in}}%
\pgfpathlineto{\pgfqpoint{3.276402in}{1.674341in}}%
\pgfpathlineto{\pgfqpoint{3.284155in}{1.666275in}}%
\pgfpathlineto{\pgfqpoint{3.297634in}{1.664976in}}%
\pgfpathlineto{\pgfqpoint{3.302196in}{1.659557in}}%
\pgfpathlineto{\pgfqpoint{3.297564in}{1.655376in}}%
\pgfpathlineto{\pgfqpoint{3.293392in}{1.640582in}}%
\pgfpathlineto{\pgfqpoint{3.283615in}{1.623960in}}%
\pgfpathlineto{\pgfqpoint{3.290178in}{1.616731in}}%
\pgfpathlineto{\pgfqpoint{3.284807in}{1.608986in}}%
\pgfpathlineto{\pgfqpoint{3.287837in}{1.592001in}}%
\pgfpathlineto{\pgfqpoint{3.297508in}{1.583085in}}%
\pgfpathlineto{\pgfqpoint{3.320457in}{1.568600in}}%
\pgfpathlineto{\pgfqpoint{3.301976in}{1.548162in}}%
\pgfpathlineto{\pgfqpoint{3.301717in}{1.540404in}}%
\pgfpathlineto{\pgfqpoint{3.286671in}{1.530400in}}%
\pgfpathlineto{\pgfqpoint{3.270035in}{1.528522in}}%
\pgfpathlineto{\pgfqpoint{3.265920in}{1.519827in}}%
\pgfpathlineto{\pgfqpoint{3.219645in}{1.509814in}}%
\pgfpathlineto{\pgfqpoint{3.165648in}{1.498939in}}%
\pgfpathlineto{\pgfqpoint{3.128532in}{1.492293in}}%
\pgfpathlineto{\pgfqpoint{3.045737in}{1.476948in}}%
\pgfusepath{stroke}%
\end{pgfscope}%
\begin{pgfscope}%
\pgfpathrectangle{\pgfqpoint{0.100000in}{0.100000in}}{\pgfqpoint{3.608454in}{2.310000in}}%
\pgfusepath{clip}%
\pgfsetbuttcap%
\pgfsetroundjoin%
\pgfsetlinewidth{0.010037pt}%
\definecolor{currentstroke}{rgb}{1.000000,1.000000,1.000000}%
\pgfsetstrokecolor{currentstroke}%
\pgfsetdash{}{0pt}%
\pgfpathmoveto{\pgfqpoint{3.359557in}{1.729620in}}%
\pgfpathlineto{\pgfqpoint{3.439208in}{1.747720in}}%
\pgfpathlineto{\pgfqpoint{3.456549in}{1.750456in}}%
\pgfpathlineto{\pgfqpoint{3.468026in}{1.705331in}}%
\pgfpathlineto{\pgfqpoint{3.466209in}{1.697251in}}%
\pgfpathlineto{\pgfqpoint{3.429341in}{1.682977in}}%
\pgfpathlineto{\pgfqpoint{3.407332in}{1.677989in}}%
\pgfpathlineto{\pgfqpoint{3.397980in}{1.666798in}}%
\pgfpathlineto{\pgfqpoint{3.369219in}{1.646587in}}%
\pgfpathlineto{\pgfqpoint{3.363164in}{1.653431in}}%
\pgfpathlineto{\pgfqpoint{3.375437in}{1.665355in}}%
\pgfpathlineto{\pgfqpoint{3.369731in}{1.670722in}}%
\pgfpathlineto{\pgfqpoint{3.359557in}{1.729620in}}%
\pgfusepath{stroke}%
\end{pgfscope}%
\begin{pgfscope}%
\pgfpathrectangle{\pgfqpoint{0.100000in}{0.100000in}}{\pgfqpoint{3.608454in}{2.310000in}}%
\pgfusepath{clip}%
\pgfsetbuttcap%
\pgfsetroundjoin%
\pgfsetlinewidth{0.010037pt}%
\definecolor{currentstroke}{rgb}{1.000000,1.000000,1.000000}%
\pgfsetstrokecolor{currentstroke}%
\pgfsetdash{}{0pt}%
\pgfpathmoveto{\pgfqpoint{3.456549in}{1.750456in}}%
\pgfpathlineto{\pgfqpoint{3.480196in}{1.757459in}}%
\pgfpathlineto{\pgfqpoint{3.487074in}{1.740513in}}%
\pgfpathlineto{\pgfqpoint{3.494962in}{1.736335in}}%
\pgfpathlineto{\pgfqpoint{3.485228in}{1.729197in}}%
\pgfpathlineto{\pgfqpoint{3.486456in}{1.708233in}}%
\pgfpathlineto{\pgfqpoint{3.466209in}{1.697251in}}%
\pgfpathlineto{\pgfqpoint{3.468026in}{1.705331in}}%
\pgfpathlineto{\pgfqpoint{3.456549in}{1.750456in}}%
\pgfusepath{stroke}%
\end{pgfscope}%
\begin{pgfscope}%
\pgfpathrectangle{\pgfqpoint{0.100000in}{0.100000in}}{\pgfqpoint{3.608454in}{2.310000in}}%
\pgfusepath{clip}%
\pgfsetbuttcap%
\pgfsetroundjoin%
\pgfsetlinewidth{0.010037pt}%
\definecolor{currentstroke}{rgb}{1.000000,1.000000,1.000000}%
\pgfsetstrokecolor{currentstroke}%
\pgfsetdash{}{0pt}%
\pgfpathmoveto{\pgfqpoint{3.284005in}{1.523011in}}%
\pgfpathlineto{\pgfqpoint{3.286671in}{1.530400in}}%
\pgfpathlineto{\pgfqpoint{3.301717in}{1.540404in}}%
\pgfpathlineto{\pgfqpoint{3.301976in}{1.548162in}}%
\pgfpathlineto{\pgfqpoint{3.320457in}{1.568600in}}%
\pgfpathlineto{\pgfqpoint{3.297508in}{1.583085in}}%
\pgfpathlineto{\pgfqpoint{3.287837in}{1.592001in}}%
\pgfpathlineto{\pgfqpoint{3.284807in}{1.608986in}}%
\pgfpathlineto{\pgfqpoint{3.290178in}{1.616731in}}%
\pgfpathlineto{\pgfqpoint{3.283615in}{1.623960in}}%
\pgfpathlineto{\pgfqpoint{3.293392in}{1.640582in}}%
\pgfpathlineto{\pgfqpoint{3.297564in}{1.655376in}}%
\pgfpathlineto{\pgfqpoint{3.302196in}{1.659557in}}%
\pgfpathlineto{\pgfqpoint{3.334450in}{1.648555in}}%
\pgfpathlineto{\pgfqpoint{3.355075in}{1.643233in}}%
\pgfpathlineto{\pgfqpoint{3.354047in}{1.619061in}}%
\pgfpathlineto{\pgfqpoint{3.341519in}{1.600731in}}%
\pgfpathlineto{\pgfqpoint{3.351842in}{1.598035in}}%
\pgfpathlineto{\pgfqpoint{3.362561in}{1.590168in}}%
\pgfpathlineto{\pgfqpoint{3.363194in}{1.568654in}}%
\pgfpathlineto{\pgfqpoint{3.359949in}{1.553423in}}%
\pgfpathlineto{\pgfqpoint{3.362108in}{1.540913in}}%
\pgfpathlineto{\pgfqpoint{3.343475in}{1.498697in}}%
\pgfpathlineto{\pgfqpoint{3.336779in}{1.478984in}}%
\pgfpathlineto{\pgfqpoint{3.327429in}{1.488566in}}%
\pgfpathlineto{\pgfqpoint{3.315030in}{1.486933in}}%
\pgfpathlineto{\pgfqpoint{3.284000in}{1.504862in}}%
\pgfpathlineto{\pgfqpoint{3.280825in}{1.514462in}}%
\pgfpathlineto{\pgfqpoint{3.284005in}{1.523011in}}%
\pgfusepath{stroke}%
\end{pgfscope}%
\begin{pgfscope}%
\pgfpathrectangle{\pgfqpoint{0.100000in}{0.100000in}}{\pgfqpoint{3.608454in}{2.310000in}}%
\pgfusepath{clip}%
\pgfsetbuttcap%
\pgfsetroundjoin%
\pgfsetlinewidth{0.010037pt}%
\definecolor{currentstroke}{rgb}{1.000000,1.000000,1.000000}%
\pgfsetstrokecolor{currentstroke}%
\pgfsetdash{}{0pt}%
\pgfpathmoveto{\pgfqpoint{2.541376in}{1.261906in}}%
\pgfpathlineto{\pgfqpoint{2.539926in}{1.281490in}}%
\pgfpathlineto{\pgfqpoint{2.545564in}{1.291207in}}%
\pgfpathlineto{\pgfqpoint{2.541848in}{1.295967in}}%
\pgfpathlineto{\pgfqpoint{2.555400in}{1.311555in}}%
\pgfpathlineto{\pgfqpoint{2.554624in}{1.315312in}}%
\pgfpathlineto{\pgfqpoint{2.567791in}{1.342507in}}%
\pgfpathlineto{\pgfqpoint{2.565129in}{1.355623in}}%
\pgfpathlineto{\pgfqpoint{2.555581in}{1.369354in}}%
\pgfpathlineto{\pgfqpoint{2.562209in}{1.386059in}}%
\pgfpathlineto{\pgfqpoint{2.553618in}{1.495995in}}%
\pgfpathlineto{\pgfqpoint{2.547432in}{1.573118in}}%
\pgfpathlineto{\pgfqpoint{2.555978in}{1.566728in}}%
\pgfpathlineto{\pgfqpoint{2.565512in}{1.566900in}}%
\pgfpathlineto{\pgfqpoint{2.588031in}{1.579945in}}%
\pgfpathlineto{\pgfqpoint{2.657102in}{1.586389in}}%
\pgfpathlineto{\pgfqpoint{2.708226in}{1.591778in}}%
\pgfpathlineto{\pgfqpoint{2.708677in}{1.586775in}}%
\pgfpathlineto{\pgfqpoint{2.720348in}{1.481089in}}%
\pgfpathlineto{\pgfqpoint{2.730393in}{1.382751in}}%
\pgfpathlineto{\pgfqpoint{2.726063in}{1.378135in}}%
\pgfpathlineto{\pgfqpoint{2.732714in}{1.366675in}}%
\pgfpathlineto{\pgfqpoint{2.732680in}{1.358417in}}%
\pgfpathlineto{\pgfqpoint{2.723189in}{1.356342in}}%
\pgfpathlineto{\pgfqpoint{2.712569in}{1.348381in}}%
\pgfpathlineto{\pgfqpoint{2.705378in}{1.351515in}}%
\pgfpathlineto{\pgfqpoint{2.694616in}{1.346419in}}%
\pgfpathlineto{\pgfqpoint{2.697946in}{1.336184in}}%
\pgfpathlineto{\pgfqpoint{2.686884in}{1.325901in}}%
\pgfpathlineto{\pgfqpoint{2.683829in}{1.314006in}}%
\pgfpathlineto{\pgfqpoint{2.676223in}{1.312054in}}%
\pgfpathlineto{\pgfqpoint{2.670547in}{1.303046in}}%
\pgfpathlineto{\pgfqpoint{2.671290in}{1.293974in}}%
\pgfpathlineto{\pgfqpoint{2.664612in}{1.287594in}}%
\pgfpathlineto{\pgfqpoint{2.654564in}{1.288579in}}%
\pgfpathlineto{\pgfqpoint{2.646926in}{1.298367in}}%
\pgfpathlineto{\pgfqpoint{2.633983in}{1.288941in}}%
\pgfpathlineto{\pgfqpoint{2.634892in}{1.280705in}}%
\pgfpathlineto{\pgfqpoint{2.620497in}{1.275892in}}%
\pgfpathlineto{\pgfqpoint{2.616881in}{1.281949in}}%
\pgfpathlineto{\pgfqpoint{2.602868in}{1.275081in}}%
\pgfpathlineto{\pgfqpoint{2.597748in}{1.265248in}}%
\pgfpathlineto{\pgfqpoint{2.580872in}{1.275331in}}%
\pgfpathlineto{\pgfqpoint{2.554124in}{1.268651in}}%
\pgfpathlineto{\pgfqpoint{2.541376in}{1.261906in}}%
\pgfusepath{stroke}%
\end{pgfscope}%
\begin{pgfscope}%
\pgfpathrectangle{\pgfqpoint{0.100000in}{0.100000in}}{\pgfqpoint{3.608454in}{2.310000in}}%
\pgfusepath{clip}%
\pgfsetbuttcap%
\pgfsetroundjoin%
\pgfsetlinewidth{0.010037pt}%
\definecolor{currentstroke}{rgb}{1.000000,1.000000,1.000000}%
\pgfsetstrokecolor{currentstroke}%
\pgfsetdash{}{0pt}%
\pgfpathmoveto{\pgfqpoint{0.655019in}{1.776916in}}%
\pgfpathlineto{\pgfqpoint{0.693698in}{1.766153in}}%
\pgfpathlineto{\pgfqpoint{0.757813in}{1.749621in}}%
\pgfpathlineto{\pgfqpoint{0.824275in}{1.732591in}}%
\pgfpathlineto{\pgfqpoint{0.885533in}{1.717818in}}%
\pgfpathlineto{\pgfqpoint{0.938464in}{1.705834in}}%
\pgfpathlineto{\pgfqpoint{0.995094in}{1.693381in}}%
\pgfpathlineto{\pgfqpoint{0.978732in}{1.616166in}}%
\pgfpathlineto{\pgfqpoint{0.964156in}{1.547595in}}%
\pgfpathlineto{\pgfqpoint{0.940244in}{1.436956in}}%
\pgfpathlineto{\pgfqpoint{0.922291in}{1.353447in}}%
\pgfpathlineto{\pgfqpoint{0.912593in}{1.306845in}}%
\pgfpathlineto{\pgfqpoint{0.900156in}{1.246360in}}%
\pgfpathlineto{\pgfqpoint{0.891552in}{1.234089in}}%
\pgfpathlineto{\pgfqpoint{0.884626in}{1.233679in}}%
\pgfpathlineto{\pgfqpoint{0.879687in}{1.244390in}}%
\pgfpathlineto{\pgfqpoint{0.868345in}{1.248280in}}%
\pgfpathlineto{\pgfqpoint{0.856124in}{1.246917in}}%
\pgfpathlineto{\pgfqpoint{0.852655in}{1.238156in}}%
\pgfpathlineto{\pgfqpoint{0.852293in}{1.207731in}}%
\pgfpathlineto{\pgfqpoint{0.848633in}{1.200785in}}%
\pgfpathlineto{\pgfqpoint{0.850725in}{1.176360in}}%
\pgfpathlineto{\pgfqpoint{0.843082in}{1.160088in}}%
\pgfpathlineto{\pgfqpoint{0.781106in}{1.255521in}}%
\pgfpathlineto{\pgfqpoint{0.720071in}{1.348255in}}%
\pgfpathlineto{\pgfqpoint{0.688300in}{1.396883in}}%
\pgfpathlineto{\pgfqpoint{0.661796in}{1.438856in}}%
\pgfpathlineto{\pgfqpoint{0.628400in}{1.490776in}}%
\pgfpathlineto{\pgfqpoint{0.590596in}{1.549016in}}%
\pgfpathlineto{\pgfqpoint{0.606139in}{1.604297in}}%
\pgfpathlineto{\pgfqpoint{0.637416in}{1.715162in}}%
\pgfpathlineto{\pgfqpoint{0.655019in}{1.776916in}}%
\pgfusepath{stroke}%
\end{pgfscope}%
\begin{pgfscope}%
\pgfpathrectangle{\pgfqpoint{0.100000in}{0.100000in}}{\pgfqpoint{3.608454in}{2.310000in}}%
\pgfusepath{clip}%
\pgfsetbuttcap%
\pgfsetroundjoin%
\pgfsetlinewidth{0.010037pt}%
\definecolor{currentstroke}{rgb}{1.000000,1.000000,1.000000}%
\pgfsetstrokecolor{currentstroke}%
\pgfsetdash{}{0pt}%
\pgfpathmoveto{\pgfqpoint{0.995094in}{1.693381in}}%
\pgfpathlineto{\pgfqpoint{1.055637in}{1.681339in}}%
\pgfpathlineto{\pgfqpoint{1.167721in}{1.659701in}}%
\pgfpathlineto{\pgfqpoint{1.153691in}{1.581874in}}%
\pgfpathlineto{\pgfqpoint{1.215397in}{1.571309in}}%
\pgfpathlineto{\pgfqpoint{1.271614in}{1.562313in}}%
\pgfpathlineto{\pgfqpoint{1.261846in}{1.500784in}}%
\pgfpathlineto{\pgfqpoint{1.251495in}{1.434433in}}%
\pgfpathlineto{\pgfqpoint{1.237637in}{1.347314in}}%
\pgfpathlineto{\pgfqpoint{1.237278in}{1.340011in}}%
\pgfpathlineto{\pgfqpoint{1.222893in}{1.249723in}}%
\pgfpathlineto{\pgfqpoint{1.163676in}{1.258901in}}%
\pgfpathlineto{\pgfqpoint{1.133501in}{1.264961in}}%
\pgfpathlineto{\pgfqpoint{1.024434in}{1.284134in}}%
\pgfpathlineto{\pgfqpoint{0.983361in}{1.292233in}}%
\pgfpathlineto{\pgfqpoint{0.912593in}{1.306845in}}%
\pgfpathlineto{\pgfqpoint{0.922291in}{1.353447in}}%
\pgfpathlineto{\pgfqpoint{0.940244in}{1.436956in}}%
\pgfpathlineto{\pgfqpoint{0.964156in}{1.547595in}}%
\pgfpathlineto{\pgfqpoint{0.978732in}{1.616166in}}%
\pgfpathlineto{\pgfqpoint{0.995094in}{1.693381in}}%
\pgfusepath{stroke}%
\end{pgfscope}%
\begin{pgfscope}%
\pgfpathrectangle{\pgfqpoint{0.100000in}{0.100000in}}{\pgfqpoint{3.608454in}{2.310000in}}%
\pgfusepath{clip}%
\pgfsetbuttcap%
\pgfsetroundjoin%
\pgfsetlinewidth{0.010037pt}%
\definecolor{currentstroke}{rgb}{1.000000,1.000000,1.000000}%
\pgfsetstrokecolor{currentstroke}%
\pgfsetdash{}{0pt}%
\pgfpathmoveto{\pgfqpoint{0.418363in}{1.850356in}}%
\pgfpathlineto{\pgfqpoint{0.439875in}{1.842853in}}%
\pgfpathlineto{\pgfqpoint{0.526049in}{1.816024in}}%
\pgfpathlineto{\pgfqpoint{0.605321in}{1.791073in}}%
\pgfpathlineto{\pgfqpoint{0.655019in}{1.776916in}}%
\pgfpathlineto{\pgfqpoint{0.637416in}{1.715162in}}%
\pgfpathlineto{\pgfqpoint{0.606139in}{1.604297in}}%
\pgfpathlineto{\pgfqpoint{0.590596in}{1.549016in}}%
\pgfpathlineto{\pgfqpoint{0.628400in}{1.490776in}}%
\pgfpathlineto{\pgfqpoint{0.661796in}{1.438856in}}%
\pgfpathlineto{\pgfqpoint{0.688300in}{1.396883in}}%
\pgfpathlineto{\pgfqpoint{0.720071in}{1.348255in}}%
\pgfpathlineto{\pgfqpoint{0.781106in}{1.255521in}}%
\pgfpathlineto{\pgfqpoint{0.843082in}{1.160088in}}%
\pgfpathlineto{\pgfqpoint{0.840594in}{1.150624in}}%
\pgfpathlineto{\pgfqpoint{0.848070in}{1.135552in}}%
\pgfpathlineto{\pgfqpoint{0.849548in}{1.114934in}}%
\pgfpathlineto{\pgfqpoint{0.860405in}{1.104932in}}%
\pgfpathlineto{\pgfqpoint{0.860834in}{1.096874in}}%
\pgfpathlineto{\pgfqpoint{0.841395in}{1.087727in}}%
\pgfpathlineto{\pgfqpoint{0.832133in}{1.078570in}}%
\pgfpathlineto{\pgfqpoint{0.829245in}{1.058349in}}%
\pgfpathlineto{\pgfqpoint{0.824557in}{1.047321in}}%
\pgfpathlineto{\pgfqpoint{0.814680in}{1.038025in}}%
\pgfpathlineto{\pgfqpoint{0.804856in}{1.013854in}}%
\pgfpathlineto{\pgfqpoint{0.818669in}{1.001263in}}%
\pgfpathlineto{\pgfqpoint{0.816886in}{0.990883in}}%
\pgfpathlineto{\pgfqpoint{0.805560in}{0.983661in}}%
\pgfpathlineto{\pgfqpoint{0.797751in}{0.984935in}}%
\pgfpathlineto{\pgfqpoint{0.705575in}{0.997727in}}%
\pgfpathlineto{\pgfqpoint{0.637495in}{1.007385in}}%
\pgfpathlineto{\pgfqpoint{0.640472in}{1.018376in}}%
\pgfpathlineto{\pgfqpoint{0.636040in}{1.036625in}}%
\pgfpathlineto{\pgfqpoint{0.635593in}{1.055032in}}%
\pgfpathlineto{\pgfqpoint{0.632697in}{1.065804in}}%
\pgfpathlineto{\pgfqpoint{0.623775in}{1.081205in}}%
\pgfpathlineto{\pgfqpoint{0.598111in}{1.116738in}}%
\pgfpathlineto{\pgfqpoint{0.589705in}{1.121092in}}%
\pgfpathlineto{\pgfqpoint{0.578875in}{1.121023in}}%
\pgfpathlineto{\pgfqpoint{0.581349in}{1.132233in}}%
\pgfpathlineto{\pgfqpoint{0.576221in}{1.146250in}}%
\pgfpathlineto{\pgfqpoint{0.551032in}{1.153217in}}%
\pgfpathlineto{\pgfqpoint{0.535677in}{1.166118in}}%
\pgfpathlineto{\pgfqpoint{0.534404in}{1.174011in}}%
\pgfpathlineto{\pgfqpoint{0.516724in}{1.193554in}}%
\pgfpathlineto{\pgfqpoint{0.499894in}{1.197294in}}%
\pgfpathlineto{\pgfqpoint{0.484301in}{1.207225in}}%
\pgfpathlineto{\pgfqpoint{0.463793in}{1.210661in}}%
\pgfpathlineto{\pgfqpoint{0.455031in}{1.223907in}}%
\pgfpathlineto{\pgfqpoint{0.463352in}{1.244873in}}%
\pgfpathlineto{\pgfqpoint{0.460807in}{1.249588in}}%
\pgfpathlineto{\pgfqpoint{0.467737in}{1.267066in}}%
\pgfpathlineto{\pgfqpoint{0.455409in}{1.276367in}}%
\pgfpathlineto{\pgfqpoint{0.459410in}{1.293230in}}%
\pgfpathlineto{\pgfqpoint{0.452825in}{1.297541in}}%
\pgfpathlineto{\pgfqpoint{0.447140in}{1.313479in}}%
\pgfpathlineto{\pgfqpoint{0.440280in}{1.318340in}}%
\pgfpathlineto{\pgfqpoint{0.439763in}{1.329853in}}%
\pgfpathlineto{\pgfqpoint{0.434389in}{1.337971in}}%
\pgfpathlineto{\pgfqpoint{0.426289in}{1.365312in}}%
\pgfpathlineto{\pgfqpoint{0.417427in}{1.378407in}}%
\pgfpathlineto{\pgfqpoint{0.419354in}{1.400637in}}%
\pgfpathlineto{\pgfqpoint{0.429775in}{1.402874in}}%
\pgfpathlineto{\pgfqpoint{0.436561in}{1.414932in}}%
\pgfpathlineto{\pgfqpoint{0.432461in}{1.428009in}}%
\pgfpathlineto{\pgfqpoint{0.421377in}{1.430192in}}%
\pgfpathlineto{\pgfqpoint{0.415847in}{1.436309in}}%
\pgfpathlineto{\pgfqpoint{0.406892in}{1.458808in}}%
\pgfpathlineto{\pgfqpoint{0.411078in}{1.466888in}}%
\pgfpathlineto{\pgfqpoint{0.408109in}{1.481933in}}%
\pgfpathlineto{\pgfqpoint{0.414682in}{1.501419in}}%
\pgfpathlineto{\pgfqpoint{0.422383in}{1.494243in}}%
\pgfpathlineto{\pgfqpoint{0.418888in}{1.485771in}}%
\pgfpathlineto{\pgfqpoint{0.432233in}{1.472341in}}%
\pgfpathlineto{\pgfqpoint{0.431381in}{1.492284in}}%
\pgfpathlineto{\pgfqpoint{0.425663in}{1.497632in}}%
\pgfpathlineto{\pgfqpoint{0.432199in}{1.515167in}}%
\pgfpathlineto{\pgfqpoint{0.456921in}{1.512375in}}%
\pgfpathlineto{\pgfqpoint{0.453594in}{1.518895in}}%
\pgfpathlineto{\pgfqpoint{0.437278in}{1.518254in}}%
\pgfpathlineto{\pgfqpoint{0.429516in}{1.528130in}}%
\pgfpathlineto{\pgfqpoint{0.419745in}{1.519360in}}%
\pgfpathlineto{\pgfqpoint{0.414561in}{1.504701in}}%
\pgfpathlineto{\pgfqpoint{0.400724in}{1.524426in}}%
\pgfpathlineto{\pgfqpoint{0.395395in}{1.528017in}}%
\pgfpathlineto{\pgfqpoint{0.397425in}{1.549485in}}%
\pgfpathlineto{\pgfqpoint{0.393207in}{1.562145in}}%
\pgfpathlineto{\pgfqpoint{0.385562in}{1.574039in}}%
\pgfpathlineto{\pgfqpoint{0.369905in}{1.610482in}}%
\pgfpathlineto{\pgfqpoint{0.375024in}{1.618541in}}%
\pgfpathlineto{\pgfqpoint{0.374964in}{1.644010in}}%
\pgfpathlineto{\pgfqpoint{0.383414in}{1.658215in}}%
\pgfpathlineto{\pgfqpoint{0.385352in}{1.680415in}}%
\pgfpathlineto{\pgfqpoint{0.377384in}{1.705851in}}%
\pgfpathlineto{\pgfqpoint{0.366813in}{1.722042in}}%
\pgfpathlineto{\pgfqpoint{0.368698in}{1.736652in}}%
\pgfpathlineto{\pgfqpoint{0.398383in}{1.772011in}}%
\pgfpathlineto{\pgfqpoint{0.399858in}{1.784073in}}%
\pgfpathlineto{\pgfqpoint{0.413234in}{1.807081in}}%
\pgfpathlineto{\pgfqpoint{0.415090in}{1.828890in}}%
\pgfpathlineto{\pgfqpoint{0.410789in}{1.834455in}}%
\pgfpathlineto{\pgfqpoint{0.418363in}{1.850356in}}%
\pgfusepath{stroke}%
\end{pgfscope}%
\begin{pgfscope}%
\pgfpathrectangle{\pgfqpoint{0.100000in}{0.100000in}}{\pgfqpoint{3.608454in}{2.310000in}}%
\pgfusepath{clip}%
\pgfsetbuttcap%
\pgfsetroundjoin%
\pgfsetlinewidth{0.010037pt}%
\definecolor{currentstroke}{rgb}{1.000000,1.000000,1.000000}%
\pgfsetstrokecolor{currentstroke}%
\pgfsetdash{}{0pt}%
\pgfpathmoveto{\pgfqpoint{2.730393in}{1.382751in}}%
\pgfpathlineto{\pgfqpoint{2.720348in}{1.481089in}}%
\pgfpathlineto{\pgfqpoint{2.708677in}{1.586775in}}%
\pgfpathlineto{\pgfqpoint{2.785105in}{1.598157in}}%
\pgfpathlineto{\pgfqpoint{2.805389in}{1.592878in}}%
\pgfpathlineto{\pgfqpoint{2.815117in}{1.587147in}}%
\pgfpathlineto{\pgfqpoint{2.827308in}{1.588736in}}%
\pgfpathlineto{\pgfqpoint{2.843381in}{1.579248in}}%
\pgfpathlineto{\pgfqpoint{2.873320in}{1.593370in}}%
\pgfpathlineto{\pgfqpoint{2.889841in}{1.593841in}}%
\pgfpathlineto{\pgfqpoint{2.909135in}{1.615366in}}%
\pgfpathlineto{\pgfqpoint{2.928765in}{1.628459in}}%
\pgfpathlineto{\pgfqpoint{2.954949in}{1.643541in}}%
\pgfpathlineto{\pgfqpoint{2.971804in}{1.538119in}}%
\pgfpathlineto{\pgfqpoint{2.963973in}{1.531348in}}%
\pgfpathlineto{\pgfqpoint{2.969029in}{1.525129in}}%
\pgfpathlineto{\pgfqpoint{2.971043in}{1.511497in}}%
\pgfpathlineto{\pgfqpoint{2.967120in}{1.498692in}}%
\pgfpathlineto{\pgfqpoint{2.966601in}{1.479976in}}%
\pgfpathlineto{\pgfqpoint{2.962917in}{1.455616in}}%
\pgfpathlineto{\pgfqpoint{2.944276in}{1.433890in}}%
\pgfpathlineto{\pgfqpoint{2.936429in}{1.429278in}}%
\pgfpathlineto{\pgfqpoint{2.930326in}{1.433249in}}%
\pgfpathlineto{\pgfqpoint{2.915224in}{1.412541in}}%
\pgfpathlineto{\pgfqpoint{2.916638in}{1.392602in}}%
\pgfpathlineto{\pgfqpoint{2.899819in}{1.397397in}}%
\pgfpathlineto{\pgfqpoint{2.891855in}{1.377484in}}%
\pgfpathlineto{\pgfqpoint{2.895881in}{1.363373in}}%
\pgfpathlineto{\pgfqpoint{2.889579in}{1.361288in}}%
\pgfpathlineto{\pgfqpoint{2.888697in}{1.350139in}}%
\pgfpathlineto{\pgfqpoint{2.873228in}{1.345641in}}%
\pgfpathlineto{\pgfqpoint{2.865190in}{1.354643in}}%
\pgfpathlineto{\pgfqpoint{2.854847in}{1.358134in}}%
\pgfpathlineto{\pgfqpoint{2.852417in}{1.367262in}}%
\pgfpathlineto{\pgfqpoint{2.843065in}{1.365657in}}%
\pgfpathlineto{\pgfqpoint{2.836938in}{1.357261in}}%
\pgfpathlineto{\pgfqpoint{2.828173in}{1.354308in}}%
\pgfpathlineto{\pgfqpoint{2.812691in}{1.360250in}}%
\pgfpathlineto{\pgfqpoint{2.804128in}{1.353181in}}%
\pgfpathlineto{\pgfqpoint{2.791958in}{1.361544in}}%
\pgfpathlineto{\pgfqpoint{2.768610in}{1.364112in}}%
\pgfpathlineto{\pgfqpoint{2.761568in}{1.379305in}}%
\pgfpathlineto{\pgfqpoint{2.752646in}{1.386037in}}%
\pgfpathlineto{\pgfqpoint{2.743996in}{1.381728in}}%
\pgfpathlineto{\pgfqpoint{2.730393in}{1.382751in}}%
\pgfusepath{stroke}%
\end{pgfscope}%
\begin{pgfscope}%
\pgfpathrectangle{\pgfqpoint{0.100000in}{0.100000in}}{\pgfqpoint{3.608454in}{2.310000in}}%
\pgfusepath{clip}%
\pgfsetbuttcap%
\pgfsetroundjoin%
\pgfsetlinewidth{0.010037pt}%
\definecolor{currentstroke}{rgb}{1.000000,1.000000,1.000000}%
\pgfsetstrokecolor{currentstroke}%
\pgfsetdash{}{0pt}%
\pgfpathmoveto{\pgfqpoint{2.477364in}{1.192570in}}%
\pgfpathlineto{\pgfqpoint{2.468410in}{1.199844in}}%
\pgfpathlineto{\pgfqpoint{2.461108in}{1.196383in}}%
\pgfpathlineto{\pgfqpoint{2.451774in}{1.213813in}}%
\pgfpathlineto{\pgfqpoint{2.456543in}{1.224768in}}%
\pgfpathlineto{\pgfqpoint{2.449670in}{1.237099in}}%
\pgfpathlineto{\pgfqpoint{2.449301in}{1.246806in}}%
\pgfpathlineto{\pgfqpoint{2.440014in}{1.250265in}}%
\pgfpathlineto{\pgfqpoint{2.435708in}{1.257580in}}%
\pgfpathlineto{\pgfqpoint{2.417550in}{1.266715in}}%
\pgfpathlineto{\pgfqpoint{2.401780in}{1.277970in}}%
\pgfpathlineto{\pgfqpoint{2.394453in}{1.286467in}}%
\pgfpathlineto{\pgfqpoint{2.394301in}{1.296831in}}%
\pgfpathlineto{\pgfqpoint{2.404106in}{1.316728in}}%
\pgfpathlineto{\pgfqpoint{2.407122in}{1.331947in}}%
\pgfpathlineto{\pgfqpoint{2.399134in}{1.340547in}}%
\pgfpathlineto{\pgfqpoint{2.388555in}{1.343787in}}%
\pgfpathlineto{\pgfqpoint{2.375719in}{1.336702in}}%
\pgfpathlineto{\pgfqpoint{2.370287in}{1.348866in}}%
\pgfpathlineto{\pgfqpoint{2.367545in}{1.365354in}}%
\pgfpathlineto{\pgfqpoint{2.348630in}{1.380053in}}%
\pgfpathlineto{\pgfqpoint{2.344850in}{1.386576in}}%
\pgfpathlineto{\pgfqpoint{2.327573in}{1.401342in}}%
\pgfpathlineto{\pgfqpoint{2.322158in}{1.412077in}}%
\pgfpathlineto{\pgfqpoint{2.317277in}{1.433365in}}%
\pgfpathlineto{\pgfqpoint{2.320622in}{1.452274in}}%
\pgfpathlineto{\pgfqpoint{2.325089in}{1.454910in}}%
\pgfpathlineto{\pgfqpoint{2.324285in}{1.470738in}}%
\pgfpathlineto{\pgfqpoint{2.336891in}{1.475440in}}%
\pgfpathlineto{\pgfqpoint{2.340698in}{1.489644in}}%
\pgfpathlineto{\pgfqpoint{2.347951in}{1.499205in}}%
\pgfpathlineto{\pgfqpoint{2.347584in}{1.511368in}}%
\pgfpathlineto{\pgfqpoint{2.338595in}{1.521042in}}%
\pgfpathlineto{\pgfqpoint{2.340730in}{1.534582in}}%
\pgfpathlineto{\pgfqpoint{2.364032in}{1.540471in}}%
\pgfpathlineto{\pgfqpoint{2.381909in}{1.551220in}}%
\pgfpathlineto{\pgfqpoint{2.383785in}{1.564755in}}%
\pgfpathlineto{\pgfqpoint{2.390006in}{1.569016in}}%
\pgfpathlineto{\pgfqpoint{2.392386in}{1.583234in}}%
\pgfpathlineto{\pgfqpoint{2.390473in}{1.592626in}}%
\pgfpathlineto{\pgfqpoint{2.378251in}{1.600427in}}%
\pgfpathlineto{\pgfqpoint{2.373331in}{1.612057in}}%
\pgfpathlineto{\pgfqpoint{2.361269in}{1.623298in}}%
\pgfpathlineto{\pgfqpoint{2.460390in}{1.627522in}}%
\pgfpathlineto{\pgfqpoint{2.526890in}{1.632246in}}%
\pgfpathlineto{\pgfqpoint{2.525674in}{1.618258in}}%
\pgfpathlineto{\pgfqpoint{2.537007in}{1.598963in}}%
\pgfpathlineto{\pgfqpoint{2.541765in}{1.582478in}}%
\pgfpathlineto{\pgfqpoint{2.547432in}{1.573118in}}%
\pgfpathlineto{\pgfqpoint{2.553618in}{1.495995in}}%
\pgfpathlineto{\pgfqpoint{2.562209in}{1.386059in}}%
\pgfpathlineto{\pgfqpoint{2.555581in}{1.369354in}}%
\pgfpathlineto{\pgfqpoint{2.565129in}{1.355623in}}%
\pgfpathlineto{\pgfqpoint{2.567791in}{1.342507in}}%
\pgfpathlineto{\pgfqpoint{2.554624in}{1.315312in}}%
\pgfpathlineto{\pgfqpoint{2.555400in}{1.311555in}}%
\pgfpathlineto{\pgfqpoint{2.541848in}{1.295967in}}%
\pgfpathlineto{\pgfqpoint{2.545564in}{1.291207in}}%
\pgfpathlineto{\pgfqpoint{2.539926in}{1.281490in}}%
\pgfpathlineto{\pgfqpoint{2.541376in}{1.261906in}}%
\pgfpathlineto{\pgfqpoint{2.534528in}{1.249888in}}%
\pgfpathlineto{\pgfqpoint{2.540098in}{1.235688in}}%
\pgfpathlineto{\pgfqpoint{2.516761in}{1.227971in}}%
\pgfpathlineto{\pgfqpoint{2.514610in}{1.219580in}}%
\pgfpathlineto{\pgfqpoint{2.520983in}{1.208943in}}%
\pgfpathlineto{\pgfqpoint{2.518051in}{1.202010in}}%
\pgfpathlineto{\pgfqpoint{2.493026in}{1.210574in}}%
\pgfpathlineto{\pgfqpoint{2.484770in}{1.211436in}}%
\pgfpathlineto{\pgfqpoint{2.474476in}{1.198417in}}%
\pgfpathlineto{\pgfqpoint{2.477364in}{1.192570in}}%
\pgfusepath{stroke}%
\end{pgfscope}%
\begin{pgfscope}%
\pgfpathrectangle{\pgfqpoint{0.100000in}{0.100000in}}{\pgfqpoint{3.608454in}{2.310000in}}%
\pgfusepath{clip}%
\pgfsetbuttcap%
\pgfsetroundjoin%
\pgfsetlinewidth{0.010037pt}%
\definecolor{currentstroke}{rgb}{1.000000,1.000000,1.000000}%
\pgfsetstrokecolor{currentstroke}%
\pgfsetdash{}{0pt}%
\pgfpathmoveto{\pgfqpoint{3.205784in}{1.431832in}}%
\pgfpathlineto{\pgfqpoint{3.198897in}{1.442058in}}%
\pgfpathlineto{\pgfqpoint{3.202783in}{1.447781in}}%
\pgfpathlineto{\pgfqpoint{3.212300in}{1.441354in}}%
\pgfpathlineto{\pgfqpoint{3.205784in}{1.431832in}}%
\pgfusepath{stroke}%
\end{pgfscope}%
\begin{pgfscope}%
\pgfpathrectangle{\pgfqpoint{0.100000in}{0.100000in}}{\pgfqpoint{3.608454in}{2.310000in}}%
\pgfusepath{clip}%
\pgfsetbuttcap%
\pgfsetroundjoin%
\pgfsetlinewidth{0.010037pt}%
\definecolor{currentstroke}{rgb}{1.000000,1.000000,1.000000}%
\pgfsetstrokecolor{currentstroke}%
\pgfsetdash{}{0pt}%
\pgfpathmoveto{\pgfqpoint{3.265920in}{1.519827in}}%
\pgfpathlineto{\pgfqpoint{3.270035in}{1.528522in}}%
\pgfpathlineto{\pgfqpoint{3.286671in}{1.530400in}}%
\pgfpathlineto{\pgfqpoint{3.284005in}{1.523011in}}%
\pgfpathlineto{\pgfqpoint{3.278517in}{1.513571in}}%
\pgfpathlineto{\pgfqpoint{3.282236in}{1.502324in}}%
\pgfpathlineto{\pgfqpoint{3.296914in}{1.488849in}}%
\pgfpathlineto{\pgfqpoint{3.300321in}{1.474658in}}%
\pgfpathlineto{\pgfqpoint{3.317232in}{1.456987in}}%
\pgfpathlineto{\pgfqpoint{3.323867in}{1.457744in}}%
\pgfpathlineto{\pgfqpoint{3.332131in}{1.431216in}}%
\pgfpathlineto{\pgfqpoint{3.330776in}{1.430952in}}%
\pgfpathlineto{\pgfqpoint{3.329269in}{1.430656in}}%
\pgfpathlineto{\pgfqpoint{3.292410in}{1.423600in}}%
\pgfpathlineto{\pgfqpoint{3.272690in}{1.493733in}}%
\pgfpathlineto{\pgfqpoint{3.265920in}{1.519827in}}%
\pgfusepath{stroke}%
\end{pgfscope}%
\begin{pgfscope}%
\pgfpathrectangle{\pgfqpoint{0.100000in}{0.100000in}}{\pgfqpoint{3.608454in}{2.310000in}}%
\pgfusepath{clip}%
\pgfsetbuttcap%
\pgfsetroundjoin%
\pgfsetlinewidth{0.010037pt}%
\definecolor{currentstroke}{rgb}{1.000000,1.000000,1.000000}%
\pgfsetstrokecolor{currentstroke}%
\pgfsetdash{}{0pt}%
\pgfpathmoveto{\pgfqpoint{2.921647in}{1.282104in}}%
\pgfpathlineto{\pgfqpoint{2.910000in}{1.282532in}}%
\pgfpathlineto{\pgfqpoint{2.899268in}{1.289925in}}%
\pgfpathlineto{\pgfqpoint{2.884769in}{1.312394in}}%
\pgfpathlineto{\pgfqpoint{2.872427in}{1.324293in}}%
\pgfpathlineto{\pgfqpoint{2.875653in}{1.333484in}}%
\pgfpathlineto{\pgfqpoint{2.873228in}{1.345641in}}%
\pgfpathlineto{\pgfqpoint{2.888697in}{1.350139in}}%
\pgfpathlineto{\pgfqpoint{2.889579in}{1.361288in}}%
\pgfpathlineto{\pgfqpoint{2.895881in}{1.363373in}}%
\pgfpathlineto{\pgfqpoint{2.891855in}{1.377484in}}%
\pgfpathlineto{\pgfqpoint{2.899819in}{1.397397in}}%
\pgfpathlineto{\pgfqpoint{2.916638in}{1.392602in}}%
\pgfpathlineto{\pgfqpoint{2.915224in}{1.412541in}}%
\pgfpathlineto{\pgfqpoint{2.930326in}{1.433249in}}%
\pgfpathlineto{\pgfqpoint{2.936429in}{1.429278in}}%
\pgfpathlineto{\pgfqpoint{2.944276in}{1.433890in}}%
\pgfpathlineto{\pgfqpoint{2.962917in}{1.455616in}}%
\pgfpathlineto{\pgfqpoint{2.966601in}{1.479976in}}%
\pgfpathlineto{\pgfqpoint{2.967120in}{1.498692in}}%
\pgfpathlineto{\pgfqpoint{2.971043in}{1.511497in}}%
\pgfpathlineto{\pgfqpoint{2.969029in}{1.525129in}}%
\pgfpathlineto{\pgfqpoint{2.963973in}{1.531348in}}%
\pgfpathlineto{\pgfqpoint{2.971804in}{1.538119in}}%
\pgfpathlineto{\pgfqpoint{2.983156in}{1.466600in}}%
\pgfpathlineto{\pgfqpoint{3.045737in}{1.476948in}}%
\pgfpathlineto{\pgfqpoint{3.052224in}{1.436127in}}%
\pgfpathlineto{\pgfqpoint{3.074880in}{1.463071in}}%
\pgfpathlineto{\pgfqpoint{3.080218in}{1.460384in}}%
\pgfpathlineto{\pgfqpoint{3.088071in}{1.475996in}}%
\pgfpathlineto{\pgfqpoint{3.098951in}{1.471101in}}%
\pgfpathlineto{\pgfqpoint{3.108437in}{1.472032in}}%
\pgfpathlineto{\pgfqpoint{3.114696in}{1.482885in}}%
\pgfpathlineto{\pgfqpoint{3.123780in}{1.488922in}}%
\pgfpathlineto{\pgfqpoint{3.136410in}{1.483605in}}%
\pgfpathlineto{\pgfqpoint{3.144720in}{1.485442in}}%
\pgfpathlineto{\pgfqpoint{3.156626in}{1.464838in}}%
\pgfpathlineto{\pgfqpoint{3.153183in}{1.449239in}}%
\pgfpathlineto{\pgfqpoint{3.122088in}{1.466802in}}%
\pgfpathlineto{\pgfqpoint{3.116263in}{1.452384in}}%
\pgfpathlineto{\pgfqpoint{3.118185in}{1.445751in}}%
\pgfpathlineto{\pgfqpoint{3.104950in}{1.423326in}}%
\pgfpathlineto{\pgfqpoint{3.098705in}{1.418866in}}%
\pgfpathlineto{\pgfqpoint{3.095884in}{1.408977in}}%
\pgfpathlineto{\pgfqpoint{3.085235in}{1.410012in}}%
\pgfpathlineto{\pgfqpoint{3.077588in}{1.383022in}}%
\pgfpathlineto{\pgfqpoint{3.073308in}{1.376829in}}%
\pgfpathlineto{\pgfqpoint{3.062309in}{1.378894in}}%
\pgfpathlineto{\pgfqpoint{3.051058in}{1.387407in}}%
\pgfpathlineto{\pgfqpoint{3.050668in}{1.374341in}}%
\pgfpathlineto{\pgfqpoint{3.039777in}{1.352383in}}%
\pgfpathlineto{\pgfqpoint{3.037068in}{1.336751in}}%
\pgfpathlineto{\pgfqpoint{3.028711in}{1.326341in}}%
\pgfpathlineto{\pgfqpoint{3.022375in}{1.309715in}}%
\pgfpathlineto{\pgfqpoint{3.022064in}{1.294324in}}%
\pgfpathlineto{\pgfqpoint{3.012280in}{1.291449in}}%
\pgfpathlineto{\pgfqpoint{3.001180in}{1.282737in}}%
\pgfpathlineto{\pgfqpoint{2.990493in}{1.281077in}}%
\pgfpathlineto{\pgfqpoint{2.988030in}{1.273703in}}%
\pgfpathlineto{\pgfqpoint{2.970810in}{1.266151in}}%
\pgfpathlineto{\pgfqpoint{2.961186in}{1.272585in}}%
\pgfpathlineto{\pgfqpoint{2.950430in}{1.260367in}}%
\pgfpathlineto{\pgfqpoint{2.931940in}{1.263968in}}%
\pgfpathlineto{\pgfqpoint{2.920600in}{1.276793in}}%
\pgfpathlineto{\pgfqpoint{2.921647in}{1.282104in}}%
\pgfusepath{stroke}%
\end{pgfscope}%
\begin{pgfscope}%
\pgfpathrectangle{\pgfqpoint{0.100000in}{0.100000in}}{\pgfqpoint{3.608454in}{2.310000in}}%
\pgfusepath{clip}%
\pgfsetbuttcap%
\pgfsetroundjoin%
\pgfsetlinewidth{0.010037pt}%
\definecolor{currentstroke}{rgb}{1.000000,1.000000,1.000000}%
\pgfsetstrokecolor{currentstroke}%
\pgfsetdash{}{0pt}%
\pgfpathmoveto{\pgfqpoint{3.329269in}{1.430656in}}%
\pgfpathlineto{\pgfqpoint{3.328797in}{1.416244in}}%
\pgfpathlineto{\pgfqpoint{3.323250in}{1.409159in}}%
\pgfpathlineto{\pgfqpoint{3.319710in}{1.393431in}}%
\pgfpathlineto{\pgfqpoint{3.303727in}{1.386181in}}%
\pgfpathlineto{\pgfqpoint{3.290316in}{1.384068in}}%
\pgfpathlineto{\pgfqpoint{3.294224in}{1.394412in}}%
\pgfpathlineto{\pgfqpoint{3.273575in}{1.403080in}}%
\pgfpathlineto{\pgfqpoint{3.256800in}{1.413932in}}%
\pgfpathlineto{\pgfqpoint{3.267185in}{1.430155in}}%
\pgfpathlineto{\pgfqpoint{3.258862in}{1.442765in}}%
\pgfpathlineto{\pgfqpoint{3.259842in}{1.453641in}}%
\pgfpathlineto{\pgfqpoint{3.253824in}{1.456627in}}%
\pgfpathlineto{\pgfqpoint{3.248940in}{1.474298in}}%
\pgfpathlineto{\pgfqpoint{3.260489in}{1.498252in}}%
\pgfpathlineto{\pgfqpoint{3.251813in}{1.502164in}}%
\pgfpathlineto{\pgfqpoint{3.249564in}{1.490346in}}%
\pgfpathlineto{\pgfqpoint{3.237178in}{1.487056in}}%
\pgfpathlineto{\pgfqpoint{3.238481in}{1.448173in}}%
\pgfpathlineto{\pgfqpoint{3.236216in}{1.435660in}}%
\pgfpathlineto{\pgfqpoint{3.242428in}{1.417824in}}%
\pgfpathlineto{\pgfqpoint{3.251985in}{1.409243in}}%
\pgfpathlineto{\pgfqpoint{3.247651in}{1.403855in}}%
\pgfpathlineto{\pgfqpoint{3.257407in}{1.395982in}}%
\pgfpathlineto{\pgfqpoint{3.260933in}{1.383203in}}%
\pgfpathlineto{\pgfqpoint{3.243052in}{1.393785in}}%
\pgfpathlineto{\pgfqpoint{3.231736in}{1.392411in}}%
\pgfpathlineto{\pgfqpoint{3.222950in}{1.403297in}}%
\pgfpathlineto{\pgfqpoint{3.214003in}{1.404357in}}%
\pgfpathlineto{\pgfqpoint{3.201288in}{1.398904in}}%
\pgfpathlineto{\pgfqpoint{3.196357in}{1.405699in}}%
\pgfpathlineto{\pgfqpoint{3.202838in}{1.419960in}}%
\pgfpathlineto{\pgfqpoint{3.205784in}{1.431832in}}%
\pgfpathlineto{\pgfqpoint{3.212300in}{1.441354in}}%
\pgfpathlineto{\pgfqpoint{3.202783in}{1.447781in}}%
\pgfpathlineto{\pgfqpoint{3.198897in}{1.442058in}}%
\pgfpathlineto{\pgfqpoint{3.189390in}{1.447867in}}%
\pgfpathlineto{\pgfqpoint{3.172555in}{1.451737in}}%
\pgfpathlineto{\pgfqpoint{3.174081in}{1.460247in}}%
\pgfpathlineto{\pgfqpoint{3.166450in}{1.465183in}}%
\pgfpathlineto{\pgfqpoint{3.156626in}{1.464838in}}%
\pgfpathlineto{\pgfqpoint{3.144720in}{1.485442in}}%
\pgfpathlineto{\pgfqpoint{3.136410in}{1.483605in}}%
\pgfpathlineto{\pgfqpoint{3.123780in}{1.488922in}}%
\pgfpathlineto{\pgfqpoint{3.114696in}{1.482885in}}%
\pgfpathlineto{\pgfqpoint{3.108437in}{1.472032in}}%
\pgfpathlineto{\pgfqpoint{3.098951in}{1.471101in}}%
\pgfpathlineto{\pgfqpoint{3.088071in}{1.475996in}}%
\pgfpathlineto{\pgfqpoint{3.080218in}{1.460384in}}%
\pgfpathlineto{\pgfqpoint{3.074880in}{1.463071in}}%
\pgfpathlineto{\pgfqpoint{3.052224in}{1.436127in}}%
\pgfpathlineto{\pgfqpoint{3.045737in}{1.476948in}}%
\pgfpathlineto{\pgfqpoint{3.128532in}{1.492293in}}%
\pgfpathlineto{\pgfqpoint{3.165648in}{1.498939in}}%
\pgfpathlineto{\pgfqpoint{3.219645in}{1.509814in}}%
\pgfpathlineto{\pgfqpoint{3.265920in}{1.519827in}}%
\pgfpathlineto{\pgfqpoint{3.272690in}{1.493733in}}%
\pgfpathlineto{\pgfqpoint{3.292410in}{1.423600in}}%
\pgfpathlineto{\pgfqpoint{3.329269in}{1.430656in}}%
\pgfusepath{stroke}%
\end{pgfscope}%
\begin{pgfscope}%
\pgfpathrectangle{\pgfqpoint{0.100000in}{0.100000in}}{\pgfqpoint{3.608454in}{2.310000in}}%
\pgfusepath{clip}%
\pgfsetbuttcap%
\pgfsetroundjoin%
\pgfsetlinewidth{0.010037pt}%
\definecolor{currentstroke}{rgb}{1.000000,1.000000,1.000000}%
\pgfsetstrokecolor{currentstroke}%
\pgfsetdash{}{0pt}%
\pgfpathmoveto{\pgfqpoint{1.662648in}{1.197867in}}%
\pgfpathlineto{\pgfqpoint{1.602243in}{1.203632in}}%
\pgfpathlineto{\pgfqpoint{1.539577in}{1.209185in}}%
\pgfpathlineto{\pgfqpoint{1.467164in}{1.216739in}}%
\pgfpathlineto{\pgfqpoint{1.359578in}{1.229535in}}%
\pgfpathlineto{\pgfqpoint{1.321411in}{1.235415in}}%
\pgfpathlineto{\pgfqpoint{1.222893in}{1.249723in}}%
\pgfpathlineto{\pgfqpoint{1.237278in}{1.340011in}}%
\pgfpathlineto{\pgfqpoint{1.237637in}{1.347314in}}%
\pgfpathlineto{\pgfqpoint{1.251495in}{1.434433in}}%
\pgfpathlineto{\pgfqpoint{1.261846in}{1.500784in}}%
\pgfpathlineto{\pgfqpoint{1.271614in}{1.562313in}}%
\pgfpathlineto{\pgfqpoint{1.338353in}{1.552733in}}%
\pgfpathlineto{\pgfqpoint{1.400578in}{1.543802in}}%
\pgfpathlineto{\pgfqpoint{1.514964in}{1.529684in}}%
\pgfpathlineto{\pgfqpoint{1.567441in}{1.524892in}}%
\pgfpathlineto{\pgfqpoint{1.650576in}{1.516826in}}%
\pgfpathlineto{\pgfqpoint{1.686540in}{1.513906in}}%
\pgfpathlineto{\pgfqpoint{1.680190in}{1.435164in}}%
\pgfpathlineto{\pgfqpoint{1.674454in}{1.359348in}}%
\pgfpathlineto{\pgfqpoint{1.669837in}{1.297616in}}%
\pgfpathlineto{\pgfqpoint{1.662648in}{1.197867in}}%
\pgfusepath{stroke}%
\end{pgfscope}%
\begin{pgfscope}%
\pgfpathrectangle{\pgfqpoint{0.100000in}{0.100000in}}{\pgfqpoint{3.608454in}{2.310000in}}%
\pgfusepath{clip}%
\pgfsetbuttcap%
\pgfsetroundjoin%
\pgfsetlinewidth{0.010037pt}%
\definecolor{currentstroke}{rgb}{1.000000,1.000000,1.000000}%
\pgfsetstrokecolor{currentstroke}%
\pgfsetdash{}{0pt}%
\pgfpathmoveto{\pgfqpoint{2.461763in}{1.153105in}}%
\pgfpathlineto{\pgfqpoint{2.463799in}{1.162233in}}%
\pgfpathlineto{\pgfqpoint{2.474354in}{1.160178in}}%
\pgfpathlineto{\pgfqpoint{2.477364in}{1.192570in}}%
\pgfpathlineto{\pgfqpoint{2.474476in}{1.198417in}}%
\pgfpathlineto{\pgfqpoint{2.484770in}{1.211436in}}%
\pgfpathlineto{\pgfqpoint{2.493026in}{1.210574in}}%
\pgfpathlineto{\pgfqpoint{2.518051in}{1.202010in}}%
\pgfpathlineto{\pgfqpoint{2.520983in}{1.208943in}}%
\pgfpathlineto{\pgfqpoint{2.514610in}{1.219580in}}%
\pgfpathlineto{\pgfqpoint{2.516761in}{1.227971in}}%
\pgfpathlineto{\pgfqpoint{2.540098in}{1.235688in}}%
\pgfpathlineto{\pgfqpoint{2.534528in}{1.249888in}}%
\pgfpathlineto{\pgfqpoint{2.541376in}{1.261906in}}%
\pgfpathlineto{\pgfqpoint{2.554124in}{1.268651in}}%
\pgfpathlineto{\pgfqpoint{2.580872in}{1.275331in}}%
\pgfpathlineto{\pgfqpoint{2.597748in}{1.265248in}}%
\pgfpathlineto{\pgfqpoint{2.602868in}{1.275081in}}%
\pgfpathlineto{\pgfqpoint{2.616881in}{1.281949in}}%
\pgfpathlineto{\pgfqpoint{2.620497in}{1.275892in}}%
\pgfpathlineto{\pgfqpoint{2.634892in}{1.280705in}}%
\pgfpathlineto{\pgfqpoint{2.633983in}{1.288941in}}%
\pgfpathlineto{\pgfqpoint{2.646926in}{1.298367in}}%
\pgfpathlineto{\pgfqpoint{2.654564in}{1.288579in}}%
\pgfpathlineto{\pgfqpoint{2.664612in}{1.287594in}}%
\pgfpathlineto{\pgfqpoint{2.671290in}{1.293974in}}%
\pgfpathlineto{\pgfqpoint{2.670547in}{1.303046in}}%
\pgfpathlineto{\pgfqpoint{2.676223in}{1.312054in}}%
\pgfpathlineto{\pgfqpoint{2.683829in}{1.314006in}}%
\pgfpathlineto{\pgfqpoint{2.686884in}{1.325901in}}%
\pgfpathlineto{\pgfqpoint{2.697946in}{1.336184in}}%
\pgfpathlineto{\pgfqpoint{2.694616in}{1.346419in}}%
\pgfpathlineto{\pgfqpoint{2.705378in}{1.351515in}}%
\pgfpathlineto{\pgfqpoint{2.712569in}{1.348381in}}%
\pgfpathlineto{\pgfqpoint{2.723189in}{1.356342in}}%
\pgfpathlineto{\pgfqpoint{2.732680in}{1.358417in}}%
\pgfpathlineto{\pgfqpoint{2.732714in}{1.366675in}}%
\pgfpathlineto{\pgfqpoint{2.726063in}{1.378135in}}%
\pgfpathlineto{\pgfqpoint{2.730393in}{1.382751in}}%
\pgfpathlineto{\pgfqpoint{2.743996in}{1.381728in}}%
\pgfpathlineto{\pgfqpoint{2.752646in}{1.386037in}}%
\pgfpathlineto{\pgfqpoint{2.761568in}{1.379305in}}%
\pgfpathlineto{\pgfqpoint{2.768610in}{1.364112in}}%
\pgfpathlineto{\pgfqpoint{2.791958in}{1.361544in}}%
\pgfpathlineto{\pgfqpoint{2.804128in}{1.353181in}}%
\pgfpathlineto{\pgfqpoint{2.812691in}{1.360250in}}%
\pgfpathlineto{\pgfqpoint{2.828173in}{1.354308in}}%
\pgfpathlineto{\pgfqpoint{2.836938in}{1.357261in}}%
\pgfpathlineto{\pgfqpoint{2.843065in}{1.365657in}}%
\pgfpathlineto{\pgfqpoint{2.852417in}{1.367262in}}%
\pgfpathlineto{\pgfqpoint{2.854847in}{1.358134in}}%
\pgfpathlineto{\pgfqpoint{2.865190in}{1.354643in}}%
\pgfpathlineto{\pgfqpoint{2.873228in}{1.345641in}}%
\pgfpathlineto{\pgfqpoint{2.875653in}{1.333484in}}%
\pgfpathlineto{\pgfqpoint{2.872427in}{1.324293in}}%
\pgfpathlineto{\pgfqpoint{2.884769in}{1.312394in}}%
\pgfpathlineto{\pgfqpoint{2.899268in}{1.289925in}}%
\pgfpathlineto{\pgfqpoint{2.910000in}{1.282532in}}%
\pgfpathlineto{\pgfqpoint{2.921647in}{1.282104in}}%
\pgfpathlineto{\pgfqpoint{2.900141in}{1.257427in}}%
\pgfpathlineto{\pgfqpoint{2.878972in}{1.242490in}}%
\pgfpathlineto{\pgfqpoint{2.871181in}{1.230629in}}%
\pgfpathlineto{\pgfqpoint{2.871313in}{1.224191in}}%
\pgfpathlineto{\pgfqpoint{2.859847in}{1.219258in}}%
\pgfpathlineto{\pgfqpoint{2.856561in}{1.209974in}}%
\pgfpathlineto{\pgfqpoint{2.832674in}{1.200623in}}%
\pgfpathlineto{\pgfqpoint{2.824207in}{1.194550in}}%
\pgfpathlineto{\pgfqpoint{2.823062in}{1.193254in}}%
\pgfpathlineto{\pgfqpoint{2.754353in}{1.186739in}}%
\pgfpathlineto{\pgfqpoint{2.712861in}{1.183250in}}%
\pgfpathlineto{\pgfqpoint{2.644773in}{1.179396in}}%
\pgfpathlineto{\pgfqpoint{2.559948in}{1.171005in}}%
\pgfpathlineto{\pgfqpoint{2.545934in}{1.172945in}}%
\pgfpathlineto{\pgfqpoint{2.548849in}{1.158640in}}%
\pgfpathlineto{\pgfqpoint{2.461763in}{1.153105in}}%
\pgfusepath{stroke}%
\end{pgfscope}%
\begin{pgfscope}%
\pgfpathrectangle{\pgfqpoint{0.100000in}{0.100000in}}{\pgfqpoint{3.608454in}{2.310000in}}%
\pgfusepath{clip}%
\pgfsetbuttcap%
\pgfsetroundjoin%
\pgfsetlinewidth{0.010037pt}%
\definecolor{currentstroke}{rgb}{1.000000,1.000000,1.000000}%
\pgfsetstrokecolor{currentstroke}%
\pgfsetdash{}{0pt}%
\pgfpathmoveto{\pgfqpoint{1.662648in}{1.197867in}}%
\pgfpathlineto{\pgfqpoint{1.669837in}{1.297616in}}%
\pgfpathlineto{\pgfqpoint{1.674454in}{1.359348in}}%
\pgfpathlineto{\pgfqpoint{1.680190in}{1.435164in}}%
\pgfpathlineto{\pgfqpoint{1.759716in}{1.429580in}}%
\pgfpathlineto{\pgfqpoint{1.860711in}{1.424063in}}%
\pgfpathlineto{\pgfqpoint{1.929402in}{1.421428in}}%
\pgfpathlineto{\pgfqpoint{1.997690in}{1.419335in}}%
\pgfpathlineto{\pgfqpoint{2.059498in}{1.418379in}}%
\pgfpathlineto{\pgfqpoint{2.088084in}{1.418676in}}%
\pgfpathlineto{\pgfqpoint{2.100667in}{1.408407in}}%
\pgfpathlineto{\pgfqpoint{2.110523in}{1.410477in}}%
\pgfpathlineto{\pgfqpoint{2.110832in}{1.401521in}}%
\pgfpathlineto{\pgfqpoint{2.103512in}{1.386028in}}%
\pgfpathlineto{\pgfqpoint{2.104306in}{1.376240in}}%
\pgfpathlineto{\pgfqpoint{2.111314in}{1.369785in}}%
\pgfpathlineto{\pgfqpoint{2.117751in}{1.356337in}}%
\pgfpathlineto{\pgfqpoint{2.130810in}{1.348616in}}%
\pgfpathlineto{\pgfqpoint{2.130373in}{1.297922in}}%
\pgfpathlineto{\pgfqpoint{2.130758in}{1.181350in}}%
\pgfpathlineto{\pgfqpoint{2.043230in}{1.181753in}}%
\pgfpathlineto{\pgfqpoint{1.951062in}{1.183366in}}%
\pgfpathlineto{\pgfqpoint{1.883211in}{1.185760in}}%
\pgfpathlineto{\pgfqpoint{1.826575in}{1.187945in}}%
\pgfpathlineto{\pgfqpoint{1.731159in}{1.193571in}}%
\pgfpathlineto{\pgfqpoint{1.662648in}{1.197867in}}%
\pgfusepath{stroke}%
\end{pgfscope}%
\begin{pgfscope}%
\pgfpathrectangle{\pgfqpoint{0.100000in}{0.100000in}}{\pgfqpoint{3.608454in}{2.310000in}}%
\pgfusepath{clip}%
\pgfsetbuttcap%
\pgfsetroundjoin%
\pgfsetlinewidth{0.010037pt}%
\definecolor{currentstroke}{rgb}{1.000000,1.000000,1.000000}%
\pgfsetstrokecolor{currentstroke}%
\pgfsetdash{}{0pt}%
\pgfpathmoveto{\pgfqpoint{2.950506in}{1.211075in}}%
\pgfpathlineto{\pgfqpoint{2.853565in}{1.197457in}}%
\pgfpathlineto{\pgfqpoint{2.824207in}{1.194550in}}%
\pgfpathlineto{\pgfqpoint{2.832674in}{1.200623in}}%
\pgfpathlineto{\pgfqpoint{2.856561in}{1.209974in}}%
\pgfpathlineto{\pgfqpoint{2.859847in}{1.219258in}}%
\pgfpathlineto{\pgfqpoint{2.871313in}{1.224191in}}%
\pgfpathlineto{\pgfqpoint{2.871181in}{1.230629in}}%
\pgfpathlineto{\pgfqpoint{2.878972in}{1.242490in}}%
\pgfpathlineto{\pgfqpoint{2.900141in}{1.257427in}}%
\pgfpathlineto{\pgfqpoint{2.921647in}{1.282104in}}%
\pgfpathlineto{\pgfqpoint{2.920600in}{1.276793in}}%
\pgfpathlineto{\pgfqpoint{2.931940in}{1.263968in}}%
\pgfpathlineto{\pgfqpoint{2.950430in}{1.260367in}}%
\pgfpathlineto{\pgfqpoint{2.961186in}{1.272585in}}%
\pgfpathlineto{\pgfqpoint{2.970810in}{1.266151in}}%
\pgfpathlineto{\pgfqpoint{2.988030in}{1.273703in}}%
\pgfpathlineto{\pgfqpoint{2.990493in}{1.281077in}}%
\pgfpathlineto{\pgfqpoint{3.001180in}{1.282737in}}%
\pgfpathlineto{\pgfqpoint{3.012280in}{1.291449in}}%
\pgfpathlineto{\pgfqpoint{3.022064in}{1.294324in}}%
\pgfpathlineto{\pgfqpoint{3.022375in}{1.309715in}}%
\pgfpathlineto{\pgfqpoint{3.028711in}{1.326341in}}%
\pgfpathlineto{\pgfqpoint{3.037068in}{1.336751in}}%
\pgfpathlineto{\pgfqpoint{3.039777in}{1.352383in}}%
\pgfpathlineto{\pgfqpoint{3.050668in}{1.374341in}}%
\pgfpathlineto{\pgfqpoint{3.051058in}{1.387407in}}%
\pgfpathlineto{\pgfqpoint{3.062309in}{1.378894in}}%
\pgfpathlineto{\pgfqpoint{3.073308in}{1.376829in}}%
\pgfpathlineto{\pgfqpoint{3.077588in}{1.383022in}}%
\pgfpathlineto{\pgfqpoint{3.085235in}{1.410012in}}%
\pgfpathlineto{\pgfqpoint{3.095884in}{1.408977in}}%
\pgfpathlineto{\pgfqpoint{3.098705in}{1.418866in}}%
\pgfpathlineto{\pgfqpoint{3.104950in}{1.423326in}}%
\pgfpathlineto{\pgfqpoint{3.118185in}{1.445751in}}%
\pgfpathlineto{\pgfqpoint{3.116263in}{1.452384in}}%
\pgfpathlineto{\pgfqpoint{3.122088in}{1.466802in}}%
\pgfpathlineto{\pgfqpoint{3.153183in}{1.449239in}}%
\pgfpathlineto{\pgfqpoint{3.156626in}{1.464838in}}%
\pgfpathlineto{\pgfqpoint{3.166450in}{1.465183in}}%
\pgfpathlineto{\pgfqpoint{3.174081in}{1.460247in}}%
\pgfpathlineto{\pgfqpoint{3.172555in}{1.451737in}}%
\pgfpathlineto{\pgfqpoint{3.189390in}{1.447867in}}%
\pgfpathlineto{\pgfqpoint{3.198897in}{1.442058in}}%
\pgfpathlineto{\pgfqpoint{3.205784in}{1.431832in}}%
\pgfpathlineto{\pgfqpoint{3.202838in}{1.419960in}}%
\pgfpathlineto{\pgfqpoint{3.196888in}{1.418988in}}%
\pgfpathlineto{\pgfqpoint{3.193438in}{1.401073in}}%
\pgfpathlineto{\pgfqpoint{3.200996in}{1.394069in}}%
\pgfpathlineto{\pgfqpoint{3.211636in}{1.399731in}}%
\pgfpathlineto{\pgfqpoint{3.221508in}{1.387775in}}%
\pgfpathlineto{\pgfqpoint{3.243566in}{1.385626in}}%
\pgfpathlineto{\pgfqpoint{3.247366in}{1.378749in}}%
\pgfpathlineto{\pgfqpoint{3.267778in}{1.372060in}}%
\pgfpathlineto{\pgfqpoint{3.265259in}{1.364168in}}%
\pgfpathlineto{\pgfqpoint{3.267915in}{1.341861in}}%
\pgfpathlineto{\pgfqpoint{3.276154in}{1.333438in}}%
\pgfpathlineto{\pgfqpoint{3.263997in}{1.329883in}}%
\pgfpathlineto{\pgfqpoint{3.265608in}{1.320360in}}%
\pgfpathlineto{\pgfqpoint{3.278578in}{1.312286in}}%
\pgfpathlineto{\pgfqpoint{3.279747in}{1.304321in}}%
\pgfpathlineto{\pgfqpoint{3.272425in}{1.298336in}}%
\pgfpathlineto{\pgfqpoint{3.257650in}{1.312515in}}%
\pgfpathlineto{\pgfqpoint{3.256190in}{1.302170in}}%
\pgfpathlineto{\pgfqpoint{3.268558in}{1.297251in}}%
\pgfpathlineto{\pgfqpoint{3.269796in}{1.292163in}}%
\pgfpathlineto{\pgfqpoint{3.286760in}{1.298887in}}%
\pgfpathlineto{\pgfqpoint{3.299759in}{1.300665in}}%
\pgfpathlineto{\pgfqpoint{3.313076in}{1.273798in}}%
\pgfpathlineto{\pgfqpoint{3.311590in}{1.273507in}}%
\pgfpathlineto{\pgfqpoint{3.305573in}{1.272263in}}%
\pgfpathlineto{\pgfqpoint{3.303799in}{1.271893in}}%
\pgfpathlineto{\pgfqpoint{3.302627in}{1.271664in}}%
\pgfpathlineto{\pgfqpoint{3.223341in}{1.255217in}}%
\pgfpathlineto{\pgfqpoint{3.152426in}{1.240702in}}%
\pgfpathlineto{\pgfqpoint{3.054377in}{1.223812in}}%
\pgfpathlineto{\pgfqpoint{2.971105in}{1.212813in}}%
\pgfpathlineto{\pgfqpoint{2.950506in}{1.211075in}}%
\pgfusepath{stroke}%
\end{pgfscope}%
\begin{pgfscope}%
\pgfpathrectangle{\pgfqpoint{0.100000in}{0.100000in}}{\pgfqpoint{3.608454in}{2.310000in}}%
\pgfusepath{clip}%
\pgfsetbuttcap%
\pgfsetroundjoin%
\pgfsetlinewidth{0.010037pt}%
\definecolor{currentstroke}{rgb}{1.000000,1.000000,1.000000}%
\pgfsetstrokecolor{currentstroke}%
\pgfsetdash{}{0pt}%
\pgfpathmoveto{\pgfqpoint{3.303727in}{1.386181in}}%
\pgfpathlineto{\pgfqpoint{3.319710in}{1.393431in}}%
\pgfpathlineto{\pgfqpoint{3.306946in}{1.356091in}}%
\pgfpathlineto{\pgfqpoint{3.298505in}{1.336318in}}%
\pgfpathlineto{\pgfqpoint{3.292424in}{1.347508in}}%
\pgfpathlineto{\pgfqpoint{3.303237in}{1.374298in}}%
\pgfpathlineto{\pgfqpoint{3.303727in}{1.386181in}}%
\pgfusepath{stroke}%
\end{pgfscope}%
\begin{pgfscope}%
\pgfpathrectangle{\pgfqpoint{0.100000in}{0.100000in}}{\pgfqpoint{3.608454in}{2.310000in}}%
\pgfusepath{clip}%
\pgfsetbuttcap%
\pgfsetroundjoin%
\pgfsetlinewidth{0.010037pt}%
\definecolor{currentstroke}{rgb}{1.000000,1.000000,1.000000}%
\pgfsetstrokecolor{currentstroke}%
\pgfsetdash{}{0pt}%
\pgfpathmoveto{\pgfqpoint{2.461763in}{1.153105in}}%
\pgfpathlineto{\pgfqpoint{2.457903in}{1.152546in}}%
\pgfpathlineto{\pgfqpoint{2.454263in}{1.152291in}}%
\pgfpathlineto{\pgfqpoint{2.455821in}{1.141115in}}%
\pgfpathlineto{\pgfqpoint{2.446430in}{1.129849in}}%
\pgfpathlineto{\pgfqpoint{2.444579in}{1.112222in}}%
\pgfpathlineto{\pgfqpoint{2.402597in}{1.109117in}}%
\pgfpathlineto{\pgfqpoint{2.406257in}{1.117394in}}%
\pgfpathlineto{\pgfqpoint{2.421394in}{1.132509in}}%
\pgfpathlineto{\pgfqpoint{2.421813in}{1.141271in}}%
\pgfpathlineto{\pgfqpoint{2.415111in}{1.149567in}}%
\pgfpathlineto{\pgfqpoint{2.352626in}{1.146263in}}%
\pgfpathlineto{\pgfqpoint{2.243380in}{1.142795in}}%
\pgfpathlineto{\pgfqpoint{2.179448in}{1.141630in}}%
\pgfpathlineto{\pgfqpoint{2.131124in}{1.141195in}}%
\pgfpathlineto{\pgfqpoint{2.130758in}{1.181350in}}%
\pgfpathlineto{\pgfqpoint{2.130373in}{1.297922in}}%
\pgfpathlineto{\pgfqpoint{2.130810in}{1.348616in}}%
\pgfpathlineto{\pgfqpoint{2.117751in}{1.356337in}}%
\pgfpathlineto{\pgfqpoint{2.111314in}{1.369785in}}%
\pgfpathlineto{\pgfqpoint{2.104306in}{1.376240in}}%
\pgfpathlineto{\pgfqpoint{2.103512in}{1.386028in}}%
\pgfpathlineto{\pgfqpoint{2.110832in}{1.401521in}}%
\pgfpathlineto{\pgfqpoint{2.110523in}{1.410477in}}%
\pgfpathlineto{\pgfqpoint{2.100667in}{1.408407in}}%
\pgfpathlineto{\pgfqpoint{2.088084in}{1.418676in}}%
\pgfpathlineto{\pgfqpoint{2.077998in}{1.436700in}}%
\pgfpathlineto{\pgfqpoint{2.069537in}{1.445017in}}%
\pgfpathlineto{\pgfqpoint{2.060697in}{1.465454in}}%
\pgfpathlineto{\pgfqpoint{2.152519in}{1.464021in}}%
\pgfpathlineto{\pgfqpoint{2.243792in}{1.467029in}}%
\pgfpathlineto{\pgfqpoint{2.302315in}{1.470405in}}%
\pgfpathlineto{\pgfqpoint{2.320622in}{1.452274in}}%
\pgfpathlineto{\pgfqpoint{2.317277in}{1.433365in}}%
\pgfpathlineto{\pgfqpoint{2.322158in}{1.412077in}}%
\pgfpathlineto{\pgfqpoint{2.327573in}{1.401342in}}%
\pgfpathlineto{\pgfqpoint{2.344850in}{1.386576in}}%
\pgfpathlineto{\pgfqpoint{2.348630in}{1.380053in}}%
\pgfpathlineto{\pgfqpoint{2.367545in}{1.365354in}}%
\pgfpathlineto{\pgfqpoint{2.370287in}{1.348866in}}%
\pgfpathlineto{\pgfqpoint{2.375719in}{1.336702in}}%
\pgfpathlineto{\pgfqpoint{2.388555in}{1.343787in}}%
\pgfpathlineto{\pgfqpoint{2.399134in}{1.340547in}}%
\pgfpathlineto{\pgfqpoint{2.407122in}{1.331947in}}%
\pgfpathlineto{\pgfqpoint{2.404106in}{1.316728in}}%
\pgfpathlineto{\pgfqpoint{2.394301in}{1.296831in}}%
\pgfpathlineto{\pgfqpoint{2.394453in}{1.286467in}}%
\pgfpathlineto{\pgfqpoint{2.401780in}{1.277970in}}%
\pgfpathlineto{\pgfqpoint{2.417550in}{1.266715in}}%
\pgfpathlineto{\pgfqpoint{2.435708in}{1.257580in}}%
\pgfpathlineto{\pgfqpoint{2.440014in}{1.250265in}}%
\pgfpathlineto{\pgfqpoint{2.449301in}{1.246806in}}%
\pgfpathlineto{\pgfqpoint{2.449670in}{1.237099in}}%
\pgfpathlineto{\pgfqpoint{2.456543in}{1.224768in}}%
\pgfpathlineto{\pgfqpoint{2.451774in}{1.213813in}}%
\pgfpathlineto{\pgfqpoint{2.461108in}{1.196383in}}%
\pgfpathlineto{\pgfqpoint{2.468410in}{1.199844in}}%
\pgfpathlineto{\pgfqpoint{2.477364in}{1.192570in}}%
\pgfpathlineto{\pgfqpoint{2.474354in}{1.160178in}}%
\pgfpathlineto{\pgfqpoint{2.463799in}{1.162233in}}%
\pgfpathlineto{\pgfqpoint{2.461763in}{1.153105in}}%
\pgfusepath{stroke}%
\end{pgfscope}%
\begin{pgfscope}%
\pgfpathrectangle{\pgfqpoint{0.100000in}{0.100000in}}{\pgfqpoint{3.608454in}{2.310000in}}%
\pgfusepath{clip}%
\pgfsetbuttcap%
\pgfsetroundjoin%
\pgfsetlinewidth{0.010037pt}%
\definecolor{currentstroke}{rgb}{1.000000,1.000000,1.000000}%
\pgfsetstrokecolor{currentstroke}%
\pgfsetdash{}{0pt}%
\pgfpathmoveto{\pgfqpoint{0.843082in}{1.160088in}}%
\pgfpathlineto{\pgfqpoint{0.850725in}{1.176360in}}%
\pgfpathlineto{\pgfqpoint{0.848633in}{1.200785in}}%
\pgfpathlineto{\pgfqpoint{0.852293in}{1.207731in}}%
\pgfpathlineto{\pgfqpoint{0.852655in}{1.238156in}}%
\pgfpathlineto{\pgfqpoint{0.856124in}{1.246917in}}%
\pgfpathlineto{\pgfqpoint{0.868345in}{1.248280in}}%
\pgfpathlineto{\pgfqpoint{0.879687in}{1.244390in}}%
\pgfpathlineto{\pgfqpoint{0.884626in}{1.233679in}}%
\pgfpathlineto{\pgfqpoint{0.891552in}{1.234089in}}%
\pgfpathlineto{\pgfqpoint{0.900156in}{1.246360in}}%
\pgfpathlineto{\pgfqpoint{0.912593in}{1.306845in}}%
\pgfpathlineto{\pgfqpoint{0.983361in}{1.292233in}}%
\pgfpathlineto{\pgfqpoint{1.024434in}{1.284134in}}%
\pgfpathlineto{\pgfqpoint{1.133501in}{1.264961in}}%
\pgfpathlineto{\pgfqpoint{1.163676in}{1.258901in}}%
\pgfpathlineto{\pgfqpoint{1.222893in}{1.249723in}}%
\pgfpathlineto{\pgfqpoint{1.210751in}{1.171542in}}%
\pgfpathlineto{\pgfqpoint{1.198122in}{1.089991in}}%
\pgfpathlineto{\pgfqpoint{1.183580in}{0.998244in}}%
\pgfpathlineto{\pgfqpoint{1.167224in}{0.892929in}}%
\pgfpathlineto{\pgfqpoint{1.154011in}{0.806432in}}%
\pgfpathlineto{\pgfqpoint{1.059373in}{0.821460in}}%
\pgfpathlineto{\pgfqpoint{1.017791in}{0.828581in}}%
\pgfpathlineto{\pgfqpoint{0.999216in}{0.839702in}}%
\pgfpathlineto{\pgfqpoint{0.878105in}{0.912775in}}%
\pgfpathlineto{\pgfqpoint{0.787216in}{0.968230in}}%
\pgfpathlineto{\pgfqpoint{0.790243in}{0.978058in}}%
\pgfpathlineto{\pgfqpoint{0.797751in}{0.984935in}}%
\pgfpathlineto{\pgfqpoint{0.805560in}{0.983661in}}%
\pgfpathlineto{\pgfqpoint{0.816886in}{0.990883in}}%
\pgfpathlineto{\pgfqpoint{0.818669in}{1.001263in}}%
\pgfpathlineto{\pgfqpoint{0.804856in}{1.013854in}}%
\pgfpathlineto{\pgfqpoint{0.814680in}{1.038025in}}%
\pgfpathlineto{\pgfqpoint{0.824557in}{1.047321in}}%
\pgfpathlineto{\pgfqpoint{0.829245in}{1.058349in}}%
\pgfpathlineto{\pgfqpoint{0.832133in}{1.078570in}}%
\pgfpathlineto{\pgfqpoint{0.841395in}{1.087727in}}%
\pgfpathlineto{\pgfqpoint{0.860834in}{1.096874in}}%
\pgfpathlineto{\pgfqpoint{0.860405in}{1.104932in}}%
\pgfpathlineto{\pgfqpoint{0.849548in}{1.114934in}}%
\pgfpathlineto{\pgfqpoint{0.848070in}{1.135552in}}%
\pgfpathlineto{\pgfqpoint{0.840594in}{1.150624in}}%
\pgfpathlineto{\pgfqpoint{0.843082in}{1.160088in}}%
\pgfusepath{stroke}%
\end{pgfscope}%
\begin{pgfscope}%
\pgfpathrectangle{\pgfqpoint{0.100000in}{0.100000in}}{\pgfqpoint{3.608454in}{2.310000in}}%
\pgfusepath{clip}%
\pgfsetbuttcap%
\pgfsetroundjoin%
\pgfsetlinewidth{0.010037pt}%
\definecolor{currentstroke}{rgb}{1.000000,1.000000,1.000000}%
\pgfsetstrokecolor{currentstroke}%
\pgfsetdash{}{0pt}%
\pgfpathmoveto{\pgfqpoint{2.141369in}{0.914830in}}%
\pgfpathlineto{\pgfqpoint{2.123825in}{0.920256in}}%
\pgfpathlineto{\pgfqpoint{2.101346in}{0.937472in}}%
\pgfpathlineto{\pgfqpoint{2.091360in}{0.941172in}}%
\pgfpathlineto{\pgfqpoint{2.085016in}{0.933736in}}%
\pgfpathlineto{\pgfqpoint{2.077003in}{0.933361in}}%
\pgfpathlineto{\pgfqpoint{2.066866in}{0.939673in}}%
\pgfpathlineto{\pgfqpoint{2.050959in}{0.931585in}}%
\pgfpathlineto{\pgfqpoint{2.045296in}{0.935568in}}%
\pgfpathlineto{\pgfqpoint{2.029643in}{0.926132in}}%
\pgfpathlineto{\pgfqpoint{2.013130in}{0.927873in}}%
\pgfpathlineto{\pgfqpoint{2.000393in}{0.934005in}}%
\pgfpathlineto{\pgfqpoint{1.991473in}{0.931697in}}%
\pgfpathlineto{\pgfqpoint{1.977195in}{0.940369in}}%
\pgfpathlineto{\pgfqpoint{1.969253in}{0.925074in}}%
\pgfpathlineto{\pgfqpoint{1.961130in}{0.938228in}}%
\pgfpathlineto{\pgfqpoint{1.945087in}{0.933125in}}%
\pgfpathlineto{\pgfqpoint{1.931055in}{0.945619in}}%
\pgfpathlineto{\pgfqpoint{1.918793in}{0.935548in}}%
\pgfpathlineto{\pgfqpoint{1.912652in}{0.944802in}}%
\pgfpathlineto{\pgfqpoint{1.903801in}{0.947808in}}%
\pgfpathlineto{\pgfqpoint{1.898410in}{0.956730in}}%
\pgfpathlineto{\pgfqpoint{1.886834in}{0.959269in}}%
\pgfpathlineto{\pgfqpoint{1.880154in}{0.952553in}}%
\pgfpathlineto{\pgfqpoint{1.868804in}{0.961257in}}%
\pgfpathlineto{\pgfqpoint{1.863513in}{0.959270in}}%
\pgfpathlineto{\pgfqpoint{1.844702in}{0.966322in}}%
\pgfpathlineto{\pgfqpoint{1.832921in}{0.967109in}}%
\pgfpathlineto{\pgfqpoint{1.827647in}{0.982085in}}%
\pgfpathlineto{\pgfqpoint{1.818203in}{0.980220in}}%
\pgfpathlineto{\pgfqpoint{1.807385in}{0.983916in}}%
\pgfpathlineto{\pgfqpoint{1.800261in}{0.981780in}}%
\pgfpathlineto{\pgfqpoint{1.784985in}{0.998601in}}%
\pgfpathlineto{\pgfqpoint{1.780728in}{0.997501in}}%
\pgfpathlineto{\pgfqpoint{1.784594in}{1.065714in}}%
\pgfpathlineto{\pgfqpoint{1.788806in}{1.150144in}}%
\pgfpathlineto{\pgfqpoint{1.719672in}{1.154011in}}%
\pgfpathlineto{\pgfqpoint{1.651450in}{1.159169in}}%
\pgfpathlineto{\pgfqpoint{1.598743in}{1.163746in}}%
\pgfpathlineto{\pgfqpoint{1.602243in}{1.203632in}}%
\pgfpathlineto{\pgfqpoint{1.662648in}{1.197867in}}%
\pgfpathlineto{\pgfqpoint{1.731159in}{1.193571in}}%
\pgfpathlineto{\pgfqpoint{1.826575in}{1.187945in}}%
\pgfpathlineto{\pgfqpoint{1.883211in}{1.185760in}}%
\pgfpathlineto{\pgfqpoint{1.951062in}{1.183366in}}%
\pgfpathlineto{\pgfqpoint{2.043230in}{1.181753in}}%
\pgfpathlineto{\pgfqpoint{2.130758in}{1.181350in}}%
\pgfpathlineto{\pgfqpoint{2.131124in}{1.141195in}}%
\pgfpathlineto{\pgfqpoint{2.136037in}{1.110943in}}%
\pgfpathlineto{\pgfqpoint{2.143669in}{1.055067in}}%
\pgfpathlineto{\pgfqpoint{2.142095in}{0.959646in}}%
\pgfpathlineto{\pgfqpoint{2.141369in}{0.914830in}}%
\pgfusepath{stroke}%
\end{pgfscope}%
\begin{pgfscope}%
\pgfpathrectangle{\pgfqpoint{0.100000in}{0.100000in}}{\pgfqpoint{3.608454in}{2.310000in}}%
\pgfusepath{clip}%
\pgfsetbuttcap%
\pgfsetroundjoin%
\pgfsetlinewidth{0.010037pt}%
\definecolor{currentstroke}{rgb}{1.000000,1.000000,1.000000}%
\pgfsetstrokecolor{currentstroke}%
\pgfsetdash{}{0pt}%
\pgfpathmoveto{\pgfqpoint{2.798250in}{1.063020in}}%
\pgfpathlineto{\pgfqpoint{2.798302in}{1.080700in}}%
\pgfpathlineto{\pgfqpoint{2.813666in}{1.087489in}}%
\pgfpathlineto{\pgfqpoint{2.814310in}{1.098342in}}%
\pgfpathlineto{\pgfqpoint{2.828067in}{1.111600in}}%
\pgfpathlineto{\pgfqpoint{2.845262in}{1.114316in}}%
\pgfpathlineto{\pgfqpoint{2.860379in}{1.128927in}}%
\pgfpathlineto{\pgfqpoint{2.876166in}{1.135457in}}%
\pgfpathlineto{\pgfqpoint{2.888015in}{1.155063in}}%
\pgfpathlineto{\pgfqpoint{2.898807in}{1.153640in}}%
\pgfpathlineto{\pgfqpoint{2.921614in}{1.171504in}}%
\pgfpathlineto{\pgfqpoint{2.933658in}{1.171840in}}%
\pgfpathlineto{\pgfqpoint{2.938743in}{1.185452in}}%
\pgfpathlineto{\pgfqpoint{2.948316in}{1.194928in}}%
\pgfpathlineto{\pgfqpoint{2.950506in}{1.211075in}}%
\pgfpathlineto{\pgfqpoint{2.971105in}{1.212813in}}%
\pgfpathlineto{\pgfqpoint{3.054377in}{1.223812in}}%
\pgfpathlineto{\pgfqpoint{3.152426in}{1.240702in}}%
\pgfpathlineto{\pgfqpoint{3.223341in}{1.255217in}}%
\pgfpathlineto{\pgfqpoint{3.302627in}{1.271664in}}%
\pgfpathlineto{\pgfqpoint{3.325283in}{1.240558in}}%
\pgfpathlineto{\pgfqpoint{3.313011in}{1.242544in}}%
\pgfpathlineto{\pgfqpoint{3.297258in}{1.238706in}}%
\pgfpathlineto{\pgfqpoint{3.281735in}{1.222854in}}%
\pgfpathlineto{\pgfqpoint{3.270579in}{1.223994in}}%
\pgfpathlineto{\pgfqpoint{3.269168in}{1.214579in}}%
\pgfpathlineto{\pgfqpoint{3.289332in}{1.222025in}}%
\pgfpathlineto{\pgfqpoint{3.309624in}{1.225095in}}%
\pgfpathlineto{\pgfqpoint{3.317158in}{1.221010in}}%
\pgfpathlineto{\pgfqpoint{3.327297in}{1.225667in}}%
\pgfpathlineto{\pgfqpoint{3.332529in}{1.222404in}}%
\pgfpathlineto{\pgfqpoint{3.337160in}{1.206867in}}%
\pgfpathlineto{\pgfqpoint{3.327523in}{1.202055in}}%
\pgfpathlineto{\pgfqpoint{3.320908in}{1.183109in}}%
\pgfpathlineto{\pgfqpoint{3.313973in}{1.175732in}}%
\pgfpathlineto{\pgfqpoint{3.292714in}{1.177346in}}%
\pgfpathlineto{\pgfqpoint{3.293856in}{1.188471in}}%
\pgfpathlineto{\pgfqpoint{3.282246in}{1.183611in}}%
\pgfpathlineto{\pgfqpoint{3.279721in}{1.174333in}}%
\pgfpathlineto{\pgfqpoint{3.287679in}{1.164722in}}%
\pgfpathlineto{\pgfqpoint{3.290383in}{1.148409in}}%
\pgfpathlineto{\pgfqpoint{3.277360in}{1.139115in}}%
\pgfpathlineto{\pgfqpoint{3.291444in}{1.135846in}}%
\pgfpathlineto{\pgfqpoint{3.297817in}{1.142680in}}%
\pgfpathlineto{\pgfqpoint{3.311909in}{1.143512in}}%
\pgfpathlineto{\pgfqpoint{3.304667in}{1.128754in}}%
\pgfpathlineto{\pgfqpoint{3.295798in}{1.121643in}}%
\pgfpathlineto{\pgfqpoint{3.268968in}{1.114523in}}%
\pgfpathlineto{\pgfqpoint{3.241422in}{1.089539in}}%
\pgfpathlineto{\pgfqpoint{3.224436in}{1.064923in}}%
\pgfpathlineto{\pgfqpoint{3.224336in}{1.054925in}}%
\pgfpathlineto{\pgfqpoint{3.217551in}{1.041130in}}%
\pgfpathlineto{\pgfqpoint{3.182729in}{1.032053in}}%
\pgfpathlineto{\pgfqpoint{3.099917in}{1.091416in}}%
\pgfpathlineto{\pgfqpoint{3.026989in}{1.080701in}}%
\pgfpathlineto{\pgfqpoint{3.026378in}{1.090595in}}%
\pgfpathlineto{\pgfqpoint{3.015293in}{1.101737in}}%
\pgfpathlineto{\pgfqpoint{3.006902in}{1.104464in}}%
\pgfpathlineto{\pgfqpoint{2.927665in}{1.096231in}}%
\pgfpathlineto{\pgfqpoint{2.909453in}{1.090054in}}%
\pgfpathlineto{\pgfqpoint{2.876568in}{1.073696in}}%
\pgfpathlineto{\pgfqpoint{2.848149in}{1.069168in}}%
\pgfpathlineto{\pgfqpoint{2.798250in}{1.063020in}}%
\pgfusepath{stroke}%
\end{pgfscope}%
\begin{pgfscope}%
\pgfpathrectangle{\pgfqpoint{0.100000in}{0.100000in}}{\pgfqpoint{3.608454in}{2.310000in}}%
\pgfusepath{clip}%
\pgfsetbuttcap%
\pgfsetroundjoin%
\pgfsetlinewidth{0.010037pt}%
\definecolor{currentstroke}{rgb}{1.000000,1.000000,1.000000}%
\pgfsetstrokecolor{currentstroke}%
\pgfsetdash{}{0pt}%
\pgfpathmoveto{\pgfqpoint{2.798250in}{1.063020in}}%
\pgfpathlineto{\pgfqpoint{2.715320in}{1.053912in}}%
\pgfpathlineto{\pgfqpoint{2.639455in}{1.047089in}}%
\pgfpathlineto{\pgfqpoint{2.561482in}{1.042047in}}%
\pgfpathlineto{\pgfqpoint{2.548098in}{1.040107in}}%
\pgfpathlineto{\pgfqpoint{2.495536in}{1.036064in}}%
\pgfpathlineto{\pgfqpoint{2.411349in}{1.031151in}}%
\pgfpathlineto{\pgfqpoint{2.426327in}{1.043584in}}%
\pgfpathlineto{\pgfqpoint{2.422910in}{1.056671in}}%
\pgfpathlineto{\pgfqpoint{2.426220in}{1.072535in}}%
\pgfpathlineto{\pgfqpoint{2.431278in}{1.080004in}}%
\pgfpathlineto{\pgfqpoint{2.431088in}{1.090395in}}%
\pgfpathlineto{\pgfqpoint{2.444556in}{1.096930in}}%
\pgfpathlineto{\pgfqpoint{2.444579in}{1.112222in}}%
\pgfpathlineto{\pgfqpoint{2.446430in}{1.129849in}}%
\pgfpathlineto{\pgfqpoint{2.455821in}{1.141115in}}%
\pgfpathlineto{\pgfqpoint{2.454263in}{1.152291in}}%
\pgfpathlineto{\pgfqpoint{2.457903in}{1.152546in}}%
\pgfpathlineto{\pgfqpoint{2.461763in}{1.153105in}}%
\pgfpathlineto{\pgfqpoint{2.548849in}{1.158640in}}%
\pgfpathlineto{\pgfqpoint{2.545934in}{1.172945in}}%
\pgfpathlineto{\pgfqpoint{2.559948in}{1.171005in}}%
\pgfpathlineto{\pgfqpoint{2.644773in}{1.179396in}}%
\pgfpathlineto{\pgfqpoint{2.712861in}{1.183250in}}%
\pgfpathlineto{\pgfqpoint{2.754353in}{1.186739in}}%
\pgfpathlineto{\pgfqpoint{2.823062in}{1.193254in}}%
\pgfpathlineto{\pgfqpoint{2.824207in}{1.194550in}}%
\pgfpathlineto{\pgfqpoint{2.853565in}{1.197457in}}%
\pgfpathlineto{\pgfqpoint{2.950506in}{1.211075in}}%
\pgfpathlineto{\pgfqpoint{2.948316in}{1.194928in}}%
\pgfpathlineto{\pgfqpoint{2.938743in}{1.185452in}}%
\pgfpathlineto{\pgfqpoint{2.933658in}{1.171840in}}%
\pgfpathlineto{\pgfqpoint{2.921614in}{1.171504in}}%
\pgfpathlineto{\pgfqpoint{2.898807in}{1.153640in}}%
\pgfpathlineto{\pgfqpoint{2.888015in}{1.155063in}}%
\pgfpathlineto{\pgfqpoint{2.876166in}{1.135457in}}%
\pgfpathlineto{\pgfqpoint{2.860379in}{1.128927in}}%
\pgfpathlineto{\pgfqpoint{2.845262in}{1.114316in}}%
\pgfpathlineto{\pgfqpoint{2.828067in}{1.111600in}}%
\pgfpathlineto{\pgfqpoint{2.814310in}{1.098342in}}%
\pgfpathlineto{\pgfqpoint{2.813666in}{1.087489in}}%
\pgfpathlineto{\pgfqpoint{2.798302in}{1.080700in}}%
\pgfpathlineto{\pgfqpoint{2.798250in}{1.063020in}}%
\pgfusepath{stroke}%
\end{pgfscope}%
\begin{pgfscope}%
\pgfpathrectangle{\pgfqpoint{0.100000in}{0.100000in}}{\pgfqpoint{3.608454in}{2.310000in}}%
\pgfusepath{clip}%
\pgfsetbuttcap%
\pgfsetroundjoin%
\pgfsetlinewidth{0.010037pt}%
\definecolor{currentstroke}{rgb}{1.000000,1.000000,1.000000}%
\pgfsetstrokecolor{currentstroke}%
\pgfsetdash{}{0pt}%
\pgfpathmoveto{\pgfqpoint{1.327465in}{0.817498in}}%
\pgfpathlineto{\pgfqpoint{1.321550in}{0.827003in}}%
\pgfpathlineto{\pgfqpoint{1.323997in}{0.835189in}}%
\pgfpathlineto{\pgfqpoint{1.442877in}{0.821252in}}%
\pgfpathlineto{\pgfqpoint{1.563330in}{0.809245in}}%
\pgfpathlineto{\pgfqpoint{1.566842in}{0.849964in}}%
\pgfpathlineto{\pgfqpoint{1.584946in}{1.026087in}}%
\pgfpathlineto{\pgfqpoint{1.596844in}{1.163852in}}%
\pgfpathlineto{\pgfqpoint{1.598743in}{1.163746in}}%
\pgfpathlineto{\pgfqpoint{1.719672in}{1.154011in}}%
\pgfpathlineto{\pgfqpoint{1.788806in}{1.150144in}}%
\pgfpathlineto{\pgfqpoint{1.780728in}{0.997501in}}%
\pgfpathlineto{\pgfqpoint{1.784985in}{0.998601in}}%
\pgfpathlineto{\pgfqpoint{1.800261in}{0.981780in}}%
\pgfpathlineto{\pgfqpoint{1.807385in}{0.983916in}}%
\pgfpathlineto{\pgfqpoint{1.818203in}{0.980220in}}%
\pgfpathlineto{\pgfqpoint{1.827647in}{0.982085in}}%
\pgfpathlineto{\pgfqpoint{1.832921in}{0.967109in}}%
\pgfpathlineto{\pgfqpoint{1.844702in}{0.966322in}}%
\pgfpathlineto{\pgfqpoint{1.863513in}{0.959270in}}%
\pgfpathlineto{\pgfqpoint{1.868804in}{0.961257in}}%
\pgfpathlineto{\pgfqpoint{1.880154in}{0.952553in}}%
\pgfpathlineto{\pgfqpoint{1.886834in}{0.959269in}}%
\pgfpathlineto{\pgfqpoint{1.898410in}{0.956730in}}%
\pgfpathlineto{\pgfqpoint{1.903801in}{0.947808in}}%
\pgfpathlineto{\pgfqpoint{1.912652in}{0.944802in}}%
\pgfpathlineto{\pgfqpoint{1.918793in}{0.935548in}}%
\pgfpathlineto{\pgfqpoint{1.931055in}{0.945619in}}%
\pgfpathlineto{\pgfqpoint{1.945087in}{0.933125in}}%
\pgfpathlineto{\pgfqpoint{1.961130in}{0.938228in}}%
\pgfpathlineto{\pgfqpoint{1.969253in}{0.925074in}}%
\pgfpathlineto{\pgfqpoint{1.977195in}{0.940369in}}%
\pgfpathlineto{\pgfqpoint{1.991473in}{0.931697in}}%
\pgfpathlineto{\pgfqpoint{2.000393in}{0.934005in}}%
\pgfpathlineto{\pgfqpoint{2.013130in}{0.927873in}}%
\pgfpathlineto{\pgfqpoint{2.029643in}{0.926132in}}%
\pgfpathlineto{\pgfqpoint{2.045296in}{0.935568in}}%
\pgfpathlineto{\pgfqpoint{2.050959in}{0.931585in}}%
\pgfpathlineto{\pgfqpoint{2.066866in}{0.939673in}}%
\pgfpathlineto{\pgfqpoint{2.077003in}{0.933361in}}%
\pgfpathlineto{\pgfqpoint{2.085016in}{0.933736in}}%
\pgfpathlineto{\pgfqpoint{2.091360in}{0.941172in}}%
\pgfpathlineto{\pgfqpoint{2.101346in}{0.937472in}}%
\pgfpathlineto{\pgfqpoint{2.123825in}{0.920256in}}%
\pgfpathlineto{\pgfqpoint{2.141369in}{0.914830in}}%
\pgfpathlineto{\pgfqpoint{2.148404in}{0.908183in}}%
\pgfpathlineto{\pgfqpoint{2.157210in}{0.911803in}}%
\pgfpathlineto{\pgfqpoint{2.170554in}{0.909031in}}%
\pgfpathlineto{\pgfqpoint{2.171929in}{0.784895in}}%
\pgfpathlineto{\pgfqpoint{2.181197in}{0.777032in}}%
\pgfpathlineto{\pgfqpoint{2.188625in}{0.762543in}}%
\pgfpathlineto{\pgfqpoint{2.186014in}{0.752849in}}%
\pgfpathlineto{\pgfqpoint{2.191847in}{0.748745in}}%
\pgfpathlineto{\pgfqpoint{2.196588in}{0.730869in}}%
\pgfpathlineto{\pgfqpoint{2.205667in}{0.721296in}}%
\pgfpathlineto{\pgfqpoint{2.207701in}{0.700970in}}%
\pgfpathlineto{\pgfqpoint{2.206122in}{0.692365in}}%
\pgfpathlineto{\pgfqpoint{2.193781in}{0.669590in}}%
\pgfpathlineto{\pgfqpoint{2.192362in}{0.653216in}}%
\pgfpathlineto{\pgfqpoint{2.196548in}{0.649798in}}%
\pgfpathlineto{\pgfqpoint{2.196912in}{0.635452in}}%
\pgfpathlineto{\pgfqpoint{2.192657in}{0.626436in}}%
\pgfpathlineto{\pgfqpoint{2.186039in}{0.623040in}}%
\pgfpathlineto{\pgfqpoint{2.179559in}{0.611258in}}%
\pgfpathlineto{\pgfqpoint{2.187818in}{0.599847in}}%
\pgfpathlineto{\pgfqpoint{2.171786in}{0.599622in}}%
\pgfpathlineto{\pgfqpoint{2.128920in}{0.580009in}}%
\pgfpathlineto{\pgfqpoint{2.137104in}{0.591747in}}%
\pgfpathlineto{\pgfqpoint{2.127193in}{0.598075in}}%
\pgfpathlineto{\pgfqpoint{2.125123in}{0.608834in}}%
\pgfpathlineto{\pgfqpoint{2.105780in}{0.590093in}}%
\pgfpathlineto{\pgfqpoint{2.114364in}{0.577284in}}%
\pgfpathlineto{\pgfqpoint{2.102118in}{0.560974in}}%
\pgfpathlineto{\pgfqpoint{2.095528in}{0.561316in}}%
\pgfpathlineto{\pgfqpoint{2.089328in}{0.543553in}}%
\pgfpathlineto{\pgfqpoint{2.069722in}{0.529582in}}%
\pgfpathlineto{\pgfqpoint{2.051405in}{0.524547in}}%
\pgfpathlineto{\pgfqpoint{2.019457in}{0.511448in}}%
\pgfpathlineto{\pgfqpoint{2.016147in}{0.518755in}}%
\pgfpathlineto{\pgfqpoint{1.996246in}{0.512029in}}%
\pgfpathlineto{\pgfqpoint{2.008486in}{0.500577in}}%
\pgfpathlineto{\pgfqpoint{1.989179in}{0.490629in}}%
\pgfpathlineto{\pgfqpoint{1.968354in}{0.475608in}}%
\pgfpathlineto{\pgfqpoint{1.963353in}{0.482590in}}%
\pgfpathlineto{\pgfqpoint{1.946310in}{0.472125in}}%
\pgfpathlineto{\pgfqpoint{1.963014in}{0.468151in}}%
\pgfpathlineto{\pgfqpoint{1.950420in}{0.451696in}}%
\pgfpathlineto{\pgfqpoint{1.944272in}{0.456593in}}%
\pgfpathlineto{\pgfqpoint{1.936015in}{0.448727in}}%
\pgfpathlineto{\pgfqpoint{1.946306in}{0.441857in}}%
\pgfpathlineto{\pgfqpoint{1.940221in}{0.431815in}}%
\pgfpathlineto{\pgfqpoint{1.932743in}{0.408043in}}%
\pgfpathlineto{\pgfqpoint{1.921956in}{0.385147in}}%
\pgfpathlineto{\pgfqpoint{1.926799in}{0.370139in}}%
\pgfpathlineto{\pgfqpoint{1.935030in}{0.334756in}}%
\pgfpathlineto{\pgfqpoint{1.935771in}{0.320436in}}%
\pgfpathlineto{\pgfqpoint{1.948489in}{0.301654in}}%
\pgfpathlineto{\pgfqpoint{1.938682in}{0.302751in}}%
\pgfpathlineto{\pgfqpoint{1.929156in}{0.293276in}}%
\pgfpathlineto{\pgfqpoint{1.918000in}{0.300723in}}%
\pgfpathlineto{\pgfqpoint{1.913996in}{0.308152in}}%
\pgfpathlineto{\pgfqpoint{1.898118in}{0.311608in}}%
\pgfpathlineto{\pgfqpoint{1.873865in}{0.312032in}}%
\pgfpathlineto{\pgfqpoint{1.855994in}{0.326101in}}%
\pgfpathlineto{\pgfqpoint{1.839768in}{0.328451in}}%
\pgfpathlineto{\pgfqpoint{1.829925in}{0.339632in}}%
\pgfpathlineto{\pgfqpoint{1.809309in}{0.344132in}}%
\pgfpathlineto{\pgfqpoint{1.798036in}{0.380084in}}%
\pgfpathlineto{\pgfqpoint{1.786510in}{0.394486in}}%
\pgfpathlineto{\pgfqpoint{1.788468in}{0.408170in}}%
\pgfpathlineto{\pgfqpoint{1.781321in}{0.418171in}}%
\pgfpathlineto{\pgfqpoint{1.785806in}{0.431841in}}%
\pgfpathlineto{\pgfqpoint{1.782106in}{0.441858in}}%
\pgfpathlineto{\pgfqpoint{1.770524in}{0.446396in}}%
\pgfpathlineto{\pgfqpoint{1.759694in}{0.457947in}}%
\pgfpathlineto{\pgfqpoint{1.755748in}{0.473412in}}%
\pgfpathlineto{\pgfqpoint{1.745500in}{0.487466in}}%
\pgfpathlineto{\pgfqpoint{1.731859in}{0.498402in}}%
\pgfpathlineto{\pgfqpoint{1.719573in}{0.529770in}}%
\pgfpathlineto{\pgfqpoint{1.709652in}{0.564077in}}%
\pgfpathlineto{\pgfqpoint{1.701513in}{0.577641in}}%
\pgfpathlineto{\pgfqpoint{1.687401in}{0.589061in}}%
\pgfpathlineto{\pgfqpoint{1.683874in}{0.597350in}}%
\pgfpathlineto{\pgfqpoint{1.670661in}{0.602497in}}%
\pgfpathlineto{\pgfqpoint{1.657620in}{0.624460in}}%
\pgfpathlineto{\pgfqpoint{1.645710in}{0.622774in}}%
\pgfpathlineto{\pgfqpoint{1.633616in}{0.628259in}}%
\pgfpathlineto{\pgfqpoint{1.616474in}{0.627199in}}%
\pgfpathlineto{\pgfqpoint{1.599059in}{0.636251in}}%
\pgfpathlineto{\pgfqpoint{1.594148in}{0.627642in}}%
\pgfpathlineto{\pgfqpoint{1.573796in}{0.627427in}}%
\pgfpathlineto{\pgfqpoint{1.563435in}{0.611106in}}%
\pgfpathlineto{\pgfqpoint{1.554423in}{0.590899in}}%
\pgfpathlineto{\pgfqpoint{1.535259in}{0.569207in}}%
\pgfpathlineto{\pgfqpoint{1.513549in}{0.578727in}}%
\pgfpathlineto{\pgfqpoint{1.510468in}{0.585018in}}%
\pgfpathlineto{\pgfqpoint{1.497274in}{0.589816in}}%
\pgfpathlineto{\pgfqpoint{1.493221in}{0.596380in}}%
\pgfpathlineto{\pgfqpoint{1.475705in}{0.603005in}}%
\pgfpathlineto{\pgfqpoint{1.465907in}{0.616550in}}%
\pgfpathlineto{\pgfqpoint{1.454464in}{0.623084in}}%
\pgfpathlineto{\pgfqpoint{1.444618in}{0.634488in}}%
\pgfpathlineto{\pgfqpoint{1.436949in}{0.653800in}}%
\pgfpathlineto{\pgfqpoint{1.437807in}{0.680190in}}%
\pgfpathlineto{\pgfqpoint{1.428809in}{0.693526in}}%
\pgfpathlineto{\pgfqpoint{1.427774in}{0.707969in}}%
\pgfpathlineto{\pgfqpoint{1.407810in}{0.729594in}}%
\pgfpathlineto{\pgfqpoint{1.396194in}{0.734211in}}%
\pgfpathlineto{\pgfqpoint{1.383867in}{0.754389in}}%
\pgfpathlineto{\pgfqpoint{1.373364in}{0.762406in}}%
\pgfpathlineto{\pgfqpoint{1.360017in}{0.781928in}}%
\pgfpathlineto{\pgfqpoint{1.346361in}{0.790377in}}%
\pgfpathlineto{\pgfqpoint{1.337418in}{0.812026in}}%
\pgfpathlineto{\pgfqpoint{1.327465in}{0.817498in}}%
\pgfpathlineto{\pgfqpoint{1.327465in}{0.817498in}}%
\pgfusepath{stroke}%
\end{pgfscope}%
\begin{pgfscope}%
\pgfpathrectangle{\pgfqpoint{0.100000in}{0.100000in}}{\pgfqpoint{3.608454in}{2.310000in}}%
\pgfusepath{clip}%
\pgfsetbuttcap%
\pgfsetroundjoin%
\pgfsetlinewidth{0.010037pt}%
\definecolor{currentstroke}{rgb}{1.000000,1.000000,1.000000}%
\pgfsetstrokecolor{currentstroke}%
\pgfsetdash{}{0pt}%
\pgfpathmoveto{\pgfqpoint{1.222893in}{1.249723in}}%
\pgfpathlineto{\pgfqpoint{1.321411in}{1.235415in}}%
\pgfpathlineto{\pgfqpoint{1.359578in}{1.229535in}}%
\pgfpathlineto{\pgfqpoint{1.467164in}{1.216739in}}%
\pgfpathlineto{\pgfqpoint{1.539577in}{1.209185in}}%
\pgfpathlineto{\pgfqpoint{1.602243in}{1.203632in}}%
\pgfpathlineto{\pgfqpoint{1.598743in}{1.163746in}}%
\pgfpathlineto{\pgfqpoint{1.596844in}{1.163852in}}%
\pgfpathlineto{\pgfqpoint{1.591104in}{1.095394in}}%
\pgfpathlineto{\pgfqpoint{1.584946in}{1.026087in}}%
\pgfpathlineto{\pgfqpoint{1.577823in}{0.953532in}}%
\pgfpathlineto{\pgfqpoint{1.566842in}{0.849964in}}%
\pgfpathlineto{\pgfqpoint{1.563330in}{0.809245in}}%
\pgfpathlineto{\pgfqpoint{1.498773in}{0.815759in}}%
\pgfpathlineto{\pgfqpoint{1.442877in}{0.821252in}}%
\pgfpathlineto{\pgfqpoint{1.365572in}{0.830032in}}%
\pgfpathlineto{\pgfqpoint{1.323997in}{0.835189in}}%
\pgfpathlineto{\pgfqpoint{1.321550in}{0.827003in}}%
\pgfpathlineto{\pgfqpoint{1.327465in}{0.817498in}}%
\pgfpathlineto{\pgfqpoint{1.277499in}{0.823987in}}%
\pgfpathlineto{\pgfqpoint{1.215853in}{0.832835in}}%
\pgfpathlineto{\pgfqpoint{1.210243in}{0.797968in}}%
\pgfpathlineto{\pgfqpoint{1.154011in}{0.806432in}}%
\pgfpathlineto{\pgfqpoint{1.167224in}{0.892929in}}%
\pgfpathlineto{\pgfqpoint{1.183580in}{0.998244in}}%
\pgfpathlineto{\pgfqpoint{1.198122in}{1.089991in}}%
\pgfpathlineto{\pgfqpoint{1.210751in}{1.171542in}}%
\pgfpathlineto{\pgfqpoint{1.222893in}{1.249723in}}%
\pgfusepath{stroke}%
\end{pgfscope}%
\begin{pgfscope}%
\pgfpathrectangle{\pgfqpoint{0.100000in}{0.100000in}}{\pgfqpoint{3.608454in}{2.310000in}}%
\pgfusepath{clip}%
\pgfsetbuttcap%
\pgfsetroundjoin%
\pgfsetlinewidth{0.010037pt}%
\definecolor{currentstroke}{rgb}{1.000000,1.000000,1.000000}%
\pgfsetstrokecolor{currentstroke}%
\pgfsetdash{}{0pt}%
\pgfpathmoveto{\pgfqpoint{2.789338in}{0.742945in}}%
\pgfpathlineto{\pgfqpoint{2.695351in}{0.732506in}}%
\pgfpathlineto{\pgfqpoint{2.612503in}{0.726074in}}%
\pgfpathlineto{\pgfqpoint{2.611467in}{0.715921in}}%
\pgfpathlineto{\pgfqpoint{2.619077in}{0.706258in}}%
\pgfpathlineto{\pgfqpoint{2.628365in}{0.700581in}}%
\pgfpathlineto{\pgfqpoint{2.628224in}{0.685633in}}%
\pgfpathlineto{\pgfqpoint{2.625763in}{0.675646in}}%
\pgfpathlineto{\pgfqpoint{2.617565in}{0.668449in}}%
\pgfpathlineto{\pgfqpoint{2.606139in}{0.669211in}}%
\pgfpathlineto{\pgfqpoint{2.595336in}{0.678128in}}%
\pgfpathlineto{\pgfqpoint{2.593410in}{0.693997in}}%
\pgfpathlineto{\pgfqpoint{2.585364in}{0.703227in}}%
\pgfpathlineto{\pgfqpoint{2.579887in}{0.670195in}}%
\pgfpathlineto{\pgfqpoint{2.561278in}{0.673334in}}%
\pgfpathlineto{\pgfqpoint{2.554807in}{0.731095in}}%
\pgfpathlineto{\pgfqpoint{2.547798in}{0.791968in}}%
\pgfpathlineto{\pgfqpoint{2.550401in}{0.879794in}}%
\pgfpathlineto{\pgfqpoint{2.552017in}{0.940061in}}%
\pgfpathlineto{\pgfqpoint{2.555453in}{1.032015in}}%
\pgfpathlineto{\pgfqpoint{2.548098in}{1.040107in}}%
\pgfpathlineto{\pgfqpoint{2.561482in}{1.042047in}}%
\pgfpathlineto{\pgfqpoint{2.639455in}{1.047089in}}%
\pgfpathlineto{\pgfqpoint{2.715320in}{1.053912in}}%
\pgfpathlineto{\pgfqpoint{2.735280in}{0.983998in}}%
\pgfpathlineto{\pgfqpoint{2.748892in}{0.932681in}}%
\pgfpathlineto{\pgfqpoint{2.761093in}{0.889718in}}%
\pgfpathlineto{\pgfqpoint{2.768151in}{0.872419in}}%
\pgfpathlineto{\pgfqpoint{2.779261in}{0.856321in}}%
\pgfpathlineto{\pgfqpoint{2.777467in}{0.848128in}}%
\pgfpathlineto{\pgfqpoint{2.785404in}{0.844137in}}%
\pgfpathlineto{\pgfqpoint{2.776025in}{0.831718in}}%
\pgfpathlineto{\pgfqpoint{2.772856in}{0.809602in}}%
\pgfpathlineto{\pgfqpoint{2.782023in}{0.783735in}}%
\pgfpathlineto{\pgfqpoint{2.780750in}{0.757706in}}%
\pgfpathlineto{\pgfqpoint{2.789338in}{0.742945in}}%
\pgfusepath{stroke}%
\end{pgfscope}%
\begin{pgfscope}%
\pgfpathrectangle{\pgfqpoint{0.100000in}{0.100000in}}{\pgfqpoint{3.608454in}{2.310000in}}%
\pgfusepath{clip}%
\pgfsetbuttcap%
\pgfsetroundjoin%
\pgfsetlinewidth{0.010037pt}%
\definecolor{currentstroke}{rgb}{1.000000,1.000000,1.000000}%
\pgfsetstrokecolor{currentstroke}%
\pgfsetdash{}{0pt}%
\pgfpathmoveto{\pgfqpoint{2.561278in}{0.673334in}}%
\pgfpathlineto{\pgfqpoint{2.542181in}{0.667869in}}%
\pgfpathlineto{\pgfqpoint{2.527946in}{0.671378in}}%
\pgfpathlineto{\pgfqpoint{2.501511in}{0.662970in}}%
\pgfpathlineto{\pgfqpoint{2.481574in}{0.652140in}}%
\pgfpathlineto{\pgfqpoint{2.473394in}{0.671709in}}%
\pgfpathlineto{\pgfqpoint{2.464961in}{0.679911in}}%
\pgfpathlineto{\pgfqpoint{2.460893in}{0.689726in}}%
\pgfpathlineto{\pgfqpoint{2.466856in}{0.716291in}}%
\pgfpathlineto{\pgfqpoint{2.410356in}{0.712812in}}%
\pgfpathlineto{\pgfqpoint{2.337092in}{0.709636in}}%
\pgfpathlineto{\pgfqpoint{2.341451in}{0.716248in}}%
\pgfpathlineto{\pgfqpoint{2.336063in}{0.728737in}}%
\pgfpathlineto{\pgfqpoint{2.341432in}{0.731282in}}%
\pgfpathlineto{\pgfqpoint{2.340241in}{0.743282in}}%
\pgfpathlineto{\pgfqpoint{2.349598in}{0.754918in}}%
\pgfpathlineto{\pgfqpoint{2.354717in}{0.777508in}}%
\pgfpathlineto{\pgfqpoint{2.371843in}{0.792391in}}%
\pgfpathlineto{\pgfqpoint{2.365495in}{0.806845in}}%
\pgfpathlineto{\pgfqpoint{2.377639in}{0.815165in}}%
\pgfpathlineto{\pgfqpoint{2.367164in}{0.830227in}}%
\pgfpathlineto{\pgfqpoint{2.359582in}{0.864708in}}%
\pgfpathlineto{\pgfqpoint{2.362458in}{0.870969in}}%
\pgfpathlineto{\pgfqpoint{2.366997in}{0.882978in}}%
\pgfpathlineto{\pgfqpoint{2.362772in}{0.895582in}}%
\pgfpathlineto{\pgfqpoint{2.364838in}{0.901316in}}%
\pgfpathlineto{\pgfqpoint{2.357405in}{0.922979in}}%
\pgfpathlineto{\pgfqpoint{2.379349in}{0.961875in}}%
\pgfpathlineto{\pgfqpoint{2.380370in}{0.972228in}}%
\pgfpathlineto{\pgfqpoint{2.390941in}{0.979756in}}%
\pgfpathlineto{\pgfqpoint{2.402220in}{1.004607in}}%
\pgfpathlineto{\pgfqpoint{2.401684in}{1.014709in}}%
\pgfpathlineto{\pgfqpoint{2.415733in}{1.025035in}}%
\pgfpathlineto{\pgfqpoint{2.411349in}{1.031151in}}%
\pgfpathlineto{\pgfqpoint{2.495536in}{1.036064in}}%
\pgfpathlineto{\pgfqpoint{2.548098in}{1.040107in}}%
\pgfpathlineto{\pgfqpoint{2.555453in}{1.032015in}}%
\pgfpathlineto{\pgfqpoint{2.552017in}{0.940061in}}%
\pgfpathlineto{\pgfqpoint{2.550401in}{0.879794in}}%
\pgfpathlineto{\pgfqpoint{2.547798in}{0.791968in}}%
\pgfpathlineto{\pgfqpoint{2.554807in}{0.731095in}}%
\pgfpathlineto{\pgfqpoint{2.561278in}{0.673334in}}%
\pgfusepath{stroke}%
\end{pgfscope}%
\begin{pgfscope}%
\pgfpathrectangle{\pgfqpoint{0.100000in}{0.100000in}}{\pgfqpoint{3.608454in}{2.310000in}}%
\pgfusepath{clip}%
\pgfsetbuttcap%
\pgfsetroundjoin%
\pgfsetlinewidth{0.010037pt}%
\definecolor{currentstroke}{rgb}{1.000000,1.000000,1.000000}%
\pgfsetstrokecolor{currentstroke}%
\pgfsetdash{}{0pt}%
\pgfpathmoveto{\pgfqpoint{2.715320in}{1.053912in}}%
\pgfpathlineto{\pgfqpoint{2.798250in}{1.063020in}}%
\pgfpathlineto{\pgfqpoint{2.848149in}{1.069168in}}%
\pgfpathlineto{\pgfqpoint{2.876568in}{1.073696in}}%
\pgfpathlineto{\pgfqpoint{2.864796in}{1.055128in}}%
\pgfpathlineto{\pgfqpoint{2.864803in}{1.046374in}}%
\pgfpathlineto{\pgfqpoint{2.876921in}{1.041653in}}%
\pgfpathlineto{\pgfqpoint{2.885162in}{1.034031in}}%
\pgfpathlineto{\pgfqpoint{2.897615in}{1.033089in}}%
\pgfpathlineto{\pgfqpoint{2.909255in}{1.011633in}}%
\pgfpathlineto{\pgfqpoint{2.921882in}{0.996503in}}%
\pgfpathlineto{\pgfqpoint{2.944306in}{0.983786in}}%
\pgfpathlineto{\pgfqpoint{2.950679in}{0.973643in}}%
\pgfpathlineto{\pgfqpoint{2.970676in}{0.962674in}}%
\pgfpathlineto{\pgfqpoint{2.972180in}{0.954719in}}%
\pgfpathlineto{\pgfqpoint{2.988402in}{0.938901in}}%
\pgfpathlineto{\pgfqpoint{3.000827in}{0.934397in}}%
\pgfpathlineto{\pgfqpoint{3.009626in}{0.919461in}}%
\pgfpathlineto{\pgfqpoint{3.013522in}{0.902673in}}%
\pgfpathlineto{\pgfqpoint{3.026428in}{0.896184in}}%
\pgfpathlineto{\pgfqpoint{3.036819in}{0.878138in}}%
\pgfpathlineto{\pgfqpoint{3.040169in}{0.864866in}}%
\pgfpathlineto{\pgfqpoint{3.055395in}{0.859218in}}%
\pgfpathlineto{\pgfqpoint{3.042619in}{0.834754in}}%
\pgfpathlineto{\pgfqpoint{3.037811in}{0.820297in}}%
\pgfpathlineto{\pgfqpoint{3.040948in}{0.813520in}}%
\pgfpathlineto{\pgfqpoint{3.035470in}{0.802258in}}%
\pgfpathlineto{\pgfqpoint{3.033142in}{0.786687in}}%
\pgfpathlineto{\pgfqpoint{3.023401in}{0.783787in}}%
\pgfpathlineto{\pgfqpoint{3.028545in}{0.769530in}}%
\pgfpathlineto{\pgfqpoint{3.028216in}{0.751443in}}%
\pgfpathlineto{\pgfqpoint{3.010902in}{0.752121in}}%
\pgfpathlineto{\pgfqpoint{2.998840in}{0.755421in}}%
\pgfpathlineto{\pgfqpoint{2.993620in}{0.749287in}}%
\pgfpathlineto{\pgfqpoint{2.997515in}{0.734795in}}%
\pgfpathlineto{\pgfqpoint{2.996668in}{0.717946in}}%
\pgfpathlineto{\pgfqpoint{2.989059in}{0.716657in}}%
\pgfpathlineto{\pgfqpoint{2.982903in}{0.732393in}}%
\pgfpathlineto{\pgfqpoint{2.920222in}{0.728215in}}%
\pgfpathlineto{\pgfqpoint{2.841186in}{0.723903in}}%
\pgfpathlineto{\pgfqpoint{2.801326in}{0.721097in}}%
\pgfpathlineto{\pgfqpoint{2.789338in}{0.742945in}}%
\pgfpathlineto{\pgfqpoint{2.780750in}{0.757706in}}%
\pgfpathlineto{\pgfqpoint{2.782023in}{0.783735in}}%
\pgfpathlineto{\pgfqpoint{2.772856in}{0.809602in}}%
\pgfpathlineto{\pgfqpoint{2.776025in}{0.831718in}}%
\pgfpathlineto{\pgfqpoint{2.785404in}{0.844137in}}%
\pgfpathlineto{\pgfqpoint{2.777467in}{0.848128in}}%
\pgfpathlineto{\pgfqpoint{2.779261in}{0.856321in}}%
\pgfpathlineto{\pgfqpoint{2.768151in}{0.872419in}}%
\pgfpathlineto{\pgfqpoint{2.761093in}{0.889718in}}%
\pgfpathlineto{\pgfqpoint{2.748892in}{0.932681in}}%
\pgfpathlineto{\pgfqpoint{2.735280in}{0.983998in}}%
\pgfpathlineto{\pgfqpoint{2.715320in}{1.053912in}}%
\pgfusepath{stroke}%
\end{pgfscope}%
\begin{pgfscope}%
\pgfpathrectangle{\pgfqpoint{0.100000in}{0.100000in}}{\pgfqpoint{3.608454in}{2.310000in}}%
\pgfusepath{clip}%
\pgfsetbuttcap%
\pgfsetroundjoin%
\pgfsetlinewidth{0.010037pt}%
\definecolor{currentstroke}{rgb}{1.000000,1.000000,1.000000}%
\pgfsetstrokecolor{currentstroke}%
\pgfsetdash{}{0pt}%
\pgfpathmoveto{\pgfqpoint{2.876568in}{1.073696in}}%
\pgfpathlineto{\pgfqpoint{2.909453in}{1.090054in}}%
\pgfpathlineto{\pgfqpoint{2.927665in}{1.096231in}}%
\pgfpathlineto{\pgfqpoint{3.006902in}{1.104464in}}%
\pgfpathlineto{\pgfqpoint{3.015293in}{1.101737in}}%
\pgfpathlineto{\pgfqpoint{3.026378in}{1.090595in}}%
\pgfpathlineto{\pgfqpoint{3.026989in}{1.080701in}}%
\pgfpathlineto{\pgfqpoint{3.099917in}{1.091416in}}%
\pgfpathlineto{\pgfqpoint{3.182729in}{1.032053in}}%
\pgfpathlineto{\pgfqpoint{3.167199in}{1.015867in}}%
\pgfpathlineto{\pgfqpoint{3.145890in}{0.978275in}}%
\pgfpathlineto{\pgfqpoint{3.151955in}{0.970196in}}%
\pgfpathlineto{\pgfqpoint{3.140588in}{0.954516in}}%
\pgfpathlineto{\pgfqpoint{3.129280in}{0.952741in}}%
\pgfpathlineto{\pgfqpoint{3.129238in}{0.943317in}}%
\pgfpathlineto{\pgfqpoint{3.121059in}{0.933457in}}%
\pgfpathlineto{\pgfqpoint{3.110879in}{0.931458in}}%
\pgfpathlineto{\pgfqpoint{3.113100in}{0.922697in}}%
\pgfpathlineto{\pgfqpoint{3.107420in}{0.915971in}}%
\pgfpathlineto{\pgfqpoint{3.093832in}{0.910161in}}%
\pgfpathlineto{\pgfqpoint{3.077270in}{0.896945in}}%
\pgfpathlineto{\pgfqpoint{3.080396in}{0.888394in}}%
\pgfpathlineto{\pgfqpoint{3.069973in}{0.883025in}}%
\pgfpathlineto{\pgfqpoint{3.061181in}{0.888670in}}%
\pgfpathlineto{\pgfqpoint{3.054753in}{0.864128in}}%
\pgfpathlineto{\pgfqpoint{3.040169in}{0.864866in}}%
\pgfpathlineto{\pgfqpoint{3.036819in}{0.878138in}}%
\pgfpathlineto{\pgfqpoint{3.026428in}{0.896184in}}%
\pgfpathlineto{\pgfqpoint{3.013522in}{0.902673in}}%
\pgfpathlineto{\pgfqpoint{3.009626in}{0.919461in}}%
\pgfpathlineto{\pgfqpoint{3.000827in}{0.934397in}}%
\pgfpathlineto{\pgfqpoint{2.988402in}{0.938901in}}%
\pgfpathlineto{\pgfqpoint{2.972180in}{0.954719in}}%
\pgfpathlineto{\pgfqpoint{2.970676in}{0.962674in}}%
\pgfpathlineto{\pgfqpoint{2.950679in}{0.973643in}}%
\pgfpathlineto{\pgfqpoint{2.944306in}{0.983786in}}%
\pgfpathlineto{\pgfqpoint{2.921882in}{0.996503in}}%
\pgfpathlineto{\pgfqpoint{2.909255in}{1.011633in}}%
\pgfpathlineto{\pgfqpoint{2.897615in}{1.033089in}}%
\pgfpathlineto{\pgfqpoint{2.885162in}{1.034031in}}%
\pgfpathlineto{\pgfqpoint{2.876921in}{1.041653in}}%
\pgfpathlineto{\pgfqpoint{2.864803in}{1.046374in}}%
\pgfpathlineto{\pgfqpoint{2.864796in}{1.055128in}}%
\pgfpathlineto{\pgfqpoint{2.876568in}{1.073696in}}%
\pgfusepath{stroke}%
\end{pgfscope}%
\begin{pgfscope}%
\pgfpathrectangle{\pgfqpoint{0.100000in}{0.100000in}}{\pgfqpoint{3.608454in}{2.310000in}}%
\pgfusepath{clip}%
\pgfsetbuttcap%
\pgfsetroundjoin%
\pgfsetlinewidth{0.010037pt}%
\definecolor{currentstroke}{rgb}{1.000000,1.000000,1.000000}%
\pgfsetstrokecolor{currentstroke}%
\pgfsetdash{}{0pt}%
\pgfpathmoveto{\pgfqpoint{2.131124in}{1.141195in}}%
\pgfpathlineto{\pgfqpoint{2.179448in}{1.141630in}}%
\pgfpathlineto{\pgfqpoint{2.243380in}{1.142795in}}%
\pgfpathlineto{\pgfqpoint{2.352626in}{1.146263in}}%
\pgfpathlineto{\pgfqpoint{2.415111in}{1.149567in}}%
\pgfpathlineto{\pgfqpoint{2.421813in}{1.141271in}}%
\pgfpathlineto{\pgfqpoint{2.421394in}{1.132509in}}%
\pgfpathlineto{\pgfqpoint{2.406257in}{1.117394in}}%
\pgfpathlineto{\pgfqpoint{2.402597in}{1.109117in}}%
\pgfpathlineto{\pgfqpoint{2.444579in}{1.112222in}}%
\pgfpathlineto{\pgfqpoint{2.444556in}{1.096930in}}%
\pgfpathlineto{\pgfqpoint{2.431088in}{1.090395in}}%
\pgfpathlineto{\pgfqpoint{2.431278in}{1.080004in}}%
\pgfpathlineto{\pgfqpoint{2.426220in}{1.072535in}}%
\pgfpathlineto{\pgfqpoint{2.422910in}{1.056671in}}%
\pgfpathlineto{\pgfqpoint{2.426327in}{1.043584in}}%
\pgfpathlineto{\pgfqpoint{2.411349in}{1.031151in}}%
\pgfpathlineto{\pgfqpoint{2.415733in}{1.025035in}}%
\pgfpathlineto{\pgfqpoint{2.401684in}{1.014709in}}%
\pgfpathlineto{\pgfqpoint{2.402220in}{1.004607in}}%
\pgfpathlineto{\pgfqpoint{2.390941in}{0.979756in}}%
\pgfpathlineto{\pgfqpoint{2.380370in}{0.972228in}}%
\pgfpathlineto{\pgfqpoint{2.379349in}{0.961875in}}%
\pgfpathlineto{\pgfqpoint{2.357405in}{0.922979in}}%
\pgfpathlineto{\pgfqpoint{2.364838in}{0.901316in}}%
\pgfpathlineto{\pgfqpoint{2.362772in}{0.895582in}}%
\pgfpathlineto{\pgfqpoint{2.366997in}{0.882978in}}%
\pgfpathlineto{\pgfqpoint{2.362458in}{0.870969in}}%
\pgfpathlineto{\pgfqpoint{2.302450in}{0.868492in}}%
\pgfpathlineto{\pgfqpoint{2.224548in}{0.867208in}}%
\pgfpathlineto{\pgfqpoint{2.170816in}{0.866723in}}%
\pgfpathlineto{\pgfqpoint{2.170554in}{0.909031in}}%
\pgfpathlineto{\pgfqpoint{2.157210in}{0.911803in}}%
\pgfpathlineto{\pgfqpoint{2.148404in}{0.908183in}}%
\pgfpathlineto{\pgfqpoint{2.141369in}{0.914830in}}%
\pgfpathlineto{\pgfqpoint{2.142095in}{0.959646in}}%
\pgfpathlineto{\pgfqpoint{2.143669in}{1.055067in}}%
\pgfpathlineto{\pgfqpoint{2.136037in}{1.110943in}}%
\pgfpathlineto{\pgfqpoint{2.131124in}{1.141195in}}%
\pgfusepath{stroke}%
\end{pgfscope}%
\begin{pgfscope}%
\pgfpathrectangle{\pgfqpoint{0.100000in}{0.100000in}}{\pgfqpoint{3.608454in}{2.310000in}}%
\pgfusepath{clip}%
\pgfsetbuttcap%
\pgfsetroundjoin%
\pgfsetlinewidth{0.010037pt}%
\definecolor{currentstroke}{rgb}{1.000000,1.000000,1.000000}%
\pgfsetstrokecolor{currentstroke}%
\pgfsetdash{}{0pt}%
\pgfpathmoveto{\pgfqpoint{2.170816in}{0.866723in}}%
\pgfpathlineto{\pgfqpoint{2.224548in}{0.867208in}}%
\pgfpathlineto{\pgfqpoint{2.302450in}{0.868492in}}%
\pgfpathlineto{\pgfqpoint{2.362458in}{0.870969in}}%
\pgfpathlineto{\pgfqpoint{2.359582in}{0.864708in}}%
\pgfpathlineto{\pgfqpoint{2.367164in}{0.830227in}}%
\pgfpathlineto{\pgfqpoint{2.377639in}{0.815165in}}%
\pgfpathlineto{\pgfqpoint{2.365495in}{0.806845in}}%
\pgfpathlineto{\pgfqpoint{2.371843in}{0.792391in}}%
\pgfpathlineto{\pgfqpoint{2.354717in}{0.777508in}}%
\pgfpathlineto{\pgfqpoint{2.349598in}{0.754918in}}%
\pgfpathlineto{\pgfqpoint{2.340241in}{0.743282in}}%
\pgfpathlineto{\pgfqpoint{2.341432in}{0.731282in}}%
\pgfpathlineto{\pgfqpoint{2.336063in}{0.728737in}}%
\pgfpathlineto{\pgfqpoint{2.341451in}{0.716248in}}%
\pgfpathlineto{\pgfqpoint{2.337092in}{0.709636in}}%
\pgfpathlineto{\pgfqpoint{2.410356in}{0.712812in}}%
\pgfpathlineto{\pgfqpoint{2.466856in}{0.716291in}}%
\pgfpathlineto{\pgfqpoint{2.460893in}{0.689726in}}%
\pgfpathlineto{\pgfqpoint{2.464961in}{0.679911in}}%
\pgfpathlineto{\pgfqpoint{2.473394in}{0.671709in}}%
\pgfpathlineto{\pgfqpoint{2.481574in}{0.652140in}}%
\pgfpathlineto{\pgfqpoint{2.455735in}{0.656649in}}%
\pgfpathlineto{\pgfqpoint{2.446204in}{0.664060in}}%
\pgfpathlineto{\pgfqpoint{2.434847in}{0.664413in}}%
\pgfpathlineto{\pgfqpoint{2.422908in}{0.648183in}}%
\pgfpathlineto{\pgfqpoint{2.425288in}{0.640793in}}%
\pgfpathlineto{\pgfqpoint{2.463429in}{0.636322in}}%
\pgfpathlineto{\pgfqpoint{2.473414in}{0.627827in}}%
\pgfpathlineto{\pgfqpoint{2.489663in}{0.623456in}}%
\pgfpathlineto{\pgfqpoint{2.482567in}{0.613398in}}%
\pgfpathlineto{\pgfqpoint{2.470611in}{0.604627in}}%
\pgfpathlineto{\pgfqpoint{2.502169in}{0.584902in}}%
\pgfpathlineto{\pgfqpoint{2.516869in}{0.581820in}}%
\pgfpathlineto{\pgfqpoint{2.522982in}{0.565696in}}%
\pgfpathlineto{\pgfqpoint{2.509142in}{0.562730in}}%
\pgfpathlineto{\pgfqpoint{2.483128in}{0.583009in}}%
\pgfpathlineto{\pgfqpoint{2.472958in}{0.585828in}}%
\pgfpathlineto{\pgfqpoint{2.468036in}{0.593839in}}%
\pgfpathlineto{\pgfqpoint{2.452950in}{0.590551in}}%
\pgfpathlineto{\pgfqpoint{2.448226in}{0.580228in}}%
\pgfpathlineto{\pgfqpoint{2.451218in}{0.568725in}}%
\pgfpathlineto{\pgfqpoint{2.441061in}{0.561938in}}%
\pgfpathlineto{\pgfqpoint{2.431792in}{0.578660in}}%
\pgfpathlineto{\pgfqpoint{2.415557in}{0.573674in}}%
\pgfpathlineto{\pgfqpoint{2.409474in}{0.563696in}}%
\pgfpathlineto{\pgfqpoint{2.397944in}{0.566534in}}%
\pgfpathlineto{\pgfqpoint{2.390573in}{0.579135in}}%
\pgfpathlineto{\pgfqpoint{2.371016in}{0.583295in}}%
\pgfpathlineto{\pgfqpoint{2.367323in}{0.589852in}}%
\pgfpathlineto{\pgfqpoint{2.346947in}{0.601273in}}%
\pgfpathlineto{\pgfqpoint{2.341889in}{0.611259in}}%
\pgfpathlineto{\pgfqpoint{2.324822in}{0.607178in}}%
\pgfpathlineto{\pgfqpoint{2.326990in}{0.616320in}}%
\pgfpathlineto{\pgfqpoint{2.305777in}{0.606942in}}%
\pgfpathlineto{\pgfqpoint{2.311475in}{0.597211in}}%
\pgfpathlineto{\pgfqpoint{2.295094in}{0.591458in}}%
\pgfpathlineto{\pgfqpoint{2.273396in}{0.594620in}}%
\pgfpathlineto{\pgfqpoint{2.229480in}{0.609663in}}%
\pgfpathlineto{\pgfqpoint{2.195594in}{0.606677in}}%
\pgfpathlineto{\pgfqpoint{2.192657in}{0.626436in}}%
\pgfpathlineto{\pgfqpoint{2.196912in}{0.635452in}}%
\pgfpathlineto{\pgfqpoint{2.196548in}{0.649798in}}%
\pgfpathlineto{\pgfqpoint{2.192362in}{0.653216in}}%
\pgfpathlineto{\pgfqpoint{2.193781in}{0.669590in}}%
\pgfpathlineto{\pgfqpoint{2.206122in}{0.692365in}}%
\pgfpathlineto{\pgfqpoint{2.207701in}{0.700970in}}%
\pgfpathlineto{\pgfqpoint{2.205667in}{0.721296in}}%
\pgfpathlineto{\pgfqpoint{2.196588in}{0.730869in}}%
\pgfpathlineto{\pgfqpoint{2.191847in}{0.748745in}}%
\pgfpathlineto{\pgfqpoint{2.186014in}{0.752849in}}%
\pgfpathlineto{\pgfqpoint{2.188625in}{0.762543in}}%
\pgfpathlineto{\pgfqpoint{2.181197in}{0.777032in}}%
\pgfpathlineto{\pgfqpoint{2.171929in}{0.784895in}}%
\pgfpathlineto{\pgfqpoint{2.170816in}{0.866723in}}%
\pgfusepath{stroke}%
\end{pgfscope}%
\begin{pgfscope}%
\pgfpathrectangle{\pgfqpoint{0.100000in}{0.100000in}}{\pgfqpoint{3.608454in}{2.310000in}}%
\pgfusepath{clip}%
\pgfsetbuttcap%
\pgfsetroundjoin%
\pgfsetlinewidth{0.010037pt}%
\definecolor{currentstroke}{rgb}{1.000000,1.000000,1.000000}%
\pgfsetstrokecolor{currentstroke}%
\pgfsetdash{}{0pt}%
\pgfpathmoveto{\pgfqpoint{2.314523in}{0.596466in}}%
\pgfpathlineto{\pgfqpoint{2.322296in}{0.601090in}}%
\pgfpathlineto{\pgfqpoint{2.331731in}{0.595623in}}%
\pgfpathlineto{\pgfqpoint{2.326469in}{0.588102in}}%
\pgfpathlineto{\pgfqpoint{2.314523in}{0.596466in}}%
\pgfusepath{stroke}%
\end{pgfscope}%
\begin{pgfscope}%
\pgfpathrectangle{\pgfqpoint{0.100000in}{0.100000in}}{\pgfqpoint{3.608454in}{2.310000in}}%
\pgfusepath{clip}%
\pgfsetbuttcap%
\pgfsetroundjoin%
\pgfsetlinewidth{0.010037pt}%
\definecolor{currentstroke}{rgb}{1.000000,1.000000,1.000000}%
\pgfsetstrokecolor{currentstroke}%
\pgfsetdash{}{0pt}%
\pgfpathmoveto{\pgfqpoint{2.628224in}{0.685633in}}%
\pgfpathlineto{\pgfqpoint{2.628365in}{0.700581in}}%
\pgfpathlineto{\pgfqpoint{2.619077in}{0.706258in}}%
\pgfpathlineto{\pgfqpoint{2.611467in}{0.715921in}}%
\pgfpathlineto{\pgfqpoint{2.612503in}{0.726074in}}%
\pgfpathlineto{\pgfqpoint{2.695351in}{0.732506in}}%
\pgfpathlineto{\pgfqpoint{2.789338in}{0.742945in}}%
\pgfpathlineto{\pgfqpoint{2.801326in}{0.721097in}}%
\pgfpathlineto{\pgfqpoint{2.841186in}{0.723903in}}%
\pgfpathlineto{\pgfqpoint{2.920222in}{0.728215in}}%
\pgfpathlineto{\pgfqpoint{2.982903in}{0.732393in}}%
\pgfpathlineto{\pgfqpoint{2.989059in}{0.716657in}}%
\pgfpathlineto{\pgfqpoint{2.996668in}{0.717946in}}%
\pgfpathlineto{\pgfqpoint{2.997515in}{0.734795in}}%
\pgfpathlineto{\pgfqpoint{2.993620in}{0.749287in}}%
\pgfpathlineto{\pgfqpoint{2.998840in}{0.755421in}}%
\pgfpathlineto{\pgfqpoint{3.010902in}{0.752121in}}%
\pgfpathlineto{\pgfqpoint{3.028216in}{0.751443in}}%
\pgfpathlineto{\pgfqpoint{3.030878in}{0.738496in}}%
\pgfpathlineto{\pgfqpoint{3.036204in}{0.731098in}}%
\pgfpathlineto{\pgfqpoint{3.040364in}{0.714912in}}%
\pgfpathlineto{\pgfqpoint{3.053250in}{0.689862in}}%
\pgfpathlineto{\pgfqpoint{3.053315in}{0.683088in}}%
\pgfpathlineto{\pgfqpoint{3.072361in}{0.653675in}}%
\pgfpathlineto{\pgfqpoint{3.074204in}{0.647599in}}%
\pgfpathlineto{\pgfqpoint{3.102776in}{0.606058in}}%
\pgfpathlineto{\pgfqpoint{3.098249in}{0.605387in}}%
\pgfpathlineto{\pgfqpoint{3.110295in}{0.575840in}}%
\pgfpathlineto{\pgfqpoint{3.143385in}{0.524475in}}%
\pgfpathlineto{\pgfqpoint{3.162598in}{0.486707in}}%
\pgfpathlineto{\pgfqpoint{3.169019in}{0.480095in}}%
\pgfpathlineto{\pgfqpoint{3.180022in}{0.456389in}}%
\pgfpathlineto{\pgfqpoint{3.183839in}{0.418430in}}%
\pgfpathlineto{\pgfqpoint{3.185359in}{0.390048in}}%
\pgfpathlineto{\pgfqpoint{3.183524in}{0.371875in}}%
\pgfpathlineto{\pgfqpoint{3.177624in}{0.358884in}}%
\pgfpathlineto{\pgfqpoint{3.180365in}{0.341848in}}%
\pgfpathlineto{\pgfqpoint{3.174013in}{0.328363in}}%
\pgfpathlineto{\pgfqpoint{3.155154in}{0.317306in}}%
\pgfpathlineto{\pgfqpoint{3.142888in}{0.318113in}}%
\pgfpathlineto{\pgfqpoint{3.134923in}{0.312333in}}%
\pgfpathlineto{\pgfqpoint{3.132641in}{0.327766in}}%
\pgfpathlineto{\pgfqpoint{3.119456in}{0.331840in}}%
\pgfpathlineto{\pgfqpoint{3.107637in}{0.353374in}}%
\pgfpathlineto{\pgfqpoint{3.106303in}{0.363180in}}%
\pgfpathlineto{\pgfqpoint{3.085149in}{0.369136in}}%
\pgfpathlineto{\pgfqpoint{3.071519in}{0.367846in}}%
\pgfpathlineto{\pgfqpoint{3.063806in}{0.382075in}}%
\pgfpathlineto{\pgfqpoint{3.054976in}{0.407650in}}%
\pgfpathlineto{\pgfqpoint{3.042722in}{0.412834in}}%
\pgfpathlineto{\pgfqpoint{3.036079in}{0.427490in}}%
\pgfpathlineto{\pgfqpoint{3.019832in}{0.436127in}}%
\pgfpathlineto{\pgfqpoint{3.010441in}{0.446931in}}%
\pgfpathlineto{\pgfqpoint{2.994241in}{0.476304in}}%
\pgfpathlineto{\pgfqpoint{2.992335in}{0.496978in}}%
\pgfpathlineto{\pgfqpoint{3.001163in}{0.510205in}}%
\pgfpathlineto{\pgfqpoint{3.000281in}{0.519397in}}%
\pgfpathlineto{\pgfqpoint{2.990079in}{0.520379in}}%
\pgfpathlineto{\pgfqpoint{2.981556in}{0.526772in}}%
\pgfpathlineto{\pgfqpoint{2.976889in}{0.518966in}}%
\pgfpathlineto{\pgfqpoint{2.985069in}{0.512605in}}%
\pgfpathlineto{\pgfqpoint{2.979115in}{0.501166in}}%
\pgfpathlineto{\pgfqpoint{2.969483in}{0.510660in}}%
\pgfpathlineto{\pgfqpoint{2.970600in}{0.536987in}}%
\pgfpathlineto{\pgfqpoint{2.975246in}{0.558347in}}%
\pgfpathlineto{\pgfqpoint{2.972858in}{0.594997in}}%
\pgfpathlineto{\pgfqpoint{2.963257in}{0.603726in}}%
\pgfpathlineto{\pgfqpoint{2.958429in}{0.614928in}}%
\pgfpathlineto{\pgfqpoint{2.941890in}{0.614679in}}%
\pgfpathlineto{\pgfqpoint{2.925541in}{0.633111in}}%
\pgfpathlineto{\pgfqpoint{2.914580in}{0.638635in}}%
\pgfpathlineto{\pgfqpoint{2.911329in}{0.650308in}}%
\pgfpathlineto{\pgfqpoint{2.900608in}{0.654407in}}%
\pgfpathlineto{\pgfqpoint{2.891722in}{0.667295in}}%
\pgfpathlineto{\pgfqpoint{2.868250in}{0.677815in}}%
\pgfpathlineto{\pgfqpoint{2.850001in}{0.678084in}}%
\pgfpathlineto{\pgfqpoint{2.842057in}{0.674044in}}%
\pgfpathlineto{\pgfqpoint{2.844031in}{0.661431in}}%
\pgfpathlineto{\pgfqpoint{2.835747in}{0.662028in}}%
\pgfpathlineto{\pgfqpoint{2.810265in}{0.644393in}}%
\pgfpathlineto{\pgfqpoint{2.779657in}{0.637403in}}%
\pgfpathlineto{\pgfqpoint{2.779155in}{0.646052in}}%
\pgfpathlineto{\pgfqpoint{2.772372in}{0.654504in}}%
\pgfpathlineto{\pgfqpoint{2.754172in}{0.666173in}}%
\pgfpathlineto{\pgfqpoint{2.728049in}{0.678137in}}%
\pgfpathlineto{\pgfqpoint{2.699674in}{0.684471in}}%
\pgfpathlineto{\pgfqpoint{2.694348in}{0.693083in}}%
\pgfpathlineto{\pgfqpoint{2.684117in}{0.685900in}}%
\pgfpathlineto{\pgfqpoint{2.671797in}{0.684316in}}%
\pgfpathlineto{\pgfqpoint{2.628914in}{0.673017in}}%
\pgfpathlineto{\pgfqpoint{2.628224in}{0.685633in}}%
\pgfusepath{stroke}%
\end{pgfscope}%
\begin{pgfscope}%
\pgfpathrectangle{\pgfqpoint{0.100000in}{0.100000in}}{\pgfqpoint{3.608454in}{2.310000in}}%
\pgfusepath{clip}%
\pgfsetbuttcap%
\pgfsetroundjoin%
\pgfsetlinewidth{0.010037pt}%
\definecolor{currentstroke}{rgb}{1.000000,1.000000,1.000000}%
\pgfsetstrokecolor{currentstroke}%
\pgfsetdash{}{0pt}%
\pgfpathmoveto{\pgfqpoint{3.106940in}{0.606649in}}%
\pgfpathlineto{\pgfqpoint{3.118035in}{0.593742in}}%
\pgfpathlineto{\pgfqpoint{3.105219in}{0.592930in}}%
\pgfpathlineto{\pgfqpoint{3.106940in}{0.606649in}}%
\pgfusepath{stroke}%
\end{pgfscope}%
\begin{pgfscope}%
\pgfpathrectangle{\pgfqpoint{0.100000in}{0.100000in}}{\pgfqpoint{3.608454in}{2.310000in}}%
\pgfusepath{clip}%
\pgfsetbuttcap%
\pgfsetroundjoin%
\pgfsetlinewidth{0.010037pt}%
\definecolor{currentstroke}{rgb}{1.000000,1.000000,1.000000}%
\pgfsetstrokecolor{currentstroke}%
\pgfsetdash{}{0pt}%
\pgfpathmoveto{\pgfqpoint{2.453791in}{2.080299in}}%
\pgfpathlineto{\pgfqpoint{2.447771in}{2.068582in}}%
\pgfpathlineto{\pgfqpoint{2.433410in}{2.061744in}}%
\pgfpathlineto{\pgfqpoint{2.427189in}{2.052541in}}%
\pgfpathlineto{\pgfqpoint{2.419903in}{2.059169in}}%
\pgfpathlineto{\pgfqpoint{2.453791in}{2.080299in}}%
\pgfusepath{stroke}%
\end{pgfscope}%
\begin{pgfscope}%
\pgfpathrectangle{\pgfqpoint{0.100000in}{0.100000in}}{\pgfqpoint{3.608454in}{2.310000in}}%
\pgfusepath{clip}%
\pgfsetbuttcap%
\pgfsetroundjoin%
\pgfsetlinewidth{0.010037pt}%
\definecolor{currentstroke}{rgb}{1.000000,1.000000,1.000000}%
\pgfsetstrokecolor{currentstroke}%
\pgfsetdash{}{0pt}%
\pgfpathmoveto{\pgfqpoint{2.458673in}{2.009787in}}%
\pgfpathlineto{\pgfqpoint{2.473368in}{2.023495in}}%
\pgfpathlineto{\pgfqpoint{2.496033in}{2.027084in}}%
\pgfpathlineto{\pgfqpoint{2.489790in}{2.017551in}}%
\pgfpathlineto{\pgfqpoint{2.474252in}{2.003764in}}%
\pgfpathlineto{\pgfqpoint{2.465168in}{1.986068in}}%
\pgfpathlineto{\pgfqpoint{2.461321in}{1.995666in}}%
\pgfpathlineto{\pgfqpoint{2.454489in}{1.997044in}}%
\pgfpathlineto{\pgfqpoint{2.453803in}{2.005709in}}%
\pgfpathlineto{\pgfqpoint{2.458673in}{2.009787in}}%
\pgfusepath{stroke}%
\end{pgfscope}%
\begin{pgfscope}%
\pgfpathrectangle{\pgfqpoint{0.100000in}{0.100000in}}{\pgfqpoint{3.608454in}{2.310000in}}%
\pgfusepath{clip}%
\pgfsetbuttcap%
\pgfsetroundjoin%
\pgfsetlinewidth{0.010037pt}%
\definecolor{currentstroke}{rgb}{1.000000,1.000000,1.000000}%
\pgfsetstrokecolor{currentstroke}%
\pgfsetdash{}{0pt}%
\pgfpathmoveto{\pgfqpoint{2.517200in}{1.842472in}}%
\pgfpathlineto{\pgfqpoint{2.513287in}{1.846815in}}%
\pgfpathlineto{\pgfqpoint{2.517401in}{1.859041in}}%
\pgfpathlineto{\pgfqpoint{2.505188in}{1.859820in}}%
\pgfpathlineto{\pgfqpoint{2.508426in}{1.870361in}}%
\pgfpathlineto{\pgfqpoint{2.506412in}{1.887150in}}%
\pgfpathlineto{\pgfqpoint{2.495454in}{1.892970in}}%
\pgfpathlineto{\pgfqpoint{2.483987in}{1.904750in}}%
\pgfpathlineto{\pgfqpoint{2.466317in}{1.908198in}}%
\pgfpathlineto{\pgfqpoint{2.449120in}{1.908100in}}%
\pgfpathlineto{\pgfqpoint{2.432207in}{1.916460in}}%
\pgfpathlineto{\pgfqpoint{2.375622in}{1.928645in}}%
\pgfpathlineto{\pgfqpoint{2.369439in}{1.941557in}}%
\pgfpathlineto{\pgfqpoint{2.358409in}{1.945963in}}%
\pgfpathlineto{\pgfqpoint{2.379243in}{1.955848in}}%
\pgfpathlineto{\pgfqpoint{2.391013in}{1.968181in}}%
\pgfpathlineto{\pgfqpoint{2.412935in}{1.971526in}}%
\pgfpathlineto{\pgfqpoint{2.426425in}{1.984088in}}%
\pgfpathlineto{\pgfqpoint{2.433502in}{1.984591in}}%
\pgfpathlineto{\pgfqpoint{2.438904in}{1.993550in}}%
\pgfpathlineto{\pgfqpoint{2.452070in}{2.004163in}}%
\pgfpathlineto{\pgfqpoint{2.453218in}{1.996662in}}%
\pgfpathlineto{\pgfqpoint{2.459151in}{1.995113in}}%
\pgfpathlineto{\pgfqpoint{2.463615in}{1.986167in}}%
\pgfpathlineto{\pgfqpoint{2.464419in}{1.970914in}}%
\pgfpathlineto{\pgfqpoint{2.477835in}{1.980029in}}%
\pgfpathlineto{\pgfqpoint{2.493468in}{1.981933in}}%
\pgfpathlineto{\pgfqpoint{2.506813in}{1.977158in}}%
\pgfpathlineto{\pgfqpoint{2.524910in}{1.952330in}}%
\pgfpathlineto{\pgfqpoint{2.544676in}{1.956298in}}%
\pgfpathlineto{\pgfqpoint{2.552707in}{1.949654in}}%
\pgfpathlineto{\pgfqpoint{2.565625in}{1.949060in}}%
\pgfpathlineto{\pgfqpoint{2.574207in}{1.960977in}}%
\pgfpathlineto{\pgfqpoint{2.590492in}{1.971511in}}%
\pgfpathlineto{\pgfqpoint{2.606143in}{1.974810in}}%
\pgfpathlineto{\pgfqpoint{2.625564in}{1.975153in}}%
\pgfpathlineto{\pgfqpoint{2.639756in}{1.983291in}}%
\pgfpathlineto{\pgfqpoint{2.651338in}{1.979543in}}%
\pgfpathlineto{\pgfqpoint{2.653805in}{1.962333in}}%
\pgfpathlineto{\pgfqpoint{2.678653in}{1.959638in}}%
\pgfpathlineto{\pgfqpoint{2.692149in}{1.967694in}}%
\pgfpathlineto{\pgfqpoint{2.701505in}{1.949546in}}%
\pgfpathlineto{\pgfqpoint{2.719357in}{1.928560in}}%
\pgfpathlineto{\pgfqpoint{2.694371in}{1.928682in}}%
\pgfpathlineto{\pgfqpoint{2.686477in}{1.926088in}}%
\pgfpathlineto{\pgfqpoint{2.675655in}{1.929459in}}%
\pgfpathlineto{\pgfqpoint{2.674927in}{1.914931in}}%
\pgfpathlineto{\pgfqpoint{2.655278in}{1.926284in}}%
\pgfpathlineto{\pgfqpoint{2.630009in}{1.929754in}}%
\pgfpathlineto{\pgfqpoint{2.623052in}{1.918697in}}%
\pgfpathlineto{\pgfqpoint{2.589991in}{1.913321in}}%
\pgfpathlineto{\pgfqpoint{2.586187in}{1.903954in}}%
\pgfpathlineto{\pgfqpoint{2.563180in}{1.901144in}}%
\pgfpathlineto{\pgfqpoint{2.556225in}{1.891622in}}%
\pgfpathlineto{\pgfqpoint{2.544055in}{1.889069in}}%
\pgfpathlineto{\pgfqpoint{2.534335in}{1.866491in}}%
\pgfpathlineto{\pgfqpoint{2.522029in}{1.844624in}}%
\pgfpathlineto{\pgfqpoint{2.517200in}{1.842472in}}%
\pgfusepath{stroke}%
\end{pgfscope}%
\begin{pgfscope}%
\pgfpathrectangle{\pgfqpoint{0.100000in}{0.100000in}}{\pgfqpoint{3.608454in}{2.310000in}}%
\pgfusepath{clip}%
\pgfsetbuttcap%
\pgfsetroundjoin%
\pgfsetlinewidth{0.010037pt}%
\definecolor{currentstroke}{rgb}{1.000000,1.000000,1.000000}%
\pgfsetstrokecolor{currentstroke}%
\pgfsetdash{}{0pt}%
\pgfpathmoveto{\pgfqpoint{2.588031in}{1.579945in}}%
\pgfpathlineto{\pgfqpoint{2.599764in}{1.592295in}}%
\pgfpathlineto{\pgfqpoint{2.605110in}{1.610201in}}%
\pgfpathlineto{\pgfqpoint{2.611467in}{1.620579in}}%
\pgfpathlineto{\pgfqpoint{2.615365in}{1.634703in}}%
\pgfpathlineto{\pgfqpoint{2.616576in}{1.662865in}}%
\pgfpathlineto{\pgfqpoint{2.610707in}{1.689854in}}%
\pgfpathlineto{\pgfqpoint{2.591312in}{1.731182in}}%
\pgfpathlineto{\pgfqpoint{2.596565in}{1.744177in}}%
\pgfpathlineto{\pgfqpoint{2.589723in}{1.762138in}}%
\pgfpathlineto{\pgfqpoint{2.601460in}{1.786979in}}%
\pgfpathlineto{\pgfqpoint{2.599547in}{1.814676in}}%
\pgfpathlineto{\pgfqpoint{2.607723in}{1.818162in}}%
\pgfpathlineto{\pgfqpoint{2.608750in}{1.831341in}}%
\pgfpathlineto{\pgfqpoint{2.623262in}{1.839799in}}%
\pgfpathlineto{\pgfqpoint{2.631469in}{1.838434in}}%
\pgfpathlineto{\pgfqpoint{2.633776in}{1.824300in}}%
\pgfpathlineto{\pgfqpoint{2.640182in}{1.823746in}}%
\pgfpathlineto{\pgfqpoint{2.646027in}{1.844133in}}%
\pgfpathlineto{\pgfqpoint{2.644040in}{1.859989in}}%
\pgfpathlineto{\pgfqpoint{2.647841in}{1.869093in}}%
\pgfpathlineto{\pgfqpoint{2.668387in}{1.878499in}}%
\pgfpathlineto{\pgfqpoint{2.659024in}{1.881847in}}%
\pgfpathlineto{\pgfqpoint{2.655974in}{1.889927in}}%
\pgfpathlineto{\pgfqpoint{2.662719in}{1.904115in}}%
\pgfpathlineto{\pgfqpoint{2.676020in}{1.909016in}}%
\pgfpathlineto{\pgfqpoint{2.691465in}{1.900645in}}%
\pgfpathlineto{\pgfqpoint{2.706027in}{1.900540in}}%
\pgfpathlineto{\pgfqpoint{2.712790in}{1.890761in}}%
\pgfpathlineto{\pgfqpoint{2.722988in}{1.891440in}}%
\pgfpathlineto{\pgfqpoint{2.754470in}{1.877846in}}%
\pgfpathlineto{\pgfqpoint{2.760759in}{1.864686in}}%
\pgfpathlineto{\pgfqpoint{2.753890in}{1.860128in}}%
\pgfpathlineto{\pgfqpoint{2.756007in}{1.850265in}}%
\pgfpathlineto{\pgfqpoint{2.762793in}{1.845879in}}%
\pgfpathlineto{\pgfqpoint{2.766575in}{1.833755in}}%
\pgfpathlineto{\pgfqpoint{2.766045in}{1.804218in}}%
\pgfpathlineto{\pgfqpoint{2.757086in}{1.797162in}}%
\pgfpathlineto{\pgfqpoint{2.755075in}{1.781646in}}%
\pgfpathlineto{\pgfqpoint{2.738454in}{1.767308in}}%
\pgfpathlineto{\pgfqpoint{2.739441in}{1.749952in}}%
\pgfpathlineto{\pgfqpoint{2.754008in}{1.743872in}}%
\pgfpathlineto{\pgfqpoint{2.770441in}{1.765527in}}%
\pgfpathlineto{\pgfqpoint{2.771800in}{1.773381in}}%
\pgfpathlineto{\pgfqpoint{2.792295in}{1.786435in}}%
\pgfpathlineto{\pgfqpoint{2.805328in}{1.780388in}}%
\pgfpathlineto{\pgfqpoint{2.813507in}{1.766723in}}%
\pgfpathlineto{\pgfqpoint{2.826701in}{1.719275in}}%
\pgfpathlineto{\pgfqpoint{2.833693in}{1.704263in}}%
\pgfpathlineto{\pgfqpoint{2.831504in}{1.698058in}}%
\pgfpathlineto{\pgfqpoint{2.831717in}{1.676927in}}%
\pgfpathlineto{\pgfqpoint{2.818998in}{1.678955in}}%
\pgfpathlineto{\pgfqpoint{2.811814in}{1.663165in}}%
\pgfpathlineto{\pgfqpoint{2.810819in}{1.652430in}}%
\pgfpathlineto{\pgfqpoint{2.801211in}{1.645539in}}%
\pgfpathlineto{\pgfqpoint{2.799076in}{1.624604in}}%
\pgfpathlineto{\pgfqpoint{2.785105in}{1.598157in}}%
\pgfpathlineto{\pgfqpoint{2.708677in}{1.586775in}}%
\pgfpathlineto{\pgfqpoint{2.708226in}{1.591778in}}%
\pgfpathlineto{\pgfqpoint{2.657102in}{1.586389in}}%
\pgfpathlineto{\pgfqpoint{2.588031in}{1.579945in}}%
\pgfusepath{stroke}%
\end{pgfscope}%
\begin{pgfscope}%
\pgfpathrectangle{\pgfqpoint{0.100000in}{0.100000in}}{\pgfqpoint{3.608454in}{2.310000in}}%
\pgfusepath{clip}%
\pgfsetbuttcap%
\pgfsetroundjoin%
\pgfsetlinewidth{0.010037pt}%
\definecolor{currentstroke}{rgb}{1.000000,1.000000,1.000000}%
\pgfsetstrokecolor{currentstroke}%
\pgfsetdash{}{0pt}%
\pgfpathmoveto{\pgfqpoint{0.000000in}{0.000000in}}%
\pgfusepath{stroke}%
\end{pgfscope}%
\begin{pgfscope}%
\pgfpathrectangle{\pgfqpoint{0.100000in}{0.100000in}}{\pgfqpoint{3.608454in}{2.310000in}}%
\pgfusepath{clip}%
\pgfsetbuttcap%
\pgfsetroundjoin%
\pgfsetlinewidth{0.010037pt}%
\definecolor{currentstroke}{rgb}{1.000000,1.000000,1.000000}%
\pgfsetstrokecolor{currentstroke}%
\pgfsetdash{}{0pt}%
\pgfusepath{stroke}%
\end{pgfscope}%
\begin{pgfscope}%
\pgfpathrectangle{\pgfqpoint{0.100000in}{0.100000in}}{\pgfqpoint{3.608454in}{2.310000in}}%
\pgfusepath{clip}%
\pgfsetbuttcap%
\pgfsetroundjoin%
\pgfsetlinewidth{0.010037pt}%
\definecolor{currentstroke}{rgb}{1.000000,1.000000,1.000000}%
\pgfsetstrokecolor{currentstroke}%
\pgfsetdash{}{0pt}%
\pgfusepath{stroke}%
\end{pgfscope}%
\begin{pgfscope}%
\pgfpathrectangle{\pgfqpoint{0.100000in}{0.100000in}}{\pgfqpoint{3.608454in}{2.310000in}}%
\pgfusepath{clip}%
\pgfsetbuttcap%
\pgfsetroundjoin%
\pgfsetlinewidth{0.010037pt}%
\definecolor{currentstroke}{rgb}{1.000000,1.000000,1.000000}%
\pgfsetstrokecolor{currentstroke}%
\pgfsetdash{}{0pt}%
\pgfusepath{stroke}%
\end{pgfscope}%
\begin{pgfscope}%
\pgfpathrectangle{\pgfqpoint{0.100000in}{0.100000in}}{\pgfqpoint{3.608454in}{2.310000in}}%
\pgfusepath{clip}%
\pgfsetbuttcap%
\pgfsetroundjoin%
\pgfsetlinewidth{0.010037pt}%
\definecolor{currentstroke}{rgb}{1.000000,1.000000,1.000000}%
\pgfsetstrokecolor{currentstroke}%
\pgfsetdash{}{0pt}%
\pgfusepath{stroke}%
\end{pgfscope}%
\begin{pgfscope}%
\pgfpathrectangle{\pgfqpoint{0.100000in}{0.100000in}}{\pgfqpoint{3.608454in}{2.310000in}}%
\pgfusepath{clip}%
\pgfsetbuttcap%
\pgfsetroundjoin%
\pgfsetlinewidth{0.010037pt}%
\definecolor{currentstroke}{rgb}{1.000000,1.000000,1.000000}%
\pgfsetstrokecolor{currentstroke}%
\pgfsetdash{}{0pt}%
\pgfusepath{stroke}%
\end{pgfscope}%
\begin{pgfscope}%
\pgfpathrectangle{\pgfqpoint{0.100000in}{0.100000in}}{\pgfqpoint{3.608454in}{2.310000in}}%
\pgfusepath{clip}%
\pgfsetbuttcap%
\pgfsetroundjoin%
\pgfsetlinewidth{0.010037pt}%
\definecolor{currentstroke}{rgb}{1.000000,1.000000,1.000000}%
\pgfsetstrokecolor{currentstroke}%
\pgfsetdash{}{0pt}%
\pgfusepath{stroke}%
\end{pgfscope}%
\begin{pgfscope}%
\pgfpathrectangle{\pgfqpoint{0.100000in}{0.100000in}}{\pgfqpoint{3.608454in}{2.310000in}}%
\pgfusepath{clip}%
\pgfsetbuttcap%
\pgfsetroundjoin%
\pgfsetlinewidth{0.010037pt}%
\definecolor{currentstroke}{rgb}{1.000000,1.000000,1.000000}%
\pgfsetstrokecolor{currentstroke}%
\pgfsetdash{}{0pt}%
\pgfusepath{stroke}%
\end{pgfscope}%
\begin{pgfscope}%
\pgfpathrectangle{\pgfqpoint{0.100000in}{0.100000in}}{\pgfqpoint{3.608454in}{2.310000in}}%
\pgfusepath{clip}%
\pgfsetbuttcap%
\pgfsetroundjoin%
\pgfsetlinewidth{0.010037pt}%
\definecolor{currentstroke}{rgb}{1.000000,1.000000,1.000000}%
\pgfsetstrokecolor{currentstroke}%
\pgfsetdash{}{0pt}%
\pgfusepath{stroke}%
\end{pgfscope}%
\begin{pgfscope}%
\pgfpathrectangle{\pgfqpoint{0.100000in}{0.100000in}}{\pgfqpoint{3.608454in}{2.310000in}}%
\pgfusepath{clip}%
\pgfsetbuttcap%
\pgfsetroundjoin%
\pgfsetlinewidth{0.010037pt}%
\definecolor{currentstroke}{rgb}{1.000000,1.000000,1.000000}%
\pgfsetstrokecolor{currentstroke}%
\pgfsetdash{}{0pt}%
\pgfusepath{stroke}%
\end{pgfscope}%
\begin{pgfscope}%
\pgfpathrectangle{\pgfqpoint{0.100000in}{0.100000in}}{\pgfqpoint{3.608454in}{2.310000in}}%
\pgfusepath{clip}%
\pgfsetbuttcap%
\pgfsetroundjoin%
\pgfsetlinewidth{0.010037pt}%
\definecolor{currentstroke}{rgb}{1.000000,1.000000,1.000000}%
\pgfsetstrokecolor{currentstroke}%
\pgfsetdash{}{0pt}%
\pgfusepath{stroke}%
\end{pgfscope}%
\begin{pgfscope}%
\pgfpathrectangle{\pgfqpoint{0.100000in}{0.100000in}}{\pgfqpoint{3.608454in}{2.310000in}}%
\pgfusepath{clip}%
\pgfsetbuttcap%
\pgfsetroundjoin%
\pgfsetlinewidth{0.010037pt}%
\definecolor{currentstroke}{rgb}{1.000000,1.000000,1.000000}%
\pgfsetstrokecolor{currentstroke}%
\pgfsetdash{}{0pt}%
\pgfusepath{stroke}%
\end{pgfscope}%
\begin{pgfscope}%
\pgfpathrectangle{\pgfqpoint{0.100000in}{0.100000in}}{\pgfqpoint{3.608454in}{2.310000in}}%
\pgfusepath{clip}%
\pgfsetbuttcap%
\pgfsetroundjoin%
\pgfsetlinewidth{0.010037pt}%
\definecolor{currentstroke}{rgb}{1.000000,1.000000,1.000000}%
\pgfsetstrokecolor{currentstroke}%
\pgfsetdash{}{0pt}%
\pgfusepath{stroke}%
\end{pgfscope}%
\begin{pgfscope}%
\pgfpathrectangle{\pgfqpoint{0.100000in}{0.100000in}}{\pgfqpoint{3.608454in}{2.310000in}}%
\pgfusepath{clip}%
\pgfsetbuttcap%
\pgfsetroundjoin%
\pgfsetlinewidth{0.010037pt}%
\definecolor{currentstroke}{rgb}{1.000000,1.000000,1.000000}%
\pgfsetstrokecolor{currentstroke}%
\pgfsetdash{}{0pt}%
\pgfusepath{stroke}%
\end{pgfscope}%
\begin{pgfscope}%
\pgfpathrectangle{\pgfqpoint{0.100000in}{0.100000in}}{\pgfqpoint{3.608454in}{2.310000in}}%
\pgfusepath{clip}%
\pgfsetbuttcap%
\pgfsetroundjoin%
\pgfsetlinewidth{0.010037pt}%
\definecolor{currentstroke}{rgb}{1.000000,1.000000,1.000000}%
\pgfsetstrokecolor{currentstroke}%
\pgfsetdash{}{0pt}%
\pgfusepath{stroke}%
\end{pgfscope}%
\begin{pgfscope}%
\pgfpathrectangle{\pgfqpoint{0.100000in}{0.100000in}}{\pgfqpoint{3.608454in}{2.310000in}}%
\pgfusepath{clip}%
\pgfsetbuttcap%
\pgfsetroundjoin%
\pgfsetlinewidth{0.010037pt}%
\definecolor{currentstroke}{rgb}{1.000000,1.000000,1.000000}%
\pgfsetstrokecolor{currentstroke}%
\pgfsetdash{}{0pt}%
\pgfusepath{stroke}%
\end{pgfscope}%
\begin{pgfscope}%
\pgfpathrectangle{\pgfqpoint{0.100000in}{0.100000in}}{\pgfqpoint{3.608454in}{2.310000in}}%
\pgfusepath{clip}%
\pgfsetbuttcap%
\pgfsetroundjoin%
\pgfsetlinewidth{0.010037pt}%
\definecolor{currentstroke}{rgb}{1.000000,1.000000,1.000000}%
\pgfsetstrokecolor{currentstroke}%
\pgfsetdash{}{0pt}%
\pgfusepath{stroke}%
\end{pgfscope}%
\begin{pgfscope}%
\pgfpathrectangle{\pgfqpoint{0.100000in}{0.100000in}}{\pgfqpoint{3.608454in}{2.310000in}}%
\pgfusepath{clip}%
\pgfsetbuttcap%
\pgfsetroundjoin%
\pgfsetlinewidth{0.010037pt}%
\definecolor{currentstroke}{rgb}{1.000000,1.000000,1.000000}%
\pgfsetstrokecolor{currentstroke}%
\pgfsetdash{}{0pt}%
\pgfusepath{stroke}%
\end{pgfscope}%
\begin{pgfscope}%
\pgfpathrectangle{\pgfqpoint{0.100000in}{0.100000in}}{\pgfqpoint{3.608454in}{2.310000in}}%
\pgfusepath{clip}%
\pgfsetbuttcap%
\pgfsetroundjoin%
\pgfsetlinewidth{0.010037pt}%
\definecolor{currentstroke}{rgb}{1.000000,1.000000,1.000000}%
\pgfsetstrokecolor{currentstroke}%
\pgfsetdash{}{0pt}%
\pgfusepath{stroke}%
\end{pgfscope}%
\begin{pgfscope}%
\pgfpathrectangle{\pgfqpoint{0.100000in}{0.100000in}}{\pgfqpoint{3.608454in}{2.310000in}}%
\pgfusepath{clip}%
\pgfsetbuttcap%
\pgfsetroundjoin%
\pgfsetlinewidth{0.010037pt}%
\definecolor{currentstroke}{rgb}{1.000000,1.000000,1.000000}%
\pgfsetstrokecolor{currentstroke}%
\pgfsetdash{}{0pt}%
\pgfusepath{stroke}%
\end{pgfscope}%
\begin{pgfscope}%
\pgfpathrectangle{\pgfqpoint{0.100000in}{0.100000in}}{\pgfqpoint{3.608454in}{2.310000in}}%
\pgfusepath{clip}%
\pgfsetbuttcap%
\pgfsetroundjoin%
\pgfsetlinewidth{0.010037pt}%
\definecolor{currentstroke}{rgb}{1.000000,1.000000,1.000000}%
\pgfsetstrokecolor{currentstroke}%
\pgfsetdash{}{0pt}%
\pgfusepath{stroke}%
\end{pgfscope}%
\begin{pgfscope}%
\pgfpathrectangle{\pgfqpoint{0.100000in}{0.100000in}}{\pgfqpoint{3.608454in}{2.310000in}}%
\pgfusepath{clip}%
\pgfsetbuttcap%
\pgfsetroundjoin%
\pgfsetlinewidth{0.010037pt}%
\definecolor{currentstroke}{rgb}{1.000000,1.000000,1.000000}%
\pgfsetstrokecolor{currentstroke}%
\pgfsetdash{}{0pt}%
\pgfusepath{stroke}%
\end{pgfscope}%
\begin{pgfscope}%
\pgfpathrectangle{\pgfqpoint{0.100000in}{0.100000in}}{\pgfqpoint{3.608454in}{2.310000in}}%
\pgfusepath{clip}%
\pgfsetbuttcap%
\pgfsetroundjoin%
\pgfsetlinewidth{0.010037pt}%
\definecolor{currentstroke}{rgb}{1.000000,1.000000,1.000000}%
\pgfsetstrokecolor{currentstroke}%
\pgfsetdash{}{0pt}%
\pgfusepath{stroke}%
\end{pgfscope}%
\begin{pgfscope}%
\pgfpathrectangle{\pgfqpoint{0.100000in}{0.100000in}}{\pgfqpoint{3.608454in}{2.310000in}}%
\pgfusepath{clip}%
\pgfsetbuttcap%
\pgfsetroundjoin%
\pgfsetlinewidth{0.010037pt}%
\definecolor{currentstroke}{rgb}{1.000000,1.000000,1.000000}%
\pgfsetstrokecolor{currentstroke}%
\pgfsetdash{}{0pt}%
\pgfusepath{stroke}%
\end{pgfscope}%
\begin{pgfscope}%
\pgfpathrectangle{\pgfqpoint{0.100000in}{0.100000in}}{\pgfqpoint{3.608454in}{2.310000in}}%
\pgfusepath{clip}%
\pgfsetbuttcap%
\pgfsetroundjoin%
\pgfsetlinewidth{0.010037pt}%
\definecolor{currentstroke}{rgb}{1.000000,1.000000,1.000000}%
\pgfsetstrokecolor{currentstroke}%
\pgfsetdash{}{0pt}%
\pgfusepath{stroke}%
\end{pgfscope}%
\begin{pgfscope}%
\pgfpathrectangle{\pgfqpoint{0.100000in}{0.100000in}}{\pgfqpoint{3.608454in}{2.310000in}}%
\pgfusepath{clip}%
\pgfsetbuttcap%
\pgfsetroundjoin%
\pgfsetlinewidth{0.010037pt}%
\definecolor{currentstroke}{rgb}{1.000000,1.000000,1.000000}%
\pgfsetstrokecolor{currentstroke}%
\pgfsetdash{}{0pt}%
\pgfusepath{stroke}%
\end{pgfscope}%
\begin{pgfscope}%
\pgfpathrectangle{\pgfqpoint{0.100000in}{0.100000in}}{\pgfqpoint{3.608454in}{2.310000in}}%
\pgfusepath{clip}%
\pgfsetbuttcap%
\pgfsetroundjoin%
\pgfsetlinewidth{0.010037pt}%
\definecolor{currentstroke}{rgb}{1.000000,1.000000,1.000000}%
\pgfsetstrokecolor{currentstroke}%
\pgfsetdash{}{0pt}%
\pgfusepath{stroke}%
\end{pgfscope}%
\begin{pgfscope}%
\pgfpathrectangle{\pgfqpoint{0.100000in}{0.100000in}}{\pgfqpoint{3.608454in}{2.310000in}}%
\pgfusepath{clip}%
\pgfsetbuttcap%
\pgfsetroundjoin%
\pgfsetlinewidth{0.010037pt}%
\definecolor{currentstroke}{rgb}{1.000000,1.000000,1.000000}%
\pgfsetstrokecolor{currentstroke}%
\pgfsetdash{}{0pt}%
\pgfusepath{stroke}%
\end{pgfscope}%
\begin{pgfscope}%
\pgfpathrectangle{\pgfqpoint{0.100000in}{0.100000in}}{\pgfqpoint{3.608454in}{2.310000in}}%
\pgfusepath{clip}%
\pgfsetbuttcap%
\pgfsetmiterjoin%
\pgfsetlinewidth{0.000000pt}%
\definecolor{currentstroke}{rgb}{0.000000,0.000000,0.000000}%
\pgfsetstrokecolor{currentstroke}%
\pgfsetstrokeopacity{0.000000}%
\pgfsetdash{}{0pt}%
\pgfpathmoveto{\pgfqpoint{0.100000in}{0.100000in}}%
\pgfpathlineto{\pgfqpoint{3.708454in}{0.100000in}}%
\pgfpathlineto{\pgfqpoint{3.708454in}{2.410000in}}%
\pgfpathlineto{\pgfqpoint{0.100000in}{2.410000in}}%
\pgfpathlineto{\pgfqpoint{0.100000in}{0.100000in}}%
\pgfpathclose%
\pgfusepath{}%
\end{pgfscope}%
\begin{pgfscope}%
\pgfpathrectangle{\pgfqpoint{0.100000in}{0.100000in}}{\pgfqpoint{3.608454in}{2.310000in}}%
\pgfusepath{clip}%
\pgfsetbuttcap%
\pgfsetmiterjoin%
\definecolor{currentfill}{rgb}{0.000000,0.364706,0.817647}%
\pgfsetfillcolor{currentfill}%
\pgfsetlinewidth{0.000000pt}%
\definecolor{currentstroke}{rgb}{0.000000,0.000000,0.000000}%
\pgfsetstrokecolor{currentstroke}%
\pgfsetstrokeopacity{0.000000}%
\pgfsetdash{}{0pt}%
\pgfpathmoveto{\pgfqpoint{2.124622in}{1.510289in}}%
\pgfpathlineto{\pgfqpoint{2.097514in}{1.510367in}}%
\pgfpathlineto{\pgfqpoint{2.097619in}{1.537721in}}%
\pgfpathlineto{\pgfqpoint{2.077274in}{1.537859in}}%
\pgfpathlineto{\pgfqpoint{2.073777in}{1.545495in}}%
\pgfpathlineto{\pgfqpoint{2.073952in}{1.566162in}}%
\pgfpathlineto{\pgfqpoint{2.121832in}{1.566001in}}%
\pgfpathlineto{\pgfqpoint{2.121931in}{1.545408in}}%
\pgfpathlineto{\pgfqpoint{2.124501in}{1.545418in}}%
\pgfpathlineto{\pgfqpoint{2.124622in}{1.510289in}}%
\pgfpathclose%
\pgfusepath{fill}%
\end{pgfscope}%
\begin{pgfscope}%
\pgfpathrectangle{\pgfqpoint{0.100000in}{0.100000in}}{\pgfqpoint{3.608454in}{2.310000in}}%
\pgfusepath{clip}%
\pgfsetbuttcap%
\pgfsetmiterjoin%
\definecolor{currentfill}{rgb}{0.000000,0.850980,0.574510}%
\pgfsetfillcolor{currentfill}%
\pgfsetlinewidth{0.000000pt}%
\definecolor{currentstroke}{rgb}{0.000000,0.000000,0.000000}%
\pgfsetstrokecolor{currentstroke}%
\pgfsetstrokeopacity{0.000000}%
\pgfsetdash{}{0pt}%
\pgfpathmoveto{\pgfqpoint{2.724860in}{1.203323in}}%
\pgfpathlineto{\pgfqpoint{2.731198in}{1.201857in}}%
\pgfpathlineto{\pgfqpoint{2.747524in}{1.217198in}}%
\pgfpathlineto{\pgfqpoint{2.742647in}{1.220616in}}%
\pgfpathlineto{\pgfqpoint{2.742337in}{1.226077in}}%
\pgfpathlineto{\pgfqpoint{2.752556in}{1.236924in}}%
\pgfpathlineto{\pgfqpoint{2.761764in}{1.247341in}}%
\pgfpathlineto{\pgfqpoint{2.765354in}{1.245580in}}%
\pgfpathlineto{\pgfqpoint{2.780361in}{1.233265in}}%
\pgfpathlineto{\pgfqpoint{2.778898in}{1.226398in}}%
\pgfpathlineto{\pgfqpoint{2.781636in}{1.214245in}}%
\pgfpathlineto{\pgfqpoint{2.779295in}{1.206052in}}%
\pgfpathlineto{\pgfqpoint{2.783865in}{1.195392in}}%
\pgfpathlineto{\pgfqpoint{2.789452in}{1.189777in}}%
\pgfpathlineto{\pgfqpoint{2.788196in}{1.183891in}}%
\pgfpathlineto{\pgfqpoint{2.783231in}{1.181649in}}%
\pgfpathlineto{\pgfqpoint{2.785832in}{1.164847in}}%
\pgfpathlineto{\pgfqpoint{2.781495in}{1.155039in}}%
\pgfpathlineto{\pgfqpoint{2.771140in}{1.159342in}}%
\pgfpathlineto{\pgfqpoint{2.761563in}{1.168904in}}%
\pgfpathlineto{\pgfqpoint{2.763981in}{1.171120in}}%
\pgfpathlineto{\pgfqpoint{2.758225in}{1.180770in}}%
\pgfpathlineto{\pgfqpoint{2.754147in}{1.186574in}}%
\pgfpathlineto{\pgfqpoint{2.722959in}{1.184982in}}%
\pgfpathlineto{\pgfqpoint{2.724860in}{1.203323in}}%
\pgfpathclose%
\pgfusepath{fill}%
\end{pgfscope}%
\begin{pgfscope}%
\pgfpathrectangle{\pgfqpoint{0.100000in}{0.100000in}}{\pgfqpoint{3.608454in}{2.310000in}}%
\pgfusepath{clip}%
\pgfsetbuttcap%
\pgfsetmiterjoin%
\definecolor{currentfill}{rgb}{0.000000,0.635294,0.682353}%
\pgfsetfillcolor{currentfill}%
\pgfsetlinewidth{0.000000pt}%
\definecolor{currentstroke}{rgb}{0.000000,0.000000,0.000000}%
\pgfsetstrokecolor{currentstroke}%
\pgfsetstrokeopacity{0.000000}%
\pgfsetdash{}{0pt}%
\pgfpathmoveto{\pgfqpoint{1.745610in}{0.979809in}}%
\pgfpathlineto{\pgfqpoint{1.743217in}{0.941925in}}%
\pgfpathlineto{\pgfqpoint{1.708835in}{0.943953in}}%
\pgfpathlineto{\pgfqpoint{1.711338in}{0.981859in}}%
\pgfpathlineto{\pgfqpoint{1.745610in}{0.979809in}}%
\pgfpathclose%
\pgfusepath{fill}%
\end{pgfscope}%
\begin{pgfscope}%
\pgfpathrectangle{\pgfqpoint{0.100000in}{0.100000in}}{\pgfqpoint{3.608454in}{2.310000in}}%
\pgfusepath{clip}%
\pgfsetbuttcap%
\pgfsetmiterjoin%
\definecolor{currentfill}{rgb}{0.000000,0.713725,0.643137}%
\pgfsetfillcolor{currentfill}%
\pgfsetlinewidth{0.000000pt}%
\definecolor{currentstroke}{rgb}{0.000000,0.000000,0.000000}%
\pgfsetstrokecolor{currentstroke}%
\pgfsetstrokeopacity{0.000000}%
\pgfsetdash{}{0pt}%
\pgfpathmoveto{\pgfqpoint{1.670049in}{1.298787in}}%
\pgfpathlineto{\pgfqpoint{1.670021in}{1.298287in}}%
\pgfpathlineto{\pgfqpoint{1.666954in}{1.256981in}}%
\pgfpathlineto{\pgfqpoint{1.666412in}{1.249548in}}%
\pgfpathlineto{\pgfqpoint{1.601877in}{1.254860in}}%
\pgfpathlineto{\pgfqpoint{1.540663in}{1.260697in}}%
\pgfpathlineto{\pgfqpoint{1.541303in}{1.267794in}}%
\pgfpathlineto{\pgfqpoint{1.506183in}{1.284643in}}%
\pgfpathlineto{\pgfqpoint{1.483235in}{1.285604in}}%
\pgfpathlineto{\pgfqpoint{1.481289in}{1.288540in}}%
\pgfpathlineto{\pgfqpoint{1.484320in}{1.315446in}}%
\pgfpathlineto{\pgfqpoint{1.491038in}{1.314728in}}%
\pgfpathlineto{\pgfqpoint{1.493221in}{1.335267in}}%
\pgfpathlineto{\pgfqpoint{1.547834in}{1.329763in}}%
\pgfpathlineto{\pgfqpoint{1.581597in}{1.326016in}}%
\pgfpathlineto{\pgfqpoint{1.579920in}{1.306214in}}%
\pgfpathlineto{\pgfqpoint{1.626877in}{1.302163in}}%
\pgfpathlineto{\pgfqpoint{1.670049in}{1.298787in}}%
\pgfpathclose%
\pgfusepath{fill}%
\end{pgfscope}%
\begin{pgfscope}%
\pgfpathrectangle{\pgfqpoint{0.100000in}{0.100000in}}{\pgfqpoint{3.608454in}{2.310000in}}%
\pgfusepath{clip}%
\pgfsetbuttcap%
\pgfsetmiterjoin%
\definecolor{currentfill}{rgb}{0.000000,0.431373,0.784314}%
\pgfsetfillcolor{currentfill}%
\pgfsetlinewidth{0.000000pt}%
\definecolor{currentstroke}{rgb}{0.000000,0.000000,0.000000}%
\pgfsetstrokecolor{currentstroke}%
\pgfsetstrokeopacity{0.000000}%
\pgfsetdash{}{0pt}%
\pgfpathmoveto{\pgfqpoint{2.103336in}{1.338237in}}%
\pgfpathlineto{\pgfqpoint{2.096838in}{1.338167in}}%
\pgfpathlineto{\pgfqpoint{2.095756in}{1.349849in}}%
\pgfpathlineto{\pgfqpoint{2.095818in}{1.372775in}}%
\pgfpathlineto{\pgfqpoint{2.071958in}{1.372885in}}%
\pgfpathlineto{\pgfqpoint{2.072459in}{1.391367in}}%
\pgfpathlineto{\pgfqpoint{2.086114in}{1.391298in}}%
\pgfpathlineto{\pgfqpoint{2.099783in}{1.391231in}}%
\pgfpathlineto{\pgfqpoint{2.103446in}{1.388933in}}%
\pgfpathlineto{\pgfqpoint{2.100521in}{1.381808in}}%
\pgfpathlineto{\pgfqpoint{2.107556in}{1.374791in}}%
\pgfpathlineto{\pgfqpoint{2.116197in}{1.360877in}}%
\pgfpathlineto{\pgfqpoint{2.112893in}{1.355625in}}%
\pgfpathlineto{\pgfqpoint{2.112374in}{1.338940in}}%
\pgfpathlineto{\pgfqpoint{2.103336in}{1.338237in}}%
\pgfpathclose%
\pgfusepath{fill}%
\end{pgfscope}%
\begin{pgfscope}%
\pgfpathrectangle{\pgfqpoint{0.100000in}{0.100000in}}{\pgfqpoint{3.608454in}{2.310000in}}%
\pgfusepath{clip}%
\pgfsetbuttcap%
\pgfsetmiterjoin%
\definecolor{currentfill}{rgb}{0.000000,0.517647,0.741176}%
\pgfsetfillcolor{currentfill}%
\pgfsetlinewidth{0.000000pt}%
\definecolor{currentstroke}{rgb}{0.000000,0.000000,0.000000}%
\pgfsetstrokecolor{currentstroke}%
\pgfsetstrokeopacity{0.000000}%
\pgfsetdash{}{0pt}%
\pgfpathmoveto{\pgfqpoint{2.760784in}{0.972182in}}%
\pgfpathlineto{\pgfqpoint{2.738788in}{0.969584in}}%
\pgfpathlineto{\pgfqpoint{2.737463in}{0.974424in}}%
\pgfpathlineto{\pgfqpoint{2.728983in}{0.972881in}}%
\pgfpathlineto{\pgfqpoint{2.715165in}{0.972447in}}%
\pgfpathlineto{\pgfqpoint{2.706101in}{0.993510in}}%
\pgfpathlineto{\pgfqpoint{2.710524in}{0.995150in}}%
\pgfpathlineto{\pgfqpoint{2.718615in}{1.004553in}}%
\pgfpathlineto{\pgfqpoint{2.721545in}{1.014250in}}%
\pgfpathlineto{\pgfqpoint{2.725316in}{1.017812in}}%
\pgfpathlineto{\pgfqpoint{2.723904in}{1.022807in}}%
\pgfpathlineto{\pgfqpoint{2.754206in}{1.025936in}}%
\pgfpathlineto{\pgfqpoint{2.753682in}{1.036469in}}%
\pgfpathlineto{\pgfqpoint{2.746037in}{1.035657in}}%
\pgfpathlineto{\pgfqpoint{2.747587in}{1.043751in}}%
\pgfpathlineto{\pgfqpoint{2.752122in}{1.044231in}}%
\pgfpathlineto{\pgfqpoint{2.756131in}{1.058082in}}%
\pgfpathlineto{\pgfqpoint{2.779057in}{1.060719in}}%
\pgfpathlineto{\pgfqpoint{2.780160in}{1.050381in}}%
\pgfpathlineto{\pgfqpoint{2.782890in}{1.048115in}}%
\pgfpathlineto{\pgfqpoint{2.779095in}{1.040030in}}%
\pgfpathlineto{\pgfqpoint{2.780892in}{1.025876in}}%
\pgfpathlineto{\pgfqpoint{2.784817in}{0.988730in}}%
\pgfpathlineto{\pgfqpoint{2.767572in}{0.987144in}}%
\pgfpathlineto{\pgfqpoint{2.765064in}{0.976399in}}%
\pgfpathlineto{\pgfqpoint{2.760784in}{0.972182in}}%
\pgfpathclose%
\pgfusepath{fill}%
\end{pgfscope}%
\begin{pgfscope}%
\pgfpathrectangle{\pgfqpoint{0.100000in}{0.100000in}}{\pgfqpoint{3.608454in}{2.310000in}}%
\pgfusepath{clip}%
\pgfsetbuttcap%
\pgfsetmiterjoin%
\definecolor{currentfill}{rgb}{0.000000,0.431373,0.784314}%
\pgfsetfillcolor{currentfill}%
\pgfsetlinewidth{0.000000pt}%
\definecolor{currentstroke}{rgb}{0.000000,0.000000,0.000000}%
\pgfsetstrokecolor{currentstroke}%
\pgfsetstrokeopacity{0.000000}%
\pgfsetdash{}{0pt}%
\pgfpathmoveto{\pgfqpoint{1.716695in}{1.867335in}}%
\pgfpathlineto{\pgfqpoint{1.714041in}{1.832923in}}%
\pgfpathlineto{\pgfqpoint{1.660454in}{1.837408in}}%
\pgfpathlineto{\pgfqpoint{1.661648in}{1.851140in}}%
\pgfpathlineto{\pgfqpoint{1.667609in}{1.909112in}}%
\pgfpathlineto{\pgfqpoint{1.664670in}{1.909384in}}%
\pgfpathlineto{\pgfqpoint{1.667012in}{1.936010in}}%
\pgfpathlineto{\pgfqpoint{1.670909in}{1.935733in}}%
\pgfpathlineto{\pgfqpoint{1.673140in}{1.963430in}}%
\pgfpathlineto{\pgfqpoint{1.718524in}{1.959750in}}%
\pgfpathlineto{\pgfqpoint{1.721015in}{1.959541in}}%
\pgfpathlineto{\pgfqpoint{1.719106in}{1.931926in}}%
\pgfpathlineto{\pgfqpoint{1.721787in}{1.931717in}}%
\pgfpathlineto{\pgfqpoint{1.720336in}{1.913435in}}%
\pgfpathlineto{\pgfqpoint{1.716695in}{1.867335in}}%
\pgfpathclose%
\pgfusepath{fill}%
\end{pgfscope}%
\begin{pgfscope}%
\pgfpathrectangle{\pgfqpoint{0.100000in}{0.100000in}}{\pgfqpoint{3.608454in}{2.310000in}}%
\pgfusepath{clip}%
\pgfsetbuttcap%
\pgfsetmiterjoin%
\definecolor{currentfill}{rgb}{0.000000,0.647059,0.676471}%
\pgfsetfillcolor{currentfill}%
\pgfsetlinewidth{0.000000pt}%
\definecolor{currentstroke}{rgb}{0.000000,0.000000,0.000000}%
\pgfsetstrokecolor{currentstroke}%
\pgfsetstrokeopacity{0.000000}%
\pgfsetdash{}{0pt}%
\pgfpathmoveto{\pgfqpoint{3.298349in}{1.383037in}}%
\pgfpathlineto{\pgfqpoint{3.304739in}{1.388327in}}%
\pgfpathlineto{\pgfqpoint{3.327417in}{1.395926in}}%
\pgfpathlineto{\pgfqpoint{3.322504in}{1.381252in}}%
\pgfpathlineto{\pgfqpoint{3.318671in}{1.381735in}}%
\pgfpathlineto{\pgfqpoint{3.314476in}{1.374248in}}%
\pgfpathlineto{\pgfqpoint{3.312084in}{1.358499in}}%
\pgfpathlineto{\pgfqpoint{3.313688in}{1.355781in}}%
\pgfpathlineto{\pgfqpoint{3.309070in}{1.336218in}}%
\pgfpathlineto{\pgfqpoint{3.307238in}{1.324528in}}%
\pgfpathlineto{\pgfqpoint{3.302528in}{1.316927in}}%
\pgfpathlineto{\pgfqpoint{3.298046in}{1.315923in}}%
\pgfpathlineto{\pgfqpoint{3.292330in}{1.326062in}}%
\pgfpathlineto{\pgfqpoint{3.292495in}{1.350793in}}%
\pgfpathlineto{\pgfqpoint{3.295985in}{1.361073in}}%
\pgfpathlineto{\pgfqpoint{3.296227in}{1.370099in}}%
\pgfpathlineto{\pgfqpoint{3.301816in}{1.371758in}}%
\pgfpathlineto{\pgfqpoint{3.303055in}{1.378748in}}%
\pgfpathlineto{\pgfqpoint{3.298349in}{1.383037in}}%
\pgfpathclose%
\pgfusepath{fill}%
\end{pgfscope}%
\begin{pgfscope}%
\pgfpathrectangle{\pgfqpoint{0.100000in}{0.100000in}}{\pgfqpoint{3.608454in}{2.310000in}}%
\pgfusepath{clip}%
\pgfsetbuttcap%
\pgfsetmiterjoin%
\definecolor{currentfill}{rgb}{0.000000,0.603922,0.698039}%
\pgfsetfillcolor{currentfill}%
\pgfsetlinewidth{0.000000pt}%
\definecolor{currentstroke}{rgb}{0.000000,0.000000,0.000000}%
\pgfsetstrokecolor{currentstroke}%
\pgfsetstrokeopacity{0.000000}%
\pgfsetdash{}{0pt}%
\pgfpathmoveto{\pgfqpoint{2.488486in}{1.795943in}}%
\pgfpathlineto{\pgfqpoint{2.468096in}{1.794922in}}%
\pgfpathlineto{\pgfqpoint{2.467658in}{1.801883in}}%
\pgfpathlineto{\pgfqpoint{2.432784in}{1.799789in}}%
\pgfpathlineto{\pgfqpoint{2.430551in}{1.834455in}}%
\pgfpathlineto{\pgfqpoint{2.419246in}{1.833798in}}%
\pgfpathlineto{\pgfqpoint{2.417621in}{1.861430in}}%
\pgfpathlineto{\pgfqpoint{2.445567in}{1.863054in}}%
\pgfpathlineto{\pgfqpoint{2.445861in}{1.856178in}}%
\pgfpathlineto{\pgfqpoint{2.459694in}{1.857127in}}%
\pgfpathlineto{\pgfqpoint{2.460529in}{1.843339in}}%
\pgfpathlineto{\pgfqpoint{2.463191in}{1.836540in}}%
\pgfpathlineto{\pgfqpoint{2.472001in}{1.837128in}}%
\pgfpathlineto{\pgfqpoint{2.473187in}{1.816374in}}%
\pgfpathlineto{\pgfqpoint{2.486722in}{1.817142in}}%
\pgfpathlineto{\pgfqpoint{2.488486in}{1.795943in}}%
\pgfpathclose%
\pgfusepath{fill}%
\end{pgfscope}%
\begin{pgfscope}%
\pgfpathrectangle{\pgfqpoint{0.100000in}{0.100000in}}{\pgfqpoint{3.608454in}{2.310000in}}%
\pgfusepath{clip}%
\pgfsetbuttcap%
\pgfsetmiterjoin%
\definecolor{currentfill}{rgb}{0.000000,0.556863,0.721569}%
\pgfsetfillcolor{currentfill}%
\pgfsetlinewidth{0.000000pt}%
\definecolor{currentstroke}{rgb}{0.000000,0.000000,0.000000}%
\pgfsetstrokecolor{currentstroke}%
\pgfsetstrokeopacity{0.000000}%
\pgfsetdash{}{0pt}%
\pgfpathmoveto{\pgfqpoint{3.384285in}{1.735025in}}%
\pgfpathlineto{\pgfqpoint{3.359316in}{1.729735in}}%
\pgfpathlineto{\pgfqpoint{3.358769in}{1.729605in}}%
\pgfpathlineto{\pgfqpoint{3.357447in}{1.732265in}}%
\pgfpathlineto{\pgfqpoint{3.358955in}{1.786375in}}%
\pgfpathlineto{\pgfqpoint{3.356483in}{1.793242in}}%
\pgfpathlineto{\pgfqpoint{3.348859in}{1.830296in}}%
\pgfpathlineto{\pgfqpoint{3.370781in}{1.834556in}}%
\pgfpathlineto{\pgfqpoint{3.374364in}{1.831752in}}%
\pgfpathlineto{\pgfqpoint{3.374319in}{1.820118in}}%
\pgfpathlineto{\pgfqpoint{3.367008in}{1.818791in}}%
\pgfpathlineto{\pgfqpoint{3.369972in}{1.804097in}}%
\pgfpathlineto{\pgfqpoint{3.374916in}{1.805051in}}%
\pgfpathlineto{\pgfqpoint{3.377987in}{1.790511in}}%
\pgfpathlineto{\pgfqpoint{3.373418in}{1.786388in}}%
\pgfpathlineto{\pgfqpoint{3.378669in}{1.782710in}}%
\pgfpathlineto{\pgfqpoint{3.380209in}{1.761732in}}%
\pgfpathlineto{\pgfqpoint{3.384285in}{1.735025in}}%
\pgfpathclose%
\pgfusepath{fill}%
\end{pgfscope}%
\begin{pgfscope}%
\pgfpathrectangle{\pgfqpoint{0.100000in}{0.100000in}}{\pgfqpoint{3.608454in}{2.310000in}}%
\pgfusepath{clip}%
\pgfsetbuttcap%
\pgfsetmiterjoin%
\definecolor{currentfill}{rgb}{0.000000,0.752941,0.623529}%
\pgfsetfillcolor{currentfill}%
\pgfsetlinewidth{0.000000pt}%
\definecolor{currentstroke}{rgb}{0.000000,0.000000,0.000000}%
\pgfsetstrokecolor{currentstroke}%
\pgfsetstrokeopacity{0.000000}%
\pgfsetdash{}{0pt}%
\pgfpathmoveto{\pgfqpoint{2.895571in}{1.311319in}}%
\pgfpathlineto{\pgfqpoint{2.890443in}{1.304155in}}%
\pgfpathlineto{\pgfqpoint{2.884060in}{1.308388in}}%
\pgfpathlineto{\pgfqpoint{2.885378in}{1.313095in}}%
\pgfpathlineto{\pgfqpoint{2.873043in}{1.323790in}}%
\pgfpathlineto{\pgfqpoint{2.874015in}{1.332782in}}%
\pgfpathlineto{\pgfqpoint{2.861780in}{1.330306in}}%
\pgfpathlineto{\pgfqpoint{2.855274in}{1.324556in}}%
\pgfpathlineto{\pgfqpoint{2.857831in}{1.319548in}}%
\pgfpathlineto{\pgfqpoint{2.850698in}{1.310252in}}%
\pgfpathlineto{\pgfqpoint{2.840076in}{1.309133in}}%
\pgfpathlineto{\pgfqpoint{2.834982in}{1.312343in}}%
\pgfpathlineto{\pgfqpoint{2.824999in}{1.309058in}}%
\pgfpathlineto{\pgfqpoint{2.811045in}{1.319515in}}%
\pgfpathlineto{\pgfqpoint{2.818045in}{1.327218in}}%
\pgfpathlineto{\pgfqpoint{2.817804in}{1.333230in}}%
\pgfpathlineto{\pgfqpoint{2.822416in}{1.337584in}}%
\pgfpathlineto{\pgfqpoint{2.826872in}{1.332234in}}%
\pgfpathlineto{\pgfqpoint{2.834430in}{1.334621in}}%
\pgfpathlineto{\pgfqpoint{2.833231in}{1.340587in}}%
\pgfpathlineto{\pgfqpoint{2.835681in}{1.347268in}}%
\pgfpathlineto{\pgfqpoint{2.841606in}{1.355934in}}%
\pgfpathlineto{\pgfqpoint{2.842704in}{1.366486in}}%
\pgfpathlineto{\pgfqpoint{2.850733in}{1.370027in}}%
\pgfpathlineto{\pgfqpoint{2.854867in}{1.358657in}}%
\pgfpathlineto{\pgfqpoint{2.861822in}{1.358808in}}%
\pgfpathlineto{\pgfqpoint{2.864952in}{1.371264in}}%
\pgfpathlineto{\pgfqpoint{2.864652in}{1.379289in}}%
\pgfpathlineto{\pgfqpoint{2.869245in}{1.379552in}}%
\pgfpathlineto{\pgfqpoint{2.869516in}{1.374371in}}%
\pgfpathlineto{\pgfqpoint{2.875638in}{1.374674in}}%
\pgfpathlineto{\pgfqpoint{2.877293in}{1.367722in}}%
\pgfpathlineto{\pgfqpoint{2.884573in}{1.368240in}}%
\pgfpathlineto{\pgfqpoint{2.885129in}{1.361064in}}%
\pgfpathlineto{\pgfqpoint{2.893788in}{1.362802in}}%
\pgfpathlineto{\pgfqpoint{2.905122in}{1.355064in}}%
\pgfpathlineto{\pgfqpoint{2.906738in}{1.347296in}}%
\pgfpathlineto{\pgfqpoint{2.898128in}{1.341949in}}%
\pgfpathlineto{\pgfqpoint{2.894444in}{1.332079in}}%
\pgfpathlineto{\pgfqpoint{2.894629in}{1.327193in}}%
\pgfpathlineto{\pgfqpoint{2.900793in}{1.318863in}}%
\pgfpathlineto{\pgfqpoint{2.895571in}{1.311319in}}%
\pgfpathclose%
\pgfusepath{fill}%
\end{pgfscope}%
\begin{pgfscope}%
\pgfpathrectangle{\pgfqpoint{0.100000in}{0.100000in}}{\pgfqpoint{3.608454in}{2.310000in}}%
\pgfusepath{clip}%
\pgfsetbuttcap%
\pgfsetmiterjoin%
\definecolor{currentfill}{rgb}{0.000000,0.623529,0.688235}%
\pgfsetfillcolor{currentfill}%
\pgfsetlinewidth{0.000000pt}%
\definecolor{currentstroke}{rgb}{0.000000,0.000000,0.000000}%
\pgfsetstrokecolor{currentstroke}%
\pgfsetstrokeopacity{0.000000}%
\pgfsetdash{}{0pt}%
\pgfpathmoveto{\pgfqpoint{1.413364in}{1.067432in}}%
\pgfpathlineto{\pgfqpoint{1.379252in}{1.071453in}}%
\pgfpathlineto{\pgfqpoint{1.385971in}{1.127718in}}%
\pgfpathlineto{\pgfqpoint{1.377042in}{1.134592in}}%
\pgfpathlineto{\pgfqpoint{1.378347in}{1.144995in}}%
\pgfpathlineto{\pgfqpoint{1.388440in}{1.147038in}}%
\pgfpathlineto{\pgfqpoint{1.347890in}{1.152024in}}%
\pgfpathlineto{\pgfqpoint{1.350154in}{1.169097in}}%
\pgfpathlineto{\pgfqpoint{1.303419in}{1.175695in}}%
\pgfpathlineto{\pgfqpoint{1.310197in}{1.221478in}}%
\pgfpathlineto{\pgfqpoint{1.322223in}{1.228643in}}%
\pgfpathlineto{\pgfqpoint{1.324687in}{1.235061in}}%
\pgfpathlineto{\pgfqpoint{1.383780in}{1.226748in}}%
\pgfpathlineto{\pgfqpoint{1.462544in}{1.217413in}}%
\pgfpathlineto{\pgfqpoint{1.462472in}{1.211222in}}%
\pgfpathlineto{\pgfqpoint{1.458207in}{1.195789in}}%
\pgfpathlineto{\pgfqpoint{1.450272in}{1.192914in}}%
\pgfpathlineto{\pgfqpoint{1.450548in}{1.172051in}}%
\pgfpathlineto{\pgfqpoint{1.448410in}{1.160662in}}%
\pgfpathlineto{\pgfqpoint{1.445381in}{1.160260in}}%
\pgfpathlineto{\pgfqpoint{1.442215in}{1.150974in}}%
\pgfpathlineto{\pgfqpoint{1.431076in}{1.139578in}}%
\pgfpathlineto{\pgfqpoint{1.421911in}{1.141159in}}%
\pgfpathlineto{\pgfqpoint{1.413364in}{1.067432in}}%
\pgfpathclose%
\pgfusepath{fill}%
\end{pgfscope}%
\begin{pgfscope}%
\pgfpathrectangle{\pgfqpoint{0.100000in}{0.100000in}}{\pgfqpoint{3.608454in}{2.310000in}}%
\pgfusepath{clip}%
\pgfsetbuttcap%
\pgfsetmiterjoin%
\definecolor{currentfill}{rgb}{0.000000,0.847059,0.576471}%
\pgfsetfillcolor{currentfill}%
\pgfsetlinewidth{0.000000pt}%
\definecolor{currentstroke}{rgb}{0.000000,0.000000,0.000000}%
\pgfsetstrokecolor{currentstroke}%
\pgfsetstrokeopacity{0.000000}%
\pgfsetdash{}{0pt}%
\pgfpathmoveto{\pgfqpoint{2.407413in}{1.214350in}}%
\pgfpathlineto{\pgfqpoint{2.386546in}{1.213497in}}%
\pgfpathlineto{\pgfqpoint{2.386664in}{1.209942in}}%
\pgfpathlineto{\pgfqpoint{2.375052in}{1.209261in}}%
\pgfpathlineto{\pgfqpoint{2.372317in}{1.217256in}}%
\pgfpathlineto{\pgfqpoint{2.371740in}{1.235384in}}%
\pgfpathlineto{\pgfqpoint{2.358074in}{1.234831in}}%
\pgfpathlineto{\pgfqpoint{2.347612in}{1.233131in}}%
\pgfpathlineto{\pgfqpoint{2.347355in}{1.241765in}}%
\pgfpathlineto{\pgfqpoint{2.350544in}{1.245278in}}%
\pgfpathlineto{\pgfqpoint{2.349233in}{1.281942in}}%
\pgfpathlineto{\pgfqpoint{2.368816in}{1.282820in}}%
\pgfpathlineto{\pgfqpoint{2.381801in}{1.267561in}}%
\pgfpathlineto{\pgfqpoint{2.392056in}{1.271141in}}%
\pgfpathlineto{\pgfqpoint{2.397241in}{1.275533in}}%
\pgfpathlineto{\pgfqpoint{2.404880in}{1.275417in}}%
\pgfpathlineto{\pgfqpoint{2.409704in}{1.273645in}}%
\pgfpathlineto{\pgfqpoint{2.418250in}{1.266369in}}%
\pgfpathlineto{\pgfqpoint{2.421525in}{1.260176in}}%
\pgfpathlineto{\pgfqpoint{2.428315in}{1.262231in}}%
\pgfpathlineto{\pgfqpoint{2.434982in}{1.258000in}}%
\pgfpathlineto{\pgfqpoint{2.445368in}{1.248025in}}%
\pgfpathlineto{\pgfqpoint{2.449790in}{1.246442in}}%
\pgfpathlineto{\pgfqpoint{2.452467in}{1.239392in}}%
\pgfpathlineto{\pgfqpoint{2.449905in}{1.236650in}}%
\pgfpathlineto{\pgfqpoint{2.444353in}{1.238826in}}%
\pgfpathlineto{\pgfqpoint{2.410673in}{1.236890in}}%
\pgfpathlineto{\pgfqpoint{2.411944in}{1.214363in}}%
\pgfpathlineto{\pgfqpoint{2.407413in}{1.214350in}}%
\pgfpathclose%
\pgfusepath{fill}%
\end{pgfscope}%
\begin{pgfscope}%
\pgfpathrectangle{\pgfqpoint{0.100000in}{0.100000in}}{\pgfqpoint{3.608454in}{2.310000in}}%
\pgfusepath{clip}%
\pgfsetbuttcap%
\pgfsetmiterjoin%
\definecolor{currentfill}{rgb}{0.000000,0.647059,0.676471}%
\pgfsetfillcolor{currentfill}%
\pgfsetlinewidth{0.000000pt}%
\definecolor{currentstroke}{rgb}{0.000000,0.000000,0.000000}%
\pgfsetstrokecolor{currentstroke}%
\pgfsetstrokeopacity{0.000000}%
\pgfsetdash{}{0pt}%
\pgfpathmoveto{\pgfqpoint{2.751930in}{0.830708in}}%
\pgfpathlineto{\pgfqpoint{2.751810in}{0.820152in}}%
\pgfpathlineto{\pgfqpoint{2.741287in}{0.817640in}}%
\pgfpathlineto{\pgfqpoint{2.737520in}{0.807875in}}%
\pgfpathlineto{\pgfqpoint{2.728545in}{0.806963in}}%
\pgfpathlineto{\pgfqpoint{2.727912in}{0.813855in}}%
\pgfpathlineto{\pgfqpoint{2.721575in}{0.813214in}}%
\pgfpathlineto{\pgfqpoint{2.720311in}{0.819497in}}%
\pgfpathlineto{\pgfqpoint{2.713352in}{0.819097in}}%
\pgfpathlineto{\pgfqpoint{2.711685in}{0.841617in}}%
\pgfpathlineto{\pgfqpoint{2.708963in}{0.847375in}}%
\pgfpathlineto{\pgfqpoint{2.717264in}{0.854842in}}%
\pgfpathlineto{\pgfqpoint{2.723195in}{0.855536in}}%
\pgfpathlineto{\pgfqpoint{2.722392in}{0.862394in}}%
\pgfpathlineto{\pgfqpoint{2.729130in}{0.864274in}}%
\pgfpathlineto{\pgfqpoint{2.728245in}{0.872298in}}%
\pgfpathlineto{\pgfqpoint{2.734886in}{0.875470in}}%
\pgfpathlineto{\pgfqpoint{2.755362in}{0.877808in}}%
\pgfpathlineto{\pgfqpoint{2.765201in}{0.880136in}}%
\pgfpathlineto{\pgfqpoint{2.768066in}{0.872209in}}%
\pgfpathlineto{\pgfqpoint{2.776169in}{0.862453in}}%
\pgfpathlineto{\pgfqpoint{2.772613in}{0.859054in}}%
\pgfpathlineto{\pgfqpoint{2.754330in}{0.856732in}}%
\pgfpathlineto{\pgfqpoint{2.755065in}{0.852187in}}%
\pgfpathlineto{\pgfqpoint{2.748134in}{0.851370in}}%
\pgfpathlineto{\pgfqpoint{2.749632in}{0.837509in}}%
\pgfpathlineto{\pgfqpoint{2.751930in}{0.830708in}}%
\pgfpathclose%
\pgfusepath{fill}%
\end{pgfscope}%
\begin{pgfscope}%
\pgfpathrectangle{\pgfqpoint{0.100000in}{0.100000in}}{\pgfqpoint{3.608454in}{2.310000in}}%
\pgfusepath{clip}%
\pgfsetbuttcap%
\pgfsetmiterjoin%
\definecolor{currentfill}{rgb}{0.000000,0.474510,0.762745}%
\pgfsetfillcolor{currentfill}%
\pgfsetlinewidth{0.000000pt}%
\definecolor{currentstroke}{rgb}{0.000000,0.000000,0.000000}%
\pgfsetstrokecolor{currentstroke}%
\pgfsetstrokeopacity{0.000000}%
\pgfsetdash{}{0pt}%
\pgfpathmoveto{\pgfqpoint{0.975129in}{1.855275in}}%
\pgfpathlineto{\pgfqpoint{0.973897in}{1.850229in}}%
\pgfpathlineto{\pgfqpoint{0.967810in}{1.841171in}}%
\pgfpathlineto{\pgfqpoint{0.968859in}{1.837111in}}%
\pgfpathlineto{\pgfqpoint{0.963091in}{1.830834in}}%
\pgfpathlineto{\pgfqpoint{0.950896in}{1.777540in}}%
\pgfpathlineto{\pgfqpoint{0.953662in}{1.776726in}}%
\pgfpathlineto{\pgfqpoint{0.937739in}{1.705964in}}%
\pgfpathlineto{\pgfqpoint{0.861529in}{1.723655in}}%
\pgfpathlineto{\pgfqpoint{0.824452in}{1.732928in}}%
\pgfpathlineto{\pgfqpoint{0.824016in}{1.733109in}}%
\pgfpathlineto{\pgfqpoint{0.855951in}{1.862354in}}%
\pgfpathlineto{\pgfqpoint{0.860241in}{1.876128in}}%
\pgfpathlineto{\pgfqpoint{0.863735in}{1.876688in}}%
\pgfpathlineto{\pgfqpoint{0.867712in}{1.868584in}}%
\pgfpathlineto{\pgfqpoint{0.875796in}{1.867766in}}%
\pgfpathlineto{\pgfqpoint{0.879074in}{1.881212in}}%
\pgfpathlineto{\pgfqpoint{0.886843in}{1.879251in}}%
\pgfpathlineto{\pgfqpoint{0.890663in}{1.885325in}}%
\pgfpathlineto{\pgfqpoint{0.895063in}{1.884202in}}%
\pgfpathlineto{\pgfqpoint{0.896690in}{1.890839in}}%
\pgfpathlineto{\pgfqpoint{0.901859in}{1.889570in}}%
\pgfpathlineto{\pgfqpoint{0.905607in}{1.902658in}}%
\pgfpathlineto{\pgfqpoint{0.910516in}{1.911417in}}%
\pgfpathlineto{\pgfqpoint{0.919645in}{1.913720in}}%
\pgfpathlineto{\pgfqpoint{0.916455in}{1.900234in}}%
\pgfpathlineto{\pgfqpoint{0.913150in}{1.901010in}}%
\pgfpathlineto{\pgfqpoint{0.909972in}{1.887658in}}%
\pgfpathlineto{\pgfqpoint{0.915455in}{1.886073in}}%
\pgfpathlineto{\pgfqpoint{0.917584in}{1.892808in}}%
\pgfpathlineto{\pgfqpoint{0.949827in}{1.885175in}}%
\pgfpathlineto{\pgfqpoint{0.953430in}{1.884971in}}%
\pgfpathlineto{\pgfqpoint{0.964170in}{1.890384in}}%
\pgfpathlineto{\pgfqpoint{0.969814in}{1.885604in}}%
\pgfpathlineto{\pgfqpoint{0.968833in}{1.878204in}}%
\pgfpathlineto{\pgfqpoint{0.975886in}{1.873015in}}%
\pgfpathlineto{\pgfqpoint{0.972986in}{1.865357in}}%
\pgfpathlineto{\pgfqpoint{0.975129in}{1.855275in}}%
\pgfpathclose%
\pgfusepath{fill}%
\end{pgfscope}%
\begin{pgfscope}%
\pgfpathrectangle{\pgfqpoint{0.100000in}{0.100000in}}{\pgfqpoint{3.608454in}{2.310000in}}%
\pgfusepath{clip}%
\pgfsetbuttcap%
\pgfsetmiterjoin%
\definecolor{currentfill}{rgb}{0.000000,0.607843,0.696078}%
\pgfsetfillcolor{currentfill}%
\pgfsetlinewidth{0.000000pt}%
\definecolor{currentstroke}{rgb}{0.000000,0.000000,0.000000}%
\pgfsetstrokecolor{currentstroke}%
\pgfsetstrokeopacity{0.000000}%
\pgfsetdash{}{0pt}%
\pgfpathmoveto{\pgfqpoint{2.272778in}{0.593877in}}%
\pgfpathlineto{\pgfqpoint{2.264140in}{0.596489in}}%
\pgfpathlineto{\pgfqpoint{2.246043in}{0.604519in}}%
\pgfpathlineto{\pgfqpoint{2.232132in}{0.608488in}}%
\pgfpathlineto{\pgfqpoint{2.221611in}{0.607031in}}%
\pgfpathlineto{\pgfqpoint{2.210990in}{0.607615in}}%
\pgfpathlineto{\pgfqpoint{2.192546in}{0.604264in}}%
\pgfpathlineto{\pgfqpoint{2.187595in}{0.601129in}}%
\pgfpathlineto{\pgfqpoint{2.181395in}{0.610955in}}%
\pgfpathlineto{\pgfqpoint{2.188702in}{0.619056in}}%
\pgfpathlineto{\pgfqpoint{2.190675in}{0.624923in}}%
\pgfpathlineto{\pgfqpoint{2.196138in}{0.630275in}}%
\pgfpathlineto{\pgfqpoint{2.196201in}{0.649140in}}%
\pgfpathlineto{\pgfqpoint{2.191936in}{0.652528in}}%
\pgfpathlineto{\pgfqpoint{2.196505in}{0.661211in}}%
\pgfpathlineto{\pgfqpoint{2.193927in}{0.668423in}}%
\pgfpathlineto{\pgfqpoint{2.197286in}{0.677217in}}%
\pgfpathlineto{\pgfqpoint{2.200790in}{0.680240in}}%
\pgfpathlineto{\pgfqpoint{2.204721in}{0.696052in}}%
\pgfpathlineto{\pgfqpoint{2.207696in}{0.700477in}}%
\pgfpathlineto{\pgfqpoint{2.204984in}{0.705199in}}%
\pgfpathlineto{\pgfqpoint{2.205558in}{0.720837in}}%
\pgfpathlineto{\pgfqpoint{2.206617in}{0.728100in}}%
\pgfpathlineto{\pgfqpoint{2.212945in}{0.728230in}}%
\pgfpathlineto{\pgfqpoint{2.216369in}{0.735263in}}%
\pgfpathlineto{\pgfqpoint{2.226879in}{0.735487in}}%
\pgfpathlineto{\pgfqpoint{2.240752in}{0.735755in}}%
\pgfpathlineto{\pgfqpoint{2.244244in}{0.734424in}}%
\pgfpathlineto{\pgfqpoint{2.251222in}{0.732248in}}%
\pgfpathlineto{\pgfqpoint{2.254758in}{0.727139in}}%
\pgfpathlineto{\pgfqpoint{2.255860in}{0.698248in}}%
\pgfpathlineto{\pgfqpoint{2.271317in}{0.699125in}}%
\pgfpathlineto{\pgfqpoint{2.272007in}{0.674022in}}%
\pgfpathlineto{\pgfqpoint{2.269704in}{0.666105in}}%
\pgfpathlineto{\pgfqpoint{2.271892in}{0.630559in}}%
\pgfpathlineto{\pgfqpoint{2.272778in}{0.593877in}}%
\pgfpathclose%
\pgfusepath{fill}%
\end{pgfscope}%
\begin{pgfscope}%
\pgfpathrectangle{\pgfqpoint{0.100000in}{0.100000in}}{\pgfqpoint{3.608454in}{2.310000in}}%
\pgfusepath{clip}%
\pgfsetbuttcap%
\pgfsetmiterjoin%
\definecolor{currentfill}{rgb}{0.000000,0.850980,0.574510}%
\pgfsetfillcolor{currentfill}%
\pgfsetlinewidth{0.000000pt}%
\definecolor{currentstroke}{rgb}{0.000000,0.000000,0.000000}%
\pgfsetstrokecolor{currentstroke}%
\pgfsetstrokeopacity{0.000000}%
\pgfsetdash{}{0pt}%
\pgfpathmoveto{\pgfqpoint{0.792611in}{1.237224in}}%
\pgfpathlineto{\pgfqpoint{0.786789in}{1.237324in}}%
\pgfpathlineto{\pgfqpoint{0.765878in}{1.242208in}}%
\pgfpathlineto{\pgfqpoint{0.692238in}{1.260226in}}%
\pgfpathlineto{\pgfqpoint{0.668493in}{1.266395in}}%
\pgfpathlineto{\pgfqpoint{0.645007in}{1.271759in}}%
\pgfpathlineto{\pgfqpoint{0.648713in}{1.279430in}}%
\pgfpathlineto{\pgfqpoint{0.647370in}{1.296028in}}%
\pgfpathlineto{\pgfqpoint{0.649389in}{1.301776in}}%
\pgfpathlineto{\pgfqpoint{0.647795in}{1.312826in}}%
\pgfpathlineto{\pgfqpoint{0.650356in}{1.315914in}}%
\pgfpathlineto{\pgfqpoint{0.644042in}{1.329081in}}%
\pgfpathlineto{\pgfqpoint{0.644925in}{1.333356in}}%
\pgfpathlineto{\pgfqpoint{0.641670in}{1.344858in}}%
\pgfpathlineto{\pgfqpoint{0.642344in}{1.357942in}}%
\pgfpathlineto{\pgfqpoint{0.645186in}{1.361551in}}%
\pgfpathlineto{\pgfqpoint{0.643441in}{1.374382in}}%
\pgfpathlineto{\pgfqpoint{0.636091in}{1.384439in}}%
\pgfpathlineto{\pgfqpoint{0.632393in}{1.385724in}}%
\pgfpathlineto{\pgfqpoint{0.632454in}{1.400962in}}%
\pgfpathlineto{\pgfqpoint{0.628697in}{1.402513in}}%
\pgfpathlineto{\pgfqpoint{0.632163in}{1.410294in}}%
\pgfpathlineto{\pgfqpoint{0.627031in}{1.414709in}}%
\pgfpathlineto{\pgfqpoint{0.624814in}{1.421196in}}%
\pgfpathlineto{\pgfqpoint{0.619173in}{1.425619in}}%
\pgfpathlineto{\pgfqpoint{0.618630in}{1.433861in}}%
\pgfpathlineto{\pgfqpoint{0.616083in}{1.438599in}}%
\pgfpathlineto{\pgfqpoint{0.607467in}{1.441382in}}%
\pgfpathlineto{\pgfqpoint{0.612588in}{1.444548in}}%
\pgfpathlineto{\pgfqpoint{0.613111in}{1.453862in}}%
\pgfpathlineto{\pgfqpoint{0.609428in}{1.457861in}}%
\pgfpathlineto{\pgfqpoint{0.610025in}{1.468870in}}%
\pgfpathlineto{\pgfqpoint{0.597769in}{1.478364in}}%
\pgfpathlineto{\pgfqpoint{0.596004in}{1.489421in}}%
\pgfpathlineto{\pgfqpoint{0.597596in}{1.492295in}}%
\pgfpathlineto{\pgfqpoint{0.603382in}{1.497016in}}%
\pgfpathlineto{\pgfqpoint{0.606996in}{1.503863in}}%
\pgfpathlineto{\pgfqpoint{0.604651in}{1.513031in}}%
\pgfpathlineto{\pgfqpoint{0.609163in}{1.520157in}}%
\pgfpathlineto{\pgfqpoint{0.620520in}{1.503015in}}%
\pgfpathlineto{\pgfqpoint{0.628251in}{1.491065in}}%
\pgfpathlineto{\pgfqpoint{0.661400in}{1.439746in}}%
\pgfpathlineto{\pgfqpoint{0.688870in}{1.397366in}}%
\pgfpathlineto{\pgfqpoint{0.720005in}{1.349216in}}%
\pgfpathlineto{\pgfqpoint{0.780428in}{1.255922in}}%
\pgfpathlineto{\pgfqpoint{0.792611in}{1.237224in}}%
\pgfpathclose%
\pgfusepath{fill}%
\end{pgfscope}%
\begin{pgfscope}%
\pgfpathrectangle{\pgfqpoint{0.100000in}{0.100000in}}{\pgfqpoint{3.608454in}{2.310000in}}%
\pgfusepath{clip}%
\pgfsetbuttcap%
\pgfsetmiterjoin%
\definecolor{currentfill}{rgb}{0.000000,0.666667,0.666667}%
\pgfsetfillcolor{currentfill}%
\pgfsetlinewidth{0.000000pt}%
\definecolor{currentstroke}{rgb}{0.000000,0.000000,0.000000}%
\pgfsetstrokecolor{currentstroke}%
\pgfsetstrokeopacity{0.000000}%
\pgfsetdash{}{0pt}%
\pgfpathmoveto{\pgfqpoint{2.910462in}{0.795395in}}%
\pgfpathlineto{\pgfqpoint{2.894796in}{0.794134in}}%
\pgfpathlineto{\pgfqpoint{2.891131in}{0.799434in}}%
\pgfpathlineto{\pgfqpoint{2.885535in}{0.800050in}}%
\pgfpathlineto{\pgfqpoint{2.884962in}{0.810094in}}%
\pgfpathlineto{\pgfqpoint{2.887379in}{0.811341in}}%
\pgfpathlineto{\pgfqpoint{2.884148in}{0.822135in}}%
\pgfpathlineto{\pgfqpoint{2.904439in}{0.824809in}}%
\pgfpathlineto{\pgfqpoint{2.917347in}{0.821092in}}%
\pgfpathlineto{\pgfqpoint{2.918401in}{0.811582in}}%
\pgfpathlineto{\pgfqpoint{2.914876in}{0.807165in}}%
\pgfpathlineto{\pgfqpoint{2.912735in}{0.796174in}}%
\pgfpathlineto{\pgfqpoint{2.910462in}{0.795395in}}%
\pgfpathclose%
\pgfusepath{fill}%
\end{pgfscope}%
\begin{pgfscope}%
\pgfpathrectangle{\pgfqpoint{0.100000in}{0.100000in}}{\pgfqpoint{3.608454in}{2.310000in}}%
\pgfusepath{clip}%
\pgfsetbuttcap%
\pgfsetmiterjoin%
\definecolor{currentfill}{rgb}{0.000000,0.788235,0.605882}%
\pgfsetfillcolor{currentfill}%
\pgfsetlinewidth{0.000000pt}%
\definecolor{currentstroke}{rgb}{0.000000,0.000000,0.000000}%
\pgfsetstrokecolor{currentstroke}%
\pgfsetstrokeopacity{0.000000}%
\pgfsetdash{}{0pt}%
\pgfpathmoveto{\pgfqpoint{0.465159in}{1.900474in}}%
\pgfpathlineto{\pgfqpoint{0.478228in}{1.889504in}}%
\pgfpathlineto{\pgfqpoint{0.483261in}{1.890283in}}%
\pgfpathlineto{\pgfqpoint{0.491499in}{1.886385in}}%
\pgfpathlineto{\pgfqpoint{0.495724in}{1.888283in}}%
\pgfpathlineto{\pgfqpoint{0.503506in}{1.885466in}}%
\pgfpathlineto{\pgfqpoint{0.516046in}{1.883757in}}%
\pgfpathlineto{\pgfqpoint{0.532832in}{1.891746in}}%
\pgfpathlineto{\pgfqpoint{0.538396in}{1.890064in}}%
\pgfpathlineto{\pgfqpoint{0.543027in}{1.893318in}}%
\pgfpathlineto{\pgfqpoint{0.549594in}{1.891292in}}%
\pgfpathlineto{\pgfqpoint{0.526008in}{1.816628in}}%
\pgfpathlineto{\pgfqpoint{0.456983in}{1.837942in}}%
\pgfpathlineto{\pgfqpoint{0.439878in}{1.843077in}}%
\pgfpathlineto{\pgfqpoint{0.441504in}{1.854677in}}%
\pgfpathlineto{\pgfqpoint{0.447509in}{1.860643in}}%
\pgfpathlineto{\pgfqpoint{0.444876in}{1.871177in}}%
\pgfpathlineto{\pgfqpoint{0.437497in}{1.874169in}}%
\pgfpathlineto{\pgfqpoint{0.441915in}{1.883291in}}%
\pgfpathlineto{\pgfqpoint{0.448521in}{1.882309in}}%
\pgfpathlineto{\pgfqpoint{0.455038in}{1.891185in}}%
\pgfpathlineto{\pgfqpoint{0.465159in}{1.900474in}}%
\pgfpathclose%
\pgfusepath{fill}%
\end{pgfscope}%
\begin{pgfscope}%
\pgfpathrectangle{\pgfqpoint{0.100000in}{0.100000in}}{\pgfqpoint{3.608454in}{2.310000in}}%
\pgfusepath{clip}%
\pgfsetbuttcap%
\pgfsetmiterjoin%
\definecolor{currentfill}{rgb}{0.000000,0.662745,0.668627}%
\pgfsetfillcolor{currentfill}%
\pgfsetlinewidth{0.000000pt}%
\definecolor{currentstroke}{rgb}{0.000000,0.000000,0.000000}%
\pgfsetstrokecolor{currentstroke}%
\pgfsetstrokeopacity{0.000000}%
\pgfsetdash{}{0pt}%
\pgfpathmoveto{\pgfqpoint{2.734886in}{0.875470in}}%
\pgfpathlineto{\pgfqpoint{2.728245in}{0.872298in}}%
\pgfpathlineto{\pgfqpoint{2.729130in}{0.864274in}}%
\pgfpathlineto{\pgfqpoint{2.722392in}{0.862394in}}%
\pgfpathlineto{\pgfqpoint{2.723195in}{0.855536in}}%
\pgfpathlineto{\pgfqpoint{2.717264in}{0.854842in}}%
\pgfpathlineto{\pgfqpoint{2.715600in}{0.875575in}}%
\pgfpathlineto{\pgfqpoint{2.682505in}{0.872318in}}%
\pgfpathlineto{\pgfqpoint{2.676747in}{0.876531in}}%
\pgfpathlineto{\pgfqpoint{2.672768in}{0.884091in}}%
\pgfpathlineto{\pgfqpoint{2.670991in}{0.904805in}}%
\pgfpathlineto{\pgfqpoint{2.673491in}{0.910352in}}%
\pgfpathlineto{\pgfqpoint{2.679726in}{0.915262in}}%
\pgfpathlineto{\pgfqpoint{2.677615in}{0.922675in}}%
\pgfpathlineto{\pgfqpoint{2.688862in}{0.935943in}}%
\pgfpathlineto{\pgfqpoint{2.686513in}{0.939240in}}%
\pgfpathlineto{\pgfqpoint{2.690775in}{0.946928in}}%
\pgfpathlineto{\pgfqpoint{2.697028in}{0.947184in}}%
\pgfpathlineto{\pgfqpoint{2.701405in}{0.940587in}}%
\pgfpathlineto{\pgfqpoint{2.714599in}{0.941836in}}%
\pgfpathlineto{\pgfqpoint{2.717216in}{0.935165in}}%
\pgfpathlineto{\pgfqpoint{2.725270in}{0.935776in}}%
\pgfpathlineto{\pgfqpoint{2.727775in}{0.904950in}}%
\pgfpathlineto{\pgfqpoint{2.731773in}{0.905420in}}%
\pgfpathlineto{\pgfqpoint{2.734886in}{0.875470in}}%
\pgfpathclose%
\pgfusepath{fill}%
\end{pgfscope}%
\begin{pgfscope}%
\pgfpathrectangle{\pgfqpoint{0.100000in}{0.100000in}}{\pgfqpoint{3.608454in}{2.310000in}}%
\pgfusepath{clip}%
\pgfsetbuttcap%
\pgfsetmiterjoin%
\definecolor{currentfill}{rgb}{0.000000,0.533333,0.733333}%
\pgfsetfillcolor{currentfill}%
\pgfsetlinewidth{0.000000pt}%
\definecolor{currentstroke}{rgb}{0.000000,0.000000,0.000000}%
\pgfsetstrokecolor{currentstroke}%
\pgfsetstrokeopacity{0.000000}%
\pgfsetdash{}{0pt}%
\pgfpathmoveto{\pgfqpoint{2.460863in}{1.433950in}}%
\pgfpathlineto{\pgfqpoint{2.459428in}{1.452340in}}%
\pgfpathlineto{\pgfqpoint{2.452287in}{1.455302in}}%
\pgfpathlineto{\pgfqpoint{2.450616in}{1.476514in}}%
\pgfpathlineto{\pgfqpoint{2.458726in}{1.477228in}}%
\pgfpathlineto{\pgfqpoint{2.467299in}{1.483195in}}%
\pgfpathlineto{\pgfqpoint{2.470141in}{1.490384in}}%
\pgfpathlineto{\pgfqpoint{2.468185in}{1.518191in}}%
\pgfpathlineto{\pgfqpoint{2.488674in}{1.519766in}}%
\pgfpathlineto{\pgfqpoint{2.508522in}{1.521688in}}%
\pgfpathlineto{\pgfqpoint{2.509486in}{1.512307in}}%
\pgfpathlineto{\pgfqpoint{2.512434in}{1.482660in}}%
\pgfpathlineto{\pgfqpoint{2.498999in}{1.481611in}}%
\pgfpathlineto{\pgfqpoint{2.500893in}{1.455137in}}%
\pgfpathlineto{\pgfqpoint{2.493952in}{1.454636in}}%
\pgfpathlineto{\pgfqpoint{2.484922in}{1.436043in}}%
\pgfpathlineto{\pgfqpoint{2.460863in}{1.433950in}}%
\pgfpathclose%
\pgfusepath{fill}%
\end{pgfscope}%
\begin{pgfscope}%
\pgfpathrectangle{\pgfqpoint{0.100000in}{0.100000in}}{\pgfqpoint{3.608454in}{2.310000in}}%
\pgfusepath{clip}%
\pgfsetbuttcap%
\pgfsetmiterjoin%
\definecolor{currentfill}{rgb}{0.000000,0.643137,0.678431}%
\pgfsetfillcolor{currentfill}%
\pgfsetlinewidth{0.000000pt}%
\definecolor{currentstroke}{rgb}{0.000000,0.000000,0.000000}%
\pgfsetstrokecolor{currentstroke}%
\pgfsetstrokeopacity{0.000000}%
\pgfsetdash{}{0pt}%
\pgfpathmoveto{\pgfqpoint{2.322557in}{0.730252in}}%
\pgfpathlineto{\pgfqpoint{2.333480in}{0.724393in}}%
\pgfpathlineto{\pgfqpoint{2.330790in}{0.720030in}}%
\pgfpathlineto{\pgfqpoint{2.330722in}{0.712813in}}%
\pgfpathlineto{\pgfqpoint{2.325328in}{0.707042in}}%
\pgfpathlineto{\pgfqpoint{2.325022in}{0.696944in}}%
\pgfpathlineto{\pgfqpoint{2.296084in}{0.696005in}}%
\pgfpathlineto{\pgfqpoint{2.292928in}{0.702917in}}%
\pgfpathlineto{\pgfqpoint{2.285347in}{0.707925in}}%
\pgfpathlineto{\pgfqpoint{2.278789in}{0.700272in}}%
\pgfpathlineto{\pgfqpoint{2.271317in}{0.699125in}}%
\pgfpathlineto{\pgfqpoint{2.255860in}{0.698248in}}%
\pgfpathlineto{\pgfqpoint{2.254758in}{0.727139in}}%
\pgfpathlineto{\pgfqpoint{2.251222in}{0.732248in}}%
\pgfpathlineto{\pgfqpoint{2.244244in}{0.734424in}}%
\pgfpathlineto{\pgfqpoint{2.248864in}{0.736251in}}%
\pgfpathlineto{\pgfqpoint{2.261771in}{0.748472in}}%
\pgfpathlineto{\pgfqpoint{2.256019in}{0.754662in}}%
\pgfpathlineto{\pgfqpoint{2.248851in}{0.756615in}}%
\pgfpathlineto{\pgfqpoint{2.244140in}{0.763307in}}%
\pgfpathlineto{\pgfqpoint{2.268095in}{0.763939in}}%
\pgfpathlineto{\pgfqpoint{2.267964in}{0.770915in}}%
\pgfpathlineto{\pgfqpoint{2.285332in}{0.771355in}}%
\pgfpathlineto{\pgfqpoint{2.286412in}{0.757594in}}%
\pgfpathlineto{\pgfqpoint{2.291967in}{0.750029in}}%
\pgfpathlineto{\pgfqpoint{2.299838in}{0.745415in}}%
\pgfpathlineto{\pgfqpoint{2.300756in}{0.740640in}}%
\pgfpathlineto{\pgfqpoint{2.311015in}{0.733955in}}%
\pgfpathlineto{\pgfqpoint{2.322557in}{0.730252in}}%
\pgfpathclose%
\pgfusepath{fill}%
\end{pgfscope}%
\begin{pgfscope}%
\pgfpathrectangle{\pgfqpoint{0.100000in}{0.100000in}}{\pgfqpoint{3.608454in}{2.310000in}}%
\pgfusepath{clip}%
\pgfsetbuttcap%
\pgfsetmiterjoin%
\definecolor{currentfill}{rgb}{0.000000,0.623529,0.688235}%
\pgfsetfillcolor{currentfill}%
\pgfsetlinewidth{0.000000pt}%
\definecolor{currentstroke}{rgb}{0.000000,0.000000,0.000000}%
\pgfsetstrokecolor{currentstroke}%
\pgfsetstrokeopacity{0.000000}%
\pgfsetdash{}{0pt}%
\pgfpathmoveto{\pgfqpoint{2.537345in}{1.268351in}}%
\pgfpathlineto{\pgfqpoint{2.534479in}{1.269453in}}%
\pgfpathlineto{\pgfqpoint{2.519677in}{1.268419in}}%
\pgfpathlineto{\pgfqpoint{2.517928in}{1.295771in}}%
\pgfpathlineto{\pgfqpoint{2.497373in}{1.294428in}}%
\pgfpathlineto{\pgfqpoint{2.495702in}{1.322010in}}%
\pgfpathlineto{\pgfqpoint{2.523132in}{1.323450in}}%
\pgfpathlineto{\pgfqpoint{2.521612in}{1.327886in}}%
\pgfpathlineto{\pgfqpoint{2.521388in}{1.342978in}}%
\pgfpathlineto{\pgfqpoint{2.542895in}{1.344839in}}%
\pgfpathlineto{\pgfqpoint{2.544363in}{1.322774in}}%
\pgfpathlineto{\pgfqpoint{2.542054in}{1.316771in}}%
\pgfpathlineto{\pgfqpoint{2.541523in}{1.307274in}}%
\pgfpathlineto{\pgfqpoint{2.543804in}{1.301922in}}%
\pgfpathlineto{\pgfqpoint{2.541416in}{1.297901in}}%
\pgfpathlineto{\pgfqpoint{2.545025in}{1.287628in}}%
\pgfpathlineto{\pgfqpoint{2.537345in}{1.268351in}}%
\pgfpathclose%
\pgfusepath{fill}%
\end{pgfscope}%
\begin{pgfscope}%
\pgfpathrectangle{\pgfqpoint{0.100000in}{0.100000in}}{\pgfqpoint{3.608454in}{2.310000in}}%
\pgfusepath{clip}%
\pgfsetbuttcap%
\pgfsetmiterjoin%
\definecolor{currentfill}{rgb}{0.000000,0.741176,0.629412}%
\pgfsetfillcolor{currentfill}%
\pgfsetlinewidth{0.000000pt}%
\definecolor{currentstroke}{rgb}{0.000000,0.000000,0.000000}%
\pgfsetstrokecolor{currentstroke}%
\pgfsetstrokeopacity{0.000000}%
\pgfsetdash{}{0pt}%
\pgfpathmoveto{\pgfqpoint{2.920044in}{1.373705in}}%
\pgfpathlineto{\pgfqpoint{2.909141in}{1.370242in}}%
\pgfpathlineto{\pgfqpoint{2.903048in}{1.365422in}}%
\pgfpathlineto{\pgfqpoint{2.905122in}{1.355064in}}%
\pgfpathlineto{\pgfqpoint{2.893788in}{1.362802in}}%
\pgfpathlineto{\pgfqpoint{2.885129in}{1.361064in}}%
\pgfpathlineto{\pgfqpoint{2.884573in}{1.368240in}}%
\pgfpathlineto{\pgfqpoint{2.877293in}{1.367722in}}%
\pgfpathlineto{\pgfqpoint{2.875638in}{1.374674in}}%
\pgfpathlineto{\pgfqpoint{2.869516in}{1.374371in}}%
\pgfpathlineto{\pgfqpoint{2.869245in}{1.379552in}}%
\pgfpathlineto{\pgfqpoint{2.876585in}{1.380575in}}%
\pgfpathlineto{\pgfqpoint{2.875723in}{1.395664in}}%
\pgfpathlineto{\pgfqpoint{2.882640in}{1.396033in}}%
\pgfpathlineto{\pgfqpoint{2.881668in}{1.410091in}}%
\pgfpathlineto{\pgfqpoint{2.883523in}{1.417325in}}%
\pgfpathlineto{\pgfqpoint{2.880748in}{1.424222in}}%
\pgfpathlineto{\pgfqpoint{2.880330in}{1.431219in}}%
\pgfpathlineto{\pgfqpoint{2.887216in}{1.431682in}}%
\pgfpathlineto{\pgfqpoint{2.886708in}{1.438806in}}%
\pgfpathlineto{\pgfqpoint{2.893522in}{1.439324in}}%
\pgfpathlineto{\pgfqpoint{2.894009in}{1.432157in}}%
\pgfpathlineto{\pgfqpoint{2.906878in}{1.433220in}}%
\pgfpathlineto{\pgfqpoint{2.907865in}{1.422811in}}%
\pgfpathlineto{\pgfqpoint{2.916394in}{1.420191in}}%
\pgfpathlineto{\pgfqpoint{2.915344in}{1.412906in}}%
\pgfpathlineto{\pgfqpoint{2.916733in}{1.399937in}}%
\pgfpathlineto{\pgfqpoint{2.910035in}{1.386821in}}%
\pgfpathlineto{\pgfqpoint{2.920044in}{1.373705in}}%
\pgfpathclose%
\pgfusepath{fill}%
\end{pgfscope}%
\begin{pgfscope}%
\pgfpathrectangle{\pgfqpoint{0.100000in}{0.100000in}}{\pgfqpoint{3.608454in}{2.310000in}}%
\pgfusepath{clip}%
\pgfsetbuttcap%
\pgfsetmiterjoin%
\definecolor{currentfill}{rgb}{0.000000,0.443137,0.778431}%
\pgfsetfillcolor{currentfill}%
\pgfsetlinewidth{0.000000pt}%
\definecolor{currentstroke}{rgb}{0.000000,0.000000,0.000000}%
\pgfsetstrokecolor{currentstroke}%
\pgfsetstrokeopacity{0.000000}%
\pgfsetdash{}{0pt}%
\pgfpathmoveto{\pgfqpoint{1.893270in}{1.388628in}}%
\pgfpathlineto{\pgfqpoint{1.928379in}{1.387333in}}%
\pgfpathlineto{\pgfqpoint{1.927666in}{1.366756in}}%
\pgfpathlineto{\pgfqpoint{1.927400in}{1.359864in}}%
\pgfpathlineto{\pgfqpoint{1.893147in}{1.361110in}}%
\pgfpathlineto{\pgfqpoint{1.892889in}{1.354245in}}%
\pgfpathlineto{\pgfqpoint{1.858721in}{1.355685in}}%
\pgfpathlineto{\pgfqpoint{1.860470in}{1.390012in}}%
\pgfpathlineto{\pgfqpoint{1.893270in}{1.388628in}}%
\pgfpathclose%
\pgfusepath{fill}%
\end{pgfscope}%
\begin{pgfscope}%
\pgfpathrectangle{\pgfqpoint{0.100000in}{0.100000in}}{\pgfqpoint{3.608454in}{2.310000in}}%
\pgfusepath{clip}%
\pgfsetbuttcap%
\pgfsetmiterjoin%
\definecolor{currentfill}{rgb}{0.000000,0.341176,0.829412}%
\pgfsetfillcolor{currentfill}%
\pgfsetlinewidth{0.000000pt}%
\definecolor{currentstroke}{rgb}{0.000000,0.000000,0.000000}%
\pgfsetstrokecolor{currentstroke}%
\pgfsetstrokeopacity{0.000000}%
\pgfsetdash{}{0pt}%
\pgfpathmoveto{\pgfqpoint{3.214106in}{1.265774in}}%
\pgfpathlineto{\pgfqpoint{3.210178in}{1.266083in}}%
\pgfpathlineto{\pgfqpoint{3.209973in}{1.277484in}}%
\pgfpathlineto{\pgfqpoint{3.202830in}{1.275174in}}%
\pgfpathlineto{\pgfqpoint{3.197116in}{1.277528in}}%
\pgfpathlineto{\pgfqpoint{3.189244in}{1.282467in}}%
\pgfpathlineto{\pgfqpoint{3.181322in}{1.282092in}}%
\pgfpathlineto{\pgfqpoint{3.178373in}{1.294019in}}%
\pgfpathlineto{\pgfqpoint{3.184110in}{1.299032in}}%
\pgfpathlineto{\pgfqpoint{3.155665in}{1.301994in}}%
\pgfpathlineto{\pgfqpoint{3.154386in}{1.307543in}}%
\pgfpathlineto{\pgfqpoint{3.159486in}{1.315467in}}%
\pgfpathlineto{\pgfqpoint{3.159782in}{1.331466in}}%
\pgfpathlineto{\pgfqpoint{3.153524in}{1.337961in}}%
\pgfpathlineto{\pgfqpoint{3.157131in}{1.351262in}}%
\pgfpathlineto{\pgfqpoint{3.165049in}{1.347998in}}%
\pgfpathlineto{\pgfqpoint{3.169983in}{1.342070in}}%
\pgfpathlineto{\pgfqpoint{3.176043in}{1.340697in}}%
\pgfpathlineto{\pgfqpoint{3.178627in}{1.363629in}}%
\pgfpathlineto{\pgfqpoint{3.195177in}{1.357724in}}%
\pgfpathlineto{\pgfqpoint{3.197288in}{1.352635in}}%
\pgfpathlineto{\pgfqpoint{3.202566in}{1.350706in}}%
\pgfpathlineto{\pgfqpoint{3.210448in}{1.365277in}}%
\pgfpathlineto{\pgfqpoint{3.217018in}{1.367687in}}%
\pgfpathlineto{\pgfqpoint{3.222015in}{1.359018in}}%
\pgfpathlineto{\pgfqpoint{3.228182in}{1.355944in}}%
\pgfpathlineto{\pgfqpoint{3.236077in}{1.357801in}}%
\pgfpathlineto{\pgfqpoint{3.240232in}{1.353321in}}%
\pgfpathlineto{\pgfqpoint{3.248164in}{1.344820in}}%
\pgfpathlineto{\pgfqpoint{3.249638in}{1.333200in}}%
\pgfpathlineto{\pgfqpoint{3.236059in}{1.328225in}}%
\pgfpathlineto{\pgfqpoint{3.240056in}{1.312499in}}%
\pgfpathlineto{\pgfqpoint{3.234147in}{1.313361in}}%
\pgfpathlineto{\pgfqpoint{3.224896in}{1.300661in}}%
\pgfpathlineto{\pgfqpoint{3.236491in}{1.297792in}}%
\pgfpathlineto{\pgfqpoint{3.240057in}{1.290021in}}%
\pgfpathlineto{\pgfqpoint{3.214106in}{1.265774in}}%
\pgfpathclose%
\pgfusepath{fill}%
\end{pgfscope}%
\begin{pgfscope}%
\pgfpathrectangle{\pgfqpoint{0.100000in}{0.100000in}}{\pgfqpoint{3.608454in}{2.310000in}}%
\pgfusepath{clip}%
\pgfsetbuttcap%
\pgfsetmiterjoin%
\definecolor{currentfill}{rgb}{0.000000,0.603922,0.698039}%
\pgfsetfillcolor{currentfill}%
\pgfsetlinewidth{0.000000pt}%
\definecolor{currentstroke}{rgb}{0.000000,0.000000,0.000000}%
\pgfsetstrokecolor{currentstroke}%
\pgfsetstrokeopacity{0.000000}%
\pgfsetdash{}{0pt}%
\pgfpathmoveto{\pgfqpoint{1.997875in}{0.723825in}}%
\pgfpathlineto{\pgfqpoint{1.981851in}{0.714949in}}%
\pgfpathlineto{\pgfqpoint{1.965169in}{0.705757in}}%
\pgfpathlineto{\pgfqpoint{1.951524in}{0.729506in}}%
\pgfpathlineto{\pgfqpoint{1.947012in}{0.726789in}}%
\pgfpathlineto{\pgfqpoint{1.942049in}{0.732981in}}%
\pgfpathlineto{\pgfqpoint{1.924502in}{0.764640in}}%
\pgfpathlineto{\pgfqpoint{1.919219in}{0.761725in}}%
\pgfpathlineto{\pgfqpoint{1.903937in}{0.789730in}}%
\pgfpathlineto{\pgfqpoint{1.913591in}{0.794961in}}%
\pgfpathlineto{\pgfqpoint{1.929646in}{0.804003in}}%
\pgfpathlineto{\pgfqpoint{1.934684in}{0.799256in}}%
\pgfpathlineto{\pgfqpoint{1.940012in}{0.801698in}}%
\pgfpathlineto{\pgfqpoint{1.966416in}{0.807704in}}%
\pgfpathlineto{\pgfqpoint{1.975845in}{0.790506in}}%
\pgfpathlineto{\pgfqpoint{1.978879in}{0.792186in}}%
\pgfpathlineto{\pgfqpoint{1.990434in}{0.771318in}}%
\pgfpathlineto{\pgfqpoint{1.975832in}{0.763164in}}%
\pgfpathlineto{\pgfqpoint{1.997875in}{0.723825in}}%
\pgfpathclose%
\pgfusepath{fill}%
\end{pgfscope}%
\begin{pgfscope}%
\pgfpathrectangle{\pgfqpoint{0.100000in}{0.100000in}}{\pgfqpoint{3.608454in}{2.310000in}}%
\pgfusepath{clip}%
\pgfsetbuttcap%
\pgfsetmiterjoin%
\definecolor{currentfill}{rgb}{0.000000,0.592157,0.703922}%
\pgfsetfillcolor{currentfill}%
\pgfsetlinewidth{0.000000pt}%
\definecolor{currentstroke}{rgb}{0.000000,0.000000,0.000000}%
\pgfsetstrokecolor{currentstroke}%
\pgfsetstrokeopacity{0.000000}%
\pgfsetdash{}{0pt}%
\pgfpathmoveto{\pgfqpoint{2.978628in}{1.415804in}}%
\pgfpathlineto{\pgfqpoint{2.973331in}{1.416066in}}%
\pgfpathlineto{\pgfqpoint{2.965224in}{1.427879in}}%
\pgfpathlineto{\pgfqpoint{2.966480in}{1.429783in}}%
\pgfpathlineto{\pgfqpoint{2.980855in}{1.444545in}}%
\pgfpathlineto{\pgfqpoint{2.987075in}{1.444310in}}%
\pgfpathlineto{\pgfqpoint{2.989302in}{1.450132in}}%
\pgfpathlineto{\pgfqpoint{2.986701in}{1.454490in}}%
\pgfpathlineto{\pgfqpoint{2.992195in}{1.461918in}}%
\pgfpathlineto{\pgfqpoint{2.989320in}{1.467586in}}%
\pgfpathlineto{\pgfqpoint{3.045971in}{1.477046in}}%
\pgfpathlineto{\pgfqpoint{3.052407in}{1.435951in}}%
\pgfpathlineto{\pgfqpoint{3.039311in}{1.439843in}}%
\pgfpathlineto{\pgfqpoint{3.032463in}{1.435406in}}%
\pgfpathlineto{\pgfqpoint{3.025784in}{1.440073in}}%
\pgfpathlineto{\pgfqpoint{3.019335in}{1.434599in}}%
\pgfpathlineto{\pgfqpoint{3.010773in}{1.432733in}}%
\pgfpathlineto{\pgfqpoint{3.008758in}{1.422029in}}%
\pgfpathlineto{\pgfqpoint{3.004458in}{1.420613in}}%
\pgfpathlineto{\pgfqpoint{2.986872in}{1.422883in}}%
\pgfpathlineto{\pgfqpoint{2.978628in}{1.415804in}}%
\pgfpathclose%
\pgfusepath{fill}%
\end{pgfscope}%
\begin{pgfscope}%
\pgfpathrectangle{\pgfqpoint{0.100000in}{0.100000in}}{\pgfqpoint{3.608454in}{2.310000in}}%
\pgfusepath{clip}%
\pgfsetbuttcap%
\pgfsetmiterjoin%
\definecolor{currentfill}{rgb}{0.000000,0.607843,0.696078}%
\pgfsetfillcolor{currentfill}%
\pgfsetlinewidth{0.000000pt}%
\definecolor{currentstroke}{rgb}{0.000000,0.000000,0.000000}%
\pgfsetstrokecolor{currentstroke}%
\pgfsetstrokeopacity{0.000000}%
\pgfsetdash{}{0pt}%
\pgfpathmoveto{\pgfqpoint{3.192630in}{1.620121in}}%
\pgfpathlineto{\pgfqpoint{3.195360in}{1.613257in}}%
\pgfpathlineto{\pgfqpoint{3.203466in}{1.605423in}}%
\pgfpathlineto{\pgfqpoint{3.206320in}{1.596141in}}%
\pgfpathlineto{\pgfqpoint{3.211501in}{1.598059in}}%
\pgfpathlineto{\pgfqpoint{3.213947in}{1.593817in}}%
\pgfpathlineto{\pgfqpoint{3.196959in}{1.580835in}}%
\pgfpathlineto{\pgfqpoint{3.184525in}{1.574369in}}%
\pgfpathlineto{\pgfqpoint{3.176843in}{1.578501in}}%
\pgfpathlineto{\pgfqpoint{3.157897in}{1.576523in}}%
\pgfpathlineto{\pgfqpoint{3.155139in}{1.587716in}}%
\pgfpathlineto{\pgfqpoint{3.165006in}{1.605498in}}%
\pgfpathlineto{\pgfqpoint{3.173979in}{1.610853in}}%
\pgfpathlineto{\pgfqpoint{3.174106in}{1.615801in}}%
\pgfpathlineto{\pgfqpoint{3.183685in}{1.619914in}}%
\pgfpathlineto{\pgfqpoint{3.192630in}{1.620121in}}%
\pgfpathclose%
\pgfusepath{fill}%
\end{pgfscope}%
\begin{pgfscope}%
\pgfpathrectangle{\pgfqpoint{0.100000in}{0.100000in}}{\pgfqpoint{3.608454in}{2.310000in}}%
\pgfusepath{clip}%
\pgfsetbuttcap%
\pgfsetmiterjoin%
\definecolor{currentfill}{rgb}{0.000000,0.784314,0.607843}%
\pgfsetfillcolor{currentfill}%
\pgfsetlinewidth{0.000000pt}%
\definecolor{currentstroke}{rgb}{0.000000,0.000000,0.000000}%
\pgfsetstrokecolor{currentstroke}%
\pgfsetstrokeopacity{0.000000}%
\pgfsetdash{}{0pt}%
\pgfpathmoveto{\pgfqpoint{2.549469in}{0.776821in}}%
\pgfpathlineto{\pgfqpoint{2.547754in}{0.792351in}}%
\pgfpathlineto{\pgfqpoint{2.549490in}{0.847081in}}%
\pgfpathlineto{\pgfqpoint{2.550184in}{0.875112in}}%
\pgfpathlineto{\pgfqpoint{2.550318in}{0.880073in}}%
\pgfpathlineto{\pgfqpoint{2.561497in}{0.881275in}}%
\pgfpathlineto{\pgfqpoint{2.564922in}{0.885225in}}%
\pgfpathlineto{\pgfqpoint{2.570549in}{0.885685in}}%
\pgfpathlineto{\pgfqpoint{2.575754in}{0.891315in}}%
\pgfpathlineto{\pgfqpoint{2.582729in}{0.895488in}}%
\pgfpathlineto{\pgfqpoint{2.583933in}{0.884712in}}%
\pgfpathlineto{\pgfqpoint{2.591982in}{0.885391in}}%
\pgfpathlineto{\pgfqpoint{2.591160in}{0.879866in}}%
\pgfpathlineto{\pgfqpoint{2.586494in}{0.875360in}}%
\pgfpathlineto{\pgfqpoint{2.586111in}{0.864841in}}%
\pgfpathlineto{\pgfqpoint{2.589627in}{0.852306in}}%
\pgfpathlineto{\pgfqpoint{2.588335in}{0.845764in}}%
\pgfpathlineto{\pgfqpoint{2.594215in}{0.842750in}}%
\pgfpathlineto{\pgfqpoint{2.607870in}{0.843949in}}%
\pgfpathlineto{\pgfqpoint{2.609026in}{0.830093in}}%
\pgfpathlineto{\pgfqpoint{2.612796in}{0.827022in}}%
\pgfpathlineto{\pgfqpoint{2.609428in}{0.825937in}}%
\pgfpathlineto{\pgfqpoint{2.610325in}{0.816266in}}%
\pgfpathlineto{\pgfqpoint{2.603471in}{0.815667in}}%
\pgfpathlineto{\pgfqpoint{2.604373in}{0.805606in}}%
\pgfpathlineto{\pgfqpoint{2.601434in}{0.804239in}}%
\pgfpathlineto{\pgfqpoint{2.574126in}{0.801999in}}%
\pgfpathlineto{\pgfqpoint{2.567924in}{0.787803in}}%
\pgfpathlineto{\pgfqpoint{2.573394in}{0.783619in}}%
\pgfpathlineto{\pgfqpoint{2.574889in}{0.778819in}}%
\pgfpathlineto{\pgfqpoint{2.549469in}{0.776821in}}%
\pgfpathclose%
\pgfusepath{fill}%
\end{pgfscope}%
\begin{pgfscope}%
\pgfpathrectangle{\pgfqpoint{0.100000in}{0.100000in}}{\pgfqpoint{3.608454in}{2.310000in}}%
\pgfusepath{clip}%
\pgfsetbuttcap%
\pgfsetmiterjoin%
\definecolor{currentfill}{rgb}{0.000000,0.800000,0.600000}%
\pgfsetfillcolor{currentfill}%
\pgfsetlinewidth{0.000000pt}%
\definecolor{currentstroke}{rgb}{0.000000,0.000000,0.000000}%
\pgfsetstrokecolor{currentstroke}%
\pgfsetstrokeopacity{0.000000}%
\pgfsetdash{}{0pt}%
\pgfpathmoveto{\pgfqpoint{2.361668in}{0.958818in}}%
\pgfpathlineto{\pgfqpoint{2.356734in}{0.953261in}}%
\pgfpathlineto{\pgfqpoint{2.361253in}{0.945877in}}%
\pgfpathlineto{\pgfqpoint{2.349626in}{0.945687in}}%
\pgfpathlineto{\pgfqpoint{2.341075in}{0.950676in}}%
\pgfpathlineto{\pgfqpoint{2.339665in}{0.959110in}}%
\pgfpathlineto{\pgfqpoint{2.335508in}{0.959654in}}%
\pgfpathlineto{\pgfqpoint{2.333604in}{0.966070in}}%
\pgfpathlineto{\pgfqpoint{2.322365in}{0.966382in}}%
\pgfpathlineto{\pgfqpoint{2.322083in}{0.986111in}}%
\pgfpathlineto{\pgfqpoint{2.326278in}{0.991006in}}%
\pgfpathlineto{\pgfqpoint{2.326016in}{0.997543in}}%
\pgfpathlineto{\pgfqpoint{2.321823in}{1.001132in}}%
\pgfpathlineto{\pgfqpoint{2.321574in}{1.020760in}}%
\pgfpathlineto{\pgfqpoint{2.314414in}{1.020654in}}%
\pgfpathlineto{\pgfqpoint{2.314200in}{1.029274in}}%
\pgfpathlineto{\pgfqpoint{2.320051in}{1.032034in}}%
\pgfpathlineto{\pgfqpoint{2.328311in}{1.029359in}}%
\pgfpathlineto{\pgfqpoint{2.328181in}{1.034618in}}%
\pgfpathlineto{\pgfqpoint{2.335892in}{1.034757in}}%
\pgfpathlineto{\pgfqpoint{2.342574in}{1.032044in}}%
\pgfpathlineto{\pgfqpoint{2.347372in}{1.026940in}}%
\pgfpathlineto{\pgfqpoint{2.356606in}{1.028702in}}%
\pgfpathlineto{\pgfqpoint{2.356746in}{1.021522in}}%
\pgfpathlineto{\pgfqpoint{2.360312in}{1.018136in}}%
\pgfpathlineto{\pgfqpoint{2.360588in}{1.007650in}}%
\pgfpathlineto{\pgfqpoint{2.364049in}{1.007736in}}%
\pgfpathlineto{\pgfqpoint{2.364599in}{0.986870in}}%
\pgfpathlineto{\pgfqpoint{2.368151in}{0.986916in}}%
\pgfpathlineto{\pgfqpoint{2.368481in}{0.976354in}}%
\pgfpathlineto{\pgfqpoint{2.364963in}{0.976228in}}%
\pgfpathlineto{\pgfqpoint{2.363796in}{0.962251in}}%
\pgfpathlineto{\pgfqpoint{2.361668in}{0.958818in}}%
\pgfpathclose%
\pgfusepath{fill}%
\end{pgfscope}%
\begin{pgfscope}%
\pgfpathrectangle{\pgfqpoint{0.100000in}{0.100000in}}{\pgfqpoint{3.608454in}{2.310000in}}%
\pgfusepath{clip}%
\pgfsetbuttcap%
\pgfsetmiterjoin%
\definecolor{currentfill}{rgb}{0.000000,0.403922,0.798039}%
\pgfsetfillcolor{currentfill}%
\pgfsetlinewidth{0.000000pt}%
\definecolor{currentstroke}{rgb}{0.000000,0.000000,0.000000}%
\pgfsetstrokecolor{currentstroke}%
\pgfsetstrokeopacity{0.000000}%
\pgfsetdash{}{0pt}%
\pgfpathmoveto{\pgfqpoint{1.670049in}{1.298787in}}%
\pgfpathlineto{\pgfqpoint{1.626877in}{1.302163in}}%
\pgfpathlineto{\pgfqpoint{1.579920in}{1.306214in}}%
\pgfpathlineto{\pgfqpoint{1.581597in}{1.326016in}}%
\pgfpathlineto{\pgfqpoint{1.602094in}{1.324825in}}%
\pgfpathlineto{\pgfqpoint{1.603799in}{1.337764in}}%
\pgfpathlineto{\pgfqpoint{1.606968in}{1.372189in}}%
\pgfpathlineto{\pgfqpoint{1.610603in}{1.406638in}}%
\pgfpathlineto{\pgfqpoint{1.677720in}{1.401672in}}%
\pgfpathlineto{\pgfqpoint{1.677687in}{1.401208in}}%
\pgfpathlineto{\pgfqpoint{1.672634in}{1.332579in}}%
\pgfpathlineto{\pgfqpoint{1.670049in}{1.298787in}}%
\pgfpathclose%
\pgfusepath{fill}%
\end{pgfscope}%
\begin{pgfscope}%
\pgfpathrectangle{\pgfqpoint{0.100000in}{0.100000in}}{\pgfqpoint{3.608454in}{2.310000in}}%
\pgfusepath{clip}%
\pgfsetbuttcap%
\pgfsetmiterjoin%
\definecolor{currentfill}{rgb}{0.000000,0.760784,0.619608}%
\pgfsetfillcolor{currentfill}%
\pgfsetlinewidth{0.000000pt}%
\definecolor{currentstroke}{rgb}{0.000000,0.000000,0.000000}%
\pgfsetstrokecolor{currentstroke}%
\pgfsetstrokeopacity{0.000000}%
\pgfsetdash{}{0pt}%
\pgfpathmoveto{\pgfqpoint{2.396088in}{0.992607in}}%
\pgfpathlineto{\pgfqpoint{2.395789in}{0.983518in}}%
\pgfpathlineto{\pgfqpoint{2.384480in}{0.979541in}}%
\pgfpathlineto{\pgfqpoint{2.384214in}{0.972642in}}%
\pgfpathlineto{\pgfqpoint{2.372289in}{0.959362in}}%
\pgfpathlineto{\pgfqpoint{2.372346in}{0.959262in}}%
\pgfpathlineto{\pgfqpoint{2.361668in}{0.958818in}}%
\pgfpathlineto{\pgfqpoint{2.363796in}{0.962251in}}%
\pgfpathlineto{\pgfqpoint{2.364963in}{0.976228in}}%
\pgfpathlineto{\pgfqpoint{2.368481in}{0.976354in}}%
\pgfpathlineto{\pgfqpoint{2.368151in}{0.986916in}}%
\pgfpathlineto{\pgfqpoint{2.364599in}{0.986870in}}%
\pgfpathlineto{\pgfqpoint{2.364049in}{1.007736in}}%
\pgfpathlineto{\pgfqpoint{2.360588in}{1.007650in}}%
\pgfpathlineto{\pgfqpoint{2.360312in}{1.018136in}}%
\pgfpathlineto{\pgfqpoint{2.356746in}{1.021522in}}%
\pgfpathlineto{\pgfqpoint{2.356606in}{1.028702in}}%
\pgfpathlineto{\pgfqpoint{2.363393in}{1.028869in}}%
\pgfpathlineto{\pgfqpoint{2.362249in}{1.063873in}}%
\pgfpathlineto{\pgfqpoint{2.389852in}{1.065601in}}%
\pgfpathlineto{\pgfqpoint{2.410748in}{1.065840in}}%
\pgfpathlineto{\pgfqpoint{2.420217in}{1.066185in}}%
\pgfpathlineto{\pgfqpoint{2.424647in}{1.062261in}}%
\pgfpathlineto{\pgfqpoint{2.420989in}{1.051604in}}%
\pgfpathlineto{\pgfqpoint{2.425407in}{1.049147in}}%
\pgfpathlineto{\pgfqpoint{2.424139in}{1.041016in}}%
\pgfpathlineto{\pgfqpoint{2.419408in}{1.040872in}}%
\pgfpathlineto{\pgfqpoint{2.417154in}{1.033431in}}%
\pgfpathlineto{\pgfqpoint{2.412041in}{1.029240in}}%
\pgfpathlineto{\pgfqpoint{2.415331in}{1.023934in}}%
\pgfpathlineto{\pgfqpoint{2.404445in}{1.017531in}}%
\pgfpathlineto{\pgfqpoint{2.401872in}{1.004184in}}%
\pgfpathlineto{\pgfqpoint{2.394572in}{0.999746in}}%
\pgfpathlineto{\pgfqpoint{2.396088in}{0.992607in}}%
\pgfpathclose%
\pgfusepath{fill}%
\end{pgfscope}%
\begin{pgfscope}%
\pgfpathrectangle{\pgfqpoint{0.100000in}{0.100000in}}{\pgfqpoint{3.608454in}{2.310000in}}%
\pgfusepath{clip}%
\pgfsetbuttcap%
\pgfsetmiterjoin%
\definecolor{currentfill}{rgb}{0.000000,0.470588,0.764706}%
\pgfsetfillcolor{currentfill}%
\pgfsetlinewidth{0.000000pt}%
\definecolor{currentstroke}{rgb}{0.000000,0.000000,0.000000}%
\pgfsetstrokecolor{currentstroke}%
\pgfsetstrokeopacity{0.000000}%
\pgfsetdash{}{0pt}%
\pgfpathmoveto{\pgfqpoint{2.720376in}{1.354526in}}%
\pgfpathlineto{\pgfqpoint{2.710112in}{1.347801in}}%
\pgfpathlineto{\pgfqpoint{2.708117in}{1.365276in}}%
\pgfpathlineto{\pgfqpoint{2.693342in}{1.363630in}}%
\pgfpathlineto{\pgfqpoint{2.691293in}{1.385879in}}%
\pgfpathlineto{\pgfqpoint{2.684115in}{1.380130in}}%
\pgfpathlineto{\pgfqpoint{2.676781in}{1.379202in}}%
\pgfpathlineto{\pgfqpoint{2.675138in}{1.396451in}}%
\pgfpathlineto{\pgfqpoint{2.677647in}{1.404907in}}%
\pgfpathlineto{\pgfqpoint{2.697801in}{1.407041in}}%
\pgfpathlineto{\pgfqpoint{2.697158in}{1.412713in}}%
\pgfpathlineto{\pgfqpoint{2.726448in}{1.415659in}}%
\pgfpathlineto{\pgfqpoint{2.729836in}{1.382914in}}%
\pgfpathlineto{\pgfqpoint{2.725607in}{1.378219in}}%
\pgfpathlineto{\pgfqpoint{2.729795in}{1.373635in}}%
\pgfpathlineto{\pgfqpoint{2.727957in}{1.367391in}}%
\pgfpathlineto{\pgfqpoint{2.733971in}{1.365444in}}%
\pgfpathlineto{\pgfqpoint{2.728448in}{1.358041in}}%
\pgfpathlineto{\pgfqpoint{2.720376in}{1.354526in}}%
\pgfpathclose%
\pgfusepath{fill}%
\end{pgfscope}%
\begin{pgfscope}%
\pgfpathrectangle{\pgfqpoint{0.100000in}{0.100000in}}{\pgfqpoint{3.608454in}{2.310000in}}%
\pgfusepath{clip}%
\pgfsetbuttcap%
\pgfsetmiterjoin%
\definecolor{currentfill}{rgb}{0.000000,0.670588,0.664706}%
\pgfsetfillcolor{currentfill}%
\pgfsetlinewidth{0.000000pt}%
\definecolor{currentstroke}{rgb}{0.000000,0.000000,0.000000}%
\pgfsetstrokecolor{currentstroke}%
\pgfsetstrokeopacity{0.000000}%
\pgfsetdash{}{0pt}%
\pgfpathmoveto{\pgfqpoint{1.708835in}{0.943953in}}%
\pgfpathlineto{\pgfqpoint{1.743217in}{0.941925in}}%
\pgfpathlineto{\pgfqpoint{1.774163in}{0.940149in}}%
\pgfpathlineto{\pgfqpoint{1.797246in}{0.938094in}}%
\pgfpathlineto{\pgfqpoint{1.810392in}{0.932854in}}%
\pgfpathlineto{\pgfqpoint{1.811548in}{0.930044in}}%
\pgfpathlineto{\pgfqpoint{1.810377in}{0.903512in}}%
\pgfpathlineto{\pgfqpoint{1.808729in}{0.868432in}}%
\pgfpathlineto{\pgfqpoint{1.763926in}{0.871040in}}%
\pgfpathlineto{\pgfqpoint{1.738977in}{0.872740in}}%
\pgfpathlineto{\pgfqpoint{1.741190in}{0.907227in}}%
\pgfpathlineto{\pgfqpoint{1.706689in}{0.909371in}}%
\pgfpathlineto{\pgfqpoint{1.708835in}{0.943953in}}%
\pgfpathclose%
\pgfusepath{fill}%
\end{pgfscope}%
\begin{pgfscope}%
\pgfpathrectangle{\pgfqpoint{0.100000in}{0.100000in}}{\pgfqpoint{3.608454in}{2.310000in}}%
\pgfusepath{clip}%
\pgfsetbuttcap%
\pgfsetmiterjoin%
\definecolor{currentfill}{rgb}{0.000000,0.835294,0.582353}%
\pgfsetfillcolor{currentfill}%
\pgfsetlinewidth{0.000000pt}%
\definecolor{currentstroke}{rgb}{0.000000,0.000000,0.000000}%
\pgfsetstrokecolor{currentstroke}%
\pgfsetstrokeopacity{0.000000}%
\pgfsetdash{}{0pt}%
\pgfpathmoveto{\pgfqpoint{2.336789in}{0.709454in}}%
\pgfpathlineto{\pgfqpoint{2.330722in}{0.712813in}}%
\pgfpathlineto{\pgfqpoint{2.330790in}{0.720030in}}%
\pgfpathlineto{\pgfqpoint{2.333480in}{0.724393in}}%
\pgfpathlineto{\pgfqpoint{2.322557in}{0.730252in}}%
\pgfpathlineto{\pgfqpoint{2.319974in}{0.746106in}}%
\pgfpathlineto{\pgfqpoint{2.325362in}{0.751678in}}%
\pgfpathlineto{\pgfqpoint{2.330853in}{0.768113in}}%
\pgfpathlineto{\pgfqpoint{2.334485in}{0.767523in}}%
\pgfpathlineto{\pgfqpoint{2.338396in}{0.773253in}}%
\pgfpathlineto{\pgfqpoint{2.338322in}{0.779971in}}%
\pgfpathlineto{\pgfqpoint{2.343388in}{0.795089in}}%
\pgfpathlineto{\pgfqpoint{2.342946in}{0.805709in}}%
\pgfpathlineto{\pgfqpoint{2.354522in}{0.806095in}}%
\pgfpathlineto{\pgfqpoint{2.366944in}{0.810346in}}%
\pgfpathlineto{\pgfqpoint{2.365200in}{0.801106in}}%
\pgfpathlineto{\pgfqpoint{2.374090in}{0.798791in}}%
\pgfpathlineto{\pgfqpoint{2.371635in}{0.792020in}}%
\pgfpathlineto{\pgfqpoint{2.364252in}{0.787285in}}%
\pgfpathlineto{\pgfqpoint{2.363437in}{0.782994in}}%
\pgfpathlineto{\pgfqpoint{2.353993in}{0.777417in}}%
\pgfpathlineto{\pgfqpoint{2.357930in}{0.763773in}}%
\pgfpathlineto{\pgfqpoint{2.367747in}{0.759421in}}%
\pgfpathlineto{\pgfqpoint{2.395986in}{0.760766in}}%
\pgfpathlineto{\pgfqpoint{2.403042in}{0.761124in}}%
\pgfpathlineto{\pgfqpoint{2.404037in}{0.740243in}}%
\pgfpathlineto{\pgfqpoint{2.377876in}{0.738907in}}%
\pgfpathlineto{\pgfqpoint{2.371609in}{0.736436in}}%
\pgfpathlineto{\pgfqpoint{2.362838in}{0.740120in}}%
\pgfpathlineto{\pgfqpoint{2.347315in}{0.735413in}}%
\pgfpathlineto{\pgfqpoint{2.346226in}{0.727949in}}%
\pgfpathlineto{\pgfqpoint{2.341299in}{0.728481in}}%
\pgfpathlineto{\pgfqpoint{2.337261in}{0.720197in}}%
\pgfpathlineto{\pgfqpoint{2.341562in}{0.715016in}}%
\pgfpathlineto{\pgfqpoint{2.336789in}{0.709454in}}%
\pgfpathclose%
\pgfusepath{fill}%
\end{pgfscope}%
\begin{pgfscope}%
\pgfpathrectangle{\pgfqpoint{0.100000in}{0.100000in}}{\pgfqpoint{3.608454in}{2.310000in}}%
\pgfusepath{clip}%
\pgfsetbuttcap%
\pgfsetmiterjoin%
\definecolor{currentfill}{rgb}{0.000000,0.533333,0.733333}%
\pgfsetfillcolor{currentfill}%
\pgfsetlinewidth{0.000000pt}%
\definecolor{currentstroke}{rgb}{0.000000,0.000000,0.000000}%
\pgfsetstrokecolor{currentstroke}%
\pgfsetstrokeopacity{0.000000}%
\pgfsetdash{}{0pt}%
\pgfpathmoveto{\pgfqpoint{2.368642in}{1.363523in}}%
\pgfpathlineto{\pgfqpoint{2.340427in}{1.362513in}}%
\pgfpathlineto{\pgfqpoint{2.340562in}{1.355530in}}%
\pgfpathlineto{\pgfqpoint{2.335984in}{1.355415in}}%
\pgfpathlineto{\pgfqpoint{2.326861in}{1.355261in}}%
\pgfpathlineto{\pgfqpoint{2.324527in}{1.369104in}}%
\pgfpathlineto{\pgfqpoint{2.307376in}{1.370072in}}%
\pgfpathlineto{\pgfqpoint{2.306650in}{1.395239in}}%
\pgfpathlineto{\pgfqpoint{2.298632in}{1.395025in}}%
\pgfpathlineto{\pgfqpoint{2.298225in}{1.417924in}}%
\pgfpathlineto{\pgfqpoint{2.291417in}{1.417786in}}%
\pgfpathlineto{\pgfqpoint{2.290718in}{1.442111in}}%
\pgfpathlineto{\pgfqpoint{2.318087in}{1.442426in}}%
\pgfpathlineto{\pgfqpoint{2.317746in}{1.438593in}}%
\pgfpathlineto{\pgfqpoint{2.353565in}{1.439495in}}%
\pgfpathlineto{\pgfqpoint{2.354524in}{1.412008in}}%
\pgfpathlineto{\pgfqpoint{2.375453in}{1.412534in}}%
\pgfpathlineto{\pgfqpoint{2.371479in}{1.401180in}}%
\pgfpathlineto{\pgfqpoint{2.376576in}{1.388834in}}%
\pgfpathlineto{\pgfqpoint{2.374602in}{1.377391in}}%
\pgfpathlineto{\pgfqpoint{2.354879in}{1.376769in}}%
\pgfpathlineto{\pgfqpoint{2.363949in}{1.370060in}}%
\pgfpathlineto{\pgfqpoint{2.368642in}{1.363523in}}%
\pgfpathclose%
\pgfusepath{fill}%
\end{pgfscope}%
\begin{pgfscope}%
\pgfpathrectangle{\pgfqpoint{0.100000in}{0.100000in}}{\pgfqpoint{3.608454in}{2.310000in}}%
\pgfusepath{clip}%
\pgfsetbuttcap%
\pgfsetmiterjoin%
\definecolor{currentfill}{rgb}{0.000000,0.858824,0.570588}%
\pgfsetfillcolor{currentfill}%
\pgfsetlinewidth{0.000000pt}%
\definecolor{currentstroke}{rgb}{0.000000,0.000000,0.000000}%
\pgfsetstrokecolor{currentstroke}%
\pgfsetstrokeopacity{0.000000}%
\pgfsetdash{}{0pt}%
\pgfpathmoveto{\pgfqpoint{1.007808in}{2.262053in}}%
\pgfpathlineto{\pgfqpoint{0.998620in}{2.223097in}}%
\pgfpathlineto{\pgfqpoint{0.960593in}{2.232110in}}%
\pgfpathlineto{\pgfqpoint{0.967184in}{2.259137in}}%
\pgfpathlineto{\pgfqpoint{0.954484in}{2.262189in}}%
\pgfpathlineto{\pgfqpoint{0.957430in}{2.274082in}}%
\pgfpathlineto{\pgfqpoint{1.007808in}{2.262053in}}%
\pgfpathclose%
\pgfusepath{fill}%
\end{pgfscope}%
\begin{pgfscope}%
\pgfpathrectangle{\pgfqpoint{0.100000in}{0.100000in}}{\pgfqpoint{3.608454in}{2.310000in}}%
\pgfusepath{clip}%
\pgfsetbuttcap%
\pgfsetmiterjoin%
\definecolor{currentfill}{rgb}{0.000000,0.537255,0.731373}%
\pgfsetfillcolor{currentfill}%
\pgfsetlinewidth{0.000000pt}%
\definecolor{currentstroke}{rgb}{0.000000,0.000000,0.000000}%
\pgfsetstrokecolor{currentstroke}%
\pgfsetstrokeopacity{0.000000}%
\pgfsetdash{}{0pt}%
\pgfpathmoveto{\pgfqpoint{2.799445in}{1.624219in}}%
\pgfpathlineto{\pgfqpoint{2.792334in}{1.628535in}}%
\pgfpathlineto{\pgfqpoint{2.764552in}{1.624067in}}%
\pgfpathlineto{\pgfqpoint{2.743822in}{1.620714in}}%
\pgfpathlineto{\pgfqpoint{2.739935in}{1.648352in}}%
\pgfpathlineto{\pgfqpoint{2.735608in}{1.675908in}}%
\pgfpathlineto{\pgfqpoint{2.723746in}{1.674452in}}%
\pgfpathlineto{\pgfqpoint{2.720350in}{1.701323in}}%
\pgfpathlineto{\pgfqpoint{2.720255in}{1.702142in}}%
\pgfpathlineto{\pgfqpoint{2.745404in}{1.705512in}}%
\pgfpathlineto{\pgfqpoint{2.744382in}{1.712367in}}%
\pgfpathlineto{\pgfqpoint{2.771368in}{1.715974in}}%
\pgfpathlineto{\pgfqpoint{2.777478in}{1.717721in}}%
\pgfpathlineto{\pgfqpoint{2.776445in}{1.724555in}}%
\pgfpathlineto{\pgfqpoint{2.790730in}{1.726819in}}%
\pgfpathlineto{\pgfqpoint{2.786333in}{1.754076in}}%
\pgfpathlineto{\pgfqpoint{2.815217in}{1.759142in}}%
\pgfpathlineto{\pgfqpoint{2.817633in}{1.749127in}}%
\pgfpathlineto{\pgfqpoint{2.821949in}{1.738559in}}%
\pgfpathlineto{\pgfqpoint{2.825537in}{1.722337in}}%
\pgfpathlineto{\pgfqpoint{2.828453in}{1.713989in}}%
\pgfpathlineto{\pgfqpoint{2.833465in}{1.706759in}}%
\pgfpathlineto{\pgfqpoint{2.831898in}{1.690908in}}%
\pgfpathlineto{\pgfqpoint{2.833245in}{1.687359in}}%
\pgfpathlineto{\pgfqpoint{2.831702in}{1.674798in}}%
\pgfpathlineto{\pgfqpoint{2.820874in}{1.679961in}}%
\pgfpathlineto{\pgfqpoint{2.815658in}{1.676414in}}%
\pgfpathlineto{\pgfqpoint{2.812091in}{1.663621in}}%
\pgfpathlineto{\pgfqpoint{2.813316in}{1.659807in}}%
\pgfpathlineto{\pgfqpoint{2.810815in}{1.651423in}}%
\pgfpathlineto{\pgfqpoint{2.803471in}{1.647833in}}%
\pgfpathlineto{\pgfqpoint{2.800576in}{1.641151in}}%
\pgfpathlineto{\pgfqpoint{2.801791in}{1.629277in}}%
\pgfpathlineto{\pgfqpoint{2.799445in}{1.624219in}}%
\pgfpathclose%
\pgfusepath{fill}%
\end{pgfscope}%
\begin{pgfscope}%
\pgfpathrectangle{\pgfqpoint{0.100000in}{0.100000in}}{\pgfqpoint{3.608454in}{2.310000in}}%
\pgfusepath{clip}%
\pgfsetbuttcap%
\pgfsetmiterjoin%
\definecolor{currentfill}{rgb}{0.000000,0.603922,0.698039}%
\pgfsetfillcolor{currentfill}%
\pgfsetlinewidth{0.000000pt}%
\definecolor{currentstroke}{rgb}{0.000000,0.000000,0.000000}%
\pgfsetstrokecolor{currentstroke}%
\pgfsetstrokeopacity{0.000000}%
\pgfsetdash{}{0pt}%
\pgfpathmoveto{\pgfqpoint{0.763231in}{2.326563in}}%
\pgfpathlineto{\pgfqpoint{0.765166in}{2.300496in}}%
\pgfpathlineto{\pgfqpoint{0.760555in}{2.298682in}}%
\pgfpathlineto{\pgfqpoint{0.752119in}{2.299760in}}%
\pgfpathlineto{\pgfqpoint{0.728782in}{2.306722in}}%
\pgfpathlineto{\pgfqpoint{0.672017in}{2.324335in}}%
\pgfpathlineto{\pgfqpoint{0.672629in}{2.335131in}}%
\pgfpathlineto{\pgfqpoint{0.667385in}{2.333960in}}%
\pgfpathlineto{\pgfqpoint{0.662898in}{2.347488in}}%
\pgfpathlineto{\pgfqpoint{0.667167in}{2.355917in}}%
\pgfpathlineto{\pgfqpoint{0.717719in}{2.339867in}}%
\pgfpathlineto{\pgfqpoint{0.763231in}{2.326563in}}%
\pgfpathclose%
\pgfusepath{fill}%
\end{pgfscope}%
\begin{pgfscope}%
\pgfpathrectangle{\pgfqpoint{0.100000in}{0.100000in}}{\pgfqpoint{3.608454in}{2.310000in}}%
\pgfusepath{clip}%
\pgfsetbuttcap%
\pgfsetmiterjoin%
\definecolor{currentfill}{rgb}{0.000000,0.984314,0.507843}%
\pgfsetfillcolor{currentfill}%
\pgfsetlinewidth{0.000000pt}%
\definecolor{currentstroke}{rgb}{0.000000,0.000000,0.000000}%
\pgfsetstrokecolor{currentstroke}%
\pgfsetstrokeopacity{0.000000}%
\pgfsetdash{}{0pt}%
\pgfpathmoveto{\pgfqpoint{1.223167in}{1.249876in}}%
\pgfpathlineto{\pgfqpoint{1.210980in}{1.171960in}}%
\pgfpathlineto{\pgfqpoint{1.300870in}{1.158512in}}%
\pgfpathlineto{\pgfqpoint{1.321043in}{1.155597in}}%
\pgfpathlineto{\pgfqpoint{1.313512in}{1.101203in}}%
\pgfpathlineto{\pgfqpoint{1.239395in}{1.111911in}}%
\pgfpathlineto{\pgfqpoint{1.235282in}{1.084666in}}%
\pgfpathlineto{\pgfqpoint{1.198276in}{1.090371in}}%
\pgfpathlineto{\pgfqpoint{1.193614in}{1.060581in}}%
\pgfpathlineto{\pgfqpoint{1.183780in}{0.997802in}}%
\pgfpathlineto{\pgfqpoint{1.164171in}{1.000833in}}%
\pgfpathlineto{\pgfqpoint{1.154422in}{0.995185in}}%
\pgfpathlineto{\pgfqpoint{1.150108in}{0.990389in}}%
\pgfpathlineto{\pgfqpoint{1.145330in}{0.990568in}}%
\pgfpathlineto{\pgfqpoint{1.133001in}{0.981587in}}%
\pgfpathlineto{\pgfqpoint{1.125550in}{0.990854in}}%
\pgfpathlineto{\pgfqpoint{1.119039in}{0.992055in}}%
\pgfpathlineto{\pgfqpoint{1.124557in}{1.025178in}}%
\pgfpathlineto{\pgfqpoint{1.075926in}{1.033485in}}%
\pgfpathlineto{\pgfqpoint{1.079528in}{1.054057in}}%
\pgfpathlineto{\pgfqpoint{1.116943in}{1.267777in}}%
\pgfpathlineto{\pgfqpoint{1.163571in}{1.259392in}}%
\pgfpathlineto{\pgfqpoint{1.223167in}{1.249876in}}%
\pgfpathclose%
\pgfusepath{fill}%
\end{pgfscope}%
\begin{pgfscope}%
\pgfpathrectangle{\pgfqpoint{0.100000in}{0.100000in}}{\pgfqpoint{3.608454in}{2.310000in}}%
\pgfusepath{clip}%
\pgfsetbuttcap%
\pgfsetmiterjoin%
\definecolor{currentfill}{rgb}{0.000000,0.541176,0.729412}%
\pgfsetfillcolor{currentfill}%
\pgfsetlinewidth{0.000000pt}%
\definecolor{currentstroke}{rgb}{0.000000,0.000000,0.000000}%
\pgfsetstrokecolor{currentstroke}%
\pgfsetstrokeopacity{0.000000}%
\pgfsetdash{}{0pt}%
\pgfpathmoveto{\pgfqpoint{2.269571in}{1.909501in}}%
\pgfpathlineto{\pgfqpoint{2.297058in}{1.910416in}}%
\pgfpathlineto{\pgfqpoint{2.331620in}{1.911661in}}%
\pgfpathlineto{\pgfqpoint{2.332207in}{1.897841in}}%
\pgfpathlineto{\pgfqpoint{2.345925in}{1.898515in}}%
\pgfpathlineto{\pgfqpoint{2.348110in}{1.850512in}}%
\pgfpathlineto{\pgfqpoint{2.334351in}{1.850010in}}%
\pgfpathlineto{\pgfqpoint{2.334657in}{1.843071in}}%
\pgfpathlineto{\pgfqpoint{2.300225in}{1.841636in}}%
\pgfpathlineto{\pgfqpoint{2.300433in}{1.834842in}}%
\pgfpathlineto{\pgfqpoint{2.266038in}{1.833884in}}%
\pgfpathlineto{\pgfqpoint{2.265057in}{1.868094in}}%
\pgfpathlineto{\pgfqpoint{2.271935in}{1.868329in}}%
\pgfpathlineto{\pgfqpoint{2.270918in}{1.895690in}}%
\pgfpathlineto{\pgfqpoint{2.269571in}{1.909501in}}%
\pgfpathclose%
\pgfusepath{fill}%
\end{pgfscope}%
\begin{pgfscope}%
\pgfpathrectangle{\pgfqpoint{0.100000in}{0.100000in}}{\pgfqpoint{3.608454in}{2.310000in}}%
\pgfusepath{clip}%
\pgfsetbuttcap%
\pgfsetmiterjoin%
\definecolor{currentfill}{rgb}{0.000000,0.674510,0.662745}%
\pgfsetfillcolor{currentfill}%
\pgfsetlinewidth{0.000000pt}%
\definecolor{currentstroke}{rgb}{0.000000,0.000000,0.000000}%
\pgfsetstrokecolor{currentstroke}%
\pgfsetstrokeopacity{0.000000}%
\pgfsetdash{}{0pt}%
\pgfpathmoveto{\pgfqpoint{1.724900in}{0.803670in}}%
\pgfpathlineto{\pgfqpoint{1.690000in}{0.806092in}}%
\pgfpathlineto{\pgfqpoint{1.692995in}{0.840738in}}%
\pgfpathlineto{\pgfqpoint{1.695447in}{0.875598in}}%
\pgfpathlineto{\pgfqpoint{1.704446in}{0.875520in}}%
\pgfpathlineto{\pgfqpoint{1.706689in}{0.909371in}}%
\pgfpathlineto{\pgfqpoint{1.741190in}{0.907227in}}%
\pgfpathlineto{\pgfqpoint{1.738977in}{0.872740in}}%
\pgfpathlineto{\pgfqpoint{1.729883in}{0.873343in}}%
\pgfpathlineto{\pgfqpoint{1.724900in}{0.803670in}}%
\pgfpathclose%
\pgfusepath{fill}%
\end{pgfscope}%
\begin{pgfscope}%
\pgfpathrectangle{\pgfqpoint{0.100000in}{0.100000in}}{\pgfqpoint{3.608454in}{2.310000in}}%
\pgfusepath{clip}%
\pgfsetbuttcap%
\pgfsetmiterjoin%
\definecolor{currentfill}{rgb}{0.000000,0.827451,0.586275}%
\pgfsetfillcolor{currentfill}%
\pgfsetlinewidth{0.000000pt}%
\definecolor{currentstroke}{rgb}{0.000000,0.000000,0.000000}%
\pgfsetstrokecolor{currentstroke}%
\pgfsetstrokeopacity{0.000000}%
\pgfsetdash{}{0pt}%
\pgfpathmoveto{\pgfqpoint{2.474841in}{0.898468in}}%
\pgfpathlineto{\pgfqpoint{2.483949in}{0.899006in}}%
\pgfpathlineto{\pgfqpoint{2.485523in}{0.870923in}}%
\pgfpathlineto{\pgfqpoint{2.458164in}{0.869135in}}%
\pgfpathlineto{\pgfqpoint{2.458381in}{0.865552in}}%
\pgfpathlineto{\pgfqpoint{2.442684in}{0.864247in}}%
\pgfpathlineto{\pgfqpoint{2.417323in}{0.873285in}}%
\pgfpathlineto{\pgfqpoint{2.410750in}{0.879777in}}%
\pgfpathlineto{\pgfqpoint{2.417347in}{0.883236in}}%
\pgfpathlineto{\pgfqpoint{2.416578in}{0.896534in}}%
\pgfpathlineto{\pgfqpoint{2.419871in}{0.899622in}}%
\pgfpathlineto{\pgfqpoint{2.426856in}{0.899501in}}%
\pgfpathlineto{\pgfqpoint{2.455570in}{0.891636in}}%
\pgfpathlineto{\pgfqpoint{2.458419in}{0.897453in}}%
\pgfpathlineto{\pgfqpoint{2.474841in}{0.898468in}}%
\pgfpathclose%
\pgfusepath{fill}%
\end{pgfscope}%
\begin{pgfscope}%
\pgfpathrectangle{\pgfqpoint{0.100000in}{0.100000in}}{\pgfqpoint{3.608454in}{2.310000in}}%
\pgfusepath{clip}%
\pgfsetbuttcap%
\pgfsetmiterjoin%
\definecolor{currentfill}{rgb}{0.000000,0.552941,0.723529}%
\pgfsetfillcolor{currentfill}%
\pgfsetlinewidth{0.000000pt}%
\definecolor{currentstroke}{rgb}{0.000000,0.000000,0.000000}%
\pgfsetstrokecolor{currentstroke}%
\pgfsetstrokeopacity{0.000000}%
\pgfsetdash{}{0pt}%
\pgfpathmoveto{\pgfqpoint{2.348213in}{1.508885in}}%
\pgfpathlineto{\pgfqpoint{2.323109in}{1.507938in}}%
\pgfpathlineto{\pgfqpoint{2.322856in}{1.514824in}}%
\pgfpathlineto{\pgfqpoint{2.302293in}{1.514165in}}%
\pgfpathlineto{\pgfqpoint{2.288611in}{1.513740in}}%
\pgfpathlineto{\pgfqpoint{2.287673in}{1.541193in}}%
\pgfpathlineto{\pgfqpoint{2.294578in}{1.541485in}}%
\pgfpathlineto{\pgfqpoint{2.293482in}{1.569183in}}%
\pgfpathlineto{\pgfqpoint{2.320902in}{1.570069in}}%
\pgfpathlineto{\pgfqpoint{2.320640in}{1.576968in}}%
\pgfpathlineto{\pgfqpoint{2.348123in}{1.578046in}}%
\pgfpathlineto{\pgfqpoint{2.348758in}{1.564236in}}%
\pgfpathlineto{\pgfqpoint{2.349290in}{1.550458in}}%
\pgfpathlineto{\pgfqpoint{2.356154in}{1.550723in}}%
\pgfpathlineto{\pgfqpoint{2.356534in}{1.539366in}}%
\pgfpathlineto{\pgfqpoint{2.341290in}{1.535608in}}%
\pgfpathlineto{\pgfqpoint{2.337899in}{1.521676in}}%
\pgfpathlineto{\pgfqpoint{2.345525in}{1.515337in}}%
\pgfpathlineto{\pgfqpoint{2.348213in}{1.508885in}}%
\pgfpathclose%
\pgfusepath{fill}%
\end{pgfscope}%
\begin{pgfscope}%
\pgfpathrectangle{\pgfqpoint{0.100000in}{0.100000in}}{\pgfqpoint{3.608454in}{2.310000in}}%
\pgfusepath{clip}%
\pgfsetbuttcap%
\pgfsetmiterjoin%
\definecolor{currentfill}{rgb}{0.000000,0.513725,0.743137}%
\pgfsetfillcolor{currentfill}%
\pgfsetlinewidth{0.000000pt}%
\definecolor{currentstroke}{rgb}{0.000000,0.000000,0.000000}%
\pgfsetstrokecolor{currentstroke}%
\pgfsetstrokeopacity{0.000000}%
\pgfsetdash{}{0pt}%
\pgfpathmoveto{\pgfqpoint{2.099037in}{1.876260in}}%
\pgfpathlineto{\pgfqpoint{2.099439in}{1.847533in}}%
\pgfpathlineto{\pgfqpoint{2.044431in}{1.847876in}}%
\pgfpathlineto{\pgfqpoint{2.037488in}{1.848022in}}%
\pgfpathlineto{\pgfqpoint{2.037019in}{1.861775in}}%
\pgfpathlineto{\pgfqpoint{2.037372in}{1.889582in}}%
\pgfpathlineto{\pgfqpoint{2.036803in}{1.896501in}}%
\pgfpathlineto{\pgfqpoint{2.036301in}{1.917281in}}%
\pgfpathlineto{\pgfqpoint{2.036642in}{1.945064in}}%
\pgfpathlineto{\pgfqpoint{2.042594in}{1.944974in}}%
\pgfpathlineto{\pgfqpoint{2.042559in}{1.951936in}}%
\pgfpathlineto{\pgfqpoint{2.097851in}{1.951504in}}%
\pgfpathlineto{\pgfqpoint{2.098775in}{1.916792in}}%
\pgfpathlineto{\pgfqpoint{2.099037in}{1.876260in}}%
\pgfpathclose%
\pgfusepath{fill}%
\end{pgfscope}%
\begin{pgfscope}%
\pgfpathrectangle{\pgfqpoint{0.100000in}{0.100000in}}{\pgfqpoint{3.608454in}{2.310000in}}%
\pgfusepath{clip}%
\pgfsetbuttcap%
\pgfsetmiterjoin%
\definecolor{currentfill}{rgb}{0.000000,0.549020,0.725490}%
\pgfsetfillcolor{currentfill}%
\pgfsetlinewidth{0.000000pt}%
\definecolor{currentstroke}{rgb}{0.000000,0.000000,0.000000}%
\pgfsetstrokecolor{currentstroke}%
\pgfsetstrokeopacity{0.000000}%
\pgfsetdash{}{0pt}%
\pgfpathmoveto{\pgfqpoint{2.817692in}{1.389496in}}%
\pgfpathlineto{\pgfqpoint{2.804293in}{1.384898in}}%
\pgfpathlineto{\pgfqpoint{2.788388in}{1.383110in}}%
\pgfpathlineto{\pgfqpoint{2.786603in}{1.400891in}}%
\pgfpathlineto{\pgfqpoint{2.777953in}{1.400497in}}%
\pgfpathlineto{\pgfqpoint{2.776762in}{1.425313in}}%
\pgfpathlineto{\pgfqpoint{2.795500in}{1.426123in}}%
\pgfpathlineto{\pgfqpoint{2.794880in}{1.439319in}}%
\pgfpathlineto{\pgfqpoint{2.819277in}{1.440727in}}%
\pgfpathlineto{\pgfqpoint{2.820238in}{1.426570in}}%
\pgfpathlineto{\pgfqpoint{2.815226in}{1.414852in}}%
\pgfpathlineto{\pgfqpoint{2.816007in}{1.404773in}}%
\pgfpathlineto{\pgfqpoint{2.818458in}{1.403774in}}%
\pgfpathlineto{\pgfqpoint{2.817692in}{1.389496in}}%
\pgfpathclose%
\pgfusepath{fill}%
\end{pgfscope}%
\begin{pgfscope}%
\pgfpathrectangle{\pgfqpoint{0.100000in}{0.100000in}}{\pgfqpoint{3.608454in}{2.310000in}}%
\pgfusepath{clip}%
\pgfsetbuttcap%
\pgfsetmiterjoin%
\definecolor{currentfill}{rgb}{0.000000,0.670588,0.664706}%
\pgfsetfillcolor{currentfill}%
\pgfsetlinewidth{0.000000pt}%
\definecolor{currentstroke}{rgb}{0.000000,0.000000,0.000000}%
\pgfsetstrokecolor{currentstroke}%
\pgfsetstrokeopacity{0.000000}%
\pgfsetdash{}{0pt}%
\pgfpathmoveto{\pgfqpoint{3.596315in}{2.093414in}}%
\pgfpathlineto{\pgfqpoint{3.619860in}{2.089848in}}%
\pgfpathlineto{\pgfqpoint{3.617148in}{2.084070in}}%
\pgfpathlineto{\pgfqpoint{3.624313in}{2.074071in}}%
\pgfpathlineto{\pgfqpoint{3.623386in}{2.068081in}}%
\pgfpathlineto{\pgfqpoint{3.634273in}{2.058968in}}%
\pgfpathlineto{\pgfqpoint{3.636385in}{2.065102in}}%
\pgfpathlineto{\pgfqpoint{3.643097in}{2.064783in}}%
\pgfpathlineto{\pgfqpoint{3.657645in}{2.048872in}}%
\pgfpathlineto{\pgfqpoint{3.660907in}{2.041358in}}%
\pgfpathlineto{\pgfqpoint{3.656677in}{2.034155in}}%
\pgfpathlineto{\pgfqpoint{3.653836in}{2.025053in}}%
\pgfpathlineto{\pgfqpoint{3.642794in}{2.024053in}}%
\pgfpathlineto{\pgfqpoint{3.643360in}{2.017403in}}%
\pgfpathlineto{\pgfqpoint{3.634988in}{2.017176in}}%
\pgfpathlineto{\pgfqpoint{3.633250in}{2.008704in}}%
\pgfpathlineto{\pgfqpoint{3.624318in}{2.007131in}}%
\pgfpathlineto{\pgfqpoint{3.621576in}{1.995429in}}%
\pgfpathlineto{\pgfqpoint{3.618035in}{1.994535in}}%
\pgfpathlineto{\pgfqpoint{3.610482in}{2.006964in}}%
\pgfpathlineto{\pgfqpoint{3.597567in}{2.032110in}}%
\pgfpathlineto{\pgfqpoint{3.603808in}{2.035391in}}%
\pgfpathlineto{\pgfqpoint{3.592928in}{2.056963in}}%
\pgfpathlineto{\pgfqpoint{3.598893in}{2.059799in}}%
\pgfpathlineto{\pgfqpoint{3.584395in}{2.086358in}}%
\pgfpathlineto{\pgfqpoint{3.596315in}{2.093414in}}%
\pgfpathclose%
\pgfusepath{fill}%
\end{pgfscope}%
\begin{pgfscope}%
\pgfpathrectangle{\pgfqpoint{0.100000in}{0.100000in}}{\pgfqpoint{3.608454in}{2.310000in}}%
\pgfusepath{clip}%
\pgfsetbuttcap%
\pgfsetmiterjoin%
\definecolor{currentfill}{rgb}{0.000000,0.529412,0.735294}%
\pgfsetfillcolor{currentfill}%
\pgfsetlinewidth{0.000000pt}%
\definecolor{currentstroke}{rgb}{0.000000,0.000000,0.000000}%
\pgfsetstrokecolor{currentstroke}%
\pgfsetstrokeopacity{0.000000}%
\pgfsetdash{}{0pt}%
\pgfpathmoveto{\pgfqpoint{2.833971in}{1.586462in}}%
\pgfpathlineto{\pgfqpoint{2.818905in}{1.581252in}}%
\pgfpathlineto{\pgfqpoint{2.796288in}{1.578242in}}%
\pgfpathlineto{\pgfqpoint{2.791356in}{1.581015in}}%
\pgfpathlineto{\pgfqpoint{2.794328in}{1.554792in}}%
\pgfpathlineto{\pgfqpoint{2.767138in}{1.551464in}}%
\pgfpathlineto{\pgfqpoint{2.746637in}{1.548837in}}%
\pgfpathlineto{\pgfqpoint{2.744178in}{1.569410in}}%
\pgfpathlineto{\pgfqpoint{2.737491in}{1.568598in}}%
\pgfpathlineto{\pgfqpoint{2.734410in}{1.575123in}}%
\pgfpathlineto{\pgfqpoint{2.731544in}{1.590059in}}%
\pgfpathlineto{\pgfqpoint{2.768691in}{1.595961in}}%
\pgfpathlineto{\pgfqpoint{2.764552in}{1.624067in}}%
\pgfpathlineto{\pgfqpoint{2.792334in}{1.628535in}}%
\pgfpathlineto{\pgfqpoint{2.799445in}{1.624219in}}%
\pgfpathlineto{\pgfqpoint{2.792747in}{1.616003in}}%
\pgfpathlineto{\pgfqpoint{2.786784in}{1.605434in}}%
\pgfpathlineto{\pgfqpoint{2.786107in}{1.596021in}}%
\pgfpathlineto{\pgfqpoint{2.793257in}{1.596373in}}%
\pgfpathlineto{\pgfqpoint{2.810519in}{1.591373in}}%
\pgfpathlineto{\pgfqpoint{2.817182in}{1.586236in}}%
\pgfpathlineto{\pgfqpoint{2.831148in}{1.589793in}}%
\pgfpathlineto{\pgfqpoint{2.833971in}{1.586462in}}%
\pgfpathclose%
\pgfusepath{fill}%
\end{pgfscope}%
\begin{pgfscope}%
\pgfpathrectangle{\pgfqpoint{0.100000in}{0.100000in}}{\pgfqpoint{3.608454in}{2.310000in}}%
\pgfusepath{clip}%
\pgfsetbuttcap%
\pgfsetmiterjoin%
\definecolor{currentfill}{rgb}{0.000000,0.556863,0.721569}%
\pgfsetfillcolor{currentfill}%
\pgfsetlinewidth{0.000000pt}%
\definecolor{currentstroke}{rgb}{0.000000,0.000000,0.000000}%
\pgfsetstrokecolor{currentstroke}%
\pgfsetstrokeopacity{0.000000}%
\pgfsetdash{}{0pt}%
\pgfpathmoveto{\pgfqpoint{2.380907in}{1.900334in}}%
\pgfpathlineto{\pgfqpoint{2.382793in}{1.866283in}}%
\pgfpathlineto{\pgfqpoint{2.383604in}{1.852652in}}%
\pgfpathlineto{\pgfqpoint{2.348110in}{1.850512in}}%
\pgfpathlineto{\pgfqpoint{2.345925in}{1.898515in}}%
\pgfpathlineto{\pgfqpoint{2.332207in}{1.897841in}}%
\pgfpathlineto{\pgfqpoint{2.331620in}{1.911661in}}%
\pgfpathlineto{\pgfqpoint{2.297058in}{1.910416in}}%
\pgfpathlineto{\pgfqpoint{2.295263in}{1.958166in}}%
\pgfpathlineto{\pgfqpoint{2.305024in}{1.961414in}}%
\pgfpathlineto{\pgfqpoint{2.313626in}{1.967707in}}%
\pgfpathlineto{\pgfqpoint{2.319119in}{1.967177in}}%
\pgfpathlineto{\pgfqpoint{2.326155in}{1.974190in}}%
\pgfpathlineto{\pgfqpoint{2.332310in}{1.976143in}}%
\pgfpathlineto{\pgfqpoint{2.338487in}{1.970865in}}%
\pgfpathlineto{\pgfqpoint{2.331663in}{1.959426in}}%
\pgfpathlineto{\pgfqpoint{2.331044in}{1.945780in}}%
\pgfpathlineto{\pgfqpoint{2.342366in}{1.952806in}}%
\pgfpathlineto{\pgfqpoint{2.350608in}{1.946791in}}%
\pgfpathlineto{\pgfqpoint{2.351813in}{1.919479in}}%
\pgfpathlineto{\pgfqpoint{2.358761in}{1.919792in}}%
\pgfpathlineto{\pgfqpoint{2.359091in}{1.912916in}}%
\pgfpathlineto{\pgfqpoint{2.365878in}{1.913327in}}%
\pgfpathlineto{\pgfqpoint{2.366647in}{1.899528in}}%
\pgfpathlineto{\pgfqpoint{2.380907in}{1.900334in}}%
\pgfpathclose%
\pgfusepath{fill}%
\end{pgfscope}%
\begin{pgfscope}%
\pgfpathrectangle{\pgfqpoint{0.100000in}{0.100000in}}{\pgfqpoint{3.608454in}{2.310000in}}%
\pgfusepath{clip}%
\pgfsetbuttcap%
\pgfsetmiterjoin%
\definecolor{currentfill}{rgb}{0.000000,0.611765,0.694118}%
\pgfsetfillcolor{currentfill}%
\pgfsetlinewidth{0.000000pt}%
\definecolor{currentstroke}{rgb}{0.000000,0.000000,0.000000}%
\pgfsetstrokecolor{currentstroke}%
\pgfsetstrokeopacity{0.000000}%
\pgfsetdash{}{0pt}%
\pgfpathmoveto{\pgfqpoint{2.634617in}{1.461203in}}%
\pgfpathlineto{\pgfqpoint{2.634909in}{1.458446in}}%
\pgfpathlineto{\pgfqpoint{2.607655in}{1.455712in}}%
\pgfpathlineto{\pgfqpoint{2.607399in}{1.458524in}}%
\pgfpathlineto{\pgfqpoint{2.583559in}{1.456426in}}%
\pgfpathlineto{\pgfqpoint{2.585430in}{1.435815in}}%
\pgfpathlineto{\pgfqpoint{2.565557in}{1.434073in}}%
\pgfpathlineto{\pgfqpoint{2.565221in}{1.447932in}}%
\pgfpathlineto{\pgfqpoint{2.557572in}{1.448924in}}%
\pgfpathlineto{\pgfqpoint{2.559183in}{1.428045in}}%
\pgfpathlineto{\pgfqpoint{2.534753in}{1.425828in}}%
\pgfpathlineto{\pgfqpoint{2.532347in}{1.453054in}}%
\pgfpathlineto{\pgfqpoint{2.531157in}{1.473630in}}%
\pgfpathlineto{\pgfqpoint{2.555646in}{1.476018in}}%
\pgfpathlineto{\pgfqpoint{2.554067in}{1.495355in}}%
\pgfpathlineto{\pgfqpoint{2.579594in}{1.497487in}}%
\pgfpathlineto{\pgfqpoint{2.578828in}{1.505444in}}%
\pgfpathlineto{\pgfqpoint{2.585581in}{1.506060in}}%
\pgfpathlineto{\pgfqpoint{2.588380in}{1.512229in}}%
\pgfpathlineto{\pgfqpoint{2.609172in}{1.513996in}}%
\pgfpathlineto{\pgfqpoint{2.610441in}{1.500202in}}%
\pgfpathlineto{\pgfqpoint{2.614139in}{1.497101in}}%
\pgfpathlineto{\pgfqpoint{2.623185in}{1.497972in}}%
\pgfpathlineto{\pgfqpoint{2.625004in}{1.477436in}}%
\pgfpathlineto{\pgfqpoint{2.633399in}{1.473636in}}%
\pgfpathlineto{\pgfqpoint{2.634617in}{1.461203in}}%
\pgfpathclose%
\pgfusepath{fill}%
\end{pgfscope}%
\begin{pgfscope}%
\pgfpathrectangle{\pgfqpoint{0.100000in}{0.100000in}}{\pgfqpoint{3.608454in}{2.310000in}}%
\pgfusepath{clip}%
\pgfsetbuttcap%
\pgfsetmiterjoin%
\definecolor{currentfill}{rgb}{0.000000,0.788235,0.605882}%
\pgfsetfillcolor{currentfill}%
\pgfsetlinewidth{0.000000pt}%
\definecolor{currentstroke}{rgb}{0.000000,0.000000,0.000000}%
\pgfsetstrokecolor{currentstroke}%
\pgfsetstrokeopacity{0.000000}%
\pgfsetdash{}{0pt}%
\pgfpathmoveto{\pgfqpoint{2.202317in}{0.720035in}}%
\pgfpathlineto{\pgfqpoint{2.181123in}{0.718301in}}%
\pgfpathlineto{\pgfqpoint{2.166253in}{0.713450in}}%
\pgfpathlineto{\pgfqpoint{2.152123in}{0.723725in}}%
\pgfpathlineto{\pgfqpoint{2.145077in}{0.730301in}}%
\pgfpathlineto{\pgfqpoint{2.140576in}{0.738107in}}%
\pgfpathlineto{\pgfqpoint{2.125247in}{0.741762in}}%
\pgfpathlineto{\pgfqpoint{2.115910in}{0.747279in}}%
\pgfpathlineto{\pgfqpoint{2.111897in}{0.751748in}}%
\pgfpathlineto{\pgfqpoint{2.108517in}{0.767596in}}%
\pgfpathlineto{\pgfqpoint{2.111050in}{0.772661in}}%
\pgfpathlineto{\pgfqpoint{2.143853in}{0.772670in}}%
\pgfpathlineto{\pgfqpoint{2.139795in}{0.782977in}}%
\pgfpathlineto{\pgfqpoint{2.173280in}{0.783716in}}%
\pgfpathlineto{\pgfqpoint{2.182637in}{0.773027in}}%
\pgfpathlineto{\pgfqpoint{2.186530in}{0.767441in}}%
\pgfpathlineto{\pgfqpoint{2.187852in}{0.760329in}}%
\pgfpathlineto{\pgfqpoint{2.185571in}{0.753524in}}%
\pgfpathlineto{\pgfqpoint{2.191816in}{0.743213in}}%
\pgfpathlineto{\pgfqpoint{2.197023in}{0.737625in}}%
\pgfpathlineto{\pgfqpoint{2.196211in}{0.730685in}}%
\pgfpathlineto{\pgfqpoint{2.202317in}{0.720035in}}%
\pgfpathclose%
\pgfusepath{fill}%
\end{pgfscope}%
\begin{pgfscope}%
\pgfpathrectangle{\pgfqpoint{0.100000in}{0.100000in}}{\pgfqpoint{3.608454in}{2.310000in}}%
\pgfusepath{clip}%
\pgfsetbuttcap%
\pgfsetmiterjoin%
\definecolor{currentfill}{rgb}{0.000000,0.725490,0.637255}%
\pgfsetfillcolor{currentfill}%
\pgfsetlinewidth{0.000000pt}%
\definecolor{currentstroke}{rgb}{0.000000,0.000000,0.000000}%
\pgfsetstrokecolor{currentstroke}%
\pgfsetstrokeopacity{0.000000}%
\pgfsetdash{}{0pt}%
\pgfpathmoveto{\pgfqpoint{3.147047in}{0.962987in}}%
\pgfpathlineto{\pgfqpoint{3.134423in}{0.971285in}}%
\pgfpathlineto{\pgfqpoint{3.115214in}{0.974844in}}%
\pgfpathlineto{\pgfqpoint{3.097264in}{0.986382in}}%
\pgfpathlineto{\pgfqpoint{3.088906in}{0.984668in}}%
\pgfpathlineto{\pgfqpoint{3.094287in}{1.003646in}}%
\pgfpathlineto{\pgfqpoint{3.092298in}{1.007042in}}%
\pgfpathlineto{\pgfqpoint{3.096860in}{1.011889in}}%
\pgfpathlineto{\pgfqpoint{3.091334in}{1.021253in}}%
\pgfpathlineto{\pgfqpoint{3.096046in}{1.025469in}}%
\pgfpathlineto{\pgfqpoint{3.076612in}{1.036110in}}%
\pgfpathlineto{\pgfqpoint{3.078779in}{1.041854in}}%
\pgfpathlineto{\pgfqpoint{3.073425in}{1.048823in}}%
\pgfpathlineto{\pgfqpoint{3.065576in}{1.050744in}}%
\pgfpathlineto{\pgfqpoint{3.079593in}{1.063190in}}%
\pgfpathlineto{\pgfqpoint{3.095167in}{1.065172in}}%
\pgfpathlineto{\pgfqpoint{3.108787in}{1.052476in}}%
\pgfpathlineto{\pgfqpoint{3.116266in}{1.079750in}}%
\pgfpathlineto{\pgfqpoint{3.145061in}{1.058883in}}%
\pgfpathlineto{\pgfqpoint{3.185535in}{1.030022in}}%
\pgfpathlineto{\pgfqpoint{3.175718in}{1.024509in}}%
\pgfpathlineto{\pgfqpoint{3.169582in}{1.018517in}}%
\pgfpathlineto{\pgfqpoint{3.162304in}{1.007975in}}%
\pgfpathlineto{\pgfqpoint{3.155468in}{0.994985in}}%
\pgfpathlineto{\pgfqpoint{3.152816in}{0.985825in}}%
\pgfpathlineto{\pgfqpoint{3.152996in}{0.969508in}}%
\pgfpathlineto{\pgfqpoint{3.147047in}{0.962987in}}%
\pgfpathclose%
\pgfusepath{fill}%
\end{pgfscope}%
\begin{pgfscope}%
\pgfpathrectangle{\pgfqpoint{0.100000in}{0.100000in}}{\pgfqpoint{3.608454in}{2.310000in}}%
\pgfusepath{clip}%
\pgfsetbuttcap%
\pgfsetmiterjoin%
\definecolor{currentfill}{rgb}{0.000000,0.796078,0.601961}%
\pgfsetfillcolor{currentfill}%
\pgfsetlinewidth{0.000000pt}%
\definecolor{currentstroke}{rgb}{0.000000,0.000000,0.000000}%
\pgfsetstrokecolor{currentstroke}%
\pgfsetstrokeopacity{0.000000}%
\pgfsetdash{}{0pt}%
\pgfpathmoveto{\pgfqpoint{2.611522in}{1.917144in}}%
\pgfpathlineto{\pgfqpoint{2.605875in}{1.974408in}}%
\pgfpathlineto{\pgfqpoint{2.625687in}{1.975346in}}%
\pgfpathlineto{\pgfqpoint{2.636425in}{1.981889in}}%
\pgfpathlineto{\pgfqpoint{2.653804in}{1.986050in}}%
\pgfpathlineto{\pgfqpoint{2.651220in}{1.978697in}}%
\pgfpathlineto{\pgfqpoint{2.652749in}{1.962987in}}%
\pgfpathlineto{\pgfqpoint{2.664373in}{1.960858in}}%
\pgfpathlineto{\pgfqpoint{2.672255in}{1.965247in}}%
\pgfpathlineto{\pgfqpoint{2.677033in}{1.960330in}}%
\pgfpathlineto{\pgfqpoint{2.686113in}{1.967978in}}%
\pgfpathlineto{\pgfqpoint{2.702927in}{1.970218in}}%
\pgfpathlineto{\pgfqpoint{2.701836in}{1.963213in}}%
\pgfpathlineto{\pgfqpoint{2.704934in}{1.955834in}}%
\pgfpathlineto{\pgfqpoint{2.703228in}{1.945319in}}%
\pgfpathlineto{\pgfqpoint{2.711142in}{1.941318in}}%
\pgfpathlineto{\pgfqpoint{2.714091in}{1.933704in}}%
\pgfpathlineto{\pgfqpoint{2.724603in}{1.932074in}}%
\pgfpathlineto{\pgfqpoint{2.735969in}{1.940892in}}%
\pgfpathlineto{\pgfqpoint{2.742949in}{1.933510in}}%
\pgfpathlineto{\pgfqpoint{2.741240in}{1.928257in}}%
\pgfpathlineto{\pgfqpoint{2.733601in}{1.928846in}}%
\pgfpathlineto{\pgfqpoint{2.725778in}{1.927125in}}%
\pgfpathlineto{\pgfqpoint{2.720067in}{1.929651in}}%
\pgfpathlineto{\pgfqpoint{2.706856in}{1.928149in}}%
\pgfpathlineto{\pgfqpoint{2.693720in}{1.923132in}}%
\pgfpathlineto{\pgfqpoint{2.684872in}{1.925062in}}%
\pgfpathlineto{\pgfqpoint{2.682294in}{1.929947in}}%
\pgfpathlineto{\pgfqpoint{2.675578in}{1.929116in}}%
\pgfpathlineto{\pgfqpoint{2.673548in}{1.913649in}}%
\pgfpathlineto{\pgfqpoint{2.664389in}{1.919036in}}%
\pgfpathlineto{\pgfqpoint{2.658981in}{1.924851in}}%
\pgfpathlineto{\pgfqpoint{2.650299in}{1.927597in}}%
\pgfpathlineto{\pgfqpoint{2.635790in}{1.929952in}}%
\pgfpathlineto{\pgfqpoint{2.629423in}{1.928847in}}%
\pgfpathlineto{\pgfqpoint{2.622249in}{1.917995in}}%
\pgfpathlineto{\pgfqpoint{2.611522in}{1.917144in}}%
\pgfpathclose%
\pgfusepath{fill}%
\end{pgfscope}%
\begin{pgfscope}%
\pgfpathrectangle{\pgfqpoint{0.100000in}{0.100000in}}{\pgfqpoint{3.608454in}{2.310000in}}%
\pgfusepath{clip}%
\pgfsetbuttcap%
\pgfsetmiterjoin%
\definecolor{currentfill}{rgb}{0.000000,0.662745,0.668627}%
\pgfsetfillcolor{currentfill}%
\pgfsetlinewidth{0.000000pt}%
\definecolor{currentstroke}{rgb}{0.000000,0.000000,0.000000}%
\pgfsetstrokecolor{currentstroke}%
\pgfsetstrokeopacity{0.000000}%
\pgfsetdash{}{0pt}%
\pgfpathmoveto{\pgfqpoint{2.482016in}{1.066685in}}%
\pgfpathlineto{\pgfqpoint{2.463288in}{1.065984in}}%
\pgfpathlineto{\pgfqpoint{2.461856in}{1.079471in}}%
\pgfpathlineto{\pgfqpoint{2.453807in}{1.083497in}}%
\pgfpathlineto{\pgfqpoint{2.441604in}{1.083493in}}%
\pgfpathlineto{\pgfqpoint{2.434616in}{1.075007in}}%
\pgfpathlineto{\pgfqpoint{2.432045in}{1.079916in}}%
\pgfpathlineto{\pgfqpoint{2.431130in}{1.090720in}}%
\pgfpathlineto{\pgfqpoint{2.439328in}{1.092627in}}%
\pgfpathlineto{\pgfqpoint{2.442664in}{1.097005in}}%
\pgfpathlineto{\pgfqpoint{2.442276in}{1.103468in}}%
\pgfpathlineto{\pgfqpoint{2.445220in}{1.109326in}}%
\pgfpathlineto{\pgfqpoint{2.446875in}{1.118895in}}%
\pgfpathlineto{\pgfqpoint{2.452158in}{1.122751in}}%
\pgfpathlineto{\pgfqpoint{2.445412in}{1.130104in}}%
\pgfpathlineto{\pgfqpoint{2.455346in}{1.132898in}}%
\pgfpathlineto{\pgfqpoint{2.454276in}{1.139497in}}%
\pgfpathlineto{\pgfqpoint{2.453898in}{1.152266in}}%
\pgfpathlineto{\pgfqpoint{2.457328in}{1.152433in}}%
\pgfpathlineto{\pgfqpoint{2.461626in}{1.152803in}}%
\pgfpathlineto{\pgfqpoint{2.465696in}{1.145939in}}%
\pgfpathlineto{\pgfqpoint{2.462058in}{1.133851in}}%
\pgfpathlineto{\pgfqpoint{2.458843in}{1.129936in}}%
\pgfpathlineto{\pgfqpoint{2.479806in}{1.130609in}}%
\pgfpathlineto{\pgfqpoint{2.480144in}{1.116463in}}%
\pgfpathlineto{\pgfqpoint{2.469563in}{1.097602in}}%
\pgfpathlineto{\pgfqpoint{2.487860in}{1.090523in}}%
\pgfpathlineto{\pgfqpoint{2.488608in}{1.069834in}}%
\pgfpathlineto{\pgfqpoint{2.482016in}{1.066685in}}%
\pgfpathclose%
\pgfusepath{fill}%
\end{pgfscope}%
\begin{pgfscope}%
\pgfpathrectangle{\pgfqpoint{0.100000in}{0.100000in}}{\pgfqpoint{3.608454in}{2.310000in}}%
\pgfusepath{clip}%
\pgfsetbuttcap%
\pgfsetmiterjoin%
\definecolor{currentfill}{rgb}{0.000000,0.470588,0.764706}%
\pgfsetfillcolor{currentfill}%
\pgfsetlinewidth{0.000000pt}%
\definecolor{currentstroke}{rgb}{0.000000,0.000000,0.000000}%
\pgfsetstrokecolor{currentstroke}%
\pgfsetstrokeopacity{0.000000}%
\pgfsetdash{}{0pt}%
\pgfpathmoveto{\pgfqpoint{2.970938in}{1.132406in}}%
\pgfpathlineto{\pgfqpoint{2.978683in}{1.148660in}}%
\pgfpathlineto{\pgfqpoint{2.973998in}{1.150189in}}%
\pgfpathlineto{\pgfqpoint{2.967129in}{1.148531in}}%
\pgfpathlineto{\pgfqpoint{2.956204in}{1.154739in}}%
\pgfpathlineto{\pgfqpoint{2.949055in}{1.160645in}}%
\pgfpathlineto{\pgfqpoint{2.955589in}{1.174438in}}%
\pgfpathlineto{\pgfqpoint{2.963826in}{1.175254in}}%
\pgfpathlineto{\pgfqpoint{2.972912in}{1.172989in}}%
\pgfpathlineto{\pgfqpoint{2.978982in}{1.167861in}}%
\pgfpathlineto{\pgfqpoint{2.987201in}{1.172347in}}%
\pgfpathlineto{\pgfqpoint{2.993174in}{1.172573in}}%
\pgfpathlineto{\pgfqpoint{3.006498in}{1.176840in}}%
\pgfpathlineto{\pgfqpoint{3.018531in}{1.178366in}}%
\pgfpathlineto{\pgfqpoint{3.020036in}{1.162708in}}%
\pgfpathlineto{\pgfqpoint{3.018191in}{1.148162in}}%
\pgfpathlineto{\pgfqpoint{3.022456in}{1.135384in}}%
\pgfpathlineto{\pgfqpoint{3.008399in}{1.134076in}}%
\pgfpathlineto{\pgfqpoint{3.007763in}{1.136339in}}%
\pgfpathlineto{\pgfqpoint{2.970938in}{1.132406in}}%
\pgfpathclose%
\pgfusepath{fill}%
\end{pgfscope}%
\begin{pgfscope}%
\pgfpathrectangle{\pgfqpoint{0.100000in}{0.100000in}}{\pgfqpoint{3.608454in}{2.310000in}}%
\pgfusepath{clip}%
\pgfsetbuttcap%
\pgfsetmiterjoin%
\definecolor{currentfill}{rgb}{0.000000,0.517647,0.741176}%
\pgfsetfillcolor{currentfill}%
\pgfsetlinewidth{0.000000pt}%
\definecolor{currentstroke}{rgb}{0.000000,0.000000,0.000000}%
\pgfsetstrokecolor{currentstroke}%
\pgfsetstrokeopacity{0.000000}%
\pgfsetdash{}{0pt}%
\pgfpathmoveto{\pgfqpoint{2.801132in}{1.502933in}}%
\pgfpathlineto{\pgfqpoint{2.795078in}{1.502065in}}%
\pgfpathlineto{\pgfqpoint{2.773279in}{1.502104in}}%
\pgfpathlineto{\pgfqpoint{2.767138in}{1.551464in}}%
\pgfpathlineto{\pgfqpoint{2.794328in}{1.554792in}}%
\pgfpathlineto{\pgfqpoint{2.791356in}{1.581015in}}%
\pgfpathlineto{\pgfqpoint{2.796288in}{1.578242in}}%
\pgfpathlineto{\pgfqpoint{2.818905in}{1.581252in}}%
\pgfpathlineto{\pgfqpoint{2.825402in}{1.579904in}}%
\pgfpathlineto{\pgfqpoint{2.831040in}{1.546066in}}%
\pgfpathlineto{\pgfqpoint{2.814322in}{1.543583in}}%
\pgfpathlineto{\pgfqpoint{2.817395in}{1.520810in}}%
\pgfpathlineto{\pgfqpoint{2.799219in}{1.517163in}}%
\pgfpathlineto{\pgfqpoint{2.801132in}{1.502933in}}%
\pgfpathclose%
\pgfusepath{fill}%
\end{pgfscope}%
\begin{pgfscope}%
\pgfpathrectangle{\pgfqpoint{0.100000in}{0.100000in}}{\pgfqpoint{3.608454in}{2.310000in}}%
\pgfusepath{clip}%
\pgfsetbuttcap%
\pgfsetmiterjoin%
\definecolor{currentfill}{rgb}{0.000000,0.709804,0.645098}%
\pgfsetfillcolor{currentfill}%
\pgfsetlinewidth{0.000000pt}%
\definecolor{currentstroke}{rgb}{0.000000,0.000000,0.000000}%
\pgfsetstrokecolor{currentstroke}%
\pgfsetstrokeopacity{0.000000}%
\pgfsetdash{}{0pt}%
\pgfpathmoveto{\pgfqpoint{3.436777in}{1.747450in}}%
\pgfpathlineto{\pgfqpoint{3.384285in}{1.735025in}}%
\pgfpathlineto{\pgfqpoint{3.380209in}{1.761732in}}%
\pgfpathlineto{\pgfqpoint{3.378669in}{1.782710in}}%
\pgfpathlineto{\pgfqpoint{3.373418in}{1.786388in}}%
\pgfpathlineto{\pgfqpoint{3.377987in}{1.790511in}}%
\pgfpathlineto{\pgfqpoint{3.414849in}{1.798333in}}%
\pgfpathlineto{\pgfqpoint{3.419015in}{1.794601in}}%
\pgfpathlineto{\pgfqpoint{3.418010in}{1.787350in}}%
\pgfpathlineto{\pgfqpoint{3.420462in}{1.768872in}}%
\pgfpathlineto{\pgfqpoint{3.427810in}{1.762986in}}%
\pgfpathlineto{\pgfqpoint{3.426466in}{1.757361in}}%
\pgfpathlineto{\pgfqpoint{3.434236in}{1.757543in}}%
\pgfpathlineto{\pgfqpoint{3.436777in}{1.747450in}}%
\pgfpathclose%
\pgfusepath{fill}%
\end{pgfscope}%
\begin{pgfscope}%
\pgfpathrectangle{\pgfqpoint{0.100000in}{0.100000in}}{\pgfqpoint{3.608454in}{2.310000in}}%
\pgfusepath{clip}%
\pgfsetbuttcap%
\pgfsetmiterjoin%
\definecolor{currentfill}{rgb}{0.000000,0.470588,0.764706}%
\pgfsetfillcolor{currentfill}%
\pgfsetlinewidth{0.000000pt}%
\definecolor{currentstroke}{rgb}{0.000000,0.000000,0.000000}%
\pgfsetstrokecolor{currentstroke}%
\pgfsetstrokeopacity{0.000000}%
\pgfsetdash{}{0pt}%
\pgfpathmoveto{\pgfqpoint{1.938292in}{1.256226in}}%
\pgfpathlineto{\pgfqpoint{1.938903in}{1.276826in}}%
\pgfpathlineto{\pgfqpoint{1.925218in}{1.277253in}}%
\pgfpathlineto{\pgfqpoint{1.926190in}{1.311718in}}%
\pgfpathlineto{\pgfqpoint{1.960306in}{1.310682in}}%
\pgfpathlineto{\pgfqpoint{1.987581in}{1.310037in}}%
\pgfpathlineto{\pgfqpoint{1.987451in}{1.303171in}}%
\pgfpathlineto{\pgfqpoint{1.994270in}{1.303029in}}%
\pgfpathlineto{\pgfqpoint{1.992866in}{1.289314in}}%
\pgfpathlineto{\pgfqpoint{1.992256in}{1.268523in}}%
\pgfpathlineto{\pgfqpoint{1.972846in}{1.269115in}}%
\pgfpathlineto{\pgfqpoint{1.972560in}{1.255285in}}%
\pgfpathlineto{\pgfqpoint{1.938292in}{1.256226in}}%
\pgfpathclose%
\pgfusepath{fill}%
\end{pgfscope}%
\begin{pgfscope}%
\pgfpathrectangle{\pgfqpoint{0.100000in}{0.100000in}}{\pgfqpoint{3.608454in}{2.310000in}}%
\pgfusepath{clip}%
\pgfsetbuttcap%
\pgfsetmiterjoin%
\definecolor{currentfill}{rgb}{0.000000,0.431373,0.784314}%
\pgfsetfillcolor{currentfill}%
\pgfsetlinewidth{0.000000pt}%
\definecolor{currentstroke}{rgb}{0.000000,0.000000,0.000000}%
\pgfsetstrokecolor{currentstroke}%
\pgfsetstrokeopacity{0.000000}%
\pgfsetdash{}{0pt}%
\pgfpathmoveto{\pgfqpoint{1.719085in}{1.966639in}}%
\pgfpathlineto{\pgfqpoint{1.718524in}{1.959750in}}%
\pgfpathlineto{\pgfqpoint{1.673140in}{1.963430in}}%
\pgfpathlineto{\pgfqpoint{1.635984in}{1.966799in}}%
\pgfpathlineto{\pgfqpoint{1.638609in}{1.994610in}}%
\pgfpathlineto{\pgfqpoint{1.635493in}{1.994898in}}%
\pgfpathlineto{\pgfqpoint{1.638167in}{2.022688in}}%
\pgfpathlineto{\pgfqpoint{1.668766in}{2.019808in}}%
\pgfpathlineto{\pgfqpoint{1.671188in}{2.047334in}}%
\pgfpathlineto{\pgfqpoint{1.695708in}{2.045165in}}%
\pgfpathlineto{\pgfqpoint{1.700751in}{2.053267in}}%
\pgfpathlineto{\pgfqpoint{1.710047in}{2.050726in}}%
\pgfpathlineto{\pgfqpoint{1.716886in}{2.052052in}}%
\pgfpathlineto{\pgfqpoint{1.715648in}{2.039417in}}%
\pgfpathlineto{\pgfqpoint{1.718555in}{2.035319in}}%
\pgfpathlineto{\pgfqpoint{1.716982in}{2.015681in}}%
\pgfpathlineto{\pgfqpoint{1.720286in}{2.015429in}}%
\pgfpathlineto{\pgfqpoint{1.718098in}{1.987796in}}%
\pgfpathlineto{\pgfqpoint{1.720698in}{1.987602in}}%
\pgfpathlineto{\pgfqpoint{1.719085in}{1.966639in}}%
\pgfpathclose%
\pgfusepath{fill}%
\end{pgfscope}%
\begin{pgfscope}%
\pgfpathrectangle{\pgfqpoint{0.100000in}{0.100000in}}{\pgfqpoint{3.608454in}{2.310000in}}%
\pgfusepath{clip}%
\pgfsetbuttcap%
\pgfsetmiterjoin%
\definecolor{currentfill}{rgb}{0.000000,0.368627,0.815686}%
\pgfsetfillcolor{currentfill}%
\pgfsetlinewidth{0.000000pt}%
\definecolor{currentstroke}{rgb}{0.000000,0.000000,0.000000}%
\pgfsetstrokecolor{currentstroke}%
\pgfsetstrokeopacity{0.000000}%
\pgfsetdash{}{0pt}%
\pgfpathmoveto{\pgfqpoint{0.934361in}{0.385742in}}%
\pgfpathlineto{\pgfqpoint{0.935737in}{0.388869in}}%
\pgfpathlineto{\pgfqpoint{0.937444in}{0.389743in}}%
\pgfpathlineto{\pgfqpoint{0.940486in}{0.388027in}}%
\pgfpathlineto{\pgfqpoint{0.940016in}{0.386614in}}%
\pgfpathlineto{\pgfqpoint{0.942968in}{0.383955in}}%
\pgfpathlineto{\pgfqpoint{0.936317in}{0.384289in}}%
\pgfpathlineto{\pgfqpoint{0.934361in}{0.385742in}}%
\pgfpathclose%
\pgfusepath{fill}%
\end{pgfscope}%
\begin{pgfscope}%
\pgfpathrectangle{\pgfqpoint{0.100000in}{0.100000in}}{\pgfqpoint{3.608454in}{2.310000in}}%
\pgfusepath{clip}%
\pgfsetbuttcap%
\pgfsetmiterjoin%
\definecolor{currentfill}{rgb}{0.000000,0.368627,0.815686}%
\pgfsetfillcolor{currentfill}%
\pgfsetlinewidth{0.000000pt}%
\definecolor{currentstroke}{rgb}{0.000000,0.000000,0.000000}%
\pgfsetstrokecolor{currentstroke}%
\pgfsetstrokeopacity{0.000000}%
\pgfsetdash{}{0pt}%
\pgfpathmoveto{\pgfqpoint{0.911727in}{0.383608in}}%
\pgfpathlineto{\pgfqpoint{0.914559in}{0.385588in}}%
\pgfpathlineto{\pgfqpoint{0.917234in}{0.385734in}}%
\pgfpathlineto{\pgfqpoint{0.927384in}{0.388734in}}%
\pgfpathlineto{\pgfqpoint{0.928405in}{0.390996in}}%
\pgfpathlineto{\pgfqpoint{0.930625in}{0.389633in}}%
\pgfpathlineto{\pgfqpoint{0.930606in}{0.387345in}}%
\pgfpathlineto{\pgfqpoint{0.924754in}{0.385978in}}%
\pgfpathlineto{\pgfqpoint{0.920480in}{0.383050in}}%
\pgfpathlineto{\pgfqpoint{0.918337in}{0.383207in}}%
\pgfpathlineto{\pgfqpoint{0.917432in}{0.381621in}}%
\pgfpathlineto{\pgfqpoint{0.912959in}{0.382564in}}%
\pgfpathlineto{\pgfqpoint{0.911727in}{0.383608in}}%
\pgfpathclose%
\pgfusepath{fill}%
\end{pgfscope}%
\begin{pgfscope}%
\pgfpathrectangle{\pgfqpoint{0.100000in}{0.100000in}}{\pgfqpoint{3.608454in}{2.310000in}}%
\pgfusepath{clip}%
\pgfsetbuttcap%
\pgfsetmiterjoin%
\definecolor{currentfill}{rgb}{0.000000,0.368627,0.815686}%
\pgfsetfillcolor{currentfill}%
\pgfsetlinewidth{0.000000pt}%
\definecolor{currentstroke}{rgb}{0.000000,0.000000,0.000000}%
\pgfsetstrokecolor{currentstroke}%
\pgfsetstrokeopacity{0.000000}%
\pgfsetdash{}{0pt}%
\pgfpathmoveto{\pgfqpoint{0.918694in}{0.390985in}}%
\pgfpathlineto{\pgfqpoint{0.922043in}{0.397743in}}%
\pgfpathlineto{\pgfqpoint{0.922804in}{0.395503in}}%
\pgfpathlineto{\pgfqpoint{0.918694in}{0.390985in}}%
\pgfpathclose%
\pgfusepath{fill}%
\end{pgfscope}%
\begin{pgfscope}%
\pgfpathrectangle{\pgfqpoint{0.100000in}{0.100000in}}{\pgfqpoint{3.608454in}{2.310000in}}%
\pgfusepath{clip}%
\pgfsetbuttcap%
\pgfsetmiterjoin%
\definecolor{currentfill}{rgb}{0.000000,0.368627,0.815686}%
\pgfsetfillcolor{currentfill}%
\pgfsetlinewidth{0.000000pt}%
\definecolor{currentstroke}{rgb}{0.000000,0.000000,0.000000}%
\pgfsetstrokecolor{currentstroke}%
\pgfsetstrokeopacity{0.000000}%
\pgfsetdash{}{0pt}%
\pgfpathmoveto{\pgfqpoint{0.979580in}{0.452265in}}%
\pgfpathlineto{\pgfqpoint{0.983915in}{0.449521in}}%
\pgfpathlineto{\pgfqpoint{0.986635in}{0.453959in}}%
\pgfpathlineto{\pgfqpoint{0.997579in}{0.447110in}}%
\pgfpathlineto{\pgfqpoint{0.996218in}{0.444880in}}%
\pgfpathlineto{\pgfqpoint{1.002410in}{0.441046in}}%
\pgfpathlineto{\pgfqpoint{1.019938in}{0.430914in}}%
\pgfpathlineto{\pgfqpoint{1.017896in}{0.425014in}}%
\pgfpathlineto{\pgfqpoint{1.014239in}{0.423539in}}%
\pgfpathlineto{\pgfqpoint{1.015064in}{0.421049in}}%
\pgfpathlineto{\pgfqpoint{1.012216in}{0.417308in}}%
\pgfpathlineto{\pgfqpoint{1.015043in}{0.415179in}}%
\pgfpathlineto{\pgfqpoint{1.018284in}{0.417154in}}%
\pgfpathlineto{\pgfqpoint{1.018953in}{0.415851in}}%
\pgfpathlineto{\pgfqpoint{1.022698in}{0.413642in}}%
\pgfpathlineto{\pgfqpoint{1.021385in}{0.411488in}}%
\pgfpathlineto{\pgfqpoint{1.023662in}{0.410127in}}%
\pgfpathlineto{\pgfqpoint{1.022393in}{0.408007in}}%
\pgfpathlineto{\pgfqpoint{1.024400in}{0.406832in}}%
\pgfpathlineto{\pgfqpoint{1.019250in}{0.398079in}}%
\pgfpathlineto{\pgfqpoint{1.020524in}{0.396033in}}%
\pgfpathlineto{\pgfqpoint{1.017303in}{0.390307in}}%
\pgfpathlineto{\pgfqpoint{1.027182in}{0.384801in}}%
\pgfpathlineto{\pgfqpoint{1.016055in}{0.364729in}}%
\pgfpathlineto{\pgfqpoint{1.004608in}{0.344079in}}%
\pgfpathlineto{\pgfqpoint{0.997055in}{0.330455in}}%
\pgfpathlineto{\pgfqpoint{0.993111in}{0.336827in}}%
\pgfpathlineto{\pgfqpoint{0.996872in}{0.337073in}}%
\pgfpathlineto{\pgfqpoint{0.997859in}{0.339919in}}%
\pgfpathlineto{\pgfqpoint{0.995627in}{0.341376in}}%
\pgfpathlineto{\pgfqpoint{0.995160in}{0.339365in}}%
\pgfpathlineto{\pgfqpoint{0.990363in}{0.340190in}}%
\pgfpathlineto{\pgfqpoint{0.988940in}{0.342964in}}%
\pgfpathlineto{\pgfqpoint{0.979357in}{0.351803in}}%
\pgfpathlineto{\pgfqpoint{0.975077in}{0.352806in}}%
\pgfpathlineto{\pgfqpoint{0.969268in}{0.355424in}}%
\pgfpathlineto{\pgfqpoint{0.963603in}{0.357154in}}%
\pgfpathlineto{\pgfqpoint{0.962385in}{0.358792in}}%
\pgfpathlineto{\pgfqpoint{0.962290in}{0.362125in}}%
\pgfpathlineto{\pgfqpoint{0.961040in}{0.364752in}}%
\pgfpathlineto{\pgfqpoint{0.957865in}{0.367144in}}%
\pgfpathlineto{\pgfqpoint{0.958163in}{0.370011in}}%
\pgfpathlineto{\pgfqpoint{0.956916in}{0.372017in}}%
\pgfpathlineto{\pgfqpoint{0.959883in}{0.375282in}}%
\pgfpathlineto{\pgfqpoint{0.960378in}{0.377067in}}%
\pgfpathlineto{\pgfqpoint{0.958642in}{0.378412in}}%
\pgfpathlineto{\pgfqpoint{0.956108in}{0.376672in}}%
\pgfpathlineto{\pgfqpoint{0.953407in}{0.377974in}}%
\pgfpathlineto{\pgfqpoint{0.952146in}{0.377061in}}%
\pgfpathlineto{\pgfqpoint{0.949343in}{0.382116in}}%
\pgfpathlineto{\pgfqpoint{0.947228in}{0.383148in}}%
\pgfpathlineto{\pgfqpoint{0.949313in}{0.385149in}}%
\pgfpathlineto{\pgfqpoint{0.949856in}{0.389896in}}%
\pgfpathlineto{\pgfqpoint{0.949598in}{0.393135in}}%
\pgfpathlineto{\pgfqpoint{0.945394in}{0.394108in}}%
\pgfpathlineto{\pgfqpoint{0.947389in}{0.396132in}}%
\pgfpathlineto{\pgfqpoint{0.942948in}{0.398138in}}%
\pgfpathlineto{\pgfqpoint{0.943546in}{0.402478in}}%
\pgfpathlineto{\pgfqpoint{0.941428in}{0.403633in}}%
\pgfpathlineto{\pgfqpoint{0.940901in}{0.405559in}}%
\pgfpathlineto{\pgfqpoint{0.938307in}{0.405202in}}%
\pgfpathlineto{\pgfqpoint{0.937360in}{0.403019in}}%
\pgfpathlineto{\pgfqpoint{0.933986in}{0.407135in}}%
\pgfpathlineto{\pgfqpoint{0.935525in}{0.410297in}}%
\pgfpathlineto{\pgfqpoint{0.930641in}{0.409132in}}%
\pgfpathlineto{\pgfqpoint{0.924871in}{0.408830in}}%
\pgfpathlineto{\pgfqpoint{0.926703in}{0.411773in}}%
\pgfpathlineto{\pgfqpoint{0.928482in}{0.410664in}}%
\pgfpathlineto{\pgfqpoint{0.928700in}{0.414487in}}%
\pgfpathlineto{\pgfqpoint{0.926958in}{0.414550in}}%
\pgfpathlineto{\pgfqpoint{0.922100in}{0.411389in}}%
\pgfpathlineto{\pgfqpoint{0.918534in}{0.408022in}}%
\pgfpathlineto{\pgfqpoint{0.919688in}{0.407473in}}%
\pgfpathlineto{\pgfqpoint{0.921994in}{0.409270in}}%
\pgfpathlineto{\pgfqpoint{0.923912in}{0.407532in}}%
\pgfpathlineto{\pgfqpoint{0.921409in}{0.405514in}}%
\pgfpathlineto{\pgfqpoint{0.919314in}{0.405095in}}%
\pgfpathlineto{\pgfqpoint{0.916088in}{0.406347in}}%
\pgfpathlineto{\pgfqpoint{0.917614in}{0.402235in}}%
\pgfpathlineto{\pgfqpoint{0.919867in}{0.400817in}}%
\pgfpathlineto{\pgfqpoint{0.913772in}{0.398471in}}%
\pgfpathlineto{\pgfqpoint{0.912875in}{0.397284in}}%
\pgfpathlineto{\pgfqpoint{0.916456in}{0.393000in}}%
\pgfpathlineto{\pgfqpoint{0.915325in}{0.388126in}}%
\pgfpathlineto{\pgfqpoint{0.913318in}{0.388120in}}%
\pgfpathlineto{\pgfqpoint{0.911580in}{0.390660in}}%
\pgfpathlineto{\pgfqpoint{0.912461in}{0.393615in}}%
\pgfpathlineto{\pgfqpoint{0.910952in}{0.394996in}}%
\pgfpathlineto{\pgfqpoint{0.908291in}{0.391990in}}%
\pgfpathlineto{\pgfqpoint{0.906140in}{0.393309in}}%
\pgfpathlineto{\pgfqpoint{0.908663in}{0.396558in}}%
\pgfpathlineto{\pgfqpoint{0.913452in}{0.403668in}}%
\pgfpathlineto{\pgfqpoint{0.912214in}{0.404513in}}%
\pgfpathlineto{\pgfqpoint{0.917407in}{0.411973in}}%
\pgfpathlineto{\pgfqpoint{0.916327in}{0.412673in}}%
\pgfpathlineto{\pgfqpoint{0.918151in}{0.415401in}}%
\pgfpathlineto{\pgfqpoint{0.921451in}{0.413351in}}%
\pgfpathlineto{\pgfqpoint{0.926490in}{0.421058in}}%
\pgfpathlineto{\pgfqpoint{0.930988in}{0.427301in}}%
\pgfpathlineto{\pgfqpoint{0.945694in}{0.417578in}}%
\pgfpathlineto{\pgfqpoint{0.946494in}{0.418805in}}%
\pgfpathlineto{\pgfqpoint{0.949551in}{0.416794in}}%
\pgfpathlineto{\pgfqpoint{0.960218in}{0.433861in}}%
\pgfpathlineto{\pgfqpoint{0.961359in}{0.436162in}}%
\pgfpathlineto{\pgfqpoint{0.967777in}{0.431965in}}%
\pgfpathlineto{\pgfqpoint{0.977356in}{0.447349in}}%
\pgfpathlineto{\pgfqpoint{0.979580in}{0.452265in}}%
\pgfpathclose%
\pgfusepath{fill}%
\end{pgfscope}%
\begin{pgfscope}%
\pgfpathrectangle{\pgfqpoint{0.100000in}{0.100000in}}{\pgfqpoint{3.608454in}{2.310000in}}%
\pgfusepath{clip}%
\pgfsetbuttcap%
\pgfsetmiterjoin%
\definecolor{currentfill}{rgb}{0.000000,0.941176,0.529412}%
\pgfsetfillcolor{currentfill}%
\pgfsetlinewidth{0.000000pt}%
\definecolor{currentstroke}{rgb}{0.000000,0.000000,0.000000}%
\pgfsetstrokecolor{currentstroke}%
\pgfsetstrokeopacity{0.000000}%
\pgfsetdash{}{0pt}%
\pgfpathmoveto{\pgfqpoint{2.923822in}{1.280887in}}%
\pgfpathlineto{\pgfqpoint{2.927867in}{1.284414in}}%
\pgfpathlineto{\pgfqpoint{2.929879in}{1.293625in}}%
\pgfpathlineto{\pgfqpoint{2.935674in}{1.297744in}}%
\pgfpathlineto{\pgfqpoint{2.931210in}{1.301107in}}%
\pgfpathlineto{\pgfqpoint{2.946226in}{1.306488in}}%
\pgfpathlineto{\pgfqpoint{2.952046in}{1.304075in}}%
\pgfpathlineto{\pgfqpoint{2.961369in}{1.291611in}}%
\pgfpathlineto{\pgfqpoint{2.967623in}{1.287269in}}%
\pgfpathlineto{\pgfqpoint{2.965505in}{1.286020in}}%
\pgfpathlineto{\pgfqpoint{2.960443in}{1.270838in}}%
\pgfpathlineto{\pgfqpoint{2.950534in}{1.260497in}}%
\pgfpathlineto{\pgfqpoint{2.942904in}{1.258852in}}%
\pgfpathlineto{\pgfqpoint{2.935481in}{1.264222in}}%
\pgfpathlineto{\pgfqpoint{2.926305in}{1.268226in}}%
\pgfpathlineto{\pgfqpoint{2.923526in}{1.275100in}}%
\pgfpathlineto{\pgfqpoint{2.923822in}{1.280887in}}%
\pgfpathclose%
\pgfusepath{fill}%
\end{pgfscope}%
\begin{pgfscope}%
\pgfpathrectangle{\pgfqpoint{0.100000in}{0.100000in}}{\pgfqpoint{3.608454in}{2.310000in}}%
\pgfusepath{clip}%
\pgfsetbuttcap%
\pgfsetmiterjoin%
\definecolor{currentfill}{rgb}{0.000000,0.443137,0.778431}%
\pgfsetfillcolor{currentfill}%
\pgfsetlinewidth{0.000000pt}%
\definecolor{currentstroke}{rgb}{0.000000,0.000000,0.000000}%
\pgfsetstrokecolor{currentstroke}%
\pgfsetstrokeopacity{0.000000}%
\pgfsetdash{}{0pt}%
\pgfpathmoveto{\pgfqpoint{3.019784in}{0.555093in}}%
\pgfpathlineto{\pgfqpoint{3.017489in}{0.571785in}}%
\pgfpathlineto{\pgfqpoint{3.006224in}{0.574274in}}%
\pgfpathlineto{\pgfqpoint{3.001283in}{0.581332in}}%
\pgfpathlineto{\pgfqpoint{3.006411in}{0.592107in}}%
\pgfpathlineto{\pgfqpoint{2.994626in}{0.604123in}}%
\pgfpathlineto{\pgfqpoint{3.039989in}{0.610672in}}%
\pgfpathlineto{\pgfqpoint{3.040352in}{0.617807in}}%
\pgfpathlineto{\pgfqpoint{3.037437in}{0.635958in}}%
\pgfpathlineto{\pgfqpoint{3.040376in}{0.630406in}}%
\pgfpathlineto{\pgfqpoint{3.046213in}{0.629170in}}%
\pgfpathlineto{\pgfqpoint{3.049410in}{0.623089in}}%
\pgfpathlineto{\pgfqpoint{3.060692in}{0.615927in}}%
\pgfpathlineto{\pgfqpoint{3.065010in}{0.604224in}}%
\pgfpathlineto{\pgfqpoint{3.071577in}{0.605034in}}%
\pgfpathlineto{\pgfqpoint{3.078312in}{0.602797in}}%
\pgfpathlineto{\pgfqpoint{3.080990in}{0.605844in}}%
\pgfpathlineto{\pgfqpoint{3.090875in}{0.590177in}}%
\pgfpathlineto{\pgfqpoint{3.099275in}{0.583081in}}%
\pgfpathlineto{\pgfqpoint{3.099105in}{0.576837in}}%
\pgfpathlineto{\pgfqpoint{3.102843in}{0.570453in}}%
\pgfpathlineto{\pgfqpoint{3.110913in}{0.514275in}}%
\pgfpathlineto{\pgfqpoint{3.091930in}{0.511405in}}%
\pgfpathlineto{\pgfqpoint{3.090421in}{0.517236in}}%
\pgfpathlineto{\pgfqpoint{3.084501in}{0.525466in}}%
\pgfpathlineto{\pgfqpoint{3.077818in}{0.527142in}}%
\pgfpathlineto{\pgfqpoint{3.076539in}{0.531713in}}%
\pgfpathlineto{\pgfqpoint{3.070237in}{0.538537in}}%
\pgfpathlineto{\pgfqpoint{3.059031in}{0.547094in}}%
\pgfpathlineto{\pgfqpoint{3.055318in}{0.556018in}}%
\pgfpathlineto{\pgfqpoint{3.048333in}{0.554969in}}%
\pgfpathlineto{\pgfqpoint{3.047294in}{0.561961in}}%
\pgfpathlineto{\pgfqpoint{3.026291in}{0.558732in}}%
\pgfpathlineto{\pgfqpoint{3.019784in}{0.555093in}}%
\pgfpathclose%
\pgfusepath{fill}%
\end{pgfscope}%
\begin{pgfscope}%
\pgfpathrectangle{\pgfqpoint{0.100000in}{0.100000in}}{\pgfqpoint{3.608454in}{2.310000in}}%
\pgfusepath{clip}%
\pgfsetbuttcap%
\pgfsetmiterjoin%
\definecolor{currentfill}{rgb}{0.000000,0.474510,0.762745}%
\pgfsetfillcolor{currentfill}%
\pgfsetlinewidth{0.000000pt}%
\definecolor{currentstroke}{rgb}{0.000000,0.000000,0.000000}%
\pgfsetstrokecolor{currentstroke}%
\pgfsetstrokeopacity{0.000000}%
\pgfsetdash{}{0pt}%
\pgfpathmoveto{\pgfqpoint{1.723873in}{1.194151in}}%
\pgfpathlineto{\pgfqpoint{1.731511in}{1.193664in}}%
\pgfpathlineto{\pgfqpoint{1.728388in}{1.154348in}}%
\pgfpathlineto{\pgfqpoint{1.685892in}{1.157286in}}%
\pgfpathlineto{\pgfqpoint{1.683362in}{1.122232in}}%
\pgfpathlineto{\pgfqpoint{1.648937in}{1.124836in}}%
\pgfpathlineto{\pgfqpoint{1.651744in}{1.159945in}}%
\pgfpathlineto{\pgfqpoint{1.598511in}{1.164402in}}%
\pgfpathlineto{\pgfqpoint{1.602002in}{1.203777in}}%
\pgfpathlineto{\pgfqpoint{1.621102in}{1.201730in}}%
\pgfpathlineto{\pgfqpoint{1.662385in}{1.198210in}}%
\pgfpathlineto{\pgfqpoint{1.664847in}{1.229439in}}%
\pgfpathlineto{\pgfqpoint{1.695230in}{1.227106in}}%
\pgfpathlineto{\pgfqpoint{1.693075in}{1.196095in}}%
\pgfpathlineto{\pgfqpoint{1.723873in}{1.194151in}}%
\pgfpathclose%
\pgfusepath{fill}%
\end{pgfscope}%
\begin{pgfscope}%
\pgfpathrectangle{\pgfqpoint{0.100000in}{0.100000in}}{\pgfqpoint{3.608454in}{2.310000in}}%
\pgfusepath{clip}%
\pgfsetbuttcap%
\pgfsetmiterjoin%
\definecolor{currentfill}{rgb}{0.000000,0.607843,0.696078}%
\pgfsetfillcolor{currentfill}%
\pgfsetlinewidth{0.000000pt}%
\definecolor{currentstroke}{rgb}{0.000000,0.000000,0.000000}%
\pgfsetstrokecolor{currentstroke}%
\pgfsetstrokeopacity{0.000000}%
\pgfsetdash{}{0pt}%
\pgfpathmoveto{\pgfqpoint{2.642473in}{1.232080in}}%
\pgfpathlineto{\pgfqpoint{2.647960in}{1.220595in}}%
\pgfpathlineto{\pgfqpoint{2.639824in}{1.220359in}}%
\pgfpathlineto{\pgfqpoint{2.634641in}{1.217991in}}%
\pgfpathlineto{\pgfqpoint{2.631929in}{1.205576in}}%
\pgfpathlineto{\pgfqpoint{2.620792in}{1.209072in}}%
\pgfpathlineto{\pgfqpoint{2.607679in}{1.208265in}}%
\pgfpathlineto{\pgfqpoint{2.595443in}{1.205574in}}%
\pgfpathlineto{\pgfqpoint{2.589411in}{1.214310in}}%
\pgfpathlineto{\pgfqpoint{2.585378in}{1.221768in}}%
\pgfpathlineto{\pgfqpoint{2.589982in}{1.231634in}}%
\pgfpathlineto{\pgfqpoint{2.586421in}{1.235460in}}%
\pgfpathlineto{\pgfqpoint{2.585128in}{1.245024in}}%
\pgfpathlineto{\pgfqpoint{2.577210in}{1.248460in}}%
\pgfpathlineto{\pgfqpoint{2.585709in}{1.266483in}}%
\pgfpathlineto{\pgfqpoint{2.586412in}{1.272897in}}%
\pgfpathlineto{\pgfqpoint{2.601681in}{1.267043in}}%
\pgfpathlineto{\pgfqpoint{2.602387in}{1.273595in}}%
\pgfpathlineto{\pgfqpoint{2.617290in}{1.282808in}}%
\pgfpathlineto{\pgfqpoint{2.620883in}{1.276667in}}%
\pgfpathlineto{\pgfqpoint{2.627054in}{1.277303in}}%
\pgfpathlineto{\pgfqpoint{2.627023in}{1.272051in}}%
\pgfpathlineto{\pgfqpoint{2.625972in}{1.268011in}}%
\pgfpathlineto{\pgfqpoint{2.630443in}{1.257485in}}%
\pgfpathlineto{\pgfqpoint{2.628815in}{1.247944in}}%
\pgfpathlineto{\pgfqpoint{2.632971in}{1.237093in}}%
\pgfpathlineto{\pgfqpoint{2.641029in}{1.235231in}}%
\pgfpathlineto{\pgfqpoint{2.642473in}{1.232080in}}%
\pgfpathclose%
\pgfusepath{fill}%
\end{pgfscope}%
\begin{pgfscope}%
\pgfpathrectangle{\pgfqpoint{0.100000in}{0.100000in}}{\pgfqpoint{3.608454in}{2.310000in}}%
\pgfusepath{clip}%
\pgfsetbuttcap%
\pgfsetmiterjoin%
\definecolor{currentfill}{rgb}{0.000000,0.415686,0.792157}%
\pgfsetfillcolor{currentfill}%
\pgfsetlinewidth{0.000000pt}%
\definecolor{currentstroke}{rgb}{0.000000,0.000000,0.000000}%
\pgfsetstrokecolor{currentstroke}%
\pgfsetstrokeopacity{0.000000}%
\pgfsetdash{}{0pt}%
\pgfpathmoveto{\pgfqpoint{2.218568in}{1.724511in}}%
\pgfpathlineto{\pgfqpoint{2.238913in}{1.724992in}}%
\pgfpathlineto{\pgfqpoint{2.252724in}{1.724220in}}%
\pgfpathlineto{\pgfqpoint{2.253472in}{1.697776in}}%
\pgfpathlineto{\pgfqpoint{2.220391in}{1.696881in}}%
\pgfpathlineto{\pgfqpoint{2.184501in}{1.696223in}}%
\pgfpathlineto{\pgfqpoint{2.184096in}{1.723863in}}%
\pgfpathlineto{\pgfqpoint{2.218568in}{1.724511in}}%
\pgfpathclose%
\pgfusepath{fill}%
\end{pgfscope}%
\begin{pgfscope}%
\pgfpathrectangle{\pgfqpoint{0.100000in}{0.100000in}}{\pgfqpoint{3.608454in}{2.310000in}}%
\pgfusepath{clip}%
\pgfsetbuttcap%
\pgfsetmiterjoin%
\definecolor{currentfill}{rgb}{0.000000,0.207843,0.896078}%
\pgfsetfillcolor{currentfill}%
\pgfsetlinewidth{0.000000pt}%
\definecolor{currentstroke}{rgb}{0.000000,0.000000,0.000000}%
\pgfsetstrokecolor{currentstroke}%
\pgfsetstrokeopacity{0.000000}%
\pgfsetdash{}{0pt}%
\pgfpathmoveto{\pgfqpoint{1.989587in}{1.638444in}}%
\pgfpathlineto{\pgfqpoint{1.981253in}{1.641675in}}%
\pgfpathlineto{\pgfqpoint{1.981953in}{1.671006in}}%
\pgfpathlineto{\pgfqpoint{1.968146in}{1.671305in}}%
\pgfpathlineto{\pgfqpoint{1.968709in}{1.697472in}}%
\pgfpathlineto{\pgfqpoint{1.956891in}{1.697804in}}%
\pgfpathlineto{\pgfqpoint{1.957610in}{1.725492in}}%
\pgfpathlineto{\pgfqpoint{1.998767in}{1.724416in}}%
\pgfpathlineto{\pgfqpoint{2.023670in}{1.724067in}}%
\pgfpathlineto{\pgfqpoint{2.023239in}{1.696380in}}%
\pgfpathlineto{\pgfqpoint{2.014851in}{1.696524in}}%
\pgfpathlineto{\pgfqpoint{2.015077in}{1.691199in}}%
\pgfpathlineto{\pgfqpoint{2.019001in}{1.688123in}}%
\pgfpathlineto{\pgfqpoint{2.018524in}{1.680571in}}%
\pgfpathlineto{\pgfqpoint{2.021565in}{1.674298in}}%
\pgfpathlineto{\pgfqpoint{2.022652in}{1.663329in}}%
\pgfpathlineto{\pgfqpoint{2.002309in}{1.663696in}}%
\pgfpathlineto{\pgfqpoint{2.001679in}{1.633660in}}%
\pgfpathlineto{\pgfqpoint{1.989587in}{1.638444in}}%
\pgfpathclose%
\pgfusepath{fill}%
\end{pgfscope}%
\begin{pgfscope}%
\pgfpathrectangle{\pgfqpoint{0.100000in}{0.100000in}}{\pgfqpoint{3.608454in}{2.310000in}}%
\pgfusepath{clip}%
\pgfsetbuttcap%
\pgfsetmiterjoin%
\definecolor{currentfill}{rgb}{0.000000,0.490196,0.754902}%
\pgfsetfillcolor{currentfill}%
\pgfsetlinewidth{0.000000pt}%
\definecolor{currentstroke}{rgb}{0.000000,0.000000,0.000000}%
\pgfsetstrokecolor{currentstroke}%
\pgfsetstrokeopacity{0.000000}%
\pgfsetdash{}{0pt}%
\pgfpathmoveto{\pgfqpoint{2.018170in}{1.419481in}}%
\pgfpathlineto{\pgfqpoint{2.031759in}{1.419257in}}%
\pgfpathlineto{\pgfqpoint{2.031304in}{1.384920in}}%
\pgfpathlineto{\pgfqpoint{1.962491in}{1.386340in}}%
\pgfpathlineto{\pgfqpoint{1.963356in}{1.420723in}}%
\pgfpathlineto{\pgfqpoint{2.018170in}{1.419481in}}%
\pgfpathclose%
\pgfusepath{fill}%
\end{pgfscope}%
\begin{pgfscope}%
\pgfpathrectangle{\pgfqpoint{0.100000in}{0.100000in}}{\pgfqpoint{3.608454in}{2.310000in}}%
\pgfusepath{clip}%
\pgfsetbuttcap%
\pgfsetmiterjoin%
\definecolor{currentfill}{rgb}{0.000000,0.592157,0.703922}%
\pgfsetfillcolor{currentfill}%
\pgfsetlinewidth{0.000000pt}%
\definecolor{currentstroke}{rgb}{0.000000,0.000000,0.000000}%
\pgfsetstrokecolor{currentstroke}%
\pgfsetstrokeopacity{0.000000}%
\pgfsetdash{}{0pt}%
\pgfpathmoveto{\pgfqpoint{2.396575in}{1.440884in}}%
\pgfpathlineto{\pgfqpoint{2.381360in}{1.440493in}}%
\pgfpathlineto{\pgfqpoint{2.381085in}{1.447410in}}%
\pgfpathlineto{\pgfqpoint{2.353369in}{1.446719in}}%
\pgfpathlineto{\pgfqpoint{2.352440in}{1.474762in}}%
\pgfpathlineto{\pgfqpoint{2.359331in}{1.474799in}}%
\pgfpathlineto{\pgfqpoint{2.358001in}{1.509016in}}%
\pgfpathlineto{\pgfqpoint{2.378619in}{1.509638in}}%
\pgfpathlineto{\pgfqpoint{2.378376in}{1.516534in}}%
\pgfpathlineto{\pgfqpoint{2.405330in}{1.517808in}}%
\pgfpathlineto{\pgfqpoint{2.406029in}{1.504012in}}%
\pgfpathlineto{\pgfqpoint{2.407342in}{1.476448in}}%
\pgfpathlineto{\pgfqpoint{2.414274in}{1.476742in}}%
\pgfpathlineto{\pgfqpoint{2.414833in}{1.467947in}}%
\pgfpathlineto{\pgfqpoint{2.412968in}{1.462469in}}%
\pgfpathlineto{\pgfqpoint{2.405497in}{1.456494in}}%
\pgfpathlineto{\pgfqpoint{2.401121in}{1.445010in}}%
\pgfpathlineto{\pgfqpoint{2.396575in}{1.440884in}}%
\pgfpathclose%
\pgfusepath{fill}%
\end{pgfscope}%
\begin{pgfscope}%
\pgfpathrectangle{\pgfqpoint{0.100000in}{0.100000in}}{\pgfqpoint{3.608454in}{2.310000in}}%
\pgfusepath{clip}%
\pgfsetbuttcap%
\pgfsetmiterjoin%
\definecolor{currentfill}{rgb}{0.000000,0.572549,0.713725}%
\pgfsetfillcolor{currentfill}%
\pgfsetlinewidth{0.000000pt}%
\definecolor{currentstroke}{rgb}{0.000000,0.000000,0.000000}%
\pgfsetstrokecolor{currentstroke}%
\pgfsetstrokeopacity{0.000000}%
\pgfsetdash{}{0pt}%
\pgfpathmoveto{\pgfqpoint{2.729948in}{1.258824in}}%
\pgfpathlineto{\pgfqpoint{2.730780in}{1.248329in}}%
\pgfpathlineto{\pgfqpoint{2.722412in}{1.251628in}}%
\pgfpathlineto{\pgfqpoint{2.718150in}{1.247622in}}%
\pgfpathlineto{\pgfqpoint{2.711292in}{1.250974in}}%
\pgfpathlineto{\pgfqpoint{2.704025in}{1.249697in}}%
\pgfpathlineto{\pgfqpoint{2.701382in}{1.255633in}}%
\pgfpathlineto{\pgfqpoint{2.694535in}{1.255275in}}%
\pgfpathlineto{\pgfqpoint{2.694291in}{1.266816in}}%
\pgfpathlineto{\pgfqpoint{2.691418in}{1.266749in}}%
\pgfpathlineto{\pgfqpoint{2.684273in}{1.275179in}}%
\pgfpathlineto{\pgfqpoint{2.690035in}{1.277558in}}%
\pgfpathlineto{\pgfqpoint{2.698255in}{1.290856in}}%
\pgfpathlineto{\pgfqpoint{2.704197in}{1.287428in}}%
\pgfpathlineto{\pgfqpoint{2.718323in}{1.291481in}}%
\pgfpathlineto{\pgfqpoint{2.719889in}{1.285770in}}%
\pgfpathlineto{\pgfqpoint{2.727518in}{1.286122in}}%
\pgfpathlineto{\pgfqpoint{2.729699in}{1.283399in}}%
\pgfpathlineto{\pgfqpoint{2.729948in}{1.258824in}}%
\pgfpathclose%
\pgfusepath{fill}%
\end{pgfscope}%
\begin{pgfscope}%
\pgfpathrectangle{\pgfqpoint{0.100000in}{0.100000in}}{\pgfqpoint{3.608454in}{2.310000in}}%
\pgfusepath{clip}%
\pgfsetbuttcap%
\pgfsetmiterjoin%
\definecolor{currentfill}{rgb}{0.000000,0.529412,0.735294}%
\pgfsetfillcolor{currentfill}%
\pgfsetlinewidth{0.000000pt}%
\definecolor{currentstroke}{rgb}{0.000000,0.000000,0.000000}%
\pgfsetstrokecolor{currentstroke}%
\pgfsetstrokeopacity{0.000000}%
\pgfsetdash{}{0pt}%
\pgfpathmoveto{\pgfqpoint{2.868390in}{1.010555in}}%
\pgfpathlineto{\pgfqpoint{2.871596in}{1.014173in}}%
\pgfpathlineto{\pgfqpoint{2.886412in}{1.013136in}}%
\pgfpathlineto{\pgfqpoint{2.885344in}{1.008668in}}%
\pgfpathlineto{\pgfqpoint{2.897463in}{0.997613in}}%
\pgfpathlineto{\pgfqpoint{2.908530in}{0.995600in}}%
\pgfpathlineto{\pgfqpoint{2.902891in}{0.990956in}}%
\pgfpathlineto{\pgfqpoint{2.896953in}{0.978821in}}%
\pgfpathlineto{\pgfqpoint{2.896071in}{0.971539in}}%
\pgfpathlineto{\pgfqpoint{2.889208in}{0.971223in}}%
\pgfpathlineto{\pgfqpoint{2.877953in}{0.974806in}}%
\pgfpathlineto{\pgfqpoint{2.870396in}{0.968762in}}%
\pgfpathlineto{\pgfqpoint{2.864261in}{0.976571in}}%
\pgfpathlineto{\pgfqpoint{2.850509in}{0.983786in}}%
\pgfpathlineto{\pgfqpoint{2.844938in}{0.981568in}}%
\pgfpathlineto{\pgfqpoint{2.837086in}{0.989053in}}%
\pgfpathlineto{\pgfqpoint{2.839244in}{0.999162in}}%
\pgfpathlineto{\pgfqpoint{2.843269in}{0.994686in}}%
\pgfpathlineto{\pgfqpoint{2.856780in}{0.993705in}}%
\pgfpathlineto{\pgfqpoint{2.859121in}{0.988728in}}%
\pgfpathlineto{\pgfqpoint{2.868081in}{0.996615in}}%
\pgfpathlineto{\pgfqpoint{2.869364in}{1.002057in}}%
\pgfpathlineto{\pgfqpoint{2.865578in}{1.008074in}}%
\pgfpathlineto{\pgfqpoint{2.868390in}{1.010555in}}%
\pgfpathclose%
\pgfusepath{fill}%
\end{pgfscope}%
\begin{pgfscope}%
\pgfpathrectangle{\pgfqpoint{0.100000in}{0.100000in}}{\pgfqpoint{3.608454in}{2.310000in}}%
\pgfusepath{clip}%
\pgfsetbuttcap%
\pgfsetmiterjoin%
\definecolor{currentfill}{rgb}{0.000000,0.364706,0.817647}%
\pgfsetfillcolor{currentfill}%
\pgfsetlinewidth{0.000000pt}%
\definecolor{currentstroke}{rgb}{0.000000,0.000000,0.000000}%
\pgfsetstrokecolor{currentstroke}%
\pgfsetstrokeopacity{0.000000}%
\pgfsetdash{}{0pt}%
\pgfpathmoveto{\pgfqpoint{1.858155in}{1.589614in}}%
\pgfpathlineto{\pgfqpoint{1.831612in}{1.590776in}}%
\pgfpathlineto{\pgfqpoint{1.803419in}{1.592288in}}%
\pgfpathlineto{\pgfqpoint{1.804075in}{1.619812in}}%
\pgfpathlineto{\pgfqpoint{1.805638in}{1.647292in}}%
\pgfpathlineto{\pgfqpoint{1.805743in}{1.664546in}}%
\pgfpathlineto{\pgfqpoint{1.860373in}{1.661671in}}%
\pgfpathlineto{\pgfqpoint{1.858155in}{1.589614in}}%
\pgfpathclose%
\pgfusepath{fill}%
\end{pgfscope}%
\begin{pgfscope}%
\pgfpathrectangle{\pgfqpoint{0.100000in}{0.100000in}}{\pgfqpoint{3.608454in}{2.310000in}}%
\pgfusepath{clip}%
\pgfsetbuttcap%
\pgfsetmiterjoin%
\definecolor{currentfill}{rgb}{0.000000,0.325490,0.837255}%
\pgfsetfillcolor{currentfill}%
\pgfsetlinewidth{0.000000pt}%
\definecolor{currentstroke}{rgb}{0.000000,0.000000,0.000000}%
\pgfsetstrokecolor{currentstroke}%
\pgfsetstrokeopacity{0.000000}%
\pgfsetdash{}{0pt}%
\pgfpathmoveto{\pgfqpoint{2.039668in}{1.566696in}}%
\pgfpathlineto{\pgfqpoint{2.043988in}{1.561009in}}%
\pgfpathlineto{\pgfqpoint{2.040513in}{1.552171in}}%
\pgfpathlineto{\pgfqpoint{2.021430in}{1.552509in}}%
\pgfpathlineto{\pgfqpoint{2.021531in}{1.557101in}}%
\pgfpathlineto{\pgfqpoint{1.994270in}{1.557682in}}%
\pgfpathlineto{\pgfqpoint{1.987553in}{1.557852in}}%
\pgfpathlineto{\pgfqpoint{1.988196in}{1.585344in}}%
\pgfpathlineto{\pgfqpoint{1.999669in}{1.585100in}}%
\pgfpathlineto{\pgfqpoint{2.000030in}{1.598859in}}%
\pgfpathlineto{\pgfqpoint{2.005666in}{1.599833in}}%
\pgfpathlineto{\pgfqpoint{2.027350in}{1.599345in}}%
\pgfpathlineto{\pgfqpoint{2.027934in}{1.590280in}}%
\pgfpathlineto{\pgfqpoint{2.036665in}{1.576098in}}%
\pgfpathlineto{\pgfqpoint{2.040288in}{1.575134in}}%
\pgfpathlineto{\pgfqpoint{2.039668in}{1.566696in}}%
\pgfpathclose%
\pgfusepath{fill}%
\end{pgfscope}%
\begin{pgfscope}%
\pgfpathrectangle{\pgfqpoint{0.100000in}{0.100000in}}{\pgfqpoint{3.608454in}{2.310000in}}%
\pgfusepath{clip}%
\pgfsetbuttcap%
\pgfsetmiterjoin%
\definecolor{currentfill}{rgb}{0.000000,0.352941,0.823529}%
\pgfsetfillcolor{currentfill}%
\pgfsetlinewidth{0.000000pt}%
\definecolor{currentstroke}{rgb}{0.000000,0.000000,0.000000}%
\pgfsetstrokecolor{currentstroke}%
\pgfsetstrokeopacity{0.000000}%
\pgfsetdash{}{0pt}%
\pgfpathmoveto{\pgfqpoint{2.258667in}{1.780519in}}%
\pgfpathlineto{\pgfqpoint{2.247155in}{1.782906in}}%
\pgfpathlineto{\pgfqpoint{2.241126in}{1.786250in}}%
\pgfpathlineto{\pgfqpoint{2.231530in}{1.786839in}}%
\pgfpathlineto{\pgfqpoint{2.231724in}{1.779996in}}%
\pgfpathlineto{\pgfqpoint{2.217797in}{1.777469in}}%
\pgfpathlineto{\pgfqpoint{2.217878in}{1.773970in}}%
\pgfpathlineto{\pgfqpoint{2.204133in}{1.773697in}}%
\pgfpathlineto{\pgfqpoint{2.204050in}{1.779422in}}%
\pgfpathlineto{\pgfqpoint{2.190305in}{1.779149in}}%
\pgfpathlineto{\pgfqpoint{2.168433in}{1.778817in}}%
\pgfpathlineto{\pgfqpoint{2.169520in}{1.792650in}}%
\pgfpathlineto{\pgfqpoint{2.162655in}{1.792587in}}%
\pgfpathlineto{\pgfqpoint{2.162319in}{1.813349in}}%
\pgfpathlineto{\pgfqpoint{2.148622in}{1.813278in}}%
\pgfpathlineto{\pgfqpoint{2.148148in}{1.837483in}}%
\pgfpathlineto{\pgfqpoint{2.156256in}{1.841131in}}%
\pgfpathlineto{\pgfqpoint{2.160005in}{1.848588in}}%
\pgfpathlineto{\pgfqpoint{2.176492in}{1.838639in}}%
\pgfpathlineto{\pgfqpoint{2.186149in}{1.838564in}}%
\pgfpathlineto{\pgfqpoint{2.190082in}{1.834950in}}%
\pgfpathlineto{\pgfqpoint{2.189353in}{1.873824in}}%
\pgfpathlineto{\pgfqpoint{2.226726in}{1.874225in}}%
\pgfpathlineto{\pgfqpoint{2.224641in}{1.861874in}}%
\pgfpathlineto{\pgfqpoint{2.230979in}{1.861346in}}%
\pgfpathlineto{\pgfqpoint{2.238078in}{1.851709in}}%
\pgfpathlineto{\pgfqpoint{2.237996in}{1.848119in}}%
\pgfpathlineto{\pgfqpoint{2.232161in}{1.839416in}}%
\pgfpathlineto{\pgfqpoint{2.232360in}{1.832947in}}%
\pgfpathlineto{\pgfqpoint{2.266038in}{1.833884in}}%
\pgfpathlineto{\pgfqpoint{2.267401in}{1.826914in}}%
\pgfpathlineto{\pgfqpoint{2.268903in}{1.780718in}}%
\pgfpathlineto{\pgfqpoint{2.258667in}{1.780519in}}%
\pgfpathclose%
\pgfusepath{fill}%
\end{pgfscope}%
\begin{pgfscope}%
\pgfpathrectangle{\pgfqpoint{0.100000in}{0.100000in}}{\pgfqpoint{3.608454in}{2.310000in}}%
\pgfusepath{clip}%
\pgfsetbuttcap%
\pgfsetmiterjoin%
\definecolor{currentfill}{rgb}{0.000000,0.482353,0.758824}%
\pgfsetfillcolor{currentfill}%
\pgfsetlinewidth{0.000000pt}%
\definecolor{currentstroke}{rgb}{0.000000,0.000000,0.000000}%
\pgfsetstrokecolor{currentstroke}%
\pgfsetstrokeopacity{0.000000}%
\pgfsetdash{}{0pt}%
\pgfpathmoveto{\pgfqpoint{2.424922in}{0.713419in}}%
\pgfpathlineto{\pgfqpoint{2.432451in}{0.691509in}}%
\pgfpathlineto{\pgfqpoint{2.435277in}{0.651886in}}%
\pgfpathlineto{\pgfqpoint{2.431222in}{0.657222in}}%
\pgfpathlineto{\pgfqpoint{2.419484in}{0.654679in}}%
\pgfpathlineto{\pgfqpoint{2.415540in}{0.649843in}}%
\pgfpathlineto{\pgfqpoint{2.407987in}{0.645878in}}%
\pgfpathlineto{\pgfqpoint{2.389280in}{0.640751in}}%
\pgfpathlineto{\pgfqpoint{2.386115in}{0.636858in}}%
\pgfpathlineto{\pgfqpoint{2.383161in}{0.637744in}}%
\pgfpathlineto{\pgfqpoint{2.368222in}{0.632799in}}%
\pgfpathlineto{\pgfqpoint{2.351410in}{0.638299in}}%
\pgfpathlineto{\pgfqpoint{2.350223in}{0.648417in}}%
\pgfpathlineto{\pgfqpoint{2.344148in}{0.649181in}}%
\pgfpathlineto{\pgfqpoint{2.339922in}{0.654462in}}%
\pgfpathlineto{\pgfqpoint{2.338216in}{0.665019in}}%
\pgfpathlineto{\pgfqpoint{2.331431in}{0.670010in}}%
\pgfpathlineto{\pgfqpoint{2.329801in}{0.674430in}}%
\pgfpathlineto{\pgfqpoint{2.331030in}{0.685463in}}%
\pgfpathlineto{\pgfqpoint{2.325022in}{0.696944in}}%
\pgfpathlineto{\pgfqpoint{2.325328in}{0.707042in}}%
\pgfpathlineto{\pgfqpoint{2.330722in}{0.712813in}}%
\pgfpathlineto{\pgfqpoint{2.336789in}{0.709454in}}%
\pgfpathlineto{\pgfqpoint{2.392212in}{0.711750in}}%
\pgfpathlineto{\pgfqpoint{2.424922in}{0.713419in}}%
\pgfpathclose%
\pgfusepath{fill}%
\end{pgfscope}%
\begin{pgfscope}%
\pgfpathrectangle{\pgfqpoint{0.100000in}{0.100000in}}{\pgfqpoint{3.608454in}{2.310000in}}%
\pgfusepath{clip}%
\pgfsetbuttcap%
\pgfsetmiterjoin%
\definecolor{currentfill}{rgb}{0.000000,0.592157,0.703922}%
\pgfsetfillcolor{currentfill}%
\pgfsetlinewidth{0.000000pt}%
\definecolor{currentstroke}{rgb}{0.000000,0.000000,0.000000}%
\pgfsetstrokecolor{currentstroke}%
\pgfsetstrokeopacity{0.000000}%
\pgfsetdash{}{0pt}%
\pgfpathmoveto{\pgfqpoint{1.895416in}{1.185262in}}%
\pgfpathlineto{\pgfqpoint{1.854168in}{1.187060in}}%
\pgfpathlineto{\pgfqpoint{1.855479in}{1.217455in}}%
\pgfpathlineto{\pgfqpoint{1.821457in}{1.218805in}}%
\pgfpathlineto{\pgfqpoint{1.821831in}{1.246792in}}%
\pgfpathlineto{\pgfqpoint{1.855927in}{1.245032in}}%
\pgfpathlineto{\pgfqpoint{1.856262in}{1.252307in}}%
\pgfpathlineto{\pgfqpoint{1.890012in}{1.250882in}}%
\pgfpathlineto{\pgfqpoint{1.890205in}{1.243605in}}%
\pgfpathlineto{\pgfqpoint{1.889149in}{1.216054in}}%
\pgfpathlineto{\pgfqpoint{1.896376in}{1.215785in}}%
\pgfpathlineto{\pgfqpoint{1.895416in}{1.185262in}}%
\pgfpathclose%
\pgfusepath{fill}%
\end{pgfscope}%
\begin{pgfscope}%
\pgfpathrectangle{\pgfqpoint{0.100000in}{0.100000in}}{\pgfqpoint{3.608454in}{2.310000in}}%
\pgfusepath{clip}%
\pgfsetbuttcap%
\pgfsetmiterjoin%
\definecolor{currentfill}{rgb}{0.000000,0.427451,0.786275}%
\pgfsetfillcolor{currentfill}%
\pgfsetlinewidth{0.000000pt}%
\definecolor{currentstroke}{rgb}{0.000000,0.000000,0.000000}%
\pgfsetstrokecolor{currentstroke}%
\pgfsetstrokeopacity{0.000000}%
\pgfsetdash{}{0pt}%
\pgfpathmoveto{\pgfqpoint{2.499443in}{1.713874in}}%
\pgfpathlineto{\pgfqpoint{2.485729in}{1.712888in}}%
\pgfpathlineto{\pgfqpoint{2.485186in}{1.719799in}}%
\pgfpathlineto{\pgfqpoint{2.450492in}{1.717550in}}%
\pgfpathlineto{\pgfqpoint{2.436777in}{1.717456in}}%
\pgfpathlineto{\pgfqpoint{2.436156in}{1.727226in}}%
\pgfpathlineto{\pgfqpoint{2.440017in}{1.736230in}}%
\pgfpathlineto{\pgfqpoint{2.439466in}{1.744650in}}%
\pgfpathlineto{\pgfqpoint{2.414986in}{1.743088in}}%
\pgfpathlineto{\pgfqpoint{2.413744in}{1.763959in}}%
\pgfpathlineto{\pgfqpoint{2.434927in}{1.765090in}}%
\pgfpathlineto{\pgfqpoint{2.432784in}{1.799789in}}%
\pgfpathlineto{\pgfqpoint{2.467658in}{1.801883in}}%
\pgfpathlineto{\pgfqpoint{2.468096in}{1.794922in}}%
\pgfpathlineto{\pgfqpoint{2.488486in}{1.795943in}}%
\pgfpathlineto{\pgfqpoint{2.491553in}{1.796332in}}%
\pgfpathlineto{\pgfqpoint{2.493489in}{1.769026in}}%
\pgfpathlineto{\pgfqpoint{2.501961in}{1.769586in}}%
\pgfpathlineto{\pgfqpoint{2.504166in}{1.741961in}}%
\pgfpathlineto{\pgfqpoint{2.497309in}{1.741424in}}%
\pgfpathlineto{\pgfqpoint{2.499443in}{1.713874in}}%
\pgfpathclose%
\pgfusepath{fill}%
\end{pgfscope}%
\begin{pgfscope}%
\pgfpathrectangle{\pgfqpoint{0.100000in}{0.100000in}}{\pgfqpoint{3.608454in}{2.310000in}}%
\pgfusepath{clip}%
\pgfsetbuttcap%
\pgfsetmiterjoin%
\definecolor{currentfill}{rgb}{0.000000,0.364706,0.817647}%
\pgfsetfillcolor{currentfill}%
\pgfsetlinewidth{0.000000pt}%
\definecolor{currentstroke}{rgb}{0.000000,0.000000,0.000000}%
\pgfsetstrokecolor{currentstroke}%
\pgfsetstrokeopacity{0.000000}%
\pgfsetdash{}{0pt}%
\pgfpathmoveto{\pgfqpoint{3.147047in}{0.962987in}}%
\pgfpathlineto{\pgfqpoint{3.144033in}{0.954523in}}%
\pgfpathlineto{\pgfqpoint{3.134057in}{0.954802in}}%
\pgfpathlineto{\pgfqpoint{3.127866in}{0.947248in}}%
\pgfpathlineto{\pgfqpoint{3.131225in}{0.944243in}}%
\pgfpathlineto{\pgfqpoint{3.123999in}{0.938426in}}%
\pgfpathlineto{\pgfqpoint{3.122700in}{0.934730in}}%
\pgfpathlineto{\pgfqpoint{3.114162in}{0.927752in}}%
\pgfpathlineto{\pgfqpoint{3.113774in}{0.923333in}}%
\pgfpathlineto{\pgfqpoint{3.106625in}{0.916011in}}%
\pgfpathlineto{\pgfqpoint{3.099237in}{0.913335in}}%
\pgfpathlineto{\pgfqpoint{3.086793in}{0.902214in}}%
\pgfpathlineto{\pgfqpoint{3.077208in}{0.901258in}}%
\pgfpathlineto{\pgfqpoint{3.071225in}{0.906814in}}%
\pgfpathlineto{\pgfqpoint{3.066440in}{0.907291in}}%
\pgfpathlineto{\pgfqpoint{3.061084in}{0.913825in}}%
\pgfpathlineto{\pgfqpoint{3.051321in}{0.914811in}}%
\pgfpathlineto{\pgfqpoint{3.044811in}{0.924182in}}%
\pgfpathlineto{\pgfqpoint{3.038279in}{0.928637in}}%
\pgfpathlineto{\pgfqpoint{3.030466in}{0.937480in}}%
\pgfpathlineto{\pgfqpoint{3.047350in}{0.952216in}}%
\pgfpathlineto{\pgfqpoint{3.064743in}{0.967659in}}%
\pgfpathlineto{\pgfqpoint{3.070676in}{0.962442in}}%
\pgfpathlineto{\pgfqpoint{3.081498in}{0.967552in}}%
\pgfpathlineto{\pgfqpoint{3.081713in}{0.979204in}}%
\pgfpathlineto{\pgfqpoint{3.088906in}{0.984668in}}%
\pgfpathlineto{\pgfqpoint{3.097264in}{0.986382in}}%
\pgfpathlineto{\pgfqpoint{3.115214in}{0.974844in}}%
\pgfpathlineto{\pgfqpoint{3.134423in}{0.971285in}}%
\pgfpathlineto{\pgfqpoint{3.147047in}{0.962987in}}%
\pgfpathclose%
\pgfusepath{fill}%
\end{pgfscope}%
\begin{pgfscope}%
\pgfpathrectangle{\pgfqpoint{0.100000in}{0.100000in}}{\pgfqpoint{3.608454in}{2.310000in}}%
\pgfusepath{clip}%
\pgfsetbuttcap%
\pgfsetmiterjoin%
\definecolor{currentfill}{rgb}{0.000000,0.356863,0.821569}%
\pgfsetfillcolor{currentfill}%
\pgfsetlinewidth{0.000000pt}%
\definecolor{currentstroke}{rgb}{0.000000,0.000000,0.000000}%
\pgfsetstrokecolor{currentstroke}%
\pgfsetstrokeopacity{0.000000}%
\pgfsetdash{}{0pt}%
\pgfpathmoveto{\pgfqpoint{3.191049in}{1.170900in}}%
\pgfpathlineto{\pgfqpoint{3.187783in}{1.156766in}}%
\pgfpathlineto{\pgfqpoint{3.188860in}{1.151551in}}%
\pgfpathlineto{\pgfqpoint{3.180850in}{1.147206in}}%
\pgfpathlineto{\pgfqpoint{3.173892in}{1.141190in}}%
\pgfpathlineto{\pgfqpoint{3.168634in}{1.140787in}}%
\pgfpathlineto{\pgfqpoint{3.164716in}{1.144245in}}%
\pgfpathlineto{\pgfqpoint{3.160885in}{1.137946in}}%
\pgfpathlineto{\pgfqpoint{3.149312in}{1.137485in}}%
\pgfpathlineto{\pgfqpoint{3.131128in}{1.128295in}}%
\pgfpathlineto{\pgfqpoint{3.121994in}{1.132796in}}%
\pgfpathlineto{\pgfqpoint{3.110689in}{1.147372in}}%
\pgfpathlineto{\pgfqpoint{3.110510in}{1.150879in}}%
\pgfpathlineto{\pgfqpoint{3.097477in}{1.148400in}}%
\pgfpathlineto{\pgfqpoint{3.093928in}{1.174159in}}%
\pgfpathlineto{\pgfqpoint{3.113209in}{1.177564in}}%
\pgfpathlineto{\pgfqpoint{3.110627in}{1.182204in}}%
\pgfpathlineto{\pgfqpoint{3.106519in}{1.208514in}}%
\pgfpathlineto{\pgfqpoint{3.113032in}{1.209517in}}%
\pgfpathlineto{\pgfqpoint{3.109908in}{1.233039in}}%
\pgfpathlineto{\pgfqpoint{3.152496in}{1.240772in}}%
\pgfpathlineto{\pgfqpoint{3.150361in}{1.233286in}}%
\pgfpathlineto{\pgfqpoint{3.154402in}{1.212187in}}%
\pgfpathlineto{\pgfqpoint{3.160658in}{1.212280in}}%
\pgfpathlineto{\pgfqpoint{3.165876in}{1.221007in}}%
\pgfpathlineto{\pgfqpoint{3.177079in}{1.221508in}}%
\pgfpathlineto{\pgfqpoint{3.185638in}{1.219606in}}%
\pgfpathlineto{\pgfqpoint{3.175634in}{1.186715in}}%
\pgfpathlineto{\pgfqpoint{3.181671in}{1.180528in}}%
\pgfpathlineto{\pgfqpoint{3.186933in}{1.171546in}}%
\pgfpathlineto{\pgfqpoint{3.191049in}{1.170900in}}%
\pgfpathclose%
\pgfusepath{fill}%
\end{pgfscope}%
\begin{pgfscope}%
\pgfpathrectangle{\pgfqpoint{0.100000in}{0.100000in}}{\pgfqpoint{3.608454in}{2.310000in}}%
\pgfusepath{clip}%
\pgfsetbuttcap%
\pgfsetmiterjoin%
\definecolor{currentfill}{rgb}{0.000000,0.505882,0.747059}%
\pgfsetfillcolor{currentfill}%
\pgfsetlinewidth{0.000000pt}%
\definecolor{currentstroke}{rgb}{0.000000,0.000000,0.000000}%
\pgfsetstrokecolor{currentstroke}%
\pgfsetstrokeopacity{0.000000}%
\pgfsetdash{}{0pt}%
\pgfpathmoveto{\pgfqpoint{1.881222in}{1.423490in}}%
\pgfpathlineto{\pgfqpoint{1.853837in}{1.424663in}}%
\pgfpathlineto{\pgfqpoint{1.856208in}{1.476676in}}%
\pgfpathlineto{\pgfqpoint{1.841990in}{1.478161in}}%
\pgfpathlineto{\pgfqpoint{1.842917in}{1.507856in}}%
\pgfpathlineto{\pgfqpoint{1.855925in}{1.507277in}}%
\pgfpathlineto{\pgfqpoint{1.857457in}{1.534695in}}%
\pgfpathlineto{\pgfqpoint{1.884664in}{1.533508in}}%
\pgfpathlineto{\pgfqpoint{1.883282in}{1.506053in}}%
\pgfpathlineto{\pgfqpoint{1.884864in}{1.505981in}}%
\pgfpathlineto{\pgfqpoint{1.883772in}{1.478524in}}%
\pgfpathlineto{\pgfqpoint{1.881222in}{1.423490in}}%
\pgfpathclose%
\pgfusepath{fill}%
\end{pgfscope}%
\begin{pgfscope}%
\pgfpathrectangle{\pgfqpoint{0.100000in}{0.100000in}}{\pgfqpoint{3.608454in}{2.310000in}}%
\pgfusepath{clip}%
\pgfsetbuttcap%
\pgfsetmiterjoin%
\definecolor{currentfill}{rgb}{0.000000,0.423529,0.788235}%
\pgfsetfillcolor{currentfill}%
\pgfsetlinewidth{0.000000pt}%
\definecolor{currentstroke}{rgb}{0.000000,0.000000,0.000000}%
\pgfsetstrokecolor{currentstroke}%
\pgfsetstrokeopacity{0.000000}%
\pgfsetdash{}{0pt}%
\pgfpathmoveto{\pgfqpoint{0.836288in}{2.219545in}}%
\pgfpathlineto{\pgfqpoint{0.829989in}{2.221231in}}%
\pgfpathlineto{\pgfqpoint{0.822836in}{2.217028in}}%
\pgfpathlineto{\pgfqpoint{0.820986in}{2.210320in}}%
\pgfpathlineto{\pgfqpoint{0.814986in}{2.205959in}}%
\pgfpathlineto{\pgfqpoint{0.810303in}{2.197629in}}%
\pgfpathlineto{\pgfqpoint{0.807006in}{2.198546in}}%
\pgfpathlineto{\pgfqpoint{0.798560in}{2.193839in}}%
\pgfpathlineto{\pgfqpoint{0.796683in}{2.187200in}}%
\pgfpathlineto{\pgfqpoint{0.778779in}{2.191862in}}%
\pgfpathlineto{\pgfqpoint{0.774962in}{2.186885in}}%
\pgfpathlineto{\pgfqpoint{0.769524in}{2.184730in}}%
\pgfpathlineto{\pgfqpoint{0.767233in}{2.177075in}}%
\pgfpathlineto{\pgfqpoint{0.763587in}{2.181602in}}%
\pgfpathlineto{\pgfqpoint{0.748201in}{2.185881in}}%
\pgfpathlineto{\pgfqpoint{0.742925in}{2.192347in}}%
\pgfpathlineto{\pgfqpoint{0.735660in}{2.196356in}}%
\pgfpathlineto{\pgfqpoint{0.733242in}{2.201231in}}%
\pgfpathlineto{\pgfqpoint{0.724213in}{2.207197in}}%
\pgfpathlineto{\pgfqpoint{0.722390in}{2.214825in}}%
\pgfpathlineto{\pgfqpoint{0.718104in}{2.222638in}}%
\pgfpathlineto{\pgfqpoint{0.722977in}{2.230297in}}%
\pgfpathlineto{\pgfqpoint{0.721489in}{2.246797in}}%
\pgfpathlineto{\pgfqpoint{0.727318in}{2.257197in}}%
\pgfpathlineto{\pgfqpoint{0.738201in}{2.260046in}}%
\pgfpathlineto{\pgfqpoint{0.741543in}{2.265250in}}%
\pgfpathlineto{\pgfqpoint{0.736966in}{2.278404in}}%
\pgfpathlineto{\pgfqpoint{0.740255in}{2.288121in}}%
\pgfpathlineto{\pgfqpoint{0.752759in}{2.291869in}}%
\pgfpathlineto{\pgfqpoint{0.760304in}{2.288239in}}%
\pgfpathlineto{\pgfqpoint{0.760184in}{2.281456in}}%
\pgfpathlineto{\pgfqpoint{0.765081in}{2.268518in}}%
\pgfpathlineto{\pgfqpoint{0.771087in}{2.257398in}}%
\pgfpathlineto{\pgfqpoint{0.769468in}{2.254538in}}%
\pgfpathlineto{\pgfqpoint{0.780160in}{2.239854in}}%
\pgfpathlineto{\pgfqpoint{0.780931in}{2.234819in}}%
\pgfpathlineto{\pgfqpoint{0.790372in}{2.232053in}}%
\pgfpathlineto{\pgfqpoint{0.791522in}{2.239767in}}%
\pgfpathlineto{\pgfqpoint{0.800018in}{2.241064in}}%
\pgfpathlineto{\pgfqpoint{0.802328in}{2.232125in}}%
\pgfpathlineto{\pgfqpoint{0.815034in}{2.234984in}}%
\pgfpathlineto{\pgfqpoint{0.819597in}{2.231350in}}%
\pgfpathlineto{\pgfqpoint{0.832614in}{2.236217in}}%
\pgfpathlineto{\pgfqpoint{0.836249in}{2.233779in}}%
\pgfpathlineto{\pgfqpoint{0.836288in}{2.219545in}}%
\pgfpathclose%
\pgfusepath{fill}%
\end{pgfscope}%
\begin{pgfscope}%
\pgfpathrectangle{\pgfqpoint{0.100000in}{0.100000in}}{\pgfqpoint{3.608454in}{2.310000in}}%
\pgfusepath{clip}%
\pgfsetbuttcap%
\pgfsetmiterjoin%
\definecolor{currentfill}{rgb}{0.000000,0.631373,0.684314}%
\pgfsetfillcolor{currentfill}%
\pgfsetlinewidth{0.000000pt}%
\definecolor{currentstroke}{rgb}{0.000000,0.000000,0.000000}%
\pgfsetstrokecolor{currentstroke}%
\pgfsetstrokeopacity{0.000000}%
\pgfsetdash{}{0pt}%
\pgfpathmoveto{\pgfqpoint{2.485040in}{0.651337in}}%
\pgfpathlineto{\pgfqpoint{2.478549in}{0.654208in}}%
\pgfpathlineto{\pgfqpoint{2.472822in}{0.672176in}}%
\pgfpathlineto{\pgfqpoint{2.466743in}{0.676753in}}%
\pgfpathlineto{\pgfqpoint{2.460731in}{0.688378in}}%
\pgfpathlineto{\pgfqpoint{2.461305in}{0.698015in}}%
\pgfpathlineto{\pgfqpoint{2.465118in}{0.704833in}}%
\pgfpathlineto{\pgfqpoint{2.467199in}{0.715885in}}%
\pgfpathlineto{\pgfqpoint{2.472257in}{0.716197in}}%
\pgfpathlineto{\pgfqpoint{2.482623in}{0.717728in}}%
\pgfpathlineto{\pgfqpoint{2.493658in}{0.717377in}}%
\pgfpathlineto{\pgfqpoint{2.494101in}{0.710117in}}%
\pgfpathlineto{\pgfqpoint{2.528702in}{0.712404in}}%
\pgfpathlineto{\pgfqpoint{2.525381in}{0.754028in}}%
\pgfpathlineto{\pgfqpoint{2.551986in}{0.756063in}}%
\pgfpathlineto{\pgfqpoint{2.556138in}{0.721372in}}%
\pgfpathlineto{\pgfqpoint{2.561760in}{0.672494in}}%
\pgfpathlineto{\pgfqpoint{2.556857in}{0.667105in}}%
\pgfpathlineto{\pgfqpoint{2.547096in}{0.670673in}}%
\pgfpathlineto{\pgfqpoint{2.541328in}{0.668013in}}%
\pgfpathlineto{\pgfqpoint{2.525544in}{0.670845in}}%
\pgfpathlineto{\pgfqpoint{2.515743in}{0.668541in}}%
\pgfpathlineto{\pgfqpoint{2.497858in}{0.661623in}}%
\pgfpathlineto{\pgfqpoint{2.485040in}{0.651337in}}%
\pgfpathclose%
\pgfusepath{fill}%
\end{pgfscope}%
\begin{pgfscope}%
\pgfpathrectangle{\pgfqpoint{0.100000in}{0.100000in}}{\pgfqpoint{3.608454in}{2.310000in}}%
\pgfusepath{clip}%
\pgfsetbuttcap%
\pgfsetmiterjoin%
\definecolor{currentfill}{rgb}{0.000000,0.690196,0.654902}%
\pgfsetfillcolor{currentfill}%
\pgfsetlinewidth{0.000000pt}%
\definecolor{currentstroke}{rgb}{0.000000,0.000000,0.000000}%
\pgfsetstrokecolor{currentstroke}%
\pgfsetstrokeopacity{0.000000}%
\pgfsetdash{}{0pt}%
\pgfpathmoveto{\pgfqpoint{3.568956in}{1.730471in}}%
\pgfpathlineto{\pgfqpoint{3.576435in}{1.723494in}}%
\pgfpathlineto{\pgfqpoint{3.574246in}{1.719609in}}%
\pgfpathlineto{\pgfqpoint{3.565620in}{1.718239in}}%
\pgfpathlineto{\pgfqpoint{3.570850in}{1.726339in}}%
\pgfpathlineto{\pgfqpoint{3.568956in}{1.730471in}}%
\pgfpathclose%
\pgfusepath{fill}%
\end{pgfscope}%
\begin{pgfscope}%
\pgfpathrectangle{\pgfqpoint{0.100000in}{0.100000in}}{\pgfqpoint{3.608454in}{2.310000in}}%
\pgfusepath{clip}%
\pgfsetbuttcap%
\pgfsetmiterjoin%
\definecolor{currentfill}{rgb}{0.000000,0.525490,0.737255}%
\pgfsetfillcolor{currentfill}%
\pgfsetlinewidth{0.000000pt}%
\definecolor{currentstroke}{rgb}{0.000000,0.000000,0.000000}%
\pgfsetstrokecolor{currentstroke}%
\pgfsetstrokeopacity{0.000000}%
\pgfsetdash{}{0pt}%
\pgfpathmoveto{\pgfqpoint{3.024683in}{1.232048in}}%
\pgfpathlineto{\pgfqpoint{3.011272in}{1.247768in}}%
\pgfpathlineto{\pgfqpoint{3.005275in}{1.242589in}}%
\pgfpathlineto{\pgfqpoint{2.993913in}{1.252683in}}%
\pgfpathlineto{\pgfqpoint{2.990512in}{1.259171in}}%
\pgfpathlineto{\pgfqpoint{2.994849in}{1.262443in}}%
\pgfpathlineto{\pgfqpoint{2.983343in}{1.270885in}}%
\pgfpathlineto{\pgfqpoint{2.992726in}{1.278073in}}%
\pgfpathlineto{\pgfqpoint{2.991086in}{1.284432in}}%
\pgfpathlineto{\pgfqpoint{3.000083in}{1.309688in}}%
\pgfpathlineto{\pgfqpoint{3.013170in}{1.308307in}}%
\pgfpathlineto{\pgfqpoint{3.022678in}{1.310484in}}%
\pgfpathlineto{\pgfqpoint{3.033267in}{1.304579in}}%
\pgfpathlineto{\pgfqpoint{3.045914in}{1.315133in}}%
\pgfpathlineto{\pgfqpoint{3.050832in}{1.323908in}}%
\pgfpathlineto{\pgfqpoint{3.060124in}{1.322440in}}%
\pgfpathlineto{\pgfqpoint{3.060878in}{1.314437in}}%
\pgfpathlineto{\pgfqpoint{3.074021in}{1.306357in}}%
\pgfpathlineto{\pgfqpoint{3.067687in}{1.299266in}}%
\pgfpathlineto{\pgfqpoint{3.064782in}{1.300929in}}%
\pgfpathlineto{\pgfqpoint{3.057612in}{1.293606in}}%
\pgfpathlineto{\pgfqpoint{3.055501in}{1.285981in}}%
\pgfpathlineto{\pgfqpoint{3.057510in}{1.278132in}}%
\pgfpathlineto{\pgfqpoint{3.069985in}{1.273083in}}%
\pgfpathlineto{\pgfqpoint{3.074664in}{1.267099in}}%
\pgfpathlineto{\pgfqpoint{3.074241in}{1.252094in}}%
\pgfpathlineto{\pgfqpoint{3.067580in}{1.245522in}}%
\pgfpathlineto{\pgfqpoint{3.036908in}{1.247132in}}%
\pgfpathlineto{\pgfqpoint{3.032148in}{1.244073in}}%
\pgfpathlineto{\pgfqpoint{3.024683in}{1.232048in}}%
\pgfpathclose%
\pgfusepath{fill}%
\end{pgfscope}%
\begin{pgfscope}%
\pgfpathrectangle{\pgfqpoint{0.100000in}{0.100000in}}{\pgfqpoint{3.608454in}{2.310000in}}%
\pgfusepath{clip}%
\pgfsetbuttcap%
\pgfsetmiterjoin%
\definecolor{currentfill}{rgb}{0.000000,0.560784,0.719608}%
\pgfsetfillcolor{currentfill}%
\pgfsetlinewidth{0.000000pt}%
\definecolor{currentstroke}{rgb}{0.000000,0.000000,0.000000}%
\pgfsetstrokecolor{currentstroke}%
\pgfsetstrokeopacity{0.000000}%
\pgfsetdash{}{0pt}%
\pgfpathmoveto{\pgfqpoint{1.267460in}{1.535366in}}%
\pgfpathlineto{\pgfqpoint{1.261144in}{1.538773in}}%
\pgfpathlineto{\pgfqpoint{1.262249in}{1.546769in}}%
\pgfpathlineto{\pgfqpoint{1.260798in}{1.552822in}}%
\pgfpathlineto{\pgfqpoint{1.249496in}{1.554569in}}%
\pgfpathlineto{\pgfqpoint{1.248664in}{1.549321in}}%
\pgfpathlineto{\pgfqpoint{1.242723in}{1.546157in}}%
\pgfpathlineto{\pgfqpoint{1.241154in}{1.553467in}}%
\pgfpathlineto{\pgfqpoint{1.220332in}{1.551758in}}%
\pgfpathlineto{\pgfqpoint{1.213136in}{1.556950in}}%
\pgfpathlineto{\pgfqpoint{1.215522in}{1.571272in}}%
\pgfpathlineto{\pgfqpoint{1.186062in}{1.576054in}}%
\pgfpathlineto{\pgfqpoint{1.154007in}{1.581932in}}%
\pgfpathlineto{\pgfqpoint{1.157543in}{1.601688in}}%
\pgfpathlineto{\pgfqpoint{1.146992in}{1.596437in}}%
\pgfpathlineto{\pgfqpoint{1.143956in}{1.603549in}}%
\pgfpathlineto{\pgfqpoint{1.144102in}{1.610430in}}%
\pgfpathlineto{\pgfqpoint{1.136553in}{1.614347in}}%
\pgfpathlineto{\pgfqpoint{1.132825in}{1.619941in}}%
\pgfpathlineto{\pgfqpoint{1.138324in}{1.627635in}}%
\pgfpathlineto{\pgfqpoint{1.142480in}{1.636751in}}%
\pgfpathlineto{\pgfqpoint{1.141060in}{1.650733in}}%
\pgfpathlineto{\pgfqpoint{1.139137in}{1.653332in}}%
\pgfpathlineto{\pgfqpoint{1.142493in}{1.660281in}}%
\pgfpathlineto{\pgfqpoint{1.141273in}{1.664762in}}%
\pgfpathlineto{\pgfqpoint{1.167956in}{1.660095in}}%
\pgfpathlineto{\pgfqpoint{1.182268in}{1.739464in}}%
\pgfpathlineto{\pgfqpoint{1.186356in}{1.762535in}}%
\pgfpathlineto{\pgfqpoint{1.199226in}{1.758954in}}%
\pgfpathlineto{\pgfqpoint{1.198346in}{1.753929in}}%
\pgfpathlineto{\pgfqpoint{1.211732in}{1.751597in}}%
\pgfpathlineto{\pgfqpoint{1.212511in}{1.756000in}}%
\pgfpathlineto{\pgfqpoint{1.225832in}{1.753673in}}%
\pgfpathlineto{\pgfqpoint{1.226972in}{1.760454in}}%
\pgfpathlineto{\pgfqpoint{1.243269in}{1.757690in}}%
\pgfpathlineto{\pgfqpoint{1.244733in}{1.764439in}}%
\pgfpathlineto{\pgfqpoint{1.261563in}{1.761526in}}%
\pgfpathlineto{\pgfqpoint{1.263755in}{1.753309in}}%
\pgfpathlineto{\pgfqpoint{1.262997in}{1.742086in}}%
\pgfpathlineto{\pgfqpoint{1.264849in}{1.733953in}}%
\pgfpathlineto{\pgfqpoint{1.270191in}{1.723778in}}%
\pgfpathlineto{\pgfqpoint{1.275945in}{1.717107in}}%
\pgfpathlineto{\pgfqpoint{1.282465in}{1.703069in}}%
\pgfpathlineto{\pgfqpoint{1.290897in}{1.695072in}}%
\pgfpathlineto{\pgfqpoint{1.287915in}{1.674697in}}%
\pgfpathlineto{\pgfqpoint{1.287325in}{1.661183in}}%
\pgfpathlineto{\pgfqpoint{1.347746in}{1.651996in}}%
\pgfpathlineto{\pgfqpoint{1.375467in}{1.648053in}}%
\pgfpathlineto{\pgfqpoint{1.374676in}{1.634274in}}%
\pgfpathlineto{\pgfqpoint{1.369658in}{1.600576in}}%
\pgfpathlineto{\pgfqpoint{1.345097in}{1.604208in}}%
\pgfpathlineto{\pgfqpoint{1.342910in}{1.583654in}}%
\pgfpathlineto{\pgfqpoint{1.338390in}{1.552623in}}%
\pgfpathlineto{\pgfqpoint{1.302483in}{1.557690in}}%
\pgfpathlineto{\pgfqpoint{1.271578in}{1.562500in}}%
\pgfpathlineto{\pgfqpoint{1.267460in}{1.535366in}}%
\pgfpathclose%
\pgfusepath{fill}%
\end{pgfscope}%
\begin{pgfscope}%
\pgfpathrectangle{\pgfqpoint{0.100000in}{0.100000in}}{\pgfqpoint{3.608454in}{2.310000in}}%
\pgfusepath{clip}%
\pgfsetbuttcap%
\pgfsetmiterjoin%
\definecolor{currentfill}{rgb}{0.000000,0.439216,0.780392}%
\pgfsetfillcolor{currentfill}%
\pgfsetlinewidth{0.000000pt}%
\definecolor{currentstroke}{rgb}{0.000000,0.000000,0.000000}%
\pgfsetstrokecolor{currentstroke}%
\pgfsetstrokeopacity{0.000000}%
\pgfsetdash{}{0pt}%
\pgfpathmoveto{\pgfqpoint{2.520499in}{1.715478in}}%
\pgfpathlineto{\pgfqpoint{2.499443in}{1.713874in}}%
\pgfpathlineto{\pgfqpoint{2.497309in}{1.741424in}}%
\pgfpathlineto{\pgfqpoint{2.504166in}{1.741961in}}%
\pgfpathlineto{\pgfqpoint{2.501961in}{1.769586in}}%
\pgfpathlineto{\pgfqpoint{2.510762in}{1.770224in}}%
\pgfpathlineto{\pgfqpoint{2.510243in}{1.777129in}}%
\pgfpathlineto{\pgfqpoint{2.529824in}{1.778725in}}%
\pgfpathlineto{\pgfqpoint{2.532257in}{1.767832in}}%
\pgfpathlineto{\pgfqpoint{2.525321in}{1.760582in}}%
\pgfpathlineto{\pgfqpoint{2.521671in}{1.742517in}}%
\pgfpathlineto{\pgfqpoint{2.524790in}{1.727344in}}%
\pgfpathlineto{\pgfqpoint{2.520499in}{1.715478in}}%
\pgfpathclose%
\pgfusepath{fill}%
\end{pgfscope}%
\begin{pgfscope}%
\pgfpathrectangle{\pgfqpoint{0.100000in}{0.100000in}}{\pgfqpoint{3.608454in}{2.310000in}}%
\pgfusepath{clip}%
\pgfsetbuttcap%
\pgfsetmiterjoin%
\definecolor{currentfill}{rgb}{0.000000,0.600000,0.700000}%
\pgfsetfillcolor{currentfill}%
\pgfsetlinewidth{0.000000pt}%
\definecolor{currentstroke}{rgb}{0.000000,0.000000,0.000000}%
\pgfsetstrokecolor{currentstroke}%
\pgfsetstrokeopacity{0.000000}%
\pgfsetdash{}{0pt}%
\pgfpathmoveto{\pgfqpoint{1.576920in}{0.953485in}}%
\pgfpathlineto{\pgfqpoint{1.580563in}{0.992038in}}%
\pgfpathlineto{\pgfqpoint{1.614395in}{0.989155in}}%
\pgfpathlineto{\pgfqpoint{1.617270in}{1.023437in}}%
\pgfpathlineto{\pgfqpoint{1.651469in}{1.020799in}}%
\pgfpathlineto{\pgfqpoint{1.648845in}{0.986397in}}%
\pgfpathlineto{\pgfqpoint{1.642812in}{0.986873in}}%
\pgfpathlineto{\pgfqpoint{1.640100in}{0.948213in}}%
\pgfpathlineto{\pgfqpoint{1.576920in}{0.953485in}}%
\pgfpathclose%
\pgfusepath{fill}%
\end{pgfscope}%
\begin{pgfscope}%
\pgfpathrectangle{\pgfqpoint{0.100000in}{0.100000in}}{\pgfqpoint{3.608454in}{2.310000in}}%
\pgfusepath{clip}%
\pgfsetbuttcap%
\pgfsetmiterjoin%
\definecolor{currentfill}{rgb}{0.000000,0.792157,0.603922}%
\pgfsetfillcolor{currentfill}%
\pgfsetlinewidth{0.000000pt}%
\definecolor{currentstroke}{rgb}{0.000000,0.000000,0.000000}%
\pgfsetstrokecolor{currentstroke}%
\pgfsetstrokeopacity{0.000000}%
\pgfsetdash{}{0pt}%
\pgfpathmoveto{\pgfqpoint{0.557819in}{2.120432in}}%
\pgfpathlineto{\pgfqpoint{0.538739in}{2.126084in}}%
\pgfpathlineto{\pgfqpoint{0.525797in}{2.131179in}}%
\pgfpathlineto{\pgfqpoint{0.534887in}{2.152356in}}%
\pgfpathlineto{\pgfqpoint{0.534769in}{2.164166in}}%
\pgfpathlineto{\pgfqpoint{0.540718in}{2.157678in}}%
\pgfpathlineto{\pgfqpoint{0.547629in}{2.160026in}}%
\pgfpathlineto{\pgfqpoint{0.555579in}{2.159647in}}%
\pgfpathlineto{\pgfqpoint{0.552341in}{2.166159in}}%
\pgfpathlineto{\pgfqpoint{0.542113in}{2.164053in}}%
\pgfpathlineto{\pgfqpoint{0.534487in}{2.174240in}}%
\pgfpathlineto{\pgfqpoint{0.538478in}{2.184732in}}%
\pgfpathlineto{\pgfqpoint{0.548088in}{2.185551in}}%
\pgfpathlineto{\pgfqpoint{0.547127in}{2.201192in}}%
\pgfpathlineto{\pgfqpoint{0.543078in}{2.205524in}}%
\pgfpathlineto{\pgfqpoint{0.545221in}{2.218928in}}%
\pgfpathlineto{\pgfqpoint{0.549339in}{2.216444in}}%
\pgfpathlineto{\pgfqpoint{0.559409in}{2.219592in}}%
\pgfpathlineto{\pgfqpoint{0.547351in}{2.228134in}}%
\pgfpathlineto{\pgfqpoint{0.549613in}{2.249680in}}%
\pgfpathlineto{\pgfqpoint{0.547461in}{2.260324in}}%
\pgfpathlineto{\pgfqpoint{0.549782in}{2.270356in}}%
\pgfpathlineto{\pgfqpoint{0.565186in}{2.263908in}}%
\pgfpathlineto{\pgfqpoint{0.593017in}{2.254746in}}%
\pgfpathlineto{\pgfqpoint{0.586531in}{2.235034in}}%
\pgfpathlineto{\pgfqpoint{0.583086in}{2.221440in}}%
\pgfpathlineto{\pgfqpoint{0.597993in}{2.216816in}}%
\pgfpathlineto{\pgfqpoint{0.598094in}{2.209351in}}%
\pgfpathlineto{\pgfqpoint{0.593046in}{2.193971in}}%
\pgfpathlineto{\pgfqpoint{0.582091in}{2.197407in}}%
\pgfpathlineto{\pgfqpoint{0.572748in}{2.166216in}}%
\pgfpathlineto{\pgfqpoint{0.631314in}{2.147810in}}%
\pgfpathlineto{\pgfqpoint{0.623310in}{2.122649in}}%
\pgfpathlineto{\pgfqpoint{0.615173in}{2.117075in}}%
\pgfpathlineto{\pgfqpoint{0.609486in}{2.121452in}}%
\pgfpathlineto{\pgfqpoint{0.593066in}{2.119148in}}%
\pgfpathlineto{\pgfqpoint{0.589293in}{2.135786in}}%
\pgfpathlineto{\pgfqpoint{0.578043in}{2.148170in}}%
\pgfpathlineto{\pgfqpoint{0.566644in}{2.148266in}}%
\pgfpathlineto{\pgfqpoint{0.557819in}{2.120432in}}%
\pgfpathclose%
\pgfusepath{fill}%
\end{pgfscope}%
\begin{pgfscope}%
\pgfpathrectangle{\pgfqpoint{0.100000in}{0.100000in}}{\pgfqpoint{3.608454in}{2.310000in}}%
\pgfusepath{clip}%
\pgfsetbuttcap%
\pgfsetmiterjoin%
\definecolor{currentfill}{rgb}{0.000000,0.839216,0.580392}%
\pgfsetfillcolor{currentfill}%
\pgfsetlinewidth{0.000000pt}%
\definecolor{currentstroke}{rgb}{0.000000,0.000000,0.000000}%
\pgfsetstrokecolor{currentstroke}%
\pgfsetstrokeopacity{0.000000}%
\pgfsetdash{}{0pt}%
\pgfpathmoveto{\pgfqpoint{0.954967in}{2.104591in}}%
\pgfpathlineto{\pgfqpoint{0.949294in}{2.080795in}}%
\pgfpathlineto{\pgfqpoint{0.933769in}{2.084665in}}%
\pgfpathlineto{\pgfqpoint{0.926998in}{2.079195in}}%
\pgfpathlineto{\pgfqpoint{0.909838in}{2.083326in}}%
\pgfpathlineto{\pgfqpoint{0.907603in}{2.074322in}}%
\pgfpathlineto{\pgfqpoint{0.899012in}{2.076012in}}%
\pgfpathlineto{\pgfqpoint{0.898511in}{2.086388in}}%
\pgfpathlineto{\pgfqpoint{0.892918in}{2.093767in}}%
\pgfpathlineto{\pgfqpoint{0.890708in}{2.100407in}}%
\pgfpathlineto{\pgfqpoint{0.882536in}{2.100932in}}%
\pgfpathlineto{\pgfqpoint{0.878570in}{2.104452in}}%
\pgfpathlineto{\pgfqpoint{0.868296in}{2.100626in}}%
\pgfpathlineto{\pgfqpoint{0.849927in}{2.099792in}}%
\pgfpathlineto{\pgfqpoint{0.849420in}{2.114482in}}%
\pgfpathlineto{\pgfqpoint{0.851613in}{2.114354in}}%
\pgfpathlineto{\pgfqpoint{0.863100in}{2.115696in}}%
\pgfpathlineto{\pgfqpoint{0.868405in}{2.124498in}}%
\pgfpathlineto{\pgfqpoint{0.875234in}{2.151233in}}%
\pgfpathlineto{\pgfqpoint{0.923573in}{2.138917in}}%
\pgfpathlineto{\pgfqpoint{0.921056in}{2.128678in}}%
\pgfpathlineto{\pgfqpoint{0.933521in}{2.118231in}}%
\pgfpathlineto{\pgfqpoint{0.950036in}{2.114158in}}%
\pgfpathlineto{\pgfqpoint{0.956600in}{2.111418in}}%
\pgfpathlineto{\pgfqpoint{0.954967in}{2.104591in}}%
\pgfpathclose%
\pgfusepath{fill}%
\end{pgfscope}%
\begin{pgfscope}%
\pgfpathrectangle{\pgfqpoint{0.100000in}{0.100000in}}{\pgfqpoint{3.608454in}{2.310000in}}%
\pgfusepath{clip}%
\pgfsetbuttcap%
\pgfsetmiterjoin%
\definecolor{currentfill}{rgb}{0.000000,0.698039,0.650980}%
\pgfsetfillcolor{currentfill}%
\pgfsetlinewidth{0.000000pt}%
\definecolor{currentstroke}{rgb}{0.000000,0.000000,0.000000}%
\pgfsetstrokecolor{currentstroke}%
\pgfsetstrokeopacity{0.000000}%
\pgfsetdash{}{0pt}%
\pgfpathmoveto{\pgfqpoint{0.877772in}{0.912470in}}%
\pgfpathlineto{\pgfqpoint{0.885192in}{0.948795in}}%
\pgfpathlineto{\pgfqpoint{0.934927in}{0.938858in}}%
\pgfpathlineto{\pgfqpoint{1.001421in}{0.926345in}}%
\pgfpathlineto{\pgfqpoint{1.075355in}{0.913752in}}%
\pgfpathlineto{\pgfqpoint{1.074172in}{0.906976in}}%
\pgfpathlineto{\pgfqpoint{1.092774in}{0.903733in}}%
\pgfpathlineto{\pgfqpoint{1.167197in}{0.891666in}}%
\pgfpathlineto{\pgfqpoint{1.153619in}{0.805601in}}%
\pgfpathlineto{\pgfqpoint{1.089707in}{0.815956in}}%
\pgfpathlineto{\pgfqpoint{1.017594in}{0.828333in}}%
\pgfpathlineto{\pgfqpoint{0.999378in}{0.839223in}}%
\pgfpathlineto{\pgfqpoint{0.900066in}{0.898746in}}%
\pgfpathlineto{\pgfqpoint{0.877772in}{0.912470in}}%
\pgfpathclose%
\pgfusepath{fill}%
\end{pgfscope}%
\begin{pgfscope}%
\pgfpathrectangle{\pgfqpoint{0.100000in}{0.100000in}}{\pgfqpoint{3.608454in}{2.310000in}}%
\pgfusepath{clip}%
\pgfsetbuttcap%
\pgfsetmiterjoin%
\definecolor{currentfill}{rgb}{0.000000,0.439216,0.780392}%
\pgfsetfillcolor{currentfill}%
\pgfsetlinewidth{0.000000pt}%
\definecolor{currentstroke}{rgb}{0.000000,0.000000,0.000000}%
\pgfsetstrokecolor{currentstroke}%
\pgfsetstrokeopacity{0.000000}%
\pgfsetdash{}{0pt}%
\pgfpathmoveto{\pgfqpoint{3.153317in}{1.449269in}}%
\pgfpathlineto{\pgfqpoint{3.152017in}{1.444895in}}%
\pgfpathlineto{\pgfqpoint{3.144962in}{1.435381in}}%
\pgfpathlineto{\pgfqpoint{3.141078in}{1.437828in}}%
\pgfpathlineto{\pgfqpoint{3.125843in}{1.434080in}}%
\pgfpathlineto{\pgfqpoint{3.119217in}{1.439782in}}%
\pgfpathlineto{\pgfqpoint{3.117131in}{1.434042in}}%
\pgfpathlineto{\pgfqpoint{3.111250in}{1.435310in}}%
\pgfpathlineto{\pgfqpoint{3.105576in}{1.439244in}}%
\pgfpathlineto{\pgfqpoint{3.087837in}{1.442494in}}%
\pgfpathlineto{\pgfqpoint{3.082421in}{1.444639in}}%
\pgfpathlineto{\pgfqpoint{3.088891in}{1.453283in}}%
\pgfpathlineto{\pgfqpoint{3.091266in}{1.461236in}}%
\pgfpathlineto{\pgfqpoint{3.097628in}{1.471273in}}%
\pgfpathlineto{\pgfqpoint{3.109226in}{1.471916in}}%
\pgfpathlineto{\pgfqpoint{3.117245in}{1.469315in}}%
\pgfpathlineto{\pgfqpoint{3.153317in}{1.449269in}}%
\pgfpathclose%
\pgfusepath{fill}%
\end{pgfscope}%
\begin{pgfscope}%
\pgfpathrectangle{\pgfqpoint{0.100000in}{0.100000in}}{\pgfqpoint{3.608454in}{2.310000in}}%
\pgfusepath{clip}%
\pgfsetbuttcap%
\pgfsetmiterjoin%
\definecolor{currentfill}{rgb}{0.000000,0.501961,0.749020}%
\pgfsetfillcolor{currentfill}%
\pgfsetlinewidth{0.000000pt}%
\definecolor{currentstroke}{rgb}{0.000000,0.000000,0.000000}%
\pgfsetstrokecolor{currentstroke}%
\pgfsetstrokeopacity{0.000000}%
\pgfsetdash{}{0pt}%
\pgfpathmoveto{\pgfqpoint{1.965169in}{0.705757in}}%
\pgfpathlineto{\pgfqpoint{1.952059in}{0.698355in}}%
\pgfpathlineto{\pgfqpoint{1.947915in}{0.687499in}}%
\pgfpathlineto{\pgfqpoint{1.926963in}{0.697491in}}%
\pgfpathlineto{\pgfqpoint{1.913119in}{0.701138in}}%
\pgfpathlineto{\pgfqpoint{1.907744in}{0.711220in}}%
\pgfpathlineto{\pgfqpoint{1.871649in}{0.712052in}}%
\pgfpathlineto{\pgfqpoint{1.866023in}{0.717001in}}%
\pgfpathlineto{\pgfqpoint{1.863869in}{0.728404in}}%
\pgfpathlineto{\pgfqpoint{1.884281in}{0.742465in}}%
\pgfpathlineto{\pgfqpoint{1.919219in}{0.761725in}}%
\pgfpathlineto{\pgfqpoint{1.924502in}{0.764640in}}%
\pgfpathlineto{\pgfqpoint{1.942049in}{0.732981in}}%
\pgfpathlineto{\pgfqpoint{1.947012in}{0.726789in}}%
\pgfpathlineto{\pgfqpoint{1.951524in}{0.729506in}}%
\pgfpathlineto{\pgfqpoint{1.965169in}{0.705757in}}%
\pgfpathclose%
\pgfusepath{fill}%
\end{pgfscope}%
\begin{pgfscope}%
\pgfpathrectangle{\pgfqpoint{0.100000in}{0.100000in}}{\pgfqpoint{3.608454in}{2.310000in}}%
\pgfusepath{clip}%
\pgfsetbuttcap%
\pgfsetmiterjoin%
\definecolor{currentfill}{rgb}{0.000000,0.800000,0.600000}%
\pgfsetfillcolor{currentfill}%
\pgfsetlinewidth{0.000000pt}%
\definecolor{currentstroke}{rgb}{0.000000,0.000000,0.000000}%
\pgfsetstrokecolor{currentstroke}%
\pgfsetstrokeopacity{0.000000}%
\pgfsetdash{}{0pt}%
\pgfpathmoveto{\pgfqpoint{1.805378in}{0.356576in}}%
\pgfpathlineto{\pgfqpoint{1.817184in}{0.364409in}}%
\pgfpathlineto{\pgfqpoint{1.821677in}{0.373163in}}%
\pgfpathlineto{\pgfqpoint{1.859982in}{0.371433in}}%
\pgfpathlineto{\pgfqpoint{1.891484in}{0.370110in}}%
\pgfpathlineto{\pgfqpoint{1.891060in}{0.356772in}}%
\pgfpathlineto{\pgfqpoint{1.899944in}{0.355045in}}%
\pgfpathlineto{\pgfqpoint{1.930477in}{0.354309in}}%
\pgfpathlineto{\pgfqpoint{1.932540in}{0.346180in}}%
\pgfpathlineto{\pgfqpoint{1.930157in}{0.343218in}}%
\pgfpathlineto{\pgfqpoint{1.936961in}{0.334615in}}%
\pgfpathlineto{\pgfqpoint{1.936672in}{0.327595in}}%
\pgfpathlineto{\pgfqpoint{1.940091in}{0.321228in}}%
\pgfpathlineto{\pgfqpoint{1.940482in}{0.313585in}}%
\pgfpathlineto{\pgfqpoint{1.950108in}{0.309249in}}%
\pgfpathlineto{\pgfqpoint{1.950498in}{0.302891in}}%
\pgfpathlineto{\pgfqpoint{1.934207in}{0.298800in}}%
\pgfpathlineto{\pgfqpoint{1.929813in}{0.292862in}}%
\pgfpathlineto{\pgfqpoint{1.923685in}{0.296794in}}%
\pgfpathlineto{\pgfqpoint{1.911823in}{0.308379in}}%
\pgfpathlineto{\pgfqpoint{1.896783in}{0.312095in}}%
\pgfpathlineto{\pgfqpoint{1.883885in}{0.310848in}}%
\pgfpathlineto{\pgfqpoint{1.870440in}{0.313575in}}%
\pgfpathlineto{\pgfqpoint{1.855483in}{0.326486in}}%
\pgfpathlineto{\pgfqpoint{1.840888in}{0.327815in}}%
\pgfpathlineto{\pgfqpoint{1.830852in}{0.339080in}}%
\pgfpathlineto{\pgfqpoint{1.828260in}{0.338364in}}%
\pgfpathlineto{\pgfqpoint{1.809077in}{0.344635in}}%
\pgfpathlineto{\pgfqpoint{1.810652in}{0.348669in}}%
\pgfpathlineto{\pgfqpoint{1.805378in}{0.356576in}}%
\pgfpathclose%
\pgfusepath{fill}%
\end{pgfscope}%
\begin{pgfscope}%
\pgfpathrectangle{\pgfqpoint{0.100000in}{0.100000in}}{\pgfqpoint{3.608454in}{2.310000in}}%
\pgfusepath{clip}%
\pgfsetbuttcap%
\pgfsetmiterjoin%
\definecolor{currentfill}{rgb}{0.000000,0.525490,0.737255}%
\pgfsetfillcolor{currentfill}%
\pgfsetlinewidth{0.000000pt}%
\definecolor{currentstroke}{rgb}{0.000000,0.000000,0.000000}%
\pgfsetstrokecolor{currentstroke}%
\pgfsetstrokeopacity{0.000000}%
\pgfsetdash{}{0pt}%
\pgfpathmoveto{\pgfqpoint{2.495702in}{1.322010in}}%
\pgfpathlineto{\pgfqpoint{2.494687in}{1.339376in}}%
\pgfpathlineto{\pgfqpoint{2.467492in}{1.337391in}}%
\pgfpathlineto{\pgfqpoint{2.467926in}{1.330458in}}%
\pgfpathlineto{\pgfqpoint{2.460771in}{1.330448in}}%
\pgfpathlineto{\pgfqpoint{2.458623in}{1.367966in}}%
\pgfpathlineto{\pgfqpoint{2.465371in}{1.368430in}}%
\pgfpathlineto{\pgfqpoint{2.506265in}{1.370984in}}%
\pgfpathlineto{\pgfqpoint{2.506514in}{1.367545in}}%
\pgfpathlineto{\pgfqpoint{2.550764in}{1.371543in}}%
\pgfpathlineto{\pgfqpoint{2.558058in}{1.370400in}}%
\pgfpathlineto{\pgfqpoint{2.562133in}{1.362791in}}%
\pgfpathlineto{\pgfqpoint{2.562407in}{1.357174in}}%
\pgfpathlineto{\pgfqpoint{2.566176in}{1.350839in}}%
\pgfpathlineto{\pgfqpoint{2.565996in}{1.346868in}}%
\pgfpathlineto{\pgfqpoint{2.542895in}{1.344839in}}%
\pgfpathlineto{\pgfqpoint{2.521388in}{1.342978in}}%
\pgfpathlineto{\pgfqpoint{2.521612in}{1.327886in}}%
\pgfpathlineto{\pgfqpoint{2.523132in}{1.323450in}}%
\pgfpathlineto{\pgfqpoint{2.495702in}{1.322010in}}%
\pgfpathclose%
\pgfusepath{fill}%
\end{pgfscope}%
\begin{pgfscope}%
\pgfpathrectangle{\pgfqpoint{0.100000in}{0.100000in}}{\pgfqpoint{3.608454in}{2.310000in}}%
\pgfusepath{clip}%
\pgfsetbuttcap%
\pgfsetmiterjoin%
\definecolor{currentfill}{rgb}{0.000000,0.568627,0.715686}%
\pgfsetfillcolor{currentfill}%
\pgfsetlinewidth{0.000000pt}%
\definecolor{currentstroke}{rgb}{0.000000,0.000000,0.000000}%
\pgfsetstrokecolor{currentstroke}%
\pgfsetstrokeopacity{0.000000}%
\pgfsetdash{}{0pt}%
\pgfpathmoveto{\pgfqpoint{3.037191in}{1.156638in}}%
\pgfpathlineto{\pgfqpoint{3.024951in}{1.162919in}}%
\pgfpathlineto{\pgfqpoint{3.020036in}{1.162708in}}%
\pgfpathlineto{\pgfqpoint{3.018531in}{1.178366in}}%
\pgfpathlineto{\pgfqpoint{3.006498in}{1.176840in}}%
\pgfpathlineto{\pgfqpoint{3.004020in}{1.198138in}}%
\pgfpathlineto{\pgfqpoint{2.996900in}{1.203053in}}%
\pgfpathlineto{\pgfqpoint{2.995823in}{1.206373in}}%
\pgfpathlineto{\pgfqpoint{2.999075in}{1.216208in}}%
\pgfpathlineto{\pgfqpoint{3.017299in}{1.218774in}}%
\pgfpathlineto{\pgfqpoint{3.037358in}{1.220901in}}%
\pgfpathlineto{\pgfqpoint{3.054182in}{1.223554in}}%
\pgfpathlineto{\pgfqpoint{3.060084in}{1.181854in}}%
\pgfpathlineto{\pgfqpoint{3.049089in}{1.181346in}}%
\pgfpathlineto{\pgfqpoint{3.039491in}{1.174669in}}%
\pgfpathlineto{\pgfqpoint{3.039762in}{1.164964in}}%
\pgfpathlineto{\pgfqpoint{3.035098in}{1.163147in}}%
\pgfpathlineto{\pgfqpoint{3.037191in}{1.156638in}}%
\pgfpathclose%
\pgfusepath{fill}%
\end{pgfscope}%
\begin{pgfscope}%
\pgfpathrectangle{\pgfqpoint{0.100000in}{0.100000in}}{\pgfqpoint{3.608454in}{2.310000in}}%
\pgfusepath{clip}%
\pgfsetbuttcap%
\pgfsetmiterjoin%
\definecolor{currentfill}{rgb}{0.000000,0.607843,0.696078}%
\pgfsetfillcolor{currentfill}%
\pgfsetlinewidth{0.000000pt}%
\definecolor{currentstroke}{rgb}{0.000000,0.000000,0.000000}%
\pgfsetstrokecolor{currentstroke}%
\pgfsetstrokeopacity{0.000000}%
\pgfsetdash{}{0pt}%
\pgfpathmoveto{\pgfqpoint{3.358684in}{1.520578in}}%
\pgfpathlineto{\pgfqpoint{3.353587in}{1.521095in}}%
\pgfpathlineto{\pgfqpoint{3.348171in}{1.542140in}}%
\pgfpathlineto{\pgfqpoint{3.330022in}{1.567556in}}%
\pgfpathlineto{\pgfqpoint{3.327731in}{1.570641in}}%
\pgfpathlineto{\pgfqpoint{3.339300in}{1.581420in}}%
\pgfpathlineto{\pgfqpoint{3.345889in}{1.596696in}}%
\pgfpathlineto{\pgfqpoint{3.350906in}{1.598176in}}%
\pgfpathlineto{\pgfqpoint{3.359788in}{1.596585in}}%
\pgfpathlineto{\pgfqpoint{3.363035in}{1.588297in}}%
\pgfpathlineto{\pgfqpoint{3.363118in}{1.574131in}}%
\pgfpathlineto{\pgfqpoint{3.365563in}{1.546396in}}%
\pgfpathlineto{\pgfqpoint{3.364102in}{1.538156in}}%
\pgfpathlineto{\pgfqpoint{3.358684in}{1.520578in}}%
\pgfpathclose%
\pgfusepath{fill}%
\end{pgfscope}%
\begin{pgfscope}%
\pgfpathrectangle{\pgfqpoint{0.100000in}{0.100000in}}{\pgfqpoint{3.608454in}{2.310000in}}%
\pgfusepath{clip}%
\pgfsetbuttcap%
\pgfsetmiterjoin%
\definecolor{currentfill}{rgb}{0.000000,0.419608,0.790196}%
\pgfsetfillcolor{currentfill}%
\pgfsetlinewidth{0.000000pt}%
\definecolor{currentstroke}{rgb}{0.000000,0.000000,0.000000}%
\pgfsetstrokecolor{currentstroke}%
\pgfsetstrokeopacity{0.000000}%
\pgfsetdash{}{0pt}%
\pgfpathmoveto{\pgfqpoint{1.939662in}{1.829664in}}%
\pgfpathlineto{\pgfqpoint{1.938531in}{1.788233in}}%
\pgfpathlineto{\pgfqpoint{1.897415in}{1.789939in}}%
\pgfpathlineto{\pgfqpoint{1.898285in}{1.810809in}}%
\pgfpathlineto{\pgfqpoint{1.849510in}{1.813058in}}%
\pgfpathlineto{\pgfqpoint{1.850854in}{1.840684in}}%
\pgfpathlineto{\pgfqpoint{1.843275in}{1.841097in}}%
\pgfpathlineto{\pgfqpoint{1.845659in}{1.896366in}}%
\pgfpathlineto{\pgfqpoint{1.885060in}{1.894397in}}%
\pgfpathlineto{\pgfqpoint{1.941869in}{1.891889in}}%
\pgfpathlineto{\pgfqpoint{1.939662in}{1.829664in}}%
\pgfpathclose%
\pgfusepath{fill}%
\end{pgfscope}%
\begin{pgfscope}%
\pgfpathrectangle{\pgfqpoint{0.100000in}{0.100000in}}{\pgfqpoint{3.608454in}{2.310000in}}%
\pgfusepath{clip}%
\pgfsetbuttcap%
\pgfsetmiterjoin%
\definecolor{currentfill}{rgb}{0.000000,0.600000,0.700000}%
\pgfsetfillcolor{currentfill}%
\pgfsetlinewidth{0.000000pt}%
\definecolor{currentstroke}{rgb}{0.000000,0.000000,0.000000}%
\pgfsetstrokecolor{currentstroke}%
\pgfsetstrokeopacity{0.000000}%
\pgfsetdash{}{0pt}%
\pgfpathmoveto{\pgfqpoint{2.173280in}{0.783716in}}%
\pgfpathlineto{\pgfqpoint{2.139795in}{0.782977in}}%
\pgfpathlineto{\pgfqpoint{2.143853in}{0.772670in}}%
\pgfpathlineto{\pgfqpoint{2.111050in}{0.772661in}}%
\pgfpathlineto{\pgfqpoint{2.107805in}{0.772710in}}%
\pgfpathlineto{\pgfqpoint{2.107625in}{0.827720in}}%
\pgfpathlineto{\pgfqpoint{2.096510in}{0.830344in}}%
\pgfpathlineto{\pgfqpoint{2.096641in}{0.856769in}}%
\pgfpathlineto{\pgfqpoint{2.125475in}{0.856951in}}%
\pgfpathlineto{\pgfqpoint{2.126427in}{0.854941in}}%
\pgfpathlineto{\pgfqpoint{2.170684in}{0.855388in}}%
\pgfpathlineto{\pgfqpoint{2.171444in}{0.784825in}}%
\pgfpathlineto{\pgfqpoint{2.173280in}{0.783716in}}%
\pgfpathclose%
\pgfusepath{fill}%
\end{pgfscope}%
\begin{pgfscope}%
\pgfpathrectangle{\pgfqpoint{0.100000in}{0.100000in}}{\pgfqpoint{3.608454in}{2.310000in}}%
\pgfusepath{clip}%
\pgfsetbuttcap%
\pgfsetmiterjoin%
\definecolor{currentfill}{rgb}{0.000000,0.827451,0.586275}%
\pgfsetfillcolor{currentfill}%
\pgfsetlinewidth{0.000000pt}%
\definecolor{currentstroke}{rgb}{0.000000,0.000000,0.000000}%
\pgfsetstrokecolor{currentstroke}%
\pgfsetstrokeopacity{0.000000}%
\pgfsetdash{}{0pt}%
\pgfpathmoveto{\pgfqpoint{0.976297in}{0.641465in}}%
\pgfpathlineto{\pgfqpoint{0.971445in}{0.636313in}}%
\pgfpathlineto{\pgfqpoint{0.973534in}{0.634639in}}%
\pgfpathlineto{\pgfqpoint{0.971825in}{0.632484in}}%
\pgfpathlineto{\pgfqpoint{0.974033in}{0.630789in}}%
\pgfpathlineto{\pgfqpoint{0.972344in}{0.628719in}}%
\pgfpathlineto{\pgfqpoint{0.974510in}{0.627028in}}%
\pgfpathlineto{\pgfqpoint{0.972240in}{0.625296in}}%
\pgfpathlineto{\pgfqpoint{0.967277in}{0.618907in}}%
\pgfpathlineto{\pgfqpoint{0.971558in}{0.615632in}}%
\pgfpathlineto{\pgfqpoint{0.969910in}{0.613404in}}%
\pgfpathlineto{\pgfqpoint{0.971435in}{0.612298in}}%
\pgfpathlineto{\pgfqpoint{0.964303in}{0.603251in}}%
\pgfpathlineto{\pgfqpoint{0.961995in}{0.601456in}}%
\pgfpathlineto{\pgfqpoint{0.955683in}{0.606572in}}%
\pgfpathlineto{\pgfqpoint{0.952323in}{0.602325in}}%
\pgfpathlineto{\pgfqpoint{0.946326in}{0.607097in}}%
\pgfpathlineto{\pgfqpoint{0.940885in}{0.601118in}}%
\pgfpathlineto{\pgfqpoint{0.936932in}{0.604209in}}%
\pgfpathlineto{\pgfqpoint{0.938627in}{0.606373in}}%
\pgfpathlineto{\pgfqpoint{0.936545in}{0.608078in}}%
\pgfpathlineto{\pgfqpoint{0.938173in}{0.610139in}}%
\pgfpathlineto{\pgfqpoint{0.934029in}{0.613492in}}%
\pgfpathlineto{\pgfqpoint{0.930663in}{0.609300in}}%
\pgfpathlineto{\pgfqpoint{0.926478in}{0.612658in}}%
\pgfpathlineto{\pgfqpoint{0.929884in}{0.616921in}}%
\pgfpathlineto{\pgfqpoint{0.921542in}{0.623737in}}%
\pgfpathlineto{\pgfqpoint{0.914790in}{0.615304in}}%
\pgfpathlineto{\pgfqpoint{0.904546in}{0.623960in}}%
\pgfpathlineto{\pgfqpoint{0.901313in}{0.620035in}}%
\pgfpathlineto{\pgfqpoint{0.895296in}{0.625196in}}%
\pgfpathlineto{\pgfqpoint{0.886576in}{0.615036in}}%
\pgfpathlineto{\pgfqpoint{0.884580in}{0.616759in}}%
\pgfpathlineto{\pgfqpoint{0.882807in}{0.614773in}}%
\pgfpathlineto{\pgfqpoint{0.870840in}{0.625356in}}%
\pgfpathlineto{\pgfqpoint{0.853614in}{0.640940in}}%
\pgfpathlineto{\pgfqpoint{0.845457in}{0.648542in}}%
\pgfpathlineto{\pgfqpoint{0.849066in}{0.652382in}}%
\pgfpathlineto{\pgfqpoint{0.847118in}{0.654238in}}%
\pgfpathlineto{\pgfqpoint{0.850812in}{0.658129in}}%
\pgfpathlineto{\pgfqpoint{0.848280in}{0.660517in}}%
\pgfpathlineto{\pgfqpoint{0.855680in}{0.668231in}}%
\pgfpathlineto{\pgfqpoint{0.855048in}{0.668896in}}%
\pgfpathlineto{\pgfqpoint{0.862419in}{0.676643in}}%
\pgfpathlineto{\pgfqpoint{0.863943in}{0.679699in}}%
\pgfpathlineto{\pgfqpoint{0.867167in}{0.677080in}}%
\pgfpathlineto{\pgfqpoint{0.868881in}{0.673318in}}%
\pgfpathlineto{\pgfqpoint{0.866790in}{0.671377in}}%
\pgfpathlineto{\pgfqpoint{0.864000in}{0.670092in}}%
\pgfpathlineto{\pgfqpoint{0.862316in}{0.667806in}}%
\pgfpathlineto{\pgfqpoint{0.859533in}{0.667863in}}%
\pgfpathlineto{\pgfqpoint{0.858589in}{0.662654in}}%
\pgfpathlineto{\pgfqpoint{0.861141in}{0.660797in}}%
\pgfpathlineto{\pgfqpoint{0.865258in}{0.656612in}}%
\pgfpathlineto{\pgfqpoint{0.868725in}{0.654793in}}%
\pgfpathlineto{\pgfqpoint{0.868360in}{0.652816in}}%
\pgfpathlineto{\pgfqpoint{0.870641in}{0.649968in}}%
\pgfpathlineto{\pgfqpoint{0.873832in}{0.649148in}}%
\pgfpathlineto{\pgfqpoint{0.875177in}{0.647013in}}%
\pgfpathlineto{\pgfqpoint{0.873862in}{0.644273in}}%
\pgfpathlineto{\pgfqpoint{0.879648in}{0.646559in}}%
\pgfpathlineto{\pgfqpoint{0.883020in}{0.648596in}}%
\pgfpathlineto{\pgfqpoint{0.884548in}{0.645448in}}%
\pgfpathlineto{\pgfqpoint{0.888156in}{0.644278in}}%
\pgfpathlineto{\pgfqpoint{0.888264in}{0.647509in}}%
\pgfpathlineto{\pgfqpoint{0.885826in}{0.651944in}}%
\pgfpathlineto{\pgfqpoint{0.882189in}{0.653954in}}%
\pgfpathlineto{\pgfqpoint{0.885088in}{0.656299in}}%
\pgfpathlineto{\pgfqpoint{0.885204in}{0.660563in}}%
\pgfpathlineto{\pgfqpoint{0.886202in}{0.663888in}}%
\pgfpathlineto{\pgfqpoint{0.884247in}{0.666557in}}%
\pgfpathlineto{\pgfqpoint{0.886027in}{0.671213in}}%
\pgfpathlineto{\pgfqpoint{0.888440in}{0.671601in}}%
\pgfpathlineto{\pgfqpoint{0.890129in}{0.670612in}}%
\pgfpathlineto{\pgfqpoint{0.889298in}{0.664217in}}%
\pgfpathlineto{\pgfqpoint{0.886556in}{0.660681in}}%
\pgfpathlineto{\pgfqpoint{0.886858in}{0.653674in}}%
\pgfpathlineto{\pgfqpoint{0.891233in}{0.651664in}}%
\pgfpathlineto{\pgfqpoint{0.892089in}{0.649695in}}%
\pgfpathlineto{\pgfqpoint{0.891874in}{0.645714in}}%
\pgfpathlineto{\pgfqpoint{0.897602in}{0.640349in}}%
\pgfpathlineto{\pgfqpoint{0.900250in}{0.642484in}}%
\pgfpathlineto{\pgfqpoint{0.901360in}{0.646766in}}%
\pgfpathlineto{\pgfqpoint{0.898368in}{0.648567in}}%
\pgfpathlineto{\pgfqpoint{0.897419in}{0.652291in}}%
\pgfpathlineto{\pgfqpoint{0.893693in}{0.653344in}}%
\pgfpathlineto{\pgfqpoint{0.891532in}{0.652766in}}%
\pgfpathlineto{\pgfqpoint{0.889610in}{0.655212in}}%
\pgfpathlineto{\pgfqpoint{0.889414in}{0.661725in}}%
\pgfpathlineto{\pgfqpoint{0.891701in}{0.663459in}}%
\pgfpathlineto{\pgfqpoint{0.897189in}{0.666384in}}%
\pgfpathlineto{\pgfqpoint{0.896490in}{0.669062in}}%
\pgfpathlineto{\pgfqpoint{0.890332in}{0.672207in}}%
\pgfpathlineto{\pgfqpoint{0.891668in}{0.673920in}}%
\pgfpathlineto{\pgfqpoint{0.887467in}{0.677256in}}%
\pgfpathlineto{\pgfqpoint{0.884828in}{0.680240in}}%
\pgfpathlineto{\pgfqpoint{0.881627in}{0.686225in}}%
\pgfpathlineto{\pgfqpoint{0.885981in}{0.691916in}}%
\pgfpathlineto{\pgfqpoint{0.888139in}{0.698608in}}%
\pgfpathlineto{\pgfqpoint{0.888448in}{0.701521in}}%
\pgfpathlineto{\pgfqpoint{0.887314in}{0.706916in}}%
\pgfpathlineto{\pgfqpoint{0.887408in}{0.709755in}}%
\pgfpathlineto{\pgfqpoint{0.886741in}{0.720168in}}%
\pgfpathlineto{\pgfqpoint{0.894146in}{0.712972in}}%
\pgfpathlineto{\pgfqpoint{0.898458in}{0.717465in}}%
\pgfpathlineto{\pgfqpoint{0.913662in}{0.703154in}}%
\pgfpathlineto{\pgfqpoint{0.915131in}{0.704745in}}%
\pgfpathlineto{\pgfqpoint{0.921351in}{0.699050in}}%
\pgfpathlineto{\pgfqpoint{0.919895in}{0.697448in}}%
\pgfpathlineto{\pgfqpoint{0.924102in}{0.693649in}}%
\pgfpathlineto{\pgfqpoint{0.925549in}{0.695258in}}%
\pgfpathlineto{\pgfqpoint{0.929228in}{0.691967in}}%
\pgfpathlineto{\pgfqpoint{0.927789in}{0.690351in}}%
\pgfpathlineto{\pgfqpoint{0.936363in}{0.682792in}}%
\pgfpathlineto{\pgfqpoint{0.937775in}{0.684411in}}%
\pgfpathlineto{\pgfqpoint{0.941559in}{0.681130in}}%
\pgfpathlineto{\pgfqpoint{0.940158in}{0.679506in}}%
\pgfpathlineto{\pgfqpoint{0.954477in}{0.667338in}}%
\pgfpathlineto{\pgfqpoint{0.963280in}{0.660093in}}%
\pgfpathlineto{\pgfqpoint{0.961484in}{0.657889in}}%
\pgfpathlineto{\pgfqpoint{0.969832in}{0.651052in}}%
\pgfpathlineto{\pgfqpoint{0.968183in}{0.648986in}}%
\pgfpathlineto{\pgfqpoint{0.976829in}{0.642165in}}%
\pgfpathlineto{\pgfqpoint{0.976297in}{0.641465in}}%
\pgfpathclose%
\pgfusepath{fill}%
\end{pgfscope}%
\begin{pgfscope}%
\pgfpathrectangle{\pgfqpoint{0.100000in}{0.100000in}}{\pgfqpoint{3.608454in}{2.310000in}}%
\pgfusepath{clip}%
\pgfsetbuttcap%
\pgfsetmiterjoin%
\definecolor{currentfill}{rgb}{0.000000,0.796078,0.601961}%
\pgfsetfillcolor{currentfill}%
\pgfsetlinewidth{0.000000pt}%
\definecolor{currentstroke}{rgb}{0.000000,0.000000,0.000000}%
\pgfsetstrokecolor{currentstroke}%
\pgfsetstrokeopacity{0.000000}%
\pgfsetdash{}{0pt}%
\pgfpathmoveto{\pgfqpoint{1.801456in}{1.757724in}}%
\pgfpathlineto{\pgfqpoint{1.800009in}{1.732253in}}%
\pgfpathlineto{\pgfqpoint{1.800957in}{1.721865in}}%
\pgfpathlineto{\pgfqpoint{1.793255in}{1.722620in}}%
\pgfpathlineto{\pgfqpoint{1.787775in}{1.725501in}}%
\pgfpathlineto{\pgfqpoint{1.775409in}{1.724494in}}%
\pgfpathlineto{\pgfqpoint{1.768321in}{1.732235in}}%
\pgfpathlineto{\pgfqpoint{1.760148in}{1.734344in}}%
\pgfpathlineto{\pgfqpoint{1.762841in}{1.760148in}}%
\pgfpathlineto{\pgfqpoint{1.801456in}{1.757724in}}%
\pgfpathclose%
\pgfusepath{fill}%
\end{pgfscope}%
\begin{pgfscope}%
\pgfpathrectangle{\pgfqpoint{0.100000in}{0.100000in}}{\pgfqpoint{3.608454in}{2.310000in}}%
\pgfusepath{clip}%
\pgfsetbuttcap%
\pgfsetmiterjoin%
\definecolor{currentfill}{rgb}{0.000000,0.462745,0.768627}%
\pgfsetfillcolor{currentfill}%
\pgfsetlinewidth{0.000000pt}%
\definecolor{currentstroke}{rgb}{0.000000,0.000000,0.000000}%
\pgfsetstrokecolor{currentstroke}%
\pgfsetstrokeopacity{0.000000}%
\pgfsetdash{}{0pt}%
\pgfpathmoveto{\pgfqpoint{2.387375in}{1.567036in}}%
\pgfpathlineto{\pgfqpoint{2.380372in}{1.562314in}}%
\pgfpathlineto{\pgfqpoint{2.376029in}{1.564684in}}%
\pgfpathlineto{\pgfqpoint{2.348758in}{1.564236in}}%
\pgfpathlineto{\pgfqpoint{2.348123in}{1.578046in}}%
\pgfpathlineto{\pgfqpoint{2.346907in}{1.605656in}}%
\pgfpathlineto{\pgfqpoint{2.360507in}{1.606246in}}%
\pgfpathlineto{\pgfqpoint{2.360207in}{1.613167in}}%
\pgfpathlineto{\pgfqpoint{2.371380in}{1.613652in}}%
\pgfpathlineto{\pgfqpoint{2.374950in}{1.609365in}}%
\pgfpathlineto{\pgfqpoint{2.376710in}{1.601691in}}%
\pgfpathlineto{\pgfqpoint{2.381293in}{1.599262in}}%
\pgfpathlineto{\pgfqpoint{2.390701in}{1.593517in}}%
\pgfpathlineto{\pgfqpoint{2.392472in}{1.584764in}}%
\pgfpathlineto{\pgfqpoint{2.390834in}{1.569129in}}%
\pgfpathlineto{\pgfqpoint{2.387375in}{1.567036in}}%
\pgfpathclose%
\pgfusepath{fill}%
\end{pgfscope}%
\begin{pgfscope}%
\pgfpathrectangle{\pgfqpoint{0.100000in}{0.100000in}}{\pgfqpoint{3.608454in}{2.310000in}}%
\pgfusepath{clip}%
\pgfsetbuttcap%
\pgfsetmiterjoin%
\definecolor{currentfill}{rgb}{0.000000,0.364706,0.817647}%
\pgfsetfillcolor{currentfill}%
\pgfsetlinewidth{0.000000pt}%
\definecolor{currentstroke}{rgb}{0.000000,0.000000,0.000000}%
\pgfsetstrokecolor{currentstroke}%
\pgfsetstrokeopacity{0.000000}%
\pgfsetdash{}{0pt}%
\pgfpathmoveto{\pgfqpoint{1.853837in}{1.424663in}}%
\pgfpathlineto{\pgfqpoint{1.881222in}{1.423490in}}%
\pgfpathlineto{\pgfqpoint{1.894644in}{1.422961in}}%
\pgfpathlineto{\pgfqpoint{1.893270in}{1.388628in}}%
\pgfpathlineto{\pgfqpoint{1.860470in}{1.390012in}}%
\pgfpathlineto{\pgfqpoint{1.824934in}{1.391676in}}%
\pgfpathlineto{\pgfqpoint{1.826801in}{1.425965in}}%
\pgfpathlineto{\pgfqpoint{1.853837in}{1.424663in}}%
\pgfpathclose%
\pgfusepath{fill}%
\end{pgfscope}%
\begin{pgfscope}%
\pgfpathrectangle{\pgfqpoint{0.100000in}{0.100000in}}{\pgfqpoint{3.608454in}{2.310000in}}%
\pgfusepath{clip}%
\pgfsetbuttcap%
\pgfsetmiterjoin%
\definecolor{currentfill}{rgb}{0.000000,0.494118,0.752941}%
\pgfsetfillcolor{currentfill}%
\pgfsetlinewidth{0.000000pt}%
\definecolor{currentstroke}{rgb}{0.000000,0.000000,0.000000}%
\pgfsetstrokecolor{currentstroke}%
\pgfsetstrokeopacity{0.000000}%
\pgfsetdash{}{0pt}%
\pgfpathmoveto{\pgfqpoint{2.374602in}{1.377391in}}%
\pgfpathlineto{\pgfqpoint{2.376576in}{1.388834in}}%
\pgfpathlineto{\pgfqpoint{2.371479in}{1.401180in}}%
\pgfpathlineto{\pgfqpoint{2.375453in}{1.412534in}}%
\pgfpathlineto{\pgfqpoint{2.354524in}{1.412008in}}%
\pgfpathlineto{\pgfqpoint{2.353565in}{1.439495in}}%
\pgfpathlineto{\pgfqpoint{2.353369in}{1.446719in}}%
\pgfpathlineto{\pgfqpoint{2.381085in}{1.447410in}}%
\pgfpathlineto{\pgfqpoint{2.381360in}{1.440493in}}%
\pgfpathlineto{\pgfqpoint{2.396575in}{1.440884in}}%
\pgfpathlineto{\pgfqpoint{2.389932in}{1.435148in}}%
\pgfpathlineto{\pgfqpoint{2.391369in}{1.431061in}}%
\pgfpathlineto{\pgfqpoint{2.397384in}{1.430530in}}%
\pgfpathlineto{\pgfqpoint{2.409257in}{1.435587in}}%
\pgfpathlineto{\pgfqpoint{2.410721in}{1.410092in}}%
\pgfpathlineto{\pgfqpoint{2.415980in}{1.389550in}}%
\pgfpathlineto{\pgfqpoint{2.402134in}{1.388642in}}%
\pgfpathlineto{\pgfqpoint{2.403570in}{1.368279in}}%
\pgfpathlineto{\pgfqpoint{2.393816in}{1.360813in}}%
\pgfpathlineto{\pgfqpoint{2.383166in}{1.360351in}}%
\pgfpathlineto{\pgfqpoint{2.376719in}{1.362091in}}%
\pgfpathlineto{\pgfqpoint{2.374602in}{1.377391in}}%
\pgfpathclose%
\pgfusepath{fill}%
\end{pgfscope}%
\begin{pgfscope}%
\pgfpathrectangle{\pgfqpoint{0.100000in}{0.100000in}}{\pgfqpoint{3.608454in}{2.310000in}}%
\pgfusepath{clip}%
\pgfsetbuttcap%
\pgfsetmiterjoin%
\definecolor{currentfill}{rgb}{0.000000,0.639216,0.680392}%
\pgfsetfillcolor{currentfill}%
\pgfsetlinewidth{0.000000pt}%
\definecolor{currentstroke}{rgb}{0.000000,0.000000,0.000000}%
\pgfsetstrokecolor{currentstroke}%
\pgfsetstrokeopacity{0.000000}%
\pgfsetdash{}{0pt}%
\pgfpathmoveto{\pgfqpoint{3.177265in}{1.821081in}}%
\pgfpathlineto{\pgfqpoint{3.173354in}{1.830395in}}%
\pgfpathlineto{\pgfqpoint{3.177753in}{1.838567in}}%
\pgfpathlineto{\pgfqpoint{3.175101in}{1.843253in}}%
\pgfpathlineto{\pgfqpoint{3.166217in}{1.845402in}}%
\pgfpathlineto{\pgfqpoint{3.162379in}{1.854420in}}%
\pgfpathlineto{\pgfqpoint{3.171115in}{1.863602in}}%
\pgfpathlineto{\pgfqpoint{3.169608in}{1.868162in}}%
\pgfpathlineto{\pgfqpoint{3.181998in}{1.877869in}}%
\pgfpathlineto{\pgfqpoint{3.186045in}{1.883942in}}%
\pgfpathlineto{\pgfqpoint{3.187678in}{1.891082in}}%
\pgfpathlineto{\pgfqpoint{3.198874in}{1.908865in}}%
\pgfpathlineto{\pgfqpoint{3.215748in}{1.929678in}}%
\pgfpathlineto{\pgfqpoint{3.223456in}{1.936959in}}%
\pgfpathlineto{\pgfqpoint{3.230888in}{1.941273in}}%
\pgfpathlineto{\pgfqpoint{3.236756in}{1.940890in}}%
\pgfpathlineto{\pgfqpoint{3.242124in}{1.938669in}}%
\pgfpathlineto{\pgfqpoint{3.262366in}{1.879593in}}%
\pgfpathlineto{\pgfqpoint{3.263105in}{1.873951in}}%
\pgfpathlineto{\pgfqpoint{3.245888in}{1.867665in}}%
\pgfpathlineto{\pgfqpoint{3.234624in}{1.863539in}}%
\pgfpathlineto{\pgfqpoint{3.227847in}{1.865817in}}%
\pgfpathlineto{\pgfqpoint{3.239480in}{1.829258in}}%
\pgfpathlineto{\pgfqpoint{3.219039in}{1.808893in}}%
\pgfpathlineto{\pgfqpoint{3.205664in}{1.810106in}}%
\pgfpathlineto{\pgfqpoint{3.200994in}{1.826823in}}%
\pgfpathlineto{\pgfqpoint{3.186645in}{1.825233in}}%
\pgfpathlineto{\pgfqpoint{3.177265in}{1.821081in}}%
\pgfpathclose%
\pgfusepath{fill}%
\end{pgfscope}%
\begin{pgfscope}%
\pgfpathrectangle{\pgfqpoint{0.100000in}{0.100000in}}{\pgfqpoint{3.608454in}{2.310000in}}%
\pgfusepath{clip}%
\pgfsetbuttcap%
\pgfsetmiterjoin%
\definecolor{currentfill}{rgb}{0.000000,0.647059,0.676471}%
\pgfsetfillcolor{currentfill}%
\pgfsetlinewidth{0.000000pt}%
\definecolor{currentstroke}{rgb}{0.000000,0.000000,0.000000}%
\pgfsetstrokecolor{currentstroke}%
\pgfsetstrokeopacity{0.000000}%
\pgfsetdash{}{0pt}%
\pgfpathmoveto{\pgfqpoint{2.268476in}{1.251773in}}%
\pgfpathlineto{\pgfqpoint{2.260153in}{1.246809in}}%
\pgfpathlineto{\pgfqpoint{2.255244in}{1.254033in}}%
\pgfpathlineto{\pgfqpoint{2.227021in}{1.254178in}}%
\pgfpathlineto{\pgfqpoint{2.227204in}{1.266822in}}%
\pgfpathlineto{\pgfqpoint{2.225942in}{1.276603in}}%
\pgfpathlineto{\pgfqpoint{2.226458in}{1.316665in}}%
\pgfpathlineto{\pgfqpoint{2.240062in}{1.316251in}}%
\pgfpathlineto{\pgfqpoint{2.240098in}{1.309360in}}%
\pgfpathlineto{\pgfqpoint{2.253822in}{1.296576in}}%
\pgfpathlineto{\pgfqpoint{2.267557in}{1.296520in}}%
\pgfpathlineto{\pgfqpoint{2.267441in}{1.289875in}}%
\pgfpathlineto{\pgfqpoint{2.274008in}{1.289693in}}%
\pgfpathlineto{\pgfqpoint{2.280744in}{1.286387in}}%
\pgfpathlineto{\pgfqpoint{2.282058in}{1.264607in}}%
\pgfpathlineto{\pgfqpoint{2.268296in}{1.264580in}}%
\pgfpathlineto{\pgfqpoint{2.268476in}{1.251773in}}%
\pgfpathclose%
\pgfusepath{fill}%
\end{pgfscope}%
\begin{pgfscope}%
\pgfpathrectangle{\pgfqpoint{0.100000in}{0.100000in}}{\pgfqpoint{3.608454in}{2.310000in}}%
\pgfusepath{clip}%
\pgfsetbuttcap%
\pgfsetmiterjoin%
\definecolor{currentfill}{rgb}{0.000000,0.709804,0.645098}%
\pgfsetfillcolor{currentfill}%
\pgfsetlinewidth{0.000000pt}%
\definecolor{currentstroke}{rgb}{0.000000,0.000000,0.000000}%
\pgfsetstrokecolor{currentstroke}%
\pgfsetstrokeopacity{0.000000}%
\pgfsetdash{}{0pt}%
\pgfpathmoveto{\pgfqpoint{2.434616in}{1.075007in}}%
\pgfpathlineto{\pgfqpoint{2.426249in}{1.075405in}}%
\pgfpathlineto{\pgfqpoint{2.427918in}{1.068838in}}%
\pgfpathlineto{\pgfqpoint{2.420217in}{1.066185in}}%
\pgfpathlineto{\pgfqpoint{2.410748in}{1.065840in}}%
\pgfpathlineto{\pgfqpoint{2.408349in}{1.109943in}}%
\pgfpathlineto{\pgfqpoint{2.406026in}{1.117238in}}%
\pgfpathlineto{\pgfqpoint{2.411938in}{1.125058in}}%
\pgfpathlineto{\pgfqpoint{2.420239in}{1.132410in}}%
\pgfpathlineto{\pgfqpoint{2.421101in}{1.141357in}}%
\pgfpathlineto{\pgfqpoint{2.410619in}{1.149798in}}%
\pgfpathlineto{\pgfqpoint{2.414649in}{1.160533in}}%
\pgfpathlineto{\pgfqpoint{2.426666in}{1.161074in}}%
\pgfpathlineto{\pgfqpoint{2.426853in}{1.142086in}}%
\pgfpathlineto{\pgfqpoint{2.440887in}{1.142638in}}%
\pgfpathlineto{\pgfqpoint{2.450987in}{1.144593in}}%
\pgfpathlineto{\pgfqpoint{2.454276in}{1.139497in}}%
\pgfpathlineto{\pgfqpoint{2.455346in}{1.132898in}}%
\pgfpathlineto{\pgfqpoint{2.445412in}{1.130104in}}%
\pgfpathlineto{\pgfqpoint{2.452158in}{1.122751in}}%
\pgfpathlineto{\pgfqpoint{2.446875in}{1.118895in}}%
\pgfpathlineto{\pgfqpoint{2.445220in}{1.109326in}}%
\pgfpathlineto{\pgfqpoint{2.442276in}{1.103468in}}%
\pgfpathlineto{\pgfqpoint{2.442664in}{1.097005in}}%
\pgfpathlineto{\pgfqpoint{2.439328in}{1.092627in}}%
\pgfpathlineto{\pgfqpoint{2.431130in}{1.090720in}}%
\pgfpathlineto{\pgfqpoint{2.432045in}{1.079916in}}%
\pgfpathlineto{\pgfqpoint{2.434616in}{1.075007in}}%
\pgfpathclose%
\pgfusepath{fill}%
\end{pgfscope}%
\begin{pgfscope}%
\pgfpathrectangle{\pgfqpoint{0.100000in}{0.100000in}}{\pgfqpoint{3.608454in}{2.310000in}}%
\pgfusepath{clip}%
\pgfsetbuttcap%
\pgfsetmiterjoin%
\definecolor{currentfill}{rgb}{0.000000,0.552941,0.723529}%
\pgfsetfillcolor{currentfill}%
\pgfsetlinewidth{0.000000pt}%
\definecolor{currentstroke}{rgb}{0.000000,0.000000,0.000000}%
\pgfsetstrokecolor{currentstroke}%
\pgfsetstrokeopacity{0.000000}%
\pgfsetdash{}{0pt}%
\pgfpathmoveto{\pgfqpoint{2.994468in}{0.861760in}}%
\pgfpathlineto{\pgfqpoint{2.991213in}{0.868134in}}%
\pgfpathlineto{\pgfqpoint{2.962520in}{0.870538in}}%
\pgfpathlineto{\pgfqpoint{2.954746in}{0.867311in}}%
\pgfpathlineto{\pgfqpoint{2.950142in}{0.872561in}}%
\pgfpathlineto{\pgfqpoint{2.951258in}{0.878733in}}%
\pgfpathlineto{\pgfqpoint{2.943431in}{0.882150in}}%
\pgfpathlineto{\pgfqpoint{2.932829in}{0.881990in}}%
\pgfpathlineto{\pgfqpoint{2.935640in}{0.887719in}}%
\pgfpathlineto{\pgfqpoint{2.943276in}{0.894232in}}%
\pgfpathlineto{\pgfqpoint{2.940759in}{0.897010in}}%
\pgfpathlineto{\pgfqpoint{2.944266in}{0.903605in}}%
\pgfpathlineto{\pgfqpoint{2.950977in}{0.910809in}}%
\pgfpathlineto{\pgfqpoint{2.959107in}{0.909040in}}%
\pgfpathlineto{\pgfqpoint{2.962570in}{0.910798in}}%
\pgfpathlineto{\pgfqpoint{2.965871in}{0.919067in}}%
\pgfpathlineto{\pgfqpoint{2.980224in}{0.924118in}}%
\pgfpathlineto{\pgfqpoint{2.986683in}{0.921479in}}%
\pgfpathlineto{\pgfqpoint{2.999992in}{0.934399in}}%
\pgfpathlineto{\pgfqpoint{3.003733in}{0.931980in}}%
\pgfpathlineto{\pgfqpoint{3.003885in}{0.926116in}}%
\pgfpathlineto{\pgfqpoint{3.010014in}{0.919433in}}%
\pgfpathlineto{\pgfqpoint{3.011589in}{0.908504in}}%
\pgfpathlineto{\pgfqpoint{3.015378in}{0.900512in}}%
\pgfpathlineto{\pgfqpoint{3.006499in}{0.892231in}}%
\pgfpathlineto{\pgfqpoint{3.008209in}{0.883079in}}%
\pgfpathlineto{\pgfqpoint{3.014555in}{0.877148in}}%
\pgfpathlineto{\pgfqpoint{3.016488in}{0.872188in}}%
\pgfpathlineto{\pgfqpoint{2.994468in}{0.861760in}}%
\pgfpathclose%
\pgfusepath{fill}%
\end{pgfscope}%
\begin{pgfscope}%
\pgfpathrectangle{\pgfqpoint{0.100000in}{0.100000in}}{\pgfqpoint{3.608454in}{2.310000in}}%
\pgfusepath{clip}%
\pgfsetbuttcap%
\pgfsetmiterjoin%
\definecolor{currentfill}{rgb}{0.000000,0.631373,0.684314}%
\pgfsetfillcolor{currentfill}%
\pgfsetlinewidth{0.000000pt}%
\definecolor{currentstroke}{rgb}{0.000000,0.000000,0.000000}%
\pgfsetstrokecolor{currentstroke}%
\pgfsetstrokeopacity{0.000000}%
\pgfsetdash{}{0pt}%
\pgfpathmoveto{\pgfqpoint{1.438374in}{0.821513in}}%
\pgfpathlineto{\pgfqpoint{1.435480in}{0.791748in}}%
\pgfpathlineto{\pgfqpoint{1.429579in}{0.734241in}}%
\pgfpathlineto{\pgfqpoint{1.426837in}{0.715528in}}%
\pgfpathlineto{\pgfqpoint{1.422243in}{0.713258in}}%
\pgfpathlineto{\pgfqpoint{1.413714in}{0.723890in}}%
\pgfpathlineto{\pgfqpoint{1.407758in}{0.729084in}}%
\pgfpathlineto{\pgfqpoint{1.396001in}{0.734660in}}%
\pgfpathlineto{\pgfqpoint{1.395928in}{0.737162in}}%
\pgfpathlineto{\pgfqpoint{1.386189in}{0.746307in}}%
\pgfpathlineto{\pgfqpoint{1.383892in}{0.754060in}}%
\pgfpathlineto{\pgfqpoint{1.373173in}{0.762125in}}%
\pgfpathlineto{\pgfqpoint{1.362767in}{0.779271in}}%
\pgfpathlineto{\pgfqpoint{1.354361in}{0.783205in}}%
\pgfpathlineto{\pgfqpoint{1.346623in}{0.789592in}}%
\pgfpathlineto{\pgfqpoint{1.337607in}{0.811852in}}%
\pgfpathlineto{\pgfqpoint{1.328213in}{0.817202in}}%
\pgfpathlineto{\pgfqpoint{1.276635in}{0.824008in}}%
\pgfpathlineto{\pgfqpoint{1.285282in}{0.888829in}}%
\pgfpathlineto{\pgfqpoint{1.287163in}{0.902597in}}%
\pgfpathlineto{\pgfqpoint{1.314425in}{0.898871in}}%
\pgfpathlineto{\pgfqpoint{1.314868in}{0.902303in}}%
\pgfpathlineto{\pgfqpoint{1.353414in}{0.915777in}}%
\pgfpathlineto{\pgfqpoint{1.350085in}{0.908598in}}%
\pgfpathlineto{\pgfqpoint{1.340564in}{0.833080in}}%
\pgfpathlineto{\pgfqpoint{1.404085in}{0.825375in}}%
\pgfpathlineto{\pgfqpoint{1.438374in}{0.821513in}}%
\pgfpathclose%
\pgfusepath{fill}%
\end{pgfscope}%
\begin{pgfscope}%
\pgfpathrectangle{\pgfqpoint{0.100000in}{0.100000in}}{\pgfqpoint{3.608454in}{2.310000in}}%
\pgfusepath{clip}%
\pgfsetbuttcap%
\pgfsetmiterjoin%
\definecolor{currentfill}{rgb}{0.000000,0.564706,0.717647}%
\pgfsetfillcolor{currentfill}%
\pgfsetlinewidth{0.000000pt}%
\definecolor{currentstroke}{rgb}{0.000000,0.000000,0.000000}%
\pgfsetstrokecolor{currentstroke}%
\pgfsetstrokeopacity{0.000000}%
\pgfsetdash{}{0pt}%
\pgfpathmoveto{\pgfqpoint{2.765201in}{0.880136in}}%
\pgfpathlineto{\pgfqpoint{2.760875in}{0.889582in}}%
\pgfpathlineto{\pgfqpoint{2.782280in}{0.891779in}}%
\pgfpathlineto{\pgfqpoint{2.779766in}{0.919893in}}%
\pgfpathlineto{\pgfqpoint{2.802911in}{0.922322in}}%
\pgfpathlineto{\pgfqpoint{2.802438in}{0.905479in}}%
\pgfpathlineto{\pgfqpoint{2.805495in}{0.895708in}}%
\pgfpathlineto{\pgfqpoint{2.814267in}{0.888921in}}%
\pgfpathlineto{\pgfqpoint{2.821573in}{0.886623in}}%
\pgfpathlineto{\pgfqpoint{2.812799in}{0.870744in}}%
\pgfpathlineto{\pgfqpoint{2.800195in}{0.867107in}}%
\pgfpathlineto{\pgfqpoint{2.799223in}{0.863800in}}%
\pgfpathlineto{\pgfqpoint{2.801485in}{0.843067in}}%
\pgfpathlineto{\pgfqpoint{2.783833in}{0.840916in}}%
\pgfpathlineto{\pgfqpoint{2.775432in}{0.832235in}}%
\pgfpathlineto{\pgfqpoint{2.776392in}{0.826653in}}%
\pgfpathlineto{\pgfqpoint{2.767748in}{0.825602in}}%
\pgfpathlineto{\pgfqpoint{2.762136in}{0.831926in}}%
\pgfpathlineto{\pgfqpoint{2.751930in}{0.830708in}}%
\pgfpathlineto{\pgfqpoint{2.749632in}{0.837509in}}%
\pgfpathlineto{\pgfqpoint{2.748134in}{0.851370in}}%
\pgfpathlineto{\pgfqpoint{2.755065in}{0.852187in}}%
\pgfpathlineto{\pgfqpoint{2.754330in}{0.856732in}}%
\pgfpathlineto{\pgfqpoint{2.772613in}{0.859054in}}%
\pgfpathlineto{\pgfqpoint{2.776169in}{0.862453in}}%
\pgfpathlineto{\pgfqpoint{2.768066in}{0.872209in}}%
\pgfpathlineto{\pgfqpoint{2.765201in}{0.880136in}}%
\pgfpathclose%
\pgfusepath{fill}%
\end{pgfscope}%
\begin{pgfscope}%
\pgfpathrectangle{\pgfqpoint{0.100000in}{0.100000in}}{\pgfqpoint{3.608454in}{2.310000in}}%
\pgfusepath{clip}%
\pgfsetbuttcap%
\pgfsetmiterjoin%
\definecolor{currentfill}{rgb}{0.000000,0.494118,0.752941}%
\pgfsetfillcolor{currentfill}%
\pgfsetlinewidth{0.000000pt}%
\definecolor{currentstroke}{rgb}{0.000000,0.000000,0.000000}%
\pgfsetstrokecolor{currentstroke}%
\pgfsetstrokeopacity{0.000000}%
\pgfsetdash{}{0pt}%
\pgfpathmoveto{\pgfqpoint{2.682505in}{0.872318in}}%
\pgfpathlineto{\pgfqpoint{2.680466in}{0.868375in}}%
\pgfpathlineto{\pgfqpoint{2.660419in}{0.866374in}}%
\pgfpathlineto{\pgfqpoint{2.660793in}{0.862897in}}%
\pgfpathlineto{\pgfqpoint{2.647203in}{0.861864in}}%
\pgfpathlineto{\pgfqpoint{2.640493in}{0.861198in}}%
\pgfpathlineto{\pgfqpoint{2.639182in}{0.874953in}}%
\pgfpathlineto{\pgfqpoint{2.648704in}{0.875737in}}%
\pgfpathlineto{\pgfqpoint{2.647138in}{0.892648in}}%
\pgfpathlineto{\pgfqpoint{2.640666in}{0.895938in}}%
\pgfpathlineto{\pgfqpoint{2.636464in}{0.900943in}}%
\pgfpathlineto{\pgfqpoint{2.635819in}{0.907320in}}%
\pgfpathlineto{\pgfqpoint{2.625197in}{0.912348in}}%
\pgfpathlineto{\pgfqpoint{2.617467in}{0.921208in}}%
\pgfpathlineto{\pgfqpoint{2.618130in}{0.927031in}}%
\pgfpathlineto{\pgfqpoint{2.623213in}{0.935505in}}%
\pgfpathlineto{\pgfqpoint{2.628827in}{0.937113in}}%
\pgfpathlineto{\pgfqpoint{2.631914in}{0.941959in}}%
\pgfpathlineto{\pgfqpoint{2.632392in}{0.948970in}}%
\pgfpathlineto{\pgfqpoint{2.635731in}{0.956087in}}%
\pgfpathlineto{\pgfqpoint{2.627618in}{0.959226in}}%
\pgfpathlineto{\pgfqpoint{2.625447in}{0.963052in}}%
\pgfpathlineto{\pgfqpoint{2.623076in}{0.990093in}}%
\pgfpathlineto{\pgfqpoint{2.622963in}{0.991238in}}%
\pgfpathlineto{\pgfqpoint{2.657488in}{0.993619in}}%
\pgfpathlineto{\pgfqpoint{2.664313in}{0.994103in}}%
\pgfpathlineto{\pgfqpoint{2.677206in}{0.979100in}}%
\pgfpathlineto{\pgfqpoint{2.675715in}{0.974764in}}%
\pgfpathlineto{\pgfqpoint{2.678957in}{0.970072in}}%
\pgfpathlineto{\pgfqpoint{2.684888in}{0.971060in}}%
\pgfpathlineto{\pgfqpoint{2.689264in}{0.964347in}}%
\pgfpathlineto{\pgfqpoint{2.694783in}{0.960330in}}%
\pgfpathlineto{\pgfqpoint{2.696367in}{0.955422in}}%
\pgfpathlineto{\pgfqpoint{2.690775in}{0.946928in}}%
\pgfpathlineto{\pgfqpoint{2.686513in}{0.939240in}}%
\pgfpathlineto{\pgfqpoint{2.688862in}{0.935943in}}%
\pgfpathlineto{\pgfqpoint{2.677615in}{0.922675in}}%
\pgfpathlineto{\pgfqpoint{2.679726in}{0.915262in}}%
\pgfpathlineto{\pgfqpoint{2.673491in}{0.910352in}}%
\pgfpathlineto{\pgfqpoint{2.670991in}{0.904805in}}%
\pgfpathlineto{\pgfqpoint{2.672768in}{0.884091in}}%
\pgfpathlineto{\pgfqpoint{2.676747in}{0.876531in}}%
\pgfpathlineto{\pgfqpoint{2.682505in}{0.872318in}}%
\pgfpathclose%
\pgfusepath{fill}%
\end{pgfscope}%
\begin{pgfscope}%
\pgfpathrectangle{\pgfqpoint{0.100000in}{0.100000in}}{\pgfqpoint{3.608454in}{2.310000in}}%
\pgfusepath{clip}%
\pgfsetbuttcap%
\pgfsetmiterjoin%
\definecolor{currentfill}{rgb}{0.000000,0.458824,0.770588}%
\pgfsetfillcolor{currentfill}%
\pgfsetlinewidth{0.000000pt}%
\definecolor{currentstroke}{rgb}{0.000000,0.000000,0.000000}%
\pgfsetstrokecolor{currentstroke}%
\pgfsetstrokeopacity{0.000000}%
\pgfsetdash{}{0pt}%
\pgfpathmoveto{\pgfqpoint{2.032756in}{2.134709in}}%
\pgfpathlineto{\pgfqpoint{2.032428in}{2.112015in}}%
\pgfpathlineto{\pgfqpoint{2.033124in}{2.098129in}}%
\pgfpathlineto{\pgfqpoint{1.991978in}{2.098902in}}%
\pgfpathlineto{\pgfqpoint{1.993550in}{2.092414in}}%
\pgfpathlineto{\pgfqpoint{1.991966in}{2.086284in}}%
\pgfpathlineto{\pgfqpoint{1.993953in}{2.077757in}}%
\pgfpathlineto{\pgfqpoint{1.992478in}{2.070813in}}%
\pgfpathlineto{\pgfqpoint{1.951799in}{2.072044in}}%
\pgfpathlineto{\pgfqpoint{1.931043in}{2.072760in}}%
\pgfpathlineto{\pgfqpoint{1.931568in}{2.086739in}}%
\pgfpathlineto{\pgfqpoint{1.930695in}{2.100743in}}%
\pgfpathlineto{\pgfqpoint{1.896095in}{2.102164in}}%
\pgfpathlineto{\pgfqpoint{1.895146in}{2.116164in}}%
\pgfpathlineto{\pgfqpoint{1.896146in}{2.138772in}}%
\pgfpathlineto{\pgfqpoint{1.946219in}{2.136867in}}%
\pgfpathlineto{\pgfqpoint{1.989377in}{2.135633in}}%
\pgfpathlineto{\pgfqpoint{2.032756in}{2.134709in}}%
\pgfpathclose%
\pgfusepath{fill}%
\end{pgfscope}%
\begin{pgfscope}%
\pgfpathrectangle{\pgfqpoint{0.100000in}{0.100000in}}{\pgfqpoint{3.608454in}{2.310000in}}%
\pgfusepath{clip}%
\pgfsetbuttcap%
\pgfsetmiterjoin%
\definecolor{currentfill}{rgb}{0.000000,0.592157,0.703922}%
\pgfsetfillcolor{currentfill}%
\pgfsetlinewidth{0.000000pt}%
\definecolor{currentstroke}{rgb}{0.000000,0.000000,0.000000}%
\pgfsetstrokecolor{currentstroke}%
\pgfsetstrokeopacity{0.000000}%
\pgfsetdash{}{0pt}%
\pgfpathmoveto{\pgfqpoint{1.805378in}{0.356576in}}%
\pgfpathlineto{\pgfqpoint{1.803098in}{0.369070in}}%
\pgfpathlineto{\pgfqpoint{1.799243in}{0.378831in}}%
\pgfpathlineto{\pgfqpoint{1.791844in}{0.386485in}}%
\pgfpathlineto{\pgfqpoint{1.787001in}{0.397445in}}%
\pgfpathlineto{\pgfqpoint{1.789197in}{0.406267in}}%
\pgfpathlineto{\pgfqpoint{1.784237in}{0.419600in}}%
\pgfpathlineto{\pgfqpoint{1.786564in}{0.431356in}}%
\pgfpathlineto{\pgfqpoint{1.783350in}{0.432395in}}%
\pgfpathlineto{\pgfqpoint{1.782660in}{0.441230in}}%
\pgfpathlineto{\pgfqpoint{1.770087in}{0.446743in}}%
\pgfpathlineto{\pgfqpoint{1.764022in}{0.455756in}}%
\pgfpathlineto{\pgfqpoint{1.759560in}{0.458124in}}%
\pgfpathlineto{\pgfqpoint{1.758847in}{0.466132in}}%
\pgfpathlineto{\pgfqpoint{1.752364in}{0.474146in}}%
\pgfpathlineto{\pgfqpoint{1.746457in}{0.486543in}}%
\pgfpathlineto{\pgfqpoint{1.737580in}{0.491300in}}%
\pgfpathlineto{\pgfqpoint{1.744553in}{0.490983in}}%
\pgfpathlineto{\pgfqpoint{1.795506in}{0.488863in}}%
\pgfpathlineto{\pgfqpoint{1.795200in}{0.474847in}}%
\pgfpathlineto{\pgfqpoint{1.836840in}{0.475148in}}%
\pgfpathlineto{\pgfqpoint{1.834828in}{0.418520in}}%
\pgfpathlineto{\pgfqpoint{1.852427in}{0.418217in}}%
\pgfpathlineto{\pgfqpoint{1.856313in}{0.408047in}}%
\pgfpathlineto{\pgfqpoint{1.857651in}{0.393438in}}%
\pgfpathlineto{\pgfqpoint{1.861148in}{0.393352in}}%
\pgfpathlineto{\pgfqpoint{1.859982in}{0.371433in}}%
\pgfpathlineto{\pgfqpoint{1.821677in}{0.373163in}}%
\pgfpathlineto{\pgfqpoint{1.817184in}{0.364409in}}%
\pgfpathlineto{\pgfqpoint{1.805378in}{0.356576in}}%
\pgfpathclose%
\pgfusepath{fill}%
\end{pgfscope}%
\begin{pgfscope}%
\pgfpathrectangle{\pgfqpoint{0.100000in}{0.100000in}}{\pgfqpoint{3.608454in}{2.310000in}}%
\pgfusepath{clip}%
\pgfsetbuttcap%
\pgfsetmiterjoin%
\definecolor{currentfill}{rgb}{0.000000,0.596078,0.701961}%
\pgfsetfillcolor{currentfill}%
\pgfsetlinewidth{0.000000pt}%
\definecolor{currentstroke}{rgb}{0.000000,0.000000,0.000000}%
\pgfsetstrokecolor{currentstroke}%
\pgfsetstrokeopacity{0.000000}%
\pgfsetdash{}{0pt}%
\pgfpathmoveto{\pgfqpoint{2.983960in}{0.768316in}}%
\pgfpathlineto{\pgfqpoint{2.977965in}{0.773429in}}%
\pgfpathlineto{\pgfqpoint{2.977554in}{0.783183in}}%
\pgfpathlineto{\pgfqpoint{2.983748in}{0.789742in}}%
\pgfpathlineto{\pgfqpoint{2.990030in}{0.793119in}}%
\pgfpathlineto{\pgfqpoint{2.986023in}{0.797823in}}%
\pgfpathlineto{\pgfqpoint{2.994265in}{0.796799in}}%
\pgfpathlineto{\pgfqpoint{3.002283in}{0.799408in}}%
\pgfpathlineto{\pgfqpoint{3.007825in}{0.797264in}}%
\pgfpathlineto{\pgfqpoint{3.013726in}{0.809687in}}%
\pgfpathlineto{\pgfqpoint{3.009224in}{0.814410in}}%
\pgfpathlineto{\pgfqpoint{3.015577in}{0.818559in}}%
\pgfpathlineto{\pgfqpoint{3.019148in}{0.828947in}}%
\pgfpathlineto{\pgfqpoint{3.023851in}{0.824906in}}%
\pgfpathlineto{\pgfqpoint{3.035156in}{0.827393in}}%
\pgfpathlineto{\pgfqpoint{3.040729in}{0.821509in}}%
\pgfpathlineto{\pgfqpoint{3.042224in}{0.818383in}}%
\pgfpathlineto{\pgfqpoint{3.038091in}{0.806725in}}%
\pgfpathlineto{\pgfqpoint{3.040053in}{0.796620in}}%
\pgfpathlineto{\pgfqpoint{3.031881in}{0.783770in}}%
\pgfpathlineto{\pgfqpoint{3.030760in}{0.775066in}}%
\pgfpathlineto{\pgfqpoint{3.020788in}{0.780741in}}%
\pgfpathlineto{\pgfqpoint{3.005604in}{0.782838in}}%
\pgfpathlineto{\pgfqpoint{2.996804in}{0.773153in}}%
\pgfpathlineto{\pgfqpoint{2.988034in}{0.774133in}}%
\pgfpathlineto{\pgfqpoint{2.983960in}{0.768316in}}%
\pgfpathclose%
\pgfusepath{fill}%
\end{pgfscope}%
\begin{pgfscope}%
\pgfpathrectangle{\pgfqpoint{0.100000in}{0.100000in}}{\pgfqpoint{3.608454in}{2.310000in}}%
\pgfusepath{clip}%
\pgfsetbuttcap%
\pgfsetmiterjoin%
\definecolor{currentfill}{rgb}{0.000000,0.364706,0.817647}%
\pgfsetfillcolor{currentfill}%
\pgfsetlinewidth{0.000000pt}%
\definecolor{currentstroke}{rgb}{0.000000,0.000000,0.000000}%
\pgfsetstrokecolor{currentstroke}%
\pgfsetstrokeopacity{0.000000}%
\pgfsetdash{}{0pt}%
\pgfpathmoveto{\pgfqpoint{1.571835in}{1.569055in}}%
\pgfpathlineto{\pgfqpoint{1.580064in}{1.651542in}}%
\pgfpathlineto{\pgfqpoint{1.583107in}{1.682203in}}%
\pgfpathlineto{\pgfqpoint{1.643808in}{1.676503in}}%
\pgfpathlineto{\pgfqpoint{1.655854in}{1.675393in}}%
\pgfpathlineto{\pgfqpoint{1.655483in}{1.658354in}}%
\pgfpathlineto{\pgfqpoint{1.653148in}{1.631106in}}%
\pgfpathlineto{\pgfqpoint{1.654543in}{1.630993in}}%
\pgfpathlineto{\pgfqpoint{1.652617in}{1.603533in}}%
\pgfpathlineto{\pgfqpoint{1.655822in}{1.596346in}}%
\pgfpathlineto{\pgfqpoint{1.654522in}{1.581965in}}%
\pgfpathlineto{\pgfqpoint{1.656181in}{1.573226in}}%
\pgfpathlineto{\pgfqpoint{1.654457in}{1.551499in}}%
\pgfpathlineto{\pgfqpoint{1.611187in}{1.555209in}}%
\pgfpathlineto{\pgfqpoint{1.610795in}{1.551818in}}%
\pgfpathlineto{\pgfqpoint{1.570505in}{1.555595in}}%
\pgfpathlineto{\pgfqpoint{1.571835in}{1.569055in}}%
\pgfpathclose%
\pgfusepath{fill}%
\end{pgfscope}%
\begin{pgfscope}%
\pgfpathrectangle{\pgfqpoint{0.100000in}{0.100000in}}{\pgfqpoint{3.608454in}{2.310000in}}%
\pgfusepath{clip}%
\pgfsetbuttcap%
\pgfsetmiterjoin%
\definecolor{currentfill}{rgb}{0.000000,0.937255,0.531373}%
\pgfsetfillcolor{currentfill}%
\pgfsetlinewidth{0.000000pt}%
\definecolor{currentstroke}{rgb}{0.000000,0.000000,0.000000}%
\pgfsetstrokecolor{currentstroke}%
\pgfsetstrokeopacity{0.000000}%
\pgfsetdash{}{0pt}%
\pgfpathmoveto{\pgfqpoint{2.393376in}{0.840961in}}%
\pgfpathlineto{\pgfqpoint{2.393171in}{0.837522in}}%
\pgfpathlineto{\pgfqpoint{2.384135in}{0.837054in}}%
\pgfpathlineto{\pgfqpoint{2.384564in}{0.826478in}}%
\pgfpathlineto{\pgfqpoint{2.378401in}{0.826111in}}%
\pgfpathlineto{\pgfqpoint{2.377206in}{0.836777in}}%
\pgfpathlineto{\pgfqpoint{2.371879in}{0.836595in}}%
\pgfpathlineto{\pgfqpoint{2.366196in}{0.837164in}}%
\pgfpathlineto{\pgfqpoint{2.364353in}{0.841406in}}%
\pgfpathlineto{\pgfqpoint{2.370206in}{0.848561in}}%
\pgfpathlineto{\pgfqpoint{2.363126in}{0.851002in}}%
\pgfpathlineto{\pgfqpoint{2.363366in}{0.855813in}}%
\pgfpathlineto{\pgfqpoint{2.368994in}{0.861333in}}%
\pgfpathlineto{\pgfqpoint{2.367614in}{0.868244in}}%
\pgfpathlineto{\pgfqpoint{2.362240in}{0.870230in}}%
\pgfpathlineto{\pgfqpoint{2.366771in}{0.881093in}}%
\pgfpathlineto{\pgfqpoint{2.366452in}{0.887645in}}%
\pgfpathlineto{\pgfqpoint{2.362887in}{0.893699in}}%
\pgfpathlineto{\pgfqpoint{2.362531in}{0.900120in}}%
\pgfpathlineto{\pgfqpoint{2.356517in}{0.914773in}}%
\pgfpathlineto{\pgfqpoint{2.362379in}{0.917175in}}%
\pgfpathlineto{\pgfqpoint{2.356747in}{0.922152in}}%
\pgfpathlineto{\pgfqpoint{2.361535in}{0.931873in}}%
\pgfpathlineto{\pgfqpoint{2.370618in}{0.931411in}}%
\pgfpathlineto{\pgfqpoint{2.365618in}{0.938943in}}%
\pgfpathlineto{\pgfqpoint{2.369367in}{0.943780in}}%
\pgfpathlineto{\pgfqpoint{2.364126in}{0.947506in}}%
\pgfpathlineto{\pgfqpoint{2.376340in}{0.951894in}}%
\pgfpathlineto{\pgfqpoint{2.377959in}{0.956142in}}%
\pgfpathlineto{\pgfqpoint{2.372346in}{0.959262in}}%
\pgfpathlineto{\pgfqpoint{2.372289in}{0.959362in}}%
\pgfpathlineto{\pgfqpoint{2.391745in}{0.960160in}}%
\pgfpathlineto{\pgfqpoint{2.392184in}{0.949754in}}%
\pgfpathlineto{\pgfqpoint{2.405968in}{0.950277in}}%
\pgfpathlineto{\pgfqpoint{2.406598in}{0.936460in}}%
\pgfpathlineto{\pgfqpoint{2.408409in}{0.898240in}}%
\pgfpathlineto{\pgfqpoint{2.401625in}{0.897879in}}%
\pgfpathlineto{\pgfqpoint{2.401868in}{0.893227in}}%
\pgfpathlineto{\pgfqpoint{2.391193in}{0.892674in}}%
\pgfpathlineto{\pgfqpoint{2.395715in}{0.879087in}}%
\pgfpathlineto{\pgfqpoint{2.396470in}{0.865278in}}%
\pgfpathlineto{\pgfqpoint{2.389913in}{0.858059in}}%
\pgfpathlineto{\pgfqpoint{2.393315in}{0.855864in}}%
\pgfpathlineto{\pgfqpoint{2.393376in}{0.840961in}}%
\pgfpathclose%
\pgfusepath{fill}%
\end{pgfscope}%
\begin{pgfscope}%
\pgfpathrectangle{\pgfqpoint{0.100000in}{0.100000in}}{\pgfqpoint{3.608454in}{2.310000in}}%
\pgfusepath{clip}%
\pgfsetbuttcap%
\pgfsetmiterjoin%
\definecolor{currentfill}{rgb}{0.000000,0.733333,0.633333}%
\pgfsetfillcolor{currentfill}%
\pgfsetlinewidth{0.000000pt}%
\definecolor{currentstroke}{rgb}{0.000000,0.000000,0.000000}%
\pgfsetstrokecolor{currentstroke}%
\pgfsetstrokeopacity{0.000000}%
\pgfsetdash{}{0pt}%
\pgfpathmoveto{\pgfqpoint{0.480357in}{2.020001in}}%
\pgfpathlineto{\pgfqpoint{0.492958in}{2.047705in}}%
\pgfpathlineto{\pgfqpoint{0.494929in}{2.056692in}}%
\pgfpathlineto{\pgfqpoint{0.510061in}{2.083946in}}%
\pgfpathlineto{\pgfqpoint{0.514910in}{2.098885in}}%
\pgfpathlineto{\pgfqpoint{0.524396in}{2.121590in}}%
\pgfpathlineto{\pgfqpoint{0.525797in}{2.131179in}}%
\pgfpathlineto{\pgfqpoint{0.538739in}{2.126084in}}%
\pgfpathlineto{\pgfqpoint{0.557819in}{2.120432in}}%
\pgfpathlineto{\pgfqpoint{0.556089in}{2.115055in}}%
\pgfpathlineto{\pgfqpoint{0.551015in}{2.110628in}}%
\pgfpathlineto{\pgfqpoint{0.550017in}{2.103751in}}%
\pgfpathlineto{\pgfqpoint{0.543848in}{2.096029in}}%
\pgfpathlineto{\pgfqpoint{0.538605in}{2.079539in}}%
\pgfpathlineto{\pgfqpoint{0.521499in}{2.085093in}}%
\pgfpathlineto{\pgfqpoint{0.518018in}{2.074545in}}%
\pgfpathlineto{\pgfqpoint{0.521244in}{2.073470in}}%
\pgfpathlineto{\pgfqpoint{0.512881in}{2.048028in}}%
\pgfpathlineto{\pgfqpoint{0.519006in}{2.044536in}}%
\pgfpathlineto{\pgfqpoint{0.512206in}{2.022672in}}%
\pgfpathlineto{\pgfqpoint{0.505570in}{2.024850in}}%
\pgfpathlineto{\pgfqpoint{0.504063in}{2.019289in}}%
\pgfpathlineto{\pgfqpoint{0.498912in}{2.014543in}}%
\pgfpathlineto{\pgfqpoint{0.480357in}{2.020001in}}%
\pgfpathclose%
\pgfusepath{fill}%
\end{pgfscope}%
\begin{pgfscope}%
\pgfpathrectangle{\pgfqpoint{0.100000in}{0.100000in}}{\pgfqpoint{3.608454in}{2.310000in}}%
\pgfusepath{clip}%
\pgfsetbuttcap%
\pgfsetmiterjoin%
\definecolor{currentfill}{rgb}{0.000000,0.827451,0.586275}%
\pgfsetfillcolor{currentfill}%
\pgfsetlinewidth{0.000000pt}%
\definecolor{currentstroke}{rgb}{0.000000,0.000000,0.000000}%
\pgfsetstrokecolor{currentstroke}%
\pgfsetstrokeopacity{0.000000}%
\pgfsetdash{}{0pt}%
\pgfpathmoveto{\pgfqpoint{0.375397in}{0.571894in}}%
\pgfpathlineto{\pgfqpoint{0.374645in}{0.576884in}}%
\pgfpathlineto{\pgfqpoint{0.375793in}{0.577086in}}%
\pgfpathlineto{\pgfqpoint{0.375397in}{0.571894in}}%
\pgfpathclose%
\pgfusepath{fill}%
\end{pgfscope}%
\begin{pgfscope}%
\pgfpathrectangle{\pgfqpoint{0.100000in}{0.100000in}}{\pgfqpoint{3.608454in}{2.310000in}}%
\pgfusepath{clip}%
\pgfsetbuttcap%
\pgfsetmiterjoin%
\definecolor{currentfill}{rgb}{0.000000,0.827451,0.586275}%
\pgfsetfillcolor{currentfill}%
\pgfsetlinewidth{0.000000pt}%
\definecolor{currentstroke}{rgb}{0.000000,0.000000,0.000000}%
\pgfsetstrokecolor{currentstroke}%
\pgfsetstrokeopacity{0.000000}%
\pgfsetdash{}{0pt}%
\pgfpathmoveto{\pgfqpoint{0.409283in}{0.538102in}}%
\pgfpathlineto{\pgfqpoint{0.408841in}{0.542212in}}%
\pgfpathlineto{\pgfqpoint{0.410919in}{0.540658in}}%
\pgfpathlineto{\pgfqpoint{0.412176in}{0.541182in}}%
\pgfpathlineto{\pgfqpoint{0.411443in}{0.546443in}}%
\pgfpathlineto{\pgfqpoint{0.412787in}{0.546879in}}%
\pgfpathlineto{\pgfqpoint{0.415145in}{0.543509in}}%
\pgfpathlineto{\pgfqpoint{0.414661in}{0.542347in}}%
\pgfpathlineto{\pgfqpoint{0.415729in}{0.539263in}}%
\pgfpathlineto{\pgfqpoint{0.410458in}{0.539171in}}%
\pgfpathlineto{\pgfqpoint{0.409283in}{0.538102in}}%
\pgfpathclose%
\pgfusepath{fill}%
\end{pgfscope}%
\begin{pgfscope}%
\pgfpathrectangle{\pgfqpoint{0.100000in}{0.100000in}}{\pgfqpoint{3.608454in}{2.310000in}}%
\pgfusepath{clip}%
\pgfsetbuttcap%
\pgfsetmiterjoin%
\definecolor{currentfill}{rgb}{0.000000,0.827451,0.586275}%
\pgfsetfillcolor{currentfill}%
\pgfsetlinewidth{0.000000pt}%
\definecolor{currentstroke}{rgb}{0.000000,0.000000,0.000000}%
\pgfsetstrokecolor{currentstroke}%
\pgfsetstrokeopacity{0.000000}%
\pgfsetdash{}{0pt}%
\pgfpathmoveto{\pgfqpoint{0.420907in}{0.523072in}}%
\pgfpathlineto{\pgfqpoint{0.420394in}{0.526482in}}%
\pgfpathlineto{\pgfqpoint{0.425789in}{0.526498in}}%
\pgfpathlineto{\pgfqpoint{0.427996in}{0.528757in}}%
\pgfpathlineto{\pgfqpoint{0.429378in}{0.528347in}}%
\pgfpathlineto{\pgfqpoint{0.431337in}{0.526034in}}%
\pgfpathlineto{\pgfqpoint{0.429229in}{0.525606in}}%
\pgfpathlineto{\pgfqpoint{0.428074in}{0.524089in}}%
\pgfpathlineto{\pgfqpoint{0.429993in}{0.521225in}}%
\pgfpathlineto{\pgfqpoint{0.430475in}{0.519141in}}%
\pgfpathlineto{\pgfqpoint{0.428313in}{0.519015in}}%
\pgfpathlineto{\pgfqpoint{0.427037in}{0.520588in}}%
\pgfpathlineto{\pgfqpoint{0.423736in}{0.522918in}}%
\pgfpathlineto{\pgfqpoint{0.420907in}{0.523072in}}%
\pgfpathclose%
\pgfusepath{fill}%
\end{pgfscope}%
\begin{pgfscope}%
\pgfpathrectangle{\pgfqpoint{0.100000in}{0.100000in}}{\pgfqpoint{3.608454in}{2.310000in}}%
\pgfusepath{clip}%
\pgfsetbuttcap%
\pgfsetmiterjoin%
\definecolor{currentfill}{rgb}{0.000000,0.827451,0.586275}%
\pgfsetfillcolor{currentfill}%
\pgfsetlinewidth{0.000000pt}%
\definecolor{currentstroke}{rgb}{0.000000,0.000000,0.000000}%
\pgfsetstrokecolor{currentstroke}%
\pgfsetstrokeopacity{0.000000}%
\pgfsetdash{}{0pt}%
\pgfpathmoveto{\pgfqpoint{0.417557in}{0.530445in}}%
\pgfpathlineto{\pgfqpoint{0.417550in}{0.532867in}}%
\pgfpathlineto{\pgfqpoint{0.412931in}{0.535157in}}%
\pgfpathlineto{\pgfqpoint{0.414886in}{0.536599in}}%
\pgfpathlineto{\pgfqpoint{0.417623in}{0.533713in}}%
\pgfpathlineto{\pgfqpoint{0.421451in}{0.532232in}}%
\pgfpathlineto{\pgfqpoint{0.423992in}{0.533859in}}%
\pgfpathlineto{\pgfqpoint{0.424555in}{0.531246in}}%
\pgfpathlineto{\pgfqpoint{0.421848in}{0.531121in}}%
\pgfpathlineto{\pgfqpoint{0.419783in}{0.529378in}}%
\pgfpathlineto{\pgfqpoint{0.417557in}{0.530445in}}%
\pgfpathclose%
\pgfusepath{fill}%
\end{pgfscope}%
\begin{pgfscope}%
\pgfpathrectangle{\pgfqpoint{0.100000in}{0.100000in}}{\pgfqpoint{3.608454in}{2.310000in}}%
\pgfusepath{clip}%
\pgfsetbuttcap%
\pgfsetmiterjoin%
\definecolor{currentfill}{rgb}{0.000000,0.827451,0.586275}%
\pgfsetfillcolor{currentfill}%
\pgfsetlinewidth{0.000000pt}%
\definecolor{currentstroke}{rgb}{0.000000,0.000000,0.000000}%
\pgfsetstrokecolor{currentstroke}%
\pgfsetstrokeopacity{0.000000}%
\pgfsetdash{}{0pt}%
\pgfpathmoveto{\pgfqpoint{0.389946in}{0.575682in}}%
\pgfpathlineto{\pgfqpoint{0.389432in}{0.578330in}}%
\pgfpathlineto{\pgfqpoint{0.392641in}{0.578848in}}%
\pgfpathlineto{\pgfqpoint{0.393455in}{0.577648in}}%
\pgfpathlineto{\pgfqpoint{0.392694in}{0.575332in}}%
\pgfpathlineto{\pgfqpoint{0.389946in}{0.575682in}}%
\pgfpathclose%
\pgfusepath{fill}%
\end{pgfscope}%
\begin{pgfscope}%
\pgfpathrectangle{\pgfqpoint{0.100000in}{0.100000in}}{\pgfqpoint{3.608454in}{2.310000in}}%
\pgfusepath{clip}%
\pgfsetbuttcap%
\pgfsetmiterjoin%
\definecolor{currentfill}{rgb}{0.000000,0.827451,0.586275}%
\pgfsetfillcolor{currentfill}%
\pgfsetlinewidth{0.000000pt}%
\definecolor{currentstroke}{rgb}{0.000000,0.000000,0.000000}%
\pgfsetstrokecolor{currentstroke}%
\pgfsetstrokeopacity{0.000000}%
\pgfsetdash{}{0pt}%
\pgfpathmoveto{\pgfqpoint{0.432376in}{0.516958in}}%
\pgfpathlineto{\pgfqpoint{0.431679in}{0.519437in}}%
\pgfpathlineto{\pgfqpoint{0.436217in}{0.517285in}}%
\pgfpathlineto{\pgfqpoint{0.434890in}{0.515977in}}%
\pgfpathlineto{\pgfqpoint{0.432376in}{0.516958in}}%
\pgfpathclose%
\pgfusepath{fill}%
\end{pgfscope}%
\begin{pgfscope}%
\pgfpathrectangle{\pgfqpoint{0.100000in}{0.100000in}}{\pgfqpoint{3.608454in}{2.310000in}}%
\pgfusepath{clip}%
\pgfsetbuttcap%
\pgfsetmiterjoin%
\definecolor{currentfill}{rgb}{0.000000,0.827451,0.586275}%
\pgfsetfillcolor{currentfill}%
\pgfsetlinewidth{0.000000pt}%
\definecolor{currentstroke}{rgb}{0.000000,0.000000,0.000000}%
\pgfsetstrokecolor{currentstroke}%
\pgfsetstrokeopacity{0.000000}%
\pgfsetdash{}{0pt}%
\pgfpathmoveto{\pgfqpoint{0.378939in}{0.591628in}}%
\pgfpathlineto{\pgfqpoint{0.379145in}{0.593528in}}%
\pgfpathlineto{\pgfqpoint{0.381451in}{0.592746in}}%
\pgfpathlineto{\pgfqpoint{0.380887in}{0.590969in}}%
\pgfpathlineto{\pgfqpoint{0.378939in}{0.591628in}}%
\pgfpathclose%
\pgfusepath{fill}%
\end{pgfscope}%
\begin{pgfscope}%
\pgfpathrectangle{\pgfqpoint{0.100000in}{0.100000in}}{\pgfqpoint{3.608454in}{2.310000in}}%
\pgfusepath{clip}%
\pgfsetbuttcap%
\pgfsetmiterjoin%
\definecolor{currentfill}{rgb}{0.000000,0.827451,0.586275}%
\pgfsetfillcolor{currentfill}%
\pgfsetlinewidth{0.000000pt}%
\definecolor{currentstroke}{rgb}{0.000000,0.000000,0.000000}%
\pgfsetstrokecolor{currentstroke}%
\pgfsetstrokeopacity{0.000000}%
\pgfsetdash{}{0pt}%
\pgfpathmoveto{\pgfqpoint{0.436549in}{0.519300in}}%
\pgfpathlineto{\pgfqpoint{0.436374in}{0.521884in}}%
\pgfpathlineto{\pgfqpoint{0.438554in}{0.523113in}}%
\pgfpathlineto{\pgfqpoint{0.439378in}{0.520715in}}%
\pgfpathlineto{\pgfqpoint{0.436549in}{0.519300in}}%
\pgfpathclose%
\pgfusepath{fill}%
\end{pgfscope}%
\begin{pgfscope}%
\pgfpathrectangle{\pgfqpoint{0.100000in}{0.100000in}}{\pgfqpoint{3.608454in}{2.310000in}}%
\pgfusepath{clip}%
\pgfsetbuttcap%
\pgfsetmiterjoin%
\definecolor{currentfill}{rgb}{0.000000,0.827451,0.586275}%
\pgfsetfillcolor{currentfill}%
\pgfsetlinewidth{0.000000pt}%
\definecolor{currentstroke}{rgb}{0.000000,0.000000,0.000000}%
\pgfsetstrokecolor{currentstroke}%
\pgfsetstrokeopacity{0.000000}%
\pgfsetdash{}{0pt}%
\pgfpathmoveto{\pgfqpoint{0.365932in}{0.607246in}}%
\pgfpathlineto{\pgfqpoint{0.366004in}{0.609622in}}%
\pgfpathlineto{\pgfqpoint{0.369338in}{0.609095in}}%
\pgfpathlineto{\pgfqpoint{0.371550in}{0.607762in}}%
\pgfpathlineto{\pgfqpoint{0.374186in}{0.608460in}}%
\pgfpathlineto{\pgfqpoint{0.375594in}{0.606756in}}%
\pgfpathlineto{\pgfqpoint{0.371710in}{0.606454in}}%
\pgfpathlineto{\pgfqpoint{0.370752in}{0.604856in}}%
\pgfpathlineto{\pgfqpoint{0.368700in}{0.607913in}}%
\pgfpathlineto{\pgfqpoint{0.365932in}{0.607246in}}%
\pgfpathclose%
\pgfusepath{fill}%
\end{pgfscope}%
\begin{pgfscope}%
\pgfpathrectangle{\pgfqpoint{0.100000in}{0.100000in}}{\pgfqpoint{3.608454in}{2.310000in}}%
\pgfusepath{clip}%
\pgfsetbuttcap%
\pgfsetmiterjoin%
\definecolor{currentfill}{rgb}{0.000000,0.827451,0.586275}%
\pgfsetfillcolor{currentfill}%
\pgfsetlinewidth{0.000000pt}%
\definecolor{currentstroke}{rgb}{0.000000,0.000000,0.000000}%
\pgfsetstrokecolor{currentstroke}%
\pgfsetstrokeopacity{0.000000}%
\pgfsetdash{}{0pt}%
\pgfpathmoveto{\pgfqpoint{0.480883in}{0.478483in}}%
\pgfpathlineto{\pgfqpoint{0.480302in}{0.480201in}}%
\pgfpathlineto{\pgfqpoint{0.484465in}{0.480081in}}%
\pgfpathlineto{\pgfqpoint{0.485718in}{0.477885in}}%
\pgfpathlineto{\pgfqpoint{0.484820in}{0.476951in}}%
\pgfpathlineto{\pgfqpoint{0.480883in}{0.478483in}}%
\pgfpathclose%
\pgfusepath{fill}%
\end{pgfscope}%
\begin{pgfscope}%
\pgfpathrectangle{\pgfqpoint{0.100000in}{0.100000in}}{\pgfqpoint{3.608454in}{2.310000in}}%
\pgfusepath{clip}%
\pgfsetbuttcap%
\pgfsetmiterjoin%
\definecolor{currentfill}{rgb}{0.000000,0.827451,0.586275}%
\pgfsetfillcolor{currentfill}%
\pgfsetlinewidth{0.000000pt}%
\definecolor{currentstroke}{rgb}{0.000000,0.000000,0.000000}%
\pgfsetstrokecolor{currentstroke}%
\pgfsetstrokeopacity{0.000000}%
\pgfsetdash{}{0pt}%
\pgfpathmoveto{\pgfqpoint{0.344293in}{0.666210in}}%
\pgfpathlineto{\pgfqpoint{0.343816in}{0.669244in}}%
\pgfpathlineto{\pgfqpoint{0.344625in}{0.670409in}}%
\pgfpathlineto{\pgfqpoint{0.346843in}{0.669814in}}%
\pgfpathlineto{\pgfqpoint{0.346503in}{0.668280in}}%
\pgfpathlineto{\pgfqpoint{0.344293in}{0.666210in}}%
\pgfpathclose%
\pgfusepath{fill}%
\end{pgfscope}%
\begin{pgfscope}%
\pgfpathrectangle{\pgfqpoint{0.100000in}{0.100000in}}{\pgfqpoint{3.608454in}{2.310000in}}%
\pgfusepath{clip}%
\pgfsetbuttcap%
\pgfsetmiterjoin%
\definecolor{currentfill}{rgb}{0.000000,0.827451,0.586275}%
\pgfsetfillcolor{currentfill}%
\pgfsetlinewidth{0.000000pt}%
\definecolor{currentstroke}{rgb}{0.000000,0.000000,0.000000}%
\pgfsetstrokecolor{currentstroke}%
\pgfsetstrokeopacity{0.000000}%
\pgfsetdash{}{0pt}%
\pgfpathmoveto{\pgfqpoint{0.506806in}{0.461794in}}%
\pgfpathlineto{\pgfqpoint{0.507677in}{0.463806in}}%
\pgfpathlineto{\pgfqpoint{0.511473in}{0.462938in}}%
\pgfpathlineto{\pgfqpoint{0.510522in}{0.460924in}}%
\pgfpathlineto{\pgfqpoint{0.506806in}{0.461794in}}%
\pgfpathclose%
\pgfusepath{fill}%
\end{pgfscope}%
\begin{pgfscope}%
\pgfpathrectangle{\pgfqpoint{0.100000in}{0.100000in}}{\pgfqpoint{3.608454in}{2.310000in}}%
\pgfusepath{clip}%
\pgfsetbuttcap%
\pgfsetmiterjoin%
\definecolor{currentfill}{rgb}{0.000000,0.827451,0.586275}%
\pgfsetfillcolor{currentfill}%
\pgfsetlinewidth{0.000000pt}%
\definecolor{currentstroke}{rgb}{0.000000,0.000000,0.000000}%
\pgfsetstrokecolor{currentstroke}%
\pgfsetstrokeopacity{0.000000}%
\pgfsetdash{}{0pt}%
\pgfpathmoveto{\pgfqpoint{0.349460in}{0.680635in}}%
\pgfpathlineto{\pgfqpoint{0.349236in}{0.682931in}}%
\pgfpathlineto{\pgfqpoint{0.345635in}{0.686541in}}%
\pgfpathlineto{\pgfqpoint{0.348354in}{0.687196in}}%
\pgfpathlineto{\pgfqpoint{0.346567in}{0.689336in}}%
\pgfpathlineto{\pgfqpoint{0.347637in}{0.690563in}}%
\pgfpathlineto{\pgfqpoint{0.350022in}{0.690962in}}%
\pgfpathlineto{\pgfqpoint{0.351654in}{0.687198in}}%
\pgfpathlineto{\pgfqpoint{0.353748in}{0.683821in}}%
\pgfpathlineto{\pgfqpoint{0.353994in}{0.680863in}}%
\pgfpathlineto{\pgfqpoint{0.352204in}{0.679863in}}%
\pgfpathlineto{\pgfqpoint{0.349460in}{0.680635in}}%
\pgfpathclose%
\pgfusepath{fill}%
\end{pgfscope}%
\begin{pgfscope}%
\pgfpathrectangle{\pgfqpoint{0.100000in}{0.100000in}}{\pgfqpoint{3.608454in}{2.310000in}}%
\pgfusepath{clip}%
\pgfsetbuttcap%
\pgfsetmiterjoin%
\definecolor{currentfill}{rgb}{0.000000,0.827451,0.586275}%
\pgfsetfillcolor{currentfill}%
\pgfsetlinewidth{0.000000pt}%
\definecolor{currentstroke}{rgb}{0.000000,0.000000,0.000000}%
\pgfsetstrokecolor{currentstroke}%
\pgfsetstrokeopacity{0.000000}%
\pgfsetdash{}{0pt}%
\pgfpathmoveto{\pgfqpoint{0.556547in}{0.436939in}}%
\pgfpathlineto{\pgfqpoint{0.555985in}{0.439788in}}%
\pgfpathlineto{\pgfqpoint{0.557578in}{0.440221in}}%
\pgfpathlineto{\pgfqpoint{0.561322in}{0.439031in}}%
\pgfpathlineto{\pgfqpoint{0.563123in}{0.437521in}}%
\pgfpathlineto{\pgfqpoint{0.565195in}{0.437153in}}%
\pgfpathlineto{\pgfqpoint{0.566888in}{0.435455in}}%
\pgfpathlineto{\pgfqpoint{0.569188in}{0.437256in}}%
\pgfpathlineto{\pgfqpoint{0.572877in}{0.438751in}}%
\pgfpathlineto{\pgfqpoint{0.575932in}{0.435745in}}%
\pgfpathlineto{\pgfqpoint{0.574169in}{0.442233in}}%
\pgfpathlineto{\pgfqpoint{0.576542in}{0.443329in}}%
\pgfpathlineto{\pgfqpoint{0.579175in}{0.442765in}}%
\pgfpathlineto{\pgfqpoint{0.583426in}{0.440455in}}%
\pgfpathlineto{\pgfqpoint{0.582448in}{0.438627in}}%
\pgfpathlineto{\pgfqpoint{0.584561in}{0.436213in}}%
\pgfpathlineto{\pgfqpoint{0.586754in}{0.436810in}}%
\pgfpathlineto{\pgfqpoint{0.586008in}{0.433801in}}%
\pgfpathlineto{\pgfqpoint{0.582411in}{0.433473in}}%
\pgfpathlineto{\pgfqpoint{0.580948in}{0.429661in}}%
\pgfpathlineto{\pgfqpoint{0.574828in}{0.431855in}}%
\pgfpathlineto{\pgfqpoint{0.572584in}{0.432073in}}%
\pgfpathlineto{\pgfqpoint{0.571112in}{0.430136in}}%
\pgfpathlineto{\pgfqpoint{0.569377in}{0.433000in}}%
\pgfpathlineto{\pgfqpoint{0.567679in}{0.433841in}}%
\pgfpathlineto{\pgfqpoint{0.562725in}{0.434707in}}%
\pgfpathlineto{\pgfqpoint{0.560495in}{0.435926in}}%
\pgfpathlineto{\pgfqpoint{0.556547in}{0.436939in}}%
\pgfpathclose%
\pgfusepath{fill}%
\end{pgfscope}%
\begin{pgfscope}%
\pgfpathrectangle{\pgfqpoint{0.100000in}{0.100000in}}{\pgfqpoint{3.608454in}{2.310000in}}%
\pgfusepath{clip}%
\pgfsetbuttcap%
\pgfsetmiterjoin%
\definecolor{currentfill}{rgb}{0.000000,0.827451,0.586275}%
\pgfsetfillcolor{currentfill}%
\pgfsetlinewidth{0.000000pt}%
\definecolor{currentstroke}{rgb}{0.000000,0.000000,0.000000}%
\pgfsetstrokecolor{currentstroke}%
\pgfsetstrokeopacity{0.000000}%
\pgfsetdash{}{0pt}%
\pgfpathmoveto{\pgfqpoint{0.523468in}{0.452981in}}%
\pgfpathlineto{\pgfqpoint{0.523015in}{0.455703in}}%
\pgfpathlineto{\pgfqpoint{0.525225in}{0.455976in}}%
\pgfpathlineto{\pgfqpoint{0.525280in}{0.453766in}}%
\pgfpathlineto{\pgfqpoint{0.523468in}{0.452981in}}%
\pgfpathclose%
\pgfusepath{fill}%
\end{pgfscope}%
\begin{pgfscope}%
\pgfpathrectangle{\pgfqpoint{0.100000in}{0.100000in}}{\pgfqpoint{3.608454in}{2.310000in}}%
\pgfusepath{clip}%
\pgfsetbuttcap%
\pgfsetmiterjoin%
\definecolor{currentfill}{rgb}{0.000000,0.827451,0.586275}%
\pgfsetfillcolor{currentfill}%
\pgfsetlinewidth{0.000000pt}%
\definecolor{currentstroke}{rgb}{0.000000,0.000000,0.000000}%
\pgfsetstrokecolor{currentstroke}%
\pgfsetstrokeopacity{0.000000}%
\pgfsetdash{}{0pt}%
\pgfpathmoveto{\pgfqpoint{0.532273in}{0.445921in}}%
\pgfpathlineto{\pgfqpoint{0.532818in}{0.446891in}}%
\pgfpathlineto{\pgfqpoint{0.538427in}{0.447049in}}%
\pgfpathlineto{\pgfqpoint{0.541753in}{0.449221in}}%
\pgfpathlineto{\pgfqpoint{0.546018in}{0.448961in}}%
\pgfpathlineto{\pgfqpoint{0.548669in}{0.445595in}}%
\pgfpathlineto{\pgfqpoint{0.551094in}{0.448774in}}%
\pgfpathlineto{\pgfqpoint{0.553730in}{0.449651in}}%
\pgfpathlineto{\pgfqpoint{0.556658in}{0.449030in}}%
\pgfpathlineto{\pgfqpoint{0.558881in}{0.447475in}}%
\pgfpathlineto{\pgfqpoint{0.559652in}{0.444922in}}%
\pgfpathlineto{\pgfqpoint{0.559201in}{0.443279in}}%
\pgfpathlineto{\pgfqpoint{0.556808in}{0.441953in}}%
\pgfpathlineto{\pgfqpoint{0.548779in}{0.444295in}}%
\pgfpathlineto{\pgfqpoint{0.543555in}{0.442792in}}%
\pgfpathlineto{\pgfqpoint{0.541563in}{0.443699in}}%
\pgfpathlineto{\pgfqpoint{0.532273in}{0.445921in}}%
\pgfpathclose%
\pgfusepath{fill}%
\end{pgfscope}%
\begin{pgfscope}%
\pgfpathrectangle{\pgfqpoint{0.100000in}{0.100000in}}{\pgfqpoint{3.608454in}{2.310000in}}%
\pgfusepath{clip}%
\pgfsetbuttcap%
\pgfsetmiterjoin%
\definecolor{currentfill}{rgb}{0.000000,0.827451,0.586275}%
\pgfsetfillcolor{currentfill}%
\pgfsetlinewidth{0.000000pt}%
\definecolor{currentstroke}{rgb}{0.000000,0.000000,0.000000}%
\pgfsetstrokecolor{currentstroke}%
\pgfsetstrokeopacity{0.000000}%
\pgfsetdash{}{0pt}%
\pgfpathmoveto{\pgfqpoint{0.467157in}{0.485560in}}%
\pgfpathlineto{\pgfqpoint{0.465100in}{0.487760in}}%
\pgfpathlineto{\pgfqpoint{0.463877in}{0.490340in}}%
\pgfpathlineto{\pgfqpoint{0.461910in}{0.492281in}}%
\pgfpathlineto{\pgfqpoint{0.462381in}{0.494960in}}%
\pgfpathlineto{\pgfqpoint{0.466531in}{0.490940in}}%
\pgfpathlineto{\pgfqpoint{0.468975in}{0.485923in}}%
\pgfpathlineto{\pgfqpoint{0.467157in}{0.485560in}}%
\pgfpathclose%
\pgfusepath{fill}%
\end{pgfscope}%
\begin{pgfscope}%
\pgfpathrectangle{\pgfqpoint{0.100000in}{0.100000in}}{\pgfqpoint{3.608454in}{2.310000in}}%
\pgfusepath{clip}%
\pgfsetbuttcap%
\pgfsetmiterjoin%
\definecolor{currentfill}{rgb}{0.000000,0.827451,0.586275}%
\pgfsetfillcolor{currentfill}%
\pgfsetlinewidth{0.000000pt}%
\definecolor{currentstroke}{rgb}{0.000000,0.000000,0.000000}%
\pgfsetstrokecolor{currentstroke}%
\pgfsetstrokeopacity{0.000000}%
\pgfsetdash{}{0pt}%
\pgfpathmoveto{\pgfqpoint{0.456040in}{0.499043in}}%
\pgfpathlineto{\pgfqpoint{0.454673in}{0.501491in}}%
\pgfpathlineto{\pgfqpoint{0.451914in}{0.502371in}}%
\pgfpathlineto{\pgfqpoint{0.451045in}{0.506296in}}%
\pgfpathlineto{\pgfqpoint{0.454103in}{0.504564in}}%
\pgfpathlineto{\pgfqpoint{0.455369in}{0.502279in}}%
\pgfpathlineto{\pgfqpoint{0.457486in}{0.502766in}}%
\pgfpathlineto{\pgfqpoint{0.461584in}{0.502365in}}%
\pgfpathlineto{\pgfqpoint{0.463798in}{0.503703in}}%
\pgfpathlineto{\pgfqpoint{0.465884in}{0.502834in}}%
\pgfpathlineto{\pgfqpoint{0.466975in}{0.500406in}}%
\pgfpathlineto{\pgfqpoint{0.466593in}{0.498917in}}%
\pgfpathlineto{\pgfqpoint{0.464307in}{0.497731in}}%
\pgfpathlineto{\pgfqpoint{0.462347in}{0.499412in}}%
\pgfpathlineto{\pgfqpoint{0.461307in}{0.496536in}}%
\pgfpathlineto{\pgfqpoint{0.459853in}{0.497056in}}%
\pgfpathlineto{\pgfqpoint{0.457626in}{0.499895in}}%
\pgfpathlineto{\pgfqpoint{0.456040in}{0.499043in}}%
\pgfpathclose%
\pgfusepath{fill}%
\end{pgfscope}%
\begin{pgfscope}%
\pgfpathrectangle{\pgfqpoint{0.100000in}{0.100000in}}{\pgfqpoint{3.608454in}{2.310000in}}%
\pgfusepath{clip}%
\pgfsetbuttcap%
\pgfsetmiterjoin%
\definecolor{currentfill}{rgb}{0.000000,0.776471,0.611765}%
\pgfsetfillcolor{currentfill}%
\pgfsetlinewidth{0.000000pt}%
\definecolor{currentstroke}{rgb}{0.000000,0.000000,0.000000}%
\pgfsetstrokecolor{currentstroke}%
\pgfsetstrokeopacity{0.000000}%
\pgfsetdash{}{0pt}%
\pgfpathmoveto{\pgfqpoint{2.231585in}{0.901788in}}%
\pgfpathlineto{\pgfqpoint{2.224654in}{0.900658in}}%
\pgfpathlineto{\pgfqpoint{2.207240in}{0.900881in}}%
\pgfpathlineto{\pgfqpoint{2.207261in}{0.903173in}}%
\pgfpathlineto{\pgfqpoint{2.191374in}{0.903396in}}%
\pgfpathlineto{\pgfqpoint{2.183884in}{0.913246in}}%
\pgfpathlineto{\pgfqpoint{2.175110in}{0.918121in}}%
\pgfpathlineto{\pgfqpoint{2.175503in}{0.924555in}}%
\pgfpathlineto{\pgfqpoint{2.174128in}{0.929756in}}%
\pgfpathlineto{\pgfqpoint{2.169177in}{0.935220in}}%
\pgfpathlineto{\pgfqpoint{2.166804in}{0.946907in}}%
\pgfpathlineto{\pgfqpoint{2.165791in}{0.959401in}}%
\pgfpathlineto{\pgfqpoint{2.156410in}{0.959401in}}%
\pgfpathlineto{\pgfqpoint{2.155587in}{0.972321in}}%
\pgfpathlineto{\pgfqpoint{2.176495in}{0.972103in}}%
\pgfpathlineto{\pgfqpoint{2.176519in}{0.997121in}}%
\pgfpathlineto{\pgfqpoint{2.180404in}{1.000901in}}%
\pgfpathlineto{\pgfqpoint{2.190786in}{1.003624in}}%
\pgfpathlineto{\pgfqpoint{2.211366in}{1.003801in}}%
\pgfpathlineto{\pgfqpoint{2.218257in}{1.006289in}}%
\pgfpathlineto{\pgfqpoint{2.238992in}{1.006643in}}%
\pgfpathlineto{\pgfqpoint{2.238537in}{0.995926in}}%
\pgfpathlineto{\pgfqpoint{2.245351in}{0.995915in}}%
\pgfpathlineto{\pgfqpoint{2.249833in}{0.992475in}}%
\pgfpathlineto{\pgfqpoint{2.250959in}{0.985619in}}%
\pgfpathlineto{\pgfqpoint{2.258094in}{0.983370in}}%
\pgfpathlineto{\pgfqpoint{2.259281in}{0.978738in}}%
\pgfpathlineto{\pgfqpoint{2.259184in}{0.957741in}}%
\pgfpathlineto{\pgfqpoint{2.245528in}{0.957904in}}%
\pgfpathlineto{\pgfqpoint{2.245189in}{0.934157in}}%
\pgfpathlineto{\pgfqpoint{2.241965in}{0.929323in}}%
\pgfpathlineto{\pgfqpoint{2.231849in}{0.927551in}}%
\pgfpathlineto{\pgfqpoint{2.231585in}{0.901788in}}%
\pgfpathclose%
\pgfusepath{fill}%
\end{pgfscope}%
\begin{pgfscope}%
\pgfpathrectangle{\pgfqpoint{0.100000in}{0.100000in}}{\pgfqpoint{3.608454in}{2.310000in}}%
\pgfusepath{clip}%
\pgfsetbuttcap%
\pgfsetmiterjoin%
\definecolor{currentfill}{rgb}{0.000000,0.721569,0.639216}%
\pgfsetfillcolor{currentfill}%
\pgfsetlinewidth{0.000000pt}%
\definecolor{currentstroke}{rgb}{0.000000,0.000000,0.000000}%
\pgfsetstrokecolor{currentstroke}%
\pgfsetstrokeopacity{0.000000}%
\pgfsetdash{}{0pt}%
\pgfpathmoveto{\pgfqpoint{0.695564in}{1.998853in}}%
\pgfpathlineto{\pgfqpoint{0.656754in}{2.010285in}}%
\pgfpathlineto{\pgfqpoint{0.620648in}{2.021281in}}%
\pgfpathlineto{\pgfqpoint{0.623427in}{2.025636in}}%
\pgfpathlineto{\pgfqpoint{0.622394in}{2.035924in}}%
\pgfpathlineto{\pgfqpoint{0.630035in}{2.038493in}}%
\pgfpathlineto{\pgfqpoint{0.629691in}{2.048489in}}%
\pgfpathlineto{\pgfqpoint{0.634662in}{2.052316in}}%
\pgfpathlineto{\pgfqpoint{0.636428in}{2.062774in}}%
\pgfpathlineto{\pgfqpoint{0.628826in}{2.076535in}}%
\pgfpathlineto{\pgfqpoint{0.630988in}{2.086576in}}%
\pgfpathlineto{\pgfqpoint{0.634943in}{2.089043in}}%
\pgfpathlineto{\pgfqpoint{0.641984in}{2.086967in}}%
\pgfpathlineto{\pgfqpoint{0.653479in}{2.086201in}}%
\pgfpathlineto{\pgfqpoint{0.650732in}{2.091686in}}%
\pgfpathlineto{\pgfqpoint{0.656547in}{2.111660in}}%
\pgfpathlineto{\pgfqpoint{0.661241in}{2.110265in}}%
\pgfpathlineto{\pgfqpoint{0.719253in}{2.092888in}}%
\pgfpathlineto{\pgfqpoint{0.749167in}{2.084439in}}%
\pgfpathlineto{\pgfqpoint{0.744663in}{2.068773in}}%
\pgfpathlineto{\pgfqpoint{0.737141in}{2.068948in}}%
\pgfpathlineto{\pgfqpoint{0.723978in}{2.065521in}}%
\pgfpathlineto{\pgfqpoint{0.713071in}{2.066429in}}%
\pgfpathlineto{\pgfqpoint{0.707552in}{2.064386in}}%
\pgfpathlineto{\pgfqpoint{0.710428in}{2.058581in}}%
\pgfpathlineto{\pgfqpoint{0.710674in}{2.050338in}}%
\pgfpathlineto{\pgfqpoint{0.702993in}{2.049946in}}%
\pgfpathlineto{\pgfqpoint{0.695191in}{2.034615in}}%
\pgfpathlineto{\pgfqpoint{0.697834in}{2.026472in}}%
\pgfpathlineto{\pgfqpoint{0.692055in}{2.010580in}}%
\pgfpathlineto{\pgfqpoint{0.692032in}{2.003466in}}%
\pgfpathlineto{\pgfqpoint{0.695564in}{1.998853in}}%
\pgfpathclose%
\pgfusepath{fill}%
\end{pgfscope}%
\begin{pgfscope}%
\pgfpathrectangle{\pgfqpoint{0.100000in}{0.100000in}}{\pgfqpoint{3.608454in}{2.310000in}}%
\pgfusepath{clip}%
\pgfsetbuttcap%
\pgfsetmiterjoin%
\definecolor{currentfill}{rgb}{0.000000,0.505882,0.747059}%
\pgfsetfillcolor{currentfill}%
\pgfsetlinewidth{0.000000pt}%
\definecolor{currentstroke}{rgb}{0.000000,0.000000,0.000000}%
\pgfsetstrokecolor{currentstroke}%
\pgfsetstrokeopacity{0.000000}%
\pgfsetdash{}{0pt}%
\pgfpathmoveto{\pgfqpoint{1.748391in}{1.048651in}}%
\pgfpathlineto{\pgfqpoint{1.713023in}{1.050837in}}%
\pgfpathlineto{\pgfqpoint{1.678367in}{1.053341in}}%
\pgfpathlineto{\pgfqpoint{1.680896in}{1.087866in}}%
\pgfpathlineto{\pgfqpoint{1.683362in}{1.122232in}}%
\pgfpathlineto{\pgfqpoint{1.752524in}{1.117682in}}%
\pgfpathlineto{\pgfqpoint{1.787003in}{1.115616in}}%
\pgfpathlineto{\pgfqpoint{1.786242in}{1.101806in}}%
\pgfpathlineto{\pgfqpoint{1.785102in}{1.081107in}}%
\pgfpathlineto{\pgfqpoint{1.750416in}{1.083136in}}%
\pgfpathlineto{\pgfqpoint{1.748391in}{1.048651in}}%
\pgfpathclose%
\pgfusepath{fill}%
\end{pgfscope}%
\begin{pgfscope}%
\pgfpathrectangle{\pgfqpoint{0.100000in}{0.100000in}}{\pgfqpoint{3.608454in}{2.310000in}}%
\pgfusepath{clip}%
\pgfsetbuttcap%
\pgfsetmiterjoin%
\definecolor{currentfill}{rgb}{0.000000,0.654902,0.672549}%
\pgfsetfillcolor{currentfill}%
\pgfsetlinewidth{0.000000pt}%
\definecolor{currentstroke}{rgb}{0.000000,0.000000,0.000000}%
\pgfsetstrokecolor{currentstroke}%
\pgfsetstrokeopacity{0.000000}%
\pgfsetdash{}{0pt}%
\pgfpathmoveto{\pgfqpoint{1.884281in}{0.742465in}}%
\pgfpathlineto{\pgfqpoint{1.863869in}{0.728404in}}%
\pgfpathlineto{\pgfqpoint{1.866023in}{0.717001in}}%
\pgfpathlineto{\pgfqpoint{1.871649in}{0.712052in}}%
\pgfpathlineto{\pgfqpoint{1.870905in}{0.703446in}}%
\pgfpathlineto{\pgfqpoint{1.835408in}{0.704889in}}%
\pgfpathlineto{\pgfqpoint{1.826743in}{0.706842in}}%
\pgfpathlineto{\pgfqpoint{1.828725in}{0.748261in}}%
\pgfpathlineto{\pgfqpoint{1.818920in}{0.750454in}}%
\pgfpathlineto{\pgfqpoint{1.813328in}{0.744891in}}%
\pgfpathlineto{\pgfqpoint{1.804967in}{0.750514in}}%
\pgfpathlineto{\pgfqpoint{1.792433in}{0.750538in}}%
\pgfpathlineto{\pgfqpoint{1.786260in}{0.758457in}}%
\pgfpathlineto{\pgfqpoint{1.788954in}{0.799743in}}%
\pgfpathlineto{\pgfqpoint{1.842140in}{0.796889in}}%
\pgfpathlineto{\pgfqpoint{1.858178in}{0.766124in}}%
\pgfpathlineto{\pgfqpoint{1.870147in}{0.767494in}}%
\pgfpathlineto{\pgfqpoint{1.884281in}{0.742465in}}%
\pgfpathclose%
\pgfusepath{fill}%
\end{pgfscope}%
\begin{pgfscope}%
\pgfpathrectangle{\pgfqpoint{0.100000in}{0.100000in}}{\pgfqpoint{3.608454in}{2.310000in}}%
\pgfusepath{clip}%
\pgfsetbuttcap%
\pgfsetmiterjoin%
\definecolor{currentfill}{rgb}{0.000000,0.905882,0.547059}%
\pgfsetfillcolor{currentfill}%
\pgfsetlinewidth{0.000000pt}%
\definecolor{currentstroke}{rgb}{0.000000,0.000000,0.000000}%
\pgfsetstrokecolor{currentstroke}%
\pgfsetstrokeopacity{0.000000}%
\pgfsetdash{}{0pt}%
\pgfpathmoveto{\pgfqpoint{2.910647in}{0.649877in}}%
\pgfpathlineto{\pgfqpoint{2.900528in}{0.654212in}}%
\pgfpathlineto{\pgfqpoint{2.893044in}{0.665397in}}%
\pgfpathlineto{\pgfqpoint{2.882864in}{0.671272in}}%
\pgfpathlineto{\pgfqpoint{2.871699in}{0.674920in}}%
\pgfpathlineto{\pgfqpoint{2.866769in}{0.678489in}}%
\pgfpathlineto{\pgfqpoint{2.870453in}{0.691569in}}%
\pgfpathlineto{\pgfqpoint{2.878191in}{0.698489in}}%
\pgfpathlineto{\pgfqpoint{2.876684in}{0.701966in}}%
\pgfpathlineto{\pgfqpoint{2.881876in}{0.708207in}}%
\pgfpathlineto{\pgfqpoint{2.880133in}{0.713968in}}%
\pgfpathlineto{\pgfqpoint{2.888448in}{0.720934in}}%
\pgfpathlineto{\pgfqpoint{2.887245in}{0.726366in}}%
\pgfpathlineto{\pgfqpoint{2.957900in}{0.730883in}}%
\pgfpathlineto{\pgfqpoint{2.966459in}{0.731456in}}%
\pgfpathlineto{\pgfqpoint{2.971461in}{0.696012in}}%
\pgfpathlineto{\pgfqpoint{2.965631in}{0.682186in}}%
\pgfpathlineto{\pgfqpoint{2.965236in}{0.672986in}}%
\pgfpathlineto{\pgfqpoint{2.962277in}{0.669805in}}%
\pgfpathlineto{\pgfqpoint{2.955684in}{0.671251in}}%
\pgfpathlineto{\pgfqpoint{2.951329in}{0.676818in}}%
\pgfpathlineto{\pgfqpoint{2.946140in}{0.674377in}}%
\pgfpathlineto{\pgfqpoint{2.944970in}{0.667447in}}%
\pgfpathlineto{\pgfqpoint{2.915694in}{0.663204in}}%
\pgfpathlineto{\pgfqpoint{2.913541in}{0.650279in}}%
\pgfpathlineto{\pgfqpoint{2.910647in}{0.649877in}}%
\pgfpathclose%
\pgfusepath{fill}%
\end{pgfscope}%
\begin{pgfscope}%
\pgfpathrectangle{\pgfqpoint{0.100000in}{0.100000in}}{\pgfqpoint{3.608454in}{2.310000in}}%
\pgfusepath{clip}%
\pgfsetbuttcap%
\pgfsetmiterjoin%
\definecolor{currentfill}{rgb}{0.000000,0.686275,0.656863}%
\pgfsetfillcolor{currentfill}%
\pgfsetlinewidth{0.000000pt}%
\definecolor{currentstroke}{rgb}{0.000000,0.000000,0.000000}%
\pgfsetstrokecolor{currentstroke}%
\pgfsetstrokeopacity{0.000000}%
\pgfsetdash{}{0pt}%
\pgfpathmoveto{\pgfqpoint{2.543041in}{1.178054in}}%
\pgfpathlineto{\pgfqpoint{2.549236in}{1.162164in}}%
\pgfpathlineto{\pgfqpoint{2.548186in}{1.158629in}}%
\pgfpathlineto{\pgfqpoint{2.520509in}{1.156888in}}%
\pgfpathlineto{\pgfqpoint{2.519321in}{1.176572in}}%
\pgfpathlineto{\pgfqpoint{2.543041in}{1.178054in}}%
\pgfpathclose%
\pgfusepath{fill}%
\end{pgfscope}%
\begin{pgfscope}%
\pgfpathrectangle{\pgfqpoint{0.100000in}{0.100000in}}{\pgfqpoint{3.608454in}{2.310000in}}%
\pgfusepath{clip}%
\pgfsetbuttcap%
\pgfsetmiterjoin%
\definecolor{currentfill}{rgb}{0.000000,0.729412,0.635294}%
\pgfsetfillcolor{currentfill}%
\pgfsetlinewidth{0.000000pt}%
\definecolor{currentstroke}{rgb}{0.000000,0.000000,0.000000}%
\pgfsetstrokecolor{currentstroke}%
\pgfsetstrokeopacity{0.000000}%
\pgfsetdash{}{0pt}%
\pgfpathmoveto{\pgfqpoint{2.213676in}{1.928846in}}%
\pgfpathlineto{\pgfqpoint{2.214175in}{1.908040in}}%
\pgfpathlineto{\pgfqpoint{2.193329in}{1.907307in}}%
\pgfpathlineto{\pgfqpoint{2.193274in}{1.914649in}}%
\pgfpathlineto{\pgfqpoint{2.172320in}{1.914091in}}%
\pgfpathlineto{\pgfqpoint{2.172440in}{1.907196in}}%
\pgfpathlineto{\pgfqpoint{2.141306in}{1.906925in}}%
\pgfpathlineto{\pgfqpoint{2.142555in}{1.917955in}}%
\pgfpathlineto{\pgfqpoint{2.138352in}{1.920870in}}%
\pgfpathlineto{\pgfqpoint{2.132410in}{1.918482in}}%
\pgfpathlineto{\pgfqpoint{2.118937in}{1.925718in}}%
\pgfpathlineto{\pgfqpoint{2.118427in}{1.958403in}}%
\pgfpathlineto{\pgfqpoint{2.125367in}{1.958378in}}%
\pgfpathlineto{\pgfqpoint{2.124628in}{2.006820in}}%
\pgfpathlineto{\pgfqpoint{2.138354in}{2.007024in}}%
\pgfpathlineto{\pgfqpoint{2.138015in}{2.041726in}}%
\pgfpathlineto{\pgfqpoint{2.172503in}{2.042117in}}%
\pgfpathlineto{\pgfqpoint{2.172471in}{2.046322in}}%
\pgfpathlineto{\pgfqpoint{2.209745in}{2.046401in}}%
\pgfpathlineto{\pgfqpoint{2.211106in}{2.032691in}}%
\pgfpathlineto{\pgfqpoint{2.211215in}{2.006566in}}%
\pgfpathlineto{\pgfqpoint{2.212580in}{1.977269in}}%
\pgfpathlineto{\pgfqpoint{2.212805in}{1.945663in}}%
\pgfpathlineto{\pgfqpoint{2.213676in}{1.928846in}}%
\pgfpathclose%
\pgfusepath{fill}%
\end{pgfscope}%
\begin{pgfscope}%
\pgfpathrectangle{\pgfqpoint{0.100000in}{0.100000in}}{\pgfqpoint{3.608454in}{2.310000in}}%
\pgfusepath{clip}%
\pgfsetbuttcap%
\pgfsetmiterjoin%
\definecolor{currentfill}{rgb}{0.000000,0.521569,0.739216}%
\pgfsetfillcolor{currentfill}%
\pgfsetlinewidth{0.000000pt}%
\definecolor{currentstroke}{rgb}{0.000000,0.000000,0.000000}%
\pgfsetstrokecolor{currentstroke}%
\pgfsetstrokeopacity{0.000000}%
\pgfsetdash{}{0pt}%
\pgfpathmoveto{\pgfqpoint{1.598608in}{1.839877in}}%
\pgfpathlineto{\pgfqpoint{1.543825in}{1.845887in}}%
\pgfpathlineto{\pgfqpoint{1.546917in}{1.873677in}}%
\pgfpathlineto{\pgfqpoint{1.549754in}{1.873361in}}%
\pgfpathlineto{\pgfqpoint{1.553831in}{1.907730in}}%
\pgfpathlineto{\pgfqpoint{1.556192in}{1.907466in}}%
\pgfpathlineto{\pgfqpoint{1.559222in}{1.935142in}}%
\pgfpathlineto{\pgfqpoint{1.562162in}{1.934829in}}%
\pgfpathlineto{\pgfqpoint{1.582466in}{1.932625in}}%
\pgfpathlineto{\pgfqpoint{1.581763in}{1.926153in}}%
\pgfpathlineto{\pgfqpoint{1.588691in}{1.925458in}}%
\pgfpathlineto{\pgfqpoint{1.587963in}{1.918755in}}%
\pgfpathlineto{\pgfqpoint{1.602057in}{1.917205in}}%
\pgfpathlineto{\pgfqpoint{1.606786in}{1.914838in}}%
\pgfpathlineto{\pgfqpoint{1.664670in}{1.909384in}}%
\pgfpathlineto{\pgfqpoint{1.667609in}{1.909112in}}%
\pgfpathlineto{\pgfqpoint{1.661648in}{1.851140in}}%
\pgfpathlineto{\pgfqpoint{1.601301in}{1.856823in}}%
\pgfpathlineto{\pgfqpoint{1.598608in}{1.839877in}}%
\pgfpathclose%
\pgfusepath{fill}%
\end{pgfscope}%
\begin{pgfscope}%
\pgfpathrectangle{\pgfqpoint{0.100000in}{0.100000in}}{\pgfqpoint{3.608454in}{2.310000in}}%
\pgfusepath{clip}%
\pgfsetbuttcap%
\pgfsetmiterjoin%
\definecolor{currentfill}{rgb}{0.000000,0.768627,0.615686}%
\pgfsetfillcolor{currentfill}%
\pgfsetlinewidth{0.000000pt}%
\definecolor{currentstroke}{rgb}{0.000000,0.000000,0.000000}%
\pgfsetstrokecolor{currentstroke}%
\pgfsetstrokeopacity{0.000000}%
\pgfsetdash{}{0pt}%
\pgfpathmoveto{\pgfqpoint{2.618822in}{1.745974in}}%
\pgfpathlineto{\pgfqpoint{2.596450in}{1.744135in}}%
\pgfpathlineto{\pgfqpoint{2.593518in}{1.755812in}}%
\pgfpathlineto{\pgfqpoint{2.590087in}{1.761390in}}%
\pgfpathlineto{\pgfqpoint{2.594695in}{1.768652in}}%
\pgfpathlineto{\pgfqpoint{2.601730in}{1.786502in}}%
\pgfpathlineto{\pgfqpoint{2.602311in}{1.800403in}}%
\pgfpathlineto{\pgfqpoint{2.625838in}{1.802295in}}%
\pgfpathlineto{\pgfqpoint{2.628418in}{1.774736in}}%
\pgfpathlineto{\pgfqpoint{2.615872in}{1.773777in}}%
\pgfpathlineto{\pgfqpoint{2.618822in}{1.745974in}}%
\pgfpathclose%
\pgfusepath{fill}%
\end{pgfscope}%
\begin{pgfscope}%
\pgfpathrectangle{\pgfqpoint{0.100000in}{0.100000in}}{\pgfqpoint{3.608454in}{2.310000in}}%
\pgfusepath{clip}%
\pgfsetbuttcap%
\pgfsetmiterjoin%
\definecolor{currentfill}{rgb}{0.000000,0.682353,0.658824}%
\pgfsetfillcolor{currentfill}%
\pgfsetlinewidth{0.000000pt}%
\definecolor{currentstroke}{rgb}{0.000000,0.000000,0.000000}%
\pgfsetstrokecolor{currentstroke}%
\pgfsetstrokeopacity{0.000000}%
\pgfsetdash{}{0pt}%
\pgfpathmoveto{\pgfqpoint{1.194163in}{1.451751in}}%
\pgfpathlineto{\pgfqpoint{1.192855in}{1.443180in}}%
\pgfpathlineto{\pgfqpoint{1.188609in}{1.438112in}}%
\pgfpathlineto{\pgfqpoint{1.189342in}{1.433929in}}%
\pgfpathlineto{\pgfqpoint{1.182094in}{1.423685in}}%
\pgfpathlineto{\pgfqpoint{1.177987in}{1.406555in}}%
\pgfpathlineto{\pgfqpoint{1.181137in}{1.398548in}}%
\pgfpathlineto{\pgfqpoint{1.179215in}{1.395609in}}%
\pgfpathlineto{\pgfqpoint{1.183363in}{1.379999in}}%
\pgfpathlineto{\pgfqpoint{1.181575in}{1.376785in}}%
\pgfpathlineto{\pgfqpoint{1.141855in}{1.383558in}}%
\pgfpathlineto{\pgfqpoint{1.103543in}{1.390370in}}%
\pgfpathlineto{\pgfqpoint{1.103689in}{1.391163in}}%
\pgfpathlineto{\pgfqpoint{1.111444in}{1.431696in}}%
\pgfpathlineto{\pgfqpoint{1.077995in}{1.437925in}}%
\pgfpathlineto{\pgfqpoint{1.068484in}{1.440713in}}%
\pgfpathlineto{\pgfqpoint{1.072686in}{1.463362in}}%
\pgfpathlineto{\pgfqpoint{1.078936in}{1.465774in}}%
\pgfpathlineto{\pgfqpoint{1.089487in}{1.463907in}}%
\pgfpathlineto{\pgfqpoint{1.092637in}{1.469533in}}%
\pgfpathlineto{\pgfqpoint{1.095705in}{1.487818in}}%
\pgfpathlineto{\pgfqpoint{1.104984in}{1.493265in}}%
\pgfpathlineto{\pgfqpoint{1.101981in}{1.496172in}}%
\pgfpathlineto{\pgfqpoint{1.148850in}{1.487724in}}%
\pgfpathlineto{\pgfqpoint{1.169072in}{1.483661in}}%
\pgfpathlineto{\pgfqpoint{1.206990in}{1.477188in}}%
\pgfpathlineto{\pgfqpoint{1.199560in}{1.472867in}}%
\pgfpathlineto{\pgfqpoint{1.199187in}{1.467135in}}%
\pgfpathlineto{\pgfqpoint{1.194698in}{1.460074in}}%
\pgfpathlineto{\pgfqpoint{1.194163in}{1.451751in}}%
\pgfpathclose%
\pgfusepath{fill}%
\end{pgfscope}%
\begin{pgfscope}%
\pgfpathrectangle{\pgfqpoint{0.100000in}{0.100000in}}{\pgfqpoint{3.608454in}{2.310000in}}%
\pgfusepath{clip}%
\pgfsetbuttcap%
\pgfsetmiterjoin%
\definecolor{currentfill}{rgb}{0.000000,0.945098,0.527451}%
\pgfsetfillcolor{currentfill}%
\pgfsetlinewidth{0.000000pt}%
\definecolor{currentstroke}{rgb}{0.000000,0.000000,0.000000}%
\pgfsetstrokecolor{currentstroke}%
\pgfsetstrokeopacity{0.000000}%
\pgfsetdash{}{0pt}%
\pgfpathmoveto{\pgfqpoint{0.456983in}{1.837942in}}%
\pgfpathlineto{\pgfqpoint{0.452336in}{1.831556in}}%
\pgfpathlineto{\pgfqpoint{0.442382in}{1.828434in}}%
\pgfpathlineto{\pgfqpoint{0.441986in}{1.818955in}}%
\pgfpathlineto{\pgfqpoint{0.436698in}{1.812439in}}%
\pgfpathlineto{\pgfqpoint{0.438276in}{1.805265in}}%
\pgfpathlineto{\pgfqpoint{0.437233in}{1.797086in}}%
\pgfpathlineto{\pgfqpoint{0.433774in}{1.794008in}}%
\pgfpathlineto{\pgfqpoint{0.427599in}{1.795918in}}%
\pgfpathlineto{\pgfqpoint{0.429636in}{1.802194in}}%
\pgfpathlineto{\pgfqpoint{0.413061in}{1.807713in}}%
\pgfpathlineto{\pgfqpoint{0.414756in}{1.827792in}}%
\pgfpathlineto{\pgfqpoint{0.410295in}{1.835155in}}%
\pgfpathlineto{\pgfqpoint{0.417067in}{1.845633in}}%
\pgfpathlineto{\pgfqpoint{0.418133in}{1.851097in}}%
\pgfpathlineto{\pgfqpoint{0.412925in}{1.860819in}}%
\pgfpathlineto{\pgfqpoint{0.414420in}{1.867082in}}%
\pgfpathlineto{\pgfqpoint{0.414023in}{1.879041in}}%
\pgfpathlineto{\pgfqpoint{0.418855in}{1.891440in}}%
\pgfpathlineto{\pgfqpoint{0.422736in}{1.897083in}}%
\pgfpathlineto{\pgfqpoint{0.423493in}{1.904021in}}%
\pgfpathlineto{\pgfqpoint{0.420036in}{1.911157in}}%
\pgfpathlineto{\pgfqpoint{0.420507in}{1.920108in}}%
\pgfpathlineto{\pgfqpoint{0.427334in}{1.927274in}}%
\pgfpathlineto{\pgfqpoint{0.439586in}{1.923181in}}%
\pgfpathlineto{\pgfqpoint{0.443285in}{1.914029in}}%
\pgfpathlineto{\pgfqpoint{0.438998in}{1.899672in}}%
\pgfpathlineto{\pgfqpoint{0.447212in}{1.898187in}}%
\pgfpathlineto{\pgfqpoint{0.456477in}{1.902564in}}%
\pgfpathlineto{\pgfqpoint{0.465159in}{1.900474in}}%
\pgfpathlineto{\pgfqpoint{0.455038in}{1.891185in}}%
\pgfpathlineto{\pgfqpoint{0.448521in}{1.882309in}}%
\pgfpathlineto{\pgfqpoint{0.441915in}{1.883291in}}%
\pgfpathlineto{\pgfqpoint{0.437497in}{1.874169in}}%
\pgfpathlineto{\pgfqpoint{0.444876in}{1.871177in}}%
\pgfpathlineto{\pgfqpoint{0.447509in}{1.860643in}}%
\pgfpathlineto{\pgfqpoint{0.441504in}{1.854677in}}%
\pgfpathlineto{\pgfqpoint{0.439878in}{1.843077in}}%
\pgfpathlineto{\pgfqpoint{0.456983in}{1.837942in}}%
\pgfpathclose%
\pgfusepath{fill}%
\end{pgfscope}%
\begin{pgfscope}%
\pgfpathrectangle{\pgfqpoint{0.100000in}{0.100000in}}{\pgfqpoint{3.608454in}{2.310000in}}%
\pgfusepath{clip}%
\pgfsetbuttcap%
\pgfsetmiterjoin%
\definecolor{currentfill}{rgb}{0.000000,0.639216,0.680392}%
\pgfsetfillcolor{currentfill}%
\pgfsetlinewidth{0.000000pt}%
\definecolor{currentstroke}{rgb}{0.000000,0.000000,0.000000}%
\pgfsetstrokecolor{currentstroke}%
\pgfsetstrokeopacity{0.000000}%
\pgfsetdash{}{0pt}%
\pgfpathmoveto{\pgfqpoint{2.720165in}{1.477649in}}%
\pgfpathlineto{\pgfqpoint{2.722450in}{1.453627in}}%
\pgfpathlineto{\pgfqpoint{2.698943in}{1.451000in}}%
\pgfpathlineto{\pgfqpoint{2.699598in}{1.433892in}}%
\pgfpathlineto{\pgfqpoint{2.694747in}{1.433296in}}%
\pgfpathlineto{\pgfqpoint{2.676915in}{1.431338in}}%
\pgfpathlineto{\pgfqpoint{2.676851in}{1.443992in}}%
\pgfpathlineto{\pgfqpoint{2.659619in}{1.442049in}}%
\pgfpathlineto{\pgfqpoint{2.657467in}{1.463709in}}%
\pgfpathlineto{\pgfqpoint{2.653860in}{1.497742in}}%
\pgfpathlineto{\pgfqpoint{2.667350in}{1.499209in}}%
\pgfpathlineto{\pgfqpoint{2.678668in}{1.500441in}}%
\pgfpathlineto{\pgfqpoint{2.679497in}{1.493619in}}%
\pgfpathlineto{\pgfqpoint{2.694173in}{1.495218in}}%
\pgfpathlineto{\pgfqpoint{2.695336in}{1.474638in}}%
\pgfpathlineto{\pgfqpoint{2.720165in}{1.477649in}}%
\pgfpathclose%
\pgfusepath{fill}%
\end{pgfscope}%
\begin{pgfscope}%
\pgfpathrectangle{\pgfqpoint{0.100000in}{0.100000in}}{\pgfqpoint{3.608454in}{2.310000in}}%
\pgfusepath{clip}%
\pgfsetbuttcap%
\pgfsetmiterjoin%
\definecolor{currentfill}{rgb}{0.000000,0.635294,0.682353}%
\pgfsetfillcolor{currentfill}%
\pgfsetlinewidth{0.000000pt}%
\definecolor{currentstroke}{rgb}{0.000000,0.000000,0.000000}%
\pgfsetstrokecolor{currentstroke}%
\pgfsetstrokeopacity{0.000000}%
\pgfsetdash{}{0pt}%
\pgfpathmoveto{\pgfqpoint{2.809988in}{1.252552in}}%
\pgfpathlineto{\pgfqpoint{2.801049in}{1.243631in}}%
\pgfpathlineto{\pgfqpoint{2.797121in}{1.243378in}}%
\pgfpathlineto{\pgfqpoint{2.788710in}{1.248796in}}%
\pgfpathlineto{\pgfqpoint{2.781851in}{1.240501in}}%
\pgfpathlineto{\pgfqpoint{2.780361in}{1.233265in}}%
\pgfpathlineto{\pgfqpoint{2.765354in}{1.245580in}}%
\pgfpathlineto{\pgfqpoint{2.767417in}{1.258344in}}%
\pgfpathlineto{\pgfqpoint{2.773095in}{1.263245in}}%
\pgfpathlineto{\pgfqpoint{2.759930in}{1.280063in}}%
\pgfpathlineto{\pgfqpoint{2.774318in}{1.293523in}}%
\pgfpathlineto{\pgfqpoint{2.791570in}{1.288910in}}%
\pgfpathlineto{\pgfqpoint{2.801059in}{1.283583in}}%
\pgfpathlineto{\pgfqpoint{2.810307in}{1.281879in}}%
\pgfpathlineto{\pgfqpoint{2.804559in}{1.276722in}}%
\pgfpathlineto{\pgfqpoint{2.800586in}{1.266861in}}%
\pgfpathlineto{\pgfqpoint{2.802573in}{1.261875in}}%
\pgfpathlineto{\pgfqpoint{2.809988in}{1.252552in}}%
\pgfpathclose%
\pgfusepath{fill}%
\end{pgfscope}%
\begin{pgfscope}%
\pgfpathrectangle{\pgfqpoint{0.100000in}{0.100000in}}{\pgfqpoint{3.608454in}{2.310000in}}%
\pgfusepath{clip}%
\pgfsetbuttcap%
\pgfsetmiterjoin%
\definecolor{currentfill}{rgb}{0.000000,0.964706,0.517647}%
\pgfsetfillcolor{currentfill}%
\pgfsetlinewidth{0.000000pt}%
\definecolor{currentstroke}{rgb}{0.000000,0.000000,0.000000}%
\pgfsetstrokecolor{currentstroke}%
\pgfsetstrokeopacity{0.000000}%
\pgfsetdash{}{0pt}%
\pgfpathmoveto{\pgfqpoint{2.182637in}{0.773027in}}%
\pgfpathlineto{\pgfqpoint{2.212262in}{0.773539in}}%
\pgfpathlineto{\pgfqpoint{2.212402in}{0.763123in}}%
\pgfpathlineto{\pgfqpoint{2.219367in}{0.763229in}}%
\pgfpathlineto{\pgfqpoint{2.219726in}{0.749195in}}%
\pgfpathlineto{\pgfqpoint{2.226572in}{0.749343in}}%
\pgfpathlineto{\pgfqpoint{2.226879in}{0.735487in}}%
\pgfpathlineto{\pgfqpoint{2.216369in}{0.735263in}}%
\pgfpathlineto{\pgfqpoint{2.212945in}{0.728230in}}%
\pgfpathlineto{\pgfqpoint{2.206617in}{0.728100in}}%
\pgfpathlineto{\pgfqpoint{2.205558in}{0.720837in}}%
\pgfpathlineto{\pgfqpoint{2.202317in}{0.720035in}}%
\pgfpathlineto{\pgfqpoint{2.196211in}{0.730685in}}%
\pgfpathlineto{\pgfqpoint{2.197023in}{0.737625in}}%
\pgfpathlineto{\pgfqpoint{2.191816in}{0.743213in}}%
\pgfpathlineto{\pgfqpoint{2.185571in}{0.753524in}}%
\pgfpathlineto{\pgfqpoint{2.187852in}{0.760329in}}%
\pgfpathlineto{\pgfqpoint{2.186530in}{0.767441in}}%
\pgfpathlineto{\pgfqpoint{2.182637in}{0.773027in}}%
\pgfpathclose%
\pgfusepath{fill}%
\end{pgfscope}%
\begin{pgfscope}%
\pgfpathrectangle{\pgfqpoint{0.100000in}{0.100000in}}{\pgfqpoint{3.608454in}{2.310000in}}%
\pgfusepath{clip}%
\pgfsetbuttcap%
\pgfsetmiterjoin%
\definecolor{currentfill}{rgb}{0.000000,0.486275,0.756863}%
\pgfsetfillcolor{currentfill}%
\pgfsetlinewidth{0.000000pt}%
\definecolor{currentstroke}{rgb}{0.000000,0.000000,0.000000}%
\pgfsetstrokecolor{currentstroke}%
\pgfsetstrokeopacity{0.000000}%
\pgfsetdash{}{0pt}%
\pgfpathmoveto{\pgfqpoint{1.732544in}{1.486727in}}%
\pgfpathlineto{\pgfqpoint{1.684512in}{1.490228in}}%
\pgfpathlineto{\pgfqpoint{1.684833in}{1.494331in}}%
\pgfpathlineto{\pgfqpoint{1.649013in}{1.497206in}}%
\pgfpathlineto{\pgfqpoint{1.650561in}{1.517162in}}%
\pgfpathlineto{\pgfqpoint{1.653438in}{1.520561in}}%
\pgfpathlineto{\pgfqpoint{1.654575in}{1.534236in}}%
\pgfpathlineto{\pgfqpoint{1.687474in}{1.531563in}}%
\pgfpathlineto{\pgfqpoint{1.688556in}{1.545253in}}%
\pgfpathlineto{\pgfqpoint{1.692731in}{1.544894in}}%
\pgfpathlineto{\pgfqpoint{1.726904in}{1.542427in}}%
\pgfpathlineto{\pgfqpoint{1.735045in}{1.541820in}}%
\pgfpathlineto{\pgfqpoint{1.733109in}{1.514401in}}%
\pgfpathlineto{\pgfqpoint{1.734320in}{1.514317in}}%
\pgfpathlineto{\pgfqpoint{1.732544in}{1.486727in}}%
\pgfpathclose%
\pgfusepath{fill}%
\end{pgfscope}%
\begin{pgfscope}%
\pgfpathrectangle{\pgfqpoint{0.100000in}{0.100000in}}{\pgfqpoint{3.608454in}{2.310000in}}%
\pgfusepath{clip}%
\pgfsetbuttcap%
\pgfsetmiterjoin%
\definecolor{currentfill}{rgb}{0.000000,0.596078,0.701961}%
\pgfsetfillcolor{currentfill}%
\pgfsetlinewidth{0.000000pt}%
\definecolor{currentstroke}{rgb}{0.000000,0.000000,0.000000}%
\pgfsetstrokecolor{currentstroke}%
\pgfsetstrokeopacity{0.000000}%
\pgfsetdash{}{0pt}%
\pgfpathmoveto{\pgfqpoint{2.610086in}{1.175304in}}%
\pgfpathlineto{\pgfqpoint{2.607679in}{1.208265in}}%
\pgfpathlineto{\pgfqpoint{2.620792in}{1.209072in}}%
\pgfpathlineto{\pgfqpoint{2.631929in}{1.205576in}}%
\pgfpathlineto{\pgfqpoint{2.634641in}{1.217991in}}%
\pgfpathlineto{\pgfqpoint{2.639824in}{1.220359in}}%
\pgfpathlineto{\pgfqpoint{2.647960in}{1.220595in}}%
\pgfpathlineto{\pgfqpoint{2.642473in}{1.232080in}}%
\pgfpathlineto{\pgfqpoint{2.661850in}{1.235103in}}%
\pgfpathlineto{\pgfqpoint{2.669571in}{1.226318in}}%
\pgfpathlineto{\pgfqpoint{2.669252in}{1.213390in}}%
\pgfpathlineto{\pgfqpoint{2.664358in}{1.203424in}}%
\pgfpathlineto{\pgfqpoint{2.677443in}{1.189404in}}%
\pgfpathlineto{\pgfqpoint{2.678723in}{1.180579in}}%
\pgfpathlineto{\pgfqpoint{2.656122in}{1.179965in}}%
\pgfpathlineto{\pgfqpoint{2.639785in}{1.178722in}}%
\pgfpathlineto{\pgfqpoint{2.610086in}{1.175304in}}%
\pgfpathclose%
\pgfusepath{fill}%
\end{pgfscope}%
\begin{pgfscope}%
\pgfpathrectangle{\pgfqpoint{0.100000in}{0.100000in}}{\pgfqpoint{3.608454in}{2.310000in}}%
\pgfusepath{clip}%
\pgfsetbuttcap%
\pgfsetmiterjoin%
\definecolor{currentfill}{rgb}{0.000000,0.564706,0.717647}%
\pgfsetfillcolor{currentfill}%
\pgfsetlinewidth{0.000000pt}%
\definecolor{currentstroke}{rgb}{0.000000,0.000000,0.000000}%
\pgfsetstrokecolor{currentstroke}%
\pgfsetstrokeopacity{0.000000}%
\pgfsetdash{}{0pt}%
\pgfpathmoveto{\pgfqpoint{0.975129in}{1.855275in}}%
\pgfpathlineto{\pgfqpoint{0.981526in}{1.858299in}}%
\pgfpathlineto{\pgfqpoint{0.983732in}{1.852523in}}%
\pgfpathlineto{\pgfqpoint{0.991130in}{1.850884in}}%
\pgfpathlineto{\pgfqpoint{0.999848in}{1.846301in}}%
\pgfpathlineto{\pgfqpoint{1.012788in}{1.842181in}}%
\pgfpathlineto{\pgfqpoint{1.013709in}{1.836495in}}%
\pgfpathlineto{\pgfqpoint{1.027024in}{1.825238in}}%
\pgfpathlineto{\pgfqpoint{1.034613in}{1.812423in}}%
\pgfpathlineto{\pgfqpoint{1.038897in}{1.813232in}}%
\pgfpathlineto{\pgfqpoint{1.041894in}{1.801489in}}%
\pgfpathlineto{\pgfqpoint{1.040506in}{1.794892in}}%
\pgfpathlineto{\pgfqpoint{1.055937in}{1.791680in}}%
\pgfpathlineto{\pgfqpoint{1.054623in}{1.785268in}}%
\pgfpathlineto{\pgfqpoint{1.074676in}{1.781213in}}%
\pgfpathlineto{\pgfqpoint{1.071945in}{1.767715in}}%
\pgfpathlineto{\pgfqpoint{1.058573in}{1.770425in}}%
\pgfpathlineto{\pgfqpoint{1.053439in}{1.743275in}}%
\pgfpathlineto{\pgfqpoint{1.055372in}{1.735866in}}%
\pgfpathlineto{\pgfqpoint{1.051274in}{1.732636in}}%
\pgfpathlineto{\pgfqpoint{1.044474in}{1.739072in}}%
\pgfpathlineto{\pgfqpoint{1.038549in}{1.738796in}}%
\pgfpathlineto{\pgfqpoint{1.041428in}{1.752757in}}%
\pgfpathlineto{\pgfqpoint{1.044792in}{1.752067in}}%
\pgfpathlineto{\pgfqpoint{1.050343in}{1.779280in}}%
\pgfpathlineto{\pgfqpoint{1.033321in}{1.782826in}}%
\pgfpathlineto{\pgfqpoint{1.027689in}{1.755650in}}%
\pgfpathlineto{\pgfqpoint{1.023421in}{1.749624in}}%
\pgfpathlineto{\pgfqpoint{1.013857in}{1.751700in}}%
\pgfpathlineto{\pgfqpoint{1.010156in}{1.733884in}}%
\pgfpathlineto{\pgfqpoint{1.002092in}{1.733915in}}%
\pgfpathlineto{\pgfqpoint{1.000366in}{1.726400in}}%
\pgfpathlineto{\pgfqpoint{0.988011in}{1.729060in}}%
\pgfpathlineto{\pgfqpoint{0.981161in}{1.696275in}}%
\pgfpathlineto{\pgfqpoint{0.949739in}{1.703791in}}%
\pgfpathlineto{\pgfqpoint{0.937739in}{1.705964in}}%
\pgfpathlineto{\pgfqpoint{0.953662in}{1.776726in}}%
\pgfpathlineto{\pgfqpoint{0.950896in}{1.777540in}}%
\pgfpathlineto{\pgfqpoint{0.963091in}{1.830834in}}%
\pgfpathlineto{\pgfqpoint{0.968859in}{1.837111in}}%
\pgfpathlineto{\pgfqpoint{0.967810in}{1.841171in}}%
\pgfpathlineto{\pgfqpoint{0.973897in}{1.850229in}}%
\pgfpathlineto{\pgfqpoint{0.975129in}{1.855275in}}%
\pgfpathclose%
\pgfusepath{fill}%
\end{pgfscope}%
\begin{pgfscope}%
\pgfpathrectangle{\pgfqpoint{0.100000in}{0.100000in}}{\pgfqpoint{3.608454in}{2.310000in}}%
\pgfusepath{clip}%
\pgfsetbuttcap%
\pgfsetmiterjoin%
\definecolor{currentfill}{rgb}{0.000000,0.909804,0.545098}%
\pgfsetfillcolor{currentfill}%
\pgfsetlinewidth{0.000000pt}%
\definecolor{currentstroke}{rgb}{0.000000,0.000000,0.000000}%
\pgfsetstrokecolor{currentstroke}%
\pgfsetstrokeopacity{0.000000}%
\pgfsetdash{}{0pt}%
\pgfpathmoveto{\pgfqpoint{0.494828in}{1.757961in}}%
\pgfpathlineto{\pgfqpoint{0.489189in}{1.745779in}}%
\pgfpathlineto{\pgfqpoint{0.485004in}{1.744684in}}%
\pgfpathlineto{\pgfqpoint{0.475334in}{1.730099in}}%
\pgfpathlineto{\pgfqpoint{0.468741in}{1.723271in}}%
\pgfpathlineto{\pgfqpoint{0.469494in}{1.715995in}}%
\pgfpathlineto{\pgfqpoint{0.458930in}{1.713085in}}%
\pgfpathlineto{\pgfqpoint{0.454296in}{1.709584in}}%
\pgfpathlineto{\pgfqpoint{0.448092in}{1.709303in}}%
\pgfpathlineto{\pgfqpoint{0.442506in}{1.705022in}}%
\pgfpathlineto{\pgfqpoint{0.441132in}{1.699873in}}%
\pgfpathlineto{\pgfqpoint{0.444780in}{1.696306in}}%
\pgfpathlineto{\pgfqpoint{0.442069in}{1.688570in}}%
\pgfpathlineto{\pgfqpoint{0.441276in}{1.674970in}}%
\pgfpathlineto{\pgfqpoint{0.406038in}{1.686194in}}%
\pgfpathlineto{\pgfqpoint{0.406652in}{1.688065in}}%
\pgfpathlineto{\pgfqpoint{0.379033in}{1.697172in}}%
\pgfpathlineto{\pgfqpoint{0.376641in}{1.705935in}}%
\pgfpathlineto{\pgfqpoint{0.372917in}{1.709765in}}%
\pgfpathlineto{\pgfqpoint{0.365968in}{1.722668in}}%
\pgfpathlineto{\pgfqpoint{0.368189in}{1.726700in}}%
\pgfpathlineto{\pgfqpoint{0.367902in}{1.737177in}}%
\pgfpathlineto{\pgfqpoint{0.378459in}{1.750564in}}%
\pgfpathlineto{\pgfqpoint{0.394964in}{1.767791in}}%
\pgfpathlineto{\pgfqpoint{0.399441in}{1.775849in}}%
\pgfpathlineto{\pgfqpoint{0.397857in}{1.779799in}}%
\pgfpathlineto{\pgfqpoint{0.406645in}{1.794088in}}%
\pgfpathlineto{\pgfqpoint{0.413061in}{1.807713in}}%
\pgfpathlineto{\pgfqpoint{0.429636in}{1.802194in}}%
\pgfpathlineto{\pgfqpoint{0.427599in}{1.795918in}}%
\pgfpathlineto{\pgfqpoint{0.433774in}{1.794008in}}%
\pgfpathlineto{\pgfqpoint{0.442827in}{1.791103in}}%
\pgfpathlineto{\pgfqpoint{0.443568in}{1.772460in}}%
\pgfpathlineto{\pgfqpoint{0.449403in}{1.767395in}}%
\pgfpathlineto{\pgfqpoint{0.449915in}{1.763500in}}%
\pgfpathlineto{\pgfqpoint{0.457583in}{1.760831in}}%
\pgfpathlineto{\pgfqpoint{0.459959in}{1.754157in}}%
\pgfpathlineto{\pgfqpoint{0.468505in}{1.753515in}}%
\pgfpathlineto{\pgfqpoint{0.466061in}{1.760966in}}%
\pgfpathlineto{\pgfqpoint{0.469737in}{1.765482in}}%
\pgfpathlineto{\pgfqpoint{0.477526in}{1.765034in}}%
\pgfpathlineto{\pgfqpoint{0.494062in}{1.772814in}}%
\pgfpathlineto{\pgfqpoint{0.498294in}{1.770162in}}%
\pgfpathlineto{\pgfqpoint{0.494828in}{1.757961in}}%
\pgfpathclose%
\pgfusepath{fill}%
\end{pgfscope}%
\begin{pgfscope}%
\pgfpathrectangle{\pgfqpoint{0.100000in}{0.100000in}}{\pgfqpoint{3.608454in}{2.310000in}}%
\pgfusepath{clip}%
\pgfsetbuttcap%
\pgfsetmiterjoin%
\definecolor{currentfill}{rgb}{0.000000,0.603922,0.698039}%
\pgfsetfillcolor{currentfill}%
\pgfsetlinewidth{0.000000pt}%
\definecolor{currentstroke}{rgb}{0.000000,0.000000,0.000000}%
\pgfsetstrokecolor{currentstroke}%
\pgfsetstrokeopacity{0.000000}%
\pgfsetdash{}{0pt}%
\pgfpathmoveto{\pgfqpoint{2.246959in}{1.366834in}}%
\pgfpathlineto{\pgfqpoint{2.241843in}{1.368445in}}%
\pgfpathlineto{\pgfqpoint{2.238456in}{1.359055in}}%
\pgfpathlineto{\pgfqpoint{2.233381in}{1.359171in}}%
\pgfpathlineto{\pgfqpoint{2.230926in}{1.365957in}}%
\pgfpathlineto{\pgfqpoint{2.224430in}{1.367560in}}%
\pgfpathlineto{\pgfqpoint{2.221539in}{1.373866in}}%
\pgfpathlineto{\pgfqpoint{2.215503in}{1.374729in}}%
\pgfpathlineto{\pgfqpoint{2.211917in}{1.379066in}}%
\pgfpathlineto{\pgfqpoint{2.212943in}{1.384548in}}%
\pgfpathlineto{\pgfqpoint{2.211089in}{1.394791in}}%
\pgfpathlineto{\pgfqpoint{2.206226in}{1.396171in}}%
\pgfpathlineto{\pgfqpoint{2.206018in}{1.417015in}}%
\pgfpathlineto{\pgfqpoint{2.205934in}{1.422201in}}%
\pgfpathlineto{\pgfqpoint{2.236527in}{1.423164in}}%
\pgfpathlineto{\pgfqpoint{2.237055in}{1.423181in}}%
\pgfpathlineto{\pgfqpoint{2.237052in}{1.396510in}}%
\pgfpathlineto{\pgfqpoint{2.247305in}{1.396610in}}%
\pgfpathlineto{\pgfqpoint{2.246959in}{1.366834in}}%
\pgfpathclose%
\pgfusepath{fill}%
\end{pgfscope}%
\begin{pgfscope}%
\pgfpathrectangle{\pgfqpoint{0.100000in}{0.100000in}}{\pgfqpoint{3.608454in}{2.310000in}}%
\pgfusepath{clip}%
\pgfsetbuttcap%
\pgfsetmiterjoin%
\definecolor{currentfill}{rgb}{0.000000,0.572549,0.713725}%
\pgfsetfillcolor{currentfill}%
\pgfsetlinewidth{0.000000pt}%
\definecolor{currentstroke}{rgb}{0.000000,0.000000,0.000000}%
\pgfsetstrokecolor{currentstroke}%
\pgfsetstrokeopacity{0.000000}%
\pgfsetdash{}{0pt}%
\pgfpathmoveto{\pgfqpoint{1.857172in}{0.996308in}}%
\pgfpathlineto{\pgfqpoint{1.856890in}{0.989419in}}%
\pgfpathlineto{\pgfqpoint{1.867617in}{0.988984in}}%
\pgfpathlineto{\pgfqpoint{1.867064in}{0.975207in}}%
\pgfpathlineto{\pgfqpoint{1.870481in}{0.975087in}}%
\pgfpathlineto{\pgfqpoint{1.869882in}{0.961093in}}%
\pgfpathlineto{\pgfqpoint{1.864563in}{0.959406in}}%
\pgfpathlineto{\pgfqpoint{1.852516in}{0.961694in}}%
\pgfpathlineto{\pgfqpoint{1.847621in}{0.966431in}}%
\pgfpathlineto{\pgfqpoint{1.832097in}{0.967259in}}%
\pgfpathlineto{\pgfqpoint{1.831143in}{0.977031in}}%
\pgfpathlineto{\pgfqpoint{1.831022in}{0.989497in}}%
\pgfpathlineto{\pgfqpoint{1.835171in}{0.996785in}}%
\pgfpathlineto{\pgfqpoint{1.857172in}{0.996308in}}%
\pgfpathclose%
\pgfusepath{fill}%
\end{pgfscope}%
\begin{pgfscope}%
\pgfpathrectangle{\pgfqpoint{0.100000in}{0.100000in}}{\pgfqpoint{3.608454in}{2.310000in}}%
\pgfusepath{clip}%
\pgfsetbuttcap%
\pgfsetmiterjoin%
\definecolor{currentfill}{rgb}{0.000000,0.533333,0.733333}%
\pgfsetfillcolor{currentfill}%
\pgfsetlinewidth{0.000000pt}%
\definecolor{currentstroke}{rgb}{0.000000,0.000000,0.000000}%
\pgfsetstrokecolor{currentstroke}%
\pgfsetstrokeopacity{0.000000}%
\pgfsetdash{}{0pt}%
\pgfpathmoveto{\pgfqpoint{2.046244in}{1.696054in}}%
\pgfpathlineto{\pgfqpoint{2.023239in}{1.696380in}}%
\pgfpathlineto{\pgfqpoint{2.023670in}{1.724067in}}%
\pgfpathlineto{\pgfqpoint{2.024621in}{1.779133in}}%
\pgfpathlineto{\pgfqpoint{2.024747in}{1.786093in}}%
\pgfpathlineto{\pgfqpoint{2.072659in}{1.785449in}}%
\pgfpathlineto{\pgfqpoint{2.073124in}{1.778481in}}%
\pgfpathlineto{\pgfqpoint{2.073025in}{1.750974in}}%
\pgfpathlineto{\pgfqpoint{2.046268in}{1.751283in}}%
\pgfpathlineto{\pgfqpoint{2.046612in}{1.723715in}}%
\pgfpathlineto{\pgfqpoint{2.046244in}{1.696054in}}%
\pgfpathclose%
\pgfusepath{fill}%
\end{pgfscope}%
\begin{pgfscope}%
\pgfpathrectangle{\pgfqpoint{0.100000in}{0.100000in}}{\pgfqpoint{3.608454in}{2.310000in}}%
\pgfusepath{clip}%
\pgfsetbuttcap%
\pgfsetmiterjoin%
\definecolor{currentfill}{rgb}{0.000000,0.490196,0.754902}%
\pgfsetfillcolor{currentfill}%
\pgfsetlinewidth{0.000000pt}%
\definecolor{currentstroke}{rgb}{0.000000,0.000000,0.000000}%
\pgfsetstrokecolor{currentstroke}%
\pgfsetstrokeopacity{0.000000}%
\pgfsetdash{}{0pt}%
\pgfpathmoveto{\pgfqpoint{1.461309in}{1.650537in}}%
\pgfpathlineto{\pgfqpoint{1.457886in}{1.623314in}}%
\pgfpathlineto{\pgfqpoint{1.451617in}{1.568947in}}%
\pgfpathlineto{\pgfqpoint{1.436631in}{1.570812in}}%
\pgfpathlineto{\pgfqpoint{1.432841in}{1.539719in}}%
\pgfpathlineto{\pgfqpoint{1.401111in}{1.544070in}}%
\pgfpathlineto{\pgfqpoint{1.338390in}{1.552623in}}%
\pgfpathlineto{\pgfqpoint{1.342910in}{1.583654in}}%
\pgfpathlineto{\pgfqpoint{1.345097in}{1.604208in}}%
\pgfpathlineto{\pgfqpoint{1.369658in}{1.600576in}}%
\pgfpathlineto{\pgfqpoint{1.374676in}{1.634274in}}%
\pgfpathlineto{\pgfqpoint{1.375467in}{1.648053in}}%
\pgfpathlineto{\pgfqpoint{1.377332in}{1.661590in}}%
\pgfpathlineto{\pgfqpoint{1.429302in}{1.654336in}}%
\pgfpathlineto{\pgfqpoint{1.461309in}{1.650537in}}%
\pgfpathclose%
\pgfusepath{fill}%
\end{pgfscope}%
\begin{pgfscope}%
\pgfpathrectangle{\pgfqpoint{0.100000in}{0.100000in}}{\pgfqpoint{3.608454in}{2.310000in}}%
\pgfusepath{clip}%
\pgfsetbuttcap%
\pgfsetmiterjoin%
\definecolor{currentfill}{rgb}{0.000000,0.301961,0.849020}%
\pgfsetfillcolor{currentfill}%
\pgfsetlinewidth{0.000000pt}%
\definecolor{currentstroke}{rgb}{0.000000,0.000000,0.000000}%
\pgfsetstrokecolor{currentstroke}%
\pgfsetstrokeopacity{0.000000}%
\pgfsetdash{}{0pt}%
\pgfpathmoveto{\pgfqpoint{1.963356in}{1.420723in}}%
\pgfpathlineto{\pgfqpoint{1.964117in}{1.448246in}}%
\pgfpathlineto{\pgfqpoint{1.965564in}{1.503301in}}%
\pgfpathlineto{\pgfqpoint{1.992970in}{1.502588in}}%
\pgfpathlineto{\pgfqpoint{2.019500in}{1.502021in}}%
\pgfpathlineto{\pgfqpoint{2.019155in}{1.481378in}}%
\pgfpathlineto{\pgfqpoint{2.018170in}{1.419481in}}%
\pgfpathlineto{\pgfqpoint{1.963356in}{1.420723in}}%
\pgfpathclose%
\pgfusepath{fill}%
\end{pgfscope}%
\begin{pgfscope}%
\pgfpathrectangle{\pgfqpoint{0.100000in}{0.100000in}}{\pgfqpoint{3.608454in}{2.310000in}}%
\pgfusepath{clip}%
\pgfsetbuttcap%
\pgfsetmiterjoin%
\definecolor{currentfill}{rgb}{0.000000,0.458824,0.770588}%
\pgfsetfillcolor{currentfill}%
\pgfsetlinewidth{0.000000pt}%
\definecolor{currentstroke}{rgb}{0.000000,0.000000,0.000000}%
\pgfsetstrokecolor{currentstroke}%
\pgfsetstrokeopacity{0.000000}%
\pgfsetdash{}{0pt}%
\pgfpathmoveto{\pgfqpoint{1.960057in}{1.995378in}}%
\pgfpathlineto{\pgfqpoint{1.946235in}{1.995849in}}%
\pgfpathlineto{\pgfqpoint{1.945081in}{2.002729in}}%
\pgfpathlineto{\pgfqpoint{1.946085in}{2.030374in}}%
\pgfpathlineto{\pgfqpoint{1.951584in}{2.030204in}}%
\pgfpathlineto{\pgfqpoint{1.952493in}{2.057946in}}%
\pgfpathlineto{\pgfqpoint{1.951799in}{2.072044in}}%
\pgfpathlineto{\pgfqpoint{1.992478in}{2.070813in}}%
\pgfpathlineto{\pgfqpoint{1.992147in}{2.069188in}}%
\pgfpathlineto{\pgfqpoint{2.026639in}{2.068568in}}%
\pgfpathlineto{\pgfqpoint{2.027218in}{2.042403in}}%
\pgfpathlineto{\pgfqpoint{2.034152in}{2.042276in}}%
\pgfpathlineto{\pgfqpoint{2.034103in}{2.035335in}}%
\pgfpathlineto{\pgfqpoint{2.061781in}{2.035022in}}%
\pgfpathlineto{\pgfqpoint{2.061858in}{2.042025in}}%
\pgfpathlineto{\pgfqpoint{2.068716in}{2.041994in}}%
\pgfpathlineto{\pgfqpoint{2.068749in}{2.048976in}}%
\pgfpathlineto{\pgfqpoint{2.075577in}{2.048780in}}%
\pgfpathlineto{\pgfqpoint{2.075622in}{2.027945in}}%
\pgfpathlineto{\pgfqpoint{2.076864in}{2.014087in}}%
\pgfpathlineto{\pgfqpoint{2.049145in}{2.014243in}}%
\pgfpathlineto{\pgfqpoint{2.006954in}{2.014902in}}%
\pgfpathlineto{\pgfqpoint{2.007393in}{1.994001in}}%
\pgfpathlineto{\pgfqpoint{1.960057in}{1.995378in}}%
\pgfpathclose%
\pgfusepath{fill}%
\end{pgfscope}%
\begin{pgfscope}%
\pgfpathrectangle{\pgfqpoint{0.100000in}{0.100000in}}{\pgfqpoint{3.608454in}{2.310000in}}%
\pgfusepath{clip}%
\pgfsetbuttcap%
\pgfsetmiterjoin%
\definecolor{currentfill}{rgb}{0.000000,0.454902,0.772549}%
\pgfsetfillcolor{currentfill}%
\pgfsetlinewidth{0.000000pt}%
\definecolor{currentstroke}{rgb}{0.000000,0.000000,0.000000}%
\pgfsetstrokecolor{currentstroke}%
\pgfsetstrokeopacity{0.000000}%
\pgfsetdash{}{0pt}%
\pgfpathmoveto{\pgfqpoint{2.100603in}{1.750815in}}%
\pgfpathlineto{\pgfqpoint{2.073025in}{1.750974in}}%
\pgfpathlineto{\pgfqpoint{2.073124in}{1.778481in}}%
\pgfpathlineto{\pgfqpoint{2.086396in}{1.778361in}}%
\pgfpathlineto{\pgfqpoint{2.086458in}{1.790823in}}%
\pgfpathlineto{\pgfqpoint{2.079554in}{1.795195in}}%
\pgfpathlineto{\pgfqpoint{2.079580in}{1.806133in}}%
\pgfpathlineto{\pgfqpoint{2.120405in}{1.806180in}}%
\pgfpathlineto{\pgfqpoint{2.135067in}{1.806256in}}%
\pgfpathlineto{\pgfqpoint{2.135124in}{1.792351in}}%
\pgfpathlineto{\pgfqpoint{2.127691in}{1.792284in}}%
\pgfpathlineto{\pgfqpoint{2.128088in}{1.771585in}}%
\pgfpathlineto{\pgfqpoint{2.118639in}{1.771585in}}%
\pgfpathlineto{\pgfqpoint{2.114392in}{1.774876in}}%
\pgfpathlineto{\pgfqpoint{2.114393in}{1.757738in}}%
\pgfpathlineto{\pgfqpoint{2.100687in}{1.757698in}}%
\pgfpathlineto{\pgfqpoint{2.100603in}{1.750815in}}%
\pgfpathclose%
\pgfusepath{fill}%
\end{pgfscope}%
\begin{pgfscope}%
\pgfpathrectangle{\pgfqpoint{0.100000in}{0.100000in}}{\pgfqpoint{3.608454in}{2.310000in}}%
\pgfusepath{clip}%
\pgfsetbuttcap%
\pgfsetmiterjoin%
\definecolor{currentfill}{rgb}{0.000000,0.796078,0.601961}%
\pgfsetfillcolor{currentfill}%
\pgfsetlinewidth{0.000000pt}%
\definecolor{currentstroke}{rgb}{0.000000,0.000000,0.000000}%
\pgfsetstrokecolor{currentstroke}%
\pgfsetstrokeopacity{0.000000}%
\pgfsetdash{}{0pt}%
\pgfpathmoveto{\pgfqpoint{2.318616in}{2.073687in}}%
\pgfpathlineto{\pgfqpoint{2.328248in}{2.078608in}}%
\pgfpathlineto{\pgfqpoint{2.333925in}{2.066565in}}%
\pgfpathlineto{\pgfqpoint{2.354068in}{2.067548in}}%
\pgfpathlineto{\pgfqpoint{2.366718in}{2.069902in}}%
\pgfpathlineto{\pgfqpoint{2.372839in}{2.068040in}}%
\pgfpathlineto{\pgfqpoint{2.374681in}{2.063610in}}%
\pgfpathlineto{\pgfqpoint{2.381315in}{2.060550in}}%
\pgfpathlineto{\pgfqpoint{2.386736in}{2.063805in}}%
\pgfpathlineto{\pgfqpoint{2.394855in}{2.063381in}}%
\pgfpathlineto{\pgfqpoint{2.394022in}{2.058782in}}%
\pgfpathlineto{\pgfqpoint{2.376788in}{2.048043in}}%
\pgfpathlineto{\pgfqpoint{2.365376in}{2.042875in}}%
\pgfpathlineto{\pgfqpoint{2.346958in}{2.036101in}}%
\pgfpathlineto{\pgfqpoint{2.336678in}{2.029365in}}%
\pgfpathlineto{\pgfqpoint{2.321610in}{2.015873in}}%
\pgfpathlineto{\pgfqpoint{2.320258in}{2.049681in}}%
\pgfpathlineto{\pgfqpoint{2.318616in}{2.073687in}}%
\pgfpathclose%
\pgfusepath{fill}%
\end{pgfscope}%
\begin{pgfscope}%
\pgfpathrectangle{\pgfqpoint{0.100000in}{0.100000in}}{\pgfqpoint{3.608454in}{2.310000in}}%
\pgfusepath{clip}%
\pgfsetbuttcap%
\pgfsetmiterjoin%
\definecolor{currentfill}{rgb}{0.000000,0.737255,0.631373}%
\pgfsetfillcolor{currentfill}%
\pgfsetlinewidth{0.000000pt}%
\definecolor{currentstroke}{rgb}{0.000000,0.000000,0.000000}%
\pgfsetstrokecolor{currentstroke}%
\pgfsetstrokeopacity{0.000000}%
\pgfsetdash{}{0pt}%
\pgfpathmoveto{\pgfqpoint{2.842704in}{1.366486in}}%
\pgfpathlineto{\pgfqpoint{2.836077in}{1.357238in}}%
\pgfpathlineto{\pgfqpoint{2.829317in}{1.356171in}}%
\pgfpathlineto{\pgfqpoint{2.825080in}{1.387318in}}%
\pgfpathlineto{\pgfqpoint{2.817692in}{1.389496in}}%
\pgfpathlineto{\pgfqpoint{2.818458in}{1.403774in}}%
\pgfpathlineto{\pgfqpoint{2.816007in}{1.404773in}}%
\pgfpathlineto{\pgfqpoint{2.815226in}{1.414852in}}%
\pgfpathlineto{\pgfqpoint{2.820238in}{1.426570in}}%
\pgfpathlineto{\pgfqpoint{2.837395in}{1.428067in}}%
\pgfpathlineto{\pgfqpoint{2.852553in}{1.427030in}}%
\pgfpathlineto{\pgfqpoint{2.853130in}{1.419128in}}%
\pgfpathlineto{\pgfqpoint{2.867100in}{1.421144in}}%
\pgfpathlineto{\pgfqpoint{2.867003in}{1.423229in}}%
\pgfpathlineto{\pgfqpoint{2.880748in}{1.424222in}}%
\pgfpathlineto{\pgfqpoint{2.883523in}{1.417325in}}%
\pgfpathlineto{\pgfqpoint{2.881668in}{1.410091in}}%
\pgfpathlineto{\pgfqpoint{2.882640in}{1.396033in}}%
\pgfpathlineto{\pgfqpoint{2.875723in}{1.395664in}}%
\pgfpathlineto{\pgfqpoint{2.876585in}{1.380575in}}%
\pgfpathlineto{\pgfqpoint{2.869245in}{1.379552in}}%
\pgfpathlineto{\pgfqpoint{2.864652in}{1.379289in}}%
\pgfpathlineto{\pgfqpoint{2.864952in}{1.371264in}}%
\pgfpathlineto{\pgfqpoint{2.861822in}{1.358808in}}%
\pgfpathlineto{\pgfqpoint{2.854867in}{1.358657in}}%
\pgfpathlineto{\pgfqpoint{2.850733in}{1.370027in}}%
\pgfpathlineto{\pgfqpoint{2.842704in}{1.366486in}}%
\pgfpathclose%
\pgfusepath{fill}%
\end{pgfscope}%
\begin{pgfscope}%
\pgfpathrectangle{\pgfqpoint{0.100000in}{0.100000in}}{\pgfqpoint{3.608454in}{2.310000in}}%
\pgfusepath{clip}%
\pgfsetbuttcap%
\pgfsetmiterjoin%
\definecolor{currentfill}{rgb}{0.000000,0.537255,0.731373}%
\pgfsetfillcolor{currentfill}%
\pgfsetlinewidth{0.000000pt}%
\definecolor{currentstroke}{rgb}{0.000000,0.000000,0.000000}%
\pgfsetstrokecolor{currentstroke}%
\pgfsetstrokeopacity{0.000000}%
\pgfsetdash{}{0pt}%
\pgfpathmoveto{\pgfqpoint{2.402794in}{1.625171in}}%
\pgfpathlineto{\pgfqpoint{2.404561in}{1.600729in}}%
\pgfpathlineto{\pgfqpoint{2.381293in}{1.599262in}}%
\pgfpathlineto{\pgfqpoint{2.376710in}{1.601691in}}%
\pgfpathlineto{\pgfqpoint{2.374950in}{1.609365in}}%
\pgfpathlineto{\pgfqpoint{2.371380in}{1.613652in}}%
\pgfpathlineto{\pgfqpoint{2.360207in}{1.613167in}}%
\pgfpathlineto{\pgfqpoint{2.360507in}{1.606246in}}%
\pgfpathlineto{\pgfqpoint{2.346907in}{1.605656in}}%
\pgfpathlineto{\pgfqpoint{2.306014in}{1.604048in}}%
\pgfpathlineto{\pgfqpoint{2.304826in}{1.617701in}}%
\pgfpathlineto{\pgfqpoint{2.303199in}{1.666191in}}%
\pgfpathlineto{\pgfqpoint{2.301639in}{1.699356in}}%
\pgfpathlineto{\pgfqpoint{2.324231in}{1.700249in}}%
\pgfpathlineto{\pgfqpoint{2.325187in}{1.694132in}}%
\pgfpathlineto{\pgfqpoint{2.325644in}{1.688292in}}%
\pgfpathlineto{\pgfqpoint{2.334254in}{1.681222in}}%
\pgfpathlineto{\pgfqpoint{2.327890in}{1.671376in}}%
\pgfpathlineto{\pgfqpoint{2.330505in}{1.653247in}}%
\pgfpathlineto{\pgfqpoint{2.333060in}{1.651644in}}%
\pgfpathlineto{\pgfqpoint{2.336274in}{1.640299in}}%
\pgfpathlineto{\pgfqpoint{2.342986in}{1.636283in}}%
\pgfpathlineto{\pgfqpoint{2.356020in}{1.633465in}}%
\pgfpathlineto{\pgfqpoint{2.361136in}{1.623207in}}%
\pgfpathlineto{\pgfqpoint{2.402794in}{1.625171in}}%
\pgfpathclose%
\pgfusepath{fill}%
\end{pgfscope}%
\begin{pgfscope}%
\pgfpathrectangle{\pgfqpoint{0.100000in}{0.100000in}}{\pgfqpoint{3.608454in}{2.310000in}}%
\pgfusepath{clip}%
\pgfsetbuttcap%
\pgfsetmiterjoin%
\definecolor{currentfill}{rgb}{0.000000,0.576471,0.711765}%
\pgfsetfillcolor{currentfill}%
\pgfsetlinewidth{0.000000pt}%
\definecolor{currentstroke}{rgb}{0.000000,0.000000,0.000000}%
\pgfsetstrokecolor{currentstroke}%
\pgfsetstrokeopacity{0.000000}%
\pgfsetdash{}{0pt}%
\pgfpathmoveto{\pgfqpoint{3.105607in}{1.274029in}}%
\pgfpathlineto{\pgfqpoint{3.099433in}{1.272390in}}%
\pgfpathlineto{\pgfqpoint{3.095309in}{1.277259in}}%
\pgfpathlineto{\pgfqpoint{3.089238in}{1.276542in}}%
\pgfpathlineto{\pgfqpoint{3.081510in}{1.270060in}}%
\pgfpathlineto{\pgfqpoint{3.074664in}{1.267099in}}%
\pgfpathlineto{\pgfqpoint{3.069985in}{1.273083in}}%
\pgfpathlineto{\pgfqpoint{3.057510in}{1.278132in}}%
\pgfpathlineto{\pgfqpoint{3.055501in}{1.285981in}}%
\pgfpathlineto{\pgfqpoint{3.057612in}{1.293606in}}%
\pgfpathlineto{\pgfqpoint{3.064782in}{1.300929in}}%
\pgfpathlineto{\pgfqpoint{3.067687in}{1.299266in}}%
\pgfpathlineto{\pgfqpoint{3.074021in}{1.306357in}}%
\pgfpathlineto{\pgfqpoint{3.075683in}{1.311970in}}%
\pgfpathlineto{\pgfqpoint{3.081757in}{1.317949in}}%
\pgfpathlineto{\pgfqpoint{3.084921in}{1.329255in}}%
\pgfpathlineto{\pgfqpoint{3.092431in}{1.332135in}}%
\pgfpathlineto{\pgfqpoint{3.098792in}{1.325084in}}%
\pgfpathlineto{\pgfqpoint{3.103381in}{1.324700in}}%
\pgfpathlineto{\pgfqpoint{3.110914in}{1.313734in}}%
\pgfpathlineto{\pgfqpoint{3.115324in}{1.315206in}}%
\pgfpathlineto{\pgfqpoint{3.125802in}{1.307188in}}%
\pgfpathlineto{\pgfqpoint{3.131753in}{1.305762in}}%
\pgfpathlineto{\pgfqpoint{3.128475in}{1.293170in}}%
\pgfpathlineto{\pgfqpoint{3.120247in}{1.288198in}}%
\pgfpathlineto{\pgfqpoint{3.117815in}{1.273052in}}%
\pgfpathlineto{\pgfqpoint{3.105607in}{1.274029in}}%
\pgfpathclose%
\pgfusepath{fill}%
\end{pgfscope}%
\begin{pgfscope}%
\pgfpathrectangle{\pgfqpoint{0.100000in}{0.100000in}}{\pgfqpoint{3.608454in}{2.310000in}}%
\pgfusepath{clip}%
\pgfsetbuttcap%
\pgfsetmiterjoin%
\definecolor{currentfill}{rgb}{0.000000,0.909804,0.545098}%
\pgfsetfillcolor{currentfill}%
\pgfsetlinewidth{0.000000pt}%
\definecolor{currentstroke}{rgb}{0.000000,0.000000,0.000000}%
\pgfsetstrokecolor{currentstroke}%
\pgfsetstrokeopacity{0.000000}%
\pgfsetdash{}{0pt}%
\pgfpathmoveto{\pgfqpoint{0.616244in}{1.922404in}}%
\pgfpathlineto{\pgfqpoint{0.668725in}{1.906595in}}%
\pgfpathlineto{\pgfqpoint{0.695319in}{1.898980in}}%
\pgfpathlineto{\pgfqpoint{0.686047in}{1.866076in}}%
\pgfpathlineto{\pgfqpoint{0.684062in}{1.866632in}}%
\pgfpathlineto{\pgfqpoint{0.674134in}{1.833816in}}%
\pgfpathlineto{\pgfqpoint{0.706614in}{1.825006in}}%
\pgfpathlineto{\pgfqpoint{0.691075in}{1.767312in}}%
\pgfpathlineto{\pgfqpoint{0.654872in}{1.777406in}}%
\pgfpathlineto{\pgfqpoint{0.605093in}{1.791654in}}%
\pgfpathlineto{\pgfqpoint{0.621562in}{1.848713in}}%
\pgfpathlineto{\pgfqpoint{0.595602in}{1.856614in}}%
\pgfpathlineto{\pgfqpoint{0.609494in}{1.902912in}}%
\pgfpathlineto{\pgfqpoint{0.616244in}{1.922404in}}%
\pgfpathclose%
\pgfusepath{fill}%
\end{pgfscope}%
\begin{pgfscope}%
\pgfpathrectangle{\pgfqpoint{0.100000in}{0.100000in}}{\pgfqpoint{3.608454in}{2.310000in}}%
\pgfusepath{clip}%
\pgfsetbuttcap%
\pgfsetmiterjoin%
\definecolor{currentfill}{rgb}{0.000000,0.403922,0.798039}%
\pgfsetfillcolor{currentfill}%
\pgfsetlinewidth{0.000000pt}%
\definecolor{currentstroke}{rgb}{0.000000,0.000000,0.000000}%
\pgfsetstrokecolor{currentstroke}%
\pgfsetstrokeopacity{0.000000}%
\pgfsetdash{}{0pt}%
\pgfpathmoveto{\pgfqpoint{2.139356in}{1.083433in}}%
\pgfpathlineto{\pgfqpoint{2.134726in}{1.115276in}}%
\pgfpathlineto{\pgfqpoint{2.106071in}{1.115212in}}%
\pgfpathlineto{\pgfqpoint{2.106820in}{1.155478in}}%
\pgfpathlineto{\pgfqpoint{2.131026in}{1.155307in}}%
\pgfpathlineto{\pgfqpoint{2.130979in}{1.163114in}}%
\pgfpathlineto{\pgfqpoint{2.165811in}{1.161887in}}%
\pgfpathlineto{\pgfqpoint{2.165488in}{1.142173in}}%
\pgfpathlineto{\pgfqpoint{2.178864in}{1.142325in}}%
\pgfpathlineto{\pgfqpoint{2.178906in}{1.127188in}}%
\pgfpathlineto{\pgfqpoint{2.192391in}{1.127027in}}%
\pgfpathlineto{\pgfqpoint{2.194911in}{1.119004in}}%
\pgfpathlineto{\pgfqpoint{2.199417in}{1.113130in}}%
\pgfpathlineto{\pgfqpoint{2.206212in}{1.109649in}}%
\pgfpathlineto{\pgfqpoint{2.206038in}{1.100680in}}%
\pgfpathlineto{\pgfqpoint{2.201478in}{1.093804in}}%
\pgfpathlineto{\pgfqpoint{2.201915in}{1.084462in}}%
\pgfpathlineto{\pgfqpoint{2.176762in}{1.083808in}}%
\pgfpathlineto{\pgfqpoint{2.163020in}{1.082921in}}%
\pgfpathlineto{\pgfqpoint{2.139356in}{1.083433in}}%
\pgfpathclose%
\pgfusepath{fill}%
\end{pgfscope}%
\begin{pgfscope}%
\pgfpathrectangle{\pgfqpoint{0.100000in}{0.100000in}}{\pgfqpoint{3.608454in}{2.310000in}}%
\pgfusepath{clip}%
\pgfsetbuttcap%
\pgfsetmiterjoin%
\definecolor{currentfill}{rgb}{0.000000,0.533333,0.733333}%
\pgfsetfillcolor{currentfill}%
\pgfsetlinewidth{0.000000pt}%
\definecolor{currentstroke}{rgb}{0.000000,0.000000,0.000000}%
\pgfsetstrokecolor{currentstroke}%
\pgfsetstrokeopacity{0.000000}%
\pgfsetdash{}{0pt}%
\pgfpathmoveto{\pgfqpoint{3.081713in}{0.979204in}}%
\pgfpathlineto{\pgfqpoint{3.081498in}{0.967552in}}%
\pgfpathlineto{\pgfqpoint{3.070676in}{0.962442in}}%
\pgfpathlineto{\pgfqpoint{3.064743in}{0.967659in}}%
\pgfpathlineto{\pgfqpoint{3.047350in}{0.952216in}}%
\pgfpathlineto{\pgfqpoint{3.034445in}{0.961587in}}%
\pgfpathlineto{\pgfqpoint{3.018836in}{0.966601in}}%
\pgfpathlineto{\pgfqpoint{3.010939in}{0.970937in}}%
\pgfpathlineto{\pgfqpoint{3.005840in}{0.971004in}}%
\pgfpathlineto{\pgfqpoint{3.016045in}{0.985596in}}%
\pgfpathlineto{\pgfqpoint{2.999338in}{0.990301in}}%
\pgfpathlineto{\pgfqpoint{2.988314in}{0.999433in}}%
\pgfpathlineto{\pgfqpoint{2.983807in}{0.993720in}}%
\pgfpathlineto{\pgfqpoint{2.971145in}{0.996030in}}%
\pgfpathlineto{\pgfqpoint{2.966134in}{1.004041in}}%
\pgfpathlineto{\pgfqpoint{2.958909in}{1.001900in}}%
\pgfpathlineto{\pgfqpoint{2.963991in}{1.017466in}}%
\pgfpathlineto{\pgfqpoint{2.961263in}{1.022919in}}%
\pgfpathlineto{\pgfqpoint{2.962380in}{1.029555in}}%
\pgfpathlineto{\pgfqpoint{2.968957in}{1.035717in}}%
\pgfpathlineto{\pgfqpoint{2.975999in}{1.050234in}}%
\pgfpathlineto{\pgfqpoint{2.982517in}{1.046654in}}%
\pgfpathlineto{\pgfqpoint{2.990720in}{1.049995in}}%
\pgfpathlineto{\pgfqpoint{2.989806in}{1.055375in}}%
\pgfpathlineto{\pgfqpoint{3.025339in}{1.058515in}}%
\pgfpathlineto{\pgfqpoint{3.026275in}{1.051810in}}%
\pgfpathlineto{\pgfqpoint{3.034331in}{1.046602in}}%
\pgfpathlineto{\pgfqpoint{3.030270in}{1.030329in}}%
\pgfpathlineto{\pgfqpoint{3.041428in}{1.022979in}}%
\pgfpathlineto{\pgfqpoint{3.046504in}{1.026091in}}%
\pgfpathlineto{\pgfqpoint{3.049502in}{1.021849in}}%
\pgfpathlineto{\pgfqpoint{3.049719in}{1.006843in}}%
\pgfpathlineto{\pgfqpoint{3.052596in}{1.003504in}}%
\pgfpathlineto{\pgfqpoint{3.051594in}{0.997455in}}%
\pgfpathlineto{\pgfqpoint{3.059264in}{0.991086in}}%
\pgfpathlineto{\pgfqpoint{3.065412in}{0.981312in}}%
\pgfpathlineto{\pgfqpoint{3.073391in}{0.976759in}}%
\pgfpathlineto{\pgfqpoint{3.081713in}{0.979204in}}%
\pgfpathclose%
\pgfusepath{fill}%
\end{pgfscope}%
\begin{pgfscope}%
\pgfpathrectangle{\pgfqpoint{0.100000in}{0.100000in}}{\pgfqpoint{3.608454in}{2.310000in}}%
\pgfusepath{clip}%
\pgfsetbuttcap%
\pgfsetmiterjoin%
\definecolor{currentfill}{rgb}{0.000000,0.811765,0.594118}%
\pgfsetfillcolor{currentfill}%
\pgfsetlinewidth{0.000000pt}%
\definecolor{currentstroke}{rgb}{0.000000,0.000000,0.000000}%
\pgfsetstrokecolor{currentstroke}%
\pgfsetstrokeopacity{0.000000}%
\pgfsetdash{}{0pt}%
\pgfpathmoveto{\pgfqpoint{2.876403in}{1.073700in}}%
\pgfpathlineto{\pgfqpoint{2.876200in}{1.069013in}}%
\pgfpathlineto{\pgfqpoint{2.869918in}{1.063176in}}%
\pgfpathlineto{\pgfqpoint{2.864683in}{1.055271in}}%
\pgfpathlineto{\pgfqpoint{2.863588in}{1.049271in}}%
\pgfpathlineto{\pgfqpoint{2.850394in}{1.050302in}}%
\pgfpathlineto{\pgfqpoint{2.847354in}{1.055558in}}%
\pgfpathlineto{\pgfqpoint{2.843160in}{1.055221in}}%
\pgfpathlineto{\pgfqpoint{2.842470in}{1.059416in}}%
\pgfpathlineto{\pgfqpoint{2.850115in}{1.070383in}}%
\pgfpathlineto{\pgfqpoint{2.840206in}{1.081206in}}%
\pgfpathlineto{\pgfqpoint{2.834354in}{1.080705in}}%
\pgfpathlineto{\pgfqpoint{2.835727in}{1.088275in}}%
\pgfpathlineto{\pgfqpoint{2.837326in}{1.091183in}}%
\pgfpathlineto{\pgfqpoint{2.851025in}{1.094550in}}%
\pgfpathlineto{\pgfqpoint{2.858178in}{1.097707in}}%
\pgfpathlineto{\pgfqpoint{2.867388in}{1.085037in}}%
\pgfpathlineto{\pgfqpoint{2.872545in}{1.080983in}}%
\pgfpathlineto{\pgfqpoint{2.876403in}{1.073700in}}%
\pgfpathclose%
\pgfusepath{fill}%
\end{pgfscope}%
\begin{pgfscope}%
\pgfpathrectangle{\pgfqpoint{0.100000in}{0.100000in}}{\pgfqpoint{3.608454in}{2.310000in}}%
\pgfusepath{clip}%
\pgfsetbuttcap%
\pgfsetmiterjoin%
\definecolor{currentfill}{rgb}{0.000000,0.768627,0.615686}%
\pgfsetfillcolor{currentfill}%
\pgfsetlinewidth{0.000000pt}%
\definecolor{currentstroke}{rgb}{0.000000,0.000000,0.000000}%
\pgfsetstrokecolor{currentstroke}%
\pgfsetstrokeopacity{0.000000}%
\pgfsetdash{}{0pt}%
\pgfpathmoveto{\pgfqpoint{2.288166in}{1.125922in}}%
\pgfpathlineto{\pgfqpoint{2.286012in}{1.115699in}}%
\pgfpathlineto{\pgfqpoint{2.286066in}{1.109865in}}%
\pgfpathlineto{\pgfqpoint{2.279199in}{1.109799in}}%
\pgfpathlineto{\pgfqpoint{2.279258in}{1.102907in}}%
\pgfpathlineto{\pgfqpoint{2.272396in}{1.102853in}}%
\pgfpathlineto{\pgfqpoint{2.272381in}{1.109733in}}%
\pgfpathlineto{\pgfqpoint{2.245103in}{1.109604in}}%
\pgfpathlineto{\pgfqpoint{2.241638in}{1.113006in}}%
\pgfpathlineto{\pgfqpoint{2.241311in}{1.142821in}}%
\pgfpathlineto{\pgfqpoint{2.243255in}{1.143432in}}%
\pgfpathlineto{\pgfqpoint{2.248459in}{1.143536in}}%
\pgfpathlineto{\pgfqpoint{2.248336in}{1.167927in}}%
\pgfpathlineto{\pgfqpoint{2.289673in}{1.168159in}}%
\pgfpathlineto{\pgfqpoint{2.289909in}{1.144788in}}%
\pgfpathlineto{\pgfqpoint{2.288004in}{1.144708in}}%
\pgfpathlineto{\pgfqpoint{2.288166in}{1.125922in}}%
\pgfpathclose%
\pgfusepath{fill}%
\end{pgfscope}%
\begin{pgfscope}%
\pgfpathrectangle{\pgfqpoint{0.100000in}{0.100000in}}{\pgfqpoint{3.608454in}{2.310000in}}%
\pgfusepath{clip}%
\pgfsetbuttcap%
\pgfsetmiterjoin%
\definecolor{currentfill}{rgb}{0.000000,0.847059,0.576471}%
\pgfsetfillcolor{currentfill}%
\pgfsetlinewidth{0.000000pt}%
\definecolor{currentstroke}{rgb}{0.000000,0.000000,0.000000}%
\pgfsetstrokecolor{currentstroke}%
\pgfsetstrokeopacity{0.000000}%
\pgfsetdash{}{0pt}%
\pgfpathmoveto{\pgfqpoint{3.197116in}{1.277528in}}%
\pgfpathlineto{\pgfqpoint{3.199968in}{1.263498in}}%
\pgfpathlineto{\pgfqpoint{3.195545in}{1.249102in}}%
\pgfpathlineto{\pgfqpoint{3.187275in}{1.247460in}}%
\pgfpathlineto{\pgfqpoint{3.152496in}{1.240772in}}%
\pgfpathlineto{\pgfqpoint{3.109908in}{1.233039in}}%
\pgfpathlineto{\pgfqpoint{3.104858in}{1.232143in}}%
\pgfpathlineto{\pgfqpoint{3.105607in}{1.274029in}}%
\pgfpathlineto{\pgfqpoint{3.117815in}{1.273052in}}%
\pgfpathlineto{\pgfqpoint{3.120247in}{1.288198in}}%
\pgfpathlineto{\pgfqpoint{3.128475in}{1.293170in}}%
\pgfpathlineto{\pgfqpoint{3.131753in}{1.305762in}}%
\pgfpathlineto{\pgfqpoint{3.140579in}{1.302693in}}%
\pgfpathlineto{\pgfqpoint{3.149648in}{1.328164in}}%
\pgfpathlineto{\pgfqpoint{3.149422in}{1.332537in}}%
\pgfpathlineto{\pgfqpoint{3.153524in}{1.337961in}}%
\pgfpathlineto{\pgfqpoint{3.159782in}{1.331466in}}%
\pgfpathlineto{\pgfqpoint{3.159486in}{1.315467in}}%
\pgfpathlineto{\pgfqpoint{3.154386in}{1.307543in}}%
\pgfpathlineto{\pgfqpoint{3.155665in}{1.301994in}}%
\pgfpathlineto{\pgfqpoint{3.184110in}{1.299032in}}%
\pgfpathlineto{\pgfqpoint{3.178373in}{1.294019in}}%
\pgfpathlineto{\pgfqpoint{3.181322in}{1.282092in}}%
\pgfpathlineto{\pgfqpoint{3.189244in}{1.282467in}}%
\pgfpathlineto{\pgfqpoint{3.197116in}{1.277528in}}%
\pgfpathclose%
\pgfusepath{fill}%
\end{pgfscope}%
\begin{pgfscope}%
\pgfpathrectangle{\pgfqpoint{0.100000in}{0.100000in}}{\pgfqpoint{3.608454in}{2.310000in}}%
\pgfusepath{clip}%
\pgfsetbuttcap%
\pgfsetmiterjoin%
\definecolor{currentfill}{rgb}{0.000000,0.713725,0.643137}%
\pgfsetfillcolor{currentfill}%
\pgfsetlinewidth{0.000000pt}%
\definecolor{currentstroke}{rgb}{0.000000,0.000000,0.000000}%
\pgfsetstrokecolor{currentstroke}%
\pgfsetstrokeopacity{0.000000}%
\pgfsetdash{}{0pt}%
\pgfpathmoveto{\pgfqpoint{1.034613in}{1.812423in}}%
\pgfpathlineto{\pgfqpoint{1.027024in}{1.825238in}}%
\pgfpathlineto{\pgfqpoint{1.013709in}{1.836495in}}%
\pgfpathlineto{\pgfqpoint{1.012788in}{1.842181in}}%
\pgfpathlineto{\pgfqpoint{0.999848in}{1.846301in}}%
\pgfpathlineto{\pgfqpoint{0.991130in}{1.850884in}}%
\pgfpathlineto{\pgfqpoint{0.983732in}{1.852523in}}%
\pgfpathlineto{\pgfqpoint{0.981526in}{1.858299in}}%
\pgfpathlineto{\pgfqpoint{0.975129in}{1.855275in}}%
\pgfpathlineto{\pgfqpoint{0.972986in}{1.865357in}}%
\pgfpathlineto{\pgfqpoint{0.975886in}{1.873015in}}%
\pgfpathlineto{\pgfqpoint{0.968833in}{1.878204in}}%
\pgfpathlineto{\pgfqpoint{0.969814in}{1.885604in}}%
\pgfpathlineto{\pgfqpoint{0.964170in}{1.890384in}}%
\pgfpathlineto{\pgfqpoint{0.969006in}{1.895451in}}%
\pgfpathlineto{\pgfqpoint{0.970494in}{1.901193in}}%
\pgfpathlineto{\pgfqpoint{0.967552in}{1.905305in}}%
\pgfpathlineto{\pgfqpoint{0.969772in}{1.911307in}}%
\pgfpathlineto{\pgfqpoint{0.976747in}{1.912867in}}%
\pgfpathlineto{\pgfqpoint{0.981235in}{1.920433in}}%
\pgfpathlineto{\pgfqpoint{0.989615in}{1.915702in}}%
\pgfpathlineto{\pgfqpoint{1.004445in}{1.925280in}}%
\pgfpathlineto{\pgfqpoint{1.009235in}{1.946422in}}%
\pgfpathlineto{\pgfqpoint{1.012082in}{1.949489in}}%
\pgfpathlineto{\pgfqpoint{1.018689in}{1.956057in}}%
\pgfpathlineto{\pgfqpoint{1.011182in}{1.972459in}}%
\pgfpathlineto{\pgfqpoint{1.018436in}{1.970411in}}%
\pgfpathlineto{\pgfqpoint{1.025284in}{1.975956in}}%
\pgfpathlineto{\pgfqpoint{1.028023in}{1.976205in}}%
\pgfpathlineto{\pgfqpoint{1.034877in}{1.968576in}}%
\pgfpathlineto{\pgfqpoint{1.048671in}{1.970878in}}%
\pgfpathlineto{\pgfqpoint{1.053520in}{1.975559in}}%
\pgfpathlineto{\pgfqpoint{1.061696in}{1.979359in}}%
\pgfpathlineto{\pgfqpoint{1.067306in}{1.970565in}}%
\pgfpathlineto{\pgfqpoint{1.064456in}{1.964544in}}%
\pgfpathlineto{\pgfqpoint{1.067866in}{1.953218in}}%
\pgfpathlineto{\pgfqpoint{1.066629in}{1.948325in}}%
\pgfpathlineto{\pgfqpoint{1.073148in}{1.929599in}}%
\pgfpathlineto{\pgfqpoint{1.077934in}{1.924256in}}%
\pgfpathlineto{\pgfqpoint{1.074674in}{1.909170in}}%
\pgfpathlineto{\pgfqpoint{1.079391in}{1.904532in}}%
\pgfpathlineto{\pgfqpoint{1.088102in}{1.901286in}}%
\pgfpathlineto{\pgfqpoint{1.091986in}{1.894613in}}%
\pgfpathlineto{\pgfqpoint{1.094071in}{1.877735in}}%
\pgfpathlineto{\pgfqpoint{1.093092in}{1.871853in}}%
\pgfpathlineto{\pgfqpoint{1.100477in}{1.862882in}}%
\pgfpathlineto{\pgfqpoint{1.102494in}{1.863989in}}%
\pgfpathlineto{\pgfqpoint{1.100225in}{1.852591in}}%
\pgfpathlineto{\pgfqpoint{1.066665in}{1.859594in}}%
\pgfpathlineto{\pgfqpoint{1.063909in}{1.846064in}}%
\pgfpathlineto{\pgfqpoint{1.061051in}{1.840480in}}%
\pgfpathlineto{\pgfqpoint{1.057180in}{1.823892in}}%
\pgfpathlineto{\pgfqpoint{1.049281in}{1.819620in}}%
\pgfpathlineto{\pgfqpoint{1.037674in}{1.816193in}}%
\pgfpathlineto{\pgfqpoint{1.034613in}{1.812423in}}%
\pgfpathclose%
\pgfusepath{fill}%
\end{pgfscope}%
\begin{pgfscope}%
\pgfpathrectangle{\pgfqpoint{0.100000in}{0.100000in}}{\pgfqpoint{3.608454in}{2.310000in}}%
\pgfusepath{clip}%
\pgfsetbuttcap%
\pgfsetmiterjoin%
\definecolor{currentfill}{rgb}{0.000000,0.666667,0.666667}%
\pgfsetfillcolor{currentfill}%
\pgfsetlinewidth{0.000000pt}%
\definecolor{currentstroke}{rgb}{0.000000,0.000000,0.000000}%
\pgfsetstrokecolor{currentstroke}%
\pgfsetstrokeopacity{0.000000}%
\pgfsetdash{}{0pt}%
\pgfpathmoveto{\pgfqpoint{1.183026in}{1.348649in}}%
\pgfpathlineto{\pgfqpoint{1.177278in}{1.346155in}}%
\pgfpathlineto{\pgfqpoint{1.170873in}{1.338075in}}%
\pgfpathlineto{\pgfqpoint{1.168304in}{1.339366in}}%
\pgfpathlineto{\pgfqpoint{1.159137in}{1.332405in}}%
\pgfpathlineto{\pgfqpoint{1.150177in}{1.332380in}}%
\pgfpathlineto{\pgfqpoint{1.141396in}{1.315621in}}%
\pgfpathlineto{\pgfqpoint{1.134341in}{1.312779in}}%
\pgfpathlineto{\pgfqpoint{1.130720in}{1.308550in}}%
\pgfpathlineto{\pgfqpoint{1.104844in}{1.313028in}}%
\pgfpathlineto{\pgfqpoint{1.024916in}{1.327610in}}%
\pgfpathlineto{\pgfqpoint{1.005205in}{1.331999in}}%
\pgfpathlineto{\pgfqpoint{1.008850in}{1.352444in}}%
\pgfpathlineto{\pgfqpoint{1.015614in}{1.351012in}}%
\pgfpathlineto{\pgfqpoint{1.016803in}{1.357032in}}%
\pgfpathlineto{\pgfqpoint{1.023679in}{1.356359in}}%
\pgfpathlineto{\pgfqpoint{1.026917in}{1.376417in}}%
\pgfpathlineto{\pgfqpoint{1.032898in}{1.379983in}}%
\pgfpathlineto{\pgfqpoint{1.038444in}{1.389127in}}%
\pgfpathlineto{\pgfqpoint{1.034771in}{1.396805in}}%
\pgfpathlineto{\pgfqpoint{1.030622in}{1.400197in}}%
\pgfpathlineto{\pgfqpoint{1.031123in}{1.409823in}}%
\pgfpathlineto{\pgfqpoint{1.036855in}{1.417471in}}%
\pgfpathlineto{\pgfqpoint{1.045136in}{1.415471in}}%
\pgfpathlineto{\pgfqpoint{1.051426in}{1.418402in}}%
\pgfpathlineto{\pgfqpoint{1.053099in}{1.429145in}}%
\pgfpathlineto{\pgfqpoint{1.060054in}{1.435184in}}%
\pgfpathlineto{\pgfqpoint{1.064107in}{1.434490in}}%
\pgfpathlineto{\pgfqpoint{1.068484in}{1.440713in}}%
\pgfpathlineto{\pgfqpoint{1.077995in}{1.437925in}}%
\pgfpathlineto{\pgfqpoint{1.111444in}{1.431696in}}%
\pgfpathlineto{\pgfqpoint{1.103689in}{1.391163in}}%
\pgfpathlineto{\pgfqpoint{1.076288in}{1.396187in}}%
\pgfpathlineto{\pgfqpoint{1.069417in}{1.390629in}}%
\pgfpathlineto{\pgfqpoint{1.065800in}{1.369302in}}%
\pgfpathlineto{\pgfqpoint{1.136475in}{1.356525in}}%
\pgfpathlineto{\pgfqpoint{1.183026in}{1.348649in}}%
\pgfpathclose%
\pgfusepath{fill}%
\end{pgfscope}%
\begin{pgfscope}%
\pgfpathrectangle{\pgfqpoint{0.100000in}{0.100000in}}{\pgfqpoint{3.608454in}{2.310000in}}%
\pgfusepath{clip}%
\pgfsetbuttcap%
\pgfsetmiterjoin%
\definecolor{currentfill}{rgb}{0.000000,0.631373,0.684314}%
\pgfsetfillcolor{currentfill}%
\pgfsetlinewidth{0.000000pt}%
\definecolor{currentstroke}{rgb}{0.000000,0.000000,0.000000}%
\pgfsetstrokecolor{currentstroke}%
\pgfsetstrokeopacity{0.000000}%
\pgfsetdash{}{0pt}%
\pgfpathmoveto{\pgfqpoint{2.102150in}{1.181491in}}%
\pgfpathlineto{\pgfqpoint{2.081042in}{1.181525in}}%
\pgfpathlineto{\pgfqpoint{2.079502in}{1.177066in}}%
\pgfpathlineto{\pgfqpoint{2.079386in}{1.149757in}}%
\pgfpathlineto{\pgfqpoint{2.055192in}{1.149919in}}%
\pgfpathlineto{\pgfqpoint{2.055083in}{1.136146in}}%
\pgfpathlineto{\pgfqpoint{2.043055in}{1.136287in}}%
\pgfpathlineto{\pgfqpoint{2.043578in}{1.181785in}}%
\pgfpathlineto{\pgfqpoint{2.045886in}{1.181780in}}%
\pgfpathlineto{\pgfqpoint{2.046716in}{1.239913in}}%
\pgfpathlineto{\pgfqpoint{2.101288in}{1.239463in}}%
\pgfpathlineto{\pgfqpoint{2.102082in}{1.211883in}}%
\pgfpathlineto{\pgfqpoint{2.102150in}{1.181491in}}%
\pgfpathclose%
\pgfusepath{fill}%
\end{pgfscope}%
\begin{pgfscope}%
\pgfpathrectangle{\pgfqpoint{0.100000in}{0.100000in}}{\pgfqpoint{3.608454in}{2.310000in}}%
\pgfusepath{clip}%
\pgfsetbuttcap%
\pgfsetmiterjoin%
\definecolor{currentfill}{rgb}{0.000000,0.654902,0.672549}%
\pgfsetfillcolor{currentfill}%
\pgfsetlinewidth{0.000000pt}%
\definecolor{currentstroke}{rgb}{0.000000,0.000000,0.000000}%
\pgfsetstrokecolor{currentstroke}%
\pgfsetstrokeopacity{0.000000}%
\pgfsetdash{}{0pt}%
\pgfpathmoveto{\pgfqpoint{1.411450in}{2.059486in}}%
\pgfpathlineto{\pgfqpoint{1.406013in}{2.044122in}}%
\pgfpathlineto{\pgfqpoint{1.405800in}{2.033757in}}%
\pgfpathlineto{\pgfqpoint{1.407470in}{2.024546in}}%
\pgfpathlineto{\pgfqpoint{1.403304in}{2.014885in}}%
\pgfpathlineto{\pgfqpoint{1.405583in}{2.011993in}}%
\pgfpathlineto{\pgfqpoint{1.407997in}{2.003940in}}%
\pgfpathlineto{\pgfqpoint{1.364475in}{2.009800in}}%
\pgfpathlineto{\pgfqpoint{1.323552in}{2.016288in}}%
\pgfpathlineto{\pgfqpoint{1.322819in}{2.011773in}}%
\pgfpathlineto{\pgfqpoint{1.303565in}{2.014897in}}%
\pgfpathlineto{\pgfqpoint{1.303017in}{2.028995in}}%
\pgfpathlineto{\pgfqpoint{1.305651in}{2.030921in}}%
\pgfpathlineto{\pgfqpoint{1.309306in}{2.053731in}}%
\pgfpathlineto{\pgfqpoint{1.303684in}{2.058169in}}%
\pgfpathlineto{\pgfqpoint{1.296843in}{2.059295in}}%
\pgfpathlineto{\pgfqpoint{1.290585in}{2.063837in}}%
\pgfpathlineto{\pgfqpoint{1.293100in}{2.076361in}}%
\pgfpathlineto{\pgfqpoint{1.309109in}{2.082641in}}%
\pgfpathlineto{\pgfqpoint{1.312471in}{2.096193in}}%
\pgfpathlineto{\pgfqpoint{1.328054in}{2.095490in}}%
\pgfpathlineto{\pgfqpoint{1.332093in}{2.093598in}}%
\pgfpathlineto{\pgfqpoint{1.342674in}{2.097531in}}%
\pgfpathlineto{\pgfqpoint{1.350200in}{2.095416in}}%
\pgfpathlineto{\pgfqpoint{1.358594in}{2.095823in}}%
\pgfpathlineto{\pgfqpoint{1.362402in}{2.089710in}}%
\pgfpathlineto{\pgfqpoint{1.367934in}{2.081540in}}%
\pgfpathlineto{\pgfqpoint{1.377859in}{2.077250in}}%
\pgfpathlineto{\pgfqpoint{1.385004in}{2.076421in}}%
\pgfpathlineto{\pgfqpoint{1.389242in}{2.072875in}}%
\pgfpathlineto{\pgfqpoint{1.404841in}{2.072375in}}%
\pgfpathlineto{\pgfqpoint{1.411723in}{2.068893in}}%
\pgfpathlineto{\pgfqpoint{1.411450in}{2.059486in}}%
\pgfpathclose%
\pgfusepath{fill}%
\end{pgfscope}%
\begin{pgfscope}%
\pgfpathrectangle{\pgfqpoint{0.100000in}{0.100000in}}{\pgfqpoint{3.608454in}{2.310000in}}%
\pgfusepath{clip}%
\pgfsetbuttcap%
\pgfsetmiterjoin%
\definecolor{currentfill}{rgb}{0.000000,0.521569,0.739216}%
\pgfsetfillcolor{currentfill}%
\pgfsetlinewidth{0.000000pt}%
\definecolor{currentstroke}{rgb}{0.000000,0.000000,0.000000}%
\pgfsetstrokecolor{currentstroke}%
\pgfsetstrokeopacity{0.000000}%
\pgfsetdash{}{0pt}%
\pgfpathmoveto{\pgfqpoint{2.002513in}{1.074759in}}%
\pgfpathlineto{\pgfqpoint{2.002912in}{1.098705in}}%
\pgfpathlineto{\pgfqpoint{2.003198in}{1.115990in}}%
\pgfpathlineto{\pgfqpoint{1.990571in}{1.116149in}}%
\pgfpathlineto{\pgfqpoint{1.990717in}{1.123080in}}%
\pgfpathlineto{\pgfqpoint{1.983944in}{1.123221in}}%
\pgfpathlineto{\pgfqpoint{1.984100in}{1.130082in}}%
\pgfpathlineto{\pgfqpoint{1.977272in}{1.130240in}}%
\pgfpathlineto{\pgfqpoint{1.977572in}{1.143978in}}%
\pgfpathlineto{\pgfqpoint{1.985971in}{1.147335in}}%
\pgfpathlineto{\pgfqpoint{1.984832in}{1.151230in}}%
\pgfpathlineto{\pgfqpoint{1.950456in}{1.151489in}}%
\pgfpathlineto{\pgfqpoint{1.951284in}{1.183534in}}%
\pgfpathlineto{\pgfqpoint{1.971148in}{1.183063in}}%
\pgfpathlineto{\pgfqpoint{2.043578in}{1.181785in}}%
\pgfpathlineto{\pgfqpoint{2.043055in}{1.136287in}}%
\pgfpathlineto{\pgfqpoint{2.055083in}{1.136146in}}%
\pgfpathlineto{\pgfqpoint{2.055192in}{1.149919in}}%
\pgfpathlineto{\pgfqpoint{2.079386in}{1.149757in}}%
\pgfpathlineto{\pgfqpoint{2.079502in}{1.177066in}}%
\pgfpathlineto{\pgfqpoint{2.081042in}{1.181525in}}%
\pgfpathlineto{\pgfqpoint{2.102150in}{1.181491in}}%
\pgfpathlineto{\pgfqpoint{2.106837in}{1.177041in}}%
\pgfpathlineto{\pgfqpoint{2.106820in}{1.155478in}}%
\pgfpathlineto{\pgfqpoint{2.106071in}{1.115212in}}%
\pgfpathlineto{\pgfqpoint{2.099224in}{1.115204in}}%
\pgfpathlineto{\pgfqpoint{2.099232in}{1.108302in}}%
\pgfpathlineto{\pgfqpoint{2.093490in}{1.108334in}}%
\pgfpathlineto{\pgfqpoint{2.088937in}{1.099591in}}%
\pgfpathlineto{\pgfqpoint{2.089737in}{1.087648in}}%
\pgfpathlineto{\pgfqpoint{2.076132in}{1.086986in}}%
\pgfpathlineto{\pgfqpoint{2.068981in}{1.083880in}}%
\pgfpathlineto{\pgfqpoint{2.065063in}{1.091216in}}%
\pgfpathlineto{\pgfqpoint{2.057637in}{1.091246in}}%
\pgfpathlineto{\pgfqpoint{2.030288in}{1.091590in}}%
\pgfpathlineto{\pgfqpoint{2.030056in}{1.074337in}}%
\pgfpathlineto{\pgfqpoint{2.002513in}{1.074759in}}%
\pgfpathclose%
\pgfusepath{fill}%
\end{pgfscope}%
\begin{pgfscope}%
\pgfpathrectangle{\pgfqpoint{0.100000in}{0.100000in}}{\pgfqpoint{3.608454in}{2.310000in}}%
\pgfusepath{clip}%
\pgfsetbuttcap%
\pgfsetmiterjoin%
\definecolor{currentfill}{rgb}{0.000000,0.796078,0.601961}%
\pgfsetfillcolor{currentfill}%
\pgfsetlinewidth{0.000000pt}%
\definecolor{currentstroke}{rgb}{0.000000,0.000000,0.000000}%
\pgfsetstrokecolor{currentstroke}%
\pgfsetstrokeopacity{0.000000}%
\pgfsetdash{}{0pt}%
\pgfpathmoveto{\pgfqpoint{1.841646in}{1.759396in}}%
\pgfpathlineto{\pgfqpoint{1.843095in}{1.749824in}}%
\pgfpathlineto{\pgfqpoint{1.835653in}{1.753170in}}%
\pgfpathlineto{\pgfqpoint{1.830770in}{1.752552in}}%
\pgfpathlineto{\pgfqpoint{1.825951in}{1.758455in}}%
\pgfpathlineto{\pgfqpoint{1.801576in}{1.759791in}}%
\pgfpathlineto{\pgfqpoint{1.801456in}{1.757724in}}%
\pgfpathlineto{\pgfqpoint{1.762841in}{1.760148in}}%
\pgfpathlineto{\pgfqpoint{1.756032in}{1.760574in}}%
\pgfpathlineto{\pgfqpoint{1.757629in}{1.774393in}}%
\pgfpathlineto{\pgfqpoint{1.759531in}{1.801865in}}%
\pgfpathlineto{\pgfqpoint{1.760814in}{1.806325in}}%
\pgfpathlineto{\pgfqpoint{1.767130in}{1.804067in}}%
\pgfpathlineto{\pgfqpoint{1.778292in}{1.808049in}}%
\pgfpathlineto{\pgfqpoint{1.781507in}{1.805854in}}%
\pgfpathlineto{\pgfqpoint{1.787641in}{1.810142in}}%
\pgfpathlineto{\pgfqpoint{1.795074in}{1.805632in}}%
\pgfpathlineto{\pgfqpoint{1.803220in}{1.813204in}}%
\pgfpathlineto{\pgfqpoint{1.802764in}{1.815650in}}%
\pgfpathlineto{\pgfqpoint{1.843805in}{1.813366in}}%
\pgfpathlineto{\pgfqpoint{1.842265in}{1.785638in}}%
\pgfpathlineto{\pgfqpoint{1.841646in}{1.759396in}}%
\pgfpathclose%
\pgfusepath{fill}%
\end{pgfscope}%
\begin{pgfscope}%
\pgfpathrectangle{\pgfqpoint{0.100000in}{0.100000in}}{\pgfqpoint{3.608454in}{2.310000in}}%
\pgfusepath{clip}%
\pgfsetbuttcap%
\pgfsetmiterjoin%
\definecolor{currentfill}{rgb}{0.000000,0.850980,0.574510}%
\pgfsetfillcolor{currentfill}%
\pgfsetlinewidth{0.000000pt}%
\definecolor{currentstroke}{rgb}{0.000000,0.000000,0.000000}%
\pgfsetstrokecolor{currentstroke}%
\pgfsetstrokeopacity{0.000000}%
\pgfsetdash{}{0pt}%
\pgfpathmoveto{\pgfqpoint{2.404037in}{0.740243in}}%
\pgfpathlineto{\pgfqpoint{2.403042in}{0.761124in}}%
\pgfpathlineto{\pgfqpoint{2.395986in}{0.760766in}}%
\pgfpathlineto{\pgfqpoint{2.395693in}{0.767722in}}%
\pgfpathlineto{\pgfqpoint{2.426622in}{0.769362in}}%
\pgfpathlineto{\pgfqpoint{2.436793in}{0.774096in}}%
\pgfpathlineto{\pgfqpoint{2.447021in}{0.775387in}}%
\pgfpathlineto{\pgfqpoint{2.447975in}{0.756350in}}%
\pgfpathlineto{\pgfqpoint{2.449720in}{0.745908in}}%
\pgfpathlineto{\pgfqpoint{2.444160in}{0.745936in}}%
\pgfpathlineto{\pgfqpoint{2.439716in}{0.742196in}}%
\pgfpathlineto{\pgfqpoint{2.404037in}{0.740243in}}%
\pgfpathclose%
\pgfusepath{fill}%
\end{pgfscope}%
\begin{pgfscope}%
\pgfpathrectangle{\pgfqpoint{0.100000in}{0.100000in}}{\pgfqpoint{3.608454in}{2.310000in}}%
\pgfusepath{clip}%
\pgfsetbuttcap%
\pgfsetmiterjoin%
\definecolor{currentfill}{rgb}{0.000000,0.427451,0.786275}%
\pgfsetfillcolor{currentfill}%
\pgfsetlinewidth{0.000000pt}%
\definecolor{currentstroke}{rgb}{0.000000,0.000000,0.000000}%
\pgfsetstrokecolor{currentstroke}%
\pgfsetstrokeopacity{0.000000}%
\pgfsetdash{}{0pt}%
\pgfpathmoveto{\pgfqpoint{1.644176in}{0.728867in}}%
\pgfpathlineto{\pgfqpoint{1.601306in}{0.732857in}}%
\pgfpathlineto{\pgfqpoint{1.599669in}{0.742990in}}%
\pgfpathlineto{\pgfqpoint{1.584520in}{0.753390in}}%
\pgfpathlineto{\pgfqpoint{1.578119in}{0.751288in}}%
\pgfpathlineto{\pgfqpoint{1.580562in}{0.779720in}}%
\pgfpathlineto{\pgfqpoint{1.578420in}{0.779909in}}%
\pgfpathlineto{\pgfqpoint{1.581375in}{0.814353in}}%
\pgfpathlineto{\pgfqpoint{1.563559in}{0.816012in}}%
\pgfpathlineto{\pgfqpoint{1.566607in}{0.850506in}}%
\pgfpathlineto{\pgfqpoint{1.624223in}{0.845743in}}%
\pgfpathlineto{\pgfqpoint{1.658603in}{0.843265in}}%
\pgfpathlineto{\pgfqpoint{1.655619in}{0.808558in}}%
\pgfpathlineto{\pgfqpoint{1.650167in}{0.808902in}}%
\pgfpathlineto{\pgfqpoint{1.644176in}{0.728867in}}%
\pgfpathclose%
\pgfusepath{fill}%
\end{pgfscope}%
\begin{pgfscope}%
\pgfpathrectangle{\pgfqpoint{0.100000in}{0.100000in}}{\pgfqpoint{3.608454in}{2.310000in}}%
\pgfusepath{clip}%
\pgfsetbuttcap%
\pgfsetmiterjoin%
\definecolor{currentfill}{rgb}{0.000000,0.521569,0.739216}%
\pgfsetfillcolor{currentfill}%
\pgfsetlinewidth{0.000000pt}%
\definecolor{currentstroke}{rgb}{0.000000,0.000000,0.000000}%
\pgfsetstrokecolor{currentstroke}%
\pgfsetstrokeopacity{0.000000}%
\pgfsetdash{}{0pt}%
\pgfpathmoveto{\pgfqpoint{2.433309in}{1.321900in}}%
\pgfpathlineto{\pgfqpoint{2.404058in}{1.320682in}}%
\pgfpathlineto{\pgfqpoint{2.402286in}{1.326271in}}%
\pgfpathlineto{\pgfqpoint{2.407445in}{1.335820in}}%
\pgfpathlineto{\pgfqpoint{2.398503in}{1.340979in}}%
\pgfpathlineto{\pgfqpoint{2.386232in}{1.344093in}}%
\pgfpathlineto{\pgfqpoint{2.380869in}{1.336490in}}%
\pgfpathlineto{\pgfqpoint{2.373739in}{1.339722in}}%
\pgfpathlineto{\pgfqpoint{2.369779in}{1.350114in}}%
\pgfpathlineto{\pgfqpoint{2.371543in}{1.352918in}}%
\pgfpathlineto{\pgfqpoint{2.368642in}{1.363523in}}%
\pgfpathlineto{\pgfqpoint{2.363949in}{1.370060in}}%
\pgfpathlineto{\pgfqpoint{2.354879in}{1.376769in}}%
\pgfpathlineto{\pgfqpoint{2.374602in}{1.377391in}}%
\pgfpathlineto{\pgfqpoint{2.376719in}{1.362091in}}%
\pgfpathlineto{\pgfqpoint{2.383166in}{1.360351in}}%
\pgfpathlineto{\pgfqpoint{2.393816in}{1.360813in}}%
\pgfpathlineto{\pgfqpoint{2.403570in}{1.368279in}}%
\pgfpathlineto{\pgfqpoint{2.402134in}{1.388642in}}%
\pgfpathlineto{\pgfqpoint{2.415980in}{1.389550in}}%
\pgfpathlineto{\pgfqpoint{2.439852in}{1.391129in}}%
\pgfpathlineto{\pgfqpoint{2.440862in}{1.377267in}}%
\pgfpathlineto{\pgfqpoint{2.464725in}{1.378777in}}%
\pgfpathlineto{\pgfqpoint{2.465371in}{1.368430in}}%
\pgfpathlineto{\pgfqpoint{2.458623in}{1.367966in}}%
\pgfpathlineto{\pgfqpoint{2.460771in}{1.330448in}}%
\pgfpathlineto{\pgfqpoint{2.439649in}{1.329250in}}%
\pgfpathlineto{\pgfqpoint{2.440188in}{1.322366in}}%
\pgfpathlineto{\pgfqpoint{2.433309in}{1.321900in}}%
\pgfpathclose%
\pgfusepath{fill}%
\end{pgfscope}%
\begin{pgfscope}%
\pgfpathrectangle{\pgfqpoint{0.100000in}{0.100000in}}{\pgfqpoint{3.608454in}{2.310000in}}%
\pgfusepath{clip}%
\pgfsetbuttcap%
\pgfsetmiterjoin%
\definecolor{currentfill}{rgb}{0.000000,0.882353,0.558824}%
\pgfsetfillcolor{currentfill}%
\pgfsetlinewidth{0.000000pt}%
\definecolor{currentstroke}{rgb}{0.000000,0.000000,0.000000}%
\pgfsetstrokecolor{currentstroke}%
\pgfsetstrokeopacity{0.000000}%
\pgfsetdash{}{0pt}%
\pgfpathmoveto{\pgfqpoint{2.342635in}{0.869578in}}%
\pgfpathlineto{\pgfqpoint{2.302086in}{0.868312in}}%
\pgfpathlineto{\pgfqpoint{2.299404in}{0.870120in}}%
\pgfpathlineto{\pgfqpoint{2.297141in}{0.880461in}}%
\pgfpathlineto{\pgfqpoint{2.297645in}{0.883786in}}%
\pgfpathlineto{\pgfqpoint{2.307139in}{0.891576in}}%
\pgfpathlineto{\pgfqpoint{2.306319in}{0.899456in}}%
\pgfpathlineto{\pgfqpoint{2.341651in}{0.899996in}}%
\pgfpathlineto{\pgfqpoint{2.342635in}{0.869578in}}%
\pgfpathclose%
\pgfusepath{fill}%
\end{pgfscope}%
\begin{pgfscope}%
\pgfpathrectangle{\pgfqpoint{0.100000in}{0.100000in}}{\pgfqpoint{3.608454in}{2.310000in}}%
\pgfusepath{clip}%
\pgfsetbuttcap%
\pgfsetmiterjoin%
\definecolor{currentfill}{rgb}{0.000000,0.403922,0.798039}%
\pgfsetfillcolor{currentfill}%
\pgfsetlinewidth{0.000000pt}%
\definecolor{currentstroke}{rgb}{0.000000,0.000000,0.000000}%
\pgfsetstrokecolor{currentstroke}%
\pgfsetstrokeopacity{0.000000}%
\pgfsetdash{}{0pt}%
\pgfpathmoveto{\pgfqpoint{1.726904in}{1.542427in}}%
\pgfpathlineto{\pgfqpoint{1.692731in}{1.544894in}}%
\pgfpathlineto{\pgfqpoint{1.694791in}{1.572316in}}%
\pgfpathlineto{\pgfqpoint{1.689939in}{1.572676in}}%
\pgfpathlineto{\pgfqpoint{1.691656in}{1.593744in}}%
\pgfpathlineto{\pgfqpoint{1.693738in}{1.600532in}}%
\pgfpathlineto{\pgfqpoint{1.695984in}{1.627874in}}%
\pgfpathlineto{\pgfqpoint{1.694322in}{1.628015in}}%
\pgfpathlineto{\pgfqpoint{1.696414in}{1.655020in}}%
\pgfpathlineto{\pgfqpoint{1.696825in}{1.671982in}}%
\pgfpathlineto{\pgfqpoint{1.746207in}{1.668241in}}%
\pgfpathlineto{\pgfqpoint{1.805743in}{1.664546in}}%
\pgfpathlineto{\pgfqpoint{1.805638in}{1.647292in}}%
\pgfpathlineto{\pgfqpoint{1.804075in}{1.619812in}}%
\pgfpathlineto{\pgfqpoint{1.803419in}{1.592288in}}%
\pgfpathlineto{\pgfqpoint{1.763649in}{1.594885in}}%
\pgfpathlineto{\pgfqpoint{1.762096in}{1.567407in}}%
\pgfpathlineto{\pgfqpoint{1.728889in}{1.569871in}}%
\pgfpathlineto{\pgfqpoint{1.726904in}{1.542427in}}%
\pgfpathclose%
\pgfusepath{fill}%
\end{pgfscope}%
\begin{pgfscope}%
\pgfpathrectangle{\pgfqpoint{0.100000in}{0.100000in}}{\pgfqpoint{3.608454in}{2.310000in}}%
\pgfusepath{clip}%
\pgfsetbuttcap%
\pgfsetmiterjoin%
\definecolor{currentfill}{rgb}{0.000000,0.639216,0.680392}%
\pgfsetfillcolor{currentfill}%
\pgfsetlinewidth{0.000000pt}%
\definecolor{currentstroke}{rgb}{0.000000,0.000000,0.000000}%
\pgfsetstrokecolor{currentstroke}%
\pgfsetstrokeopacity{0.000000}%
\pgfsetdash{}{0pt}%
\pgfpathmoveto{\pgfqpoint{2.731544in}{1.590059in}}%
\pgfpathlineto{\pgfqpoint{2.707766in}{1.586556in}}%
\pgfpathlineto{\pgfqpoint{2.706079in}{1.591455in}}%
\pgfpathlineto{\pgfqpoint{2.703247in}{1.616030in}}%
\pgfpathlineto{\pgfqpoint{2.710082in}{1.616628in}}%
\pgfpathlineto{\pgfqpoint{2.706426in}{1.644227in}}%
\pgfpathlineto{\pgfqpoint{2.739935in}{1.648352in}}%
\pgfpathlineto{\pgfqpoint{2.743822in}{1.620714in}}%
\pgfpathlineto{\pgfqpoint{2.764552in}{1.624067in}}%
\pgfpathlineto{\pgfqpoint{2.768691in}{1.595961in}}%
\pgfpathlineto{\pgfqpoint{2.731544in}{1.590059in}}%
\pgfpathclose%
\pgfusepath{fill}%
\end{pgfscope}%
\begin{pgfscope}%
\pgfpathrectangle{\pgfqpoint{0.100000in}{0.100000in}}{\pgfqpoint{3.608454in}{2.310000in}}%
\pgfusepath{clip}%
\pgfsetbuttcap%
\pgfsetmiterjoin%
\definecolor{currentfill}{rgb}{0.000000,0.470588,0.764706}%
\pgfsetfillcolor{currentfill}%
\pgfsetlinewidth{0.000000pt}%
\definecolor{currentstroke}{rgb}{0.000000,0.000000,0.000000}%
\pgfsetstrokecolor{currentstroke}%
\pgfsetstrokeopacity{0.000000}%
\pgfsetdash{}{0pt}%
\pgfpathmoveto{\pgfqpoint{3.269282in}{1.579037in}}%
\pgfpathlineto{\pgfqpoint{3.244201in}{1.592250in}}%
\pgfpathlineto{\pgfqpoint{3.251074in}{1.598426in}}%
\pgfpathlineto{\pgfqpoint{3.234122in}{1.609162in}}%
\pgfpathlineto{\pgfqpoint{3.248038in}{1.619832in}}%
\pgfpathlineto{\pgfqpoint{3.244631in}{1.624617in}}%
\pgfpathlineto{\pgfqpoint{3.248335in}{1.629894in}}%
\pgfpathlineto{\pgfqpoint{3.252558in}{1.629161in}}%
\pgfpathlineto{\pgfqpoint{3.262436in}{1.621007in}}%
\pgfpathlineto{\pgfqpoint{3.258629in}{1.617098in}}%
\pgfpathlineto{\pgfqpoint{3.266342in}{1.608149in}}%
\pgfpathlineto{\pgfqpoint{3.276398in}{1.614345in}}%
\pgfpathlineto{\pgfqpoint{3.284382in}{1.624557in}}%
\pgfpathlineto{\pgfqpoint{3.289903in}{1.618119in}}%
\pgfpathlineto{\pgfqpoint{3.288120in}{1.611032in}}%
\pgfpathlineto{\pgfqpoint{3.284662in}{1.609183in}}%
\pgfpathlineto{\pgfqpoint{3.286080in}{1.595797in}}%
\pgfpathlineto{\pgfqpoint{3.275573in}{1.583652in}}%
\pgfpathlineto{\pgfqpoint{3.269282in}{1.579037in}}%
\pgfpathclose%
\pgfusepath{fill}%
\end{pgfscope}%
\begin{pgfscope}%
\pgfpathrectangle{\pgfqpoint{0.100000in}{0.100000in}}{\pgfqpoint{3.608454in}{2.310000in}}%
\pgfusepath{clip}%
\pgfsetbuttcap%
\pgfsetmiterjoin%
\definecolor{currentfill}{rgb}{0.000000,0.556863,0.721569}%
\pgfsetfillcolor{currentfill}%
\pgfsetlinewidth{0.000000pt}%
\definecolor{currentstroke}{rgb}{0.000000,0.000000,0.000000}%
\pgfsetstrokecolor{currentstroke}%
\pgfsetstrokeopacity{0.000000}%
\pgfsetdash{}{0pt}%
\pgfpathmoveto{\pgfqpoint{1.826743in}{0.706842in}}%
\pgfpathlineto{\pgfqpoint{1.835408in}{0.704889in}}%
\pgfpathlineto{\pgfqpoint{1.833962in}{0.671147in}}%
\pgfpathlineto{\pgfqpoint{1.810600in}{0.672318in}}%
\pgfpathlineto{\pgfqpoint{1.798187in}{0.672904in}}%
\pgfpathlineto{\pgfqpoint{1.799907in}{0.708077in}}%
\pgfpathlineto{\pgfqpoint{1.826743in}{0.706842in}}%
\pgfpathclose%
\pgfusepath{fill}%
\end{pgfscope}%
\begin{pgfscope}%
\pgfpathrectangle{\pgfqpoint{0.100000in}{0.100000in}}{\pgfqpoint{3.608454in}{2.310000in}}%
\pgfusepath{clip}%
\pgfsetbuttcap%
\pgfsetmiterjoin%
\definecolor{currentfill}{rgb}{0.000000,0.627451,0.686275}%
\pgfsetfillcolor{currentfill}%
\pgfsetlinewidth{0.000000pt}%
\definecolor{currentstroke}{rgb}{0.000000,0.000000,0.000000}%
\pgfsetstrokecolor{currentstroke}%
\pgfsetstrokeopacity{0.000000}%
\pgfsetdash{}{0pt}%
\pgfpathmoveto{\pgfqpoint{2.236527in}{1.423164in}}%
\pgfpathlineto{\pgfqpoint{2.205934in}{1.422201in}}%
\pgfpathlineto{\pgfqpoint{2.204405in}{1.465435in}}%
\pgfpathlineto{\pgfqpoint{2.244011in}{1.467014in}}%
\pgfpathlineto{\pgfqpoint{2.246988in}{1.453136in}}%
\pgfpathlineto{\pgfqpoint{2.246319in}{1.447619in}}%
\pgfpathlineto{\pgfqpoint{2.235992in}{1.447300in}}%
\pgfpathlineto{\pgfqpoint{2.236527in}{1.423164in}}%
\pgfpathclose%
\pgfusepath{fill}%
\end{pgfscope}%
\begin{pgfscope}%
\pgfpathrectangle{\pgfqpoint{0.100000in}{0.100000in}}{\pgfqpoint{3.608454in}{2.310000in}}%
\pgfusepath{clip}%
\pgfsetbuttcap%
\pgfsetmiterjoin%
\definecolor{currentfill}{rgb}{0.000000,0.513725,0.743137}%
\pgfsetfillcolor{currentfill}%
\pgfsetlinewidth{0.000000pt}%
\definecolor{currentstroke}{rgb}{0.000000,0.000000,0.000000}%
\pgfsetstrokecolor{currentstroke}%
\pgfsetstrokeopacity{0.000000}%
\pgfsetdash{}{0pt}%
\pgfpathmoveto{\pgfqpoint{3.115324in}{1.315206in}}%
\pgfpathlineto{\pgfqpoint{3.110914in}{1.313734in}}%
\pgfpathlineto{\pgfqpoint{3.103381in}{1.324700in}}%
\pgfpathlineto{\pgfqpoint{3.098792in}{1.325084in}}%
\pgfpathlineto{\pgfqpoint{3.092431in}{1.332135in}}%
\pgfpathlineto{\pgfqpoint{3.084921in}{1.329255in}}%
\pgfpathlineto{\pgfqpoint{3.081757in}{1.317949in}}%
\pgfpathlineto{\pgfqpoint{3.075683in}{1.311970in}}%
\pgfpathlineto{\pgfqpoint{3.074021in}{1.306357in}}%
\pgfpathlineto{\pgfqpoint{3.060878in}{1.314437in}}%
\pgfpathlineto{\pgfqpoint{3.060124in}{1.322440in}}%
\pgfpathlineto{\pgfqpoint{3.050832in}{1.323908in}}%
\pgfpathlineto{\pgfqpoint{3.045914in}{1.315133in}}%
\pgfpathlineto{\pgfqpoint{3.033267in}{1.304579in}}%
\pgfpathlineto{\pgfqpoint{3.022678in}{1.310484in}}%
\pgfpathlineto{\pgfqpoint{3.025954in}{1.320474in}}%
\pgfpathlineto{\pgfqpoint{3.037059in}{1.337144in}}%
\pgfpathlineto{\pgfqpoint{3.038592in}{1.342913in}}%
\pgfpathlineto{\pgfqpoint{3.039437in}{1.352438in}}%
\pgfpathlineto{\pgfqpoint{3.046369in}{1.360627in}}%
\pgfpathlineto{\pgfqpoint{3.044609in}{1.363268in}}%
\pgfpathlineto{\pgfqpoint{3.050299in}{1.374365in}}%
\pgfpathlineto{\pgfqpoint{3.050650in}{1.387261in}}%
\pgfpathlineto{\pgfqpoint{3.058026in}{1.385268in}}%
\pgfpathlineto{\pgfqpoint{3.062954in}{1.378601in}}%
\pgfpathlineto{\pgfqpoint{3.073594in}{1.376808in}}%
\pgfpathlineto{\pgfqpoint{3.078058in}{1.382882in}}%
\pgfpathlineto{\pgfqpoint{3.110807in}{1.366936in}}%
\pgfpathlineto{\pgfqpoint{3.110708in}{1.357752in}}%
\pgfpathlineto{\pgfqpoint{3.107539in}{1.353570in}}%
\pgfpathlineto{\pgfqpoint{3.120128in}{1.335691in}}%
\pgfpathlineto{\pgfqpoint{3.121341in}{1.327517in}}%
\pgfpathlineto{\pgfqpoint{3.116356in}{1.322582in}}%
\pgfpathlineto{\pgfqpoint{3.115324in}{1.315206in}}%
\pgfpathclose%
\pgfusepath{fill}%
\end{pgfscope}%
\begin{pgfscope}%
\pgfpathrectangle{\pgfqpoint{0.100000in}{0.100000in}}{\pgfqpoint{3.608454in}{2.310000in}}%
\pgfusepath{clip}%
\pgfsetbuttcap%
\pgfsetmiterjoin%
\definecolor{currentfill}{rgb}{0.000000,0.556863,0.721569}%
\pgfsetfillcolor{currentfill}%
\pgfsetlinewidth{0.000000pt}%
\definecolor{currentstroke}{rgb}{0.000000,0.000000,0.000000}%
\pgfsetstrokecolor{currentstroke}%
\pgfsetstrokeopacity{0.000000}%
\pgfsetdash{}{0pt}%
\pgfpathmoveto{\pgfqpoint{1.810600in}{0.672318in}}%
\pgfpathlineto{\pgfqpoint{1.809932in}{0.655294in}}%
\pgfpathlineto{\pgfqpoint{1.778772in}{0.657165in}}%
\pgfpathlineto{\pgfqpoint{1.753800in}{0.658495in}}%
\pgfpathlineto{\pgfqpoint{1.755704in}{0.692007in}}%
\pgfpathlineto{\pgfqpoint{1.757470in}{0.722100in}}%
\pgfpathlineto{\pgfqpoint{1.792400in}{0.720174in}}%
\pgfpathlineto{\pgfqpoint{1.791780in}{0.708494in}}%
\pgfpathlineto{\pgfqpoint{1.799907in}{0.708077in}}%
\pgfpathlineto{\pgfqpoint{1.798187in}{0.672904in}}%
\pgfpathlineto{\pgfqpoint{1.810600in}{0.672318in}}%
\pgfpathclose%
\pgfusepath{fill}%
\end{pgfscope}%
\begin{pgfscope}%
\pgfpathrectangle{\pgfqpoint{0.100000in}{0.100000in}}{\pgfqpoint{3.608454in}{2.310000in}}%
\pgfusepath{clip}%
\pgfsetbuttcap%
\pgfsetmiterjoin%
\definecolor{currentfill}{rgb}{0.000000,0.521569,0.739216}%
\pgfsetfillcolor{currentfill}%
\pgfsetlinewidth{0.000000pt}%
\definecolor{currentstroke}{rgb}{0.000000,0.000000,0.000000}%
\pgfsetstrokecolor{currentstroke}%
\pgfsetstrokeopacity{0.000000}%
\pgfsetdash{}{0pt}%
\pgfpathmoveto{\pgfqpoint{3.185535in}{1.030022in}}%
\pgfpathlineto{\pgfqpoint{3.145061in}{1.058883in}}%
\pgfpathlineto{\pgfqpoint{3.150261in}{1.071349in}}%
\pgfpathlineto{\pgfqpoint{3.156843in}{1.077749in}}%
\pgfpathlineto{\pgfqpoint{3.156648in}{1.092471in}}%
\pgfpathlineto{\pgfqpoint{3.148549in}{1.102631in}}%
\pgfpathlineto{\pgfqpoint{3.154787in}{1.104979in}}%
\pgfpathlineto{\pgfqpoint{3.174249in}{1.109020in}}%
\pgfpathlineto{\pgfqpoint{3.182406in}{1.101945in}}%
\pgfpathlineto{\pgfqpoint{3.194034in}{1.089009in}}%
\pgfpathlineto{\pgfqpoint{3.200611in}{1.103096in}}%
\pgfpathlineto{\pgfqpoint{3.228319in}{1.108263in}}%
\pgfpathlineto{\pgfqpoint{3.237815in}{1.092298in}}%
\pgfpathlineto{\pgfqpoint{3.243003in}{1.088398in}}%
\pgfpathlineto{\pgfqpoint{3.231067in}{1.072331in}}%
\pgfpathlineto{\pgfqpoint{3.226781in}{1.061809in}}%
\pgfpathlineto{\pgfqpoint{3.223101in}{1.039243in}}%
\pgfpathlineto{\pgfqpoint{3.209280in}{1.039463in}}%
\pgfpathlineto{\pgfqpoint{3.196751in}{1.036307in}}%
\pgfpathlineto{\pgfqpoint{3.185535in}{1.030022in}}%
\pgfpathclose%
\pgfusepath{fill}%
\end{pgfscope}%
\begin{pgfscope}%
\pgfpathrectangle{\pgfqpoint{0.100000in}{0.100000in}}{\pgfqpoint{3.608454in}{2.310000in}}%
\pgfusepath{clip}%
\pgfsetbuttcap%
\pgfsetmiterjoin%
\definecolor{currentfill}{rgb}{0.000000,0.490196,0.754902}%
\pgfsetfillcolor{currentfill}%
\pgfsetlinewidth{0.000000pt}%
\definecolor{currentstroke}{rgb}{0.000000,0.000000,0.000000}%
\pgfsetstrokecolor{currentstroke}%
\pgfsetstrokeopacity{0.000000}%
\pgfsetdash{}{0pt}%
\pgfpathmoveto{\pgfqpoint{0.628251in}{1.491065in}}%
\pgfpathlineto{\pgfqpoint{0.620520in}{1.503015in}}%
\pgfpathlineto{\pgfqpoint{0.623308in}{1.518183in}}%
\pgfpathlineto{\pgfqpoint{0.620041in}{1.519423in}}%
\pgfpathlineto{\pgfqpoint{0.623821in}{1.536529in}}%
\pgfpathlineto{\pgfqpoint{0.631100in}{1.537480in}}%
\pgfpathlineto{\pgfqpoint{0.632981in}{1.544578in}}%
\pgfpathlineto{\pgfqpoint{0.606969in}{1.551881in}}%
\pgfpathlineto{\pgfqpoint{0.607050in}{1.554247in}}%
\pgfpathlineto{\pgfqpoint{0.592721in}{1.558163in}}%
\pgfpathlineto{\pgfqpoint{0.599878in}{1.583444in}}%
\pgfpathlineto{\pgfqpoint{0.605938in}{1.604509in}}%
\pgfpathlineto{\pgfqpoint{0.622947in}{1.663573in}}%
\pgfpathlineto{\pgfqpoint{0.637413in}{1.715698in}}%
\pgfpathlineto{\pgfqpoint{0.654872in}{1.777406in}}%
\pgfpathlineto{\pgfqpoint{0.691075in}{1.767312in}}%
\pgfpathlineto{\pgfqpoint{0.693113in}{1.766738in}}%
\pgfpathlineto{\pgfqpoint{0.681032in}{1.722500in}}%
\pgfpathlineto{\pgfqpoint{0.671164in}{1.687956in}}%
\pgfpathlineto{\pgfqpoint{0.672323in}{1.687633in}}%
\pgfpathlineto{\pgfqpoint{0.662884in}{1.653976in}}%
\pgfpathlineto{\pgfqpoint{0.661398in}{1.654341in}}%
\pgfpathlineto{\pgfqpoint{0.650489in}{1.614771in}}%
\pgfpathlineto{\pgfqpoint{0.693366in}{1.603165in}}%
\pgfpathlineto{\pgfqpoint{0.755783in}{1.587192in}}%
\pgfpathlineto{\pgfqpoint{0.758589in}{1.580048in}}%
\pgfpathlineto{\pgfqpoint{0.753615in}{1.567143in}}%
\pgfpathlineto{\pgfqpoint{0.753176in}{1.557006in}}%
\pgfpathlineto{\pgfqpoint{0.749642in}{1.550822in}}%
\pgfpathlineto{\pgfqpoint{0.738497in}{1.547647in}}%
\pgfpathlineto{\pgfqpoint{0.730097in}{1.540859in}}%
\pgfpathlineto{\pgfqpoint{0.730363in}{1.531286in}}%
\pgfpathlineto{\pgfqpoint{0.725481in}{1.527601in}}%
\pgfpathlineto{\pgfqpoint{0.724276in}{1.521184in}}%
\pgfpathlineto{\pgfqpoint{0.718561in}{1.521037in}}%
\pgfpathlineto{\pgfqpoint{0.655392in}{1.537674in}}%
\pgfpathlineto{\pgfqpoint{0.647659in}{1.529088in}}%
\pgfpathlineto{\pgfqpoint{0.645777in}{1.522062in}}%
\pgfpathlineto{\pgfqpoint{0.652474in}{1.520325in}}%
\pgfpathlineto{\pgfqpoint{0.643201in}{1.487011in}}%
\pgfpathlineto{\pgfqpoint{0.628251in}{1.491065in}}%
\pgfpathclose%
\pgfusepath{fill}%
\end{pgfscope}%
\begin{pgfscope}%
\pgfpathrectangle{\pgfqpoint{0.100000in}{0.100000in}}{\pgfqpoint{3.608454in}{2.310000in}}%
\pgfusepath{clip}%
\pgfsetbuttcap%
\pgfsetmiterjoin%
\definecolor{currentfill}{rgb}{0.000000,0.470588,0.764706}%
\pgfsetfillcolor{currentfill}%
\pgfsetlinewidth{0.000000pt}%
\definecolor{currentstroke}{rgb}{0.000000,0.000000,0.000000}%
\pgfsetstrokecolor{currentstroke}%
\pgfsetstrokeopacity{0.000000}%
\pgfsetdash{}{0pt}%
\pgfpathmoveto{\pgfqpoint{1.889333in}{1.734748in}}%
\pgfpathlineto{\pgfqpoint{1.882514in}{1.735032in}}%
\pgfpathlineto{\pgfqpoint{1.883458in}{1.755759in}}%
\pgfpathlineto{\pgfqpoint{1.896283in}{1.755231in}}%
\pgfpathlineto{\pgfqpoint{1.897415in}{1.789939in}}%
\pgfpathlineto{\pgfqpoint{1.938531in}{1.788233in}}%
\pgfpathlineto{\pgfqpoint{1.945390in}{1.787969in}}%
\pgfpathlineto{\pgfqpoint{1.945343in}{1.781068in}}%
\pgfpathlineto{\pgfqpoint{1.944449in}{1.753368in}}%
\pgfpathlineto{\pgfqpoint{1.917260in}{1.754415in}}%
\pgfpathlineto{\pgfqpoint{1.916523in}{1.733891in}}%
\pgfpathlineto{\pgfqpoint{1.889333in}{1.734748in}}%
\pgfpathclose%
\pgfusepath{fill}%
\end{pgfscope}%
\begin{pgfscope}%
\pgfpathrectangle{\pgfqpoint{0.100000in}{0.100000in}}{\pgfqpoint{3.608454in}{2.310000in}}%
\pgfusepath{clip}%
\pgfsetbuttcap%
\pgfsetmiterjoin%
\definecolor{currentfill}{rgb}{0.000000,0.647059,0.676471}%
\pgfsetfillcolor{currentfill}%
\pgfsetlinewidth{0.000000pt}%
\definecolor{currentstroke}{rgb}{0.000000,0.000000,0.000000}%
\pgfsetstrokecolor{currentstroke}%
\pgfsetstrokeopacity{0.000000}%
\pgfsetdash{}{0pt}%
\pgfpathmoveto{\pgfqpoint{3.262571in}{1.349107in}}%
\pgfpathlineto{\pgfqpoint{3.255578in}{1.348500in}}%
\pgfpathlineto{\pgfqpoint{3.250135in}{1.356031in}}%
\pgfpathlineto{\pgfqpoint{3.244523in}{1.358215in}}%
\pgfpathlineto{\pgfqpoint{3.240232in}{1.353321in}}%
\pgfpathlineto{\pgfqpoint{3.236077in}{1.357801in}}%
\pgfpathlineto{\pgfqpoint{3.228182in}{1.355944in}}%
\pgfpathlineto{\pgfqpoint{3.222015in}{1.359018in}}%
\pgfpathlineto{\pgfqpoint{3.217018in}{1.367687in}}%
\pgfpathlineto{\pgfqpoint{3.209188in}{1.375295in}}%
\pgfpathlineto{\pgfqpoint{3.214965in}{1.384823in}}%
\pgfpathlineto{\pgfqpoint{3.216571in}{1.392855in}}%
\pgfpathlineto{\pgfqpoint{3.225674in}{1.386154in}}%
\pgfpathlineto{\pgfqpoint{3.242132in}{1.387601in}}%
\pgfpathlineto{\pgfqpoint{3.253966in}{1.377009in}}%
\pgfpathlineto{\pgfqpoint{3.269135in}{1.372212in}}%
\pgfpathlineto{\pgfqpoint{3.266819in}{1.363023in}}%
\pgfpathlineto{\pgfqpoint{3.268336in}{1.356399in}}%
\pgfpathlineto{\pgfqpoint{3.266028in}{1.348977in}}%
\pgfpathlineto{\pgfqpoint{3.262571in}{1.349107in}}%
\pgfpathclose%
\pgfusepath{fill}%
\end{pgfscope}%
\begin{pgfscope}%
\pgfpathrectangle{\pgfqpoint{0.100000in}{0.100000in}}{\pgfqpoint{3.608454in}{2.310000in}}%
\pgfusepath{clip}%
\pgfsetbuttcap%
\pgfsetmiterjoin%
\definecolor{currentfill}{rgb}{0.000000,0.498039,0.750980}%
\pgfsetfillcolor{currentfill}%
\pgfsetlinewidth{0.000000pt}%
\definecolor{currentstroke}{rgb}{0.000000,0.000000,0.000000}%
\pgfsetstrokecolor{currentstroke}%
\pgfsetstrokeopacity{0.000000}%
\pgfsetdash{}{0pt}%
\pgfpathmoveto{\pgfqpoint{1.938292in}{1.256226in}}%
\pgfpathlineto{\pgfqpoint{1.938084in}{1.242176in}}%
\pgfpathlineto{\pgfqpoint{1.890205in}{1.243605in}}%
\pgfpathlineto{\pgfqpoint{1.890012in}{1.250882in}}%
\pgfpathlineto{\pgfqpoint{1.891797in}{1.319820in}}%
\pgfpathlineto{\pgfqpoint{1.892401in}{1.333545in}}%
\pgfpathlineto{\pgfqpoint{1.926584in}{1.332368in}}%
\pgfpathlineto{\pgfqpoint{1.926190in}{1.311718in}}%
\pgfpathlineto{\pgfqpoint{1.925218in}{1.277253in}}%
\pgfpathlineto{\pgfqpoint{1.938903in}{1.276826in}}%
\pgfpathlineto{\pgfqpoint{1.938292in}{1.256226in}}%
\pgfpathclose%
\pgfusepath{fill}%
\end{pgfscope}%
\begin{pgfscope}%
\pgfpathrectangle{\pgfqpoint{0.100000in}{0.100000in}}{\pgfqpoint{3.608454in}{2.310000in}}%
\pgfusepath{clip}%
\pgfsetbuttcap%
\pgfsetmiterjoin%
\definecolor{currentfill}{rgb}{0.000000,0.631373,0.684314}%
\pgfsetfillcolor{currentfill}%
\pgfsetlinewidth{0.000000pt}%
\definecolor{currentstroke}{rgb}{0.000000,0.000000,0.000000}%
\pgfsetstrokecolor{currentstroke}%
\pgfsetstrokeopacity{0.000000}%
\pgfsetdash{}{0pt}%
\pgfpathmoveto{\pgfqpoint{2.856403in}{0.756814in}}%
\pgfpathlineto{\pgfqpoint{2.831112in}{0.753917in}}%
\pgfpathlineto{\pgfqpoint{2.819817in}{0.752621in}}%
\pgfpathlineto{\pgfqpoint{2.818573in}{0.766684in}}%
\pgfpathlineto{\pgfqpoint{2.811462in}{0.766109in}}%
\pgfpathlineto{\pgfqpoint{2.810152in}{0.780002in}}%
\pgfpathlineto{\pgfqpoint{2.820126in}{0.780667in}}%
\pgfpathlineto{\pgfqpoint{2.825145in}{0.783605in}}%
\pgfpathlineto{\pgfqpoint{2.821443in}{0.790502in}}%
\pgfpathlineto{\pgfqpoint{2.821071in}{0.796344in}}%
\pgfpathlineto{\pgfqpoint{2.814786in}{0.800405in}}%
\pgfpathlineto{\pgfqpoint{2.809408in}{0.807037in}}%
\pgfpathlineto{\pgfqpoint{2.808283in}{0.818788in}}%
\pgfpathlineto{\pgfqpoint{2.853827in}{0.823500in}}%
\pgfpathlineto{\pgfqpoint{2.854299in}{0.818491in}}%
\pgfpathlineto{\pgfqpoint{2.862925in}{0.815951in}}%
\pgfpathlineto{\pgfqpoint{2.865088in}{0.801722in}}%
\pgfpathlineto{\pgfqpoint{2.870385in}{0.802308in}}%
\pgfpathlineto{\pgfqpoint{2.875554in}{0.798633in}}%
\pgfpathlineto{\pgfqpoint{2.877599in}{0.779820in}}%
\pgfpathlineto{\pgfqpoint{2.854226in}{0.777256in}}%
\pgfpathlineto{\pgfqpoint{2.856403in}{0.756814in}}%
\pgfpathclose%
\pgfusepath{fill}%
\end{pgfscope}%
\begin{pgfscope}%
\pgfpathrectangle{\pgfqpoint{0.100000in}{0.100000in}}{\pgfqpoint{3.608454in}{2.310000in}}%
\pgfusepath{clip}%
\pgfsetbuttcap%
\pgfsetmiterjoin%
\definecolor{currentfill}{rgb}{0.000000,0.619608,0.690196}%
\pgfsetfillcolor{currentfill}%
\pgfsetlinewidth{0.000000pt}%
\definecolor{currentstroke}{rgb}{0.000000,0.000000,0.000000}%
\pgfsetstrokecolor{currentstroke}%
\pgfsetstrokeopacity{0.000000}%
\pgfsetdash{}{0pt}%
\pgfpathmoveto{\pgfqpoint{2.631561in}{1.887489in}}%
\pgfpathlineto{\pgfqpoint{2.627523in}{1.889173in}}%
\pgfpathlineto{\pgfqpoint{2.629695in}{1.902348in}}%
\pgfpathlineto{\pgfqpoint{2.633726in}{1.900071in}}%
\pgfpathlineto{\pgfqpoint{2.635138in}{1.890834in}}%
\pgfpathlineto{\pgfqpoint{2.631561in}{1.887489in}}%
\pgfpathclose%
\pgfusepath{fill}%
\end{pgfscope}%
\begin{pgfscope}%
\pgfpathrectangle{\pgfqpoint{0.100000in}{0.100000in}}{\pgfqpoint{3.608454in}{2.310000in}}%
\pgfusepath{clip}%
\pgfsetbuttcap%
\pgfsetmiterjoin%
\definecolor{currentfill}{rgb}{0.000000,0.619608,0.690196}%
\pgfsetfillcolor{currentfill}%
\pgfsetlinewidth{0.000000pt}%
\definecolor{currentstroke}{rgb}{0.000000,0.000000,0.000000}%
\pgfsetstrokecolor{currentstroke}%
\pgfsetstrokeopacity{0.000000}%
\pgfsetdash{}{0pt}%
\pgfpathmoveto{\pgfqpoint{2.706092in}{1.900170in}}%
\pgfpathlineto{\pgfqpoint{2.703507in}{1.899857in}}%
\pgfpathlineto{\pgfqpoint{2.707810in}{1.866142in}}%
\pgfpathlineto{\pgfqpoint{2.680662in}{1.863242in}}%
\pgfpathlineto{\pgfqpoint{2.681411in}{1.856342in}}%
\pgfpathlineto{\pgfqpoint{2.674595in}{1.855705in}}%
\pgfpathlineto{\pgfqpoint{2.654113in}{1.853564in}}%
\pgfpathlineto{\pgfqpoint{2.653419in}{1.860520in}}%
\pgfpathlineto{\pgfqpoint{2.644203in}{1.859587in}}%
\pgfpathlineto{\pgfqpoint{2.648170in}{1.869172in}}%
\pgfpathlineto{\pgfqpoint{2.656698in}{1.874267in}}%
\pgfpathlineto{\pgfqpoint{2.661603in}{1.879931in}}%
\pgfpathlineto{\pgfqpoint{2.658078in}{1.883394in}}%
\pgfpathlineto{\pgfqpoint{2.656143in}{1.891002in}}%
\pgfpathlineto{\pgfqpoint{2.664544in}{1.902846in}}%
\pgfpathlineto{\pgfqpoint{2.671993in}{1.905904in}}%
\pgfpathlineto{\pgfqpoint{2.673199in}{1.909457in}}%
\pgfpathlineto{\pgfqpoint{2.691843in}{1.900527in}}%
\pgfpathlineto{\pgfqpoint{2.698822in}{1.902391in}}%
\pgfpathlineto{\pgfqpoint{2.706092in}{1.900170in}}%
\pgfpathclose%
\pgfusepath{fill}%
\end{pgfscope}%
\begin{pgfscope}%
\pgfpathrectangle{\pgfqpoint{0.100000in}{0.100000in}}{\pgfqpoint{3.608454in}{2.310000in}}%
\pgfusepath{clip}%
\pgfsetbuttcap%
\pgfsetmiterjoin%
\definecolor{currentfill}{rgb}{0.000000,0.670588,0.664706}%
\pgfsetfillcolor{currentfill}%
\pgfsetlinewidth{0.000000pt}%
\definecolor{currentstroke}{rgb}{0.000000,0.000000,0.000000}%
\pgfsetstrokecolor{currentstroke}%
\pgfsetstrokeopacity{0.000000}%
\pgfsetdash{}{0pt}%
\pgfpathmoveto{\pgfqpoint{2.878161in}{0.725802in}}%
\pgfpathlineto{\pgfqpoint{2.887245in}{0.726366in}}%
\pgfpathlineto{\pgfqpoint{2.888448in}{0.720934in}}%
\pgfpathlineto{\pgfqpoint{2.880133in}{0.713968in}}%
\pgfpathlineto{\pgfqpoint{2.881876in}{0.708207in}}%
\pgfpathlineto{\pgfqpoint{2.876684in}{0.701966in}}%
\pgfpathlineto{\pgfqpoint{2.878191in}{0.698489in}}%
\pgfpathlineto{\pgfqpoint{2.870453in}{0.691569in}}%
\pgfpathlineto{\pgfqpoint{2.866769in}{0.678489in}}%
\pgfpathlineto{\pgfqpoint{2.854948in}{0.675589in}}%
\pgfpathlineto{\pgfqpoint{2.847384in}{0.676389in}}%
\pgfpathlineto{\pgfqpoint{2.841677in}{0.670279in}}%
\pgfpathlineto{\pgfqpoint{2.844033in}{0.665474in}}%
\pgfpathlineto{\pgfqpoint{2.829477in}{0.667841in}}%
\pgfpathlineto{\pgfqpoint{2.797575in}{0.664470in}}%
\pgfpathlineto{\pgfqpoint{2.788514in}{0.665867in}}%
\pgfpathlineto{\pgfqpoint{2.787376in}{0.676778in}}%
\pgfpathlineto{\pgfqpoint{2.769691in}{0.676534in}}%
\pgfpathlineto{\pgfqpoint{2.766984in}{0.705619in}}%
\pgfpathlineto{\pgfqpoint{2.760006in}{0.704933in}}%
\pgfpathlineto{\pgfqpoint{2.760979in}{0.694551in}}%
\pgfpathlineto{\pgfqpoint{2.735873in}{0.692265in}}%
\pgfpathlineto{\pgfqpoint{2.726751in}{0.687235in}}%
\pgfpathlineto{\pgfqpoint{2.735056in}{0.696150in}}%
\pgfpathlineto{\pgfqpoint{2.730084in}{0.707397in}}%
\pgfpathlineto{\pgfqpoint{2.733506in}{0.712730in}}%
\pgfpathlineto{\pgfqpoint{2.721186in}{0.711980in}}%
\pgfpathlineto{\pgfqpoint{2.719091in}{0.734938in}}%
\pgfpathlineto{\pgfqpoint{2.789288in}{0.742878in}}%
\pgfpathlineto{\pgfqpoint{2.794346in}{0.734429in}}%
\pgfpathlineto{\pgfqpoint{2.795289in}{0.729104in}}%
\pgfpathlineto{\pgfqpoint{2.801217in}{0.720980in}}%
\pgfpathlineto{\pgfqpoint{2.834431in}{0.723057in}}%
\pgfpathlineto{\pgfqpoint{2.878161in}{0.725802in}}%
\pgfpathclose%
\pgfusepath{fill}%
\end{pgfscope}%
\begin{pgfscope}%
\pgfpathrectangle{\pgfqpoint{0.100000in}{0.100000in}}{\pgfqpoint{3.608454in}{2.310000in}}%
\pgfusepath{clip}%
\pgfsetbuttcap%
\pgfsetmiterjoin%
\definecolor{currentfill}{rgb}{0.000000,0.431373,0.784314}%
\pgfsetfillcolor{currentfill}%
\pgfsetlinewidth{0.000000pt}%
\definecolor{currentstroke}{rgb}{0.000000,0.000000,0.000000}%
\pgfsetstrokecolor{currentstroke}%
\pgfsetstrokeopacity{0.000000}%
\pgfsetdash{}{0pt}%
\pgfpathmoveto{\pgfqpoint{1.718524in}{1.959750in}}%
\pgfpathlineto{\pgfqpoint{1.719085in}{1.966639in}}%
\pgfpathlineto{\pgfqpoint{1.739705in}{1.965094in}}%
\pgfpathlineto{\pgfqpoint{1.739184in}{1.958172in}}%
\pgfpathlineto{\pgfqpoint{1.762088in}{1.956478in}}%
\pgfpathlineto{\pgfqpoint{1.760656in}{1.935885in}}%
\pgfpathlineto{\pgfqpoint{1.774399in}{1.934946in}}%
\pgfpathlineto{\pgfqpoint{1.775193in}{1.927930in}}%
\pgfpathlineto{\pgfqpoint{1.759368in}{1.916843in}}%
\pgfpathlineto{\pgfqpoint{1.756437in}{1.912852in}}%
\pgfpathlineto{\pgfqpoint{1.744713in}{1.908214in}}%
\pgfpathlineto{\pgfqpoint{1.730862in}{1.910799in}}%
\pgfpathlineto{\pgfqpoint{1.726123in}{1.914246in}}%
\pgfpathlineto{\pgfqpoint{1.720336in}{1.913435in}}%
\pgfpathlineto{\pgfqpoint{1.721787in}{1.931717in}}%
\pgfpathlineto{\pgfqpoint{1.719106in}{1.931926in}}%
\pgfpathlineto{\pgfqpoint{1.721015in}{1.959541in}}%
\pgfpathlineto{\pgfqpoint{1.718524in}{1.959750in}}%
\pgfpathclose%
\pgfusepath{fill}%
\end{pgfscope}%
\begin{pgfscope}%
\pgfpathrectangle{\pgfqpoint{0.100000in}{0.100000in}}{\pgfqpoint{3.608454in}{2.310000in}}%
\pgfusepath{clip}%
\pgfsetbuttcap%
\pgfsetmiterjoin%
\definecolor{currentfill}{rgb}{0.000000,0.364706,0.817647}%
\pgfsetfillcolor{currentfill}%
\pgfsetlinewidth{0.000000pt}%
\definecolor{currentstroke}{rgb}{0.000000,0.000000,0.000000}%
\pgfsetstrokecolor{currentstroke}%
\pgfsetstrokeopacity{0.000000}%
\pgfsetdash{}{0pt}%
\pgfpathmoveto{\pgfqpoint{1.670021in}{1.298287in}}%
\pgfpathlineto{\pgfqpoint{1.670049in}{1.298787in}}%
\pgfpathlineto{\pgfqpoint{1.672634in}{1.332579in}}%
\pgfpathlineto{\pgfqpoint{1.707133in}{1.330203in}}%
\pgfpathlineto{\pgfqpoint{1.756162in}{1.326912in}}%
\pgfpathlineto{\pgfqpoint{1.754213in}{1.292521in}}%
\pgfpathlineto{\pgfqpoint{1.699539in}{1.296124in}}%
\pgfpathlineto{\pgfqpoint{1.670021in}{1.298287in}}%
\pgfpathclose%
\pgfusepath{fill}%
\end{pgfscope}%
\begin{pgfscope}%
\pgfpathrectangle{\pgfqpoint{0.100000in}{0.100000in}}{\pgfqpoint{3.608454in}{2.310000in}}%
\pgfusepath{clip}%
\pgfsetbuttcap%
\pgfsetmiterjoin%
\definecolor{currentfill}{rgb}{0.000000,0.454902,0.772549}%
\pgfsetfillcolor{currentfill}%
\pgfsetlinewidth{0.000000pt}%
\definecolor{currentstroke}{rgb}{0.000000,0.000000,0.000000}%
\pgfsetstrokecolor{currentstroke}%
\pgfsetstrokeopacity{0.000000}%
\pgfsetdash{}{0pt}%
\pgfpathmoveto{\pgfqpoint{2.084392in}{1.676317in}}%
\pgfpathlineto{\pgfqpoint{2.084271in}{1.648958in}}%
\pgfpathlineto{\pgfqpoint{2.056914in}{1.649093in}}%
\pgfpathlineto{\pgfqpoint{2.017541in}{1.649592in}}%
\pgfpathlineto{\pgfqpoint{2.019889in}{1.653420in}}%
\pgfpathlineto{\pgfqpoint{2.019320in}{1.660716in}}%
\pgfpathlineto{\pgfqpoint{2.022652in}{1.663329in}}%
\pgfpathlineto{\pgfqpoint{2.021565in}{1.674298in}}%
\pgfpathlineto{\pgfqpoint{2.018524in}{1.680571in}}%
\pgfpathlineto{\pgfqpoint{2.019001in}{1.688123in}}%
\pgfpathlineto{\pgfqpoint{2.015077in}{1.691199in}}%
\pgfpathlineto{\pgfqpoint{2.014851in}{1.696524in}}%
\pgfpathlineto{\pgfqpoint{2.023239in}{1.696380in}}%
\pgfpathlineto{\pgfqpoint{2.046244in}{1.696054in}}%
\pgfpathlineto{\pgfqpoint{2.057273in}{1.695912in}}%
\pgfpathlineto{\pgfqpoint{2.057034in}{1.676704in}}%
\pgfpathlineto{\pgfqpoint{2.084392in}{1.676317in}}%
\pgfpathclose%
\pgfusepath{fill}%
\end{pgfscope}%
\begin{pgfscope}%
\pgfpathrectangle{\pgfqpoint{0.100000in}{0.100000in}}{\pgfqpoint{3.608454in}{2.310000in}}%
\pgfusepath{clip}%
\pgfsetbuttcap%
\pgfsetmiterjoin%
\definecolor{currentfill}{rgb}{0.000000,0.439216,0.780392}%
\pgfsetfillcolor{currentfill}%
\pgfsetlinewidth{0.000000pt}%
\definecolor{currentstroke}{rgb}{0.000000,0.000000,0.000000}%
\pgfsetstrokecolor{currentstroke}%
\pgfsetstrokeopacity{0.000000}%
\pgfsetdash{}{0pt}%
\pgfpathmoveto{\pgfqpoint{3.420582in}{1.833115in}}%
\pgfpathlineto{\pgfqpoint{3.422865in}{1.818650in}}%
\pgfpathlineto{\pgfqpoint{3.426046in}{1.819395in}}%
\pgfpathlineto{\pgfqpoint{3.430894in}{1.808006in}}%
\pgfpathlineto{\pgfqpoint{3.434992in}{1.802755in}}%
\pgfpathlineto{\pgfqpoint{3.414849in}{1.798333in}}%
\pgfpathlineto{\pgfqpoint{3.377987in}{1.790511in}}%
\pgfpathlineto{\pgfqpoint{3.374916in}{1.805051in}}%
\pgfpathlineto{\pgfqpoint{3.369972in}{1.804097in}}%
\pgfpathlineto{\pgfqpoint{3.367008in}{1.818791in}}%
\pgfpathlineto{\pgfqpoint{3.374319in}{1.820118in}}%
\pgfpathlineto{\pgfqpoint{3.374364in}{1.831752in}}%
\pgfpathlineto{\pgfqpoint{3.382485in}{1.831082in}}%
\pgfpathlineto{\pgfqpoint{3.396386in}{1.835470in}}%
\pgfpathlineto{\pgfqpoint{3.401069in}{1.830482in}}%
\pgfpathlineto{\pgfqpoint{3.412904in}{1.835488in}}%
\pgfpathlineto{\pgfqpoint{3.420582in}{1.833115in}}%
\pgfpathclose%
\pgfusepath{fill}%
\end{pgfscope}%
\begin{pgfscope}%
\pgfpathrectangle{\pgfqpoint{0.100000in}{0.100000in}}{\pgfqpoint{3.608454in}{2.310000in}}%
\pgfusepath{clip}%
\pgfsetbuttcap%
\pgfsetmiterjoin%
\definecolor{currentfill}{rgb}{0.000000,0.403922,0.798039}%
\pgfsetfillcolor{currentfill}%
\pgfsetlinewidth{0.000000pt}%
\definecolor{currentstroke}{rgb}{0.000000,0.000000,0.000000}%
\pgfsetstrokecolor{currentstroke}%
\pgfsetstrokeopacity{0.000000}%
\pgfsetdash{}{0pt}%
\pgfpathmoveto{\pgfqpoint{1.357331in}{0.522521in}}%
\pgfpathlineto{\pgfqpoint{1.350294in}{0.534978in}}%
\pgfpathlineto{\pgfqpoint{1.341731in}{0.539953in}}%
\pgfpathlineto{\pgfqpoint{1.343780in}{0.546075in}}%
\pgfpathlineto{\pgfqpoint{1.344206in}{0.557855in}}%
\pgfpathlineto{\pgfqpoint{1.348048in}{0.565231in}}%
\pgfpathlineto{\pgfqpoint{1.356655in}{0.558900in}}%
\pgfpathlineto{\pgfqpoint{1.368635in}{0.560958in}}%
\pgfpathlineto{\pgfqpoint{1.371871in}{0.559296in}}%
\pgfpathlineto{\pgfqpoint{1.370602in}{0.544014in}}%
\pgfpathlineto{\pgfqpoint{1.366190in}{0.537636in}}%
\pgfpathlineto{\pgfqpoint{1.369907in}{0.525625in}}%
\pgfpathlineto{\pgfqpoint{1.369026in}{0.517805in}}%
\pgfpathlineto{\pgfqpoint{1.357331in}{0.522521in}}%
\pgfpathclose%
\pgfusepath{fill}%
\end{pgfscope}%
\begin{pgfscope}%
\pgfpathrectangle{\pgfqpoint{0.100000in}{0.100000in}}{\pgfqpoint{3.608454in}{2.310000in}}%
\pgfusepath{clip}%
\pgfsetbuttcap%
\pgfsetmiterjoin%
\definecolor{currentfill}{rgb}{0.000000,0.521569,0.739216}%
\pgfsetfillcolor{currentfill}%
\pgfsetlinewidth{0.000000pt}%
\definecolor{currentstroke}{rgb}{0.000000,0.000000,0.000000}%
\pgfsetstrokecolor{currentstroke}%
\pgfsetstrokeopacity{0.000000}%
\pgfsetdash{}{0pt}%
\pgfpathmoveto{\pgfqpoint{0.575043in}{1.085831in}}%
\pgfpathlineto{\pgfqpoint{0.566660in}{1.090402in}}%
\pgfpathlineto{\pgfqpoint{0.566840in}{1.100543in}}%
\pgfpathlineto{\pgfqpoint{0.574437in}{1.095234in}}%
\pgfpathlineto{\pgfqpoint{0.575043in}{1.085831in}}%
\pgfpathclose%
\pgfusepath{fill}%
\end{pgfscope}%
\begin{pgfscope}%
\pgfpathrectangle{\pgfqpoint{0.100000in}{0.100000in}}{\pgfqpoint{3.608454in}{2.310000in}}%
\pgfusepath{clip}%
\pgfsetbuttcap%
\pgfsetmiterjoin%
\definecolor{currentfill}{rgb}{0.000000,0.521569,0.739216}%
\pgfsetfillcolor{currentfill}%
\pgfsetlinewidth{0.000000pt}%
\definecolor{currentstroke}{rgb}{0.000000,0.000000,0.000000}%
\pgfsetstrokecolor{currentstroke}%
\pgfsetstrokeopacity{0.000000}%
\pgfsetdash{}{0pt}%
\pgfpathmoveto{\pgfqpoint{0.623749in}{1.080645in}}%
\pgfpathlineto{\pgfqpoint{0.614676in}{1.095570in}}%
\pgfpathlineto{\pgfqpoint{0.606778in}{1.102859in}}%
\pgfpathlineto{\pgfqpoint{0.598881in}{1.114997in}}%
\pgfpathlineto{\pgfqpoint{0.593720in}{1.119151in}}%
\pgfpathlineto{\pgfqpoint{0.585255in}{1.116816in}}%
\pgfpathlineto{\pgfqpoint{0.578599in}{1.121274in}}%
\pgfpathlineto{\pgfqpoint{0.581839in}{1.128446in}}%
\pgfpathlineto{\pgfqpoint{0.579418in}{1.141321in}}%
\pgfpathlineto{\pgfqpoint{0.576385in}{1.146196in}}%
\pgfpathlineto{\pgfqpoint{0.563289in}{1.149271in}}%
\pgfpathlineto{\pgfqpoint{0.558328in}{1.148636in}}%
\pgfpathlineto{\pgfqpoint{0.535794in}{1.165997in}}%
\pgfpathlineto{\pgfqpoint{0.534238in}{1.176440in}}%
\pgfpathlineto{\pgfqpoint{0.524045in}{1.187970in}}%
\pgfpathlineto{\pgfqpoint{0.527921in}{1.193369in}}%
\pgfpathlineto{\pgfqpoint{0.537333in}{1.227732in}}%
\pgfpathlineto{\pgfqpoint{0.547247in}{1.223234in}}%
\pgfpathlineto{\pgfqpoint{0.547967in}{1.217656in}}%
\pgfpathlineto{\pgfqpoint{0.598059in}{1.204495in}}%
\pgfpathlineto{\pgfqpoint{0.649520in}{1.191532in}}%
\pgfpathlineto{\pgfqpoint{0.668493in}{1.266395in}}%
\pgfpathlineto{\pgfqpoint{0.692238in}{1.260226in}}%
\pgfpathlineto{\pgfqpoint{0.765878in}{1.242208in}}%
\pgfpathlineto{\pgfqpoint{0.786789in}{1.237324in}}%
\pgfpathlineto{\pgfqpoint{0.792611in}{1.237224in}}%
\pgfpathlineto{\pgfqpoint{0.842429in}{1.160488in}}%
\pgfpathlineto{\pgfqpoint{0.840274in}{1.151795in}}%
\pgfpathlineto{\pgfqpoint{0.843543in}{1.141162in}}%
\pgfpathlineto{\pgfqpoint{0.847819in}{1.135796in}}%
\pgfpathlineto{\pgfqpoint{0.848041in}{1.126325in}}%
\pgfpathlineto{\pgfqpoint{0.851945in}{1.113704in}}%
\pgfpathlineto{\pgfqpoint{0.862231in}{1.099514in}}%
\pgfpathlineto{\pgfqpoint{0.868397in}{1.094701in}}%
\pgfpathlineto{\pgfqpoint{0.895182in}{1.086583in}}%
\pgfpathlineto{\pgfqpoint{0.900636in}{1.091652in}}%
\pgfpathlineto{\pgfqpoint{0.914070in}{1.089859in}}%
\pgfpathlineto{\pgfqpoint{0.908989in}{1.065097in}}%
\pgfpathlineto{\pgfqpoint{0.899087in}{1.016739in}}%
\pgfpathlineto{\pgfqpoint{0.858565in}{1.025184in}}%
\pgfpathlineto{\pgfqpoint{0.859997in}{1.031950in}}%
\pgfpathlineto{\pgfqpoint{0.839911in}{1.036150in}}%
\pgfpathlineto{\pgfqpoint{0.832599in}{1.002476in}}%
\pgfpathlineto{\pgfqpoint{0.807337in}{1.008296in}}%
\pgfpathlineto{\pgfqpoint{0.805414in}{1.014538in}}%
\pgfpathlineto{\pgfqpoint{0.809954in}{1.027225in}}%
\pgfpathlineto{\pgfqpoint{0.809266in}{1.037990in}}%
\pgfpathlineto{\pgfqpoint{0.787832in}{1.044715in}}%
\pgfpathlineto{\pgfqpoint{0.721788in}{1.059750in}}%
\pgfpathlineto{\pgfqpoint{0.660882in}{1.074556in}}%
\pgfpathlineto{\pgfqpoint{0.647432in}{1.078308in}}%
\pgfpathlineto{\pgfqpoint{0.640257in}{1.084919in}}%
\pgfpathlineto{\pgfqpoint{0.631612in}{1.088268in}}%
\pgfpathlineto{\pgfqpoint{0.626202in}{1.085453in}}%
\pgfpathlineto{\pgfqpoint{0.623749in}{1.080645in}}%
\pgfpathclose%
\pgfusepath{fill}%
\end{pgfscope}%
\begin{pgfscope}%
\pgfpathrectangle{\pgfqpoint{0.100000in}{0.100000in}}{\pgfqpoint{3.608454in}{2.310000in}}%
\pgfusepath{clip}%
\pgfsetbuttcap%
\pgfsetmiterjoin%
\definecolor{currentfill}{rgb}{0.000000,0.764706,0.617647}%
\pgfsetfillcolor{currentfill}%
\pgfsetlinewidth{0.000000pt}%
\definecolor{currentstroke}{rgb}{0.000000,0.000000,0.000000}%
\pgfsetstrokecolor{currentstroke}%
\pgfsetstrokeopacity{0.000000}%
\pgfsetdash{}{0pt}%
\pgfpathmoveto{\pgfqpoint{2.669844in}{0.772749in}}%
\pgfpathlineto{\pgfqpoint{2.683521in}{0.774139in}}%
\pgfpathlineto{\pgfqpoint{2.682532in}{0.784446in}}%
\pgfpathlineto{\pgfqpoint{2.685966in}{0.784799in}}%
\pgfpathlineto{\pgfqpoint{2.683683in}{0.809309in}}%
\pgfpathlineto{\pgfqpoint{2.685856in}{0.816401in}}%
\pgfpathlineto{\pgfqpoint{2.692745in}{0.817134in}}%
\pgfpathlineto{\pgfqpoint{2.693497in}{0.810278in}}%
\pgfpathlineto{\pgfqpoint{2.714111in}{0.812491in}}%
\pgfpathlineto{\pgfqpoint{2.713352in}{0.819097in}}%
\pgfpathlineto{\pgfqpoint{2.720311in}{0.819497in}}%
\pgfpathlineto{\pgfqpoint{2.721575in}{0.813214in}}%
\pgfpathlineto{\pgfqpoint{2.727912in}{0.813855in}}%
\pgfpathlineto{\pgfqpoint{2.728545in}{0.806963in}}%
\pgfpathlineto{\pgfqpoint{2.737520in}{0.807875in}}%
\pgfpathlineto{\pgfqpoint{2.737974in}{0.800305in}}%
\pgfpathlineto{\pgfqpoint{2.733549in}{0.786465in}}%
\pgfpathlineto{\pgfqpoint{2.706700in}{0.783758in}}%
\pgfpathlineto{\pgfqpoint{2.704075in}{0.776471in}}%
\pgfpathlineto{\pgfqpoint{2.708723in}{0.733977in}}%
\pgfpathlineto{\pgfqpoint{2.674593in}{0.730790in}}%
\pgfpathlineto{\pgfqpoint{2.673739in}{0.731186in}}%
\pgfpathlineto{\pgfqpoint{2.669844in}{0.772749in}}%
\pgfpathclose%
\pgfusepath{fill}%
\end{pgfscope}%
\begin{pgfscope}%
\pgfpathrectangle{\pgfqpoint{0.100000in}{0.100000in}}{\pgfqpoint{3.608454in}{2.310000in}}%
\pgfusepath{clip}%
\pgfsetbuttcap%
\pgfsetmiterjoin%
\definecolor{currentfill}{rgb}{0.000000,0.431373,0.784314}%
\pgfsetfillcolor{currentfill}%
\pgfsetlinewidth{0.000000pt}%
\definecolor{currentstroke}{rgb}{0.000000,0.000000,0.000000}%
\pgfsetstrokecolor{currentstroke}%
\pgfsetstrokeopacity{0.000000}%
\pgfsetdash{}{0pt}%
\pgfpathmoveto{\pgfqpoint{2.420217in}{1.066185in}}%
\pgfpathlineto{\pgfqpoint{2.427918in}{1.068838in}}%
\pgfpathlineto{\pgfqpoint{2.426249in}{1.075405in}}%
\pgfpathlineto{\pgfqpoint{2.434616in}{1.075007in}}%
\pgfpathlineto{\pgfqpoint{2.441604in}{1.083493in}}%
\pgfpathlineto{\pgfqpoint{2.453807in}{1.083497in}}%
\pgfpathlineto{\pgfqpoint{2.461856in}{1.079471in}}%
\pgfpathlineto{\pgfqpoint{2.463288in}{1.065984in}}%
\pgfpathlineto{\pgfqpoint{2.482016in}{1.066685in}}%
\pgfpathlineto{\pgfqpoint{2.483127in}{1.034794in}}%
\pgfpathlineto{\pgfqpoint{2.473129in}{1.034113in}}%
\pgfpathlineto{\pgfqpoint{2.473863in}{1.022295in}}%
\pgfpathlineto{\pgfqpoint{2.477335in}{1.022510in}}%
\pgfpathlineto{\pgfqpoint{2.478692in}{1.001827in}}%
\pgfpathlineto{\pgfqpoint{2.482530in}{0.995135in}}%
\pgfpathlineto{\pgfqpoint{2.467215in}{0.998833in}}%
\pgfpathlineto{\pgfqpoint{2.420148in}{0.996176in}}%
\pgfpathlineto{\pgfqpoint{2.413440in}{0.992269in}}%
\pgfpathlineto{\pgfqpoint{2.413949in}{0.985540in}}%
\pgfpathlineto{\pgfqpoint{2.407574in}{0.985155in}}%
\pgfpathlineto{\pgfqpoint{2.396088in}{0.992607in}}%
\pgfpathlineto{\pgfqpoint{2.394572in}{0.999746in}}%
\pgfpathlineto{\pgfqpoint{2.401872in}{1.004184in}}%
\pgfpathlineto{\pgfqpoint{2.404445in}{1.017531in}}%
\pgfpathlineto{\pgfqpoint{2.415331in}{1.023934in}}%
\pgfpathlineto{\pgfqpoint{2.412041in}{1.029240in}}%
\pgfpathlineto{\pgfqpoint{2.417154in}{1.033431in}}%
\pgfpathlineto{\pgfqpoint{2.419408in}{1.040872in}}%
\pgfpathlineto{\pgfqpoint{2.424139in}{1.041016in}}%
\pgfpathlineto{\pgfqpoint{2.425407in}{1.049147in}}%
\pgfpathlineto{\pgfqpoint{2.420989in}{1.051604in}}%
\pgfpathlineto{\pgfqpoint{2.424647in}{1.062261in}}%
\pgfpathlineto{\pgfqpoint{2.420217in}{1.066185in}}%
\pgfpathclose%
\pgfusepath{fill}%
\end{pgfscope}%
\begin{pgfscope}%
\pgfpathrectangle{\pgfqpoint{0.100000in}{0.100000in}}{\pgfqpoint{3.608454in}{2.310000in}}%
\pgfusepath{clip}%
\pgfsetbuttcap%
\pgfsetmiterjoin%
\definecolor{currentfill}{rgb}{0.000000,0.937255,0.531373}%
\pgfsetfillcolor{currentfill}%
\pgfsetlinewidth{0.000000pt}%
\definecolor{currentstroke}{rgb}{0.000000,0.000000,0.000000}%
\pgfsetstrokecolor{currentstroke}%
\pgfsetstrokeopacity{0.000000}%
\pgfsetdash{}{0pt}%
\pgfpathmoveto{\pgfqpoint{2.826981in}{1.247494in}}%
\pgfpathlineto{\pgfqpoint{2.816644in}{1.252449in}}%
\pgfpathlineto{\pgfqpoint{2.809988in}{1.252552in}}%
\pgfpathlineto{\pgfqpoint{2.802573in}{1.261875in}}%
\pgfpathlineto{\pgfqpoint{2.800586in}{1.266861in}}%
\pgfpathlineto{\pgfqpoint{2.804559in}{1.276722in}}%
\pgfpathlineto{\pgfqpoint{2.810307in}{1.281879in}}%
\pgfpathlineto{\pgfqpoint{2.814875in}{1.291146in}}%
\pgfpathlineto{\pgfqpoint{2.820760in}{1.291911in}}%
\pgfpathlineto{\pgfqpoint{2.825277in}{1.301789in}}%
\pgfpathlineto{\pgfqpoint{2.824999in}{1.309058in}}%
\pgfpathlineto{\pgfqpoint{2.834982in}{1.312343in}}%
\pgfpathlineto{\pgfqpoint{2.840076in}{1.309133in}}%
\pgfpathlineto{\pgfqpoint{2.850698in}{1.310252in}}%
\pgfpathlineto{\pgfqpoint{2.855369in}{1.302588in}}%
\pgfpathlineto{\pgfqpoint{2.847809in}{1.300549in}}%
\pgfpathlineto{\pgfqpoint{2.837862in}{1.289352in}}%
\pgfpathlineto{\pgfqpoint{2.839446in}{1.282704in}}%
\pgfpathlineto{\pgfqpoint{2.849787in}{1.279563in}}%
\pgfpathlineto{\pgfqpoint{2.860524in}{1.271364in}}%
\pgfpathlineto{\pgfqpoint{2.851560in}{1.271717in}}%
\pgfpathlineto{\pgfqpoint{2.850587in}{1.262241in}}%
\pgfpathlineto{\pgfqpoint{2.841576in}{1.260538in}}%
\pgfpathlineto{\pgfqpoint{2.833788in}{1.254755in}}%
\pgfpathlineto{\pgfqpoint{2.830889in}{1.256941in}}%
\pgfpathlineto{\pgfqpoint{2.825042in}{1.252479in}}%
\pgfpathlineto{\pgfqpoint{2.826981in}{1.247494in}}%
\pgfpathclose%
\pgfusepath{fill}%
\end{pgfscope}%
\begin{pgfscope}%
\pgfpathrectangle{\pgfqpoint{0.100000in}{0.100000in}}{\pgfqpoint{3.608454in}{2.310000in}}%
\pgfusepath{clip}%
\pgfsetbuttcap%
\pgfsetmiterjoin%
\definecolor{currentfill}{rgb}{0.000000,0.678431,0.660784}%
\pgfsetfillcolor{currentfill}%
\pgfsetlinewidth{0.000000pt}%
\definecolor{currentstroke}{rgb}{0.000000,0.000000,0.000000}%
\pgfsetstrokecolor{currentstroke}%
\pgfsetstrokeopacity{0.000000}%
\pgfsetdash{}{0pt}%
\pgfpathmoveto{\pgfqpoint{3.127726in}{0.403873in}}%
\pgfpathlineto{\pgfqpoint{3.133900in}{0.367620in}}%
\pgfpathlineto{\pgfqpoint{3.099940in}{0.362168in}}%
\pgfpathlineto{\pgfqpoint{3.093963in}{0.366222in}}%
\pgfpathlineto{\pgfqpoint{3.076512in}{0.366235in}}%
\pgfpathlineto{\pgfqpoint{3.068767in}{0.369764in}}%
\pgfpathlineto{\pgfqpoint{3.063694in}{0.379993in}}%
\pgfpathlineto{\pgfqpoint{3.058788in}{0.396588in}}%
\pgfpathlineto{\pgfqpoint{3.055285in}{0.402609in}}%
\pgfpathlineto{\pgfqpoint{3.046296in}{0.409678in}}%
\pgfpathlineto{\pgfqpoint{3.037249in}{0.409747in}}%
\pgfpathlineto{\pgfqpoint{3.032422in}{0.421403in}}%
\pgfpathlineto{\pgfqpoint{3.036040in}{0.423473in}}%
\pgfpathlineto{\pgfqpoint{3.037010in}{0.432074in}}%
\pgfpathlineto{\pgfqpoint{3.072400in}{0.437158in}}%
\pgfpathlineto{\pgfqpoint{3.075618in}{0.416692in}}%
\pgfpathlineto{\pgfqpoint{3.096468in}{0.420138in}}%
\pgfpathlineto{\pgfqpoint{3.099914in}{0.399053in}}%
\pgfpathlineto{\pgfqpoint{3.127726in}{0.403873in}}%
\pgfpathclose%
\pgfusepath{fill}%
\end{pgfscope}%
\begin{pgfscope}%
\pgfpathrectangle{\pgfqpoint{0.100000in}{0.100000in}}{\pgfqpoint{3.608454in}{2.310000in}}%
\pgfusepath{clip}%
\pgfsetbuttcap%
\pgfsetmiterjoin%
\definecolor{currentfill}{rgb}{0.000000,0.427451,0.786275}%
\pgfsetfillcolor{currentfill}%
\pgfsetlinewidth{0.000000pt}%
\definecolor{currentstroke}{rgb}{0.000000,0.000000,0.000000}%
\pgfsetstrokecolor{currentstroke}%
\pgfsetstrokeopacity{0.000000}%
\pgfsetdash{}{0pt}%
\pgfpathmoveto{\pgfqpoint{0.508309in}{1.308909in}}%
\pgfpathlineto{\pgfqpoint{0.468241in}{1.320678in}}%
\pgfpathlineto{\pgfqpoint{0.438474in}{1.329831in}}%
\pgfpathlineto{\pgfqpoint{0.433387in}{1.338858in}}%
\pgfpathlineto{\pgfqpoint{0.433303in}{1.348727in}}%
\pgfpathlineto{\pgfqpoint{0.430147in}{1.350732in}}%
\pgfpathlineto{\pgfqpoint{0.425066in}{1.366592in}}%
\pgfpathlineto{\pgfqpoint{0.420021in}{1.371609in}}%
\pgfpathlineto{\pgfqpoint{0.416793in}{1.379069in}}%
\pgfpathlineto{\pgfqpoint{0.418273in}{1.400974in}}%
\pgfpathlineto{\pgfqpoint{0.425399in}{1.400805in}}%
\pgfpathlineto{\pgfqpoint{0.429729in}{1.404759in}}%
\pgfpathlineto{\pgfqpoint{0.434644in}{1.414437in}}%
\pgfpathlineto{\pgfqpoint{0.433004in}{1.425553in}}%
\pgfpathlineto{\pgfqpoint{0.418942in}{1.431643in}}%
\pgfpathlineto{\pgfqpoint{0.413490in}{1.439050in}}%
\pgfpathlineto{\pgfqpoint{0.411162in}{1.446533in}}%
\pgfpathlineto{\pgfqpoint{0.411566in}{1.453040in}}%
\pgfpathlineto{\pgfqpoint{0.422155in}{1.452127in}}%
\pgfpathlineto{\pgfqpoint{0.422616in}{1.464232in}}%
\pgfpathlineto{\pgfqpoint{0.424905in}{1.469119in}}%
\pgfpathlineto{\pgfqpoint{0.432540in}{1.470642in}}%
\pgfpathlineto{\pgfqpoint{0.441374in}{1.466010in}}%
\pgfpathlineto{\pgfqpoint{0.445670in}{1.467188in}}%
\pgfpathlineto{\pgfqpoint{0.469193in}{1.459876in}}%
\pgfpathlineto{\pgfqpoint{0.470704in}{1.451769in}}%
\pgfpathlineto{\pgfqpoint{0.465818in}{1.444533in}}%
\pgfpathlineto{\pgfqpoint{0.465842in}{1.434190in}}%
\pgfpathlineto{\pgfqpoint{0.476039in}{1.431016in}}%
\pgfpathlineto{\pgfqpoint{0.475697in}{1.423778in}}%
\pgfpathlineto{\pgfqpoint{0.471004in}{1.413448in}}%
\pgfpathlineto{\pgfqpoint{0.474622in}{1.404992in}}%
\pgfpathlineto{\pgfqpoint{0.486008in}{1.393700in}}%
\pgfpathlineto{\pgfqpoint{0.499634in}{1.368954in}}%
\pgfpathlineto{\pgfqpoint{0.496557in}{1.356685in}}%
\pgfpathlineto{\pgfqpoint{0.492171in}{1.356794in}}%
\pgfpathlineto{\pgfqpoint{0.488481in}{1.345190in}}%
\pgfpathlineto{\pgfqpoint{0.493971in}{1.332372in}}%
\pgfpathlineto{\pgfqpoint{0.498712in}{1.326448in}}%
\pgfpathlineto{\pgfqpoint{0.503087in}{1.324821in}}%
\pgfpathlineto{\pgfqpoint{0.508309in}{1.308909in}}%
\pgfpathclose%
\pgfusepath{fill}%
\end{pgfscope}%
\begin{pgfscope}%
\pgfpathrectangle{\pgfqpoint{0.100000in}{0.100000in}}{\pgfqpoint{3.608454in}{2.310000in}}%
\pgfusepath{clip}%
\pgfsetbuttcap%
\pgfsetmiterjoin%
\definecolor{currentfill}{rgb}{0.000000,0.364706,0.817647}%
\pgfsetfillcolor{currentfill}%
\pgfsetlinewidth{0.000000pt}%
\definecolor{currentstroke}{rgb}{0.000000,0.000000,0.000000}%
\pgfsetstrokecolor{currentstroke}%
\pgfsetstrokeopacity{0.000000}%
\pgfsetdash{}{0pt}%
\pgfpathmoveto{\pgfqpoint{1.082084in}{0.210879in}}%
\pgfpathlineto{\pgfqpoint{1.080487in}{0.211762in}}%
\pgfpathlineto{\pgfqpoint{1.079012in}{0.214678in}}%
\pgfpathlineto{\pgfqpoint{1.073207in}{0.214262in}}%
\pgfpathlineto{\pgfqpoint{1.071517in}{0.212096in}}%
\pgfpathlineto{\pgfqpoint{1.069651in}{0.211374in}}%
\pgfpathlineto{\pgfqpoint{1.067915in}{0.213771in}}%
\pgfpathlineto{\pgfqpoint{1.064056in}{0.213557in}}%
\pgfpathlineto{\pgfqpoint{1.066029in}{0.216398in}}%
\pgfpathlineto{\pgfqpoint{1.068443in}{0.216507in}}%
\pgfpathlineto{\pgfqpoint{1.067670in}{0.218823in}}%
\pgfpathlineto{\pgfqpoint{1.067684in}{0.221839in}}%
\pgfpathlineto{\pgfqpoint{1.068559in}{0.223929in}}%
\pgfpathlineto{\pgfqpoint{1.067133in}{0.226051in}}%
\pgfpathlineto{\pgfqpoint{1.068011in}{0.227410in}}%
\pgfpathlineto{\pgfqpoint{1.070899in}{0.228330in}}%
\pgfpathlineto{\pgfqpoint{1.071501in}{0.231767in}}%
\pgfpathlineto{\pgfqpoint{1.068523in}{0.231093in}}%
\pgfpathlineto{\pgfqpoint{1.068667in}{0.235893in}}%
\pgfpathlineto{\pgfqpoint{1.067104in}{0.238920in}}%
\pgfpathlineto{\pgfqpoint{1.067789in}{0.242382in}}%
\pgfpathlineto{\pgfqpoint{1.067367in}{0.244154in}}%
\pgfpathlineto{\pgfqpoint{1.065528in}{0.244746in}}%
\pgfpathlineto{\pgfqpoint{1.061828in}{0.247605in}}%
\pgfpathlineto{\pgfqpoint{1.061314in}{0.250040in}}%
\pgfpathlineto{\pgfqpoint{1.062351in}{0.251212in}}%
\pgfpathlineto{\pgfqpoint{1.061213in}{0.253041in}}%
\pgfpathlineto{\pgfqpoint{1.062006in}{0.255168in}}%
\pgfpathlineto{\pgfqpoint{1.061710in}{0.260818in}}%
\pgfpathlineto{\pgfqpoint{1.062710in}{0.261351in}}%
\pgfpathlineto{\pgfqpoint{1.063933in}{0.265257in}}%
\pgfpathlineto{\pgfqpoint{1.061849in}{0.265034in}}%
\pgfpathlineto{\pgfqpoint{1.062080in}{0.267893in}}%
\pgfpathlineto{\pgfqpoint{1.064059in}{0.276502in}}%
\pgfpathlineto{\pgfqpoint{1.064827in}{0.282112in}}%
\pgfpathlineto{\pgfqpoint{1.062541in}{0.281407in}}%
\pgfpathlineto{\pgfqpoint{1.062850in}{0.278917in}}%
\pgfpathlineto{\pgfqpoint{1.061375in}{0.278552in}}%
\pgfpathlineto{\pgfqpoint{1.061437in}{0.275880in}}%
\pgfpathlineto{\pgfqpoint{1.059564in}{0.269832in}}%
\pgfpathlineto{\pgfqpoint{1.059156in}{0.263608in}}%
\pgfpathlineto{\pgfqpoint{1.057760in}{0.259141in}}%
\pgfpathlineto{\pgfqpoint{1.056698in}{0.253549in}}%
\pgfpathlineto{\pgfqpoint{1.054233in}{0.250993in}}%
\pgfpathlineto{\pgfqpoint{1.052026in}{0.254747in}}%
\pgfpathlineto{\pgfqpoint{1.052338in}{0.257691in}}%
\pgfpathlineto{\pgfqpoint{1.050359in}{0.258935in}}%
\pgfpathlineto{\pgfqpoint{1.045539in}{0.260432in}}%
\pgfpathlineto{\pgfqpoint{1.048292in}{0.263786in}}%
\pgfpathlineto{\pgfqpoint{1.048887in}{0.269049in}}%
\pgfpathlineto{\pgfqpoint{1.048759in}{0.271883in}}%
\pgfpathlineto{\pgfqpoint{1.045964in}{0.271941in}}%
\pgfpathlineto{\pgfqpoint{1.044789in}{0.274978in}}%
\pgfpathlineto{\pgfqpoint{1.043552in}{0.275786in}}%
\pgfpathlineto{\pgfqpoint{1.044076in}{0.279452in}}%
\pgfpathlineto{\pgfqpoint{1.040938in}{0.280527in}}%
\pgfpathlineto{\pgfqpoint{1.038980in}{0.283005in}}%
\pgfpathlineto{\pgfqpoint{1.037883in}{0.281325in}}%
\pgfpathlineto{\pgfqpoint{1.040576in}{0.278511in}}%
\pgfpathlineto{\pgfqpoint{1.040732in}{0.274755in}}%
\pgfpathlineto{\pgfqpoint{1.043542in}{0.270946in}}%
\pgfpathlineto{\pgfqpoint{1.040999in}{0.270087in}}%
\pgfpathlineto{\pgfqpoint{1.042205in}{0.268402in}}%
\pgfpathlineto{\pgfqpoint{1.044236in}{0.268878in}}%
\pgfpathlineto{\pgfqpoint{1.043891in}{0.261248in}}%
\pgfpathlineto{\pgfqpoint{1.039661in}{0.261025in}}%
\pgfpathlineto{\pgfqpoint{1.040400in}{0.263502in}}%
\pgfpathlineto{\pgfqpoint{1.037395in}{0.262751in}}%
\pgfpathlineto{\pgfqpoint{1.033326in}{0.260670in}}%
\pgfpathlineto{\pgfqpoint{1.033334in}{0.264446in}}%
\pgfpathlineto{\pgfqpoint{1.032173in}{0.267309in}}%
\pgfpathlineto{\pgfqpoint{1.029605in}{0.268061in}}%
\pgfpathlineto{\pgfqpoint{1.025973in}{0.276463in}}%
\pgfpathlineto{\pgfqpoint{1.026377in}{0.278219in}}%
\pgfpathlineto{\pgfqpoint{1.024142in}{0.282697in}}%
\pgfpathlineto{\pgfqpoint{1.024419in}{0.288147in}}%
\pgfpathlineto{\pgfqpoint{1.022440in}{0.293955in}}%
\pgfpathlineto{\pgfqpoint{1.018989in}{0.298401in}}%
\pgfpathlineto{\pgfqpoint{1.017453in}{0.301716in}}%
\pgfpathlineto{\pgfqpoint{1.012827in}{0.307991in}}%
\pgfpathlineto{\pgfqpoint{1.008318in}{0.316535in}}%
\pgfpathlineto{\pgfqpoint{1.011237in}{0.316069in}}%
\pgfpathlineto{\pgfqpoint{1.015372in}{0.318061in}}%
\pgfpathlineto{\pgfqpoint{1.015809in}{0.323540in}}%
\pgfpathlineto{\pgfqpoint{1.016829in}{0.325554in}}%
\pgfpathlineto{\pgfqpoint{1.011575in}{0.323679in}}%
\pgfpathlineto{\pgfqpoint{1.009826in}{0.326603in}}%
\pgfpathlineto{\pgfqpoint{1.006161in}{0.325696in}}%
\pgfpathlineto{\pgfqpoint{1.005156in}{0.323648in}}%
\pgfpathlineto{\pgfqpoint{1.000051in}{0.326986in}}%
\pgfpathlineto{\pgfqpoint{0.997055in}{0.330455in}}%
\pgfpathlineto{\pgfqpoint{1.004608in}{0.344079in}}%
\pgfpathlineto{\pgfqpoint{1.009161in}{0.338796in}}%
\pgfpathlineto{\pgfqpoint{1.012055in}{0.339391in}}%
\pgfpathlineto{\pgfqpoint{1.015589in}{0.334105in}}%
\pgfpathlineto{\pgfqpoint{1.021713in}{0.335942in}}%
\pgfpathlineto{\pgfqpoint{1.027923in}{0.333090in}}%
\pgfpathlineto{\pgfqpoint{1.029124in}{0.331408in}}%
\pgfpathlineto{\pgfqpoint{1.024345in}{0.326303in}}%
\pgfpathlineto{\pgfqpoint{1.025043in}{0.322901in}}%
\pgfpathlineto{\pgfqpoint{1.028234in}{0.318423in}}%
\pgfpathlineto{\pgfqpoint{1.027542in}{0.314145in}}%
\pgfpathlineto{\pgfqpoint{1.033899in}{0.293981in}}%
\pgfpathlineto{\pgfqpoint{1.032438in}{0.287802in}}%
\pgfpathlineto{\pgfqpoint{1.030648in}{0.284485in}}%
\pgfpathlineto{\pgfqpoint{1.031736in}{0.284026in}}%
\pgfpathlineto{\pgfqpoint{1.035421in}{0.285214in}}%
\pgfpathlineto{\pgfqpoint{1.041713in}{0.285984in}}%
\pgfpathlineto{\pgfqpoint{1.046557in}{0.284922in}}%
\pgfpathlineto{\pgfqpoint{1.049608in}{0.287280in}}%
\pgfpathlineto{\pgfqpoint{1.051192in}{0.290700in}}%
\pgfpathlineto{\pgfqpoint{1.054086in}{0.290900in}}%
\pgfpathlineto{\pgfqpoint{1.058160in}{0.294663in}}%
\pgfpathlineto{\pgfqpoint{1.061301in}{0.293799in}}%
\pgfpathlineto{\pgfqpoint{1.068256in}{0.294295in}}%
\pgfpathlineto{\pgfqpoint{1.069607in}{0.293456in}}%
\pgfpathlineto{\pgfqpoint{1.071897in}{0.287157in}}%
\pgfpathlineto{\pgfqpoint{1.071746in}{0.284296in}}%
\pgfpathlineto{\pgfqpoint{1.069635in}{0.282181in}}%
\pgfpathlineto{\pgfqpoint{1.069594in}{0.277534in}}%
\pgfpathlineto{\pgfqpoint{1.072600in}{0.274944in}}%
\pgfpathlineto{\pgfqpoint{1.072637in}{0.271740in}}%
\pgfpathlineto{\pgfqpoint{1.074084in}{0.269420in}}%
\pgfpathlineto{\pgfqpoint{1.073539in}{0.266568in}}%
\pgfpathlineto{\pgfqpoint{1.073899in}{0.262067in}}%
\pgfpathlineto{\pgfqpoint{1.077857in}{0.256052in}}%
\pgfpathlineto{\pgfqpoint{1.078297in}{0.250259in}}%
\pgfpathlineto{\pgfqpoint{1.080003in}{0.246871in}}%
\pgfpathlineto{\pgfqpoint{1.078915in}{0.244757in}}%
\pgfpathlineto{\pgfqpoint{1.079981in}{0.235727in}}%
\pgfpathlineto{\pgfqpoint{1.079690in}{0.232569in}}%
\pgfpathlineto{\pgfqpoint{1.080751in}{0.226649in}}%
\pgfpathlineto{\pgfqpoint{1.081182in}{0.216121in}}%
\pgfpathlineto{\pgfqpoint{1.082084in}{0.210879in}}%
\pgfpathclose%
\pgfusepath{fill}%
\end{pgfscope}%
\begin{pgfscope}%
\pgfpathrectangle{\pgfqpoint{0.100000in}{0.100000in}}{\pgfqpoint{3.608454in}{2.310000in}}%
\pgfusepath{clip}%
\pgfsetbuttcap%
\pgfsetmiterjoin%
\definecolor{currentfill}{rgb}{0.000000,0.364706,0.817647}%
\pgfsetfillcolor{currentfill}%
\pgfsetlinewidth{0.000000pt}%
\definecolor{currentstroke}{rgb}{0.000000,0.000000,0.000000}%
\pgfsetstrokecolor{currentstroke}%
\pgfsetstrokeopacity{0.000000}%
\pgfsetdash{}{0pt}%
\pgfpathmoveto{\pgfqpoint{1.046772in}{0.216481in}}%
\pgfpathlineto{\pgfqpoint{1.049473in}{0.221987in}}%
\pgfpathlineto{\pgfqpoint{1.051491in}{0.223563in}}%
\pgfpathlineto{\pgfqpoint{1.053235in}{0.226047in}}%
\pgfpathlineto{\pgfqpoint{1.053669in}{0.234273in}}%
\pgfpathlineto{\pgfqpoint{1.055631in}{0.239433in}}%
\pgfpathlineto{\pgfqpoint{1.056457in}{0.243986in}}%
\pgfpathlineto{\pgfqpoint{1.058002in}{0.245736in}}%
\pgfpathlineto{\pgfqpoint{1.057331in}{0.248039in}}%
\pgfpathlineto{\pgfqpoint{1.057875in}{0.253294in}}%
\pgfpathlineto{\pgfqpoint{1.059869in}{0.251052in}}%
\pgfpathlineto{\pgfqpoint{1.059422in}{0.246810in}}%
\pgfpathlineto{\pgfqpoint{1.061246in}{0.245939in}}%
\pgfpathlineto{\pgfqpoint{1.065594in}{0.242418in}}%
\pgfpathlineto{\pgfqpoint{1.066004in}{0.236870in}}%
\pgfpathlineto{\pgfqpoint{1.065636in}{0.230146in}}%
\pgfpathlineto{\pgfqpoint{1.064005in}{0.231270in}}%
\pgfpathlineto{\pgfqpoint{1.064341in}{0.239679in}}%
\pgfpathlineto{\pgfqpoint{1.062886in}{0.239330in}}%
\pgfpathlineto{\pgfqpoint{1.060924in}{0.235715in}}%
\pgfpathlineto{\pgfqpoint{1.062789in}{0.235055in}}%
\pgfpathlineto{\pgfqpoint{1.062708in}{0.232793in}}%
\pgfpathlineto{\pgfqpoint{1.061477in}{0.229384in}}%
\pgfpathlineto{\pgfqpoint{1.062971in}{0.228484in}}%
\pgfpathlineto{\pgfqpoint{1.062440in}{0.223965in}}%
\pgfpathlineto{\pgfqpoint{1.059310in}{0.224988in}}%
\pgfpathlineto{\pgfqpoint{1.061475in}{0.221010in}}%
\pgfpathlineto{\pgfqpoint{1.059297in}{0.219771in}}%
\pgfpathlineto{\pgfqpoint{1.058337in}{0.221103in}}%
\pgfpathlineto{\pgfqpoint{1.056457in}{0.219297in}}%
\pgfpathlineto{\pgfqpoint{1.053559in}{0.218893in}}%
\pgfpathlineto{\pgfqpoint{1.050357in}{0.216700in}}%
\pgfpathlineto{\pgfqpoint{1.046772in}{0.216481in}}%
\pgfpathclose%
\pgfusepath{fill}%
\end{pgfscope}%
\begin{pgfscope}%
\pgfpathrectangle{\pgfqpoint{0.100000in}{0.100000in}}{\pgfqpoint{3.608454in}{2.310000in}}%
\pgfusepath{clip}%
\pgfsetbuttcap%
\pgfsetmiterjoin%
\definecolor{currentfill}{rgb}{0.000000,0.364706,0.817647}%
\pgfsetfillcolor{currentfill}%
\pgfsetlinewidth{0.000000pt}%
\definecolor{currentstroke}{rgb}{0.000000,0.000000,0.000000}%
\pgfsetstrokecolor{currentstroke}%
\pgfsetstrokeopacity{0.000000}%
\pgfsetdash{}{0pt}%
\pgfpathmoveto{\pgfqpoint{1.031898in}{0.249479in}}%
\pgfpathlineto{\pgfqpoint{1.031917in}{0.252366in}}%
\pgfpathlineto{\pgfqpoint{1.033185in}{0.256544in}}%
\pgfpathlineto{\pgfqpoint{1.037420in}{0.256450in}}%
\pgfpathlineto{\pgfqpoint{1.039053in}{0.255636in}}%
\pgfpathlineto{\pgfqpoint{1.040317in}{0.257461in}}%
\pgfpathlineto{\pgfqpoint{1.041757in}{0.255742in}}%
\pgfpathlineto{\pgfqpoint{1.045494in}{0.257061in}}%
\pgfpathlineto{\pgfqpoint{1.048001in}{0.253354in}}%
\pgfpathlineto{\pgfqpoint{1.048579in}{0.249693in}}%
\pgfpathlineto{\pgfqpoint{1.052991in}{0.247043in}}%
\pgfpathlineto{\pgfqpoint{1.054422in}{0.245099in}}%
\pgfpathlineto{\pgfqpoint{1.053698in}{0.242447in}}%
\pgfpathlineto{\pgfqpoint{1.052033in}{0.241302in}}%
\pgfpathlineto{\pgfqpoint{1.050384in}{0.238572in}}%
\pgfpathlineto{\pgfqpoint{1.047084in}{0.239857in}}%
\pgfpathlineto{\pgfqpoint{1.044554in}{0.244818in}}%
\pgfpathlineto{\pgfqpoint{1.042604in}{0.247304in}}%
\pgfpathlineto{\pgfqpoint{1.041099in}{0.250574in}}%
\pgfpathlineto{\pgfqpoint{1.037351in}{0.246961in}}%
\pgfpathlineto{\pgfqpoint{1.031898in}{0.249479in}}%
\pgfpathclose%
\pgfusepath{fill}%
\end{pgfscope}%
\begin{pgfscope}%
\pgfpathrectangle{\pgfqpoint{0.100000in}{0.100000in}}{\pgfqpoint{3.608454in}{2.310000in}}%
\pgfusepath{clip}%
\pgfsetbuttcap%
\pgfsetmiterjoin%
\definecolor{currentfill}{rgb}{0.000000,0.482353,0.758824}%
\pgfsetfillcolor{currentfill}%
\pgfsetlinewidth{0.000000pt}%
\definecolor{currentstroke}{rgb}{0.000000,0.000000,0.000000}%
\pgfsetstrokecolor{currentstroke}%
\pgfsetstrokeopacity{0.000000}%
\pgfsetdash{}{0pt}%
\pgfpathmoveto{\pgfqpoint{3.344220in}{1.667169in}}%
\pgfpathlineto{\pgfqpoint{3.344000in}{1.680413in}}%
\pgfpathlineto{\pgfqpoint{3.366550in}{1.688953in}}%
\pgfpathlineto{\pgfqpoint{3.369634in}{1.670838in}}%
\pgfpathlineto{\pgfqpoint{3.375111in}{1.665432in}}%
\pgfpathlineto{\pgfqpoint{3.362988in}{1.653409in}}%
\pgfpathlineto{\pgfqpoint{3.369189in}{1.645562in}}%
\pgfpathlineto{\pgfqpoint{3.372252in}{1.639364in}}%
\pgfpathlineto{\pgfqpoint{3.379684in}{1.642975in}}%
\pgfpathlineto{\pgfqpoint{3.395641in}{1.645504in}}%
\pgfpathlineto{\pgfqpoint{3.404650in}{1.652435in}}%
\pgfpathlineto{\pgfqpoint{3.420950in}{1.656635in}}%
\pgfpathlineto{\pgfqpoint{3.428750in}{1.659772in}}%
\pgfpathlineto{\pgfqpoint{3.439665in}{1.671536in}}%
\pgfpathlineto{\pgfqpoint{3.441705in}{1.675928in}}%
\pgfpathlineto{\pgfqpoint{3.448912in}{1.669829in}}%
\pgfpathlineto{\pgfqpoint{3.454605in}{1.672223in}}%
\pgfpathlineto{\pgfqpoint{3.458853in}{1.668402in}}%
\pgfpathlineto{\pgfqpoint{3.465859in}{1.676647in}}%
\pgfpathlineto{\pgfqpoint{3.471136in}{1.676378in}}%
\pgfpathlineto{\pgfqpoint{3.438928in}{1.650105in}}%
\pgfpathlineto{\pgfqpoint{3.418976in}{1.636559in}}%
\pgfpathlineto{\pgfqpoint{3.411351in}{1.630261in}}%
\pgfpathlineto{\pgfqpoint{3.399067in}{1.623101in}}%
\pgfpathlineto{\pgfqpoint{3.394098in}{1.622800in}}%
\pgfpathlineto{\pgfqpoint{3.382218in}{1.615957in}}%
\pgfpathlineto{\pgfqpoint{3.369738in}{1.613583in}}%
\pgfpathlineto{\pgfqpoint{3.361694in}{1.608438in}}%
\pgfpathlineto{\pgfqpoint{3.353119in}{1.610609in}}%
\pgfpathlineto{\pgfqpoint{3.350697in}{1.605499in}}%
\pgfpathlineto{\pgfqpoint{3.343167in}{1.599400in}}%
\pgfpathlineto{\pgfqpoint{3.344569in}{1.612129in}}%
\pgfpathlineto{\pgfqpoint{3.352075in}{1.614317in}}%
\pgfpathlineto{\pgfqpoint{3.355371in}{1.641928in}}%
\pgfpathlineto{\pgfqpoint{3.353043in}{1.654103in}}%
\pgfpathlineto{\pgfqpoint{3.346226in}{1.661401in}}%
\pgfpathlineto{\pgfqpoint{3.344220in}{1.667169in}}%
\pgfpathclose%
\pgfusepath{fill}%
\end{pgfscope}%
\begin{pgfscope}%
\pgfpathrectangle{\pgfqpoint{0.100000in}{0.100000in}}{\pgfqpoint{3.608454in}{2.310000in}}%
\pgfusepath{clip}%
\pgfsetbuttcap%
\pgfsetmiterjoin%
\definecolor{currentfill}{rgb}{0.000000,0.478431,0.760784}%
\pgfsetfillcolor{currentfill}%
\pgfsetlinewidth{0.000000pt}%
\definecolor{currentstroke}{rgb}{0.000000,0.000000,0.000000}%
\pgfsetstrokecolor{currentstroke}%
\pgfsetstrokeopacity{0.000000}%
\pgfsetdash{}{0pt}%
\pgfpathmoveto{\pgfqpoint{1.157736in}{1.964285in}}%
\pgfpathlineto{\pgfqpoint{1.146770in}{1.973621in}}%
\pgfpathlineto{\pgfqpoint{1.147664in}{1.994687in}}%
\pgfpathlineto{\pgfqpoint{1.143228in}{2.002013in}}%
\pgfpathlineto{\pgfqpoint{1.145294in}{2.011278in}}%
\pgfpathlineto{\pgfqpoint{1.148015in}{2.013748in}}%
\pgfpathlineto{\pgfqpoint{1.158552in}{2.015787in}}%
\pgfpathlineto{\pgfqpoint{1.162054in}{2.020418in}}%
\pgfpathlineto{\pgfqpoint{1.164980in}{2.032458in}}%
\pgfpathlineto{\pgfqpoint{1.160126in}{2.041022in}}%
\pgfpathlineto{\pgfqpoint{1.153421in}{2.042231in}}%
\pgfpathlineto{\pgfqpoint{1.155484in}{2.052872in}}%
\pgfpathlineto{\pgfqpoint{1.141533in}{2.055406in}}%
\pgfpathlineto{\pgfqpoint{1.146831in}{2.082471in}}%
\pgfpathlineto{\pgfqpoint{1.132903in}{2.085429in}}%
\pgfpathlineto{\pgfqpoint{1.137742in}{2.109618in}}%
\pgfpathlineto{\pgfqpoint{1.134932in}{2.119015in}}%
\pgfpathlineto{\pgfqpoint{1.135799in}{2.128987in}}%
\pgfpathlineto{\pgfqpoint{1.140741in}{2.130368in}}%
\pgfpathlineto{\pgfqpoint{1.143835in}{2.140807in}}%
\pgfpathlineto{\pgfqpoint{1.149026in}{2.145360in}}%
\pgfpathlineto{\pgfqpoint{1.150853in}{2.134585in}}%
\pgfpathlineto{\pgfqpoint{1.148672in}{2.124295in}}%
\pgfpathlineto{\pgfqpoint{1.152365in}{2.116451in}}%
\pgfpathlineto{\pgfqpoint{1.165476in}{2.115944in}}%
\pgfpathlineto{\pgfqpoint{1.173281in}{2.111946in}}%
\pgfpathlineto{\pgfqpoint{1.173331in}{2.107985in}}%
\pgfpathlineto{\pgfqpoint{1.179403in}{2.102869in}}%
\pgfpathlineto{\pgfqpoint{1.191810in}{2.101317in}}%
\pgfpathlineto{\pgfqpoint{1.187119in}{2.075916in}}%
\pgfpathlineto{\pgfqpoint{1.192015in}{2.070842in}}%
\pgfpathlineto{\pgfqpoint{1.199816in}{2.068342in}}%
\pgfpathlineto{\pgfqpoint{1.196622in}{2.051472in}}%
\pgfpathlineto{\pgfqpoint{1.203686in}{2.050112in}}%
\pgfpathlineto{\pgfqpoint{1.203339in}{2.044520in}}%
\pgfpathlineto{\pgfqpoint{1.210034in}{2.040419in}}%
\pgfpathlineto{\pgfqpoint{1.213019in}{2.026776in}}%
\pgfpathlineto{\pgfqpoint{1.216917in}{2.019801in}}%
\pgfpathlineto{\pgfqpoint{1.219896in}{2.007840in}}%
\pgfpathlineto{\pgfqpoint{1.224493in}{2.008821in}}%
\pgfpathlineto{\pgfqpoint{1.228711in}{2.003790in}}%
\pgfpathlineto{\pgfqpoint{1.224264in}{1.998009in}}%
\pgfpathlineto{\pgfqpoint{1.225529in}{1.987742in}}%
\pgfpathlineto{\pgfqpoint{1.210052in}{1.990143in}}%
\pgfpathlineto{\pgfqpoint{1.208958in}{1.985065in}}%
\pgfpathlineto{\pgfqpoint{1.203995in}{1.980934in}}%
\pgfpathlineto{\pgfqpoint{1.199884in}{1.972842in}}%
\pgfpathlineto{\pgfqpoint{1.192609in}{1.969852in}}%
\pgfpathlineto{\pgfqpoint{1.185357in}{1.963304in}}%
\pgfpathlineto{\pgfqpoint{1.176413in}{1.966592in}}%
\pgfpathlineto{\pgfqpoint{1.173618in}{1.970061in}}%
\pgfpathlineto{\pgfqpoint{1.165857in}{1.971182in}}%
\pgfpathlineto{\pgfqpoint{1.157736in}{1.964285in}}%
\pgfpathclose%
\pgfusepath{fill}%
\end{pgfscope}%
\begin{pgfscope}%
\pgfpathrectangle{\pgfqpoint{0.100000in}{0.100000in}}{\pgfqpoint{3.608454in}{2.310000in}}%
\pgfusepath{clip}%
\pgfsetbuttcap%
\pgfsetmiterjoin%
\definecolor{currentfill}{rgb}{0.000000,0.372549,0.813725}%
\pgfsetfillcolor{currentfill}%
\pgfsetlinewidth{0.000000pt}%
\definecolor{currentstroke}{rgb}{0.000000,0.000000,0.000000}%
\pgfsetstrokecolor{currentstroke}%
\pgfsetstrokeopacity{0.000000}%
\pgfsetdash{}{0pt}%
\pgfpathmoveto{\pgfqpoint{1.964117in}{1.448246in}}%
\pgfpathlineto{\pgfqpoint{1.936670in}{1.449037in}}%
\pgfpathlineto{\pgfqpoint{1.937428in}{1.476563in}}%
\pgfpathlineto{\pgfqpoint{1.938437in}{1.514106in}}%
\pgfpathlineto{\pgfqpoint{1.942299in}{1.518835in}}%
\pgfpathlineto{\pgfqpoint{1.952480in}{1.526297in}}%
\pgfpathlineto{\pgfqpoint{1.966281in}{1.530825in}}%
\pgfpathlineto{\pgfqpoint{1.965564in}{1.503301in}}%
\pgfpathlineto{\pgfqpoint{1.964117in}{1.448246in}}%
\pgfpathclose%
\pgfusepath{fill}%
\end{pgfscope}%
\begin{pgfscope}%
\pgfpathrectangle{\pgfqpoint{0.100000in}{0.100000in}}{\pgfqpoint{3.608454in}{2.310000in}}%
\pgfusepath{clip}%
\pgfsetbuttcap%
\pgfsetmiterjoin%
\definecolor{currentfill}{rgb}{0.000000,0.678431,0.660784}%
\pgfsetfillcolor{currentfill}%
\pgfsetlinewidth{0.000000pt}%
\definecolor{currentstroke}{rgb}{0.000000,0.000000,0.000000}%
\pgfsetstrokecolor{currentstroke}%
\pgfsetstrokeopacity{0.000000}%
\pgfsetdash{}{0pt}%
\pgfpathmoveto{\pgfqpoint{0.737189in}{2.001359in}}%
\pgfpathlineto{\pgfqpoint{0.734618in}{1.987950in}}%
\pgfpathlineto{\pgfqpoint{0.723444in}{1.948508in}}%
\pgfpathlineto{\pgfqpoint{0.710132in}{1.952205in}}%
\pgfpathlineto{\pgfqpoint{0.711901in}{1.958542in}}%
\pgfpathlineto{\pgfqpoint{0.706253in}{1.964177in}}%
\pgfpathlineto{\pgfqpoint{0.686387in}{1.969750in}}%
\pgfpathlineto{\pgfqpoint{0.693514in}{1.994463in}}%
\pgfpathlineto{\pgfqpoint{0.695564in}{1.998853in}}%
\pgfpathlineto{\pgfqpoint{0.692032in}{2.003466in}}%
\pgfpathlineto{\pgfqpoint{0.692055in}{2.010580in}}%
\pgfpathlineto{\pgfqpoint{0.697834in}{2.026472in}}%
\pgfpathlineto{\pgfqpoint{0.695191in}{2.034615in}}%
\pgfpathlineto{\pgfqpoint{0.702993in}{2.049946in}}%
\pgfpathlineto{\pgfqpoint{0.710674in}{2.050338in}}%
\pgfpathlineto{\pgfqpoint{0.710428in}{2.058581in}}%
\pgfpathlineto{\pgfqpoint{0.707552in}{2.064386in}}%
\pgfpathlineto{\pgfqpoint{0.713071in}{2.066429in}}%
\pgfpathlineto{\pgfqpoint{0.723978in}{2.065521in}}%
\pgfpathlineto{\pgfqpoint{0.737141in}{2.068948in}}%
\pgfpathlineto{\pgfqpoint{0.724757in}{2.026561in}}%
\pgfpathlineto{\pgfqpoint{0.731442in}{2.024554in}}%
\pgfpathlineto{\pgfqpoint{0.729534in}{2.017963in}}%
\pgfpathlineto{\pgfqpoint{0.736122in}{2.015946in}}%
\pgfpathlineto{\pgfqpoint{0.730749in}{2.003174in}}%
\pgfpathlineto{\pgfqpoint{0.737189in}{2.001359in}}%
\pgfpathclose%
\pgfusepath{fill}%
\end{pgfscope}%
\begin{pgfscope}%
\pgfpathrectangle{\pgfqpoint{0.100000in}{0.100000in}}{\pgfqpoint{3.608454in}{2.310000in}}%
\pgfusepath{clip}%
\pgfsetbuttcap%
\pgfsetmiterjoin%
\definecolor{currentfill}{rgb}{0.000000,0.521569,0.739216}%
\pgfsetfillcolor{currentfill}%
\pgfsetlinewidth{0.000000pt}%
\definecolor{currentstroke}{rgb}{0.000000,0.000000,0.000000}%
\pgfsetstrokecolor{currentstroke}%
\pgfsetstrokeopacity{0.000000}%
\pgfsetdash{}{0pt}%
\pgfpathmoveto{\pgfqpoint{1.792481in}{1.427797in}}%
\pgfpathlineto{\pgfqpoint{1.724173in}{1.432221in}}%
\pgfpathlineto{\pgfqpoint{1.726124in}{1.459646in}}%
\pgfpathlineto{\pgfqpoint{1.682362in}{1.462771in}}%
\pgfpathlineto{\pgfqpoint{1.684512in}{1.490228in}}%
\pgfpathlineto{\pgfqpoint{1.732544in}{1.486727in}}%
\pgfpathlineto{\pgfqpoint{1.760676in}{1.485050in}}%
\pgfpathlineto{\pgfqpoint{1.758946in}{1.457374in}}%
\pgfpathlineto{\pgfqpoint{1.793800in}{1.455314in}}%
\pgfpathlineto{\pgfqpoint{1.792481in}{1.427797in}}%
\pgfpathclose%
\pgfusepath{fill}%
\end{pgfscope}%
\begin{pgfscope}%
\pgfpathrectangle{\pgfqpoint{0.100000in}{0.100000in}}{\pgfqpoint{3.608454in}{2.310000in}}%
\pgfusepath{clip}%
\pgfsetbuttcap%
\pgfsetmiterjoin%
\definecolor{currentfill}{rgb}{0.000000,0.505882,0.747059}%
\pgfsetfillcolor{currentfill}%
\pgfsetlinewidth{0.000000pt}%
\definecolor{currentstroke}{rgb}{0.000000,0.000000,0.000000}%
\pgfsetstrokecolor{currentstroke}%
\pgfsetstrokeopacity{0.000000}%
\pgfsetdash{}{0pt}%
\pgfpathmoveto{\pgfqpoint{1.223167in}{1.249876in}}%
\pgfpathlineto{\pgfqpoint{1.163571in}{1.259392in}}%
\pgfpathlineto{\pgfqpoint{1.116943in}{1.267777in}}%
\pgfpathlineto{\pgfqpoint{1.078577in}{1.274542in}}%
\pgfpathlineto{\pgfqpoint{1.087638in}{1.275829in}}%
\pgfpathlineto{\pgfqpoint{1.091741in}{1.280236in}}%
\pgfpathlineto{\pgfqpoint{1.099188in}{1.278642in}}%
\pgfpathlineto{\pgfqpoint{1.109948in}{1.282813in}}%
\pgfpathlineto{\pgfqpoint{1.113627in}{1.291412in}}%
\pgfpathlineto{\pgfqpoint{1.122123in}{1.294263in}}%
\pgfpathlineto{\pgfqpoint{1.130720in}{1.308550in}}%
\pgfpathlineto{\pgfqpoint{1.134341in}{1.312779in}}%
\pgfpathlineto{\pgfqpoint{1.141396in}{1.315621in}}%
\pgfpathlineto{\pgfqpoint{1.150177in}{1.332380in}}%
\pgfpathlineto{\pgfqpoint{1.159137in}{1.332405in}}%
\pgfpathlineto{\pgfqpoint{1.168304in}{1.339366in}}%
\pgfpathlineto{\pgfqpoint{1.170873in}{1.338075in}}%
\pgfpathlineto{\pgfqpoint{1.177278in}{1.346155in}}%
\pgfpathlineto{\pgfqpoint{1.183026in}{1.348649in}}%
\pgfpathlineto{\pgfqpoint{1.186021in}{1.352257in}}%
\pgfpathlineto{\pgfqpoint{1.183516in}{1.359664in}}%
\pgfpathlineto{\pgfqpoint{1.181575in}{1.376785in}}%
\pgfpathlineto{\pgfqpoint{1.240533in}{1.367269in}}%
\pgfpathlineto{\pgfqpoint{1.237792in}{1.349736in}}%
\pgfpathlineto{\pgfqpoint{1.237566in}{1.340907in}}%
\pgfpathlineto{\pgfqpoint{1.234133in}{1.318758in}}%
\pgfpathlineto{\pgfqpoint{1.248158in}{1.317741in}}%
\pgfpathlineto{\pgfqpoint{1.282717in}{1.312481in}}%
\pgfpathlineto{\pgfqpoint{1.284736in}{1.306313in}}%
\pgfpathlineto{\pgfqpoint{1.298775in}{1.307121in}}%
\pgfpathlineto{\pgfqpoint{1.305818in}{1.299625in}}%
\pgfpathlineto{\pgfqpoint{1.303897in}{1.295339in}}%
\pgfpathlineto{\pgfqpoint{1.293840in}{1.286121in}}%
\pgfpathlineto{\pgfqpoint{1.291651in}{1.276850in}}%
\pgfpathlineto{\pgfqpoint{1.286167in}{1.271003in}}%
\pgfpathlineto{\pgfqpoint{1.280183in}{1.269745in}}%
\pgfpathlineto{\pgfqpoint{1.272888in}{1.260292in}}%
\pgfpathlineto{\pgfqpoint{1.268205in}{1.247503in}}%
\pgfpathlineto{\pgfqpoint{1.264762in}{1.243586in}}%
\pgfpathlineto{\pgfqpoint{1.223167in}{1.249876in}}%
\pgfpathclose%
\pgfusepath{fill}%
\end{pgfscope}%
\begin{pgfscope}%
\pgfpathrectangle{\pgfqpoint{0.100000in}{0.100000in}}{\pgfqpoint{3.608454in}{2.310000in}}%
\pgfusepath{clip}%
\pgfsetbuttcap%
\pgfsetmiterjoin%
\definecolor{currentfill}{rgb}{0.000000,0.423529,0.788235}%
\pgfsetfillcolor{currentfill}%
\pgfsetlinewidth{0.000000pt}%
\definecolor{currentstroke}{rgb}{0.000000,0.000000,0.000000}%
\pgfsetstrokecolor{currentstroke}%
\pgfsetstrokeopacity{0.000000}%
\pgfsetdash{}{0pt}%
\pgfpathmoveto{\pgfqpoint{2.669947in}{1.049451in}}%
\pgfpathlineto{\pgfqpoint{2.668245in}{1.039940in}}%
\pgfpathlineto{\pgfqpoint{2.670176in}{1.036633in}}%
\pgfpathlineto{\pgfqpoint{2.669558in}{1.020872in}}%
\pgfpathlineto{\pgfqpoint{2.673359in}{1.013797in}}%
\pgfpathlineto{\pgfqpoint{2.666467in}{1.008380in}}%
\pgfpathlineto{\pgfqpoint{2.659190in}{1.009555in}}%
\pgfpathlineto{\pgfqpoint{2.657488in}{0.993619in}}%
\pgfpathlineto{\pgfqpoint{2.622963in}{0.991238in}}%
\pgfpathlineto{\pgfqpoint{2.623076in}{0.990093in}}%
\pgfpathlineto{\pgfqpoint{2.595636in}{0.988190in}}%
\pgfpathlineto{\pgfqpoint{2.593933in}{1.008935in}}%
\pgfpathlineto{\pgfqpoint{2.599703in}{1.020694in}}%
\pgfpathlineto{\pgfqpoint{2.599104in}{1.027892in}}%
\pgfpathlineto{\pgfqpoint{2.607222in}{1.025513in}}%
\pgfpathlineto{\pgfqpoint{2.612717in}{1.030317in}}%
\pgfpathlineto{\pgfqpoint{2.611745in}{1.044790in}}%
\pgfpathlineto{\pgfqpoint{2.636003in}{1.046348in}}%
\pgfpathlineto{\pgfqpoint{2.634590in}{1.067850in}}%
\pgfpathlineto{\pgfqpoint{2.639911in}{1.067474in}}%
\pgfpathlineto{\pgfqpoint{2.648663in}{1.074175in}}%
\pgfpathlineto{\pgfqpoint{2.648588in}{1.077186in}}%
\pgfpathlineto{\pgfqpoint{2.662888in}{1.068219in}}%
\pgfpathlineto{\pgfqpoint{2.665061in}{1.060409in}}%
\pgfpathlineto{\pgfqpoint{2.668473in}{1.060141in}}%
\pgfpathlineto{\pgfqpoint{2.669947in}{1.049451in}}%
\pgfpathclose%
\pgfusepath{fill}%
\end{pgfscope}%
\begin{pgfscope}%
\pgfpathrectangle{\pgfqpoint{0.100000in}{0.100000in}}{\pgfqpoint{3.608454in}{2.310000in}}%
\pgfusepath{clip}%
\pgfsetbuttcap%
\pgfsetmiterjoin%
\definecolor{currentfill}{rgb}{0.000000,0.349020,0.825490}%
\pgfsetfillcolor{currentfill}%
\pgfsetlinewidth{0.000000pt}%
\definecolor{currentstroke}{rgb}{0.000000,0.000000,0.000000}%
\pgfsetstrokecolor{currentstroke}%
\pgfsetstrokeopacity{0.000000}%
\pgfsetdash{}{0pt}%
\pgfpathmoveto{\pgfqpoint{1.255196in}{1.457818in}}%
\pgfpathlineto{\pgfqpoint{1.267460in}{1.535366in}}%
\pgfpathlineto{\pgfqpoint{1.271578in}{1.562500in}}%
\pgfpathlineto{\pgfqpoint{1.302483in}{1.557690in}}%
\pgfpathlineto{\pgfqpoint{1.338390in}{1.552623in}}%
\pgfpathlineto{\pgfqpoint{1.401111in}{1.544070in}}%
\pgfpathlineto{\pgfqpoint{1.400660in}{1.537925in}}%
\pgfpathlineto{\pgfqpoint{1.411585in}{1.530462in}}%
\pgfpathlineto{\pgfqpoint{1.412111in}{1.525836in}}%
\pgfpathlineto{\pgfqpoint{1.405998in}{1.509570in}}%
\pgfpathlineto{\pgfqpoint{1.407647in}{1.498770in}}%
\pgfpathlineto{\pgfqpoint{1.407785in}{1.490488in}}%
\pgfpathlineto{\pgfqpoint{1.404073in}{1.463909in}}%
\pgfpathlineto{\pgfqpoint{1.403922in}{1.457766in}}%
\pgfpathlineto{\pgfqpoint{1.379506in}{1.460495in}}%
\pgfpathlineto{\pgfqpoint{1.380971in}{1.474066in}}%
\pgfpathlineto{\pgfqpoint{1.364204in}{1.476264in}}%
\pgfpathlineto{\pgfqpoint{1.362220in}{1.462510in}}%
\pgfpathlineto{\pgfqpoint{1.355468in}{1.463655in}}%
\pgfpathlineto{\pgfqpoint{1.354499in}{1.456659in}}%
\pgfpathlineto{\pgfqpoint{1.324000in}{1.460674in}}%
\pgfpathlineto{\pgfqpoint{1.322533in}{1.450461in}}%
\pgfpathlineto{\pgfqpoint{1.285446in}{1.455840in}}%
\pgfpathlineto{\pgfqpoint{1.284941in}{1.452441in}}%
\pgfpathlineto{\pgfqpoint{1.255196in}{1.457818in}}%
\pgfpathclose%
\pgfusepath{fill}%
\end{pgfscope}%
\begin{pgfscope}%
\pgfpathrectangle{\pgfqpoint{0.100000in}{0.100000in}}{\pgfqpoint{3.608454in}{2.310000in}}%
\pgfusepath{clip}%
\pgfsetbuttcap%
\pgfsetmiterjoin%
\definecolor{currentfill}{rgb}{0.000000,0.717647,0.641176}%
\pgfsetfillcolor{currentfill}%
\pgfsetlinewidth{0.000000pt}%
\definecolor{currentstroke}{rgb}{0.000000,0.000000,0.000000}%
\pgfsetstrokecolor{currentstroke}%
\pgfsetstrokeopacity{0.000000}%
\pgfsetdash{}{0pt}%
\pgfpathmoveto{\pgfqpoint{1.709808in}{0.564368in}}%
\pgfpathlineto{\pgfqpoint{1.748728in}{0.562219in}}%
\pgfpathlineto{\pgfqpoint{1.744553in}{0.490983in}}%
\pgfpathlineto{\pgfqpoint{1.737580in}{0.491300in}}%
\pgfpathlineto{\pgfqpoint{1.732274in}{0.498247in}}%
\pgfpathlineto{\pgfqpoint{1.729908in}{0.510537in}}%
\pgfpathlineto{\pgfqpoint{1.725203in}{0.520544in}}%
\pgfpathlineto{\pgfqpoint{1.726134in}{0.523344in}}%
\pgfpathlineto{\pgfqpoint{1.719531in}{0.529734in}}%
\pgfpathlineto{\pgfqpoint{1.716987in}{0.543109in}}%
\pgfpathlineto{\pgfqpoint{1.710227in}{0.553019in}}%
\pgfpathlineto{\pgfqpoint{1.709808in}{0.564368in}}%
\pgfpathclose%
\pgfusepath{fill}%
\end{pgfscope}%
\begin{pgfscope}%
\pgfpathrectangle{\pgfqpoint{0.100000in}{0.100000in}}{\pgfqpoint{3.608454in}{2.310000in}}%
\pgfusepath{clip}%
\pgfsetbuttcap%
\pgfsetmiterjoin%
\definecolor{currentfill}{rgb}{0.000000,0.631373,0.684314}%
\pgfsetfillcolor{currentfill}%
\pgfsetlinewidth{0.000000pt}%
\definecolor{currentstroke}{rgb}{0.000000,0.000000,0.000000}%
\pgfsetstrokecolor{currentstroke}%
\pgfsetstrokeopacity{0.000000}%
\pgfsetdash{}{0pt}%
\pgfpathmoveto{\pgfqpoint{2.111050in}{0.772661in}}%
\pgfpathlineto{\pgfqpoint{2.108517in}{0.767596in}}%
\pgfpathlineto{\pgfqpoint{2.111897in}{0.751748in}}%
\pgfpathlineto{\pgfqpoint{2.115910in}{0.747279in}}%
\pgfpathlineto{\pgfqpoint{2.106573in}{0.739232in}}%
\pgfpathlineto{\pgfqpoint{2.099586in}{0.743769in}}%
\pgfpathlineto{\pgfqpoint{2.096621in}{0.750671in}}%
\pgfpathlineto{\pgfqpoint{2.088232in}{0.752553in}}%
\pgfpathlineto{\pgfqpoint{2.062519in}{0.748614in}}%
\pgfpathlineto{\pgfqpoint{2.056507in}{0.745639in}}%
\pgfpathlineto{\pgfqpoint{2.058523in}{0.753655in}}%
\pgfpathlineto{\pgfqpoint{2.048151in}{0.760702in}}%
\pgfpathlineto{\pgfqpoint{2.047446in}{0.765794in}}%
\pgfpathlineto{\pgfqpoint{2.042310in}{0.767446in}}%
\pgfpathlineto{\pgfqpoint{2.034850in}{0.788420in}}%
\pgfpathlineto{\pgfqpoint{2.023652in}{0.805506in}}%
\pgfpathlineto{\pgfqpoint{2.013777in}{0.811827in}}%
\pgfpathlineto{\pgfqpoint{2.034479in}{0.813847in}}%
\pgfpathlineto{\pgfqpoint{2.034871in}{0.852102in}}%
\pgfpathlineto{\pgfqpoint{2.044379in}{0.851895in}}%
\pgfpathlineto{\pgfqpoint{2.050385in}{0.846426in}}%
\pgfpathlineto{\pgfqpoint{2.053670in}{0.846942in}}%
\pgfpathlineto{\pgfqpoint{2.064298in}{0.842428in}}%
\pgfpathlineto{\pgfqpoint{2.062420in}{0.861532in}}%
\pgfpathlineto{\pgfqpoint{2.086189in}{0.861574in}}%
\pgfpathlineto{\pgfqpoint{2.095246in}{0.861494in}}%
\pgfpathlineto{\pgfqpoint{2.096641in}{0.856769in}}%
\pgfpathlineto{\pgfqpoint{2.096510in}{0.830344in}}%
\pgfpathlineto{\pgfqpoint{2.107625in}{0.827720in}}%
\pgfpathlineto{\pgfqpoint{2.107805in}{0.772710in}}%
\pgfpathlineto{\pgfqpoint{2.111050in}{0.772661in}}%
\pgfpathclose%
\pgfusepath{fill}%
\end{pgfscope}%
\begin{pgfscope}%
\pgfpathrectangle{\pgfqpoint{0.100000in}{0.100000in}}{\pgfqpoint{3.608454in}{2.310000in}}%
\pgfusepath{clip}%
\pgfsetbuttcap%
\pgfsetmiterjoin%
\definecolor{currentfill}{rgb}{0.000000,0.600000,0.700000}%
\pgfsetfillcolor{currentfill}%
\pgfsetlinewidth{0.000000pt}%
\definecolor{currentstroke}{rgb}{0.000000,0.000000,0.000000}%
\pgfsetstrokecolor{currentstroke}%
\pgfsetstrokeopacity{0.000000}%
\pgfsetdash{}{0pt}%
\pgfpathmoveto{\pgfqpoint{0.469229in}{1.459834in}}%
\pgfpathlineto{\pgfqpoint{0.465525in}{1.465978in}}%
\pgfpathlineto{\pgfqpoint{0.471844in}{1.486991in}}%
\pgfpathlineto{\pgfqpoint{0.471692in}{1.491720in}}%
\pgfpathlineto{\pgfqpoint{0.476612in}{1.508153in}}%
\pgfpathlineto{\pgfqpoint{0.471851in}{1.509726in}}%
\pgfpathlineto{\pgfqpoint{0.471900in}{1.514999in}}%
\pgfpathlineto{\pgfqpoint{0.477138in}{1.516663in}}%
\pgfpathlineto{\pgfqpoint{0.480918in}{1.524945in}}%
\pgfpathlineto{\pgfqpoint{0.474910in}{1.526813in}}%
\pgfpathlineto{\pgfqpoint{0.479765in}{1.542960in}}%
\pgfpathlineto{\pgfqpoint{0.465266in}{1.547913in}}%
\pgfpathlineto{\pgfqpoint{0.460110in}{1.545792in}}%
\pgfpathlineto{\pgfqpoint{0.455284in}{1.549289in}}%
\pgfpathlineto{\pgfqpoint{0.454748in}{1.561199in}}%
\pgfpathlineto{\pgfqpoint{0.452551in}{1.565246in}}%
\pgfpathlineto{\pgfqpoint{0.452005in}{1.577321in}}%
\pgfpathlineto{\pgfqpoint{0.446281in}{1.581110in}}%
\pgfpathlineto{\pgfqpoint{0.450911in}{1.584637in}}%
\pgfpathlineto{\pgfqpoint{0.480562in}{1.575547in}}%
\pgfpathlineto{\pgfqpoint{0.485256in}{1.568714in}}%
\pgfpathlineto{\pgfqpoint{0.484414in}{1.564406in}}%
\pgfpathlineto{\pgfqpoint{0.489967in}{1.557121in}}%
\pgfpathlineto{\pgfqpoint{0.498193in}{1.555922in}}%
\pgfpathlineto{\pgfqpoint{0.502169in}{1.569148in}}%
\pgfpathlineto{\pgfqpoint{0.507112in}{1.573994in}}%
\pgfpathlineto{\pgfqpoint{0.514461in}{1.575916in}}%
\pgfpathlineto{\pgfqpoint{0.520214in}{1.588877in}}%
\pgfpathlineto{\pgfqpoint{0.532464in}{1.597564in}}%
\pgfpathlineto{\pgfqpoint{0.538914in}{1.596597in}}%
\pgfpathlineto{\pgfqpoint{0.556871in}{1.597127in}}%
\pgfpathlineto{\pgfqpoint{0.563417in}{1.600581in}}%
\pgfpathlineto{\pgfqpoint{0.567874in}{1.599010in}}%
\pgfpathlineto{\pgfqpoint{0.570419in}{1.591928in}}%
\pgfpathlineto{\pgfqpoint{0.599878in}{1.583444in}}%
\pgfpathlineto{\pgfqpoint{0.592721in}{1.558163in}}%
\pgfpathlineto{\pgfqpoint{0.607050in}{1.554247in}}%
\pgfpathlineto{\pgfqpoint{0.606969in}{1.551881in}}%
\pgfpathlineto{\pgfqpoint{0.632981in}{1.544578in}}%
\pgfpathlineto{\pgfqpoint{0.631100in}{1.537480in}}%
\pgfpathlineto{\pgfqpoint{0.623821in}{1.536529in}}%
\pgfpathlineto{\pgfqpoint{0.620041in}{1.519423in}}%
\pgfpathlineto{\pgfqpoint{0.623308in}{1.518183in}}%
\pgfpathlineto{\pgfqpoint{0.620520in}{1.503015in}}%
\pgfpathlineto{\pgfqpoint{0.609163in}{1.520157in}}%
\pgfpathlineto{\pgfqpoint{0.604651in}{1.513031in}}%
\pgfpathlineto{\pgfqpoint{0.606996in}{1.503863in}}%
\pgfpathlineto{\pgfqpoint{0.603382in}{1.497016in}}%
\pgfpathlineto{\pgfqpoint{0.597596in}{1.492295in}}%
\pgfpathlineto{\pgfqpoint{0.592836in}{1.501027in}}%
\pgfpathlineto{\pgfqpoint{0.583557in}{1.498623in}}%
\pgfpathlineto{\pgfqpoint{0.577147in}{1.506768in}}%
\pgfpathlineto{\pgfqpoint{0.566537in}{1.504878in}}%
\pgfpathlineto{\pgfqpoint{0.550096in}{1.491227in}}%
\pgfpathlineto{\pgfqpoint{0.537164in}{1.478720in}}%
\pgfpathlineto{\pgfqpoint{0.526181in}{1.471930in}}%
\pgfpathlineto{\pgfqpoint{0.515295in}{1.495292in}}%
\pgfpathlineto{\pgfqpoint{0.508101in}{1.469498in}}%
\pgfpathlineto{\pgfqpoint{0.504262in}{1.472446in}}%
\pgfpathlineto{\pgfqpoint{0.489989in}{1.471204in}}%
\pgfpathlineto{\pgfqpoint{0.469229in}{1.459834in}}%
\pgfpathclose%
\pgfusepath{fill}%
\end{pgfscope}%
\begin{pgfscope}%
\pgfpathrectangle{\pgfqpoint{0.100000in}{0.100000in}}{\pgfqpoint{3.608454in}{2.310000in}}%
\pgfusepath{clip}%
\pgfsetbuttcap%
\pgfsetmiterjoin%
\definecolor{currentfill}{rgb}{0.000000,0.478431,0.760784}%
\pgfsetfillcolor{currentfill}%
\pgfsetlinewidth{0.000000pt}%
\definecolor{currentstroke}{rgb}{0.000000,0.000000,0.000000}%
\pgfsetstrokecolor{currentstroke}%
\pgfsetstrokeopacity{0.000000}%
\pgfsetdash{}{0pt}%
\pgfpathmoveto{\pgfqpoint{1.886398in}{1.921692in}}%
\pgfpathlineto{\pgfqpoint{1.884519in}{1.921788in}}%
\pgfpathlineto{\pgfqpoint{1.885796in}{1.949443in}}%
\pgfpathlineto{\pgfqpoint{1.863197in}{1.950599in}}%
\pgfpathlineto{\pgfqpoint{1.864737in}{1.978421in}}%
\pgfpathlineto{\pgfqpoint{1.862851in}{1.978498in}}%
\pgfpathlineto{\pgfqpoint{1.864232in}{2.006135in}}%
\pgfpathlineto{\pgfqpoint{1.875835in}{2.005576in}}%
\pgfpathlineto{\pgfqpoint{1.917334in}{2.003750in}}%
\pgfpathlineto{\pgfqpoint{1.918791in}{1.996815in}}%
\pgfpathlineto{\pgfqpoint{1.946235in}{1.995849in}}%
\pgfpathlineto{\pgfqpoint{1.960057in}{1.995378in}}%
\pgfpathlineto{\pgfqpoint{1.960762in}{1.974477in}}%
\pgfpathlineto{\pgfqpoint{1.959926in}{1.946668in}}%
\pgfpathlineto{\pgfqpoint{1.940671in}{1.947331in}}%
\pgfpathlineto{\pgfqpoint{1.939707in}{1.919594in}}%
\pgfpathlineto{\pgfqpoint{1.886398in}{1.921692in}}%
\pgfpathclose%
\pgfusepath{fill}%
\end{pgfscope}%
\begin{pgfscope}%
\pgfpathrectangle{\pgfqpoint{0.100000in}{0.100000in}}{\pgfqpoint{3.608454in}{2.310000in}}%
\pgfusepath{clip}%
\pgfsetbuttcap%
\pgfsetmiterjoin%
\definecolor{currentfill}{rgb}{0.000000,0.529412,0.735294}%
\pgfsetfillcolor{currentfill}%
\pgfsetlinewidth{0.000000pt}%
\definecolor{currentstroke}{rgb}{0.000000,0.000000,0.000000}%
\pgfsetstrokecolor{currentstroke}%
\pgfsetstrokeopacity{0.000000}%
\pgfsetdash{}{0pt}%
\pgfpathmoveto{\pgfqpoint{3.420582in}{1.833115in}}%
\pgfpathlineto{\pgfqpoint{3.412904in}{1.835488in}}%
\pgfpathlineto{\pgfqpoint{3.401069in}{1.830482in}}%
\pgfpathlineto{\pgfqpoint{3.396386in}{1.835470in}}%
\pgfpathlineto{\pgfqpoint{3.382485in}{1.831082in}}%
\pgfpathlineto{\pgfqpoint{3.374364in}{1.831752in}}%
\pgfpathlineto{\pgfqpoint{3.370781in}{1.834556in}}%
\pgfpathlineto{\pgfqpoint{3.375120in}{1.840353in}}%
\pgfpathlineto{\pgfqpoint{3.375426in}{1.849437in}}%
\pgfpathlineto{\pgfqpoint{3.371356in}{1.850568in}}%
\pgfpathlineto{\pgfqpoint{3.372767in}{1.866271in}}%
\pgfpathlineto{\pgfqpoint{3.365489in}{1.867233in}}%
\pgfpathlineto{\pgfqpoint{3.365877in}{1.874467in}}%
\pgfpathlineto{\pgfqpoint{3.354162in}{1.876806in}}%
\pgfpathlineto{\pgfqpoint{3.358819in}{1.880708in}}%
\pgfpathlineto{\pgfqpoint{3.363839in}{1.892385in}}%
\pgfpathlineto{\pgfqpoint{3.353338in}{1.893246in}}%
\pgfpathlineto{\pgfqpoint{3.350605in}{1.900061in}}%
\pgfpathlineto{\pgfqpoint{3.349620in}{1.909475in}}%
\pgfpathlineto{\pgfqpoint{3.352226in}{1.924022in}}%
\pgfpathlineto{\pgfqpoint{3.361540in}{1.921976in}}%
\pgfpathlineto{\pgfqpoint{3.363137in}{1.929015in}}%
\pgfpathlineto{\pgfqpoint{3.370017in}{1.926361in}}%
\pgfpathlineto{\pgfqpoint{3.371539in}{1.933332in}}%
\pgfpathlineto{\pgfqpoint{3.384688in}{1.930010in}}%
\pgfpathlineto{\pgfqpoint{3.381069in}{1.911013in}}%
\pgfpathlineto{\pgfqpoint{3.384973in}{1.910295in}}%
\pgfpathlineto{\pgfqpoint{3.398006in}{1.914180in}}%
\pgfpathlineto{\pgfqpoint{3.396256in}{1.920610in}}%
\pgfpathlineto{\pgfqpoint{3.399587in}{1.926596in}}%
\pgfpathlineto{\pgfqpoint{3.405777in}{1.928235in}}%
\pgfpathlineto{\pgfqpoint{3.410255in}{1.935145in}}%
\pgfpathlineto{\pgfqpoint{3.415486in}{1.930416in}}%
\pgfpathlineto{\pgfqpoint{3.422059in}{1.929884in}}%
\pgfpathlineto{\pgfqpoint{3.423697in}{1.926204in}}%
\pgfpathlineto{\pgfqpoint{3.431085in}{1.927713in}}%
\pgfpathlineto{\pgfqpoint{3.433355in}{1.925260in}}%
\pgfpathlineto{\pgfqpoint{3.438972in}{1.915003in}}%
\pgfpathlineto{\pgfqpoint{3.442370in}{1.903310in}}%
\pgfpathlineto{\pgfqpoint{3.431626in}{1.898622in}}%
\pgfpathlineto{\pgfqpoint{3.437941in}{1.884217in}}%
\pgfpathlineto{\pgfqpoint{3.431173in}{1.881530in}}%
\pgfpathlineto{\pgfqpoint{3.432335in}{1.876511in}}%
\pgfpathlineto{\pgfqpoint{3.424616in}{1.868740in}}%
\pgfpathlineto{\pgfqpoint{3.421501in}{1.872886in}}%
\pgfpathlineto{\pgfqpoint{3.415970in}{1.857152in}}%
\pgfpathlineto{\pgfqpoint{3.413847in}{1.847236in}}%
\pgfpathlineto{\pgfqpoint{3.421042in}{1.838069in}}%
\pgfpathlineto{\pgfqpoint{3.420582in}{1.833115in}}%
\pgfpathclose%
\pgfusepath{fill}%
\end{pgfscope}%
\begin{pgfscope}%
\pgfpathrectangle{\pgfqpoint{0.100000in}{0.100000in}}{\pgfqpoint{3.608454in}{2.310000in}}%
\pgfusepath{clip}%
\pgfsetbuttcap%
\pgfsetmiterjoin%
\definecolor{currentfill}{rgb}{0.000000,0.666667,0.666667}%
\pgfsetfillcolor{currentfill}%
\pgfsetlinewidth{0.000000pt}%
\definecolor{currentstroke}{rgb}{0.000000,0.000000,0.000000}%
\pgfsetstrokecolor{currentstroke}%
\pgfsetstrokeopacity{0.000000}%
\pgfsetdash{}{0pt}%
\pgfpathmoveto{\pgfqpoint{3.043474in}{1.583306in}}%
\pgfpathlineto{\pgfqpoint{3.040035in}{1.579416in}}%
\pgfpathlineto{\pgfqpoint{3.028102in}{1.574866in}}%
\pgfpathlineto{\pgfqpoint{3.020349in}{1.574020in}}%
\pgfpathlineto{\pgfqpoint{3.013646in}{1.587974in}}%
\pgfpathlineto{\pgfqpoint{2.995768in}{1.584878in}}%
\pgfpathlineto{\pgfqpoint{2.991632in}{1.609704in}}%
\pgfpathlineto{\pgfqpoint{2.984410in}{1.607734in}}%
\pgfpathlineto{\pgfqpoint{2.961131in}{1.605506in}}%
\pgfpathlineto{\pgfqpoint{2.955120in}{1.642813in}}%
\pgfpathlineto{\pgfqpoint{2.973125in}{1.655287in}}%
\pgfpathlineto{\pgfqpoint{2.989337in}{1.668857in}}%
\pgfpathlineto{\pgfqpoint{2.999470in}{1.675879in}}%
\pgfpathlineto{\pgfqpoint{3.011435in}{1.687328in}}%
\pgfpathlineto{\pgfqpoint{3.021181in}{1.698369in}}%
\pgfpathlineto{\pgfqpoint{3.027192in}{1.702502in}}%
\pgfpathlineto{\pgfqpoint{3.032000in}{1.700788in}}%
\pgfpathlineto{\pgfqpoint{3.039413in}{1.658799in}}%
\pgfpathlineto{\pgfqpoint{3.047674in}{1.660178in}}%
\pgfpathlineto{\pgfqpoint{3.050043in}{1.647962in}}%
\pgfpathlineto{\pgfqpoint{3.048278in}{1.646809in}}%
\pgfpathlineto{\pgfqpoint{3.050721in}{1.630639in}}%
\pgfpathlineto{\pgfqpoint{3.051910in}{1.617021in}}%
\pgfpathlineto{\pgfqpoint{3.045476in}{1.613591in}}%
\pgfpathlineto{\pgfqpoint{3.046524in}{1.607153in}}%
\pgfpathlineto{\pgfqpoint{3.039983in}{1.605278in}}%
\pgfpathlineto{\pgfqpoint{3.043474in}{1.583306in}}%
\pgfpathclose%
\pgfusepath{fill}%
\end{pgfscope}%
\begin{pgfscope}%
\pgfpathrectangle{\pgfqpoint{0.100000in}{0.100000in}}{\pgfqpoint{3.608454in}{2.310000in}}%
\pgfusepath{clip}%
\pgfsetbuttcap%
\pgfsetmiterjoin%
\definecolor{currentfill}{rgb}{0.000000,0.333333,0.833333}%
\pgfsetfillcolor{currentfill}%
\pgfsetlinewidth{0.000000pt}%
\definecolor{currentstroke}{rgb}{0.000000,0.000000,0.000000}%
\pgfsetstrokecolor{currentstroke}%
\pgfsetstrokeopacity{0.000000}%
\pgfsetdash{}{0pt}%
\pgfpathmoveto{\pgfqpoint{1.992970in}{1.502588in}}%
\pgfpathlineto{\pgfqpoint{1.965564in}{1.503301in}}%
\pgfpathlineto{\pgfqpoint{1.966281in}{1.530825in}}%
\pgfpathlineto{\pgfqpoint{1.952480in}{1.526297in}}%
\pgfpathlineto{\pgfqpoint{1.952621in}{1.531206in}}%
\pgfpathlineto{\pgfqpoint{1.946376in}{1.531379in}}%
\pgfpathlineto{\pgfqpoint{1.946070in}{1.524600in}}%
\pgfpathlineto{\pgfqpoint{1.925386in}{1.523978in}}%
\pgfpathlineto{\pgfqpoint{1.922421in}{1.521780in}}%
\pgfpathlineto{\pgfqpoint{1.912100in}{1.522170in}}%
\pgfpathlineto{\pgfqpoint{1.911513in}{1.532437in}}%
\pgfpathlineto{\pgfqpoint{1.912258in}{1.559919in}}%
\pgfpathlineto{\pgfqpoint{1.912766in}{1.573654in}}%
\pgfpathlineto{\pgfqpoint{1.939982in}{1.572847in}}%
\pgfpathlineto{\pgfqpoint{1.939698in}{1.559064in}}%
\pgfpathlineto{\pgfqpoint{1.987553in}{1.557852in}}%
\pgfpathlineto{\pgfqpoint{1.994270in}{1.557682in}}%
\pgfpathlineto{\pgfqpoint{1.992970in}{1.502588in}}%
\pgfpathclose%
\pgfusepath{fill}%
\end{pgfscope}%
\begin{pgfscope}%
\pgfpathrectangle{\pgfqpoint{0.100000in}{0.100000in}}{\pgfqpoint{3.608454in}{2.310000in}}%
\pgfusepath{clip}%
\pgfsetbuttcap%
\pgfsetmiterjoin%
\definecolor{currentfill}{rgb}{0.000000,0.462745,0.768627}%
\pgfsetfillcolor{currentfill}%
\pgfsetlinewidth{0.000000pt}%
\definecolor{currentstroke}{rgb}{0.000000,0.000000,0.000000}%
\pgfsetstrokecolor{currentstroke}%
\pgfsetstrokeopacity{0.000000}%
\pgfsetdash{}{0pt}%
\pgfpathmoveto{\pgfqpoint{2.115120in}{1.723283in}}%
\pgfpathlineto{\pgfqpoint{2.114776in}{1.743905in}}%
\pgfpathlineto{\pgfqpoint{2.101083in}{1.743915in}}%
\pgfpathlineto{\pgfqpoint{2.100603in}{1.750815in}}%
\pgfpathlineto{\pgfqpoint{2.100687in}{1.757698in}}%
\pgfpathlineto{\pgfqpoint{2.114393in}{1.757738in}}%
\pgfpathlineto{\pgfqpoint{2.114392in}{1.774876in}}%
\pgfpathlineto{\pgfqpoint{2.118639in}{1.771585in}}%
\pgfpathlineto{\pgfqpoint{2.128088in}{1.771585in}}%
\pgfpathlineto{\pgfqpoint{2.167464in}{1.771934in}}%
\pgfpathlineto{\pgfqpoint{2.168433in}{1.778817in}}%
\pgfpathlineto{\pgfqpoint{2.190305in}{1.779149in}}%
\pgfpathlineto{\pgfqpoint{2.190701in}{1.751566in}}%
\pgfpathlineto{\pgfqpoint{2.176891in}{1.751349in}}%
\pgfpathlineto{\pgfqpoint{2.177278in}{1.723751in}}%
\pgfpathlineto{\pgfqpoint{2.115120in}{1.723283in}}%
\pgfpathclose%
\pgfusepath{fill}%
\end{pgfscope}%
\begin{pgfscope}%
\pgfpathrectangle{\pgfqpoint{0.100000in}{0.100000in}}{\pgfqpoint{3.608454in}{2.310000in}}%
\pgfusepath{clip}%
\pgfsetbuttcap%
\pgfsetmiterjoin%
\definecolor{currentfill}{rgb}{0.000000,0.529412,0.735294}%
\pgfsetfillcolor{currentfill}%
\pgfsetlinewidth{0.000000pt}%
\definecolor{currentstroke}{rgb}{0.000000,0.000000,0.000000}%
\pgfsetstrokecolor{currentstroke}%
\pgfsetstrokeopacity{0.000000}%
\pgfsetdash{}{0pt}%
\pgfpathmoveto{\pgfqpoint{2.248598in}{1.467244in}}%
\pgfpathlineto{\pgfqpoint{2.247888in}{1.491592in}}%
\pgfpathlineto{\pgfqpoint{2.220496in}{1.490915in}}%
\pgfpathlineto{\pgfqpoint{2.219965in}{1.511654in}}%
\pgfpathlineto{\pgfqpoint{2.206358in}{1.511376in}}%
\pgfpathlineto{\pgfqpoint{2.205803in}{1.538835in}}%
\pgfpathlineto{\pgfqpoint{2.266869in}{1.540448in}}%
\pgfpathlineto{\pgfqpoint{2.287673in}{1.541193in}}%
\pgfpathlineto{\pgfqpoint{2.288611in}{1.513740in}}%
\pgfpathlineto{\pgfqpoint{2.302293in}{1.514165in}}%
\pgfpathlineto{\pgfqpoint{2.303069in}{1.470734in}}%
\pgfpathlineto{\pgfqpoint{2.290253in}{1.469670in}}%
\pgfpathlineto{\pgfqpoint{2.248598in}{1.467244in}}%
\pgfpathclose%
\pgfusepath{fill}%
\end{pgfscope}%
\begin{pgfscope}%
\pgfpathrectangle{\pgfqpoint{0.100000in}{0.100000in}}{\pgfqpoint{3.608454in}{2.310000in}}%
\pgfusepath{clip}%
\pgfsetbuttcap%
\pgfsetmiterjoin%
\definecolor{currentfill}{rgb}{0.000000,0.627451,0.686275}%
\pgfsetfillcolor{currentfill}%
\pgfsetlinewidth{0.000000pt}%
\definecolor{currentstroke}{rgb}{0.000000,0.000000,0.000000}%
\pgfsetstrokecolor{currentstroke}%
\pgfsetstrokeopacity{0.000000}%
\pgfsetdash{}{0pt}%
\pgfpathmoveto{\pgfqpoint{1.141273in}{1.664762in}}%
\pgfpathlineto{\pgfqpoint{1.142493in}{1.660281in}}%
\pgfpathlineto{\pgfqpoint{1.139137in}{1.653332in}}%
\pgfpathlineto{\pgfqpoint{1.141060in}{1.650733in}}%
\pgfpathlineto{\pgfqpoint{1.142480in}{1.636751in}}%
\pgfpathlineto{\pgfqpoint{1.138324in}{1.627635in}}%
\pgfpathlineto{\pgfqpoint{1.132825in}{1.619941in}}%
\pgfpathlineto{\pgfqpoint{1.121923in}{1.621467in}}%
\pgfpathlineto{\pgfqpoint{1.117387in}{1.618454in}}%
\pgfpathlineto{\pgfqpoint{1.110470in}{1.623700in}}%
\pgfpathlineto{\pgfqpoint{1.101371in}{1.618838in}}%
\pgfpathlineto{\pgfqpoint{1.089153in}{1.621193in}}%
\pgfpathlineto{\pgfqpoint{1.070321in}{1.603891in}}%
\pgfpathlineto{\pgfqpoint{1.051249in}{1.601475in}}%
\pgfpathlineto{\pgfqpoint{0.978586in}{1.616384in}}%
\pgfpathlineto{\pgfqpoint{0.994964in}{1.693281in}}%
\pgfpathlineto{\pgfqpoint{1.004416in}{1.690817in}}%
\pgfpathlineto{\pgfqpoint{1.054992in}{1.681239in}}%
\pgfpathlineto{\pgfqpoint{1.060120in}{1.706792in}}%
\pgfpathlineto{\pgfqpoint{1.080051in}{1.702943in}}%
\pgfpathlineto{\pgfqpoint{1.081382in}{1.709710in}}%
\pgfpathlineto{\pgfqpoint{1.090342in}{1.707956in}}%
\pgfpathlineto{\pgfqpoint{1.091640in}{1.714726in}}%
\pgfpathlineto{\pgfqpoint{1.098147in}{1.713450in}}%
\pgfpathlineto{\pgfqpoint{1.101946in}{1.704218in}}%
\pgfpathlineto{\pgfqpoint{1.109727in}{1.693683in}}%
\pgfpathlineto{\pgfqpoint{1.117899in}{1.692151in}}%
\pgfpathlineto{\pgfqpoint{1.122487in}{1.688955in}}%
\pgfpathlineto{\pgfqpoint{1.126068in}{1.701146in}}%
\pgfpathlineto{\pgfqpoint{1.141976in}{1.698141in}}%
\pgfpathlineto{\pgfqpoint{1.143732in}{1.692681in}}%
\pgfpathlineto{\pgfqpoint{1.141935in}{1.686685in}}%
\pgfpathlineto{\pgfqpoint{1.137407in}{1.682303in}}%
\pgfpathlineto{\pgfqpoint{1.137582in}{1.673009in}}%
\pgfpathlineto{\pgfqpoint{1.141528in}{1.670743in}}%
\pgfpathlineto{\pgfqpoint{1.141273in}{1.664762in}}%
\pgfpathclose%
\pgfusepath{fill}%
\end{pgfscope}%
\begin{pgfscope}%
\pgfpathrectangle{\pgfqpoint{0.100000in}{0.100000in}}{\pgfqpoint{3.608454in}{2.310000in}}%
\pgfusepath{clip}%
\pgfsetbuttcap%
\pgfsetmiterjoin%
\definecolor{currentfill}{rgb}{0.000000,0.713725,0.643137}%
\pgfsetfillcolor{currentfill}%
\pgfsetlinewidth{0.000000pt}%
\definecolor{currentstroke}{rgb}{0.000000,0.000000,0.000000}%
\pgfsetstrokecolor{currentstroke}%
\pgfsetstrokeopacity{0.000000}%
\pgfsetdash{}{0pt}%
\pgfpathmoveto{\pgfqpoint{2.811045in}{1.319515in}}%
\pgfpathlineto{\pgfqpoint{2.804061in}{1.326162in}}%
\pgfpathlineto{\pgfqpoint{2.796848in}{1.326453in}}%
\pgfpathlineto{\pgfqpoint{2.787097in}{1.337635in}}%
\pgfpathlineto{\pgfqpoint{2.789928in}{1.341177in}}%
\pgfpathlineto{\pgfqpoint{2.786720in}{1.347970in}}%
\pgfpathlineto{\pgfqpoint{2.779558in}{1.348364in}}%
\pgfpathlineto{\pgfqpoint{2.772188in}{1.346181in}}%
\pgfpathlineto{\pgfqpoint{2.768060in}{1.364802in}}%
\pgfpathlineto{\pgfqpoint{2.777359in}{1.362093in}}%
\pgfpathlineto{\pgfqpoint{2.784131in}{1.363939in}}%
\pgfpathlineto{\pgfqpoint{2.790830in}{1.362794in}}%
\pgfpathlineto{\pgfqpoint{2.798224in}{1.355170in}}%
\pgfpathlineto{\pgfqpoint{2.804813in}{1.353816in}}%
\pgfpathlineto{\pgfqpoint{2.807405in}{1.358642in}}%
\pgfpathlineto{\pgfqpoint{2.812960in}{1.360831in}}%
\pgfpathlineto{\pgfqpoint{2.829317in}{1.356171in}}%
\pgfpathlineto{\pgfqpoint{2.836077in}{1.357238in}}%
\pgfpathlineto{\pgfqpoint{2.842704in}{1.366486in}}%
\pgfpathlineto{\pgfqpoint{2.841606in}{1.355934in}}%
\pgfpathlineto{\pgfqpoint{2.835681in}{1.347268in}}%
\pgfpathlineto{\pgfqpoint{2.833231in}{1.340587in}}%
\pgfpathlineto{\pgfqpoint{2.834430in}{1.334621in}}%
\pgfpathlineto{\pgfqpoint{2.826872in}{1.332234in}}%
\pgfpathlineto{\pgfqpoint{2.822416in}{1.337584in}}%
\pgfpathlineto{\pgfqpoint{2.817804in}{1.333230in}}%
\pgfpathlineto{\pgfqpoint{2.818045in}{1.327218in}}%
\pgfpathlineto{\pgfqpoint{2.811045in}{1.319515in}}%
\pgfpathclose%
\pgfusepath{fill}%
\end{pgfscope}%
\begin{pgfscope}%
\pgfpathrectangle{\pgfqpoint{0.100000in}{0.100000in}}{\pgfqpoint{3.608454in}{2.310000in}}%
\pgfusepath{clip}%
\pgfsetbuttcap%
\pgfsetmiterjoin%
\definecolor{currentfill}{rgb}{0.000000,0.560784,0.719608}%
\pgfsetfillcolor{currentfill}%
\pgfsetlinewidth{0.000000pt}%
\definecolor{currentstroke}{rgb}{0.000000,0.000000,0.000000}%
\pgfsetstrokecolor{currentstroke}%
\pgfsetstrokeopacity{0.000000}%
\pgfsetdash{}{0pt}%
\pgfpathmoveto{\pgfqpoint{2.839212in}{0.870674in}}%
\pgfpathlineto{\pgfqpoint{2.834772in}{0.868856in}}%
\pgfpathlineto{\pgfqpoint{2.830726in}{0.859944in}}%
\pgfpathlineto{\pgfqpoint{2.822592in}{0.856703in}}%
\pgfpathlineto{\pgfqpoint{2.816271in}{0.861716in}}%
\pgfpathlineto{\pgfqpoint{2.812799in}{0.870744in}}%
\pgfpathlineto{\pgfqpoint{2.821573in}{0.886623in}}%
\pgfpathlineto{\pgfqpoint{2.827685in}{0.882746in}}%
\pgfpathlineto{\pgfqpoint{2.830174in}{0.891948in}}%
\pgfpathlineto{\pgfqpoint{2.830655in}{0.902493in}}%
\pgfpathlineto{\pgfqpoint{2.835752in}{0.904367in}}%
\pgfpathlineto{\pgfqpoint{2.833490in}{0.924483in}}%
\pgfpathlineto{\pgfqpoint{2.844937in}{0.925826in}}%
\pgfpathlineto{\pgfqpoint{2.848183in}{0.924494in}}%
\pgfpathlineto{\pgfqpoint{2.849065in}{0.920723in}}%
\pgfpathlineto{\pgfqpoint{2.866524in}{0.926111in}}%
\pgfpathlineto{\pgfqpoint{2.874063in}{0.928161in}}%
\pgfpathlineto{\pgfqpoint{2.881402in}{0.908328in}}%
\pgfpathlineto{\pgfqpoint{2.878409in}{0.905708in}}%
\pgfpathlineto{\pgfqpoint{2.893588in}{0.882495in}}%
\pgfpathlineto{\pgfqpoint{2.901077in}{0.870542in}}%
\pgfpathlineto{\pgfqpoint{2.898491in}{0.872568in}}%
\pgfpathlineto{\pgfqpoint{2.890867in}{0.862703in}}%
\pgfpathlineto{\pgfqpoint{2.887187in}{0.854824in}}%
\pgfpathlineto{\pgfqpoint{2.893436in}{0.849951in}}%
\pgfpathlineto{\pgfqpoint{2.891100in}{0.843787in}}%
\pgfpathlineto{\pgfqpoint{2.872842in}{0.842457in}}%
\pgfpathlineto{\pgfqpoint{2.871284in}{0.855909in}}%
\pgfpathlineto{\pgfqpoint{2.855107in}{0.854082in}}%
\pgfpathlineto{\pgfqpoint{2.854444in}{0.865881in}}%
\pgfpathlineto{\pgfqpoint{2.851231in}{0.869522in}}%
\pgfpathlineto{\pgfqpoint{2.839212in}{0.870674in}}%
\pgfpathclose%
\pgfusepath{fill}%
\end{pgfscope}%
\begin{pgfscope}%
\pgfpathrectangle{\pgfqpoint{0.100000in}{0.100000in}}{\pgfqpoint{3.608454in}{2.310000in}}%
\pgfusepath{clip}%
\pgfsetbuttcap%
\pgfsetmiterjoin%
\definecolor{currentfill}{rgb}{0.000000,0.419608,0.790196}%
\pgfsetfillcolor{currentfill}%
\pgfsetlinewidth{0.000000pt}%
\definecolor{currentstroke}{rgb}{0.000000,0.000000,0.000000}%
\pgfsetstrokecolor{currentstroke}%
\pgfsetstrokeopacity{0.000000}%
\pgfsetdash{}{0pt}%
\pgfpathmoveto{\pgfqpoint{1.981983in}{1.828333in}}%
\pgfpathlineto{\pgfqpoint{1.939662in}{1.829664in}}%
\pgfpathlineto{\pgfqpoint{1.941869in}{1.891889in}}%
\pgfpathlineto{\pgfqpoint{1.983409in}{1.890641in}}%
\pgfpathlineto{\pgfqpoint{2.020229in}{1.889881in}}%
\pgfpathlineto{\pgfqpoint{2.018973in}{1.880728in}}%
\pgfpathlineto{\pgfqpoint{2.015031in}{1.874970in}}%
\pgfpathlineto{\pgfqpoint{2.004519in}{1.867267in}}%
\pgfpathlineto{\pgfqpoint{2.003389in}{1.863983in}}%
\pgfpathlineto{\pgfqpoint{2.012969in}{1.848244in}}%
\pgfpathlineto{\pgfqpoint{2.021780in}{1.845315in}}%
\pgfpathlineto{\pgfqpoint{2.024583in}{1.841364in}}%
\pgfpathlineto{\pgfqpoint{1.995378in}{1.841947in}}%
\pgfpathlineto{\pgfqpoint{1.994509in}{1.839579in}}%
\pgfpathlineto{\pgfqpoint{1.982265in}{1.839933in}}%
\pgfpathlineto{\pgfqpoint{1.981983in}{1.828333in}}%
\pgfpathclose%
\pgfusepath{fill}%
\end{pgfscope}%
\begin{pgfscope}%
\pgfpathrectangle{\pgfqpoint{0.100000in}{0.100000in}}{\pgfqpoint{3.608454in}{2.310000in}}%
\pgfusepath{clip}%
\pgfsetbuttcap%
\pgfsetmiterjoin%
\definecolor{currentfill}{rgb}{0.000000,0.682353,0.658824}%
\pgfsetfillcolor{currentfill}%
\pgfsetlinewidth{0.000000pt}%
\definecolor{currentstroke}{rgb}{0.000000,0.000000,0.000000}%
\pgfsetstrokecolor{currentstroke}%
\pgfsetstrokeopacity{0.000000}%
\pgfsetdash{}{0pt}%
\pgfpathmoveto{\pgfqpoint{2.971951in}{1.538265in}}%
\pgfpathlineto{\pgfqpoint{2.963825in}{1.532458in}}%
\pgfpathlineto{\pgfqpoint{2.952109in}{1.531961in}}%
\pgfpathlineto{\pgfqpoint{2.948242in}{1.534882in}}%
\pgfpathlineto{\pgfqpoint{2.947280in}{1.541412in}}%
\pgfpathlineto{\pgfqpoint{2.937196in}{1.539921in}}%
\pgfpathlineto{\pgfqpoint{2.934085in}{1.560268in}}%
\pgfpathlineto{\pgfqpoint{2.939221in}{1.561031in}}%
\pgfpathlineto{\pgfqpoint{2.932782in}{1.601168in}}%
\pgfpathlineto{\pgfqpoint{2.961131in}{1.605506in}}%
\pgfpathlineto{\pgfqpoint{2.984410in}{1.607734in}}%
\pgfpathlineto{\pgfqpoint{2.991632in}{1.609704in}}%
\pgfpathlineto{\pgfqpoint{2.995768in}{1.584878in}}%
\pgfpathlineto{\pgfqpoint{2.988249in}{1.569893in}}%
\pgfpathlineto{\pgfqpoint{2.990516in}{1.558619in}}%
\pgfpathlineto{\pgfqpoint{2.969282in}{1.554858in}}%
\pgfpathlineto{\pgfqpoint{2.971951in}{1.538265in}}%
\pgfpathclose%
\pgfusepath{fill}%
\end{pgfscope}%
\begin{pgfscope}%
\pgfpathrectangle{\pgfqpoint{0.100000in}{0.100000in}}{\pgfqpoint{3.608454in}{2.310000in}}%
\pgfusepath{clip}%
\pgfsetbuttcap%
\pgfsetmiterjoin%
\definecolor{currentfill}{rgb}{0.000000,0.443137,0.778431}%
\pgfsetfillcolor{currentfill}%
\pgfsetlinewidth{0.000000pt}%
\definecolor{currentstroke}{rgb}{0.000000,0.000000,0.000000}%
\pgfsetstrokecolor{currentstroke}%
\pgfsetstrokeopacity{0.000000}%
\pgfsetdash{}{0pt}%
\pgfpathmoveto{\pgfqpoint{1.643808in}{1.676503in}}%
\pgfpathlineto{\pgfqpoint{1.583107in}{1.682203in}}%
\pgfpathlineto{\pgfqpoint{1.586973in}{1.721839in}}%
\pgfpathlineto{\pgfqpoint{1.592333in}{1.775294in}}%
\pgfpathlineto{\pgfqpoint{1.598608in}{1.839877in}}%
\pgfpathlineto{\pgfqpoint{1.601301in}{1.856823in}}%
\pgfpathlineto{\pgfqpoint{1.661648in}{1.851140in}}%
\pgfpathlineto{\pgfqpoint{1.660454in}{1.837408in}}%
\pgfpathlineto{\pgfqpoint{1.714041in}{1.832923in}}%
\pgfpathlineto{\pgfqpoint{1.710799in}{1.791190in}}%
\pgfpathlineto{\pgfqpoint{1.707249in}{1.750308in}}%
\pgfpathlineto{\pgfqpoint{1.704828in}{1.727690in}}%
\pgfpathlineto{\pgfqpoint{1.697918in}{1.727667in}}%
\pgfpathlineto{\pgfqpoint{1.695616in}{1.726789in}}%
\pgfpathlineto{\pgfqpoint{1.659487in}{1.729827in}}%
\pgfpathlineto{\pgfqpoint{1.654175in}{1.729234in}}%
\pgfpathlineto{\pgfqpoint{1.648016in}{1.724287in}}%
\pgfpathlineto{\pgfqpoint{1.643808in}{1.676503in}}%
\pgfpathclose%
\pgfusepath{fill}%
\end{pgfscope}%
\begin{pgfscope}%
\pgfpathrectangle{\pgfqpoint{0.100000in}{0.100000in}}{\pgfqpoint{3.608454in}{2.310000in}}%
\pgfusepath{clip}%
\pgfsetbuttcap%
\pgfsetmiterjoin%
\definecolor{currentfill}{rgb}{0.000000,0.768627,0.615686}%
\pgfsetfillcolor{currentfill}%
\pgfsetlinewidth{0.000000pt}%
\definecolor{currentstroke}{rgb}{0.000000,0.000000,0.000000}%
\pgfsetstrokecolor{currentstroke}%
\pgfsetstrokeopacity{0.000000}%
\pgfsetdash{}{0pt}%
\pgfpathmoveto{\pgfqpoint{2.867721in}{1.199764in}}%
\pgfpathlineto{\pgfqpoint{2.859737in}{1.192156in}}%
\pgfpathlineto{\pgfqpoint{2.856311in}{1.184739in}}%
\pgfpathlineto{\pgfqpoint{2.851139in}{1.181704in}}%
\pgfpathlineto{\pgfqpoint{2.846368in}{1.183090in}}%
\pgfpathlineto{\pgfqpoint{2.836596in}{1.177507in}}%
\pgfpathlineto{\pgfqpoint{2.827195in}{1.174638in}}%
\pgfpathlineto{\pgfqpoint{2.814063in}{1.180206in}}%
\pgfpathlineto{\pgfqpoint{2.810250in}{1.179188in}}%
\pgfpathlineto{\pgfqpoint{2.804581in}{1.191394in}}%
\pgfpathlineto{\pgfqpoint{2.806029in}{1.196000in}}%
\pgfpathlineto{\pgfqpoint{2.812724in}{1.202234in}}%
\pgfpathlineto{\pgfqpoint{2.815210in}{1.209500in}}%
\pgfpathlineto{\pgfqpoint{2.831285in}{1.222519in}}%
\pgfpathlineto{\pgfqpoint{2.833882in}{1.215723in}}%
\pgfpathlineto{\pgfqpoint{2.833845in}{1.206501in}}%
\pgfpathlineto{\pgfqpoint{2.837048in}{1.201444in}}%
\pgfpathlineto{\pgfqpoint{2.828062in}{1.198088in}}%
\pgfpathlineto{\pgfqpoint{2.827558in}{1.194906in}}%
\pgfpathlineto{\pgfqpoint{2.867721in}{1.199764in}}%
\pgfpathclose%
\pgfusepath{fill}%
\end{pgfscope}%
\begin{pgfscope}%
\pgfpathrectangle{\pgfqpoint{0.100000in}{0.100000in}}{\pgfqpoint{3.608454in}{2.310000in}}%
\pgfusepath{clip}%
\pgfsetbuttcap%
\pgfsetmiterjoin%
\definecolor{currentfill}{rgb}{0.000000,0.784314,0.607843}%
\pgfsetfillcolor{currentfill}%
\pgfsetlinewidth{0.000000pt}%
\definecolor{currentstroke}{rgb}{0.000000,0.000000,0.000000}%
\pgfsetstrokecolor{currentstroke}%
\pgfsetstrokeopacity{0.000000}%
\pgfsetdash{}{0pt}%
\pgfpathmoveto{\pgfqpoint{2.843160in}{1.055221in}}%
\pgfpathlineto{\pgfqpoint{2.833248in}{1.051039in}}%
\pgfpathlineto{\pgfqpoint{2.830948in}{1.045651in}}%
\pgfpathlineto{\pgfqpoint{2.825426in}{1.046467in}}%
\pgfpathlineto{\pgfqpoint{2.820073in}{1.037913in}}%
\pgfpathlineto{\pgfqpoint{2.812171in}{1.037511in}}%
\pgfpathlineto{\pgfqpoint{2.814885in}{1.050035in}}%
\pgfpathlineto{\pgfqpoint{2.808303in}{1.061863in}}%
\pgfpathlineto{\pgfqpoint{2.798375in}{1.062974in}}%
\pgfpathlineto{\pgfqpoint{2.798205in}{1.081874in}}%
\pgfpathlineto{\pgfqpoint{2.802072in}{1.085804in}}%
\pgfpathlineto{\pgfqpoint{2.807766in}{1.084602in}}%
\pgfpathlineto{\pgfqpoint{2.814347in}{1.089123in}}%
\pgfpathlineto{\pgfqpoint{2.819581in}{1.083755in}}%
\pgfpathlineto{\pgfqpoint{2.827437in}{1.087830in}}%
\pgfpathlineto{\pgfqpoint{2.835727in}{1.088275in}}%
\pgfpathlineto{\pgfqpoint{2.834354in}{1.080705in}}%
\pgfpathlineto{\pgfqpoint{2.840206in}{1.081206in}}%
\pgfpathlineto{\pgfqpoint{2.850115in}{1.070383in}}%
\pgfpathlineto{\pgfqpoint{2.842470in}{1.059416in}}%
\pgfpathlineto{\pgfqpoint{2.843160in}{1.055221in}}%
\pgfpathclose%
\pgfusepath{fill}%
\end{pgfscope}%
\begin{pgfscope}%
\pgfpathrectangle{\pgfqpoint{0.100000in}{0.100000in}}{\pgfqpoint{3.608454in}{2.310000in}}%
\pgfusepath{clip}%
\pgfsetbuttcap%
\pgfsetmiterjoin%
\definecolor{currentfill}{rgb}{0.000000,0.478431,0.760784}%
\pgfsetfillcolor{currentfill}%
\pgfsetlinewidth{0.000000pt}%
\definecolor{currentstroke}{rgb}{0.000000,0.000000,0.000000}%
\pgfsetstrokecolor{currentstroke}%
\pgfsetstrokeopacity{0.000000}%
\pgfsetdash{}{0pt}%
\pgfpathmoveto{\pgfqpoint{2.872479in}{1.551436in}}%
\pgfpathlineto{\pgfqpoint{2.869199in}{1.556678in}}%
\pgfpathlineto{\pgfqpoint{2.868477in}{1.562448in}}%
\pgfpathlineto{\pgfqpoint{2.873516in}{1.568184in}}%
\pgfpathlineto{\pgfqpoint{2.879290in}{1.569013in}}%
\pgfpathlineto{\pgfqpoint{2.878552in}{1.574944in}}%
\pgfpathlineto{\pgfqpoint{2.884117in}{1.575727in}}%
\pgfpathlineto{\pgfqpoint{2.883315in}{1.581642in}}%
\pgfpathlineto{\pgfqpoint{2.877765in}{1.580902in}}%
\pgfpathlineto{\pgfqpoint{2.876219in}{1.593055in}}%
\pgfpathlineto{\pgfqpoint{2.889848in}{1.593630in}}%
\pgfpathlineto{\pgfqpoint{2.897159in}{1.600476in}}%
\pgfpathlineto{\pgfqpoint{2.905051in}{1.610697in}}%
\pgfpathlineto{\pgfqpoint{2.928077in}{1.628592in}}%
\pgfpathlineto{\pgfqpoint{2.934104in}{1.630691in}}%
\pgfpathlineto{\pgfqpoint{2.955120in}{1.642813in}}%
\pgfpathlineto{\pgfqpoint{2.961131in}{1.605506in}}%
\pgfpathlineto{\pgfqpoint{2.932782in}{1.601168in}}%
\pgfpathlineto{\pgfqpoint{2.939221in}{1.561031in}}%
\pgfpathlineto{\pgfqpoint{2.934085in}{1.560268in}}%
\pgfpathlineto{\pgfqpoint{2.914636in}{1.557327in}}%
\pgfpathlineto{\pgfqpoint{2.915375in}{1.550896in}}%
\pgfpathlineto{\pgfqpoint{2.901826in}{1.549489in}}%
\pgfpathlineto{\pgfqpoint{2.901003in}{1.555251in}}%
\pgfpathlineto{\pgfqpoint{2.872479in}{1.551436in}}%
\pgfpathclose%
\pgfusepath{fill}%
\end{pgfscope}%
\begin{pgfscope}%
\pgfpathrectangle{\pgfqpoint{0.100000in}{0.100000in}}{\pgfqpoint{3.608454in}{2.310000in}}%
\pgfusepath{clip}%
\pgfsetbuttcap%
\pgfsetmiterjoin%
\definecolor{currentfill}{rgb}{0.000000,0.639216,0.680392}%
\pgfsetfillcolor{currentfill}%
\pgfsetlinewidth{0.000000pt}%
\definecolor{currentstroke}{rgb}{0.000000,0.000000,0.000000}%
\pgfsetstrokecolor{currentstroke}%
\pgfsetstrokeopacity{0.000000}%
\pgfsetdash{}{0pt}%
\pgfpathmoveto{\pgfqpoint{0.737189in}{2.001359in}}%
\pgfpathlineto{\pgfqpoint{0.730749in}{2.003174in}}%
\pgfpathlineto{\pgfqpoint{0.736122in}{2.015946in}}%
\pgfpathlineto{\pgfqpoint{0.729534in}{2.017963in}}%
\pgfpathlineto{\pgfqpoint{0.731442in}{2.024554in}}%
\pgfpathlineto{\pgfqpoint{0.724757in}{2.026561in}}%
\pgfpathlineto{\pgfqpoint{0.737141in}{2.068948in}}%
\pgfpathlineto{\pgfqpoint{0.744663in}{2.068773in}}%
\pgfpathlineto{\pgfqpoint{0.749167in}{2.084439in}}%
\pgfpathlineto{\pgfqpoint{0.761418in}{2.129669in}}%
\pgfpathlineto{\pgfqpoint{0.774845in}{2.127276in}}%
\pgfpathlineto{\pgfqpoint{0.782427in}{2.131997in}}%
\pgfpathlineto{\pgfqpoint{0.784697in}{2.127385in}}%
\pgfpathlineto{\pgfqpoint{0.790335in}{2.130635in}}%
\pgfpathlineto{\pgfqpoint{0.823108in}{2.121596in}}%
\pgfpathlineto{\pgfqpoint{0.851613in}{2.114354in}}%
\pgfpathlineto{\pgfqpoint{0.849420in}{2.114482in}}%
\pgfpathlineto{\pgfqpoint{0.849927in}{2.099792in}}%
\pgfpathlineto{\pgfqpoint{0.868296in}{2.100626in}}%
\pgfpathlineto{\pgfqpoint{0.864637in}{2.088901in}}%
\pgfpathlineto{\pgfqpoint{0.871344in}{2.087217in}}%
\pgfpathlineto{\pgfqpoint{0.872020in}{2.076043in}}%
\pgfpathlineto{\pgfqpoint{0.875412in}{2.075250in}}%
\pgfpathlineto{\pgfqpoint{0.869214in}{2.048910in}}%
\pgfpathlineto{\pgfqpoint{0.849176in}{2.054195in}}%
\pgfpathlineto{\pgfqpoint{0.846645in}{2.043323in}}%
\pgfpathlineto{\pgfqpoint{0.841907in}{2.040973in}}%
\pgfpathlineto{\pgfqpoint{0.838162in}{2.031358in}}%
\pgfpathlineto{\pgfqpoint{0.835457in}{2.032072in}}%
\pgfpathlineto{\pgfqpoint{0.831103in}{2.015303in}}%
\pgfpathlineto{\pgfqpoint{0.825859in}{2.013085in}}%
\pgfpathlineto{\pgfqpoint{0.813450in}{2.016488in}}%
\pgfpathlineto{\pgfqpoint{0.811964in}{2.010809in}}%
\pgfpathlineto{\pgfqpoint{0.797207in}{2.013901in}}%
\pgfpathlineto{\pgfqpoint{0.796435in}{2.001902in}}%
\pgfpathlineto{\pgfqpoint{0.802344in}{2.000341in}}%
\pgfpathlineto{\pgfqpoint{0.799750in}{1.984499in}}%
\pgfpathlineto{\pgfqpoint{0.737189in}{2.001359in}}%
\pgfpathclose%
\pgfusepath{fill}%
\end{pgfscope}%
\begin{pgfscope}%
\pgfpathrectangle{\pgfqpoint{0.100000in}{0.100000in}}{\pgfqpoint{3.608454in}{2.310000in}}%
\pgfusepath{clip}%
\pgfsetbuttcap%
\pgfsetmiterjoin%
\definecolor{currentfill}{rgb}{0.000000,0.670588,0.664706}%
\pgfsetfillcolor{currentfill}%
\pgfsetlinewidth{0.000000pt}%
\definecolor{currentstroke}{rgb}{0.000000,0.000000,0.000000}%
\pgfsetstrokecolor{currentstroke}%
\pgfsetstrokeopacity{0.000000}%
\pgfsetdash{}{0pt}%
\pgfpathmoveto{\pgfqpoint{1.601877in}{1.254860in}}%
\pgfpathlineto{\pgfqpoint{1.596724in}{1.204258in}}%
\pgfpathlineto{\pgfqpoint{1.538731in}{1.209408in}}%
\pgfpathlineto{\pgfqpoint{1.513551in}{1.211748in}}%
\pgfpathlineto{\pgfqpoint{1.466660in}{1.216957in}}%
\pgfpathlineto{\pgfqpoint{1.469341in}{1.240362in}}%
\pgfpathlineto{\pgfqpoint{1.471264in}{1.248705in}}%
\pgfpathlineto{\pgfqpoint{1.469792in}{1.261228in}}%
\pgfpathlineto{\pgfqpoint{1.463748in}{1.269672in}}%
\pgfpathlineto{\pgfqpoint{1.450742in}{1.268060in}}%
\pgfpathlineto{\pgfqpoint{1.457227in}{1.286192in}}%
\pgfpathlineto{\pgfqpoint{1.454881in}{1.289984in}}%
\pgfpathlineto{\pgfqpoint{1.458524in}{1.289130in}}%
\pgfpathlineto{\pgfqpoint{1.470019in}{1.293572in}}%
\pgfpathlineto{\pgfqpoint{1.474843in}{1.297491in}}%
\pgfpathlineto{\pgfqpoint{1.481289in}{1.288540in}}%
\pgfpathlineto{\pgfqpoint{1.483235in}{1.285604in}}%
\pgfpathlineto{\pgfqpoint{1.506183in}{1.284643in}}%
\pgfpathlineto{\pgfqpoint{1.541303in}{1.267794in}}%
\pgfpathlineto{\pgfqpoint{1.540663in}{1.260697in}}%
\pgfpathlineto{\pgfqpoint{1.601877in}{1.254860in}}%
\pgfpathclose%
\pgfusepath{fill}%
\end{pgfscope}%
\begin{pgfscope}%
\pgfpathrectangle{\pgfqpoint{0.100000in}{0.100000in}}{\pgfqpoint{3.608454in}{2.310000in}}%
\pgfusepath{clip}%
\pgfsetbuttcap%
\pgfsetmiterjoin%
\definecolor{currentfill}{rgb}{0.000000,0.380392,0.809804}%
\pgfsetfillcolor{currentfill}%
\pgfsetlinewidth{0.000000pt}%
\definecolor{currentstroke}{rgb}{0.000000,0.000000,0.000000}%
\pgfsetstrokecolor{currentstroke}%
\pgfsetstrokeopacity{0.000000}%
\pgfsetdash{}{0pt}%
\pgfpathmoveto{\pgfqpoint{2.347355in}{1.241765in}}%
\pgfpathlineto{\pgfqpoint{2.337514in}{1.241571in}}%
\pgfpathlineto{\pgfqpoint{2.337381in}{1.248460in}}%
\pgfpathlineto{\pgfqpoint{2.323611in}{1.248060in}}%
\pgfpathlineto{\pgfqpoint{2.322748in}{1.276780in}}%
\pgfpathlineto{\pgfqpoint{2.315298in}{1.276599in}}%
\pgfpathlineto{\pgfqpoint{2.314915in}{1.287249in}}%
\pgfpathlineto{\pgfqpoint{2.313973in}{1.320063in}}%
\pgfpathlineto{\pgfqpoint{2.319587in}{1.318096in}}%
\pgfpathlineto{\pgfqpoint{2.327621in}{1.321058in}}%
\pgfpathlineto{\pgfqpoint{2.327371in}{1.331768in}}%
\pgfpathlineto{\pgfqpoint{2.336652in}{1.332026in}}%
\pgfpathlineto{\pgfqpoint{2.335984in}{1.355415in}}%
\pgfpathlineto{\pgfqpoint{2.340562in}{1.355530in}}%
\pgfpathlineto{\pgfqpoint{2.340427in}{1.362513in}}%
\pgfpathlineto{\pgfqpoint{2.368642in}{1.363523in}}%
\pgfpathlineto{\pgfqpoint{2.371543in}{1.352918in}}%
\pgfpathlineto{\pgfqpoint{2.369779in}{1.350114in}}%
\pgfpathlineto{\pgfqpoint{2.373739in}{1.339722in}}%
\pgfpathlineto{\pgfqpoint{2.380869in}{1.336490in}}%
\pgfpathlineto{\pgfqpoint{2.386232in}{1.344093in}}%
\pgfpathlineto{\pgfqpoint{2.398503in}{1.340979in}}%
\pgfpathlineto{\pgfqpoint{2.407445in}{1.335820in}}%
\pgfpathlineto{\pgfqpoint{2.402286in}{1.326271in}}%
\pgfpathlineto{\pgfqpoint{2.404058in}{1.320682in}}%
\pgfpathlineto{\pgfqpoint{2.433309in}{1.321900in}}%
\pgfpathlineto{\pgfqpoint{2.435549in}{1.287521in}}%
\pgfpathlineto{\pgfqpoint{2.414877in}{1.286654in}}%
\pgfpathlineto{\pgfqpoint{2.415279in}{1.279749in}}%
\pgfpathlineto{\pgfqpoint{2.404880in}{1.275417in}}%
\pgfpathlineto{\pgfqpoint{2.397241in}{1.275533in}}%
\pgfpathlineto{\pgfqpoint{2.392056in}{1.271141in}}%
\pgfpathlineto{\pgfqpoint{2.381801in}{1.267561in}}%
\pgfpathlineto{\pgfqpoint{2.368816in}{1.282820in}}%
\pgfpathlineto{\pgfqpoint{2.349233in}{1.281942in}}%
\pgfpathlineto{\pgfqpoint{2.350544in}{1.245278in}}%
\pgfpathlineto{\pgfqpoint{2.347355in}{1.241765in}}%
\pgfpathclose%
\pgfusepath{fill}%
\end{pgfscope}%
\begin{pgfscope}%
\pgfpathrectangle{\pgfqpoint{0.100000in}{0.100000in}}{\pgfqpoint{3.608454in}{2.310000in}}%
\pgfusepath{clip}%
\pgfsetbuttcap%
\pgfsetmiterjoin%
\definecolor{currentfill}{rgb}{0.000000,0.301961,0.849020}%
\pgfsetfillcolor{currentfill}%
\pgfsetlinewidth{0.000000pt}%
\definecolor{currentstroke}{rgb}{0.000000,0.000000,0.000000}%
\pgfsetstrokecolor{currentstroke}%
\pgfsetstrokeopacity{0.000000}%
\pgfsetdash{}{0pt}%
\pgfpathmoveto{\pgfqpoint{2.130827in}{1.263573in}}%
\pgfpathlineto{\pgfqpoint{2.101955in}{1.263584in}}%
\pgfpathlineto{\pgfqpoint{2.102752in}{1.281057in}}%
\pgfpathlineto{\pgfqpoint{2.102737in}{1.315533in}}%
\pgfpathlineto{\pgfqpoint{2.103312in}{1.318980in}}%
\pgfpathlineto{\pgfqpoint{2.103336in}{1.338237in}}%
\pgfpathlineto{\pgfqpoint{2.112374in}{1.338940in}}%
\pgfpathlineto{\pgfqpoint{2.112893in}{1.355625in}}%
\pgfpathlineto{\pgfqpoint{2.116197in}{1.360877in}}%
\pgfpathlineto{\pgfqpoint{2.107556in}{1.374791in}}%
\pgfpathlineto{\pgfqpoint{2.100521in}{1.381808in}}%
\pgfpathlineto{\pgfqpoint{2.131100in}{1.381764in}}%
\pgfpathlineto{\pgfqpoint{2.131152in}{1.375734in}}%
\pgfpathlineto{\pgfqpoint{2.154965in}{1.375795in}}%
\pgfpathlineto{\pgfqpoint{2.155044in}{1.381527in}}%
\pgfpathlineto{\pgfqpoint{2.182439in}{1.381621in}}%
\pgfpathlineto{\pgfqpoint{2.182827in}{1.356530in}}%
\pgfpathlineto{\pgfqpoint{2.188988in}{1.359699in}}%
\pgfpathlineto{\pgfqpoint{2.198627in}{1.357636in}}%
\pgfpathlineto{\pgfqpoint{2.199219in}{1.335898in}}%
\pgfpathlineto{\pgfqpoint{2.198808in}{1.305319in}}%
\pgfpathlineto{\pgfqpoint{2.164617in}{1.305750in}}%
\pgfpathlineto{\pgfqpoint{2.164228in}{1.277853in}}%
\pgfpathlineto{\pgfqpoint{2.165424in}{1.263791in}}%
\pgfpathlineto{\pgfqpoint{2.154102in}{1.264813in}}%
\pgfpathlineto{\pgfqpoint{2.130815in}{1.265402in}}%
\pgfpathlineto{\pgfqpoint{2.130827in}{1.263573in}}%
\pgfpathclose%
\pgfusepath{fill}%
\end{pgfscope}%
\begin{pgfscope}%
\pgfpathrectangle{\pgfqpoint{0.100000in}{0.100000in}}{\pgfqpoint{3.608454in}{2.310000in}}%
\pgfusepath{clip}%
\pgfsetbuttcap%
\pgfsetmiterjoin%
\definecolor{currentfill}{rgb}{0.000000,0.611765,0.694118}%
\pgfsetfillcolor{currentfill}%
\pgfsetlinewidth{0.000000pt}%
\definecolor{currentstroke}{rgb}{0.000000,0.000000,0.000000}%
\pgfsetstrokecolor{currentstroke}%
\pgfsetstrokeopacity{0.000000}%
\pgfsetdash{}{0pt}%
\pgfpathmoveto{\pgfqpoint{1.522840in}{2.169802in}}%
\pgfpathlineto{\pgfqpoint{1.578022in}{2.163356in}}%
\pgfpathlineto{\pgfqpoint{1.578243in}{2.156172in}}%
\pgfpathlineto{\pgfqpoint{1.575926in}{2.135412in}}%
\pgfpathlineto{\pgfqpoint{1.578597in}{2.128096in}}%
\pgfpathlineto{\pgfqpoint{1.534783in}{2.133134in}}%
\pgfpathlineto{\pgfqpoint{1.523357in}{2.134506in}}%
\pgfpathlineto{\pgfqpoint{1.525884in}{2.155265in}}%
\pgfpathlineto{\pgfqpoint{1.519036in}{2.156103in}}%
\pgfpathlineto{\pgfqpoint{1.522840in}{2.169802in}}%
\pgfpathclose%
\pgfusepath{fill}%
\end{pgfscope}%
\begin{pgfscope}%
\pgfpathrectangle{\pgfqpoint{0.100000in}{0.100000in}}{\pgfqpoint{3.608454in}{2.310000in}}%
\pgfusepath{clip}%
\pgfsetbuttcap%
\pgfsetmiterjoin%
\definecolor{currentfill}{rgb}{0.000000,0.419608,0.790196}%
\pgfsetfillcolor{currentfill}%
\pgfsetlinewidth{0.000000pt}%
\definecolor{currentstroke}{rgb}{0.000000,0.000000,0.000000}%
\pgfsetstrokecolor{currentstroke}%
\pgfsetstrokeopacity{0.000000}%
\pgfsetdash{}{0pt}%
\pgfpathmoveto{\pgfqpoint{3.062330in}{0.671592in}}%
\pgfpathlineto{\pgfqpoint{3.055227in}{0.666862in}}%
\pgfpathlineto{\pgfqpoint{3.041523in}{0.664546in}}%
\pgfpathlineto{\pgfqpoint{3.039796in}{0.675377in}}%
\pgfpathlineto{\pgfqpoint{3.035034in}{0.681162in}}%
\pgfpathlineto{\pgfqpoint{3.019158in}{0.678567in}}%
\pgfpathlineto{\pgfqpoint{3.011416in}{0.670241in}}%
\pgfpathlineto{\pgfqpoint{3.004229in}{0.666876in}}%
\pgfpathlineto{\pgfqpoint{2.999393in}{0.700496in}}%
\pgfpathlineto{\pgfqpoint{2.971461in}{0.696012in}}%
\pgfpathlineto{\pgfqpoint{2.966459in}{0.731456in}}%
\pgfpathlineto{\pgfqpoint{2.983300in}{0.732552in}}%
\pgfpathlineto{\pgfqpoint{2.985291in}{0.750228in}}%
\pgfpathlineto{\pgfqpoint{2.966799in}{0.748504in}}%
\pgfpathlineto{\pgfqpoint{2.964698in}{0.765826in}}%
\pgfpathlineto{\pgfqpoint{2.983960in}{0.768316in}}%
\pgfpathlineto{\pgfqpoint{2.988034in}{0.774133in}}%
\pgfpathlineto{\pgfqpoint{2.996804in}{0.773153in}}%
\pgfpathlineto{\pgfqpoint{3.005604in}{0.782838in}}%
\pgfpathlineto{\pgfqpoint{3.020788in}{0.780741in}}%
\pgfpathlineto{\pgfqpoint{3.030760in}{0.775066in}}%
\pgfpathlineto{\pgfqpoint{3.034315in}{0.767379in}}%
\pgfpathlineto{\pgfqpoint{3.032529in}{0.754216in}}%
\pgfpathlineto{\pgfqpoint{3.035355in}{0.750593in}}%
\pgfpathlineto{\pgfqpoint{3.035557in}{0.741809in}}%
\pgfpathlineto{\pgfqpoint{3.041274in}{0.727152in}}%
\pgfpathlineto{\pgfqpoint{3.042489in}{0.719793in}}%
\pgfpathlineto{\pgfqpoint{3.052515in}{0.693343in}}%
\pgfpathlineto{\pgfqpoint{3.062330in}{0.671592in}}%
\pgfpathclose%
\pgfusepath{fill}%
\end{pgfscope}%
\begin{pgfscope}%
\pgfpathrectangle{\pgfqpoint{0.100000in}{0.100000in}}{\pgfqpoint{3.608454in}{2.310000in}}%
\pgfusepath{clip}%
\pgfsetbuttcap%
\pgfsetmiterjoin%
\definecolor{currentfill}{rgb}{0.000000,0.721569,0.639216}%
\pgfsetfillcolor{currentfill}%
\pgfsetlinewidth{0.000000pt}%
\definecolor{currentstroke}{rgb}{0.000000,0.000000,0.000000}%
\pgfsetstrokecolor{currentstroke}%
\pgfsetstrokeopacity{0.000000}%
\pgfsetdash{}{0pt}%
\pgfpathmoveto{\pgfqpoint{1.590719in}{1.104616in}}%
\pgfpathlineto{\pgfqpoint{1.589893in}{1.095406in}}%
\pgfpathlineto{\pgfqpoint{1.585141in}{1.042675in}}%
\pgfpathlineto{\pgfqpoint{1.569604in}{1.044093in}}%
\pgfpathlineto{\pgfqpoint{1.568889in}{1.037039in}}%
\pgfpathlineto{\pgfqpoint{1.562067in}{1.037697in}}%
\pgfpathlineto{\pgfqpoint{1.561377in}{1.030839in}}%
\pgfpathlineto{\pgfqpoint{1.554565in}{1.031506in}}%
\pgfpathlineto{\pgfqpoint{1.553889in}{1.024630in}}%
\pgfpathlineto{\pgfqpoint{1.540290in}{1.025994in}}%
\pgfpathlineto{\pgfqpoint{1.539569in}{1.019066in}}%
\pgfpathlineto{\pgfqpoint{1.523829in}{1.020634in}}%
\pgfpathlineto{\pgfqpoint{1.512139in}{1.021816in}}%
\pgfpathlineto{\pgfqpoint{1.516667in}{1.064131in}}%
\pgfpathlineto{\pgfqpoint{1.534671in}{1.070306in}}%
\pgfpathlineto{\pgfqpoint{1.548858in}{1.068831in}}%
\pgfpathlineto{\pgfqpoint{1.550211in}{1.083114in}}%
\pgfpathlineto{\pgfqpoint{1.544598in}{1.083669in}}%
\pgfpathlineto{\pgfqpoint{1.531357in}{1.115154in}}%
\pgfpathlineto{\pgfqpoint{1.528359in}{1.113844in}}%
\pgfpathlineto{\pgfqpoint{1.506250in}{1.116084in}}%
\pgfpathlineto{\pgfqpoint{1.508110in}{1.125109in}}%
\pgfpathlineto{\pgfqpoint{1.508502in}{1.136752in}}%
\pgfpathlineto{\pgfqpoint{1.505428in}{1.150972in}}%
\pgfpathlineto{\pgfqpoint{1.532570in}{1.148149in}}%
\pgfpathlineto{\pgfqpoint{1.532225in}{1.144723in}}%
\pgfpathlineto{\pgfqpoint{1.545891in}{1.143339in}}%
\pgfpathlineto{\pgfqpoint{1.545193in}{1.136393in}}%
\pgfpathlineto{\pgfqpoint{1.572537in}{1.133851in}}%
\pgfpathlineto{\pgfqpoint{1.569376in}{1.106608in}}%
\pgfpathlineto{\pgfqpoint{1.590719in}{1.104616in}}%
\pgfpathclose%
\pgfusepath{fill}%
\end{pgfscope}%
\begin{pgfscope}%
\pgfpathrectangle{\pgfqpoint{0.100000in}{0.100000in}}{\pgfqpoint{3.608454in}{2.310000in}}%
\pgfusepath{clip}%
\pgfsetbuttcap%
\pgfsetmiterjoin%
\definecolor{currentfill}{rgb}{0.000000,0.631373,0.684314}%
\pgfsetfillcolor{currentfill}%
\pgfsetlinewidth{0.000000pt}%
\definecolor{currentstroke}{rgb}{0.000000,0.000000,0.000000}%
\pgfsetstrokecolor{currentstroke}%
\pgfsetstrokeopacity{0.000000}%
\pgfsetdash{}{0pt}%
\pgfpathmoveto{\pgfqpoint{3.130816in}{1.638979in}}%
\pgfpathlineto{\pgfqpoint{3.123343in}{1.674373in}}%
\pgfpathlineto{\pgfqpoint{3.115261in}{1.672766in}}%
\pgfpathlineto{\pgfqpoint{3.109751in}{1.709791in}}%
\pgfpathlineto{\pgfqpoint{3.111728in}{1.718970in}}%
\pgfpathlineto{\pgfqpoint{3.141312in}{1.724493in}}%
\pgfpathlineto{\pgfqpoint{3.142173in}{1.718464in}}%
\pgfpathlineto{\pgfqpoint{3.156855in}{1.718354in}}%
\pgfpathlineto{\pgfqpoint{3.155225in}{1.724423in}}%
\pgfpathlineto{\pgfqpoint{3.173622in}{1.728849in}}%
\pgfpathlineto{\pgfqpoint{3.167852in}{1.733612in}}%
\pgfpathlineto{\pgfqpoint{3.190729in}{1.738354in}}%
\pgfpathlineto{\pgfqpoint{3.194930in}{1.721757in}}%
\pgfpathlineto{\pgfqpoint{3.196916in}{1.713203in}}%
\pgfpathlineto{\pgfqpoint{3.187096in}{1.712911in}}%
\pgfpathlineto{\pgfqpoint{3.180627in}{1.708614in}}%
\pgfpathlineto{\pgfqpoint{3.183986in}{1.686549in}}%
\pgfpathlineto{\pgfqpoint{3.162699in}{1.682336in}}%
\pgfpathlineto{\pgfqpoint{3.172073in}{1.651526in}}%
\pgfpathlineto{\pgfqpoint{3.167806in}{1.647018in}}%
\pgfpathlineto{\pgfqpoint{3.130816in}{1.638979in}}%
\pgfpathclose%
\pgfusepath{fill}%
\end{pgfscope}%
\begin{pgfscope}%
\pgfpathrectangle{\pgfqpoint{0.100000in}{0.100000in}}{\pgfqpoint{3.608454in}{2.310000in}}%
\pgfusepath{clip}%
\pgfsetbuttcap%
\pgfsetmiterjoin%
\definecolor{currentfill}{rgb}{0.000000,0.796078,0.601961}%
\pgfsetfillcolor{currentfill}%
\pgfsetlinewidth{0.000000pt}%
\definecolor{currentstroke}{rgb}{0.000000,0.000000,0.000000}%
\pgfsetstrokecolor{currentstroke}%
\pgfsetstrokeopacity{0.000000}%
\pgfsetdash{}{0pt}%
\pgfpathmoveto{\pgfqpoint{0.800486in}{1.923543in}}%
\pgfpathlineto{\pgfqpoint{0.805150in}{1.923524in}}%
\pgfpathlineto{\pgfqpoint{0.814693in}{1.932896in}}%
\pgfpathlineto{\pgfqpoint{0.818538in}{1.934210in}}%
\pgfpathlineto{\pgfqpoint{0.839459in}{1.928946in}}%
\pgfpathlineto{\pgfqpoint{0.844130in}{1.923115in}}%
\pgfpathlineto{\pgfqpoint{0.842438in}{1.916478in}}%
\pgfpathlineto{\pgfqpoint{0.857097in}{1.912795in}}%
\pgfpathlineto{\pgfqpoint{0.860164in}{1.926417in}}%
\pgfpathlineto{\pgfqpoint{0.866297in}{1.932367in}}%
\pgfpathlineto{\pgfqpoint{0.875640in}{1.944745in}}%
\pgfpathlineto{\pgfqpoint{0.884981in}{1.950005in}}%
\pgfpathlineto{\pgfqpoint{0.899991in}{1.946190in}}%
\pgfpathlineto{\pgfqpoint{0.896730in}{1.932620in}}%
\pgfpathlineto{\pgfqpoint{0.901098in}{1.927625in}}%
\pgfpathlineto{\pgfqpoint{0.898851in}{1.918021in}}%
\pgfpathlineto{\pgfqpoint{0.904574in}{1.912817in}}%
\pgfpathlineto{\pgfqpoint{0.910516in}{1.911417in}}%
\pgfpathlineto{\pgfqpoint{0.905607in}{1.902658in}}%
\pgfpathlineto{\pgfqpoint{0.901859in}{1.889570in}}%
\pgfpathlineto{\pgfqpoint{0.896690in}{1.890839in}}%
\pgfpathlineto{\pgfqpoint{0.895063in}{1.884202in}}%
\pgfpathlineto{\pgfqpoint{0.890663in}{1.885325in}}%
\pgfpathlineto{\pgfqpoint{0.886843in}{1.879251in}}%
\pgfpathlineto{\pgfqpoint{0.879074in}{1.881212in}}%
\pgfpathlineto{\pgfqpoint{0.875796in}{1.867766in}}%
\pgfpathlineto{\pgfqpoint{0.867712in}{1.868584in}}%
\pgfpathlineto{\pgfqpoint{0.863735in}{1.876688in}}%
\pgfpathlineto{\pgfqpoint{0.860241in}{1.876128in}}%
\pgfpathlineto{\pgfqpoint{0.855951in}{1.862354in}}%
\pgfpathlineto{\pgfqpoint{0.824016in}{1.733109in}}%
\pgfpathlineto{\pgfqpoint{0.757203in}{1.749818in}}%
\pgfpathlineto{\pgfqpoint{0.761786in}{1.771448in}}%
\pgfpathlineto{\pgfqpoint{0.774453in}{1.820402in}}%
\pgfpathlineto{\pgfqpoint{0.773761in}{1.820585in}}%
\pgfpathlineto{\pgfqpoint{0.790924in}{1.886917in}}%
\pgfpathlineto{\pgfqpoint{0.796408in}{1.906884in}}%
\pgfpathlineto{\pgfqpoint{0.800486in}{1.923543in}}%
\pgfpathclose%
\pgfusepath{fill}%
\end{pgfscope}%
\begin{pgfscope}%
\pgfpathrectangle{\pgfqpoint{0.100000in}{0.100000in}}{\pgfqpoint{3.608454in}{2.310000in}}%
\pgfusepath{clip}%
\pgfsetbuttcap%
\pgfsetmiterjoin%
\definecolor{currentfill}{rgb}{0.000000,0.635294,0.682353}%
\pgfsetfillcolor{currentfill}%
\pgfsetlinewidth{0.000000pt}%
\definecolor{currentstroke}{rgb}{0.000000,0.000000,0.000000}%
\pgfsetstrokecolor{currentstroke}%
\pgfsetstrokeopacity{0.000000}%
\pgfsetdash{}{0pt}%
\pgfpathmoveto{\pgfqpoint{2.862925in}{0.815951in}}%
\pgfpathlineto{\pgfqpoint{2.854299in}{0.818491in}}%
\pgfpathlineto{\pgfqpoint{2.853827in}{0.823500in}}%
\pgfpathlineto{\pgfqpoint{2.847520in}{0.837484in}}%
\pgfpathlineto{\pgfqpoint{2.844094in}{0.841666in}}%
\pgfpathlineto{\pgfqpoint{2.847916in}{0.849972in}}%
\pgfpathlineto{\pgfqpoint{2.855107in}{0.854082in}}%
\pgfpathlineto{\pgfqpoint{2.871284in}{0.855909in}}%
\pgfpathlineto{\pgfqpoint{2.872842in}{0.842457in}}%
\pgfpathlineto{\pgfqpoint{2.891100in}{0.843787in}}%
\pgfpathlineto{\pgfqpoint{2.894217in}{0.842012in}}%
\pgfpathlineto{\pgfqpoint{2.895939in}{0.833978in}}%
\pgfpathlineto{\pgfqpoint{2.900223in}{0.831517in}}%
\pgfpathlineto{\pgfqpoint{2.904439in}{0.824809in}}%
\pgfpathlineto{\pgfqpoint{2.884148in}{0.822135in}}%
\pgfpathlineto{\pgfqpoint{2.875215in}{0.821548in}}%
\pgfpathlineto{\pgfqpoint{2.875686in}{0.817596in}}%
\pgfpathlineto{\pgfqpoint{2.862925in}{0.815951in}}%
\pgfpathclose%
\pgfusepath{fill}%
\end{pgfscope}%
\begin{pgfscope}%
\pgfpathrectangle{\pgfqpoint{0.100000in}{0.100000in}}{\pgfqpoint{3.608454in}{2.310000in}}%
\pgfusepath{clip}%
\pgfsetbuttcap%
\pgfsetmiterjoin%
\definecolor{currentfill}{rgb}{0.000000,0.517647,0.741176}%
\pgfsetfillcolor{currentfill}%
\pgfsetlinewidth{0.000000pt}%
\definecolor{currentstroke}{rgb}{0.000000,0.000000,0.000000}%
\pgfsetstrokecolor{currentstroke}%
\pgfsetstrokeopacity{0.000000}%
\pgfsetdash{}{0pt}%
\pgfpathmoveto{\pgfqpoint{1.793434in}{1.427744in}}%
\pgfpathlineto{\pgfqpoint{1.792481in}{1.427797in}}%
\pgfpathlineto{\pgfqpoint{1.793800in}{1.455314in}}%
\pgfpathlineto{\pgfqpoint{1.758946in}{1.457374in}}%
\pgfpathlineto{\pgfqpoint{1.760676in}{1.485050in}}%
\pgfpathlineto{\pgfqpoint{1.793857in}{1.483052in}}%
\pgfpathlineto{\pgfqpoint{1.795374in}{1.510393in}}%
\pgfpathlineto{\pgfqpoint{1.842917in}{1.507856in}}%
\pgfpathlineto{\pgfqpoint{1.841990in}{1.478161in}}%
\pgfpathlineto{\pgfqpoint{1.856208in}{1.476676in}}%
\pgfpathlineto{\pgfqpoint{1.853837in}{1.424663in}}%
\pgfpathlineto{\pgfqpoint{1.826801in}{1.425965in}}%
\pgfpathlineto{\pgfqpoint{1.793434in}{1.427744in}}%
\pgfpathclose%
\pgfusepath{fill}%
\end{pgfscope}%
\begin{pgfscope}%
\pgfpathrectangle{\pgfqpoint{0.100000in}{0.100000in}}{\pgfqpoint{3.608454in}{2.310000in}}%
\pgfusepath{clip}%
\pgfsetbuttcap%
\pgfsetmiterjoin%
\definecolor{currentfill}{rgb}{0.000000,0.341176,0.829412}%
\pgfsetfillcolor{currentfill}%
\pgfsetlinewidth{0.000000pt}%
\definecolor{currentstroke}{rgb}{0.000000,0.000000,0.000000}%
\pgfsetstrokecolor{currentstroke}%
\pgfsetstrokeopacity{0.000000}%
\pgfsetdash{}{0pt}%
\pgfpathmoveto{\pgfqpoint{1.538610in}{1.726627in}}%
\pgfpathlineto{\pgfqpoint{1.586973in}{1.721839in}}%
\pgfpathlineto{\pgfqpoint{1.583107in}{1.682203in}}%
\pgfpathlineto{\pgfqpoint{1.580064in}{1.651542in}}%
\pgfpathlineto{\pgfqpoint{1.531372in}{1.656425in}}%
\pgfpathlineto{\pgfqpoint{1.532653in}{1.671867in}}%
\pgfpathlineto{\pgfqpoint{1.538610in}{1.726627in}}%
\pgfpathclose%
\pgfusepath{fill}%
\end{pgfscope}%
\begin{pgfscope}%
\pgfpathrectangle{\pgfqpoint{0.100000in}{0.100000in}}{\pgfqpoint{3.608454in}{2.310000in}}%
\pgfusepath{clip}%
\pgfsetbuttcap%
\pgfsetmiterjoin%
\definecolor{currentfill}{rgb}{0.000000,0.392157,0.803922}%
\pgfsetfillcolor{currentfill}%
\pgfsetlinewidth{0.000000pt}%
\definecolor{currentstroke}{rgb}{0.000000,0.000000,0.000000}%
\pgfsetstrokecolor{currentstroke}%
\pgfsetstrokeopacity{0.000000}%
\pgfsetdash{}{0pt}%
\pgfpathmoveto{\pgfqpoint{0.429033in}{1.527264in}}%
\pgfpathlineto{\pgfqpoint{0.432587in}{1.529018in}}%
\pgfpathlineto{\pgfqpoint{0.433169in}{1.538861in}}%
\pgfpathlineto{\pgfqpoint{0.431374in}{1.550880in}}%
\pgfpathlineto{\pgfqpoint{0.425226in}{1.565689in}}%
\pgfpathlineto{\pgfqpoint{0.427977in}{1.570534in}}%
\pgfpathlineto{\pgfqpoint{0.438495in}{1.570367in}}%
\pgfpathlineto{\pgfqpoint{0.446440in}{1.577149in}}%
\pgfpathlineto{\pgfqpoint{0.446281in}{1.581110in}}%
\pgfpathlineto{\pgfqpoint{0.452005in}{1.577321in}}%
\pgfpathlineto{\pgfqpoint{0.452551in}{1.565246in}}%
\pgfpathlineto{\pgfqpoint{0.454748in}{1.561199in}}%
\pgfpathlineto{\pgfqpoint{0.455284in}{1.549289in}}%
\pgfpathlineto{\pgfqpoint{0.460110in}{1.545792in}}%
\pgfpathlineto{\pgfqpoint{0.465266in}{1.547913in}}%
\pgfpathlineto{\pgfqpoint{0.479765in}{1.542960in}}%
\pgfpathlineto{\pgfqpoint{0.474910in}{1.526813in}}%
\pgfpathlineto{\pgfqpoint{0.480918in}{1.524945in}}%
\pgfpathlineto{\pgfqpoint{0.477138in}{1.516663in}}%
\pgfpathlineto{\pgfqpoint{0.471900in}{1.514999in}}%
\pgfpathlineto{\pgfqpoint{0.471851in}{1.509726in}}%
\pgfpathlineto{\pgfqpoint{0.476612in}{1.508153in}}%
\pgfpathlineto{\pgfqpoint{0.471692in}{1.491720in}}%
\pgfpathlineto{\pgfqpoint{0.471844in}{1.486991in}}%
\pgfpathlineto{\pgfqpoint{0.465525in}{1.465978in}}%
\pgfpathlineto{\pgfqpoint{0.469229in}{1.459834in}}%
\pgfpathlineto{\pgfqpoint{0.469193in}{1.459876in}}%
\pgfpathlineto{\pgfqpoint{0.445670in}{1.467188in}}%
\pgfpathlineto{\pgfqpoint{0.441374in}{1.466010in}}%
\pgfpathlineto{\pgfqpoint{0.432540in}{1.470642in}}%
\pgfpathlineto{\pgfqpoint{0.424905in}{1.469119in}}%
\pgfpathlineto{\pgfqpoint{0.422616in}{1.464232in}}%
\pgfpathlineto{\pgfqpoint{0.422155in}{1.452127in}}%
\pgfpathlineto{\pgfqpoint{0.411566in}{1.453040in}}%
\pgfpathlineto{\pgfqpoint{0.411162in}{1.446533in}}%
\pgfpathlineto{\pgfqpoint{0.406795in}{1.454089in}}%
\pgfpathlineto{\pgfqpoint{0.406757in}{1.459003in}}%
\pgfpathlineto{\pgfqpoint{0.410394in}{1.466958in}}%
\pgfpathlineto{\pgfqpoint{0.410387in}{1.477682in}}%
\pgfpathlineto{\pgfqpoint{0.407424in}{1.482308in}}%
\pgfpathlineto{\pgfqpoint{0.411743in}{1.490790in}}%
\pgfpathlineto{\pgfqpoint{0.413684in}{1.501555in}}%
\pgfpathlineto{\pgfqpoint{0.420869in}{1.501805in}}%
\pgfpathlineto{\pgfqpoint{0.418472in}{1.484385in}}%
\pgfpathlineto{\pgfqpoint{0.423806in}{1.481158in}}%
\pgfpathlineto{\pgfqpoint{0.428604in}{1.473723in}}%
\pgfpathlineto{\pgfqpoint{0.431932in}{1.475097in}}%
\pgfpathlineto{\pgfqpoint{0.432089in}{1.485171in}}%
\pgfpathlineto{\pgfqpoint{0.424993in}{1.499408in}}%
\pgfpathlineto{\pgfqpoint{0.428477in}{1.506604in}}%
\pgfpathlineto{\pgfqpoint{0.424773in}{1.512850in}}%
\pgfpathlineto{\pgfqpoint{0.435086in}{1.517401in}}%
\pgfpathlineto{\pgfqpoint{0.433615in}{1.522888in}}%
\pgfpathlineto{\pgfqpoint{0.429033in}{1.527264in}}%
\pgfpathclose%
\pgfusepath{fill}%
\end{pgfscope}%
\begin{pgfscope}%
\pgfpathrectangle{\pgfqpoint{0.100000in}{0.100000in}}{\pgfqpoint{3.608454in}{2.310000in}}%
\pgfusepath{clip}%
\pgfsetbuttcap%
\pgfsetmiterjoin%
\definecolor{currentfill}{rgb}{0.000000,0.392157,0.803922}%
\pgfsetfillcolor{currentfill}%
\pgfsetlinewidth{0.000000pt}%
\definecolor{currentstroke}{rgb}{0.000000,0.000000,0.000000}%
\pgfsetstrokecolor{currentstroke}%
\pgfsetstrokeopacity{0.000000}%
\pgfsetdash{}{0pt}%
\pgfpathmoveto{\pgfqpoint{0.422980in}{1.525827in}}%
\pgfpathlineto{\pgfqpoint{0.419309in}{1.509781in}}%
\pgfpathlineto{\pgfqpoint{0.414223in}{1.505038in}}%
\pgfpathlineto{\pgfqpoint{0.408463in}{1.513232in}}%
\pgfpathlineto{\pgfqpoint{0.405158in}{1.513449in}}%
\pgfpathlineto{\pgfqpoint{0.399143in}{1.525491in}}%
\pgfpathlineto{\pgfqpoint{0.393060in}{1.528517in}}%
\pgfpathlineto{\pgfqpoint{0.396639in}{1.537475in}}%
\pgfpathlineto{\pgfqpoint{0.396978in}{1.549508in}}%
\pgfpathlineto{\pgfqpoint{0.403577in}{1.549156in}}%
\pgfpathlineto{\pgfqpoint{0.410415in}{1.537899in}}%
\pgfpathlineto{\pgfqpoint{0.419978in}{1.533247in}}%
\pgfpathlineto{\pgfqpoint{0.422980in}{1.525827in}}%
\pgfpathclose%
\pgfusepath{fill}%
\end{pgfscope}%
\begin{pgfscope}%
\pgfpathrectangle{\pgfqpoint{0.100000in}{0.100000in}}{\pgfqpoint{3.608454in}{2.310000in}}%
\pgfusepath{clip}%
\pgfsetbuttcap%
\pgfsetmiterjoin%
\definecolor{currentfill}{rgb}{0.000000,0.592157,0.703922}%
\pgfsetfillcolor{currentfill}%
\pgfsetlinewidth{0.000000pt}%
\definecolor{currentstroke}{rgb}{0.000000,0.000000,0.000000}%
\pgfsetstrokecolor{currentstroke}%
\pgfsetstrokeopacity{0.000000}%
\pgfsetdash{}{0pt}%
\pgfpathmoveto{\pgfqpoint{2.410650in}{1.552521in}}%
\pgfpathlineto{\pgfqpoint{2.391551in}{1.551544in}}%
\pgfpathlineto{\pgfqpoint{2.387375in}{1.567036in}}%
\pgfpathlineto{\pgfqpoint{2.390834in}{1.569129in}}%
\pgfpathlineto{\pgfqpoint{2.392472in}{1.584764in}}%
\pgfpathlineto{\pgfqpoint{2.390701in}{1.593517in}}%
\pgfpathlineto{\pgfqpoint{2.381293in}{1.599262in}}%
\pgfpathlineto{\pgfqpoint{2.404561in}{1.600729in}}%
\pgfpathlineto{\pgfqpoint{2.402794in}{1.625171in}}%
\pgfpathlineto{\pgfqpoint{2.435480in}{1.626662in}}%
\pgfpathlineto{\pgfqpoint{2.473885in}{1.628654in}}%
\pgfpathlineto{\pgfqpoint{2.475849in}{1.601806in}}%
\pgfpathlineto{\pgfqpoint{2.482705in}{1.602290in}}%
\pgfpathlineto{\pgfqpoint{2.482413in}{1.595355in}}%
\pgfpathlineto{\pgfqpoint{2.484797in}{1.561007in}}%
\pgfpathlineto{\pgfqpoint{2.451488in}{1.558550in}}%
\pgfpathlineto{\pgfqpoint{2.451750in}{1.555104in}}%
\pgfpathlineto{\pgfqpoint{2.410650in}{1.552521in}}%
\pgfpathclose%
\pgfusepath{fill}%
\end{pgfscope}%
\begin{pgfscope}%
\pgfpathrectangle{\pgfqpoint{0.100000in}{0.100000in}}{\pgfqpoint{3.608454in}{2.310000in}}%
\pgfusepath{clip}%
\pgfsetbuttcap%
\pgfsetmiterjoin%
\definecolor{currentfill}{rgb}{0.000000,0.462745,0.768627}%
\pgfsetfillcolor{currentfill}%
\pgfsetlinewidth{0.000000pt}%
\definecolor{currentstroke}{rgb}{0.000000,0.000000,0.000000}%
\pgfsetstrokecolor{currentstroke}%
\pgfsetstrokeopacity{0.000000}%
\pgfsetdash{}{0pt}%
\pgfpathmoveto{\pgfqpoint{2.583434in}{1.352723in}}%
\pgfpathlineto{\pgfqpoint{2.590232in}{1.353266in}}%
\pgfpathlineto{\pgfqpoint{2.590013in}{1.348185in}}%
\pgfpathlineto{\pgfqpoint{2.584141in}{1.345785in}}%
\pgfpathlineto{\pgfqpoint{2.582510in}{1.335565in}}%
\pgfpathlineto{\pgfqpoint{2.585800in}{1.324135in}}%
\pgfpathlineto{\pgfqpoint{2.572272in}{1.322088in}}%
\pgfpathlineto{\pgfqpoint{2.565371in}{1.317560in}}%
\pgfpathlineto{\pgfqpoint{2.560427in}{1.323864in}}%
\pgfpathlineto{\pgfqpoint{2.544363in}{1.322774in}}%
\pgfpathlineto{\pgfqpoint{2.542895in}{1.344839in}}%
\pgfpathlineto{\pgfqpoint{2.565996in}{1.346868in}}%
\pgfpathlineto{\pgfqpoint{2.566176in}{1.350839in}}%
\pgfpathlineto{\pgfqpoint{2.583434in}{1.352723in}}%
\pgfpathclose%
\pgfusepath{fill}%
\end{pgfscope}%
\begin{pgfscope}%
\pgfpathrectangle{\pgfqpoint{0.100000in}{0.100000in}}{\pgfqpoint{3.608454in}{2.310000in}}%
\pgfusepath{clip}%
\pgfsetbuttcap%
\pgfsetmiterjoin%
\definecolor{currentfill}{rgb}{0.000000,0.701961,0.649020}%
\pgfsetfillcolor{currentfill}%
\pgfsetlinewidth{0.000000pt}%
\definecolor{currentstroke}{rgb}{0.000000,0.000000,0.000000}%
\pgfsetstrokecolor{currentstroke}%
\pgfsetstrokeopacity{0.000000}%
\pgfsetdash{}{0pt}%
\pgfpathmoveto{\pgfqpoint{2.062941in}{0.991185in}}%
\pgfpathlineto{\pgfqpoint{2.062908in}{0.984282in}}%
\pgfpathlineto{\pgfqpoint{2.055870in}{0.984320in}}%
\pgfpathlineto{\pgfqpoint{2.055634in}{0.956575in}}%
\pgfpathlineto{\pgfqpoint{2.041706in}{0.956720in}}%
\pgfpathlineto{\pgfqpoint{2.014398in}{0.957120in}}%
\pgfpathlineto{\pgfqpoint{2.014740in}{0.977835in}}%
\pgfpathlineto{\pgfqpoint{2.007778in}{0.978034in}}%
\pgfpathlineto{\pgfqpoint{2.007887in}{0.984833in}}%
\pgfpathlineto{\pgfqpoint{1.987413in}{0.985257in}}%
\pgfpathlineto{\pgfqpoint{1.987552in}{0.992168in}}%
\pgfpathlineto{\pgfqpoint{1.980692in}{0.992304in}}%
\pgfpathlineto{\pgfqpoint{1.981120in}{1.013018in}}%
\pgfpathlineto{\pgfqpoint{1.984031in}{1.018827in}}%
\pgfpathlineto{\pgfqpoint{1.988856in}{1.020093in}}%
\pgfpathlineto{\pgfqpoint{1.994760in}{1.012926in}}%
\pgfpathlineto{\pgfqpoint{1.999760in}{1.016432in}}%
\pgfpathlineto{\pgfqpoint{2.009922in}{1.016922in}}%
\pgfpathlineto{\pgfqpoint{2.010212in}{1.033190in}}%
\pgfpathlineto{\pgfqpoint{2.013392in}{1.033129in}}%
\pgfpathlineto{\pgfqpoint{2.013608in}{1.046935in}}%
\pgfpathlineto{\pgfqpoint{2.043347in}{1.046522in}}%
\pgfpathlineto{\pgfqpoint{2.043077in}{1.035565in}}%
\pgfpathlineto{\pgfqpoint{2.036171in}{1.027587in}}%
\pgfpathlineto{\pgfqpoint{2.035603in}{0.998326in}}%
\pgfpathlineto{\pgfqpoint{2.049296in}{0.998169in}}%
\pgfpathlineto{\pgfqpoint{2.049224in}{0.991276in}}%
\pgfpathlineto{\pgfqpoint{2.062941in}{0.991185in}}%
\pgfpathclose%
\pgfusepath{fill}%
\end{pgfscope}%
\begin{pgfscope}%
\pgfpathrectangle{\pgfqpoint{0.100000in}{0.100000in}}{\pgfqpoint{3.608454in}{2.310000in}}%
\pgfusepath{clip}%
\pgfsetbuttcap%
\pgfsetmiterjoin%
\definecolor{currentfill}{rgb}{0.000000,0.580392,0.709804}%
\pgfsetfillcolor{currentfill}%
\pgfsetlinewidth{0.000000pt}%
\definecolor{currentstroke}{rgb}{0.000000,0.000000,0.000000}%
\pgfsetstrokecolor{currentstroke}%
\pgfsetstrokeopacity{0.000000}%
\pgfsetdash{}{0pt}%
\pgfpathmoveto{\pgfqpoint{3.152496in}{1.240772in}}%
\pgfpathlineto{\pgfqpoint{3.187275in}{1.247460in}}%
\pgfpathlineto{\pgfqpoint{3.188912in}{1.234693in}}%
\pgfpathlineto{\pgfqpoint{3.185638in}{1.219606in}}%
\pgfpathlineto{\pgfqpoint{3.177079in}{1.221508in}}%
\pgfpathlineto{\pgfqpoint{3.165876in}{1.221007in}}%
\pgfpathlineto{\pgfqpoint{3.160658in}{1.212280in}}%
\pgfpathlineto{\pgfqpoint{3.154402in}{1.212187in}}%
\pgfpathlineto{\pgfqpoint{3.150361in}{1.233286in}}%
\pgfpathlineto{\pgfqpoint{3.152496in}{1.240772in}}%
\pgfpathclose%
\pgfusepath{fill}%
\end{pgfscope}%
\begin{pgfscope}%
\pgfpathrectangle{\pgfqpoint{0.100000in}{0.100000in}}{\pgfqpoint{3.608454in}{2.310000in}}%
\pgfusepath{clip}%
\pgfsetbuttcap%
\pgfsetmiterjoin%
\definecolor{currentfill}{rgb}{0.000000,0.792157,0.603922}%
\pgfsetfillcolor{currentfill}%
\pgfsetlinewidth{0.000000pt}%
\definecolor{currentstroke}{rgb}{0.000000,0.000000,0.000000}%
\pgfsetstrokecolor{currentstroke}%
\pgfsetstrokeopacity{0.000000}%
\pgfsetdash{}{0pt}%
\pgfpathmoveto{\pgfqpoint{1.183780in}{0.997802in}}%
\pgfpathlineto{\pgfqpoint{1.176800in}{0.953138in}}%
\pgfpathlineto{\pgfqpoint{1.167197in}{0.891666in}}%
\pgfpathlineto{\pgfqpoint{1.092774in}{0.903733in}}%
\pgfpathlineto{\pgfqpoint{1.074172in}{0.906976in}}%
\pgfpathlineto{\pgfqpoint{1.075355in}{0.913752in}}%
\pgfpathlineto{\pgfqpoint{1.084657in}{0.967062in}}%
\pgfpathlineto{\pgfqpoint{1.087666in}{0.981611in}}%
\pgfpathlineto{\pgfqpoint{1.096233in}{0.984306in}}%
\pgfpathlineto{\pgfqpoint{1.106555in}{0.985188in}}%
\pgfpathlineto{\pgfqpoint{1.117590in}{0.983356in}}%
\pgfpathlineto{\pgfqpoint{1.119039in}{0.992055in}}%
\pgfpathlineto{\pgfqpoint{1.125550in}{0.990854in}}%
\pgfpathlineto{\pgfqpoint{1.133001in}{0.981587in}}%
\pgfpathlineto{\pgfqpoint{1.145330in}{0.990568in}}%
\pgfpathlineto{\pgfqpoint{1.150108in}{0.990389in}}%
\pgfpathlineto{\pgfqpoint{1.154422in}{0.995185in}}%
\pgfpathlineto{\pgfqpoint{1.164171in}{1.000833in}}%
\pgfpathlineto{\pgfqpoint{1.183780in}{0.997802in}}%
\pgfpathclose%
\pgfusepath{fill}%
\end{pgfscope}%
\begin{pgfscope}%
\pgfpathrectangle{\pgfqpoint{0.100000in}{0.100000in}}{\pgfqpoint{3.608454in}{2.310000in}}%
\pgfusepath{clip}%
\pgfsetbuttcap%
\pgfsetmiterjoin%
\definecolor{currentfill}{rgb}{0.000000,0.760784,0.619608}%
\pgfsetfillcolor{currentfill}%
\pgfsetlinewidth{0.000000pt}%
\definecolor{currentstroke}{rgb}{0.000000,0.000000,0.000000}%
\pgfsetstrokecolor{currentstroke}%
\pgfsetstrokeopacity{0.000000}%
\pgfsetdash{}{0pt}%
\pgfpathmoveto{\pgfqpoint{2.056507in}{0.745639in}}%
\pgfpathlineto{\pgfqpoint{2.061056in}{0.729018in}}%
\pgfpathlineto{\pgfqpoint{2.057335in}{0.727032in}}%
\pgfpathlineto{\pgfqpoint{2.053852in}{0.716384in}}%
\pgfpathlineto{\pgfqpoint{2.054567in}{0.712981in}}%
\pgfpathlineto{\pgfqpoint{2.040106in}{0.712927in}}%
\pgfpathlineto{\pgfqpoint{2.021875in}{0.703710in}}%
\pgfpathlineto{\pgfqpoint{2.016298in}{0.728776in}}%
\pgfpathlineto{\pgfqpoint{2.016972in}{0.734385in}}%
\pgfpathlineto{\pgfqpoint{1.997875in}{0.723825in}}%
\pgfpathlineto{\pgfqpoint{1.975832in}{0.763164in}}%
\pgfpathlineto{\pgfqpoint{1.990434in}{0.771318in}}%
\pgfpathlineto{\pgfqpoint{1.978879in}{0.792186in}}%
\pgfpathlineto{\pgfqpoint{2.013777in}{0.811827in}}%
\pgfpathlineto{\pgfqpoint{2.023652in}{0.805506in}}%
\pgfpathlineto{\pgfqpoint{2.034850in}{0.788420in}}%
\pgfpathlineto{\pgfqpoint{2.042310in}{0.767446in}}%
\pgfpathlineto{\pgfqpoint{2.047446in}{0.765794in}}%
\pgfpathlineto{\pgfqpoint{2.048151in}{0.760702in}}%
\pgfpathlineto{\pgfqpoint{2.058523in}{0.753655in}}%
\pgfpathlineto{\pgfqpoint{2.056507in}{0.745639in}}%
\pgfpathclose%
\pgfusepath{fill}%
\end{pgfscope}%
\begin{pgfscope}%
\pgfpathrectangle{\pgfqpoint{0.100000in}{0.100000in}}{\pgfqpoint{3.608454in}{2.310000in}}%
\pgfusepath{clip}%
\pgfsetbuttcap%
\pgfsetmiterjoin%
\definecolor{currentfill}{rgb}{0.000000,0.466667,0.766667}%
\pgfsetfillcolor{currentfill}%
\pgfsetlinewidth{0.000000pt}%
\definecolor{currentstroke}{rgb}{0.000000,0.000000,0.000000}%
\pgfsetstrokecolor{currentstroke}%
\pgfsetstrokeopacity{0.000000}%
\pgfsetdash{}{0pt}%
\pgfpathmoveto{\pgfqpoint{1.951799in}{2.072044in}}%
\pgfpathlineto{\pgfqpoint{1.952493in}{2.057946in}}%
\pgfpathlineto{\pgfqpoint{1.951584in}{2.030204in}}%
\pgfpathlineto{\pgfqpoint{1.946085in}{2.030374in}}%
\pgfpathlineto{\pgfqpoint{1.918355in}{2.031373in}}%
\pgfpathlineto{\pgfqpoint{1.917529in}{2.045335in}}%
\pgfpathlineto{\pgfqpoint{1.876114in}{2.047161in}}%
\pgfpathlineto{\pgfqpoint{1.848523in}{2.048601in}}%
\pgfpathlineto{\pgfqpoint{1.847326in}{2.062601in}}%
\pgfpathlineto{\pgfqpoint{1.848899in}{2.090580in}}%
\pgfpathlineto{\pgfqpoint{1.867703in}{2.089585in}}%
\pgfpathlineto{\pgfqpoint{1.895459in}{2.088252in}}%
\pgfpathlineto{\pgfqpoint{1.896095in}{2.102164in}}%
\pgfpathlineto{\pgfqpoint{1.930695in}{2.100743in}}%
\pgfpathlineto{\pgfqpoint{1.931568in}{2.086739in}}%
\pgfpathlineto{\pgfqpoint{1.931043in}{2.072760in}}%
\pgfpathlineto{\pgfqpoint{1.951799in}{2.072044in}}%
\pgfpathclose%
\pgfusepath{fill}%
\end{pgfscope}%
\begin{pgfscope}%
\pgfpathrectangle{\pgfqpoint{0.100000in}{0.100000in}}{\pgfqpoint{3.608454in}{2.310000in}}%
\pgfusepath{clip}%
\pgfsetbuttcap%
\pgfsetmiterjoin%
\definecolor{currentfill}{rgb}{0.000000,0.654902,0.672549}%
\pgfsetfillcolor{currentfill}%
\pgfsetlinewidth{0.000000pt}%
\definecolor{currentstroke}{rgb}{0.000000,0.000000,0.000000}%
\pgfsetstrokecolor{currentstroke}%
\pgfsetstrokeopacity{0.000000}%
\pgfsetdash{}{0pt}%
\pgfpathmoveto{\pgfqpoint{2.452854in}{0.589739in}}%
\pgfpathlineto{\pgfqpoint{2.452510in}{0.580153in}}%
\pgfpathlineto{\pgfqpoint{2.447694in}{0.578303in}}%
\pgfpathlineto{\pgfqpoint{2.452088in}{0.571325in}}%
\pgfpathlineto{\pgfqpoint{2.450488in}{0.567507in}}%
\pgfpathlineto{\pgfqpoint{2.441407in}{0.560816in}}%
\pgfpathlineto{\pgfqpoint{2.437817in}{0.565645in}}%
\pgfpathlineto{\pgfqpoint{2.435406in}{0.574170in}}%
\pgfpathlineto{\pgfqpoint{2.431293in}{0.576903in}}%
\pgfpathlineto{\pgfqpoint{2.428352in}{0.570405in}}%
\pgfpathlineto{\pgfqpoint{2.421368in}{0.576035in}}%
\pgfpathlineto{\pgfqpoint{2.411945in}{0.568422in}}%
\pgfpathlineto{\pgfqpoint{2.411612in}{0.562816in}}%
\pgfpathlineto{\pgfqpoint{2.405732in}{0.562002in}}%
\pgfpathlineto{\pgfqpoint{2.397767in}{0.557468in}}%
\pgfpathlineto{\pgfqpoint{2.389463in}{0.565906in}}%
\pgfpathlineto{\pgfqpoint{2.380049in}{0.566008in}}%
\pgfpathlineto{\pgfqpoint{2.367175in}{0.570371in}}%
\pgfpathlineto{\pgfqpoint{2.363150in}{0.573969in}}%
\pgfpathlineto{\pgfqpoint{2.369537in}{0.580614in}}%
\pgfpathlineto{\pgfqpoint{2.366976in}{0.588543in}}%
\pgfpathlineto{\pgfqpoint{2.347924in}{0.595251in}}%
\pgfpathlineto{\pgfqpoint{2.341614in}{0.599740in}}%
\pgfpathlineto{\pgfqpoint{2.341637in}{0.608452in}}%
\pgfpathlineto{\pgfqpoint{2.330757in}{0.608626in}}%
\pgfpathlineto{\pgfqpoint{2.323648in}{0.606472in}}%
\pgfpathlineto{\pgfqpoint{2.335052in}{0.617153in}}%
\pgfpathlineto{\pgfqpoint{2.340753in}{0.625307in}}%
\pgfpathlineto{\pgfqpoint{2.351303in}{0.626566in}}%
\pgfpathlineto{\pgfqpoint{2.358501in}{0.617112in}}%
\pgfpathlineto{\pgfqpoint{2.357771in}{0.613019in}}%
\pgfpathlineto{\pgfqpoint{2.372197in}{0.609200in}}%
\pgfpathlineto{\pgfqpoint{2.378486in}{0.611623in}}%
\pgfpathlineto{\pgfqpoint{2.376744in}{0.617237in}}%
\pgfpathlineto{\pgfqpoint{2.371581in}{0.618635in}}%
\pgfpathlineto{\pgfqpoint{2.365882in}{0.630690in}}%
\pgfpathlineto{\pgfqpoint{2.368222in}{0.632799in}}%
\pgfpathlineto{\pgfqpoint{2.383161in}{0.637744in}}%
\pgfpathlineto{\pgfqpoint{2.386115in}{0.636858in}}%
\pgfpathlineto{\pgfqpoint{2.386695in}{0.629800in}}%
\pgfpathlineto{\pgfqpoint{2.392099in}{0.624223in}}%
\pgfpathlineto{\pgfqpoint{2.416070in}{0.624294in}}%
\pgfpathlineto{\pgfqpoint{2.420890in}{0.618181in}}%
\pgfpathlineto{\pgfqpoint{2.428180in}{0.614445in}}%
\pgfpathlineto{\pgfqpoint{2.430825in}{0.609081in}}%
\pgfpathlineto{\pgfqpoint{2.440322in}{0.609463in}}%
\pgfpathlineto{\pgfqpoint{2.443757in}{0.607123in}}%
\pgfpathlineto{\pgfqpoint{2.442006in}{0.599263in}}%
\pgfpathlineto{\pgfqpoint{2.445754in}{0.592557in}}%
\pgfpathlineto{\pgfqpoint{2.452854in}{0.589739in}}%
\pgfpathclose%
\pgfusepath{fill}%
\end{pgfscope}%
\begin{pgfscope}%
\pgfpathrectangle{\pgfqpoint{0.100000in}{0.100000in}}{\pgfqpoint{3.608454in}{2.310000in}}%
\pgfusepath{clip}%
\pgfsetbuttcap%
\pgfsetmiterjoin%
\definecolor{currentfill}{rgb}{0.000000,0.568627,0.715686}%
\pgfsetfillcolor{currentfill}%
\pgfsetlinewidth{0.000000pt}%
\definecolor{currentstroke}{rgb}{0.000000,0.000000,0.000000}%
\pgfsetstrokecolor{currentstroke}%
\pgfsetstrokeopacity{0.000000}%
\pgfsetdash{}{0pt}%
\pgfpathmoveto{\pgfqpoint{1.782538in}{1.034546in}}%
\pgfpathlineto{\pgfqpoint{1.785102in}{1.081107in}}%
\pgfpathlineto{\pgfqpoint{1.786242in}{1.101806in}}%
\pgfpathlineto{\pgfqpoint{1.792285in}{1.107896in}}%
\pgfpathlineto{\pgfqpoint{1.798623in}{1.104938in}}%
\pgfpathlineto{\pgfqpoint{1.804299in}{1.099161in}}%
\pgfpathlineto{\pgfqpoint{1.814486in}{1.099133in}}%
\pgfpathlineto{\pgfqpoint{1.818722in}{1.107444in}}%
\pgfpathlineto{\pgfqpoint{1.826418in}{1.110191in}}%
\pgfpathlineto{\pgfqpoint{1.825972in}{1.094304in}}%
\pgfpathlineto{\pgfqpoint{1.873684in}{1.092215in}}%
\pgfpathlineto{\pgfqpoint{1.873435in}{1.071561in}}%
\pgfpathlineto{\pgfqpoint{1.872137in}{1.035604in}}%
\pgfpathlineto{\pgfqpoint{1.863914in}{1.037452in}}%
\pgfpathlineto{\pgfqpoint{1.796992in}{1.040622in}}%
\pgfpathlineto{\pgfqpoint{1.796623in}{1.033765in}}%
\pgfpathlineto{\pgfqpoint{1.782538in}{1.034546in}}%
\pgfpathclose%
\pgfusepath{fill}%
\end{pgfscope}%
\begin{pgfscope}%
\pgfpathrectangle{\pgfqpoint{0.100000in}{0.100000in}}{\pgfqpoint{3.608454in}{2.310000in}}%
\pgfusepath{clip}%
\pgfsetbuttcap%
\pgfsetmiterjoin%
\definecolor{currentfill}{rgb}{0.000000,0.682353,0.658824}%
\pgfsetfillcolor{currentfill}%
\pgfsetlinewidth{0.000000pt}%
\definecolor{currentstroke}{rgb}{0.000000,0.000000,0.000000}%
\pgfsetstrokecolor{currentstroke}%
\pgfsetstrokeopacity{0.000000}%
\pgfsetdash{}{0pt}%
\pgfpathmoveto{\pgfqpoint{3.008758in}{1.422029in}}%
\pgfpathlineto{\pgfqpoint{3.021096in}{1.413057in}}%
\pgfpathlineto{\pgfqpoint{3.018667in}{1.406375in}}%
\pgfpathlineto{\pgfqpoint{3.019598in}{1.394420in}}%
\pgfpathlineto{\pgfqpoint{3.012891in}{1.393149in}}%
\pgfpathlineto{\pgfqpoint{3.010884in}{1.388904in}}%
\pgfpathlineto{\pgfqpoint{3.001117in}{1.390072in}}%
\pgfpathlineto{\pgfqpoint{2.997346in}{1.397787in}}%
\pgfpathlineto{\pgfqpoint{2.989078in}{1.401477in}}%
\pgfpathlineto{\pgfqpoint{2.980404in}{1.409102in}}%
\pgfpathlineto{\pgfqpoint{2.978628in}{1.415804in}}%
\pgfpathlineto{\pgfqpoint{2.986872in}{1.422883in}}%
\pgfpathlineto{\pgfqpoint{3.004458in}{1.420613in}}%
\pgfpathlineto{\pgfqpoint{3.008758in}{1.422029in}}%
\pgfpathclose%
\pgfusepath{fill}%
\end{pgfscope}%
\begin{pgfscope}%
\pgfpathrectangle{\pgfqpoint{0.100000in}{0.100000in}}{\pgfqpoint{3.608454in}{2.310000in}}%
\pgfusepath{clip}%
\pgfsetbuttcap%
\pgfsetmiterjoin%
\definecolor{currentfill}{rgb}{0.000000,0.478431,0.760784}%
\pgfsetfillcolor{currentfill}%
\pgfsetlinewidth{0.000000pt}%
\definecolor{currentstroke}{rgb}{0.000000,0.000000,0.000000}%
\pgfsetstrokecolor{currentstroke}%
\pgfsetstrokeopacity{0.000000}%
\pgfsetdash{}{0pt}%
\pgfpathmoveto{\pgfqpoint{2.814875in}{1.291146in}}%
\pgfpathlineto{\pgfqpoint{2.810307in}{1.281879in}}%
\pgfpathlineto{\pgfqpoint{2.801059in}{1.283583in}}%
\pgfpathlineto{\pgfqpoint{2.791570in}{1.288910in}}%
\pgfpathlineto{\pgfqpoint{2.774318in}{1.293523in}}%
\pgfpathlineto{\pgfqpoint{2.759930in}{1.280063in}}%
\pgfpathlineto{\pgfqpoint{2.757404in}{1.276520in}}%
\pgfpathlineto{\pgfqpoint{2.740147in}{1.292411in}}%
\pgfpathlineto{\pgfqpoint{2.733018in}{1.291688in}}%
\pgfpathlineto{\pgfqpoint{2.727518in}{1.286122in}}%
\pgfpathlineto{\pgfqpoint{2.719889in}{1.285770in}}%
\pgfpathlineto{\pgfqpoint{2.718323in}{1.291481in}}%
\pgfpathlineto{\pgfqpoint{2.724571in}{1.300390in}}%
\pgfpathlineto{\pgfqpoint{2.726400in}{1.307715in}}%
\pgfpathlineto{\pgfqpoint{2.726160in}{1.316456in}}%
\pgfpathlineto{\pgfqpoint{2.730992in}{1.324932in}}%
\pgfpathlineto{\pgfqpoint{2.724859in}{1.333241in}}%
\pgfpathlineto{\pgfqpoint{2.722897in}{1.338929in}}%
\pgfpathlineto{\pgfqpoint{2.718732in}{1.341235in}}%
\pgfpathlineto{\pgfqpoint{2.726746in}{1.347306in}}%
\pgfpathlineto{\pgfqpoint{2.735344in}{1.352880in}}%
\pgfpathlineto{\pgfqpoint{2.736943in}{1.344983in}}%
\pgfpathlineto{\pgfqpoint{2.750154in}{1.334910in}}%
\pgfpathlineto{\pgfqpoint{2.760419in}{1.343088in}}%
\pgfpathlineto{\pgfqpoint{2.772188in}{1.346181in}}%
\pgfpathlineto{\pgfqpoint{2.779558in}{1.348364in}}%
\pgfpathlineto{\pgfqpoint{2.786720in}{1.347970in}}%
\pgfpathlineto{\pgfqpoint{2.789928in}{1.341177in}}%
\pgfpathlineto{\pgfqpoint{2.787097in}{1.337635in}}%
\pgfpathlineto{\pgfqpoint{2.796848in}{1.326453in}}%
\pgfpathlineto{\pgfqpoint{2.784277in}{1.310449in}}%
\pgfpathlineto{\pgfqpoint{2.793467in}{1.297141in}}%
\pgfpathlineto{\pgfqpoint{2.806402in}{1.296003in}}%
\pgfpathlineto{\pgfqpoint{2.814875in}{1.291146in}}%
\pgfpathclose%
\pgfusepath{fill}%
\end{pgfscope}%
\begin{pgfscope}%
\pgfpathrectangle{\pgfqpoint{0.100000in}{0.100000in}}{\pgfqpoint{3.608454in}{2.310000in}}%
\pgfusepath{clip}%
\pgfsetbuttcap%
\pgfsetmiterjoin%
\definecolor{currentfill}{rgb}{0.000000,0.619608,0.690196}%
\pgfsetfillcolor{currentfill}%
\pgfsetlinewidth{0.000000pt}%
\definecolor{currentstroke}{rgb}{0.000000,0.000000,0.000000}%
\pgfsetstrokecolor{currentstroke}%
\pgfsetstrokeopacity{0.000000}%
\pgfsetdash{}{0pt}%
\pgfpathmoveto{\pgfqpoint{2.657467in}{1.463709in}}%
\pgfpathlineto{\pgfqpoint{2.634617in}{1.461203in}}%
\pgfpathlineto{\pgfqpoint{2.633399in}{1.473636in}}%
\pgfpathlineto{\pgfqpoint{2.625004in}{1.477436in}}%
\pgfpathlineto{\pgfqpoint{2.623185in}{1.497972in}}%
\pgfpathlineto{\pgfqpoint{2.614139in}{1.497101in}}%
\pgfpathlineto{\pgfqpoint{2.610441in}{1.500202in}}%
\pgfpathlineto{\pgfqpoint{2.609172in}{1.513996in}}%
\pgfpathlineto{\pgfqpoint{2.588380in}{1.512229in}}%
\pgfpathlineto{\pgfqpoint{2.586580in}{1.532701in}}%
\pgfpathlineto{\pgfqpoint{2.614051in}{1.535158in}}%
\pgfpathlineto{\pgfqpoint{2.615905in}{1.514539in}}%
\pgfpathlineto{\pgfqpoint{2.633720in}{1.516230in}}%
\pgfpathlineto{\pgfqpoint{2.633101in}{1.523054in}}%
\pgfpathlineto{\pgfqpoint{2.646259in}{1.524584in}}%
\pgfpathlineto{\pgfqpoint{2.649395in}{1.497238in}}%
\pgfpathlineto{\pgfqpoint{2.653860in}{1.497742in}}%
\pgfpathlineto{\pgfqpoint{2.657467in}{1.463709in}}%
\pgfpathclose%
\pgfusepath{fill}%
\end{pgfscope}%
\begin{pgfscope}%
\pgfpathrectangle{\pgfqpoint{0.100000in}{0.100000in}}{\pgfqpoint{3.608454in}{2.310000in}}%
\pgfusepath{clip}%
\pgfsetbuttcap%
\pgfsetmiterjoin%
\definecolor{currentfill}{rgb}{0.000000,0.713725,0.643137}%
\pgfsetfillcolor{currentfill}%
\pgfsetlinewidth{0.000000pt}%
\definecolor{currentstroke}{rgb}{0.000000,0.000000,0.000000}%
\pgfsetstrokecolor{currentstroke}%
\pgfsetstrokeopacity{0.000000}%
\pgfsetdash{}{0pt}%
\pgfpathmoveto{\pgfqpoint{2.356606in}{1.028702in}}%
\pgfpathlineto{\pgfqpoint{2.347372in}{1.026940in}}%
\pgfpathlineto{\pgfqpoint{2.342574in}{1.032044in}}%
\pgfpathlineto{\pgfqpoint{2.335892in}{1.034757in}}%
\pgfpathlineto{\pgfqpoint{2.328181in}{1.034618in}}%
\pgfpathlineto{\pgfqpoint{2.328311in}{1.029359in}}%
\pgfpathlineto{\pgfqpoint{2.320051in}{1.032034in}}%
\pgfpathlineto{\pgfqpoint{2.314200in}{1.029274in}}%
\pgfpathlineto{\pgfqpoint{2.309489in}{1.028555in}}%
\pgfpathlineto{\pgfqpoint{2.300760in}{1.031780in}}%
\pgfpathlineto{\pgfqpoint{2.293540in}{1.031438in}}%
\pgfpathlineto{\pgfqpoint{2.293330in}{1.054914in}}%
\pgfpathlineto{\pgfqpoint{2.310094in}{1.055119in}}%
\pgfpathlineto{\pgfqpoint{2.313444in}{1.062948in}}%
\pgfpathlineto{\pgfqpoint{2.313347in}{1.069167in}}%
\pgfpathlineto{\pgfqpoint{2.340700in}{1.069582in}}%
\pgfpathlineto{\pgfqpoint{2.340391in}{1.078576in}}%
\pgfpathlineto{\pgfqpoint{2.343663in}{1.088662in}}%
\pgfpathlineto{\pgfqpoint{2.350507in}{1.098801in}}%
\pgfpathlineto{\pgfqpoint{2.361164in}{1.098864in}}%
\pgfpathlineto{\pgfqpoint{2.362249in}{1.063873in}}%
\pgfpathlineto{\pgfqpoint{2.363393in}{1.028869in}}%
\pgfpathlineto{\pgfqpoint{2.356606in}{1.028702in}}%
\pgfpathclose%
\pgfusepath{fill}%
\end{pgfscope}%
\begin{pgfscope}%
\pgfpathrectangle{\pgfqpoint{0.100000in}{0.100000in}}{\pgfqpoint{3.608454in}{2.310000in}}%
\pgfusepath{clip}%
\pgfsetbuttcap%
\pgfsetmiterjoin%
\definecolor{currentfill}{rgb}{0.000000,0.337255,0.831373}%
\pgfsetfillcolor{currentfill}%
\pgfsetlinewidth{0.000000pt}%
\definecolor{currentstroke}{rgb}{0.000000,0.000000,0.000000}%
\pgfsetstrokecolor{currentstroke}%
\pgfsetstrokeopacity{0.000000}%
\pgfsetdash{}{0pt}%
\pgfpathmoveto{\pgfqpoint{1.844714in}{0.901800in}}%
\pgfpathlineto{\pgfqpoint{1.846276in}{0.936427in}}%
\pgfpathlineto{\pgfqpoint{1.847621in}{0.966431in}}%
\pgfpathlineto{\pgfqpoint{1.852516in}{0.961694in}}%
\pgfpathlineto{\pgfqpoint{1.864563in}{0.959406in}}%
\pgfpathlineto{\pgfqpoint{1.869882in}{0.961093in}}%
\pgfpathlineto{\pgfqpoint{1.877716in}{0.953284in}}%
\pgfpathlineto{\pgfqpoint{1.883512in}{0.956032in}}%
\pgfpathlineto{\pgfqpoint{1.885921in}{0.960469in}}%
\pgfpathlineto{\pgfqpoint{1.898581in}{0.956614in}}%
\pgfpathlineto{\pgfqpoint{1.906317in}{0.946847in}}%
\pgfpathlineto{\pgfqpoint{1.913020in}{0.946294in}}%
\pgfpathlineto{\pgfqpoint{1.910664in}{0.938402in}}%
\pgfpathlineto{\pgfqpoint{1.909495in}{0.904858in}}%
\pgfpathlineto{\pgfqpoint{1.880200in}{0.905864in}}%
\pgfpathlineto{\pgfqpoint{1.880039in}{0.900259in}}%
\pgfpathlineto{\pgfqpoint{1.844714in}{0.901800in}}%
\pgfpathclose%
\pgfusepath{fill}%
\end{pgfscope}%
\begin{pgfscope}%
\pgfpathrectangle{\pgfqpoint{0.100000in}{0.100000in}}{\pgfqpoint{3.608454in}{2.310000in}}%
\pgfusepath{clip}%
\pgfsetbuttcap%
\pgfsetmiterjoin%
\definecolor{currentfill}{rgb}{0.000000,0.627451,0.686275}%
\pgfsetfillcolor{currentfill}%
\pgfsetlinewidth{0.000000pt}%
\definecolor{currentstroke}{rgb}{0.000000,0.000000,0.000000}%
\pgfsetstrokecolor{currentstroke}%
\pgfsetstrokeopacity{0.000000}%
\pgfsetdash{}{0pt}%
\pgfpathmoveto{\pgfqpoint{2.227204in}{1.266822in}}%
\pgfpathlineto{\pgfqpoint{2.199885in}{1.267142in}}%
\pgfpathlineto{\pgfqpoint{2.199587in}{1.254527in}}%
\pgfpathlineto{\pgfqpoint{2.196058in}{1.254577in}}%
\pgfpathlineto{\pgfqpoint{2.192435in}{1.247763in}}%
\pgfpathlineto{\pgfqpoint{2.180991in}{1.248016in}}%
\pgfpathlineto{\pgfqpoint{2.181129in}{1.252610in}}%
\pgfpathlineto{\pgfqpoint{2.165075in}{1.253073in}}%
\pgfpathlineto{\pgfqpoint{2.165424in}{1.263791in}}%
\pgfpathlineto{\pgfqpoint{2.164228in}{1.277853in}}%
\pgfpathlineto{\pgfqpoint{2.164617in}{1.305750in}}%
\pgfpathlineto{\pgfqpoint{2.198808in}{1.305319in}}%
\pgfpathlineto{\pgfqpoint{2.199219in}{1.335898in}}%
\pgfpathlineto{\pgfqpoint{2.226654in}{1.335261in}}%
\pgfpathlineto{\pgfqpoint{2.226458in}{1.316665in}}%
\pgfpathlineto{\pgfqpoint{2.225942in}{1.276603in}}%
\pgfpathlineto{\pgfqpoint{2.227204in}{1.266822in}}%
\pgfpathclose%
\pgfusepath{fill}%
\end{pgfscope}%
\begin{pgfscope}%
\pgfpathrectangle{\pgfqpoint{0.100000in}{0.100000in}}{\pgfqpoint{3.608454in}{2.310000in}}%
\pgfusepath{clip}%
\pgfsetbuttcap%
\pgfsetmiterjoin%
\definecolor{currentfill}{rgb}{0.000000,0.674510,0.662745}%
\pgfsetfillcolor{currentfill}%
\pgfsetlinewidth{0.000000pt}%
\definecolor{currentstroke}{rgb}{0.000000,0.000000,0.000000}%
\pgfsetstrokecolor{currentstroke}%
\pgfsetstrokeopacity{0.000000}%
\pgfsetdash{}{0pt}%
\pgfpathmoveto{\pgfqpoint{0.499716in}{1.160624in}}%
\pgfpathlineto{\pgfqpoint{0.490524in}{1.164969in}}%
\pgfpathlineto{\pgfqpoint{0.493536in}{1.171597in}}%
\pgfpathlineto{\pgfqpoint{0.503853in}{1.164499in}}%
\pgfpathlineto{\pgfqpoint{0.499716in}{1.160624in}}%
\pgfpathclose%
\pgfusepath{fill}%
\end{pgfscope}%
\begin{pgfscope}%
\pgfpathrectangle{\pgfqpoint{0.100000in}{0.100000in}}{\pgfqpoint{3.608454in}{2.310000in}}%
\pgfusepath{clip}%
\pgfsetbuttcap%
\pgfsetmiterjoin%
\definecolor{currentfill}{rgb}{0.000000,0.674510,0.662745}%
\pgfsetfillcolor{currentfill}%
\pgfsetlinewidth{0.000000pt}%
\definecolor{currentstroke}{rgb}{0.000000,0.000000,0.000000}%
\pgfsetstrokecolor{currentstroke}%
\pgfsetstrokeopacity{0.000000}%
\pgfsetdash{}{0pt}%
\pgfpathmoveto{\pgfqpoint{0.473802in}{1.162641in}}%
\pgfpathlineto{\pgfqpoint{0.468569in}{1.171802in}}%
\pgfpathlineto{\pgfqpoint{0.474440in}{1.173296in}}%
\pgfpathlineto{\pgfqpoint{0.483849in}{1.167231in}}%
\pgfpathlineto{\pgfqpoint{0.482039in}{1.164194in}}%
\pgfpathlineto{\pgfqpoint{0.473802in}{1.162641in}}%
\pgfpathclose%
\pgfusepath{fill}%
\end{pgfscope}%
\begin{pgfscope}%
\pgfpathrectangle{\pgfqpoint{0.100000in}{0.100000in}}{\pgfqpoint{3.608454in}{2.310000in}}%
\pgfusepath{clip}%
\pgfsetbuttcap%
\pgfsetmiterjoin%
\definecolor{currentfill}{rgb}{0.000000,0.674510,0.662745}%
\pgfsetfillcolor{currentfill}%
\pgfsetlinewidth{0.000000pt}%
\definecolor{currentstroke}{rgb}{0.000000,0.000000,0.000000}%
\pgfsetstrokecolor{currentstroke}%
\pgfsetstrokeopacity{0.000000}%
\pgfsetdash{}{0pt}%
\pgfpathmoveto{\pgfqpoint{0.537333in}{1.227732in}}%
\pgfpathlineto{\pgfqpoint{0.527921in}{1.193369in}}%
\pgfpathlineto{\pgfqpoint{0.524045in}{1.187970in}}%
\pgfpathlineto{\pgfqpoint{0.516406in}{1.193887in}}%
\pgfpathlineto{\pgfqpoint{0.508386in}{1.194136in}}%
\pgfpathlineto{\pgfqpoint{0.499604in}{1.197528in}}%
\pgfpathlineto{\pgfqpoint{0.493082in}{1.203453in}}%
\pgfpathlineto{\pgfqpoint{0.479238in}{1.208664in}}%
\pgfpathlineto{\pgfqpoint{0.464027in}{1.210445in}}%
\pgfpathlineto{\pgfqpoint{0.462181in}{1.217763in}}%
\pgfpathlineto{\pgfqpoint{0.454827in}{1.224336in}}%
\pgfpathlineto{\pgfqpoint{0.460433in}{1.233714in}}%
\pgfpathlineto{\pgfqpoint{0.459304in}{1.237786in}}%
\pgfpathlineto{\pgfqpoint{0.463136in}{1.244048in}}%
\pgfpathlineto{\pgfqpoint{0.460490in}{1.249495in}}%
\pgfpathlineto{\pgfqpoint{0.464659in}{1.256580in}}%
\pgfpathlineto{\pgfqpoint{0.467179in}{1.267774in}}%
\pgfpathlineto{\pgfqpoint{0.460247in}{1.270875in}}%
\pgfpathlineto{\pgfqpoint{0.454219in}{1.280640in}}%
\pgfpathlineto{\pgfqpoint{0.458963in}{1.289133in}}%
\pgfpathlineto{\pgfqpoint{0.457966in}{1.295353in}}%
\pgfpathlineto{\pgfqpoint{0.452130in}{1.297956in}}%
\pgfpathlineto{\pgfqpoint{0.446221in}{1.313638in}}%
\pgfpathlineto{\pgfqpoint{0.439451in}{1.318713in}}%
\pgfpathlineto{\pgfqpoint{0.438474in}{1.329831in}}%
\pgfpathlineto{\pgfqpoint{0.468241in}{1.320678in}}%
\pgfpathlineto{\pgfqpoint{0.508309in}{1.308909in}}%
\pgfpathlineto{\pgfqpoint{0.509532in}{1.308555in}}%
\pgfpathlineto{\pgfqpoint{0.505758in}{1.295236in}}%
\pgfpathlineto{\pgfqpoint{0.512448in}{1.293350in}}%
\pgfpathlineto{\pgfqpoint{0.510504in}{1.286661in}}%
\pgfpathlineto{\pgfqpoint{0.514797in}{1.280686in}}%
\pgfpathlineto{\pgfqpoint{0.521448in}{1.276381in}}%
\pgfpathlineto{\pgfqpoint{0.519557in}{1.269679in}}%
\pgfpathlineto{\pgfqpoint{0.523955in}{1.268417in}}%
\pgfpathlineto{\pgfqpoint{0.522085in}{1.261761in}}%
\pgfpathlineto{\pgfqpoint{0.530958in}{1.259199in}}%
\pgfpathlineto{\pgfqpoint{0.529070in}{1.252510in}}%
\pgfpathlineto{\pgfqpoint{0.536255in}{1.250941in}}%
\pgfpathlineto{\pgfqpoint{0.533846in}{1.244020in}}%
\pgfpathlineto{\pgfqpoint{0.539132in}{1.241671in}}%
\pgfpathlineto{\pgfqpoint{0.537333in}{1.227732in}}%
\pgfpathclose%
\pgfusepath{fill}%
\end{pgfscope}%
\begin{pgfscope}%
\pgfpathrectangle{\pgfqpoint{0.100000in}{0.100000in}}{\pgfqpoint{3.608454in}{2.310000in}}%
\pgfusepath{clip}%
\pgfsetbuttcap%
\pgfsetmiterjoin%
\definecolor{currentfill}{rgb}{0.000000,0.415686,0.792157}%
\pgfsetfillcolor{currentfill}%
\pgfsetlinewidth{0.000000pt}%
\definecolor{currentstroke}{rgb}{0.000000,0.000000,0.000000}%
\pgfsetstrokecolor{currentstroke}%
\pgfsetstrokeopacity{0.000000}%
\pgfsetdash{}{0pt}%
\pgfpathmoveto{\pgfqpoint{3.370781in}{1.834556in}}%
\pgfpathlineto{\pgfqpoint{3.348859in}{1.830296in}}%
\pgfpathlineto{\pgfqpoint{3.344700in}{1.848799in}}%
\pgfpathlineto{\pgfqpoint{3.333464in}{1.857292in}}%
\pgfpathlineto{\pgfqpoint{3.332649in}{1.866641in}}%
\pgfpathlineto{\pgfqpoint{3.324985in}{1.884098in}}%
\pgfpathlineto{\pgfqpoint{3.324873in}{1.895982in}}%
\pgfpathlineto{\pgfqpoint{3.327709in}{1.903478in}}%
\pgfpathlineto{\pgfqpoint{3.324626in}{1.910739in}}%
\pgfpathlineto{\pgfqpoint{3.323703in}{1.920977in}}%
\pgfpathlineto{\pgfqpoint{3.316899in}{1.928969in}}%
\pgfpathlineto{\pgfqpoint{3.316446in}{1.944043in}}%
\pgfpathlineto{\pgfqpoint{3.313379in}{1.946070in}}%
\pgfpathlineto{\pgfqpoint{3.314088in}{1.952894in}}%
\pgfpathlineto{\pgfqpoint{3.312126in}{1.959898in}}%
\pgfpathlineto{\pgfqpoint{3.351137in}{1.969833in}}%
\pgfpathlineto{\pgfqpoint{3.355196in}{1.970349in}}%
\pgfpathlineto{\pgfqpoint{3.359777in}{1.957182in}}%
\pgfpathlineto{\pgfqpoint{3.358298in}{1.952554in}}%
\pgfpathlineto{\pgfqpoint{3.368238in}{1.950258in}}%
\pgfpathlineto{\pgfqpoint{3.366415in}{1.942507in}}%
\pgfpathlineto{\pgfqpoint{3.373495in}{1.940927in}}%
\pgfpathlineto{\pgfqpoint{3.371539in}{1.933332in}}%
\pgfpathlineto{\pgfqpoint{3.370017in}{1.926361in}}%
\pgfpathlineto{\pgfqpoint{3.363137in}{1.929015in}}%
\pgfpathlineto{\pgfqpoint{3.361540in}{1.921976in}}%
\pgfpathlineto{\pgfqpoint{3.352226in}{1.924022in}}%
\pgfpathlineto{\pgfqpoint{3.349620in}{1.909475in}}%
\pgfpathlineto{\pgfqpoint{3.350605in}{1.900061in}}%
\pgfpathlineto{\pgfqpoint{3.353338in}{1.893246in}}%
\pgfpathlineto{\pgfqpoint{3.363839in}{1.892385in}}%
\pgfpathlineto{\pgfqpoint{3.358819in}{1.880708in}}%
\pgfpathlineto{\pgfqpoint{3.354162in}{1.876806in}}%
\pgfpathlineto{\pgfqpoint{3.365877in}{1.874467in}}%
\pgfpathlineto{\pgfqpoint{3.365489in}{1.867233in}}%
\pgfpathlineto{\pgfqpoint{3.372767in}{1.866271in}}%
\pgfpathlineto{\pgfqpoint{3.371356in}{1.850568in}}%
\pgfpathlineto{\pgfqpoint{3.375426in}{1.849437in}}%
\pgfpathlineto{\pgfqpoint{3.375120in}{1.840353in}}%
\pgfpathlineto{\pgfqpoint{3.370781in}{1.834556in}}%
\pgfpathclose%
\pgfusepath{fill}%
\end{pgfscope}%
\begin{pgfscope}%
\pgfpathrectangle{\pgfqpoint{0.100000in}{0.100000in}}{\pgfqpoint{3.608454in}{2.310000in}}%
\pgfusepath{clip}%
\pgfsetbuttcap%
\pgfsetmiterjoin%
\definecolor{currentfill}{rgb}{0.000000,0.745098,0.627451}%
\pgfsetfillcolor{currentfill}%
\pgfsetlinewidth{0.000000pt}%
\definecolor{currentstroke}{rgb}{0.000000,0.000000,0.000000}%
\pgfsetstrokecolor{currentstroke}%
\pgfsetstrokeopacity{0.000000}%
\pgfsetdash{}{0pt}%
\pgfpathmoveto{\pgfqpoint{2.342946in}{0.805709in}}%
\pgfpathlineto{\pgfqpoint{2.340198in}{0.812791in}}%
\pgfpathlineto{\pgfqpoint{2.344014in}{0.814843in}}%
\pgfpathlineto{\pgfqpoint{2.342569in}{0.828917in}}%
\pgfpathlineto{\pgfqpoint{2.344269in}{0.832393in}}%
\pgfpathlineto{\pgfqpoint{2.365973in}{0.833290in}}%
\pgfpathlineto{\pgfqpoint{2.371879in}{0.836595in}}%
\pgfpathlineto{\pgfqpoint{2.377206in}{0.836777in}}%
\pgfpathlineto{\pgfqpoint{2.378401in}{0.826111in}}%
\pgfpathlineto{\pgfqpoint{2.384564in}{0.826478in}}%
\pgfpathlineto{\pgfqpoint{2.384135in}{0.837054in}}%
\pgfpathlineto{\pgfqpoint{2.393171in}{0.837522in}}%
\pgfpathlineto{\pgfqpoint{2.393376in}{0.840961in}}%
\pgfpathlineto{\pgfqpoint{2.404935in}{0.832654in}}%
\pgfpathlineto{\pgfqpoint{2.401669in}{0.830402in}}%
\pgfpathlineto{\pgfqpoint{2.399936in}{0.821863in}}%
\pgfpathlineto{\pgfqpoint{2.395425in}{0.819431in}}%
\pgfpathlineto{\pgfqpoint{2.394278in}{0.809641in}}%
\pgfpathlineto{\pgfqpoint{2.396163in}{0.790004in}}%
\pgfpathlineto{\pgfqpoint{2.395693in}{0.767722in}}%
\pgfpathlineto{\pgfqpoint{2.395986in}{0.760766in}}%
\pgfpathlineto{\pgfqpoint{2.367747in}{0.759421in}}%
\pgfpathlineto{\pgfqpoint{2.357930in}{0.763773in}}%
\pgfpathlineto{\pgfqpoint{2.353993in}{0.777417in}}%
\pgfpathlineto{\pgfqpoint{2.363437in}{0.782994in}}%
\pgfpathlineto{\pgfqpoint{2.364252in}{0.787285in}}%
\pgfpathlineto{\pgfqpoint{2.371635in}{0.792020in}}%
\pgfpathlineto{\pgfqpoint{2.374090in}{0.798791in}}%
\pgfpathlineto{\pgfqpoint{2.365200in}{0.801106in}}%
\pgfpathlineto{\pgfqpoint{2.366944in}{0.810346in}}%
\pgfpathlineto{\pgfqpoint{2.354522in}{0.806095in}}%
\pgfpathlineto{\pgfqpoint{2.342946in}{0.805709in}}%
\pgfpathclose%
\pgfusepath{fill}%
\end{pgfscope}%
\begin{pgfscope}%
\pgfpathrectangle{\pgfqpoint{0.100000in}{0.100000in}}{\pgfqpoint{3.608454in}{2.310000in}}%
\pgfusepath{clip}%
\pgfsetbuttcap%
\pgfsetmiterjoin%
\definecolor{currentfill}{rgb}{0.000000,0.760784,0.619608}%
\pgfsetfillcolor{currentfill}%
\pgfsetlinewidth{0.000000pt}%
\definecolor{currentstroke}{rgb}{0.000000,0.000000,0.000000}%
\pgfsetstrokecolor{currentstroke}%
\pgfsetstrokeopacity{0.000000}%
\pgfsetdash{}{0pt}%
\pgfpathmoveto{\pgfqpoint{0.509532in}{1.308555in}}%
\pgfpathlineto{\pgfqpoint{0.508309in}{1.308909in}}%
\pgfpathlineto{\pgfqpoint{0.503087in}{1.324821in}}%
\pgfpathlineto{\pgfqpoint{0.498712in}{1.326448in}}%
\pgfpathlineto{\pgfqpoint{0.493971in}{1.332372in}}%
\pgfpathlineto{\pgfqpoint{0.488481in}{1.345190in}}%
\pgfpathlineto{\pgfqpoint{0.492171in}{1.356794in}}%
\pgfpathlineto{\pgfqpoint{0.496557in}{1.356685in}}%
\pgfpathlineto{\pgfqpoint{0.499634in}{1.368954in}}%
\pgfpathlineto{\pgfqpoint{0.486008in}{1.393700in}}%
\pgfpathlineto{\pgfqpoint{0.520707in}{1.412942in}}%
\pgfpathlineto{\pgfqpoint{0.533257in}{1.413969in}}%
\pgfpathlineto{\pgfqpoint{0.548331in}{1.412190in}}%
\pgfpathlineto{\pgfqpoint{0.570915in}{1.425063in}}%
\pgfpathlineto{\pgfqpoint{0.577622in}{1.423252in}}%
\pgfpathlineto{\pgfqpoint{0.583298in}{1.427984in}}%
\pgfpathlineto{\pgfqpoint{0.584724in}{1.432975in}}%
\pgfpathlineto{\pgfqpoint{0.605864in}{1.444998in}}%
\pgfpathlineto{\pgfqpoint{0.607467in}{1.441382in}}%
\pgfpathlineto{\pgfqpoint{0.616083in}{1.438599in}}%
\pgfpathlineto{\pgfqpoint{0.618630in}{1.433861in}}%
\pgfpathlineto{\pgfqpoint{0.619173in}{1.425619in}}%
\pgfpathlineto{\pgfqpoint{0.624814in}{1.421196in}}%
\pgfpathlineto{\pgfqpoint{0.627031in}{1.414709in}}%
\pgfpathlineto{\pgfqpoint{0.632163in}{1.410294in}}%
\pgfpathlineto{\pgfqpoint{0.628697in}{1.402513in}}%
\pgfpathlineto{\pgfqpoint{0.632454in}{1.400962in}}%
\pgfpathlineto{\pgfqpoint{0.632393in}{1.385724in}}%
\pgfpathlineto{\pgfqpoint{0.636091in}{1.384439in}}%
\pgfpathlineto{\pgfqpoint{0.643441in}{1.374382in}}%
\pgfpathlineto{\pgfqpoint{0.645186in}{1.361551in}}%
\pgfpathlineto{\pgfqpoint{0.642344in}{1.357942in}}%
\pgfpathlineto{\pgfqpoint{0.641670in}{1.344858in}}%
\pgfpathlineto{\pgfqpoint{0.644925in}{1.333356in}}%
\pgfpathlineto{\pgfqpoint{0.644042in}{1.329081in}}%
\pgfpathlineto{\pgfqpoint{0.650356in}{1.315914in}}%
\pgfpathlineto{\pgfqpoint{0.647795in}{1.312826in}}%
\pgfpathlineto{\pgfqpoint{0.649389in}{1.301776in}}%
\pgfpathlineto{\pgfqpoint{0.647370in}{1.296028in}}%
\pgfpathlineto{\pgfqpoint{0.648713in}{1.279430in}}%
\pgfpathlineto{\pgfqpoint{0.645007in}{1.271759in}}%
\pgfpathlineto{\pgfqpoint{0.606581in}{1.281871in}}%
\pgfpathlineto{\pgfqpoint{0.550101in}{1.297200in}}%
\pgfpathlineto{\pgfqpoint{0.509532in}{1.308555in}}%
\pgfpathclose%
\pgfusepath{fill}%
\end{pgfscope}%
\begin{pgfscope}%
\pgfpathrectangle{\pgfqpoint{0.100000in}{0.100000in}}{\pgfqpoint{3.608454in}{2.310000in}}%
\pgfusepath{clip}%
\pgfsetbuttcap%
\pgfsetmiterjoin%
\definecolor{currentfill}{rgb}{0.000000,0.564706,0.717647}%
\pgfsetfillcolor{currentfill}%
\pgfsetlinewidth{0.000000pt}%
\definecolor{currentstroke}{rgb}{0.000000,0.000000,0.000000}%
\pgfsetstrokecolor{currentstroke}%
\pgfsetstrokeopacity{0.000000}%
\pgfsetdash{}{0pt}%
\pgfpathmoveto{\pgfqpoint{1.007808in}{2.262053in}}%
\pgfpathlineto{\pgfqpoint{1.050960in}{2.251947in}}%
\pgfpathlineto{\pgfqpoint{1.075629in}{2.246526in}}%
\pgfpathlineto{\pgfqpoint{1.074794in}{2.233716in}}%
\pgfpathlineto{\pgfqpoint{1.075971in}{2.228263in}}%
\pgfpathlineto{\pgfqpoint{1.072055in}{2.221371in}}%
\pgfpathlineto{\pgfqpoint{1.061432in}{2.221573in}}%
\pgfpathlineto{\pgfqpoint{1.061972in}{2.214474in}}%
\pgfpathlineto{\pgfqpoint{1.056030in}{2.187273in}}%
\pgfpathlineto{\pgfqpoint{1.047061in}{2.189256in}}%
\pgfpathlineto{\pgfqpoint{1.043899in}{2.172864in}}%
\pgfpathlineto{\pgfqpoint{1.036396in}{2.174561in}}%
\pgfpathlineto{\pgfqpoint{1.033549in}{2.168208in}}%
\pgfpathlineto{\pgfqpoint{1.027896in}{2.166208in}}%
\pgfpathlineto{\pgfqpoint{1.015079in}{2.170487in}}%
\pgfpathlineto{\pgfqpoint{1.015735in}{2.176278in}}%
\pgfpathlineto{\pgfqpoint{1.012094in}{2.185135in}}%
\pgfpathlineto{\pgfqpoint{1.011635in}{2.199142in}}%
\pgfpathlineto{\pgfqpoint{0.994462in}{2.195826in}}%
\pgfpathlineto{\pgfqpoint{0.993310in}{2.200722in}}%
\pgfpathlineto{\pgfqpoint{0.998620in}{2.223097in}}%
\pgfpathlineto{\pgfqpoint{1.007808in}{2.262053in}}%
\pgfpathclose%
\pgfusepath{fill}%
\end{pgfscope}%
\begin{pgfscope}%
\pgfpathrectangle{\pgfqpoint{0.100000in}{0.100000in}}{\pgfqpoint{3.608454in}{2.310000in}}%
\pgfusepath{clip}%
\pgfsetbuttcap%
\pgfsetmiterjoin%
\definecolor{currentfill}{rgb}{0.000000,0.631373,0.684314}%
\pgfsetfillcolor{currentfill}%
\pgfsetlinewidth{0.000000pt}%
\definecolor{currentstroke}{rgb}{0.000000,0.000000,0.000000}%
\pgfsetstrokecolor{currentstroke}%
\pgfsetstrokeopacity{0.000000}%
\pgfsetdash{}{0pt}%
\pgfpathmoveto{\pgfqpoint{2.515690in}{0.998385in}}%
\pgfpathlineto{\pgfqpoint{2.515503in}{1.018069in}}%
\pgfpathlineto{\pgfqpoint{2.512709in}{1.025980in}}%
\pgfpathlineto{\pgfqpoint{2.508135in}{1.025655in}}%
\pgfpathlineto{\pgfqpoint{2.507414in}{1.036455in}}%
\pgfpathlineto{\pgfqpoint{2.509787in}{1.036620in}}%
\pgfpathlineto{\pgfqpoint{2.508734in}{1.056586in}}%
\pgfpathlineto{\pgfqpoint{2.512736in}{1.056859in}}%
\pgfpathlineto{\pgfqpoint{2.517201in}{1.066138in}}%
\pgfpathlineto{\pgfqpoint{2.521534in}{1.068158in}}%
\pgfpathlineto{\pgfqpoint{2.535012in}{1.069064in}}%
\pgfpathlineto{\pgfqpoint{2.534914in}{1.072050in}}%
\pgfpathlineto{\pgfqpoint{2.550818in}{1.070538in}}%
\pgfpathlineto{\pgfqpoint{2.556697in}{1.071568in}}%
\pgfpathlineto{\pgfqpoint{2.560066in}{1.064151in}}%
\pgfpathlineto{\pgfqpoint{2.561622in}{1.041248in}}%
\pgfpathlineto{\pgfqpoint{2.547499in}{1.040339in}}%
\pgfpathlineto{\pgfqpoint{2.554995in}{1.031701in}}%
\pgfpathlineto{\pgfqpoint{2.553769in}{0.997504in}}%
\pgfpathlineto{\pgfqpoint{2.528607in}{0.995823in}}%
\pgfpathlineto{\pgfqpoint{2.528379in}{0.999297in}}%
\pgfpathlineto{\pgfqpoint{2.515690in}{0.998385in}}%
\pgfpathclose%
\pgfusepath{fill}%
\end{pgfscope}%
\begin{pgfscope}%
\pgfpathrectangle{\pgfqpoint{0.100000in}{0.100000in}}{\pgfqpoint{3.608454in}{2.310000in}}%
\pgfusepath{clip}%
\pgfsetbuttcap%
\pgfsetmiterjoin%
\definecolor{currentfill}{rgb}{0.000000,0.992157,0.503922}%
\pgfsetfillcolor{currentfill}%
\pgfsetlinewidth{0.000000pt}%
\definecolor{currentstroke}{rgb}{0.000000,0.000000,0.000000}%
\pgfsetstrokecolor{currentstroke}%
\pgfsetstrokeopacity{0.000000}%
\pgfsetdash{}{0pt}%
\pgfpathmoveto{\pgfqpoint{0.647635in}{2.272305in}}%
\pgfpathlineto{\pgfqpoint{0.648220in}{2.267247in}}%
\pgfpathlineto{\pgfqpoint{0.639739in}{2.262486in}}%
\pgfpathlineto{\pgfqpoint{0.635231in}{2.255252in}}%
\pgfpathlineto{\pgfqpoint{0.623636in}{2.252362in}}%
\pgfpathlineto{\pgfqpoint{0.595222in}{2.261545in}}%
\pgfpathlineto{\pgfqpoint{0.593017in}{2.254746in}}%
\pgfpathlineto{\pgfqpoint{0.565186in}{2.263908in}}%
\pgfpathlineto{\pgfqpoint{0.549782in}{2.270356in}}%
\pgfpathlineto{\pgfqpoint{0.551434in}{2.287191in}}%
\pgfpathlineto{\pgfqpoint{0.549994in}{2.294555in}}%
\pgfpathlineto{\pgfqpoint{0.544785in}{2.302357in}}%
\pgfpathlineto{\pgfqpoint{0.544644in}{2.308702in}}%
\pgfpathlineto{\pgfqpoint{0.548618in}{2.325631in}}%
\pgfpathlineto{\pgfqpoint{0.557234in}{2.339937in}}%
\pgfpathlineto{\pgfqpoint{0.566992in}{2.328262in}}%
\pgfpathlineto{\pgfqpoint{0.573972in}{2.323699in}}%
\pgfpathlineto{\pgfqpoint{0.582903in}{2.313221in}}%
\pgfpathlineto{\pgfqpoint{0.591705in}{2.308913in}}%
\pgfpathlineto{\pgfqpoint{0.598948in}{2.307790in}}%
\pgfpathlineto{\pgfqpoint{0.604150in}{2.302970in}}%
\pgfpathlineto{\pgfqpoint{0.620717in}{2.296209in}}%
\pgfpathlineto{\pgfqpoint{0.627218in}{2.298959in}}%
\pgfpathlineto{\pgfqpoint{0.630726in}{2.290276in}}%
\pgfpathlineto{\pgfqpoint{0.646614in}{2.290179in}}%
\pgfpathlineto{\pgfqpoint{0.648837in}{2.285967in}}%
\pgfpathlineto{\pgfqpoint{0.646729in}{2.278658in}}%
\pgfpathlineto{\pgfqpoint{0.647635in}{2.272305in}}%
\pgfpathclose%
\pgfusepath{fill}%
\end{pgfscope}%
\begin{pgfscope}%
\pgfpathrectangle{\pgfqpoint{0.100000in}{0.100000in}}{\pgfqpoint{3.608454in}{2.310000in}}%
\pgfusepath{clip}%
\pgfsetbuttcap%
\pgfsetmiterjoin%
\definecolor{currentfill}{rgb}{0.000000,0.639216,0.680392}%
\pgfsetfillcolor{currentfill}%
\pgfsetlinewidth{0.000000pt}%
\definecolor{currentstroke}{rgb}{0.000000,0.000000,0.000000}%
\pgfsetstrokecolor{currentstroke}%
\pgfsetstrokeopacity{0.000000}%
\pgfsetdash{}{0pt}%
\pgfpathmoveto{\pgfqpoint{2.726071in}{1.419220in}}%
\pgfpathlineto{\pgfqpoint{2.726448in}{1.415659in}}%
\pgfpathlineto{\pgfqpoint{2.697158in}{1.412713in}}%
\pgfpathlineto{\pgfqpoint{2.694747in}{1.433296in}}%
\pgfpathlineto{\pgfqpoint{2.699598in}{1.433892in}}%
\pgfpathlineto{\pgfqpoint{2.698943in}{1.451000in}}%
\pgfpathlineto{\pgfqpoint{2.722450in}{1.453627in}}%
\pgfpathlineto{\pgfqpoint{2.726071in}{1.419220in}}%
\pgfpathclose%
\pgfusepath{fill}%
\end{pgfscope}%
\begin{pgfscope}%
\pgfpathrectangle{\pgfqpoint{0.100000in}{0.100000in}}{\pgfqpoint{3.608454in}{2.310000in}}%
\pgfusepath{clip}%
\pgfsetbuttcap%
\pgfsetmiterjoin%
\definecolor{currentfill}{rgb}{0.000000,0.529412,0.735294}%
\pgfsetfillcolor{currentfill}%
\pgfsetlinewidth{0.000000pt}%
\definecolor{currentstroke}{rgb}{0.000000,0.000000,0.000000}%
\pgfsetstrokecolor{currentstroke}%
\pgfsetstrokeopacity{0.000000}%
\pgfsetdash{}{0pt}%
\pgfpathmoveto{\pgfqpoint{2.266038in}{1.833884in}}%
\pgfpathlineto{\pgfqpoint{2.232360in}{1.832947in}}%
\pgfpathlineto{\pgfqpoint{2.232161in}{1.839416in}}%
\pgfpathlineto{\pgfqpoint{2.237996in}{1.848119in}}%
\pgfpathlineto{\pgfqpoint{2.238078in}{1.851709in}}%
\pgfpathlineto{\pgfqpoint{2.230979in}{1.861346in}}%
\pgfpathlineto{\pgfqpoint{2.224641in}{1.861874in}}%
\pgfpathlineto{\pgfqpoint{2.226726in}{1.874225in}}%
\pgfpathlineto{\pgfqpoint{2.229731in}{1.876962in}}%
\pgfpathlineto{\pgfqpoint{2.233501in}{1.887214in}}%
\pgfpathlineto{\pgfqpoint{2.237467in}{1.890615in}}%
\pgfpathlineto{\pgfqpoint{2.253310in}{1.897735in}}%
\pgfpathlineto{\pgfqpoint{2.256290in}{1.902446in}}%
\pgfpathlineto{\pgfqpoint{2.256111in}{1.909051in}}%
\pgfpathlineto{\pgfqpoint{2.269571in}{1.909501in}}%
\pgfpathlineto{\pgfqpoint{2.270918in}{1.895690in}}%
\pgfpathlineto{\pgfqpoint{2.271935in}{1.868329in}}%
\pgfpathlineto{\pgfqpoint{2.265057in}{1.868094in}}%
\pgfpathlineto{\pgfqpoint{2.266038in}{1.833884in}}%
\pgfpathclose%
\pgfusepath{fill}%
\end{pgfscope}%
\begin{pgfscope}%
\pgfpathrectangle{\pgfqpoint{0.100000in}{0.100000in}}{\pgfqpoint{3.608454in}{2.310000in}}%
\pgfusepath{clip}%
\pgfsetbuttcap%
\pgfsetmiterjoin%
\definecolor{currentfill}{rgb}{0.000000,0.560784,0.719608}%
\pgfsetfillcolor{currentfill}%
\pgfsetlinewidth{0.000000pt}%
\definecolor{currentstroke}{rgb}{0.000000,0.000000,0.000000}%
\pgfsetstrokecolor{currentstroke}%
\pgfsetstrokeopacity{0.000000}%
\pgfsetdash{}{0pt}%
\pgfpathmoveto{\pgfqpoint{2.442684in}{0.864247in}}%
\pgfpathlineto{\pgfqpoint{2.438751in}{0.854506in}}%
\pgfpathlineto{\pgfqpoint{2.428148in}{0.846306in}}%
\pgfpathlineto{\pgfqpoint{2.412628in}{0.839279in}}%
\pgfpathlineto{\pgfqpoint{2.404935in}{0.832654in}}%
\pgfpathlineto{\pgfqpoint{2.393376in}{0.840961in}}%
\pgfpathlineto{\pgfqpoint{2.393315in}{0.855864in}}%
\pgfpathlineto{\pgfqpoint{2.389913in}{0.858059in}}%
\pgfpathlineto{\pgfqpoint{2.396470in}{0.865278in}}%
\pgfpathlineto{\pgfqpoint{2.395715in}{0.879087in}}%
\pgfpathlineto{\pgfqpoint{2.391193in}{0.892674in}}%
\pgfpathlineto{\pgfqpoint{2.401868in}{0.893227in}}%
\pgfpathlineto{\pgfqpoint{2.401625in}{0.897879in}}%
\pgfpathlineto{\pgfqpoint{2.408409in}{0.898240in}}%
\pgfpathlineto{\pgfqpoint{2.408616in}{0.893304in}}%
\pgfpathlineto{\pgfqpoint{2.416578in}{0.896534in}}%
\pgfpathlineto{\pgfqpoint{2.417347in}{0.883236in}}%
\pgfpathlineto{\pgfqpoint{2.410750in}{0.879777in}}%
\pgfpathlineto{\pgfqpoint{2.417323in}{0.873285in}}%
\pgfpathlineto{\pgfqpoint{2.442684in}{0.864247in}}%
\pgfpathclose%
\pgfusepath{fill}%
\end{pgfscope}%
\begin{pgfscope}%
\pgfpathrectangle{\pgfqpoint{0.100000in}{0.100000in}}{\pgfqpoint{3.608454in}{2.310000in}}%
\pgfusepath{clip}%
\pgfsetbuttcap%
\pgfsetmiterjoin%
\definecolor{currentfill}{rgb}{0.000000,0.741176,0.629412}%
\pgfsetfillcolor{currentfill}%
\pgfsetlinewidth{0.000000pt}%
\definecolor{currentstroke}{rgb}{0.000000,0.000000,0.000000}%
\pgfsetstrokecolor{currentstroke}%
\pgfsetstrokeopacity{0.000000}%
\pgfsetdash{}{0pt}%
\pgfpathmoveto{\pgfqpoint{2.172440in}{1.907196in}}%
\pgfpathlineto{\pgfqpoint{2.172623in}{1.893374in}}%
\pgfpathlineto{\pgfqpoint{2.174916in}{1.893366in}}%
\pgfpathlineto{\pgfqpoint{2.175363in}{1.880599in}}%
\pgfpathlineto{\pgfqpoint{2.142690in}{1.880249in}}%
\pgfpathlineto{\pgfqpoint{2.147223in}{1.876558in}}%
\pgfpathlineto{\pgfqpoint{2.099037in}{1.876260in}}%
\pgfpathlineto{\pgfqpoint{2.098775in}{1.916792in}}%
\pgfpathlineto{\pgfqpoint{2.097851in}{1.951504in}}%
\pgfpathlineto{\pgfqpoint{2.097902in}{1.958459in}}%
\pgfpathlineto{\pgfqpoint{2.118427in}{1.958403in}}%
\pgfpathlineto{\pgfqpoint{2.118937in}{1.925718in}}%
\pgfpathlineto{\pgfqpoint{2.132410in}{1.918482in}}%
\pgfpathlineto{\pgfqpoint{2.138352in}{1.920870in}}%
\pgfpathlineto{\pgfqpoint{2.142555in}{1.917955in}}%
\pgfpathlineto{\pgfqpoint{2.141306in}{1.906925in}}%
\pgfpathlineto{\pgfqpoint{2.172440in}{1.907196in}}%
\pgfpathclose%
\pgfusepath{fill}%
\end{pgfscope}%
\begin{pgfscope}%
\pgfpathrectangle{\pgfqpoint{0.100000in}{0.100000in}}{\pgfqpoint{3.608454in}{2.310000in}}%
\pgfusepath{clip}%
\pgfsetbuttcap%
\pgfsetmiterjoin%
\definecolor{currentfill}{rgb}{0.000000,0.745098,0.627451}%
\pgfsetfillcolor{currentfill}%
\pgfsetlinewidth{0.000000pt}%
\definecolor{currentstroke}{rgb}{0.000000,0.000000,0.000000}%
\pgfsetstrokecolor{currentstroke}%
\pgfsetstrokeopacity{0.000000}%
\pgfsetdash{}{0pt}%
\pgfpathmoveto{\pgfqpoint{2.231585in}{0.901788in}}%
\pgfpathlineto{\pgfqpoint{2.231529in}{0.896018in}}%
\pgfpathlineto{\pgfqpoint{2.240853in}{0.895980in}}%
\pgfpathlineto{\pgfqpoint{2.240807in}{0.867381in}}%
\pgfpathlineto{\pgfqpoint{2.186423in}{0.866575in}}%
\pgfpathlineto{\pgfqpoint{2.182676in}{0.872245in}}%
\pgfpathlineto{\pgfqpoint{2.183288in}{0.876547in}}%
\pgfpathlineto{\pgfqpoint{2.189397in}{0.878197in}}%
\pgfpathlineto{\pgfqpoint{2.195345in}{0.891554in}}%
\pgfpathlineto{\pgfqpoint{2.190887in}{0.898874in}}%
\pgfpathlineto{\pgfqpoint{2.191374in}{0.903396in}}%
\pgfpathlineto{\pgfqpoint{2.207261in}{0.903173in}}%
\pgfpathlineto{\pgfqpoint{2.207240in}{0.900881in}}%
\pgfpathlineto{\pgfqpoint{2.224654in}{0.900658in}}%
\pgfpathlineto{\pgfqpoint{2.231585in}{0.901788in}}%
\pgfpathclose%
\pgfusepath{fill}%
\end{pgfscope}%
\begin{pgfscope}%
\pgfpathrectangle{\pgfqpoint{0.100000in}{0.100000in}}{\pgfqpoint{3.608454in}{2.310000in}}%
\pgfusepath{clip}%
\pgfsetbuttcap%
\pgfsetmiterjoin%
\definecolor{currentfill}{rgb}{0.000000,0.282353,0.858824}%
\pgfsetfillcolor{currentfill}%
\pgfsetlinewidth{0.000000pt}%
\definecolor{currentstroke}{rgb}{0.000000,0.000000,0.000000}%
\pgfsetstrokecolor{currentstroke}%
\pgfsetstrokeopacity{0.000000}%
\pgfsetdash{}{0pt}%
\pgfpathmoveto{\pgfqpoint{2.454210in}{1.655405in}}%
\pgfpathlineto{\pgfqpoint{2.433615in}{1.653903in}}%
\pgfpathlineto{\pgfqpoint{2.406360in}{1.653264in}}%
\pgfpathlineto{\pgfqpoint{2.406543in}{1.649815in}}%
\pgfpathlineto{\pgfqpoint{2.372395in}{1.647911in}}%
\pgfpathlineto{\pgfqpoint{2.370685in}{1.678611in}}%
\pgfpathlineto{\pgfqpoint{2.363976in}{1.679112in}}%
\pgfpathlineto{\pgfqpoint{2.357123in}{1.675633in}}%
\pgfpathlineto{\pgfqpoint{2.356072in}{1.695517in}}%
\pgfpathlineto{\pgfqpoint{2.355403in}{1.705801in}}%
\pgfpathlineto{\pgfqpoint{2.376100in}{1.706906in}}%
\pgfpathlineto{\pgfqpoint{2.375661in}{1.713794in}}%
\pgfpathlineto{\pgfqpoint{2.405834in}{1.715444in}}%
\pgfpathlineto{\pgfqpoint{2.416473in}{1.716189in}}%
\pgfpathlineto{\pgfqpoint{2.414986in}{1.743088in}}%
\pgfpathlineto{\pgfqpoint{2.439466in}{1.744650in}}%
\pgfpathlineto{\pgfqpoint{2.440017in}{1.736230in}}%
\pgfpathlineto{\pgfqpoint{2.436156in}{1.727226in}}%
\pgfpathlineto{\pgfqpoint{2.436777in}{1.717456in}}%
\pgfpathlineto{\pgfqpoint{2.450492in}{1.717550in}}%
\pgfpathlineto{\pgfqpoint{2.454210in}{1.655405in}}%
\pgfpathclose%
\pgfusepath{fill}%
\end{pgfscope}%
\begin{pgfscope}%
\pgfpathrectangle{\pgfqpoint{0.100000in}{0.100000in}}{\pgfqpoint{3.608454in}{2.310000in}}%
\pgfusepath{clip}%
\pgfsetbuttcap%
\pgfsetmiterjoin%
\definecolor{currentfill}{rgb}{0.000000,0.388235,0.805882}%
\pgfsetfillcolor{currentfill}%
\pgfsetlinewidth{0.000000pt}%
\definecolor{currentstroke}{rgb}{0.000000,0.000000,0.000000}%
\pgfsetstrokecolor{currentstroke}%
\pgfsetstrokeopacity{0.000000}%
\pgfsetdash{}{0pt}%
\pgfpathmoveto{\pgfqpoint{1.278870in}{1.882470in}}%
\pgfpathlineto{\pgfqpoint{1.317490in}{1.876575in}}%
\pgfpathlineto{\pgfqpoint{1.318547in}{1.875904in}}%
\pgfpathlineto{\pgfqpoint{1.364865in}{1.868873in}}%
\pgfpathlineto{\pgfqpoint{1.383661in}{1.866297in}}%
\pgfpathlineto{\pgfqpoint{1.384830in}{1.855882in}}%
\pgfpathlineto{\pgfqpoint{1.389944in}{1.843550in}}%
\pgfpathlineto{\pgfqpoint{1.396824in}{1.840362in}}%
\pgfpathlineto{\pgfqpoint{1.406562in}{1.831802in}}%
\pgfpathlineto{\pgfqpoint{1.412375in}{1.821775in}}%
\pgfpathlineto{\pgfqpoint{1.417230in}{1.817131in}}%
\pgfpathlineto{\pgfqpoint{1.419171in}{1.807060in}}%
\pgfpathlineto{\pgfqpoint{1.417395in}{1.794618in}}%
\pgfpathlineto{\pgfqpoint{1.338318in}{1.806097in}}%
\pgfpathlineto{\pgfqpoint{1.337257in}{1.799111in}}%
\pgfpathlineto{\pgfqpoint{1.323648in}{1.801191in}}%
\pgfpathlineto{\pgfqpoint{1.322614in}{1.794241in}}%
\pgfpathlineto{\pgfqpoint{1.315604in}{1.795301in}}%
\pgfpathlineto{\pgfqpoint{1.314743in}{1.788523in}}%
\pgfpathlineto{\pgfqpoint{1.304569in}{1.790081in}}%
\pgfpathlineto{\pgfqpoint{1.303529in}{1.783157in}}%
\pgfpathlineto{\pgfqpoint{1.291011in}{1.784895in}}%
\pgfpathlineto{\pgfqpoint{1.285167in}{1.795237in}}%
\pgfpathlineto{\pgfqpoint{1.279919in}{1.798565in}}%
\pgfpathlineto{\pgfqpoint{1.273821in}{1.796453in}}%
\pgfpathlineto{\pgfqpoint{1.271394in}{1.792453in}}%
\pgfpathlineto{\pgfqpoint{1.264050in}{1.788274in}}%
\pgfpathlineto{\pgfqpoint{1.260477in}{1.800714in}}%
\pgfpathlineto{\pgfqpoint{1.254530in}{1.801740in}}%
\pgfpathlineto{\pgfqpoint{1.251746in}{1.806946in}}%
\pgfpathlineto{\pgfqpoint{1.253366in}{1.816700in}}%
\pgfpathlineto{\pgfqpoint{1.250644in}{1.821729in}}%
\pgfpathlineto{\pgfqpoint{1.249318in}{1.830814in}}%
\pgfpathlineto{\pgfqpoint{1.243670in}{1.842268in}}%
\pgfpathlineto{\pgfqpoint{1.245880in}{1.848777in}}%
\pgfpathlineto{\pgfqpoint{1.241305in}{1.854999in}}%
\pgfpathlineto{\pgfqpoint{1.224995in}{1.857794in}}%
\pgfpathlineto{\pgfqpoint{1.226134in}{1.864328in}}%
\pgfpathlineto{\pgfqpoint{1.204580in}{1.868127in}}%
\pgfpathlineto{\pgfqpoint{1.209266in}{1.894331in}}%
\pgfpathlineto{\pgfqpoint{1.228507in}{1.890208in}}%
\pgfpathlineto{\pgfqpoint{1.278870in}{1.882470in}}%
\pgfpathclose%
\pgfusepath{fill}%
\end{pgfscope}%
\begin{pgfscope}%
\pgfpathrectangle{\pgfqpoint{0.100000in}{0.100000in}}{\pgfqpoint{3.608454in}{2.310000in}}%
\pgfusepath{clip}%
\pgfsetbuttcap%
\pgfsetmiterjoin%
\definecolor{currentfill}{rgb}{0.000000,0.470588,0.764706}%
\pgfsetfillcolor{currentfill}%
\pgfsetlinewidth{0.000000pt}%
\definecolor{currentstroke}{rgb}{0.000000,0.000000,0.000000}%
\pgfsetstrokecolor{currentstroke}%
\pgfsetstrokeopacity{0.000000}%
\pgfsetdash{}{0pt}%
\pgfpathmoveto{\pgfqpoint{2.551959in}{0.940228in}}%
\pgfpathlineto{\pgfqpoint{2.546945in}{0.939900in}}%
\pgfpathlineto{\pgfqpoint{2.544297in}{0.934956in}}%
\pgfpathlineto{\pgfqpoint{2.535053in}{0.931443in}}%
\pgfpathlineto{\pgfqpoint{2.521327in}{0.932876in}}%
\pgfpathlineto{\pgfqpoint{2.520583in}{0.943229in}}%
\pgfpathlineto{\pgfqpoint{2.506684in}{0.942391in}}%
\pgfpathlineto{\pgfqpoint{2.506964in}{0.937791in}}%
\pgfpathlineto{\pgfqpoint{2.489985in}{0.934122in}}%
\pgfpathlineto{\pgfqpoint{2.469213in}{0.932802in}}%
\pgfpathlineto{\pgfqpoint{2.468504in}{0.944341in}}%
\pgfpathlineto{\pgfqpoint{2.466935in}{0.967664in}}%
\pgfpathlineto{\pgfqpoint{2.484251in}{0.968666in}}%
\pgfpathlineto{\pgfqpoint{2.483100in}{0.985946in}}%
\pgfpathlineto{\pgfqpoint{2.503832in}{0.987282in}}%
\pgfpathlineto{\pgfqpoint{2.512839in}{0.989031in}}%
\pgfpathlineto{\pgfqpoint{2.512313in}{0.997036in}}%
\pgfpathlineto{\pgfqpoint{2.515690in}{0.998385in}}%
\pgfpathlineto{\pgfqpoint{2.528379in}{0.999297in}}%
\pgfpathlineto{\pgfqpoint{2.528607in}{0.995823in}}%
\pgfpathlineto{\pgfqpoint{2.553769in}{0.997504in}}%
\pgfpathlineto{\pgfqpoint{2.553484in}{0.986177in}}%
\pgfpathlineto{\pgfqpoint{2.551959in}{0.940228in}}%
\pgfpathclose%
\pgfusepath{fill}%
\end{pgfscope}%
\begin{pgfscope}%
\pgfpathrectangle{\pgfqpoint{0.100000in}{0.100000in}}{\pgfqpoint{3.608454in}{2.310000in}}%
\pgfusepath{clip}%
\pgfsetbuttcap%
\pgfsetmiterjoin%
\definecolor{currentfill}{rgb}{0.000000,0.670588,0.664706}%
\pgfsetfillcolor{currentfill}%
\pgfsetlinewidth{0.000000pt}%
\definecolor{currentstroke}{rgb}{0.000000,0.000000,0.000000}%
\pgfsetstrokecolor{currentstroke}%
\pgfsetstrokeopacity{0.000000}%
\pgfsetdash{}{0pt}%
\pgfpathmoveto{\pgfqpoint{3.458788in}{2.127000in}}%
\pgfpathlineto{\pgfqpoint{3.457596in}{2.136892in}}%
\pgfpathlineto{\pgfqpoint{3.481080in}{2.206952in}}%
\pgfpathlineto{\pgfqpoint{3.491936in}{2.206386in}}%
\pgfpathlineto{\pgfqpoint{3.494641in}{2.194072in}}%
\pgfpathlineto{\pgfqpoint{3.504213in}{2.190328in}}%
\pgfpathlineto{\pgfqpoint{3.518463in}{2.203609in}}%
\pgfpathlineto{\pgfqpoint{3.528567in}{2.206604in}}%
\pgfpathlineto{\pgfqpoint{3.527922in}{2.211919in}}%
\pgfpathlineto{\pgfqpoint{3.535054in}{2.214135in}}%
\pgfpathlineto{\pgfqpoint{3.552913in}{2.206672in}}%
\pgfpathlineto{\pgfqpoint{3.558787in}{2.200716in}}%
\pgfpathlineto{\pgfqpoint{3.564560in}{2.199223in}}%
\pgfpathlineto{\pgfqpoint{3.575635in}{2.163717in}}%
\pgfpathlineto{\pgfqpoint{3.591481in}{2.113768in}}%
\pgfpathlineto{\pgfqpoint{3.591685in}{2.108859in}}%
\pgfpathlineto{\pgfqpoint{3.596315in}{2.093414in}}%
\pgfpathlineto{\pgfqpoint{3.584395in}{2.086358in}}%
\pgfpathlineto{\pgfqpoint{3.566650in}{2.075414in}}%
\pgfpathlineto{\pgfqpoint{3.565274in}{2.075679in}}%
\pgfpathlineto{\pgfqpoint{3.557203in}{2.103454in}}%
\pgfpathlineto{\pgfqpoint{3.546699in}{2.136666in}}%
\pgfpathlineto{\pgfqpoint{3.526210in}{2.131697in}}%
\pgfpathlineto{\pgfqpoint{3.522068in}{2.145170in}}%
\pgfpathlineto{\pgfqpoint{3.474712in}{2.131510in}}%
\pgfpathlineto{\pgfqpoint{3.458788in}{2.127000in}}%
\pgfpathclose%
\pgfusepath{fill}%
\end{pgfscope}%
\begin{pgfscope}%
\pgfpathrectangle{\pgfqpoint{0.100000in}{0.100000in}}{\pgfqpoint{3.608454in}{2.310000in}}%
\pgfusepath{clip}%
\pgfsetbuttcap%
\pgfsetmiterjoin%
\definecolor{currentfill}{rgb}{0.000000,0.482353,0.758824}%
\pgfsetfillcolor{currentfill}%
\pgfsetlinewidth{0.000000pt}%
\definecolor{currentstroke}{rgb}{0.000000,0.000000,0.000000}%
\pgfsetstrokecolor{currentstroke}%
\pgfsetstrokeopacity{0.000000}%
\pgfsetdash{}{0pt}%
\pgfpathmoveto{\pgfqpoint{2.714189in}{1.531011in}}%
\pgfpathlineto{\pgfqpoint{2.707766in}{1.586556in}}%
\pgfpathlineto{\pgfqpoint{2.731544in}{1.590059in}}%
\pgfpathlineto{\pgfqpoint{2.734410in}{1.575123in}}%
\pgfpathlineto{\pgfqpoint{2.737491in}{1.568598in}}%
\pgfpathlineto{\pgfqpoint{2.744178in}{1.569410in}}%
\pgfpathlineto{\pgfqpoint{2.746637in}{1.548837in}}%
\pgfpathlineto{\pgfqpoint{2.739953in}{1.548013in}}%
\pgfpathlineto{\pgfqpoint{2.741540in}{1.534284in}}%
\pgfpathlineto{\pgfqpoint{2.714189in}{1.531011in}}%
\pgfpathclose%
\pgfusepath{fill}%
\end{pgfscope}%
\begin{pgfscope}%
\pgfpathrectangle{\pgfqpoint{0.100000in}{0.100000in}}{\pgfqpoint{3.608454in}{2.310000in}}%
\pgfusepath{clip}%
\pgfsetbuttcap%
\pgfsetmiterjoin%
\definecolor{currentfill}{rgb}{0.000000,0.650980,0.674510}%
\pgfsetfillcolor{currentfill}%
\pgfsetlinewidth{0.000000pt}%
\definecolor{currentstroke}{rgb}{0.000000,0.000000,0.000000}%
\pgfsetstrokecolor{currentstroke}%
\pgfsetstrokeopacity{0.000000}%
\pgfsetdash{}{0pt}%
\pgfpathmoveto{\pgfqpoint{1.303419in}{1.175695in}}%
\pgfpathlineto{\pgfqpoint{1.300870in}{1.158512in}}%
\pgfpathlineto{\pgfqpoint{1.210980in}{1.171960in}}%
\pgfpathlineto{\pgfqpoint{1.223167in}{1.249876in}}%
\pgfpathlineto{\pgfqpoint{1.264762in}{1.243586in}}%
\pgfpathlineto{\pgfqpoint{1.268205in}{1.247503in}}%
\pgfpathlineto{\pgfqpoint{1.272888in}{1.260292in}}%
\pgfpathlineto{\pgfqpoint{1.280183in}{1.269745in}}%
\pgfpathlineto{\pgfqpoint{1.286167in}{1.271003in}}%
\pgfpathlineto{\pgfqpoint{1.291651in}{1.276850in}}%
\pgfpathlineto{\pgfqpoint{1.293840in}{1.286121in}}%
\pgfpathlineto{\pgfqpoint{1.303897in}{1.295339in}}%
\pgfpathlineto{\pgfqpoint{1.305818in}{1.299625in}}%
\pgfpathlineto{\pgfqpoint{1.314836in}{1.308720in}}%
\pgfpathlineto{\pgfqpoint{1.325950in}{1.311919in}}%
\pgfpathlineto{\pgfqpoint{1.327961in}{1.309954in}}%
\pgfpathlineto{\pgfqpoint{1.329166in}{1.295757in}}%
\pgfpathlineto{\pgfqpoint{1.325433in}{1.268700in}}%
\pgfpathlineto{\pgfqpoint{1.347428in}{1.265725in}}%
\pgfpathlineto{\pgfqpoint{1.347103in}{1.263333in}}%
\pgfpathlineto{\pgfqpoint{1.375196in}{1.259964in}}%
\pgfpathlineto{\pgfqpoint{1.373495in}{1.246788in}}%
\pgfpathlineto{\pgfqpoint{1.383780in}{1.226748in}}%
\pgfpathlineto{\pgfqpoint{1.324687in}{1.235061in}}%
\pgfpathlineto{\pgfqpoint{1.322223in}{1.228643in}}%
\pgfpathlineto{\pgfqpoint{1.310197in}{1.221478in}}%
\pgfpathlineto{\pgfqpoint{1.303419in}{1.175695in}}%
\pgfpathclose%
\pgfusepath{fill}%
\end{pgfscope}%
\begin{pgfscope}%
\pgfpathrectangle{\pgfqpoint{0.100000in}{0.100000in}}{\pgfqpoint{3.608454in}{2.310000in}}%
\pgfusepath{clip}%
\pgfsetbuttcap%
\pgfsetmiterjoin%
\definecolor{currentfill}{rgb}{0.000000,0.760784,0.619608}%
\pgfsetfillcolor{currentfill}%
\pgfsetlinewidth{0.000000pt}%
\definecolor{currentstroke}{rgb}{0.000000,0.000000,0.000000}%
\pgfsetstrokecolor{currentstroke}%
\pgfsetstrokeopacity{0.000000}%
\pgfsetdash{}{0pt}%
\pgfpathmoveto{\pgfqpoint{2.420115in}{2.049767in}}%
\pgfpathlineto{\pgfqpoint{2.414517in}{2.052884in}}%
\pgfpathlineto{\pgfqpoint{2.417513in}{2.057828in}}%
\pgfpathlineto{\pgfqpoint{2.438907in}{2.070311in}}%
\pgfpathlineto{\pgfqpoint{2.446529in}{2.077612in}}%
\pgfpathlineto{\pgfqpoint{2.455726in}{2.081379in}}%
\pgfpathlineto{\pgfqpoint{2.450327in}{2.070558in}}%
\pgfpathlineto{\pgfqpoint{2.444429in}{2.065660in}}%
\pgfpathlineto{\pgfqpoint{2.430864in}{2.058882in}}%
\pgfpathlineto{\pgfqpoint{2.433030in}{2.055826in}}%
\pgfpathlineto{\pgfqpoint{2.420115in}{2.049767in}}%
\pgfpathclose%
\pgfusepath{fill}%
\end{pgfscope}%
\begin{pgfscope}%
\pgfpathrectangle{\pgfqpoint{0.100000in}{0.100000in}}{\pgfqpoint{3.608454in}{2.310000in}}%
\pgfusepath{clip}%
\pgfsetbuttcap%
\pgfsetmiterjoin%
\definecolor{currentfill}{rgb}{0.000000,0.760784,0.619608}%
\pgfsetfillcolor{currentfill}%
\pgfsetlinewidth{0.000000pt}%
\definecolor{currentstroke}{rgb}{0.000000,0.000000,0.000000}%
\pgfsetstrokecolor{currentstroke}%
\pgfsetstrokeopacity{0.000000}%
\pgfsetdash{}{0pt}%
\pgfpathmoveto{\pgfqpoint{2.485656in}{1.981514in}}%
\pgfpathlineto{\pgfqpoint{2.486521in}{1.969842in}}%
\pgfpathlineto{\pgfqpoint{2.489953in}{1.963218in}}%
\pgfpathlineto{\pgfqpoint{2.483177in}{1.962621in}}%
\pgfpathlineto{\pgfqpoint{2.484672in}{1.942048in}}%
\pgfpathlineto{\pgfqpoint{2.436761in}{1.938649in}}%
\pgfpathlineto{\pgfqpoint{2.435480in}{1.959378in}}%
\pgfpathlineto{\pgfqpoint{2.442299in}{1.959772in}}%
\pgfpathlineto{\pgfqpoint{2.438144in}{1.966329in}}%
\pgfpathlineto{\pgfqpoint{2.436662in}{1.987686in}}%
\pgfpathlineto{\pgfqpoint{2.438695in}{1.993269in}}%
\pgfpathlineto{\pgfqpoint{2.450631in}{2.004021in}}%
\pgfpathlineto{\pgfqpoint{2.454724in}{2.005602in}}%
\pgfpathlineto{\pgfqpoint{2.462516in}{2.016619in}}%
\pgfpathlineto{\pgfqpoint{2.472955in}{2.023672in}}%
\pgfpathlineto{\pgfqpoint{2.487300in}{2.027022in}}%
\pgfpathlineto{\pgfqpoint{2.495202in}{2.026998in}}%
\pgfpathlineto{\pgfqpoint{2.500592in}{2.021625in}}%
\pgfpathlineto{\pgfqpoint{2.489998in}{2.020382in}}%
\pgfpathlineto{\pgfqpoint{2.488577in}{2.015517in}}%
\pgfpathlineto{\pgfqpoint{2.482568in}{2.011992in}}%
\pgfpathlineto{\pgfqpoint{2.473948in}{2.003663in}}%
\pgfpathlineto{\pgfqpoint{2.473922in}{1.999171in}}%
\pgfpathlineto{\pgfqpoint{2.463650in}{1.984401in}}%
\pgfpathlineto{\pgfqpoint{2.462316in}{1.973885in}}%
\pgfpathlineto{\pgfqpoint{2.465262in}{1.970695in}}%
\pgfpathlineto{\pgfqpoint{2.474637in}{1.982064in}}%
\pgfpathlineto{\pgfqpoint{2.479239in}{1.985245in}}%
\pgfpathlineto{\pgfqpoint{2.485656in}{1.981514in}}%
\pgfpathclose%
\pgfusepath{fill}%
\end{pgfscope}%
\begin{pgfscope}%
\pgfpathrectangle{\pgfqpoint{0.100000in}{0.100000in}}{\pgfqpoint{3.608454in}{2.310000in}}%
\pgfusepath{clip}%
\pgfsetbuttcap%
\pgfsetmiterjoin%
\definecolor{currentfill}{rgb}{0.000000,0.372549,0.813725}%
\pgfsetfillcolor{currentfill}%
\pgfsetlinewidth{0.000000pt}%
\definecolor{currentstroke}{rgb}{0.000000,0.000000,0.000000}%
\pgfsetstrokecolor{currentstroke}%
\pgfsetstrokeopacity{0.000000}%
\pgfsetdash{}{0pt}%
\pgfpathmoveto{\pgfqpoint{1.481289in}{1.288540in}}%
\pgfpathlineto{\pgfqpoint{1.474843in}{1.297491in}}%
\pgfpathlineto{\pgfqpoint{1.470019in}{1.293572in}}%
\pgfpathlineto{\pgfqpoint{1.458524in}{1.289130in}}%
\pgfpathlineto{\pgfqpoint{1.454881in}{1.289984in}}%
\pgfpathlineto{\pgfqpoint{1.448675in}{1.296375in}}%
\pgfpathlineto{\pgfqpoint{1.444613in}{1.311073in}}%
\pgfpathlineto{\pgfqpoint{1.426840in}{1.337144in}}%
\pgfpathlineto{\pgfqpoint{1.435665in}{1.345243in}}%
\pgfpathlineto{\pgfqpoint{1.436154in}{1.349875in}}%
\pgfpathlineto{\pgfqpoint{1.431712in}{1.356218in}}%
\pgfpathlineto{\pgfqpoint{1.470992in}{1.351885in}}%
\pgfpathlineto{\pgfqpoint{1.474789in}{1.385857in}}%
\pgfpathlineto{\pgfqpoint{1.552744in}{1.377402in}}%
\pgfpathlineto{\pgfqpoint{1.550467in}{1.357020in}}%
\pgfpathlineto{\pgfqpoint{1.547834in}{1.329763in}}%
\pgfpathlineto{\pgfqpoint{1.493221in}{1.335267in}}%
\pgfpathlineto{\pgfqpoint{1.491038in}{1.314728in}}%
\pgfpathlineto{\pgfqpoint{1.484320in}{1.315446in}}%
\pgfpathlineto{\pgfqpoint{1.481289in}{1.288540in}}%
\pgfpathclose%
\pgfusepath{fill}%
\end{pgfscope}%
\begin{pgfscope}%
\pgfpathrectangle{\pgfqpoint{0.100000in}{0.100000in}}{\pgfqpoint{3.608454in}{2.310000in}}%
\pgfusepath{clip}%
\pgfsetbuttcap%
\pgfsetmiterjoin%
\definecolor{currentfill}{rgb}{0.000000,0.717647,0.641176}%
\pgfsetfillcolor{currentfill}%
\pgfsetlinewidth{0.000000pt}%
\definecolor{currentstroke}{rgb}{0.000000,0.000000,0.000000}%
\pgfsetstrokecolor{currentstroke}%
\pgfsetstrokeopacity{0.000000}%
\pgfsetdash{}{0pt}%
\pgfpathmoveto{\pgfqpoint{2.977955in}{1.500863in}}%
\pgfpathlineto{\pgfqpoint{2.956251in}{1.497406in}}%
\pgfpathlineto{\pgfqpoint{2.954062in}{1.518144in}}%
\pgfpathlineto{\pgfqpoint{2.949522in}{1.517575in}}%
\pgfpathlineto{\pgfqpoint{2.949140in}{1.528048in}}%
\pgfpathlineto{\pgfqpoint{2.952109in}{1.531961in}}%
\pgfpathlineto{\pgfqpoint{2.963825in}{1.532458in}}%
\pgfpathlineto{\pgfqpoint{2.971951in}{1.538265in}}%
\pgfpathlineto{\pgfqpoint{2.977955in}{1.500863in}}%
\pgfpathclose%
\pgfusepath{fill}%
\end{pgfscope}%
\begin{pgfscope}%
\pgfpathrectangle{\pgfqpoint{0.100000in}{0.100000in}}{\pgfqpoint{3.608454in}{2.310000in}}%
\pgfusepath{clip}%
\pgfsetbuttcap%
\pgfsetmiterjoin%
\definecolor{currentfill}{rgb}{0.000000,0.760784,0.619608}%
\pgfsetfillcolor{currentfill}%
\pgfsetlinewidth{0.000000pt}%
\definecolor{currentstroke}{rgb}{0.000000,0.000000,0.000000}%
\pgfsetstrokecolor{currentstroke}%
\pgfsetstrokeopacity{0.000000}%
\pgfsetdash{}{0pt}%
\pgfpathmoveto{\pgfqpoint{2.573691in}{1.108364in}}%
\pgfpathlineto{\pgfqpoint{2.573068in}{1.094430in}}%
\pgfpathlineto{\pgfqpoint{2.579650in}{1.089158in}}%
\pgfpathlineto{\pgfqpoint{2.575892in}{1.080309in}}%
\pgfpathlineto{\pgfqpoint{2.559562in}{1.077181in}}%
\pgfpathlineto{\pgfqpoint{2.556697in}{1.071568in}}%
\pgfpathlineto{\pgfqpoint{2.550818in}{1.070538in}}%
\pgfpathlineto{\pgfqpoint{2.534914in}{1.072050in}}%
\pgfpathlineto{\pgfqpoint{2.531082in}{1.076803in}}%
\pgfpathlineto{\pgfqpoint{2.525076in}{1.078728in}}%
\pgfpathlineto{\pgfqpoint{2.517692in}{1.084251in}}%
\pgfpathlineto{\pgfqpoint{2.517006in}{1.100169in}}%
\pgfpathlineto{\pgfqpoint{2.519200in}{1.103140in}}%
\pgfpathlineto{\pgfqpoint{2.544235in}{1.104329in}}%
\pgfpathlineto{\pgfqpoint{2.552190in}{1.107568in}}%
\pgfpathlineto{\pgfqpoint{2.553934in}{1.104609in}}%
\pgfpathlineto{\pgfqpoint{2.561312in}{1.107350in}}%
\pgfpathlineto{\pgfqpoint{2.573691in}{1.108364in}}%
\pgfpathclose%
\pgfusepath{fill}%
\end{pgfscope}%
\begin{pgfscope}%
\pgfpathrectangle{\pgfqpoint{0.100000in}{0.100000in}}{\pgfqpoint{3.608454in}{2.310000in}}%
\pgfusepath{clip}%
\pgfsetbuttcap%
\pgfsetmiterjoin%
\definecolor{currentfill}{rgb}{0.000000,0.662745,0.668627}%
\pgfsetfillcolor{currentfill}%
\pgfsetlinewidth{0.000000pt}%
\definecolor{currentstroke}{rgb}{0.000000,0.000000,0.000000}%
\pgfsetstrokecolor{currentstroke}%
\pgfsetstrokeopacity{0.000000}%
\pgfsetdash{}{0pt}%
\pgfpathmoveto{\pgfqpoint{1.518670in}{0.813225in}}%
\pgfpathlineto{\pgfqpoint{1.498395in}{0.815195in}}%
\pgfpathlineto{\pgfqpoint{1.443108in}{0.820998in}}%
\pgfpathlineto{\pgfqpoint{1.447334in}{0.862167in}}%
\pgfpathlineto{\pgfqpoint{1.413830in}{0.865768in}}%
\pgfpathlineto{\pgfqpoint{1.417927in}{0.900770in}}%
\pgfpathlineto{\pgfqpoint{1.420339in}{0.900498in}}%
\pgfpathlineto{\pgfqpoint{1.421822in}{0.913943in}}%
\pgfpathlineto{\pgfqpoint{1.449063in}{0.911411in}}%
\pgfpathlineto{\pgfqpoint{1.452184in}{0.931773in}}%
\pgfpathlineto{\pgfqpoint{1.458090in}{0.986273in}}%
\pgfpathlineto{\pgfqpoint{1.465068in}{0.985473in}}%
\pgfpathlineto{\pgfqpoint{1.464338in}{0.978551in}}%
\pgfpathlineto{\pgfqpoint{1.505497in}{0.973876in}}%
\pgfpathlineto{\pgfqpoint{1.506226in}{0.980812in}}%
\pgfpathlineto{\pgfqpoint{1.519945in}{0.979398in}}%
\pgfpathlineto{\pgfqpoint{1.526759in}{0.978642in}}%
\pgfpathlineto{\pgfqpoint{1.524717in}{0.957918in}}%
\pgfpathlineto{\pgfqpoint{1.532929in}{0.957126in}}%
\pgfpathlineto{\pgfqpoint{1.531592in}{0.944243in}}%
\pgfpathlineto{\pgfqpoint{1.545317in}{0.943021in}}%
\pgfpathlineto{\pgfqpoint{1.544636in}{0.936153in}}%
\pgfpathlineto{\pgfqpoint{1.530771in}{0.937426in}}%
\pgfpathlineto{\pgfqpoint{1.529406in}{0.923693in}}%
\pgfpathlineto{\pgfqpoint{1.526429in}{0.923948in}}%
\pgfpathlineto{\pgfqpoint{1.523143in}{0.889921in}}%
\pgfpathlineto{\pgfqpoint{1.519940in}{0.890229in}}%
\pgfpathlineto{\pgfqpoint{1.516573in}{0.855187in}}%
\pgfpathlineto{\pgfqpoint{1.522650in}{0.854609in}}%
\pgfpathlineto{\pgfqpoint{1.518670in}{0.813225in}}%
\pgfpathclose%
\pgfusepath{fill}%
\end{pgfscope}%
\begin{pgfscope}%
\pgfpathrectangle{\pgfqpoint{0.100000in}{0.100000in}}{\pgfqpoint{3.608454in}{2.310000in}}%
\pgfusepath{clip}%
\pgfsetbuttcap%
\pgfsetmiterjoin%
\definecolor{currentfill}{rgb}{0.000000,0.654902,0.672549}%
\pgfsetfillcolor{currentfill}%
\pgfsetlinewidth{0.000000pt}%
\definecolor{currentstroke}{rgb}{0.000000,0.000000,0.000000}%
\pgfsetstrokecolor{currentstroke}%
\pgfsetstrokeopacity{0.000000}%
\pgfsetdash{}{0pt}%
\pgfpathmoveto{\pgfqpoint{1.411936in}{2.184794in}}%
\pgfpathlineto{\pgfqpoint{1.407741in}{2.164527in}}%
\pgfpathlineto{\pgfqpoint{1.404613in}{2.164996in}}%
\pgfpathlineto{\pgfqpoint{1.402592in}{2.151198in}}%
\pgfpathlineto{\pgfqpoint{1.399456in}{2.139818in}}%
\pgfpathlineto{\pgfqpoint{1.396285in}{2.142117in}}%
\pgfpathlineto{\pgfqpoint{1.392437in}{2.124428in}}%
\pgfpathlineto{\pgfqpoint{1.389566in}{2.105191in}}%
\pgfpathlineto{\pgfqpoint{1.380515in}{2.107723in}}%
\pgfpathlineto{\pgfqpoint{1.378739in}{2.102305in}}%
\pgfpathlineto{\pgfqpoint{1.364847in}{2.104580in}}%
\pgfpathlineto{\pgfqpoint{1.362402in}{2.089710in}}%
\pgfpathlineto{\pgfqpoint{1.358594in}{2.095823in}}%
\pgfpathlineto{\pgfqpoint{1.350200in}{2.095416in}}%
\pgfpathlineto{\pgfqpoint{1.342674in}{2.097531in}}%
\pgfpathlineto{\pgfqpoint{1.332093in}{2.093598in}}%
\pgfpathlineto{\pgfqpoint{1.328054in}{2.095490in}}%
\pgfpathlineto{\pgfqpoint{1.333325in}{2.126759in}}%
\pgfpathlineto{\pgfqpoint{1.323124in}{2.128216in}}%
\pgfpathlineto{\pgfqpoint{1.324278in}{2.135234in}}%
\pgfpathlineto{\pgfqpoint{1.318309in}{2.136230in}}%
\pgfpathlineto{\pgfqpoint{1.319123in}{2.142944in}}%
\pgfpathlineto{\pgfqpoint{1.278167in}{2.149917in}}%
\pgfpathlineto{\pgfqpoint{1.276965in}{2.143064in}}%
\pgfpathlineto{\pgfqpoint{1.266372in}{2.144919in}}%
\pgfpathlineto{\pgfqpoint{1.265169in}{2.138097in}}%
\pgfpathlineto{\pgfqpoint{1.234540in}{2.143503in}}%
\pgfpathlineto{\pgfqpoint{1.235802in}{2.150383in}}%
\pgfpathlineto{\pgfqpoint{1.242847in}{2.149099in}}%
\pgfpathlineto{\pgfqpoint{1.254169in}{2.210455in}}%
\pgfpathlineto{\pgfqpoint{1.308831in}{2.201031in}}%
\pgfpathlineto{\pgfqpoint{1.372332in}{2.190734in}}%
\pgfpathlineto{\pgfqpoint{1.411936in}{2.184794in}}%
\pgfpathclose%
\pgfusepath{fill}%
\end{pgfscope}%
\begin{pgfscope}%
\pgfpathrectangle{\pgfqpoint{0.100000in}{0.100000in}}{\pgfqpoint{3.608454in}{2.310000in}}%
\pgfusepath{clip}%
\pgfsetbuttcap%
\pgfsetmiterjoin%
\definecolor{currentfill}{rgb}{0.000000,0.541176,0.729412}%
\pgfsetfillcolor{currentfill}%
\pgfsetlinewidth{0.000000pt}%
\definecolor{currentstroke}{rgb}{0.000000,0.000000,0.000000}%
\pgfsetstrokecolor{currentstroke}%
\pgfsetstrokeopacity{0.000000}%
\pgfsetdash{}{0pt}%
\pgfpathmoveto{\pgfqpoint{3.286885in}{1.169498in}}%
\pgfpathlineto{\pgfqpoint{3.287098in}{1.161529in}}%
\pgfpathlineto{\pgfqpoint{3.273900in}{1.157290in}}%
\pgfpathlineto{\pgfqpoint{3.265343in}{1.162433in}}%
\pgfpathlineto{\pgfqpoint{3.248882in}{1.168666in}}%
\pgfpathlineto{\pgfqpoint{3.248243in}{1.176631in}}%
\pgfpathlineto{\pgfqpoint{3.250804in}{1.179929in}}%
\pgfpathlineto{\pgfqpoint{3.246587in}{1.185703in}}%
\pgfpathlineto{\pgfqpoint{3.244554in}{1.193790in}}%
\pgfpathlineto{\pgfqpoint{3.232624in}{1.197760in}}%
\pgfpathlineto{\pgfqpoint{3.231820in}{1.204552in}}%
\pgfpathlineto{\pgfqpoint{3.226514in}{1.211542in}}%
\pgfpathlineto{\pgfqpoint{3.230833in}{1.216563in}}%
\pgfpathlineto{\pgfqpoint{3.238791in}{1.214139in}}%
\pgfpathlineto{\pgfqpoint{3.239989in}{1.209565in}}%
\pgfpathlineto{\pgfqpoint{3.247206in}{1.209157in}}%
\pgfpathlineto{\pgfqpoint{3.253779in}{1.205127in}}%
\pgfpathlineto{\pgfqpoint{3.258356in}{1.209245in}}%
\pgfpathlineto{\pgfqpoint{3.262155in}{1.203038in}}%
\pgfpathlineto{\pgfqpoint{3.265874in}{1.209619in}}%
\pgfpathlineto{\pgfqpoint{3.269662in}{1.209194in}}%
\pgfpathlineto{\pgfqpoint{3.272014in}{1.215831in}}%
\pgfpathlineto{\pgfqpoint{3.282173in}{1.217740in}}%
\pgfpathlineto{\pgfqpoint{3.289627in}{1.222451in}}%
\pgfpathlineto{\pgfqpoint{3.293663in}{1.220232in}}%
\pgfpathlineto{\pgfqpoint{3.301128in}{1.225302in}}%
\pgfpathlineto{\pgfqpoint{3.310778in}{1.227788in}}%
\pgfpathlineto{\pgfqpoint{3.318505in}{1.221809in}}%
\pgfpathlineto{\pgfqpoint{3.321566in}{1.228593in}}%
\pgfpathlineto{\pgfqpoint{3.327494in}{1.228385in}}%
\pgfpathlineto{\pgfqpoint{3.332202in}{1.223602in}}%
\pgfpathlineto{\pgfqpoint{3.337579in}{1.209759in}}%
\pgfpathlineto{\pgfqpoint{3.336122in}{1.199785in}}%
\pgfpathlineto{\pgfqpoint{3.328969in}{1.197798in}}%
\pgfpathlineto{\pgfqpoint{3.323739in}{1.188050in}}%
\pgfpathlineto{\pgfqpoint{3.323024in}{1.183127in}}%
\pgfpathlineto{\pgfqpoint{3.315978in}{1.174823in}}%
\pgfpathlineto{\pgfqpoint{3.307317in}{1.174597in}}%
\pgfpathlineto{\pgfqpoint{3.297293in}{1.177393in}}%
\pgfpathlineto{\pgfqpoint{3.295040in}{1.173975in}}%
\pgfpathlineto{\pgfqpoint{3.285428in}{1.174379in}}%
\pgfpathlineto{\pgfqpoint{3.286885in}{1.169498in}}%
\pgfpathclose%
\pgfusepath{fill}%
\end{pgfscope}%
\begin{pgfscope}%
\pgfpathrectangle{\pgfqpoint{0.100000in}{0.100000in}}{\pgfqpoint{3.608454in}{2.310000in}}%
\pgfusepath{clip}%
\pgfsetbuttcap%
\pgfsetmiterjoin%
\definecolor{currentfill}{rgb}{0.000000,0.541176,0.729412}%
\pgfsetfillcolor{currentfill}%
\pgfsetlinewidth{0.000000pt}%
\definecolor{currentstroke}{rgb}{0.000000,0.000000,0.000000}%
\pgfsetstrokecolor{currentstroke}%
\pgfsetstrokeopacity{0.000000}%
\pgfsetdash{}{0pt}%
\pgfpathmoveto{\pgfqpoint{3.324933in}{1.250084in}}%
\pgfpathlineto{\pgfqpoint{3.338120in}{1.231991in}}%
\pgfpathlineto{\pgfqpoint{3.331617in}{1.232924in}}%
\pgfpathlineto{\pgfqpoint{3.324107in}{1.249673in}}%
\pgfpathlineto{\pgfqpoint{3.324933in}{1.250084in}}%
\pgfpathclose%
\pgfusepath{fill}%
\end{pgfscope}%
\begin{pgfscope}%
\pgfpathrectangle{\pgfqpoint{0.100000in}{0.100000in}}{\pgfqpoint{3.608454in}{2.310000in}}%
\pgfusepath{clip}%
\pgfsetbuttcap%
\pgfsetmiterjoin%
\definecolor{currentfill}{rgb}{0.000000,0.505882,0.747059}%
\pgfsetfillcolor{currentfill}%
\pgfsetlinewidth{0.000000pt}%
\definecolor{currentstroke}{rgb}{0.000000,0.000000,0.000000}%
\pgfsetstrokecolor{currentstroke}%
\pgfsetstrokeopacity{0.000000}%
\pgfsetdash{}{0pt}%
\pgfpathmoveto{\pgfqpoint{1.923573in}{0.599973in}}%
\pgfpathlineto{\pgfqpoint{1.925990in}{0.598363in}}%
\pgfpathlineto{\pgfqpoint{1.945778in}{0.610158in}}%
\pgfpathlineto{\pgfqpoint{1.957611in}{0.597327in}}%
\pgfpathlineto{\pgfqpoint{1.953286in}{0.593627in}}%
\pgfpathlineto{\pgfqpoint{1.950391in}{0.577980in}}%
\pgfpathlineto{\pgfqpoint{1.923720in}{0.556660in}}%
\pgfpathlineto{\pgfqpoint{1.915913in}{0.565873in}}%
\pgfpathlineto{\pgfqpoint{1.907128in}{0.576661in}}%
\pgfpathlineto{\pgfqpoint{1.923573in}{0.599973in}}%
\pgfpathclose%
\pgfusepath{fill}%
\end{pgfscope}%
\begin{pgfscope}%
\pgfpathrectangle{\pgfqpoint{0.100000in}{0.100000in}}{\pgfqpoint{3.608454in}{2.310000in}}%
\pgfusepath{clip}%
\pgfsetbuttcap%
\pgfsetmiterjoin%
\definecolor{currentfill}{rgb}{0.000000,0.623529,0.688235}%
\pgfsetfillcolor{currentfill}%
\pgfsetlinewidth{0.000000pt}%
\definecolor{currentstroke}{rgb}{0.000000,0.000000,0.000000}%
\pgfsetstrokecolor{currentstroke}%
\pgfsetstrokeopacity{0.000000}%
\pgfsetdash{}{0pt}%
\pgfpathmoveto{\pgfqpoint{1.888636in}{1.700186in}}%
\pgfpathlineto{\pgfqpoint{1.956891in}{1.697804in}}%
\pgfpathlineto{\pgfqpoint{1.968709in}{1.697472in}}%
\pgfpathlineto{\pgfqpoint{1.968146in}{1.671305in}}%
\pgfpathlineto{\pgfqpoint{1.928813in}{1.672403in}}%
\pgfpathlineto{\pgfqpoint{1.926808in}{1.664949in}}%
\pgfpathlineto{\pgfqpoint{1.927556in}{1.657034in}}%
\pgfpathlineto{\pgfqpoint{1.925767in}{1.649254in}}%
\pgfpathlineto{\pgfqpoint{1.922640in}{1.646360in}}%
\pgfpathlineto{\pgfqpoint{1.914725in}{1.650247in}}%
\pgfpathlineto{\pgfqpoint{1.905777in}{1.655731in}}%
\pgfpathlineto{\pgfqpoint{1.904126in}{1.659838in}}%
\pgfpathlineto{\pgfqpoint{1.902600in}{1.663282in}}%
\pgfpathlineto{\pgfqpoint{1.893266in}{1.666452in}}%
\pgfpathlineto{\pgfqpoint{1.888697in}{1.672141in}}%
\pgfpathlineto{\pgfqpoint{1.884096in}{1.673023in}}%
\pgfpathlineto{\pgfqpoint{1.881583in}{1.680508in}}%
\pgfpathlineto{\pgfqpoint{1.872226in}{1.686715in}}%
\pgfpathlineto{\pgfqpoint{1.867137in}{1.695946in}}%
\pgfpathlineto{\pgfqpoint{1.862167in}{1.696367in}}%
\pgfpathlineto{\pgfqpoint{1.859723in}{1.701494in}}%
\pgfpathlineto{\pgfqpoint{1.888636in}{1.700186in}}%
\pgfpathclose%
\pgfusepath{fill}%
\end{pgfscope}%
\begin{pgfscope}%
\pgfpathrectangle{\pgfqpoint{0.100000in}{0.100000in}}{\pgfqpoint{3.608454in}{2.310000in}}%
\pgfusepath{clip}%
\pgfsetbuttcap%
\pgfsetmiterjoin%
\definecolor{currentfill}{rgb}{0.000000,0.486275,0.756863}%
\pgfsetfillcolor{currentfill}%
\pgfsetlinewidth{0.000000pt}%
\definecolor{currentstroke}{rgb}{0.000000,0.000000,0.000000}%
\pgfsetstrokecolor{currentstroke}%
\pgfsetstrokeopacity{0.000000}%
\pgfsetdash{}{0pt}%
\pgfpathmoveto{\pgfqpoint{2.688201in}{1.338715in}}%
\pgfpathlineto{\pgfqpoint{2.674475in}{1.337189in}}%
\pgfpathlineto{\pgfqpoint{2.669050in}{1.334303in}}%
\pgfpathlineto{\pgfqpoint{2.667619in}{1.346789in}}%
\pgfpathlineto{\pgfqpoint{2.662944in}{1.348688in}}%
\pgfpathlineto{\pgfqpoint{2.643662in}{1.346766in}}%
\pgfpathlineto{\pgfqpoint{2.641627in}{1.364614in}}%
\pgfpathlineto{\pgfqpoint{2.638885in}{1.368939in}}%
\pgfpathlineto{\pgfqpoint{2.653405in}{1.370465in}}%
\pgfpathlineto{\pgfqpoint{2.650805in}{1.393553in}}%
\pgfpathlineto{\pgfqpoint{2.675138in}{1.396451in}}%
\pgfpathlineto{\pgfqpoint{2.676781in}{1.379202in}}%
\pgfpathlineto{\pgfqpoint{2.684115in}{1.380130in}}%
\pgfpathlineto{\pgfqpoint{2.691293in}{1.385879in}}%
\pgfpathlineto{\pgfqpoint{2.693342in}{1.363630in}}%
\pgfpathlineto{\pgfqpoint{2.683930in}{1.355595in}}%
\pgfpathlineto{\pgfqpoint{2.679482in}{1.355109in}}%
\pgfpathlineto{\pgfqpoint{2.680182in}{1.348259in}}%
\pgfpathlineto{\pgfqpoint{2.687551in}{1.344405in}}%
\pgfpathlineto{\pgfqpoint{2.688201in}{1.338715in}}%
\pgfpathclose%
\pgfusepath{fill}%
\end{pgfscope}%
\begin{pgfscope}%
\pgfpathrectangle{\pgfqpoint{0.100000in}{0.100000in}}{\pgfqpoint{3.608454in}{2.310000in}}%
\pgfusepath{clip}%
\pgfsetbuttcap%
\pgfsetmiterjoin%
\definecolor{currentfill}{rgb}{0.000000,0.596078,0.701961}%
\pgfsetfillcolor{currentfill}%
\pgfsetlinewidth{0.000000pt}%
\definecolor{currentstroke}{rgb}{0.000000,0.000000,0.000000}%
\pgfsetstrokecolor{currentstroke}%
\pgfsetstrokeopacity{0.000000}%
\pgfsetdash{}{0pt}%
\pgfpathmoveto{\pgfqpoint{2.043578in}{1.181785in}}%
\pgfpathlineto{\pgfqpoint{1.971148in}{1.183063in}}%
\pgfpathlineto{\pgfqpoint{1.971690in}{1.220737in}}%
\pgfpathlineto{\pgfqpoint{2.011063in}{1.220009in}}%
\pgfpathlineto{\pgfqpoint{2.012052in}{1.268224in}}%
\pgfpathlineto{\pgfqpoint{2.022297in}{1.268017in}}%
\pgfpathlineto{\pgfqpoint{2.022450in}{1.274884in}}%
\pgfpathlineto{\pgfqpoint{2.047190in}{1.274443in}}%
\pgfpathlineto{\pgfqpoint{2.047067in}{1.264108in}}%
\pgfpathlineto{\pgfqpoint{2.046716in}{1.239913in}}%
\pgfpathlineto{\pgfqpoint{2.045886in}{1.181780in}}%
\pgfpathlineto{\pgfqpoint{2.043578in}{1.181785in}}%
\pgfpathclose%
\pgfusepath{fill}%
\end{pgfscope}%
\begin{pgfscope}%
\pgfpathrectangle{\pgfqpoint{0.100000in}{0.100000in}}{\pgfqpoint{3.608454in}{2.310000in}}%
\pgfusepath{clip}%
\pgfsetbuttcap%
\pgfsetmiterjoin%
\definecolor{currentfill}{rgb}{0.000000,0.901961,0.549020}%
\pgfsetfillcolor{currentfill}%
\pgfsetlinewidth{0.000000pt}%
\definecolor{currentstroke}{rgb}{0.000000,0.000000,0.000000}%
\pgfsetstrokecolor{currentstroke}%
\pgfsetstrokeopacity{0.000000}%
\pgfsetdash{}{0pt}%
\pgfpathmoveto{\pgfqpoint{2.867721in}{1.199764in}}%
\pgfpathlineto{\pgfqpoint{2.827558in}{1.194906in}}%
\pgfpathlineto{\pgfqpoint{2.828062in}{1.198088in}}%
\pgfpathlineto{\pgfqpoint{2.837048in}{1.201444in}}%
\pgfpathlineto{\pgfqpoint{2.856686in}{1.210206in}}%
\pgfpathlineto{\pgfqpoint{2.859365in}{1.219154in}}%
\pgfpathlineto{\pgfqpoint{2.871264in}{1.223809in}}%
\pgfpathlineto{\pgfqpoint{2.870950in}{1.230065in}}%
\pgfpathlineto{\pgfqpoint{2.878218in}{1.237030in}}%
\pgfpathlineto{\pgfqpoint{2.878670in}{1.243251in}}%
\pgfpathlineto{\pgfqpoint{2.887645in}{1.250541in}}%
\pgfpathlineto{\pgfqpoint{2.902185in}{1.260526in}}%
\pgfpathlineto{\pgfqpoint{2.908022in}{1.254600in}}%
\pgfpathlineto{\pgfqpoint{2.916003in}{1.242873in}}%
\pgfpathlineto{\pgfqpoint{2.904928in}{1.234962in}}%
\pgfpathlineto{\pgfqpoint{2.907489in}{1.230091in}}%
\pgfpathlineto{\pgfqpoint{2.901040in}{1.226866in}}%
\pgfpathlineto{\pgfqpoint{2.886736in}{1.225061in}}%
\pgfpathlineto{\pgfqpoint{2.880082in}{1.220503in}}%
\pgfpathlineto{\pgfqpoint{2.876442in}{1.212071in}}%
\pgfpathlineto{\pgfqpoint{2.869228in}{1.206518in}}%
\pgfpathlineto{\pgfqpoint{2.867721in}{1.199764in}}%
\pgfpathclose%
\pgfusepath{fill}%
\end{pgfscope}%
\begin{pgfscope}%
\pgfpathrectangle{\pgfqpoint{0.100000in}{0.100000in}}{\pgfqpoint{3.608454in}{2.310000in}}%
\pgfusepath{clip}%
\pgfsetbuttcap%
\pgfsetmiterjoin%
\definecolor{currentfill}{rgb}{0.000000,0.505882,0.747059}%
\pgfsetfillcolor{currentfill}%
\pgfsetlinewidth{0.000000pt}%
\definecolor{currentstroke}{rgb}{0.000000,0.000000,0.000000}%
\pgfsetstrokecolor{currentstroke}%
\pgfsetstrokeopacity{0.000000}%
\pgfsetdash{}{0pt}%
\pgfpathmoveto{\pgfqpoint{1.683362in}{1.122232in}}%
\pgfpathlineto{\pgfqpoint{1.680896in}{1.087866in}}%
\pgfpathlineto{\pgfqpoint{1.646284in}{1.091077in}}%
\pgfpathlineto{\pgfqpoint{1.589893in}{1.095406in}}%
\pgfpathlineto{\pgfqpoint{1.590719in}{1.104616in}}%
\pgfpathlineto{\pgfqpoint{1.596009in}{1.164627in}}%
\pgfpathlineto{\pgfqpoint{1.598511in}{1.164402in}}%
\pgfpathlineto{\pgfqpoint{1.651744in}{1.159945in}}%
\pgfpathlineto{\pgfqpoint{1.648937in}{1.124836in}}%
\pgfpathlineto{\pgfqpoint{1.683362in}{1.122232in}}%
\pgfpathclose%
\pgfusepath{fill}%
\end{pgfscope}%
\begin{pgfscope}%
\pgfpathrectangle{\pgfqpoint{0.100000in}{0.100000in}}{\pgfqpoint{3.608454in}{2.310000in}}%
\pgfusepath{clip}%
\pgfsetbuttcap%
\pgfsetmiterjoin%
\definecolor{currentfill}{rgb}{0.000000,0.447059,0.776471}%
\pgfsetfillcolor{currentfill}%
\pgfsetlinewidth{0.000000pt}%
\definecolor{currentstroke}{rgb}{0.000000,0.000000,0.000000}%
\pgfsetstrokecolor{currentstroke}%
\pgfsetstrokeopacity{0.000000}%
\pgfsetdash{}{0pt}%
\pgfpathmoveto{\pgfqpoint{1.907060in}{0.482191in}}%
\pgfpathlineto{\pgfqpoint{1.900788in}{0.477356in}}%
\pgfpathlineto{\pgfqpoint{1.902126in}{0.472699in}}%
\pgfpathlineto{\pgfqpoint{1.836840in}{0.475148in}}%
\pgfpathlineto{\pgfqpoint{1.839000in}{0.522528in}}%
\pgfpathlineto{\pgfqpoint{1.840202in}{0.558059in}}%
\pgfpathlineto{\pgfqpoint{1.797569in}{0.560061in}}%
\pgfpathlineto{\pgfqpoint{1.799798in}{0.602976in}}%
\pgfpathlineto{\pgfqpoint{1.786496in}{0.603609in}}%
\pgfpathlineto{\pgfqpoint{1.787654in}{0.626021in}}%
\pgfpathlineto{\pgfqpoint{1.817099in}{0.623588in}}%
\pgfpathlineto{\pgfqpoint{1.834695in}{0.613722in}}%
\pgfpathlineto{\pgfqpoint{1.835751in}{0.642247in}}%
\pgfpathlineto{\pgfqpoint{1.858713in}{0.641352in}}%
\pgfpathlineto{\pgfqpoint{1.870125in}{0.624788in}}%
\pgfpathlineto{\pgfqpoint{1.878483in}{0.632521in}}%
\pgfpathlineto{\pgfqpoint{1.896461in}{0.616726in}}%
\pgfpathlineto{\pgfqpoint{1.898375in}{0.608966in}}%
\pgfpathlineto{\pgfqpoint{1.907278in}{0.617155in}}%
\pgfpathlineto{\pgfqpoint{1.918807in}{0.601623in}}%
\pgfpathlineto{\pgfqpoint{1.923573in}{0.599973in}}%
\pgfpathlineto{\pgfqpoint{1.907128in}{0.576661in}}%
\pgfpathlineto{\pgfqpoint{1.915913in}{0.565873in}}%
\pgfpathlineto{\pgfqpoint{1.882610in}{0.539717in}}%
\pgfpathlineto{\pgfqpoint{1.895145in}{0.523860in}}%
\pgfpathlineto{\pgfqpoint{1.889121in}{0.521871in}}%
\pgfpathlineto{\pgfqpoint{1.907719in}{0.482706in}}%
\pgfpathlineto{\pgfqpoint{1.907060in}{0.482191in}}%
\pgfpathclose%
\pgfusepath{fill}%
\end{pgfscope}%
\begin{pgfscope}%
\pgfpathrectangle{\pgfqpoint{0.100000in}{0.100000in}}{\pgfqpoint{3.608454in}{2.310000in}}%
\pgfusepath{clip}%
\pgfsetbuttcap%
\pgfsetmiterjoin%
\definecolor{currentfill}{rgb}{0.000000,0.572549,0.713725}%
\pgfsetfillcolor{currentfill}%
\pgfsetlinewidth{0.000000pt}%
\definecolor{currentstroke}{rgb}{0.000000,0.000000,0.000000}%
\pgfsetstrokecolor{currentstroke}%
\pgfsetstrokeopacity{0.000000}%
\pgfsetdash{}{0pt}%
\pgfpathmoveto{\pgfqpoint{1.831143in}{0.977031in}}%
\pgfpathlineto{\pgfqpoint{1.820478in}{0.987070in}}%
\pgfpathlineto{\pgfqpoint{1.818337in}{0.980492in}}%
\pgfpathlineto{\pgfqpoint{1.811536in}{0.984126in}}%
\pgfpathlineto{\pgfqpoint{1.805846in}{0.981264in}}%
\pgfpathlineto{\pgfqpoint{1.799655in}{0.981871in}}%
\pgfpathlineto{\pgfqpoint{1.785322in}{0.998334in}}%
\pgfpathlineto{\pgfqpoint{1.780677in}{0.997355in}}%
\pgfpathlineto{\pgfqpoint{1.781300in}{1.012079in}}%
\pgfpathlineto{\pgfqpoint{1.782538in}{1.034546in}}%
\pgfpathlineto{\pgfqpoint{1.796623in}{1.033765in}}%
\pgfpathlineto{\pgfqpoint{1.796992in}{1.040622in}}%
\pgfpathlineto{\pgfqpoint{1.863914in}{1.037452in}}%
\pgfpathlineto{\pgfqpoint{1.872137in}{1.035604in}}%
\pgfpathlineto{\pgfqpoint{1.871452in}{1.016408in}}%
\pgfpathlineto{\pgfqpoint{1.858037in}{1.016994in}}%
\pgfpathlineto{\pgfqpoint{1.857172in}{0.996308in}}%
\pgfpathlineto{\pgfqpoint{1.835171in}{0.996785in}}%
\pgfpathlineto{\pgfqpoint{1.831022in}{0.989497in}}%
\pgfpathlineto{\pgfqpoint{1.831143in}{0.977031in}}%
\pgfpathclose%
\pgfusepath{fill}%
\end{pgfscope}%
\begin{pgfscope}%
\pgfpathrectangle{\pgfqpoint{0.100000in}{0.100000in}}{\pgfqpoint{3.608454in}{2.310000in}}%
\pgfusepath{clip}%
\pgfsetbuttcap%
\pgfsetmiterjoin%
\definecolor{currentfill}{rgb}{0.000000,0.349020,0.825490}%
\pgfsetfillcolor{currentfill}%
\pgfsetlinewidth{0.000000pt}%
\definecolor{currentstroke}{rgb}{0.000000,0.000000,0.000000}%
\pgfsetstrokecolor{currentstroke}%
\pgfsetstrokeopacity{0.000000}%
\pgfsetdash{}{0pt}%
\pgfpathmoveto{\pgfqpoint{2.023407in}{1.319659in}}%
\pgfpathlineto{\pgfqpoint{2.023077in}{1.302462in}}%
\pgfpathlineto{\pgfqpoint{1.994270in}{1.303029in}}%
\pgfpathlineto{\pgfqpoint{1.987451in}{1.303171in}}%
\pgfpathlineto{\pgfqpoint{1.987581in}{1.310037in}}%
\pgfpathlineto{\pgfqpoint{1.960306in}{1.310682in}}%
\pgfpathlineto{\pgfqpoint{1.961886in}{1.365738in}}%
\pgfpathlineto{\pgfqpoint{1.962491in}{1.386340in}}%
\pgfpathlineto{\pgfqpoint{2.031304in}{1.384920in}}%
\pgfpathlineto{\pgfqpoint{2.043672in}{1.384754in}}%
\pgfpathlineto{\pgfqpoint{2.043381in}{1.357123in}}%
\pgfpathlineto{\pgfqpoint{2.034708in}{1.355023in}}%
\pgfpathlineto{\pgfqpoint{2.030596in}{1.356610in}}%
\pgfpathlineto{\pgfqpoint{2.021694in}{1.353951in}}%
\pgfpathlineto{\pgfqpoint{2.021423in}{1.343723in}}%
\pgfpathlineto{\pgfqpoint{2.014646in}{1.343860in}}%
\pgfpathlineto{\pgfqpoint{2.014334in}{1.326679in}}%
\pgfpathlineto{\pgfqpoint{2.021173in}{1.326536in}}%
\pgfpathlineto{\pgfqpoint{2.023407in}{1.319659in}}%
\pgfpathclose%
\pgfusepath{fill}%
\end{pgfscope}%
\begin{pgfscope}%
\pgfpathrectangle{\pgfqpoint{0.100000in}{0.100000in}}{\pgfqpoint{3.608454in}{2.310000in}}%
\pgfusepath{clip}%
\pgfsetbuttcap%
\pgfsetmiterjoin%
\definecolor{currentfill}{rgb}{0.000000,0.482353,0.758824}%
\pgfsetfillcolor{currentfill}%
\pgfsetlinewidth{0.000000pt}%
\definecolor{currentstroke}{rgb}{0.000000,0.000000,0.000000}%
\pgfsetstrokecolor{currentstroke}%
\pgfsetstrokeopacity{0.000000}%
\pgfsetdash{}{0pt}%
\pgfpathmoveto{\pgfqpoint{2.986399in}{1.102377in}}%
\pgfpathlineto{\pgfqpoint{2.947537in}{1.098430in}}%
\pgfpathlineto{\pgfqpoint{2.946773in}{1.103937in}}%
\pgfpathlineto{\pgfqpoint{2.932868in}{1.114078in}}%
\pgfpathlineto{\pgfqpoint{2.926537in}{1.111891in}}%
\pgfpathlineto{\pgfqpoint{2.924776in}{1.115579in}}%
\pgfpathlineto{\pgfqpoint{2.926944in}{1.122067in}}%
\pgfpathlineto{\pgfqpoint{2.930944in}{1.123301in}}%
\pgfpathlineto{\pgfqpoint{2.943450in}{1.124830in}}%
\pgfpathlineto{\pgfqpoint{2.951598in}{1.127420in}}%
\pgfpathlineto{\pgfqpoint{2.956399in}{1.133737in}}%
\pgfpathlineto{\pgfqpoint{2.963707in}{1.130755in}}%
\pgfpathlineto{\pgfqpoint{2.970938in}{1.132406in}}%
\pgfpathlineto{\pgfqpoint{3.007763in}{1.136339in}}%
\pgfpathlineto{\pgfqpoint{3.008399in}{1.134076in}}%
\pgfpathlineto{\pgfqpoint{3.009870in}{1.122844in}}%
\pgfpathlineto{\pgfqpoint{3.007642in}{1.113989in}}%
\pgfpathlineto{\pgfqpoint{3.008536in}{1.104446in}}%
\pgfpathlineto{\pgfqpoint{2.986399in}{1.102377in}}%
\pgfpathclose%
\pgfusepath{fill}%
\end{pgfscope}%
\begin{pgfscope}%
\pgfpathrectangle{\pgfqpoint{0.100000in}{0.100000in}}{\pgfqpoint{3.608454in}{2.310000in}}%
\pgfusepath{clip}%
\pgfsetbuttcap%
\pgfsetmiterjoin%
\definecolor{currentfill}{rgb}{0.000000,0.631373,0.684314}%
\pgfsetfillcolor{currentfill}%
\pgfsetlinewidth{0.000000pt}%
\definecolor{currentstroke}{rgb}{0.000000,0.000000,0.000000}%
\pgfsetstrokecolor{currentstroke}%
\pgfsetstrokeopacity{0.000000}%
\pgfsetdash{}{0pt}%
\pgfpathmoveto{\pgfqpoint{2.922696in}{0.996327in}}%
\pgfpathlineto{\pgfqpoint{2.916400in}{0.998276in}}%
\pgfpathlineto{\pgfqpoint{2.908530in}{0.995600in}}%
\pgfpathlineto{\pgfqpoint{2.897463in}{0.997613in}}%
\pgfpathlineto{\pgfqpoint{2.885344in}{1.008668in}}%
\pgfpathlineto{\pgfqpoint{2.886412in}{1.013136in}}%
\pgfpathlineto{\pgfqpoint{2.871596in}{1.014173in}}%
\pgfpathlineto{\pgfqpoint{2.868390in}{1.010555in}}%
\pgfpathlineto{\pgfqpoint{2.864861in}{1.018162in}}%
\pgfpathlineto{\pgfqpoint{2.864516in}{1.026478in}}%
\pgfpathlineto{\pgfqpoint{2.859003in}{1.029918in}}%
\pgfpathlineto{\pgfqpoint{2.863172in}{1.045964in}}%
\pgfpathlineto{\pgfqpoint{2.864794in}{1.047108in}}%
\pgfpathlineto{\pgfqpoint{2.872404in}{1.042068in}}%
\pgfpathlineto{\pgfqpoint{2.876324in}{1.042276in}}%
\pgfpathlineto{\pgfqpoint{2.884093in}{1.034847in}}%
\pgfpathlineto{\pgfqpoint{2.888625in}{1.033073in}}%
\pgfpathlineto{\pgfqpoint{2.895066in}{1.035008in}}%
\pgfpathlineto{\pgfqpoint{2.907921in}{1.017944in}}%
\pgfpathlineto{\pgfqpoint{2.910642in}{1.010268in}}%
\pgfpathlineto{\pgfqpoint{2.919954in}{1.001789in}}%
\pgfpathlineto{\pgfqpoint{2.922696in}{0.996327in}}%
\pgfpathclose%
\pgfusepath{fill}%
\end{pgfscope}%
\begin{pgfscope}%
\pgfpathrectangle{\pgfqpoint{0.100000in}{0.100000in}}{\pgfqpoint{3.608454in}{2.310000in}}%
\pgfusepath{clip}%
\pgfsetbuttcap%
\pgfsetmiterjoin%
\definecolor{currentfill}{rgb}{0.000000,0.647059,0.676471}%
\pgfsetfillcolor{currentfill}%
\pgfsetlinewidth{0.000000pt}%
\definecolor{currentstroke}{rgb}{0.000000,0.000000,0.000000}%
\pgfsetstrokecolor{currentstroke}%
\pgfsetstrokeopacity{0.000000}%
\pgfsetdash{}{0pt}%
\pgfpathmoveto{\pgfqpoint{0.960301in}{1.531851in}}%
\pgfpathlineto{\pgfqpoint{0.954290in}{1.503771in}}%
\pgfpathlineto{\pgfqpoint{0.939931in}{1.436894in}}%
\pgfpathlineto{\pgfqpoint{0.882599in}{1.449402in}}%
\pgfpathlineto{\pgfqpoint{0.836825in}{1.499170in}}%
\pgfpathlineto{\pgfqpoint{0.795453in}{1.508995in}}%
\pgfpathlineto{\pgfqpoint{0.752002in}{1.519867in}}%
\pgfpathlineto{\pgfqpoint{0.724276in}{1.521184in}}%
\pgfpathlineto{\pgfqpoint{0.725481in}{1.527601in}}%
\pgfpathlineto{\pgfqpoint{0.730363in}{1.531286in}}%
\pgfpathlineto{\pgfqpoint{0.730097in}{1.540859in}}%
\pgfpathlineto{\pgfqpoint{0.738497in}{1.547647in}}%
\pgfpathlineto{\pgfqpoint{0.749642in}{1.550822in}}%
\pgfpathlineto{\pgfqpoint{0.753176in}{1.557006in}}%
\pgfpathlineto{\pgfqpoint{0.753615in}{1.567143in}}%
\pgfpathlineto{\pgfqpoint{0.758589in}{1.580048in}}%
\pgfpathlineto{\pgfqpoint{0.755783in}{1.587192in}}%
\pgfpathlineto{\pgfqpoint{0.785394in}{1.632033in}}%
\pgfpathlineto{\pgfqpoint{0.798688in}{1.628778in}}%
\pgfpathlineto{\pgfqpoint{0.805480in}{1.656197in}}%
\pgfpathlineto{\pgfqpoint{0.824452in}{1.732928in}}%
\pgfpathlineto{\pgfqpoint{0.861529in}{1.723655in}}%
\pgfpathlineto{\pgfqpoint{0.937739in}{1.705964in}}%
\pgfpathlineto{\pgfqpoint{0.949739in}{1.703791in}}%
\pgfpathlineto{\pgfqpoint{0.981161in}{1.696275in}}%
\pgfpathlineto{\pgfqpoint{0.994964in}{1.693281in}}%
\pgfpathlineto{\pgfqpoint{0.978586in}{1.616384in}}%
\pgfpathlineto{\pgfqpoint{0.960301in}{1.531851in}}%
\pgfpathclose%
\pgfusepath{fill}%
\end{pgfscope}%
\begin{pgfscope}%
\pgfpathrectangle{\pgfqpoint{0.100000in}{0.100000in}}{\pgfqpoint{3.608454in}{2.310000in}}%
\pgfusepath{clip}%
\pgfsetbuttcap%
\pgfsetmiterjoin%
\definecolor{currentfill}{rgb}{0.000000,0.796078,0.601961}%
\pgfsetfillcolor{currentfill}%
\pgfsetlinewidth{0.000000pt}%
\definecolor{currentstroke}{rgb}{0.000000,0.000000,0.000000}%
\pgfsetstrokecolor{currentstroke}%
\pgfsetstrokeopacity{0.000000}%
\pgfsetdash{}{0pt}%
\pgfpathmoveto{\pgfqpoint{1.859723in}{1.701494in}}%
\pgfpathlineto{\pgfqpoint{1.862167in}{1.696367in}}%
\pgfpathlineto{\pgfqpoint{1.867137in}{1.695946in}}%
\pgfpathlineto{\pgfqpoint{1.872226in}{1.686715in}}%
\pgfpathlineto{\pgfqpoint{1.881583in}{1.680508in}}%
\pgfpathlineto{\pgfqpoint{1.884096in}{1.673023in}}%
\pgfpathlineto{\pgfqpoint{1.888697in}{1.672141in}}%
\pgfpathlineto{\pgfqpoint{1.893266in}{1.666452in}}%
\pgfpathlineto{\pgfqpoint{1.902600in}{1.663282in}}%
\pgfpathlineto{\pgfqpoint{1.904126in}{1.659838in}}%
\pgfpathlineto{\pgfqpoint{1.860373in}{1.661671in}}%
\pgfpathlineto{\pgfqpoint{1.805743in}{1.664546in}}%
\pgfpathlineto{\pgfqpoint{1.805445in}{1.674881in}}%
\pgfpathlineto{\pgfqpoint{1.806996in}{1.702523in}}%
\pgfpathlineto{\pgfqpoint{1.807130in}{1.721260in}}%
\pgfpathlineto{\pgfqpoint{1.817373in}{1.718961in}}%
\pgfpathlineto{\pgfqpoint{1.829398in}{1.718027in}}%
\pgfpathlineto{\pgfqpoint{1.842006in}{1.721715in}}%
\pgfpathlineto{\pgfqpoint{1.841021in}{1.702423in}}%
\pgfpathlineto{\pgfqpoint{1.859723in}{1.701494in}}%
\pgfpathclose%
\pgfusepath{fill}%
\end{pgfscope}%
\begin{pgfscope}%
\pgfpathrectangle{\pgfqpoint{0.100000in}{0.100000in}}{\pgfqpoint{3.608454in}{2.310000in}}%
\pgfusepath{clip}%
\pgfsetbuttcap%
\pgfsetmiterjoin%
\definecolor{currentfill}{rgb}{0.000000,0.709804,0.645098}%
\pgfsetfillcolor{currentfill}%
\pgfsetlinewidth{0.000000pt}%
\definecolor{currentstroke}{rgb}{0.000000,0.000000,0.000000}%
\pgfsetstrokecolor{currentstroke}%
\pgfsetstrokeopacity{0.000000}%
\pgfsetdash{}{0pt}%
\pgfpathmoveto{\pgfqpoint{2.811462in}{0.766109in}}%
\pgfpathlineto{\pgfqpoint{2.792796in}{0.763762in}}%
\pgfpathlineto{\pgfqpoint{2.794089in}{0.749184in}}%
\pgfpathlineto{\pgfqpoint{2.786852in}{0.748618in}}%
\pgfpathlineto{\pgfqpoint{2.781206in}{0.755220in}}%
\pgfpathlineto{\pgfqpoint{2.779513in}{0.764362in}}%
\pgfpathlineto{\pgfqpoint{2.781911in}{0.785540in}}%
\pgfpathlineto{\pgfqpoint{2.780161in}{0.791452in}}%
\pgfpathlineto{\pgfqpoint{2.774910in}{0.796997in}}%
\pgfpathlineto{\pgfqpoint{2.755229in}{0.795820in}}%
\pgfpathlineto{\pgfqpoint{2.755934in}{0.788907in}}%
\pgfpathlineto{\pgfqpoint{2.733549in}{0.786465in}}%
\pgfpathlineto{\pgfqpoint{2.737974in}{0.800305in}}%
\pgfpathlineto{\pgfqpoint{2.737520in}{0.807875in}}%
\pgfpathlineto{\pgfqpoint{2.741287in}{0.817640in}}%
\pgfpathlineto{\pgfqpoint{2.751810in}{0.820152in}}%
\pgfpathlineto{\pgfqpoint{2.751930in}{0.830708in}}%
\pgfpathlineto{\pgfqpoint{2.762136in}{0.831926in}}%
\pgfpathlineto{\pgfqpoint{2.767748in}{0.825602in}}%
\pgfpathlineto{\pgfqpoint{2.776392in}{0.826653in}}%
\pgfpathlineto{\pgfqpoint{2.778092in}{0.819839in}}%
\pgfpathlineto{\pgfqpoint{2.786918in}{0.816730in}}%
\pgfpathlineto{\pgfqpoint{2.808283in}{0.818788in}}%
\pgfpathlineto{\pgfqpoint{2.809408in}{0.807037in}}%
\pgfpathlineto{\pgfqpoint{2.814786in}{0.800405in}}%
\pgfpathlineto{\pgfqpoint{2.821071in}{0.796344in}}%
\pgfpathlineto{\pgfqpoint{2.821443in}{0.790502in}}%
\pgfpathlineto{\pgfqpoint{2.825145in}{0.783605in}}%
\pgfpathlineto{\pgfqpoint{2.820126in}{0.780667in}}%
\pgfpathlineto{\pgfqpoint{2.810152in}{0.780002in}}%
\pgfpathlineto{\pgfqpoint{2.811462in}{0.766109in}}%
\pgfpathclose%
\pgfusepath{fill}%
\end{pgfscope}%
\begin{pgfscope}%
\pgfpathrectangle{\pgfqpoint{0.100000in}{0.100000in}}{\pgfqpoint{3.608454in}{2.310000in}}%
\pgfusepath{clip}%
\pgfsetbuttcap%
\pgfsetmiterjoin%
\definecolor{currentfill}{rgb}{0.000000,0.635294,0.682353}%
\pgfsetfillcolor{currentfill}%
\pgfsetlinewidth{0.000000pt}%
\definecolor{currentstroke}{rgb}{0.000000,0.000000,0.000000}%
\pgfsetstrokecolor{currentstroke}%
\pgfsetstrokeopacity{0.000000}%
\pgfsetdash{}{0pt}%
\pgfpathmoveto{\pgfqpoint{1.088551in}{1.993499in}}%
\pgfpathlineto{\pgfqpoint{1.086269in}{1.988945in}}%
\pgfpathlineto{\pgfqpoint{1.076963in}{1.988148in}}%
\pgfpathlineto{\pgfqpoint{1.069612in}{1.982716in}}%
\pgfpathlineto{\pgfqpoint{1.061696in}{1.979359in}}%
\pgfpathlineto{\pgfqpoint{1.053520in}{1.975559in}}%
\pgfpathlineto{\pgfqpoint{1.048671in}{1.970878in}}%
\pgfpathlineto{\pgfqpoint{1.034877in}{1.968576in}}%
\pgfpathlineto{\pgfqpoint{1.028023in}{1.976205in}}%
\pgfpathlineto{\pgfqpoint{1.025284in}{1.975956in}}%
\pgfpathlineto{\pgfqpoint{1.030719in}{1.983835in}}%
\pgfpathlineto{\pgfqpoint{1.028886in}{1.992775in}}%
\pgfpathlineto{\pgfqpoint{1.033186in}{1.997675in}}%
\pgfpathlineto{\pgfqpoint{1.040540in}{1.999733in}}%
\pgfpathlineto{\pgfqpoint{1.037550in}{2.013584in}}%
\pgfpathlineto{\pgfqpoint{1.042048in}{2.022019in}}%
\pgfpathlineto{\pgfqpoint{1.042825in}{2.030433in}}%
\pgfpathlineto{\pgfqpoint{1.049410in}{2.040224in}}%
\pgfpathlineto{\pgfqpoint{1.056865in}{2.058071in}}%
\pgfpathlineto{\pgfqpoint{1.055012in}{2.060043in}}%
\pgfpathlineto{\pgfqpoint{1.041035in}{2.060695in}}%
\pgfpathlineto{\pgfqpoint{1.041857in}{2.066086in}}%
\pgfpathlineto{\pgfqpoint{1.027942in}{2.078990in}}%
\pgfpathlineto{\pgfqpoint{1.027701in}{2.087646in}}%
\pgfpathlineto{\pgfqpoint{1.023030in}{2.093793in}}%
\pgfpathlineto{\pgfqpoint{1.014258in}{2.117180in}}%
\pgfpathlineto{\pgfqpoint{1.003956in}{2.123021in}}%
\pgfpathlineto{\pgfqpoint{1.003713in}{2.127225in}}%
\pgfpathlineto{\pgfqpoint{0.994972in}{2.136779in}}%
\pgfpathlineto{\pgfqpoint{1.001953in}{2.138275in}}%
\pgfpathlineto{\pgfqpoint{1.011451in}{2.134988in}}%
\pgfpathlineto{\pgfqpoint{1.020317in}{2.134665in}}%
\pgfpathlineto{\pgfqpoint{1.034696in}{2.123940in}}%
\pgfpathlineto{\pgfqpoint{1.033798in}{2.117353in}}%
\pgfpathlineto{\pgfqpoint{1.039741in}{2.112438in}}%
\pgfpathlineto{\pgfqpoint{1.048078in}{2.111068in}}%
\pgfpathlineto{\pgfqpoint{1.060904in}{2.103004in}}%
\pgfpathlineto{\pgfqpoint{1.068378in}{2.094075in}}%
\pgfpathlineto{\pgfqpoint{1.072235in}{2.094680in}}%
\pgfpathlineto{\pgfqpoint{1.085713in}{2.091817in}}%
\pgfpathlineto{\pgfqpoint{1.092364in}{2.093954in}}%
\pgfpathlineto{\pgfqpoint{1.091464in}{2.102885in}}%
\pgfpathlineto{\pgfqpoint{1.093234in}{2.110956in}}%
\pgfpathlineto{\pgfqpoint{1.090372in}{2.118413in}}%
\pgfpathlineto{\pgfqpoint{1.093543in}{2.127613in}}%
\pgfpathlineto{\pgfqpoint{1.118004in}{2.122768in}}%
\pgfpathlineto{\pgfqpoint{1.111233in}{2.089809in}}%
\pgfpathlineto{\pgfqpoint{1.119922in}{2.088067in}}%
\pgfpathlineto{\pgfqpoint{1.114370in}{2.060903in}}%
\pgfpathlineto{\pgfqpoint{1.108312in}{2.062127in}}%
\pgfpathlineto{\pgfqpoint{1.103510in}{2.056028in}}%
\pgfpathlineto{\pgfqpoint{1.092981in}{2.055856in}}%
\pgfpathlineto{\pgfqpoint{1.092070in}{2.051353in}}%
\pgfpathlineto{\pgfqpoint{1.085312in}{2.052758in}}%
\pgfpathlineto{\pgfqpoint{1.078479in}{2.042039in}}%
\pgfpathlineto{\pgfqpoint{1.077963in}{2.034819in}}%
\pgfpathlineto{\pgfqpoint{1.081054in}{2.029697in}}%
\pgfpathlineto{\pgfqpoint{1.081296in}{2.021948in}}%
\pgfpathlineto{\pgfqpoint{1.076933in}{2.013289in}}%
\pgfpathlineto{\pgfqpoint{1.076031in}{2.007092in}}%
\pgfpathlineto{\pgfqpoint{1.080846in}{2.002972in}}%
\pgfpathlineto{\pgfqpoint{1.081829in}{1.997042in}}%
\pgfpathlineto{\pgfqpoint{1.088551in}{1.993499in}}%
\pgfpathclose%
\pgfusepath{fill}%
\end{pgfscope}%
\begin{pgfscope}%
\pgfpathrectangle{\pgfqpoint{0.100000in}{0.100000in}}{\pgfqpoint{3.608454in}{2.310000in}}%
\pgfusepath{clip}%
\pgfsetbuttcap%
\pgfsetmiterjoin%
\definecolor{currentfill}{rgb}{0.000000,0.701961,0.649020}%
\pgfsetfillcolor{currentfill}%
\pgfsetlinewidth{0.000000pt}%
\definecolor{currentstroke}{rgb}{0.000000,0.000000,0.000000}%
\pgfsetstrokecolor{currentstroke}%
\pgfsetstrokeopacity{0.000000}%
\pgfsetdash{}{0pt}%
\pgfpathmoveto{\pgfqpoint{2.193329in}{1.907307in}}%
\pgfpathlineto{\pgfqpoint{2.214175in}{1.908040in}}%
\pgfpathlineto{\pgfqpoint{2.213676in}{1.928846in}}%
\pgfpathlineto{\pgfqpoint{2.255541in}{1.929754in}}%
\pgfpathlineto{\pgfqpoint{2.256111in}{1.909051in}}%
\pgfpathlineto{\pgfqpoint{2.256290in}{1.902446in}}%
\pgfpathlineto{\pgfqpoint{2.253310in}{1.897735in}}%
\pgfpathlineto{\pgfqpoint{2.237467in}{1.890615in}}%
\pgfpathlineto{\pgfqpoint{2.233501in}{1.887214in}}%
\pgfpathlineto{\pgfqpoint{2.229731in}{1.876962in}}%
\pgfpathlineto{\pgfqpoint{2.226726in}{1.874225in}}%
\pgfpathlineto{\pgfqpoint{2.189353in}{1.873824in}}%
\pgfpathlineto{\pgfqpoint{2.188794in}{1.893513in}}%
\pgfpathlineto{\pgfqpoint{2.193529in}{1.893561in}}%
\pgfpathlineto{\pgfqpoint{2.193329in}{1.907307in}}%
\pgfpathclose%
\pgfusepath{fill}%
\end{pgfscope}%
\begin{pgfscope}%
\pgfpathrectangle{\pgfqpoint{0.100000in}{0.100000in}}{\pgfqpoint{3.608454in}{2.310000in}}%
\pgfusepath{clip}%
\pgfsetbuttcap%
\pgfsetmiterjoin%
\definecolor{currentfill}{rgb}{0.000000,0.592157,0.703922}%
\pgfsetfillcolor{currentfill}%
\pgfsetlinewidth{0.000000pt}%
\definecolor{currentstroke}{rgb}{0.000000,0.000000,0.000000}%
\pgfsetstrokecolor{currentstroke}%
\pgfsetstrokeopacity{0.000000}%
\pgfsetdash{}{0pt}%
\pgfpathmoveto{\pgfqpoint{1.854168in}{1.187060in}}%
\pgfpathlineto{\pgfqpoint{1.825437in}{1.188388in}}%
\pgfpathlineto{\pgfqpoint{1.820080in}{1.188664in}}%
\pgfpathlineto{\pgfqpoint{1.821457in}{1.218805in}}%
\pgfpathlineto{\pgfqpoint{1.855479in}{1.217455in}}%
\pgfpathlineto{\pgfqpoint{1.854168in}{1.187060in}}%
\pgfpathclose%
\pgfusepath{fill}%
\end{pgfscope}%
\begin{pgfscope}%
\pgfpathrectangle{\pgfqpoint{0.100000in}{0.100000in}}{\pgfqpoint{3.608454in}{2.310000in}}%
\pgfusepath{clip}%
\pgfsetbuttcap%
\pgfsetmiterjoin%
\definecolor{currentfill}{rgb}{0.000000,0.701961,0.649020}%
\pgfsetfillcolor{currentfill}%
\pgfsetlinewidth{0.000000pt}%
\definecolor{currentstroke}{rgb}{0.000000,0.000000,0.000000}%
\pgfsetstrokecolor{currentstroke}%
\pgfsetstrokeopacity{0.000000}%
\pgfsetdash{}{0pt}%
\pgfpathmoveto{\pgfqpoint{2.589411in}{1.214310in}}%
\pgfpathlineto{\pgfqpoint{2.577947in}{1.209252in}}%
\pgfpathlineto{\pgfqpoint{2.567654in}{1.211922in}}%
\pgfpathlineto{\pgfqpoint{2.566571in}{1.201683in}}%
\pgfpathlineto{\pgfqpoint{2.563459in}{1.197714in}}%
\pgfpathlineto{\pgfqpoint{2.555635in}{1.195943in}}%
\pgfpathlineto{\pgfqpoint{2.539380in}{1.187318in}}%
\pgfpathlineto{\pgfqpoint{2.533377in}{1.204248in}}%
\pgfpathlineto{\pgfqpoint{2.535791in}{1.209919in}}%
\pgfpathlineto{\pgfqpoint{2.533129in}{1.219044in}}%
\pgfpathlineto{\pgfqpoint{2.523585in}{1.228784in}}%
\pgfpathlineto{\pgfqpoint{2.528115in}{1.232902in}}%
\pgfpathlineto{\pgfqpoint{2.540147in}{1.235271in}}%
\pgfpathlineto{\pgfqpoint{2.534558in}{1.250105in}}%
\pgfpathlineto{\pgfqpoint{2.541890in}{1.261426in}}%
\pgfpathlineto{\pgfqpoint{2.549388in}{1.262674in}}%
\pgfpathlineto{\pgfqpoint{2.546912in}{1.269077in}}%
\pgfpathlineto{\pgfqpoint{2.553706in}{1.268477in}}%
\pgfpathlineto{\pgfqpoint{2.562848in}{1.271335in}}%
\pgfpathlineto{\pgfqpoint{2.567447in}{1.266096in}}%
\pgfpathlineto{\pgfqpoint{2.570793in}{1.273710in}}%
\pgfpathlineto{\pgfqpoint{2.578231in}{1.275886in}}%
\pgfpathlineto{\pgfqpoint{2.586412in}{1.272897in}}%
\pgfpathlineto{\pgfqpoint{2.585709in}{1.266483in}}%
\pgfpathlineto{\pgfqpoint{2.577210in}{1.248460in}}%
\pgfpathlineto{\pgfqpoint{2.585128in}{1.245024in}}%
\pgfpathlineto{\pgfqpoint{2.586421in}{1.235460in}}%
\pgfpathlineto{\pgfqpoint{2.589982in}{1.231634in}}%
\pgfpathlineto{\pgfqpoint{2.585378in}{1.221768in}}%
\pgfpathlineto{\pgfqpoint{2.589411in}{1.214310in}}%
\pgfpathclose%
\pgfusepath{fill}%
\end{pgfscope}%
\begin{pgfscope}%
\pgfpathrectangle{\pgfqpoint{0.100000in}{0.100000in}}{\pgfqpoint{3.608454in}{2.310000in}}%
\pgfusepath{clip}%
\pgfsetbuttcap%
\pgfsetmiterjoin%
\definecolor{currentfill}{rgb}{0.000000,0.588235,0.705882}%
\pgfsetfillcolor{currentfill}%
\pgfsetlinewidth{0.000000pt}%
\definecolor{currentstroke}{rgb}{0.000000,0.000000,0.000000}%
\pgfsetstrokecolor{currentstroke}%
\pgfsetstrokeopacity{0.000000}%
\pgfsetdash{}{0pt}%
\pgfpathmoveto{\pgfqpoint{1.638167in}{2.022688in}}%
\pgfpathlineto{\pgfqpoint{1.635493in}{1.994898in}}%
\pgfpathlineto{\pgfqpoint{1.638609in}{1.994610in}}%
\pgfpathlineto{\pgfqpoint{1.635984in}{1.966799in}}%
\pgfpathlineto{\pgfqpoint{1.625532in}{1.967784in}}%
\pgfpathlineto{\pgfqpoint{1.624849in}{1.960734in}}%
\pgfpathlineto{\pgfqpoint{1.611520in}{1.962062in}}%
\pgfpathlineto{\pgfqpoint{1.612308in}{1.970040in}}%
\pgfpathlineto{\pgfqpoint{1.595495in}{1.971747in}}%
\pgfpathlineto{\pgfqpoint{1.592373in}{1.975555in}}%
\pgfpathlineto{\pgfqpoint{1.582076in}{1.976639in}}%
\pgfpathlineto{\pgfqpoint{1.583733in}{1.990580in}}%
\pgfpathlineto{\pgfqpoint{1.560751in}{1.993122in}}%
\pgfpathlineto{\pgfqpoint{1.558985in}{1.997946in}}%
\pgfpathlineto{\pgfqpoint{1.552079in}{1.998723in}}%
\pgfpathlineto{\pgfqpoint{1.552609in}{2.003322in}}%
\pgfpathlineto{\pgfqpoint{1.545674in}{2.004107in}}%
\pgfpathlineto{\pgfqpoint{1.547528in}{2.020313in}}%
\pgfpathlineto{\pgfqpoint{1.543137in}{2.020826in}}%
\pgfpathlineto{\pgfqpoint{1.546296in}{2.048425in}}%
\pgfpathlineto{\pgfqpoint{1.548607in}{2.048174in}}%
\pgfpathlineto{\pgfqpoint{1.551007in}{2.068911in}}%
\pgfpathlineto{\pgfqpoint{1.557865in}{2.068157in}}%
\pgfpathlineto{\pgfqpoint{1.557068in}{2.061233in}}%
\pgfpathlineto{\pgfqpoint{1.570807in}{2.059681in}}%
\pgfpathlineto{\pgfqpoint{1.570016in}{2.052759in}}%
\pgfpathlineto{\pgfqpoint{1.590639in}{2.050503in}}%
\pgfpathlineto{\pgfqpoint{1.589896in}{2.043588in}}%
\pgfpathlineto{\pgfqpoint{1.594243in}{2.043152in}}%
\pgfpathlineto{\pgfqpoint{1.592763in}{2.029281in}}%
\pgfpathlineto{\pgfqpoint{1.613315in}{2.027119in}}%
\pgfpathlineto{\pgfqpoint{1.617764in}{2.024744in}}%
\pgfpathlineto{\pgfqpoint{1.638167in}{2.022688in}}%
\pgfpathclose%
\pgfusepath{fill}%
\end{pgfscope}%
\begin{pgfscope}%
\pgfpathrectangle{\pgfqpoint{0.100000in}{0.100000in}}{\pgfqpoint{3.608454in}{2.310000in}}%
\pgfusepath{clip}%
\pgfsetbuttcap%
\pgfsetmiterjoin%
\definecolor{currentfill}{rgb}{0.000000,0.647059,0.676471}%
\pgfsetfillcolor{currentfill}%
\pgfsetlinewidth{0.000000pt}%
\definecolor{currentstroke}{rgb}{0.000000,0.000000,0.000000}%
\pgfsetstrokecolor{currentstroke}%
\pgfsetstrokeopacity{0.000000}%
\pgfsetdash{}{0pt}%
\pgfpathmoveto{\pgfqpoint{2.534753in}{1.425828in}}%
\pgfpathlineto{\pgfqpoint{2.534185in}{1.410410in}}%
\pgfpathlineto{\pgfqpoint{2.528514in}{1.407327in}}%
\pgfpathlineto{\pgfqpoint{2.503724in}{1.405402in}}%
\pgfpathlineto{\pgfqpoint{2.506265in}{1.370984in}}%
\pgfpathlineto{\pgfqpoint{2.465371in}{1.368430in}}%
\pgfpathlineto{\pgfqpoint{2.464725in}{1.378777in}}%
\pgfpathlineto{\pgfqpoint{2.471713in}{1.379003in}}%
\pgfpathlineto{\pgfqpoint{2.470095in}{1.403294in}}%
\pgfpathlineto{\pgfqpoint{2.463205in}{1.402934in}}%
\pgfpathlineto{\pgfqpoint{2.462297in}{1.414401in}}%
\pgfpathlineto{\pgfqpoint{2.457647in}{1.415090in}}%
\pgfpathlineto{\pgfqpoint{2.457108in}{1.423280in}}%
\pgfpathlineto{\pgfqpoint{2.461611in}{1.423638in}}%
\pgfpathlineto{\pgfqpoint{2.460863in}{1.433950in}}%
\pgfpathlineto{\pgfqpoint{2.484922in}{1.436043in}}%
\pgfpathlineto{\pgfqpoint{2.493952in}{1.454636in}}%
\pgfpathlineto{\pgfqpoint{2.500893in}{1.455137in}}%
\pgfpathlineto{\pgfqpoint{2.498999in}{1.481611in}}%
\pgfpathlineto{\pgfqpoint{2.512434in}{1.482660in}}%
\pgfpathlineto{\pgfqpoint{2.509486in}{1.512307in}}%
\pgfpathlineto{\pgfqpoint{2.516336in}{1.513050in}}%
\pgfpathlineto{\pgfqpoint{2.520205in}{1.472936in}}%
\pgfpathlineto{\pgfqpoint{2.531157in}{1.473630in}}%
\pgfpathlineto{\pgfqpoint{2.532347in}{1.453054in}}%
\pgfpathlineto{\pgfqpoint{2.534753in}{1.425828in}}%
\pgfpathclose%
\pgfusepath{fill}%
\end{pgfscope}%
\begin{pgfscope}%
\pgfpathrectangle{\pgfqpoint{0.100000in}{0.100000in}}{\pgfqpoint{3.608454in}{2.310000in}}%
\pgfusepath{clip}%
\pgfsetbuttcap%
\pgfsetmiterjoin%
\definecolor{currentfill}{rgb}{0.000000,0.478431,0.760784}%
\pgfsetfillcolor{currentfill}%
\pgfsetlinewidth{0.000000pt}%
\definecolor{currentstroke}{rgb}{0.000000,0.000000,0.000000}%
\pgfsetstrokecolor{currentstroke}%
\pgfsetstrokeopacity{0.000000}%
\pgfsetdash{}{0pt}%
\pgfpathmoveto{\pgfqpoint{1.638167in}{2.022688in}}%
\pgfpathlineto{\pgfqpoint{1.617764in}{2.024744in}}%
\pgfpathlineto{\pgfqpoint{1.613315in}{2.027119in}}%
\pgfpathlineto{\pgfqpoint{1.592763in}{2.029281in}}%
\pgfpathlineto{\pgfqpoint{1.594243in}{2.043152in}}%
\pgfpathlineto{\pgfqpoint{1.589896in}{2.043588in}}%
\pgfpathlineto{\pgfqpoint{1.590639in}{2.050503in}}%
\pgfpathlineto{\pgfqpoint{1.570016in}{2.052759in}}%
\pgfpathlineto{\pgfqpoint{1.570807in}{2.059681in}}%
\pgfpathlineto{\pgfqpoint{1.557068in}{2.061233in}}%
\pgfpathlineto{\pgfqpoint{1.557865in}{2.068157in}}%
\pgfpathlineto{\pgfqpoint{1.560678in}{2.074073in}}%
\pgfpathlineto{\pgfqpoint{1.562458in}{2.089780in}}%
\pgfpathlineto{\pgfqpoint{1.569948in}{2.088334in}}%
\pgfpathlineto{\pgfqpoint{1.582871in}{2.091842in}}%
\pgfpathlineto{\pgfqpoint{1.590181in}{2.089744in}}%
\pgfpathlineto{\pgfqpoint{1.601867in}{2.089959in}}%
\pgfpathlineto{\pgfqpoint{1.607083in}{2.083595in}}%
\pgfpathlineto{\pgfqpoint{1.619223in}{2.081466in}}%
\pgfpathlineto{\pgfqpoint{1.623124in}{2.077735in}}%
\pgfpathlineto{\pgfqpoint{1.632154in}{2.074810in}}%
\pgfpathlineto{\pgfqpoint{1.643336in}{2.085862in}}%
\pgfpathlineto{\pgfqpoint{1.647283in}{2.084906in}}%
\pgfpathlineto{\pgfqpoint{1.648734in}{2.078432in}}%
\pgfpathlineto{\pgfqpoint{1.653010in}{2.075697in}}%
\pgfpathlineto{\pgfqpoint{1.664466in}{2.078762in}}%
\pgfpathlineto{\pgfqpoint{1.669040in}{2.084553in}}%
\pgfpathlineto{\pgfqpoint{1.686894in}{2.083116in}}%
\pgfpathlineto{\pgfqpoint{1.690580in}{2.101766in}}%
\pgfpathlineto{\pgfqpoint{1.687537in}{2.102006in}}%
\pgfpathlineto{\pgfqpoint{1.688717in}{2.115878in}}%
\pgfpathlineto{\pgfqpoint{1.723190in}{2.112972in}}%
\pgfpathlineto{\pgfqpoint{1.737002in}{2.111989in}}%
\pgfpathlineto{\pgfqpoint{1.735922in}{2.098000in}}%
\pgfpathlineto{\pgfqpoint{1.738640in}{2.097811in}}%
\pgfpathlineto{\pgfqpoint{1.736501in}{2.069931in}}%
\pgfpathlineto{\pgfqpoint{1.739211in}{2.069716in}}%
\pgfpathlineto{\pgfqpoint{1.738152in}{2.055718in}}%
\pgfpathlineto{\pgfqpoint{1.710634in}{2.057876in}}%
\pgfpathlineto{\pgfqpoint{1.710047in}{2.050726in}}%
\pgfpathlineto{\pgfqpoint{1.700751in}{2.053267in}}%
\pgfpathlineto{\pgfqpoint{1.695708in}{2.045165in}}%
\pgfpathlineto{\pgfqpoint{1.671188in}{2.047334in}}%
\pgfpathlineto{\pgfqpoint{1.668766in}{2.019808in}}%
\pgfpathlineto{\pgfqpoint{1.638167in}{2.022688in}}%
\pgfpathclose%
\pgfusepath{fill}%
\end{pgfscope}%
\begin{pgfscope}%
\pgfpathrectangle{\pgfqpoint{0.100000in}{0.100000in}}{\pgfqpoint{3.608454in}{2.310000in}}%
\pgfusepath{clip}%
\pgfsetbuttcap%
\pgfsetmiterjoin%
\definecolor{currentfill}{rgb}{0.000000,0.690196,0.654902}%
\pgfsetfillcolor{currentfill}%
\pgfsetlinewidth{0.000000pt}%
\definecolor{currentstroke}{rgb}{0.000000,0.000000,0.000000}%
\pgfsetstrokecolor{currentstroke}%
\pgfsetstrokeopacity{0.000000}%
\pgfsetdash{}{0pt}%
\pgfpathmoveto{\pgfqpoint{3.535136in}{1.728569in}}%
\pgfpathlineto{\pgfqpoint{3.546887in}{1.721036in}}%
\pgfpathlineto{\pgfqpoint{3.532181in}{1.716409in}}%
\pgfpathlineto{\pgfqpoint{3.530762in}{1.723556in}}%
\pgfpathlineto{\pgfqpoint{3.535136in}{1.728569in}}%
\pgfpathclose%
\pgfusepath{fill}%
\end{pgfscope}%
\begin{pgfscope}%
\pgfpathrectangle{\pgfqpoint{0.100000in}{0.100000in}}{\pgfqpoint{3.608454in}{2.310000in}}%
\pgfusepath{clip}%
\pgfsetbuttcap%
\pgfsetmiterjoin%
\definecolor{currentfill}{rgb}{0.000000,0.494118,0.752941}%
\pgfsetfillcolor{currentfill}%
\pgfsetlinewidth{0.000000pt}%
\definecolor{currentstroke}{rgb}{0.000000,0.000000,0.000000}%
\pgfsetstrokecolor{currentstroke}%
\pgfsetstrokeopacity{0.000000}%
\pgfsetdash{}{0pt}%
\pgfpathmoveto{\pgfqpoint{3.153524in}{1.337961in}}%
\pgfpathlineto{\pgfqpoint{3.149422in}{1.332537in}}%
\pgfpathlineto{\pgfqpoint{3.149648in}{1.328164in}}%
\pgfpathlineto{\pgfqpoint{3.140579in}{1.302693in}}%
\pgfpathlineto{\pgfqpoint{3.131753in}{1.305762in}}%
\pgfpathlineto{\pgfqpoint{3.125802in}{1.307188in}}%
\pgfpathlineto{\pgfqpoint{3.115324in}{1.315206in}}%
\pgfpathlineto{\pgfqpoint{3.116356in}{1.322582in}}%
\pgfpathlineto{\pgfqpoint{3.121341in}{1.327517in}}%
\pgfpathlineto{\pgfqpoint{3.120128in}{1.335691in}}%
\pgfpathlineto{\pgfqpoint{3.107539in}{1.353570in}}%
\pgfpathlineto{\pgfqpoint{3.110708in}{1.357752in}}%
\pgfpathlineto{\pgfqpoint{3.110807in}{1.366936in}}%
\pgfpathlineto{\pgfqpoint{3.115799in}{1.372179in}}%
\pgfpathlineto{\pgfqpoint{3.123623in}{1.385333in}}%
\pgfpathlineto{\pgfqpoint{3.125203in}{1.394430in}}%
\pgfpathlineto{\pgfqpoint{3.129858in}{1.404228in}}%
\pgfpathlineto{\pgfqpoint{3.137714in}{1.398047in}}%
\pgfpathlineto{\pgfqpoint{3.141594in}{1.398020in}}%
\pgfpathlineto{\pgfqpoint{3.153282in}{1.414127in}}%
\pgfpathlineto{\pgfqpoint{3.158541in}{1.406177in}}%
\pgfpathlineto{\pgfqpoint{3.169602in}{1.394541in}}%
\pgfpathlineto{\pgfqpoint{3.175771in}{1.395430in}}%
\pgfpathlineto{\pgfqpoint{3.161085in}{1.369292in}}%
\pgfpathlineto{\pgfqpoint{3.168135in}{1.369114in}}%
\pgfpathlineto{\pgfqpoint{3.178627in}{1.363629in}}%
\pgfpathlineto{\pgfqpoint{3.176043in}{1.340697in}}%
\pgfpathlineto{\pgfqpoint{3.169983in}{1.342070in}}%
\pgfpathlineto{\pgfqpoint{3.165049in}{1.347998in}}%
\pgfpathlineto{\pgfqpoint{3.157131in}{1.351262in}}%
\pgfpathlineto{\pgfqpoint{3.153524in}{1.337961in}}%
\pgfpathclose%
\pgfusepath{fill}%
\end{pgfscope}%
\begin{pgfscope}%
\pgfpathrectangle{\pgfqpoint{0.100000in}{0.100000in}}{\pgfqpoint{3.608454in}{2.310000in}}%
\pgfusepath{clip}%
\pgfsetbuttcap%
\pgfsetmiterjoin%
\definecolor{currentfill}{rgb}{0.000000,0.427451,0.786275}%
\pgfsetfillcolor{currentfill}%
\pgfsetlinewidth{0.000000pt}%
\definecolor{currentstroke}{rgb}{0.000000,0.000000,0.000000}%
\pgfsetstrokecolor{currentstroke}%
\pgfsetstrokeopacity{0.000000}%
\pgfsetdash{}{0pt}%
\pgfpathmoveto{\pgfqpoint{1.713023in}{1.050837in}}%
\pgfpathlineto{\pgfqpoint{1.710432in}{1.016551in}}%
\pgfpathlineto{\pgfqpoint{1.685686in}{1.018201in}}%
\pgfpathlineto{\pgfqpoint{1.651469in}{1.020799in}}%
\pgfpathlineto{\pgfqpoint{1.617270in}{1.023437in}}%
\pgfpathlineto{\pgfqpoint{1.583688in}{1.026354in}}%
\pgfpathlineto{\pgfqpoint{1.585141in}{1.042675in}}%
\pgfpathlineto{\pgfqpoint{1.589893in}{1.095406in}}%
\pgfpathlineto{\pgfqpoint{1.646284in}{1.091077in}}%
\pgfpathlineto{\pgfqpoint{1.680896in}{1.087866in}}%
\pgfpathlineto{\pgfqpoint{1.678367in}{1.053341in}}%
\pgfpathlineto{\pgfqpoint{1.713023in}{1.050837in}}%
\pgfpathclose%
\pgfusepath{fill}%
\end{pgfscope}%
\begin{pgfscope}%
\pgfpathrectangle{\pgfqpoint{0.100000in}{0.100000in}}{\pgfqpoint{3.608454in}{2.310000in}}%
\pgfusepath{clip}%
\pgfsetbuttcap%
\pgfsetmiterjoin%
\definecolor{currentfill}{rgb}{0.000000,0.470588,0.764706}%
\pgfsetfillcolor{currentfill}%
\pgfsetlinewidth{0.000000pt}%
\definecolor{currentstroke}{rgb}{0.000000,0.000000,0.000000}%
\pgfsetstrokecolor{currentstroke}%
\pgfsetstrokeopacity{0.000000}%
\pgfsetdash{}{0pt}%
\pgfpathmoveto{\pgfqpoint{2.281822in}{0.951221in}}%
\pgfpathlineto{\pgfqpoint{2.271324in}{0.957636in}}%
\pgfpathlineto{\pgfqpoint{2.259184in}{0.957741in}}%
\pgfpathlineto{\pgfqpoint{2.259281in}{0.978738in}}%
\pgfpathlineto{\pgfqpoint{2.258094in}{0.983370in}}%
\pgfpathlineto{\pgfqpoint{2.250959in}{0.985619in}}%
\pgfpathlineto{\pgfqpoint{2.249833in}{0.992475in}}%
\pgfpathlineto{\pgfqpoint{2.245351in}{0.995915in}}%
\pgfpathlineto{\pgfqpoint{2.238537in}{0.995926in}}%
\pgfpathlineto{\pgfqpoint{2.238992in}{1.006643in}}%
\pgfpathlineto{\pgfqpoint{2.218257in}{1.006289in}}%
\pgfpathlineto{\pgfqpoint{2.216940in}{1.014382in}}%
\pgfpathlineto{\pgfqpoint{2.226237in}{1.022520in}}%
\pgfpathlineto{\pgfqpoint{2.241035in}{1.037964in}}%
\pgfpathlineto{\pgfqpoint{2.245679in}{1.038542in}}%
\pgfpathlineto{\pgfqpoint{2.245419in}{1.061527in}}%
\pgfpathlineto{\pgfqpoint{2.269359in}{1.061552in}}%
\pgfpathlineto{\pgfqpoint{2.269431in}{1.054652in}}%
\pgfpathlineto{\pgfqpoint{2.293330in}{1.054914in}}%
\pgfpathlineto{\pgfqpoint{2.293540in}{1.031438in}}%
\pgfpathlineto{\pgfqpoint{2.300760in}{1.031780in}}%
\pgfpathlineto{\pgfqpoint{2.309489in}{1.028555in}}%
\pgfpathlineto{\pgfqpoint{2.314200in}{1.029274in}}%
\pgfpathlineto{\pgfqpoint{2.314414in}{1.020654in}}%
\pgfpathlineto{\pgfqpoint{2.321574in}{1.020760in}}%
\pgfpathlineto{\pgfqpoint{2.321823in}{1.001132in}}%
\pgfpathlineto{\pgfqpoint{2.326016in}{0.997543in}}%
\pgfpathlineto{\pgfqpoint{2.326278in}{0.991006in}}%
\pgfpathlineto{\pgfqpoint{2.322083in}{0.986111in}}%
\pgfpathlineto{\pgfqpoint{2.289234in}{0.985811in}}%
\pgfpathlineto{\pgfqpoint{2.289533in}{0.964998in}}%
\pgfpathlineto{\pgfqpoint{2.288597in}{0.951635in}}%
\pgfpathlineto{\pgfqpoint{2.281822in}{0.951221in}}%
\pgfpathclose%
\pgfusepath{fill}%
\end{pgfscope}%
\begin{pgfscope}%
\pgfpathrectangle{\pgfqpoint{0.100000in}{0.100000in}}{\pgfqpoint{3.608454in}{2.310000in}}%
\pgfusepath{clip}%
\pgfsetbuttcap%
\pgfsetmiterjoin%
\definecolor{currentfill}{rgb}{0.000000,0.611765,0.694118}%
\pgfsetfillcolor{currentfill}%
\pgfsetlinewidth{0.000000pt}%
\definecolor{currentstroke}{rgb}{0.000000,0.000000,0.000000}%
\pgfsetstrokecolor{currentstroke}%
\pgfsetstrokeopacity{0.000000}%
\pgfsetdash{}{0pt}%
\pgfpathmoveto{\pgfqpoint{1.553831in}{1.907730in}}%
\pgfpathlineto{\pgfqpoint{1.549754in}{1.873361in}}%
\pgfpathlineto{\pgfqpoint{1.546917in}{1.873677in}}%
\pgfpathlineto{\pgfqpoint{1.543825in}{1.845887in}}%
\pgfpathlineto{\pgfqpoint{1.494987in}{1.851589in}}%
\pgfpathlineto{\pgfqpoint{1.488701in}{1.851762in}}%
\pgfpathlineto{\pgfqpoint{1.475384in}{1.853420in}}%
\pgfpathlineto{\pgfqpoint{1.477206in}{1.868092in}}%
\pgfpathlineto{\pgfqpoint{1.449171in}{1.871656in}}%
\pgfpathlineto{\pgfqpoint{1.453065in}{1.884941in}}%
\pgfpathlineto{\pgfqpoint{1.456350in}{1.910291in}}%
\pgfpathlineto{\pgfqpoint{1.446358in}{1.912489in}}%
\pgfpathlineto{\pgfqpoint{1.446859in}{1.927356in}}%
\pgfpathlineto{\pgfqpoint{1.434747in}{1.935937in}}%
\pgfpathlineto{\pgfqpoint{1.421111in}{1.937794in}}%
\pgfpathlineto{\pgfqpoint{1.422053in}{1.944628in}}%
\pgfpathlineto{\pgfqpoint{1.417308in}{1.945271in}}%
\pgfpathlineto{\pgfqpoint{1.421242in}{1.955525in}}%
\pgfpathlineto{\pgfqpoint{1.414681in}{1.966931in}}%
\pgfpathlineto{\pgfqpoint{1.410750in}{1.980307in}}%
\pgfpathlineto{\pgfqpoint{1.406757in}{1.983167in}}%
\pgfpathlineto{\pgfqpoint{1.409028in}{1.997695in}}%
\pgfpathlineto{\pgfqpoint{1.407997in}{2.003940in}}%
\pgfpathlineto{\pgfqpoint{1.405583in}{2.011993in}}%
\pgfpathlineto{\pgfqpoint{1.469044in}{2.003996in}}%
\pgfpathlineto{\pgfqpoint{1.468884in}{2.002862in}}%
\pgfpathlineto{\pgfqpoint{1.503208in}{1.998589in}}%
\pgfpathlineto{\pgfqpoint{1.501201in}{1.997658in}}%
\pgfpathlineto{\pgfqpoint{1.497843in}{1.970127in}}%
\pgfpathlineto{\pgfqpoint{1.495943in}{1.970351in}}%
\pgfpathlineto{\pgfqpoint{1.492673in}{1.942949in}}%
\pgfpathlineto{\pgfqpoint{1.490580in}{1.943193in}}%
\pgfpathlineto{\pgfqpoint{1.487149in}{1.915660in}}%
\pgfpathlineto{\pgfqpoint{1.553831in}{1.907730in}}%
\pgfpathclose%
\pgfusepath{fill}%
\end{pgfscope}%
\begin{pgfscope}%
\pgfpathrectangle{\pgfqpoint{0.100000in}{0.100000in}}{\pgfqpoint{3.608454in}{2.310000in}}%
\pgfusepath{clip}%
\pgfsetbuttcap%
\pgfsetmiterjoin%
\definecolor{currentfill}{rgb}{0.000000,0.776471,0.611765}%
\pgfsetfillcolor{currentfill}%
\pgfsetlinewidth{0.000000pt}%
\definecolor{currentstroke}{rgb}{0.000000,0.000000,0.000000}%
\pgfsetstrokecolor{currentstroke}%
\pgfsetstrokeopacity{0.000000}%
\pgfsetdash{}{0pt}%
\pgfpathmoveto{\pgfqpoint{2.978722in}{0.804662in}}%
\pgfpathlineto{\pgfqpoint{2.970476in}{0.810318in}}%
\pgfpathlineto{\pgfqpoint{2.953500in}{0.817959in}}%
\pgfpathlineto{\pgfqpoint{2.929182in}{0.813904in}}%
\pgfpathlineto{\pgfqpoint{2.926243in}{0.822499in}}%
\pgfpathlineto{\pgfqpoint{2.917347in}{0.821092in}}%
\pgfpathlineto{\pgfqpoint{2.904439in}{0.824809in}}%
\pgfpathlineto{\pgfqpoint{2.900223in}{0.831517in}}%
\pgfpathlineto{\pgfqpoint{2.895939in}{0.833978in}}%
\pgfpathlineto{\pgfqpoint{2.894217in}{0.842012in}}%
\pgfpathlineto{\pgfqpoint{2.891100in}{0.843787in}}%
\pgfpathlineto{\pgfqpoint{2.893436in}{0.849951in}}%
\pgfpathlineto{\pgfqpoint{2.887187in}{0.854824in}}%
\pgfpathlineto{\pgfqpoint{2.890867in}{0.862703in}}%
\pgfpathlineto{\pgfqpoint{2.898491in}{0.872568in}}%
\pgfpathlineto{\pgfqpoint{2.901077in}{0.870542in}}%
\pgfpathlineto{\pgfqpoint{2.913824in}{0.850086in}}%
\pgfpathlineto{\pgfqpoint{2.922187in}{0.854988in}}%
\pgfpathlineto{\pgfqpoint{2.930056in}{0.865279in}}%
\pgfpathlineto{\pgfqpoint{2.928863in}{0.867621in}}%
\pgfpathlineto{\pgfqpoint{2.932829in}{0.881990in}}%
\pgfpathlineto{\pgfqpoint{2.943431in}{0.882150in}}%
\pgfpathlineto{\pgfqpoint{2.951258in}{0.878733in}}%
\pgfpathlineto{\pgfqpoint{2.950142in}{0.872561in}}%
\pgfpathlineto{\pgfqpoint{2.954746in}{0.867311in}}%
\pgfpathlineto{\pgfqpoint{2.962520in}{0.870538in}}%
\pgfpathlineto{\pgfqpoint{2.967645in}{0.858971in}}%
\pgfpathlineto{\pgfqpoint{2.967579in}{0.838638in}}%
\pgfpathlineto{\pgfqpoint{2.975754in}{0.837437in}}%
\pgfpathlineto{\pgfqpoint{2.979065in}{0.834250in}}%
\pgfpathlineto{\pgfqpoint{2.975294in}{0.828495in}}%
\pgfpathlineto{\pgfqpoint{2.978722in}{0.804662in}}%
\pgfpathclose%
\pgfusepath{fill}%
\end{pgfscope}%
\begin{pgfscope}%
\pgfpathrectangle{\pgfqpoint{0.100000in}{0.100000in}}{\pgfqpoint{3.608454in}{2.310000in}}%
\pgfusepath{clip}%
\pgfsetbuttcap%
\pgfsetmiterjoin%
\definecolor{currentfill}{rgb}{0.000000,0.458824,0.770588}%
\pgfsetfillcolor{currentfill}%
\pgfsetlinewidth{0.000000pt}%
\definecolor{currentstroke}{rgb}{0.000000,0.000000,0.000000}%
\pgfsetstrokecolor{currentstroke}%
\pgfsetstrokeopacity{0.000000}%
\pgfsetdash{}{0pt}%
\pgfpathmoveto{\pgfqpoint{3.269493in}{1.294748in}}%
\pgfpathlineto{\pgfqpoint{3.266535in}{1.293082in}}%
\pgfpathlineto{\pgfqpoint{3.248223in}{1.267423in}}%
\pgfpathlineto{\pgfqpoint{3.248646in}{1.259422in}}%
\pgfpathlineto{\pgfqpoint{3.221614in}{1.254572in}}%
\pgfpathlineto{\pgfqpoint{3.214106in}{1.265774in}}%
\pgfpathlineto{\pgfqpoint{3.240057in}{1.290021in}}%
\pgfpathlineto{\pgfqpoint{3.236491in}{1.297792in}}%
\pgfpathlineto{\pgfqpoint{3.224896in}{1.300661in}}%
\pgfpathlineto{\pgfqpoint{3.234147in}{1.313361in}}%
\pgfpathlineto{\pgfqpoint{3.240056in}{1.312499in}}%
\pgfpathlineto{\pgfqpoint{3.236059in}{1.328225in}}%
\pgfpathlineto{\pgfqpoint{3.249638in}{1.333200in}}%
\pgfpathlineto{\pgfqpoint{3.248164in}{1.344820in}}%
\pgfpathlineto{\pgfqpoint{3.240232in}{1.353321in}}%
\pgfpathlineto{\pgfqpoint{3.244523in}{1.358215in}}%
\pgfpathlineto{\pgfqpoint{3.250135in}{1.356031in}}%
\pgfpathlineto{\pgfqpoint{3.255578in}{1.348500in}}%
\pgfpathlineto{\pgfqpoint{3.262571in}{1.349107in}}%
\pgfpathlineto{\pgfqpoint{3.267940in}{1.341269in}}%
\pgfpathlineto{\pgfqpoint{3.273671in}{1.340685in}}%
\pgfpathlineto{\pgfqpoint{3.276753in}{1.332527in}}%
\pgfpathlineto{\pgfqpoint{3.273954in}{1.327806in}}%
\pgfpathlineto{\pgfqpoint{3.268757in}{1.331100in}}%
\pgfpathlineto{\pgfqpoint{3.270994in}{1.316465in}}%
\pgfpathlineto{\pgfqpoint{3.280061in}{1.309564in}}%
\pgfpathlineto{\pgfqpoint{3.278895in}{1.302449in}}%
\pgfpathlineto{\pgfqpoint{3.274672in}{1.301154in}}%
\pgfpathlineto{\pgfqpoint{3.269493in}{1.294748in}}%
\pgfpathclose%
\pgfusepath{fill}%
\end{pgfscope}%
\begin{pgfscope}%
\pgfpathrectangle{\pgfqpoint{0.100000in}{0.100000in}}{\pgfqpoint{3.608454in}{2.310000in}}%
\pgfusepath{clip}%
\pgfsetbuttcap%
\pgfsetmiterjoin%
\definecolor{currentfill}{rgb}{0.000000,0.447059,0.776471}%
\pgfsetfillcolor{currentfill}%
\pgfsetlinewidth{0.000000pt}%
\definecolor{currentstroke}{rgb}{0.000000,0.000000,0.000000}%
\pgfsetstrokecolor{currentstroke}%
\pgfsetstrokeopacity{0.000000}%
\pgfsetdash{}{0pt}%
\pgfpathmoveto{\pgfqpoint{2.844033in}{0.665474in}}%
\pgfpathlineto{\pgfqpoint{2.844132in}{0.660367in}}%
\pgfpathlineto{\pgfqpoint{2.835832in}{0.661752in}}%
\pgfpathlineto{\pgfqpoint{2.827136in}{0.657037in}}%
\pgfpathlineto{\pgfqpoint{2.809853in}{0.644061in}}%
\pgfpathlineto{\pgfqpoint{2.799185in}{0.640400in}}%
\pgfpathlineto{\pgfqpoint{2.793287in}{0.640173in}}%
\pgfpathlineto{\pgfqpoint{2.783205in}{0.636006in}}%
\pgfpathlineto{\pgfqpoint{2.779514in}{0.637469in}}%
\pgfpathlineto{\pgfqpoint{2.779023in}{0.645884in}}%
\pgfpathlineto{\pgfqpoint{2.769443in}{0.656196in}}%
\pgfpathlineto{\pgfqpoint{2.764564in}{0.656808in}}%
\pgfpathlineto{\pgfqpoint{2.757459in}{0.662818in}}%
\pgfpathlineto{\pgfqpoint{2.742727in}{0.670681in}}%
\pgfpathlineto{\pgfqpoint{2.727447in}{0.677658in}}%
\pgfpathlineto{\pgfqpoint{2.726751in}{0.687235in}}%
\pgfpathlineto{\pgfqpoint{2.735873in}{0.692265in}}%
\pgfpathlineto{\pgfqpoint{2.760979in}{0.694551in}}%
\pgfpathlineto{\pgfqpoint{2.760006in}{0.704933in}}%
\pgfpathlineto{\pgfqpoint{2.766984in}{0.705619in}}%
\pgfpathlineto{\pgfqpoint{2.769691in}{0.676534in}}%
\pgfpathlineto{\pgfqpoint{2.787376in}{0.676778in}}%
\pgfpathlineto{\pgfqpoint{2.788514in}{0.665867in}}%
\pgfpathlineto{\pgfqpoint{2.797575in}{0.664470in}}%
\pgfpathlineto{\pgfqpoint{2.829477in}{0.667841in}}%
\pgfpathlineto{\pgfqpoint{2.844033in}{0.665474in}}%
\pgfpathclose%
\pgfusepath{fill}%
\end{pgfscope}%
\begin{pgfscope}%
\pgfpathrectangle{\pgfqpoint{0.100000in}{0.100000in}}{\pgfqpoint{3.608454in}{2.310000in}}%
\pgfusepath{clip}%
\pgfsetbuttcap%
\pgfsetmiterjoin%
\definecolor{currentfill}{rgb}{0.000000,0.737255,0.631373}%
\pgfsetfillcolor{currentfill}%
\pgfsetlinewidth{0.000000pt}%
\definecolor{currentstroke}{rgb}{0.000000,0.000000,0.000000}%
\pgfsetstrokecolor{currentstroke}%
\pgfsetstrokeopacity{0.000000}%
\pgfsetdash{}{0pt}%
\pgfpathmoveto{\pgfqpoint{3.074241in}{1.252094in}}%
\pgfpathlineto{\pgfqpoint{3.073773in}{1.226802in}}%
\pgfpathlineto{\pgfqpoint{3.054182in}{1.223554in}}%
\pgfpathlineto{\pgfqpoint{3.037358in}{1.220901in}}%
\pgfpathlineto{\pgfqpoint{3.017299in}{1.218774in}}%
\pgfpathlineto{\pgfqpoint{3.015913in}{1.223529in}}%
\pgfpathlineto{\pgfqpoint{3.025017in}{1.229246in}}%
\pgfpathlineto{\pgfqpoint{3.024683in}{1.232048in}}%
\pgfpathlineto{\pgfqpoint{3.032148in}{1.244073in}}%
\pgfpathlineto{\pgfqpoint{3.036908in}{1.247132in}}%
\pgfpathlineto{\pgfqpoint{3.067580in}{1.245522in}}%
\pgfpathlineto{\pgfqpoint{3.074241in}{1.252094in}}%
\pgfpathclose%
\pgfusepath{fill}%
\end{pgfscope}%
\begin{pgfscope}%
\pgfpathrectangle{\pgfqpoint{0.100000in}{0.100000in}}{\pgfqpoint{3.608454in}{2.310000in}}%
\pgfusepath{clip}%
\pgfsetbuttcap%
\pgfsetmiterjoin%
\definecolor{currentfill}{rgb}{0.000000,0.721569,0.639216}%
\pgfsetfillcolor{currentfill}%
\pgfsetlinewidth{0.000000pt}%
\definecolor{currentstroke}{rgb}{0.000000,0.000000,0.000000}%
\pgfsetstrokecolor{currentstroke}%
\pgfsetstrokeopacity{0.000000}%
\pgfsetdash{}{0pt}%
\pgfpathmoveto{\pgfqpoint{1.523829in}{1.020634in}}%
\pgfpathlineto{\pgfqpoint{1.519945in}{0.979398in}}%
\pgfpathlineto{\pgfqpoint{1.506226in}{0.980812in}}%
\pgfpathlineto{\pgfqpoint{1.505497in}{0.973876in}}%
\pgfpathlineto{\pgfqpoint{1.464338in}{0.978551in}}%
\pgfpathlineto{\pgfqpoint{1.465068in}{0.985473in}}%
\pgfpathlineto{\pgfqpoint{1.458090in}{0.986273in}}%
\pgfpathlineto{\pgfqpoint{1.460404in}{1.006662in}}%
\pgfpathlineto{\pgfqpoint{1.432998in}{1.009699in}}%
\pgfpathlineto{\pgfqpoint{1.435358in}{1.030206in}}%
\pgfpathlineto{\pgfqpoint{1.436822in}{1.030070in}}%
\pgfpathlineto{\pgfqpoint{1.440664in}{1.064305in}}%
\pgfpathlineto{\pgfqpoint{1.442195in}{1.078025in}}%
\pgfpathlineto{\pgfqpoint{1.517207in}{1.069928in}}%
\pgfpathlineto{\pgfqpoint{1.516667in}{1.064131in}}%
\pgfpathlineto{\pgfqpoint{1.512139in}{1.021816in}}%
\pgfpathlineto{\pgfqpoint{1.523829in}{1.020634in}}%
\pgfpathclose%
\pgfusepath{fill}%
\end{pgfscope}%
\begin{pgfscope}%
\pgfpathrectangle{\pgfqpoint{0.100000in}{0.100000in}}{\pgfqpoint{3.608454in}{2.310000in}}%
\pgfusepath{clip}%
\pgfsetbuttcap%
\pgfsetmiterjoin%
\definecolor{currentfill}{rgb}{0.000000,0.352941,0.823529}%
\pgfsetfillcolor{currentfill}%
\pgfsetlinewidth{0.000000pt}%
\definecolor{currentstroke}{rgb}{0.000000,0.000000,0.000000}%
\pgfsetstrokecolor{currentstroke}%
\pgfsetstrokeopacity{0.000000}%
\pgfsetdash{}{0pt}%
\pgfpathmoveto{\pgfqpoint{3.530808in}{1.730215in}}%
\pgfpathlineto{\pgfqpoint{3.529206in}{1.730530in}}%
\pgfpathlineto{\pgfqpoint{3.531215in}{1.735382in}}%
\pgfpathlineto{\pgfqpoint{3.526862in}{1.745686in}}%
\pgfpathlineto{\pgfqpoint{3.518613in}{1.736187in}}%
\pgfpathlineto{\pgfqpoint{3.513375in}{1.745585in}}%
\pgfpathlineto{\pgfqpoint{3.499874in}{1.758602in}}%
\pgfpathlineto{\pgfqpoint{3.495434in}{1.768247in}}%
\pgfpathlineto{\pgfqpoint{3.481572in}{1.755490in}}%
\pgfpathlineto{\pgfqpoint{3.479909in}{1.757791in}}%
\pgfpathlineto{\pgfqpoint{3.456367in}{1.750710in}}%
\pgfpathlineto{\pgfqpoint{3.436777in}{1.747450in}}%
\pgfpathlineto{\pgfqpoint{3.434236in}{1.757543in}}%
\pgfpathlineto{\pgfqpoint{3.426466in}{1.757361in}}%
\pgfpathlineto{\pgfqpoint{3.427810in}{1.762986in}}%
\pgfpathlineto{\pgfqpoint{3.420462in}{1.768872in}}%
\pgfpathlineto{\pgfqpoint{3.418010in}{1.787350in}}%
\pgfpathlineto{\pgfqpoint{3.419015in}{1.794601in}}%
\pgfpathlineto{\pgfqpoint{3.414849in}{1.798333in}}%
\pgfpathlineto{\pgfqpoint{3.434992in}{1.802755in}}%
\pgfpathlineto{\pgfqpoint{3.471016in}{1.810988in}}%
\pgfpathlineto{\pgfqpoint{3.477550in}{1.822939in}}%
\pgfpathlineto{\pgfqpoint{3.487553in}{1.830826in}}%
\pgfpathlineto{\pgfqpoint{3.494171in}{1.831598in}}%
\pgfpathlineto{\pgfqpoint{3.501460in}{1.818237in}}%
\pgfpathlineto{\pgfqpoint{3.511382in}{1.816812in}}%
\pgfpathlineto{\pgfqpoint{3.507584in}{1.811929in}}%
\pgfpathlineto{\pgfqpoint{3.498548in}{1.805841in}}%
\pgfpathlineto{\pgfqpoint{3.500415in}{1.803471in}}%
\pgfpathlineto{\pgfqpoint{3.494052in}{1.793429in}}%
\pgfpathlineto{\pgfqpoint{3.496118in}{1.786298in}}%
\pgfpathlineto{\pgfqpoint{3.502227in}{1.787338in}}%
\pgfpathlineto{\pgfqpoint{3.513856in}{1.782235in}}%
\pgfpathlineto{\pgfqpoint{3.520466in}{1.774625in}}%
\pgfpathlineto{\pgfqpoint{3.519312in}{1.765836in}}%
\pgfpathlineto{\pgfqpoint{3.529828in}{1.763255in}}%
\pgfpathlineto{\pgfqpoint{3.531850in}{1.755199in}}%
\pgfpathlineto{\pgfqpoint{3.540059in}{1.751718in}}%
\pgfpathlineto{\pgfqpoint{3.547764in}{1.752941in}}%
\pgfpathlineto{\pgfqpoint{3.566780in}{1.763298in}}%
\pgfpathlineto{\pgfqpoint{3.569480in}{1.754165in}}%
\pgfpathlineto{\pgfqpoint{3.542220in}{1.740804in}}%
\pgfpathlineto{\pgfqpoint{3.540912in}{1.736515in}}%
\pgfpathlineto{\pgfqpoint{3.530808in}{1.730215in}}%
\pgfpathclose%
\pgfusepath{fill}%
\end{pgfscope}%
\begin{pgfscope}%
\pgfpathrectangle{\pgfqpoint{0.100000in}{0.100000in}}{\pgfqpoint{3.608454in}{2.310000in}}%
\pgfusepath{clip}%
\pgfsetbuttcap%
\pgfsetmiterjoin%
\definecolor{currentfill}{rgb}{0.000000,0.827451,0.586275}%
\pgfsetfillcolor{currentfill}%
\pgfsetlinewidth{0.000000pt}%
\definecolor{currentstroke}{rgb}{0.000000,0.000000,0.000000}%
\pgfsetstrokecolor{currentstroke}%
\pgfsetstrokeopacity{0.000000}%
\pgfsetdash{}{0pt}%
\pgfpathmoveto{\pgfqpoint{0.737385in}{0.650268in}}%
\pgfpathlineto{\pgfqpoint{0.738818in}{0.653776in}}%
\pgfpathlineto{\pgfqpoint{0.738884in}{0.657312in}}%
\pgfpathlineto{\pgfqpoint{0.737658in}{0.662474in}}%
\pgfpathlineto{\pgfqpoint{0.738418in}{0.665625in}}%
\pgfpathlineto{\pgfqpoint{0.736741in}{0.671963in}}%
\pgfpathlineto{\pgfqpoint{0.734805in}{0.674605in}}%
\pgfpathlineto{\pgfqpoint{0.729012in}{0.675764in}}%
\pgfpathlineto{\pgfqpoint{0.728310in}{0.678102in}}%
\pgfpathlineto{\pgfqpoint{0.728866in}{0.682342in}}%
\pgfpathlineto{\pgfqpoint{0.731069in}{0.684627in}}%
\pgfpathlineto{\pgfqpoint{0.733927in}{0.685526in}}%
\pgfpathlineto{\pgfqpoint{0.736427in}{0.687461in}}%
\pgfpathlineto{\pgfqpoint{0.737267in}{0.686485in}}%
\pgfpathlineto{\pgfqpoint{0.736221in}{0.683768in}}%
\pgfpathlineto{\pgfqpoint{0.738962in}{0.676360in}}%
\pgfpathlineto{\pgfqpoint{0.740767in}{0.675472in}}%
\pgfpathlineto{\pgfqpoint{0.746631in}{0.673961in}}%
\pgfpathlineto{\pgfqpoint{0.747962in}{0.671116in}}%
\pgfpathlineto{\pgfqpoint{0.747857in}{0.668032in}}%
\pgfpathlineto{\pgfqpoint{0.745828in}{0.665738in}}%
\pgfpathlineto{\pgfqpoint{0.747273in}{0.661310in}}%
\pgfpathlineto{\pgfqpoint{0.747476in}{0.658092in}}%
\pgfpathlineto{\pgfqpoint{0.749575in}{0.653902in}}%
\pgfpathlineto{\pgfqpoint{0.752259in}{0.652120in}}%
\pgfpathlineto{\pgfqpoint{0.753791in}{0.649042in}}%
\pgfpathlineto{\pgfqpoint{0.752833in}{0.647643in}}%
\pgfpathlineto{\pgfqpoint{0.749601in}{0.646731in}}%
\pgfpathlineto{\pgfqpoint{0.748012in}{0.649373in}}%
\pgfpathlineto{\pgfqpoint{0.745554in}{0.651542in}}%
\pgfpathlineto{\pgfqpoint{0.742391in}{0.652099in}}%
\pgfpathlineto{\pgfqpoint{0.737385in}{0.650268in}}%
\pgfpathclose%
\pgfusepath{fill}%
\end{pgfscope}%
\begin{pgfscope}%
\pgfpathrectangle{\pgfqpoint{0.100000in}{0.100000in}}{\pgfqpoint{3.608454in}{2.310000in}}%
\pgfusepath{clip}%
\pgfsetbuttcap%
\pgfsetmiterjoin%
\definecolor{currentfill}{rgb}{0.000000,0.827451,0.586275}%
\pgfsetfillcolor{currentfill}%
\pgfsetlinewidth{0.000000pt}%
\definecolor{currentstroke}{rgb}{0.000000,0.000000,0.000000}%
\pgfsetstrokecolor{currentstroke}%
\pgfsetstrokeopacity{0.000000}%
\pgfsetdash{}{0pt}%
\pgfpathmoveto{\pgfqpoint{0.863943in}{0.679699in}}%
\pgfpathlineto{\pgfqpoint{0.862419in}{0.676643in}}%
\pgfpathlineto{\pgfqpoint{0.855048in}{0.668896in}}%
\pgfpathlineto{\pgfqpoint{0.855680in}{0.668231in}}%
\pgfpathlineto{\pgfqpoint{0.848280in}{0.660517in}}%
\pgfpathlineto{\pgfqpoint{0.850812in}{0.658129in}}%
\pgfpathlineto{\pgfqpoint{0.847118in}{0.654238in}}%
\pgfpathlineto{\pgfqpoint{0.849066in}{0.652382in}}%
\pgfpathlineto{\pgfqpoint{0.845457in}{0.648542in}}%
\pgfpathlineto{\pgfqpoint{0.853614in}{0.640940in}}%
\pgfpathlineto{\pgfqpoint{0.870840in}{0.625356in}}%
\pgfpathlineto{\pgfqpoint{0.882807in}{0.614773in}}%
\pgfpathlineto{\pgfqpoint{0.884580in}{0.616759in}}%
\pgfpathlineto{\pgfqpoint{0.886576in}{0.615036in}}%
\pgfpathlineto{\pgfqpoint{0.888579in}{0.613334in}}%
\pgfpathlineto{\pgfqpoint{0.883388in}{0.607238in}}%
\pgfpathlineto{\pgfqpoint{0.881527in}{0.608877in}}%
\pgfpathlineto{\pgfqpoint{0.874796in}{0.601023in}}%
\pgfpathlineto{\pgfqpoint{0.872825in}{0.602697in}}%
\pgfpathlineto{\pgfqpoint{0.869420in}{0.598690in}}%
\pgfpathlineto{\pgfqpoint{0.867382in}{0.600383in}}%
\pgfpathlineto{\pgfqpoint{0.853750in}{0.584249in}}%
\pgfpathlineto{\pgfqpoint{0.856024in}{0.582334in}}%
\pgfpathlineto{\pgfqpoint{0.850691in}{0.576622in}}%
\pgfpathlineto{\pgfqpoint{0.848929in}{0.578173in}}%
\pgfpathlineto{\pgfqpoint{0.847165in}{0.576109in}}%
\pgfpathlineto{\pgfqpoint{0.845387in}{0.577604in}}%
\pgfpathlineto{\pgfqpoint{0.843712in}{0.575563in}}%
\pgfpathlineto{\pgfqpoint{0.839724in}{0.579022in}}%
\pgfpathlineto{\pgfqpoint{0.838193in}{0.577309in}}%
\pgfpathlineto{\pgfqpoint{0.836205in}{0.579038in}}%
\pgfpathlineto{\pgfqpoint{0.831146in}{0.572881in}}%
\pgfpathlineto{\pgfqpoint{0.829339in}{0.574413in}}%
\pgfpathlineto{\pgfqpoint{0.824219in}{0.568707in}}%
\pgfpathlineto{\pgfqpoint{0.812726in}{0.579011in}}%
\pgfpathlineto{\pgfqpoint{0.814280in}{0.580959in}}%
\pgfpathlineto{\pgfqpoint{0.812542in}{0.582561in}}%
\pgfpathlineto{\pgfqpoint{0.814282in}{0.584612in}}%
\pgfpathlineto{\pgfqpoint{0.812396in}{0.586335in}}%
\pgfpathlineto{\pgfqpoint{0.814106in}{0.588275in}}%
\pgfpathlineto{\pgfqpoint{0.812082in}{0.590252in}}%
\pgfpathlineto{\pgfqpoint{0.818724in}{0.592259in}}%
\pgfpathlineto{\pgfqpoint{0.818875in}{0.594792in}}%
\pgfpathlineto{\pgfqpoint{0.816853in}{0.596475in}}%
\pgfpathlineto{\pgfqpoint{0.818540in}{0.598214in}}%
\pgfpathlineto{\pgfqpoint{0.821571in}{0.596416in}}%
\pgfpathlineto{\pgfqpoint{0.820998in}{0.592054in}}%
\pgfpathlineto{\pgfqpoint{0.823073in}{0.586861in}}%
\pgfpathlineto{\pgfqpoint{0.829242in}{0.582849in}}%
\pgfpathlineto{\pgfqpoint{0.831926in}{0.582649in}}%
\pgfpathlineto{\pgfqpoint{0.836286in}{0.583592in}}%
\pgfpathlineto{\pgfqpoint{0.839848in}{0.584937in}}%
\pgfpathlineto{\pgfqpoint{0.841321in}{0.586303in}}%
\pgfpathlineto{\pgfqpoint{0.844401in}{0.593367in}}%
\pgfpathlineto{\pgfqpoint{0.847330in}{0.597336in}}%
\pgfpathlineto{\pgfqpoint{0.848129in}{0.603625in}}%
\pgfpathlineto{\pgfqpoint{0.845939in}{0.606018in}}%
\pgfpathlineto{\pgfqpoint{0.848681in}{0.607604in}}%
\pgfpathlineto{\pgfqpoint{0.851564in}{0.603948in}}%
\pgfpathlineto{\pgfqpoint{0.856618in}{0.604295in}}%
\pgfpathlineto{\pgfqpoint{0.858912in}{0.606588in}}%
\pgfpathlineto{\pgfqpoint{0.860360in}{0.609378in}}%
\pgfpathlineto{\pgfqpoint{0.859005in}{0.614759in}}%
\pgfpathlineto{\pgfqpoint{0.854769in}{0.613300in}}%
\pgfpathlineto{\pgfqpoint{0.852074in}{0.613498in}}%
\pgfpathlineto{\pgfqpoint{0.851004in}{0.616189in}}%
\pgfpathlineto{\pgfqpoint{0.847338in}{0.616442in}}%
\pgfpathlineto{\pgfqpoint{0.843639in}{0.617453in}}%
\pgfpathlineto{\pgfqpoint{0.840644in}{0.617353in}}%
\pgfpathlineto{\pgfqpoint{0.831603in}{0.615210in}}%
\pgfpathlineto{\pgfqpoint{0.833848in}{0.619214in}}%
\pgfpathlineto{\pgfqpoint{0.834393in}{0.625065in}}%
\pgfpathlineto{\pgfqpoint{0.831455in}{0.625440in}}%
\pgfpathlineto{\pgfqpoint{0.831880in}{0.622625in}}%
\pgfpathlineto{\pgfqpoint{0.829678in}{0.621808in}}%
\pgfpathlineto{\pgfqpoint{0.829726in}{0.624782in}}%
\pgfpathlineto{\pgfqpoint{0.828624in}{0.627896in}}%
\pgfpathlineto{\pgfqpoint{0.822394in}{0.633936in}}%
\pgfpathlineto{\pgfqpoint{0.815333in}{0.636154in}}%
\pgfpathlineto{\pgfqpoint{0.813484in}{0.637133in}}%
\pgfpathlineto{\pgfqpoint{0.805786in}{0.649884in}}%
\pgfpathlineto{\pgfqpoint{0.804377in}{0.652802in}}%
\pgfpathlineto{\pgfqpoint{0.804337in}{0.655328in}}%
\pgfpathlineto{\pgfqpoint{0.806949in}{0.658877in}}%
\pgfpathlineto{\pgfqpoint{0.809295in}{0.660793in}}%
\pgfpathlineto{\pgfqpoint{0.809700in}{0.668848in}}%
\pgfpathlineto{\pgfqpoint{0.813648in}{0.667058in}}%
\pgfpathlineto{\pgfqpoint{0.818571in}{0.669230in}}%
\pgfpathlineto{\pgfqpoint{0.813780in}{0.677479in}}%
\pgfpathlineto{\pgfqpoint{0.811754in}{0.680104in}}%
\pgfpathlineto{\pgfqpoint{0.810600in}{0.686441in}}%
\pgfpathlineto{\pgfqpoint{0.809362in}{0.690533in}}%
\pgfpathlineto{\pgfqpoint{0.811858in}{0.693024in}}%
\pgfpathlineto{\pgfqpoint{0.822834in}{0.689897in}}%
\pgfpathlineto{\pgfqpoint{0.823528in}{0.687901in}}%
\pgfpathlineto{\pgfqpoint{0.827552in}{0.687211in}}%
\pgfpathlineto{\pgfqpoint{0.827406in}{0.689342in}}%
\pgfpathlineto{\pgfqpoint{0.838119in}{0.687842in}}%
\pgfpathlineto{\pgfqpoint{0.840746in}{0.682200in}}%
\pgfpathlineto{\pgfqpoint{0.842608in}{0.680338in}}%
\pgfpathlineto{\pgfqpoint{0.844777in}{0.680608in}}%
\pgfpathlineto{\pgfqpoint{0.843442in}{0.683454in}}%
\pgfpathlineto{\pgfqpoint{0.846822in}{0.686101in}}%
\pgfpathlineto{\pgfqpoint{0.855456in}{0.683423in}}%
\pgfpathlineto{\pgfqpoint{0.860934in}{0.681377in}}%
\pgfpathlineto{\pgfqpoint{0.863943in}{0.679699in}}%
\pgfpathclose%
\pgfusepath{fill}%
\end{pgfscope}%
\begin{pgfscope}%
\pgfpathrectangle{\pgfqpoint{0.100000in}{0.100000in}}{\pgfqpoint{3.608454in}{2.310000in}}%
\pgfusepath{clip}%
\pgfsetbuttcap%
\pgfsetmiterjoin%
\definecolor{currentfill}{rgb}{0.000000,0.796078,0.601961}%
\pgfsetfillcolor{currentfill}%
\pgfsetlinewidth{0.000000pt}%
\definecolor{currentstroke}{rgb}{0.000000,0.000000,0.000000}%
\pgfsetstrokecolor{currentstroke}%
\pgfsetstrokeopacity{0.000000}%
\pgfsetdash{}{0pt}%
\pgfpathmoveto{\pgfqpoint{1.805743in}{1.664546in}}%
\pgfpathlineto{\pgfqpoint{1.746207in}{1.668241in}}%
\pgfpathlineto{\pgfqpoint{1.748309in}{1.699187in}}%
\pgfpathlineto{\pgfqpoint{1.750543in}{1.731001in}}%
\pgfpathlineto{\pgfqpoint{1.760148in}{1.734344in}}%
\pgfpathlineto{\pgfqpoint{1.768321in}{1.732235in}}%
\pgfpathlineto{\pgfqpoint{1.775409in}{1.724494in}}%
\pgfpathlineto{\pgfqpoint{1.787775in}{1.725501in}}%
\pgfpathlineto{\pgfqpoint{1.793255in}{1.722620in}}%
\pgfpathlineto{\pgfqpoint{1.800957in}{1.721865in}}%
\pgfpathlineto{\pgfqpoint{1.807130in}{1.721260in}}%
\pgfpathlineto{\pgfqpoint{1.806996in}{1.702523in}}%
\pgfpathlineto{\pgfqpoint{1.805445in}{1.674881in}}%
\pgfpathlineto{\pgfqpoint{1.805743in}{1.664546in}}%
\pgfpathclose%
\pgfusepath{fill}%
\end{pgfscope}%
\begin{pgfscope}%
\pgfpathrectangle{\pgfqpoint{0.100000in}{0.100000in}}{\pgfqpoint{3.608454in}{2.310000in}}%
\pgfusepath{clip}%
\pgfsetbuttcap%
\pgfsetmiterjoin%
\definecolor{currentfill}{rgb}{0.000000,0.607843,0.696078}%
\pgfsetfillcolor{currentfill}%
\pgfsetlinewidth{0.000000pt}%
\definecolor{currentstroke}{rgb}{0.000000,0.000000,0.000000}%
\pgfsetstrokecolor{currentstroke}%
\pgfsetstrokeopacity{0.000000}%
\pgfsetdash{}{0pt}%
\pgfpathmoveto{\pgfqpoint{0.895815in}{1.966858in}}%
\pgfpathlineto{\pgfqpoint{0.891257in}{1.963356in}}%
\pgfpathlineto{\pgfqpoint{0.890375in}{1.956866in}}%
\pgfpathlineto{\pgfqpoint{0.884981in}{1.950005in}}%
\pgfpathlineto{\pgfqpoint{0.875640in}{1.944745in}}%
\pgfpathlineto{\pgfqpoint{0.866297in}{1.932367in}}%
\pgfpathlineto{\pgfqpoint{0.860164in}{1.926417in}}%
\pgfpathlineto{\pgfqpoint{0.857097in}{1.912795in}}%
\pgfpathlineto{\pgfqpoint{0.842438in}{1.916478in}}%
\pgfpathlineto{\pgfqpoint{0.844130in}{1.923115in}}%
\pgfpathlineto{\pgfqpoint{0.839459in}{1.928946in}}%
\pgfpathlineto{\pgfqpoint{0.818538in}{1.934210in}}%
\pgfpathlineto{\pgfqpoint{0.814693in}{1.932896in}}%
\pgfpathlineto{\pgfqpoint{0.805150in}{1.923524in}}%
\pgfpathlineto{\pgfqpoint{0.800486in}{1.923543in}}%
\pgfpathlineto{\pgfqpoint{0.785890in}{1.927298in}}%
\pgfpathlineto{\pgfqpoint{0.792882in}{1.935165in}}%
\pgfpathlineto{\pgfqpoint{0.793934in}{1.941045in}}%
\pgfpathlineto{\pgfqpoint{0.803077in}{1.950037in}}%
\pgfpathlineto{\pgfqpoint{0.801936in}{1.953999in}}%
\pgfpathlineto{\pgfqpoint{0.793832in}{1.958919in}}%
\pgfpathlineto{\pgfqpoint{0.793910in}{1.961931in}}%
\pgfpathlineto{\pgfqpoint{0.807519in}{1.962306in}}%
\pgfpathlineto{\pgfqpoint{0.807404in}{1.969358in}}%
\pgfpathlineto{\pgfqpoint{0.813336in}{1.971025in}}%
\pgfpathlineto{\pgfqpoint{0.814622in}{1.975723in}}%
\pgfpathlineto{\pgfqpoint{0.808158in}{1.981911in}}%
\pgfpathlineto{\pgfqpoint{0.799750in}{1.984499in}}%
\pgfpathlineto{\pgfqpoint{0.802344in}{2.000341in}}%
\pgfpathlineto{\pgfqpoint{0.796435in}{2.001902in}}%
\pgfpathlineto{\pgfqpoint{0.797207in}{2.013901in}}%
\pgfpathlineto{\pgfqpoint{0.811964in}{2.010809in}}%
\pgfpathlineto{\pgfqpoint{0.813450in}{2.016488in}}%
\pgfpathlineto{\pgfqpoint{0.825859in}{2.013085in}}%
\pgfpathlineto{\pgfqpoint{0.831103in}{2.015303in}}%
\pgfpathlineto{\pgfqpoint{0.835457in}{2.032072in}}%
\pgfpathlineto{\pgfqpoint{0.838162in}{2.031358in}}%
\pgfpathlineto{\pgfqpoint{0.841907in}{2.040973in}}%
\pgfpathlineto{\pgfqpoint{0.846645in}{2.043323in}}%
\pgfpathlineto{\pgfqpoint{0.858767in}{2.040242in}}%
\pgfpathlineto{\pgfqpoint{0.853215in}{2.027577in}}%
\pgfpathlineto{\pgfqpoint{0.855049in}{2.021169in}}%
\pgfpathlineto{\pgfqpoint{0.853023in}{2.013264in}}%
\pgfpathlineto{\pgfqpoint{0.854053in}{1.998430in}}%
\pgfpathlineto{\pgfqpoint{0.859301in}{1.991402in}}%
\pgfpathlineto{\pgfqpoint{0.859744in}{1.982888in}}%
\pgfpathlineto{\pgfqpoint{0.871160in}{1.980064in}}%
\pgfpathlineto{\pgfqpoint{0.869552in}{1.973496in}}%
\pgfpathlineto{\pgfqpoint{0.895815in}{1.966858in}}%
\pgfpathclose%
\pgfusepath{fill}%
\end{pgfscope}%
\begin{pgfscope}%
\pgfpathrectangle{\pgfqpoint{0.100000in}{0.100000in}}{\pgfqpoint{3.608454in}{2.310000in}}%
\pgfusepath{clip}%
\pgfsetbuttcap%
\pgfsetmiterjoin%
\definecolor{currentfill}{rgb}{0.000000,0.364706,0.817647}%
\pgfsetfillcolor{currentfill}%
\pgfsetlinewidth{0.000000pt}%
\definecolor{currentstroke}{rgb}{0.000000,0.000000,0.000000}%
\pgfsetstrokecolor{currentstroke}%
\pgfsetstrokeopacity{0.000000}%
\pgfsetdash{}{0pt}%
\pgfpathmoveto{\pgfqpoint{1.791350in}{1.359031in}}%
\pgfpathlineto{\pgfqpoint{1.789033in}{1.324734in}}%
\pgfpathlineto{\pgfqpoint{1.756162in}{1.326912in}}%
\pgfpathlineto{\pgfqpoint{1.707133in}{1.330203in}}%
\pgfpathlineto{\pgfqpoint{1.709976in}{1.364447in}}%
\pgfpathlineto{\pgfqpoint{1.715265in}{1.364101in}}%
\pgfpathlineto{\pgfqpoint{1.717833in}{1.398242in}}%
\pgfpathlineto{\pgfqpoint{1.791324in}{1.393448in}}%
\pgfpathlineto{\pgfqpoint{1.792462in}{1.393414in}}%
\pgfpathlineto{\pgfqpoint{1.790342in}{1.359067in}}%
\pgfpathlineto{\pgfqpoint{1.791350in}{1.359031in}}%
\pgfpathclose%
\pgfusepath{fill}%
\end{pgfscope}%
\begin{pgfscope}%
\pgfpathrectangle{\pgfqpoint{0.100000in}{0.100000in}}{\pgfqpoint{3.608454in}{2.310000in}}%
\pgfusepath{clip}%
\pgfsetbuttcap%
\pgfsetmiterjoin%
\definecolor{currentfill}{rgb}{0.000000,0.486275,0.756863}%
\pgfsetfillcolor{currentfill}%
\pgfsetlinewidth{0.000000pt}%
\definecolor{currentstroke}{rgb}{0.000000,0.000000,0.000000}%
\pgfsetstrokecolor{currentstroke}%
\pgfsetstrokeopacity{0.000000}%
\pgfsetdash{}{0pt}%
\pgfpathmoveto{\pgfqpoint{2.464725in}{1.378777in}}%
\pgfpathlineto{\pgfqpoint{2.440862in}{1.377267in}}%
\pgfpathlineto{\pgfqpoint{2.439852in}{1.391129in}}%
\pgfpathlineto{\pgfqpoint{2.415980in}{1.389550in}}%
\pgfpathlineto{\pgfqpoint{2.410721in}{1.410092in}}%
\pgfpathlineto{\pgfqpoint{2.409257in}{1.435587in}}%
\pgfpathlineto{\pgfqpoint{2.410568in}{1.438329in}}%
\pgfpathlineto{\pgfqpoint{2.418888in}{1.437396in}}%
\pgfpathlineto{\pgfqpoint{2.427687in}{1.440218in}}%
\pgfpathlineto{\pgfqpoint{2.432920in}{1.438056in}}%
\pgfpathlineto{\pgfqpoint{2.431914in}{1.453643in}}%
\pgfpathlineto{\pgfqpoint{2.452287in}{1.455302in}}%
\pgfpathlineto{\pgfqpoint{2.459428in}{1.452340in}}%
\pgfpathlineto{\pgfqpoint{2.460863in}{1.433950in}}%
\pgfpathlineto{\pgfqpoint{2.461611in}{1.423638in}}%
\pgfpathlineto{\pgfqpoint{2.457108in}{1.423280in}}%
\pgfpathlineto{\pgfqpoint{2.457647in}{1.415090in}}%
\pgfpathlineto{\pgfqpoint{2.462297in}{1.414401in}}%
\pgfpathlineto{\pgfqpoint{2.463205in}{1.402934in}}%
\pgfpathlineto{\pgfqpoint{2.470095in}{1.403294in}}%
\pgfpathlineto{\pgfqpoint{2.471713in}{1.379003in}}%
\pgfpathlineto{\pgfqpoint{2.464725in}{1.378777in}}%
\pgfpathclose%
\pgfusepath{fill}%
\end{pgfscope}%
\begin{pgfscope}%
\pgfpathrectangle{\pgfqpoint{0.100000in}{0.100000in}}{\pgfqpoint{3.608454in}{2.310000in}}%
\pgfusepath{clip}%
\pgfsetbuttcap%
\pgfsetmiterjoin%
\definecolor{currentfill}{rgb}{0.000000,0.556863,0.721569}%
\pgfsetfillcolor{currentfill}%
\pgfsetlinewidth{0.000000pt}%
\definecolor{currentstroke}{rgb}{0.000000,0.000000,0.000000}%
\pgfsetstrokecolor{currentstroke}%
\pgfsetstrokeopacity{0.000000}%
\pgfsetdash{}{0pt}%
\pgfpathmoveto{\pgfqpoint{3.344220in}{1.667169in}}%
\pgfpathlineto{\pgfqpoint{3.332839in}{1.649737in}}%
\pgfpathlineto{\pgfqpoint{3.302353in}{1.660141in}}%
\pgfpathlineto{\pgfqpoint{3.298905in}{1.664653in}}%
\pgfpathlineto{\pgfqpoint{3.292838in}{1.664005in}}%
\pgfpathlineto{\pgfqpoint{3.282783in}{1.667998in}}%
\pgfpathlineto{\pgfqpoint{3.277583in}{1.676067in}}%
\pgfpathlineto{\pgfqpoint{3.274806in}{1.685909in}}%
\pgfpathlineto{\pgfqpoint{3.267733in}{1.692380in}}%
\pgfpathlineto{\pgfqpoint{3.302028in}{1.725960in}}%
\pgfpathlineto{\pgfqpoint{3.311248in}{1.723631in}}%
\pgfpathlineto{\pgfqpoint{3.324878in}{1.725364in}}%
\pgfpathlineto{\pgfqpoint{3.325396in}{1.731486in}}%
\pgfpathlineto{\pgfqpoint{3.333233in}{1.729959in}}%
\pgfpathlineto{\pgfqpoint{3.335015in}{1.726252in}}%
\pgfpathlineto{\pgfqpoint{3.347348in}{1.723269in}}%
\pgfpathlineto{\pgfqpoint{3.358373in}{1.723673in}}%
\pgfpathlineto{\pgfqpoint{3.358769in}{1.729605in}}%
\pgfpathlineto{\pgfqpoint{3.359316in}{1.729735in}}%
\pgfpathlineto{\pgfqpoint{3.366550in}{1.688953in}}%
\pgfpathlineto{\pgfqpoint{3.344000in}{1.680413in}}%
\pgfpathlineto{\pgfqpoint{3.344220in}{1.667169in}}%
\pgfpathclose%
\pgfusepath{fill}%
\end{pgfscope}%
\begin{pgfscope}%
\pgfpathrectangle{\pgfqpoint{0.100000in}{0.100000in}}{\pgfqpoint{3.608454in}{2.310000in}}%
\pgfusepath{clip}%
\pgfsetbuttcap%
\pgfsetmiterjoin%
\definecolor{currentfill}{rgb}{0.000000,0.674510,0.662745}%
\pgfsetfillcolor{currentfill}%
\pgfsetlinewidth{0.000000pt}%
\definecolor{currentstroke}{rgb}{0.000000,0.000000,0.000000}%
\pgfsetstrokecolor{currentstroke}%
\pgfsetstrokeopacity{0.000000}%
\pgfsetdash{}{0pt}%
\pgfpathmoveto{\pgfqpoint{1.443108in}{0.820998in}}%
\pgfpathlineto{\pgfqpoint{1.438374in}{0.821513in}}%
\pgfpathlineto{\pgfqpoint{1.404085in}{0.825375in}}%
\pgfpathlineto{\pgfqpoint{1.340564in}{0.833080in}}%
\pgfpathlineto{\pgfqpoint{1.350085in}{0.908598in}}%
\pgfpathlineto{\pgfqpoint{1.353414in}{0.915777in}}%
\pgfpathlineto{\pgfqpoint{1.356403in}{0.942467in}}%
\pgfpathlineto{\pgfqpoint{1.355496in}{0.949729in}}%
\pgfpathlineto{\pgfqpoint{1.357228in}{0.963939in}}%
\pgfpathlineto{\pgfqpoint{1.378188in}{0.960531in}}%
\pgfpathlineto{\pgfqpoint{1.379878in}{0.974364in}}%
\pgfpathlineto{\pgfqpoint{1.388260in}{0.973319in}}%
\pgfpathlineto{\pgfqpoint{1.392442in}{1.007538in}}%
\pgfpathlineto{\pgfqpoint{1.432205in}{1.002872in}}%
\pgfpathlineto{\pgfqpoint{1.432998in}{1.009699in}}%
\pgfpathlineto{\pgfqpoint{1.460404in}{1.006662in}}%
\pgfpathlineto{\pgfqpoint{1.458090in}{0.986273in}}%
\pgfpathlineto{\pgfqpoint{1.452184in}{0.931773in}}%
\pgfpathlineto{\pgfqpoint{1.449063in}{0.911411in}}%
\pgfpathlineto{\pgfqpoint{1.421822in}{0.913943in}}%
\pgfpathlineto{\pgfqpoint{1.420339in}{0.900498in}}%
\pgfpathlineto{\pgfqpoint{1.417927in}{0.900770in}}%
\pgfpathlineto{\pgfqpoint{1.413830in}{0.865768in}}%
\pgfpathlineto{\pgfqpoint{1.447334in}{0.862167in}}%
\pgfpathlineto{\pgfqpoint{1.443108in}{0.820998in}}%
\pgfpathclose%
\pgfusepath{fill}%
\end{pgfscope}%
\begin{pgfscope}%
\pgfpathrectangle{\pgfqpoint{0.100000in}{0.100000in}}{\pgfqpoint{3.608454in}{2.310000in}}%
\pgfusepath{clip}%
\pgfsetbuttcap%
\pgfsetmiterjoin%
\definecolor{currentfill}{rgb}{0.000000,0.627451,0.686275}%
\pgfsetfillcolor{currentfill}%
\pgfsetlinewidth{0.000000pt}%
\definecolor{currentstroke}{rgb}{0.000000,0.000000,0.000000}%
\pgfsetstrokecolor{currentstroke}%
\pgfsetstrokeopacity{0.000000}%
\pgfsetdash{}{0pt}%
\pgfpathmoveto{\pgfqpoint{2.853827in}{0.823500in}}%
\pgfpathlineto{\pgfqpoint{2.808283in}{0.818788in}}%
\pgfpathlineto{\pgfqpoint{2.786918in}{0.816730in}}%
\pgfpathlineto{\pgfqpoint{2.778092in}{0.819839in}}%
\pgfpathlineto{\pgfqpoint{2.776392in}{0.826653in}}%
\pgfpathlineto{\pgfqpoint{2.775432in}{0.832235in}}%
\pgfpathlineto{\pgfqpoint{2.783833in}{0.840916in}}%
\pgfpathlineto{\pgfqpoint{2.801485in}{0.843067in}}%
\pgfpathlineto{\pgfqpoint{2.799223in}{0.863800in}}%
\pgfpathlineto{\pgfqpoint{2.800195in}{0.867107in}}%
\pgfpathlineto{\pgfqpoint{2.812799in}{0.870744in}}%
\pgfpathlineto{\pgfqpoint{2.816271in}{0.861716in}}%
\pgfpathlineto{\pgfqpoint{2.822592in}{0.856703in}}%
\pgfpathlineto{\pgfqpoint{2.830726in}{0.859944in}}%
\pgfpathlineto{\pgfqpoint{2.834772in}{0.868856in}}%
\pgfpathlineto{\pgfqpoint{2.839212in}{0.870674in}}%
\pgfpathlineto{\pgfqpoint{2.851231in}{0.869522in}}%
\pgfpathlineto{\pgfqpoint{2.854444in}{0.865881in}}%
\pgfpathlineto{\pgfqpoint{2.855107in}{0.854082in}}%
\pgfpathlineto{\pgfqpoint{2.847916in}{0.849972in}}%
\pgfpathlineto{\pgfqpoint{2.844094in}{0.841666in}}%
\pgfpathlineto{\pgfqpoint{2.847520in}{0.837484in}}%
\pgfpathlineto{\pgfqpoint{2.853827in}{0.823500in}}%
\pgfpathclose%
\pgfusepath{fill}%
\end{pgfscope}%
\begin{pgfscope}%
\pgfpathrectangle{\pgfqpoint{0.100000in}{0.100000in}}{\pgfqpoint{3.608454in}{2.310000in}}%
\pgfusepath{clip}%
\pgfsetbuttcap%
\pgfsetmiterjoin%
\definecolor{currentfill}{rgb}{0.000000,0.843137,0.578431}%
\pgfsetfillcolor{currentfill}%
\pgfsetlinewidth{0.000000pt}%
\definecolor{currentstroke}{rgb}{0.000000,0.000000,0.000000}%
\pgfsetstrokecolor{currentstroke}%
\pgfsetstrokeopacity{0.000000}%
\pgfsetdash{}{0pt}%
\pgfpathmoveto{\pgfqpoint{2.756041in}{1.789197in}}%
\pgfpathlineto{\pgfqpoint{2.738013in}{1.786827in}}%
\pgfpathlineto{\pgfqpoint{2.729298in}{1.855252in}}%
\pgfpathlineto{\pgfqpoint{2.728097in}{1.868854in}}%
\pgfpathlineto{\pgfqpoint{2.707810in}{1.866142in}}%
\pgfpathlineto{\pgfqpoint{2.703507in}{1.899857in}}%
\pgfpathlineto{\pgfqpoint{2.706092in}{1.900170in}}%
\pgfpathlineto{\pgfqpoint{2.712457in}{1.891199in}}%
\pgfpathlineto{\pgfqpoint{2.723164in}{1.891127in}}%
\pgfpathlineto{\pgfqpoint{2.731458in}{1.886319in}}%
\pgfpathlineto{\pgfqpoint{2.749951in}{1.882041in}}%
\pgfpathlineto{\pgfqpoint{2.753811in}{1.874987in}}%
\pgfpathlineto{\pgfqpoint{2.760125in}{1.868000in}}%
\pgfpathlineto{\pgfqpoint{2.764215in}{1.858878in}}%
\pgfpathlineto{\pgfqpoint{2.754329in}{1.860433in}}%
\pgfpathlineto{\pgfqpoint{2.755769in}{1.850748in}}%
\pgfpathlineto{\pgfqpoint{2.762952in}{1.847681in}}%
\pgfpathlineto{\pgfqpoint{2.767020in}{1.834435in}}%
\pgfpathlineto{\pgfqpoint{2.765648in}{1.826069in}}%
\pgfpathlineto{\pgfqpoint{2.767351in}{1.804576in}}%
\pgfpathlineto{\pgfqpoint{2.762782in}{1.798736in}}%
\pgfpathlineto{\pgfqpoint{2.757883in}{1.797924in}}%
\pgfpathlineto{\pgfqpoint{2.756041in}{1.789197in}}%
\pgfpathclose%
\pgfusepath{fill}%
\end{pgfscope}%
\begin{pgfscope}%
\pgfpathrectangle{\pgfqpoint{0.100000in}{0.100000in}}{\pgfqpoint{3.608454in}{2.310000in}}%
\pgfusepath{clip}%
\pgfsetbuttcap%
\pgfsetmiterjoin%
\definecolor{currentfill}{rgb}{0.000000,0.513725,0.743137}%
\pgfsetfillcolor{currentfill}%
\pgfsetlinewidth{0.000000pt}%
\definecolor{currentstroke}{rgb}{0.000000,0.000000,0.000000}%
\pgfsetstrokecolor{currentstroke}%
\pgfsetstrokeopacity{0.000000}%
\pgfsetdash{}{0pt}%
\pgfpathmoveto{\pgfqpoint{2.773279in}{1.502104in}}%
\pgfpathlineto{\pgfqpoint{2.766536in}{1.500999in}}%
\pgfpathlineto{\pgfqpoint{2.767029in}{1.496892in}}%
\pgfpathlineto{\pgfqpoint{2.746418in}{1.494296in}}%
\pgfpathlineto{\pgfqpoint{2.747392in}{1.486269in}}%
\pgfpathlineto{\pgfqpoint{2.741875in}{1.483652in}}%
\pgfpathlineto{\pgfqpoint{2.719775in}{1.481004in}}%
\pgfpathlineto{\pgfqpoint{2.714189in}{1.531011in}}%
\pgfpathlineto{\pgfqpoint{2.741540in}{1.534284in}}%
\pgfpathlineto{\pgfqpoint{2.739953in}{1.548013in}}%
\pgfpathlineto{\pgfqpoint{2.746637in}{1.548837in}}%
\pgfpathlineto{\pgfqpoint{2.767138in}{1.551464in}}%
\pgfpathlineto{\pgfqpoint{2.773279in}{1.502104in}}%
\pgfpathclose%
\pgfusepath{fill}%
\end{pgfscope}%
\begin{pgfscope}%
\pgfpathrectangle{\pgfqpoint{0.100000in}{0.100000in}}{\pgfqpoint{3.608454in}{2.310000in}}%
\pgfusepath{clip}%
\pgfsetbuttcap%
\pgfsetmiterjoin%
\definecolor{currentfill}{rgb}{0.000000,0.768627,0.615686}%
\pgfsetfillcolor{currentfill}%
\pgfsetlinewidth{0.000000pt}%
\definecolor{currentstroke}{rgb}{0.000000,0.000000,0.000000}%
\pgfsetstrokecolor{currentstroke}%
\pgfsetstrokeopacity{0.000000}%
\pgfsetdash{}{0pt}%
\pgfpathmoveto{\pgfqpoint{3.130816in}{1.638979in}}%
\pgfpathlineto{\pgfqpoint{3.131860in}{1.634049in}}%
\pgfpathlineto{\pgfqpoint{3.109153in}{1.629395in}}%
\pgfpathlineto{\pgfqpoint{3.105576in}{1.628770in}}%
\pgfpathlineto{\pgfqpoint{3.094591in}{1.638194in}}%
\pgfpathlineto{\pgfqpoint{3.078140in}{1.636089in}}%
\pgfpathlineto{\pgfqpoint{3.050721in}{1.630639in}}%
\pgfpathlineto{\pgfqpoint{3.048278in}{1.646809in}}%
\pgfpathlineto{\pgfqpoint{3.050043in}{1.647962in}}%
\pgfpathlineto{\pgfqpoint{3.047674in}{1.660178in}}%
\pgfpathlineto{\pgfqpoint{3.039413in}{1.658799in}}%
\pgfpathlineto{\pgfqpoint{3.032000in}{1.700788in}}%
\pgfpathlineto{\pgfqpoint{3.036017in}{1.700877in}}%
\pgfpathlineto{\pgfqpoint{3.042918in}{1.694827in}}%
\pgfpathlineto{\pgfqpoint{3.049340in}{1.697252in}}%
\pgfpathlineto{\pgfqpoint{3.064906in}{1.706470in}}%
\pgfpathlineto{\pgfqpoint{3.066495in}{1.705513in}}%
\pgfpathlineto{\pgfqpoint{3.102306in}{1.712116in}}%
\pgfpathlineto{\pgfqpoint{3.109751in}{1.709791in}}%
\pgfpathlineto{\pgfqpoint{3.115261in}{1.672766in}}%
\pgfpathlineto{\pgfqpoint{3.123343in}{1.674373in}}%
\pgfpathlineto{\pgfqpoint{3.130816in}{1.638979in}}%
\pgfpathclose%
\pgfusepath{fill}%
\end{pgfscope}%
\begin{pgfscope}%
\pgfpathrectangle{\pgfqpoint{0.100000in}{0.100000in}}{\pgfqpoint{3.608454in}{2.310000in}}%
\pgfusepath{clip}%
\pgfsetbuttcap%
\pgfsetmiterjoin%
\definecolor{currentfill}{rgb}{0.000000,0.552941,0.723529}%
\pgfsetfillcolor{currentfill}%
\pgfsetlinewidth{0.000000pt}%
\definecolor{currentstroke}{rgb}{0.000000,0.000000,0.000000}%
\pgfsetstrokecolor{currentstroke}%
\pgfsetstrokeopacity{0.000000}%
\pgfsetdash{}{0pt}%
\pgfpathmoveto{\pgfqpoint{2.682759in}{1.669371in}}%
\pgfpathlineto{\pgfqpoint{2.696454in}{1.670883in}}%
\pgfpathlineto{\pgfqpoint{2.693396in}{1.698299in}}%
\pgfpathlineto{\pgfqpoint{2.720350in}{1.701323in}}%
\pgfpathlineto{\pgfqpoint{2.723746in}{1.674452in}}%
\pgfpathlineto{\pgfqpoint{2.735608in}{1.675908in}}%
\pgfpathlineto{\pgfqpoint{2.739935in}{1.648352in}}%
\pgfpathlineto{\pgfqpoint{2.706426in}{1.644227in}}%
\pgfpathlineto{\pgfqpoint{2.685918in}{1.641892in}}%
\pgfpathlineto{\pgfqpoint{2.682759in}{1.669371in}}%
\pgfpathclose%
\pgfusepath{fill}%
\end{pgfscope}%
\begin{pgfscope}%
\pgfpathrectangle{\pgfqpoint{0.100000in}{0.100000in}}{\pgfqpoint{3.608454in}{2.310000in}}%
\pgfusepath{clip}%
\pgfsetbuttcap%
\pgfsetmiterjoin%
\definecolor{currentfill}{rgb}{0.000000,0.447059,0.776471}%
\pgfsetfillcolor{currentfill}%
\pgfsetlinewidth{0.000000pt}%
\definecolor{currentstroke}{rgb}{0.000000,0.000000,0.000000}%
\pgfsetstrokecolor{currentstroke}%
\pgfsetstrokeopacity{0.000000}%
\pgfsetdash{}{0pt}%
\pgfpathmoveto{\pgfqpoint{1.992478in}{2.070813in}}%
\pgfpathlineto{\pgfqpoint{1.993953in}{2.077757in}}%
\pgfpathlineto{\pgfqpoint{1.991966in}{2.086284in}}%
\pgfpathlineto{\pgfqpoint{1.993550in}{2.092414in}}%
\pgfpathlineto{\pgfqpoint{1.991978in}{2.098902in}}%
\pgfpathlineto{\pgfqpoint{2.033124in}{2.098129in}}%
\pgfpathlineto{\pgfqpoint{2.074816in}{2.097239in}}%
\pgfpathlineto{\pgfqpoint{2.075068in}{2.055715in}}%
\pgfpathlineto{\pgfqpoint{2.075577in}{2.048780in}}%
\pgfpathlineto{\pgfqpoint{2.068749in}{2.048976in}}%
\pgfpathlineto{\pgfqpoint{2.068716in}{2.041994in}}%
\pgfpathlineto{\pgfqpoint{2.061858in}{2.042025in}}%
\pgfpathlineto{\pgfqpoint{2.061781in}{2.035022in}}%
\pgfpathlineto{\pgfqpoint{2.034103in}{2.035335in}}%
\pgfpathlineto{\pgfqpoint{2.034152in}{2.042276in}}%
\pgfpathlineto{\pgfqpoint{2.027218in}{2.042403in}}%
\pgfpathlineto{\pgfqpoint{2.026639in}{2.068568in}}%
\pgfpathlineto{\pgfqpoint{1.992147in}{2.069188in}}%
\pgfpathlineto{\pgfqpoint{1.992478in}{2.070813in}}%
\pgfpathclose%
\pgfusepath{fill}%
\end{pgfscope}%
\begin{pgfscope}%
\pgfpathrectangle{\pgfqpoint{0.100000in}{0.100000in}}{\pgfqpoint{3.608454in}{2.310000in}}%
\pgfusepath{clip}%
\pgfsetbuttcap%
\pgfsetmiterjoin%
\definecolor{currentfill}{rgb}{0.000000,0.439216,0.780392}%
\pgfsetfillcolor{currentfill}%
\pgfsetlinewidth{0.000000pt}%
\definecolor{currentstroke}{rgb}{0.000000,0.000000,0.000000}%
\pgfsetstrokecolor{currentstroke}%
\pgfsetstrokeopacity{0.000000}%
\pgfsetdash{}{0pt}%
\pgfpathmoveto{\pgfqpoint{3.181076in}{1.501996in}}%
\pgfpathlineto{\pgfqpoint{3.166603in}{1.499150in}}%
\pgfpathlineto{\pgfqpoint{3.162472in}{1.516430in}}%
\pgfpathlineto{\pgfqpoint{3.146592in}{1.539583in}}%
\pgfpathlineto{\pgfqpoint{3.143924in}{1.538541in}}%
\pgfpathlineto{\pgfqpoint{3.138889in}{1.546016in}}%
\pgfpathlineto{\pgfqpoint{3.133310in}{1.544667in}}%
\pgfpathlineto{\pgfqpoint{3.130234in}{1.548464in}}%
\pgfpathlineto{\pgfqpoint{3.133732in}{1.555574in}}%
\pgfpathlineto{\pgfqpoint{3.131483in}{1.559276in}}%
\pgfpathlineto{\pgfqpoint{3.139700in}{1.576466in}}%
\pgfpathlineto{\pgfqpoint{3.155139in}{1.587716in}}%
\pgfpathlineto{\pgfqpoint{3.157897in}{1.576523in}}%
\pgfpathlineto{\pgfqpoint{3.176843in}{1.578501in}}%
\pgfpathlineto{\pgfqpoint{3.184525in}{1.574369in}}%
\pgfpathlineto{\pgfqpoint{3.196959in}{1.580835in}}%
\pgfpathlineto{\pgfqpoint{3.208359in}{1.574869in}}%
\pgfpathlineto{\pgfqpoint{3.201240in}{1.566916in}}%
\pgfpathlineto{\pgfqpoint{3.212062in}{1.546662in}}%
\pgfpathlineto{\pgfqpoint{3.204878in}{1.540940in}}%
\pgfpathlineto{\pgfqpoint{3.208552in}{1.535223in}}%
\pgfpathlineto{\pgfqpoint{3.216418in}{1.535851in}}%
\pgfpathlineto{\pgfqpoint{3.221777in}{1.528537in}}%
\pgfpathlineto{\pgfqpoint{3.227209in}{1.526852in}}%
\pgfpathlineto{\pgfqpoint{3.239245in}{1.513965in}}%
\pgfpathlineto{\pgfqpoint{3.181076in}{1.501996in}}%
\pgfpathclose%
\pgfusepath{fill}%
\end{pgfscope}%
\begin{pgfscope}%
\pgfpathrectangle{\pgfqpoint{0.100000in}{0.100000in}}{\pgfqpoint{3.608454in}{2.310000in}}%
\pgfusepath{clip}%
\pgfsetbuttcap%
\pgfsetmiterjoin%
\definecolor{currentfill}{rgb}{0.000000,0.278431,0.860784}%
\pgfsetfillcolor{currentfill}%
\pgfsetlinewidth{0.000000pt}%
\definecolor{currentstroke}{rgb}{0.000000,0.000000,0.000000}%
\pgfsetstrokecolor{currentstroke}%
\pgfsetstrokeopacity{0.000000}%
\pgfsetdash{}{0pt}%
\pgfpathmoveto{\pgfqpoint{1.654457in}{1.551499in}}%
\pgfpathlineto{\pgfqpoint{1.655781in}{1.547900in}}%
\pgfpathlineto{\pgfqpoint{1.654575in}{1.534236in}}%
\pgfpathlineto{\pgfqpoint{1.653438in}{1.520561in}}%
\pgfpathlineto{\pgfqpoint{1.650561in}{1.517162in}}%
\pgfpathlineto{\pgfqpoint{1.595804in}{1.522015in}}%
\pgfpathlineto{\pgfqpoint{1.533391in}{1.528263in}}%
\pgfpathlineto{\pgfqpoint{1.494704in}{1.532257in}}%
\pgfpathlineto{\pgfqpoint{1.500517in}{1.584019in}}%
\pgfpathlineto{\pgfqpoint{1.537205in}{1.579689in}}%
\pgfpathlineto{\pgfqpoint{1.536486in}{1.572765in}}%
\pgfpathlineto{\pgfqpoint{1.571835in}{1.569055in}}%
\pgfpathlineto{\pgfqpoint{1.570505in}{1.555595in}}%
\pgfpathlineto{\pgfqpoint{1.610795in}{1.551818in}}%
\pgfpathlineto{\pgfqpoint{1.611187in}{1.555209in}}%
\pgfpathlineto{\pgfqpoint{1.654457in}{1.551499in}}%
\pgfpathclose%
\pgfusepath{fill}%
\end{pgfscope}%
\begin{pgfscope}%
\pgfpathrectangle{\pgfqpoint{0.100000in}{0.100000in}}{\pgfqpoint{3.608454in}{2.310000in}}%
\pgfusepath{clip}%
\pgfsetbuttcap%
\pgfsetmiterjoin%
\definecolor{currentfill}{rgb}{0.000000,0.713725,0.643137}%
\pgfsetfillcolor{currentfill}%
\pgfsetlinewidth{0.000000pt}%
\definecolor{currentstroke}{rgb}{0.000000,0.000000,0.000000}%
\pgfsetstrokecolor{currentstroke}%
\pgfsetstrokeopacity{0.000000}%
\pgfsetdash{}{0pt}%
\pgfpathmoveto{\pgfqpoint{1.421683in}{1.386315in}}%
\pgfpathlineto{\pgfqpoint{1.421188in}{1.383353in}}%
\pgfpathlineto{\pgfqpoint{1.430521in}{1.376213in}}%
\pgfpathlineto{\pgfqpoint{1.436137in}{1.369904in}}%
\pgfpathlineto{\pgfqpoint{1.436653in}{1.364650in}}%
\pgfpathlineto{\pgfqpoint{1.431712in}{1.356218in}}%
\pgfpathlineto{\pgfqpoint{1.436154in}{1.349875in}}%
\pgfpathlineto{\pgfqpoint{1.435665in}{1.345243in}}%
\pgfpathlineto{\pgfqpoint{1.426840in}{1.337144in}}%
\pgfpathlineto{\pgfqpoint{1.422616in}{1.335846in}}%
\pgfpathlineto{\pgfqpoint{1.412097in}{1.337063in}}%
\pgfpathlineto{\pgfqpoint{1.406993in}{1.344082in}}%
\pgfpathlineto{\pgfqpoint{1.404621in}{1.362130in}}%
\pgfpathlineto{\pgfqpoint{1.405672in}{1.369838in}}%
\pgfpathlineto{\pgfqpoint{1.412208in}{1.374642in}}%
\pgfpathlineto{\pgfqpoint{1.403546in}{1.376980in}}%
\pgfpathlineto{\pgfqpoint{1.396241in}{1.384895in}}%
\pgfpathlineto{\pgfqpoint{1.398138in}{1.389426in}}%
\pgfpathlineto{\pgfqpoint{1.421683in}{1.386315in}}%
\pgfpathclose%
\pgfusepath{fill}%
\end{pgfscope}%
\begin{pgfscope}%
\pgfpathrectangle{\pgfqpoint{0.100000in}{0.100000in}}{\pgfqpoint{3.608454in}{2.310000in}}%
\pgfusepath{clip}%
\pgfsetbuttcap%
\pgfsetmiterjoin%
\definecolor{currentfill}{rgb}{0.000000,0.572549,0.713725}%
\pgfsetfillcolor{currentfill}%
\pgfsetlinewidth{0.000000pt}%
\definecolor{currentstroke}{rgb}{0.000000,0.000000,0.000000}%
\pgfsetstrokecolor{currentstroke}%
\pgfsetstrokeopacity{0.000000}%
\pgfsetdash{}{0pt}%
\pgfpathmoveto{\pgfqpoint{2.517006in}{1.100169in}}%
\pgfpathlineto{\pgfqpoint{2.510594in}{1.099857in}}%
\pgfpathlineto{\pgfqpoint{2.510006in}{1.121373in}}%
\pgfpathlineto{\pgfqpoint{2.512080in}{1.125849in}}%
\pgfpathlineto{\pgfqpoint{2.519951in}{1.126253in}}%
\pgfpathlineto{\pgfqpoint{2.518781in}{1.156793in}}%
\pgfpathlineto{\pgfqpoint{2.520509in}{1.156888in}}%
\pgfpathlineto{\pgfqpoint{2.548186in}{1.158629in}}%
\pgfpathlineto{\pgfqpoint{2.548297in}{1.153932in}}%
\pgfpathlineto{\pgfqpoint{2.555391in}{1.144733in}}%
\pgfpathlineto{\pgfqpoint{2.556313in}{1.139100in}}%
\pgfpathlineto{\pgfqpoint{2.570262in}{1.139592in}}%
\pgfpathlineto{\pgfqpoint{2.579740in}{1.137359in}}%
\pgfpathlineto{\pgfqpoint{2.583376in}{1.141723in}}%
\pgfpathlineto{\pgfqpoint{2.583459in}{1.148544in}}%
\pgfpathlineto{\pgfqpoint{2.597965in}{1.148745in}}%
\pgfpathlineto{\pgfqpoint{2.606315in}{1.148332in}}%
\pgfpathlineto{\pgfqpoint{2.604989in}{1.140525in}}%
\pgfpathlineto{\pgfqpoint{2.606773in}{1.136211in}}%
\pgfpathlineto{\pgfqpoint{2.605651in}{1.120589in}}%
\pgfpathlineto{\pgfqpoint{2.584331in}{1.121514in}}%
\pgfpathlineto{\pgfqpoint{2.573891in}{1.113688in}}%
\pgfpathlineto{\pgfqpoint{2.573691in}{1.108364in}}%
\pgfpathlineto{\pgfqpoint{2.561312in}{1.107350in}}%
\pgfpathlineto{\pgfqpoint{2.553934in}{1.104609in}}%
\pgfpathlineto{\pgfqpoint{2.552190in}{1.107568in}}%
\pgfpathlineto{\pgfqpoint{2.544235in}{1.104329in}}%
\pgfpathlineto{\pgfqpoint{2.519200in}{1.103140in}}%
\pgfpathlineto{\pgfqpoint{2.517006in}{1.100169in}}%
\pgfpathclose%
\pgfusepath{fill}%
\end{pgfscope}%
\begin{pgfscope}%
\pgfpathrectangle{\pgfqpoint{0.100000in}{0.100000in}}{\pgfqpoint{3.608454in}{2.310000in}}%
\pgfusepath{clip}%
\pgfsetbuttcap%
\pgfsetmiterjoin%
\definecolor{currentfill}{rgb}{0.000000,0.360784,0.819608}%
\pgfsetfillcolor{currentfill}%
\pgfsetlinewidth{0.000000pt}%
\definecolor{currentstroke}{rgb}{0.000000,0.000000,0.000000}%
\pgfsetstrokecolor{currentstroke}%
\pgfsetstrokeopacity{0.000000}%
\pgfsetdash{}{0pt}%
\pgfpathmoveto{\pgfqpoint{1.027182in}{0.384801in}}%
\pgfpathlineto{\pgfqpoint{1.017303in}{0.390307in}}%
\pgfpathlineto{\pgfqpoint{1.020524in}{0.396033in}}%
\pgfpathlineto{\pgfqpoint{1.019250in}{0.398079in}}%
\pgfpathlineto{\pgfqpoint{1.024400in}{0.406832in}}%
\pgfpathlineto{\pgfqpoint{1.022393in}{0.408007in}}%
\pgfpathlineto{\pgfqpoint{1.023662in}{0.410127in}}%
\pgfpathlineto{\pgfqpoint{1.021385in}{0.411488in}}%
\pgfpathlineto{\pgfqpoint{1.022698in}{0.413642in}}%
\pgfpathlineto{\pgfqpoint{1.018953in}{0.415851in}}%
\pgfpathlineto{\pgfqpoint{1.018284in}{0.417154in}}%
\pgfpathlineto{\pgfqpoint{1.015043in}{0.415179in}}%
\pgfpathlineto{\pgfqpoint{1.012216in}{0.417308in}}%
\pgfpathlineto{\pgfqpoint{1.015064in}{0.421049in}}%
\pgfpathlineto{\pgfqpoint{1.014239in}{0.423539in}}%
\pgfpathlineto{\pgfqpoint{1.017896in}{0.425014in}}%
\pgfpathlineto{\pgfqpoint{1.019938in}{0.430914in}}%
\pgfpathlineto{\pgfqpoint{1.002410in}{0.441046in}}%
\pgfpathlineto{\pgfqpoint{0.996218in}{0.444880in}}%
\pgfpathlineto{\pgfqpoint{0.997579in}{0.447110in}}%
\pgfpathlineto{\pgfqpoint{0.986635in}{0.453959in}}%
\pgfpathlineto{\pgfqpoint{0.983915in}{0.449521in}}%
\pgfpathlineto{\pgfqpoint{0.979580in}{0.452265in}}%
\pgfpathlineto{\pgfqpoint{0.986452in}{0.463301in}}%
\pgfpathlineto{\pgfqpoint{0.972860in}{0.471797in}}%
\pgfpathlineto{\pgfqpoint{0.968975in}{0.474416in}}%
\pgfpathlineto{\pgfqpoint{0.966684in}{0.470967in}}%
\pgfpathlineto{\pgfqpoint{0.950166in}{0.482212in}}%
\pgfpathlineto{\pgfqpoint{0.940198in}{0.483801in}}%
\pgfpathlineto{\pgfqpoint{0.915352in}{0.487908in}}%
\pgfpathlineto{\pgfqpoint{0.914054in}{0.486250in}}%
\pgfpathlineto{\pgfqpoint{0.902205in}{0.495355in}}%
\pgfpathlineto{\pgfqpoint{0.894650in}{0.485392in}}%
\pgfpathlineto{\pgfqpoint{0.892390in}{0.483966in}}%
\pgfpathlineto{\pgfqpoint{0.888822in}{0.486756in}}%
\pgfpathlineto{\pgfqpoint{0.887256in}{0.484755in}}%
\pgfpathlineto{\pgfqpoint{0.883274in}{0.487912in}}%
\pgfpathlineto{\pgfqpoint{0.880946in}{0.484915in}}%
\pgfpathlineto{\pgfqpoint{0.878680in}{0.486403in}}%
\pgfpathlineto{\pgfqpoint{0.867623in}{0.494883in}}%
\pgfpathlineto{\pgfqpoint{0.854046in}{0.505533in}}%
\pgfpathlineto{\pgfqpoint{0.845087in}{0.513012in}}%
\pgfpathlineto{\pgfqpoint{0.846928in}{0.515267in}}%
\pgfpathlineto{\pgfqpoint{0.831654in}{0.528086in}}%
\pgfpathlineto{\pgfqpoint{0.830011in}{0.526141in}}%
\pgfpathlineto{\pgfqpoint{0.821946in}{0.532976in}}%
\pgfpathlineto{\pgfqpoint{0.819842in}{0.530526in}}%
\pgfpathlineto{\pgfqpoint{0.812784in}{0.536599in}}%
\pgfpathlineto{\pgfqpoint{0.802699in}{0.545595in}}%
\pgfpathlineto{\pgfqpoint{0.807682in}{0.551257in}}%
\pgfpathlineto{\pgfqpoint{0.806268in}{0.552525in}}%
\pgfpathlineto{\pgfqpoint{0.813107in}{0.560166in}}%
\pgfpathlineto{\pgfqpoint{0.813592in}{0.559787in}}%
\pgfpathlineto{\pgfqpoint{0.820507in}{0.567425in}}%
\pgfpathlineto{\pgfqpoint{0.823025in}{0.569214in}}%
\pgfpathlineto{\pgfqpoint{0.824219in}{0.568707in}}%
\pgfpathlineto{\pgfqpoint{0.829339in}{0.574413in}}%
\pgfpathlineto{\pgfqpoint{0.831146in}{0.572881in}}%
\pgfpathlineto{\pgfqpoint{0.836205in}{0.579038in}}%
\pgfpathlineto{\pgfqpoint{0.838193in}{0.577309in}}%
\pgfpathlineto{\pgfqpoint{0.839724in}{0.579022in}}%
\pgfpathlineto{\pgfqpoint{0.843712in}{0.575563in}}%
\pgfpathlineto{\pgfqpoint{0.845387in}{0.577604in}}%
\pgfpathlineto{\pgfqpoint{0.847165in}{0.576109in}}%
\pgfpathlineto{\pgfqpoint{0.848929in}{0.578173in}}%
\pgfpathlineto{\pgfqpoint{0.850691in}{0.576622in}}%
\pgfpathlineto{\pgfqpoint{0.856024in}{0.582334in}}%
\pgfpathlineto{\pgfqpoint{0.853750in}{0.584249in}}%
\pgfpathlineto{\pgfqpoint{0.867382in}{0.600383in}}%
\pgfpathlineto{\pgfqpoint{0.869420in}{0.598690in}}%
\pgfpathlineto{\pgfqpoint{0.872825in}{0.602697in}}%
\pgfpathlineto{\pgfqpoint{0.874796in}{0.601023in}}%
\pgfpathlineto{\pgfqpoint{0.881527in}{0.608877in}}%
\pgfpathlineto{\pgfqpoint{0.883388in}{0.607238in}}%
\pgfpathlineto{\pgfqpoint{0.888579in}{0.613334in}}%
\pgfpathlineto{\pgfqpoint{0.886576in}{0.615036in}}%
\pgfpathlineto{\pgfqpoint{0.895296in}{0.625196in}}%
\pgfpathlineto{\pgfqpoint{0.901313in}{0.620035in}}%
\pgfpathlineto{\pgfqpoint{0.904546in}{0.623960in}}%
\pgfpathlineto{\pgfqpoint{0.914790in}{0.615304in}}%
\pgfpathlineto{\pgfqpoint{0.921542in}{0.623737in}}%
\pgfpathlineto{\pgfqpoint{0.929884in}{0.616921in}}%
\pgfpathlineto{\pgfqpoint{0.926478in}{0.612658in}}%
\pgfpathlineto{\pgfqpoint{0.930663in}{0.609300in}}%
\pgfpathlineto{\pgfqpoint{0.934029in}{0.613492in}}%
\pgfpathlineto{\pgfqpoint{0.938173in}{0.610139in}}%
\pgfpathlineto{\pgfqpoint{0.936545in}{0.608078in}}%
\pgfpathlineto{\pgfqpoint{0.938627in}{0.606373in}}%
\pgfpathlineto{\pgfqpoint{0.936932in}{0.604209in}}%
\pgfpathlineto{\pgfqpoint{0.940885in}{0.601118in}}%
\pgfpathlineto{\pgfqpoint{0.946326in}{0.607097in}}%
\pgfpathlineto{\pgfqpoint{0.952323in}{0.602325in}}%
\pgfpathlineto{\pgfqpoint{0.955683in}{0.606572in}}%
\pgfpathlineto{\pgfqpoint{0.961995in}{0.601456in}}%
\pgfpathlineto{\pgfqpoint{0.964303in}{0.603251in}}%
\pgfpathlineto{\pgfqpoint{0.971435in}{0.612298in}}%
\pgfpathlineto{\pgfqpoint{0.969910in}{0.613404in}}%
\pgfpathlineto{\pgfqpoint{0.971558in}{0.615632in}}%
\pgfpathlineto{\pgfqpoint{0.967277in}{0.618907in}}%
\pgfpathlineto{\pgfqpoint{0.972240in}{0.625296in}}%
\pgfpathlineto{\pgfqpoint{0.974510in}{0.627028in}}%
\pgfpathlineto{\pgfqpoint{0.972344in}{0.628719in}}%
\pgfpathlineto{\pgfqpoint{0.974033in}{0.630789in}}%
\pgfpathlineto{\pgfqpoint{0.971825in}{0.632484in}}%
\pgfpathlineto{\pgfqpoint{0.973534in}{0.634639in}}%
\pgfpathlineto{\pgfqpoint{0.971445in}{0.636313in}}%
\pgfpathlineto{\pgfqpoint{0.976297in}{0.641465in}}%
\pgfpathlineto{\pgfqpoint{0.997892in}{0.624739in}}%
\pgfpathlineto{\pgfqpoint{1.026608in}{0.603847in}}%
\pgfpathlineto{\pgfqpoint{1.052163in}{0.586365in}}%
\pgfpathlineto{\pgfqpoint{1.056440in}{0.584227in}}%
\pgfpathlineto{\pgfqpoint{1.066205in}{0.578011in}}%
\pgfpathlineto{\pgfqpoint{1.074192in}{0.589904in}}%
\pgfpathlineto{\pgfqpoint{1.083089in}{0.584700in}}%
\pgfpathlineto{\pgfqpoint{1.110388in}{0.567904in}}%
\pgfpathlineto{\pgfqpoint{1.124077in}{0.559590in}}%
\pgfpathlineto{\pgfqpoint{1.114317in}{0.541984in}}%
\pgfpathlineto{\pgfqpoint{1.101033in}{0.518021in}}%
\pgfpathlineto{\pgfqpoint{1.087919in}{0.494364in}}%
\pgfpathlineto{\pgfqpoint{1.075939in}{0.472754in}}%
\pgfpathlineto{\pgfqpoint{1.063110in}{0.449610in}}%
\pgfpathlineto{\pgfqpoint{1.050408in}{0.426699in}}%
\pgfpathlineto{\pgfqpoint{1.037825in}{0.403999in}}%
\pgfpathlineto{\pgfqpoint{1.027182in}{0.384801in}}%
\pgfpathclose%
\pgfusepath{fill}%
\end{pgfscope}%
\begin{pgfscope}%
\pgfpathrectangle{\pgfqpoint{0.100000in}{0.100000in}}{\pgfqpoint{3.608454in}{2.310000in}}%
\pgfusepath{clip}%
\pgfsetbuttcap%
\pgfsetmiterjoin%
\definecolor{currentfill}{rgb}{0.000000,0.580392,0.709804}%
\pgfsetfillcolor{currentfill}%
\pgfsetlinewidth{0.000000pt}%
\definecolor{currentstroke}{rgb}{0.000000,0.000000,0.000000}%
\pgfsetstrokecolor{currentstroke}%
\pgfsetstrokeopacity{0.000000}%
\pgfsetdash{}{0pt}%
\pgfpathmoveto{\pgfqpoint{2.289673in}{1.168159in}}%
\pgfpathlineto{\pgfqpoint{2.248336in}{1.167927in}}%
\pgfpathlineto{\pgfqpoint{2.211773in}{1.168036in}}%
\pgfpathlineto{\pgfqpoint{2.211810in}{1.181895in}}%
\pgfpathlineto{\pgfqpoint{2.194584in}{1.181869in}}%
\pgfpathlineto{\pgfqpoint{2.194554in}{1.204465in}}%
\pgfpathlineto{\pgfqpoint{2.193249in}{1.204475in}}%
\pgfpathlineto{\pgfqpoint{2.193719in}{1.240848in}}%
\pgfpathlineto{\pgfqpoint{2.192435in}{1.247763in}}%
\pgfpathlineto{\pgfqpoint{2.196058in}{1.254577in}}%
\pgfpathlineto{\pgfqpoint{2.199587in}{1.254527in}}%
\pgfpathlineto{\pgfqpoint{2.199885in}{1.267142in}}%
\pgfpathlineto{\pgfqpoint{2.227204in}{1.266822in}}%
\pgfpathlineto{\pgfqpoint{2.227021in}{1.254178in}}%
\pgfpathlineto{\pgfqpoint{2.255244in}{1.254033in}}%
\pgfpathlineto{\pgfqpoint{2.260153in}{1.246809in}}%
\pgfpathlineto{\pgfqpoint{2.268476in}{1.251773in}}%
\pgfpathlineto{\pgfqpoint{2.268526in}{1.240287in}}%
\pgfpathlineto{\pgfqpoint{2.278985in}{1.235414in}}%
\pgfpathlineto{\pgfqpoint{2.279053in}{1.231943in}}%
\pgfpathlineto{\pgfqpoint{2.280237in}{1.188871in}}%
\pgfpathlineto{\pgfqpoint{2.290504in}{1.188921in}}%
\pgfpathlineto{\pgfqpoint{2.289673in}{1.168159in}}%
\pgfpathclose%
\pgfusepath{fill}%
\end{pgfscope}%
\begin{pgfscope}%
\pgfpathrectangle{\pgfqpoint{0.100000in}{0.100000in}}{\pgfqpoint{3.608454in}{2.310000in}}%
\pgfusepath{clip}%
\pgfsetbuttcap%
\pgfsetmiterjoin%
\definecolor{currentfill}{rgb}{0.000000,0.800000,0.600000}%
\pgfsetfillcolor{currentfill}%
\pgfsetlinewidth{0.000000pt}%
\definecolor{currentstroke}{rgb}{0.000000,0.000000,0.000000}%
\pgfsetstrokecolor{currentstroke}%
\pgfsetstrokeopacity{0.000000}%
\pgfsetdash{}{0pt}%
\pgfpathmoveto{\pgfqpoint{2.175503in}{0.924555in}}%
\pgfpathlineto{\pgfqpoint{2.170065in}{0.926960in}}%
\pgfpathlineto{\pgfqpoint{2.157026in}{0.928424in}}%
\pgfpathlineto{\pgfqpoint{2.147313in}{0.939341in}}%
\pgfpathlineto{\pgfqpoint{2.141202in}{0.939334in}}%
\pgfpathlineto{\pgfqpoint{2.140773in}{0.915277in}}%
\pgfpathlineto{\pgfqpoint{2.123559in}{0.920377in}}%
\pgfpathlineto{\pgfqpoint{2.115386in}{0.923776in}}%
\pgfpathlineto{\pgfqpoint{2.109858in}{0.929360in}}%
\pgfpathlineto{\pgfqpoint{2.092356in}{0.940891in}}%
\pgfpathlineto{\pgfqpoint{2.087358in}{0.933839in}}%
\pgfpathlineto{\pgfqpoint{2.079921in}{0.933374in}}%
\pgfpathlineto{\pgfqpoint{2.070424in}{0.935180in}}%
\pgfpathlineto{\pgfqpoint{2.068217in}{0.938848in}}%
\pgfpathlineto{\pgfqpoint{2.056583in}{0.934051in}}%
\pgfpathlineto{\pgfqpoint{2.051407in}{0.934194in}}%
\pgfpathlineto{\pgfqpoint{2.045055in}{0.940136in}}%
\pgfpathlineto{\pgfqpoint{2.045185in}{0.954206in}}%
\pgfpathlineto{\pgfqpoint{2.041706in}{0.956720in}}%
\pgfpathlineto{\pgfqpoint{2.055634in}{0.956575in}}%
\pgfpathlineto{\pgfqpoint{2.055870in}{0.984320in}}%
\pgfpathlineto{\pgfqpoint{2.062908in}{0.984282in}}%
\pgfpathlineto{\pgfqpoint{2.062941in}{0.991185in}}%
\pgfpathlineto{\pgfqpoint{2.073230in}{0.991141in}}%
\pgfpathlineto{\pgfqpoint{2.073255in}{0.998037in}}%
\pgfpathlineto{\pgfqpoint{2.110850in}{0.997922in}}%
\pgfpathlineto{\pgfqpoint{2.110869in}{0.984129in}}%
\pgfpathlineto{\pgfqpoint{2.141998in}{0.984277in}}%
\pgfpathlineto{\pgfqpoint{2.142324in}{1.001839in}}%
\pgfpathlineto{\pgfqpoint{2.151260in}{0.999398in}}%
\pgfpathlineto{\pgfqpoint{2.170781in}{0.999186in}}%
\pgfpathlineto{\pgfqpoint{2.176519in}{0.997121in}}%
\pgfpathlineto{\pgfqpoint{2.176495in}{0.972103in}}%
\pgfpathlineto{\pgfqpoint{2.155587in}{0.972321in}}%
\pgfpathlineto{\pgfqpoint{2.156410in}{0.959401in}}%
\pgfpathlineto{\pgfqpoint{2.165791in}{0.959401in}}%
\pgfpathlineto{\pgfqpoint{2.166804in}{0.946907in}}%
\pgfpathlineto{\pgfqpoint{2.169177in}{0.935220in}}%
\pgfpathlineto{\pgfqpoint{2.174128in}{0.929756in}}%
\pgfpathlineto{\pgfqpoint{2.175503in}{0.924555in}}%
\pgfpathclose%
\pgfusepath{fill}%
\end{pgfscope}%
\begin{pgfscope}%
\pgfpathrectangle{\pgfqpoint{0.100000in}{0.100000in}}{\pgfqpoint{3.608454in}{2.310000in}}%
\pgfusepath{clip}%
\pgfsetbuttcap%
\pgfsetmiterjoin%
\definecolor{currentfill}{rgb}{0.000000,0.282353,0.858824}%
\pgfsetfillcolor{currentfill}%
\pgfsetlinewidth{0.000000pt}%
\definecolor{currentstroke}{rgb}{0.000000,0.000000,0.000000}%
\pgfsetstrokecolor{currentstroke}%
\pgfsetstrokeopacity{0.000000}%
\pgfsetdash{}{0pt}%
\pgfpathmoveto{\pgfqpoint{2.138316in}{1.510307in}}%
\pgfpathlineto{\pgfqpoint{2.124622in}{1.510289in}}%
\pgfpathlineto{\pgfqpoint{2.124501in}{1.545418in}}%
\pgfpathlineto{\pgfqpoint{2.121931in}{1.545408in}}%
\pgfpathlineto{\pgfqpoint{2.121832in}{1.566001in}}%
\pgfpathlineto{\pgfqpoint{2.156052in}{1.566272in}}%
\pgfpathlineto{\pgfqpoint{2.155774in}{1.593709in}}%
\pgfpathlineto{\pgfqpoint{2.196902in}{1.594193in}}%
\pgfpathlineto{\pgfqpoint{2.210467in}{1.594411in}}%
\pgfpathlineto{\pgfqpoint{2.210987in}{1.567017in}}%
\pgfpathlineto{\pgfqpoint{2.204148in}{1.566920in}}%
\pgfpathlineto{\pgfqpoint{2.204489in}{1.546166in}}%
\pgfpathlineto{\pgfqpoint{2.205803in}{1.538835in}}%
\pgfpathlineto{\pgfqpoint{2.206358in}{1.511376in}}%
\pgfpathlineto{\pgfqpoint{2.192689in}{1.511198in}}%
\pgfpathlineto{\pgfqpoint{2.172375in}{1.510519in}}%
\pgfpathlineto{\pgfqpoint{2.138316in}{1.510307in}}%
\pgfpathclose%
\pgfusepath{fill}%
\end{pgfscope}%
\begin{pgfscope}%
\pgfpathrectangle{\pgfqpoint{0.100000in}{0.100000in}}{\pgfqpoint{3.608454in}{2.310000in}}%
\pgfusepath{clip}%
\pgfsetbuttcap%
\pgfsetmiterjoin%
\definecolor{currentfill}{rgb}{0.000000,0.505882,0.747059}%
\pgfsetfillcolor{currentfill}%
\pgfsetlinewidth{0.000000pt}%
\definecolor{currentstroke}{rgb}{0.000000,0.000000,0.000000}%
\pgfsetstrokecolor{currentstroke}%
\pgfsetstrokeopacity{0.000000}%
\pgfsetdash{}{0pt}%
\pgfpathmoveto{\pgfqpoint{1.782538in}{1.034546in}}%
\pgfpathlineto{\pgfqpoint{1.781300in}{1.012079in}}%
\pgfpathlineto{\pgfqpoint{1.754289in}{1.013713in}}%
\pgfpathlineto{\pgfqpoint{1.746176in}{1.014219in}}%
\pgfpathlineto{\pgfqpoint{1.748391in}{1.048651in}}%
\pgfpathlineto{\pgfqpoint{1.750416in}{1.083136in}}%
\pgfpathlineto{\pgfqpoint{1.785102in}{1.081107in}}%
\pgfpathlineto{\pgfqpoint{1.782538in}{1.034546in}}%
\pgfpathclose%
\pgfusepath{fill}%
\end{pgfscope}%
\begin{pgfscope}%
\pgfpathrectangle{\pgfqpoint{0.100000in}{0.100000in}}{\pgfqpoint{3.608454in}{2.310000in}}%
\pgfusepath{clip}%
\pgfsetbuttcap%
\pgfsetmiterjoin%
\definecolor{currentfill}{rgb}{0.000000,0.486275,0.756863}%
\pgfsetfillcolor{currentfill}%
\pgfsetlinewidth{0.000000pt}%
\definecolor{currentstroke}{rgb}{0.000000,0.000000,0.000000}%
\pgfsetstrokecolor{currentstroke}%
\pgfsetstrokeopacity{0.000000}%
\pgfsetdash{}{0pt}%
\pgfpathmoveto{\pgfqpoint{1.893386in}{1.070793in}}%
\pgfpathlineto{\pgfqpoint{1.873435in}{1.071561in}}%
\pgfpathlineto{\pgfqpoint{1.873684in}{1.092215in}}%
\pgfpathlineto{\pgfqpoint{1.825972in}{1.094304in}}%
\pgfpathlineto{\pgfqpoint{1.826418in}{1.110191in}}%
\pgfpathlineto{\pgfqpoint{1.826923in}{1.122162in}}%
\pgfpathlineto{\pgfqpoint{1.854121in}{1.120697in}}%
\pgfpathlineto{\pgfqpoint{1.855366in}{1.155288in}}%
\pgfpathlineto{\pgfqpoint{1.844616in}{1.170342in}}%
\pgfpathlineto{\pgfqpoint{1.837450in}{1.170858in}}%
\pgfpathlineto{\pgfqpoint{1.830118in}{1.178718in}}%
\pgfpathlineto{\pgfqpoint{1.825437in}{1.188388in}}%
\pgfpathlineto{\pgfqpoint{1.854168in}{1.187060in}}%
\pgfpathlineto{\pgfqpoint{1.895416in}{1.185262in}}%
\pgfpathlineto{\pgfqpoint{1.951284in}{1.183534in}}%
\pgfpathlineto{\pgfqpoint{1.950456in}{1.151489in}}%
\pgfpathlineto{\pgfqpoint{1.949595in}{1.117555in}}%
\pgfpathlineto{\pgfqpoint{1.901714in}{1.119083in}}%
\pgfpathlineto{\pgfqpoint{1.900701in}{1.084287in}}%
\pgfpathlineto{\pgfqpoint{1.893880in}{1.084537in}}%
\pgfpathlineto{\pgfqpoint{1.893386in}{1.070793in}}%
\pgfpathclose%
\pgfusepath{fill}%
\end{pgfscope}%
\begin{pgfscope}%
\pgfpathrectangle{\pgfqpoint{0.100000in}{0.100000in}}{\pgfqpoint{3.608454in}{2.310000in}}%
\pgfusepath{clip}%
\pgfsetbuttcap%
\pgfsetmiterjoin%
\definecolor{currentfill}{rgb}{0.000000,0.678431,0.660784}%
\pgfsetfillcolor{currentfill}%
\pgfsetlinewidth{0.000000pt}%
\definecolor{currentstroke}{rgb}{0.000000,0.000000,0.000000}%
\pgfsetstrokecolor{currentstroke}%
\pgfsetstrokeopacity{0.000000}%
\pgfsetdash{}{0pt}%
\pgfpathmoveto{\pgfqpoint{1.068484in}{1.440713in}}%
\pgfpathlineto{\pgfqpoint{1.064107in}{1.434490in}}%
\pgfpathlineto{\pgfqpoint{1.060054in}{1.435184in}}%
\pgfpathlineto{\pgfqpoint{1.053099in}{1.429145in}}%
\pgfpathlineto{\pgfqpoint{1.051426in}{1.418402in}}%
\pgfpathlineto{\pgfqpoint{1.045136in}{1.415471in}}%
\pgfpathlineto{\pgfqpoint{1.036855in}{1.417471in}}%
\pgfpathlineto{\pgfqpoint{1.031123in}{1.409823in}}%
\pgfpathlineto{\pgfqpoint{1.000345in}{1.416064in}}%
\pgfpathlineto{\pgfqpoint{0.938194in}{1.428823in}}%
\pgfpathlineto{\pgfqpoint{0.939931in}{1.436894in}}%
\pgfpathlineto{\pgfqpoint{0.954290in}{1.503771in}}%
\pgfpathlineto{\pgfqpoint{0.960301in}{1.531851in}}%
\pgfpathlineto{\pgfqpoint{1.034454in}{1.516384in}}%
\pgfpathlineto{\pgfqpoint{1.061605in}{1.511149in}}%
\pgfpathlineto{\pgfqpoint{1.072847in}{1.517633in}}%
\pgfpathlineto{\pgfqpoint{1.077321in}{1.529704in}}%
\pgfpathlineto{\pgfqpoint{1.075863in}{1.537379in}}%
\pgfpathlineto{\pgfqpoint{1.077577in}{1.552756in}}%
\pgfpathlineto{\pgfqpoint{1.087943in}{1.551215in}}%
\pgfpathlineto{\pgfqpoint{1.090523in}{1.546657in}}%
\pgfpathlineto{\pgfqpoint{1.105133in}{1.552914in}}%
\pgfpathlineto{\pgfqpoint{1.116941in}{1.552445in}}%
\pgfpathlineto{\pgfqpoint{1.112401in}{1.545424in}}%
\pgfpathlineto{\pgfqpoint{1.120187in}{1.531417in}}%
\pgfpathlineto{\pgfqpoint{1.129198in}{1.527981in}}%
\pgfpathlineto{\pgfqpoint{1.131819in}{1.519099in}}%
\pgfpathlineto{\pgfqpoint{1.128853in}{1.510727in}}%
\pgfpathlineto{\pgfqpoint{1.137161in}{1.499935in}}%
\pgfpathlineto{\pgfqpoint{1.136590in}{1.496867in}}%
\pgfpathlineto{\pgfqpoint{1.148005in}{1.494799in}}%
\pgfpathlineto{\pgfqpoint{1.148850in}{1.487724in}}%
\pgfpathlineto{\pgfqpoint{1.101981in}{1.496172in}}%
\pgfpathlineto{\pgfqpoint{1.104984in}{1.493265in}}%
\pgfpathlineto{\pgfqpoint{1.095705in}{1.487818in}}%
\pgfpathlineto{\pgfqpoint{1.092637in}{1.469533in}}%
\pgfpathlineto{\pgfqpoint{1.089487in}{1.463907in}}%
\pgfpathlineto{\pgfqpoint{1.078936in}{1.465774in}}%
\pgfpathlineto{\pgfqpoint{1.072686in}{1.463362in}}%
\pgfpathlineto{\pgfqpoint{1.068484in}{1.440713in}}%
\pgfpathclose%
\pgfusepath{fill}%
\end{pgfscope}%
\begin{pgfscope}%
\pgfpathrectangle{\pgfqpoint{0.100000in}{0.100000in}}{\pgfqpoint{3.608454in}{2.310000in}}%
\pgfusepath{clip}%
\pgfsetbuttcap%
\pgfsetmiterjoin%
\definecolor{currentfill}{rgb}{0.000000,0.560784,0.719608}%
\pgfsetfillcolor{currentfill}%
\pgfsetlinewidth{0.000000pt}%
\definecolor{currentstroke}{rgb}{0.000000,0.000000,0.000000}%
\pgfsetstrokecolor{currentstroke}%
\pgfsetstrokeopacity{0.000000}%
\pgfsetdash{}{0pt}%
\pgfpathmoveto{\pgfqpoint{2.831040in}{1.546066in}}%
\pgfpathlineto{\pgfqpoint{2.825402in}{1.579904in}}%
\pgfpathlineto{\pgfqpoint{2.818905in}{1.581252in}}%
\pgfpathlineto{\pgfqpoint{2.833971in}{1.586462in}}%
\pgfpathlineto{\pgfqpoint{2.840470in}{1.580694in}}%
\pgfpathlineto{\pgfqpoint{2.847896in}{1.579132in}}%
\pgfpathlineto{\pgfqpoint{2.860187in}{1.585033in}}%
\pgfpathlineto{\pgfqpoint{2.873522in}{1.593422in}}%
\pgfpathlineto{\pgfqpoint{2.876219in}{1.593055in}}%
\pgfpathlineto{\pgfqpoint{2.877765in}{1.580902in}}%
\pgfpathlineto{\pgfqpoint{2.883315in}{1.581642in}}%
\pgfpathlineto{\pgfqpoint{2.884117in}{1.575727in}}%
\pgfpathlineto{\pgfqpoint{2.878552in}{1.574944in}}%
\pgfpathlineto{\pgfqpoint{2.879290in}{1.569013in}}%
\pgfpathlineto{\pgfqpoint{2.873516in}{1.568184in}}%
\pgfpathlineto{\pgfqpoint{2.868477in}{1.562448in}}%
\pgfpathlineto{\pgfqpoint{2.869199in}{1.556678in}}%
\pgfpathlineto{\pgfqpoint{2.853507in}{1.554639in}}%
\pgfpathlineto{\pgfqpoint{2.854535in}{1.548999in}}%
\pgfpathlineto{\pgfqpoint{2.831040in}{1.546066in}}%
\pgfpathclose%
\pgfusepath{fill}%
\end{pgfscope}%
\begin{pgfscope}%
\pgfpathrectangle{\pgfqpoint{0.100000in}{0.100000in}}{\pgfqpoint{3.608454in}{2.310000in}}%
\pgfusepath{clip}%
\pgfsetbuttcap%
\pgfsetmiterjoin%
\definecolor{currentfill}{rgb}{0.000000,0.611765,0.694118}%
\pgfsetfillcolor{currentfill}%
\pgfsetlinewidth{0.000000pt}%
\definecolor{currentstroke}{rgb}{0.000000,0.000000,0.000000}%
\pgfsetstrokecolor{currentstroke}%
\pgfsetstrokeopacity{0.000000}%
\pgfsetdash{}{0pt}%
\pgfpathmoveto{\pgfqpoint{1.611520in}{1.962062in}}%
\pgfpathlineto{\pgfqpoint{1.606786in}{1.914838in}}%
\pgfpathlineto{\pgfqpoint{1.602057in}{1.917205in}}%
\pgfpathlineto{\pgfqpoint{1.587963in}{1.918755in}}%
\pgfpathlineto{\pgfqpoint{1.588691in}{1.925458in}}%
\pgfpathlineto{\pgfqpoint{1.581763in}{1.926153in}}%
\pgfpathlineto{\pgfqpoint{1.582466in}{1.932625in}}%
\pgfpathlineto{\pgfqpoint{1.562162in}{1.934829in}}%
\pgfpathlineto{\pgfqpoint{1.564057in}{1.962319in}}%
\pgfpathlineto{\pgfqpoint{1.573489in}{1.961301in}}%
\pgfpathlineto{\pgfqpoint{1.574985in}{1.975097in}}%
\pgfpathlineto{\pgfqpoint{1.582076in}{1.976639in}}%
\pgfpathlineto{\pgfqpoint{1.592373in}{1.975555in}}%
\pgfpathlineto{\pgfqpoint{1.595495in}{1.971747in}}%
\pgfpathlineto{\pgfqpoint{1.612308in}{1.970040in}}%
\pgfpathlineto{\pgfqpoint{1.611520in}{1.962062in}}%
\pgfpathclose%
\pgfusepath{fill}%
\end{pgfscope}%
\begin{pgfscope}%
\pgfpathrectangle{\pgfqpoint{0.100000in}{0.100000in}}{\pgfqpoint{3.608454in}{2.310000in}}%
\pgfusepath{clip}%
\pgfsetbuttcap%
\pgfsetmiterjoin%
\definecolor{currentfill}{rgb}{0.000000,0.560784,0.719608}%
\pgfsetfillcolor{currentfill}%
\pgfsetlinewidth{0.000000pt}%
\definecolor{currentstroke}{rgb}{0.000000,0.000000,0.000000}%
\pgfsetstrokecolor{currentstroke}%
\pgfsetstrokeopacity{0.000000}%
\pgfsetdash{}{0pt}%
\pgfpathmoveto{\pgfqpoint{3.298349in}{1.383037in}}%
\pgfpathlineto{\pgfqpoint{3.291222in}{1.379099in}}%
\pgfpathlineto{\pgfqpoint{3.288344in}{1.383217in}}%
\pgfpathlineto{\pgfqpoint{3.291045in}{1.389057in}}%
\pgfpathlineto{\pgfqpoint{3.289366in}{1.396494in}}%
\pgfpathlineto{\pgfqpoint{3.281967in}{1.394943in}}%
\pgfpathlineto{\pgfqpoint{3.282347in}{1.403895in}}%
\pgfpathlineto{\pgfqpoint{3.281272in}{1.408524in}}%
\pgfpathlineto{\pgfqpoint{3.284538in}{1.423341in}}%
\pgfpathlineto{\pgfqpoint{3.290599in}{1.431089in}}%
\pgfpathlineto{\pgfqpoint{3.274847in}{1.487194in}}%
\pgfpathlineto{\pgfqpoint{3.281434in}{1.488221in}}%
\pgfpathlineto{\pgfqpoint{3.288460in}{1.495799in}}%
\pgfpathlineto{\pgfqpoint{3.293773in}{1.492637in}}%
\pgfpathlineto{\pgfqpoint{3.298424in}{1.484780in}}%
\pgfpathlineto{\pgfqpoint{3.300159in}{1.474127in}}%
\pgfpathlineto{\pgfqpoint{3.304555in}{1.471386in}}%
\pgfpathlineto{\pgfqpoint{3.308496in}{1.463629in}}%
\pgfpathlineto{\pgfqpoint{3.317538in}{1.456649in}}%
\pgfpathlineto{\pgfqpoint{3.324360in}{1.455596in}}%
\pgfpathlineto{\pgfqpoint{3.331882in}{1.431286in}}%
\pgfpathlineto{\pgfqpoint{3.329200in}{1.401923in}}%
\pgfpathlineto{\pgfqpoint{3.327417in}{1.395926in}}%
\pgfpathlineto{\pgfqpoint{3.304739in}{1.388327in}}%
\pgfpathlineto{\pgfqpoint{3.298349in}{1.383037in}}%
\pgfpathclose%
\pgfusepath{fill}%
\end{pgfscope}%
\begin{pgfscope}%
\pgfpathrectangle{\pgfqpoint{0.100000in}{0.100000in}}{\pgfqpoint{3.608454in}{2.310000in}}%
\pgfusepath{clip}%
\pgfsetbuttcap%
\pgfsetmiterjoin%
\definecolor{currentfill}{rgb}{0.000000,0.443137,0.778431}%
\pgfsetfillcolor{currentfill}%
\pgfsetlinewidth{0.000000pt}%
\definecolor{currentstroke}{rgb}{0.000000,0.000000,0.000000}%
\pgfsetstrokecolor{currentstroke}%
\pgfsetstrokeopacity{0.000000}%
\pgfsetdash{}{0pt}%
\pgfpathmoveto{\pgfqpoint{2.115090in}{1.593467in}}%
\pgfpathlineto{\pgfqpoint{2.067449in}{1.593735in}}%
\pgfpathlineto{\pgfqpoint{2.067679in}{1.614586in}}%
\pgfpathlineto{\pgfqpoint{2.063540in}{1.621479in}}%
\pgfpathlineto{\pgfqpoint{2.056624in}{1.621513in}}%
\pgfpathlineto{\pgfqpoint{2.056914in}{1.649093in}}%
\pgfpathlineto{\pgfqpoint{2.084271in}{1.648958in}}%
\pgfpathlineto{\pgfqpoint{2.111776in}{1.648912in}}%
\pgfpathlineto{\pgfqpoint{2.111772in}{1.614418in}}%
\pgfpathlineto{\pgfqpoint{2.115152in}{1.614399in}}%
\pgfpathlineto{\pgfqpoint{2.115090in}{1.593467in}}%
\pgfpathclose%
\pgfusepath{fill}%
\end{pgfscope}%
\begin{pgfscope}%
\pgfpathrectangle{\pgfqpoint{0.100000in}{0.100000in}}{\pgfqpoint{3.608454in}{2.310000in}}%
\pgfusepath{clip}%
\pgfsetbuttcap%
\pgfsetmiterjoin%
\definecolor{currentfill}{rgb}{0.000000,0.698039,0.650980}%
\pgfsetfillcolor{currentfill}%
\pgfsetlinewidth{0.000000pt}%
\definecolor{currentstroke}{rgb}{0.000000,0.000000,0.000000}%
\pgfsetstrokecolor{currentstroke}%
\pgfsetstrokeopacity{0.000000}%
\pgfsetdash{}{0pt}%
\pgfpathmoveto{\pgfqpoint{1.194163in}{1.451751in}}%
\pgfpathlineto{\pgfqpoint{1.249416in}{1.442841in}}%
\pgfpathlineto{\pgfqpoint{1.253203in}{1.445105in}}%
\pgfpathlineto{\pgfqpoint{1.251645in}{1.434868in}}%
\pgfpathlineto{\pgfqpoint{1.240533in}{1.367269in}}%
\pgfpathlineto{\pgfqpoint{1.181575in}{1.376785in}}%
\pgfpathlineto{\pgfqpoint{1.183363in}{1.379999in}}%
\pgfpathlineto{\pgfqpoint{1.179215in}{1.395609in}}%
\pgfpathlineto{\pgfqpoint{1.181137in}{1.398548in}}%
\pgfpathlineto{\pgfqpoint{1.177987in}{1.406555in}}%
\pgfpathlineto{\pgfqpoint{1.182094in}{1.423685in}}%
\pgfpathlineto{\pgfqpoint{1.189342in}{1.433929in}}%
\pgfpathlineto{\pgfqpoint{1.188609in}{1.438112in}}%
\pgfpathlineto{\pgfqpoint{1.192855in}{1.443180in}}%
\pgfpathlineto{\pgfqpoint{1.194163in}{1.451751in}}%
\pgfpathclose%
\pgfusepath{fill}%
\end{pgfscope}%
\begin{pgfscope}%
\pgfpathrectangle{\pgfqpoint{0.100000in}{0.100000in}}{\pgfqpoint{3.608454in}{2.310000in}}%
\pgfusepath{clip}%
\pgfsetbuttcap%
\pgfsetmiterjoin%
\definecolor{currentfill}{rgb}{0.000000,0.564706,0.717647}%
\pgfsetfillcolor{currentfill}%
\pgfsetlinewidth{0.000000pt}%
\definecolor{currentstroke}{rgb}{0.000000,0.000000,0.000000}%
\pgfsetstrokecolor{currentstroke}%
\pgfsetstrokeopacity{0.000000}%
\pgfsetdash{}{0pt}%
\pgfpathmoveto{\pgfqpoint{2.780892in}{1.025876in}}%
\pgfpathlineto{\pgfqpoint{2.779095in}{1.040030in}}%
\pgfpathlineto{\pgfqpoint{2.782890in}{1.048115in}}%
\pgfpathlineto{\pgfqpoint{2.780160in}{1.050381in}}%
\pgfpathlineto{\pgfqpoint{2.779057in}{1.060719in}}%
\pgfpathlineto{\pgfqpoint{2.798375in}{1.062974in}}%
\pgfpathlineto{\pgfqpoint{2.808303in}{1.061863in}}%
\pgfpathlineto{\pgfqpoint{2.814885in}{1.050035in}}%
\pgfpathlineto{\pgfqpoint{2.812171in}{1.037511in}}%
\pgfpathlineto{\pgfqpoint{2.810594in}{1.033632in}}%
\pgfpathlineto{\pgfqpoint{2.801686in}{1.031063in}}%
\pgfpathlineto{\pgfqpoint{2.799254in}{1.027968in}}%
\pgfpathlineto{\pgfqpoint{2.780892in}{1.025876in}}%
\pgfpathclose%
\pgfusepath{fill}%
\end{pgfscope}%
\begin{pgfscope}%
\pgfpathrectangle{\pgfqpoint{0.100000in}{0.100000in}}{\pgfqpoint{3.608454in}{2.310000in}}%
\pgfusepath{clip}%
\pgfsetbuttcap%
\pgfsetmiterjoin%
\definecolor{currentfill}{rgb}{0.000000,0.917647,0.541176}%
\pgfsetfillcolor{currentfill}%
\pgfsetlinewidth{0.000000pt}%
\definecolor{currentstroke}{rgb}{0.000000,0.000000,0.000000}%
\pgfsetstrokecolor{currentstroke}%
\pgfsetstrokeopacity{0.000000}%
\pgfsetdash{}{0pt}%
\pgfpathmoveto{\pgfqpoint{0.427334in}{1.927274in}}%
\pgfpathlineto{\pgfqpoint{0.440684in}{1.949326in}}%
\pgfpathlineto{\pgfqpoint{0.456163in}{1.968351in}}%
\pgfpathlineto{\pgfqpoint{0.467729in}{1.989714in}}%
\pgfpathlineto{\pgfqpoint{0.479271in}{1.985850in}}%
\pgfpathlineto{\pgfqpoint{0.487714in}{1.989945in}}%
\pgfpathlineto{\pgfqpoint{0.498412in}{1.984367in}}%
\pgfpathlineto{\pgfqpoint{0.501091in}{1.975934in}}%
\pgfpathlineto{\pgfqpoint{0.509770in}{1.969016in}}%
\pgfpathlineto{\pgfqpoint{0.521206in}{1.965301in}}%
\pgfpathlineto{\pgfqpoint{0.517042in}{1.952171in}}%
\pgfpathlineto{\pgfqpoint{0.517068in}{1.946672in}}%
\pgfpathlineto{\pgfqpoint{0.537209in}{1.940659in}}%
\pgfpathlineto{\pgfqpoint{0.534705in}{1.932520in}}%
\pgfpathlineto{\pgfqpoint{0.568235in}{1.922352in}}%
\pgfpathlineto{\pgfqpoint{0.572489in}{1.914545in}}%
\pgfpathlineto{\pgfqpoint{0.572519in}{1.906519in}}%
\pgfpathlineto{\pgfqpoint{0.565189in}{1.899442in}}%
\pgfpathlineto{\pgfqpoint{0.562002in}{1.893971in}}%
\pgfpathlineto{\pgfqpoint{0.551234in}{1.896670in}}%
\pgfpathlineto{\pgfqpoint{0.549594in}{1.891292in}}%
\pgfpathlineto{\pgfqpoint{0.543027in}{1.893318in}}%
\pgfpathlineto{\pgfqpoint{0.538396in}{1.890064in}}%
\pgfpathlineto{\pgfqpoint{0.532832in}{1.891746in}}%
\pgfpathlineto{\pgfqpoint{0.516046in}{1.883757in}}%
\pgfpathlineto{\pgfqpoint{0.503506in}{1.885466in}}%
\pgfpathlineto{\pgfqpoint{0.495724in}{1.888283in}}%
\pgfpathlineto{\pgfqpoint{0.491499in}{1.886385in}}%
\pgfpathlineto{\pgfqpoint{0.483261in}{1.890283in}}%
\pgfpathlineto{\pgfqpoint{0.478228in}{1.889504in}}%
\pgfpathlineto{\pgfqpoint{0.465159in}{1.900474in}}%
\pgfpathlineto{\pgfqpoint{0.456477in}{1.902564in}}%
\pgfpathlineto{\pgfqpoint{0.447212in}{1.898187in}}%
\pgfpathlineto{\pgfqpoint{0.438998in}{1.899672in}}%
\pgfpathlineto{\pgfqpoint{0.443285in}{1.914029in}}%
\pgfpathlineto{\pgfqpoint{0.439586in}{1.923181in}}%
\pgfpathlineto{\pgfqpoint{0.427334in}{1.927274in}}%
\pgfpathclose%
\pgfusepath{fill}%
\end{pgfscope}%
\begin{pgfscope}%
\pgfpathrectangle{\pgfqpoint{0.100000in}{0.100000in}}{\pgfqpoint{3.608454in}{2.310000in}}%
\pgfusepath{clip}%
\pgfsetbuttcap%
\pgfsetmiterjoin%
\definecolor{currentfill}{rgb}{0.000000,0.666667,0.666667}%
\pgfsetfillcolor{currentfill}%
\pgfsetlinewidth{0.000000pt}%
\definecolor{currentstroke}{rgb}{0.000000,0.000000,0.000000}%
\pgfsetstrokecolor{currentstroke}%
\pgfsetstrokeopacity{0.000000}%
\pgfsetdash{}{0pt}%
\pgfpathmoveto{\pgfqpoint{3.106445in}{0.607053in}}%
\pgfpathlineto{\pgfqpoint{3.090036in}{0.604447in}}%
\pgfpathlineto{\pgfqpoint{3.090875in}{0.590177in}}%
\pgfpathlineto{\pgfqpoint{3.080990in}{0.605844in}}%
\pgfpathlineto{\pgfqpoint{3.078312in}{0.602797in}}%
\pgfpathlineto{\pgfqpoint{3.071577in}{0.605034in}}%
\pgfpathlineto{\pgfqpoint{3.065010in}{0.604224in}}%
\pgfpathlineto{\pgfqpoint{3.060692in}{0.615927in}}%
\pgfpathlineto{\pgfqpoint{3.049410in}{0.623089in}}%
\pgfpathlineto{\pgfqpoint{3.046213in}{0.629170in}}%
\pgfpathlineto{\pgfqpoint{3.040376in}{0.630406in}}%
\pgfpathlineto{\pgfqpoint{3.037437in}{0.635958in}}%
\pgfpathlineto{\pgfqpoint{3.029457in}{0.642400in}}%
\pgfpathlineto{\pgfqpoint{3.025702in}{0.651277in}}%
\pgfpathlineto{\pgfqpoint{3.020704in}{0.653264in}}%
\pgfpathlineto{\pgfqpoint{3.006579in}{0.647225in}}%
\pgfpathlineto{\pgfqpoint{3.004229in}{0.666876in}}%
\pgfpathlineto{\pgfqpoint{3.011416in}{0.670241in}}%
\pgfpathlineto{\pgfqpoint{3.019158in}{0.678567in}}%
\pgfpathlineto{\pgfqpoint{3.035034in}{0.681162in}}%
\pgfpathlineto{\pgfqpoint{3.039796in}{0.675377in}}%
\pgfpathlineto{\pgfqpoint{3.041523in}{0.664546in}}%
\pgfpathlineto{\pgfqpoint{3.055227in}{0.666862in}}%
\pgfpathlineto{\pgfqpoint{3.062330in}{0.671592in}}%
\pgfpathlineto{\pgfqpoint{3.078208in}{0.644633in}}%
\pgfpathlineto{\pgfqpoint{3.086273in}{0.632061in}}%
\pgfpathlineto{\pgfqpoint{3.106445in}{0.607053in}}%
\pgfpathclose%
\pgfusepath{fill}%
\end{pgfscope}%
\begin{pgfscope}%
\pgfpathrectangle{\pgfqpoint{0.100000in}{0.100000in}}{\pgfqpoint{3.608454in}{2.310000in}}%
\pgfusepath{clip}%
\pgfsetbuttcap%
\pgfsetmiterjoin%
\definecolor{currentfill}{rgb}{0.000000,0.549020,0.725490}%
\pgfsetfillcolor{currentfill}%
\pgfsetlinewidth{0.000000pt}%
\definecolor{currentstroke}{rgb}{0.000000,0.000000,0.000000}%
\pgfsetstrokecolor{currentstroke}%
\pgfsetstrokeopacity{0.000000}%
\pgfsetdash{}{0pt}%
\pgfpathmoveto{\pgfqpoint{0.914070in}{1.089859in}}%
\pgfpathlineto{\pgfqpoint{0.900636in}{1.091652in}}%
\pgfpathlineto{\pgfqpoint{0.895182in}{1.086583in}}%
\pgfpathlineto{\pgfqpoint{0.868397in}{1.094701in}}%
\pgfpathlineto{\pgfqpoint{0.862231in}{1.099514in}}%
\pgfpathlineto{\pgfqpoint{0.851945in}{1.113704in}}%
\pgfpathlineto{\pgfqpoint{0.848041in}{1.126325in}}%
\pgfpathlineto{\pgfqpoint{0.847819in}{1.135796in}}%
\pgfpathlineto{\pgfqpoint{0.843543in}{1.141162in}}%
\pgfpathlineto{\pgfqpoint{0.840274in}{1.151795in}}%
\pgfpathlineto{\pgfqpoint{0.842429in}{1.160488in}}%
\pgfpathlineto{\pgfqpoint{0.792611in}{1.237224in}}%
\pgfpathlineto{\pgfqpoint{0.780428in}{1.255922in}}%
\pgfpathlineto{\pgfqpoint{0.720005in}{1.349216in}}%
\pgfpathlineto{\pgfqpoint{0.688870in}{1.397366in}}%
\pgfpathlineto{\pgfqpoint{0.661400in}{1.439746in}}%
\pgfpathlineto{\pgfqpoint{0.665974in}{1.438337in}}%
\pgfpathlineto{\pgfqpoint{0.717194in}{1.472445in}}%
\pgfpathlineto{\pgfqpoint{0.695807in}{1.514308in}}%
\pgfpathlineto{\pgfqpoint{0.697419in}{1.520493in}}%
\pgfpathlineto{\pgfqpoint{0.718561in}{1.521037in}}%
\pgfpathlineto{\pgfqpoint{0.724276in}{1.521184in}}%
\pgfpathlineto{\pgfqpoint{0.752002in}{1.519867in}}%
\pgfpathlineto{\pgfqpoint{0.795453in}{1.508995in}}%
\pgfpathlineto{\pgfqpoint{0.836825in}{1.499170in}}%
\pgfpathlineto{\pgfqpoint{0.882599in}{1.449402in}}%
\pgfpathlineto{\pgfqpoint{0.939931in}{1.436894in}}%
\pgfpathlineto{\pgfqpoint{0.938194in}{1.428823in}}%
\pgfpathlineto{\pgfqpoint{0.925741in}{1.369767in}}%
\pgfpathlineto{\pgfqpoint{0.912295in}{1.307222in}}%
\pgfpathlineto{\pgfqpoint{0.983512in}{1.292540in}}%
\pgfpathlineto{\pgfqpoint{1.005762in}{1.288169in}}%
\pgfpathlineto{\pgfqpoint{0.994707in}{1.264046in}}%
\pgfpathlineto{\pgfqpoint{0.991172in}{1.251855in}}%
\pgfpathlineto{\pgfqpoint{0.991491in}{1.243095in}}%
\pgfpathlineto{\pgfqpoint{0.986930in}{1.238736in}}%
\pgfpathlineto{\pgfqpoint{0.978736in}{1.235692in}}%
\pgfpathlineto{\pgfqpoint{0.961661in}{1.233095in}}%
\pgfpathlineto{\pgfqpoint{0.951426in}{1.230336in}}%
\pgfpathlineto{\pgfqpoint{0.949023in}{1.226105in}}%
\pgfpathlineto{\pgfqpoint{0.943617in}{1.227923in}}%
\pgfpathlineto{\pgfqpoint{0.940117in}{1.223821in}}%
\pgfpathlineto{\pgfqpoint{0.941162in}{1.217180in}}%
\pgfpathlineto{\pgfqpoint{0.937616in}{1.205156in}}%
\pgfpathlineto{\pgfqpoint{0.933245in}{1.183825in}}%
\pgfpathlineto{\pgfqpoint{0.914070in}{1.089859in}}%
\pgfpathclose%
\pgfusepath{fill}%
\end{pgfscope}%
\begin{pgfscope}%
\pgfpathrectangle{\pgfqpoint{0.100000in}{0.100000in}}{\pgfqpoint{3.608454in}{2.310000in}}%
\pgfusepath{clip}%
\pgfsetbuttcap%
\pgfsetmiterjoin%
\definecolor{currentfill}{rgb}{0.000000,0.552941,0.723529}%
\pgfsetfillcolor{currentfill}%
\pgfsetlinewidth{0.000000pt}%
\definecolor{currentstroke}{rgb}{0.000000,0.000000,0.000000}%
\pgfsetstrokecolor{currentstroke}%
\pgfsetstrokeopacity{0.000000}%
\pgfsetdash{}{0pt}%
\pgfpathmoveto{\pgfqpoint{2.956251in}{1.497406in}}%
\pgfpathlineto{\pgfqpoint{2.935644in}{1.495065in}}%
\pgfpathlineto{\pgfqpoint{2.928425in}{1.497499in}}%
\pgfpathlineto{\pgfqpoint{2.908656in}{1.495198in}}%
\pgfpathlineto{\pgfqpoint{2.909284in}{1.489578in}}%
\pgfpathlineto{\pgfqpoint{2.878220in}{1.486351in}}%
\pgfpathlineto{\pgfqpoint{2.876912in}{1.491928in}}%
\pgfpathlineto{\pgfqpoint{2.874225in}{1.518225in}}%
\pgfpathlineto{\pgfqpoint{2.871819in}{1.517489in}}%
\pgfpathlineto{\pgfqpoint{2.870685in}{1.525264in}}%
\pgfpathlineto{\pgfqpoint{2.876286in}{1.526111in}}%
\pgfpathlineto{\pgfqpoint{2.872479in}{1.551436in}}%
\pgfpathlineto{\pgfqpoint{2.901003in}{1.555251in}}%
\pgfpathlineto{\pgfqpoint{2.901826in}{1.549489in}}%
\pgfpathlineto{\pgfqpoint{2.915375in}{1.550896in}}%
\pgfpathlineto{\pgfqpoint{2.914636in}{1.557327in}}%
\pgfpathlineto{\pgfqpoint{2.934085in}{1.560268in}}%
\pgfpathlineto{\pgfqpoint{2.937196in}{1.539921in}}%
\pgfpathlineto{\pgfqpoint{2.947280in}{1.541412in}}%
\pgfpathlineto{\pgfqpoint{2.948242in}{1.534882in}}%
\pgfpathlineto{\pgfqpoint{2.952109in}{1.531961in}}%
\pgfpathlineto{\pgfqpoint{2.949140in}{1.528048in}}%
\pgfpathlineto{\pgfqpoint{2.949522in}{1.517575in}}%
\pgfpathlineto{\pgfqpoint{2.954062in}{1.518144in}}%
\pgfpathlineto{\pgfqpoint{2.956251in}{1.497406in}}%
\pgfpathclose%
\pgfusepath{fill}%
\end{pgfscope}%
\begin{pgfscope}%
\pgfpathrectangle{\pgfqpoint{0.100000in}{0.100000in}}{\pgfqpoint{3.608454in}{2.310000in}}%
\pgfusepath{clip}%
\pgfsetbuttcap%
\pgfsetmiterjoin%
\definecolor{currentfill}{rgb}{0.000000,0.658824,0.670588}%
\pgfsetfillcolor{currentfill}%
\pgfsetlinewidth{0.000000pt}%
\definecolor{currentstroke}{rgb}{0.000000,0.000000,0.000000}%
\pgfsetstrokecolor{currentstroke}%
\pgfsetstrokeopacity{0.000000}%
\pgfsetdash{}{0pt}%
\pgfpathmoveto{\pgfqpoint{2.996900in}{1.203053in}}%
\pgfpathlineto{\pgfqpoint{3.004020in}{1.198138in}}%
\pgfpathlineto{\pgfqpoint{3.006498in}{1.176840in}}%
\pgfpathlineto{\pgfqpoint{2.993174in}{1.172573in}}%
\pgfpathlineto{\pgfqpoint{2.987201in}{1.172347in}}%
\pgfpathlineto{\pgfqpoint{2.978982in}{1.167861in}}%
\pgfpathlineto{\pgfqpoint{2.972912in}{1.172989in}}%
\pgfpathlineto{\pgfqpoint{2.963826in}{1.175254in}}%
\pgfpathlineto{\pgfqpoint{2.968126in}{1.182900in}}%
\pgfpathlineto{\pgfqpoint{2.955269in}{1.192553in}}%
\pgfpathlineto{\pgfqpoint{2.948361in}{1.195220in}}%
\pgfpathlineto{\pgfqpoint{2.950303in}{1.201290in}}%
\pgfpathlineto{\pgfqpoint{2.950033in}{1.210864in}}%
\pgfpathlineto{\pgfqpoint{2.970559in}{1.212860in}}%
\pgfpathlineto{\pgfqpoint{2.970773in}{1.210739in}}%
\pgfpathlineto{\pgfqpoint{2.980945in}{1.198576in}}%
\pgfpathlineto{\pgfqpoint{2.987922in}{1.204284in}}%
\pgfpathlineto{\pgfqpoint{2.996900in}{1.203053in}}%
\pgfpathclose%
\pgfusepath{fill}%
\end{pgfscope}%
\begin{pgfscope}%
\pgfpathrectangle{\pgfqpoint{0.100000in}{0.100000in}}{\pgfqpoint{3.608454in}{2.310000in}}%
\pgfusepath{clip}%
\pgfsetbuttcap%
\pgfsetmiterjoin%
\definecolor{currentfill}{rgb}{0.000000,0.364706,0.817647}%
\pgfsetfillcolor{currentfill}%
\pgfsetlinewidth{0.000000pt}%
\definecolor{currentstroke}{rgb}{0.000000,0.000000,0.000000}%
\pgfsetstrokecolor{currentstroke}%
\pgfsetstrokeopacity{0.000000}%
\pgfsetdash{}{0pt}%
\pgfpathmoveto{\pgfqpoint{1.655822in}{1.596346in}}%
\pgfpathlineto{\pgfqpoint{1.691656in}{1.593744in}}%
\pgfpathlineto{\pgfqpoint{1.689939in}{1.572676in}}%
\pgfpathlineto{\pgfqpoint{1.694791in}{1.572316in}}%
\pgfpathlineto{\pgfqpoint{1.692731in}{1.544894in}}%
\pgfpathlineto{\pgfqpoint{1.688556in}{1.545253in}}%
\pgfpathlineto{\pgfqpoint{1.687474in}{1.531563in}}%
\pgfpathlineto{\pgfqpoint{1.654575in}{1.534236in}}%
\pgfpathlineto{\pgfqpoint{1.655781in}{1.547900in}}%
\pgfpathlineto{\pgfqpoint{1.654457in}{1.551499in}}%
\pgfpathlineto{\pgfqpoint{1.656181in}{1.573226in}}%
\pgfpathlineto{\pgfqpoint{1.654522in}{1.581965in}}%
\pgfpathlineto{\pgfqpoint{1.655822in}{1.596346in}}%
\pgfpathclose%
\pgfusepath{fill}%
\end{pgfscope}%
\begin{pgfscope}%
\pgfpathrectangle{\pgfqpoint{0.100000in}{0.100000in}}{\pgfqpoint{3.608454in}{2.310000in}}%
\pgfusepath{clip}%
\pgfsetbuttcap%
\pgfsetmiterjoin%
\definecolor{currentfill}{rgb}{0.000000,0.615686,0.692157}%
\pgfsetfillcolor{currentfill}%
\pgfsetlinewidth{0.000000pt}%
\definecolor{currentstroke}{rgb}{0.000000,0.000000,0.000000}%
\pgfsetstrokecolor{currentstroke}%
\pgfsetstrokeopacity{0.000000}%
\pgfsetdash{}{0pt}%
\pgfpathmoveto{\pgfqpoint{2.722959in}{1.184982in}}%
\pgfpathlineto{\pgfqpoint{2.722324in}{1.177007in}}%
\pgfpathlineto{\pgfqpoint{2.713741in}{1.174022in}}%
\pgfpathlineto{\pgfqpoint{2.710970in}{1.165941in}}%
\pgfpathlineto{\pgfqpoint{2.705300in}{1.170782in}}%
\pgfpathlineto{\pgfqpoint{2.696790in}{1.173951in}}%
\pgfpathlineto{\pgfqpoint{2.690027in}{1.171438in}}%
\pgfpathlineto{\pgfqpoint{2.690217in}{1.180528in}}%
\pgfpathlineto{\pgfqpoint{2.678723in}{1.180579in}}%
\pgfpathlineto{\pgfqpoint{2.677443in}{1.189404in}}%
\pgfpathlineto{\pgfqpoint{2.664358in}{1.203424in}}%
\pgfpathlineto{\pgfqpoint{2.669252in}{1.213390in}}%
\pgfpathlineto{\pgfqpoint{2.669571in}{1.226318in}}%
\pgfpathlineto{\pgfqpoint{2.661850in}{1.235103in}}%
\pgfpathlineto{\pgfqpoint{2.665870in}{1.236725in}}%
\pgfpathlineto{\pgfqpoint{2.667849in}{1.244767in}}%
\pgfpathlineto{\pgfqpoint{2.692446in}{1.245040in}}%
\pgfpathlineto{\pgfqpoint{2.690908in}{1.235308in}}%
\pgfpathlineto{\pgfqpoint{2.692552in}{1.225974in}}%
\pgfpathlineto{\pgfqpoint{2.696385in}{1.221833in}}%
\pgfpathlineto{\pgfqpoint{2.703168in}{1.221287in}}%
\pgfpathlineto{\pgfqpoint{2.705489in}{1.213403in}}%
\pgfpathlineto{\pgfqpoint{2.709289in}{1.208321in}}%
\pgfpathlineto{\pgfqpoint{2.723224in}{1.208766in}}%
\pgfpathlineto{\pgfqpoint{2.724860in}{1.203323in}}%
\pgfpathlineto{\pgfqpoint{2.722959in}{1.184982in}}%
\pgfpathclose%
\pgfusepath{fill}%
\end{pgfscope}%
\begin{pgfscope}%
\pgfpathrectangle{\pgfqpoint{0.100000in}{0.100000in}}{\pgfqpoint{3.608454in}{2.310000in}}%
\pgfusepath{clip}%
\pgfsetbuttcap%
\pgfsetmiterjoin%
\definecolor{currentfill}{rgb}{0.000000,0.505882,0.747059}%
\pgfsetfillcolor{currentfill}%
\pgfsetlinewidth{0.000000pt}%
\definecolor{currentstroke}{rgb}{0.000000,0.000000,0.000000}%
\pgfsetstrokecolor{currentstroke}%
\pgfsetstrokeopacity{0.000000}%
\pgfsetdash{}{0pt}%
\pgfpathmoveto{\pgfqpoint{3.206484in}{1.172807in}}%
\pgfpathlineto{\pgfqpoint{3.191049in}{1.170900in}}%
\pgfpathlineto{\pgfqpoint{3.186933in}{1.171546in}}%
\pgfpathlineto{\pgfqpoint{3.181671in}{1.180528in}}%
\pgfpathlineto{\pgfqpoint{3.175634in}{1.186715in}}%
\pgfpathlineto{\pgfqpoint{3.185638in}{1.219606in}}%
\pgfpathlineto{\pgfqpoint{3.195538in}{1.217166in}}%
\pgfpathlineto{\pgfqpoint{3.206943in}{1.219779in}}%
\pgfpathlineto{\pgfqpoint{3.216999in}{1.216602in}}%
\pgfpathlineto{\pgfqpoint{3.218220in}{1.212547in}}%
\pgfpathlineto{\pgfqpoint{3.226514in}{1.211542in}}%
\pgfpathlineto{\pgfqpoint{3.231820in}{1.204552in}}%
\pgfpathlineto{\pgfqpoint{3.232624in}{1.197760in}}%
\pgfpathlineto{\pgfqpoint{3.229923in}{1.198333in}}%
\pgfpathlineto{\pgfqpoint{3.206484in}{1.172807in}}%
\pgfpathclose%
\pgfusepath{fill}%
\end{pgfscope}%
\begin{pgfscope}%
\pgfpathrectangle{\pgfqpoint{0.100000in}{0.100000in}}{\pgfqpoint{3.608454in}{2.310000in}}%
\pgfusepath{clip}%
\pgfsetbuttcap%
\pgfsetmiterjoin%
\definecolor{currentfill}{rgb}{0.000000,0.568627,0.715686}%
\pgfsetfillcolor{currentfill}%
\pgfsetlinewidth{0.000000pt}%
\definecolor{currentstroke}{rgb}{0.000000,0.000000,0.000000}%
\pgfsetstrokecolor{currentstroke}%
\pgfsetstrokeopacity{0.000000}%
\pgfsetdash{}{0pt}%
\pgfpathmoveto{\pgfqpoint{1.051274in}{1.732636in}}%
\pgfpathlineto{\pgfqpoint{1.054206in}{1.729137in}}%
\pgfpathlineto{\pgfqpoint{1.064066in}{1.727041in}}%
\pgfpathlineto{\pgfqpoint{1.060120in}{1.706792in}}%
\pgfpathlineto{\pgfqpoint{1.054992in}{1.681239in}}%
\pgfpathlineto{\pgfqpoint{1.004416in}{1.690817in}}%
\pgfpathlineto{\pgfqpoint{0.994964in}{1.693281in}}%
\pgfpathlineto{\pgfqpoint{0.981161in}{1.696275in}}%
\pgfpathlineto{\pgfqpoint{0.988011in}{1.729060in}}%
\pgfpathlineto{\pgfqpoint{1.000366in}{1.726400in}}%
\pgfpathlineto{\pgfqpoint{1.002092in}{1.733915in}}%
\pgfpathlineto{\pgfqpoint{1.010156in}{1.733884in}}%
\pgfpathlineto{\pgfqpoint{1.013857in}{1.751700in}}%
\pgfpathlineto{\pgfqpoint{1.023421in}{1.749624in}}%
\pgfpathlineto{\pgfqpoint{1.027689in}{1.755650in}}%
\pgfpathlineto{\pgfqpoint{1.033321in}{1.782826in}}%
\pgfpathlineto{\pgfqpoint{1.050343in}{1.779280in}}%
\pgfpathlineto{\pgfqpoint{1.044792in}{1.752067in}}%
\pgfpathlineto{\pgfqpoint{1.041428in}{1.752757in}}%
\pgfpathlineto{\pgfqpoint{1.038549in}{1.738796in}}%
\pgfpathlineto{\pgfqpoint{1.044474in}{1.739072in}}%
\pgfpathlineto{\pgfqpoint{1.051274in}{1.732636in}}%
\pgfpathclose%
\pgfusepath{fill}%
\end{pgfscope}%
\begin{pgfscope}%
\pgfpathrectangle{\pgfqpoint{0.100000in}{0.100000in}}{\pgfqpoint{3.608454in}{2.310000in}}%
\pgfusepath{clip}%
\pgfsetbuttcap%
\pgfsetmiterjoin%
\definecolor{currentfill}{rgb}{0.000000,0.443137,0.778431}%
\pgfsetfillcolor{currentfill}%
\pgfsetlinewidth{0.000000pt}%
\definecolor{currentstroke}{rgb}{0.000000,0.000000,0.000000}%
\pgfsetstrokecolor{currentstroke}%
\pgfsetstrokeopacity{0.000000}%
\pgfsetdash{}{0pt}%
\pgfpathmoveto{\pgfqpoint{2.723904in}{1.022807in}}%
\pgfpathlineto{\pgfqpoint{2.715608in}{1.053498in}}%
\pgfpathlineto{\pgfqpoint{2.698781in}{1.052062in}}%
\pgfpathlineto{\pgfqpoint{2.696968in}{1.070359in}}%
\pgfpathlineto{\pgfqpoint{2.693282in}{1.075443in}}%
\pgfpathlineto{\pgfqpoint{2.694900in}{1.080379in}}%
\pgfpathlineto{\pgfqpoint{2.693772in}{1.094220in}}%
\pgfpathlineto{\pgfqpoint{2.697680in}{1.093767in}}%
\pgfpathlineto{\pgfqpoint{2.714180in}{1.097018in}}%
\pgfpathlineto{\pgfqpoint{2.722462in}{1.100629in}}%
\pgfpathlineto{\pgfqpoint{2.724406in}{1.094850in}}%
\pgfpathlineto{\pgfqpoint{2.736951in}{1.085202in}}%
\pgfpathlineto{\pgfqpoint{2.739787in}{1.093594in}}%
\pgfpathlineto{\pgfqpoint{2.749929in}{1.091001in}}%
\pgfpathlineto{\pgfqpoint{2.755774in}{1.081517in}}%
\pgfpathlineto{\pgfqpoint{2.751742in}{1.068739in}}%
\pgfpathlineto{\pgfqpoint{2.756131in}{1.058082in}}%
\pgfpathlineto{\pgfqpoint{2.752122in}{1.044231in}}%
\pgfpathlineto{\pgfqpoint{2.747587in}{1.043751in}}%
\pgfpathlineto{\pgfqpoint{2.746037in}{1.035657in}}%
\pgfpathlineto{\pgfqpoint{2.753682in}{1.036469in}}%
\pgfpathlineto{\pgfqpoint{2.754206in}{1.025936in}}%
\pgfpathlineto{\pgfqpoint{2.723904in}{1.022807in}}%
\pgfpathclose%
\pgfusepath{fill}%
\end{pgfscope}%
\begin{pgfscope}%
\pgfpathrectangle{\pgfqpoint{0.100000in}{0.100000in}}{\pgfqpoint{3.608454in}{2.310000in}}%
\pgfusepath{clip}%
\pgfsetbuttcap%
\pgfsetmiterjoin%
\definecolor{currentfill}{rgb}{0.000000,0.466667,0.766667}%
\pgfsetfillcolor{currentfill}%
\pgfsetlinewidth{0.000000pt}%
\definecolor{currentstroke}{rgb}{0.000000,0.000000,0.000000}%
\pgfsetstrokecolor{currentstroke}%
\pgfsetstrokeopacity{0.000000}%
\pgfsetdash{}{0pt}%
\pgfpathmoveto{\pgfqpoint{2.667350in}{1.499209in}}%
\pgfpathlineto{\pgfqpoint{2.653860in}{1.497742in}}%
\pgfpathlineto{\pgfqpoint{2.649395in}{1.497238in}}%
\pgfpathlineto{\pgfqpoint{2.646259in}{1.524584in}}%
\pgfpathlineto{\pgfqpoint{2.633101in}{1.523054in}}%
\pgfpathlineto{\pgfqpoint{2.633720in}{1.516230in}}%
\pgfpathlineto{\pgfqpoint{2.615905in}{1.514539in}}%
\pgfpathlineto{\pgfqpoint{2.614051in}{1.535158in}}%
\pgfpathlineto{\pgfqpoint{2.638514in}{1.537637in}}%
\pgfpathlineto{\pgfqpoint{2.636183in}{1.558287in}}%
\pgfpathlineto{\pgfqpoint{2.660092in}{1.560783in}}%
\pgfpathlineto{\pgfqpoint{2.661344in}{1.549614in}}%
\pgfpathlineto{\pgfqpoint{2.660236in}{1.540245in}}%
\pgfpathlineto{\pgfqpoint{2.661863in}{1.526375in}}%
\pgfpathlineto{\pgfqpoint{2.664197in}{1.526649in}}%
\pgfpathlineto{\pgfqpoint{2.667350in}{1.499209in}}%
\pgfpathclose%
\pgfusepath{fill}%
\end{pgfscope}%
\begin{pgfscope}%
\pgfpathrectangle{\pgfqpoint{0.100000in}{0.100000in}}{\pgfqpoint{3.608454in}{2.310000in}}%
\pgfusepath{clip}%
\pgfsetbuttcap%
\pgfsetmiterjoin%
\definecolor{currentfill}{rgb}{0.000000,0.690196,0.654902}%
\pgfsetfillcolor{currentfill}%
\pgfsetlinewidth{0.000000pt}%
\definecolor{currentstroke}{rgb}{0.000000,0.000000,0.000000}%
\pgfsetstrokecolor{currentstroke}%
\pgfsetstrokeopacity{0.000000}%
\pgfsetdash{}{0pt}%
\pgfpathmoveto{\pgfqpoint{0.429033in}{1.527264in}}%
\pgfpathlineto{\pgfqpoint{0.422980in}{1.525827in}}%
\pgfpathlineto{\pgfqpoint{0.419978in}{1.533247in}}%
\pgfpathlineto{\pgfqpoint{0.410415in}{1.537899in}}%
\pgfpathlineto{\pgfqpoint{0.403577in}{1.549156in}}%
\pgfpathlineto{\pgfqpoint{0.396978in}{1.549508in}}%
\pgfpathlineto{\pgfqpoint{0.393248in}{1.563598in}}%
\pgfpathlineto{\pgfqpoint{0.387572in}{1.570213in}}%
\pgfpathlineto{\pgfqpoint{0.382759in}{1.579907in}}%
\pgfpathlineto{\pgfqpoint{0.380186in}{1.589872in}}%
\pgfpathlineto{\pgfqpoint{0.372415in}{1.602800in}}%
\pgfpathlineto{\pgfqpoint{0.369555in}{1.609925in}}%
\pgfpathlineto{\pgfqpoint{0.374700in}{1.618945in}}%
\pgfpathlineto{\pgfqpoint{0.373785in}{1.631630in}}%
\pgfpathlineto{\pgfqpoint{0.374411in}{1.644351in}}%
\pgfpathlineto{\pgfqpoint{0.377409in}{1.651582in}}%
\pgfpathlineto{\pgfqpoint{0.382897in}{1.658786in}}%
\pgfpathlineto{\pgfqpoint{0.384563in}{1.669377in}}%
\pgfpathlineto{\pgfqpoint{0.384726in}{1.681179in}}%
\pgfpathlineto{\pgfqpoint{0.379033in}{1.697172in}}%
\pgfpathlineto{\pgfqpoint{0.406652in}{1.688065in}}%
\pgfpathlineto{\pgfqpoint{0.406038in}{1.686194in}}%
\pgfpathlineto{\pgfqpoint{0.441276in}{1.674970in}}%
\pgfpathlineto{\pgfqpoint{0.436761in}{1.661477in}}%
\pgfpathlineto{\pgfqpoint{0.437233in}{1.653925in}}%
\pgfpathlineto{\pgfqpoint{0.434583in}{1.644086in}}%
\pgfpathlineto{\pgfqpoint{0.443292in}{1.641381in}}%
\pgfpathlineto{\pgfqpoint{0.440384in}{1.631531in}}%
\pgfpathlineto{\pgfqpoint{0.435002in}{1.619433in}}%
\pgfpathlineto{\pgfqpoint{0.440654in}{1.612562in}}%
\pgfpathlineto{\pgfqpoint{0.447344in}{1.609339in}}%
\pgfpathlineto{\pgfqpoint{0.445039in}{1.597163in}}%
\pgfpathlineto{\pgfqpoint{0.448894in}{1.592474in}}%
\pgfpathlineto{\pgfqpoint{0.450911in}{1.584637in}}%
\pgfpathlineto{\pgfqpoint{0.446281in}{1.581110in}}%
\pgfpathlineto{\pgfqpoint{0.446440in}{1.577149in}}%
\pgfpathlineto{\pgfqpoint{0.438495in}{1.570367in}}%
\pgfpathlineto{\pgfqpoint{0.427977in}{1.570534in}}%
\pgfpathlineto{\pgfqpoint{0.425226in}{1.565689in}}%
\pgfpathlineto{\pgfqpoint{0.431374in}{1.550880in}}%
\pgfpathlineto{\pgfqpoint{0.433169in}{1.538861in}}%
\pgfpathlineto{\pgfqpoint{0.432587in}{1.529018in}}%
\pgfpathlineto{\pgfqpoint{0.429033in}{1.527264in}}%
\pgfpathclose%
\pgfusepath{fill}%
\end{pgfscope}%
\begin{pgfscope}%
\pgfpathrectangle{\pgfqpoint{0.100000in}{0.100000in}}{\pgfqpoint{3.608454in}{2.310000in}}%
\pgfusepath{clip}%
\pgfsetbuttcap%
\pgfsetmiterjoin%
\definecolor{currentfill}{rgb}{0.000000,0.686275,0.656863}%
\pgfsetfillcolor{currentfill}%
\pgfsetlinewidth{0.000000pt}%
\definecolor{currentstroke}{rgb}{0.000000,0.000000,0.000000}%
\pgfsetstrokecolor{currentstroke}%
\pgfsetstrokeopacity{0.000000}%
\pgfsetdash{}{0pt}%
\pgfpathmoveto{\pgfqpoint{2.929879in}{1.293625in}}%
\pgfpathlineto{\pgfqpoint{2.916184in}{1.292250in}}%
\pgfpathlineto{\pgfqpoint{2.909448in}{1.297631in}}%
\pgfpathlineto{\pgfqpoint{2.908038in}{1.306264in}}%
\pgfpathlineto{\pgfqpoint{2.902369in}{1.314765in}}%
\pgfpathlineto{\pgfqpoint{2.895571in}{1.311319in}}%
\pgfpathlineto{\pgfqpoint{2.900793in}{1.318863in}}%
\pgfpathlineto{\pgfqpoint{2.894629in}{1.327193in}}%
\pgfpathlineto{\pgfqpoint{2.894444in}{1.332079in}}%
\pgfpathlineto{\pgfqpoint{2.898128in}{1.341949in}}%
\pgfpathlineto{\pgfqpoint{2.906738in}{1.347296in}}%
\pgfpathlineto{\pgfqpoint{2.905122in}{1.355064in}}%
\pgfpathlineto{\pgfqpoint{2.903048in}{1.365422in}}%
\pgfpathlineto{\pgfqpoint{2.909141in}{1.370242in}}%
\pgfpathlineto{\pgfqpoint{2.920044in}{1.373705in}}%
\pgfpathlineto{\pgfqpoint{2.930069in}{1.365085in}}%
\pgfpathlineto{\pgfqpoint{2.936140in}{1.370645in}}%
\pgfpathlineto{\pgfqpoint{2.940346in}{1.366252in}}%
\pgfpathlineto{\pgfqpoint{2.957191in}{1.367025in}}%
\pgfpathlineto{\pgfqpoint{2.965422in}{1.379505in}}%
\pgfpathlineto{\pgfqpoint{2.972924in}{1.373592in}}%
\pgfpathlineto{\pgfqpoint{2.976752in}{1.368166in}}%
\pgfpathlineto{\pgfqpoint{2.975847in}{1.360884in}}%
\pgfpathlineto{\pgfqpoint{2.958041in}{1.346022in}}%
\pgfpathlineto{\pgfqpoint{2.953019in}{1.337104in}}%
\pgfpathlineto{\pgfqpoint{2.952368in}{1.321633in}}%
\pgfpathlineto{\pgfqpoint{2.944352in}{1.321971in}}%
\pgfpathlineto{\pgfqpoint{2.941096in}{1.316634in}}%
\pgfpathlineto{\pgfqpoint{2.946226in}{1.306488in}}%
\pgfpathlineto{\pgfqpoint{2.931210in}{1.301107in}}%
\pgfpathlineto{\pgfqpoint{2.935674in}{1.297744in}}%
\pgfpathlineto{\pgfqpoint{2.929879in}{1.293625in}}%
\pgfpathclose%
\pgfusepath{fill}%
\end{pgfscope}%
\begin{pgfscope}%
\pgfpathrectangle{\pgfqpoint{0.100000in}{0.100000in}}{\pgfqpoint{3.608454in}{2.310000in}}%
\pgfusepath{clip}%
\pgfsetbuttcap%
\pgfsetmiterjoin%
\definecolor{currentfill}{rgb}{0.000000,0.490196,0.754902}%
\pgfsetfillcolor{currentfill}%
\pgfsetlinewidth{0.000000pt}%
\definecolor{currentstroke}{rgb}{0.000000,0.000000,0.000000}%
\pgfsetstrokecolor{currentstroke}%
\pgfsetstrokeopacity{0.000000}%
\pgfsetdash{}{0pt}%
\pgfpathmoveto{\pgfqpoint{1.538610in}{1.726627in}}%
\pgfpathlineto{\pgfqpoint{1.532653in}{1.671867in}}%
\pgfpathlineto{\pgfqpoint{1.531372in}{1.656425in}}%
\pgfpathlineto{\pgfqpoint{1.508517in}{1.658724in}}%
\pgfpathlineto{\pgfqpoint{1.507108in}{1.644984in}}%
\pgfpathlineto{\pgfqpoint{1.504277in}{1.645283in}}%
\pgfpathlineto{\pgfqpoint{1.500514in}{1.634480in}}%
\pgfpathlineto{\pgfqpoint{1.486783in}{1.637214in}}%
\pgfpathlineto{\pgfqpoint{1.487410in}{1.642958in}}%
\pgfpathlineto{\pgfqpoint{1.491303in}{1.646739in}}%
\pgfpathlineto{\pgfqpoint{1.461309in}{1.650537in}}%
\pgfpathlineto{\pgfqpoint{1.429302in}{1.654336in}}%
\pgfpathlineto{\pgfqpoint{1.377332in}{1.661590in}}%
\pgfpathlineto{\pgfqpoint{1.376226in}{1.661744in}}%
\pgfpathlineto{\pgfqpoint{1.379922in}{1.688990in}}%
\pgfpathlineto{\pgfqpoint{1.382345in}{1.688646in}}%
\pgfpathlineto{\pgfqpoint{1.386059in}{1.715853in}}%
\pgfpathlineto{\pgfqpoint{1.388966in}{1.743023in}}%
\pgfpathlineto{\pgfqpoint{1.388262in}{1.745418in}}%
\pgfpathlineto{\pgfqpoint{1.412417in}{1.742036in}}%
\pgfpathlineto{\pgfqpoint{1.441739in}{1.737485in}}%
\pgfpathlineto{\pgfqpoint{1.474731in}{1.733639in}}%
\pgfpathlineto{\pgfqpoint{1.538610in}{1.726627in}}%
\pgfpathclose%
\pgfusepath{fill}%
\end{pgfscope}%
\begin{pgfscope}%
\pgfpathrectangle{\pgfqpoint{0.100000in}{0.100000in}}{\pgfqpoint{3.608454in}{2.310000in}}%
\pgfusepath{clip}%
\pgfsetbuttcap%
\pgfsetmiterjoin%
\definecolor{currentfill}{rgb}{0.000000,0.372549,0.813725}%
\pgfsetfillcolor{currentfill}%
\pgfsetlinewidth{0.000000pt}%
\definecolor{currentstroke}{rgb}{0.000000,0.000000,0.000000}%
\pgfsetstrokecolor{currentstroke}%
\pgfsetstrokeopacity{0.000000}%
\pgfsetdash{}{0pt}%
\pgfpathmoveto{\pgfqpoint{2.088125in}{1.418722in}}%
\pgfpathlineto{\pgfqpoint{2.086211in}{1.418729in}}%
\pgfpathlineto{\pgfqpoint{2.031759in}{1.419257in}}%
\pgfpathlineto{\pgfqpoint{2.018170in}{1.419481in}}%
\pgfpathlineto{\pgfqpoint{2.019155in}{1.481378in}}%
\pgfpathlineto{\pgfqpoint{2.056854in}{1.480886in}}%
\pgfpathlineto{\pgfqpoint{2.053550in}{1.477162in}}%
\pgfpathlineto{\pgfqpoint{2.059082in}{1.471267in}}%
\pgfpathlineto{\pgfqpoint{2.060821in}{1.465149in}}%
\pgfpathlineto{\pgfqpoint{2.069544in}{1.443619in}}%
\pgfpathlineto{\pgfqpoint{2.077207in}{1.438452in}}%
\pgfpathlineto{\pgfqpoint{2.077853in}{1.433443in}}%
\pgfpathlineto{\pgfqpoint{2.082351in}{1.428944in}}%
\pgfpathlineto{\pgfqpoint{2.081282in}{1.423410in}}%
\pgfpathlineto{\pgfqpoint{2.088125in}{1.418722in}}%
\pgfpathclose%
\pgfusepath{fill}%
\end{pgfscope}%
\begin{pgfscope}%
\pgfpathrectangle{\pgfqpoint{0.100000in}{0.100000in}}{\pgfqpoint{3.608454in}{2.310000in}}%
\pgfusepath{clip}%
\pgfsetbuttcap%
\pgfsetmiterjoin%
\definecolor{currentfill}{rgb}{0.000000,0.619608,0.690196}%
\pgfsetfillcolor{currentfill}%
\pgfsetlinewidth{0.000000pt}%
\definecolor{currentstroke}{rgb}{0.000000,0.000000,0.000000}%
\pgfsetstrokecolor{currentstroke}%
\pgfsetstrokeopacity{0.000000}%
\pgfsetdash{}{0pt}%
\pgfpathmoveto{\pgfqpoint{1.624223in}{0.845743in}}%
\pgfpathlineto{\pgfqpoint{1.566607in}{0.850506in}}%
\pgfpathlineto{\pgfqpoint{1.563559in}{0.816012in}}%
\pgfpathlineto{\pgfqpoint{1.562944in}{0.809149in}}%
\pgfpathlineto{\pgfqpoint{1.518670in}{0.813225in}}%
\pgfpathlineto{\pgfqpoint{1.522650in}{0.854609in}}%
\pgfpathlineto{\pgfqpoint{1.516573in}{0.855187in}}%
\pgfpathlineto{\pgfqpoint{1.519940in}{0.890229in}}%
\pgfpathlineto{\pgfqpoint{1.523143in}{0.889921in}}%
\pgfpathlineto{\pgfqpoint{1.526429in}{0.923948in}}%
\pgfpathlineto{\pgfqpoint{1.529406in}{0.923693in}}%
\pgfpathlineto{\pgfqpoint{1.530771in}{0.937426in}}%
\pgfpathlineto{\pgfqpoint{1.544636in}{0.936153in}}%
\pgfpathlineto{\pgfqpoint{1.574789in}{0.933425in}}%
\pgfpathlineto{\pgfqpoint{1.573241in}{0.919040in}}%
\pgfpathlineto{\pgfqpoint{1.638056in}{0.913724in}}%
\pgfpathlineto{\pgfqpoint{1.635380in}{0.879674in}}%
\pgfpathlineto{\pgfqpoint{1.626580in}{0.880308in}}%
\pgfpathlineto{\pgfqpoint{1.624223in}{0.845743in}}%
\pgfpathclose%
\pgfusepath{fill}%
\end{pgfscope}%
\begin{pgfscope}%
\pgfpathrectangle{\pgfqpoint{0.100000in}{0.100000in}}{\pgfqpoint{3.608454in}{2.310000in}}%
\pgfusepath{clip}%
\pgfsetbuttcap%
\pgfsetmiterjoin%
\definecolor{currentfill}{rgb}{0.000000,0.388235,0.805882}%
\pgfsetfillcolor{currentfill}%
\pgfsetlinewidth{0.000000pt}%
\definecolor{currentstroke}{rgb}{0.000000,0.000000,0.000000}%
\pgfsetstrokecolor{currentstroke}%
\pgfsetstrokeopacity{0.000000}%
\pgfsetdash{}{0pt}%
\pgfpathmoveto{\pgfqpoint{2.547881in}{1.571936in}}%
\pgfpathlineto{\pgfqpoint{2.550384in}{1.539587in}}%
\pgfpathlineto{\pgfqpoint{2.521564in}{1.536820in}}%
\pgfpathlineto{\pgfqpoint{2.522216in}{1.529995in}}%
\pgfpathlineto{\pgfqpoint{2.508442in}{1.528613in}}%
\pgfpathlineto{\pgfqpoint{2.508522in}{1.521688in}}%
\pgfpathlineto{\pgfqpoint{2.488674in}{1.519766in}}%
\pgfpathlineto{\pgfqpoint{2.484797in}{1.561007in}}%
\pgfpathlineto{\pgfqpoint{2.482413in}{1.595355in}}%
\pgfpathlineto{\pgfqpoint{2.482705in}{1.602290in}}%
\pgfpathlineto{\pgfqpoint{2.475849in}{1.601806in}}%
\pgfpathlineto{\pgfqpoint{2.473885in}{1.628654in}}%
\pgfpathlineto{\pgfqpoint{2.526558in}{1.632431in}}%
\pgfpathlineto{\pgfqpoint{2.525980in}{1.616699in}}%
\pgfpathlineto{\pgfqpoint{2.528499in}{1.610045in}}%
\pgfpathlineto{\pgfqpoint{2.536221in}{1.600158in}}%
\pgfpathlineto{\pgfqpoint{2.541972in}{1.582322in}}%
\pgfpathlineto{\pgfqpoint{2.547881in}{1.571936in}}%
\pgfpathclose%
\pgfusepath{fill}%
\end{pgfscope}%
\begin{pgfscope}%
\pgfpathrectangle{\pgfqpoint{0.100000in}{0.100000in}}{\pgfqpoint{3.608454in}{2.310000in}}%
\pgfusepath{clip}%
\pgfsetbuttcap%
\pgfsetmiterjoin%
\definecolor{currentfill}{rgb}{0.000000,0.490196,0.754902}%
\pgfsetfillcolor{currentfill}%
\pgfsetlinewidth{0.000000pt}%
\definecolor{currentstroke}{rgb}{0.000000,0.000000,0.000000}%
\pgfsetstrokecolor{currentstroke}%
\pgfsetstrokeopacity{0.000000}%
\pgfsetdash{}{0pt}%
\pgfpathmoveto{\pgfqpoint{1.962491in}{1.386340in}}%
\pgfpathlineto{\pgfqpoint{1.961886in}{1.365738in}}%
\pgfpathlineto{\pgfqpoint{1.927666in}{1.366756in}}%
\pgfpathlineto{\pgfqpoint{1.928379in}{1.387333in}}%
\pgfpathlineto{\pgfqpoint{1.893270in}{1.388628in}}%
\pgfpathlineto{\pgfqpoint{1.894644in}{1.422961in}}%
\pgfpathlineto{\pgfqpoint{1.908588in}{1.422459in}}%
\pgfpathlineto{\pgfqpoint{1.963356in}{1.420723in}}%
\pgfpathlineto{\pgfqpoint{1.962491in}{1.386340in}}%
\pgfpathclose%
\pgfusepath{fill}%
\end{pgfscope}%
\begin{pgfscope}%
\pgfpathrectangle{\pgfqpoint{0.100000in}{0.100000in}}{\pgfqpoint{3.608454in}{2.310000in}}%
\pgfusepath{clip}%
\pgfsetbuttcap%
\pgfsetmiterjoin%
\definecolor{currentfill}{rgb}{0.000000,0.698039,0.650980}%
\pgfsetfillcolor{currentfill}%
\pgfsetlinewidth{0.000000pt}%
\definecolor{currentstroke}{rgb}{0.000000,0.000000,0.000000}%
\pgfsetstrokecolor{currentstroke}%
\pgfsetstrokeopacity{0.000000}%
\pgfsetdash{}{0pt}%
\pgfpathmoveto{\pgfqpoint{1.691656in}{1.593744in}}%
\pgfpathlineto{\pgfqpoint{1.655822in}{1.596346in}}%
\pgfpathlineto{\pgfqpoint{1.652617in}{1.603533in}}%
\pgfpathlineto{\pgfqpoint{1.654543in}{1.630993in}}%
\pgfpathlineto{\pgfqpoint{1.653148in}{1.631106in}}%
\pgfpathlineto{\pgfqpoint{1.655483in}{1.658354in}}%
\pgfpathlineto{\pgfqpoint{1.655854in}{1.675393in}}%
\pgfpathlineto{\pgfqpoint{1.643808in}{1.676503in}}%
\pgfpathlineto{\pgfqpoint{1.648016in}{1.724287in}}%
\pgfpathlineto{\pgfqpoint{1.654175in}{1.729234in}}%
\pgfpathlineto{\pgfqpoint{1.659487in}{1.729827in}}%
\pgfpathlineto{\pgfqpoint{1.695616in}{1.726789in}}%
\pgfpathlineto{\pgfqpoint{1.697918in}{1.727667in}}%
\pgfpathlineto{\pgfqpoint{1.696530in}{1.709950in}}%
\pgfpathlineto{\pgfqpoint{1.697612in}{1.702965in}}%
\pgfpathlineto{\pgfqpoint{1.748309in}{1.699187in}}%
\pgfpathlineto{\pgfqpoint{1.746207in}{1.668241in}}%
\pgfpathlineto{\pgfqpoint{1.696825in}{1.671982in}}%
\pgfpathlineto{\pgfqpoint{1.696414in}{1.655020in}}%
\pgfpathlineto{\pgfqpoint{1.694322in}{1.628015in}}%
\pgfpathlineto{\pgfqpoint{1.695984in}{1.627874in}}%
\pgfpathlineto{\pgfqpoint{1.693738in}{1.600532in}}%
\pgfpathlineto{\pgfqpoint{1.691656in}{1.593744in}}%
\pgfpathclose%
\pgfusepath{fill}%
\end{pgfscope}%
\begin{pgfscope}%
\pgfpathrectangle{\pgfqpoint{0.100000in}{0.100000in}}{\pgfqpoint{3.608454in}{2.310000in}}%
\pgfusepath{clip}%
\pgfsetbuttcap%
\pgfsetmiterjoin%
\definecolor{currentfill}{rgb}{0.000000,0.721569,0.639216}%
\pgfsetfillcolor{currentfill}%
\pgfsetlinewidth{0.000000pt}%
\definecolor{currentstroke}{rgb}{0.000000,0.000000,0.000000}%
\pgfsetstrokecolor{currentstroke}%
\pgfsetstrokeopacity{0.000000}%
\pgfsetdash{}{0pt}%
\pgfpathmoveto{\pgfqpoint{2.136953in}{2.110228in}}%
\pgfpathlineto{\pgfqpoint{2.145780in}{2.110086in}}%
\pgfpathlineto{\pgfqpoint{2.146826in}{2.106493in}}%
\pgfpathlineto{\pgfqpoint{2.168582in}{2.104438in}}%
\pgfpathlineto{\pgfqpoint{2.170811in}{2.095732in}}%
\pgfpathlineto{\pgfqpoint{2.188133in}{2.098370in}}%
\pgfpathlineto{\pgfqpoint{2.188169in}{2.101766in}}%
\pgfpathlineto{\pgfqpoint{2.201785in}{2.106387in}}%
\pgfpathlineto{\pgfqpoint{2.208113in}{2.105322in}}%
\pgfpathlineto{\pgfqpoint{2.208627in}{2.060466in}}%
\pgfpathlineto{\pgfqpoint{2.209745in}{2.046401in}}%
\pgfpathlineto{\pgfqpoint{2.172471in}{2.046322in}}%
\pgfpathlineto{\pgfqpoint{2.172503in}{2.042117in}}%
\pgfpathlineto{\pgfqpoint{2.138015in}{2.041726in}}%
\pgfpathlineto{\pgfqpoint{2.137206in}{2.083483in}}%
\pgfpathlineto{\pgfqpoint{2.136953in}{2.110228in}}%
\pgfpathclose%
\pgfusepath{fill}%
\end{pgfscope}%
\begin{pgfscope}%
\pgfpathrectangle{\pgfqpoint{0.100000in}{0.100000in}}{\pgfqpoint{3.608454in}{2.310000in}}%
\pgfusepath{clip}%
\pgfsetbuttcap%
\pgfsetmiterjoin%
\definecolor{currentfill}{rgb}{0.000000,0.545098,0.727451}%
\pgfsetfillcolor{currentfill}%
\pgfsetlinewidth{0.000000pt}%
\definecolor{currentstroke}{rgb}{0.000000,0.000000,0.000000}%
\pgfsetstrokecolor{currentstroke}%
\pgfsetstrokeopacity{0.000000}%
\pgfsetdash{}{0pt}%
\pgfpathmoveto{\pgfqpoint{2.554067in}{1.495355in}}%
\pgfpathlineto{\pgfqpoint{2.555646in}{1.476018in}}%
\pgfpathlineto{\pgfqpoint{2.531157in}{1.473630in}}%
\pgfpathlineto{\pgfqpoint{2.520205in}{1.472936in}}%
\pgfpathlineto{\pgfqpoint{2.516336in}{1.513050in}}%
\pgfpathlineto{\pgfqpoint{2.509486in}{1.512307in}}%
\pgfpathlineto{\pgfqpoint{2.508522in}{1.521688in}}%
\pgfpathlineto{\pgfqpoint{2.508442in}{1.528613in}}%
\pgfpathlineto{\pgfqpoint{2.522216in}{1.529995in}}%
\pgfpathlineto{\pgfqpoint{2.521564in}{1.536820in}}%
\pgfpathlineto{\pgfqpoint{2.550384in}{1.539587in}}%
\pgfpathlineto{\pgfqpoint{2.554067in}{1.495355in}}%
\pgfpathclose%
\pgfusepath{fill}%
\end{pgfscope}%
\begin{pgfscope}%
\pgfpathrectangle{\pgfqpoint{0.100000in}{0.100000in}}{\pgfqpoint{3.608454in}{2.310000in}}%
\pgfusepath{clip}%
\pgfsetbuttcap%
\pgfsetmiterjoin%
\definecolor{currentfill}{rgb}{0.000000,0.537255,0.731373}%
\pgfsetfillcolor{currentfill}%
\pgfsetlinewidth{0.000000pt}%
\definecolor{currentstroke}{rgb}{0.000000,0.000000,0.000000}%
\pgfsetstrokecolor{currentstroke}%
\pgfsetstrokeopacity{0.000000}%
\pgfsetdash{}{0pt}%
\pgfpathmoveto{\pgfqpoint{2.452854in}{0.589739in}}%
\pgfpathlineto{\pgfqpoint{2.445754in}{0.592557in}}%
\pgfpathlineto{\pgfqpoint{2.442006in}{0.599263in}}%
\pgfpathlineto{\pgfqpoint{2.443757in}{0.607123in}}%
\pgfpathlineto{\pgfqpoint{2.440322in}{0.609463in}}%
\pgfpathlineto{\pgfqpoint{2.430825in}{0.609081in}}%
\pgfpathlineto{\pgfqpoint{2.428180in}{0.614445in}}%
\pgfpathlineto{\pgfqpoint{2.420890in}{0.618181in}}%
\pgfpathlineto{\pgfqpoint{2.416070in}{0.624294in}}%
\pgfpathlineto{\pgfqpoint{2.392099in}{0.624223in}}%
\pgfpathlineto{\pgfqpoint{2.386695in}{0.629800in}}%
\pgfpathlineto{\pgfqpoint{2.386115in}{0.636858in}}%
\pgfpathlineto{\pgfqpoint{2.389280in}{0.640751in}}%
\pgfpathlineto{\pgfqpoint{2.407987in}{0.645878in}}%
\pgfpathlineto{\pgfqpoint{2.415540in}{0.649843in}}%
\pgfpathlineto{\pgfqpoint{2.419484in}{0.654679in}}%
\pgfpathlineto{\pgfqpoint{2.431222in}{0.657222in}}%
\pgfpathlineto{\pgfqpoint{2.435277in}{0.651886in}}%
\pgfpathlineto{\pgfqpoint{2.432451in}{0.691509in}}%
\pgfpathlineto{\pgfqpoint{2.424922in}{0.713419in}}%
\pgfpathlineto{\pgfqpoint{2.467199in}{0.715885in}}%
\pgfpathlineto{\pgfqpoint{2.465118in}{0.704833in}}%
\pgfpathlineto{\pgfqpoint{2.461305in}{0.698015in}}%
\pgfpathlineto{\pgfqpoint{2.460731in}{0.688378in}}%
\pgfpathlineto{\pgfqpoint{2.466743in}{0.676753in}}%
\pgfpathlineto{\pgfqpoint{2.472822in}{0.672176in}}%
\pgfpathlineto{\pgfqpoint{2.478549in}{0.654208in}}%
\pgfpathlineto{\pgfqpoint{2.485040in}{0.651337in}}%
\pgfpathlineto{\pgfqpoint{2.478638in}{0.648879in}}%
\pgfpathlineto{\pgfqpoint{2.472203in}{0.638511in}}%
\pgfpathlineto{\pgfqpoint{2.465852in}{0.639118in}}%
\pgfpathlineto{\pgfqpoint{2.464257in}{0.631545in}}%
\pgfpathlineto{\pgfqpoint{2.470761in}{0.631547in}}%
\pgfpathlineto{\pgfqpoint{2.473516in}{0.626309in}}%
\pgfpathlineto{\pgfqpoint{2.481391in}{0.627004in}}%
\pgfpathlineto{\pgfqpoint{2.482013in}{0.636124in}}%
\pgfpathlineto{\pgfqpoint{2.487826in}{0.640181in}}%
\pgfpathlineto{\pgfqpoint{2.500761in}{0.631538in}}%
\pgfpathlineto{\pgfqpoint{2.493847in}{0.616542in}}%
\pgfpathlineto{\pgfqpoint{2.478334in}{0.616245in}}%
\pgfpathlineto{\pgfqpoint{2.480862in}{0.608985in}}%
\pgfpathlineto{\pgfqpoint{2.477552in}{0.599501in}}%
\pgfpathlineto{\pgfqpoint{2.485322in}{0.596196in}}%
\pgfpathlineto{\pgfqpoint{2.492517in}{0.589207in}}%
\pgfpathlineto{\pgfqpoint{2.511782in}{0.586087in}}%
\pgfpathlineto{\pgfqpoint{2.520554in}{0.574448in}}%
\pgfpathlineto{\pgfqpoint{2.526501in}{0.573415in}}%
\pgfpathlineto{\pgfqpoint{2.522992in}{0.564279in}}%
\pgfpathlineto{\pgfqpoint{2.516518in}{0.560200in}}%
\pgfpathlineto{\pgfqpoint{2.511813in}{0.562316in}}%
\pgfpathlineto{\pgfqpoint{2.503209in}{0.561121in}}%
\pgfpathlineto{\pgfqpoint{2.491659in}{0.574580in}}%
\pgfpathlineto{\pgfqpoint{2.475551in}{0.580137in}}%
\pgfpathlineto{\pgfqpoint{2.466715in}{0.580608in}}%
\pgfpathlineto{\pgfqpoint{2.468264in}{0.585589in}}%
\pgfpathlineto{\pgfqpoint{2.466744in}{0.592994in}}%
\pgfpathlineto{\pgfqpoint{2.452854in}{0.589739in}}%
\pgfpathclose%
\pgfusepath{fill}%
\end{pgfscope}%
\begin{pgfscope}%
\pgfpathrectangle{\pgfqpoint{0.100000in}{0.100000in}}{\pgfqpoint{3.608454in}{2.310000in}}%
\pgfusepath{clip}%
\pgfsetbuttcap%
\pgfsetmiterjoin%
\definecolor{currentfill}{rgb}{0.000000,0.607843,0.696078}%
\pgfsetfillcolor{currentfill}%
\pgfsetlinewidth{0.000000pt}%
\definecolor{currentstroke}{rgb}{0.000000,0.000000,0.000000}%
\pgfsetstrokecolor{currentstroke}%
\pgfsetstrokeopacity{0.000000}%
\pgfsetdash{}{0pt}%
\pgfpathmoveto{\pgfqpoint{2.306014in}{1.604048in}}%
\pgfpathlineto{\pgfqpoint{2.264937in}{1.602761in}}%
\pgfpathlineto{\pgfqpoint{2.237469in}{1.601952in}}%
\pgfpathlineto{\pgfqpoint{2.237674in}{1.595106in}}%
\pgfpathlineto{\pgfqpoint{2.223956in}{1.594717in}}%
\pgfpathlineto{\pgfqpoint{2.223427in}{1.615413in}}%
\pgfpathlineto{\pgfqpoint{2.222007in}{1.615392in}}%
\pgfpathlineto{\pgfqpoint{2.221382in}{1.649990in}}%
\pgfpathlineto{\pgfqpoint{2.276089in}{1.651404in}}%
\pgfpathlineto{\pgfqpoint{2.275677in}{1.665325in}}%
\pgfpathlineto{\pgfqpoint{2.303199in}{1.666191in}}%
\pgfpathlineto{\pgfqpoint{2.304826in}{1.617701in}}%
\pgfpathlineto{\pgfqpoint{2.306014in}{1.604048in}}%
\pgfpathclose%
\pgfusepath{fill}%
\end{pgfscope}%
\begin{pgfscope}%
\pgfpathrectangle{\pgfqpoint{0.100000in}{0.100000in}}{\pgfqpoint{3.608454in}{2.310000in}}%
\pgfusepath{clip}%
\pgfsetbuttcap%
\pgfsetmiterjoin%
\definecolor{currentfill}{rgb}{0.000000,0.333333,0.833333}%
\pgfsetfillcolor{currentfill}%
\pgfsetlinewidth{0.000000pt}%
\definecolor{currentstroke}{rgb}{0.000000,0.000000,0.000000}%
\pgfsetstrokecolor{currentstroke}%
\pgfsetstrokeopacity{0.000000}%
\pgfsetdash{}{0pt}%
\pgfpathmoveto{\pgfqpoint{2.866524in}{0.926111in}}%
\pgfpathlineto{\pgfqpoint{2.849065in}{0.920723in}}%
\pgfpathlineto{\pgfqpoint{2.848183in}{0.924494in}}%
\pgfpathlineto{\pgfqpoint{2.848096in}{0.929927in}}%
\pgfpathlineto{\pgfqpoint{2.843454in}{0.939687in}}%
\pgfpathlineto{\pgfqpoint{2.838874in}{0.945035in}}%
\pgfpathlineto{\pgfqpoint{2.830443in}{0.932055in}}%
\pgfpathlineto{\pgfqpoint{2.818507in}{0.934685in}}%
\pgfpathlineto{\pgfqpoint{2.809234in}{0.933588in}}%
\pgfpathlineto{\pgfqpoint{2.807171in}{0.925663in}}%
\pgfpathlineto{\pgfqpoint{2.802892in}{0.925244in}}%
\pgfpathlineto{\pgfqpoint{2.796122in}{0.932528in}}%
\pgfpathlineto{\pgfqpoint{2.791999in}{0.944257in}}%
\pgfpathlineto{\pgfqpoint{2.778830in}{0.942946in}}%
\pgfpathlineto{\pgfqpoint{2.773147in}{0.947109in}}%
\pgfpathlineto{\pgfqpoint{2.771603in}{0.963496in}}%
\pgfpathlineto{\pgfqpoint{2.762398in}{0.964970in}}%
\pgfpathlineto{\pgfqpoint{2.760784in}{0.972182in}}%
\pgfpathlineto{\pgfqpoint{2.765064in}{0.976399in}}%
\pgfpathlineto{\pgfqpoint{2.767572in}{0.987144in}}%
\pgfpathlineto{\pgfqpoint{2.784817in}{0.988730in}}%
\pgfpathlineto{\pgfqpoint{2.780892in}{1.025876in}}%
\pgfpathlineto{\pgfqpoint{2.799254in}{1.027968in}}%
\pgfpathlineto{\pgfqpoint{2.801686in}{1.031063in}}%
\pgfpathlineto{\pgfqpoint{2.810594in}{1.033632in}}%
\pgfpathlineto{\pgfqpoint{2.810860in}{1.028741in}}%
\pgfpathlineto{\pgfqpoint{2.818142in}{1.023149in}}%
\pgfpathlineto{\pgfqpoint{2.825899in}{1.020792in}}%
\pgfpathlineto{\pgfqpoint{2.830646in}{1.010243in}}%
\pgfpathlineto{\pgfqpoint{2.827384in}{1.003120in}}%
\pgfpathlineto{\pgfqpoint{2.836248in}{0.996454in}}%
\pgfpathlineto{\pgfqpoint{2.839244in}{0.999162in}}%
\pgfpathlineto{\pgfqpoint{2.837086in}{0.989053in}}%
\pgfpathlineto{\pgfqpoint{2.844938in}{0.981568in}}%
\pgfpathlineto{\pgfqpoint{2.850509in}{0.983786in}}%
\pgfpathlineto{\pgfqpoint{2.864261in}{0.976571in}}%
\pgfpathlineto{\pgfqpoint{2.870396in}{0.968762in}}%
\pgfpathlineto{\pgfqpoint{2.880855in}{0.952903in}}%
\pgfpathlineto{\pgfqpoint{2.864668in}{0.946906in}}%
\pgfpathlineto{\pgfqpoint{2.866524in}{0.926111in}}%
\pgfpathclose%
\pgfusepath{fill}%
\end{pgfscope}%
\begin{pgfscope}%
\pgfpathrectangle{\pgfqpoint{0.100000in}{0.100000in}}{\pgfqpoint{3.608454in}{2.310000in}}%
\pgfusepath{clip}%
\pgfsetbuttcap%
\pgfsetmiterjoin%
\definecolor{currentfill}{rgb}{0.000000,0.803922,0.598039}%
\pgfsetfillcolor{currentfill}%
\pgfsetlinewidth{0.000000pt}%
\definecolor{currentstroke}{rgb}{0.000000,0.000000,0.000000}%
\pgfsetstrokecolor{currentstroke}%
\pgfsetstrokeopacity{0.000000}%
\pgfsetdash{}{0pt}%
\pgfpathmoveto{\pgfqpoint{2.882548in}{1.076650in}}%
\pgfpathlineto{\pgfqpoint{2.876403in}{1.073700in}}%
\pgfpathlineto{\pgfqpoint{2.872545in}{1.080983in}}%
\pgfpathlineto{\pgfqpoint{2.867388in}{1.085037in}}%
\pgfpathlineto{\pgfqpoint{2.858178in}{1.097707in}}%
\pgfpathlineto{\pgfqpoint{2.851025in}{1.094550in}}%
\pgfpathlineto{\pgfqpoint{2.837326in}{1.091183in}}%
\pgfpathlineto{\pgfqpoint{2.835727in}{1.088275in}}%
\pgfpathlineto{\pgfqpoint{2.827437in}{1.087830in}}%
\pgfpathlineto{\pgfqpoint{2.819581in}{1.083755in}}%
\pgfpathlineto{\pgfqpoint{2.814347in}{1.089123in}}%
\pgfpathlineto{\pgfqpoint{2.813980in}{1.097663in}}%
\pgfpathlineto{\pgfqpoint{2.817068in}{1.103142in}}%
\pgfpathlineto{\pgfqpoint{2.828227in}{1.112365in}}%
\pgfpathlineto{\pgfqpoint{2.835111in}{1.113784in}}%
\pgfpathlineto{\pgfqpoint{2.845692in}{1.114629in}}%
\pgfpathlineto{\pgfqpoint{2.859932in}{1.127064in}}%
\pgfpathlineto{\pgfqpoint{2.864797in}{1.125728in}}%
\pgfpathlineto{\pgfqpoint{2.866841in}{1.114054in}}%
\pgfpathlineto{\pgfqpoint{2.884882in}{1.098977in}}%
\pgfpathlineto{\pgfqpoint{2.885380in}{1.093119in}}%
\pgfpathlineto{\pgfqpoint{2.879644in}{1.078954in}}%
\pgfpathlineto{\pgfqpoint{2.882548in}{1.076650in}}%
\pgfpathclose%
\pgfusepath{fill}%
\end{pgfscope}%
\begin{pgfscope}%
\pgfpathrectangle{\pgfqpoint{0.100000in}{0.100000in}}{\pgfqpoint{3.608454in}{2.310000in}}%
\pgfusepath{clip}%
\pgfsetbuttcap%
\pgfsetmiterjoin%
\definecolor{currentfill}{rgb}{0.000000,0.419608,0.790196}%
\pgfsetfillcolor{currentfill}%
\pgfsetlinewidth{0.000000pt}%
\definecolor{currentstroke}{rgb}{0.000000,0.000000,0.000000}%
\pgfsetstrokecolor{currentstroke}%
\pgfsetstrokeopacity{0.000000}%
\pgfsetdash{}{0pt}%
\pgfpathmoveto{\pgfqpoint{2.433615in}{1.653903in}}%
\pgfpathlineto{\pgfqpoint{2.435480in}{1.626662in}}%
\pgfpathlineto{\pgfqpoint{2.402794in}{1.625171in}}%
\pgfpathlineto{\pgfqpoint{2.361136in}{1.623207in}}%
\pgfpathlineto{\pgfqpoint{2.356020in}{1.633465in}}%
\pgfpathlineto{\pgfqpoint{2.342986in}{1.636283in}}%
\pgfpathlineto{\pgfqpoint{2.336274in}{1.640299in}}%
\pgfpathlineto{\pgfqpoint{2.333060in}{1.651644in}}%
\pgfpathlineto{\pgfqpoint{2.330505in}{1.653247in}}%
\pgfpathlineto{\pgfqpoint{2.327890in}{1.671376in}}%
\pgfpathlineto{\pgfqpoint{2.334254in}{1.681222in}}%
\pgfpathlineto{\pgfqpoint{2.325644in}{1.688292in}}%
\pgfpathlineto{\pgfqpoint{2.325187in}{1.694132in}}%
\pgfpathlineto{\pgfqpoint{2.356072in}{1.695517in}}%
\pgfpathlineto{\pgfqpoint{2.357123in}{1.675633in}}%
\pgfpathlineto{\pgfqpoint{2.363976in}{1.679112in}}%
\pgfpathlineto{\pgfqpoint{2.370685in}{1.678611in}}%
\pgfpathlineto{\pgfqpoint{2.372395in}{1.647911in}}%
\pgfpathlineto{\pgfqpoint{2.406543in}{1.649815in}}%
\pgfpathlineto{\pgfqpoint{2.406360in}{1.653264in}}%
\pgfpathlineto{\pgfqpoint{2.433615in}{1.653903in}}%
\pgfpathclose%
\pgfusepath{fill}%
\end{pgfscope}%
\begin{pgfscope}%
\pgfpathrectangle{\pgfqpoint{0.100000in}{0.100000in}}{\pgfqpoint{3.608454in}{2.310000in}}%
\pgfusepath{clip}%
\pgfsetbuttcap%
\pgfsetmiterjoin%
\definecolor{currentfill}{rgb}{0.000000,0.807843,0.596078}%
\pgfsetfillcolor{currentfill}%
\pgfsetlinewidth{0.000000pt}%
\definecolor{currentstroke}{rgb}{0.000000,0.000000,0.000000}%
\pgfsetstrokecolor{currentstroke}%
\pgfsetstrokeopacity{0.000000}%
\pgfsetdash{}{0pt}%
\pgfpathmoveto{\pgfqpoint{2.288166in}{1.125922in}}%
\pgfpathlineto{\pgfqpoint{2.288004in}{1.144708in}}%
\pgfpathlineto{\pgfqpoint{2.289909in}{1.144788in}}%
\pgfpathlineto{\pgfqpoint{2.289673in}{1.168159in}}%
\pgfpathlineto{\pgfqpoint{2.290504in}{1.188921in}}%
\pgfpathlineto{\pgfqpoint{2.280237in}{1.188871in}}%
\pgfpathlineto{\pgfqpoint{2.279053in}{1.231943in}}%
\pgfpathlineto{\pgfqpoint{2.310037in}{1.232515in}}%
\pgfpathlineto{\pgfqpoint{2.310515in}{1.218707in}}%
\pgfpathlineto{\pgfqpoint{2.346916in}{1.219387in}}%
\pgfpathlineto{\pgfqpoint{2.347208in}{1.211310in}}%
\pgfpathlineto{\pgfqpoint{2.350740in}{1.202506in}}%
\pgfpathlineto{\pgfqpoint{2.357801in}{1.198383in}}%
\pgfpathlineto{\pgfqpoint{2.357912in}{1.194600in}}%
\pgfpathlineto{\pgfqpoint{2.345378in}{1.193313in}}%
\pgfpathlineto{\pgfqpoint{2.345750in}{1.172705in}}%
\pgfpathlineto{\pgfqpoint{2.352645in}{1.172837in}}%
\pgfpathlineto{\pgfqpoint{2.353073in}{1.147045in}}%
\pgfpathlineto{\pgfqpoint{2.335244in}{1.146265in}}%
\pgfpathlineto{\pgfqpoint{2.332518in}{1.146200in}}%
\pgfpathlineto{\pgfqpoint{2.332736in}{1.133366in}}%
\pgfpathlineto{\pgfqpoint{2.325633in}{1.133246in}}%
\pgfpathlineto{\pgfqpoint{2.325882in}{1.126330in}}%
\pgfpathlineto{\pgfqpoint{2.288166in}{1.125922in}}%
\pgfpathclose%
\pgfusepath{fill}%
\end{pgfscope}%
\begin{pgfscope}%
\pgfpathrectangle{\pgfqpoint{0.100000in}{0.100000in}}{\pgfqpoint{3.608454in}{2.310000in}}%
\pgfusepath{clip}%
\pgfsetbuttcap%
\pgfsetmiterjoin%
\definecolor{currentfill}{rgb}{0.000000,0.509804,0.745098}%
\pgfsetfillcolor{currentfill}%
\pgfsetlinewidth{0.000000pt}%
\definecolor{currentstroke}{rgb}{0.000000,0.000000,0.000000}%
\pgfsetstrokecolor{currentstroke}%
\pgfsetstrokeopacity{0.000000}%
\pgfsetdash{}{0pt}%
\pgfpathmoveto{\pgfqpoint{1.202252in}{1.853038in}}%
\pgfpathlineto{\pgfqpoint{1.204580in}{1.868127in}}%
\pgfpathlineto{\pgfqpoint{1.226134in}{1.864328in}}%
\pgfpathlineto{\pgfqpoint{1.224995in}{1.857794in}}%
\pgfpathlineto{\pgfqpoint{1.241305in}{1.854999in}}%
\pgfpathlineto{\pgfqpoint{1.245880in}{1.848777in}}%
\pgfpathlineto{\pgfqpoint{1.243670in}{1.842268in}}%
\pgfpathlineto{\pgfqpoint{1.249318in}{1.830814in}}%
\pgfpathlineto{\pgfqpoint{1.250644in}{1.821729in}}%
\pgfpathlineto{\pgfqpoint{1.253366in}{1.816700in}}%
\pgfpathlineto{\pgfqpoint{1.251746in}{1.806946in}}%
\pgfpathlineto{\pgfqpoint{1.244733in}{1.764439in}}%
\pgfpathlineto{\pgfqpoint{1.243269in}{1.757690in}}%
\pgfpathlineto{\pgfqpoint{1.226972in}{1.760454in}}%
\pgfpathlineto{\pgfqpoint{1.225832in}{1.753673in}}%
\pgfpathlineto{\pgfqpoint{1.212511in}{1.756000in}}%
\pgfpathlineto{\pgfqpoint{1.211732in}{1.751597in}}%
\pgfpathlineto{\pgfqpoint{1.198346in}{1.753929in}}%
\pgfpathlineto{\pgfqpoint{1.199226in}{1.758954in}}%
\pgfpathlineto{\pgfqpoint{1.186356in}{1.762535in}}%
\pgfpathlineto{\pgfqpoint{1.188888in}{1.777013in}}%
\pgfpathlineto{\pgfqpoint{1.180737in}{1.781578in}}%
\pgfpathlineto{\pgfqpoint{1.179359in}{1.788436in}}%
\pgfpathlineto{\pgfqpoint{1.170551in}{1.790061in}}%
\pgfpathlineto{\pgfqpoint{1.174897in}{1.813546in}}%
\pgfpathlineto{\pgfqpoint{1.186131in}{1.812373in}}%
\pgfpathlineto{\pgfqpoint{1.195518in}{1.814680in}}%
\pgfpathlineto{\pgfqpoint{1.202252in}{1.853038in}}%
\pgfpathclose%
\pgfusepath{fill}%
\end{pgfscope}%
\begin{pgfscope}%
\pgfpathrectangle{\pgfqpoint{0.100000in}{0.100000in}}{\pgfqpoint{3.608454in}{2.310000in}}%
\pgfusepath{clip}%
\pgfsetbuttcap%
\pgfsetmiterjoin%
\definecolor{currentfill}{rgb}{0.000000,0.572549,0.713725}%
\pgfsetfillcolor{currentfill}%
\pgfsetlinewidth{0.000000pt}%
\definecolor{currentstroke}{rgb}{0.000000,0.000000,0.000000}%
\pgfsetstrokecolor{currentstroke}%
\pgfsetstrokeopacity{0.000000}%
\pgfsetdash{}{0pt}%
\pgfpathmoveto{\pgfqpoint{1.873743in}{0.830250in}}%
\pgfpathlineto{\pgfqpoint{1.873122in}{0.813238in}}%
\pgfpathlineto{\pgfqpoint{1.842140in}{0.796889in}}%
\pgfpathlineto{\pgfqpoint{1.788954in}{0.799743in}}%
\pgfpathlineto{\pgfqpoint{1.759459in}{0.801389in}}%
\pgfpathlineto{\pgfqpoint{1.763926in}{0.871040in}}%
\pgfpathlineto{\pgfqpoint{1.808729in}{0.868432in}}%
\pgfpathlineto{\pgfqpoint{1.868391in}{0.865576in}}%
\pgfpathlineto{\pgfqpoint{1.866975in}{0.830682in}}%
\pgfpathlineto{\pgfqpoint{1.873743in}{0.830250in}}%
\pgfpathclose%
\pgfusepath{fill}%
\end{pgfscope}%
\begin{pgfscope}%
\pgfpathrectangle{\pgfqpoint{0.100000in}{0.100000in}}{\pgfqpoint{3.608454in}{2.310000in}}%
\pgfusepath{clip}%
\pgfsetbuttcap%
\pgfsetmiterjoin%
\definecolor{currentfill}{rgb}{0.000000,0.666667,0.666667}%
\pgfsetfillcolor{currentfill}%
\pgfsetlinewidth{0.000000pt}%
\definecolor{currentstroke}{rgb}{0.000000,0.000000,0.000000}%
\pgfsetstrokecolor{currentstroke}%
\pgfsetstrokeopacity{0.000000}%
\pgfsetdash{}{0pt}%
\pgfpathmoveto{\pgfqpoint{1.103689in}{1.391163in}}%
\pgfpathlineto{\pgfqpoint{1.103543in}{1.390370in}}%
\pgfpathlineto{\pgfqpoint{1.141855in}{1.383558in}}%
\pgfpathlineto{\pgfqpoint{1.181575in}{1.376785in}}%
\pgfpathlineto{\pgfqpoint{1.183516in}{1.359664in}}%
\pgfpathlineto{\pgfqpoint{1.186021in}{1.352257in}}%
\pgfpathlineto{\pgfqpoint{1.183026in}{1.348649in}}%
\pgfpathlineto{\pgfqpoint{1.136475in}{1.356525in}}%
\pgfpathlineto{\pgfqpoint{1.065800in}{1.369302in}}%
\pgfpathlineto{\pgfqpoint{1.069417in}{1.390629in}}%
\pgfpathlineto{\pgfqpoint{1.076288in}{1.396187in}}%
\pgfpathlineto{\pgfqpoint{1.103689in}{1.391163in}}%
\pgfpathclose%
\pgfusepath{fill}%
\end{pgfscope}%
\begin{pgfscope}%
\pgfpathrectangle{\pgfqpoint{0.100000in}{0.100000in}}{\pgfqpoint{3.608454in}{2.310000in}}%
\pgfusepath{clip}%
\pgfsetbuttcap%
\pgfsetmiterjoin%
\definecolor{currentfill}{rgb}{0.000000,0.643137,0.678431}%
\pgfsetfillcolor{currentfill}%
\pgfsetlinewidth{0.000000pt}%
\definecolor{currentstroke}{rgb}{0.000000,0.000000,0.000000}%
\pgfsetstrokecolor{currentstroke}%
\pgfsetstrokeopacity{0.000000}%
\pgfsetdash{}{0pt}%
\pgfpathmoveto{\pgfqpoint{2.123559in}{0.920377in}}%
\pgfpathlineto{\pgfqpoint{2.140773in}{0.915277in}}%
\pgfpathlineto{\pgfqpoint{2.141202in}{0.939334in}}%
\pgfpathlineto{\pgfqpoint{2.147313in}{0.939341in}}%
\pgfpathlineto{\pgfqpoint{2.157026in}{0.928424in}}%
\pgfpathlineto{\pgfqpoint{2.170065in}{0.926960in}}%
\pgfpathlineto{\pgfqpoint{2.175503in}{0.924555in}}%
\pgfpathlineto{\pgfqpoint{2.175110in}{0.918121in}}%
\pgfpathlineto{\pgfqpoint{2.183884in}{0.913246in}}%
\pgfpathlineto{\pgfqpoint{2.191374in}{0.903396in}}%
\pgfpathlineto{\pgfqpoint{2.190887in}{0.898874in}}%
\pgfpathlineto{\pgfqpoint{2.195345in}{0.891554in}}%
\pgfpathlineto{\pgfqpoint{2.189397in}{0.878197in}}%
\pgfpathlineto{\pgfqpoint{2.183288in}{0.876547in}}%
\pgfpathlineto{\pgfqpoint{2.182676in}{0.872245in}}%
\pgfpathlineto{\pgfqpoint{2.186423in}{0.866575in}}%
\pgfpathlineto{\pgfqpoint{2.170562in}{0.866376in}}%
\pgfpathlineto{\pgfqpoint{2.170684in}{0.855388in}}%
\pgfpathlineto{\pgfqpoint{2.126427in}{0.854941in}}%
\pgfpathlineto{\pgfqpoint{2.125475in}{0.856951in}}%
\pgfpathlineto{\pgfqpoint{2.096641in}{0.856769in}}%
\pgfpathlineto{\pgfqpoint{2.095246in}{0.861494in}}%
\pgfpathlineto{\pgfqpoint{2.086189in}{0.861574in}}%
\pgfpathlineto{\pgfqpoint{2.086328in}{0.894484in}}%
\pgfpathlineto{\pgfqpoint{2.100603in}{0.896104in}}%
\pgfpathlineto{\pgfqpoint{2.115003in}{0.895858in}}%
\pgfpathlineto{\pgfqpoint{2.123587in}{0.890672in}}%
\pgfpathlineto{\pgfqpoint{2.123559in}{0.920377in}}%
\pgfpathclose%
\pgfusepath{fill}%
\end{pgfscope}%
\begin{pgfscope}%
\pgfpathrectangle{\pgfqpoint{0.100000in}{0.100000in}}{\pgfqpoint{3.608454in}{2.310000in}}%
\pgfusepath{clip}%
\pgfsetbuttcap%
\pgfsetmiterjoin%
\definecolor{currentfill}{rgb}{0.000000,0.890196,0.554902}%
\pgfsetfillcolor{currentfill}%
\pgfsetlinewidth{0.000000pt}%
\definecolor{currentstroke}{rgb}{0.000000,0.000000,0.000000}%
\pgfsetstrokecolor{currentstroke}%
\pgfsetstrokeopacity{0.000000}%
\pgfsetdash{}{0pt}%
\pgfpathmoveto{\pgfqpoint{3.001283in}{0.581332in}}%
\pgfpathlineto{\pgfqpoint{2.990145in}{0.581978in}}%
\pgfpathlineto{\pgfqpoint{2.974786in}{0.579836in}}%
\pgfpathlineto{\pgfqpoint{2.972832in}{0.594908in}}%
\pgfpathlineto{\pgfqpoint{2.963258in}{0.602568in}}%
\pgfpathlineto{\pgfqpoint{2.965965in}{0.605834in}}%
\pgfpathlineto{\pgfqpoint{2.973159in}{0.605106in}}%
\pgfpathlineto{\pgfqpoint{2.978125in}{0.608670in}}%
\pgfpathlineto{\pgfqpoint{2.976237in}{0.622163in}}%
\pgfpathlineto{\pgfqpoint{2.985432in}{0.623511in}}%
\pgfpathlineto{\pgfqpoint{2.982209in}{0.644864in}}%
\pgfpathlineto{\pgfqpoint{2.995671in}{0.646679in}}%
\pgfpathlineto{\pgfqpoint{2.996330in}{0.642369in}}%
\pgfpathlineto{\pgfqpoint{3.006579in}{0.647225in}}%
\pgfpathlineto{\pgfqpoint{3.020704in}{0.653264in}}%
\pgfpathlineto{\pgfqpoint{3.025702in}{0.651277in}}%
\pgfpathlineto{\pgfqpoint{3.029457in}{0.642400in}}%
\pgfpathlineto{\pgfqpoint{3.037437in}{0.635958in}}%
\pgfpathlineto{\pgfqpoint{3.040352in}{0.617807in}}%
\pgfpathlineto{\pgfqpoint{3.039989in}{0.610672in}}%
\pgfpathlineto{\pgfqpoint{2.994626in}{0.604123in}}%
\pgfpathlineto{\pgfqpoint{3.006411in}{0.592107in}}%
\pgfpathlineto{\pgfqpoint{3.001283in}{0.581332in}}%
\pgfpathclose%
\pgfusepath{fill}%
\end{pgfscope}%
\begin{pgfscope}%
\pgfpathrectangle{\pgfqpoint{0.100000in}{0.100000in}}{\pgfqpoint{3.608454in}{2.310000in}}%
\pgfusepath{clip}%
\pgfsetbuttcap%
\pgfsetmiterjoin%
\definecolor{currentfill}{rgb}{0.000000,0.615686,0.692157}%
\pgfsetfillcolor{currentfill}%
\pgfsetlinewidth{0.000000pt}%
\definecolor{currentstroke}{rgb}{0.000000,0.000000,0.000000}%
\pgfsetstrokecolor{currentstroke}%
\pgfsetstrokeopacity{0.000000}%
\pgfsetdash{}{0pt}%
\pgfpathmoveto{\pgfqpoint{2.002513in}{1.074759in}}%
\pgfpathlineto{\pgfqpoint{2.030056in}{1.074337in}}%
\pgfpathlineto{\pgfqpoint{2.030288in}{1.091590in}}%
\pgfpathlineto{\pgfqpoint{2.057637in}{1.091246in}}%
\pgfpathlineto{\pgfqpoint{2.057547in}{1.080899in}}%
\pgfpathlineto{\pgfqpoint{2.060962in}{1.080885in}}%
\pgfpathlineto{\pgfqpoint{2.060863in}{1.067101in}}%
\pgfpathlineto{\pgfqpoint{2.053817in}{1.067155in}}%
\pgfpathlineto{\pgfqpoint{2.053684in}{1.060206in}}%
\pgfpathlineto{\pgfqpoint{2.050250in}{1.053326in}}%
\pgfpathlineto{\pgfqpoint{2.043425in}{1.053406in}}%
\pgfpathlineto{\pgfqpoint{2.043347in}{1.046522in}}%
\pgfpathlineto{\pgfqpoint{2.013608in}{1.046935in}}%
\pgfpathlineto{\pgfqpoint{2.010159in}{1.057965in}}%
\pgfpathlineto{\pgfqpoint{2.001980in}{1.055914in}}%
\pgfpathlineto{\pgfqpoint{2.002513in}{1.074759in}}%
\pgfpathclose%
\pgfusepath{fill}%
\end{pgfscope}%
\begin{pgfscope}%
\pgfpathrectangle{\pgfqpoint{0.100000in}{0.100000in}}{\pgfqpoint{3.608454in}{2.310000in}}%
\pgfusepath{clip}%
\pgfsetbuttcap%
\pgfsetmiterjoin%
\definecolor{currentfill}{rgb}{0.000000,0.721569,0.639216}%
\pgfsetfillcolor{currentfill}%
\pgfsetlinewidth{0.000000pt}%
\definecolor{currentstroke}{rgb}{0.000000,0.000000,0.000000}%
\pgfsetstrokecolor{currentstroke}%
\pgfsetstrokeopacity{0.000000}%
\pgfsetdash{}{0pt}%
\pgfpathmoveto{\pgfqpoint{2.949055in}{1.160645in}}%
\pgfpathlineto{\pgfqpoint{2.942760in}{1.162773in}}%
\pgfpathlineto{\pgfqpoint{2.938556in}{1.155078in}}%
\pgfpathlineto{\pgfqpoint{2.934882in}{1.156995in}}%
\pgfpathlineto{\pgfqpoint{2.930144in}{1.169183in}}%
\pgfpathlineto{\pgfqpoint{2.933147in}{1.171253in}}%
\pgfpathlineto{\pgfqpoint{2.936507in}{1.182790in}}%
\pgfpathlineto{\pgfqpoint{2.929794in}{1.191773in}}%
\pgfpathlineto{\pgfqpoint{2.933739in}{1.199536in}}%
\pgfpathlineto{\pgfqpoint{2.934516in}{1.206979in}}%
\pgfpathlineto{\pgfqpoint{2.940357in}{1.211530in}}%
\pgfpathlineto{\pgfqpoint{2.951682in}{1.213011in}}%
\pgfpathlineto{\pgfqpoint{2.950033in}{1.210864in}}%
\pgfpathlineto{\pgfqpoint{2.950303in}{1.201290in}}%
\pgfpathlineto{\pgfqpoint{2.948361in}{1.195220in}}%
\pgfpathlineto{\pgfqpoint{2.955269in}{1.192553in}}%
\pgfpathlineto{\pgfqpoint{2.968126in}{1.182900in}}%
\pgfpathlineto{\pgfqpoint{2.963826in}{1.175254in}}%
\pgfpathlineto{\pgfqpoint{2.955589in}{1.174438in}}%
\pgfpathlineto{\pgfqpoint{2.949055in}{1.160645in}}%
\pgfpathclose%
\pgfusepath{fill}%
\end{pgfscope}%
\begin{pgfscope}%
\pgfpathrectangle{\pgfqpoint{0.100000in}{0.100000in}}{\pgfqpoint{3.608454in}{2.310000in}}%
\pgfusepath{clip}%
\pgfsetbuttcap%
\pgfsetmiterjoin%
\definecolor{currentfill}{rgb}{0.000000,0.419608,0.790196}%
\pgfsetfillcolor{currentfill}%
\pgfsetlinewidth{0.000000pt}%
\definecolor{currentstroke}{rgb}{0.000000,0.000000,0.000000}%
\pgfsetstrokecolor{currentstroke}%
\pgfsetstrokeopacity{0.000000}%
\pgfsetdash{}{0pt}%
\pgfpathmoveto{\pgfqpoint{1.812489in}{1.842876in}}%
\pgfpathlineto{\pgfqpoint{1.843275in}{1.841097in}}%
\pgfpathlineto{\pgfqpoint{1.850854in}{1.840684in}}%
\pgfpathlineto{\pgfqpoint{1.849510in}{1.813058in}}%
\pgfpathlineto{\pgfqpoint{1.843805in}{1.813366in}}%
\pgfpathlineto{\pgfqpoint{1.802764in}{1.815650in}}%
\pgfpathlineto{\pgfqpoint{1.802162in}{1.823198in}}%
\pgfpathlineto{\pgfqpoint{1.809782in}{1.824532in}}%
\pgfpathlineto{\pgfqpoint{1.808240in}{1.831561in}}%
\pgfpathlineto{\pgfqpoint{1.812489in}{1.842876in}}%
\pgfpathclose%
\pgfusepath{fill}%
\end{pgfscope}%
\begin{pgfscope}%
\pgfpathrectangle{\pgfqpoint{0.100000in}{0.100000in}}{\pgfqpoint{3.608454in}{2.310000in}}%
\pgfusepath{clip}%
\pgfsetbuttcap%
\pgfsetmiterjoin%
\definecolor{currentfill}{rgb}{0.000000,0.666667,0.666667}%
\pgfsetfillcolor{currentfill}%
\pgfsetlinewidth{0.000000pt}%
\definecolor{currentstroke}{rgb}{0.000000,0.000000,0.000000}%
\pgfsetstrokecolor{currentstroke}%
\pgfsetstrokeopacity{0.000000}%
\pgfsetdash{}{0pt}%
\pgfpathmoveto{\pgfqpoint{1.498395in}{0.815195in}}%
\pgfpathlineto{\pgfqpoint{1.518670in}{0.813225in}}%
\pgfpathlineto{\pgfqpoint{1.562944in}{0.809149in}}%
\pgfpathlineto{\pgfqpoint{1.563559in}{0.816012in}}%
\pgfpathlineto{\pgfqpoint{1.581375in}{0.814353in}}%
\pgfpathlineto{\pgfqpoint{1.578420in}{0.779909in}}%
\pgfpathlineto{\pgfqpoint{1.580562in}{0.779720in}}%
\pgfpathlineto{\pgfqpoint{1.578119in}{0.751288in}}%
\pgfpathlineto{\pgfqpoint{1.574315in}{0.749565in}}%
\pgfpathlineto{\pgfqpoint{1.567213in}{0.756734in}}%
\pgfpathlineto{\pgfqpoint{1.562114in}{0.758888in}}%
\pgfpathlineto{\pgfqpoint{1.518651in}{0.714404in}}%
\pgfpathlineto{\pgfqpoint{1.486032in}{0.744751in}}%
\pgfpathlineto{\pgfqpoint{1.498395in}{0.815195in}}%
\pgfpathclose%
\pgfusepath{fill}%
\end{pgfscope}%
\begin{pgfscope}%
\pgfpathrectangle{\pgfqpoint{0.100000in}{0.100000in}}{\pgfqpoint{3.608454in}{2.310000in}}%
\pgfusepath{clip}%
\pgfsetbuttcap%
\pgfsetmiterjoin%
\definecolor{currentfill}{rgb}{0.000000,0.345098,0.827451}%
\pgfsetfillcolor{currentfill}%
\pgfsetlinewidth{0.000000pt}%
\definecolor{currentstroke}{rgb}{0.000000,0.000000,0.000000}%
\pgfsetstrokecolor{currentstroke}%
\pgfsetstrokeopacity{0.000000}%
\pgfsetdash{}{0pt}%
\pgfpathmoveto{\pgfqpoint{3.271268in}{1.222696in}}%
\pgfpathlineto{\pgfqpoint{3.265839in}{1.231923in}}%
\pgfpathlineto{\pgfqpoint{3.266234in}{1.243000in}}%
\pgfpathlineto{\pgfqpoint{3.259748in}{1.247063in}}%
\pgfpathlineto{\pgfqpoint{3.249894in}{1.248153in}}%
\pgfpathlineto{\pgfqpoint{3.248646in}{1.259422in}}%
\pgfpathlineto{\pgfqpoint{3.248223in}{1.267423in}}%
\pgfpathlineto{\pgfqpoint{3.266535in}{1.293082in}}%
\pgfpathlineto{\pgfqpoint{3.269493in}{1.294748in}}%
\pgfpathlineto{\pgfqpoint{3.275926in}{1.293375in}}%
\pgfpathlineto{\pgfqpoint{3.281080in}{1.300393in}}%
\pgfpathlineto{\pgfqpoint{3.293773in}{1.298157in}}%
\pgfpathlineto{\pgfqpoint{3.299684in}{1.300538in}}%
\pgfpathlineto{\pgfqpoint{3.305220in}{1.288703in}}%
\pgfpathlineto{\pgfqpoint{3.312824in}{1.275983in}}%
\pgfpathlineto{\pgfqpoint{3.324933in}{1.250084in}}%
\pgfpathlineto{\pgfqpoint{3.324107in}{1.249673in}}%
\pgfpathlineto{\pgfqpoint{3.316533in}{1.263940in}}%
\pgfpathlineto{\pgfqpoint{3.313407in}{1.258882in}}%
\pgfpathlineto{\pgfqpoint{3.326492in}{1.238123in}}%
\pgfpathlineto{\pgfqpoint{3.319080in}{1.242583in}}%
\pgfpathlineto{\pgfqpoint{3.311911in}{1.242756in}}%
\pgfpathlineto{\pgfqpoint{3.299704in}{1.233838in}}%
\pgfpathlineto{\pgfqpoint{3.289367in}{1.230321in}}%
\pgfpathlineto{\pgfqpoint{3.283250in}{1.223621in}}%
\pgfpathlineto{\pgfqpoint{3.271268in}{1.222696in}}%
\pgfpathclose%
\pgfusepath{fill}%
\end{pgfscope}%
\begin{pgfscope}%
\pgfpathrectangle{\pgfqpoint{0.100000in}{0.100000in}}{\pgfqpoint{3.608454in}{2.310000in}}%
\pgfusepath{clip}%
\pgfsetbuttcap%
\pgfsetmiterjoin%
\definecolor{currentfill}{rgb}{0.000000,0.776471,0.611765}%
\pgfsetfillcolor{currentfill}%
\pgfsetlinewidth{0.000000pt}%
\definecolor{currentstroke}{rgb}{0.000000,0.000000,0.000000}%
\pgfsetstrokecolor{currentstroke}%
\pgfsetstrokeopacity{0.000000}%
\pgfsetdash{}{0pt}%
\pgfpathmoveto{\pgfqpoint{2.002912in}{1.098705in}}%
\pgfpathlineto{\pgfqpoint{1.969632in}{1.099419in}}%
\pgfpathlineto{\pgfqpoint{1.956052in}{1.103169in}}%
\pgfpathlineto{\pgfqpoint{1.956419in}{1.117360in}}%
\pgfpathlineto{\pgfqpoint{1.949595in}{1.117555in}}%
\pgfpathlineto{\pgfqpoint{1.950456in}{1.151489in}}%
\pgfpathlineto{\pgfqpoint{1.984832in}{1.151230in}}%
\pgfpathlineto{\pgfqpoint{1.985971in}{1.147335in}}%
\pgfpathlineto{\pgfqpoint{1.977572in}{1.143978in}}%
\pgfpathlineto{\pgfqpoint{1.977272in}{1.130240in}}%
\pgfpathlineto{\pgfqpoint{1.984100in}{1.130082in}}%
\pgfpathlineto{\pgfqpoint{1.983944in}{1.123221in}}%
\pgfpathlineto{\pgfqpoint{1.990717in}{1.123080in}}%
\pgfpathlineto{\pgfqpoint{1.990571in}{1.116149in}}%
\pgfpathlineto{\pgfqpoint{2.003198in}{1.115990in}}%
\pgfpathlineto{\pgfqpoint{2.002912in}{1.098705in}}%
\pgfpathclose%
\pgfusepath{fill}%
\end{pgfscope}%
\begin{pgfscope}%
\pgfpathrectangle{\pgfqpoint{0.100000in}{0.100000in}}{\pgfqpoint{3.608454in}{2.310000in}}%
\pgfusepath{clip}%
\pgfsetbuttcap%
\pgfsetmiterjoin%
\definecolor{currentfill}{rgb}{0.000000,0.439216,0.780392}%
\pgfsetfillcolor{currentfill}%
\pgfsetlinewidth{0.000000pt}%
\definecolor{currentstroke}{rgb}{0.000000,0.000000,0.000000}%
\pgfsetstrokecolor{currentstroke}%
\pgfsetstrokeopacity{0.000000}%
\pgfsetdash{}{0pt}%
\pgfpathmoveto{\pgfqpoint{1.857457in}{1.534695in}}%
\pgfpathlineto{\pgfqpoint{1.855925in}{1.507277in}}%
\pgfpathlineto{\pgfqpoint{1.842917in}{1.507856in}}%
\pgfpathlineto{\pgfqpoint{1.795374in}{1.510393in}}%
\pgfpathlineto{\pgfqpoint{1.796934in}{1.565224in}}%
\pgfpathlineto{\pgfqpoint{1.762096in}{1.567407in}}%
\pgfpathlineto{\pgfqpoint{1.763649in}{1.594885in}}%
\pgfpathlineto{\pgfqpoint{1.803419in}{1.592288in}}%
\pgfpathlineto{\pgfqpoint{1.831612in}{1.590776in}}%
\pgfpathlineto{\pgfqpoint{1.830208in}{1.563450in}}%
\pgfpathlineto{\pgfqpoint{1.858208in}{1.562101in}}%
\pgfpathlineto{\pgfqpoint{1.857457in}{1.534695in}}%
\pgfpathclose%
\pgfusepath{fill}%
\end{pgfscope}%
\begin{pgfscope}%
\pgfpathrectangle{\pgfqpoint{0.100000in}{0.100000in}}{\pgfqpoint{3.608454in}{2.310000in}}%
\pgfusepath{clip}%
\pgfsetbuttcap%
\pgfsetmiterjoin%
\definecolor{currentfill}{rgb}{0.000000,0.576471,0.711765}%
\pgfsetfillcolor{currentfill}%
\pgfsetlinewidth{0.000000pt}%
\definecolor{currentstroke}{rgb}{0.000000,0.000000,0.000000}%
\pgfsetstrokecolor{currentstroke}%
\pgfsetstrokeopacity{0.000000}%
\pgfsetdash{}{0pt}%
\pgfpathmoveto{\pgfqpoint{2.410748in}{1.065840in}}%
\pgfpathlineto{\pgfqpoint{2.389852in}{1.065601in}}%
\pgfpathlineto{\pgfqpoint{2.362249in}{1.063873in}}%
\pgfpathlineto{\pgfqpoint{2.361164in}{1.098864in}}%
\pgfpathlineto{\pgfqpoint{2.350507in}{1.098801in}}%
\pgfpathlineto{\pgfqpoint{2.340339in}{1.098431in}}%
\pgfpathlineto{\pgfqpoint{2.339731in}{1.125367in}}%
\pgfpathlineto{\pgfqpoint{2.345455in}{1.126590in}}%
\pgfpathlineto{\pgfqpoint{2.343960in}{1.136912in}}%
\pgfpathlineto{\pgfqpoint{2.335244in}{1.146265in}}%
\pgfpathlineto{\pgfqpoint{2.353073in}{1.147045in}}%
\pgfpathlineto{\pgfqpoint{2.410619in}{1.149798in}}%
\pgfpathlineto{\pgfqpoint{2.421101in}{1.141357in}}%
\pgfpathlineto{\pgfqpoint{2.420239in}{1.132410in}}%
\pgfpathlineto{\pgfqpoint{2.411938in}{1.125058in}}%
\pgfpathlineto{\pgfqpoint{2.406026in}{1.117238in}}%
\pgfpathlineto{\pgfqpoint{2.408349in}{1.109943in}}%
\pgfpathlineto{\pgfqpoint{2.410748in}{1.065840in}}%
\pgfpathclose%
\pgfusepath{fill}%
\end{pgfscope}%
\begin{pgfscope}%
\pgfpathrectangle{\pgfqpoint{0.100000in}{0.100000in}}{\pgfqpoint{3.608454in}{2.310000in}}%
\pgfusepath{clip}%
\pgfsetbuttcap%
\pgfsetmiterjoin%
\definecolor{currentfill}{rgb}{0.000000,0.650980,0.674510}%
\pgfsetfillcolor{currentfill}%
\pgfsetlinewidth{0.000000pt}%
\definecolor{currentstroke}{rgb}{0.000000,0.000000,0.000000}%
\pgfsetstrokecolor{currentstroke}%
\pgfsetstrokeopacity{0.000000}%
\pgfsetdash{}{0pt}%
\pgfpathmoveto{\pgfqpoint{1.383780in}{1.226748in}}%
\pgfpathlineto{\pgfqpoint{1.373495in}{1.246788in}}%
\pgfpathlineto{\pgfqpoint{1.375196in}{1.259964in}}%
\pgfpathlineto{\pgfqpoint{1.347103in}{1.263333in}}%
\pgfpathlineto{\pgfqpoint{1.347428in}{1.265725in}}%
\pgfpathlineto{\pgfqpoint{1.349065in}{1.285654in}}%
\pgfpathlineto{\pgfqpoint{1.352269in}{1.306231in}}%
\pgfpathlineto{\pgfqpoint{1.360950in}{1.306443in}}%
\pgfpathlineto{\pgfqpoint{1.365837in}{1.343240in}}%
\pgfpathlineto{\pgfqpoint{1.412097in}{1.337063in}}%
\pgfpathlineto{\pgfqpoint{1.422616in}{1.335846in}}%
\pgfpathlineto{\pgfqpoint{1.426840in}{1.337144in}}%
\pgfpathlineto{\pgfqpoint{1.444613in}{1.311073in}}%
\pgfpathlineto{\pgfqpoint{1.448675in}{1.296375in}}%
\pgfpathlineto{\pgfqpoint{1.454881in}{1.289984in}}%
\pgfpathlineto{\pgfqpoint{1.457227in}{1.286192in}}%
\pgfpathlineto{\pgfqpoint{1.450742in}{1.268060in}}%
\pgfpathlineto{\pgfqpoint{1.463748in}{1.269672in}}%
\pgfpathlineto{\pgfqpoint{1.469792in}{1.261228in}}%
\pgfpathlineto{\pgfqpoint{1.471264in}{1.248705in}}%
\pgfpathlineto{\pgfqpoint{1.469341in}{1.240362in}}%
\pgfpathlineto{\pgfqpoint{1.466660in}{1.216957in}}%
\pgfpathlineto{\pgfqpoint{1.462544in}{1.217413in}}%
\pgfpathlineto{\pgfqpoint{1.383780in}{1.226748in}}%
\pgfpathclose%
\pgfusepath{fill}%
\end{pgfscope}%
\begin{pgfscope}%
\pgfpathrectangle{\pgfqpoint{0.100000in}{0.100000in}}{\pgfqpoint{3.608454in}{2.310000in}}%
\pgfusepath{clip}%
\pgfsetbuttcap%
\pgfsetmiterjoin%
\definecolor{currentfill}{rgb}{0.000000,0.427451,0.786275}%
\pgfsetfillcolor{currentfill}%
\pgfsetlinewidth{0.000000pt}%
\definecolor{currentstroke}{rgb}{0.000000,0.000000,0.000000}%
\pgfsetstrokecolor{currentstroke}%
\pgfsetstrokeopacity{0.000000}%
\pgfsetdash{}{0pt}%
\pgfpathmoveto{\pgfqpoint{2.204405in}{1.465435in}}%
\pgfpathlineto{\pgfqpoint{2.193447in}{1.465250in}}%
\pgfpathlineto{\pgfqpoint{2.192689in}{1.511198in}}%
\pgfpathlineto{\pgfqpoint{2.206358in}{1.511376in}}%
\pgfpathlineto{\pgfqpoint{2.219965in}{1.511654in}}%
\pgfpathlineto{\pgfqpoint{2.220496in}{1.490915in}}%
\pgfpathlineto{\pgfqpoint{2.247888in}{1.491592in}}%
\pgfpathlineto{\pgfqpoint{2.248598in}{1.467244in}}%
\pgfpathlineto{\pgfqpoint{2.244011in}{1.467014in}}%
\pgfpathlineto{\pgfqpoint{2.204405in}{1.465435in}}%
\pgfpathclose%
\pgfusepath{fill}%
\end{pgfscope}%
\begin{pgfscope}%
\pgfpathrectangle{\pgfqpoint{0.100000in}{0.100000in}}{\pgfqpoint{3.608454in}{2.310000in}}%
\pgfusepath{clip}%
\pgfsetbuttcap%
\pgfsetmiterjoin%
\definecolor{currentfill}{rgb}{0.000000,0.180392,0.909804}%
\pgfsetfillcolor{currentfill}%
\pgfsetlinewidth{0.000000pt}%
\definecolor{currentstroke}{rgb}{0.000000,0.000000,0.000000}%
\pgfsetstrokecolor{currentstroke}%
\pgfsetstrokeopacity{0.000000}%
\pgfsetdash{}{0pt}%
\pgfpathmoveto{\pgfqpoint{1.945778in}{0.610158in}}%
\pgfpathlineto{\pgfqpoint{1.925990in}{0.598363in}}%
\pgfpathlineto{\pgfqpoint{1.923573in}{0.599973in}}%
\pgfpathlineto{\pgfqpoint{1.918807in}{0.601623in}}%
\pgfpathlineto{\pgfqpoint{1.907278in}{0.617155in}}%
\pgfpathlineto{\pgfqpoint{1.898375in}{0.608966in}}%
\pgfpathlineto{\pgfqpoint{1.896461in}{0.616726in}}%
\pgfpathlineto{\pgfqpoint{1.878483in}{0.632521in}}%
\pgfpathlineto{\pgfqpoint{1.870125in}{0.624788in}}%
\pgfpathlineto{\pgfqpoint{1.858713in}{0.641352in}}%
\pgfpathlineto{\pgfqpoint{1.859593in}{0.670192in}}%
\pgfpathlineto{\pgfqpoint{1.871994in}{0.669892in}}%
\pgfpathlineto{\pgfqpoint{1.891457in}{0.663121in}}%
\pgfpathlineto{\pgfqpoint{1.892292in}{0.669010in}}%
\pgfpathlineto{\pgfqpoint{1.903598in}{0.691454in}}%
\pgfpathlineto{\pgfqpoint{1.913119in}{0.701138in}}%
\pgfpathlineto{\pgfqpoint{1.926963in}{0.697491in}}%
\pgfpathlineto{\pgfqpoint{1.947915in}{0.687499in}}%
\pgfpathlineto{\pgfqpoint{1.952059in}{0.698355in}}%
\pgfpathlineto{\pgfqpoint{1.965169in}{0.705757in}}%
\pgfpathlineto{\pgfqpoint{1.981851in}{0.714949in}}%
\pgfpathlineto{\pgfqpoint{1.984720in}{0.711352in}}%
\pgfpathlineto{\pgfqpoint{1.986509in}{0.700655in}}%
\pgfpathlineto{\pgfqpoint{1.993815in}{0.692311in}}%
\pgfpathlineto{\pgfqpoint{1.995500in}{0.684736in}}%
\pgfpathlineto{\pgfqpoint{1.971709in}{0.671338in}}%
\pgfpathlineto{\pgfqpoint{1.975419in}{0.668492in}}%
\pgfpathlineto{\pgfqpoint{1.978079in}{0.661346in}}%
\pgfpathlineto{\pgfqpoint{1.986455in}{0.652084in}}%
\pgfpathlineto{\pgfqpoint{1.993561in}{0.650091in}}%
\pgfpathlineto{\pgfqpoint{1.987116in}{0.645931in}}%
\pgfpathlineto{\pgfqpoint{1.982721in}{0.639405in}}%
\pgfpathlineto{\pgfqpoint{1.979102in}{0.635917in}}%
\pgfpathlineto{\pgfqpoint{1.966633in}{0.631014in}}%
\pgfpathlineto{\pgfqpoint{1.945778in}{0.610158in}}%
\pgfpathclose%
\pgfusepath{fill}%
\end{pgfscope}%
\begin{pgfscope}%
\pgfpathrectangle{\pgfqpoint{0.100000in}{0.100000in}}{\pgfqpoint{3.608454in}{2.310000in}}%
\pgfusepath{clip}%
\pgfsetbuttcap%
\pgfsetmiterjoin%
\definecolor{currentfill}{rgb}{0.000000,0.690196,0.654902}%
\pgfsetfillcolor{currentfill}%
\pgfsetlinewidth{0.000000pt}%
\definecolor{currentstroke}{rgb}{0.000000,0.000000,0.000000}%
\pgfsetstrokecolor{currentstroke}%
\pgfsetstrokeopacity{0.000000}%
\pgfsetdash{}{0pt}%
\pgfpathmoveto{\pgfqpoint{2.288408in}{0.781863in}}%
\pgfpathlineto{\pgfqpoint{2.287535in}{0.813168in}}%
\pgfpathlineto{\pgfqpoint{2.280375in}{0.819890in}}%
\pgfpathlineto{\pgfqpoint{2.279806in}{0.840743in}}%
\pgfpathlineto{\pgfqpoint{2.272892in}{0.840599in}}%
\pgfpathlineto{\pgfqpoint{2.267654in}{0.847225in}}%
\pgfpathlineto{\pgfqpoint{2.258858in}{0.847301in}}%
\pgfpathlineto{\pgfqpoint{2.258381in}{0.867551in}}%
\pgfpathlineto{\pgfqpoint{2.302086in}{0.868312in}}%
\pgfpathlineto{\pgfqpoint{2.342635in}{0.869578in}}%
\pgfpathlineto{\pgfqpoint{2.344274in}{0.869656in}}%
\pgfpathlineto{\pgfqpoint{2.340706in}{0.860346in}}%
\pgfpathlineto{\pgfqpoint{2.337645in}{0.859095in}}%
\pgfpathlineto{\pgfqpoint{2.331640in}{0.845019in}}%
\pgfpathlineto{\pgfqpoint{2.334864in}{0.835567in}}%
\pgfpathlineto{\pgfqpoint{2.344959in}{0.835935in}}%
\pgfpathlineto{\pgfqpoint{2.344269in}{0.832393in}}%
\pgfpathlineto{\pgfqpoint{2.342569in}{0.828917in}}%
\pgfpathlineto{\pgfqpoint{2.344014in}{0.814843in}}%
\pgfpathlineto{\pgfqpoint{2.340198in}{0.812791in}}%
\pgfpathlineto{\pgfqpoint{2.342946in}{0.805709in}}%
\pgfpathlineto{\pgfqpoint{2.343388in}{0.795089in}}%
\pgfpathlineto{\pgfqpoint{2.338322in}{0.779971in}}%
\pgfpathlineto{\pgfqpoint{2.336629in}{0.786857in}}%
\pgfpathlineto{\pgfqpoint{2.324523in}{0.778629in}}%
\pgfpathlineto{\pgfqpoint{2.319400in}{0.785321in}}%
\pgfpathlineto{\pgfqpoint{2.315987in}{0.782756in}}%
\pgfpathlineto{\pgfqpoint{2.288408in}{0.781863in}}%
\pgfpathclose%
\pgfusepath{fill}%
\end{pgfscope}%
\begin{pgfscope}%
\pgfpathrectangle{\pgfqpoint{0.100000in}{0.100000in}}{\pgfqpoint{3.608454in}{2.310000in}}%
\pgfusepath{clip}%
\pgfsetbuttcap%
\pgfsetmiterjoin%
\definecolor{currentfill}{rgb}{0.000000,0.776471,0.611765}%
\pgfsetfillcolor{currentfill}%
\pgfsetlinewidth{0.000000pt}%
\definecolor{currentstroke}{rgb}{0.000000,0.000000,0.000000}%
\pgfsetstrokecolor{currentstroke}%
\pgfsetstrokeopacity{0.000000}%
\pgfsetdash{}{0pt}%
\pgfpathmoveto{\pgfqpoint{2.611522in}{1.917144in}}%
\pgfpathlineto{\pgfqpoint{2.589883in}{1.913322in}}%
\pgfpathlineto{\pgfqpoint{2.586900in}{1.909857in}}%
\pgfpathlineto{\pgfqpoint{2.586216in}{1.900993in}}%
\pgfpathlineto{\pgfqpoint{2.576778in}{1.896148in}}%
\pgfpathlineto{\pgfqpoint{2.572058in}{1.888198in}}%
\pgfpathlineto{\pgfqpoint{2.566665in}{1.889771in}}%
\pgfpathlineto{\pgfqpoint{2.572918in}{1.898549in}}%
\pgfpathlineto{\pgfqpoint{2.575837in}{1.905644in}}%
\pgfpathlineto{\pgfqpoint{2.562752in}{1.899838in}}%
\pgfpathlineto{\pgfqpoint{2.559587in}{1.892271in}}%
\pgfpathlineto{\pgfqpoint{2.552985in}{1.887697in}}%
\pgfpathlineto{\pgfqpoint{2.547425in}{1.892494in}}%
\pgfpathlineto{\pgfqpoint{2.541575in}{1.885933in}}%
\pgfpathlineto{\pgfqpoint{2.537283in}{1.876850in}}%
\pgfpathlineto{\pgfqpoint{2.533757in}{1.876549in}}%
\pgfpathlineto{\pgfqpoint{2.531516in}{1.904166in}}%
\pgfpathlineto{\pgfqpoint{2.528620in}{1.910846in}}%
\pgfpathlineto{\pgfqpoint{2.514867in}{1.909757in}}%
\pgfpathlineto{\pgfqpoint{2.513243in}{1.930478in}}%
\pgfpathlineto{\pgfqpoint{2.485703in}{1.928298in}}%
\pgfpathlineto{\pgfqpoint{2.484672in}{1.942048in}}%
\pgfpathlineto{\pgfqpoint{2.483177in}{1.962621in}}%
\pgfpathlineto{\pgfqpoint{2.489953in}{1.963218in}}%
\pgfpathlineto{\pgfqpoint{2.486521in}{1.969842in}}%
\pgfpathlineto{\pgfqpoint{2.485656in}{1.981514in}}%
\pgfpathlineto{\pgfqpoint{2.493445in}{1.982175in}}%
\pgfpathlineto{\pgfqpoint{2.508763in}{1.975427in}}%
\pgfpathlineto{\pgfqpoint{2.511603in}{1.969882in}}%
\pgfpathlineto{\pgfqpoint{2.526869in}{1.951632in}}%
\pgfpathlineto{\pgfqpoint{2.538957in}{1.952508in}}%
\pgfpathlineto{\pgfqpoint{2.544000in}{1.955889in}}%
\pgfpathlineto{\pgfqpoint{2.552498in}{1.948973in}}%
\pgfpathlineto{\pgfqpoint{2.564052in}{1.959757in}}%
\pgfpathlineto{\pgfqpoint{2.568367in}{1.953169in}}%
\pgfpathlineto{\pgfqpoint{2.574148in}{1.960469in}}%
\pgfpathlineto{\pgfqpoint{2.592197in}{1.970111in}}%
\pgfpathlineto{\pgfqpoint{2.605875in}{1.974408in}}%
\pgfpathlineto{\pgfqpoint{2.611522in}{1.917144in}}%
\pgfpathclose%
\pgfusepath{fill}%
\end{pgfscope}%
\begin{pgfscope}%
\pgfpathrectangle{\pgfqpoint{0.100000in}{0.100000in}}{\pgfqpoint{3.608454in}{2.310000in}}%
\pgfusepath{clip}%
\pgfsetbuttcap%
\pgfsetmiterjoin%
\definecolor{currentfill}{rgb}{0.000000,0.537255,0.731373}%
\pgfsetfillcolor{currentfill}%
\pgfsetlinewidth{0.000000pt}%
\definecolor{currentstroke}{rgb}{0.000000,0.000000,0.000000}%
\pgfsetstrokecolor{currentstroke}%
\pgfsetstrokeopacity{0.000000}%
\pgfsetdash{}{0pt}%
\pgfpathmoveto{\pgfqpoint{3.284382in}{1.624557in}}%
\pgfpathlineto{\pgfqpoint{3.276398in}{1.614345in}}%
\pgfpathlineto{\pgfqpoint{3.266342in}{1.608149in}}%
\pgfpathlineto{\pgfqpoint{3.258629in}{1.617098in}}%
\pgfpathlineto{\pgfqpoint{3.262436in}{1.621007in}}%
\pgfpathlineto{\pgfqpoint{3.252558in}{1.629161in}}%
\pgfpathlineto{\pgfqpoint{3.248335in}{1.629894in}}%
\pgfpathlineto{\pgfqpoint{3.244631in}{1.624617in}}%
\pgfpathlineto{\pgfqpoint{3.248038in}{1.619832in}}%
\pgfpathlineto{\pgfqpoint{3.234122in}{1.609162in}}%
\pgfpathlineto{\pgfqpoint{3.221218in}{1.609411in}}%
\pgfpathlineto{\pgfqpoint{3.217800in}{1.603350in}}%
\pgfpathlineto{\pgfqpoint{3.217800in}{1.596744in}}%
\pgfpathlineto{\pgfqpoint{3.213947in}{1.593817in}}%
\pgfpathlineto{\pgfqpoint{3.211501in}{1.598059in}}%
\pgfpathlineto{\pgfqpoint{3.206320in}{1.596141in}}%
\pgfpathlineto{\pgfqpoint{3.203466in}{1.605423in}}%
\pgfpathlineto{\pgfqpoint{3.195360in}{1.613257in}}%
\pgfpathlineto{\pgfqpoint{3.192630in}{1.620121in}}%
\pgfpathlineto{\pgfqpoint{3.195407in}{1.620855in}}%
\pgfpathlineto{\pgfqpoint{3.203745in}{1.634729in}}%
\pgfpathlineto{\pgfqpoint{3.209405in}{1.635945in}}%
\pgfpathlineto{\pgfqpoint{3.210875in}{1.652502in}}%
\pgfpathlineto{\pgfqpoint{3.210379in}{1.663470in}}%
\pgfpathlineto{\pgfqpoint{3.215154in}{1.664914in}}%
\pgfpathlineto{\pgfqpoint{3.253079in}{1.672195in}}%
\pgfpathlineto{\pgfqpoint{3.245833in}{1.699582in}}%
\pgfpathlineto{\pgfqpoint{3.252945in}{1.701154in}}%
\pgfpathlineto{\pgfqpoint{3.258600in}{1.697493in}}%
\pgfpathlineto{\pgfqpoint{3.260640in}{1.692420in}}%
\pgfpathlineto{\pgfqpoint{3.267733in}{1.692380in}}%
\pgfpathlineto{\pgfqpoint{3.274806in}{1.685909in}}%
\pgfpathlineto{\pgfqpoint{3.277583in}{1.676067in}}%
\pgfpathlineto{\pgfqpoint{3.282783in}{1.667998in}}%
\pgfpathlineto{\pgfqpoint{3.292838in}{1.664005in}}%
\pgfpathlineto{\pgfqpoint{3.298905in}{1.664653in}}%
\pgfpathlineto{\pgfqpoint{3.302353in}{1.660141in}}%
\pgfpathlineto{\pgfqpoint{3.295750in}{1.652950in}}%
\pgfpathlineto{\pgfqpoint{3.294577in}{1.644062in}}%
\pgfpathlineto{\pgfqpoint{3.284382in}{1.624557in}}%
\pgfpathclose%
\pgfusepath{fill}%
\end{pgfscope}%
\begin{pgfscope}%
\pgfpathrectangle{\pgfqpoint{0.100000in}{0.100000in}}{\pgfqpoint{3.608454in}{2.310000in}}%
\pgfusepath{clip}%
\pgfsetbuttcap%
\pgfsetmiterjoin%
\definecolor{currentfill}{rgb}{0.000000,0.364706,0.817647}%
\pgfsetfillcolor{currentfill}%
\pgfsetlinewidth{0.000000pt}%
\definecolor{currentstroke}{rgb}{0.000000,0.000000,0.000000}%
\pgfsetstrokecolor{currentstroke}%
\pgfsetstrokeopacity{0.000000}%
\pgfsetdash{}{0pt}%
\pgfpathmoveto{\pgfqpoint{1.858721in}{1.355685in}}%
\pgfpathlineto{\pgfqpoint{1.892889in}{1.354245in}}%
\pgfpathlineto{\pgfqpoint{1.892401in}{1.333545in}}%
\pgfpathlineto{\pgfqpoint{1.891797in}{1.319820in}}%
\pgfpathlineto{\pgfqpoint{1.824072in}{1.322793in}}%
\pgfpathlineto{\pgfqpoint{1.789033in}{1.324734in}}%
\pgfpathlineto{\pgfqpoint{1.791350in}{1.359031in}}%
\pgfpathlineto{\pgfqpoint{1.858721in}{1.355685in}}%
\pgfpathclose%
\pgfusepath{fill}%
\end{pgfscope}%
\begin{pgfscope}%
\pgfpathrectangle{\pgfqpoint{0.100000in}{0.100000in}}{\pgfqpoint{3.608454in}{2.310000in}}%
\pgfusepath{clip}%
\pgfsetbuttcap%
\pgfsetmiterjoin%
\definecolor{currentfill}{rgb}{0.000000,0.827451,0.586275}%
\pgfsetfillcolor{currentfill}%
\pgfsetlinewidth{0.000000pt}%
\definecolor{currentstroke}{rgb}{0.000000,0.000000,0.000000}%
\pgfsetstrokecolor{currentstroke}%
\pgfsetstrokeopacity{0.000000}%
\pgfsetdash{}{0pt}%
\pgfpathmoveto{\pgfqpoint{0.592949in}{0.431644in}}%
\pgfpathlineto{\pgfqpoint{0.591752in}{0.435712in}}%
\pgfpathlineto{\pgfqpoint{0.594116in}{0.436502in}}%
\pgfpathlineto{\pgfqpoint{0.596254in}{0.435454in}}%
\pgfpathlineto{\pgfqpoint{0.596813in}{0.431481in}}%
\pgfpathlineto{\pgfqpoint{0.595244in}{0.430874in}}%
\pgfpathlineto{\pgfqpoint{0.592949in}{0.431644in}}%
\pgfpathclose%
\pgfusepath{fill}%
\end{pgfscope}%
\begin{pgfscope}%
\pgfpathrectangle{\pgfqpoint{0.100000in}{0.100000in}}{\pgfqpoint{3.608454in}{2.310000in}}%
\pgfusepath{clip}%
\pgfsetbuttcap%
\pgfsetmiterjoin%
\definecolor{currentfill}{rgb}{0.000000,0.827451,0.586275}%
\pgfsetfillcolor{currentfill}%
\pgfsetlinewidth{0.000000pt}%
\definecolor{currentstroke}{rgb}{0.000000,0.000000,0.000000}%
\pgfsetstrokecolor{currentstroke}%
\pgfsetstrokeopacity{0.000000}%
\pgfsetdash{}{0pt}%
\pgfpathmoveto{\pgfqpoint{0.598016in}{0.429456in}}%
\pgfpathlineto{\pgfqpoint{0.599382in}{0.432202in}}%
\pgfpathlineto{\pgfqpoint{0.600924in}{0.431678in}}%
\pgfpathlineto{\pgfqpoint{0.601971in}{0.429670in}}%
\pgfpathlineto{\pgfqpoint{0.598016in}{0.429456in}}%
\pgfpathclose%
\pgfusepath{fill}%
\end{pgfscope}%
\begin{pgfscope}%
\pgfpathrectangle{\pgfqpoint{0.100000in}{0.100000in}}{\pgfqpoint{3.608454in}{2.310000in}}%
\pgfusepath{clip}%
\pgfsetbuttcap%
\pgfsetmiterjoin%
\definecolor{currentfill}{rgb}{0.000000,0.827451,0.586275}%
\pgfsetfillcolor{currentfill}%
\pgfsetlinewidth{0.000000pt}%
\definecolor{currentstroke}{rgb}{0.000000,0.000000,0.000000}%
\pgfsetstrokecolor{currentstroke}%
\pgfsetstrokeopacity{0.000000}%
\pgfsetdash{}{0pt}%
\pgfpathmoveto{\pgfqpoint{0.679037in}{0.397130in}}%
\pgfpathlineto{\pgfqpoint{0.677306in}{0.394866in}}%
\pgfpathlineto{\pgfqpoint{0.676295in}{0.398443in}}%
\pgfpathlineto{\pgfqpoint{0.673518in}{0.397756in}}%
\pgfpathlineto{\pgfqpoint{0.676763in}{0.402425in}}%
\pgfpathlineto{\pgfqpoint{0.680361in}{0.402575in}}%
\pgfpathlineto{\pgfqpoint{0.679428in}{0.400744in}}%
\pgfpathlineto{\pgfqpoint{0.680747in}{0.399061in}}%
\pgfpathlineto{\pgfqpoint{0.683725in}{0.397976in}}%
\pgfpathlineto{\pgfqpoint{0.681665in}{0.395162in}}%
\pgfpathlineto{\pgfqpoint{0.680358in}{0.397560in}}%
\pgfpathlineto{\pgfqpoint{0.679037in}{0.397130in}}%
\pgfpathclose%
\pgfusepath{fill}%
\end{pgfscope}%
\begin{pgfscope}%
\pgfpathrectangle{\pgfqpoint{0.100000in}{0.100000in}}{\pgfqpoint{3.608454in}{2.310000in}}%
\pgfusepath{clip}%
\pgfsetbuttcap%
\pgfsetmiterjoin%
\definecolor{currentfill}{rgb}{0.000000,0.827451,0.586275}%
\pgfsetfillcolor{currentfill}%
\pgfsetlinewidth{0.000000pt}%
\definecolor{currentstroke}{rgb}{0.000000,0.000000,0.000000}%
\pgfsetstrokecolor{currentstroke}%
\pgfsetstrokeopacity{0.000000}%
\pgfsetdash{}{0pt}%
\pgfpathmoveto{\pgfqpoint{0.650134in}{0.407855in}}%
\pgfpathlineto{\pgfqpoint{0.650363in}{0.410897in}}%
\pgfpathlineto{\pgfqpoint{0.653136in}{0.411103in}}%
\pgfpathlineto{\pgfqpoint{0.652771in}{0.408615in}}%
\pgfpathlineto{\pgfqpoint{0.650134in}{0.407855in}}%
\pgfpathclose%
\pgfusepath{fill}%
\end{pgfscope}%
\begin{pgfscope}%
\pgfpathrectangle{\pgfqpoint{0.100000in}{0.100000in}}{\pgfqpoint{3.608454in}{2.310000in}}%
\pgfusepath{clip}%
\pgfsetbuttcap%
\pgfsetmiterjoin%
\definecolor{currentfill}{rgb}{0.000000,0.827451,0.586275}%
\pgfsetfillcolor{currentfill}%
\pgfsetlinewidth{0.000000pt}%
\definecolor{currentstroke}{rgb}{0.000000,0.000000,0.000000}%
\pgfsetstrokecolor{currentstroke}%
\pgfsetstrokeopacity{0.000000}%
\pgfsetdash{}{0pt}%
\pgfpathmoveto{\pgfqpoint{0.676688in}{0.386431in}}%
\pgfpathlineto{\pgfqpoint{0.683893in}{0.390538in}}%
\pgfpathlineto{\pgfqpoint{0.685712in}{0.388279in}}%
\pgfpathlineto{\pgfqpoint{0.683374in}{0.387317in}}%
\pgfpathlineto{\pgfqpoint{0.676688in}{0.386431in}}%
\pgfpathclose%
\pgfusepath{fill}%
\end{pgfscope}%
\begin{pgfscope}%
\pgfpathrectangle{\pgfqpoint{0.100000in}{0.100000in}}{\pgfqpoint{3.608454in}{2.310000in}}%
\pgfusepath{clip}%
\pgfsetbuttcap%
\pgfsetmiterjoin%
\definecolor{currentfill}{rgb}{0.000000,0.827451,0.586275}%
\pgfsetfillcolor{currentfill}%
\pgfsetlinewidth{0.000000pt}%
\definecolor{currentstroke}{rgb}{0.000000,0.000000,0.000000}%
\pgfsetstrokecolor{currentstroke}%
\pgfsetstrokeopacity{0.000000}%
\pgfsetdash{}{0pt}%
\pgfpathmoveto{\pgfqpoint{0.698256in}{0.395576in}}%
\pgfpathlineto{\pgfqpoint{0.696998in}{0.397305in}}%
\pgfpathlineto{\pgfqpoint{0.701060in}{0.399308in}}%
\pgfpathlineto{\pgfqpoint{0.699966in}{0.402043in}}%
\pgfpathlineto{\pgfqpoint{0.697786in}{0.402037in}}%
\pgfpathlineto{\pgfqpoint{0.696446in}{0.403146in}}%
\pgfpathlineto{\pgfqpoint{0.691042in}{0.403748in}}%
\pgfpathlineto{\pgfqpoint{0.686539in}{0.403836in}}%
\pgfpathlineto{\pgfqpoint{0.684692in}{0.405317in}}%
\pgfpathlineto{\pgfqpoint{0.682874in}{0.404197in}}%
\pgfpathlineto{\pgfqpoint{0.680377in}{0.404275in}}%
\pgfpathlineto{\pgfqpoint{0.680706in}{0.406672in}}%
\pgfpathlineto{\pgfqpoint{0.675294in}{0.406801in}}%
\pgfpathlineto{\pgfqpoint{0.672429in}{0.407601in}}%
\pgfpathlineto{\pgfqpoint{0.669990in}{0.409978in}}%
\pgfpathlineto{\pgfqpoint{0.672102in}{0.411831in}}%
\pgfpathlineto{\pgfqpoint{0.676413in}{0.414087in}}%
\pgfpathlineto{\pgfqpoint{0.674276in}{0.416221in}}%
\pgfpathlineto{\pgfqpoint{0.672028in}{0.416339in}}%
\pgfpathlineto{\pgfqpoint{0.667446in}{0.412698in}}%
\pgfpathlineto{\pgfqpoint{0.664420in}{0.412115in}}%
\pgfpathlineto{\pgfqpoint{0.661755in}{0.412674in}}%
\pgfpathlineto{\pgfqpoint{0.660202in}{0.409906in}}%
\pgfpathlineto{\pgfqpoint{0.657982in}{0.409975in}}%
\pgfpathlineto{\pgfqpoint{0.652958in}{0.413217in}}%
\pgfpathlineto{\pgfqpoint{0.652756in}{0.415126in}}%
\pgfpathlineto{\pgfqpoint{0.654783in}{0.415315in}}%
\pgfpathlineto{\pgfqpoint{0.655313in}{0.417608in}}%
\pgfpathlineto{\pgfqpoint{0.654939in}{0.420817in}}%
\pgfpathlineto{\pgfqpoint{0.650913in}{0.416098in}}%
\pgfpathlineto{\pgfqpoint{0.650688in}{0.414349in}}%
\pgfpathlineto{\pgfqpoint{0.647513in}{0.414740in}}%
\pgfpathlineto{\pgfqpoint{0.645402in}{0.416502in}}%
\pgfpathlineto{\pgfqpoint{0.646738in}{0.419922in}}%
\pgfpathlineto{\pgfqpoint{0.645503in}{0.423049in}}%
\pgfpathlineto{\pgfqpoint{0.644228in}{0.422685in}}%
\pgfpathlineto{\pgfqpoint{0.644008in}{0.418053in}}%
\pgfpathlineto{\pgfqpoint{0.636879in}{0.418407in}}%
\pgfpathlineto{\pgfqpoint{0.638459in}{0.416408in}}%
\pgfpathlineto{\pgfqpoint{0.636161in}{0.415808in}}%
\pgfpathlineto{\pgfqpoint{0.635340in}{0.417833in}}%
\pgfpathlineto{\pgfqpoint{0.633246in}{0.416909in}}%
\pgfpathlineto{\pgfqpoint{0.629966in}{0.419033in}}%
\pgfpathlineto{\pgfqpoint{0.625891in}{0.422864in}}%
\pgfpathlineto{\pgfqpoint{0.621156in}{0.424136in}}%
\pgfpathlineto{\pgfqpoint{0.617033in}{0.423185in}}%
\pgfpathlineto{\pgfqpoint{0.614359in}{0.424662in}}%
\pgfpathlineto{\pgfqpoint{0.613000in}{0.429251in}}%
\pgfpathlineto{\pgfqpoint{0.614183in}{0.431739in}}%
\pgfpathlineto{\pgfqpoint{0.619670in}{0.431167in}}%
\pgfpathlineto{\pgfqpoint{0.626353in}{0.433684in}}%
\pgfpathlineto{\pgfqpoint{0.627719in}{0.431363in}}%
\pgfpathlineto{\pgfqpoint{0.630050in}{0.430192in}}%
\pgfpathlineto{\pgfqpoint{0.631614in}{0.430608in}}%
\pgfpathlineto{\pgfqpoint{0.636883in}{0.428698in}}%
\pgfpathlineto{\pgfqpoint{0.640000in}{0.425903in}}%
\pgfpathlineto{\pgfqpoint{0.638907in}{0.421869in}}%
\pgfpathlineto{\pgfqpoint{0.640826in}{0.421267in}}%
\pgfpathlineto{\pgfqpoint{0.643623in}{0.424929in}}%
\pgfpathlineto{\pgfqpoint{0.647952in}{0.423982in}}%
\pgfpathlineto{\pgfqpoint{0.649233in}{0.421491in}}%
\pgfpathlineto{\pgfqpoint{0.651625in}{0.421666in}}%
\pgfpathlineto{\pgfqpoint{0.667655in}{0.424721in}}%
\pgfpathlineto{\pgfqpoint{0.673676in}{0.424199in}}%
\pgfpathlineto{\pgfqpoint{0.676435in}{0.424381in}}%
\pgfpathlineto{\pgfqpoint{0.682895in}{0.421450in}}%
\pgfpathlineto{\pgfqpoint{0.686743in}{0.417731in}}%
\pgfpathlineto{\pgfqpoint{0.685490in}{0.416186in}}%
\pgfpathlineto{\pgfqpoint{0.685449in}{0.411697in}}%
\pgfpathlineto{\pgfqpoint{0.687707in}{0.412299in}}%
\pgfpathlineto{\pgfqpoint{0.691151in}{0.410326in}}%
\pgfpathlineto{\pgfqpoint{0.690257in}{0.408099in}}%
\pgfpathlineto{\pgfqpoint{0.691712in}{0.406391in}}%
\pgfpathlineto{\pgfqpoint{0.693589in}{0.407652in}}%
\pgfpathlineto{\pgfqpoint{0.692837in}{0.410463in}}%
\pgfpathlineto{\pgfqpoint{0.691256in}{0.412047in}}%
\pgfpathlineto{\pgfqpoint{0.692624in}{0.413404in}}%
\pgfpathlineto{\pgfqpoint{0.696548in}{0.415075in}}%
\pgfpathlineto{\pgfqpoint{0.699114in}{0.416794in}}%
\pgfpathlineto{\pgfqpoint{0.705119in}{0.417836in}}%
\pgfpathlineto{\pgfqpoint{0.708125in}{0.417265in}}%
\pgfpathlineto{\pgfqpoint{0.711215in}{0.417647in}}%
\pgfpathlineto{\pgfqpoint{0.714242in}{0.417166in}}%
\pgfpathlineto{\pgfqpoint{0.721904in}{0.414573in}}%
\pgfpathlineto{\pgfqpoint{0.721574in}{0.415801in}}%
\pgfpathlineto{\pgfqpoint{0.725765in}{0.415354in}}%
\pgfpathlineto{\pgfqpoint{0.727666in}{0.413737in}}%
\pgfpathlineto{\pgfqpoint{0.725147in}{0.414065in}}%
\pgfpathlineto{\pgfqpoint{0.723593in}{0.412237in}}%
\pgfpathlineto{\pgfqpoint{0.720711in}{0.414781in}}%
\pgfpathlineto{\pgfqpoint{0.719039in}{0.412825in}}%
\pgfpathlineto{\pgfqpoint{0.715284in}{0.416046in}}%
\pgfpathlineto{\pgfqpoint{0.713680in}{0.414200in}}%
\pgfpathlineto{\pgfqpoint{0.710015in}{0.417271in}}%
\pgfpathlineto{\pgfqpoint{0.706826in}{0.413694in}}%
\pgfpathlineto{\pgfqpoint{0.707923in}{0.412710in}}%
\pgfpathlineto{\pgfqpoint{0.701711in}{0.405421in}}%
\pgfpathlineto{\pgfqpoint{0.699223in}{0.404064in}}%
\pgfpathlineto{\pgfqpoint{0.702855in}{0.400929in}}%
\pgfpathlineto{\pgfqpoint{0.698256in}{0.395576in}}%
\pgfpathclose%
\pgfusepath{fill}%
\end{pgfscope}%
\begin{pgfscope}%
\pgfpathrectangle{\pgfqpoint{0.100000in}{0.100000in}}{\pgfqpoint{3.608454in}{2.310000in}}%
\pgfusepath{clip}%
\pgfsetbuttcap%
\pgfsetmiterjoin%
\definecolor{currentfill}{rgb}{0.000000,0.478431,0.760784}%
\pgfsetfillcolor{currentfill}%
\pgfsetlinewidth{0.000000pt}%
\definecolor{currentstroke}{rgb}{0.000000,0.000000,0.000000}%
\pgfsetstrokecolor{currentstroke}%
\pgfsetstrokeopacity{0.000000}%
\pgfsetdash{}{0pt}%
\pgfpathmoveto{\pgfqpoint{1.886398in}{1.921692in}}%
\pgfpathlineto{\pgfqpoint{1.885060in}{1.894397in}}%
\pgfpathlineto{\pgfqpoint{1.845659in}{1.896366in}}%
\pgfpathlineto{\pgfqpoint{1.802493in}{1.899008in}}%
\pgfpathlineto{\pgfqpoint{1.796957in}{1.908802in}}%
\pgfpathlineto{\pgfqpoint{1.796911in}{1.916921in}}%
\pgfpathlineto{\pgfqpoint{1.800710in}{1.922898in}}%
\pgfpathlineto{\pgfqpoint{1.801655in}{1.932573in}}%
\pgfpathlineto{\pgfqpoint{1.799830in}{1.937828in}}%
\pgfpathlineto{\pgfqpoint{1.802847in}{1.946100in}}%
\pgfpathlineto{\pgfqpoint{1.801850in}{1.950975in}}%
\pgfpathlineto{\pgfqpoint{1.796923in}{1.954510in}}%
\pgfpathlineto{\pgfqpoint{1.828664in}{1.952547in}}%
\pgfpathlineto{\pgfqpoint{1.863197in}{1.950599in}}%
\pgfpathlineto{\pgfqpoint{1.885796in}{1.949443in}}%
\pgfpathlineto{\pgfqpoint{1.884519in}{1.921788in}}%
\pgfpathlineto{\pgfqpoint{1.886398in}{1.921692in}}%
\pgfpathclose%
\pgfusepath{fill}%
\end{pgfscope}%
\begin{pgfscope}%
\pgfpathrectangle{\pgfqpoint{0.100000in}{0.100000in}}{\pgfqpoint{3.608454in}{2.310000in}}%
\pgfusepath{clip}%
\pgfsetbuttcap%
\pgfsetmiterjoin%
\definecolor{currentfill}{rgb}{0.000000,0.356863,0.821569}%
\pgfsetfillcolor{currentfill}%
\pgfsetlinewidth{0.000000pt}%
\definecolor{currentstroke}{rgb}{0.000000,0.000000,0.000000}%
\pgfsetstrokecolor{currentstroke}%
\pgfsetstrokeopacity{0.000000}%
\pgfsetdash{}{0pt}%
\pgfpathmoveto{\pgfqpoint{1.790983in}{1.190359in}}%
\pgfpathlineto{\pgfqpoint{1.788727in}{1.150707in}}%
\pgfpathlineto{\pgfqpoint{1.728388in}{1.154348in}}%
\pgfpathlineto{\pgfqpoint{1.731511in}{1.193664in}}%
\pgfpathlineto{\pgfqpoint{1.723873in}{1.194151in}}%
\pgfpathlineto{\pgfqpoint{1.725856in}{1.224921in}}%
\pgfpathlineto{\pgfqpoint{1.751890in}{1.223219in}}%
\pgfpathlineto{\pgfqpoint{1.752303in}{1.230115in}}%
\pgfpathlineto{\pgfqpoint{1.786488in}{1.228084in}}%
\pgfpathlineto{\pgfqpoint{1.787082in}{1.221084in}}%
\pgfpathlineto{\pgfqpoint{1.785504in}{1.190712in}}%
\pgfpathlineto{\pgfqpoint{1.790983in}{1.190359in}}%
\pgfpathclose%
\pgfusepath{fill}%
\end{pgfscope}%
\begin{pgfscope}%
\pgfpathrectangle{\pgfqpoint{0.100000in}{0.100000in}}{\pgfqpoint{3.608454in}{2.310000in}}%
\pgfusepath{clip}%
\pgfsetbuttcap%
\pgfsetmiterjoin%
\definecolor{currentfill}{rgb}{0.000000,0.525490,0.737255}%
\pgfsetfillcolor{currentfill}%
\pgfsetlinewidth{0.000000pt}%
\definecolor{currentstroke}{rgb}{0.000000,0.000000,0.000000}%
\pgfsetstrokecolor{currentstroke}%
\pgfsetstrokeopacity{0.000000}%
\pgfsetdash{}{0pt}%
\pgfpathmoveto{\pgfqpoint{2.084392in}{1.676317in}}%
\pgfpathlineto{\pgfqpoint{2.057034in}{1.676704in}}%
\pgfpathlineto{\pgfqpoint{2.057273in}{1.695912in}}%
\pgfpathlineto{\pgfqpoint{2.046244in}{1.696054in}}%
\pgfpathlineto{\pgfqpoint{2.046612in}{1.723715in}}%
\pgfpathlineto{\pgfqpoint{2.046268in}{1.751283in}}%
\pgfpathlineto{\pgfqpoint{2.073025in}{1.750974in}}%
\pgfpathlineto{\pgfqpoint{2.100603in}{1.750815in}}%
\pgfpathlineto{\pgfqpoint{2.101083in}{1.743915in}}%
\pgfpathlineto{\pgfqpoint{2.114776in}{1.743905in}}%
\pgfpathlineto{\pgfqpoint{2.115120in}{1.723283in}}%
\pgfpathlineto{\pgfqpoint{2.115164in}{1.695734in}}%
\pgfpathlineto{\pgfqpoint{2.111687in}{1.695719in}}%
\pgfpathlineto{\pgfqpoint{2.084490in}{1.695746in}}%
\pgfpathlineto{\pgfqpoint{2.084392in}{1.676317in}}%
\pgfpathclose%
\pgfusepath{fill}%
\end{pgfscope}%
\begin{pgfscope}%
\pgfpathrectangle{\pgfqpoint{0.100000in}{0.100000in}}{\pgfqpoint{3.608454in}{2.310000in}}%
\pgfusepath{clip}%
\pgfsetbuttcap%
\pgfsetmiterjoin%
\definecolor{currentfill}{rgb}{0.000000,0.427451,0.786275}%
\pgfsetfillcolor{currentfill}%
\pgfsetlinewidth{0.000000pt}%
\definecolor{currentstroke}{rgb}{0.000000,0.000000,0.000000}%
\pgfsetstrokecolor{currentstroke}%
\pgfsetstrokeopacity{0.000000}%
\pgfsetdash{}{0pt}%
\pgfpathmoveto{\pgfqpoint{2.138550in}{1.463912in}}%
\pgfpathlineto{\pgfqpoint{2.128892in}{1.463893in}}%
\pgfpathlineto{\pgfqpoint{2.060821in}{1.465149in}}%
\pgfpathlineto{\pgfqpoint{2.059082in}{1.471267in}}%
\pgfpathlineto{\pgfqpoint{2.053550in}{1.477162in}}%
\pgfpathlineto{\pgfqpoint{2.056854in}{1.480886in}}%
\pgfpathlineto{\pgfqpoint{2.058174in}{1.490199in}}%
\pgfpathlineto{\pgfqpoint{2.083784in}{1.490019in}}%
\pgfpathlineto{\pgfqpoint{2.083971in}{1.510379in}}%
\pgfpathlineto{\pgfqpoint{2.097514in}{1.510367in}}%
\pgfpathlineto{\pgfqpoint{2.124622in}{1.510289in}}%
\pgfpathlineto{\pgfqpoint{2.138316in}{1.510307in}}%
\pgfpathlineto{\pgfqpoint{2.138550in}{1.463912in}}%
\pgfpathclose%
\pgfusepath{fill}%
\end{pgfscope}%
\begin{pgfscope}%
\pgfpathrectangle{\pgfqpoint{0.100000in}{0.100000in}}{\pgfqpoint{3.608454in}{2.310000in}}%
\pgfusepath{clip}%
\pgfsetbuttcap%
\pgfsetmiterjoin%
\definecolor{currentfill}{rgb}{0.000000,0.556863,0.721569}%
\pgfsetfillcolor{currentfill}%
\pgfsetlinewidth{0.000000pt}%
\definecolor{currentstroke}{rgb}{0.000000,0.000000,0.000000}%
\pgfsetstrokecolor{currentstroke}%
\pgfsetstrokeopacity{0.000000}%
\pgfsetdash{}{0pt}%
\pgfpathmoveto{\pgfqpoint{1.941489in}{0.354044in}}%
\pgfpathlineto{\pgfqpoint{1.939142in}{0.354101in}}%
\pgfpathlineto{\pgfqpoint{1.934663in}{0.372389in}}%
\pgfpathlineto{\pgfqpoint{1.935228in}{0.399177in}}%
\pgfpathlineto{\pgfqpoint{1.938012in}{0.416549in}}%
\pgfpathlineto{\pgfqpoint{1.956579in}{0.451354in}}%
\pgfpathlineto{\pgfqpoint{1.961353in}{0.455090in}}%
\pgfpathlineto{\pgfqpoint{1.966578in}{0.467186in}}%
\pgfpathlineto{\pgfqpoint{1.973872in}{0.476940in}}%
\pgfpathlineto{\pgfqpoint{1.976992in}{0.475338in}}%
\pgfpathlineto{\pgfqpoint{1.975277in}{0.471174in}}%
\pgfpathlineto{\pgfqpoint{1.962738in}{0.456890in}}%
\pgfpathlineto{\pgfqpoint{1.952264in}{0.440704in}}%
\pgfpathlineto{\pgfqpoint{1.945455in}{0.427440in}}%
\pgfpathlineto{\pgfqpoint{1.940167in}{0.414451in}}%
\pgfpathlineto{\pgfqpoint{1.937055in}{0.402500in}}%
\pgfpathlineto{\pgfqpoint{1.935766in}{0.383962in}}%
\pgfpathlineto{\pgfqpoint{1.938010in}{0.367087in}}%
\pgfpathlineto{\pgfqpoint{1.941489in}{0.354044in}}%
\pgfpathclose%
\pgfusepath{fill}%
\end{pgfscope}%
\begin{pgfscope}%
\pgfpathrectangle{\pgfqpoint{0.100000in}{0.100000in}}{\pgfqpoint{3.608454in}{2.310000in}}%
\pgfusepath{clip}%
\pgfsetbuttcap%
\pgfsetmiterjoin%
\definecolor{currentfill}{rgb}{0.000000,0.556863,0.721569}%
\pgfsetfillcolor{currentfill}%
\pgfsetlinewidth{0.000000pt}%
\definecolor{currentstroke}{rgb}{0.000000,0.000000,0.000000}%
\pgfsetstrokecolor{currentstroke}%
\pgfsetstrokeopacity{0.000000}%
\pgfsetdash{}{0pt}%
\pgfpathmoveto{\pgfqpoint{1.930477in}{0.354309in}}%
\pgfpathlineto{\pgfqpoint{1.899944in}{0.355045in}}%
\pgfpathlineto{\pgfqpoint{1.891060in}{0.356772in}}%
\pgfpathlineto{\pgfqpoint{1.891484in}{0.370110in}}%
\pgfpathlineto{\pgfqpoint{1.859982in}{0.371433in}}%
\pgfpathlineto{\pgfqpoint{1.861148in}{0.393352in}}%
\pgfpathlineto{\pgfqpoint{1.857651in}{0.393438in}}%
\pgfpathlineto{\pgfqpoint{1.856313in}{0.408047in}}%
\pgfpathlineto{\pgfqpoint{1.852427in}{0.418217in}}%
\pgfpathlineto{\pgfqpoint{1.834828in}{0.418520in}}%
\pgfpathlineto{\pgfqpoint{1.836840in}{0.475148in}}%
\pgfpathlineto{\pgfqpoint{1.902126in}{0.472699in}}%
\pgfpathlineto{\pgfqpoint{1.900788in}{0.477356in}}%
\pgfpathlineto{\pgfqpoint{1.907060in}{0.482191in}}%
\pgfpathlineto{\pgfqpoint{1.924683in}{0.477880in}}%
\pgfpathlineto{\pgfqpoint{1.938836in}{0.498314in}}%
\pgfpathlineto{\pgfqpoint{1.954355in}{0.511185in}}%
\pgfpathlineto{\pgfqpoint{1.962232in}{0.510712in}}%
\pgfpathlineto{\pgfqpoint{1.966594in}{0.506597in}}%
\pgfpathlineto{\pgfqpoint{1.973994in}{0.506774in}}%
\pgfpathlineto{\pgfqpoint{1.981043in}{0.498344in}}%
\pgfpathlineto{\pgfqpoint{1.979461in}{0.488391in}}%
\pgfpathlineto{\pgfqpoint{1.978605in}{0.483866in}}%
\pgfpathlineto{\pgfqpoint{1.969758in}{0.476284in}}%
\pgfpathlineto{\pgfqpoint{1.965121in}{0.476210in}}%
\pgfpathlineto{\pgfqpoint{1.961299in}{0.468390in}}%
\pgfpathlineto{\pgfqpoint{1.951003in}{0.452565in}}%
\pgfpathlineto{\pgfqpoint{1.945104in}{0.457357in}}%
\pgfpathlineto{\pgfqpoint{1.938160in}{0.454348in}}%
\pgfpathlineto{\pgfqpoint{1.938434in}{0.445675in}}%
\pgfpathlineto{\pgfqpoint{1.946055in}{0.442493in}}%
\pgfpathlineto{\pgfqpoint{1.941015in}{0.432395in}}%
\pgfpathlineto{\pgfqpoint{1.934658in}{0.413330in}}%
\pgfpathlineto{\pgfqpoint{1.931292in}{0.396174in}}%
\pgfpathlineto{\pgfqpoint{1.923085in}{0.385830in}}%
\pgfpathlineto{\pgfqpoint{1.922580in}{0.373704in}}%
\pgfpathlineto{\pgfqpoint{1.928475in}{0.367207in}}%
\pgfpathlineto{\pgfqpoint{1.930477in}{0.354309in}}%
\pgfpathclose%
\pgfusepath{fill}%
\end{pgfscope}%
\begin{pgfscope}%
\pgfpathrectangle{\pgfqpoint{0.100000in}{0.100000in}}{\pgfqpoint{3.608454in}{2.310000in}}%
\pgfusepath{clip}%
\pgfsetbuttcap%
\pgfsetmiterjoin%
\definecolor{currentfill}{rgb}{0.000000,0.470588,0.764706}%
\pgfsetfillcolor{currentfill}%
\pgfsetlinewidth{0.000000pt}%
\definecolor{currentstroke}{rgb}{0.000000,0.000000,0.000000}%
\pgfsetstrokecolor{currentstroke}%
\pgfsetstrokeopacity{0.000000}%
\pgfsetdash{}{0pt}%
\pgfpathmoveto{\pgfqpoint{2.023670in}{1.724067in}}%
\pgfpathlineto{\pgfqpoint{1.998767in}{1.724416in}}%
\pgfpathlineto{\pgfqpoint{1.999302in}{1.751993in}}%
\pgfpathlineto{\pgfqpoint{1.944686in}{1.753364in}}%
\pgfpathlineto{\pgfqpoint{1.944449in}{1.753368in}}%
\pgfpathlineto{\pgfqpoint{1.945343in}{1.781068in}}%
\pgfpathlineto{\pgfqpoint{1.986361in}{1.779846in}}%
\pgfpathlineto{\pgfqpoint{2.024621in}{1.779133in}}%
\pgfpathlineto{\pgfqpoint{2.023670in}{1.724067in}}%
\pgfpathclose%
\pgfusepath{fill}%
\end{pgfscope}%
\begin{pgfscope}%
\pgfpathrectangle{\pgfqpoint{0.100000in}{0.100000in}}{\pgfqpoint{3.608454in}{2.310000in}}%
\pgfusepath{clip}%
\pgfsetbuttcap%
\pgfsetmiterjoin%
\definecolor{currentfill}{rgb}{0.000000,0.862745,0.568627}%
\pgfsetfillcolor{currentfill}%
\pgfsetlinewidth{0.000000pt}%
\definecolor{currentstroke}{rgb}{0.000000,0.000000,0.000000}%
\pgfsetstrokecolor{currentstroke}%
\pgfsetstrokeopacity{0.000000}%
\pgfsetdash{}{0pt}%
\pgfpathmoveto{\pgfqpoint{2.313347in}{1.069167in}}%
\pgfpathlineto{\pgfqpoint{2.313444in}{1.062948in}}%
\pgfpathlineto{\pgfqpoint{2.310094in}{1.055119in}}%
\pgfpathlineto{\pgfqpoint{2.293330in}{1.054914in}}%
\pgfpathlineto{\pgfqpoint{2.269431in}{1.054652in}}%
\pgfpathlineto{\pgfqpoint{2.269359in}{1.061552in}}%
\pgfpathlineto{\pgfqpoint{2.245419in}{1.061527in}}%
\pgfpathlineto{\pgfqpoint{2.247654in}{1.068396in}}%
\pgfpathlineto{\pgfqpoint{2.247558in}{1.082128in}}%
\pgfpathlineto{\pgfqpoint{2.259040in}{1.082097in}}%
\pgfpathlineto{\pgfqpoint{2.258979in}{1.087834in}}%
\pgfpathlineto{\pgfqpoint{2.272664in}{1.088030in}}%
\pgfpathlineto{\pgfqpoint{2.272396in}{1.102853in}}%
\pgfpathlineto{\pgfqpoint{2.279258in}{1.102907in}}%
\pgfpathlineto{\pgfqpoint{2.279199in}{1.109799in}}%
\pgfpathlineto{\pgfqpoint{2.286066in}{1.109865in}}%
\pgfpathlineto{\pgfqpoint{2.286012in}{1.115699in}}%
\pgfpathlineto{\pgfqpoint{2.295109in}{1.108191in}}%
\pgfpathlineto{\pgfqpoint{2.293303in}{1.102454in}}%
\pgfpathlineto{\pgfqpoint{2.302022in}{1.100303in}}%
\pgfpathlineto{\pgfqpoint{2.308868in}{1.093507in}}%
\pgfpathlineto{\pgfqpoint{2.307635in}{1.090812in}}%
\pgfpathlineto{\pgfqpoint{2.313117in}{1.083232in}}%
\pgfpathlineto{\pgfqpoint{2.313347in}{1.069167in}}%
\pgfpathclose%
\pgfusepath{fill}%
\end{pgfscope}%
\begin{pgfscope}%
\pgfpathrectangle{\pgfqpoint{0.100000in}{0.100000in}}{\pgfqpoint{3.608454in}{2.310000in}}%
\pgfusepath{clip}%
\pgfsetbuttcap%
\pgfsetmiterjoin%
\definecolor{currentfill}{rgb}{0.000000,0.568627,0.715686}%
\pgfsetfillcolor{currentfill}%
\pgfsetlinewidth{0.000000pt}%
\definecolor{currentstroke}{rgb}{0.000000,0.000000,0.000000}%
\pgfsetstrokecolor{currentstroke}%
\pgfsetstrokeopacity{0.000000}%
\pgfsetdash{}{0pt}%
\pgfpathmoveto{\pgfqpoint{2.958909in}{1.001900in}}%
\pgfpathlineto{\pgfqpoint{2.949644in}{0.997974in}}%
\pgfpathlineto{\pgfqpoint{2.953478in}{0.990526in}}%
\pgfpathlineto{\pgfqpoint{2.950161in}{0.981612in}}%
\pgfpathlineto{\pgfqpoint{2.956878in}{0.975335in}}%
\pgfpathlineto{\pgfqpoint{2.954883in}{0.971559in}}%
\pgfpathlineto{\pgfqpoint{2.948243in}{0.978295in}}%
\pgfpathlineto{\pgfqpoint{2.941645in}{0.986529in}}%
\pgfpathlineto{\pgfqpoint{2.922696in}{0.996327in}}%
\pgfpathlineto{\pgfqpoint{2.919954in}{1.001789in}}%
\pgfpathlineto{\pgfqpoint{2.910642in}{1.010268in}}%
\pgfpathlineto{\pgfqpoint{2.907921in}{1.017944in}}%
\pgfpathlineto{\pgfqpoint{2.895066in}{1.035008in}}%
\pgfpathlineto{\pgfqpoint{2.888625in}{1.033073in}}%
\pgfpathlineto{\pgfqpoint{2.884093in}{1.034847in}}%
\pgfpathlineto{\pgfqpoint{2.876324in}{1.042276in}}%
\pgfpathlineto{\pgfqpoint{2.872404in}{1.042068in}}%
\pgfpathlineto{\pgfqpoint{2.864794in}{1.047108in}}%
\pgfpathlineto{\pgfqpoint{2.863588in}{1.049271in}}%
\pgfpathlineto{\pgfqpoint{2.864683in}{1.055271in}}%
\pgfpathlineto{\pgfqpoint{2.869918in}{1.063176in}}%
\pgfpathlineto{\pgfqpoint{2.876200in}{1.069013in}}%
\pgfpathlineto{\pgfqpoint{2.876403in}{1.073700in}}%
\pgfpathlineto{\pgfqpoint{2.882548in}{1.076650in}}%
\pgfpathlineto{\pgfqpoint{2.903153in}{1.086612in}}%
\pgfpathlineto{\pgfqpoint{2.922867in}{1.095204in}}%
\pgfpathlineto{\pgfqpoint{2.931616in}{1.096854in}}%
\pgfpathlineto{\pgfqpoint{2.934795in}{1.069673in}}%
\pgfpathlineto{\pgfqpoint{2.951248in}{1.056924in}}%
\pgfpathlineto{\pgfqpoint{2.957922in}{1.053177in}}%
\pgfpathlineto{\pgfqpoint{2.975999in}{1.050234in}}%
\pgfpathlineto{\pgfqpoint{2.968957in}{1.035717in}}%
\pgfpathlineto{\pgfqpoint{2.962380in}{1.029555in}}%
\pgfpathlineto{\pgfqpoint{2.961263in}{1.022919in}}%
\pgfpathlineto{\pgfqpoint{2.963991in}{1.017466in}}%
\pgfpathlineto{\pgfqpoint{2.958909in}{1.001900in}}%
\pgfpathclose%
\pgfusepath{fill}%
\end{pgfscope}%
\begin{pgfscope}%
\pgfpathrectangle{\pgfqpoint{0.100000in}{0.100000in}}{\pgfqpoint{3.608454in}{2.310000in}}%
\pgfusepath{clip}%
\pgfsetbuttcap%
\pgfsetmiterjoin%
\definecolor{currentfill}{rgb}{0.000000,0.407843,0.796078}%
\pgfsetfillcolor{currentfill}%
\pgfsetlinewidth{0.000000pt}%
\definecolor{currentstroke}{rgb}{0.000000,0.000000,0.000000}%
\pgfsetstrokecolor{currentstroke}%
\pgfsetstrokeopacity{0.000000}%
\pgfsetdash{}{0pt}%
\pgfpathmoveto{\pgfqpoint{2.684273in}{1.275179in}}%
\pgfpathlineto{\pgfqpoint{2.678952in}{1.278333in}}%
\pgfpathlineto{\pgfqpoint{2.675593in}{1.285078in}}%
\pgfpathlineto{\pgfqpoint{2.669652in}{1.289191in}}%
\pgfpathlineto{\pgfqpoint{2.655765in}{1.288068in}}%
\pgfpathlineto{\pgfqpoint{2.649917in}{1.290900in}}%
\pgfpathlineto{\pgfqpoint{2.648711in}{1.297118in}}%
\pgfpathlineto{\pgfqpoint{2.643610in}{1.301513in}}%
\pgfpathlineto{\pgfqpoint{2.640611in}{1.294225in}}%
\pgfpathlineto{\pgfqpoint{2.636960in}{1.295013in}}%
\pgfpathlineto{\pgfqpoint{2.636434in}{1.301852in}}%
\pgfpathlineto{\pgfqpoint{2.629550in}{1.301264in}}%
\pgfpathlineto{\pgfqpoint{2.629213in}{1.305901in}}%
\pgfpathlineto{\pgfqpoint{2.622478in}{1.305021in}}%
\pgfpathlineto{\pgfqpoint{2.621408in}{1.315375in}}%
\pgfpathlineto{\pgfqpoint{2.644392in}{1.317477in}}%
\pgfpathlineto{\pgfqpoint{2.641760in}{1.344149in}}%
\pgfpathlineto{\pgfqpoint{2.643662in}{1.346766in}}%
\pgfpathlineto{\pgfqpoint{2.662944in}{1.348688in}}%
\pgfpathlineto{\pgfqpoint{2.667619in}{1.346789in}}%
\pgfpathlineto{\pgfqpoint{2.669050in}{1.334303in}}%
\pgfpathlineto{\pgfqpoint{2.674475in}{1.337189in}}%
\pgfpathlineto{\pgfqpoint{2.688201in}{1.338715in}}%
\pgfpathlineto{\pgfqpoint{2.697081in}{1.338065in}}%
\pgfpathlineto{\pgfqpoint{2.697292in}{1.333124in}}%
\pgfpathlineto{\pgfqpoint{2.707002in}{1.331992in}}%
\pgfpathlineto{\pgfqpoint{2.713052in}{1.339727in}}%
\pgfpathlineto{\pgfqpoint{2.718732in}{1.341235in}}%
\pgfpathlineto{\pgfqpoint{2.722897in}{1.338929in}}%
\pgfpathlineto{\pgfqpoint{2.724859in}{1.333241in}}%
\pgfpathlineto{\pgfqpoint{2.730992in}{1.324932in}}%
\pgfpathlineto{\pgfqpoint{2.726160in}{1.316456in}}%
\pgfpathlineto{\pgfqpoint{2.726400in}{1.307715in}}%
\pgfpathlineto{\pgfqpoint{2.724571in}{1.300390in}}%
\pgfpathlineto{\pgfqpoint{2.718323in}{1.291481in}}%
\pgfpathlineto{\pgfqpoint{2.704197in}{1.287428in}}%
\pgfpathlineto{\pgfqpoint{2.698255in}{1.290856in}}%
\pgfpathlineto{\pgfqpoint{2.690035in}{1.277558in}}%
\pgfpathlineto{\pgfqpoint{2.684273in}{1.275179in}}%
\pgfpathclose%
\pgfusepath{fill}%
\end{pgfscope}%
\begin{pgfscope}%
\pgfpathrectangle{\pgfqpoint{0.100000in}{0.100000in}}{\pgfqpoint{3.608454in}{2.310000in}}%
\pgfusepath{clip}%
\pgfsetbuttcap%
\pgfsetmiterjoin%
\definecolor{currentfill}{rgb}{0.000000,0.368627,0.815686}%
\pgfsetfillcolor{currentfill}%
\pgfsetlinewidth{0.000000pt}%
\definecolor{currentstroke}{rgb}{0.000000,0.000000,0.000000}%
\pgfsetstrokecolor{currentstroke}%
\pgfsetstrokeopacity{0.000000}%
\pgfsetdash{}{0pt}%
\pgfpathmoveto{\pgfqpoint{1.912766in}{1.573654in}}%
\pgfpathlineto{\pgfqpoint{1.912258in}{1.559919in}}%
\pgfpathlineto{\pgfqpoint{1.884919in}{1.560911in}}%
\pgfpathlineto{\pgfqpoint{1.885929in}{1.588436in}}%
\pgfpathlineto{\pgfqpoint{1.858155in}{1.589614in}}%
\pgfpathlineto{\pgfqpoint{1.860373in}{1.661671in}}%
\pgfpathlineto{\pgfqpoint{1.904126in}{1.659838in}}%
\pgfpathlineto{\pgfqpoint{1.905777in}{1.655731in}}%
\pgfpathlineto{\pgfqpoint{1.914725in}{1.650247in}}%
\pgfpathlineto{\pgfqpoint{1.914002in}{1.614954in}}%
\pgfpathlineto{\pgfqpoint{1.912766in}{1.573654in}}%
\pgfpathclose%
\pgfusepath{fill}%
\end{pgfscope}%
\begin{pgfscope}%
\pgfpathrectangle{\pgfqpoint{0.100000in}{0.100000in}}{\pgfqpoint{3.608454in}{2.310000in}}%
\pgfusepath{clip}%
\pgfsetbuttcap%
\pgfsetmiterjoin%
\definecolor{currentfill}{rgb}{0.000000,0.458824,0.770588}%
\pgfsetfillcolor{currentfill}%
\pgfsetlinewidth{0.000000pt}%
\definecolor{currentstroke}{rgb}{0.000000,0.000000,0.000000}%
\pgfsetstrokecolor{currentstroke}%
\pgfsetstrokeopacity{0.000000}%
\pgfsetdash{}{0pt}%
\pgfpathmoveto{\pgfqpoint{2.719775in}{1.481004in}}%
\pgfpathlineto{\pgfqpoint{2.720165in}{1.477649in}}%
\pgfpathlineto{\pgfqpoint{2.695336in}{1.474638in}}%
\pgfpathlineto{\pgfqpoint{2.694173in}{1.495218in}}%
\pgfpathlineto{\pgfqpoint{2.679497in}{1.493619in}}%
\pgfpathlineto{\pgfqpoint{2.678668in}{1.500441in}}%
\pgfpathlineto{\pgfqpoint{2.667350in}{1.499209in}}%
\pgfpathlineto{\pgfqpoint{2.664197in}{1.526649in}}%
\pgfpathlineto{\pgfqpoint{2.661863in}{1.526375in}}%
\pgfpathlineto{\pgfqpoint{2.660236in}{1.540245in}}%
\pgfpathlineto{\pgfqpoint{2.661344in}{1.549614in}}%
\pgfpathlineto{\pgfqpoint{2.668362in}{1.548067in}}%
\pgfpathlineto{\pgfqpoint{2.688778in}{1.550091in}}%
\pgfpathlineto{\pgfqpoint{2.684290in}{1.589030in}}%
\pgfpathlineto{\pgfqpoint{2.706079in}{1.591455in}}%
\pgfpathlineto{\pgfqpoint{2.707766in}{1.586556in}}%
\pgfpathlineto{\pgfqpoint{2.714189in}{1.531011in}}%
\pgfpathlineto{\pgfqpoint{2.719775in}{1.481004in}}%
\pgfpathclose%
\pgfusepath{fill}%
\end{pgfscope}%
\begin{pgfscope}%
\pgfpathrectangle{\pgfqpoint{0.100000in}{0.100000in}}{\pgfqpoint{3.608454in}{2.310000in}}%
\pgfusepath{clip}%
\pgfsetbuttcap%
\pgfsetmiterjoin%
\definecolor{currentfill}{rgb}{0.000000,1.000000,0.500000}%
\pgfsetfillcolor{currentfill}%
\pgfsetlinewidth{0.000000pt}%
\definecolor{currentstroke}{rgb}{0.000000,0.000000,0.000000}%
\pgfsetstrokecolor{currentstroke}%
\pgfsetstrokeopacity{0.000000}%
\pgfsetdash{}{0pt}%
\pgfpathmoveto{\pgfqpoint{2.837048in}{1.201444in}}%
\pgfpathlineto{\pgfqpoint{2.833845in}{1.206501in}}%
\pgfpathlineto{\pgfqpoint{2.833882in}{1.215723in}}%
\pgfpathlineto{\pgfqpoint{2.831285in}{1.222519in}}%
\pgfpathlineto{\pgfqpoint{2.826345in}{1.240166in}}%
\pgfpathlineto{\pgfqpoint{2.826981in}{1.247494in}}%
\pgfpathlineto{\pgfqpoint{2.825042in}{1.252479in}}%
\pgfpathlineto{\pgfqpoint{2.830889in}{1.256941in}}%
\pgfpathlineto{\pgfqpoint{2.833788in}{1.254755in}}%
\pgfpathlineto{\pgfqpoint{2.841576in}{1.260538in}}%
\pgfpathlineto{\pgfqpoint{2.850587in}{1.262241in}}%
\pgfpathlineto{\pgfqpoint{2.851560in}{1.271717in}}%
\pgfpathlineto{\pgfqpoint{2.860524in}{1.271364in}}%
\pgfpathlineto{\pgfqpoint{2.872542in}{1.266952in}}%
\pgfpathlineto{\pgfqpoint{2.874724in}{1.257922in}}%
\pgfpathlineto{\pgfqpoint{2.880671in}{1.252069in}}%
\pgfpathlineto{\pgfqpoint{2.887645in}{1.250541in}}%
\pgfpathlineto{\pgfqpoint{2.878670in}{1.243251in}}%
\pgfpathlineto{\pgfqpoint{2.878218in}{1.237030in}}%
\pgfpathlineto{\pgfqpoint{2.870950in}{1.230065in}}%
\pgfpathlineto{\pgfqpoint{2.871264in}{1.223809in}}%
\pgfpathlineto{\pgfqpoint{2.859365in}{1.219154in}}%
\pgfpathlineto{\pgfqpoint{2.856686in}{1.210206in}}%
\pgfpathlineto{\pgfqpoint{2.837048in}{1.201444in}}%
\pgfpathclose%
\pgfusepath{fill}%
\end{pgfscope}%
\begin{pgfscope}%
\pgfpathrectangle{\pgfqpoint{0.100000in}{0.100000in}}{\pgfqpoint{3.608454in}{2.310000in}}%
\pgfusepath{clip}%
\pgfsetbuttcap%
\pgfsetmiterjoin%
\definecolor{currentfill}{rgb}{0.000000,0.564706,0.717647}%
\pgfsetfillcolor{currentfill}%
\pgfsetlinewidth{0.000000pt}%
\definecolor{currentstroke}{rgb}{0.000000,0.000000,0.000000}%
\pgfsetstrokecolor{currentstroke}%
\pgfsetstrokeopacity{0.000000}%
\pgfsetdash{}{0pt}%
\pgfpathmoveto{\pgfqpoint{2.989806in}{1.055375in}}%
\pgfpathlineto{\pgfqpoint{2.990720in}{1.049995in}}%
\pgfpathlineto{\pgfqpoint{2.982517in}{1.046654in}}%
\pgfpathlineto{\pgfqpoint{2.975999in}{1.050234in}}%
\pgfpathlineto{\pgfqpoint{2.957922in}{1.053177in}}%
\pgfpathlineto{\pgfqpoint{2.951248in}{1.056924in}}%
\pgfpathlineto{\pgfqpoint{2.934795in}{1.069673in}}%
\pgfpathlineto{\pgfqpoint{2.931616in}{1.096854in}}%
\pgfpathlineto{\pgfqpoint{2.922867in}{1.095204in}}%
\pgfpathlineto{\pgfqpoint{2.922309in}{1.102781in}}%
\pgfpathlineto{\pgfqpoint{2.926537in}{1.111891in}}%
\pgfpathlineto{\pgfqpoint{2.932868in}{1.114078in}}%
\pgfpathlineto{\pgfqpoint{2.946773in}{1.103937in}}%
\pgfpathlineto{\pgfqpoint{2.947537in}{1.098430in}}%
\pgfpathlineto{\pgfqpoint{2.986399in}{1.102377in}}%
\pgfpathlineto{\pgfqpoint{2.985051in}{1.092680in}}%
\pgfpathlineto{\pgfqpoint{2.980220in}{1.091029in}}%
\pgfpathlineto{\pgfqpoint{2.984590in}{1.076632in}}%
\pgfpathlineto{\pgfqpoint{2.986559in}{1.063157in}}%
\pgfpathlineto{\pgfqpoint{2.989806in}{1.055375in}}%
\pgfpathclose%
\pgfusepath{fill}%
\end{pgfscope}%
\begin{pgfscope}%
\pgfpathrectangle{\pgfqpoint{0.100000in}{0.100000in}}{\pgfqpoint{3.608454in}{2.310000in}}%
\pgfusepath{clip}%
\pgfsetbuttcap%
\pgfsetmiterjoin%
\definecolor{currentfill}{rgb}{0.000000,0.827451,0.586275}%
\pgfsetfillcolor{currentfill}%
\pgfsetlinewidth{0.000000pt}%
\definecolor{currentstroke}{rgb}{0.000000,0.000000,0.000000}%
\pgfsetstrokecolor{currentstroke}%
\pgfsetstrokeopacity{0.000000}%
\pgfsetdash{}{0pt}%
\pgfpathmoveto{\pgfqpoint{1.037345in}{0.183729in}}%
\pgfpathlineto{\pgfqpoint{1.037725in}{0.186685in}}%
\pgfpathlineto{\pgfqpoint{1.039756in}{0.184678in}}%
\pgfpathlineto{\pgfqpoint{1.037345in}{0.183729in}}%
\pgfpathclose%
\pgfusepath{fill}%
\end{pgfscope}%
\begin{pgfscope}%
\pgfpathrectangle{\pgfqpoint{0.100000in}{0.100000in}}{\pgfqpoint{3.608454in}{2.310000in}}%
\pgfusepath{clip}%
\pgfsetbuttcap%
\pgfsetmiterjoin%
\definecolor{currentfill}{rgb}{0.000000,0.827451,0.586275}%
\pgfsetfillcolor{currentfill}%
\pgfsetlinewidth{0.000000pt}%
\definecolor{currentstroke}{rgb}{0.000000,0.000000,0.000000}%
\pgfsetstrokecolor{currentstroke}%
\pgfsetstrokeopacity{0.000000}%
\pgfsetdash{}{0pt}%
\pgfpathmoveto{\pgfqpoint{1.035979in}{0.195213in}}%
\pgfpathlineto{\pgfqpoint{1.035168in}{0.197283in}}%
\pgfpathlineto{\pgfqpoint{1.034662in}{0.206903in}}%
\pgfpathlineto{\pgfqpoint{1.035313in}{0.207991in}}%
\pgfpathlineto{\pgfqpoint{1.034954in}{0.212110in}}%
\pgfpathlineto{\pgfqpoint{1.035841in}{0.213389in}}%
\pgfpathlineto{\pgfqpoint{1.034185in}{0.215042in}}%
\pgfpathlineto{\pgfqpoint{1.034010in}{0.217336in}}%
\pgfpathlineto{\pgfqpoint{1.037398in}{0.218597in}}%
\pgfpathlineto{\pgfqpoint{1.038802in}{0.221050in}}%
\pgfpathlineto{\pgfqpoint{1.037009in}{0.222070in}}%
\pgfpathlineto{\pgfqpoint{1.036796in}{0.223966in}}%
\pgfpathlineto{\pgfqpoint{1.038521in}{0.226447in}}%
\pgfpathlineto{\pgfqpoint{1.036445in}{0.231584in}}%
\pgfpathlineto{\pgfqpoint{1.041639in}{0.235717in}}%
\pgfpathlineto{\pgfqpoint{1.043821in}{0.231654in}}%
\pgfpathlineto{\pgfqpoint{1.048322in}{0.228489in}}%
\pgfpathlineto{\pgfqpoint{1.045161in}{0.227298in}}%
\pgfpathlineto{\pgfqpoint{1.046187in}{0.224668in}}%
\pgfpathlineto{\pgfqpoint{1.045919in}{0.221757in}}%
\pgfpathlineto{\pgfqpoint{1.044595in}{0.217845in}}%
\pgfpathlineto{\pgfqpoint{1.044114in}{0.214196in}}%
\pgfpathlineto{\pgfqpoint{1.042540in}{0.210865in}}%
\pgfpathlineto{\pgfqpoint{1.043020in}{0.209624in}}%
\pgfpathlineto{\pgfqpoint{1.041842in}{0.206198in}}%
\pgfpathlineto{\pgfqpoint{1.039406in}{0.205292in}}%
\pgfpathlineto{\pgfqpoint{1.040077in}{0.203620in}}%
\pgfpathlineto{\pgfqpoint{1.037635in}{0.200811in}}%
\pgfpathlineto{\pgfqpoint{1.035979in}{0.195213in}}%
\pgfpathclose%
\pgfusepath{fill}%
\end{pgfscope}%
\begin{pgfscope}%
\pgfpathrectangle{\pgfqpoint{0.100000in}{0.100000in}}{\pgfqpoint{3.608454in}{2.310000in}}%
\pgfusepath{clip}%
\pgfsetbuttcap%
\pgfsetmiterjoin%
\definecolor{currentfill}{rgb}{0.000000,0.827451,0.586275}%
\pgfsetfillcolor{currentfill}%
\pgfsetlinewidth{0.000000pt}%
\definecolor{currentstroke}{rgb}{0.000000,0.000000,0.000000}%
\pgfsetstrokecolor{currentstroke}%
\pgfsetstrokeopacity{0.000000}%
\pgfsetdash{}{0pt}%
\pgfpathmoveto{\pgfqpoint{1.029215in}{0.224020in}}%
\pgfpathlineto{\pgfqpoint{1.032249in}{0.226756in}}%
\pgfpathlineto{\pgfqpoint{1.032356in}{0.230382in}}%
\pgfpathlineto{\pgfqpoint{1.033329in}{0.232191in}}%
\pgfpathlineto{\pgfqpoint{1.035308in}{0.231235in}}%
\pgfpathlineto{\pgfqpoint{1.034682in}{0.226435in}}%
\pgfpathlineto{\pgfqpoint{1.033872in}{0.224494in}}%
\pgfpathlineto{\pgfqpoint{1.031932in}{0.223434in}}%
\pgfpathlineto{\pgfqpoint{1.029215in}{0.224020in}}%
\pgfpathclose%
\pgfusepath{fill}%
\end{pgfscope}%
\begin{pgfscope}%
\pgfpathrectangle{\pgfqpoint{0.100000in}{0.100000in}}{\pgfqpoint{3.608454in}{2.310000in}}%
\pgfusepath{clip}%
\pgfsetbuttcap%
\pgfsetmiterjoin%
\definecolor{currentfill}{rgb}{0.000000,0.827451,0.586275}%
\pgfsetfillcolor{currentfill}%
\pgfsetlinewidth{0.000000pt}%
\definecolor{currentstroke}{rgb}{0.000000,0.000000,0.000000}%
\pgfsetstrokecolor{currentstroke}%
\pgfsetstrokeopacity{0.000000}%
\pgfsetdash{}{0pt}%
\pgfpathmoveto{\pgfqpoint{1.031898in}{0.249479in}}%
\pgfpathlineto{\pgfqpoint{1.037351in}{0.246961in}}%
\pgfpathlineto{\pgfqpoint{1.041099in}{0.250574in}}%
\pgfpathlineto{\pgfqpoint{1.042604in}{0.247304in}}%
\pgfpathlineto{\pgfqpoint{1.044554in}{0.244818in}}%
\pgfpathlineto{\pgfqpoint{1.047084in}{0.239857in}}%
\pgfpathlineto{\pgfqpoint{1.050384in}{0.238572in}}%
\pgfpathlineto{\pgfqpoint{1.051184in}{0.237535in}}%
\pgfpathlineto{\pgfqpoint{1.049176in}{0.229939in}}%
\pgfpathlineto{\pgfqpoint{1.046352in}{0.230420in}}%
\pgfpathlineto{\pgfqpoint{1.040103in}{0.241756in}}%
\pgfpathlineto{\pgfqpoint{1.039585in}{0.239633in}}%
\pgfpathlineto{\pgfqpoint{1.040523in}{0.237630in}}%
\pgfpathlineto{\pgfqpoint{1.039041in}{0.234732in}}%
\pgfpathlineto{\pgfqpoint{1.036002in}{0.232709in}}%
\pgfpathlineto{\pgfqpoint{1.034369in}{0.234138in}}%
\pgfpathlineto{\pgfqpoint{1.033903in}{0.237752in}}%
\pgfpathlineto{\pgfqpoint{1.034739in}{0.241664in}}%
\pgfpathlineto{\pgfqpoint{1.031898in}{0.249479in}}%
\pgfpathclose%
\pgfusepath{fill}%
\end{pgfscope}%
\begin{pgfscope}%
\pgfpathrectangle{\pgfqpoint{0.100000in}{0.100000in}}{\pgfqpoint{3.608454in}{2.310000in}}%
\pgfusepath{clip}%
\pgfsetbuttcap%
\pgfsetmiterjoin%
\definecolor{currentfill}{rgb}{0.000000,0.827451,0.586275}%
\pgfsetfillcolor{currentfill}%
\pgfsetlinewidth{0.000000pt}%
\definecolor{currentstroke}{rgb}{0.000000,0.000000,0.000000}%
\pgfsetstrokecolor{currentstroke}%
\pgfsetstrokeopacity{0.000000}%
\pgfsetdash{}{0pt}%
\pgfpathmoveto{\pgfqpoint{1.062273in}{0.185108in}}%
\pgfpathlineto{\pgfqpoint{1.060831in}{0.189361in}}%
\pgfpathlineto{\pgfqpoint{1.064048in}{0.191263in}}%
\pgfpathlineto{\pgfqpoint{1.066461in}{0.190126in}}%
\pgfpathlineto{\pgfqpoint{1.067412in}{0.187983in}}%
\pgfpathlineto{\pgfqpoint{1.065144in}{0.184845in}}%
\pgfpathlineto{\pgfqpoint{1.062273in}{0.185108in}}%
\pgfpathclose%
\pgfusepath{fill}%
\end{pgfscope}%
\begin{pgfscope}%
\pgfpathrectangle{\pgfqpoint{0.100000in}{0.100000in}}{\pgfqpoint{3.608454in}{2.310000in}}%
\pgfusepath{clip}%
\pgfsetbuttcap%
\pgfsetmiterjoin%
\definecolor{currentfill}{rgb}{0.000000,0.827451,0.586275}%
\pgfsetfillcolor{currentfill}%
\pgfsetlinewidth{0.000000pt}%
\definecolor{currentstroke}{rgb}{0.000000,0.000000,0.000000}%
\pgfsetstrokecolor{currentstroke}%
\pgfsetstrokeopacity{0.000000}%
\pgfsetdash{}{0pt}%
\pgfpathmoveto{\pgfqpoint{1.053060in}{0.199716in}}%
\pgfpathlineto{\pgfqpoint{1.049731in}{0.199089in}}%
\pgfpathlineto{\pgfqpoint{1.049501in}{0.194791in}}%
\pgfpathlineto{\pgfqpoint{1.047203in}{0.194053in}}%
\pgfpathlineto{\pgfqpoint{1.047508in}{0.191183in}}%
\pgfpathlineto{\pgfqpoint{1.045251in}{0.188644in}}%
\pgfpathlineto{\pgfqpoint{1.044390in}{0.191995in}}%
\pgfpathlineto{\pgfqpoint{1.042486in}{0.190633in}}%
\pgfpathlineto{\pgfqpoint{1.041396in}{0.191487in}}%
\pgfpathlineto{\pgfqpoint{1.042996in}{0.194874in}}%
\pgfpathlineto{\pgfqpoint{1.045532in}{0.196881in}}%
\pgfpathlineto{\pgfqpoint{1.047284in}{0.196813in}}%
\pgfpathlineto{\pgfqpoint{1.048561in}{0.199863in}}%
\pgfpathlineto{\pgfqpoint{1.047065in}{0.200403in}}%
\pgfpathlineto{\pgfqpoint{1.046014in}{0.202505in}}%
\pgfpathlineto{\pgfqpoint{1.046696in}{0.204982in}}%
\pgfpathlineto{\pgfqpoint{1.047059in}{0.209727in}}%
\pgfpathlineto{\pgfqpoint{1.050063in}{0.211128in}}%
\pgfpathlineto{\pgfqpoint{1.054012in}{0.204553in}}%
\pgfpathlineto{\pgfqpoint{1.054925in}{0.201896in}}%
\pgfpathlineto{\pgfqpoint{1.053060in}{0.199716in}}%
\pgfpathclose%
\pgfusepath{fill}%
\end{pgfscope}%
\begin{pgfscope}%
\pgfpathrectangle{\pgfqpoint{0.100000in}{0.100000in}}{\pgfqpoint{3.608454in}{2.310000in}}%
\pgfusepath{clip}%
\pgfsetbuttcap%
\pgfsetmiterjoin%
\definecolor{currentfill}{rgb}{0.000000,0.827451,0.586275}%
\pgfsetfillcolor{currentfill}%
\pgfsetlinewidth{0.000000pt}%
\definecolor{currentstroke}{rgb}{0.000000,0.000000,0.000000}%
\pgfsetstrokecolor{currentstroke}%
\pgfsetstrokeopacity{0.000000}%
\pgfsetdash{}{0pt}%
\pgfpathmoveto{\pgfqpoint{1.064056in}{0.213557in}}%
\pgfpathlineto{\pgfqpoint{1.067915in}{0.213771in}}%
\pgfpathlineto{\pgfqpoint{1.069651in}{0.211374in}}%
\pgfpathlineto{\pgfqpoint{1.071517in}{0.212096in}}%
\pgfpathlineto{\pgfqpoint{1.073207in}{0.214262in}}%
\pgfpathlineto{\pgfqpoint{1.079012in}{0.214678in}}%
\pgfpathlineto{\pgfqpoint{1.080487in}{0.211762in}}%
\pgfpathlineto{\pgfqpoint{1.082084in}{0.210879in}}%
\pgfpathlineto{\pgfqpoint{1.082183in}{0.207755in}}%
\pgfpathlineto{\pgfqpoint{1.079181in}{0.204297in}}%
\pgfpathlineto{\pgfqpoint{1.082604in}{0.201601in}}%
\pgfpathlineto{\pgfqpoint{1.080183in}{0.196814in}}%
\pgfpathlineto{\pgfqpoint{1.082403in}{0.194305in}}%
\pgfpathlineto{\pgfqpoint{1.081247in}{0.188005in}}%
\pgfpathlineto{\pgfqpoint{1.083332in}{0.187330in}}%
\pgfpathlineto{\pgfqpoint{1.085981in}{0.184345in}}%
\pgfpathlineto{\pgfqpoint{1.090049in}{0.178164in}}%
\pgfpathlineto{\pgfqpoint{1.089922in}{0.175542in}}%
\pgfpathlineto{\pgfqpoint{1.087857in}{0.173091in}}%
\pgfpathlineto{\pgfqpoint{1.085139in}{0.174080in}}%
\pgfpathlineto{\pgfqpoint{1.081805in}{0.173502in}}%
\pgfpathlineto{\pgfqpoint{1.082038in}{0.171601in}}%
\pgfpathlineto{\pgfqpoint{1.079838in}{0.170598in}}%
\pgfpathlineto{\pgfqpoint{1.078315in}{0.173884in}}%
\pgfpathlineto{\pgfqpoint{1.075481in}{0.173524in}}%
\pgfpathlineto{\pgfqpoint{1.073238in}{0.172315in}}%
\pgfpathlineto{\pgfqpoint{1.072944in}{0.170978in}}%
\pgfpathlineto{\pgfqpoint{1.070260in}{0.170374in}}%
\pgfpathlineto{\pgfqpoint{1.070834in}{0.168040in}}%
\pgfpathlineto{\pgfqpoint{1.069320in}{0.165633in}}%
\pgfpathlineto{\pgfqpoint{1.067405in}{0.165345in}}%
\pgfpathlineto{\pgfqpoint{1.065482in}{0.169372in}}%
\pgfpathlineto{\pgfqpoint{1.069168in}{0.169743in}}%
\pgfpathlineto{\pgfqpoint{1.070313in}{0.171536in}}%
\pgfpathlineto{\pgfqpoint{1.072838in}{0.172914in}}%
\pgfpathlineto{\pgfqpoint{1.074001in}{0.176696in}}%
\pgfpathlineto{\pgfqpoint{1.075966in}{0.177890in}}%
\pgfpathlineto{\pgfqpoint{1.073883in}{0.179967in}}%
\pgfpathlineto{\pgfqpoint{1.073154in}{0.177269in}}%
\pgfpathlineto{\pgfqpoint{1.069781in}{0.176867in}}%
\pgfpathlineto{\pgfqpoint{1.068604in}{0.173680in}}%
\pgfpathlineto{\pgfqpoint{1.065898in}{0.174566in}}%
\pgfpathlineto{\pgfqpoint{1.065624in}{0.179594in}}%
\pgfpathlineto{\pgfqpoint{1.063296in}{0.179530in}}%
\pgfpathlineto{\pgfqpoint{1.064425in}{0.183198in}}%
\pgfpathlineto{\pgfqpoint{1.069941in}{0.185345in}}%
\pgfpathlineto{\pgfqpoint{1.069443in}{0.182567in}}%
\pgfpathlineto{\pgfqpoint{1.070640in}{0.180586in}}%
\pgfpathlineto{\pgfqpoint{1.071934in}{0.187975in}}%
\pgfpathlineto{\pgfqpoint{1.073144in}{0.190094in}}%
\pgfpathlineto{\pgfqpoint{1.076606in}{0.193050in}}%
\pgfpathlineto{\pgfqpoint{1.072321in}{0.195142in}}%
\pgfpathlineto{\pgfqpoint{1.072458in}{0.197025in}}%
\pgfpathlineto{\pgfqpoint{1.070874in}{0.199642in}}%
\pgfpathlineto{\pgfqpoint{1.070121in}{0.204047in}}%
\pgfpathlineto{\pgfqpoint{1.071795in}{0.204349in}}%
\pgfpathlineto{\pgfqpoint{1.071791in}{0.206742in}}%
\pgfpathlineto{\pgfqpoint{1.069524in}{0.207403in}}%
\pgfpathlineto{\pgfqpoint{1.067853in}{0.209464in}}%
\pgfpathlineto{\pgfqpoint{1.067534in}{0.211631in}}%
\pgfpathlineto{\pgfqpoint{1.065578in}{0.211187in}}%
\pgfpathlineto{\pgfqpoint{1.064056in}{0.213557in}}%
\pgfpathclose%
\pgfusepath{fill}%
\end{pgfscope}%
\begin{pgfscope}%
\pgfpathrectangle{\pgfqpoint{0.100000in}{0.100000in}}{\pgfqpoint{3.608454in}{2.310000in}}%
\pgfusepath{clip}%
\pgfsetbuttcap%
\pgfsetmiterjoin%
\definecolor{currentfill}{rgb}{0.000000,0.827451,0.586275}%
\pgfsetfillcolor{currentfill}%
\pgfsetlinewidth{0.000000pt}%
\definecolor{currentstroke}{rgb}{0.000000,0.000000,0.000000}%
\pgfsetstrokecolor{currentstroke}%
\pgfsetstrokeopacity{0.000000}%
\pgfsetdash{}{0pt}%
\pgfpathmoveto{\pgfqpoint{1.053511in}{0.195656in}}%
\pgfpathlineto{\pgfqpoint{1.056239in}{0.203313in}}%
\pgfpathlineto{\pgfqpoint{1.055808in}{0.208025in}}%
\pgfpathlineto{\pgfqpoint{1.056404in}{0.208822in}}%
\pgfpathlineto{\pgfqpoint{1.055449in}{0.213304in}}%
\pgfpathlineto{\pgfqpoint{1.057580in}{0.213988in}}%
\pgfpathlineto{\pgfqpoint{1.064166in}{0.208247in}}%
\pgfpathlineto{\pgfqpoint{1.067947in}{0.205632in}}%
\pgfpathlineto{\pgfqpoint{1.068462in}{0.202960in}}%
\pgfpathlineto{\pgfqpoint{1.067374in}{0.200817in}}%
\pgfpathlineto{\pgfqpoint{1.069430in}{0.198670in}}%
\pgfpathlineto{\pgfqpoint{1.071014in}{0.192172in}}%
\pgfpathlineto{\pgfqpoint{1.066528in}{0.191514in}}%
\pgfpathlineto{\pgfqpoint{1.064349in}{0.192706in}}%
\pgfpathlineto{\pgfqpoint{1.066503in}{0.196582in}}%
\pgfpathlineto{\pgfqpoint{1.064154in}{0.196268in}}%
\pgfpathlineto{\pgfqpoint{1.062590in}{0.194222in}}%
\pgfpathlineto{\pgfqpoint{1.059615in}{0.193121in}}%
\pgfpathlineto{\pgfqpoint{1.057190in}{0.195536in}}%
\pgfpathlineto{\pgfqpoint{1.055925in}{0.194675in}}%
\pgfpathlineto{\pgfqpoint{1.053511in}{0.195656in}}%
\pgfpathclose%
\pgfusepath{fill}%
\end{pgfscope}%
\begin{pgfscope}%
\pgfpathrectangle{\pgfqpoint{0.100000in}{0.100000in}}{\pgfqpoint{3.608454in}{2.310000in}}%
\pgfusepath{clip}%
\pgfsetbuttcap%
\pgfsetmiterjoin%
\definecolor{currentfill}{rgb}{0.000000,0.694118,0.652941}%
\pgfsetfillcolor{currentfill}%
\pgfsetlinewidth{0.000000pt}%
\definecolor{currentstroke}{rgb}{0.000000,0.000000,0.000000}%
\pgfsetstrokecolor{currentstroke}%
\pgfsetstrokeopacity{0.000000}%
\pgfsetdash{}{0pt}%
\pgfpathmoveto{\pgfqpoint{1.157736in}{1.964285in}}%
\pgfpathlineto{\pgfqpoint{1.140024in}{1.967674in}}%
\pgfpathlineto{\pgfqpoint{1.137125in}{1.963079in}}%
\pgfpathlineto{\pgfqpoint{1.128763in}{1.959930in}}%
\pgfpathlineto{\pgfqpoint{1.126712in}{1.970229in}}%
\pgfpathlineto{\pgfqpoint{1.113810in}{1.978599in}}%
\pgfpathlineto{\pgfqpoint{1.110227in}{1.984421in}}%
\pgfpathlineto{\pgfqpoint{1.103252in}{1.984857in}}%
\pgfpathlineto{\pgfqpoint{1.092311in}{1.978823in}}%
\pgfpathlineto{\pgfqpoint{1.088551in}{1.993499in}}%
\pgfpathlineto{\pgfqpoint{1.081829in}{1.997042in}}%
\pgfpathlineto{\pgfqpoint{1.080846in}{2.002972in}}%
\pgfpathlineto{\pgfqpoint{1.076031in}{2.007092in}}%
\pgfpathlineto{\pgfqpoint{1.076933in}{2.013289in}}%
\pgfpathlineto{\pgfqpoint{1.081296in}{2.021948in}}%
\pgfpathlineto{\pgfqpoint{1.081054in}{2.029697in}}%
\pgfpathlineto{\pgfqpoint{1.077963in}{2.034819in}}%
\pgfpathlineto{\pgfqpoint{1.078479in}{2.042039in}}%
\pgfpathlineto{\pgfqpoint{1.085312in}{2.052758in}}%
\pgfpathlineto{\pgfqpoint{1.092070in}{2.051353in}}%
\pgfpathlineto{\pgfqpoint{1.092981in}{2.055856in}}%
\pgfpathlineto{\pgfqpoint{1.103510in}{2.056028in}}%
\pgfpathlineto{\pgfqpoint{1.108312in}{2.062127in}}%
\pgfpathlineto{\pgfqpoint{1.114370in}{2.060903in}}%
\pgfpathlineto{\pgfqpoint{1.119922in}{2.088067in}}%
\pgfpathlineto{\pgfqpoint{1.111233in}{2.089809in}}%
\pgfpathlineto{\pgfqpoint{1.118004in}{2.122768in}}%
\pgfpathlineto{\pgfqpoint{1.134932in}{2.119015in}}%
\pgfpathlineto{\pgfqpoint{1.137742in}{2.109618in}}%
\pgfpathlineto{\pgfqpoint{1.132903in}{2.085429in}}%
\pgfpathlineto{\pgfqpoint{1.146831in}{2.082471in}}%
\pgfpathlineto{\pgfqpoint{1.141533in}{2.055406in}}%
\pgfpathlineto{\pgfqpoint{1.155484in}{2.052872in}}%
\pgfpathlineto{\pgfqpoint{1.153421in}{2.042231in}}%
\pgfpathlineto{\pgfqpoint{1.160126in}{2.041022in}}%
\pgfpathlineto{\pgfqpoint{1.164980in}{2.032458in}}%
\pgfpathlineto{\pgfqpoint{1.162054in}{2.020418in}}%
\pgfpathlineto{\pgfqpoint{1.158552in}{2.015787in}}%
\pgfpathlineto{\pgfqpoint{1.148015in}{2.013748in}}%
\pgfpathlineto{\pgfqpoint{1.145294in}{2.011278in}}%
\pgfpathlineto{\pgfqpoint{1.143228in}{2.002013in}}%
\pgfpathlineto{\pgfqpoint{1.147664in}{1.994687in}}%
\pgfpathlineto{\pgfqpoint{1.146770in}{1.973621in}}%
\pgfpathlineto{\pgfqpoint{1.157736in}{1.964285in}}%
\pgfpathclose%
\pgfusepath{fill}%
\end{pgfscope}%
\begin{pgfscope}%
\pgfpathrectangle{\pgfqpoint{0.100000in}{0.100000in}}{\pgfqpoint{3.608454in}{2.310000in}}%
\pgfusepath{clip}%
\pgfsetbuttcap%
\pgfsetmiterjoin%
\definecolor{currentfill}{rgb}{0.000000,0.435294,0.782353}%
\pgfsetfillcolor{currentfill}%
\pgfsetlinewidth{0.000000pt}%
\definecolor{currentstroke}{rgb}{0.000000,0.000000,0.000000}%
\pgfsetstrokecolor{currentstroke}%
\pgfsetstrokeopacity{0.000000}%
\pgfsetdash{}{0pt}%
\pgfpathmoveto{\pgfqpoint{1.937428in}{1.476563in}}%
\pgfpathlineto{\pgfqpoint{1.883772in}{1.478524in}}%
\pgfpathlineto{\pgfqpoint{1.884864in}{1.505981in}}%
\pgfpathlineto{\pgfqpoint{1.883282in}{1.506053in}}%
\pgfpathlineto{\pgfqpoint{1.884664in}{1.533508in}}%
\pgfpathlineto{\pgfqpoint{1.911513in}{1.532437in}}%
\pgfpathlineto{\pgfqpoint{1.912100in}{1.522170in}}%
\pgfpathlineto{\pgfqpoint{1.922421in}{1.521780in}}%
\pgfpathlineto{\pgfqpoint{1.925386in}{1.523978in}}%
\pgfpathlineto{\pgfqpoint{1.946070in}{1.524600in}}%
\pgfpathlineto{\pgfqpoint{1.946376in}{1.531379in}}%
\pgfpathlineto{\pgfqpoint{1.952621in}{1.531206in}}%
\pgfpathlineto{\pgfqpoint{1.952480in}{1.526297in}}%
\pgfpathlineto{\pgfqpoint{1.942299in}{1.518835in}}%
\pgfpathlineto{\pgfqpoint{1.938437in}{1.514106in}}%
\pgfpathlineto{\pgfqpoint{1.937428in}{1.476563in}}%
\pgfpathclose%
\pgfusepath{fill}%
\end{pgfscope}%
\begin{pgfscope}%
\pgfpathrectangle{\pgfqpoint{0.100000in}{0.100000in}}{\pgfqpoint{3.608454in}{2.310000in}}%
\pgfusepath{clip}%
\pgfsetbuttcap%
\pgfsetmiterjoin%
\definecolor{currentfill}{rgb}{0.000000,0.678431,0.660784}%
\pgfsetfillcolor{currentfill}%
\pgfsetlinewidth{0.000000pt}%
\definecolor{currentstroke}{rgb}{0.000000,0.000000,0.000000}%
\pgfsetstrokecolor{currentstroke}%
\pgfsetstrokeopacity{0.000000}%
\pgfsetdash{}{0pt}%
\pgfpathmoveto{\pgfqpoint{0.799750in}{1.984499in}}%
\pgfpathlineto{\pgfqpoint{0.808158in}{1.981911in}}%
\pgfpathlineto{\pgfqpoint{0.814622in}{1.975723in}}%
\pgfpathlineto{\pgfqpoint{0.813336in}{1.971025in}}%
\pgfpathlineto{\pgfqpoint{0.807404in}{1.969358in}}%
\pgfpathlineto{\pgfqpoint{0.807519in}{1.962306in}}%
\pgfpathlineto{\pgfqpoint{0.793910in}{1.961931in}}%
\pgfpathlineto{\pgfqpoint{0.793832in}{1.958919in}}%
\pgfpathlineto{\pgfqpoint{0.801936in}{1.953999in}}%
\pgfpathlineto{\pgfqpoint{0.803077in}{1.950037in}}%
\pgfpathlineto{\pgfqpoint{0.793934in}{1.941045in}}%
\pgfpathlineto{\pgfqpoint{0.792882in}{1.935165in}}%
\pgfpathlineto{\pgfqpoint{0.785890in}{1.927298in}}%
\pgfpathlineto{\pgfqpoint{0.800486in}{1.923543in}}%
\pgfpathlineto{\pgfqpoint{0.796408in}{1.906884in}}%
\pgfpathlineto{\pgfqpoint{0.764102in}{1.916069in}}%
\pgfpathlineto{\pgfqpoint{0.762277in}{1.909354in}}%
\pgfpathlineto{\pgfqpoint{0.715935in}{1.921940in}}%
\pgfpathlineto{\pgfqpoint{0.723444in}{1.948508in}}%
\pgfpathlineto{\pgfqpoint{0.734618in}{1.987950in}}%
\pgfpathlineto{\pgfqpoint{0.737189in}{2.001359in}}%
\pgfpathlineto{\pgfqpoint{0.799750in}{1.984499in}}%
\pgfpathclose%
\pgfusepath{fill}%
\end{pgfscope}%
\begin{pgfscope}%
\pgfpathrectangle{\pgfqpoint{0.100000in}{0.100000in}}{\pgfqpoint{3.608454in}{2.310000in}}%
\pgfusepath{clip}%
\pgfsetbuttcap%
\pgfsetmiterjoin%
\definecolor{currentfill}{rgb}{0.000000,0.698039,0.650980}%
\pgfsetfillcolor{currentfill}%
\pgfsetlinewidth{0.000000pt}%
\definecolor{currentstroke}{rgb}{0.000000,0.000000,0.000000}%
\pgfsetstrokecolor{currentstroke}%
\pgfsetstrokeopacity{0.000000}%
\pgfsetdash{}{0pt}%
\pgfpathmoveto{\pgfqpoint{1.267460in}{1.535366in}}%
\pgfpathlineto{\pgfqpoint{1.255196in}{1.457818in}}%
\pgfpathlineto{\pgfqpoint{1.253203in}{1.445105in}}%
\pgfpathlineto{\pgfqpoint{1.249416in}{1.442841in}}%
\pgfpathlineto{\pgfqpoint{1.194163in}{1.451751in}}%
\pgfpathlineto{\pgfqpoint{1.194698in}{1.460074in}}%
\pgfpathlineto{\pgfqpoint{1.199187in}{1.467135in}}%
\pgfpathlineto{\pgfqpoint{1.199560in}{1.472867in}}%
\pgfpathlineto{\pgfqpoint{1.206990in}{1.477188in}}%
\pgfpathlineto{\pgfqpoint{1.169072in}{1.483661in}}%
\pgfpathlineto{\pgfqpoint{1.148850in}{1.487724in}}%
\pgfpathlineto{\pgfqpoint{1.148005in}{1.494799in}}%
\pgfpathlineto{\pgfqpoint{1.158113in}{1.555519in}}%
\pgfpathlineto{\pgfqpoint{1.167943in}{1.559368in}}%
\pgfpathlineto{\pgfqpoint{1.173397in}{1.557849in}}%
\pgfpathlineto{\pgfqpoint{1.190708in}{1.558686in}}%
\pgfpathlineto{\pgfqpoint{1.198089in}{1.561222in}}%
\pgfpathlineto{\pgfqpoint{1.203263in}{1.558508in}}%
\pgfpathlineto{\pgfqpoint{1.213136in}{1.556950in}}%
\pgfpathlineto{\pgfqpoint{1.220332in}{1.551758in}}%
\pgfpathlineto{\pgfqpoint{1.241154in}{1.553467in}}%
\pgfpathlineto{\pgfqpoint{1.242723in}{1.546157in}}%
\pgfpathlineto{\pgfqpoint{1.248664in}{1.549321in}}%
\pgfpathlineto{\pgfqpoint{1.249496in}{1.554569in}}%
\pgfpathlineto{\pgfqpoint{1.260798in}{1.552822in}}%
\pgfpathlineto{\pgfqpoint{1.262249in}{1.546769in}}%
\pgfpathlineto{\pgfqpoint{1.261144in}{1.538773in}}%
\pgfpathlineto{\pgfqpoint{1.267460in}{1.535366in}}%
\pgfpathclose%
\pgfusepath{fill}%
\end{pgfscope}%
\begin{pgfscope}%
\pgfpathrectangle{\pgfqpoint{0.100000in}{0.100000in}}{\pgfqpoint{3.608454in}{2.310000in}}%
\pgfusepath{clip}%
\pgfsetbuttcap%
\pgfsetmiterjoin%
\definecolor{currentfill}{rgb}{0.000000,0.631373,0.684314}%
\pgfsetfillcolor{currentfill}%
\pgfsetlinewidth{0.000000pt}%
\definecolor{currentstroke}{rgb}{0.000000,0.000000,0.000000}%
\pgfsetstrokecolor{currentstroke}%
\pgfsetstrokeopacity{0.000000}%
\pgfsetdash{}{0pt}%
\pgfpathmoveto{\pgfqpoint{2.106820in}{1.155478in}}%
\pgfpathlineto{\pgfqpoint{2.106837in}{1.177041in}}%
\pgfpathlineto{\pgfqpoint{2.102150in}{1.181491in}}%
\pgfpathlineto{\pgfqpoint{2.102082in}{1.211883in}}%
\pgfpathlineto{\pgfqpoint{2.101288in}{1.239463in}}%
\pgfpathlineto{\pgfqpoint{2.101955in}{1.263584in}}%
\pgfpathlineto{\pgfqpoint{2.130827in}{1.263573in}}%
\pgfpathlineto{\pgfqpoint{2.130830in}{1.210394in}}%
\pgfpathlineto{\pgfqpoint{2.164491in}{1.209439in}}%
\pgfpathlineto{\pgfqpoint{2.166408in}{1.204778in}}%
\pgfpathlineto{\pgfqpoint{2.165811in}{1.161887in}}%
\pgfpathlineto{\pgfqpoint{2.130979in}{1.163114in}}%
\pgfpathlineto{\pgfqpoint{2.131026in}{1.155307in}}%
\pgfpathlineto{\pgfqpoint{2.106820in}{1.155478in}}%
\pgfpathclose%
\pgfusepath{fill}%
\end{pgfscope}%
\begin{pgfscope}%
\pgfpathrectangle{\pgfqpoint{0.100000in}{0.100000in}}{\pgfqpoint{3.608454in}{2.310000in}}%
\pgfusepath{clip}%
\pgfsetbuttcap%
\pgfsetmiterjoin%
\definecolor{currentfill}{rgb}{0.000000,0.333333,0.833333}%
\pgfsetfillcolor{currentfill}%
\pgfsetlinewidth{0.000000pt}%
\definecolor{currentstroke}{rgb}{0.000000,0.000000,0.000000}%
\pgfsetstrokecolor{currentstroke}%
\pgfsetstrokeopacity{0.000000}%
\pgfsetdash{}{0pt}%
\pgfpathmoveto{\pgfqpoint{3.080944in}{0.255920in}}%
\pgfpathlineto{\pgfqpoint{3.081905in}{0.258847in}}%
\pgfpathlineto{\pgfqpoint{3.093100in}{0.264126in}}%
\pgfpathlineto{\pgfqpoint{3.097421in}{0.270578in}}%
\pgfpathlineto{\pgfqpoint{3.104568in}{0.275364in}}%
\pgfpathlineto{\pgfqpoint{3.114895in}{0.274075in}}%
\pgfpathlineto{\pgfqpoint{3.119568in}{0.271124in}}%
\pgfpathlineto{\pgfqpoint{3.110703in}{0.266525in}}%
\pgfpathlineto{\pgfqpoint{3.103147in}{0.268644in}}%
\pgfpathlineto{\pgfqpoint{3.101531in}{0.264974in}}%
\pgfpathlineto{\pgfqpoint{3.089987in}{0.258378in}}%
\pgfpathlineto{\pgfqpoint{3.080944in}{0.255920in}}%
\pgfpathclose%
\pgfusepath{fill}%
\end{pgfscope}%
\begin{pgfscope}%
\pgfpathrectangle{\pgfqpoint{0.100000in}{0.100000in}}{\pgfqpoint{3.608454in}{2.310000in}}%
\pgfusepath{clip}%
\pgfsetbuttcap%
\pgfsetmiterjoin%
\definecolor{currentfill}{rgb}{0.000000,0.333333,0.833333}%
\pgfsetfillcolor{currentfill}%
\pgfsetlinewidth{0.000000pt}%
\definecolor{currentstroke}{rgb}{0.000000,0.000000,0.000000}%
\pgfsetstrokecolor{currentstroke}%
\pgfsetstrokeopacity{0.000000}%
\pgfsetdash{}{0pt}%
\pgfpathmoveto{\pgfqpoint{3.099940in}{0.362168in}}%
\pgfpathlineto{\pgfqpoint{3.133900in}{0.367620in}}%
\pgfpathlineto{\pgfqpoint{3.127726in}{0.403873in}}%
\pgfpathlineto{\pgfqpoint{3.126693in}{0.409803in}}%
\pgfpathlineto{\pgfqpoint{3.173551in}{0.417385in}}%
\pgfpathlineto{\pgfqpoint{3.184616in}{0.418072in}}%
\pgfpathlineto{\pgfqpoint{3.185207in}{0.398995in}}%
\pgfpathlineto{\pgfqpoint{3.187948in}{0.377138in}}%
\pgfpathlineto{\pgfqpoint{3.187162in}{0.370691in}}%
\pgfpathlineto{\pgfqpoint{3.180144in}{0.367508in}}%
\pgfpathlineto{\pgfqpoint{3.176678in}{0.346943in}}%
\pgfpathlineto{\pgfqpoint{3.178542in}{0.336148in}}%
\pgfpathlineto{\pgfqpoint{3.168359in}{0.323529in}}%
\pgfpathlineto{\pgfqpoint{3.163072in}{0.325621in}}%
\pgfpathlineto{\pgfqpoint{3.160642in}{0.320071in}}%
\pgfpathlineto{\pgfqpoint{3.154594in}{0.316205in}}%
\pgfpathlineto{\pgfqpoint{3.143564in}{0.317485in}}%
\pgfpathlineto{\pgfqpoint{3.140330in}{0.313854in}}%
\pgfpathlineto{\pgfqpoint{3.126999in}{0.309857in}}%
\pgfpathlineto{\pgfqpoint{3.119526in}{0.317390in}}%
\pgfpathlineto{\pgfqpoint{3.119713in}{0.326025in}}%
\pgfpathlineto{\pgfqpoint{3.116825in}{0.336890in}}%
\pgfpathlineto{\pgfqpoint{3.113204in}{0.339888in}}%
\pgfpathlineto{\pgfqpoint{3.109115in}{0.349472in}}%
\pgfpathlineto{\pgfqpoint{3.101613in}{0.355282in}}%
\pgfpathlineto{\pgfqpoint{3.099940in}{0.362168in}}%
\pgfpathclose%
\pgfusepath{fill}%
\end{pgfscope}%
\begin{pgfscope}%
\pgfpathrectangle{\pgfqpoint{0.100000in}{0.100000in}}{\pgfqpoint{3.608454in}{2.310000in}}%
\pgfusepath{clip}%
\pgfsetbuttcap%
\pgfsetmiterjoin%
\definecolor{currentfill}{rgb}{0.000000,0.698039,0.650980}%
\pgfsetfillcolor{currentfill}%
\pgfsetlinewidth{0.000000pt}%
\definecolor{currentstroke}{rgb}{0.000000,0.000000,0.000000}%
\pgfsetstrokecolor{currentstroke}%
\pgfsetstrokeopacity{0.000000}%
\pgfsetdash{}{0pt}%
\pgfpathmoveto{\pgfqpoint{2.547754in}{0.792351in}}%
\pgfpathlineto{\pgfqpoint{2.536679in}{0.788374in}}%
\pgfpathlineto{\pgfqpoint{2.518619in}{0.784958in}}%
\pgfpathlineto{\pgfqpoint{2.516272in}{0.816527in}}%
\pgfpathlineto{\pgfqpoint{2.488805in}{0.814690in}}%
\pgfpathlineto{\pgfqpoint{2.485523in}{0.870923in}}%
\pgfpathlineto{\pgfqpoint{2.519263in}{0.872679in}}%
\pgfpathlineto{\pgfqpoint{2.550184in}{0.875112in}}%
\pgfpathlineto{\pgfqpoint{2.549490in}{0.847081in}}%
\pgfpathlineto{\pgfqpoint{2.547754in}{0.792351in}}%
\pgfpathclose%
\pgfusepath{fill}%
\end{pgfscope}%
\begin{pgfscope}%
\pgfpathrectangle{\pgfqpoint{0.100000in}{0.100000in}}{\pgfqpoint{3.608454in}{2.310000in}}%
\pgfusepath{clip}%
\pgfsetbuttcap%
\pgfsetmiterjoin%
\definecolor{currentfill}{rgb}{0.000000,0.490196,0.754902}%
\pgfsetfillcolor{currentfill}%
\pgfsetlinewidth{0.000000pt}%
\definecolor{currentstroke}{rgb}{0.000000,0.000000,0.000000}%
\pgfsetstrokecolor{currentstroke}%
\pgfsetstrokeopacity{0.000000}%
\pgfsetdash{}{0pt}%
\pgfpathmoveto{\pgfqpoint{1.244902in}{0.674701in}}%
\pgfpathlineto{\pgfqpoint{1.249588in}{0.682742in}}%
\pgfpathlineto{\pgfqpoint{1.263659in}{0.679913in}}%
\pgfpathlineto{\pgfqpoint{1.258636in}{0.675464in}}%
\pgfpathlineto{\pgfqpoint{1.252077in}{0.677773in}}%
\pgfpathlineto{\pgfqpoint{1.244902in}{0.674701in}}%
\pgfpathclose%
\pgfusepath{fill}%
\end{pgfscope}%
\begin{pgfscope}%
\pgfpathrectangle{\pgfqpoint{0.100000in}{0.100000in}}{\pgfqpoint{3.608454in}{2.310000in}}%
\pgfusepath{clip}%
\pgfsetbuttcap%
\pgfsetmiterjoin%
\definecolor{currentfill}{rgb}{0.000000,0.490196,0.754902}%
\pgfsetfillcolor{currentfill}%
\pgfsetlinewidth{0.000000pt}%
\definecolor{currentstroke}{rgb}{0.000000,0.000000,0.000000}%
\pgfsetstrokecolor{currentstroke}%
\pgfsetstrokeopacity{0.000000}%
\pgfsetdash{}{0pt}%
\pgfpathmoveto{\pgfqpoint{1.293790in}{0.642427in}}%
\pgfpathlineto{\pgfqpoint{1.285285in}{0.654449in}}%
\pgfpathlineto{\pgfqpoint{1.281587in}{0.664883in}}%
\pgfpathlineto{\pgfqpoint{1.282478in}{0.669549in}}%
\pgfpathlineto{\pgfqpoint{1.292812in}{0.673886in}}%
\pgfpathlineto{\pgfqpoint{1.304918in}{0.671461in}}%
\pgfpathlineto{\pgfqpoint{1.308962in}{0.666408in}}%
\pgfpathlineto{\pgfqpoint{1.315946in}{0.662668in}}%
\pgfpathlineto{\pgfqpoint{1.318323in}{0.655591in}}%
\pgfpathlineto{\pgfqpoint{1.315484in}{0.650135in}}%
\pgfpathlineto{\pgfqpoint{1.309751in}{0.648208in}}%
\pgfpathlineto{\pgfqpoint{1.305345in}{0.642514in}}%
\pgfpathlineto{\pgfqpoint{1.293790in}{0.642427in}}%
\pgfpathclose%
\pgfusepath{fill}%
\end{pgfscope}%
\begin{pgfscope}%
\pgfpathrectangle{\pgfqpoint{0.100000in}{0.100000in}}{\pgfqpoint{3.608454in}{2.310000in}}%
\pgfusepath{clip}%
\pgfsetbuttcap%
\pgfsetmiterjoin%
\definecolor{currentfill}{rgb}{0.000000,0.478431,0.760784}%
\pgfsetfillcolor{currentfill}%
\pgfsetlinewidth{0.000000pt}%
\definecolor{currentstroke}{rgb}{0.000000,0.000000,0.000000}%
\pgfsetstrokecolor{currentstroke}%
\pgfsetstrokeopacity{0.000000}%
\pgfsetdash{}{0pt}%
\pgfpathmoveto{\pgfqpoint{1.946085in}{2.030374in}}%
\pgfpathlineto{\pgfqpoint{1.945081in}{2.002729in}}%
\pgfpathlineto{\pgfqpoint{1.946235in}{1.995849in}}%
\pgfpathlineto{\pgfqpoint{1.918791in}{1.996815in}}%
\pgfpathlineto{\pgfqpoint{1.917334in}{2.003750in}}%
\pgfpathlineto{\pgfqpoint{1.918355in}{2.031373in}}%
\pgfpathlineto{\pgfqpoint{1.946085in}{2.030374in}}%
\pgfpathclose%
\pgfusepath{fill}%
\end{pgfscope}%
\begin{pgfscope}%
\pgfpathrectangle{\pgfqpoint{0.100000in}{0.100000in}}{\pgfqpoint{3.608454in}{2.310000in}}%
\pgfusepath{clip}%
\pgfsetbuttcap%
\pgfsetmiterjoin%
\definecolor{currentfill}{rgb}{0.000000,0.439216,0.780392}%
\pgfsetfillcolor{currentfill}%
\pgfsetlinewidth{0.000000pt}%
\definecolor{currentstroke}{rgb}{0.000000,0.000000,0.000000}%
\pgfsetstrokecolor{currentstroke}%
\pgfsetstrokeopacity{0.000000}%
\pgfsetdash{}{0pt}%
\pgfpathmoveto{\pgfqpoint{1.474731in}{1.733639in}}%
\pgfpathlineto{\pgfqpoint{1.441739in}{1.737485in}}%
\pgfpathlineto{\pgfqpoint{1.412417in}{1.742036in}}%
\pgfpathlineto{\pgfqpoint{1.417395in}{1.794618in}}%
\pgfpathlineto{\pgfqpoint{1.419171in}{1.807060in}}%
\pgfpathlineto{\pgfqpoint{1.417230in}{1.817131in}}%
\pgfpathlineto{\pgfqpoint{1.412375in}{1.821775in}}%
\pgfpathlineto{\pgfqpoint{1.406562in}{1.831802in}}%
\pgfpathlineto{\pgfqpoint{1.396824in}{1.840362in}}%
\pgfpathlineto{\pgfqpoint{1.389944in}{1.843550in}}%
\pgfpathlineto{\pgfqpoint{1.384830in}{1.855882in}}%
\pgfpathlineto{\pgfqpoint{1.383661in}{1.866297in}}%
\pgfpathlineto{\pgfqpoint{1.364865in}{1.868873in}}%
\pgfpathlineto{\pgfqpoint{1.372850in}{1.883266in}}%
\pgfpathlineto{\pgfqpoint{1.375793in}{1.885029in}}%
\pgfpathlineto{\pgfqpoint{1.342451in}{1.889782in}}%
\pgfpathlineto{\pgfqpoint{1.345983in}{1.906551in}}%
\pgfpathlineto{\pgfqpoint{1.348402in}{1.908679in}}%
\pgfpathlineto{\pgfqpoint{1.366491in}{1.905688in}}%
\pgfpathlineto{\pgfqpoint{1.380767in}{1.908242in}}%
\pgfpathlineto{\pgfqpoint{1.383830in}{1.929135in}}%
\pgfpathlineto{\pgfqpoint{1.386336in}{1.938080in}}%
\pgfpathlineto{\pgfqpoint{1.395427in}{1.936746in}}%
\pgfpathlineto{\pgfqpoint{1.400977in}{1.942930in}}%
\pgfpathlineto{\pgfqpoint{1.407750in}{1.941975in}}%
\pgfpathlineto{\pgfqpoint{1.408401in}{1.946543in}}%
\pgfpathlineto{\pgfqpoint{1.417308in}{1.945271in}}%
\pgfpathlineto{\pgfqpoint{1.422053in}{1.944628in}}%
\pgfpathlineto{\pgfqpoint{1.421111in}{1.937794in}}%
\pgfpathlineto{\pgfqpoint{1.434747in}{1.935937in}}%
\pgfpathlineto{\pgfqpoint{1.446859in}{1.927356in}}%
\pgfpathlineto{\pgfqpoint{1.446358in}{1.912489in}}%
\pgfpathlineto{\pgfqpoint{1.456350in}{1.910291in}}%
\pgfpathlineto{\pgfqpoint{1.453065in}{1.884941in}}%
\pgfpathlineto{\pgfqpoint{1.449171in}{1.871656in}}%
\pgfpathlineto{\pgfqpoint{1.477206in}{1.868092in}}%
\pgfpathlineto{\pgfqpoint{1.475384in}{1.853420in}}%
\pgfpathlineto{\pgfqpoint{1.488701in}{1.851762in}}%
\pgfpathlineto{\pgfqpoint{1.486605in}{1.827974in}}%
\pgfpathlineto{\pgfqpoint{1.481253in}{1.786760in}}%
\pgfpathlineto{\pgfqpoint{1.474731in}{1.733639in}}%
\pgfpathclose%
\pgfusepath{fill}%
\end{pgfscope}%
\begin{pgfscope}%
\pgfpathrectangle{\pgfqpoint{0.100000in}{0.100000in}}{\pgfqpoint{3.608454in}{2.310000in}}%
\pgfusepath{clip}%
\pgfsetbuttcap%
\pgfsetmiterjoin%
\definecolor{currentfill}{rgb}{0.000000,0.603922,0.698039}%
\pgfsetfillcolor{currentfill}%
\pgfsetlinewidth{0.000000pt}%
\definecolor{currentstroke}{rgb}{0.000000,0.000000,0.000000}%
\pgfsetstrokecolor{currentstroke}%
\pgfsetstrokeopacity{0.000000}%
\pgfsetdash{}{0pt}%
\pgfpathmoveto{\pgfqpoint{2.184096in}{1.723863in}}%
\pgfpathlineto{\pgfqpoint{2.177278in}{1.723751in}}%
\pgfpathlineto{\pgfqpoint{2.176891in}{1.751349in}}%
\pgfpathlineto{\pgfqpoint{2.190701in}{1.751566in}}%
\pgfpathlineto{\pgfqpoint{2.190305in}{1.779149in}}%
\pgfpathlineto{\pgfqpoint{2.204050in}{1.779422in}}%
\pgfpathlineto{\pgfqpoint{2.204133in}{1.773697in}}%
\pgfpathlineto{\pgfqpoint{2.217878in}{1.773970in}}%
\pgfpathlineto{\pgfqpoint{2.218568in}{1.724511in}}%
\pgfpathlineto{\pgfqpoint{2.184096in}{1.723863in}}%
\pgfpathclose%
\pgfusepath{fill}%
\end{pgfscope}%
\begin{pgfscope}%
\pgfpathrectangle{\pgfqpoint{0.100000in}{0.100000in}}{\pgfqpoint{3.608454in}{2.310000in}}%
\pgfusepath{clip}%
\pgfsetbuttcap%
\pgfsetmiterjoin%
\definecolor{currentfill}{rgb}{0.000000,0.364706,0.817647}%
\pgfsetfillcolor{currentfill}%
\pgfsetlinewidth{0.000000pt}%
\definecolor{currentstroke}{rgb}{0.000000,0.000000,0.000000}%
\pgfsetstrokecolor{currentstroke}%
\pgfsetstrokeopacity{0.000000}%
\pgfsetdash{}{0pt}%
\pgfpathmoveto{\pgfqpoint{1.789033in}{1.324734in}}%
\pgfpathlineto{\pgfqpoint{1.824072in}{1.322793in}}%
\pgfpathlineto{\pgfqpoint{1.822395in}{1.288502in}}%
\pgfpathlineto{\pgfqpoint{1.782549in}{1.290657in}}%
\pgfpathlineto{\pgfqpoint{1.754213in}{1.292521in}}%
\pgfpathlineto{\pgfqpoint{1.756162in}{1.326912in}}%
\pgfpathlineto{\pgfqpoint{1.789033in}{1.324734in}}%
\pgfpathclose%
\pgfusepath{fill}%
\end{pgfscope}%
\begin{pgfscope}%
\pgfpathrectangle{\pgfqpoint{0.100000in}{0.100000in}}{\pgfqpoint{3.608454in}{2.310000in}}%
\pgfusepath{clip}%
\pgfsetbuttcap%
\pgfsetmiterjoin%
\definecolor{currentfill}{rgb}{0.000000,0.619608,0.690196}%
\pgfsetfillcolor{currentfill}%
\pgfsetlinewidth{0.000000pt}%
\definecolor{currentstroke}{rgb}{0.000000,0.000000,0.000000}%
\pgfsetstrokecolor{currentstroke}%
\pgfsetstrokeopacity{0.000000}%
\pgfsetdash{}{0pt}%
\pgfpathmoveto{\pgfqpoint{1.720336in}{1.913435in}}%
\pgfpathlineto{\pgfqpoint{1.726123in}{1.914246in}}%
\pgfpathlineto{\pgfqpoint{1.730862in}{1.910799in}}%
\pgfpathlineto{\pgfqpoint{1.744713in}{1.908214in}}%
\pgfpathlineto{\pgfqpoint{1.756437in}{1.912852in}}%
\pgfpathlineto{\pgfqpoint{1.759368in}{1.916843in}}%
\pgfpathlineto{\pgfqpoint{1.775193in}{1.927930in}}%
\pgfpathlineto{\pgfqpoint{1.783790in}{1.937380in}}%
\pgfpathlineto{\pgfqpoint{1.787591in}{1.935130in}}%
\pgfpathlineto{\pgfqpoint{1.799830in}{1.937828in}}%
\pgfpathlineto{\pgfqpoint{1.801655in}{1.932573in}}%
\pgfpathlineto{\pgfqpoint{1.800710in}{1.922898in}}%
\pgfpathlineto{\pgfqpoint{1.796911in}{1.916921in}}%
\pgfpathlineto{\pgfqpoint{1.796957in}{1.908802in}}%
\pgfpathlineto{\pgfqpoint{1.802493in}{1.899008in}}%
\pgfpathlineto{\pgfqpoint{1.808345in}{1.892715in}}%
\pgfpathlineto{\pgfqpoint{1.811895in}{1.879653in}}%
\pgfpathlineto{\pgfqpoint{1.807298in}{1.876074in}}%
\pgfpathlineto{\pgfqpoint{1.801983in}{1.865242in}}%
\pgfpathlineto{\pgfqpoint{1.808999in}{1.861084in}}%
\pgfpathlineto{\pgfqpoint{1.773598in}{1.863245in}}%
\pgfpathlineto{\pgfqpoint{1.716695in}{1.867335in}}%
\pgfpathlineto{\pgfqpoint{1.720336in}{1.913435in}}%
\pgfpathclose%
\pgfusepath{fill}%
\end{pgfscope}%
\begin{pgfscope}%
\pgfpathrectangle{\pgfqpoint{0.100000in}{0.100000in}}{\pgfqpoint{3.608454in}{2.310000in}}%
\pgfusepath{clip}%
\pgfsetbuttcap%
\pgfsetmiterjoin%
\definecolor{currentfill}{rgb}{0.000000,0.407843,0.796078}%
\pgfsetfillcolor{currentfill}%
\pgfsetlinewidth{0.000000pt}%
\definecolor{currentstroke}{rgb}{0.000000,0.000000,0.000000}%
\pgfsetstrokecolor{currentstroke}%
\pgfsetstrokeopacity{0.000000}%
\pgfsetdash{}{0pt}%
\pgfpathmoveto{\pgfqpoint{1.347428in}{1.265725in}}%
\pgfpathlineto{\pgfqpoint{1.325433in}{1.268700in}}%
\pgfpathlineto{\pgfqpoint{1.329166in}{1.295757in}}%
\pgfpathlineto{\pgfqpoint{1.327961in}{1.309954in}}%
\pgfpathlineto{\pgfqpoint{1.325950in}{1.311919in}}%
\pgfpathlineto{\pgfqpoint{1.328108in}{1.320759in}}%
\pgfpathlineto{\pgfqpoint{1.328146in}{1.332828in}}%
\pgfpathlineto{\pgfqpoint{1.325556in}{1.338863in}}%
\pgfpathlineto{\pgfqpoint{1.333852in}{1.337711in}}%
\pgfpathlineto{\pgfqpoint{1.343717in}{1.409436in}}%
\pgfpathlineto{\pgfqpoint{1.350593in}{1.411552in}}%
\pgfpathlineto{\pgfqpoint{1.352191in}{1.404691in}}%
\pgfpathlineto{\pgfqpoint{1.355918in}{1.399849in}}%
\pgfpathlineto{\pgfqpoint{1.369057in}{1.398085in}}%
\pgfpathlineto{\pgfqpoint{1.377365in}{1.386918in}}%
\pgfpathlineto{\pgfqpoint{1.383841in}{1.385038in}}%
\pgfpathlineto{\pgfqpoint{1.387549in}{1.389567in}}%
\pgfpathlineto{\pgfqpoint{1.393282in}{1.388976in}}%
\pgfpathlineto{\pgfqpoint{1.396241in}{1.384895in}}%
\pgfpathlineto{\pgfqpoint{1.403546in}{1.376980in}}%
\pgfpathlineto{\pgfqpoint{1.412208in}{1.374642in}}%
\pgfpathlineto{\pgfqpoint{1.405672in}{1.369838in}}%
\pgfpathlineto{\pgfqpoint{1.404621in}{1.362130in}}%
\pgfpathlineto{\pgfqpoint{1.406993in}{1.344082in}}%
\pgfpathlineto{\pgfqpoint{1.412097in}{1.337063in}}%
\pgfpathlineto{\pgfqpoint{1.365837in}{1.343240in}}%
\pgfpathlineto{\pgfqpoint{1.360950in}{1.306443in}}%
\pgfpathlineto{\pgfqpoint{1.352269in}{1.306231in}}%
\pgfpathlineto{\pgfqpoint{1.349065in}{1.285654in}}%
\pgfpathlineto{\pgfqpoint{1.347428in}{1.265725in}}%
\pgfpathclose%
\pgfusepath{fill}%
\end{pgfscope}%
\begin{pgfscope}%
\pgfpathrectangle{\pgfqpoint{0.100000in}{0.100000in}}{\pgfqpoint{3.608454in}{2.310000in}}%
\pgfusepath{clip}%
\pgfsetbuttcap%
\pgfsetmiterjoin%
\definecolor{currentfill}{rgb}{0.000000,0.796078,0.601961}%
\pgfsetfillcolor{currentfill}%
\pgfsetlinewidth{0.000000pt}%
\definecolor{currentstroke}{rgb}{0.000000,0.000000,0.000000}%
\pgfsetstrokecolor{currentstroke}%
\pgfsetstrokeopacity{0.000000}%
\pgfsetdash{}{0pt}%
\pgfpathmoveto{\pgfqpoint{1.716695in}{1.867335in}}%
\pgfpathlineto{\pgfqpoint{1.773598in}{1.863245in}}%
\pgfpathlineto{\pgfqpoint{1.808999in}{1.861084in}}%
\pgfpathlineto{\pgfqpoint{1.812344in}{1.853586in}}%
\pgfpathlineto{\pgfqpoint{1.809666in}{1.848703in}}%
\pgfpathlineto{\pgfqpoint{1.812489in}{1.842876in}}%
\pgfpathlineto{\pgfqpoint{1.808240in}{1.831561in}}%
\pgfpathlineto{\pgfqpoint{1.809782in}{1.824532in}}%
\pgfpathlineto{\pgfqpoint{1.802162in}{1.823198in}}%
\pgfpathlineto{\pgfqpoint{1.802764in}{1.815650in}}%
\pgfpathlineto{\pgfqpoint{1.803220in}{1.813204in}}%
\pgfpathlineto{\pgfqpoint{1.795074in}{1.805632in}}%
\pgfpathlineto{\pgfqpoint{1.787641in}{1.810142in}}%
\pgfpathlineto{\pgfqpoint{1.781507in}{1.805854in}}%
\pgfpathlineto{\pgfqpoint{1.778292in}{1.808049in}}%
\pgfpathlineto{\pgfqpoint{1.767130in}{1.804067in}}%
\pgfpathlineto{\pgfqpoint{1.760814in}{1.806325in}}%
\pgfpathlineto{\pgfqpoint{1.754993in}{1.802225in}}%
\pgfpathlineto{\pgfqpoint{1.747956in}{1.802589in}}%
\pgfpathlineto{\pgfqpoint{1.737625in}{1.793677in}}%
\pgfpathlineto{\pgfqpoint{1.729800in}{1.795494in}}%
\pgfpathlineto{\pgfqpoint{1.723276in}{1.792199in}}%
\pgfpathlineto{\pgfqpoint{1.710799in}{1.791190in}}%
\pgfpathlineto{\pgfqpoint{1.714041in}{1.832923in}}%
\pgfpathlineto{\pgfqpoint{1.716695in}{1.867335in}}%
\pgfpathclose%
\pgfusepath{fill}%
\end{pgfscope}%
\begin{pgfscope}%
\pgfpathrectangle{\pgfqpoint{0.100000in}{0.100000in}}{\pgfqpoint{3.608454in}{2.310000in}}%
\pgfusepath{clip}%
\pgfsetbuttcap%
\pgfsetmiterjoin%
\definecolor{currentfill}{rgb}{0.000000,0.490196,0.754902}%
\pgfsetfillcolor{currentfill}%
\pgfsetlinewidth{0.000000pt}%
\definecolor{currentstroke}{rgb}{0.000000,0.000000,0.000000}%
\pgfsetstrokecolor{currentstroke}%
\pgfsetstrokeopacity{0.000000}%
\pgfsetdash{}{0pt}%
\pgfpathmoveto{\pgfqpoint{2.835111in}{1.113784in}}%
\pgfpathlineto{\pgfqpoint{2.828227in}{1.112365in}}%
\pgfpathlineto{\pgfqpoint{2.817068in}{1.103142in}}%
\pgfpathlineto{\pgfqpoint{2.805797in}{1.111156in}}%
\pgfpathlineto{\pgfqpoint{2.790002in}{1.116268in}}%
\pgfpathlineto{\pgfqpoint{2.775615in}{1.111731in}}%
\pgfpathlineto{\pgfqpoint{2.772428in}{1.117636in}}%
\pgfpathlineto{\pgfqpoint{2.765621in}{1.120136in}}%
\pgfpathlineto{\pgfqpoint{2.761095in}{1.125193in}}%
\pgfpathlineto{\pgfqpoint{2.766985in}{1.132686in}}%
\pgfpathlineto{\pgfqpoint{2.763659in}{1.139195in}}%
\pgfpathlineto{\pgfqpoint{2.750357in}{1.150558in}}%
\pgfpathlineto{\pgfqpoint{2.751113in}{1.160516in}}%
\pgfpathlineto{\pgfqpoint{2.761563in}{1.168904in}}%
\pgfpathlineto{\pgfqpoint{2.771140in}{1.159342in}}%
\pgfpathlineto{\pgfqpoint{2.781495in}{1.155039in}}%
\pgfpathlineto{\pgfqpoint{2.785832in}{1.164847in}}%
\pgfpathlineto{\pgfqpoint{2.783231in}{1.181649in}}%
\pgfpathlineto{\pgfqpoint{2.788196in}{1.183891in}}%
\pgfpathlineto{\pgfqpoint{2.789452in}{1.189777in}}%
\pgfpathlineto{\pgfqpoint{2.804581in}{1.191394in}}%
\pgfpathlineto{\pgfqpoint{2.810250in}{1.179188in}}%
\pgfpathlineto{\pgfqpoint{2.814063in}{1.180206in}}%
\pgfpathlineto{\pgfqpoint{2.827195in}{1.174638in}}%
\pgfpathlineto{\pgfqpoint{2.824792in}{1.161761in}}%
\pgfpathlineto{\pgfqpoint{2.831674in}{1.145493in}}%
\pgfpathlineto{\pgfqpoint{2.823619in}{1.137722in}}%
\pgfpathlineto{\pgfqpoint{2.829784in}{1.127369in}}%
\pgfpathlineto{\pgfqpoint{2.833923in}{1.123330in}}%
\pgfpathlineto{\pgfqpoint{2.835111in}{1.113784in}}%
\pgfpathclose%
\pgfusepath{fill}%
\end{pgfscope}%
\begin{pgfscope}%
\pgfpathrectangle{\pgfqpoint{0.100000in}{0.100000in}}{\pgfqpoint{3.608454in}{2.310000in}}%
\pgfusepath{clip}%
\pgfsetbuttcap%
\pgfsetmiterjoin%
\definecolor{currentfill}{rgb}{0.000000,0.364706,0.817647}%
\pgfsetfillcolor{currentfill}%
\pgfsetlinewidth{0.000000pt}%
\definecolor{currentstroke}{rgb}{0.000000,0.000000,0.000000}%
\pgfsetstrokecolor{currentstroke}%
\pgfsetstrokeopacity{0.000000}%
\pgfsetdash{}{0pt}%
\pgfpathmoveto{\pgfqpoint{1.793434in}{1.427744in}}%
\pgfpathlineto{\pgfqpoint{1.791324in}{1.393448in}}%
\pgfpathlineto{\pgfqpoint{1.717833in}{1.398242in}}%
\pgfpathlineto{\pgfqpoint{1.716323in}{1.398352in}}%
\pgfpathlineto{\pgfqpoint{1.718916in}{1.432564in}}%
\pgfpathlineto{\pgfqpoint{1.724173in}{1.432221in}}%
\pgfpathlineto{\pgfqpoint{1.792481in}{1.427797in}}%
\pgfpathlineto{\pgfqpoint{1.793434in}{1.427744in}}%
\pgfpathclose%
\pgfusepath{fill}%
\end{pgfscope}%
\begin{pgfscope}%
\pgfpathrectangle{\pgfqpoint{0.100000in}{0.100000in}}{\pgfqpoint{3.608454in}{2.310000in}}%
\pgfusepath{clip}%
\pgfsetbuttcap%
\pgfsetmiterjoin%
\definecolor{currentfill}{rgb}{0.000000,0.607843,0.696078}%
\pgfsetfillcolor{currentfill}%
\pgfsetlinewidth{0.000000pt}%
\definecolor{currentstroke}{rgb}{0.000000,0.000000,0.000000}%
\pgfsetstrokecolor{currentstroke}%
\pgfsetstrokeopacity{0.000000}%
\pgfsetdash{}{0pt}%
\pgfpathmoveto{\pgfqpoint{3.083124in}{1.089029in}}%
\pgfpathlineto{\pgfqpoint{3.099339in}{1.091621in}}%
\pgfpathlineto{\pgfqpoint{3.116266in}{1.079750in}}%
\pgfpathlineto{\pgfqpoint{3.108787in}{1.052476in}}%
\pgfpathlineto{\pgfqpoint{3.095167in}{1.065172in}}%
\pgfpathlineto{\pgfqpoint{3.079593in}{1.063190in}}%
\pgfpathlineto{\pgfqpoint{3.065576in}{1.050744in}}%
\pgfpathlineto{\pgfqpoint{3.060908in}{1.061022in}}%
\pgfpathlineto{\pgfqpoint{3.054685in}{1.068940in}}%
\pgfpathlineto{\pgfqpoint{3.042269in}{1.083215in}}%
\pgfpathlineto{\pgfqpoint{3.083124in}{1.089029in}}%
\pgfpathclose%
\pgfusepath{fill}%
\end{pgfscope}%
\begin{pgfscope}%
\pgfpathrectangle{\pgfqpoint{0.100000in}{0.100000in}}{\pgfqpoint{3.608454in}{2.310000in}}%
\pgfusepath{clip}%
\pgfsetbuttcap%
\pgfsetmiterjoin%
\definecolor{currentfill}{rgb}{0.000000,0.713725,0.643137}%
\pgfsetfillcolor{currentfill}%
\pgfsetlinewidth{0.000000pt}%
\definecolor{currentstroke}{rgb}{0.000000,0.000000,0.000000}%
\pgfsetstrokecolor{currentstroke}%
\pgfsetstrokeopacity{0.000000}%
\pgfsetdash{}{0pt}%
\pgfpathmoveto{\pgfqpoint{1.748728in}{0.562219in}}%
\pgfpathlineto{\pgfqpoint{1.709808in}{0.564368in}}%
\pgfpathlineto{\pgfqpoint{1.702681in}{0.571973in}}%
\pgfpathlineto{\pgfqpoint{1.701376in}{0.578127in}}%
\pgfpathlineto{\pgfqpoint{1.687298in}{0.588703in}}%
\pgfpathlineto{\pgfqpoint{1.682525in}{0.597073in}}%
\pgfpathlineto{\pgfqpoint{1.678141in}{0.597930in}}%
\pgfpathlineto{\pgfqpoint{1.672342in}{0.605442in}}%
\pgfpathlineto{\pgfqpoint{1.667527in}{0.607537in}}%
\pgfpathlineto{\pgfqpoint{1.661236in}{0.620787in}}%
\pgfpathlineto{\pgfqpoint{1.653250in}{0.625462in}}%
\pgfpathlineto{\pgfqpoint{1.648358in}{0.623948in}}%
\pgfpathlineto{\pgfqpoint{1.623671in}{0.627789in}}%
\pgfpathlineto{\pgfqpoint{1.617429in}{0.627185in}}%
\pgfpathlineto{\pgfqpoint{1.599353in}{0.636098in}}%
\pgfpathlineto{\pgfqpoint{1.583492in}{0.651351in}}%
\pgfpathlineto{\pgfqpoint{1.585018in}{0.669731in}}%
\pgfpathlineto{\pgfqpoint{1.600425in}{0.668531in}}%
\pgfpathlineto{\pgfqpoint{1.602461in}{0.693594in}}%
\pgfpathlineto{\pgfqpoint{1.616475in}{0.692376in}}%
\pgfpathlineto{\pgfqpoint{1.616885in}{0.697027in}}%
\pgfpathlineto{\pgfqpoint{1.644365in}{0.694777in}}%
\pgfpathlineto{\pgfqpoint{1.648797in}{0.693179in}}%
\pgfpathlineto{\pgfqpoint{1.649534in}{0.685985in}}%
\pgfpathlineto{\pgfqpoint{1.642706in}{0.676280in}}%
\pgfpathlineto{\pgfqpoint{1.647908in}{0.670831in}}%
\pgfpathlineto{\pgfqpoint{1.640690in}{0.665716in}}%
\pgfpathlineto{\pgfqpoint{1.695622in}{0.661844in}}%
\pgfpathlineto{\pgfqpoint{1.713553in}{0.660732in}}%
\pgfpathlineto{\pgfqpoint{1.710253in}{0.607672in}}%
\pgfpathlineto{\pgfqpoint{1.751098in}{0.605189in}}%
\pgfpathlineto{\pgfqpoint{1.748728in}{0.562219in}}%
\pgfpathclose%
\pgfusepath{fill}%
\end{pgfscope}%
\begin{pgfscope}%
\pgfpathrectangle{\pgfqpoint{0.100000in}{0.100000in}}{\pgfqpoint{3.608454in}{2.310000in}}%
\pgfusepath{clip}%
\pgfsetbuttcap%
\pgfsetmiterjoin%
\definecolor{currentfill}{rgb}{0.000000,0.654902,0.672549}%
\pgfsetfillcolor{currentfill}%
\pgfsetlinewidth{0.000000pt}%
\definecolor{currentstroke}{rgb}{0.000000,0.000000,0.000000}%
\pgfsetstrokecolor{currentstroke}%
\pgfsetstrokeopacity{0.000000}%
\pgfsetdash{}{0pt}%
\pgfpathmoveto{\pgfqpoint{3.051009in}{1.477958in}}%
\pgfpathlineto{\pgfqpoint{3.051273in}{1.485944in}}%
\pgfpathlineto{\pgfqpoint{3.048216in}{1.491838in}}%
\pgfpathlineto{\pgfqpoint{3.054176in}{1.507940in}}%
\pgfpathlineto{\pgfqpoint{3.061504in}{1.519627in}}%
\pgfpathlineto{\pgfqpoint{3.068856in}{1.543568in}}%
\pgfpathlineto{\pgfqpoint{3.072318in}{1.562246in}}%
\pgfpathlineto{\pgfqpoint{3.098915in}{1.567130in}}%
\pgfpathlineto{\pgfqpoint{3.107213in}{1.564628in}}%
\pgfpathlineto{\pgfqpoint{3.111267in}{1.560885in}}%
\pgfpathlineto{\pgfqpoint{3.110368in}{1.555879in}}%
\pgfpathlineto{\pgfqpoint{3.115439in}{1.550679in}}%
\pgfpathlineto{\pgfqpoint{3.110584in}{1.534960in}}%
\pgfpathlineto{\pgfqpoint{3.112902in}{1.529498in}}%
\pgfpathlineto{\pgfqpoint{3.119792in}{1.526024in}}%
\pgfpathlineto{\pgfqpoint{3.116662in}{1.519085in}}%
\pgfpathlineto{\pgfqpoint{3.117244in}{1.511717in}}%
\pgfpathlineto{\pgfqpoint{3.114952in}{1.497798in}}%
\pgfpathlineto{\pgfqpoint{3.111561in}{1.488897in}}%
\pgfpathlineto{\pgfqpoint{3.051009in}{1.477958in}}%
\pgfpathclose%
\pgfusepath{fill}%
\end{pgfscope}%
\begin{pgfscope}%
\pgfpathrectangle{\pgfqpoint{0.100000in}{0.100000in}}{\pgfqpoint{3.608454in}{2.310000in}}%
\pgfusepath{clip}%
\pgfsetbuttcap%
\pgfsetmiterjoin%
\definecolor{currentfill}{rgb}{0.000000,0.400000,0.800000}%
\pgfsetfillcolor{currentfill}%
\pgfsetlinewidth{0.000000pt}%
\definecolor{currentstroke}{rgb}{0.000000,0.000000,0.000000}%
\pgfsetstrokecolor{currentstroke}%
\pgfsetstrokeopacity{0.000000}%
\pgfsetdash{}{0pt}%
\pgfpathmoveto{\pgfqpoint{2.335984in}{1.355415in}}%
\pgfpathlineto{\pgfqpoint{2.336652in}{1.332026in}}%
\pgfpathlineto{\pgfqpoint{2.327371in}{1.331768in}}%
\pgfpathlineto{\pgfqpoint{2.327621in}{1.321058in}}%
\pgfpathlineto{\pgfqpoint{2.319587in}{1.318096in}}%
\pgfpathlineto{\pgfqpoint{2.313973in}{1.320063in}}%
\pgfpathlineto{\pgfqpoint{2.313325in}{1.348137in}}%
\pgfpathlineto{\pgfqpoint{2.284099in}{1.347536in}}%
\pgfpathlineto{\pgfqpoint{2.284008in}{1.361432in}}%
\pgfpathlineto{\pgfqpoint{2.271150in}{1.361570in}}%
\pgfpathlineto{\pgfqpoint{2.271164in}{1.369568in}}%
\pgfpathlineto{\pgfqpoint{2.307376in}{1.370072in}}%
\pgfpathlineto{\pgfqpoint{2.324527in}{1.369104in}}%
\pgfpathlineto{\pgfqpoint{2.326861in}{1.355261in}}%
\pgfpathlineto{\pgfqpoint{2.335984in}{1.355415in}}%
\pgfpathclose%
\pgfusepath{fill}%
\end{pgfscope}%
\begin{pgfscope}%
\pgfpathrectangle{\pgfqpoint{0.100000in}{0.100000in}}{\pgfqpoint{3.608454in}{2.310000in}}%
\pgfusepath{clip}%
\pgfsetbuttcap%
\pgfsetmiterjoin%
\definecolor{currentfill}{rgb}{0.000000,0.537255,0.731373}%
\pgfsetfillcolor{currentfill}%
\pgfsetlinewidth{0.000000pt}%
\definecolor{currentstroke}{rgb}{0.000000,0.000000,0.000000}%
\pgfsetstrokecolor{currentstroke}%
\pgfsetstrokeopacity{0.000000}%
\pgfsetdash{}{0pt}%
\pgfpathmoveto{\pgfqpoint{2.182439in}{1.381621in}}%
\pgfpathlineto{\pgfqpoint{2.155044in}{1.381527in}}%
\pgfpathlineto{\pgfqpoint{2.154965in}{1.375795in}}%
\pgfpathlineto{\pgfqpoint{2.131152in}{1.375734in}}%
\pgfpathlineto{\pgfqpoint{2.131100in}{1.381764in}}%
\pgfpathlineto{\pgfqpoint{2.100521in}{1.381808in}}%
\pgfpathlineto{\pgfqpoint{2.103446in}{1.388933in}}%
\pgfpathlineto{\pgfqpoint{2.099783in}{1.391231in}}%
\pgfpathlineto{\pgfqpoint{2.086114in}{1.391298in}}%
\pgfpathlineto{\pgfqpoint{2.086211in}{1.418729in}}%
\pgfpathlineto{\pgfqpoint{2.088125in}{1.418722in}}%
\pgfpathlineto{\pgfqpoint{2.094548in}{1.411120in}}%
\pgfpathlineto{\pgfqpoint{2.101635in}{1.407776in}}%
\pgfpathlineto{\pgfqpoint{2.107194in}{1.411974in}}%
\pgfpathlineto{\pgfqpoint{2.104194in}{1.428993in}}%
\pgfpathlineto{\pgfqpoint{2.130612in}{1.428716in}}%
\pgfpathlineto{\pgfqpoint{2.130676in}{1.421843in}}%
\pgfpathlineto{\pgfqpoint{2.154074in}{1.421652in}}%
\pgfpathlineto{\pgfqpoint{2.154147in}{1.429674in}}%
\pgfpathlineto{\pgfqpoint{2.181491in}{1.429705in}}%
\pgfpathlineto{\pgfqpoint{2.181819in}{1.416001in}}%
\pgfpathlineto{\pgfqpoint{2.182439in}{1.381621in}}%
\pgfpathclose%
\pgfusepath{fill}%
\end{pgfscope}%
\begin{pgfscope}%
\pgfpathrectangle{\pgfqpoint{0.100000in}{0.100000in}}{\pgfqpoint{3.608454in}{2.310000in}}%
\pgfusepath{clip}%
\pgfsetbuttcap%
\pgfsetmiterjoin%
\definecolor{currentfill}{rgb}{0.000000,0.521569,0.739216}%
\pgfsetfillcolor{currentfill}%
\pgfsetlinewidth{0.000000pt}%
\definecolor{currentstroke}{rgb}{0.000000,0.000000,0.000000}%
\pgfsetstrokecolor{currentstroke}%
\pgfsetstrokeopacity{0.000000}%
\pgfsetdash{}{0pt}%
\pgfpathmoveto{\pgfqpoint{2.461626in}{1.152803in}}%
\pgfpathlineto{\pgfqpoint{2.457328in}{1.152433in}}%
\pgfpathlineto{\pgfqpoint{2.453898in}{1.152266in}}%
\pgfpathlineto{\pgfqpoint{2.454276in}{1.139497in}}%
\pgfpathlineto{\pgfqpoint{2.450987in}{1.144593in}}%
\pgfpathlineto{\pgfqpoint{2.440887in}{1.142638in}}%
\pgfpathlineto{\pgfqpoint{2.426853in}{1.142086in}}%
\pgfpathlineto{\pgfqpoint{2.426666in}{1.161074in}}%
\pgfpathlineto{\pgfqpoint{2.414649in}{1.160533in}}%
\pgfpathlineto{\pgfqpoint{2.413654in}{1.167791in}}%
\pgfpathlineto{\pgfqpoint{2.406457in}{1.183183in}}%
\pgfpathlineto{\pgfqpoint{2.415125in}{1.196618in}}%
\pgfpathlineto{\pgfqpoint{2.408011in}{1.196285in}}%
\pgfpathlineto{\pgfqpoint{2.407413in}{1.214350in}}%
\pgfpathlineto{\pgfqpoint{2.411944in}{1.214363in}}%
\pgfpathlineto{\pgfqpoint{2.410673in}{1.236890in}}%
\pgfpathlineto{\pgfqpoint{2.444353in}{1.238826in}}%
\pgfpathlineto{\pgfqpoint{2.449905in}{1.236650in}}%
\pgfpathlineto{\pgfqpoint{2.457038in}{1.223290in}}%
\pgfpathlineto{\pgfqpoint{2.451445in}{1.214650in}}%
\pgfpathlineto{\pgfqpoint{2.460932in}{1.196917in}}%
\pgfpathlineto{\pgfqpoint{2.467123in}{1.192232in}}%
\pgfpathlineto{\pgfqpoint{2.472760in}{1.194489in}}%
\pgfpathlineto{\pgfqpoint{2.477226in}{1.192084in}}%
\pgfpathlineto{\pgfqpoint{2.479456in}{1.190579in}}%
\pgfpathlineto{\pgfqpoint{2.477580in}{1.181419in}}%
\pgfpathlineto{\pgfqpoint{2.478809in}{1.176580in}}%
\pgfpathlineto{\pgfqpoint{2.474357in}{1.170803in}}%
\pgfpathlineto{\pgfqpoint{2.477166in}{1.167044in}}%
\pgfpathlineto{\pgfqpoint{2.473332in}{1.159111in}}%
\pgfpathlineto{\pgfqpoint{2.464384in}{1.163029in}}%
\pgfpathlineto{\pgfqpoint{2.461626in}{1.152803in}}%
\pgfpathclose%
\pgfusepath{fill}%
\end{pgfscope}%
\begin{pgfscope}%
\pgfpathrectangle{\pgfqpoint{0.100000in}{0.100000in}}{\pgfqpoint{3.608454in}{2.310000in}}%
\pgfusepath{clip}%
\pgfsetbuttcap%
\pgfsetmiterjoin%
\definecolor{currentfill}{rgb}{0.000000,0.686275,0.656863}%
\pgfsetfillcolor{currentfill}%
\pgfsetlinewidth{0.000000pt}%
\definecolor{currentstroke}{rgb}{0.000000,0.000000,0.000000}%
\pgfsetstrokecolor{currentstroke}%
\pgfsetstrokeopacity{0.000000}%
\pgfsetdash{}{0pt}%
\pgfpathmoveto{\pgfqpoint{2.547754in}{0.792351in}}%
\pgfpathlineto{\pgfqpoint{2.549469in}{0.776821in}}%
\pgfpathlineto{\pgfqpoint{2.551986in}{0.756063in}}%
\pgfpathlineto{\pgfqpoint{2.525381in}{0.754028in}}%
\pgfpathlineto{\pgfqpoint{2.487539in}{0.751603in}}%
\pgfpathlineto{\pgfqpoint{2.485619in}{0.780455in}}%
\pgfpathlineto{\pgfqpoint{2.468672in}{0.778139in}}%
\pgfpathlineto{\pgfqpoint{2.467492in}{0.796662in}}%
\pgfpathlineto{\pgfqpoint{2.462239in}{0.799019in}}%
\pgfpathlineto{\pgfqpoint{2.461411in}{0.812905in}}%
\pgfpathlineto{\pgfqpoint{2.488805in}{0.814690in}}%
\pgfpathlineto{\pgfqpoint{2.516272in}{0.816527in}}%
\pgfpathlineto{\pgfqpoint{2.518619in}{0.784958in}}%
\pgfpathlineto{\pgfqpoint{2.536679in}{0.788374in}}%
\pgfpathlineto{\pgfqpoint{2.547754in}{0.792351in}}%
\pgfpathclose%
\pgfusepath{fill}%
\end{pgfscope}%
\begin{pgfscope}%
\pgfpathrectangle{\pgfqpoint{0.100000in}{0.100000in}}{\pgfqpoint{3.608454in}{2.310000in}}%
\pgfusepath{clip}%
\pgfsetbuttcap%
\pgfsetmiterjoin%
\definecolor{currentfill}{rgb}{0.000000,0.611765,0.694118}%
\pgfsetfillcolor{currentfill}%
\pgfsetlinewidth{0.000000pt}%
\definecolor{currentstroke}{rgb}{0.000000,0.000000,0.000000}%
\pgfsetstrokecolor{currentstroke}%
\pgfsetstrokeopacity{0.000000}%
\pgfsetdash{}{0pt}%
\pgfpathmoveto{\pgfqpoint{1.557865in}{2.068157in}}%
\pgfpathlineto{\pgfqpoint{1.551007in}{2.068911in}}%
\pgfpathlineto{\pgfqpoint{1.548607in}{2.048174in}}%
\pgfpathlineto{\pgfqpoint{1.546296in}{2.048425in}}%
\pgfpathlineto{\pgfqpoint{1.543137in}{2.020826in}}%
\pgfpathlineto{\pgfqpoint{1.520131in}{2.023482in}}%
\pgfpathlineto{\pgfqpoint{1.519298in}{2.016515in}}%
\pgfpathlineto{\pgfqpoint{1.512427in}{2.017326in}}%
\pgfpathlineto{\pgfqpoint{1.513276in}{2.024315in}}%
\pgfpathlineto{\pgfqpoint{1.506437in}{2.025137in}}%
\pgfpathlineto{\pgfqpoint{1.501940in}{2.025689in}}%
\pgfpathlineto{\pgfqpoint{1.503629in}{2.039468in}}%
\pgfpathlineto{\pgfqpoint{1.496777in}{2.040294in}}%
\pgfpathlineto{\pgfqpoint{1.500521in}{2.053891in}}%
\pgfpathlineto{\pgfqpoint{1.503848in}{2.080824in}}%
\pgfpathlineto{\pgfqpoint{1.498672in}{2.081480in}}%
\pgfpathlineto{\pgfqpoint{1.496354in}{2.088951in}}%
\pgfpathlineto{\pgfqpoint{1.498726in}{2.097310in}}%
\pgfpathlineto{\pgfqpoint{1.507101in}{2.093171in}}%
\pgfpathlineto{\pgfqpoint{1.527435in}{2.089391in}}%
\pgfpathlineto{\pgfqpoint{1.529442in}{2.105984in}}%
\pgfpathlineto{\pgfqpoint{1.531517in}{2.105747in}}%
\pgfpathlineto{\pgfqpoint{1.534783in}{2.133134in}}%
\pgfpathlineto{\pgfqpoint{1.578597in}{2.128096in}}%
\pgfpathlineto{\pgfqpoint{1.589978in}{2.126848in}}%
\pgfpathlineto{\pgfqpoint{1.589220in}{2.119949in}}%
\pgfpathlineto{\pgfqpoint{1.596092in}{2.119196in}}%
\pgfpathlineto{\pgfqpoint{1.595340in}{2.112255in}}%
\pgfpathlineto{\pgfqpoint{1.626106in}{2.109098in}}%
\pgfpathlineto{\pgfqpoint{1.623124in}{2.077735in}}%
\pgfpathlineto{\pgfqpoint{1.619223in}{2.081466in}}%
\pgfpathlineto{\pgfqpoint{1.607083in}{2.083595in}}%
\pgfpathlineto{\pgfqpoint{1.601867in}{2.089959in}}%
\pgfpathlineto{\pgfqpoint{1.590181in}{2.089744in}}%
\pgfpathlineto{\pgfqpoint{1.582871in}{2.091842in}}%
\pgfpathlineto{\pgfqpoint{1.569948in}{2.088334in}}%
\pgfpathlineto{\pgfqpoint{1.562458in}{2.089780in}}%
\pgfpathlineto{\pgfqpoint{1.560678in}{2.074073in}}%
\pgfpathlineto{\pgfqpoint{1.557865in}{2.068157in}}%
\pgfpathclose%
\pgfusepath{fill}%
\end{pgfscope}%
\begin{pgfscope}%
\pgfpathrectangle{\pgfqpoint{0.100000in}{0.100000in}}{\pgfqpoint{3.608454in}{2.310000in}}%
\pgfusepath{clip}%
\pgfsetbuttcap%
\pgfsetmiterjoin%
\definecolor{currentfill}{rgb}{0.000000,0.854902,0.572549}%
\pgfsetfillcolor{currentfill}%
\pgfsetlinewidth{0.000000pt}%
\definecolor{currentstroke}{rgb}{0.000000,0.000000,0.000000}%
\pgfsetstrokecolor{currentstroke}%
\pgfsetstrokeopacity{0.000000}%
\pgfsetdash{}{0pt}%
\pgfpathmoveto{\pgfqpoint{1.795506in}{0.488863in}}%
\pgfpathlineto{\pgfqpoint{1.744553in}{0.490983in}}%
\pgfpathlineto{\pgfqpoint{1.748728in}{0.562219in}}%
\pgfpathlineto{\pgfqpoint{1.797569in}{0.560061in}}%
\pgfpathlineto{\pgfqpoint{1.796154in}{0.523908in}}%
\pgfpathlineto{\pgfqpoint{1.797091in}{0.523880in}}%
\pgfpathlineto{\pgfqpoint{1.795506in}{0.488863in}}%
\pgfpathclose%
\pgfusepath{fill}%
\end{pgfscope}%
\begin{pgfscope}%
\pgfpathrectangle{\pgfqpoint{0.100000in}{0.100000in}}{\pgfqpoint{3.608454in}{2.310000in}}%
\pgfusepath{clip}%
\pgfsetbuttcap%
\pgfsetmiterjoin%
\definecolor{currentfill}{rgb}{0.000000,0.666667,0.666667}%
\pgfsetfillcolor{currentfill}%
\pgfsetlinewidth{0.000000pt}%
\definecolor{currentstroke}{rgb}{0.000000,0.000000,0.000000}%
\pgfsetstrokecolor{currentstroke}%
\pgfsetstrokeopacity{0.000000}%
\pgfsetdash{}{0pt}%
\pgfpathmoveto{\pgfqpoint{1.443108in}{0.820998in}}%
\pgfpathlineto{\pgfqpoint{1.498395in}{0.815195in}}%
\pgfpathlineto{\pgfqpoint{1.486032in}{0.744751in}}%
\pgfpathlineto{\pgfqpoint{1.426837in}{0.715528in}}%
\pgfpathlineto{\pgfqpoint{1.429579in}{0.734241in}}%
\pgfpathlineto{\pgfqpoint{1.435480in}{0.791748in}}%
\pgfpathlineto{\pgfqpoint{1.438374in}{0.821513in}}%
\pgfpathlineto{\pgfqpoint{1.443108in}{0.820998in}}%
\pgfpathclose%
\pgfusepath{fill}%
\end{pgfscope}%
\begin{pgfscope}%
\pgfpathrectangle{\pgfqpoint{0.100000in}{0.100000in}}{\pgfqpoint{3.608454in}{2.310000in}}%
\pgfusepath{clip}%
\pgfsetbuttcap%
\pgfsetmiterjoin%
\definecolor{currentfill}{rgb}{0.000000,0.721569,0.639216}%
\pgfsetfillcolor{currentfill}%
\pgfsetlinewidth{0.000000pt}%
\definecolor{currentstroke}{rgb}{0.000000,0.000000,0.000000}%
\pgfsetstrokecolor{currentstroke}%
\pgfsetstrokeopacity{0.000000}%
\pgfsetdash{}{0pt}%
\pgfpathmoveto{\pgfqpoint{1.596724in}{1.204258in}}%
\pgfpathlineto{\pgfqpoint{1.602002in}{1.203777in}}%
\pgfpathlineto{\pgfqpoint{1.598511in}{1.164402in}}%
\pgfpathlineto{\pgfqpoint{1.596009in}{1.164627in}}%
\pgfpathlineto{\pgfqpoint{1.590719in}{1.104616in}}%
\pgfpathlineto{\pgfqpoint{1.569376in}{1.106608in}}%
\pgfpathlineto{\pgfqpoint{1.572537in}{1.133851in}}%
\pgfpathlineto{\pgfqpoint{1.545193in}{1.136393in}}%
\pgfpathlineto{\pgfqpoint{1.545891in}{1.143339in}}%
\pgfpathlineto{\pgfqpoint{1.532225in}{1.144723in}}%
\pgfpathlineto{\pgfqpoint{1.532570in}{1.148149in}}%
\pgfpathlineto{\pgfqpoint{1.538731in}{1.209408in}}%
\pgfpathlineto{\pgfqpoint{1.596724in}{1.204258in}}%
\pgfpathclose%
\pgfusepath{fill}%
\end{pgfscope}%
\begin{pgfscope}%
\pgfpathrectangle{\pgfqpoint{0.100000in}{0.100000in}}{\pgfqpoint{3.608454in}{2.310000in}}%
\pgfusepath{clip}%
\pgfsetbuttcap%
\pgfsetmiterjoin%
\definecolor{currentfill}{rgb}{0.000000,0.827451,0.586275}%
\pgfsetfillcolor{currentfill}%
\pgfsetlinewidth{0.000000pt}%
\definecolor{currentstroke}{rgb}{0.000000,0.000000,0.000000}%
\pgfsetstrokecolor{currentstroke}%
\pgfsetstrokeopacity{0.000000}%
\pgfsetdash{}{0pt}%
\pgfpathmoveto{\pgfqpoint{0.655369in}{0.626110in}}%
\pgfpathlineto{\pgfqpoint{0.652280in}{0.629547in}}%
\pgfpathlineto{\pgfqpoint{0.652480in}{0.635784in}}%
\pgfpathlineto{\pgfqpoint{0.651564in}{0.638676in}}%
\pgfpathlineto{\pgfqpoint{0.654993in}{0.639490in}}%
\pgfpathlineto{\pgfqpoint{0.653508in}{0.635458in}}%
\pgfpathlineto{\pgfqpoint{0.653829in}{0.632282in}}%
\pgfpathlineto{\pgfqpoint{0.656074in}{0.628020in}}%
\pgfpathlineto{\pgfqpoint{0.655369in}{0.626110in}}%
\pgfpathclose%
\pgfusepath{fill}%
\end{pgfscope}%
\begin{pgfscope}%
\pgfpathrectangle{\pgfqpoint{0.100000in}{0.100000in}}{\pgfqpoint{3.608454in}{2.310000in}}%
\pgfusepath{clip}%
\pgfsetbuttcap%
\pgfsetmiterjoin%
\definecolor{currentfill}{rgb}{0.000000,0.827451,0.586275}%
\pgfsetfillcolor{currentfill}%
\pgfsetlinewidth{0.000000pt}%
\definecolor{currentstroke}{rgb}{0.000000,0.000000,0.000000}%
\pgfsetstrokecolor{currentstroke}%
\pgfsetstrokeopacity{0.000000}%
\pgfsetdash{}{0pt}%
\pgfpathmoveto{\pgfqpoint{0.705347in}{0.551735in}}%
\pgfpathlineto{\pgfqpoint{0.705017in}{0.553533in}}%
\pgfpathlineto{\pgfqpoint{0.702425in}{0.555599in}}%
\pgfpathlineto{\pgfqpoint{0.701422in}{0.558633in}}%
\pgfpathlineto{\pgfqpoint{0.701257in}{0.561453in}}%
\pgfpathlineto{\pgfqpoint{0.699278in}{0.565095in}}%
\pgfpathlineto{\pgfqpoint{0.699632in}{0.569687in}}%
\pgfpathlineto{\pgfqpoint{0.700852in}{0.572122in}}%
\pgfpathlineto{\pgfqpoint{0.702514in}{0.571731in}}%
\pgfpathlineto{\pgfqpoint{0.706248in}{0.566407in}}%
\pgfpathlineto{\pgfqpoint{0.710913in}{0.566989in}}%
\pgfpathlineto{\pgfqpoint{0.713889in}{0.566223in}}%
\pgfpathlineto{\pgfqpoint{0.714374in}{0.564683in}}%
\pgfpathlineto{\pgfqpoint{0.716943in}{0.563500in}}%
\pgfpathlineto{\pgfqpoint{0.716755in}{0.561123in}}%
\pgfpathlineto{\pgfqpoint{0.719375in}{0.558723in}}%
\pgfpathlineto{\pgfqpoint{0.719272in}{0.555594in}}%
\pgfpathlineto{\pgfqpoint{0.717154in}{0.553156in}}%
\pgfpathlineto{\pgfqpoint{0.715054in}{0.551897in}}%
\pgfpathlineto{\pgfqpoint{0.714314in}{0.546900in}}%
\pgfpathlineto{\pgfqpoint{0.711935in}{0.547634in}}%
\pgfpathlineto{\pgfqpoint{0.708937in}{0.550219in}}%
\pgfpathlineto{\pgfqpoint{0.705347in}{0.551735in}}%
\pgfpathclose%
\pgfusepath{fill}%
\end{pgfscope}%
\begin{pgfscope}%
\pgfpathrectangle{\pgfqpoint{0.100000in}{0.100000in}}{\pgfqpoint{3.608454in}{2.310000in}}%
\pgfusepath{clip}%
\pgfsetbuttcap%
\pgfsetmiterjoin%
\definecolor{currentfill}{rgb}{0.000000,0.827451,0.586275}%
\pgfsetfillcolor{currentfill}%
\pgfsetlinewidth{0.000000pt}%
\definecolor{currentstroke}{rgb}{0.000000,0.000000,0.000000}%
\pgfsetstrokecolor{currentstroke}%
\pgfsetstrokeopacity{0.000000}%
\pgfsetdash{}{0pt}%
\pgfpathmoveto{\pgfqpoint{0.894650in}{0.485392in}}%
\pgfpathlineto{\pgfqpoint{0.878747in}{0.464023in}}%
\pgfpathlineto{\pgfqpoint{0.878374in}{0.466226in}}%
\pgfpathlineto{\pgfqpoint{0.874246in}{0.469508in}}%
\pgfpathlineto{\pgfqpoint{0.870359in}{0.464352in}}%
\pgfpathlineto{\pgfqpoint{0.869217in}{0.461424in}}%
\pgfpathlineto{\pgfqpoint{0.864761in}{0.455531in}}%
\pgfpathlineto{\pgfqpoint{0.853555in}{0.463966in}}%
\pgfpathlineto{\pgfqpoint{0.836300in}{0.477446in}}%
\pgfpathlineto{\pgfqpoint{0.822897in}{0.488128in}}%
\pgfpathlineto{\pgfqpoint{0.824629in}{0.490280in}}%
\pgfpathlineto{\pgfqpoint{0.820776in}{0.493434in}}%
\pgfpathlineto{\pgfqpoint{0.819027in}{0.491316in}}%
\pgfpathlineto{\pgfqpoint{0.816987in}{0.492818in}}%
\pgfpathlineto{\pgfqpoint{0.815390in}{0.490823in}}%
\pgfpathlineto{\pgfqpoint{0.813395in}{0.492489in}}%
\pgfpathlineto{\pgfqpoint{0.814957in}{0.494457in}}%
\pgfpathlineto{\pgfqpoint{0.804074in}{0.503387in}}%
\pgfpathlineto{\pgfqpoint{0.802432in}{0.501473in}}%
\pgfpathlineto{\pgfqpoint{0.800941in}{0.502662in}}%
\pgfpathlineto{\pgfqpoint{0.799211in}{0.500636in}}%
\pgfpathlineto{\pgfqpoint{0.797356in}{0.502139in}}%
\pgfpathlineto{\pgfqpoint{0.795190in}{0.500666in}}%
\pgfpathlineto{\pgfqpoint{0.791902in}{0.496847in}}%
\pgfpathlineto{\pgfqpoint{0.789763in}{0.498705in}}%
\pgfpathlineto{\pgfqpoint{0.788112in}{0.496794in}}%
\pgfpathlineto{\pgfqpoint{0.786206in}{0.498438in}}%
\pgfpathlineto{\pgfqpoint{0.782945in}{0.494696in}}%
\pgfpathlineto{\pgfqpoint{0.781042in}{0.496290in}}%
\pgfpathlineto{\pgfqpoint{0.778697in}{0.494994in}}%
\pgfpathlineto{\pgfqpoint{0.775519in}{0.491429in}}%
\pgfpathlineto{\pgfqpoint{0.773696in}{0.492989in}}%
\pgfpathlineto{\pgfqpoint{0.770402in}{0.489224in}}%
\pgfpathlineto{\pgfqpoint{0.767845in}{0.491365in}}%
\pgfpathlineto{\pgfqpoint{0.764542in}{0.487613in}}%
\pgfpathlineto{\pgfqpoint{0.762636in}{0.489281in}}%
\pgfpathlineto{\pgfqpoint{0.759307in}{0.485494in}}%
\pgfpathlineto{\pgfqpoint{0.756995in}{0.484116in}}%
\pgfpathlineto{\pgfqpoint{0.755169in}{0.485784in}}%
\pgfpathlineto{\pgfqpoint{0.752873in}{0.484455in}}%
\pgfpathlineto{\pgfqpoint{0.749547in}{0.480704in}}%
\pgfpathlineto{\pgfqpoint{0.747682in}{0.482304in}}%
\pgfpathlineto{\pgfqpoint{0.746173in}{0.480590in}}%
\pgfpathlineto{\pgfqpoint{0.744357in}{0.482168in}}%
\pgfpathlineto{\pgfqpoint{0.740722in}{0.478712in}}%
\pgfpathlineto{\pgfqpoint{0.740538in}{0.478732in}}%
\pgfpathlineto{\pgfqpoint{0.737273in}{0.481661in}}%
\pgfpathlineto{\pgfqpoint{0.735963in}{0.479492in}}%
\pgfpathlineto{\pgfqpoint{0.734892in}{0.478188in}}%
\pgfpathlineto{\pgfqpoint{0.727742in}{0.479830in}}%
\pgfpathlineto{\pgfqpoint{0.728503in}{0.483239in}}%
\pgfpathlineto{\pgfqpoint{0.731793in}{0.484480in}}%
\pgfpathlineto{\pgfqpoint{0.735386in}{0.488777in}}%
\pgfpathlineto{\pgfqpoint{0.736702in}{0.491209in}}%
\pgfpathlineto{\pgfqpoint{0.736725in}{0.494624in}}%
\pgfpathlineto{\pgfqpoint{0.737998in}{0.497489in}}%
\pgfpathlineto{\pgfqpoint{0.739339in}{0.496648in}}%
\pgfpathlineto{\pgfqpoint{0.741092in}{0.498876in}}%
\pgfpathlineto{\pgfqpoint{0.745616in}{0.498996in}}%
\pgfpathlineto{\pgfqpoint{0.746961in}{0.501426in}}%
\pgfpathlineto{\pgfqpoint{0.748923in}{0.507053in}}%
\pgfpathlineto{\pgfqpoint{0.749487in}{0.509973in}}%
\pgfpathlineto{\pgfqpoint{0.750888in}{0.512260in}}%
\pgfpathlineto{\pgfqpoint{0.751044in}{0.515415in}}%
\pgfpathlineto{\pgfqpoint{0.754037in}{0.517140in}}%
\pgfpathlineto{\pgfqpoint{0.753941in}{0.521042in}}%
\pgfpathlineto{\pgfqpoint{0.752010in}{0.522096in}}%
\pgfpathlineto{\pgfqpoint{0.748553in}{0.518499in}}%
\pgfpathlineto{\pgfqpoint{0.745838in}{0.518264in}}%
\pgfpathlineto{\pgfqpoint{0.744687in}{0.519881in}}%
\pgfpathlineto{\pgfqpoint{0.738995in}{0.520517in}}%
\pgfpathlineto{\pgfqpoint{0.733381in}{0.523327in}}%
\pgfpathlineto{\pgfqpoint{0.730861in}{0.525258in}}%
\pgfpathlineto{\pgfqpoint{0.728015in}{0.528863in}}%
\pgfpathlineto{\pgfqpoint{0.727312in}{0.530894in}}%
\pgfpathlineto{\pgfqpoint{0.728494in}{0.533400in}}%
\pgfpathlineto{\pgfqpoint{0.730236in}{0.534212in}}%
\pgfpathlineto{\pgfqpoint{0.729733in}{0.537152in}}%
\pgfpathlineto{\pgfqpoint{0.730007in}{0.540442in}}%
\pgfpathlineto{\pgfqpoint{0.731612in}{0.540941in}}%
\pgfpathlineto{\pgfqpoint{0.730213in}{0.545263in}}%
\pgfpathlineto{\pgfqpoint{0.731563in}{0.546123in}}%
\pgfpathlineto{\pgfqpoint{0.728919in}{0.548253in}}%
\pgfpathlineto{\pgfqpoint{0.728522in}{0.551246in}}%
\pgfpathlineto{\pgfqpoint{0.729262in}{0.553763in}}%
\pgfpathlineto{\pgfqpoint{0.731148in}{0.553340in}}%
\pgfpathlineto{\pgfqpoint{0.731806in}{0.555149in}}%
\pgfpathlineto{\pgfqpoint{0.729059in}{0.555932in}}%
\pgfpathlineto{\pgfqpoint{0.728087in}{0.558977in}}%
\pgfpathlineto{\pgfqpoint{0.731970in}{0.558050in}}%
\pgfpathlineto{\pgfqpoint{0.736704in}{0.560032in}}%
\pgfpathlineto{\pgfqpoint{0.739886in}{0.560095in}}%
\pgfpathlineto{\pgfqpoint{0.738328in}{0.563808in}}%
\pgfpathlineto{\pgfqpoint{0.739124in}{0.565657in}}%
\pgfpathlineto{\pgfqpoint{0.741979in}{0.566643in}}%
\pgfpathlineto{\pgfqpoint{0.742241in}{0.570962in}}%
\pgfpathlineto{\pgfqpoint{0.739760in}{0.569380in}}%
\pgfpathlineto{\pgfqpoint{0.738485in}{0.573179in}}%
\pgfpathlineto{\pgfqpoint{0.741686in}{0.576289in}}%
\pgfpathlineto{\pgfqpoint{0.739656in}{0.579992in}}%
\pgfpathlineto{\pgfqpoint{0.741135in}{0.581696in}}%
\pgfpathlineto{\pgfqpoint{0.743178in}{0.581393in}}%
\pgfpathlineto{\pgfqpoint{0.744280in}{0.583186in}}%
\pgfpathlineto{\pgfqpoint{0.743630in}{0.585155in}}%
\pgfpathlineto{\pgfqpoint{0.740303in}{0.585479in}}%
\pgfpathlineto{\pgfqpoint{0.741843in}{0.588299in}}%
\pgfpathlineto{\pgfqpoint{0.746374in}{0.589014in}}%
\pgfpathlineto{\pgfqpoint{0.747420in}{0.592192in}}%
\pgfpathlineto{\pgfqpoint{0.750741in}{0.589193in}}%
\pgfpathlineto{\pgfqpoint{0.753431in}{0.589913in}}%
\pgfpathlineto{\pgfqpoint{0.754707in}{0.592801in}}%
\pgfpathlineto{\pgfqpoint{0.757034in}{0.594176in}}%
\pgfpathlineto{\pgfqpoint{0.767817in}{0.596834in}}%
\pgfpathlineto{\pgfqpoint{0.772792in}{0.597266in}}%
\pgfpathlineto{\pgfqpoint{0.774461in}{0.595375in}}%
\pgfpathlineto{\pgfqpoint{0.777267in}{0.596507in}}%
\pgfpathlineto{\pgfqpoint{0.778048in}{0.599219in}}%
\pgfpathlineto{\pgfqpoint{0.783684in}{0.603781in}}%
\pgfpathlineto{\pgfqpoint{0.789079in}{0.605814in}}%
\pgfpathlineto{\pgfqpoint{0.794866in}{0.605498in}}%
\pgfpathlineto{\pgfqpoint{0.798355in}{0.602512in}}%
\pgfpathlineto{\pgfqpoint{0.799353in}{0.597971in}}%
\pgfpathlineto{\pgfqpoint{0.799751in}{0.593446in}}%
\pgfpathlineto{\pgfqpoint{0.801076in}{0.590791in}}%
\pgfpathlineto{\pgfqpoint{0.803852in}{0.589451in}}%
\pgfpathlineto{\pgfqpoint{0.808679in}{0.590732in}}%
\pgfpathlineto{\pgfqpoint{0.812082in}{0.590252in}}%
\pgfpathlineto{\pgfqpoint{0.814106in}{0.588275in}}%
\pgfpathlineto{\pgfqpoint{0.812396in}{0.586335in}}%
\pgfpathlineto{\pgfqpoint{0.814282in}{0.584612in}}%
\pgfpathlineto{\pgfqpoint{0.812542in}{0.582561in}}%
\pgfpathlineto{\pgfqpoint{0.814280in}{0.580959in}}%
\pgfpathlineto{\pgfqpoint{0.812726in}{0.579011in}}%
\pgfpathlineto{\pgfqpoint{0.824219in}{0.568707in}}%
\pgfpathlineto{\pgfqpoint{0.823025in}{0.569214in}}%
\pgfpathlineto{\pgfqpoint{0.820507in}{0.567425in}}%
\pgfpathlineto{\pgfqpoint{0.813592in}{0.559787in}}%
\pgfpathlineto{\pgfqpoint{0.813107in}{0.560166in}}%
\pgfpathlineto{\pgfqpoint{0.806268in}{0.552525in}}%
\pgfpathlineto{\pgfqpoint{0.807682in}{0.551257in}}%
\pgfpathlineto{\pgfqpoint{0.802699in}{0.545595in}}%
\pgfpathlineto{\pgfqpoint{0.812784in}{0.536599in}}%
\pgfpathlineto{\pgfqpoint{0.819842in}{0.530526in}}%
\pgfpathlineto{\pgfqpoint{0.821946in}{0.532976in}}%
\pgfpathlineto{\pgfqpoint{0.830011in}{0.526141in}}%
\pgfpathlineto{\pgfqpoint{0.831654in}{0.528086in}}%
\pgfpathlineto{\pgfqpoint{0.846928in}{0.515267in}}%
\pgfpathlineto{\pgfqpoint{0.845087in}{0.513012in}}%
\pgfpathlineto{\pgfqpoint{0.854046in}{0.505533in}}%
\pgfpathlineto{\pgfqpoint{0.867623in}{0.494883in}}%
\pgfpathlineto{\pgfqpoint{0.878680in}{0.486403in}}%
\pgfpathlineto{\pgfqpoint{0.880946in}{0.484915in}}%
\pgfpathlineto{\pgfqpoint{0.883274in}{0.487912in}}%
\pgfpathlineto{\pgfqpoint{0.887256in}{0.484755in}}%
\pgfpathlineto{\pgfqpoint{0.888822in}{0.486756in}}%
\pgfpathlineto{\pgfqpoint{0.892390in}{0.483966in}}%
\pgfpathlineto{\pgfqpoint{0.894650in}{0.485392in}}%
\pgfpathclose%
\pgfusepath{fill}%
\end{pgfscope}%
\begin{pgfscope}%
\pgfpathrectangle{\pgfqpoint{0.100000in}{0.100000in}}{\pgfqpoint{3.608454in}{2.310000in}}%
\pgfusepath{clip}%
\pgfsetbuttcap%
\pgfsetmiterjoin%
\definecolor{currentfill}{rgb}{0.000000,0.713725,0.643137}%
\pgfsetfillcolor{currentfill}%
\pgfsetlinewidth{0.000000pt}%
\definecolor{currentstroke}{rgb}{0.000000,0.000000,0.000000}%
\pgfsetstrokecolor{currentstroke}%
\pgfsetstrokeopacity{0.000000}%
\pgfsetdash{}{0pt}%
\pgfpathmoveto{\pgfqpoint{1.012082in}{1.949489in}}%
\pgfpathlineto{\pgfqpoint{1.009235in}{1.946422in}}%
\pgfpathlineto{\pgfqpoint{1.004445in}{1.925280in}}%
\pgfpathlineto{\pgfqpoint{0.989615in}{1.915702in}}%
\pgfpathlineto{\pgfqpoint{0.981235in}{1.920433in}}%
\pgfpathlineto{\pgfqpoint{0.976747in}{1.912867in}}%
\pgfpathlineto{\pgfqpoint{0.969772in}{1.911307in}}%
\pgfpathlineto{\pgfqpoint{0.967552in}{1.905305in}}%
\pgfpathlineto{\pgfqpoint{0.970494in}{1.901193in}}%
\pgfpathlineto{\pgfqpoint{0.969006in}{1.895451in}}%
\pgfpathlineto{\pgfqpoint{0.964170in}{1.890384in}}%
\pgfpathlineto{\pgfqpoint{0.953430in}{1.884971in}}%
\pgfpathlineto{\pgfqpoint{0.949827in}{1.885175in}}%
\pgfpathlineto{\pgfqpoint{0.917584in}{1.892808in}}%
\pgfpathlineto{\pgfqpoint{0.915455in}{1.886073in}}%
\pgfpathlineto{\pgfqpoint{0.909972in}{1.887658in}}%
\pgfpathlineto{\pgfqpoint{0.913150in}{1.901010in}}%
\pgfpathlineto{\pgfqpoint{0.916455in}{1.900234in}}%
\pgfpathlineto{\pgfqpoint{0.919645in}{1.913720in}}%
\pgfpathlineto{\pgfqpoint{0.910516in}{1.911417in}}%
\pgfpathlineto{\pgfqpoint{0.904574in}{1.912817in}}%
\pgfpathlineto{\pgfqpoint{0.898851in}{1.918021in}}%
\pgfpathlineto{\pgfqpoint{0.901098in}{1.927625in}}%
\pgfpathlineto{\pgfqpoint{0.896730in}{1.932620in}}%
\pgfpathlineto{\pgfqpoint{0.899991in}{1.946190in}}%
\pgfpathlineto{\pgfqpoint{0.884981in}{1.950005in}}%
\pgfpathlineto{\pgfqpoint{0.890375in}{1.956866in}}%
\pgfpathlineto{\pgfqpoint{0.891257in}{1.963356in}}%
\pgfpathlineto{\pgfqpoint{0.895815in}{1.966858in}}%
\pgfpathlineto{\pgfqpoint{0.899963in}{1.971006in}}%
\pgfpathlineto{\pgfqpoint{0.904646in}{1.980156in}}%
\pgfpathlineto{\pgfqpoint{0.923373in}{1.975654in}}%
\pgfpathlineto{\pgfqpoint{0.923969in}{1.964164in}}%
\pgfpathlineto{\pgfqpoint{0.931228in}{1.960671in}}%
\pgfpathlineto{\pgfqpoint{0.942731in}{1.965121in}}%
\pgfpathlineto{\pgfqpoint{1.012082in}{1.949489in}}%
\pgfpathclose%
\pgfusepath{fill}%
\end{pgfscope}%
\begin{pgfscope}%
\pgfpathrectangle{\pgfqpoint{0.100000in}{0.100000in}}{\pgfqpoint{3.608454in}{2.310000in}}%
\pgfusepath{clip}%
\pgfsetbuttcap%
\pgfsetmiterjoin%
\definecolor{currentfill}{rgb}{0.000000,0.709804,0.645098}%
\pgfsetfillcolor{currentfill}%
\pgfsetlinewidth{0.000000pt}%
\definecolor{currentstroke}{rgb}{0.000000,0.000000,0.000000}%
\pgfsetstrokecolor{currentstroke}%
\pgfsetstrokeopacity{0.000000}%
\pgfsetdash{}{0pt}%
\pgfpathmoveto{\pgfqpoint{2.477226in}{1.192084in}}%
\pgfpathlineto{\pgfqpoint{2.472760in}{1.194489in}}%
\pgfpathlineto{\pgfqpoint{2.467123in}{1.192232in}}%
\pgfpathlineto{\pgfqpoint{2.460932in}{1.196917in}}%
\pgfpathlineto{\pgfqpoint{2.451445in}{1.214650in}}%
\pgfpathlineto{\pgfqpoint{2.457038in}{1.223290in}}%
\pgfpathlineto{\pgfqpoint{2.449905in}{1.236650in}}%
\pgfpathlineto{\pgfqpoint{2.452467in}{1.239392in}}%
\pgfpathlineto{\pgfqpoint{2.449790in}{1.246442in}}%
\pgfpathlineto{\pgfqpoint{2.445368in}{1.248025in}}%
\pgfpathlineto{\pgfqpoint{2.434982in}{1.258000in}}%
\pgfpathlineto{\pgfqpoint{2.428315in}{1.262231in}}%
\pgfpathlineto{\pgfqpoint{2.421525in}{1.260176in}}%
\pgfpathlineto{\pgfqpoint{2.418250in}{1.266369in}}%
\pgfpathlineto{\pgfqpoint{2.409704in}{1.273645in}}%
\pgfpathlineto{\pgfqpoint{2.404880in}{1.275417in}}%
\pgfpathlineto{\pgfqpoint{2.415279in}{1.279749in}}%
\pgfpathlineto{\pgfqpoint{2.414877in}{1.286654in}}%
\pgfpathlineto{\pgfqpoint{2.435549in}{1.287521in}}%
\pgfpathlineto{\pgfqpoint{2.470034in}{1.289131in}}%
\pgfpathlineto{\pgfqpoint{2.470338in}{1.282183in}}%
\pgfpathlineto{\pgfqpoint{2.497937in}{1.284068in}}%
\pgfpathlineto{\pgfqpoint{2.499015in}{1.266823in}}%
\pgfpathlineto{\pgfqpoint{2.500587in}{1.242556in}}%
\pgfpathlineto{\pgfqpoint{2.501888in}{1.221865in}}%
\pgfpathlineto{\pgfqpoint{2.489480in}{1.220900in}}%
\pgfpathlineto{\pgfqpoint{2.488838in}{1.212212in}}%
\pgfpathlineto{\pgfqpoint{2.482512in}{1.210574in}}%
\pgfpathlineto{\pgfqpoint{2.473800in}{1.196498in}}%
\pgfpathlineto{\pgfqpoint{2.477226in}{1.192084in}}%
\pgfpathclose%
\pgfusepath{fill}%
\end{pgfscope}%
\begin{pgfscope}%
\pgfpathrectangle{\pgfqpoint{0.100000in}{0.100000in}}{\pgfqpoint{3.608454in}{2.310000in}}%
\pgfusepath{clip}%
\pgfsetbuttcap%
\pgfsetmiterjoin%
\definecolor{currentfill}{rgb}{0.000000,0.454902,0.772549}%
\pgfsetfillcolor{currentfill}%
\pgfsetlinewidth{0.000000pt}%
\definecolor{currentstroke}{rgb}{0.000000,0.000000,0.000000}%
\pgfsetstrokecolor{currentstroke}%
\pgfsetstrokeopacity{0.000000}%
\pgfsetdash{}{0pt}%
\pgfpathmoveto{\pgfqpoint{3.245459in}{1.515286in}}%
\pgfpathlineto{\pgfqpoint{3.239646in}{1.514041in}}%
\pgfpathlineto{\pgfqpoint{3.239245in}{1.513965in}}%
\pgfpathlineto{\pgfqpoint{3.227209in}{1.526852in}}%
\pgfpathlineto{\pgfqpoint{3.221777in}{1.528537in}}%
\pgfpathlineto{\pgfqpoint{3.216418in}{1.535851in}}%
\pgfpathlineto{\pgfqpoint{3.208552in}{1.535223in}}%
\pgfpathlineto{\pgfqpoint{3.204878in}{1.540940in}}%
\pgfpathlineto{\pgfqpoint{3.212062in}{1.546662in}}%
\pgfpathlineto{\pgfqpoint{3.201240in}{1.566916in}}%
\pgfpathlineto{\pgfqpoint{3.208359in}{1.574869in}}%
\pgfpathlineto{\pgfqpoint{3.196959in}{1.580835in}}%
\pgfpathlineto{\pgfqpoint{3.213947in}{1.593817in}}%
\pgfpathlineto{\pgfqpoint{3.217800in}{1.596744in}}%
\pgfpathlineto{\pgfqpoint{3.217800in}{1.603350in}}%
\pgfpathlineto{\pgfqpoint{3.221218in}{1.609411in}}%
\pgfpathlineto{\pgfqpoint{3.234122in}{1.609162in}}%
\pgfpathlineto{\pgfqpoint{3.251074in}{1.598426in}}%
\pgfpathlineto{\pgfqpoint{3.244201in}{1.592250in}}%
\pgfpathlineto{\pgfqpoint{3.269282in}{1.579037in}}%
\pgfpathlineto{\pgfqpoint{3.262873in}{1.561037in}}%
\pgfpathlineto{\pgfqpoint{3.250514in}{1.547606in}}%
\pgfpathlineto{\pgfqpoint{3.252158in}{1.542409in}}%
\pgfpathlineto{\pgfqpoint{3.250059in}{1.534993in}}%
\pgfpathlineto{\pgfqpoint{3.250934in}{1.527034in}}%
\pgfpathlineto{\pgfqpoint{3.248921in}{1.519116in}}%
\pgfpathlineto{\pgfqpoint{3.245459in}{1.515286in}}%
\pgfpathclose%
\pgfusepath{fill}%
\end{pgfscope}%
\begin{pgfscope}%
\pgfpathrectangle{\pgfqpoint{0.100000in}{0.100000in}}{\pgfqpoint{3.608454in}{2.310000in}}%
\pgfusepath{clip}%
\pgfsetbuttcap%
\pgfsetmiterjoin%
\definecolor{currentfill}{rgb}{0.000000,0.592157,0.703922}%
\pgfsetfillcolor{currentfill}%
\pgfsetlinewidth{0.000000pt}%
\definecolor{currentstroke}{rgb}{0.000000,0.000000,0.000000}%
\pgfsetstrokecolor{currentstroke}%
\pgfsetstrokeopacity{0.000000}%
\pgfsetdash{}{0pt}%
\pgfpathmoveto{\pgfqpoint{2.887245in}{0.726366in}}%
\pgfpathlineto{\pgfqpoint{2.878161in}{0.725802in}}%
\pgfpathlineto{\pgfqpoint{2.874929in}{0.755916in}}%
\pgfpathlineto{\pgfqpoint{2.856733in}{0.753994in}}%
\pgfpathlineto{\pgfqpoint{2.856403in}{0.756814in}}%
\pgfpathlineto{\pgfqpoint{2.854226in}{0.777256in}}%
\pgfpathlineto{\pgfqpoint{2.877599in}{0.779820in}}%
\pgfpathlineto{\pgfqpoint{2.875554in}{0.798633in}}%
\pgfpathlineto{\pgfqpoint{2.870385in}{0.802308in}}%
\pgfpathlineto{\pgfqpoint{2.865088in}{0.801722in}}%
\pgfpathlineto{\pgfqpoint{2.862925in}{0.815951in}}%
\pgfpathlineto{\pgfqpoint{2.875686in}{0.817596in}}%
\pgfpathlineto{\pgfqpoint{2.875215in}{0.821548in}}%
\pgfpathlineto{\pgfqpoint{2.884148in}{0.822135in}}%
\pgfpathlineto{\pgfqpoint{2.887379in}{0.811341in}}%
\pgfpathlineto{\pgfqpoint{2.884962in}{0.810094in}}%
\pgfpathlineto{\pgfqpoint{2.885535in}{0.800050in}}%
\pgfpathlineto{\pgfqpoint{2.891131in}{0.799434in}}%
\pgfpathlineto{\pgfqpoint{2.894796in}{0.794134in}}%
\pgfpathlineto{\pgfqpoint{2.910462in}{0.795395in}}%
\pgfpathlineto{\pgfqpoint{2.911798in}{0.789691in}}%
\pgfpathlineto{\pgfqpoint{2.920094in}{0.780812in}}%
\pgfpathlineto{\pgfqpoint{2.920160in}{0.773477in}}%
\pgfpathlineto{\pgfqpoint{2.925268in}{0.774204in}}%
\pgfpathlineto{\pgfqpoint{2.928578in}{0.749275in}}%
\pgfpathlineto{\pgfqpoint{2.937427in}{0.747697in}}%
\pgfpathlineto{\pgfqpoint{2.943923in}{0.740325in}}%
\pgfpathlineto{\pgfqpoint{2.955222in}{0.740323in}}%
\pgfpathlineto{\pgfqpoint{2.957900in}{0.730883in}}%
\pgfpathlineto{\pgfqpoint{2.887245in}{0.726366in}}%
\pgfpathclose%
\pgfusepath{fill}%
\end{pgfscope}%
\begin{pgfscope}%
\pgfpathrectangle{\pgfqpoint{0.100000in}{0.100000in}}{\pgfqpoint{3.608454in}{2.310000in}}%
\pgfusepath{clip}%
\pgfsetbuttcap%
\pgfsetmiterjoin%
\definecolor{currentfill}{rgb}{0.000000,0.827451,0.586275}%
\pgfsetfillcolor{currentfill}%
\pgfsetlinewidth{0.000000pt}%
\definecolor{currentstroke}{rgb}{0.000000,0.000000,0.000000}%
\pgfsetstrokecolor{currentstroke}%
\pgfsetstrokeopacity{0.000000}%
\pgfsetdash{}{0pt}%
\pgfpathmoveto{\pgfqpoint{2.416578in}{0.896534in}}%
\pgfpathlineto{\pgfqpoint{2.408616in}{0.893304in}}%
\pgfpathlineto{\pgfqpoint{2.408409in}{0.898240in}}%
\pgfpathlineto{\pgfqpoint{2.406598in}{0.936460in}}%
\pgfpathlineto{\pgfqpoint{2.424326in}{0.937210in}}%
\pgfpathlineto{\pgfqpoint{2.416509in}{0.931498in}}%
\pgfpathlineto{\pgfqpoint{2.427713in}{0.930421in}}%
\pgfpathlineto{\pgfqpoint{2.427361in}{0.937371in}}%
\pgfpathlineto{\pgfqpoint{2.441137in}{0.939280in}}%
\pgfpathlineto{\pgfqpoint{2.440822in}{0.945153in}}%
\pgfpathlineto{\pgfqpoint{2.451281in}{0.943313in}}%
\pgfpathlineto{\pgfqpoint{2.468504in}{0.944341in}}%
\pgfpathlineto{\pgfqpoint{2.469213in}{0.932802in}}%
\pgfpathlineto{\pgfqpoint{2.470515in}{0.912039in}}%
\pgfpathlineto{\pgfqpoint{2.474111in}{0.910147in}}%
\pgfpathlineto{\pgfqpoint{2.474841in}{0.898468in}}%
\pgfpathlineto{\pgfqpoint{2.458419in}{0.897453in}}%
\pgfpathlineto{\pgfqpoint{2.455570in}{0.891636in}}%
\pgfpathlineto{\pgfqpoint{2.426856in}{0.899501in}}%
\pgfpathlineto{\pgfqpoint{2.419871in}{0.899622in}}%
\pgfpathlineto{\pgfqpoint{2.416578in}{0.896534in}}%
\pgfpathclose%
\pgfusepath{fill}%
\end{pgfscope}%
\begin{pgfscope}%
\pgfpathrectangle{\pgfqpoint{0.100000in}{0.100000in}}{\pgfqpoint{3.608454in}{2.310000in}}%
\pgfusepath{clip}%
\pgfsetbuttcap%
\pgfsetmiterjoin%
\definecolor{currentfill}{rgb}{0.000000,0.333333,0.833333}%
\pgfsetfillcolor{currentfill}%
\pgfsetlinewidth{0.000000pt}%
\definecolor{currentstroke}{rgb}{0.000000,0.000000,0.000000}%
\pgfsetstrokecolor{currentstroke}%
\pgfsetstrokeopacity{0.000000}%
\pgfsetdash{}{0pt}%
\pgfpathmoveto{\pgfqpoint{1.913591in}{0.794961in}}%
\pgfpathlineto{\pgfqpoint{1.908521in}{0.806762in}}%
\pgfpathlineto{\pgfqpoint{1.900998in}{0.829156in}}%
\pgfpathlineto{\pgfqpoint{1.873743in}{0.830250in}}%
\pgfpathlineto{\pgfqpoint{1.866975in}{0.830682in}}%
\pgfpathlineto{\pgfqpoint{1.868391in}{0.865576in}}%
\pgfpathlineto{\pgfqpoint{1.878275in}{0.865031in}}%
\pgfpathlineto{\pgfqpoint{1.878486in}{0.869447in}}%
\pgfpathlineto{\pgfqpoint{1.912098in}{0.867745in}}%
\pgfpathlineto{\pgfqpoint{1.913431in}{0.902088in}}%
\pgfpathlineto{\pgfqpoint{1.909495in}{0.904858in}}%
\pgfpathlineto{\pgfqpoint{1.910664in}{0.938402in}}%
\pgfpathlineto{\pgfqpoint{1.917943in}{0.934949in}}%
\pgfpathlineto{\pgfqpoint{1.931087in}{0.945835in}}%
\pgfpathlineto{\pgfqpoint{1.938143in}{0.938336in}}%
\pgfpathlineto{\pgfqpoint{1.943227in}{0.939444in}}%
\pgfpathlineto{\pgfqpoint{1.941990in}{0.901235in}}%
\pgfpathlineto{\pgfqpoint{1.948881in}{0.900975in}}%
\pgfpathlineto{\pgfqpoint{1.946927in}{0.865843in}}%
\pgfpathlineto{\pgfqpoint{1.971315in}{0.864806in}}%
\pgfpathlineto{\pgfqpoint{1.970138in}{0.830158in}}%
\pgfpathlineto{\pgfqpoint{1.966891in}{0.830295in}}%
\pgfpathlineto{\pgfqpoint{1.966416in}{0.807704in}}%
\pgfpathlineto{\pgfqpoint{1.940012in}{0.801698in}}%
\pgfpathlineto{\pgfqpoint{1.934684in}{0.799256in}}%
\pgfpathlineto{\pgfqpoint{1.929646in}{0.804003in}}%
\pgfpathlineto{\pgfqpoint{1.913591in}{0.794961in}}%
\pgfpathclose%
\pgfusepath{fill}%
\end{pgfscope}%
\begin{pgfscope}%
\pgfpathrectangle{\pgfqpoint{0.100000in}{0.100000in}}{\pgfqpoint{3.608454in}{2.310000in}}%
\pgfusepath{clip}%
\pgfsetbuttcap%
\pgfsetmiterjoin%
\definecolor{currentfill}{rgb}{0.000000,0.454902,0.772549}%
\pgfsetfillcolor{currentfill}%
\pgfsetlinewidth{0.000000pt}%
\definecolor{currentstroke}{rgb}{0.000000,0.000000,0.000000}%
\pgfsetstrokecolor{currentstroke}%
\pgfsetstrokeopacity{0.000000}%
\pgfsetdash{}{0pt}%
\pgfpathmoveto{\pgfqpoint{2.698781in}{1.052062in}}%
\pgfpathlineto{\pgfqpoint{2.669947in}{1.049451in}}%
\pgfpathlineto{\pgfqpoint{2.668473in}{1.060141in}}%
\pgfpathlineto{\pgfqpoint{2.665061in}{1.060409in}}%
\pgfpathlineto{\pgfqpoint{2.662888in}{1.068219in}}%
\pgfpathlineto{\pgfqpoint{2.648588in}{1.077186in}}%
\pgfpathlineto{\pgfqpoint{2.644297in}{1.082383in}}%
\pgfpathlineto{\pgfqpoint{2.643683in}{1.102257in}}%
\pgfpathlineto{\pgfqpoint{2.651600in}{1.103191in}}%
\pgfpathlineto{\pgfqpoint{2.654621in}{1.100124in}}%
\pgfpathlineto{\pgfqpoint{2.668887in}{1.104876in}}%
\pgfpathlineto{\pgfqpoint{2.681536in}{1.102539in}}%
\pgfpathlineto{\pgfqpoint{2.685775in}{1.104241in}}%
\pgfpathlineto{\pgfqpoint{2.688788in}{1.097666in}}%
\pgfpathlineto{\pgfqpoint{2.693772in}{1.094220in}}%
\pgfpathlineto{\pgfqpoint{2.694900in}{1.080379in}}%
\pgfpathlineto{\pgfqpoint{2.693282in}{1.075443in}}%
\pgfpathlineto{\pgfqpoint{2.696968in}{1.070359in}}%
\pgfpathlineto{\pgfqpoint{2.698781in}{1.052062in}}%
\pgfpathclose%
\pgfusepath{fill}%
\end{pgfscope}%
\begin{pgfscope}%
\pgfpathrectangle{\pgfqpoint{0.100000in}{0.100000in}}{\pgfqpoint{3.608454in}{2.310000in}}%
\pgfusepath{clip}%
\pgfsetbuttcap%
\pgfsetmiterjoin%
\definecolor{currentfill}{rgb}{0.000000,0.647059,0.676471}%
\pgfsetfillcolor{currentfill}%
\pgfsetlinewidth{0.000000pt}%
\definecolor{currentstroke}{rgb}{0.000000,0.000000,0.000000}%
\pgfsetstrokecolor{currentstroke}%
\pgfsetstrokeopacity{0.000000}%
\pgfsetdash{}{0pt}%
\pgfpathmoveto{\pgfqpoint{2.550184in}{0.875112in}}%
\pgfpathlineto{\pgfqpoint{2.519263in}{0.872679in}}%
\pgfpathlineto{\pgfqpoint{2.517297in}{0.901251in}}%
\pgfpathlineto{\pgfqpoint{2.526832in}{0.901949in}}%
\pgfpathlineto{\pgfqpoint{2.525476in}{0.919309in}}%
\pgfpathlineto{\pgfqpoint{2.534527in}{0.922423in}}%
\pgfpathlineto{\pgfqpoint{2.537884in}{0.927921in}}%
\pgfpathlineto{\pgfqpoint{2.535053in}{0.931443in}}%
\pgfpathlineto{\pgfqpoint{2.544297in}{0.934956in}}%
\pgfpathlineto{\pgfqpoint{2.546945in}{0.939900in}}%
\pgfpathlineto{\pgfqpoint{2.551959in}{0.940228in}}%
\pgfpathlineto{\pgfqpoint{2.551506in}{0.923399in}}%
\pgfpathlineto{\pgfqpoint{2.580198in}{0.924865in}}%
\pgfpathlineto{\pgfqpoint{2.582729in}{0.895488in}}%
\pgfpathlineto{\pgfqpoint{2.575754in}{0.891315in}}%
\pgfpathlineto{\pgfqpoint{2.570549in}{0.885685in}}%
\pgfpathlineto{\pgfqpoint{2.564922in}{0.885225in}}%
\pgfpathlineto{\pgfqpoint{2.561497in}{0.881275in}}%
\pgfpathlineto{\pgfqpoint{2.550318in}{0.880073in}}%
\pgfpathlineto{\pgfqpoint{2.550184in}{0.875112in}}%
\pgfpathclose%
\pgfusepath{fill}%
\end{pgfscope}%
\begin{pgfscope}%
\pgfpathrectangle{\pgfqpoint{0.100000in}{0.100000in}}{\pgfqpoint{3.608454in}{2.310000in}}%
\pgfusepath{clip}%
\pgfsetbuttcap%
\pgfsetmiterjoin%
\definecolor{currentfill}{rgb}{0.000000,0.725490,0.637255}%
\pgfsetfillcolor{currentfill}%
\pgfsetlinewidth{0.000000pt}%
\definecolor{currentstroke}{rgb}{0.000000,0.000000,0.000000}%
\pgfsetstrokecolor{currentstroke}%
\pgfsetstrokeopacity{0.000000}%
\pgfsetdash{}{0pt}%
\pgfpathmoveto{\pgfqpoint{3.215154in}{1.664914in}}%
\pgfpathlineto{\pgfqpoint{3.207689in}{1.691369in}}%
\pgfpathlineto{\pgfqpoint{3.183986in}{1.686549in}}%
\pgfpathlineto{\pgfqpoint{3.180627in}{1.708614in}}%
\pgfpathlineto{\pgfqpoint{3.187096in}{1.712911in}}%
\pgfpathlineto{\pgfqpoint{3.196916in}{1.713203in}}%
\pgfpathlineto{\pgfqpoint{3.194930in}{1.721757in}}%
\pgfpathlineto{\pgfqpoint{3.217063in}{1.727096in}}%
\pgfpathlineto{\pgfqpoint{3.220803in}{1.715267in}}%
\pgfpathlineto{\pgfqpoint{3.232758in}{1.716945in}}%
\pgfpathlineto{\pgfqpoint{3.233713in}{1.712856in}}%
\pgfpathlineto{\pgfqpoint{3.246255in}{1.715536in}}%
\pgfpathlineto{\pgfqpoint{3.248659in}{1.703631in}}%
\pgfpathlineto{\pgfqpoint{3.252945in}{1.701154in}}%
\pgfpathlineto{\pgfqpoint{3.245833in}{1.699582in}}%
\pgfpathlineto{\pgfqpoint{3.253079in}{1.672195in}}%
\pgfpathlineto{\pgfqpoint{3.215154in}{1.664914in}}%
\pgfpathclose%
\pgfusepath{fill}%
\end{pgfscope}%
\begin{pgfscope}%
\pgfpathrectangle{\pgfqpoint{0.100000in}{0.100000in}}{\pgfqpoint{3.608454in}{2.310000in}}%
\pgfusepath{clip}%
\pgfsetbuttcap%
\pgfsetmiterjoin%
\definecolor{currentfill}{rgb}{0.000000,0.521569,0.739216}%
\pgfsetfillcolor{currentfill}%
\pgfsetlinewidth{0.000000pt}%
\definecolor{currentstroke}{rgb}{0.000000,0.000000,0.000000}%
\pgfsetstrokecolor{currentstroke}%
\pgfsetstrokeopacity{0.000000}%
\pgfsetdash{}{0pt}%
\pgfpathmoveto{\pgfqpoint{2.020851in}{0.893088in}}%
\pgfpathlineto{\pgfqpoint{2.014996in}{0.896820in}}%
\pgfpathlineto{\pgfqpoint{1.984737in}{0.897974in}}%
\pgfpathlineto{\pgfqpoint{1.977949in}{0.898974in}}%
\pgfpathlineto{\pgfqpoint{1.978804in}{0.941263in}}%
\pgfpathlineto{\pgfqpoint{1.982724in}{0.934735in}}%
\pgfpathlineto{\pgfqpoint{1.993908in}{0.931626in}}%
\pgfpathlineto{\pgfqpoint{2.000025in}{0.945069in}}%
\pgfpathlineto{\pgfqpoint{2.014398in}{0.957120in}}%
\pgfpathlineto{\pgfqpoint{2.041706in}{0.956720in}}%
\pgfpathlineto{\pgfqpoint{2.045185in}{0.954206in}}%
\pgfpathlineto{\pgfqpoint{2.045055in}{0.940136in}}%
\pgfpathlineto{\pgfqpoint{2.051407in}{0.934194in}}%
\pgfpathlineto{\pgfqpoint{2.056583in}{0.934051in}}%
\pgfpathlineto{\pgfqpoint{2.050404in}{0.930484in}}%
\pgfpathlineto{\pgfqpoint{2.049865in}{0.897339in}}%
\pgfpathlineto{\pgfqpoint{2.020851in}{0.893088in}}%
\pgfpathclose%
\pgfusepath{fill}%
\end{pgfscope}%
\begin{pgfscope}%
\pgfpathrectangle{\pgfqpoint{0.100000in}{0.100000in}}{\pgfqpoint{3.608454in}{2.310000in}}%
\pgfusepath{clip}%
\pgfsetbuttcap%
\pgfsetmiterjoin%
\definecolor{currentfill}{rgb}{0.000000,0.564706,0.717647}%
\pgfsetfillcolor{currentfill}%
\pgfsetlinewidth{0.000000pt}%
\definecolor{currentstroke}{rgb}{0.000000,0.000000,0.000000}%
\pgfsetstrokecolor{currentstroke}%
\pgfsetstrokeopacity{0.000000}%
\pgfsetdash{}{0pt}%
\pgfpathmoveto{\pgfqpoint{2.561760in}{0.672494in}}%
\pgfpathlineto{\pgfqpoint{2.556138in}{0.721372in}}%
\pgfpathlineto{\pgfqpoint{2.551986in}{0.756063in}}%
\pgfpathlineto{\pgfqpoint{2.549469in}{0.776821in}}%
\pgfpathlineto{\pgfqpoint{2.574889in}{0.778819in}}%
\pgfpathlineto{\pgfqpoint{2.573394in}{0.783619in}}%
\pgfpathlineto{\pgfqpoint{2.567924in}{0.787803in}}%
\pgfpathlineto{\pgfqpoint{2.574126in}{0.801999in}}%
\pgfpathlineto{\pgfqpoint{2.601434in}{0.804239in}}%
\pgfpathlineto{\pgfqpoint{2.604373in}{0.805606in}}%
\pgfpathlineto{\pgfqpoint{2.603471in}{0.815667in}}%
\pgfpathlineto{\pgfqpoint{2.610325in}{0.816266in}}%
\pgfpathlineto{\pgfqpoint{2.609428in}{0.825937in}}%
\pgfpathlineto{\pgfqpoint{2.612796in}{0.827022in}}%
\pgfpathlineto{\pgfqpoint{2.622706in}{0.822514in}}%
\pgfpathlineto{\pgfqpoint{2.634007in}{0.811449in}}%
\pgfpathlineto{\pgfqpoint{2.652230in}{0.813085in}}%
\pgfpathlineto{\pgfqpoint{2.653773in}{0.795845in}}%
\pgfpathlineto{\pgfqpoint{2.613728in}{0.792291in}}%
\pgfpathlineto{\pgfqpoint{2.613565in}{0.781771in}}%
\pgfpathlineto{\pgfqpoint{2.610143in}{0.781433in}}%
\pgfpathlineto{\pgfqpoint{2.611445in}{0.760754in}}%
\pgfpathlineto{\pgfqpoint{2.605232in}{0.759125in}}%
\pgfpathlineto{\pgfqpoint{2.601460in}{0.754867in}}%
\pgfpathlineto{\pgfqpoint{2.602701in}{0.749237in}}%
\pgfpathlineto{\pgfqpoint{2.609757in}{0.745194in}}%
\pgfpathlineto{\pgfqpoint{2.612513in}{0.725646in}}%
\pgfpathlineto{\pgfqpoint{2.610908in}{0.714979in}}%
\pgfpathlineto{\pgfqpoint{2.618895in}{0.705619in}}%
\pgfpathlineto{\pgfqpoint{2.627748in}{0.701087in}}%
\pgfpathlineto{\pgfqpoint{2.628964in}{0.696389in}}%
\pgfpathlineto{\pgfqpoint{2.626017in}{0.689071in}}%
\pgfpathlineto{\pgfqpoint{2.632066in}{0.682154in}}%
\pgfpathlineto{\pgfqpoint{2.627322in}{0.678308in}}%
\pgfpathlineto{\pgfqpoint{2.622714in}{0.669017in}}%
\pgfpathlineto{\pgfqpoint{2.613254in}{0.665766in}}%
\pgfpathlineto{\pgfqpoint{2.606414in}{0.667773in}}%
\pgfpathlineto{\pgfqpoint{2.595006in}{0.677521in}}%
\pgfpathlineto{\pgfqpoint{2.592621in}{0.682803in}}%
\pgfpathlineto{\pgfqpoint{2.594483in}{0.687365in}}%
\pgfpathlineto{\pgfqpoint{2.592065in}{0.696291in}}%
\pgfpathlineto{\pgfqpoint{2.588127in}{0.698733in}}%
\pgfpathlineto{\pgfqpoint{2.583318in}{0.695147in}}%
\pgfpathlineto{\pgfqpoint{2.581260in}{0.684211in}}%
\pgfpathlineto{\pgfqpoint{2.581606in}{0.673454in}}%
\pgfpathlineto{\pgfqpoint{2.580066in}{0.669312in}}%
\pgfpathlineto{\pgfqpoint{2.564924in}{0.674357in}}%
\pgfpathlineto{\pgfqpoint{2.561760in}{0.672494in}}%
\pgfpathclose%
\pgfusepath{fill}%
\end{pgfscope}%
\begin{pgfscope}%
\pgfpathrectangle{\pgfqpoint{0.100000in}{0.100000in}}{\pgfqpoint{3.608454in}{2.310000in}}%
\pgfusepath{clip}%
\pgfsetbuttcap%
\pgfsetmiterjoin%
\definecolor{currentfill}{rgb}{0.000000,0.545098,0.727451}%
\pgfsetfillcolor{currentfill}%
\pgfsetlinewidth{0.000000pt}%
\definecolor{currentstroke}{rgb}{0.000000,0.000000,0.000000}%
\pgfsetstrokecolor{currentstroke}%
\pgfsetstrokeopacity{0.000000}%
\pgfsetdash{}{0pt}%
\pgfpathmoveto{\pgfqpoint{1.440664in}{1.064305in}}%
\pgfpathlineto{\pgfqpoint{1.436822in}{1.030070in}}%
\pgfpathlineto{\pgfqpoint{1.435358in}{1.030206in}}%
\pgfpathlineto{\pgfqpoint{1.432998in}{1.009699in}}%
\pgfpathlineto{\pgfqpoint{1.432205in}{1.002872in}}%
\pgfpathlineto{\pgfqpoint{1.392442in}{1.007538in}}%
\pgfpathlineto{\pgfqpoint{1.360437in}{1.011466in}}%
\pgfpathlineto{\pgfqpoint{1.362154in}{1.025644in}}%
\pgfpathlineto{\pgfqpoint{1.340304in}{1.035758in}}%
\pgfpathlineto{\pgfqpoint{1.333401in}{1.040515in}}%
\pgfpathlineto{\pgfqpoint{1.278910in}{1.047826in}}%
\pgfpathlineto{\pgfqpoint{1.236876in}{1.053946in}}%
\pgfpathlineto{\pgfqpoint{1.193614in}{1.060581in}}%
\pgfpathlineto{\pgfqpoint{1.198276in}{1.090371in}}%
\pgfpathlineto{\pgfqpoint{1.235282in}{1.084666in}}%
\pgfpathlineto{\pgfqpoint{1.239395in}{1.111911in}}%
\pgfpathlineto{\pgfqpoint{1.313512in}{1.101203in}}%
\pgfpathlineto{\pgfqpoint{1.321043in}{1.155597in}}%
\pgfpathlineto{\pgfqpoint{1.300870in}{1.158512in}}%
\pgfpathlineto{\pgfqpoint{1.303419in}{1.175695in}}%
\pgfpathlineto{\pgfqpoint{1.350154in}{1.169097in}}%
\pgfpathlineto{\pgfqpoint{1.347890in}{1.152024in}}%
\pgfpathlineto{\pgfqpoint{1.388440in}{1.147038in}}%
\pgfpathlineto{\pgfqpoint{1.378347in}{1.144995in}}%
\pgfpathlineto{\pgfqpoint{1.377042in}{1.134592in}}%
\pgfpathlineto{\pgfqpoint{1.385971in}{1.127718in}}%
\pgfpathlineto{\pgfqpoint{1.379252in}{1.071453in}}%
\pgfpathlineto{\pgfqpoint{1.413364in}{1.067432in}}%
\pgfpathlineto{\pgfqpoint{1.440664in}{1.064305in}}%
\pgfpathclose%
\pgfusepath{fill}%
\end{pgfscope}%
\begin{pgfscope}%
\pgfpathrectangle{\pgfqpoint{0.100000in}{0.100000in}}{\pgfqpoint{3.608454in}{2.310000in}}%
\pgfusepath{clip}%
\pgfsetbuttcap%
\pgfsetmiterjoin%
\definecolor{currentfill}{rgb}{0.000000,0.627451,0.686275}%
\pgfsetfillcolor{currentfill}%
\pgfsetlinewidth{0.000000pt}%
\definecolor{currentstroke}{rgb}{0.000000,0.000000,0.000000}%
\pgfsetstrokecolor{currentstroke}%
\pgfsetstrokeopacity{0.000000}%
\pgfsetdash{}{0pt}%
\pgfpathmoveto{\pgfqpoint{1.993561in}{0.650091in}}%
\pgfpathlineto{\pgfqpoint{1.986455in}{0.652084in}}%
\pgfpathlineto{\pgfqpoint{1.978079in}{0.661346in}}%
\pgfpathlineto{\pgfqpoint{1.975419in}{0.668492in}}%
\pgfpathlineto{\pgfqpoint{1.971709in}{0.671338in}}%
\pgfpathlineto{\pgfqpoint{1.995500in}{0.684736in}}%
\pgfpathlineto{\pgfqpoint{1.993815in}{0.692311in}}%
\pgfpathlineto{\pgfqpoint{1.986509in}{0.700655in}}%
\pgfpathlineto{\pgfqpoint{1.984720in}{0.711352in}}%
\pgfpathlineto{\pgfqpoint{1.981851in}{0.714949in}}%
\pgfpathlineto{\pgfqpoint{1.997875in}{0.723825in}}%
\pgfpathlineto{\pgfqpoint{2.016972in}{0.734385in}}%
\pgfpathlineto{\pgfqpoint{2.016298in}{0.728776in}}%
\pgfpathlineto{\pgfqpoint{2.021875in}{0.703710in}}%
\pgfpathlineto{\pgfqpoint{2.026699in}{0.691630in}}%
\pgfpathlineto{\pgfqpoint{2.047653in}{0.694679in}}%
\pgfpathlineto{\pgfqpoint{2.049731in}{0.675985in}}%
\pgfpathlineto{\pgfqpoint{2.051242in}{0.645234in}}%
\pgfpathlineto{\pgfqpoint{2.031567in}{0.644083in}}%
\pgfpathlineto{\pgfqpoint{2.031314in}{0.647827in}}%
\pgfpathlineto{\pgfqpoint{2.025563in}{0.657506in}}%
\pgfpathlineto{\pgfqpoint{2.016277in}{0.656724in}}%
\pgfpathlineto{\pgfqpoint{2.008653in}{0.653539in}}%
\pgfpathlineto{\pgfqpoint{1.993561in}{0.650091in}}%
\pgfpathclose%
\pgfusepath{fill}%
\end{pgfscope}%
\begin{pgfscope}%
\pgfpathrectangle{\pgfqpoint{0.100000in}{0.100000in}}{\pgfqpoint{3.608454in}{2.310000in}}%
\pgfusepath{clip}%
\pgfsetbuttcap%
\pgfsetmiterjoin%
\definecolor{currentfill}{rgb}{0.000000,0.850980,0.574510}%
\pgfsetfillcolor{currentfill}%
\pgfsetlinewidth{0.000000pt}%
\definecolor{currentstroke}{rgb}{0.000000,0.000000,0.000000}%
\pgfsetstrokecolor{currentstroke}%
\pgfsetstrokeopacity{0.000000}%
\pgfsetdash{}{0pt}%
\pgfpathmoveto{\pgfqpoint{2.991086in}{1.284432in}}%
\pgfpathlineto{\pgfqpoint{2.974604in}{1.294585in}}%
\pgfpathlineto{\pgfqpoint{2.968540in}{1.292624in}}%
\pgfpathlineto{\pgfqpoint{2.967623in}{1.287269in}}%
\pgfpathlineto{\pgfqpoint{2.961369in}{1.291611in}}%
\pgfpathlineto{\pgfqpoint{2.952046in}{1.304075in}}%
\pgfpathlineto{\pgfqpoint{2.946226in}{1.306488in}}%
\pgfpathlineto{\pgfqpoint{2.941096in}{1.316634in}}%
\pgfpathlineto{\pgfqpoint{2.944352in}{1.321971in}}%
\pgfpathlineto{\pgfqpoint{2.952368in}{1.321633in}}%
\pgfpathlineto{\pgfqpoint{2.953019in}{1.337104in}}%
\pgfpathlineto{\pgfqpoint{2.958041in}{1.346022in}}%
\pgfpathlineto{\pgfqpoint{2.968495in}{1.342247in}}%
\pgfpathlineto{\pgfqpoint{2.973652in}{1.345179in}}%
\pgfpathlineto{\pgfqpoint{2.981470in}{1.336849in}}%
\pgfpathlineto{\pgfqpoint{2.992200in}{1.336598in}}%
\pgfpathlineto{\pgfqpoint{3.006750in}{1.353994in}}%
\pgfpathlineto{\pgfqpoint{3.011961in}{1.351516in}}%
\pgfpathlineto{\pgfqpoint{3.013175in}{1.342782in}}%
\pgfpathlineto{\pgfqpoint{3.020093in}{1.338488in}}%
\pgfpathlineto{\pgfqpoint{3.026105in}{1.338632in}}%
\pgfpathlineto{\pgfqpoint{3.038592in}{1.342913in}}%
\pgfpathlineto{\pgfqpoint{3.037059in}{1.337144in}}%
\pgfpathlineto{\pgfqpoint{3.025954in}{1.320474in}}%
\pgfpathlineto{\pgfqpoint{3.022678in}{1.310484in}}%
\pgfpathlineto{\pgfqpoint{3.013170in}{1.308307in}}%
\pgfpathlineto{\pgfqpoint{3.000083in}{1.309688in}}%
\pgfpathlineto{\pgfqpoint{2.991086in}{1.284432in}}%
\pgfpathclose%
\pgfusepath{fill}%
\end{pgfscope}%
\begin{pgfscope}%
\pgfpathrectangle{\pgfqpoint{0.100000in}{0.100000in}}{\pgfqpoint{3.608454in}{2.310000in}}%
\pgfusepath{clip}%
\pgfsetbuttcap%
\pgfsetmiterjoin%
\definecolor{currentfill}{rgb}{0.000000,0.592157,0.703922}%
\pgfsetfillcolor{currentfill}%
\pgfsetlinewidth{0.000000pt}%
\definecolor{currentstroke}{rgb}{0.000000,0.000000,0.000000}%
\pgfsetstrokecolor{currentstroke}%
\pgfsetstrokeopacity{0.000000}%
\pgfsetdash{}{0pt}%
\pgfpathmoveto{\pgfqpoint{3.302028in}{1.725960in}}%
\pgfpathlineto{\pgfqpoint{3.267733in}{1.692380in}}%
\pgfpathlineto{\pgfqpoint{3.260640in}{1.692420in}}%
\pgfpathlineto{\pgfqpoint{3.258600in}{1.697493in}}%
\pgfpathlineto{\pgfqpoint{3.252945in}{1.701154in}}%
\pgfpathlineto{\pgfqpoint{3.248659in}{1.703631in}}%
\pgfpathlineto{\pgfqpoint{3.246255in}{1.715536in}}%
\pgfpathlineto{\pgfqpoint{3.233713in}{1.712856in}}%
\pgfpathlineto{\pgfqpoint{3.232758in}{1.716945in}}%
\pgfpathlineto{\pgfqpoint{3.220803in}{1.715267in}}%
\pgfpathlineto{\pgfqpoint{3.217063in}{1.727096in}}%
\pgfpathlineto{\pgfqpoint{3.210494in}{1.750645in}}%
\pgfpathlineto{\pgfqpoint{3.243936in}{1.759541in}}%
\pgfpathlineto{\pgfqpoint{3.246838in}{1.769301in}}%
\pgfpathlineto{\pgfqpoint{3.252161in}{1.774686in}}%
\pgfpathlineto{\pgfqpoint{3.264610in}{1.770699in}}%
\pgfpathlineto{\pgfqpoint{3.264917in}{1.776723in}}%
\pgfpathlineto{\pgfqpoint{3.272069in}{1.775491in}}%
\pgfpathlineto{\pgfqpoint{3.279173in}{1.774369in}}%
\pgfpathlineto{\pgfqpoint{3.279470in}{1.768032in}}%
\pgfpathlineto{\pgfqpoint{3.283745in}{1.758930in}}%
\pgfpathlineto{\pgfqpoint{3.281055in}{1.749476in}}%
\pgfpathlineto{\pgfqpoint{3.287975in}{1.743465in}}%
\pgfpathlineto{\pgfqpoint{3.300176in}{1.740144in}}%
\pgfpathlineto{\pgfqpoint{3.296564in}{1.727321in}}%
\pgfpathlineto{\pgfqpoint{3.302028in}{1.725960in}}%
\pgfpathclose%
\pgfusepath{fill}%
\end{pgfscope}%
\begin{pgfscope}%
\pgfpathrectangle{\pgfqpoint{0.100000in}{0.100000in}}{\pgfqpoint{3.608454in}{2.310000in}}%
\pgfusepath{clip}%
\pgfsetbuttcap%
\pgfsetmiterjoin%
\definecolor{currentfill}{rgb}{0.000000,0.478431,0.760784}%
\pgfsetfillcolor{currentfill}%
\pgfsetlinewidth{0.000000pt}%
\definecolor{currentstroke}{rgb}{0.000000,0.000000,0.000000}%
\pgfsetstrokecolor{currentstroke}%
\pgfsetstrokeopacity{0.000000}%
\pgfsetdash{}{0pt}%
\pgfpathmoveto{\pgfqpoint{1.917334in}{2.003750in}}%
\pgfpathlineto{\pgfqpoint{1.875835in}{2.005576in}}%
\pgfpathlineto{\pgfqpoint{1.877148in}{2.033168in}}%
\pgfpathlineto{\pgfqpoint{1.876114in}{2.047161in}}%
\pgfpathlineto{\pgfqpoint{1.917529in}{2.045335in}}%
\pgfpathlineto{\pgfqpoint{1.918355in}{2.031373in}}%
\pgfpathlineto{\pgfqpoint{1.917334in}{2.003750in}}%
\pgfpathclose%
\pgfusepath{fill}%
\end{pgfscope}%
\begin{pgfscope}%
\pgfpathrectangle{\pgfqpoint{0.100000in}{0.100000in}}{\pgfqpoint{3.608454in}{2.310000in}}%
\pgfusepath{clip}%
\pgfsetbuttcap%
\pgfsetmiterjoin%
\definecolor{currentfill}{rgb}{0.000000,0.576471,0.711765}%
\pgfsetfillcolor{currentfill}%
\pgfsetlinewidth{0.000000pt}%
\definecolor{currentstroke}{rgb}{0.000000,0.000000,0.000000}%
\pgfsetstrokecolor{currentstroke}%
\pgfsetstrokeopacity{0.000000}%
\pgfsetdash{}{0pt}%
\pgfpathmoveto{\pgfqpoint{1.933314in}{1.014281in}}%
\pgfpathlineto{\pgfqpoint{1.932881in}{1.000512in}}%
\pgfpathlineto{\pgfqpoint{1.905475in}{1.001362in}}%
\pgfpathlineto{\pgfqpoint{1.905912in}{1.015133in}}%
\pgfpathlineto{\pgfqpoint{1.871452in}{1.016408in}}%
\pgfpathlineto{\pgfqpoint{1.872137in}{1.035604in}}%
\pgfpathlineto{\pgfqpoint{1.873435in}{1.071561in}}%
\pgfpathlineto{\pgfqpoint{1.893386in}{1.070793in}}%
\pgfpathlineto{\pgfqpoint{1.893326in}{1.056979in}}%
\pgfpathlineto{\pgfqpoint{1.911760in}{1.056365in}}%
\pgfpathlineto{\pgfqpoint{1.923203in}{1.052514in}}%
\pgfpathlineto{\pgfqpoint{1.934200in}{1.052365in}}%
\pgfpathlineto{\pgfqpoint{1.933314in}{1.014281in}}%
\pgfpathclose%
\pgfusepath{fill}%
\end{pgfscope}%
\begin{pgfscope}%
\pgfpathrectangle{\pgfqpoint{0.100000in}{0.100000in}}{\pgfqpoint{3.608454in}{2.310000in}}%
\pgfusepath{clip}%
\pgfsetbuttcap%
\pgfsetmiterjoin%
\definecolor{currentfill}{rgb}{0.000000,0.458824,0.770588}%
\pgfsetfillcolor{currentfill}%
\pgfsetlinewidth{0.000000pt}%
\definecolor{currentstroke}{rgb}{0.000000,0.000000,0.000000}%
\pgfsetstrokecolor{currentstroke}%
\pgfsetstrokeopacity{0.000000}%
\pgfsetdash{}{0pt}%
\pgfpathmoveto{\pgfqpoint{2.165979in}{1.695985in}}%
\pgfpathlineto{\pgfqpoint{2.166463in}{1.649078in}}%
\pgfpathlineto{\pgfqpoint{2.111776in}{1.648912in}}%
\pgfpathlineto{\pgfqpoint{2.084271in}{1.648958in}}%
\pgfpathlineto{\pgfqpoint{2.084392in}{1.676317in}}%
\pgfpathlineto{\pgfqpoint{2.084490in}{1.695746in}}%
\pgfpathlineto{\pgfqpoint{2.111687in}{1.695719in}}%
\pgfpathlineto{\pgfqpoint{2.111765in}{1.676267in}}%
\pgfpathlineto{\pgfqpoint{2.138939in}{1.676359in}}%
\pgfpathlineto{\pgfqpoint{2.138833in}{1.695825in}}%
\pgfpathlineto{\pgfqpoint{2.165979in}{1.695985in}}%
\pgfpathclose%
\pgfusepath{fill}%
\end{pgfscope}%
\begin{pgfscope}%
\pgfpathrectangle{\pgfqpoint{0.100000in}{0.100000in}}{\pgfqpoint{3.608454in}{2.310000in}}%
\pgfusepath{clip}%
\pgfsetbuttcap%
\pgfsetmiterjoin%
\definecolor{currentfill}{rgb}{0.000000,0.713725,0.643137}%
\pgfsetfillcolor{currentfill}%
\pgfsetlinewidth{0.000000pt}%
\definecolor{currentstroke}{rgb}{0.000000,0.000000,0.000000}%
\pgfsetstrokecolor{currentstroke}%
\pgfsetstrokeopacity{0.000000}%
\pgfsetdash{}{0pt}%
\pgfpathmoveto{\pgfqpoint{2.824999in}{1.309058in}}%
\pgfpathlineto{\pgfqpoint{2.825277in}{1.301789in}}%
\pgfpathlineto{\pgfqpoint{2.820760in}{1.291911in}}%
\pgfpathlineto{\pgfqpoint{2.814875in}{1.291146in}}%
\pgfpathlineto{\pgfqpoint{2.806402in}{1.296003in}}%
\pgfpathlineto{\pgfqpoint{2.793467in}{1.297141in}}%
\pgfpathlineto{\pgfqpoint{2.784277in}{1.310449in}}%
\pgfpathlineto{\pgfqpoint{2.796848in}{1.326453in}}%
\pgfpathlineto{\pgfqpoint{2.804061in}{1.326162in}}%
\pgfpathlineto{\pgfqpoint{2.811045in}{1.319515in}}%
\pgfpathlineto{\pgfqpoint{2.824999in}{1.309058in}}%
\pgfpathclose%
\pgfusepath{fill}%
\end{pgfscope}%
\begin{pgfscope}%
\pgfpathrectangle{\pgfqpoint{0.100000in}{0.100000in}}{\pgfqpoint{3.608454in}{2.310000in}}%
\pgfusepath{clip}%
\pgfsetbuttcap%
\pgfsetmiterjoin%
\definecolor{currentfill}{rgb}{0.000000,0.607843,0.696078}%
\pgfsetfillcolor{currentfill}%
\pgfsetlinewidth{0.000000pt}%
\definecolor{currentstroke}{rgb}{0.000000,0.000000,0.000000}%
\pgfsetstrokecolor{currentstroke}%
\pgfsetstrokeopacity{0.000000}%
\pgfsetdash{}{0pt}%
\pgfpathmoveto{\pgfqpoint{2.586580in}{1.532701in}}%
\pgfpathlineto{\pgfqpoint{2.588380in}{1.512229in}}%
\pgfpathlineto{\pgfqpoint{2.585581in}{1.506060in}}%
\pgfpathlineto{\pgfqpoint{2.578828in}{1.505444in}}%
\pgfpathlineto{\pgfqpoint{2.579594in}{1.497487in}}%
\pgfpathlineto{\pgfqpoint{2.554067in}{1.495355in}}%
\pgfpathlineto{\pgfqpoint{2.550384in}{1.539587in}}%
\pgfpathlineto{\pgfqpoint{2.547881in}{1.571936in}}%
\pgfpathlineto{\pgfqpoint{2.554797in}{1.567328in}}%
\pgfpathlineto{\pgfqpoint{2.566678in}{1.566887in}}%
\pgfpathlineto{\pgfqpoint{2.582676in}{1.575047in}}%
\pgfpathlineto{\pgfqpoint{2.586580in}{1.532701in}}%
\pgfpathclose%
\pgfusepath{fill}%
\end{pgfscope}%
\begin{pgfscope}%
\pgfpathrectangle{\pgfqpoint{0.100000in}{0.100000in}}{\pgfqpoint{3.608454in}{2.310000in}}%
\pgfusepath{clip}%
\pgfsetbuttcap%
\pgfsetmiterjoin%
\definecolor{currentfill}{rgb}{0.000000,0.478431,0.760784}%
\pgfsetfillcolor{currentfill}%
\pgfsetlinewidth{0.000000pt}%
\definecolor{currentstroke}{rgb}{0.000000,0.000000,0.000000}%
\pgfsetstrokecolor{currentstroke}%
\pgfsetstrokeopacity{0.000000}%
\pgfsetdash{}{0pt}%
\pgfpathmoveto{\pgfqpoint{1.844714in}{0.901800in}}%
\pgfpathlineto{\pgfqpoint{1.880039in}{0.900259in}}%
\pgfpathlineto{\pgfqpoint{1.880200in}{0.905864in}}%
\pgfpathlineto{\pgfqpoint{1.909495in}{0.904858in}}%
\pgfpathlineto{\pgfqpoint{1.913431in}{0.902088in}}%
\pgfpathlineto{\pgfqpoint{1.912098in}{0.867745in}}%
\pgfpathlineto{\pgfqpoint{1.878486in}{0.869447in}}%
\pgfpathlineto{\pgfqpoint{1.878275in}{0.865031in}}%
\pgfpathlineto{\pgfqpoint{1.868391in}{0.865576in}}%
\pgfpathlineto{\pgfqpoint{1.808729in}{0.868432in}}%
\pgfpathlineto{\pgfqpoint{1.810377in}{0.903512in}}%
\pgfpathlineto{\pgfqpoint{1.844714in}{0.901800in}}%
\pgfpathclose%
\pgfusepath{fill}%
\end{pgfscope}%
\begin{pgfscope}%
\pgfpathrectangle{\pgfqpoint{0.100000in}{0.100000in}}{\pgfqpoint{3.608454in}{2.310000in}}%
\pgfusepath{clip}%
\pgfsetbuttcap%
\pgfsetmiterjoin%
\definecolor{currentfill}{rgb}{0.000000,0.572549,0.713725}%
\pgfsetfillcolor{currentfill}%
\pgfsetlinewidth{0.000000pt}%
\definecolor{currentstroke}{rgb}{0.000000,0.000000,0.000000}%
\pgfsetstrokecolor{currentstroke}%
\pgfsetstrokeopacity{0.000000}%
\pgfsetdash{}{0pt}%
\pgfpathmoveto{\pgfqpoint{2.307376in}{1.370072in}}%
\pgfpathlineto{\pgfqpoint{2.271164in}{1.369568in}}%
\pgfpathlineto{\pgfqpoint{2.271150in}{1.361570in}}%
\pgfpathlineto{\pgfqpoint{2.264061in}{1.361534in}}%
\pgfpathlineto{\pgfqpoint{2.246959in}{1.366834in}}%
\pgfpathlineto{\pgfqpoint{2.247305in}{1.396610in}}%
\pgfpathlineto{\pgfqpoint{2.237052in}{1.396510in}}%
\pgfpathlineto{\pgfqpoint{2.237055in}{1.423181in}}%
\pgfpathlineto{\pgfqpoint{2.267438in}{1.423982in}}%
\pgfpathlineto{\pgfqpoint{2.267659in}{1.417346in}}%
\pgfpathlineto{\pgfqpoint{2.291417in}{1.417786in}}%
\pgfpathlineto{\pgfqpoint{2.298225in}{1.417924in}}%
\pgfpathlineto{\pgfqpoint{2.298632in}{1.395025in}}%
\pgfpathlineto{\pgfqpoint{2.306650in}{1.395239in}}%
\pgfpathlineto{\pgfqpoint{2.307376in}{1.370072in}}%
\pgfpathclose%
\pgfusepath{fill}%
\end{pgfscope}%
\begin{pgfscope}%
\pgfpathrectangle{\pgfqpoint{0.100000in}{0.100000in}}{\pgfqpoint{3.608454in}{2.310000in}}%
\pgfusepath{clip}%
\pgfsetbuttcap%
\pgfsetmiterjoin%
\definecolor{currentfill}{rgb}{0.000000,0.666667,0.666667}%
\pgfsetfillcolor{currentfill}%
\pgfsetlinewidth{0.000000pt}%
\definecolor{currentstroke}{rgb}{0.000000,0.000000,0.000000}%
\pgfsetstrokecolor{currentstroke}%
\pgfsetstrokeopacity{0.000000}%
\pgfsetdash{}{0pt}%
\pgfpathmoveto{\pgfqpoint{2.043077in}{1.035565in}}%
\pgfpathlineto{\pgfqpoint{2.043347in}{1.046522in}}%
\pgfpathlineto{\pgfqpoint{2.043425in}{1.053406in}}%
\pgfpathlineto{\pgfqpoint{2.050250in}{1.053326in}}%
\pgfpathlineto{\pgfqpoint{2.053684in}{1.060206in}}%
\pgfpathlineto{\pgfqpoint{2.053817in}{1.067155in}}%
\pgfpathlineto{\pgfqpoint{2.060863in}{1.067101in}}%
\pgfpathlineto{\pgfqpoint{2.060962in}{1.080885in}}%
\pgfpathlineto{\pgfqpoint{2.057547in}{1.080899in}}%
\pgfpathlineto{\pgfqpoint{2.057637in}{1.091246in}}%
\pgfpathlineto{\pgfqpoint{2.065063in}{1.091216in}}%
\pgfpathlineto{\pgfqpoint{2.068981in}{1.083880in}}%
\pgfpathlineto{\pgfqpoint{2.076132in}{1.086986in}}%
\pgfpathlineto{\pgfqpoint{2.089737in}{1.087648in}}%
\pgfpathlineto{\pgfqpoint{2.088937in}{1.099591in}}%
\pgfpathlineto{\pgfqpoint{2.093490in}{1.108334in}}%
\pgfpathlineto{\pgfqpoint{2.099232in}{1.108302in}}%
\pgfpathlineto{\pgfqpoint{2.099224in}{1.115204in}}%
\pgfpathlineto{\pgfqpoint{2.106071in}{1.115212in}}%
\pgfpathlineto{\pgfqpoint{2.134726in}{1.115276in}}%
\pgfpathlineto{\pgfqpoint{2.139356in}{1.083433in}}%
\pgfpathlineto{\pgfqpoint{2.140727in}{1.073891in}}%
\pgfpathlineto{\pgfqpoint{2.098590in}{1.073824in}}%
\pgfpathlineto{\pgfqpoint{2.098258in}{1.064900in}}%
\pgfpathlineto{\pgfqpoint{2.101602in}{1.057429in}}%
\pgfpathlineto{\pgfqpoint{2.096410in}{1.052482in}}%
\pgfpathlineto{\pgfqpoint{2.091976in}{1.044393in}}%
\pgfpathlineto{\pgfqpoint{2.081865in}{1.049099in}}%
\pgfpathlineto{\pgfqpoint{2.069896in}{1.042778in}}%
\pgfpathlineto{\pgfqpoint{2.058740in}{1.038661in}}%
\pgfpathlineto{\pgfqpoint{2.050477in}{1.038266in}}%
\pgfpathlineto{\pgfqpoint{2.043077in}{1.035565in}}%
\pgfpathclose%
\pgfusepath{fill}%
\end{pgfscope}%
\begin{pgfscope}%
\pgfpathrectangle{\pgfqpoint{0.100000in}{0.100000in}}{\pgfqpoint{3.608454in}{2.310000in}}%
\pgfusepath{clip}%
\pgfsetbuttcap%
\pgfsetmiterjoin%
\definecolor{currentfill}{rgb}{0.000000,0.627451,0.686275}%
\pgfsetfillcolor{currentfill}%
\pgfsetlinewidth{0.000000pt}%
\definecolor{currentstroke}{rgb}{0.000000,0.000000,0.000000}%
\pgfsetstrokecolor{currentstroke}%
\pgfsetstrokeopacity{0.000000}%
\pgfsetdash{}{0pt}%
\pgfpathmoveto{\pgfqpoint{2.088125in}{1.418722in}}%
\pgfpathlineto{\pgfqpoint{2.081282in}{1.423410in}}%
\pgfpathlineto{\pgfqpoint{2.082351in}{1.428944in}}%
\pgfpathlineto{\pgfqpoint{2.077853in}{1.433443in}}%
\pgfpathlineto{\pgfqpoint{2.077207in}{1.438452in}}%
\pgfpathlineto{\pgfqpoint{2.069544in}{1.443619in}}%
\pgfpathlineto{\pgfqpoint{2.060821in}{1.465149in}}%
\pgfpathlineto{\pgfqpoint{2.128892in}{1.463893in}}%
\pgfpathlineto{\pgfqpoint{2.130971in}{1.449316in}}%
\pgfpathlineto{\pgfqpoint{2.130612in}{1.428716in}}%
\pgfpathlineto{\pgfqpoint{2.104194in}{1.428993in}}%
\pgfpathlineto{\pgfqpoint{2.107194in}{1.411974in}}%
\pgfpathlineto{\pgfqpoint{2.101635in}{1.407776in}}%
\pgfpathlineto{\pgfqpoint{2.094548in}{1.411120in}}%
\pgfpathlineto{\pgfqpoint{2.088125in}{1.418722in}}%
\pgfpathclose%
\pgfusepath{fill}%
\end{pgfscope}%
\begin{pgfscope}%
\pgfpathrectangle{\pgfqpoint{0.100000in}{0.100000in}}{\pgfqpoint{3.608454in}{2.310000in}}%
\pgfusepath{clip}%
\pgfsetbuttcap%
\pgfsetmiterjoin%
\definecolor{currentfill}{rgb}{0.000000,0.733333,0.633333}%
\pgfsetfillcolor{currentfill}%
\pgfsetlinewidth{0.000000pt}%
\definecolor{currentstroke}{rgb}{0.000000,0.000000,0.000000}%
\pgfsetstrokecolor{currentstroke}%
\pgfsetstrokeopacity{0.000000}%
\pgfsetdash{}{0pt}%
\pgfpathmoveto{\pgfqpoint{2.722462in}{1.100629in}}%
\pgfpathlineto{\pgfqpoint{2.714180in}{1.097018in}}%
\pgfpathlineto{\pgfqpoint{2.697680in}{1.093767in}}%
\pgfpathlineto{\pgfqpoint{2.693772in}{1.094220in}}%
\pgfpathlineto{\pgfqpoint{2.688788in}{1.097666in}}%
\pgfpathlineto{\pgfqpoint{2.685775in}{1.104241in}}%
\pgfpathlineto{\pgfqpoint{2.685202in}{1.110108in}}%
\pgfpathlineto{\pgfqpoint{2.689025in}{1.121466in}}%
\pgfpathlineto{\pgfqpoint{2.681516in}{1.127762in}}%
\pgfpathlineto{\pgfqpoint{2.678100in}{1.135237in}}%
\pgfpathlineto{\pgfqpoint{2.683825in}{1.139915in}}%
\pgfpathlineto{\pgfqpoint{2.690045in}{1.139163in}}%
\pgfpathlineto{\pgfqpoint{2.693376in}{1.142517in}}%
\pgfpathlineto{\pgfqpoint{2.695803in}{1.138527in}}%
\pgfpathlineto{\pgfqpoint{2.701451in}{1.138405in}}%
\pgfpathlineto{\pgfqpoint{2.703971in}{1.134522in}}%
\pgfpathlineto{\pgfqpoint{2.712831in}{1.140583in}}%
\pgfpathlineto{\pgfqpoint{2.722754in}{1.138754in}}%
\pgfpathlineto{\pgfqpoint{2.729200in}{1.134130in}}%
\pgfpathlineto{\pgfqpoint{2.732727in}{1.126386in}}%
\pgfpathlineto{\pgfqpoint{2.731314in}{1.113832in}}%
\pgfpathlineto{\pgfqpoint{2.722462in}{1.100629in}}%
\pgfpathclose%
\pgfusepath{fill}%
\end{pgfscope}%
\begin{pgfscope}%
\pgfpathrectangle{\pgfqpoint{0.100000in}{0.100000in}}{\pgfqpoint{3.608454in}{2.310000in}}%
\pgfusepath{clip}%
\pgfsetbuttcap%
\pgfsetmiterjoin%
\definecolor{currentfill}{rgb}{0.000000,0.627451,0.686275}%
\pgfsetfillcolor{currentfill}%
\pgfsetlinewidth{0.000000pt}%
\definecolor{currentstroke}{rgb}{0.000000,0.000000,0.000000}%
\pgfsetstrokecolor{currentstroke}%
\pgfsetstrokeopacity{0.000000}%
\pgfsetdash{}{0pt}%
\pgfpathmoveto{\pgfqpoint{1.411936in}{2.184794in}}%
\pgfpathlineto{\pgfqpoint{1.473904in}{2.176145in}}%
\pgfpathlineto{\pgfqpoint{1.522840in}{2.169802in}}%
\pgfpathlineto{\pgfqpoint{1.519036in}{2.156103in}}%
\pgfpathlineto{\pgfqpoint{1.525884in}{2.155265in}}%
\pgfpathlineto{\pgfqpoint{1.523357in}{2.134506in}}%
\pgfpathlineto{\pgfqpoint{1.534783in}{2.133134in}}%
\pgfpathlineto{\pgfqpoint{1.531517in}{2.105747in}}%
\pgfpathlineto{\pgfqpoint{1.529442in}{2.105984in}}%
\pgfpathlineto{\pgfqpoint{1.527435in}{2.089391in}}%
\pgfpathlineto{\pgfqpoint{1.507101in}{2.093171in}}%
\pgfpathlineto{\pgfqpoint{1.498726in}{2.097310in}}%
\pgfpathlineto{\pgfqpoint{1.496354in}{2.088951in}}%
\pgfpathlineto{\pgfqpoint{1.492326in}{2.087153in}}%
\pgfpathlineto{\pgfqpoint{1.487031in}{2.078600in}}%
\pgfpathlineto{\pgfqpoint{1.476331in}{2.073148in}}%
\pgfpathlineto{\pgfqpoint{1.470183in}{2.074395in}}%
\pgfpathlineto{\pgfqpoint{1.468292in}{2.069917in}}%
\pgfpathlineto{\pgfqpoint{1.449108in}{2.071253in}}%
\pgfpathlineto{\pgfqpoint{1.440516in}{2.074870in}}%
\pgfpathlineto{\pgfqpoint{1.437875in}{2.069615in}}%
\pgfpathlineto{\pgfqpoint{1.428142in}{2.073167in}}%
\pgfpathlineto{\pgfqpoint{1.422183in}{2.067656in}}%
\pgfpathlineto{\pgfqpoint{1.414702in}{2.064290in}}%
\pgfpathlineto{\pgfqpoint{1.411450in}{2.059486in}}%
\pgfpathlineto{\pgfqpoint{1.411723in}{2.068893in}}%
\pgfpathlineto{\pgfqpoint{1.404841in}{2.072375in}}%
\pgfpathlineto{\pgfqpoint{1.389242in}{2.072875in}}%
\pgfpathlineto{\pgfqpoint{1.385004in}{2.076421in}}%
\pgfpathlineto{\pgfqpoint{1.377859in}{2.077250in}}%
\pgfpathlineto{\pgfqpoint{1.367934in}{2.081540in}}%
\pgfpathlineto{\pgfqpoint{1.362402in}{2.089710in}}%
\pgfpathlineto{\pgfqpoint{1.364847in}{2.104580in}}%
\pgfpathlineto{\pgfqpoint{1.378739in}{2.102305in}}%
\pgfpathlineto{\pgfqpoint{1.380515in}{2.107723in}}%
\pgfpathlineto{\pgfqpoint{1.389566in}{2.105191in}}%
\pgfpathlineto{\pgfqpoint{1.392437in}{2.124428in}}%
\pgfpathlineto{\pgfqpoint{1.396285in}{2.142117in}}%
\pgfpathlineto{\pgfqpoint{1.399456in}{2.139818in}}%
\pgfpathlineto{\pgfqpoint{1.402592in}{2.151198in}}%
\pgfpathlineto{\pgfqpoint{1.404613in}{2.164996in}}%
\pgfpathlineto{\pgfqpoint{1.407741in}{2.164527in}}%
\pgfpathlineto{\pgfqpoint{1.411936in}{2.184794in}}%
\pgfpathclose%
\pgfusepath{fill}%
\end{pgfscope}%
\begin{pgfscope}%
\pgfpathrectangle{\pgfqpoint{0.100000in}{0.100000in}}{\pgfqpoint{3.608454in}{2.310000in}}%
\pgfusepath{clip}%
\pgfsetbuttcap%
\pgfsetmiterjoin%
\definecolor{currentfill}{rgb}{0.000000,0.572549,0.713725}%
\pgfsetfillcolor{currentfill}%
\pgfsetlinewidth{0.000000pt}%
\definecolor{currentstroke}{rgb}{0.000000,0.000000,0.000000}%
\pgfsetstrokecolor{currentstroke}%
\pgfsetstrokeopacity{0.000000}%
\pgfsetdash{}{0pt}%
\pgfpathmoveto{\pgfqpoint{3.150978in}{0.513727in}}%
\pgfpathlineto{\pgfqpoint{3.118731in}{0.508688in}}%
\pgfpathlineto{\pgfqpoint{3.117607in}{0.515401in}}%
\pgfpathlineto{\pgfqpoint{3.110913in}{0.514275in}}%
\pgfpathlineto{\pgfqpoint{3.102843in}{0.570453in}}%
\pgfpathlineto{\pgfqpoint{3.099105in}{0.576837in}}%
\pgfpathlineto{\pgfqpoint{3.099275in}{0.583081in}}%
\pgfpathlineto{\pgfqpoint{3.090875in}{0.590177in}}%
\pgfpathlineto{\pgfqpoint{3.090036in}{0.604447in}}%
\pgfpathlineto{\pgfqpoint{3.106445in}{0.607053in}}%
\pgfpathlineto{\pgfqpoint{3.119882in}{0.592598in}}%
\pgfpathlineto{\pgfqpoint{3.125061in}{0.582905in}}%
\pgfpathlineto{\pgfqpoint{3.120768in}{0.574633in}}%
\pgfpathlineto{\pgfqpoint{3.121728in}{0.567411in}}%
\pgfpathlineto{\pgfqpoint{3.127445in}{0.552410in}}%
\pgfpathlineto{\pgfqpoint{3.144148in}{0.527548in}}%
\pgfpathlineto{\pgfqpoint{3.150978in}{0.513727in}}%
\pgfpathclose%
\pgfusepath{fill}%
\end{pgfscope}%
\begin{pgfscope}%
\pgfpathrectangle{\pgfqpoint{0.100000in}{0.100000in}}{\pgfqpoint{3.608454in}{2.310000in}}%
\pgfusepath{clip}%
\pgfsetbuttcap%
\pgfsetmiterjoin%
\definecolor{currentfill}{rgb}{0.000000,0.674510,0.662745}%
\pgfsetfillcolor{currentfill}%
\pgfsetlinewidth{0.000000pt}%
\definecolor{currentstroke}{rgb}{0.000000,0.000000,0.000000}%
\pgfsetstrokecolor{currentstroke}%
\pgfsetstrokeopacity{0.000000}%
\pgfsetdash{}{0pt}%
\pgfpathmoveto{\pgfqpoint{2.218257in}{1.006289in}}%
\pgfpathlineto{\pgfqpoint{2.211366in}{1.003801in}}%
\pgfpathlineto{\pgfqpoint{2.190786in}{1.003624in}}%
\pgfpathlineto{\pgfqpoint{2.190828in}{1.025374in}}%
\pgfpathlineto{\pgfqpoint{2.166850in}{1.025625in}}%
\pgfpathlineto{\pgfqpoint{2.162454in}{1.027888in}}%
\pgfpathlineto{\pgfqpoint{2.162548in}{1.034836in}}%
\pgfpathlineto{\pgfqpoint{2.166153in}{1.040540in}}%
\pgfpathlineto{\pgfqpoint{2.166477in}{1.058801in}}%
\pgfpathlineto{\pgfqpoint{2.169375in}{1.063715in}}%
\pgfpathlineto{\pgfqpoint{2.166287in}{1.069101in}}%
\pgfpathlineto{\pgfqpoint{2.166395in}{1.075969in}}%
\pgfpathlineto{\pgfqpoint{2.172137in}{1.075918in}}%
\pgfpathlineto{\pgfqpoint{2.176762in}{1.083808in}}%
\pgfpathlineto{\pgfqpoint{2.201915in}{1.084462in}}%
\pgfpathlineto{\pgfqpoint{2.201928in}{1.082158in}}%
\pgfpathlineto{\pgfqpoint{2.247558in}{1.082128in}}%
\pgfpathlineto{\pgfqpoint{2.247654in}{1.068396in}}%
\pgfpathlineto{\pgfqpoint{2.245419in}{1.061527in}}%
\pgfpathlineto{\pgfqpoint{2.245679in}{1.038542in}}%
\pgfpathlineto{\pgfqpoint{2.241035in}{1.037964in}}%
\pgfpathlineto{\pgfqpoint{2.226237in}{1.022520in}}%
\pgfpathlineto{\pgfqpoint{2.216940in}{1.014382in}}%
\pgfpathlineto{\pgfqpoint{2.218257in}{1.006289in}}%
\pgfpathclose%
\pgfusepath{fill}%
\end{pgfscope}%
\begin{pgfscope}%
\pgfpathrectangle{\pgfqpoint{0.100000in}{0.100000in}}{\pgfqpoint{3.608454in}{2.310000in}}%
\pgfusepath{clip}%
\pgfsetbuttcap%
\pgfsetmiterjoin%
\definecolor{currentfill}{rgb}{0.000000,0.411765,0.794118}%
\pgfsetfillcolor{currentfill}%
\pgfsetlinewidth{0.000000pt}%
\definecolor{currentstroke}{rgb}{0.000000,0.000000,0.000000}%
\pgfsetstrokecolor{currentstroke}%
\pgfsetstrokeopacity{0.000000}%
\pgfsetdash{}{0pt}%
\pgfpathmoveto{\pgfqpoint{1.075355in}{0.913752in}}%
\pgfpathlineto{\pgfqpoint{1.001421in}{0.926345in}}%
\pgfpathlineto{\pgfqpoint{0.934927in}{0.938858in}}%
\pgfpathlineto{\pgfqpoint{0.885192in}{0.948795in}}%
\pgfpathlineto{\pgfqpoint{0.899087in}{1.016739in}}%
\pgfpathlineto{\pgfqpoint{0.908989in}{1.065097in}}%
\pgfpathlineto{\pgfqpoint{0.947034in}{1.057499in}}%
\pgfpathlineto{\pgfqpoint{0.975479in}{1.042443in}}%
\pgfpathlineto{\pgfqpoint{0.985165in}{1.053983in}}%
\pgfpathlineto{\pgfqpoint{1.012835in}{1.044864in}}%
\pgfpathlineto{\pgfqpoint{1.027746in}{1.042105in}}%
\pgfpathlineto{\pgfqpoint{1.031005in}{1.053711in}}%
\pgfpathlineto{\pgfqpoint{1.015879in}{1.056517in}}%
\pgfpathlineto{\pgfqpoint{1.022683in}{1.074223in}}%
\pgfpathlineto{\pgfqpoint{1.031414in}{1.079449in}}%
\pgfpathlineto{\pgfqpoint{1.039948in}{1.073968in}}%
\pgfpathlineto{\pgfqpoint{1.049066in}{1.073939in}}%
\pgfpathlineto{\pgfqpoint{1.055533in}{1.068101in}}%
\pgfpathlineto{\pgfqpoint{1.063896in}{1.064728in}}%
\pgfpathlineto{\pgfqpoint{1.073305in}{1.055337in}}%
\pgfpathlineto{\pgfqpoint{1.079528in}{1.054057in}}%
\pgfpathlineto{\pgfqpoint{1.075926in}{1.033485in}}%
\pgfpathlineto{\pgfqpoint{1.124557in}{1.025178in}}%
\pgfpathlineto{\pgfqpoint{1.119039in}{0.992055in}}%
\pgfpathlineto{\pgfqpoint{1.117590in}{0.983356in}}%
\pgfpathlineto{\pgfqpoint{1.106555in}{0.985188in}}%
\pgfpathlineto{\pgfqpoint{1.096233in}{0.984306in}}%
\pgfpathlineto{\pgfqpoint{1.087666in}{0.981611in}}%
\pgfpathlineto{\pgfqpoint{1.084657in}{0.967062in}}%
\pgfpathlineto{\pgfqpoint{1.075355in}{0.913752in}}%
\pgfpathclose%
\pgfusepath{fill}%
\end{pgfscope}%
\begin{pgfscope}%
\pgfpathrectangle{\pgfqpoint{0.100000in}{0.100000in}}{\pgfqpoint{3.608454in}{2.310000in}}%
\pgfusepath{clip}%
\pgfsetbuttcap%
\pgfsetmiterjoin%
\definecolor{currentfill}{rgb}{0.000000,0.447059,0.776471}%
\pgfsetfillcolor{currentfill}%
\pgfsetlinewidth{0.000000pt}%
\definecolor{currentstroke}{rgb}{0.000000,0.000000,0.000000}%
\pgfsetstrokecolor{currentstroke}%
\pgfsetstrokeopacity{0.000000}%
\pgfsetdash{}{0pt}%
\pgfpathmoveto{\pgfqpoint{2.032756in}{2.134709in}}%
\pgfpathlineto{\pgfqpoint{2.098692in}{2.134056in}}%
\pgfpathlineto{\pgfqpoint{2.098766in}{2.165029in}}%
\pgfpathlineto{\pgfqpoint{2.103732in}{2.162515in}}%
\pgfpathlineto{\pgfqpoint{2.109017in}{2.163883in}}%
\pgfpathlineto{\pgfqpoint{2.115955in}{2.157789in}}%
\pgfpathlineto{\pgfqpoint{2.123516in}{2.124861in}}%
\pgfpathlineto{\pgfqpoint{2.123134in}{2.116351in}}%
\pgfpathlineto{\pgfqpoint{2.128642in}{2.111522in}}%
\pgfpathlineto{\pgfqpoint{2.136953in}{2.110228in}}%
\pgfpathlineto{\pgfqpoint{2.137206in}{2.083483in}}%
\pgfpathlineto{\pgfqpoint{2.095560in}{2.083291in}}%
\pgfpathlineto{\pgfqpoint{2.095537in}{2.097237in}}%
\pgfpathlineto{\pgfqpoint{2.074816in}{2.097239in}}%
\pgfpathlineto{\pgfqpoint{2.033124in}{2.098129in}}%
\pgfpathlineto{\pgfqpoint{2.032428in}{2.112015in}}%
\pgfpathlineto{\pgfqpoint{2.032756in}{2.134709in}}%
\pgfpathclose%
\pgfusepath{fill}%
\end{pgfscope}%
\begin{pgfscope}%
\pgfpathrectangle{\pgfqpoint{0.100000in}{0.100000in}}{\pgfqpoint{3.608454in}{2.310000in}}%
\pgfusepath{clip}%
\pgfsetbuttcap%
\pgfsetmiterjoin%
\definecolor{currentfill}{rgb}{0.000000,0.576471,0.711765}%
\pgfsetfillcolor{currentfill}%
\pgfsetlinewidth{0.000000pt}%
\definecolor{currentstroke}{rgb}{0.000000,0.000000,0.000000}%
\pgfsetstrokecolor{currentstroke}%
\pgfsetstrokeopacity{0.000000}%
\pgfsetdash{}{0pt}%
\pgfpathmoveto{\pgfqpoint{1.001953in}{2.138275in}}%
\pgfpathlineto{\pgfqpoint{0.996502in}{2.145181in}}%
\pgfpathlineto{\pgfqpoint{1.000929in}{2.147949in}}%
\pgfpathlineto{\pgfqpoint{1.000904in}{2.156205in}}%
\pgfpathlineto{\pgfqpoint{0.996594in}{2.161889in}}%
\pgfpathlineto{\pgfqpoint{0.988985in}{2.182170in}}%
\pgfpathlineto{\pgfqpoint{0.993310in}{2.200722in}}%
\pgfpathlineto{\pgfqpoint{0.994462in}{2.195826in}}%
\pgfpathlineto{\pgfqpoint{1.011635in}{2.199142in}}%
\pgfpathlineto{\pgfqpoint{1.012094in}{2.185135in}}%
\pgfpathlineto{\pgfqpoint{1.015735in}{2.176278in}}%
\pgfpathlineto{\pgfqpoint{1.015079in}{2.170487in}}%
\pgfpathlineto{\pgfqpoint{1.027896in}{2.166208in}}%
\pgfpathlineto{\pgfqpoint{1.033549in}{2.168208in}}%
\pgfpathlineto{\pgfqpoint{1.036396in}{2.174561in}}%
\pgfpathlineto{\pgfqpoint{1.043899in}{2.172864in}}%
\pgfpathlineto{\pgfqpoint{1.047061in}{2.189256in}}%
\pgfpathlineto{\pgfqpoint{1.056030in}{2.187273in}}%
\pgfpathlineto{\pgfqpoint{1.061972in}{2.214474in}}%
\pgfpathlineto{\pgfqpoint{1.061432in}{2.221573in}}%
\pgfpathlineto{\pgfqpoint{1.072055in}{2.221371in}}%
\pgfpathlineto{\pgfqpoint{1.075971in}{2.228263in}}%
\pgfpathlineto{\pgfqpoint{1.074794in}{2.233716in}}%
\pgfpathlineto{\pgfqpoint{1.075629in}{2.246526in}}%
\pgfpathlineto{\pgfqpoint{1.109601in}{2.239048in}}%
\pgfpathlineto{\pgfqpoint{1.109470in}{2.229400in}}%
\pgfpathlineto{\pgfqpoint{1.112359in}{2.223590in}}%
\pgfpathlineto{\pgfqpoint{1.118945in}{2.225334in}}%
\pgfpathlineto{\pgfqpoint{1.123862in}{2.213141in}}%
\pgfpathlineto{\pgfqpoint{1.119541in}{2.205769in}}%
\pgfpathlineto{\pgfqpoint{1.133270in}{2.197089in}}%
\pgfpathlineto{\pgfqpoint{1.131142in}{2.190334in}}%
\pgfpathlineto{\pgfqpoint{1.136982in}{2.181768in}}%
\pgfpathlineto{\pgfqpoint{1.134122in}{2.180102in}}%
\pgfpathlineto{\pgfqpoint{1.140465in}{2.171053in}}%
\pgfpathlineto{\pgfqpoint{1.139517in}{2.164977in}}%
\pgfpathlineto{\pgfqpoint{1.150149in}{2.159708in}}%
\pgfpathlineto{\pgfqpoint{1.154266in}{2.150464in}}%
\pgfpathlineto{\pgfqpoint{1.149026in}{2.145360in}}%
\pgfpathlineto{\pgfqpoint{1.143835in}{2.140807in}}%
\pgfpathlineto{\pgfqpoint{1.140741in}{2.130368in}}%
\pgfpathlineto{\pgfqpoint{1.135799in}{2.128987in}}%
\pgfpathlineto{\pgfqpoint{1.134932in}{2.119015in}}%
\pgfpathlineto{\pgfqpoint{1.118004in}{2.122768in}}%
\pgfpathlineto{\pgfqpoint{1.093543in}{2.127613in}}%
\pgfpathlineto{\pgfqpoint{1.090372in}{2.118413in}}%
\pgfpathlineto{\pgfqpoint{1.093234in}{2.110956in}}%
\pgfpathlineto{\pgfqpoint{1.091464in}{2.102885in}}%
\pgfpathlineto{\pgfqpoint{1.092364in}{2.093954in}}%
\pgfpathlineto{\pgfqpoint{1.085713in}{2.091817in}}%
\pgfpathlineto{\pgfqpoint{1.072235in}{2.094680in}}%
\pgfpathlineto{\pgfqpoint{1.068378in}{2.094075in}}%
\pgfpathlineto{\pgfqpoint{1.060904in}{2.103004in}}%
\pgfpathlineto{\pgfqpoint{1.048078in}{2.111068in}}%
\pgfpathlineto{\pgfqpoint{1.039741in}{2.112438in}}%
\pgfpathlineto{\pgfqpoint{1.033798in}{2.117353in}}%
\pgfpathlineto{\pgfqpoint{1.034696in}{2.123940in}}%
\pgfpathlineto{\pgfqpoint{1.020317in}{2.134665in}}%
\pgfpathlineto{\pgfqpoint{1.011451in}{2.134988in}}%
\pgfpathlineto{\pgfqpoint{1.001953in}{2.138275in}}%
\pgfpathclose%
\pgfusepath{fill}%
\end{pgfscope}%
\begin{pgfscope}%
\pgfpathrectangle{\pgfqpoint{0.100000in}{0.100000in}}{\pgfqpoint{3.608454in}{2.310000in}}%
\pgfusepath{clip}%
\pgfsetbuttcap%
\pgfsetmiterjoin%
\definecolor{currentfill}{rgb}{0.000000,0.666667,0.666667}%
\pgfsetfillcolor{currentfill}%
\pgfsetlinewidth{0.000000pt}%
\definecolor{currentstroke}{rgb}{0.000000,0.000000,0.000000}%
\pgfsetstrokecolor{currentstroke}%
\pgfsetstrokeopacity{0.000000}%
\pgfsetdash{}{0pt}%
\pgfpathmoveto{\pgfqpoint{1.583492in}{0.651351in}}%
\pgfpathlineto{\pgfqpoint{1.599353in}{0.636098in}}%
\pgfpathlineto{\pgfqpoint{1.593544in}{0.627739in}}%
\pgfpathlineto{\pgfqpoint{1.573827in}{0.627230in}}%
\pgfpathlineto{\pgfqpoint{1.567953in}{0.615750in}}%
\pgfpathlineto{\pgfqpoint{1.563163in}{0.610468in}}%
\pgfpathlineto{\pgfqpoint{1.559813in}{0.597646in}}%
\pgfpathlineto{\pgfqpoint{1.546004in}{0.580048in}}%
\pgfpathlineto{\pgfqpoint{1.539619in}{0.575381in}}%
\pgfpathlineto{\pgfqpoint{1.538032in}{0.569527in}}%
\pgfpathlineto{\pgfqpoint{1.529354in}{0.569946in}}%
\pgfpathlineto{\pgfqpoint{1.514568in}{0.578283in}}%
\pgfpathlineto{\pgfqpoint{1.508590in}{0.586063in}}%
\pgfpathlineto{\pgfqpoint{1.497563in}{0.588987in}}%
\pgfpathlineto{\pgfqpoint{1.493508in}{0.596165in}}%
\pgfpathlineto{\pgfqpoint{1.480379in}{0.599940in}}%
\pgfpathlineto{\pgfqpoint{1.469478in}{0.607802in}}%
\pgfpathlineto{\pgfqpoint{1.463905in}{0.617758in}}%
\pgfpathlineto{\pgfqpoint{1.457236in}{0.620260in}}%
\pgfpathlineto{\pgfqpoint{1.445703in}{0.631256in}}%
\pgfpathlineto{\pgfqpoint{1.443059in}{0.642585in}}%
\pgfpathlineto{\pgfqpoint{1.437063in}{0.653883in}}%
\pgfpathlineto{\pgfqpoint{1.436330in}{0.666295in}}%
\pgfpathlineto{\pgfqpoint{1.437587in}{0.679889in}}%
\pgfpathlineto{\pgfqpoint{1.428449in}{0.693343in}}%
\pgfpathlineto{\pgfqpoint{1.428817in}{0.701695in}}%
\pgfpathlineto{\pgfqpoint{1.425660in}{0.710940in}}%
\pgfpathlineto{\pgfqpoint{1.422243in}{0.713258in}}%
\pgfpathlineto{\pgfqpoint{1.426837in}{0.715528in}}%
\pgfpathlineto{\pgfqpoint{1.486032in}{0.744751in}}%
\pgfpathlineto{\pgfqpoint{1.518651in}{0.714404in}}%
\pgfpathlineto{\pgfqpoint{1.583492in}{0.651351in}}%
\pgfpathclose%
\pgfusepath{fill}%
\end{pgfscope}%
\begin{pgfscope}%
\pgfpathrectangle{\pgfqpoint{0.100000in}{0.100000in}}{\pgfqpoint{3.608454in}{2.310000in}}%
\pgfusepath{clip}%
\pgfsetbuttcap%
\pgfsetmiterjoin%
\definecolor{currentfill}{rgb}{0.000000,0.639216,0.680392}%
\pgfsetfillcolor{currentfill}%
\pgfsetlinewidth{0.000000pt}%
\definecolor{currentstroke}{rgb}{0.000000,0.000000,0.000000}%
\pgfsetstrokecolor{currentstroke}%
\pgfsetstrokeopacity{0.000000}%
\pgfsetdash{}{0pt}%
\pgfpathmoveto{\pgfqpoint{2.715165in}{0.972447in}}%
\pgfpathlineto{\pgfqpoint{2.705960in}{0.970065in}}%
\pgfpathlineto{\pgfqpoint{2.694783in}{0.960330in}}%
\pgfpathlineto{\pgfqpoint{2.689264in}{0.964347in}}%
\pgfpathlineto{\pgfqpoint{2.684888in}{0.971060in}}%
\pgfpathlineto{\pgfqpoint{2.678957in}{0.970072in}}%
\pgfpathlineto{\pgfqpoint{2.675715in}{0.974764in}}%
\pgfpathlineto{\pgfqpoint{2.677206in}{0.979100in}}%
\pgfpathlineto{\pgfqpoint{2.664313in}{0.994103in}}%
\pgfpathlineto{\pgfqpoint{2.657488in}{0.993619in}}%
\pgfpathlineto{\pgfqpoint{2.659190in}{1.009555in}}%
\pgfpathlineto{\pgfqpoint{2.666467in}{1.008380in}}%
\pgfpathlineto{\pgfqpoint{2.673359in}{1.013797in}}%
\pgfpathlineto{\pgfqpoint{2.669558in}{1.020872in}}%
\pgfpathlineto{\pgfqpoint{2.670176in}{1.036633in}}%
\pgfpathlineto{\pgfqpoint{2.668245in}{1.039940in}}%
\pgfpathlineto{\pgfqpoint{2.669947in}{1.049451in}}%
\pgfpathlineto{\pgfqpoint{2.698781in}{1.052062in}}%
\pgfpathlineto{\pgfqpoint{2.715608in}{1.053498in}}%
\pgfpathlineto{\pgfqpoint{2.723904in}{1.022807in}}%
\pgfpathlineto{\pgfqpoint{2.725316in}{1.017812in}}%
\pgfpathlineto{\pgfqpoint{2.721545in}{1.014250in}}%
\pgfpathlineto{\pgfqpoint{2.718615in}{1.004553in}}%
\pgfpathlineto{\pgfqpoint{2.710524in}{0.995150in}}%
\pgfpathlineto{\pgfqpoint{2.706101in}{0.993510in}}%
\pgfpathlineto{\pgfqpoint{2.715165in}{0.972447in}}%
\pgfpathclose%
\pgfusepath{fill}%
\end{pgfscope}%
\begin{pgfscope}%
\pgfpathrectangle{\pgfqpoint{0.100000in}{0.100000in}}{\pgfqpoint{3.608454in}{2.310000in}}%
\pgfusepath{clip}%
\pgfsetbuttcap%
\pgfsetmiterjoin%
\definecolor{currentfill}{rgb}{0.000000,0.945098,0.527451}%
\pgfsetfillcolor{currentfill}%
\pgfsetlinewidth{0.000000pt}%
\definecolor{currentstroke}{rgb}{0.000000,0.000000,0.000000}%
\pgfsetstrokecolor{currentstroke}%
\pgfsetstrokeopacity{0.000000}%
\pgfsetdash{}{0pt}%
\pgfpathmoveto{\pgfqpoint{0.605938in}{1.604509in}}%
\pgfpathlineto{\pgfqpoint{0.600479in}{1.609616in}}%
\pgfpathlineto{\pgfqpoint{0.604334in}{1.622797in}}%
\pgfpathlineto{\pgfqpoint{0.600541in}{1.629953in}}%
\pgfpathlineto{\pgfqpoint{0.601624in}{1.635606in}}%
\pgfpathlineto{\pgfqpoint{0.594599in}{1.640016in}}%
\pgfpathlineto{\pgfqpoint{0.584658in}{1.656876in}}%
\pgfpathlineto{\pgfqpoint{0.574470in}{1.662350in}}%
\pgfpathlineto{\pgfqpoint{0.566606in}{1.658522in}}%
\pgfpathlineto{\pgfqpoint{0.558505in}{1.659196in}}%
\pgfpathlineto{\pgfqpoint{0.555953in}{1.662869in}}%
\pgfpathlineto{\pgfqpoint{0.560209in}{1.677284in}}%
\pgfpathlineto{\pgfqpoint{0.544816in}{1.681752in}}%
\pgfpathlineto{\pgfqpoint{0.555668in}{1.716515in}}%
\pgfpathlineto{\pgfqpoint{0.561272in}{1.737785in}}%
\pgfpathlineto{\pgfqpoint{0.494828in}{1.757961in}}%
\pgfpathlineto{\pgfqpoint{0.498294in}{1.770162in}}%
\pgfpathlineto{\pgfqpoint{0.494062in}{1.772814in}}%
\pgfpathlineto{\pgfqpoint{0.477526in}{1.765034in}}%
\pgfpathlineto{\pgfqpoint{0.469737in}{1.765482in}}%
\pgfpathlineto{\pgfqpoint{0.466061in}{1.760966in}}%
\pgfpathlineto{\pgfqpoint{0.468505in}{1.753515in}}%
\pgfpathlineto{\pgfqpoint{0.459959in}{1.754157in}}%
\pgfpathlineto{\pgfqpoint{0.457583in}{1.760831in}}%
\pgfpathlineto{\pgfqpoint{0.449915in}{1.763500in}}%
\pgfpathlineto{\pgfqpoint{0.449403in}{1.767395in}}%
\pgfpathlineto{\pgfqpoint{0.443568in}{1.772460in}}%
\pgfpathlineto{\pgfqpoint{0.442827in}{1.791103in}}%
\pgfpathlineto{\pgfqpoint{0.433774in}{1.794008in}}%
\pgfpathlineto{\pgfqpoint{0.437233in}{1.797086in}}%
\pgfpathlineto{\pgfqpoint{0.438276in}{1.805265in}}%
\pgfpathlineto{\pgfqpoint{0.436698in}{1.812439in}}%
\pgfpathlineto{\pgfqpoint{0.441986in}{1.818955in}}%
\pgfpathlineto{\pgfqpoint{0.442382in}{1.828434in}}%
\pgfpathlineto{\pgfqpoint{0.452336in}{1.831556in}}%
\pgfpathlineto{\pgfqpoint{0.456983in}{1.837942in}}%
\pgfpathlineto{\pgfqpoint{0.526008in}{1.816628in}}%
\pgfpathlineto{\pgfqpoint{0.549594in}{1.891292in}}%
\pgfpathlineto{\pgfqpoint{0.551234in}{1.896670in}}%
\pgfpathlineto{\pgfqpoint{0.562002in}{1.893971in}}%
\pgfpathlineto{\pgfqpoint{0.565189in}{1.899442in}}%
\pgfpathlineto{\pgfqpoint{0.572519in}{1.906519in}}%
\pgfpathlineto{\pgfqpoint{0.572489in}{1.914545in}}%
\pgfpathlineto{\pgfqpoint{0.568235in}{1.922352in}}%
\pgfpathlineto{\pgfqpoint{0.571017in}{1.931177in}}%
\pgfpathlineto{\pgfqpoint{0.579411in}{1.933417in}}%
\pgfpathlineto{\pgfqpoint{0.616244in}{1.922404in}}%
\pgfpathlineto{\pgfqpoint{0.609494in}{1.902912in}}%
\pgfpathlineto{\pgfqpoint{0.595602in}{1.856614in}}%
\pgfpathlineto{\pgfqpoint{0.621562in}{1.848713in}}%
\pgfpathlineto{\pgfqpoint{0.605093in}{1.791654in}}%
\pgfpathlineto{\pgfqpoint{0.654872in}{1.777406in}}%
\pgfpathlineto{\pgfqpoint{0.637413in}{1.715698in}}%
\pgfpathlineto{\pgfqpoint{0.622947in}{1.663573in}}%
\pgfpathlineto{\pgfqpoint{0.605938in}{1.604509in}}%
\pgfpathclose%
\pgfusepath{fill}%
\end{pgfscope}%
\begin{pgfscope}%
\pgfpathrectangle{\pgfqpoint{0.100000in}{0.100000in}}{\pgfqpoint{3.608454in}{2.310000in}}%
\pgfusepath{clip}%
\pgfsetbuttcap%
\pgfsetmiterjoin%
\definecolor{currentfill}{rgb}{0.000000,0.580392,0.709804}%
\pgfsetfillcolor{currentfill}%
\pgfsetlinewidth{0.000000pt}%
\definecolor{currentstroke}{rgb}{0.000000,0.000000,0.000000}%
\pgfsetstrokecolor{currentstroke}%
\pgfsetstrokeopacity{0.000000}%
\pgfsetdash{}{0pt}%
\pgfpathmoveto{\pgfqpoint{2.636183in}{1.558287in}}%
\pgfpathlineto{\pgfqpoint{2.638514in}{1.537637in}}%
\pgfpathlineto{\pgfqpoint{2.614051in}{1.535158in}}%
\pgfpathlineto{\pgfqpoint{2.586580in}{1.532701in}}%
\pgfpathlineto{\pgfqpoint{2.582676in}{1.575047in}}%
\pgfpathlineto{\pgfqpoint{2.596356in}{1.586648in}}%
\pgfpathlineto{\pgfqpoint{2.601592in}{1.595258in}}%
\pgfpathlineto{\pgfqpoint{2.605772in}{1.609448in}}%
\pgfpathlineto{\pgfqpoint{2.612010in}{1.620134in}}%
\pgfpathlineto{\pgfqpoint{2.620303in}{1.620947in}}%
\pgfpathlineto{\pgfqpoint{2.621651in}{1.607409in}}%
\pgfpathlineto{\pgfqpoint{2.648568in}{1.609914in}}%
\pgfpathlineto{\pgfqpoint{2.650897in}{1.588647in}}%
\pgfpathlineto{\pgfqpoint{2.649386in}{1.585311in}}%
\pgfpathlineto{\pgfqpoint{2.633441in}{1.583836in}}%
\pgfpathlineto{\pgfqpoint{2.636183in}{1.558287in}}%
\pgfpathclose%
\pgfusepath{fill}%
\end{pgfscope}%
\begin{pgfscope}%
\pgfpathrectangle{\pgfqpoint{0.100000in}{0.100000in}}{\pgfqpoint{3.608454in}{2.310000in}}%
\pgfusepath{clip}%
\pgfsetbuttcap%
\pgfsetmiterjoin%
\definecolor{currentfill}{rgb}{0.000000,0.525490,0.737255}%
\pgfsetfillcolor{currentfill}%
\pgfsetlinewidth{0.000000pt}%
\definecolor{currentstroke}{rgb}{0.000000,0.000000,0.000000}%
\pgfsetstrokecolor{currentstroke}%
\pgfsetstrokeopacity{0.000000}%
\pgfsetdash{}{0pt}%
\pgfpathmoveto{\pgfqpoint{2.395693in}{0.767722in}}%
\pgfpathlineto{\pgfqpoint{2.396163in}{0.790004in}}%
\pgfpathlineto{\pgfqpoint{2.394278in}{0.809641in}}%
\pgfpathlineto{\pgfqpoint{2.395425in}{0.819431in}}%
\pgfpathlineto{\pgfqpoint{2.399936in}{0.821863in}}%
\pgfpathlineto{\pgfqpoint{2.401669in}{0.830402in}}%
\pgfpathlineto{\pgfqpoint{2.404935in}{0.832654in}}%
\pgfpathlineto{\pgfqpoint{2.412628in}{0.839279in}}%
\pgfpathlineto{\pgfqpoint{2.428148in}{0.846306in}}%
\pgfpathlineto{\pgfqpoint{2.438751in}{0.854506in}}%
\pgfpathlineto{\pgfqpoint{2.442684in}{0.864247in}}%
\pgfpathlineto{\pgfqpoint{2.458381in}{0.865552in}}%
\pgfpathlineto{\pgfqpoint{2.458164in}{0.869135in}}%
\pgfpathlineto{\pgfqpoint{2.485523in}{0.870923in}}%
\pgfpathlineto{\pgfqpoint{2.488805in}{0.814690in}}%
\pgfpathlineto{\pgfqpoint{2.461411in}{0.812905in}}%
\pgfpathlineto{\pgfqpoint{2.462239in}{0.799019in}}%
\pgfpathlineto{\pgfqpoint{2.467492in}{0.796662in}}%
\pgfpathlineto{\pgfqpoint{2.468672in}{0.778139in}}%
\pgfpathlineto{\pgfqpoint{2.447021in}{0.775387in}}%
\pgfpathlineto{\pgfqpoint{2.436793in}{0.774096in}}%
\pgfpathlineto{\pgfqpoint{2.426622in}{0.769362in}}%
\pgfpathlineto{\pgfqpoint{2.395693in}{0.767722in}}%
\pgfpathclose%
\pgfusepath{fill}%
\end{pgfscope}%
\begin{pgfscope}%
\pgfpathrectangle{\pgfqpoint{0.100000in}{0.100000in}}{\pgfqpoint{3.608454in}{2.310000in}}%
\pgfusepath{clip}%
\pgfsetbuttcap%
\pgfsetmiterjoin%
\definecolor{currentfill}{rgb}{0.000000,0.666667,0.666667}%
\pgfsetfillcolor{currentfill}%
\pgfsetlinewidth{0.000000pt}%
\definecolor{currentstroke}{rgb}{0.000000,0.000000,0.000000}%
\pgfsetstrokecolor{currentstroke}%
\pgfsetstrokeopacity{0.000000}%
\pgfsetdash{}{0pt}%
\pgfpathmoveto{\pgfqpoint{2.244140in}{0.763307in}}%
\pgfpathlineto{\pgfqpoint{2.248851in}{0.756615in}}%
\pgfpathlineto{\pgfqpoint{2.256019in}{0.754662in}}%
\pgfpathlineto{\pgfqpoint{2.261771in}{0.748472in}}%
\pgfpathlineto{\pgfqpoint{2.248864in}{0.736251in}}%
\pgfpathlineto{\pgfqpoint{2.244244in}{0.734424in}}%
\pgfpathlineto{\pgfqpoint{2.240752in}{0.735755in}}%
\pgfpathlineto{\pgfqpoint{2.226879in}{0.735487in}}%
\pgfpathlineto{\pgfqpoint{2.226572in}{0.749343in}}%
\pgfpathlineto{\pgfqpoint{2.219726in}{0.749195in}}%
\pgfpathlineto{\pgfqpoint{2.219367in}{0.763229in}}%
\pgfpathlineto{\pgfqpoint{2.212402in}{0.763123in}}%
\pgfpathlineto{\pgfqpoint{2.212262in}{0.773539in}}%
\pgfpathlineto{\pgfqpoint{2.182637in}{0.773027in}}%
\pgfpathlineto{\pgfqpoint{2.173280in}{0.783716in}}%
\pgfpathlineto{\pgfqpoint{2.171444in}{0.784825in}}%
\pgfpathlineto{\pgfqpoint{2.170684in}{0.855388in}}%
\pgfpathlineto{\pgfqpoint{2.170562in}{0.866376in}}%
\pgfpathlineto{\pgfqpoint{2.186423in}{0.866575in}}%
\pgfpathlineto{\pgfqpoint{2.240807in}{0.867381in}}%
\pgfpathlineto{\pgfqpoint{2.258381in}{0.867551in}}%
\pgfpathlineto{\pgfqpoint{2.258858in}{0.847301in}}%
\pgfpathlineto{\pgfqpoint{2.251967in}{0.847074in}}%
\pgfpathlineto{\pgfqpoint{2.248708in}{0.834367in}}%
\pgfpathlineto{\pgfqpoint{2.249143in}{0.822767in}}%
\pgfpathlineto{\pgfqpoint{2.255994in}{0.822878in}}%
\pgfpathlineto{\pgfqpoint{2.256664in}{0.805503in}}%
\pgfpathlineto{\pgfqpoint{2.254067in}{0.798428in}}%
\pgfpathlineto{\pgfqpoint{2.245714in}{0.798315in}}%
\pgfpathlineto{\pgfqpoint{2.249736in}{0.781810in}}%
\pgfpathlineto{\pgfqpoint{2.246031in}{0.774913in}}%
\pgfpathlineto{\pgfqpoint{2.244140in}{0.763307in}}%
\pgfpathclose%
\pgfusepath{fill}%
\end{pgfscope}%
\begin{pgfscope}%
\pgfpathrectangle{\pgfqpoint{0.100000in}{0.100000in}}{\pgfqpoint{3.608454in}{2.310000in}}%
\pgfusepath{clip}%
\pgfsetbuttcap%
\pgfsetmiterjoin%
\definecolor{currentfill}{rgb}{0.000000,0.541176,0.729412}%
\pgfsetfillcolor{currentfill}%
\pgfsetlinewidth{0.000000pt}%
\definecolor{currentstroke}{rgb}{0.000000,0.000000,0.000000}%
\pgfsetstrokecolor{currentstroke}%
\pgfsetstrokeopacity{0.000000}%
\pgfsetdash{}{0pt}%
\pgfpathmoveto{\pgfqpoint{2.859932in}{1.127064in}}%
\pgfpathlineto{\pgfqpoint{2.845692in}{1.114629in}}%
\pgfpathlineto{\pgfqpoint{2.835111in}{1.113784in}}%
\pgfpathlineto{\pgfqpoint{2.833923in}{1.123330in}}%
\pgfpathlineto{\pgfqpoint{2.829784in}{1.127369in}}%
\pgfpathlineto{\pgfqpoint{2.823619in}{1.137722in}}%
\pgfpathlineto{\pgfqpoint{2.831674in}{1.145493in}}%
\pgfpathlineto{\pgfqpoint{2.824792in}{1.161761in}}%
\pgfpathlineto{\pgfqpoint{2.827195in}{1.174638in}}%
\pgfpathlineto{\pgfqpoint{2.836596in}{1.177507in}}%
\pgfpathlineto{\pgfqpoint{2.846368in}{1.183090in}}%
\pgfpathlineto{\pgfqpoint{2.851139in}{1.181704in}}%
\pgfpathlineto{\pgfqpoint{2.853813in}{1.173625in}}%
\pgfpathlineto{\pgfqpoint{2.859396in}{1.177849in}}%
\pgfpathlineto{\pgfqpoint{2.865237in}{1.171385in}}%
\pgfpathlineto{\pgfqpoint{2.860386in}{1.164003in}}%
\pgfpathlineto{\pgfqpoint{2.868331in}{1.159628in}}%
\pgfpathlineto{\pgfqpoint{2.879988in}{1.149617in}}%
\pgfpathlineto{\pgfqpoint{2.877093in}{1.136816in}}%
\pgfpathlineto{\pgfqpoint{2.865280in}{1.133278in}}%
\pgfpathlineto{\pgfqpoint{2.859932in}{1.127064in}}%
\pgfpathclose%
\pgfusepath{fill}%
\end{pgfscope}%
\begin{pgfscope}%
\pgfpathrectangle{\pgfqpoint{0.100000in}{0.100000in}}{\pgfqpoint{3.608454in}{2.310000in}}%
\pgfusepath{clip}%
\pgfsetbuttcap%
\pgfsetmiterjoin%
\definecolor{currentfill}{rgb}{0.000000,0.439216,0.780392}%
\pgfsetfillcolor{currentfill}%
\pgfsetlinewidth{0.000000pt}%
\definecolor{currentstroke}{rgb}{0.000000,0.000000,0.000000}%
\pgfsetstrokecolor{currentstroke}%
\pgfsetstrokeopacity{0.000000}%
\pgfsetdash{}{0pt}%
\pgfpathmoveto{\pgfqpoint{1.291011in}{1.784895in}}%
\pgfpathlineto{\pgfqpoint{1.303529in}{1.783157in}}%
\pgfpathlineto{\pgfqpoint{1.304569in}{1.790081in}}%
\pgfpathlineto{\pgfqpoint{1.314743in}{1.788523in}}%
\pgfpathlineto{\pgfqpoint{1.315604in}{1.795301in}}%
\pgfpathlineto{\pgfqpoint{1.322614in}{1.794241in}}%
\pgfpathlineto{\pgfqpoint{1.323648in}{1.801191in}}%
\pgfpathlineto{\pgfqpoint{1.337257in}{1.799111in}}%
\pgfpathlineto{\pgfqpoint{1.338318in}{1.806097in}}%
\pgfpathlineto{\pgfqpoint{1.417395in}{1.794618in}}%
\pgfpathlineto{\pgfqpoint{1.412417in}{1.742036in}}%
\pgfpathlineto{\pgfqpoint{1.388262in}{1.745418in}}%
\pgfpathlineto{\pgfqpoint{1.384511in}{1.743611in}}%
\pgfpathlineto{\pgfqpoint{1.352534in}{1.748041in}}%
\pgfpathlineto{\pgfqpoint{1.342072in}{1.748707in}}%
\pgfpathlineto{\pgfqpoint{1.322608in}{1.757434in}}%
\pgfpathlineto{\pgfqpoint{1.323149in}{1.760841in}}%
\pgfpathlineto{\pgfqpoint{1.310078in}{1.766325in}}%
\pgfpathlineto{\pgfqpoint{1.296615in}{1.768471in}}%
\pgfpathlineto{\pgfqpoint{1.297605in}{1.774698in}}%
\pgfpathlineto{\pgfqpoint{1.291011in}{1.784895in}}%
\pgfpathclose%
\pgfusepath{fill}%
\end{pgfscope}%
\begin{pgfscope}%
\pgfpathrectangle{\pgfqpoint{0.100000in}{0.100000in}}{\pgfqpoint{3.608454in}{2.310000in}}%
\pgfusepath{clip}%
\pgfsetbuttcap%
\pgfsetmiterjoin%
\definecolor{currentfill}{rgb}{0.000000,0.411765,0.794118}%
\pgfsetfillcolor{currentfill}%
\pgfsetlinewidth{0.000000pt}%
\definecolor{currentstroke}{rgb}{0.000000,0.000000,0.000000}%
\pgfsetstrokecolor{currentstroke}%
\pgfsetstrokeopacity{0.000000}%
\pgfsetdash{}{0pt}%
\pgfpathmoveto{\pgfqpoint{2.522980in}{1.660007in}}%
\pgfpathlineto{\pgfqpoint{2.467930in}{1.655975in}}%
\pgfpathlineto{\pgfqpoint{2.454210in}{1.655405in}}%
\pgfpathlineto{\pgfqpoint{2.450492in}{1.717550in}}%
\pgfpathlineto{\pgfqpoint{2.485186in}{1.719799in}}%
\pgfpathlineto{\pgfqpoint{2.485729in}{1.712888in}}%
\pgfpathlineto{\pgfqpoint{2.499443in}{1.713874in}}%
\pgfpathlineto{\pgfqpoint{2.520499in}{1.715478in}}%
\pgfpathlineto{\pgfqpoint{2.520327in}{1.709190in}}%
\pgfpathlineto{\pgfqpoint{2.517313in}{1.703277in}}%
\pgfpathlineto{\pgfqpoint{2.515606in}{1.691180in}}%
\pgfpathlineto{\pgfqpoint{2.522980in}{1.660007in}}%
\pgfpathclose%
\pgfusepath{fill}%
\end{pgfscope}%
\begin{pgfscope}%
\pgfpathrectangle{\pgfqpoint{0.100000in}{0.100000in}}{\pgfqpoint{3.608454in}{2.310000in}}%
\pgfusepath{clip}%
\pgfsetbuttcap%
\pgfsetmiterjoin%
\definecolor{currentfill}{rgb}{0.000000,0.419608,0.790196}%
\pgfsetfillcolor{currentfill}%
\pgfsetlinewidth{0.000000pt}%
\definecolor{currentstroke}{rgb}{0.000000,0.000000,0.000000}%
\pgfsetstrokecolor{currentstroke}%
\pgfsetstrokeopacity{0.000000}%
\pgfsetdash{}{0pt}%
\pgfpathmoveto{\pgfqpoint{0.721788in}{1.059750in}}%
\pgfpathlineto{\pgfqpoint{0.715639in}{1.032543in}}%
\pgfpathlineto{\pgfqpoint{0.714200in}{1.032869in}}%
\pgfpathlineto{\pgfqpoint{0.705665in}{0.997629in}}%
\pgfpathlineto{\pgfqpoint{0.637733in}{1.007252in}}%
\pgfpathlineto{\pgfqpoint{0.638583in}{1.013548in}}%
\pgfpathlineto{\pgfqpoint{0.632350in}{1.019382in}}%
\pgfpathlineto{\pgfqpoint{0.635822in}{1.034911in}}%
\pgfpathlineto{\pgfqpoint{0.636748in}{1.045662in}}%
\pgfpathlineto{\pgfqpoint{0.634791in}{1.058910in}}%
\pgfpathlineto{\pgfqpoint{0.628784in}{1.074472in}}%
\pgfpathlineto{\pgfqpoint{0.623749in}{1.080645in}}%
\pgfpathlineto{\pgfqpoint{0.626202in}{1.085453in}}%
\pgfpathlineto{\pgfqpoint{0.631612in}{1.088268in}}%
\pgfpathlineto{\pgfqpoint{0.640257in}{1.084919in}}%
\pgfpathlineto{\pgfqpoint{0.647432in}{1.078308in}}%
\pgfpathlineto{\pgfqpoint{0.660882in}{1.074556in}}%
\pgfpathlineto{\pgfqpoint{0.721788in}{1.059750in}}%
\pgfpathclose%
\pgfusepath{fill}%
\end{pgfscope}%
\begin{pgfscope}%
\pgfpathrectangle{\pgfqpoint{0.100000in}{0.100000in}}{\pgfqpoint{3.608454in}{2.310000in}}%
\pgfusepath{clip}%
\pgfsetbuttcap%
\pgfsetmiterjoin%
\definecolor{currentfill}{rgb}{0.000000,0.713725,0.643137}%
\pgfsetfillcolor{currentfill}%
\pgfsetlinewidth{0.000000pt}%
\definecolor{currentstroke}{rgb}{0.000000,0.000000,0.000000}%
\pgfsetstrokecolor{currentstroke}%
\pgfsetstrokeopacity{0.000000}%
\pgfsetdash{}{0pt}%
\pgfpathmoveto{\pgfqpoint{2.380907in}{1.900334in}}%
\pgfpathlineto{\pgfqpoint{2.366647in}{1.899528in}}%
\pgfpathlineto{\pgfqpoint{2.365878in}{1.913327in}}%
\pgfpathlineto{\pgfqpoint{2.359091in}{1.912916in}}%
\pgfpathlineto{\pgfqpoint{2.358761in}{1.919792in}}%
\pgfpathlineto{\pgfqpoint{2.351813in}{1.919479in}}%
\pgfpathlineto{\pgfqpoint{2.350608in}{1.946791in}}%
\pgfpathlineto{\pgfqpoint{2.356731in}{1.945339in}}%
\pgfpathlineto{\pgfqpoint{2.361222in}{1.948386in}}%
\pgfpathlineto{\pgfqpoint{2.378630in}{1.955337in}}%
\pgfpathlineto{\pgfqpoint{2.391147in}{1.967641in}}%
\pgfpathlineto{\pgfqpoint{2.411560in}{1.970914in}}%
\pgfpathlineto{\pgfqpoint{2.420196in}{1.976038in}}%
\pgfpathlineto{\pgfqpoint{2.425368in}{1.983108in}}%
\pgfpathlineto{\pgfqpoint{2.433892in}{1.984602in}}%
\pgfpathlineto{\pgfqpoint{2.436662in}{1.987686in}}%
\pgfpathlineto{\pgfqpoint{2.438144in}{1.966329in}}%
\pgfpathlineto{\pgfqpoint{2.442299in}{1.959772in}}%
\pgfpathlineto{\pgfqpoint{2.435480in}{1.959378in}}%
\pgfpathlineto{\pgfqpoint{2.436761in}{1.938649in}}%
\pgfpathlineto{\pgfqpoint{2.438472in}{1.913015in}}%
\pgfpathlineto{\pgfqpoint{2.432701in}{1.915926in}}%
\pgfpathlineto{\pgfqpoint{2.385839in}{1.925963in}}%
\pgfpathlineto{\pgfqpoint{2.387268in}{1.900687in}}%
\pgfpathlineto{\pgfqpoint{2.380907in}{1.900334in}}%
\pgfpathclose%
\pgfusepath{fill}%
\end{pgfscope}%
\begin{pgfscope}%
\pgfpathrectangle{\pgfqpoint{0.100000in}{0.100000in}}{\pgfqpoint{3.608454in}{2.310000in}}%
\pgfusepath{clip}%
\pgfsetbuttcap%
\pgfsetmiterjoin%
\definecolor{currentfill}{rgb}{0.000000,0.403922,0.798039}%
\pgfsetfillcolor{currentfill}%
\pgfsetlinewidth{0.000000pt}%
\definecolor{currentstroke}{rgb}{0.000000,0.000000,0.000000}%
\pgfsetstrokecolor{currentstroke}%
\pgfsetstrokeopacity{0.000000}%
\pgfsetdash{}{0pt}%
\pgfpathmoveto{\pgfqpoint{1.577074in}{1.409804in}}%
\pgfpathlineto{\pgfqpoint{1.610603in}{1.406638in}}%
\pgfpathlineto{\pgfqpoint{1.606968in}{1.372189in}}%
\pgfpathlineto{\pgfqpoint{1.603799in}{1.337764in}}%
\pgfpathlineto{\pgfqpoint{1.602094in}{1.324825in}}%
\pgfpathlineto{\pgfqpoint{1.581597in}{1.326016in}}%
\pgfpathlineto{\pgfqpoint{1.547834in}{1.329763in}}%
\pgfpathlineto{\pgfqpoint{1.550467in}{1.357020in}}%
\pgfpathlineto{\pgfqpoint{1.571014in}{1.354845in}}%
\pgfpathlineto{\pgfqpoint{1.577074in}{1.409804in}}%
\pgfpathclose%
\pgfusepath{fill}%
\end{pgfscope}%
\begin{pgfscope}%
\pgfpathrectangle{\pgfqpoint{0.100000in}{0.100000in}}{\pgfqpoint{3.608454in}{2.310000in}}%
\pgfusepath{clip}%
\pgfsetbuttcap%
\pgfsetmiterjoin%
\definecolor{currentfill}{rgb}{0.000000,0.537255,0.731373}%
\pgfsetfillcolor{currentfill}%
\pgfsetlinewidth{0.000000pt}%
\definecolor{currentstroke}{rgb}{0.000000,0.000000,0.000000}%
\pgfsetstrokecolor{currentstroke}%
\pgfsetstrokeopacity{0.000000}%
\pgfsetdash{}{0pt}%
\pgfpathmoveto{\pgfqpoint{2.101955in}{1.263584in}}%
\pgfpathlineto{\pgfqpoint{2.101288in}{1.239463in}}%
\pgfpathlineto{\pgfqpoint{2.046716in}{1.239913in}}%
\pgfpathlineto{\pgfqpoint{2.047067in}{1.264108in}}%
\pgfpathlineto{\pgfqpoint{2.074465in}{1.263717in}}%
\pgfpathlineto{\pgfqpoint{2.075326in}{1.295024in}}%
\pgfpathlineto{\pgfqpoint{2.075903in}{1.319066in}}%
\pgfpathlineto{\pgfqpoint{2.103312in}{1.318980in}}%
\pgfpathlineto{\pgfqpoint{2.102737in}{1.315533in}}%
\pgfpathlineto{\pgfqpoint{2.102752in}{1.281057in}}%
\pgfpathlineto{\pgfqpoint{2.101955in}{1.263584in}}%
\pgfpathclose%
\pgfusepath{fill}%
\end{pgfscope}%
\begin{pgfscope}%
\pgfpathrectangle{\pgfqpoint{0.100000in}{0.100000in}}{\pgfqpoint{3.608454in}{2.310000in}}%
\pgfusepath{clip}%
\pgfsetbuttcap%
\pgfsetmiterjoin%
\definecolor{currentfill}{rgb}{0.000000,0.517647,0.741176}%
\pgfsetfillcolor{currentfill}%
\pgfsetlinewidth{0.000000pt}%
\definecolor{currentstroke}{rgb}{0.000000,0.000000,0.000000}%
\pgfsetstrokecolor{currentstroke}%
\pgfsetstrokeopacity{0.000000}%
\pgfsetdash{}{0pt}%
\pgfpathmoveto{\pgfqpoint{2.356072in}{1.695517in}}%
\pgfpathlineto{\pgfqpoint{2.325187in}{1.694132in}}%
\pgfpathlineto{\pgfqpoint{2.324231in}{1.700249in}}%
\pgfpathlineto{\pgfqpoint{2.301639in}{1.699356in}}%
\pgfpathlineto{\pgfqpoint{2.294773in}{1.699116in}}%
\pgfpathlineto{\pgfqpoint{2.293781in}{1.726557in}}%
\pgfpathlineto{\pgfqpoint{2.319291in}{1.727547in}}%
\pgfpathlineto{\pgfqpoint{2.310801in}{1.738104in}}%
\pgfpathlineto{\pgfqpoint{2.316678in}{1.738350in}}%
\pgfpathlineto{\pgfqpoint{2.317255in}{1.744761in}}%
\pgfpathlineto{\pgfqpoint{2.336281in}{1.746036in}}%
\pgfpathlineto{\pgfqpoint{2.339782in}{1.753146in}}%
\pgfpathlineto{\pgfqpoint{2.373549in}{1.754537in}}%
\pgfpathlineto{\pgfqpoint{2.373196in}{1.761948in}}%
\pgfpathlineto{\pgfqpoint{2.396439in}{1.763257in}}%
\pgfpathlineto{\pgfqpoint{2.389968in}{1.750300in}}%
\pgfpathlineto{\pgfqpoint{2.394954in}{1.739243in}}%
\pgfpathlineto{\pgfqpoint{2.395001in}{1.732218in}}%
\pgfpathlineto{\pgfqpoint{2.400634in}{1.726438in}}%
\pgfpathlineto{\pgfqpoint{2.405834in}{1.715444in}}%
\pgfpathlineto{\pgfqpoint{2.375661in}{1.713794in}}%
\pgfpathlineto{\pgfqpoint{2.376100in}{1.706906in}}%
\pgfpathlineto{\pgfqpoint{2.355403in}{1.705801in}}%
\pgfpathlineto{\pgfqpoint{2.356072in}{1.695517in}}%
\pgfpathclose%
\pgfusepath{fill}%
\end{pgfscope}%
\begin{pgfscope}%
\pgfpathrectangle{\pgfqpoint{0.100000in}{0.100000in}}{\pgfqpoint{3.608454in}{2.310000in}}%
\pgfusepath{clip}%
\pgfsetbuttcap%
\pgfsetmiterjoin%
\definecolor{currentfill}{rgb}{0.000000,0.603922,0.698039}%
\pgfsetfillcolor{currentfill}%
\pgfsetlinewidth{0.000000pt}%
\definecolor{currentstroke}{rgb}{0.000000,0.000000,0.000000}%
\pgfsetstrokecolor{currentstroke}%
\pgfsetstrokeopacity{0.000000}%
\pgfsetdash{}{0pt}%
\pgfpathmoveto{\pgfqpoint{1.234540in}{2.143503in}}%
\pgfpathlineto{\pgfqpoint{1.232482in}{2.132070in}}%
\pgfpathlineto{\pgfqpoint{1.202148in}{2.137522in}}%
\pgfpathlineto{\pgfqpoint{1.200330in}{2.146187in}}%
\pgfpathlineto{\pgfqpoint{1.194140in}{2.150958in}}%
\pgfpathlineto{\pgfqpoint{1.150149in}{2.159708in}}%
\pgfpathlineto{\pgfqpoint{1.139517in}{2.164977in}}%
\pgfpathlineto{\pgfqpoint{1.140465in}{2.171053in}}%
\pgfpathlineto{\pgfqpoint{1.134122in}{2.180102in}}%
\pgfpathlineto{\pgfqpoint{1.136982in}{2.181768in}}%
\pgfpathlineto{\pgfqpoint{1.131142in}{2.190334in}}%
\pgfpathlineto{\pgfqpoint{1.133270in}{2.197089in}}%
\pgfpathlineto{\pgfqpoint{1.119541in}{2.205769in}}%
\pgfpathlineto{\pgfqpoint{1.123862in}{2.213141in}}%
\pgfpathlineto{\pgfqpoint{1.118945in}{2.225334in}}%
\pgfpathlineto{\pgfqpoint{1.112359in}{2.223590in}}%
\pgfpathlineto{\pgfqpoint{1.109470in}{2.229400in}}%
\pgfpathlineto{\pgfqpoint{1.109601in}{2.239048in}}%
\pgfpathlineto{\pgfqpoint{1.157261in}{2.229106in}}%
\pgfpathlineto{\pgfqpoint{1.193179in}{2.222014in}}%
\pgfpathlineto{\pgfqpoint{1.254169in}{2.210455in}}%
\pgfpathlineto{\pgfqpoint{1.242847in}{2.149099in}}%
\pgfpathlineto{\pgfqpoint{1.235802in}{2.150383in}}%
\pgfpathlineto{\pgfqpoint{1.234540in}{2.143503in}}%
\pgfpathclose%
\pgfusepath{fill}%
\end{pgfscope}%
\begin{pgfscope}%
\pgfpathrectangle{\pgfqpoint{0.100000in}{0.100000in}}{\pgfqpoint{3.608454in}{2.310000in}}%
\pgfusepath{clip}%
\pgfsetbuttcap%
\pgfsetmiterjoin%
\definecolor{currentfill}{rgb}{0.000000,0.380392,0.809804}%
\pgfsetfillcolor{currentfill}%
\pgfsetlinewidth{0.000000pt}%
\definecolor{currentstroke}{rgb}{0.000000,0.000000,0.000000}%
\pgfsetstrokecolor{currentstroke}%
\pgfsetstrokeopacity{0.000000}%
\pgfsetdash{}{0pt}%
\pgfpathmoveto{\pgfqpoint{1.807069in}{2.050968in}}%
\pgfpathlineto{\pgfqpoint{1.738152in}{2.055718in}}%
\pgfpathlineto{\pgfqpoint{1.739211in}{2.069716in}}%
\pgfpathlineto{\pgfqpoint{1.736501in}{2.069931in}}%
\pgfpathlineto{\pgfqpoint{1.738640in}{2.097811in}}%
\pgfpathlineto{\pgfqpoint{1.735922in}{2.098000in}}%
\pgfpathlineto{\pgfqpoint{1.737002in}{2.111989in}}%
\pgfpathlineto{\pgfqpoint{1.723190in}{2.112972in}}%
\pgfpathlineto{\pgfqpoint{1.724331in}{2.126897in}}%
\pgfpathlineto{\pgfqpoint{1.728606in}{2.126565in}}%
\pgfpathlineto{\pgfqpoint{1.729180in}{2.133512in}}%
\pgfpathlineto{\pgfqpoint{1.736050in}{2.132965in}}%
\pgfpathlineto{\pgfqpoint{1.737269in}{2.148337in}}%
\pgfpathlineto{\pgfqpoint{1.790039in}{2.144561in}}%
\pgfpathlineto{\pgfqpoint{1.833904in}{2.141855in}}%
\pgfpathlineto{\pgfqpoint{1.832606in}{2.119399in}}%
\pgfpathlineto{\pgfqpoint{1.833772in}{2.105376in}}%
\pgfpathlineto{\pgfqpoint{1.826839in}{2.105738in}}%
\pgfpathlineto{\pgfqpoint{1.828130in}{2.091731in}}%
\pgfpathlineto{\pgfqpoint{1.826539in}{2.063794in}}%
\pgfpathlineto{\pgfqpoint{1.827828in}{2.049744in}}%
\pgfpathlineto{\pgfqpoint{1.807069in}{2.050968in}}%
\pgfpathclose%
\pgfusepath{fill}%
\end{pgfscope}%
\begin{pgfscope}%
\pgfpathrectangle{\pgfqpoint{0.100000in}{0.100000in}}{\pgfqpoint{3.608454in}{2.310000in}}%
\pgfusepath{clip}%
\pgfsetbuttcap%
\pgfsetmiterjoin%
\definecolor{currentfill}{rgb}{0.000000,0.639216,0.680392}%
\pgfsetfillcolor{currentfill}%
\pgfsetlinewidth{0.000000pt}%
\definecolor{currentstroke}{rgb}{0.000000,0.000000,0.000000}%
\pgfsetstrokecolor{currentstroke}%
\pgfsetstrokeopacity{0.000000}%
\pgfsetdash{}{0pt}%
\pgfpathmoveto{\pgfqpoint{1.889333in}{1.734748in}}%
\pgfpathlineto{\pgfqpoint{1.888636in}{1.700186in}}%
\pgfpathlineto{\pgfqpoint{1.859723in}{1.701494in}}%
\pgfpathlineto{\pgfqpoint{1.841021in}{1.702423in}}%
\pgfpathlineto{\pgfqpoint{1.842006in}{1.721715in}}%
\pgfpathlineto{\pgfqpoint{1.829398in}{1.718027in}}%
\pgfpathlineto{\pgfqpoint{1.817373in}{1.718961in}}%
\pgfpathlineto{\pgfqpoint{1.807130in}{1.721260in}}%
\pgfpathlineto{\pgfqpoint{1.800957in}{1.721865in}}%
\pgfpathlineto{\pgfqpoint{1.800009in}{1.732253in}}%
\pgfpathlineto{\pgfqpoint{1.801456in}{1.757724in}}%
\pgfpathlineto{\pgfqpoint{1.801576in}{1.759791in}}%
\pgfpathlineto{\pgfqpoint{1.825951in}{1.758455in}}%
\pgfpathlineto{\pgfqpoint{1.830770in}{1.752552in}}%
\pgfpathlineto{\pgfqpoint{1.835653in}{1.753170in}}%
\pgfpathlineto{\pgfqpoint{1.843095in}{1.749824in}}%
\pgfpathlineto{\pgfqpoint{1.841646in}{1.759396in}}%
\pgfpathlineto{\pgfqpoint{1.846482in}{1.757163in}}%
\pgfpathlineto{\pgfqpoint{1.883458in}{1.755759in}}%
\pgfpathlineto{\pgfqpoint{1.882514in}{1.735032in}}%
\pgfpathlineto{\pgfqpoint{1.889333in}{1.734748in}}%
\pgfpathclose%
\pgfusepath{fill}%
\end{pgfscope}%
\begin{pgfscope}%
\pgfpathrectangle{\pgfqpoint{0.100000in}{0.100000in}}{\pgfqpoint{3.608454in}{2.310000in}}%
\pgfusepath{clip}%
\pgfsetbuttcap%
\pgfsetmiterjoin%
\definecolor{currentfill}{rgb}{0.000000,0.925490,0.537255}%
\pgfsetfillcolor{currentfill}%
\pgfsetlinewidth{0.000000pt}%
\definecolor{currentstroke}{rgb}{0.000000,0.000000,0.000000}%
\pgfsetstrokecolor{currentstroke}%
\pgfsetstrokeopacity{0.000000}%
\pgfsetdash{}{0pt}%
\pgfpathmoveto{\pgfqpoint{2.895571in}{1.311319in}}%
\pgfpathlineto{\pgfqpoint{2.902369in}{1.314765in}}%
\pgfpathlineto{\pgfqpoint{2.908038in}{1.306264in}}%
\pgfpathlineto{\pgfqpoint{2.909448in}{1.297631in}}%
\pgfpathlineto{\pgfqpoint{2.916184in}{1.292250in}}%
\pgfpathlineto{\pgfqpoint{2.929879in}{1.293625in}}%
\pgfpathlineto{\pgfqpoint{2.927867in}{1.284414in}}%
\pgfpathlineto{\pgfqpoint{2.923822in}{1.280887in}}%
\pgfpathlineto{\pgfqpoint{2.921023in}{1.282534in}}%
\pgfpathlineto{\pgfqpoint{2.902185in}{1.260526in}}%
\pgfpathlineto{\pgfqpoint{2.887645in}{1.250541in}}%
\pgfpathlineto{\pgfqpoint{2.880671in}{1.252069in}}%
\pgfpathlineto{\pgfqpoint{2.874724in}{1.257922in}}%
\pgfpathlineto{\pgfqpoint{2.872542in}{1.266952in}}%
\pgfpathlineto{\pgfqpoint{2.860524in}{1.271364in}}%
\pgfpathlineto{\pgfqpoint{2.849787in}{1.279563in}}%
\pgfpathlineto{\pgfqpoint{2.839446in}{1.282704in}}%
\pgfpathlineto{\pgfqpoint{2.837862in}{1.289352in}}%
\pgfpathlineto{\pgfqpoint{2.847809in}{1.300549in}}%
\pgfpathlineto{\pgfqpoint{2.855369in}{1.302588in}}%
\pgfpathlineto{\pgfqpoint{2.850698in}{1.310252in}}%
\pgfpathlineto{\pgfqpoint{2.857831in}{1.319548in}}%
\pgfpathlineto{\pgfqpoint{2.855274in}{1.324556in}}%
\pgfpathlineto{\pgfqpoint{2.861780in}{1.330306in}}%
\pgfpathlineto{\pgfqpoint{2.874015in}{1.332782in}}%
\pgfpathlineto{\pgfqpoint{2.873043in}{1.323790in}}%
\pgfpathlineto{\pgfqpoint{2.885378in}{1.313095in}}%
\pgfpathlineto{\pgfqpoint{2.884060in}{1.308388in}}%
\pgfpathlineto{\pgfqpoint{2.890443in}{1.304155in}}%
\pgfpathlineto{\pgfqpoint{2.895571in}{1.311319in}}%
\pgfpathclose%
\pgfusepath{fill}%
\end{pgfscope}%
\begin{pgfscope}%
\pgfpathrectangle{\pgfqpoint{0.100000in}{0.100000in}}{\pgfqpoint{3.608454in}{2.310000in}}%
\pgfusepath{clip}%
\pgfsetbuttcap%
\pgfsetmiterjoin%
\definecolor{currentfill}{rgb}{0.000000,0.462745,0.768627}%
\pgfsetfillcolor{currentfill}%
\pgfsetlinewidth{0.000000pt}%
\definecolor{currentstroke}{rgb}{0.000000,0.000000,0.000000}%
\pgfsetstrokecolor{currentstroke}%
\pgfsetstrokeopacity{0.000000}%
\pgfsetdash{}{0pt}%
\pgfpathmoveto{\pgfqpoint{2.193901in}{1.649520in}}%
\pgfpathlineto{\pgfqpoint{2.166463in}{1.649078in}}%
\pgfpathlineto{\pgfqpoint{2.165979in}{1.695985in}}%
\pgfpathlineto{\pgfqpoint{2.184501in}{1.696223in}}%
\pgfpathlineto{\pgfqpoint{2.220391in}{1.696881in}}%
\pgfpathlineto{\pgfqpoint{2.221382in}{1.649990in}}%
\pgfpathlineto{\pgfqpoint{2.193901in}{1.649520in}}%
\pgfpathclose%
\pgfusepath{fill}%
\end{pgfscope}%
\begin{pgfscope}%
\pgfpathrectangle{\pgfqpoint{0.100000in}{0.100000in}}{\pgfqpoint{3.608454in}{2.310000in}}%
\pgfusepath{clip}%
\pgfsetbuttcap%
\pgfsetmiterjoin%
\definecolor{currentfill}{rgb}{0.000000,0.768627,0.615686}%
\pgfsetfillcolor{currentfill}%
\pgfsetlinewidth{0.000000pt}%
\definecolor{currentstroke}{rgb}{0.000000,0.000000,0.000000}%
\pgfsetstrokecolor{currentstroke}%
\pgfsetstrokeopacity{0.000000}%
\pgfsetdash{}{0pt}%
\pgfpathmoveto{\pgfqpoint{2.936140in}{1.370645in}}%
\pgfpathlineto{\pgfqpoint{2.930069in}{1.365085in}}%
\pgfpathlineto{\pgfqpoint{2.920044in}{1.373705in}}%
\pgfpathlineto{\pgfqpoint{2.910035in}{1.386821in}}%
\pgfpathlineto{\pgfqpoint{2.916733in}{1.399937in}}%
\pgfpathlineto{\pgfqpoint{2.915344in}{1.412906in}}%
\pgfpathlineto{\pgfqpoint{2.916394in}{1.420191in}}%
\pgfpathlineto{\pgfqpoint{2.907865in}{1.422811in}}%
\pgfpathlineto{\pgfqpoint{2.906878in}{1.433220in}}%
\pgfpathlineto{\pgfqpoint{2.907685in}{1.436860in}}%
\pgfpathlineto{\pgfqpoint{2.914756in}{1.436835in}}%
\pgfpathlineto{\pgfqpoint{2.914676in}{1.445590in}}%
\pgfpathlineto{\pgfqpoint{2.927738in}{1.446958in}}%
\pgfpathlineto{\pgfqpoint{2.938939in}{1.450447in}}%
\pgfpathlineto{\pgfqpoint{2.941441in}{1.448688in}}%
\pgfpathlineto{\pgfqpoint{2.954171in}{1.450091in}}%
\pgfpathlineto{\pgfqpoint{2.951761in}{1.443677in}}%
\pgfpathlineto{\pgfqpoint{2.962415in}{1.436530in}}%
\pgfpathlineto{\pgfqpoint{2.966480in}{1.429783in}}%
\pgfpathlineto{\pgfqpoint{2.965224in}{1.427879in}}%
\pgfpathlineto{\pgfqpoint{2.973331in}{1.416066in}}%
\pgfpathlineto{\pgfqpoint{2.967704in}{1.409901in}}%
\pgfpathlineto{\pgfqpoint{2.961168in}{1.406181in}}%
\pgfpathlineto{\pgfqpoint{2.954499in}{1.407993in}}%
\pgfpathlineto{\pgfqpoint{2.943680in}{1.394391in}}%
\pgfpathlineto{\pgfqpoint{2.936799in}{1.397031in}}%
\pgfpathlineto{\pgfqpoint{2.932362in}{1.392633in}}%
\pgfpathlineto{\pgfqpoint{2.933951in}{1.375085in}}%
\pgfpathlineto{\pgfqpoint{2.936140in}{1.370645in}}%
\pgfpathclose%
\pgfusepath{fill}%
\end{pgfscope}%
\begin{pgfscope}%
\pgfpathrectangle{\pgfqpoint{0.100000in}{0.100000in}}{\pgfqpoint{3.608454in}{2.310000in}}%
\pgfusepath{clip}%
\pgfsetbuttcap%
\pgfsetmiterjoin%
\definecolor{currentfill}{rgb}{0.000000,0.431373,0.784314}%
\pgfsetfillcolor{currentfill}%
\pgfsetlinewidth{0.000000pt}%
\definecolor{currentstroke}{rgb}{0.000000,0.000000,0.000000}%
\pgfsetstrokecolor{currentstroke}%
\pgfsetstrokeopacity{0.000000}%
\pgfsetdash{}{0pt}%
\pgfpathmoveto{\pgfqpoint{2.258667in}{1.780519in}}%
\pgfpathlineto{\pgfqpoint{2.268903in}{1.780718in}}%
\pgfpathlineto{\pgfqpoint{2.267401in}{1.826914in}}%
\pgfpathlineto{\pgfqpoint{2.266038in}{1.833884in}}%
\pgfpathlineto{\pgfqpoint{2.300433in}{1.834842in}}%
\pgfpathlineto{\pgfqpoint{2.300225in}{1.841636in}}%
\pgfpathlineto{\pgfqpoint{2.334657in}{1.843071in}}%
\pgfpathlineto{\pgfqpoint{2.337567in}{1.780978in}}%
\pgfpathlineto{\pgfqpoint{2.344426in}{1.781264in}}%
\pgfpathlineto{\pgfqpoint{2.344731in}{1.774369in}}%
\pgfpathlineto{\pgfqpoint{2.372249in}{1.775882in}}%
\pgfpathlineto{\pgfqpoint{2.373196in}{1.761948in}}%
\pgfpathlineto{\pgfqpoint{2.373549in}{1.754537in}}%
\pgfpathlineto{\pgfqpoint{2.339782in}{1.753146in}}%
\pgfpathlineto{\pgfqpoint{2.336281in}{1.746036in}}%
\pgfpathlineto{\pgfqpoint{2.317255in}{1.744761in}}%
\pgfpathlineto{\pgfqpoint{2.316678in}{1.738350in}}%
\pgfpathlineto{\pgfqpoint{2.310801in}{1.738104in}}%
\pgfpathlineto{\pgfqpoint{2.319291in}{1.727547in}}%
\pgfpathlineto{\pgfqpoint{2.293781in}{1.726557in}}%
\pgfpathlineto{\pgfqpoint{2.273836in}{1.725888in}}%
\pgfpathlineto{\pgfqpoint{2.272950in}{1.753323in}}%
\pgfpathlineto{\pgfqpoint{2.285525in}{1.753739in}}%
\pgfpathlineto{\pgfqpoint{2.281318in}{1.764757in}}%
\pgfpathlineto{\pgfqpoint{2.273713in}{1.770000in}}%
\pgfpathlineto{\pgfqpoint{2.264630in}{1.772555in}}%
\pgfpathlineto{\pgfqpoint{2.258667in}{1.780519in}}%
\pgfpathclose%
\pgfusepath{fill}%
\end{pgfscope}%
\begin{pgfscope}%
\pgfpathrectangle{\pgfqpoint{0.100000in}{0.100000in}}{\pgfqpoint{3.608454in}{2.310000in}}%
\pgfusepath{clip}%
\pgfsetbuttcap%
\pgfsetmiterjoin%
\definecolor{currentfill}{rgb}{0.000000,0.678431,0.660784}%
\pgfsetfillcolor{currentfill}%
\pgfsetlinewidth{0.000000pt}%
\definecolor{currentstroke}{rgb}{0.000000,0.000000,0.000000}%
\pgfsetstrokecolor{currentstroke}%
\pgfsetstrokeopacity{0.000000}%
\pgfsetdash{}{0pt}%
\pgfpathmoveto{\pgfqpoint{1.938836in}{0.498314in}}%
\pgfpathlineto{\pgfqpoint{1.924683in}{0.477880in}}%
\pgfpathlineto{\pgfqpoint{1.907060in}{0.482191in}}%
\pgfpathlineto{\pgfqpoint{1.907719in}{0.482706in}}%
\pgfpathlineto{\pgfqpoint{1.889121in}{0.521871in}}%
\pgfpathlineto{\pgfqpoint{1.895145in}{0.523860in}}%
\pgfpathlineto{\pgfqpoint{1.882610in}{0.539717in}}%
\pgfpathlineto{\pgfqpoint{1.915913in}{0.565873in}}%
\pgfpathlineto{\pgfqpoint{1.923720in}{0.556660in}}%
\pgfpathlineto{\pgfqpoint{1.913527in}{0.548723in}}%
\pgfpathlineto{\pgfqpoint{1.925742in}{0.532805in}}%
\pgfpathlineto{\pgfqpoint{1.911040in}{0.521565in}}%
\pgfpathlineto{\pgfqpoint{1.916523in}{0.511256in}}%
\pgfpathlineto{\pgfqpoint{1.926187in}{0.508771in}}%
\pgfpathlineto{\pgfqpoint{1.926522in}{0.503678in}}%
\pgfpathlineto{\pgfqpoint{1.938836in}{0.498314in}}%
\pgfpathclose%
\pgfusepath{fill}%
\end{pgfscope}%
\begin{pgfscope}%
\pgfpathrectangle{\pgfqpoint{0.100000in}{0.100000in}}{\pgfqpoint{3.608454in}{2.310000in}}%
\pgfusepath{clip}%
\pgfsetbuttcap%
\pgfsetmiterjoin%
\definecolor{currentfill}{rgb}{0.000000,0.596078,0.701961}%
\pgfsetfillcolor{currentfill}%
\pgfsetlinewidth{0.000000pt}%
\definecolor{currentstroke}{rgb}{0.000000,0.000000,0.000000}%
\pgfsetstrokecolor{currentstroke}%
\pgfsetstrokeopacity{0.000000}%
\pgfsetdash{}{0pt}%
\pgfpathmoveto{\pgfqpoint{2.999992in}{0.934399in}}%
\pgfpathlineto{\pgfqpoint{2.986683in}{0.921479in}}%
\pgfpathlineto{\pgfqpoint{2.980224in}{0.924118in}}%
\pgfpathlineto{\pgfqpoint{2.965871in}{0.919067in}}%
\pgfpathlineto{\pgfqpoint{2.962570in}{0.910798in}}%
\pgfpathlineto{\pgfqpoint{2.959107in}{0.909040in}}%
\pgfpathlineto{\pgfqpoint{2.950977in}{0.910809in}}%
\pgfpathlineto{\pgfqpoint{2.944266in}{0.903605in}}%
\pgfpathlineto{\pgfqpoint{2.937410in}{0.908377in}}%
\pgfpathlineto{\pgfqpoint{2.937493in}{0.915146in}}%
\pgfpathlineto{\pgfqpoint{2.933286in}{0.922296in}}%
\pgfpathlineto{\pgfqpoint{2.925307in}{0.930408in}}%
\pgfpathlineto{\pgfqpoint{2.919634in}{0.933284in}}%
\pgfpathlineto{\pgfqpoint{2.918368in}{0.938239in}}%
\pgfpathlineto{\pgfqpoint{2.909363in}{0.953452in}}%
\pgfpathlineto{\pgfqpoint{2.910547in}{0.958327in}}%
\pgfpathlineto{\pgfqpoint{2.915431in}{0.959683in}}%
\pgfpathlineto{\pgfqpoint{2.918991in}{0.967320in}}%
\pgfpathlineto{\pgfqpoint{2.922909in}{0.970556in}}%
\pgfpathlineto{\pgfqpoint{2.932420in}{0.972207in}}%
\pgfpathlineto{\pgfqpoint{2.942345in}{0.978613in}}%
\pgfpathlineto{\pgfqpoint{2.948243in}{0.978295in}}%
\pgfpathlineto{\pgfqpoint{2.954883in}{0.971559in}}%
\pgfpathlineto{\pgfqpoint{2.956878in}{0.975335in}}%
\pgfpathlineto{\pgfqpoint{2.950161in}{0.981612in}}%
\pgfpathlineto{\pgfqpoint{2.953478in}{0.990526in}}%
\pgfpathlineto{\pgfqpoint{2.949644in}{0.997974in}}%
\pgfpathlineto{\pgfqpoint{2.958909in}{1.001900in}}%
\pgfpathlineto{\pgfqpoint{2.966134in}{1.004041in}}%
\pgfpathlineto{\pgfqpoint{2.971145in}{0.996030in}}%
\pgfpathlineto{\pgfqpoint{2.983807in}{0.993720in}}%
\pgfpathlineto{\pgfqpoint{2.988314in}{0.999433in}}%
\pgfpathlineto{\pgfqpoint{2.999338in}{0.990301in}}%
\pgfpathlineto{\pgfqpoint{3.016045in}{0.985596in}}%
\pgfpathlineto{\pgfqpoint{3.005840in}{0.971004in}}%
\pgfpathlineto{\pgfqpoint{2.984051in}{0.944289in}}%
\pgfpathlineto{\pgfqpoint{2.988444in}{0.938374in}}%
\pgfpathlineto{\pgfqpoint{2.999992in}{0.934399in}}%
\pgfpathclose%
\pgfusepath{fill}%
\end{pgfscope}%
\begin{pgfscope}%
\pgfpathrectangle{\pgfqpoint{0.100000in}{0.100000in}}{\pgfqpoint{3.608454in}{2.310000in}}%
\pgfusepath{clip}%
\pgfsetbuttcap%
\pgfsetmiterjoin%
\definecolor{currentfill}{rgb}{0.000000,0.466667,0.766667}%
\pgfsetfillcolor{currentfill}%
\pgfsetlinewidth{0.000000pt}%
\definecolor{currentstroke}{rgb}{0.000000,0.000000,0.000000}%
\pgfsetstrokecolor{currentstroke}%
\pgfsetstrokeopacity{0.000000}%
\pgfsetdash{}{0pt}%
\pgfpathmoveto{\pgfqpoint{3.243003in}{1.088398in}}%
\pgfpathlineto{\pgfqpoint{3.237815in}{1.092298in}}%
\pgfpathlineto{\pgfqpoint{3.228319in}{1.108263in}}%
\pgfpathlineto{\pgfqpoint{3.227150in}{1.124441in}}%
\pgfpathlineto{\pgfqpoint{3.219619in}{1.130949in}}%
\pgfpathlineto{\pgfqpoint{3.216293in}{1.140248in}}%
\pgfpathlineto{\pgfqpoint{3.211683in}{1.141958in}}%
\pgfpathlineto{\pgfqpoint{3.206484in}{1.172807in}}%
\pgfpathlineto{\pgfqpoint{3.229923in}{1.198333in}}%
\pgfpathlineto{\pgfqpoint{3.232624in}{1.197760in}}%
\pgfpathlineto{\pgfqpoint{3.244554in}{1.193790in}}%
\pgfpathlineto{\pgfqpoint{3.246587in}{1.185703in}}%
\pgfpathlineto{\pgfqpoint{3.250804in}{1.179929in}}%
\pgfpathlineto{\pgfqpoint{3.248243in}{1.176631in}}%
\pgfpathlineto{\pgfqpoint{3.248882in}{1.168666in}}%
\pgfpathlineto{\pgfqpoint{3.265343in}{1.162433in}}%
\pgfpathlineto{\pgfqpoint{3.273900in}{1.157290in}}%
\pgfpathlineto{\pgfqpoint{3.287098in}{1.161529in}}%
\pgfpathlineto{\pgfqpoint{3.286885in}{1.169498in}}%
\pgfpathlineto{\pgfqpoint{3.295207in}{1.169584in}}%
\pgfpathlineto{\pgfqpoint{3.297149in}{1.166013in}}%
\pgfpathlineto{\pgfqpoint{3.293259in}{1.151590in}}%
\pgfpathlineto{\pgfqpoint{3.288603in}{1.146209in}}%
\pgfpathlineto{\pgfqpoint{3.292278in}{1.142409in}}%
\pgfpathlineto{\pgfqpoint{3.303080in}{1.146473in}}%
\pgfpathlineto{\pgfqpoint{3.308404in}{1.145005in}}%
\pgfpathlineto{\pgfqpoint{3.312700in}{1.148475in}}%
\pgfpathlineto{\pgfqpoint{3.320379in}{1.148245in}}%
\pgfpathlineto{\pgfqpoint{3.305824in}{1.120574in}}%
\pgfpathlineto{\pgfqpoint{3.300774in}{1.117204in}}%
\pgfpathlineto{\pgfqpoint{3.293113in}{1.118696in}}%
\pgfpathlineto{\pgfqpoint{3.285161in}{1.117485in}}%
\pgfpathlineto{\pgfqpoint{3.263408in}{1.107685in}}%
\pgfpathlineto{\pgfqpoint{3.243003in}{1.088398in}}%
\pgfpathclose%
\pgfusepath{fill}%
\end{pgfscope}%
\begin{pgfscope}%
\pgfpathrectangle{\pgfqpoint{0.100000in}{0.100000in}}{\pgfqpoint{3.608454in}{2.310000in}}%
\pgfusepath{clip}%
\pgfsetbuttcap%
\pgfsetmiterjoin%
\definecolor{currentfill}{rgb}{0.000000,0.509804,0.745098}%
\pgfsetfillcolor{currentfill}%
\pgfsetlinewidth{0.000000pt}%
\definecolor{currentstroke}{rgb}{0.000000,0.000000,0.000000}%
\pgfsetstrokecolor{currentstroke}%
\pgfsetstrokeopacity{0.000000}%
\pgfsetdash{}{0pt}%
\pgfpathmoveto{\pgfqpoint{1.757470in}{0.722100in}}%
\pgfpathlineto{\pgfqpoint{1.755704in}{0.692007in}}%
\pgfpathlineto{\pgfqpoint{1.697817in}{0.695195in}}%
\pgfpathlineto{\pgfqpoint{1.699678in}{0.725186in}}%
\pgfpathlineto{\pgfqpoint{1.678348in}{0.726376in}}%
\pgfpathlineto{\pgfqpoint{1.681284in}{0.762092in}}%
\pgfpathlineto{\pgfqpoint{1.684601in}{0.806456in}}%
\pgfpathlineto{\pgfqpoint{1.690000in}{0.806092in}}%
\pgfpathlineto{\pgfqpoint{1.724900in}{0.803670in}}%
\pgfpathlineto{\pgfqpoint{1.753852in}{0.801689in}}%
\pgfpathlineto{\pgfqpoint{1.751590in}{0.761847in}}%
\pgfpathlineto{\pgfqpoint{1.759948in}{0.761279in}}%
\pgfpathlineto{\pgfqpoint{1.757470in}{0.722100in}}%
\pgfpathclose%
\pgfusepath{fill}%
\end{pgfscope}%
\begin{pgfscope}%
\pgfpathrectangle{\pgfqpoint{0.100000in}{0.100000in}}{\pgfqpoint{3.608454in}{2.310000in}}%
\pgfusepath{clip}%
\pgfsetbuttcap%
\pgfsetmiterjoin%
\definecolor{currentfill}{rgb}{0.000000,0.505882,0.747059}%
\pgfsetfillcolor{currentfill}%
\pgfsetlinewidth{0.000000pt}%
\definecolor{currentstroke}{rgb}{0.000000,0.000000,0.000000}%
\pgfsetstrokecolor{currentstroke}%
\pgfsetstrokeopacity{0.000000}%
\pgfsetdash{}{0pt}%
\pgfpathmoveto{\pgfqpoint{2.839244in}{0.999162in}}%
\pgfpathlineto{\pgfqpoint{2.836248in}{0.996454in}}%
\pgfpathlineto{\pgfqpoint{2.827384in}{1.003120in}}%
\pgfpathlineto{\pgfqpoint{2.830646in}{1.010243in}}%
\pgfpathlineto{\pgfqpoint{2.825899in}{1.020792in}}%
\pgfpathlineto{\pgfqpoint{2.818142in}{1.023149in}}%
\pgfpathlineto{\pgfqpoint{2.810860in}{1.028741in}}%
\pgfpathlineto{\pgfqpoint{2.810594in}{1.033632in}}%
\pgfpathlineto{\pgfqpoint{2.812171in}{1.037511in}}%
\pgfpathlineto{\pgfqpoint{2.820073in}{1.037913in}}%
\pgfpathlineto{\pgfqpoint{2.825426in}{1.046467in}}%
\pgfpathlineto{\pgfqpoint{2.830948in}{1.045651in}}%
\pgfpathlineto{\pgfqpoint{2.833248in}{1.051039in}}%
\pgfpathlineto{\pgfqpoint{2.843160in}{1.055221in}}%
\pgfpathlineto{\pgfqpoint{2.847354in}{1.055558in}}%
\pgfpathlineto{\pgfqpoint{2.850394in}{1.050302in}}%
\pgfpathlineto{\pgfqpoint{2.863588in}{1.049271in}}%
\pgfpathlineto{\pgfqpoint{2.864794in}{1.047108in}}%
\pgfpathlineto{\pgfqpoint{2.863172in}{1.045964in}}%
\pgfpathlineto{\pgfqpoint{2.859003in}{1.029918in}}%
\pgfpathlineto{\pgfqpoint{2.864516in}{1.026478in}}%
\pgfpathlineto{\pgfqpoint{2.864861in}{1.018162in}}%
\pgfpathlineto{\pgfqpoint{2.868390in}{1.010555in}}%
\pgfpathlineto{\pgfqpoint{2.865578in}{1.008074in}}%
\pgfpathlineto{\pgfqpoint{2.869364in}{1.002057in}}%
\pgfpathlineto{\pgfqpoint{2.868081in}{0.996615in}}%
\pgfpathlineto{\pgfqpoint{2.859121in}{0.988728in}}%
\pgfpathlineto{\pgfqpoint{2.856780in}{0.993705in}}%
\pgfpathlineto{\pgfqpoint{2.843269in}{0.994686in}}%
\pgfpathlineto{\pgfqpoint{2.839244in}{0.999162in}}%
\pgfpathclose%
\pgfusepath{fill}%
\end{pgfscope}%
\begin{pgfscope}%
\pgfpathrectangle{\pgfqpoint{0.100000in}{0.100000in}}{\pgfqpoint{3.608454in}{2.310000in}}%
\pgfusepath{clip}%
\pgfsetbuttcap%
\pgfsetmiterjoin%
\definecolor{currentfill}{rgb}{0.000000,0.545098,0.727451}%
\pgfsetfillcolor{currentfill}%
\pgfsetlinewidth{0.000000pt}%
\definecolor{currentstroke}{rgb}{0.000000,0.000000,0.000000}%
\pgfsetstrokecolor{currentstroke}%
\pgfsetstrokeopacity{0.000000}%
\pgfsetdash{}{0pt}%
\pgfpathmoveto{\pgfqpoint{2.075326in}{1.295024in}}%
\pgfpathlineto{\pgfqpoint{2.047953in}{1.295247in}}%
\pgfpathlineto{\pgfqpoint{2.048478in}{1.319317in}}%
\pgfpathlineto{\pgfqpoint{2.023407in}{1.319659in}}%
\pgfpathlineto{\pgfqpoint{2.021173in}{1.326536in}}%
\pgfpathlineto{\pgfqpoint{2.014334in}{1.326679in}}%
\pgfpathlineto{\pgfqpoint{2.014646in}{1.343860in}}%
\pgfpathlineto{\pgfqpoint{2.021423in}{1.343723in}}%
\pgfpathlineto{\pgfqpoint{2.021694in}{1.353951in}}%
\pgfpathlineto{\pgfqpoint{2.030596in}{1.356610in}}%
\pgfpathlineto{\pgfqpoint{2.034708in}{1.355023in}}%
\pgfpathlineto{\pgfqpoint{2.043381in}{1.357123in}}%
\pgfpathlineto{\pgfqpoint{2.043672in}{1.384754in}}%
\pgfpathlineto{\pgfqpoint{2.058697in}{1.384598in}}%
\pgfpathlineto{\pgfqpoint{2.058778in}{1.391470in}}%
\pgfpathlineto{\pgfqpoint{2.072459in}{1.391367in}}%
\pgfpathlineto{\pgfqpoint{2.071958in}{1.372885in}}%
\pgfpathlineto{\pgfqpoint{2.095818in}{1.372775in}}%
\pgfpathlineto{\pgfqpoint{2.095756in}{1.349849in}}%
\pgfpathlineto{\pgfqpoint{2.096838in}{1.338167in}}%
\pgfpathlineto{\pgfqpoint{2.103336in}{1.338237in}}%
\pgfpathlineto{\pgfqpoint{2.103312in}{1.318980in}}%
\pgfpathlineto{\pgfqpoint{2.075903in}{1.319066in}}%
\pgfpathlineto{\pgfqpoint{2.075326in}{1.295024in}}%
\pgfpathclose%
\pgfusepath{fill}%
\end{pgfscope}%
\begin{pgfscope}%
\pgfpathrectangle{\pgfqpoint{0.100000in}{0.100000in}}{\pgfqpoint{3.608454in}{2.310000in}}%
\pgfusepath{clip}%
\pgfsetbuttcap%
\pgfsetmiterjoin%
\definecolor{currentfill}{rgb}{0.000000,0.478431,0.760784}%
\pgfsetfillcolor{currentfill}%
\pgfsetlinewidth{0.000000pt}%
\definecolor{currentstroke}{rgb}{0.000000,0.000000,0.000000}%
\pgfsetstrokecolor{currentstroke}%
\pgfsetstrokeopacity{0.000000}%
\pgfsetdash{}{0pt}%
\pgfpathmoveto{\pgfqpoint{2.024967in}{0.596552in}}%
\pgfpathlineto{\pgfqpoint{2.015636in}{0.587127in}}%
\pgfpathlineto{\pgfqpoint{2.012654in}{0.578657in}}%
\pgfpathlineto{\pgfqpoint{1.992065in}{0.566132in}}%
\pgfpathlineto{\pgfqpoint{1.997763in}{0.573057in}}%
\pgfpathlineto{\pgfqpoint{1.981908in}{0.592727in}}%
\pgfpathlineto{\pgfqpoint{1.976342in}{0.597283in}}%
\pgfpathlineto{\pgfqpoint{1.957611in}{0.597327in}}%
\pgfpathlineto{\pgfqpoint{1.945778in}{0.610158in}}%
\pgfpathlineto{\pgfqpoint{1.966633in}{0.631014in}}%
\pgfpathlineto{\pgfqpoint{1.979102in}{0.635917in}}%
\pgfpathlineto{\pgfqpoint{1.982721in}{0.639405in}}%
\pgfpathlineto{\pgfqpoint{1.990391in}{0.638498in}}%
\pgfpathlineto{\pgfqpoint{2.000273in}{0.618823in}}%
\pgfpathlineto{\pgfqpoint{2.008648in}{0.612091in}}%
\pgfpathlineto{\pgfqpoint{2.014382in}{0.611205in}}%
\pgfpathlineto{\pgfqpoint{2.019290in}{0.599307in}}%
\pgfpathlineto{\pgfqpoint{2.024967in}{0.596552in}}%
\pgfpathclose%
\pgfusepath{fill}%
\end{pgfscope}%
\begin{pgfscope}%
\pgfpathrectangle{\pgfqpoint{0.100000in}{0.100000in}}{\pgfqpoint{3.608454in}{2.310000in}}%
\pgfusepath{clip}%
\pgfsetbuttcap%
\pgfsetmiterjoin%
\definecolor{currentfill}{rgb}{0.000000,0.647059,0.676471}%
\pgfsetfillcolor{currentfill}%
\pgfsetlinewidth{0.000000pt}%
\definecolor{currentstroke}{rgb}{0.000000,0.000000,0.000000}%
\pgfsetstrokecolor{currentstroke}%
\pgfsetstrokeopacity{0.000000}%
\pgfsetdash{}{0pt}%
\pgfpathmoveto{\pgfqpoint{0.755783in}{1.587192in}}%
\pgfpathlineto{\pgfqpoint{0.693366in}{1.603165in}}%
\pgfpathlineto{\pgfqpoint{0.650489in}{1.614771in}}%
\pgfpathlineto{\pgfqpoint{0.661398in}{1.654341in}}%
\pgfpathlineto{\pgfqpoint{0.662884in}{1.653976in}}%
\pgfpathlineto{\pgfqpoint{0.672323in}{1.687633in}}%
\pgfpathlineto{\pgfqpoint{0.671164in}{1.687956in}}%
\pgfpathlineto{\pgfqpoint{0.681032in}{1.722500in}}%
\pgfpathlineto{\pgfqpoint{0.693113in}{1.766738in}}%
\pgfpathlineto{\pgfqpoint{0.728707in}{1.756911in}}%
\pgfpathlineto{\pgfqpoint{0.757203in}{1.749818in}}%
\pgfpathlineto{\pgfqpoint{0.824016in}{1.733109in}}%
\pgfpathlineto{\pgfqpoint{0.824452in}{1.732928in}}%
\pgfpathlineto{\pgfqpoint{0.805480in}{1.656197in}}%
\pgfpathlineto{\pgfqpoint{0.798688in}{1.628778in}}%
\pgfpathlineto{\pgfqpoint{0.785394in}{1.632033in}}%
\pgfpathlineto{\pgfqpoint{0.755783in}{1.587192in}}%
\pgfpathclose%
\pgfusepath{fill}%
\end{pgfscope}%
\begin{pgfscope}%
\pgfpathrectangle{\pgfqpoint{0.100000in}{0.100000in}}{\pgfqpoint{3.608454in}{2.310000in}}%
\pgfusepath{clip}%
\pgfsetbuttcap%
\pgfsetmiterjoin%
\definecolor{currentfill}{rgb}{0.000000,0.329412,0.835294}%
\pgfsetfillcolor{currentfill}%
\pgfsetlinewidth{0.000000pt}%
\definecolor{currentstroke}{rgb}{0.000000,0.000000,0.000000}%
\pgfsetstrokecolor{currentstroke}%
\pgfsetstrokeopacity{0.000000}%
\pgfsetdash{}{0pt}%
\pgfpathmoveto{\pgfqpoint{1.989587in}{1.638444in}}%
\pgfpathlineto{\pgfqpoint{1.988757in}{1.605977in}}%
\pgfpathlineto{\pgfqpoint{1.968280in}{1.606524in}}%
\pgfpathlineto{\pgfqpoint{1.968471in}{1.613417in}}%
\pgfpathlineto{\pgfqpoint{1.914002in}{1.614954in}}%
\pgfpathlineto{\pgfqpoint{1.914725in}{1.650247in}}%
\pgfpathlineto{\pgfqpoint{1.922640in}{1.646360in}}%
\pgfpathlineto{\pgfqpoint{1.925767in}{1.649254in}}%
\pgfpathlineto{\pgfqpoint{1.927556in}{1.657034in}}%
\pgfpathlineto{\pgfqpoint{1.926808in}{1.664949in}}%
\pgfpathlineto{\pgfqpoint{1.928813in}{1.672403in}}%
\pgfpathlineto{\pgfqpoint{1.968146in}{1.671305in}}%
\pgfpathlineto{\pgfqpoint{1.981953in}{1.671006in}}%
\pgfpathlineto{\pgfqpoint{1.981253in}{1.641675in}}%
\pgfpathlineto{\pgfqpoint{1.989587in}{1.638444in}}%
\pgfpathclose%
\pgfusepath{fill}%
\end{pgfscope}%
\begin{pgfscope}%
\pgfpathrectangle{\pgfqpoint{0.100000in}{0.100000in}}{\pgfqpoint{3.608454in}{2.310000in}}%
\pgfusepath{clip}%
\pgfsetbuttcap%
\pgfsetmiterjoin%
\definecolor{currentfill}{rgb}{0.000000,0.545098,0.727451}%
\pgfsetfillcolor{currentfill}%
\pgfsetlinewidth{0.000000pt}%
\definecolor{currentstroke}{rgb}{0.000000,0.000000,0.000000}%
\pgfsetstrokecolor{currentstroke}%
\pgfsetstrokeopacity{0.000000}%
\pgfsetdash{}{0pt}%
\pgfpathmoveto{\pgfqpoint{3.191049in}{1.170900in}}%
\pgfpathlineto{\pgfqpoint{3.206484in}{1.172807in}}%
\pgfpathlineto{\pgfqpoint{3.211683in}{1.141958in}}%
\pgfpathlineto{\pgfqpoint{3.216293in}{1.140248in}}%
\pgfpathlineto{\pgfqpoint{3.219619in}{1.130949in}}%
\pgfpathlineto{\pgfqpoint{3.227150in}{1.124441in}}%
\pgfpathlineto{\pgfqpoint{3.228319in}{1.108263in}}%
\pgfpathlineto{\pgfqpoint{3.200611in}{1.103096in}}%
\pgfpathlineto{\pgfqpoint{3.194034in}{1.089009in}}%
\pgfpathlineto{\pgfqpoint{3.182406in}{1.101945in}}%
\pgfpathlineto{\pgfqpoint{3.174249in}{1.109020in}}%
\pgfpathlineto{\pgfqpoint{3.162834in}{1.117655in}}%
\pgfpathlineto{\pgfqpoint{3.160679in}{1.126465in}}%
\pgfpathlineto{\pgfqpoint{3.160885in}{1.137946in}}%
\pgfpathlineto{\pgfqpoint{3.164716in}{1.144245in}}%
\pgfpathlineto{\pgfqpoint{3.168634in}{1.140787in}}%
\pgfpathlineto{\pgfqpoint{3.173892in}{1.141190in}}%
\pgfpathlineto{\pgfqpoint{3.180850in}{1.147206in}}%
\pgfpathlineto{\pgfqpoint{3.188860in}{1.151551in}}%
\pgfpathlineto{\pgfqpoint{3.187783in}{1.156766in}}%
\pgfpathlineto{\pgfqpoint{3.191049in}{1.170900in}}%
\pgfpathclose%
\pgfusepath{fill}%
\end{pgfscope}%
\begin{pgfscope}%
\pgfpathrectangle{\pgfqpoint{0.100000in}{0.100000in}}{\pgfqpoint{3.608454in}{2.310000in}}%
\pgfusepath{clip}%
\pgfsetbuttcap%
\pgfsetmiterjoin%
\definecolor{currentfill}{rgb}{0.000000,0.576471,0.711765}%
\pgfsetfillcolor{currentfill}%
\pgfsetlinewidth{0.000000pt}%
\definecolor{currentstroke}{rgb}{0.000000,0.000000,0.000000}%
\pgfsetstrokecolor{currentstroke}%
\pgfsetstrokeopacity{0.000000}%
\pgfsetdash{}{0pt}%
\pgfpathmoveto{\pgfqpoint{2.786852in}{0.748618in}}%
\pgfpathlineto{\pgfqpoint{2.789288in}{0.742878in}}%
\pgfpathlineto{\pgfqpoint{2.719091in}{0.734938in}}%
\pgfpathlineto{\pgfqpoint{2.708723in}{0.733977in}}%
\pgfpathlineto{\pgfqpoint{2.704075in}{0.776471in}}%
\pgfpathlineto{\pgfqpoint{2.706700in}{0.783758in}}%
\pgfpathlineto{\pgfqpoint{2.733549in}{0.786465in}}%
\pgfpathlineto{\pgfqpoint{2.755934in}{0.788907in}}%
\pgfpathlineto{\pgfqpoint{2.755229in}{0.795820in}}%
\pgfpathlineto{\pgfqpoint{2.774910in}{0.796997in}}%
\pgfpathlineto{\pgfqpoint{2.780161in}{0.791452in}}%
\pgfpathlineto{\pgfqpoint{2.781911in}{0.785540in}}%
\pgfpathlineto{\pgfqpoint{2.779513in}{0.764362in}}%
\pgfpathlineto{\pgfqpoint{2.781206in}{0.755220in}}%
\pgfpathlineto{\pgfqpoint{2.786852in}{0.748618in}}%
\pgfpathclose%
\pgfusepath{fill}%
\end{pgfscope}%
\begin{pgfscope}%
\pgfpathrectangle{\pgfqpoint{0.100000in}{0.100000in}}{\pgfqpoint{3.608454in}{2.310000in}}%
\pgfusepath{clip}%
\pgfsetbuttcap%
\pgfsetmiterjoin%
\definecolor{currentfill}{rgb}{0.000000,0.654902,0.672549}%
\pgfsetfillcolor{currentfill}%
\pgfsetlinewidth{0.000000pt}%
\definecolor{currentstroke}{rgb}{0.000000,0.000000,0.000000}%
\pgfsetstrokecolor{currentstroke}%
\pgfsetstrokeopacity{0.000000}%
\pgfsetdash{}{0pt}%
\pgfpathmoveto{\pgfqpoint{0.875234in}{2.151233in}}%
\pgfpathlineto{\pgfqpoint{0.821826in}{2.165426in}}%
\pgfpathlineto{\pgfqpoint{0.836363in}{2.217995in}}%
\pgfpathlineto{\pgfqpoint{0.843741in}{2.216808in}}%
\pgfpathlineto{\pgfqpoint{0.857008in}{2.211536in}}%
\pgfpathlineto{\pgfqpoint{0.862221in}{2.202501in}}%
\pgfpathlineto{\pgfqpoint{0.873890in}{2.207982in}}%
\pgfpathlineto{\pgfqpoint{0.878250in}{2.201957in}}%
\pgfpathlineto{\pgfqpoint{0.879920in}{2.194564in}}%
\pgfpathlineto{\pgfqpoint{0.891121in}{2.195618in}}%
\pgfpathlineto{\pgfqpoint{0.894472in}{2.191910in}}%
\pgfpathlineto{\pgfqpoint{0.902377in}{2.193011in}}%
\pgfpathlineto{\pgfqpoint{0.908021in}{2.186968in}}%
\pgfpathlineto{\pgfqpoint{0.913019in}{2.206570in}}%
\pgfpathlineto{\pgfqpoint{0.916424in}{2.212894in}}%
\pgfpathlineto{\pgfqpoint{0.926800in}{2.253370in}}%
\pgfpathlineto{\pgfqpoint{0.920139in}{2.255027in}}%
\pgfpathlineto{\pgfqpoint{0.925222in}{2.260906in}}%
\pgfpathlineto{\pgfqpoint{0.927075in}{2.268237in}}%
\pgfpathlineto{\pgfqpoint{0.937150in}{2.279229in}}%
\pgfpathlineto{\pgfqpoint{0.957430in}{2.274082in}}%
\pgfpathlineto{\pgfqpoint{0.954484in}{2.262189in}}%
\pgfpathlineto{\pgfqpoint{0.967184in}{2.259137in}}%
\pgfpathlineto{\pgfqpoint{0.960593in}{2.232110in}}%
\pgfpathlineto{\pgfqpoint{0.998620in}{2.223097in}}%
\pgfpathlineto{\pgfqpoint{0.993310in}{2.200722in}}%
\pgfpathlineto{\pgfqpoint{0.988985in}{2.182170in}}%
\pgfpathlineto{\pgfqpoint{0.996594in}{2.161889in}}%
\pgfpathlineto{\pgfqpoint{1.000904in}{2.156205in}}%
\pgfpathlineto{\pgfqpoint{1.000929in}{2.147949in}}%
\pgfpathlineto{\pgfqpoint{0.996502in}{2.145181in}}%
\pgfpathlineto{\pgfqpoint{1.001953in}{2.138275in}}%
\pgfpathlineto{\pgfqpoint{0.994972in}{2.136779in}}%
\pgfpathlineto{\pgfqpoint{1.003713in}{2.127225in}}%
\pgfpathlineto{\pgfqpoint{1.003956in}{2.123021in}}%
\pgfpathlineto{\pgfqpoint{1.014258in}{2.117180in}}%
\pgfpathlineto{\pgfqpoint{1.023030in}{2.093793in}}%
\pgfpathlineto{\pgfqpoint{1.027701in}{2.087646in}}%
\pgfpathlineto{\pgfqpoint{0.993104in}{2.095521in}}%
\pgfpathlineto{\pgfqpoint{0.954967in}{2.104591in}}%
\pgfpathlineto{\pgfqpoint{0.956600in}{2.111418in}}%
\pgfpathlineto{\pgfqpoint{0.950036in}{2.114158in}}%
\pgfpathlineto{\pgfqpoint{0.933521in}{2.118231in}}%
\pgfpathlineto{\pgfqpoint{0.921056in}{2.128678in}}%
\pgfpathlineto{\pgfqpoint{0.923573in}{2.138917in}}%
\pgfpathlineto{\pgfqpoint{0.875234in}{2.151233in}}%
\pgfpathclose%
\pgfusepath{fill}%
\end{pgfscope}%
\begin{pgfscope}%
\pgfpathrectangle{\pgfqpoint{0.100000in}{0.100000in}}{\pgfqpoint{3.608454in}{2.310000in}}%
\pgfusepath{clip}%
\pgfsetbuttcap%
\pgfsetmiterjoin%
\definecolor{currentfill}{rgb}{0.000000,0.682353,0.658824}%
\pgfsetfillcolor{currentfill}%
\pgfsetlinewidth{0.000000pt}%
\definecolor{currentstroke}{rgb}{0.000000,0.000000,0.000000}%
\pgfsetstrokecolor{currentstroke}%
\pgfsetstrokeopacity{0.000000}%
\pgfsetdash{}{0pt}%
\pgfpathmoveto{\pgfqpoint{2.187595in}{0.601129in}}%
\pgfpathlineto{\pgfqpoint{2.169818in}{0.598417in}}%
\pgfpathlineto{\pgfqpoint{2.151781in}{0.590398in}}%
\pgfpathlineto{\pgfqpoint{2.151337in}{0.616492in}}%
\pgfpathlineto{\pgfqpoint{2.145490in}{0.616619in}}%
\pgfpathlineto{\pgfqpoint{2.145181in}{0.634496in}}%
\pgfpathlineto{\pgfqpoint{2.134719in}{0.634369in}}%
\pgfpathlineto{\pgfqpoint{2.125194in}{0.664549in}}%
\pgfpathlineto{\pgfqpoint{2.109149in}{0.671435in}}%
\pgfpathlineto{\pgfqpoint{2.105170in}{0.678118in}}%
\pgfpathlineto{\pgfqpoint{2.098891in}{0.678184in}}%
\pgfpathlineto{\pgfqpoint{2.097748in}{0.687767in}}%
\pgfpathlineto{\pgfqpoint{2.093068in}{0.691255in}}%
\pgfpathlineto{\pgfqpoint{2.110954in}{0.708601in}}%
\pgfpathlineto{\pgfqpoint{2.117525in}{0.716961in}}%
\pgfpathlineto{\pgfqpoint{2.123775in}{0.713712in}}%
\pgfpathlineto{\pgfqpoint{2.132588in}{0.712274in}}%
\pgfpathlineto{\pgfqpoint{2.143868in}{0.708025in}}%
\pgfpathlineto{\pgfqpoint{2.166253in}{0.713450in}}%
\pgfpathlineto{\pgfqpoint{2.181123in}{0.718301in}}%
\pgfpathlineto{\pgfqpoint{2.202317in}{0.720035in}}%
\pgfpathlineto{\pgfqpoint{2.205558in}{0.720837in}}%
\pgfpathlineto{\pgfqpoint{2.204984in}{0.705199in}}%
\pgfpathlineto{\pgfqpoint{2.207696in}{0.700477in}}%
\pgfpathlineto{\pgfqpoint{2.204721in}{0.696052in}}%
\pgfpathlineto{\pgfqpoint{2.200790in}{0.680240in}}%
\pgfpathlineto{\pgfqpoint{2.197286in}{0.677217in}}%
\pgfpathlineto{\pgfqpoint{2.193927in}{0.668423in}}%
\pgfpathlineto{\pgfqpoint{2.196505in}{0.661211in}}%
\pgfpathlineto{\pgfqpoint{2.191936in}{0.652528in}}%
\pgfpathlineto{\pgfqpoint{2.196201in}{0.649140in}}%
\pgfpathlineto{\pgfqpoint{2.196138in}{0.630275in}}%
\pgfpathlineto{\pgfqpoint{2.190675in}{0.624923in}}%
\pgfpathlineto{\pgfqpoint{2.188702in}{0.619056in}}%
\pgfpathlineto{\pgfqpoint{2.181395in}{0.610955in}}%
\pgfpathlineto{\pgfqpoint{2.187595in}{0.601129in}}%
\pgfpathclose%
\pgfusepath{fill}%
\end{pgfscope}%
\begin{pgfscope}%
\pgfpathrectangle{\pgfqpoint{0.100000in}{0.100000in}}{\pgfqpoint{3.608454in}{2.310000in}}%
\pgfusepath{clip}%
\pgfsetbuttcap%
\pgfsetmiterjoin%
\definecolor{currentfill}{rgb}{0.000000,0.301961,0.849020}%
\pgfsetfillcolor{currentfill}%
\pgfsetlinewidth{0.000000pt}%
\definecolor{currentstroke}{rgb}{0.000000,0.000000,0.000000}%
\pgfsetstrokecolor{currentstroke}%
\pgfsetstrokeopacity{0.000000}%
\pgfsetdash{}{0pt}%
\pgfpathmoveto{\pgfqpoint{3.217018in}{1.367687in}}%
\pgfpathlineto{\pgfqpoint{3.210448in}{1.365277in}}%
\pgfpathlineto{\pgfqpoint{3.202566in}{1.350706in}}%
\pgfpathlineto{\pgfqpoint{3.197288in}{1.352635in}}%
\pgfpathlineto{\pgfqpoint{3.195177in}{1.357724in}}%
\pgfpathlineto{\pgfqpoint{3.178627in}{1.363629in}}%
\pgfpathlineto{\pgfqpoint{3.168135in}{1.369114in}}%
\pgfpathlineto{\pgfqpoint{3.161085in}{1.369292in}}%
\pgfpathlineto{\pgfqpoint{3.175771in}{1.395430in}}%
\pgfpathlineto{\pgfqpoint{3.175115in}{1.399614in}}%
\pgfpathlineto{\pgfqpoint{3.182634in}{1.411460in}}%
\pgfpathlineto{\pgfqpoint{3.194258in}{1.405953in}}%
\pgfpathlineto{\pgfqpoint{3.196379in}{1.396687in}}%
\pgfpathlineto{\pgfqpoint{3.199440in}{1.393885in}}%
\pgfpathlineto{\pgfqpoint{3.209863in}{1.398797in}}%
\pgfpathlineto{\pgfqpoint{3.214679in}{1.396873in}}%
\pgfpathlineto{\pgfqpoint{3.216571in}{1.392855in}}%
\pgfpathlineto{\pgfqpoint{3.214965in}{1.384823in}}%
\pgfpathlineto{\pgfqpoint{3.209188in}{1.375295in}}%
\pgfpathlineto{\pgfqpoint{3.217018in}{1.367687in}}%
\pgfpathclose%
\pgfusepath{fill}%
\end{pgfscope}%
\begin{pgfscope}%
\pgfpathrectangle{\pgfqpoint{0.100000in}{0.100000in}}{\pgfqpoint{3.608454in}{2.310000in}}%
\pgfusepath{clip}%
\pgfsetbuttcap%
\pgfsetmiterjoin%
\definecolor{currentfill}{rgb}{0.000000,0.333333,0.833333}%
\pgfsetfillcolor{currentfill}%
\pgfsetlinewidth{0.000000pt}%
\definecolor{currentstroke}{rgb}{0.000000,0.000000,0.000000}%
\pgfsetstrokecolor{currentstroke}%
\pgfsetstrokeopacity{0.000000}%
\pgfsetdash{}{0pt}%
\pgfpathmoveto{\pgfqpoint{1.821457in}{1.218805in}}%
\pgfpathlineto{\pgfqpoint{1.820080in}{1.188664in}}%
\pgfpathlineto{\pgfqpoint{1.790983in}{1.190359in}}%
\pgfpathlineto{\pgfqpoint{1.785504in}{1.190712in}}%
\pgfpathlineto{\pgfqpoint{1.787082in}{1.221084in}}%
\pgfpathlineto{\pgfqpoint{1.786488in}{1.228084in}}%
\pgfpathlineto{\pgfqpoint{1.752303in}{1.230115in}}%
\pgfpathlineto{\pgfqpoint{1.754126in}{1.258033in}}%
\pgfpathlineto{\pgfqpoint{1.754193in}{1.271800in}}%
\pgfpathlineto{\pgfqpoint{1.781396in}{1.270043in}}%
\pgfpathlineto{\pgfqpoint{1.782549in}{1.290657in}}%
\pgfpathlineto{\pgfqpoint{1.822395in}{1.288502in}}%
\pgfpathlineto{\pgfqpoint{1.823241in}{1.288458in}}%
\pgfpathlineto{\pgfqpoint{1.821521in}{1.253978in}}%
\pgfpathlineto{\pgfqpoint{1.821831in}{1.246792in}}%
\pgfpathlineto{\pgfqpoint{1.821457in}{1.218805in}}%
\pgfpathclose%
\pgfusepath{fill}%
\end{pgfscope}%
\begin{pgfscope}%
\pgfpathrectangle{\pgfqpoint{0.100000in}{0.100000in}}{\pgfqpoint{3.608454in}{2.310000in}}%
\pgfusepath{clip}%
\pgfsetbuttcap%
\pgfsetmiterjoin%
\definecolor{currentfill}{rgb}{0.000000,0.674510,0.662745}%
\pgfsetfillcolor{currentfill}%
\pgfsetlinewidth{0.000000pt}%
\definecolor{currentstroke}{rgb}{0.000000,0.000000,0.000000}%
\pgfsetstrokecolor{currentstroke}%
\pgfsetstrokeopacity{0.000000}%
\pgfsetdash{}{0pt}%
\pgfpathmoveto{\pgfqpoint{1.759459in}{0.801389in}}%
\pgfpathlineto{\pgfqpoint{1.753852in}{0.801689in}}%
\pgfpathlineto{\pgfqpoint{1.724900in}{0.803670in}}%
\pgfpathlineto{\pgfqpoint{1.729883in}{0.873343in}}%
\pgfpathlineto{\pgfqpoint{1.738977in}{0.872740in}}%
\pgfpathlineto{\pgfqpoint{1.763926in}{0.871040in}}%
\pgfpathlineto{\pgfqpoint{1.759459in}{0.801389in}}%
\pgfpathclose%
\pgfusepath{fill}%
\end{pgfscope}%
\begin{pgfscope}%
\pgfpathrectangle{\pgfqpoint{0.100000in}{0.100000in}}{\pgfqpoint{3.608454in}{2.310000in}}%
\pgfusepath{clip}%
\pgfsetbuttcap%
\pgfsetmiterjoin%
\definecolor{currentfill}{rgb}{0.000000,0.333333,0.833333}%
\pgfsetfillcolor{currentfill}%
\pgfsetlinewidth{0.000000pt}%
\definecolor{currentstroke}{rgb}{0.000000,0.000000,0.000000}%
\pgfsetstrokecolor{currentstroke}%
\pgfsetstrokeopacity{0.000000}%
\pgfsetdash{}{0pt}%
\pgfpathmoveto{\pgfqpoint{1.752303in}{1.230115in}}%
\pgfpathlineto{\pgfqpoint{1.751890in}{1.223219in}}%
\pgfpathlineto{\pgfqpoint{1.725856in}{1.224921in}}%
\pgfpathlineto{\pgfqpoint{1.724493in}{1.225005in}}%
\pgfpathlineto{\pgfqpoint{1.726337in}{1.252520in}}%
\pgfpathlineto{\pgfqpoint{1.666954in}{1.256981in}}%
\pgfpathlineto{\pgfqpoint{1.670021in}{1.298287in}}%
\pgfpathlineto{\pgfqpoint{1.699539in}{1.296124in}}%
\pgfpathlineto{\pgfqpoint{1.754213in}{1.292521in}}%
\pgfpathlineto{\pgfqpoint{1.782549in}{1.290657in}}%
\pgfpathlineto{\pgfqpoint{1.781396in}{1.270043in}}%
\pgfpathlineto{\pgfqpoint{1.754193in}{1.271800in}}%
\pgfpathlineto{\pgfqpoint{1.754126in}{1.258033in}}%
\pgfpathlineto{\pgfqpoint{1.752303in}{1.230115in}}%
\pgfpathclose%
\pgfusepath{fill}%
\end{pgfscope}%
\begin{pgfscope}%
\pgfpathrectangle{\pgfqpoint{0.100000in}{0.100000in}}{\pgfqpoint{3.608454in}{2.310000in}}%
\pgfusepath{clip}%
\pgfsetbuttcap%
\pgfsetmiterjoin%
\definecolor{currentfill}{rgb}{0.000000,0.521569,0.739216}%
\pgfsetfillcolor{currentfill}%
\pgfsetlinewidth{0.000000pt}%
\definecolor{currentstroke}{rgb}{0.000000,0.000000,0.000000}%
\pgfsetstrokecolor{currentstroke}%
\pgfsetstrokeopacity{0.000000}%
\pgfsetdash{}{0pt}%
\pgfpathmoveto{\pgfqpoint{2.280744in}{1.286387in}}%
\pgfpathlineto{\pgfqpoint{2.314915in}{1.287249in}}%
\pgfpathlineto{\pgfqpoint{2.315298in}{1.276599in}}%
\pgfpathlineto{\pgfqpoint{2.322748in}{1.276780in}}%
\pgfpathlineto{\pgfqpoint{2.323611in}{1.248060in}}%
\pgfpathlineto{\pgfqpoint{2.337381in}{1.248460in}}%
\pgfpathlineto{\pgfqpoint{2.337514in}{1.241571in}}%
\pgfpathlineto{\pgfqpoint{2.347355in}{1.241765in}}%
\pgfpathlineto{\pgfqpoint{2.347612in}{1.233131in}}%
\pgfpathlineto{\pgfqpoint{2.358074in}{1.234831in}}%
\pgfpathlineto{\pgfqpoint{2.371740in}{1.235384in}}%
\pgfpathlineto{\pgfqpoint{2.372317in}{1.217256in}}%
\pgfpathlineto{\pgfqpoint{2.375052in}{1.209261in}}%
\pgfpathlineto{\pgfqpoint{2.375334in}{1.200990in}}%
\pgfpathlineto{\pgfqpoint{2.373090in}{1.191695in}}%
\pgfpathlineto{\pgfqpoint{2.361419in}{1.191422in}}%
\pgfpathlineto{\pgfqpoint{2.357912in}{1.194600in}}%
\pgfpathlineto{\pgfqpoint{2.357801in}{1.198383in}}%
\pgfpathlineto{\pgfqpoint{2.350740in}{1.202506in}}%
\pgfpathlineto{\pgfqpoint{2.347208in}{1.211310in}}%
\pgfpathlineto{\pgfqpoint{2.346916in}{1.219387in}}%
\pgfpathlineto{\pgfqpoint{2.310515in}{1.218707in}}%
\pgfpathlineto{\pgfqpoint{2.310037in}{1.232515in}}%
\pgfpathlineto{\pgfqpoint{2.279053in}{1.231943in}}%
\pgfpathlineto{\pgfqpoint{2.278985in}{1.235414in}}%
\pgfpathlineto{\pgfqpoint{2.268526in}{1.240287in}}%
\pgfpathlineto{\pgfqpoint{2.268476in}{1.251773in}}%
\pgfpathlineto{\pgfqpoint{2.268296in}{1.264580in}}%
\pgfpathlineto{\pgfqpoint{2.282058in}{1.264607in}}%
\pgfpathlineto{\pgfqpoint{2.280744in}{1.286387in}}%
\pgfpathclose%
\pgfusepath{fill}%
\end{pgfscope}%
\begin{pgfscope}%
\pgfpathrectangle{\pgfqpoint{0.100000in}{0.100000in}}{\pgfqpoint{3.608454in}{2.310000in}}%
\pgfusepath{clip}%
\pgfsetbuttcap%
\pgfsetmiterjoin%
\definecolor{currentfill}{rgb}{0.000000,0.682353,0.658824}%
\pgfsetfillcolor{currentfill}%
\pgfsetlinewidth{0.000000pt}%
\definecolor{currentstroke}{rgb}{0.000000,0.000000,0.000000}%
\pgfsetstrokecolor{currentstroke}%
\pgfsetstrokeopacity{0.000000}%
\pgfsetdash{}{0pt}%
\pgfpathmoveto{\pgfqpoint{2.831112in}{0.753917in}}%
\pgfpathlineto{\pgfqpoint{2.834431in}{0.723057in}}%
\pgfpathlineto{\pgfqpoint{2.801217in}{0.720980in}}%
\pgfpathlineto{\pgfqpoint{2.795289in}{0.729104in}}%
\pgfpathlineto{\pgfqpoint{2.794346in}{0.734429in}}%
\pgfpathlineto{\pgfqpoint{2.789288in}{0.742878in}}%
\pgfpathlineto{\pgfqpoint{2.786852in}{0.748618in}}%
\pgfpathlineto{\pgfqpoint{2.794089in}{0.749184in}}%
\pgfpathlineto{\pgfqpoint{2.792796in}{0.763762in}}%
\pgfpathlineto{\pgfqpoint{2.811462in}{0.766109in}}%
\pgfpathlineto{\pgfqpoint{2.818573in}{0.766684in}}%
\pgfpathlineto{\pgfqpoint{2.819817in}{0.752621in}}%
\pgfpathlineto{\pgfqpoint{2.831112in}{0.753917in}}%
\pgfpathclose%
\pgfusepath{fill}%
\end{pgfscope}%
\begin{pgfscope}%
\pgfpathrectangle{\pgfqpoint{0.100000in}{0.100000in}}{\pgfqpoint{3.608454in}{2.310000in}}%
\pgfusepath{clip}%
\pgfsetbuttcap%
\pgfsetmiterjoin%
\definecolor{currentfill}{rgb}{0.000000,0.431373,0.784314}%
\pgfsetfillcolor{currentfill}%
\pgfsetlinewidth{0.000000pt}%
\definecolor{currentstroke}{rgb}{0.000000,0.000000,0.000000}%
\pgfsetstrokecolor{currentstroke}%
\pgfsetstrokeopacity{0.000000}%
\pgfsetdash{}{0pt}%
\pgfpathmoveto{\pgfqpoint{2.024967in}{0.596552in}}%
\pgfpathlineto{\pgfqpoint{2.019290in}{0.599307in}}%
\pgfpathlineto{\pgfqpoint{2.014382in}{0.611205in}}%
\pgfpathlineto{\pgfqpoint{2.008648in}{0.612091in}}%
\pgfpathlineto{\pgfqpoint{2.000273in}{0.618823in}}%
\pgfpathlineto{\pgfqpoint{1.990391in}{0.638498in}}%
\pgfpathlineto{\pgfqpoint{1.982721in}{0.639405in}}%
\pgfpathlineto{\pgfqpoint{1.987116in}{0.645931in}}%
\pgfpathlineto{\pgfqpoint{1.993561in}{0.650091in}}%
\pgfpathlineto{\pgfqpoint{2.008653in}{0.653539in}}%
\pgfpathlineto{\pgfqpoint{2.016277in}{0.656724in}}%
\pgfpathlineto{\pgfqpoint{2.025563in}{0.657506in}}%
\pgfpathlineto{\pgfqpoint{2.031314in}{0.647827in}}%
\pgfpathlineto{\pgfqpoint{2.031567in}{0.644083in}}%
\pgfpathlineto{\pgfqpoint{2.051242in}{0.645234in}}%
\pgfpathlineto{\pgfqpoint{2.049731in}{0.675985in}}%
\pgfpathlineto{\pgfqpoint{2.065563in}{0.666170in}}%
\pgfpathlineto{\pgfqpoint{2.081886in}{0.666080in}}%
\pgfpathlineto{\pgfqpoint{2.084351in}{0.694073in}}%
\pgfpathlineto{\pgfqpoint{2.089584in}{0.697805in}}%
\pgfpathlineto{\pgfqpoint{2.093068in}{0.691255in}}%
\pgfpathlineto{\pgfqpoint{2.097748in}{0.687767in}}%
\pgfpathlineto{\pgfqpoint{2.098891in}{0.678184in}}%
\pgfpathlineto{\pgfqpoint{2.105170in}{0.678118in}}%
\pgfpathlineto{\pgfqpoint{2.109149in}{0.671435in}}%
\pgfpathlineto{\pgfqpoint{2.125194in}{0.664549in}}%
\pgfpathlineto{\pgfqpoint{2.134719in}{0.634369in}}%
\pgfpathlineto{\pgfqpoint{2.145181in}{0.634496in}}%
\pgfpathlineto{\pgfqpoint{2.145490in}{0.616619in}}%
\pgfpathlineto{\pgfqpoint{2.151337in}{0.616492in}}%
\pgfpathlineto{\pgfqpoint{2.151781in}{0.590398in}}%
\pgfpathlineto{\pgfqpoint{2.133973in}{0.582254in}}%
\pgfpathlineto{\pgfqpoint{2.137484in}{0.591301in}}%
\pgfpathlineto{\pgfqpoint{2.124852in}{0.587349in}}%
\pgfpathlineto{\pgfqpoint{2.128361in}{0.601241in}}%
\pgfpathlineto{\pgfqpoint{2.125042in}{0.608247in}}%
\pgfpathlineto{\pgfqpoint{2.119618in}{0.605904in}}%
\pgfpathlineto{\pgfqpoint{2.115649in}{0.599011in}}%
\pgfpathlineto{\pgfqpoint{2.109756in}{0.601205in}}%
\pgfpathlineto{\pgfqpoint{2.105684in}{0.596360in}}%
\pgfpathlineto{\pgfqpoint{2.105660in}{0.588738in}}%
\pgfpathlineto{\pgfqpoint{2.113115in}{0.585047in}}%
\pgfpathlineto{\pgfqpoint{2.114156in}{0.572287in}}%
\pgfpathlineto{\pgfqpoint{2.125821in}{0.572042in}}%
\pgfpathlineto{\pgfqpoint{2.098432in}{0.551471in}}%
\pgfpathlineto{\pgfqpoint{2.102863in}{0.561208in}}%
\pgfpathlineto{\pgfqpoint{2.090532in}{0.582574in}}%
\pgfpathlineto{\pgfqpoint{2.091703in}{0.589805in}}%
\pgfpathlineto{\pgfqpoint{2.088573in}{0.592795in}}%
\pgfpathlineto{\pgfqpoint{2.077251in}{0.591805in}}%
\pgfpathlineto{\pgfqpoint{2.074520in}{0.580523in}}%
\pgfpathlineto{\pgfqpoint{2.068450in}{0.580491in}}%
\pgfpathlineto{\pgfqpoint{2.066381in}{0.572643in}}%
\pgfpathlineto{\pgfqpoint{2.061080in}{0.568664in}}%
\pgfpathlineto{\pgfqpoint{2.054185in}{0.571463in}}%
\pgfpathlineto{\pgfqpoint{2.047518in}{0.566563in}}%
\pgfpathlineto{\pgfqpoint{2.040450in}{0.572451in}}%
\pgfpathlineto{\pgfqpoint{2.037581in}{0.580386in}}%
\pgfpathlineto{\pgfqpoint{2.032880in}{0.583631in}}%
\pgfpathlineto{\pgfqpoint{2.032870in}{0.593014in}}%
\pgfpathlineto{\pgfqpoint{2.024967in}{0.596552in}}%
\pgfpathclose%
\pgfusepath{fill}%
\end{pgfscope}%
\begin{pgfscope}%
\pgfpathrectangle{\pgfqpoint{0.100000in}{0.100000in}}{\pgfqpoint{3.608454in}{2.310000in}}%
\pgfusepath{clip}%
\pgfsetbuttcap%
\pgfsetmiterjoin%
\definecolor{currentfill}{rgb}{0.000000,0.439216,0.780392}%
\pgfsetfillcolor{currentfill}%
\pgfsetlinewidth{0.000000pt}%
\definecolor{currentstroke}{rgb}{0.000000,0.000000,0.000000}%
\pgfsetstrokecolor{currentstroke}%
\pgfsetstrokeopacity{0.000000}%
\pgfsetdash{}{0pt}%
\pgfpathmoveto{\pgfqpoint{1.938143in}{0.938336in}}%
\pgfpathlineto{\pgfqpoint{1.931087in}{0.945835in}}%
\pgfpathlineto{\pgfqpoint{1.917943in}{0.934949in}}%
\pgfpathlineto{\pgfqpoint{1.910664in}{0.938402in}}%
\pgfpathlineto{\pgfqpoint{1.913020in}{0.946294in}}%
\pgfpathlineto{\pgfqpoint{1.906317in}{0.946847in}}%
\pgfpathlineto{\pgfqpoint{1.898581in}{0.956614in}}%
\pgfpathlineto{\pgfqpoint{1.885921in}{0.960469in}}%
\pgfpathlineto{\pgfqpoint{1.883512in}{0.956032in}}%
\pgfpathlineto{\pgfqpoint{1.877716in}{0.953284in}}%
\pgfpathlineto{\pgfqpoint{1.869882in}{0.961093in}}%
\pgfpathlineto{\pgfqpoint{1.870481in}{0.975087in}}%
\pgfpathlineto{\pgfqpoint{1.867064in}{0.975207in}}%
\pgfpathlineto{\pgfqpoint{1.867617in}{0.988984in}}%
\pgfpathlineto{\pgfqpoint{1.856890in}{0.989419in}}%
\pgfpathlineto{\pgfqpoint{1.857172in}{0.996308in}}%
\pgfpathlineto{\pgfqpoint{1.858037in}{1.016994in}}%
\pgfpathlineto{\pgfqpoint{1.871452in}{1.016408in}}%
\pgfpathlineto{\pgfqpoint{1.905912in}{1.015133in}}%
\pgfpathlineto{\pgfqpoint{1.905475in}{1.001362in}}%
\pgfpathlineto{\pgfqpoint{1.932881in}{1.000512in}}%
\pgfpathlineto{\pgfqpoint{1.939766in}{1.000299in}}%
\pgfpathlineto{\pgfqpoint{1.938143in}{0.938336in}}%
\pgfpathclose%
\pgfusepath{fill}%
\end{pgfscope}%
\begin{pgfscope}%
\pgfpathrectangle{\pgfqpoint{0.100000in}{0.100000in}}{\pgfqpoint{3.608454in}{2.310000in}}%
\pgfusepath{clip}%
\pgfsetbuttcap%
\pgfsetmiterjoin%
\definecolor{currentfill}{rgb}{0.000000,0.498039,0.750980}%
\pgfsetfillcolor{currentfill}%
\pgfsetlinewidth{0.000000pt}%
\definecolor{currentstroke}{rgb}{0.000000,0.000000,0.000000}%
\pgfsetstrokecolor{currentstroke}%
\pgfsetstrokeopacity{0.000000}%
\pgfsetdash{}{0pt}%
\pgfpathmoveto{\pgfqpoint{3.358769in}{1.729605in}}%
\pgfpathlineto{\pgfqpoint{3.358373in}{1.723673in}}%
\pgfpathlineto{\pgfqpoint{3.347348in}{1.723269in}}%
\pgfpathlineto{\pgfqpoint{3.335015in}{1.726252in}}%
\pgfpathlineto{\pgfqpoint{3.333233in}{1.729959in}}%
\pgfpathlineto{\pgfqpoint{3.325396in}{1.731486in}}%
\pgfpathlineto{\pgfqpoint{3.324878in}{1.725364in}}%
\pgfpathlineto{\pgfqpoint{3.311248in}{1.723631in}}%
\pgfpathlineto{\pgfqpoint{3.302028in}{1.725960in}}%
\pgfpathlineto{\pgfqpoint{3.296564in}{1.727321in}}%
\pgfpathlineto{\pgfqpoint{3.300176in}{1.740144in}}%
\pgfpathlineto{\pgfqpoint{3.287975in}{1.743465in}}%
\pgfpathlineto{\pgfqpoint{3.281055in}{1.749476in}}%
\pgfpathlineto{\pgfqpoint{3.283745in}{1.758930in}}%
\pgfpathlineto{\pgfqpoint{3.279470in}{1.768032in}}%
\pgfpathlineto{\pgfqpoint{3.279173in}{1.774369in}}%
\pgfpathlineto{\pgfqpoint{3.291128in}{1.772508in}}%
\pgfpathlineto{\pgfqpoint{3.300308in}{1.775508in}}%
\pgfpathlineto{\pgfqpoint{3.309890in}{1.786988in}}%
\pgfpathlineto{\pgfqpoint{3.287566in}{1.849431in}}%
\pgfpathlineto{\pgfqpoint{3.296075in}{1.852707in}}%
\pgfpathlineto{\pgfqpoint{3.332649in}{1.866641in}}%
\pgfpathlineto{\pgfqpoint{3.333464in}{1.857292in}}%
\pgfpathlineto{\pgfqpoint{3.344700in}{1.848799in}}%
\pgfpathlineto{\pgfqpoint{3.348859in}{1.830296in}}%
\pgfpathlineto{\pgfqpoint{3.356483in}{1.793242in}}%
\pgfpathlineto{\pgfqpoint{3.358955in}{1.786375in}}%
\pgfpathlineto{\pgfqpoint{3.357447in}{1.732265in}}%
\pgfpathlineto{\pgfqpoint{3.358769in}{1.729605in}}%
\pgfpathclose%
\pgfusepath{fill}%
\end{pgfscope}%
\begin{pgfscope}%
\pgfpathrectangle{\pgfqpoint{0.100000in}{0.100000in}}{\pgfqpoint{3.608454in}{2.310000in}}%
\pgfusepath{clip}%
\pgfsetbuttcap%
\pgfsetmiterjoin%
\definecolor{currentfill}{rgb}{0.000000,0.466667,0.766667}%
\pgfsetfillcolor{currentfill}%
\pgfsetlinewidth{0.000000pt}%
\definecolor{currentstroke}{rgb}{0.000000,0.000000,0.000000}%
\pgfsetstrokecolor{currentstroke}%
\pgfsetstrokeopacity{0.000000}%
\pgfsetdash{}{0pt}%
\pgfpathmoveto{\pgfqpoint{1.364865in}{1.868873in}}%
\pgfpathlineto{\pgfqpoint{1.318547in}{1.875904in}}%
\pgfpathlineto{\pgfqpoint{1.317490in}{1.876575in}}%
\pgfpathlineto{\pgfqpoint{1.278870in}{1.882470in}}%
\pgfpathlineto{\pgfqpoint{1.281086in}{1.895779in}}%
\pgfpathlineto{\pgfqpoint{1.266400in}{1.898228in}}%
\pgfpathlineto{\pgfqpoint{1.268744in}{1.912078in}}%
\pgfpathlineto{\pgfqpoint{1.275970in}{1.910859in}}%
\pgfpathlineto{\pgfqpoint{1.278245in}{1.924515in}}%
\pgfpathlineto{\pgfqpoint{1.292921in}{1.929061in}}%
\pgfpathlineto{\pgfqpoint{1.299700in}{1.927978in}}%
\pgfpathlineto{\pgfqpoint{1.301786in}{1.941660in}}%
\pgfpathlineto{\pgfqpoint{1.307165in}{1.954995in}}%
\pgfpathlineto{\pgfqpoint{1.311739in}{1.954347in}}%
\pgfpathlineto{\pgfqpoint{1.313126in}{1.960922in}}%
\pgfpathlineto{\pgfqpoint{1.302894in}{1.962642in}}%
\pgfpathlineto{\pgfqpoint{1.301322in}{1.969905in}}%
\pgfpathlineto{\pgfqpoint{1.302428in}{1.976762in}}%
\pgfpathlineto{\pgfqpoint{1.316114in}{1.974564in}}%
\pgfpathlineto{\pgfqpoint{1.322819in}{2.011773in}}%
\pgfpathlineto{\pgfqpoint{1.323552in}{2.016288in}}%
\pgfpathlineto{\pgfqpoint{1.364475in}{2.009800in}}%
\pgfpathlineto{\pgfqpoint{1.407997in}{2.003940in}}%
\pgfpathlineto{\pgfqpoint{1.409028in}{1.997695in}}%
\pgfpathlineto{\pgfqpoint{1.406757in}{1.983167in}}%
\pgfpathlineto{\pgfqpoint{1.410750in}{1.980307in}}%
\pgfpathlineto{\pgfqpoint{1.414681in}{1.966931in}}%
\pgfpathlineto{\pgfqpoint{1.421242in}{1.955525in}}%
\pgfpathlineto{\pgfqpoint{1.417308in}{1.945271in}}%
\pgfpathlineto{\pgfqpoint{1.408401in}{1.946543in}}%
\pgfpathlineto{\pgfqpoint{1.407750in}{1.941975in}}%
\pgfpathlineto{\pgfqpoint{1.400977in}{1.942930in}}%
\pgfpathlineto{\pgfqpoint{1.395427in}{1.936746in}}%
\pgfpathlineto{\pgfqpoint{1.386336in}{1.938080in}}%
\pgfpathlineto{\pgfqpoint{1.383830in}{1.929135in}}%
\pgfpathlineto{\pgfqpoint{1.380767in}{1.908242in}}%
\pgfpathlineto{\pgfqpoint{1.366491in}{1.905688in}}%
\pgfpathlineto{\pgfqpoint{1.348402in}{1.908679in}}%
\pgfpathlineto{\pgfqpoint{1.345983in}{1.906551in}}%
\pgfpathlineto{\pgfqpoint{1.342451in}{1.889782in}}%
\pgfpathlineto{\pgfqpoint{1.375793in}{1.885029in}}%
\pgfpathlineto{\pgfqpoint{1.372850in}{1.883266in}}%
\pgfpathlineto{\pgfqpoint{1.364865in}{1.868873in}}%
\pgfpathclose%
\pgfusepath{fill}%
\end{pgfscope}%
\begin{pgfscope}%
\pgfpathrectangle{\pgfqpoint{0.100000in}{0.100000in}}{\pgfqpoint{3.608454in}{2.310000in}}%
\pgfusepath{clip}%
\pgfsetbuttcap%
\pgfsetmiterjoin%
\definecolor{currentfill}{rgb}{0.000000,0.647059,0.676471}%
\pgfsetfillcolor{currentfill}%
\pgfsetlinewidth{0.000000pt}%
\definecolor{currentstroke}{rgb}{0.000000,0.000000,0.000000}%
\pgfsetstrokecolor{currentstroke}%
\pgfsetstrokeopacity{0.000000}%
\pgfsetdash{}{0pt}%
\pgfpathmoveto{\pgfqpoint{2.893588in}{0.882495in}}%
\pgfpathlineto{\pgfqpoint{2.878409in}{0.905708in}}%
\pgfpathlineto{\pgfqpoint{2.881402in}{0.908328in}}%
\pgfpathlineto{\pgfqpoint{2.874063in}{0.928161in}}%
\pgfpathlineto{\pgfqpoint{2.866524in}{0.926111in}}%
\pgfpathlineto{\pgfqpoint{2.864668in}{0.946906in}}%
\pgfpathlineto{\pgfqpoint{2.880855in}{0.952903in}}%
\pgfpathlineto{\pgfqpoint{2.889760in}{0.943778in}}%
\pgfpathlineto{\pgfqpoint{2.898524in}{0.954072in}}%
\pgfpathlineto{\pgfqpoint{2.909363in}{0.953452in}}%
\pgfpathlineto{\pgfqpoint{2.918368in}{0.938239in}}%
\pgfpathlineto{\pgfqpoint{2.911845in}{0.934034in}}%
\pgfpathlineto{\pgfqpoint{2.910167in}{0.929143in}}%
\pgfpathlineto{\pgfqpoint{2.905128in}{0.926113in}}%
\pgfpathlineto{\pgfqpoint{2.899562in}{0.918887in}}%
\pgfpathlineto{\pgfqpoint{2.899980in}{0.912387in}}%
\pgfpathlineto{\pgfqpoint{2.903959in}{0.905713in}}%
\pgfpathlineto{\pgfqpoint{2.909397in}{0.902008in}}%
\pgfpathlineto{\pgfqpoint{2.910191in}{0.894634in}}%
\pgfpathlineto{\pgfqpoint{2.893588in}{0.882495in}}%
\pgfpathclose%
\pgfusepath{fill}%
\end{pgfscope}%
\begin{pgfscope}%
\pgfpathrectangle{\pgfqpoint{0.100000in}{0.100000in}}{\pgfqpoint{3.608454in}{2.310000in}}%
\pgfusepath{clip}%
\pgfsetbuttcap%
\pgfsetmiterjoin%
\definecolor{currentfill}{rgb}{0.000000,0.611765,0.694118}%
\pgfsetfillcolor{currentfill}%
\pgfsetlinewidth{0.000000pt}%
\definecolor{currentstroke}{rgb}{0.000000,0.000000,0.000000}%
\pgfsetstrokecolor{currentstroke}%
\pgfsetstrokeopacity{0.000000}%
\pgfsetdash{}{0pt}%
\pgfpathmoveto{\pgfqpoint{1.496354in}{2.088951in}}%
\pgfpathlineto{\pgfqpoint{1.498672in}{2.081480in}}%
\pgfpathlineto{\pgfqpoint{1.503848in}{2.080824in}}%
\pgfpathlineto{\pgfqpoint{1.500521in}{2.053891in}}%
\pgfpathlineto{\pgfqpoint{1.496777in}{2.040294in}}%
\pgfpathlineto{\pgfqpoint{1.503629in}{2.039468in}}%
\pgfpathlineto{\pgfqpoint{1.501940in}{2.025689in}}%
\pgfpathlineto{\pgfqpoint{1.506437in}{2.025137in}}%
\pgfpathlineto{\pgfqpoint{1.503208in}{1.998589in}}%
\pgfpathlineto{\pgfqpoint{1.468884in}{2.002862in}}%
\pgfpathlineto{\pgfqpoint{1.469044in}{2.003996in}}%
\pgfpathlineto{\pgfqpoint{1.405583in}{2.011993in}}%
\pgfpathlineto{\pgfqpoint{1.403304in}{2.014885in}}%
\pgfpathlineto{\pgfqpoint{1.407470in}{2.024546in}}%
\pgfpathlineto{\pgfqpoint{1.405800in}{2.033757in}}%
\pgfpathlineto{\pgfqpoint{1.406013in}{2.044122in}}%
\pgfpathlineto{\pgfqpoint{1.411450in}{2.059486in}}%
\pgfpathlineto{\pgfqpoint{1.414702in}{2.064290in}}%
\pgfpathlineto{\pgfqpoint{1.422183in}{2.067656in}}%
\pgfpathlineto{\pgfqpoint{1.428142in}{2.073167in}}%
\pgfpathlineto{\pgfqpoint{1.437875in}{2.069615in}}%
\pgfpathlineto{\pgfqpoint{1.440516in}{2.074870in}}%
\pgfpathlineto{\pgfqpoint{1.449108in}{2.071253in}}%
\pgfpathlineto{\pgfqpoint{1.468292in}{2.069917in}}%
\pgfpathlineto{\pgfqpoint{1.470183in}{2.074395in}}%
\pgfpathlineto{\pgfqpoint{1.476331in}{2.073148in}}%
\pgfpathlineto{\pgfqpoint{1.487031in}{2.078600in}}%
\pgfpathlineto{\pgfqpoint{1.492326in}{2.087153in}}%
\pgfpathlineto{\pgfqpoint{1.496354in}{2.088951in}}%
\pgfpathclose%
\pgfusepath{fill}%
\end{pgfscope}%
\begin{pgfscope}%
\pgfpathrectangle{\pgfqpoint{0.100000in}{0.100000in}}{\pgfqpoint{3.608454in}{2.310000in}}%
\pgfusepath{clip}%
\pgfsetbuttcap%
\pgfsetmiterjoin%
\definecolor{currentfill}{rgb}{0.000000,0.423529,0.788235}%
\pgfsetfillcolor{currentfill}%
\pgfsetlinewidth{0.000000pt}%
\definecolor{currentstroke}{rgb}{0.000000,0.000000,0.000000}%
\pgfsetstrokecolor{currentstroke}%
\pgfsetstrokeopacity{0.000000}%
\pgfsetdash{}{0pt}%
\pgfpathmoveto{\pgfqpoint{3.153317in}{1.449269in}}%
\pgfpathlineto{\pgfqpoint{3.117245in}{1.469315in}}%
\pgfpathlineto{\pgfqpoint{3.109226in}{1.471916in}}%
\pgfpathlineto{\pgfqpoint{3.110506in}{1.479451in}}%
\pgfpathlineto{\pgfqpoint{3.116008in}{1.487243in}}%
\pgfpathlineto{\pgfqpoint{3.111561in}{1.488897in}}%
\pgfpathlineto{\pgfqpoint{3.114952in}{1.497798in}}%
\pgfpathlineto{\pgfqpoint{3.117244in}{1.511717in}}%
\pgfpathlineto{\pgfqpoint{3.116662in}{1.519085in}}%
\pgfpathlineto{\pgfqpoint{3.119792in}{1.526024in}}%
\pgfpathlineto{\pgfqpoint{3.133137in}{1.523448in}}%
\pgfpathlineto{\pgfqpoint{3.137378in}{1.521010in}}%
\pgfpathlineto{\pgfqpoint{3.143924in}{1.538541in}}%
\pgfpathlineto{\pgfqpoint{3.146592in}{1.539583in}}%
\pgfpathlineto{\pgfqpoint{3.162472in}{1.516430in}}%
\pgfpathlineto{\pgfqpoint{3.166603in}{1.499150in}}%
\pgfpathlineto{\pgfqpoint{3.161516in}{1.490087in}}%
\pgfpathlineto{\pgfqpoint{3.159520in}{1.465911in}}%
\pgfpathlineto{\pgfqpoint{3.156244in}{1.464323in}}%
\pgfpathlineto{\pgfqpoint{3.153317in}{1.449269in}}%
\pgfpathclose%
\pgfusepath{fill}%
\end{pgfscope}%
\begin{pgfscope}%
\pgfpathrectangle{\pgfqpoint{0.100000in}{0.100000in}}{\pgfqpoint{3.608454in}{2.310000in}}%
\pgfusepath{clip}%
\pgfsetbuttcap%
\pgfsetmiterjoin%
\definecolor{currentfill}{rgb}{0.000000,0.827451,0.586275}%
\pgfsetfillcolor{currentfill}%
\pgfsetlinewidth{0.000000pt}%
\definecolor{currentstroke}{rgb}{0.000000,0.000000,0.000000}%
\pgfsetstrokecolor{currentstroke}%
\pgfsetstrokeopacity{0.000000}%
\pgfsetdash{}{0pt}%
\pgfpathmoveto{\pgfqpoint{0.735331in}{0.472883in}}%
\pgfpathlineto{\pgfqpoint{0.738160in}{0.475971in}}%
\pgfpathlineto{\pgfqpoint{0.743525in}{0.474782in}}%
\pgfpathlineto{\pgfqpoint{0.740586in}{0.472887in}}%
\pgfpathlineto{\pgfqpoint{0.735331in}{0.472883in}}%
\pgfpathclose%
\pgfusepath{fill}%
\end{pgfscope}%
\begin{pgfscope}%
\pgfpathrectangle{\pgfqpoint{0.100000in}{0.100000in}}{\pgfqpoint{3.608454in}{2.310000in}}%
\pgfusepath{clip}%
\pgfsetbuttcap%
\pgfsetmiterjoin%
\definecolor{currentfill}{rgb}{0.000000,0.827451,0.586275}%
\pgfsetfillcolor{currentfill}%
\pgfsetlinewidth{0.000000pt}%
\definecolor{currentstroke}{rgb}{0.000000,0.000000,0.000000}%
\pgfsetstrokecolor{currentstroke}%
\pgfsetstrokeopacity{0.000000}%
\pgfsetdash{}{0pt}%
\pgfpathmoveto{\pgfqpoint{0.864761in}{0.455531in}}%
\pgfpathlineto{\pgfqpoint{0.863736in}{0.453184in}}%
\pgfpathlineto{\pgfqpoint{0.857604in}{0.444984in}}%
\pgfpathlineto{\pgfqpoint{0.856061in}{0.446120in}}%
\pgfpathlineto{\pgfqpoint{0.849915in}{0.438009in}}%
\pgfpathlineto{\pgfqpoint{0.848446in}{0.439078in}}%
\pgfpathlineto{\pgfqpoint{0.842589in}{0.431260in}}%
\pgfpathlineto{\pgfqpoint{0.841100in}{0.432388in}}%
\pgfpathlineto{\pgfqpoint{0.839547in}{0.430349in}}%
\pgfpathlineto{\pgfqpoint{0.837531in}{0.431857in}}%
\pgfpathlineto{\pgfqpoint{0.834469in}{0.427785in}}%
\pgfpathlineto{\pgfqpoint{0.832475in}{0.429280in}}%
\pgfpathlineto{\pgfqpoint{0.830966in}{0.427222in}}%
\pgfpathlineto{\pgfqpoint{0.827494in}{0.429779in}}%
\pgfpathlineto{\pgfqpoint{0.824498in}{0.425742in}}%
\pgfpathlineto{\pgfqpoint{0.820512in}{0.428750in}}%
\pgfpathlineto{\pgfqpoint{0.817435in}{0.424709in}}%
\pgfpathlineto{\pgfqpoint{0.818042in}{0.424203in}}%
\pgfpathlineto{\pgfqpoint{0.812034in}{0.416427in}}%
\pgfpathlineto{\pgfqpoint{0.811149in}{0.413923in}}%
\pgfpathlineto{\pgfqpoint{0.815079in}{0.410909in}}%
\pgfpathlineto{\pgfqpoint{0.811859in}{0.406633in}}%
\pgfpathlineto{\pgfqpoint{0.809601in}{0.406866in}}%
\pgfpathlineto{\pgfqpoint{0.808196in}{0.405860in}}%
\pgfpathlineto{\pgfqpoint{0.804050in}{0.407001in}}%
\pgfpathlineto{\pgfqpoint{0.800501in}{0.409249in}}%
\pgfpathlineto{\pgfqpoint{0.798949in}{0.408471in}}%
\pgfpathlineto{\pgfqpoint{0.797005in}{0.409992in}}%
\pgfpathlineto{\pgfqpoint{0.794932in}{0.410045in}}%
\pgfpathlineto{\pgfqpoint{0.789736in}{0.403394in}}%
\pgfpathlineto{\pgfqpoint{0.786031in}{0.403802in}}%
\pgfpathlineto{\pgfqpoint{0.784195in}{0.405252in}}%
\pgfpathlineto{\pgfqpoint{0.781062in}{0.403097in}}%
\pgfpathlineto{\pgfqpoint{0.777277in}{0.402089in}}%
\pgfpathlineto{\pgfqpoint{0.777113in}{0.402156in}}%
\pgfpathlineto{\pgfqpoint{0.777652in}{0.403284in}}%
\pgfpathlineto{\pgfqpoint{0.773253in}{0.404050in}}%
\pgfpathlineto{\pgfqpoint{0.771380in}{0.405564in}}%
\pgfpathlineto{\pgfqpoint{0.768215in}{0.404903in}}%
\pgfpathlineto{\pgfqpoint{0.766993in}{0.403276in}}%
\pgfpathlineto{\pgfqpoint{0.765003in}{0.403792in}}%
\pgfpathlineto{\pgfqpoint{0.761612in}{0.402864in}}%
\pgfpathlineto{\pgfqpoint{0.760327in}{0.401259in}}%
\pgfpathlineto{\pgfqpoint{0.764308in}{0.397983in}}%
\pgfpathlineto{\pgfqpoint{0.759811in}{0.395476in}}%
\pgfpathlineto{\pgfqpoint{0.758622in}{0.397717in}}%
\pgfpathlineto{\pgfqpoint{0.754992in}{0.396142in}}%
\pgfpathlineto{\pgfqpoint{0.753722in}{0.397568in}}%
\pgfpathlineto{\pgfqpoint{0.743096in}{0.399381in}}%
\pgfpathlineto{\pgfqpoint{0.741766in}{0.396518in}}%
\pgfpathlineto{\pgfqpoint{0.739332in}{0.398261in}}%
\pgfpathlineto{\pgfqpoint{0.739251in}{0.400641in}}%
\pgfpathlineto{\pgfqpoint{0.736510in}{0.401202in}}%
\pgfpathlineto{\pgfqpoint{0.736143in}{0.398969in}}%
\pgfpathlineto{\pgfqpoint{0.734429in}{0.397349in}}%
\pgfpathlineto{\pgfqpoint{0.730331in}{0.398641in}}%
\pgfpathlineto{\pgfqpoint{0.727510in}{0.401229in}}%
\pgfpathlineto{\pgfqpoint{0.723876in}{0.399715in}}%
\pgfpathlineto{\pgfqpoint{0.726580in}{0.395915in}}%
\pgfpathlineto{\pgfqpoint{0.720792in}{0.394591in}}%
\pgfpathlineto{\pgfqpoint{0.718356in}{0.391937in}}%
\pgfpathlineto{\pgfqpoint{0.717498in}{0.394606in}}%
\pgfpathlineto{\pgfqpoint{0.714729in}{0.394241in}}%
\pgfpathlineto{\pgfqpoint{0.712644in}{0.395735in}}%
\pgfpathlineto{\pgfqpoint{0.710848in}{0.395775in}}%
\pgfpathlineto{\pgfqpoint{0.705221in}{0.398449in}}%
\pgfpathlineto{\pgfqpoint{0.702788in}{0.397277in}}%
\pgfpathlineto{\pgfqpoint{0.700932in}{0.397327in}}%
\pgfpathlineto{\pgfqpoint{0.698256in}{0.395576in}}%
\pgfpathlineto{\pgfqpoint{0.702855in}{0.400929in}}%
\pgfpathlineto{\pgfqpoint{0.699223in}{0.404064in}}%
\pgfpathlineto{\pgfqpoint{0.701711in}{0.405421in}}%
\pgfpathlineto{\pgfqpoint{0.707923in}{0.412710in}}%
\pgfpathlineto{\pgfqpoint{0.706826in}{0.413694in}}%
\pgfpathlineto{\pgfqpoint{0.710015in}{0.417271in}}%
\pgfpathlineto{\pgfqpoint{0.713680in}{0.414200in}}%
\pgfpathlineto{\pgfqpoint{0.715284in}{0.416046in}}%
\pgfpathlineto{\pgfqpoint{0.719039in}{0.412825in}}%
\pgfpathlineto{\pgfqpoint{0.720711in}{0.414781in}}%
\pgfpathlineto{\pgfqpoint{0.723593in}{0.412237in}}%
\pgfpathlineto{\pgfqpoint{0.725147in}{0.414065in}}%
\pgfpathlineto{\pgfqpoint{0.727666in}{0.413737in}}%
\pgfpathlineto{\pgfqpoint{0.730412in}{0.411143in}}%
\pgfpathlineto{\pgfqpoint{0.733440in}{0.415672in}}%
\pgfpathlineto{\pgfqpoint{0.739202in}{0.417909in}}%
\pgfpathlineto{\pgfqpoint{0.745161in}{0.418898in}}%
\pgfpathlineto{\pgfqpoint{0.748345in}{0.417706in}}%
\pgfpathlineto{\pgfqpoint{0.751670in}{0.419346in}}%
\pgfpathlineto{\pgfqpoint{0.756304in}{0.418993in}}%
\pgfpathlineto{\pgfqpoint{0.757747in}{0.421251in}}%
\pgfpathlineto{\pgfqpoint{0.768447in}{0.430721in}}%
\pgfpathlineto{\pgfqpoint{0.770758in}{0.430833in}}%
\pgfpathlineto{\pgfqpoint{0.770614in}{0.432762in}}%
\pgfpathlineto{\pgfqpoint{0.774098in}{0.436410in}}%
\pgfpathlineto{\pgfqpoint{0.779568in}{0.437653in}}%
\pgfpathlineto{\pgfqpoint{0.790893in}{0.428785in}}%
\pgfpathlineto{\pgfqpoint{0.796011in}{0.435184in}}%
\pgfpathlineto{\pgfqpoint{0.787216in}{0.442053in}}%
\pgfpathlineto{\pgfqpoint{0.772048in}{0.445375in}}%
\pgfpathlineto{\pgfqpoint{0.769184in}{0.446551in}}%
\pgfpathlineto{\pgfqpoint{0.768571in}{0.448855in}}%
\pgfpathlineto{\pgfqpoint{0.769448in}{0.452003in}}%
\pgfpathlineto{\pgfqpoint{0.768194in}{0.453396in}}%
\pgfpathlineto{\pgfqpoint{0.771817in}{0.456869in}}%
\pgfpathlineto{\pgfqpoint{0.770582in}{0.457460in}}%
\pgfpathlineto{\pgfqpoint{0.763393in}{0.455439in}}%
\pgfpathlineto{\pgfqpoint{0.761052in}{0.449862in}}%
\pgfpathlineto{\pgfqpoint{0.759380in}{0.447981in}}%
\pgfpathlineto{\pgfqpoint{0.757127in}{0.448341in}}%
\pgfpathlineto{\pgfqpoint{0.755962in}{0.451093in}}%
\pgfpathlineto{\pgfqpoint{0.756950in}{0.451960in}}%
\pgfpathlineto{\pgfqpoint{0.756995in}{0.455328in}}%
\pgfpathlineto{\pgfqpoint{0.758416in}{0.463072in}}%
\pgfpathlineto{\pgfqpoint{0.757580in}{0.464630in}}%
\pgfpathlineto{\pgfqpoint{0.759553in}{0.467203in}}%
\pgfpathlineto{\pgfqpoint{0.757448in}{0.468248in}}%
\pgfpathlineto{\pgfqpoint{0.756266in}{0.467141in}}%
\pgfpathlineto{\pgfqpoint{0.752660in}{0.466199in}}%
\pgfpathlineto{\pgfqpoint{0.752967in}{0.468537in}}%
\pgfpathlineto{\pgfqpoint{0.751538in}{0.473475in}}%
\pgfpathlineto{\pgfqpoint{0.753567in}{0.476503in}}%
\pgfpathlineto{\pgfqpoint{0.750953in}{0.477091in}}%
\pgfpathlineto{\pgfqpoint{0.745672in}{0.477306in}}%
\pgfpathlineto{\pgfqpoint{0.740722in}{0.478712in}}%
\pgfpathlineto{\pgfqpoint{0.744357in}{0.482168in}}%
\pgfpathlineto{\pgfqpoint{0.746173in}{0.480590in}}%
\pgfpathlineto{\pgfqpoint{0.747682in}{0.482304in}}%
\pgfpathlineto{\pgfqpoint{0.749547in}{0.480704in}}%
\pgfpathlineto{\pgfqpoint{0.752873in}{0.484455in}}%
\pgfpathlineto{\pgfqpoint{0.755169in}{0.485784in}}%
\pgfpathlineto{\pgfqpoint{0.756995in}{0.484116in}}%
\pgfpathlineto{\pgfqpoint{0.759307in}{0.485494in}}%
\pgfpathlineto{\pgfqpoint{0.762636in}{0.489281in}}%
\pgfpathlineto{\pgfqpoint{0.764542in}{0.487613in}}%
\pgfpathlineto{\pgfqpoint{0.767845in}{0.491365in}}%
\pgfpathlineto{\pgfqpoint{0.770402in}{0.489224in}}%
\pgfpathlineto{\pgfqpoint{0.773696in}{0.492989in}}%
\pgfpathlineto{\pgfqpoint{0.775519in}{0.491429in}}%
\pgfpathlineto{\pgfqpoint{0.778697in}{0.494994in}}%
\pgfpathlineto{\pgfqpoint{0.781042in}{0.496290in}}%
\pgfpathlineto{\pgfqpoint{0.782945in}{0.494696in}}%
\pgfpathlineto{\pgfqpoint{0.786206in}{0.498438in}}%
\pgfpathlineto{\pgfqpoint{0.788112in}{0.496794in}}%
\pgfpathlineto{\pgfqpoint{0.789763in}{0.498705in}}%
\pgfpathlineto{\pgfqpoint{0.791902in}{0.496847in}}%
\pgfpathlineto{\pgfqpoint{0.795190in}{0.500666in}}%
\pgfpathlineto{\pgfqpoint{0.797356in}{0.502139in}}%
\pgfpathlineto{\pgfqpoint{0.799211in}{0.500636in}}%
\pgfpathlineto{\pgfqpoint{0.800941in}{0.502662in}}%
\pgfpathlineto{\pgfqpoint{0.802432in}{0.501473in}}%
\pgfpathlineto{\pgfqpoint{0.804074in}{0.503387in}}%
\pgfpathlineto{\pgfqpoint{0.814957in}{0.494457in}}%
\pgfpathlineto{\pgfqpoint{0.813395in}{0.492489in}}%
\pgfpathlineto{\pgfqpoint{0.815390in}{0.490823in}}%
\pgfpathlineto{\pgfqpoint{0.816987in}{0.492818in}}%
\pgfpathlineto{\pgfqpoint{0.819027in}{0.491316in}}%
\pgfpathlineto{\pgfqpoint{0.820776in}{0.493434in}}%
\pgfpathlineto{\pgfqpoint{0.824629in}{0.490280in}}%
\pgfpathlineto{\pgfqpoint{0.822897in}{0.488128in}}%
\pgfpathlineto{\pgfqpoint{0.836300in}{0.477446in}}%
\pgfpathlineto{\pgfqpoint{0.853555in}{0.463966in}}%
\pgfpathlineto{\pgfqpoint{0.864761in}{0.455531in}}%
\pgfpathclose%
\pgfusepath{fill}%
\end{pgfscope}%
\begin{pgfscope}%
\pgfpathrectangle{\pgfqpoint{0.100000in}{0.100000in}}{\pgfqpoint{3.608454in}{2.310000in}}%
\pgfusepath{clip}%
\pgfsetbuttcap%
\pgfsetmiterjoin%
\definecolor{currentfill}{rgb}{0.000000,0.827451,0.586275}%
\pgfsetfillcolor{currentfill}%
\pgfsetlinewidth{0.000000pt}%
\definecolor{currentstroke}{rgb}{0.000000,0.000000,0.000000}%
\pgfsetstrokecolor{currentstroke}%
\pgfsetstrokeopacity{0.000000}%
\pgfsetdash{}{0pt}%
\pgfpathmoveto{\pgfqpoint{0.735963in}{0.479492in}}%
\pgfpathlineto{\pgfqpoint{0.737273in}{0.481661in}}%
\pgfpathlineto{\pgfqpoint{0.740538in}{0.478732in}}%
\pgfpathlineto{\pgfqpoint{0.735963in}{0.479492in}}%
\pgfpathclose%
\pgfusepath{fill}%
\end{pgfscope}%
\begin{pgfscope}%
\pgfpathrectangle{\pgfqpoint{0.100000in}{0.100000in}}{\pgfqpoint{3.608454in}{2.310000in}}%
\pgfusepath{clip}%
\pgfsetbuttcap%
\pgfsetmiterjoin%
\definecolor{currentfill}{rgb}{0.000000,0.694118,0.652941}%
\pgfsetfillcolor{currentfill}%
\pgfsetlinewidth{0.000000pt}%
\definecolor{currentstroke}{rgb}{0.000000,0.000000,0.000000}%
\pgfsetstrokecolor{currentstroke}%
\pgfsetstrokeopacity{0.000000}%
\pgfsetdash{}{0pt}%
\pgfpathmoveto{\pgfqpoint{1.128763in}{1.959930in}}%
\pgfpathlineto{\pgfqpoint{1.126427in}{1.948422in}}%
\pgfpathlineto{\pgfqpoint{1.131774in}{1.943515in}}%
\pgfpathlineto{\pgfqpoint{1.135886in}{1.943764in}}%
\pgfpathlineto{\pgfqpoint{1.137082in}{1.935988in}}%
\pgfpathlineto{\pgfqpoint{1.133161in}{1.915579in}}%
\pgfpathlineto{\pgfqpoint{1.139942in}{1.914263in}}%
\pgfpathlineto{\pgfqpoint{1.138589in}{1.907407in}}%
\pgfpathlineto{\pgfqpoint{1.146489in}{1.905873in}}%
\pgfpathlineto{\pgfqpoint{1.145219in}{1.892062in}}%
\pgfpathlineto{\pgfqpoint{1.150185in}{1.887641in}}%
\pgfpathlineto{\pgfqpoint{1.165886in}{1.884470in}}%
\pgfpathlineto{\pgfqpoint{1.164838in}{1.878825in}}%
\pgfpathlineto{\pgfqpoint{1.181983in}{1.875682in}}%
\pgfpathlineto{\pgfqpoint{1.177604in}{1.866287in}}%
\pgfpathlineto{\pgfqpoint{1.166980in}{1.866417in}}%
\pgfpathlineto{\pgfqpoint{1.159714in}{1.863660in}}%
\pgfpathlineto{\pgfqpoint{1.157878in}{1.868529in}}%
\pgfpathlineto{\pgfqpoint{1.144630in}{1.867633in}}%
\pgfpathlineto{\pgfqpoint{1.134588in}{1.873151in}}%
\pgfpathlineto{\pgfqpoint{1.130114in}{1.870874in}}%
\pgfpathlineto{\pgfqpoint{1.127333in}{1.864844in}}%
\pgfpathlineto{\pgfqpoint{1.123023in}{1.868301in}}%
\pgfpathlineto{\pgfqpoint{1.105971in}{1.872044in}}%
\pgfpathlineto{\pgfqpoint{1.102494in}{1.863989in}}%
\pgfpathlineto{\pgfqpoint{1.100477in}{1.862882in}}%
\pgfpathlineto{\pgfqpoint{1.093092in}{1.871853in}}%
\pgfpathlineto{\pgfqpoint{1.094071in}{1.877735in}}%
\pgfpathlineto{\pgfqpoint{1.091986in}{1.894613in}}%
\pgfpathlineto{\pgfqpoint{1.088102in}{1.901286in}}%
\pgfpathlineto{\pgfqpoint{1.079391in}{1.904532in}}%
\pgfpathlineto{\pgfqpoint{1.074674in}{1.909170in}}%
\pgfpathlineto{\pgfqpoint{1.077934in}{1.924256in}}%
\pgfpathlineto{\pgfqpoint{1.073148in}{1.929599in}}%
\pgfpathlineto{\pgfqpoint{1.066629in}{1.948325in}}%
\pgfpathlineto{\pgfqpoint{1.067866in}{1.953218in}}%
\pgfpathlineto{\pgfqpoint{1.064456in}{1.964544in}}%
\pgfpathlineto{\pgfqpoint{1.067306in}{1.970565in}}%
\pgfpathlineto{\pgfqpoint{1.061696in}{1.979359in}}%
\pgfpathlineto{\pgfqpoint{1.069612in}{1.982716in}}%
\pgfpathlineto{\pgfqpoint{1.076963in}{1.988148in}}%
\pgfpathlineto{\pgfqpoint{1.086269in}{1.988945in}}%
\pgfpathlineto{\pgfqpoint{1.088551in}{1.993499in}}%
\pgfpathlineto{\pgfqpoint{1.092311in}{1.978823in}}%
\pgfpathlineto{\pgfqpoint{1.103252in}{1.984857in}}%
\pgfpathlineto{\pgfqpoint{1.110227in}{1.984421in}}%
\pgfpathlineto{\pgfqpoint{1.113810in}{1.978599in}}%
\pgfpathlineto{\pgfqpoint{1.126712in}{1.970229in}}%
\pgfpathlineto{\pgfqpoint{1.128763in}{1.959930in}}%
\pgfpathclose%
\pgfusepath{fill}%
\end{pgfscope}%
\begin{pgfscope}%
\pgfpathrectangle{\pgfqpoint{0.100000in}{0.100000in}}{\pgfqpoint{3.608454in}{2.310000in}}%
\pgfusepath{clip}%
\pgfsetbuttcap%
\pgfsetmiterjoin%
\definecolor{currentfill}{rgb}{0.000000,0.607843,0.696078}%
\pgfsetfillcolor{currentfill}%
\pgfsetlinewidth{0.000000pt}%
\definecolor{currentstroke}{rgb}{0.000000,0.000000,0.000000}%
\pgfsetstrokecolor{currentstroke}%
\pgfsetstrokeopacity{0.000000}%
\pgfsetdash{}{0pt}%
\pgfpathmoveto{\pgfqpoint{0.869214in}{2.048910in}}%
\pgfpathlineto{\pgfqpoint{0.906027in}{2.039422in}}%
\pgfpathlineto{\pgfqpoint{0.909977in}{2.024514in}}%
\pgfpathlineto{\pgfqpoint{0.914485in}{2.022985in}}%
\pgfpathlineto{\pgfqpoint{0.921587in}{2.014227in}}%
\pgfpathlineto{\pgfqpoint{0.923281in}{2.004122in}}%
\pgfpathlineto{\pgfqpoint{0.916109in}{1.996717in}}%
\pgfpathlineto{\pgfqpoint{0.904646in}{1.980156in}}%
\pgfpathlineto{\pgfqpoint{0.899963in}{1.971006in}}%
\pgfpathlineto{\pgfqpoint{0.895815in}{1.966858in}}%
\pgfpathlineto{\pgfqpoint{0.869552in}{1.973496in}}%
\pgfpathlineto{\pgfqpoint{0.871160in}{1.980064in}}%
\pgfpathlineto{\pgfqpoint{0.859744in}{1.982888in}}%
\pgfpathlineto{\pgfqpoint{0.859301in}{1.991402in}}%
\pgfpathlineto{\pgfqpoint{0.854053in}{1.998430in}}%
\pgfpathlineto{\pgfqpoint{0.853023in}{2.013264in}}%
\pgfpathlineto{\pgfqpoint{0.855049in}{2.021169in}}%
\pgfpathlineto{\pgfqpoint{0.853215in}{2.027577in}}%
\pgfpathlineto{\pgfqpoint{0.858767in}{2.040242in}}%
\pgfpathlineto{\pgfqpoint{0.846645in}{2.043323in}}%
\pgfpathlineto{\pgfqpoint{0.849176in}{2.054195in}}%
\pgfpathlineto{\pgfqpoint{0.869214in}{2.048910in}}%
\pgfpathclose%
\pgfusepath{fill}%
\end{pgfscope}%
\begin{pgfscope}%
\pgfpathrectangle{\pgfqpoint{0.100000in}{0.100000in}}{\pgfqpoint{3.608454in}{2.310000in}}%
\pgfusepath{clip}%
\pgfsetbuttcap%
\pgfsetmiterjoin%
\definecolor{currentfill}{rgb}{0.000000,0.364706,0.817647}%
\pgfsetfillcolor{currentfill}%
\pgfsetlinewidth{0.000000pt}%
\definecolor{currentstroke}{rgb}{0.000000,0.000000,0.000000}%
\pgfsetstrokecolor{currentstroke}%
\pgfsetstrokeopacity{0.000000}%
\pgfsetdash{}{0pt}%
\pgfpathmoveto{\pgfqpoint{1.707133in}{1.330203in}}%
\pgfpathlineto{\pgfqpoint{1.672634in}{1.332579in}}%
\pgfpathlineto{\pgfqpoint{1.677687in}{1.401208in}}%
\pgfpathlineto{\pgfqpoint{1.716323in}{1.398352in}}%
\pgfpathlineto{\pgfqpoint{1.717833in}{1.398242in}}%
\pgfpathlineto{\pgfqpoint{1.715265in}{1.364101in}}%
\pgfpathlineto{\pgfqpoint{1.709976in}{1.364447in}}%
\pgfpathlineto{\pgfqpoint{1.707133in}{1.330203in}}%
\pgfpathclose%
\pgfusepath{fill}%
\end{pgfscope}%
\begin{pgfscope}%
\pgfpathrectangle{\pgfqpoint{0.100000in}{0.100000in}}{\pgfqpoint{3.608454in}{2.310000in}}%
\pgfusepath{clip}%
\pgfsetbuttcap%
\pgfsetmiterjoin%
\definecolor{currentfill}{rgb}{0.000000,0.243137,0.878431}%
\pgfsetfillcolor{currentfill}%
\pgfsetlinewidth{0.000000pt}%
\definecolor{currentstroke}{rgb}{0.000000,0.000000,0.000000}%
\pgfsetstrokecolor{currentstroke}%
\pgfsetstrokeopacity{0.000000}%
\pgfsetdash{}{0pt}%
\pgfpathmoveto{\pgfqpoint{2.034871in}{0.852102in}}%
\pgfpathlineto{\pgfqpoint{2.034479in}{0.813847in}}%
\pgfpathlineto{\pgfqpoint{2.013777in}{0.811827in}}%
\pgfpathlineto{\pgfqpoint{1.978879in}{0.792186in}}%
\pgfpathlineto{\pgfqpoint{1.975845in}{0.790506in}}%
\pgfpathlineto{\pgfqpoint{1.966416in}{0.807704in}}%
\pgfpathlineto{\pgfqpoint{1.966891in}{0.830295in}}%
\pgfpathlineto{\pgfqpoint{1.970138in}{0.830158in}}%
\pgfpathlineto{\pgfqpoint{1.971315in}{0.864806in}}%
\pgfpathlineto{\pgfqpoint{1.946927in}{0.865843in}}%
\pgfpathlineto{\pgfqpoint{1.948881in}{0.900975in}}%
\pgfpathlineto{\pgfqpoint{1.941990in}{0.901235in}}%
\pgfpathlineto{\pgfqpoint{1.943227in}{0.939444in}}%
\pgfpathlineto{\pgfqpoint{1.945662in}{0.932105in}}%
\pgfpathlineto{\pgfqpoint{1.950349in}{0.931638in}}%
\pgfpathlineto{\pgfqpoint{1.958950in}{0.938059in}}%
\pgfpathlineto{\pgfqpoint{1.963046in}{0.937524in}}%
\pgfpathlineto{\pgfqpoint{1.961356in}{0.929934in}}%
\pgfpathlineto{\pgfqpoint{1.964587in}{0.923699in}}%
\pgfpathlineto{\pgfqpoint{1.968361in}{0.923919in}}%
\pgfpathlineto{\pgfqpoint{1.978804in}{0.941263in}}%
\pgfpathlineto{\pgfqpoint{1.977949in}{0.898974in}}%
\pgfpathlineto{\pgfqpoint{1.984737in}{0.897974in}}%
\pgfpathlineto{\pgfqpoint{2.014996in}{0.896820in}}%
\pgfpathlineto{\pgfqpoint{2.020851in}{0.893088in}}%
\pgfpathlineto{\pgfqpoint{2.020147in}{0.852551in}}%
\pgfpathlineto{\pgfqpoint{2.034871in}{0.852102in}}%
\pgfpathclose%
\pgfusepath{fill}%
\end{pgfscope}%
\begin{pgfscope}%
\pgfpathrectangle{\pgfqpoint{0.100000in}{0.100000in}}{\pgfqpoint{3.608454in}{2.310000in}}%
\pgfusepath{clip}%
\pgfsetbuttcap%
\pgfsetmiterjoin%
\definecolor{currentfill}{rgb}{0.000000,0.545098,0.727451}%
\pgfsetfillcolor{currentfill}%
\pgfsetlinewidth{0.000000pt}%
\definecolor{currentstroke}{rgb}{0.000000,0.000000,0.000000}%
\pgfsetstrokecolor{currentstroke}%
\pgfsetstrokeopacity{0.000000}%
\pgfsetdash{}{0pt}%
\pgfpathmoveto{\pgfqpoint{1.978804in}{0.941263in}}%
\pgfpathlineto{\pgfqpoint{1.968361in}{0.923919in}}%
\pgfpathlineto{\pgfqpoint{1.964587in}{0.923699in}}%
\pgfpathlineto{\pgfqpoint{1.961356in}{0.929934in}}%
\pgfpathlineto{\pgfqpoint{1.963046in}{0.937524in}}%
\pgfpathlineto{\pgfqpoint{1.958950in}{0.938059in}}%
\pgfpathlineto{\pgfqpoint{1.950349in}{0.931638in}}%
\pgfpathlineto{\pgfqpoint{1.945662in}{0.932105in}}%
\pgfpathlineto{\pgfqpoint{1.943227in}{0.939444in}}%
\pgfpathlineto{\pgfqpoint{1.938143in}{0.938336in}}%
\pgfpathlineto{\pgfqpoint{1.939766in}{1.000299in}}%
\pgfpathlineto{\pgfqpoint{1.932881in}{1.000512in}}%
\pgfpathlineto{\pgfqpoint{1.933314in}{1.014281in}}%
\pgfpathlineto{\pgfqpoint{1.981120in}{1.013018in}}%
\pgfpathlineto{\pgfqpoint{1.980692in}{0.992304in}}%
\pgfpathlineto{\pgfqpoint{1.987552in}{0.992168in}}%
\pgfpathlineto{\pgfqpoint{1.987413in}{0.985257in}}%
\pgfpathlineto{\pgfqpoint{2.007887in}{0.984833in}}%
\pgfpathlineto{\pgfqpoint{2.007778in}{0.978034in}}%
\pgfpathlineto{\pgfqpoint{2.014740in}{0.977835in}}%
\pgfpathlineto{\pgfqpoint{2.014398in}{0.957120in}}%
\pgfpathlineto{\pgfqpoint{2.000025in}{0.945069in}}%
\pgfpathlineto{\pgfqpoint{1.993908in}{0.931626in}}%
\pgfpathlineto{\pgfqpoint{1.982724in}{0.934735in}}%
\pgfpathlineto{\pgfqpoint{1.978804in}{0.941263in}}%
\pgfpathclose%
\pgfusepath{fill}%
\end{pgfscope}%
\begin{pgfscope}%
\pgfpathrectangle{\pgfqpoint{0.100000in}{0.100000in}}{\pgfqpoint{3.608454in}{2.310000in}}%
\pgfusepath{clip}%
\pgfsetbuttcap%
\pgfsetmiterjoin%
\definecolor{currentfill}{rgb}{0.000000,0.674510,0.662745}%
\pgfsetfillcolor{currentfill}%
\pgfsetlinewidth{0.000000pt}%
\definecolor{currentstroke}{rgb}{0.000000,0.000000,0.000000}%
\pgfsetstrokecolor{currentstroke}%
\pgfsetstrokeopacity{0.000000}%
\pgfsetdash{}{0pt}%
\pgfpathmoveto{\pgfqpoint{1.355496in}{0.949729in}}%
\pgfpathlineto{\pgfqpoint{1.356403in}{0.942467in}}%
\pgfpathlineto{\pgfqpoint{1.353414in}{0.915777in}}%
\pgfpathlineto{\pgfqpoint{1.314868in}{0.902303in}}%
\pgfpathlineto{\pgfqpoint{1.314425in}{0.898871in}}%
\pgfpathlineto{\pgfqpoint{1.287163in}{0.902597in}}%
\pgfpathlineto{\pgfqpoint{1.285282in}{0.888829in}}%
\pgfpathlineto{\pgfqpoint{1.264784in}{0.891644in}}%
\pgfpathlineto{\pgfqpoint{1.256509in}{0.894758in}}%
\pgfpathlineto{\pgfqpoint{1.259177in}{0.906354in}}%
\pgfpathlineto{\pgfqpoint{1.256598in}{0.912844in}}%
\pgfpathlineto{\pgfqpoint{1.257767in}{0.919573in}}%
\pgfpathlineto{\pgfqpoint{1.253634in}{0.922827in}}%
\pgfpathlineto{\pgfqpoint{1.255093in}{0.930635in}}%
\pgfpathlineto{\pgfqpoint{1.253754in}{0.941144in}}%
\pgfpathlineto{\pgfqpoint{1.245517in}{0.942241in}}%
\pgfpathlineto{\pgfqpoint{1.248648in}{0.963976in}}%
\pgfpathlineto{\pgfqpoint{1.252648in}{0.963109in}}%
\pgfpathlineto{\pgfqpoint{1.318032in}{0.954222in}}%
\pgfpathlineto{\pgfqpoint{1.355496in}{0.949729in}}%
\pgfpathclose%
\pgfusepath{fill}%
\end{pgfscope}%
\begin{pgfscope}%
\pgfpathrectangle{\pgfqpoint{0.100000in}{0.100000in}}{\pgfqpoint{3.608454in}{2.310000in}}%
\pgfusepath{clip}%
\pgfsetbuttcap%
\pgfsetmiterjoin%
\definecolor{currentfill}{rgb}{0.000000,0.329412,0.835294}%
\pgfsetfillcolor{currentfill}%
\pgfsetlinewidth{0.000000pt}%
\definecolor{currentstroke}{rgb}{0.000000,0.000000,0.000000}%
\pgfsetstrokecolor{currentstroke}%
\pgfsetstrokeopacity{0.000000}%
\pgfsetdash{}{0pt}%
\pgfpathmoveto{\pgfqpoint{2.042559in}{1.951936in}}%
\pgfpathlineto{\pgfqpoint{2.042594in}{1.944974in}}%
\pgfpathlineto{\pgfqpoint{2.036642in}{1.945064in}}%
\pgfpathlineto{\pgfqpoint{2.036301in}{1.917281in}}%
\pgfpathlineto{\pgfqpoint{2.036803in}{1.896501in}}%
\pgfpathlineto{\pgfqpoint{2.019610in}{1.896790in}}%
\pgfpathlineto{\pgfqpoint{2.020229in}{1.889881in}}%
\pgfpathlineto{\pgfqpoint{1.983409in}{1.890641in}}%
\pgfpathlineto{\pgfqpoint{1.981775in}{1.895615in}}%
\pgfpathlineto{\pgfqpoint{1.981211in}{1.918353in}}%
\pgfpathlineto{\pgfqpoint{1.981936in}{1.946008in}}%
\pgfpathlineto{\pgfqpoint{1.959926in}{1.946668in}}%
\pgfpathlineto{\pgfqpoint{1.960762in}{1.974477in}}%
\pgfpathlineto{\pgfqpoint{1.960057in}{1.995378in}}%
\pgfpathlineto{\pgfqpoint{2.007393in}{1.994001in}}%
\pgfpathlineto{\pgfqpoint{2.006954in}{2.014902in}}%
\pgfpathlineto{\pgfqpoint{2.049145in}{2.014243in}}%
\pgfpathlineto{\pgfqpoint{2.048844in}{1.986478in}}%
\pgfpathlineto{\pgfqpoint{2.041913in}{1.986517in}}%
\pgfpathlineto{\pgfqpoint{2.042858in}{1.972645in}}%
\pgfpathlineto{\pgfqpoint{2.042559in}{1.951936in}}%
\pgfpathclose%
\pgfusepath{fill}%
\end{pgfscope}%
\begin{pgfscope}%
\pgfpathrectangle{\pgfqpoint{0.100000in}{0.100000in}}{\pgfqpoint{3.608454in}{2.310000in}}%
\pgfusepath{clip}%
\pgfsetbuttcap%
\pgfsetmiterjoin%
\definecolor{currentfill}{rgb}{0.000000,0.854902,0.572549}%
\pgfsetfillcolor{currentfill}%
\pgfsetlinewidth{0.000000pt}%
\definecolor{currentstroke}{rgb}{0.000000,0.000000,0.000000}%
\pgfsetstrokecolor{currentstroke}%
\pgfsetstrokeopacity{0.000000}%
\pgfsetdash{}{0pt}%
\pgfpathmoveto{\pgfqpoint{2.738013in}{1.786827in}}%
\pgfpathlineto{\pgfqpoint{2.697115in}{1.781860in}}%
\pgfpathlineto{\pgfqpoint{2.683319in}{1.780357in}}%
\pgfpathlineto{\pgfqpoint{2.680280in}{1.807943in}}%
\pgfpathlineto{\pgfqpoint{2.677389in}{1.835395in}}%
\pgfpathlineto{\pgfqpoint{2.674595in}{1.855705in}}%
\pgfpathlineto{\pgfqpoint{2.681411in}{1.856342in}}%
\pgfpathlineto{\pgfqpoint{2.680662in}{1.863242in}}%
\pgfpathlineto{\pgfqpoint{2.707810in}{1.866142in}}%
\pgfpathlineto{\pgfqpoint{2.728097in}{1.868854in}}%
\pgfpathlineto{\pgfqpoint{2.729298in}{1.855252in}}%
\pgfpathlineto{\pgfqpoint{2.738013in}{1.786827in}}%
\pgfpathclose%
\pgfusepath{fill}%
\end{pgfscope}%
\begin{pgfscope}%
\pgfpathrectangle{\pgfqpoint{0.100000in}{0.100000in}}{\pgfqpoint{3.608454in}{2.310000in}}%
\pgfusepath{clip}%
\pgfsetbuttcap%
\pgfsetmiterjoin%
\definecolor{currentfill}{rgb}{0.000000,0.815686,0.592157}%
\pgfsetfillcolor{currentfill}%
\pgfsetlinewidth{0.000000pt}%
\definecolor{currentstroke}{rgb}{0.000000,0.000000,0.000000}%
\pgfsetstrokecolor{currentstroke}%
\pgfsetstrokeopacity{0.000000}%
\pgfsetdash{}{0pt}%
\pgfpathmoveto{\pgfqpoint{3.069493in}{0.458265in}}%
\pgfpathlineto{\pgfqpoint{3.072400in}{0.437158in}}%
\pgfpathlineto{\pgfqpoint{3.037010in}{0.432074in}}%
\pgfpathlineto{\pgfqpoint{3.036461in}{0.440065in}}%
\pgfpathlineto{\pgfqpoint{3.030688in}{0.445529in}}%
\pgfpathlineto{\pgfqpoint{3.027406in}{0.441545in}}%
\pgfpathlineto{\pgfqpoint{3.030649in}{0.432623in}}%
\pgfpathlineto{\pgfqpoint{3.024312in}{0.431402in}}%
\pgfpathlineto{\pgfqpoint{3.021799in}{0.431432in}}%
\pgfpathlineto{\pgfqpoint{3.006353in}{0.451398in}}%
\pgfpathlineto{\pgfqpoint{2.995081in}{0.469136in}}%
\pgfpathlineto{\pgfqpoint{2.984705in}{0.479666in}}%
\pgfpathlineto{\pgfqpoint{2.986857in}{0.487228in}}%
\pgfpathlineto{\pgfqpoint{2.992126in}{0.496993in}}%
\pgfpathlineto{\pgfqpoint{3.027503in}{0.502091in}}%
\pgfpathlineto{\pgfqpoint{3.030929in}{0.477492in}}%
\pgfpathlineto{\pgfqpoint{3.065776in}{0.482786in}}%
\pgfpathlineto{\pgfqpoint{3.069493in}{0.458265in}}%
\pgfpathclose%
\pgfusepath{fill}%
\end{pgfscope}%
\begin{pgfscope}%
\pgfpathrectangle{\pgfqpoint{0.100000in}{0.100000in}}{\pgfqpoint{3.608454in}{2.310000in}}%
\pgfusepath{clip}%
\pgfsetbuttcap%
\pgfsetmiterjoin%
\definecolor{currentfill}{rgb}{0.000000,0.662745,0.668627}%
\pgfsetfillcolor{currentfill}%
\pgfsetlinewidth{0.000000pt}%
\definecolor{currentstroke}{rgb}{0.000000,0.000000,0.000000}%
\pgfsetstrokecolor{currentstroke}%
\pgfsetstrokeopacity{0.000000}%
\pgfsetdash{}{0pt}%
\pgfpathmoveto{\pgfqpoint{3.046524in}{1.607153in}}%
\pgfpathlineto{\pgfqpoint{3.045476in}{1.613591in}}%
\pgfpathlineto{\pgfqpoint{3.051910in}{1.617021in}}%
\pgfpathlineto{\pgfqpoint{3.050721in}{1.630639in}}%
\pgfpathlineto{\pgfqpoint{3.078140in}{1.636089in}}%
\pgfpathlineto{\pgfqpoint{3.094591in}{1.638194in}}%
\pgfpathlineto{\pgfqpoint{3.105576in}{1.628770in}}%
\pgfpathlineto{\pgfqpoint{3.109153in}{1.629395in}}%
\pgfpathlineto{\pgfqpoint{3.110757in}{1.621055in}}%
\pgfpathlineto{\pgfqpoint{3.106858in}{1.608202in}}%
\pgfpathlineto{\pgfqpoint{3.068655in}{1.604850in}}%
\pgfpathlineto{\pgfqpoint{3.054170in}{1.610224in}}%
\pgfpathlineto{\pgfqpoint{3.046524in}{1.607153in}}%
\pgfpathclose%
\pgfusepath{fill}%
\end{pgfscope}%
\begin{pgfscope}%
\pgfpathrectangle{\pgfqpoint{0.100000in}{0.100000in}}{\pgfqpoint{3.608454in}{2.310000in}}%
\pgfusepath{clip}%
\pgfsetbuttcap%
\pgfsetmiterjoin%
\definecolor{currentfill}{rgb}{0.000000,0.603922,0.698039}%
\pgfsetfillcolor{currentfill}%
\pgfsetlinewidth{0.000000pt}%
\definecolor{currentstroke}{rgb}{0.000000,0.000000,0.000000}%
\pgfsetstrokecolor{currentstroke}%
\pgfsetstrokeopacity{0.000000}%
\pgfsetdash{}{0pt}%
\pgfpathmoveto{\pgfqpoint{2.583434in}{1.352723in}}%
\pgfpathlineto{\pgfqpoint{2.566176in}{1.350839in}}%
\pgfpathlineto{\pgfqpoint{2.562407in}{1.357174in}}%
\pgfpathlineto{\pgfqpoint{2.562133in}{1.362791in}}%
\pgfpathlineto{\pgfqpoint{2.558058in}{1.370400in}}%
\pgfpathlineto{\pgfqpoint{2.560407in}{1.373999in}}%
\pgfpathlineto{\pgfqpoint{2.558464in}{1.382618in}}%
\pgfpathlineto{\pgfqpoint{2.562752in}{1.385889in}}%
\pgfpathlineto{\pgfqpoint{2.559183in}{1.428045in}}%
\pgfpathlineto{\pgfqpoint{2.557572in}{1.448924in}}%
\pgfpathlineto{\pgfqpoint{2.565221in}{1.447932in}}%
\pgfpathlineto{\pgfqpoint{2.565557in}{1.434073in}}%
\pgfpathlineto{\pgfqpoint{2.585430in}{1.435815in}}%
\pgfpathlineto{\pgfqpoint{2.583559in}{1.456426in}}%
\pgfpathlineto{\pgfqpoint{2.607399in}{1.458524in}}%
\pgfpathlineto{\pgfqpoint{2.607655in}{1.455712in}}%
\pgfpathlineto{\pgfqpoint{2.612117in}{1.412594in}}%
\pgfpathlineto{\pgfqpoint{2.614715in}{1.406613in}}%
\pgfpathlineto{\pgfqpoint{2.613308in}{1.399991in}}%
\pgfpathlineto{\pgfqpoint{2.597905in}{1.398852in}}%
\pgfpathlineto{\pgfqpoint{2.598672in}{1.388501in}}%
\pgfpathlineto{\pgfqpoint{2.591897in}{1.387949in}}%
\pgfpathlineto{\pgfqpoint{2.593047in}{1.374201in}}%
\pgfpathlineto{\pgfqpoint{2.581657in}{1.373507in}}%
\pgfpathlineto{\pgfqpoint{2.583434in}{1.352723in}}%
\pgfpathclose%
\pgfusepath{fill}%
\end{pgfscope}%
\begin{pgfscope}%
\pgfpathrectangle{\pgfqpoint{0.100000in}{0.100000in}}{\pgfqpoint{3.608454in}{2.310000in}}%
\pgfusepath{clip}%
\pgfsetbuttcap%
\pgfsetmiterjoin%
\definecolor{currentfill}{rgb}{0.000000,0.341176,0.829412}%
\pgfsetfillcolor{currentfill}%
\pgfsetlinewidth{0.000000pt}%
\definecolor{currentstroke}{rgb}{0.000000,0.000000,0.000000}%
\pgfsetstrokecolor{currentstroke}%
\pgfsetstrokeopacity{0.000000}%
\pgfsetdash{}{0pt}%
\pgfpathmoveto{\pgfqpoint{1.580064in}{1.651542in}}%
\pgfpathlineto{\pgfqpoint{1.571835in}{1.569055in}}%
\pgfpathlineto{\pgfqpoint{1.536486in}{1.572765in}}%
\pgfpathlineto{\pgfqpoint{1.537205in}{1.579689in}}%
\pgfpathlineto{\pgfqpoint{1.500517in}{1.584019in}}%
\pgfpathlineto{\pgfqpoint{1.503839in}{1.609878in}}%
\pgfpathlineto{\pgfqpoint{1.507108in}{1.644984in}}%
\pgfpathlineto{\pgfqpoint{1.508517in}{1.658724in}}%
\pgfpathlineto{\pgfqpoint{1.531372in}{1.656425in}}%
\pgfpathlineto{\pgfqpoint{1.580064in}{1.651542in}}%
\pgfpathclose%
\pgfusepath{fill}%
\end{pgfscope}%
\begin{pgfscope}%
\pgfpathrectangle{\pgfqpoint{0.100000in}{0.100000in}}{\pgfqpoint{3.608454in}{2.310000in}}%
\pgfusepath{clip}%
\pgfsetbuttcap%
\pgfsetmiterjoin%
\definecolor{currentfill}{rgb}{0.000000,0.615686,0.692157}%
\pgfsetfillcolor{currentfill}%
\pgfsetlinewidth{0.000000pt}%
\definecolor{currentstroke}{rgb}{0.000000,0.000000,0.000000}%
\pgfsetstrokecolor{currentstroke}%
\pgfsetstrokeopacity{0.000000}%
\pgfsetdash{}{0pt}%
\pgfpathmoveto{\pgfqpoint{1.651469in}{1.020799in}}%
\pgfpathlineto{\pgfqpoint{1.685686in}{1.018201in}}%
\pgfpathlineto{\pgfqpoint{1.683230in}{0.983785in}}%
\pgfpathlineto{\pgfqpoint{1.711338in}{0.981859in}}%
\pgfpathlineto{\pgfqpoint{1.708835in}{0.943953in}}%
\pgfpathlineto{\pgfqpoint{1.640100in}{0.948213in}}%
\pgfpathlineto{\pgfqpoint{1.642812in}{0.986873in}}%
\pgfpathlineto{\pgfqpoint{1.648845in}{0.986397in}}%
\pgfpathlineto{\pgfqpoint{1.651469in}{1.020799in}}%
\pgfpathclose%
\pgfusepath{fill}%
\end{pgfscope}%
\begin{pgfscope}%
\pgfpathrectangle{\pgfqpoint{0.100000in}{0.100000in}}{\pgfqpoint{3.608454in}{2.310000in}}%
\pgfusepath{clip}%
\pgfsetbuttcap%
\pgfsetmiterjoin%
\definecolor{currentfill}{rgb}{0.000000,0.156863,0.921569}%
\pgfsetfillcolor{currentfill}%
\pgfsetlinewidth{0.000000pt}%
\definecolor{currentstroke}{rgb}{0.000000,0.000000,0.000000}%
\pgfsetstrokecolor{currentstroke}%
\pgfsetstrokeopacity{0.000000}%
\pgfsetdash{}{0pt}%
\pgfpathmoveto{\pgfqpoint{1.396241in}{1.384895in}}%
\pgfpathlineto{\pgfqpoint{1.393282in}{1.388976in}}%
\pgfpathlineto{\pgfqpoint{1.387549in}{1.389567in}}%
\pgfpathlineto{\pgfqpoint{1.383841in}{1.385038in}}%
\pgfpathlineto{\pgfqpoint{1.377365in}{1.386918in}}%
\pgfpathlineto{\pgfqpoint{1.369057in}{1.398085in}}%
\pgfpathlineto{\pgfqpoint{1.355918in}{1.399849in}}%
\pgfpathlineto{\pgfqpoint{1.352191in}{1.404691in}}%
\pgfpathlineto{\pgfqpoint{1.350593in}{1.411552in}}%
\pgfpathlineto{\pgfqpoint{1.346932in}{1.416687in}}%
\pgfpathlineto{\pgfqpoint{1.349566in}{1.420460in}}%
\pgfpathlineto{\pgfqpoint{1.290367in}{1.428918in}}%
\pgfpathlineto{\pgfqpoint{1.251645in}{1.434868in}}%
\pgfpathlineto{\pgfqpoint{1.253203in}{1.445105in}}%
\pgfpathlineto{\pgfqpoint{1.255196in}{1.457818in}}%
\pgfpathlineto{\pgfqpoint{1.284941in}{1.452441in}}%
\pgfpathlineto{\pgfqpoint{1.285446in}{1.455840in}}%
\pgfpathlineto{\pgfqpoint{1.322533in}{1.450461in}}%
\pgfpathlineto{\pgfqpoint{1.324000in}{1.460674in}}%
\pgfpathlineto{\pgfqpoint{1.354499in}{1.456659in}}%
\pgfpathlineto{\pgfqpoint{1.355468in}{1.463655in}}%
\pgfpathlineto{\pgfqpoint{1.362220in}{1.462510in}}%
\pgfpathlineto{\pgfqpoint{1.364204in}{1.476264in}}%
\pgfpathlineto{\pgfqpoint{1.380971in}{1.474066in}}%
\pgfpathlineto{\pgfqpoint{1.379506in}{1.460495in}}%
\pgfpathlineto{\pgfqpoint{1.403922in}{1.457766in}}%
\pgfpathlineto{\pgfqpoint{1.426636in}{1.453930in}}%
\pgfpathlineto{\pgfqpoint{1.435412in}{1.444400in}}%
\pgfpathlineto{\pgfqpoint{1.437752in}{1.434786in}}%
\pgfpathlineto{\pgfqpoint{1.444005in}{1.434780in}}%
\pgfpathlineto{\pgfqpoint{1.451897in}{1.428337in}}%
\pgfpathlineto{\pgfqpoint{1.439656in}{1.412644in}}%
\pgfpathlineto{\pgfqpoint{1.424568in}{1.406130in}}%
\pgfpathlineto{\pgfqpoint{1.423585in}{1.392666in}}%
\pgfpathlineto{\pgfqpoint{1.421683in}{1.386315in}}%
\pgfpathlineto{\pgfqpoint{1.398138in}{1.389426in}}%
\pgfpathlineto{\pgfqpoint{1.396241in}{1.384895in}}%
\pgfpathclose%
\pgfusepath{fill}%
\end{pgfscope}%
\begin{pgfscope}%
\pgfpathrectangle{\pgfqpoint{0.100000in}{0.100000in}}{\pgfqpoint{3.608454in}{2.310000in}}%
\pgfusepath{clip}%
\pgfsetbuttcap%
\pgfsetmiterjoin%
\definecolor{currentfill}{rgb}{0.000000,0.603922,0.698039}%
\pgfsetfillcolor{currentfill}%
\pgfsetlinewidth{0.000000pt}%
\definecolor{currentstroke}{rgb}{0.000000,0.000000,0.000000}%
\pgfsetstrokecolor{currentstroke}%
\pgfsetstrokeopacity{0.000000}%
\pgfsetdash{}{0pt}%
\pgfpathmoveto{\pgfqpoint{3.091930in}{0.511405in}}%
\pgfpathlineto{\pgfqpoint{3.088462in}{0.501164in}}%
\pgfpathlineto{\pgfqpoint{3.097174in}{0.491133in}}%
\pgfpathlineto{\pgfqpoint{3.100740in}{0.491314in}}%
\pgfpathlineto{\pgfqpoint{3.109319in}{0.479362in}}%
\pgfpathlineto{\pgfqpoint{3.095343in}{0.476565in}}%
\pgfpathlineto{\pgfqpoint{3.096465in}{0.469557in}}%
\pgfpathlineto{\pgfqpoint{3.089414in}{0.468427in}}%
\pgfpathlineto{\pgfqpoint{3.090529in}{0.461360in}}%
\pgfpathlineto{\pgfqpoint{3.069493in}{0.458265in}}%
\pgfpathlineto{\pgfqpoint{3.065776in}{0.482786in}}%
\pgfpathlineto{\pgfqpoint{3.030929in}{0.477492in}}%
\pgfpathlineto{\pgfqpoint{3.027503in}{0.502091in}}%
\pgfpathlineto{\pgfqpoint{3.021367in}{0.543921in}}%
\pgfpathlineto{\pgfqpoint{3.016866in}{0.550384in}}%
\pgfpathlineto{\pgfqpoint{3.019784in}{0.555093in}}%
\pgfpathlineto{\pgfqpoint{3.026291in}{0.558732in}}%
\pgfpathlineto{\pgfqpoint{3.047294in}{0.561961in}}%
\pgfpathlineto{\pgfqpoint{3.048333in}{0.554969in}}%
\pgfpathlineto{\pgfqpoint{3.055318in}{0.556018in}}%
\pgfpathlineto{\pgfqpoint{3.059031in}{0.547094in}}%
\pgfpathlineto{\pgfqpoint{3.070237in}{0.538537in}}%
\pgfpathlineto{\pgfqpoint{3.076539in}{0.531713in}}%
\pgfpathlineto{\pgfqpoint{3.077818in}{0.527142in}}%
\pgfpathlineto{\pgfqpoint{3.084501in}{0.525466in}}%
\pgfpathlineto{\pgfqpoint{3.090421in}{0.517236in}}%
\pgfpathlineto{\pgfqpoint{3.091930in}{0.511405in}}%
\pgfpathclose%
\pgfusepath{fill}%
\end{pgfscope}%
\begin{pgfscope}%
\pgfpathrectangle{\pgfqpoint{0.100000in}{0.100000in}}{\pgfqpoint{3.608454in}{2.310000in}}%
\pgfusepath{clip}%
\pgfsetbuttcap%
\pgfsetmiterjoin%
\definecolor{currentfill}{rgb}{0.000000,0.835294,0.582353}%
\pgfsetfillcolor{currentfill}%
\pgfsetlinewidth{0.000000pt}%
\definecolor{currentstroke}{rgb}{0.000000,0.000000,0.000000}%
\pgfsetstrokecolor{currentstroke}%
\pgfsetstrokeopacity{0.000000}%
\pgfsetdash{}{0pt}%
\pgfpathmoveto{\pgfqpoint{0.509532in}{1.308555in}}%
\pgfpathlineto{\pgfqpoint{0.550101in}{1.297200in}}%
\pgfpathlineto{\pgfqpoint{0.606581in}{1.281871in}}%
\pgfpathlineto{\pgfqpoint{0.645007in}{1.271759in}}%
\pgfpathlineto{\pgfqpoint{0.668493in}{1.266395in}}%
\pgfpathlineto{\pgfqpoint{0.649520in}{1.191532in}}%
\pgfpathlineto{\pgfqpoint{0.598059in}{1.204495in}}%
\pgfpathlineto{\pgfqpoint{0.547967in}{1.217656in}}%
\pgfpathlineto{\pgfqpoint{0.547247in}{1.223234in}}%
\pgfpathlineto{\pgfqpoint{0.537333in}{1.227732in}}%
\pgfpathlineto{\pgfqpoint{0.539132in}{1.241671in}}%
\pgfpathlineto{\pgfqpoint{0.533846in}{1.244020in}}%
\pgfpathlineto{\pgfqpoint{0.536255in}{1.250941in}}%
\pgfpathlineto{\pgfqpoint{0.529070in}{1.252510in}}%
\pgfpathlineto{\pgfqpoint{0.530958in}{1.259199in}}%
\pgfpathlineto{\pgfqpoint{0.522085in}{1.261761in}}%
\pgfpathlineto{\pgfqpoint{0.523955in}{1.268417in}}%
\pgfpathlineto{\pgfqpoint{0.519557in}{1.269679in}}%
\pgfpathlineto{\pgfqpoint{0.521448in}{1.276381in}}%
\pgfpathlineto{\pgfqpoint{0.514797in}{1.280686in}}%
\pgfpathlineto{\pgfqpoint{0.510504in}{1.286661in}}%
\pgfpathlineto{\pgfqpoint{0.512448in}{1.293350in}}%
\pgfpathlineto{\pgfqpoint{0.505758in}{1.295236in}}%
\pgfpathlineto{\pgfqpoint{0.509532in}{1.308555in}}%
\pgfpathclose%
\pgfusepath{fill}%
\end{pgfscope}%
\begin{pgfscope}%
\pgfpathrectangle{\pgfqpoint{0.100000in}{0.100000in}}{\pgfqpoint{3.608454in}{2.310000in}}%
\pgfusepath{clip}%
\pgfsetbuttcap%
\pgfsetmiterjoin%
\definecolor{currentfill}{rgb}{0.000000,0.427451,0.786275}%
\pgfsetfillcolor{currentfill}%
\pgfsetlinewidth{0.000000pt}%
\definecolor{currentstroke}{rgb}{0.000000,0.000000,0.000000}%
\pgfsetstrokecolor{currentstroke}%
\pgfsetstrokeopacity{0.000000}%
\pgfsetdash{}{0pt}%
\pgfpathmoveto{\pgfqpoint{2.193447in}{1.465250in}}%
\pgfpathlineto{\pgfqpoint{2.180394in}{1.464846in}}%
\pgfpathlineto{\pgfqpoint{2.138550in}{1.463912in}}%
\pgfpathlineto{\pgfqpoint{2.138316in}{1.510307in}}%
\pgfpathlineto{\pgfqpoint{2.172375in}{1.510519in}}%
\pgfpathlineto{\pgfqpoint{2.192689in}{1.511198in}}%
\pgfpathlineto{\pgfqpoint{2.193447in}{1.465250in}}%
\pgfpathclose%
\pgfusepath{fill}%
\end{pgfscope}%
\begin{pgfscope}%
\pgfpathrectangle{\pgfqpoint{0.100000in}{0.100000in}}{\pgfqpoint{3.608454in}{2.310000in}}%
\pgfusepath{clip}%
\pgfsetbuttcap%
\pgfsetmiterjoin%
\definecolor{currentfill}{rgb}{0.000000,0.709804,0.645098}%
\pgfsetfillcolor{currentfill}%
\pgfsetlinewidth{0.000000pt}%
\definecolor{currentstroke}{rgb}{0.000000,0.000000,0.000000}%
\pgfsetstrokecolor{currentstroke}%
\pgfsetstrokeopacity{0.000000}%
\pgfsetdash{}{0pt}%
\pgfpathmoveto{\pgfqpoint{2.724860in}{1.203323in}}%
\pgfpathlineto{\pgfqpoint{2.723224in}{1.208766in}}%
\pgfpathlineto{\pgfqpoint{2.709289in}{1.208321in}}%
\pgfpathlineto{\pgfqpoint{2.705489in}{1.213403in}}%
\pgfpathlineto{\pgfqpoint{2.703168in}{1.221287in}}%
\pgfpathlineto{\pgfqpoint{2.696385in}{1.221833in}}%
\pgfpathlineto{\pgfqpoint{2.692552in}{1.225974in}}%
\pgfpathlineto{\pgfqpoint{2.690908in}{1.235308in}}%
\pgfpathlineto{\pgfqpoint{2.692446in}{1.245040in}}%
\pgfpathlineto{\pgfqpoint{2.693473in}{1.249077in}}%
\pgfpathlineto{\pgfqpoint{2.704025in}{1.249697in}}%
\pgfpathlineto{\pgfqpoint{2.711292in}{1.250974in}}%
\pgfpathlineto{\pgfqpoint{2.718150in}{1.247622in}}%
\pgfpathlineto{\pgfqpoint{2.722412in}{1.251628in}}%
\pgfpathlineto{\pgfqpoint{2.730780in}{1.248329in}}%
\pgfpathlineto{\pgfqpoint{2.729948in}{1.258824in}}%
\pgfpathlineto{\pgfqpoint{2.741472in}{1.258776in}}%
\pgfpathlineto{\pgfqpoint{2.742755in}{1.250848in}}%
\pgfpathlineto{\pgfqpoint{2.751350in}{1.246500in}}%
\pgfpathlineto{\pgfqpoint{2.752556in}{1.236924in}}%
\pgfpathlineto{\pgfqpoint{2.742337in}{1.226077in}}%
\pgfpathlineto{\pgfqpoint{2.742647in}{1.220616in}}%
\pgfpathlineto{\pgfqpoint{2.747524in}{1.217198in}}%
\pgfpathlineto{\pgfqpoint{2.731198in}{1.201857in}}%
\pgfpathlineto{\pgfqpoint{2.724860in}{1.203323in}}%
\pgfpathclose%
\pgfusepath{fill}%
\end{pgfscope}%
\begin{pgfscope}%
\pgfpathrectangle{\pgfqpoint{0.100000in}{0.100000in}}{\pgfqpoint{3.608454in}{2.310000in}}%
\pgfusepath{clip}%
\pgfsetbuttcap%
\pgfsetmiterjoin%
\definecolor{currentfill}{rgb}{0.000000,0.505882,0.747059}%
\pgfsetfillcolor{currentfill}%
\pgfsetlinewidth{0.000000pt}%
\definecolor{currentstroke}{rgb}{0.000000,0.000000,0.000000}%
\pgfsetstrokecolor{currentstroke}%
\pgfsetstrokeopacity{0.000000}%
\pgfsetdash{}{0pt}%
\pgfpathmoveto{\pgfqpoint{1.711338in}{0.981859in}}%
\pgfpathlineto{\pgfqpoint{1.683230in}{0.983785in}}%
\pgfpathlineto{\pgfqpoint{1.685686in}{1.018201in}}%
\pgfpathlineto{\pgfqpoint{1.710432in}{1.016551in}}%
\pgfpathlineto{\pgfqpoint{1.713023in}{1.050837in}}%
\pgfpathlineto{\pgfqpoint{1.748391in}{1.048651in}}%
\pgfpathlineto{\pgfqpoint{1.746176in}{1.014219in}}%
\pgfpathlineto{\pgfqpoint{1.754289in}{1.013713in}}%
\pgfpathlineto{\pgfqpoint{1.752117in}{0.979371in}}%
\pgfpathlineto{\pgfqpoint{1.745610in}{0.979809in}}%
\pgfpathlineto{\pgfqpoint{1.711338in}{0.981859in}}%
\pgfpathclose%
\pgfusepath{fill}%
\end{pgfscope}%
\begin{pgfscope}%
\pgfpathrectangle{\pgfqpoint{0.100000in}{0.100000in}}{\pgfqpoint{3.608454in}{2.310000in}}%
\pgfusepath{clip}%
\pgfsetbuttcap%
\pgfsetmiterjoin%
\definecolor{currentfill}{rgb}{0.000000,0.650980,0.674510}%
\pgfsetfillcolor{currentfill}%
\pgfsetlinewidth{0.000000pt}%
\definecolor{currentstroke}{rgb}{0.000000,0.000000,0.000000}%
\pgfsetstrokecolor{currentstroke}%
\pgfsetstrokeopacity{0.000000}%
\pgfsetdash{}{0pt}%
\pgfpathmoveto{\pgfqpoint{1.666412in}{1.249548in}}%
\pgfpathlineto{\pgfqpoint{1.664847in}{1.229439in}}%
\pgfpathlineto{\pgfqpoint{1.662385in}{1.198210in}}%
\pgfpathlineto{\pgfqpoint{1.621102in}{1.201730in}}%
\pgfpathlineto{\pgfqpoint{1.602002in}{1.203777in}}%
\pgfpathlineto{\pgfqpoint{1.596724in}{1.204258in}}%
\pgfpathlineto{\pgfqpoint{1.601877in}{1.254860in}}%
\pgfpathlineto{\pgfqpoint{1.666412in}{1.249548in}}%
\pgfpathclose%
\pgfusepath{fill}%
\end{pgfscope}%
\begin{pgfscope}%
\pgfpathrectangle{\pgfqpoint{0.100000in}{0.100000in}}{\pgfqpoint{3.608454in}{2.310000in}}%
\pgfusepath{clip}%
\pgfsetbuttcap%
\pgfsetmiterjoin%
\definecolor{currentfill}{rgb}{0.000000,0.309804,0.845098}%
\pgfsetfillcolor{currentfill}%
\pgfsetlinewidth{0.000000pt}%
\definecolor{currentstroke}{rgb}{0.000000,0.000000,0.000000}%
\pgfsetstrokecolor{currentstroke}%
\pgfsetstrokeopacity{0.000000}%
\pgfsetdash{}{0pt}%
\pgfpathmoveto{\pgfqpoint{2.685775in}{1.104241in}}%
\pgfpathlineto{\pgfqpoint{2.681536in}{1.102539in}}%
\pgfpathlineto{\pgfqpoint{2.668887in}{1.104876in}}%
\pgfpathlineto{\pgfqpoint{2.654621in}{1.100124in}}%
\pgfpathlineto{\pgfqpoint{2.651600in}{1.103191in}}%
\pgfpathlineto{\pgfqpoint{2.643683in}{1.102257in}}%
\pgfpathlineto{\pgfqpoint{2.632392in}{1.103973in}}%
\pgfpathlineto{\pgfqpoint{2.605693in}{1.111968in}}%
\pgfpathlineto{\pgfqpoint{2.605651in}{1.120589in}}%
\pgfpathlineto{\pgfqpoint{2.606773in}{1.136211in}}%
\pgfpathlineto{\pgfqpoint{2.604989in}{1.140525in}}%
\pgfpathlineto{\pgfqpoint{2.606315in}{1.148332in}}%
\pgfpathlineto{\pgfqpoint{2.597965in}{1.148745in}}%
\pgfpathlineto{\pgfqpoint{2.599908in}{1.153357in}}%
\pgfpathlineto{\pgfqpoint{2.607032in}{1.160645in}}%
\pgfpathlineto{\pgfqpoint{2.606663in}{1.174945in}}%
\pgfpathlineto{\pgfqpoint{2.610086in}{1.175304in}}%
\pgfpathlineto{\pgfqpoint{2.639785in}{1.178722in}}%
\pgfpathlineto{\pgfqpoint{2.656122in}{1.179965in}}%
\pgfpathlineto{\pgfqpoint{2.678723in}{1.180579in}}%
\pgfpathlineto{\pgfqpoint{2.690217in}{1.180528in}}%
\pgfpathlineto{\pgfqpoint{2.690027in}{1.171438in}}%
\pgfpathlineto{\pgfqpoint{2.689646in}{1.154594in}}%
\pgfpathlineto{\pgfqpoint{2.693285in}{1.148452in}}%
\pgfpathlineto{\pgfqpoint{2.693376in}{1.142517in}}%
\pgfpathlineto{\pgfqpoint{2.690045in}{1.139163in}}%
\pgfpathlineto{\pgfqpoint{2.683825in}{1.139915in}}%
\pgfpathlineto{\pgfqpoint{2.678100in}{1.135237in}}%
\pgfpathlineto{\pgfqpoint{2.681516in}{1.127762in}}%
\pgfpathlineto{\pgfqpoint{2.689025in}{1.121466in}}%
\pgfpathlineto{\pgfqpoint{2.685202in}{1.110108in}}%
\pgfpathlineto{\pgfqpoint{2.685775in}{1.104241in}}%
\pgfpathclose%
\pgfusepath{fill}%
\end{pgfscope}%
\begin{pgfscope}%
\pgfpathrectangle{\pgfqpoint{0.100000in}{0.100000in}}{\pgfqpoint{3.608454in}{2.310000in}}%
\pgfusepath{clip}%
\pgfsetbuttcap%
\pgfsetmiterjoin%
\definecolor{currentfill}{rgb}{0.000000,0.752941,0.623529}%
\pgfsetfillcolor{currentfill}%
\pgfsetlinewidth{0.000000pt}%
\definecolor{currentstroke}{rgb}{0.000000,0.000000,0.000000}%
\pgfsetstrokecolor{currentstroke}%
\pgfsetstrokeopacity{0.000000}%
\pgfsetdash{}{0pt}%
\pgfpathmoveto{\pgfqpoint{0.486008in}{1.393700in}}%
\pgfpathlineto{\pgfqpoint{0.474622in}{1.404992in}}%
\pgfpathlineto{\pgfqpoint{0.471004in}{1.413448in}}%
\pgfpathlineto{\pgfqpoint{0.475697in}{1.423778in}}%
\pgfpathlineto{\pgfqpoint{0.476039in}{1.431016in}}%
\pgfpathlineto{\pgfqpoint{0.465842in}{1.434190in}}%
\pgfpathlineto{\pgfqpoint{0.465818in}{1.444533in}}%
\pgfpathlineto{\pgfqpoint{0.470704in}{1.451769in}}%
\pgfpathlineto{\pgfqpoint{0.469193in}{1.459876in}}%
\pgfpathlineto{\pgfqpoint{0.469229in}{1.459834in}}%
\pgfpathlineto{\pgfqpoint{0.489989in}{1.471204in}}%
\pgfpathlineto{\pgfqpoint{0.504262in}{1.472446in}}%
\pgfpathlineto{\pgfqpoint{0.508101in}{1.469498in}}%
\pgfpathlineto{\pgfqpoint{0.515295in}{1.495292in}}%
\pgfpathlineto{\pgfqpoint{0.526181in}{1.471930in}}%
\pgfpathlineto{\pgfqpoint{0.537164in}{1.478720in}}%
\pgfpathlineto{\pgfqpoint{0.550096in}{1.491227in}}%
\pgfpathlineto{\pgfqpoint{0.566537in}{1.504878in}}%
\pgfpathlineto{\pgfqpoint{0.577147in}{1.506768in}}%
\pgfpathlineto{\pgfqpoint{0.583557in}{1.498623in}}%
\pgfpathlineto{\pgfqpoint{0.592836in}{1.501027in}}%
\pgfpathlineto{\pgfqpoint{0.597596in}{1.492295in}}%
\pgfpathlineto{\pgfqpoint{0.596004in}{1.489421in}}%
\pgfpathlineto{\pgfqpoint{0.597769in}{1.478364in}}%
\pgfpathlineto{\pgfqpoint{0.610025in}{1.468870in}}%
\pgfpathlineto{\pgfqpoint{0.609428in}{1.457861in}}%
\pgfpathlineto{\pgfqpoint{0.613111in}{1.453862in}}%
\pgfpathlineto{\pgfqpoint{0.612588in}{1.444548in}}%
\pgfpathlineto{\pgfqpoint{0.607467in}{1.441382in}}%
\pgfpathlineto{\pgfqpoint{0.605864in}{1.444998in}}%
\pgfpathlineto{\pgfqpoint{0.584724in}{1.432975in}}%
\pgfpathlineto{\pgfqpoint{0.583298in}{1.427984in}}%
\pgfpathlineto{\pgfqpoint{0.577622in}{1.423252in}}%
\pgfpathlineto{\pgfqpoint{0.570915in}{1.425063in}}%
\pgfpathlineto{\pgfqpoint{0.548331in}{1.412190in}}%
\pgfpathlineto{\pgfqpoint{0.533257in}{1.413969in}}%
\pgfpathlineto{\pgfqpoint{0.520707in}{1.412942in}}%
\pgfpathlineto{\pgfqpoint{0.486008in}{1.393700in}}%
\pgfpathclose%
\pgfusepath{fill}%
\end{pgfscope}%
\begin{pgfscope}%
\pgfpathrectangle{\pgfqpoint{0.100000in}{0.100000in}}{\pgfqpoint{3.608454in}{2.310000in}}%
\pgfusepath{clip}%
\pgfsetbuttcap%
\pgfsetmiterjoin%
\definecolor{currentfill}{rgb}{0.000000,0.431373,0.784314}%
\pgfsetfillcolor{currentfill}%
\pgfsetlinewidth{0.000000pt}%
\definecolor{currentstroke}{rgb}{0.000000,0.000000,0.000000}%
\pgfsetstrokecolor{currentstroke}%
\pgfsetstrokeopacity{0.000000}%
\pgfsetdash{}{0pt}%
\pgfpathmoveto{\pgfqpoint{1.664670in}{1.909384in}}%
\pgfpathlineto{\pgfqpoint{1.606786in}{1.914838in}}%
\pgfpathlineto{\pgfqpoint{1.611520in}{1.962062in}}%
\pgfpathlineto{\pgfqpoint{1.624849in}{1.960734in}}%
\pgfpathlineto{\pgfqpoint{1.625532in}{1.967784in}}%
\pgfpathlineto{\pgfqpoint{1.635984in}{1.966799in}}%
\pgfpathlineto{\pgfqpoint{1.673140in}{1.963430in}}%
\pgfpathlineto{\pgfqpoint{1.670909in}{1.935733in}}%
\pgfpathlineto{\pgfqpoint{1.667012in}{1.936010in}}%
\pgfpathlineto{\pgfqpoint{1.664670in}{1.909384in}}%
\pgfpathclose%
\pgfusepath{fill}%
\end{pgfscope}%
\begin{pgfscope}%
\pgfpathrectangle{\pgfqpoint{0.100000in}{0.100000in}}{\pgfqpoint{3.608454in}{2.310000in}}%
\pgfusepath{clip}%
\pgfsetbuttcap%
\pgfsetmiterjoin%
\definecolor{currentfill}{rgb}{0.000000,0.521569,0.739216}%
\pgfsetfillcolor{currentfill}%
\pgfsetlinewidth{0.000000pt}%
\definecolor{currentstroke}{rgb}{0.000000,0.000000,0.000000}%
\pgfsetstrokecolor{currentstroke}%
\pgfsetstrokeopacity{0.000000}%
\pgfsetdash{}{0pt}%
\pgfpathmoveto{\pgfqpoint{3.282347in}{1.403895in}}%
\pgfpathlineto{\pgfqpoint{3.278306in}{1.404567in}}%
\pgfpathlineto{\pgfqpoint{3.267145in}{1.400703in}}%
\pgfpathlineto{\pgfqpoint{3.253828in}{1.416317in}}%
\pgfpathlineto{\pgfqpoint{3.255954in}{1.419237in}}%
\pgfpathlineto{\pgfqpoint{3.254537in}{1.428747in}}%
\pgfpathlineto{\pgfqpoint{3.261997in}{1.431811in}}%
\pgfpathlineto{\pgfqpoint{3.248879in}{1.439534in}}%
\pgfpathlineto{\pgfqpoint{3.254321in}{1.452877in}}%
\pgfpathlineto{\pgfqpoint{3.250356in}{1.457926in}}%
\pgfpathlineto{\pgfqpoint{3.250047in}{1.464723in}}%
\pgfpathlineto{\pgfqpoint{3.246140in}{1.469571in}}%
\pgfpathlineto{\pgfqpoint{3.249034in}{1.483689in}}%
\pgfpathlineto{\pgfqpoint{3.256327in}{1.489994in}}%
\pgfpathlineto{\pgfqpoint{3.266239in}{1.490550in}}%
\pgfpathlineto{\pgfqpoint{3.273108in}{1.493337in}}%
\pgfpathlineto{\pgfqpoint{3.274847in}{1.487194in}}%
\pgfpathlineto{\pgfqpoint{3.290599in}{1.431089in}}%
\pgfpathlineto{\pgfqpoint{3.284538in}{1.423341in}}%
\pgfpathlineto{\pgfqpoint{3.281272in}{1.408524in}}%
\pgfpathlineto{\pgfqpoint{3.282347in}{1.403895in}}%
\pgfpathclose%
\pgfusepath{fill}%
\end{pgfscope}%
\begin{pgfscope}%
\pgfpathrectangle{\pgfqpoint{0.100000in}{0.100000in}}{\pgfqpoint{3.608454in}{2.310000in}}%
\pgfusepath{clip}%
\pgfsetbuttcap%
\pgfsetmiterjoin%
\definecolor{currentfill}{rgb}{0.000000,0.443137,0.778431}%
\pgfsetfillcolor{currentfill}%
\pgfsetlinewidth{0.000000pt}%
\definecolor{currentstroke}{rgb}{0.000000,0.000000,0.000000}%
\pgfsetstrokecolor{currentstroke}%
\pgfsetstrokeopacity{0.000000}%
\pgfsetdash{}{0pt}%
\pgfpathmoveto{\pgfqpoint{2.529824in}{1.778725in}}%
\pgfpathlineto{\pgfqpoint{2.510243in}{1.777129in}}%
\pgfpathlineto{\pgfqpoint{2.510762in}{1.770224in}}%
\pgfpathlineto{\pgfqpoint{2.501961in}{1.769586in}}%
\pgfpathlineto{\pgfqpoint{2.493489in}{1.769026in}}%
\pgfpathlineto{\pgfqpoint{2.491553in}{1.796332in}}%
\pgfpathlineto{\pgfqpoint{2.488486in}{1.795943in}}%
\pgfpathlineto{\pgfqpoint{2.486722in}{1.817142in}}%
\pgfpathlineto{\pgfqpoint{2.473187in}{1.816374in}}%
\pgfpathlineto{\pgfqpoint{2.472001in}{1.837128in}}%
\pgfpathlineto{\pgfqpoint{2.463191in}{1.836540in}}%
\pgfpathlineto{\pgfqpoint{2.460529in}{1.843339in}}%
\pgfpathlineto{\pgfqpoint{2.459694in}{1.857127in}}%
\pgfpathlineto{\pgfqpoint{2.473620in}{1.857983in}}%
\pgfpathlineto{\pgfqpoint{2.480448in}{1.858288in}}%
\pgfpathlineto{\pgfqpoint{2.481279in}{1.844556in}}%
\pgfpathlineto{\pgfqpoint{2.488144in}{1.844796in}}%
\pgfpathlineto{\pgfqpoint{2.492447in}{1.838114in}}%
\pgfpathlineto{\pgfqpoint{2.492843in}{1.831235in}}%
\pgfpathlineto{\pgfqpoint{2.513293in}{1.828191in}}%
\pgfpathlineto{\pgfqpoint{2.502443in}{1.807863in}}%
\pgfpathlineto{\pgfqpoint{2.500063in}{1.795356in}}%
\pgfpathlineto{\pgfqpoint{2.505725in}{1.792829in}}%
\pgfpathlineto{\pgfqpoint{2.509822in}{1.800026in}}%
\pgfpathlineto{\pgfqpoint{2.514631in}{1.802211in}}%
\pgfpathlineto{\pgfqpoint{2.522578in}{1.818649in}}%
\pgfpathlineto{\pgfqpoint{2.529869in}{1.821187in}}%
\pgfpathlineto{\pgfqpoint{2.533595in}{1.825201in}}%
\pgfpathlineto{\pgfqpoint{2.540562in}{1.840166in}}%
\pgfpathlineto{\pgfqpoint{2.541258in}{1.846577in}}%
\pgfpathlineto{\pgfqpoint{2.552435in}{1.851077in}}%
\pgfpathlineto{\pgfqpoint{2.552379in}{1.842511in}}%
\pgfpathlineto{\pgfqpoint{2.545403in}{1.830966in}}%
\pgfpathlineto{\pgfqpoint{2.545230in}{1.823543in}}%
\pgfpathlineto{\pgfqpoint{2.541717in}{1.820773in}}%
\pgfpathlineto{\pgfqpoint{2.533894in}{1.801143in}}%
\pgfpathlineto{\pgfqpoint{2.529824in}{1.778725in}}%
\pgfpathclose%
\pgfusepath{fill}%
\end{pgfscope}%
\begin{pgfscope}%
\pgfpathrectangle{\pgfqpoint{0.100000in}{0.100000in}}{\pgfqpoint{3.608454in}{2.310000in}}%
\pgfusepath{clip}%
\pgfsetbuttcap%
\pgfsetmiterjoin%
\definecolor{currentfill}{rgb}{0.000000,0.654902,0.672549}%
\pgfsetfillcolor{currentfill}%
\pgfsetlinewidth{0.000000pt}%
\definecolor{currentstroke}{rgb}{0.000000,0.000000,0.000000}%
\pgfsetstrokecolor{currentstroke}%
\pgfsetstrokeopacity{0.000000}%
\pgfsetdash{}{0pt}%
\pgfpathmoveto{\pgfqpoint{3.143924in}{1.538541in}}%
\pgfpathlineto{\pgfqpoint{3.137378in}{1.521010in}}%
\pgfpathlineto{\pgfqpoint{3.133137in}{1.523448in}}%
\pgfpathlineto{\pgfqpoint{3.119792in}{1.526024in}}%
\pgfpathlineto{\pgfqpoint{3.112902in}{1.529498in}}%
\pgfpathlineto{\pgfqpoint{3.110584in}{1.534960in}}%
\pgfpathlineto{\pgfqpoint{3.115439in}{1.550679in}}%
\pgfpathlineto{\pgfqpoint{3.110368in}{1.555879in}}%
\pgfpathlineto{\pgfqpoint{3.111267in}{1.560885in}}%
\pgfpathlineto{\pgfqpoint{3.107213in}{1.564628in}}%
\pgfpathlineto{\pgfqpoint{3.098915in}{1.567130in}}%
\pgfpathlineto{\pgfqpoint{3.072318in}{1.562246in}}%
\pgfpathlineto{\pgfqpoint{3.069571in}{1.576329in}}%
\pgfpathlineto{\pgfqpoint{3.045363in}{1.572457in}}%
\pgfpathlineto{\pgfqpoint{3.043474in}{1.583306in}}%
\pgfpathlineto{\pgfqpoint{3.039983in}{1.605278in}}%
\pgfpathlineto{\pgfqpoint{3.046524in}{1.607153in}}%
\pgfpathlineto{\pgfqpoint{3.054170in}{1.610224in}}%
\pgfpathlineto{\pgfqpoint{3.068655in}{1.604850in}}%
\pgfpathlineto{\pgfqpoint{3.106858in}{1.608202in}}%
\pgfpathlineto{\pgfqpoint{3.111013in}{1.603891in}}%
\pgfpathlineto{\pgfqpoint{3.125203in}{1.608488in}}%
\pgfpathlineto{\pgfqpoint{3.128786in}{1.603472in}}%
\pgfpathlineto{\pgfqpoint{3.133554in}{1.604605in}}%
\pgfpathlineto{\pgfqpoint{3.144862in}{1.595009in}}%
\pgfpathlineto{\pgfqpoint{3.154199in}{1.598846in}}%
\pgfpathlineto{\pgfqpoint{3.165006in}{1.605498in}}%
\pgfpathlineto{\pgfqpoint{3.155139in}{1.587716in}}%
\pgfpathlineto{\pgfqpoint{3.139700in}{1.576466in}}%
\pgfpathlineto{\pgfqpoint{3.131483in}{1.559276in}}%
\pgfpathlineto{\pgfqpoint{3.133732in}{1.555574in}}%
\pgfpathlineto{\pgfqpoint{3.130234in}{1.548464in}}%
\pgfpathlineto{\pgfqpoint{3.133310in}{1.544667in}}%
\pgfpathlineto{\pgfqpoint{3.138889in}{1.546016in}}%
\pgfpathlineto{\pgfqpoint{3.143924in}{1.538541in}}%
\pgfpathclose%
\pgfusepath{fill}%
\end{pgfscope}%
\begin{pgfscope}%
\pgfpathrectangle{\pgfqpoint{0.100000in}{0.100000in}}{\pgfqpoint{3.608454in}{2.310000in}}%
\pgfusepath{clip}%
\pgfsetbuttcap%
\pgfsetmiterjoin%
\definecolor{currentfill}{rgb}{0.000000,0.345098,0.827451}%
\pgfsetfillcolor{currentfill}%
\pgfsetlinewidth{0.000000pt}%
\definecolor{currentstroke}{rgb}{0.000000,0.000000,0.000000}%
\pgfsetstrokecolor{currentstroke}%
\pgfsetstrokeopacity{0.000000}%
\pgfsetdash{}{0pt}%
\pgfpathmoveto{\pgfqpoint{2.196902in}{1.594193in}}%
\pgfpathlineto{\pgfqpoint{2.155774in}{1.593709in}}%
\pgfpathlineto{\pgfqpoint{2.115090in}{1.593467in}}%
\pgfpathlineto{\pgfqpoint{2.115152in}{1.614399in}}%
\pgfpathlineto{\pgfqpoint{2.111772in}{1.614418in}}%
\pgfpathlineto{\pgfqpoint{2.111776in}{1.648912in}}%
\pgfpathlineto{\pgfqpoint{2.166463in}{1.649078in}}%
\pgfpathlineto{\pgfqpoint{2.193901in}{1.649520in}}%
\pgfpathlineto{\pgfqpoint{2.194446in}{1.614873in}}%
\pgfpathlineto{\pgfqpoint{2.196575in}{1.614898in}}%
\pgfpathlineto{\pgfqpoint{2.196902in}{1.594193in}}%
\pgfpathclose%
\pgfusepath{fill}%
\end{pgfscope}%
\begin{pgfscope}%
\pgfpathrectangle{\pgfqpoint{0.100000in}{0.100000in}}{\pgfqpoint{3.608454in}{2.310000in}}%
\pgfusepath{clip}%
\pgfsetbuttcap%
\pgfsetmiterjoin%
\definecolor{currentfill}{rgb}{0.000000,0.447059,0.776471}%
\pgfsetfillcolor{currentfill}%
\pgfsetlinewidth{0.000000pt}%
\definecolor{currentstroke}{rgb}{0.000000,0.000000,0.000000}%
\pgfsetstrokecolor{currentstroke}%
\pgfsetstrokeopacity{0.000000}%
\pgfsetdash{}{0pt}%
\pgfpathmoveto{\pgfqpoint{3.466647in}{1.697217in}}%
\pgfpathlineto{\pgfqpoint{3.461826in}{1.697323in}}%
\pgfpathlineto{\pgfqpoint{3.452537in}{1.691767in}}%
\pgfpathlineto{\pgfqpoint{3.447030in}{1.691452in}}%
\pgfpathlineto{\pgfqpoint{3.437379in}{1.684748in}}%
\pgfpathlineto{\pgfqpoint{3.433364in}{1.685208in}}%
\pgfpathlineto{\pgfqpoint{3.406113in}{1.677314in}}%
\pgfpathlineto{\pgfqpoint{3.369189in}{1.645562in}}%
\pgfpathlineto{\pgfqpoint{3.362988in}{1.653409in}}%
\pgfpathlineto{\pgfqpoint{3.375111in}{1.665432in}}%
\pgfpathlineto{\pgfqpoint{3.369634in}{1.670838in}}%
\pgfpathlineto{\pgfqpoint{3.366550in}{1.688953in}}%
\pgfpathlineto{\pgfqpoint{3.359316in}{1.729735in}}%
\pgfpathlineto{\pgfqpoint{3.384285in}{1.735025in}}%
\pgfpathlineto{\pgfqpoint{3.436777in}{1.747450in}}%
\pgfpathlineto{\pgfqpoint{3.456367in}{1.750710in}}%
\pgfpathlineto{\pgfqpoint{3.464103in}{1.723936in}}%
\pgfpathlineto{\pgfqpoint{3.468281in}{1.705449in}}%
\pgfpathlineto{\pgfqpoint{3.466647in}{1.697217in}}%
\pgfpathclose%
\pgfusepath{fill}%
\end{pgfscope}%
\begin{pgfscope}%
\pgfpathrectangle{\pgfqpoint{0.100000in}{0.100000in}}{\pgfqpoint{3.608454in}{2.310000in}}%
\pgfusepath{clip}%
\pgfsetbuttcap%
\pgfsetmiterjoin%
\definecolor{currentfill}{rgb}{0.000000,0.580392,0.709804}%
\pgfsetfillcolor{currentfill}%
\pgfsetlinewidth{0.000000pt}%
\definecolor{currentstroke}{rgb}{0.000000,0.000000,0.000000}%
\pgfsetstrokecolor{currentstroke}%
\pgfsetstrokeopacity{0.000000}%
\pgfsetdash{}{0pt}%
\pgfpathmoveto{\pgfqpoint{2.353369in}{1.446719in}}%
\pgfpathlineto{\pgfqpoint{2.353565in}{1.439495in}}%
\pgfpathlineto{\pgfqpoint{2.317746in}{1.438593in}}%
\pgfpathlineto{\pgfqpoint{2.318087in}{1.442426in}}%
\pgfpathlineto{\pgfqpoint{2.290718in}{1.442111in}}%
\pgfpathlineto{\pgfqpoint{2.290253in}{1.469670in}}%
\pgfpathlineto{\pgfqpoint{2.303069in}{1.470734in}}%
\pgfpathlineto{\pgfqpoint{2.302293in}{1.514165in}}%
\pgfpathlineto{\pgfqpoint{2.322856in}{1.514824in}}%
\pgfpathlineto{\pgfqpoint{2.323109in}{1.507938in}}%
\pgfpathlineto{\pgfqpoint{2.348213in}{1.508885in}}%
\pgfpathlineto{\pgfqpoint{2.358001in}{1.509016in}}%
\pgfpathlineto{\pgfqpoint{2.359331in}{1.474799in}}%
\pgfpathlineto{\pgfqpoint{2.352440in}{1.474762in}}%
\pgfpathlineto{\pgfqpoint{2.353369in}{1.446719in}}%
\pgfpathclose%
\pgfusepath{fill}%
\end{pgfscope}%
\begin{pgfscope}%
\pgfpathrectangle{\pgfqpoint{0.100000in}{0.100000in}}{\pgfqpoint{3.608454in}{2.310000in}}%
\pgfusepath{clip}%
\pgfsetbuttcap%
\pgfsetmiterjoin%
\definecolor{currentfill}{rgb}{0.000000,0.635294,0.682353}%
\pgfsetfillcolor{currentfill}%
\pgfsetlinewidth{0.000000pt}%
\definecolor{currentstroke}{rgb}{0.000000,0.000000,0.000000}%
\pgfsetstrokecolor{currentstroke}%
\pgfsetstrokeopacity{0.000000}%
\pgfsetdash{}{0pt}%
\pgfpathmoveto{\pgfqpoint{2.798375in}{1.062974in}}%
\pgfpathlineto{\pgfqpoint{2.779057in}{1.060719in}}%
\pgfpathlineto{\pgfqpoint{2.756131in}{1.058082in}}%
\pgfpathlineto{\pgfqpoint{2.751742in}{1.068739in}}%
\pgfpathlineto{\pgfqpoint{2.755774in}{1.081517in}}%
\pgfpathlineto{\pgfqpoint{2.749929in}{1.091001in}}%
\pgfpathlineto{\pgfqpoint{2.755630in}{1.093871in}}%
\pgfpathlineto{\pgfqpoint{2.763588in}{1.107918in}}%
\pgfpathlineto{\pgfqpoint{2.762622in}{1.113861in}}%
\pgfpathlineto{\pgfqpoint{2.765621in}{1.120136in}}%
\pgfpathlineto{\pgfqpoint{2.772428in}{1.117636in}}%
\pgfpathlineto{\pgfqpoint{2.775615in}{1.111731in}}%
\pgfpathlineto{\pgfqpoint{2.790002in}{1.116268in}}%
\pgfpathlineto{\pgfqpoint{2.805797in}{1.111156in}}%
\pgfpathlineto{\pgfqpoint{2.817068in}{1.103142in}}%
\pgfpathlineto{\pgfqpoint{2.813980in}{1.097663in}}%
\pgfpathlineto{\pgfqpoint{2.814347in}{1.089123in}}%
\pgfpathlineto{\pgfqpoint{2.807766in}{1.084602in}}%
\pgfpathlineto{\pgfqpoint{2.802072in}{1.085804in}}%
\pgfpathlineto{\pgfqpoint{2.798205in}{1.081874in}}%
\pgfpathlineto{\pgfqpoint{2.798375in}{1.062974in}}%
\pgfpathclose%
\pgfusepath{fill}%
\end{pgfscope}%
\begin{pgfscope}%
\pgfpathrectangle{\pgfqpoint{0.100000in}{0.100000in}}{\pgfqpoint{3.608454in}{2.310000in}}%
\pgfusepath{clip}%
\pgfsetbuttcap%
\pgfsetmiterjoin%
\definecolor{currentfill}{rgb}{0.000000,0.682353,0.658824}%
\pgfsetfillcolor{currentfill}%
\pgfsetlinewidth{0.000000pt}%
\definecolor{currentstroke}{rgb}{0.000000,0.000000,0.000000}%
\pgfsetstrokecolor{currentstroke}%
\pgfsetstrokeopacity{0.000000}%
\pgfsetdash{}{0pt}%
\pgfpathmoveto{\pgfqpoint{1.516667in}{1.064131in}}%
\pgfpathlineto{\pgfqpoint{1.517207in}{1.069928in}}%
\pgfpathlineto{\pgfqpoint{1.442195in}{1.078025in}}%
\pgfpathlineto{\pgfqpoint{1.440664in}{1.064305in}}%
\pgfpathlineto{\pgfqpoint{1.413364in}{1.067432in}}%
\pgfpathlineto{\pgfqpoint{1.421911in}{1.141159in}}%
\pgfpathlineto{\pgfqpoint{1.431076in}{1.139578in}}%
\pgfpathlineto{\pgfqpoint{1.442215in}{1.150974in}}%
\pgfpathlineto{\pgfqpoint{1.445381in}{1.160260in}}%
\pgfpathlineto{\pgfqpoint{1.448410in}{1.160662in}}%
\pgfpathlineto{\pgfqpoint{1.450548in}{1.172051in}}%
\pgfpathlineto{\pgfqpoint{1.450272in}{1.192914in}}%
\pgfpathlineto{\pgfqpoint{1.458207in}{1.195789in}}%
\pgfpathlineto{\pgfqpoint{1.462472in}{1.211222in}}%
\pgfpathlineto{\pgfqpoint{1.462544in}{1.217413in}}%
\pgfpathlineto{\pgfqpoint{1.466660in}{1.216957in}}%
\pgfpathlineto{\pgfqpoint{1.513551in}{1.211748in}}%
\pgfpathlineto{\pgfqpoint{1.538731in}{1.209408in}}%
\pgfpathlineto{\pgfqpoint{1.532570in}{1.148149in}}%
\pgfpathlineto{\pgfqpoint{1.505428in}{1.150972in}}%
\pgfpathlineto{\pgfqpoint{1.508502in}{1.136752in}}%
\pgfpathlineto{\pgfqpoint{1.508110in}{1.125109in}}%
\pgfpathlineto{\pgfqpoint{1.506250in}{1.116084in}}%
\pgfpathlineto{\pgfqpoint{1.528359in}{1.113844in}}%
\pgfpathlineto{\pgfqpoint{1.531357in}{1.115154in}}%
\pgfpathlineto{\pgfqpoint{1.544598in}{1.083669in}}%
\pgfpathlineto{\pgfqpoint{1.550211in}{1.083114in}}%
\pgfpathlineto{\pgfqpoint{1.548858in}{1.068831in}}%
\pgfpathlineto{\pgfqpoint{1.534671in}{1.070306in}}%
\pgfpathlineto{\pgfqpoint{1.516667in}{1.064131in}}%
\pgfpathclose%
\pgfusepath{fill}%
\end{pgfscope}%
\begin{pgfscope}%
\pgfpathrectangle{\pgfqpoint{0.100000in}{0.100000in}}{\pgfqpoint{3.608454in}{2.310000in}}%
\pgfusepath{clip}%
\pgfsetbuttcap%
\pgfsetmiterjoin%
\definecolor{currentfill}{rgb}{0.000000,0.411765,0.794118}%
\pgfsetfillcolor{currentfill}%
\pgfsetlinewidth{0.000000pt}%
\definecolor{currentstroke}{rgb}{0.000000,0.000000,0.000000}%
\pgfsetstrokecolor{currentstroke}%
\pgfsetstrokeopacity{0.000000}%
\pgfsetdash{}{0pt}%
\pgfpathmoveto{\pgfqpoint{2.266869in}{1.540448in}}%
\pgfpathlineto{\pgfqpoint{2.205803in}{1.538835in}}%
\pgfpathlineto{\pgfqpoint{2.204489in}{1.546166in}}%
\pgfpathlineto{\pgfqpoint{2.204148in}{1.566920in}}%
\pgfpathlineto{\pgfqpoint{2.210987in}{1.567017in}}%
\pgfpathlineto{\pgfqpoint{2.210467in}{1.594411in}}%
\pgfpathlineto{\pgfqpoint{2.223956in}{1.594717in}}%
\pgfpathlineto{\pgfqpoint{2.237674in}{1.595106in}}%
\pgfpathlineto{\pgfqpoint{2.237469in}{1.601952in}}%
\pgfpathlineto{\pgfqpoint{2.264937in}{1.602761in}}%
\pgfpathlineto{\pgfqpoint{2.266869in}{1.540448in}}%
\pgfpathclose%
\pgfusepath{fill}%
\end{pgfscope}%
\begin{pgfscope}%
\pgfpathrectangle{\pgfqpoint{0.100000in}{0.100000in}}{\pgfqpoint{3.608454in}{2.310000in}}%
\pgfusepath{clip}%
\pgfsetbuttcap%
\pgfsetmiterjoin%
\definecolor{currentfill}{rgb}{0.000000,0.647059,0.676471}%
\pgfsetfillcolor{currentfill}%
\pgfsetlinewidth{0.000000pt}%
\definecolor{currentstroke}{rgb}{0.000000,0.000000,0.000000}%
\pgfsetstrokecolor{currentstroke}%
\pgfsetstrokeopacity{0.000000}%
\pgfsetdash{}{0pt}%
\pgfpathmoveto{\pgfqpoint{0.718561in}{1.521037in}}%
\pgfpathlineto{\pgfqpoint{0.697419in}{1.520493in}}%
\pgfpathlineto{\pgfqpoint{0.695807in}{1.514308in}}%
\pgfpathlineto{\pgfqpoint{0.717194in}{1.472445in}}%
\pgfpathlineto{\pgfqpoint{0.665974in}{1.438337in}}%
\pgfpathlineto{\pgfqpoint{0.661400in}{1.439746in}}%
\pgfpathlineto{\pgfqpoint{0.628251in}{1.491065in}}%
\pgfpathlineto{\pgfqpoint{0.643201in}{1.487011in}}%
\pgfpathlineto{\pgfqpoint{0.652474in}{1.520325in}}%
\pgfpathlineto{\pgfqpoint{0.645777in}{1.522062in}}%
\pgfpathlineto{\pgfqpoint{0.647659in}{1.529088in}}%
\pgfpathlineto{\pgfqpoint{0.655392in}{1.537674in}}%
\pgfpathlineto{\pgfqpoint{0.718561in}{1.521037in}}%
\pgfpathclose%
\pgfusepath{fill}%
\end{pgfscope}%
\begin{pgfscope}%
\pgfpathrectangle{\pgfqpoint{0.100000in}{0.100000in}}{\pgfqpoint{3.608454in}{2.310000in}}%
\pgfusepath{clip}%
\pgfsetbuttcap%
\pgfsetmiterjoin%
\definecolor{currentfill}{rgb}{0.000000,0.737255,0.631373}%
\pgfsetfillcolor{currentfill}%
\pgfsetlinewidth{0.000000pt}%
\definecolor{currentstroke}{rgb}{0.000000,0.000000,0.000000}%
\pgfsetstrokecolor{currentstroke}%
\pgfsetstrokeopacity{0.000000}%
\pgfsetdash{}{0pt}%
\pgfpathmoveto{\pgfqpoint{2.996900in}{1.203053in}}%
\pgfpathlineto{\pgfqpoint{2.987922in}{1.204284in}}%
\pgfpathlineto{\pgfqpoint{2.980945in}{1.198576in}}%
\pgfpathlineto{\pgfqpoint{2.970773in}{1.210739in}}%
\pgfpathlineto{\pgfqpoint{2.970559in}{1.212860in}}%
\pgfpathlineto{\pgfqpoint{2.950033in}{1.210864in}}%
\pgfpathlineto{\pgfqpoint{2.951682in}{1.213011in}}%
\pgfpathlineto{\pgfqpoint{2.957411in}{1.217461in}}%
\pgfpathlineto{\pgfqpoint{2.958209in}{1.221651in}}%
\pgfpathlineto{\pgfqpoint{2.974075in}{1.228469in}}%
\pgfpathlineto{\pgfqpoint{2.984220in}{1.230496in}}%
\pgfpathlineto{\pgfqpoint{2.986436in}{1.233543in}}%
\pgfpathlineto{\pgfqpoint{3.005275in}{1.242589in}}%
\pgfpathlineto{\pgfqpoint{3.011272in}{1.247768in}}%
\pgfpathlineto{\pgfqpoint{3.024683in}{1.232048in}}%
\pgfpathlineto{\pgfqpoint{3.025017in}{1.229246in}}%
\pgfpathlineto{\pgfqpoint{3.015913in}{1.223529in}}%
\pgfpathlineto{\pgfqpoint{3.017299in}{1.218774in}}%
\pgfpathlineto{\pgfqpoint{2.999075in}{1.216208in}}%
\pgfpathlineto{\pgfqpoint{2.995823in}{1.206373in}}%
\pgfpathlineto{\pgfqpoint{2.996900in}{1.203053in}}%
\pgfpathclose%
\pgfusepath{fill}%
\end{pgfscope}%
\begin{pgfscope}%
\pgfpathrectangle{\pgfqpoint{0.100000in}{0.100000in}}{\pgfqpoint{3.608454in}{2.310000in}}%
\pgfusepath{clip}%
\pgfsetbuttcap%
\pgfsetmiterjoin%
\definecolor{currentfill}{rgb}{0.000000,0.678431,0.660784}%
\pgfsetfillcolor{currentfill}%
\pgfsetlinewidth{0.000000pt}%
\definecolor{currentstroke}{rgb}{0.000000,0.000000,0.000000}%
\pgfsetstrokecolor{currentstroke}%
\pgfsetstrokeopacity{0.000000}%
\pgfsetdash{}{0pt}%
\pgfpathmoveto{\pgfqpoint{2.869199in}{1.556678in}}%
\pgfpathlineto{\pgfqpoint{2.872479in}{1.551436in}}%
\pgfpathlineto{\pgfqpoint{2.876286in}{1.526111in}}%
\pgfpathlineto{\pgfqpoint{2.870685in}{1.525264in}}%
\pgfpathlineto{\pgfqpoint{2.871819in}{1.517489in}}%
\pgfpathlineto{\pgfqpoint{2.862804in}{1.514836in}}%
\pgfpathlineto{\pgfqpoint{2.846858in}{1.512561in}}%
\pgfpathlineto{\pgfqpoint{2.848845in}{1.496480in}}%
\pgfpathlineto{\pgfqpoint{2.831782in}{1.495222in}}%
\pgfpathlineto{\pgfqpoint{2.831500in}{1.499705in}}%
\pgfpathlineto{\pgfqpoint{2.811880in}{1.499432in}}%
\pgfpathlineto{\pgfqpoint{2.811270in}{1.504385in}}%
\pgfpathlineto{\pgfqpoint{2.801132in}{1.502933in}}%
\pgfpathlineto{\pgfqpoint{2.799219in}{1.517163in}}%
\pgfpathlineto{\pgfqpoint{2.817395in}{1.520810in}}%
\pgfpathlineto{\pgfqpoint{2.814322in}{1.543583in}}%
\pgfpathlineto{\pgfqpoint{2.831040in}{1.546066in}}%
\pgfpathlineto{\pgfqpoint{2.854535in}{1.548999in}}%
\pgfpathlineto{\pgfqpoint{2.853507in}{1.554639in}}%
\pgfpathlineto{\pgfqpoint{2.869199in}{1.556678in}}%
\pgfpathclose%
\pgfusepath{fill}%
\end{pgfscope}%
\begin{pgfscope}%
\pgfpathrectangle{\pgfqpoint{0.100000in}{0.100000in}}{\pgfqpoint{3.608454in}{2.310000in}}%
\pgfusepath{clip}%
\pgfsetbuttcap%
\pgfsetmiterjoin%
\definecolor{currentfill}{rgb}{0.000000,0.878431,0.560784}%
\pgfsetfillcolor{currentfill}%
\pgfsetlinewidth{0.000000pt}%
\definecolor{currentstroke}{rgb}{0.000000,0.000000,0.000000}%
\pgfsetstrokecolor{currentstroke}%
\pgfsetstrokeopacity{0.000000}%
\pgfsetdash{}{0pt}%
\pgfpathmoveto{\pgfqpoint{2.722462in}{1.100629in}}%
\pgfpathlineto{\pgfqpoint{2.731314in}{1.113832in}}%
\pgfpathlineto{\pgfqpoint{2.732727in}{1.126386in}}%
\pgfpathlineto{\pgfqpoint{2.729200in}{1.134130in}}%
\pgfpathlineto{\pgfqpoint{2.729930in}{1.144288in}}%
\pgfpathlineto{\pgfqpoint{2.735267in}{1.144371in}}%
\pgfpathlineto{\pgfqpoint{2.738900in}{1.152461in}}%
\pgfpathlineto{\pgfqpoint{2.735445in}{1.160714in}}%
\pgfpathlineto{\pgfqpoint{2.734679in}{1.168795in}}%
\pgfpathlineto{\pgfqpoint{2.738211in}{1.181250in}}%
\pgfpathlineto{\pgfqpoint{2.744380in}{1.183721in}}%
\pgfpathlineto{\pgfqpoint{2.758225in}{1.180770in}}%
\pgfpathlineto{\pgfqpoint{2.763981in}{1.171120in}}%
\pgfpathlineto{\pgfqpoint{2.761563in}{1.168904in}}%
\pgfpathlineto{\pgfqpoint{2.751113in}{1.160516in}}%
\pgfpathlineto{\pgfqpoint{2.750357in}{1.150558in}}%
\pgfpathlineto{\pgfqpoint{2.763659in}{1.139195in}}%
\pgfpathlineto{\pgfqpoint{2.766985in}{1.132686in}}%
\pgfpathlineto{\pgfqpoint{2.761095in}{1.125193in}}%
\pgfpathlineto{\pgfqpoint{2.765621in}{1.120136in}}%
\pgfpathlineto{\pgfqpoint{2.762622in}{1.113861in}}%
\pgfpathlineto{\pgfqpoint{2.763588in}{1.107918in}}%
\pgfpathlineto{\pgfqpoint{2.755630in}{1.093871in}}%
\pgfpathlineto{\pgfqpoint{2.749929in}{1.091001in}}%
\pgfpathlineto{\pgfqpoint{2.739787in}{1.093594in}}%
\pgfpathlineto{\pgfqpoint{2.736951in}{1.085202in}}%
\pgfpathlineto{\pgfqpoint{2.724406in}{1.094850in}}%
\pgfpathlineto{\pgfqpoint{2.722462in}{1.100629in}}%
\pgfpathclose%
\pgfusepath{fill}%
\end{pgfscope}%
\begin{pgfscope}%
\pgfpathrectangle{\pgfqpoint{0.100000in}{0.100000in}}{\pgfqpoint{3.608454in}{2.310000in}}%
\pgfusepath{clip}%
\pgfsetbuttcap%
\pgfsetmiterjoin%
\definecolor{currentfill}{rgb}{0.000000,0.690196,0.654902}%
\pgfsetfillcolor{currentfill}%
\pgfsetlinewidth{0.000000pt}%
\definecolor{currentstroke}{rgb}{0.000000,0.000000,0.000000}%
\pgfsetstrokecolor{currentstroke}%
\pgfsetstrokeopacity{0.000000}%
\pgfsetdash{}{0pt}%
\pgfpathmoveto{\pgfqpoint{3.271268in}{1.222696in}}%
\pgfpathlineto{\pgfqpoint{3.272014in}{1.215831in}}%
\pgfpathlineto{\pgfqpoint{3.269662in}{1.209194in}}%
\pgfpathlineto{\pgfqpoint{3.265874in}{1.209619in}}%
\pgfpathlineto{\pgfqpoint{3.262155in}{1.203038in}}%
\pgfpathlineto{\pgfqpoint{3.258356in}{1.209245in}}%
\pgfpathlineto{\pgfqpoint{3.253779in}{1.205127in}}%
\pgfpathlineto{\pgfqpoint{3.247206in}{1.209157in}}%
\pgfpathlineto{\pgfqpoint{3.239989in}{1.209565in}}%
\pgfpathlineto{\pgfqpoint{3.238791in}{1.214139in}}%
\pgfpathlineto{\pgfqpoint{3.230833in}{1.216563in}}%
\pgfpathlineto{\pgfqpoint{3.226514in}{1.211542in}}%
\pgfpathlineto{\pgfqpoint{3.218220in}{1.212547in}}%
\pgfpathlineto{\pgfqpoint{3.216999in}{1.216602in}}%
\pgfpathlineto{\pgfqpoint{3.206943in}{1.219779in}}%
\pgfpathlineto{\pgfqpoint{3.195538in}{1.217166in}}%
\pgfpathlineto{\pgfqpoint{3.185638in}{1.219606in}}%
\pgfpathlineto{\pgfqpoint{3.188912in}{1.234693in}}%
\pgfpathlineto{\pgfqpoint{3.187275in}{1.247460in}}%
\pgfpathlineto{\pgfqpoint{3.195545in}{1.249102in}}%
\pgfpathlineto{\pgfqpoint{3.199968in}{1.263498in}}%
\pgfpathlineto{\pgfqpoint{3.197116in}{1.277528in}}%
\pgfpathlineto{\pgfqpoint{3.202830in}{1.275174in}}%
\pgfpathlineto{\pgfqpoint{3.209973in}{1.277484in}}%
\pgfpathlineto{\pgfqpoint{3.210178in}{1.266083in}}%
\pgfpathlineto{\pgfqpoint{3.214106in}{1.265774in}}%
\pgfpathlineto{\pgfqpoint{3.221614in}{1.254572in}}%
\pgfpathlineto{\pgfqpoint{3.248646in}{1.259422in}}%
\pgfpathlineto{\pgfqpoint{3.249894in}{1.248153in}}%
\pgfpathlineto{\pgfqpoint{3.259748in}{1.247063in}}%
\pgfpathlineto{\pgfqpoint{3.266234in}{1.243000in}}%
\pgfpathlineto{\pgfqpoint{3.265839in}{1.231923in}}%
\pgfpathlineto{\pgfqpoint{3.271268in}{1.222696in}}%
\pgfpathclose%
\pgfusepath{fill}%
\end{pgfscope}%
\begin{pgfscope}%
\pgfpathrectangle{\pgfqpoint{0.100000in}{0.100000in}}{\pgfqpoint{3.608454in}{2.310000in}}%
\pgfusepath{clip}%
\pgfsetbuttcap%
\pgfsetmiterjoin%
\definecolor{currentfill}{rgb}{0.000000,0.486275,0.756863}%
\pgfsetfillcolor{currentfill}%
\pgfsetlinewidth{0.000000pt}%
\definecolor{currentstroke}{rgb}{0.000000,0.000000,0.000000}%
\pgfsetstrokecolor{currentstroke}%
\pgfsetstrokeopacity{0.000000}%
\pgfsetdash{}{0pt}%
\pgfpathmoveto{\pgfqpoint{1.826418in}{1.110191in}}%
\pgfpathlineto{\pgfqpoint{1.818722in}{1.107444in}}%
\pgfpathlineto{\pgfqpoint{1.814486in}{1.099133in}}%
\pgfpathlineto{\pgfqpoint{1.804299in}{1.099161in}}%
\pgfpathlineto{\pgfqpoint{1.798623in}{1.104938in}}%
\pgfpathlineto{\pgfqpoint{1.792285in}{1.107896in}}%
\pgfpathlineto{\pgfqpoint{1.786242in}{1.101806in}}%
\pgfpathlineto{\pgfqpoint{1.787003in}{1.115616in}}%
\pgfpathlineto{\pgfqpoint{1.788727in}{1.150707in}}%
\pgfpathlineto{\pgfqpoint{1.790983in}{1.190359in}}%
\pgfpathlineto{\pgfqpoint{1.820080in}{1.188664in}}%
\pgfpathlineto{\pgfqpoint{1.825437in}{1.188388in}}%
\pgfpathlineto{\pgfqpoint{1.830118in}{1.178718in}}%
\pgfpathlineto{\pgfqpoint{1.837450in}{1.170858in}}%
\pgfpathlineto{\pgfqpoint{1.844616in}{1.170342in}}%
\pgfpathlineto{\pgfqpoint{1.855366in}{1.155288in}}%
\pgfpathlineto{\pgfqpoint{1.854121in}{1.120697in}}%
\pgfpathlineto{\pgfqpoint{1.826923in}{1.122162in}}%
\pgfpathlineto{\pgfqpoint{1.826418in}{1.110191in}}%
\pgfpathclose%
\pgfusepath{fill}%
\end{pgfscope}%
\begin{pgfscope}%
\pgfpathrectangle{\pgfqpoint{0.100000in}{0.100000in}}{\pgfqpoint{3.608454in}{2.310000in}}%
\pgfusepath{clip}%
\pgfsetbuttcap%
\pgfsetmiterjoin%
\definecolor{currentfill}{rgb}{0.000000,0.313725,0.843137}%
\pgfsetfillcolor{currentfill}%
\pgfsetlinewidth{0.000000pt}%
\definecolor{currentstroke}{rgb}{0.000000,0.000000,0.000000}%
\pgfsetstrokecolor{currentstroke}%
\pgfsetstrokeopacity{0.000000}%
\pgfsetdash{}{0pt}%
\pgfpathmoveto{\pgfqpoint{3.083124in}{1.089029in}}%
\pgfpathlineto{\pgfqpoint{3.042269in}{1.083215in}}%
\pgfpathlineto{\pgfqpoint{3.054685in}{1.068940in}}%
\pgfpathlineto{\pgfqpoint{3.046213in}{1.063185in}}%
\pgfpathlineto{\pgfqpoint{3.047179in}{1.057968in}}%
\pgfpathlineto{\pgfqpoint{3.040979in}{1.056871in}}%
\pgfpathlineto{\pgfqpoint{3.034032in}{1.059624in}}%
\pgfpathlineto{\pgfqpoint{3.026275in}{1.051810in}}%
\pgfpathlineto{\pgfqpoint{3.025339in}{1.058515in}}%
\pgfpathlineto{\pgfqpoint{2.989806in}{1.055375in}}%
\pgfpathlineto{\pgfqpoint{2.986559in}{1.063157in}}%
\pgfpathlineto{\pgfqpoint{2.984590in}{1.076632in}}%
\pgfpathlineto{\pgfqpoint{2.980220in}{1.091029in}}%
\pgfpathlineto{\pgfqpoint{2.985051in}{1.092680in}}%
\pgfpathlineto{\pgfqpoint{2.986399in}{1.102377in}}%
\pgfpathlineto{\pgfqpoint{3.008536in}{1.104446in}}%
\pgfpathlineto{\pgfqpoint{3.007642in}{1.113989in}}%
\pgfpathlineto{\pgfqpoint{3.009870in}{1.122844in}}%
\pgfpathlineto{\pgfqpoint{3.008399in}{1.134076in}}%
\pgfpathlineto{\pgfqpoint{3.022456in}{1.135384in}}%
\pgfpathlineto{\pgfqpoint{3.018191in}{1.148162in}}%
\pgfpathlineto{\pgfqpoint{3.020036in}{1.162708in}}%
\pgfpathlineto{\pgfqpoint{3.024951in}{1.162919in}}%
\pgfpathlineto{\pgfqpoint{3.037191in}{1.156638in}}%
\pgfpathlineto{\pgfqpoint{3.046178in}{1.152956in}}%
\pgfpathlineto{\pgfqpoint{3.055178in}{1.146538in}}%
\pgfpathlineto{\pgfqpoint{3.057775in}{1.140860in}}%
\pgfpathlineto{\pgfqpoint{3.062964in}{1.138236in}}%
\pgfpathlineto{\pgfqpoint{3.067859in}{1.131866in}}%
\pgfpathlineto{\pgfqpoint{3.066320in}{1.122502in}}%
\pgfpathlineto{\pgfqpoint{3.070229in}{1.116768in}}%
\pgfpathlineto{\pgfqpoint{3.070169in}{1.111400in}}%
\pgfpathlineto{\pgfqpoint{3.078110in}{1.113287in}}%
\pgfpathlineto{\pgfqpoint{3.083344in}{1.107498in}}%
\pgfpathlineto{\pgfqpoint{3.086909in}{1.097540in}}%
\pgfpathlineto{\pgfqpoint{3.083124in}{1.089029in}}%
\pgfpathclose%
\pgfusepath{fill}%
\end{pgfscope}%
\begin{pgfscope}%
\pgfpathrectangle{\pgfqpoint{0.100000in}{0.100000in}}{\pgfqpoint{3.608454in}{2.310000in}}%
\pgfusepath{clip}%
\pgfsetbuttcap%
\pgfsetmiterjoin%
\definecolor{currentfill}{rgb}{0.000000,0.439216,0.780392}%
\pgfsetfillcolor{currentfill}%
\pgfsetlinewidth{0.000000pt}%
\definecolor{currentstroke}{rgb}{0.000000,0.000000,0.000000}%
\pgfsetstrokecolor{currentstroke}%
\pgfsetstrokeopacity{0.000000}%
\pgfsetdash{}{0pt}%
\pgfpathmoveto{\pgfqpoint{2.034871in}{0.852102in}}%
\pgfpathlineto{\pgfqpoint{2.020147in}{0.852551in}}%
\pgfpathlineto{\pgfqpoint{2.020851in}{0.893088in}}%
\pgfpathlineto{\pgfqpoint{2.049865in}{0.897339in}}%
\pgfpathlineto{\pgfqpoint{2.050404in}{0.930484in}}%
\pgfpathlineto{\pgfqpoint{2.056583in}{0.934051in}}%
\pgfpathlineto{\pgfqpoint{2.068217in}{0.938848in}}%
\pgfpathlineto{\pgfqpoint{2.070424in}{0.935180in}}%
\pgfpathlineto{\pgfqpoint{2.079921in}{0.933374in}}%
\pgfpathlineto{\pgfqpoint{2.087358in}{0.933839in}}%
\pgfpathlineto{\pgfqpoint{2.092356in}{0.940891in}}%
\pgfpathlineto{\pgfqpoint{2.109858in}{0.929360in}}%
\pgfpathlineto{\pgfqpoint{2.115386in}{0.923776in}}%
\pgfpathlineto{\pgfqpoint{2.123559in}{0.920377in}}%
\pgfpathlineto{\pgfqpoint{2.123587in}{0.890672in}}%
\pgfpathlineto{\pgfqpoint{2.115003in}{0.895858in}}%
\pgfpathlineto{\pgfqpoint{2.100603in}{0.896104in}}%
\pgfpathlineto{\pgfqpoint{2.086328in}{0.894484in}}%
\pgfpathlineto{\pgfqpoint{2.086189in}{0.861574in}}%
\pgfpathlineto{\pgfqpoint{2.062420in}{0.861532in}}%
\pgfpathlineto{\pgfqpoint{2.064298in}{0.842428in}}%
\pgfpathlineto{\pgfqpoint{2.053670in}{0.846942in}}%
\pgfpathlineto{\pgfqpoint{2.050385in}{0.846426in}}%
\pgfpathlineto{\pgfqpoint{2.044379in}{0.851895in}}%
\pgfpathlineto{\pgfqpoint{2.034871in}{0.852102in}}%
\pgfpathclose%
\pgfusepath{fill}%
\end{pgfscope}%
\begin{pgfscope}%
\pgfpathrectangle{\pgfqpoint{0.100000in}{0.100000in}}{\pgfqpoint{3.608454in}{2.310000in}}%
\pgfusepath{clip}%
\pgfsetbuttcap%
\pgfsetmiterjoin%
\definecolor{currentfill}{rgb}{0.000000,0.956863,0.521569}%
\pgfsetfillcolor{currentfill}%
\pgfsetlinewidth{0.000000pt}%
\definecolor{currentstroke}{rgb}{0.000000,0.000000,0.000000}%
\pgfsetstrokecolor{currentstroke}%
\pgfsetstrokeopacity{0.000000}%
\pgfsetdash{}{0pt}%
\pgfpathmoveto{\pgfqpoint{2.961168in}{1.406181in}}%
\pgfpathlineto{\pgfqpoint{2.965023in}{1.393727in}}%
\pgfpathlineto{\pgfqpoint{2.962379in}{1.389071in}}%
\pgfpathlineto{\pgfqpoint{2.967663in}{1.383858in}}%
\pgfpathlineto{\pgfqpoint{2.965422in}{1.379505in}}%
\pgfpathlineto{\pgfqpoint{2.957191in}{1.367025in}}%
\pgfpathlineto{\pgfqpoint{2.940346in}{1.366252in}}%
\pgfpathlineto{\pgfqpoint{2.936140in}{1.370645in}}%
\pgfpathlineto{\pgfqpoint{2.933951in}{1.375085in}}%
\pgfpathlineto{\pgfqpoint{2.932362in}{1.392633in}}%
\pgfpathlineto{\pgfqpoint{2.936799in}{1.397031in}}%
\pgfpathlineto{\pgfqpoint{2.943680in}{1.394391in}}%
\pgfpathlineto{\pgfqpoint{2.954499in}{1.407993in}}%
\pgfpathlineto{\pgfqpoint{2.961168in}{1.406181in}}%
\pgfpathclose%
\pgfusepath{fill}%
\end{pgfscope}%
\begin{pgfscope}%
\pgfpathrectangle{\pgfqpoint{0.100000in}{0.100000in}}{\pgfqpoint{3.608454in}{2.310000in}}%
\pgfusepath{clip}%
\pgfsetbuttcap%
\pgfsetmiterjoin%
\definecolor{currentfill}{rgb}{0.000000,0.427451,0.786275}%
\pgfsetfillcolor{currentfill}%
\pgfsetlinewidth{0.000000pt}%
\definecolor{currentstroke}{rgb}{0.000000,0.000000,0.000000}%
\pgfsetstrokecolor{currentstroke}%
\pgfsetstrokeopacity{0.000000}%
\pgfsetdash{}{0pt}%
\pgfpathmoveto{\pgfqpoint{1.623124in}{2.077735in}}%
\pgfpathlineto{\pgfqpoint{1.626106in}{2.109098in}}%
\pgfpathlineto{\pgfqpoint{1.630877in}{2.157836in}}%
\pgfpathlineto{\pgfqpoint{1.689092in}{2.152335in}}%
\pgfpathlineto{\pgfqpoint{1.737269in}{2.148337in}}%
\pgfpathlineto{\pgfqpoint{1.736050in}{2.132965in}}%
\pgfpathlineto{\pgfqpoint{1.729180in}{2.133512in}}%
\pgfpathlineto{\pgfqpoint{1.728606in}{2.126565in}}%
\pgfpathlineto{\pgfqpoint{1.724331in}{2.126897in}}%
\pgfpathlineto{\pgfqpoint{1.723190in}{2.112972in}}%
\pgfpathlineto{\pgfqpoint{1.688717in}{2.115878in}}%
\pgfpathlineto{\pgfqpoint{1.687537in}{2.102006in}}%
\pgfpathlineto{\pgfqpoint{1.690580in}{2.101766in}}%
\pgfpathlineto{\pgfqpoint{1.686894in}{2.083116in}}%
\pgfpathlineto{\pgfqpoint{1.669040in}{2.084553in}}%
\pgfpathlineto{\pgfqpoint{1.664466in}{2.078762in}}%
\pgfpathlineto{\pgfqpoint{1.653010in}{2.075697in}}%
\pgfpathlineto{\pgfqpoint{1.648734in}{2.078432in}}%
\pgfpathlineto{\pgfqpoint{1.647283in}{2.084906in}}%
\pgfpathlineto{\pgfqpoint{1.643336in}{2.085862in}}%
\pgfpathlineto{\pgfqpoint{1.632154in}{2.074810in}}%
\pgfpathlineto{\pgfqpoint{1.623124in}{2.077735in}}%
\pgfpathclose%
\pgfusepath{fill}%
\end{pgfscope}%
\begin{pgfscope}%
\pgfpathrectangle{\pgfqpoint{0.100000in}{0.100000in}}{\pgfqpoint{3.608454in}{2.310000in}}%
\pgfusepath{clip}%
\pgfsetbuttcap%
\pgfsetmiterjoin%
\definecolor{currentfill}{rgb}{0.000000,0.505882,0.747059}%
\pgfsetfillcolor{currentfill}%
\pgfsetlinewidth{0.000000pt}%
\definecolor{currentstroke}{rgb}{0.000000,0.000000,0.000000}%
\pgfsetstrokecolor{currentstroke}%
\pgfsetstrokeopacity{0.000000}%
\pgfsetdash{}{0pt}%
\pgfpathmoveto{\pgfqpoint{1.576920in}{0.953485in}}%
\pgfpathlineto{\pgfqpoint{1.640100in}{0.948213in}}%
\pgfpathlineto{\pgfqpoint{1.708835in}{0.943953in}}%
\pgfpathlineto{\pgfqpoint{1.706689in}{0.909371in}}%
\pgfpathlineto{\pgfqpoint{1.704446in}{0.875520in}}%
\pgfpathlineto{\pgfqpoint{1.695447in}{0.875598in}}%
\pgfpathlineto{\pgfqpoint{1.635380in}{0.879674in}}%
\pgfpathlineto{\pgfqpoint{1.638056in}{0.913724in}}%
\pgfpathlineto{\pgfqpoint{1.573241in}{0.919040in}}%
\pgfpathlineto{\pgfqpoint{1.574789in}{0.933425in}}%
\pgfpathlineto{\pgfqpoint{1.576920in}{0.953485in}}%
\pgfpathclose%
\pgfusepath{fill}%
\end{pgfscope}%
\begin{pgfscope}%
\pgfpathrectangle{\pgfqpoint{0.100000in}{0.100000in}}{\pgfqpoint{3.608454in}{2.310000in}}%
\pgfusepath{clip}%
\pgfsetbuttcap%
\pgfsetmiterjoin%
\definecolor{currentfill}{rgb}{0.000000,0.584314,0.707843}%
\pgfsetfillcolor{currentfill}%
\pgfsetlinewidth{0.000000pt}%
\definecolor{currentstroke}{rgb}{0.000000,0.000000,0.000000}%
\pgfsetstrokecolor{currentstroke}%
\pgfsetstrokeopacity{0.000000}%
\pgfsetdash{}{0pt}%
\pgfpathmoveto{\pgfqpoint{3.145061in}{1.058883in}}%
\pgfpathlineto{\pgfqpoint{3.116266in}{1.079750in}}%
\pgfpathlineto{\pgfqpoint{3.099339in}{1.091621in}}%
\pgfpathlineto{\pgfqpoint{3.083124in}{1.089029in}}%
\pgfpathlineto{\pgfqpoint{3.086909in}{1.097540in}}%
\pgfpathlineto{\pgfqpoint{3.083344in}{1.107498in}}%
\pgfpathlineto{\pgfqpoint{3.078110in}{1.113287in}}%
\pgfpathlineto{\pgfqpoint{3.070169in}{1.111400in}}%
\pgfpathlineto{\pgfqpoint{3.070229in}{1.116768in}}%
\pgfpathlineto{\pgfqpoint{3.066320in}{1.122502in}}%
\pgfpathlineto{\pgfqpoint{3.067859in}{1.131866in}}%
\pgfpathlineto{\pgfqpoint{3.062964in}{1.138236in}}%
\pgfpathlineto{\pgfqpoint{3.057775in}{1.140860in}}%
\pgfpathlineto{\pgfqpoint{3.097477in}{1.148400in}}%
\pgfpathlineto{\pgfqpoint{3.110510in}{1.150879in}}%
\pgfpathlineto{\pgfqpoint{3.110689in}{1.147372in}}%
\pgfpathlineto{\pgfqpoint{3.121994in}{1.132796in}}%
\pgfpathlineto{\pgfqpoint{3.131128in}{1.128295in}}%
\pgfpathlineto{\pgfqpoint{3.149312in}{1.137485in}}%
\pgfpathlineto{\pgfqpoint{3.160885in}{1.137946in}}%
\pgfpathlineto{\pgfqpoint{3.160679in}{1.126465in}}%
\pgfpathlineto{\pgfqpoint{3.162834in}{1.117655in}}%
\pgfpathlineto{\pgfqpoint{3.174249in}{1.109020in}}%
\pgfpathlineto{\pgfqpoint{3.154787in}{1.104979in}}%
\pgfpathlineto{\pgfqpoint{3.148549in}{1.102631in}}%
\pgfpathlineto{\pgfqpoint{3.156648in}{1.092471in}}%
\pgfpathlineto{\pgfqpoint{3.156843in}{1.077749in}}%
\pgfpathlineto{\pgfqpoint{3.150261in}{1.071349in}}%
\pgfpathlineto{\pgfqpoint{3.145061in}{1.058883in}}%
\pgfpathclose%
\pgfusepath{fill}%
\end{pgfscope}%
\begin{pgfscope}%
\pgfpathrectangle{\pgfqpoint{0.100000in}{0.100000in}}{\pgfqpoint{3.608454in}{2.310000in}}%
\pgfusepath{clip}%
\pgfsetbuttcap%
\pgfsetmiterjoin%
\definecolor{currentfill}{rgb}{0.000000,0.435294,0.782353}%
\pgfsetfillcolor{currentfill}%
\pgfsetlinewidth{0.000000pt}%
\definecolor{currentstroke}{rgb}{0.000000,0.000000,0.000000}%
\pgfsetstrokecolor{currentstroke}%
\pgfsetstrokeopacity{0.000000}%
\pgfsetdash{}{0pt}%
\pgfpathmoveto{\pgfqpoint{1.964117in}{1.448246in}}%
\pgfpathlineto{\pgfqpoint{1.963356in}{1.420723in}}%
\pgfpathlineto{\pgfqpoint{1.908588in}{1.422459in}}%
\pgfpathlineto{\pgfqpoint{1.909610in}{1.449945in}}%
\pgfpathlineto{\pgfqpoint{1.936670in}{1.449037in}}%
\pgfpathlineto{\pgfqpoint{1.964117in}{1.448246in}}%
\pgfpathclose%
\pgfusepath{fill}%
\end{pgfscope}%
\begin{pgfscope}%
\pgfpathrectangle{\pgfqpoint{0.100000in}{0.100000in}}{\pgfqpoint{3.608454in}{2.310000in}}%
\pgfusepath{clip}%
\pgfsetbuttcap%
\pgfsetmiterjoin%
\definecolor{currentfill}{rgb}{0.000000,0.717647,0.641176}%
\pgfsetfillcolor{currentfill}%
\pgfsetlinewidth{0.000000pt}%
\definecolor{currentstroke}{rgb}{0.000000,0.000000,0.000000}%
\pgfsetstrokecolor{currentstroke}%
\pgfsetstrokeopacity{0.000000}%
\pgfsetdash{}{0pt}%
\pgfpathmoveto{\pgfqpoint{2.510006in}{1.121373in}}%
\pgfpathlineto{\pgfqpoint{2.498683in}{1.126546in}}%
\pgfpathlineto{\pgfqpoint{2.492139in}{1.131174in}}%
\pgfpathlineto{\pgfqpoint{2.479806in}{1.130609in}}%
\pgfpathlineto{\pgfqpoint{2.458843in}{1.129936in}}%
\pgfpathlineto{\pgfqpoint{2.462058in}{1.133851in}}%
\pgfpathlineto{\pgfqpoint{2.465696in}{1.145939in}}%
\pgfpathlineto{\pgfqpoint{2.461626in}{1.152803in}}%
\pgfpathlineto{\pgfqpoint{2.464384in}{1.163029in}}%
\pgfpathlineto{\pgfqpoint{2.473332in}{1.159111in}}%
\pgfpathlineto{\pgfqpoint{2.477166in}{1.167044in}}%
\pgfpathlineto{\pgfqpoint{2.474357in}{1.170803in}}%
\pgfpathlineto{\pgfqpoint{2.478809in}{1.176580in}}%
\pgfpathlineto{\pgfqpoint{2.498496in}{1.176925in}}%
\pgfpathlineto{\pgfqpoint{2.497579in}{1.190610in}}%
\pgfpathlineto{\pgfqpoint{2.518313in}{1.191718in}}%
\pgfpathlineto{\pgfqpoint{2.519321in}{1.176572in}}%
\pgfpathlineto{\pgfqpoint{2.520509in}{1.156888in}}%
\pgfpathlineto{\pgfqpoint{2.518781in}{1.156793in}}%
\pgfpathlineto{\pgfqpoint{2.519951in}{1.126253in}}%
\pgfpathlineto{\pgfqpoint{2.512080in}{1.125849in}}%
\pgfpathlineto{\pgfqpoint{2.510006in}{1.121373in}}%
\pgfpathclose%
\pgfusepath{fill}%
\end{pgfscope}%
\begin{pgfscope}%
\pgfpathrectangle{\pgfqpoint{0.100000in}{0.100000in}}{\pgfqpoint{3.608454in}{2.310000in}}%
\pgfusepath{clip}%
\pgfsetbuttcap%
\pgfsetmiterjoin%
\definecolor{currentfill}{rgb}{0.000000,0.474510,0.762745}%
\pgfsetfillcolor{currentfill}%
\pgfsetlinewidth{0.000000pt}%
\definecolor{currentstroke}{rgb}{0.000000,0.000000,0.000000}%
\pgfsetstrokecolor{currentstroke}%
\pgfsetstrokeopacity{0.000000}%
\pgfsetdash{}{0pt}%
\pgfpathmoveto{\pgfqpoint{1.401111in}{1.544070in}}%
\pgfpathlineto{\pgfqpoint{1.432841in}{1.539719in}}%
\pgfpathlineto{\pgfqpoint{1.436631in}{1.570812in}}%
\pgfpathlineto{\pgfqpoint{1.451617in}{1.568947in}}%
\pgfpathlineto{\pgfqpoint{1.457886in}{1.623314in}}%
\pgfpathlineto{\pgfqpoint{1.461309in}{1.650537in}}%
\pgfpathlineto{\pgfqpoint{1.491303in}{1.646739in}}%
\pgfpathlineto{\pgfqpoint{1.487410in}{1.642958in}}%
\pgfpathlineto{\pgfqpoint{1.486783in}{1.637214in}}%
\pgfpathlineto{\pgfqpoint{1.500514in}{1.634480in}}%
\pgfpathlineto{\pgfqpoint{1.504277in}{1.645283in}}%
\pgfpathlineto{\pgfqpoint{1.507108in}{1.644984in}}%
\pgfpathlineto{\pgfqpoint{1.503839in}{1.609878in}}%
\pgfpathlineto{\pgfqpoint{1.500517in}{1.584019in}}%
\pgfpathlineto{\pgfqpoint{1.494704in}{1.532257in}}%
\pgfpathlineto{\pgfqpoint{1.440561in}{1.538635in}}%
\pgfpathlineto{\pgfqpoint{1.440255in}{1.533604in}}%
\pgfpathlineto{\pgfqpoint{1.446700in}{1.523501in}}%
\pgfpathlineto{\pgfqpoint{1.455629in}{1.496066in}}%
\pgfpathlineto{\pgfqpoint{1.447426in}{1.486216in}}%
\pgfpathlineto{\pgfqpoint{1.444777in}{1.488383in}}%
\pgfpathlineto{\pgfqpoint{1.436560in}{1.486098in}}%
\pgfpathlineto{\pgfqpoint{1.426803in}{1.487168in}}%
\pgfpathlineto{\pgfqpoint{1.410993in}{1.493320in}}%
\pgfpathlineto{\pgfqpoint{1.407647in}{1.498770in}}%
\pgfpathlineto{\pgfqpoint{1.405998in}{1.509570in}}%
\pgfpathlineto{\pgfqpoint{1.412111in}{1.525836in}}%
\pgfpathlineto{\pgfqpoint{1.411585in}{1.530462in}}%
\pgfpathlineto{\pgfqpoint{1.400660in}{1.537925in}}%
\pgfpathlineto{\pgfqpoint{1.401111in}{1.544070in}}%
\pgfpathclose%
\pgfusepath{fill}%
\end{pgfscope}%
\begin{pgfscope}%
\pgfpathrectangle{\pgfqpoint{0.100000in}{0.100000in}}{\pgfqpoint{3.608454in}{2.310000in}}%
\pgfusepath{clip}%
\pgfsetbuttcap%
\pgfsetmiterjoin%
\definecolor{currentfill}{rgb}{0.000000,0.749020,0.625490}%
\pgfsetfillcolor{currentfill}%
\pgfsetlinewidth{0.000000pt}%
\definecolor{currentstroke}{rgb}{0.000000,0.000000,0.000000}%
\pgfsetstrokecolor{currentstroke}%
\pgfsetstrokeopacity{0.000000}%
\pgfsetdash{}{0pt}%
\pgfpathmoveto{\pgfqpoint{2.815217in}{1.759142in}}%
\pgfpathlineto{\pgfqpoint{2.786333in}{1.754076in}}%
\pgfpathlineto{\pgfqpoint{2.790730in}{1.726819in}}%
\pgfpathlineto{\pgfqpoint{2.776445in}{1.724555in}}%
\pgfpathlineto{\pgfqpoint{2.777478in}{1.717721in}}%
\pgfpathlineto{\pgfqpoint{2.771368in}{1.715974in}}%
\pgfpathlineto{\pgfqpoint{2.744382in}{1.712367in}}%
\pgfpathlineto{\pgfqpoint{2.745404in}{1.705512in}}%
\pgfpathlineto{\pgfqpoint{2.720255in}{1.702142in}}%
\pgfpathlineto{\pgfqpoint{2.717010in}{1.728719in}}%
\pgfpathlineto{\pgfqpoint{2.703428in}{1.727130in}}%
\pgfpathlineto{\pgfqpoint{2.697115in}{1.781860in}}%
\pgfpathlineto{\pgfqpoint{2.738013in}{1.786827in}}%
\pgfpathlineto{\pgfqpoint{2.756041in}{1.789197in}}%
\pgfpathlineto{\pgfqpoint{2.756426in}{1.784651in}}%
\pgfpathlineto{\pgfqpoint{2.746228in}{1.773703in}}%
\pgfpathlineto{\pgfqpoint{2.742788in}{1.773732in}}%
\pgfpathlineto{\pgfqpoint{2.738774in}{1.765761in}}%
\pgfpathlineto{\pgfqpoint{2.738244in}{1.753504in}}%
\pgfpathlineto{\pgfqpoint{2.741189in}{1.748238in}}%
\pgfpathlineto{\pgfqpoint{2.756060in}{1.743595in}}%
\pgfpathlineto{\pgfqpoint{2.775084in}{1.776539in}}%
\pgfpathlineto{\pgfqpoint{2.786967in}{1.781353in}}%
\pgfpathlineto{\pgfqpoint{2.792933in}{1.786497in}}%
\pgfpathlineto{\pgfqpoint{2.804037in}{1.782140in}}%
\pgfpathlineto{\pgfqpoint{2.812223in}{1.770016in}}%
\pgfpathlineto{\pgfqpoint{2.815217in}{1.759142in}}%
\pgfpathclose%
\pgfusepath{fill}%
\end{pgfscope}%
\begin{pgfscope}%
\pgfpathrectangle{\pgfqpoint{0.100000in}{0.100000in}}{\pgfqpoint{3.608454in}{2.310000in}}%
\pgfusepath{clip}%
\pgfsetbuttcap%
\pgfsetmiterjoin%
\definecolor{currentfill}{rgb}{0.000000,0.701961,0.649020}%
\pgfsetfillcolor{currentfill}%
\pgfsetlinewidth{0.000000pt}%
\definecolor{currentstroke}{rgb}{0.000000,0.000000,0.000000}%
\pgfsetstrokecolor{currentstroke}%
\pgfsetstrokeopacity{0.000000}%
\pgfsetdash{}{0pt}%
\pgfpathmoveto{\pgfqpoint{3.050650in}{1.387261in}}%
\pgfpathlineto{\pgfqpoint{3.050299in}{1.374365in}}%
\pgfpathlineto{\pgfqpoint{3.044609in}{1.363268in}}%
\pgfpathlineto{\pgfqpoint{3.046369in}{1.360627in}}%
\pgfpathlineto{\pgfqpoint{3.039437in}{1.352438in}}%
\pgfpathlineto{\pgfqpoint{3.038592in}{1.342913in}}%
\pgfpathlineto{\pgfqpoint{3.026105in}{1.338632in}}%
\pgfpathlineto{\pgfqpoint{3.020093in}{1.338488in}}%
\pgfpathlineto{\pgfqpoint{3.013175in}{1.342782in}}%
\pgfpathlineto{\pgfqpoint{3.011961in}{1.351516in}}%
\pgfpathlineto{\pgfqpoint{3.012394in}{1.360390in}}%
\pgfpathlineto{\pgfqpoint{3.016902in}{1.365346in}}%
\pgfpathlineto{\pgfqpoint{3.018910in}{1.376644in}}%
\pgfpathlineto{\pgfqpoint{3.010884in}{1.388904in}}%
\pgfpathlineto{\pgfqpoint{3.012891in}{1.393149in}}%
\pgfpathlineto{\pgfqpoint{3.019598in}{1.394420in}}%
\pgfpathlineto{\pgfqpoint{3.018667in}{1.406375in}}%
\pgfpathlineto{\pgfqpoint{3.021096in}{1.413057in}}%
\pgfpathlineto{\pgfqpoint{3.008758in}{1.422029in}}%
\pgfpathlineto{\pgfqpoint{3.010773in}{1.432733in}}%
\pgfpathlineto{\pgfqpoint{3.019335in}{1.434599in}}%
\pgfpathlineto{\pgfqpoint{3.025784in}{1.440073in}}%
\pgfpathlineto{\pgfqpoint{3.032463in}{1.435406in}}%
\pgfpathlineto{\pgfqpoint{3.039311in}{1.439843in}}%
\pgfpathlineto{\pgfqpoint{3.052407in}{1.435951in}}%
\pgfpathlineto{\pgfqpoint{3.060760in}{1.438938in}}%
\pgfpathlineto{\pgfqpoint{3.063986in}{1.437446in}}%
\pgfpathlineto{\pgfqpoint{3.062165in}{1.429526in}}%
\pgfpathlineto{\pgfqpoint{3.065468in}{1.428454in}}%
\pgfpathlineto{\pgfqpoint{3.063909in}{1.418899in}}%
\pgfpathlineto{\pgfqpoint{3.055324in}{1.410700in}}%
\pgfpathlineto{\pgfqpoint{3.056652in}{1.404504in}}%
\pgfpathlineto{\pgfqpoint{3.050650in}{1.387261in}}%
\pgfpathclose%
\pgfusepath{fill}%
\end{pgfscope}%
\begin{pgfscope}%
\pgfpathrectangle{\pgfqpoint{0.100000in}{0.100000in}}{\pgfqpoint{3.608454in}{2.310000in}}%
\pgfusepath{clip}%
\pgfsetbuttcap%
\pgfsetmiterjoin%
\definecolor{currentfill}{rgb}{0.000000,0.600000,0.700000}%
\pgfsetfillcolor{currentfill}%
\pgfsetlinewidth{0.000000pt}%
\definecolor{currentstroke}{rgb}{0.000000,0.000000,0.000000}%
\pgfsetstrokecolor{currentstroke}%
\pgfsetstrokeopacity{0.000000}%
\pgfsetdash{}{0pt}%
\pgfpathmoveto{\pgfqpoint{3.217063in}{1.727096in}}%
\pgfpathlineto{\pgfqpoint{3.194930in}{1.721757in}}%
\pgfpathlineto{\pgfqpoint{3.190729in}{1.738354in}}%
\pgfpathlineto{\pgfqpoint{3.167852in}{1.733612in}}%
\pgfpathlineto{\pgfqpoint{3.162874in}{1.739825in}}%
\pgfpathlineto{\pgfqpoint{3.158361in}{1.759713in}}%
\pgfpathlineto{\pgfqpoint{3.158574in}{1.768645in}}%
\pgfpathlineto{\pgfqpoint{3.153298in}{1.788862in}}%
\pgfpathlineto{\pgfqpoint{3.157336in}{1.795352in}}%
\pgfpathlineto{\pgfqpoint{3.167511in}{1.806116in}}%
\pgfpathlineto{\pgfqpoint{3.177989in}{1.809166in}}%
\pgfpathlineto{\pgfqpoint{3.177265in}{1.821081in}}%
\pgfpathlineto{\pgfqpoint{3.186645in}{1.825233in}}%
\pgfpathlineto{\pgfqpoint{3.200994in}{1.826823in}}%
\pgfpathlineto{\pgfqpoint{3.205664in}{1.810106in}}%
\pgfpathlineto{\pgfqpoint{3.219039in}{1.808893in}}%
\pgfpathlineto{\pgfqpoint{3.239480in}{1.829258in}}%
\pgfpathlineto{\pgfqpoint{3.227847in}{1.865817in}}%
\pgfpathlineto{\pgfqpoint{3.234624in}{1.863539in}}%
\pgfpathlineto{\pgfqpoint{3.245888in}{1.867665in}}%
\pgfpathlineto{\pgfqpoint{3.260468in}{1.823521in}}%
\pgfpathlineto{\pgfqpoint{3.257858in}{1.810998in}}%
\pgfpathlineto{\pgfqpoint{3.267521in}{1.808820in}}%
\pgfpathlineto{\pgfqpoint{3.270444in}{1.800501in}}%
\pgfpathlineto{\pgfqpoint{3.267628in}{1.792368in}}%
\pgfpathlineto{\pgfqpoint{3.271017in}{1.786225in}}%
\pgfpathlineto{\pgfqpoint{3.272069in}{1.775491in}}%
\pgfpathlineto{\pgfqpoint{3.264917in}{1.776723in}}%
\pgfpathlineto{\pgfqpoint{3.264610in}{1.770699in}}%
\pgfpathlineto{\pgfqpoint{3.252161in}{1.774686in}}%
\pgfpathlineto{\pgfqpoint{3.246838in}{1.769301in}}%
\pgfpathlineto{\pgfqpoint{3.243936in}{1.759541in}}%
\pgfpathlineto{\pgfqpoint{3.210494in}{1.750645in}}%
\pgfpathlineto{\pgfqpoint{3.217063in}{1.727096in}}%
\pgfpathclose%
\pgfusepath{fill}%
\end{pgfscope}%
\begin{pgfscope}%
\pgfpathrectangle{\pgfqpoint{0.100000in}{0.100000in}}{\pgfqpoint{3.608454in}{2.310000in}}%
\pgfusepath{clip}%
\pgfsetbuttcap%
\pgfsetmiterjoin%
\definecolor{currentfill}{rgb}{0.000000,0.443137,0.778431}%
\pgfsetfillcolor{currentfill}%
\pgfsetlinewidth{0.000000pt}%
\definecolor{currentstroke}{rgb}{0.000000,0.000000,0.000000}%
\pgfsetstrokecolor{currentstroke}%
\pgfsetstrokeopacity{0.000000}%
\pgfsetdash{}{0pt}%
\pgfpathmoveto{\pgfqpoint{2.115090in}{1.593467in}}%
\pgfpathlineto{\pgfqpoint{2.155774in}{1.593709in}}%
\pgfpathlineto{\pgfqpoint{2.156052in}{1.566272in}}%
\pgfpathlineto{\pgfqpoint{2.121832in}{1.566001in}}%
\pgfpathlineto{\pgfqpoint{2.073952in}{1.566162in}}%
\pgfpathlineto{\pgfqpoint{2.067148in}{1.566188in}}%
\pgfpathlineto{\pgfqpoint{2.067449in}{1.593735in}}%
\pgfpathlineto{\pgfqpoint{2.115090in}{1.593467in}}%
\pgfpathclose%
\pgfusepath{fill}%
\end{pgfscope}%
\begin{pgfscope}%
\pgfpathrectangle{\pgfqpoint{0.100000in}{0.100000in}}{\pgfqpoint{3.608454in}{2.310000in}}%
\pgfusepath{clip}%
\pgfsetbuttcap%
\pgfsetmiterjoin%
\definecolor{currentfill}{rgb}{0.000000,0.603922,0.698039}%
\pgfsetfillcolor{currentfill}%
\pgfsetlinewidth{0.000000pt}%
\definecolor{currentstroke}{rgb}{0.000000,0.000000,0.000000}%
\pgfsetstrokecolor{currentstroke}%
\pgfsetstrokeopacity{0.000000}%
\pgfsetdash{}{0pt}%
\pgfpathmoveto{\pgfqpoint{2.647203in}{0.861864in}}%
\pgfpathlineto{\pgfqpoint{2.648743in}{0.852393in}}%
\pgfpathlineto{\pgfqpoint{2.656320in}{0.836856in}}%
\pgfpathlineto{\pgfqpoint{2.650875in}{0.827115in}}%
\pgfpathlineto{\pgfqpoint{2.652230in}{0.813085in}}%
\pgfpathlineto{\pgfqpoint{2.634007in}{0.811449in}}%
\pgfpathlineto{\pgfqpoint{2.622706in}{0.822514in}}%
\pgfpathlineto{\pgfqpoint{2.612796in}{0.827022in}}%
\pgfpathlineto{\pgfqpoint{2.609026in}{0.830093in}}%
\pgfpathlineto{\pgfqpoint{2.607870in}{0.843949in}}%
\pgfpathlineto{\pgfqpoint{2.594215in}{0.842750in}}%
\pgfpathlineto{\pgfqpoint{2.588335in}{0.845764in}}%
\pgfpathlineto{\pgfqpoint{2.589627in}{0.852306in}}%
\pgfpathlineto{\pgfqpoint{2.586111in}{0.864841in}}%
\pgfpathlineto{\pgfqpoint{2.586494in}{0.875360in}}%
\pgfpathlineto{\pgfqpoint{2.591160in}{0.879866in}}%
\pgfpathlineto{\pgfqpoint{2.591982in}{0.885391in}}%
\pgfpathlineto{\pgfqpoint{2.583933in}{0.884712in}}%
\pgfpathlineto{\pgfqpoint{2.582729in}{0.895488in}}%
\pgfpathlineto{\pgfqpoint{2.580198in}{0.924865in}}%
\pgfpathlineto{\pgfqpoint{2.591702in}{0.925524in}}%
\pgfpathlineto{\pgfqpoint{2.593439in}{0.932696in}}%
\pgfpathlineto{\pgfqpoint{2.610666in}{0.933472in}}%
\pgfpathlineto{\pgfqpoint{2.618130in}{0.927031in}}%
\pgfpathlineto{\pgfqpoint{2.617467in}{0.921208in}}%
\pgfpathlineto{\pgfqpoint{2.625197in}{0.912348in}}%
\pgfpathlineto{\pgfqpoint{2.635819in}{0.907320in}}%
\pgfpathlineto{\pgfqpoint{2.636464in}{0.900943in}}%
\pgfpathlineto{\pgfqpoint{2.640666in}{0.895938in}}%
\pgfpathlineto{\pgfqpoint{2.647138in}{0.892648in}}%
\pgfpathlineto{\pgfqpoint{2.648704in}{0.875737in}}%
\pgfpathlineto{\pgfqpoint{2.639182in}{0.874953in}}%
\pgfpathlineto{\pgfqpoint{2.640493in}{0.861198in}}%
\pgfpathlineto{\pgfqpoint{2.647203in}{0.861864in}}%
\pgfpathclose%
\pgfusepath{fill}%
\end{pgfscope}%
\begin{pgfscope}%
\pgfpathrectangle{\pgfqpoint{0.100000in}{0.100000in}}{\pgfqpoint{3.608454in}{2.310000in}}%
\pgfusepath{clip}%
\pgfsetbuttcap%
\pgfsetmiterjoin%
\definecolor{currentfill}{rgb}{0.000000,0.321569,0.839216}%
\pgfsetfillcolor{currentfill}%
\pgfsetlinewidth{0.000000pt}%
\definecolor{currentstroke}{rgb}{0.000000,0.000000,0.000000}%
\pgfsetstrokecolor{currentstroke}%
\pgfsetstrokeopacity{0.000000}%
\pgfsetdash{}{0pt}%
\pgfpathmoveto{\pgfqpoint{2.223956in}{1.594717in}}%
\pgfpathlineto{\pgfqpoint{2.210467in}{1.594411in}}%
\pgfpathlineto{\pgfqpoint{2.196902in}{1.594193in}}%
\pgfpathlineto{\pgfqpoint{2.196575in}{1.614898in}}%
\pgfpathlineto{\pgfqpoint{2.194446in}{1.614873in}}%
\pgfpathlineto{\pgfqpoint{2.193901in}{1.649520in}}%
\pgfpathlineto{\pgfqpoint{2.221382in}{1.649990in}}%
\pgfpathlineto{\pgfqpoint{2.222007in}{1.615392in}}%
\pgfpathlineto{\pgfqpoint{2.223427in}{1.615413in}}%
\pgfpathlineto{\pgfqpoint{2.223956in}{1.594717in}}%
\pgfpathclose%
\pgfusepath{fill}%
\end{pgfscope}%
\begin{pgfscope}%
\pgfpathrectangle{\pgfqpoint{0.100000in}{0.100000in}}{\pgfqpoint{3.608454in}{2.310000in}}%
\pgfusepath{clip}%
\pgfsetbuttcap%
\pgfsetmiterjoin%
\definecolor{currentfill}{rgb}{0.000000,0.694118,0.652941}%
\pgfsetfillcolor{currentfill}%
\pgfsetlinewidth{0.000000pt}%
\definecolor{currentstroke}{rgb}{0.000000,0.000000,0.000000}%
\pgfsetstrokecolor{currentstroke}%
\pgfsetstrokeopacity{0.000000}%
\pgfsetdash{}{0pt}%
\pgfpathmoveto{\pgfqpoint{2.720255in}{1.702142in}}%
\pgfpathlineto{\pgfqpoint{2.720350in}{1.701323in}}%
\pgfpathlineto{\pgfqpoint{2.693396in}{1.698299in}}%
\pgfpathlineto{\pgfqpoint{2.689796in}{1.725591in}}%
\pgfpathlineto{\pgfqpoint{2.676133in}{1.724089in}}%
\pgfpathlineto{\pgfqpoint{2.669984in}{1.779116in}}%
\pgfpathlineto{\pgfqpoint{2.683319in}{1.780357in}}%
\pgfpathlineto{\pgfqpoint{2.697115in}{1.781860in}}%
\pgfpathlineto{\pgfqpoint{2.703428in}{1.727130in}}%
\pgfpathlineto{\pgfqpoint{2.717010in}{1.728719in}}%
\pgfpathlineto{\pgfqpoint{2.720255in}{1.702142in}}%
\pgfpathclose%
\pgfusepath{fill}%
\end{pgfscope}%
\begin{pgfscope}%
\pgfpathrectangle{\pgfqpoint{0.100000in}{0.100000in}}{\pgfqpoint{3.608454in}{2.310000in}}%
\pgfusepath{clip}%
\pgfsetbuttcap%
\pgfsetmiterjoin%
\definecolor{currentfill}{rgb}{0.000000,0.635294,0.682353}%
\pgfsetfillcolor{currentfill}%
\pgfsetlinewidth{0.000000pt}%
\definecolor{currentstroke}{rgb}{0.000000,0.000000,0.000000}%
\pgfsetstrokecolor{currentstroke}%
\pgfsetstrokeopacity{0.000000}%
\pgfsetdash{}{0pt}%
\pgfpathmoveto{\pgfqpoint{2.882548in}{1.076650in}}%
\pgfpathlineto{\pgfqpoint{2.879644in}{1.078954in}}%
\pgfpathlineto{\pgfqpoint{2.885380in}{1.093119in}}%
\pgfpathlineto{\pgfqpoint{2.884882in}{1.098977in}}%
\pgfpathlineto{\pgfqpoint{2.866841in}{1.114054in}}%
\pgfpathlineto{\pgfqpoint{2.864797in}{1.125728in}}%
\pgfpathlineto{\pgfqpoint{2.859932in}{1.127064in}}%
\pgfpathlineto{\pgfqpoint{2.865280in}{1.133278in}}%
\pgfpathlineto{\pgfqpoint{2.877093in}{1.136816in}}%
\pgfpathlineto{\pgfqpoint{2.879988in}{1.149617in}}%
\pgfpathlineto{\pgfqpoint{2.895600in}{1.161365in}}%
\pgfpathlineto{\pgfqpoint{2.898012in}{1.153842in}}%
\pgfpathlineto{\pgfqpoint{2.904596in}{1.155570in}}%
\pgfpathlineto{\pgfqpoint{2.905552in}{1.151887in}}%
\pgfpathlineto{\pgfqpoint{2.916363in}{1.143909in}}%
\pgfpathlineto{\pgfqpoint{2.925369in}{1.125057in}}%
\pgfpathlineto{\pgfqpoint{2.930944in}{1.123301in}}%
\pgfpathlineto{\pgfqpoint{2.926944in}{1.122067in}}%
\pgfpathlineto{\pgfqpoint{2.924776in}{1.115579in}}%
\pgfpathlineto{\pgfqpoint{2.926537in}{1.111891in}}%
\pgfpathlineto{\pgfqpoint{2.922309in}{1.102781in}}%
\pgfpathlineto{\pgfqpoint{2.922867in}{1.095204in}}%
\pgfpathlineto{\pgfqpoint{2.903153in}{1.086612in}}%
\pgfpathlineto{\pgfqpoint{2.882548in}{1.076650in}}%
\pgfpathclose%
\pgfusepath{fill}%
\end{pgfscope}%
\begin{pgfscope}%
\pgfpathrectangle{\pgfqpoint{0.100000in}{0.100000in}}{\pgfqpoint{3.608454in}{2.310000in}}%
\pgfusepath{clip}%
\pgfsetbuttcap%
\pgfsetmiterjoin%
\definecolor{currentfill}{rgb}{0.000000,0.764706,0.617647}%
\pgfsetfillcolor{currentfill}%
\pgfsetlinewidth{0.000000pt}%
\definecolor{currentstroke}{rgb}{0.000000,0.000000,0.000000}%
\pgfsetstrokecolor{currentstroke}%
\pgfsetstrokeopacity{0.000000}%
\pgfsetdash{}{0pt}%
\pgfpathmoveto{\pgfqpoint{2.272396in}{1.102853in}}%
\pgfpathlineto{\pgfqpoint{2.272664in}{1.088030in}}%
\pgfpathlineto{\pgfqpoint{2.258979in}{1.087834in}}%
\pgfpathlineto{\pgfqpoint{2.259040in}{1.082097in}}%
\pgfpathlineto{\pgfqpoint{2.247558in}{1.082128in}}%
\pgfpathlineto{\pgfqpoint{2.201928in}{1.082158in}}%
\pgfpathlineto{\pgfqpoint{2.201915in}{1.084462in}}%
\pgfpathlineto{\pgfqpoint{2.201478in}{1.093804in}}%
\pgfpathlineto{\pgfqpoint{2.206038in}{1.100680in}}%
\pgfpathlineto{\pgfqpoint{2.206212in}{1.109649in}}%
\pgfpathlineto{\pgfqpoint{2.199417in}{1.113130in}}%
\pgfpathlineto{\pgfqpoint{2.194911in}{1.119004in}}%
\pgfpathlineto{\pgfqpoint{2.192391in}{1.127027in}}%
\pgfpathlineto{\pgfqpoint{2.178906in}{1.127188in}}%
\pgfpathlineto{\pgfqpoint{2.178864in}{1.142325in}}%
\pgfpathlineto{\pgfqpoint{2.243255in}{1.143432in}}%
\pgfpathlineto{\pgfqpoint{2.241311in}{1.142821in}}%
\pgfpathlineto{\pgfqpoint{2.241638in}{1.113006in}}%
\pgfpathlineto{\pgfqpoint{2.245103in}{1.109604in}}%
\pgfpathlineto{\pgfqpoint{2.272381in}{1.109733in}}%
\pgfpathlineto{\pgfqpoint{2.272396in}{1.102853in}}%
\pgfpathclose%
\pgfusepath{fill}%
\end{pgfscope}%
\begin{pgfscope}%
\pgfpathrectangle{\pgfqpoint{0.100000in}{0.100000in}}{\pgfqpoint{3.608454in}{2.310000in}}%
\pgfusepath{clip}%
\pgfsetbuttcap%
\pgfsetmiterjoin%
\definecolor{currentfill}{rgb}{0.000000,0.447059,0.776471}%
\pgfsetfillcolor{currentfill}%
\pgfsetlinewidth{0.000000pt}%
\definecolor{currentstroke}{rgb}{0.000000,0.000000,0.000000}%
\pgfsetstrokecolor{currentstroke}%
\pgfsetstrokeopacity{0.000000}%
\pgfsetdash{}{0pt}%
\pgfpathmoveto{\pgfqpoint{2.526558in}{1.632431in}}%
\pgfpathlineto{\pgfqpoint{2.473885in}{1.628654in}}%
\pgfpathlineto{\pgfqpoint{2.435480in}{1.626662in}}%
\pgfpathlineto{\pgfqpoint{2.433615in}{1.653903in}}%
\pgfpathlineto{\pgfqpoint{2.454210in}{1.655405in}}%
\pgfpathlineto{\pgfqpoint{2.467930in}{1.655975in}}%
\pgfpathlineto{\pgfqpoint{2.522980in}{1.660007in}}%
\pgfpathlineto{\pgfqpoint{2.527343in}{1.655339in}}%
\pgfpathlineto{\pgfqpoint{2.524832in}{1.642175in}}%
\pgfpathlineto{\pgfqpoint{2.526558in}{1.632431in}}%
\pgfpathclose%
\pgfusepath{fill}%
\end{pgfscope}%
\begin{pgfscope}%
\pgfpathrectangle{\pgfqpoint{0.100000in}{0.100000in}}{\pgfqpoint{3.608454in}{2.310000in}}%
\pgfusepath{clip}%
\pgfsetbuttcap%
\pgfsetmiterjoin%
\definecolor{currentfill}{rgb}{0.000000,0.501961,0.749020}%
\pgfsetfillcolor{currentfill}%
\pgfsetlinewidth{0.000000pt}%
\definecolor{currentstroke}{rgb}{0.000000,0.000000,0.000000}%
\pgfsetstrokecolor{currentstroke}%
\pgfsetstrokeopacity{0.000000}%
\pgfsetdash{}{0pt}%
\pgfpathmoveto{\pgfqpoint{1.774163in}{0.940149in}}%
\pgfpathlineto{\pgfqpoint{1.743217in}{0.941925in}}%
\pgfpathlineto{\pgfqpoint{1.745610in}{0.979809in}}%
\pgfpathlineto{\pgfqpoint{1.752117in}{0.979371in}}%
\pgfpathlineto{\pgfqpoint{1.754289in}{1.013713in}}%
\pgfpathlineto{\pgfqpoint{1.781300in}{1.012079in}}%
\pgfpathlineto{\pgfqpoint{1.780677in}{0.997355in}}%
\pgfpathlineto{\pgfqpoint{1.779205in}{0.970733in}}%
\pgfpathlineto{\pgfqpoint{1.775993in}{0.971342in}}%
\pgfpathlineto{\pgfqpoint{1.774163in}{0.940149in}}%
\pgfpathclose%
\pgfusepath{fill}%
\end{pgfscope}%
\begin{pgfscope}%
\pgfpathrectangle{\pgfqpoint{0.100000in}{0.100000in}}{\pgfqpoint{3.608454in}{2.310000in}}%
\pgfusepath{clip}%
\pgfsetbuttcap%
\pgfsetmiterjoin%
\definecolor{currentfill}{rgb}{0.000000,0.811765,0.594118}%
\pgfsetfillcolor{currentfill}%
\pgfsetlinewidth{0.000000pt}%
\definecolor{currentstroke}{rgb}{0.000000,0.000000,0.000000}%
\pgfsetstrokecolor{currentstroke}%
\pgfsetstrokeopacity{0.000000}%
\pgfsetdash{}{0pt}%
\pgfpathmoveto{\pgfqpoint{2.338322in}{0.779971in}}%
\pgfpathlineto{\pgfqpoint{2.338396in}{0.773253in}}%
\pgfpathlineto{\pgfqpoint{2.334485in}{0.767523in}}%
\pgfpathlineto{\pgfqpoint{2.330853in}{0.768113in}}%
\pgfpathlineto{\pgfqpoint{2.325362in}{0.751678in}}%
\pgfpathlineto{\pgfqpoint{2.319974in}{0.746106in}}%
\pgfpathlineto{\pgfqpoint{2.322557in}{0.730252in}}%
\pgfpathlineto{\pgfqpoint{2.311015in}{0.733955in}}%
\pgfpathlineto{\pgfqpoint{2.300756in}{0.740640in}}%
\pgfpathlineto{\pgfqpoint{2.299838in}{0.745415in}}%
\pgfpathlineto{\pgfqpoint{2.291967in}{0.750029in}}%
\pgfpathlineto{\pgfqpoint{2.286412in}{0.757594in}}%
\pgfpathlineto{\pgfqpoint{2.285332in}{0.771355in}}%
\pgfpathlineto{\pgfqpoint{2.288408in}{0.781863in}}%
\pgfpathlineto{\pgfqpoint{2.315987in}{0.782756in}}%
\pgfpathlineto{\pgfqpoint{2.319400in}{0.785321in}}%
\pgfpathlineto{\pgfqpoint{2.324523in}{0.778629in}}%
\pgfpathlineto{\pgfqpoint{2.336629in}{0.786857in}}%
\pgfpathlineto{\pgfqpoint{2.338322in}{0.779971in}}%
\pgfpathclose%
\pgfusepath{fill}%
\end{pgfscope}%
\begin{pgfscope}%
\pgfpathrectangle{\pgfqpoint{0.100000in}{0.100000in}}{\pgfqpoint{3.608454in}{2.310000in}}%
\pgfusepath{clip}%
\pgfsetbuttcap%
\pgfsetmiterjoin%
\definecolor{currentfill}{rgb}{0.000000,0.654902,0.672549}%
\pgfsetfillcolor{currentfill}%
\pgfsetlinewidth{0.000000pt}%
\definecolor{currentstroke}{rgb}{0.000000,0.000000,0.000000}%
\pgfsetstrokecolor{currentstroke}%
\pgfsetstrokeopacity{0.000000}%
\pgfsetdash{}{0pt}%
\pgfpathmoveto{\pgfqpoint{2.618130in}{0.927031in}}%
\pgfpathlineto{\pgfqpoint{2.610666in}{0.933472in}}%
\pgfpathlineto{\pgfqpoint{2.593439in}{0.932696in}}%
\pgfpathlineto{\pgfqpoint{2.591702in}{0.925524in}}%
\pgfpathlineto{\pgfqpoint{2.580198in}{0.924865in}}%
\pgfpathlineto{\pgfqpoint{2.551506in}{0.923399in}}%
\pgfpathlineto{\pgfqpoint{2.551959in}{0.940228in}}%
\pgfpathlineto{\pgfqpoint{2.553484in}{0.986177in}}%
\pgfpathlineto{\pgfqpoint{2.595636in}{0.988190in}}%
\pgfpathlineto{\pgfqpoint{2.623076in}{0.990093in}}%
\pgfpathlineto{\pgfqpoint{2.625447in}{0.963052in}}%
\pgfpathlineto{\pgfqpoint{2.627618in}{0.959226in}}%
\pgfpathlineto{\pgfqpoint{2.635731in}{0.956087in}}%
\pgfpathlineto{\pgfqpoint{2.632392in}{0.948970in}}%
\pgfpathlineto{\pgfqpoint{2.631914in}{0.941959in}}%
\pgfpathlineto{\pgfqpoint{2.628827in}{0.937113in}}%
\pgfpathlineto{\pgfqpoint{2.623213in}{0.935505in}}%
\pgfpathlineto{\pgfqpoint{2.618130in}{0.927031in}}%
\pgfpathclose%
\pgfusepath{fill}%
\end{pgfscope}%
\begin{pgfscope}%
\pgfpathrectangle{\pgfqpoint{0.100000in}{0.100000in}}{\pgfqpoint{3.608454in}{2.310000in}}%
\pgfusepath{clip}%
\pgfsetbuttcap%
\pgfsetmiterjoin%
\definecolor{currentfill}{rgb}{0.000000,0.525490,0.737255}%
\pgfsetfillcolor{currentfill}%
\pgfsetlinewidth{0.000000pt}%
\definecolor{currentstroke}{rgb}{0.000000,0.000000,0.000000}%
\pgfsetstrokecolor{currentstroke}%
\pgfsetstrokeopacity{0.000000}%
\pgfsetdash{}{0pt}%
\pgfpathmoveto{\pgfqpoint{2.949055in}{1.160645in}}%
\pgfpathlineto{\pgfqpoint{2.956204in}{1.154739in}}%
\pgfpathlineto{\pgfqpoint{2.967129in}{1.148531in}}%
\pgfpathlineto{\pgfqpoint{2.973998in}{1.150189in}}%
\pgfpathlineto{\pgfqpoint{2.978683in}{1.148660in}}%
\pgfpathlineto{\pgfqpoint{2.970938in}{1.132406in}}%
\pgfpathlineto{\pgfqpoint{2.963707in}{1.130755in}}%
\pgfpathlineto{\pgfqpoint{2.956399in}{1.133737in}}%
\pgfpathlineto{\pgfqpoint{2.951598in}{1.127420in}}%
\pgfpathlineto{\pgfqpoint{2.943450in}{1.124830in}}%
\pgfpathlineto{\pgfqpoint{2.930944in}{1.123301in}}%
\pgfpathlineto{\pgfqpoint{2.925369in}{1.125057in}}%
\pgfpathlineto{\pgfqpoint{2.916363in}{1.143909in}}%
\pgfpathlineto{\pgfqpoint{2.905552in}{1.151887in}}%
\pgfpathlineto{\pgfqpoint{2.904596in}{1.155570in}}%
\pgfpathlineto{\pgfqpoint{2.909797in}{1.165077in}}%
\pgfpathlineto{\pgfqpoint{2.920582in}{1.172131in}}%
\pgfpathlineto{\pgfqpoint{2.930144in}{1.169183in}}%
\pgfpathlineto{\pgfqpoint{2.934882in}{1.156995in}}%
\pgfpathlineto{\pgfqpoint{2.938556in}{1.155078in}}%
\pgfpathlineto{\pgfqpoint{2.942760in}{1.162773in}}%
\pgfpathlineto{\pgfqpoint{2.949055in}{1.160645in}}%
\pgfpathclose%
\pgfusepath{fill}%
\end{pgfscope}%
\begin{pgfscope}%
\pgfpathrectangle{\pgfqpoint{0.100000in}{0.100000in}}{\pgfqpoint{3.608454in}{2.310000in}}%
\pgfusepath{clip}%
\pgfsetbuttcap%
\pgfsetmiterjoin%
\definecolor{currentfill}{rgb}{0.000000,0.603922,0.698039}%
\pgfsetfillcolor{currentfill}%
\pgfsetlinewidth{0.000000pt}%
\definecolor{currentstroke}{rgb}{0.000000,0.000000,0.000000}%
\pgfsetstrokecolor{currentstroke}%
\pgfsetstrokeopacity{0.000000}%
\pgfsetdash{}{0pt}%
\pgfpathmoveto{\pgfqpoint{0.650466in}{2.312218in}}%
\pgfpathlineto{\pgfqpoint{0.640068in}{2.318805in}}%
\pgfpathlineto{\pgfqpoint{0.635270in}{2.326403in}}%
\pgfpathlineto{\pgfqpoint{0.636669in}{2.333571in}}%
\pgfpathlineto{\pgfqpoint{0.643160in}{2.329044in}}%
\pgfpathlineto{\pgfqpoint{0.650928in}{2.336446in}}%
\pgfpathlineto{\pgfqpoint{0.659748in}{2.329464in}}%
\pgfpathlineto{\pgfqpoint{0.650466in}{2.312218in}}%
\pgfpathclose%
\pgfusepath{fill}%
\end{pgfscope}%
\begin{pgfscope}%
\pgfpathrectangle{\pgfqpoint{0.100000in}{0.100000in}}{\pgfqpoint{3.608454in}{2.310000in}}%
\pgfusepath{clip}%
\pgfsetbuttcap%
\pgfsetmiterjoin%
\definecolor{currentfill}{rgb}{0.000000,0.556863,0.721569}%
\pgfsetfillcolor{currentfill}%
\pgfsetlinewidth{0.000000pt}%
\definecolor{currentstroke}{rgb}{0.000000,0.000000,0.000000}%
\pgfsetstrokecolor{currentstroke}%
\pgfsetstrokeopacity{0.000000}%
\pgfsetdash{}{0pt}%
\pgfpathmoveto{\pgfqpoint{1.644176in}{0.728867in}}%
\pgfpathlineto{\pgfqpoint{1.678348in}{0.726376in}}%
\pgfpathlineto{\pgfqpoint{1.699678in}{0.725186in}}%
\pgfpathlineto{\pgfqpoint{1.697817in}{0.695195in}}%
\pgfpathlineto{\pgfqpoint{1.695622in}{0.661844in}}%
\pgfpathlineto{\pgfqpoint{1.640690in}{0.665716in}}%
\pgfpathlineto{\pgfqpoint{1.647908in}{0.670831in}}%
\pgfpathlineto{\pgfqpoint{1.642706in}{0.676280in}}%
\pgfpathlineto{\pgfqpoint{1.649534in}{0.685985in}}%
\pgfpathlineto{\pgfqpoint{1.648797in}{0.693179in}}%
\pgfpathlineto{\pgfqpoint{1.644365in}{0.694777in}}%
\pgfpathlineto{\pgfqpoint{1.616885in}{0.697027in}}%
\pgfpathlineto{\pgfqpoint{1.616475in}{0.692376in}}%
\pgfpathlineto{\pgfqpoint{1.602461in}{0.693594in}}%
\pgfpathlineto{\pgfqpoint{1.600425in}{0.668531in}}%
\pgfpathlineto{\pgfqpoint{1.585018in}{0.669731in}}%
\pgfpathlineto{\pgfqpoint{1.583492in}{0.651351in}}%
\pgfpathlineto{\pgfqpoint{1.518651in}{0.714404in}}%
\pgfpathlineto{\pgfqpoint{1.562114in}{0.758888in}}%
\pgfpathlineto{\pgfqpoint{1.567213in}{0.756734in}}%
\pgfpathlineto{\pgfqpoint{1.574315in}{0.749565in}}%
\pgfpathlineto{\pgfqpoint{1.578119in}{0.751288in}}%
\pgfpathlineto{\pgfqpoint{1.584520in}{0.753390in}}%
\pgfpathlineto{\pgfqpoint{1.599669in}{0.742990in}}%
\pgfpathlineto{\pgfqpoint{1.601306in}{0.732857in}}%
\pgfpathlineto{\pgfqpoint{1.644176in}{0.728867in}}%
\pgfpathclose%
\pgfusepath{fill}%
\end{pgfscope}%
\begin{pgfscope}%
\pgfpathrectangle{\pgfqpoint{0.100000in}{0.100000in}}{\pgfqpoint{3.608454in}{2.310000in}}%
\pgfusepath{clip}%
\pgfsetbuttcap%
\pgfsetmiterjoin%
\definecolor{currentfill}{rgb}{0.000000,0.745098,0.627451}%
\pgfsetfillcolor{currentfill}%
\pgfsetlinewidth{0.000000pt}%
\definecolor{currentstroke}{rgb}{0.000000,0.000000,0.000000}%
\pgfsetstrokecolor{currentstroke}%
\pgfsetstrokeopacity{0.000000}%
\pgfsetdash{}{0pt}%
\pgfpathmoveto{\pgfqpoint{2.297141in}{0.880461in}}%
\pgfpathlineto{\pgfqpoint{2.299404in}{0.870120in}}%
\pgfpathlineto{\pgfqpoint{2.302086in}{0.868312in}}%
\pgfpathlineto{\pgfqpoint{2.258381in}{0.867551in}}%
\pgfpathlineto{\pgfqpoint{2.240807in}{0.867381in}}%
\pgfpathlineto{\pgfqpoint{2.240853in}{0.895980in}}%
\pgfpathlineto{\pgfqpoint{2.231529in}{0.896018in}}%
\pgfpathlineto{\pgfqpoint{2.231585in}{0.901788in}}%
\pgfpathlineto{\pgfqpoint{2.231849in}{0.927551in}}%
\pgfpathlineto{\pgfqpoint{2.241965in}{0.929323in}}%
\pgfpathlineto{\pgfqpoint{2.245189in}{0.934157in}}%
\pgfpathlineto{\pgfqpoint{2.245528in}{0.957904in}}%
\pgfpathlineto{\pgfqpoint{2.259184in}{0.957741in}}%
\pgfpathlineto{\pgfqpoint{2.271324in}{0.957636in}}%
\pgfpathlineto{\pgfqpoint{2.281822in}{0.951221in}}%
\pgfpathlineto{\pgfqpoint{2.272954in}{0.950762in}}%
\pgfpathlineto{\pgfqpoint{2.274161in}{0.944083in}}%
\pgfpathlineto{\pgfqpoint{2.282183in}{0.933447in}}%
\pgfpathlineto{\pgfqpoint{2.284812in}{0.906848in}}%
\pgfpathlineto{\pgfqpoint{2.280766in}{0.896956in}}%
\pgfpathlineto{\pgfqpoint{2.282169in}{0.889435in}}%
\pgfpathlineto{\pgfqpoint{2.286293in}{0.889498in}}%
\pgfpathlineto{\pgfqpoint{2.297141in}{0.880461in}}%
\pgfpathclose%
\pgfusepath{fill}%
\end{pgfscope}%
\begin{pgfscope}%
\pgfpathrectangle{\pgfqpoint{0.100000in}{0.100000in}}{\pgfqpoint{3.608454in}{2.310000in}}%
\pgfusepath{clip}%
\pgfsetbuttcap%
\pgfsetmiterjoin%
\definecolor{currentfill}{rgb}{0.000000,0.521569,0.739216}%
\pgfsetfillcolor{currentfill}%
\pgfsetlinewidth{0.000000pt}%
\definecolor{currentstroke}{rgb}{0.000000,0.000000,0.000000}%
\pgfsetstrokecolor{currentstroke}%
\pgfsetstrokeopacity{0.000000}%
\pgfsetdash{}{0pt}%
\pgfpathmoveto{\pgfqpoint{1.795374in}{1.510393in}}%
\pgfpathlineto{\pgfqpoint{1.793857in}{1.483052in}}%
\pgfpathlineto{\pgfqpoint{1.760676in}{1.485050in}}%
\pgfpathlineto{\pgfqpoint{1.732544in}{1.486727in}}%
\pgfpathlineto{\pgfqpoint{1.734320in}{1.514317in}}%
\pgfpathlineto{\pgfqpoint{1.733109in}{1.514401in}}%
\pgfpathlineto{\pgfqpoint{1.735045in}{1.541820in}}%
\pgfpathlineto{\pgfqpoint{1.726904in}{1.542427in}}%
\pgfpathlineto{\pgfqpoint{1.728889in}{1.569871in}}%
\pgfpathlineto{\pgfqpoint{1.762096in}{1.567407in}}%
\pgfpathlineto{\pgfqpoint{1.796934in}{1.565224in}}%
\pgfpathlineto{\pgfqpoint{1.795374in}{1.510393in}}%
\pgfpathclose%
\pgfusepath{fill}%
\end{pgfscope}%
\begin{pgfscope}%
\pgfpathrectangle{\pgfqpoint{0.100000in}{0.100000in}}{\pgfqpoint{3.608454in}{2.310000in}}%
\pgfusepath{clip}%
\pgfsetbuttcap%
\pgfsetmiterjoin%
\definecolor{currentfill}{rgb}{0.000000,0.513725,0.743137}%
\pgfsetfillcolor{currentfill}%
\pgfsetlinewidth{0.000000pt}%
\definecolor{currentstroke}{rgb}{0.000000,0.000000,0.000000}%
\pgfsetstrokecolor{currentstroke}%
\pgfsetstrokeopacity{0.000000}%
\pgfsetdash{}{0pt}%
\pgfpathmoveto{\pgfqpoint{2.794880in}{1.439319in}}%
\pgfpathlineto{\pgfqpoint{2.795500in}{1.426123in}}%
\pgfpathlineto{\pgfqpoint{2.776762in}{1.425313in}}%
\pgfpathlineto{\pgfqpoint{2.757924in}{1.424641in}}%
\pgfpathlineto{\pgfqpoint{2.746401in}{1.421668in}}%
\pgfpathlineto{\pgfqpoint{2.726071in}{1.419220in}}%
\pgfpathlineto{\pgfqpoint{2.722450in}{1.453627in}}%
\pgfpathlineto{\pgfqpoint{2.720165in}{1.477649in}}%
\pgfpathlineto{\pgfqpoint{2.719775in}{1.481004in}}%
\pgfpathlineto{\pgfqpoint{2.741875in}{1.483652in}}%
\pgfpathlineto{\pgfqpoint{2.747392in}{1.486269in}}%
\pgfpathlineto{\pgfqpoint{2.746418in}{1.494296in}}%
\pgfpathlineto{\pgfqpoint{2.767029in}{1.496892in}}%
\pgfpathlineto{\pgfqpoint{2.766536in}{1.500999in}}%
\pgfpathlineto{\pgfqpoint{2.773279in}{1.502104in}}%
\pgfpathlineto{\pgfqpoint{2.795078in}{1.502065in}}%
\pgfpathlineto{\pgfqpoint{2.795940in}{1.480285in}}%
\pgfpathlineto{\pgfqpoint{2.799407in}{1.480427in}}%
\pgfpathlineto{\pgfqpoint{2.800713in}{1.463456in}}%
\pgfpathlineto{\pgfqpoint{2.798306in}{1.443885in}}%
\pgfpathlineto{\pgfqpoint{2.794880in}{1.439319in}}%
\pgfpathclose%
\pgfusepath{fill}%
\end{pgfscope}%
\begin{pgfscope}%
\pgfpathrectangle{\pgfqpoint{0.100000in}{0.100000in}}{\pgfqpoint{3.608454in}{2.310000in}}%
\pgfusepath{clip}%
\pgfsetbuttcap%
\pgfsetmiterjoin%
\definecolor{currentfill}{rgb}{0.000000,0.584314,0.707843}%
\pgfsetfillcolor{currentfill}%
\pgfsetlinewidth{0.000000pt}%
\definecolor{currentstroke}{rgb}{0.000000,0.000000,0.000000}%
\pgfsetstrokecolor{currentstroke}%
\pgfsetstrokeopacity{0.000000}%
\pgfsetdash{}{0pt}%
\pgfpathmoveto{\pgfqpoint{2.290718in}{1.442111in}}%
\pgfpathlineto{\pgfqpoint{2.291417in}{1.417786in}}%
\pgfpathlineto{\pgfqpoint{2.267659in}{1.417346in}}%
\pgfpathlineto{\pgfqpoint{2.267438in}{1.423982in}}%
\pgfpathlineto{\pgfqpoint{2.237055in}{1.423181in}}%
\pgfpathlineto{\pgfqpoint{2.236527in}{1.423164in}}%
\pgfpathlineto{\pgfqpoint{2.235992in}{1.447300in}}%
\pgfpathlineto{\pgfqpoint{2.246319in}{1.447619in}}%
\pgfpathlineto{\pgfqpoint{2.246988in}{1.453136in}}%
\pgfpathlineto{\pgfqpoint{2.244011in}{1.467014in}}%
\pgfpathlineto{\pgfqpoint{2.248598in}{1.467244in}}%
\pgfpathlineto{\pgfqpoint{2.290253in}{1.469670in}}%
\pgfpathlineto{\pgfqpoint{2.290718in}{1.442111in}}%
\pgfpathclose%
\pgfusepath{fill}%
\end{pgfscope}%
\begin{pgfscope}%
\pgfpathrectangle{\pgfqpoint{0.100000in}{0.100000in}}{\pgfqpoint{3.608454in}{2.310000in}}%
\pgfusepath{clip}%
\pgfsetbuttcap%
\pgfsetmiterjoin%
\definecolor{currentfill}{rgb}{0.000000,0.709804,0.645098}%
\pgfsetfillcolor{currentfill}%
\pgfsetlinewidth{0.000000pt}%
\definecolor{currentstroke}{rgb}{0.000000,0.000000,0.000000}%
\pgfsetstrokecolor{currentstroke}%
\pgfsetstrokeopacity{0.000000}%
\pgfsetdash{}{0pt}%
\pgfpathmoveto{\pgfqpoint{2.190786in}{1.003624in}}%
\pgfpathlineto{\pgfqpoint{2.180404in}{1.000901in}}%
\pgfpathlineto{\pgfqpoint{2.176519in}{0.997121in}}%
\pgfpathlineto{\pgfqpoint{2.170781in}{0.999186in}}%
\pgfpathlineto{\pgfqpoint{2.151260in}{0.999398in}}%
\pgfpathlineto{\pgfqpoint{2.142324in}{1.001839in}}%
\pgfpathlineto{\pgfqpoint{2.141998in}{0.984277in}}%
\pgfpathlineto{\pgfqpoint{2.110869in}{0.984129in}}%
\pgfpathlineto{\pgfqpoint{2.110850in}{0.997922in}}%
\pgfpathlineto{\pgfqpoint{2.073255in}{0.998037in}}%
\pgfpathlineto{\pgfqpoint{2.073230in}{0.991141in}}%
\pgfpathlineto{\pgfqpoint{2.062941in}{0.991185in}}%
\pgfpathlineto{\pgfqpoint{2.049224in}{0.991276in}}%
\pgfpathlineto{\pgfqpoint{2.049296in}{0.998169in}}%
\pgfpathlineto{\pgfqpoint{2.035603in}{0.998326in}}%
\pgfpathlineto{\pgfqpoint{2.036171in}{1.027587in}}%
\pgfpathlineto{\pgfqpoint{2.043077in}{1.035565in}}%
\pgfpathlineto{\pgfqpoint{2.050477in}{1.038266in}}%
\pgfpathlineto{\pgfqpoint{2.058740in}{1.038661in}}%
\pgfpathlineto{\pgfqpoint{2.069896in}{1.042778in}}%
\pgfpathlineto{\pgfqpoint{2.081865in}{1.049099in}}%
\pgfpathlineto{\pgfqpoint{2.091976in}{1.044393in}}%
\pgfpathlineto{\pgfqpoint{2.096410in}{1.052482in}}%
\pgfpathlineto{\pgfqpoint{2.101602in}{1.057429in}}%
\pgfpathlineto{\pgfqpoint{2.098258in}{1.064900in}}%
\pgfpathlineto{\pgfqpoint{2.098590in}{1.073824in}}%
\pgfpathlineto{\pgfqpoint{2.140727in}{1.073891in}}%
\pgfpathlineto{\pgfqpoint{2.139356in}{1.083433in}}%
\pgfpathlineto{\pgfqpoint{2.163020in}{1.082921in}}%
\pgfpathlineto{\pgfqpoint{2.176762in}{1.083808in}}%
\pgfpathlineto{\pgfqpoint{2.172137in}{1.075918in}}%
\pgfpathlineto{\pgfqpoint{2.166395in}{1.075969in}}%
\pgfpathlineto{\pgfqpoint{2.166287in}{1.069101in}}%
\pgfpathlineto{\pgfqpoint{2.169375in}{1.063715in}}%
\pgfpathlineto{\pgfqpoint{2.166477in}{1.058801in}}%
\pgfpathlineto{\pgfqpoint{2.166153in}{1.040540in}}%
\pgfpathlineto{\pgfqpoint{2.162548in}{1.034836in}}%
\pgfpathlineto{\pgfqpoint{2.162454in}{1.027888in}}%
\pgfpathlineto{\pgfqpoint{2.166850in}{1.025625in}}%
\pgfpathlineto{\pgfqpoint{2.190828in}{1.025374in}}%
\pgfpathlineto{\pgfqpoint{2.190786in}{1.003624in}}%
\pgfpathclose%
\pgfusepath{fill}%
\end{pgfscope}%
\begin{pgfscope}%
\pgfpathrectangle{\pgfqpoint{0.100000in}{0.100000in}}{\pgfqpoint{3.608454in}{2.310000in}}%
\pgfusepath{clip}%
\pgfsetbuttcap%
\pgfsetmiterjoin%
\definecolor{currentfill}{rgb}{0.000000,0.623529,0.688235}%
\pgfsetfillcolor{currentfill}%
\pgfsetlinewidth{0.000000pt}%
\definecolor{currentstroke}{rgb}{0.000000,0.000000,0.000000}%
\pgfsetstrokecolor{currentstroke}%
\pgfsetstrokeopacity{0.000000}%
\pgfsetdash{}{0pt}%
\pgfpathmoveto{\pgfqpoint{3.573009in}{1.957078in}}%
\pgfpathlineto{\pgfqpoint{3.580481in}{1.956181in}}%
\pgfpathlineto{\pgfqpoint{3.577419in}{1.950172in}}%
\pgfpathlineto{\pgfqpoint{3.573009in}{1.957078in}}%
\pgfpathclose%
\pgfusepath{fill}%
\end{pgfscope}%
\begin{pgfscope}%
\pgfpathrectangle{\pgfqpoint{0.100000in}{0.100000in}}{\pgfqpoint{3.608454in}{2.310000in}}%
\pgfusepath{clip}%
\pgfsetbuttcap%
\pgfsetmiterjoin%
\definecolor{currentfill}{rgb}{0.000000,0.623529,0.688235}%
\pgfsetfillcolor{currentfill}%
\pgfsetlinewidth{0.000000pt}%
\definecolor{currentstroke}{rgb}{0.000000,0.000000,0.000000}%
\pgfsetstrokecolor{currentstroke}%
\pgfsetstrokeopacity{0.000000}%
\pgfsetdash{}{0pt}%
\pgfpathmoveto{\pgfqpoint{3.550714in}{1.937007in}}%
\pgfpathlineto{\pgfqpoint{3.553924in}{1.946267in}}%
\pgfpathlineto{\pgfqpoint{3.547311in}{1.956692in}}%
\pgfpathlineto{\pgfqpoint{3.543555in}{1.955713in}}%
\pgfpathlineto{\pgfqpoint{3.538805in}{1.962876in}}%
\pgfpathlineto{\pgfqpoint{3.535140in}{1.964201in}}%
\pgfpathlineto{\pgfqpoint{3.535820in}{1.974394in}}%
\pgfpathlineto{\pgfqpoint{3.534775in}{1.988683in}}%
\pgfpathlineto{\pgfqpoint{3.532167in}{1.996604in}}%
\pgfpathlineto{\pgfqpoint{3.539818in}{1.996897in}}%
\pgfpathlineto{\pgfqpoint{3.527215in}{2.022326in}}%
\pgfpathlineto{\pgfqpoint{3.514278in}{2.013412in}}%
\pgfpathlineto{\pgfqpoint{3.494071in}{2.051630in}}%
\pgfpathlineto{\pgfqpoint{3.496496in}{2.059377in}}%
\pgfpathlineto{\pgfqpoint{3.490714in}{2.062918in}}%
\pgfpathlineto{\pgfqpoint{3.491275in}{2.076574in}}%
\pgfpathlineto{\pgfqpoint{3.487328in}{2.085547in}}%
\pgfpathlineto{\pgfqpoint{3.478067in}{2.117502in}}%
\pgfpathlineto{\pgfqpoint{3.474712in}{2.131510in}}%
\pgfpathlineto{\pgfqpoint{3.522068in}{2.145170in}}%
\pgfpathlineto{\pgfqpoint{3.526210in}{2.131697in}}%
\pgfpathlineto{\pgfqpoint{3.546699in}{2.136666in}}%
\pgfpathlineto{\pgfqpoint{3.557203in}{2.103454in}}%
\pgfpathlineto{\pgfqpoint{3.565274in}{2.075679in}}%
\pgfpathlineto{\pgfqpoint{3.566650in}{2.075414in}}%
\pgfpathlineto{\pgfqpoint{3.584395in}{2.086358in}}%
\pgfpathlineto{\pgfqpoint{3.598893in}{2.059799in}}%
\pgfpathlineto{\pgfqpoint{3.592928in}{2.056963in}}%
\pgfpathlineto{\pgfqpoint{3.603808in}{2.035391in}}%
\pgfpathlineto{\pgfqpoint{3.597567in}{2.032110in}}%
\pgfpathlineto{\pgfqpoint{3.610482in}{2.006964in}}%
\pgfpathlineto{\pgfqpoint{3.618035in}{1.994535in}}%
\pgfpathlineto{\pgfqpoint{3.613315in}{1.987169in}}%
\pgfpathlineto{\pgfqpoint{3.607145in}{1.996865in}}%
\pgfpathlineto{\pgfqpoint{3.599803in}{1.992874in}}%
\pgfpathlineto{\pgfqpoint{3.608045in}{1.984589in}}%
\pgfpathlineto{\pgfqpoint{3.601528in}{1.973738in}}%
\pgfpathlineto{\pgfqpoint{3.594431in}{1.979179in}}%
\pgfpathlineto{\pgfqpoint{3.594740in}{1.984518in}}%
\pgfpathlineto{\pgfqpoint{3.589736in}{1.989331in}}%
\pgfpathlineto{\pgfqpoint{3.587038in}{1.983926in}}%
\pgfpathlineto{\pgfqpoint{3.586028in}{1.973437in}}%
\pgfpathlineto{\pgfqpoint{3.576059in}{1.975363in}}%
\pgfpathlineto{\pgfqpoint{3.564314in}{1.981978in}}%
\pgfpathlineto{\pgfqpoint{3.561572in}{1.978778in}}%
\pgfpathlineto{\pgfqpoint{3.565543in}{1.970193in}}%
\pgfpathlineto{\pgfqpoint{3.562497in}{1.956796in}}%
\pgfpathlineto{\pgfqpoint{3.565408in}{1.948473in}}%
\pgfpathlineto{\pgfqpoint{3.560888in}{1.939136in}}%
\pgfpathlineto{\pgfqpoint{3.550714in}{1.937007in}}%
\pgfpathclose%
\pgfusepath{fill}%
\end{pgfscope}%
\begin{pgfscope}%
\pgfpathrectangle{\pgfqpoint{0.100000in}{0.100000in}}{\pgfqpoint{3.608454in}{2.310000in}}%
\pgfusepath{clip}%
\pgfsetbuttcap%
\pgfsetmiterjoin%
\definecolor{currentfill}{rgb}{0.000000,0.564706,0.717647}%
\pgfsetfillcolor{currentfill}%
\pgfsetlinewidth{0.000000pt}%
\definecolor{currentstroke}{rgb}{0.000000,0.000000,0.000000}%
\pgfsetstrokecolor{currentstroke}%
\pgfsetstrokeopacity{0.000000}%
\pgfsetdash{}{0pt}%
\pgfpathmoveto{\pgfqpoint{1.859593in}{0.670192in}}%
\pgfpathlineto{\pgfqpoint{1.833962in}{0.671147in}}%
\pgfpathlineto{\pgfqpoint{1.835408in}{0.704889in}}%
\pgfpathlineto{\pgfqpoint{1.870905in}{0.703446in}}%
\pgfpathlineto{\pgfqpoint{1.871649in}{0.712052in}}%
\pgfpathlineto{\pgfqpoint{1.907744in}{0.711220in}}%
\pgfpathlineto{\pgfqpoint{1.913119in}{0.701138in}}%
\pgfpathlineto{\pgfqpoint{1.903598in}{0.691454in}}%
\pgfpathlineto{\pgfqpoint{1.892292in}{0.669010in}}%
\pgfpathlineto{\pgfqpoint{1.891457in}{0.663121in}}%
\pgfpathlineto{\pgfqpoint{1.871994in}{0.669892in}}%
\pgfpathlineto{\pgfqpoint{1.859593in}{0.670192in}}%
\pgfpathclose%
\pgfusepath{fill}%
\end{pgfscope}%
\begin{pgfscope}%
\pgfpathrectangle{\pgfqpoint{0.100000in}{0.100000in}}{\pgfqpoint{3.608454in}{2.310000in}}%
\pgfusepath{clip}%
\pgfsetbuttcap%
\pgfsetmiterjoin%
\definecolor{currentfill}{rgb}{0.000000,0.803922,0.598039}%
\pgfsetfillcolor{currentfill}%
\pgfsetlinewidth{0.000000pt}%
\definecolor{currentstroke}{rgb}{0.000000,0.000000,0.000000}%
\pgfsetstrokecolor{currentstroke}%
\pgfsetstrokeopacity{0.000000}%
\pgfsetdash{}{0pt}%
\pgfpathmoveto{\pgfqpoint{2.362240in}{0.870230in}}%
\pgfpathlineto{\pgfqpoint{2.367614in}{0.868244in}}%
\pgfpathlineto{\pgfqpoint{2.368994in}{0.861333in}}%
\pgfpathlineto{\pgfqpoint{2.363366in}{0.855813in}}%
\pgfpathlineto{\pgfqpoint{2.363126in}{0.851002in}}%
\pgfpathlineto{\pgfqpoint{2.370206in}{0.848561in}}%
\pgfpathlineto{\pgfqpoint{2.364353in}{0.841406in}}%
\pgfpathlineto{\pgfqpoint{2.366196in}{0.837164in}}%
\pgfpathlineto{\pgfqpoint{2.371879in}{0.836595in}}%
\pgfpathlineto{\pgfqpoint{2.365973in}{0.833290in}}%
\pgfpathlineto{\pgfqpoint{2.344269in}{0.832393in}}%
\pgfpathlineto{\pgfqpoint{2.344959in}{0.835935in}}%
\pgfpathlineto{\pgfqpoint{2.334864in}{0.835567in}}%
\pgfpathlineto{\pgfqpoint{2.331640in}{0.845019in}}%
\pgfpathlineto{\pgfqpoint{2.337645in}{0.859095in}}%
\pgfpathlineto{\pgfqpoint{2.340706in}{0.860346in}}%
\pgfpathlineto{\pgfqpoint{2.344274in}{0.869656in}}%
\pgfpathlineto{\pgfqpoint{2.362240in}{0.870230in}}%
\pgfpathclose%
\pgfusepath{fill}%
\end{pgfscope}%
\begin{pgfscope}%
\pgfpathrectangle{\pgfqpoint{0.100000in}{0.100000in}}{\pgfqpoint{3.608454in}{2.310000in}}%
\pgfusepath{clip}%
\pgfsetbuttcap%
\pgfsetmiterjoin%
\definecolor{currentfill}{rgb}{0.000000,0.749020,0.625490}%
\pgfsetfillcolor{currentfill}%
\pgfsetlinewidth{0.000000pt}%
\definecolor{currentstroke}{rgb}{0.000000,0.000000,0.000000}%
\pgfsetstrokecolor{currentstroke}%
\pgfsetstrokeopacity{0.000000}%
\pgfsetdash{}{0pt}%
\pgfpathmoveto{\pgfqpoint{2.093068in}{0.691255in}}%
\pgfpathlineto{\pgfqpoint{2.089584in}{0.697805in}}%
\pgfpathlineto{\pgfqpoint{2.084351in}{0.694073in}}%
\pgfpathlineto{\pgfqpoint{2.081886in}{0.666080in}}%
\pgfpathlineto{\pgfqpoint{2.065563in}{0.666170in}}%
\pgfpathlineto{\pgfqpoint{2.049731in}{0.675985in}}%
\pgfpathlineto{\pgfqpoint{2.047653in}{0.694679in}}%
\pgfpathlineto{\pgfqpoint{2.026699in}{0.691630in}}%
\pgfpathlineto{\pgfqpoint{2.021875in}{0.703710in}}%
\pgfpathlineto{\pgfqpoint{2.040106in}{0.712927in}}%
\pgfpathlineto{\pgfqpoint{2.054567in}{0.712981in}}%
\pgfpathlineto{\pgfqpoint{2.053852in}{0.716384in}}%
\pgfpathlineto{\pgfqpoint{2.057335in}{0.727032in}}%
\pgfpathlineto{\pgfqpoint{2.061056in}{0.729018in}}%
\pgfpathlineto{\pgfqpoint{2.056507in}{0.745639in}}%
\pgfpathlineto{\pgfqpoint{2.062519in}{0.748614in}}%
\pgfpathlineto{\pgfqpoint{2.088232in}{0.752553in}}%
\pgfpathlineto{\pgfqpoint{2.096621in}{0.750671in}}%
\pgfpathlineto{\pgfqpoint{2.099586in}{0.743769in}}%
\pgfpathlineto{\pgfqpoint{2.106573in}{0.739232in}}%
\pgfpathlineto{\pgfqpoint{2.115910in}{0.747279in}}%
\pgfpathlineto{\pgfqpoint{2.125247in}{0.741762in}}%
\pgfpathlineto{\pgfqpoint{2.140576in}{0.738107in}}%
\pgfpathlineto{\pgfqpoint{2.145077in}{0.730301in}}%
\pgfpathlineto{\pgfqpoint{2.152123in}{0.723725in}}%
\pgfpathlineto{\pgfqpoint{2.166253in}{0.713450in}}%
\pgfpathlineto{\pgfqpoint{2.143868in}{0.708025in}}%
\pgfpathlineto{\pgfqpoint{2.132588in}{0.712274in}}%
\pgfpathlineto{\pgfqpoint{2.123775in}{0.713712in}}%
\pgfpathlineto{\pgfqpoint{2.117525in}{0.716961in}}%
\pgfpathlineto{\pgfqpoint{2.110954in}{0.708601in}}%
\pgfpathlineto{\pgfqpoint{2.093068in}{0.691255in}}%
\pgfpathclose%
\pgfusepath{fill}%
\end{pgfscope}%
\begin{pgfscope}%
\pgfpathrectangle{\pgfqpoint{0.100000in}{0.100000in}}{\pgfqpoint{3.608454in}{2.310000in}}%
\pgfusepath{clip}%
\pgfsetbuttcap%
\pgfsetmiterjoin%
\definecolor{currentfill}{rgb}{0.000000,0.654902,0.672549}%
\pgfsetfillcolor{currentfill}%
\pgfsetlinewidth{0.000000pt}%
\definecolor{currentstroke}{rgb}{0.000000,0.000000,0.000000}%
\pgfsetstrokecolor{currentstroke}%
\pgfsetstrokeopacity{0.000000}%
\pgfsetdash{}{0pt}%
\pgfpathmoveto{\pgfqpoint{2.904596in}{1.155570in}}%
\pgfpathlineto{\pgfqpoint{2.898012in}{1.153842in}}%
\pgfpathlineto{\pgfqpoint{2.895600in}{1.161365in}}%
\pgfpathlineto{\pgfqpoint{2.879988in}{1.149617in}}%
\pgfpathlineto{\pgfqpoint{2.868331in}{1.159628in}}%
\pgfpathlineto{\pgfqpoint{2.860386in}{1.164003in}}%
\pgfpathlineto{\pgfqpoint{2.865237in}{1.171385in}}%
\pgfpathlineto{\pgfqpoint{2.859396in}{1.177849in}}%
\pgfpathlineto{\pgfqpoint{2.853813in}{1.173625in}}%
\pgfpathlineto{\pgfqpoint{2.851139in}{1.181704in}}%
\pgfpathlineto{\pgfqpoint{2.856311in}{1.184739in}}%
\pgfpathlineto{\pgfqpoint{2.859737in}{1.192156in}}%
\pgfpathlineto{\pgfqpoint{2.867721in}{1.199764in}}%
\pgfpathlineto{\pgfqpoint{2.869228in}{1.206518in}}%
\pgfpathlineto{\pgfqpoint{2.876442in}{1.212071in}}%
\pgfpathlineto{\pgfqpoint{2.880082in}{1.220503in}}%
\pgfpathlineto{\pgfqpoint{2.886736in}{1.225061in}}%
\pgfpathlineto{\pgfqpoint{2.901040in}{1.226866in}}%
\pgfpathlineto{\pgfqpoint{2.907547in}{1.214437in}}%
\pgfpathlineto{\pgfqpoint{2.920870in}{1.223465in}}%
\pgfpathlineto{\pgfqpoint{2.921871in}{1.228104in}}%
\pgfpathlineto{\pgfqpoint{2.933591in}{1.232567in}}%
\pgfpathlineto{\pgfqpoint{2.937030in}{1.237689in}}%
\pgfpathlineto{\pgfqpoint{2.943559in}{1.240996in}}%
\pgfpathlineto{\pgfqpoint{2.960170in}{1.246122in}}%
\pgfpathlineto{\pgfqpoint{2.968379in}{1.238517in}}%
\pgfpathlineto{\pgfqpoint{2.974075in}{1.228469in}}%
\pgfpathlineto{\pgfqpoint{2.958209in}{1.221651in}}%
\pgfpathlineto{\pgfqpoint{2.957411in}{1.217461in}}%
\pgfpathlineto{\pgfqpoint{2.951682in}{1.213011in}}%
\pgfpathlineto{\pgfqpoint{2.940357in}{1.211530in}}%
\pgfpathlineto{\pgfqpoint{2.934516in}{1.206979in}}%
\pgfpathlineto{\pgfqpoint{2.933739in}{1.199536in}}%
\pgfpathlineto{\pgfqpoint{2.929794in}{1.191773in}}%
\pgfpathlineto{\pgfqpoint{2.936507in}{1.182790in}}%
\pgfpathlineto{\pgfqpoint{2.933147in}{1.171253in}}%
\pgfpathlineto{\pgfqpoint{2.930144in}{1.169183in}}%
\pgfpathlineto{\pgfqpoint{2.920582in}{1.172131in}}%
\pgfpathlineto{\pgfqpoint{2.909797in}{1.165077in}}%
\pgfpathlineto{\pgfqpoint{2.904596in}{1.155570in}}%
\pgfpathclose%
\pgfusepath{fill}%
\end{pgfscope}%
\begin{pgfscope}%
\pgfpathrectangle{\pgfqpoint{0.100000in}{0.100000in}}{\pgfqpoint{3.608454in}{2.310000in}}%
\pgfusepath{clip}%
\pgfsetbuttcap%
\pgfsetmiterjoin%
\definecolor{currentfill}{rgb}{0.000000,0.701961,0.649020}%
\pgfsetfillcolor{currentfill}%
\pgfsetlinewidth{0.000000pt}%
\definecolor{currentstroke}{rgb}{0.000000,0.000000,0.000000}%
\pgfsetstrokecolor{currentstroke}%
\pgfsetstrokeopacity{0.000000}%
\pgfsetdash{}{0pt}%
\pgfpathmoveto{\pgfqpoint{2.537283in}{1.876850in}}%
\pgfpathlineto{\pgfqpoint{2.528827in}{1.854971in}}%
\pgfpathlineto{\pgfqpoint{2.521970in}{1.844929in}}%
\pgfpathlineto{\pgfqpoint{2.520382in}{1.829615in}}%
\pgfpathlineto{\pgfqpoint{2.513293in}{1.828191in}}%
\pgfpathlineto{\pgfqpoint{2.492843in}{1.831235in}}%
\pgfpathlineto{\pgfqpoint{2.492447in}{1.838114in}}%
\pgfpathlineto{\pgfqpoint{2.488144in}{1.844796in}}%
\pgfpathlineto{\pgfqpoint{2.481279in}{1.844556in}}%
\pgfpathlineto{\pgfqpoint{2.480448in}{1.858288in}}%
\pgfpathlineto{\pgfqpoint{2.473620in}{1.857983in}}%
\pgfpathlineto{\pgfqpoint{2.471787in}{1.885405in}}%
\pgfpathlineto{\pgfqpoint{2.457894in}{1.884450in}}%
\pgfpathlineto{\pgfqpoint{2.456551in}{1.904950in}}%
\pgfpathlineto{\pgfqpoint{2.438472in}{1.913015in}}%
\pgfpathlineto{\pgfqpoint{2.436761in}{1.938649in}}%
\pgfpathlineto{\pgfqpoint{2.484672in}{1.942048in}}%
\pgfpathlineto{\pgfqpoint{2.485703in}{1.928298in}}%
\pgfpathlineto{\pgfqpoint{2.513243in}{1.930478in}}%
\pgfpathlineto{\pgfqpoint{2.514867in}{1.909757in}}%
\pgfpathlineto{\pgfqpoint{2.528620in}{1.910846in}}%
\pgfpathlineto{\pgfqpoint{2.531516in}{1.904166in}}%
\pgfpathlineto{\pgfqpoint{2.533757in}{1.876549in}}%
\pgfpathlineto{\pgfqpoint{2.537283in}{1.876850in}}%
\pgfpathclose%
\pgfusepath{fill}%
\end{pgfscope}%
\begin{pgfscope}%
\pgfpathrectangle{\pgfqpoint{0.100000in}{0.100000in}}{\pgfqpoint{3.608454in}{2.310000in}}%
\pgfusepath{clip}%
\pgfsetbuttcap%
\pgfsetmiterjoin%
\definecolor{currentfill}{rgb}{0.000000,0.545098,0.727451}%
\pgfsetfillcolor{currentfill}%
\pgfsetlinewidth{0.000000pt}%
\definecolor{currentstroke}{rgb}{0.000000,0.000000,0.000000}%
\pgfsetstrokecolor{currentstroke}%
\pgfsetstrokeopacity{0.000000}%
\pgfsetdash{}{0pt}%
\pgfpathmoveto{\pgfqpoint{2.507414in}{1.036455in}}%
\pgfpathlineto{\pgfqpoint{2.483127in}{1.034794in}}%
\pgfpathlineto{\pgfqpoint{2.482016in}{1.066685in}}%
\pgfpathlineto{\pgfqpoint{2.488608in}{1.069834in}}%
\pgfpathlineto{\pgfqpoint{2.487860in}{1.090523in}}%
\pgfpathlineto{\pgfqpoint{2.469563in}{1.097602in}}%
\pgfpathlineto{\pgfqpoint{2.480144in}{1.116463in}}%
\pgfpathlineto{\pgfqpoint{2.479806in}{1.130609in}}%
\pgfpathlineto{\pgfqpoint{2.492139in}{1.131174in}}%
\pgfpathlineto{\pgfqpoint{2.498683in}{1.126546in}}%
\pgfpathlineto{\pgfqpoint{2.510006in}{1.121373in}}%
\pgfpathlineto{\pgfqpoint{2.510594in}{1.099857in}}%
\pgfpathlineto{\pgfqpoint{2.517006in}{1.100169in}}%
\pgfpathlineto{\pgfqpoint{2.517692in}{1.084251in}}%
\pgfpathlineto{\pgfqpoint{2.525076in}{1.078728in}}%
\pgfpathlineto{\pgfqpoint{2.531082in}{1.076803in}}%
\pgfpathlineto{\pgfqpoint{2.534914in}{1.072050in}}%
\pgfpathlineto{\pgfqpoint{2.535012in}{1.069064in}}%
\pgfpathlineto{\pgfqpoint{2.521534in}{1.068158in}}%
\pgfpathlineto{\pgfqpoint{2.517201in}{1.066138in}}%
\pgfpathlineto{\pgfqpoint{2.512736in}{1.056859in}}%
\pgfpathlineto{\pgfqpoint{2.508734in}{1.056586in}}%
\pgfpathlineto{\pgfqpoint{2.509787in}{1.036620in}}%
\pgfpathlineto{\pgfqpoint{2.507414in}{1.036455in}}%
\pgfpathclose%
\pgfusepath{fill}%
\end{pgfscope}%
\begin{pgfscope}%
\pgfpathrectangle{\pgfqpoint{0.100000in}{0.100000in}}{\pgfqpoint{3.608454in}{2.310000in}}%
\pgfusepath{clip}%
\pgfsetbuttcap%
\pgfsetmiterjoin%
\definecolor{currentfill}{rgb}{0.000000,0.490196,0.754902}%
\pgfsetfillcolor{currentfill}%
\pgfsetlinewidth{0.000000pt}%
\definecolor{currentstroke}{rgb}{0.000000,0.000000,0.000000}%
\pgfsetstrokecolor{currentstroke}%
\pgfsetstrokeopacity{0.000000}%
\pgfsetdash{}{0pt}%
\pgfpathmoveto{\pgfqpoint{1.960306in}{1.310682in}}%
\pgfpathlineto{\pgfqpoint{1.926190in}{1.311718in}}%
\pgfpathlineto{\pgfqpoint{1.926584in}{1.332368in}}%
\pgfpathlineto{\pgfqpoint{1.892401in}{1.333545in}}%
\pgfpathlineto{\pgfqpoint{1.892889in}{1.354245in}}%
\pgfpathlineto{\pgfqpoint{1.893147in}{1.361110in}}%
\pgfpathlineto{\pgfqpoint{1.927400in}{1.359864in}}%
\pgfpathlineto{\pgfqpoint{1.927666in}{1.366756in}}%
\pgfpathlineto{\pgfqpoint{1.961886in}{1.365738in}}%
\pgfpathlineto{\pgfqpoint{1.960306in}{1.310682in}}%
\pgfpathclose%
\pgfusepath{fill}%
\end{pgfscope}%
\begin{pgfscope}%
\pgfpathrectangle{\pgfqpoint{0.100000in}{0.100000in}}{\pgfqpoint{3.608454in}{2.310000in}}%
\pgfusepath{clip}%
\pgfsetbuttcap%
\pgfsetmiterjoin%
\definecolor{currentfill}{rgb}{0.000000,0.494118,0.752941}%
\pgfsetfillcolor{currentfill}%
\pgfsetlinewidth{0.000000pt}%
\definecolor{currentstroke}{rgb}{0.000000,0.000000,0.000000}%
\pgfsetstrokecolor{currentstroke}%
\pgfsetstrokeopacity{0.000000}%
\pgfsetdash{}{0pt}%
\pgfpathmoveto{\pgfqpoint{2.548186in}{1.158629in}}%
\pgfpathlineto{\pgfqpoint{2.549236in}{1.162164in}}%
\pgfpathlineto{\pgfqpoint{2.543041in}{1.178054in}}%
\pgfpathlineto{\pgfqpoint{2.539380in}{1.187318in}}%
\pgfpathlineto{\pgfqpoint{2.555635in}{1.195943in}}%
\pgfpathlineto{\pgfqpoint{2.563459in}{1.197714in}}%
\pgfpathlineto{\pgfqpoint{2.566571in}{1.201683in}}%
\pgfpathlineto{\pgfqpoint{2.567654in}{1.211922in}}%
\pgfpathlineto{\pgfqpoint{2.577947in}{1.209252in}}%
\pgfpathlineto{\pgfqpoint{2.589411in}{1.214310in}}%
\pgfpathlineto{\pgfqpoint{2.595443in}{1.205574in}}%
\pgfpathlineto{\pgfqpoint{2.607679in}{1.208265in}}%
\pgfpathlineto{\pgfqpoint{2.610086in}{1.175304in}}%
\pgfpathlineto{\pgfqpoint{2.606663in}{1.174945in}}%
\pgfpathlineto{\pgfqpoint{2.607032in}{1.160645in}}%
\pgfpathlineto{\pgfqpoint{2.599908in}{1.153357in}}%
\pgfpathlineto{\pgfqpoint{2.597965in}{1.148745in}}%
\pgfpathlineto{\pgfqpoint{2.583459in}{1.148544in}}%
\pgfpathlineto{\pgfqpoint{2.583376in}{1.141723in}}%
\pgfpathlineto{\pgfqpoint{2.579740in}{1.137359in}}%
\pgfpathlineto{\pgfqpoint{2.570262in}{1.139592in}}%
\pgfpathlineto{\pgfqpoint{2.556313in}{1.139100in}}%
\pgfpathlineto{\pgfqpoint{2.555391in}{1.144733in}}%
\pgfpathlineto{\pgfqpoint{2.548297in}{1.153932in}}%
\pgfpathlineto{\pgfqpoint{2.548186in}{1.158629in}}%
\pgfpathclose%
\pgfusepath{fill}%
\end{pgfscope}%
\begin{pgfscope}%
\pgfpathrectangle{\pgfqpoint{0.100000in}{0.100000in}}{\pgfqpoint{3.608454in}{2.310000in}}%
\pgfusepath{clip}%
\pgfsetbuttcap%
\pgfsetmiterjoin%
\definecolor{currentfill}{rgb}{0.000000,0.317647,0.841176}%
\pgfsetfillcolor{currentfill}%
\pgfsetlinewidth{0.000000pt}%
\definecolor{currentstroke}{rgb}{0.000000,0.000000,0.000000}%
\pgfsetstrokecolor{currentstroke}%
\pgfsetstrokeopacity{0.000000}%
\pgfsetdash{}{0pt}%
\pgfpathmoveto{\pgfqpoint{3.245459in}{1.515286in}}%
\pgfpathlineto{\pgfqpoint{3.266966in}{1.519876in}}%
\pgfpathlineto{\pgfqpoint{3.271013in}{1.528198in}}%
\pgfpathlineto{\pgfqpoint{3.275612in}{1.531070in}}%
\pgfpathlineto{\pgfqpoint{3.284318in}{1.531602in}}%
\pgfpathlineto{\pgfqpoint{3.284745in}{1.526913in}}%
\pgfpathlineto{\pgfqpoint{3.281349in}{1.515956in}}%
\pgfpathlineto{\pgfqpoint{3.285117in}{1.511332in}}%
\pgfpathlineto{\pgfqpoint{3.285323in}{1.505755in}}%
\pgfpathlineto{\pgfqpoint{3.282148in}{1.502290in}}%
\pgfpathlineto{\pgfqpoint{3.288460in}{1.495799in}}%
\pgfpathlineto{\pgfqpoint{3.281434in}{1.488221in}}%
\pgfpathlineto{\pgfqpoint{3.274847in}{1.487194in}}%
\pgfpathlineto{\pgfqpoint{3.273108in}{1.493337in}}%
\pgfpathlineto{\pgfqpoint{3.266239in}{1.490550in}}%
\pgfpathlineto{\pgfqpoint{3.256327in}{1.489994in}}%
\pgfpathlineto{\pgfqpoint{3.258440in}{1.494561in}}%
\pgfpathlineto{\pgfqpoint{3.257846in}{1.502321in}}%
\pgfpathlineto{\pgfqpoint{3.251965in}{1.502268in}}%
\pgfpathlineto{\pgfqpoint{3.239646in}{1.514041in}}%
\pgfpathlineto{\pgfqpoint{3.245459in}{1.515286in}}%
\pgfpathclose%
\pgfusepath{fill}%
\end{pgfscope}%
\begin{pgfscope}%
\pgfpathrectangle{\pgfqpoint{0.100000in}{0.100000in}}{\pgfqpoint{3.608454in}{2.310000in}}%
\pgfusepath{clip}%
\pgfsetbuttcap%
\pgfsetmiterjoin%
\definecolor{currentfill}{rgb}{0.000000,0.364706,0.817647}%
\pgfsetfillcolor{currentfill}%
\pgfsetlinewidth{0.000000pt}%
\definecolor{currentstroke}{rgb}{0.000000,0.000000,0.000000}%
\pgfsetstrokecolor{currentstroke}%
\pgfsetstrokeopacity{0.000000}%
\pgfsetdash{}{0pt}%
\pgfpathmoveto{\pgfqpoint{1.791350in}{1.359031in}}%
\pgfpathlineto{\pgfqpoint{1.790342in}{1.359067in}}%
\pgfpathlineto{\pgfqpoint{1.792462in}{1.393414in}}%
\pgfpathlineto{\pgfqpoint{1.824934in}{1.391676in}}%
\pgfpathlineto{\pgfqpoint{1.860470in}{1.390012in}}%
\pgfpathlineto{\pgfqpoint{1.858721in}{1.355685in}}%
\pgfpathlineto{\pgfqpoint{1.791350in}{1.359031in}}%
\pgfpathclose%
\pgfusepath{fill}%
\end{pgfscope}%
\begin{pgfscope}%
\pgfpathrectangle{\pgfqpoint{0.100000in}{0.100000in}}{\pgfqpoint{3.608454in}{2.310000in}}%
\pgfusepath{clip}%
\pgfsetbuttcap%
\pgfsetmiterjoin%
\definecolor{currentfill}{rgb}{0.000000,0.478431,0.760784}%
\pgfsetfillcolor{currentfill}%
\pgfsetlinewidth{0.000000pt}%
\definecolor{currentstroke}{rgb}{0.000000,0.000000,0.000000}%
\pgfsetstrokecolor{currentstroke}%
\pgfsetstrokeopacity{0.000000}%
\pgfsetdash{}{0pt}%
\pgfpathmoveto{\pgfqpoint{1.895416in}{1.185262in}}%
\pgfpathlineto{\pgfqpoint{1.896376in}{1.215785in}}%
\pgfpathlineto{\pgfqpoint{1.889149in}{1.216054in}}%
\pgfpathlineto{\pgfqpoint{1.890205in}{1.243605in}}%
\pgfpathlineto{\pgfqpoint{1.938084in}{1.242176in}}%
\pgfpathlineto{\pgfqpoint{1.938292in}{1.256226in}}%
\pgfpathlineto{\pgfqpoint{1.972560in}{1.255285in}}%
\pgfpathlineto{\pgfqpoint{1.972846in}{1.269115in}}%
\pgfpathlineto{\pgfqpoint{1.992256in}{1.268523in}}%
\pgfpathlineto{\pgfqpoint{2.012052in}{1.268224in}}%
\pgfpathlineto{\pgfqpoint{2.011063in}{1.220009in}}%
\pgfpathlineto{\pgfqpoint{1.971690in}{1.220737in}}%
\pgfpathlineto{\pgfqpoint{1.971148in}{1.183063in}}%
\pgfpathlineto{\pgfqpoint{1.951284in}{1.183534in}}%
\pgfpathlineto{\pgfqpoint{1.895416in}{1.185262in}}%
\pgfpathclose%
\pgfusepath{fill}%
\end{pgfscope}%
\begin{pgfscope}%
\pgfpathrectangle{\pgfqpoint{0.100000in}{0.100000in}}{\pgfqpoint{3.608454in}{2.310000in}}%
\pgfusepath{clip}%
\pgfsetbuttcap%
\pgfsetmiterjoin%
\definecolor{currentfill}{rgb}{0.000000,0.470588,0.764706}%
\pgfsetfillcolor{currentfill}%
\pgfsetlinewidth{0.000000pt}%
\definecolor{currentstroke}{rgb}{0.000000,0.000000,0.000000}%
\pgfsetstrokecolor{currentstroke}%
\pgfsetstrokeopacity{0.000000}%
\pgfsetdash{}{0pt}%
\pgfpathmoveto{\pgfqpoint{1.811548in}{0.930044in}}%
\pgfpathlineto{\pgfqpoint{1.810392in}{0.932854in}}%
\pgfpathlineto{\pgfqpoint{1.797246in}{0.938094in}}%
\pgfpathlineto{\pgfqpoint{1.774163in}{0.940149in}}%
\pgfpathlineto{\pgfqpoint{1.775993in}{0.971342in}}%
\pgfpathlineto{\pgfqpoint{1.779205in}{0.970733in}}%
\pgfpathlineto{\pgfqpoint{1.780677in}{0.997355in}}%
\pgfpathlineto{\pgfqpoint{1.785322in}{0.998334in}}%
\pgfpathlineto{\pgfqpoint{1.799655in}{0.981871in}}%
\pgfpathlineto{\pgfqpoint{1.805846in}{0.981264in}}%
\pgfpathlineto{\pgfqpoint{1.811536in}{0.984126in}}%
\pgfpathlineto{\pgfqpoint{1.818337in}{0.980492in}}%
\pgfpathlineto{\pgfqpoint{1.820478in}{0.987070in}}%
\pgfpathlineto{\pgfqpoint{1.831143in}{0.977031in}}%
\pgfpathlineto{\pgfqpoint{1.832097in}{0.967259in}}%
\pgfpathlineto{\pgfqpoint{1.847621in}{0.966431in}}%
\pgfpathlineto{\pgfqpoint{1.846276in}{0.936427in}}%
\pgfpathlineto{\pgfqpoint{1.811868in}{0.938019in}}%
\pgfpathlineto{\pgfqpoint{1.811548in}{0.930044in}}%
\pgfpathclose%
\pgfusepath{fill}%
\end{pgfscope}%
\begin{pgfscope}%
\pgfpathrectangle{\pgfqpoint{0.100000in}{0.100000in}}{\pgfqpoint{3.608454in}{2.310000in}}%
\pgfusepath{clip}%
\pgfsetbuttcap%
\pgfsetmiterjoin%
\definecolor{currentfill}{rgb}{0.000000,0.537255,0.731373}%
\pgfsetfillcolor{currentfill}%
\pgfsetlinewidth{0.000000pt}%
\definecolor{currentstroke}{rgb}{0.000000,0.000000,0.000000}%
\pgfsetstrokecolor{currentstroke}%
\pgfsetstrokeopacity{0.000000}%
\pgfsetdash{}{0pt}%
\pgfpathmoveto{\pgfqpoint{0.906140in}{0.393309in}}%
\pgfpathlineto{\pgfqpoint{0.902140in}{0.395283in}}%
\pgfpathlineto{\pgfqpoint{0.900276in}{0.397748in}}%
\pgfpathlineto{\pgfqpoint{0.897685in}{0.398394in}}%
\pgfpathlineto{\pgfqpoint{0.898240in}{0.401261in}}%
\pgfpathlineto{\pgfqpoint{0.897317in}{0.402234in}}%
\pgfpathlineto{\pgfqpoint{0.894866in}{0.400086in}}%
\pgfpathlineto{\pgfqpoint{0.892669in}{0.399270in}}%
\pgfpathlineto{\pgfqpoint{0.891166in}{0.400988in}}%
\pgfpathlineto{\pgfqpoint{0.889388in}{0.400332in}}%
\pgfpathlineto{\pgfqpoint{0.888087in}{0.396647in}}%
\pgfpathlineto{\pgfqpoint{0.885872in}{0.401723in}}%
\pgfpathlineto{\pgfqpoint{0.885272in}{0.397710in}}%
\pgfpathlineto{\pgfqpoint{0.881126in}{0.397511in}}%
\pgfpathlineto{\pgfqpoint{0.880525in}{0.396050in}}%
\pgfpathlineto{\pgfqpoint{0.876957in}{0.396163in}}%
\pgfpathlineto{\pgfqpoint{0.872853in}{0.397496in}}%
\pgfpathlineto{\pgfqpoint{0.866722in}{0.396548in}}%
\pgfpathlineto{\pgfqpoint{0.864048in}{0.397682in}}%
\pgfpathlineto{\pgfqpoint{0.863015in}{0.395943in}}%
\pgfpathlineto{\pgfqpoint{0.860626in}{0.397438in}}%
\pgfpathlineto{\pgfqpoint{0.860500in}{0.399959in}}%
\pgfpathlineto{\pgfqpoint{0.857986in}{0.398830in}}%
\pgfpathlineto{\pgfqpoint{0.854460in}{0.399934in}}%
\pgfpathlineto{\pgfqpoint{0.854188in}{0.405776in}}%
\pgfpathlineto{\pgfqpoint{0.857163in}{0.407525in}}%
\pgfpathlineto{\pgfqpoint{0.861136in}{0.406969in}}%
\pgfpathlineto{\pgfqpoint{0.862984in}{0.405050in}}%
\pgfpathlineto{\pgfqpoint{0.865125in}{0.405991in}}%
\pgfpathlineto{\pgfqpoint{0.867877in}{0.404847in}}%
\pgfpathlineto{\pgfqpoint{0.868825in}{0.406515in}}%
\pgfpathlineto{\pgfqpoint{0.874638in}{0.407994in}}%
\pgfpathlineto{\pgfqpoint{0.872011in}{0.409155in}}%
\pgfpathlineto{\pgfqpoint{0.868887in}{0.408859in}}%
\pgfpathlineto{\pgfqpoint{0.865823in}{0.407518in}}%
\pgfpathlineto{\pgfqpoint{0.863498in}{0.411733in}}%
\pgfpathlineto{\pgfqpoint{0.863048in}{0.414812in}}%
\pgfpathlineto{\pgfqpoint{0.869799in}{0.420224in}}%
\pgfpathlineto{\pgfqpoint{0.874638in}{0.421627in}}%
\pgfpathlineto{\pgfqpoint{0.880313in}{0.424977in}}%
\pgfpathlineto{\pgfqpoint{0.883440in}{0.427931in}}%
\pgfpathlineto{\pgfqpoint{0.884480in}{0.433448in}}%
\pgfpathlineto{\pgfqpoint{0.887231in}{0.433726in}}%
\pgfpathlineto{\pgfqpoint{0.889978in}{0.432755in}}%
\pgfpathlineto{\pgfqpoint{0.896244in}{0.433381in}}%
\pgfpathlineto{\pgfqpoint{0.902329in}{0.432923in}}%
\pgfpathlineto{\pgfqpoint{0.902058in}{0.429768in}}%
\pgfpathlineto{\pgfqpoint{0.903873in}{0.427224in}}%
\pgfpathlineto{\pgfqpoint{0.909287in}{0.427185in}}%
\pgfpathlineto{\pgfqpoint{0.907671in}{0.433445in}}%
\pgfpathlineto{\pgfqpoint{0.909967in}{0.434891in}}%
\pgfpathlineto{\pgfqpoint{0.905039in}{0.438552in}}%
\pgfpathlineto{\pgfqpoint{0.903341in}{0.440916in}}%
\pgfpathlineto{\pgfqpoint{0.897892in}{0.441660in}}%
\pgfpathlineto{\pgfqpoint{0.894137in}{0.440185in}}%
\pgfpathlineto{\pgfqpoint{0.888440in}{0.441543in}}%
\pgfpathlineto{\pgfqpoint{0.882939in}{0.440631in}}%
\pgfpathlineto{\pgfqpoint{0.881028in}{0.436394in}}%
\pgfpathlineto{\pgfqpoint{0.879917in}{0.438011in}}%
\pgfpathlineto{\pgfqpoint{0.876955in}{0.437975in}}%
\pgfpathlineto{\pgfqpoint{0.871270in}{0.436691in}}%
\pgfpathlineto{\pgfqpoint{0.868136in}{0.433288in}}%
\pgfpathlineto{\pgfqpoint{0.865686in}{0.432241in}}%
\pgfpathlineto{\pgfqpoint{0.862763in}{0.432076in}}%
\pgfpathlineto{\pgfqpoint{0.860841in}{0.433543in}}%
\pgfpathlineto{\pgfqpoint{0.860190in}{0.428478in}}%
\pgfpathlineto{\pgfqpoint{0.855724in}{0.426050in}}%
\pgfpathlineto{\pgfqpoint{0.853611in}{0.426386in}}%
\pgfpathlineto{\pgfqpoint{0.851971in}{0.427933in}}%
\pgfpathlineto{\pgfqpoint{0.849085in}{0.424656in}}%
\pgfpathlineto{\pgfqpoint{0.846580in}{0.424147in}}%
\pgfpathlineto{\pgfqpoint{0.840273in}{0.426897in}}%
\pgfpathlineto{\pgfqpoint{0.838838in}{0.425763in}}%
\pgfpathlineto{\pgfqpoint{0.836699in}{0.425925in}}%
\pgfpathlineto{\pgfqpoint{0.835510in}{0.423824in}}%
\pgfpathlineto{\pgfqpoint{0.831134in}{0.425525in}}%
\pgfpathlineto{\pgfqpoint{0.826995in}{0.422232in}}%
\pgfpathlineto{\pgfqpoint{0.824051in}{0.421527in}}%
\pgfpathlineto{\pgfqpoint{0.825544in}{0.418315in}}%
\pgfpathlineto{\pgfqpoint{0.828849in}{0.415327in}}%
\pgfpathlineto{\pgfqpoint{0.831233in}{0.409187in}}%
\pgfpathlineto{\pgfqpoint{0.831242in}{0.406233in}}%
\pgfpathlineto{\pgfqpoint{0.828822in}{0.408037in}}%
\pgfpathlineto{\pgfqpoint{0.826771in}{0.405229in}}%
\pgfpathlineto{\pgfqpoint{0.822419in}{0.408486in}}%
\pgfpathlineto{\pgfqpoint{0.820929in}{0.406477in}}%
\pgfpathlineto{\pgfqpoint{0.815079in}{0.410909in}}%
\pgfpathlineto{\pgfqpoint{0.811149in}{0.413923in}}%
\pgfpathlineto{\pgfqpoint{0.812034in}{0.416427in}}%
\pgfpathlineto{\pgfqpoint{0.818042in}{0.424203in}}%
\pgfpathlineto{\pgfqpoint{0.817435in}{0.424709in}}%
\pgfpathlineto{\pgfqpoint{0.820512in}{0.428750in}}%
\pgfpathlineto{\pgfqpoint{0.824498in}{0.425742in}}%
\pgfpathlineto{\pgfqpoint{0.827494in}{0.429779in}}%
\pgfpathlineto{\pgfqpoint{0.830966in}{0.427222in}}%
\pgfpathlineto{\pgfqpoint{0.832475in}{0.429280in}}%
\pgfpathlineto{\pgfqpoint{0.834469in}{0.427785in}}%
\pgfpathlineto{\pgfqpoint{0.837531in}{0.431857in}}%
\pgfpathlineto{\pgfqpoint{0.839547in}{0.430349in}}%
\pgfpathlineto{\pgfqpoint{0.841100in}{0.432388in}}%
\pgfpathlineto{\pgfqpoint{0.842589in}{0.431260in}}%
\pgfpathlineto{\pgfqpoint{0.848446in}{0.439078in}}%
\pgfpathlineto{\pgfqpoint{0.849915in}{0.438009in}}%
\pgfpathlineto{\pgfqpoint{0.856061in}{0.446120in}}%
\pgfpathlineto{\pgfqpoint{0.857604in}{0.444984in}}%
\pgfpathlineto{\pgfqpoint{0.863736in}{0.453184in}}%
\pgfpathlineto{\pgfqpoint{0.864761in}{0.455531in}}%
\pgfpathlineto{\pgfqpoint{0.869217in}{0.461424in}}%
\pgfpathlineto{\pgfqpoint{0.870359in}{0.464352in}}%
\pgfpathlineto{\pgfqpoint{0.874246in}{0.469508in}}%
\pgfpathlineto{\pgfqpoint{0.878374in}{0.466226in}}%
\pgfpathlineto{\pgfqpoint{0.878747in}{0.464023in}}%
\pgfpathlineto{\pgfqpoint{0.894650in}{0.485392in}}%
\pgfpathlineto{\pgfqpoint{0.902205in}{0.495355in}}%
\pgfpathlineto{\pgfqpoint{0.914054in}{0.486250in}}%
\pgfpathlineto{\pgfqpoint{0.915352in}{0.487908in}}%
\pgfpathlineto{\pgfqpoint{0.940198in}{0.483801in}}%
\pgfpathlineto{\pgfqpoint{0.950166in}{0.482212in}}%
\pgfpathlineto{\pgfqpoint{0.966684in}{0.470967in}}%
\pgfpathlineto{\pgfqpoint{0.968975in}{0.474416in}}%
\pgfpathlineto{\pgfqpoint{0.972860in}{0.471797in}}%
\pgfpathlineto{\pgfqpoint{0.986452in}{0.463301in}}%
\pgfpathlineto{\pgfqpoint{0.979580in}{0.452265in}}%
\pgfpathlineto{\pgfqpoint{0.977356in}{0.447349in}}%
\pgfpathlineto{\pgfqpoint{0.967777in}{0.431965in}}%
\pgfpathlineto{\pgfqpoint{0.961359in}{0.436162in}}%
\pgfpathlineto{\pgfqpoint{0.960218in}{0.433861in}}%
\pgfpathlineto{\pgfqpoint{0.949551in}{0.416794in}}%
\pgfpathlineto{\pgfqpoint{0.946494in}{0.418805in}}%
\pgfpathlineto{\pgfqpoint{0.945694in}{0.417578in}}%
\pgfpathlineto{\pgfqpoint{0.930988in}{0.427301in}}%
\pgfpathlineto{\pgfqpoint{0.926490in}{0.421058in}}%
\pgfpathlineto{\pgfqpoint{0.921451in}{0.413351in}}%
\pgfpathlineto{\pgfqpoint{0.918151in}{0.415401in}}%
\pgfpathlineto{\pgfqpoint{0.916327in}{0.412673in}}%
\pgfpathlineto{\pgfqpoint{0.917407in}{0.411973in}}%
\pgfpathlineto{\pgfqpoint{0.912214in}{0.404513in}}%
\pgfpathlineto{\pgfqpoint{0.913452in}{0.403668in}}%
\pgfpathlineto{\pgfqpoint{0.908663in}{0.396558in}}%
\pgfpathlineto{\pgfqpoint{0.906140in}{0.393309in}}%
\pgfpathclose%
\pgfusepath{fill}%
\end{pgfscope}%
\begin{pgfscope}%
\pgfpathrectangle{\pgfqpoint{0.100000in}{0.100000in}}{\pgfqpoint{3.608454in}{2.310000in}}%
\pgfusepath{clip}%
\pgfsetbuttcap%
\pgfsetmiterjoin%
\definecolor{currentfill}{rgb}{0.000000,0.874510,0.562745}%
\pgfsetfillcolor{currentfill}%
\pgfsetlinewidth{0.000000pt}%
\definecolor{currentstroke}{rgb}{0.000000,0.000000,0.000000}%
\pgfsetstrokecolor{currentstroke}%
\pgfsetstrokeopacity{0.000000}%
\pgfsetdash{}{0pt}%
\pgfpathmoveto{\pgfqpoint{2.901040in}{1.226866in}}%
\pgfpathlineto{\pgfqpoint{2.907489in}{1.230091in}}%
\pgfpathlineto{\pgfqpoint{2.904928in}{1.234962in}}%
\pgfpathlineto{\pgfqpoint{2.916003in}{1.242873in}}%
\pgfpathlineto{\pgfqpoint{2.908022in}{1.254600in}}%
\pgfpathlineto{\pgfqpoint{2.902185in}{1.260526in}}%
\pgfpathlineto{\pgfqpoint{2.921023in}{1.282534in}}%
\pgfpathlineto{\pgfqpoint{2.923822in}{1.280887in}}%
\pgfpathlineto{\pgfqpoint{2.923526in}{1.275100in}}%
\pgfpathlineto{\pgfqpoint{2.926305in}{1.268226in}}%
\pgfpathlineto{\pgfqpoint{2.935481in}{1.264222in}}%
\pgfpathlineto{\pgfqpoint{2.942904in}{1.258852in}}%
\pgfpathlineto{\pgfqpoint{2.950534in}{1.260497in}}%
\pgfpathlineto{\pgfqpoint{2.960443in}{1.270838in}}%
\pgfpathlineto{\pgfqpoint{2.965505in}{1.286020in}}%
\pgfpathlineto{\pgfqpoint{2.967623in}{1.287269in}}%
\pgfpathlineto{\pgfqpoint{2.968540in}{1.292624in}}%
\pgfpathlineto{\pgfqpoint{2.974604in}{1.294585in}}%
\pgfpathlineto{\pgfqpoint{2.991086in}{1.284432in}}%
\pgfpathlineto{\pgfqpoint{2.992726in}{1.278073in}}%
\pgfpathlineto{\pgfqpoint{2.983343in}{1.270885in}}%
\pgfpathlineto{\pgfqpoint{2.994849in}{1.262443in}}%
\pgfpathlineto{\pgfqpoint{2.990512in}{1.259171in}}%
\pgfpathlineto{\pgfqpoint{2.993913in}{1.252683in}}%
\pgfpathlineto{\pgfqpoint{3.005275in}{1.242589in}}%
\pgfpathlineto{\pgfqpoint{2.986436in}{1.233543in}}%
\pgfpathlineto{\pgfqpoint{2.984220in}{1.230496in}}%
\pgfpathlineto{\pgfqpoint{2.974075in}{1.228469in}}%
\pgfpathlineto{\pgfqpoint{2.968379in}{1.238517in}}%
\pgfpathlineto{\pgfqpoint{2.960170in}{1.246122in}}%
\pgfpathlineto{\pgfqpoint{2.943559in}{1.240996in}}%
\pgfpathlineto{\pgfqpoint{2.937030in}{1.237689in}}%
\pgfpathlineto{\pgfqpoint{2.933591in}{1.232567in}}%
\pgfpathlineto{\pgfqpoint{2.921871in}{1.228104in}}%
\pgfpathlineto{\pgfqpoint{2.920870in}{1.223465in}}%
\pgfpathlineto{\pgfqpoint{2.907547in}{1.214437in}}%
\pgfpathlineto{\pgfqpoint{2.901040in}{1.226866in}}%
\pgfpathclose%
\pgfusepath{fill}%
\end{pgfscope}%
\begin{pgfscope}%
\pgfpathrectangle{\pgfqpoint{0.100000in}{0.100000in}}{\pgfqpoint{3.608454in}{2.310000in}}%
\pgfusepath{clip}%
\pgfsetbuttcap%
\pgfsetmiterjoin%
\definecolor{currentfill}{rgb}{0.000000,0.478431,0.760784}%
\pgfsetfillcolor{currentfill}%
\pgfsetlinewidth{0.000000pt}%
\definecolor{currentstroke}{rgb}{0.000000,0.000000,0.000000}%
\pgfsetstrokecolor{currentstroke}%
\pgfsetstrokeopacity{0.000000}%
\pgfsetdash{}{0pt}%
\pgfpathmoveto{\pgfqpoint{1.864232in}{2.006135in}}%
\pgfpathlineto{\pgfqpoint{1.862851in}{1.978498in}}%
\pgfpathlineto{\pgfqpoint{1.864737in}{1.978421in}}%
\pgfpathlineto{\pgfqpoint{1.863197in}{1.950599in}}%
\pgfpathlineto{\pgfqpoint{1.828664in}{1.952547in}}%
\pgfpathlineto{\pgfqpoint{1.830555in}{1.980267in}}%
\pgfpathlineto{\pgfqpoint{1.828413in}{1.980392in}}%
\pgfpathlineto{\pgfqpoint{1.830021in}{2.007934in}}%
\pgfpathlineto{\pgfqpoint{1.864232in}{2.006135in}}%
\pgfpathclose%
\pgfusepath{fill}%
\end{pgfscope}%
\begin{pgfscope}%
\pgfpathrectangle{\pgfqpoint{0.100000in}{0.100000in}}{\pgfqpoint{3.608454in}{2.310000in}}%
\pgfusepath{clip}%
\pgfsetbuttcap%
\pgfsetmiterjoin%
\definecolor{currentfill}{rgb}{0.000000,0.505882,0.747059}%
\pgfsetfillcolor{currentfill}%
\pgfsetlinewidth{0.000000pt}%
\definecolor{currentstroke}{rgb}{0.000000,0.000000,0.000000}%
\pgfsetstrokecolor{currentstroke}%
\pgfsetstrokeopacity{0.000000}%
\pgfsetdash{}{0pt}%
\pgfpathmoveto{\pgfqpoint{1.728388in}{1.154348in}}%
\pgfpathlineto{\pgfqpoint{1.788727in}{1.150707in}}%
\pgfpathlineto{\pgfqpoint{1.787003in}{1.115616in}}%
\pgfpathlineto{\pgfqpoint{1.752524in}{1.117682in}}%
\pgfpathlineto{\pgfqpoint{1.683362in}{1.122232in}}%
\pgfpathlineto{\pgfqpoint{1.685892in}{1.157286in}}%
\pgfpathlineto{\pgfqpoint{1.728388in}{1.154348in}}%
\pgfpathclose%
\pgfusepath{fill}%
\end{pgfscope}%
\begin{pgfscope}%
\pgfpathrectangle{\pgfqpoint{0.100000in}{0.100000in}}{\pgfqpoint{3.608454in}{2.310000in}}%
\pgfusepath{clip}%
\pgfsetbuttcap%
\pgfsetmiterjoin%
\definecolor{currentfill}{rgb}{0.000000,0.482353,0.758824}%
\pgfsetfillcolor{currentfill}%
\pgfsetlinewidth{0.000000pt}%
\definecolor{currentstroke}{rgb}{0.000000,0.000000,0.000000}%
\pgfsetstrokecolor{currentstroke}%
\pgfsetstrokeopacity{0.000000}%
\pgfsetdash{}{0pt}%
\pgfpathmoveto{\pgfqpoint{2.615996in}{1.634560in}}%
\pgfpathlineto{\pgfqpoint{2.617722in}{1.657792in}}%
\pgfpathlineto{\pgfqpoint{2.615667in}{1.673618in}}%
\pgfpathlineto{\pgfqpoint{2.611600in}{1.688088in}}%
\pgfpathlineto{\pgfqpoint{2.600135in}{1.708748in}}%
\pgfpathlineto{\pgfqpoint{2.592486in}{1.725771in}}%
\pgfpathlineto{\pgfqpoint{2.591900in}{1.732408in}}%
\pgfpathlineto{\pgfqpoint{2.596450in}{1.744135in}}%
\pgfpathlineto{\pgfqpoint{2.618822in}{1.745974in}}%
\pgfpathlineto{\pgfqpoint{2.645951in}{1.748717in}}%
\pgfpathlineto{\pgfqpoint{2.648810in}{1.721346in}}%
\pgfpathlineto{\pgfqpoint{2.676133in}{1.724089in}}%
\pgfpathlineto{\pgfqpoint{2.689796in}{1.725591in}}%
\pgfpathlineto{\pgfqpoint{2.693396in}{1.698299in}}%
\pgfpathlineto{\pgfqpoint{2.696454in}{1.670883in}}%
\pgfpathlineto{\pgfqpoint{2.682759in}{1.669371in}}%
\pgfpathlineto{\pgfqpoint{2.655531in}{1.666251in}}%
\pgfpathlineto{\pgfqpoint{2.658513in}{1.638934in}}%
\pgfpathlineto{\pgfqpoint{2.615996in}{1.634560in}}%
\pgfpathclose%
\pgfusepath{fill}%
\end{pgfscope}%
\begin{pgfscope}%
\pgfpathrectangle{\pgfqpoint{0.100000in}{0.100000in}}{\pgfqpoint{3.608454in}{2.310000in}}%
\pgfusepath{clip}%
\pgfsetbuttcap%
\pgfsetmiterjoin%
\definecolor{currentfill}{rgb}{0.000000,0.372549,0.813725}%
\pgfsetfillcolor{currentfill}%
\pgfsetlinewidth{0.000000pt}%
\definecolor{currentstroke}{rgb}{0.000000,0.000000,0.000000}%
\pgfsetstrokecolor{currentstroke}%
\pgfsetstrokeopacity{0.000000}%
\pgfsetdash{}{0pt}%
\pgfpathmoveto{\pgfqpoint{1.148005in}{1.494799in}}%
\pgfpathlineto{\pgfqpoint{1.136590in}{1.496867in}}%
\pgfpathlineto{\pgfqpoint{1.137161in}{1.499935in}}%
\pgfpathlineto{\pgfqpoint{1.128853in}{1.510727in}}%
\pgfpathlineto{\pgfqpoint{1.131819in}{1.519099in}}%
\pgfpathlineto{\pgfqpoint{1.129198in}{1.527981in}}%
\pgfpathlineto{\pgfqpoint{1.120187in}{1.531417in}}%
\pgfpathlineto{\pgfqpoint{1.112401in}{1.545424in}}%
\pgfpathlineto{\pgfqpoint{1.116941in}{1.552445in}}%
\pgfpathlineto{\pgfqpoint{1.105133in}{1.552914in}}%
\pgfpathlineto{\pgfqpoint{1.090523in}{1.546657in}}%
\pgfpathlineto{\pgfqpoint{1.087943in}{1.551215in}}%
\pgfpathlineto{\pgfqpoint{1.077577in}{1.552756in}}%
\pgfpathlineto{\pgfqpoint{1.075863in}{1.537379in}}%
\pgfpathlineto{\pgfqpoint{1.077321in}{1.529704in}}%
\pgfpathlineto{\pgfqpoint{1.072847in}{1.517633in}}%
\pgfpathlineto{\pgfqpoint{1.061605in}{1.511149in}}%
\pgfpathlineto{\pgfqpoint{1.034454in}{1.516384in}}%
\pgfpathlineto{\pgfqpoint{0.960301in}{1.531851in}}%
\pgfpathlineto{\pgfqpoint{0.978586in}{1.616384in}}%
\pgfpathlineto{\pgfqpoint{1.051249in}{1.601475in}}%
\pgfpathlineto{\pgfqpoint{1.070321in}{1.603891in}}%
\pgfpathlineto{\pgfqpoint{1.089153in}{1.621193in}}%
\pgfpathlineto{\pgfqpoint{1.101371in}{1.618838in}}%
\pgfpathlineto{\pgfqpoint{1.110470in}{1.623700in}}%
\pgfpathlineto{\pgfqpoint{1.117387in}{1.618454in}}%
\pgfpathlineto{\pgfqpoint{1.121923in}{1.621467in}}%
\pgfpathlineto{\pgfqpoint{1.132825in}{1.619941in}}%
\pgfpathlineto{\pgfqpoint{1.136553in}{1.614347in}}%
\pgfpathlineto{\pgfqpoint{1.144102in}{1.610430in}}%
\pgfpathlineto{\pgfqpoint{1.143956in}{1.603549in}}%
\pgfpathlineto{\pgfqpoint{1.146992in}{1.596437in}}%
\pgfpathlineto{\pgfqpoint{1.157543in}{1.601688in}}%
\pgfpathlineto{\pgfqpoint{1.154007in}{1.581932in}}%
\pgfpathlineto{\pgfqpoint{1.186062in}{1.576054in}}%
\pgfpathlineto{\pgfqpoint{1.215522in}{1.571272in}}%
\pgfpathlineto{\pgfqpoint{1.213136in}{1.556950in}}%
\pgfpathlineto{\pgfqpoint{1.203263in}{1.558508in}}%
\pgfpathlineto{\pgfqpoint{1.198089in}{1.561222in}}%
\pgfpathlineto{\pgfqpoint{1.190708in}{1.558686in}}%
\pgfpathlineto{\pgfqpoint{1.173397in}{1.557849in}}%
\pgfpathlineto{\pgfqpoint{1.167943in}{1.559368in}}%
\pgfpathlineto{\pgfqpoint{1.158113in}{1.555519in}}%
\pgfpathlineto{\pgfqpoint{1.148005in}{1.494799in}}%
\pgfpathclose%
\pgfusepath{fill}%
\end{pgfscope}%
\begin{pgfscope}%
\pgfpathrectangle{\pgfqpoint{0.100000in}{0.100000in}}{\pgfqpoint{3.608454in}{2.310000in}}%
\pgfusepath{clip}%
\pgfsetbuttcap%
\pgfsetmiterjoin%
\definecolor{currentfill}{rgb}{0.000000,0.286275,0.856863}%
\pgfsetfillcolor{currentfill}%
\pgfsetlinewidth{0.000000pt}%
\definecolor{currentstroke}{rgb}{0.000000,0.000000,0.000000}%
\pgfsetstrokecolor{currentstroke}%
\pgfsetstrokeopacity{0.000000}%
\pgfsetdash{}{0pt}%
\pgfpathmoveto{\pgfqpoint{2.056854in}{1.480886in}}%
\pgfpathlineto{\pgfqpoint{2.019155in}{1.481378in}}%
\pgfpathlineto{\pgfqpoint{2.019500in}{1.502021in}}%
\pgfpathlineto{\pgfqpoint{1.992970in}{1.502588in}}%
\pgfpathlineto{\pgfqpoint{1.994270in}{1.557682in}}%
\pgfpathlineto{\pgfqpoint{2.021531in}{1.557101in}}%
\pgfpathlineto{\pgfqpoint{2.021430in}{1.552509in}}%
\pgfpathlineto{\pgfqpoint{2.040513in}{1.552171in}}%
\pgfpathlineto{\pgfqpoint{2.043988in}{1.561009in}}%
\pgfpathlineto{\pgfqpoint{2.039668in}{1.566696in}}%
\pgfpathlineto{\pgfqpoint{2.067148in}{1.566188in}}%
\pgfpathlineto{\pgfqpoint{2.073952in}{1.566162in}}%
\pgfpathlineto{\pgfqpoint{2.073777in}{1.545495in}}%
\pgfpathlineto{\pgfqpoint{2.077274in}{1.537859in}}%
\pgfpathlineto{\pgfqpoint{2.097619in}{1.537721in}}%
\pgfpathlineto{\pgfqpoint{2.097514in}{1.510367in}}%
\pgfpathlineto{\pgfqpoint{2.083971in}{1.510379in}}%
\pgfpathlineto{\pgfqpoint{2.083784in}{1.490019in}}%
\pgfpathlineto{\pgfqpoint{2.058174in}{1.490199in}}%
\pgfpathlineto{\pgfqpoint{2.056854in}{1.480886in}}%
\pgfpathclose%
\pgfusepath{fill}%
\end{pgfscope}%
\begin{pgfscope}%
\pgfpathrectangle{\pgfqpoint{0.100000in}{0.100000in}}{\pgfqpoint{3.608454in}{2.310000in}}%
\pgfusepath{clip}%
\pgfsetbuttcap%
\pgfsetmiterjoin%
\definecolor{currentfill}{rgb}{0.000000,0.490196,0.754902}%
\pgfsetfillcolor{currentfill}%
\pgfsetlinewidth{0.000000pt}%
\definecolor{currentstroke}{rgb}{0.000000,0.000000,0.000000}%
\pgfsetstrokecolor{currentstroke}%
\pgfsetstrokeopacity{0.000000}%
\pgfsetdash{}{0pt}%
\pgfpathmoveto{\pgfqpoint{2.293781in}{1.726557in}}%
\pgfpathlineto{\pgfqpoint{2.294773in}{1.699116in}}%
\pgfpathlineto{\pgfqpoint{2.253472in}{1.697776in}}%
\pgfpathlineto{\pgfqpoint{2.252724in}{1.724220in}}%
\pgfpathlineto{\pgfqpoint{2.238913in}{1.724992in}}%
\pgfpathlineto{\pgfqpoint{2.218568in}{1.724511in}}%
\pgfpathlineto{\pgfqpoint{2.217878in}{1.773970in}}%
\pgfpathlineto{\pgfqpoint{2.217797in}{1.777469in}}%
\pgfpathlineto{\pgfqpoint{2.231724in}{1.779996in}}%
\pgfpathlineto{\pgfqpoint{2.231530in}{1.786839in}}%
\pgfpathlineto{\pgfqpoint{2.241126in}{1.786250in}}%
\pgfpathlineto{\pgfqpoint{2.247155in}{1.782906in}}%
\pgfpathlineto{\pgfqpoint{2.258667in}{1.780519in}}%
\pgfpathlineto{\pgfqpoint{2.264630in}{1.772555in}}%
\pgfpathlineto{\pgfqpoint{2.273713in}{1.770000in}}%
\pgfpathlineto{\pgfqpoint{2.281318in}{1.764757in}}%
\pgfpathlineto{\pgfqpoint{2.285525in}{1.753739in}}%
\pgfpathlineto{\pgfqpoint{2.272950in}{1.753323in}}%
\pgfpathlineto{\pgfqpoint{2.273836in}{1.725888in}}%
\pgfpathlineto{\pgfqpoint{2.293781in}{1.726557in}}%
\pgfpathclose%
\pgfusepath{fill}%
\end{pgfscope}%
\begin{pgfscope}%
\pgfpathrectangle{\pgfqpoint{0.100000in}{0.100000in}}{\pgfqpoint{3.608454in}{2.310000in}}%
\pgfusepath{clip}%
\pgfsetbuttcap%
\pgfsetmiterjoin%
\definecolor{currentfill}{rgb}{0.000000,0.470588,0.764706}%
\pgfsetfillcolor{currentfill}%
\pgfsetlinewidth{0.000000pt}%
\definecolor{currentstroke}{rgb}{0.000000,0.000000,0.000000}%
\pgfsetstrokecolor{currentstroke}%
\pgfsetstrokeopacity{0.000000}%
\pgfsetdash{}{0pt}%
\pgfpathmoveto{\pgfqpoint{1.944686in}{1.753364in}}%
\pgfpathlineto{\pgfqpoint{1.999302in}{1.751993in}}%
\pgfpathlineto{\pgfqpoint{1.998767in}{1.724416in}}%
\pgfpathlineto{\pgfqpoint{1.957610in}{1.725492in}}%
\pgfpathlineto{\pgfqpoint{1.943727in}{1.725985in}}%
\pgfpathlineto{\pgfqpoint{1.944686in}{1.753364in}}%
\pgfpathclose%
\pgfusepath{fill}%
\end{pgfscope}%
\begin{pgfscope}%
\pgfpathrectangle{\pgfqpoint{0.100000in}{0.100000in}}{\pgfqpoint{3.608454in}{2.310000in}}%
\pgfusepath{clip}%
\pgfsetbuttcap%
\pgfsetmiterjoin%
\definecolor{currentfill}{rgb}{0.000000,0.501961,0.749020}%
\pgfsetfillcolor{currentfill}%
\pgfsetlinewidth{0.000000pt}%
\definecolor{currentstroke}{rgb}{0.000000,0.000000,0.000000}%
\pgfsetstrokecolor{currentstroke}%
\pgfsetstrokeopacity{0.000000}%
\pgfsetdash{}{0pt}%
\pgfpathmoveto{\pgfqpoint{2.086211in}{1.418729in}}%
\pgfpathlineto{\pgfqpoint{2.086114in}{1.391298in}}%
\pgfpathlineto{\pgfqpoint{2.072459in}{1.391367in}}%
\pgfpathlineto{\pgfqpoint{2.058778in}{1.391470in}}%
\pgfpathlineto{\pgfqpoint{2.058697in}{1.384598in}}%
\pgfpathlineto{\pgfqpoint{2.043672in}{1.384754in}}%
\pgfpathlineto{\pgfqpoint{2.031304in}{1.384920in}}%
\pgfpathlineto{\pgfqpoint{2.031759in}{1.419257in}}%
\pgfpathlineto{\pgfqpoint{2.086211in}{1.418729in}}%
\pgfpathclose%
\pgfusepath{fill}%
\end{pgfscope}%
\begin{pgfscope}%
\pgfpathrectangle{\pgfqpoint{0.100000in}{0.100000in}}{\pgfqpoint{3.608454in}{2.310000in}}%
\pgfusepath{clip}%
\pgfsetbuttcap%
\pgfsetmiterjoin%
\definecolor{currentfill}{rgb}{0.000000,0.458824,0.770588}%
\pgfsetfillcolor{currentfill}%
\pgfsetlinewidth{0.000000pt}%
\definecolor{currentstroke}{rgb}{0.000000,0.000000,0.000000}%
\pgfsetstrokecolor{currentstroke}%
\pgfsetstrokeopacity{0.000000}%
\pgfsetdash{}{0pt}%
\pgfpathmoveto{\pgfqpoint{2.414986in}{1.743088in}}%
\pgfpathlineto{\pgfqpoint{2.416473in}{1.716189in}}%
\pgfpathlineto{\pgfqpoint{2.405834in}{1.715444in}}%
\pgfpathlineto{\pgfqpoint{2.400634in}{1.726438in}}%
\pgfpathlineto{\pgfqpoint{2.395001in}{1.732218in}}%
\pgfpathlineto{\pgfqpoint{2.394954in}{1.739243in}}%
\pgfpathlineto{\pgfqpoint{2.389968in}{1.750300in}}%
\pgfpathlineto{\pgfqpoint{2.396439in}{1.763257in}}%
\pgfpathlineto{\pgfqpoint{2.373196in}{1.761948in}}%
\pgfpathlineto{\pgfqpoint{2.372249in}{1.775882in}}%
\pgfpathlineto{\pgfqpoint{2.344731in}{1.774369in}}%
\pgfpathlineto{\pgfqpoint{2.344426in}{1.781264in}}%
\pgfpathlineto{\pgfqpoint{2.337567in}{1.780978in}}%
\pgfpathlineto{\pgfqpoint{2.334657in}{1.843071in}}%
\pgfpathlineto{\pgfqpoint{2.334351in}{1.850010in}}%
\pgfpathlineto{\pgfqpoint{2.348110in}{1.850512in}}%
\pgfpathlineto{\pgfqpoint{2.383604in}{1.852652in}}%
\pgfpathlineto{\pgfqpoint{2.382793in}{1.866283in}}%
\pgfpathlineto{\pgfqpoint{2.417000in}{1.868413in}}%
\pgfpathlineto{\pgfqpoint{2.417621in}{1.861430in}}%
\pgfpathlineto{\pgfqpoint{2.419246in}{1.833798in}}%
\pgfpathlineto{\pgfqpoint{2.430551in}{1.834455in}}%
\pgfpathlineto{\pgfqpoint{2.432784in}{1.799789in}}%
\pgfpathlineto{\pgfqpoint{2.434927in}{1.765090in}}%
\pgfpathlineto{\pgfqpoint{2.413744in}{1.763959in}}%
\pgfpathlineto{\pgfqpoint{2.414986in}{1.743088in}}%
\pgfpathclose%
\pgfusepath{fill}%
\end{pgfscope}%
\begin{pgfscope}%
\pgfpathrectangle{\pgfqpoint{0.100000in}{0.100000in}}{\pgfqpoint{3.608454in}{2.310000in}}%
\pgfusepath{clip}%
\pgfsetbuttcap%
\pgfsetmiterjoin%
\definecolor{currentfill}{rgb}{0.000000,0.894118,0.552941}%
\pgfsetfillcolor{currentfill}%
\pgfsetlinewidth{0.000000pt}%
\definecolor{currentstroke}{rgb}{0.000000,0.000000,0.000000}%
\pgfsetstrokecolor{currentstroke}%
\pgfsetstrokeopacity{0.000000}%
\pgfsetdash{}{0pt}%
\pgfpathmoveto{\pgfqpoint{0.531886in}{1.661402in}}%
\pgfpathlineto{\pgfqpoint{0.522050in}{1.660114in}}%
\pgfpathlineto{\pgfqpoint{0.515901in}{1.652222in}}%
\pgfpathlineto{\pgfqpoint{0.512624in}{1.653371in}}%
\pgfpathlineto{\pgfqpoint{0.504486in}{1.647574in}}%
\pgfpathlineto{\pgfqpoint{0.490498in}{1.651799in}}%
\pgfpathlineto{\pgfqpoint{0.488413in}{1.645315in}}%
\pgfpathlineto{\pgfqpoint{0.436761in}{1.661477in}}%
\pgfpathlineto{\pgfqpoint{0.441276in}{1.674970in}}%
\pgfpathlineto{\pgfqpoint{0.442069in}{1.688570in}}%
\pgfpathlineto{\pgfqpoint{0.444780in}{1.696306in}}%
\pgfpathlineto{\pgfqpoint{0.441132in}{1.699873in}}%
\pgfpathlineto{\pgfqpoint{0.442506in}{1.705022in}}%
\pgfpathlineto{\pgfqpoint{0.448092in}{1.709303in}}%
\pgfpathlineto{\pgfqpoint{0.454296in}{1.709584in}}%
\pgfpathlineto{\pgfqpoint{0.458930in}{1.713085in}}%
\pgfpathlineto{\pgfqpoint{0.469494in}{1.715995in}}%
\pgfpathlineto{\pgfqpoint{0.468741in}{1.723271in}}%
\pgfpathlineto{\pgfqpoint{0.475334in}{1.730099in}}%
\pgfpathlineto{\pgfqpoint{0.485004in}{1.744684in}}%
\pgfpathlineto{\pgfqpoint{0.489189in}{1.745779in}}%
\pgfpathlineto{\pgfqpoint{0.494828in}{1.757961in}}%
\pgfpathlineto{\pgfqpoint{0.561272in}{1.737785in}}%
\pgfpathlineto{\pgfqpoint{0.555668in}{1.716515in}}%
\pgfpathlineto{\pgfqpoint{0.544816in}{1.681752in}}%
\pgfpathlineto{\pgfqpoint{0.535033in}{1.684697in}}%
\pgfpathlineto{\pgfqpoint{0.534540in}{1.676936in}}%
\pgfpathlineto{\pgfqpoint{0.540383in}{1.672372in}}%
\pgfpathlineto{\pgfqpoint{0.537339in}{1.664898in}}%
\pgfpathlineto{\pgfqpoint{0.531886in}{1.661402in}}%
\pgfpathclose%
\pgfusepath{fill}%
\end{pgfscope}%
\begin{pgfscope}%
\pgfpathrectangle{\pgfqpoint{0.100000in}{0.100000in}}{\pgfqpoint{3.608454in}{2.310000in}}%
\pgfusepath{clip}%
\pgfsetbuttcap%
\pgfsetmiterjoin%
\definecolor{currentfill}{rgb}{0.000000,0.415686,0.792157}%
\pgfsetfillcolor{currentfill}%
\pgfsetlinewidth{0.000000pt}%
\definecolor{currentstroke}{rgb}{0.000000,0.000000,0.000000}%
\pgfsetstrokecolor{currentstroke}%
\pgfsetstrokeopacity{0.000000}%
\pgfsetdash{}{0pt}%
\pgfpathmoveto{\pgfqpoint{2.688201in}{1.338715in}}%
\pgfpathlineto{\pgfqpoint{2.687551in}{1.344405in}}%
\pgfpathlineto{\pgfqpoint{2.680182in}{1.348259in}}%
\pgfpathlineto{\pgfqpoint{2.679482in}{1.355109in}}%
\pgfpathlineto{\pgfqpoint{2.683930in}{1.355595in}}%
\pgfpathlineto{\pgfqpoint{2.693342in}{1.363630in}}%
\pgfpathlineto{\pgfqpoint{2.708117in}{1.365276in}}%
\pgfpathlineto{\pgfqpoint{2.710112in}{1.347801in}}%
\pgfpathlineto{\pgfqpoint{2.720376in}{1.354526in}}%
\pgfpathlineto{\pgfqpoint{2.726746in}{1.347306in}}%
\pgfpathlineto{\pgfqpoint{2.718732in}{1.341235in}}%
\pgfpathlineto{\pgfqpoint{2.713052in}{1.339727in}}%
\pgfpathlineto{\pgfqpoint{2.707002in}{1.331992in}}%
\pgfpathlineto{\pgfqpoint{2.697292in}{1.333124in}}%
\pgfpathlineto{\pgfqpoint{2.697081in}{1.338065in}}%
\pgfpathlineto{\pgfqpoint{2.688201in}{1.338715in}}%
\pgfpathclose%
\pgfusepath{fill}%
\end{pgfscope}%
\begin{pgfscope}%
\pgfpathrectangle{\pgfqpoint{0.100000in}{0.100000in}}{\pgfqpoint{3.608454in}{2.310000in}}%
\pgfusepath{clip}%
\pgfsetbuttcap%
\pgfsetmiterjoin%
\definecolor{currentfill}{rgb}{0.000000,0.435294,0.782353}%
\pgfsetfillcolor{currentfill}%
\pgfsetlinewidth{0.000000pt}%
\definecolor{currentstroke}{rgb}{0.000000,0.000000,0.000000}%
\pgfsetstrokecolor{currentstroke}%
\pgfsetstrokeopacity{0.000000}%
\pgfsetdash{}{0pt}%
\pgfpathmoveto{\pgfqpoint{1.912258in}{1.559919in}}%
\pgfpathlineto{\pgfqpoint{1.911513in}{1.532437in}}%
\pgfpathlineto{\pgfqpoint{1.884664in}{1.533508in}}%
\pgfpathlineto{\pgfqpoint{1.857457in}{1.534695in}}%
\pgfpathlineto{\pgfqpoint{1.858208in}{1.562101in}}%
\pgfpathlineto{\pgfqpoint{1.830208in}{1.563450in}}%
\pgfpathlineto{\pgfqpoint{1.831612in}{1.590776in}}%
\pgfpathlineto{\pgfqpoint{1.858155in}{1.589614in}}%
\pgfpathlineto{\pgfqpoint{1.885929in}{1.588436in}}%
\pgfpathlineto{\pgfqpoint{1.884919in}{1.560911in}}%
\pgfpathlineto{\pgfqpoint{1.912258in}{1.559919in}}%
\pgfpathclose%
\pgfusepath{fill}%
\end{pgfscope}%
\begin{pgfscope}%
\pgfpathrectangle{\pgfqpoint{0.100000in}{0.100000in}}{\pgfqpoint{3.608454in}{2.310000in}}%
\pgfusepath{clip}%
\pgfsetbuttcap%
\pgfsetmiterjoin%
\definecolor{currentfill}{rgb}{0.000000,0.313725,0.843137}%
\pgfsetfillcolor{currentfill}%
\pgfsetlinewidth{0.000000pt}%
\definecolor{currentstroke}{rgb}{0.000000,0.000000,0.000000}%
\pgfsetstrokecolor{currentstroke}%
\pgfsetstrokeopacity{0.000000}%
\pgfsetdash{}{0pt}%
\pgfpathmoveto{\pgfqpoint{1.968280in}{1.606524in}}%
\pgfpathlineto{\pgfqpoint{1.988757in}{1.605977in}}%
\pgfpathlineto{\pgfqpoint{1.989587in}{1.638444in}}%
\pgfpathlineto{\pgfqpoint{2.001679in}{1.633660in}}%
\pgfpathlineto{\pgfqpoint{2.002309in}{1.663696in}}%
\pgfpathlineto{\pgfqpoint{2.022652in}{1.663329in}}%
\pgfpathlineto{\pgfqpoint{2.019320in}{1.660716in}}%
\pgfpathlineto{\pgfqpoint{2.019889in}{1.653420in}}%
\pgfpathlineto{\pgfqpoint{2.017541in}{1.649592in}}%
\pgfpathlineto{\pgfqpoint{2.056914in}{1.649093in}}%
\pgfpathlineto{\pgfqpoint{2.056624in}{1.621513in}}%
\pgfpathlineto{\pgfqpoint{2.063540in}{1.621479in}}%
\pgfpathlineto{\pgfqpoint{2.067679in}{1.614586in}}%
\pgfpathlineto{\pgfqpoint{2.067449in}{1.593735in}}%
\pgfpathlineto{\pgfqpoint{2.067148in}{1.566188in}}%
\pgfpathlineto{\pgfqpoint{2.039668in}{1.566696in}}%
\pgfpathlineto{\pgfqpoint{2.040288in}{1.575134in}}%
\pgfpathlineto{\pgfqpoint{2.036665in}{1.576098in}}%
\pgfpathlineto{\pgfqpoint{2.027934in}{1.590280in}}%
\pgfpathlineto{\pgfqpoint{2.027350in}{1.599345in}}%
\pgfpathlineto{\pgfqpoint{2.005666in}{1.599833in}}%
\pgfpathlineto{\pgfqpoint{2.000030in}{1.598859in}}%
\pgfpathlineto{\pgfqpoint{1.999669in}{1.585100in}}%
\pgfpathlineto{\pgfqpoint{1.988196in}{1.585344in}}%
\pgfpathlineto{\pgfqpoint{1.967706in}{1.585865in}}%
\pgfpathlineto{\pgfqpoint{1.968280in}{1.606524in}}%
\pgfpathclose%
\pgfusepath{fill}%
\end{pgfscope}%
\begin{pgfscope}%
\pgfpathrectangle{\pgfqpoint{0.100000in}{0.100000in}}{\pgfqpoint{3.608454in}{2.310000in}}%
\pgfusepath{clip}%
\pgfsetbuttcap%
\pgfsetmiterjoin%
\definecolor{currentfill}{rgb}{0.000000,0.650980,0.674510}%
\pgfsetfillcolor{currentfill}%
\pgfsetlinewidth{0.000000pt}%
\definecolor{currentstroke}{rgb}{0.000000,0.000000,0.000000}%
\pgfsetstrokecolor{currentstroke}%
\pgfsetstrokeopacity{0.000000}%
\pgfsetdash{}{0pt}%
\pgfpathmoveto{\pgfqpoint{3.051321in}{0.914811in}}%
\pgfpathlineto{\pgfqpoint{3.047304in}{0.907519in}}%
\pgfpathlineto{\pgfqpoint{3.038277in}{0.916697in}}%
\pgfpathlineto{\pgfqpoint{3.022585in}{0.898037in}}%
\pgfpathlineto{\pgfqpoint{3.015378in}{0.900512in}}%
\pgfpathlineto{\pgfqpoint{3.011589in}{0.908504in}}%
\pgfpathlineto{\pgfqpoint{3.010014in}{0.919433in}}%
\pgfpathlineto{\pgfqpoint{3.003885in}{0.926116in}}%
\pgfpathlineto{\pgfqpoint{3.003733in}{0.931980in}}%
\pgfpathlineto{\pgfqpoint{2.999992in}{0.934399in}}%
\pgfpathlineto{\pgfqpoint{2.988444in}{0.938374in}}%
\pgfpathlineto{\pgfqpoint{2.984051in}{0.944289in}}%
\pgfpathlineto{\pgfqpoint{3.005840in}{0.971004in}}%
\pgfpathlineto{\pgfqpoint{3.010939in}{0.970937in}}%
\pgfpathlineto{\pgfqpoint{3.018836in}{0.966601in}}%
\pgfpathlineto{\pgfqpoint{3.034445in}{0.961587in}}%
\pgfpathlineto{\pgfqpoint{3.047350in}{0.952216in}}%
\pgfpathlineto{\pgfqpoint{3.030466in}{0.937480in}}%
\pgfpathlineto{\pgfqpoint{3.038279in}{0.928637in}}%
\pgfpathlineto{\pgfqpoint{3.044811in}{0.924182in}}%
\pgfpathlineto{\pgfqpoint{3.051321in}{0.914811in}}%
\pgfpathclose%
\pgfusepath{fill}%
\end{pgfscope}%
\begin{pgfscope}%
\pgfpathrectangle{\pgfqpoint{0.100000in}{0.100000in}}{\pgfqpoint{3.608454in}{2.310000in}}%
\pgfusepath{clip}%
\pgfsetbuttcap%
\pgfsetmiterjoin%
\definecolor{currentfill}{rgb}{0.000000,0.478431,0.760784}%
\pgfsetfillcolor{currentfill}%
\pgfsetlinewidth{0.000000pt}%
\definecolor{currentstroke}{rgb}{0.000000,0.000000,0.000000}%
\pgfsetstrokecolor{currentstroke}%
\pgfsetstrokeopacity{0.000000}%
\pgfsetdash{}{0pt}%
\pgfpathmoveto{\pgfqpoint{1.046091in}{0.176022in}}%
\pgfpathlineto{\pgfqpoint{1.046239in}{0.179505in}}%
\pgfpathlineto{\pgfqpoint{1.048650in}{0.177612in}}%
\pgfpathlineto{\pgfqpoint{1.046091in}{0.176022in}}%
\pgfpathclose%
\pgfusepath{fill}%
\end{pgfscope}%
\begin{pgfscope}%
\pgfpathrectangle{\pgfqpoint{0.100000in}{0.100000in}}{\pgfqpoint{3.608454in}{2.310000in}}%
\pgfusepath{clip}%
\pgfsetbuttcap%
\pgfsetmiterjoin%
\definecolor{currentfill}{rgb}{0.000000,0.478431,0.760784}%
\pgfsetfillcolor{currentfill}%
\pgfsetlinewidth{0.000000pt}%
\definecolor{currentstroke}{rgb}{0.000000,0.000000,0.000000}%
\pgfsetstrokecolor{currentstroke}%
\pgfsetstrokeopacity{0.000000}%
\pgfsetdash{}{0pt}%
\pgfpathmoveto{\pgfqpoint{1.042447in}{0.162554in}}%
\pgfpathlineto{\pgfqpoint{1.043313in}{0.165192in}}%
\pgfpathlineto{\pgfqpoint{1.044966in}{0.165456in}}%
\pgfpathlineto{\pgfqpoint{1.044338in}{0.161816in}}%
\pgfpathlineto{\pgfqpoint{1.042447in}{0.162554in}}%
\pgfpathclose%
\pgfusepath{fill}%
\end{pgfscope}%
\begin{pgfscope}%
\pgfpathrectangle{\pgfqpoint{0.100000in}{0.100000in}}{\pgfqpoint{3.608454in}{2.310000in}}%
\pgfusepath{clip}%
\pgfsetbuttcap%
\pgfsetmiterjoin%
\definecolor{currentfill}{rgb}{0.000000,0.478431,0.760784}%
\pgfsetfillcolor{currentfill}%
\pgfsetlinewidth{0.000000pt}%
\definecolor{currentstroke}{rgb}{0.000000,0.000000,0.000000}%
\pgfsetstrokecolor{currentstroke}%
\pgfsetstrokeopacity{0.000000}%
\pgfsetdash{}{0pt}%
\pgfpathmoveto{\pgfqpoint{1.045943in}{0.145022in}}%
\pgfpathlineto{\pgfqpoint{1.044467in}{0.150857in}}%
\pgfpathlineto{\pgfqpoint{1.045363in}{0.152527in}}%
\pgfpathlineto{\pgfqpoint{1.043425in}{0.154046in}}%
\pgfpathlineto{\pgfqpoint{1.045163in}{0.157939in}}%
\pgfpathlineto{\pgfqpoint{1.045050in}{0.160374in}}%
\pgfpathlineto{\pgfqpoint{1.047193in}{0.161279in}}%
\pgfpathlineto{\pgfqpoint{1.047667in}{0.157760in}}%
\pgfpathlineto{\pgfqpoint{1.045744in}{0.156113in}}%
\pgfpathlineto{\pgfqpoint{1.046983in}{0.154946in}}%
\pgfpathlineto{\pgfqpoint{1.047011in}{0.149220in}}%
\pgfpathlineto{\pgfqpoint{1.049976in}{0.148483in}}%
\pgfpathlineto{\pgfqpoint{1.045943in}{0.145022in}}%
\pgfpathclose%
\pgfusepath{fill}%
\end{pgfscope}%
\begin{pgfscope}%
\pgfpathrectangle{\pgfqpoint{0.100000in}{0.100000in}}{\pgfqpoint{3.608454in}{2.310000in}}%
\pgfusepath{clip}%
\pgfsetbuttcap%
\pgfsetmiterjoin%
\definecolor{currentfill}{rgb}{0.000000,0.478431,0.760784}%
\pgfsetfillcolor{currentfill}%
\pgfsetlinewidth{0.000000pt}%
\definecolor{currentstroke}{rgb}{0.000000,0.000000,0.000000}%
\pgfsetstrokecolor{currentstroke}%
\pgfsetstrokeopacity{0.000000}%
\pgfsetdash{}{0pt}%
\pgfpathmoveto{\pgfqpoint{1.071810in}{0.150553in}}%
\pgfpathlineto{\pgfqpoint{1.072025in}{0.147586in}}%
\pgfpathlineto{\pgfqpoint{1.070746in}{0.144824in}}%
\pgfpathlineto{\pgfqpoint{1.069328in}{0.144797in}}%
\pgfpathlineto{\pgfqpoint{1.068043in}{0.147923in}}%
\pgfpathlineto{\pgfqpoint{1.069455in}{0.148655in}}%
\pgfpathlineto{\pgfqpoint{1.069622in}{0.151729in}}%
\pgfpathlineto{\pgfqpoint{1.071810in}{0.150553in}}%
\pgfpathclose%
\pgfusepath{fill}%
\end{pgfscope}%
\begin{pgfscope}%
\pgfpathrectangle{\pgfqpoint{0.100000in}{0.100000in}}{\pgfqpoint{3.608454in}{2.310000in}}%
\pgfusepath{clip}%
\pgfsetbuttcap%
\pgfsetmiterjoin%
\definecolor{currentfill}{rgb}{0.000000,0.478431,0.760784}%
\pgfsetfillcolor{currentfill}%
\pgfsetlinewidth{0.000000pt}%
\definecolor{currentstroke}{rgb}{0.000000,0.000000,0.000000}%
\pgfsetstrokecolor{currentstroke}%
\pgfsetstrokeopacity{0.000000}%
\pgfsetdash{}{0pt}%
\pgfpathmoveto{\pgfqpoint{1.065482in}{0.169372in}}%
\pgfpathlineto{\pgfqpoint{1.067405in}{0.165345in}}%
\pgfpathlineto{\pgfqpoint{1.069320in}{0.165633in}}%
\pgfpathlineto{\pgfqpoint{1.070834in}{0.168040in}}%
\pgfpathlineto{\pgfqpoint{1.070260in}{0.170374in}}%
\pgfpathlineto{\pgfqpoint{1.072944in}{0.170978in}}%
\pgfpathlineto{\pgfqpoint{1.073238in}{0.172315in}}%
\pgfpathlineto{\pgfqpoint{1.075481in}{0.173524in}}%
\pgfpathlineto{\pgfqpoint{1.078315in}{0.173884in}}%
\pgfpathlineto{\pgfqpoint{1.079838in}{0.170598in}}%
\pgfpathlineto{\pgfqpoint{1.082038in}{0.171601in}}%
\pgfpathlineto{\pgfqpoint{1.081805in}{0.173502in}}%
\pgfpathlineto{\pgfqpoint{1.085139in}{0.174080in}}%
\pgfpathlineto{\pgfqpoint{1.087857in}{0.173091in}}%
\pgfpathlineto{\pgfqpoint{1.089922in}{0.175542in}}%
\pgfpathlineto{\pgfqpoint{1.090049in}{0.178164in}}%
\pgfpathlineto{\pgfqpoint{1.093019in}{0.176002in}}%
\pgfpathlineto{\pgfqpoint{1.096013in}{0.170850in}}%
\pgfpathlineto{\pgfqpoint{1.096478in}{0.167420in}}%
\pgfpathlineto{\pgfqpoint{1.098614in}{0.165154in}}%
\pgfpathlineto{\pgfqpoint{1.100805in}{0.164982in}}%
\pgfpathlineto{\pgfqpoint{1.101142in}{0.162059in}}%
\pgfpathlineto{\pgfqpoint{1.100030in}{0.159422in}}%
\pgfpathlineto{\pgfqpoint{1.096111in}{0.156382in}}%
\pgfpathlineto{\pgfqpoint{1.095842in}{0.153748in}}%
\pgfpathlineto{\pgfqpoint{1.094580in}{0.151215in}}%
\pgfpathlineto{\pgfqpoint{1.094242in}{0.146717in}}%
\pgfpathlineto{\pgfqpoint{1.093131in}{0.142662in}}%
\pgfpathlineto{\pgfqpoint{1.090632in}{0.141147in}}%
\pgfpathlineto{\pgfqpoint{1.084308in}{0.135553in}}%
\pgfpathlineto{\pgfqpoint{1.081972in}{0.136078in}}%
\pgfpathlineto{\pgfqpoint{1.079389in}{0.134752in}}%
\pgfpathlineto{\pgfqpoint{1.077089in}{0.136527in}}%
\pgfpathlineto{\pgfqpoint{1.074765in}{0.136279in}}%
\pgfpathlineto{\pgfqpoint{1.076502in}{0.144158in}}%
\pgfpathlineto{\pgfqpoint{1.076368in}{0.146457in}}%
\pgfpathlineto{\pgfqpoint{1.079906in}{0.148920in}}%
\pgfpathlineto{\pgfqpoint{1.081825in}{0.148421in}}%
\pgfpathlineto{\pgfqpoint{1.081856in}{0.151528in}}%
\pgfpathlineto{\pgfqpoint{1.083827in}{0.155130in}}%
\pgfpathlineto{\pgfqpoint{1.085699in}{0.162221in}}%
\pgfpathlineto{\pgfqpoint{1.083172in}{0.161819in}}%
\pgfpathlineto{\pgfqpoint{1.083697in}{0.159212in}}%
\pgfpathlineto{\pgfqpoint{1.082643in}{0.155646in}}%
\pgfpathlineto{\pgfqpoint{1.078109in}{0.149553in}}%
\pgfpathlineto{\pgfqpoint{1.074518in}{0.148168in}}%
\pgfpathlineto{\pgfqpoint{1.073929in}{0.150391in}}%
\pgfpathlineto{\pgfqpoint{1.070564in}{0.153361in}}%
\pgfpathlineto{\pgfqpoint{1.067603in}{0.152672in}}%
\pgfpathlineto{\pgfqpoint{1.065718in}{0.151421in}}%
\pgfpathlineto{\pgfqpoint{1.067029in}{0.156508in}}%
\pgfpathlineto{\pgfqpoint{1.068970in}{0.159001in}}%
\pgfpathlineto{\pgfqpoint{1.070918in}{0.159489in}}%
\pgfpathlineto{\pgfqpoint{1.072200in}{0.162116in}}%
\pgfpathlineto{\pgfqpoint{1.074427in}{0.163792in}}%
\pgfpathlineto{\pgfqpoint{1.074481in}{0.165957in}}%
\pgfpathlineto{\pgfqpoint{1.077573in}{0.166662in}}%
\pgfpathlineto{\pgfqpoint{1.077538in}{0.168796in}}%
\pgfpathlineto{\pgfqpoint{1.074777in}{0.170290in}}%
\pgfpathlineto{\pgfqpoint{1.071926in}{0.169175in}}%
\pgfpathlineto{\pgfqpoint{1.073508in}{0.167354in}}%
\pgfpathlineto{\pgfqpoint{1.070989in}{0.164095in}}%
\pgfpathlineto{\pgfqpoint{1.067088in}{0.160807in}}%
\pgfpathlineto{\pgfqpoint{1.065060in}{0.165053in}}%
\pgfpathlineto{\pgfqpoint{1.065482in}{0.169372in}}%
\pgfpathclose%
\pgfusepath{fill}%
\end{pgfscope}%
\begin{pgfscope}%
\pgfpathrectangle{\pgfqpoint{0.100000in}{0.100000in}}{\pgfqpoint{3.608454in}{2.310000in}}%
\pgfusepath{clip}%
\pgfsetbuttcap%
\pgfsetmiterjoin%
\definecolor{currentfill}{rgb}{0.000000,0.478431,0.760784}%
\pgfsetfillcolor{currentfill}%
\pgfsetlinewidth{0.000000pt}%
\definecolor{currentstroke}{rgb}{0.000000,0.000000,0.000000}%
\pgfsetstrokecolor{currentstroke}%
\pgfsetstrokeopacity{0.000000}%
\pgfsetdash{}{0pt}%
\pgfpathmoveto{\pgfqpoint{1.041943in}{0.169830in}}%
\pgfpathlineto{\pgfqpoint{1.039521in}{0.170855in}}%
\pgfpathlineto{\pgfqpoint{1.042951in}{0.172853in}}%
\pgfpathlineto{\pgfqpoint{1.043959in}{0.170752in}}%
\pgfpathlineto{\pgfqpoint{1.047379in}{0.171633in}}%
\pgfpathlineto{\pgfqpoint{1.048310in}{0.169620in}}%
\pgfpathlineto{\pgfqpoint{1.045621in}{0.169580in}}%
\pgfpathlineto{\pgfqpoint{1.044188in}{0.167269in}}%
\pgfpathlineto{\pgfqpoint{1.041932in}{0.166429in}}%
\pgfpathlineto{\pgfqpoint{1.040306in}{0.167304in}}%
\pgfpathlineto{\pgfqpoint{1.041943in}{0.169830in}}%
\pgfpathclose%
\pgfusepath{fill}%
\end{pgfscope}%
\begin{pgfscope}%
\pgfpathrectangle{\pgfqpoint{0.100000in}{0.100000in}}{\pgfqpoint{3.608454in}{2.310000in}}%
\pgfusepath{clip}%
\pgfsetbuttcap%
\pgfsetmiterjoin%
\definecolor{currentfill}{rgb}{0.000000,0.478431,0.760784}%
\pgfsetfillcolor{currentfill}%
\pgfsetlinewidth{0.000000pt}%
\definecolor{currentstroke}{rgb}{0.000000,0.000000,0.000000}%
\pgfsetstrokecolor{currentstroke}%
\pgfsetstrokeopacity{0.000000}%
\pgfsetdash{}{0pt}%
\pgfpathmoveto{\pgfqpoint{1.051425in}{0.157660in}}%
\pgfpathlineto{\pgfqpoint{1.051517in}{0.160010in}}%
\pgfpathlineto{\pgfqpoint{1.046502in}{0.163773in}}%
\pgfpathlineto{\pgfqpoint{1.050660in}{0.165847in}}%
\pgfpathlineto{\pgfqpoint{1.049847in}{0.167856in}}%
\pgfpathlineto{\pgfqpoint{1.050846in}{0.170035in}}%
\pgfpathlineto{\pgfqpoint{1.048993in}{0.170836in}}%
\pgfpathlineto{\pgfqpoint{1.047832in}{0.173344in}}%
\pgfpathlineto{\pgfqpoint{1.051430in}{0.175827in}}%
\pgfpathlineto{\pgfqpoint{1.051845in}{0.178802in}}%
\pgfpathlineto{\pgfqpoint{1.053869in}{0.179753in}}%
\pgfpathlineto{\pgfqpoint{1.053809in}{0.182723in}}%
\pgfpathlineto{\pgfqpoint{1.051267in}{0.183712in}}%
\pgfpathlineto{\pgfqpoint{1.045462in}{0.182905in}}%
\pgfpathlineto{\pgfqpoint{1.047279in}{0.185118in}}%
\pgfpathlineto{\pgfqpoint{1.049511in}{0.186205in}}%
\pgfpathlineto{\pgfqpoint{1.052044in}{0.186328in}}%
\pgfpathlineto{\pgfqpoint{1.051609in}{0.190204in}}%
\pgfpathlineto{\pgfqpoint{1.052862in}{0.192676in}}%
\pgfpathlineto{\pgfqpoint{1.054788in}{0.191850in}}%
\pgfpathlineto{\pgfqpoint{1.059113in}{0.188805in}}%
\pgfpathlineto{\pgfqpoint{1.059157in}{0.184625in}}%
\pgfpathlineto{\pgfqpoint{1.057111in}{0.183014in}}%
\pgfpathlineto{\pgfqpoint{1.057987in}{0.179356in}}%
\pgfpathlineto{\pgfqpoint{1.059495in}{0.179632in}}%
\pgfpathlineto{\pgfqpoint{1.061360in}{0.177220in}}%
\pgfpathlineto{\pgfqpoint{1.062648in}{0.172710in}}%
\pgfpathlineto{\pgfqpoint{1.062360in}{0.170277in}}%
\pgfpathlineto{\pgfqpoint{1.060929in}{0.168990in}}%
\pgfpathlineto{\pgfqpoint{1.062884in}{0.167409in}}%
\pgfpathlineto{\pgfqpoint{1.062703in}{0.164914in}}%
\pgfpathlineto{\pgfqpoint{1.060871in}{0.164522in}}%
\pgfpathlineto{\pgfqpoint{1.059752in}{0.166916in}}%
\pgfpathlineto{\pgfqpoint{1.058308in}{0.164637in}}%
\pgfpathlineto{\pgfqpoint{1.061166in}{0.162445in}}%
\pgfpathlineto{\pgfqpoint{1.061172in}{0.160447in}}%
\pgfpathlineto{\pgfqpoint{1.062512in}{0.157360in}}%
\pgfpathlineto{\pgfqpoint{1.060934in}{0.153696in}}%
\pgfpathlineto{\pgfqpoint{1.064018in}{0.154722in}}%
\pgfpathlineto{\pgfqpoint{1.063461in}{0.152526in}}%
\pgfpathlineto{\pgfqpoint{1.061554in}{0.151190in}}%
\pgfpathlineto{\pgfqpoint{1.060026in}{0.151947in}}%
\pgfpathlineto{\pgfqpoint{1.058945in}{0.148455in}}%
\pgfpathlineto{\pgfqpoint{1.061641in}{0.148773in}}%
\pgfpathlineto{\pgfqpoint{1.059952in}{0.145277in}}%
\pgfpathlineto{\pgfqpoint{1.059401in}{0.142596in}}%
\pgfpathlineto{\pgfqpoint{1.057983in}{0.140914in}}%
\pgfpathlineto{\pgfqpoint{1.055057in}{0.141929in}}%
\pgfpathlineto{\pgfqpoint{1.054548in}{0.147234in}}%
\pgfpathlineto{\pgfqpoint{1.052684in}{0.150955in}}%
\pgfpathlineto{\pgfqpoint{1.052499in}{0.153302in}}%
\pgfpathlineto{\pgfqpoint{1.054358in}{0.157385in}}%
\pgfpathlineto{\pgfqpoint{1.051425in}{0.157660in}}%
\pgfpathclose%
\pgfusepath{fill}%
\end{pgfscope}%
\begin{pgfscope}%
\pgfpathrectangle{\pgfqpoint{0.100000in}{0.100000in}}{\pgfqpoint{3.608454in}{2.310000in}}%
\pgfusepath{clip}%
\pgfsetbuttcap%
\pgfsetmiterjoin%
\definecolor{currentfill}{rgb}{0.000000,0.682353,0.658824}%
\pgfsetfillcolor{currentfill}%
\pgfsetlinewidth{0.000000pt}%
\definecolor{currentstroke}{rgb}{0.000000,0.000000,0.000000}%
\pgfsetstrokecolor{currentstroke}%
\pgfsetstrokeopacity{0.000000}%
\pgfsetdash{}{0pt}%
\pgfpathmoveto{\pgfqpoint{0.875234in}{2.151233in}}%
\pgfpathlineto{\pgfqpoint{0.868405in}{2.124498in}}%
\pgfpathlineto{\pgfqpoint{0.863100in}{2.115696in}}%
\pgfpathlineto{\pgfqpoint{0.851613in}{2.114354in}}%
\pgfpathlineto{\pgfqpoint{0.823108in}{2.121596in}}%
\pgfpathlineto{\pgfqpoint{0.790335in}{2.130635in}}%
\pgfpathlineto{\pgfqpoint{0.784697in}{2.127385in}}%
\pgfpathlineto{\pgfqpoint{0.782427in}{2.131997in}}%
\pgfpathlineto{\pgfqpoint{0.774845in}{2.127276in}}%
\pgfpathlineto{\pgfqpoint{0.761418in}{2.129669in}}%
\pgfpathlineto{\pgfqpoint{0.758684in}{2.134138in}}%
\pgfpathlineto{\pgfqpoint{0.762232in}{2.143437in}}%
\pgfpathlineto{\pgfqpoint{0.763393in}{2.155197in}}%
\pgfpathlineto{\pgfqpoint{0.762473in}{2.166902in}}%
\pgfpathlineto{\pgfqpoint{0.767233in}{2.177075in}}%
\pgfpathlineto{\pgfqpoint{0.769524in}{2.184730in}}%
\pgfpathlineto{\pgfqpoint{0.774962in}{2.186885in}}%
\pgfpathlineto{\pgfqpoint{0.778779in}{2.191862in}}%
\pgfpathlineto{\pgfqpoint{0.796683in}{2.187200in}}%
\pgfpathlineto{\pgfqpoint{0.798560in}{2.193839in}}%
\pgfpathlineto{\pgfqpoint{0.807006in}{2.198546in}}%
\pgfpathlineto{\pgfqpoint{0.810303in}{2.197629in}}%
\pgfpathlineto{\pgfqpoint{0.814986in}{2.205959in}}%
\pgfpathlineto{\pgfqpoint{0.820986in}{2.210320in}}%
\pgfpathlineto{\pgfqpoint{0.822836in}{2.217028in}}%
\pgfpathlineto{\pgfqpoint{0.829989in}{2.221231in}}%
\pgfpathlineto{\pgfqpoint{0.836288in}{2.219545in}}%
\pgfpathlineto{\pgfqpoint{0.836363in}{2.217995in}}%
\pgfpathlineto{\pgfqpoint{0.821826in}{2.165426in}}%
\pgfpathlineto{\pgfqpoint{0.875234in}{2.151233in}}%
\pgfpathclose%
\pgfusepath{fill}%
\end{pgfscope}%
\begin{pgfscope}%
\pgfpathrectangle{\pgfqpoint{0.100000in}{0.100000in}}{\pgfqpoint{3.608454in}{2.310000in}}%
\pgfusepath{clip}%
\pgfsetbuttcap%
\pgfsetmiterjoin%
\definecolor{currentfill}{rgb}{0.000000,0.709804,0.645098}%
\pgfsetfillcolor{currentfill}%
\pgfsetlinewidth{0.000000pt}%
\definecolor{currentstroke}{rgb}{0.000000,0.000000,0.000000}%
\pgfsetstrokecolor{currentstroke}%
\pgfsetstrokeopacity{0.000000}%
\pgfsetdash{}{0pt}%
\pgfpathmoveto{\pgfqpoint{2.966480in}{1.429783in}}%
\pgfpathlineto{\pgfqpoint{2.962415in}{1.436530in}}%
\pgfpathlineto{\pgfqpoint{2.951761in}{1.443677in}}%
\pgfpathlineto{\pgfqpoint{2.954171in}{1.450091in}}%
\pgfpathlineto{\pgfqpoint{2.941441in}{1.448688in}}%
\pgfpathlineto{\pgfqpoint{2.938939in}{1.450447in}}%
\pgfpathlineto{\pgfqpoint{2.935564in}{1.458054in}}%
\pgfpathlineto{\pgfqpoint{2.933975in}{1.470774in}}%
\pgfpathlineto{\pgfqpoint{2.938569in}{1.471342in}}%
\pgfpathlineto{\pgfqpoint{2.935644in}{1.495065in}}%
\pgfpathlineto{\pgfqpoint{2.956251in}{1.497406in}}%
\pgfpathlineto{\pgfqpoint{2.977955in}{1.500863in}}%
\pgfpathlineto{\pgfqpoint{2.983438in}{1.466654in}}%
\pgfpathlineto{\pgfqpoint{2.989320in}{1.467586in}}%
\pgfpathlineto{\pgfqpoint{2.992195in}{1.461918in}}%
\pgfpathlineto{\pgfqpoint{2.986701in}{1.454490in}}%
\pgfpathlineto{\pgfqpoint{2.989302in}{1.450132in}}%
\pgfpathlineto{\pgfqpoint{2.987075in}{1.444310in}}%
\pgfpathlineto{\pgfqpoint{2.980855in}{1.444545in}}%
\pgfpathlineto{\pgfqpoint{2.966480in}{1.429783in}}%
\pgfpathclose%
\pgfusepath{fill}%
\end{pgfscope}%
\begin{pgfscope}%
\pgfpathrectangle{\pgfqpoint{0.100000in}{0.100000in}}{\pgfqpoint{3.608454in}{2.310000in}}%
\pgfusepath{clip}%
\pgfsetbuttcap%
\pgfsetmiterjoin%
\definecolor{currentfill}{rgb}{0.000000,0.411765,0.794118}%
\pgfsetfillcolor{currentfill}%
\pgfsetlinewidth{0.000000pt}%
\definecolor{currentstroke}{rgb}{0.000000,0.000000,0.000000}%
\pgfsetstrokecolor{currentstroke}%
\pgfsetstrokeopacity{0.000000}%
\pgfsetdash{}{0pt}%
\pgfpathmoveto{\pgfqpoint{3.466647in}{1.697217in}}%
\pgfpathlineto{\pgfqpoint{3.468281in}{1.705449in}}%
\pgfpathlineto{\pgfqpoint{3.464103in}{1.723936in}}%
\pgfpathlineto{\pgfqpoint{3.456367in}{1.750710in}}%
\pgfpathlineto{\pgfqpoint{3.479909in}{1.757791in}}%
\pgfpathlineto{\pgfqpoint{3.481572in}{1.755490in}}%
\pgfpathlineto{\pgfqpoint{3.495434in}{1.768247in}}%
\pgfpathlineto{\pgfqpoint{3.499874in}{1.758602in}}%
\pgfpathlineto{\pgfqpoint{3.513375in}{1.745585in}}%
\pgfpathlineto{\pgfqpoint{3.518613in}{1.736187in}}%
\pgfpathlineto{\pgfqpoint{3.513987in}{1.733317in}}%
\pgfpathlineto{\pgfqpoint{3.514812in}{1.727935in}}%
\pgfpathlineto{\pgfqpoint{3.510151in}{1.723018in}}%
\pgfpathlineto{\pgfqpoint{3.494576in}{1.717460in}}%
\pgfpathlineto{\pgfqpoint{3.494096in}{1.728843in}}%
\pgfpathlineto{\pgfqpoint{3.490331in}{1.736004in}}%
\pgfpathlineto{\pgfqpoint{3.485585in}{1.730146in}}%
\pgfpathlineto{\pgfqpoint{3.485574in}{1.722330in}}%
\pgfpathlineto{\pgfqpoint{3.488922in}{1.715480in}}%
\pgfpathlineto{\pgfqpoint{3.487219in}{1.706707in}}%
\pgfpathlineto{\pgfqpoint{3.482998in}{1.705766in}}%
\pgfpathlineto{\pgfqpoint{3.466647in}{1.697217in}}%
\pgfpathclose%
\pgfusepath{fill}%
\end{pgfscope}%
\begin{pgfscope}%
\pgfpathrectangle{\pgfqpoint{0.100000in}{0.100000in}}{\pgfqpoint{3.608454in}{2.310000in}}%
\pgfusepath{clip}%
\pgfsetbuttcap%
\pgfsetmiterjoin%
\definecolor{currentfill}{rgb}{0.000000,0.458824,0.770588}%
\pgfsetfillcolor{currentfill}%
\pgfsetlinewidth{0.000000pt}%
\definecolor{currentstroke}{rgb}{0.000000,0.000000,0.000000}%
\pgfsetstrokecolor{currentstroke}%
\pgfsetstrokeopacity{0.000000}%
\pgfsetdash{}{0pt}%
\pgfpathmoveto{\pgfqpoint{2.148148in}{1.837483in}}%
\pgfpathlineto{\pgfqpoint{2.141258in}{1.840834in}}%
\pgfpathlineto{\pgfqpoint{2.120021in}{1.840682in}}%
\pgfpathlineto{\pgfqpoint{2.120016in}{1.847577in}}%
\pgfpathlineto{\pgfqpoint{2.099439in}{1.847533in}}%
\pgfpathlineto{\pgfqpoint{2.099037in}{1.876260in}}%
\pgfpathlineto{\pgfqpoint{2.147223in}{1.876558in}}%
\pgfpathlineto{\pgfqpoint{2.142690in}{1.880249in}}%
\pgfpathlineto{\pgfqpoint{2.175363in}{1.880599in}}%
\pgfpathlineto{\pgfqpoint{2.174916in}{1.893366in}}%
\pgfpathlineto{\pgfqpoint{2.172623in}{1.893374in}}%
\pgfpathlineto{\pgfqpoint{2.172440in}{1.907196in}}%
\pgfpathlineto{\pgfqpoint{2.172320in}{1.914091in}}%
\pgfpathlineto{\pgfqpoint{2.193274in}{1.914649in}}%
\pgfpathlineto{\pgfqpoint{2.193329in}{1.907307in}}%
\pgfpathlineto{\pgfqpoint{2.193529in}{1.893561in}}%
\pgfpathlineto{\pgfqpoint{2.188794in}{1.893513in}}%
\pgfpathlineto{\pgfqpoint{2.189353in}{1.873824in}}%
\pgfpathlineto{\pgfqpoint{2.190082in}{1.834950in}}%
\pgfpathlineto{\pgfqpoint{2.186149in}{1.838564in}}%
\pgfpathlineto{\pgfqpoint{2.176492in}{1.838639in}}%
\pgfpathlineto{\pgfqpoint{2.160005in}{1.848588in}}%
\pgfpathlineto{\pgfqpoint{2.156256in}{1.841131in}}%
\pgfpathlineto{\pgfqpoint{2.148148in}{1.837483in}}%
\pgfpathclose%
\pgfusepath{fill}%
\end{pgfscope}%
\begin{pgfscope}%
\pgfpathrectangle{\pgfqpoint{0.100000in}{0.100000in}}{\pgfqpoint{3.608454in}{2.310000in}}%
\pgfusepath{clip}%
\pgfsetbuttcap%
\pgfsetmiterjoin%
\definecolor{currentfill}{rgb}{0.000000,0.678431,0.660784}%
\pgfsetfillcolor{currentfill}%
\pgfsetlinewidth{0.000000pt}%
\definecolor{currentstroke}{rgb}{0.000000,0.000000,0.000000}%
\pgfsetstrokecolor{currentstroke}%
\pgfsetstrokeopacity{0.000000}%
\pgfsetdash{}{0pt}%
\pgfpathmoveto{\pgfqpoint{1.141273in}{1.664762in}}%
\pgfpathlineto{\pgfqpoint{1.141528in}{1.670743in}}%
\pgfpathlineto{\pgfqpoint{1.137582in}{1.673009in}}%
\pgfpathlineto{\pgfqpoint{1.137407in}{1.682303in}}%
\pgfpathlineto{\pgfqpoint{1.141935in}{1.686685in}}%
\pgfpathlineto{\pgfqpoint{1.143732in}{1.692681in}}%
\pgfpathlineto{\pgfqpoint{1.141976in}{1.698141in}}%
\pgfpathlineto{\pgfqpoint{1.126068in}{1.701146in}}%
\pgfpathlineto{\pgfqpoint{1.126636in}{1.717223in}}%
\pgfpathlineto{\pgfqpoint{1.125111in}{1.725638in}}%
\pgfpathlineto{\pgfqpoint{1.117055in}{1.725647in}}%
\pgfpathlineto{\pgfqpoint{1.116145in}{1.733195in}}%
\pgfpathlineto{\pgfqpoint{1.117477in}{1.741580in}}%
\pgfpathlineto{\pgfqpoint{1.124259in}{1.750336in}}%
\pgfpathlineto{\pgfqpoint{1.182268in}{1.739464in}}%
\pgfpathlineto{\pgfqpoint{1.167956in}{1.660095in}}%
\pgfpathlineto{\pgfqpoint{1.141273in}{1.664762in}}%
\pgfpathclose%
\pgfusepath{fill}%
\end{pgfscope}%
\begin{pgfscope}%
\pgfpathrectangle{\pgfqpoint{0.100000in}{0.100000in}}{\pgfqpoint{3.608454in}{2.310000in}}%
\pgfusepath{clip}%
\pgfsetbuttcap%
\pgfsetmiterjoin%
\definecolor{currentfill}{rgb}{0.000000,0.713725,0.643137}%
\pgfsetfillcolor{currentfill}%
\pgfsetlinewidth{0.000000pt}%
\definecolor{currentstroke}{rgb}{0.000000,0.000000,0.000000}%
\pgfsetstrokecolor{currentstroke}%
\pgfsetstrokeopacity{0.000000}%
\pgfsetdash{}{0pt}%
\pgfpathmoveto{\pgfqpoint{2.193249in}{1.204475in}}%
\pgfpathlineto{\pgfqpoint{2.166408in}{1.204778in}}%
\pgfpathlineto{\pgfqpoint{2.164491in}{1.209439in}}%
\pgfpathlineto{\pgfqpoint{2.130830in}{1.210394in}}%
\pgfpathlineto{\pgfqpoint{2.130827in}{1.263573in}}%
\pgfpathlineto{\pgfqpoint{2.130815in}{1.265402in}}%
\pgfpathlineto{\pgfqpoint{2.154102in}{1.264813in}}%
\pgfpathlineto{\pgfqpoint{2.165424in}{1.263791in}}%
\pgfpathlineto{\pgfqpoint{2.165075in}{1.253073in}}%
\pgfpathlineto{\pgfqpoint{2.181129in}{1.252610in}}%
\pgfpathlineto{\pgfqpoint{2.180991in}{1.248016in}}%
\pgfpathlineto{\pgfqpoint{2.192435in}{1.247763in}}%
\pgfpathlineto{\pgfqpoint{2.193719in}{1.240848in}}%
\pgfpathlineto{\pgfqpoint{2.193249in}{1.204475in}}%
\pgfpathclose%
\pgfusepath{fill}%
\end{pgfscope}%
\begin{pgfscope}%
\pgfpathrectangle{\pgfqpoint{0.100000in}{0.100000in}}{\pgfqpoint{3.608454in}{2.310000in}}%
\pgfusepath{clip}%
\pgfsetbuttcap%
\pgfsetmiterjoin%
\definecolor{currentfill}{rgb}{0.000000,0.552941,0.723529}%
\pgfsetfillcolor{currentfill}%
\pgfsetlinewidth{0.000000pt}%
\definecolor{currentstroke}{rgb}{0.000000,0.000000,0.000000}%
\pgfsetstrokecolor{currentstroke}%
\pgfsetstrokeopacity{0.000000}%
\pgfsetdash{}{0pt}%
\pgfpathmoveto{\pgfqpoint{3.097477in}{1.148400in}}%
\pgfpathlineto{\pgfqpoint{3.057775in}{1.140860in}}%
\pgfpathlineto{\pgfqpoint{3.055178in}{1.146538in}}%
\pgfpathlineto{\pgfqpoint{3.046178in}{1.152956in}}%
\pgfpathlineto{\pgfqpoint{3.037191in}{1.156638in}}%
\pgfpathlineto{\pgfqpoint{3.035098in}{1.163147in}}%
\pgfpathlineto{\pgfqpoint{3.039762in}{1.164964in}}%
\pgfpathlineto{\pgfqpoint{3.039491in}{1.174669in}}%
\pgfpathlineto{\pgfqpoint{3.049089in}{1.181346in}}%
\pgfpathlineto{\pgfqpoint{3.060084in}{1.181854in}}%
\pgfpathlineto{\pgfqpoint{3.054182in}{1.223554in}}%
\pgfpathlineto{\pgfqpoint{3.073773in}{1.226802in}}%
\pgfpathlineto{\pgfqpoint{3.074241in}{1.252094in}}%
\pgfpathlineto{\pgfqpoint{3.074664in}{1.267099in}}%
\pgfpathlineto{\pgfqpoint{3.081510in}{1.270060in}}%
\pgfpathlineto{\pgfqpoint{3.089238in}{1.276542in}}%
\pgfpathlineto{\pgfqpoint{3.095309in}{1.277259in}}%
\pgfpathlineto{\pgfqpoint{3.099433in}{1.272390in}}%
\pgfpathlineto{\pgfqpoint{3.105607in}{1.274029in}}%
\pgfpathlineto{\pgfqpoint{3.104858in}{1.232143in}}%
\pgfpathlineto{\pgfqpoint{3.109908in}{1.233039in}}%
\pgfpathlineto{\pgfqpoint{3.113032in}{1.209517in}}%
\pgfpathlineto{\pgfqpoint{3.106519in}{1.208514in}}%
\pgfpathlineto{\pgfqpoint{3.110627in}{1.182204in}}%
\pgfpathlineto{\pgfqpoint{3.113209in}{1.177564in}}%
\pgfpathlineto{\pgfqpoint{3.093928in}{1.174159in}}%
\pgfpathlineto{\pgfqpoint{3.097477in}{1.148400in}}%
\pgfpathclose%
\pgfusepath{fill}%
\end{pgfscope}%
\begin{pgfscope}%
\pgfpathrectangle{\pgfqpoint{0.100000in}{0.100000in}}{\pgfqpoint{3.608454in}{2.310000in}}%
\pgfusepath{clip}%
\pgfsetbuttcap%
\pgfsetmiterjoin%
\definecolor{currentfill}{rgb}{0.000000,0.490196,0.754902}%
\pgfsetfillcolor{currentfill}%
\pgfsetlinewidth{0.000000pt}%
\definecolor{currentstroke}{rgb}{0.000000,0.000000,0.000000}%
\pgfsetstrokecolor{currentstroke}%
\pgfsetstrokeopacity{0.000000}%
\pgfsetdash{}{0pt}%
\pgfpathmoveto{\pgfqpoint{1.380303in}{0.447788in}}%
\pgfpathlineto{\pgfqpoint{1.383896in}{0.457637in}}%
\pgfpathlineto{\pgfqpoint{1.384080in}{0.463783in}}%
\pgfpathlineto{\pgfqpoint{1.393078in}{0.457489in}}%
\pgfpathlineto{\pgfqpoint{1.394896in}{0.449171in}}%
\pgfpathlineto{\pgfqpoint{1.389591in}{0.443421in}}%
\pgfpathlineto{\pgfqpoint{1.380303in}{0.447788in}}%
\pgfpathclose%
\pgfusepath{fill}%
\end{pgfscope}%
\begin{pgfscope}%
\pgfpathrectangle{\pgfqpoint{0.100000in}{0.100000in}}{\pgfqpoint{3.608454in}{2.310000in}}%
\pgfusepath{clip}%
\pgfsetbuttcap%
\pgfsetmiterjoin%
\definecolor{currentfill}{rgb}{0.000000,0.490196,0.754902}%
\pgfsetfillcolor{currentfill}%
\pgfsetlinewidth{0.000000pt}%
\definecolor{currentstroke}{rgb}{0.000000,0.000000,0.000000}%
\pgfsetstrokecolor{currentstroke}%
\pgfsetstrokeopacity{0.000000}%
\pgfsetdash{}{0pt}%
\pgfpathmoveto{\pgfqpoint{1.385807in}{0.418653in}}%
\pgfpathlineto{\pgfqpoint{1.386519in}{0.422119in}}%
\pgfpathlineto{\pgfqpoint{1.397550in}{0.420104in}}%
\pgfpathlineto{\pgfqpoint{1.394836in}{0.415075in}}%
\pgfpathlineto{\pgfqpoint{1.385807in}{0.418653in}}%
\pgfpathclose%
\pgfusepath{fill}%
\end{pgfscope}%
\begin{pgfscope}%
\pgfpathrectangle{\pgfqpoint{0.100000in}{0.100000in}}{\pgfqpoint{3.608454in}{2.310000in}}%
\pgfusepath{clip}%
\pgfsetbuttcap%
\pgfsetmiterjoin%
\definecolor{currentfill}{rgb}{0.000000,0.490196,0.754902}%
\pgfsetfillcolor{currentfill}%
\pgfsetlinewidth{0.000000pt}%
\definecolor{currentstroke}{rgb}{0.000000,0.000000,0.000000}%
\pgfsetstrokecolor{currentstroke}%
\pgfsetstrokeopacity{0.000000}%
\pgfsetdash{}{0pt}%
\pgfpathmoveto{\pgfqpoint{1.407961in}{0.408746in}}%
\pgfpathlineto{\pgfqpoint{1.405253in}{0.414109in}}%
\pgfpathlineto{\pgfqpoint{1.413113in}{0.426464in}}%
\pgfpathlineto{\pgfqpoint{1.404816in}{0.436091in}}%
\pgfpathlineto{\pgfqpoint{1.404807in}{0.443981in}}%
\pgfpathlineto{\pgfqpoint{1.412629in}{0.450646in}}%
\pgfpathlineto{\pgfqpoint{1.418487in}{0.448805in}}%
\pgfpathlineto{\pgfqpoint{1.419893in}{0.442848in}}%
\pgfpathlineto{\pgfqpoint{1.418178in}{0.434960in}}%
\pgfpathlineto{\pgfqpoint{1.432543in}{0.428456in}}%
\pgfpathlineto{\pgfqpoint{1.436516in}{0.412431in}}%
\pgfpathlineto{\pgfqpoint{1.441689in}{0.404958in}}%
\pgfpathlineto{\pgfqpoint{1.438189in}{0.398895in}}%
\pgfpathlineto{\pgfqpoint{1.431764in}{0.398037in}}%
\pgfpathlineto{\pgfqpoint{1.421058in}{0.403867in}}%
\pgfpathlineto{\pgfqpoint{1.407961in}{0.408746in}}%
\pgfpathclose%
\pgfusepath{fill}%
\end{pgfscope}%
\begin{pgfscope}%
\pgfpathrectangle{\pgfqpoint{0.100000in}{0.100000in}}{\pgfqpoint{3.608454in}{2.310000in}}%
\pgfusepath{clip}%
\pgfsetbuttcap%
\pgfsetmiterjoin%
\definecolor{currentfill}{rgb}{0.000000,0.490196,0.754902}%
\pgfsetfillcolor{currentfill}%
\pgfsetlinewidth{0.000000pt}%
\definecolor{currentstroke}{rgb}{0.000000,0.000000,0.000000}%
\pgfsetstrokecolor{currentstroke}%
\pgfsetstrokeopacity{0.000000}%
\pgfsetdash{}{0pt}%
\pgfpathmoveto{\pgfqpoint{1.406047in}{0.473337in}}%
\pgfpathlineto{\pgfqpoint{1.416814in}{0.465266in}}%
\pgfpathlineto{\pgfqpoint{1.412882in}{0.461231in}}%
\pgfpathlineto{\pgfqpoint{1.404816in}{0.462303in}}%
\pgfpathlineto{\pgfqpoint{1.398464in}{0.468448in}}%
\pgfpathlineto{\pgfqpoint{1.393702in}{0.476372in}}%
\pgfpathlineto{\pgfqpoint{1.378874in}{0.487793in}}%
\pgfpathlineto{\pgfqpoint{1.389247in}{0.492428in}}%
\pgfpathlineto{\pgfqpoint{1.406047in}{0.473337in}}%
\pgfpathclose%
\pgfusepath{fill}%
\end{pgfscope}%
\begin{pgfscope}%
\pgfpathrectangle{\pgfqpoint{0.100000in}{0.100000in}}{\pgfqpoint{3.608454in}{2.310000in}}%
\pgfusepath{clip}%
\pgfsetbuttcap%
\pgfsetmiterjoin%
\definecolor{currentfill}{rgb}{0.000000,0.694118,0.652941}%
\pgfsetfillcolor{currentfill}%
\pgfsetlinewidth{0.000000pt}%
\definecolor{currentstroke}{rgb}{0.000000,0.000000,0.000000}%
\pgfsetstrokecolor{currentstroke}%
\pgfsetstrokeopacity{0.000000}%
\pgfsetdash{}{0pt}%
\pgfpathmoveto{\pgfqpoint{1.576920in}{0.953485in}}%
\pgfpathlineto{\pgfqpoint{1.574789in}{0.933425in}}%
\pgfpathlineto{\pgfqpoint{1.544636in}{0.936153in}}%
\pgfpathlineto{\pgfqpoint{1.545317in}{0.943021in}}%
\pgfpathlineto{\pgfqpoint{1.531592in}{0.944243in}}%
\pgfpathlineto{\pgfqpoint{1.532929in}{0.957126in}}%
\pgfpathlineto{\pgfqpoint{1.524717in}{0.957918in}}%
\pgfpathlineto{\pgfqpoint{1.526759in}{0.978642in}}%
\pgfpathlineto{\pgfqpoint{1.519945in}{0.979398in}}%
\pgfpathlineto{\pgfqpoint{1.523829in}{1.020634in}}%
\pgfpathlineto{\pgfqpoint{1.539569in}{1.019066in}}%
\pgfpathlineto{\pgfqpoint{1.540290in}{1.025994in}}%
\pgfpathlineto{\pgfqpoint{1.553889in}{1.024630in}}%
\pgfpathlineto{\pgfqpoint{1.554565in}{1.031506in}}%
\pgfpathlineto{\pgfqpoint{1.561377in}{1.030839in}}%
\pgfpathlineto{\pgfqpoint{1.562067in}{1.037697in}}%
\pgfpathlineto{\pgfqpoint{1.568889in}{1.037039in}}%
\pgfpathlineto{\pgfqpoint{1.569604in}{1.044093in}}%
\pgfpathlineto{\pgfqpoint{1.585141in}{1.042675in}}%
\pgfpathlineto{\pgfqpoint{1.583688in}{1.026354in}}%
\pgfpathlineto{\pgfqpoint{1.617270in}{1.023437in}}%
\pgfpathlineto{\pgfqpoint{1.614395in}{0.989155in}}%
\pgfpathlineto{\pgfqpoint{1.580563in}{0.992038in}}%
\pgfpathlineto{\pgfqpoint{1.576920in}{0.953485in}}%
\pgfpathclose%
\pgfusepath{fill}%
\end{pgfscope}%
\begin{pgfscope}%
\pgfpathrectangle{\pgfqpoint{0.100000in}{0.100000in}}{\pgfqpoint{3.608454in}{2.310000in}}%
\pgfusepath{clip}%
\pgfsetbuttcap%
\pgfsetmiterjoin%
\definecolor{currentfill}{rgb}{0.000000,0.611765,0.694118}%
\pgfsetfillcolor{currentfill}%
\pgfsetlinewidth{0.000000pt}%
\definecolor{currentstroke}{rgb}{0.000000,0.000000,0.000000}%
\pgfsetstrokecolor{currentstroke}%
\pgfsetstrokeopacity{0.000000}%
\pgfsetdash{}{0pt}%
\pgfpathmoveto{\pgfqpoint{3.088906in}{0.984668in}}%
\pgfpathlineto{\pgfqpoint{3.081713in}{0.979204in}}%
\pgfpathlineto{\pgfqpoint{3.073391in}{0.976759in}}%
\pgfpathlineto{\pgfqpoint{3.065412in}{0.981312in}}%
\pgfpathlineto{\pgfqpoint{3.059264in}{0.991086in}}%
\pgfpathlineto{\pgfqpoint{3.051594in}{0.997455in}}%
\pgfpathlineto{\pgfqpoint{3.052596in}{1.003504in}}%
\pgfpathlineto{\pgfqpoint{3.049719in}{1.006843in}}%
\pgfpathlineto{\pgfqpoint{3.049502in}{1.021849in}}%
\pgfpathlineto{\pgfqpoint{3.046504in}{1.026091in}}%
\pgfpathlineto{\pgfqpoint{3.041428in}{1.022979in}}%
\pgfpathlineto{\pgfqpoint{3.030270in}{1.030329in}}%
\pgfpathlineto{\pgfqpoint{3.034331in}{1.046602in}}%
\pgfpathlineto{\pgfqpoint{3.026275in}{1.051810in}}%
\pgfpathlineto{\pgfqpoint{3.034032in}{1.059624in}}%
\pgfpathlineto{\pgfqpoint{3.040979in}{1.056871in}}%
\pgfpathlineto{\pgfqpoint{3.047179in}{1.057968in}}%
\pgfpathlineto{\pgfqpoint{3.046213in}{1.063185in}}%
\pgfpathlineto{\pgfqpoint{3.054685in}{1.068940in}}%
\pgfpathlineto{\pgfqpoint{3.060908in}{1.061022in}}%
\pgfpathlineto{\pgfqpoint{3.065576in}{1.050744in}}%
\pgfpathlineto{\pgfqpoint{3.073425in}{1.048823in}}%
\pgfpathlineto{\pgfqpoint{3.078779in}{1.041854in}}%
\pgfpathlineto{\pgfqpoint{3.076612in}{1.036110in}}%
\pgfpathlineto{\pgfqpoint{3.096046in}{1.025469in}}%
\pgfpathlineto{\pgfqpoint{3.091334in}{1.021253in}}%
\pgfpathlineto{\pgfqpoint{3.096860in}{1.011889in}}%
\pgfpathlineto{\pgfqpoint{3.092298in}{1.007042in}}%
\pgfpathlineto{\pgfqpoint{3.094287in}{1.003646in}}%
\pgfpathlineto{\pgfqpoint{3.088906in}{0.984668in}}%
\pgfpathclose%
\pgfusepath{fill}%
\end{pgfscope}%
\begin{pgfscope}%
\pgfpathrectangle{\pgfqpoint{0.100000in}{0.100000in}}{\pgfqpoint{3.608454in}{2.310000in}}%
\pgfusepath{clip}%
\pgfsetbuttcap%
\pgfsetmiterjoin%
\definecolor{currentfill}{rgb}{0.000000,0.611765,0.694118}%
\pgfsetfillcolor{currentfill}%
\pgfsetlinewidth{0.000000pt}%
\definecolor{currentstroke}{rgb}{0.000000,0.000000,0.000000}%
\pgfsetstrokecolor{currentstroke}%
\pgfsetstrokeopacity{0.000000}%
\pgfsetdash{}{0pt}%
\pgfpathmoveto{\pgfqpoint{1.562162in}{1.934829in}}%
\pgfpathlineto{\pgfqpoint{1.559222in}{1.935142in}}%
\pgfpathlineto{\pgfqpoint{1.556192in}{1.907466in}}%
\pgfpathlineto{\pgfqpoint{1.553831in}{1.907730in}}%
\pgfpathlineto{\pgfqpoint{1.487149in}{1.915660in}}%
\pgfpathlineto{\pgfqpoint{1.490580in}{1.943193in}}%
\pgfpathlineto{\pgfqpoint{1.492673in}{1.942949in}}%
\pgfpathlineto{\pgfqpoint{1.495943in}{1.970351in}}%
\pgfpathlineto{\pgfqpoint{1.497843in}{1.970127in}}%
\pgfpathlineto{\pgfqpoint{1.501201in}{1.997658in}}%
\pgfpathlineto{\pgfqpoint{1.503208in}{1.998589in}}%
\pgfpathlineto{\pgfqpoint{1.506437in}{2.025137in}}%
\pgfpathlineto{\pgfqpoint{1.513276in}{2.024315in}}%
\pgfpathlineto{\pgfqpoint{1.512427in}{2.017326in}}%
\pgfpathlineto{\pgfqpoint{1.519298in}{2.016515in}}%
\pgfpathlineto{\pgfqpoint{1.520131in}{2.023482in}}%
\pgfpathlineto{\pgfqpoint{1.543137in}{2.020826in}}%
\pgfpathlineto{\pgfqpoint{1.547528in}{2.020313in}}%
\pgfpathlineto{\pgfqpoint{1.545674in}{2.004107in}}%
\pgfpathlineto{\pgfqpoint{1.552609in}{2.003322in}}%
\pgfpathlineto{\pgfqpoint{1.552079in}{1.998723in}}%
\pgfpathlineto{\pgfqpoint{1.558985in}{1.997946in}}%
\pgfpathlineto{\pgfqpoint{1.560751in}{1.993122in}}%
\pgfpathlineto{\pgfqpoint{1.583733in}{1.990580in}}%
\pgfpathlineto{\pgfqpoint{1.582076in}{1.976639in}}%
\pgfpathlineto{\pgfqpoint{1.574985in}{1.975097in}}%
\pgfpathlineto{\pgfqpoint{1.573489in}{1.961301in}}%
\pgfpathlineto{\pgfqpoint{1.564057in}{1.962319in}}%
\pgfpathlineto{\pgfqpoint{1.562162in}{1.934829in}}%
\pgfpathclose%
\pgfusepath{fill}%
\end{pgfscope}%
\begin{pgfscope}%
\pgfpathrectangle{\pgfqpoint{0.100000in}{0.100000in}}{\pgfqpoint{3.608454in}{2.310000in}}%
\pgfusepath{clip}%
\pgfsetbuttcap%
\pgfsetmiterjoin%
\definecolor{currentfill}{rgb}{0.000000,0.431373,0.784314}%
\pgfsetfillcolor{currentfill}%
\pgfsetlinewidth{0.000000pt}%
\definecolor{currentstroke}{rgb}{0.000000,0.000000,0.000000}%
\pgfsetstrokecolor{currentstroke}%
\pgfsetstrokeopacity{0.000000}%
\pgfsetdash{}{0pt}%
\pgfpathmoveto{\pgfqpoint{2.013559in}{0.517174in}}%
\pgfpathlineto{\pgfqpoint{2.014894in}{0.542939in}}%
\pgfpathlineto{\pgfqpoint{1.992065in}{0.566132in}}%
\pgfpathlineto{\pgfqpoint{2.012654in}{0.578657in}}%
\pgfpathlineto{\pgfqpoint{2.015636in}{0.587127in}}%
\pgfpathlineto{\pgfqpoint{2.024967in}{0.596552in}}%
\pgfpathlineto{\pgfqpoint{2.032870in}{0.593014in}}%
\pgfpathlineto{\pgfqpoint{2.032880in}{0.583631in}}%
\pgfpathlineto{\pgfqpoint{2.037581in}{0.580386in}}%
\pgfpathlineto{\pgfqpoint{2.040450in}{0.572451in}}%
\pgfpathlineto{\pgfqpoint{2.047518in}{0.566563in}}%
\pgfpathlineto{\pgfqpoint{2.054185in}{0.571463in}}%
\pgfpathlineto{\pgfqpoint{2.061080in}{0.568664in}}%
\pgfpathlineto{\pgfqpoint{2.066381in}{0.572643in}}%
\pgfpathlineto{\pgfqpoint{2.068450in}{0.580491in}}%
\pgfpathlineto{\pgfqpoint{2.074520in}{0.580523in}}%
\pgfpathlineto{\pgfqpoint{2.077251in}{0.591805in}}%
\pgfpathlineto{\pgfqpoint{2.088573in}{0.592795in}}%
\pgfpathlineto{\pgfqpoint{2.091703in}{0.589805in}}%
\pgfpathlineto{\pgfqpoint{2.090532in}{0.582574in}}%
\pgfpathlineto{\pgfqpoint{2.102863in}{0.561208in}}%
\pgfpathlineto{\pgfqpoint{2.098432in}{0.551471in}}%
\pgfpathlineto{\pgfqpoint{2.079822in}{0.534485in}}%
\pgfpathlineto{\pgfqpoint{2.074339in}{0.533228in}}%
\pgfpathlineto{\pgfqpoint{2.049544in}{0.518585in}}%
\pgfpathlineto{\pgfqpoint{2.036951in}{0.514000in}}%
\pgfpathlineto{\pgfqpoint{2.033736in}{0.517714in}}%
\pgfpathlineto{\pgfqpoint{2.020863in}{0.511930in}}%
\pgfpathlineto{\pgfqpoint{2.021021in}{0.521517in}}%
\pgfpathlineto{\pgfqpoint{2.013559in}{0.517174in}}%
\pgfpathclose%
\pgfusepath{fill}%
\end{pgfscope}%
\begin{pgfscope}%
\pgfpathrectangle{\pgfqpoint{0.100000in}{0.100000in}}{\pgfqpoint{3.608454in}{2.310000in}}%
\pgfusepath{clip}%
\pgfsetbuttcap%
\pgfsetmiterjoin%
\definecolor{currentfill}{rgb}{0.000000,0.686275,0.656863}%
\pgfsetfillcolor{currentfill}%
\pgfsetlinewidth{0.000000pt}%
\definecolor{currentstroke}{rgb}{0.000000,0.000000,0.000000}%
\pgfsetstrokecolor{currentstroke}%
\pgfsetstrokeopacity{0.000000}%
\pgfsetdash{}{0pt}%
\pgfpathmoveto{\pgfqpoint{3.111250in}{1.435310in}}%
\pgfpathlineto{\pgfqpoint{3.103999in}{1.422556in}}%
\pgfpathlineto{\pgfqpoint{3.099019in}{1.419473in}}%
\pgfpathlineto{\pgfqpoint{3.095711in}{1.408904in}}%
\pgfpathlineto{\pgfqpoint{3.086938in}{1.414331in}}%
\pgfpathlineto{\pgfqpoint{3.084328in}{1.406748in}}%
\pgfpathlineto{\pgfqpoint{3.063909in}{1.418899in}}%
\pgfpathlineto{\pgfqpoint{3.065468in}{1.428454in}}%
\pgfpathlineto{\pgfqpoint{3.062165in}{1.429526in}}%
\pgfpathlineto{\pgfqpoint{3.063986in}{1.437446in}}%
\pgfpathlineto{\pgfqpoint{3.060760in}{1.438938in}}%
\pgfpathlineto{\pgfqpoint{3.052407in}{1.435951in}}%
\pgfpathlineto{\pgfqpoint{3.045971in}{1.477046in}}%
\pgfpathlineto{\pgfqpoint{3.051009in}{1.477958in}}%
\pgfpathlineto{\pgfqpoint{3.111561in}{1.488897in}}%
\pgfpathlineto{\pgfqpoint{3.116008in}{1.487243in}}%
\pgfpathlineto{\pgfqpoint{3.110506in}{1.479451in}}%
\pgfpathlineto{\pgfqpoint{3.109226in}{1.471916in}}%
\pgfpathlineto{\pgfqpoint{3.097628in}{1.471273in}}%
\pgfpathlineto{\pgfqpoint{3.091266in}{1.461236in}}%
\pgfpathlineto{\pgfqpoint{3.088891in}{1.453283in}}%
\pgfpathlineto{\pgfqpoint{3.082421in}{1.444639in}}%
\pgfpathlineto{\pgfqpoint{3.087837in}{1.442494in}}%
\pgfpathlineto{\pgfqpoint{3.105576in}{1.439244in}}%
\pgfpathlineto{\pgfqpoint{3.111250in}{1.435310in}}%
\pgfpathclose%
\pgfusepath{fill}%
\end{pgfscope}%
\begin{pgfscope}%
\pgfpathrectangle{\pgfqpoint{0.100000in}{0.100000in}}{\pgfqpoint{3.608454in}{2.310000in}}%
\pgfusepath{clip}%
\pgfsetbuttcap%
\pgfsetmiterjoin%
\definecolor{currentfill}{rgb}{0.000000,0.611765,0.694118}%
\pgfsetfillcolor{currentfill}%
\pgfsetlinewidth{0.000000pt}%
\definecolor{currentstroke}{rgb}{0.000000,0.000000,0.000000}%
\pgfsetstrokecolor{currentstroke}%
\pgfsetstrokeopacity{0.000000}%
\pgfsetdash{}{0pt}%
\pgfpathmoveto{\pgfqpoint{3.279173in}{1.774369in}}%
\pgfpathlineto{\pgfqpoint{3.272069in}{1.775491in}}%
\pgfpathlineto{\pgfqpoint{3.271017in}{1.786225in}}%
\pgfpathlineto{\pgfqpoint{3.267628in}{1.792368in}}%
\pgfpathlineto{\pgfqpoint{3.270444in}{1.800501in}}%
\pgfpathlineto{\pgfqpoint{3.267521in}{1.808820in}}%
\pgfpathlineto{\pgfqpoint{3.257858in}{1.810998in}}%
\pgfpathlineto{\pgfqpoint{3.260468in}{1.823521in}}%
\pgfpathlineto{\pgfqpoint{3.245888in}{1.867665in}}%
\pgfpathlineto{\pgfqpoint{3.263105in}{1.873951in}}%
\pgfpathlineto{\pgfqpoint{3.276788in}{1.878840in}}%
\pgfpathlineto{\pgfqpoint{3.280930in}{1.867522in}}%
\pgfpathlineto{\pgfqpoint{3.277231in}{1.863034in}}%
\pgfpathlineto{\pgfqpoint{3.286142in}{1.855787in}}%
\pgfpathlineto{\pgfqpoint{3.295690in}{1.856839in}}%
\pgfpathlineto{\pgfqpoint{3.296075in}{1.852707in}}%
\pgfpathlineto{\pgfqpoint{3.287566in}{1.849431in}}%
\pgfpathlineto{\pgfqpoint{3.309890in}{1.786988in}}%
\pgfpathlineto{\pgfqpoint{3.300308in}{1.775508in}}%
\pgfpathlineto{\pgfqpoint{3.291128in}{1.772508in}}%
\pgfpathlineto{\pgfqpoint{3.279173in}{1.774369in}}%
\pgfpathclose%
\pgfusepath{fill}%
\end{pgfscope}%
\begin{pgfscope}%
\pgfpathrectangle{\pgfqpoint{0.100000in}{0.100000in}}{\pgfqpoint{3.608454in}{2.310000in}}%
\pgfusepath{clip}%
\pgfsetbuttcap%
\pgfsetmiterjoin%
\definecolor{currentfill}{rgb}{0.000000,0.517647,0.741176}%
\pgfsetfillcolor{currentfill}%
\pgfsetlinewidth{0.000000pt}%
\definecolor{currentstroke}{rgb}{0.000000,0.000000,0.000000}%
\pgfsetstrokecolor{currentstroke}%
\pgfsetstrokeopacity{0.000000}%
\pgfsetdash{}{0pt}%
\pgfpathmoveto{\pgfqpoint{2.558058in}{1.370400in}}%
\pgfpathlineto{\pgfqpoint{2.550764in}{1.371543in}}%
\pgfpathlineto{\pgfqpoint{2.506514in}{1.367545in}}%
\pgfpathlineto{\pgfqpoint{2.506265in}{1.370984in}}%
\pgfpathlineto{\pgfqpoint{2.503724in}{1.405402in}}%
\pgfpathlineto{\pgfqpoint{2.528514in}{1.407327in}}%
\pgfpathlineto{\pgfqpoint{2.534185in}{1.410410in}}%
\pgfpathlineto{\pgfqpoint{2.534753in}{1.425828in}}%
\pgfpathlineto{\pgfqpoint{2.559183in}{1.428045in}}%
\pgfpathlineto{\pgfqpoint{2.562752in}{1.385889in}}%
\pgfpathlineto{\pgfqpoint{2.558464in}{1.382618in}}%
\pgfpathlineto{\pgfqpoint{2.560407in}{1.373999in}}%
\pgfpathlineto{\pgfqpoint{2.558058in}{1.370400in}}%
\pgfpathclose%
\pgfusepath{fill}%
\end{pgfscope}%
\begin{pgfscope}%
\pgfpathrectangle{\pgfqpoint{0.100000in}{0.100000in}}{\pgfqpoint{3.608454in}{2.310000in}}%
\pgfusepath{clip}%
\pgfsetbuttcap%
\pgfsetmiterjoin%
\definecolor{currentfill}{rgb}{0.000000,0.419608,0.790196}%
\pgfsetfillcolor{currentfill}%
\pgfsetlinewidth{0.000000pt}%
\definecolor{currentstroke}{rgb}{0.000000,0.000000,0.000000}%
\pgfsetstrokecolor{currentstroke}%
\pgfsetstrokeopacity{0.000000}%
\pgfsetdash{}{0pt}%
\pgfpathmoveto{\pgfqpoint{1.843275in}{1.841097in}}%
\pgfpathlineto{\pgfqpoint{1.812489in}{1.842876in}}%
\pgfpathlineto{\pgfqpoint{1.809666in}{1.848703in}}%
\pgfpathlineto{\pgfqpoint{1.812344in}{1.853586in}}%
\pgfpathlineto{\pgfqpoint{1.808999in}{1.861084in}}%
\pgfpathlineto{\pgfqpoint{1.801983in}{1.865242in}}%
\pgfpathlineto{\pgfqpoint{1.807298in}{1.876074in}}%
\pgfpathlineto{\pgfqpoint{1.811895in}{1.879653in}}%
\pgfpathlineto{\pgfqpoint{1.808345in}{1.892715in}}%
\pgfpathlineto{\pgfqpoint{1.802493in}{1.899008in}}%
\pgfpathlineto{\pgfqpoint{1.845659in}{1.896366in}}%
\pgfpathlineto{\pgfqpoint{1.843275in}{1.841097in}}%
\pgfpathclose%
\pgfusepath{fill}%
\end{pgfscope}%
\begin{pgfscope}%
\pgfpathrectangle{\pgfqpoint{0.100000in}{0.100000in}}{\pgfqpoint{3.608454in}{2.310000in}}%
\pgfusepath{clip}%
\pgfsetbuttcap%
\pgfsetmiterjoin%
\definecolor{currentfill}{rgb}{0.000000,0.827451,0.586275}%
\pgfsetfillcolor{currentfill}%
\pgfsetlinewidth{0.000000pt}%
\definecolor{currentstroke}{rgb}{0.000000,0.000000,0.000000}%
\pgfsetstrokecolor{currentstroke}%
\pgfsetstrokeopacity{0.000000}%
\pgfsetdash{}{0pt}%
\pgfpathmoveto{\pgfqpoint{0.787216in}{0.442053in}}%
\pgfpathlineto{\pgfqpoint{0.796011in}{0.435184in}}%
\pgfpathlineto{\pgfqpoint{0.790893in}{0.428785in}}%
\pgfpathlineto{\pgfqpoint{0.779568in}{0.437653in}}%
\pgfpathlineto{\pgfqpoint{0.783683in}{0.438305in}}%
\pgfpathlineto{\pgfqpoint{0.787048in}{0.440340in}}%
\pgfpathlineto{\pgfqpoint{0.787216in}{0.442053in}}%
\pgfpathclose%
\pgfusepath{fill}%
\end{pgfscope}%
\begin{pgfscope}%
\pgfpathrectangle{\pgfqpoint{0.100000in}{0.100000in}}{\pgfqpoint{3.608454in}{2.310000in}}%
\pgfusepath{clip}%
\pgfsetbuttcap%
\pgfsetmiterjoin%
\definecolor{currentfill}{rgb}{0.000000,0.796078,0.601961}%
\pgfsetfillcolor{currentfill}%
\pgfsetlinewidth{0.000000pt}%
\definecolor{currentstroke}{rgb}{0.000000,0.000000,0.000000}%
\pgfsetstrokecolor{currentstroke}%
\pgfsetstrokeopacity{0.000000}%
\pgfsetdash{}{0pt}%
\pgfpathmoveto{\pgfqpoint{1.760148in}{1.734344in}}%
\pgfpathlineto{\pgfqpoint{1.750543in}{1.731001in}}%
\pgfpathlineto{\pgfqpoint{1.748309in}{1.699187in}}%
\pgfpathlineto{\pgfqpoint{1.697612in}{1.702965in}}%
\pgfpathlineto{\pgfqpoint{1.696530in}{1.709950in}}%
\pgfpathlineto{\pgfqpoint{1.697918in}{1.727667in}}%
\pgfpathlineto{\pgfqpoint{1.704828in}{1.727690in}}%
\pgfpathlineto{\pgfqpoint{1.707249in}{1.750308in}}%
\pgfpathlineto{\pgfqpoint{1.710799in}{1.791190in}}%
\pgfpathlineto{\pgfqpoint{1.723276in}{1.792199in}}%
\pgfpathlineto{\pgfqpoint{1.729800in}{1.795494in}}%
\pgfpathlineto{\pgfqpoint{1.737625in}{1.793677in}}%
\pgfpathlineto{\pgfqpoint{1.747956in}{1.802589in}}%
\pgfpathlineto{\pgfqpoint{1.754993in}{1.802225in}}%
\pgfpathlineto{\pgfqpoint{1.760814in}{1.806325in}}%
\pgfpathlineto{\pgfqpoint{1.759531in}{1.801865in}}%
\pgfpathlineto{\pgfqpoint{1.757629in}{1.774393in}}%
\pgfpathlineto{\pgfqpoint{1.756032in}{1.760574in}}%
\pgfpathlineto{\pgfqpoint{1.762841in}{1.760148in}}%
\pgfpathlineto{\pgfqpoint{1.760148in}{1.734344in}}%
\pgfpathclose%
\pgfusepath{fill}%
\end{pgfscope}%
\begin{pgfscope}%
\pgfpathrectangle{\pgfqpoint{0.100000in}{0.100000in}}{\pgfqpoint{3.608454in}{2.310000in}}%
\pgfusepath{clip}%
\pgfsetbuttcap%
\pgfsetmiterjoin%
\definecolor{currentfill}{rgb}{0.000000,0.541176,0.729412}%
\pgfsetfillcolor{currentfill}%
\pgfsetlinewidth{0.000000pt}%
\definecolor{currentstroke}{rgb}{0.000000,0.000000,0.000000}%
\pgfsetstrokecolor{currentstroke}%
\pgfsetstrokeopacity{0.000000}%
\pgfsetdash{}{0pt}%
\pgfpathmoveto{\pgfqpoint{2.406029in}{1.504012in}}%
\pgfpathlineto{\pgfqpoint{2.405330in}{1.517808in}}%
\pgfpathlineto{\pgfqpoint{2.378376in}{1.516534in}}%
\pgfpathlineto{\pgfqpoint{2.378619in}{1.509638in}}%
\pgfpathlineto{\pgfqpoint{2.358001in}{1.509016in}}%
\pgfpathlineto{\pgfqpoint{2.348213in}{1.508885in}}%
\pgfpathlineto{\pgfqpoint{2.345525in}{1.515337in}}%
\pgfpathlineto{\pgfqpoint{2.337899in}{1.521676in}}%
\pgfpathlineto{\pgfqpoint{2.341290in}{1.535608in}}%
\pgfpathlineto{\pgfqpoint{2.356534in}{1.539366in}}%
\pgfpathlineto{\pgfqpoint{2.356154in}{1.550723in}}%
\pgfpathlineto{\pgfqpoint{2.349290in}{1.550458in}}%
\pgfpathlineto{\pgfqpoint{2.348758in}{1.564236in}}%
\pgfpathlineto{\pgfqpoint{2.376029in}{1.564684in}}%
\pgfpathlineto{\pgfqpoint{2.380372in}{1.562314in}}%
\pgfpathlineto{\pgfqpoint{2.387375in}{1.567036in}}%
\pgfpathlineto{\pgfqpoint{2.391551in}{1.551544in}}%
\pgfpathlineto{\pgfqpoint{2.410650in}{1.552521in}}%
\pgfpathlineto{\pgfqpoint{2.412483in}{1.524944in}}%
\pgfpathlineto{\pgfqpoint{2.425489in}{1.525655in}}%
\pgfpathlineto{\pgfqpoint{2.426710in}{1.505121in}}%
\pgfpathlineto{\pgfqpoint{2.406029in}{1.504012in}}%
\pgfpathclose%
\pgfusepath{fill}%
\end{pgfscope}%
\begin{pgfscope}%
\pgfpathrectangle{\pgfqpoint{0.100000in}{0.100000in}}{\pgfqpoint{3.608454in}{2.310000in}}%
\pgfusepath{clip}%
\pgfsetbuttcap%
\pgfsetmiterjoin%
\definecolor{currentfill}{rgb}{0.000000,0.494118,0.752941}%
\pgfsetfillcolor{currentfill}%
\pgfsetlinewidth{0.000000pt}%
\definecolor{currentstroke}{rgb}{0.000000,0.000000,0.000000}%
\pgfsetstrokecolor{currentstroke}%
\pgfsetstrokeopacity{0.000000}%
\pgfsetdash{}{0pt}%
\pgfpathmoveto{\pgfqpoint{2.708723in}{0.733977in}}%
\pgfpathlineto{\pgfqpoint{2.719091in}{0.734938in}}%
\pgfpathlineto{\pgfqpoint{2.721186in}{0.711980in}}%
\pgfpathlineto{\pgfqpoint{2.733506in}{0.712730in}}%
\pgfpathlineto{\pgfqpoint{2.730084in}{0.707397in}}%
\pgfpathlineto{\pgfqpoint{2.735056in}{0.696150in}}%
\pgfpathlineto{\pgfqpoint{2.726751in}{0.687235in}}%
\pgfpathlineto{\pgfqpoint{2.727447in}{0.677658in}}%
\pgfpathlineto{\pgfqpoint{2.710794in}{0.682159in}}%
\pgfpathlineto{\pgfqpoint{2.699063in}{0.683673in}}%
\pgfpathlineto{\pgfqpoint{2.682757in}{0.683563in}}%
\pgfpathlineto{\pgfqpoint{2.666620in}{0.680661in}}%
\pgfpathlineto{\pgfqpoint{2.645613in}{0.674473in}}%
\pgfpathlineto{\pgfqpoint{2.635436in}{0.672947in}}%
\pgfpathlineto{\pgfqpoint{2.622714in}{0.669017in}}%
\pgfpathlineto{\pgfqpoint{2.627322in}{0.678308in}}%
\pgfpathlineto{\pgfqpoint{2.632066in}{0.682154in}}%
\pgfpathlineto{\pgfqpoint{2.626017in}{0.689071in}}%
\pgfpathlineto{\pgfqpoint{2.628964in}{0.696389in}}%
\pgfpathlineto{\pgfqpoint{2.627748in}{0.701087in}}%
\pgfpathlineto{\pgfqpoint{2.618895in}{0.705619in}}%
\pgfpathlineto{\pgfqpoint{2.610908in}{0.714979in}}%
\pgfpathlineto{\pgfqpoint{2.612513in}{0.725646in}}%
\pgfpathlineto{\pgfqpoint{2.674593in}{0.730790in}}%
\pgfpathlineto{\pgfqpoint{2.708723in}{0.733977in}}%
\pgfpathclose%
\pgfusepath{fill}%
\end{pgfscope}%
\begin{pgfscope}%
\pgfpathrectangle{\pgfqpoint{0.100000in}{0.100000in}}{\pgfqpoint{3.608454in}{2.310000in}}%
\pgfusepath{clip}%
\pgfsetbuttcap%
\pgfsetmiterjoin%
\definecolor{currentfill}{rgb}{0.000000,0.364706,0.817647}%
\pgfsetfillcolor{currentfill}%
\pgfsetlinewidth{0.000000pt}%
\definecolor{currentstroke}{rgb}{0.000000,0.000000,0.000000}%
\pgfsetstrokecolor{currentstroke}%
\pgfsetstrokeopacity{0.000000}%
\pgfsetdash{}{0pt}%
\pgfpathmoveto{\pgfqpoint{1.826801in}{1.425965in}}%
\pgfpathlineto{\pgfqpoint{1.824934in}{1.391676in}}%
\pgfpathlineto{\pgfqpoint{1.792462in}{1.393414in}}%
\pgfpathlineto{\pgfqpoint{1.791324in}{1.393448in}}%
\pgfpathlineto{\pgfqpoint{1.793434in}{1.427744in}}%
\pgfpathlineto{\pgfqpoint{1.826801in}{1.425965in}}%
\pgfpathclose%
\pgfusepath{fill}%
\end{pgfscope}%
\begin{pgfscope}%
\pgfpathrectangle{\pgfqpoint{0.100000in}{0.100000in}}{\pgfqpoint{3.608454in}{2.310000in}}%
\pgfusepath{clip}%
\pgfsetbuttcap%
\pgfsetmiterjoin%
\definecolor{currentfill}{rgb}{0.000000,0.815686,0.592157}%
\pgfsetfillcolor{currentfill}%
\pgfsetlinewidth{0.000000pt}%
\definecolor{currentstroke}{rgb}{0.000000,0.000000,0.000000}%
\pgfsetstrokecolor{currentstroke}%
\pgfsetstrokeopacity{0.000000}%
\pgfsetdash{}{0pt}%
\pgfpathmoveto{\pgfqpoint{2.350507in}{1.098801in}}%
\pgfpathlineto{\pgfqpoint{2.343663in}{1.088662in}}%
\pgfpathlineto{\pgfqpoint{2.340391in}{1.078576in}}%
\pgfpathlineto{\pgfqpoint{2.340700in}{1.069582in}}%
\pgfpathlineto{\pgfqpoint{2.313347in}{1.069167in}}%
\pgfpathlineto{\pgfqpoint{2.313117in}{1.083232in}}%
\pgfpathlineto{\pgfqpoint{2.307635in}{1.090812in}}%
\pgfpathlineto{\pgfqpoint{2.308868in}{1.093507in}}%
\pgfpathlineto{\pgfqpoint{2.302022in}{1.100303in}}%
\pgfpathlineto{\pgfqpoint{2.293303in}{1.102454in}}%
\pgfpathlineto{\pgfqpoint{2.295109in}{1.108191in}}%
\pgfpathlineto{\pgfqpoint{2.286012in}{1.115699in}}%
\pgfpathlineto{\pgfqpoint{2.288166in}{1.125922in}}%
\pgfpathlineto{\pgfqpoint{2.325882in}{1.126330in}}%
\pgfpathlineto{\pgfqpoint{2.325633in}{1.133246in}}%
\pgfpathlineto{\pgfqpoint{2.332736in}{1.133366in}}%
\pgfpathlineto{\pgfqpoint{2.332518in}{1.146200in}}%
\pgfpathlineto{\pgfqpoint{2.335244in}{1.146265in}}%
\pgfpathlineto{\pgfqpoint{2.343960in}{1.136912in}}%
\pgfpathlineto{\pgfqpoint{2.345455in}{1.126590in}}%
\pgfpathlineto{\pgfqpoint{2.339731in}{1.125367in}}%
\pgfpathlineto{\pgfqpoint{2.340339in}{1.098431in}}%
\pgfpathlineto{\pgfqpoint{2.350507in}{1.098801in}}%
\pgfpathclose%
\pgfusepath{fill}%
\end{pgfscope}%
\begin{pgfscope}%
\pgfpathrectangle{\pgfqpoint{0.100000in}{0.100000in}}{\pgfqpoint{3.608454in}{2.310000in}}%
\pgfusepath{clip}%
\pgfsetbuttcap%
\pgfsetmiterjoin%
\definecolor{currentfill}{rgb}{0.000000,0.411765,0.794118}%
\pgfsetfillcolor{currentfill}%
\pgfsetlinewidth{0.000000pt}%
\definecolor{currentstroke}{rgb}{0.000000,0.000000,0.000000}%
\pgfsetstrokecolor{currentstroke}%
\pgfsetstrokeopacity{0.000000}%
\pgfsetdash{}{0pt}%
\pgfpathmoveto{\pgfqpoint{2.313973in}{1.320063in}}%
\pgfpathlineto{\pgfqpoint{2.314915in}{1.287249in}}%
\pgfpathlineto{\pgfqpoint{2.280744in}{1.286387in}}%
\pgfpathlineto{\pgfqpoint{2.274008in}{1.289693in}}%
\pgfpathlineto{\pgfqpoint{2.267441in}{1.289875in}}%
\pgfpathlineto{\pgfqpoint{2.267557in}{1.296520in}}%
\pgfpathlineto{\pgfqpoint{2.253822in}{1.296576in}}%
\pgfpathlineto{\pgfqpoint{2.240098in}{1.309360in}}%
\pgfpathlineto{\pgfqpoint{2.240062in}{1.316251in}}%
\pgfpathlineto{\pgfqpoint{2.226458in}{1.316665in}}%
\pgfpathlineto{\pgfqpoint{2.226654in}{1.335261in}}%
\pgfpathlineto{\pgfqpoint{2.226701in}{1.338754in}}%
\pgfpathlineto{\pgfqpoint{2.233532in}{1.346189in}}%
\pgfpathlineto{\pgfqpoint{2.232957in}{1.349258in}}%
\pgfpathlineto{\pgfqpoint{2.238456in}{1.359055in}}%
\pgfpathlineto{\pgfqpoint{2.241843in}{1.368445in}}%
\pgfpathlineto{\pgfqpoint{2.246959in}{1.366834in}}%
\pgfpathlineto{\pgfqpoint{2.264061in}{1.361534in}}%
\pgfpathlineto{\pgfqpoint{2.271150in}{1.361570in}}%
\pgfpathlineto{\pgfqpoint{2.284008in}{1.361432in}}%
\pgfpathlineto{\pgfqpoint{2.284099in}{1.347536in}}%
\pgfpathlineto{\pgfqpoint{2.313325in}{1.348137in}}%
\pgfpathlineto{\pgfqpoint{2.313973in}{1.320063in}}%
\pgfpathclose%
\pgfusepath{fill}%
\end{pgfscope}%
\begin{pgfscope}%
\pgfpathrectangle{\pgfqpoint{0.100000in}{0.100000in}}{\pgfqpoint{3.608454in}{2.310000in}}%
\pgfusepath{clip}%
\pgfsetbuttcap%
\pgfsetmiterjoin%
\definecolor{currentfill}{rgb}{0.000000,0.521569,0.739216}%
\pgfsetfillcolor{currentfill}%
\pgfsetlinewidth{0.000000pt}%
\definecolor{currentstroke}{rgb}{0.000000,0.000000,0.000000}%
\pgfsetstrokecolor{currentstroke}%
\pgfsetstrokeopacity{0.000000}%
\pgfsetdash{}{0pt}%
\pgfpathmoveto{\pgfqpoint{2.986023in}{0.797823in}}%
\pgfpathlineto{\pgfqpoint{2.983849in}{0.804589in}}%
\pgfpathlineto{\pgfqpoint{2.978722in}{0.804662in}}%
\pgfpathlineto{\pgfqpoint{2.975294in}{0.828495in}}%
\pgfpathlineto{\pgfqpoint{2.979065in}{0.834250in}}%
\pgfpathlineto{\pgfqpoint{2.975754in}{0.837437in}}%
\pgfpathlineto{\pgfqpoint{2.967579in}{0.838638in}}%
\pgfpathlineto{\pgfqpoint{2.967645in}{0.858971in}}%
\pgfpathlineto{\pgfqpoint{2.962520in}{0.870538in}}%
\pgfpathlineto{\pgfqpoint{2.991213in}{0.868134in}}%
\pgfpathlineto{\pgfqpoint{2.994468in}{0.861760in}}%
\pgfpathlineto{\pgfqpoint{2.996708in}{0.858258in}}%
\pgfpathlineto{\pgfqpoint{3.005118in}{0.855085in}}%
\pgfpathlineto{\pgfqpoint{3.011213in}{0.848731in}}%
\pgfpathlineto{\pgfqpoint{3.020884in}{0.848437in}}%
\pgfpathlineto{\pgfqpoint{3.025040in}{0.841844in}}%
\pgfpathlineto{\pgfqpoint{3.033775in}{0.835716in}}%
\pgfpathlineto{\pgfqpoint{3.039659in}{0.833647in}}%
\pgfpathlineto{\pgfqpoint{3.044197in}{0.827018in}}%
\pgfpathlineto{\pgfqpoint{3.040729in}{0.821509in}}%
\pgfpathlineto{\pgfqpoint{3.035156in}{0.827393in}}%
\pgfpathlineto{\pgfqpoint{3.023851in}{0.824906in}}%
\pgfpathlineto{\pgfqpoint{3.019148in}{0.828947in}}%
\pgfpathlineto{\pgfqpoint{3.015577in}{0.818559in}}%
\pgfpathlineto{\pgfqpoint{3.009224in}{0.814410in}}%
\pgfpathlineto{\pgfqpoint{3.013726in}{0.809687in}}%
\pgfpathlineto{\pgfqpoint{3.007825in}{0.797264in}}%
\pgfpathlineto{\pgfqpoint{3.002283in}{0.799408in}}%
\pgfpathlineto{\pgfqpoint{2.994265in}{0.796799in}}%
\pgfpathlineto{\pgfqpoint{2.986023in}{0.797823in}}%
\pgfpathclose%
\pgfusepath{fill}%
\end{pgfscope}%
\begin{pgfscope}%
\pgfpathrectangle{\pgfqpoint{0.100000in}{0.100000in}}{\pgfqpoint{3.608454in}{2.310000in}}%
\pgfusepath{clip}%
\pgfsetbuttcap%
\pgfsetmiterjoin%
\definecolor{currentfill}{rgb}{0.000000,0.419608,0.790196}%
\pgfsetfillcolor{currentfill}%
\pgfsetlinewidth{0.000000pt}%
\definecolor{currentstroke}{rgb}{0.000000,0.000000,0.000000}%
\pgfsetstrokecolor{currentstroke}%
\pgfsetstrokeopacity{0.000000}%
\pgfsetdash{}{0pt}%
\pgfpathmoveto{\pgfqpoint{2.025168in}{1.813601in}}%
\pgfpathlineto{\pgfqpoint{2.024747in}{1.786093in}}%
\pgfpathlineto{\pgfqpoint{2.024621in}{1.779133in}}%
\pgfpathlineto{\pgfqpoint{1.986361in}{1.779846in}}%
\pgfpathlineto{\pgfqpoint{1.945343in}{1.781068in}}%
\pgfpathlineto{\pgfqpoint{1.945390in}{1.787969in}}%
\pgfpathlineto{\pgfqpoint{1.938531in}{1.788233in}}%
\pgfpathlineto{\pgfqpoint{1.939662in}{1.829664in}}%
\pgfpathlineto{\pgfqpoint{1.981983in}{1.828333in}}%
\pgfpathlineto{\pgfqpoint{2.001203in}{1.827773in}}%
\pgfpathlineto{\pgfqpoint{2.000983in}{1.813993in}}%
\pgfpathlineto{\pgfqpoint{2.025168in}{1.813601in}}%
\pgfpathclose%
\pgfusepath{fill}%
\end{pgfscope}%
\begin{pgfscope}%
\pgfpathrectangle{\pgfqpoint{0.100000in}{0.100000in}}{\pgfqpoint{3.608454in}{2.310000in}}%
\pgfusepath{clip}%
\pgfsetbuttcap%
\pgfsetmiterjoin%
\definecolor{currentfill}{rgb}{0.000000,0.678431,0.660784}%
\pgfsetfillcolor{currentfill}%
\pgfsetlinewidth{0.000000pt}%
\definecolor{currentstroke}{rgb}{0.000000,0.000000,0.000000}%
\pgfsetstrokecolor{currentstroke}%
\pgfsetstrokeopacity{0.000000}%
\pgfsetdash{}{0pt}%
\pgfpathmoveto{\pgfqpoint{1.919219in}{0.761725in}}%
\pgfpathlineto{\pgfqpoint{1.884281in}{0.742465in}}%
\pgfpathlineto{\pgfqpoint{1.870147in}{0.767494in}}%
\pgfpathlineto{\pgfqpoint{1.858178in}{0.766124in}}%
\pgfpathlineto{\pgfqpoint{1.842140in}{0.796889in}}%
\pgfpathlineto{\pgfqpoint{1.873122in}{0.813238in}}%
\pgfpathlineto{\pgfqpoint{1.873743in}{0.830250in}}%
\pgfpathlineto{\pgfqpoint{1.900998in}{0.829156in}}%
\pgfpathlineto{\pgfqpoint{1.908521in}{0.806762in}}%
\pgfpathlineto{\pgfqpoint{1.913591in}{0.794961in}}%
\pgfpathlineto{\pgfqpoint{1.903937in}{0.789730in}}%
\pgfpathlineto{\pgfqpoint{1.919219in}{0.761725in}}%
\pgfpathclose%
\pgfusepath{fill}%
\end{pgfscope}%
\begin{pgfscope}%
\pgfpathrectangle{\pgfqpoint{0.100000in}{0.100000in}}{\pgfqpoint{3.608454in}{2.310000in}}%
\pgfusepath{clip}%
\pgfsetbuttcap%
\pgfsetmiterjoin%
\definecolor{currentfill}{rgb}{0.000000,0.580392,0.709804}%
\pgfsetfillcolor{currentfill}%
\pgfsetlinewidth{0.000000pt}%
\definecolor{currentstroke}{rgb}{0.000000,0.000000,0.000000}%
\pgfsetstrokecolor{currentstroke}%
\pgfsetstrokeopacity{0.000000}%
\pgfsetdash{}{0pt}%
\pgfpathmoveto{\pgfqpoint{3.111250in}{1.435310in}}%
\pgfpathlineto{\pgfqpoint{3.117131in}{1.434042in}}%
\pgfpathlineto{\pgfqpoint{3.119217in}{1.439782in}}%
\pgfpathlineto{\pgfqpoint{3.125843in}{1.434080in}}%
\pgfpathlineto{\pgfqpoint{3.127584in}{1.429231in}}%
\pgfpathlineto{\pgfqpoint{3.124006in}{1.419453in}}%
\pgfpathlineto{\pgfqpoint{3.131195in}{1.415085in}}%
\pgfpathlineto{\pgfqpoint{3.129858in}{1.404228in}}%
\pgfpathlineto{\pgfqpoint{3.125203in}{1.394430in}}%
\pgfpathlineto{\pgfqpoint{3.123623in}{1.385333in}}%
\pgfpathlineto{\pgfqpoint{3.115799in}{1.372179in}}%
\pgfpathlineto{\pgfqpoint{3.110807in}{1.366936in}}%
\pgfpathlineto{\pgfqpoint{3.078058in}{1.382882in}}%
\pgfpathlineto{\pgfqpoint{3.073594in}{1.376808in}}%
\pgfpathlineto{\pgfqpoint{3.062954in}{1.378601in}}%
\pgfpathlineto{\pgfqpoint{3.058026in}{1.385268in}}%
\pgfpathlineto{\pgfqpoint{3.050650in}{1.387261in}}%
\pgfpathlineto{\pgfqpoint{3.056652in}{1.404504in}}%
\pgfpathlineto{\pgfqpoint{3.055324in}{1.410700in}}%
\pgfpathlineto{\pgfqpoint{3.063909in}{1.418899in}}%
\pgfpathlineto{\pgfqpoint{3.084328in}{1.406748in}}%
\pgfpathlineto{\pgfqpoint{3.086938in}{1.414331in}}%
\pgfpathlineto{\pgfqpoint{3.095711in}{1.408904in}}%
\pgfpathlineto{\pgfqpoint{3.099019in}{1.419473in}}%
\pgfpathlineto{\pgfqpoint{3.103999in}{1.422556in}}%
\pgfpathlineto{\pgfqpoint{3.111250in}{1.435310in}}%
\pgfpathclose%
\pgfusepath{fill}%
\end{pgfscope}%
\begin{pgfscope}%
\pgfpathrectangle{\pgfqpoint{0.100000in}{0.100000in}}{\pgfqpoint{3.608454in}{2.310000in}}%
\pgfusepath{clip}%
\pgfsetbuttcap%
\pgfsetmiterjoin%
\definecolor{currentfill}{rgb}{0.000000,0.741176,0.629412}%
\pgfsetfillcolor{currentfill}%
\pgfsetlinewidth{0.000000pt}%
\definecolor{currentstroke}{rgb}{0.000000,0.000000,0.000000}%
\pgfsetstrokecolor{currentstroke}%
\pgfsetstrokeopacity{0.000000}%
\pgfsetdash{}{0pt}%
\pgfpathmoveto{\pgfqpoint{2.414649in}{1.160533in}}%
\pgfpathlineto{\pgfqpoint{2.410619in}{1.149798in}}%
\pgfpathlineto{\pgfqpoint{2.353073in}{1.147045in}}%
\pgfpathlineto{\pgfqpoint{2.352645in}{1.172837in}}%
\pgfpathlineto{\pgfqpoint{2.345750in}{1.172705in}}%
\pgfpathlineto{\pgfqpoint{2.345378in}{1.193313in}}%
\pgfpathlineto{\pgfqpoint{2.357912in}{1.194600in}}%
\pgfpathlineto{\pgfqpoint{2.361419in}{1.191422in}}%
\pgfpathlineto{\pgfqpoint{2.373090in}{1.191695in}}%
\pgfpathlineto{\pgfqpoint{2.375334in}{1.200990in}}%
\pgfpathlineto{\pgfqpoint{2.375052in}{1.209261in}}%
\pgfpathlineto{\pgfqpoint{2.386664in}{1.209942in}}%
\pgfpathlineto{\pgfqpoint{2.386546in}{1.213497in}}%
\pgfpathlineto{\pgfqpoint{2.407413in}{1.214350in}}%
\pgfpathlineto{\pgfqpoint{2.408011in}{1.196285in}}%
\pgfpathlineto{\pgfqpoint{2.415125in}{1.196618in}}%
\pgfpathlineto{\pgfqpoint{2.406457in}{1.183183in}}%
\pgfpathlineto{\pgfqpoint{2.413654in}{1.167791in}}%
\pgfpathlineto{\pgfqpoint{2.414649in}{1.160533in}}%
\pgfpathclose%
\pgfusepath{fill}%
\end{pgfscope}%
\begin{pgfscope}%
\pgfpathrectangle{\pgfqpoint{0.100000in}{0.100000in}}{\pgfqpoint{3.608454in}{2.310000in}}%
\pgfusepath{clip}%
\pgfsetbuttcap%
\pgfsetmiterjoin%
\definecolor{currentfill}{rgb}{0.000000,0.611765,0.694118}%
\pgfsetfillcolor{currentfill}%
\pgfsetlinewidth{0.000000pt}%
\definecolor{currentstroke}{rgb}{0.000000,0.000000,0.000000}%
\pgfsetstrokecolor{currentstroke}%
\pgfsetstrokeopacity{0.000000}%
\pgfsetdash{}{0pt}%
\pgfpathmoveto{\pgfqpoint{1.578022in}{2.163356in}}%
\pgfpathlineto{\pgfqpoint{1.630877in}{2.157836in}}%
\pgfpathlineto{\pgfqpoint{1.626106in}{2.109098in}}%
\pgfpathlineto{\pgfqpoint{1.595340in}{2.112255in}}%
\pgfpathlineto{\pgfqpoint{1.596092in}{2.119196in}}%
\pgfpathlineto{\pgfqpoint{1.589220in}{2.119949in}}%
\pgfpathlineto{\pgfqpoint{1.589978in}{2.126848in}}%
\pgfpathlineto{\pgfqpoint{1.578597in}{2.128096in}}%
\pgfpathlineto{\pgfqpoint{1.575926in}{2.135412in}}%
\pgfpathlineto{\pgfqpoint{1.578243in}{2.156172in}}%
\pgfpathlineto{\pgfqpoint{1.578022in}{2.163356in}}%
\pgfpathclose%
\pgfusepath{fill}%
\end{pgfscope}%
\begin{pgfscope}%
\pgfpathrectangle{\pgfqpoint{0.100000in}{0.100000in}}{\pgfqpoint{3.608454in}{2.310000in}}%
\pgfusepath{clip}%
\pgfsetbuttcap%
\pgfsetmiterjoin%
\definecolor{currentfill}{rgb}{0.000000,0.584314,0.707843}%
\pgfsetfillcolor{currentfill}%
\pgfsetlinewidth{0.000000pt}%
\definecolor{currentstroke}{rgb}{0.000000,0.000000,0.000000}%
\pgfsetstrokecolor{currentstroke}%
\pgfsetstrokeopacity{0.000000}%
\pgfsetdash{}{0pt}%
\pgfpathmoveto{\pgfqpoint{2.118427in}{1.958403in}}%
\pgfpathlineto{\pgfqpoint{2.097902in}{1.958459in}}%
\pgfpathlineto{\pgfqpoint{2.097851in}{1.951504in}}%
\pgfpathlineto{\pgfqpoint{2.042559in}{1.951936in}}%
\pgfpathlineto{\pgfqpoint{2.042858in}{1.972645in}}%
\pgfpathlineto{\pgfqpoint{2.041913in}{1.986517in}}%
\pgfpathlineto{\pgfqpoint{2.048844in}{1.986478in}}%
\pgfpathlineto{\pgfqpoint{2.049145in}{2.014243in}}%
\pgfpathlineto{\pgfqpoint{2.076864in}{2.014087in}}%
\pgfpathlineto{\pgfqpoint{2.075622in}{2.027945in}}%
\pgfpathlineto{\pgfqpoint{2.075577in}{2.048780in}}%
\pgfpathlineto{\pgfqpoint{2.075068in}{2.055715in}}%
\pgfpathlineto{\pgfqpoint{2.074816in}{2.097239in}}%
\pgfpathlineto{\pgfqpoint{2.095537in}{2.097237in}}%
\pgfpathlineto{\pgfqpoint{2.095560in}{2.083291in}}%
\pgfpathlineto{\pgfqpoint{2.137206in}{2.083483in}}%
\pgfpathlineto{\pgfqpoint{2.138015in}{2.041726in}}%
\pgfpathlineto{\pgfqpoint{2.138354in}{2.007024in}}%
\pgfpathlineto{\pgfqpoint{2.124628in}{2.006820in}}%
\pgfpathlineto{\pgfqpoint{2.125367in}{1.958378in}}%
\pgfpathlineto{\pgfqpoint{2.118427in}{1.958403in}}%
\pgfpathclose%
\pgfusepath{fill}%
\end{pgfscope}%
\begin{pgfscope}%
\pgfpathrectangle{\pgfqpoint{0.100000in}{0.100000in}}{\pgfqpoint{3.608454in}{2.310000in}}%
\pgfusepath{clip}%
\pgfsetbuttcap%
\pgfsetmiterjoin%
\definecolor{currentfill}{rgb}{0.000000,0.666667,0.666667}%
\pgfsetfillcolor{currentfill}%
\pgfsetlinewidth{0.000000pt}%
\definecolor{currentstroke}{rgb}{0.000000,0.000000,0.000000}%
\pgfsetstrokecolor{currentstroke}%
\pgfsetstrokeopacity{0.000000}%
\pgfsetdash{}{0pt}%
\pgfpathmoveto{\pgfqpoint{2.272778in}{0.593877in}}%
\pgfpathlineto{\pgfqpoint{2.271892in}{0.630559in}}%
\pgfpathlineto{\pgfqpoint{2.269704in}{0.666105in}}%
\pgfpathlineto{\pgfqpoint{2.272007in}{0.674022in}}%
\pgfpathlineto{\pgfqpoint{2.271317in}{0.699125in}}%
\pgfpathlineto{\pgfqpoint{2.278789in}{0.700272in}}%
\pgfpathlineto{\pgfqpoint{2.285347in}{0.707925in}}%
\pgfpathlineto{\pgfqpoint{2.292928in}{0.702917in}}%
\pgfpathlineto{\pgfqpoint{2.296084in}{0.696005in}}%
\pgfpathlineto{\pgfqpoint{2.325022in}{0.696944in}}%
\pgfpathlineto{\pgfqpoint{2.331030in}{0.685463in}}%
\pgfpathlineto{\pgfqpoint{2.329801in}{0.674430in}}%
\pgfpathlineto{\pgfqpoint{2.331431in}{0.670010in}}%
\pgfpathlineto{\pgfqpoint{2.338216in}{0.665019in}}%
\pgfpathlineto{\pgfqpoint{2.339922in}{0.654462in}}%
\pgfpathlineto{\pgfqpoint{2.344148in}{0.649181in}}%
\pgfpathlineto{\pgfqpoint{2.350223in}{0.648417in}}%
\pgfpathlineto{\pgfqpoint{2.351410in}{0.638299in}}%
\pgfpathlineto{\pgfqpoint{2.368222in}{0.632799in}}%
\pgfpathlineto{\pgfqpoint{2.365882in}{0.630690in}}%
\pgfpathlineto{\pgfqpoint{2.371581in}{0.618635in}}%
\pgfpathlineto{\pgfqpoint{2.376744in}{0.617237in}}%
\pgfpathlineto{\pgfqpoint{2.378486in}{0.611623in}}%
\pgfpathlineto{\pgfqpoint{2.372197in}{0.609200in}}%
\pgfpathlineto{\pgfqpoint{2.357771in}{0.613019in}}%
\pgfpathlineto{\pgfqpoint{2.358501in}{0.617112in}}%
\pgfpathlineto{\pgfqpoint{2.351303in}{0.626566in}}%
\pgfpathlineto{\pgfqpoint{2.340753in}{0.625307in}}%
\pgfpathlineto{\pgfqpoint{2.335052in}{0.617153in}}%
\pgfpathlineto{\pgfqpoint{2.323648in}{0.606472in}}%
\pgfpathlineto{\pgfqpoint{2.323245in}{0.615819in}}%
\pgfpathlineto{\pgfqpoint{2.312586in}{0.611195in}}%
\pgfpathlineto{\pgfqpoint{2.305299in}{0.605725in}}%
\pgfpathlineto{\pgfqpoint{2.307552in}{0.598624in}}%
\pgfpathlineto{\pgfqpoint{2.315848in}{0.597697in}}%
\pgfpathlineto{\pgfqpoint{2.326725in}{0.599490in}}%
\pgfpathlineto{\pgfqpoint{2.336035in}{0.593898in}}%
\pgfpathlineto{\pgfqpoint{2.328591in}{0.587197in}}%
\pgfpathlineto{\pgfqpoint{2.311674in}{0.595112in}}%
\pgfpathlineto{\pgfqpoint{2.304943in}{0.595012in}}%
\pgfpathlineto{\pgfqpoint{2.294524in}{0.590708in}}%
\pgfpathlineto{\pgfqpoint{2.272778in}{0.593877in}}%
\pgfpathclose%
\pgfusepath{fill}%
\end{pgfscope}%
\begin{pgfscope}%
\pgfpathrectangle{\pgfqpoint{0.100000in}{0.100000in}}{\pgfqpoint{3.608454in}{2.310000in}}%
\pgfusepath{clip}%
\pgfsetbuttcap%
\pgfsetmiterjoin%
\definecolor{currentfill}{rgb}{0.000000,0.682353,0.658824}%
\pgfsetfillcolor{currentfill}%
\pgfsetlinewidth{0.000000pt}%
\definecolor{currentstroke}{rgb}{0.000000,0.000000,0.000000}%
\pgfsetstrokecolor{currentstroke}%
\pgfsetstrokeopacity{0.000000}%
\pgfsetdash{}{0pt}%
\pgfpathmoveto{\pgfqpoint{1.836840in}{0.475148in}}%
\pgfpathlineto{\pgfqpoint{1.795200in}{0.474847in}}%
\pgfpathlineto{\pgfqpoint{1.795506in}{0.488863in}}%
\pgfpathlineto{\pgfqpoint{1.797091in}{0.523880in}}%
\pgfpathlineto{\pgfqpoint{1.796154in}{0.523908in}}%
\pgfpathlineto{\pgfqpoint{1.797569in}{0.560061in}}%
\pgfpathlineto{\pgfqpoint{1.840202in}{0.558059in}}%
\pgfpathlineto{\pgfqpoint{1.839000in}{0.522528in}}%
\pgfpathlineto{\pgfqpoint{1.836840in}{0.475148in}}%
\pgfpathclose%
\pgfusepath{fill}%
\end{pgfscope}%
\begin{pgfscope}%
\pgfpathrectangle{\pgfqpoint{0.100000in}{0.100000in}}{\pgfqpoint{3.608454in}{2.310000in}}%
\pgfusepath{clip}%
\pgfsetbuttcap%
\pgfsetmiterjoin%
\definecolor{currentfill}{rgb}{0.000000,0.588235,0.705882}%
\pgfsetfillcolor{currentfill}%
\pgfsetlinewidth{0.000000pt}%
\definecolor{currentstroke}{rgb}{0.000000,0.000000,0.000000}%
\pgfsetstrokecolor{currentstroke}%
\pgfsetstrokeopacity{0.000000}%
\pgfsetdash{}{0pt}%
\pgfpathmoveto{\pgfqpoint{2.966459in}{0.731456in}}%
\pgfpathlineto{\pgfqpoint{2.957900in}{0.730883in}}%
\pgfpathlineto{\pgfqpoint{2.955222in}{0.740323in}}%
\pgfpathlineto{\pgfqpoint{2.943923in}{0.740325in}}%
\pgfpathlineto{\pgfqpoint{2.937427in}{0.747697in}}%
\pgfpathlineto{\pgfqpoint{2.928578in}{0.749275in}}%
\pgfpathlineto{\pgfqpoint{2.925268in}{0.774204in}}%
\pgfpathlineto{\pgfqpoint{2.920160in}{0.773477in}}%
\pgfpathlineto{\pgfqpoint{2.920094in}{0.780812in}}%
\pgfpathlineto{\pgfqpoint{2.911798in}{0.789691in}}%
\pgfpathlineto{\pgfqpoint{2.910462in}{0.795395in}}%
\pgfpathlineto{\pgfqpoint{2.912735in}{0.796174in}}%
\pgfpathlineto{\pgfqpoint{2.914876in}{0.807165in}}%
\pgfpathlineto{\pgfqpoint{2.918401in}{0.811582in}}%
\pgfpathlineto{\pgfqpoint{2.917347in}{0.821092in}}%
\pgfpathlineto{\pgfqpoint{2.926243in}{0.822499in}}%
\pgfpathlineto{\pgfqpoint{2.929182in}{0.813904in}}%
\pgfpathlineto{\pgfqpoint{2.953500in}{0.817959in}}%
\pgfpathlineto{\pgfqpoint{2.970476in}{0.810318in}}%
\pgfpathlineto{\pgfqpoint{2.978722in}{0.804662in}}%
\pgfpathlineto{\pgfqpoint{2.983849in}{0.804589in}}%
\pgfpathlineto{\pgfqpoint{2.986023in}{0.797823in}}%
\pgfpathlineto{\pgfqpoint{2.990030in}{0.793119in}}%
\pgfpathlineto{\pgfqpoint{2.983748in}{0.789742in}}%
\pgfpathlineto{\pgfqpoint{2.977554in}{0.783183in}}%
\pgfpathlineto{\pgfqpoint{2.977965in}{0.773429in}}%
\pgfpathlineto{\pgfqpoint{2.983960in}{0.768316in}}%
\pgfpathlineto{\pgfqpoint{2.964698in}{0.765826in}}%
\pgfpathlineto{\pgfqpoint{2.966799in}{0.748504in}}%
\pgfpathlineto{\pgfqpoint{2.985291in}{0.750228in}}%
\pgfpathlineto{\pgfqpoint{2.983300in}{0.732552in}}%
\pgfpathlineto{\pgfqpoint{2.966459in}{0.731456in}}%
\pgfpathclose%
\pgfusepath{fill}%
\end{pgfscope}%
\begin{pgfscope}%
\pgfpathrectangle{\pgfqpoint{0.100000in}{0.100000in}}{\pgfqpoint{3.608454in}{2.310000in}}%
\pgfusepath{clip}%
\pgfsetbuttcap%
\pgfsetmiterjoin%
\definecolor{currentfill}{rgb}{0.000000,0.670588,0.664706}%
\pgfsetfillcolor{currentfill}%
\pgfsetlinewidth{0.000000pt}%
\definecolor{currentstroke}{rgb}{0.000000,0.000000,0.000000}%
\pgfsetstrokecolor{currentstroke}%
\pgfsetstrokeopacity{0.000000}%
\pgfsetdash{}{0pt}%
\pgfpathmoveto{\pgfqpoint{0.579411in}{1.933417in}}%
\pgfpathlineto{\pgfqpoint{0.571017in}{1.931177in}}%
\pgfpathlineto{\pgfqpoint{0.568235in}{1.922352in}}%
\pgfpathlineto{\pgfqpoint{0.534705in}{1.932520in}}%
\pgfpathlineto{\pgfqpoint{0.537209in}{1.940659in}}%
\pgfpathlineto{\pgfqpoint{0.517068in}{1.946672in}}%
\pgfpathlineto{\pgfqpoint{0.517042in}{1.952171in}}%
\pgfpathlineto{\pgfqpoint{0.521206in}{1.965301in}}%
\pgfpathlineto{\pgfqpoint{0.509770in}{1.969016in}}%
\pgfpathlineto{\pgfqpoint{0.501091in}{1.975934in}}%
\pgfpathlineto{\pgfqpoint{0.498412in}{1.984367in}}%
\pgfpathlineto{\pgfqpoint{0.487714in}{1.989945in}}%
\pgfpathlineto{\pgfqpoint{0.479271in}{1.985850in}}%
\pgfpathlineto{\pgfqpoint{0.467729in}{1.989714in}}%
\pgfpathlineto{\pgfqpoint{0.475658in}{2.007052in}}%
\pgfpathlineto{\pgfqpoint{0.480357in}{2.020001in}}%
\pgfpathlineto{\pgfqpoint{0.498912in}{2.014543in}}%
\pgfpathlineto{\pgfqpoint{0.504063in}{2.019289in}}%
\pgfpathlineto{\pgfqpoint{0.505570in}{2.024850in}}%
\pgfpathlineto{\pgfqpoint{0.512206in}{2.022672in}}%
\pgfpathlineto{\pgfqpoint{0.519006in}{2.044536in}}%
\pgfpathlineto{\pgfqpoint{0.512881in}{2.048028in}}%
\pgfpathlineto{\pgfqpoint{0.521244in}{2.073470in}}%
\pgfpathlineto{\pgfqpoint{0.518018in}{2.074545in}}%
\pgfpathlineto{\pgfqpoint{0.521499in}{2.085093in}}%
\pgfpathlineto{\pgfqpoint{0.538605in}{2.079539in}}%
\pgfpathlineto{\pgfqpoint{0.543848in}{2.096029in}}%
\pgfpathlineto{\pgfqpoint{0.561348in}{2.090351in}}%
\pgfpathlineto{\pgfqpoint{0.566154in}{2.086439in}}%
\pgfpathlineto{\pgfqpoint{0.568024in}{2.080965in}}%
\pgfpathlineto{\pgfqpoint{0.576179in}{2.072349in}}%
\pgfpathlineto{\pgfqpoint{0.575705in}{2.064800in}}%
\pgfpathlineto{\pgfqpoint{0.572541in}{2.061484in}}%
\pgfpathlineto{\pgfqpoint{0.575268in}{2.054730in}}%
\pgfpathlineto{\pgfqpoint{0.580214in}{2.049876in}}%
\pgfpathlineto{\pgfqpoint{0.581481in}{2.042623in}}%
\pgfpathlineto{\pgfqpoint{0.587828in}{2.036379in}}%
\pgfpathlineto{\pgfqpoint{0.623427in}{2.025636in}}%
\pgfpathlineto{\pgfqpoint{0.620648in}{2.021281in}}%
\pgfpathlineto{\pgfqpoint{0.616042in}{2.014692in}}%
\pgfpathlineto{\pgfqpoint{0.613506in}{2.004223in}}%
\pgfpathlineto{\pgfqpoint{0.608267in}{1.997300in}}%
\pgfpathlineto{\pgfqpoint{0.605387in}{1.987819in}}%
\pgfpathlineto{\pgfqpoint{0.605729in}{1.977984in}}%
\pgfpathlineto{\pgfqpoint{0.603007in}{1.965607in}}%
\pgfpathlineto{\pgfqpoint{0.595256in}{1.957298in}}%
\pgfpathlineto{\pgfqpoint{0.586510in}{1.951287in}}%
\pgfpathlineto{\pgfqpoint{0.585151in}{1.943983in}}%
\pgfpathlineto{\pgfqpoint{0.579411in}{1.933417in}}%
\pgfpathclose%
\pgfusepath{fill}%
\end{pgfscope}%
\begin{pgfscope}%
\pgfpathrectangle{\pgfqpoint{0.100000in}{0.100000in}}{\pgfqpoint{3.608454in}{2.310000in}}%
\pgfusepath{clip}%
\pgfsetbuttcap%
\pgfsetmiterjoin%
\definecolor{currentfill}{rgb}{0.000000,0.545098,0.727451}%
\pgfsetfillcolor{currentfill}%
\pgfsetlinewidth{0.000000pt}%
\definecolor{currentstroke}{rgb}{0.000000,0.000000,0.000000}%
\pgfsetstrokecolor{currentstroke}%
\pgfsetstrokeopacity{0.000000}%
\pgfsetdash{}{0pt}%
\pgfpathmoveto{\pgfqpoint{2.765201in}{0.880136in}}%
\pgfpathlineto{\pgfqpoint{2.755362in}{0.877808in}}%
\pgfpathlineto{\pgfqpoint{2.734886in}{0.875470in}}%
\pgfpathlineto{\pgfqpoint{2.731773in}{0.905420in}}%
\pgfpathlineto{\pgfqpoint{2.727775in}{0.904950in}}%
\pgfpathlineto{\pgfqpoint{2.725270in}{0.935776in}}%
\pgfpathlineto{\pgfqpoint{2.717216in}{0.935165in}}%
\pgfpathlineto{\pgfqpoint{2.714599in}{0.941836in}}%
\pgfpathlineto{\pgfqpoint{2.701405in}{0.940587in}}%
\pgfpathlineto{\pgfqpoint{2.697028in}{0.947184in}}%
\pgfpathlineto{\pgfqpoint{2.690775in}{0.946928in}}%
\pgfpathlineto{\pgfqpoint{2.696367in}{0.955422in}}%
\pgfpathlineto{\pgfqpoint{2.694783in}{0.960330in}}%
\pgfpathlineto{\pgfqpoint{2.705960in}{0.970065in}}%
\pgfpathlineto{\pgfqpoint{2.715165in}{0.972447in}}%
\pgfpathlineto{\pgfqpoint{2.728983in}{0.972881in}}%
\pgfpathlineto{\pgfqpoint{2.737463in}{0.974424in}}%
\pgfpathlineto{\pgfqpoint{2.738788in}{0.969584in}}%
\pgfpathlineto{\pgfqpoint{2.760784in}{0.972182in}}%
\pgfpathlineto{\pgfqpoint{2.762398in}{0.964970in}}%
\pgfpathlineto{\pgfqpoint{2.771603in}{0.963496in}}%
\pgfpathlineto{\pgfqpoint{2.773147in}{0.947109in}}%
\pgfpathlineto{\pgfqpoint{2.778830in}{0.942946in}}%
\pgfpathlineto{\pgfqpoint{2.791999in}{0.944257in}}%
\pgfpathlineto{\pgfqpoint{2.796122in}{0.932528in}}%
\pgfpathlineto{\pgfqpoint{2.802892in}{0.925244in}}%
\pgfpathlineto{\pgfqpoint{2.802911in}{0.922322in}}%
\pgfpathlineto{\pgfqpoint{2.779766in}{0.919893in}}%
\pgfpathlineto{\pgfqpoint{2.782280in}{0.891779in}}%
\pgfpathlineto{\pgfqpoint{2.760875in}{0.889582in}}%
\pgfpathlineto{\pgfqpoint{2.765201in}{0.880136in}}%
\pgfpathclose%
\pgfusepath{fill}%
\end{pgfscope}%
\begin{pgfscope}%
\pgfpathrectangle{\pgfqpoint{0.100000in}{0.100000in}}{\pgfqpoint{3.608454in}{2.310000in}}%
\pgfusepath{clip}%
\pgfsetbuttcap%
\pgfsetmiterjoin%
\definecolor{currentfill}{rgb}{0.000000,0.556863,0.721569}%
\pgfsetfillcolor{currentfill}%
\pgfsetlinewidth{0.000000pt}%
\definecolor{currentstroke}{rgb}{0.000000,0.000000,0.000000}%
\pgfsetstrokecolor{currentstroke}%
\pgfsetstrokeopacity{0.000000}%
\pgfsetdash{}{0pt}%
\pgfpathmoveto{\pgfqpoint{1.305818in}{1.299625in}}%
\pgfpathlineto{\pgfqpoint{1.298775in}{1.307121in}}%
\pgfpathlineto{\pgfqpoint{1.284736in}{1.306313in}}%
\pgfpathlineto{\pgfqpoint{1.282717in}{1.312481in}}%
\pgfpathlineto{\pgfqpoint{1.248158in}{1.317741in}}%
\pgfpathlineto{\pgfqpoint{1.234133in}{1.318758in}}%
\pgfpathlineto{\pgfqpoint{1.237566in}{1.340907in}}%
\pgfpathlineto{\pgfqpoint{1.237792in}{1.349736in}}%
\pgfpathlineto{\pgfqpoint{1.240533in}{1.367269in}}%
\pgfpathlineto{\pgfqpoint{1.251645in}{1.434868in}}%
\pgfpathlineto{\pgfqpoint{1.290367in}{1.428918in}}%
\pgfpathlineto{\pgfqpoint{1.349566in}{1.420460in}}%
\pgfpathlineto{\pgfqpoint{1.346932in}{1.416687in}}%
\pgfpathlineto{\pgfqpoint{1.350593in}{1.411552in}}%
\pgfpathlineto{\pgfqpoint{1.343717in}{1.409436in}}%
\pgfpathlineto{\pgfqpoint{1.333852in}{1.337711in}}%
\pgfpathlineto{\pgfqpoint{1.325556in}{1.338863in}}%
\pgfpathlineto{\pgfqpoint{1.328146in}{1.332828in}}%
\pgfpathlineto{\pgfqpoint{1.328108in}{1.320759in}}%
\pgfpathlineto{\pgfqpoint{1.325950in}{1.311919in}}%
\pgfpathlineto{\pgfqpoint{1.314836in}{1.308720in}}%
\pgfpathlineto{\pgfqpoint{1.305818in}{1.299625in}}%
\pgfpathclose%
\pgfusepath{fill}%
\end{pgfscope}%
\begin{pgfscope}%
\pgfpathrectangle{\pgfqpoint{0.100000in}{0.100000in}}{\pgfqpoint{3.608454in}{2.310000in}}%
\pgfusepath{clip}%
\pgfsetbuttcap%
\pgfsetmiterjoin%
\definecolor{currentfill}{rgb}{0.000000,0.827451,0.586275}%
\pgfsetfillcolor{currentfill}%
\pgfsetlinewidth{0.000000pt}%
\definecolor{currentstroke}{rgb}{0.000000,0.000000,0.000000}%
\pgfsetstrokecolor{currentstroke}%
\pgfsetstrokeopacity{0.000000}%
\pgfsetdash{}{0pt}%
\pgfpathmoveto{\pgfqpoint{1.124077in}{0.559590in}}%
\pgfpathlineto{\pgfqpoint{1.110388in}{0.567904in}}%
\pgfpathlineto{\pgfqpoint{1.083089in}{0.584700in}}%
\pgfpathlineto{\pgfqpoint{1.074192in}{0.589904in}}%
\pgfpathlineto{\pgfqpoint{1.066205in}{0.578011in}}%
\pgfpathlineto{\pgfqpoint{1.056440in}{0.584227in}}%
\pgfpathlineto{\pgfqpoint{1.052163in}{0.586365in}}%
\pgfpathlineto{\pgfqpoint{1.026608in}{0.603847in}}%
\pgfpathlineto{\pgfqpoint{0.997892in}{0.624739in}}%
\pgfpathlineto{\pgfqpoint{0.976297in}{0.641465in}}%
\pgfpathlineto{\pgfqpoint{0.976829in}{0.642165in}}%
\pgfpathlineto{\pgfqpoint{0.968183in}{0.648986in}}%
\pgfpathlineto{\pgfqpoint{0.969832in}{0.651052in}}%
\pgfpathlineto{\pgfqpoint{0.961484in}{0.657889in}}%
\pgfpathlineto{\pgfqpoint{0.963280in}{0.660093in}}%
\pgfpathlineto{\pgfqpoint{0.954477in}{0.667338in}}%
\pgfpathlineto{\pgfqpoint{0.940158in}{0.679506in}}%
\pgfpathlineto{\pgfqpoint{0.941559in}{0.681130in}}%
\pgfpathlineto{\pgfqpoint{0.937775in}{0.684411in}}%
\pgfpathlineto{\pgfqpoint{0.936363in}{0.682792in}}%
\pgfpathlineto{\pgfqpoint{0.927789in}{0.690351in}}%
\pgfpathlineto{\pgfqpoint{0.929228in}{0.691967in}}%
\pgfpathlineto{\pgfqpoint{0.925549in}{0.695258in}}%
\pgfpathlineto{\pgfqpoint{0.924102in}{0.693649in}}%
\pgfpathlineto{\pgfqpoint{0.919895in}{0.697448in}}%
\pgfpathlineto{\pgfqpoint{0.921351in}{0.699050in}}%
\pgfpathlineto{\pgfqpoint{0.915131in}{0.704745in}}%
\pgfpathlineto{\pgfqpoint{0.913662in}{0.703154in}}%
\pgfpathlineto{\pgfqpoint{0.898458in}{0.717465in}}%
\pgfpathlineto{\pgfqpoint{0.894146in}{0.712972in}}%
\pgfpathlineto{\pgfqpoint{0.886741in}{0.720168in}}%
\pgfpathlineto{\pgfqpoint{0.883880in}{0.727259in}}%
\pgfpathlineto{\pgfqpoint{0.884750in}{0.730564in}}%
\pgfpathlineto{\pgfqpoint{0.884173in}{0.733970in}}%
\pgfpathlineto{\pgfqpoint{0.881741in}{0.738770in}}%
\pgfpathlineto{\pgfqpoint{0.884973in}{0.737504in}}%
\pgfpathlineto{\pgfqpoint{0.888394in}{0.737302in}}%
\pgfpathlineto{\pgfqpoint{0.891146in}{0.738897in}}%
\pgfpathlineto{\pgfqpoint{0.897248in}{0.744755in}}%
\pgfpathlineto{\pgfqpoint{0.899391in}{0.744053in}}%
\pgfpathlineto{\pgfqpoint{0.904078in}{0.738998in}}%
\pgfpathlineto{\pgfqpoint{0.914603in}{0.732006in}}%
\pgfpathlineto{\pgfqpoint{0.917793in}{0.730737in}}%
\pgfpathlineto{\pgfqpoint{0.924697in}{0.730323in}}%
\pgfpathlineto{\pgfqpoint{0.934212in}{0.732247in}}%
\pgfpathlineto{\pgfqpoint{0.938989in}{0.735849in}}%
\pgfpathlineto{\pgfqpoint{0.940161in}{0.737749in}}%
\pgfpathlineto{\pgfqpoint{0.943533in}{0.740578in}}%
\pgfpathlineto{\pgfqpoint{0.945793in}{0.741496in}}%
\pgfpathlineto{\pgfqpoint{0.952899in}{0.741432in}}%
\pgfpathlineto{\pgfqpoint{0.955066in}{0.743565in}}%
\pgfpathlineto{\pgfqpoint{0.959761in}{0.743939in}}%
\pgfpathlineto{\pgfqpoint{0.965368in}{0.745161in}}%
\pgfpathlineto{\pgfqpoint{0.968573in}{0.740621in}}%
\pgfpathlineto{\pgfqpoint{0.974554in}{0.739226in}}%
\pgfpathlineto{\pgfqpoint{0.985159in}{0.739890in}}%
\pgfpathlineto{\pgfqpoint{0.995988in}{0.741765in}}%
\pgfpathlineto{\pgfqpoint{0.999534in}{0.740148in}}%
\pgfpathlineto{\pgfqpoint{0.999964in}{0.737676in}}%
\pgfpathlineto{\pgfqpoint{1.004192in}{0.734276in}}%
\pgfpathlineto{\pgfqpoint{1.008232in}{0.732003in}}%
\pgfpathlineto{\pgfqpoint{1.011631in}{0.730934in}}%
\pgfpathlineto{\pgfqpoint{1.018774in}{0.731477in}}%
\pgfpathlineto{\pgfqpoint{1.028230in}{0.734980in}}%
\pgfpathlineto{\pgfqpoint{1.031476in}{0.735474in}}%
\pgfpathlineto{\pgfqpoint{1.033074in}{0.730543in}}%
\pgfpathlineto{\pgfqpoint{1.032211in}{0.727412in}}%
\pgfpathlineto{\pgfqpoint{1.034678in}{0.727466in}}%
\pgfpathlineto{\pgfqpoint{1.036005in}{0.725068in}}%
\pgfpathlineto{\pgfqpoint{1.036019in}{0.722231in}}%
\pgfpathlineto{\pgfqpoint{1.030729in}{0.720182in}}%
\pgfpathlineto{\pgfqpoint{1.026664in}{0.721652in}}%
\pgfpathlineto{\pgfqpoint{1.027301in}{0.717005in}}%
\pgfpathlineto{\pgfqpoint{1.031426in}{0.715395in}}%
\pgfpathlineto{\pgfqpoint{1.033829in}{0.717336in}}%
\pgfpathlineto{\pgfqpoint{1.039293in}{0.717228in}}%
\pgfpathlineto{\pgfqpoint{1.042163in}{0.713874in}}%
\pgfpathlineto{\pgfqpoint{1.038932in}{0.710194in}}%
\pgfpathlineto{\pgfqpoint{1.038327in}{0.708052in}}%
\pgfpathlineto{\pgfqpoint{1.040735in}{0.706277in}}%
\pgfpathlineto{\pgfqpoint{1.041143in}{0.703816in}}%
\pgfpathlineto{\pgfqpoint{1.044548in}{0.704736in}}%
\pgfpathlineto{\pgfqpoint{1.049310in}{0.702353in}}%
\pgfpathlineto{\pgfqpoint{1.052657in}{0.701365in}}%
\pgfpathlineto{\pgfqpoint{1.054785in}{0.696074in}}%
\pgfpathlineto{\pgfqpoint{1.056904in}{0.696328in}}%
\pgfpathlineto{\pgfqpoint{1.058734in}{0.693326in}}%
\pgfpathlineto{\pgfqpoint{1.052615in}{0.689071in}}%
\pgfpathlineto{\pgfqpoint{1.050791in}{0.686844in}}%
\pgfpathlineto{\pgfqpoint{1.057270in}{0.682397in}}%
\pgfpathlineto{\pgfqpoint{1.056506in}{0.680783in}}%
\pgfpathlineto{\pgfqpoint{1.053335in}{0.680298in}}%
\pgfpathlineto{\pgfqpoint{1.058156in}{0.676641in}}%
\pgfpathlineto{\pgfqpoint{1.058655in}{0.673622in}}%
\pgfpathlineto{\pgfqpoint{1.061740in}{0.674956in}}%
\pgfpathlineto{\pgfqpoint{1.066675in}{0.673514in}}%
\pgfpathlineto{\pgfqpoint{1.067308in}{0.671656in}}%
\pgfpathlineto{\pgfqpoint{1.064920in}{0.669776in}}%
\pgfpathlineto{\pgfqpoint{1.069974in}{0.668359in}}%
\pgfpathlineto{\pgfqpoint{1.073388in}{0.668942in}}%
\pgfpathlineto{\pgfqpoint{1.077468in}{0.664660in}}%
\pgfpathlineto{\pgfqpoint{1.078685in}{0.664612in}}%
\pgfpathlineto{\pgfqpoint{1.079652in}{0.659570in}}%
\pgfpathlineto{\pgfqpoint{1.082792in}{0.657912in}}%
\pgfpathlineto{\pgfqpoint{1.082915in}{0.655259in}}%
\pgfpathlineto{\pgfqpoint{1.085209in}{0.653941in}}%
\pgfpathlineto{\pgfqpoint{1.085707in}{0.649109in}}%
\pgfpathlineto{\pgfqpoint{1.092374in}{0.641077in}}%
\pgfpathlineto{\pgfqpoint{1.095064in}{0.640992in}}%
\pgfpathlineto{\pgfqpoint{1.098844in}{0.638529in}}%
\pgfpathlineto{\pgfqpoint{1.101206in}{0.635275in}}%
\pgfpathlineto{\pgfqpoint{1.104111in}{0.633554in}}%
\pgfpathlineto{\pgfqpoint{1.106568in}{0.626099in}}%
\pgfpathlineto{\pgfqpoint{1.108579in}{0.623361in}}%
\pgfpathlineto{\pgfqpoint{1.113936in}{0.621999in}}%
\pgfpathlineto{\pgfqpoint{1.120984in}{0.619591in}}%
\pgfpathlineto{\pgfqpoint{1.123831in}{0.619357in}}%
\pgfpathlineto{\pgfqpoint{1.127068in}{0.617346in}}%
\pgfpathlineto{\pgfqpoint{1.130555in}{0.612352in}}%
\pgfpathlineto{\pgfqpoint{1.132716in}{0.606694in}}%
\pgfpathlineto{\pgfqpoint{1.132681in}{0.605047in}}%
\pgfpathlineto{\pgfqpoint{1.134831in}{0.600830in}}%
\pgfpathlineto{\pgfqpoint{1.137136in}{0.597981in}}%
\pgfpathlineto{\pgfqpoint{1.137981in}{0.595190in}}%
\pgfpathlineto{\pgfqpoint{1.141010in}{0.591617in}}%
\pgfpathlineto{\pgfqpoint{1.139745in}{0.587853in}}%
\pgfpathlineto{\pgfqpoint{1.124077in}{0.559590in}}%
\pgfpathclose%
\pgfusepath{fill}%
\end{pgfscope}%
\begin{pgfscope}%
\pgfpathrectangle{\pgfqpoint{0.100000in}{0.100000in}}{\pgfqpoint{3.608454in}{2.310000in}}%
\pgfusepath{clip}%
\pgfsetbuttcap%
\pgfsetmiterjoin%
\definecolor{currentfill}{rgb}{0.000000,0.776471,0.611765}%
\pgfsetfillcolor{currentfill}%
\pgfsetlinewidth{0.000000pt}%
\definecolor{currentstroke}{rgb}{0.000000,0.000000,0.000000}%
\pgfsetstrokecolor{currentstroke}%
\pgfsetstrokeopacity{0.000000}%
\pgfsetdash{}{0pt}%
\pgfpathmoveto{\pgfqpoint{2.523585in}{1.228784in}}%
\pgfpathlineto{\pgfqpoint{2.519923in}{1.230014in}}%
\pgfpathlineto{\pgfqpoint{2.519105in}{1.243921in}}%
\pgfpathlineto{\pgfqpoint{2.500587in}{1.242556in}}%
\pgfpathlineto{\pgfqpoint{2.499015in}{1.266823in}}%
\pgfpathlineto{\pgfqpoint{2.519677in}{1.268419in}}%
\pgfpathlineto{\pgfqpoint{2.534479in}{1.269453in}}%
\pgfpathlineto{\pgfqpoint{2.537345in}{1.268351in}}%
\pgfpathlineto{\pgfqpoint{2.541890in}{1.261426in}}%
\pgfpathlineto{\pgfqpoint{2.534558in}{1.250105in}}%
\pgfpathlineto{\pgfqpoint{2.540147in}{1.235271in}}%
\pgfpathlineto{\pgfqpoint{2.528115in}{1.232902in}}%
\pgfpathlineto{\pgfqpoint{2.523585in}{1.228784in}}%
\pgfpathclose%
\pgfusepath{fill}%
\end{pgfscope}%
\begin{pgfscope}%
\pgfpathrectangle{\pgfqpoint{0.100000in}{0.100000in}}{\pgfqpoint{3.608454in}{2.310000in}}%
\pgfusepath{clip}%
\pgfsetbuttcap%
\pgfsetmiterjoin%
\definecolor{currentfill}{rgb}{0.000000,0.643137,0.678431}%
\pgfsetfillcolor{currentfill}%
\pgfsetlinewidth{0.000000pt}%
\definecolor{currentstroke}{rgb}{0.000000,0.000000,0.000000}%
\pgfsetstrokecolor{currentstroke}%
\pgfsetstrokeopacity{0.000000}%
\pgfsetdash{}{0pt}%
\pgfpathmoveto{\pgfqpoint{2.468504in}{0.944341in}}%
\pgfpathlineto{\pgfqpoint{2.451281in}{0.943313in}}%
\pgfpathlineto{\pgfqpoint{2.440822in}{0.945153in}}%
\pgfpathlineto{\pgfqpoint{2.441137in}{0.939280in}}%
\pgfpathlineto{\pgfqpoint{2.427361in}{0.937371in}}%
\pgfpathlineto{\pgfqpoint{2.427713in}{0.930421in}}%
\pgfpathlineto{\pgfqpoint{2.416509in}{0.931498in}}%
\pgfpathlineto{\pgfqpoint{2.424326in}{0.937210in}}%
\pgfpathlineto{\pgfqpoint{2.406598in}{0.936460in}}%
\pgfpathlineto{\pgfqpoint{2.405968in}{0.950277in}}%
\pgfpathlineto{\pgfqpoint{2.392184in}{0.949754in}}%
\pgfpathlineto{\pgfqpoint{2.391745in}{0.960160in}}%
\pgfpathlineto{\pgfqpoint{2.372289in}{0.959362in}}%
\pgfpathlineto{\pgfqpoint{2.384214in}{0.972642in}}%
\pgfpathlineto{\pgfqpoint{2.384480in}{0.979541in}}%
\pgfpathlineto{\pgfqpoint{2.395789in}{0.983518in}}%
\pgfpathlineto{\pgfqpoint{2.396088in}{0.992607in}}%
\pgfpathlineto{\pgfqpoint{2.407574in}{0.985155in}}%
\pgfpathlineto{\pgfqpoint{2.413949in}{0.985540in}}%
\pgfpathlineto{\pgfqpoint{2.413440in}{0.992269in}}%
\pgfpathlineto{\pgfqpoint{2.420148in}{0.996176in}}%
\pgfpathlineto{\pgfqpoint{2.467215in}{0.998833in}}%
\pgfpathlineto{\pgfqpoint{2.482530in}{0.995135in}}%
\pgfpathlineto{\pgfqpoint{2.483100in}{0.985946in}}%
\pgfpathlineto{\pgfqpoint{2.484251in}{0.968666in}}%
\pgfpathlineto{\pgfqpoint{2.466935in}{0.967664in}}%
\pgfpathlineto{\pgfqpoint{2.468504in}{0.944341in}}%
\pgfpathclose%
\pgfusepath{fill}%
\end{pgfscope}%
\begin{pgfscope}%
\pgfpathrectangle{\pgfqpoint{0.100000in}{0.100000in}}{\pgfqpoint{3.608454in}{2.310000in}}%
\pgfusepath{clip}%
\pgfsetbuttcap%
\pgfsetmiterjoin%
\definecolor{currentfill}{rgb}{0.000000,0.427451,0.786275}%
\pgfsetfillcolor{currentfill}%
\pgfsetlinewidth{0.000000pt}%
\definecolor{currentstroke}{rgb}{0.000000,0.000000,0.000000}%
\pgfsetstrokecolor{currentstroke}%
\pgfsetstrokeopacity{0.000000}%
\pgfsetdash{}{0pt}%
\pgfpathmoveto{\pgfqpoint{3.506058in}{1.889401in}}%
\pgfpathlineto{\pgfqpoint{3.496870in}{1.892724in}}%
\pgfpathlineto{\pgfqpoint{3.498018in}{1.895957in}}%
\pgfpathlineto{\pgfqpoint{3.491360in}{1.900167in}}%
\pgfpathlineto{\pgfqpoint{3.485597in}{1.897568in}}%
\pgfpathlineto{\pgfqpoint{3.484176in}{1.904080in}}%
\pgfpathlineto{\pgfqpoint{3.474707in}{1.903468in}}%
\pgfpathlineto{\pgfqpoint{3.465497in}{1.899405in}}%
\pgfpathlineto{\pgfqpoint{3.454146in}{1.936802in}}%
\pgfpathlineto{\pgfqpoint{3.442334in}{1.974192in}}%
\pgfpathlineto{\pgfqpoint{3.428979in}{2.013989in}}%
\pgfpathlineto{\pgfqpoint{3.435281in}{2.018401in}}%
\pgfpathlineto{\pgfqpoint{3.440288in}{2.012007in}}%
\pgfpathlineto{\pgfqpoint{3.443602in}{2.017361in}}%
\pgfpathlineto{\pgfqpoint{3.442691in}{2.028080in}}%
\pgfpathlineto{\pgfqpoint{3.447268in}{2.025988in}}%
\pgfpathlineto{\pgfqpoint{3.451169in}{2.029587in}}%
\pgfpathlineto{\pgfqpoint{3.444152in}{2.035398in}}%
\pgfpathlineto{\pgfqpoint{3.446349in}{2.043593in}}%
\pgfpathlineto{\pgfqpoint{3.457409in}{2.056674in}}%
\pgfpathlineto{\pgfqpoint{3.455828in}{2.064013in}}%
\pgfpathlineto{\pgfqpoint{3.460024in}{2.069258in}}%
\pgfpathlineto{\pgfqpoint{3.460838in}{2.075413in}}%
\pgfpathlineto{\pgfqpoint{3.455091in}{2.080073in}}%
\pgfpathlineto{\pgfqpoint{3.456695in}{2.090917in}}%
\pgfpathlineto{\pgfqpoint{3.452786in}{2.093614in}}%
\pgfpathlineto{\pgfqpoint{3.454249in}{2.105442in}}%
\pgfpathlineto{\pgfqpoint{3.460446in}{2.114346in}}%
\pgfpathlineto{\pgfqpoint{3.458788in}{2.127000in}}%
\pgfpathlineto{\pgfqpoint{3.474712in}{2.131510in}}%
\pgfpathlineto{\pgfqpoint{3.478067in}{2.117502in}}%
\pgfpathlineto{\pgfqpoint{3.487328in}{2.085547in}}%
\pgfpathlineto{\pgfqpoint{3.491275in}{2.076574in}}%
\pgfpathlineto{\pgfqpoint{3.490714in}{2.062918in}}%
\pgfpathlineto{\pgfqpoint{3.496496in}{2.059377in}}%
\pgfpathlineto{\pgfqpoint{3.494071in}{2.051630in}}%
\pgfpathlineto{\pgfqpoint{3.514278in}{2.013412in}}%
\pgfpathlineto{\pgfqpoint{3.527215in}{2.022326in}}%
\pgfpathlineto{\pgfqpoint{3.539818in}{1.996897in}}%
\pgfpathlineto{\pgfqpoint{3.532167in}{1.996604in}}%
\pgfpathlineto{\pgfqpoint{3.534775in}{1.988683in}}%
\pgfpathlineto{\pgfqpoint{3.535820in}{1.974394in}}%
\pgfpathlineto{\pgfqpoint{3.535140in}{1.964201in}}%
\pgfpathlineto{\pgfqpoint{3.538805in}{1.962876in}}%
\pgfpathlineto{\pgfqpoint{3.543555in}{1.955713in}}%
\pgfpathlineto{\pgfqpoint{3.547311in}{1.956692in}}%
\pgfpathlineto{\pgfqpoint{3.553924in}{1.946267in}}%
\pgfpathlineto{\pgfqpoint{3.550714in}{1.937007in}}%
\pgfpathlineto{\pgfqpoint{3.547033in}{1.936458in}}%
\pgfpathlineto{\pgfqpoint{3.546482in}{1.926493in}}%
\pgfpathlineto{\pgfqpoint{3.535703in}{1.920334in}}%
\pgfpathlineto{\pgfqpoint{3.534745in}{1.914657in}}%
\pgfpathlineto{\pgfqpoint{3.529959in}{1.911075in}}%
\pgfpathlineto{\pgfqpoint{3.516703in}{1.917002in}}%
\pgfpathlineto{\pgfqpoint{3.510197in}{1.910231in}}%
\pgfpathlineto{\pgfqpoint{3.508494in}{1.902315in}}%
\pgfpathlineto{\pgfqpoint{3.514068in}{1.894022in}}%
\pgfpathlineto{\pgfqpoint{3.506058in}{1.889401in}}%
\pgfpathclose%
\pgfusepath{fill}%
\end{pgfscope}%
\begin{pgfscope}%
\pgfpathrectangle{\pgfqpoint{0.100000in}{0.100000in}}{\pgfqpoint{3.608454in}{2.310000in}}%
\pgfusepath{clip}%
\pgfsetbuttcap%
\pgfsetmiterjoin%
\definecolor{currentfill}{rgb}{0.000000,0.615686,0.692157}%
\pgfsetfillcolor{currentfill}%
\pgfsetlinewidth{0.000000pt}%
\definecolor{currentstroke}{rgb}{0.000000,0.000000,0.000000}%
\pgfsetstrokecolor{currentstroke}%
\pgfsetstrokeopacity{0.000000}%
\pgfsetdash{}{0pt}%
\pgfpathmoveto{\pgfqpoint{2.519263in}{0.872679in}}%
\pgfpathlineto{\pgfqpoint{2.485523in}{0.870923in}}%
\pgfpathlineto{\pgfqpoint{2.483949in}{0.899006in}}%
\pgfpathlineto{\pgfqpoint{2.474841in}{0.898468in}}%
\pgfpathlineto{\pgfqpoint{2.474111in}{0.910147in}}%
\pgfpathlineto{\pgfqpoint{2.470515in}{0.912039in}}%
\pgfpathlineto{\pgfqpoint{2.469213in}{0.932802in}}%
\pgfpathlineto{\pgfqpoint{2.489985in}{0.934122in}}%
\pgfpathlineto{\pgfqpoint{2.506964in}{0.937791in}}%
\pgfpathlineto{\pgfqpoint{2.506684in}{0.942391in}}%
\pgfpathlineto{\pgfqpoint{2.520583in}{0.943229in}}%
\pgfpathlineto{\pgfqpoint{2.521327in}{0.932876in}}%
\pgfpathlineto{\pgfqpoint{2.535053in}{0.931443in}}%
\pgfpathlineto{\pgfqpoint{2.537884in}{0.927921in}}%
\pgfpathlineto{\pgfqpoint{2.534527in}{0.922423in}}%
\pgfpathlineto{\pgfqpoint{2.525476in}{0.919309in}}%
\pgfpathlineto{\pgfqpoint{2.526832in}{0.901949in}}%
\pgfpathlineto{\pgfqpoint{2.517297in}{0.901251in}}%
\pgfpathlineto{\pgfqpoint{2.519263in}{0.872679in}}%
\pgfpathclose%
\pgfusepath{fill}%
\end{pgfscope}%
\begin{pgfscope}%
\pgfpathrectangle{\pgfqpoint{0.100000in}{0.100000in}}{\pgfqpoint{3.608454in}{2.310000in}}%
\pgfusepath{clip}%
\pgfsetbuttcap%
\pgfsetmiterjoin%
\definecolor{currentfill}{rgb}{0.000000,0.403922,0.798039}%
\pgfsetfillcolor{currentfill}%
\pgfsetlinewidth{0.000000pt}%
\definecolor{currentstroke}{rgb}{0.000000,0.000000,0.000000}%
\pgfsetstrokecolor{currentstroke}%
\pgfsetstrokeopacity{0.000000}%
\pgfsetdash{}{0pt}%
\pgfpathmoveto{\pgfqpoint{0.623427in}{2.025636in}}%
\pgfpathlineto{\pgfqpoint{0.587828in}{2.036379in}}%
\pgfpathlineto{\pgfqpoint{0.581481in}{2.042623in}}%
\pgfpathlineto{\pgfqpoint{0.580214in}{2.049876in}}%
\pgfpathlineto{\pgfqpoint{0.575268in}{2.054730in}}%
\pgfpathlineto{\pgfqpoint{0.572541in}{2.061484in}}%
\pgfpathlineto{\pgfqpoint{0.575705in}{2.064800in}}%
\pgfpathlineto{\pgfqpoint{0.576179in}{2.072349in}}%
\pgfpathlineto{\pgfqpoint{0.568024in}{2.080965in}}%
\pgfpathlineto{\pgfqpoint{0.566154in}{2.086439in}}%
\pgfpathlineto{\pgfqpoint{0.561348in}{2.090351in}}%
\pgfpathlineto{\pgfqpoint{0.543848in}{2.096029in}}%
\pgfpathlineto{\pgfqpoint{0.550017in}{2.103751in}}%
\pgfpathlineto{\pgfqpoint{0.551015in}{2.110628in}}%
\pgfpathlineto{\pgfqpoint{0.556089in}{2.115055in}}%
\pgfpathlineto{\pgfqpoint{0.557819in}{2.120432in}}%
\pgfpathlineto{\pgfqpoint{0.566644in}{2.148266in}}%
\pgfpathlineto{\pgfqpoint{0.578043in}{2.148170in}}%
\pgfpathlineto{\pgfqpoint{0.589293in}{2.135786in}}%
\pgfpathlineto{\pgfqpoint{0.593066in}{2.119148in}}%
\pgfpathlineto{\pgfqpoint{0.609486in}{2.121452in}}%
\pgfpathlineto{\pgfqpoint{0.615173in}{2.117075in}}%
\pgfpathlineto{\pgfqpoint{0.623310in}{2.122649in}}%
\pgfpathlineto{\pgfqpoint{0.631314in}{2.147810in}}%
\pgfpathlineto{\pgfqpoint{0.669147in}{2.136525in}}%
\pgfpathlineto{\pgfqpoint{0.661241in}{2.110265in}}%
\pgfpathlineto{\pgfqpoint{0.656547in}{2.111660in}}%
\pgfpathlineto{\pgfqpoint{0.650732in}{2.091686in}}%
\pgfpathlineto{\pgfqpoint{0.653479in}{2.086201in}}%
\pgfpathlineto{\pgfqpoint{0.641984in}{2.086967in}}%
\pgfpathlineto{\pgfqpoint{0.634943in}{2.089043in}}%
\pgfpathlineto{\pgfqpoint{0.630988in}{2.086576in}}%
\pgfpathlineto{\pgfqpoint{0.628826in}{2.076535in}}%
\pgfpathlineto{\pgfqpoint{0.636428in}{2.062774in}}%
\pgfpathlineto{\pgfqpoint{0.634662in}{2.052316in}}%
\pgfpathlineto{\pgfqpoint{0.629691in}{2.048489in}}%
\pgfpathlineto{\pgfqpoint{0.630035in}{2.038493in}}%
\pgfpathlineto{\pgfqpoint{0.622394in}{2.035924in}}%
\pgfpathlineto{\pgfqpoint{0.623427in}{2.025636in}}%
\pgfpathclose%
\pgfusepath{fill}%
\end{pgfscope}%
\begin{pgfscope}%
\pgfpathrectangle{\pgfqpoint{0.100000in}{0.100000in}}{\pgfqpoint{3.608454in}{2.310000in}}%
\pgfusepath{clip}%
\pgfsetbuttcap%
\pgfsetmiterjoin%
\definecolor{currentfill}{rgb}{0.000000,0.466667,0.766667}%
\pgfsetfillcolor{currentfill}%
\pgfsetlinewidth{0.000000pt}%
\definecolor{currentstroke}{rgb}{0.000000,0.000000,0.000000}%
\pgfsetstrokecolor{currentstroke}%
\pgfsetstrokeopacity{0.000000}%
\pgfsetdash{}{0pt}%
\pgfpathmoveto{\pgfqpoint{3.027503in}{0.502091in}}%
\pgfpathlineto{\pgfqpoint{2.992126in}{0.496993in}}%
\pgfpathlineto{\pgfqpoint{2.998004in}{0.510943in}}%
\pgfpathlineto{\pgfqpoint{2.992287in}{0.512337in}}%
\pgfpathlineto{\pgfqpoint{2.985816in}{0.507630in}}%
\pgfpathlineto{\pgfqpoint{2.984712in}{0.500436in}}%
\pgfpathlineto{\pgfqpoint{2.980383in}{0.498839in}}%
\pgfpathlineto{\pgfqpoint{2.968704in}{0.512418in}}%
\pgfpathlineto{\pgfqpoint{2.971482in}{0.527436in}}%
\pgfpathlineto{\pgfqpoint{2.969072in}{0.537771in}}%
\pgfpathlineto{\pgfqpoint{2.971652in}{0.541641in}}%
\pgfpathlineto{\pgfqpoint{2.975315in}{0.570540in}}%
\pgfpathlineto{\pgfqpoint{2.974786in}{0.579836in}}%
\pgfpathlineto{\pgfqpoint{2.990145in}{0.581978in}}%
\pgfpathlineto{\pgfqpoint{3.001283in}{0.581332in}}%
\pgfpathlineto{\pgfqpoint{3.006224in}{0.574274in}}%
\pgfpathlineto{\pgfqpoint{3.017489in}{0.571785in}}%
\pgfpathlineto{\pgfqpoint{3.019784in}{0.555093in}}%
\pgfpathlineto{\pgfqpoint{3.016866in}{0.550384in}}%
\pgfpathlineto{\pgfqpoint{3.021367in}{0.543921in}}%
\pgfpathlineto{\pgfqpoint{3.027503in}{0.502091in}}%
\pgfpathclose%
\pgfusepath{fill}%
\end{pgfscope}%
\begin{pgfscope}%
\pgfpathrectangle{\pgfqpoint{0.100000in}{0.100000in}}{\pgfqpoint{3.608454in}{2.310000in}}%
\pgfusepath{clip}%
\pgfsetbuttcap%
\pgfsetmiterjoin%
\definecolor{currentfill}{rgb}{0.000000,0.290196,0.854902}%
\pgfsetfillcolor{currentfill}%
\pgfsetlinewidth{0.000000pt}%
\definecolor{currentstroke}{rgb}{0.000000,0.000000,0.000000}%
\pgfsetstrokecolor{currentstroke}%
\pgfsetstrokeopacity{0.000000}%
\pgfsetdash{}{0pt}%
\pgfpathmoveto{\pgfqpoint{3.420582in}{1.833115in}}%
\pgfpathlineto{\pgfqpoint{3.421042in}{1.838069in}}%
\pgfpathlineto{\pgfqpoint{3.413847in}{1.847236in}}%
\pgfpathlineto{\pgfqpoint{3.415970in}{1.857152in}}%
\pgfpathlineto{\pgfqpoint{3.421501in}{1.872886in}}%
\pgfpathlineto{\pgfqpoint{3.424616in}{1.868740in}}%
\pgfpathlineto{\pgfqpoint{3.432335in}{1.876511in}}%
\pgfpathlineto{\pgfqpoint{3.431173in}{1.881530in}}%
\pgfpathlineto{\pgfqpoint{3.437941in}{1.884217in}}%
\pgfpathlineto{\pgfqpoint{3.431626in}{1.898622in}}%
\pgfpathlineto{\pgfqpoint{3.442370in}{1.903310in}}%
\pgfpathlineto{\pgfqpoint{3.438972in}{1.915003in}}%
\pgfpathlineto{\pgfqpoint{3.433355in}{1.925260in}}%
\pgfpathlineto{\pgfqpoint{3.440922in}{1.921126in}}%
\pgfpathlineto{\pgfqpoint{3.441678in}{1.930119in}}%
\pgfpathlineto{\pgfqpoint{3.453336in}{1.932723in}}%
\pgfpathlineto{\pgfqpoint{3.454146in}{1.936802in}}%
\pgfpathlineto{\pgfqpoint{3.465497in}{1.899405in}}%
\pgfpathlineto{\pgfqpoint{3.474707in}{1.903468in}}%
\pgfpathlineto{\pgfqpoint{3.484176in}{1.904080in}}%
\pgfpathlineto{\pgfqpoint{3.485597in}{1.897568in}}%
\pgfpathlineto{\pgfqpoint{3.491360in}{1.900167in}}%
\pgfpathlineto{\pgfqpoint{3.498018in}{1.895957in}}%
\pgfpathlineto{\pgfqpoint{3.496870in}{1.892724in}}%
\pgfpathlineto{\pgfqpoint{3.506058in}{1.889401in}}%
\pgfpathlineto{\pgfqpoint{3.507318in}{1.881477in}}%
\pgfpathlineto{\pgfqpoint{3.505212in}{1.874640in}}%
\pgfpathlineto{\pgfqpoint{3.499936in}{1.871265in}}%
\pgfpathlineto{\pgfqpoint{3.499231in}{1.853431in}}%
\pgfpathlineto{\pgfqpoint{3.493789in}{1.834815in}}%
\pgfpathlineto{\pgfqpoint{3.494171in}{1.831598in}}%
\pgfpathlineto{\pgfqpoint{3.487553in}{1.830826in}}%
\pgfpathlineto{\pgfqpoint{3.477550in}{1.822939in}}%
\pgfpathlineto{\pgfqpoint{3.471016in}{1.810988in}}%
\pgfpathlineto{\pgfqpoint{3.434992in}{1.802755in}}%
\pgfpathlineto{\pgfqpoint{3.430894in}{1.808006in}}%
\pgfpathlineto{\pgfqpoint{3.426046in}{1.819395in}}%
\pgfpathlineto{\pgfqpoint{3.422865in}{1.818650in}}%
\pgfpathlineto{\pgfqpoint{3.420582in}{1.833115in}}%
\pgfpathclose%
\pgfusepath{fill}%
\end{pgfscope}%
\begin{pgfscope}%
\pgfpathrectangle{\pgfqpoint{0.100000in}{0.100000in}}{\pgfqpoint{3.608454in}{2.310000in}}%
\pgfusepath{clip}%
\pgfsetbuttcap%
\pgfsetmiterjoin%
\definecolor{currentfill}{rgb}{0.000000,0.698039,0.650980}%
\pgfsetfillcolor{currentfill}%
\pgfsetlinewidth{0.000000pt}%
\definecolor{currentstroke}{rgb}{0.000000,0.000000,0.000000}%
\pgfsetstrokecolor{currentstroke}%
\pgfsetstrokeopacity{0.000000}%
\pgfsetdash{}{0pt}%
\pgfpathmoveto{\pgfqpoint{2.478809in}{1.176580in}}%
\pgfpathlineto{\pgfqpoint{2.477580in}{1.181419in}}%
\pgfpathlineto{\pgfqpoint{2.479456in}{1.190579in}}%
\pgfpathlineto{\pgfqpoint{2.477226in}{1.192084in}}%
\pgfpathlineto{\pgfqpoint{2.473800in}{1.196498in}}%
\pgfpathlineto{\pgfqpoint{2.482512in}{1.210574in}}%
\pgfpathlineto{\pgfqpoint{2.488838in}{1.212212in}}%
\pgfpathlineto{\pgfqpoint{2.489480in}{1.220900in}}%
\pgfpathlineto{\pgfqpoint{2.501888in}{1.221865in}}%
\pgfpathlineto{\pgfqpoint{2.500587in}{1.242556in}}%
\pgfpathlineto{\pgfqpoint{2.519105in}{1.243921in}}%
\pgfpathlineto{\pgfqpoint{2.519923in}{1.230014in}}%
\pgfpathlineto{\pgfqpoint{2.523585in}{1.228784in}}%
\pgfpathlineto{\pgfqpoint{2.533129in}{1.219044in}}%
\pgfpathlineto{\pgfqpoint{2.535791in}{1.209919in}}%
\pgfpathlineto{\pgfqpoint{2.533377in}{1.204248in}}%
\pgfpathlineto{\pgfqpoint{2.539380in}{1.187318in}}%
\pgfpathlineto{\pgfqpoint{2.543041in}{1.178054in}}%
\pgfpathlineto{\pgfqpoint{2.519321in}{1.176572in}}%
\pgfpathlineto{\pgfqpoint{2.518313in}{1.191718in}}%
\pgfpathlineto{\pgfqpoint{2.497579in}{1.190610in}}%
\pgfpathlineto{\pgfqpoint{2.498496in}{1.176925in}}%
\pgfpathlineto{\pgfqpoint{2.478809in}{1.176580in}}%
\pgfpathclose%
\pgfusepath{fill}%
\end{pgfscope}%
\begin{pgfscope}%
\pgfpathrectangle{\pgfqpoint{0.100000in}{0.100000in}}{\pgfqpoint{3.608454in}{2.310000in}}%
\pgfusepath{clip}%
\pgfsetbuttcap%
\pgfsetmiterjoin%
\definecolor{currentfill}{rgb}{0.000000,0.505882,0.747059}%
\pgfsetfillcolor{currentfill}%
\pgfsetlinewidth{0.000000pt}%
\definecolor{currentstroke}{rgb}{0.000000,0.000000,0.000000}%
\pgfsetstrokecolor{currentstroke}%
\pgfsetstrokeopacity{0.000000}%
\pgfsetdash{}{0pt}%
\pgfpathmoveto{\pgfqpoint{3.165006in}{1.605498in}}%
\pgfpathlineto{\pgfqpoint{3.154199in}{1.598846in}}%
\pgfpathlineto{\pgfqpoint{3.144862in}{1.595009in}}%
\pgfpathlineto{\pgfqpoint{3.133554in}{1.604605in}}%
\pgfpathlineto{\pgfqpoint{3.128786in}{1.603472in}}%
\pgfpathlineto{\pgfqpoint{3.125203in}{1.608488in}}%
\pgfpathlineto{\pgfqpoint{3.111013in}{1.603891in}}%
\pgfpathlineto{\pgfqpoint{3.106858in}{1.608202in}}%
\pgfpathlineto{\pgfqpoint{3.110757in}{1.621055in}}%
\pgfpathlineto{\pgfqpoint{3.109153in}{1.629395in}}%
\pgfpathlineto{\pgfqpoint{3.131860in}{1.634049in}}%
\pgfpathlineto{\pgfqpoint{3.130816in}{1.638979in}}%
\pgfpathlineto{\pgfqpoint{3.167806in}{1.647018in}}%
\pgfpathlineto{\pgfqpoint{3.172073in}{1.651526in}}%
\pgfpathlineto{\pgfqpoint{3.162699in}{1.682336in}}%
\pgfpathlineto{\pgfqpoint{3.183986in}{1.686549in}}%
\pgfpathlineto{\pgfqpoint{3.207689in}{1.691369in}}%
\pgfpathlineto{\pgfqpoint{3.215154in}{1.664914in}}%
\pgfpathlineto{\pgfqpoint{3.210379in}{1.663470in}}%
\pgfpathlineto{\pgfqpoint{3.210875in}{1.652502in}}%
\pgfpathlineto{\pgfqpoint{3.209405in}{1.635945in}}%
\pgfpathlineto{\pgfqpoint{3.203745in}{1.634729in}}%
\pgfpathlineto{\pgfqpoint{3.195407in}{1.620855in}}%
\pgfpathlineto{\pgfqpoint{3.192630in}{1.620121in}}%
\pgfpathlineto{\pgfqpoint{3.183685in}{1.619914in}}%
\pgfpathlineto{\pgfqpoint{3.174106in}{1.615801in}}%
\pgfpathlineto{\pgfqpoint{3.173979in}{1.610853in}}%
\pgfpathlineto{\pgfqpoint{3.165006in}{1.605498in}}%
\pgfpathclose%
\pgfusepath{fill}%
\end{pgfscope}%
\begin{pgfscope}%
\pgfpathrectangle{\pgfqpoint{0.100000in}{0.100000in}}{\pgfqpoint{3.608454in}{2.310000in}}%
\pgfusepath{clip}%
\pgfsetbuttcap%
\pgfsetmiterjoin%
\definecolor{currentfill}{rgb}{0.000000,0.823529,0.588235}%
\pgfsetfillcolor{currentfill}%
\pgfsetlinewidth{0.000000pt}%
\definecolor{currentstroke}{rgb}{0.000000,0.000000,0.000000}%
\pgfsetstrokecolor{currentstroke}%
\pgfsetstrokeopacity{0.000000}%
\pgfsetdash{}{0pt}%
\pgfpathmoveto{\pgfqpoint{2.674593in}{0.730790in}}%
\pgfpathlineto{\pgfqpoint{2.612513in}{0.725646in}}%
\pgfpathlineto{\pgfqpoint{2.609757in}{0.745194in}}%
\pgfpathlineto{\pgfqpoint{2.602701in}{0.749237in}}%
\pgfpathlineto{\pgfqpoint{2.601460in}{0.754867in}}%
\pgfpathlineto{\pgfqpoint{2.605232in}{0.759125in}}%
\pgfpathlineto{\pgfqpoint{2.611445in}{0.760754in}}%
\pgfpathlineto{\pgfqpoint{2.610143in}{0.781433in}}%
\pgfpathlineto{\pgfqpoint{2.613565in}{0.781771in}}%
\pgfpathlineto{\pgfqpoint{2.613728in}{0.792291in}}%
\pgfpathlineto{\pgfqpoint{2.653773in}{0.795845in}}%
\pgfpathlineto{\pgfqpoint{2.655184in}{0.780136in}}%
\pgfpathlineto{\pgfqpoint{2.660506in}{0.772013in}}%
\pgfpathlineto{\pgfqpoint{2.669844in}{0.772749in}}%
\pgfpathlineto{\pgfqpoint{2.673739in}{0.731186in}}%
\pgfpathlineto{\pgfqpoint{2.674593in}{0.730790in}}%
\pgfpathclose%
\pgfusepath{fill}%
\end{pgfscope}%
\begin{pgfscope}%
\pgfpathrectangle{\pgfqpoint{0.100000in}{0.100000in}}{\pgfqpoint{3.608454in}{2.310000in}}%
\pgfusepath{clip}%
\pgfsetbuttcap%
\pgfsetmiterjoin%
\definecolor{currentfill}{rgb}{0.000000,0.674510,0.662745}%
\pgfsetfillcolor{currentfill}%
\pgfsetlinewidth{0.000000pt}%
\definecolor{currentstroke}{rgb}{0.000000,0.000000,0.000000}%
\pgfsetstrokecolor{currentstroke}%
\pgfsetstrokeopacity{0.000000}%
\pgfsetdash{}{0pt}%
\pgfpathmoveto{\pgfqpoint{1.245517in}{0.942241in}}%
\pgfpathlineto{\pgfqpoint{1.176800in}{0.953138in}}%
\pgfpathlineto{\pgfqpoint{1.183780in}{0.997802in}}%
\pgfpathlineto{\pgfqpoint{1.193614in}{1.060581in}}%
\pgfpathlineto{\pgfqpoint{1.236876in}{1.053946in}}%
\pgfpathlineto{\pgfqpoint{1.278910in}{1.047826in}}%
\pgfpathlineto{\pgfqpoint{1.333401in}{1.040515in}}%
\pgfpathlineto{\pgfqpoint{1.340304in}{1.035758in}}%
\pgfpathlineto{\pgfqpoint{1.362154in}{1.025644in}}%
\pgfpathlineto{\pgfqpoint{1.360437in}{1.011466in}}%
\pgfpathlineto{\pgfqpoint{1.392442in}{1.007538in}}%
\pgfpathlineto{\pgfqpoint{1.388260in}{0.973319in}}%
\pgfpathlineto{\pgfqpoint{1.379878in}{0.974364in}}%
\pgfpathlineto{\pgfqpoint{1.378188in}{0.960531in}}%
\pgfpathlineto{\pgfqpoint{1.357228in}{0.963939in}}%
\pgfpathlineto{\pgfqpoint{1.355496in}{0.949729in}}%
\pgfpathlineto{\pgfqpoint{1.318032in}{0.954222in}}%
\pgfpathlineto{\pgfqpoint{1.252648in}{0.963109in}}%
\pgfpathlineto{\pgfqpoint{1.248648in}{0.963976in}}%
\pgfpathlineto{\pgfqpoint{1.245517in}{0.942241in}}%
\pgfpathclose%
\pgfusepath{fill}%
\end{pgfscope}%
\begin{pgfscope}%
\pgfpathrectangle{\pgfqpoint{0.100000in}{0.100000in}}{\pgfqpoint{3.608454in}{2.310000in}}%
\pgfusepath{clip}%
\pgfsetbuttcap%
\pgfsetmiterjoin%
\definecolor{currentfill}{rgb}{0.000000,0.670588,0.664706}%
\pgfsetfillcolor{currentfill}%
\pgfsetlinewidth{0.000000pt}%
\definecolor{currentstroke}{rgb}{0.000000,0.000000,0.000000}%
\pgfsetstrokecolor{currentstroke}%
\pgfsetstrokeopacity{0.000000}%
\pgfsetdash{}{0pt}%
\pgfpathmoveto{\pgfqpoint{2.676133in}{1.724089in}}%
\pgfpathlineto{\pgfqpoint{2.648810in}{1.721346in}}%
\pgfpathlineto{\pgfqpoint{2.645951in}{1.748717in}}%
\pgfpathlineto{\pgfqpoint{2.618822in}{1.745974in}}%
\pgfpathlineto{\pgfqpoint{2.615872in}{1.773777in}}%
\pgfpathlineto{\pgfqpoint{2.628418in}{1.774736in}}%
\pgfpathlineto{\pgfqpoint{2.625838in}{1.802295in}}%
\pgfpathlineto{\pgfqpoint{2.680280in}{1.807943in}}%
\pgfpathlineto{\pgfqpoint{2.683319in}{1.780357in}}%
\pgfpathlineto{\pgfqpoint{2.669984in}{1.779116in}}%
\pgfpathlineto{\pgfqpoint{2.676133in}{1.724089in}}%
\pgfpathclose%
\pgfusepath{fill}%
\end{pgfscope}%
\begin{pgfscope}%
\pgfpathrectangle{\pgfqpoint{0.100000in}{0.100000in}}{\pgfqpoint{3.608454in}{2.310000in}}%
\pgfusepath{clip}%
\pgfsetbuttcap%
\pgfsetmiterjoin%
\definecolor{currentfill}{rgb}{0.000000,0.411765,0.794118}%
\pgfsetfillcolor{currentfill}%
\pgfsetlinewidth{0.000000pt}%
\definecolor{currentstroke}{rgb}{0.000000,0.000000,0.000000}%
\pgfsetstrokecolor{currentstroke}%
\pgfsetstrokeopacity{0.000000}%
\pgfsetdash{}{0pt}%
\pgfpathmoveto{\pgfqpoint{2.575892in}{1.080309in}}%
\pgfpathlineto{\pgfqpoint{2.579650in}{1.089158in}}%
\pgfpathlineto{\pgfqpoint{2.573068in}{1.094430in}}%
\pgfpathlineto{\pgfqpoint{2.573691in}{1.108364in}}%
\pgfpathlineto{\pgfqpoint{2.573891in}{1.113688in}}%
\pgfpathlineto{\pgfqpoint{2.584331in}{1.121514in}}%
\pgfpathlineto{\pgfqpoint{2.605651in}{1.120589in}}%
\pgfpathlineto{\pgfqpoint{2.605693in}{1.111968in}}%
\pgfpathlineto{\pgfqpoint{2.632392in}{1.103973in}}%
\pgfpathlineto{\pgfqpoint{2.643683in}{1.102257in}}%
\pgfpathlineto{\pgfqpoint{2.644297in}{1.082383in}}%
\pgfpathlineto{\pgfqpoint{2.648588in}{1.077186in}}%
\pgfpathlineto{\pgfqpoint{2.648663in}{1.074175in}}%
\pgfpathlineto{\pgfqpoint{2.639911in}{1.067474in}}%
\pgfpathlineto{\pgfqpoint{2.634590in}{1.067850in}}%
\pgfpathlineto{\pgfqpoint{2.636003in}{1.046348in}}%
\pgfpathlineto{\pgfqpoint{2.611745in}{1.044790in}}%
\pgfpathlineto{\pgfqpoint{2.610505in}{1.048649in}}%
\pgfpathlineto{\pgfqpoint{2.608880in}{1.074941in}}%
\pgfpathlineto{\pgfqpoint{2.607711in}{1.078660in}}%
\pgfpathlineto{\pgfqpoint{2.601253in}{1.080261in}}%
\pgfpathlineto{\pgfqpoint{2.585581in}{1.074349in}}%
\pgfpathlineto{\pgfqpoint{2.580566in}{1.079844in}}%
\pgfpathlineto{\pgfqpoint{2.575892in}{1.080309in}}%
\pgfpathclose%
\pgfusepath{fill}%
\end{pgfscope}%
\begin{pgfscope}%
\pgfpathrectangle{\pgfqpoint{0.100000in}{0.100000in}}{\pgfqpoint{3.608454in}{2.310000in}}%
\pgfusepath{clip}%
\pgfsetbuttcap%
\pgfsetmiterjoin%
\definecolor{currentfill}{rgb}{0.000000,0.949020,0.525490}%
\pgfsetfillcolor{currentfill}%
\pgfsetlinewidth{0.000000pt}%
\definecolor{currentstroke}{rgb}{0.000000,0.000000,0.000000}%
\pgfsetstrokecolor{currentstroke}%
\pgfsetstrokeopacity{0.000000}%
\pgfsetdash{}{0pt}%
\pgfpathmoveto{\pgfqpoint{3.010884in}{1.388904in}}%
\pgfpathlineto{\pgfqpoint{3.018910in}{1.376644in}}%
\pgfpathlineto{\pgfqpoint{3.016902in}{1.365346in}}%
\pgfpathlineto{\pgfqpoint{3.012394in}{1.360390in}}%
\pgfpathlineto{\pgfqpoint{3.011961in}{1.351516in}}%
\pgfpathlineto{\pgfqpoint{3.006750in}{1.353994in}}%
\pgfpathlineto{\pgfqpoint{2.992200in}{1.336598in}}%
\pgfpathlineto{\pgfqpoint{2.981470in}{1.336849in}}%
\pgfpathlineto{\pgfqpoint{2.973652in}{1.345179in}}%
\pgfpathlineto{\pgfqpoint{2.968495in}{1.342247in}}%
\pgfpathlineto{\pgfqpoint{2.958041in}{1.346022in}}%
\pgfpathlineto{\pgfqpoint{2.975847in}{1.360884in}}%
\pgfpathlineto{\pgfqpoint{2.976752in}{1.368166in}}%
\pgfpathlineto{\pgfqpoint{2.972924in}{1.373592in}}%
\pgfpathlineto{\pgfqpoint{2.965422in}{1.379505in}}%
\pgfpathlineto{\pgfqpoint{2.967663in}{1.383858in}}%
\pgfpathlineto{\pgfqpoint{2.962379in}{1.389071in}}%
\pgfpathlineto{\pgfqpoint{2.965023in}{1.393727in}}%
\pgfpathlineto{\pgfqpoint{2.961168in}{1.406181in}}%
\pgfpathlineto{\pgfqpoint{2.967704in}{1.409901in}}%
\pgfpathlineto{\pgfqpoint{2.973331in}{1.416066in}}%
\pgfpathlineto{\pgfqpoint{2.978628in}{1.415804in}}%
\pgfpathlineto{\pgfqpoint{2.980404in}{1.409102in}}%
\pgfpathlineto{\pgfqpoint{2.989078in}{1.401477in}}%
\pgfpathlineto{\pgfqpoint{2.997346in}{1.397787in}}%
\pgfpathlineto{\pgfqpoint{3.001117in}{1.390072in}}%
\pgfpathlineto{\pgfqpoint{3.010884in}{1.388904in}}%
\pgfpathclose%
\pgfusepath{fill}%
\end{pgfscope}%
\begin{pgfscope}%
\pgfpathrectangle{\pgfqpoint{0.100000in}{0.100000in}}{\pgfqpoint{3.608454in}{2.310000in}}%
\pgfusepath{clip}%
\pgfsetbuttcap%
\pgfsetmiterjoin%
\definecolor{currentfill}{rgb}{0.000000,0.784314,0.607843}%
\pgfsetfillcolor{currentfill}%
\pgfsetlinewidth{0.000000pt}%
\definecolor{currentstroke}{rgb}{0.000000,0.000000,0.000000}%
\pgfsetstrokecolor{currentstroke}%
\pgfsetstrokeopacity{0.000000}%
\pgfsetdash{}{0pt}%
\pgfpathmoveto{\pgfqpoint{0.983512in}{1.292540in}}%
\pgfpathlineto{\pgfqpoint{0.912295in}{1.307222in}}%
\pgfpathlineto{\pgfqpoint{0.925741in}{1.369767in}}%
\pgfpathlineto{\pgfqpoint{0.938194in}{1.428823in}}%
\pgfpathlineto{\pgfqpoint{1.000345in}{1.416064in}}%
\pgfpathlineto{\pgfqpoint{1.031123in}{1.409823in}}%
\pgfpathlineto{\pgfqpoint{1.030622in}{1.400197in}}%
\pgfpathlineto{\pgfqpoint{1.034771in}{1.396805in}}%
\pgfpathlineto{\pgfqpoint{1.038444in}{1.389127in}}%
\pgfpathlineto{\pgfqpoint{1.032898in}{1.379983in}}%
\pgfpathlineto{\pgfqpoint{1.026917in}{1.376417in}}%
\pgfpathlineto{\pgfqpoint{1.023679in}{1.356359in}}%
\pgfpathlineto{\pgfqpoint{1.016803in}{1.357032in}}%
\pgfpathlineto{\pgfqpoint{1.015614in}{1.351012in}}%
\pgfpathlineto{\pgfqpoint{1.008850in}{1.352444in}}%
\pgfpathlineto{\pgfqpoint{1.005205in}{1.331999in}}%
\pgfpathlineto{\pgfqpoint{0.991804in}{1.334673in}}%
\pgfpathlineto{\pgfqpoint{0.983512in}{1.292540in}}%
\pgfpathclose%
\pgfusepath{fill}%
\end{pgfscope}%
\begin{pgfscope}%
\pgfpathrectangle{\pgfqpoint{0.100000in}{0.100000in}}{\pgfqpoint{3.608454in}{2.310000in}}%
\pgfusepath{clip}%
\pgfsetbuttcap%
\pgfsetmiterjoin%
\definecolor{currentfill}{rgb}{0.000000,0.490196,0.754902}%
\pgfsetfillcolor{currentfill}%
\pgfsetlinewidth{0.000000pt}%
\definecolor{currentstroke}{rgb}{0.000000,0.000000,0.000000}%
\pgfsetstrokecolor{currentstroke}%
\pgfsetstrokeopacity{0.000000}%
\pgfsetdash{}{0pt}%
\pgfpathmoveto{\pgfqpoint{2.684273in}{1.275179in}}%
\pgfpathlineto{\pgfqpoint{2.691418in}{1.266749in}}%
\pgfpathlineto{\pgfqpoint{2.694291in}{1.266816in}}%
\pgfpathlineto{\pgfqpoint{2.694535in}{1.255275in}}%
\pgfpathlineto{\pgfqpoint{2.701382in}{1.255633in}}%
\pgfpathlineto{\pgfqpoint{2.704025in}{1.249697in}}%
\pgfpathlineto{\pgfqpoint{2.693473in}{1.249077in}}%
\pgfpathlineto{\pgfqpoint{2.692446in}{1.245040in}}%
\pgfpathlineto{\pgfqpoint{2.667849in}{1.244767in}}%
\pgfpathlineto{\pgfqpoint{2.665870in}{1.236725in}}%
\pgfpathlineto{\pgfqpoint{2.661850in}{1.235103in}}%
\pgfpathlineto{\pgfqpoint{2.642473in}{1.232080in}}%
\pgfpathlineto{\pgfqpoint{2.641029in}{1.235231in}}%
\pgfpathlineto{\pgfqpoint{2.632971in}{1.237093in}}%
\pgfpathlineto{\pgfqpoint{2.628815in}{1.247944in}}%
\pgfpathlineto{\pgfqpoint{2.630443in}{1.257485in}}%
\pgfpathlineto{\pgfqpoint{2.625972in}{1.268011in}}%
\pgfpathlineto{\pgfqpoint{2.627023in}{1.272051in}}%
\pgfpathlineto{\pgfqpoint{2.630595in}{1.278625in}}%
\pgfpathlineto{\pgfqpoint{2.635552in}{1.279697in}}%
\pgfpathlineto{\pgfqpoint{2.634127in}{1.288599in}}%
\pgfpathlineto{\pgfqpoint{2.636960in}{1.295013in}}%
\pgfpathlineto{\pgfqpoint{2.640611in}{1.294225in}}%
\pgfpathlineto{\pgfqpoint{2.643610in}{1.301513in}}%
\pgfpathlineto{\pgfqpoint{2.648711in}{1.297118in}}%
\pgfpathlineto{\pgfqpoint{2.649917in}{1.290900in}}%
\pgfpathlineto{\pgfqpoint{2.655765in}{1.288068in}}%
\pgfpathlineto{\pgfqpoint{2.669652in}{1.289191in}}%
\pgfpathlineto{\pgfqpoint{2.675593in}{1.285078in}}%
\pgfpathlineto{\pgfqpoint{2.678952in}{1.278333in}}%
\pgfpathlineto{\pgfqpoint{2.684273in}{1.275179in}}%
\pgfpathclose%
\pgfusepath{fill}%
\end{pgfscope}%
\begin{pgfscope}%
\pgfpathrectangle{\pgfqpoint{0.100000in}{0.100000in}}{\pgfqpoint{3.608454in}{2.310000in}}%
\pgfusepath{clip}%
\pgfsetbuttcap%
\pgfsetmiterjoin%
\definecolor{currentfill}{rgb}{0.000000,0.478431,0.760784}%
\pgfsetfillcolor{currentfill}%
\pgfsetlinewidth{0.000000pt}%
\definecolor{currentstroke}{rgb}{0.000000,0.000000,0.000000}%
\pgfsetstrokecolor{currentstroke}%
\pgfsetstrokeopacity{0.000000}%
\pgfsetdash{}{0pt}%
\pgfpathmoveto{\pgfqpoint{1.896146in}{2.138772in}}%
\pgfpathlineto{\pgfqpoint{1.895146in}{2.116164in}}%
\pgfpathlineto{\pgfqpoint{1.896095in}{2.102164in}}%
\pgfpathlineto{\pgfqpoint{1.895459in}{2.088252in}}%
\pgfpathlineto{\pgfqpoint{1.867703in}{2.089585in}}%
\pgfpathlineto{\pgfqpoint{1.868375in}{2.103508in}}%
\pgfpathlineto{\pgfqpoint{1.833772in}{2.105376in}}%
\pgfpathlineto{\pgfqpoint{1.832606in}{2.119399in}}%
\pgfpathlineto{\pgfqpoint{1.833904in}{2.141855in}}%
\pgfpathlineto{\pgfqpoint{1.896146in}{2.138772in}}%
\pgfpathclose%
\pgfusepath{fill}%
\end{pgfscope}%
\begin{pgfscope}%
\pgfpathrectangle{\pgfqpoint{0.100000in}{0.100000in}}{\pgfqpoint{3.608454in}{2.310000in}}%
\pgfusepath{clip}%
\pgfsetbuttcap%
\pgfsetmiterjoin%
\definecolor{currentfill}{rgb}{0.000000,0.490196,0.754902}%
\pgfsetfillcolor{currentfill}%
\pgfsetlinewidth{0.000000pt}%
\definecolor{currentstroke}{rgb}{0.000000,0.000000,0.000000}%
\pgfsetstrokecolor{currentstroke}%
\pgfsetstrokeopacity{0.000000}%
\pgfsetdash{}{0pt}%
\pgfpathmoveto{\pgfqpoint{2.013608in}{1.046935in}}%
\pgfpathlineto{\pgfqpoint{2.013392in}{1.033129in}}%
\pgfpathlineto{\pgfqpoint{2.010212in}{1.033190in}}%
\pgfpathlineto{\pgfqpoint{2.009922in}{1.016922in}}%
\pgfpathlineto{\pgfqpoint{1.999760in}{1.016432in}}%
\pgfpathlineto{\pgfqpoint{1.994760in}{1.012926in}}%
\pgfpathlineto{\pgfqpoint{1.988856in}{1.020093in}}%
\pgfpathlineto{\pgfqpoint{1.984031in}{1.018827in}}%
\pgfpathlineto{\pgfqpoint{1.981120in}{1.013018in}}%
\pgfpathlineto{\pgfqpoint{1.933314in}{1.014281in}}%
\pgfpathlineto{\pgfqpoint{1.934200in}{1.052365in}}%
\pgfpathlineto{\pgfqpoint{1.923203in}{1.052514in}}%
\pgfpathlineto{\pgfqpoint{1.911760in}{1.056365in}}%
\pgfpathlineto{\pgfqpoint{1.893326in}{1.056979in}}%
\pgfpathlineto{\pgfqpoint{1.893386in}{1.070793in}}%
\pgfpathlineto{\pgfqpoint{1.893880in}{1.084537in}}%
\pgfpathlineto{\pgfqpoint{1.900701in}{1.084287in}}%
\pgfpathlineto{\pgfqpoint{1.901714in}{1.119083in}}%
\pgfpathlineto{\pgfqpoint{1.949595in}{1.117555in}}%
\pgfpathlineto{\pgfqpoint{1.956419in}{1.117360in}}%
\pgfpathlineto{\pgfqpoint{1.956052in}{1.103169in}}%
\pgfpathlineto{\pgfqpoint{1.969632in}{1.099419in}}%
\pgfpathlineto{\pgfqpoint{2.002912in}{1.098705in}}%
\pgfpathlineto{\pgfqpoint{2.002513in}{1.074759in}}%
\pgfpathlineto{\pgfqpoint{2.001980in}{1.055914in}}%
\pgfpathlineto{\pgfqpoint{2.010159in}{1.057965in}}%
\pgfpathlineto{\pgfqpoint{2.013608in}{1.046935in}}%
\pgfpathclose%
\pgfusepath{fill}%
\end{pgfscope}%
\begin{pgfscope}%
\pgfpathrectangle{\pgfqpoint{0.100000in}{0.100000in}}{\pgfqpoint{3.608454in}{2.310000in}}%
\pgfusepath{clip}%
\pgfsetbuttcap%
\pgfsetmiterjoin%
\definecolor{currentfill}{rgb}{0.000000,0.482353,0.758824}%
\pgfsetfillcolor{currentfill}%
\pgfsetlinewidth{0.000000pt}%
\definecolor{currentstroke}{rgb}{0.000000,0.000000,0.000000}%
\pgfsetstrokecolor{currentstroke}%
\pgfsetstrokeopacity{0.000000}%
\pgfsetdash{}{0pt}%
\pgfpathmoveto{\pgfqpoint{2.111687in}{1.695719in}}%
\pgfpathlineto{\pgfqpoint{2.115164in}{1.695734in}}%
\pgfpathlineto{\pgfqpoint{2.115120in}{1.723283in}}%
\pgfpathlineto{\pgfqpoint{2.177278in}{1.723751in}}%
\pgfpathlineto{\pgfqpoint{2.184096in}{1.723863in}}%
\pgfpathlineto{\pgfqpoint{2.184501in}{1.696223in}}%
\pgfpathlineto{\pgfqpoint{2.165979in}{1.695985in}}%
\pgfpathlineto{\pgfqpoint{2.138833in}{1.695825in}}%
\pgfpathlineto{\pgfqpoint{2.138939in}{1.676359in}}%
\pgfpathlineto{\pgfqpoint{2.111765in}{1.676267in}}%
\pgfpathlineto{\pgfqpoint{2.111687in}{1.695719in}}%
\pgfpathclose%
\pgfusepath{fill}%
\end{pgfscope}%
\begin{pgfscope}%
\pgfpathrectangle{\pgfqpoint{0.100000in}{0.100000in}}{\pgfqpoint{3.608454in}{2.310000in}}%
\pgfusepath{clip}%
\pgfsetbuttcap%
\pgfsetmiterjoin%
\definecolor{currentfill}{rgb}{0.000000,0.823529,0.588235}%
\pgfsetfillcolor{currentfill}%
\pgfsetlinewidth{0.000000pt}%
\definecolor{currentstroke}{rgb}{0.000000,0.000000,0.000000}%
\pgfsetstrokecolor{currentstroke}%
\pgfsetstrokeopacity{0.000000}%
\pgfsetdash{}{0pt}%
\pgfpathmoveto{\pgfqpoint{2.804581in}{1.191394in}}%
\pgfpathlineto{\pgfqpoint{2.789452in}{1.189777in}}%
\pgfpathlineto{\pgfqpoint{2.783865in}{1.195392in}}%
\pgfpathlineto{\pgfqpoint{2.779295in}{1.206052in}}%
\pgfpathlineto{\pgfqpoint{2.781636in}{1.214245in}}%
\pgfpathlineto{\pgfqpoint{2.778898in}{1.226398in}}%
\pgfpathlineto{\pgfqpoint{2.780361in}{1.233265in}}%
\pgfpathlineto{\pgfqpoint{2.781851in}{1.240501in}}%
\pgfpathlineto{\pgfqpoint{2.788710in}{1.248796in}}%
\pgfpathlineto{\pgfqpoint{2.797121in}{1.243378in}}%
\pgfpathlineto{\pgfqpoint{2.801049in}{1.243631in}}%
\pgfpathlineto{\pgfqpoint{2.809988in}{1.252552in}}%
\pgfpathlineto{\pgfqpoint{2.816644in}{1.252449in}}%
\pgfpathlineto{\pgfqpoint{2.826981in}{1.247494in}}%
\pgfpathlineto{\pgfqpoint{2.826345in}{1.240166in}}%
\pgfpathlineto{\pgfqpoint{2.831285in}{1.222519in}}%
\pgfpathlineto{\pgfqpoint{2.815210in}{1.209500in}}%
\pgfpathlineto{\pgfqpoint{2.812724in}{1.202234in}}%
\pgfpathlineto{\pgfqpoint{2.806029in}{1.196000in}}%
\pgfpathlineto{\pgfqpoint{2.804581in}{1.191394in}}%
\pgfpathclose%
\pgfusepath{fill}%
\end{pgfscope}%
\begin{pgfscope}%
\pgfpathrectangle{\pgfqpoint{0.100000in}{0.100000in}}{\pgfqpoint{3.608454in}{2.310000in}}%
\pgfusepath{clip}%
\pgfsetbuttcap%
\pgfsetmiterjoin%
\definecolor{currentfill}{rgb}{0.000000,0.615686,0.692157}%
\pgfsetfillcolor{currentfill}%
\pgfsetlinewidth{0.000000pt}%
\definecolor{currentstroke}{rgb}{0.000000,0.000000,0.000000}%
\pgfsetstrokecolor{currentstroke}%
\pgfsetstrokeopacity{0.000000}%
\pgfsetdash{}{0pt}%
\pgfpathmoveto{\pgfqpoint{2.765354in}{1.245580in}}%
\pgfpathlineto{\pgfqpoint{2.761764in}{1.247341in}}%
\pgfpathlineto{\pgfqpoint{2.752556in}{1.236924in}}%
\pgfpathlineto{\pgfqpoint{2.751350in}{1.246500in}}%
\pgfpathlineto{\pgfqpoint{2.742755in}{1.250848in}}%
\pgfpathlineto{\pgfqpoint{2.741472in}{1.258776in}}%
\pgfpathlineto{\pgfqpoint{2.729948in}{1.258824in}}%
\pgfpathlineto{\pgfqpoint{2.729699in}{1.283399in}}%
\pgfpathlineto{\pgfqpoint{2.727518in}{1.286122in}}%
\pgfpathlineto{\pgfqpoint{2.733018in}{1.291688in}}%
\pgfpathlineto{\pgfqpoint{2.740147in}{1.292411in}}%
\pgfpathlineto{\pgfqpoint{2.757404in}{1.276520in}}%
\pgfpathlineto{\pgfqpoint{2.759930in}{1.280063in}}%
\pgfpathlineto{\pgfqpoint{2.773095in}{1.263245in}}%
\pgfpathlineto{\pgfqpoint{2.767417in}{1.258344in}}%
\pgfpathlineto{\pgfqpoint{2.765354in}{1.245580in}}%
\pgfpathclose%
\pgfusepath{fill}%
\end{pgfscope}%
\begin{pgfscope}%
\pgfpathrectangle{\pgfqpoint{0.100000in}{0.100000in}}{\pgfqpoint{3.608454in}{2.310000in}}%
\pgfusepath{clip}%
\pgfsetbuttcap%
\pgfsetmiterjoin%
\definecolor{currentfill}{rgb}{0.000000,0.403922,0.798039}%
\pgfsetfillcolor{currentfill}%
\pgfsetlinewidth{0.000000pt}%
\definecolor{currentstroke}{rgb}{0.000000,0.000000,0.000000}%
\pgfsetstrokecolor{currentstroke}%
\pgfsetstrokeopacity{0.000000}%
\pgfsetdash{}{0pt}%
\pgfpathmoveto{\pgfqpoint{1.682362in}{1.462771in}}%
\pgfpathlineto{\pgfqpoint{1.677720in}{1.401672in}}%
\pgfpathlineto{\pgfqpoint{1.610603in}{1.406638in}}%
\pgfpathlineto{\pgfqpoint{1.577074in}{1.409804in}}%
\pgfpathlineto{\pgfqpoint{1.580447in}{1.444028in}}%
\pgfpathlineto{\pgfqpoint{1.594536in}{1.442692in}}%
\pgfpathlineto{\pgfqpoint{1.598652in}{1.483744in}}%
\pgfpathlineto{\pgfqpoint{1.591803in}{1.484401in}}%
\pgfpathlineto{\pgfqpoint{1.595804in}{1.522015in}}%
\pgfpathlineto{\pgfqpoint{1.650561in}{1.517162in}}%
\pgfpathlineto{\pgfqpoint{1.649013in}{1.497206in}}%
\pgfpathlineto{\pgfqpoint{1.684833in}{1.494331in}}%
\pgfpathlineto{\pgfqpoint{1.684512in}{1.490228in}}%
\pgfpathlineto{\pgfqpoint{1.682362in}{1.462771in}}%
\pgfpathclose%
\pgfusepath{fill}%
\end{pgfscope}%
\begin{pgfscope}%
\pgfpathrectangle{\pgfqpoint{0.100000in}{0.100000in}}{\pgfqpoint{3.608454in}{2.310000in}}%
\pgfusepath{clip}%
\pgfsetbuttcap%
\pgfsetmiterjoin%
\definecolor{currentfill}{rgb}{0.000000,0.470588,0.764706}%
\pgfsetfillcolor{currentfill}%
\pgfsetlinewidth{0.000000pt}%
\definecolor{currentstroke}{rgb}{0.000000,0.000000,0.000000}%
\pgfsetstrokecolor{currentstroke}%
\pgfsetstrokeopacity{0.000000}%
\pgfsetdash{}{0pt}%
\pgfpathmoveto{\pgfqpoint{1.846276in}{0.936427in}}%
\pgfpathlineto{\pgfqpoint{1.844714in}{0.901800in}}%
\pgfpathlineto{\pgfqpoint{1.810377in}{0.903512in}}%
\pgfpathlineto{\pgfqpoint{1.811548in}{0.930044in}}%
\pgfpathlineto{\pgfqpoint{1.811868in}{0.938019in}}%
\pgfpathlineto{\pgfqpoint{1.846276in}{0.936427in}}%
\pgfpathclose%
\pgfusepath{fill}%
\end{pgfscope}%
\begin{pgfscope}%
\pgfpathrectangle{\pgfqpoint{0.100000in}{0.100000in}}{\pgfqpoint{3.608454in}{2.310000in}}%
\pgfusepath{clip}%
\pgfsetbuttcap%
\pgfsetmiterjoin%
\definecolor{currentfill}{rgb}{0.000000,0.588235,0.705882}%
\pgfsetfillcolor{currentfill}%
\pgfsetlinewidth{0.000000pt}%
\definecolor{currentstroke}{rgb}{0.000000,0.000000,0.000000}%
\pgfsetstrokecolor{currentstroke}%
\pgfsetstrokeopacity{0.000000}%
\pgfsetdash{}{0pt}%
\pgfpathmoveto{\pgfqpoint{2.238456in}{1.359055in}}%
\pgfpathlineto{\pgfqpoint{2.232957in}{1.349258in}}%
\pgfpathlineto{\pgfqpoint{2.233532in}{1.346189in}}%
\pgfpathlineto{\pgfqpoint{2.226701in}{1.338754in}}%
\pgfpathlineto{\pgfqpoint{2.226654in}{1.335261in}}%
\pgfpathlineto{\pgfqpoint{2.199219in}{1.335898in}}%
\pgfpathlineto{\pgfqpoint{2.198627in}{1.357636in}}%
\pgfpathlineto{\pgfqpoint{2.188988in}{1.359699in}}%
\pgfpathlineto{\pgfqpoint{2.182827in}{1.356530in}}%
\pgfpathlineto{\pgfqpoint{2.182439in}{1.381621in}}%
\pgfpathlineto{\pgfqpoint{2.181819in}{1.416001in}}%
\pgfpathlineto{\pgfqpoint{2.206018in}{1.417015in}}%
\pgfpathlineto{\pgfqpoint{2.206226in}{1.396171in}}%
\pgfpathlineto{\pgfqpoint{2.211089in}{1.394791in}}%
\pgfpathlineto{\pgfqpoint{2.212943in}{1.384548in}}%
\pgfpathlineto{\pgfqpoint{2.211917in}{1.379066in}}%
\pgfpathlineto{\pgfqpoint{2.215503in}{1.374729in}}%
\pgfpathlineto{\pgfqpoint{2.221539in}{1.373866in}}%
\pgfpathlineto{\pgfqpoint{2.224430in}{1.367560in}}%
\pgfpathlineto{\pgfqpoint{2.230926in}{1.365957in}}%
\pgfpathlineto{\pgfqpoint{2.233381in}{1.359171in}}%
\pgfpathlineto{\pgfqpoint{2.238456in}{1.359055in}}%
\pgfpathclose%
\pgfusepath{fill}%
\end{pgfscope}%
\begin{pgfscope}%
\pgfpathrectangle{\pgfqpoint{0.100000in}{0.100000in}}{\pgfqpoint{3.608454in}{2.310000in}}%
\pgfusepath{clip}%
\pgfsetbuttcap%
\pgfsetmiterjoin%
\definecolor{currentfill}{rgb}{0.000000,0.356863,0.821569}%
\pgfsetfillcolor{currentfill}%
\pgfsetlinewidth{0.000000pt}%
\definecolor{currentstroke}{rgb}{0.000000,0.000000,0.000000}%
\pgfsetstrokecolor{currentstroke}%
\pgfsetstrokeopacity{0.000000}%
\pgfsetdash{}{0pt}%
\pgfpathmoveto{\pgfqpoint{2.886708in}{1.438806in}}%
\pgfpathlineto{\pgfqpoint{2.887216in}{1.431682in}}%
\pgfpathlineto{\pgfqpoint{2.880330in}{1.431219in}}%
\pgfpathlineto{\pgfqpoint{2.880748in}{1.424222in}}%
\pgfpathlineto{\pgfqpoint{2.867003in}{1.423229in}}%
\pgfpathlineto{\pgfqpoint{2.867100in}{1.421144in}}%
\pgfpathlineto{\pgfqpoint{2.853130in}{1.419128in}}%
\pgfpathlineto{\pgfqpoint{2.852553in}{1.427030in}}%
\pgfpathlineto{\pgfqpoint{2.837395in}{1.428067in}}%
\pgfpathlineto{\pgfqpoint{2.820238in}{1.426570in}}%
\pgfpathlineto{\pgfqpoint{2.819277in}{1.440727in}}%
\pgfpathlineto{\pgfqpoint{2.794880in}{1.439319in}}%
\pgfpathlineto{\pgfqpoint{2.798306in}{1.443885in}}%
\pgfpathlineto{\pgfqpoint{2.800713in}{1.463456in}}%
\pgfpathlineto{\pgfqpoint{2.799407in}{1.480427in}}%
\pgfpathlineto{\pgfqpoint{2.795940in}{1.480285in}}%
\pgfpathlineto{\pgfqpoint{2.795078in}{1.502065in}}%
\pgfpathlineto{\pgfqpoint{2.801132in}{1.502933in}}%
\pgfpathlineto{\pgfqpoint{2.811270in}{1.504385in}}%
\pgfpathlineto{\pgfqpoint{2.811880in}{1.499432in}}%
\pgfpathlineto{\pgfqpoint{2.831500in}{1.499705in}}%
\pgfpathlineto{\pgfqpoint{2.831782in}{1.495222in}}%
\pgfpathlineto{\pgfqpoint{2.848845in}{1.496480in}}%
\pgfpathlineto{\pgfqpoint{2.846858in}{1.512561in}}%
\pgfpathlineto{\pgfqpoint{2.862804in}{1.514836in}}%
\pgfpathlineto{\pgfqpoint{2.871819in}{1.517489in}}%
\pgfpathlineto{\pgfqpoint{2.874225in}{1.518225in}}%
\pgfpathlineto{\pgfqpoint{2.876912in}{1.491928in}}%
\pgfpathlineto{\pgfqpoint{2.878220in}{1.486351in}}%
\pgfpathlineto{\pgfqpoint{2.879950in}{1.469292in}}%
\pgfpathlineto{\pgfqpoint{2.878244in}{1.466104in}}%
\pgfpathlineto{\pgfqpoint{2.864318in}{1.465519in}}%
\pgfpathlineto{\pgfqpoint{2.864747in}{1.458144in}}%
\pgfpathlineto{\pgfqpoint{2.869385in}{1.458446in}}%
\pgfpathlineto{\pgfqpoint{2.872185in}{1.451622in}}%
\pgfpathlineto{\pgfqpoint{2.872947in}{1.440094in}}%
\pgfpathlineto{\pgfqpoint{2.886708in}{1.438806in}}%
\pgfpathclose%
\pgfusepath{fill}%
\end{pgfscope}%
\begin{pgfscope}%
\pgfpathrectangle{\pgfqpoint{0.100000in}{0.100000in}}{\pgfqpoint{3.608454in}{2.310000in}}%
\pgfusepath{clip}%
\pgfsetbuttcap%
\pgfsetmiterjoin%
\definecolor{currentfill}{rgb}{0.000000,0.627451,0.686275}%
\pgfsetfillcolor{currentfill}%
\pgfsetlinewidth{0.000000pt}%
\definecolor{currentstroke}{rgb}{0.000000,0.000000,0.000000}%
\pgfsetstrokecolor{currentstroke}%
\pgfsetstrokeopacity{0.000000}%
\pgfsetdash{}{0pt}%
\pgfpathmoveto{\pgfqpoint{2.130612in}{1.428716in}}%
\pgfpathlineto{\pgfqpoint{2.130971in}{1.449316in}}%
\pgfpathlineto{\pgfqpoint{2.128892in}{1.463893in}}%
\pgfpathlineto{\pgfqpoint{2.138550in}{1.463912in}}%
\pgfpathlineto{\pgfqpoint{2.180394in}{1.464846in}}%
\pgfpathlineto{\pgfqpoint{2.181491in}{1.429705in}}%
\pgfpathlineto{\pgfqpoint{2.154147in}{1.429674in}}%
\pgfpathlineto{\pgfqpoint{2.154074in}{1.421652in}}%
\pgfpathlineto{\pgfqpoint{2.130676in}{1.421843in}}%
\pgfpathlineto{\pgfqpoint{2.130612in}{1.428716in}}%
\pgfpathclose%
\pgfusepath{fill}%
\end{pgfscope}%
\begin{pgfscope}%
\pgfpathrectangle{\pgfqpoint{0.100000in}{0.100000in}}{\pgfqpoint{3.608454in}{2.310000in}}%
\pgfusepath{clip}%
\pgfsetbuttcap%
\pgfsetmiterjoin%
\definecolor{currentfill}{rgb}{0.000000,0.427451,0.786275}%
\pgfsetfillcolor{currentfill}%
\pgfsetlinewidth{0.000000pt}%
\definecolor{currentstroke}{rgb}{0.000000,0.000000,0.000000}%
\pgfsetstrokecolor{currentstroke}%
\pgfsetstrokeopacity{0.000000}%
\pgfsetdash{}{0pt}%
\pgfpathmoveto{\pgfqpoint{2.073124in}{1.778481in}}%
\pgfpathlineto{\pgfqpoint{2.072659in}{1.785449in}}%
\pgfpathlineto{\pgfqpoint{2.024747in}{1.786093in}}%
\pgfpathlineto{\pgfqpoint{2.025168in}{1.813601in}}%
\pgfpathlineto{\pgfqpoint{2.025490in}{1.836778in}}%
\pgfpathlineto{\pgfqpoint{2.034935in}{1.834889in}}%
\pgfpathlineto{\pgfqpoint{2.041094in}{1.830482in}}%
\pgfpathlineto{\pgfqpoint{2.044201in}{1.834191in}}%
\pgfpathlineto{\pgfqpoint{2.044431in}{1.847876in}}%
\pgfpathlineto{\pgfqpoint{2.099439in}{1.847533in}}%
\pgfpathlineto{\pgfqpoint{2.120016in}{1.847577in}}%
\pgfpathlineto{\pgfqpoint{2.120021in}{1.840682in}}%
\pgfpathlineto{\pgfqpoint{2.120405in}{1.806180in}}%
\pgfpathlineto{\pgfqpoint{2.079580in}{1.806133in}}%
\pgfpathlineto{\pgfqpoint{2.079554in}{1.795195in}}%
\pgfpathlineto{\pgfqpoint{2.086458in}{1.790823in}}%
\pgfpathlineto{\pgfqpoint{2.086396in}{1.778361in}}%
\pgfpathlineto{\pgfqpoint{2.073124in}{1.778481in}}%
\pgfpathclose%
\pgfusepath{fill}%
\end{pgfscope}%
\begin{pgfscope}%
\pgfpathrectangle{\pgfqpoint{0.100000in}{0.100000in}}{\pgfqpoint{3.608454in}{2.310000in}}%
\pgfusepath{clip}%
\pgfsetbuttcap%
\pgfsetmiterjoin%
\definecolor{currentfill}{rgb}{0.000000,0.705882,0.647059}%
\pgfsetfillcolor{currentfill}%
\pgfsetlinewidth{0.000000pt}%
\definecolor{currentstroke}{rgb}{0.000000,0.000000,0.000000}%
\pgfsetstrokecolor{currentstroke}%
\pgfsetstrokeopacity{0.000000}%
\pgfsetdash{}{0pt}%
\pgfpathmoveto{\pgfqpoint{2.248336in}{1.167927in}}%
\pgfpathlineto{\pgfqpoint{2.248459in}{1.143536in}}%
\pgfpathlineto{\pgfqpoint{2.243255in}{1.143432in}}%
\pgfpathlineto{\pgfqpoint{2.178864in}{1.142325in}}%
\pgfpathlineto{\pgfqpoint{2.165488in}{1.142173in}}%
\pgfpathlineto{\pgfqpoint{2.165811in}{1.161887in}}%
\pgfpathlineto{\pgfqpoint{2.166408in}{1.204778in}}%
\pgfpathlineto{\pgfqpoint{2.193249in}{1.204475in}}%
\pgfpathlineto{\pgfqpoint{2.194554in}{1.204465in}}%
\pgfpathlineto{\pgfqpoint{2.194584in}{1.181869in}}%
\pgfpathlineto{\pgfqpoint{2.211810in}{1.181895in}}%
\pgfpathlineto{\pgfqpoint{2.211773in}{1.168036in}}%
\pgfpathlineto{\pgfqpoint{2.248336in}{1.167927in}}%
\pgfpathclose%
\pgfusepath{fill}%
\end{pgfscope}%
\begin{pgfscope}%
\pgfpathrectangle{\pgfqpoint{0.100000in}{0.100000in}}{\pgfqpoint{3.608454in}{2.310000in}}%
\pgfusepath{clip}%
\pgfsetbuttcap%
\pgfsetmiterjoin%
\definecolor{currentfill}{rgb}{0.000000,0.517647,0.741176}%
\pgfsetfillcolor{currentfill}%
\pgfsetlinewidth{0.000000pt}%
\definecolor{currentstroke}{rgb}{0.000000,0.000000,0.000000}%
\pgfsetstrokecolor{currentstroke}%
\pgfsetstrokeopacity{0.000000}%
\pgfsetdash{}{0pt}%
\pgfpathmoveto{\pgfqpoint{2.722959in}{1.184982in}}%
\pgfpathlineto{\pgfqpoint{2.754147in}{1.186574in}}%
\pgfpathlineto{\pgfqpoint{2.758225in}{1.180770in}}%
\pgfpathlineto{\pgfqpoint{2.744380in}{1.183721in}}%
\pgfpathlineto{\pgfqpoint{2.738211in}{1.181250in}}%
\pgfpathlineto{\pgfqpoint{2.734679in}{1.168795in}}%
\pgfpathlineto{\pgfqpoint{2.735445in}{1.160714in}}%
\pgfpathlineto{\pgfqpoint{2.738900in}{1.152461in}}%
\pgfpathlineto{\pgfqpoint{2.735267in}{1.144371in}}%
\pgfpathlineto{\pgfqpoint{2.729930in}{1.144288in}}%
\pgfpathlineto{\pgfqpoint{2.729200in}{1.134130in}}%
\pgfpathlineto{\pgfqpoint{2.722754in}{1.138754in}}%
\pgfpathlineto{\pgfqpoint{2.712831in}{1.140583in}}%
\pgfpathlineto{\pgfqpoint{2.703971in}{1.134522in}}%
\pgfpathlineto{\pgfqpoint{2.701451in}{1.138405in}}%
\pgfpathlineto{\pgfqpoint{2.695803in}{1.138527in}}%
\pgfpathlineto{\pgfqpoint{2.693376in}{1.142517in}}%
\pgfpathlineto{\pgfqpoint{2.693285in}{1.148452in}}%
\pgfpathlineto{\pgfqpoint{2.689646in}{1.154594in}}%
\pgfpathlineto{\pgfqpoint{2.690027in}{1.171438in}}%
\pgfpathlineto{\pgfqpoint{2.696790in}{1.173951in}}%
\pgfpathlineto{\pgfqpoint{2.705300in}{1.170782in}}%
\pgfpathlineto{\pgfqpoint{2.710970in}{1.165941in}}%
\pgfpathlineto{\pgfqpoint{2.713741in}{1.174022in}}%
\pgfpathlineto{\pgfqpoint{2.722324in}{1.177007in}}%
\pgfpathlineto{\pgfqpoint{2.722959in}{1.184982in}}%
\pgfpathclose%
\pgfusepath{fill}%
\end{pgfscope}%
\begin{pgfscope}%
\pgfpathrectangle{\pgfqpoint{0.100000in}{0.100000in}}{\pgfqpoint{3.608454in}{2.310000in}}%
\pgfusepath{clip}%
\pgfsetbuttcap%
\pgfsetmiterjoin%
\definecolor{currentfill}{rgb}{0.000000,0.341176,0.829412}%
\pgfsetfillcolor{currentfill}%
\pgfsetlinewidth{0.000000pt}%
\definecolor{currentstroke}{rgb}{0.000000,0.000000,0.000000}%
\pgfsetstrokecolor{currentstroke}%
\pgfsetstrokeopacity{0.000000}%
\pgfsetdash{}{0pt}%
\pgfpathmoveto{\pgfqpoint{1.586973in}{1.721839in}}%
\pgfpathlineto{\pgfqpoint{1.538610in}{1.726627in}}%
\pgfpathlineto{\pgfqpoint{1.474731in}{1.733639in}}%
\pgfpathlineto{\pgfqpoint{1.481253in}{1.786760in}}%
\pgfpathlineto{\pgfqpoint{1.486605in}{1.827974in}}%
\pgfpathlineto{\pgfqpoint{1.488701in}{1.851762in}}%
\pgfpathlineto{\pgfqpoint{1.494987in}{1.851589in}}%
\pgfpathlineto{\pgfqpoint{1.543825in}{1.845887in}}%
\pgfpathlineto{\pgfqpoint{1.598608in}{1.839877in}}%
\pgfpathlineto{\pgfqpoint{1.592333in}{1.775294in}}%
\pgfpathlineto{\pgfqpoint{1.586973in}{1.721839in}}%
\pgfpathclose%
\pgfusepath{fill}%
\end{pgfscope}%
\begin{pgfscope}%
\pgfpathrectangle{\pgfqpoint{0.100000in}{0.100000in}}{\pgfqpoint{3.608454in}{2.310000in}}%
\pgfusepath{clip}%
\pgfsetbuttcap%
\pgfsetmiterjoin%
\definecolor{currentfill}{rgb}{0.000000,0.776471,0.611765}%
\pgfsetfillcolor{currentfill}%
\pgfsetlinewidth{0.000000pt}%
\definecolor{currentstroke}{rgb}{0.000000,0.000000,0.000000}%
\pgfsetstrokecolor{currentstroke}%
\pgfsetstrokeopacity{0.000000}%
\pgfsetdash{}{0pt}%
\pgfpathmoveto{\pgfqpoint{1.365564in}{0.266179in}}%
\pgfpathlineto{\pgfqpoint{1.364279in}{0.272148in}}%
\pgfpathlineto{\pgfqpoint{1.358083in}{0.284462in}}%
\pgfpathlineto{\pgfqpoint{1.360076in}{0.291758in}}%
\pgfpathlineto{\pgfqpoint{1.373325in}{0.304088in}}%
\pgfpathlineto{\pgfqpoint{1.377160in}{0.317108in}}%
\pgfpathlineto{\pgfqpoint{1.381805in}{0.326247in}}%
\pgfpathlineto{\pgfqpoint{1.380932in}{0.332508in}}%
\pgfpathlineto{\pgfqpoint{1.385909in}{0.340951in}}%
\pgfpathlineto{\pgfqpoint{1.410377in}{0.343837in}}%
\pgfpathlineto{\pgfqpoint{1.412863in}{0.346234in}}%
\pgfpathlineto{\pgfqpoint{1.413879in}{0.356269in}}%
\pgfpathlineto{\pgfqpoint{1.419360in}{0.363772in}}%
\pgfpathlineto{\pgfqpoint{1.424674in}{0.362812in}}%
\pgfpathlineto{\pgfqpoint{1.428368in}{0.354664in}}%
\pgfpathlineto{\pgfqpoint{1.431876in}{0.341640in}}%
\pgfpathlineto{\pgfqpoint{1.439179in}{0.333553in}}%
\pgfpathlineto{\pgfqpoint{1.446187in}{0.319277in}}%
\pgfpathlineto{\pgfqpoint{1.449694in}{0.299992in}}%
\pgfpathlineto{\pgfqpoint{1.442998in}{0.291249in}}%
\pgfpathlineto{\pgfqpoint{1.448609in}{0.287822in}}%
\pgfpathlineto{\pgfqpoint{1.445177in}{0.280721in}}%
\pgfpathlineto{\pgfqpoint{1.446258in}{0.271967in}}%
\pgfpathlineto{\pgfqpoint{1.449844in}{0.264785in}}%
\pgfpathlineto{\pgfqpoint{1.446708in}{0.261925in}}%
\pgfpathlineto{\pgfqpoint{1.430233in}{0.260876in}}%
\pgfpathlineto{\pgfqpoint{1.416961in}{0.263446in}}%
\pgfpathlineto{\pgfqpoint{1.406310in}{0.270716in}}%
\pgfpathlineto{\pgfqpoint{1.400271in}{0.270027in}}%
\pgfpathlineto{\pgfqpoint{1.387382in}{0.272320in}}%
\pgfpathlineto{\pgfqpoint{1.365564in}{0.266179in}}%
\pgfpathclose%
\pgfusepath{fill}%
\end{pgfscope}%
\begin{pgfscope}%
\pgfpathrectangle{\pgfqpoint{0.100000in}{0.100000in}}{\pgfqpoint{3.608454in}{2.310000in}}%
\pgfusepath{clip}%
\pgfsetbuttcap%
\pgfsetmiterjoin%
\definecolor{currentfill}{rgb}{0.000000,0.368627,0.815686}%
\pgfsetfillcolor{currentfill}%
\pgfsetlinewidth{0.000000pt}%
\definecolor{currentstroke}{rgb}{0.000000,0.000000,0.000000}%
\pgfsetstrokecolor{currentstroke}%
\pgfsetstrokeopacity{0.000000}%
\pgfsetdash{}{0pt}%
\pgfpathmoveto{\pgfqpoint{2.168433in}{1.778817in}}%
\pgfpathlineto{\pgfqpoint{2.167464in}{1.771934in}}%
\pgfpathlineto{\pgfqpoint{2.128088in}{1.771585in}}%
\pgfpathlineto{\pgfqpoint{2.127691in}{1.792284in}}%
\pgfpathlineto{\pgfqpoint{2.135124in}{1.792351in}}%
\pgfpathlineto{\pgfqpoint{2.135067in}{1.806256in}}%
\pgfpathlineto{\pgfqpoint{2.120405in}{1.806180in}}%
\pgfpathlineto{\pgfqpoint{2.120021in}{1.840682in}}%
\pgfpathlineto{\pgfqpoint{2.141258in}{1.840834in}}%
\pgfpathlineto{\pgfqpoint{2.148148in}{1.837483in}}%
\pgfpathlineto{\pgfqpoint{2.148622in}{1.813278in}}%
\pgfpathlineto{\pgfqpoint{2.162319in}{1.813349in}}%
\pgfpathlineto{\pgfqpoint{2.162655in}{1.792587in}}%
\pgfpathlineto{\pgfqpoint{2.169520in}{1.792650in}}%
\pgfpathlineto{\pgfqpoint{2.168433in}{1.778817in}}%
\pgfpathclose%
\pgfusepath{fill}%
\end{pgfscope}%
\begin{pgfscope}%
\pgfpathrectangle{\pgfqpoint{0.100000in}{0.100000in}}{\pgfqpoint{3.608454in}{2.310000in}}%
\pgfusepath{clip}%
\pgfsetbuttcap%
\pgfsetmiterjoin%
\definecolor{currentfill}{rgb}{0.000000,0.454902,0.772549}%
\pgfsetfillcolor{currentfill}%
\pgfsetlinewidth{0.000000pt}%
\definecolor{currentstroke}{rgb}{0.000000,0.000000,0.000000}%
\pgfsetstrokecolor{currentstroke}%
\pgfsetstrokeopacity{0.000000}%
\pgfsetdash{}{0pt}%
\pgfpathmoveto{\pgfqpoint{2.636960in}{1.295013in}}%
\pgfpathlineto{\pgfqpoint{2.634127in}{1.288599in}}%
\pgfpathlineto{\pgfqpoint{2.635552in}{1.279697in}}%
\pgfpathlineto{\pgfqpoint{2.630595in}{1.278625in}}%
\pgfpathlineto{\pgfqpoint{2.627023in}{1.272051in}}%
\pgfpathlineto{\pgfqpoint{2.627054in}{1.277303in}}%
\pgfpathlineto{\pgfqpoint{2.620883in}{1.276667in}}%
\pgfpathlineto{\pgfqpoint{2.617290in}{1.282808in}}%
\pgfpathlineto{\pgfqpoint{2.602387in}{1.273595in}}%
\pgfpathlineto{\pgfqpoint{2.601681in}{1.267043in}}%
\pgfpathlineto{\pgfqpoint{2.586412in}{1.272897in}}%
\pgfpathlineto{\pgfqpoint{2.578231in}{1.275886in}}%
\pgfpathlineto{\pgfqpoint{2.570793in}{1.273710in}}%
\pgfpathlineto{\pgfqpoint{2.567447in}{1.266096in}}%
\pgfpathlineto{\pgfqpoint{2.562848in}{1.271335in}}%
\pgfpathlineto{\pgfqpoint{2.553706in}{1.268477in}}%
\pgfpathlineto{\pgfqpoint{2.546912in}{1.269077in}}%
\pgfpathlineto{\pgfqpoint{2.549388in}{1.262674in}}%
\pgfpathlineto{\pgfqpoint{2.541890in}{1.261426in}}%
\pgfpathlineto{\pgfqpoint{2.537345in}{1.268351in}}%
\pgfpathlineto{\pgfqpoint{2.545025in}{1.287628in}}%
\pgfpathlineto{\pgfqpoint{2.541416in}{1.297901in}}%
\pgfpathlineto{\pgfqpoint{2.543804in}{1.301922in}}%
\pgfpathlineto{\pgfqpoint{2.541523in}{1.307274in}}%
\pgfpathlineto{\pgfqpoint{2.542054in}{1.316771in}}%
\pgfpathlineto{\pgfqpoint{2.544363in}{1.322774in}}%
\pgfpathlineto{\pgfqpoint{2.560427in}{1.323864in}}%
\pgfpathlineto{\pgfqpoint{2.565371in}{1.317560in}}%
\pgfpathlineto{\pgfqpoint{2.572272in}{1.322088in}}%
\pgfpathlineto{\pgfqpoint{2.585800in}{1.324135in}}%
\pgfpathlineto{\pgfqpoint{2.594128in}{1.324280in}}%
\pgfpathlineto{\pgfqpoint{2.608516in}{1.322156in}}%
\pgfpathlineto{\pgfqpoint{2.610839in}{1.324868in}}%
\pgfpathlineto{\pgfqpoint{2.620401in}{1.325720in}}%
\pgfpathlineto{\pgfqpoint{2.621408in}{1.315375in}}%
\pgfpathlineto{\pgfqpoint{2.622478in}{1.305021in}}%
\pgfpathlineto{\pgfqpoint{2.629213in}{1.305901in}}%
\pgfpathlineto{\pgfqpoint{2.629550in}{1.301264in}}%
\pgfpathlineto{\pgfqpoint{2.636434in}{1.301852in}}%
\pgfpathlineto{\pgfqpoint{2.636960in}{1.295013in}}%
\pgfpathclose%
\pgfusepath{fill}%
\end{pgfscope}%
\begin{pgfscope}%
\pgfpathrectangle{\pgfqpoint{0.100000in}{0.100000in}}{\pgfqpoint{3.608454in}{2.310000in}}%
\pgfusepath{clip}%
\pgfsetbuttcap%
\pgfsetmiterjoin%
\definecolor{currentfill}{rgb}{0.000000,0.721569,0.639216}%
\pgfsetfillcolor{currentfill}%
\pgfsetlinewidth{0.000000pt}%
\definecolor{currentstroke}{rgb}{0.000000,0.000000,0.000000}%
\pgfsetstrokecolor{currentstroke}%
\pgfsetstrokeopacity{0.000000}%
\pgfsetdash{}{0pt}%
\pgfpathmoveto{\pgfqpoint{2.295263in}{1.958166in}}%
\pgfpathlineto{\pgfqpoint{2.297058in}{1.910416in}}%
\pgfpathlineto{\pgfqpoint{2.269571in}{1.909501in}}%
\pgfpathlineto{\pgfqpoint{2.256111in}{1.909051in}}%
\pgfpathlineto{\pgfqpoint{2.255541in}{1.929754in}}%
\pgfpathlineto{\pgfqpoint{2.213676in}{1.928846in}}%
\pgfpathlineto{\pgfqpoint{2.212805in}{1.945663in}}%
\pgfpathlineto{\pgfqpoint{2.212580in}{1.977269in}}%
\pgfpathlineto{\pgfqpoint{2.211215in}{2.006566in}}%
\pgfpathlineto{\pgfqpoint{2.211106in}{2.032691in}}%
\pgfpathlineto{\pgfqpoint{2.209745in}{2.046401in}}%
\pgfpathlineto{\pgfqpoint{2.208627in}{2.060466in}}%
\pgfpathlineto{\pgfqpoint{2.208113in}{2.105322in}}%
\pgfpathlineto{\pgfqpoint{2.213604in}{2.105132in}}%
\pgfpathlineto{\pgfqpoint{2.227390in}{2.098684in}}%
\pgfpathlineto{\pgfqpoint{2.232324in}{2.099097in}}%
\pgfpathlineto{\pgfqpoint{2.231425in}{2.090560in}}%
\pgfpathlineto{\pgfqpoint{2.239307in}{2.091665in}}%
\pgfpathlineto{\pgfqpoint{2.247149in}{2.073655in}}%
\pgfpathlineto{\pgfqpoint{2.251836in}{2.075702in}}%
\pgfpathlineto{\pgfqpoint{2.251160in}{2.083595in}}%
\pgfpathlineto{\pgfqpoint{2.263561in}{2.085282in}}%
\pgfpathlineto{\pgfqpoint{2.266399in}{2.077856in}}%
\pgfpathlineto{\pgfqpoint{2.274138in}{2.073437in}}%
\pgfpathlineto{\pgfqpoint{2.282131in}{2.073030in}}%
\pgfpathlineto{\pgfqpoint{2.282543in}{2.066327in}}%
\pgfpathlineto{\pgfqpoint{2.297364in}{2.061677in}}%
\pgfpathlineto{\pgfqpoint{2.307312in}{2.064841in}}%
\pgfpathlineto{\pgfqpoint{2.318616in}{2.073687in}}%
\pgfpathlineto{\pgfqpoint{2.320258in}{2.049681in}}%
\pgfpathlineto{\pgfqpoint{2.321610in}{2.015873in}}%
\pgfpathlineto{\pgfqpoint{2.309782in}{2.000701in}}%
\pgfpathlineto{\pgfqpoint{2.298141in}{1.987730in}}%
\pgfpathlineto{\pgfqpoint{2.293110in}{1.984555in}}%
\pgfpathlineto{\pgfqpoint{2.282209in}{1.972690in}}%
\pgfpathlineto{\pgfqpoint{2.265774in}{1.959901in}}%
\pgfpathlineto{\pgfqpoint{2.269412in}{1.953632in}}%
\pgfpathlineto{\pgfqpoint{2.274053in}{1.951302in}}%
\pgfpathlineto{\pgfqpoint{2.282372in}{1.952765in}}%
\pgfpathlineto{\pgfqpoint{2.295263in}{1.958166in}}%
\pgfpathclose%
\pgfusepath{fill}%
\end{pgfscope}%
\begin{pgfscope}%
\pgfpathrectangle{\pgfqpoint{0.100000in}{0.100000in}}{\pgfqpoint{3.608454in}{2.310000in}}%
\pgfusepath{clip}%
\pgfsetbuttcap%
\pgfsetmiterjoin%
\definecolor{currentfill}{rgb}{0.000000,0.372549,0.813725}%
\pgfsetfillcolor{currentfill}%
\pgfsetlinewidth{0.000000pt}%
\definecolor{currentstroke}{rgb}{0.000000,0.000000,0.000000}%
\pgfsetstrokecolor{currentstroke}%
\pgfsetstrokeopacity{0.000000}%
\pgfsetdash{}{0pt}%
\pgfpathmoveto{\pgfqpoint{2.287673in}{1.541193in}}%
\pgfpathlineto{\pgfqpoint{2.266869in}{1.540448in}}%
\pgfpathlineto{\pgfqpoint{2.264937in}{1.602761in}}%
\pgfpathlineto{\pgfqpoint{2.306014in}{1.604048in}}%
\pgfpathlineto{\pgfqpoint{2.346907in}{1.605656in}}%
\pgfpathlineto{\pgfqpoint{2.348123in}{1.578046in}}%
\pgfpathlineto{\pgfqpoint{2.320640in}{1.576968in}}%
\pgfpathlineto{\pgfqpoint{2.320902in}{1.570069in}}%
\pgfpathlineto{\pgfqpoint{2.293482in}{1.569183in}}%
\pgfpathlineto{\pgfqpoint{2.294578in}{1.541485in}}%
\pgfpathlineto{\pgfqpoint{2.287673in}{1.541193in}}%
\pgfpathclose%
\pgfusepath{fill}%
\end{pgfscope}%
\begin{pgfscope}%
\pgfpathrectangle{\pgfqpoint{0.100000in}{0.100000in}}{\pgfqpoint{3.608454in}{2.310000in}}%
\pgfusepath{clip}%
\pgfsetbuttcap%
\pgfsetmiterjoin%
\definecolor{currentfill}{rgb}{0.000000,0.549020,0.725490}%
\pgfsetfillcolor{currentfill}%
\pgfsetlinewidth{0.000000pt}%
\definecolor{currentstroke}{rgb}{0.000000,0.000000,0.000000}%
\pgfsetstrokecolor{currentstroke}%
\pgfsetstrokeopacity{0.000000}%
\pgfsetdash{}{0pt}%
\pgfpathmoveto{\pgfqpoint{3.039659in}{0.833647in}}%
\pgfpathlineto{\pgfqpoint{3.033775in}{0.835716in}}%
\pgfpathlineto{\pgfqpoint{3.025040in}{0.841844in}}%
\pgfpathlineto{\pgfqpoint{3.020884in}{0.848437in}}%
\pgfpathlineto{\pgfqpoint{3.011213in}{0.848731in}}%
\pgfpathlineto{\pgfqpoint{3.005118in}{0.855085in}}%
\pgfpathlineto{\pgfqpoint{2.996708in}{0.858258in}}%
\pgfpathlineto{\pgfqpoint{2.994468in}{0.861760in}}%
\pgfpathlineto{\pgfqpoint{3.016488in}{0.872188in}}%
\pgfpathlineto{\pgfqpoint{3.014555in}{0.877148in}}%
\pgfpathlineto{\pgfqpoint{3.008209in}{0.883079in}}%
\pgfpathlineto{\pgfqpoint{3.006499in}{0.892231in}}%
\pgfpathlineto{\pgfqpoint{3.015378in}{0.900512in}}%
\pgfpathlineto{\pgfqpoint{3.022585in}{0.898037in}}%
\pgfpathlineto{\pgfqpoint{3.038277in}{0.916697in}}%
\pgfpathlineto{\pgfqpoint{3.047304in}{0.907519in}}%
\pgfpathlineto{\pgfqpoint{3.051321in}{0.914811in}}%
\pgfpathlineto{\pgfqpoint{3.061084in}{0.913825in}}%
\pgfpathlineto{\pgfqpoint{3.066440in}{0.907291in}}%
\pgfpathlineto{\pgfqpoint{3.071225in}{0.906814in}}%
\pgfpathlineto{\pgfqpoint{3.077208in}{0.901258in}}%
\pgfpathlineto{\pgfqpoint{3.081530in}{0.894395in}}%
\pgfpathlineto{\pgfqpoint{3.080811in}{0.888841in}}%
\pgfpathlineto{\pgfqpoint{3.069735in}{0.881587in}}%
\pgfpathlineto{\pgfqpoint{3.064847in}{0.872766in}}%
\pgfpathlineto{\pgfqpoint{3.055250in}{0.864816in}}%
\pgfpathlineto{\pgfqpoint{3.057975in}{0.856666in}}%
\pgfpathlineto{\pgfqpoint{3.043134in}{0.834472in}}%
\pgfpathlineto{\pgfqpoint{3.039659in}{0.833647in}}%
\pgfpathclose%
\pgfusepath{fill}%
\end{pgfscope}%
\begin{pgfscope}%
\pgfpathrectangle{\pgfqpoint{0.100000in}{0.100000in}}{\pgfqpoint{3.608454in}{2.310000in}}%
\pgfusepath{clip}%
\pgfsetbuttcap%
\pgfsetmiterjoin%
\definecolor{currentfill}{rgb}{0.000000,0.909804,0.545098}%
\pgfsetfillcolor{currentfill}%
\pgfsetlinewidth{0.000000pt}%
\definecolor{currentstroke}{rgb}{0.000000,0.000000,0.000000}%
\pgfsetstrokecolor{currentstroke}%
\pgfsetstrokeopacity{0.000000}%
\pgfsetdash{}{0pt}%
\pgfpathmoveto{\pgfqpoint{0.715935in}{1.921940in}}%
\pgfpathlineto{\pgfqpoint{0.762277in}{1.909354in}}%
\pgfpathlineto{\pgfqpoint{0.764102in}{1.916069in}}%
\pgfpathlineto{\pgfqpoint{0.796408in}{1.906884in}}%
\pgfpathlineto{\pgfqpoint{0.790924in}{1.886917in}}%
\pgfpathlineto{\pgfqpoint{0.773761in}{1.820585in}}%
\pgfpathlineto{\pgfqpoint{0.774453in}{1.820402in}}%
\pgfpathlineto{\pgfqpoint{0.761786in}{1.771448in}}%
\pgfpathlineto{\pgfqpoint{0.757203in}{1.749818in}}%
\pgfpathlineto{\pgfqpoint{0.728707in}{1.756911in}}%
\pgfpathlineto{\pgfqpoint{0.693113in}{1.766738in}}%
\pgfpathlineto{\pgfqpoint{0.691075in}{1.767312in}}%
\pgfpathlineto{\pgfqpoint{0.706614in}{1.825006in}}%
\pgfpathlineto{\pgfqpoint{0.674134in}{1.833816in}}%
\pgfpathlineto{\pgfqpoint{0.684062in}{1.866632in}}%
\pgfpathlineto{\pgfqpoint{0.686047in}{1.866076in}}%
\pgfpathlineto{\pgfqpoint{0.695319in}{1.898980in}}%
\pgfpathlineto{\pgfqpoint{0.697108in}{1.905730in}}%
\pgfpathlineto{\pgfqpoint{0.703763in}{1.903845in}}%
\pgfpathlineto{\pgfqpoint{0.709441in}{1.923753in}}%
\pgfpathlineto{\pgfqpoint{0.715935in}{1.921940in}}%
\pgfpathclose%
\pgfusepath{fill}%
\end{pgfscope}%
\begin{pgfscope}%
\pgfpathrectangle{\pgfqpoint{0.100000in}{0.100000in}}{\pgfqpoint{3.608454in}{2.310000in}}%
\pgfusepath{clip}%
\pgfsetbuttcap%
\pgfsetmiterjoin%
\definecolor{currentfill}{rgb}{0.000000,0.568627,0.715686}%
\pgfsetfillcolor{currentfill}%
\pgfsetlinewidth{0.000000pt}%
\definecolor{currentstroke}{rgb}{0.000000,0.000000,0.000000}%
\pgfsetstrokecolor{currentstroke}%
\pgfsetstrokeopacity{0.000000}%
\pgfsetdash{}{0pt}%
\pgfpathmoveto{\pgfqpoint{3.127726in}{0.403873in}}%
\pgfpathlineto{\pgfqpoint{3.099914in}{0.399053in}}%
\pgfpathlineto{\pgfqpoint{3.096468in}{0.420138in}}%
\pgfpathlineto{\pgfqpoint{3.075618in}{0.416692in}}%
\pgfpathlineto{\pgfqpoint{3.072400in}{0.437158in}}%
\pgfpathlineto{\pgfqpoint{3.069493in}{0.458265in}}%
\pgfpathlineto{\pgfqpoint{3.090529in}{0.461360in}}%
\pgfpathlineto{\pgfqpoint{3.089414in}{0.468427in}}%
\pgfpathlineto{\pgfqpoint{3.096465in}{0.469557in}}%
\pgfpathlineto{\pgfqpoint{3.095343in}{0.476565in}}%
\pgfpathlineto{\pgfqpoint{3.109319in}{0.479362in}}%
\pgfpathlineto{\pgfqpoint{3.100740in}{0.491314in}}%
\pgfpathlineto{\pgfqpoint{3.097174in}{0.491133in}}%
\pgfpathlineto{\pgfqpoint{3.088462in}{0.501164in}}%
\pgfpathlineto{\pgfqpoint{3.091930in}{0.511405in}}%
\pgfpathlineto{\pgfqpoint{3.110913in}{0.514275in}}%
\pgfpathlineto{\pgfqpoint{3.117607in}{0.515401in}}%
\pgfpathlineto{\pgfqpoint{3.118731in}{0.508688in}}%
\pgfpathlineto{\pgfqpoint{3.150978in}{0.513727in}}%
\pgfpathlineto{\pgfqpoint{3.156513in}{0.502751in}}%
\pgfpathlineto{\pgfqpoint{3.167005in}{0.486501in}}%
\pgfpathlineto{\pgfqpoint{3.173443in}{0.474576in}}%
\pgfpathlineto{\pgfqpoint{3.181412in}{0.456522in}}%
\pgfpathlineto{\pgfqpoint{3.183733in}{0.440357in}}%
\pgfpathlineto{\pgfqpoint{3.184616in}{0.418072in}}%
\pgfpathlineto{\pgfqpoint{3.173551in}{0.417385in}}%
\pgfpathlineto{\pgfqpoint{3.126693in}{0.409803in}}%
\pgfpathlineto{\pgfqpoint{3.127726in}{0.403873in}}%
\pgfpathclose%
\pgfusepath{fill}%
\end{pgfscope}%
\begin{pgfscope}%
\pgfpathrectangle{\pgfqpoint{0.100000in}{0.100000in}}{\pgfqpoint{3.608454in}{2.310000in}}%
\pgfusepath{clip}%
\pgfsetbuttcap%
\pgfsetmiterjoin%
\definecolor{currentfill}{rgb}{0.000000,0.521569,0.739216}%
\pgfsetfillcolor{currentfill}%
\pgfsetlinewidth{0.000000pt}%
\definecolor{currentstroke}{rgb}{0.000000,0.000000,0.000000}%
\pgfsetstrokecolor{currentstroke}%
\pgfsetstrokeopacity{0.000000}%
\pgfsetdash{}{0pt}%
\pgfpathmoveto{\pgfqpoint{2.621408in}{1.315375in}}%
\pgfpathlineto{\pgfqpoint{2.620401in}{1.325720in}}%
\pgfpathlineto{\pgfqpoint{2.610839in}{1.324868in}}%
\pgfpathlineto{\pgfqpoint{2.608516in}{1.322156in}}%
\pgfpathlineto{\pgfqpoint{2.594128in}{1.324280in}}%
\pgfpathlineto{\pgfqpoint{2.585800in}{1.324135in}}%
\pgfpathlineto{\pgfqpoint{2.582510in}{1.335565in}}%
\pgfpathlineto{\pgfqpoint{2.584141in}{1.345785in}}%
\pgfpathlineto{\pgfqpoint{2.590013in}{1.348185in}}%
\pgfpathlineto{\pgfqpoint{2.590232in}{1.353266in}}%
\pgfpathlineto{\pgfqpoint{2.583434in}{1.352723in}}%
\pgfpathlineto{\pgfqpoint{2.581657in}{1.373507in}}%
\pgfpathlineto{\pgfqpoint{2.593047in}{1.374201in}}%
\pgfpathlineto{\pgfqpoint{2.591897in}{1.387949in}}%
\pgfpathlineto{\pgfqpoint{2.598672in}{1.388501in}}%
\pgfpathlineto{\pgfqpoint{2.597905in}{1.398852in}}%
\pgfpathlineto{\pgfqpoint{2.613308in}{1.399991in}}%
\pgfpathlineto{\pgfqpoint{2.616632in}{1.400311in}}%
\pgfpathlineto{\pgfqpoint{2.617652in}{1.390119in}}%
\pgfpathlineto{\pgfqpoint{2.632842in}{1.391462in}}%
\pgfpathlineto{\pgfqpoint{2.635715in}{1.368602in}}%
\pgfpathlineto{\pgfqpoint{2.638885in}{1.368939in}}%
\pgfpathlineto{\pgfqpoint{2.641627in}{1.364614in}}%
\pgfpathlineto{\pgfqpoint{2.643662in}{1.346766in}}%
\pgfpathlineto{\pgfqpoint{2.641760in}{1.344149in}}%
\pgfpathlineto{\pgfqpoint{2.644392in}{1.317477in}}%
\pgfpathlineto{\pgfqpoint{2.621408in}{1.315375in}}%
\pgfpathclose%
\pgfusepath{fill}%
\end{pgfscope}%
\begin{pgfscope}%
\pgfpathrectangle{\pgfqpoint{0.100000in}{0.100000in}}{\pgfqpoint{3.608454in}{2.310000in}}%
\pgfusepath{clip}%
\pgfsetbuttcap%
\pgfsetmiterjoin%
\definecolor{currentfill}{rgb}{0.000000,0.905882,0.547059}%
\pgfsetfillcolor{currentfill}%
\pgfsetlinewidth{0.000000pt}%
\definecolor{currentstroke}{rgb}{0.000000,0.000000,0.000000}%
\pgfsetstrokecolor{currentstroke}%
\pgfsetstrokeopacity{0.000000}%
\pgfsetdash{}{0pt}%
\pgfpathmoveto{\pgfqpoint{2.258858in}{0.847301in}}%
\pgfpathlineto{\pgfqpoint{2.267654in}{0.847225in}}%
\pgfpathlineto{\pgfqpoint{2.272892in}{0.840599in}}%
\pgfpathlineto{\pgfqpoint{2.279806in}{0.840743in}}%
\pgfpathlineto{\pgfqpoint{2.280375in}{0.819890in}}%
\pgfpathlineto{\pgfqpoint{2.287535in}{0.813168in}}%
\pgfpathlineto{\pgfqpoint{2.288408in}{0.781863in}}%
\pgfpathlineto{\pgfqpoint{2.285332in}{0.771355in}}%
\pgfpathlineto{\pgfqpoint{2.267964in}{0.770915in}}%
\pgfpathlineto{\pgfqpoint{2.268095in}{0.763939in}}%
\pgfpathlineto{\pgfqpoint{2.244140in}{0.763307in}}%
\pgfpathlineto{\pgfqpoint{2.246031in}{0.774913in}}%
\pgfpathlineto{\pgfqpoint{2.249736in}{0.781810in}}%
\pgfpathlineto{\pgfqpoint{2.245714in}{0.798315in}}%
\pgfpathlineto{\pgfqpoint{2.254067in}{0.798428in}}%
\pgfpathlineto{\pgfqpoint{2.256664in}{0.805503in}}%
\pgfpathlineto{\pgfqpoint{2.255994in}{0.822878in}}%
\pgfpathlineto{\pgfqpoint{2.249143in}{0.822767in}}%
\pgfpathlineto{\pgfqpoint{2.248708in}{0.834367in}}%
\pgfpathlineto{\pgfqpoint{2.251967in}{0.847074in}}%
\pgfpathlineto{\pgfqpoint{2.258858in}{0.847301in}}%
\pgfpathclose%
\pgfusepath{fill}%
\end{pgfscope}%
\begin{pgfscope}%
\pgfpathrectangle{\pgfqpoint{0.100000in}{0.100000in}}{\pgfqpoint{3.608454in}{2.310000in}}%
\pgfusepath{clip}%
\pgfsetbuttcap%
\pgfsetmiterjoin%
\definecolor{currentfill}{rgb}{0.000000,0.603922,0.698039}%
\pgfsetfillcolor{currentfill}%
\pgfsetlinewidth{0.000000pt}%
\definecolor{currentstroke}{rgb}{0.000000,0.000000,0.000000}%
\pgfsetstrokecolor{currentstroke}%
\pgfsetstrokeopacity{0.000000}%
\pgfsetdash{}{0pt}%
\pgfpathmoveto{\pgfqpoint{1.291011in}{1.784895in}}%
\pgfpathlineto{\pgfqpoint{1.297605in}{1.774698in}}%
\pgfpathlineto{\pgfqpoint{1.296615in}{1.768471in}}%
\pgfpathlineto{\pgfqpoint{1.310078in}{1.766325in}}%
\pgfpathlineto{\pgfqpoint{1.323149in}{1.760841in}}%
\pgfpathlineto{\pgfqpoint{1.322608in}{1.757434in}}%
\pgfpathlineto{\pgfqpoint{1.342072in}{1.748707in}}%
\pgfpathlineto{\pgfqpoint{1.352534in}{1.748041in}}%
\pgfpathlineto{\pgfqpoint{1.384511in}{1.743611in}}%
\pgfpathlineto{\pgfqpoint{1.388262in}{1.745418in}}%
\pgfpathlineto{\pgfqpoint{1.388966in}{1.743023in}}%
\pgfpathlineto{\pgfqpoint{1.386059in}{1.715853in}}%
\pgfpathlineto{\pgfqpoint{1.382345in}{1.688646in}}%
\pgfpathlineto{\pgfqpoint{1.379922in}{1.688990in}}%
\pgfpathlineto{\pgfqpoint{1.376226in}{1.661744in}}%
\pgfpathlineto{\pgfqpoint{1.377332in}{1.661590in}}%
\pgfpathlineto{\pgfqpoint{1.375467in}{1.648053in}}%
\pgfpathlineto{\pgfqpoint{1.347746in}{1.651996in}}%
\pgfpathlineto{\pgfqpoint{1.287325in}{1.661183in}}%
\pgfpathlineto{\pgfqpoint{1.287915in}{1.674697in}}%
\pgfpathlineto{\pgfqpoint{1.290897in}{1.695072in}}%
\pgfpathlineto{\pgfqpoint{1.282465in}{1.703069in}}%
\pgfpathlineto{\pgfqpoint{1.275945in}{1.717107in}}%
\pgfpathlineto{\pgfqpoint{1.270191in}{1.723778in}}%
\pgfpathlineto{\pgfqpoint{1.264849in}{1.733953in}}%
\pgfpathlineto{\pgfqpoint{1.262997in}{1.742086in}}%
\pgfpathlineto{\pgfqpoint{1.263755in}{1.753309in}}%
\pgfpathlineto{\pgfqpoint{1.261563in}{1.761526in}}%
\pgfpathlineto{\pgfqpoint{1.244733in}{1.764439in}}%
\pgfpathlineto{\pgfqpoint{1.251746in}{1.806946in}}%
\pgfpathlineto{\pgfqpoint{1.254530in}{1.801740in}}%
\pgfpathlineto{\pgfqpoint{1.260477in}{1.800714in}}%
\pgfpathlineto{\pgfqpoint{1.264050in}{1.788274in}}%
\pgfpathlineto{\pgfqpoint{1.271394in}{1.792453in}}%
\pgfpathlineto{\pgfqpoint{1.273821in}{1.796453in}}%
\pgfpathlineto{\pgfqpoint{1.279919in}{1.798565in}}%
\pgfpathlineto{\pgfqpoint{1.285167in}{1.795237in}}%
\pgfpathlineto{\pgfqpoint{1.291011in}{1.784895in}}%
\pgfpathclose%
\pgfusepath{fill}%
\end{pgfscope}%
\begin{pgfscope}%
\pgfpathrectangle{\pgfqpoint{0.100000in}{0.100000in}}{\pgfqpoint{3.608454in}{2.310000in}}%
\pgfusepath{clip}%
\pgfsetbuttcap%
\pgfsetmiterjoin%
\definecolor{currentfill}{rgb}{0.000000,0.588235,0.705882}%
\pgfsetfillcolor{currentfill}%
\pgfsetlinewidth{0.000000pt}%
\definecolor{currentstroke}{rgb}{0.000000,0.000000,0.000000}%
\pgfsetstrokecolor{currentstroke}%
\pgfsetstrokeopacity{0.000000}%
\pgfsetdash{}{0pt}%
\pgfpathmoveto{\pgfqpoint{2.417621in}{1.861430in}}%
\pgfpathlineto{\pgfqpoint{2.417000in}{1.868413in}}%
\pgfpathlineto{\pgfqpoint{2.382793in}{1.866283in}}%
\pgfpathlineto{\pgfqpoint{2.380907in}{1.900334in}}%
\pgfpathlineto{\pgfqpoint{2.387268in}{1.900687in}}%
\pgfpathlineto{\pgfqpoint{2.385839in}{1.925963in}}%
\pgfpathlineto{\pgfqpoint{2.432701in}{1.915926in}}%
\pgfpathlineto{\pgfqpoint{2.438472in}{1.913015in}}%
\pgfpathlineto{\pgfqpoint{2.456551in}{1.904950in}}%
\pgfpathlineto{\pgfqpoint{2.457894in}{1.884450in}}%
\pgfpathlineto{\pgfqpoint{2.471787in}{1.885405in}}%
\pgfpathlineto{\pgfqpoint{2.473620in}{1.857983in}}%
\pgfpathlineto{\pgfqpoint{2.459694in}{1.857127in}}%
\pgfpathlineto{\pgfqpoint{2.445861in}{1.856178in}}%
\pgfpathlineto{\pgfqpoint{2.445567in}{1.863054in}}%
\pgfpathlineto{\pgfqpoint{2.417621in}{1.861430in}}%
\pgfpathclose%
\pgfusepath{fill}%
\end{pgfscope}%
\begin{pgfscope}%
\pgfpathrectangle{\pgfqpoint{0.100000in}{0.100000in}}{\pgfqpoint{3.608454in}{2.310000in}}%
\pgfusepath{clip}%
\pgfsetbuttcap%
\pgfsetmiterjoin%
\definecolor{currentfill}{rgb}{0.000000,0.701961,0.649020}%
\pgfsetfillcolor{currentfill}%
\pgfsetlinewidth{0.000000pt}%
\definecolor{currentstroke}{rgb}{0.000000,0.000000,0.000000}%
\pgfsetstrokecolor{currentstroke}%
\pgfsetstrokeopacity{0.000000}%
\pgfsetdash{}{0pt}%
\pgfpathmoveto{\pgfqpoint{1.126068in}{1.701146in}}%
\pgfpathlineto{\pgfqpoint{1.122487in}{1.688955in}}%
\pgfpathlineto{\pgfqpoint{1.117899in}{1.692151in}}%
\pgfpathlineto{\pgfqpoint{1.109727in}{1.693683in}}%
\pgfpathlineto{\pgfqpoint{1.101946in}{1.704218in}}%
\pgfpathlineto{\pgfqpoint{1.098147in}{1.713450in}}%
\pgfpathlineto{\pgfqpoint{1.091640in}{1.714726in}}%
\pgfpathlineto{\pgfqpoint{1.090342in}{1.707956in}}%
\pgfpathlineto{\pgfqpoint{1.081382in}{1.709710in}}%
\pgfpathlineto{\pgfqpoint{1.080051in}{1.702943in}}%
\pgfpathlineto{\pgfqpoint{1.060120in}{1.706792in}}%
\pgfpathlineto{\pgfqpoint{1.064066in}{1.727041in}}%
\pgfpathlineto{\pgfqpoint{1.054206in}{1.729137in}}%
\pgfpathlineto{\pgfqpoint{1.051274in}{1.732636in}}%
\pgfpathlineto{\pgfqpoint{1.055372in}{1.735866in}}%
\pgfpathlineto{\pgfqpoint{1.053439in}{1.743275in}}%
\pgfpathlineto{\pgfqpoint{1.058573in}{1.770425in}}%
\pgfpathlineto{\pgfqpoint{1.071945in}{1.767715in}}%
\pgfpathlineto{\pgfqpoint{1.074676in}{1.781213in}}%
\pgfpathlineto{\pgfqpoint{1.054623in}{1.785268in}}%
\pgfpathlineto{\pgfqpoint{1.055937in}{1.791680in}}%
\pgfpathlineto{\pgfqpoint{1.040506in}{1.794892in}}%
\pgfpathlineto{\pgfqpoint{1.041894in}{1.801489in}}%
\pgfpathlineto{\pgfqpoint{1.038897in}{1.813232in}}%
\pgfpathlineto{\pgfqpoint{1.034613in}{1.812423in}}%
\pgfpathlineto{\pgfqpoint{1.037674in}{1.816193in}}%
\pgfpathlineto{\pgfqpoint{1.049281in}{1.819620in}}%
\pgfpathlineto{\pgfqpoint{1.057180in}{1.823892in}}%
\pgfpathlineto{\pgfqpoint{1.061051in}{1.840480in}}%
\pgfpathlineto{\pgfqpoint{1.063909in}{1.846064in}}%
\pgfpathlineto{\pgfqpoint{1.066665in}{1.859594in}}%
\pgfpathlineto{\pgfqpoint{1.100225in}{1.852591in}}%
\pgfpathlineto{\pgfqpoint{1.102494in}{1.863989in}}%
\pgfpathlineto{\pgfqpoint{1.105971in}{1.872044in}}%
\pgfpathlineto{\pgfqpoint{1.123023in}{1.868301in}}%
\pgfpathlineto{\pgfqpoint{1.127333in}{1.864844in}}%
\pgfpathlineto{\pgfqpoint{1.130114in}{1.870874in}}%
\pgfpathlineto{\pgfqpoint{1.134588in}{1.873151in}}%
\pgfpathlineto{\pgfqpoint{1.144630in}{1.867633in}}%
\pgfpathlineto{\pgfqpoint{1.157878in}{1.868529in}}%
\pgfpathlineto{\pgfqpoint{1.159714in}{1.863660in}}%
\pgfpathlineto{\pgfqpoint{1.166980in}{1.866417in}}%
\pgfpathlineto{\pgfqpoint{1.177604in}{1.866287in}}%
\pgfpathlineto{\pgfqpoint{1.181983in}{1.875682in}}%
\pgfpathlineto{\pgfqpoint{1.187877in}{1.878039in}}%
\pgfpathlineto{\pgfqpoint{1.191695in}{1.873351in}}%
\pgfpathlineto{\pgfqpoint{1.198018in}{1.855881in}}%
\pgfpathlineto{\pgfqpoint{1.202252in}{1.853038in}}%
\pgfpathlineto{\pgfqpoint{1.195518in}{1.814680in}}%
\pgfpathlineto{\pgfqpoint{1.186131in}{1.812373in}}%
\pgfpathlineto{\pgfqpoint{1.174897in}{1.813546in}}%
\pgfpathlineto{\pgfqpoint{1.170551in}{1.790061in}}%
\pgfpathlineto{\pgfqpoint{1.179359in}{1.788436in}}%
\pgfpathlineto{\pgfqpoint{1.180737in}{1.781578in}}%
\pgfpathlineto{\pgfqpoint{1.188888in}{1.777013in}}%
\pgfpathlineto{\pgfqpoint{1.186356in}{1.762535in}}%
\pgfpathlineto{\pgfqpoint{1.182268in}{1.739464in}}%
\pgfpathlineto{\pgfqpoint{1.124259in}{1.750336in}}%
\pgfpathlineto{\pgfqpoint{1.117477in}{1.741580in}}%
\pgfpathlineto{\pgfqpoint{1.116145in}{1.733195in}}%
\pgfpathlineto{\pgfqpoint{1.117055in}{1.725647in}}%
\pgfpathlineto{\pgfqpoint{1.125111in}{1.725638in}}%
\pgfpathlineto{\pgfqpoint{1.126636in}{1.717223in}}%
\pgfpathlineto{\pgfqpoint{1.126068in}{1.701146in}}%
\pgfpathclose%
\pgfusepath{fill}%
\end{pgfscope}%
\begin{pgfscope}%
\pgfpathrectangle{\pgfqpoint{0.100000in}{0.100000in}}{\pgfqpoint{3.608454in}{2.310000in}}%
\pgfusepath{clip}%
\pgfsetbuttcap%
\pgfsetmiterjoin%
\definecolor{currentfill}{rgb}{0.000000,0.701961,0.649020}%
\pgfsetfillcolor{currentfill}%
\pgfsetlinewidth{0.000000pt}%
\definecolor{currentstroke}{rgb}{0.000000,0.000000,0.000000}%
\pgfsetstrokecolor{currentstroke}%
\pgfsetstrokeopacity{0.000000}%
\pgfsetdash{}{0pt}%
\pgfpathmoveto{\pgfqpoint{0.718104in}{2.222638in}}%
\pgfpathlineto{\pgfqpoint{0.722390in}{2.214825in}}%
\pgfpathlineto{\pgfqpoint{0.724213in}{2.207197in}}%
\pgfpathlineto{\pgfqpoint{0.733242in}{2.201231in}}%
\pgfpathlineto{\pgfqpoint{0.735660in}{2.196356in}}%
\pgfpathlineto{\pgfqpoint{0.742925in}{2.192347in}}%
\pgfpathlineto{\pgfqpoint{0.748201in}{2.185881in}}%
\pgfpathlineto{\pgfqpoint{0.763587in}{2.181602in}}%
\pgfpathlineto{\pgfqpoint{0.767233in}{2.177075in}}%
\pgfpathlineto{\pgfqpoint{0.762473in}{2.166902in}}%
\pgfpathlineto{\pgfqpoint{0.763393in}{2.155197in}}%
\pgfpathlineto{\pgfqpoint{0.762232in}{2.143437in}}%
\pgfpathlineto{\pgfqpoint{0.758684in}{2.134138in}}%
\pgfpathlineto{\pgfqpoint{0.761418in}{2.129669in}}%
\pgfpathlineto{\pgfqpoint{0.749167in}{2.084439in}}%
\pgfpathlineto{\pgfqpoint{0.719253in}{2.092888in}}%
\pgfpathlineto{\pgfqpoint{0.661241in}{2.110265in}}%
\pgfpathlineto{\pgfqpoint{0.669147in}{2.136525in}}%
\pgfpathlineto{\pgfqpoint{0.675514in}{2.134601in}}%
\pgfpathlineto{\pgfqpoint{0.675711in}{2.144078in}}%
\pgfpathlineto{\pgfqpoint{0.682016in}{2.153295in}}%
\pgfpathlineto{\pgfqpoint{0.682663in}{2.166183in}}%
\pgfpathlineto{\pgfqpoint{0.680298in}{2.173569in}}%
\pgfpathlineto{\pgfqpoint{0.697139in}{2.190297in}}%
\pgfpathlineto{\pgfqpoint{0.699138in}{2.201963in}}%
\pgfpathlineto{\pgfqpoint{0.694658in}{2.204301in}}%
\pgfpathlineto{\pgfqpoint{0.694418in}{2.209381in}}%
\pgfpathlineto{\pgfqpoint{0.701083in}{2.215781in}}%
\pgfpathlineto{\pgfqpoint{0.712161in}{2.222097in}}%
\pgfpathlineto{\pgfqpoint{0.718104in}{2.222638in}}%
\pgfpathclose%
\pgfusepath{fill}%
\end{pgfscope}%
\begin{pgfscope}%
\pgfpathrectangle{\pgfqpoint{0.100000in}{0.100000in}}{\pgfqpoint{3.608454in}{2.310000in}}%
\pgfusepath{clip}%
\pgfsetbuttcap%
\pgfsetmiterjoin%
\definecolor{currentfill}{rgb}{0.000000,0.776471,0.611765}%
\pgfsetfillcolor{currentfill}%
\pgfsetlinewidth{0.000000pt}%
\definecolor{currentstroke}{rgb}{0.000000,0.000000,0.000000}%
\pgfsetstrokecolor{currentstroke}%
\pgfsetstrokeopacity{0.000000}%
\pgfsetdash{}{0pt}%
\pgfpathmoveto{\pgfqpoint{0.538914in}{1.596597in}}%
\pgfpathlineto{\pgfqpoint{0.532464in}{1.597564in}}%
\pgfpathlineto{\pgfqpoint{0.520214in}{1.588877in}}%
\pgfpathlineto{\pgfqpoint{0.514461in}{1.575916in}}%
\pgfpathlineto{\pgfqpoint{0.507112in}{1.573994in}}%
\pgfpathlineto{\pgfqpoint{0.502169in}{1.569148in}}%
\pgfpathlineto{\pgfqpoint{0.498193in}{1.555922in}}%
\pgfpathlineto{\pgfqpoint{0.489967in}{1.557121in}}%
\pgfpathlineto{\pgfqpoint{0.484414in}{1.564406in}}%
\pgfpathlineto{\pgfqpoint{0.485256in}{1.568714in}}%
\pgfpathlineto{\pgfqpoint{0.480562in}{1.575547in}}%
\pgfpathlineto{\pgfqpoint{0.450911in}{1.584637in}}%
\pgfpathlineto{\pgfqpoint{0.448894in}{1.592474in}}%
\pgfpathlineto{\pgfqpoint{0.445039in}{1.597163in}}%
\pgfpathlineto{\pgfqpoint{0.447344in}{1.609339in}}%
\pgfpathlineto{\pgfqpoint{0.440654in}{1.612562in}}%
\pgfpathlineto{\pgfqpoint{0.435002in}{1.619433in}}%
\pgfpathlineto{\pgfqpoint{0.440384in}{1.631531in}}%
\pgfpathlineto{\pgfqpoint{0.443292in}{1.641381in}}%
\pgfpathlineto{\pgfqpoint{0.434583in}{1.644086in}}%
\pgfpathlineto{\pgfqpoint{0.437233in}{1.653925in}}%
\pgfpathlineto{\pgfqpoint{0.436761in}{1.661477in}}%
\pgfpathlineto{\pgfqpoint{0.488413in}{1.645315in}}%
\pgfpathlineto{\pgfqpoint{0.490498in}{1.651799in}}%
\pgfpathlineto{\pgfqpoint{0.504486in}{1.647574in}}%
\pgfpathlineto{\pgfqpoint{0.512624in}{1.653371in}}%
\pgfpathlineto{\pgfqpoint{0.515901in}{1.652222in}}%
\pgfpathlineto{\pgfqpoint{0.522050in}{1.660114in}}%
\pgfpathlineto{\pgfqpoint{0.531886in}{1.661402in}}%
\pgfpathlineto{\pgfqpoint{0.534394in}{1.657082in}}%
\pgfpathlineto{\pgfqpoint{0.529796in}{1.650793in}}%
\pgfpathlineto{\pgfqpoint{0.526964in}{1.639730in}}%
\pgfpathlineto{\pgfqpoint{0.529553in}{1.637642in}}%
\pgfpathlineto{\pgfqpoint{0.540950in}{1.606159in}}%
\pgfpathlineto{\pgfqpoint{0.538914in}{1.596597in}}%
\pgfpathclose%
\pgfusepath{fill}%
\end{pgfscope}%
\begin{pgfscope}%
\pgfpathrectangle{\pgfqpoint{0.100000in}{0.100000in}}{\pgfqpoint{3.608454in}{2.310000in}}%
\pgfusepath{clip}%
\pgfsetbuttcap%
\pgfsetmiterjoin%
\definecolor{currentfill}{rgb}{0.000000,0.596078,0.701961}%
\pgfsetfillcolor{currentfill}%
\pgfsetlinewidth{0.000000pt}%
\definecolor{currentstroke}{rgb}{0.000000,0.000000,0.000000}%
\pgfsetstrokecolor{currentstroke}%
\pgfsetstrokeopacity{0.000000}%
\pgfsetdash{}{0pt}%
\pgfpathmoveto{\pgfqpoint{2.075326in}{1.295024in}}%
\pgfpathlineto{\pgfqpoint{2.074465in}{1.263717in}}%
\pgfpathlineto{\pgfqpoint{2.047067in}{1.264108in}}%
\pgfpathlineto{\pgfqpoint{2.047190in}{1.274443in}}%
\pgfpathlineto{\pgfqpoint{2.022450in}{1.274884in}}%
\pgfpathlineto{\pgfqpoint{2.022297in}{1.268017in}}%
\pgfpathlineto{\pgfqpoint{2.012052in}{1.268224in}}%
\pgfpathlineto{\pgfqpoint{1.992256in}{1.268523in}}%
\pgfpathlineto{\pgfqpoint{1.992866in}{1.289314in}}%
\pgfpathlineto{\pgfqpoint{1.994270in}{1.303029in}}%
\pgfpathlineto{\pgfqpoint{2.023077in}{1.302462in}}%
\pgfpathlineto{\pgfqpoint{2.023407in}{1.319659in}}%
\pgfpathlineto{\pgfqpoint{2.048478in}{1.319317in}}%
\pgfpathlineto{\pgfqpoint{2.047953in}{1.295247in}}%
\pgfpathlineto{\pgfqpoint{2.075326in}{1.295024in}}%
\pgfpathclose%
\pgfusepath{fill}%
\end{pgfscope}%
\begin{pgfscope}%
\pgfpathrectangle{\pgfqpoint{0.100000in}{0.100000in}}{\pgfqpoint{3.608454in}{2.310000in}}%
\pgfusepath{clip}%
\pgfsetbuttcap%
\pgfsetmiterjoin%
\definecolor{currentfill}{rgb}{0.000000,0.462745,0.768627}%
\pgfsetfillcolor{currentfill}%
\pgfsetlinewidth{0.000000pt}%
\definecolor{currentstroke}{rgb}{0.000000,0.000000,0.000000}%
\pgfsetstrokecolor{currentstroke}%
\pgfsetstrokeopacity{0.000000}%
\pgfsetdash{}{0pt}%
\pgfpathmoveto{\pgfqpoint{1.303565in}{2.014897in}}%
\pgfpathlineto{\pgfqpoint{1.299591in}{2.015508in}}%
\pgfpathlineto{\pgfqpoint{1.292244in}{2.021526in}}%
\pgfpathlineto{\pgfqpoint{1.264971in}{2.020151in}}%
\pgfpathlineto{\pgfqpoint{1.259993in}{2.025470in}}%
\pgfpathlineto{\pgfqpoint{1.251605in}{2.040078in}}%
\pgfpathlineto{\pgfqpoint{1.244581in}{2.048177in}}%
\pgfpathlineto{\pgfqpoint{1.231866in}{2.052930in}}%
\pgfpathlineto{\pgfqpoint{1.211322in}{2.055765in}}%
\pgfpathlineto{\pgfqpoint{1.210055in}{2.048940in}}%
\pgfpathlineto{\pgfqpoint{1.203686in}{2.050112in}}%
\pgfpathlineto{\pgfqpoint{1.196622in}{2.051472in}}%
\pgfpathlineto{\pgfqpoint{1.199816in}{2.068342in}}%
\pgfpathlineto{\pgfqpoint{1.192015in}{2.070842in}}%
\pgfpathlineto{\pgfqpoint{1.187119in}{2.075916in}}%
\pgfpathlineto{\pgfqpoint{1.191810in}{2.101317in}}%
\pgfpathlineto{\pgfqpoint{1.179403in}{2.102869in}}%
\pgfpathlineto{\pgfqpoint{1.173331in}{2.107985in}}%
\pgfpathlineto{\pgfqpoint{1.173281in}{2.111946in}}%
\pgfpathlineto{\pgfqpoint{1.165476in}{2.115944in}}%
\pgfpathlineto{\pgfqpoint{1.152365in}{2.116451in}}%
\pgfpathlineto{\pgfqpoint{1.148672in}{2.124295in}}%
\pgfpathlineto{\pgfqpoint{1.150853in}{2.134585in}}%
\pgfpathlineto{\pgfqpoint{1.149026in}{2.145360in}}%
\pgfpathlineto{\pgfqpoint{1.154266in}{2.150464in}}%
\pgfpathlineto{\pgfqpoint{1.150149in}{2.159708in}}%
\pgfpathlineto{\pgfqpoint{1.194140in}{2.150958in}}%
\pgfpathlineto{\pgfqpoint{1.200330in}{2.146187in}}%
\pgfpathlineto{\pgfqpoint{1.202148in}{2.137522in}}%
\pgfpathlineto{\pgfqpoint{1.232482in}{2.132070in}}%
\pgfpathlineto{\pgfqpoint{1.234540in}{2.143503in}}%
\pgfpathlineto{\pgfqpoint{1.265169in}{2.138097in}}%
\pgfpathlineto{\pgfqpoint{1.266372in}{2.144919in}}%
\pgfpathlineto{\pgfqpoint{1.276965in}{2.143064in}}%
\pgfpathlineto{\pgfqpoint{1.278167in}{2.149917in}}%
\pgfpathlineto{\pgfqpoint{1.319123in}{2.142944in}}%
\pgfpathlineto{\pgfqpoint{1.318309in}{2.136230in}}%
\pgfpathlineto{\pgfqpoint{1.324278in}{2.135234in}}%
\pgfpathlineto{\pgfqpoint{1.323124in}{2.128216in}}%
\pgfpathlineto{\pgfqpoint{1.333325in}{2.126759in}}%
\pgfpathlineto{\pgfqpoint{1.328054in}{2.095490in}}%
\pgfpathlineto{\pgfqpoint{1.312471in}{2.096193in}}%
\pgfpathlineto{\pgfqpoint{1.309109in}{2.082641in}}%
\pgfpathlineto{\pgfqpoint{1.293100in}{2.076361in}}%
\pgfpathlineto{\pgfqpoint{1.290585in}{2.063837in}}%
\pgfpathlineto{\pgfqpoint{1.296843in}{2.059295in}}%
\pgfpathlineto{\pgfqpoint{1.303684in}{2.058169in}}%
\pgfpathlineto{\pgfqpoint{1.309306in}{2.053731in}}%
\pgfpathlineto{\pgfqpoint{1.305651in}{2.030921in}}%
\pgfpathlineto{\pgfqpoint{1.303017in}{2.028995in}}%
\pgfpathlineto{\pgfqpoint{1.303565in}{2.014897in}}%
\pgfpathclose%
\pgfusepath{fill}%
\end{pgfscope}%
\begin{pgfscope}%
\pgfpathrectangle{\pgfqpoint{0.100000in}{0.100000in}}{\pgfqpoint{3.608454in}{2.310000in}}%
\pgfusepath{clip}%
\pgfsetbuttcap%
\pgfsetmiterjoin%
\definecolor{currentfill}{rgb}{0.000000,0.462745,0.768627}%
\pgfsetfillcolor{currentfill}%
\pgfsetlinewidth{0.000000pt}%
\definecolor{currentstroke}{rgb}{0.000000,0.000000,0.000000}%
\pgfsetstrokecolor{currentstroke}%
\pgfsetstrokeopacity{0.000000}%
\pgfsetdash{}{0pt}%
\pgfpathmoveto{\pgfqpoint{1.807069in}{2.050968in}}%
\pgfpathlineto{\pgfqpoint{1.827828in}{2.049744in}}%
\pgfpathlineto{\pgfqpoint{1.826539in}{2.063794in}}%
\pgfpathlineto{\pgfqpoint{1.828130in}{2.091731in}}%
\pgfpathlineto{\pgfqpoint{1.826839in}{2.105738in}}%
\pgfpathlineto{\pgfqpoint{1.833772in}{2.105376in}}%
\pgfpathlineto{\pgfqpoint{1.868375in}{2.103508in}}%
\pgfpathlineto{\pgfqpoint{1.867703in}{2.089585in}}%
\pgfpathlineto{\pgfqpoint{1.848899in}{2.090580in}}%
\pgfpathlineto{\pgfqpoint{1.847326in}{2.062601in}}%
\pgfpathlineto{\pgfqpoint{1.848523in}{2.048601in}}%
\pgfpathlineto{\pgfqpoint{1.876114in}{2.047161in}}%
\pgfpathlineto{\pgfqpoint{1.877148in}{2.033168in}}%
\pgfpathlineto{\pgfqpoint{1.875835in}{2.005576in}}%
\pgfpathlineto{\pgfqpoint{1.864232in}{2.006135in}}%
\pgfpathlineto{\pgfqpoint{1.830021in}{2.007934in}}%
\pgfpathlineto{\pgfqpoint{1.799746in}{2.009759in}}%
\pgfpathlineto{\pgfqpoint{1.801529in}{2.037408in}}%
\pgfpathlineto{\pgfqpoint{1.806219in}{2.037109in}}%
\pgfpathlineto{\pgfqpoint{1.807069in}{2.050968in}}%
\pgfpathclose%
\pgfusepath{fill}%
\end{pgfscope}%
\begin{pgfscope}%
\pgfpathrectangle{\pgfqpoint{0.100000in}{0.100000in}}{\pgfqpoint{3.608454in}{2.310000in}}%
\pgfusepath{clip}%
\pgfsetbuttcap%
\pgfsetmiterjoin%
\definecolor{currentfill}{rgb}{0.000000,0.580392,0.709804}%
\pgfsetfillcolor{currentfill}%
\pgfsetlinewidth{0.000000pt}%
\definecolor{currentstroke}{rgb}{0.000000,0.000000,0.000000}%
\pgfsetstrokecolor{currentstroke}%
\pgfsetstrokeopacity{0.000000}%
\pgfsetdash{}{0pt}%
\pgfpathmoveto{\pgfqpoint{2.409257in}{1.435587in}}%
\pgfpathlineto{\pgfqpoint{2.397384in}{1.430530in}}%
\pgfpathlineto{\pgfqpoint{2.391369in}{1.431061in}}%
\pgfpathlineto{\pgfqpoint{2.389932in}{1.435148in}}%
\pgfpathlineto{\pgfqpoint{2.396575in}{1.440884in}}%
\pgfpathlineto{\pgfqpoint{2.401121in}{1.445010in}}%
\pgfpathlineto{\pgfqpoint{2.405497in}{1.456494in}}%
\pgfpathlineto{\pgfqpoint{2.412968in}{1.462469in}}%
\pgfpathlineto{\pgfqpoint{2.414833in}{1.467947in}}%
\pgfpathlineto{\pgfqpoint{2.414274in}{1.476742in}}%
\pgfpathlineto{\pgfqpoint{2.407342in}{1.476448in}}%
\pgfpathlineto{\pgfqpoint{2.406029in}{1.504012in}}%
\pgfpathlineto{\pgfqpoint{2.426710in}{1.505121in}}%
\pgfpathlineto{\pgfqpoint{2.425489in}{1.525655in}}%
\pgfpathlineto{\pgfqpoint{2.412483in}{1.524944in}}%
\pgfpathlineto{\pgfqpoint{2.410650in}{1.552521in}}%
\pgfpathlineto{\pgfqpoint{2.451750in}{1.555104in}}%
\pgfpathlineto{\pgfqpoint{2.451488in}{1.558550in}}%
\pgfpathlineto{\pgfqpoint{2.484797in}{1.561007in}}%
\pgfpathlineto{\pgfqpoint{2.488674in}{1.519766in}}%
\pgfpathlineto{\pgfqpoint{2.468185in}{1.518191in}}%
\pgfpathlineto{\pgfqpoint{2.470141in}{1.490384in}}%
\pgfpathlineto{\pgfqpoint{2.467299in}{1.483195in}}%
\pgfpathlineto{\pgfqpoint{2.458726in}{1.477228in}}%
\pgfpathlineto{\pgfqpoint{2.450616in}{1.476514in}}%
\pgfpathlineto{\pgfqpoint{2.452287in}{1.455302in}}%
\pgfpathlineto{\pgfqpoint{2.431914in}{1.453643in}}%
\pgfpathlineto{\pgfqpoint{2.432920in}{1.438056in}}%
\pgfpathlineto{\pgfqpoint{2.427687in}{1.440218in}}%
\pgfpathlineto{\pgfqpoint{2.418888in}{1.437396in}}%
\pgfpathlineto{\pgfqpoint{2.410568in}{1.438329in}}%
\pgfpathlineto{\pgfqpoint{2.409257in}{1.435587in}}%
\pgfpathclose%
\pgfusepath{fill}%
\end{pgfscope}%
\begin{pgfscope}%
\pgfpathrectangle{\pgfqpoint{0.100000in}{0.100000in}}{\pgfqpoint{3.608454in}{2.310000in}}%
\pgfusepath{clip}%
\pgfsetbuttcap%
\pgfsetmiterjoin%
\definecolor{currentfill}{rgb}{0.000000,0.560784,0.719608}%
\pgfsetfillcolor{currentfill}%
\pgfsetlinewidth{0.000000pt}%
\definecolor{currentstroke}{rgb}{0.000000,0.000000,0.000000}%
\pgfsetstrokecolor{currentstroke}%
\pgfsetstrokeopacity{0.000000}%
\pgfsetdash{}{0pt}%
\pgfpathmoveto{\pgfqpoint{1.843805in}{1.813366in}}%
\pgfpathlineto{\pgfqpoint{1.849510in}{1.813058in}}%
\pgfpathlineto{\pgfqpoint{1.898285in}{1.810809in}}%
\pgfpathlineto{\pgfqpoint{1.897415in}{1.789939in}}%
\pgfpathlineto{\pgfqpoint{1.896283in}{1.755231in}}%
\pgfpathlineto{\pgfqpoint{1.883458in}{1.755759in}}%
\pgfpathlineto{\pgfqpoint{1.846482in}{1.757163in}}%
\pgfpathlineto{\pgfqpoint{1.841646in}{1.759396in}}%
\pgfpathlineto{\pgfqpoint{1.842265in}{1.785638in}}%
\pgfpathlineto{\pgfqpoint{1.843805in}{1.813366in}}%
\pgfpathclose%
\pgfusepath{fill}%
\end{pgfscope}%
\begin{pgfscope}%
\pgfpathrectangle{\pgfqpoint{0.100000in}{0.100000in}}{\pgfqpoint{3.608454in}{2.310000in}}%
\pgfusepath{clip}%
\pgfsetbuttcap%
\pgfsetmiterjoin%
\definecolor{currentfill}{rgb}{0.000000,0.796078,0.601961}%
\pgfsetfillcolor{currentfill}%
\pgfsetlinewidth{0.000000pt}%
\definecolor{currentstroke}{rgb}{0.000000,0.000000,0.000000}%
\pgfsetstrokecolor{currentstroke}%
\pgfsetstrokeopacity{0.000000}%
\pgfsetdash{}{0pt}%
\pgfpathmoveto{\pgfqpoint{2.336789in}{0.709454in}}%
\pgfpathlineto{\pgfqpoint{2.341562in}{0.715016in}}%
\pgfpathlineto{\pgfqpoint{2.337261in}{0.720197in}}%
\pgfpathlineto{\pgfqpoint{2.341299in}{0.728481in}}%
\pgfpathlineto{\pgfqpoint{2.346226in}{0.727949in}}%
\pgfpathlineto{\pgfqpoint{2.347315in}{0.735413in}}%
\pgfpathlineto{\pgfqpoint{2.362838in}{0.740120in}}%
\pgfpathlineto{\pgfqpoint{2.371609in}{0.736436in}}%
\pgfpathlineto{\pgfqpoint{2.377876in}{0.738907in}}%
\pgfpathlineto{\pgfqpoint{2.404037in}{0.740243in}}%
\pgfpathlineto{\pgfqpoint{2.439716in}{0.742196in}}%
\pgfpathlineto{\pgfqpoint{2.444160in}{0.745936in}}%
\pgfpathlineto{\pgfqpoint{2.449720in}{0.745908in}}%
\pgfpathlineto{\pgfqpoint{2.458936in}{0.746389in}}%
\pgfpathlineto{\pgfqpoint{2.458660in}{0.749893in}}%
\pgfpathlineto{\pgfqpoint{2.470222in}{0.750540in}}%
\pgfpathlineto{\pgfqpoint{2.472257in}{0.716197in}}%
\pgfpathlineto{\pgfqpoint{2.467199in}{0.715885in}}%
\pgfpathlineto{\pgfqpoint{2.424922in}{0.713419in}}%
\pgfpathlineto{\pgfqpoint{2.392212in}{0.711750in}}%
\pgfpathlineto{\pgfqpoint{2.336789in}{0.709454in}}%
\pgfpathclose%
\pgfusepath{fill}%
\end{pgfscope}%
\begin{pgfscope}%
\pgfpathrectangle{\pgfqpoint{0.100000in}{0.100000in}}{\pgfqpoint{3.608454in}{2.310000in}}%
\pgfusepath{clip}%
\pgfsetbuttcap%
\pgfsetmiterjoin%
\definecolor{currentfill}{rgb}{0.000000,0.439216,0.780392}%
\pgfsetfillcolor{currentfill}%
\pgfsetlinewidth{0.000000pt}%
\definecolor{currentstroke}{rgb}{0.000000,0.000000,0.000000}%
\pgfsetstrokecolor{currentstroke}%
\pgfsetstrokeopacity{0.000000}%
\pgfsetdash{}{0pt}%
\pgfpathmoveto{\pgfqpoint{2.221382in}{1.649990in}}%
\pgfpathlineto{\pgfqpoint{2.220391in}{1.696881in}}%
\pgfpathlineto{\pgfqpoint{2.253472in}{1.697776in}}%
\pgfpathlineto{\pgfqpoint{2.294773in}{1.699116in}}%
\pgfpathlineto{\pgfqpoint{2.301639in}{1.699356in}}%
\pgfpathlineto{\pgfqpoint{2.303199in}{1.666191in}}%
\pgfpathlineto{\pgfqpoint{2.275677in}{1.665325in}}%
\pgfpathlineto{\pgfqpoint{2.276089in}{1.651404in}}%
\pgfpathlineto{\pgfqpoint{2.221382in}{1.649990in}}%
\pgfpathclose%
\pgfusepath{fill}%
\end{pgfscope}%
\begin{pgfscope}%
\pgfpathrectangle{\pgfqpoint{0.100000in}{0.100000in}}{\pgfqpoint{3.608454in}{2.310000in}}%
\pgfusepath{clip}%
\pgfsetbuttcap%
\pgfsetmiterjoin%
\definecolor{currentfill}{rgb}{0.000000,0.529412,0.735294}%
\pgfsetfillcolor{currentfill}%
\pgfsetlinewidth{0.000000pt}%
\definecolor{currentstroke}{rgb}{0.000000,0.000000,0.000000}%
\pgfsetstrokecolor{currentstroke}%
\pgfsetstrokeopacity{0.000000}%
\pgfsetdash{}{0pt}%
\pgfpathmoveto{\pgfqpoint{3.109751in}{1.709791in}}%
\pgfpathlineto{\pgfqpoint{3.102306in}{1.712116in}}%
\pgfpathlineto{\pgfqpoint{3.066495in}{1.705513in}}%
\pgfpathlineto{\pgfqpoint{3.064906in}{1.706470in}}%
\pgfpathlineto{\pgfqpoint{3.049340in}{1.697252in}}%
\pgfpathlineto{\pgfqpoint{3.042918in}{1.694827in}}%
\pgfpathlineto{\pgfqpoint{3.036017in}{1.700877in}}%
\pgfpathlineto{\pgfqpoint{3.032000in}{1.700788in}}%
\pgfpathlineto{\pgfqpoint{3.027192in}{1.702502in}}%
\pgfpathlineto{\pgfqpoint{3.030633in}{1.712923in}}%
\pgfpathlineto{\pgfqpoint{3.037383in}{1.717726in}}%
\pgfpathlineto{\pgfqpoint{3.040436in}{1.722640in}}%
\pgfpathlineto{\pgfqpoint{3.033554in}{1.734692in}}%
\pgfpathlineto{\pgfqpoint{3.028625in}{1.736179in}}%
\pgfpathlineto{\pgfqpoint{3.027760in}{1.742504in}}%
\pgfpathlineto{\pgfqpoint{3.022366in}{1.756709in}}%
\pgfpathlineto{\pgfqpoint{3.033985in}{1.763896in}}%
\pgfpathlineto{\pgfqpoint{3.044744in}{1.769031in}}%
\pgfpathlineto{\pgfqpoint{3.058321in}{1.773021in}}%
\pgfpathlineto{\pgfqpoint{3.081941in}{1.776783in}}%
\pgfpathlineto{\pgfqpoint{3.094629in}{1.776945in}}%
\pgfpathlineto{\pgfqpoint{3.106459in}{1.771415in}}%
\pgfpathlineto{\pgfqpoint{3.117744in}{1.776426in}}%
\pgfpathlineto{\pgfqpoint{3.141833in}{1.780817in}}%
\pgfpathlineto{\pgfqpoint{3.153298in}{1.788862in}}%
\pgfpathlineto{\pgfqpoint{3.158574in}{1.768645in}}%
\pgfpathlineto{\pgfqpoint{3.158361in}{1.759713in}}%
\pgfpathlineto{\pgfqpoint{3.162874in}{1.739825in}}%
\pgfpathlineto{\pgfqpoint{3.167852in}{1.733612in}}%
\pgfpathlineto{\pgfqpoint{3.173622in}{1.728849in}}%
\pgfpathlineto{\pgfqpoint{3.155225in}{1.724423in}}%
\pgfpathlineto{\pgfqpoint{3.156855in}{1.718354in}}%
\pgfpathlineto{\pgfqpoint{3.142173in}{1.718464in}}%
\pgfpathlineto{\pgfqpoint{3.141312in}{1.724493in}}%
\pgfpathlineto{\pgfqpoint{3.111728in}{1.718970in}}%
\pgfpathlineto{\pgfqpoint{3.109751in}{1.709791in}}%
\pgfpathclose%
\pgfusepath{fill}%
\end{pgfscope}%
\begin{pgfscope}%
\pgfpathrectangle{\pgfqpoint{0.100000in}{0.100000in}}{\pgfqpoint{3.608454in}{2.310000in}}%
\pgfusepath{clip}%
\pgfsetbuttcap%
\pgfsetmiterjoin%
\definecolor{currentfill}{rgb}{0.000000,0.466667,0.766667}%
\pgfsetfillcolor{currentfill}%
\pgfsetlinewidth{0.000000pt}%
\definecolor{currentstroke}{rgb}{0.000000,0.000000,0.000000}%
\pgfsetstrokecolor{currentstroke}%
\pgfsetstrokeopacity{0.000000}%
\pgfsetdash{}{0pt}%
\pgfpathmoveto{\pgfqpoint{2.989320in}{1.467586in}}%
\pgfpathlineto{\pgfqpoint{2.983438in}{1.466654in}}%
\pgfpathlineto{\pgfqpoint{2.977955in}{1.500863in}}%
\pgfpathlineto{\pgfqpoint{2.971951in}{1.538265in}}%
\pgfpathlineto{\pgfqpoint{2.969282in}{1.554858in}}%
\pgfpathlineto{\pgfqpoint{2.990516in}{1.558619in}}%
\pgfpathlineto{\pgfqpoint{2.988249in}{1.569893in}}%
\pgfpathlineto{\pgfqpoint{2.995768in}{1.584878in}}%
\pgfpathlineto{\pgfqpoint{3.013646in}{1.587974in}}%
\pgfpathlineto{\pgfqpoint{3.020349in}{1.574020in}}%
\pgfpathlineto{\pgfqpoint{3.028102in}{1.574866in}}%
\pgfpathlineto{\pgfqpoint{3.040035in}{1.579416in}}%
\pgfpathlineto{\pgfqpoint{3.043474in}{1.583306in}}%
\pgfpathlineto{\pgfqpoint{3.045363in}{1.572457in}}%
\pgfpathlineto{\pgfqpoint{3.069571in}{1.576329in}}%
\pgfpathlineto{\pgfqpoint{3.072318in}{1.562246in}}%
\pgfpathlineto{\pgfqpoint{3.068856in}{1.543568in}}%
\pgfpathlineto{\pgfqpoint{3.061504in}{1.519627in}}%
\pgfpathlineto{\pgfqpoint{3.054176in}{1.507940in}}%
\pgfpathlineto{\pgfqpoint{3.048216in}{1.491838in}}%
\pgfpathlineto{\pgfqpoint{3.051273in}{1.485944in}}%
\pgfpathlineto{\pgfqpoint{3.051009in}{1.477958in}}%
\pgfpathlineto{\pgfqpoint{3.045971in}{1.477046in}}%
\pgfpathlineto{\pgfqpoint{2.989320in}{1.467586in}}%
\pgfpathclose%
\pgfusepath{fill}%
\end{pgfscope}%
\begin{pgfscope}%
\pgfpathrectangle{\pgfqpoint{0.100000in}{0.100000in}}{\pgfqpoint{3.608454in}{2.310000in}}%
\pgfusepath{clip}%
\pgfsetbuttcap%
\pgfsetmiterjoin%
\definecolor{currentfill}{rgb}{0.000000,0.596078,0.701961}%
\pgfsetfillcolor{currentfill}%
\pgfsetlinewidth{0.000000pt}%
\definecolor{currentstroke}{rgb}{0.000000,0.000000,0.000000}%
\pgfsetstrokecolor{currentstroke}%
\pgfsetstrokeopacity{0.000000}%
\pgfsetdash{}{0pt}%
\pgfpathmoveto{\pgfqpoint{1.787654in}{0.626021in}}%
\pgfpathlineto{\pgfqpoint{1.781527in}{0.626427in}}%
\pgfpathlineto{\pgfqpoint{1.782309in}{0.639635in}}%
\pgfpathlineto{\pgfqpoint{1.777624in}{0.639888in}}%
\pgfpathlineto{\pgfqpoint{1.778772in}{0.657165in}}%
\pgfpathlineto{\pgfqpoint{1.809932in}{0.655294in}}%
\pgfpathlineto{\pgfqpoint{1.810600in}{0.672318in}}%
\pgfpathlineto{\pgfqpoint{1.833962in}{0.671147in}}%
\pgfpathlineto{\pgfqpoint{1.859593in}{0.670192in}}%
\pgfpathlineto{\pgfqpoint{1.858713in}{0.641352in}}%
\pgfpathlineto{\pgfqpoint{1.835751in}{0.642247in}}%
\pgfpathlineto{\pgfqpoint{1.834695in}{0.613722in}}%
\pgfpathlineto{\pgfqpoint{1.817099in}{0.623588in}}%
\pgfpathlineto{\pgfqpoint{1.787654in}{0.626021in}}%
\pgfpathclose%
\pgfusepath{fill}%
\end{pgfscope}%
\begin{pgfscope}%
\pgfpathrectangle{\pgfqpoint{0.100000in}{0.100000in}}{\pgfqpoint{3.608454in}{2.310000in}}%
\pgfusepath{clip}%
\pgfsetbuttcap%
\pgfsetmiterjoin%
\definecolor{currentfill}{rgb}{0.000000,0.619608,0.690196}%
\pgfsetfillcolor{currentfill}%
\pgfsetlinewidth{0.000000pt}%
\definecolor{currentstroke}{rgb}{0.000000,0.000000,0.000000}%
\pgfsetstrokecolor{currentstroke}%
\pgfsetstrokeopacity{0.000000}%
\pgfsetdash{}{0pt}%
\pgfpathmoveto{\pgfqpoint{2.870396in}{0.968762in}}%
\pgfpathlineto{\pgfqpoint{2.877953in}{0.974806in}}%
\pgfpathlineto{\pgfqpoint{2.889208in}{0.971223in}}%
\pgfpathlineto{\pgfqpoint{2.896071in}{0.971539in}}%
\pgfpathlineto{\pgfqpoint{2.896953in}{0.978821in}}%
\pgfpathlineto{\pgfqpoint{2.902891in}{0.990956in}}%
\pgfpathlineto{\pgfqpoint{2.908530in}{0.995600in}}%
\pgfpathlineto{\pgfqpoint{2.916400in}{0.998276in}}%
\pgfpathlineto{\pgfqpoint{2.922696in}{0.996327in}}%
\pgfpathlineto{\pgfqpoint{2.941645in}{0.986529in}}%
\pgfpathlineto{\pgfqpoint{2.948243in}{0.978295in}}%
\pgfpathlineto{\pgfqpoint{2.942345in}{0.978613in}}%
\pgfpathlineto{\pgfqpoint{2.932420in}{0.972207in}}%
\pgfpathlineto{\pgfqpoint{2.922909in}{0.970556in}}%
\pgfpathlineto{\pgfqpoint{2.918991in}{0.967320in}}%
\pgfpathlineto{\pgfqpoint{2.915431in}{0.959683in}}%
\pgfpathlineto{\pgfqpoint{2.910547in}{0.958327in}}%
\pgfpathlineto{\pgfqpoint{2.909363in}{0.953452in}}%
\pgfpathlineto{\pgfqpoint{2.898524in}{0.954072in}}%
\pgfpathlineto{\pgfqpoint{2.889760in}{0.943778in}}%
\pgfpathlineto{\pgfqpoint{2.880855in}{0.952903in}}%
\pgfpathlineto{\pgfqpoint{2.870396in}{0.968762in}}%
\pgfpathclose%
\pgfusepath{fill}%
\end{pgfscope}%
\begin{pgfscope}%
\pgfpathrectangle{\pgfqpoint{0.100000in}{0.100000in}}{\pgfqpoint{3.608454in}{2.310000in}}%
\pgfusepath{clip}%
\pgfsetbuttcap%
\pgfsetmiterjoin%
\definecolor{currentfill}{rgb}{0.000000,0.627451,0.686275}%
\pgfsetfillcolor{currentfill}%
\pgfsetlinewidth{0.000000pt}%
\definecolor{currentstroke}{rgb}{0.000000,0.000000,0.000000}%
\pgfsetstrokecolor{currentstroke}%
\pgfsetstrokeopacity{0.000000}%
\pgfsetdash{}{0pt}%
\pgfpathmoveto{\pgfqpoint{2.205934in}{1.422201in}}%
\pgfpathlineto{\pgfqpoint{2.206018in}{1.417015in}}%
\pgfpathlineto{\pgfqpoint{2.181819in}{1.416001in}}%
\pgfpathlineto{\pgfqpoint{2.181491in}{1.429705in}}%
\pgfpathlineto{\pgfqpoint{2.180394in}{1.464846in}}%
\pgfpathlineto{\pgfqpoint{2.193447in}{1.465250in}}%
\pgfpathlineto{\pgfqpoint{2.204405in}{1.465435in}}%
\pgfpathlineto{\pgfqpoint{2.205934in}{1.422201in}}%
\pgfpathclose%
\pgfusepath{fill}%
\end{pgfscope}%
\begin{pgfscope}%
\pgfpathrectangle{\pgfqpoint{0.100000in}{0.100000in}}{\pgfqpoint{3.608454in}{2.310000in}}%
\pgfusepath{clip}%
\pgfsetbuttcap%
\pgfsetmiterjoin%
\definecolor{currentfill}{rgb}{0.000000,0.615686,0.692157}%
\pgfsetfillcolor{currentfill}%
\pgfsetlinewidth{0.000000pt}%
\definecolor{currentstroke}{rgb}{0.000000,0.000000,0.000000}%
\pgfsetstrokecolor{currentstroke}%
\pgfsetstrokeopacity{0.000000}%
\pgfsetdash{}{0pt}%
\pgfpathmoveto{\pgfqpoint{1.826743in}{0.706842in}}%
\pgfpathlineto{\pgfqpoint{1.799907in}{0.708077in}}%
\pgfpathlineto{\pgfqpoint{1.791780in}{0.708494in}}%
\pgfpathlineto{\pgfqpoint{1.792400in}{0.720174in}}%
\pgfpathlineto{\pgfqpoint{1.757470in}{0.722100in}}%
\pgfpathlineto{\pgfqpoint{1.759948in}{0.761279in}}%
\pgfpathlineto{\pgfqpoint{1.751590in}{0.761847in}}%
\pgfpathlineto{\pgfqpoint{1.753852in}{0.801689in}}%
\pgfpathlineto{\pgfqpoint{1.759459in}{0.801389in}}%
\pgfpathlineto{\pgfqpoint{1.788954in}{0.799743in}}%
\pgfpathlineto{\pgfqpoint{1.786260in}{0.758457in}}%
\pgfpathlineto{\pgfqpoint{1.792433in}{0.750538in}}%
\pgfpathlineto{\pgfqpoint{1.804967in}{0.750514in}}%
\pgfpathlineto{\pgfqpoint{1.813328in}{0.744891in}}%
\pgfpathlineto{\pgfqpoint{1.818920in}{0.750454in}}%
\pgfpathlineto{\pgfqpoint{1.828725in}{0.748261in}}%
\pgfpathlineto{\pgfqpoint{1.826743in}{0.706842in}}%
\pgfpathclose%
\pgfusepath{fill}%
\end{pgfscope}%
\begin{pgfscope}%
\pgfpathrectangle{\pgfqpoint{0.100000in}{0.100000in}}{\pgfqpoint{3.608454in}{2.310000in}}%
\pgfusepath{clip}%
\pgfsetbuttcap%
\pgfsetmiterjoin%
\definecolor{currentfill}{rgb}{0.000000,0.639216,0.680392}%
\pgfsetfillcolor{currentfill}%
\pgfsetlinewidth{0.000000pt}%
\definecolor{currentstroke}{rgb}{0.000000,0.000000,0.000000}%
\pgfsetstrokecolor{currentstroke}%
\pgfsetstrokeopacity{0.000000}%
\pgfsetdash{}{0pt}%
\pgfpathmoveto{\pgfqpoint{2.472257in}{0.716197in}}%
\pgfpathlineto{\pgfqpoint{2.470222in}{0.750540in}}%
\pgfpathlineto{\pgfqpoint{2.458660in}{0.749893in}}%
\pgfpathlineto{\pgfqpoint{2.458936in}{0.746389in}}%
\pgfpathlineto{\pgfqpoint{2.449720in}{0.745908in}}%
\pgfpathlineto{\pgfqpoint{2.447975in}{0.756350in}}%
\pgfpathlineto{\pgfqpoint{2.447021in}{0.775387in}}%
\pgfpathlineto{\pgfqpoint{2.468672in}{0.778139in}}%
\pgfpathlineto{\pgfqpoint{2.485619in}{0.780455in}}%
\pgfpathlineto{\pgfqpoint{2.487539in}{0.751603in}}%
\pgfpathlineto{\pgfqpoint{2.525381in}{0.754028in}}%
\pgfpathlineto{\pgfqpoint{2.528702in}{0.712404in}}%
\pgfpathlineto{\pgfqpoint{2.494101in}{0.710117in}}%
\pgfpathlineto{\pgfqpoint{2.493658in}{0.717377in}}%
\pgfpathlineto{\pgfqpoint{2.482623in}{0.717728in}}%
\pgfpathlineto{\pgfqpoint{2.472257in}{0.716197in}}%
\pgfpathclose%
\pgfusepath{fill}%
\end{pgfscope}%
\begin{pgfscope}%
\pgfpathrectangle{\pgfqpoint{0.100000in}{0.100000in}}{\pgfqpoint{3.608454in}{2.310000in}}%
\pgfusepath{clip}%
\pgfsetbuttcap%
\pgfsetmiterjoin%
\definecolor{currentfill}{rgb}{0.000000,0.196078,0.901961}%
\pgfsetfillcolor{currentfill}%
\pgfsetlinewidth{0.000000pt}%
\definecolor{currentstroke}{rgb}{0.000000,0.000000,0.000000}%
\pgfsetstrokecolor{currentstroke}%
\pgfsetstrokeopacity{0.000000}%
\pgfsetdash{}{0pt}%
\pgfpathmoveto{\pgfqpoint{3.129858in}{1.404228in}}%
\pgfpathlineto{\pgfqpoint{3.131195in}{1.415085in}}%
\pgfpathlineto{\pgfqpoint{3.124006in}{1.419453in}}%
\pgfpathlineto{\pgfqpoint{3.127584in}{1.429231in}}%
\pgfpathlineto{\pgfqpoint{3.125843in}{1.434080in}}%
\pgfpathlineto{\pgfqpoint{3.141078in}{1.437828in}}%
\pgfpathlineto{\pgfqpoint{3.144962in}{1.435381in}}%
\pgfpathlineto{\pgfqpoint{3.152017in}{1.444895in}}%
\pgfpathlineto{\pgfqpoint{3.153317in}{1.449269in}}%
\pgfpathlineto{\pgfqpoint{3.156244in}{1.464323in}}%
\pgfpathlineto{\pgfqpoint{3.159520in}{1.465911in}}%
\pgfpathlineto{\pgfqpoint{3.161516in}{1.490087in}}%
\pgfpathlineto{\pgfqpoint{3.166603in}{1.499150in}}%
\pgfpathlineto{\pgfqpoint{3.181076in}{1.501996in}}%
\pgfpathlineto{\pgfqpoint{3.177742in}{1.492668in}}%
\pgfpathlineto{\pgfqpoint{3.186509in}{1.489666in}}%
\pgfpathlineto{\pgfqpoint{3.191082in}{1.485805in}}%
\pgfpathlineto{\pgfqpoint{3.188742in}{1.472751in}}%
\pgfpathlineto{\pgfqpoint{3.192873in}{1.468222in}}%
\pgfpathlineto{\pgfqpoint{3.201406in}{1.464696in}}%
\pgfpathlineto{\pgfqpoint{3.204352in}{1.460668in}}%
\pgfpathlineto{\pgfqpoint{3.223628in}{1.451194in}}%
\pgfpathlineto{\pgfqpoint{3.226611in}{1.445349in}}%
\pgfpathlineto{\pgfqpoint{3.225469in}{1.437928in}}%
\pgfpathlineto{\pgfqpoint{3.227384in}{1.433116in}}%
\pgfpathlineto{\pgfqpoint{3.238143in}{1.432326in}}%
\pgfpathlineto{\pgfqpoint{3.241599in}{1.419010in}}%
\pgfpathlineto{\pgfqpoint{3.251419in}{1.409012in}}%
\pgfpathlineto{\pgfqpoint{3.254865in}{1.395654in}}%
\pgfpathlineto{\pgfqpoint{3.260051in}{1.390426in}}%
\pgfpathlineto{\pgfqpoint{3.258704in}{1.383890in}}%
\pgfpathlineto{\pgfqpoint{3.252557in}{1.388608in}}%
\pgfpathlineto{\pgfqpoint{3.237222in}{1.393476in}}%
\pgfpathlineto{\pgfqpoint{3.230597in}{1.391657in}}%
\pgfpathlineto{\pgfqpoint{3.214679in}{1.396873in}}%
\pgfpathlineto{\pgfqpoint{3.209863in}{1.398797in}}%
\pgfpathlineto{\pgfqpoint{3.199440in}{1.393885in}}%
\pgfpathlineto{\pgfqpoint{3.196379in}{1.396687in}}%
\pgfpathlineto{\pgfqpoint{3.194258in}{1.405953in}}%
\pgfpathlineto{\pgfqpoint{3.182634in}{1.411460in}}%
\pgfpathlineto{\pgfqpoint{3.175115in}{1.399614in}}%
\pgfpathlineto{\pgfqpoint{3.175771in}{1.395430in}}%
\pgfpathlineto{\pgfqpoint{3.169602in}{1.394541in}}%
\pgfpathlineto{\pgfqpoint{3.158541in}{1.406177in}}%
\pgfpathlineto{\pgfqpoint{3.153282in}{1.414127in}}%
\pgfpathlineto{\pgfqpoint{3.141594in}{1.398020in}}%
\pgfpathlineto{\pgfqpoint{3.137714in}{1.398047in}}%
\pgfpathlineto{\pgfqpoint{3.129858in}{1.404228in}}%
\pgfpathclose%
\pgfusepath{fill}%
\end{pgfscope}%
\begin{pgfscope}%
\pgfpathrectangle{\pgfqpoint{0.100000in}{0.100000in}}{\pgfqpoint{3.608454in}{2.310000in}}%
\pgfusepath{clip}%
\pgfsetbuttcap%
\pgfsetmiterjoin%
\definecolor{currentfill}{rgb}{0.000000,0.952941,0.523529}%
\pgfsetfillcolor{currentfill}%
\pgfsetlinewidth{0.000000pt}%
\definecolor{currentstroke}{rgb}{0.000000,0.000000,0.000000}%
\pgfsetstrokecolor{currentstroke}%
\pgfsetstrokeopacity{0.000000}%
\pgfsetdash{}{0pt}%
\pgfpathmoveto{\pgfqpoint{0.605938in}{1.604509in}}%
\pgfpathlineto{\pgfqpoint{0.599878in}{1.583444in}}%
\pgfpathlineto{\pgfqpoint{0.570419in}{1.591928in}}%
\pgfpathlineto{\pgfqpoint{0.567874in}{1.599010in}}%
\pgfpathlineto{\pgfqpoint{0.563417in}{1.600581in}}%
\pgfpathlineto{\pgfqpoint{0.556871in}{1.597127in}}%
\pgfpathlineto{\pgfqpoint{0.538914in}{1.596597in}}%
\pgfpathlineto{\pgfqpoint{0.540950in}{1.606159in}}%
\pgfpathlineto{\pgfqpoint{0.529553in}{1.637642in}}%
\pgfpathlineto{\pgfqpoint{0.526964in}{1.639730in}}%
\pgfpathlineto{\pgfqpoint{0.529796in}{1.650793in}}%
\pgfpathlineto{\pgfqpoint{0.534394in}{1.657082in}}%
\pgfpathlineto{\pgfqpoint{0.531886in}{1.661402in}}%
\pgfpathlineto{\pgfqpoint{0.537339in}{1.664898in}}%
\pgfpathlineto{\pgfqpoint{0.540383in}{1.672372in}}%
\pgfpathlineto{\pgfqpoint{0.534540in}{1.676936in}}%
\pgfpathlineto{\pgfqpoint{0.535033in}{1.684697in}}%
\pgfpathlineto{\pgfqpoint{0.544816in}{1.681752in}}%
\pgfpathlineto{\pgfqpoint{0.560209in}{1.677284in}}%
\pgfpathlineto{\pgfqpoint{0.555953in}{1.662869in}}%
\pgfpathlineto{\pgfqpoint{0.558505in}{1.659196in}}%
\pgfpathlineto{\pgfqpoint{0.566606in}{1.658522in}}%
\pgfpathlineto{\pgfqpoint{0.574470in}{1.662350in}}%
\pgfpathlineto{\pgfqpoint{0.584658in}{1.656876in}}%
\pgfpathlineto{\pgfqpoint{0.594599in}{1.640016in}}%
\pgfpathlineto{\pgfqpoint{0.601624in}{1.635606in}}%
\pgfpathlineto{\pgfqpoint{0.600541in}{1.629953in}}%
\pgfpathlineto{\pgfqpoint{0.604334in}{1.622797in}}%
\pgfpathlineto{\pgfqpoint{0.600479in}{1.609616in}}%
\pgfpathlineto{\pgfqpoint{0.605938in}{1.604509in}}%
\pgfpathclose%
\pgfusepath{fill}%
\end{pgfscope}%
\begin{pgfscope}%
\pgfpathrectangle{\pgfqpoint{0.100000in}{0.100000in}}{\pgfqpoint{3.608454in}{2.310000in}}%
\pgfusepath{clip}%
\pgfsetbuttcap%
\pgfsetmiterjoin%
\definecolor{currentfill}{rgb}{0.000000,0.470588,0.764706}%
\pgfsetfillcolor{currentfill}%
\pgfsetlinewidth{0.000000pt}%
\definecolor{currentstroke}{rgb}{0.000000,0.000000,0.000000}%
\pgfsetstrokecolor{currentstroke}%
\pgfsetstrokeopacity{0.000000}%
\pgfsetdash{}{0pt}%
\pgfpathmoveto{\pgfqpoint{1.888636in}{1.700186in}}%
\pgfpathlineto{\pgfqpoint{1.889333in}{1.734748in}}%
\pgfpathlineto{\pgfqpoint{1.916523in}{1.733891in}}%
\pgfpathlineto{\pgfqpoint{1.917260in}{1.754415in}}%
\pgfpathlineto{\pgfqpoint{1.944449in}{1.753368in}}%
\pgfpathlineto{\pgfqpoint{1.944686in}{1.753364in}}%
\pgfpathlineto{\pgfqpoint{1.943727in}{1.725985in}}%
\pgfpathlineto{\pgfqpoint{1.957610in}{1.725492in}}%
\pgfpathlineto{\pgfqpoint{1.956891in}{1.697804in}}%
\pgfpathlineto{\pgfqpoint{1.888636in}{1.700186in}}%
\pgfpathclose%
\pgfusepath{fill}%
\end{pgfscope}%
\begin{pgfscope}%
\pgfpathrectangle{\pgfqpoint{0.100000in}{0.100000in}}{\pgfqpoint{3.608454in}{2.310000in}}%
\pgfusepath{clip}%
\pgfsetbuttcap%
\pgfsetmiterjoin%
\definecolor{currentfill}{rgb}{0.000000,0.909804,0.545098}%
\pgfsetfillcolor{currentfill}%
\pgfsetlinewidth{0.000000pt}%
\definecolor{currentstroke}{rgb}{0.000000,0.000000,0.000000}%
\pgfsetstrokecolor{currentstroke}%
\pgfsetstrokeopacity{0.000000}%
\pgfsetdash{}{0pt}%
\pgfpathmoveto{\pgfqpoint{0.937150in}{2.279229in}}%
\pgfpathlineto{\pgfqpoint{0.927075in}{2.268237in}}%
\pgfpathlineto{\pgfqpoint{0.925222in}{2.260906in}}%
\pgfpathlineto{\pgfqpoint{0.920139in}{2.255027in}}%
\pgfpathlineto{\pgfqpoint{0.926800in}{2.253370in}}%
\pgfpathlineto{\pgfqpoint{0.916424in}{2.212894in}}%
\pgfpathlineto{\pgfqpoint{0.913019in}{2.206570in}}%
\pgfpathlineto{\pgfqpoint{0.908021in}{2.186968in}}%
\pgfpathlineto{\pgfqpoint{0.902377in}{2.193011in}}%
\pgfpathlineto{\pgfqpoint{0.894472in}{2.191910in}}%
\pgfpathlineto{\pgfqpoint{0.891121in}{2.195618in}}%
\pgfpathlineto{\pgfqpoint{0.879920in}{2.194564in}}%
\pgfpathlineto{\pgfqpoint{0.878250in}{2.201957in}}%
\pgfpathlineto{\pgfqpoint{0.873890in}{2.207982in}}%
\pgfpathlineto{\pgfqpoint{0.862221in}{2.202501in}}%
\pgfpathlineto{\pgfqpoint{0.857008in}{2.211536in}}%
\pgfpathlineto{\pgfqpoint{0.843741in}{2.216808in}}%
\pgfpathlineto{\pgfqpoint{0.836363in}{2.217995in}}%
\pgfpathlineto{\pgfqpoint{0.836288in}{2.219545in}}%
\pgfpathlineto{\pgfqpoint{0.836249in}{2.233779in}}%
\pgfpathlineto{\pgfqpoint{0.832614in}{2.236217in}}%
\pgfpathlineto{\pgfqpoint{0.819597in}{2.231350in}}%
\pgfpathlineto{\pgfqpoint{0.815034in}{2.234984in}}%
\pgfpathlineto{\pgfqpoint{0.802328in}{2.232125in}}%
\pgfpathlineto{\pgfqpoint{0.800018in}{2.241064in}}%
\pgfpathlineto{\pgfqpoint{0.791522in}{2.239767in}}%
\pgfpathlineto{\pgfqpoint{0.790372in}{2.232053in}}%
\pgfpathlineto{\pgfqpoint{0.780931in}{2.234819in}}%
\pgfpathlineto{\pgfqpoint{0.780160in}{2.239854in}}%
\pgfpathlineto{\pgfqpoint{0.769468in}{2.254538in}}%
\pgfpathlineto{\pgfqpoint{0.771087in}{2.257398in}}%
\pgfpathlineto{\pgfqpoint{0.765081in}{2.268518in}}%
\pgfpathlineto{\pgfqpoint{0.760184in}{2.281456in}}%
\pgfpathlineto{\pgfqpoint{0.760304in}{2.288239in}}%
\pgfpathlineto{\pgfqpoint{0.757956in}{2.296770in}}%
\pgfpathlineto{\pgfqpoint{0.760555in}{2.298682in}}%
\pgfpathlineto{\pgfqpoint{0.765166in}{2.300496in}}%
\pgfpathlineto{\pgfqpoint{0.763231in}{2.326563in}}%
\pgfpathlineto{\pgfqpoint{0.801462in}{2.315507in}}%
\pgfpathlineto{\pgfqpoint{0.847920in}{2.302634in}}%
\pgfpathlineto{\pgfqpoint{0.882493in}{2.293352in}}%
\pgfpathlineto{\pgfqpoint{0.937150in}{2.279229in}}%
\pgfpathclose%
\pgfusepath{fill}%
\end{pgfscope}%
\begin{pgfscope}%
\pgfpathrectangle{\pgfqpoint{0.100000in}{0.100000in}}{\pgfqpoint{3.608454in}{2.310000in}}%
\pgfusepath{clip}%
\pgfsetbuttcap%
\pgfsetmiterjoin%
\definecolor{currentfill}{rgb}{0.000000,0.580392,0.709804}%
\pgfsetfillcolor{currentfill}%
\pgfsetlinewidth{0.000000pt}%
\definecolor{currentstroke}{rgb}{0.000000,0.000000,0.000000}%
\pgfsetstrokecolor{currentstroke}%
\pgfsetstrokeopacity{0.000000}%
\pgfsetdash{}{0pt}%
\pgfpathmoveto{\pgfqpoint{2.821573in}{0.886623in}}%
\pgfpathlineto{\pgfqpoint{2.814267in}{0.888921in}}%
\pgfpathlineto{\pgfqpoint{2.805495in}{0.895708in}}%
\pgfpathlineto{\pgfqpoint{2.802438in}{0.905479in}}%
\pgfpathlineto{\pgfqpoint{2.802911in}{0.922322in}}%
\pgfpathlineto{\pgfqpoint{2.802892in}{0.925244in}}%
\pgfpathlineto{\pgfqpoint{2.807171in}{0.925663in}}%
\pgfpathlineto{\pgfqpoint{2.809234in}{0.933588in}}%
\pgfpathlineto{\pgfqpoint{2.818507in}{0.934685in}}%
\pgfpathlineto{\pgfqpoint{2.830443in}{0.932055in}}%
\pgfpathlineto{\pgfqpoint{2.838874in}{0.945035in}}%
\pgfpathlineto{\pgfqpoint{2.843454in}{0.939687in}}%
\pgfpathlineto{\pgfqpoint{2.848096in}{0.929927in}}%
\pgfpathlineto{\pgfqpoint{2.848183in}{0.924494in}}%
\pgfpathlineto{\pgfqpoint{2.844937in}{0.925826in}}%
\pgfpathlineto{\pgfqpoint{2.833490in}{0.924483in}}%
\pgfpathlineto{\pgfqpoint{2.835752in}{0.904367in}}%
\pgfpathlineto{\pgfqpoint{2.830655in}{0.902493in}}%
\pgfpathlineto{\pgfqpoint{2.830174in}{0.891948in}}%
\pgfpathlineto{\pgfqpoint{2.827685in}{0.882746in}}%
\pgfpathlineto{\pgfqpoint{2.821573in}{0.886623in}}%
\pgfpathclose%
\pgfusepath{fill}%
\end{pgfscope}%
\begin{pgfscope}%
\pgfpathrectangle{\pgfqpoint{0.100000in}{0.100000in}}{\pgfqpoint{3.608454in}{2.310000in}}%
\pgfusepath{clip}%
\pgfsetbuttcap%
\pgfsetmiterjoin%
\definecolor{currentfill}{rgb}{0.000000,0.286275,0.856863}%
\pgfsetfillcolor{currentfill}%
\pgfsetlinewidth{0.000000pt}%
\definecolor{currentstroke}{rgb}{0.000000,0.000000,0.000000}%
\pgfsetstrokecolor{currentstroke}%
\pgfsetstrokeopacity{0.000000}%
\pgfsetdash{}{0pt}%
\pgfpathmoveto{\pgfqpoint{3.238143in}{1.432326in}}%
\pgfpathlineto{\pgfqpoint{3.227384in}{1.433116in}}%
\pgfpathlineto{\pgfqpoint{3.225469in}{1.437928in}}%
\pgfpathlineto{\pgfqpoint{3.226611in}{1.445349in}}%
\pgfpathlineto{\pgfqpoint{3.223628in}{1.451194in}}%
\pgfpathlineto{\pgfqpoint{3.204352in}{1.460668in}}%
\pgfpathlineto{\pgfqpoint{3.201406in}{1.464696in}}%
\pgfpathlineto{\pgfqpoint{3.192873in}{1.468222in}}%
\pgfpathlineto{\pgfqpoint{3.188742in}{1.472751in}}%
\pgfpathlineto{\pgfqpoint{3.191082in}{1.485805in}}%
\pgfpathlineto{\pgfqpoint{3.186509in}{1.489666in}}%
\pgfpathlineto{\pgfqpoint{3.177742in}{1.492668in}}%
\pgfpathlineto{\pgfqpoint{3.181076in}{1.501996in}}%
\pgfpathlineto{\pgfqpoint{3.239245in}{1.513965in}}%
\pgfpathlineto{\pgfqpoint{3.239646in}{1.514041in}}%
\pgfpathlineto{\pgfqpoint{3.251965in}{1.502268in}}%
\pgfpathlineto{\pgfqpoint{3.252119in}{1.493455in}}%
\pgfpathlineto{\pgfqpoint{3.248435in}{1.488117in}}%
\pgfpathlineto{\pgfqpoint{3.238301in}{1.480374in}}%
\pgfpathlineto{\pgfqpoint{3.235613in}{1.470733in}}%
\pgfpathlineto{\pgfqpoint{3.238603in}{1.459477in}}%
\pgfpathlineto{\pgfqpoint{3.239193in}{1.448194in}}%
\pgfpathlineto{\pgfqpoint{3.235959in}{1.434274in}}%
\pgfpathlineto{\pgfqpoint{3.238143in}{1.432326in}}%
\pgfpathclose%
\pgfusepath{fill}%
\end{pgfscope}%
\begin{pgfscope}%
\pgfpathrectangle{\pgfqpoint{0.100000in}{0.100000in}}{\pgfqpoint{3.608454in}{2.310000in}}%
\pgfusepath{clip}%
\pgfsetbuttcap%
\pgfsetmiterjoin%
\definecolor{currentfill}{rgb}{0.000000,0.592157,0.703922}%
\pgfsetfillcolor{currentfill}%
\pgfsetlinewidth{0.000000pt}%
\definecolor{currentstroke}{rgb}{0.000000,0.000000,0.000000}%
\pgfsetstrokecolor{currentstroke}%
\pgfsetstrokeopacity{0.000000}%
\pgfsetdash{}{0pt}%
\pgfpathmoveto{\pgfqpoint{1.821831in}{1.246792in}}%
\pgfpathlineto{\pgfqpoint{1.821521in}{1.253978in}}%
\pgfpathlineto{\pgfqpoint{1.823241in}{1.288458in}}%
\pgfpathlineto{\pgfqpoint{1.822395in}{1.288502in}}%
\pgfpathlineto{\pgfqpoint{1.824072in}{1.322793in}}%
\pgfpathlineto{\pgfqpoint{1.891797in}{1.319820in}}%
\pgfpathlineto{\pgfqpoint{1.890012in}{1.250882in}}%
\pgfpathlineto{\pgfqpoint{1.856262in}{1.252307in}}%
\pgfpathlineto{\pgfqpoint{1.855927in}{1.245032in}}%
\pgfpathlineto{\pgfqpoint{1.821831in}{1.246792in}}%
\pgfpathclose%
\pgfusepath{fill}%
\end{pgfscope}%
\begin{pgfscope}%
\pgfpathrectangle{\pgfqpoint{0.100000in}{0.100000in}}{\pgfqpoint{3.608454in}{2.310000in}}%
\pgfusepath{clip}%
\pgfsetbuttcap%
\pgfsetmiterjoin%
\definecolor{currentfill}{rgb}{0.000000,0.678431,0.660784}%
\pgfsetfillcolor{currentfill}%
\pgfsetlinewidth{0.000000pt}%
\definecolor{currentstroke}{rgb}{0.000000,0.000000,0.000000}%
\pgfsetstrokecolor{currentstroke}%
\pgfsetstrokeopacity{0.000000}%
\pgfsetdash{}{0pt}%
\pgfpathmoveto{\pgfqpoint{2.625838in}{1.802295in}}%
\pgfpathlineto{\pgfqpoint{2.602311in}{1.800403in}}%
\pgfpathlineto{\pgfqpoint{2.600324in}{1.814817in}}%
\pgfpathlineto{\pgfqpoint{2.608715in}{1.818757in}}%
\pgfpathlineto{\pgfqpoint{2.608533in}{1.830495in}}%
\pgfpathlineto{\pgfqpoint{2.613100in}{1.831951in}}%
\pgfpathlineto{\pgfqpoint{2.615491in}{1.837204in}}%
\pgfpathlineto{\pgfqpoint{2.623020in}{1.837013in}}%
\pgfpathlineto{\pgfqpoint{2.625532in}{1.845348in}}%
\pgfpathlineto{\pgfqpoint{2.632241in}{1.851902in}}%
\pgfpathlineto{\pgfqpoint{2.634383in}{1.847792in}}%
\pgfpathlineto{\pgfqpoint{2.634627in}{1.836092in}}%
\pgfpathlineto{\pgfqpoint{2.632415in}{1.832536in}}%
\pgfpathlineto{\pgfqpoint{2.633665in}{1.824033in}}%
\pgfpathlineto{\pgfqpoint{2.640260in}{1.822509in}}%
\pgfpathlineto{\pgfqpoint{2.645608in}{1.835898in}}%
\pgfpathlineto{\pgfqpoint{2.646411in}{1.852152in}}%
\pgfpathlineto{\pgfqpoint{2.644203in}{1.859587in}}%
\pgfpathlineto{\pgfqpoint{2.653419in}{1.860520in}}%
\pgfpathlineto{\pgfqpoint{2.654113in}{1.853564in}}%
\pgfpathlineto{\pgfqpoint{2.674595in}{1.855705in}}%
\pgfpathlineto{\pgfqpoint{2.677389in}{1.835395in}}%
\pgfpathlineto{\pgfqpoint{2.680280in}{1.807943in}}%
\pgfpathlineto{\pgfqpoint{2.625838in}{1.802295in}}%
\pgfpathclose%
\pgfusepath{fill}%
\end{pgfscope}%
\begin{pgfscope}%
\pgfpathrectangle{\pgfqpoint{0.100000in}{0.100000in}}{\pgfqpoint{3.608454in}{2.310000in}}%
\pgfusepath{clip}%
\pgfsetbuttcap%
\pgfsetmiterjoin%
\definecolor{currentfill}{rgb}{0.000000,0.305882,0.847059}%
\pgfsetfillcolor{currentfill}%
\pgfsetlinewidth{0.000000pt}%
\definecolor{currentstroke}{rgb}{0.000000,0.000000,0.000000}%
\pgfsetstrokecolor{currentstroke}%
\pgfsetstrokeopacity{0.000000}%
\pgfsetdash{}{0pt}%
\pgfpathmoveto{\pgfqpoint{1.830021in}{2.007934in}}%
\pgfpathlineto{\pgfqpoint{1.828413in}{1.980392in}}%
\pgfpathlineto{\pgfqpoint{1.830555in}{1.980267in}}%
\pgfpathlineto{\pgfqpoint{1.828664in}{1.952547in}}%
\pgfpathlineto{\pgfqpoint{1.796923in}{1.954510in}}%
\pgfpathlineto{\pgfqpoint{1.801850in}{1.950975in}}%
\pgfpathlineto{\pgfqpoint{1.802847in}{1.946100in}}%
\pgfpathlineto{\pgfqpoint{1.799830in}{1.937828in}}%
\pgfpathlineto{\pgfqpoint{1.787591in}{1.935130in}}%
\pgfpathlineto{\pgfqpoint{1.783790in}{1.937380in}}%
\pgfpathlineto{\pgfqpoint{1.775193in}{1.927930in}}%
\pgfpathlineto{\pgfqpoint{1.774399in}{1.934946in}}%
\pgfpathlineto{\pgfqpoint{1.760656in}{1.935885in}}%
\pgfpathlineto{\pgfqpoint{1.762088in}{1.956478in}}%
\pgfpathlineto{\pgfqpoint{1.739184in}{1.958172in}}%
\pgfpathlineto{\pgfqpoint{1.739705in}{1.965094in}}%
\pgfpathlineto{\pgfqpoint{1.719085in}{1.966639in}}%
\pgfpathlineto{\pgfqpoint{1.720698in}{1.987602in}}%
\pgfpathlineto{\pgfqpoint{1.718098in}{1.987796in}}%
\pgfpathlineto{\pgfqpoint{1.720286in}{2.015429in}}%
\pgfpathlineto{\pgfqpoint{1.716982in}{2.015681in}}%
\pgfpathlineto{\pgfqpoint{1.718555in}{2.035319in}}%
\pgfpathlineto{\pgfqpoint{1.715648in}{2.039417in}}%
\pgfpathlineto{\pgfqpoint{1.716886in}{2.052052in}}%
\pgfpathlineto{\pgfqpoint{1.710047in}{2.050726in}}%
\pgfpathlineto{\pgfqpoint{1.710634in}{2.057876in}}%
\pgfpathlineto{\pgfqpoint{1.738152in}{2.055718in}}%
\pgfpathlineto{\pgfqpoint{1.807069in}{2.050968in}}%
\pgfpathlineto{\pgfqpoint{1.806219in}{2.037109in}}%
\pgfpathlineto{\pgfqpoint{1.801529in}{2.037408in}}%
\pgfpathlineto{\pgfqpoint{1.799746in}{2.009759in}}%
\pgfpathlineto{\pgfqpoint{1.830021in}{2.007934in}}%
\pgfpathclose%
\pgfusepath{fill}%
\end{pgfscope}%
\begin{pgfscope}%
\pgfpathrectangle{\pgfqpoint{0.100000in}{0.100000in}}{\pgfqpoint{3.608454in}{2.310000in}}%
\pgfusepath{clip}%
\pgfsetbuttcap%
\pgfsetmiterjoin%
\definecolor{currentfill}{rgb}{0.000000,0.423529,0.788235}%
\pgfsetfillcolor{currentfill}%
\pgfsetlinewidth{0.000000pt}%
\definecolor{currentstroke}{rgb}{0.000000,0.000000,0.000000}%
\pgfsetstrokecolor{currentstroke}%
\pgfsetstrokeopacity{0.000000}%
\pgfsetdash{}{0pt}%
\pgfpathmoveto{\pgfqpoint{3.285323in}{1.505755in}}%
\pgfpathlineto{\pgfqpoint{3.285117in}{1.511332in}}%
\pgfpathlineto{\pgfqpoint{3.281349in}{1.515956in}}%
\pgfpathlineto{\pgfqpoint{3.284745in}{1.526913in}}%
\pgfpathlineto{\pgfqpoint{3.284318in}{1.531602in}}%
\pgfpathlineto{\pgfqpoint{3.275612in}{1.531070in}}%
\pgfpathlineto{\pgfqpoint{3.271013in}{1.528198in}}%
\pgfpathlineto{\pgfqpoint{3.266966in}{1.519876in}}%
\pgfpathlineto{\pgfqpoint{3.245459in}{1.515286in}}%
\pgfpathlineto{\pgfqpoint{3.248921in}{1.519116in}}%
\pgfpathlineto{\pgfqpoint{3.250934in}{1.527034in}}%
\pgfpathlineto{\pgfqpoint{3.250059in}{1.534993in}}%
\pgfpathlineto{\pgfqpoint{3.252158in}{1.542409in}}%
\pgfpathlineto{\pgfqpoint{3.250514in}{1.547606in}}%
\pgfpathlineto{\pgfqpoint{3.262873in}{1.561037in}}%
\pgfpathlineto{\pgfqpoint{3.269282in}{1.579037in}}%
\pgfpathlineto{\pgfqpoint{3.275573in}{1.583652in}}%
\pgfpathlineto{\pgfqpoint{3.286080in}{1.595797in}}%
\pgfpathlineto{\pgfqpoint{3.294732in}{1.592344in}}%
\pgfpathlineto{\pgfqpoint{3.297442in}{1.582891in}}%
\pgfpathlineto{\pgfqpoint{3.303159in}{1.582589in}}%
\pgfpathlineto{\pgfqpoint{3.313073in}{1.572820in}}%
\pgfpathlineto{\pgfqpoint{3.330022in}{1.567556in}}%
\pgfpathlineto{\pgfqpoint{3.348171in}{1.542140in}}%
\pgfpathlineto{\pgfqpoint{3.353587in}{1.521095in}}%
\pgfpathlineto{\pgfqpoint{3.358684in}{1.520578in}}%
\pgfpathlineto{\pgfqpoint{3.354142in}{1.510014in}}%
\pgfpathlineto{\pgfqpoint{3.344357in}{1.498736in}}%
\pgfpathlineto{\pgfqpoint{3.337885in}{1.476407in}}%
\pgfpathlineto{\pgfqpoint{3.334527in}{1.471502in}}%
\pgfpathlineto{\pgfqpoint{3.328098in}{1.470069in}}%
\pgfpathlineto{\pgfqpoint{3.329144in}{1.488418in}}%
\pgfpathlineto{\pgfqpoint{3.310575in}{1.490636in}}%
\pgfpathlineto{\pgfqpoint{3.301327in}{1.493916in}}%
\pgfpathlineto{\pgfqpoint{3.288516in}{1.501805in}}%
\pgfpathlineto{\pgfqpoint{3.285323in}{1.505755in}}%
\pgfpathclose%
\pgfusepath{fill}%
\end{pgfscope}%
\begin{pgfscope}%
\pgfpathrectangle{\pgfqpoint{0.100000in}{0.100000in}}{\pgfqpoint{3.608454in}{2.310000in}}%
\pgfusepath{clip}%
\pgfsetbuttcap%
\pgfsetmiterjoin%
\definecolor{currentfill}{rgb}{0.000000,0.333333,0.833333}%
\pgfsetfillcolor{currentfill}%
\pgfsetlinewidth{0.000000pt}%
\definecolor{currentstroke}{rgb}{0.000000,0.000000,0.000000}%
\pgfsetstrokecolor{currentstroke}%
\pgfsetstrokeopacity{0.000000}%
\pgfsetdash{}{0pt}%
\pgfpathmoveto{\pgfqpoint{1.725856in}{1.224921in}}%
\pgfpathlineto{\pgfqpoint{1.723873in}{1.194151in}}%
\pgfpathlineto{\pgfqpoint{1.693075in}{1.196095in}}%
\pgfpathlineto{\pgfqpoint{1.695230in}{1.227106in}}%
\pgfpathlineto{\pgfqpoint{1.664847in}{1.229439in}}%
\pgfpathlineto{\pgfqpoint{1.666412in}{1.249548in}}%
\pgfpathlineto{\pgfqpoint{1.666954in}{1.256981in}}%
\pgfpathlineto{\pgfqpoint{1.726337in}{1.252520in}}%
\pgfpathlineto{\pgfqpoint{1.724493in}{1.225005in}}%
\pgfpathlineto{\pgfqpoint{1.725856in}{1.224921in}}%
\pgfpathclose%
\pgfusepath{fill}%
\end{pgfscope}%
\begin{pgfscope}%
\pgfpathrectangle{\pgfqpoint{0.100000in}{0.100000in}}{\pgfqpoint{3.608454in}{2.310000in}}%
\pgfusepath{clip}%
\pgfsetbuttcap%
\pgfsetmiterjoin%
\definecolor{currentfill}{rgb}{0.000000,0.596078,0.701961}%
\pgfsetfillcolor{currentfill}%
\pgfsetlinewidth{0.000000pt}%
\definecolor{currentstroke}{rgb}{0.000000,0.000000,0.000000}%
\pgfsetstrokecolor{currentstroke}%
\pgfsetstrokeopacity{0.000000}%
\pgfsetdash{}{0pt}%
\pgfpathmoveto{\pgfqpoint{2.856403in}{0.756814in}}%
\pgfpathlineto{\pgfqpoint{2.856733in}{0.753994in}}%
\pgfpathlineto{\pgfqpoint{2.874929in}{0.755916in}}%
\pgfpathlineto{\pgfqpoint{2.878161in}{0.725802in}}%
\pgfpathlineto{\pgfqpoint{2.834431in}{0.723057in}}%
\pgfpathlineto{\pgfqpoint{2.831112in}{0.753917in}}%
\pgfpathlineto{\pgfqpoint{2.856403in}{0.756814in}}%
\pgfpathclose%
\pgfusepath{fill}%
\end{pgfscope}%
\begin{pgfscope}%
\pgfpathrectangle{\pgfqpoint{0.100000in}{0.100000in}}{\pgfqpoint{3.608454in}{2.310000in}}%
\pgfusepath{clip}%
\pgfsetbuttcap%
\pgfsetmiterjoin%
\definecolor{currentfill}{rgb}{0.000000,0.552941,0.723529}%
\pgfsetfillcolor{currentfill}%
\pgfsetlinewidth{0.000000pt}%
\definecolor{currentstroke}{rgb}{0.000000,0.000000,0.000000}%
\pgfsetstrokecolor{currentstroke}%
\pgfsetstrokeopacity{0.000000}%
\pgfsetdash{}{0pt}%
\pgfpathmoveto{\pgfqpoint{2.519677in}{1.268419in}}%
\pgfpathlineto{\pgfqpoint{2.499015in}{1.266823in}}%
\pgfpathlineto{\pgfqpoint{2.497937in}{1.284068in}}%
\pgfpathlineto{\pgfqpoint{2.470338in}{1.282183in}}%
\pgfpathlineto{\pgfqpoint{2.470034in}{1.289131in}}%
\pgfpathlineto{\pgfqpoint{2.435549in}{1.287521in}}%
\pgfpathlineto{\pgfqpoint{2.433309in}{1.321900in}}%
\pgfpathlineto{\pgfqpoint{2.440188in}{1.322366in}}%
\pgfpathlineto{\pgfqpoint{2.439649in}{1.329250in}}%
\pgfpathlineto{\pgfqpoint{2.460771in}{1.330448in}}%
\pgfpathlineto{\pgfqpoint{2.467926in}{1.330458in}}%
\pgfpathlineto{\pgfqpoint{2.467492in}{1.337391in}}%
\pgfpathlineto{\pgfqpoint{2.494687in}{1.339376in}}%
\pgfpathlineto{\pgfqpoint{2.495702in}{1.322010in}}%
\pgfpathlineto{\pgfqpoint{2.497373in}{1.294428in}}%
\pgfpathlineto{\pgfqpoint{2.517928in}{1.295771in}}%
\pgfpathlineto{\pgfqpoint{2.519677in}{1.268419in}}%
\pgfpathclose%
\pgfusepath{fill}%
\end{pgfscope}%
\begin{pgfscope}%
\pgfpathrectangle{\pgfqpoint{0.100000in}{0.100000in}}{\pgfqpoint{3.608454in}{2.310000in}}%
\pgfusepath{clip}%
\pgfsetbuttcap%
\pgfsetmiterjoin%
\definecolor{currentfill}{rgb}{0.000000,0.647059,0.676471}%
\pgfsetfillcolor{currentfill}%
\pgfsetlinewidth{0.000000pt}%
\definecolor{currentstroke}{rgb}{0.000000,0.000000,0.000000}%
\pgfsetstrokecolor{currentstroke}%
\pgfsetstrokeopacity{0.000000}%
\pgfsetdash{}{0pt}%
\pgfpathmoveto{\pgfqpoint{2.669844in}{0.772749in}}%
\pgfpathlineto{\pgfqpoint{2.660506in}{0.772013in}}%
\pgfpathlineto{\pgfqpoint{2.655184in}{0.780136in}}%
\pgfpathlineto{\pgfqpoint{2.653773in}{0.795845in}}%
\pgfpathlineto{\pgfqpoint{2.652230in}{0.813085in}}%
\pgfpathlineto{\pgfqpoint{2.650875in}{0.827115in}}%
\pgfpathlineto{\pgfqpoint{2.656320in}{0.836856in}}%
\pgfpathlineto{\pgfqpoint{2.648743in}{0.852393in}}%
\pgfpathlineto{\pgfqpoint{2.647203in}{0.861864in}}%
\pgfpathlineto{\pgfqpoint{2.660793in}{0.862897in}}%
\pgfpathlineto{\pgfqpoint{2.660419in}{0.866374in}}%
\pgfpathlineto{\pgfqpoint{2.680466in}{0.868375in}}%
\pgfpathlineto{\pgfqpoint{2.682505in}{0.872318in}}%
\pgfpathlineto{\pgfqpoint{2.715600in}{0.875575in}}%
\pgfpathlineto{\pgfqpoint{2.717264in}{0.854842in}}%
\pgfpathlineto{\pgfqpoint{2.708963in}{0.847375in}}%
\pgfpathlineto{\pgfqpoint{2.711685in}{0.841617in}}%
\pgfpathlineto{\pgfqpoint{2.713352in}{0.819097in}}%
\pgfpathlineto{\pgfqpoint{2.714111in}{0.812491in}}%
\pgfpathlineto{\pgfqpoint{2.693497in}{0.810278in}}%
\pgfpathlineto{\pgfqpoint{2.692745in}{0.817134in}}%
\pgfpathlineto{\pgfqpoint{2.685856in}{0.816401in}}%
\pgfpathlineto{\pgfqpoint{2.683683in}{0.809309in}}%
\pgfpathlineto{\pgfqpoint{2.685966in}{0.784799in}}%
\pgfpathlineto{\pgfqpoint{2.682532in}{0.784446in}}%
\pgfpathlineto{\pgfqpoint{2.683521in}{0.774139in}}%
\pgfpathlineto{\pgfqpoint{2.669844in}{0.772749in}}%
\pgfpathclose%
\pgfusepath{fill}%
\end{pgfscope}%
\begin{pgfscope}%
\pgfpathrectangle{\pgfqpoint{0.100000in}{0.100000in}}{\pgfqpoint{3.608454in}{2.310000in}}%
\pgfusepath{clip}%
\pgfsetbuttcap%
\pgfsetmiterjoin%
\definecolor{currentfill}{rgb}{0.000000,0.368627,0.815686}%
\pgfsetfillcolor{currentfill}%
\pgfsetlinewidth{0.000000pt}%
\definecolor{currentstroke}{rgb}{0.000000,0.000000,0.000000}%
\pgfsetstrokecolor{currentstroke}%
\pgfsetstrokeopacity{0.000000}%
\pgfsetdash{}{0pt}%
\pgfpathmoveto{\pgfqpoint{0.658484in}{2.308460in}}%
\pgfpathlineto{\pgfqpoint{0.663035in}{2.299847in}}%
\pgfpathlineto{\pgfqpoint{0.657090in}{2.299503in}}%
\pgfpathlineto{\pgfqpoint{0.654017in}{2.294863in}}%
\pgfpathlineto{\pgfqpoint{0.658915in}{2.282360in}}%
\pgfpathlineto{\pgfqpoint{0.663385in}{2.275790in}}%
\pgfpathlineto{\pgfqpoint{0.662941in}{2.269541in}}%
\pgfpathlineto{\pgfqpoint{0.657685in}{2.267521in}}%
\pgfpathlineto{\pgfqpoint{0.656759in}{2.274209in}}%
\pgfpathlineto{\pgfqpoint{0.651475in}{2.279265in}}%
\pgfpathlineto{\pgfqpoint{0.653699in}{2.284940in}}%
\pgfpathlineto{\pgfqpoint{0.647875in}{2.296843in}}%
\pgfpathlineto{\pgfqpoint{0.658484in}{2.308460in}}%
\pgfpathclose%
\pgfusepath{fill}%
\end{pgfscope}%
\begin{pgfscope}%
\pgfpathrectangle{\pgfqpoint{0.100000in}{0.100000in}}{\pgfqpoint{3.608454in}{2.310000in}}%
\pgfusepath{clip}%
\pgfsetbuttcap%
\pgfsetmiterjoin%
\definecolor{currentfill}{rgb}{0.000000,0.368627,0.815686}%
\pgfsetfillcolor{currentfill}%
\pgfsetlinewidth{0.000000pt}%
\definecolor{currentstroke}{rgb}{0.000000,0.000000,0.000000}%
\pgfsetstrokecolor{currentstroke}%
\pgfsetstrokeopacity{0.000000}%
\pgfsetdash{}{0pt}%
\pgfpathmoveto{\pgfqpoint{0.669147in}{2.136525in}}%
\pgfpathlineto{\pgfqpoint{0.631314in}{2.147810in}}%
\pgfpathlineto{\pgfqpoint{0.572748in}{2.166216in}}%
\pgfpathlineto{\pgfqpoint{0.582091in}{2.197407in}}%
\pgfpathlineto{\pgfqpoint{0.593046in}{2.193971in}}%
\pgfpathlineto{\pgfqpoint{0.598094in}{2.209351in}}%
\pgfpathlineto{\pgfqpoint{0.597993in}{2.216816in}}%
\pgfpathlineto{\pgfqpoint{0.583086in}{2.221440in}}%
\pgfpathlineto{\pgfqpoint{0.586531in}{2.235034in}}%
\pgfpathlineto{\pgfqpoint{0.593017in}{2.254746in}}%
\pgfpathlineto{\pgfqpoint{0.595222in}{2.261545in}}%
\pgfpathlineto{\pgfqpoint{0.623636in}{2.252362in}}%
\pgfpathlineto{\pgfqpoint{0.635231in}{2.255252in}}%
\pgfpathlineto{\pgfqpoint{0.639739in}{2.262486in}}%
\pgfpathlineto{\pgfqpoint{0.648220in}{2.267247in}}%
\pgfpathlineto{\pgfqpoint{0.647635in}{2.272305in}}%
\pgfpathlineto{\pgfqpoint{0.652787in}{2.268601in}}%
\pgfpathlineto{\pgfqpoint{0.651158in}{2.248072in}}%
\pgfpathlineto{\pgfqpoint{0.656337in}{2.252428in}}%
\pgfpathlineto{\pgfqpoint{0.657168in}{2.258990in}}%
\pgfpathlineto{\pgfqpoint{0.664956in}{2.268322in}}%
\pgfpathlineto{\pgfqpoint{0.670631in}{2.272725in}}%
\pgfpathlineto{\pgfqpoint{0.666622in}{2.280032in}}%
\pgfpathlineto{\pgfqpoint{0.667712in}{2.286805in}}%
\pgfpathlineto{\pgfqpoint{0.660569in}{2.294722in}}%
\pgfpathlineto{\pgfqpoint{0.669465in}{2.296315in}}%
\pgfpathlineto{\pgfqpoint{0.662370in}{2.309255in}}%
\pgfpathlineto{\pgfqpoint{0.674084in}{2.320271in}}%
\pgfpathlineto{\pgfqpoint{0.672017in}{2.324335in}}%
\pgfpathlineto{\pgfqpoint{0.728782in}{2.306722in}}%
\pgfpathlineto{\pgfqpoint{0.752119in}{2.299760in}}%
\pgfpathlineto{\pgfqpoint{0.760555in}{2.298682in}}%
\pgfpathlineto{\pgfqpoint{0.757956in}{2.296770in}}%
\pgfpathlineto{\pgfqpoint{0.760304in}{2.288239in}}%
\pgfpathlineto{\pgfqpoint{0.752759in}{2.291869in}}%
\pgfpathlineto{\pgfqpoint{0.740255in}{2.288121in}}%
\pgfpathlineto{\pgfqpoint{0.736966in}{2.278404in}}%
\pgfpathlineto{\pgfqpoint{0.741543in}{2.265250in}}%
\pgfpathlineto{\pgfqpoint{0.738201in}{2.260046in}}%
\pgfpathlineto{\pgfqpoint{0.727318in}{2.257197in}}%
\pgfpathlineto{\pgfqpoint{0.721489in}{2.246797in}}%
\pgfpathlineto{\pgfqpoint{0.722977in}{2.230297in}}%
\pgfpathlineto{\pgfqpoint{0.718104in}{2.222638in}}%
\pgfpathlineto{\pgfqpoint{0.712161in}{2.222097in}}%
\pgfpathlineto{\pgfqpoint{0.701083in}{2.215781in}}%
\pgfpathlineto{\pgfqpoint{0.694418in}{2.209381in}}%
\pgfpathlineto{\pgfqpoint{0.694658in}{2.204301in}}%
\pgfpathlineto{\pgfqpoint{0.699138in}{2.201963in}}%
\pgfpathlineto{\pgfqpoint{0.697139in}{2.190297in}}%
\pgfpathlineto{\pgfqpoint{0.680298in}{2.173569in}}%
\pgfpathlineto{\pgfqpoint{0.682663in}{2.166183in}}%
\pgfpathlineto{\pgfqpoint{0.682016in}{2.153295in}}%
\pgfpathlineto{\pgfqpoint{0.675711in}{2.144078in}}%
\pgfpathlineto{\pgfqpoint{0.675514in}{2.134601in}}%
\pgfpathlineto{\pgfqpoint{0.669147in}{2.136525in}}%
\pgfpathclose%
\pgfusepath{fill}%
\end{pgfscope}%
\begin{pgfscope}%
\pgfpathrectangle{\pgfqpoint{0.100000in}{0.100000in}}{\pgfqpoint{3.608454in}{2.310000in}}%
\pgfusepath{clip}%
\pgfsetbuttcap%
\pgfsetmiterjoin%
\definecolor{currentfill}{rgb}{0.000000,0.541176,0.729412}%
\pgfsetfillcolor{currentfill}%
\pgfsetlinewidth{0.000000pt}%
\definecolor{currentstroke}{rgb}{0.000000,0.000000,0.000000}%
\pgfsetstrokecolor{currentstroke}%
\pgfsetstrokeopacity{0.000000}%
\pgfsetdash{}{0pt}%
\pgfpathmoveto{\pgfqpoint{2.615996in}{1.634560in}}%
\pgfpathlineto{\pgfqpoint{2.658513in}{1.638934in}}%
\pgfpathlineto{\pgfqpoint{2.655531in}{1.666251in}}%
\pgfpathlineto{\pgfqpoint{2.682759in}{1.669371in}}%
\pgfpathlineto{\pgfqpoint{2.685918in}{1.641892in}}%
\pgfpathlineto{\pgfqpoint{2.706426in}{1.644227in}}%
\pgfpathlineto{\pgfqpoint{2.710082in}{1.616628in}}%
\pgfpathlineto{\pgfqpoint{2.703247in}{1.616030in}}%
\pgfpathlineto{\pgfqpoint{2.648568in}{1.609914in}}%
\pgfpathlineto{\pgfqpoint{2.621651in}{1.607409in}}%
\pgfpathlineto{\pgfqpoint{2.620303in}{1.620947in}}%
\pgfpathlineto{\pgfqpoint{2.612010in}{1.620134in}}%
\pgfpathlineto{\pgfqpoint{2.615996in}{1.634560in}}%
\pgfpathclose%
\pgfusepath{fill}%
\end{pgfscope}%
\begin{pgfscope}%
\pgfpathrectangle{\pgfqpoint{0.100000in}{0.100000in}}{\pgfqpoint{3.608454in}{2.310000in}}%
\pgfusepath{clip}%
\pgfsetbuttcap%
\pgfsetmiterjoin%
\definecolor{currentfill}{rgb}{0.000000,0.654902,0.672549}%
\pgfsetfillcolor{currentfill}%
\pgfsetlinewidth{0.000000pt}%
\definecolor{currentstroke}{rgb}{0.000000,0.000000,0.000000}%
\pgfsetstrokecolor{currentstroke}%
\pgfsetstrokeopacity{0.000000}%
\pgfsetdash{}{0pt}%
\pgfpathmoveto{\pgfqpoint{1.303565in}{2.014897in}}%
\pgfpathlineto{\pgfqpoint{1.322819in}{2.011773in}}%
\pgfpathlineto{\pgfqpoint{1.316114in}{1.974564in}}%
\pgfpathlineto{\pgfqpoint{1.302428in}{1.976762in}}%
\pgfpathlineto{\pgfqpoint{1.268315in}{1.982523in}}%
\pgfpathlineto{\pgfqpoint{1.261011in}{1.981452in}}%
\pgfpathlineto{\pgfqpoint{1.225529in}{1.987742in}}%
\pgfpathlineto{\pgfqpoint{1.224264in}{1.998009in}}%
\pgfpathlineto{\pgfqpoint{1.228711in}{2.003790in}}%
\pgfpathlineto{\pgfqpoint{1.224493in}{2.008821in}}%
\pgfpathlineto{\pgfqpoint{1.219896in}{2.007840in}}%
\pgfpathlineto{\pgfqpoint{1.216917in}{2.019801in}}%
\pgfpathlineto{\pgfqpoint{1.213019in}{2.026776in}}%
\pgfpathlineto{\pgfqpoint{1.210034in}{2.040419in}}%
\pgfpathlineto{\pgfqpoint{1.203339in}{2.044520in}}%
\pgfpathlineto{\pgfqpoint{1.203686in}{2.050112in}}%
\pgfpathlineto{\pgfqpoint{1.210055in}{2.048940in}}%
\pgfpathlineto{\pgfqpoint{1.211322in}{2.055765in}}%
\pgfpathlineto{\pgfqpoint{1.231866in}{2.052930in}}%
\pgfpathlineto{\pgfqpoint{1.244581in}{2.048177in}}%
\pgfpathlineto{\pgfqpoint{1.251605in}{2.040078in}}%
\pgfpathlineto{\pgfqpoint{1.259993in}{2.025470in}}%
\pgfpathlineto{\pgfqpoint{1.264971in}{2.020151in}}%
\pgfpathlineto{\pgfqpoint{1.292244in}{2.021526in}}%
\pgfpathlineto{\pgfqpoint{1.299591in}{2.015508in}}%
\pgfpathlineto{\pgfqpoint{1.303565in}{2.014897in}}%
\pgfpathclose%
\pgfusepath{fill}%
\end{pgfscope}%
\begin{pgfscope}%
\pgfpathrectangle{\pgfqpoint{0.100000in}{0.100000in}}{\pgfqpoint{3.608454in}{2.310000in}}%
\pgfusepath{clip}%
\pgfsetbuttcap%
\pgfsetmiterjoin%
\definecolor{currentfill}{rgb}{0.000000,0.380392,0.809804}%
\pgfsetfillcolor{currentfill}%
\pgfsetlinewidth{0.000000pt}%
\definecolor{currentstroke}{rgb}{0.000000,0.000000,0.000000}%
\pgfsetstrokecolor{currentstroke}%
\pgfsetstrokeopacity{0.000000}%
\pgfsetdash{}{0pt}%
\pgfpathmoveto{\pgfqpoint{3.343167in}{1.599400in}}%
\pgfpathlineto{\pgfqpoint{3.345889in}{1.596696in}}%
\pgfpathlineto{\pgfqpoint{3.339300in}{1.581420in}}%
\pgfpathlineto{\pgfqpoint{3.327731in}{1.570641in}}%
\pgfpathlineto{\pgfqpoint{3.330022in}{1.567556in}}%
\pgfpathlineto{\pgfqpoint{3.313073in}{1.572820in}}%
\pgfpathlineto{\pgfqpoint{3.303159in}{1.582589in}}%
\pgfpathlineto{\pgfqpoint{3.297442in}{1.582891in}}%
\pgfpathlineto{\pgfqpoint{3.294732in}{1.592344in}}%
\pgfpathlineto{\pgfqpoint{3.286080in}{1.595797in}}%
\pgfpathlineto{\pgfqpoint{3.284662in}{1.609183in}}%
\pgfpathlineto{\pgfqpoint{3.288120in}{1.611032in}}%
\pgfpathlineto{\pgfqpoint{3.289903in}{1.618119in}}%
\pgfpathlineto{\pgfqpoint{3.284382in}{1.624557in}}%
\pgfpathlineto{\pgfqpoint{3.294577in}{1.644062in}}%
\pgfpathlineto{\pgfqpoint{3.295750in}{1.652950in}}%
\pgfpathlineto{\pgfqpoint{3.302353in}{1.660141in}}%
\pgfpathlineto{\pgfqpoint{3.332839in}{1.649737in}}%
\pgfpathlineto{\pgfqpoint{3.344220in}{1.667169in}}%
\pgfpathlineto{\pgfqpoint{3.346226in}{1.661401in}}%
\pgfpathlineto{\pgfqpoint{3.353043in}{1.654103in}}%
\pgfpathlineto{\pgfqpoint{3.355371in}{1.641928in}}%
\pgfpathlineto{\pgfqpoint{3.352075in}{1.614317in}}%
\pgfpathlineto{\pgfqpoint{3.344569in}{1.612129in}}%
\pgfpathlineto{\pgfqpoint{3.343167in}{1.599400in}}%
\pgfpathclose%
\pgfusepath{fill}%
\end{pgfscope}%
\begin{pgfscope}%
\pgfpathrectangle{\pgfqpoint{0.100000in}{0.100000in}}{\pgfqpoint{3.608454in}{2.310000in}}%
\pgfusepath{clip}%
\pgfsetbuttcap%
\pgfsetmiterjoin%
\definecolor{currentfill}{rgb}{0.000000,0.670588,0.664706}%
\pgfsetfillcolor{currentfill}%
\pgfsetlinewidth{0.000000pt}%
\definecolor{currentstroke}{rgb}{0.000000,0.000000,0.000000}%
\pgfsetstrokecolor{currentstroke}%
\pgfsetstrokeopacity{0.000000}%
\pgfsetdash{}{0pt}%
\pgfpathmoveto{\pgfqpoint{2.893588in}{0.882495in}}%
\pgfpathlineto{\pgfqpoint{2.910191in}{0.894634in}}%
\pgfpathlineto{\pgfqpoint{2.909397in}{0.902008in}}%
\pgfpathlineto{\pgfqpoint{2.903959in}{0.905713in}}%
\pgfpathlineto{\pgfqpoint{2.899980in}{0.912387in}}%
\pgfpathlineto{\pgfqpoint{2.899562in}{0.918887in}}%
\pgfpathlineto{\pgfqpoint{2.905128in}{0.926113in}}%
\pgfpathlineto{\pgfqpoint{2.910167in}{0.929143in}}%
\pgfpathlineto{\pgfqpoint{2.911845in}{0.934034in}}%
\pgfpathlineto{\pgfqpoint{2.918368in}{0.938239in}}%
\pgfpathlineto{\pgfqpoint{2.919634in}{0.933284in}}%
\pgfpathlineto{\pgfqpoint{2.925307in}{0.930408in}}%
\pgfpathlineto{\pgfqpoint{2.933286in}{0.922296in}}%
\pgfpathlineto{\pgfqpoint{2.937493in}{0.915146in}}%
\pgfpathlineto{\pgfqpoint{2.937410in}{0.908377in}}%
\pgfpathlineto{\pgfqpoint{2.944266in}{0.903605in}}%
\pgfpathlineto{\pgfqpoint{2.940759in}{0.897010in}}%
\pgfpathlineto{\pgfqpoint{2.943276in}{0.894232in}}%
\pgfpathlineto{\pgfqpoint{2.935640in}{0.887719in}}%
\pgfpathlineto{\pgfqpoint{2.932829in}{0.881990in}}%
\pgfpathlineto{\pgfqpoint{2.928863in}{0.867621in}}%
\pgfpathlineto{\pgfqpoint{2.930056in}{0.865279in}}%
\pgfpathlineto{\pgfqpoint{2.922187in}{0.854988in}}%
\pgfpathlineto{\pgfqpoint{2.913824in}{0.850086in}}%
\pgfpathlineto{\pgfqpoint{2.901077in}{0.870542in}}%
\pgfpathlineto{\pgfqpoint{2.893588in}{0.882495in}}%
\pgfpathclose%
\pgfusepath{fill}%
\end{pgfscope}%
\begin{pgfscope}%
\pgfpathrectangle{\pgfqpoint{0.100000in}{0.100000in}}{\pgfqpoint{3.608454in}{2.310000in}}%
\pgfusepath{clip}%
\pgfsetbuttcap%
\pgfsetmiterjoin%
\definecolor{currentfill}{rgb}{0.000000,0.509804,0.745098}%
\pgfsetfillcolor{currentfill}%
\pgfsetlinewidth{0.000000pt}%
\definecolor{currentstroke}{rgb}{0.000000,0.000000,0.000000}%
\pgfsetstrokecolor{currentstroke}%
\pgfsetstrokeopacity{0.000000}%
\pgfsetdash{}{0pt}%
\pgfpathmoveto{\pgfqpoint{2.515690in}{0.998385in}}%
\pgfpathlineto{\pgfqpoint{2.512313in}{0.997036in}}%
\pgfpathlineto{\pgfqpoint{2.512839in}{0.989031in}}%
\pgfpathlineto{\pgfqpoint{2.503832in}{0.987282in}}%
\pgfpathlineto{\pgfqpoint{2.483100in}{0.985946in}}%
\pgfpathlineto{\pgfqpoint{2.482530in}{0.995135in}}%
\pgfpathlineto{\pgfqpoint{2.478692in}{1.001827in}}%
\pgfpathlineto{\pgfqpoint{2.477335in}{1.022510in}}%
\pgfpathlineto{\pgfqpoint{2.473863in}{1.022295in}}%
\pgfpathlineto{\pgfqpoint{2.473129in}{1.034113in}}%
\pgfpathlineto{\pgfqpoint{2.483127in}{1.034794in}}%
\pgfpathlineto{\pgfqpoint{2.507414in}{1.036455in}}%
\pgfpathlineto{\pgfqpoint{2.508135in}{1.025655in}}%
\pgfpathlineto{\pgfqpoint{2.512709in}{1.025980in}}%
\pgfpathlineto{\pgfqpoint{2.515503in}{1.018069in}}%
\pgfpathlineto{\pgfqpoint{2.515690in}{0.998385in}}%
\pgfpathclose%
\pgfusepath{fill}%
\end{pgfscope}%
\begin{pgfscope}%
\pgfpathrectangle{\pgfqpoint{0.100000in}{0.100000in}}{\pgfqpoint{3.608454in}{2.310000in}}%
\pgfusepath{clip}%
\pgfsetbuttcap%
\pgfsetmiterjoin%
\definecolor{currentfill}{rgb}{0.000000,0.564706,0.717647}%
\pgfsetfillcolor{currentfill}%
\pgfsetlinewidth{0.000000pt}%
\definecolor{currentstroke}{rgb}{0.000000,0.000000,0.000000}%
\pgfsetstrokecolor{currentstroke}%
\pgfsetstrokeopacity{0.000000}%
\pgfsetdash{}{0pt}%
\pgfpathmoveto{\pgfqpoint{2.595636in}{0.988190in}}%
\pgfpathlineto{\pgfqpoint{2.553484in}{0.986177in}}%
\pgfpathlineto{\pgfqpoint{2.553769in}{0.997504in}}%
\pgfpathlineto{\pgfqpoint{2.554995in}{1.031701in}}%
\pgfpathlineto{\pgfqpoint{2.547499in}{1.040339in}}%
\pgfpathlineto{\pgfqpoint{2.561622in}{1.041248in}}%
\pgfpathlineto{\pgfqpoint{2.560066in}{1.064151in}}%
\pgfpathlineto{\pgfqpoint{2.556697in}{1.071568in}}%
\pgfpathlineto{\pgfqpoint{2.559562in}{1.077181in}}%
\pgfpathlineto{\pgfqpoint{2.575892in}{1.080309in}}%
\pgfpathlineto{\pgfqpoint{2.580566in}{1.079844in}}%
\pgfpathlineto{\pgfqpoint{2.585581in}{1.074349in}}%
\pgfpathlineto{\pgfqpoint{2.601253in}{1.080261in}}%
\pgfpathlineto{\pgfqpoint{2.607711in}{1.078660in}}%
\pgfpathlineto{\pgfqpoint{2.608880in}{1.074941in}}%
\pgfpathlineto{\pgfqpoint{2.610505in}{1.048649in}}%
\pgfpathlineto{\pgfqpoint{2.611745in}{1.044790in}}%
\pgfpathlineto{\pgfqpoint{2.612717in}{1.030317in}}%
\pgfpathlineto{\pgfqpoint{2.607222in}{1.025513in}}%
\pgfpathlineto{\pgfqpoint{2.599104in}{1.027892in}}%
\pgfpathlineto{\pgfqpoint{2.599703in}{1.020694in}}%
\pgfpathlineto{\pgfqpoint{2.593933in}{1.008935in}}%
\pgfpathlineto{\pgfqpoint{2.595636in}{0.988190in}}%
\pgfpathclose%
\pgfusepath{fill}%
\end{pgfscope}%
\begin{pgfscope}%
\pgfpathrectangle{\pgfqpoint{0.100000in}{0.100000in}}{\pgfqpoint{3.608454in}{2.310000in}}%
\pgfusepath{clip}%
\pgfsetbuttcap%
\pgfsetmiterjoin%
\definecolor{currentfill}{rgb}{0.000000,0.549020,0.725490}%
\pgfsetfillcolor{currentfill}%
\pgfsetlinewidth{0.000000pt}%
\definecolor{currentstroke}{rgb}{0.000000,0.000000,0.000000}%
\pgfsetstrokecolor{currentstroke}%
\pgfsetstrokeopacity{0.000000}%
\pgfsetdash{}{0pt}%
\pgfpathmoveto{\pgfqpoint{1.975277in}{0.471174in}}%
\pgfpathlineto{\pgfqpoint{1.976992in}{0.475338in}}%
\pgfpathlineto{\pgfqpoint{1.978118in}{0.480166in}}%
\pgfpathlineto{\pgfqpoint{1.984185in}{0.481426in}}%
\pgfpathlineto{\pgfqpoint{1.994367in}{0.489945in}}%
\pgfpathlineto{\pgfqpoint{1.996428in}{0.487179in}}%
\pgfpathlineto{\pgfqpoint{1.984727in}{0.479179in}}%
\pgfpathlineto{\pgfqpoint{1.975277in}{0.471174in}}%
\pgfpathclose%
\pgfusepath{fill}%
\end{pgfscope}%
\begin{pgfscope}%
\pgfpathrectangle{\pgfqpoint{0.100000in}{0.100000in}}{\pgfqpoint{3.608454in}{2.310000in}}%
\pgfusepath{clip}%
\pgfsetbuttcap%
\pgfsetmiterjoin%
\definecolor{currentfill}{rgb}{0.000000,0.549020,0.725490}%
\pgfsetfillcolor{currentfill}%
\pgfsetlinewidth{0.000000pt}%
\definecolor{currentstroke}{rgb}{0.000000,0.000000,0.000000}%
\pgfsetstrokecolor{currentstroke}%
\pgfsetstrokeopacity{0.000000}%
\pgfsetdash{}{0pt}%
\pgfpathmoveto{\pgfqpoint{2.013559in}{0.517174in}}%
\pgfpathlineto{\pgfqpoint{2.002833in}{0.511803in}}%
\pgfpathlineto{\pgfqpoint{1.996542in}{0.517508in}}%
\pgfpathlineto{\pgfqpoint{1.995156in}{0.524126in}}%
\pgfpathlineto{\pgfqpoint{1.989574in}{0.521754in}}%
\pgfpathlineto{\pgfqpoint{1.994165in}{0.512797in}}%
\pgfpathlineto{\pgfqpoint{2.008150in}{0.500988in}}%
\pgfpathlineto{\pgfqpoint{1.995428in}{0.494588in}}%
\pgfpathlineto{\pgfqpoint{1.992059in}{0.491469in}}%
\pgfpathlineto{\pgfqpoint{1.981043in}{0.498344in}}%
\pgfpathlineto{\pgfqpoint{1.973994in}{0.506774in}}%
\pgfpathlineto{\pgfqpoint{1.966594in}{0.506597in}}%
\pgfpathlineto{\pgfqpoint{1.962232in}{0.510712in}}%
\pgfpathlineto{\pgfqpoint{1.954355in}{0.511185in}}%
\pgfpathlineto{\pgfqpoint{1.938836in}{0.498314in}}%
\pgfpathlineto{\pgfqpoint{1.926522in}{0.503678in}}%
\pgfpathlineto{\pgfqpoint{1.926187in}{0.508771in}}%
\pgfpathlineto{\pgfqpoint{1.916523in}{0.511256in}}%
\pgfpathlineto{\pgfqpoint{1.911040in}{0.521565in}}%
\pgfpathlineto{\pgfqpoint{1.925742in}{0.532805in}}%
\pgfpathlineto{\pgfqpoint{1.913527in}{0.548723in}}%
\pgfpathlineto{\pgfqpoint{1.923720in}{0.556660in}}%
\pgfpathlineto{\pgfqpoint{1.950391in}{0.577980in}}%
\pgfpathlineto{\pgfqpoint{1.953286in}{0.593627in}}%
\pgfpathlineto{\pgfqpoint{1.957611in}{0.597327in}}%
\pgfpathlineto{\pgfqpoint{1.976342in}{0.597283in}}%
\pgfpathlineto{\pgfqpoint{1.981908in}{0.592727in}}%
\pgfpathlineto{\pgfqpoint{1.997763in}{0.573057in}}%
\pgfpathlineto{\pgfqpoint{1.992065in}{0.566132in}}%
\pgfpathlineto{\pgfqpoint{2.014894in}{0.542939in}}%
\pgfpathlineto{\pgfqpoint{2.013559in}{0.517174in}}%
\pgfpathclose%
\pgfusepath{fill}%
\end{pgfscope}%
\begin{pgfscope}%
\pgfpathrectangle{\pgfqpoint{0.100000in}{0.100000in}}{\pgfqpoint{3.608454in}{2.310000in}}%
\pgfusepath{clip}%
\pgfsetbuttcap%
\pgfsetmiterjoin%
\definecolor{currentfill}{rgb}{0.000000,0.882353,0.558824}%
\pgfsetfillcolor{currentfill}%
\pgfsetlinewidth{0.000000pt}%
\definecolor{currentstroke}{rgb}{0.000000,0.000000,0.000000}%
\pgfsetstrokecolor{currentstroke}%
\pgfsetstrokeopacity{0.000000}%
\pgfsetdash{}{0pt}%
\pgfpathmoveto{\pgfqpoint{2.362240in}{0.870230in}}%
\pgfpathlineto{\pgfqpoint{2.344274in}{0.869656in}}%
\pgfpathlineto{\pgfqpoint{2.342635in}{0.869578in}}%
\pgfpathlineto{\pgfqpoint{2.341651in}{0.899996in}}%
\pgfpathlineto{\pgfqpoint{2.306319in}{0.899456in}}%
\pgfpathlineto{\pgfqpoint{2.307139in}{0.891576in}}%
\pgfpathlineto{\pgfqpoint{2.297645in}{0.883786in}}%
\pgfpathlineto{\pgfqpoint{2.297141in}{0.880461in}}%
\pgfpathlineto{\pgfqpoint{2.286293in}{0.889498in}}%
\pgfpathlineto{\pgfqpoint{2.282169in}{0.889435in}}%
\pgfpathlineto{\pgfqpoint{2.280766in}{0.896956in}}%
\pgfpathlineto{\pgfqpoint{2.284812in}{0.906848in}}%
\pgfpathlineto{\pgfqpoint{2.282183in}{0.933447in}}%
\pgfpathlineto{\pgfqpoint{2.274161in}{0.944083in}}%
\pgfpathlineto{\pgfqpoint{2.272954in}{0.950762in}}%
\pgfpathlineto{\pgfqpoint{2.281822in}{0.951221in}}%
\pgfpathlineto{\pgfqpoint{2.288597in}{0.951635in}}%
\pgfpathlineto{\pgfqpoint{2.289533in}{0.964998in}}%
\pgfpathlineto{\pgfqpoint{2.289234in}{0.985811in}}%
\pgfpathlineto{\pgfqpoint{2.322083in}{0.986111in}}%
\pgfpathlineto{\pgfqpoint{2.322365in}{0.966382in}}%
\pgfpathlineto{\pgfqpoint{2.333604in}{0.966070in}}%
\pgfpathlineto{\pgfqpoint{2.335508in}{0.959654in}}%
\pgfpathlineto{\pgfqpoint{2.339665in}{0.959110in}}%
\pgfpathlineto{\pgfqpoint{2.341075in}{0.950676in}}%
\pgfpathlineto{\pgfqpoint{2.349626in}{0.945687in}}%
\pgfpathlineto{\pgfqpoint{2.361253in}{0.945877in}}%
\pgfpathlineto{\pgfqpoint{2.356734in}{0.953261in}}%
\pgfpathlineto{\pgfqpoint{2.361668in}{0.958818in}}%
\pgfpathlineto{\pgfqpoint{2.372346in}{0.959262in}}%
\pgfpathlineto{\pgfqpoint{2.377959in}{0.956142in}}%
\pgfpathlineto{\pgfqpoint{2.376340in}{0.951894in}}%
\pgfpathlineto{\pgfqpoint{2.364126in}{0.947506in}}%
\pgfpathlineto{\pgfqpoint{2.369367in}{0.943780in}}%
\pgfpathlineto{\pgfqpoint{2.365618in}{0.938943in}}%
\pgfpathlineto{\pgfqpoint{2.370618in}{0.931411in}}%
\pgfpathlineto{\pgfqpoint{2.361535in}{0.931873in}}%
\pgfpathlineto{\pgfqpoint{2.356747in}{0.922152in}}%
\pgfpathlineto{\pgfqpoint{2.362379in}{0.917175in}}%
\pgfpathlineto{\pgfqpoint{2.356517in}{0.914773in}}%
\pgfpathlineto{\pgfqpoint{2.362531in}{0.900120in}}%
\pgfpathlineto{\pgfqpoint{2.362887in}{0.893699in}}%
\pgfpathlineto{\pgfqpoint{2.366452in}{0.887645in}}%
\pgfpathlineto{\pgfqpoint{2.366771in}{0.881093in}}%
\pgfpathlineto{\pgfqpoint{2.362240in}{0.870230in}}%
\pgfpathclose%
\pgfusepath{fill}%
\end{pgfscope}%
\begin{pgfscope}%
\pgfpathrectangle{\pgfqpoint{0.100000in}{0.100000in}}{\pgfqpoint{3.608454in}{2.310000in}}%
\pgfusepath{clip}%
\pgfsetbuttcap%
\pgfsetmiterjoin%
\definecolor{currentfill}{rgb}{0.000000,0.686275,0.656863}%
\pgfsetfillcolor{currentfill}%
\pgfsetlinewidth{0.000000pt}%
\definecolor{currentstroke}{rgb}{0.000000,0.000000,0.000000}%
\pgfsetstrokecolor{currentstroke}%
\pgfsetstrokeopacity{0.000000}%
\pgfsetdash{}{0pt}%
\pgfpathmoveto{\pgfqpoint{3.332649in}{1.866641in}}%
\pgfpathlineto{\pgfqpoint{3.296075in}{1.852707in}}%
\pgfpathlineto{\pgfqpoint{3.295690in}{1.856839in}}%
\pgfpathlineto{\pgfqpoint{3.286142in}{1.855787in}}%
\pgfpathlineto{\pgfqpoint{3.277231in}{1.863034in}}%
\pgfpathlineto{\pgfqpoint{3.280930in}{1.867522in}}%
\pgfpathlineto{\pgfqpoint{3.276788in}{1.878840in}}%
\pgfpathlineto{\pgfqpoint{3.263105in}{1.873951in}}%
\pgfpathlineto{\pgfqpoint{3.262366in}{1.879593in}}%
\pgfpathlineto{\pgfqpoint{3.242124in}{1.938669in}}%
\pgfpathlineto{\pgfqpoint{3.236756in}{1.940890in}}%
\pgfpathlineto{\pgfqpoint{3.268366in}{1.947882in}}%
\pgfpathlineto{\pgfqpoint{3.312126in}{1.959898in}}%
\pgfpathlineto{\pgfqpoint{3.314088in}{1.952894in}}%
\pgfpathlineto{\pgfqpoint{3.313379in}{1.946070in}}%
\pgfpathlineto{\pgfqpoint{3.316446in}{1.944043in}}%
\pgfpathlineto{\pgfqpoint{3.316899in}{1.928969in}}%
\pgfpathlineto{\pgfqpoint{3.323703in}{1.920977in}}%
\pgfpathlineto{\pgfqpoint{3.324626in}{1.910739in}}%
\pgfpathlineto{\pgfqpoint{3.327709in}{1.903478in}}%
\pgfpathlineto{\pgfqpoint{3.324873in}{1.895982in}}%
\pgfpathlineto{\pgfqpoint{3.324985in}{1.884098in}}%
\pgfpathlineto{\pgfqpoint{3.332649in}{1.866641in}}%
\pgfpathclose%
\pgfusepath{fill}%
\end{pgfscope}%
\begin{pgfscope}%
\pgfpathrectangle{\pgfqpoint{0.100000in}{0.100000in}}{\pgfqpoint{3.608454in}{2.310000in}}%
\pgfusepath{clip}%
\pgfsetbuttcap%
\pgfsetmiterjoin%
\definecolor{currentfill}{rgb}{0.000000,0.556863,0.721569}%
\pgfsetfillcolor{currentfill}%
\pgfsetlinewidth{0.000000pt}%
\definecolor{currentstroke}{rgb}{0.000000,0.000000,0.000000}%
\pgfsetstrokecolor{currentstroke}%
\pgfsetstrokeopacity{0.000000}%
\pgfsetdash{}{0pt}%
\pgfpathmoveto{\pgfqpoint{1.678348in}{0.726376in}}%
\pgfpathlineto{\pgfqpoint{1.644176in}{0.728867in}}%
\pgfpathlineto{\pgfqpoint{1.650167in}{0.808902in}}%
\pgfpathlineto{\pgfqpoint{1.655619in}{0.808558in}}%
\pgfpathlineto{\pgfqpoint{1.658603in}{0.843265in}}%
\pgfpathlineto{\pgfqpoint{1.624223in}{0.845743in}}%
\pgfpathlineto{\pgfqpoint{1.626580in}{0.880308in}}%
\pgfpathlineto{\pgfqpoint{1.635380in}{0.879674in}}%
\pgfpathlineto{\pgfqpoint{1.695447in}{0.875598in}}%
\pgfpathlineto{\pgfqpoint{1.692995in}{0.840738in}}%
\pgfpathlineto{\pgfqpoint{1.690000in}{0.806092in}}%
\pgfpathlineto{\pgfqpoint{1.684601in}{0.806456in}}%
\pgfpathlineto{\pgfqpoint{1.681284in}{0.762092in}}%
\pgfpathlineto{\pgfqpoint{1.678348in}{0.726376in}}%
\pgfpathclose%
\pgfusepath{fill}%
\end{pgfscope}%
\begin{pgfscope}%
\pgfpathrectangle{\pgfqpoint{0.100000in}{0.100000in}}{\pgfqpoint{3.608454in}{2.310000in}}%
\pgfusepath{clip}%
\pgfsetbuttcap%
\pgfsetmiterjoin%
\definecolor{currentfill}{rgb}{0.000000,0.827451,0.586275}%
\pgfsetfillcolor{currentfill}%
\pgfsetlinewidth{0.000000pt}%
\definecolor{currentstroke}{rgb}{0.000000,0.000000,0.000000}%
\pgfsetstrokecolor{currentstroke}%
\pgfsetstrokeopacity{0.000000}%
\pgfsetdash{}{0pt}%
\pgfpathmoveto{\pgfqpoint{0.779307in}{0.362404in}}%
\pgfpathlineto{\pgfqpoint{0.779048in}{0.364878in}}%
\pgfpathlineto{\pgfqpoint{0.781778in}{0.364420in}}%
\pgfpathlineto{\pgfqpoint{0.784239in}{0.359816in}}%
\pgfpathlineto{\pgfqpoint{0.779307in}{0.362404in}}%
\pgfpathclose%
\pgfusepath{fill}%
\end{pgfscope}%
\begin{pgfscope}%
\pgfpathrectangle{\pgfqpoint{0.100000in}{0.100000in}}{\pgfqpoint{3.608454in}{2.310000in}}%
\pgfusepath{clip}%
\pgfsetbuttcap%
\pgfsetmiterjoin%
\definecolor{currentfill}{rgb}{0.000000,0.827451,0.586275}%
\pgfsetfillcolor{currentfill}%
\pgfsetlinewidth{0.000000pt}%
\definecolor{currentstroke}{rgb}{0.000000,0.000000,0.000000}%
\pgfsetstrokecolor{currentstroke}%
\pgfsetstrokeopacity{0.000000}%
\pgfsetdash{}{0pt}%
\pgfpathmoveto{\pgfqpoint{0.771381in}{0.365345in}}%
\pgfpathlineto{\pgfqpoint{0.773191in}{0.367159in}}%
\pgfpathlineto{\pgfqpoint{0.777403in}{0.367463in}}%
\pgfpathlineto{\pgfqpoint{0.776562in}{0.365684in}}%
\pgfpathlineto{\pgfqpoint{0.771381in}{0.365345in}}%
\pgfpathclose%
\pgfusepath{fill}%
\end{pgfscope}%
\begin{pgfscope}%
\pgfpathrectangle{\pgfqpoint{0.100000in}{0.100000in}}{\pgfqpoint{3.608454in}{2.310000in}}%
\pgfusepath{clip}%
\pgfsetbuttcap%
\pgfsetmiterjoin%
\definecolor{currentfill}{rgb}{0.000000,0.827451,0.586275}%
\pgfsetfillcolor{currentfill}%
\pgfsetlinewidth{0.000000pt}%
\definecolor{currentstroke}{rgb}{0.000000,0.000000,0.000000}%
\pgfsetstrokecolor{currentstroke}%
\pgfsetstrokeopacity{0.000000}%
\pgfsetdash{}{0pt}%
\pgfpathmoveto{\pgfqpoint{0.764308in}{0.397983in}}%
\pgfpathlineto{\pgfqpoint{0.760327in}{0.401259in}}%
\pgfpathlineto{\pgfqpoint{0.761612in}{0.402864in}}%
\pgfpathlineto{\pgfqpoint{0.765003in}{0.403792in}}%
\pgfpathlineto{\pgfqpoint{0.766993in}{0.403276in}}%
\pgfpathlineto{\pgfqpoint{0.768215in}{0.404903in}}%
\pgfpathlineto{\pgfqpoint{0.771380in}{0.405564in}}%
\pgfpathlineto{\pgfqpoint{0.773253in}{0.404050in}}%
\pgfpathlineto{\pgfqpoint{0.777652in}{0.403284in}}%
\pgfpathlineto{\pgfqpoint{0.777113in}{0.402156in}}%
\pgfpathlineto{\pgfqpoint{0.772733in}{0.402967in}}%
\pgfpathlineto{\pgfqpoint{0.767839in}{0.402629in}}%
\pgfpathlineto{\pgfqpoint{0.767479in}{0.399727in}}%
\pgfpathlineto{\pgfqpoint{0.764308in}{0.397983in}}%
\pgfpathclose%
\pgfusepath{fill}%
\end{pgfscope}%
\begin{pgfscope}%
\pgfpathrectangle{\pgfqpoint{0.100000in}{0.100000in}}{\pgfqpoint{3.608454in}{2.310000in}}%
\pgfusepath{clip}%
\pgfsetbuttcap%
\pgfsetmiterjoin%
\definecolor{currentfill}{rgb}{0.000000,0.827451,0.586275}%
\pgfsetfillcolor{currentfill}%
\pgfsetlinewidth{0.000000pt}%
\definecolor{currentstroke}{rgb}{0.000000,0.000000,0.000000}%
\pgfsetstrokecolor{currentstroke}%
\pgfsetstrokeopacity{0.000000}%
\pgfsetdash{}{0pt}%
\pgfpathmoveto{\pgfqpoint{0.833159in}{0.390083in}}%
\pgfpathlineto{\pgfqpoint{0.831710in}{0.387820in}}%
\pgfpathlineto{\pgfqpoint{0.834237in}{0.387444in}}%
\pgfpathlineto{\pgfqpoint{0.834627in}{0.385754in}}%
\pgfpathlineto{\pgfqpoint{0.836745in}{0.385201in}}%
\pgfpathlineto{\pgfqpoint{0.837269in}{0.381931in}}%
\pgfpathlineto{\pgfqpoint{0.835377in}{0.380107in}}%
\pgfpathlineto{\pgfqpoint{0.832974in}{0.379753in}}%
\pgfpathlineto{\pgfqpoint{0.827040in}{0.383181in}}%
\pgfpathlineto{\pgfqpoint{0.822060in}{0.382223in}}%
\pgfpathlineto{\pgfqpoint{0.818304in}{0.385966in}}%
\pgfpathlineto{\pgfqpoint{0.817822in}{0.388703in}}%
\pgfpathlineto{\pgfqpoint{0.819792in}{0.390087in}}%
\pgfpathlineto{\pgfqpoint{0.823330in}{0.389436in}}%
\pgfpathlineto{\pgfqpoint{0.823863in}{0.391532in}}%
\pgfpathlineto{\pgfqpoint{0.826549in}{0.387455in}}%
\pgfpathlineto{\pgfqpoint{0.828702in}{0.391954in}}%
\pgfpathlineto{\pgfqpoint{0.832601in}{0.391664in}}%
\pgfpathlineto{\pgfqpoint{0.834521in}{0.394379in}}%
\pgfpathlineto{\pgfqpoint{0.835969in}{0.392848in}}%
\pgfpathlineto{\pgfqpoint{0.836007in}{0.391059in}}%
\pgfpathlineto{\pgfqpoint{0.833159in}{0.390083in}}%
\pgfpathclose%
\pgfusepath{fill}%
\end{pgfscope}%
\begin{pgfscope}%
\pgfpathrectangle{\pgfqpoint{0.100000in}{0.100000in}}{\pgfqpoint{3.608454in}{2.310000in}}%
\pgfusepath{clip}%
\pgfsetbuttcap%
\pgfsetmiterjoin%
\definecolor{currentfill}{rgb}{0.000000,0.827451,0.586275}%
\pgfsetfillcolor{currentfill}%
\pgfsetlinewidth{0.000000pt}%
\definecolor{currentstroke}{rgb}{0.000000,0.000000,0.000000}%
\pgfsetstrokecolor{currentstroke}%
\pgfsetstrokeopacity{0.000000}%
\pgfsetdash{}{0pt}%
\pgfpathmoveto{\pgfqpoint{0.824220in}{0.377345in}}%
\pgfpathlineto{\pgfqpoint{0.824880in}{0.374202in}}%
\pgfpathlineto{\pgfqpoint{0.822730in}{0.374227in}}%
\pgfpathlineto{\pgfqpoint{0.819725in}{0.370152in}}%
\pgfpathlineto{\pgfqpoint{0.823688in}{0.368053in}}%
\pgfpathlineto{\pgfqpoint{0.818723in}{0.365758in}}%
\pgfpathlineto{\pgfqpoint{0.816629in}{0.366618in}}%
\pgfpathlineto{\pgfqpoint{0.815348in}{0.369220in}}%
\pgfpathlineto{\pgfqpoint{0.813787in}{0.368393in}}%
\pgfpathlineto{\pgfqpoint{0.813447in}{0.366368in}}%
\pgfpathlineto{\pgfqpoint{0.811430in}{0.365767in}}%
\pgfpathlineto{\pgfqpoint{0.803663in}{0.368137in}}%
\pgfpathlineto{\pgfqpoint{0.804036in}{0.365808in}}%
\pgfpathlineto{\pgfqpoint{0.801470in}{0.365130in}}%
\pgfpathlineto{\pgfqpoint{0.797979in}{0.368165in}}%
\pgfpathlineto{\pgfqpoint{0.796078in}{0.368564in}}%
\pgfpathlineto{\pgfqpoint{0.791325in}{0.365921in}}%
\pgfpathlineto{\pgfqpoint{0.786527in}{0.365586in}}%
\pgfpathlineto{\pgfqpoint{0.787213in}{0.367927in}}%
\pgfpathlineto{\pgfqpoint{0.792491in}{0.369127in}}%
\pgfpathlineto{\pgfqpoint{0.792763in}{0.372545in}}%
\pgfpathlineto{\pgfqpoint{0.789745in}{0.372900in}}%
\pgfpathlineto{\pgfqpoint{0.786003in}{0.370933in}}%
\pgfpathlineto{\pgfqpoint{0.784989in}{0.376973in}}%
\pgfpathlineto{\pgfqpoint{0.787640in}{0.381629in}}%
\pgfpathlineto{\pgfqpoint{0.786375in}{0.384997in}}%
\pgfpathlineto{\pgfqpoint{0.792271in}{0.388607in}}%
\pgfpathlineto{\pgfqpoint{0.796911in}{0.388654in}}%
\pgfpathlineto{\pgfqpoint{0.799734in}{0.387836in}}%
\pgfpathlineto{\pgfqpoint{0.801573in}{0.385770in}}%
\pgfpathlineto{\pgfqpoint{0.799429in}{0.380434in}}%
\pgfpathlineto{\pgfqpoint{0.803009in}{0.382399in}}%
\pgfpathlineto{\pgfqpoint{0.804618in}{0.387969in}}%
\pgfpathlineto{\pgfqpoint{0.807031in}{0.389222in}}%
\pgfpathlineto{\pgfqpoint{0.809285in}{0.388827in}}%
\pgfpathlineto{\pgfqpoint{0.809591in}{0.384951in}}%
\pgfpathlineto{\pgfqpoint{0.811916in}{0.384419in}}%
\pgfpathlineto{\pgfqpoint{0.812085in}{0.387811in}}%
\pgfpathlineto{\pgfqpoint{0.816863in}{0.385495in}}%
\pgfpathlineto{\pgfqpoint{0.818200in}{0.383011in}}%
\pgfpathlineto{\pgfqpoint{0.820527in}{0.380666in}}%
\pgfpathlineto{\pgfqpoint{0.818276in}{0.378662in}}%
\pgfpathlineto{\pgfqpoint{0.824441in}{0.378263in}}%
\pgfpathlineto{\pgfqpoint{0.824220in}{0.377345in}}%
\pgfpathclose%
\pgfusepath{fill}%
\end{pgfscope}%
\begin{pgfscope}%
\pgfpathrectangle{\pgfqpoint{0.100000in}{0.100000in}}{\pgfqpoint{3.608454in}{2.310000in}}%
\pgfusepath{clip}%
\pgfsetbuttcap%
\pgfsetmiterjoin%
\definecolor{currentfill}{rgb}{0.000000,0.827451,0.586275}%
\pgfsetfillcolor{currentfill}%
\pgfsetlinewidth{0.000000pt}%
\definecolor{currentstroke}{rgb}{0.000000,0.000000,0.000000}%
\pgfsetstrokecolor{currentstroke}%
\pgfsetstrokeopacity{0.000000}%
\pgfsetdash{}{0pt}%
\pgfpathmoveto{\pgfqpoint{0.831242in}{0.406233in}}%
\pgfpathlineto{\pgfqpoint{0.829444in}{0.406598in}}%
\pgfpathlineto{\pgfqpoint{0.826973in}{0.404781in}}%
\pgfpathlineto{\pgfqpoint{0.821584in}{0.404692in}}%
\pgfpathlineto{\pgfqpoint{0.820126in}{0.405792in}}%
\pgfpathlineto{\pgfqpoint{0.817124in}{0.404505in}}%
\pgfpathlineto{\pgfqpoint{0.814345in}{0.404191in}}%
\pgfpathlineto{\pgfqpoint{0.812992in}{0.402013in}}%
\pgfpathlineto{\pgfqpoint{0.811166in}{0.403256in}}%
\pgfpathlineto{\pgfqpoint{0.810441in}{0.400723in}}%
\pgfpathlineto{\pgfqpoint{0.803324in}{0.399430in}}%
\pgfpathlineto{\pgfqpoint{0.800398in}{0.399952in}}%
\pgfpathlineto{\pgfqpoint{0.797438in}{0.402092in}}%
\pgfpathlineto{\pgfqpoint{0.794014in}{0.402817in}}%
\pgfpathlineto{\pgfqpoint{0.793204in}{0.401308in}}%
\pgfpathlineto{\pgfqpoint{0.788959in}{0.401438in}}%
\pgfpathlineto{\pgfqpoint{0.786334in}{0.400255in}}%
\pgfpathlineto{\pgfqpoint{0.785696in}{0.403237in}}%
\pgfpathlineto{\pgfqpoint{0.783361in}{0.401415in}}%
\pgfpathlineto{\pgfqpoint{0.780693in}{0.402239in}}%
\pgfpathlineto{\pgfqpoint{0.779576in}{0.400441in}}%
\pgfpathlineto{\pgfqpoint{0.777277in}{0.402089in}}%
\pgfpathlineto{\pgfqpoint{0.781062in}{0.403097in}}%
\pgfpathlineto{\pgfqpoint{0.784195in}{0.405252in}}%
\pgfpathlineto{\pgfqpoint{0.786031in}{0.403802in}}%
\pgfpathlineto{\pgfqpoint{0.789736in}{0.403394in}}%
\pgfpathlineto{\pgfqpoint{0.794932in}{0.410045in}}%
\pgfpathlineto{\pgfqpoint{0.797005in}{0.409992in}}%
\pgfpathlineto{\pgfqpoint{0.798949in}{0.408471in}}%
\pgfpathlineto{\pgfqpoint{0.800501in}{0.409249in}}%
\pgfpathlineto{\pgfqpoint{0.804050in}{0.407001in}}%
\pgfpathlineto{\pgfqpoint{0.808196in}{0.405860in}}%
\pgfpathlineto{\pgfqpoint{0.809601in}{0.406866in}}%
\pgfpathlineto{\pgfqpoint{0.811859in}{0.406633in}}%
\pgfpathlineto{\pgfqpoint{0.815079in}{0.410909in}}%
\pgfpathlineto{\pgfqpoint{0.820929in}{0.406477in}}%
\pgfpathlineto{\pgfqpoint{0.822419in}{0.408486in}}%
\pgfpathlineto{\pgfqpoint{0.826771in}{0.405229in}}%
\pgfpathlineto{\pgfqpoint{0.828822in}{0.408037in}}%
\pgfpathlineto{\pgfqpoint{0.831242in}{0.406233in}}%
\pgfpathclose%
\pgfusepath{fill}%
\end{pgfscope}%
\begin{pgfscope}%
\pgfpathrectangle{\pgfqpoint{0.100000in}{0.100000in}}{\pgfqpoint{3.608454in}{2.310000in}}%
\pgfusepath{clip}%
\pgfsetbuttcap%
\pgfsetmiterjoin%
\definecolor{currentfill}{rgb}{0.000000,0.454902,0.772549}%
\pgfsetfillcolor{currentfill}%
\pgfsetlinewidth{0.000000pt}%
\definecolor{currentstroke}{rgb}{0.000000,0.000000,0.000000}%
\pgfsetstrokecolor{currentstroke}%
\pgfsetstrokeopacity{0.000000}%
\pgfsetdash{}{0pt}%
\pgfpathmoveto{\pgfqpoint{1.204580in}{1.868127in}}%
\pgfpathlineto{\pgfqpoint{1.202252in}{1.853038in}}%
\pgfpathlineto{\pgfqpoint{1.198018in}{1.855881in}}%
\pgfpathlineto{\pgfqpoint{1.191695in}{1.873351in}}%
\pgfpathlineto{\pgfqpoint{1.187877in}{1.878039in}}%
\pgfpathlineto{\pgfqpoint{1.181983in}{1.875682in}}%
\pgfpathlineto{\pgfqpoint{1.164838in}{1.878825in}}%
\pgfpathlineto{\pgfqpoint{1.165886in}{1.884470in}}%
\pgfpathlineto{\pgfqpoint{1.150185in}{1.887641in}}%
\pgfpathlineto{\pgfqpoint{1.145219in}{1.892062in}}%
\pgfpathlineto{\pgfqpoint{1.146489in}{1.905873in}}%
\pgfpathlineto{\pgfqpoint{1.138589in}{1.907407in}}%
\pgfpathlineto{\pgfqpoint{1.139942in}{1.914263in}}%
\pgfpathlineto{\pgfqpoint{1.133161in}{1.915579in}}%
\pgfpathlineto{\pgfqpoint{1.137082in}{1.935988in}}%
\pgfpathlineto{\pgfqpoint{1.135886in}{1.943764in}}%
\pgfpathlineto{\pgfqpoint{1.131774in}{1.943515in}}%
\pgfpathlineto{\pgfqpoint{1.126427in}{1.948422in}}%
\pgfpathlineto{\pgfqpoint{1.128763in}{1.959930in}}%
\pgfpathlineto{\pgfqpoint{1.137125in}{1.963079in}}%
\pgfpathlineto{\pgfqpoint{1.140024in}{1.967674in}}%
\pgfpathlineto{\pgfqpoint{1.157736in}{1.964285in}}%
\pgfpathlineto{\pgfqpoint{1.165857in}{1.971182in}}%
\pgfpathlineto{\pgfqpoint{1.173618in}{1.970061in}}%
\pgfpathlineto{\pgfqpoint{1.176413in}{1.966592in}}%
\pgfpathlineto{\pgfqpoint{1.185357in}{1.963304in}}%
\pgfpathlineto{\pgfqpoint{1.192609in}{1.969852in}}%
\pgfpathlineto{\pgfqpoint{1.199884in}{1.972842in}}%
\pgfpathlineto{\pgfqpoint{1.203995in}{1.980934in}}%
\pgfpathlineto{\pgfqpoint{1.208958in}{1.985065in}}%
\pgfpathlineto{\pgfqpoint{1.210052in}{1.990143in}}%
\pgfpathlineto{\pgfqpoint{1.225529in}{1.987742in}}%
\pgfpathlineto{\pgfqpoint{1.261011in}{1.981452in}}%
\pgfpathlineto{\pgfqpoint{1.268315in}{1.982523in}}%
\pgfpathlineto{\pgfqpoint{1.302428in}{1.976762in}}%
\pgfpathlineto{\pgfqpoint{1.301322in}{1.969905in}}%
\pgfpathlineto{\pgfqpoint{1.302894in}{1.962642in}}%
\pgfpathlineto{\pgfqpoint{1.313126in}{1.960922in}}%
\pgfpathlineto{\pgfqpoint{1.311739in}{1.954347in}}%
\pgfpathlineto{\pgfqpoint{1.307165in}{1.954995in}}%
\pgfpathlineto{\pgfqpoint{1.301786in}{1.941660in}}%
\pgfpathlineto{\pgfqpoint{1.299700in}{1.927978in}}%
\pgfpathlineto{\pgfqpoint{1.292921in}{1.929061in}}%
\pgfpathlineto{\pgfqpoint{1.278245in}{1.924515in}}%
\pgfpathlineto{\pgfqpoint{1.275970in}{1.910859in}}%
\pgfpathlineto{\pgfqpoint{1.268744in}{1.912078in}}%
\pgfpathlineto{\pgfqpoint{1.266400in}{1.898228in}}%
\pgfpathlineto{\pgfqpoint{1.281086in}{1.895779in}}%
\pgfpathlineto{\pgfqpoint{1.278870in}{1.882470in}}%
\pgfpathlineto{\pgfqpoint{1.228507in}{1.890208in}}%
\pgfpathlineto{\pgfqpoint{1.209266in}{1.894331in}}%
\pgfpathlineto{\pgfqpoint{1.204580in}{1.868127in}}%
\pgfpathclose%
\pgfusepath{fill}%
\end{pgfscope}%
\begin{pgfscope}%
\pgfpathrectangle{\pgfqpoint{0.100000in}{0.100000in}}{\pgfqpoint{3.608454in}{2.310000in}}%
\pgfusepath{clip}%
\pgfsetbuttcap%
\pgfsetmiterjoin%
\definecolor{currentfill}{rgb}{0.000000,0.509804,0.745098}%
\pgfsetfillcolor{currentfill}%
\pgfsetlinewidth{0.000000pt}%
\definecolor{currentstroke}{rgb}{0.000000,0.000000,0.000000}%
\pgfsetstrokecolor{currentstroke}%
\pgfsetstrokeopacity{0.000000}%
\pgfsetdash{}{0pt}%
\pgfpathmoveto{\pgfqpoint{1.908588in}{1.422459in}}%
\pgfpathlineto{\pgfqpoint{1.894644in}{1.422961in}}%
\pgfpathlineto{\pgfqpoint{1.881222in}{1.423490in}}%
\pgfpathlineto{\pgfqpoint{1.883772in}{1.478524in}}%
\pgfpathlineto{\pgfqpoint{1.937428in}{1.476563in}}%
\pgfpathlineto{\pgfqpoint{1.936670in}{1.449037in}}%
\pgfpathlineto{\pgfqpoint{1.909610in}{1.449945in}}%
\pgfpathlineto{\pgfqpoint{1.908588in}{1.422459in}}%
\pgfpathclose%
\pgfusepath{fill}%
\end{pgfscope}%
\begin{pgfscope}%
\pgfpathrectangle{\pgfqpoint{0.100000in}{0.100000in}}{\pgfqpoint{3.608454in}{2.310000in}}%
\pgfusepath{clip}%
\pgfsetbuttcap%
\pgfsetmiterjoin%
\definecolor{currentfill}{rgb}{0.000000,0.576471,0.711765}%
\pgfsetfillcolor{currentfill}%
\pgfsetlinewidth{0.000000pt}%
\definecolor{currentstroke}{rgb}{0.000000,0.000000,0.000000}%
\pgfsetstrokecolor{currentstroke}%
\pgfsetstrokeopacity{0.000000}%
\pgfsetdash{}{0pt}%
\pgfpathmoveto{\pgfqpoint{0.695319in}{1.898980in}}%
\pgfpathlineto{\pgfqpoint{0.668725in}{1.906595in}}%
\pgfpathlineto{\pgfqpoint{0.616244in}{1.922404in}}%
\pgfpathlineto{\pgfqpoint{0.579411in}{1.933417in}}%
\pgfpathlineto{\pgfqpoint{0.585151in}{1.943983in}}%
\pgfpathlineto{\pgfqpoint{0.586510in}{1.951287in}}%
\pgfpathlineto{\pgfqpoint{0.595256in}{1.957298in}}%
\pgfpathlineto{\pgfqpoint{0.603007in}{1.965607in}}%
\pgfpathlineto{\pgfqpoint{0.605729in}{1.977984in}}%
\pgfpathlineto{\pgfqpoint{0.605387in}{1.987819in}}%
\pgfpathlineto{\pgfqpoint{0.608267in}{1.997300in}}%
\pgfpathlineto{\pgfqpoint{0.613506in}{2.004223in}}%
\pgfpathlineto{\pgfqpoint{0.616042in}{2.014692in}}%
\pgfpathlineto{\pgfqpoint{0.620648in}{2.021281in}}%
\pgfpathlineto{\pgfqpoint{0.656754in}{2.010285in}}%
\pgfpathlineto{\pgfqpoint{0.695564in}{1.998853in}}%
\pgfpathlineto{\pgfqpoint{0.693514in}{1.994463in}}%
\pgfpathlineto{\pgfqpoint{0.686387in}{1.969750in}}%
\pgfpathlineto{\pgfqpoint{0.706253in}{1.964177in}}%
\pgfpathlineto{\pgfqpoint{0.711901in}{1.958542in}}%
\pgfpathlineto{\pgfqpoint{0.710132in}{1.952205in}}%
\pgfpathlineto{\pgfqpoint{0.723444in}{1.948508in}}%
\pgfpathlineto{\pgfqpoint{0.715935in}{1.921940in}}%
\pgfpathlineto{\pgfqpoint{0.709441in}{1.923753in}}%
\pgfpathlineto{\pgfqpoint{0.703763in}{1.903845in}}%
\pgfpathlineto{\pgfqpoint{0.697108in}{1.905730in}}%
\pgfpathlineto{\pgfqpoint{0.695319in}{1.898980in}}%
\pgfpathclose%
\pgfusepath{fill}%
\end{pgfscope}%
\begin{pgfscope}%
\pgfpathrectangle{\pgfqpoint{0.100000in}{0.100000in}}{\pgfqpoint{3.608454in}{2.310000in}}%
\pgfusepath{clip}%
\pgfsetbuttcap%
\pgfsetmiterjoin%
\definecolor{currentfill}{rgb}{0.000000,0.423529,0.788235}%
\pgfsetfillcolor{currentfill}%
\pgfsetlinewidth{0.000000pt}%
\definecolor{currentstroke}{rgb}{0.000000,0.000000,0.000000}%
\pgfsetstrokecolor{currentstroke}%
\pgfsetstrokeopacity{0.000000}%
\pgfsetdash{}{0pt}%
\pgfpathmoveto{\pgfqpoint{1.716323in}{1.398352in}}%
\pgfpathlineto{\pgfqpoint{1.677687in}{1.401208in}}%
\pgfpathlineto{\pgfqpoint{1.677720in}{1.401672in}}%
\pgfpathlineto{\pgfqpoint{1.682362in}{1.462771in}}%
\pgfpathlineto{\pgfqpoint{1.726124in}{1.459646in}}%
\pgfpathlineto{\pgfqpoint{1.724173in}{1.432221in}}%
\pgfpathlineto{\pgfqpoint{1.718916in}{1.432564in}}%
\pgfpathlineto{\pgfqpoint{1.716323in}{1.398352in}}%
\pgfpathclose%
\pgfusepath{fill}%
\end{pgfscope}%
\begin{pgfscope}%
\pgfpathrectangle{\pgfqpoint{0.100000in}{0.100000in}}{\pgfqpoint{3.608454in}{2.310000in}}%
\pgfusepath{clip}%
\pgfsetbuttcap%
\pgfsetmiterjoin%
\definecolor{currentfill}{rgb}{0.000000,0.815686,0.592157}%
\pgfsetfillcolor{currentfill}%
\pgfsetlinewidth{0.000000pt}%
\definecolor{currentstroke}{rgb}{0.000000,0.000000,0.000000}%
\pgfsetstrokecolor{currentstroke}%
\pgfsetstrokeopacity{0.000000}%
\pgfsetdash{}{0pt}%
\pgfpathmoveto{\pgfqpoint{0.908989in}{1.065097in}}%
\pgfpathlineto{\pgfqpoint{0.914070in}{1.089859in}}%
\pgfpathlineto{\pgfqpoint{0.933245in}{1.183825in}}%
\pgfpathlineto{\pgfqpoint{0.937616in}{1.205156in}}%
\pgfpathlineto{\pgfqpoint{0.941162in}{1.217180in}}%
\pgfpathlineto{\pgfqpoint{0.940117in}{1.223821in}}%
\pgfpathlineto{\pgfqpoint{0.943617in}{1.227923in}}%
\pgfpathlineto{\pgfqpoint{0.949023in}{1.226105in}}%
\pgfpathlineto{\pgfqpoint{0.951426in}{1.230336in}}%
\pgfpathlineto{\pgfqpoint{0.961661in}{1.233095in}}%
\pgfpathlineto{\pgfqpoint{0.978736in}{1.235692in}}%
\pgfpathlineto{\pgfqpoint{0.986930in}{1.238736in}}%
\pgfpathlineto{\pgfqpoint{0.991491in}{1.243095in}}%
\pgfpathlineto{\pgfqpoint{0.991172in}{1.251855in}}%
\pgfpathlineto{\pgfqpoint{0.994707in}{1.264046in}}%
\pgfpathlineto{\pgfqpoint{1.005762in}{1.288169in}}%
\pgfpathlineto{\pgfqpoint{0.983512in}{1.292540in}}%
\pgfpathlineto{\pgfqpoint{0.991804in}{1.334673in}}%
\pgfpathlineto{\pgfqpoint{1.005205in}{1.331999in}}%
\pgfpathlineto{\pgfqpoint{1.024916in}{1.327610in}}%
\pgfpathlineto{\pgfqpoint{1.104844in}{1.313028in}}%
\pgfpathlineto{\pgfqpoint{1.130720in}{1.308550in}}%
\pgfpathlineto{\pgfqpoint{1.122123in}{1.294263in}}%
\pgfpathlineto{\pgfqpoint{1.113627in}{1.291412in}}%
\pgfpathlineto{\pgfqpoint{1.109948in}{1.282813in}}%
\pgfpathlineto{\pgfqpoint{1.099188in}{1.278642in}}%
\pgfpathlineto{\pgfqpoint{1.091741in}{1.280236in}}%
\pgfpathlineto{\pgfqpoint{1.087638in}{1.275829in}}%
\pgfpathlineto{\pgfqpoint{1.078577in}{1.274542in}}%
\pgfpathlineto{\pgfqpoint{1.116943in}{1.267777in}}%
\pgfpathlineto{\pgfqpoint{1.079528in}{1.054057in}}%
\pgfpathlineto{\pgfqpoint{1.073305in}{1.055337in}}%
\pgfpathlineto{\pgfqpoint{1.063896in}{1.064728in}}%
\pgfpathlineto{\pgfqpoint{1.055533in}{1.068101in}}%
\pgfpathlineto{\pgfqpoint{1.049066in}{1.073939in}}%
\pgfpathlineto{\pgfqpoint{1.039948in}{1.073968in}}%
\pgfpathlineto{\pgfqpoint{1.031414in}{1.079449in}}%
\pgfpathlineto{\pgfqpoint{1.022683in}{1.074223in}}%
\pgfpathlineto{\pgfqpoint{1.015879in}{1.056517in}}%
\pgfpathlineto{\pgfqpoint{1.031005in}{1.053711in}}%
\pgfpathlineto{\pgfqpoint{1.027746in}{1.042105in}}%
\pgfpathlineto{\pgfqpoint{1.012835in}{1.044864in}}%
\pgfpathlineto{\pgfqpoint{0.985165in}{1.053983in}}%
\pgfpathlineto{\pgfqpoint{0.975479in}{1.042443in}}%
\pgfpathlineto{\pgfqpoint{0.947034in}{1.057499in}}%
\pgfpathlineto{\pgfqpoint{0.908989in}{1.065097in}}%
\pgfpathclose%
\pgfusepath{fill}%
\end{pgfscope}%
\begin{pgfscope}%
\pgfpathrectangle{\pgfqpoint{0.100000in}{0.100000in}}{\pgfqpoint{3.608454in}{2.310000in}}%
\pgfusepath{clip}%
\pgfsetbuttcap%
\pgfsetmiterjoin%
\definecolor{currentfill}{rgb}{0.000000,0.219608,0.890196}%
\pgfsetfillcolor{currentfill}%
\pgfsetlinewidth{0.000000pt}%
\definecolor{currentstroke}{rgb}{0.000000,0.000000,0.000000}%
\pgfsetstrokecolor{currentstroke}%
\pgfsetstrokeopacity{0.000000}%
\pgfsetdash{}{0pt}%
\pgfpathmoveto{\pgfqpoint{1.577074in}{1.409804in}}%
\pgfpathlineto{\pgfqpoint{1.571014in}{1.354845in}}%
\pgfpathlineto{\pgfqpoint{1.550467in}{1.357020in}}%
\pgfpathlineto{\pgfqpoint{1.552744in}{1.377402in}}%
\pgfpathlineto{\pgfqpoint{1.474789in}{1.385857in}}%
\pgfpathlineto{\pgfqpoint{1.470992in}{1.351885in}}%
\pgfpathlineto{\pgfqpoint{1.431712in}{1.356218in}}%
\pgfpathlineto{\pgfqpoint{1.436653in}{1.364650in}}%
\pgfpathlineto{\pgfqpoint{1.436137in}{1.369904in}}%
\pgfpathlineto{\pgfqpoint{1.430521in}{1.376213in}}%
\pgfpathlineto{\pgfqpoint{1.421188in}{1.383353in}}%
\pgfpathlineto{\pgfqpoint{1.421683in}{1.386315in}}%
\pgfpathlineto{\pgfqpoint{1.423585in}{1.392666in}}%
\pgfpathlineto{\pgfqpoint{1.424568in}{1.406130in}}%
\pgfpathlineto{\pgfqpoint{1.439656in}{1.412644in}}%
\pgfpathlineto{\pgfqpoint{1.451897in}{1.428337in}}%
\pgfpathlineto{\pgfqpoint{1.444005in}{1.434780in}}%
\pgfpathlineto{\pgfqpoint{1.450957in}{1.440841in}}%
\pgfpathlineto{\pgfqpoint{1.459571in}{1.445102in}}%
\pgfpathlineto{\pgfqpoint{1.461324in}{1.456249in}}%
\pgfpathlineto{\pgfqpoint{1.464403in}{1.459400in}}%
\pgfpathlineto{\pgfqpoint{1.465579in}{1.476906in}}%
\pgfpathlineto{\pgfqpoint{1.501446in}{1.472932in}}%
\pgfpathlineto{\pgfqpoint{1.499141in}{1.452348in}}%
\pgfpathlineto{\pgfqpoint{1.553648in}{1.446628in}}%
\pgfpathlineto{\pgfqpoint{1.580447in}{1.444028in}}%
\pgfpathlineto{\pgfqpoint{1.577074in}{1.409804in}}%
\pgfpathclose%
\pgfusepath{fill}%
\end{pgfscope}%
\begin{pgfscope}%
\pgfpathrectangle{\pgfqpoint{0.100000in}{0.100000in}}{\pgfqpoint{3.608454in}{2.310000in}}%
\pgfusepath{clip}%
\pgfsetbuttcap%
\pgfsetmiterjoin%
\definecolor{currentfill}{rgb}{0.000000,0.701961,0.649020}%
\pgfsetfillcolor{currentfill}%
\pgfsetlinewidth{0.000000pt}%
\definecolor{currentstroke}{rgb}{0.000000,0.000000,0.000000}%
\pgfsetstrokecolor{currentstroke}%
\pgfsetstrokeopacity{0.000000}%
\pgfsetdash{}{0pt}%
\pgfpathmoveto{\pgfqpoint{1.797569in}{0.560061in}}%
\pgfpathlineto{\pgfqpoint{1.748728in}{0.562219in}}%
\pgfpathlineto{\pgfqpoint{1.751098in}{0.605189in}}%
\pgfpathlineto{\pgfqpoint{1.710253in}{0.607672in}}%
\pgfpathlineto{\pgfqpoint{1.713553in}{0.660732in}}%
\pgfpathlineto{\pgfqpoint{1.695622in}{0.661844in}}%
\pgfpathlineto{\pgfqpoint{1.697817in}{0.695195in}}%
\pgfpathlineto{\pgfqpoint{1.755704in}{0.692007in}}%
\pgfpathlineto{\pgfqpoint{1.753800in}{0.658495in}}%
\pgfpathlineto{\pgfqpoint{1.778772in}{0.657165in}}%
\pgfpathlineto{\pgfqpoint{1.777624in}{0.639888in}}%
\pgfpathlineto{\pgfqpoint{1.782309in}{0.639635in}}%
\pgfpathlineto{\pgfqpoint{1.781527in}{0.626427in}}%
\pgfpathlineto{\pgfqpoint{1.787654in}{0.626021in}}%
\pgfpathlineto{\pgfqpoint{1.786496in}{0.603609in}}%
\pgfpathlineto{\pgfqpoint{1.799798in}{0.602976in}}%
\pgfpathlineto{\pgfqpoint{1.797569in}{0.560061in}}%
\pgfpathclose%
\pgfusepath{fill}%
\end{pgfscope}%
\begin{pgfscope}%
\pgfpathrectangle{\pgfqpoint{0.100000in}{0.100000in}}{\pgfqpoint{3.608454in}{2.310000in}}%
\pgfusepath{clip}%
\pgfsetbuttcap%
\pgfsetmiterjoin%
\definecolor{currentfill}{rgb}{0.000000,0.674510,0.662745}%
\pgfsetfillcolor{currentfill}%
\pgfsetlinewidth{0.000000pt}%
\definecolor{currentstroke}{rgb}{0.000000,0.000000,0.000000}%
\pgfsetstrokecolor{currentstroke}%
\pgfsetstrokeopacity{0.000000}%
\pgfsetdash{}{0pt}%
\pgfpathmoveto{\pgfqpoint{1.167197in}{0.891666in}}%
\pgfpathlineto{\pgfqpoint{1.176800in}{0.953138in}}%
\pgfpathlineto{\pgfqpoint{1.245517in}{0.942241in}}%
\pgfpathlineto{\pgfqpoint{1.253754in}{0.941144in}}%
\pgfpathlineto{\pgfqpoint{1.255093in}{0.930635in}}%
\pgfpathlineto{\pgfqpoint{1.253634in}{0.922827in}}%
\pgfpathlineto{\pgfqpoint{1.257767in}{0.919573in}}%
\pgfpathlineto{\pgfqpoint{1.256598in}{0.912844in}}%
\pgfpathlineto{\pgfqpoint{1.259177in}{0.906354in}}%
\pgfpathlineto{\pgfqpoint{1.256509in}{0.894758in}}%
\pgfpathlineto{\pgfqpoint{1.264784in}{0.891644in}}%
\pgfpathlineto{\pgfqpoint{1.285282in}{0.888829in}}%
\pgfpathlineto{\pgfqpoint{1.276635in}{0.824008in}}%
\pgfpathlineto{\pgfqpoint{1.215509in}{0.832612in}}%
\pgfpathlineto{\pgfqpoint{1.210301in}{0.797108in}}%
\pgfpathlineto{\pgfqpoint{1.153619in}{0.805601in}}%
\pgfpathlineto{\pgfqpoint{1.167197in}{0.891666in}}%
\pgfpathclose%
\pgfusepath{fill}%
\end{pgfscope}%
\begin{pgfscope}%
\pgfpathrectangle{\pgfqpoint{0.100000in}{0.100000in}}{\pgfqpoint{3.608454in}{2.310000in}}%
\pgfusepath{clip}%
\pgfsetbuttcap%
\pgfsetmiterjoin%
\definecolor{currentfill}{rgb}{0.000000,0.478431,0.760784}%
\pgfsetfillcolor{currentfill}%
\pgfsetlinewidth{0.000000pt}%
\definecolor{currentstroke}{rgb}{0.000000,0.000000,0.000000}%
\pgfsetstrokecolor{currentstroke}%
\pgfsetstrokeopacity{0.000000}%
\pgfsetdash{}{0pt}%
\pgfpathmoveto{\pgfqpoint{1.983409in}{1.890641in}}%
\pgfpathlineto{\pgfqpoint{1.941869in}{1.891889in}}%
\pgfpathlineto{\pgfqpoint{1.885060in}{1.894397in}}%
\pgfpathlineto{\pgfqpoint{1.886398in}{1.921692in}}%
\pgfpathlineto{\pgfqpoint{1.939707in}{1.919594in}}%
\pgfpathlineto{\pgfqpoint{1.940671in}{1.947331in}}%
\pgfpathlineto{\pgfqpoint{1.959926in}{1.946668in}}%
\pgfpathlineto{\pgfqpoint{1.981936in}{1.946008in}}%
\pgfpathlineto{\pgfqpoint{1.981211in}{1.918353in}}%
\pgfpathlineto{\pgfqpoint{1.981775in}{1.895615in}}%
\pgfpathlineto{\pgfqpoint{1.983409in}{1.890641in}}%
\pgfpathclose%
\pgfusepath{fill}%
\end{pgfscope}%
\begin{pgfscope}%
\pgfpathrectangle{\pgfqpoint{0.100000in}{0.100000in}}{\pgfqpoint{3.608454in}{2.310000in}}%
\pgfusepath{clip}%
\pgfsetbuttcap%
\pgfsetmiterjoin%
\definecolor{currentfill}{rgb}{0.000000,0.556863,0.721569}%
\pgfsetfillcolor{currentfill}%
\pgfsetlinewidth{0.000000pt}%
\definecolor{currentstroke}{rgb}{0.000000,0.000000,0.000000}%
\pgfsetstrokecolor{currentstroke}%
\pgfsetstrokeopacity{0.000000}%
\pgfsetdash{}{0pt}%
\pgfpathmoveto{\pgfqpoint{3.371539in}{1.933332in}}%
\pgfpathlineto{\pgfqpoint{3.373495in}{1.940927in}}%
\pgfpathlineto{\pgfqpoint{3.366415in}{1.942507in}}%
\pgfpathlineto{\pgfqpoint{3.368238in}{1.950258in}}%
\pgfpathlineto{\pgfqpoint{3.358298in}{1.952554in}}%
\pgfpathlineto{\pgfqpoint{3.359777in}{1.957182in}}%
\pgfpathlineto{\pgfqpoint{3.355196in}{1.970349in}}%
\pgfpathlineto{\pgfqpoint{3.368542in}{1.973373in}}%
\pgfpathlineto{\pgfqpoint{3.412317in}{1.985521in}}%
\pgfpathlineto{\pgfqpoint{3.415090in}{2.005571in}}%
\pgfpathlineto{\pgfqpoint{3.427643in}{2.008700in}}%
\pgfpathlineto{\pgfqpoint{3.428979in}{2.013989in}}%
\pgfpathlineto{\pgfqpoint{3.442334in}{1.974192in}}%
\pgfpathlineto{\pgfqpoint{3.454146in}{1.936802in}}%
\pgfpathlineto{\pgfqpoint{3.453336in}{1.932723in}}%
\pgfpathlineto{\pgfqpoint{3.441678in}{1.930119in}}%
\pgfpathlineto{\pgfqpoint{3.440922in}{1.921126in}}%
\pgfpathlineto{\pgfqpoint{3.433355in}{1.925260in}}%
\pgfpathlineto{\pgfqpoint{3.431085in}{1.927713in}}%
\pgfpathlineto{\pgfqpoint{3.423697in}{1.926204in}}%
\pgfpathlineto{\pgfqpoint{3.422059in}{1.929884in}}%
\pgfpathlineto{\pgfqpoint{3.415486in}{1.930416in}}%
\pgfpathlineto{\pgfqpoint{3.410255in}{1.935145in}}%
\pgfpathlineto{\pgfqpoint{3.405777in}{1.928235in}}%
\pgfpathlineto{\pgfqpoint{3.399587in}{1.926596in}}%
\pgfpathlineto{\pgfqpoint{3.396256in}{1.920610in}}%
\pgfpathlineto{\pgfqpoint{3.398006in}{1.914180in}}%
\pgfpathlineto{\pgfqpoint{3.384973in}{1.910295in}}%
\pgfpathlineto{\pgfqpoint{3.381069in}{1.911013in}}%
\pgfpathlineto{\pgfqpoint{3.384688in}{1.930010in}}%
\pgfpathlineto{\pgfqpoint{3.371539in}{1.933332in}}%
\pgfpathclose%
\pgfusepath{fill}%
\end{pgfscope}%
\begin{pgfscope}%
\pgfpathrectangle{\pgfqpoint{0.100000in}{0.100000in}}{\pgfqpoint{3.608454in}{2.310000in}}%
\pgfusepath{clip}%
\pgfsetbuttcap%
\pgfsetmiterjoin%
\definecolor{currentfill}{rgb}{0.000000,0.325490,0.837255}%
\pgfsetfillcolor{currentfill}%
\pgfsetlinewidth{0.000000pt}%
\definecolor{currentstroke}{rgb}{0.000000,0.000000,0.000000}%
\pgfsetstrokecolor{currentstroke}%
\pgfsetstrokeopacity{0.000000}%
\pgfsetdash{}{0pt}%
\pgfpathmoveto{\pgfqpoint{1.912766in}{1.573654in}}%
\pgfpathlineto{\pgfqpoint{1.914002in}{1.614954in}}%
\pgfpathlineto{\pgfqpoint{1.968471in}{1.613417in}}%
\pgfpathlineto{\pgfqpoint{1.968280in}{1.606524in}}%
\pgfpathlineto{\pgfqpoint{1.967706in}{1.585865in}}%
\pgfpathlineto{\pgfqpoint{1.988196in}{1.585344in}}%
\pgfpathlineto{\pgfqpoint{1.987553in}{1.557852in}}%
\pgfpathlineto{\pgfqpoint{1.939698in}{1.559064in}}%
\pgfpathlineto{\pgfqpoint{1.939982in}{1.572847in}}%
\pgfpathlineto{\pgfqpoint{1.912766in}{1.573654in}}%
\pgfpathclose%
\pgfusepath{fill}%
\end{pgfscope}%
\begin{pgfscope}%
\pgfpathrectangle{\pgfqpoint{0.100000in}{0.100000in}}{\pgfqpoint{3.608454in}{2.310000in}}%
\pgfusepath{clip}%
\pgfsetbuttcap%
\pgfsetmiterjoin%
\definecolor{currentfill}{rgb}{0.000000,0.450980,0.774510}%
\pgfsetfillcolor{currentfill}%
\pgfsetlinewidth{0.000000pt}%
\definecolor{currentstroke}{rgb}{0.000000,0.000000,0.000000}%
\pgfsetstrokecolor{currentstroke}%
\pgfsetstrokeopacity{0.000000}%
\pgfsetdash{}{0pt}%
\pgfpathmoveto{\pgfqpoint{1.580447in}{1.444028in}}%
\pgfpathlineto{\pgfqpoint{1.553648in}{1.446628in}}%
\pgfpathlineto{\pgfqpoint{1.499141in}{1.452348in}}%
\pgfpathlineto{\pgfqpoint{1.501446in}{1.472932in}}%
\pgfpathlineto{\pgfqpoint{1.465579in}{1.476906in}}%
\pgfpathlineto{\pgfqpoint{1.459655in}{1.485514in}}%
\pgfpathlineto{\pgfqpoint{1.455629in}{1.496066in}}%
\pgfpathlineto{\pgfqpoint{1.446700in}{1.523501in}}%
\pgfpathlineto{\pgfqpoint{1.440255in}{1.533604in}}%
\pgfpathlineto{\pgfqpoint{1.440561in}{1.538635in}}%
\pgfpathlineto{\pgfqpoint{1.494704in}{1.532257in}}%
\pgfpathlineto{\pgfqpoint{1.533391in}{1.528263in}}%
\pgfpathlineto{\pgfqpoint{1.595804in}{1.522015in}}%
\pgfpathlineto{\pgfqpoint{1.591803in}{1.484401in}}%
\pgfpathlineto{\pgfqpoint{1.598652in}{1.483744in}}%
\pgfpathlineto{\pgfqpoint{1.594536in}{1.442692in}}%
\pgfpathlineto{\pgfqpoint{1.580447in}{1.444028in}}%
\pgfpathclose%
\pgfusepath{fill}%
\end{pgfscope}%
\begin{pgfscope}%
\pgfpathrectangle{\pgfqpoint{0.100000in}{0.100000in}}{\pgfqpoint{3.608454in}{2.310000in}}%
\pgfusepath{clip}%
\pgfsetbuttcap%
\pgfsetmiterjoin%
\definecolor{currentfill}{rgb}{0.000000,0.741176,0.629412}%
\pgfsetfillcolor{currentfill}%
\pgfsetlinewidth{0.000000pt}%
\definecolor{currentstroke}{rgb}{0.000000,0.000000,0.000000}%
\pgfsetstrokecolor{currentstroke}%
\pgfsetstrokeopacity{0.000000}%
\pgfsetdash{}{0pt}%
\pgfpathmoveto{\pgfqpoint{0.869214in}{2.048910in}}%
\pgfpathlineto{\pgfqpoint{0.875412in}{2.075250in}}%
\pgfpathlineto{\pgfqpoint{0.872020in}{2.076043in}}%
\pgfpathlineto{\pgfqpoint{0.871344in}{2.087217in}}%
\pgfpathlineto{\pgfqpoint{0.864637in}{2.088901in}}%
\pgfpathlineto{\pgfqpoint{0.868296in}{2.100626in}}%
\pgfpathlineto{\pgfqpoint{0.878570in}{2.104452in}}%
\pgfpathlineto{\pgfqpoint{0.882536in}{2.100932in}}%
\pgfpathlineto{\pgfqpoint{0.890708in}{2.100407in}}%
\pgfpathlineto{\pgfqpoint{0.892918in}{2.093767in}}%
\pgfpathlineto{\pgfqpoint{0.898511in}{2.086388in}}%
\pgfpathlineto{\pgfqpoint{0.899012in}{2.076012in}}%
\pgfpathlineto{\pgfqpoint{0.907603in}{2.074322in}}%
\pgfpathlineto{\pgfqpoint{0.909838in}{2.083326in}}%
\pgfpathlineto{\pgfqpoint{0.926998in}{2.079195in}}%
\pgfpathlineto{\pgfqpoint{0.933769in}{2.084665in}}%
\pgfpathlineto{\pgfqpoint{0.949294in}{2.080795in}}%
\pgfpathlineto{\pgfqpoint{0.954967in}{2.104591in}}%
\pgfpathlineto{\pgfqpoint{0.993104in}{2.095521in}}%
\pgfpathlineto{\pgfqpoint{1.027701in}{2.087646in}}%
\pgfpathlineto{\pgfqpoint{1.027942in}{2.078990in}}%
\pgfpathlineto{\pgfqpoint{1.041857in}{2.066086in}}%
\pgfpathlineto{\pgfqpoint{1.041035in}{2.060695in}}%
\pgfpathlineto{\pgfqpoint{1.055012in}{2.060043in}}%
\pgfpathlineto{\pgfqpoint{1.056865in}{2.058071in}}%
\pgfpathlineto{\pgfqpoint{1.049410in}{2.040224in}}%
\pgfpathlineto{\pgfqpoint{1.042825in}{2.030433in}}%
\pgfpathlineto{\pgfqpoint{1.042048in}{2.022019in}}%
\pgfpathlineto{\pgfqpoint{1.037550in}{2.013584in}}%
\pgfpathlineto{\pgfqpoint{1.040540in}{1.999733in}}%
\pgfpathlineto{\pgfqpoint{1.033186in}{1.997675in}}%
\pgfpathlineto{\pgfqpoint{1.028886in}{1.992775in}}%
\pgfpathlineto{\pgfqpoint{1.030719in}{1.983835in}}%
\pgfpathlineto{\pgfqpoint{1.025284in}{1.975956in}}%
\pgfpathlineto{\pgfqpoint{1.018436in}{1.970411in}}%
\pgfpathlineto{\pgfqpoint{1.011182in}{1.972459in}}%
\pgfpathlineto{\pgfqpoint{1.018689in}{1.956057in}}%
\pgfpathlineto{\pgfqpoint{1.012082in}{1.949489in}}%
\pgfpathlineto{\pgfqpoint{0.942731in}{1.965121in}}%
\pgfpathlineto{\pgfqpoint{0.931228in}{1.960671in}}%
\pgfpathlineto{\pgfqpoint{0.923969in}{1.964164in}}%
\pgfpathlineto{\pgfqpoint{0.923373in}{1.975654in}}%
\pgfpathlineto{\pgfqpoint{0.904646in}{1.980156in}}%
\pgfpathlineto{\pgfqpoint{0.916109in}{1.996717in}}%
\pgfpathlineto{\pgfqpoint{0.923281in}{2.004122in}}%
\pgfpathlineto{\pgfqpoint{0.921587in}{2.014227in}}%
\pgfpathlineto{\pgfqpoint{0.914485in}{2.022985in}}%
\pgfpathlineto{\pgfqpoint{0.909977in}{2.024514in}}%
\pgfpathlineto{\pgfqpoint{0.906027in}{2.039422in}}%
\pgfpathlineto{\pgfqpoint{0.869214in}{2.048910in}}%
\pgfpathclose%
\pgfusepath{fill}%
\end{pgfscope}%
\begin{pgfscope}%
\pgfpathrectangle{\pgfqpoint{0.100000in}{0.100000in}}{\pgfqpoint{3.608454in}{2.310000in}}%
\pgfusepath{clip}%
\pgfsetbuttcap%
\pgfsetmiterjoin%
\definecolor{currentfill}{rgb}{0.000000,0.698039,0.650980}%
\pgfsetfillcolor{currentfill}%
\pgfsetlinewidth{0.000000pt}%
\definecolor{currentstroke}{rgb}{0.000000,0.000000,0.000000}%
\pgfsetstrokecolor{currentstroke}%
\pgfsetstrokeopacity{0.000000}%
\pgfsetdash{}{0pt}%
\pgfpathmoveto{\pgfqpoint{2.971461in}{0.696012in}}%
\pgfpathlineto{\pgfqpoint{2.999393in}{0.700496in}}%
\pgfpathlineto{\pgfqpoint{3.004229in}{0.666876in}}%
\pgfpathlineto{\pgfqpoint{3.006579in}{0.647225in}}%
\pgfpathlineto{\pgfqpoint{2.996330in}{0.642369in}}%
\pgfpathlineto{\pgfqpoint{2.995671in}{0.646679in}}%
\pgfpathlineto{\pgfqpoint{2.982209in}{0.644864in}}%
\pgfpathlineto{\pgfqpoint{2.985432in}{0.623511in}}%
\pgfpathlineto{\pgfqpoint{2.976237in}{0.622163in}}%
\pgfpathlineto{\pgfqpoint{2.978125in}{0.608670in}}%
\pgfpathlineto{\pgfqpoint{2.973159in}{0.605106in}}%
\pgfpathlineto{\pgfqpoint{2.965965in}{0.605834in}}%
\pgfpathlineto{\pgfqpoint{2.963258in}{0.602568in}}%
\pgfpathlineto{\pgfqpoint{2.962394in}{0.607479in}}%
\pgfpathlineto{\pgfqpoint{2.957118in}{0.615111in}}%
\pgfpathlineto{\pgfqpoint{2.942926in}{0.613795in}}%
\pgfpathlineto{\pgfqpoint{2.937858in}{0.621000in}}%
\pgfpathlineto{\pgfqpoint{2.930039in}{0.626757in}}%
\pgfpathlineto{\pgfqpoint{2.926843in}{0.632121in}}%
\pgfpathlineto{\pgfqpoint{2.921124in}{0.633342in}}%
\pgfpathlineto{\pgfqpoint{2.913615in}{0.638189in}}%
\pgfpathlineto{\pgfqpoint{2.910647in}{0.649877in}}%
\pgfpathlineto{\pgfqpoint{2.913541in}{0.650279in}}%
\pgfpathlineto{\pgfqpoint{2.915694in}{0.663204in}}%
\pgfpathlineto{\pgfqpoint{2.944970in}{0.667447in}}%
\pgfpathlineto{\pgfqpoint{2.946140in}{0.674377in}}%
\pgfpathlineto{\pgfqpoint{2.951329in}{0.676818in}}%
\pgfpathlineto{\pgfqpoint{2.955684in}{0.671251in}}%
\pgfpathlineto{\pgfqpoint{2.962277in}{0.669805in}}%
\pgfpathlineto{\pgfqpoint{2.965236in}{0.672986in}}%
\pgfpathlineto{\pgfqpoint{2.965631in}{0.682186in}}%
\pgfpathlineto{\pgfqpoint{2.971461in}{0.696012in}}%
\pgfpathclose%
\pgfusepath{fill}%
\end{pgfscope}%
\begin{pgfscope}%
\pgfpathrectangle{\pgfqpoint{0.100000in}{0.100000in}}{\pgfqpoint{3.608454in}{2.310000in}}%
\pgfusepath{clip}%
\pgfsetbuttcap%
\pgfsetmiterjoin%
\definecolor{currentfill}{rgb}{0.000000,0.427451,0.786275}%
\pgfsetfillcolor{currentfill}%
\pgfsetlinewidth{0.000000pt}%
\definecolor{currentstroke}{rgb}{0.000000,0.000000,0.000000}%
\pgfsetstrokecolor{currentstroke}%
\pgfsetstrokeopacity{0.000000}%
\pgfsetdash{}{0pt}%
\pgfpathmoveto{\pgfqpoint{2.817692in}{1.389496in}}%
\pgfpathlineto{\pgfqpoint{2.825080in}{1.387318in}}%
\pgfpathlineto{\pgfqpoint{2.829317in}{1.356171in}}%
\pgfpathlineto{\pgfqpoint{2.812960in}{1.360831in}}%
\pgfpathlineto{\pgfqpoint{2.807405in}{1.358642in}}%
\pgfpathlineto{\pgfqpoint{2.804813in}{1.353816in}}%
\pgfpathlineto{\pgfqpoint{2.798224in}{1.355170in}}%
\pgfpathlineto{\pgfqpoint{2.790830in}{1.362794in}}%
\pgfpathlineto{\pgfqpoint{2.784131in}{1.363939in}}%
\pgfpathlineto{\pgfqpoint{2.777359in}{1.362093in}}%
\pgfpathlineto{\pgfqpoint{2.768060in}{1.364802in}}%
\pgfpathlineto{\pgfqpoint{2.772188in}{1.346181in}}%
\pgfpathlineto{\pgfqpoint{2.760419in}{1.343088in}}%
\pgfpathlineto{\pgfqpoint{2.750154in}{1.334910in}}%
\pgfpathlineto{\pgfqpoint{2.736943in}{1.344983in}}%
\pgfpathlineto{\pgfqpoint{2.735344in}{1.352880in}}%
\pgfpathlineto{\pgfqpoint{2.726746in}{1.347306in}}%
\pgfpathlineto{\pgfqpoint{2.720376in}{1.354526in}}%
\pgfpathlineto{\pgfqpoint{2.728448in}{1.358041in}}%
\pgfpathlineto{\pgfqpoint{2.733971in}{1.365444in}}%
\pgfpathlineto{\pgfqpoint{2.727957in}{1.367391in}}%
\pgfpathlineto{\pgfqpoint{2.729795in}{1.373635in}}%
\pgfpathlineto{\pgfqpoint{2.725607in}{1.378219in}}%
\pgfpathlineto{\pgfqpoint{2.729836in}{1.382914in}}%
\pgfpathlineto{\pgfqpoint{2.726448in}{1.415659in}}%
\pgfpathlineto{\pgfqpoint{2.726071in}{1.419220in}}%
\pgfpathlineto{\pgfqpoint{2.746401in}{1.421668in}}%
\pgfpathlineto{\pgfqpoint{2.757924in}{1.424641in}}%
\pgfpathlineto{\pgfqpoint{2.776762in}{1.425313in}}%
\pgfpathlineto{\pgfqpoint{2.777953in}{1.400497in}}%
\pgfpathlineto{\pgfqpoint{2.786603in}{1.400891in}}%
\pgfpathlineto{\pgfqpoint{2.788388in}{1.383110in}}%
\pgfpathlineto{\pgfqpoint{2.804293in}{1.384898in}}%
\pgfpathlineto{\pgfqpoint{2.817692in}{1.389496in}}%
\pgfpathclose%
\pgfusepath{fill}%
\end{pgfscope}%
\begin{pgfscope}%
\pgfpathrectangle{\pgfqpoint{0.100000in}{0.100000in}}{\pgfqpoint{3.608454in}{2.310000in}}%
\pgfusepath{clip}%
\pgfsetbuttcap%
\pgfsetmiterjoin%
\definecolor{currentfill}{rgb}{0.000000,0.611765,0.694118}%
\pgfsetfillcolor{currentfill}%
\pgfsetlinewidth{0.000000pt}%
\definecolor{currentstroke}{rgb}{0.000000,0.000000,0.000000}%
\pgfsetstrokecolor{currentstroke}%
\pgfsetstrokeopacity{0.000000}%
\pgfsetdash{}{0pt}%
\pgfpathmoveto{\pgfqpoint{2.906878in}{1.433220in}}%
\pgfpathlineto{\pgfqpoint{2.894009in}{1.432157in}}%
\pgfpathlineto{\pgfqpoint{2.893522in}{1.439324in}}%
\pgfpathlineto{\pgfqpoint{2.886708in}{1.438806in}}%
\pgfpathlineto{\pgfqpoint{2.872947in}{1.440094in}}%
\pgfpathlineto{\pgfqpoint{2.872185in}{1.451622in}}%
\pgfpathlineto{\pgfqpoint{2.869385in}{1.458446in}}%
\pgfpathlineto{\pgfqpoint{2.864747in}{1.458144in}}%
\pgfpathlineto{\pgfqpoint{2.864318in}{1.465519in}}%
\pgfpathlineto{\pgfqpoint{2.878244in}{1.466104in}}%
\pgfpathlineto{\pgfqpoint{2.879950in}{1.469292in}}%
\pgfpathlineto{\pgfqpoint{2.878220in}{1.486351in}}%
\pgfpathlineto{\pgfqpoint{2.909284in}{1.489578in}}%
\pgfpathlineto{\pgfqpoint{2.908656in}{1.495198in}}%
\pgfpathlineto{\pgfqpoint{2.928425in}{1.497499in}}%
\pgfpathlineto{\pgfqpoint{2.935644in}{1.495065in}}%
\pgfpathlineto{\pgfqpoint{2.938569in}{1.471342in}}%
\pgfpathlineto{\pgfqpoint{2.933975in}{1.470774in}}%
\pgfpathlineto{\pgfqpoint{2.935564in}{1.458054in}}%
\pgfpathlineto{\pgfqpoint{2.938939in}{1.450447in}}%
\pgfpathlineto{\pgfqpoint{2.927738in}{1.446958in}}%
\pgfpathlineto{\pgfqpoint{2.914676in}{1.445590in}}%
\pgfpathlineto{\pgfqpoint{2.914756in}{1.436835in}}%
\pgfpathlineto{\pgfqpoint{2.907685in}{1.436860in}}%
\pgfpathlineto{\pgfqpoint{2.906878in}{1.433220in}}%
\pgfpathclose%
\pgfusepath{fill}%
\end{pgfscope}%
\begin{pgfscope}%
\pgfpathrectangle{\pgfqpoint{0.100000in}{0.100000in}}{\pgfqpoint{3.608454in}{2.310000in}}%
\pgfusepath{clip}%
\pgfsetbuttcap%
\pgfsetmiterjoin%
\definecolor{currentfill}{rgb}{0.000000,0.458824,0.770588}%
\pgfsetfillcolor{currentfill}%
\pgfsetlinewidth{0.000000pt}%
\definecolor{currentstroke}{rgb}{0.000000,0.000000,0.000000}%
\pgfsetstrokecolor{currentstroke}%
\pgfsetstrokeopacity{0.000000}%
\pgfsetdash{}{0pt}%
\pgfpathmoveto{\pgfqpoint{2.025168in}{1.813601in}}%
\pgfpathlineto{\pgfqpoint{2.000983in}{1.813993in}}%
\pgfpathlineto{\pgfqpoint{2.001203in}{1.827773in}}%
\pgfpathlineto{\pgfqpoint{1.981983in}{1.828333in}}%
\pgfpathlineto{\pgfqpoint{1.982265in}{1.839933in}}%
\pgfpathlineto{\pgfqpoint{1.994509in}{1.839579in}}%
\pgfpathlineto{\pgfqpoint{1.995378in}{1.841947in}}%
\pgfpathlineto{\pgfqpoint{2.024583in}{1.841364in}}%
\pgfpathlineto{\pgfqpoint{2.021780in}{1.845315in}}%
\pgfpathlineto{\pgfqpoint{2.012969in}{1.848244in}}%
\pgfpathlineto{\pgfqpoint{2.003389in}{1.863983in}}%
\pgfpathlineto{\pgfqpoint{2.004519in}{1.867267in}}%
\pgfpathlineto{\pgfqpoint{2.015031in}{1.874970in}}%
\pgfpathlineto{\pgfqpoint{2.018973in}{1.880728in}}%
\pgfpathlineto{\pgfqpoint{2.020229in}{1.889881in}}%
\pgfpathlineto{\pgfqpoint{2.019610in}{1.896790in}}%
\pgfpathlineto{\pgfqpoint{2.036803in}{1.896501in}}%
\pgfpathlineto{\pgfqpoint{2.037372in}{1.889582in}}%
\pgfpathlineto{\pgfqpoint{2.037019in}{1.861775in}}%
\pgfpathlineto{\pgfqpoint{2.037488in}{1.848022in}}%
\pgfpathlineto{\pgfqpoint{2.044431in}{1.847876in}}%
\pgfpathlineto{\pgfqpoint{2.044201in}{1.834191in}}%
\pgfpathlineto{\pgfqpoint{2.041094in}{1.830482in}}%
\pgfpathlineto{\pgfqpoint{2.034935in}{1.834889in}}%
\pgfpathlineto{\pgfqpoint{2.025490in}{1.836778in}}%
\pgfpathlineto{\pgfqpoint{2.025168in}{1.813601in}}%
\pgfpathclose%
\pgfusepath{fill}%
\end{pgfscope}%
\begin{pgfscope}%
\pgfpathrectangle{\pgfqpoint{0.100000in}{0.100000in}}{\pgfqpoint{3.608454in}{2.310000in}}%
\pgfusepath{clip}%
\pgfsetbuttcap%
\pgfsetmiterjoin%
\definecolor{currentfill}{rgb}{0.000000,0.647059,0.676471}%
\pgfsetfillcolor{currentfill}%
\pgfsetlinewidth{0.000000pt}%
\definecolor{currentstroke}{rgb}{0.000000,0.000000,0.000000}%
\pgfsetstrokecolor{currentstroke}%
\pgfsetstrokeopacity{0.000000}%
\pgfsetdash{}{0pt}%
\pgfpathmoveto{\pgfqpoint{2.661344in}{1.549614in}}%
\pgfpathlineto{\pgfqpoint{2.660092in}{1.560783in}}%
\pgfpathlineto{\pgfqpoint{2.636183in}{1.558287in}}%
\pgfpathlineto{\pgfqpoint{2.633441in}{1.583836in}}%
\pgfpathlineto{\pgfqpoint{2.649386in}{1.585311in}}%
\pgfpathlineto{\pgfqpoint{2.650897in}{1.588647in}}%
\pgfpathlineto{\pgfqpoint{2.648568in}{1.609914in}}%
\pgfpathlineto{\pgfqpoint{2.703247in}{1.616030in}}%
\pgfpathlineto{\pgfqpoint{2.706079in}{1.591455in}}%
\pgfpathlineto{\pgfqpoint{2.684290in}{1.589030in}}%
\pgfpathlineto{\pgfqpoint{2.688778in}{1.550091in}}%
\pgfpathlineto{\pgfqpoint{2.668362in}{1.548067in}}%
\pgfpathlineto{\pgfqpoint{2.661344in}{1.549614in}}%
\pgfpathclose%
\pgfusepath{fill}%
\end{pgfscope}%
\begin{pgfscope}%
\pgfpathrectangle{\pgfqpoint{0.100000in}{0.100000in}}{\pgfqpoint{3.608454in}{2.310000in}}%
\pgfusepath{clip}%
\pgfsetbuttcap%
\pgfsetmiterjoin%
\definecolor{currentfill}{rgb}{0.000000,0.000000,1.000000}%
\pgfsetfillcolor{currentfill}%
\pgfsetlinewidth{0.000000pt}%
\definecolor{currentstroke}{rgb}{0.000000,0.000000,0.000000}%
\pgfsetstrokecolor{currentstroke}%
\pgfsetstrokeopacity{0.000000}%
\pgfsetdash{}{0pt}%
\pgfpathmoveto{\pgfqpoint{1.455629in}{1.496066in}}%
\pgfpathlineto{\pgfqpoint{1.459655in}{1.485514in}}%
\pgfpathlineto{\pgfqpoint{1.465579in}{1.476906in}}%
\pgfpathlineto{\pgfqpoint{1.464403in}{1.459400in}}%
\pgfpathlineto{\pgfqpoint{1.461324in}{1.456249in}}%
\pgfpathlineto{\pgfqpoint{1.459571in}{1.445102in}}%
\pgfpathlineto{\pgfqpoint{1.450957in}{1.440841in}}%
\pgfpathlineto{\pgfqpoint{1.444005in}{1.434780in}}%
\pgfpathlineto{\pgfqpoint{1.437752in}{1.434786in}}%
\pgfpathlineto{\pgfqpoint{1.435412in}{1.444400in}}%
\pgfpathlineto{\pgfqpoint{1.426636in}{1.453930in}}%
\pgfpathlineto{\pgfqpoint{1.403922in}{1.457766in}}%
\pgfpathlineto{\pgfqpoint{1.404073in}{1.463909in}}%
\pgfpathlineto{\pgfqpoint{1.407785in}{1.490488in}}%
\pgfpathlineto{\pgfqpoint{1.407647in}{1.498770in}}%
\pgfpathlineto{\pgfqpoint{1.410993in}{1.493320in}}%
\pgfpathlineto{\pgfqpoint{1.426803in}{1.487168in}}%
\pgfpathlineto{\pgfqpoint{1.436560in}{1.486098in}}%
\pgfpathlineto{\pgfqpoint{1.444777in}{1.488383in}}%
\pgfpathlineto{\pgfqpoint{1.447426in}{1.486216in}}%
\pgfpathlineto{\pgfqpoint{1.455629in}{1.496066in}}%
\pgfpathclose%
\pgfusepath{fill}%
\end{pgfscope}%
\begin{pgfscope}%
\pgfpathrectangle{\pgfqpoint{0.100000in}{0.100000in}}{\pgfqpoint{3.608454in}{2.310000in}}%
\pgfusepath{clip}%
\pgfsetbuttcap%
\pgfsetmiterjoin%
\definecolor{currentfill}{rgb}{0.000000,0.384314,0.807843}%
\pgfsetfillcolor{currentfill}%
\pgfsetlinewidth{0.000000pt}%
\definecolor{currentstroke}{rgb}{0.000000,0.000000,0.000000}%
\pgfsetstrokecolor{currentstroke}%
\pgfsetstrokeopacity{0.000000}%
\pgfsetdash{}{0pt}%
\pgfpathmoveto{\pgfqpoint{2.638885in}{1.368939in}}%
\pgfpathlineto{\pgfqpoint{2.635715in}{1.368602in}}%
\pgfpathlineto{\pgfqpoint{2.632842in}{1.391462in}}%
\pgfpathlineto{\pgfqpoint{2.617652in}{1.390119in}}%
\pgfpathlineto{\pgfqpoint{2.616632in}{1.400311in}}%
\pgfpathlineto{\pgfqpoint{2.613308in}{1.399991in}}%
\pgfpathlineto{\pgfqpoint{2.614715in}{1.406613in}}%
\pgfpathlineto{\pgfqpoint{2.612117in}{1.412594in}}%
\pgfpathlineto{\pgfqpoint{2.607655in}{1.455712in}}%
\pgfpathlineto{\pgfqpoint{2.634909in}{1.458446in}}%
\pgfpathlineto{\pgfqpoint{2.634617in}{1.461203in}}%
\pgfpathlineto{\pgfqpoint{2.657467in}{1.463709in}}%
\pgfpathlineto{\pgfqpoint{2.659619in}{1.442049in}}%
\pgfpathlineto{\pgfqpoint{2.676851in}{1.443992in}}%
\pgfpathlineto{\pgfqpoint{2.676915in}{1.431338in}}%
\pgfpathlineto{\pgfqpoint{2.694747in}{1.433296in}}%
\pgfpathlineto{\pgfqpoint{2.697158in}{1.412713in}}%
\pgfpathlineto{\pgfqpoint{2.697801in}{1.407041in}}%
\pgfpathlineto{\pgfqpoint{2.677647in}{1.404907in}}%
\pgfpathlineto{\pgfqpoint{2.675138in}{1.396451in}}%
\pgfpathlineto{\pgfqpoint{2.650805in}{1.393553in}}%
\pgfpathlineto{\pgfqpoint{2.653405in}{1.370465in}}%
\pgfpathlineto{\pgfqpoint{2.638885in}{1.368939in}}%
\pgfpathclose%
\pgfusepath{fill}%
\end{pgfscope}%
\begin{pgfscope}%
\pgfpathrectangle{\pgfqpoint{0.100000in}{0.100000in}}{\pgfqpoint{3.608454in}{2.310000in}}%
\pgfusepath{clip}%
\pgfsetbuttcap%
\pgfsetmiterjoin%
\definecolor{currentfill}{rgb}{0.000000,0.827451,0.586275}%
\pgfsetfillcolor{currentfill}%
\pgfsetlinewidth{0.000000pt}%
\definecolor{currentstroke}{rgb}{0.000000,0.000000,0.000000}%
\pgfsetstrokecolor{currentstroke}%
\pgfsetstrokeopacity{0.000000}%
\pgfsetdash{}{0pt}%
\pgfpathmoveto{\pgfqpoint{0.899087in}{1.016739in}}%
\pgfpathlineto{\pgfqpoint{0.885192in}{0.948795in}}%
\pgfpathlineto{\pgfqpoint{0.877772in}{0.912470in}}%
\pgfpathlineto{\pgfqpoint{0.787923in}{0.968647in}}%
\pgfpathlineto{\pgfqpoint{0.790528in}{0.978335in}}%
\pgfpathlineto{\pgfqpoint{0.797923in}{0.984723in}}%
\pgfpathlineto{\pgfqpoint{0.705665in}{0.997629in}}%
\pgfpathlineto{\pgfqpoint{0.714200in}{1.032869in}}%
\pgfpathlineto{\pgfqpoint{0.715639in}{1.032543in}}%
\pgfpathlineto{\pgfqpoint{0.721788in}{1.059750in}}%
\pgfpathlineto{\pgfqpoint{0.787832in}{1.044715in}}%
\pgfpathlineto{\pgfqpoint{0.809266in}{1.037990in}}%
\pgfpathlineto{\pgfqpoint{0.809954in}{1.027225in}}%
\pgfpathlineto{\pgfqpoint{0.805414in}{1.014538in}}%
\pgfpathlineto{\pgfqpoint{0.807337in}{1.008296in}}%
\pgfpathlineto{\pgfqpoint{0.832599in}{1.002476in}}%
\pgfpathlineto{\pgfqpoint{0.839911in}{1.036150in}}%
\pgfpathlineto{\pgfqpoint{0.859997in}{1.031950in}}%
\pgfpathlineto{\pgfqpoint{0.858565in}{1.025184in}}%
\pgfpathlineto{\pgfqpoint{0.899087in}{1.016739in}}%
\pgfpathclose%
\pgfusepath{fill}%
\end{pgfscope}%
\begin{pgfscope}%
\pgfpathrectangle{\pgfqpoint{3.319565in}{0.488889in}}{\pgfqpoint{0.360845in}{0.924000in}}%
\pgfusepath{clip}%
\pgfsetbuttcap%
\pgfsetmiterjoin%
\definecolor{currentfill}{rgb}{0.000000,0.854902,0.572549}%
\pgfsetfillcolor{currentfill}%
\pgfsetlinewidth{0.000000pt}%
\definecolor{currentstroke}{rgb}{0.000000,0.000000,0.000000}%
\pgfsetstrokecolor{currentstroke}%
\pgfsetstrokeopacity{0.000000}%
\pgfsetdash{}{0pt}%
\pgfpathmoveto{\pgfqpoint{3.581475in}{0.488889in}}%
\pgfpathlineto{\pgfqpoint{3.680411in}{0.488889in}}%
\pgfpathlineto{\pgfqpoint{3.680411in}{0.581289in}}%
\pgfpathlineto{\pgfqpoint{3.581475in}{0.581289in}}%
\pgfpathlineto{\pgfqpoint{3.581475in}{0.488889in}}%
\pgfpathclose%
\pgfusepath{fill}%
\end{pgfscope}%
\begin{pgfscope}%
\pgfpathrectangle{\pgfqpoint{3.319565in}{0.488889in}}{\pgfqpoint{0.360845in}{0.924000in}}%
\pgfusepath{clip}%
\pgfsetbuttcap%
\pgfsetmiterjoin%
\definecolor{currentfill}{rgb}{0.000000,0.815686,0.592157}%
\pgfsetfillcolor{currentfill}%
\pgfsetlinewidth{0.000000pt}%
\definecolor{currentstroke}{rgb}{0.000000,0.000000,0.000000}%
\pgfsetstrokecolor{currentstroke}%
\pgfsetstrokeopacity{0.000000}%
\pgfsetdash{}{0pt}%
\pgfpathmoveto{\pgfqpoint{3.578539in}{0.498129in}}%
\pgfpathlineto{\pgfqpoint{3.680411in}{0.498129in}}%
\pgfpathlineto{\pgfqpoint{3.680411in}{0.590529in}}%
\pgfpathlineto{\pgfqpoint{3.578539in}{0.590529in}}%
\pgfpathlineto{\pgfqpoint{3.578539in}{0.498129in}}%
\pgfpathclose%
\pgfusepath{fill}%
\end{pgfscope}%
\begin{pgfscope}%
\pgfpathrectangle{\pgfqpoint{3.319565in}{0.488889in}}{\pgfqpoint{0.360845in}{0.924000in}}%
\pgfusepath{clip}%
\pgfsetbuttcap%
\pgfsetmiterjoin%
\definecolor{currentfill}{rgb}{0.000000,0.788235,0.605882}%
\pgfsetfillcolor{currentfill}%
\pgfsetlinewidth{0.000000pt}%
\definecolor{currentstroke}{rgb}{0.000000,0.000000,0.000000}%
\pgfsetstrokecolor{currentstroke}%
\pgfsetstrokeopacity{0.000000}%
\pgfsetdash{}{0pt}%
\pgfpathmoveto{\pgfqpoint{3.576250in}{0.507369in}}%
\pgfpathlineto{\pgfqpoint{3.680411in}{0.507369in}}%
\pgfpathlineto{\pgfqpoint{3.680411in}{0.599769in}}%
\pgfpathlineto{\pgfqpoint{3.576250in}{0.599769in}}%
\pgfpathlineto{\pgfqpoint{3.576250in}{0.507369in}}%
\pgfpathclose%
\pgfusepath{fill}%
\end{pgfscope}%
\begin{pgfscope}%
\pgfpathrectangle{\pgfqpoint{3.319565in}{0.488889in}}{\pgfqpoint{0.360845in}{0.924000in}}%
\pgfusepath{clip}%
\pgfsetbuttcap%
\pgfsetmiterjoin%
\definecolor{currentfill}{rgb}{0.000000,0.760784,0.619608}%
\pgfsetfillcolor{currentfill}%
\pgfsetlinewidth{0.000000pt}%
\definecolor{currentstroke}{rgb}{0.000000,0.000000,0.000000}%
\pgfsetstrokecolor{currentstroke}%
\pgfsetstrokeopacity{0.000000}%
\pgfsetdash{}{0pt}%
\pgfpathmoveto{\pgfqpoint{3.574081in}{0.516609in}}%
\pgfpathlineto{\pgfqpoint{3.680411in}{0.516609in}}%
\pgfpathlineto{\pgfqpoint{3.680411in}{0.609009in}}%
\pgfpathlineto{\pgfqpoint{3.574081in}{0.609009in}}%
\pgfpathlineto{\pgfqpoint{3.574081in}{0.516609in}}%
\pgfpathclose%
\pgfusepath{fill}%
\end{pgfscope}%
\begin{pgfscope}%
\pgfpathrectangle{\pgfqpoint{3.319565in}{0.488889in}}{\pgfqpoint{0.360845in}{0.924000in}}%
\pgfusepath{clip}%
\pgfsetbuttcap%
\pgfsetmiterjoin%
\definecolor{currentfill}{rgb}{0.000000,0.749020,0.625490}%
\pgfsetfillcolor{currentfill}%
\pgfsetlinewidth{0.000000pt}%
\definecolor{currentstroke}{rgb}{0.000000,0.000000,0.000000}%
\pgfsetstrokecolor{currentstroke}%
\pgfsetstrokeopacity{0.000000}%
\pgfsetdash{}{0pt}%
\pgfpathmoveto{\pgfqpoint{3.573102in}{0.525849in}}%
\pgfpathlineto{\pgfqpoint{3.680411in}{0.525849in}}%
\pgfpathlineto{\pgfqpoint{3.680411in}{0.618249in}}%
\pgfpathlineto{\pgfqpoint{3.573102in}{0.618249in}}%
\pgfpathlineto{\pgfqpoint{3.573102in}{0.525849in}}%
\pgfpathclose%
\pgfusepath{fill}%
\end{pgfscope}%
\begin{pgfscope}%
\pgfpathrectangle{\pgfqpoint{3.319565in}{0.488889in}}{\pgfqpoint{0.360845in}{0.924000in}}%
\pgfusepath{clip}%
\pgfsetbuttcap%
\pgfsetmiterjoin%
\definecolor{currentfill}{rgb}{0.000000,0.713725,0.643137}%
\pgfsetfillcolor{currentfill}%
\pgfsetlinewidth{0.000000pt}%
\definecolor{currentstroke}{rgb}{0.000000,0.000000,0.000000}%
\pgfsetstrokecolor{currentstroke}%
\pgfsetstrokeopacity{0.000000}%
\pgfsetdash{}{0pt}%
\pgfpathmoveto{\pgfqpoint{3.570208in}{0.535089in}}%
\pgfpathlineto{\pgfqpoint{3.680411in}{0.535089in}}%
\pgfpathlineto{\pgfqpoint{3.680411in}{0.627489in}}%
\pgfpathlineto{\pgfqpoint{3.570208in}{0.627489in}}%
\pgfpathlineto{\pgfqpoint{3.570208in}{0.535089in}}%
\pgfpathclose%
\pgfusepath{fill}%
\end{pgfscope}%
\begin{pgfscope}%
\pgfpathrectangle{\pgfqpoint{3.319565in}{0.488889in}}{\pgfqpoint{0.360845in}{0.924000in}}%
\pgfusepath{clip}%
\pgfsetbuttcap%
\pgfsetmiterjoin%
\definecolor{currentfill}{rgb}{0.000000,0.701961,0.649020}%
\pgfsetfillcolor{currentfill}%
\pgfsetlinewidth{0.000000pt}%
\definecolor{currentstroke}{rgb}{0.000000,0.000000,0.000000}%
\pgfsetstrokecolor{currentstroke}%
\pgfsetstrokeopacity{0.000000}%
\pgfsetdash{}{0pt}%
\pgfpathmoveto{\pgfqpoint{3.569139in}{0.544329in}}%
\pgfpathlineto{\pgfqpoint{3.680411in}{0.544329in}}%
\pgfpathlineto{\pgfqpoint{3.680411in}{0.636729in}}%
\pgfpathlineto{\pgfqpoint{3.569139in}{0.636729in}}%
\pgfpathlineto{\pgfqpoint{3.569139in}{0.544329in}}%
\pgfpathclose%
\pgfusepath{fill}%
\end{pgfscope}%
\begin{pgfscope}%
\pgfpathrectangle{\pgfqpoint{3.319565in}{0.488889in}}{\pgfqpoint{0.360845in}{0.924000in}}%
\pgfusepath{clip}%
\pgfsetbuttcap%
\pgfsetmiterjoin%
\definecolor{currentfill}{rgb}{0.000000,0.690196,0.654902}%
\pgfsetfillcolor{currentfill}%
\pgfsetlinewidth{0.000000pt}%
\definecolor{currentstroke}{rgb}{0.000000,0.000000,0.000000}%
\pgfsetstrokecolor{currentstroke}%
\pgfsetstrokeopacity{0.000000}%
\pgfsetdash{}{0pt}%
\pgfpathmoveto{\pgfqpoint{3.568218in}{0.553569in}}%
\pgfpathlineto{\pgfqpoint{3.680411in}{0.553569in}}%
\pgfpathlineto{\pgfqpoint{3.680411in}{0.645969in}}%
\pgfpathlineto{\pgfqpoint{3.568218in}{0.645969in}}%
\pgfpathlineto{\pgfqpoint{3.568218in}{0.553569in}}%
\pgfpathclose%
\pgfusepath{fill}%
\end{pgfscope}%
\begin{pgfscope}%
\pgfpathrectangle{\pgfqpoint{3.319565in}{0.488889in}}{\pgfqpoint{0.360845in}{0.924000in}}%
\pgfusepath{clip}%
\pgfsetbuttcap%
\pgfsetmiterjoin%
\definecolor{currentfill}{rgb}{0.000000,0.678431,0.660784}%
\pgfsetfillcolor{currentfill}%
\pgfsetlinewidth{0.000000pt}%
\definecolor{currentstroke}{rgb}{0.000000,0.000000,0.000000}%
\pgfsetstrokecolor{currentstroke}%
\pgfsetstrokeopacity{0.000000}%
\pgfsetdash{}{0pt}%
\pgfpathmoveto{\pgfqpoint{3.567329in}{0.562809in}}%
\pgfpathlineto{\pgfqpoint{3.680411in}{0.562809in}}%
\pgfpathlineto{\pgfqpoint{3.680411in}{0.655209in}}%
\pgfpathlineto{\pgfqpoint{3.567329in}{0.655209in}}%
\pgfpathlineto{\pgfqpoint{3.567329in}{0.562809in}}%
\pgfpathclose%
\pgfusepath{fill}%
\end{pgfscope}%
\begin{pgfscope}%
\pgfpathrectangle{\pgfqpoint{3.319565in}{0.488889in}}{\pgfqpoint{0.360845in}{0.924000in}}%
\pgfusepath{clip}%
\pgfsetbuttcap%
\pgfsetmiterjoin%
\definecolor{currentfill}{rgb}{0.000000,0.670588,0.664706}%
\pgfsetfillcolor{currentfill}%
\pgfsetlinewidth{0.000000pt}%
\definecolor{currentstroke}{rgb}{0.000000,0.000000,0.000000}%
\pgfsetstrokecolor{currentstroke}%
\pgfsetstrokeopacity{0.000000}%
\pgfsetdash{}{0pt}%
\pgfpathmoveto{\pgfqpoint{3.566822in}{0.572049in}}%
\pgfpathlineto{\pgfqpoint{3.680411in}{0.572049in}}%
\pgfpathlineto{\pgfqpoint{3.680411in}{0.664449in}}%
\pgfpathlineto{\pgfqpoint{3.566822in}{0.664449in}}%
\pgfpathlineto{\pgfqpoint{3.566822in}{0.572049in}}%
\pgfpathclose%
\pgfusepath{fill}%
\end{pgfscope}%
\begin{pgfscope}%
\pgfpathrectangle{\pgfqpoint{3.319565in}{0.488889in}}{\pgfqpoint{0.360845in}{0.924000in}}%
\pgfusepath{clip}%
\pgfsetbuttcap%
\pgfsetmiterjoin%
\definecolor{currentfill}{rgb}{0.000000,0.666667,0.666667}%
\pgfsetfillcolor{currentfill}%
\pgfsetlinewidth{0.000000pt}%
\definecolor{currentstroke}{rgb}{0.000000,0.000000,0.000000}%
\pgfsetstrokecolor{currentstroke}%
\pgfsetstrokeopacity{0.000000}%
\pgfsetdash{}{0pt}%
\pgfpathmoveto{\pgfqpoint{3.566410in}{0.581289in}}%
\pgfpathlineto{\pgfqpoint{3.680411in}{0.581289in}}%
\pgfpathlineto{\pgfqpoint{3.680411in}{0.673689in}}%
\pgfpathlineto{\pgfqpoint{3.566410in}{0.673689in}}%
\pgfpathlineto{\pgfqpoint{3.566410in}{0.581289in}}%
\pgfpathclose%
\pgfusepath{fill}%
\end{pgfscope}%
\begin{pgfscope}%
\pgfpathrectangle{\pgfqpoint{3.319565in}{0.488889in}}{\pgfqpoint{0.360845in}{0.924000in}}%
\pgfusepath{clip}%
\pgfsetbuttcap%
\pgfsetmiterjoin%
\definecolor{currentfill}{rgb}{0.000000,0.654902,0.672549}%
\pgfsetfillcolor{currentfill}%
\pgfsetlinewidth{0.000000pt}%
\definecolor{currentstroke}{rgb}{0.000000,0.000000,0.000000}%
\pgfsetstrokecolor{currentstroke}%
\pgfsetstrokeopacity{0.000000}%
\pgfsetdash{}{0pt}%
\pgfpathmoveto{\pgfqpoint{3.565309in}{0.590529in}}%
\pgfpathlineto{\pgfqpoint{3.680411in}{0.590529in}}%
\pgfpathlineto{\pgfqpoint{3.680411in}{0.682929in}}%
\pgfpathlineto{\pgfqpoint{3.565309in}{0.682929in}}%
\pgfpathlineto{\pgfqpoint{3.565309in}{0.590529in}}%
\pgfpathclose%
\pgfusepath{fill}%
\end{pgfscope}%
\begin{pgfscope}%
\pgfpathrectangle{\pgfqpoint{3.319565in}{0.488889in}}{\pgfqpoint{0.360845in}{0.924000in}}%
\pgfusepath{clip}%
\pgfsetbuttcap%
\pgfsetmiterjoin%
\definecolor{currentfill}{rgb}{0.000000,0.643137,0.678431}%
\pgfsetfillcolor{currentfill}%
\pgfsetlinewidth{0.000000pt}%
\definecolor{currentstroke}{rgb}{0.000000,0.000000,0.000000}%
\pgfsetstrokecolor{currentstroke}%
\pgfsetstrokeopacity{0.000000}%
\pgfsetdash{}{0pt}%
\pgfpathmoveto{\pgfqpoint{3.564406in}{0.599769in}}%
\pgfpathlineto{\pgfqpoint{3.680411in}{0.599769in}}%
\pgfpathlineto{\pgfqpoint{3.680411in}{0.692169in}}%
\pgfpathlineto{\pgfqpoint{3.564406in}{0.692169in}}%
\pgfpathlineto{\pgfqpoint{3.564406in}{0.599769in}}%
\pgfpathclose%
\pgfusepath{fill}%
\end{pgfscope}%
\begin{pgfscope}%
\pgfpathrectangle{\pgfqpoint{3.319565in}{0.488889in}}{\pgfqpoint{0.360845in}{0.924000in}}%
\pgfusepath{clip}%
\pgfsetbuttcap%
\pgfsetmiterjoin%
\definecolor{currentfill}{rgb}{0.000000,0.631373,0.684314}%
\pgfsetfillcolor{currentfill}%
\pgfsetlinewidth{0.000000pt}%
\definecolor{currentstroke}{rgb}{0.000000,0.000000,0.000000}%
\pgfsetstrokecolor{currentstroke}%
\pgfsetstrokeopacity{0.000000}%
\pgfsetdash{}{0pt}%
\pgfpathmoveto{\pgfqpoint{3.563660in}{0.609009in}}%
\pgfpathlineto{\pgfqpoint{3.680411in}{0.609009in}}%
\pgfpathlineto{\pgfqpoint{3.680411in}{0.701409in}}%
\pgfpathlineto{\pgfqpoint{3.563660in}{0.701409in}}%
\pgfpathlineto{\pgfqpoint{3.563660in}{0.609009in}}%
\pgfpathclose%
\pgfusepath{fill}%
\end{pgfscope}%
\begin{pgfscope}%
\pgfpathrectangle{\pgfqpoint{3.319565in}{0.488889in}}{\pgfqpoint{0.360845in}{0.924000in}}%
\pgfusepath{clip}%
\pgfsetbuttcap%
\pgfsetmiterjoin%
\definecolor{currentfill}{rgb}{0.000000,0.627451,0.686275}%
\pgfsetfillcolor{currentfill}%
\pgfsetlinewidth{0.000000pt}%
\definecolor{currentstroke}{rgb}{0.000000,0.000000,0.000000}%
\pgfsetstrokecolor{currentstroke}%
\pgfsetstrokeopacity{0.000000}%
\pgfsetdash{}{0pt}%
\pgfpathmoveto{\pgfqpoint{3.563249in}{0.618249in}}%
\pgfpathlineto{\pgfqpoint{3.680411in}{0.618249in}}%
\pgfpathlineto{\pgfqpoint{3.680411in}{0.710649in}}%
\pgfpathlineto{\pgfqpoint{3.563249in}{0.710649in}}%
\pgfpathlineto{\pgfqpoint{3.563249in}{0.618249in}}%
\pgfpathclose%
\pgfusepath{fill}%
\end{pgfscope}%
\begin{pgfscope}%
\pgfpathrectangle{\pgfqpoint{3.319565in}{0.488889in}}{\pgfqpoint{0.360845in}{0.924000in}}%
\pgfusepath{clip}%
\pgfsetbuttcap%
\pgfsetmiterjoin%
\definecolor{currentfill}{rgb}{0.000000,0.607843,0.696078}%
\pgfsetfillcolor{currentfill}%
\pgfsetlinewidth{0.000000pt}%
\definecolor{currentstroke}{rgb}{0.000000,0.000000,0.000000}%
\pgfsetstrokecolor{currentstroke}%
\pgfsetstrokeopacity{0.000000}%
\pgfsetdash{}{0pt}%
\pgfpathmoveto{\pgfqpoint{3.561733in}{0.627489in}}%
\pgfpathlineto{\pgfqpoint{3.680411in}{0.627489in}}%
\pgfpathlineto{\pgfqpoint{3.680411in}{0.719889in}}%
\pgfpathlineto{\pgfqpoint{3.561733in}{0.719889in}}%
\pgfpathlineto{\pgfqpoint{3.561733in}{0.627489in}}%
\pgfpathclose%
\pgfusepath{fill}%
\end{pgfscope}%
\begin{pgfscope}%
\pgfpathrectangle{\pgfqpoint{3.319565in}{0.488889in}}{\pgfqpoint{0.360845in}{0.924000in}}%
\pgfusepath{clip}%
\pgfsetbuttcap%
\pgfsetmiterjoin%
\definecolor{currentfill}{rgb}{0.000000,0.603922,0.698039}%
\pgfsetfillcolor{currentfill}%
\pgfsetlinewidth{0.000000pt}%
\definecolor{currentstroke}{rgb}{0.000000,0.000000,0.000000}%
\pgfsetstrokecolor{currentstroke}%
\pgfsetstrokeopacity{0.000000}%
\pgfsetdash{}{0pt}%
\pgfpathmoveto{\pgfqpoint{3.561363in}{0.636729in}}%
\pgfpathlineto{\pgfqpoint{3.680411in}{0.636729in}}%
\pgfpathlineto{\pgfqpoint{3.680411in}{0.729129in}}%
\pgfpathlineto{\pgfqpoint{3.561363in}{0.729129in}}%
\pgfpathlineto{\pgfqpoint{3.561363in}{0.636729in}}%
\pgfpathclose%
\pgfusepath{fill}%
\end{pgfscope}%
\begin{pgfscope}%
\pgfpathrectangle{\pgfqpoint{3.319565in}{0.488889in}}{\pgfqpoint{0.360845in}{0.924000in}}%
\pgfusepath{clip}%
\pgfsetbuttcap%
\pgfsetmiterjoin%
\definecolor{currentfill}{rgb}{0.000000,0.600000,0.700000}%
\pgfsetfillcolor{currentfill}%
\pgfsetlinewidth{0.000000pt}%
\definecolor{currentstroke}{rgb}{0.000000,0.000000,0.000000}%
\pgfsetstrokecolor{currentstroke}%
\pgfsetstrokeopacity{0.000000}%
\pgfsetdash{}{0pt}%
\pgfpathmoveto{\pgfqpoint{3.560871in}{0.645969in}}%
\pgfpathlineto{\pgfqpoint{3.680411in}{0.645969in}}%
\pgfpathlineto{\pgfqpoint{3.680411in}{0.738369in}}%
\pgfpathlineto{\pgfqpoint{3.560871in}{0.738369in}}%
\pgfpathlineto{\pgfqpoint{3.560871in}{0.645969in}}%
\pgfpathclose%
\pgfusepath{fill}%
\end{pgfscope}%
\begin{pgfscope}%
\pgfpathrectangle{\pgfqpoint{3.319565in}{0.488889in}}{\pgfqpoint{0.360845in}{0.924000in}}%
\pgfusepath{clip}%
\pgfsetbuttcap%
\pgfsetmiterjoin%
\definecolor{currentfill}{rgb}{0.000000,0.596078,0.701961}%
\pgfsetfillcolor{currentfill}%
\pgfsetlinewidth{0.000000pt}%
\definecolor{currentstroke}{rgb}{0.000000,0.000000,0.000000}%
\pgfsetstrokecolor{currentstroke}%
\pgfsetstrokeopacity{0.000000}%
\pgfsetdash{}{0pt}%
\pgfpathmoveto{\pgfqpoint{3.560811in}{0.655209in}}%
\pgfpathlineto{\pgfqpoint{3.680411in}{0.655209in}}%
\pgfpathlineto{\pgfqpoint{3.680411in}{0.747609in}}%
\pgfpathlineto{\pgfqpoint{3.560811in}{0.747609in}}%
\pgfpathlineto{\pgfqpoint{3.560811in}{0.655209in}}%
\pgfpathclose%
\pgfusepath{fill}%
\end{pgfscope}%
\begin{pgfscope}%
\pgfpathrectangle{\pgfqpoint{3.319565in}{0.488889in}}{\pgfqpoint{0.360845in}{0.924000in}}%
\pgfusepath{clip}%
\pgfsetbuttcap%
\pgfsetmiterjoin%
\definecolor{currentfill}{rgb}{0.000000,0.584314,0.707843}%
\pgfsetfillcolor{currentfill}%
\pgfsetlinewidth{0.000000pt}%
\definecolor{currentstroke}{rgb}{0.000000,0.000000,0.000000}%
\pgfsetstrokecolor{currentstroke}%
\pgfsetstrokeopacity{0.000000}%
\pgfsetdash{}{0pt}%
\pgfpathmoveto{\pgfqpoint{3.559810in}{0.664449in}}%
\pgfpathlineto{\pgfqpoint{3.680411in}{0.664449in}}%
\pgfpathlineto{\pgfqpoint{3.680411in}{0.756849in}}%
\pgfpathlineto{\pgfqpoint{3.559810in}{0.756849in}}%
\pgfpathlineto{\pgfqpoint{3.559810in}{0.664449in}}%
\pgfpathclose%
\pgfusepath{fill}%
\end{pgfscope}%
\begin{pgfscope}%
\pgfpathrectangle{\pgfqpoint{3.319565in}{0.488889in}}{\pgfqpoint{0.360845in}{0.924000in}}%
\pgfusepath{clip}%
\pgfsetbuttcap%
\pgfsetmiterjoin%
\definecolor{currentfill}{rgb}{0.000000,0.576471,0.711765}%
\pgfsetfillcolor{currentfill}%
\pgfsetlinewidth{0.000000pt}%
\definecolor{currentstroke}{rgb}{0.000000,0.000000,0.000000}%
\pgfsetstrokecolor{currentstroke}%
\pgfsetstrokeopacity{0.000000}%
\pgfsetdash{}{0pt}%
\pgfpathmoveto{\pgfqpoint{3.559144in}{0.673689in}}%
\pgfpathlineto{\pgfqpoint{3.680411in}{0.673689in}}%
\pgfpathlineto{\pgfqpoint{3.680411in}{0.766089in}}%
\pgfpathlineto{\pgfqpoint{3.559144in}{0.766089in}}%
\pgfpathlineto{\pgfqpoint{3.559144in}{0.673689in}}%
\pgfpathclose%
\pgfusepath{fill}%
\end{pgfscope}%
\begin{pgfscope}%
\pgfpathrectangle{\pgfqpoint{3.319565in}{0.488889in}}{\pgfqpoint{0.360845in}{0.924000in}}%
\pgfusepath{clip}%
\pgfsetbuttcap%
\pgfsetmiterjoin%
\definecolor{currentfill}{rgb}{0.000000,0.568627,0.715686}%
\pgfsetfillcolor{currentfill}%
\pgfsetlinewidth{0.000000pt}%
\definecolor{currentstroke}{rgb}{0.000000,0.000000,0.000000}%
\pgfsetstrokecolor{currentstroke}%
\pgfsetstrokeopacity{0.000000}%
\pgfsetdash{}{0pt}%
\pgfpathmoveto{\pgfqpoint{3.558529in}{0.682929in}}%
\pgfpathlineto{\pgfqpoint{3.680411in}{0.682929in}}%
\pgfpathlineto{\pgfqpoint{3.680411in}{0.775329in}}%
\pgfpathlineto{\pgfqpoint{3.558529in}{0.775329in}}%
\pgfpathlineto{\pgfqpoint{3.558529in}{0.682929in}}%
\pgfpathclose%
\pgfusepath{fill}%
\end{pgfscope}%
\begin{pgfscope}%
\pgfpathrectangle{\pgfqpoint{3.319565in}{0.488889in}}{\pgfqpoint{0.360845in}{0.924000in}}%
\pgfusepath{clip}%
\pgfsetbuttcap%
\pgfsetmiterjoin%
\definecolor{currentfill}{rgb}{0.000000,0.564706,0.717647}%
\pgfsetfillcolor{currentfill}%
\pgfsetlinewidth{0.000000pt}%
\definecolor{currentstroke}{rgb}{0.000000,0.000000,0.000000}%
\pgfsetstrokecolor{currentstroke}%
\pgfsetstrokeopacity{0.000000}%
\pgfsetdash{}{0pt}%
\pgfpathmoveto{\pgfqpoint{3.558230in}{0.692169in}}%
\pgfpathlineto{\pgfqpoint{3.680411in}{0.692169in}}%
\pgfpathlineto{\pgfqpoint{3.680411in}{0.784569in}}%
\pgfpathlineto{\pgfqpoint{3.558230in}{0.784569in}}%
\pgfpathlineto{\pgfqpoint{3.558230in}{0.692169in}}%
\pgfpathclose%
\pgfusepath{fill}%
\end{pgfscope}%
\begin{pgfscope}%
\pgfpathrectangle{\pgfqpoint{3.319565in}{0.488889in}}{\pgfqpoint{0.360845in}{0.924000in}}%
\pgfusepath{clip}%
\pgfsetbuttcap%
\pgfsetmiterjoin%
\definecolor{currentfill}{rgb}{0.000000,0.560784,0.719608}%
\pgfsetfillcolor{currentfill}%
\pgfsetlinewidth{0.000000pt}%
\definecolor{currentstroke}{rgb}{0.000000,0.000000,0.000000}%
\pgfsetstrokecolor{currentstroke}%
\pgfsetstrokeopacity{0.000000}%
\pgfsetdash{}{0pt}%
\pgfpathmoveto{\pgfqpoint{3.557690in}{0.701409in}}%
\pgfpathlineto{\pgfqpoint{3.680411in}{0.701409in}}%
\pgfpathlineto{\pgfqpoint{3.680411in}{0.793809in}}%
\pgfpathlineto{\pgfqpoint{3.557690in}{0.793809in}}%
\pgfpathlineto{\pgfqpoint{3.557690in}{0.701409in}}%
\pgfpathclose%
\pgfusepath{fill}%
\end{pgfscope}%
\begin{pgfscope}%
\pgfpathrectangle{\pgfqpoint{3.319565in}{0.488889in}}{\pgfqpoint{0.360845in}{0.924000in}}%
\pgfusepath{clip}%
\pgfsetbuttcap%
\pgfsetmiterjoin%
\definecolor{currentfill}{rgb}{0.000000,0.552941,0.723529}%
\pgfsetfillcolor{currentfill}%
\pgfsetlinewidth{0.000000pt}%
\definecolor{currentstroke}{rgb}{0.000000,0.000000,0.000000}%
\pgfsetstrokecolor{currentstroke}%
\pgfsetstrokeopacity{0.000000}%
\pgfsetdash{}{0pt}%
\pgfpathmoveto{\pgfqpoint{3.557112in}{0.710649in}}%
\pgfpathlineto{\pgfqpoint{3.680411in}{0.710649in}}%
\pgfpathlineto{\pgfqpoint{3.680411in}{0.803049in}}%
\pgfpathlineto{\pgfqpoint{3.557112in}{0.803049in}}%
\pgfpathlineto{\pgfqpoint{3.557112in}{0.710649in}}%
\pgfpathclose%
\pgfusepath{fill}%
\end{pgfscope}%
\begin{pgfscope}%
\pgfpathrectangle{\pgfqpoint{3.319565in}{0.488889in}}{\pgfqpoint{0.360845in}{0.924000in}}%
\pgfusepath{clip}%
\pgfsetbuttcap%
\pgfsetmiterjoin%
\definecolor{currentfill}{rgb}{0.000000,0.549020,0.725490}%
\pgfsetfillcolor{currentfill}%
\pgfsetlinewidth{0.000000pt}%
\definecolor{currentstroke}{rgb}{0.000000,0.000000,0.000000}%
\pgfsetstrokecolor{currentstroke}%
\pgfsetstrokeopacity{0.000000}%
\pgfsetdash{}{0pt}%
\pgfpathmoveto{\pgfqpoint{3.556888in}{0.719889in}}%
\pgfpathlineto{\pgfqpoint{3.680411in}{0.719889in}}%
\pgfpathlineto{\pgfqpoint{3.680411in}{0.812289in}}%
\pgfpathlineto{\pgfqpoint{3.556888in}{0.812289in}}%
\pgfpathlineto{\pgfqpoint{3.556888in}{0.719889in}}%
\pgfpathclose%
\pgfusepath{fill}%
\end{pgfscope}%
\begin{pgfscope}%
\pgfpathrectangle{\pgfqpoint{3.319565in}{0.488889in}}{\pgfqpoint{0.360845in}{0.924000in}}%
\pgfusepath{clip}%
\pgfsetbuttcap%
\pgfsetmiterjoin%
\definecolor{currentfill}{rgb}{0.000000,0.545098,0.727451}%
\pgfsetfillcolor{currentfill}%
\pgfsetlinewidth{0.000000pt}%
\definecolor{currentstroke}{rgb}{0.000000,0.000000,0.000000}%
\pgfsetstrokecolor{currentstroke}%
\pgfsetstrokeopacity{0.000000}%
\pgfsetdash{}{0pt}%
\pgfpathmoveto{\pgfqpoint{3.556562in}{0.729129in}}%
\pgfpathlineto{\pgfqpoint{3.680411in}{0.729129in}}%
\pgfpathlineto{\pgfqpoint{3.680411in}{0.821529in}}%
\pgfpathlineto{\pgfqpoint{3.556562in}{0.821529in}}%
\pgfpathlineto{\pgfqpoint{3.556562in}{0.729129in}}%
\pgfpathclose%
\pgfusepath{fill}%
\end{pgfscope}%
\begin{pgfscope}%
\pgfpathrectangle{\pgfqpoint{3.319565in}{0.488889in}}{\pgfqpoint{0.360845in}{0.924000in}}%
\pgfusepath{clip}%
\pgfsetbuttcap%
\pgfsetmiterjoin%
\definecolor{currentfill}{rgb}{0.000000,0.537255,0.731373}%
\pgfsetfillcolor{currentfill}%
\pgfsetlinewidth{0.000000pt}%
\definecolor{currentstroke}{rgb}{0.000000,0.000000,0.000000}%
\pgfsetstrokecolor{currentstroke}%
\pgfsetstrokeopacity{0.000000}%
\pgfsetdash{}{0pt}%
\pgfpathmoveto{\pgfqpoint{3.556029in}{0.738369in}}%
\pgfpathlineto{\pgfqpoint{3.680411in}{0.738369in}}%
\pgfpathlineto{\pgfqpoint{3.680411in}{0.830769in}}%
\pgfpathlineto{\pgfqpoint{3.556029in}{0.830769in}}%
\pgfpathlineto{\pgfqpoint{3.556029in}{0.738369in}}%
\pgfpathclose%
\pgfusepath{fill}%
\end{pgfscope}%
\begin{pgfscope}%
\pgfpathrectangle{\pgfqpoint{3.319565in}{0.488889in}}{\pgfqpoint{0.360845in}{0.924000in}}%
\pgfusepath{clip}%
\pgfsetbuttcap%
\pgfsetmiterjoin%
\definecolor{currentfill}{rgb}{0.000000,0.537255,0.731373}%
\pgfsetfillcolor{currentfill}%
\pgfsetlinewidth{0.000000pt}%
\definecolor{currentstroke}{rgb}{0.000000,0.000000,0.000000}%
\pgfsetstrokecolor{currentstroke}%
\pgfsetstrokeopacity{0.000000}%
\pgfsetdash{}{0pt}%
\pgfpathmoveto{\pgfqpoint{3.555811in}{0.747609in}}%
\pgfpathlineto{\pgfqpoint{3.680411in}{0.747609in}}%
\pgfpathlineto{\pgfqpoint{3.680411in}{0.840009in}}%
\pgfpathlineto{\pgfqpoint{3.555811in}{0.840009in}}%
\pgfpathlineto{\pgfqpoint{3.555811in}{0.747609in}}%
\pgfpathclose%
\pgfusepath{fill}%
\end{pgfscope}%
\begin{pgfscope}%
\pgfpathrectangle{\pgfqpoint{3.319565in}{0.488889in}}{\pgfqpoint{0.360845in}{0.924000in}}%
\pgfusepath{clip}%
\pgfsetbuttcap%
\pgfsetmiterjoin%
\definecolor{currentfill}{rgb}{0.000000,0.537255,0.731373}%
\pgfsetfillcolor{currentfill}%
\pgfsetlinewidth{0.000000pt}%
\definecolor{currentstroke}{rgb}{0.000000,0.000000,0.000000}%
\pgfsetstrokecolor{currentstroke}%
\pgfsetstrokeopacity{0.000000}%
\pgfsetdash{}{0pt}%
\pgfpathmoveto{\pgfqpoint{3.555811in}{0.756849in}}%
\pgfpathlineto{\pgfqpoint{3.680411in}{0.756849in}}%
\pgfpathlineto{\pgfqpoint{3.680411in}{0.849249in}}%
\pgfpathlineto{\pgfqpoint{3.555811in}{0.849249in}}%
\pgfpathlineto{\pgfqpoint{3.555811in}{0.756849in}}%
\pgfpathclose%
\pgfusepath{fill}%
\end{pgfscope}%
\begin{pgfscope}%
\pgfpathrectangle{\pgfqpoint{3.319565in}{0.488889in}}{\pgfqpoint{0.360845in}{0.924000in}}%
\pgfusepath{clip}%
\pgfsetbuttcap%
\pgfsetmiterjoin%
\definecolor{currentfill}{rgb}{0.000000,0.529412,0.735294}%
\pgfsetfillcolor{currentfill}%
\pgfsetlinewidth{0.000000pt}%
\definecolor{currentstroke}{rgb}{0.000000,0.000000,0.000000}%
\pgfsetstrokecolor{currentstroke}%
\pgfsetstrokeopacity{0.000000}%
\pgfsetdash{}{0pt}%
\pgfpathmoveto{\pgfqpoint{3.555410in}{0.766089in}}%
\pgfpathlineto{\pgfqpoint{3.680411in}{0.766089in}}%
\pgfpathlineto{\pgfqpoint{3.680411in}{0.858489in}}%
\pgfpathlineto{\pgfqpoint{3.555410in}{0.858489in}}%
\pgfpathlineto{\pgfqpoint{3.555410in}{0.766089in}}%
\pgfpathclose%
\pgfusepath{fill}%
\end{pgfscope}%
\begin{pgfscope}%
\pgfpathrectangle{\pgfqpoint{3.319565in}{0.488889in}}{\pgfqpoint{0.360845in}{0.924000in}}%
\pgfusepath{clip}%
\pgfsetbuttcap%
\pgfsetmiterjoin%
\definecolor{currentfill}{rgb}{0.000000,0.525490,0.737255}%
\pgfsetfillcolor{currentfill}%
\pgfsetlinewidth{0.000000pt}%
\definecolor{currentstroke}{rgb}{0.000000,0.000000,0.000000}%
\pgfsetstrokecolor{currentstroke}%
\pgfsetstrokeopacity{0.000000}%
\pgfsetdash{}{0pt}%
\pgfpathmoveto{\pgfqpoint{3.555056in}{0.775329in}}%
\pgfpathlineto{\pgfqpoint{3.680411in}{0.775329in}}%
\pgfpathlineto{\pgfqpoint{3.680411in}{0.867729in}}%
\pgfpathlineto{\pgfqpoint{3.555056in}{0.867729in}}%
\pgfpathlineto{\pgfqpoint{3.555056in}{0.775329in}}%
\pgfpathclose%
\pgfusepath{fill}%
\end{pgfscope}%
\begin{pgfscope}%
\pgfpathrectangle{\pgfqpoint{3.319565in}{0.488889in}}{\pgfqpoint{0.360845in}{0.924000in}}%
\pgfusepath{clip}%
\pgfsetbuttcap%
\pgfsetmiterjoin%
\definecolor{currentfill}{rgb}{0.000000,0.521569,0.739216}%
\pgfsetfillcolor{currentfill}%
\pgfsetlinewidth{0.000000pt}%
\definecolor{currentstroke}{rgb}{0.000000,0.000000,0.000000}%
\pgfsetstrokecolor{currentstroke}%
\pgfsetstrokeopacity{0.000000}%
\pgfsetdash{}{0pt}%
\pgfpathmoveto{\pgfqpoint{3.554599in}{0.784569in}}%
\pgfpathlineto{\pgfqpoint{3.680411in}{0.784569in}}%
\pgfpathlineto{\pgfqpoint{3.680411in}{0.876969in}}%
\pgfpathlineto{\pgfqpoint{3.554599in}{0.876969in}}%
\pgfpathlineto{\pgfqpoint{3.554599in}{0.784569in}}%
\pgfpathclose%
\pgfusepath{fill}%
\end{pgfscope}%
\begin{pgfscope}%
\pgfpathrectangle{\pgfqpoint{3.319565in}{0.488889in}}{\pgfqpoint{0.360845in}{0.924000in}}%
\pgfusepath{clip}%
\pgfsetbuttcap%
\pgfsetmiterjoin%
\definecolor{currentfill}{rgb}{0.000000,0.521569,0.739216}%
\pgfsetfillcolor{currentfill}%
\pgfsetlinewidth{0.000000pt}%
\definecolor{currentstroke}{rgb}{0.000000,0.000000,0.000000}%
\pgfsetstrokecolor{currentstroke}%
\pgfsetstrokeopacity{0.000000}%
\pgfsetdash{}{0pt}%
\pgfpathmoveto{\pgfqpoint{3.554552in}{0.793809in}}%
\pgfpathlineto{\pgfqpoint{3.680411in}{0.793809in}}%
\pgfpathlineto{\pgfqpoint{3.680411in}{0.886209in}}%
\pgfpathlineto{\pgfqpoint{3.554552in}{0.886209in}}%
\pgfpathlineto{\pgfqpoint{3.554552in}{0.793809in}}%
\pgfpathclose%
\pgfusepath{fill}%
\end{pgfscope}%
\begin{pgfscope}%
\pgfpathrectangle{\pgfqpoint{3.319565in}{0.488889in}}{\pgfqpoint{0.360845in}{0.924000in}}%
\pgfusepath{clip}%
\pgfsetbuttcap%
\pgfsetmiterjoin%
\definecolor{currentfill}{rgb}{0.000000,0.521569,0.739216}%
\pgfsetfillcolor{currentfill}%
\pgfsetlinewidth{0.000000pt}%
\definecolor{currentstroke}{rgb}{0.000000,0.000000,0.000000}%
\pgfsetstrokecolor{currentstroke}%
\pgfsetstrokeopacity{0.000000}%
\pgfsetdash{}{0pt}%
\pgfpathmoveto{\pgfqpoint{3.554552in}{0.803049in}}%
\pgfpathlineto{\pgfqpoint{3.680411in}{0.803049in}}%
\pgfpathlineto{\pgfqpoint{3.680411in}{0.895449in}}%
\pgfpathlineto{\pgfqpoint{3.554552in}{0.895449in}}%
\pgfpathlineto{\pgfqpoint{3.554552in}{0.803049in}}%
\pgfpathclose%
\pgfusepath{fill}%
\end{pgfscope}%
\begin{pgfscope}%
\pgfpathrectangle{\pgfqpoint{3.319565in}{0.488889in}}{\pgfqpoint{0.360845in}{0.924000in}}%
\pgfusepath{clip}%
\pgfsetbuttcap%
\pgfsetmiterjoin%
\definecolor{currentfill}{rgb}{0.000000,0.521569,0.739216}%
\pgfsetfillcolor{currentfill}%
\pgfsetlinewidth{0.000000pt}%
\definecolor{currentstroke}{rgb}{0.000000,0.000000,0.000000}%
\pgfsetstrokecolor{currentstroke}%
\pgfsetstrokeopacity{0.000000}%
\pgfsetdash{}{0pt}%
\pgfpathmoveto{\pgfqpoint{3.554552in}{0.812289in}}%
\pgfpathlineto{\pgfqpoint{3.680411in}{0.812289in}}%
\pgfpathlineto{\pgfqpoint{3.680411in}{0.904689in}}%
\pgfpathlineto{\pgfqpoint{3.554552in}{0.904689in}}%
\pgfpathlineto{\pgfqpoint{3.554552in}{0.812289in}}%
\pgfpathclose%
\pgfusepath{fill}%
\end{pgfscope}%
\begin{pgfscope}%
\pgfpathrectangle{\pgfqpoint{3.319565in}{0.488889in}}{\pgfqpoint{0.360845in}{0.924000in}}%
\pgfusepath{clip}%
\pgfsetbuttcap%
\pgfsetmiterjoin%
\definecolor{currentfill}{rgb}{0.000000,0.521569,0.739216}%
\pgfsetfillcolor{currentfill}%
\pgfsetlinewidth{0.000000pt}%
\definecolor{currentstroke}{rgb}{0.000000,0.000000,0.000000}%
\pgfsetstrokecolor{currentstroke}%
\pgfsetstrokeopacity{0.000000}%
\pgfsetdash{}{0pt}%
\pgfpathmoveto{\pgfqpoint{3.554552in}{0.821529in}}%
\pgfpathlineto{\pgfqpoint{3.680411in}{0.821529in}}%
\pgfpathlineto{\pgfqpoint{3.680411in}{0.913929in}}%
\pgfpathlineto{\pgfqpoint{3.554552in}{0.913929in}}%
\pgfpathlineto{\pgfqpoint{3.554552in}{0.821529in}}%
\pgfpathclose%
\pgfusepath{fill}%
\end{pgfscope}%
\begin{pgfscope}%
\pgfpathrectangle{\pgfqpoint{3.319565in}{0.488889in}}{\pgfqpoint{0.360845in}{0.924000in}}%
\pgfusepath{clip}%
\pgfsetbuttcap%
\pgfsetmiterjoin%
\definecolor{currentfill}{rgb}{0.000000,0.521569,0.739216}%
\pgfsetfillcolor{currentfill}%
\pgfsetlinewidth{0.000000pt}%
\definecolor{currentstroke}{rgb}{0.000000,0.000000,0.000000}%
\pgfsetstrokecolor{currentstroke}%
\pgfsetstrokeopacity{0.000000}%
\pgfsetdash{}{0pt}%
\pgfpathmoveto{\pgfqpoint{3.554552in}{0.830769in}}%
\pgfpathlineto{\pgfqpoint{3.680411in}{0.830769in}}%
\pgfpathlineto{\pgfqpoint{3.680411in}{0.923169in}}%
\pgfpathlineto{\pgfqpoint{3.554552in}{0.923169in}}%
\pgfpathlineto{\pgfqpoint{3.554552in}{0.830769in}}%
\pgfpathclose%
\pgfusepath{fill}%
\end{pgfscope}%
\begin{pgfscope}%
\pgfpathrectangle{\pgfqpoint{3.319565in}{0.488889in}}{\pgfqpoint{0.360845in}{0.924000in}}%
\pgfusepath{clip}%
\pgfsetbuttcap%
\pgfsetmiterjoin%
\definecolor{currentfill}{rgb}{0.000000,0.517647,0.741176}%
\pgfsetfillcolor{currentfill}%
\pgfsetlinewidth{0.000000pt}%
\definecolor{currentstroke}{rgb}{0.000000,0.000000,0.000000}%
\pgfsetstrokecolor{currentstroke}%
\pgfsetstrokeopacity{0.000000}%
\pgfsetdash{}{0pt}%
\pgfpathmoveto{\pgfqpoint{3.554430in}{0.840009in}}%
\pgfpathlineto{\pgfqpoint{3.680411in}{0.840009in}}%
\pgfpathlineto{\pgfqpoint{3.680411in}{0.932409in}}%
\pgfpathlineto{\pgfqpoint{3.554430in}{0.932409in}}%
\pgfpathlineto{\pgfqpoint{3.554430in}{0.840009in}}%
\pgfpathclose%
\pgfusepath{fill}%
\end{pgfscope}%
\begin{pgfscope}%
\pgfpathrectangle{\pgfqpoint{3.319565in}{0.488889in}}{\pgfqpoint{0.360845in}{0.924000in}}%
\pgfusepath{clip}%
\pgfsetbuttcap%
\pgfsetmiterjoin%
\definecolor{currentfill}{rgb}{0.000000,0.505882,0.747059}%
\pgfsetfillcolor{currentfill}%
\pgfsetlinewidth{0.000000pt}%
\definecolor{currentstroke}{rgb}{0.000000,0.000000,0.000000}%
\pgfsetstrokecolor{currentstroke}%
\pgfsetstrokeopacity{0.000000}%
\pgfsetdash{}{0pt}%
\pgfpathmoveto{\pgfqpoint{3.553331in}{0.849249in}}%
\pgfpathlineto{\pgfqpoint{3.680411in}{0.849249in}}%
\pgfpathlineto{\pgfqpoint{3.680411in}{0.941649in}}%
\pgfpathlineto{\pgfqpoint{3.553331in}{0.941649in}}%
\pgfpathlineto{\pgfqpoint{3.553331in}{0.849249in}}%
\pgfpathclose%
\pgfusepath{fill}%
\end{pgfscope}%
\begin{pgfscope}%
\pgfpathrectangle{\pgfqpoint{3.319565in}{0.488889in}}{\pgfqpoint{0.360845in}{0.924000in}}%
\pgfusepath{clip}%
\pgfsetbuttcap%
\pgfsetmiterjoin%
\definecolor{currentfill}{rgb}{0.000000,0.494118,0.752941}%
\pgfsetfillcolor{currentfill}%
\pgfsetlinewidth{0.000000pt}%
\definecolor{currentstroke}{rgb}{0.000000,0.000000,0.000000}%
\pgfsetstrokecolor{currentstroke}%
\pgfsetstrokeopacity{0.000000}%
\pgfsetdash{}{0pt}%
\pgfpathmoveto{\pgfqpoint{3.552507in}{0.858489in}}%
\pgfpathlineto{\pgfqpoint{3.680411in}{0.858489in}}%
\pgfpathlineto{\pgfqpoint{3.680411in}{0.950889in}}%
\pgfpathlineto{\pgfqpoint{3.552507in}{0.950889in}}%
\pgfpathlineto{\pgfqpoint{3.552507in}{0.858489in}}%
\pgfpathclose%
\pgfusepath{fill}%
\end{pgfscope}%
\begin{pgfscope}%
\pgfpathrectangle{\pgfqpoint{3.319565in}{0.488889in}}{\pgfqpoint{0.360845in}{0.924000in}}%
\pgfusepath{clip}%
\pgfsetbuttcap%
\pgfsetmiterjoin%
\definecolor{currentfill}{rgb}{0.000000,0.490196,0.754902}%
\pgfsetfillcolor{currentfill}%
\pgfsetlinewidth{0.000000pt}%
\definecolor{currentstroke}{rgb}{0.000000,0.000000,0.000000}%
\pgfsetstrokecolor{currentstroke}%
\pgfsetstrokeopacity{0.000000}%
\pgfsetdash{}{0pt}%
\pgfpathmoveto{\pgfqpoint{3.552233in}{0.867729in}}%
\pgfpathlineto{\pgfqpoint{3.680411in}{0.867729in}}%
\pgfpathlineto{\pgfqpoint{3.680411in}{0.960129in}}%
\pgfpathlineto{\pgfqpoint{3.552233in}{0.960129in}}%
\pgfpathlineto{\pgfqpoint{3.552233in}{0.867729in}}%
\pgfpathclose%
\pgfusepath{fill}%
\end{pgfscope}%
\begin{pgfscope}%
\pgfpathrectangle{\pgfqpoint{3.319565in}{0.488889in}}{\pgfqpoint{0.360845in}{0.924000in}}%
\pgfusepath{clip}%
\pgfsetbuttcap%
\pgfsetmiterjoin%
\definecolor{currentfill}{rgb}{0.000000,0.482353,0.758824}%
\pgfsetfillcolor{currentfill}%
\pgfsetlinewidth{0.000000pt}%
\definecolor{currentstroke}{rgb}{0.000000,0.000000,0.000000}%
\pgfsetstrokecolor{currentstroke}%
\pgfsetstrokeopacity{0.000000}%
\pgfsetdash{}{0pt}%
\pgfpathmoveto{\pgfqpoint{3.551600in}{0.876969in}}%
\pgfpathlineto{\pgfqpoint{3.680411in}{0.876969in}}%
\pgfpathlineto{\pgfqpoint{3.680411in}{0.969369in}}%
\pgfpathlineto{\pgfqpoint{3.551600in}{0.969369in}}%
\pgfpathlineto{\pgfqpoint{3.551600in}{0.876969in}}%
\pgfpathclose%
\pgfusepath{fill}%
\end{pgfscope}%
\begin{pgfscope}%
\pgfpathrectangle{\pgfqpoint{3.319565in}{0.488889in}}{\pgfqpoint{0.360845in}{0.924000in}}%
\pgfusepath{clip}%
\pgfsetbuttcap%
\pgfsetmiterjoin%
\definecolor{currentfill}{rgb}{0.000000,0.482353,0.758824}%
\pgfsetfillcolor{currentfill}%
\pgfsetlinewidth{0.000000pt}%
\definecolor{currentstroke}{rgb}{0.000000,0.000000,0.000000}%
\pgfsetstrokecolor{currentstroke}%
\pgfsetstrokeopacity{0.000000}%
\pgfsetdash{}{0pt}%
\pgfpathmoveto{\pgfqpoint{3.551392in}{0.886209in}}%
\pgfpathlineto{\pgfqpoint{3.680411in}{0.886209in}}%
\pgfpathlineto{\pgfqpoint{3.680411in}{0.978609in}}%
\pgfpathlineto{\pgfqpoint{3.551392in}{0.978609in}}%
\pgfpathlineto{\pgfqpoint{3.551392in}{0.886209in}}%
\pgfpathclose%
\pgfusepath{fill}%
\end{pgfscope}%
\begin{pgfscope}%
\pgfpathrectangle{\pgfqpoint{3.319565in}{0.488889in}}{\pgfqpoint{0.360845in}{0.924000in}}%
\pgfusepath{clip}%
\pgfsetbuttcap%
\pgfsetmiterjoin%
\definecolor{currentfill}{rgb}{0.000000,0.482353,0.758824}%
\pgfsetfillcolor{currentfill}%
\pgfsetlinewidth{0.000000pt}%
\definecolor{currentstroke}{rgb}{0.000000,0.000000,0.000000}%
\pgfsetstrokecolor{currentstroke}%
\pgfsetstrokeopacity{0.000000}%
\pgfsetdash{}{0pt}%
\pgfpathmoveto{\pgfqpoint{3.551392in}{0.895449in}}%
\pgfpathlineto{\pgfqpoint{3.680411in}{0.895449in}}%
\pgfpathlineto{\pgfqpoint{3.680411in}{0.987849in}}%
\pgfpathlineto{\pgfqpoint{3.551392in}{0.987849in}}%
\pgfpathlineto{\pgfqpoint{3.551392in}{0.895449in}}%
\pgfpathclose%
\pgfusepath{fill}%
\end{pgfscope}%
\begin{pgfscope}%
\pgfpathrectangle{\pgfqpoint{3.319565in}{0.488889in}}{\pgfqpoint{0.360845in}{0.924000in}}%
\pgfusepath{clip}%
\pgfsetbuttcap%
\pgfsetmiterjoin%
\definecolor{currentfill}{rgb}{0.000000,0.482353,0.758824}%
\pgfsetfillcolor{currentfill}%
\pgfsetlinewidth{0.000000pt}%
\definecolor{currentstroke}{rgb}{0.000000,0.000000,0.000000}%
\pgfsetstrokecolor{currentstroke}%
\pgfsetstrokeopacity{0.000000}%
\pgfsetdash{}{0pt}%
\pgfpathmoveto{\pgfqpoint{3.551392in}{0.904689in}}%
\pgfpathlineto{\pgfqpoint{3.680411in}{0.904689in}}%
\pgfpathlineto{\pgfqpoint{3.680411in}{0.997089in}}%
\pgfpathlineto{\pgfqpoint{3.551392in}{0.997089in}}%
\pgfpathlineto{\pgfqpoint{3.551392in}{0.904689in}}%
\pgfpathclose%
\pgfusepath{fill}%
\end{pgfscope}%
\begin{pgfscope}%
\pgfpathrectangle{\pgfqpoint{3.319565in}{0.488889in}}{\pgfqpoint{0.360845in}{0.924000in}}%
\pgfusepath{clip}%
\pgfsetbuttcap%
\pgfsetmiterjoin%
\definecolor{currentfill}{rgb}{0.000000,0.482353,0.758824}%
\pgfsetfillcolor{currentfill}%
\pgfsetlinewidth{0.000000pt}%
\definecolor{currentstroke}{rgb}{0.000000,0.000000,0.000000}%
\pgfsetstrokecolor{currentstroke}%
\pgfsetstrokeopacity{0.000000}%
\pgfsetdash{}{0pt}%
\pgfpathmoveto{\pgfqpoint{3.551392in}{0.913929in}}%
\pgfpathlineto{\pgfqpoint{3.680411in}{0.913929in}}%
\pgfpathlineto{\pgfqpoint{3.680411in}{1.006329in}}%
\pgfpathlineto{\pgfqpoint{3.551392in}{1.006329in}}%
\pgfpathlineto{\pgfqpoint{3.551392in}{0.913929in}}%
\pgfpathclose%
\pgfusepath{fill}%
\end{pgfscope}%
\begin{pgfscope}%
\pgfpathrectangle{\pgfqpoint{3.319565in}{0.488889in}}{\pgfqpoint{0.360845in}{0.924000in}}%
\pgfusepath{clip}%
\pgfsetbuttcap%
\pgfsetmiterjoin%
\definecolor{currentfill}{rgb}{0.000000,0.478431,0.760784}%
\pgfsetfillcolor{currentfill}%
\pgfsetlinewidth{0.000000pt}%
\definecolor{currentstroke}{rgb}{0.000000,0.000000,0.000000}%
\pgfsetstrokecolor{currentstroke}%
\pgfsetstrokeopacity{0.000000}%
\pgfsetdash{}{0pt}%
\pgfpathmoveto{\pgfqpoint{3.551102in}{0.923169in}}%
\pgfpathlineto{\pgfqpoint{3.680411in}{0.923169in}}%
\pgfpathlineto{\pgfqpoint{3.680411in}{1.015569in}}%
\pgfpathlineto{\pgfqpoint{3.551102in}{1.015569in}}%
\pgfpathlineto{\pgfqpoint{3.551102in}{0.923169in}}%
\pgfpathclose%
\pgfusepath{fill}%
\end{pgfscope}%
\begin{pgfscope}%
\pgfpathrectangle{\pgfqpoint{3.319565in}{0.488889in}}{\pgfqpoint{0.360845in}{0.924000in}}%
\pgfusepath{clip}%
\pgfsetbuttcap%
\pgfsetmiterjoin%
\definecolor{currentfill}{rgb}{0.000000,0.474510,0.762745}%
\pgfsetfillcolor{currentfill}%
\pgfsetlinewidth{0.000000pt}%
\definecolor{currentstroke}{rgb}{0.000000,0.000000,0.000000}%
\pgfsetstrokecolor{currentstroke}%
\pgfsetstrokeopacity{0.000000}%
\pgfsetdash{}{0pt}%
\pgfpathmoveto{\pgfqpoint{3.550994in}{0.932409in}}%
\pgfpathlineto{\pgfqpoint{3.680411in}{0.932409in}}%
\pgfpathlineto{\pgfqpoint{3.680411in}{1.024809in}}%
\pgfpathlineto{\pgfqpoint{3.550994in}{1.024809in}}%
\pgfpathlineto{\pgfqpoint{3.550994in}{0.932409in}}%
\pgfpathclose%
\pgfusepath{fill}%
\end{pgfscope}%
\begin{pgfscope}%
\pgfpathrectangle{\pgfqpoint{3.319565in}{0.488889in}}{\pgfqpoint{0.360845in}{0.924000in}}%
\pgfusepath{clip}%
\pgfsetbuttcap%
\pgfsetmiterjoin%
\definecolor{currentfill}{rgb}{0.000000,0.466667,0.766667}%
\pgfsetfillcolor{currentfill}%
\pgfsetlinewidth{0.000000pt}%
\definecolor{currentstroke}{rgb}{0.000000,0.000000,0.000000}%
\pgfsetstrokecolor{currentstroke}%
\pgfsetstrokeopacity{0.000000}%
\pgfsetdash{}{0pt}%
\pgfpathmoveto{\pgfqpoint{3.550302in}{0.941649in}}%
\pgfpathlineto{\pgfqpoint{3.680411in}{0.941649in}}%
\pgfpathlineto{\pgfqpoint{3.680411in}{1.034049in}}%
\pgfpathlineto{\pgfqpoint{3.550302in}{1.034049in}}%
\pgfpathlineto{\pgfqpoint{3.550302in}{0.941649in}}%
\pgfpathclose%
\pgfusepath{fill}%
\end{pgfscope}%
\begin{pgfscope}%
\pgfpathrectangle{\pgfqpoint{3.319565in}{0.488889in}}{\pgfqpoint{0.360845in}{0.924000in}}%
\pgfusepath{clip}%
\pgfsetbuttcap%
\pgfsetmiterjoin%
\definecolor{currentfill}{rgb}{0.000000,0.466667,0.766667}%
\pgfsetfillcolor{currentfill}%
\pgfsetlinewidth{0.000000pt}%
\definecolor{currentstroke}{rgb}{0.000000,0.000000,0.000000}%
\pgfsetstrokecolor{currentstroke}%
\pgfsetstrokeopacity{0.000000}%
\pgfsetdash{}{0pt}%
\pgfpathmoveto{\pgfqpoint{3.550168in}{0.950889in}}%
\pgfpathlineto{\pgfqpoint{3.680411in}{0.950889in}}%
\pgfpathlineto{\pgfqpoint{3.680411in}{1.043289in}}%
\pgfpathlineto{\pgfqpoint{3.550168in}{1.043289in}}%
\pgfpathlineto{\pgfqpoint{3.550168in}{0.950889in}}%
\pgfpathclose%
\pgfusepath{fill}%
\end{pgfscope}%
\begin{pgfscope}%
\pgfpathrectangle{\pgfqpoint{3.319565in}{0.488889in}}{\pgfqpoint{0.360845in}{0.924000in}}%
\pgfusepath{clip}%
\pgfsetbuttcap%
\pgfsetmiterjoin%
\definecolor{currentfill}{rgb}{0.000000,0.462745,0.768627}%
\pgfsetfillcolor{currentfill}%
\pgfsetlinewidth{0.000000pt}%
\definecolor{currentstroke}{rgb}{0.000000,0.000000,0.000000}%
\pgfsetstrokecolor{currentstroke}%
\pgfsetstrokeopacity{0.000000}%
\pgfsetdash{}{0pt}%
\pgfpathmoveto{\pgfqpoint{3.549899in}{0.960129in}}%
\pgfpathlineto{\pgfqpoint{3.680411in}{0.960129in}}%
\pgfpathlineto{\pgfqpoint{3.680411in}{1.052529in}}%
\pgfpathlineto{\pgfqpoint{3.549899in}{1.052529in}}%
\pgfpathlineto{\pgfqpoint{3.549899in}{0.960129in}}%
\pgfpathclose%
\pgfusepath{fill}%
\end{pgfscope}%
\begin{pgfscope}%
\pgfpathrectangle{\pgfqpoint{3.319565in}{0.488889in}}{\pgfqpoint{0.360845in}{0.924000in}}%
\pgfusepath{clip}%
\pgfsetbuttcap%
\pgfsetmiterjoin%
\definecolor{currentfill}{rgb}{0.000000,0.454902,0.772549}%
\pgfsetfillcolor{currentfill}%
\pgfsetlinewidth{0.000000pt}%
\definecolor{currentstroke}{rgb}{0.000000,0.000000,0.000000}%
\pgfsetstrokecolor{currentstroke}%
\pgfsetstrokeopacity{0.000000}%
\pgfsetdash{}{0pt}%
\pgfpathmoveto{\pgfqpoint{3.549135in}{0.969369in}}%
\pgfpathlineto{\pgfqpoint{3.680411in}{0.969369in}}%
\pgfpathlineto{\pgfqpoint{3.680411in}{1.061769in}}%
\pgfpathlineto{\pgfqpoint{3.549135in}{1.061769in}}%
\pgfpathlineto{\pgfqpoint{3.549135in}{0.969369in}}%
\pgfpathclose%
\pgfusepath{fill}%
\end{pgfscope}%
\begin{pgfscope}%
\pgfpathrectangle{\pgfqpoint{3.319565in}{0.488889in}}{\pgfqpoint{0.360845in}{0.924000in}}%
\pgfusepath{clip}%
\pgfsetbuttcap%
\pgfsetmiterjoin%
\definecolor{currentfill}{rgb}{0.000000,0.447059,0.776471}%
\pgfsetfillcolor{currentfill}%
\pgfsetlinewidth{0.000000pt}%
\definecolor{currentstroke}{rgb}{0.000000,0.000000,0.000000}%
\pgfsetstrokecolor{currentstroke}%
\pgfsetstrokeopacity{0.000000}%
\pgfsetdash{}{0pt}%
\pgfpathmoveto{\pgfqpoint{3.548594in}{0.978609in}}%
\pgfpathlineto{\pgfqpoint{3.680411in}{0.978609in}}%
\pgfpathlineto{\pgfqpoint{3.680411in}{1.071009in}}%
\pgfpathlineto{\pgfqpoint{3.548594in}{1.071009in}}%
\pgfpathlineto{\pgfqpoint{3.548594in}{0.978609in}}%
\pgfpathclose%
\pgfusepath{fill}%
\end{pgfscope}%
\begin{pgfscope}%
\pgfpathrectangle{\pgfqpoint{3.319565in}{0.488889in}}{\pgfqpoint{0.360845in}{0.924000in}}%
\pgfusepath{clip}%
\pgfsetbuttcap%
\pgfsetmiterjoin%
\definecolor{currentfill}{rgb}{0.000000,0.447059,0.776471}%
\pgfsetfillcolor{currentfill}%
\pgfsetlinewidth{0.000000pt}%
\definecolor{currentstroke}{rgb}{0.000000,0.000000,0.000000}%
\pgfsetstrokecolor{currentstroke}%
\pgfsetstrokeopacity{0.000000}%
\pgfsetdash{}{0pt}%
\pgfpathmoveto{\pgfqpoint{3.548591in}{0.987849in}}%
\pgfpathlineto{\pgfqpoint{3.680411in}{0.987849in}}%
\pgfpathlineto{\pgfqpoint{3.680411in}{1.080249in}}%
\pgfpathlineto{\pgfqpoint{3.548591in}{1.080249in}}%
\pgfpathlineto{\pgfqpoint{3.548591in}{0.987849in}}%
\pgfpathclose%
\pgfusepath{fill}%
\end{pgfscope}%
\begin{pgfscope}%
\pgfpathrectangle{\pgfqpoint{3.319565in}{0.488889in}}{\pgfqpoint{0.360845in}{0.924000in}}%
\pgfusepath{clip}%
\pgfsetbuttcap%
\pgfsetmiterjoin%
\definecolor{currentfill}{rgb}{0.000000,0.443137,0.778431}%
\pgfsetfillcolor{currentfill}%
\pgfsetlinewidth{0.000000pt}%
\definecolor{currentstroke}{rgb}{0.000000,0.000000,0.000000}%
\pgfsetstrokecolor{currentstroke}%
\pgfsetstrokeopacity{0.000000}%
\pgfsetdash{}{0pt}%
\pgfpathmoveto{\pgfqpoint{3.548193in}{0.997089in}}%
\pgfpathlineto{\pgfqpoint{3.680411in}{0.997089in}}%
\pgfpathlineto{\pgfqpoint{3.680411in}{1.089489in}}%
\pgfpathlineto{\pgfqpoint{3.548193in}{1.089489in}}%
\pgfpathlineto{\pgfqpoint{3.548193in}{0.997089in}}%
\pgfpathclose%
\pgfusepath{fill}%
\end{pgfscope}%
\begin{pgfscope}%
\pgfpathrectangle{\pgfqpoint{3.319565in}{0.488889in}}{\pgfqpoint{0.360845in}{0.924000in}}%
\pgfusepath{clip}%
\pgfsetbuttcap%
\pgfsetmiterjoin%
\definecolor{currentfill}{rgb}{0.000000,0.439216,0.780392}%
\pgfsetfillcolor{currentfill}%
\pgfsetlinewidth{0.000000pt}%
\definecolor{currentstroke}{rgb}{0.000000,0.000000,0.000000}%
\pgfsetstrokecolor{currentstroke}%
\pgfsetstrokeopacity{0.000000}%
\pgfsetdash{}{0pt}%
\pgfpathmoveto{\pgfqpoint{3.548087in}{1.006329in}}%
\pgfpathlineto{\pgfqpoint{3.680411in}{1.006329in}}%
\pgfpathlineto{\pgfqpoint{3.680411in}{1.098729in}}%
\pgfpathlineto{\pgfqpoint{3.548087in}{1.098729in}}%
\pgfpathlineto{\pgfqpoint{3.548087in}{1.006329in}}%
\pgfpathclose%
\pgfusepath{fill}%
\end{pgfscope}%
\begin{pgfscope}%
\pgfpathrectangle{\pgfqpoint{3.319565in}{0.488889in}}{\pgfqpoint{0.360845in}{0.924000in}}%
\pgfusepath{clip}%
\pgfsetbuttcap%
\pgfsetmiterjoin%
\definecolor{currentfill}{rgb}{0.000000,0.431373,0.784314}%
\pgfsetfillcolor{currentfill}%
\pgfsetlinewidth{0.000000pt}%
\definecolor{currentstroke}{rgb}{0.000000,0.000000,0.000000}%
\pgfsetstrokecolor{currentstroke}%
\pgfsetstrokeopacity{0.000000}%
\pgfsetdash{}{0pt}%
\pgfpathmoveto{\pgfqpoint{3.547398in}{1.015569in}}%
\pgfpathlineto{\pgfqpoint{3.680411in}{1.015569in}}%
\pgfpathlineto{\pgfqpoint{3.680411in}{1.107969in}}%
\pgfpathlineto{\pgfqpoint{3.547398in}{1.107969in}}%
\pgfpathlineto{\pgfqpoint{3.547398in}{1.015569in}}%
\pgfpathclose%
\pgfusepath{fill}%
\end{pgfscope}%
\begin{pgfscope}%
\pgfpathrectangle{\pgfqpoint{3.319565in}{0.488889in}}{\pgfqpoint{0.360845in}{0.924000in}}%
\pgfusepath{clip}%
\pgfsetbuttcap%
\pgfsetmiterjoin%
\definecolor{currentfill}{rgb}{0.000000,0.431373,0.784314}%
\pgfsetfillcolor{currentfill}%
\pgfsetlinewidth{0.000000pt}%
\definecolor{currentstroke}{rgb}{0.000000,0.000000,0.000000}%
\pgfsetstrokecolor{currentstroke}%
\pgfsetstrokeopacity{0.000000}%
\pgfsetdash{}{0pt}%
\pgfpathmoveto{\pgfqpoint{3.547398in}{1.024809in}}%
\pgfpathlineto{\pgfqpoint{3.680411in}{1.024809in}}%
\pgfpathlineto{\pgfqpoint{3.680411in}{1.117209in}}%
\pgfpathlineto{\pgfqpoint{3.547398in}{1.117209in}}%
\pgfpathlineto{\pgfqpoint{3.547398in}{1.024809in}}%
\pgfpathclose%
\pgfusepath{fill}%
\end{pgfscope}%
\begin{pgfscope}%
\pgfpathrectangle{\pgfqpoint{3.319565in}{0.488889in}}{\pgfqpoint{0.360845in}{0.924000in}}%
\pgfusepath{clip}%
\pgfsetbuttcap%
\pgfsetmiterjoin%
\definecolor{currentfill}{rgb}{0.000000,0.431373,0.784314}%
\pgfsetfillcolor{currentfill}%
\pgfsetlinewidth{0.000000pt}%
\definecolor{currentstroke}{rgb}{0.000000,0.000000,0.000000}%
\pgfsetstrokecolor{currentstroke}%
\pgfsetstrokeopacity{0.000000}%
\pgfsetdash{}{0pt}%
\pgfpathmoveto{\pgfqpoint{3.547398in}{1.034049in}}%
\pgfpathlineto{\pgfqpoint{3.680411in}{1.034049in}}%
\pgfpathlineto{\pgfqpoint{3.680411in}{1.126449in}}%
\pgfpathlineto{\pgfqpoint{3.547398in}{1.126449in}}%
\pgfpathlineto{\pgfqpoint{3.547398in}{1.034049in}}%
\pgfpathclose%
\pgfusepath{fill}%
\end{pgfscope}%
\begin{pgfscope}%
\pgfpathrectangle{\pgfqpoint{3.319565in}{0.488889in}}{\pgfqpoint{0.360845in}{0.924000in}}%
\pgfusepath{clip}%
\pgfsetbuttcap%
\pgfsetmiterjoin%
\definecolor{currentfill}{rgb}{0.000000,0.427451,0.786275}%
\pgfsetfillcolor{currentfill}%
\pgfsetlinewidth{0.000000pt}%
\definecolor{currentstroke}{rgb}{0.000000,0.000000,0.000000}%
\pgfsetstrokecolor{currentstroke}%
\pgfsetstrokeopacity{0.000000}%
\pgfsetdash{}{0pt}%
\pgfpathmoveto{\pgfqpoint{3.547013in}{1.043289in}}%
\pgfpathlineto{\pgfqpoint{3.680411in}{1.043289in}}%
\pgfpathlineto{\pgfqpoint{3.680411in}{1.135689in}}%
\pgfpathlineto{\pgfqpoint{3.547013in}{1.135689in}}%
\pgfpathlineto{\pgfqpoint{3.547013in}{1.043289in}}%
\pgfpathclose%
\pgfusepath{fill}%
\end{pgfscope}%
\begin{pgfscope}%
\pgfpathrectangle{\pgfqpoint{3.319565in}{0.488889in}}{\pgfqpoint{0.360845in}{0.924000in}}%
\pgfusepath{clip}%
\pgfsetbuttcap%
\pgfsetmiterjoin%
\definecolor{currentfill}{rgb}{0.000000,0.427451,0.786275}%
\pgfsetfillcolor{currentfill}%
\pgfsetlinewidth{0.000000pt}%
\definecolor{currentstroke}{rgb}{0.000000,0.000000,0.000000}%
\pgfsetstrokecolor{currentstroke}%
\pgfsetstrokeopacity{0.000000}%
\pgfsetdash{}{0pt}%
\pgfpathmoveto{\pgfqpoint{3.546966in}{1.052529in}}%
\pgfpathlineto{\pgfqpoint{3.680411in}{1.052529in}}%
\pgfpathlineto{\pgfqpoint{3.680411in}{1.144929in}}%
\pgfpathlineto{\pgfqpoint{3.546966in}{1.144929in}}%
\pgfpathlineto{\pgfqpoint{3.546966in}{1.052529in}}%
\pgfpathclose%
\pgfusepath{fill}%
\end{pgfscope}%
\begin{pgfscope}%
\pgfpathrectangle{\pgfqpoint{3.319565in}{0.488889in}}{\pgfqpoint{0.360845in}{0.924000in}}%
\pgfusepath{clip}%
\pgfsetbuttcap%
\pgfsetmiterjoin%
\definecolor{currentfill}{rgb}{0.000000,0.423529,0.788235}%
\pgfsetfillcolor{currentfill}%
\pgfsetlinewidth{0.000000pt}%
\definecolor{currentstroke}{rgb}{0.000000,0.000000,0.000000}%
\pgfsetstrokecolor{currentstroke}%
\pgfsetstrokeopacity{0.000000}%
\pgfsetdash{}{0pt}%
\pgfpathmoveto{\pgfqpoint{3.546797in}{1.061769in}}%
\pgfpathlineto{\pgfqpoint{3.680411in}{1.061769in}}%
\pgfpathlineto{\pgfqpoint{3.680411in}{1.154169in}}%
\pgfpathlineto{\pgfqpoint{3.546797in}{1.154169in}}%
\pgfpathlineto{\pgfqpoint{3.546797in}{1.061769in}}%
\pgfpathclose%
\pgfusepath{fill}%
\end{pgfscope}%
\begin{pgfscope}%
\pgfpathrectangle{\pgfqpoint{3.319565in}{0.488889in}}{\pgfqpoint{0.360845in}{0.924000in}}%
\pgfusepath{clip}%
\pgfsetbuttcap%
\pgfsetmiterjoin%
\definecolor{currentfill}{rgb}{0.000000,0.423529,0.788235}%
\pgfsetfillcolor{currentfill}%
\pgfsetlinewidth{0.000000pt}%
\definecolor{currentstroke}{rgb}{0.000000,0.000000,0.000000}%
\pgfsetstrokecolor{currentstroke}%
\pgfsetstrokeopacity{0.000000}%
\pgfsetdash{}{0pt}%
\pgfpathmoveto{\pgfqpoint{3.546797in}{1.071009in}}%
\pgfpathlineto{\pgfqpoint{3.680411in}{1.071009in}}%
\pgfpathlineto{\pgfqpoint{3.680411in}{1.163409in}}%
\pgfpathlineto{\pgfqpoint{3.546797in}{1.163409in}}%
\pgfpathlineto{\pgfqpoint{3.546797in}{1.071009in}}%
\pgfpathclose%
\pgfusepath{fill}%
\end{pgfscope}%
\begin{pgfscope}%
\pgfpathrectangle{\pgfqpoint{3.319565in}{0.488889in}}{\pgfqpoint{0.360845in}{0.924000in}}%
\pgfusepath{clip}%
\pgfsetbuttcap%
\pgfsetmiterjoin%
\definecolor{currentfill}{rgb}{0.000000,0.419608,0.790196}%
\pgfsetfillcolor{currentfill}%
\pgfsetlinewidth{0.000000pt}%
\definecolor{currentstroke}{rgb}{0.000000,0.000000,0.000000}%
\pgfsetstrokecolor{currentstroke}%
\pgfsetstrokeopacity{0.000000}%
\pgfsetdash{}{0pt}%
\pgfpathmoveto{\pgfqpoint{3.546460in}{1.080249in}}%
\pgfpathlineto{\pgfqpoint{3.680411in}{1.080249in}}%
\pgfpathlineto{\pgfqpoint{3.680411in}{1.172649in}}%
\pgfpathlineto{\pgfqpoint{3.546460in}{1.172649in}}%
\pgfpathlineto{\pgfqpoint{3.546460in}{1.080249in}}%
\pgfpathclose%
\pgfusepath{fill}%
\end{pgfscope}%
\begin{pgfscope}%
\pgfpathrectangle{\pgfqpoint{3.319565in}{0.488889in}}{\pgfqpoint{0.360845in}{0.924000in}}%
\pgfusepath{clip}%
\pgfsetbuttcap%
\pgfsetmiterjoin%
\definecolor{currentfill}{rgb}{0.000000,0.419608,0.790196}%
\pgfsetfillcolor{currentfill}%
\pgfsetlinewidth{0.000000pt}%
\definecolor{currentstroke}{rgb}{0.000000,0.000000,0.000000}%
\pgfsetstrokecolor{currentstroke}%
\pgfsetstrokeopacity{0.000000}%
\pgfsetdash{}{0pt}%
\pgfpathmoveto{\pgfqpoint{3.546402in}{1.089489in}}%
\pgfpathlineto{\pgfqpoint{3.680411in}{1.089489in}}%
\pgfpathlineto{\pgfqpoint{3.680411in}{1.181889in}}%
\pgfpathlineto{\pgfqpoint{3.546402in}{1.181889in}}%
\pgfpathlineto{\pgfqpoint{3.546402in}{1.089489in}}%
\pgfpathclose%
\pgfusepath{fill}%
\end{pgfscope}%
\begin{pgfscope}%
\pgfpathrectangle{\pgfqpoint{3.319565in}{0.488889in}}{\pgfqpoint{0.360845in}{0.924000in}}%
\pgfusepath{clip}%
\pgfsetbuttcap%
\pgfsetmiterjoin%
\definecolor{currentfill}{rgb}{0.000000,0.411765,0.794118}%
\pgfsetfillcolor{currentfill}%
\pgfsetlinewidth{0.000000pt}%
\definecolor{currentstroke}{rgb}{0.000000,0.000000,0.000000}%
\pgfsetstrokecolor{currentstroke}%
\pgfsetstrokeopacity{0.000000}%
\pgfsetdash{}{0pt}%
\pgfpathmoveto{\pgfqpoint{3.545929in}{1.098729in}}%
\pgfpathlineto{\pgfqpoint{3.680411in}{1.098729in}}%
\pgfpathlineto{\pgfqpoint{3.680411in}{1.191129in}}%
\pgfpathlineto{\pgfqpoint{3.545929in}{1.191129in}}%
\pgfpathlineto{\pgfqpoint{3.545929in}{1.098729in}}%
\pgfpathclose%
\pgfusepath{fill}%
\end{pgfscope}%
\begin{pgfscope}%
\pgfpathrectangle{\pgfqpoint{3.319565in}{0.488889in}}{\pgfqpoint{0.360845in}{0.924000in}}%
\pgfusepath{clip}%
\pgfsetbuttcap%
\pgfsetmiterjoin%
\definecolor{currentfill}{rgb}{0.000000,0.411765,0.794118}%
\pgfsetfillcolor{currentfill}%
\pgfsetlinewidth{0.000000pt}%
\definecolor{currentstroke}{rgb}{0.000000,0.000000,0.000000}%
\pgfsetstrokecolor{currentstroke}%
\pgfsetstrokeopacity{0.000000}%
\pgfsetdash{}{0pt}%
\pgfpathmoveto{\pgfqpoint{3.545690in}{1.107969in}}%
\pgfpathlineto{\pgfqpoint{3.680411in}{1.107969in}}%
\pgfpathlineto{\pgfqpoint{3.680411in}{1.200369in}}%
\pgfpathlineto{\pgfqpoint{3.545690in}{1.200369in}}%
\pgfpathlineto{\pgfqpoint{3.545690in}{1.107969in}}%
\pgfpathclose%
\pgfusepath{fill}%
\end{pgfscope}%
\begin{pgfscope}%
\pgfpathrectangle{\pgfqpoint{3.319565in}{0.488889in}}{\pgfqpoint{0.360845in}{0.924000in}}%
\pgfusepath{clip}%
\pgfsetbuttcap%
\pgfsetmiterjoin%
\definecolor{currentfill}{rgb}{0.000000,0.411765,0.794118}%
\pgfsetfillcolor{currentfill}%
\pgfsetlinewidth{0.000000pt}%
\definecolor{currentstroke}{rgb}{0.000000,0.000000,0.000000}%
\pgfsetstrokecolor{currentstroke}%
\pgfsetstrokeopacity{0.000000}%
\pgfsetdash{}{0pt}%
\pgfpathmoveto{\pgfqpoint{3.545690in}{1.117209in}}%
\pgfpathlineto{\pgfqpoint{3.680411in}{1.117209in}}%
\pgfpathlineto{\pgfqpoint{3.680411in}{1.209609in}}%
\pgfpathlineto{\pgfqpoint{3.545690in}{1.209609in}}%
\pgfpathlineto{\pgfqpoint{3.545690in}{1.117209in}}%
\pgfpathclose%
\pgfusepath{fill}%
\end{pgfscope}%
\begin{pgfscope}%
\pgfpathrectangle{\pgfqpoint{3.319565in}{0.488889in}}{\pgfqpoint{0.360845in}{0.924000in}}%
\pgfusepath{clip}%
\pgfsetbuttcap%
\pgfsetmiterjoin%
\definecolor{currentfill}{rgb}{0.000000,0.403922,0.798039}%
\pgfsetfillcolor{currentfill}%
\pgfsetlinewidth{0.000000pt}%
\definecolor{currentstroke}{rgb}{0.000000,0.000000,0.000000}%
\pgfsetstrokecolor{currentstroke}%
\pgfsetstrokeopacity{0.000000}%
\pgfsetdash{}{0pt}%
\pgfpathmoveto{\pgfqpoint{3.545305in}{1.126449in}}%
\pgfpathlineto{\pgfqpoint{3.680411in}{1.126449in}}%
\pgfpathlineto{\pgfqpoint{3.680411in}{1.218849in}}%
\pgfpathlineto{\pgfqpoint{3.545305in}{1.218849in}}%
\pgfpathlineto{\pgfqpoint{3.545305in}{1.126449in}}%
\pgfpathclose%
\pgfusepath{fill}%
\end{pgfscope}%
\begin{pgfscope}%
\pgfpathrectangle{\pgfqpoint{3.319565in}{0.488889in}}{\pgfqpoint{0.360845in}{0.924000in}}%
\pgfusepath{clip}%
\pgfsetbuttcap%
\pgfsetmiterjoin%
\definecolor{currentfill}{rgb}{0.000000,0.392157,0.803922}%
\pgfsetfillcolor{currentfill}%
\pgfsetlinewidth{0.000000pt}%
\definecolor{currentstroke}{rgb}{0.000000,0.000000,0.000000}%
\pgfsetstrokecolor{currentstroke}%
\pgfsetstrokeopacity{0.000000}%
\pgfsetdash{}{0pt}%
\pgfpathmoveto{\pgfqpoint{3.544072in}{1.135689in}}%
\pgfpathlineto{\pgfqpoint{3.680411in}{1.135689in}}%
\pgfpathlineto{\pgfqpoint{3.680411in}{1.228089in}}%
\pgfpathlineto{\pgfqpoint{3.544072in}{1.228089in}}%
\pgfpathlineto{\pgfqpoint{3.544072in}{1.135689in}}%
\pgfpathclose%
\pgfusepath{fill}%
\end{pgfscope}%
\begin{pgfscope}%
\pgfpathrectangle{\pgfqpoint{3.319565in}{0.488889in}}{\pgfqpoint{0.360845in}{0.924000in}}%
\pgfusepath{clip}%
\pgfsetbuttcap%
\pgfsetmiterjoin%
\definecolor{currentfill}{rgb}{0.000000,0.392157,0.803922}%
\pgfsetfillcolor{currentfill}%
\pgfsetlinewidth{0.000000pt}%
\definecolor{currentstroke}{rgb}{0.000000,0.000000,0.000000}%
\pgfsetstrokecolor{currentstroke}%
\pgfsetstrokeopacity{0.000000}%
\pgfsetdash{}{0pt}%
\pgfpathmoveto{\pgfqpoint{3.544072in}{1.144929in}}%
\pgfpathlineto{\pgfqpoint{3.680411in}{1.144929in}}%
\pgfpathlineto{\pgfqpoint{3.680411in}{1.237329in}}%
\pgfpathlineto{\pgfqpoint{3.544072in}{1.237329in}}%
\pgfpathlineto{\pgfqpoint{3.544072in}{1.144929in}}%
\pgfpathclose%
\pgfusepath{fill}%
\end{pgfscope}%
\begin{pgfscope}%
\pgfpathrectangle{\pgfqpoint{3.319565in}{0.488889in}}{\pgfqpoint{0.360845in}{0.924000in}}%
\pgfusepath{clip}%
\pgfsetbuttcap%
\pgfsetmiterjoin%
\definecolor{currentfill}{rgb}{0.000000,0.388235,0.805882}%
\pgfsetfillcolor{currentfill}%
\pgfsetlinewidth{0.000000pt}%
\definecolor{currentstroke}{rgb}{0.000000,0.000000,0.000000}%
\pgfsetstrokecolor{currentstroke}%
\pgfsetstrokeopacity{0.000000}%
\pgfsetdash{}{0pt}%
\pgfpathmoveto{\pgfqpoint{3.543791in}{1.154169in}}%
\pgfpathlineto{\pgfqpoint{3.680411in}{1.154169in}}%
\pgfpathlineto{\pgfqpoint{3.680411in}{1.246569in}}%
\pgfpathlineto{\pgfqpoint{3.543791in}{1.246569in}}%
\pgfpathlineto{\pgfqpoint{3.543791in}{1.154169in}}%
\pgfpathclose%
\pgfusepath{fill}%
\end{pgfscope}%
\begin{pgfscope}%
\pgfpathrectangle{\pgfqpoint{3.319565in}{0.488889in}}{\pgfqpoint{0.360845in}{0.924000in}}%
\pgfusepath{clip}%
\pgfsetbuttcap%
\pgfsetmiterjoin%
\definecolor{currentfill}{rgb}{0.000000,0.388235,0.805882}%
\pgfsetfillcolor{currentfill}%
\pgfsetlinewidth{0.000000pt}%
\definecolor{currentstroke}{rgb}{0.000000,0.000000,0.000000}%
\pgfsetstrokecolor{currentstroke}%
\pgfsetstrokeopacity{0.000000}%
\pgfsetdash{}{0pt}%
\pgfpathmoveto{\pgfqpoint{3.543791in}{1.163409in}}%
\pgfpathlineto{\pgfqpoint{3.680411in}{1.163409in}}%
\pgfpathlineto{\pgfqpoint{3.680411in}{1.255809in}}%
\pgfpathlineto{\pgfqpoint{3.543791in}{1.255809in}}%
\pgfpathlineto{\pgfqpoint{3.543791in}{1.163409in}}%
\pgfpathclose%
\pgfusepath{fill}%
\end{pgfscope}%
\begin{pgfscope}%
\pgfpathrectangle{\pgfqpoint{3.319565in}{0.488889in}}{\pgfqpoint{0.360845in}{0.924000in}}%
\pgfusepath{clip}%
\pgfsetbuttcap%
\pgfsetmiterjoin%
\definecolor{currentfill}{rgb}{0.000000,0.388235,0.805882}%
\pgfsetfillcolor{currentfill}%
\pgfsetlinewidth{0.000000pt}%
\definecolor{currentstroke}{rgb}{0.000000,0.000000,0.000000}%
\pgfsetstrokecolor{currentstroke}%
\pgfsetstrokeopacity{0.000000}%
\pgfsetdash{}{0pt}%
\pgfpathmoveto{\pgfqpoint{3.543791in}{1.172649in}}%
\pgfpathlineto{\pgfqpoint{3.680411in}{1.172649in}}%
\pgfpathlineto{\pgfqpoint{3.680411in}{1.265049in}}%
\pgfpathlineto{\pgfqpoint{3.543791in}{1.265049in}}%
\pgfpathlineto{\pgfqpoint{3.543791in}{1.172649in}}%
\pgfpathclose%
\pgfusepath{fill}%
\end{pgfscope}%
\begin{pgfscope}%
\pgfpathrectangle{\pgfqpoint{3.319565in}{0.488889in}}{\pgfqpoint{0.360845in}{0.924000in}}%
\pgfusepath{clip}%
\pgfsetbuttcap%
\pgfsetmiterjoin%
\definecolor{currentfill}{rgb}{0.000000,0.380392,0.809804}%
\pgfsetfillcolor{currentfill}%
\pgfsetlinewidth{0.000000pt}%
\definecolor{currentstroke}{rgb}{0.000000,0.000000,0.000000}%
\pgfsetstrokecolor{currentstroke}%
\pgfsetstrokeopacity{0.000000}%
\pgfsetdash{}{0pt}%
\pgfpathmoveto{\pgfqpoint{3.543213in}{1.181889in}}%
\pgfpathlineto{\pgfqpoint{3.680411in}{1.181889in}}%
\pgfpathlineto{\pgfqpoint{3.680411in}{1.274289in}}%
\pgfpathlineto{\pgfqpoint{3.543213in}{1.274289in}}%
\pgfpathlineto{\pgfqpoint{3.543213in}{1.181889in}}%
\pgfpathclose%
\pgfusepath{fill}%
\end{pgfscope}%
\begin{pgfscope}%
\pgfpathrectangle{\pgfqpoint{3.319565in}{0.488889in}}{\pgfqpoint{0.360845in}{0.924000in}}%
\pgfusepath{clip}%
\pgfsetbuttcap%
\pgfsetmiterjoin%
\definecolor{currentfill}{rgb}{0.000000,0.380392,0.809804}%
\pgfsetfillcolor{currentfill}%
\pgfsetlinewidth{0.000000pt}%
\definecolor{currentstroke}{rgb}{0.000000,0.000000,0.000000}%
\pgfsetstrokecolor{currentstroke}%
\pgfsetstrokeopacity{0.000000}%
\pgfsetdash{}{0pt}%
\pgfpathmoveto{\pgfqpoint{3.543213in}{1.191129in}}%
\pgfpathlineto{\pgfqpoint{3.680411in}{1.191129in}}%
\pgfpathlineto{\pgfqpoint{3.680411in}{1.283529in}}%
\pgfpathlineto{\pgfqpoint{3.543213in}{1.283529in}}%
\pgfpathlineto{\pgfqpoint{3.543213in}{1.191129in}}%
\pgfpathclose%
\pgfusepath{fill}%
\end{pgfscope}%
\begin{pgfscope}%
\pgfpathrectangle{\pgfqpoint{3.319565in}{0.488889in}}{\pgfqpoint{0.360845in}{0.924000in}}%
\pgfusepath{clip}%
\pgfsetbuttcap%
\pgfsetmiterjoin%
\definecolor{currentfill}{rgb}{0.000000,0.380392,0.809804}%
\pgfsetfillcolor{currentfill}%
\pgfsetlinewidth{0.000000pt}%
\definecolor{currentstroke}{rgb}{0.000000,0.000000,0.000000}%
\pgfsetstrokecolor{currentstroke}%
\pgfsetstrokeopacity{0.000000}%
\pgfsetdash{}{0pt}%
\pgfpathmoveto{\pgfqpoint{3.543104in}{1.200369in}}%
\pgfpathlineto{\pgfqpoint{3.680411in}{1.200369in}}%
\pgfpathlineto{\pgfqpoint{3.680411in}{1.292769in}}%
\pgfpathlineto{\pgfqpoint{3.543104in}{1.292769in}}%
\pgfpathlineto{\pgfqpoint{3.543104in}{1.200369in}}%
\pgfpathclose%
\pgfusepath{fill}%
\end{pgfscope}%
\begin{pgfscope}%
\pgfpathrectangle{\pgfqpoint{3.319565in}{0.488889in}}{\pgfqpoint{0.360845in}{0.924000in}}%
\pgfusepath{clip}%
\pgfsetbuttcap%
\pgfsetmiterjoin%
\definecolor{currentfill}{rgb}{0.000000,0.372549,0.813725}%
\pgfsetfillcolor{currentfill}%
\pgfsetlinewidth{0.000000pt}%
\definecolor{currentstroke}{rgb}{0.000000,0.000000,0.000000}%
\pgfsetstrokecolor{currentstroke}%
\pgfsetstrokeopacity{0.000000}%
\pgfsetdash{}{0pt}%
\pgfpathmoveto{\pgfqpoint{3.542527in}{1.209609in}}%
\pgfpathlineto{\pgfqpoint{3.680411in}{1.209609in}}%
\pgfpathlineto{\pgfqpoint{3.680411in}{1.302009in}}%
\pgfpathlineto{\pgfqpoint{3.542527in}{1.302009in}}%
\pgfpathlineto{\pgfqpoint{3.542527in}{1.209609in}}%
\pgfpathclose%
\pgfusepath{fill}%
\end{pgfscope}%
\begin{pgfscope}%
\pgfpathrectangle{\pgfqpoint{3.319565in}{0.488889in}}{\pgfqpoint{0.360845in}{0.924000in}}%
\pgfusepath{clip}%
\pgfsetbuttcap%
\pgfsetmiterjoin%
\definecolor{currentfill}{rgb}{0.000000,0.368627,0.815686}%
\pgfsetfillcolor{currentfill}%
\pgfsetlinewidth{0.000000pt}%
\definecolor{currentstroke}{rgb}{0.000000,0.000000,0.000000}%
\pgfsetstrokecolor{currentstroke}%
\pgfsetstrokeopacity{0.000000}%
\pgfsetdash{}{0pt}%
\pgfpathmoveto{\pgfqpoint{3.542427in}{1.218849in}}%
\pgfpathlineto{\pgfqpoint{3.680411in}{1.218849in}}%
\pgfpathlineto{\pgfqpoint{3.680411in}{1.311249in}}%
\pgfpathlineto{\pgfqpoint{3.542427in}{1.311249in}}%
\pgfpathlineto{\pgfqpoint{3.542427in}{1.218849in}}%
\pgfpathclose%
\pgfusepath{fill}%
\end{pgfscope}%
\begin{pgfscope}%
\pgfpathrectangle{\pgfqpoint{3.319565in}{0.488889in}}{\pgfqpoint{0.360845in}{0.924000in}}%
\pgfusepath{clip}%
\pgfsetbuttcap%
\pgfsetmiterjoin%
\definecolor{currentfill}{rgb}{0.000000,0.364706,0.817647}%
\pgfsetfillcolor{currentfill}%
\pgfsetlinewidth{0.000000pt}%
\definecolor{currentstroke}{rgb}{0.000000,0.000000,0.000000}%
\pgfsetstrokecolor{currentstroke}%
\pgfsetstrokeopacity{0.000000}%
\pgfsetdash{}{0pt}%
\pgfpathmoveto{\pgfqpoint{3.541911in}{1.228089in}}%
\pgfpathlineto{\pgfqpoint{3.680411in}{1.228089in}}%
\pgfpathlineto{\pgfqpoint{3.680411in}{1.320489in}}%
\pgfpathlineto{\pgfqpoint{3.541911in}{1.320489in}}%
\pgfpathlineto{\pgfqpoint{3.541911in}{1.228089in}}%
\pgfpathclose%
\pgfusepath{fill}%
\end{pgfscope}%
\begin{pgfscope}%
\pgfpathrectangle{\pgfqpoint{3.319565in}{0.488889in}}{\pgfqpoint{0.360845in}{0.924000in}}%
\pgfusepath{clip}%
\pgfsetbuttcap%
\pgfsetmiterjoin%
\definecolor{currentfill}{rgb}{0.000000,0.356863,0.821569}%
\pgfsetfillcolor{currentfill}%
\pgfsetlinewidth{0.000000pt}%
\definecolor{currentstroke}{rgb}{0.000000,0.000000,0.000000}%
\pgfsetstrokecolor{currentstroke}%
\pgfsetstrokeopacity{0.000000}%
\pgfsetdash{}{0pt}%
\pgfpathmoveto{\pgfqpoint{3.541325in}{1.237329in}}%
\pgfpathlineto{\pgfqpoint{3.680411in}{1.237329in}}%
\pgfpathlineto{\pgfqpoint{3.680411in}{1.329729in}}%
\pgfpathlineto{\pgfqpoint{3.541325in}{1.329729in}}%
\pgfpathlineto{\pgfqpoint{3.541325in}{1.237329in}}%
\pgfpathclose%
\pgfusepath{fill}%
\end{pgfscope}%
\begin{pgfscope}%
\pgfpathrectangle{\pgfqpoint{3.319565in}{0.488889in}}{\pgfqpoint{0.360845in}{0.924000in}}%
\pgfusepath{clip}%
\pgfsetbuttcap%
\pgfsetmiterjoin%
\definecolor{currentfill}{rgb}{0.000000,0.352941,0.823529}%
\pgfsetfillcolor{currentfill}%
\pgfsetlinewidth{0.000000pt}%
\definecolor{currentstroke}{rgb}{0.000000,0.000000,0.000000}%
\pgfsetstrokecolor{currentstroke}%
\pgfsetstrokeopacity{0.000000}%
\pgfsetdash{}{0pt}%
\pgfpathmoveto{\pgfqpoint{3.540997in}{1.246569in}}%
\pgfpathlineto{\pgfqpoint{3.680411in}{1.246569in}}%
\pgfpathlineto{\pgfqpoint{3.680411in}{1.338969in}}%
\pgfpathlineto{\pgfqpoint{3.540997in}{1.338969in}}%
\pgfpathlineto{\pgfqpoint{3.540997in}{1.246569in}}%
\pgfpathclose%
\pgfusepath{fill}%
\end{pgfscope}%
\begin{pgfscope}%
\pgfpathrectangle{\pgfqpoint{3.319565in}{0.488889in}}{\pgfqpoint{0.360845in}{0.924000in}}%
\pgfusepath{clip}%
\pgfsetbuttcap%
\pgfsetmiterjoin%
\definecolor{currentfill}{rgb}{0.000000,0.352941,0.823529}%
\pgfsetfillcolor{currentfill}%
\pgfsetlinewidth{0.000000pt}%
\definecolor{currentstroke}{rgb}{0.000000,0.000000,0.000000}%
\pgfsetstrokecolor{currentstroke}%
\pgfsetstrokeopacity{0.000000}%
\pgfsetdash{}{0pt}%
\pgfpathmoveto{\pgfqpoint{3.540986in}{1.255809in}}%
\pgfpathlineto{\pgfqpoint{3.680411in}{1.255809in}}%
\pgfpathlineto{\pgfqpoint{3.680411in}{1.348209in}}%
\pgfpathlineto{\pgfqpoint{3.540986in}{1.348209in}}%
\pgfpathlineto{\pgfqpoint{3.540986in}{1.255809in}}%
\pgfpathclose%
\pgfusepath{fill}%
\end{pgfscope}%
\begin{pgfscope}%
\pgfpathrectangle{\pgfqpoint{3.319565in}{0.488889in}}{\pgfqpoint{0.360845in}{0.924000in}}%
\pgfusepath{clip}%
\pgfsetbuttcap%
\pgfsetmiterjoin%
\definecolor{currentfill}{rgb}{0.000000,0.352941,0.823529}%
\pgfsetfillcolor{currentfill}%
\pgfsetlinewidth{0.000000pt}%
\definecolor{currentstroke}{rgb}{0.000000,0.000000,0.000000}%
\pgfsetstrokecolor{currentstroke}%
\pgfsetstrokeopacity{0.000000}%
\pgfsetdash{}{0pt}%
\pgfpathmoveto{\pgfqpoint{3.540986in}{1.265049in}}%
\pgfpathlineto{\pgfqpoint{3.680411in}{1.265049in}}%
\pgfpathlineto{\pgfqpoint{3.680411in}{1.357449in}}%
\pgfpathlineto{\pgfqpoint{3.540986in}{1.357449in}}%
\pgfpathlineto{\pgfqpoint{3.540986in}{1.265049in}}%
\pgfpathclose%
\pgfusepath{fill}%
\end{pgfscope}%
\begin{pgfscope}%
\pgfpathrectangle{\pgfqpoint{3.319565in}{0.488889in}}{\pgfqpoint{0.360845in}{0.924000in}}%
\pgfusepath{clip}%
\pgfsetbuttcap%
\pgfsetmiterjoin%
\definecolor{currentfill}{rgb}{0.000000,0.341176,0.829412}%
\pgfsetfillcolor{currentfill}%
\pgfsetlinewidth{0.000000pt}%
\definecolor{currentstroke}{rgb}{0.000000,0.000000,0.000000}%
\pgfsetstrokecolor{currentstroke}%
\pgfsetstrokeopacity{0.000000}%
\pgfsetdash{}{0pt}%
\pgfpathmoveto{\pgfqpoint{3.540080in}{1.274289in}}%
\pgfpathlineto{\pgfqpoint{3.680411in}{1.274289in}}%
\pgfpathlineto{\pgfqpoint{3.680411in}{1.366689in}}%
\pgfpathlineto{\pgfqpoint{3.540080in}{1.366689in}}%
\pgfpathlineto{\pgfqpoint{3.540080in}{1.274289in}}%
\pgfpathclose%
\pgfusepath{fill}%
\end{pgfscope}%
\begin{pgfscope}%
\pgfpathrectangle{\pgfqpoint{3.319565in}{0.488889in}}{\pgfqpoint{0.360845in}{0.924000in}}%
\pgfusepath{clip}%
\pgfsetbuttcap%
\pgfsetmiterjoin%
\definecolor{currentfill}{rgb}{0.000000,0.333333,0.833333}%
\pgfsetfillcolor{currentfill}%
\pgfsetlinewidth{0.000000pt}%
\definecolor{currentstroke}{rgb}{0.000000,0.000000,0.000000}%
\pgfsetstrokecolor{currentstroke}%
\pgfsetstrokeopacity{0.000000}%
\pgfsetdash{}{0pt}%
\pgfpathmoveto{\pgfqpoint{3.539586in}{1.283529in}}%
\pgfpathlineto{\pgfqpoint{3.680411in}{1.283529in}}%
\pgfpathlineto{\pgfqpoint{3.680411in}{1.375929in}}%
\pgfpathlineto{\pgfqpoint{3.539586in}{1.375929in}}%
\pgfpathlineto{\pgfqpoint{3.539586in}{1.283529in}}%
\pgfpathclose%
\pgfusepath{fill}%
\end{pgfscope}%
\begin{pgfscope}%
\pgfpathrectangle{\pgfqpoint{3.319565in}{0.488889in}}{\pgfqpoint{0.360845in}{0.924000in}}%
\pgfusepath{clip}%
\pgfsetbuttcap%
\pgfsetmiterjoin%
\definecolor{currentfill}{rgb}{0.000000,0.333333,0.833333}%
\pgfsetfillcolor{currentfill}%
\pgfsetlinewidth{0.000000pt}%
\definecolor{currentstroke}{rgb}{0.000000,0.000000,0.000000}%
\pgfsetstrokecolor{currentstroke}%
\pgfsetstrokeopacity{0.000000}%
\pgfsetdash{}{0pt}%
\pgfpathmoveto{\pgfqpoint{3.539576in}{1.292769in}}%
\pgfpathlineto{\pgfqpoint{3.680411in}{1.292769in}}%
\pgfpathlineto{\pgfqpoint{3.680411in}{1.385169in}}%
\pgfpathlineto{\pgfqpoint{3.539576in}{1.385169in}}%
\pgfpathlineto{\pgfqpoint{3.539576in}{1.292769in}}%
\pgfpathclose%
\pgfusepath{fill}%
\end{pgfscope}%
\begin{pgfscope}%
\pgfpathrectangle{\pgfqpoint{3.319565in}{0.488889in}}{\pgfqpoint{0.360845in}{0.924000in}}%
\pgfusepath{clip}%
\pgfsetbuttcap%
\pgfsetmiterjoin%
\definecolor{currentfill}{rgb}{0.000000,0.333333,0.833333}%
\pgfsetfillcolor{currentfill}%
\pgfsetlinewidth{0.000000pt}%
\definecolor{currentstroke}{rgb}{0.000000,0.000000,0.000000}%
\pgfsetstrokecolor{currentstroke}%
\pgfsetstrokeopacity{0.000000}%
\pgfsetdash{}{0pt}%
\pgfpathmoveto{\pgfqpoint{3.539297in}{1.302009in}}%
\pgfpathlineto{\pgfqpoint{3.680411in}{1.302009in}}%
\pgfpathlineto{\pgfqpoint{3.680411in}{1.394409in}}%
\pgfpathlineto{\pgfqpoint{3.539297in}{1.394409in}}%
\pgfpathlineto{\pgfqpoint{3.539297in}{1.302009in}}%
\pgfpathclose%
\pgfusepath{fill}%
\end{pgfscope}%
\begin{pgfscope}%
\pgfpathrectangle{\pgfqpoint{3.319565in}{0.488889in}}{\pgfqpoint{0.360845in}{0.924000in}}%
\pgfusepath{clip}%
\pgfsetbuttcap%
\pgfsetmiterjoin%
\definecolor{currentfill}{rgb}{0.000000,0.333333,0.833333}%
\pgfsetfillcolor{currentfill}%
\pgfsetlinewidth{0.000000pt}%
\definecolor{currentstroke}{rgb}{0.000000,0.000000,0.000000}%
\pgfsetstrokecolor{currentstroke}%
\pgfsetstrokeopacity{0.000000}%
\pgfsetdash{}{0pt}%
\pgfpathmoveto{\pgfqpoint{3.539297in}{1.311249in}}%
\pgfpathlineto{\pgfqpoint{3.680411in}{1.311249in}}%
\pgfpathlineto{\pgfqpoint{3.680411in}{1.403649in}}%
\pgfpathlineto{\pgfqpoint{3.539297in}{1.403649in}}%
\pgfpathlineto{\pgfqpoint{3.539297in}{1.311249in}}%
\pgfpathclose%
\pgfusepath{fill}%
\end{pgfscope}%
\begin{pgfscope}%
\pgfpathrectangle{\pgfqpoint{3.319565in}{0.488889in}}{\pgfqpoint{0.360845in}{0.924000in}}%
\pgfusepath{clip}%
\pgfsetbuttcap%
\pgfsetmiterjoin%
\definecolor{currentfill}{rgb}{0.000000,0.313725,0.843137}%
\pgfsetfillcolor{currentfill}%
\pgfsetlinewidth{0.000000pt}%
\definecolor{currentstroke}{rgb}{0.000000,0.000000,0.000000}%
\pgfsetstrokecolor{currentstroke}%
\pgfsetstrokeopacity{0.000000}%
\pgfsetdash{}{0pt}%
\pgfpathmoveto{\pgfqpoint{3.537936in}{1.320489in}}%
\pgfpathlineto{\pgfqpoint{3.680411in}{1.320489in}}%
\pgfpathlineto{\pgfqpoint{3.680411in}{1.412889in}}%
\pgfpathlineto{\pgfqpoint{3.537936in}{1.412889in}}%
\pgfpathlineto{\pgfqpoint{3.537936in}{1.320489in}}%
\pgfpathclose%
\pgfusepath{fill}%
\end{pgfscope}%
\begin{pgfscope}%
\pgfpathrectangle{\pgfqpoint{3.319565in}{0.488889in}}{\pgfqpoint{0.360845in}{0.924000in}}%
\pgfusepath{clip}%
\pgfsetbuttcap%
\pgfsetmiterjoin%
\definecolor{currentfill}{rgb}{0.000000,0.301961,0.849020}%
\pgfsetfillcolor{currentfill}%
\pgfsetlinewidth{0.000000pt}%
\definecolor{currentstroke}{rgb}{0.000000,0.000000,0.000000}%
\pgfsetstrokecolor{currentstroke}%
\pgfsetstrokeopacity{0.000000}%
\pgfsetdash{}{0pt}%
\pgfpathmoveto{\pgfqpoint{3.536971in}{1.329729in}}%
\pgfpathlineto{\pgfqpoint{3.680411in}{1.329729in}}%
\pgfpathlineto{\pgfqpoint{3.680411in}{1.422129in}}%
\pgfpathlineto{\pgfqpoint{3.536971in}{1.422129in}}%
\pgfpathlineto{\pgfqpoint{3.536971in}{1.329729in}}%
\pgfpathclose%
\pgfusepath{fill}%
\end{pgfscope}%
\begin{pgfscope}%
\pgfpathrectangle{\pgfqpoint{3.319565in}{0.488889in}}{\pgfqpoint{0.360845in}{0.924000in}}%
\pgfusepath{clip}%
\pgfsetbuttcap%
\pgfsetmiterjoin%
\definecolor{currentfill}{rgb}{0.000000,0.290196,0.854902}%
\pgfsetfillcolor{currentfill}%
\pgfsetlinewidth{0.000000pt}%
\definecolor{currentstroke}{rgb}{0.000000,0.000000,0.000000}%
\pgfsetstrokecolor{currentstroke}%
\pgfsetstrokeopacity{0.000000}%
\pgfsetdash{}{0pt}%
\pgfpathmoveto{\pgfqpoint{3.535827in}{1.338969in}}%
\pgfpathlineto{\pgfqpoint{3.680411in}{1.338969in}}%
\pgfpathlineto{\pgfqpoint{3.680411in}{1.431369in}}%
\pgfpathlineto{\pgfqpoint{3.535827in}{1.431369in}}%
\pgfpathlineto{\pgfqpoint{3.535827in}{1.338969in}}%
\pgfpathclose%
\pgfusepath{fill}%
\end{pgfscope}%
\begin{pgfscope}%
\pgfpathrectangle{\pgfqpoint{3.319565in}{0.488889in}}{\pgfqpoint{0.360845in}{0.924000in}}%
\pgfusepath{clip}%
\pgfsetbuttcap%
\pgfsetmiterjoin%
\definecolor{currentfill}{rgb}{0.000000,0.286275,0.856863}%
\pgfsetfillcolor{currentfill}%
\pgfsetlinewidth{0.000000pt}%
\definecolor{currentstroke}{rgb}{0.000000,0.000000,0.000000}%
\pgfsetstrokecolor{currentstroke}%
\pgfsetstrokeopacity{0.000000}%
\pgfsetdash{}{0pt}%
\pgfpathmoveto{\pgfqpoint{3.535510in}{1.348209in}}%
\pgfpathlineto{\pgfqpoint{3.680411in}{1.348209in}}%
\pgfpathlineto{\pgfqpoint{3.680411in}{1.440609in}}%
\pgfpathlineto{\pgfqpoint{3.535510in}{1.440609in}}%
\pgfpathlineto{\pgfqpoint{3.535510in}{1.348209in}}%
\pgfpathclose%
\pgfusepath{fill}%
\end{pgfscope}%
\begin{pgfscope}%
\pgfpathrectangle{\pgfqpoint{3.319565in}{0.488889in}}{\pgfqpoint{0.360845in}{0.924000in}}%
\pgfusepath{clip}%
\pgfsetbuttcap%
\pgfsetmiterjoin%
\definecolor{currentfill}{rgb}{0.000000,0.243137,0.878431}%
\pgfsetfillcolor{currentfill}%
\pgfsetlinewidth{0.000000pt}%
\definecolor{currentstroke}{rgb}{0.000000,0.000000,0.000000}%
\pgfsetstrokecolor{currentstroke}%
\pgfsetstrokeopacity{0.000000}%
\pgfsetdash{}{0pt}%
\pgfpathmoveto{\pgfqpoint{3.532156in}{1.357449in}}%
\pgfpathlineto{\pgfqpoint{3.680411in}{1.357449in}}%
\pgfpathlineto{\pgfqpoint{3.680411in}{1.449849in}}%
\pgfpathlineto{\pgfqpoint{3.532156in}{1.449849in}}%
\pgfpathlineto{\pgfqpoint{3.532156in}{1.357449in}}%
\pgfpathclose%
\pgfusepath{fill}%
\end{pgfscope}%
\begin{pgfscope}%
\pgfpathrectangle{\pgfqpoint{3.319565in}{0.488889in}}{\pgfqpoint{0.360845in}{0.924000in}}%
\pgfusepath{clip}%
\pgfsetbuttcap%
\pgfsetmiterjoin%
\definecolor{currentfill}{rgb}{0.000000,0.243137,0.878431}%
\pgfsetfillcolor{currentfill}%
\pgfsetlinewidth{0.000000pt}%
\definecolor{currentstroke}{rgb}{0.000000,0.000000,0.000000}%
\pgfsetstrokecolor{currentstroke}%
\pgfsetstrokeopacity{0.000000}%
\pgfsetdash{}{0pt}%
\pgfpathmoveto{\pgfqpoint{3.532156in}{1.366689in}}%
\pgfpathlineto{\pgfqpoint{3.680411in}{1.366689in}}%
\pgfpathlineto{\pgfqpoint{3.680411in}{1.459089in}}%
\pgfpathlineto{\pgfqpoint{3.532156in}{1.459089in}}%
\pgfpathlineto{\pgfqpoint{3.532156in}{1.366689in}}%
\pgfpathclose%
\pgfusepath{fill}%
\end{pgfscope}%
\begin{pgfscope}%
\pgfpathrectangle{\pgfqpoint{3.319565in}{0.488889in}}{\pgfqpoint{0.360845in}{0.924000in}}%
\pgfusepath{clip}%
\pgfsetbuttcap%
\pgfsetmiterjoin%
\definecolor{currentfill}{rgb}{0.000000,0.219608,0.890196}%
\pgfsetfillcolor{currentfill}%
\pgfsetlinewidth{0.000000pt}%
\definecolor{currentstroke}{rgb}{0.000000,0.000000,0.000000}%
\pgfsetstrokecolor{currentstroke}%
\pgfsetstrokeopacity{0.000000}%
\pgfsetdash{}{0pt}%
\pgfpathmoveto{\pgfqpoint{3.530237in}{1.375929in}}%
\pgfpathlineto{\pgfqpoint{3.680411in}{1.375929in}}%
\pgfpathlineto{\pgfqpoint{3.680411in}{1.468329in}}%
\pgfpathlineto{\pgfqpoint{3.530237in}{1.468329in}}%
\pgfpathlineto{\pgfqpoint{3.530237in}{1.375929in}}%
\pgfpathclose%
\pgfusepath{fill}%
\end{pgfscope}%
\begin{pgfscope}%
\pgfpathrectangle{\pgfqpoint{3.319565in}{0.488889in}}{\pgfqpoint{0.360845in}{0.924000in}}%
\pgfusepath{clip}%
\pgfsetbuttcap%
\pgfsetmiterjoin%
\definecolor{currentfill}{rgb}{0.000000,0.196078,0.901961}%
\pgfsetfillcolor{currentfill}%
\pgfsetlinewidth{0.000000pt}%
\definecolor{currentstroke}{rgb}{0.000000,0.000000,0.000000}%
\pgfsetstrokecolor{currentstroke}%
\pgfsetstrokeopacity{0.000000}%
\pgfsetdash{}{0pt}%
\pgfpathmoveto{\pgfqpoint{3.528298in}{1.385169in}}%
\pgfpathlineto{\pgfqpoint{3.680411in}{1.385169in}}%
\pgfpathlineto{\pgfqpoint{3.680411in}{1.477569in}}%
\pgfpathlineto{\pgfqpoint{3.528298in}{1.477569in}}%
\pgfpathlineto{\pgfqpoint{3.528298in}{1.385169in}}%
\pgfpathclose%
\pgfusepath{fill}%
\end{pgfscope}%
\begin{pgfscope}%
\pgfpathrectangle{\pgfqpoint{3.319565in}{0.488889in}}{\pgfqpoint{0.360845in}{0.924000in}}%
\pgfusepath{clip}%
\pgfsetbuttcap%
\pgfsetmiterjoin%
\definecolor{currentfill}{rgb}{0.000000,0.196078,0.901961}%
\pgfsetfillcolor{currentfill}%
\pgfsetlinewidth{0.000000pt}%
\definecolor{currentstroke}{rgb}{0.000000,0.000000,0.000000}%
\pgfsetstrokecolor{currentstroke}%
\pgfsetstrokeopacity{0.000000}%
\pgfsetdash{}{0pt}%
\pgfpathmoveto{\pgfqpoint{3.528298in}{1.394409in}}%
\pgfpathlineto{\pgfqpoint{3.680411in}{1.394409in}}%
\pgfpathlineto{\pgfqpoint{3.680411in}{1.486809in}}%
\pgfpathlineto{\pgfqpoint{3.528298in}{1.486809in}}%
\pgfpathlineto{\pgfqpoint{3.528298in}{1.394409in}}%
\pgfpathclose%
\pgfusepath{fill}%
\end{pgfscope}%
\begin{pgfscope}%
\pgfpathrectangle{\pgfqpoint{3.319565in}{0.488889in}}{\pgfqpoint{0.360845in}{0.924000in}}%
\pgfusepath{clip}%
\pgfsetbuttcap%
\pgfsetmiterjoin%
\definecolor{currentfill}{rgb}{0.000000,0.000000,1.000000}%
\pgfsetfillcolor{currentfill}%
\pgfsetlinewidth{0.000000pt}%
\definecolor{currentstroke}{rgb}{0.000000,0.000000,0.000000}%
\pgfsetstrokecolor{currentstroke}%
\pgfsetstrokeopacity{0.000000}%
\pgfsetdash{}{0pt}%
\pgfpathmoveto{\pgfqpoint{3.512377in}{1.403649in}}%
\pgfpathlineto{\pgfqpoint{3.680411in}{1.403649in}}%
\pgfpathlineto{\pgfqpoint{3.680411in}{1.496049in}}%
\pgfpathlineto{\pgfqpoint{3.512377in}{1.496049in}}%
\pgfpathlineto{\pgfqpoint{3.512377in}{1.403649in}}%
\pgfpathclose%
\pgfusepath{fill}%
\end{pgfscope}%
\begin{pgfscope}%
\definecolor{textcolor}{rgb}{0.000000,0.000000,0.000000}%
\pgfsetstrokecolor{textcolor}%
\pgfsetfillcolor{textcolor}%
\pgftext[x=3.471121in,y=1.468329in,left,base]{\color{textcolor}\setmainfont{Lato}\rmfamily\fontsize{7.000000}{8.400000}\selectfont 46.6\%}%
\end{pgfscope}%
\begin{pgfscope}%
\definecolor{textcolor}{rgb}{0.000000,0.000000,0.000000}%
\pgfsetstrokecolor{textcolor}%
\pgfsetfillcolor{textcolor}%
\pgftext[x=3.471121in,y=0.359529in,left,base]{\color{textcolor}\setmainfont{Lato}\rmfamily\fontsize{7.000000}{8.400000}\selectfont 24.1\%}%
\end{pgfscope}%
\end{pgfpicture}%
\makeatother%
\endgroup%


\vspace{-4mm}
\footnotesize{Source: American Community Survey, Dorn, Author's Calculations}
\end{minipage} \hspace{3mm}
\begin{minipage}{0.24\textwidth}
\vspace{-1mm}

\footnotesize \input{text/acs_ftfy.txt}
\end{minipage}
\newpage
\begin{minipage}{0.76\textwidth} 
\normalsize \textbf{Employment Rates of Largest Commuter Zones, \input{text/acs_ftfy_year.txt}}
\vspace{-1mm}

\footnotesize \hspace{42mm} \textit{all ages} \hspace{20mm} \textit{age 25--54}\\
\rowcolors{1}{}{black!5} \setlength{\tabcolsep}{3.1pt} \color{black!90}
	{\renewcommand{\arraystretch}{1.55}
		\begin{tabular}{p{24mm} R{15mm} R{15mm} R{15mm} R{15mm}}
		 	& full-time \& full-year & employed & full-time \& full-year & employed \\
			 \input data/acs_emp_tables.tex \hline
		\end{tabular}}	
\vspace{-2mm}

\footnotesize{Source: American Community Survey, Dorn, Author's Calculations}
\vspace{6mm}

\small The share of the age 16 and older population that is employed in any amount at a point in time varies by state and over time. \input{text/state_emp.txt}
\end{minipage}
\vspace{2mm}

\normalsize \textbf{Employment Rate by State}\\
\footnotesize{\textit{employed share of population, percent}}\\
\vspace{-2mm}
\hspace{-8mm} %% Creator: Matplotlib, PGF backend
%%
%% To include the figure in your LaTeX document, write
%%   \input{<filename>.pgf}
%%
%% Make sure the required packages are loaded in your preamble
%%   \usepackage{pgf}
%%
%% Also ensure that all the required font packages are loaded; for instance,
%% the lmodern package is sometimes necessary when using math font.
%%   \usepackage{lmodern}
%%
%% Figures using additional raster images can only be included by \input if
%% they are in the same directory as the main LaTeX file. For loading figures
%% from other directories you can use the `import` package
%%   \usepackage{import}
%%
%% and then include the figures with
%%   \import{<path to file>}{<filename>.pgf}
%%
%% Matplotlib used the following preamble
%%   
%%   \usepackage{fontspec}
%%   \setmainfont{DejaVuSerif.ttf}[Path=\detokenize{/home/brian/miniconda3/lib/python3.8/site-packages/matplotlib/mpl-data/fonts/ttf/}]
%%   \setsansfont{DejaVuSans.ttf}[Path=\detokenize{/home/brian/miniconda3/lib/python3.8/site-packages/matplotlib/mpl-data/fonts/ttf/}]
%%   \setmonofont{DejaVuSansMono.ttf}[Path=\detokenize{/home/brian/miniconda3/lib/python3.8/site-packages/matplotlib/mpl-data/fonts/ttf/}]
%%   \makeatletter\@ifpackageloaded{underscore}{}{\usepackage[strings]{underscore}}\makeatother
%%
\begingroup%
\makeatletter%
\begin{pgfpicture}%
\pgfpathrectangle{\pgfpointorigin}{\pgfqpoint{6.487985in}{2.088185in}}%
\pgfusepath{use as bounding box, clip}%
\begin{pgfscope}%
\pgfsetbuttcap%
\pgfsetmiterjoin%
\definecolor{currentfill}{rgb}{1.000000,1.000000,1.000000}%
\pgfsetfillcolor{currentfill}%
\pgfsetlinewidth{0.000000pt}%
\definecolor{currentstroke}{rgb}{1.000000,1.000000,1.000000}%
\pgfsetstrokecolor{currentstroke}%
\pgfsetdash{}{0pt}%
\pgfpathmoveto{\pgfqpoint{0.000000in}{0.000000in}}%
\pgfpathlineto{\pgfqpoint{6.487985in}{0.000000in}}%
\pgfpathlineto{\pgfqpoint{6.487985in}{2.088185in}}%
\pgfpathlineto{\pgfqpoint{0.000000in}{2.088185in}}%
\pgfpathlineto{\pgfqpoint{0.000000in}{0.000000in}}%
\pgfpathclose%
\pgfusepath{fill}%
\end{pgfscope}%
\begin{pgfscope}%
\pgfpathrectangle{\pgfqpoint{0.100000in}{0.100000in}}{\pgfqpoint{2.862985in}{1.832778in}}%
\pgfusepath{clip}%
\pgfsetbuttcap%
\pgfsetmiterjoin%
\definecolor{currentfill}{rgb}{0.768090,0.244817,0.498465}%
\pgfsetfillcolor{currentfill}%
\pgfsetlinewidth{0.000000pt}%
\definecolor{currentstroke}{rgb}{0.000000,0.000000,0.000000}%
\pgfsetstrokecolor{currentstroke}%
\pgfsetstrokeopacity{0.000000}%
\pgfsetdash{}{0pt}%
\pgfpathmoveto{\pgfqpoint{1.055609in}{0.529755in}}%
\pgfpathlineto{\pgfqpoint{1.047380in}{0.529889in}}%
\pgfpathlineto{\pgfqpoint{1.041780in}{0.536088in}}%
\pgfpathlineto{\pgfqpoint{1.037972in}{0.552550in}}%
\pgfpathlineto{\pgfqpoint{1.046822in}{0.554612in}}%
\pgfpathlineto{\pgfqpoint{1.055813in}{0.552394in}}%
\pgfpathlineto{\pgfqpoint{1.064963in}{0.545909in}}%
\pgfpathlineto{\pgfqpoint{1.064959in}{0.536940in}}%
\pgfpathlineto{\pgfqpoint{1.055609in}{0.529755in}}%
\pgfpathclose%
\pgfusepath{fill}%
\end{pgfscope}%
\begin{pgfscope}%
\pgfpathrectangle{\pgfqpoint{0.100000in}{0.100000in}}{\pgfqpoint{2.862985in}{1.832778in}}%
\pgfusepath{clip}%
\pgfsetbuttcap%
\pgfsetmiterjoin%
\definecolor{currentfill}{rgb}{0.768090,0.244817,0.498465}%
\pgfsetfillcolor{currentfill}%
\pgfsetlinewidth{0.000000pt}%
\definecolor{currentstroke}{rgb}{0.000000,0.000000,0.000000}%
\pgfsetstrokecolor{currentstroke}%
\pgfsetstrokeopacity{0.000000}%
\pgfsetdash{}{0pt}%
\pgfpathmoveto{\pgfqpoint{1.105694in}{0.432023in}}%
\pgfpathlineto{\pgfqpoint{1.096488in}{0.435528in}}%
\pgfpathlineto{\pgfqpoint{1.096446in}{0.442404in}}%
\pgfpathlineto{\pgfqpoint{1.085326in}{0.448219in}}%
\pgfpathlineto{\pgfqpoint{1.086680in}{0.462860in}}%
\pgfpathlineto{\pgfqpoint{1.089901in}{0.469769in}}%
\pgfpathlineto{\pgfqpoint{1.096701in}{0.463410in}}%
\pgfpathlineto{\pgfqpoint{1.107760in}{0.464472in}}%
\pgfpathlineto{\pgfqpoint{1.104967in}{0.444236in}}%
\pgfpathlineto{\pgfqpoint{1.105694in}{0.432023in}}%
\pgfpathclose%
\pgfusepath{fill}%
\end{pgfscope}%
\begin{pgfscope}%
\pgfpathrectangle{\pgfqpoint{0.100000in}{0.100000in}}{\pgfqpoint{2.862985in}{1.832778in}}%
\pgfusepath{clip}%
\pgfsetbuttcap%
\pgfsetmiterjoin%
\definecolor{currentfill}{rgb}{0.768090,0.244817,0.498465}%
\pgfsetfillcolor{currentfill}%
\pgfsetlinewidth{0.000000pt}%
\definecolor{currentstroke}{rgb}{0.000000,0.000000,0.000000}%
\pgfsetstrokecolor{currentstroke}%
\pgfsetstrokeopacity{0.000000}%
\pgfsetdash{}{0pt}%
\pgfpathmoveto{\pgfqpoint{1.144333in}{0.386887in}}%
\pgfpathlineto{\pgfqpoint{1.132651in}{0.388876in}}%
\pgfpathlineto{\pgfqpoint{1.126191in}{0.398877in}}%
\pgfpathlineto{\pgfqpoint{1.125279in}{0.407100in}}%
\pgfpathlineto{\pgfqpoint{1.144333in}{0.386887in}}%
\pgfpathclose%
\pgfusepath{fill}%
\end{pgfscope}%
\begin{pgfscope}%
\pgfpathrectangle{\pgfqpoint{0.100000in}{0.100000in}}{\pgfqpoint{2.862985in}{1.832778in}}%
\pgfusepath{clip}%
\pgfsetbuttcap%
\pgfsetmiterjoin%
\definecolor{currentfill}{rgb}{0.768090,0.244817,0.498465}%
\pgfsetfillcolor{currentfill}%
\pgfsetlinewidth{0.000000pt}%
\definecolor{currentstroke}{rgb}{0.000000,0.000000,0.000000}%
\pgfsetstrokecolor{currentstroke}%
\pgfsetstrokeopacity{0.000000}%
\pgfsetdash{}{0pt}%
\pgfpathmoveto{\pgfqpoint{1.120394in}{0.388631in}}%
\pgfpathlineto{\pgfqpoint{1.126519in}{0.383659in}}%
\pgfpathlineto{\pgfqpoint{1.127848in}{0.375110in}}%
\pgfpathlineto{\pgfqpoint{1.115519in}{0.375583in}}%
\pgfpathlineto{\pgfqpoint{1.120394in}{0.388631in}}%
\pgfpathclose%
\pgfusepath{fill}%
\end{pgfscope}%
\begin{pgfscope}%
\pgfpathrectangle{\pgfqpoint{0.100000in}{0.100000in}}{\pgfqpoint{2.862985in}{1.832778in}}%
\pgfusepath{clip}%
\pgfsetbuttcap%
\pgfsetmiterjoin%
\definecolor{currentfill}{rgb}{0.768090,0.244817,0.498465}%
\pgfsetfillcolor{currentfill}%
\pgfsetlinewidth{0.000000pt}%
\definecolor{currentstroke}{rgb}{0.000000,0.000000,0.000000}%
\pgfsetstrokecolor{currentstroke}%
\pgfsetstrokeopacity{0.000000}%
\pgfsetdash{}{0pt}%
\pgfpathmoveto{\pgfqpoint{1.148521in}{0.340880in}}%
\pgfpathlineto{\pgfqpoint{1.136217in}{0.345303in}}%
\pgfpathlineto{\pgfqpoint{1.139744in}{0.356026in}}%
\pgfpathlineto{\pgfqpoint{1.134586in}{0.366952in}}%
\pgfpathlineto{\pgfqpoint{1.136206in}{0.374463in}}%
\pgfpathlineto{\pgfqpoint{1.146107in}{0.377326in}}%
\pgfpathlineto{\pgfqpoint{1.145916in}{0.365292in}}%
\pgfpathlineto{\pgfqpoint{1.157627in}{0.358618in}}%
\pgfpathlineto{\pgfqpoint{1.164326in}{0.342480in}}%
\pgfpathlineto{\pgfqpoint{1.156861in}{0.336266in}}%
\pgfpathlineto{\pgfqpoint{1.148521in}{0.340880in}}%
\pgfpathclose%
\pgfusepath{fill}%
\end{pgfscope}%
\begin{pgfscope}%
\pgfpathrectangle{\pgfqpoint{0.100000in}{0.100000in}}{\pgfqpoint{2.862985in}{1.832778in}}%
\pgfusepath{clip}%
\pgfsetbuttcap%
\pgfsetmiterjoin%
\definecolor{currentfill}{rgb}{0.768090,0.244817,0.498465}%
\pgfsetfillcolor{currentfill}%
\pgfsetlinewidth{0.000000pt}%
\definecolor{currentstroke}{rgb}{0.000000,0.000000,0.000000}%
\pgfsetstrokecolor{currentstroke}%
\pgfsetstrokeopacity{0.000000}%
\pgfsetdash{}{0pt}%
\pgfpathmoveto{\pgfqpoint{1.104084in}{0.231768in}}%
\pgfpathlineto{\pgfqpoint{1.098222in}{0.243877in}}%
\pgfpathlineto{\pgfqpoint{1.100818in}{0.252022in}}%
\pgfpathlineto{\pgfqpoint{1.112315in}{0.263296in}}%
\pgfpathlineto{\pgfqpoint{1.113404in}{0.271938in}}%
\pgfpathlineto{\pgfqpoint{1.120928in}{0.291072in}}%
\pgfpathlineto{\pgfqpoint{1.140718in}{0.294341in}}%
\pgfpathlineto{\pgfqpoint{1.144506in}{0.305541in}}%
\pgfpathlineto{\pgfqpoint{1.153232in}{0.301341in}}%
\pgfpathlineto{\pgfqpoint{1.170961in}{0.269908in}}%
\pgfpathlineto{\pgfqpoint{1.171185in}{0.259645in}}%
\pgfpathlineto{\pgfqpoint{1.166173in}{0.250763in}}%
\pgfpathlineto{\pgfqpoint{1.172532in}{0.230369in}}%
\pgfpathlineto{\pgfqpoint{1.168089in}{0.227081in}}%
\pgfpathlineto{\pgfqpoint{1.152955in}{0.228037in}}%
\pgfpathlineto{\pgfqpoint{1.134756in}{0.234119in}}%
\pgfpathlineto{\pgfqpoint{1.119460in}{0.236631in}}%
\pgfpathlineto{\pgfqpoint{1.115375in}{0.233063in}}%
\pgfpathlineto{\pgfqpoint{1.104084in}{0.231768in}}%
\pgfpathclose%
\pgfusepath{fill}%
\end{pgfscope}%
\begin{pgfscope}%
\pgfpathrectangle{\pgfqpoint{0.100000in}{0.100000in}}{\pgfqpoint{2.862985in}{1.832778in}}%
\pgfusepath{clip}%
\pgfsetbuttcap%
\pgfsetmiterjoin%
\definecolor{currentfill}{rgb}{0.660374,0.134144,0.588971}%
\pgfsetfillcolor{currentfill}%
\pgfsetlinewidth{0.000000pt}%
\definecolor{currentstroke}{rgb}{0.000000,0.000000,0.000000}%
\pgfsetstrokecolor{currentstroke}%
\pgfsetstrokeopacity{0.000000}%
\pgfsetdash{}{0pt}%
\pgfpathmoveto{\pgfqpoint{0.780273in}{1.824988in}}%
\pgfpathlineto{\pgfqpoint{0.765426in}{1.766159in}}%
\pgfpathlineto{\pgfqpoint{0.754869in}{1.724007in}}%
\pgfpathlineto{\pgfqpoint{0.742317in}{1.673299in}}%
\pgfpathlineto{\pgfqpoint{0.739852in}{1.660801in}}%
\pgfpathlineto{\pgfqpoint{0.741462in}{1.649338in}}%
\pgfpathlineto{\pgfqpoint{0.739329in}{1.638779in}}%
\pgfpathlineto{\pgfqpoint{0.651709in}{1.661643in}}%
\pgfpathlineto{\pgfqpoint{0.643794in}{1.658987in}}%
\pgfpathlineto{\pgfqpoint{0.637021in}{1.661292in}}%
\pgfpathlineto{\pgfqpoint{0.605660in}{1.662045in}}%
\pgfpathlineto{\pgfqpoint{0.595053in}{1.659055in}}%
\pgfpathlineto{\pgfqpoint{0.584520in}{1.660065in}}%
\pgfpathlineto{\pgfqpoint{0.580055in}{1.664746in}}%
\pgfpathlineto{\pgfqpoint{0.551979in}{1.663725in}}%
\pgfpathlineto{\pgfqpoint{0.547545in}{1.671374in}}%
\pgfpathlineto{\pgfqpoint{0.539127in}{1.675404in}}%
\pgfpathlineto{\pgfqpoint{0.526880in}{1.677860in}}%
\pgfpathlineto{\pgfqpoint{0.505740in}{1.674264in}}%
\pgfpathlineto{\pgfqpoint{0.497937in}{1.677787in}}%
\pgfpathlineto{\pgfqpoint{0.485905in}{1.687156in}}%
\pgfpathlineto{\pgfqpoint{0.489522in}{1.705519in}}%
\pgfpathlineto{\pgfqpoint{0.488254in}{1.714897in}}%
\pgfpathlineto{\pgfqpoint{0.478718in}{1.724852in}}%
\pgfpathlineto{\pgfqpoint{0.472610in}{1.724234in}}%
\pgfpathlineto{\pgfqpoint{0.468229in}{1.734304in}}%
\pgfpathlineto{\pgfqpoint{0.457849in}{1.738374in}}%
\pgfpathlineto{\pgfqpoint{0.450305in}{1.737856in}}%
\pgfpathlineto{\pgfqpoint{0.447864in}{1.748198in}}%
\pgfpathlineto{\pgfqpoint{0.455380in}{1.747098in}}%
\pgfpathlineto{\pgfqpoint{0.454783in}{1.761472in}}%
\pgfpathlineto{\pgfqpoint{0.451507in}{1.770008in}}%
\pgfpathlineto{\pgfqpoint{0.452951in}{1.780932in}}%
\pgfpathlineto{\pgfqpoint{0.456884in}{1.789149in}}%
\pgfpathlineto{\pgfqpoint{0.456662in}{1.804741in}}%
\pgfpathlineto{\pgfqpoint{0.454563in}{1.810398in}}%
\pgfpathlineto{\pgfqpoint{0.458188in}{1.828548in}}%
\pgfpathlineto{\pgfqpoint{0.457142in}{1.840246in}}%
\pgfpathlineto{\pgfqpoint{0.453517in}{1.845854in}}%
\pgfpathlineto{\pgfqpoint{0.454070in}{1.864192in}}%
\pgfpathlineto{\pgfqpoint{0.459342in}{1.877678in}}%
\pgfpathlineto{\pgfqpoint{0.465077in}{1.874378in}}%
\pgfpathlineto{\pgfqpoint{0.484064in}{1.854764in}}%
\pgfpathlineto{\pgfqpoint{0.507071in}{1.844031in}}%
\pgfpathlineto{\pgfqpoint{0.518852in}{1.842761in}}%
\pgfpathlineto{\pgfqpoint{0.529951in}{1.838191in}}%
\pgfpathlineto{\pgfqpoint{0.533101in}{1.822814in}}%
\pgfpathlineto{\pgfqpoint{0.523375in}{1.819901in}}%
\pgfpathlineto{\pgfqpoint{0.516929in}{1.807782in}}%
\pgfpathlineto{\pgfqpoint{0.524497in}{1.808455in}}%
\pgfpathlineto{\pgfqpoint{0.527549in}{1.813895in}}%
\pgfpathlineto{\pgfqpoint{0.538221in}{1.820697in}}%
\pgfpathlineto{\pgfqpoint{0.537670in}{1.810596in}}%
\pgfpathlineto{\pgfqpoint{0.530544in}{1.808967in}}%
\pgfpathlineto{\pgfqpoint{0.531727in}{1.795983in}}%
\pgfpathlineto{\pgfqpoint{0.523050in}{1.782976in}}%
\pgfpathlineto{\pgfqpoint{0.517496in}{1.789161in}}%
\pgfpathlineto{\pgfqpoint{0.500991in}{1.785798in}}%
\pgfpathlineto{\pgfqpoint{0.500123in}{1.778232in}}%
\pgfpathlineto{\pgfqpoint{0.505813in}{1.773622in}}%
\pgfpathlineto{\pgfqpoint{0.514493in}{1.773169in}}%
\pgfpathlineto{\pgfqpoint{0.526495in}{1.783019in}}%
\pgfpathlineto{\pgfqpoint{0.535994in}{1.783837in}}%
\pgfpathlineto{\pgfqpoint{0.536358in}{1.794741in}}%
\pgfpathlineto{\pgfqpoint{0.541254in}{1.810756in}}%
\pgfpathlineto{\pgfqpoint{0.552856in}{1.822666in}}%
\pgfpathlineto{\pgfqpoint{0.548988in}{1.831891in}}%
\pgfpathlineto{\pgfqpoint{0.551630in}{1.841823in}}%
\pgfpathlineto{\pgfqpoint{0.546483in}{1.851722in}}%
\pgfpathlineto{\pgfqpoint{0.553818in}{1.864381in}}%
\pgfpathlineto{\pgfqpoint{0.554917in}{1.871996in}}%
\pgfpathlineto{\pgfqpoint{0.548548in}{1.877001in}}%
\pgfpathlineto{\pgfqpoint{0.549609in}{1.889804in}}%
\pgfpathlineto{\pgfqpoint{0.625918in}{1.866620in}}%
\pgfpathlineto{\pgfqpoint{0.706952in}{1.843896in}}%
\pgfpathlineto{\pgfqpoint{0.780273in}{1.824988in}}%
\pgfpathclose%
\pgfusepath{fill}%
\end{pgfscope}%
\begin{pgfscope}%
\pgfpathrectangle{\pgfqpoint{0.100000in}{0.100000in}}{\pgfqpoint{2.862985in}{1.832778in}}%
\pgfusepath{clip}%
\pgfsetbuttcap%
\pgfsetmiterjoin%
\definecolor{currentfill}{rgb}{0.660374,0.134144,0.588971}%
\pgfsetfillcolor{currentfill}%
\pgfsetlinewidth{0.000000pt}%
\definecolor{currentstroke}{rgb}{0.000000,0.000000,0.000000}%
\pgfsetstrokecolor{currentstroke}%
\pgfsetstrokeopacity{0.000000}%
\pgfsetdash{}{0pt}%
\pgfpathmoveto{\pgfqpoint{0.537056in}{1.845512in}}%
\pgfpathlineto{\pgfqpoint{0.540813in}{1.831026in}}%
\pgfpathlineto{\pgfqpoint{0.546706in}{1.826214in}}%
\pgfpathlineto{\pgfqpoint{0.542004in}{1.820089in}}%
\pgfpathlineto{\pgfqpoint{0.537421in}{1.829064in}}%
\pgfpathlineto{\pgfqpoint{0.537056in}{1.845512in}}%
\pgfpathclose%
\pgfusepath{fill}%
\end{pgfscope}%
\begin{pgfscope}%
\pgfpathrectangle{\pgfqpoint{0.100000in}{0.100000in}}{\pgfqpoint{2.862985in}{1.832778in}}%
\pgfusepath{clip}%
\pgfsetbuttcap%
\pgfsetmiterjoin%
\definecolor{currentfill}{rgb}{0.697324,0.169573,0.560919}%
\pgfsetfillcolor{currentfill}%
\pgfsetlinewidth{0.000000pt}%
\definecolor{currentstroke}{rgb}{0.000000,0.000000,0.000000}%
\pgfsetstrokecolor{currentstroke}%
\pgfsetstrokeopacity{0.000000}%
\pgfsetdash{}{0pt}%
\pgfpathmoveto{\pgfqpoint{0.819624in}{1.815480in}}%
\pgfpathlineto{\pgfqpoint{0.901194in}{1.797175in}}%
\pgfpathlineto{\pgfqpoint{0.978016in}{1.781647in}}%
\pgfpathlineto{\pgfqpoint{1.037121in}{1.770817in}}%
\pgfpathlineto{\pgfqpoint{1.088650in}{1.762160in}}%
\pgfpathlineto{\pgfqpoint{1.140293in}{1.754214in}}%
\pgfpathlineto{\pgfqpoint{1.184272in}{1.748018in}}%
\pgfpathlineto{\pgfqpoint{1.228320in}{1.742337in}}%
\pgfpathlineto{\pgfqpoint{1.272432in}{1.737171in}}%
\pgfpathlineto{\pgfqpoint{1.314001in}{1.732781in}}%
\pgfpathlineto{\pgfqpoint{1.308284in}{1.669428in}}%
\pgfpathlineto{\pgfqpoint{1.299775in}{1.583815in}}%
\pgfpathlineto{\pgfqpoint{1.295312in}{1.539783in}}%
\pgfpathlineto{\pgfqpoint{1.288893in}{1.480441in}}%
\pgfpathlineto{\pgfqpoint{1.243453in}{1.485403in}}%
\pgfpathlineto{\pgfqpoint{1.191429in}{1.491308in}}%
\pgfpathlineto{\pgfqpoint{1.119194in}{1.501195in}}%
\pgfpathlineto{\pgfqpoint{1.086931in}{1.505797in}}%
\pgfpathlineto{\pgfqpoint{1.007445in}{1.518274in}}%
\pgfpathlineto{\pgfqpoint{0.980083in}{1.523280in}}%
\pgfpathlineto{\pgfqpoint{0.974380in}{1.490854in}}%
\pgfpathlineto{\pgfqpoint{0.971262in}{1.493168in}}%
\pgfpathlineto{\pgfqpoint{0.965403in}{1.508767in}}%
\pgfpathlineto{\pgfqpoint{0.957871in}{1.507806in}}%
\pgfpathlineto{\pgfqpoint{0.952480in}{1.499920in}}%
\pgfpathlineto{\pgfqpoint{0.941293in}{1.498866in}}%
\pgfpathlineto{\pgfqpoint{0.939045in}{1.502638in}}%
\pgfpathlineto{\pgfqpoint{0.928566in}{1.502181in}}%
\pgfpathlineto{\pgfqpoint{0.923297in}{1.505835in}}%
\pgfpathlineto{\pgfqpoint{0.915931in}{1.500183in}}%
\pgfpathlineto{\pgfqpoint{0.898038in}{1.505269in}}%
\pgfpathlineto{\pgfqpoint{0.890230in}{1.502538in}}%
\pgfpathlineto{\pgfqpoint{0.886029in}{1.514658in}}%
\pgfpathlineto{\pgfqpoint{0.885833in}{1.526190in}}%
\pgfpathlineto{\pgfqpoint{0.873499in}{1.534484in}}%
\pgfpathlineto{\pgfqpoint{0.875510in}{1.546436in}}%
\pgfpathlineto{\pgfqpoint{0.866728in}{1.566180in}}%
\pgfpathlineto{\pgfqpoint{0.867319in}{1.582964in}}%
\pgfpathlineto{\pgfqpoint{0.859772in}{1.591178in}}%
\pgfpathlineto{\pgfqpoint{0.852690in}{1.583911in}}%
\pgfpathlineto{\pgfqpoint{0.843050in}{1.579980in}}%
\pgfpathlineto{\pgfqpoint{0.835712in}{1.588082in}}%
\pgfpathlineto{\pgfqpoint{0.836905in}{1.601019in}}%
\pgfpathlineto{\pgfqpoint{0.845705in}{1.606116in}}%
\pgfpathlineto{\pgfqpoint{0.843366in}{1.614312in}}%
\pgfpathlineto{\pgfqpoint{0.858663in}{1.653985in}}%
\pgfpathlineto{\pgfqpoint{0.846601in}{1.654979in}}%
\pgfpathlineto{\pgfqpoint{0.845361in}{1.662098in}}%
\pgfpathlineto{\pgfqpoint{0.836575in}{1.668238in}}%
\pgfpathlineto{\pgfqpoint{0.837137in}{1.675101in}}%
\pgfpathlineto{\pgfqpoint{0.832502in}{1.680433in}}%
\pgfpathlineto{\pgfqpoint{0.824957in}{1.699791in}}%
\pgfpathlineto{\pgfqpoint{0.818063in}{1.703739in}}%
\pgfpathlineto{\pgfqpoint{0.813126in}{1.720446in}}%
\pgfpathlineto{\pgfqpoint{0.814291in}{1.731551in}}%
\pgfpathlineto{\pgfqpoint{0.805080in}{1.752004in}}%
\pgfpathlineto{\pgfqpoint{0.819624in}{1.815480in}}%
\pgfpathclose%
\pgfusepath{fill}%
\end{pgfscope}%
\begin{pgfscope}%
\pgfpathrectangle{\pgfqpoint{0.100000in}{0.100000in}}{\pgfqpoint{2.862985in}{1.832778in}}%
\pgfusepath{clip}%
\pgfsetbuttcap%
\pgfsetmiterjoin%
\definecolor{currentfill}{rgb}{0.805467,0.289057,0.462415}%
\pgfsetfillcolor{currentfill}%
\pgfsetlinewidth{0.000000pt}%
\definecolor{currentstroke}{rgb}{0.000000,0.000000,0.000000}%
\pgfsetstrokecolor{currentstroke}%
\pgfsetstrokeopacity{0.000000}%
\pgfsetdash{}{0pt}%
\pgfpathmoveto{\pgfqpoint{2.788801in}{1.489012in}}%
\pgfpathlineto{\pgfqpoint{2.787168in}{1.495954in}}%
\pgfpathlineto{\pgfqpoint{2.778094in}{1.502035in}}%
\pgfpathlineto{\pgfqpoint{2.754117in}{1.580468in}}%
\pgfpathlineto{\pgfqpoint{2.741192in}{1.618318in}}%
\pgfpathlineto{\pgfqpoint{2.751583in}{1.629222in}}%
\pgfpathlineto{\pgfqpoint{2.761683in}{1.655908in}}%
\pgfpathlineto{\pgfqpoint{2.766720in}{1.664437in}}%
\pgfpathlineto{\pgfqpoint{2.763152in}{1.668018in}}%
\pgfpathlineto{\pgfqpoint{2.761957in}{1.691707in}}%
\pgfpathlineto{\pgfqpoint{2.766463in}{1.698969in}}%
\pgfpathlineto{\pgfqpoint{2.764488in}{1.715849in}}%
\pgfpathlineto{\pgfqpoint{2.782474in}{1.771228in}}%
\pgfpathlineto{\pgfqpoint{2.790578in}{1.771530in}}%
\pgfpathlineto{\pgfqpoint{2.793942in}{1.761638in}}%
\pgfpathlineto{\pgfqpoint{2.801144in}{1.758798in}}%
\pgfpathlineto{\pgfqpoint{2.814749in}{1.770433in}}%
\pgfpathlineto{\pgfqpoint{2.825417in}{1.777306in}}%
\pgfpathlineto{\pgfqpoint{2.848922in}{1.765204in}}%
\pgfpathlineto{\pgfqpoint{2.870157in}{1.698008in}}%
\pgfpathlineto{\pgfqpoint{2.874194in}{1.681476in}}%
\pgfpathlineto{\pgfqpoint{2.883495in}{1.679537in}}%
\pgfpathlineto{\pgfqpoint{2.896110in}{1.668309in}}%
\pgfpathlineto{\pgfqpoint{2.895396in}{1.661767in}}%
\pgfpathlineto{\pgfqpoint{2.903989in}{1.654013in}}%
\pgfpathlineto{\pgfqpoint{2.910970in}{1.658462in}}%
\pgfpathlineto{\pgfqpoint{2.925577in}{1.641426in}}%
\pgfpathlineto{\pgfqpoint{2.919049in}{1.627792in}}%
\pgfpathlineto{\pgfqpoint{2.910294in}{1.627522in}}%
\pgfpathlineto{\pgfqpoint{2.903271in}{1.615351in}}%
\pgfpathlineto{\pgfqpoint{2.894762in}{1.613625in}}%
\pgfpathlineto{\pgfqpoint{2.888516in}{1.607145in}}%
\pgfpathlineto{\pgfqpoint{2.869860in}{1.600195in}}%
\pgfpathlineto{\pgfqpoint{2.858568in}{1.588995in}}%
\pgfpathlineto{\pgfqpoint{2.852832in}{1.597034in}}%
\pgfpathlineto{\pgfqpoint{2.847620in}{1.591272in}}%
\pgfpathlineto{\pgfqpoint{2.849047in}{1.567979in}}%
\pgfpathlineto{\pgfqpoint{2.844920in}{1.558753in}}%
\pgfpathlineto{\pgfqpoint{2.835918in}{1.561273in}}%
\pgfpathlineto{\pgfqpoint{2.834465in}{1.551814in}}%
\pgfpathlineto{\pgfqpoint{2.830576in}{1.547924in}}%
\pgfpathlineto{\pgfqpoint{2.822791in}{1.548872in}}%
\pgfpathlineto{\pgfqpoint{2.823281in}{1.540028in}}%
\pgfpathlineto{\pgfqpoint{2.811482in}{1.542540in}}%
\pgfpathlineto{\pgfqpoint{2.805038in}{1.530287in}}%
\pgfpathlineto{\pgfqpoint{2.807472in}{1.523875in}}%
\pgfpathlineto{\pgfqpoint{2.803639in}{1.513212in}}%
\pgfpathlineto{\pgfqpoint{2.797604in}{1.505369in}}%
\pgfpathlineto{\pgfqpoint{2.796086in}{1.488995in}}%
\pgfpathlineto{\pgfqpoint{2.788801in}{1.489012in}}%
\pgfpathclose%
\pgfusepath{fill}%
\end{pgfscope}%
\begin{pgfscope}%
\pgfpathrectangle{\pgfqpoint{0.100000in}{0.100000in}}{\pgfqpoint{2.862985in}{1.832778in}}%
\pgfusepath{clip}%
\pgfsetbuttcap%
\pgfsetmiterjoin%
\definecolor{currentfill}{rgb}{0.805467,0.289057,0.462415}%
\pgfsetfillcolor{currentfill}%
\pgfsetlinewidth{0.000000pt}%
\definecolor{currentstroke}{rgb}{0.000000,0.000000,0.000000}%
\pgfsetstrokecolor{currentstroke}%
\pgfsetstrokeopacity{0.000000}%
\pgfsetdash{}{0pt}%
\pgfpathmoveto{\pgfqpoint{2.873211in}{1.595412in}}%
\pgfpathlineto{\pgfqpoint{2.878538in}{1.600991in}}%
\pgfpathlineto{\pgfqpoint{2.883609in}{1.595747in}}%
\pgfpathlineto{\pgfqpoint{2.874500in}{1.588798in}}%
\pgfpathlineto{\pgfqpoint{2.873211in}{1.595412in}}%
\pgfpathclose%
\pgfusepath{fill}%
\end{pgfscope}%
\begin{pgfscope}%
\pgfpathrectangle{\pgfqpoint{0.100000in}{0.100000in}}{\pgfqpoint{2.862985in}{1.832778in}}%
\pgfusepath{clip}%
\pgfsetbuttcap%
\pgfsetmiterjoin%
\definecolor{currentfill}{rgb}{0.423689,0.000646,0.658956}%
\pgfsetfillcolor{currentfill}%
\pgfsetlinewidth{0.000000pt}%
\definecolor{currentstroke}{rgb}{0.000000,0.000000,0.000000}%
\pgfsetstrokecolor{currentstroke}%
\pgfsetstrokeopacity{0.000000}%
\pgfsetdash{}{0pt}%
\pgfpathmoveto{\pgfqpoint{1.295312in}{1.539783in}}%
\pgfpathlineto{\pgfqpoint{1.299775in}{1.583815in}}%
\pgfpathlineto{\pgfqpoint{1.308284in}{1.669428in}}%
\pgfpathlineto{\pgfqpoint{1.314001in}{1.732781in}}%
\pgfpathlineto{\pgfqpoint{1.360821in}{1.728389in}}%
\pgfpathlineto{\pgfqpoint{1.420718in}{1.723620in}}%
\pgfpathlineto{\pgfqpoint{1.475466in}{1.720093in}}%
\pgfpathlineto{\pgfqpoint{1.525039in}{1.717585in}}%
\pgfpathlineto{\pgfqpoint{1.599019in}{1.715047in}}%
\pgfpathlineto{\pgfqpoint{1.604068in}{1.694197in}}%
\pgfpathlineto{\pgfqpoint{1.602261in}{1.684232in}}%
\pgfpathlineto{\pgfqpoint{1.601666in}{1.663751in}}%
\pgfpathlineto{\pgfqpoint{1.605101in}{1.648413in}}%
\pgfpathlineto{\pgfqpoint{1.612970in}{1.625784in}}%
\pgfpathlineto{\pgfqpoint{1.612952in}{1.597294in}}%
\pgfpathlineto{\pgfqpoint{1.614408in}{1.564197in}}%
\pgfpathlineto{\pgfqpoint{1.616408in}{1.555277in}}%
\pgfpathlineto{\pgfqpoint{1.622257in}{1.545485in}}%
\pgfpathlineto{\pgfqpoint{1.624186in}{1.530226in}}%
\pgfpathlineto{\pgfqpoint{1.623353in}{1.520039in}}%
\pgfpathlineto{\pgfqpoint{1.561333in}{1.521383in}}%
\pgfpathlineto{\pgfqpoint{1.516215in}{1.523670in}}%
\pgfpathlineto{\pgfqpoint{1.450074in}{1.527191in}}%
\pgfpathlineto{\pgfqpoint{1.384849in}{1.531824in}}%
\pgfpathlineto{\pgfqpoint{1.341407in}{1.535329in}}%
\pgfpathlineto{\pgfqpoint{1.295312in}{1.539783in}}%
\pgfpathclose%
\pgfusepath{fill}%
\end{pgfscope}%
\begin{pgfscope}%
\pgfpathrectangle{\pgfqpoint{0.100000in}{0.100000in}}{\pgfqpoint{2.862985in}{1.832778in}}%
\pgfusepath{clip}%
\pgfsetbuttcap%
\pgfsetmiterjoin%
\definecolor{currentfill}{rgb}{0.471457,0.005678,0.659897}%
\pgfsetfillcolor{currentfill}%
\pgfsetlinewidth{0.000000pt}%
\definecolor{currentstroke}{rgb}{0.000000,0.000000,0.000000}%
\pgfsetstrokecolor{currentstroke}%
\pgfsetstrokeopacity{0.000000}%
\pgfsetdash{}{0pt}%
\pgfpathmoveto{\pgfqpoint{1.276565in}{1.355506in}}%
\pgfpathlineto{\pgfqpoint{1.283725in}{1.429317in}}%
\pgfpathlineto{\pgfqpoint{1.288893in}{1.480441in}}%
\pgfpathlineto{\pgfqpoint{1.295312in}{1.539783in}}%
\pgfpathlineto{\pgfqpoint{1.341407in}{1.535329in}}%
\pgfpathlineto{\pgfqpoint{1.384849in}{1.531824in}}%
\pgfpathlineto{\pgfqpoint{1.450074in}{1.527191in}}%
\pgfpathlineto{\pgfqpoint{1.516215in}{1.523670in}}%
\pgfpathlineto{\pgfqpoint{1.561333in}{1.521383in}}%
\pgfpathlineto{\pgfqpoint{1.623353in}{1.520039in}}%
\pgfpathlineto{\pgfqpoint{1.619154in}{1.507786in}}%
\pgfpathlineto{\pgfqpoint{1.610772in}{1.498167in}}%
\pgfpathlineto{\pgfqpoint{1.617191in}{1.487092in}}%
\pgfpathlineto{\pgfqpoint{1.624268in}{1.484725in}}%
\pgfpathlineto{\pgfqpoint{1.627626in}{1.478363in}}%
\pgfpathlineto{\pgfqpoint{1.626855in}{1.431914in}}%
\pgfpathlineto{\pgfqpoint{1.625565in}{1.366544in}}%
\pgfpathlineto{\pgfqpoint{1.620791in}{1.349375in}}%
\pgfpathlineto{\pgfqpoint{1.625061in}{1.339878in}}%
\pgfpathlineto{\pgfqpoint{1.616489in}{1.319470in}}%
\pgfpathlineto{\pgfqpoint{1.625519in}{1.303018in}}%
\pgfpathlineto{\pgfqpoint{1.617847in}{1.304277in}}%
\pgfpathlineto{\pgfqpoint{1.612612in}{1.314513in}}%
\pgfpathlineto{\pgfqpoint{1.593930in}{1.321529in}}%
\pgfpathlineto{\pgfqpoint{1.582148in}{1.327701in}}%
\pgfpathlineto{\pgfqpoint{1.562377in}{1.328214in}}%
\pgfpathlineto{\pgfqpoint{1.555515in}{1.322591in}}%
\pgfpathlineto{\pgfqpoint{1.533137in}{1.333663in}}%
\pgfpathlineto{\pgfqpoint{1.531420in}{1.337164in}}%
\pgfpathlineto{\pgfqpoint{1.453322in}{1.340860in}}%
\pgfpathlineto{\pgfqpoint{1.405874in}{1.343580in}}%
\pgfpathlineto{\pgfqpoint{1.366682in}{1.346649in}}%
\pgfpathlineto{\pgfqpoint{1.301922in}{1.352768in}}%
\pgfpathlineto{\pgfqpoint{1.276565in}{1.355506in}}%
\pgfpathclose%
\pgfusepath{fill}%
\end{pgfscope}%
\begin{pgfscope}%
\pgfpathrectangle{\pgfqpoint{0.100000in}{0.100000in}}{\pgfqpoint{2.862985in}{1.832778in}}%
\pgfusepath{clip}%
\pgfsetbuttcap%
\pgfsetmiterjoin%
\definecolor{currentfill}{rgb}{0.645872,0.120898,0.598867}%
\pgfsetfillcolor{currentfill}%
\pgfsetlinewidth{0.000000pt}%
\definecolor{currentstroke}{rgb}{0.000000,0.000000,0.000000}%
\pgfsetstrokecolor{currentstroke}%
\pgfsetstrokeopacity{0.000000}%
\pgfsetdash{}{0pt}%
\pgfpathmoveto{\pgfqpoint{1.264283in}{1.230524in}}%
\pgfpathlineto{\pgfqpoint{1.222647in}{1.234326in}}%
\pgfpathlineto{\pgfqpoint{1.131892in}{1.245527in}}%
\pgfpathlineto{\pgfqpoint{1.082522in}{1.252613in}}%
\pgfpathlineto{\pgfqpoint{1.029571in}{1.260214in}}%
\pgfpathlineto{\pgfqpoint{0.984967in}{1.267352in}}%
\pgfpathlineto{\pgfqpoint{0.936009in}{1.275734in}}%
\pgfpathlineto{\pgfqpoint{0.947141in}{1.337483in}}%
\pgfpathlineto{\pgfqpoint{0.958416in}{1.400800in}}%
\pgfpathlineto{\pgfqpoint{0.974380in}{1.490854in}}%
\pgfpathlineto{\pgfqpoint{0.980083in}{1.523280in}}%
\pgfpathlineto{\pgfqpoint{1.007445in}{1.518274in}}%
\pgfpathlineto{\pgfqpoint{1.086931in}{1.505797in}}%
\pgfpathlineto{\pgfqpoint{1.119194in}{1.501195in}}%
\pgfpathlineto{\pgfqpoint{1.191429in}{1.491308in}}%
\pgfpathlineto{\pgfqpoint{1.243453in}{1.485403in}}%
\pgfpathlineto{\pgfqpoint{1.288893in}{1.480441in}}%
\pgfpathlineto{\pgfqpoint{1.283725in}{1.429317in}}%
\pgfpathlineto{\pgfqpoint{1.276565in}{1.355506in}}%
\pgfpathlineto{\pgfqpoint{1.270419in}{1.292777in}}%
\pgfpathlineto{\pgfqpoint{1.264283in}{1.230524in}}%
\pgfpathclose%
\pgfusepath{fill}%
\end{pgfscope}%
\begin{pgfscope}%
\pgfpathrectangle{\pgfqpoint{0.100000in}{0.100000in}}{\pgfqpoint{2.862985in}{1.832778in}}%
\pgfusepath{clip}%
\pgfsetbuttcap%
\pgfsetmiterjoin%
\definecolor{currentfill}{rgb}{0.573632,0.060028,0.637349}%
\pgfsetfillcolor{currentfill}%
\pgfsetlinewidth{0.000000pt}%
\definecolor{currentstroke}{rgb}{0.000000,0.000000,0.000000}%
\pgfsetstrokecolor{currentstroke}%
\pgfsetstrokeopacity{0.000000}%
\pgfsetdash{}{0pt}%
\pgfpathmoveto{\pgfqpoint{2.025520in}{1.315700in}}%
\pgfpathlineto{\pgfqpoint{1.972758in}{1.311952in}}%
\pgfpathlineto{\pgfqpoint{1.894114in}{1.308600in}}%
\pgfpathlineto{\pgfqpoint{1.891125in}{1.316544in}}%
\pgfpathlineto{\pgfqpoint{1.873686in}{1.322483in}}%
\pgfpathlineto{\pgfqpoint{1.869839in}{1.333702in}}%
\pgfpathlineto{\pgfqpoint{1.868228in}{1.347582in}}%
\pgfpathlineto{\pgfqpoint{1.872158in}{1.354693in}}%
\pgfpathlineto{\pgfqpoint{1.865951in}{1.361517in}}%
\pgfpathlineto{\pgfqpoint{1.864454in}{1.369657in}}%
\pgfpathlineto{\pgfqpoint{1.862453in}{1.387663in}}%
\pgfpathlineto{\pgfqpoint{1.856505in}{1.397444in}}%
\pgfpathlineto{\pgfqpoint{1.845940in}{1.402933in}}%
\pgfpathlineto{\pgfqpoint{1.834444in}{1.411992in}}%
\pgfpathlineto{\pgfqpoint{1.828501in}{1.422710in}}%
\pgfpathlineto{\pgfqpoint{1.817838in}{1.427026in}}%
\pgfpathlineto{\pgfqpoint{1.811570in}{1.434051in}}%
\pgfpathlineto{\pgfqpoint{1.803976in}{1.435243in}}%
\pgfpathlineto{\pgfqpoint{1.790419in}{1.445673in}}%
\pgfpathlineto{\pgfqpoint{1.792621in}{1.457675in}}%
\pgfpathlineto{\pgfqpoint{1.792197in}{1.480500in}}%
\pgfpathlineto{\pgfqpoint{1.796378in}{1.486783in}}%
\pgfpathlineto{\pgfqpoint{1.792621in}{1.496270in}}%
\pgfpathlineto{\pgfqpoint{1.786004in}{1.498105in}}%
\pgfpathlineto{\pgfqpoint{1.786550in}{1.506437in}}%
\pgfpathlineto{\pgfqpoint{1.794761in}{1.519600in}}%
\pgfpathlineto{\pgfqpoint{1.811020in}{1.530009in}}%
\pgfpathlineto{\pgfqpoint{1.810007in}{1.567025in}}%
\pgfpathlineto{\pgfqpoint{1.818143in}{1.572586in}}%
\pgfpathlineto{\pgfqpoint{1.825842in}{1.568882in}}%
\pgfpathlineto{\pgfqpoint{1.841526in}{1.574304in}}%
\pgfpathlineto{\pgfqpoint{1.871034in}{1.587932in}}%
\pgfpathlineto{\pgfqpoint{1.874878in}{1.583713in}}%
\pgfpathlineto{\pgfqpoint{1.869289in}{1.564600in}}%
\pgfpathlineto{\pgfqpoint{1.877604in}{1.568794in}}%
\pgfpathlineto{\pgfqpoint{1.891845in}{1.564606in}}%
\pgfpathlineto{\pgfqpoint{1.900596in}{1.561110in}}%
\pgfpathlineto{\pgfqpoint{1.905502in}{1.550866in}}%
\pgfpathlineto{\pgfqpoint{1.950397in}{1.541198in}}%
\pgfpathlineto{\pgfqpoint{1.963816in}{1.534565in}}%
\pgfpathlineto{\pgfqpoint{1.977460in}{1.534643in}}%
\pgfpathlineto{\pgfqpoint{1.991480in}{1.531907in}}%
\pgfpathlineto{\pgfqpoint{2.000578in}{1.522561in}}%
\pgfpathlineto{\pgfqpoint{2.009272in}{1.517943in}}%
\pgfpathlineto{\pgfqpoint{2.010870in}{1.504623in}}%
\pgfpathlineto{\pgfqpoint{2.008301in}{1.496259in}}%
\pgfpathlineto{\pgfqpoint{2.017991in}{1.495641in}}%
\pgfpathlineto{\pgfqpoint{2.014727in}{1.485941in}}%
\pgfpathlineto{\pgfqpoint{2.017832in}{1.482495in}}%
\pgfpathlineto{\pgfqpoint{2.020919in}{1.473339in}}%
\pgfpathlineto{\pgfqpoint{2.011476in}{1.468491in}}%
\pgfpathlineto{\pgfqpoint{2.005993in}{1.454984in}}%
\pgfpathlineto{\pgfqpoint{2.004273in}{1.445433in}}%
\pgfpathlineto{\pgfqpoint{2.009533in}{1.443798in}}%
\pgfpathlineto{\pgfqpoint{2.016268in}{1.450960in}}%
\pgfpathlineto{\pgfqpoint{2.021994in}{1.463372in}}%
\pgfpathlineto{\pgfqpoint{2.029741in}{1.467683in}}%
\pgfpathlineto{\pgfqpoint{2.035563in}{1.462075in}}%
\pgfpathlineto{\pgfqpoint{2.029832in}{1.445104in}}%
\pgfpathlineto{\pgfqpoint{2.028034in}{1.431928in}}%
\pgfpathlineto{\pgfqpoint{2.029728in}{1.422457in}}%
\pgfpathlineto{\pgfqpoint{2.024406in}{1.417099in}}%
\pgfpathlineto{\pgfqpoint{2.021748in}{1.404007in}}%
\pgfpathlineto{\pgfqpoint{2.023920in}{1.390247in}}%
\pgfpathlineto{\pgfqpoint{2.017640in}{1.369897in}}%
\pgfpathlineto{\pgfqpoint{2.017772in}{1.359728in}}%
\pgfpathlineto{\pgfqpoint{2.022734in}{1.337693in}}%
\pgfpathlineto{\pgfqpoint{2.025951in}{1.333913in}}%
\pgfpathlineto{\pgfqpoint{2.025520in}{1.315700in}}%
\pgfpathclose%
\pgfusepath{fill}%
\end{pgfscope}%
\begin{pgfscope}%
\pgfpathrectangle{\pgfqpoint{0.100000in}{0.100000in}}{\pgfqpoint{2.862985in}{1.832778in}}%
\pgfusepath{clip}%
\pgfsetbuttcap%
\pgfsetmiterjoin%
\definecolor{currentfill}{rgb}{0.573632,0.060028,0.637349}%
\pgfsetfillcolor{currentfill}%
\pgfsetlinewidth{0.000000pt}%
\definecolor{currentstroke}{rgb}{0.000000,0.000000,0.000000}%
\pgfsetstrokecolor{currentstroke}%
\pgfsetstrokeopacity{0.000000}%
\pgfsetdash{}{0pt}%
\pgfpathmoveto{\pgfqpoint{2.045302in}{1.494291in}}%
\pgfpathlineto{\pgfqpoint{2.046480in}{1.485502in}}%
\pgfpathlineto{\pgfqpoint{2.035690in}{1.462343in}}%
\pgfpathlineto{\pgfqpoint{2.030896in}{1.469051in}}%
\pgfpathlineto{\pgfqpoint{2.045302in}{1.494291in}}%
\pgfpathclose%
\pgfusepath{fill}%
\end{pgfscope}%
\begin{pgfscope}%
\pgfpathrectangle{\pgfqpoint{0.100000in}{0.100000in}}{\pgfqpoint{2.862985in}{1.832778in}}%
\pgfusepath{clip}%
\pgfsetbuttcap%
\pgfsetmiterjoin%
\definecolor{currentfill}{rgb}{0.697324,0.169573,0.560919}%
\pgfsetfillcolor{currentfill}%
\pgfsetlinewidth{0.000000pt}%
\definecolor{currentstroke}{rgb}{0.000000,0.000000,0.000000}%
\pgfsetstrokecolor{currentstroke}%
\pgfsetstrokeopacity{0.000000}%
\pgfsetdash{}{0pt}%
\pgfpathmoveto{\pgfqpoint{0.739329in}{1.638779in}}%
\pgfpathlineto{\pgfqpoint{0.741462in}{1.649338in}}%
\pgfpathlineto{\pgfqpoint{0.739852in}{1.660801in}}%
\pgfpathlineto{\pgfqpoint{0.742317in}{1.673299in}}%
\pgfpathlineto{\pgfqpoint{0.754869in}{1.724007in}}%
\pgfpathlineto{\pgfqpoint{0.765426in}{1.766159in}}%
\pgfpathlineto{\pgfqpoint{0.780273in}{1.824988in}}%
\pgfpathlineto{\pgfqpoint{0.819624in}{1.815480in}}%
\pgfpathlineto{\pgfqpoint{0.805080in}{1.752004in}}%
\pgfpathlineto{\pgfqpoint{0.814291in}{1.731551in}}%
\pgfpathlineto{\pgfqpoint{0.813126in}{1.720446in}}%
\pgfpathlineto{\pgfqpoint{0.818063in}{1.703739in}}%
\pgfpathlineto{\pgfqpoint{0.824957in}{1.699791in}}%
\pgfpathlineto{\pgfqpoint{0.832502in}{1.680433in}}%
\pgfpathlineto{\pgfqpoint{0.837137in}{1.675101in}}%
\pgfpathlineto{\pgfqpoint{0.836575in}{1.668238in}}%
\pgfpathlineto{\pgfqpoint{0.845361in}{1.662098in}}%
\pgfpathlineto{\pgfqpoint{0.846601in}{1.654979in}}%
\pgfpathlineto{\pgfqpoint{0.858663in}{1.653985in}}%
\pgfpathlineto{\pgfqpoint{0.843366in}{1.614312in}}%
\pgfpathlineto{\pgfqpoint{0.845705in}{1.606116in}}%
\pgfpathlineto{\pgfqpoint{0.836905in}{1.601019in}}%
\pgfpathlineto{\pgfqpoint{0.835712in}{1.588082in}}%
\pgfpathlineto{\pgfqpoint{0.843050in}{1.579980in}}%
\pgfpathlineto{\pgfqpoint{0.852690in}{1.583911in}}%
\pgfpathlineto{\pgfqpoint{0.859772in}{1.591178in}}%
\pgfpathlineto{\pgfqpoint{0.867319in}{1.582964in}}%
\pgfpathlineto{\pgfqpoint{0.866728in}{1.566180in}}%
\pgfpathlineto{\pgfqpoint{0.875510in}{1.546436in}}%
\pgfpathlineto{\pgfqpoint{0.873499in}{1.534484in}}%
\pgfpathlineto{\pgfqpoint{0.885833in}{1.526190in}}%
\pgfpathlineto{\pgfqpoint{0.886029in}{1.514658in}}%
\pgfpathlineto{\pgfqpoint{0.890230in}{1.502538in}}%
\pgfpathlineto{\pgfqpoint{0.898038in}{1.505269in}}%
\pgfpathlineto{\pgfqpoint{0.915931in}{1.500183in}}%
\pgfpathlineto{\pgfqpoint{0.923297in}{1.505835in}}%
\pgfpathlineto{\pgfqpoint{0.928566in}{1.502181in}}%
\pgfpathlineto{\pgfqpoint{0.939045in}{1.502638in}}%
\pgfpathlineto{\pgfqpoint{0.941293in}{1.498866in}}%
\pgfpathlineto{\pgfqpoint{0.952480in}{1.499920in}}%
\pgfpathlineto{\pgfqpoint{0.957871in}{1.507806in}}%
\pgfpathlineto{\pgfqpoint{0.965403in}{1.508767in}}%
\pgfpathlineto{\pgfqpoint{0.971262in}{1.493168in}}%
\pgfpathlineto{\pgfqpoint{0.974380in}{1.490854in}}%
\pgfpathlineto{\pgfqpoint{0.958416in}{1.400800in}}%
\pgfpathlineto{\pgfqpoint{0.947141in}{1.337483in}}%
\pgfpathlineto{\pgfqpoint{0.858212in}{1.354650in}}%
\pgfpathlineto{\pgfqpoint{0.810177in}{1.364205in}}%
\pgfpathlineto{\pgfqpoint{0.765246in}{1.374085in}}%
\pgfpathlineto{\pgfqpoint{0.723250in}{1.383593in}}%
\pgfpathlineto{\pgfqpoint{0.674647in}{1.395314in}}%
\pgfpathlineto{\pgfqpoint{0.699717in}{1.498158in}}%
\pgfpathlineto{\pgfqpoint{0.701069in}{1.505674in}}%
\pgfpathlineto{\pgfqpoint{0.712217in}{1.525368in}}%
\pgfpathlineto{\pgfqpoint{0.700673in}{1.537205in}}%
\pgfpathlineto{\pgfqpoint{0.703070in}{1.548828in}}%
\pgfpathlineto{\pgfqpoint{0.707464in}{1.551738in}}%
\pgfpathlineto{\pgfqpoint{0.715313in}{1.563715in}}%
\pgfpathlineto{\pgfqpoint{0.726752in}{1.572002in}}%
\pgfpathlineto{\pgfqpoint{0.727377in}{1.578133in}}%
\pgfpathlineto{\pgfqpoint{0.734282in}{1.584259in}}%
\pgfpathlineto{\pgfqpoint{0.740015in}{1.595725in}}%
\pgfpathlineto{\pgfqpoint{0.751774in}{1.607856in}}%
\pgfpathlineto{\pgfqpoint{0.750939in}{1.619838in}}%
\pgfpathlineto{\pgfqpoint{0.742899in}{1.626494in}}%
\pgfpathlineto{\pgfqpoint{0.739329in}{1.638779in}}%
\pgfpathclose%
\pgfusepath{fill}%
\end{pgfscope}%
\begin{pgfscope}%
\pgfpathrectangle{\pgfqpoint{0.100000in}{0.100000in}}{\pgfqpoint{2.862985in}{1.832778in}}%
\pgfusepath{clip}%
\pgfsetbuttcap%
\pgfsetmiterjoin%
\definecolor{currentfill}{rgb}{0.600266,0.081516,0.625342}%
\pgfsetfillcolor{currentfill}%
\pgfsetlinewidth{0.000000pt}%
\definecolor{currentstroke}{rgb}{0.000000,0.000000,0.000000}%
\pgfsetstrokecolor{currentstroke}%
\pgfsetstrokeopacity{0.000000}%
\pgfsetdash{}{0pt}%
\pgfpathmoveto{\pgfqpoint{2.686015in}{1.438080in}}%
\pgfpathlineto{\pgfqpoint{2.682320in}{1.449744in}}%
\pgfpathlineto{\pgfqpoint{2.675463in}{1.485127in}}%
\pgfpathlineto{\pgfqpoint{2.666598in}{1.496013in}}%
\pgfpathlineto{\pgfqpoint{2.658822in}{1.515590in}}%
\pgfpathlineto{\pgfqpoint{2.661375in}{1.530104in}}%
\pgfpathlineto{\pgfqpoint{2.659327in}{1.540801in}}%
\pgfpathlineto{\pgfqpoint{2.652683in}{1.551308in}}%
\pgfpathlineto{\pgfqpoint{2.652410in}{1.562882in}}%
\pgfpathlineto{\pgfqpoint{2.648552in}{1.575365in}}%
\pgfpathlineto{\pgfqpoint{2.683077in}{1.583862in}}%
\pgfpathlineto{\pgfqpoint{2.727922in}{1.595966in}}%
\pgfpathlineto{\pgfqpoint{2.729711in}{1.589025in}}%
\pgfpathlineto{\pgfqpoint{2.726856in}{1.577975in}}%
\pgfpathlineto{\pgfqpoint{2.733559in}{1.569133in}}%
\pgfpathlineto{\pgfqpoint{2.729999in}{1.557930in}}%
\pgfpathlineto{\pgfqpoint{2.715852in}{1.543846in}}%
\pgfpathlineto{\pgfqpoint{2.719741in}{1.533267in}}%
\pgfpathlineto{\pgfqpoint{2.716859in}{1.510906in}}%
\pgfpathlineto{\pgfqpoint{2.712824in}{1.498298in}}%
\pgfpathlineto{\pgfqpoint{2.717972in}{1.462444in}}%
\pgfpathlineto{\pgfqpoint{2.717278in}{1.449830in}}%
\pgfpathlineto{\pgfqpoint{2.722258in}{1.445782in}}%
\pgfpathlineto{\pgfqpoint{2.686015in}{1.438080in}}%
\pgfpathclose%
\pgfusepath{fill}%
\end{pgfscope}%
\begin{pgfscope}%
\pgfpathrectangle{\pgfqpoint{0.100000in}{0.100000in}}{\pgfqpoint{2.862985in}{1.832778in}}%
\pgfusepath{clip}%
\pgfsetbuttcap%
\pgfsetmiterjoin%
\definecolor{currentfill}{rgb}{0.471457,0.005678,0.659897}%
\pgfsetfillcolor{currentfill}%
\pgfsetlinewidth{0.000000pt}%
\definecolor{currentstroke}{rgb}{0.000000,0.000000,0.000000}%
\pgfsetstrokecolor{currentstroke}%
\pgfsetstrokeopacity{0.000000}%
\pgfsetdash{}{0pt}%
\pgfpathmoveto{\pgfqpoint{1.625565in}{1.366544in}}%
\pgfpathlineto{\pgfqpoint{1.626855in}{1.431914in}}%
\pgfpathlineto{\pgfqpoint{1.627626in}{1.478363in}}%
\pgfpathlineto{\pgfqpoint{1.624268in}{1.484725in}}%
\pgfpathlineto{\pgfqpoint{1.617191in}{1.487092in}}%
\pgfpathlineto{\pgfqpoint{1.610772in}{1.498167in}}%
\pgfpathlineto{\pgfqpoint{1.619154in}{1.507786in}}%
\pgfpathlineto{\pgfqpoint{1.623353in}{1.520039in}}%
\pgfpathlineto{\pgfqpoint{1.624186in}{1.530226in}}%
\pgfpathlineto{\pgfqpoint{1.622257in}{1.545485in}}%
\pgfpathlineto{\pgfqpoint{1.616408in}{1.555277in}}%
\pgfpathlineto{\pgfqpoint{1.614408in}{1.564197in}}%
\pgfpathlineto{\pgfqpoint{1.612952in}{1.597294in}}%
\pgfpathlineto{\pgfqpoint{1.612970in}{1.625784in}}%
\pgfpathlineto{\pgfqpoint{1.605101in}{1.648413in}}%
\pgfpathlineto{\pgfqpoint{1.601666in}{1.663751in}}%
\pgfpathlineto{\pgfqpoint{1.602261in}{1.684232in}}%
\pgfpathlineto{\pgfqpoint{1.604068in}{1.694197in}}%
\pgfpathlineto{\pgfqpoint{1.599019in}{1.715047in}}%
\pgfpathlineto{\pgfqpoint{1.633393in}{1.714359in}}%
\pgfpathlineto{\pgfqpoint{1.685604in}{1.713911in}}%
\pgfpathlineto{\pgfqpoint{1.685890in}{1.737609in}}%
\pgfpathlineto{\pgfqpoint{1.699180in}{1.735000in}}%
\pgfpathlineto{\pgfqpoint{1.705549in}{1.706104in}}%
\pgfpathlineto{\pgfqpoint{1.710245in}{1.695714in}}%
\pgfpathlineto{\pgfqpoint{1.721922in}{1.695407in}}%
\pgfpathlineto{\pgfqpoint{1.724536in}{1.691882in}}%
\pgfpathlineto{\pgfqpoint{1.740825in}{1.690320in}}%
\pgfpathlineto{\pgfqpoint{1.743562in}{1.683156in}}%
\pgfpathlineto{\pgfqpoint{1.754788in}{1.684767in}}%
\pgfpathlineto{\pgfqpoint{1.763507in}{1.691467in}}%
\pgfpathlineto{\pgfqpoint{1.778542in}{1.691216in}}%
\pgfpathlineto{\pgfqpoint{1.787845in}{1.685828in}}%
\pgfpathlineto{\pgfqpoint{1.788915in}{1.680774in}}%
\pgfpathlineto{\pgfqpoint{1.797767in}{1.679717in}}%
\pgfpathlineto{\pgfqpoint{1.803544in}{1.665932in}}%
\pgfpathlineto{\pgfqpoint{1.807272in}{1.674408in}}%
\pgfpathlineto{\pgfqpoint{1.817441in}{1.674929in}}%
\pgfpathlineto{\pgfqpoint{1.820013in}{1.668328in}}%
\pgfpathlineto{\pgfqpoint{1.831455in}{1.665315in}}%
\pgfpathlineto{\pgfqpoint{1.837765in}{1.655808in}}%
\pgfpathlineto{\pgfqpoint{1.851759in}{1.658760in}}%
\pgfpathlineto{\pgfqpoint{1.867158in}{1.670427in}}%
\pgfpathlineto{\pgfqpoint{1.872772in}{1.660142in}}%
\pgfpathlineto{\pgfqpoint{1.898039in}{1.662957in}}%
\pgfpathlineto{\pgfqpoint{1.908840in}{1.655898in}}%
\pgfpathlineto{\pgfqpoint{1.915126in}{1.658418in}}%
\pgfpathlineto{\pgfqpoint{1.920167in}{1.654444in}}%
\pgfpathlineto{\pgfqpoint{1.905213in}{1.645016in}}%
\pgfpathlineto{\pgfqpoint{1.883866in}{1.636615in}}%
\pgfpathlineto{\pgfqpoint{1.862725in}{1.619824in}}%
\pgfpathlineto{\pgfqpoint{1.844396in}{1.597742in}}%
\pgfpathlineto{\pgfqpoint{1.830521in}{1.584687in}}%
\pgfpathlineto{\pgfqpoint{1.818367in}{1.575713in}}%
\pgfpathlineto{\pgfqpoint{1.810007in}{1.567025in}}%
\pgfpathlineto{\pgfqpoint{1.811020in}{1.530009in}}%
\pgfpathlineto{\pgfqpoint{1.794761in}{1.519600in}}%
\pgfpathlineto{\pgfqpoint{1.786550in}{1.506437in}}%
\pgfpathlineto{\pgfqpoint{1.786004in}{1.498105in}}%
\pgfpathlineto{\pgfqpoint{1.792621in}{1.496270in}}%
\pgfpathlineto{\pgfqpoint{1.796378in}{1.486783in}}%
\pgfpathlineto{\pgfqpoint{1.792197in}{1.480500in}}%
\pgfpathlineto{\pgfqpoint{1.792621in}{1.457675in}}%
\pgfpathlineto{\pgfqpoint{1.790419in}{1.445673in}}%
\pgfpathlineto{\pgfqpoint{1.803976in}{1.435243in}}%
\pgfpathlineto{\pgfqpoint{1.811570in}{1.434051in}}%
\pgfpathlineto{\pgfqpoint{1.817838in}{1.427026in}}%
\pgfpathlineto{\pgfqpoint{1.828501in}{1.422710in}}%
\pgfpathlineto{\pgfqpoint{1.834444in}{1.411992in}}%
\pgfpathlineto{\pgfqpoint{1.845940in}{1.402933in}}%
\pgfpathlineto{\pgfqpoint{1.856505in}{1.397444in}}%
\pgfpathlineto{\pgfqpoint{1.862453in}{1.387663in}}%
\pgfpathlineto{\pgfqpoint{1.864454in}{1.369657in}}%
\pgfpathlineto{\pgfqpoint{1.808388in}{1.367621in}}%
\pgfpathlineto{\pgfqpoint{1.760597in}{1.366622in}}%
\pgfpathlineto{\pgfqpoint{1.717055in}{1.365982in}}%
\pgfpathlineto{\pgfqpoint{1.670992in}{1.366052in}}%
\pgfpathlineto{\pgfqpoint{1.625565in}{1.366544in}}%
\pgfpathclose%
\pgfusepath{fill}%
\end{pgfscope}%
\begin{pgfscope}%
\pgfpathrectangle{\pgfqpoint{0.100000in}{0.100000in}}{\pgfqpoint{2.862985in}{1.832778in}}%
\pgfusepath{clip}%
\pgfsetbuttcap%
\pgfsetmiterjoin%
\definecolor{currentfill}{rgb}{0.727670,0.200586,0.535293}%
\pgfsetfillcolor{currentfill}%
\pgfsetlinewidth{0.000000pt}%
\definecolor{currentstroke}{rgb}{0.000000,0.000000,0.000000}%
\pgfsetstrokecolor{currentstroke}%
\pgfsetstrokeopacity{0.000000}%
\pgfsetdash{}{0pt}%
\pgfpathmoveto{\pgfqpoint{0.352592in}{1.488751in}}%
\pgfpathlineto{\pgfqpoint{0.348152in}{1.496921in}}%
\pgfpathlineto{\pgfqpoint{0.350987in}{1.517893in}}%
\pgfpathlineto{\pgfqpoint{0.356459in}{1.528853in}}%
\pgfpathlineto{\pgfqpoint{0.353724in}{1.543679in}}%
\pgfpathlineto{\pgfqpoint{0.359432in}{1.549914in}}%
\pgfpathlineto{\pgfqpoint{0.369847in}{1.566720in}}%
\pgfpathlineto{\pgfqpoint{0.378677in}{1.576864in}}%
\pgfpathlineto{\pgfqpoint{0.391585in}{1.598959in}}%
\pgfpathlineto{\pgfqpoint{0.401485in}{1.623028in}}%
\pgfpathlineto{\pgfqpoint{0.411995in}{1.645601in}}%
\pgfpathlineto{\pgfqpoint{0.414119in}{1.655029in}}%
\pgfpathlineto{\pgfqpoint{0.428617in}{1.682000in}}%
\pgfpathlineto{\pgfqpoint{0.431409in}{1.693845in}}%
\pgfpathlineto{\pgfqpoint{0.437582in}{1.706284in}}%
\pgfpathlineto{\pgfqpoint{0.439788in}{1.721458in}}%
\pgfpathlineto{\pgfqpoint{0.451573in}{1.728857in}}%
\pgfpathlineto{\pgfqpoint{0.465546in}{1.732587in}}%
\pgfpathlineto{\pgfqpoint{0.472610in}{1.724234in}}%
\pgfpathlineto{\pgfqpoint{0.478718in}{1.724852in}}%
\pgfpathlineto{\pgfqpoint{0.488254in}{1.714897in}}%
\pgfpathlineto{\pgfqpoint{0.489522in}{1.705519in}}%
\pgfpathlineto{\pgfqpoint{0.485905in}{1.687156in}}%
\pgfpathlineto{\pgfqpoint{0.497937in}{1.677787in}}%
\pgfpathlineto{\pgfqpoint{0.505740in}{1.674264in}}%
\pgfpathlineto{\pgfqpoint{0.526880in}{1.677860in}}%
\pgfpathlineto{\pgfqpoint{0.539127in}{1.675404in}}%
\pgfpathlineto{\pgfqpoint{0.547545in}{1.671374in}}%
\pgfpathlineto{\pgfqpoint{0.551979in}{1.663725in}}%
\pgfpathlineto{\pgfqpoint{0.580055in}{1.664746in}}%
\pgfpathlineto{\pgfqpoint{0.584520in}{1.660065in}}%
\pgfpathlineto{\pgfqpoint{0.595053in}{1.659055in}}%
\pgfpathlineto{\pgfqpoint{0.605660in}{1.662045in}}%
\pgfpathlineto{\pgfqpoint{0.637021in}{1.661292in}}%
\pgfpathlineto{\pgfqpoint{0.643794in}{1.658987in}}%
\pgfpathlineto{\pgfqpoint{0.651709in}{1.661643in}}%
\pgfpathlineto{\pgfqpoint{0.739329in}{1.638779in}}%
\pgfpathlineto{\pgfqpoint{0.742899in}{1.626494in}}%
\pgfpathlineto{\pgfqpoint{0.750939in}{1.619838in}}%
\pgfpathlineto{\pgfqpoint{0.751774in}{1.607856in}}%
\pgfpathlineto{\pgfqpoint{0.740015in}{1.595725in}}%
\pgfpathlineto{\pgfqpoint{0.734282in}{1.584259in}}%
\pgfpathlineto{\pgfqpoint{0.727377in}{1.578133in}}%
\pgfpathlineto{\pgfqpoint{0.726752in}{1.572002in}}%
\pgfpathlineto{\pgfqpoint{0.715313in}{1.563715in}}%
\pgfpathlineto{\pgfqpoint{0.707464in}{1.551738in}}%
\pgfpathlineto{\pgfqpoint{0.703070in}{1.548828in}}%
\pgfpathlineto{\pgfqpoint{0.700673in}{1.537205in}}%
\pgfpathlineto{\pgfqpoint{0.712217in}{1.525368in}}%
\pgfpathlineto{\pgfqpoint{0.701069in}{1.505674in}}%
\pgfpathlineto{\pgfqpoint{0.699717in}{1.498158in}}%
\pgfpathlineto{\pgfqpoint{0.674647in}{1.395314in}}%
\pgfpathlineto{\pgfqpoint{0.621916in}{1.408826in}}%
\pgfpathlineto{\pgfqpoint{0.571046in}{1.421943in}}%
\pgfpathlineto{\pgfqpoint{0.540358in}{1.430483in}}%
\pgfpathlineto{\pgfqpoint{0.500927in}{1.441715in}}%
\pgfpathlineto{\pgfqpoint{0.438032in}{1.461511in}}%
\pgfpathlineto{\pgfqpoint{0.369661in}{1.482798in}}%
\pgfpathlineto{\pgfqpoint{0.352592in}{1.488751in}}%
\pgfpathclose%
\pgfusepath{fill}%
\end{pgfscope}%
\begin{pgfscope}%
\pgfpathrectangle{\pgfqpoint{0.100000in}{0.100000in}}{\pgfqpoint{2.862985in}{1.832778in}}%
\pgfusepath{clip}%
\pgfsetbuttcap%
\pgfsetmiterjoin%
\definecolor{currentfill}{rgb}{0.600266,0.081516,0.625342}%
\pgfsetfillcolor{currentfill}%
\pgfsetlinewidth{0.000000pt}%
\definecolor{currentstroke}{rgb}{0.000000,0.000000,0.000000}%
\pgfsetstrokecolor{currentstroke}%
\pgfsetstrokeopacity{0.000000}%
\pgfsetdash{}{0pt}%
\pgfpathmoveto{\pgfqpoint{2.722258in}{1.445782in}}%
\pgfpathlineto{\pgfqpoint{2.717278in}{1.449830in}}%
\pgfpathlineto{\pgfqpoint{2.717972in}{1.462444in}}%
\pgfpathlineto{\pgfqpoint{2.712824in}{1.498298in}}%
\pgfpathlineto{\pgfqpoint{2.716859in}{1.510906in}}%
\pgfpathlineto{\pgfqpoint{2.719741in}{1.533267in}}%
\pgfpathlineto{\pgfqpoint{2.715852in}{1.543846in}}%
\pgfpathlineto{\pgfqpoint{2.729999in}{1.557930in}}%
\pgfpathlineto{\pgfqpoint{2.733559in}{1.569133in}}%
\pgfpathlineto{\pgfqpoint{2.726856in}{1.577975in}}%
\pgfpathlineto{\pgfqpoint{2.729711in}{1.589025in}}%
\pgfpathlineto{\pgfqpoint{2.727922in}{1.595966in}}%
\pgfpathlineto{\pgfqpoint{2.729459in}{1.610832in}}%
\pgfpathlineto{\pgfqpoint{2.732332in}{1.615421in}}%
\pgfpathlineto{\pgfqpoint{2.741192in}{1.618318in}}%
\pgfpathlineto{\pgfqpoint{2.754117in}{1.580468in}}%
\pgfpathlineto{\pgfqpoint{2.778094in}{1.502035in}}%
\pgfpathlineto{\pgfqpoint{2.787168in}{1.495954in}}%
\pgfpathlineto{\pgfqpoint{2.788801in}{1.489012in}}%
\pgfpathlineto{\pgfqpoint{2.793590in}{1.486207in}}%
\pgfpathlineto{\pgfqpoint{2.793226in}{1.473617in}}%
\pgfpathlineto{\pgfqpoint{2.788151in}{1.473415in}}%
\pgfpathlineto{\pgfqpoint{2.777871in}{1.465574in}}%
\pgfpathlineto{\pgfqpoint{2.774906in}{1.457709in}}%
\pgfpathlineto{\pgfqpoint{2.747390in}{1.450901in}}%
\pgfpathlineto{\pgfqpoint{2.722258in}{1.445782in}}%
\pgfpathclose%
\pgfusepath{fill}%
\end{pgfscope}%
\begin{pgfscope}%
\pgfpathrectangle{\pgfqpoint{0.100000in}{0.100000in}}{\pgfqpoint{2.862985in}{1.832778in}}%
\pgfusepath{clip}%
\pgfsetbuttcap%
\pgfsetmiterjoin%
\definecolor{currentfill}{rgb}{0.517933,0.021563,0.654109}%
\pgfsetfillcolor{currentfill}%
\pgfsetlinewidth{0.000000pt}%
\definecolor{currentstroke}{rgb}{0.000000,0.000000,0.000000}%
\pgfsetstrokecolor{currentstroke}%
\pgfsetstrokeopacity{0.000000}%
\pgfsetdash{}{0pt}%
\pgfpathmoveto{\pgfqpoint{1.861865in}{1.172908in}}%
\pgfpathlineto{\pgfqpoint{1.847339in}{1.187293in}}%
\pgfpathlineto{\pgfqpoint{1.800906in}{1.184615in}}%
\pgfpathlineto{\pgfqpoint{1.728489in}{1.182229in}}%
\pgfpathlineto{\pgfqpoint{1.655637in}{1.183365in}}%
\pgfpathlineto{\pgfqpoint{1.650523in}{1.192282in}}%
\pgfpathlineto{\pgfqpoint{1.652611in}{1.201038in}}%
\pgfpathlineto{\pgfqpoint{1.651620in}{1.216037in}}%
\pgfpathlineto{\pgfqpoint{1.647770in}{1.230565in}}%
\pgfpathlineto{\pgfqpoint{1.648217in}{1.238237in}}%
\pgfpathlineto{\pgfqpoint{1.640033in}{1.246905in}}%
\pgfpathlineto{\pgfqpoint{1.641814in}{1.258967in}}%
\pgfpathlineto{\pgfqpoint{1.629260in}{1.282791in}}%
\pgfpathlineto{\pgfqpoint{1.625519in}{1.303018in}}%
\pgfpathlineto{\pgfqpoint{1.616489in}{1.319470in}}%
\pgfpathlineto{\pgfqpoint{1.625061in}{1.339878in}}%
\pgfpathlineto{\pgfqpoint{1.620791in}{1.349375in}}%
\pgfpathlineto{\pgfqpoint{1.625565in}{1.366544in}}%
\pgfpathlineto{\pgfqpoint{1.670992in}{1.366052in}}%
\pgfpathlineto{\pgfqpoint{1.717055in}{1.365982in}}%
\pgfpathlineto{\pgfqpoint{1.760597in}{1.366622in}}%
\pgfpathlineto{\pgfqpoint{1.808388in}{1.367621in}}%
\pgfpathlineto{\pgfqpoint{1.864454in}{1.369657in}}%
\pgfpathlineto{\pgfqpoint{1.865951in}{1.361517in}}%
\pgfpathlineto{\pgfqpoint{1.872158in}{1.354693in}}%
\pgfpathlineto{\pgfqpoint{1.868228in}{1.347582in}}%
\pgfpathlineto{\pgfqpoint{1.869839in}{1.333702in}}%
\pgfpathlineto{\pgfqpoint{1.873686in}{1.322483in}}%
\pgfpathlineto{\pgfqpoint{1.891125in}{1.316544in}}%
\pgfpathlineto{\pgfqpoint{1.894114in}{1.308600in}}%
\pgfpathlineto{\pgfqpoint{1.903684in}{1.299681in}}%
\pgfpathlineto{\pgfqpoint{1.907588in}{1.290454in}}%
\pgfpathlineto{\pgfqpoint{1.917285in}{1.284265in}}%
\pgfpathlineto{\pgfqpoint{1.918802in}{1.276813in}}%
\pgfpathlineto{\pgfqpoint{1.916914in}{1.265532in}}%
\pgfpathlineto{\pgfqpoint{1.911979in}{1.262151in}}%
\pgfpathlineto{\pgfqpoint{1.910490in}{1.251413in}}%
\pgfpathlineto{\pgfqpoint{1.896306in}{1.242884in}}%
\pgfpathlineto{\pgfqpoint{1.877818in}{1.238212in}}%
\pgfpathlineto{\pgfqpoint{1.876124in}{1.227469in}}%
\pgfpathlineto{\pgfqpoint{1.883256in}{1.219794in}}%
\pgfpathlineto{\pgfqpoint{1.883547in}{1.210143in}}%
\pgfpathlineto{\pgfqpoint{1.877793in}{1.202558in}}%
\pgfpathlineto{\pgfqpoint{1.874773in}{1.191288in}}%
\pgfpathlineto{\pgfqpoint{1.864771in}{1.187558in}}%
\pgfpathlineto{\pgfqpoint{1.865408in}{1.174999in}}%
\pgfpathlineto{\pgfqpoint{1.861865in}{1.172908in}}%
\pgfpathclose%
\pgfusepath{fill}%
\end{pgfscope}%
\begin{pgfscope}%
\pgfpathrectangle{\pgfqpoint{0.100000in}{0.100000in}}{\pgfqpoint{2.862985in}{1.832778in}}%
\pgfusepath{clip}%
\pgfsetbuttcap%
\pgfsetmiterjoin%
\definecolor{currentfill}{rgb}{0.610667,0.090204,0.619951}%
\pgfsetfillcolor{currentfill}%
\pgfsetlinewidth{0.000000pt}%
\definecolor{currentstroke}{rgb}{0.000000,0.000000,0.000000}%
\pgfsetstrokecolor{currentstroke}%
\pgfsetstrokeopacity{0.000000}%
\pgfsetdash{}{0pt}%
\pgfpathmoveto{\pgfqpoint{2.763120in}{1.409488in}}%
\pgfpathlineto{\pgfqpoint{2.749362in}{1.407318in}}%
\pgfpathlineto{\pgfqpoint{2.686166in}{1.392957in}}%
\pgfpathlineto{\pgfqpoint{2.685061in}{1.394631in}}%
\pgfpathlineto{\pgfqpoint{2.686015in}{1.438080in}}%
\pgfpathlineto{\pgfqpoint{2.722258in}{1.445782in}}%
\pgfpathlineto{\pgfqpoint{2.747390in}{1.450901in}}%
\pgfpathlineto{\pgfqpoint{2.774906in}{1.457709in}}%
\pgfpathlineto{\pgfqpoint{2.777871in}{1.465574in}}%
\pgfpathlineto{\pgfqpoint{2.788151in}{1.473415in}}%
\pgfpathlineto{\pgfqpoint{2.793226in}{1.473617in}}%
\pgfpathlineto{\pgfqpoint{2.799896in}{1.462179in}}%
\pgfpathlineto{\pgfqpoint{2.793843in}{1.445482in}}%
\pgfpathlineto{\pgfqpoint{2.792955in}{1.435691in}}%
\pgfpathlineto{\pgfqpoint{2.805207in}{1.436611in}}%
\pgfpathlineto{\pgfqpoint{2.810763in}{1.431914in}}%
\pgfpathlineto{\pgfqpoint{2.823224in}{1.412717in}}%
\pgfpathlineto{\pgfqpoint{2.829417in}{1.410382in}}%
\pgfpathlineto{\pgfqpoint{2.839790in}{1.411248in}}%
\pgfpathlineto{\pgfqpoint{2.847009in}{1.417772in}}%
\pgfpathlineto{\pgfqpoint{2.851808in}{1.411950in}}%
\pgfpathlineto{\pgfqpoint{2.821631in}{1.396025in}}%
\pgfpathlineto{\pgfqpoint{2.820688in}{1.407452in}}%
\pgfpathlineto{\pgfqpoint{2.806980in}{1.389693in}}%
\pgfpathlineto{\pgfqpoint{2.802160in}{1.386633in}}%
\pgfpathlineto{\pgfqpoint{2.795438in}{1.396881in}}%
\pgfpathlineto{\pgfqpoint{2.793598in}{1.398285in}}%
\pgfpathlineto{\pgfqpoint{2.787340in}{1.401600in}}%
\pgfpathlineto{\pgfqpoint{2.781882in}{1.415045in}}%
\pgfpathlineto{\pgfqpoint{2.763120in}{1.409488in}}%
\pgfpathclose%
\pgfusepath{fill}%
\end{pgfscope}%
\begin{pgfscope}%
\pgfpathrectangle{\pgfqpoint{0.100000in}{0.100000in}}{\pgfqpoint{2.862985in}{1.832778in}}%
\pgfusepath{clip}%
\pgfsetbuttcap%
\pgfsetmiterjoin%
\definecolor{currentfill}{rgb}{0.417642,0.000564,0.658390}%
\pgfsetfillcolor{currentfill}%
\pgfsetlinewidth{0.000000pt}%
\definecolor{currentstroke}{rgb}{0.000000,0.000000,0.000000}%
\pgfsetstrokecolor{currentstroke}%
\pgfsetstrokeopacity{0.000000}%
\pgfsetdash{}{0pt}%
\pgfpathmoveto{\pgfqpoint{1.353739in}{1.159333in}}%
\pgfpathlineto{\pgfqpoint{1.358777in}{1.221808in}}%
\pgfpathlineto{\pgfqpoint{1.330243in}{1.224124in}}%
\pgfpathlineto{\pgfqpoint{1.264283in}{1.230524in}}%
\pgfpathlineto{\pgfqpoint{1.270419in}{1.292777in}}%
\pgfpathlineto{\pgfqpoint{1.276565in}{1.355506in}}%
\pgfpathlineto{\pgfqpoint{1.301922in}{1.352768in}}%
\pgfpathlineto{\pgfqpoint{1.366682in}{1.346649in}}%
\pgfpathlineto{\pgfqpoint{1.405874in}{1.343580in}}%
\pgfpathlineto{\pgfqpoint{1.453322in}{1.340860in}}%
\pgfpathlineto{\pgfqpoint{1.531420in}{1.337164in}}%
\pgfpathlineto{\pgfqpoint{1.533137in}{1.333663in}}%
\pgfpathlineto{\pgfqpoint{1.555515in}{1.322591in}}%
\pgfpathlineto{\pgfqpoint{1.562377in}{1.328214in}}%
\pgfpathlineto{\pgfqpoint{1.582148in}{1.327701in}}%
\pgfpathlineto{\pgfqpoint{1.593930in}{1.321529in}}%
\pgfpathlineto{\pgfqpoint{1.612612in}{1.314513in}}%
\pgfpathlineto{\pgfqpoint{1.617847in}{1.304277in}}%
\pgfpathlineto{\pgfqpoint{1.625519in}{1.303018in}}%
\pgfpathlineto{\pgfqpoint{1.629260in}{1.282791in}}%
\pgfpathlineto{\pgfqpoint{1.641814in}{1.258967in}}%
\pgfpathlineto{\pgfqpoint{1.640033in}{1.246905in}}%
\pgfpathlineto{\pgfqpoint{1.648217in}{1.238237in}}%
\pgfpathlineto{\pgfqpoint{1.647770in}{1.230565in}}%
\pgfpathlineto{\pgfqpoint{1.651620in}{1.216037in}}%
\pgfpathlineto{\pgfqpoint{1.652611in}{1.201038in}}%
\pgfpathlineto{\pgfqpoint{1.650523in}{1.192282in}}%
\pgfpathlineto{\pgfqpoint{1.655637in}{1.183365in}}%
\pgfpathlineto{\pgfqpoint{1.662651in}{1.167151in}}%
\pgfpathlineto{\pgfqpoint{1.669364in}{1.160551in}}%
\pgfpathlineto{\pgfqpoint{1.677367in}{1.146251in}}%
\pgfpathlineto{\pgfqpoint{1.654686in}{1.146016in}}%
\pgfpathlineto{\pgfqpoint{1.605647in}{1.146774in}}%
\pgfpathlineto{\pgfqpoint{1.551466in}{1.148434in}}%
\pgfpathlineto{\pgfqpoint{1.496966in}{1.150526in}}%
\pgfpathlineto{\pgfqpoint{1.416835in}{1.154903in}}%
\pgfpathlineto{\pgfqpoint{1.353739in}{1.159333in}}%
\pgfpathclose%
\pgfusepath{fill}%
\end{pgfscope}%
\begin{pgfscope}%
\pgfpathrectangle{\pgfqpoint{0.100000in}{0.100000in}}{\pgfqpoint{2.862985in}{1.832778in}}%
\pgfusepath{clip}%
\pgfsetbuttcap%
\pgfsetmiterjoin%
\definecolor{currentfill}{rgb}{0.752312,0.227133,0.513149}%
\pgfsetfillcolor{currentfill}%
\pgfsetlinewidth{0.000000pt}%
\definecolor{currentstroke}{rgb}{0.000000,0.000000,0.000000}%
\pgfsetstrokecolor{currentstroke}%
\pgfsetstrokeopacity{0.000000}%
\pgfsetdash{}{0pt}%
\pgfpathmoveto{\pgfqpoint{2.397084in}{1.347781in}}%
\pgfpathlineto{\pgfqpoint{2.422175in}{1.371705in}}%
\pgfpathlineto{\pgfqpoint{2.425294in}{1.380209in}}%
\pgfpathlineto{\pgfqpoint{2.432646in}{1.387490in}}%
\pgfpathlineto{\pgfqpoint{2.427129in}{1.398090in}}%
\pgfpathlineto{\pgfqpoint{2.420212in}{1.404290in}}%
\pgfpathlineto{\pgfqpoint{2.418213in}{1.415277in}}%
\pgfpathlineto{\pgfqpoint{2.443959in}{1.426558in}}%
\pgfpathlineto{\pgfqpoint{2.465274in}{1.430126in}}%
\pgfpathlineto{\pgfqpoint{2.476728in}{1.430365in}}%
\pgfpathlineto{\pgfqpoint{2.485453in}{1.426054in}}%
\pgfpathlineto{\pgfqpoint{2.493958in}{1.429903in}}%
\pgfpathlineto{\pgfqpoint{2.514702in}{1.434217in}}%
\pgfpathlineto{\pgfqpoint{2.521870in}{1.439785in}}%
\pgfpathlineto{\pgfqpoint{2.532498in}{1.452111in}}%
\pgfpathlineto{\pgfqpoint{2.542167in}{1.457558in}}%
\pgfpathlineto{\pgfqpoint{2.542849in}{1.462782in}}%
\pgfpathlineto{\pgfqpoint{2.537774in}{1.474701in}}%
\pgfpathlineto{\pgfqpoint{2.541443in}{1.481714in}}%
\pgfpathlineto{\pgfqpoint{2.536503in}{1.489256in}}%
\pgfpathlineto{\pgfqpoint{2.528939in}{1.489782in}}%
\pgfpathlineto{\pgfqpoint{2.547870in}{1.512563in}}%
\pgfpathlineto{\pgfqpoint{2.550119in}{1.521245in}}%
\pgfpathlineto{\pgfqpoint{2.565018in}{1.543386in}}%
\pgfpathlineto{\pgfqpoint{2.578850in}{1.555368in}}%
\pgfpathlineto{\pgfqpoint{2.588347in}{1.560371in}}%
\pgfpathlineto{\pgfqpoint{2.619425in}{1.567413in}}%
\pgfpathlineto{\pgfqpoint{2.648552in}{1.575365in}}%
\pgfpathlineto{\pgfqpoint{2.652410in}{1.562882in}}%
\pgfpathlineto{\pgfqpoint{2.652683in}{1.551308in}}%
\pgfpathlineto{\pgfqpoint{2.659327in}{1.540801in}}%
\pgfpathlineto{\pgfqpoint{2.661375in}{1.530104in}}%
\pgfpathlineto{\pgfqpoint{2.658822in}{1.515590in}}%
\pgfpathlineto{\pgfqpoint{2.666598in}{1.496013in}}%
\pgfpathlineto{\pgfqpoint{2.675463in}{1.485127in}}%
\pgfpathlineto{\pgfqpoint{2.682320in}{1.449744in}}%
\pgfpathlineto{\pgfqpoint{2.686015in}{1.438080in}}%
\pgfpathlineto{\pgfqpoint{2.685061in}{1.394631in}}%
\pgfpathlineto{\pgfqpoint{2.686166in}{1.392957in}}%
\pgfpathlineto{\pgfqpoint{2.694238in}{1.346227in}}%
\pgfpathlineto{\pgfqpoint{2.698765in}{1.341969in}}%
\pgfpathlineto{\pgfqpoint{2.689028in}{1.332508in}}%
\pgfpathlineto{\pgfqpoint{2.693832in}{1.327078in}}%
\pgfpathlineto{\pgfqpoint{2.689609in}{1.318865in}}%
\pgfpathlineto{\pgfqpoint{2.689645in}{1.315368in}}%
\pgfpathlineto{\pgfqpoint{2.684366in}{1.312213in}}%
\pgfpathlineto{\pgfqpoint{2.681795in}{1.305239in}}%
\pgfpathlineto{\pgfqpoint{2.682610in}{1.324417in}}%
\pgfpathlineto{\pgfqpoint{2.666246in}{1.328639in}}%
\pgfpathlineto{\pgfqpoint{2.640655in}{1.337368in}}%
\pgfpathlineto{\pgfqpoint{2.637035in}{1.341668in}}%
\pgfpathlineto{\pgfqpoint{2.626342in}{1.342698in}}%
\pgfpathlineto{\pgfqpoint{2.620190in}{1.349098in}}%
\pgfpathlineto{\pgfqpoint{2.616871in}{1.361844in}}%
\pgfpathlineto{\pgfqpoint{2.608151in}{1.363429in}}%
\pgfpathlineto{\pgfqpoint{2.602319in}{1.370116in}}%
\pgfpathlineto{\pgfqpoint{2.528076in}{1.355148in}}%
\pgfpathlineto{\pgfqpoint{2.492594in}{1.347810in}}%
\pgfpathlineto{\pgfqpoint{2.438721in}{1.337977in}}%
\pgfpathlineto{\pgfqpoint{2.399926in}{1.331432in}}%
\pgfpathlineto{\pgfqpoint{2.397084in}{1.347781in}}%
\pgfpathclose%
\pgfusepath{fill}%
\end{pgfscope}%
\begin{pgfscope}%
\pgfpathrectangle{\pgfqpoint{0.100000in}{0.100000in}}{\pgfqpoint{2.862985in}{1.832778in}}%
\pgfusepath{clip}%
\pgfsetbuttcap%
\pgfsetmiterjoin%
\definecolor{currentfill}{rgb}{0.752312,0.227133,0.513149}%
\pgfsetfillcolor{currentfill}%
\pgfsetlinewidth{0.000000pt}%
\definecolor{currentstroke}{rgb}{0.000000,0.000000,0.000000}%
\pgfsetstrokecolor{currentstroke}%
\pgfsetstrokeopacity{0.000000}%
\pgfsetdash{}{0pt}%
\pgfpathmoveto{\pgfqpoint{2.695139in}{1.301344in}}%
\pgfpathlineto{\pgfqpoint{2.683657in}{1.297767in}}%
\pgfpathlineto{\pgfqpoint{2.681726in}{1.301055in}}%
\pgfpathlineto{\pgfqpoint{2.685411in}{1.312094in}}%
\pgfpathlineto{\pgfqpoint{2.692237in}{1.313181in}}%
\pgfpathlineto{\pgfqpoint{2.691626in}{1.316656in}}%
\pgfpathlineto{\pgfqpoint{2.697754in}{1.321872in}}%
\pgfpathlineto{\pgfqpoint{2.715472in}{1.326028in}}%
\pgfpathlineto{\pgfqpoint{2.723360in}{1.332319in}}%
\pgfpathlineto{\pgfqpoint{2.741079in}{1.337561in}}%
\pgfpathlineto{\pgfqpoint{2.749161in}{1.335643in}}%
\pgfpathlineto{\pgfqpoint{2.755959in}{1.344095in}}%
\pgfpathlineto{\pgfqpoint{2.766234in}{1.345211in}}%
\pgfpathlineto{\pgfqpoint{2.748704in}{1.328725in}}%
\pgfpathlineto{\pgfqpoint{2.695139in}{1.301344in}}%
\pgfpathclose%
\pgfusepath{fill}%
\end{pgfscope}%
\begin{pgfscope}%
\pgfpathrectangle{\pgfqpoint{0.100000in}{0.100000in}}{\pgfqpoint{2.862985in}{1.832778in}}%
\pgfusepath{clip}%
\pgfsetbuttcap%
\pgfsetmiterjoin%
\definecolor{currentfill}{rgb}{0.710549,0.182868,0.550004}%
\pgfsetfillcolor{currentfill}%
\pgfsetlinewidth{0.000000pt}%
\definecolor{currentstroke}{rgb}{0.000000,0.000000,0.000000}%
\pgfsetstrokecolor{currentstroke}%
\pgfsetstrokeopacity{0.000000}%
\pgfsetdash{}{0pt}%
\pgfpathmoveto{\pgfqpoint{2.437178in}{1.192485in}}%
\pgfpathlineto{\pgfqpoint{2.387526in}{1.184275in}}%
\pgfpathlineto{\pgfqpoint{2.378519in}{1.241018in}}%
\pgfpathlineto{\pgfqpoint{2.365146in}{1.324661in}}%
\pgfpathlineto{\pgfqpoint{2.397084in}{1.347781in}}%
\pgfpathlineto{\pgfqpoint{2.399926in}{1.331432in}}%
\pgfpathlineto{\pgfqpoint{2.438721in}{1.337977in}}%
\pgfpathlineto{\pgfqpoint{2.492594in}{1.347810in}}%
\pgfpathlineto{\pgfqpoint{2.528076in}{1.355148in}}%
\pgfpathlineto{\pgfqpoint{2.602319in}{1.370116in}}%
\pgfpathlineto{\pgfqpoint{2.608151in}{1.363429in}}%
\pgfpathlineto{\pgfqpoint{2.616871in}{1.361844in}}%
\pgfpathlineto{\pgfqpoint{2.620190in}{1.349098in}}%
\pgfpathlineto{\pgfqpoint{2.626342in}{1.342698in}}%
\pgfpathlineto{\pgfqpoint{2.637035in}{1.341668in}}%
\pgfpathlineto{\pgfqpoint{2.640655in}{1.337368in}}%
\pgfpathlineto{\pgfqpoint{2.636981in}{1.334051in}}%
\pgfpathlineto{\pgfqpoint{2.633670in}{1.322314in}}%
\pgfpathlineto{\pgfqpoint{2.625913in}{1.309125in}}%
\pgfpathlineto{\pgfqpoint{2.631120in}{1.303390in}}%
\pgfpathlineto{\pgfqpoint{2.626859in}{1.297245in}}%
\pgfpathlineto{\pgfqpoint{2.629263in}{1.283769in}}%
\pgfpathlineto{\pgfqpoint{2.636936in}{1.276695in}}%
\pgfpathlineto{\pgfqpoint{2.655144in}{1.265202in}}%
\pgfpathlineto{\pgfqpoint{2.640481in}{1.248986in}}%
\pgfpathlineto{\pgfqpoint{2.640275in}{1.242831in}}%
\pgfpathlineto{\pgfqpoint{2.628338in}{1.234894in}}%
\pgfpathlineto{\pgfqpoint{2.615138in}{1.233404in}}%
\pgfpathlineto{\pgfqpoint{2.611873in}{1.226505in}}%
\pgfpathlineto{\pgfqpoint{2.575158in}{1.218561in}}%
\pgfpathlineto{\pgfqpoint{2.532316in}{1.209933in}}%
\pgfpathlineto{\pgfqpoint{2.502868in}{1.204659in}}%
\pgfpathlineto{\pgfqpoint{2.437178in}{1.192485in}}%
\pgfpathclose%
\pgfusepath{fill}%
\end{pgfscope}%
\begin{pgfscope}%
\pgfpathrectangle{\pgfqpoint{0.100000in}{0.100000in}}{\pgfqpoint{2.862985in}{1.832778in}}%
\pgfusepath{clip}%
\pgfsetbuttcap%
\pgfsetmiterjoin%
\definecolor{currentfill}{rgb}{0.650746,0.125309,0.595617}%
\pgfsetfillcolor{currentfill}%
\pgfsetlinewidth{0.000000pt}%
\definecolor{currentstroke}{rgb}{0.000000,0.000000,0.000000}%
\pgfsetstrokecolor{currentstroke}%
\pgfsetstrokeopacity{0.000000}%
\pgfsetdash{}{0pt}%
\pgfpathmoveto{\pgfqpoint{2.686166in}{1.392957in}}%
\pgfpathlineto{\pgfqpoint{2.749362in}{1.407318in}}%
\pgfpathlineto{\pgfqpoint{2.763120in}{1.409488in}}%
\pgfpathlineto{\pgfqpoint{2.772226in}{1.373686in}}%
\pgfpathlineto{\pgfqpoint{2.770785in}{1.367275in}}%
\pgfpathlineto{\pgfqpoint{2.741533in}{1.355951in}}%
\pgfpathlineto{\pgfqpoint{2.724071in}{1.351992in}}%
\pgfpathlineto{\pgfqpoint{2.716651in}{1.343114in}}%
\pgfpathlineto{\pgfqpoint{2.693832in}{1.327078in}}%
\pgfpathlineto{\pgfqpoint{2.689028in}{1.332508in}}%
\pgfpathlineto{\pgfqpoint{2.698765in}{1.341969in}}%
\pgfpathlineto{\pgfqpoint{2.694238in}{1.346227in}}%
\pgfpathlineto{\pgfqpoint{2.686166in}{1.392957in}}%
\pgfpathclose%
\pgfusepath{fill}%
\end{pgfscope}%
\begin{pgfscope}%
\pgfpathrectangle{\pgfqpoint{0.100000in}{0.100000in}}{\pgfqpoint{2.862985in}{1.832778in}}%
\pgfusepath{clip}%
\pgfsetbuttcap%
\pgfsetmiterjoin%
\definecolor{currentfill}{rgb}{0.650746,0.125309,0.595617}%
\pgfsetfillcolor{currentfill}%
\pgfsetlinewidth{0.000000pt}%
\definecolor{currentstroke}{rgb}{0.000000,0.000000,0.000000}%
\pgfsetstrokecolor{currentstroke}%
\pgfsetstrokeopacity{0.000000}%
\pgfsetdash{}{0pt}%
\pgfpathmoveto{\pgfqpoint{2.763120in}{1.409488in}}%
\pgfpathlineto{\pgfqpoint{2.781882in}{1.415045in}}%
\pgfpathlineto{\pgfqpoint{2.787340in}{1.401600in}}%
\pgfpathlineto{\pgfqpoint{2.793598in}{1.398285in}}%
\pgfpathlineto{\pgfqpoint{2.785875in}{1.392622in}}%
\pgfpathlineto{\pgfqpoint{2.786849in}{1.375988in}}%
\pgfpathlineto{\pgfqpoint{2.770785in}{1.367275in}}%
\pgfpathlineto{\pgfqpoint{2.772226in}{1.373686in}}%
\pgfpathlineto{\pgfqpoint{2.763120in}{1.409488in}}%
\pgfpathclose%
\pgfusepath{fill}%
\end{pgfscope}%
\begin{pgfscope}%
\pgfpathrectangle{\pgfqpoint{0.100000in}{0.100000in}}{\pgfqpoint{2.862985in}{1.832778in}}%
\pgfusepath{clip}%
\pgfsetbuttcap%
\pgfsetmiterjoin%
\definecolor{currentfill}{rgb}{0.660374,0.134144,0.588971}%
\pgfsetfillcolor{currentfill}%
\pgfsetlinewidth{0.000000pt}%
\definecolor{currentstroke}{rgb}{0.000000,0.000000,0.000000}%
\pgfsetstrokecolor{currentstroke}%
\pgfsetstrokeopacity{0.000000}%
\pgfsetdash{}{0pt}%
\pgfpathmoveto{\pgfqpoint{2.626223in}{1.229031in}}%
\pgfpathlineto{\pgfqpoint{2.628338in}{1.234894in}}%
\pgfpathlineto{\pgfqpoint{2.640275in}{1.242831in}}%
\pgfpathlineto{\pgfqpoint{2.640481in}{1.248986in}}%
\pgfpathlineto{\pgfqpoint{2.655144in}{1.265202in}}%
\pgfpathlineto{\pgfqpoint{2.636936in}{1.276695in}}%
\pgfpathlineto{\pgfqpoint{2.629263in}{1.283769in}}%
\pgfpathlineto{\pgfqpoint{2.626859in}{1.297245in}}%
\pgfpathlineto{\pgfqpoint{2.631120in}{1.303390in}}%
\pgfpathlineto{\pgfqpoint{2.625913in}{1.309125in}}%
\pgfpathlineto{\pgfqpoint{2.633670in}{1.322314in}}%
\pgfpathlineto{\pgfqpoint{2.636981in}{1.334051in}}%
\pgfpathlineto{\pgfqpoint{2.640655in}{1.337368in}}%
\pgfpathlineto{\pgfqpoint{2.666246in}{1.328639in}}%
\pgfpathlineto{\pgfqpoint{2.682610in}{1.324417in}}%
\pgfpathlineto{\pgfqpoint{2.681795in}{1.305239in}}%
\pgfpathlineto{\pgfqpoint{2.671854in}{1.290696in}}%
\pgfpathlineto{\pgfqpoint{2.680045in}{1.288556in}}%
\pgfpathlineto{\pgfqpoint{2.688550in}{1.282315in}}%
\pgfpathlineto{\pgfqpoint{2.689051in}{1.265245in}}%
\pgfpathlineto{\pgfqpoint{2.686477in}{1.253161in}}%
\pgfpathlineto{\pgfqpoint{2.688190in}{1.243235in}}%
\pgfpathlineto{\pgfqpoint{2.673407in}{1.209740in}}%
\pgfpathlineto{\pgfqpoint{2.668094in}{1.194100in}}%
\pgfpathlineto{\pgfqpoint{2.660676in}{1.201703in}}%
\pgfpathlineto{\pgfqpoint{2.650838in}{1.200407in}}%
\pgfpathlineto{\pgfqpoint{2.626218in}{1.214632in}}%
\pgfpathlineto{\pgfqpoint{2.623699in}{1.222249in}}%
\pgfpathlineto{\pgfqpoint{2.626223in}{1.229031in}}%
\pgfpathclose%
\pgfusepath{fill}%
\end{pgfscope}%
\begin{pgfscope}%
\pgfpathrectangle{\pgfqpoint{0.100000in}{0.100000in}}{\pgfqpoint{2.862985in}{1.832778in}}%
\pgfusepath{clip}%
\pgfsetbuttcap%
\pgfsetmiterjoin%
\definecolor{currentfill}{rgb}{0.679160,0.151848,0.575189}%
\pgfsetfillcolor{currentfill}%
\pgfsetlinewidth{0.000000pt}%
\definecolor{currentstroke}{rgb}{0.000000,0.000000,0.000000}%
\pgfsetstrokecolor{currentstroke}%
\pgfsetstrokeopacity{0.000000}%
\pgfsetdash{}{0pt}%
\pgfpathmoveto{\pgfqpoint{2.037013in}{1.021868in}}%
\pgfpathlineto{\pgfqpoint{2.035862in}{1.037407in}}%
\pgfpathlineto{\pgfqpoint{2.040336in}{1.045116in}}%
\pgfpathlineto{\pgfqpoint{2.037387in}{1.048893in}}%
\pgfpathlineto{\pgfqpoint{2.048140in}{1.061261in}}%
\pgfpathlineto{\pgfqpoint{2.047524in}{1.064241in}}%
\pgfpathlineto{\pgfqpoint{2.057971in}{1.085818in}}%
\pgfpathlineto{\pgfqpoint{2.055859in}{1.096224in}}%
\pgfpathlineto{\pgfqpoint{2.048283in}{1.107118in}}%
\pgfpathlineto{\pgfqpoint{2.053542in}{1.120373in}}%
\pgfpathlineto{\pgfqpoint{2.046726in}{1.207597in}}%
\pgfpathlineto{\pgfqpoint{2.041818in}{1.268787in}}%
\pgfpathlineto{\pgfqpoint{2.048598in}{1.263717in}}%
\pgfpathlineto{\pgfqpoint{2.056163in}{1.263854in}}%
\pgfpathlineto{\pgfqpoint{2.074029in}{1.274203in}}%
\pgfpathlineto{\pgfqpoint{2.128831in}{1.279317in}}%
\pgfpathlineto{\pgfqpoint{2.169393in}{1.283592in}}%
\pgfpathlineto{\pgfqpoint{2.169752in}{1.279623in}}%
\pgfpathlineto{\pgfqpoint{2.179011in}{1.195770in}}%
\pgfpathlineto{\pgfqpoint{2.186981in}{1.117748in}}%
\pgfpathlineto{\pgfqpoint{2.183545in}{1.114085in}}%
\pgfpathlineto{\pgfqpoint{2.188822in}{1.104993in}}%
\pgfpathlineto{\pgfqpoint{2.188795in}{1.098441in}}%
\pgfpathlineto{\pgfqpoint{2.181265in}{1.096795in}}%
\pgfpathlineto{\pgfqpoint{2.172839in}{1.090478in}}%
\pgfpathlineto{\pgfqpoint{2.167134in}{1.092965in}}%
\pgfpathlineto{\pgfqpoint{2.158595in}{1.088922in}}%
\pgfpathlineto{\pgfqpoint{2.161237in}{1.080801in}}%
\pgfpathlineto{\pgfqpoint{2.152460in}{1.072642in}}%
\pgfpathlineto{\pgfqpoint{2.150037in}{1.063205in}}%
\pgfpathlineto{\pgfqpoint{2.144002in}{1.061656in}}%
\pgfpathlineto{\pgfqpoint{2.139498in}{1.054509in}}%
\pgfpathlineto{\pgfqpoint{2.140088in}{1.047311in}}%
\pgfpathlineto{\pgfqpoint{2.134790in}{1.042250in}}%
\pgfpathlineto{\pgfqpoint{2.126817in}{1.043031in}}%
\pgfpathlineto{\pgfqpoint{2.120757in}{1.050797in}}%
\pgfpathlineto{\pgfqpoint{2.110489in}{1.043318in}}%
\pgfpathlineto{\pgfqpoint{2.111209in}{1.036783in}}%
\pgfpathlineto{\pgfqpoint{2.099789in}{1.032965in}}%
\pgfpathlineto{\pgfqpoint{2.096919in}{1.037771in}}%
\pgfpathlineto{\pgfqpoint{2.085801in}{1.032321in}}%
\pgfpathlineto{\pgfqpoint{2.081739in}{1.024520in}}%
\pgfpathlineto{\pgfqpoint{2.068349in}{1.032520in}}%
\pgfpathlineto{\pgfqpoint{2.047127in}{1.027220in}}%
\pgfpathlineto{\pgfqpoint{2.037013in}{1.021868in}}%
\pgfpathclose%
\pgfusepath{fill}%
\end{pgfscope}%
\begin{pgfscope}%
\pgfpathrectangle{\pgfqpoint{0.100000in}{0.100000in}}{\pgfqpoint{2.862985in}{1.832778in}}%
\pgfusepath{clip}%
\pgfsetbuttcap%
\pgfsetmiterjoin%
\definecolor{currentfill}{rgb}{0.727670,0.200586,0.535293}%
\pgfsetfillcolor{currentfill}%
\pgfsetlinewidth{0.000000pt}%
\definecolor{currentstroke}{rgb}{0.000000,0.000000,0.000000}%
\pgfsetstrokecolor{currentstroke}%
\pgfsetstrokeopacity{0.000000}%
\pgfsetdash{}{0pt}%
\pgfpathmoveto{\pgfqpoint{0.540358in}{1.430483in}}%
\pgfpathlineto{\pgfqpoint{0.571046in}{1.421943in}}%
\pgfpathlineto{\pgfqpoint{0.621916in}{1.408826in}}%
\pgfpathlineto{\pgfqpoint{0.674647in}{1.395314in}}%
\pgfpathlineto{\pgfqpoint{0.723250in}{1.383593in}}%
\pgfpathlineto{\pgfqpoint{0.765246in}{1.374085in}}%
\pgfpathlineto{\pgfqpoint{0.810177in}{1.364205in}}%
\pgfpathlineto{\pgfqpoint{0.797195in}{1.302942in}}%
\pgfpathlineto{\pgfqpoint{0.785630in}{1.248537in}}%
\pgfpathlineto{\pgfqpoint{0.766658in}{1.160754in}}%
\pgfpathlineto{\pgfqpoint{0.752414in}{1.094498in}}%
\pgfpathlineto{\pgfqpoint{0.744719in}{1.057523in}}%
\pgfpathlineto{\pgfqpoint{0.734852in}{1.009534in}}%
\pgfpathlineto{\pgfqpoint{0.728025in}{0.999798in}}%
\pgfpathlineto{\pgfqpoint{0.722531in}{0.999473in}}%
\pgfpathlineto{\pgfqpoint{0.718612in}{1.007971in}}%
\pgfpathlineto{\pgfqpoint{0.709613in}{1.011057in}}%
\pgfpathlineto{\pgfqpoint{0.699916in}{1.009976in}}%
\pgfpathlineto{\pgfqpoint{0.697164in}{1.003024in}}%
\pgfpathlineto{\pgfqpoint{0.696877in}{0.978886in}}%
\pgfpathlineto{\pgfqpoint{0.693973in}{0.973374in}}%
\pgfpathlineto{\pgfqpoint{0.695633in}{0.953995in}}%
\pgfpathlineto{\pgfqpoint{0.689569in}{0.941085in}}%
\pgfpathlineto{\pgfqpoint{0.640397in}{1.016802in}}%
\pgfpathlineto{\pgfqpoint{0.591971in}{1.090379in}}%
\pgfpathlineto{\pgfqpoint{0.566763in}{1.128960in}}%
\pgfpathlineto{\pgfqpoint{0.545735in}{1.162262in}}%
\pgfpathlineto{\pgfqpoint{0.519238in}{1.203456in}}%
\pgfpathlineto{\pgfqpoint{0.489244in}{1.249665in}}%
\pgfpathlineto{\pgfqpoint{0.501576in}{1.293524in}}%
\pgfpathlineto{\pgfqpoint{0.526391in}{1.381486in}}%
\pgfpathlineto{\pgfqpoint{0.540358in}{1.430483in}}%
\pgfpathclose%
\pgfusepath{fill}%
\end{pgfscope}%
\begin{pgfscope}%
\pgfpathrectangle{\pgfqpoint{0.100000in}{0.100000in}}{\pgfqpoint{2.862985in}{1.832778in}}%
\pgfusepath{clip}%
\pgfsetbuttcap%
\pgfsetmiterjoin%
\definecolor{currentfill}{rgb}{0.393304,0.001114,0.655199}%
\pgfsetfillcolor{currentfill}%
\pgfsetlinewidth{0.000000pt}%
\definecolor{currentstroke}{rgb}{0.000000,0.000000,0.000000}%
\pgfsetstrokecolor{currentstroke}%
\pgfsetstrokeopacity{0.000000}%
\pgfsetdash{}{0pt}%
\pgfpathmoveto{\pgfqpoint{0.810177in}{1.364205in}}%
\pgfpathlineto{\pgfqpoint{0.858212in}{1.354650in}}%
\pgfpathlineto{\pgfqpoint{0.947141in}{1.337483in}}%
\pgfpathlineto{\pgfqpoint{0.936009in}{1.275734in}}%
\pgfpathlineto{\pgfqpoint{0.984967in}{1.267352in}}%
\pgfpathlineto{\pgfqpoint{1.029571in}{1.260214in}}%
\pgfpathlineto{\pgfqpoint{1.021820in}{1.211396in}}%
\pgfpathlineto{\pgfqpoint{1.013608in}{1.158753in}}%
\pgfpathlineto{\pgfqpoint{1.002613in}{1.089632in}}%
\pgfpathlineto{\pgfqpoint{1.002328in}{1.083837in}}%
\pgfpathlineto{\pgfqpoint{0.990915in}{1.012202in}}%
\pgfpathlineto{\pgfqpoint{0.943932in}{1.019484in}}%
\pgfpathlineto{\pgfqpoint{0.919991in}{1.024292in}}%
\pgfpathlineto{\pgfqpoint{0.833455in}{1.039504in}}%
\pgfpathlineto{\pgfqpoint{0.800868in}{1.045930in}}%
\pgfpathlineto{\pgfqpoint{0.744719in}{1.057523in}}%
\pgfpathlineto{\pgfqpoint{0.752414in}{1.094498in}}%
\pgfpathlineto{\pgfqpoint{0.766658in}{1.160754in}}%
\pgfpathlineto{\pgfqpoint{0.785630in}{1.248537in}}%
\pgfpathlineto{\pgfqpoint{0.797195in}{1.302942in}}%
\pgfpathlineto{\pgfqpoint{0.810177in}{1.364205in}}%
\pgfpathclose%
\pgfusepath{fill}%
\end{pgfscope}%
\begin{pgfscope}%
\pgfpathrectangle{\pgfqpoint{0.100000in}{0.100000in}}{\pgfqpoint{2.862985in}{1.832778in}}%
\pgfusepath{clip}%
\pgfsetbuttcap%
\pgfsetmiterjoin%
\definecolor{currentfill}{rgb}{0.748289,0.222711,0.516834}%
\pgfsetfillcolor{currentfill}%
\pgfsetlinewidth{0.000000pt}%
\definecolor{currentstroke}{rgb}{0.000000,0.000000,0.000000}%
\pgfsetstrokecolor{currentstroke}%
\pgfsetstrokeopacity{0.000000}%
\pgfsetdash{}{0pt}%
\pgfpathmoveto{\pgfqpoint{0.352592in}{1.488751in}}%
\pgfpathlineto{\pgfqpoint{0.369661in}{1.482798in}}%
\pgfpathlineto{\pgfqpoint{0.438032in}{1.461511in}}%
\pgfpathlineto{\pgfqpoint{0.500927in}{1.441715in}}%
\pgfpathlineto{\pgfqpoint{0.540358in}{1.430483in}}%
\pgfpathlineto{\pgfqpoint{0.526391in}{1.381486in}}%
\pgfpathlineto{\pgfqpoint{0.501576in}{1.293524in}}%
\pgfpathlineto{\pgfqpoint{0.489244in}{1.249665in}}%
\pgfpathlineto{\pgfqpoint{0.519238in}{1.203456in}}%
\pgfpathlineto{\pgfqpoint{0.545735in}{1.162262in}}%
\pgfpathlineto{\pgfqpoint{0.566763in}{1.128960in}}%
\pgfpathlineto{\pgfqpoint{0.591971in}{1.090379in}}%
\pgfpathlineto{\pgfqpoint{0.640397in}{1.016802in}}%
\pgfpathlineto{\pgfqpoint{0.689569in}{0.941085in}}%
\pgfpathlineto{\pgfqpoint{0.687595in}{0.933576in}}%
\pgfpathlineto{\pgfqpoint{0.693526in}{0.921617in}}%
\pgfpathlineto{\pgfqpoint{0.694699in}{0.905259in}}%
\pgfpathlineto{\pgfqpoint{0.703313in}{0.897324in}}%
\pgfpathlineto{\pgfqpoint{0.703654in}{0.890930in}}%
\pgfpathlineto{\pgfqpoint{0.688230in}{0.883673in}}%
\pgfpathlineto{\pgfqpoint{0.680882in}{0.876407in}}%
\pgfpathlineto{\pgfqpoint{0.678591in}{0.860364in}}%
\pgfpathlineto{\pgfqpoint{0.674871in}{0.851614in}}%
\pgfpathlineto{\pgfqpoint{0.667035in}{0.844239in}}%
\pgfpathlineto{\pgfqpoint{0.659240in}{0.825061in}}%
\pgfpathlineto{\pgfqpoint{0.670200in}{0.815071in}}%
\pgfpathlineto{\pgfqpoint{0.668785in}{0.806836in}}%
\pgfpathlineto{\pgfqpoint{0.659798in}{0.801105in}}%
\pgfpathlineto{\pgfqpoint{0.653603in}{0.802117in}}%
\pgfpathlineto{\pgfqpoint{0.580470in}{0.812266in}}%
\pgfpathlineto{\pgfqpoint{0.526454in}{0.819928in}}%
\pgfpathlineto{\pgfqpoint{0.528816in}{0.828649in}}%
\pgfpathlineto{\pgfqpoint{0.525300in}{0.843128in}}%
\pgfpathlineto{\pgfqpoint{0.524945in}{0.857733in}}%
\pgfpathlineto{\pgfqpoint{0.522647in}{0.866279in}}%
\pgfpathlineto{\pgfqpoint{0.515568in}{0.878498in}}%
\pgfpathlineto{\pgfqpoint{0.495206in}{0.906691in}}%
\pgfpathlineto{\pgfqpoint{0.488537in}{0.910145in}}%
\pgfpathlineto{\pgfqpoint{0.479944in}{0.910090in}}%
\pgfpathlineto{\pgfqpoint{0.481908in}{0.918984in}}%
\pgfpathlineto{\pgfqpoint{0.477839in}{0.930106in}}%
\pgfpathlineto{\pgfqpoint{0.457853in}{0.935633in}}%
\pgfpathlineto{\pgfqpoint{0.445671in}{0.945869in}}%
\pgfpathlineto{\pgfqpoint{0.444661in}{0.952132in}}%
\pgfpathlineto{\pgfqpoint{0.430633in}{0.967637in}}%
\pgfpathlineto{\pgfqpoint{0.417280in}{0.970604in}}%
\pgfpathlineto{\pgfqpoint{0.404908in}{0.978484in}}%
\pgfpathlineto{\pgfqpoint{0.388637in}{0.981210in}}%
\pgfpathlineto{\pgfqpoint{0.381685in}{0.991719in}}%
\pgfpathlineto{\pgfqpoint{0.388287in}{1.008354in}}%
\pgfpathlineto{\pgfqpoint{0.386268in}{1.012095in}}%
\pgfpathlineto{\pgfqpoint{0.391766in}{1.025962in}}%
\pgfpathlineto{\pgfqpoint{0.381986in}{1.033342in}}%
\pgfpathlineto{\pgfqpoint{0.385159in}{1.046721in}}%
\pgfpathlineto{\pgfqpoint{0.379935in}{1.050141in}}%
\pgfpathlineto{\pgfqpoint{0.375424in}{1.062787in}}%
\pgfpathlineto{\pgfqpoint{0.369981in}{1.066644in}}%
\pgfpathlineto{\pgfqpoint{0.369572in}{1.075778in}}%
\pgfpathlineto{\pgfqpoint{0.365307in}{1.082219in}}%
\pgfpathlineto{\pgfqpoint{0.358881in}{1.103911in}}%
\pgfpathlineto{\pgfqpoint{0.351850in}{1.114301in}}%
\pgfpathlineto{\pgfqpoint{0.353378in}{1.131939in}}%
\pgfpathlineto{\pgfqpoint{0.361647in}{1.133714in}}%
\pgfpathlineto{\pgfqpoint{0.367031in}{1.143280in}}%
\pgfpathlineto{\pgfqpoint{0.363778in}{1.153656in}}%
\pgfpathlineto{\pgfqpoint{0.354984in}{1.155388in}}%
\pgfpathlineto{\pgfqpoint{0.350596in}{1.160241in}}%
\pgfpathlineto{\pgfqpoint{0.343491in}{1.178092in}}%
\pgfpathlineto{\pgfqpoint{0.346812in}{1.184503in}}%
\pgfpathlineto{\pgfqpoint{0.344457in}{1.196440in}}%
\pgfpathlineto{\pgfqpoint{0.349672in}{1.211900in}}%
\pgfpathlineto{\pgfqpoint{0.355782in}{1.206207in}}%
\pgfpathlineto{\pgfqpoint{0.353009in}{1.199485in}}%
\pgfpathlineto{\pgfqpoint{0.363597in}{1.188830in}}%
\pgfpathlineto{\pgfqpoint{0.362921in}{1.204652in}}%
\pgfpathlineto{\pgfqpoint{0.358384in}{1.208896in}}%
\pgfpathlineto{\pgfqpoint{0.363570in}{1.222808in}}%
\pgfpathlineto{\pgfqpoint{0.383185in}{1.220593in}}%
\pgfpathlineto{\pgfqpoint{0.380545in}{1.225766in}}%
\pgfpathlineto{\pgfqpoint{0.367599in}{1.225257in}}%
\pgfpathlineto{\pgfqpoint{0.361441in}{1.233093in}}%
\pgfpathlineto{\pgfqpoint{0.353689in}{1.226135in}}%
\pgfpathlineto{\pgfqpoint{0.349576in}{1.214504in}}%
\pgfpathlineto{\pgfqpoint{0.338597in}{1.230154in}}%
\pgfpathlineto{\pgfqpoint{0.334369in}{1.233003in}}%
\pgfpathlineto{\pgfqpoint{0.335980in}{1.250036in}}%
\pgfpathlineto{\pgfqpoint{0.332634in}{1.260081in}}%
\pgfpathlineto{\pgfqpoint{0.326568in}{1.269517in}}%
\pgfpathlineto{\pgfqpoint{0.314145in}{1.298432in}}%
\pgfpathlineto{\pgfqpoint{0.318207in}{1.304826in}}%
\pgfpathlineto{\pgfqpoint{0.318159in}{1.325033in}}%
\pgfpathlineto{\pgfqpoint{0.324864in}{1.336303in}}%
\pgfpathlineto{\pgfqpoint{0.326401in}{1.353918in}}%
\pgfpathlineto{\pgfqpoint{0.320079in}{1.374099in}}%
\pgfpathlineto{\pgfqpoint{0.311692in}{1.386945in}}%
\pgfpathlineto{\pgfqpoint{0.313188in}{1.398536in}}%
\pgfpathlineto{\pgfqpoint{0.336740in}{1.426591in}}%
\pgfpathlineto{\pgfqpoint{0.337910in}{1.436161in}}%
\pgfpathlineto{\pgfqpoint{0.348523in}{1.454416in}}%
\pgfpathlineto{\pgfqpoint{0.349995in}{1.471719in}}%
\pgfpathlineto{\pgfqpoint{0.346583in}{1.476134in}}%
\pgfpathlineto{\pgfqpoint{0.352592in}{1.488751in}}%
\pgfpathclose%
\pgfusepath{fill}%
\end{pgfscope}%
\begin{pgfscope}%
\pgfpathrectangle{\pgfqpoint{0.100000in}{0.100000in}}{\pgfqpoint{2.862985in}{1.832778in}}%
\pgfusepath{clip}%
\pgfsetbuttcap%
\pgfsetmiterjoin%
\definecolor{currentfill}{rgb}{0.723444,0.196158,0.538981}%
\pgfsetfillcolor{currentfill}%
\pgfsetlinewidth{0.000000pt}%
\definecolor{currentstroke}{rgb}{0.000000,0.000000,0.000000}%
\pgfsetstrokecolor{currentstroke}%
\pgfsetstrokeopacity{0.000000}%
\pgfsetdash{}{0pt}%
\pgfpathmoveto{\pgfqpoint{2.186981in}{1.117748in}}%
\pgfpathlineto{\pgfqpoint{2.179011in}{1.195770in}}%
\pgfpathlineto{\pgfqpoint{2.169752in}{1.279623in}}%
\pgfpathlineto{\pgfqpoint{2.230390in}{1.288653in}}%
\pgfpathlineto{\pgfqpoint{2.246483in}{1.284465in}}%
\pgfpathlineto{\pgfqpoint{2.254202in}{1.279918in}}%
\pgfpathlineto{\pgfqpoint{2.263874in}{1.281179in}}%
\pgfpathlineto{\pgfqpoint{2.276627in}{1.273651in}}%
\pgfpathlineto{\pgfqpoint{2.300380in}{1.284855in}}%
\pgfpathlineto{\pgfqpoint{2.313489in}{1.285229in}}%
\pgfpathlineto{\pgfqpoint{2.328797in}{1.302307in}}%
\pgfpathlineto{\pgfqpoint{2.344372in}{1.312695in}}%
\pgfpathlineto{\pgfqpoint{2.365146in}{1.324661in}}%
\pgfpathlineto{\pgfqpoint{2.378519in}{1.241018in}}%
\pgfpathlineto{\pgfqpoint{2.372306in}{1.235646in}}%
\pgfpathlineto{\pgfqpoint{2.376317in}{1.230712in}}%
\pgfpathlineto{\pgfqpoint{2.377915in}{1.219896in}}%
\pgfpathlineto{\pgfqpoint{2.374803in}{1.209737in}}%
\pgfpathlineto{\pgfqpoint{2.374391in}{1.194887in}}%
\pgfpathlineto{\pgfqpoint{2.371468in}{1.175560in}}%
\pgfpathlineto{\pgfqpoint{2.356678in}{1.158322in}}%
\pgfpathlineto{\pgfqpoint{2.350452in}{1.154663in}}%
\pgfpathlineto{\pgfqpoint{2.345610in}{1.157813in}}%
\pgfpathlineto{\pgfqpoint{2.333627in}{1.141383in}}%
\pgfpathlineto{\pgfqpoint{2.334750in}{1.125564in}}%
\pgfpathlineto{\pgfqpoint{2.321406in}{1.129368in}}%
\pgfpathlineto{\pgfqpoint{2.315087in}{1.113569in}}%
\pgfpathlineto{\pgfqpoint{2.318281in}{1.102373in}}%
\pgfpathlineto{\pgfqpoint{2.313281in}{1.100719in}}%
\pgfpathlineto{\pgfqpoint{2.312581in}{1.091873in}}%
\pgfpathlineto{\pgfqpoint{2.300308in}{1.088305in}}%
\pgfpathlineto{\pgfqpoint{2.293930in}{1.095447in}}%
\pgfpathlineto{\pgfqpoint{2.285724in}{1.098217in}}%
\pgfpathlineto{\pgfqpoint{2.283796in}{1.105459in}}%
\pgfpathlineto{\pgfqpoint{2.276376in}{1.104185in}}%
\pgfpathlineto{\pgfqpoint{2.271515in}{1.097524in}}%
\pgfpathlineto{\pgfqpoint{2.264561in}{1.095181in}}%
\pgfpathlineto{\pgfqpoint{2.252277in}{1.099895in}}%
\pgfpathlineto{\pgfqpoint{2.245483in}{1.094286in}}%
\pgfpathlineto{\pgfqpoint{2.235827in}{1.100922in}}%
\pgfpathlineto{\pgfqpoint{2.217303in}{1.102959in}}%
\pgfpathlineto{\pgfqpoint{2.211715in}{1.115014in}}%
\pgfpathlineto{\pgfqpoint{2.204636in}{1.120355in}}%
\pgfpathlineto{\pgfqpoint{2.197773in}{1.116936in}}%
\pgfpathlineto{\pgfqpoint{2.186981in}{1.117748in}}%
\pgfpathclose%
\pgfusepath{fill}%
\end{pgfscope}%
\begin{pgfscope}%
\pgfpathrectangle{\pgfqpoint{0.100000in}{0.100000in}}{\pgfqpoint{2.862985in}{1.832778in}}%
\pgfusepath{clip}%
\pgfsetbuttcap%
\pgfsetmiterjoin%
\definecolor{currentfill}{rgb}{0.645872,0.120898,0.598867}%
\pgfsetfillcolor{currentfill}%
\pgfsetlinewidth{0.000000pt}%
\definecolor{currentstroke}{rgb}{0.000000,0.000000,0.000000}%
\pgfsetstrokecolor{currentstroke}%
\pgfsetstrokeopacity{0.000000}%
\pgfsetdash{}{0pt}%
\pgfpathmoveto{\pgfqpoint{1.986225in}{0.966856in}}%
\pgfpathlineto{\pgfqpoint{1.979121in}{0.972628in}}%
\pgfpathlineto{\pgfqpoint{1.973328in}{0.969882in}}%
\pgfpathlineto{\pgfqpoint{1.965922in}{0.983711in}}%
\pgfpathlineto{\pgfqpoint{1.969705in}{0.992402in}}%
\pgfpathlineto{\pgfqpoint{1.964252in}{1.002186in}}%
\pgfpathlineto{\pgfqpoint{1.963959in}{1.009888in}}%
\pgfpathlineto{\pgfqpoint{1.956591in}{1.012632in}}%
\pgfpathlineto{\pgfqpoint{1.953174in}{1.018436in}}%
\pgfpathlineto{\pgfqpoint{1.938768in}{1.025684in}}%
\pgfpathlineto{\pgfqpoint{1.926256in}{1.034614in}}%
\pgfpathlineto{\pgfqpoint{1.920442in}{1.041355in}}%
\pgfpathlineto{\pgfqpoint{1.920322in}{1.049578in}}%
\pgfpathlineto{\pgfqpoint{1.928102in}{1.065365in}}%
\pgfpathlineto{\pgfqpoint{1.930495in}{1.077440in}}%
\pgfpathlineto{\pgfqpoint{1.924157in}{1.084263in}}%
\pgfpathlineto{\pgfqpoint{1.915763in}{1.086834in}}%
\pgfpathlineto{\pgfqpoint{1.905579in}{1.081212in}}%
\pgfpathlineto{\pgfqpoint{1.901269in}{1.090863in}}%
\pgfpathlineto{\pgfqpoint{1.899094in}{1.103945in}}%
\pgfpathlineto{\pgfqpoint{1.884086in}{1.115607in}}%
\pgfpathlineto{\pgfqpoint{1.881087in}{1.120782in}}%
\pgfpathlineto{\pgfqpoint{1.867379in}{1.132498in}}%
\pgfpathlineto{\pgfqpoint{1.863083in}{1.141015in}}%
\pgfpathlineto{\pgfqpoint{1.859210in}{1.157906in}}%
\pgfpathlineto{\pgfqpoint{1.861865in}{1.172908in}}%
\pgfpathlineto{\pgfqpoint{1.865408in}{1.174999in}}%
\pgfpathlineto{\pgfqpoint{1.864771in}{1.187558in}}%
\pgfpathlineto{\pgfqpoint{1.874773in}{1.191288in}}%
\pgfpathlineto{\pgfqpoint{1.877793in}{1.202558in}}%
\pgfpathlineto{\pgfqpoint{1.883547in}{1.210143in}}%
\pgfpathlineto{\pgfqpoint{1.883256in}{1.219794in}}%
\pgfpathlineto{\pgfqpoint{1.876124in}{1.227469in}}%
\pgfpathlineto{\pgfqpoint{1.877818in}{1.238212in}}%
\pgfpathlineto{\pgfqpoint{1.896306in}{1.242884in}}%
\pgfpathlineto{\pgfqpoint{1.910490in}{1.251413in}}%
\pgfpathlineto{\pgfqpoint{1.911979in}{1.262151in}}%
\pgfpathlineto{\pgfqpoint{1.916914in}{1.265532in}}%
\pgfpathlineto{\pgfqpoint{1.918802in}{1.276813in}}%
\pgfpathlineto{\pgfqpoint{1.917285in}{1.284265in}}%
\pgfpathlineto{\pgfqpoint{1.907588in}{1.290454in}}%
\pgfpathlineto{\pgfqpoint{1.903684in}{1.299681in}}%
\pgfpathlineto{\pgfqpoint{1.894114in}{1.308600in}}%
\pgfpathlineto{\pgfqpoint{1.972758in}{1.311952in}}%
\pgfpathlineto{\pgfqpoint{2.025520in}{1.315700in}}%
\pgfpathlineto{\pgfqpoint{2.024554in}{1.304601in}}%
\pgfpathlineto{\pgfqpoint{2.033546in}{1.289293in}}%
\pgfpathlineto{\pgfqpoint{2.037322in}{1.276213in}}%
\pgfpathlineto{\pgfqpoint{2.041818in}{1.268787in}}%
\pgfpathlineto{\pgfqpoint{2.046726in}{1.207597in}}%
\pgfpathlineto{\pgfqpoint{2.053542in}{1.120373in}}%
\pgfpathlineto{\pgfqpoint{2.048283in}{1.107118in}}%
\pgfpathlineto{\pgfqpoint{2.055859in}{1.096224in}}%
\pgfpathlineto{\pgfqpoint{2.057971in}{1.085818in}}%
\pgfpathlineto{\pgfqpoint{2.047524in}{1.064241in}}%
\pgfpathlineto{\pgfqpoint{2.048140in}{1.061261in}}%
\pgfpathlineto{\pgfqpoint{2.037387in}{1.048893in}}%
\pgfpathlineto{\pgfqpoint{2.040336in}{1.045116in}}%
\pgfpathlineto{\pgfqpoint{2.035862in}{1.037407in}}%
\pgfpathlineto{\pgfqpoint{2.037013in}{1.021868in}}%
\pgfpathlineto{\pgfqpoint{2.031580in}{1.012333in}}%
\pgfpathlineto{\pgfqpoint{2.035999in}{1.001067in}}%
\pgfpathlineto{\pgfqpoint{2.017483in}{0.994944in}}%
\pgfpathlineto{\pgfqpoint{2.015777in}{0.988286in}}%
\pgfpathlineto{\pgfqpoint{2.020833in}{0.979847in}}%
\pgfpathlineto{\pgfqpoint{2.018506in}{0.974346in}}%
\pgfpathlineto{\pgfqpoint{1.998652in}{0.981141in}}%
\pgfpathlineto{\pgfqpoint{1.992101in}{0.981825in}}%
\pgfpathlineto{\pgfqpoint{1.983934in}{0.971495in}}%
\pgfpathlineto{\pgfqpoint{1.986225in}{0.966856in}}%
\pgfpathclose%
\pgfusepath{fill}%
\end{pgfscope}%
\begin{pgfscope}%
\pgfpathrectangle{\pgfqpoint{0.100000in}{0.100000in}}{\pgfqpoint{2.862985in}{1.832778in}}%
\pgfusepath{clip}%
\pgfsetbuttcap%
\pgfsetmiterjoin%
\definecolor{currentfill}{rgb}{0.362553,0.003243,0.649245}%
\pgfsetfillcolor{currentfill}%
\pgfsetlinewidth{0.000000pt}%
\definecolor{currentstroke}{rgb}{0.000000,0.000000,0.000000}%
\pgfsetstrokecolor{currentstroke}%
\pgfsetstrokeopacity{0.000000}%
\pgfsetdash{}{0pt}%
\pgfpathmoveto{\pgfqpoint{2.564161in}{1.156689in}}%
\pgfpathlineto{\pgfqpoint{2.558697in}{1.164803in}}%
\pgfpathlineto{\pgfqpoint{2.561780in}{1.169343in}}%
\pgfpathlineto{\pgfqpoint{2.569331in}{1.164244in}}%
\pgfpathlineto{\pgfqpoint{2.564161in}{1.156689in}}%
\pgfpathclose%
\pgfusepath{fill}%
\end{pgfscope}%
\begin{pgfscope}%
\pgfpathrectangle{\pgfqpoint{0.100000in}{0.100000in}}{\pgfqpoint{2.862985in}{1.832778in}}%
\pgfusepath{clip}%
\pgfsetbuttcap%
\pgfsetmiterjoin%
\definecolor{currentfill}{rgb}{0.783383,0.262500,0.483918}%
\pgfsetfillcolor{currentfill}%
\pgfsetlinewidth{0.000000pt}%
\definecolor{currentstroke}{rgb}{0.000000,0.000000,0.000000}%
\pgfsetstrokecolor{currentstroke}%
\pgfsetstrokeopacity{0.000000}%
\pgfsetdash{}{0pt}%
\pgfpathmoveto{\pgfqpoint{2.611873in}{1.226505in}}%
\pgfpathlineto{\pgfqpoint{2.615138in}{1.233404in}}%
\pgfpathlineto{\pgfqpoint{2.628338in}{1.234894in}}%
\pgfpathlineto{\pgfqpoint{2.626223in}{1.229031in}}%
\pgfpathlineto{\pgfqpoint{2.621868in}{1.221542in}}%
\pgfpathlineto{\pgfqpoint{2.624819in}{1.212619in}}%
\pgfpathlineto{\pgfqpoint{2.636465in}{1.201927in}}%
\pgfpathlineto{\pgfqpoint{2.639168in}{1.190668in}}%
\pgfpathlineto{\pgfqpoint{2.652585in}{1.176648in}}%
\pgfpathlineto{\pgfqpoint{2.657849in}{1.177248in}}%
\pgfpathlineto{\pgfqpoint{2.664406in}{1.156201in}}%
\pgfpathlineto{\pgfqpoint{2.663331in}{1.155991in}}%
\pgfpathlineto{\pgfqpoint{2.662135in}{1.155756in}}%
\pgfpathlineto{\pgfqpoint{2.632891in}{1.150158in}}%
\pgfpathlineto{\pgfqpoint{2.617245in}{1.205802in}}%
\pgfpathlineto{\pgfqpoint{2.611873in}{1.226505in}}%
\pgfpathclose%
\pgfusepath{fill}%
\end{pgfscope}%
\begin{pgfscope}%
\pgfpathrectangle{\pgfqpoint{0.100000in}{0.100000in}}{\pgfqpoint{2.862985in}{1.832778in}}%
\pgfusepath{clip}%
\pgfsetbuttcap%
\pgfsetmiterjoin%
\definecolor{currentfill}{rgb}{0.918109,0.453603,0.339529}%
\pgfsetfillcolor{currentfill}%
\pgfsetlinewidth{0.000000pt}%
\definecolor{currentstroke}{rgb}{0.000000,0.000000,0.000000}%
\pgfsetstrokecolor{currentstroke}%
\pgfsetstrokeopacity{0.000000}%
\pgfsetdash{}{0pt}%
\pgfpathmoveto{\pgfqpoint{2.338724in}{1.037893in}}%
\pgfpathlineto{\pgfqpoint{2.329483in}{1.038233in}}%
\pgfpathlineto{\pgfqpoint{2.320968in}{1.044098in}}%
\pgfpathlineto{\pgfqpoint{2.309465in}{1.061926in}}%
\pgfpathlineto{\pgfqpoint{2.299672in}{1.071367in}}%
\pgfpathlineto{\pgfqpoint{2.302231in}{1.078659in}}%
\pgfpathlineto{\pgfqpoint{2.300308in}{1.088305in}}%
\pgfpathlineto{\pgfqpoint{2.312581in}{1.091873in}}%
\pgfpathlineto{\pgfqpoint{2.313281in}{1.100719in}}%
\pgfpathlineto{\pgfqpoint{2.318281in}{1.102373in}}%
\pgfpathlineto{\pgfqpoint{2.315087in}{1.113569in}}%
\pgfpathlineto{\pgfqpoint{2.321406in}{1.129368in}}%
\pgfpathlineto{\pgfqpoint{2.334750in}{1.125564in}}%
\pgfpathlineto{\pgfqpoint{2.333627in}{1.141383in}}%
\pgfpathlineto{\pgfqpoint{2.345610in}{1.157813in}}%
\pgfpathlineto{\pgfqpoint{2.350452in}{1.154663in}}%
\pgfpathlineto{\pgfqpoint{2.356678in}{1.158322in}}%
\pgfpathlineto{\pgfqpoint{2.371468in}{1.175560in}}%
\pgfpathlineto{\pgfqpoint{2.374391in}{1.194887in}}%
\pgfpathlineto{\pgfqpoint{2.374803in}{1.209737in}}%
\pgfpathlineto{\pgfqpoint{2.377915in}{1.219896in}}%
\pgfpathlineto{\pgfqpoint{2.376317in}{1.230712in}}%
\pgfpathlineto{\pgfqpoint{2.372306in}{1.235646in}}%
\pgfpathlineto{\pgfqpoint{2.378519in}{1.241018in}}%
\pgfpathlineto{\pgfqpoint{2.387526in}{1.184275in}}%
\pgfpathlineto{\pgfqpoint{2.437178in}{1.192485in}}%
\pgfpathlineto{\pgfqpoint{2.442325in}{1.160097in}}%
\pgfpathlineto{\pgfqpoint{2.460300in}{1.181474in}}%
\pgfpathlineto{\pgfqpoint{2.464536in}{1.179343in}}%
\pgfpathlineto{\pgfqpoint{2.470767in}{1.191730in}}%
\pgfpathlineto{\pgfqpoint{2.479399in}{1.187846in}}%
\pgfpathlineto{\pgfqpoint{2.486925in}{1.188584in}}%
\pgfpathlineto{\pgfqpoint{2.491891in}{1.197195in}}%
\pgfpathlineto{\pgfqpoint{2.499098in}{1.201985in}}%
\pgfpathlineto{\pgfqpoint{2.509119in}{1.197766in}}%
\pgfpathlineto{\pgfqpoint{2.515712in}{1.199224in}}%
\pgfpathlineto{\pgfqpoint{2.525158in}{1.182877in}}%
\pgfpathlineto{\pgfqpoint{2.522427in}{1.170500in}}%
\pgfpathlineto{\pgfqpoint{2.497755in}{1.184435in}}%
\pgfpathlineto{\pgfqpoint{2.493134in}{1.172996in}}%
\pgfpathlineto{\pgfqpoint{2.494659in}{1.167732in}}%
\pgfpathlineto{\pgfqpoint{2.484159in}{1.149941in}}%
\pgfpathlineto{\pgfqpoint{2.479203in}{1.146402in}}%
\pgfpathlineto{\pgfqpoint{2.476965in}{1.138556in}}%
\pgfpathlineto{\pgfqpoint{2.468517in}{1.139377in}}%
\pgfpathlineto{\pgfqpoint{2.462449in}{1.117963in}}%
\pgfpathlineto{\pgfqpoint{2.459053in}{1.113049in}}%
\pgfpathlineto{\pgfqpoint{2.450326in}{1.114688in}}%
\pgfpathlineto{\pgfqpoint{2.441399in}{1.121442in}}%
\pgfpathlineto{\pgfqpoint{2.441091in}{1.111076in}}%
\pgfpathlineto{\pgfqpoint{2.432449in}{1.093654in}}%
\pgfpathlineto{\pgfqpoint{2.430300in}{1.081251in}}%
\pgfpathlineto{\pgfqpoint{2.423669in}{1.072991in}}%
\pgfpathlineto{\pgfqpoint{2.418643in}{1.059800in}}%
\pgfpathlineto{\pgfqpoint{2.418396in}{1.047589in}}%
\pgfpathlineto{\pgfqpoint{2.410633in}{1.045308in}}%
\pgfpathlineto{\pgfqpoint{2.401826in}{1.038395in}}%
\pgfpathlineto{\pgfqpoint{2.393347in}{1.037079in}}%
\pgfpathlineto{\pgfqpoint{2.391393in}{1.031228in}}%
\pgfpathlineto{\pgfqpoint{2.377730in}{1.025236in}}%
\pgfpathlineto{\pgfqpoint{2.370094in}{1.030341in}}%
\pgfpathlineto{\pgfqpoint{2.361560in}{1.020647in}}%
\pgfpathlineto{\pgfqpoint{2.346890in}{1.023505in}}%
\pgfpathlineto{\pgfqpoint{2.337893in}{1.033679in}}%
\pgfpathlineto{\pgfqpoint{2.338724in}{1.037893in}}%
\pgfpathclose%
\pgfusepath{fill}%
\end{pgfscope}%
\begin{pgfscope}%
\pgfpathrectangle{\pgfqpoint{0.100000in}{0.100000in}}{\pgfqpoint{2.862985in}{1.832778in}}%
\pgfusepath{clip}%
\pgfsetbuttcap%
\pgfsetmiterjoin%
\definecolor{currentfill}{rgb}{0.579029,0.064296,0.635126}%
\pgfsetfillcolor{currentfill}%
\pgfsetlinewidth{0.000000pt}%
\definecolor{currentstroke}{rgb}{0.000000,0.000000,0.000000}%
\pgfsetstrokecolor{currentstroke}%
\pgfsetstrokeopacity{0.000000}%
\pgfsetdash{}{0pt}%
\pgfpathmoveto{\pgfqpoint{2.662135in}{1.155756in}}%
\pgfpathlineto{\pgfqpoint{2.661760in}{1.144322in}}%
\pgfpathlineto{\pgfqpoint{2.657360in}{1.138701in}}%
\pgfpathlineto{\pgfqpoint{2.654551in}{1.126221in}}%
\pgfpathlineto{\pgfqpoint{2.641870in}{1.120469in}}%
\pgfpathlineto{\pgfqpoint{2.631229in}{1.118793in}}%
\pgfpathlineto{\pgfqpoint{2.634330in}{1.127000in}}%
\pgfpathlineto{\pgfqpoint{2.617947in}{1.133877in}}%
\pgfpathlineto{\pgfqpoint{2.604638in}{1.142487in}}%
\pgfpathlineto{\pgfqpoint{2.612878in}{1.155358in}}%
\pgfpathlineto{\pgfqpoint{2.606274in}{1.165363in}}%
\pgfpathlineto{\pgfqpoint{2.607051in}{1.173993in}}%
\pgfpathlineto{\pgfqpoint{2.602276in}{1.176362in}}%
\pgfpathlineto{\pgfqpoint{2.598402in}{1.190382in}}%
\pgfpathlineto{\pgfqpoint{2.607564in}{1.209388in}}%
\pgfpathlineto{\pgfqpoint{2.600681in}{1.212492in}}%
\pgfpathlineto{\pgfqpoint{2.598896in}{1.203115in}}%
\pgfpathlineto{\pgfqpoint{2.589070in}{1.200505in}}%
\pgfpathlineto{\pgfqpoint{2.590103in}{1.169655in}}%
\pgfpathlineto{\pgfqpoint{2.588306in}{1.159726in}}%
\pgfpathlineto{\pgfqpoint{2.593235in}{1.145575in}}%
\pgfpathlineto{\pgfqpoint{2.600818in}{1.138767in}}%
\pgfpathlineto{\pgfqpoint{2.597379in}{1.134492in}}%
\pgfpathlineto{\pgfqpoint{2.605119in}{1.128246in}}%
\pgfpathlineto{\pgfqpoint{2.607917in}{1.118107in}}%
\pgfpathlineto{\pgfqpoint{2.593730in}{1.126503in}}%
\pgfpathlineto{\pgfqpoint{2.584752in}{1.125412in}}%
\pgfpathlineto{\pgfqpoint{2.577780in}{1.134049in}}%
\pgfpathlineto{\pgfqpoint{2.570682in}{1.134890in}}%
\pgfpathlineto{\pgfqpoint{2.560594in}{1.130564in}}%
\pgfpathlineto{\pgfqpoint{2.556681in}{1.135955in}}%
\pgfpathlineto{\pgfqpoint{2.561823in}{1.147270in}}%
\pgfpathlineto{\pgfqpoint{2.564161in}{1.156689in}}%
\pgfpathlineto{\pgfqpoint{2.569331in}{1.164244in}}%
\pgfpathlineto{\pgfqpoint{2.561780in}{1.169343in}}%
\pgfpathlineto{\pgfqpoint{2.558697in}{1.164803in}}%
\pgfpathlineto{\pgfqpoint{2.551154in}{1.169411in}}%
\pgfpathlineto{\pgfqpoint{2.537796in}{1.172482in}}%
\pgfpathlineto{\pgfqpoint{2.539008in}{1.179234in}}%
\pgfpathlineto{\pgfqpoint{2.532953in}{1.183150in}}%
\pgfpathlineto{\pgfqpoint{2.525158in}{1.182877in}}%
\pgfpathlineto{\pgfqpoint{2.515712in}{1.199224in}}%
\pgfpathlineto{\pgfqpoint{2.509119in}{1.197766in}}%
\pgfpathlineto{\pgfqpoint{2.499098in}{1.201985in}}%
\pgfpathlineto{\pgfqpoint{2.491891in}{1.197195in}}%
\pgfpathlineto{\pgfqpoint{2.486925in}{1.188584in}}%
\pgfpathlineto{\pgfqpoint{2.479399in}{1.187846in}}%
\pgfpathlineto{\pgfqpoint{2.470767in}{1.191730in}}%
\pgfpathlineto{\pgfqpoint{2.464536in}{1.179343in}}%
\pgfpathlineto{\pgfqpoint{2.460300in}{1.181474in}}%
\pgfpathlineto{\pgfqpoint{2.442325in}{1.160097in}}%
\pgfpathlineto{\pgfqpoint{2.437178in}{1.192485in}}%
\pgfpathlineto{\pgfqpoint{2.502868in}{1.204659in}}%
\pgfpathlineto{\pgfqpoint{2.532316in}{1.209933in}}%
\pgfpathlineto{\pgfqpoint{2.575158in}{1.218561in}}%
\pgfpathlineto{\pgfqpoint{2.611873in}{1.226505in}}%
\pgfpathlineto{\pgfqpoint{2.617245in}{1.205802in}}%
\pgfpathlineto{\pgfqpoint{2.632891in}{1.150158in}}%
\pgfpathlineto{\pgfqpoint{2.662135in}{1.155756in}}%
\pgfpathclose%
\pgfusepath{fill}%
\end{pgfscope}%
\begin{pgfscope}%
\pgfpathrectangle{\pgfqpoint{0.100000in}{0.100000in}}{\pgfqpoint{2.862985in}{1.832778in}}%
\pgfusepath{clip}%
\pgfsetbuttcap%
\pgfsetmiterjoin%
\definecolor{currentfill}{rgb}{0.453677,0.002755,0.660310}%
\pgfsetfillcolor{currentfill}%
\pgfsetlinewidth{0.000000pt}%
\definecolor{currentstroke}{rgb}{0.000000,0.000000,0.000000}%
\pgfsetstrokecolor{currentstroke}%
\pgfsetstrokeopacity{0.000000}%
\pgfsetdash{}{0pt}%
\pgfpathmoveto{\pgfqpoint{1.339821in}{0.971059in}}%
\pgfpathlineto{\pgfqpoint{1.291895in}{0.975633in}}%
\pgfpathlineto{\pgfqpoint{1.242175in}{0.980039in}}%
\pgfpathlineto{\pgfqpoint{1.184722in}{0.986032in}}%
\pgfpathlineto{\pgfqpoint{1.099362in}{0.996184in}}%
\pgfpathlineto{\pgfqpoint{1.069080in}{1.000850in}}%
\pgfpathlineto{\pgfqpoint{0.990915in}{1.012202in}}%
\pgfpathlineto{\pgfqpoint{1.002328in}{1.083837in}}%
\pgfpathlineto{\pgfqpoint{1.002613in}{1.089632in}}%
\pgfpathlineto{\pgfqpoint{1.013608in}{1.158753in}}%
\pgfpathlineto{\pgfqpoint{1.021820in}{1.211396in}}%
\pgfpathlineto{\pgfqpoint{1.029571in}{1.260214in}}%
\pgfpathlineto{\pgfqpoint{1.082522in}{1.252613in}}%
\pgfpathlineto{\pgfqpoint{1.131892in}{1.245527in}}%
\pgfpathlineto{\pgfqpoint{1.222647in}{1.234326in}}%
\pgfpathlineto{\pgfqpoint{1.264283in}{1.230524in}}%
\pgfpathlineto{\pgfqpoint{1.330243in}{1.224124in}}%
\pgfpathlineto{\pgfqpoint{1.358777in}{1.221808in}}%
\pgfpathlineto{\pgfqpoint{1.353739in}{1.159333in}}%
\pgfpathlineto{\pgfqpoint{1.349188in}{1.099180in}}%
\pgfpathlineto{\pgfqpoint{1.345525in}{1.050201in}}%
\pgfpathlineto{\pgfqpoint{1.339821in}{0.971059in}}%
\pgfpathclose%
\pgfusepath{fill}%
\end{pgfscope}%
\begin{pgfscope}%
\pgfpathrectangle{\pgfqpoint{0.100000in}{0.100000in}}{\pgfqpoint{2.862985in}{1.832778in}}%
\pgfusepath{clip}%
\pgfsetbuttcap%
\pgfsetmiterjoin%
\definecolor{currentfill}{rgb}{0.875376,0.383347,0.389976}%
\pgfsetfillcolor{currentfill}%
\pgfsetlinewidth{0.000000pt}%
\definecolor{currentstroke}{rgb}{0.000000,0.000000,0.000000}%
\pgfsetstrokecolor{currentstroke}%
\pgfsetstrokeopacity{0.000000}%
\pgfsetdash{}{0pt}%
\pgfpathmoveto{\pgfqpoint{1.973847in}{0.935544in}}%
\pgfpathlineto{\pgfqpoint{1.975463in}{0.942786in}}%
\pgfpathlineto{\pgfqpoint{1.983837in}{0.941156in}}%
\pgfpathlineto{\pgfqpoint{1.986225in}{0.966856in}}%
\pgfpathlineto{\pgfqpoint{1.983934in}{0.971495in}}%
\pgfpathlineto{\pgfqpoint{1.992101in}{0.981825in}}%
\pgfpathlineto{\pgfqpoint{1.998652in}{0.981141in}}%
\pgfpathlineto{\pgfqpoint{2.018506in}{0.974346in}}%
\pgfpathlineto{\pgfqpoint{2.020833in}{0.979847in}}%
\pgfpathlineto{\pgfqpoint{2.015777in}{0.988286in}}%
\pgfpathlineto{\pgfqpoint{2.017483in}{0.994944in}}%
\pgfpathlineto{\pgfqpoint{2.035999in}{1.001067in}}%
\pgfpathlineto{\pgfqpoint{2.031580in}{1.012333in}}%
\pgfpathlineto{\pgfqpoint{2.037013in}{1.021868in}}%
\pgfpathlineto{\pgfqpoint{2.047127in}{1.027220in}}%
\pgfpathlineto{\pgfqpoint{2.068349in}{1.032520in}}%
\pgfpathlineto{\pgfqpoint{2.081739in}{1.024520in}}%
\pgfpathlineto{\pgfqpoint{2.085801in}{1.032321in}}%
\pgfpathlineto{\pgfqpoint{2.096919in}{1.037771in}}%
\pgfpathlineto{\pgfqpoint{2.099789in}{1.032965in}}%
\pgfpathlineto{\pgfqpoint{2.111209in}{1.036783in}}%
\pgfpathlineto{\pgfqpoint{2.110489in}{1.043318in}}%
\pgfpathlineto{\pgfqpoint{2.120757in}{1.050797in}}%
\pgfpathlineto{\pgfqpoint{2.126817in}{1.043031in}}%
\pgfpathlineto{\pgfqpoint{2.134790in}{1.042250in}}%
\pgfpathlineto{\pgfqpoint{2.140088in}{1.047311in}}%
\pgfpathlineto{\pgfqpoint{2.139498in}{1.054509in}}%
\pgfpathlineto{\pgfqpoint{2.144002in}{1.061656in}}%
\pgfpathlineto{\pgfqpoint{2.150037in}{1.063205in}}%
\pgfpathlineto{\pgfqpoint{2.152460in}{1.072642in}}%
\pgfpathlineto{\pgfqpoint{2.161237in}{1.080801in}}%
\pgfpathlineto{\pgfqpoint{2.158595in}{1.088922in}}%
\pgfpathlineto{\pgfqpoint{2.167134in}{1.092965in}}%
\pgfpathlineto{\pgfqpoint{2.172839in}{1.090478in}}%
\pgfpathlineto{\pgfqpoint{2.181265in}{1.096795in}}%
\pgfpathlineto{\pgfqpoint{2.188795in}{1.098441in}}%
\pgfpathlineto{\pgfqpoint{2.188822in}{1.104993in}}%
\pgfpathlineto{\pgfqpoint{2.183545in}{1.114085in}}%
\pgfpathlineto{\pgfqpoint{2.186981in}{1.117748in}}%
\pgfpathlineto{\pgfqpoint{2.197773in}{1.116936in}}%
\pgfpathlineto{\pgfqpoint{2.204636in}{1.120355in}}%
\pgfpathlineto{\pgfqpoint{2.211715in}{1.115014in}}%
\pgfpathlineto{\pgfqpoint{2.217303in}{1.102959in}}%
\pgfpathlineto{\pgfqpoint{2.235827in}{1.100922in}}%
\pgfpathlineto{\pgfqpoint{2.245483in}{1.094286in}}%
\pgfpathlineto{\pgfqpoint{2.252277in}{1.099895in}}%
\pgfpathlineto{\pgfqpoint{2.264561in}{1.095181in}}%
\pgfpathlineto{\pgfqpoint{2.271515in}{1.097524in}}%
\pgfpathlineto{\pgfqpoint{2.276376in}{1.104185in}}%
\pgfpathlineto{\pgfqpoint{2.283796in}{1.105459in}}%
\pgfpathlineto{\pgfqpoint{2.285724in}{1.098217in}}%
\pgfpathlineto{\pgfqpoint{2.293930in}{1.095447in}}%
\pgfpathlineto{\pgfqpoint{2.300308in}{1.088305in}}%
\pgfpathlineto{\pgfqpoint{2.302231in}{1.078659in}}%
\pgfpathlineto{\pgfqpoint{2.299672in}{1.071367in}}%
\pgfpathlineto{\pgfqpoint{2.309465in}{1.061926in}}%
\pgfpathlineto{\pgfqpoint{2.320968in}{1.044098in}}%
\pgfpathlineto{\pgfqpoint{2.329483in}{1.038233in}}%
\pgfpathlineto{\pgfqpoint{2.338724in}{1.037893in}}%
\pgfpathlineto{\pgfqpoint{2.321661in}{1.018314in}}%
\pgfpathlineto{\pgfqpoint{2.304865in}{1.006463in}}%
\pgfpathlineto{\pgfqpoint{2.298683in}{0.997053in}}%
\pgfpathlineto{\pgfqpoint{2.298788in}{0.991945in}}%
\pgfpathlineto{\pgfqpoint{2.289691in}{0.988030in}}%
\pgfpathlineto{\pgfqpoint{2.287084in}{0.980665in}}%
\pgfpathlineto{\pgfqpoint{2.268132in}{0.973246in}}%
\pgfpathlineto{\pgfqpoint{2.261414in}{0.968427in}}%
\pgfpathlineto{\pgfqpoint{2.260505in}{0.967399in}}%
\pgfpathlineto{\pgfqpoint{2.205991in}{0.962230in}}%
\pgfpathlineto{\pgfqpoint{2.173070in}{0.959461in}}%
\pgfpathlineto{\pgfqpoint{2.119049in}{0.956404in}}%
\pgfpathlineto{\pgfqpoint{2.051748in}{0.949746in}}%
\pgfpathlineto{\pgfqpoint{2.040629in}{0.951286in}}%
\pgfpathlineto{\pgfqpoint{2.042942in}{0.939936in}}%
\pgfpathlineto{\pgfqpoint{1.973847in}{0.935544in}}%
\pgfpathclose%
\pgfusepath{fill}%
\end{pgfscope}%
\begin{pgfscope}%
\pgfpathrectangle{\pgfqpoint{0.100000in}{0.100000in}}{\pgfqpoint{2.862985in}{1.832778in}}%
\pgfusepath{clip}%
\pgfsetbuttcap%
\pgfsetmiterjoin%
\definecolor{currentfill}{rgb}{0.540570,0.034950,0.648640}%
\pgfsetfillcolor{currentfill}%
\pgfsetlinewidth{0.000000pt}%
\definecolor{currentstroke}{rgb}{0.000000,0.000000,0.000000}%
\pgfsetstrokecolor{currentstroke}%
\pgfsetstrokeopacity{0.000000}%
\pgfsetdash{}{0pt}%
\pgfpathmoveto{\pgfqpoint{1.339821in}{0.971059in}}%
\pgfpathlineto{\pgfqpoint{1.345525in}{1.050201in}}%
\pgfpathlineto{\pgfqpoint{1.349188in}{1.099180in}}%
\pgfpathlineto{\pgfqpoint{1.353739in}{1.159333in}}%
\pgfpathlineto{\pgfqpoint{1.416835in}{1.154903in}}%
\pgfpathlineto{\pgfqpoint{1.496966in}{1.150526in}}%
\pgfpathlineto{\pgfqpoint{1.551466in}{1.148434in}}%
\pgfpathlineto{\pgfqpoint{1.605647in}{1.146774in}}%
\pgfpathlineto{\pgfqpoint{1.654686in}{1.146016in}}%
\pgfpathlineto{\pgfqpoint{1.677367in}{1.146251in}}%
\pgfpathlineto{\pgfqpoint{1.687350in}{1.138104in}}%
\pgfpathlineto{\pgfqpoint{1.695170in}{1.139746in}}%
\pgfpathlineto{\pgfqpoint{1.695414in}{1.132640in}}%
\pgfpathlineto{\pgfqpoint{1.689607in}{1.120348in}}%
\pgfpathlineto{\pgfqpoint{1.690237in}{1.112582in}}%
\pgfpathlineto{\pgfqpoint{1.695797in}{1.107461in}}%
\pgfpathlineto{\pgfqpoint{1.700905in}{1.096790in}}%
\pgfpathlineto{\pgfqpoint{1.711265in}{1.090665in}}%
\pgfpathlineto{\pgfqpoint{1.710919in}{1.050444in}}%
\pgfpathlineto{\pgfqpoint{1.711224in}{0.957954in}}%
\pgfpathlineto{\pgfqpoint{1.641779in}{0.958274in}}%
\pgfpathlineto{\pgfqpoint{1.568651in}{0.959554in}}%
\pgfpathlineto{\pgfqpoint{1.514818in}{0.961453in}}%
\pgfpathlineto{\pgfqpoint{1.469882in}{0.963187in}}%
\pgfpathlineto{\pgfqpoint{1.394179in}{0.967651in}}%
\pgfpathlineto{\pgfqpoint{1.339821in}{0.971059in}}%
\pgfpathclose%
\pgfusepath{fill}%
\end{pgfscope}%
\begin{pgfscope}%
\pgfpathrectangle{\pgfqpoint{0.100000in}{0.100000in}}{\pgfqpoint{2.862985in}{1.832778in}}%
\pgfusepath{clip}%
\pgfsetbuttcap%
\pgfsetmiterjoin%
\definecolor{currentfill}{rgb}{0.540570,0.034950,0.648640}%
\pgfsetfillcolor{currentfill}%
\pgfsetlinewidth{0.000000pt}%
\definecolor{currentstroke}{rgb}{0.000000,0.000000,0.000000}%
\pgfsetstrokecolor{currentstroke}%
\pgfsetstrokeopacity{0.000000}%
\pgfsetdash{}{0pt}%
\pgfpathmoveto{\pgfqpoint{2.361621in}{0.981538in}}%
\pgfpathlineto{\pgfqpoint{2.284707in}{0.970734in}}%
\pgfpathlineto{\pgfqpoint{2.261414in}{0.968427in}}%
\pgfpathlineto{\pgfqpoint{2.268132in}{0.973246in}}%
\pgfpathlineto{\pgfqpoint{2.287084in}{0.980665in}}%
\pgfpathlineto{\pgfqpoint{2.289691in}{0.988030in}}%
\pgfpathlineto{\pgfqpoint{2.298788in}{0.991945in}}%
\pgfpathlineto{\pgfqpoint{2.298683in}{0.997053in}}%
\pgfpathlineto{\pgfqpoint{2.304865in}{1.006463in}}%
\pgfpathlineto{\pgfqpoint{2.321661in}{1.018314in}}%
\pgfpathlineto{\pgfqpoint{2.338724in}{1.037893in}}%
\pgfpathlineto{\pgfqpoint{2.337893in}{1.033679in}}%
\pgfpathlineto{\pgfqpoint{2.346890in}{1.023505in}}%
\pgfpathlineto{\pgfqpoint{2.361560in}{1.020647in}}%
\pgfpathlineto{\pgfqpoint{2.370094in}{1.030341in}}%
\pgfpathlineto{\pgfqpoint{2.377730in}{1.025236in}}%
\pgfpathlineto{\pgfqpoint{2.391393in}{1.031228in}}%
\pgfpathlineto{\pgfqpoint{2.393347in}{1.037079in}}%
\pgfpathlineto{\pgfqpoint{2.401826in}{1.038395in}}%
\pgfpathlineto{\pgfqpoint{2.410633in}{1.045308in}}%
\pgfpathlineto{\pgfqpoint{2.418396in}{1.047589in}}%
\pgfpathlineto{\pgfqpoint{2.418643in}{1.059800in}}%
\pgfpathlineto{\pgfqpoint{2.423669in}{1.072991in}}%
\pgfpathlineto{\pgfqpoint{2.430300in}{1.081251in}}%
\pgfpathlineto{\pgfqpoint{2.432449in}{1.093654in}}%
\pgfpathlineto{\pgfqpoint{2.441091in}{1.111076in}}%
\pgfpathlineto{\pgfqpoint{2.441399in}{1.121442in}}%
\pgfpathlineto{\pgfqpoint{2.450326in}{1.114688in}}%
\pgfpathlineto{\pgfqpoint{2.459053in}{1.113049in}}%
\pgfpathlineto{\pgfqpoint{2.462449in}{1.117963in}}%
\pgfpathlineto{\pgfqpoint{2.468517in}{1.139377in}}%
\pgfpathlineto{\pgfqpoint{2.476965in}{1.138556in}}%
\pgfpathlineto{\pgfqpoint{2.479203in}{1.146402in}}%
\pgfpathlineto{\pgfqpoint{2.484159in}{1.149941in}}%
\pgfpathlineto{\pgfqpoint{2.494659in}{1.167732in}}%
\pgfpathlineto{\pgfqpoint{2.493134in}{1.172996in}}%
\pgfpathlineto{\pgfqpoint{2.497755in}{1.184435in}}%
\pgfpathlineto{\pgfqpoint{2.522427in}{1.170500in}}%
\pgfpathlineto{\pgfqpoint{2.525158in}{1.182877in}}%
\pgfpathlineto{\pgfqpoint{2.532953in}{1.183150in}}%
\pgfpathlineto{\pgfqpoint{2.539008in}{1.179234in}}%
\pgfpathlineto{\pgfqpoint{2.537796in}{1.172482in}}%
\pgfpathlineto{\pgfqpoint{2.551154in}{1.169411in}}%
\pgfpathlineto{\pgfqpoint{2.558697in}{1.164803in}}%
\pgfpathlineto{\pgfqpoint{2.564161in}{1.156689in}}%
\pgfpathlineto{\pgfqpoint{2.561823in}{1.147270in}}%
\pgfpathlineto{\pgfqpoint{2.557103in}{1.146498in}}%
\pgfpathlineto{\pgfqpoint{2.554365in}{1.132285in}}%
\pgfpathlineto{\pgfqpoint{2.560362in}{1.126728in}}%
\pgfpathlineto{\pgfqpoint{2.568804in}{1.131220in}}%
\pgfpathlineto{\pgfqpoint{2.576637in}{1.121734in}}%
\pgfpathlineto{\pgfqpoint{2.594137in}{1.120029in}}%
\pgfpathlineto{\pgfqpoint{2.597152in}{1.114573in}}%
\pgfpathlineto{\pgfqpoint{2.613348in}{1.109266in}}%
\pgfpathlineto{\pgfqpoint{2.611349in}{1.103004in}}%
\pgfpathlineto{\pgfqpoint{2.613456in}{1.085305in}}%
\pgfpathlineto{\pgfqpoint{2.619993in}{1.078622in}}%
\pgfpathlineto{\pgfqpoint{2.610347in}{1.075802in}}%
\pgfpathlineto{\pgfqpoint{2.611626in}{1.068246in}}%
\pgfpathlineto{\pgfqpoint{2.621916in}{1.061840in}}%
\pgfpathlineto{\pgfqpoint{2.622844in}{1.055521in}}%
\pgfpathlineto{\pgfqpoint{2.617034in}{1.050772in}}%
\pgfpathlineto{\pgfqpoint{2.605312in}{1.062022in}}%
\pgfpathlineto{\pgfqpoint{2.604153in}{1.053814in}}%
\pgfpathlineto{\pgfqpoint{2.613966in}{1.049911in}}%
\pgfpathlineto{\pgfqpoint{2.614948in}{1.045875in}}%
\pgfpathlineto{\pgfqpoint{2.628408in}{1.051209in}}%
\pgfpathlineto{\pgfqpoint{2.638722in}{1.052620in}}%
\pgfpathlineto{\pgfqpoint{2.649287in}{1.031303in}}%
\pgfpathlineto{\pgfqpoint{2.648108in}{1.031073in}}%
\pgfpathlineto{\pgfqpoint{2.643334in}{1.030086in}}%
\pgfpathlineto{\pgfqpoint{2.641927in}{1.029792in}}%
\pgfpathlineto{\pgfqpoint{2.640997in}{1.029610in}}%
\pgfpathlineto{\pgfqpoint{2.578091in}{1.016561in}}%
\pgfpathlineto{\pgfqpoint{2.521826in}{1.005045in}}%
\pgfpathlineto{\pgfqpoint{2.444033in}{0.991644in}}%
\pgfpathlineto{\pgfqpoint{2.377964in}{0.982918in}}%
\pgfpathlineto{\pgfqpoint{2.361621in}{0.981538in}}%
\pgfpathclose%
\pgfusepath{fill}%
\end{pgfscope}%
\begin{pgfscope}%
\pgfpathrectangle{\pgfqpoint{0.100000in}{0.100000in}}{\pgfqpoint{2.862985in}{1.832778in}}%
\pgfusepath{clip}%
\pgfsetbuttcap%
\pgfsetmiterjoin%
\definecolor{currentfill}{rgb}{0.540570,0.034950,0.648640}%
\pgfsetfillcolor{currentfill}%
\pgfsetlinewidth{0.000000pt}%
\definecolor{currentstroke}{rgb}{0.000000,0.000000,0.000000}%
\pgfsetstrokecolor{currentstroke}%
\pgfsetstrokeopacity{0.000000}%
\pgfsetdash{}{0pt}%
\pgfpathmoveto{\pgfqpoint{2.641870in}{1.120469in}}%
\pgfpathlineto{\pgfqpoint{2.654551in}{1.126221in}}%
\pgfpathlineto{\pgfqpoint{2.644424in}{1.096596in}}%
\pgfpathlineto{\pgfqpoint{2.637727in}{1.080908in}}%
\pgfpathlineto{\pgfqpoint{2.632902in}{1.089786in}}%
\pgfpathlineto{\pgfqpoint{2.641482in}{1.111041in}}%
\pgfpathlineto{\pgfqpoint{2.641870in}{1.120469in}}%
\pgfpathclose%
\pgfusepath{fill}%
\end{pgfscope}%
\begin{pgfscope}%
\pgfpathrectangle{\pgfqpoint{0.100000in}{0.100000in}}{\pgfqpoint{2.862985in}{1.832778in}}%
\pgfusepath{clip}%
\pgfsetbuttcap%
\pgfsetmiterjoin%
\definecolor{currentfill}{rgb}{0.692840,0.165141,0.564522}%
\pgfsetfillcolor{currentfill}%
\pgfsetlinewidth{0.000000pt}%
\definecolor{currentstroke}{rgb}{0.000000,0.000000,0.000000}%
\pgfsetstrokecolor{currentstroke}%
\pgfsetstrokeopacity{0.000000}%
\pgfsetdash{}{0pt}%
\pgfpathmoveto{\pgfqpoint{1.973847in}{0.935544in}}%
\pgfpathlineto{\pgfqpoint{1.970785in}{0.935101in}}%
\pgfpathlineto{\pgfqpoint{1.967896in}{0.934898in}}%
\pgfpathlineto{\pgfqpoint{1.969133in}{0.926031in}}%
\pgfpathlineto{\pgfqpoint{1.961682in}{0.917093in}}%
\pgfpathlineto{\pgfqpoint{1.960213in}{0.903107in}}%
\pgfpathlineto{\pgfqpoint{1.926904in}{0.900644in}}%
\pgfpathlineto{\pgfqpoint{1.929808in}{0.907211in}}%
\pgfpathlineto{\pgfqpoint{1.941818in}{0.919204in}}%
\pgfpathlineto{\pgfqpoint{1.942150in}{0.926155in}}%
\pgfpathlineto{\pgfqpoint{1.936833in}{0.932737in}}%
\pgfpathlineto{\pgfqpoint{1.887257in}{0.930116in}}%
\pgfpathlineto{\pgfqpoint{1.800580in}{0.927364in}}%
\pgfpathlineto{\pgfqpoint{1.749856in}{0.926440in}}%
\pgfpathlineto{\pgfqpoint{1.711515in}{0.926094in}}%
\pgfpathlineto{\pgfqpoint{1.711224in}{0.957954in}}%
\pgfpathlineto{\pgfqpoint{1.710919in}{1.050444in}}%
\pgfpathlineto{\pgfqpoint{1.711265in}{1.090665in}}%
\pgfpathlineto{\pgfqpoint{1.700905in}{1.096790in}}%
\pgfpathlineto{\pgfqpoint{1.695797in}{1.107461in}}%
\pgfpathlineto{\pgfqpoint{1.690237in}{1.112582in}}%
\pgfpathlineto{\pgfqpoint{1.689607in}{1.120348in}}%
\pgfpathlineto{\pgfqpoint{1.695414in}{1.132640in}}%
\pgfpathlineto{\pgfqpoint{1.695170in}{1.139746in}}%
\pgfpathlineto{\pgfqpoint{1.687350in}{1.138104in}}%
\pgfpathlineto{\pgfqpoint{1.677367in}{1.146251in}}%
\pgfpathlineto{\pgfqpoint{1.669364in}{1.160551in}}%
\pgfpathlineto{\pgfqpoint{1.662651in}{1.167151in}}%
\pgfpathlineto{\pgfqpoint{1.655637in}{1.183365in}}%
\pgfpathlineto{\pgfqpoint{1.728489in}{1.182229in}}%
\pgfpathlineto{\pgfqpoint{1.800906in}{1.184615in}}%
\pgfpathlineto{\pgfqpoint{1.847339in}{1.187293in}}%
\pgfpathlineto{\pgfqpoint{1.861865in}{1.172908in}}%
\pgfpathlineto{\pgfqpoint{1.859210in}{1.157906in}}%
\pgfpathlineto{\pgfqpoint{1.863083in}{1.141015in}}%
\pgfpathlineto{\pgfqpoint{1.867379in}{1.132498in}}%
\pgfpathlineto{\pgfqpoint{1.881087in}{1.120782in}}%
\pgfpathlineto{\pgfqpoint{1.884086in}{1.115607in}}%
\pgfpathlineto{\pgfqpoint{1.899094in}{1.103945in}}%
\pgfpathlineto{\pgfqpoint{1.901269in}{1.090863in}}%
\pgfpathlineto{\pgfqpoint{1.905579in}{1.081212in}}%
\pgfpathlineto{\pgfqpoint{1.915763in}{1.086834in}}%
\pgfpathlineto{\pgfqpoint{1.924157in}{1.084263in}}%
\pgfpathlineto{\pgfqpoint{1.930495in}{1.077440in}}%
\pgfpathlineto{\pgfqpoint{1.928102in}{1.065365in}}%
\pgfpathlineto{\pgfqpoint{1.920322in}{1.049578in}}%
\pgfpathlineto{\pgfqpoint{1.920442in}{1.041355in}}%
\pgfpathlineto{\pgfqpoint{1.926256in}{1.034614in}}%
\pgfpathlineto{\pgfqpoint{1.938768in}{1.025684in}}%
\pgfpathlineto{\pgfqpoint{1.953174in}{1.018436in}}%
\pgfpathlineto{\pgfqpoint{1.956591in}{1.012632in}}%
\pgfpathlineto{\pgfqpoint{1.963959in}{1.009888in}}%
\pgfpathlineto{\pgfqpoint{1.964252in}{1.002186in}}%
\pgfpathlineto{\pgfqpoint{1.969705in}{0.992402in}}%
\pgfpathlineto{\pgfqpoint{1.965922in}{0.983711in}}%
\pgfpathlineto{\pgfqpoint{1.973328in}{0.969882in}}%
\pgfpathlineto{\pgfqpoint{1.979121in}{0.972628in}}%
\pgfpathlineto{\pgfqpoint{1.986225in}{0.966856in}}%
\pgfpathlineto{\pgfqpoint{1.983837in}{0.941156in}}%
\pgfpathlineto{\pgfqpoint{1.975463in}{0.942786in}}%
\pgfpathlineto{\pgfqpoint{1.973847in}{0.935544in}}%
\pgfpathclose%
\pgfusepath{fill}%
\end{pgfscope}%
\begin{pgfscope}%
\pgfpathrectangle{\pgfqpoint{0.100000in}{0.100000in}}{\pgfqpoint{2.862985in}{1.832778in}}%
\pgfusepath{clip}%
\pgfsetbuttcap%
\pgfsetmiterjoin%
\definecolor{currentfill}{rgb}{0.719181,0.191729,0.542663}%
\pgfsetfillcolor{currentfill}%
\pgfsetlinewidth{0.000000pt}%
\definecolor{currentstroke}{rgb}{0.000000,0.000000,0.000000}%
\pgfsetstrokecolor{currentstroke}%
\pgfsetstrokeopacity{0.000000}%
\pgfsetdash{}{0pt}%
\pgfpathmoveto{\pgfqpoint{0.689569in}{0.941085in}}%
\pgfpathlineto{\pgfqpoint{0.695633in}{0.953995in}}%
\pgfpathlineto{\pgfqpoint{0.693973in}{0.973374in}}%
\pgfpathlineto{\pgfqpoint{0.696877in}{0.978886in}}%
\pgfpathlineto{\pgfqpoint{0.697164in}{1.003024in}}%
\pgfpathlineto{\pgfqpoint{0.699916in}{1.009976in}}%
\pgfpathlineto{\pgfqpoint{0.709613in}{1.011057in}}%
\pgfpathlineto{\pgfqpoint{0.718612in}{1.007971in}}%
\pgfpathlineto{\pgfqpoint{0.722531in}{0.999473in}}%
\pgfpathlineto{\pgfqpoint{0.728025in}{0.999798in}}%
\pgfpathlineto{\pgfqpoint{0.734852in}{1.009534in}}%
\pgfpathlineto{\pgfqpoint{0.744719in}{1.057523in}}%
\pgfpathlineto{\pgfqpoint{0.800868in}{1.045930in}}%
\pgfpathlineto{\pgfqpoint{0.833455in}{1.039504in}}%
\pgfpathlineto{\pgfqpoint{0.919991in}{1.024292in}}%
\pgfpathlineto{\pgfqpoint{0.943932in}{1.019484in}}%
\pgfpathlineto{\pgfqpoint{0.990915in}{1.012202in}}%
\pgfpathlineto{\pgfqpoint{0.981281in}{0.950172in}}%
\pgfpathlineto{\pgfqpoint{0.971262in}{0.885469in}}%
\pgfpathlineto{\pgfqpoint{0.959723in}{0.812676in}}%
\pgfpathlineto{\pgfqpoint{0.946747in}{0.729118in}}%
\pgfpathlineto{\pgfqpoint{0.936264in}{0.660490in}}%
\pgfpathlineto{\pgfqpoint{0.861176in}{0.672414in}}%
\pgfpathlineto{\pgfqpoint{0.828185in}{0.678064in}}%
\pgfpathlineto{\pgfqpoint{0.813447in}{0.686887in}}%
\pgfpathlineto{\pgfqpoint{0.717356in}{0.744864in}}%
\pgfpathlineto{\pgfqpoint{0.645244in}{0.788862in}}%
\pgfpathlineto{\pgfqpoint{0.647646in}{0.796660in}}%
\pgfpathlineto{\pgfqpoint{0.653603in}{0.802117in}}%
\pgfpathlineto{\pgfqpoint{0.659798in}{0.801105in}}%
\pgfpathlineto{\pgfqpoint{0.668785in}{0.806836in}}%
\pgfpathlineto{\pgfqpoint{0.670200in}{0.815071in}}%
\pgfpathlineto{\pgfqpoint{0.659240in}{0.825061in}}%
\pgfpathlineto{\pgfqpoint{0.667035in}{0.844239in}}%
\pgfpathlineto{\pgfqpoint{0.674871in}{0.851614in}}%
\pgfpathlineto{\pgfqpoint{0.678591in}{0.860364in}}%
\pgfpathlineto{\pgfqpoint{0.680882in}{0.876407in}}%
\pgfpathlineto{\pgfqpoint{0.688230in}{0.883673in}}%
\pgfpathlineto{\pgfqpoint{0.703654in}{0.890930in}}%
\pgfpathlineto{\pgfqpoint{0.703313in}{0.897324in}}%
\pgfpathlineto{\pgfqpoint{0.694699in}{0.905259in}}%
\pgfpathlineto{\pgfqpoint{0.693526in}{0.921617in}}%
\pgfpathlineto{\pgfqpoint{0.687595in}{0.933576in}}%
\pgfpathlineto{\pgfqpoint{0.689569in}{0.941085in}}%
\pgfpathclose%
\pgfusepath{fill}%
\end{pgfscope}%
\begin{pgfscope}%
\pgfpathrectangle{\pgfqpoint{0.100000in}{0.100000in}}{\pgfqpoint{2.862985in}{1.832778in}}%
\pgfusepath{clip}%
\pgfsetbuttcap%
\pgfsetmiterjoin%
\definecolor{currentfill}{rgb}{0.697324,0.169573,0.560919}%
\pgfsetfillcolor{currentfill}%
\pgfsetlinewidth{0.000000pt}%
\definecolor{currentstroke}{rgb}{0.000000,0.000000,0.000000}%
\pgfsetstrokecolor{currentstroke}%
\pgfsetstrokeopacity{0.000000}%
\pgfsetdash{}{0pt}%
\pgfpathmoveto{\pgfqpoint{1.719643in}{0.746495in}}%
\pgfpathlineto{\pgfqpoint{1.705724in}{0.750800in}}%
\pgfpathlineto{\pgfqpoint{1.687888in}{0.764459in}}%
\pgfpathlineto{\pgfqpoint{1.679966in}{0.767394in}}%
\pgfpathlineto{\pgfqpoint{1.674932in}{0.761495in}}%
\pgfpathlineto{\pgfqpoint{1.668575in}{0.761197in}}%
\pgfpathlineto{\pgfqpoint{1.660532in}{0.766205in}}%
\pgfpathlineto{\pgfqpoint{1.647911in}{0.759788in}}%
\pgfpathlineto{\pgfqpoint{1.643418in}{0.762948in}}%
\pgfpathlineto{\pgfqpoint{1.630999in}{0.755462in}}%
\pgfpathlineto{\pgfqpoint{1.617897in}{0.756843in}}%
\pgfpathlineto{\pgfqpoint{1.607791in}{0.761708in}}%
\pgfpathlineto{\pgfqpoint{1.600714in}{0.759877in}}%
\pgfpathlineto{\pgfqpoint{1.589386in}{0.766758in}}%
\pgfpathlineto{\pgfqpoint{1.583084in}{0.754622in}}%
\pgfpathlineto{\pgfqpoint{1.576640in}{0.765059in}}%
\pgfpathlineto{\pgfqpoint{1.563911in}{0.761010in}}%
\pgfpathlineto{\pgfqpoint{1.552778in}{0.770922in}}%
\pgfpathlineto{\pgfqpoint{1.543049in}{0.762933in}}%
\pgfpathlineto{\pgfqpoint{1.538177in}{0.770275in}}%
\pgfpathlineto{\pgfqpoint{1.531154in}{0.772659in}}%
\pgfpathlineto{\pgfqpoint{1.526877in}{0.779738in}}%
\pgfpathlineto{\pgfqpoint{1.517692in}{0.781753in}}%
\pgfpathlineto{\pgfqpoint{1.512392in}{0.776425in}}%
\pgfpathlineto{\pgfqpoint{1.503387in}{0.783330in}}%
\pgfpathlineto{\pgfqpoint{1.499189in}{0.781753in}}%
\pgfpathlineto{\pgfqpoint{1.484265in}{0.787349in}}%
\pgfpathlineto{\pgfqpoint{1.474917in}{0.787973in}}%
\pgfpathlineto{\pgfqpoint{1.470733in}{0.799856in}}%
\pgfpathlineto{\pgfqpoint{1.463240in}{0.798376in}}%
\pgfpathlineto{\pgfqpoint{1.454657in}{0.801308in}}%
\pgfpathlineto{\pgfqpoint{1.449005in}{0.799614in}}%
\pgfpathlineto{\pgfqpoint{1.436884in}{0.812959in}}%
\pgfpathlineto{\pgfqpoint{1.433507in}{0.812086in}}%
\pgfpathlineto{\pgfqpoint{1.436574in}{0.866207in}}%
\pgfpathlineto{\pgfqpoint{1.439916in}{0.933195in}}%
\pgfpathlineto{\pgfqpoint{1.385064in}{0.936263in}}%
\pgfpathlineto{\pgfqpoint{1.330936in}{0.940356in}}%
\pgfpathlineto{\pgfqpoint{1.289118in}{0.943987in}}%
\pgfpathlineto{\pgfqpoint{1.291895in}{0.975633in}}%
\pgfpathlineto{\pgfqpoint{1.339821in}{0.971059in}}%
\pgfpathlineto{\pgfqpoint{1.394179in}{0.967651in}}%
\pgfpathlineto{\pgfqpoint{1.469882in}{0.963187in}}%
\pgfpathlineto{\pgfqpoint{1.514818in}{0.961453in}}%
\pgfpathlineto{\pgfqpoint{1.568651in}{0.959554in}}%
\pgfpathlineto{\pgfqpoint{1.641779in}{0.958274in}}%
\pgfpathlineto{\pgfqpoint{1.711224in}{0.957954in}}%
\pgfpathlineto{\pgfqpoint{1.711515in}{0.926094in}}%
\pgfpathlineto{\pgfqpoint{1.715413in}{0.902092in}}%
\pgfpathlineto{\pgfqpoint{1.721468in}{0.857760in}}%
\pgfpathlineto{\pgfqpoint{1.720219in}{0.782052in}}%
\pgfpathlineto{\pgfqpoint{1.719643in}{0.746495in}}%
\pgfpathclose%
\pgfusepath{fill}%
\end{pgfscope}%
\begin{pgfscope}%
\pgfpathrectangle{\pgfqpoint{0.100000in}{0.100000in}}{\pgfqpoint{2.862985in}{1.832778in}}%
\pgfusepath{clip}%
\pgfsetbuttcap%
\pgfsetmiterjoin%
\definecolor{currentfill}{rgb}{0.768090,0.244817,0.498465}%
\pgfsetfillcolor{currentfill}%
\pgfsetlinewidth{0.000000pt}%
\definecolor{currentstroke}{rgb}{0.000000,0.000000,0.000000}%
\pgfsetstrokecolor{currentstroke}%
\pgfsetstrokeopacity{0.000000}%
\pgfsetdash{}{0pt}%
\pgfpathmoveto{\pgfqpoint{2.240820in}{0.864070in}}%
\pgfpathlineto{\pgfqpoint{2.240861in}{0.878097in}}%
\pgfpathlineto{\pgfqpoint{2.253051in}{0.883484in}}%
\pgfpathlineto{\pgfqpoint{2.253562in}{0.892095in}}%
\pgfpathlineto{\pgfqpoint{2.264476in}{0.902614in}}%
\pgfpathlineto{\pgfqpoint{2.278119in}{0.904768in}}%
\pgfpathlineto{\pgfqpoint{2.290113in}{0.916361in}}%
\pgfpathlineto{\pgfqpoint{2.302639in}{0.921543in}}%
\pgfpathlineto{\pgfqpoint{2.312040in}{0.937098in}}%
\pgfpathlineto{\pgfqpoint{2.320602in}{0.935969in}}%
\pgfpathlineto{\pgfqpoint{2.338697in}{0.950142in}}%
\pgfpathlineto{\pgfqpoint{2.348253in}{0.950409in}}%
\pgfpathlineto{\pgfqpoint{2.352288in}{0.961208in}}%
\pgfpathlineto{\pgfqpoint{2.359883in}{0.968727in}}%
\pgfpathlineto{\pgfqpoint{2.361621in}{0.981538in}}%
\pgfpathlineto{\pgfqpoint{2.377964in}{0.982918in}}%
\pgfpathlineto{\pgfqpoint{2.444033in}{0.991644in}}%
\pgfpathlineto{\pgfqpoint{2.521826in}{1.005045in}}%
\pgfpathlineto{\pgfqpoint{2.578091in}{1.016561in}}%
\pgfpathlineto{\pgfqpoint{2.640997in}{1.029610in}}%
\pgfpathlineto{\pgfqpoint{2.658973in}{1.004931in}}%
\pgfpathlineto{\pgfqpoint{2.649236in}{1.006506in}}%
\pgfpathlineto{\pgfqpoint{2.636737in}{1.003461in}}%
\pgfpathlineto{\pgfqpoint{2.624421in}{0.990884in}}%
\pgfpathlineto{\pgfqpoint{2.615570in}{0.991789in}}%
\pgfpathlineto{\pgfqpoint{2.614450in}{0.984319in}}%
\pgfpathlineto{\pgfqpoint{2.630449in}{0.990226in}}%
\pgfpathlineto{\pgfqpoint{2.646548in}{0.992662in}}%
\pgfpathlineto{\pgfqpoint{2.652526in}{0.989421in}}%
\pgfpathlineto{\pgfqpoint{2.660571in}{0.993116in}}%
\pgfpathlineto{\pgfqpoint{2.664722in}{0.990527in}}%
\pgfpathlineto{\pgfqpoint{2.668396in}{0.978199in}}%
\pgfpathlineto{\pgfqpoint{2.660750in}{0.974382in}}%
\pgfpathlineto{\pgfqpoint{2.655502in}{0.959350in}}%
\pgfpathlineto{\pgfqpoint{2.650000in}{0.953497in}}%
\pgfpathlineto{\pgfqpoint{2.633132in}{0.954777in}}%
\pgfpathlineto{\pgfqpoint{2.634039in}{0.963604in}}%
\pgfpathlineto{\pgfqpoint{2.624827in}{0.959748in}}%
\pgfpathlineto{\pgfqpoint{2.622823in}{0.952387in}}%
\pgfpathlineto{\pgfqpoint{2.629137in}{0.944761in}}%
\pgfpathlineto{\pgfqpoint{2.631282in}{0.931819in}}%
\pgfpathlineto{\pgfqpoint{2.620950in}{0.924444in}}%
\pgfpathlineto{\pgfqpoint{2.632125in}{0.921851in}}%
\pgfpathlineto{\pgfqpoint{2.637181in}{0.927273in}}%
\pgfpathlineto{\pgfqpoint{2.648362in}{0.927934in}}%
\pgfpathlineto{\pgfqpoint{2.642616in}{0.916224in}}%
\pgfpathlineto{\pgfqpoint{2.635579in}{0.910582in}}%
\pgfpathlineto{\pgfqpoint{2.614292in}{0.904933in}}%
\pgfpathlineto{\pgfqpoint{2.592436in}{0.885111in}}%
\pgfpathlineto{\pgfqpoint{2.578960in}{0.865580in}}%
\pgfpathlineto{\pgfqpoint{2.578880in}{0.857648in}}%
\pgfpathlineto{\pgfqpoint{2.573497in}{0.846702in}}%
\pgfpathlineto{\pgfqpoint{2.545869in}{0.839501in}}%
\pgfpathlineto{\pgfqpoint{2.480165in}{0.886600in}}%
\pgfpathlineto{\pgfqpoint{2.422304in}{0.878098in}}%
\pgfpathlineto{\pgfqpoint{2.421819in}{0.885949in}}%
\pgfpathlineto{\pgfqpoint{2.413024in}{0.894789in}}%
\pgfpathlineto{\pgfqpoint{2.406366in}{0.896952in}}%
\pgfpathlineto{\pgfqpoint{2.343499in}{0.890420in}}%
\pgfpathlineto{\pgfqpoint{2.329049in}{0.885519in}}%
\pgfpathlineto{\pgfqpoint{2.302958in}{0.872540in}}%
\pgfpathlineto{\pgfqpoint{2.280410in}{0.868948in}}%
\pgfpathlineto{\pgfqpoint{2.240820in}{0.864070in}}%
\pgfpathclose%
\pgfusepath{fill}%
\end{pgfscope}%
\begin{pgfscope}%
\pgfpathrectangle{\pgfqpoint{0.100000in}{0.100000in}}{\pgfqpoint{2.862985in}{1.832778in}}%
\pgfusepath{clip}%
\pgfsetbuttcap%
\pgfsetmiterjoin%
\definecolor{currentfill}{rgb}{0.805467,0.289057,0.462415}%
\pgfsetfillcolor{currentfill}%
\pgfsetlinewidth{0.000000pt}%
\definecolor{currentstroke}{rgb}{0.000000,0.000000,0.000000}%
\pgfsetstrokecolor{currentstroke}%
\pgfsetstrokeopacity{0.000000}%
\pgfsetdash{}{0pt}%
\pgfpathmoveto{\pgfqpoint{2.240820in}{0.864070in}}%
\pgfpathlineto{\pgfqpoint{2.175022in}{0.856843in}}%
\pgfpathlineto{\pgfqpoint{2.114830in}{0.851430in}}%
\pgfpathlineto{\pgfqpoint{2.052965in}{0.847429in}}%
\pgfpathlineto{\pgfqpoint{2.042346in}{0.845890in}}%
\pgfpathlineto{\pgfqpoint{2.000643in}{0.842683in}}%
\pgfpathlineto{\pgfqpoint{1.933848in}{0.838785in}}%
\pgfpathlineto{\pgfqpoint{1.945732in}{0.848649in}}%
\pgfpathlineto{\pgfqpoint{1.943020in}{0.859033in}}%
\pgfpathlineto{\pgfqpoint{1.945647in}{0.871619in}}%
\pgfpathlineto{\pgfqpoint{1.949660in}{0.877545in}}%
\pgfpathlineto{\pgfqpoint{1.949509in}{0.885789in}}%
\pgfpathlineto{\pgfqpoint{1.960195in}{0.890974in}}%
\pgfpathlineto{\pgfqpoint{1.960213in}{0.903107in}}%
\pgfpathlineto{\pgfqpoint{1.961682in}{0.917093in}}%
\pgfpathlineto{\pgfqpoint{1.969133in}{0.926031in}}%
\pgfpathlineto{\pgfqpoint{1.967896in}{0.934898in}}%
\pgfpathlineto{\pgfqpoint{1.970785in}{0.935101in}}%
\pgfpathlineto{\pgfqpoint{1.973847in}{0.935544in}}%
\pgfpathlineto{\pgfqpoint{2.042942in}{0.939936in}}%
\pgfpathlineto{\pgfqpoint{2.040629in}{0.951286in}}%
\pgfpathlineto{\pgfqpoint{2.051748in}{0.949746in}}%
\pgfpathlineto{\pgfqpoint{2.119049in}{0.956404in}}%
\pgfpathlineto{\pgfqpoint{2.173070in}{0.959461in}}%
\pgfpathlineto{\pgfqpoint{2.205991in}{0.962230in}}%
\pgfpathlineto{\pgfqpoint{2.260505in}{0.967399in}}%
\pgfpathlineto{\pgfqpoint{2.261414in}{0.968427in}}%
\pgfpathlineto{\pgfqpoint{2.284707in}{0.970734in}}%
\pgfpathlineto{\pgfqpoint{2.361621in}{0.981538in}}%
\pgfpathlineto{\pgfqpoint{2.359883in}{0.968727in}}%
\pgfpathlineto{\pgfqpoint{2.352288in}{0.961208in}}%
\pgfpathlineto{\pgfqpoint{2.348253in}{0.950409in}}%
\pgfpathlineto{\pgfqpoint{2.338697in}{0.950142in}}%
\pgfpathlineto{\pgfqpoint{2.320602in}{0.935969in}}%
\pgfpathlineto{\pgfqpoint{2.312040in}{0.937098in}}%
\pgfpathlineto{\pgfqpoint{2.302639in}{0.921543in}}%
\pgfpathlineto{\pgfqpoint{2.290113in}{0.916361in}}%
\pgfpathlineto{\pgfqpoint{2.278119in}{0.904768in}}%
\pgfpathlineto{\pgfqpoint{2.264476in}{0.902614in}}%
\pgfpathlineto{\pgfqpoint{2.253562in}{0.892095in}}%
\pgfpathlineto{\pgfqpoint{2.253051in}{0.883484in}}%
\pgfpathlineto{\pgfqpoint{2.240861in}{0.878097in}}%
\pgfpathlineto{\pgfqpoint{2.240820in}{0.864070in}}%
\pgfpathclose%
\pgfusepath{fill}%
\end{pgfscope}%
\begin{pgfscope}%
\pgfpathrectangle{\pgfqpoint{0.100000in}{0.100000in}}{\pgfqpoint{2.862985in}{1.832778in}}%
\pgfusepath{clip}%
\pgfsetbuttcap%
\pgfsetmiterjoin%
\definecolor{currentfill}{rgb}{0.645872,0.120898,0.598867}%
\pgfsetfillcolor{currentfill}%
\pgfsetlinewidth{0.000000pt}%
\definecolor{currentstroke}{rgb}{0.000000,0.000000,0.000000}%
\pgfsetstrokecolor{currentstroke}%
\pgfsetstrokeopacity{0.000000}%
\pgfsetdash{}{0pt}%
\pgfpathmoveto{\pgfqpoint{1.073883in}{0.669271in}}%
\pgfpathlineto{\pgfqpoint{1.069190in}{0.676811in}}%
\pgfpathlineto{\pgfqpoint{1.071132in}{0.683307in}}%
\pgfpathlineto{\pgfqpoint{1.104118in}{0.679215in}}%
\pgfpathlineto{\pgfqpoint{1.165452in}{0.672249in}}%
\pgfpathlineto{\pgfqpoint{1.209801in}{0.667891in}}%
\pgfpathlineto{\pgfqpoint{1.261021in}{0.662722in}}%
\pgfpathlineto{\pgfqpoint{1.263807in}{0.695029in}}%
\pgfpathlineto{\pgfqpoint{1.272520in}{0.777201in}}%
\pgfpathlineto{\pgfqpoint{1.278172in}{0.834767in}}%
\pgfpathlineto{\pgfqpoint{1.283057in}{0.889756in}}%
\pgfpathlineto{\pgfqpoint{1.287612in}{0.944071in}}%
\pgfpathlineto{\pgfqpoint{1.289118in}{0.943987in}}%
\pgfpathlineto{\pgfqpoint{1.330936in}{0.940356in}}%
\pgfpathlineto{\pgfqpoint{1.385064in}{0.936263in}}%
\pgfpathlineto{\pgfqpoint{1.439916in}{0.933195in}}%
\pgfpathlineto{\pgfqpoint{1.436574in}{0.866207in}}%
\pgfpathlineto{\pgfqpoint{1.433507in}{0.812086in}}%
\pgfpathlineto{\pgfqpoint{1.436884in}{0.812959in}}%
\pgfpathlineto{\pgfqpoint{1.449005in}{0.799614in}}%
\pgfpathlineto{\pgfqpoint{1.454657in}{0.801308in}}%
\pgfpathlineto{\pgfqpoint{1.463240in}{0.798376in}}%
\pgfpathlineto{\pgfqpoint{1.470733in}{0.799856in}}%
\pgfpathlineto{\pgfqpoint{1.474917in}{0.787973in}}%
\pgfpathlineto{\pgfqpoint{1.484265in}{0.787349in}}%
\pgfpathlineto{\pgfqpoint{1.499189in}{0.781753in}}%
\pgfpathlineto{\pgfqpoint{1.503387in}{0.783330in}}%
\pgfpathlineto{\pgfqpoint{1.512392in}{0.776425in}}%
\pgfpathlineto{\pgfqpoint{1.517692in}{0.781753in}}%
\pgfpathlineto{\pgfqpoint{1.526877in}{0.779738in}}%
\pgfpathlineto{\pgfqpoint{1.531154in}{0.772659in}}%
\pgfpathlineto{\pgfqpoint{1.538177in}{0.770275in}}%
\pgfpathlineto{\pgfqpoint{1.543049in}{0.762933in}}%
\pgfpathlineto{\pgfqpoint{1.552778in}{0.770922in}}%
\pgfpathlineto{\pgfqpoint{1.563911in}{0.761010in}}%
\pgfpathlineto{\pgfqpoint{1.576640in}{0.765059in}}%
\pgfpathlineto{\pgfqpoint{1.583084in}{0.754622in}}%
\pgfpathlineto{\pgfqpoint{1.589386in}{0.766758in}}%
\pgfpathlineto{\pgfqpoint{1.600714in}{0.759877in}}%
\pgfpathlineto{\pgfqpoint{1.607791in}{0.761708in}}%
\pgfpathlineto{\pgfqpoint{1.617897in}{0.756843in}}%
\pgfpathlineto{\pgfqpoint{1.630999in}{0.755462in}}%
\pgfpathlineto{\pgfqpoint{1.643418in}{0.762948in}}%
\pgfpathlineto{\pgfqpoint{1.647911in}{0.759788in}}%
\pgfpathlineto{\pgfqpoint{1.660532in}{0.766205in}}%
\pgfpathlineto{\pgfqpoint{1.668575in}{0.761197in}}%
\pgfpathlineto{\pgfqpoint{1.674932in}{0.761495in}}%
\pgfpathlineto{\pgfqpoint{1.679966in}{0.767394in}}%
\pgfpathlineto{\pgfqpoint{1.687888in}{0.764459in}}%
\pgfpathlineto{\pgfqpoint{1.705724in}{0.750800in}}%
\pgfpathlineto{\pgfqpoint{1.719643in}{0.746495in}}%
\pgfpathlineto{\pgfqpoint{1.725225in}{0.741221in}}%
\pgfpathlineto{\pgfqpoint{1.732212in}{0.744093in}}%
\pgfpathlineto{\pgfqpoint{1.742799in}{0.741894in}}%
\pgfpathlineto{\pgfqpoint{1.743006in}{0.708326in}}%
\pgfpathlineto{\pgfqpoint{1.743890in}{0.643403in}}%
\pgfpathlineto{\pgfqpoint{1.751243in}{0.637164in}}%
\pgfpathlineto{\pgfqpoint{1.757137in}{0.625668in}}%
\pgfpathlineto{\pgfqpoint{1.755065in}{0.617977in}}%
\pgfpathlineto{\pgfqpoint{1.759693in}{0.614721in}}%
\pgfpathlineto{\pgfqpoint{1.763454in}{0.600538in}}%
\pgfpathlineto{\pgfqpoint{1.770658in}{0.592942in}}%
\pgfpathlineto{\pgfqpoint{1.772272in}{0.576816in}}%
\pgfpathlineto{\pgfqpoint{1.771019in}{0.569988in}}%
\pgfpathlineto{\pgfqpoint{1.761227in}{0.551919in}}%
\pgfpathlineto{\pgfqpoint{1.760101in}{0.538927in}}%
\pgfpathlineto{\pgfqpoint{1.763423in}{0.536215in}}%
\pgfpathlineto{\pgfqpoint{1.763712in}{0.524834in}}%
\pgfpathlineto{\pgfqpoint{1.760335in}{0.517680in}}%
\pgfpathlineto{\pgfqpoint{1.755085in}{0.514985in}}%
\pgfpathlineto{\pgfqpoint{1.749944in}{0.505637in}}%
\pgfpathlineto{\pgfqpoint{1.756496in}{0.496583in}}%
\pgfpathlineto{\pgfqpoint{1.743777in}{0.496405in}}%
\pgfpathlineto{\pgfqpoint{1.709766in}{0.480844in}}%
\pgfpathlineto{\pgfqpoint{1.716259in}{0.490157in}}%
\pgfpathlineto{\pgfqpoint{1.708396in}{0.495178in}}%
\pgfpathlineto{\pgfqpoint{1.706753in}{0.503714in}}%
\pgfpathlineto{\pgfqpoint{1.691406in}{0.488845in}}%
\pgfpathlineto{\pgfqpoint{1.698217in}{0.478682in}}%
\pgfpathlineto{\pgfqpoint{1.688501in}{0.465742in}}%
\pgfpathlineto{\pgfqpoint{1.683273in}{0.466013in}}%
\pgfpathlineto{\pgfqpoint{1.678354in}{0.451920in}}%
\pgfpathlineto{\pgfqpoint{1.662798in}{0.440835in}}%
\pgfpathlineto{\pgfqpoint{1.648265in}{0.436840in}}%
\pgfpathlineto{\pgfqpoint{1.622917in}{0.426447in}}%
\pgfpathlineto{\pgfqpoint{1.620291in}{0.432244in}}%
\pgfpathlineto{\pgfqpoint{1.604501in}{0.426908in}}%
\pgfpathlineto{\pgfqpoint{1.614213in}{0.417822in}}%
\pgfpathlineto{\pgfqpoint{1.598894in}{0.409929in}}%
\pgfpathlineto{\pgfqpoint{1.582371in}{0.398011in}}%
\pgfpathlineto{\pgfqpoint{1.578403in}{0.403551in}}%
\pgfpathlineto{\pgfqpoint{1.564881in}{0.395248in}}%
\pgfpathlineto{\pgfqpoint{1.578135in}{0.392095in}}%
\pgfpathlineto{\pgfqpoint{1.568142in}{0.379039in}}%
\pgfpathlineto{\pgfqpoint{1.563264in}{0.382924in}}%
\pgfpathlineto{\pgfqpoint{1.556713in}{0.376684in}}%
\pgfpathlineto{\pgfqpoint{1.564878in}{0.371233in}}%
\pgfpathlineto{\pgfqpoint{1.560050in}{0.363265in}}%
\pgfpathlineto{\pgfqpoint{1.554117in}{0.344404in}}%
\pgfpathlineto{\pgfqpoint{1.545559in}{0.326238in}}%
\pgfpathlineto{\pgfqpoint{1.549401in}{0.314331in}}%
\pgfpathlineto{\pgfqpoint{1.555932in}{0.286258in}}%
\pgfpathlineto{\pgfqpoint{1.556520in}{0.274896in}}%
\pgfpathlineto{\pgfqpoint{1.566610in}{0.259994in}}%
\pgfpathlineto{\pgfqpoint{1.558829in}{0.260865in}}%
\pgfpathlineto{\pgfqpoint{1.551272in}{0.253347in}}%
\pgfpathlineto{\pgfqpoint{1.542420in}{0.259256in}}%
\pgfpathlineto{\pgfqpoint{1.539243in}{0.265150in}}%
\pgfpathlineto{\pgfqpoint{1.526645in}{0.267892in}}%
\pgfpathlineto{\pgfqpoint{1.507402in}{0.268229in}}%
\pgfpathlineto{\pgfqpoint{1.493224in}{0.279391in}}%
\pgfpathlineto{\pgfqpoint{1.480350in}{0.281256in}}%
\pgfpathlineto{\pgfqpoint{1.472540in}{0.290127in}}%
\pgfpathlineto{\pgfqpoint{1.456183in}{0.293697in}}%
\pgfpathlineto{\pgfqpoint{1.447239in}{0.322221in}}%
\pgfpathlineto{\pgfqpoint{1.438094in}{0.333648in}}%
\pgfpathlineto{\pgfqpoint{1.439648in}{0.344506in}}%
\pgfpathlineto{\pgfqpoint{1.433978in}{0.352440in}}%
\pgfpathlineto{\pgfqpoint{1.437536in}{0.363286in}}%
\pgfpathlineto{\pgfqpoint{1.434601in}{0.371234in}}%
\pgfpathlineto{\pgfqpoint{1.425411in}{0.374834in}}%
\pgfpathlineto{\pgfqpoint{1.416818in}{0.383999in}}%
\pgfpathlineto{\pgfqpoint{1.413687in}{0.396269in}}%
\pgfpathlineto{\pgfqpoint{1.405556in}{0.407420in}}%
\pgfpathlineto{\pgfqpoint{1.394734in}{0.416096in}}%
\pgfpathlineto{\pgfqpoint{1.384986in}{0.440984in}}%
\pgfpathlineto{\pgfqpoint{1.377115in}{0.468203in}}%
\pgfpathlineto{\pgfqpoint{1.370657in}{0.478966in}}%
\pgfpathlineto{\pgfqpoint{1.359460in}{0.488026in}}%
\pgfpathlineto{\pgfqpoint{1.356662in}{0.494603in}}%
\pgfpathlineto{\pgfqpoint{1.346179in}{0.498686in}}%
\pgfpathlineto{\pgfqpoint{1.335831in}{0.516112in}}%
\pgfpathlineto{\pgfqpoint{1.326382in}{0.514774in}}%
\pgfpathlineto{\pgfqpoint{1.316787in}{0.519126in}}%
\pgfpathlineto{\pgfqpoint{1.303186in}{0.518285in}}%
\pgfpathlineto{\pgfqpoint{1.289369in}{0.525467in}}%
\pgfpathlineto{\pgfqpoint{1.285472in}{0.518637in}}%
\pgfpathlineto{\pgfqpoint{1.269325in}{0.518466in}}%
\pgfpathlineto{\pgfqpoint{1.261104in}{0.505517in}}%
\pgfpathlineto{\pgfqpoint{1.253954in}{0.489484in}}%
\pgfpathlineto{\pgfqpoint{1.238750in}{0.472274in}}%
\pgfpathlineto{\pgfqpoint{1.221524in}{0.479827in}}%
\pgfpathlineto{\pgfqpoint{1.219080in}{0.484818in}}%
\pgfpathlineto{\pgfqpoint{1.208611in}{0.488625in}}%
\pgfpathlineto{\pgfqpoint{1.205396in}{0.493833in}}%
\pgfpathlineto{\pgfqpoint{1.191499in}{0.499090in}}%
\pgfpathlineto{\pgfqpoint{1.183725in}{0.509836in}}%
\pgfpathlineto{\pgfqpoint{1.174646in}{0.515021in}}%
\pgfpathlineto{\pgfqpoint{1.166834in}{0.524068in}}%
\pgfpathlineto{\pgfqpoint{1.160749in}{0.539391in}}%
\pgfpathlineto{\pgfqpoint{1.161430in}{0.560329in}}%
\pgfpathlineto{\pgfqpoint{1.154291in}{0.570909in}}%
\pgfpathlineto{\pgfqpoint{1.153469in}{0.582369in}}%
\pgfpathlineto{\pgfqpoint{1.137630in}{0.599527in}}%
\pgfpathlineto{\pgfqpoint{1.128414in}{0.603189in}}%
\pgfpathlineto{\pgfqpoint{1.118634in}{0.619199in}}%
\pgfpathlineto{\pgfqpoint{1.110300in}{0.625560in}}%
\pgfpathlineto{\pgfqpoint{1.099711in}{0.641049in}}%
\pgfpathlineto{\pgfqpoint{1.088876in}{0.647753in}}%
\pgfpathlineto{\pgfqpoint{1.081780in}{0.664929in}}%
\pgfpathlineto{\pgfqpoint{1.073883in}{0.669271in}}%
\pgfpathclose%
\pgfusepath{fill}%
\end{pgfscope}%
\begin{pgfscope}%
\pgfpathrectangle{\pgfqpoint{0.100000in}{0.100000in}}{\pgfqpoint{2.862985in}{1.832778in}}%
\pgfusepath{clip}%
\pgfsetbuttcap%
\pgfsetmiterjoin%
\definecolor{currentfill}{rgb}{0.846788,0.342551,0.420579}%
\pgfsetfillcolor{currentfill}%
\pgfsetlinewidth{0.000000pt}%
\definecolor{currentstroke}{rgb}{0.000000,0.000000,0.000000}%
\pgfsetstrokecolor{currentstroke}%
\pgfsetstrokeopacity{0.000000}%
\pgfsetdash{}{0pt}%
\pgfpathmoveto{\pgfqpoint{0.990915in}{1.012202in}}%
\pgfpathlineto{\pgfqpoint{1.069080in}{1.000850in}}%
\pgfpathlineto{\pgfqpoint{1.099362in}{0.996184in}}%
\pgfpathlineto{\pgfqpoint{1.184722in}{0.986032in}}%
\pgfpathlineto{\pgfqpoint{1.242175in}{0.980039in}}%
\pgfpathlineto{\pgfqpoint{1.291895in}{0.975633in}}%
\pgfpathlineto{\pgfqpoint{1.289118in}{0.943987in}}%
\pgfpathlineto{\pgfqpoint{1.287612in}{0.944071in}}%
\pgfpathlineto{\pgfqpoint{1.283057in}{0.889756in}}%
\pgfpathlineto{\pgfqpoint{1.278172in}{0.834767in}}%
\pgfpathlineto{\pgfqpoint{1.272520in}{0.777201in}}%
\pgfpathlineto{\pgfqpoint{1.263807in}{0.695029in}}%
\pgfpathlineto{\pgfqpoint{1.261021in}{0.662722in}}%
\pgfpathlineto{\pgfqpoint{1.209801in}{0.667891in}}%
\pgfpathlineto{\pgfqpoint{1.165452in}{0.672249in}}%
\pgfpathlineto{\pgfqpoint{1.104118in}{0.679215in}}%
\pgfpathlineto{\pgfqpoint{1.071132in}{0.683307in}}%
\pgfpathlineto{\pgfqpoint{1.069190in}{0.676811in}}%
\pgfpathlineto{\pgfqpoint{1.073883in}{0.669271in}}%
\pgfpathlineto{\pgfqpoint{1.034240in}{0.674419in}}%
\pgfpathlineto{\pgfqpoint{0.985330in}{0.681439in}}%
\pgfpathlineto{\pgfqpoint{0.980878in}{0.653775in}}%
\pgfpathlineto{\pgfqpoint{0.936264in}{0.660490in}}%
\pgfpathlineto{\pgfqpoint{0.946747in}{0.729118in}}%
\pgfpathlineto{\pgfqpoint{0.959723in}{0.812676in}}%
\pgfpathlineto{\pgfqpoint{0.971262in}{0.885469in}}%
\pgfpathlineto{\pgfqpoint{0.981281in}{0.950172in}}%
\pgfpathlineto{\pgfqpoint{0.990915in}{1.012202in}}%
\pgfpathclose%
\pgfusepath{fill}%
\end{pgfscope}%
\begin{pgfscope}%
\pgfpathrectangle{\pgfqpoint{0.100000in}{0.100000in}}{\pgfqpoint{2.862985in}{1.832778in}}%
\pgfusepath{clip}%
\pgfsetbuttcap%
\pgfsetmiterjoin%
\definecolor{currentfill}{rgb}{0.833422,0.324635,0.434366}%
\pgfsetfillcolor{currentfill}%
\pgfsetlinewidth{0.000000pt}%
\definecolor{currentstroke}{rgb}{0.000000,0.000000,0.000000}%
\pgfsetstrokecolor{currentstroke}%
\pgfsetstrokeopacity{0.000000}%
\pgfsetdash{}{0pt}%
\pgfpathmoveto{\pgfqpoint{2.233749in}{0.610119in}}%
\pgfpathlineto{\pgfqpoint{2.159178in}{0.601837in}}%
\pgfpathlineto{\pgfqpoint{2.093446in}{0.596734in}}%
\pgfpathlineto{\pgfqpoint{2.092624in}{0.588678in}}%
\pgfpathlineto{\pgfqpoint{2.098661in}{0.581011in}}%
\pgfpathlineto{\pgfqpoint{2.106031in}{0.576507in}}%
\pgfpathlineto{\pgfqpoint{2.105919in}{0.564647in}}%
\pgfpathlineto{\pgfqpoint{2.103966in}{0.556723in}}%
\pgfpathlineto{\pgfqpoint{2.097462in}{0.551013in}}%
\pgfpathlineto{\pgfqpoint{2.088396in}{0.551618in}}%
\pgfpathlineto{\pgfqpoint{2.079825in}{0.558693in}}%
\pgfpathlineto{\pgfqpoint{2.078297in}{0.571283in}}%
\pgfpathlineto{\pgfqpoint{2.071913in}{0.578606in}}%
\pgfpathlineto{\pgfqpoint{2.067568in}{0.552398in}}%
\pgfpathlineto{\pgfqpoint{2.052803in}{0.554889in}}%
\pgfpathlineto{\pgfqpoint{2.047669in}{0.600717in}}%
\pgfpathlineto{\pgfqpoint{2.042108in}{0.649015in}}%
\pgfpathlineto{\pgfqpoint{2.044173in}{0.718697in}}%
\pgfpathlineto{\pgfqpoint{2.045455in}{0.766513in}}%
\pgfpathlineto{\pgfqpoint{2.048181in}{0.839470in}}%
\pgfpathlineto{\pgfqpoint{2.042346in}{0.845890in}}%
\pgfpathlineto{\pgfqpoint{2.052965in}{0.847429in}}%
\pgfpathlineto{\pgfqpoint{2.114830in}{0.851430in}}%
\pgfpathlineto{\pgfqpoint{2.175022in}{0.856843in}}%
\pgfpathlineto{\pgfqpoint{2.190858in}{0.801373in}}%
\pgfpathlineto{\pgfqpoint{2.201658in}{0.760657in}}%
\pgfpathlineto{\pgfqpoint{2.211339in}{0.726571in}}%
\pgfpathlineto{\pgfqpoint{2.216938in}{0.712845in}}%
\pgfpathlineto{\pgfqpoint{2.225753in}{0.700073in}}%
\pgfpathlineto{\pgfqpoint{2.224330in}{0.693573in}}%
\pgfpathlineto{\pgfqpoint{2.230627in}{0.690406in}}%
\pgfpathlineto{\pgfqpoint{2.223185in}{0.680553in}}%
\pgfpathlineto{\pgfqpoint{2.220672in}{0.663006in}}%
\pgfpathlineto{\pgfqpoint{2.227945in}{0.642483in}}%
\pgfpathlineto{\pgfqpoint{2.226935in}{0.621831in}}%
\pgfpathlineto{\pgfqpoint{2.233749in}{0.610119in}}%
\pgfpathclose%
\pgfusepath{fill}%
\end{pgfscope}%
\begin{pgfscope}%
\pgfpathrectangle{\pgfqpoint{0.100000in}{0.100000in}}{\pgfqpoint{2.862985in}{1.832778in}}%
\pgfusepath{clip}%
\pgfsetbuttcap%
\pgfsetmiterjoin%
\definecolor{currentfill}{rgb}{0.928329,0.472975,0.326067}%
\pgfsetfillcolor{currentfill}%
\pgfsetlinewidth{0.000000pt}%
\definecolor{currentstroke}{rgb}{0.000000,0.000000,0.000000}%
\pgfsetstrokecolor{currentstroke}%
\pgfsetstrokeopacity{0.000000}%
\pgfsetdash{}{0pt}%
\pgfpathmoveto{\pgfqpoint{2.052803in}{0.554889in}}%
\pgfpathlineto{\pgfqpoint{2.037652in}{0.550553in}}%
\pgfpathlineto{\pgfqpoint{2.026358in}{0.553337in}}%
\pgfpathlineto{\pgfqpoint{2.005384in}{0.546666in}}%
\pgfpathlineto{\pgfqpoint{1.989565in}{0.538074in}}%
\pgfpathlineto{\pgfqpoint{1.983075in}{0.553600in}}%
\pgfpathlineto{\pgfqpoint{1.976384in}{0.560107in}}%
\pgfpathlineto{\pgfqpoint{1.973157in}{0.567894in}}%
\pgfpathlineto{\pgfqpoint{1.977888in}{0.588972in}}%
\pgfpathlineto{\pgfqpoint{1.933060in}{0.586211in}}%
\pgfpathlineto{\pgfqpoint{1.874932in}{0.583691in}}%
\pgfpathlineto{\pgfqpoint{1.878390in}{0.588938in}}%
\pgfpathlineto{\pgfqpoint{1.874116in}{0.598847in}}%
\pgfpathlineto{\pgfqpoint{1.878375in}{0.600866in}}%
\pgfpathlineto{\pgfqpoint{1.877430in}{0.610386in}}%
\pgfpathlineto{\pgfqpoint{1.884854in}{0.619619in}}%
\pgfpathlineto{\pgfqpoint{1.888916in}{0.637542in}}%
\pgfpathlineto{\pgfqpoint{1.902503in}{0.649350in}}%
\pgfpathlineto{\pgfqpoint{1.897467in}{0.660818in}}%
\pgfpathlineto{\pgfqpoint{1.907102in}{0.667419in}}%
\pgfpathlineto{\pgfqpoint{1.898791in}{0.679370in}}%
\pgfpathlineto{\pgfqpoint{1.892776in}{0.706727in}}%
\pgfpathlineto{\pgfqpoint{1.895057in}{0.711695in}}%
\pgfpathlineto{\pgfqpoint{1.898659in}{0.721223in}}%
\pgfpathlineto{\pgfqpoint{1.895307in}{0.731223in}}%
\pgfpathlineto{\pgfqpoint{1.896946in}{0.735772in}}%
\pgfpathlineto{\pgfqpoint{1.891048in}{0.752960in}}%
\pgfpathlineto{\pgfqpoint{1.908459in}{0.783821in}}%
\pgfpathlineto{\pgfqpoint{1.909269in}{0.792035in}}%
\pgfpathlineto{\pgfqpoint{1.917656in}{0.798007in}}%
\pgfpathlineto{\pgfqpoint{1.926605in}{0.817725in}}%
\pgfpathlineto{\pgfqpoint{1.926180in}{0.825739in}}%
\pgfpathlineto{\pgfqpoint{1.937326in}{0.833932in}}%
\pgfpathlineto{\pgfqpoint{1.933848in}{0.838785in}}%
\pgfpathlineto{\pgfqpoint{2.000643in}{0.842683in}}%
\pgfpathlineto{\pgfqpoint{2.042346in}{0.845890in}}%
\pgfpathlineto{\pgfqpoint{2.048181in}{0.839470in}}%
\pgfpathlineto{\pgfqpoint{2.045455in}{0.766513in}}%
\pgfpathlineto{\pgfqpoint{2.044173in}{0.718697in}}%
\pgfpathlineto{\pgfqpoint{2.042108in}{0.649015in}}%
\pgfpathlineto{\pgfqpoint{2.047669in}{0.600717in}}%
\pgfpathlineto{\pgfqpoint{2.052803in}{0.554889in}}%
\pgfpathclose%
\pgfusepath{fill}%
\end{pgfscope}%
\begin{pgfscope}%
\pgfpathrectangle{\pgfqpoint{0.100000in}{0.100000in}}{\pgfqpoint{2.862985in}{1.832778in}}%
\pgfusepath{clip}%
\pgfsetbuttcap%
\pgfsetmiterjoin%
\definecolor{currentfill}{rgb}{0.731862,0.205013,0.531601}%
\pgfsetfillcolor{currentfill}%
\pgfsetlinewidth{0.000000pt}%
\definecolor{currentstroke}{rgb}{0.000000,0.000000,0.000000}%
\pgfsetstrokecolor{currentstroke}%
\pgfsetstrokeopacity{0.000000}%
\pgfsetdash{}{0pt}%
\pgfpathmoveto{\pgfqpoint{2.175022in}{0.856843in}}%
\pgfpathlineto{\pgfqpoint{2.240820in}{0.864070in}}%
\pgfpathlineto{\pgfqpoint{2.280410in}{0.868948in}}%
\pgfpathlineto{\pgfqpoint{2.302958in}{0.872540in}}%
\pgfpathlineto{\pgfqpoint{2.293618in}{0.857809in}}%
\pgfpathlineto{\pgfqpoint{2.293623in}{0.850863in}}%
\pgfpathlineto{\pgfqpoint{2.303238in}{0.847117in}}%
\pgfpathlineto{\pgfqpoint{2.309776in}{0.841070in}}%
\pgfpathlineto{\pgfqpoint{2.319657in}{0.840323in}}%
\pgfpathlineto{\pgfqpoint{2.328892in}{0.823299in}}%
\pgfpathlineto{\pgfqpoint{2.338910in}{0.811294in}}%
\pgfpathlineto{\pgfqpoint{2.356702in}{0.801205in}}%
\pgfpathlineto{\pgfqpoint{2.361758in}{0.793157in}}%
\pgfpathlineto{\pgfqpoint{2.377624in}{0.784454in}}%
\pgfpathlineto{\pgfqpoint{2.378817in}{0.778143in}}%
\pgfpathlineto{\pgfqpoint{2.391688in}{0.765593in}}%
\pgfpathlineto{\pgfqpoint{2.401546in}{0.762019in}}%
\pgfpathlineto{\pgfqpoint{2.408528in}{0.750169in}}%
\pgfpathlineto{\pgfqpoint{2.411618in}{0.736849in}}%
\pgfpathlineto{\pgfqpoint{2.421858in}{0.731700in}}%
\pgfpathlineto{\pgfqpoint{2.430102in}{0.717383in}}%
\pgfpathlineto{\pgfqpoint{2.432761in}{0.706853in}}%
\pgfpathlineto{\pgfqpoint{2.444841in}{0.702371in}}%
\pgfpathlineto{\pgfqpoint{2.434704in}{0.682961in}}%
\pgfpathlineto{\pgfqpoint{2.430889in}{0.671491in}}%
\pgfpathlineto{\pgfqpoint{2.433378in}{0.666114in}}%
\pgfpathlineto{\pgfqpoint{2.429032in}{0.657178in}}%
\pgfpathlineto{\pgfqpoint{2.427185in}{0.644824in}}%
\pgfpathlineto{\pgfqpoint{2.419456in}{0.642523in}}%
\pgfpathlineto{\pgfqpoint{2.423537in}{0.631212in}}%
\pgfpathlineto{\pgfqpoint{2.423277in}{0.616862in}}%
\pgfpathlineto{\pgfqpoint{2.409539in}{0.617400in}}%
\pgfpathlineto{\pgfqpoint{2.399969in}{0.620018in}}%
\pgfpathlineto{\pgfqpoint{2.395828in}{0.615151in}}%
\pgfpathlineto{\pgfqpoint{2.398918in}{0.603653in}}%
\pgfpathlineto{\pgfqpoint{2.398246in}{0.590285in}}%
\pgfpathlineto{\pgfqpoint{2.392209in}{0.589262in}}%
\pgfpathlineto{\pgfqpoint{2.387325in}{0.601747in}}%
\pgfpathlineto{\pgfqpoint{2.337593in}{0.598433in}}%
\pgfpathlineto{\pgfqpoint{2.274885in}{0.595011in}}%
\pgfpathlineto{\pgfqpoint{2.243260in}{0.592785in}}%
\pgfpathlineto{\pgfqpoint{2.233749in}{0.610119in}}%
\pgfpathlineto{\pgfqpoint{2.226935in}{0.621831in}}%
\pgfpathlineto{\pgfqpoint{2.227945in}{0.642483in}}%
\pgfpathlineto{\pgfqpoint{2.220672in}{0.663006in}}%
\pgfpathlineto{\pgfqpoint{2.223185in}{0.680553in}}%
\pgfpathlineto{\pgfqpoint{2.230627in}{0.690406in}}%
\pgfpathlineto{\pgfqpoint{2.224330in}{0.693573in}}%
\pgfpathlineto{\pgfqpoint{2.225753in}{0.700073in}}%
\pgfpathlineto{\pgfqpoint{2.216938in}{0.712845in}}%
\pgfpathlineto{\pgfqpoint{2.211339in}{0.726571in}}%
\pgfpathlineto{\pgfqpoint{2.201658in}{0.760657in}}%
\pgfpathlineto{\pgfqpoint{2.190858in}{0.801373in}}%
\pgfpathlineto{\pgfqpoint{2.175022in}{0.856843in}}%
\pgfpathclose%
\pgfusepath{fill}%
\end{pgfscope}%
\begin{pgfscope}%
\pgfpathrectangle{\pgfqpoint{0.100000in}{0.100000in}}{\pgfqpoint{2.862985in}{1.832778in}}%
\pgfusepath{clip}%
\pgfsetbuttcap%
\pgfsetmiterjoin%
\definecolor{currentfill}{rgb}{0.859750,0.360588,0.406917}%
\pgfsetfillcolor{currentfill}%
\pgfsetlinewidth{0.000000pt}%
\definecolor{currentstroke}{rgb}{0.000000,0.000000,0.000000}%
\pgfsetstrokecolor{currentstroke}%
\pgfsetstrokeopacity{0.000000}%
\pgfsetdash{}{0pt}%
\pgfpathmoveto{\pgfqpoint{2.302958in}{0.872540in}}%
\pgfpathlineto{\pgfqpoint{2.329049in}{0.885519in}}%
\pgfpathlineto{\pgfqpoint{2.343499in}{0.890420in}}%
\pgfpathlineto{\pgfqpoint{2.406366in}{0.896952in}}%
\pgfpathlineto{\pgfqpoint{2.413024in}{0.894789in}}%
\pgfpathlineto{\pgfqpoint{2.421819in}{0.885949in}}%
\pgfpathlineto{\pgfqpoint{2.422304in}{0.878098in}}%
\pgfpathlineto{\pgfqpoint{2.480165in}{0.886600in}}%
\pgfpathlineto{\pgfqpoint{2.545869in}{0.839501in}}%
\pgfpathlineto{\pgfqpoint{2.533547in}{0.826658in}}%
\pgfpathlineto{\pgfqpoint{2.516640in}{0.796833in}}%
\pgfpathlineto{\pgfqpoint{2.521453in}{0.790423in}}%
\pgfpathlineto{\pgfqpoint{2.512434in}{0.777981in}}%
\pgfpathlineto{\pgfqpoint{2.503462in}{0.776573in}}%
\pgfpathlineto{\pgfqpoint{2.503429in}{0.769097in}}%
\pgfpathlineto{\pgfqpoint{2.496939in}{0.761274in}}%
\pgfpathlineto{\pgfqpoint{2.488862in}{0.759687in}}%
\pgfpathlineto{\pgfqpoint{2.490625in}{0.752736in}}%
\pgfpathlineto{\pgfqpoint{2.486118in}{0.747400in}}%
\pgfpathlineto{\pgfqpoint{2.475337in}{0.742790in}}%
\pgfpathlineto{\pgfqpoint{2.462196in}{0.732304in}}%
\pgfpathlineto{\pgfqpoint{2.464677in}{0.725520in}}%
\pgfpathlineto{\pgfqpoint{2.456407in}{0.721260in}}%
\pgfpathlineto{\pgfqpoint{2.449432in}{0.725739in}}%
\pgfpathlineto{\pgfqpoint{2.444331in}{0.706267in}}%
\pgfpathlineto{\pgfqpoint{2.432761in}{0.706853in}}%
\pgfpathlineto{\pgfqpoint{2.430102in}{0.717383in}}%
\pgfpathlineto{\pgfqpoint{2.421858in}{0.731700in}}%
\pgfpathlineto{\pgfqpoint{2.411618in}{0.736849in}}%
\pgfpathlineto{\pgfqpoint{2.408528in}{0.750169in}}%
\pgfpathlineto{\pgfqpoint{2.401546in}{0.762019in}}%
\pgfpathlineto{\pgfqpoint{2.391688in}{0.765593in}}%
\pgfpathlineto{\pgfqpoint{2.378817in}{0.778143in}}%
\pgfpathlineto{\pgfqpoint{2.377624in}{0.784454in}}%
\pgfpathlineto{\pgfqpoint{2.361758in}{0.793157in}}%
\pgfpathlineto{\pgfqpoint{2.356702in}{0.801205in}}%
\pgfpathlineto{\pgfqpoint{2.338910in}{0.811294in}}%
\pgfpathlineto{\pgfqpoint{2.328892in}{0.823299in}}%
\pgfpathlineto{\pgfqpoint{2.319657in}{0.840323in}}%
\pgfpathlineto{\pgfqpoint{2.309776in}{0.841070in}}%
\pgfpathlineto{\pgfqpoint{2.303238in}{0.847117in}}%
\pgfpathlineto{\pgfqpoint{2.293623in}{0.850863in}}%
\pgfpathlineto{\pgfqpoint{2.293618in}{0.857809in}}%
\pgfpathlineto{\pgfqpoint{2.302958in}{0.872540in}}%
\pgfpathclose%
\pgfusepath{fill}%
\end{pgfscope}%
\begin{pgfscope}%
\pgfpathrectangle{\pgfqpoint{0.100000in}{0.100000in}}{\pgfqpoint{2.862985in}{1.832778in}}%
\pgfusepath{clip}%
\pgfsetbuttcap%
\pgfsetmiterjoin%
\definecolor{currentfill}{rgb}{0.856547,0.356066,0.410322}%
\pgfsetfillcolor{currentfill}%
\pgfsetlinewidth{0.000000pt}%
\definecolor{currentstroke}{rgb}{0.000000,0.000000,0.000000}%
\pgfsetstrokecolor{currentstroke}%
\pgfsetstrokeopacity{0.000000}%
\pgfsetdash{}{0pt}%
\pgfpathmoveto{\pgfqpoint{1.711515in}{0.926094in}}%
\pgfpathlineto{\pgfqpoint{1.749856in}{0.926440in}}%
\pgfpathlineto{\pgfqpoint{1.800580in}{0.927364in}}%
\pgfpathlineto{\pgfqpoint{1.887257in}{0.930116in}}%
\pgfpathlineto{\pgfqpoint{1.936833in}{0.932737in}}%
\pgfpathlineto{\pgfqpoint{1.942150in}{0.926155in}}%
\pgfpathlineto{\pgfqpoint{1.941818in}{0.919204in}}%
\pgfpathlineto{\pgfqpoint{1.929808in}{0.907211in}}%
\pgfpathlineto{\pgfqpoint{1.926904in}{0.900644in}}%
\pgfpathlineto{\pgfqpoint{1.960213in}{0.903107in}}%
\pgfpathlineto{\pgfqpoint{1.960195in}{0.890974in}}%
\pgfpathlineto{\pgfqpoint{1.949509in}{0.885789in}}%
\pgfpathlineto{\pgfqpoint{1.949660in}{0.877545in}}%
\pgfpathlineto{\pgfqpoint{1.945647in}{0.871619in}}%
\pgfpathlineto{\pgfqpoint{1.943020in}{0.859033in}}%
\pgfpathlineto{\pgfqpoint{1.945732in}{0.848649in}}%
\pgfpathlineto{\pgfqpoint{1.933848in}{0.838785in}}%
\pgfpathlineto{\pgfqpoint{1.937326in}{0.833932in}}%
\pgfpathlineto{\pgfqpoint{1.926180in}{0.825739in}}%
\pgfpathlineto{\pgfqpoint{1.926605in}{0.817725in}}%
\pgfpathlineto{\pgfqpoint{1.917656in}{0.798007in}}%
\pgfpathlineto{\pgfqpoint{1.909269in}{0.792035in}}%
\pgfpathlineto{\pgfqpoint{1.908459in}{0.783821in}}%
\pgfpathlineto{\pgfqpoint{1.891048in}{0.752960in}}%
\pgfpathlineto{\pgfqpoint{1.896946in}{0.735772in}}%
\pgfpathlineto{\pgfqpoint{1.895307in}{0.731223in}}%
\pgfpathlineto{\pgfqpoint{1.898659in}{0.721223in}}%
\pgfpathlineto{\pgfqpoint{1.895057in}{0.711695in}}%
\pgfpathlineto{\pgfqpoint{1.847446in}{0.709729in}}%
\pgfpathlineto{\pgfqpoint{1.785638in}{0.708711in}}%
\pgfpathlineto{\pgfqpoint{1.743006in}{0.708326in}}%
\pgfpathlineto{\pgfqpoint{1.742799in}{0.741894in}}%
\pgfpathlineto{\pgfqpoint{1.732212in}{0.744093in}}%
\pgfpathlineto{\pgfqpoint{1.725225in}{0.741221in}}%
\pgfpathlineto{\pgfqpoint{1.719643in}{0.746495in}}%
\pgfpathlineto{\pgfqpoint{1.720219in}{0.782052in}}%
\pgfpathlineto{\pgfqpoint{1.721468in}{0.857760in}}%
\pgfpathlineto{\pgfqpoint{1.715413in}{0.902092in}}%
\pgfpathlineto{\pgfqpoint{1.711515in}{0.926094in}}%
\pgfpathclose%
\pgfusepath{fill}%
\end{pgfscope}%
\begin{pgfscope}%
\pgfpathrectangle{\pgfqpoint{0.100000in}{0.100000in}}{\pgfqpoint{2.862985in}{1.832778in}}%
\pgfusepath{clip}%
\pgfsetbuttcap%
\pgfsetmiterjoin%
\definecolor{currentfill}{rgb}{0.823132,0.311261,0.444806}%
\pgfsetfillcolor{currentfill}%
\pgfsetlinewidth{0.000000pt}%
\definecolor{currentstroke}{rgb}{0.000000,0.000000,0.000000}%
\pgfsetstrokecolor{currentstroke}%
\pgfsetstrokeopacity{0.000000}%
\pgfsetdash{}{0pt}%
\pgfpathmoveto{\pgfqpoint{1.743006in}{0.708326in}}%
\pgfpathlineto{\pgfqpoint{1.785638in}{0.708711in}}%
\pgfpathlineto{\pgfqpoint{1.847446in}{0.709729in}}%
\pgfpathlineto{\pgfqpoint{1.895057in}{0.711695in}}%
\pgfpathlineto{\pgfqpoint{1.892776in}{0.706727in}}%
\pgfpathlineto{\pgfqpoint{1.898791in}{0.679370in}}%
\pgfpathlineto{\pgfqpoint{1.907102in}{0.667419in}}%
\pgfpathlineto{\pgfqpoint{1.897467in}{0.660818in}}%
\pgfpathlineto{\pgfqpoint{1.902503in}{0.649350in}}%
\pgfpathlineto{\pgfqpoint{1.888916in}{0.637542in}}%
\pgfpathlineto{\pgfqpoint{1.884854in}{0.619619in}}%
\pgfpathlineto{\pgfqpoint{1.877430in}{0.610386in}}%
\pgfpathlineto{\pgfqpoint{1.878375in}{0.600866in}}%
\pgfpathlineto{\pgfqpoint{1.874116in}{0.598847in}}%
\pgfpathlineto{\pgfqpoint{1.878390in}{0.588938in}}%
\pgfpathlineto{\pgfqpoint{1.874932in}{0.583691in}}%
\pgfpathlineto{\pgfqpoint{1.933060in}{0.586211in}}%
\pgfpathlineto{\pgfqpoint{1.977888in}{0.588972in}}%
\pgfpathlineto{\pgfqpoint{1.973157in}{0.567894in}}%
\pgfpathlineto{\pgfqpoint{1.976384in}{0.560107in}}%
\pgfpathlineto{\pgfqpoint{1.983075in}{0.553600in}}%
\pgfpathlineto{\pgfqpoint{1.989565in}{0.538074in}}%
\pgfpathlineto{\pgfqpoint{1.969065in}{0.541651in}}%
\pgfpathlineto{\pgfqpoint{1.961502in}{0.547531in}}%
\pgfpathlineto{\pgfqpoint{1.952492in}{0.547811in}}%
\pgfpathlineto{\pgfqpoint{1.943019in}{0.534934in}}%
\pgfpathlineto{\pgfqpoint{1.944907in}{0.529071in}}%
\pgfpathlineto{\pgfqpoint{1.975169in}{0.525523in}}%
\pgfpathlineto{\pgfqpoint{1.983091in}{0.518783in}}%
\pgfpathlineto{\pgfqpoint{1.995984in}{0.515315in}}%
\pgfpathlineto{\pgfqpoint{1.990353in}{0.507336in}}%
\pgfpathlineto{\pgfqpoint{1.980867in}{0.500376in}}%
\pgfpathlineto{\pgfqpoint{2.005906in}{0.484726in}}%
\pgfpathlineto{\pgfqpoint{2.017569in}{0.482281in}}%
\pgfpathlineto{\pgfqpoint{2.022419in}{0.469488in}}%
\pgfpathlineto{\pgfqpoint{2.011438in}{0.467135in}}%
\pgfpathlineto{\pgfqpoint{1.990799in}{0.483224in}}%
\pgfpathlineto{\pgfqpoint{1.982729in}{0.485461in}}%
\pgfpathlineto{\pgfqpoint{1.978824in}{0.491817in}}%
\pgfpathlineto{\pgfqpoint{1.966854in}{0.489208in}}%
\pgfpathlineto{\pgfqpoint{1.963107in}{0.481018in}}%
\pgfpathlineto{\pgfqpoint{1.965480in}{0.471892in}}%
\pgfpathlineto{\pgfqpoint{1.957422in}{0.466506in}}%
\pgfpathlineto{\pgfqpoint{1.950068in}{0.479774in}}%
\pgfpathlineto{\pgfqpoint{1.937187in}{0.475818in}}%
\pgfpathlineto{\pgfqpoint{1.932360in}{0.467901in}}%
\pgfpathlineto{\pgfqpoint{1.923213in}{0.470153in}}%
\pgfpathlineto{\pgfqpoint{1.917364in}{0.480150in}}%
\pgfpathlineto{\pgfqpoint{1.901847in}{0.483451in}}%
\pgfpathlineto{\pgfqpoint{1.898917in}{0.488653in}}%
\pgfpathlineto{\pgfqpoint{1.882751in}{0.497715in}}%
\pgfpathlineto{\pgfqpoint{1.878738in}{0.505638in}}%
\pgfpathlineto{\pgfqpoint{1.865197in}{0.502400in}}%
\pgfpathlineto{\pgfqpoint{1.866917in}{0.509654in}}%
\pgfpathlineto{\pgfqpoint{1.850086in}{0.502213in}}%
\pgfpathlineto{\pgfqpoint{1.854607in}{0.494492in}}%
\pgfpathlineto{\pgfqpoint{1.841610in}{0.489928in}}%
\pgfpathlineto{\pgfqpoint{1.824395in}{0.492437in}}%
\pgfpathlineto{\pgfqpoint{1.789551in}{0.504372in}}%
\pgfpathlineto{\pgfqpoint{1.762666in}{0.502003in}}%
\pgfpathlineto{\pgfqpoint{1.760335in}{0.517680in}}%
\pgfpathlineto{\pgfqpoint{1.763712in}{0.524834in}}%
\pgfpathlineto{\pgfqpoint{1.763423in}{0.536215in}}%
\pgfpathlineto{\pgfqpoint{1.760101in}{0.538927in}}%
\pgfpathlineto{\pgfqpoint{1.761227in}{0.551919in}}%
\pgfpathlineto{\pgfqpoint{1.771019in}{0.569988in}}%
\pgfpathlineto{\pgfqpoint{1.772272in}{0.576816in}}%
\pgfpathlineto{\pgfqpoint{1.770658in}{0.592942in}}%
\pgfpathlineto{\pgfqpoint{1.763454in}{0.600538in}}%
\pgfpathlineto{\pgfqpoint{1.759693in}{0.614721in}}%
\pgfpathlineto{\pgfqpoint{1.755065in}{0.617977in}}%
\pgfpathlineto{\pgfqpoint{1.757137in}{0.625668in}}%
\pgfpathlineto{\pgfqpoint{1.751243in}{0.637164in}}%
\pgfpathlineto{\pgfqpoint{1.743890in}{0.643403in}}%
\pgfpathlineto{\pgfqpoint{1.743006in}{0.708326in}}%
\pgfpathclose%
\pgfusepath{fill}%
\end{pgfscope}%
\begin{pgfscope}%
\pgfpathrectangle{\pgfqpoint{0.100000in}{0.100000in}}{\pgfqpoint{2.862985in}{1.832778in}}%
\pgfusepath{clip}%
\pgfsetbuttcap%
\pgfsetmiterjoin%
\definecolor{currentfill}{rgb}{0.823132,0.311261,0.444806}%
\pgfsetfillcolor{currentfill}%
\pgfsetlinewidth{0.000000pt}%
\definecolor{currentstroke}{rgb}{0.000000,0.000000,0.000000}%
\pgfsetstrokecolor{currentstroke}%
\pgfsetstrokeopacity{0.000000}%
\pgfsetdash{}{0pt}%
\pgfpathmoveto{\pgfqpoint{1.857025in}{0.493901in}}%
\pgfpathlineto{\pgfqpoint{1.863193in}{0.497570in}}%
\pgfpathlineto{\pgfqpoint{1.870679in}{0.493233in}}%
\pgfpathlineto{\pgfqpoint{1.866504in}{0.487265in}}%
\pgfpathlineto{\pgfqpoint{1.857025in}{0.493901in}}%
\pgfpathclose%
\pgfusepath{fill}%
\end{pgfscope}%
\begin{pgfscope}%
\pgfpathrectangle{\pgfqpoint{0.100000in}{0.100000in}}{\pgfqpoint{2.862985in}{1.832778in}}%
\pgfusepath{clip}%
\pgfsetbuttcap%
\pgfsetmiterjoin%
\definecolor{currentfill}{rgb}{0.787133,0.266922,0.480307}%
\pgfsetfillcolor{currentfill}%
\pgfsetlinewidth{0.000000pt}%
\definecolor{currentstroke}{rgb}{0.000000,0.000000,0.000000}%
\pgfsetstrokecolor{currentstroke}%
\pgfsetstrokeopacity{0.000000}%
\pgfsetdash{}{0pt}%
\pgfpathmoveto{\pgfqpoint{2.105919in}{0.564647in}}%
\pgfpathlineto{\pgfqpoint{2.106031in}{0.576507in}}%
\pgfpathlineto{\pgfqpoint{2.098661in}{0.581011in}}%
\pgfpathlineto{\pgfqpoint{2.092624in}{0.588678in}}%
\pgfpathlineto{\pgfqpoint{2.093446in}{0.596734in}}%
\pgfpathlineto{\pgfqpoint{2.159178in}{0.601837in}}%
\pgfpathlineto{\pgfqpoint{2.233749in}{0.610119in}}%
\pgfpathlineto{\pgfqpoint{2.243260in}{0.592785in}}%
\pgfpathlineto{\pgfqpoint{2.274885in}{0.595011in}}%
\pgfpathlineto{\pgfqpoint{2.337593in}{0.598433in}}%
\pgfpathlineto{\pgfqpoint{2.387325in}{0.601747in}}%
\pgfpathlineto{\pgfqpoint{2.392209in}{0.589262in}}%
\pgfpathlineto{\pgfqpoint{2.398246in}{0.590285in}}%
\pgfpathlineto{\pgfqpoint{2.398918in}{0.603653in}}%
\pgfpathlineto{\pgfqpoint{2.395828in}{0.615151in}}%
\pgfpathlineto{\pgfqpoint{2.399969in}{0.620018in}}%
\pgfpathlineto{\pgfqpoint{2.409539in}{0.617400in}}%
\pgfpathlineto{\pgfqpoint{2.423277in}{0.616862in}}%
\pgfpathlineto{\pgfqpoint{2.425389in}{0.606589in}}%
\pgfpathlineto{\pgfqpoint{2.429615in}{0.600719in}}%
\pgfpathlineto{\pgfqpoint{2.432915in}{0.587877in}}%
\pgfpathlineto{\pgfqpoint{2.443139in}{0.568003in}}%
\pgfpathlineto{\pgfqpoint{2.443190in}{0.562628in}}%
\pgfpathlineto{\pgfqpoint{2.458302in}{0.539291in}}%
\pgfpathlineto{\pgfqpoint{2.459764in}{0.534471in}}%
\pgfpathlineto{\pgfqpoint{2.482434in}{0.501512in}}%
\pgfpathlineto{\pgfqpoint{2.478842in}{0.500979in}}%
\pgfpathlineto{\pgfqpoint{2.488399in}{0.477537in}}%
\pgfpathlineto{\pgfqpoint{2.514653in}{0.436782in}}%
\pgfpathlineto{\pgfqpoint{2.529896in}{0.406818in}}%
\pgfpathlineto{\pgfqpoint{2.534991in}{0.401571in}}%
\pgfpathlineto{\pgfqpoint{2.543722in}{0.382763in}}%
\pgfpathlineto{\pgfqpoint{2.546750in}{0.352646in}}%
\pgfpathlineto{\pgfqpoint{2.547955in}{0.330127in}}%
\pgfpathlineto{\pgfqpoint{2.546500in}{0.315709in}}%
\pgfpathlineto{\pgfqpoint{2.541819in}{0.305401in}}%
\pgfpathlineto{\pgfqpoint{2.543994in}{0.291885in}}%
\pgfpathlineto{\pgfqpoint{2.538954in}{0.281186in}}%
\pgfpathlineto{\pgfqpoint{2.523991in}{0.272413in}}%
\pgfpathlineto{\pgfqpoint{2.514259in}{0.273053in}}%
\pgfpathlineto{\pgfqpoint{2.507940in}{0.268467in}}%
\pgfpathlineto{\pgfqpoint{2.506129in}{0.280712in}}%
\pgfpathlineto{\pgfqpoint{2.495668in}{0.283944in}}%
\pgfpathlineto{\pgfqpoint{2.486290in}{0.301030in}}%
\pgfpathlineto{\pgfqpoint{2.485231in}{0.308810in}}%
\pgfpathlineto{\pgfqpoint{2.468448in}{0.313535in}}%
\pgfpathlineto{\pgfqpoint{2.457633in}{0.312512in}}%
\pgfpathlineto{\pgfqpoint{2.451514in}{0.323801in}}%
\pgfpathlineto{\pgfqpoint{2.444508in}{0.344093in}}%
\pgfpathlineto{\pgfqpoint{2.434786in}{0.348206in}}%
\pgfpathlineto{\pgfqpoint{2.429515in}{0.359834in}}%
\pgfpathlineto{\pgfqpoint{2.416625in}{0.366687in}}%
\pgfpathlineto{\pgfqpoint{2.409174in}{0.375259in}}%
\pgfpathlineto{\pgfqpoint{2.396320in}{0.398564in}}%
\pgfpathlineto{\pgfqpoint{2.394809in}{0.414966in}}%
\pgfpathlineto{\pgfqpoint{2.401813in}{0.425461in}}%
\pgfpathlineto{\pgfqpoint{2.401113in}{0.432754in}}%
\pgfpathlineto{\pgfqpoint{2.393019in}{0.433533in}}%
\pgfpathlineto{\pgfqpoint{2.386256in}{0.438606in}}%
\pgfpathlineto{\pgfqpoint{2.382553in}{0.432412in}}%
\pgfpathlineto{\pgfqpoint{2.389043in}{0.427365in}}%
\pgfpathlineto{\pgfqpoint{2.384319in}{0.418289in}}%
\pgfpathlineto{\pgfqpoint{2.376677in}{0.425822in}}%
\pgfpathlineto{\pgfqpoint{2.377564in}{0.446710in}}%
\pgfpathlineto{\pgfqpoint{2.381250in}{0.463657in}}%
\pgfpathlineto{\pgfqpoint{2.379355in}{0.492736in}}%
\pgfpathlineto{\pgfqpoint{2.371738in}{0.499662in}}%
\pgfpathlineto{\pgfqpoint{2.367907in}{0.508549in}}%
\pgfpathlineto{\pgfqpoint{2.354785in}{0.508352in}}%
\pgfpathlineto{\pgfqpoint{2.341814in}{0.522976in}}%
\pgfpathlineto{\pgfqpoint{2.333117in}{0.527359in}}%
\pgfpathlineto{\pgfqpoint{2.330538in}{0.536620in}}%
\pgfpathlineto{\pgfqpoint{2.322031in}{0.539872in}}%
\pgfpathlineto{\pgfqpoint{2.314981in}{0.550098in}}%
\pgfpathlineto{\pgfqpoint{2.296358in}{0.558444in}}%
\pgfpathlineto{\pgfqpoint{2.281879in}{0.558658in}}%
\pgfpathlineto{\pgfqpoint{2.275576in}{0.555452in}}%
\pgfpathlineto{\pgfqpoint{2.277143in}{0.545445in}}%
\pgfpathlineto{\pgfqpoint{2.270570in}{0.545919in}}%
\pgfpathlineto{\pgfqpoint{2.250353in}{0.531927in}}%
\pgfpathlineto{\pgfqpoint{2.226068in}{0.526381in}}%
\pgfpathlineto{\pgfqpoint{2.225669in}{0.533243in}}%
\pgfpathlineto{\pgfqpoint{2.220287in}{0.539949in}}%
\pgfpathlineto{\pgfqpoint{2.205848in}{0.549208in}}%
\pgfpathlineto{\pgfqpoint{2.185121in}{0.558700in}}%
\pgfpathlineto{\pgfqpoint{2.162608in}{0.563725in}}%
\pgfpathlineto{\pgfqpoint{2.158382in}{0.570558in}}%
\pgfpathlineto{\pgfqpoint{2.150265in}{0.564859in}}%
\pgfpathlineto{\pgfqpoint{2.140491in}{0.563602in}}%
\pgfpathlineto{\pgfqpoint{2.106467in}{0.554637in}}%
\pgfpathlineto{\pgfqpoint{2.105919in}{0.564647in}}%
\pgfpathclose%
\pgfusepath{fill}%
\end{pgfscope}%
\begin{pgfscope}%
\pgfpathrectangle{\pgfqpoint{0.100000in}{0.100000in}}{\pgfqpoint{2.862985in}{1.832778in}}%
\pgfusepath{clip}%
\pgfsetbuttcap%
\pgfsetmiterjoin%
\definecolor{currentfill}{rgb}{0.787133,0.266922,0.480307}%
\pgfsetfillcolor{currentfill}%
\pgfsetlinewidth{0.000000pt}%
\definecolor{currentstroke}{rgb}{0.000000,0.000000,0.000000}%
\pgfsetstrokecolor{currentstroke}%
\pgfsetstrokeopacity{0.000000}%
\pgfsetdash{}{0pt}%
\pgfpathmoveto{\pgfqpoint{2.485737in}{0.501980in}}%
\pgfpathlineto{\pgfqpoint{2.494540in}{0.491740in}}%
\pgfpathlineto{\pgfqpoint{2.484372in}{0.491096in}}%
\pgfpathlineto{\pgfqpoint{2.485737in}{0.501980in}}%
\pgfpathclose%
\pgfusepath{fill}%
\end{pgfscope}%
\begin{pgfscope}%
\pgfpathrectangle{\pgfqpoint{0.100000in}{0.100000in}}{\pgfqpoint{2.862985in}{1.832778in}}%
\pgfusepath{clip}%
\pgfsetbuttcap%
\pgfsetmiterjoin%
\definecolor{currentfill}{rgb}{0.710549,0.182868,0.550004}%
\pgfsetfillcolor{currentfill}%
\pgfsetlinewidth{0.000000pt}%
\definecolor{currentstroke}{rgb}{0.000000,0.000000,0.000000}%
\pgfsetstrokecolor{currentstroke}%
\pgfsetstrokeopacity{0.000000}%
\pgfsetdash{}{0pt}%
\pgfpathmoveto{\pgfqpoint{1.967522in}{1.671190in}}%
\pgfpathlineto{\pgfqpoint{1.962746in}{1.661894in}}%
\pgfpathlineto{\pgfqpoint{1.951351in}{1.656468in}}%
\pgfpathlineto{\pgfqpoint{1.946416in}{1.649166in}}%
\pgfpathlineto{\pgfqpoint{1.940635in}{1.654425in}}%
\pgfpathlineto{\pgfqpoint{1.967522in}{1.671190in}}%
\pgfpathclose%
\pgfusepath{fill}%
\end{pgfscope}%
\begin{pgfscope}%
\pgfpathrectangle{\pgfqpoint{0.100000in}{0.100000in}}{\pgfqpoint{2.862985in}{1.832778in}}%
\pgfusepath{clip}%
\pgfsetbuttcap%
\pgfsetmiterjoin%
\definecolor{currentfill}{rgb}{0.710549,0.182868,0.550004}%
\pgfsetfillcolor{currentfill}%
\pgfsetlinewidth{0.000000pt}%
\definecolor{currentstroke}{rgb}{0.000000,0.000000,0.000000}%
\pgfsetstrokecolor{currentstroke}%
\pgfsetstrokeopacity{0.000000}%
\pgfsetdash{}{0pt}%
\pgfpathmoveto{\pgfqpoint{1.971396in}{1.615244in}}%
\pgfpathlineto{\pgfqpoint{1.983054in}{1.626120in}}%
\pgfpathlineto{\pgfqpoint{2.001037in}{1.628969in}}%
\pgfpathlineto{\pgfqpoint{1.996084in}{1.621405in}}%
\pgfpathlineto{\pgfqpoint{1.983756in}{1.610466in}}%
\pgfpathlineto{\pgfqpoint{1.976549in}{1.596426in}}%
\pgfpathlineto{\pgfqpoint{1.973496in}{1.604041in}}%
\pgfpathlineto{\pgfqpoint{1.968076in}{1.605134in}}%
\pgfpathlineto{\pgfqpoint{1.967532in}{1.612009in}}%
\pgfpathlineto{\pgfqpoint{1.971396in}{1.615244in}}%
\pgfpathclose%
\pgfusepath{fill}%
\end{pgfscope}%
\begin{pgfscope}%
\pgfpathrectangle{\pgfqpoint{0.100000in}{0.100000in}}{\pgfqpoint{2.862985in}{1.832778in}}%
\pgfusepath{clip}%
\pgfsetbuttcap%
\pgfsetmiterjoin%
\definecolor{currentfill}{rgb}{0.710549,0.182868,0.550004}%
\pgfsetfillcolor{currentfill}%
\pgfsetlinewidth{0.000000pt}%
\definecolor{currentstroke}{rgb}{0.000000,0.000000,0.000000}%
\pgfsetstrokecolor{currentstroke}%
\pgfsetstrokeopacity{0.000000}%
\pgfsetdash{}{0pt}%
\pgfpathmoveto{\pgfqpoint{2.017832in}{1.482495in}}%
\pgfpathlineto{\pgfqpoint{2.014727in}{1.485941in}}%
\pgfpathlineto{\pgfqpoint{2.017991in}{1.495641in}}%
\pgfpathlineto{\pgfqpoint{2.008301in}{1.496259in}}%
\pgfpathlineto{\pgfqpoint{2.010870in}{1.504623in}}%
\pgfpathlineto{\pgfqpoint{2.009272in}{1.517943in}}%
\pgfpathlineto{\pgfqpoint{2.000578in}{1.522561in}}%
\pgfpathlineto{\pgfqpoint{1.991480in}{1.531907in}}%
\pgfpathlineto{\pgfqpoint{1.977460in}{1.534643in}}%
\pgfpathlineto{\pgfqpoint{1.963816in}{1.534565in}}%
\pgfpathlineto{\pgfqpoint{1.950397in}{1.541198in}}%
\pgfpathlineto{\pgfqpoint{1.905502in}{1.550866in}}%
\pgfpathlineto{\pgfqpoint{1.900596in}{1.561110in}}%
\pgfpathlineto{\pgfqpoint{1.891845in}{1.564606in}}%
\pgfpathlineto{\pgfqpoint{1.908375in}{1.572449in}}%
\pgfpathlineto{\pgfqpoint{1.917713in}{1.582234in}}%
\pgfpathlineto{\pgfqpoint{1.935106in}{1.584888in}}%
\pgfpathlineto{\pgfqpoint{1.945809in}{1.594855in}}%
\pgfpathlineto{\pgfqpoint{1.951425in}{1.595254in}}%
\pgfpathlineto{\pgfqpoint{1.955710in}{1.602362in}}%
\pgfpathlineto{\pgfqpoint{1.966157in}{1.610783in}}%
\pgfpathlineto{\pgfqpoint{1.967067in}{1.604831in}}%
\pgfpathlineto{\pgfqpoint{1.971774in}{1.603602in}}%
\pgfpathlineto{\pgfqpoint{1.975316in}{1.596504in}}%
\pgfpathlineto{\pgfqpoint{1.975954in}{1.584403in}}%
\pgfpathlineto{\pgfqpoint{1.986599in}{1.591634in}}%
\pgfpathlineto{\pgfqpoint{1.999002in}{1.593145in}}%
\pgfpathlineto{\pgfqpoint{2.009590in}{1.589356in}}%
\pgfpathlineto{\pgfqpoint{2.023949in}{1.569658in}}%
\pgfpathlineto{\pgfqpoint{2.039631in}{1.572806in}}%
\pgfpathlineto{\pgfqpoint{2.046003in}{1.567535in}}%
\pgfpathlineto{\pgfqpoint{2.056252in}{1.567064in}}%
\pgfpathlineto{\pgfqpoint{2.063061in}{1.576518in}}%
\pgfpathlineto{\pgfqpoint{2.075982in}{1.584876in}}%
\pgfpathlineto{\pgfqpoint{2.088400in}{1.587493in}}%
\pgfpathlineto{\pgfqpoint{2.103808in}{1.587766in}}%
\pgfpathlineto{\pgfqpoint{2.115068in}{1.594222in}}%
\pgfpathlineto{\pgfqpoint{2.124258in}{1.591248in}}%
\pgfpathlineto{\pgfqpoint{2.126215in}{1.577594in}}%
\pgfpathlineto{\pgfqpoint{2.145930in}{1.575456in}}%
\pgfpathlineto{\pgfqpoint{2.156638in}{1.581847in}}%
\pgfpathlineto{\pgfqpoint{2.164061in}{1.567449in}}%
\pgfpathlineto{\pgfqpoint{2.178225in}{1.550799in}}%
\pgfpathlineto{\pgfqpoint{2.158401in}{1.550895in}}%
\pgfpathlineto{\pgfqpoint{2.152137in}{1.548837in}}%
\pgfpathlineto{\pgfqpoint{2.143551in}{1.551511in}}%
\pgfpathlineto{\pgfqpoint{2.142974in}{1.539985in}}%
\pgfpathlineto{\pgfqpoint{2.127384in}{1.548993in}}%
\pgfpathlineto{\pgfqpoint{2.107335in}{1.551746in}}%
\pgfpathlineto{\pgfqpoint{2.101816in}{1.542973in}}%
\pgfpathlineto{\pgfqpoint{2.075584in}{1.538707in}}%
\pgfpathlineto{\pgfqpoint{2.072566in}{1.531276in}}%
\pgfpathlineto{\pgfqpoint{2.054312in}{1.529046in}}%
\pgfpathlineto{\pgfqpoint{2.048794in}{1.521492in}}%
\pgfpathlineto{\pgfqpoint{2.039139in}{1.519466in}}%
\pgfpathlineto{\pgfqpoint{2.031427in}{1.501552in}}%
\pgfpathlineto{\pgfqpoint{2.021663in}{1.484203in}}%
\pgfpathlineto{\pgfqpoint{2.017832in}{1.482495in}}%
\pgfpathclose%
\pgfusepath{fill}%
\end{pgfscope}%
\begin{pgfscope}%
\pgfpathrectangle{\pgfqpoint{0.100000in}{0.100000in}}{\pgfqpoint{2.862985in}{1.832778in}}%
\pgfusepath{clip}%
\pgfsetbuttcap%
\pgfsetmiterjoin%
\definecolor{currentfill}{rgb}{0.710549,0.182868,0.550004}%
\pgfsetfillcolor{currentfill}%
\pgfsetlinewidth{0.000000pt}%
\definecolor{currentstroke}{rgb}{0.000000,0.000000,0.000000}%
\pgfsetstrokecolor{currentstroke}%
\pgfsetstrokeopacity{0.000000}%
\pgfsetdash{}{0pt}%
\pgfpathmoveto{\pgfqpoint{2.074029in}{1.274203in}}%
\pgfpathlineto{\pgfqpoint{2.083339in}{1.284002in}}%
\pgfpathlineto{\pgfqpoint{2.087580in}{1.298209in}}%
\pgfpathlineto{\pgfqpoint{2.092623in}{1.306443in}}%
\pgfpathlineto{\pgfqpoint{2.095716in}{1.317649in}}%
\pgfpathlineto{\pgfqpoint{2.096678in}{1.339993in}}%
\pgfpathlineto{\pgfqpoint{2.092021in}{1.361407in}}%
\pgfpathlineto{\pgfqpoint{2.076633in}{1.394197in}}%
\pgfpathlineto{\pgfqpoint{2.080800in}{1.404507in}}%
\pgfpathlineto{\pgfqpoint{2.075372in}{1.418758in}}%
\pgfpathlineto{\pgfqpoint{2.084684in}{1.438466in}}%
\pgfpathlineto{\pgfqpoint{2.083166in}{1.460442in}}%
\pgfpathlineto{\pgfqpoint{2.089653in}{1.463207in}}%
\pgfpathlineto{\pgfqpoint{2.090468in}{1.473664in}}%
\pgfpathlineto{\pgfqpoint{2.101982in}{1.480375in}}%
\pgfpathlineto{\pgfqpoint{2.108494in}{1.479291in}}%
\pgfpathlineto{\pgfqpoint{2.110324in}{1.468077in}}%
\pgfpathlineto{\pgfqpoint{2.115407in}{1.467638in}}%
\pgfpathlineto{\pgfqpoint{2.120044in}{1.483813in}}%
\pgfpathlineto{\pgfqpoint{2.118468in}{1.496393in}}%
\pgfpathlineto{\pgfqpoint{2.121483in}{1.503617in}}%
\pgfpathlineto{\pgfqpoint{2.137785in}{1.511080in}}%
\pgfpathlineto{\pgfqpoint{2.130356in}{1.513735in}}%
\pgfpathlineto{\pgfqpoint{2.127936in}{1.520146in}}%
\pgfpathlineto{\pgfqpoint{2.133288in}{1.531404in}}%
\pgfpathlineto{\pgfqpoint{2.143841in}{1.535292in}}%
\pgfpathlineto{\pgfqpoint{2.156095in}{1.528650in}}%
\pgfpathlineto{\pgfqpoint{2.167648in}{1.528567in}}%
\pgfpathlineto{\pgfqpoint{2.173015in}{1.520808in}}%
\pgfpathlineto{\pgfqpoint{2.181106in}{1.521347in}}%
\pgfpathlineto{\pgfqpoint{2.206083in}{1.510561in}}%
\pgfpathlineto{\pgfqpoint{2.211074in}{1.500120in}}%
\pgfpathlineto{\pgfqpoint{2.205624in}{1.496504in}}%
\pgfpathlineto{\pgfqpoint{2.207303in}{1.488679in}}%
\pgfpathlineto{\pgfqpoint{2.212687in}{1.485198in}}%
\pgfpathlineto{\pgfqpoint{2.215688in}{1.475579in}}%
\pgfpathlineto{\pgfqpoint{2.215268in}{1.452144in}}%
\pgfpathlineto{\pgfqpoint{2.208159in}{1.446546in}}%
\pgfpathlineto{\pgfqpoint{2.206564in}{1.434236in}}%
\pgfpathlineto{\pgfqpoint{2.193377in}{1.422859in}}%
\pgfpathlineto{\pgfqpoint{2.194159in}{1.409089in}}%
\pgfpathlineto{\pgfqpoint{2.205717in}{1.404265in}}%
\pgfpathlineto{\pgfqpoint{2.218756in}{1.421446in}}%
\pgfpathlineto{\pgfqpoint{2.219833in}{1.427678in}}%
\pgfpathlineto{\pgfqpoint{2.236094in}{1.438035in}}%
\pgfpathlineto{\pgfqpoint{2.246435in}{1.433238in}}%
\pgfpathlineto{\pgfqpoint{2.252924in}{1.422395in}}%
\pgfpathlineto{\pgfqpoint{2.263393in}{1.384749in}}%
\pgfpathlineto{\pgfqpoint{2.268940in}{1.372839in}}%
\pgfpathlineto{\pgfqpoint{2.267203in}{1.367915in}}%
\pgfpathlineto{\pgfqpoint{2.267372in}{1.351150in}}%
\pgfpathlineto{\pgfqpoint{2.257281in}{1.352759in}}%
\pgfpathlineto{\pgfqpoint{2.251581in}{1.340231in}}%
\pgfpathlineto{\pgfqpoint{2.250791in}{1.331714in}}%
\pgfpathlineto{\pgfqpoint{2.243169in}{1.326246in}}%
\pgfpathlineto{\pgfqpoint{2.241475in}{1.309636in}}%
\pgfpathlineto{\pgfqpoint{2.230390in}{1.288653in}}%
\pgfpathlineto{\pgfqpoint{2.169752in}{1.279623in}}%
\pgfpathlineto{\pgfqpoint{2.169393in}{1.283592in}}%
\pgfpathlineto{\pgfqpoint{2.128831in}{1.279317in}}%
\pgfpathlineto{\pgfqpoint{2.074029in}{1.274203in}}%
\pgfpathclose%
\pgfusepath{fill}%
\end{pgfscope}%
\begin{pgfscope}%
\pgfpathrectangle{\pgfqpoint{0.100000in}{0.100000in}}{\pgfqpoint{2.862985in}{1.832778in}}%
\pgfusepath{clip}%
\pgfsetbuttcap%
\pgfsetmiterjoin%
\definecolor{currentfill}{rgb}{0.645872,0.120898,0.598867}%
\pgfsetfillcolor{currentfill}%
\pgfsetlinewidth{0.000000pt}%
\definecolor{currentstroke}{rgb}{0.000000,0.000000,0.000000}%
\pgfsetstrokecolor{currentstroke}%
\pgfsetstrokeopacity{0.000000}%
\pgfsetdash{}{0pt}%
\pgfpathmoveto{\pgfqpoint{0.542548in}{0.431818in}}%
\pgfpathlineto{\pgfqpoint{0.541337in}{0.430752in}}%
\pgfpathlineto{\pgfqpoint{0.535781in}{0.432173in}}%
\pgfpathlineto{\pgfqpoint{0.536196in}{0.434969in}}%
\pgfpathlineto{\pgfqpoint{0.539341in}{0.436249in}}%
\pgfpathlineto{\pgfqpoint{0.543824in}{0.442694in}}%
\pgfpathlineto{\pgfqpoint{0.544569in}{0.447051in}}%
\pgfpathlineto{\pgfqpoint{0.547090in}{0.448192in}}%
\pgfpathlineto{\pgfqpoint{0.551190in}{0.448197in}}%
\pgfpathlineto{\pgfqpoint{0.555477in}{0.462333in}}%
\pgfpathlineto{\pgfqpoint{0.557157in}{0.463496in}}%
\pgfpathlineto{\pgfqpoint{0.555651in}{0.466260in}}%
\pgfpathlineto{\pgfqpoint{0.552125in}{0.463464in}}%
\pgfpathlineto{\pgfqpoint{0.541766in}{0.467233in}}%
\pgfpathlineto{\pgfqpoint{0.535583in}{0.473582in}}%
\pgfpathlineto{\pgfqpoint{0.538280in}{0.476479in}}%
\pgfpathlineto{\pgfqpoint{0.537053in}{0.480565in}}%
\pgfpathlineto{\pgfqpoint{0.537809in}{0.485392in}}%
\pgfpathlineto{\pgfqpoint{0.536479in}{0.491035in}}%
\pgfpathlineto{\pgfqpoint{0.536266in}{0.496684in}}%
\pgfpathlineto{\pgfqpoint{0.542217in}{0.495739in}}%
\pgfpathlineto{\pgfqpoint{0.543395in}{0.497404in}}%
\pgfpathlineto{\pgfqpoint{0.546355in}{0.497243in}}%
\pgfpathlineto{\pgfqpoint{0.549764in}{0.491217in}}%
\pgfpathlineto{\pgfqpoint{0.553051in}{0.486806in}}%
\pgfpathlineto{\pgfqpoint{0.550418in}{0.483935in}}%
\pgfpathlineto{\pgfqpoint{0.554970in}{0.482833in}}%
\pgfpathlineto{\pgfqpoint{0.556007in}{0.484241in}}%
\pgfpathlineto{\pgfqpoint{0.553573in}{0.489956in}}%
\pgfpathlineto{\pgfqpoint{0.548447in}{0.494692in}}%
\pgfpathlineto{\pgfqpoint{0.549074in}{0.496676in}}%
\pgfpathlineto{\pgfqpoint{0.544984in}{0.501691in}}%
\pgfpathlineto{\pgfqpoint{0.549161in}{0.502710in}}%
\pgfpathlineto{\pgfqpoint{0.545542in}{0.513448in}}%
\pgfpathlineto{\pgfqpoint{0.548484in}{0.514936in}}%
\pgfpathlineto{\pgfqpoint{0.549202in}{0.517511in}}%
\pgfpathlineto{\pgfqpoint{0.547288in}{0.520569in}}%
\pgfpathlineto{\pgfqpoint{0.551239in}{0.521335in}}%
\pgfpathlineto{\pgfqpoint{0.551948in}{0.523648in}}%
\pgfpathlineto{\pgfqpoint{0.556660in}{0.520176in}}%
\pgfpathlineto{\pgfqpoint{0.557707in}{0.523941in}}%
\pgfpathlineto{\pgfqpoint{0.560183in}{0.525597in}}%
\pgfpathlineto{\pgfqpoint{0.572233in}{0.527901in}}%
\pgfpathlineto{\pgfqpoint{0.575932in}{0.527086in}}%
\pgfpathlineto{\pgfqpoint{0.579031in}{0.531659in}}%
\pgfpathlineto{\pgfqpoint{0.586384in}{0.534658in}}%
\pgfpathlineto{\pgfqpoint{0.592047in}{0.533521in}}%
\pgfpathlineto{\pgfqpoint{0.594631in}{0.529538in}}%
\pgfpathlineto{\pgfqpoint{0.594495in}{0.523212in}}%
\pgfpathlineto{\pgfqpoint{0.597394in}{0.521778in}}%
\pgfpathlineto{\pgfqpoint{0.603306in}{0.521978in}}%
\pgfpathlineto{\pgfqpoint{0.610878in}{0.523624in}}%
\pgfpathlineto{\pgfqpoint{0.610785in}{0.521601in}}%
\pgfpathlineto{\pgfqpoint{0.620357in}{0.515832in}}%
\pgfpathlineto{\pgfqpoint{0.627052in}{0.517379in}}%
\pgfpathlineto{\pgfqpoint{0.629022in}{0.519381in}}%
\pgfpathlineto{\pgfqpoint{0.631133in}{0.524508in}}%
\pgfpathlineto{\pgfqpoint{0.633750in}{0.527792in}}%
\pgfpathlineto{\pgfqpoint{0.633984in}{0.532712in}}%
\pgfpathlineto{\pgfqpoint{0.631587in}{0.534647in}}%
\pgfpathlineto{\pgfqpoint{0.635136in}{0.536396in}}%
\pgfpathlineto{\pgfqpoint{0.641098in}{0.533768in}}%
\pgfpathlineto{\pgfqpoint{0.643040in}{0.535390in}}%
\pgfpathlineto{\pgfqpoint{0.643382in}{0.542352in}}%
\pgfpathlineto{\pgfqpoint{0.639528in}{0.540897in}}%
\pgfpathlineto{\pgfqpoint{0.636855in}{0.543309in}}%
\pgfpathlineto{\pgfqpoint{0.632516in}{0.544812in}}%
\pgfpathlineto{\pgfqpoint{0.625629in}{0.544561in}}%
\pgfpathlineto{\pgfqpoint{0.619030in}{0.546848in}}%
\pgfpathlineto{\pgfqpoint{0.618699in}{0.552545in}}%
\pgfpathlineto{\pgfqpoint{0.612797in}{0.558597in}}%
\pgfpathlineto{\pgfqpoint{0.605313in}{0.561138in}}%
\pgfpathlineto{\pgfqpoint{0.598780in}{0.572853in}}%
\pgfpathlineto{\pgfqpoint{0.599316in}{0.577189in}}%
\pgfpathlineto{\pgfqpoint{0.602833in}{0.579165in}}%
\pgfpathlineto{\pgfqpoint{0.602315in}{0.583269in}}%
\pgfpathlineto{\pgfqpoint{0.603132in}{0.587352in}}%
\pgfpathlineto{\pgfqpoint{0.605915in}{0.584557in}}%
\pgfpathlineto{\pgfqpoint{0.609590in}{0.585723in}}%
\pgfpathlineto{\pgfqpoint{0.609407in}{0.589270in}}%
\pgfpathlineto{\pgfqpoint{0.604311in}{0.596390in}}%
\pgfpathlineto{\pgfqpoint{0.602820in}{0.603734in}}%
\pgfpathlineto{\pgfqpoint{0.607085in}{0.604294in}}%
\pgfpathlineto{\pgfqpoint{0.613468in}{0.603843in}}%
\pgfpathlineto{\pgfqpoint{0.615478in}{0.601387in}}%
\pgfpathlineto{\pgfqpoint{0.621201in}{0.602266in}}%
\pgfpathlineto{\pgfqpoint{0.626520in}{0.601112in}}%
\pgfpathlineto{\pgfqpoint{0.630115in}{0.599083in}}%
\pgfpathlineto{\pgfqpoint{0.632604in}{0.600010in}}%
\pgfpathlineto{\pgfqpoint{0.639212in}{0.597983in}}%
\pgfpathlineto{\pgfqpoint{0.639273in}{0.596187in}}%
\pgfpathlineto{\pgfqpoint{0.644305in}{0.596794in}}%
\pgfpathlineto{\pgfqpoint{0.651053in}{0.592414in}}%
\pgfpathlineto{\pgfqpoint{0.650791in}{0.589506in}}%
\pgfpathlineto{\pgfqpoint{0.647580in}{0.588164in}}%
\pgfpathlineto{\pgfqpoint{0.644167in}{0.584972in}}%
\pgfpathlineto{\pgfqpoint{0.643755in}{0.580588in}}%
\pgfpathlineto{\pgfqpoint{0.651106in}{0.574997in}}%
\pgfpathlineto{\pgfqpoint{0.649930in}{0.572318in}}%
\pgfpathlineto{\pgfqpoint{0.655161in}{0.570260in}}%
\pgfpathlineto{\pgfqpoint{0.656445in}{0.566900in}}%
\pgfpathlineto{\pgfqpoint{0.663083in}{0.569570in}}%
\pgfpathlineto{\pgfqpoint{0.666007in}{0.566071in}}%
\pgfpathlineto{\pgfqpoint{0.667356in}{0.568189in}}%
\pgfpathlineto{\pgfqpoint{0.663383in}{0.575083in}}%
\pgfpathlineto{\pgfqpoint{0.664795in}{0.577148in}}%
\pgfpathlineto{\pgfqpoint{0.665421in}{0.582594in}}%
\pgfpathlineto{\pgfqpoint{0.663838in}{0.584991in}}%
\pgfpathlineto{\pgfqpoint{0.665025in}{0.588211in}}%
\pgfpathlineto{\pgfqpoint{0.668915in}{0.587908in}}%
\pgfpathlineto{\pgfqpoint{0.668110in}{0.582429in}}%
\pgfpathlineto{\pgfqpoint{0.665627in}{0.580602in}}%
\pgfpathlineto{\pgfqpoint{0.665803in}{0.573231in}}%
\pgfpathlineto{\pgfqpoint{0.670258in}{0.572350in}}%
\pgfpathlineto{\pgfqpoint{0.670312in}{0.566870in}}%
\pgfpathlineto{\pgfqpoint{0.675178in}{0.563315in}}%
\pgfpathlineto{\pgfqpoint{0.677415in}{0.565557in}}%
\pgfpathlineto{\pgfqpoint{0.674970in}{0.572380in}}%
\pgfpathlineto{\pgfqpoint{0.672442in}{0.574106in}}%
\pgfpathlineto{\pgfqpoint{0.669955in}{0.572970in}}%
\pgfpathlineto{\pgfqpoint{0.667767in}{0.574420in}}%
\pgfpathlineto{\pgfqpoint{0.668149in}{0.580725in}}%
\pgfpathlineto{\pgfqpoint{0.671800in}{0.583304in}}%
\pgfpathlineto{\pgfqpoint{0.675395in}{0.583066in}}%
\pgfpathlineto{\pgfqpoint{0.673905in}{0.586634in}}%
\pgfpathlineto{\pgfqpoint{0.668518in}{0.589366in}}%
\pgfpathlineto{\pgfqpoint{0.661383in}{0.600253in}}%
\pgfpathlineto{\pgfqpoint{0.664944in}{0.605475in}}%
\pgfpathlineto{\pgfqpoint{0.666926in}{0.610393in}}%
\pgfpathlineto{\pgfqpoint{0.667082in}{0.620269in}}%
\pgfpathlineto{\pgfqpoint{0.666319in}{0.628601in}}%
\pgfpathlineto{\pgfqpoint{0.664142in}{0.634072in}}%
\pgfpathlineto{\pgfqpoint{0.664576in}{0.639541in}}%
\pgfpathlineto{\pgfqpoint{0.670144in}{0.644285in}}%
\pgfpathlineto{\pgfqpoint{0.675681in}{0.649922in}}%
\pgfpathlineto{\pgfqpoint{0.689564in}{0.638405in}}%
\pgfpathlineto{\pgfqpoint{0.698640in}{0.637556in}}%
\pgfpathlineto{\pgfqpoint{0.706132in}{0.640424in}}%
\pgfpathlineto{\pgfqpoint{0.712708in}{0.644201in}}%
\pgfpathlineto{\pgfqpoint{0.714105in}{0.646000in}}%
\pgfpathlineto{\pgfqpoint{0.730340in}{0.650103in}}%
\pgfpathlineto{\pgfqpoint{0.731749in}{0.647098in}}%
\pgfpathlineto{\pgfqpoint{0.737266in}{0.644638in}}%
\pgfpathlineto{\pgfqpoint{0.747793in}{0.645392in}}%
\pgfpathlineto{\pgfqpoint{0.749716in}{0.644720in}}%
\pgfpathlineto{\pgfqpoint{0.754448in}{0.646515in}}%
\pgfpathlineto{\pgfqpoint{0.756350in}{0.643411in}}%
\pgfpathlineto{\pgfqpoint{0.760894in}{0.640575in}}%
\pgfpathlineto{\pgfqpoint{0.763167in}{0.640896in}}%
\pgfpathlineto{\pgfqpoint{0.767566in}{0.637488in}}%
\pgfpathlineto{\pgfqpoint{0.775000in}{0.638115in}}%
\pgfpathlineto{\pgfqpoint{0.782425in}{0.640874in}}%
\pgfpathlineto{\pgfqpoint{0.788489in}{0.632051in}}%
\pgfpathlineto{\pgfqpoint{0.788019in}{0.630071in}}%
\pgfpathlineto{\pgfqpoint{0.781328in}{0.627796in}}%
\pgfpathlineto{\pgfqpoint{0.783115in}{0.624355in}}%
\pgfpathlineto{\pgfqpoint{0.789275in}{0.627976in}}%
\pgfpathlineto{\pgfqpoint{0.791790in}{0.627100in}}%
\pgfpathlineto{\pgfqpoint{0.792860in}{0.622768in}}%
\pgfpathlineto{\pgfqpoint{0.789183in}{0.620936in}}%
\pgfpathlineto{\pgfqpoint{0.791809in}{0.615461in}}%
\pgfpathlineto{\pgfqpoint{0.795623in}{0.616386in}}%
\pgfpathlineto{\pgfqpoint{0.801157in}{0.613560in}}%
\pgfpathlineto{\pgfqpoint{0.806450in}{0.605718in}}%
\pgfpathlineto{\pgfqpoint{0.801793in}{0.603540in}}%
\pgfpathlineto{\pgfqpoint{0.805907in}{0.597761in}}%
\pgfpathlineto{\pgfqpoint{0.802441in}{0.596410in}}%
\pgfpathlineto{\pgfqpoint{0.806753in}{0.590922in}}%
\pgfpathlineto{\pgfqpoint{0.812012in}{0.591506in}}%
\pgfpathlineto{\pgfqpoint{0.814781in}{0.589037in}}%
\pgfpathlineto{\pgfqpoint{0.821785in}{0.585396in}}%
\pgfpathlineto{\pgfqpoint{0.830820in}{0.572470in}}%
\pgfpathlineto{\pgfqpoint{0.830616in}{0.569379in}}%
\pgfpathlineto{\pgfqpoint{0.844166in}{0.559477in}}%
\pgfpathlineto{\pgfqpoint{0.844458in}{0.555975in}}%
\pgfpathlineto{\pgfqpoint{0.848069in}{0.550735in}}%
\pgfpathlineto{\pgfqpoint{0.859163in}{0.548115in}}%
\pgfpathlineto{\pgfqpoint{0.863277in}{0.546252in}}%
\pgfpathlineto{\pgfqpoint{0.867148in}{0.540476in}}%
\pgfpathlineto{\pgfqpoint{0.867646in}{0.535415in}}%
\pgfpathlineto{\pgfqpoint{0.870973in}{0.531291in}}%
\pgfpathlineto{\pgfqpoint{0.871972in}{0.526640in}}%
\pgfpathlineto{\pgfqpoint{0.875090in}{0.524941in}}%
\pgfpathlineto{\pgfqpoint{0.859646in}{0.497067in}}%
\pgfpathlineto{\pgfqpoint{0.827590in}{0.439200in}}%
\pgfpathlineto{\pgfqpoint{0.780557in}{0.354280in}}%
\pgfpathlineto{\pgfqpoint{0.762143in}{0.321025in}}%
\pgfpathlineto{\pgfqpoint{0.766065in}{0.316559in}}%
\pgfpathlineto{\pgfqpoint{0.767793in}{0.317984in}}%
\pgfpathlineto{\pgfqpoint{0.771296in}{0.312785in}}%
\pgfpathlineto{\pgfqpoint{0.776159in}{0.314293in}}%
\pgfpathlineto{\pgfqpoint{0.782614in}{0.311178in}}%
\pgfpathlineto{\pgfqpoint{0.778433in}{0.306310in}}%
\pgfpathlineto{\pgfqpoint{0.781662in}{0.299741in}}%
\pgfpathlineto{\pgfqpoint{0.780995in}{0.296543in}}%
\pgfpathlineto{\pgfqpoint{0.786145in}{0.279677in}}%
\pgfpathlineto{\pgfqpoint{0.783966in}{0.272007in}}%
\pgfpathlineto{\pgfqpoint{0.793805in}{0.273906in}}%
\pgfpathlineto{\pgfqpoint{0.796221in}{0.272662in}}%
\pgfpathlineto{\pgfqpoint{0.798844in}{0.274685in}}%
\pgfpathlineto{\pgfqpoint{0.800709in}{0.278509in}}%
\pgfpathlineto{\pgfqpoint{0.803381in}{0.280701in}}%
\pgfpathlineto{\pgfqpoint{0.814784in}{0.280424in}}%
\pgfpathlineto{\pgfqpoint{0.817361in}{0.273064in}}%
\pgfpathlineto{\pgfqpoint{0.815149in}{0.266603in}}%
\pgfpathlineto{\pgfqpoint{0.817662in}{0.264471in}}%
\pgfpathlineto{\pgfqpoint{0.818983in}{0.260825in}}%
\pgfpathlineto{\pgfqpoint{0.818695in}{0.253967in}}%
\pgfpathlineto{\pgfqpoint{0.822048in}{0.249051in}}%
\pgfpathlineto{\pgfqpoint{0.823886in}{0.241346in}}%
\pgfpathlineto{\pgfqpoint{0.822942in}{0.237184in}}%
\pgfpathlineto{\pgfqpoint{0.824080in}{0.232966in}}%
\pgfpathlineto{\pgfqpoint{0.825722in}{0.208539in}}%
\pgfpathlineto{\pgfqpoint{0.823370in}{0.206615in}}%
\pgfpathlineto{\pgfqpoint{0.826472in}{0.203957in}}%
\pgfpathlineto{\pgfqpoint{0.824034in}{0.200723in}}%
\pgfpathlineto{\pgfqpoint{0.826233in}{0.198087in}}%
\pgfpathlineto{\pgfqpoint{0.824817in}{0.193586in}}%
\pgfpathlineto{\pgfqpoint{0.827984in}{0.192366in}}%
\pgfpathlineto{\pgfqpoint{0.832058in}{0.185522in}}%
\pgfpathlineto{\pgfqpoint{0.835085in}{0.183460in}}%
\pgfpathlineto{\pgfqpoint{0.835877in}{0.180445in}}%
\pgfpathlineto{\pgfqpoint{0.837631in}{0.179139in}}%
\pgfpathlineto{\pgfqpoint{0.837265in}{0.176636in}}%
\pgfpathlineto{\pgfqpoint{0.841046in}{0.174822in}}%
\pgfpathlineto{\pgfqpoint{0.840112in}{0.170006in}}%
\pgfpathlineto{\pgfqpoint{0.836858in}{0.167483in}}%
\pgfpathlineto{\pgfqpoint{0.835380in}{0.162795in}}%
\pgfpathlineto{\pgfqpoint{0.834988in}{0.156518in}}%
\pgfpathlineto{\pgfqpoint{0.833180in}{0.155917in}}%
\pgfpathlineto{\pgfqpoint{0.827271in}{0.150415in}}%
\pgfpathlineto{\pgfqpoint{0.822019in}{0.148962in}}%
\pgfpathlineto{\pgfqpoint{0.819513in}{0.151181in}}%
\pgfpathlineto{\pgfqpoint{0.821935in}{0.157071in}}%
\pgfpathlineto{\pgfqpoint{0.821195in}{0.159962in}}%
\pgfpathlineto{\pgfqpoint{0.824665in}{0.161426in}}%
\pgfpathlineto{\pgfqpoint{0.827909in}{0.169947in}}%
\pgfpathlineto{\pgfqpoint{0.826539in}{0.177225in}}%
\pgfpathlineto{\pgfqpoint{0.824346in}{0.178943in}}%
\pgfpathlineto{\pgfqpoint{0.817012in}{0.178387in}}%
\pgfpathlineto{\pgfqpoint{0.818490in}{0.176486in}}%
\pgfpathlineto{\pgfqpoint{0.813181in}{0.171028in}}%
\pgfpathlineto{\pgfqpoint{0.811471in}{0.174063in}}%
\pgfpathlineto{\pgfqpoint{0.811916in}{0.178027in}}%
\pgfpathlineto{\pgfqpoint{0.814963in}{0.177954in}}%
\pgfpathlineto{\pgfqpoint{0.817677in}{0.180866in}}%
\pgfpathlineto{\pgfqpoint{0.819282in}{0.185118in}}%
\pgfpathlineto{\pgfqpoint{0.825331in}{0.183918in}}%
\pgfpathlineto{\pgfqpoint{0.820488in}{0.186190in}}%
\pgfpathlineto{\pgfqpoint{0.820915in}{0.189233in}}%
\pgfpathlineto{\pgfqpoint{0.818899in}{0.190607in}}%
\pgfpathlineto{\pgfqpoint{0.818074in}{0.194752in}}%
\pgfpathlineto{\pgfqpoint{0.819505in}{0.196911in}}%
\pgfpathlineto{\pgfqpoint{0.816093in}{0.203617in}}%
\pgfpathlineto{\pgfqpoint{0.816381in}{0.207883in}}%
\pgfpathlineto{\pgfqpoint{0.812383in}{0.211682in}}%
\pgfpathlineto{\pgfqpoint{0.810220in}{0.214756in}}%
\pgfpathlineto{\pgfqpoint{0.813207in}{0.217685in}}%
\pgfpathlineto{\pgfqpoint{0.814313in}{0.222063in}}%
\pgfpathlineto{\pgfqpoint{0.813326in}{0.224937in}}%
\pgfpathlineto{\pgfqpoint{0.814481in}{0.228366in}}%
\pgfpathlineto{\pgfqpoint{0.813009in}{0.239543in}}%
\pgfpathlineto{\pgfqpoint{0.810807in}{0.244614in}}%
\pgfpathlineto{\pgfqpoint{0.808288in}{0.246519in}}%
\pgfpathlineto{\pgfqpoint{0.809360in}{0.253116in}}%
\pgfpathlineto{\pgfqpoint{0.808604in}{0.258268in}}%
\pgfpathlineto{\pgfqpoint{0.809974in}{0.261572in}}%
\pgfpathlineto{\pgfqpoint{0.807140in}{0.262035in}}%
\pgfpathlineto{\pgfqpoint{0.806380in}{0.254009in}}%
\pgfpathlineto{\pgfqpoint{0.803196in}{0.245055in}}%
\pgfpathlineto{\pgfqpoint{0.800639in}{0.247094in}}%
\pgfpathlineto{\pgfqpoint{0.800297in}{0.250037in}}%
\pgfpathlineto{\pgfqpoint{0.797168in}{0.254057in}}%
\pgfpathlineto{\pgfqpoint{0.798758in}{0.256688in}}%
\pgfpathlineto{\pgfqpoint{0.798264in}{0.262613in}}%
\pgfpathlineto{\pgfqpoint{0.794215in}{0.260530in}}%
\pgfpathlineto{\pgfqpoint{0.795009in}{0.257316in}}%
\pgfpathlineto{\pgfqpoint{0.794247in}{0.253236in}}%
\pgfpathlineto{\pgfqpoint{0.789702in}{0.253244in}}%
\pgfpathlineto{\pgfqpoint{0.787665in}{0.255773in}}%
\pgfpathlineto{\pgfqpoint{0.784590in}{0.256255in}}%
\pgfpathlineto{\pgfqpoint{0.782521in}{0.258978in}}%
\pgfpathlineto{\pgfqpoint{0.778961in}{0.266454in}}%
\pgfpathlineto{\pgfqpoint{0.777807in}{0.272259in}}%
\pgfpathlineto{\pgfqpoint{0.778546in}{0.274726in}}%
\pgfpathlineto{\pgfqpoint{0.776822in}{0.280160in}}%
\pgfpathlineto{\pgfqpoint{0.773967in}{0.283305in}}%
\pgfpathlineto{\pgfqpoint{0.769029in}{0.291344in}}%
\pgfpathlineto{\pgfqpoint{0.765447in}{0.298326in}}%
\pgfpathlineto{\pgfqpoint{0.769102in}{0.298431in}}%
\pgfpathlineto{\pgfqpoint{0.771252in}{0.299897in}}%
\pgfpathlineto{\pgfqpoint{0.771675in}{0.304435in}}%
\pgfpathlineto{\pgfqpoint{0.761334in}{0.304983in}}%
\pgfpathlineto{\pgfqpoint{0.752539in}{0.314807in}}%
\pgfpathlineto{\pgfqpoint{0.754852in}{0.317392in}}%
\pgfpathlineto{\pgfqpoint{0.750599in}{0.317833in}}%
\pgfpathlineto{\pgfqpoint{0.742152in}{0.326967in}}%
\pgfpathlineto{\pgfqpoint{0.728770in}{0.331547in}}%
\pgfpathlineto{\pgfqpoint{0.726968in}{0.336915in}}%
\pgfpathlineto{\pgfqpoint{0.724191in}{0.339711in}}%
\pgfpathlineto{\pgfqpoint{0.724065in}{0.342336in}}%
\pgfpathlineto{\pgfqpoint{0.721976in}{0.344286in}}%
\pgfpathlineto{\pgfqpoint{0.724717in}{0.347544in}}%
\pgfpathlineto{\pgfqpoint{0.718465in}{0.347858in}}%
\pgfpathlineto{\pgfqpoint{0.716876in}{0.352075in}}%
\pgfpathlineto{\pgfqpoint{0.714030in}{0.353870in}}%
\pgfpathlineto{\pgfqpoint{0.719729in}{0.356037in}}%
\pgfpathlineto{\pgfqpoint{0.716765in}{0.359902in}}%
\pgfpathlineto{\pgfqpoint{0.711311in}{0.361524in}}%
\pgfpathlineto{\pgfqpoint{0.712160in}{0.368663in}}%
\pgfpathlineto{\pgfqpoint{0.710596in}{0.371190in}}%
\pgfpathlineto{\pgfqpoint{0.707780in}{0.372970in}}%
\pgfpathlineto{\pgfqpoint{0.705372in}{0.371185in}}%
\pgfpathlineto{\pgfqpoint{0.703950in}{0.372864in}}%
\pgfpathlineto{\pgfqpoint{0.699800in}{0.373345in}}%
\pgfpathlineto{\pgfqpoint{0.700134in}{0.379823in}}%
\pgfpathlineto{\pgfqpoint{0.694119in}{0.375448in}}%
\pgfpathlineto{\pgfqpoint{0.694246in}{0.366430in}}%
\pgfpathlineto{\pgfqpoint{0.687788in}{0.364702in}}%
\pgfpathlineto{\pgfqpoint{0.683572in}{0.361101in}}%
\pgfpathlineto{\pgfqpoint{0.680182in}{0.360342in}}%
\pgfpathlineto{\pgfqpoint{0.676514in}{0.363971in}}%
\pgfpathlineto{\pgfqpoint{0.673248in}{0.363490in}}%
\pgfpathlineto{\pgfqpoint{0.673854in}{0.366805in}}%
\pgfpathlineto{\pgfqpoint{0.669211in}{0.365020in}}%
\pgfpathlineto{\pgfqpoint{0.664909in}{0.364820in}}%
\pgfpathlineto{\pgfqpoint{0.659874in}{0.363010in}}%
\pgfpathlineto{\pgfqpoint{0.649656in}{0.364028in}}%
\pgfpathlineto{\pgfqpoint{0.646376in}{0.363358in}}%
\pgfpathlineto{\pgfqpoint{0.644009in}{0.366339in}}%
\pgfpathlineto{\pgfqpoint{0.639640in}{0.366637in}}%
\pgfpathlineto{\pgfqpoint{0.638731in}{0.370500in}}%
\pgfpathlineto{\pgfqpoint{0.641335in}{0.372604in}}%
\pgfpathlineto{\pgfqpoint{0.644390in}{0.372406in}}%
\pgfpathlineto{\pgfqpoint{0.646940in}{0.369537in}}%
\pgfpathlineto{\pgfqpoint{0.648619in}{0.373944in}}%
\pgfpathlineto{\pgfqpoint{0.646276in}{0.378827in}}%
\pgfpathlineto{\pgfqpoint{0.651579in}{0.383125in}}%
\pgfpathlineto{\pgfqpoint{0.656893in}{0.384709in}}%
\pgfpathlineto{\pgfqpoint{0.660593in}{0.387292in}}%
\pgfpathlineto{\pgfqpoint{0.662990in}{0.390076in}}%
\pgfpathlineto{\pgfqpoint{0.664327in}{0.394638in}}%
\pgfpathlineto{\pgfqpoint{0.668552in}{0.393396in}}%
\pgfpathlineto{\pgfqpoint{0.677714in}{0.394034in}}%
\pgfpathlineto{\pgfqpoint{0.679680in}{0.389032in}}%
\pgfpathlineto{\pgfqpoint{0.682894in}{0.388814in}}%
\pgfpathlineto{\pgfqpoint{0.688742in}{0.381156in}}%
\pgfpathlineto{\pgfqpoint{0.688753in}{0.383947in}}%
\pgfpathlineto{\pgfqpoint{0.686710in}{0.384855in}}%
\pgfpathlineto{\pgfqpoint{0.682971in}{0.393996in}}%
\pgfpathlineto{\pgfqpoint{0.684790in}{0.395238in}}%
\pgfpathlineto{\pgfqpoint{0.679797in}{0.400007in}}%
\pgfpathlineto{\pgfqpoint{0.675334in}{0.400936in}}%
\pgfpathlineto{\pgfqpoint{0.671081in}{0.399287in}}%
\pgfpathlineto{\pgfqpoint{0.663798in}{0.400571in}}%
\pgfpathlineto{\pgfqpoint{0.660344in}{0.398261in}}%
\pgfpathlineto{\pgfqpoint{0.652478in}{0.396476in}}%
\pgfpathlineto{\pgfqpoint{0.651759in}{0.393907in}}%
\pgfpathlineto{\pgfqpoint{0.645716in}{0.393207in}}%
\pgfpathlineto{\pgfqpoint{0.644165in}{0.390270in}}%
\pgfpathlineto{\pgfqpoint{0.640709in}{0.388119in}}%
\pgfpathlineto{\pgfqpoint{0.636533in}{0.388489in}}%
\pgfpathlineto{\pgfqpoint{0.634448in}{0.386408in}}%
\pgfpathlineto{\pgfqpoint{0.631818in}{0.386506in}}%
\pgfpathlineto{\pgfqpoint{0.629304in}{0.388378in}}%
\pgfpathlineto{\pgfqpoint{0.624909in}{0.388044in}}%
\pgfpathlineto{\pgfqpoint{0.623874in}{0.386318in}}%
\pgfpathlineto{\pgfqpoint{0.619258in}{0.388090in}}%
\pgfpathlineto{\pgfqpoint{0.615339in}{0.383518in}}%
\pgfpathlineto{\pgfqpoint{0.619097in}{0.379095in}}%
\pgfpathlineto{\pgfqpoint{0.620543in}{0.375457in}}%
\pgfpathlineto{\pgfqpoint{0.619210in}{0.372376in}}%
\pgfpathlineto{\pgfqpoint{0.613870in}{0.370250in}}%
\pgfpathlineto{\pgfqpoint{0.610776in}{0.371980in}}%
\pgfpathlineto{\pgfqpoint{0.606882in}{0.371043in}}%
\pgfpathlineto{\pgfqpoint{0.606413in}{0.368343in}}%
\pgfpathlineto{\pgfqpoint{0.601673in}{0.369298in}}%
\pgfpathlineto{\pgfqpoint{0.602722in}{0.366600in}}%
\pgfpathlineto{\pgfqpoint{0.595731in}{0.366031in}}%
\pgfpathlineto{\pgfqpoint{0.591205in}{0.369456in}}%
\pgfpathlineto{\pgfqpoint{0.588787in}{0.367278in}}%
\pgfpathlineto{\pgfqpoint{0.582354in}{0.369597in}}%
\pgfpathlineto{\pgfqpoint{0.580752in}{0.367248in}}%
\pgfpathlineto{\pgfqpoint{0.577563in}{0.366203in}}%
\pgfpathlineto{\pgfqpoint{0.575005in}{0.368957in}}%
\pgfpathlineto{\pgfqpoint{0.566257in}{0.368272in}}%
\pgfpathlineto{\pgfqpoint{0.566266in}{0.364211in}}%
\pgfpathlineto{\pgfqpoint{0.562752in}{0.363136in}}%
\pgfpathlineto{\pgfqpoint{0.560018in}{0.365224in}}%
\pgfpathlineto{\pgfqpoint{0.556404in}{0.364064in}}%
\pgfpathlineto{\pgfqpoint{0.551727in}{0.363787in}}%
\pgfpathlineto{\pgfqpoint{0.544759in}{0.366724in}}%
\pgfpathlineto{\pgfqpoint{0.540769in}{0.363696in}}%
\pgfpathlineto{\pgfqpoint{0.535063in}{0.367743in}}%
\pgfpathlineto{\pgfqpoint{0.533197in}{0.366209in}}%
\pgfpathlineto{\pgfqpoint{0.531887in}{0.362246in}}%
\pgfpathlineto{\pgfqpoint{0.528218in}{0.359922in}}%
\pgfpathlineto{\pgfqpoint{0.522236in}{0.362582in}}%
\pgfpathlineto{\pgfqpoint{0.512326in}{0.369153in}}%
\pgfpathlineto{\pgfqpoint{0.509438in}{0.369737in}}%
\pgfpathlineto{\pgfqpoint{0.507808in}{0.368368in}}%
\pgfpathlineto{\pgfqpoint{0.503449in}{0.369299in}}%
\pgfpathlineto{\pgfqpoint{0.501322in}{0.371344in}}%
\pgfpathlineto{\pgfqpoint{0.496969in}{0.372446in}}%
\pgfpathlineto{\pgfqpoint{0.491023in}{0.372510in}}%
\pgfpathlineto{\pgfqpoint{0.488662in}{0.374950in}}%
\pgfpathlineto{\pgfqpoint{0.493187in}{0.377474in}}%
\pgfpathlineto{\pgfqpoint{0.492131in}{0.379966in}}%
\pgfpathlineto{\pgfqpoint{0.489022in}{0.379406in}}%
\pgfpathlineto{\pgfqpoint{0.487470in}{0.377447in}}%
\pgfpathlineto{\pgfqpoint{0.480192in}{0.374639in}}%
\pgfpathlineto{\pgfqpoint{0.476512in}{0.375594in}}%
\pgfpathlineto{\pgfqpoint{0.474964in}{0.381297in}}%
\pgfpathlineto{\pgfqpoint{0.470373in}{0.378845in}}%
\pgfpathlineto{\pgfqpoint{0.469039in}{0.386073in}}%
\pgfpathlineto{\pgfqpoint{0.466864in}{0.384210in}}%
\pgfpathlineto{\pgfqpoint{0.467228in}{0.381429in}}%
\pgfpathlineto{\pgfqpoint{0.462226in}{0.382163in}}%
\pgfpathlineto{\pgfqpoint{0.467429in}{0.387010in}}%
\pgfpathlineto{\pgfqpoint{0.472702in}{0.384130in}}%
\pgfpathlineto{\pgfqpoint{0.473833in}{0.385576in}}%
\pgfpathlineto{\pgfqpoint{0.479380in}{0.383843in}}%
\pgfpathlineto{\pgfqpoint{0.479033in}{0.385816in}}%
\pgfpathlineto{\pgfqpoint{0.494168in}{0.386469in}}%
\pgfpathlineto{\pgfqpoint{0.501551in}{0.383391in}}%
\pgfpathlineto{\pgfqpoint{0.504408in}{0.380355in}}%
\pgfpathlineto{\pgfqpoint{0.503561in}{0.377705in}}%
\pgfpathlineto{\pgfqpoint{0.506429in}{0.376702in}}%
\pgfpathlineto{\pgfqpoint{0.513558in}{0.380885in}}%
\pgfpathlineto{\pgfqpoint{0.522779in}{0.381153in}}%
\pgfpathlineto{\pgfqpoint{0.531155in}{0.378631in}}%
\pgfpathlineto{\pgfqpoint{0.535899in}{0.379054in}}%
\pgfpathlineto{\pgfqpoint{0.537665in}{0.375142in}}%
\pgfpathlineto{\pgfqpoint{0.540318in}{0.379455in}}%
\pgfpathlineto{\pgfqpoint{0.541950in}{0.380389in}}%
\pgfpathlineto{\pgfqpoint{0.549498in}{0.382264in}}%
\pgfpathlineto{\pgfqpoint{0.552304in}{0.381118in}}%
\pgfpathlineto{\pgfqpoint{0.555338in}{0.382479in}}%
\pgfpathlineto{\pgfqpoint{0.558958in}{0.381767in}}%
\pgfpathlineto{\pgfqpoint{0.559812in}{0.383352in}}%
\pgfpathlineto{\pgfqpoint{0.567363in}{0.390625in}}%
\pgfpathlineto{\pgfqpoint{0.570561in}{0.391703in}}%
\pgfpathlineto{\pgfqpoint{0.572451in}{0.395629in}}%
\pgfpathlineto{\pgfqpoint{0.582481in}{0.398668in}}%
\pgfpathlineto{\pgfqpoint{0.583737in}{0.400987in}}%
\pgfpathlineto{\pgfqpoint{0.569522in}{0.404566in}}%
\pgfpathlineto{\pgfqpoint{0.570081in}{0.408415in}}%
\pgfpathlineto{\pgfqpoint{0.568619in}{0.413499in}}%
\pgfpathlineto{\pgfqpoint{0.565309in}{0.413077in}}%
\pgfpathlineto{\pgfqpoint{0.563057in}{0.406479in}}%
\pgfpathlineto{\pgfqpoint{0.560443in}{0.405936in}}%
\pgfpathlineto{\pgfqpoint{0.558742in}{0.408071in}}%
\pgfpathlineto{\pgfqpoint{0.560726in}{0.417413in}}%
\pgfpathlineto{\pgfqpoint{0.558750in}{0.421328in}}%
\pgfpathlineto{\pgfqpoint{0.554922in}{0.425943in}}%
\pgfpathlineto{\pgfqpoint{0.556762in}{0.429465in}}%
\pgfpathlineto{\pgfqpoint{0.548683in}{0.429591in}}%
\pgfpathlineto{\pgfqpoint{0.548867in}{0.430958in}}%
\pgfpathlineto{\pgfqpoint{0.544046in}{0.432871in}}%
\pgfpathlineto{\pgfqpoint{0.542548in}{0.431818in}}%
\pgfpathclose%
\pgfusepath{fill}%
\end{pgfscope}%
\begin{pgfscope}%
\pgfpathrectangle{\pgfqpoint{0.100000in}{0.100000in}}{\pgfqpoint{2.862985in}{1.832778in}}%
\pgfusepath{clip}%
\pgfsetbuttcap%
\pgfsetmiterjoin%
\definecolor{currentfill}{rgb}{0.645872,0.120898,0.598867}%
\pgfsetfillcolor{currentfill}%
\pgfsetlinewidth{0.000000pt}%
\definecolor{currentstroke}{rgb}{0.000000,0.000000,0.000000}%
\pgfsetstrokecolor{currentstroke}%
\pgfsetstrokeopacity{0.000000}%
\pgfsetdash{}{0pt}%
\pgfpathmoveto{\pgfqpoint{0.528131in}{0.499862in}}%
\pgfpathlineto{\pgfqpoint{0.527132in}{0.498106in}}%
\pgfpathlineto{\pgfqpoint{0.529775in}{0.494368in}}%
\pgfpathlineto{\pgfqpoint{0.525597in}{0.490871in}}%
\pgfpathlineto{\pgfqpoint{0.525713in}{0.487863in}}%
\pgfpathlineto{\pgfqpoint{0.523824in}{0.487153in}}%
\pgfpathlineto{\pgfqpoint{0.517331in}{0.491639in}}%
\pgfpathlineto{\pgfqpoint{0.512703in}{0.501458in}}%
\pgfpathlineto{\pgfqpoint{0.512393in}{0.504295in}}%
\pgfpathlineto{\pgfqpoint{0.513583in}{0.507691in}}%
\pgfpathlineto{\pgfqpoint{0.518598in}{0.502496in}}%
\pgfpathlineto{\pgfqpoint{0.520132in}{0.503498in}}%
\pgfpathlineto{\pgfqpoint{0.524386in}{0.502545in}}%
\pgfpathlineto{\pgfqpoint{0.528131in}{0.499862in}}%
\pgfpathclose%
\pgfusepath{fill}%
\end{pgfscope}%
\begin{pgfscope}%
\pgfpathrectangle{\pgfqpoint{0.100000in}{0.100000in}}{\pgfqpoint{2.862985in}{1.832778in}}%
\pgfusepath{clip}%
\pgfsetbuttcap%
\pgfsetmiterjoin%
\definecolor{currentfill}{rgb}{0.645872,0.120898,0.598867}%
\pgfsetfillcolor{currentfill}%
\pgfsetlinewidth{0.000000pt}%
\definecolor{currentstroke}{rgb}{0.000000,0.000000,0.000000}%
\pgfsetstrokecolor{currentstroke}%
\pgfsetstrokeopacity{0.000000}%
\pgfsetdash{}{0pt}%
\pgfpathmoveto{\pgfqpoint{0.451126in}{0.386109in}}%
\pgfpathlineto{\pgfqpoint{0.446562in}{0.385368in}}%
\pgfpathlineto{\pgfqpoint{0.443260in}{0.386895in}}%
\pgfpathlineto{\pgfqpoint{0.441816in}{0.389380in}}%
\pgfpathlineto{\pgfqpoint{0.443453in}{0.392945in}}%
\pgfpathlineto{\pgfqpoint{0.447096in}{0.392002in}}%
\pgfpathlineto{\pgfqpoint{0.453086in}{0.394117in}}%
\pgfpathlineto{\pgfqpoint{0.455184in}{0.391196in}}%
\pgfpathlineto{\pgfqpoint{0.457217in}{0.391735in}}%
\pgfpathlineto{\pgfqpoint{0.462141in}{0.389917in}}%
\pgfpathlineto{\pgfqpoint{0.464268in}{0.387608in}}%
\pgfpathlineto{\pgfqpoint{0.461672in}{0.381567in}}%
\pgfpathlineto{\pgfqpoint{0.459090in}{0.380286in}}%
\pgfpathlineto{\pgfqpoint{0.456769in}{0.380976in}}%
\pgfpathlineto{\pgfqpoint{0.451126in}{0.386109in}}%
\pgfpathclose%
\pgfusepath{fill}%
\end{pgfscope}%
\begin{pgfscope}%
\pgfpathrectangle{\pgfqpoint{0.100000in}{0.100000in}}{\pgfqpoint{2.862985in}{1.832778in}}%
\pgfusepath{clip}%
\pgfsetbuttcap%
\pgfsetmiterjoin%
\definecolor{currentfill}{rgb}{0.645872,0.120898,0.598867}%
\pgfsetfillcolor{currentfill}%
\pgfsetlinewidth{0.000000pt}%
\definecolor{currentstroke}{rgb}{0.000000,0.000000,0.000000}%
\pgfsetstrokecolor{currentstroke}%
\pgfsetstrokeopacity{0.000000}%
\pgfsetdash{}{0pt}%
\pgfpathmoveto{\pgfqpoint{0.410323in}{0.391598in}}%
\pgfpathlineto{\pgfqpoint{0.403857in}{0.394682in}}%
\pgfpathlineto{\pgfqpoint{0.397873in}{0.395600in}}%
\pgfpathlineto{\pgfqpoint{0.395987in}{0.396913in}}%
\pgfpathlineto{\pgfqpoint{0.398135in}{0.399499in}}%
\pgfpathlineto{\pgfqpoint{0.401946in}{0.396269in}}%
\pgfpathlineto{\pgfqpoint{0.405142in}{0.396714in}}%
\pgfpathlineto{\pgfqpoint{0.410341in}{0.398951in}}%
\pgfpathlineto{\pgfqpoint{0.410761in}{0.401832in}}%
\pgfpathlineto{\pgfqpoint{0.417146in}{0.400406in}}%
\pgfpathlineto{\pgfqpoint{0.415651in}{0.396069in}}%
\pgfpathlineto{\pgfqpoint{0.419719in}{0.395751in}}%
\pgfpathlineto{\pgfqpoint{0.416306in}{0.390592in}}%
\pgfpathlineto{\pgfqpoint{0.413001in}{0.392264in}}%
\pgfpathlineto{\pgfqpoint{0.410323in}{0.391598in}}%
\pgfpathclose%
\pgfusepath{fill}%
\end{pgfscope}%
\begin{pgfscope}%
\pgfpathrectangle{\pgfqpoint{0.100000in}{0.100000in}}{\pgfqpoint{2.862985in}{1.832778in}}%
\pgfusepath{clip}%
\pgfsetbuttcap%
\pgfsetmiterjoin%
\definecolor{currentfill}{rgb}{0.645872,0.120898,0.598867}%
\pgfsetfillcolor{currentfill}%
\pgfsetlinewidth{0.000000pt}%
\definecolor{currentstroke}{rgb}{0.000000,0.000000,0.000000}%
\pgfsetstrokecolor{currentstroke}%
\pgfsetstrokeopacity{0.000000}%
\pgfsetdash{}{0pt}%
\pgfpathmoveto{\pgfqpoint{0.398962in}{0.402261in}}%
\pgfpathlineto{\pgfqpoint{0.396586in}{0.400969in}}%
\pgfpathlineto{\pgfqpoint{0.390110in}{0.402955in}}%
\pgfpathlineto{\pgfqpoint{0.385271in}{0.401805in}}%
\pgfpathlineto{\pgfqpoint{0.382813in}{0.405073in}}%
\pgfpathlineto{\pgfqpoint{0.383885in}{0.406569in}}%
\pgfpathlineto{\pgfqpoint{0.387552in}{0.406819in}}%
\pgfpathlineto{\pgfqpoint{0.390357in}{0.405880in}}%
\pgfpathlineto{\pgfqpoint{0.393414in}{0.407367in}}%
\pgfpathlineto{\pgfqpoint{0.398100in}{0.405378in}}%
\pgfpathlineto{\pgfqpoint{0.398962in}{0.402261in}}%
\pgfpathclose%
\pgfusepath{fill}%
\end{pgfscope}%
\begin{pgfscope}%
\pgfpathrectangle{\pgfqpoint{0.100000in}{0.100000in}}{\pgfqpoint{2.862985in}{1.832778in}}%
\pgfusepath{clip}%
\pgfsetbuttcap%
\pgfsetmiterjoin%
\definecolor{currentfill}{rgb}{0.645872,0.120898,0.598867}%
\pgfsetfillcolor{currentfill}%
\pgfsetlinewidth{0.000000pt}%
\definecolor{currentstroke}{rgb}{0.000000,0.000000,0.000000}%
\pgfsetstrokecolor{currentstroke}%
\pgfsetstrokeopacity{0.000000}%
\pgfsetdash{}{0pt}%
\pgfpathmoveto{\pgfqpoint{0.322628in}{0.450357in}}%
\pgfpathlineto{\pgfqpoint{0.322456in}{0.447153in}}%
\pgfpathlineto{\pgfqpoint{0.317729in}{0.445139in}}%
\pgfpathlineto{\pgfqpoint{0.313412in}{0.450180in}}%
\pgfpathlineto{\pgfqpoint{0.318253in}{0.451860in}}%
\pgfpathlineto{\pgfqpoint{0.322628in}{0.450357in}}%
\pgfpathclose%
\pgfusepath{fill}%
\end{pgfscope}%
\begin{pgfscope}%
\pgfpathrectangle{\pgfqpoint{0.100000in}{0.100000in}}{\pgfqpoint{2.862985in}{1.832778in}}%
\pgfusepath{clip}%
\pgfsetbuttcap%
\pgfsetmiterjoin%
\definecolor{currentfill}{rgb}{0.645872,0.120898,0.598867}%
\pgfsetfillcolor{currentfill}%
\pgfsetlinewidth{0.000000pt}%
\definecolor{currentstroke}{rgb}{0.000000,0.000000,0.000000}%
\pgfsetstrokecolor{currentstroke}%
\pgfsetstrokeopacity{0.000000}%
\pgfsetdash{}{0pt}%
\pgfpathmoveto{\pgfqpoint{0.283638in}{0.468371in}}%
\pgfpathlineto{\pgfqpoint{0.284642in}{0.469764in}}%
\pgfpathlineto{\pgfqpoint{0.290440in}{0.469524in}}%
\pgfpathlineto{\pgfqpoint{0.291458in}{0.471597in}}%
\pgfpathlineto{\pgfqpoint{0.294002in}{0.469788in}}%
\pgfpathlineto{\pgfqpoint{0.292533in}{0.466308in}}%
\pgfpathlineto{\pgfqpoint{0.290372in}{0.466014in}}%
\pgfpathlineto{\pgfqpoint{0.285207in}{0.467106in}}%
\pgfpathlineto{\pgfqpoint{0.283638in}{0.468371in}}%
\pgfpathclose%
\pgfusepath{fill}%
\end{pgfscope}%
\begin{pgfscope}%
\pgfpathrectangle{\pgfqpoint{0.100000in}{0.100000in}}{\pgfqpoint{2.862985in}{1.832778in}}%
\pgfusepath{clip}%
\pgfsetbuttcap%
\pgfsetmiterjoin%
\definecolor{currentfill}{rgb}{0.645872,0.120898,0.598867}%
\pgfsetfillcolor{currentfill}%
\pgfsetlinewidth{0.000000pt}%
\definecolor{currentstroke}{rgb}{0.000000,0.000000,0.000000}%
\pgfsetstrokecolor{currentstroke}%
\pgfsetstrokeopacity{0.000000}%
\pgfsetdash{}{0pt}%
\pgfpathmoveto{\pgfqpoint{0.275415in}{0.479156in}}%
\pgfpathlineto{\pgfqpoint{0.275063in}{0.482294in}}%
\pgfpathlineto{\pgfqpoint{0.278088in}{0.482062in}}%
\pgfpathlineto{\pgfqpoint{0.277878in}{0.486470in}}%
\pgfpathlineto{\pgfqpoint{0.280921in}{0.484091in}}%
\pgfpathlineto{\pgfqpoint{0.279868in}{0.480653in}}%
\pgfpathlineto{\pgfqpoint{0.275415in}{0.479156in}}%
\pgfpathclose%
\pgfusepath{fill}%
\end{pgfscope}%
\begin{pgfscope}%
\pgfpathrectangle{\pgfqpoint{0.100000in}{0.100000in}}{\pgfqpoint{2.862985in}{1.832778in}}%
\pgfusepath{clip}%
\pgfsetbuttcap%
\pgfsetmiterjoin%
\definecolor{currentfill}{rgb}{0.645872,0.120898,0.598867}%
\pgfsetfillcolor{currentfill}%
\pgfsetlinewidth{0.000000pt}%
\definecolor{currentstroke}{rgb}{0.000000,0.000000,0.000000}%
\pgfsetstrokecolor{currentstroke}%
\pgfsetstrokeopacity{0.000000}%
\pgfsetdash{}{0pt}%
\pgfpathmoveto{\pgfqpoint{0.541980in}{0.591016in}}%
\pgfpathlineto{\pgfqpoint{0.537154in}{0.592058in}}%
\pgfpathlineto{\pgfqpoint{0.536118in}{0.595283in}}%
\pgfpathlineto{\pgfqpoint{0.537776in}{0.598507in}}%
\pgfpathlineto{\pgfqpoint{0.541491in}{0.600190in}}%
\pgfpathlineto{\pgfqpoint{0.545744in}{0.591916in}}%
\pgfpathlineto{\pgfqpoint{0.549601in}{0.591548in}}%
\pgfpathlineto{\pgfqpoint{0.552479in}{0.588995in}}%
\pgfpathlineto{\pgfqpoint{0.552422in}{0.585482in}}%
\pgfpathlineto{\pgfqpoint{0.550318in}{0.583797in}}%
\pgfpathlineto{\pgfqpoint{0.553740in}{0.573408in}}%
\pgfpathlineto{\pgfqpoint{0.557484in}{0.569595in}}%
\pgfpathlineto{\pgfqpoint{0.553656in}{0.568385in}}%
\pgfpathlineto{\pgfqpoint{0.550557in}{0.572663in}}%
\pgfpathlineto{\pgfqpoint{0.543823in}{0.571338in}}%
\pgfpathlineto{\pgfqpoint{0.545918in}{0.575572in}}%
\pgfpathlineto{\pgfqpoint{0.544943in}{0.577258in}}%
\pgfpathlineto{\pgfqpoint{0.545514in}{0.581671in}}%
\pgfpathlineto{\pgfqpoint{0.544180in}{0.589648in}}%
\pgfpathlineto{\pgfqpoint{0.541980in}{0.591016in}}%
\pgfpathclose%
\pgfusepath{fill}%
\end{pgfscope}%
\begin{pgfscope}%
\pgfpathrectangle{\pgfqpoint{0.100000in}{0.100000in}}{\pgfqpoint{2.862985in}{1.832778in}}%
\pgfusepath{clip}%
\pgfsetbuttcap%
\pgfsetmiterjoin%
\definecolor{currentfill}{rgb}{0.645872,0.120898,0.598867}%
\pgfsetfillcolor{currentfill}%
\pgfsetlinewidth{0.000000pt}%
\definecolor{currentstroke}{rgb}{0.000000,0.000000,0.000000}%
\pgfsetstrokecolor{currentstroke}%
\pgfsetstrokeopacity{0.000000}%
\pgfsetdash{}{0pt}%
\pgfpathmoveto{\pgfqpoint{0.711519in}{0.353334in}}%
\pgfpathlineto{\pgfqpoint{0.704523in}{0.353412in}}%
\pgfpathlineto{\pgfqpoint{0.705062in}{0.356924in}}%
\pgfpathlineto{\pgfqpoint{0.707862in}{0.358180in}}%
\pgfpathlineto{\pgfqpoint{0.711519in}{0.353334in}}%
\pgfpathclose%
\pgfusepath{fill}%
\end{pgfscope}%
\begin{pgfscope}%
\pgfpathrectangle{\pgfqpoint{0.100000in}{0.100000in}}{\pgfqpoint{2.862985in}{1.832778in}}%
\pgfusepath{clip}%
\pgfsetbuttcap%
\pgfsetmiterjoin%
\definecolor{currentfill}{rgb}{0.645872,0.120898,0.598867}%
\pgfsetfillcolor{currentfill}%
\pgfsetlinewidth{0.000000pt}%
\definecolor{currentstroke}{rgb}{0.000000,0.000000,0.000000}%
\pgfsetstrokecolor{currentstroke}%
\pgfsetstrokeopacity{0.000000}%
\pgfsetdash{}{0pt}%
\pgfpathmoveto{\pgfqpoint{0.701812in}{0.356514in}}%
\pgfpathlineto{\pgfqpoint{0.696800in}{0.355301in}}%
\pgfpathlineto{\pgfqpoint{0.692038in}{0.353224in}}%
\pgfpathlineto{\pgfqpoint{0.690567in}{0.355671in}}%
\pgfpathlineto{\pgfqpoint{0.698757in}{0.357328in}}%
\pgfpathlineto{\pgfqpoint{0.701812in}{0.356514in}}%
\pgfpathclose%
\pgfusepath{fill}%
\end{pgfscope}%
\begin{pgfscope}%
\pgfpathrectangle{\pgfqpoint{0.100000in}{0.100000in}}{\pgfqpoint{2.862985in}{1.832778in}}%
\pgfusepath{clip}%
\pgfsetbuttcap%
\pgfsetmiterjoin%
\definecolor{currentfill}{rgb}{0.645872,0.120898,0.598867}%
\pgfsetfillcolor{currentfill}%
\pgfsetlinewidth{0.000000pt}%
\definecolor{currentstroke}{rgb}{0.000000,0.000000,0.000000}%
\pgfsetstrokecolor{currentstroke}%
\pgfsetstrokeopacity{0.000000}%
\pgfsetdash{}{0pt}%
\pgfpathmoveto{\pgfqpoint{0.624990in}{0.354028in}}%
\pgfpathlineto{\pgfqpoint{0.625182in}{0.351599in}}%
\pgfpathlineto{\pgfqpoint{0.622294in}{0.350090in}}%
\pgfpathlineto{\pgfqpoint{0.614235in}{0.353552in}}%
\pgfpathlineto{\pgfqpoint{0.612986in}{0.352159in}}%
\pgfpathlineto{\pgfqpoint{0.611036in}{0.358747in}}%
\pgfpathlineto{\pgfqpoint{0.615074in}{0.359664in}}%
\pgfpathlineto{\pgfqpoint{0.616940in}{0.356957in}}%
\pgfpathlineto{\pgfqpoint{0.620992in}{0.360345in}}%
\pgfpathlineto{\pgfqpoint{0.624990in}{0.354028in}}%
\pgfpathclose%
\pgfusepath{fill}%
\end{pgfscope}%
\begin{pgfscope}%
\pgfpathrectangle{\pgfqpoint{0.100000in}{0.100000in}}{\pgfqpoint{2.862985in}{1.832778in}}%
\pgfusepath{clip}%
\pgfsetbuttcap%
\pgfsetmiterjoin%
\definecolor{currentfill}{rgb}{0.645872,0.120898,0.598867}%
\pgfsetfillcolor{currentfill}%
\pgfsetlinewidth{0.000000pt}%
\definecolor{currentstroke}{rgb}{0.000000,0.000000,0.000000}%
\pgfsetstrokecolor{currentstroke}%
\pgfsetstrokeopacity{0.000000}%
\pgfsetdash{}{0pt}%
\pgfpathmoveto{\pgfqpoint{0.806774in}{0.219412in}}%
\pgfpathlineto{\pgfqpoint{0.800017in}{0.216896in}}%
\pgfpathlineto{\pgfqpoint{0.796437in}{0.216745in}}%
\pgfpathlineto{\pgfqpoint{0.799350in}{0.222082in}}%
\pgfpathlineto{\pgfqpoint{0.801840in}{0.224073in}}%
\pgfpathlineto{\pgfqpoint{0.801760in}{0.229600in}}%
\pgfpathlineto{\pgfqpoint{0.804465in}{0.239578in}}%
\pgfpathlineto{\pgfqpoint{0.806957in}{0.241473in}}%
\pgfpathlineto{\pgfqpoint{0.812710in}{0.238729in}}%
\pgfpathlineto{\pgfqpoint{0.811870in}{0.237158in}}%
\pgfpathlineto{\pgfqpoint{0.812306in}{0.229604in}}%
\pgfpathlineto{\pgfqpoint{0.810440in}{0.227566in}}%
\pgfpathlineto{\pgfqpoint{0.809640in}{0.222201in}}%
\pgfpathlineto{\pgfqpoint{0.806774in}{0.219412in}}%
\pgfpathclose%
\pgfusepath{fill}%
\end{pgfscope}%
\begin{pgfscope}%
\pgfpathrectangle{\pgfqpoint{0.100000in}{0.100000in}}{\pgfqpoint{2.862985in}{1.832778in}}%
\pgfusepath{clip}%
\pgfsetbuttcap%
\pgfsetmiterjoin%
\definecolor{currentfill}{rgb}{0.645872,0.120898,0.598867}%
\pgfsetfillcolor{currentfill}%
\pgfsetlinewidth{0.000000pt}%
\definecolor{currentstroke}{rgb}{0.000000,0.000000,0.000000}%
\pgfsetstrokecolor{currentstroke}%
\pgfsetstrokeopacity{0.000000}%
\pgfsetdash{}{0pt}%
\pgfpathmoveto{\pgfqpoint{0.791809in}{0.244146in}}%
\pgfpathlineto{\pgfqpoint{0.795768in}{0.237310in}}%
\pgfpathlineto{\pgfqpoint{0.795034in}{0.235680in}}%
\pgfpathlineto{\pgfqpoint{0.800329in}{0.234022in}}%
\pgfpathlineto{\pgfqpoint{0.798809in}{0.227738in}}%
\pgfpathlineto{\pgfqpoint{0.796503in}{0.228029in}}%
\pgfpathlineto{\pgfqpoint{0.790445in}{0.238982in}}%
\pgfpathlineto{\pgfqpoint{0.791307in}{0.234052in}}%
\pgfpathlineto{\pgfqpoint{0.790275in}{0.231213in}}%
\pgfpathlineto{\pgfqpoint{0.787709in}{0.229950in}}%
\pgfpathlineto{\pgfqpoint{0.786251in}{0.231311in}}%
\pgfpathlineto{\pgfqpoint{0.785437in}{0.237623in}}%
\pgfpathlineto{\pgfqpoint{0.785917in}{0.241717in}}%
\pgfpathlineto{\pgfqpoint{0.784208in}{0.243504in}}%
\pgfpathlineto{\pgfqpoint{0.788722in}{0.248675in}}%
\pgfpathlineto{\pgfqpoint{0.791684in}{0.250209in}}%
\pgfpathlineto{\pgfqpoint{0.792670in}{0.248193in}}%
\pgfpathlineto{\pgfqpoint{0.795760in}{0.249753in}}%
\pgfpathlineto{\pgfqpoint{0.797999in}{0.245516in}}%
\pgfpathlineto{\pgfqpoint{0.802840in}{0.240126in}}%
\pgfpathlineto{\pgfqpoint{0.802179in}{0.237684in}}%
\pgfpathlineto{\pgfqpoint{0.799491in}{0.235024in}}%
\pgfpathlineto{\pgfqpoint{0.797005in}{0.236278in}}%
\pgfpathlineto{\pgfqpoint{0.791809in}{0.244146in}}%
\pgfpathclose%
\pgfusepath{fill}%
\end{pgfscope}%
\begin{pgfscope}%
\pgfpathrectangle{\pgfqpoint{0.100000in}{0.100000in}}{\pgfqpoint{2.862985in}{1.832778in}}%
\pgfusepath{clip}%
\pgfsetbuttcap%
\pgfsetmiterjoin%
\definecolor{currentfill}{rgb}{0.645872,0.120898,0.598867}%
\pgfsetfillcolor{currentfill}%
\pgfsetlinewidth{0.000000pt}%
\definecolor{currentstroke}{rgb}{0.000000,0.000000,0.000000}%
\pgfsetstrokecolor{currentstroke}%
\pgfsetstrokeopacity{0.000000}%
\pgfsetdash{}{0pt}%
\pgfpathmoveto{\pgfqpoint{0.595353in}{0.341071in}}%
\pgfpathlineto{\pgfqpoint{0.590449in}{0.340433in}}%
\pgfpathlineto{\pgfqpoint{0.583542in}{0.337778in}}%
\pgfpathlineto{\pgfqpoint{0.581931in}{0.339221in}}%
\pgfpathlineto{\pgfqpoint{0.589146in}{0.340851in}}%
\pgfpathlineto{\pgfqpoint{0.587212in}{0.342273in}}%
\pgfpathlineto{\pgfqpoint{0.582739in}{0.343315in}}%
\pgfpathlineto{\pgfqpoint{0.581544in}{0.346649in}}%
\pgfpathlineto{\pgfqpoint{0.583795in}{0.349521in}}%
\pgfpathlineto{\pgfqpoint{0.585690in}{0.356112in}}%
\pgfpathlineto{\pgfqpoint{0.592614in}{0.357340in}}%
\pgfpathlineto{\pgfqpoint{0.595966in}{0.354492in}}%
\pgfpathlineto{\pgfqpoint{0.599115in}{0.357248in}}%
\pgfpathlineto{\pgfqpoint{0.603064in}{0.356572in}}%
\pgfpathlineto{\pgfqpoint{0.605902in}{0.351482in}}%
\pgfpathlineto{\pgfqpoint{0.608305in}{0.356006in}}%
\pgfpathlineto{\pgfqpoint{0.615126in}{0.345423in}}%
\pgfpathlineto{\pgfqpoint{0.612662in}{0.344577in}}%
\pgfpathlineto{\pgfqpoint{0.611150in}{0.342283in}}%
\pgfpathlineto{\pgfqpoint{0.614007in}{0.340190in}}%
\pgfpathlineto{\pgfqpoint{0.609576in}{0.338227in}}%
\pgfpathlineto{\pgfqpoint{0.606364in}{0.339308in}}%
\pgfpathlineto{\pgfqpoint{0.600296in}{0.342890in}}%
\pgfpathlineto{\pgfqpoint{0.595353in}{0.341071in}}%
\pgfpathclose%
\pgfusepath{fill}%
\end{pgfscope}%
\begin{pgfscope}%
\pgfpathrectangle{\pgfqpoint{0.100000in}{0.100000in}}{\pgfqpoint{2.862985in}{1.832778in}}%
\pgfusepath{clip}%
\pgfsetbuttcap%
\pgfsetmiterjoin%
\definecolor{currentfill}{rgb}{0.645872,0.120898,0.598867}%
\pgfsetfillcolor{currentfill}%
\pgfsetlinewidth{0.000000pt}%
\definecolor{currentstroke}{rgb}{0.000000,0.000000,0.000000}%
\pgfsetstrokecolor{currentstroke}%
\pgfsetstrokeopacity{0.000000}%
\pgfsetdash{}{0pt}%
\pgfpathmoveto{\pgfqpoint{0.787691in}{0.199220in}}%
\pgfpathlineto{\pgfqpoint{0.786719in}{0.208994in}}%
\pgfpathlineto{\pgfqpoint{0.789340in}{0.212583in}}%
\pgfpathlineto{\pgfqpoint{0.786538in}{0.212832in}}%
\pgfpathlineto{\pgfqpoint{0.787562in}{0.219812in}}%
\pgfpathlineto{\pgfqpoint{0.788770in}{0.223955in}}%
\pgfpathlineto{\pgfqpoint{0.787720in}{0.229522in}}%
\pgfpathlineto{\pgfqpoint{0.790287in}{0.230600in}}%
\pgfpathlineto{\pgfqpoint{0.791124in}{0.232039in}}%
\pgfpathlineto{\pgfqpoint{0.793653in}{0.231493in}}%
\pgfpathlineto{\pgfqpoint{0.795842in}{0.227033in}}%
\pgfpathlineto{\pgfqpoint{0.796249in}{0.222947in}}%
\pgfpathlineto{\pgfqpoint{0.793571in}{0.210709in}}%
\pgfpathlineto{\pgfqpoint{0.787691in}{0.199220in}}%
\pgfpathclose%
\pgfusepath{fill}%
\end{pgfscope}%
\begin{pgfscope}%
\pgfpathrectangle{\pgfqpoint{0.100000in}{0.100000in}}{\pgfqpoint{2.862985in}{1.832778in}}%
\pgfusepath{clip}%
\pgfsetbuttcap%
\pgfsetmiterjoin%
\definecolor{currentfill}{rgb}{0.645872,0.120898,0.598867}%
\pgfsetfillcolor{currentfill}%
\pgfsetlinewidth{0.000000pt}%
\definecolor{currentstroke}{rgb}{0.000000,0.000000,0.000000}%
\pgfsetstrokecolor{currentstroke}%
\pgfsetstrokeopacity{0.000000}%
\pgfsetdash{}{0pt}%
\pgfpathmoveto{\pgfqpoint{0.787110in}{0.229056in}}%
\pgfpathlineto{\pgfqpoint{0.786975in}{0.224290in}}%
\pgfpathlineto{\pgfqpoint{0.784778in}{0.222127in}}%
\pgfpathlineto{\pgfqpoint{0.782238in}{0.222926in}}%
\pgfpathlineto{\pgfqpoint{0.785448in}{0.229848in}}%
\pgfpathlineto{\pgfqpoint{0.787110in}{0.229056in}}%
\pgfpathclose%
\pgfusepath{fill}%
\end{pgfscope}%
\begin{pgfscope}%
\pgfpathrectangle{\pgfqpoint{0.100000in}{0.100000in}}{\pgfqpoint{2.862985in}{1.832778in}}%
\pgfusepath{clip}%
\pgfsetbuttcap%
\pgfsetmiterjoin%
\definecolor{currentfill}{rgb}{0.645872,0.120898,0.598867}%
\pgfsetfillcolor{currentfill}%
\pgfsetlinewidth{0.000000pt}%
\definecolor{currentstroke}{rgb}{0.000000,0.000000,0.000000}%
\pgfsetstrokecolor{currentstroke}%
\pgfsetstrokeopacity{0.000000}%
\pgfsetdash{}{0pt}%
\pgfpathmoveto{\pgfqpoint{0.814076in}{0.207898in}}%
\pgfpathlineto{\pgfqpoint{0.812202in}{0.199923in}}%
\pgfpathlineto{\pgfqpoint{0.807835in}{0.197402in}}%
\pgfpathlineto{\pgfqpoint{0.802107in}{0.199724in}}%
\pgfpathlineto{\pgfqpoint{0.803685in}{0.202824in}}%
\pgfpathlineto{\pgfqpoint{0.803653in}{0.204970in}}%
\pgfpathlineto{\pgfqpoint{0.805783in}{0.208616in}}%
\pgfpathlineto{\pgfqpoint{0.804286in}{0.207882in}}%
\pgfpathlineto{\pgfqpoint{0.805377in}{0.209663in}}%
\pgfpathlineto{\pgfqpoint{0.803304in}{0.213709in}}%
\pgfpathlineto{\pgfqpoint{0.805657in}{0.214455in}}%
\pgfpathlineto{\pgfqpoint{0.811104in}{0.209721in}}%
\pgfpathlineto{\pgfqpoint{0.814076in}{0.207898in}}%
\pgfpathclose%
\pgfusepath{fill}%
\end{pgfscope}%
\begin{pgfscope}%
\pgfpathrectangle{\pgfqpoint{0.100000in}{0.100000in}}{\pgfqpoint{2.862985in}{1.832778in}}%
\pgfusepath{clip}%
\pgfsetbuttcap%
\pgfsetmiterjoin%
\definecolor{currentfill}{rgb}{0.645872,0.120898,0.598867}%
\pgfsetfillcolor{currentfill}%
\pgfsetlinewidth{0.000000pt}%
\definecolor{currentstroke}{rgb}{0.000000,0.000000,0.000000}%
\pgfsetstrokecolor{currentstroke}%
\pgfsetstrokeopacity{0.000000}%
\pgfsetdash{}{0pt}%
\pgfpathmoveto{\pgfqpoint{0.794962in}{0.193796in}}%
\pgfpathlineto{\pgfqpoint{0.792492in}{0.192836in}}%
\pgfpathlineto{\pgfqpoint{0.793902in}{0.201252in}}%
\pgfpathlineto{\pgfqpoint{0.796866in}{0.203216in}}%
\pgfpathlineto{\pgfqpoint{0.796144in}{0.209521in}}%
\pgfpathlineto{\pgfqpoint{0.797348in}{0.212317in}}%
\pgfpathlineto{\pgfqpoint{0.799731in}{0.213365in}}%
\pgfpathlineto{\pgfqpoint{0.802232in}{0.210622in}}%
\pgfpathlineto{\pgfqpoint{0.804379in}{0.206901in}}%
\pgfpathlineto{\pgfqpoint{0.797478in}{0.199610in}}%
\pgfpathlineto{\pgfqpoint{0.794962in}{0.193796in}}%
\pgfpathclose%
\pgfusepath{fill}%
\end{pgfscope}%
\begin{pgfscope}%
\pgfpathrectangle{\pgfqpoint{0.100000in}{0.100000in}}{\pgfqpoint{2.862985in}{1.832778in}}%
\pgfusepath{clip}%
\pgfsetbuttcap%
\pgfsetmiterjoin%
\definecolor{currentfill}{rgb}{0.645872,0.120898,0.598867}%
\pgfsetfillcolor{currentfill}%
\pgfsetlinewidth{0.000000pt}%
\definecolor{currentstroke}{rgb}{0.000000,0.000000,0.000000}%
\pgfsetstrokecolor{currentstroke}%
\pgfsetstrokeopacity{0.000000}%
\pgfsetdash{}{0pt}%
\pgfpathmoveto{\pgfqpoint{0.815172in}{0.202568in}}%
\pgfpathlineto{\pgfqpoint{0.816382in}{0.196760in}}%
\pgfpathlineto{\pgfqpoint{0.810974in}{0.197425in}}%
\pgfpathlineto{\pgfqpoint{0.815172in}{0.202568in}}%
\pgfpathclose%
\pgfusepath{fill}%
\end{pgfscope}%
\begin{pgfscope}%
\pgfpathrectangle{\pgfqpoint{0.100000in}{0.100000in}}{\pgfqpoint{2.862985in}{1.832778in}}%
\pgfusepath{clip}%
\pgfsetbuttcap%
\pgfsetmiterjoin%
\definecolor{currentfill}{rgb}{0.645872,0.120898,0.598867}%
\pgfsetfillcolor{currentfill}%
\pgfsetlinewidth{0.000000pt}%
\definecolor{currentstroke}{rgb}{0.000000,0.000000,0.000000}%
\pgfsetstrokecolor{currentstroke}%
\pgfsetstrokeopacity{0.000000}%
\pgfsetdash{}{0pt}%
\pgfpathmoveto{\pgfqpoint{0.817752in}{0.193694in}}%
\pgfpathlineto{\pgfqpoint{0.818638in}{0.190111in}}%
\pgfpathlineto{\pgfqpoint{0.820335in}{0.188965in}}%
\pgfpathlineto{\pgfqpoint{0.820026in}{0.185593in}}%
\pgfpathlineto{\pgfqpoint{0.817656in}{0.184459in}}%
\pgfpathlineto{\pgfqpoint{0.815732in}{0.189493in}}%
\pgfpathlineto{\pgfqpoint{0.817752in}{0.193694in}}%
\pgfpathclose%
\pgfusepath{fill}%
\end{pgfscope}%
\begin{pgfscope}%
\pgfpathrectangle{\pgfqpoint{0.100000in}{0.100000in}}{\pgfqpoint{2.862985in}{1.832778in}}%
\pgfusepath{clip}%
\pgfsetbuttcap%
\pgfsetmiterjoin%
\definecolor{currentfill}{rgb}{0.645872,0.120898,0.598867}%
\pgfsetfillcolor{currentfill}%
\pgfsetlinewidth{0.000000pt}%
\definecolor{currentstroke}{rgb}{0.000000,0.000000,0.000000}%
\pgfsetstrokecolor{currentstroke}%
\pgfsetstrokeopacity{0.000000}%
\pgfsetdash{}{0pt}%
\pgfpathmoveto{\pgfqpoint{0.813239in}{0.195068in}}%
\pgfpathlineto{\pgfqpoint{0.811768in}{0.190907in}}%
\pgfpathlineto{\pgfqpoint{0.809484in}{0.191060in}}%
\pgfpathlineto{\pgfqpoint{0.808043in}{0.194492in}}%
\pgfpathlineto{\pgfqpoint{0.810390in}{0.196117in}}%
\pgfpathlineto{\pgfqpoint{0.813239in}{0.195068in}}%
\pgfpathclose%
\pgfusepath{fill}%
\end{pgfscope}%
\begin{pgfscope}%
\pgfpathrectangle{\pgfqpoint{0.100000in}{0.100000in}}{\pgfqpoint{2.862985in}{1.832778in}}%
\pgfusepath{clip}%
\pgfsetbuttcap%
\pgfsetmiterjoin%
\definecolor{currentfill}{rgb}{0.645872,0.120898,0.598867}%
\pgfsetfillcolor{currentfill}%
\pgfsetlinewidth{0.000000pt}%
\definecolor{currentstroke}{rgb}{0.000000,0.000000,0.000000}%
\pgfsetstrokecolor{currentstroke}%
\pgfsetstrokeopacity{0.000000}%
\pgfsetdash{}{0pt}%
\pgfpathmoveto{\pgfqpoint{0.807076in}{0.173613in}}%
\pgfpathlineto{\pgfqpoint{0.809114in}{0.171384in}}%
\pgfpathlineto{\pgfqpoint{0.809563in}{0.166686in}}%
\pgfpathlineto{\pgfqpoint{0.810636in}{0.165321in}}%
\pgfpathlineto{\pgfqpoint{0.808884in}{0.161070in}}%
\pgfpathlineto{\pgfqpoint{0.806578in}{0.158676in}}%
\pgfpathlineto{\pgfqpoint{0.805352in}{0.154236in}}%
\pgfpathlineto{\pgfqpoint{0.802832in}{0.161156in}}%
\pgfpathlineto{\pgfqpoint{0.802414in}{0.167439in}}%
\pgfpathlineto{\pgfqpoint{0.801004in}{0.170567in}}%
\pgfpathlineto{\pgfqpoint{0.796147in}{0.173102in}}%
\pgfpathlineto{\pgfqpoint{0.800703in}{0.178390in}}%
\pgfpathlineto{\pgfqpoint{0.797666in}{0.180650in}}%
\pgfpathlineto{\pgfqpoint{0.800341in}{0.182902in}}%
\pgfpathlineto{\pgfqpoint{0.803893in}{0.191549in}}%
\pgfpathlineto{\pgfqpoint{0.800212in}{0.194597in}}%
\pgfpathlineto{\pgfqpoint{0.801630in}{0.197524in}}%
\pgfpathlineto{\pgfqpoint{0.806388in}{0.194928in}}%
\pgfpathlineto{\pgfqpoint{0.806870in}{0.190831in}}%
\pgfpathlineto{\pgfqpoint{0.805070in}{0.188498in}}%
\pgfpathlineto{\pgfqpoint{0.808638in}{0.185341in}}%
\pgfpathlineto{\pgfqpoint{0.809761in}{0.180424in}}%
\pgfpathlineto{\pgfqpoint{0.809609in}{0.173528in}}%
\pgfpathlineto{\pgfqpoint{0.807076in}{0.173613in}}%
\pgfpathclose%
\pgfusepath{fill}%
\end{pgfscope}%
\begin{pgfscope}%
\pgfpathrectangle{\pgfqpoint{0.100000in}{0.100000in}}{\pgfqpoint{2.862985in}{1.832778in}}%
\pgfusepath{clip}%
\pgfsetbuttcap%
\pgfsetmiterjoin%
\definecolor{currentfill}{rgb}{0.645872,0.120898,0.598867}%
\pgfsetfillcolor{currentfill}%
\pgfsetlinewidth{0.000000pt}%
\definecolor{currentstroke}{rgb}{0.000000,0.000000,0.000000}%
\pgfsetstrokecolor{currentstroke}%
\pgfsetstrokeopacity{0.000000}%
\pgfsetdash{}{0pt}%
\pgfpathmoveto{\pgfqpoint{0.815789in}{0.190924in}}%
\pgfpathlineto{\pgfqpoint{0.814679in}{0.188532in}}%
\pgfpathlineto{\pgfqpoint{0.815722in}{0.184648in}}%
\pgfpathlineto{\pgfqpoint{0.813053in}{0.184321in}}%
\pgfpathlineto{\pgfqpoint{0.810255in}{0.186119in}}%
\pgfpathlineto{\pgfqpoint{0.811163in}{0.190080in}}%
\pgfpathlineto{\pgfqpoint{0.815789in}{0.190924in}}%
\pgfpathclose%
\pgfusepath{fill}%
\end{pgfscope}%
\begin{pgfscope}%
\pgfpathrectangle{\pgfqpoint{0.100000in}{0.100000in}}{\pgfqpoint{2.862985in}{1.832778in}}%
\pgfusepath{clip}%
\pgfsetbuttcap%
\pgfsetmiterjoin%
\definecolor{currentfill}{rgb}{0.645872,0.120898,0.598867}%
\pgfsetfillcolor{currentfill}%
\pgfsetlinewidth{0.000000pt}%
\definecolor{currentstroke}{rgb}{0.000000,0.000000,0.000000}%
\pgfsetstrokecolor{currentstroke}%
\pgfsetstrokeopacity{0.000000}%
\pgfsetdash{}{0pt}%
\pgfpathmoveto{\pgfqpoint{0.803351in}{0.191195in}}%
\pgfpathlineto{\pgfqpoint{0.797461in}{0.188553in}}%
\pgfpathlineto{\pgfqpoint{0.795459in}{0.189503in}}%
\pgfpathlineto{\pgfqpoint{0.799476in}{0.192989in}}%
\pgfpathlineto{\pgfqpoint{0.803351in}{0.191195in}}%
\pgfpathclose%
\pgfusepath{fill}%
\end{pgfscope}%
\begin{pgfscope}%
\pgfpathrectangle{\pgfqpoint{0.100000in}{0.100000in}}{\pgfqpoint{2.862985in}{1.832778in}}%
\pgfusepath{clip}%
\pgfsetbuttcap%
\pgfsetmiterjoin%
\definecolor{currentfill}{rgb}{0.645872,0.120898,0.598867}%
\pgfsetfillcolor{currentfill}%
\pgfsetlinewidth{0.000000pt}%
\definecolor{currentstroke}{rgb}{0.000000,0.000000,0.000000}%
\pgfsetstrokecolor{currentstroke}%
\pgfsetstrokeopacity{0.000000}%
\pgfsetdash{}{0pt}%
\pgfpathmoveto{\pgfqpoint{0.815854in}{0.165422in}}%
\pgfpathlineto{\pgfqpoint{0.814449in}{0.169400in}}%
\pgfpathlineto{\pgfqpoint{0.817419in}{0.170378in}}%
\pgfpathlineto{\pgfqpoint{0.818319in}{0.174668in}}%
\pgfpathlineto{\pgfqpoint{0.821538in}{0.174673in}}%
\pgfpathlineto{\pgfqpoint{0.821524in}{0.177660in}}%
\pgfpathlineto{\pgfqpoint{0.825997in}{0.176953in}}%
\pgfpathlineto{\pgfqpoint{0.826760in}{0.168856in}}%
\pgfpathlineto{\pgfqpoint{0.824165in}{0.163694in}}%
\pgfpathlineto{\pgfqpoint{0.820315in}{0.160469in}}%
\pgfpathlineto{\pgfqpoint{0.815854in}{0.165422in}}%
\pgfpathclose%
\pgfusepath{fill}%
\end{pgfscope}%
\begin{pgfscope}%
\pgfpathrectangle{\pgfqpoint{0.100000in}{0.100000in}}{\pgfqpoint{2.862985in}{1.832778in}}%
\pgfusepath{clip}%
\pgfsetbuttcap%
\pgfsetmiterjoin%
\definecolor{currentfill}{rgb}{0.645872,0.120898,0.598867}%
\pgfsetfillcolor{currentfill}%
\pgfsetlinewidth{0.000000pt}%
\definecolor{currentstroke}{rgb}{0.000000,0.000000,0.000000}%
\pgfsetstrokecolor{currentstroke}%
\pgfsetstrokeopacity{0.000000}%
\pgfsetdash{}{0pt}%
\pgfpathmoveto{\pgfqpoint{0.814114in}{0.168482in}}%
\pgfpathlineto{\pgfqpoint{0.815383in}{0.164715in}}%
\pgfpathlineto{\pgfqpoint{0.812255in}{0.164257in}}%
\pgfpathlineto{\pgfqpoint{0.814114in}{0.168482in}}%
\pgfpathclose%
\pgfusepath{fill}%
\end{pgfscope}%
\begin{pgfscope}%
\pgfpathrectangle{\pgfqpoint{0.100000in}{0.100000in}}{\pgfqpoint{2.862985in}{1.832778in}}%
\pgfusepath{clip}%
\pgfsetbuttcap%
\pgfsetmiterjoin%
\definecolor{currentfill}{rgb}{0.645872,0.120898,0.598867}%
\pgfsetfillcolor{currentfill}%
\pgfsetlinewidth{0.000000pt}%
\definecolor{currentstroke}{rgb}{0.000000,0.000000,0.000000}%
\pgfsetstrokecolor{currentstroke}%
\pgfsetstrokeopacity{0.000000}%
\pgfsetdash{}{0pt}%
\pgfpathmoveto{\pgfqpoint{0.816751in}{0.162942in}}%
\pgfpathlineto{\pgfqpoint{0.816315in}{0.158141in}}%
\pgfpathlineto{\pgfqpoint{0.813015in}{0.158538in}}%
\pgfpathlineto{\pgfqpoint{0.814083in}{0.160858in}}%
\pgfpathlineto{\pgfqpoint{0.816751in}{0.162942in}}%
\pgfpathclose%
\pgfusepath{fill}%
\end{pgfscope}%
\begin{pgfscope}%
\definecolor{textcolor}{rgb}{0.000000,0.000000,0.000000}%
\pgfsetstrokecolor{textcolor}%
\pgfsetfillcolor{textcolor}%
\pgftext[x=1.302454in,y=1.896122in,left,base]{\color{textcolor}\setmainfont{Lato}\rmfamily\fontsize{9.000000}{10.800000}\selectfont November 2023}%
\end{pgfscope}%
\begin{pgfscope}%
\pgfpathrectangle{\pgfqpoint{3.525000in}{0.100000in}}{\pgfqpoint{2.862985in}{1.832778in}}%
\pgfusepath{clip}%
\pgfsetbuttcap%
\pgfsetmiterjoin%
\definecolor{currentfill}{rgb}{0.783383,0.262500,0.483918}%
\pgfsetfillcolor{currentfill}%
\pgfsetlinewidth{0.000000pt}%
\definecolor{currentstroke}{rgb}{0.000000,0.000000,0.000000}%
\pgfsetstrokecolor{currentstroke}%
\pgfsetstrokeopacity{0.000000}%
\pgfsetdash{}{0pt}%
\pgfpathmoveto{\pgfqpoint{4.480609in}{0.529755in}}%
\pgfpathlineto{\pgfqpoint{4.472380in}{0.529889in}}%
\pgfpathlineto{\pgfqpoint{4.466780in}{0.536088in}}%
\pgfpathlineto{\pgfqpoint{4.462972in}{0.552550in}}%
\pgfpathlineto{\pgfqpoint{4.471822in}{0.554612in}}%
\pgfpathlineto{\pgfqpoint{4.480813in}{0.552394in}}%
\pgfpathlineto{\pgfqpoint{4.489963in}{0.545909in}}%
\pgfpathlineto{\pgfqpoint{4.489959in}{0.536940in}}%
\pgfpathlineto{\pgfqpoint{4.480609in}{0.529755in}}%
\pgfpathclose%
\pgfusepath{fill}%
\end{pgfscope}%
\begin{pgfscope}%
\pgfpathrectangle{\pgfqpoint{3.525000in}{0.100000in}}{\pgfqpoint{2.862985in}{1.832778in}}%
\pgfusepath{clip}%
\pgfsetbuttcap%
\pgfsetmiterjoin%
\definecolor{currentfill}{rgb}{0.783383,0.262500,0.483918}%
\pgfsetfillcolor{currentfill}%
\pgfsetlinewidth{0.000000pt}%
\definecolor{currentstroke}{rgb}{0.000000,0.000000,0.000000}%
\pgfsetstrokecolor{currentstroke}%
\pgfsetstrokeopacity{0.000000}%
\pgfsetdash{}{0pt}%
\pgfpathmoveto{\pgfqpoint{4.530694in}{0.432023in}}%
\pgfpathlineto{\pgfqpoint{4.521488in}{0.435528in}}%
\pgfpathlineto{\pgfqpoint{4.521446in}{0.442404in}}%
\pgfpathlineto{\pgfqpoint{4.510326in}{0.448219in}}%
\pgfpathlineto{\pgfqpoint{4.511680in}{0.462860in}}%
\pgfpathlineto{\pgfqpoint{4.514901in}{0.469769in}}%
\pgfpathlineto{\pgfqpoint{4.521701in}{0.463410in}}%
\pgfpathlineto{\pgfqpoint{4.532760in}{0.464472in}}%
\pgfpathlineto{\pgfqpoint{4.529967in}{0.444236in}}%
\pgfpathlineto{\pgfqpoint{4.530694in}{0.432023in}}%
\pgfpathclose%
\pgfusepath{fill}%
\end{pgfscope}%
\begin{pgfscope}%
\pgfpathrectangle{\pgfqpoint{3.525000in}{0.100000in}}{\pgfqpoint{2.862985in}{1.832778in}}%
\pgfusepath{clip}%
\pgfsetbuttcap%
\pgfsetmiterjoin%
\definecolor{currentfill}{rgb}{0.783383,0.262500,0.483918}%
\pgfsetfillcolor{currentfill}%
\pgfsetlinewidth{0.000000pt}%
\definecolor{currentstroke}{rgb}{0.000000,0.000000,0.000000}%
\pgfsetstrokecolor{currentstroke}%
\pgfsetstrokeopacity{0.000000}%
\pgfsetdash{}{0pt}%
\pgfpathmoveto{\pgfqpoint{4.569333in}{0.386887in}}%
\pgfpathlineto{\pgfqpoint{4.557651in}{0.388876in}}%
\pgfpathlineto{\pgfqpoint{4.551191in}{0.398877in}}%
\pgfpathlineto{\pgfqpoint{4.550279in}{0.407100in}}%
\pgfpathlineto{\pgfqpoint{4.569333in}{0.386887in}}%
\pgfpathclose%
\pgfusepath{fill}%
\end{pgfscope}%
\begin{pgfscope}%
\pgfpathrectangle{\pgfqpoint{3.525000in}{0.100000in}}{\pgfqpoint{2.862985in}{1.832778in}}%
\pgfusepath{clip}%
\pgfsetbuttcap%
\pgfsetmiterjoin%
\definecolor{currentfill}{rgb}{0.783383,0.262500,0.483918}%
\pgfsetfillcolor{currentfill}%
\pgfsetlinewidth{0.000000pt}%
\definecolor{currentstroke}{rgb}{0.000000,0.000000,0.000000}%
\pgfsetstrokecolor{currentstroke}%
\pgfsetstrokeopacity{0.000000}%
\pgfsetdash{}{0pt}%
\pgfpathmoveto{\pgfqpoint{4.545394in}{0.388631in}}%
\pgfpathlineto{\pgfqpoint{4.551519in}{0.383659in}}%
\pgfpathlineto{\pgfqpoint{4.552848in}{0.375110in}}%
\pgfpathlineto{\pgfqpoint{4.540519in}{0.375583in}}%
\pgfpathlineto{\pgfqpoint{4.545394in}{0.388631in}}%
\pgfpathclose%
\pgfusepath{fill}%
\end{pgfscope}%
\begin{pgfscope}%
\pgfpathrectangle{\pgfqpoint{3.525000in}{0.100000in}}{\pgfqpoint{2.862985in}{1.832778in}}%
\pgfusepath{clip}%
\pgfsetbuttcap%
\pgfsetmiterjoin%
\definecolor{currentfill}{rgb}{0.783383,0.262500,0.483918}%
\pgfsetfillcolor{currentfill}%
\pgfsetlinewidth{0.000000pt}%
\definecolor{currentstroke}{rgb}{0.000000,0.000000,0.000000}%
\pgfsetstrokecolor{currentstroke}%
\pgfsetstrokeopacity{0.000000}%
\pgfsetdash{}{0pt}%
\pgfpathmoveto{\pgfqpoint{4.573521in}{0.340880in}}%
\pgfpathlineto{\pgfqpoint{4.561217in}{0.345303in}}%
\pgfpathlineto{\pgfqpoint{4.564744in}{0.356026in}}%
\pgfpathlineto{\pgfqpoint{4.559586in}{0.366952in}}%
\pgfpathlineto{\pgfqpoint{4.561206in}{0.374463in}}%
\pgfpathlineto{\pgfqpoint{4.571107in}{0.377326in}}%
\pgfpathlineto{\pgfqpoint{4.570916in}{0.365292in}}%
\pgfpathlineto{\pgfqpoint{4.582627in}{0.358618in}}%
\pgfpathlineto{\pgfqpoint{4.589326in}{0.342480in}}%
\pgfpathlineto{\pgfqpoint{4.581861in}{0.336266in}}%
\pgfpathlineto{\pgfqpoint{4.573521in}{0.340880in}}%
\pgfpathclose%
\pgfusepath{fill}%
\end{pgfscope}%
\begin{pgfscope}%
\pgfpathrectangle{\pgfqpoint{3.525000in}{0.100000in}}{\pgfqpoint{2.862985in}{1.832778in}}%
\pgfusepath{clip}%
\pgfsetbuttcap%
\pgfsetmiterjoin%
\definecolor{currentfill}{rgb}{0.783383,0.262500,0.483918}%
\pgfsetfillcolor{currentfill}%
\pgfsetlinewidth{0.000000pt}%
\definecolor{currentstroke}{rgb}{0.000000,0.000000,0.000000}%
\pgfsetstrokecolor{currentstroke}%
\pgfsetstrokeopacity{0.000000}%
\pgfsetdash{}{0pt}%
\pgfpathmoveto{\pgfqpoint{4.529084in}{0.231768in}}%
\pgfpathlineto{\pgfqpoint{4.523222in}{0.243877in}}%
\pgfpathlineto{\pgfqpoint{4.525818in}{0.252022in}}%
\pgfpathlineto{\pgfqpoint{4.537315in}{0.263296in}}%
\pgfpathlineto{\pgfqpoint{4.538404in}{0.271938in}}%
\pgfpathlineto{\pgfqpoint{4.545928in}{0.291072in}}%
\pgfpathlineto{\pgfqpoint{4.565718in}{0.294341in}}%
\pgfpathlineto{\pgfqpoint{4.569506in}{0.305541in}}%
\pgfpathlineto{\pgfqpoint{4.578232in}{0.301341in}}%
\pgfpathlineto{\pgfqpoint{4.595961in}{0.269908in}}%
\pgfpathlineto{\pgfqpoint{4.596185in}{0.259645in}}%
\pgfpathlineto{\pgfqpoint{4.591173in}{0.250763in}}%
\pgfpathlineto{\pgfqpoint{4.597532in}{0.230369in}}%
\pgfpathlineto{\pgfqpoint{4.593089in}{0.227081in}}%
\pgfpathlineto{\pgfqpoint{4.577955in}{0.228037in}}%
\pgfpathlineto{\pgfqpoint{4.559756in}{0.234119in}}%
\pgfpathlineto{\pgfqpoint{4.544460in}{0.236631in}}%
\pgfpathlineto{\pgfqpoint{4.540375in}{0.233063in}}%
\pgfpathlineto{\pgfqpoint{4.529084in}{0.231768in}}%
\pgfpathclose%
\pgfusepath{fill}%
\end{pgfscope}%
\begin{pgfscope}%
\pgfpathrectangle{\pgfqpoint{3.525000in}{0.100000in}}{\pgfqpoint{2.862985in}{1.832778in}}%
\pgfusepath{clip}%
\pgfsetbuttcap%
\pgfsetmiterjoin%
\definecolor{currentfill}{rgb}{0.679160,0.151848,0.575189}%
\pgfsetfillcolor{currentfill}%
\pgfsetlinewidth{0.000000pt}%
\definecolor{currentstroke}{rgb}{0.000000,0.000000,0.000000}%
\pgfsetstrokecolor{currentstroke}%
\pgfsetstrokeopacity{0.000000}%
\pgfsetdash{}{0pt}%
\pgfpathmoveto{\pgfqpoint{4.205273in}{1.824988in}}%
\pgfpathlineto{\pgfqpoint{4.190426in}{1.766159in}}%
\pgfpathlineto{\pgfqpoint{4.179869in}{1.724007in}}%
\pgfpathlineto{\pgfqpoint{4.167317in}{1.673299in}}%
\pgfpathlineto{\pgfqpoint{4.164852in}{1.660801in}}%
\pgfpathlineto{\pgfqpoint{4.166462in}{1.649338in}}%
\pgfpathlineto{\pgfqpoint{4.164329in}{1.638779in}}%
\pgfpathlineto{\pgfqpoint{4.076709in}{1.661643in}}%
\pgfpathlineto{\pgfqpoint{4.068794in}{1.658987in}}%
\pgfpathlineto{\pgfqpoint{4.062021in}{1.661292in}}%
\pgfpathlineto{\pgfqpoint{4.030660in}{1.662045in}}%
\pgfpathlineto{\pgfqpoint{4.020053in}{1.659055in}}%
\pgfpathlineto{\pgfqpoint{4.009520in}{1.660065in}}%
\pgfpathlineto{\pgfqpoint{4.005055in}{1.664746in}}%
\pgfpathlineto{\pgfqpoint{3.976979in}{1.663725in}}%
\pgfpathlineto{\pgfqpoint{3.972545in}{1.671374in}}%
\pgfpathlineto{\pgfqpoint{3.964127in}{1.675404in}}%
\pgfpathlineto{\pgfqpoint{3.951880in}{1.677860in}}%
\pgfpathlineto{\pgfqpoint{3.930740in}{1.674264in}}%
\pgfpathlineto{\pgfqpoint{3.922937in}{1.677787in}}%
\pgfpathlineto{\pgfqpoint{3.910905in}{1.687156in}}%
\pgfpathlineto{\pgfqpoint{3.914522in}{1.705519in}}%
\pgfpathlineto{\pgfqpoint{3.913254in}{1.714897in}}%
\pgfpathlineto{\pgfqpoint{3.903718in}{1.724852in}}%
\pgfpathlineto{\pgfqpoint{3.897610in}{1.724234in}}%
\pgfpathlineto{\pgfqpoint{3.893229in}{1.734304in}}%
\pgfpathlineto{\pgfqpoint{3.882849in}{1.738374in}}%
\pgfpathlineto{\pgfqpoint{3.875305in}{1.737856in}}%
\pgfpathlineto{\pgfqpoint{3.872864in}{1.748198in}}%
\pgfpathlineto{\pgfqpoint{3.880380in}{1.747098in}}%
\pgfpathlineto{\pgfqpoint{3.879783in}{1.761472in}}%
\pgfpathlineto{\pgfqpoint{3.876507in}{1.770008in}}%
\pgfpathlineto{\pgfqpoint{3.877951in}{1.780932in}}%
\pgfpathlineto{\pgfqpoint{3.881884in}{1.789149in}}%
\pgfpathlineto{\pgfqpoint{3.881662in}{1.804741in}}%
\pgfpathlineto{\pgfqpoint{3.879563in}{1.810398in}}%
\pgfpathlineto{\pgfqpoint{3.883188in}{1.828548in}}%
\pgfpathlineto{\pgfqpoint{3.882142in}{1.840246in}}%
\pgfpathlineto{\pgfqpoint{3.878517in}{1.845854in}}%
\pgfpathlineto{\pgfqpoint{3.879070in}{1.864192in}}%
\pgfpathlineto{\pgfqpoint{3.884342in}{1.877678in}}%
\pgfpathlineto{\pgfqpoint{3.890077in}{1.874378in}}%
\pgfpathlineto{\pgfqpoint{3.909064in}{1.854764in}}%
\pgfpathlineto{\pgfqpoint{3.932071in}{1.844031in}}%
\pgfpathlineto{\pgfqpoint{3.943852in}{1.842761in}}%
\pgfpathlineto{\pgfqpoint{3.954951in}{1.838191in}}%
\pgfpathlineto{\pgfqpoint{3.958101in}{1.822814in}}%
\pgfpathlineto{\pgfqpoint{3.948375in}{1.819901in}}%
\pgfpathlineto{\pgfqpoint{3.941929in}{1.807782in}}%
\pgfpathlineto{\pgfqpoint{3.949497in}{1.808455in}}%
\pgfpathlineto{\pgfqpoint{3.952549in}{1.813895in}}%
\pgfpathlineto{\pgfqpoint{3.963221in}{1.820697in}}%
\pgfpathlineto{\pgfqpoint{3.962670in}{1.810596in}}%
\pgfpathlineto{\pgfqpoint{3.955544in}{1.808967in}}%
\pgfpathlineto{\pgfqpoint{3.956727in}{1.795983in}}%
\pgfpathlineto{\pgfqpoint{3.948050in}{1.782976in}}%
\pgfpathlineto{\pgfqpoint{3.942496in}{1.789161in}}%
\pgfpathlineto{\pgfqpoint{3.925991in}{1.785798in}}%
\pgfpathlineto{\pgfqpoint{3.925123in}{1.778232in}}%
\pgfpathlineto{\pgfqpoint{3.930813in}{1.773622in}}%
\pgfpathlineto{\pgfqpoint{3.939493in}{1.773169in}}%
\pgfpathlineto{\pgfqpoint{3.951495in}{1.783019in}}%
\pgfpathlineto{\pgfqpoint{3.960994in}{1.783837in}}%
\pgfpathlineto{\pgfqpoint{3.961358in}{1.794741in}}%
\pgfpathlineto{\pgfqpoint{3.966254in}{1.810756in}}%
\pgfpathlineto{\pgfqpoint{3.977856in}{1.822666in}}%
\pgfpathlineto{\pgfqpoint{3.973988in}{1.831891in}}%
\pgfpathlineto{\pgfqpoint{3.976630in}{1.841823in}}%
\pgfpathlineto{\pgfqpoint{3.971483in}{1.851722in}}%
\pgfpathlineto{\pgfqpoint{3.978818in}{1.864381in}}%
\pgfpathlineto{\pgfqpoint{3.979917in}{1.871996in}}%
\pgfpathlineto{\pgfqpoint{3.973548in}{1.877001in}}%
\pgfpathlineto{\pgfqpoint{3.974609in}{1.889804in}}%
\pgfpathlineto{\pgfqpoint{4.050918in}{1.866620in}}%
\pgfpathlineto{\pgfqpoint{4.131952in}{1.843896in}}%
\pgfpathlineto{\pgfqpoint{4.205273in}{1.824988in}}%
\pgfpathclose%
\pgfusepath{fill}%
\end{pgfscope}%
\begin{pgfscope}%
\pgfpathrectangle{\pgfqpoint{3.525000in}{0.100000in}}{\pgfqpoint{2.862985in}{1.832778in}}%
\pgfusepath{clip}%
\pgfsetbuttcap%
\pgfsetmiterjoin%
\definecolor{currentfill}{rgb}{0.679160,0.151848,0.575189}%
\pgfsetfillcolor{currentfill}%
\pgfsetlinewidth{0.000000pt}%
\definecolor{currentstroke}{rgb}{0.000000,0.000000,0.000000}%
\pgfsetstrokecolor{currentstroke}%
\pgfsetstrokeopacity{0.000000}%
\pgfsetdash{}{0pt}%
\pgfpathmoveto{\pgfqpoint{3.962056in}{1.845512in}}%
\pgfpathlineto{\pgfqpoint{3.965813in}{1.831026in}}%
\pgfpathlineto{\pgfqpoint{3.971706in}{1.826214in}}%
\pgfpathlineto{\pgfqpoint{3.967004in}{1.820089in}}%
\pgfpathlineto{\pgfqpoint{3.962421in}{1.829064in}}%
\pgfpathlineto{\pgfqpoint{3.962056in}{1.845512in}}%
\pgfpathclose%
\pgfusepath{fill}%
\end{pgfscope}%
\begin{pgfscope}%
\pgfpathrectangle{\pgfqpoint{3.525000in}{0.100000in}}{\pgfqpoint{2.862985in}{1.832778in}}%
\pgfusepath{clip}%
\pgfsetbuttcap%
\pgfsetmiterjoin%
\definecolor{currentfill}{rgb}{0.692840,0.165141,0.564522}%
\pgfsetfillcolor{currentfill}%
\pgfsetlinewidth{0.000000pt}%
\definecolor{currentstroke}{rgb}{0.000000,0.000000,0.000000}%
\pgfsetstrokecolor{currentstroke}%
\pgfsetstrokeopacity{0.000000}%
\pgfsetdash{}{0pt}%
\pgfpathmoveto{\pgfqpoint{4.244624in}{1.815480in}}%
\pgfpathlineto{\pgfqpoint{4.326194in}{1.797175in}}%
\pgfpathlineto{\pgfqpoint{4.403016in}{1.781647in}}%
\pgfpathlineto{\pgfqpoint{4.462121in}{1.770817in}}%
\pgfpathlineto{\pgfqpoint{4.513650in}{1.762160in}}%
\pgfpathlineto{\pgfqpoint{4.565293in}{1.754214in}}%
\pgfpathlineto{\pgfqpoint{4.609272in}{1.748018in}}%
\pgfpathlineto{\pgfqpoint{4.653320in}{1.742337in}}%
\pgfpathlineto{\pgfqpoint{4.697432in}{1.737171in}}%
\pgfpathlineto{\pgfqpoint{4.739001in}{1.732781in}}%
\pgfpathlineto{\pgfqpoint{4.733284in}{1.669428in}}%
\pgfpathlineto{\pgfqpoint{4.724775in}{1.583815in}}%
\pgfpathlineto{\pgfqpoint{4.720312in}{1.539783in}}%
\pgfpathlineto{\pgfqpoint{4.713893in}{1.480441in}}%
\pgfpathlineto{\pgfqpoint{4.668453in}{1.485403in}}%
\pgfpathlineto{\pgfqpoint{4.616429in}{1.491308in}}%
\pgfpathlineto{\pgfqpoint{4.544194in}{1.501195in}}%
\pgfpathlineto{\pgfqpoint{4.511931in}{1.505797in}}%
\pgfpathlineto{\pgfqpoint{4.432445in}{1.518274in}}%
\pgfpathlineto{\pgfqpoint{4.405083in}{1.523280in}}%
\pgfpathlineto{\pgfqpoint{4.399380in}{1.490854in}}%
\pgfpathlineto{\pgfqpoint{4.396262in}{1.493168in}}%
\pgfpathlineto{\pgfqpoint{4.390403in}{1.508767in}}%
\pgfpathlineto{\pgfqpoint{4.382871in}{1.507806in}}%
\pgfpathlineto{\pgfqpoint{4.377480in}{1.499920in}}%
\pgfpathlineto{\pgfqpoint{4.366293in}{1.498866in}}%
\pgfpathlineto{\pgfqpoint{4.364045in}{1.502638in}}%
\pgfpathlineto{\pgfqpoint{4.353566in}{1.502181in}}%
\pgfpathlineto{\pgfqpoint{4.348297in}{1.505835in}}%
\pgfpathlineto{\pgfqpoint{4.340931in}{1.500183in}}%
\pgfpathlineto{\pgfqpoint{4.323038in}{1.505269in}}%
\pgfpathlineto{\pgfqpoint{4.315230in}{1.502538in}}%
\pgfpathlineto{\pgfqpoint{4.311029in}{1.514658in}}%
\pgfpathlineto{\pgfqpoint{4.310833in}{1.526190in}}%
\pgfpathlineto{\pgfqpoint{4.298499in}{1.534484in}}%
\pgfpathlineto{\pgfqpoint{4.300510in}{1.546436in}}%
\pgfpathlineto{\pgfqpoint{4.291728in}{1.566180in}}%
\pgfpathlineto{\pgfqpoint{4.292319in}{1.582964in}}%
\pgfpathlineto{\pgfqpoint{4.284772in}{1.591178in}}%
\pgfpathlineto{\pgfqpoint{4.277690in}{1.583911in}}%
\pgfpathlineto{\pgfqpoint{4.268050in}{1.579980in}}%
\pgfpathlineto{\pgfqpoint{4.260712in}{1.588082in}}%
\pgfpathlineto{\pgfqpoint{4.261905in}{1.601019in}}%
\pgfpathlineto{\pgfqpoint{4.270705in}{1.606116in}}%
\pgfpathlineto{\pgfqpoint{4.268366in}{1.614312in}}%
\pgfpathlineto{\pgfqpoint{4.283663in}{1.653985in}}%
\pgfpathlineto{\pgfqpoint{4.271601in}{1.654979in}}%
\pgfpathlineto{\pgfqpoint{4.270361in}{1.662098in}}%
\pgfpathlineto{\pgfqpoint{4.261575in}{1.668238in}}%
\pgfpathlineto{\pgfqpoint{4.262137in}{1.675101in}}%
\pgfpathlineto{\pgfqpoint{4.257502in}{1.680433in}}%
\pgfpathlineto{\pgfqpoint{4.249957in}{1.699791in}}%
\pgfpathlineto{\pgfqpoint{4.243063in}{1.703739in}}%
\pgfpathlineto{\pgfqpoint{4.238126in}{1.720446in}}%
\pgfpathlineto{\pgfqpoint{4.239291in}{1.731551in}}%
\pgfpathlineto{\pgfqpoint{4.230080in}{1.752004in}}%
\pgfpathlineto{\pgfqpoint{4.244624in}{1.815480in}}%
\pgfpathclose%
\pgfusepath{fill}%
\end{pgfscope}%
\begin{pgfscope}%
\pgfpathrectangle{\pgfqpoint{3.525000in}{0.100000in}}{\pgfqpoint{2.862985in}{1.832778in}}%
\pgfusepath{clip}%
\pgfsetbuttcap%
\pgfsetmiterjoin%
\definecolor{currentfill}{rgb}{0.846788,0.342551,0.420579}%
\pgfsetfillcolor{currentfill}%
\pgfsetlinewidth{0.000000pt}%
\definecolor{currentstroke}{rgb}{0.000000,0.000000,0.000000}%
\pgfsetstrokecolor{currentstroke}%
\pgfsetstrokeopacity{0.000000}%
\pgfsetdash{}{0pt}%
\pgfpathmoveto{\pgfqpoint{6.213801in}{1.489012in}}%
\pgfpathlineto{\pgfqpoint{6.212168in}{1.495954in}}%
\pgfpathlineto{\pgfqpoint{6.203094in}{1.502035in}}%
\pgfpathlineto{\pgfqpoint{6.179117in}{1.580468in}}%
\pgfpathlineto{\pgfqpoint{6.166192in}{1.618318in}}%
\pgfpathlineto{\pgfqpoint{6.176583in}{1.629222in}}%
\pgfpathlineto{\pgfqpoint{6.186683in}{1.655908in}}%
\pgfpathlineto{\pgfqpoint{6.191720in}{1.664437in}}%
\pgfpathlineto{\pgfqpoint{6.188152in}{1.668018in}}%
\pgfpathlineto{\pgfqpoint{6.186957in}{1.691707in}}%
\pgfpathlineto{\pgfqpoint{6.191463in}{1.698969in}}%
\pgfpathlineto{\pgfqpoint{6.189488in}{1.715849in}}%
\pgfpathlineto{\pgfqpoint{6.207474in}{1.771228in}}%
\pgfpathlineto{\pgfqpoint{6.215578in}{1.771530in}}%
\pgfpathlineto{\pgfqpoint{6.218942in}{1.761638in}}%
\pgfpathlineto{\pgfqpoint{6.226144in}{1.758798in}}%
\pgfpathlineto{\pgfqpoint{6.239749in}{1.770433in}}%
\pgfpathlineto{\pgfqpoint{6.250417in}{1.777306in}}%
\pgfpathlineto{\pgfqpoint{6.273922in}{1.765204in}}%
\pgfpathlineto{\pgfqpoint{6.295157in}{1.698008in}}%
\pgfpathlineto{\pgfqpoint{6.299194in}{1.681476in}}%
\pgfpathlineto{\pgfqpoint{6.308495in}{1.679537in}}%
\pgfpathlineto{\pgfqpoint{6.321110in}{1.668309in}}%
\pgfpathlineto{\pgfqpoint{6.320396in}{1.661767in}}%
\pgfpathlineto{\pgfqpoint{6.328989in}{1.654013in}}%
\pgfpathlineto{\pgfqpoint{6.335970in}{1.658462in}}%
\pgfpathlineto{\pgfqpoint{6.350577in}{1.641426in}}%
\pgfpathlineto{\pgfqpoint{6.344049in}{1.627792in}}%
\pgfpathlineto{\pgfqpoint{6.335294in}{1.627522in}}%
\pgfpathlineto{\pgfqpoint{6.328271in}{1.615351in}}%
\pgfpathlineto{\pgfqpoint{6.319762in}{1.613625in}}%
\pgfpathlineto{\pgfqpoint{6.313516in}{1.607145in}}%
\pgfpathlineto{\pgfqpoint{6.294860in}{1.600195in}}%
\pgfpathlineto{\pgfqpoint{6.283568in}{1.588995in}}%
\pgfpathlineto{\pgfqpoint{6.277832in}{1.597034in}}%
\pgfpathlineto{\pgfqpoint{6.272620in}{1.591272in}}%
\pgfpathlineto{\pgfqpoint{6.274047in}{1.567979in}}%
\pgfpathlineto{\pgfqpoint{6.269920in}{1.558753in}}%
\pgfpathlineto{\pgfqpoint{6.260918in}{1.561273in}}%
\pgfpathlineto{\pgfqpoint{6.259465in}{1.551814in}}%
\pgfpathlineto{\pgfqpoint{6.255576in}{1.547924in}}%
\pgfpathlineto{\pgfqpoint{6.247791in}{1.548872in}}%
\pgfpathlineto{\pgfqpoint{6.248281in}{1.540028in}}%
\pgfpathlineto{\pgfqpoint{6.236482in}{1.542540in}}%
\pgfpathlineto{\pgfqpoint{6.230038in}{1.530287in}}%
\pgfpathlineto{\pgfqpoint{6.232472in}{1.523875in}}%
\pgfpathlineto{\pgfqpoint{6.228639in}{1.513212in}}%
\pgfpathlineto{\pgfqpoint{6.222604in}{1.505369in}}%
\pgfpathlineto{\pgfqpoint{6.221086in}{1.488995in}}%
\pgfpathlineto{\pgfqpoint{6.213801in}{1.489012in}}%
\pgfpathclose%
\pgfusepath{fill}%
\end{pgfscope}%
\begin{pgfscope}%
\pgfpathrectangle{\pgfqpoint{3.525000in}{0.100000in}}{\pgfqpoint{2.862985in}{1.832778in}}%
\pgfusepath{clip}%
\pgfsetbuttcap%
\pgfsetmiterjoin%
\definecolor{currentfill}{rgb}{0.846788,0.342551,0.420579}%
\pgfsetfillcolor{currentfill}%
\pgfsetlinewidth{0.000000pt}%
\definecolor{currentstroke}{rgb}{0.000000,0.000000,0.000000}%
\pgfsetstrokecolor{currentstroke}%
\pgfsetstrokeopacity{0.000000}%
\pgfsetdash{}{0pt}%
\pgfpathmoveto{\pgfqpoint{6.298211in}{1.595412in}}%
\pgfpathlineto{\pgfqpoint{6.303538in}{1.600991in}}%
\pgfpathlineto{\pgfqpoint{6.308609in}{1.595747in}}%
\pgfpathlineto{\pgfqpoint{6.299500in}{1.588798in}}%
\pgfpathlineto{\pgfqpoint{6.298211in}{1.595412in}}%
\pgfpathclose%
\pgfusepath{fill}%
\end{pgfscope}%
\begin{pgfscope}%
\pgfpathrectangle{\pgfqpoint{3.525000in}{0.100000in}}{\pgfqpoint{2.862985in}{1.832778in}}%
\pgfusepath{clip}%
\pgfsetbuttcap%
\pgfsetmiterjoin%
\definecolor{currentfill}{rgb}{0.429719,0.000831,0.659425}%
\pgfsetfillcolor{currentfill}%
\pgfsetlinewidth{0.000000pt}%
\definecolor{currentstroke}{rgb}{0.000000,0.000000,0.000000}%
\pgfsetstrokecolor{currentstroke}%
\pgfsetstrokeopacity{0.000000}%
\pgfsetdash{}{0pt}%
\pgfpathmoveto{\pgfqpoint{4.720312in}{1.539783in}}%
\pgfpathlineto{\pgfqpoint{4.724775in}{1.583815in}}%
\pgfpathlineto{\pgfqpoint{4.733284in}{1.669428in}}%
\pgfpathlineto{\pgfqpoint{4.739001in}{1.732781in}}%
\pgfpathlineto{\pgfqpoint{4.785821in}{1.728389in}}%
\pgfpathlineto{\pgfqpoint{4.845718in}{1.723620in}}%
\pgfpathlineto{\pgfqpoint{4.900466in}{1.720093in}}%
\pgfpathlineto{\pgfqpoint{4.950039in}{1.717585in}}%
\pgfpathlineto{\pgfqpoint{5.024019in}{1.715047in}}%
\pgfpathlineto{\pgfqpoint{5.029068in}{1.694197in}}%
\pgfpathlineto{\pgfqpoint{5.027261in}{1.684232in}}%
\pgfpathlineto{\pgfqpoint{5.026666in}{1.663751in}}%
\pgfpathlineto{\pgfqpoint{5.030101in}{1.648413in}}%
\pgfpathlineto{\pgfqpoint{5.037970in}{1.625784in}}%
\pgfpathlineto{\pgfqpoint{5.037952in}{1.597294in}}%
\pgfpathlineto{\pgfqpoint{5.039408in}{1.564197in}}%
\pgfpathlineto{\pgfqpoint{5.041408in}{1.555277in}}%
\pgfpathlineto{\pgfqpoint{5.047257in}{1.545485in}}%
\pgfpathlineto{\pgfqpoint{5.049186in}{1.530226in}}%
\pgfpathlineto{\pgfqpoint{5.048353in}{1.520039in}}%
\pgfpathlineto{\pgfqpoint{4.986333in}{1.521383in}}%
\pgfpathlineto{\pgfqpoint{4.941215in}{1.523670in}}%
\pgfpathlineto{\pgfqpoint{4.875074in}{1.527191in}}%
\pgfpathlineto{\pgfqpoint{4.809849in}{1.531824in}}%
\pgfpathlineto{\pgfqpoint{4.766407in}{1.535329in}}%
\pgfpathlineto{\pgfqpoint{4.720312in}{1.539783in}}%
\pgfpathclose%
\pgfusepath{fill}%
\end{pgfscope}%
\begin{pgfscope}%
\pgfpathrectangle{\pgfqpoint{3.525000in}{0.100000in}}{\pgfqpoint{2.862985in}{1.832778in}}%
\pgfusepath{clip}%
\pgfsetbuttcap%
\pgfsetmiterjoin%
\definecolor{currentfill}{rgb}{0.477344,0.006980,0.659549}%
\pgfsetfillcolor{currentfill}%
\pgfsetlinewidth{0.000000pt}%
\definecolor{currentstroke}{rgb}{0.000000,0.000000,0.000000}%
\pgfsetstrokecolor{currentstroke}%
\pgfsetstrokeopacity{0.000000}%
\pgfsetdash{}{0pt}%
\pgfpathmoveto{\pgfqpoint{4.701565in}{1.355506in}}%
\pgfpathlineto{\pgfqpoint{4.708725in}{1.429317in}}%
\pgfpathlineto{\pgfqpoint{4.713893in}{1.480441in}}%
\pgfpathlineto{\pgfqpoint{4.720312in}{1.539783in}}%
\pgfpathlineto{\pgfqpoint{4.766407in}{1.535329in}}%
\pgfpathlineto{\pgfqpoint{4.809849in}{1.531824in}}%
\pgfpathlineto{\pgfqpoint{4.875074in}{1.527191in}}%
\pgfpathlineto{\pgfqpoint{4.941215in}{1.523670in}}%
\pgfpathlineto{\pgfqpoint{4.986333in}{1.521383in}}%
\pgfpathlineto{\pgfqpoint{5.048353in}{1.520039in}}%
\pgfpathlineto{\pgfqpoint{5.044154in}{1.507786in}}%
\pgfpathlineto{\pgfqpoint{5.035772in}{1.498167in}}%
\pgfpathlineto{\pgfqpoint{5.042191in}{1.487092in}}%
\pgfpathlineto{\pgfqpoint{5.049268in}{1.484725in}}%
\pgfpathlineto{\pgfqpoint{5.052626in}{1.478363in}}%
\pgfpathlineto{\pgfqpoint{5.051855in}{1.431914in}}%
\pgfpathlineto{\pgfqpoint{5.050565in}{1.366544in}}%
\pgfpathlineto{\pgfqpoint{5.045791in}{1.349375in}}%
\pgfpathlineto{\pgfqpoint{5.050061in}{1.339878in}}%
\pgfpathlineto{\pgfqpoint{5.041489in}{1.319470in}}%
\pgfpathlineto{\pgfqpoint{5.050519in}{1.303018in}}%
\pgfpathlineto{\pgfqpoint{5.042847in}{1.304277in}}%
\pgfpathlineto{\pgfqpoint{5.037612in}{1.314513in}}%
\pgfpathlineto{\pgfqpoint{5.018930in}{1.321529in}}%
\pgfpathlineto{\pgfqpoint{5.007148in}{1.327701in}}%
\pgfpathlineto{\pgfqpoint{4.987377in}{1.328214in}}%
\pgfpathlineto{\pgfqpoint{4.980515in}{1.322591in}}%
\pgfpathlineto{\pgfqpoint{4.958137in}{1.333663in}}%
\pgfpathlineto{\pgfqpoint{4.956420in}{1.337164in}}%
\pgfpathlineto{\pgfqpoint{4.878322in}{1.340860in}}%
\pgfpathlineto{\pgfqpoint{4.830874in}{1.343580in}}%
\pgfpathlineto{\pgfqpoint{4.791682in}{1.346649in}}%
\pgfpathlineto{\pgfqpoint{4.726922in}{1.352768in}}%
\pgfpathlineto{\pgfqpoint{4.701565in}{1.355506in}}%
\pgfpathclose%
\pgfusepath{fill}%
\end{pgfscope}%
\begin{pgfscope}%
\pgfpathrectangle{\pgfqpoint{3.525000in}{0.100000in}}{\pgfqpoint{2.862985in}{1.832778in}}%
\pgfusepath{clip}%
\pgfsetbuttcap%
\pgfsetmiterjoin%
\definecolor{currentfill}{rgb}{0.669845,0.142992,0.582154}%
\pgfsetfillcolor{currentfill}%
\pgfsetlinewidth{0.000000pt}%
\definecolor{currentstroke}{rgb}{0.000000,0.000000,0.000000}%
\pgfsetstrokecolor{currentstroke}%
\pgfsetstrokeopacity{0.000000}%
\pgfsetdash{}{0pt}%
\pgfpathmoveto{\pgfqpoint{4.689283in}{1.230524in}}%
\pgfpathlineto{\pgfqpoint{4.647647in}{1.234326in}}%
\pgfpathlineto{\pgfqpoint{4.556892in}{1.245527in}}%
\pgfpathlineto{\pgfqpoint{4.507522in}{1.252613in}}%
\pgfpathlineto{\pgfqpoint{4.454571in}{1.260214in}}%
\pgfpathlineto{\pgfqpoint{4.409967in}{1.267352in}}%
\pgfpathlineto{\pgfqpoint{4.361009in}{1.275734in}}%
\pgfpathlineto{\pgfqpoint{4.372141in}{1.337483in}}%
\pgfpathlineto{\pgfqpoint{4.383416in}{1.400800in}}%
\pgfpathlineto{\pgfqpoint{4.399380in}{1.490854in}}%
\pgfpathlineto{\pgfqpoint{4.405083in}{1.523280in}}%
\pgfpathlineto{\pgfqpoint{4.432445in}{1.518274in}}%
\pgfpathlineto{\pgfqpoint{4.511931in}{1.505797in}}%
\pgfpathlineto{\pgfqpoint{4.544194in}{1.501195in}}%
\pgfpathlineto{\pgfqpoint{4.616429in}{1.491308in}}%
\pgfpathlineto{\pgfqpoint{4.668453in}{1.485403in}}%
\pgfpathlineto{\pgfqpoint{4.713893in}{1.480441in}}%
\pgfpathlineto{\pgfqpoint{4.708725in}{1.429317in}}%
\pgfpathlineto{\pgfqpoint{4.701565in}{1.355506in}}%
\pgfpathlineto{\pgfqpoint{4.695419in}{1.292777in}}%
\pgfpathlineto{\pgfqpoint{4.689283in}{1.230524in}}%
\pgfpathclose%
\pgfusepath{fill}%
\end{pgfscope}%
\begin{pgfscope}%
\pgfpathrectangle{\pgfqpoint{3.525000in}{0.100000in}}{\pgfqpoint{2.862985in}{1.832778in}}%
\pgfusepath{clip}%
\pgfsetbuttcap%
\pgfsetmiterjoin%
\definecolor{currentfill}{rgb}{0.625987,0.103312,0.611305}%
\pgfsetfillcolor{currentfill}%
\pgfsetlinewidth{0.000000pt}%
\definecolor{currentstroke}{rgb}{0.000000,0.000000,0.000000}%
\pgfsetstrokecolor{currentstroke}%
\pgfsetstrokeopacity{0.000000}%
\pgfsetdash{}{0pt}%
\pgfpathmoveto{\pgfqpoint{5.450520in}{1.315700in}}%
\pgfpathlineto{\pgfqpoint{5.397758in}{1.311952in}}%
\pgfpathlineto{\pgfqpoint{5.319114in}{1.308600in}}%
\pgfpathlineto{\pgfqpoint{5.316125in}{1.316544in}}%
\pgfpathlineto{\pgfqpoint{5.298686in}{1.322483in}}%
\pgfpathlineto{\pgfqpoint{5.294839in}{1.333702in}}%
\pgfpathlineto{\pgfqpoint{5.293228in}{1.347582in}}%
\pgfpathlineto{\pgfqpoint{5.297158in}{1.354693in}}%
\pgfpathlineto{\pgfqpoint{5.290951in}{1.361517in}}%
\pgfpathlineto{\pgfqpoint{5.289454in}{1.369657in}}%
\pgfpathlineto{\pgfqpoint{5.287453in}{1.387663in}}%
\pgfpathlineto{\pgfqpoint{5.281505in}{1.397444in}}%
\pgfpathlineto{\pgfqpoint{5.270940in}{1.402933in}}%
\pgfpathlineto{\pgfqpoint{5.259444in}{1.411992in}}%
\pgfpathlineto{\pgfqpoint{5.253501in}{1.422710in}}%
\pgfpathlineto{\pgfqpoint{5.242838in}{1.427026in}}%
\pgfpathlineto{\pgfqpoint{5.236570in}{1.434051in}}%
\pgfpathlineto{\pgfqpoint{5.228976in}{1.435243in}}%
\pgfpathlineto{\pgfqpoint{5.215419in}{1.445673in}}%
\pgfpathlineto{\pgfqpoint{5.217621in}{1.457675in}}%
\pgfpathlineto{\pgfqpoint{5.217197in}{1.480500in}}%
\pgfpathlineto{\pgfqpoint{5.221378in}{1.486783in}}%
\pgfpathlineto{\pgfqpoint{5.217621in}{1.496270in}}%
\pgfpathlineto{\pgfqpoint{5.211004in}{1.498105in}}%
\pgfpathlineto{\pgfqpoint{5.211550in}{1.506437in}}%
\pgfpathlineto{\pgfqpoint{5.219761in}{1.519600in}}%
\pgfpathlineto{\pgfqpoint{5.236020in}{1.530009in}}%
\pgfpathlineto{\pgfqpoint{5.235007in}{1.567025in}}%
\pgfpathlineto{\pgfqpoint{5.243143in}{1.572586in}}%
\pgfpathlineto{\pgfqpoint{5.250842in}{1.568882in}}%
\pgfpathlineto{\pgfqpoint{5.266526in}{1.574304in}}%
\pgfpathlineto{\pgfqpoint{5.296034in}{1.587932in}}%
\pgfpathlineto{\pgfqpoint{5.299878in}{1.583713in}}%
\pgfpathlineto{\pgfqpoint{5.294289in}{1.564600in}}%
\pgfpathlineto{\pgfqpoint{5.302604in}{1.568794in}}%
\pgfpathlineto{\pgfqpoint{5.316845in}{1.564606in}}%
\pgfpathlineto{\pgfqpoint{5.325596in}{1.561110in}}%
\pgfpathlineto{\pgfqpoint{5.330502in}{1.550866in}}%
\pgfpathlineto{\pgfqpoint{5.375397in}{1.541198in}}%
\pgfpathlineto{\pgfqpoint{5.388816in}{1.534565in}}%
\pgfpathlineto{\pgfqpoint{5.402460in}{1.534643in}}%
\pgfpathlineto{\pgfqpoint{5.416480in}{1.531907in}}%
\pgfpathlineto{\pgfqpoint{5.425578in}{1.522561in}}%
\pgfpathlineto{\pgfqpoint{5.434272in}{1.517943in}}%
\pgfpathlineto{\pgfqpoint{5.435870in}{1.504623in}}%
\pgfpathlineto{\pgfqpoint{5.433301in}{1.496259in}}%
\pgfpathlineto{\pgfqpoint{5.442991in}{1.495641in}}%
\pgfpathlineto{\pgfqpoint{5.439727in}{1.485941in}}%
\pgfpathlineto{\pgfqpoint{5.442832in}{1.482495in}}%
\pgfpathlineto{\pgfqpoint{5.445919in}{1.473339in}}%
\pgfpathlineto{\pgfqpoint{5.436476in}{1.468491in}}%
\pgfpathlineto{\pgfqpoint{5.430993in}{1.454984in}}%
\pgfpathlineto{\pgfqpoint{5.429273in}{1.445433in}}%
\pgfpathlineto{\pgfqpoint{5.434533in}{1.443798in}}%
\pgfpathlineto{\pgfqpoint{5.441268in}{1.450960in}}%
\pgfpathlineto{\pgfqpoint{5.446994in}{1.463372in}}%
\pgfpathlineto{\pgfqpoint{5.454741in}{1.467683in}}%
\pgfpathlineto{\pgfqpoint{5.460563in}{1.462075in}}%
\pgfpathlineto{\pgfqpoint{5.454832in}{1.445104in}}%
\pgfpathlineto{\pgfqpoint{5.453034in}{1.431928in}}%
\pgfpathlineto{\pgfqpoint{5.454728in}{1.422457in}}%
\pgfpathlineto{\pgfqpoint{5.449406in}{1.417099in}}%
\pgfpathlineto{\pgfqpoint{5.446748in}{1.404007in}}%
\pgfpathlineto{\pgfqpoint{5.448920in}{1.390247in}}%
\pgfpathlineto{\pgfqpoint{5.442640in}{1.369897in}}%
\pgfpathlineto{\pgfqpoint{5.442772in}{1.359728in}}%
\pgfpathlineto{\pgfqpoint{5.447734in}{1.337693in}}%
\pgfpathlineto{\pgfqpoint{5.450951in}{1.333913in}}%
\pgfpathlineto{\pgfqpoint{5.450520in}{1.315700in}}%
\pgfpathclose%
\pgfusepath{fill}%
\end{pgfscope}%
\begin{pgfscope}%
\pgfpathrectangle{\pgfqpoint{3.525000in}{0.100000in}}{\pgfqpoint{2.862985in}{1.832778in}}%
\pgfusepath{clip}%
\pgfsetbuttcap%
\pgfsetmiterjoin%
\definecolor{currentfill}{rgb}{0.625987,0.103312,0.611305}%
\pgfsetfillcolor{currentfill}%
\pgfsetlinewidth{0.000000pt}%
\definecolor{currentstroke}{rgb}{0.000000,0.000000,0.000000}%
\pgfsetstrokecolor{currentstroke}%
\pgfsetstrokeopacity{0.000000}%
\pgfsetdash{}{0pt}%
\pgfpathmoveto{\pgfqpoint{5.470302in}{1.494291in}}%
\pgfpathlineto{\pgfqpoint{5.471480in}{1.485502in}}%
\pgfpathlineto{\pgfqpoint{5.460690in}{1.462343in}}%
\pgfpathlineto{\pgfqpoint{5.455896in}{1.469051in}}%
\pgfpathlineto{\pgfqpoint{5.470302in}{1.494291in}}%
\pgfpathclose%
\pgfusepath{fill}%
\end{pgfscope}%
\begin{pgfscope}%
\pgfpathrectangle{\pgfqpoint{3.525000in}{0.100000in}}{\pgfqpoint{2.862985in}{1.832778in}}%
\pgfusepath{clip}%
\pgfsetbuttcap%
\pgfsetmiterjoin%
\definecolor{currentfill}{rgb}{0.688318,0.160709,0.568103}%
\pgfsetfillcolor{currentfill}%
\pgfsetlinewidth{0.000000pt}%
\definecolor{currentstroke}{rgb}{0.000000,0.000000,0.000000}%
\pgfsetstrokecolor{currentstroke}%
\pgfsetstrokeopacity{0.000000}%
\pgfsetdash{}{0pt}%
\pgfpathmoveto{\pgfqpoint{4.164329in}{1.638779in}}%
\pgfpathlineto{\pgfqpoint{4.166462in}{1.649338in}}%
\pgfpathlineto{\pgfqpoint{4.164852in}{1.660801in}}%
\pgfpathlineto{\pgfqpoint{4.167317in}{1.673299in}}%
\pgfpathlineto{\pgfqpoint{4.179869in}{1.724007in}}%
\pgfpathlineto{\pgfqpoint{4.190426in}{1.766159in}}%
\pgfpathlineto{\pgfqpoint{4.205273in}{1.824988in}}%
\pgfpathlineto{\pgfqpoint{4.244624in}{1.815480in}}%
\pgfpathlineto{\pgfqpoint{4.230080in}{1.752004in}}%
\pgfpathlineto{\pgfqpoint{4.239291in}{1.731551in}}%
\pgfpathlineto{\pgfqpoint{4.238126in}{1.720446in}}%
\pgfpathlineto{\pgfqpoint{4.243063in}{1.703739in}}%
\pgfpathlineto{\pgfqpoint{4.249957in}{1.699791in}}%
\pgfpathlineto{\pgfqpoint{4.257502in}{1.680433in}}%
\pgfpathlineto{\pgfqpoint{4.262137in}{1.675101in}}%
\pgfpathlineto{\pgfqpoint{4.261575in}{1.668238in}}%
\pgfpathlineto{\pgfqpoint{4.270361in}{1.662098in}}%
\pgfpathlineto{\pgfqpoint{4.271601in}{1.654979in}}%
\pgfpathlineto{\pgfqpoint{4.283663in}{1.653985in}}%
\pgfpathlineto{\pgfqpoint{4.268366in}{1.614312in}}%
\pgfpathlineto{\pgfqpoint{4.270705in}{1.606116in}}%
\pgfpathlineto{\pgfqpoint{4.261905in}{1.601019in}}%
\pgfpathlineto{\pgfqpoint{4.260712in}{1.588082in}}%
\pgfpathlineto{\pgfqpoint{4.268050in}{1.579980in}}%
\pgfpathlineto{\pgfqpoint{4.277690in}{1.583911in}}%
\pgfpathlineto{\pgfqpoint{4.284772in}{1.591178in}}%
\pgfpathlineto{\pgfqpoint{4.292319in}{1.582964in}}%
\pgfpathlineto{\pgfqpoint{4.291728in}{1.566180in}}%
\pgfpathlineto{\pgfqpoint{4.300510in}{1.546436in}}%
\pgfpathlineto{\pgfqpoint{4.298499in}{1.534484in}}%
\pgfpathlineto{\pgfqpoint{4.310833in}{1.526190in}}%
\pgfpathlineto{\pgfqpoint{4.311029in}{1.514658in}}%
\pgfpathlineto{\pgfqpoint{4.315230in}{1.502538in}}%
\pgfpathlineto{\pgfqpoint{4.323038in}{1.505269in}}%
\pgfpathlineto{\pgfqpoint{4.340931in}{1.500183in}}%
\pgfpathlineto{\pgfqpoint{4.348297in}{1.505835in}}%
\pgfpathlineto{\pgfqpoint{4.353566in}{1.502181in}}%
\pgfpathlineto{\pgfqpoint{4.364045in}{1.502638in}}%
\pgfpathlineto{\pgfqpoint{4.366293in}{1.498866in}}%
\pgfpathlineto{\pgfqpoint{4.377480in}{1.499920in}}%
\pgfpathlineto{\pgfqpoint{4.382871in}{1.507806in}}%
\pgfpathlineto{\pgfqpoint{4.390403in}{1.508767in}}%
\pgfpathlineto{\pgfqpoint{4.396262in}{1.493168in}}%
\pgfpathlineto{\pgfqpoint{4.399380in}{1.490854in}}%
\pgfpathlineto{\pgfqpoint{4.383416in}{1.400800in}}%
\pgfpathlineto{\pgfqpoint{4.372141in}{1.337483in}}%
\pgfpathlineto{\pgfqpoint{4.283212in}{1.354650in}}%
\pgfpathlineto{\pgfqpoint{4.235177in}{1.364205in}}%
\pgfpathlineto{\pgfqpoint{4.190246in}{1.374085in}}%
\pgfpathlineto{\pgfqpoint{4.148250in}{1.383593in}}%
\pgfpathlineto{\pgfqpoint{4.099647in}{1.395314in}}%
\pgfpathlineto{\pgfqpoint{4.124717in}{1.498158in}}%
\pgfpathlineto{\pgfqpoint{4.126069in}{1.505674in}}%
\pgfpathlineto{\pgfqpoint{4.137217in}{1.525368in}}%
\pgfpathlineto{\pgfqpoint{4.125673in}{1.537205in}}%
\pgfpathlineto{\pgfqpoint{4.128070in}{1.548828in}}%
\pgfpathlineto{\pgfqpoint{4.132464in}{1.551738in}}%
\pgfpathlineto{\pgfqpoint{4.140313in}{1.563715in}}%
\pgfpathlineto{\pgfqpoint{4.151752in}{1.572002in}}%
\pgfpathlineto{\pgfqpoint{4.152377in}{1.578133in}}%
\pgfpathlineto{\pgfqpoint{4.159282in}{1.584259in}}%
\pgfpathlineto{\pgfqpoint{4.165015in}{1.595725in}}%
\pgfpathlineto{\pgfqpoint{4.176774in}{1.607856in}}%
\pgfpathlineto{\pgfqpoint{4.175939in}{1.619838in}}%
\pgfpathlineto{\pgfqpoint{4.167899in}{1.626494in}}%
\pgfpathlineto{\pgfqpoint{4.164329in}{1.638779in}}%
\pgfpathclose%
\pgfusepath{fill}%
\end{pgfscope}%
\begin{pgfscope}%
\pgfpathrectangle{\pgfqpoint{3.525000in}{0.100000in}}{\pgfqpoint{2.862985in}{1.832778in}}%
\pgfusepath{clip}%
\pgfsetbuttcap%
\pgfsetmiterjoin%
\definecolor{currentfill}{rgb}{0.679160,0.151848,0.575189}%
\pgfsetfillcolor{currentfill}%
\pgfsetlinewidth{0.000000pt}%
\definecolor{currentstroke}{rgb}{0.000000,0.000000,0.000000}%
\pgfsetstrokecolor{currentstroke}%
\pgfsetstrokeopacity{0.000000}%
\pgfsetdash{}{0pt}%
\pgfpathmoveto{\pgfqpoint{6.111015in}{1.438080in}}%
\pgfpathlineto{\pgfqpoint{6.107320in}{1.449744in}}%
\pgfpathlineto{\pgfqpoint{6.100463in}{1.485127in}}%
\pgfpathlineto{\pgfqpoint{6.091598in}{1.496013in}}%
\pgfpathlineto{\pgfqpoint{6.083822in}{1.515590in}}%
\pgfpathlineto{\pgfqpoint{6.086375in}{1.530104in}}%
\pgfpathlineto{\pgfqpoint{6.084327in}{1.540801in}}%
\pgfpathlineto{\pgfqpoint{6.077683in}{1.551308in}}%
\pgfpathlineto{\pgfqpoint{6.077410in}{1.562882in}}%
\pgfpathlineto{\pgfqpoint{6.073552in}{1.575365in}}%
\pgfpathlineto{\pgfqpoint{6.108077in}{1.583862in}}%
\pgfpathlineto{\pgfqpoint{6.152922in}{1.595966in}}%
\pgfpathlineto{\pgfqpoint{6.154711in}{1.589025in}}%
\pgfpathlineto{\pgfqpoint{6.151856in}{1.577975in}}%
\pgfpathlineto{\pgfqpoint{6.158559in}{1.569133in}}%
\pgfpathlineto{\pgfqpoint{6.154999in}{1.557930in}}%
\pgfpathlineto{\pgfqpoint{6.140852in}{1.543846in}}%
\pgfpathlineto{\pgfqpoint{6.144741in}{1.533267in}}%
\pgfpathlineto{\pgfqpoint{6.141859in}{1.510906in}}%
\pgfpathlineto{\pgfqpoint{6.137824in}{1.498298in}}%
\pgfpathlineto{\pgfqpoint{6.142972in}{1.462444in}}%
\pgfpathlineto{\pgfqpoint{6.142278in}{1.449830in}}%
\pgfpathlineto{\pgfqpoint{6.147258in}{1.445782in}}%
\pgfpathlineto{\pgfqpoint{6.111015in}{1.438080in}}%
\pgfpathclose%
\pgfusepath{fill}%
\end{pgfscope}%
\begin{pgfscope}%
\pgfpathrectangle{\pgfqpoint{3.525000in}{0.100000in}}{\pgfqpoint{2.862985in}{1.832778in}}%
\pgfusepath{clip}%
\pgfsetbuttcap%
\pgfsetmiterjoin%
\definecolor{currentfill}{rgb}{0.471457,0.005678,0.659897}%
\pgfsetfillcolor{currentfill}%
\pgfsetlinewidth{0.000000pt}%
\definecolor{currentstroke}{rgb}{0.000000,0.000000,0.000000}%
\pgfsetstrokecolor{currentstroke}%
\pgfsetstrokeopacity{0.000000}%
\pgfsetdash{}{0pt}%
\pgfpathmoveto{\pgfqpoint{5.050565in}{1.366544in}}%
\pgfpathlineto{\pgfqpoint{5.051855in}{1.431914in}}%
\pgfpathlineto{\pgfqpoint{5.052626in}{1.478363in}}%
\pgfpathlineto{\pgfqpoint{5.049268in}{1.484725in}}%
\pgfpathlineto{\pgfqpoint{5.042191in}{1.487092in}}%
\pgfpathlineto{\pgfqpoint{5.035772in}{1.498167in}}%
\pgfpathlineto{\pgfqpoint{5.044154in}{1.507786in}}%
\pgfpathlineto{\pgfqpoint{5.048353in}{1.520039in}}%
\pgfpathlineto{\pgfqpoint{5.049186in}{1.530226in}}%
\pgfpathlineto{\pgfqpoint{5.047257in}{1.545485in}}%
\pgfpathlineto{\pgfqpoint{5.041408in}{1.555277in}}%
\pgfpathlineto{\pgfqpoint{5.039408in}{1.564197in}}%
\pgfpathlineto{\pgfqpoint{5.037952in}{1.597294in}}%
\pgfpathlineto{\pgfqpoint{5.037970in}{1.625784in}}%
\pgfpathlineto{\pgfqpoint{5.030101in}{1.648413in}}%
\pgfpathlineto{\pgfqpoint{5.026666in}{1.663751in}}%
\pgfpathlineto{\pgfqpoint{5.027261in}{1.684232in}}%
\pgfpathlineto{\pgfqpoint{5.029068in}{1.694197in}}%
\pgfpathlineto{\pgfqpoint{5.024019in}{1.715047in}}%
\pgfpathlineto{\pgfqpoint{5.058393in}{1.714359in}}%
\pgfpathlineto{\pgfqpoint{5.110604in}{1.713911in}}%
\pgfpathlineto{\pgfqpoint{5.110890in}{1.737609in}}%
\pgfpathlineto{\pgfqpoint{5.124180in}{1.735000in}}%
\pgfpathlineto{\pgfqpoint{5.130549in}{1.706104in}}%
\pgfpathlineto{\pgfqpoint{5.135245in}{1.695714in}}%
\pgfpathlineto{\pgfqpoint{5.146922in}{1.695407in}}%
\pgfpathlineto{\pgfqpoint{5.149536in}{1.691882in}}%
\pgfpathlineto{\pgfqpoint{5.165825in}{1.690320in}}%
\pgfpathlineto{\pgfqpoint{5.168562in}{1.683156in}}%
\pgfpathlineto{\pgfqpoint{5.179788in}{1.684767in}}%
\pgfpathlineto{\pgfqpoint{5.188507in}{1.691467in}}%
\pgfpathlineto{\pgfqpoint{5.203542in}{1.691216in}}%
\pgfpathlineto{\pgfqpoint{5.212845in}{1.685828in}}%
\pgfpathlineto{\pgfqpoint{5.213915in}{1.680774in}}%
\pgfpathlineto{\pgfqpoint{5.222767in}{1.679717in}}%
\pgfpathlineto{\pgfqpoint{5.228544in}{1.665932in}}%
\pgfpathlineto{\pgfqpoint{5.232272in}{1.674408in}}%
\pgfpathlineto{\pgfqpoint{5.242441in}{1.674929in}}%
\pgfpathlineto{\pgfqpoint{5.245013in}{1.668328in}}%
\pgfpathlineto{\pgfqpoint{5.256455in}{1.665315in}}%
\pgfpathlineto{\pgfqpoint{5.262765in}{1.655808in}}%
\pgfpathlineto{\pgfqpoint{5.276759in}{1.658760in}}%
\pgfpathlineto{\pgfqpoint{5.292158in}{1.670427in}}%
\pgfpathlineto{\pgfqpoint{5.297772in}{1.660142in}}%
\pgfpathlineto{\pgfqpoint{5.323039in}{1.662957in}}%
\pgfpathlineto{\pgfqpoint{5.333840in}{1.655898in}}%
\pgfpathlineto{\pgfqpoint{5.340126in}{1.658418in}}%
\pgfpathlineto{\pgfqpoint{5.345167in}{1.654444in}}%
\pgfpathlineto{\pgfqpoint{5.330213in}{1.645016in}}%
\pgfpathlineto{\pgfqpoint{5.308866in}{1.636615in}}%
\pgfpathlineto{\pgfqpoint{5.287725in}{1.619824in}}%
\pgfpathlineto{\pgfqpoint{5.269396in}{1.597742in}}%
\pgfpathlineto{\pgfqpoint{5.255521in}{1.584687in}}%
\pgfpathlineto{\pgfqpoint{5.243367in}{1.575713in}}%
\pgfpathlineto{\pgfqpoint{5.235007in}{1.567025in}}%
\pgfpathlineto{\pgfqpoint{5.236020in}{1.530009in}}%
\pgfpathlineto{\pgfqpoint{5.219761in}{1.519600in}}%
\pgfpathlineto{\pgfqpoint{5.211550in}{1.506437in}}%
\pgfpathlineto{\pgfqpoint{5.211004in}{1.498105in}}%
\pgfpathlineto{\pgfqpoint{5.217621in}{1.496270in}}%
\pgfpathlineto{\pgfqpoint{5.221378in}{1.486783in}}%
\pgfpathlineto{\pgfqpoint{5.217197in}{1.480500in}}%
\pgfpathlineto{\pgfqpoint{5.217621in}{1.457675in}}%
\pgfpathlineto{\pgfqpoint{5.215419in}{1.445673in}}%
\pgfpathlineto{\pgfqpoint{5.228976in}{1.435243in}}%
\pgfpathlineto{\pgfqpoint{5.236570in}{1.434051in}}%
\pgfpathlineto{\pgfqpoint{5.242838in}{1.427026in}}%
\pgfpathlineto{\pgfqpoint{5.253501in}{1.422710in}}%
\pgfpathlineto{\pgfqpoint{5.259444in}{1.411992in}}%
\pgfpathlineto{\pgfqpoint{5.270940in}{1.402933in}}%
\pgfpathlineto{\pgfqpoint{5.281505in}{1.397444in}}%
\pgfpathlineto{\pgfqpoint{5.287453in}{1.387663in}}%
\pgfpathlineto{\pgfqpoint{5.289454in}{1.369657in}}%
\pgfpathlineto{\pgfqpoint{5.233388in}{1.367621in}}%
\pgfpathlineto{\pgfqpoint{5.185597in}{1.366622in}}%
\pgfpathlineto{\pgfqpoint{5.142055in}{1.365982in}}%
\pgfpathlineto{\pgfqpoint{5.095992in}{1.366052in}}%
\pgfpathlineto{\pgfqpoint{5.050565in}{1.366544in}}%
\pgfpathclose%
\pgfusepath{fill}%
\end{pgfscope}%
\begin{pgfscope}%
\pgfpathrectangle{\pgfqpoint{3.525000in}{0.100000in}}{\pgfqpoint{2.862985in}{1.832778in}}%
\pgfusepath{clip}%
\pgfsetbuttcap%
\pgfsetmiterjoin%
\definecolor{currentfill}{rgb}{0.723444,0.196158,0.538981}%
\pgfsetfillcolor{currentfill}%
\pgfsetlinewidth{0.000000pt}%
\definecolor{currentstroke}{rgb}{0.000000,0.000000,0.000000}%
\pgfsetstrokecolor{currentstroke}%
\pgfsetstrokeopacity{0.000000}%
\pgfsetdash{}{0pt}%
\pgfpathmoveto{\pgfqpoint{3.777592in}{1.488751in}}%
\pgfpathlineto{\pgfqpoint{3.773152in}{1.496921in}}%
\pgfpathlineto{\pgfqpoint{3.775987in}{1.517893in}}%
\pgfpathlineto{\pgfqpoint{3.781459in}{1.528853in}}%
\pgfpathlineto{\pgfqpoint{3.778724in}{1.543679in}}%
\pgfpathlineto{\pgfqpoint{3.784432in}{1.549914in}}%
\pgfpathlineto{\pgfqpoint{3.794847in}{1.566720in}}%
\pgfpathlineto{\pgfqpoint{3.803677in}{1.576864in}}%
\pgfpathlineto{\pgfqpoint{3.816585in}{1.598959in}}%
\pgfpathlineto{\pgfqpoint{3.826485in}{1.623028in}}%
\pgfpathlineto{\pgfqpoint{3.836995in}{1.645601in}}%
\pgfpathlineto{\pgfqpoint{3.839119in}{1.655029in}}%
\pgfpathlineto{\pgfqpoint{3.853617in}{1.682000in}}%
\pgfpathlineto{\pgfqpoint{3.856409in}{1.693845in}}%
\pgfpathlineto{\pgfqpoint{3.862582in}{1.706284in}}%
\pgfpathlineto{\pgfqpoint{3.864788in}{1.721458in}}%
\pgfpathlineto{\pgfqpoint{3.876573in}{1.728857in}}%
\pgfpathlineto{\pgfqpoint{3.890546in}{1.732587in}}%
\pgfpathlineto{\pgfqpoint{3.897610in}{1.724234in}}%
\pgfpathlineto{\pgfqpoint{3.903718in}{1.724852in}}%
\pgfpathlineto{\pgfqpoint{3.913254in}{1.714897in}}%
\pgfpathlineto{\pgfqpoint{3.914522in}{1.705519in}}%
\pgfpathlineto{\pgfqpoint{3.910905in}{1.687156in}}%
\pgfpathlineto{\pgfqpoint{3.922937in}{1.677787in}}%
\pgfpathlineto{\pgfqpoint{3.930740in}{1.674264in}}%
\pgfpathlineto{\pgfqpoint{3.951880in}{1.677860in}}%
\pgfpathlineto{\pgfqpoint{3.964127in}{1.675404in}}%
\pgfpathlineto{\pgfqpoint{3.972545in}{1.671374in}}%
\pgfpathlineto{\pgfqpoint{3.976979in}{1.663725in}}%
\pgfpathlineto{\pgfqpoint{4.005055in}{1.664746in}}%
\pgfpathlineto{\pgfqpoint{4.009520in}{1.660065in}}%
\pgfpathlineto{\pgfqpoint{4.020053in}{1.659055in}}%
\pgfpathlineto{\pgfqpoint{4.030660in}{1.662045in}}%
\pgfpathlineto{\pgfqpoint{4.062021in}{1.661292in}}%
\pgfpathlineto{\pgfqpoint{4.068794in}{1.658987in}}%
\pgfpathlineto{\pgfqpoint{4.076709in}{1.661643in}}%
\pgfpathlineto{\pgfqpoint{4.164329in}{1.638779in}}%
\pgfpathlineto{\pgfqpoint{4.167899in}{1.626494in}}%
\pgfpathlineto{\pgfqpoint{4.175939in}{1.619838in}}%
\pgfpathlineto{\pgfqpoint{4.176774in}{1.607856in}}%
\pgfpathlineto{\pgfqpoint{4.165015in}{1.595725in}}%
\pgfpathlineto{\pgfqpoint{4.159282in}{1.584259in}}%
\pgfpathlineto{\pgfqpoint{4.152377in}{1.578133in}}%
\pgfpathlineto{\pgfqpoint{4.151752in}{1.572002in}}%
\pgfpathlineto{\pgfqpoint{4.140313in}{1.563715in}}%
\pgfpathlineto{\pgfqpoint{4.132464in}{1.551738in}}%
\pgfpathlineto{\pgfqpoint{4.128070in}{1.548828in}}%
\pgfpathlineto{\pgfqpoint{4.125673in}{1.537205in}}%
\pgfpathlineto{\pgfqpoint{4.137217in}{1.525368in}}%
\pgfpathlineto{\pgfqpoint{4.126069in}{1.505674in}}%
\pgfpathlineto{\pgfqpoint{4.124717in}{1.498158in}}%
\pgfpathlineto{\pgfqpoint{4.099647in}{1.395314in}}%
\pgfpathlineto{\pgfqpoint{4.046916in}{1.408826in}}%
\pgfpathlineto{\pgfqpoint{3.996046in}{1.421943in}}%
\pgfpathlineto{\pgfqpoint{3.965358in}{1.430483in}}%
\pgfpathlineto{\pgfqpoint{3.925927in}{1.441715in}}%
\pgfpathlineto{\pgfqpoint{3.863032in}{1.461511in}}%
\pgfpathlineto{\pgfqpoint{3.794661in}{1.482798in}}%
\pgfpathlineto{\pgfqpoint{3.777592in}{1.488751in}}%
\pgfpathclose%
\pgfusepath{fill}%
\end{pgfscope}%
\begin{pgfscope}%
\pgfpathrectangle{\pgfqpoint{3.525000in}{0.100000in}}{\pgfqpoint{2.862985in}{1.832778in}}%
\pgfusepath{clip}%
\pgfsetbuttcap%
\pgfsetmiterjoin%
\definecolor{currentfill}{rgb}{0.573632,0.060028,0.637349}%
\pgfsetfillcolor{currentfill}%
\pgfsetlinewidth{0.000000pt}%
\definecolor{currentstroke}{rgb}{0.000000,0.000000,0.000000}%
\pgfsetstrokecolor{currentstroke}%
\pgfsetstrokeopacity{0.000000}%
\pgfsetdash{}{0pt}%
\pgfpathmoveto{\pgfqpoint{6.147258in}{1.445782in}}%
\pgfpathlineto{\pgfqpoint{6.142278in}{1.449830in}}%
\pgfpathlineto{\pgfqpoint{6.142972in}{1.462444in}}%
\pgfpathlineto{\pgfqpoint{6.137824in}{1.498298in}}%
\pgfpathlineto{\pgfqpoint{6.141859in}{1.510906in}}%
\pgfpathlineto{\pgfqpoint{6.144741in}{1.533267in}}%
\pgfpathlineto{\pgfqpoint{6.140852in}{1.543846in}}%
\pgfpathlineto{\pgfqpoint{6.154999in}{1.557930in}}%
\pgfpathlineto{\pgfqpoint{6.158559in}{1.569133in}}%
\pgfpathlineto{\pgfqpoint{6.151856in}{1.577975in}}%
\pgfpathlineto{\pgfqpoint{6.154711in}{1.589025in}}%
\pgfpathlineto{\pgfqpoint{6.152922in}{1.595966in}}%
\pgfpathlineto{\pgfqpoint{6.154459in}{1.610832in}}%
\pgfpathlineto{\pgfqpoint{6.157332in}{1.615421in}}%
\pgfpathlineto{\pgfqpoint{6.166192in}{1.618318in}}%
\pgfpathlineto{\pgfqpoint{6.179117in}{1.580468in}}%
\pgfpathlineto{\pgfqpoint{6.203094in}{1.502035in}}%
\pgfpathlineto{\pgfqpoint{6.212168in}{1.495954in}}%
\pgfpathlineto{\pgfqpoint{6.213801in}{1.489012in}}%
\pgfpathlineto{\pgfqpoint{6.218590in}{1.486207in}}%
\pgfpathlineto{\pgfqpoint{6.218226in}{1.473617in}}%
\pgfpathlineto{\pgfqpoint{6.213151in}{1.473415in}}%
\pgfpathlineto{\pgfqpoint{6.202871in}{1.465574in}}%
\pgfpathlineto{\pgfqpoint{6.199906in}{1.457709in}}%
\pgfpathlineto{\pgfqpoint{6.172390in}{1.450901in}}%
\pgfpathlineto{\pgfqpoint{6.147258in}{1.445782in}}%
\pgfpathclose%
\pgfusepath{fill}%
\end{pgfscope}%
\begin{pgfscope}%
\pgfpathrectangle{\pgfqpoint{3.525000in}{0.100000in}}{\pgfqpoint{2.862985in}{1.832778in}}%
\pgfusepath{clip}%
\pgfsetbuttcap%
\pgfsetmiterjoin%
\definecolor{currentfill}{rgb}{0.471457,0.005678,0.659897}%
\pgfsetfillcolor{currentfill}%
\pgfsetlinewidth{0.000000pt}%
\definecolor{currentstroke}{rgb}{0.000000,0.000000,0.000000}%
\pgfsetstrokecolor{currentstroke}%
\pgfsetstrokeopacity{0.000000}%
\pgfsetdash{}{0pt}%
\pgfpathmoveto{\pgfqpoint{5.286865in}{1.172908in}}%
\pgfpathlineto{\pgfqpoint{5.272339in}{1.187293in}}%
\pgfpathlineto{\pgfqpoint{5.225906in}{1.184615in}}%
\pgfpathlineto{\pgfqpoint{5.153489in}{1.182229in}}%
\pgfpathlineto{\pgfqpoint{5.080637in}{1.183365in}}%
\pgfpathlineto{\pgfqpoint{5.075523in}{1.192282in}}%
\pgfpathlineto{\pgfqpoint{5.077611in}{1.201038in}}%
\pgfpathlineto{\pgfqpoint{5.076620in}{1.216037in}}%
\pgfpathlineto{\pgfqpoint{5.072770in}{1.230565in}}%
\pgfpathlineto{\pgfqpoint{5.073217in}{1.238237in}}%
\pgfpathlineto{\pgfqpoint{5.065033in}{1.246905in}}%
\pgfpathlineto{\pgfqpoint{5.066814in}{1.258967in}}%
\pgfpathlineto{\pgfqpoint{5.054260in}{1.282791in}}%
\pgfpathlineto{\pgfqpoint{5.050519in}{1.303018in}}%
\pgfpathlineto{\pgfqpoint{5.041489in}{1.319470in}}%
\pgfpathlineto{\pgfqpoint{5.050061in}{1.339878in}}%
\pgfpathlineto{\pgfqpoint{5.045791in}{1.349375in}}%
\pgfpathlineto{\pgfqpoint{5.050565in}{1.366544in}}%
\pgfpathlineto{\pgfqpoint{5.095992in}{1.366052in}}%
\pgfpathlineto{\pgfqpoint{5.142055in}{1.365982in}}%
\pgfpathlineto{\pgfqpoint{5.185597in}{1.366622in}}%
\pgfpathlineto{\pgfqpoint{5.233388in}{1.367621in}}%
\pgfpathlineto{\pgfqpoint{5.289454in}{1.369657in}}%
\pgfpathlineto{\pgfqpoint{5.290951in}{1.361517in}}%
\pgfpathlineto{\pgfqpoint{5.297158in}{1.354693in}}%
\pgfpathlineto{\pgfqpoint{5.293228in}{1.347582in}}%
\pgfpathlineto{\pgfqpoint{5.294839in}{1.333702in}}%
\pgfpathlineto{\pgfqpoint{5.298686in}{1.322483in}}%
\pgfpathlineto{\pgfqpoint{5.316125in}{1.316544in}}%
\pgfpathlineto{\pgfqpoint{5.319114in}{1.308600in}}%
\pgfpathlineto{\pgfqpoint{5.328684in}{1.299681in}}%
\pgfpathlineto{\pgfqpoint{5.332588in}{1.290454in}}%
\pgfpathlineto{\pgfqpoint{5.342285in}{1.284265in}}%
\pgfpathlineto{\pgfqpoint{5.343802in}{1.276813in}}%
\pgfpathlineto{\pgfqpoint{5.341914in}{1.265532in}}%
\pgfpathlineto{\pgfqpoint{5.336979in}{1.262151in}}%
\pgfpathlineto{\pgfqpoint{5.335490in}{1.251413in}}%
\pgfpathlineto{\pgfqpoint{5.321306in}{1.242884in}}%
\pgfpathlineto{\pgfqpoint{5.302818in}{1.238212in}}%
\pgfpathlineto{\pgfqpoint{5.301124in}{1.227469in}}%
\pgfpathlineto{\pgfqpoint{5.308256in}{1.219794in}}%
\pgfpathlineto{\pgfqpoint{5.308547in}{1.210143in}}%
\pgfpathlineto{\pgfqpoint{5.302793in}{1.202558in}}%
\pgfpathlineto{\pgfqpoint{5.299773in}{1.191288in}}%
\pgfpathlineto{\pgfqpoint{5.289771in}{1.187558in}}%
\pgfpathlineto{\pgfqpoint{5.290408in}{1.174999in}}%
\pgfpathlineto{\pgfqpoint{5.286865in}{1.172908in}}%
\pgfpathclose%
\pgfusepath{fill}%
\end{pgfscope}%
\begin{pgfscope}%
\pgfpathrectangle{\pgfqpoint{3.525000in}{0.100000in}}{\pgfqpoint{2.862985in}{1.832778in}}%
\pgfusepath{clip}%
\pgfsetbuttcap%
\pgfsetmiterjoin%
\definecolor{currentfill}{rgb}{0.640959,0.116492,0.602065}%
\pgfsetfillcolor{currentfill}%
\pgfsetlinewidth{0.000000pt}%
\definecolor{currentstroke}{rgb}{0.000000,0.000000,0.000000}%
\pgfsetstrokecolor{currentstroke}%
\pgfsetstrokeopacity{0.000000}%
\pgfsetdash{}{0pt}%
\pgfpathmoveto{\pgfqpoint{6.188120in}{1.409488in}}%
\pgfpathlineto{\pgfqpoint{6.174362in}{1.407318in}}%
\pgfpathlineto{\pgfqpoint{6.111166in}{1.392957in}}%
\pgfpathlineto{\pgfqpoint{6.110061in}{1.394631in}}%
\pgfpathlineto{\pgfqpoint{6.111015in}{1.438080in}}%
\pgfpathlineto{\pgfqpoint{6.147258in}{1.445782in}}%
\pgfpathlineto{\pgfqpoint{6.172390in}{1.450901in}}%
\pgfpathlineto{\pgfqpoint{6.199906in}{1.457709in}}%
\pgfpathlineto{\pgfqpoint{6.202871in}{1.465574in}}%
\pgfpathlineto{\pgfqpoint{6.213151in}{1.473415in}}%
\pgfpathlineto{\pgfqpoint{6.218226in}{1.473617in}}%
\pgfpathlineto{\pgfqpoint{6.224896in}{1.462179in}}%
\pgfpathlineto{\pgfqpoint{6.218843in}{1.445482in}}%
\pgfpathlineto{\pgfqpoint{6.217955in}{1.435691in}}%
\pgfpathlineto{\pgfqpoint{6.230207in}{1.436611in}}%
\pgfpathlineto{\pgfqpoint{6.235763in}{1.431914in}}%
\pgfpathlineto{\pgfqpoint{6.248224in}{1.412717in}}%
\pgfpathlineto{\pgfqpoint{6.254417in}{1.410382in}}%
\pgfpathlineto{\pgfqpoint{6.264790in}{1.411248in}}%
\pgfpathlineto{\pgfqpoint{6.272009in}{1.417772in}}%
\pgfpathlineto{\pgfqpoint{6.276808in}{1.411950in}}%
\pgfpathlineto{\pgfqpoint{6.246631in}{1.396025in}}%
\pgfpathlineto{\pgfqpoint{6.245688in}{1.407452in}}%
\pgfpathlineto{\pgfqpoint{6.231980in}{1.389693in}}%
\pgfpathlineto{\pgfqpoint{6.227160in}{1.386633in}}%
\pgfpathlineto{\pgfqpoint{6.220438in}{1.396881in}}%
\pgfpathlineto{\pgfqpoint{6.218598in}{1.398285in}}%
\pgfpathlineto{\pgfqpoint{6.212340in}{1.401600in}}%
\pgfpathlineto{\pgfqpoint{6.206882in}{1.415045in}}%
\pgfpathlineto{\pgfqpoint{6.188120in}{1.409488in}}%
\pgfpathclose%
\pgfusepath{fill}%
\end{pgfscope}%
\begin{pgfscope}%
\pgfpathrectangle{\pgfqpoint{3.525000in}{0.100000in}}{\pgfqpoint{2.862985in}{1.832778in}}%
\pgfusepath{clip}%
\pgfsetbuttcap%
\pgfsetmiterjoin%
\definecolor{currentfill}{rgb}{0.399411,0.000859,0.656133}%
\pgfsetfillcolor{currentfill}%
\pgfsetlinewidth{0.000000pt}%
\definecolor{currentstroke}{rgb}{0.000000,0.000000,0.000000}%
\pgfsetstrokecolor{currentstroke}%
\pgfsetstrokeopacity{0.000000}%
\pgfsetdash{}{0pt}%
\pgfpathmoveto{\pgfqpoint{4.778739in}{1.159333in}}%
\pgfpathlineto{\pgfqpoint{4.783777in}{1.221808in}}%
\pgfpathlineto{\pgfqpoint{4.755243in}{1.224124in}}%
\pgfpathlineto{\pgfqpoint{4.689283in}{1.230524in}}%
\pgfpathlineto{\pgfqpoint{4.695419in}{1.292777in}}%
\pgfpathlineto{\pgfqpoint{4.701565in}{1.355506in}}%
\pgfpathlineto{\pgfqpoint{4.726922in}{1.352768in}}%
\pgfpathlineto{\pgfqpoint{4.791682in}{1.346649in}}%
\pgfpathlineto{\pgfqpoint{4.830874in}{1.343580in}}%
\pgfpathlineto{\pgfqpoint{4.878322in}{1.340860in}}%
\pgfpathlineto{\pgfqpoint{4.956420in}{1.337164in}}%
\pgfpathlineto{\pgfqpoint{4.958137in}{1.333663in}}%
\pgfpathlineto{\pgfqpoint{4.980515in}{1.322591in}}%
\pgfpathlineto{\pgfqpoint{4.987377in}{1.328214in}}%
\pgfpathlineto{\pgfqpoint{5.007148in}{1.327701in}}%
\pgfpathlineto{\pgfqpoint{5.018930in}{1.321529in}}%
\pgfpathlineto{\pgfqpoint{5.037612in}{1.314513in}}%
\pgfpathlineto{\pgfqpoint{5.042847in}{1.304277in}}%
\pgfpathlineto{\pgfqpoint{5.050519in}{1.303018in}}%
\pgfpathlineto{\pgfqpoint{5.054260in}{1.282791in}}%
\pgfpathlineto{\pgfqpoint{5.066814in}{1.258967in}}%
\pgfpathlineto{\pgfqpoint{5.065033in}{1.246905in}}%
\pgfpathlineto{\pgfqpoint{5.073217in}{1.238237in}}%
\pgfpathlineto{\pgfqpoint{5.072770in}{1.230565in}}%
\pgfpathlineto{\pgfqpoint{5.076620in}{1.216037in}}%
\pgfpathlineto{\pgfqpoint{5.077611in}{1.201038in}}%
\pgfpathlineto{\pgfqpoint{5.075523in}{1.192282in}}%
\pgfpathlineto{\pgfqpoint{5.080637in}{1.183365in}}%
\pgfpathlineto{\pgfqpoint{5.087651in}{1.167151in}}%
\pgfpathlineto{\pgfqpoint{5.094364in}{1.160551in}}%
\pgfpathlineto{\pgfqpoint{5.102367in}{1.146251in}}%
\pgfpathlineto{\pgfqpoint{5.079686in}{1.146016in}}%
\pgfpathlineto{\pgfqpoint{5.030647in}{1.146774in}}%
\pgfpathlineto{\pgfqpoint{4.976466in}{1.148434in}}%
\pgfpathlineto{\pgfqpoint{4.921966in}{1.150526in}}%
\pgfpathlineto{\pgfqpoint{4.841835in}{1.154903in}}%
\pgfpathlineto{\pgfqpoint{4.778739in}{1.159333in}}%
\pgfpathclose%
\pgfusepath{fill}%
\end{pgfscope}%
\begin{pgfscope}%
\pgfpathrectangle{\pgfqpoint{3.525000in}{0.100000in}}{\pgfqpoint{2.862985in}{1.832778in}}%
\pgfusepath{clip}%
\pgfsetbuttcap%
\pgfsetmiterjoin%
\definecolor{currentfill}{rgb}{0.794549,0.275770,0.473117}%
\pgfsetfillcolor{currentfill}%
\pgfsetlinewidth{0.000000pt}%
\definecolor{currentstroke}{rgb}{0.000000,0.000000,0.000000}%
\pgfsetstrokecolor{currentstroke}%
\pgfsetstrokeopacity{0.000000}%
\pgfsetdash{}{0pt}%
\pgfpathmoveto{\pgfqpoint{5.822084in}{1.347781in}}%
\pgfpathlineto{\pgfqpoint{5.847175in}{1.371705in}}%
\pgfpathlineto{\pgfqpoint{5.850294in}{1.380209in}}%
\pgfpathlineto{\pgfqpoint{5.857646in}{1.387490in}}%
\pgfpathlineto{\pgfqpoint{5.852129in}{1.398090in}}%
\pgfpathlineto{\pgfqpoint{5.845212in}{1.404290in}}%
\pgfpathlineto{\pgfqpoint{5.843213in}{1.415277in}}%
\pgfpathlineto{\pgfqpoint{5.868959in}{1.426558in}}%
\pgfpathlineto{\pgfqpoint{5.890274in}{1.430126in}}%
\pgfpathlineto{\pgfqpoint{5.901728in}{1.430365in}}%
\pgfpathlineto{\pgfqpoint{5.910453in}{1.426054in}}%
\pgfpathlineto{\pgfqpoint{5.918958in}{1.429903in}}%
\pgfpathlineto{\pgfqpoint{5.939702in}{1.434217in}}%
\pgfpathlineto{\pgfqpoint{5.946870in}{1.439785in}}%
\pgfpathlineto{\pgfqpoint{5.957498in}{1.452111in}}%
\pgfpathlineto{\pgfqpoint{5.967167in}{1.457558in}}%
\pgfpathlineto{\pgfqpoint{5.967849in}{1.462782in}}%
\pgfpathlineto{\pgfqpoint{5.962774in}{1.474701in}}%
\pgfpathlineto{\pgfqpoint{5.966443in}{1.481714in}}%
\pgfpathlineto{\pgfqpoint{5.961503in}{1.489256in}}%
\pgfpathlineto{\pgfqpoint{5.953939in}{1.489782in}}%
\pgfpathlineto{\pgfqpoint{5.972870in}{1.512563in}}%
\pgfpathlineto{\pgfqpoint{5.975119in}{1.521245in}}%
\pgfpathlineto{\pgfqpoint{5.990018in}{1.543386in}}%
\pgfpathlineto{\pgfqpoint{6.003850in}{1.555368in}}%
\pgfpathlineto{\pgfqpoint{6.013347in}{1.560371in}}%
\pgfpathlineto{\pgfqpoint{6.044425in}{1.567413in}}%
\pgfpathlineto{\pgfqpoint{6.073552in}{1.575365in}}%
\pgfpathlineto{\pgfqpoint{6.077410in}{1.562882in}}%
\pgfpathlineto{\pgfqpoint{6.077683in}{1.551308in}}%
\pgfpathlineto{\pgfqpoint{6.084327in}{1.540801in}}%
\pgfpathlineto{\pgfqpoint{6.086375in}{1.530104in}}%
\pgfpathlineto{\pgfqpoint{6.083822in}{1.515590in}}%
\pgfpathlineto{\pgfqpoint{6.091598in}{1.496013in}}%
\pgfpathlineto{\pgfqpoint{6.100463in}{1.485127in}}%
\pgfpathlineto{\pgfqpoint{6.107320in}{1.449744in}}%
\pgfpathlineto{\pgfqpoint{6.111015in}{1.438080in}}%
\pgfpathlineto{\pgfqpoint{6.110061in}{1.394631in}}%
\pgfpathlineto{\pgfqpoint{6.111166in}{1.392957in}}%
\pgfpathlineto{\pgfqpoint{6.119238in}{1.346227in}}%
\pgfpathlineto{\pgfqpoint{6.123765in}{1.341969in}}%
\pgfpathlineto{\pgfqpoint{6.114028in}{1.332508in}}%
\pgfpathlineto{\pgfqpoint{6.118832in}{1.327078in}}%
\pgfpathlineto{\pgfqpoint{6.114609in}{1.318865in}}%
\pgfpathlineto{\pgfqpoint{6.114645in}{1.315368in}}%
\pgfpathlineto{\pgfqpoint{6.109366in}{1.312213in}}%
\pgfpathlineto{\pgfqpoint{6.106795in}{1.305239in}}%
\pgfpathlineto{\pgfqpoint{6.107610in}{1.324417in}}%
\pgfpathlineto{\pgfqpoint{6.091246in}{1.328639in}}%
\pgfpathlineto{\pgfqpoint{6.065655in}{1.337368in}}%
\pgfpathlineto{\pgfqpoint{6.062035in}{1.341668in}}%
\pgfpathlineto{\pgfqpoint{6.051342in}{1.342698in}}%
\pgfpathlineto{\pgfqpoint{6.045190in}{1.349098in}}%
\pgfpathlineto{\pgfqpoint{6.041871in}{1.361844in}}%
\pgfpathlineto{\pgfqpoint{6.033151in}{1.363429in}}%
\pgfpathlineto{\pgfqpoint{6.027319in}{1.370116in}}%
\pgfpathlineto{\pgfqpoint{5.953076in}{1.355148in}}%
\pgfpathlineto{\pgfqpoint{5.917594in}{1.347810in}}%
\pgfpathlineto{\pgfqpoint{5.863721in}{1.337977in}}%
\pgfpathlineto{\pgfqpoint{5.824926in}{1.331432in}}%
\pgfpathlineto{\pgfqpoint{5.822084in}{1.347781in}}%
\pgfpathclose%
\pgfusepath{fill}%
\end{pgfscope}%
\begin{pgfscope}%
\pgfpathrectangle{\pgfqpoint{3.525000in}{0.100000in}}{\pgfqpoint{2.862985in}{1.832778in}}%
\pgfusepath{clip}%
\pgfsetbuttcap%
\pgfsetmiterjoin%
\definecolor{currentfill}{rgb}{0.794549,0.275770,0.473117}%
\pgfsetfillcolor{currentfill}%
\pgfsetlinewidth{0.000000pt}%
\definecolor{currentstroke}{rgb}{0.000000,0.000000,0.000000}%
\pgfsetstrokecolor{currentstroke}%
\pgfsetstrokeopacity{0.000000}%
\pgfsetdash{}{0pt}%
\pgfpathmoveto{\pgfqpoint{6.120139in}{1.301344in}}%
\pgfpathlineto{\pgfqpoint{6.108657in}{1.297767in}}%
\pgfpathlineto{\pgfqpoint{6.106726in}{1.301055in}}%
\pgfpathlineto{\pgfqpoint{6.110411in}{1.312094in}}%
\pgfpathlineto{\pgfqpoint{6.117237in}{1.313181in}}%
\pgfpathlineto{\pgfqpoint{6.116626in}{1.316656in}}%
\pgfpathlineto{\pgfqpoint{6.122754in}{1.321872in}}%
\pgfpathlineto{\pgfqpoint{6.140472in}{1.326028in}}%
\pgfpathlineto{\pgfqpoint{6.148360in}{1.332319in}}%
\pgfpathlineto{\pgfqpoint{6.166079in}{1.337561in}}%
\pgfpathlineto{\pgfqpoint{6.174161in}{1.335643in}}%
\pgfpathlineto{\pgfqpoint{6.180959in}{1.344095in}}%
\pgfpathlineto{\pgfqpoint{6.191234in}{1.345211in}}%
\pgfpathlineto{\pgfqpoint{6.173704in}{1.328725in}}%
\pgfpathlineto{\pgfqpoint{6.120139in}{1.301344in}}%
\pgfpathclose%
\pgfusepath{fill}%
\end{pgfscope}%
\begin{pgfscope}%
\pgfpathrectangle{\pgfqpoint{3.525000in}{0.100000in}}{\pgfqpoint{2.862985in}{1.832778in}}%
\pgfusepath{clip}%
\pgfsetbuttcap%
\pgfsetmiterjoin%
\definecolor{currentfill}{rgb}{0.756304,0.231555,0.509468}%
\pgfsetfillcolor{currentfill}%
\pgfsetlinewidth{0.000000pt}%
\definecolor{currentstroke}{rgb}{0.000000,0.000000,0.000000}%
\pgfsetstrokecolor{currentstroke}%
\pgfsetstrokeopacity{0.000000}%
\pgfsetdash{}{0pt}%
\pgfpathmoveto{\pgfqpoint{5.862178in}{1.192485in}}%
\pgfpathlineto{\pgfqpoint{5.812526in}{1.184275in}}%
\pgfpathlineto{\pgfqpoint{5.803519in}{1.241018in}}%
\pgfpathlineto{\pgfqpoint{5.790146in}{1.324661in}}%
\pgfpathlineto{\pgfqpoint{5.822084in}{1.347781in}}%
\pgfpathlineto{\pgfqpoint{5.824926in}{1.331432in}}%
\pgfpathlineto{\pgfqpoint{5.863721in}{1.337977in}}%
\pgfpathlineto{\pgfqpoint{5.917594in}{1.347810in}}%
\pgfpathlineto{\pgfqpoint{5.953076in}{1.355148in}}%
\pgfpathlineto{\pgfqpoint{6.027319in}{1.370116in}}%
\pgfpathlineto{\pgfqpoint{6.033151in}{1.363429in}}%
\pgfpathlineto{\pgfqpoint{6.041871in}{1.361844in}}%
\pgfpathlineto{\pgfqpoint{6.045190in}{1.349098in}}%
\pgfpathlineto{\pgfqpoint{6.051342in}{1.342698in}}%
\pgfpathlineto{\pgfqpoint{6.062035in}{1.341668in}}%
\pgfpathlineto{\pgfqpoint{6.065655in}{1.337368in}}%
\pgfpathlineto{\pgfqpoint{6.061981in}{1.334051in}}%
\pgfpathlineto{\pgfqpoint{6.058670in}{1.322314in}}%
\pgfpathlineto{\pgfqpoint{6.050913in}{1.309125in}}%
\pgfpathlineto{\pgfqpoint{6.056120in}{1.303390in}}%
\pgfpathlineto{\pgfqpoint{6.051859in}{1.297245in}}%
\pgfpathlineto{\pgfqpoint{6.054263in}{1.283769in}}%
\pgfpathlineto{\pgfqpoint{6.061936in}{1.276695in}}%
\pgfpathlineto{\pgfqpoint{6.080144in}{1.265202in}}%
\pgfpathlineto{\pgfqpoint{6.065481in}{1.248986in}}%
\pgfpathlineto{\pgfqpoint{6.065275in}{1.242831in}}%
\pgfpathlineto{\pgfqpoint{6.053338in}{1.234894in}}%
\pgfpathlineto{\pgfqpoint{6.040138in}{1.233404in}}%
\pgfpathlineto{\pgfqpoint{6.036873in}{1.226505in}}%
\pgfpathlineto{\pgfqpoint{6.000158in}{1.218561in}}%
\pgfpathlineto{\pgfqpoint{5.957316in}{1.209933in}}%
\pgfpathlineto{\pgfqpoint{5.927868in}{1.204659in}}%
\pgfpathlineto{\pgfqpoint{5.862178in}{1.192485in}}%
\pgfpathclose%
\pgfusepath{fill}%
\end{pgfscope}%
\begin{pgfscope}%
\pgfpathrectangle{\pgfqpoint{3.525000in}{0.100000in}}{\pgfqpoint{2.862985in}{1.832778in}}%
\pgfusepath{clip}%
\pgfsetbuttcap%
\pgfsetmiterjoin%
\definecolor{currentfill}{rgb}{0.625987,0.103312,0.611305}%
\pgfsetfillcolor{currentfill}%
\pgfsetlinewidth{0.000000pt}%
\definecolor{currentstroke}{rgb}{0.000000,0.000000,0.000000}%
\pgfsetstrokecolor{currentstroke}%
\pgfsetstrokeopacity{0.000000}%
\pgfsetdash{}{0pt}%
\pgfpathmoveto{\pgfqpoint{6.111166in}{1.392957in}}%
\pgfpathlineto{\pgfqpoint{6.174362in}{1.407318in}}%
\pgfpathlineto{\pgfqpoint{6.188120in}{1.409488in}}%
\pgfpathlineto{\pgfqpoint{6.197226in}{1.373686in}}%
\pgfpathlineto{\pgfqpoint{6.195785in}{1.367275in}}%
\pgfpathlineto{\pgfqpoint{6.166533in}{1.355951in}}%
\pgfpathlineto{\pgfqpoint{6.149071in}{1.351992in}}%
\pgfpathlineto{\pgfqpoint{6.141651in}{1.343114in}}%
\pgfpathlineto{\pgfqpoint{6.118832in}{1.327078in}}%
\pgfpathlineto{\pgfqpoint{6.114028in}{1.332508in}}%
\pgfpathlineto{\pgfqpoint{6.123765in}{1.341969in}}%
\pgfpathlineto{\pgfqpoint{6.119238in}{1.346227in}}%
\pgfpathlineto{\pgfqpoint{6.111166in}{1.392957in}}%
\pgfpathclose%
\pgfusepath{fill}%
\end{pgfscope}%
\begin{pgfscope}%
\pgfpathrectangle{\pgfqpoint{3.525000in}{0.100000in}}{\pgfqpoint{2.862985in}{1.832778in}}%
\pgfusepath{clip}%
\pgfsetbuttcap%
\pgfsetmiterjoin%
\definecolor{currentfill}{rgb}{0.688318,0.160709,0.568103}%
\pgfsetfillcolor{currentfill}%
\pgfsetlinewidth{0.000000pt}%
\definecolor{currentstroke}{rgb}{0.000000,0.000000,0.000000}%
\pgfsetstrokecolor{currentstroke}%
\pgfsetstrokeopacity{0.000000}%
\pgfsetdash{}{0pt}%
\pgfpathmoveto{\pgfqpoint{6.188120in}{1.409488in}}%
\pgfpathlineto{\pgfqpoint{6.206882in}{1.415045in}}%
\pgfpathlineto{\pgfqpoint{6.212340in}{1.401600in}}%
\pgfpathlineto{\pgfqpoint{6.218598in}{1.398285in}}%
\pgfpathlineto{\pgfqpoint{6.210875in}{1.392622in}}%
\pgfpathlineto{\pgfqpoint{6.211849in}{1.375988in}}%
\pgfpathlineto{\pgfqpoint{6.195785in}{1.367275in}}%
\pgfpathlineto{\pgfqpoint{6.197226in}{1.373686in}}%
\pgfpathlineto{\pgfqpoint{6.188120in}{1.409488in}}%
\pgfpathclose%
\pgfusepath{fill}%
\end{pgfscope}%
\begin{pgfscope}%
\pgfpathrectangle{\pgfqpoint{3.525000in}{0.100000in}}{\pgfqpoint{2.862985in}{1.832778in}}%
\pgfusepath{clip}%
\pgfsetbuttcap%
\pgfsetmiterjoin%
\definecolor{currentfill}{rgb}{0.650746,0.125309,0.595617}%
\pgfsetfillcolor{currentfill}%
\pgfsetlinewidth{0.000000pt}%
\definecolor{currentstroke}{rgb}{0.000000,0.000000,0.000000}%
\pgfsetstrokecolor{currentstroke}%
\pgfsetstrokeopacity{0.000000}%
\pgfsetdash{}{0pt}%
\pgfpathmoveto{\pgfqpoint{6.051223in}{1.229031in}}%
\pgfpathlineto{\pgfqpoint{6.053338in}{1.234894in}}%
\pgfpathlineto{\pgfqpoint{6.065275in}{1.242831in}}%
\pgfpathlineto{\pgfqpoint{6.065481in}{1.248986in}}%
\pgfpathlineto{\pgfqpoint{6.080144in}{1.265202in}}%
\pgfpathlineto{\pgfqpoint{6.061936in}{1.276695in}}%
\pgfpathlineto{\pgfqpoint{6.054263in}{1.283769in}}%
\pgfpathlineto{\pgfqpoint{6.051859in}{1.297245in}}%
\pgfpathlineto{\pgfqpoint{6.056120in}{1.303390in}}%
\pgfpathlineto{\pgfqpoint{6.050913in}{1.309125in}}%
\pgfpathlineto{\pgfqpoint{6.058670in}{1.322314in}}%
\pgfpathlineto{\pgfqpoint{6.061981in}{1.334051in}}%
\pgfpathlineto{\pgfqpoint{6.065655in}{1.337368in}}%
\pgfpathlineto{\pgfqpoint{6.091246in}{1.328639in}}%
\pgfpathlineto{\pgfqpoint{6.107610in}{1.324417in}}%
\pgfpathlineto{\pgfqpoint{6.106795in}{1.305239in}}%
\pgfpathlineto{\pgfqpoint{6.096854in}{1.290696in}}%
\pgfpathlineto{\pgfqpoint{6.105045in}{1.288556in}}%
\pgfpathlineto{\pgfqpoint{6.113550in}{1.282315in}}%
\pgfpathlineto{\pgfqpoint{6.114051in}{1.265245in}}%
\pgfpathlineto{\pgfqpoint{6.111477in}{1.253161in}}%
\pgfpathlineto{\pgfqpoint{6.113190in}{1.243235in}}%
\pgfpathlineto{\pgfqpoint{6.098407in}{1.209740in}}%
\pgfpathlineto{\pgfqpoint{6.093094in}{1.194100in}}%
\pgfpathlineto{\pgfqpoint{6.085676in}{1.201703in}}%
\pgfpathlineto{\pgfqpoint{6.075838in}{1.200407in}}%
\pgfpathlineto{\pgfqpoint{6.051218in}{1.214632in}}%
\pgfpathlineto{\pgfqpoint{6.048699in}{1.222249in}}%
\pgfpathlineto{\pgfqpoint{6.051223in}{1.229031in}}%
\pgfpathclose%
\pgfusepath{fill}%
\end{pgfscope}%
\begin{pgfscope}%
\pgfpathrectangle{\pgfqpoint{3.525000in}{0.100000in}}{\pgfqpoint{2.862985in}{1.832778in}}%
\pgfusepath{clip}%
\pgfsetbuttcap%
\pgfsetmiterjoin%
\definecolor{currentfill}{rgb}{0.669845,0.142992,0.582154}%
\pgfsetfillcolor{currentfill}%
\pgfsetlinewidth{0.000000pt}%
\definecolor{currentstroke}{rgb}{0.000000,0.000000,0.000000}%
\pgfsetstrokecolor{currentstroke}%
\pgfsetstrokeopacity{0.000000}%
\pgfsetdash{}{0pt}%
\pgfpathmoveto{\pgfqpoint{5.462013in}{1.021868in}}%
\pgfpathlineto{\pgfqpoint{5.460862in}{1.037407in}}%
\pgfpathlineto{\pgfqpoint{5.465336in}{1.045116in}}%
\pgfpathlineto{\pgfqpoint{5.462387in}{1.048893in}}%
\pgfpathlineto{\pgfqpoint{5.473140in}{1.061261in}}%
\pgfpathlineto{\pgfqpoint{5.472524in}{1.064241in}}%
\pgfpathlineto{\pgfqpoint{5.482971in}{1.085818in}}%
\pgfpathlineto{\pgfqpoint{5.480859in}{1.096224in}}%
\pgfpathlineto{\pgfqpoint{5.473283in}{1.107118in}}%
\pgfpathlineto{\pgfqpoint{5.478542in}{1.120373in}}%
\pgfpathlineto{\pgfqpoint{5.471726in}{1.207597in}}%
\pgfpathlineto{\pgfqpoint{5.466818in}{1.268787in}}%
\pgfpathlineto{\pgfqpoint{5.473598in}{1.263717in}}%
\pgfpathlineto{\pgfqpoint{5.481163in}{1.263854in}}%
\pgfpathlineto{\pgfqpoint{5.499029in}{1.274203in}}%
\pgfpathlineto{\pgfqpoint{5.553831in}{1.279317in}}%
\pgfpathlineto{\pgfqpoint{5.594393in}{1.283592in}}%
\pgfpathlineto{\pgfqpoint{5.594752in}{1.279623in}}%
\pgfpathlineto{\pgfqpoint{5.604011in}{1.195770in}}%
\pgfpathlineto{\pgfqpoint{5.611981in}{1.117748in}}%
\pgfpathlineto{\pgfqpoint{5.608545in}{1.114085in}}%
\pgfpathlineto{\pgfqpoint{5.613822in}{1.104993in}}%
\pgfpathlineto{\pgfqpoint{5.613795in}{1.098441in}}%
\pgfpathlineto{\pgfqpoint{5.606265in}{1.096795in}}%
\pgfpathlineto{\pgfqpoint{5.597839in}{1.090478in}}%
\pgfpathlineto{\pgfqpoint{5.592134in}{1.092965in}}%
\pgfpathlineto{\pgfqpoint{5.583595in}{1.088922in}}%
\pgfpathlineto{\pgfqpoint{5.586237in}{1.080801in}}%
\pgfpathlineto{\pgfqpoint{5.577460in}{1.072642in}}%
\pgfpathlineto{\pgfqpoint{5.575037in}{1.063205in}}%
\pgfpathlineto{\pgfqpoint{5.569002in}{1.061656in}}%
\pgfpathlineto{\pgfqpoint{5.564498in}{1.054509in}}%
\pgfpathlineto{\pgfqpoint{5.565088in}{1.047311in}}%
\pgfpathlineto{\pgfqpoint{5.559790in}{1.042250in}}%
\pgfpathlineto{\pgfqpoint{5.551817in}{1.043031in}}%
\pgfpathlineto{\pgfqpoint{5.545757in}{1.050797in}}%
\pgfpathlineto{\pgfqpoint{5.535489in}{1.043318in}}%
\pgfpathlineto{\pgfqpoint{5.536209in}{1.036783in}}%
\pgfpathlineto{\pgfqpoint{5.524789in}{1.032965in}}%
\pgfpathlineto{\pgfqpoint{5.521919in}{1.037771in}}%
\pgfpathlineto{\pgfqpoint{5.510801in}{1.032321in}}%
\pgfpathlineto{\pgfqpoint{5.506739in}{1.024520in}}%
\pgfpathlineto{\pgfqpoint{5.493349in}{1.032520in}}%
\pgfpathlineto{\pgfqpoint{5.472127in}{1.027220in}}%
\pgfpathlineto{\pgfqpoint{5.462013in}{1.021868in}}%
\pgfpathclose%
\pgfusepath{fill}%
\end{pgfscope}%
\begin{pgfscope}%
\pgfpathrectangle{\pgfqpoint{3.525000in}{0.100000in}}{\pgfqpoint{2.862985in}{1.832778in}}%
\pgfusepath{clip}%
\pgfsetbuttcap%
\pgfsetmiterjoin%
\definecolor{currentfill}{rgb}{0.783383,0.262500,0.483918}%
\pgfsetfillcolor{currentfill}%
\pgfsetlinewidth{0.000000pt}%
\definecolor{currentstroke}{rgb}{0.000000,0.000000,0.000000}%
\pgfsetstrokecolor{currentstroke}%
\pgfsetstrokeopacity{0.000000}%
\pgfsetdash{}{0pt}%
\pgfpathmoveto{\pgfqpoint{3.965358in}{1.430483in}}%
\pgfpathlineto{\pgfqpoint{3.996046in}{1.421943in}}%
\pgfpathlineto{\pgfqpoint{4.046916in}{1.408826in}}%
\pgfpathlineto{\pgfqpoint{4.099647in}{1.395314in}}%
\pgfpathlineto{\pgfqpoint{4.148250in}{1.383593in}}%
\pgfpathlineto{\pgfqpoint{4.190246in}{1.374085in}}%
\pgfpathlineto{\pgfqpoint{4.235177in}{1.364205in}}%
\pgfpathlineto{\pgfqpoint{4.222195in}{1.302942in}}%
\pgfpathlineto{\pgfqpoint{4.210630in}{1.248537in}}%
\pgfpathlineto{\pgfqpoint{4.191658in}{1.160754in}}%
\pgfpathlineto{\pgfqpoint{4.177414in}{1.094498in}}%
\pgfpathlineto{\pgfqpoint{4.169719in}{1.057523in}}%
\pgfpathlineto{\pgfqpoint{4.159852in}{1.009534in}}%
\pgfpathlineto{\pgfqpoint{4.153025in}{0.999798in}}%
\pgfpathlineto{\pgfqpoint{4.147531in}{0.999473in}}%
\pgfpathlineto{\pgfqpoint{4.143612in}{1.007971in}}%
\pgfpathlineto{\pgfqpoint{4.134613in}{1.011057in}}%
\pgfpathlineto{\pgfqpoint{4.124916in}{1.009976in}}%
\pgfpathlineto{\pgfqpoint{4.122164in}{1.003024in}}%
\pgfpathlineto{\pgfqpoint{4.121877in}{0.978886in}}%
\pgfpathlineto{\pgfqpoint{4.118973in}{0.973374in}}%
\pgfpathlineto{\pgfqpoint{4.120633in}{0.953995in}}%
\pgfpathlineto{\pgfqpoint{4.114569in}{0.941085in}}%
\pgfpathlineto{\pgfqpoint{4.065397in}{1.016802in}}%
\pgfpathlineto{\pgfqpoint{4.016971in}{1.090379in}}%
\pgfpathlineto{\pgfqpoint{3.991763in}{1.128960in}}%
\pgfpathlineto{\pgfqpoint{3.970735in}{1.162262in}}%
\pgfpathlineto{\pgfqpoint{3.944238in}{1.203456in}}%
\pgfpathlineto{\pgfqpoint{3.914244in}{1.249665in}}%
\pgfpathlineto{\pgfqpoint{3.926576in}{1.293524in}}%
\pgfpathlineto{\pgfqpoint{3.951391in}{1.381486in}}%
\pgfpathlineto{\pgfqpoint{3.965358in}{1.430483in}}%
\pgfpathclose%
\pgfusepath{fill}%
\end{pgfscope}%
\begin{pgfscope}%
\pgfpathrectangle{\pgfqpoint{3.525000in}{0.100000in}}{\pgfqpoint{2.862985in}{1.832778in}}%
\pgfusepath{clip}%
\pgfsetbuttcap%
\pgfsetmiterjoin%
\definecolor{currentfill}{rgb}{0.423689,0.000646,0.658956}%
\pgfsetfillcolor{currentfill}%
\pgfsetlinewidth{0.000000pt}%
\definecolor{currentstroke}{rgb}{0.000000,0.000000,0.000000}%
\pgfsetstrokecolor{currentstroke}%
\pgfsetstrokeopacity{0.000000}%
\pgfsetdash{}{0pt}%
\pgfpathmoveto{\pgfqpoint{4.235177in}{1.364205in}}%
\pgfpathlineto{\pgfqpoint{4.283212in}{1.354650in}}%
\pgfpathlineto{\pgfqpoint{4.372141in}{1.337483in}}%
\pgfpathlineto{\pgfqpoint{4.361009in}{1.275734in}}%
\pgfpathlineto{\pgfqpoint{4.409967in}{1.267352in}}%
\pgfpathlineto{\pgfqpoint{4.454571in}{1.260214in}}%
\pgfpathlineto{\pgfqpoint{4.446820in}{1.211396in}}%
\pgfpathlineto{\pgfqpoint{4.438608in}{1.158753in}}%
\pgfpathlineto{\pgfqpoint{4.427613in}{1.089632in}}%
\pgfpathlineto{\pgfqpoint{4.427328in}{1.083837in}}%
\pgfpathlineto{\pgfqpoint{4.415915in}{1.012202in}}%
\pgfpathlineto{\pgfqpoint{4.368932in}{1.019484in}}%
\pgfpathlineto{\pgfqpoint{4.344991in}{1.024292in}}%
\pgfpathlineto{\pgfqpoint{4.258455in}{1.039504in}}%
\pgfpathlineto{\pgfqpoint{4.225868in}{1.045930in}}%
\pgfpathlineto{\pgfqpoint{4.169719in}{1.057523in}}%
\pgfpathlineto{\pgfqpoint{4.177414in}{1.094498in}}%
\pgfpathlineto{\pgfqpoint{4.191658in}{1.160754in}}%
\pgfpathlineto{\pgfqpoint{4.210630in}{1.248537in}}%
\pgfpathlineto{\pgfqpoint{4.222195in}{1.302942in}}%
\pgfpathlineto{\pgfqpoint{4.235177in}{1.364205in}}%
\pgfpathclose%
\pgfusepath{fill}%
\end{pgfscope}%
\begin{pgfscope}%
\pgfpathrectangle{\pgfqpoint{3.525000in}{0.100000in}}{\pgfqpoint{2.862985in}{1.832778in}}%
\pgfusepath{clip}%
\pgfsetbuttcap%
\pgfsetmiterjoin%
\definecolor{currentfill}{rgb}{0.736019,0.209439,0.527908}%
\pgfsetfillcolor{currentfill}%
\pgfsetlinewidth{0.000000pt}%
\definecolor{currentstroke}{rgb}{0.000000,0.000000,0.000000}%
\pgfsetstrokecolor{currentstroke}%
\pgfsetstrokeopacity{0.000000}%
\pgfsetdash{}{0pt}%
\pgfpathmoveto{\pgfqpoint{3.777592in}{1.488751in}}%
\pgfpathlineto{\pgfqpoint{3.794661in}{1.482798in}}%
\pgfpathlineto{\pgfqpoint{3.863032in}{1.461511in}}%
\pgfpathlineto{\pgfqpoint{3.925927in}{1.441715in}}%
\pgfpathlineto{\pgfqpoint{3.965358in}{1.430483in}}%
\pgfpathlineto{\pgfqpoint{3.951391in}{1.381486in}}%
\pgfpathlineto{\pgfqpoint{3.926576in}{1.293524in}}%
\pgfpathlineto{\pgfqpoint{3.914244in}{1.249665in}}%
\pgfpathlineto{\pgfqpoint{3.944238in}{1.203456in}}%
\pgfpathlineto{\pgfqpoint{3.970735in}{1.162262in}}%
\pgfpathlineto{\pgfqpoint{3.991763in}{1.128960in}}%
\pgfpathlineto{\pgfqpoint{4.016971in}{1.090379in}}%
\pgfpathlineto{\pgfqpoint{4.065397in}{1.016802in}}%
\pgfpathlineto{\pgfqpoint{4.114569in}{0.941085in}}%
\pgfpathlineto{\pgfqpoint{4.112595in}{0.933576in}}%
\pgfpathlineto{\pgfqpoint{4.118526in}{0.921617in}}%
\pgfpathlineto{\pgfqpoint{4.119699in}{0.905259in}}%
\pgfpathlineto{\pgfqpoint{4.128313in}{0.897324in}}%
\pgfpathlineto{\pgfqpoint{4.128654in}{0.890930in}}%
\pgfpathlineto{\pgfqpoint{4.113230in}{0.883673in}}%
\pgfpathlineto{\pgfqpoint{4.105882in}{0.876407in}}%
\pgfpathlineto{\pgfqpoint{4.103591in}{0.860364in}}%
\pgfpathlineto{\pgfqpoint{4.099871in}{0.851614in}}%
\pgfpathlineto{\pgfqpoint{4.092035in}{0.844239in}}%
\pgfpathlineto{\pgfqpoint{4.084240in}{0.825061in}}%
\pgfpathlineto{\pgfqpoint{4.095200in}{0.815071in}}%
\pgfpathlineto{\pgfqpoint{4.093785in}{0.806836in}}%
\pgfpathlineto{\pgfqpoint{4.084798in}{0.801105in}}%
\pgfpathlineto{\pgfqpoint{4.078603in}{0.802117in}}%
\pgfpathlineto{\pgfqpoint{4.005470in}{0.812266in}}%
\pgfpathlineto{\pgfqpoint{3.951454in}{0.819928in}}%
\pgfpathlineto{\pgfqpoint{3.953816in}{0.828649in}}%
\pgfpathlineto{\pgfqpoint{3.950300in}{0.843128in}}%
\pgfpathlineto{\pgfqpoint{3.949945in}{0.857733in}}%
\pgfpathlineto{\pgfqpoint{3.947647in}{0.866279in}}%
\pgfpathlineto{\pgfqpoint{3.940568in}{0.878498in}}%
\pgfpathlineto{\pgfqpoint{3.920206in}{0.906691in}}%
\pgfpathlineto{\pgfqpoint{3.913537in}{0.910145in}}%
\pgfpathlineto{\pgfqpoint{3.904944in}{0.910090in}}%
\pgfpathlineto{\pgfqpoint{3.906908in}{0.918984in}}%
\pgfpathlineto{\pgfqpoint{3.902839in}{0.930106in}}%
\pgfpathlineto{\pgfqpoint{3.882853in}{0.935633in}}%
\pgfpathlineto{\pgfqpoint{3.870671in}{0.945869in}}%
\pgfpathlineto{\pgfqpoint{3.869661in}{0.952132in}}%
\pgfpathlineto{\pgfqpoint{3.855633in}{0.967637in}}%
\pgfpathlineto{\pgfqpoint{3.842280in}{0.970604in}}%
\pgfpathlineto{\pgfqpoint{3.829908in}{0.978484in}}%
\pgfpathlineto{\pgfqpoint{3.813637in}{0.981210in}}%
\pgfpathlineto{\pgfqpoint{3.806685in}{0.991719in}}%
\pgfpathlineto{\pgfqpoint{3.813287in}{1.008354in}}%
\pgfpathlineto{\pgfqpoint{3.811268in}{1.012095in}}%
\pgfpathlineto{\pgfqpoint{3.816766in}{1.025962in}}%
\pgfpathlineto{\pgfqpoint{3.806986in}{1.033342in}}%
\pgfpathlineto{\pgfqpoint{3.810159in}{1.046721in}}%
\pgfpathlineto{\pgfqpoint{3.804935in}{1.050141in}}%
\pgfpathlineto{\pgfqpoint{3.800424in}{1.062787in}}%
\pgfpathlineto{\pgfqpoint{3.794981in}{1.066644in}}%
\pgfpathlineto{\pgfqpoint{3.794572in}{1.075778in}}%
\pgfpathlineto{\pgfqpoint{3.790307in}{1.082219in}}%
\pgfpathlineto{\pgfqpoint{3.783881in}{1.103911in}}%
\pgfpathlineto{\pgfqpoint{3.776850in}{1.114301in}}%
\pgfpathlineto{\pgfqpoint{3.778378in}{1.131939in}}%
\pgfpathlineto{\pgfqpoint{3.786647in}{1.133714in}}%
\pgfpathlineto{\pgfqpoint{3.792031in}{1.143280in}}%
\pgfpathlineto{\pgfqpoint{3.788778in}{1.153656in}}%
\pgfpathlineto{\pgfqpoint{3.779984in}{1.155388in}}%
\pgfpathlineto{\pgfqpoint{3.775596in}{1.160241in}}%
\pgfpathlineto{\pgfqpoint{3.768491in}{1.178092in}}%
\pgfpathlineto{\pgfqpoint{3.771812in}{1.184503in}}%
\pgfpathlineto{\pgfqpoint{3.769457in}{1.196440in}}%
\pgfpathlineto{\pgfqpoint{3.774672in}{1.211900in}}%
\pgfpathlineto{\pgfqpoint{3.780782in}{1.206207in}}%
\pgfpathlineto{\pgfqpoint{3.778009in}{1.199485in}}%
\pgfpathlineto{\pgfqpoint{3.788597in}{1.188830in}}%
\pgfpathlineto{\pgfqpoint{3.787921in}{1.204652in}}%
\pgfpathlineto{\pgfqpoint{3.783384in}{1.208896in}}%
\pgfpathlineto{\pgfqpoint{3.788570in}{1.222808in}}%
\pgfpathlineto{\pgfqpoint{3.808185in}{1.220593in}}%
\pgfpathlineto{\pgfqpoint{3.805545in}{1.225766in}}%
\pgfpathlineto{\pgfqpoint{3.792599in}{1.225257in}}%
\pgfpathlineto{\pgfqpoint{3.786441in}{1.233093in}}%
\pgfpathlineto{\pgfqpoint{3.778689in}{1.226135in}}%
\pgfpathlineto{\pgfqpoint{3.774576in}{1.214504in}}%
\pgfpathlineto{\pgfqpoint{3.763597in}{1.230154in}}%
\pgfpathlineto{\pgfqpoint{3.759369in}{1.233003in}}%
\pgfpathlineto{\pgfqpoint{3.760980in}{1.250036in}}%
\pgfpathlineto{\pgfqpoint{3.757634in}{1.260081in}}%
\pgfpathlineto{\pgfqpoint{3.751568in}{1.269517in}}%
\pgfpathlineto{\pgfqpoint{3.739145in}{1.298432in}}%
\pgfpathlineto{\pgfqpoint{3.743207in}{1.304826in}}%
\pgfpathlineto{\pgfqpoint{3.743159in}{1.325033in}}%
\pgfpathlineto{\pgfqpoint{3.749864in}{1.336303in}}%
\pgfpathlineto{\pgfqpoint{3.751401in}{1.353918in}}%
\pgfpathlineto{\pgfqpoint{3.745079in}{1.374099in}}%
\pgfpathlineto{\pgfqpoint{3.736692in}{1.386945in}}%
\pgfpathlineto{\pgfqpoint{3.738188in}{1.398536in}}%
\pgfpathlineto{\pgfqpoint{3.761740in}{1.426591in}}%
\pgfpathlineto{\pgfqpoint{3.762910in}{1.436161in}}%
\pgfpathlineto{\pgfqpoint{3.773523in}{1.454416in}}%
\pgfpathlineto{\pgfqpoint{3.774995in}{1.471719in}}%
\pgfpathlineto{\pgfqpoint{3.771583in}{1.476134in}}%
\pgfpathlineto{\pgfqpoint{3.777592in}{1.488751in}}%
\pgfpathclose%
\pgfusepath{fill}%
\end{pgfscope}%
\begin{pgfscope}%
\pgfpathrectangle{\pgfqpoint{3.525000in}{0.100000in}}{\pgfqpoint{2.862985in}{1.832778in}}%
\pgfusepath{clip}%
\pgfsetbuttcap%
\pgfsetmiterjoin%
\definecolor{currentfill}{rgb}{0.760264,0.235976,0.505794}%
\pgfsetfillcolor{currentfill}%
\pgfsetlinewidth{0.000000pt}%
\definecolor{currentstroke}{rgb}{0.000000,0.000000,0.000000}%
\pgfsetstrokecolor{currentstroke}%
\pgfsetstrokeopacity{0.000000}%
\pgfsetdash{}{0pt}%
\pgfpathmoveto{\pgfqpoint{5.611981in}{1.117748in}}%
\pgfpathlineto{\pgfqpoint{5.604011in}{1.195770in}}%
\pgfpathlineto{\pgfqpoint{5.594752in}{1.279623in}}%
\pgfpathlineto{\pgfqpoint{5.655390in}{1.288653in}}%
\pgfpathlineto{\pgfqpoint{5.671483in}{1.284465in}}%
\pgfpathlineto{\pgfqpoint{5.679202in}{1.279918in}}%
\pgfpathlineto{\pgfqpoint{5.688874in}{1.281179in}}%
\pgfpathlineto{\pgfqpoint{5.701627in}{1.273651in}}%
\pgfpathlineto{\pgfqpoint{5.725380in}{1.284855in}}%
\pgfpathlineto{\pgfqpoint{5.738489in}{1.285229in}}%
\pgfpathlineto{\pgfqpoint{5.753797in}{1.302307in}}%
\pgfpathlineto{\pgfqpoint{5.769372in}{1.312695in}}%
\pgfpathlineto{\pgfqpoint{5.790146in}{1.324661in}}%
\pgfpathlineto{\pgfqpoint{5.803519in}{1.241018in}}%
\pgfpathlineto{\pgfqpoint{5.797306in}{1.235646in}}%
\pgfpathlineto{\pgfqpoint{5.801317in}{1.230712in}}%
\pgfpathlineto{\pgfqpoint{5.802915in}{1.219896in}}%
\pgfpathlineto{\pgfqpoint{5.799803in}{1.209737in}}%
\pgfpathlineto{\pgfqpoint{5.799391in}{1.194887in}}%
\pgfpathlineto{\pgfqpoint{5.796468in}{1.175560in}}%
\pgfpathlineto{\pgfqpoint{5.781678in}{1.158322in}}%
\pgfpathlineto{\pgfqpoint{5.775452in}{1.154663in}}%
\pgfpathlineto{\pgfqpoint{5.770610in}{1.157813in}}%
\pgfpathlineto{\pgfqpoint{5.758627in}{1.141383in}}%
\pgfpathlineto{\pgfqpoint{5.759750in}{1.125564in}}%
\pgfpathlineto{\pgfqpoint{5.746406in}{1.129368in}}%
\pgfpathlineto{\pgfqpoint{5.740087in}{1.113569in}}%
\pgfpathlineto{\pgfqpoint{5.743281in}{1.102373in}}%
\pgfpathlineto{\pgfqpoint{5.738281in}{1.100719in}}%
\pgfpathlineto{\pgfqpoint{5.737581in}{1.091873in}}%
\pgfpathlineto{\pgfqpoint{5.725308in}{1.088305in}}%
\pgfpathlineto{\pgfqpoint{5.718930in}{1.095447in}}%
\pgfpathlineto{\pgfqpoint{5.710724in}{1.098217in}}%
\pgfpathlineto{\pgfqpoint{5.708796in}{1.105459in}}%
\pgfpathlineto{\pgfqpoint{5.701376in}{1.104185in}}%
\pgfpathlineto{\pgfqpoint{5.696515in}{1.097524in}}%
\pgfpathlineto{\pgfqpoint{5.689561in}{1.095181in}}%
\pgfpathlineto{\pgfqpoint{5.677277in}{1.099895in}}%
\pgfpathlineto{\pgfqpoint{5.670483in}{1.094286in}}%
\pgfpathlineto{\pgfqpoint{5.660827in}{1.100922in}}%
\pgfpathlineto{\pgfqpoint{5.642303in}{1.102959in}}%
\pgfpathlineto{\pgfqpoint{5.636715in}{1.115014in}}%
\pgfpathlineto{\pgfqpoint{5.629636in}{1.120355in}}%
\pgfpathlineto{\pgfqpoint{5.622773in}{1.116936in}}%
\pgfpathlineto{\pgfqpoint{5.611981in}{1.117748in}}%
\pgfpathclose%
\pgfusepath{fill}%
\end{pgfscope}%
\begin{pgfscope}%
\pgfpathrectangle{\pgfqpoint{3.525000in}{0.100000in}}{\pgfqpoint{2.862985in}{1.832778in}}%
\pgfusepath{clip}%
\pgfsetbuttcap%
\pgfsetmiterjoin%
\definecolor{currentfill}{rgb}{0.665129,0.138566,0.585582}%
\pgfsetfillcolor{currentfill}%
\pgfsetlinewidth{0.000000pt}%
\definecolor{currentstroke}{rgb}{0.000000,0.000000,0.000000}%
\pgfsetstrokecolor{currentstroke}%
\pgfsetstrokeopacity{0.000000}%
\pgfsetdash{}{0pt}%
\pgfpathmoveto{\pgfqpoint{5.411225in}{0.966856in}}%
\pgfpathlineto{\pgfqpoint{5.404121in}{0.972628in}}%
\pgfpathlineto{\pgfqpoint{5.398328in}{0.969882in}}%
\pgfpathlineto{\pgfqpoint{5.390922in}{0.983711in}}%
\pgfpathlineto{\pgfqpoint{5.394705in}{0.992402in}}%
\pgfpathlineto{\pgfqpoint{5.389252in}{1.002186in}}%
\pgfpathlineto{\pgfqpoint{5.388959in}{1.009888in}}%
\pgfpathlineto{\pgfqpoint{5.381591in}{1.012632in}}%
\pgfpathlineto{\pgfqpoint{5.378174in}{1.018436in}}%
\pgfpathlineto{\pgfqpoint{5.363768in}{1.025684in}}%
\pgfpathlineto{\pgfqpoint{5.351256in}{1.034614in}}%
\pgfpathlineto{\pgfqpoint{5.345442in}{1.041355in}}%
\pgfpathlineto{\pgfqpoint{5.345322in}{1.049578in}}%
\pgfpathlineto{\pgfqpoint{5.353102in}{1.065365in}}%
\pgfpathlineto{\pgfqpoint{5.355495in}{1.077440in}}%
\pgfpathlineto{\pgfqpoint{5.349157in}{1.084263in}}%
\pgfpathlineto{\pgfqpoint{5.340763in}{1.086834in}}%
\pgfpathlineto{\pgfqpoint{5.330579in}{1.081212in}}%
\pgfpathlineto{\pgfqpoint{5.326269in}{1.090863in}}%
\pgfpathlineto{\pgfqpoint{5.324094in}{1.103945in}}%
\pgfpathlineto{\pgfqpoint{5.309086in}{1.115607in}}%
\pgfpathlineto{\pgfqpoint{5.306087in}{1.120782in}}%
\pgfpathlineto{\pgfqpoint{5.292379in}{1.132498in}}%
\pgfpathlineto{\pgfqpoint{5.288083in}{1.141015in}}%
\pgfpathlineto{\pgfqpoint{5.284210in}{1.157906in}}%
\pgfpathlineto{\pgfqpoint{5.286865in}{1.172908in}}%
\pgfpathlineto{\pgfqpoint{5.290408in}{1.174999in}}%
\pgfpathlineto{\pgfqpoint{5.289771in}{1.187558in}}%
\pgfpathlineto{\pgfqpoint{5.299773in}{1.191288in}}%
\pgfpathlineto{\pgfqpoint{5.302793in}{1.202558in}}%
\pgfpathlineto{\pgfqpoint{5.308547in}{1.210143in}}%
\pgfpathlineto{\pgfqpoint{5.308256in}{1.219794in}}%
\pgfpathlineto{\pgfqpoint{5.301124in}{1.227469in}}%
\pgfpathlineto{\pgfqpoint{5.302818in}{1.238212in}}%
\pgfpathlineto{\pgfqpoint{5.321306in}{1.242884in}}%
\pgfpathlineto{\pgfqpoint{5.335490in}{1.251413in}}%
\pgfpathlineto{\pgfqpoint{5.336979in}{1.262151in}}%
\pgfpathlineto{\pgfqpoint{5.341914in}{1.265532in}}%
\pgfpathlineto{\pgfqpoint{5.343802in}{1.276813in}}%
\pgfpathlineto{\pgfqpoint{5.342285in}{1.284265in}}%
\pgfpathlineto{\pgfqpoint{5.332588in}{1.290454in}}%
\pgfpathlineto{\pgfqpoint{5.328684in}{1.299681in}}%
\pgfpathlineto{\pgfqpoint{5.319114in}{1.308600in}}%
\pgfpathlineto{\pgfqpoint{5.397758in}{1.311952in}}%
\pgfpathlineto{\pgfqpoint{5.450520in}{1.315700in}}%
\pgfpathlineto{\pgfqpoint{5.449554in}{1.304601in}}%
\pgfpathlineto{\pgfqpoint{5.458546in}{1.289293in}}%
\pgfpathlineto{\pgfqpoint{5.462322in}{1.276213in}}%
\pgfpathlineto{\pgfqpoint{5.466818in}{1.268787in}}%
\pgfpathlineto{\pgfqpoint{5.471726in}{1.207597in}}%
\pgfpathlineto{\pgfqpoint{5.478542in}{1.120373in}}%
\pgfpathlineto{\pgfqpoint{5.473283in}{1.107118in}}%
\pgfpathlineto{\pgfqpoint{5.480859in}{1.096224in}}%
\pgfpathlineto{\pgfqpoint{5.482971in}{1.085818in}}%
\pgfpathlineto{\pgfqpoint{5.472524in}{1.064241in}}%
\pgfpathlineto{\pgfqpoint{5.473140in}{1.061261in}}%
\pgfpathlineto{\pgfqpoint{5.462387in}{1.048893in}}%
\pgfpathlineto{\pgfqpoint{5.465336in}{1.045116in}}%
\pgfpathlineto{\pgfqpoint{5.460862in}{1.037407in}}%
\pgfpathlineto{\pgfqpoint{5.462013in}{1.021868in}}%
\pgfpathlineto{\pgfqpoint{5.456580in}{1.012333in}}%
\pgfpathlineto{\pgfqpoint{5.460999in}{1.001067in}}%
\pgfpathlineto{\pgfqpoint{5.442483in}{0.994944in}}%
\pgfpathlineto{\pgfqpoint{5.440777in}{0.988286in}}%
\pgfpathlineto{\pgfqpoint{5.445833in}{0.979847in}}%
\pgfpathlineto{\pgfqpoint{5.443506in}{0.974346in}}%
\pgfpathlineto{\pgfqpoint{5.423652in}{0.981141in}}%
\pgfpathlineto{\pgfqpoint{5.417101in}{0.981825in}}%
\pgfpathlineto{\pgfqpoint{5.408934in}{0.971495in}}%
\pgfpathlineto{\pgfqpoint{5.411225in}{0.966856in}}%
\pgfpathclose%
\pgfusepath{fill}%
\end{pgfscope}%
\begin{pgfscope}%
\pgfpathrectangle{\pgfqpoint{3.525000in}{0.100000in}}{\pgfqpoint{2.862985in}{1.832778in}}%
\pgfusepath{clip}%
\pgfsetbuttcap%
\pgfsetmiterjoin%
\definecolor{currentfill}{rgb}{0.429719,0.000831,0.659425}%
\pgfsetfillcolor{currentfill}%
\pgfsetlinewidth{0.000000pt}%
\definecolor{currentstroke}{rgb}{0.000000,0.000000,0.000000}%
\pgfsetstrokecolor{currentstroke}%
\pgfsetstrokeopacity{0.000000}%
\pgfsetdash{}{0pt}%
\pgfpathmoveto{\pgfqpoint{5.989161in}{1.156689in}}%
\pgfpathlineto{\pgfqpoint{5.983697in}{1.164803in}}%
\pgfpathlineto{\pgfqpoint{5.986780in}{1.169343in}}%
\pgfpathlineto{\pgfqpoint{5.994331in}{1.164244in}}%
\pgfpathlineto{\pgfqpoint{5.989161in}{1.156689in}}%
\pgfpathclose%
\pgfusepath{fill}%
\end{pgfscope}%
\begin{pgfscope}%
\pgfpathrectangle{\pgfqpoint{3.525000in}{0.100000in}}{\pgfqpoint{2.862985in}{1.832778in}}%
\pgfusepath{clip}%
\pgfsetbuttcap%
\pgfsetmiterjoin%
\definecolor{currentfill}{rgb}{0.809052,0.293491,0.458870}%
\pgfsetfillcolor{currentfill}%
\pgfsetlinewidth{0.000000pt}%
\definecolor{currentstroke}{rgb}{0.000000,0.000000,0.000000}%
\pgfsetstrokecolor{currentstroke}%
\pgfsetstrokeopacity{0.000000}%
\pgfsetdash{}{0pt}%
\pgfpathmoveto{\pgfqpoint{6.036873in}{1.226505in}}%
\pgfpathlineto{\pgfqpoint{6.040138in}{1.233404in}}%
\pgfpathlineto{\pgfqpoint{6.053338in}{1.234894in}}%
\pgfpathlineto{\pgfqpoint{6.051223in}{1.229031in}}%
\pgfpathlineto{\pgfqpoint{6.046868in}{1.221542in}}%
\pgfpathlineto{\pgfqpoint{6.049819in}{1.212619in}}%
\pgfpathlineto{\pgfqpoint{6.061465in}{1.201927in}}%
\pgfpathlineto{\pgfqpoint{6.064168in}{1.190668in}}%
\pgfpathlineto{\pgfqpoint{6.077585in}{1.176648in}}%
\pgfpathlineto{\pgfqpoint{6.082849in}{1.177248in}}%
\pgfpathlineto{\pgfqpoint{6.089406in}{1.156201in}}%
\pgfpathlineto{\pgfqpoint{6.088331in}{1.155991in}}%
\pgfpathlineto{\pgfqpoint{6.087135in}{1.155756in}}%
\pgfpathlineto{\pgfqpoint{6.057891in}{1.150158in}}%
\pgfpathlineto{\pgfqpoint{6.042245in}{1.205802in}}%
\pgfpathlineto{\pgfqpoint{6.036873in}{1.226505in}}%
\pgfpathclose%
\pgfusepath{fill}%
\end{pgfscope}%
\begin{pgfscope}%
\pgfpathrectangle{\pgfqpoint{3.525000in}{0.100000in}}{\pgfqpoint{2.862985in}{1.832778in}}%
\pgfusepath{clip}%
\pgfsetbuttcap%
\pgfsetmiterjoin%
\definecolor{currentfill}{rgb}{0.923287,0.463251,0.332801}%
\pgfsetfillcolor{currentfill}%
\pgfsetlinewidth{0.000000pt}%
\definecolor{currentstroke}{rgb}{0.000000,0.000000,0.000000}%
\pgfsetstrokecolor{currentstroke}%
\pgfsetstrokeopacity{0.000000}%
\pgfsetdash{}{0pt}%
\pgfpathmoveto{\pgfqpoint{5.763724in}{1.037893in}}%
\pgfpathlineto{\pgfqpoint{5.754483in}{1.038233in}}%
\pgfpathlineto{\pgfqpoint{5.745968in}{1.044098in}}%
\pgfpathlineto{\pgfqpoint{5.734465in}{1.061926in}}%
\pgfpathlineto{\pgfqpoint{5.724672in}{1.071367in}}%
\pgfpathlineto{\pgfqpoint{5.727231in}{1.078659in}}%
\pgfpathlineto{\pgfqpoint{5.725308in}{1.088305in}}%
\pgfpathlineto{\pgfqpoint{5.737581in}{1.091873in}}%
\pgfpathlineto{\pgfqpoint{5.738281in}{1.100719in}}%
\pgfpathlineto{\pgfqpoint{5.743281in}{1.102373in}}%
\pgfpathlineto{\pgfqpoint{5.740087in}{1.113569in}}%
\pgfpathlineto{\pgfqpoint{5.746406in}{1.129368in}}%
\pgfpathlineto{\pgfqpoint{5.759750in}{1.125564in}}%
\pgfpathlineto{\pgfqpoint{5.758627in}{1.141383in}}%
\pgfpathlineto{\pgfqpoint{5.770610in}{1.157813in}}%
\pgfpathlineto{\pgfqpoint{5.775452in}{1.154663in}}%
\pgfpathlineto{\pgfqpoint{5.781678in}{1.158322in}}%
\pgfpathlineto{\pgfqpoint{5.796468in}{1.175560in}}%
\pgfpathlineto{\pgfqpoint{5.799391in}{1.194887in}}%
\pgfpathlineto{\pgfqpoint{5.799803in}{1.209737in}}%
\pgfpathlineto{\pgfqpoint{5.802915in}{1.219896in}}%
\pgfpathlineto{\pgfqpoint{5.801317in}{1.230712in}}%
\pgfpathlineto{\pgfqpoint{5.797306in}{1.235646in}}%
\pgfpathlineto{\pgfqpoint{5.803519in}{1.241018in}}%
\pgfpathlineto{\pgfqpoint{5.812526in}{1.184275in}}%
\pgfpathlineto{\pgfqpoint{5.862178in}{1.192485in}}%
\pgfpathlineto{\pgfqpoint{5.867325in}{1.160097in}}%
\pgfpathlineto{\pgfqpoint{5.885300in}{1.181474in}}%
\pgfpathlineto{\pgfqpoint{5.889536in}{1.179343in}}%
\pgfpathlineto{\pgfqpoint{5.895767in}{1.191730in}}%
\pgfpathlineto{\pgfqpoint{5.904399in}{1.187846in}}%
\pgfpathlineto{\pgfqpoint{5.911925in}{1.188584in}}%
\pgfpathlineto{\pgfqpoint{5.916891in}{1.197195in}}%
\pgfpathlineto{\pgfqpoint{5.924098in}{1.201985in}}%
\pgfpathlineto{\pgfqpoint{5.934119in}{1.197766in}}%
\pgfpathlineto{\pgfqpoint{5.940712in}{1.199224in}}%
\pgfpathlineto{\pgfqpoint{5.950158in}{1.182877in}}%
\pgfpathlineto{\pgfqpoint{5.947427in}{1.170500in}}%
\pgfpathlineto{\pgfqpoint{5.922755in}{1.184435in}}%
\pgfpathlineto{\pgfqpoint{5.918134in}{1.172996in}}%
\pgfpathlineto{\pgfqpoint{5.919659in}{1.167732in}}%
\pgfpathlineto{\pgfqpoint{5.909159in}{1.149941in}}%
\pgfpathlineto{\pgfqpoint{5.904203in}{1.146402in}}%
\pgfpathlineto{\pgfqpoint{5.901965in}{1.138556in}}%
\pgfpathlineto{\pgfqpoint{5.893517in}{1.139377in}}%
\pgfpathlineto{\pgfqpoint{5.887449in}{1.117963in}}%
\pgfpathlineto{\pgfqpoint{5.884053in}{1.113049in}}%
\pgfpathlineto{\pgfqpoint{5.875326in}{1.114688in}}%
\pgfpathlineto{\pgfqpoint{5.866399in}{1.121442in}}%
\pgfpathlineto{\pgfqpoint{5.866091in}{1.111076in}}%
\pgfpathlineto{\pgfqpoint{5.857449in}{1.093654in}}%
\pgfpathlineto{\pgfqpoint{5.855300in}{1.081251in}}%
\pgfpathlineto{\pgfqpoint{5.848669in}{1.072991in}}%
\pgfpathlineto{\pgfqpoint{5.843643in}{1.059800in}}%
\pgfpathlineto{\pgfqpoint{5.843396in}{1.047589in}}%
\pgfpathlineto{\pgfqpoint{5.835633in}{1.045308in}}%
\pgfpathlineto{\pgfqpoint{5.826826in}{1.038395in}}%
\pgfpathlineto{\pgfqpoint{5.818347in}{1.037079in}}%
\pgfpathlineto{\pgfqpoint{5.816393in}{1.031228in}}%
\pgfpathlineto{\pgfqpoint{5.802730in}{1.025236in}}%
\pgfpathlineto{\pgfqpoint{5.795094in}{1.030341in}}%
\pgfpathlineto{\pgfqpoint{5.786560in}{1.020647in}}%
\pgfpathlineto{\pgfqpoint{5.771890in}{1.023505in}}%
\pgfpathlineto{\pgfqpoint{5.762893in}{1.033679in}}%
\pgfpathlineto{\pgfqpoint{5.763724in}{1.037893in}}%
\pgfpathclose%
\pgfusepath{fill}%
\end{pgfscope}%
\begin{pgfscope}%
\pgfpathrectangle{\pgfqpoint{3.525000in}{0.100000in}}{\pgfqpoint{2.862985in}{1.832778in}}%
\pgfusepath{clip}%
\pgfsetbuttcap%
\pgfsetmiterjoin%
\definecolor{currentfill}{rgb}{0.625987,0.103312,0.611305}%
\pgfsetfillcolor{currentfill}%
\pgfsetlinewidth{0.000000pt}%
\definecolor{currentstroke}{rgb}{0.000000,0.000000,0.000000}%
\pgfsetstrokecolor{currentstroke}%
\pgfsetstrokeopacity{0.000000}%
\pgfsetdash{}{0pt}%
\pgfpathmoveto{\pgfqpoint{6.087135in}{1.155756in}}%
\pgfpathlineto{\pgfqpoint{6.086760in}{1.144322in}}%
\pgfpathlineto{\pgfqpoint{6.082360in}{1.138701in}}%
\pgfpathlineto{\pgfqpoint{6.079551in}{1.126221in}}%
\pgfpathlineto{\pgfqpoint{6.066870in}{1.120469in}}%
\pgfpathlineto{\pgfqpoint{6.056229in}{1.118793in}}%
\pgfpathlineto{\pgfqpoint{6.059330in}{1.127000in}}%
\pgfpathlineto{\pgfqpoint{6.042947in}{1.133877in}}%
\pgfpathlineto{\pgfqpoint{6.029638in}{1.142487in}}%
\pgfpathlineto{\pgfqpoint{6.037878in}{1.155358in}}%
\pgfpathlineto{\pgfqpoint{6.031274in}{1.165363in}}%
\pgfpathlineto{\pgfqpoint{6.032051in}{1.173993in}}%
\pgfpathlineto{\pgfqpoint{6.027276in}{1.176362in}}%
\pgfpathlineto{\pgfqpoint{6.023402in}{1.190382in}}%
\pgfpathlineto{\pgfqpoint{6.032564in}{1.209388in}}%
\pgfpathlineto{\pgfqpoint{6.025681in}{1.212492in}}%
\pgfpathlineto{\pgfqpoint{6.023896in}{1.203115in}}%
\pgfpathlineto{\pgfqpoint{6.014070in}{1.200505in}}%
\pgfpathlineto{\pgfqpoint{6.015103in}{1.169655in}}%
\pgfpathlineto{\pgfqpoint{6.013306in}{1.159726in}}%
\pgfpathlineto{\pgfqpoint{6.018235in}{1.145575in}}%
\pgfpathlineto{\pgfqpoint{6.025818in}{1.138767in}}%
\pgfpathlineto{\pgfqpoint{6.022379in}{1.134492in}}%
\pgfpathlineto{\pgfqpoint{6.030119in}{1.128246in}}%
\pgfpathlineto{\pgfqpoint{6.032917in}{1.118107in}}%
\pgfpathlineto{\pgfqpoint{6.018730in}{1.126503in}}%
\pgfpathlineto{\pgfqpoint{6.009752in}{1.125412in}}%
\pgfpathlineto{\pgfqpoint{6.002780in}{1.134049in}}%
\pgfpathlineto{\pgfqpoint{5.995682in}{1.134890in}}%
\pgfpathlineto{\pgfqpoint{5.985594in}{1.130564in}}%
\pgfpathlineto{\pgfqpoint{5.981681in}{1.135955in}}%
\pgfpathlineto{\pgfqpoint{5.986823in}{1.147270in}}%
\pgfpathlineto{\pgfqpoint{5.989161in}{1.156689in}}%
\pgfpathlineto{\pgfqpoint{5.994331in}{1.164244in}}%
\pgfpathlineto{\pgfqpoint{5.986780in}{1.169343in}}%
\pgfpathlineto{\pgfqpoint{5.983697in}{1.164803in}}%
\pgfpathlineto{\pgfqpoint{5.976154in}{1.169411in}}%
\pgfpathlineto{\pgfqpoint{5.962796in}{1.172482in}}%
\pgfpathlineto{\pgfqpoint{5.964008in}{1.179234in}}%
\pgfpathlineto{\pgfqpoint{5.957953in}{1.183150in}}%
\pgfpathlineto{\pgfqpoint{5.950158in}{1.182877in}}%
\pgfpathlineto{\pgfqpoint{5.940712in}{1.199224in}}%
\pgfpathlineto{\pgfqpoint{5.934119in}{1.197766in}}%
\pgfpathlineto{\pgfqpoint{5.924098in}{1.201985in}}%
\pgfpathlineto{\pgfqpoint{5.916891in}{1.197195in}}%
\pgfpathlineto{\pgfqpoint{5.911925in}{1.188584in}}%
\pgfpathlineto{\pgfqpoint{5.904399in}{1.187846in}}%
\pgfpathlineto{\pgfqpoint{5.895767in}{1.191730in}}%
\pgfpathlineto{\pgfqpoint{5.889536in}{1.179343in}}%
\pgfpathlineto{\pgfqpoint{5.885300in}{1.181474in}}%
\pgfpathlineto{\pgfqpoint{5.867325in}{1.160097in}}%
\pgfpathlineto{\pgfqpoint{5.862178in}{1.192485in}}%
\pgfpathlineto{\pgfqpoint{5.927868in}{1.204659in}}%
\pgfpathlineto{\pgfqpoint{5.957316in}{1.209933in}}%
\pgfpathlineto{\pgfqpoint{6.000158in}{1.218561in}}%
\pgfpathlineto{\pgfqpoint{6.036873in}{1.226505in}}%
\pgfpathlineto{\pgfqpoint{6.042245in}{1.205802in}}%
\pgfpathlineto{\pgfqpoint{6.057891in}{1.150158in}}%
\pgfpathlineto{\pgfqpoint{6.087135in}{1.155756in}}%
\pgfpathclose%
\pgfusepath{fill}%
\end{pgfscope}%
\begin{pgfscope}%
\pgfpathrectangle{\pgfqpoint{3.525000in}{0.100000in}}{\pgfqpoint{2.862985in}{1.832778in}}%
\pgfusepath{clip}%
\pgfsetbuttcap%
\pgfsetmiterjoin%
\definecolor{currentfill}{rgb}{0.471457,0.005678,0.659897}%
\pgfsetfillcolor{currentfill}%
\pgfsetlinewidth{0.000000pt}%
\definecolor{currentstroke}{rgb}{0.000000,0.000000,0.000000}%
\pgfsetstrokecolor{currentstroke}%
\pgfsetstrokeopacity{0.000000}%
\pgfsetdash{}{0pt}%
\pgfpathmoveto{\pgfqpoint{4.764821in}{0.971059in}}%
\pgfpathlineto{\pgfqpoint{4.716895in}{0.975633in}}%
\pgfpathlineto{\pgfqpoint{4.667175in}{0.980039in}}%
\pgfpathlineto{\pgfqpoint{4.609722in}{0.986032in}}%
\pgfpathlineto{\pgfqpoint{4.524362in}{0.996184in}}%
\pgfpathlineto{\pgfqpoint{4.494080in}{1.000850in}}%
\pgfpathlineto{\pgfqpoint{4.415915in}{1.012202in}}%
\pgfpathlineto{\pgfqpoint{4.427328in}{1.083837in}}%
\pgfpathlineto{\pgfqpoint{4.427613in}{1.089632in}}%
\pgfpathlineto{\pgfqpoint{4.438608in}{1.158753in}}%
\pgfpathlineto{\pgfqpoint{4.446820in}{1.211396in}}%
\pgfpathlineto{\pgfqpoint{4.454571in}{1.260214in}}%
\pgfpathlineto{\pgfqpoint{4.507522in}{1.252613in}}%
\pgfpathlineto{\pgfqpoint{4.556892in}{1.245527in}}%
\pgfpathlineto{\pgfqpoint{4.647647in}{1.234326in}}%
\pgfpathlineto{\pgfqpoint{4.689283in}{1.230524in}}%
\pgfpathlineto{\pgfqpoint{4.755243in}{1.224124in}}%
\pgfpathlineto{\pgfqpoint{4.783777in}{1.221808in}}%
\pgfpathlineto{\pgfqpoint{4.778739in}{1.159333in}}%
\pgfpathlineto{\pgfqpoint{4.774188in}{1.099180in}}%
\pgfpathlineto{\pgfqpoint{4.770525in}{1.050201in}}%
\pgfpathlineto{\pgfqpoint{4.764821in}{0.971059in}}%
\pgfpathclose%
\pgfusepath{fill}%
\end{pgfscope}%
\begin{pgfscope}%
\pgfpathrectangle{\pgfqpoint{3.525000in}{0.100000in}}{\pgfqpoint{2.862985in}{1.832778in}}%
\pgfusepath{clip}%
\pgfsetbuttcap%
\pgfsetmiterjoin%
\definecolor{currentfill}{rgb}{0.856547,0.356066,0.410322}%
\pgfsetfillcolor{currentfill}%
\pgfsetlinewidth{0.000000pt}%
\definecolor{currentstroke}{rgb}{0.000000,0.000000,0.000000}%
\pgfsetstrokecolor{currentstroke}%
\pgfsetstrokeopacity{0.000000}%
\pgfsetdash{}{0pt}%
\pgfpathmoveto{\pgfqpoint{5.398847in}{0.935544in}}%
\pgfpathlineto{\pgfqpoint{5.400463in}{0.942786in}}%
\pgfpathlineto{\pgfqpoint{5.408837in}{0.941156in}}%
\pgfpathlineto{\pgfqpoint{5.411225in}{0.966856in}}%
\pgfpathlineto{\pgfqpoint{5.408934in}{0.971495in}}%
\pgfpathlineto{\pgfqpoint{5.417101in}{0.981825in}}%
\pgfpathlineto{\pgfqpoint{5.423652in}{0.981141in}}%
\pgfpathlineto{\pgfqpoint{5.443506in}{0.974346in}}%
\pgfpathlineto{\pgfqpoint{5.445833in}{0.979847in}}%
\pgfpathlineto{\pgfqpoint{5.440777in}{0.988286in}}%
\pgfpathlineto{\pgfqpoint{5.442483in}{0.994944in}}%
\pgfpathlineto{\pgfqpoint{5.460999in}{1.001067in}}%
\pgfpathlineto{\pgfqpoint{5.456580in}{1.012333in}}%
\pgfpathlineto{\pgfqpoint{5.462013in}{1.021868in}}%
\pgfpathlineto{\pgfqpoint{5.472127in}{1.027220in}}%
\pgfpathlineto{\pgfqpoint{5.493349in}{1.032520in}}%
\pgfpathlineto{\pgfqpoint{5.506739in}{1.024520in}}%
\pgfpathlineto{\pgfqpoint{5.510801in}{1.032321in}}%
\pgfpathlineto{\pgfqpoint{5.521919in}{1.037771in}}%
\pgfpathlineto{\pgfqpoint{5.524789in}{1.032965in}}%
\pgfpathlineto{\pgfqpoint{5.536209in}{1.036783in}}%
\pgfpathlineto{\pgfqpoint{5.535489in}{1.043318in}}%
\pgfpathlineto{\pgfqpoint{5.545757in}{1.050797in}}%
\pgfpathlineto{\pgfqpoint{5.551817in}{1.043031in}}%
\pgfpathlineto{\pgfqpoint{5.559790in}{1.042250in}}%
\pgfpathlineto{\pgfqpoint{5.565088in}{1.047311in}}%
\pgfpathlineto{\pgfqpoint{5.564498in}{1.054509in}}%
\pgfpathlineto{\pgfqpoint{5.569002in}{1.061656in}}%
\pgfpathlineto{\pgfqpoint{5.575037in}{1.063205in}}%
\pgfpathlineto{\pgfqpoint{5.577460in}{1.072642in}}%
\pgfpathlineto{\pgfqpoint{5.586237in}{1.080801in}}%
\pgfpathlineto{\pgfqpoint{5.583595in}{1.088922in}}%
\pgfpathlineto{\pgfqpoint{5.592134in}{1.092965in}}%
\pgfpathlineto{\pgfqpoint{5.597839in}{1.090478in}}%
\pgfpathlineto{\pgfqpoint{5.606265in}{1.096795in}}%
\pgfpathlineto{\pgfqpoint{5.613795in}{1.098441in}}%
\pgfpathlineto{\pgfqpoint{5.613822in}{1.104993in}}%
\pgfpathlineto{\pgfqpoint{5.608545in}{1.114085in}}%
\pgfpathlineto{\pgfqpoint{5.611981in}{1.117748in}}%
\pgfpathlineto{\pgfqpoint{5.622773in}{1.116936in}}%
\pgfpathlineto{\pgfqpoint{5.629636in}{1.120355in}}%
\pgfpathlineto{\pgfqpoint{5.636715in}{1.115014in}}%
\pgfpathlineto{\pgfqpoint{5.642303in}{1.102959in}}%
\pgfpathlineto{\pgfqpoint{5.660827in}{1.100922in}}%
\pgfpathlineto{\pgfqpoint{5.670483in}{1.094286in}}%
\pgfpathlineto{\pgfqpoint{5.677277in}{1.099895in}}%
\pgfpathlineto{\pgfqpoint{5.689561in}{1.095181in}}%
\pgfpathlineto{\pgfqpoint{5.696515in}{1.097524in}}%
\pgfpathlineto{\pgfqpoint{5.701376in}{1.104185in}}%
\pgfpathlineto{\pgfqpoint{5.708796in}{1.105459in}}%
\pgfpathlineto{\pgfqpoint{5.710724in}{1.098217in}}%
\pgfpathlineto{\pgfqpoint{5.718930in}{1.095447in}}%
\pgfpathlineto{\pgfqpoint{5.725308in}{1.088305in}}%
\pgfpathlineto{\pgfqpoint{5.727231in}{1.078659in}}%
\pgfpathlineto{\pgfqpoint{5.724672in}{1.071367in}}%
\pgfpathlineto{\pgfqpoint{5.734465in}{1.061926in}}%
\pgfpathlineto{\pgfqpoint{5.745968in}{1.044098in}}%
\pgfpathlineto{\pgfqpoint{5.754483in}{1.038233in}}%
\pgfpathlineto{\pgfqpoint{5.763724in}{1.037893in}}%
\pgfpathlineto{\pgfqpoint{5.746661in}{1.018314in}}%
\pgfpathlineto{\pgfqpoint{5.729865in}{1.006463in}}%
\pgfpathlineto{\pgfqpoint{5.723683in}{0.997053in}}%
\pgfpathlineto{\pgfqpoint{5.723788in}{0.991945in}}%
\pgfpathlineto{\pgfqpoint{5.714691in}{0.988030in}}%
\pgfpathlineto{\pgfqpoint{5.712084in}{0.980665in}}%
\pgfpathlineto{\pgfqpoint{5.693132in}{0.973246in}}%
\pgfpathlineto{\pgfqpoint{5.686414in}{0.968427in}}%
\pgfpathlineto{\pgfqpoint{5.685505in}{0.967399in}}%
\pgfpathlineto{\pgfqpoint{5.630991in}{0.962230in}}%
\pgfpathlineto{\pgfqpoint{5.598070in}{0.959461in}}%
\pgfpathlineto{\pgfqpoint{5.544049in}{0.956404in}}%
\pgfpathlineto{\pgfqpoint{5.476748in}{0.949746in}}%
\pgfpathlineto{\pgfqpoint{5.465629in}{0.951286in}}%
\pgfpathlineto{\pgfqpoint{5.467942in}{0.939936in}}%
\pgfpathlineto{\pgfqpoint{5.398847in}{0.935544in}}%
\pgfpathclose%
\pgfusepath{fill}%
\end{pgfscope}%
\begin{pgfscope}%
\pgfpathrectangle{\pgfqpoint{3.525000in}{0.100000in}}{\pgfqpoint{2.862985in}{1.832778in}}%
\pgfusepath{clip}%
\pgfsetbuttcap%
\pgfsetmiterjoin%
\definecolor{currentfill}{rgb}{0.551715,0.043136,0.645277}%
\pgfsetfillcolor{currentfill}%
\pgfsetlinewidth{0.000000pt}%
\definecolor{currentstroke}{rgb}{0.000000,0.000000,0.000000}%
\pgfsetstrokecolor{currentstroke}%
\pgfsetstrokeopacity{0.000000}%
\pgfsetdash{}{0pt}%
\pgfpathmoveto{\pgfqpoint{4.764821in}{0.971059in}}%
\pgfpathlineto{\pgfqpoint{4.770525in}{1.050201in}}%
\pgfpathlineto{\pgfqpoint{4.774188in}{1.099180in}}%
\pgfpathlineto{\pgfqpoint{4.778739in}{1.159333in}}%
\pgfpathlineto{\pgfqpoint{4.841835in}{1.154903in}}%
\pgfpathlineto{\pgfqpoint{4.921966in}{1.150526in}}%
\pgfpathlineto{\pgfqpoint{4.976466in}{1.148434in}}%
\pgfpathlineto{\pgfqpoint{5.030647in}{1.146774in}}%
\pgfpathlineto{\pgfqpoint{5.079686in}{1.146016in}}%
\pgfpathlineto{\pgfqpoint{5.102367in}{1.146251in}}%
\pgfpathlineto{\pgfqpoint{5.112350in}{1.138104in}}%
\pgfpathlineto{\pgfqpoint{5.120170in}{1.139746in}}%
\pgfpathlineto{\pgfqpoint{5.120414in}{1.132640in}}%
\pgfpathlineto{\pgfqpoint{5.114607in}{1.120348in}}%
\pgfpathlineto{\pgfqpoint{5.115237in}{1.112582in}}%
\pgfpathlineto{\pgfqpoint{5.120797in}{1.107461in}}%
\pgfpathlineto{\pgfqpoint{5.125905in}{1.096790in}}%
\pgfpathlineto{\pgfqpoint{5.136265in}{1.090665in}}%
\pgfpathlineto{\pgfqpoint{5.135919in}{1.050444in}}%
\pgfpathlineto{\pgfqpoint{5.136224in}{0.957954in}}%
\pgfpathlineto{\pgfqpoint{5.066779in}{0.958274in}}%
\pgfpathlineto{\pgfqpoint{4.993651in}{0.959554in}}%
\pgfpathlineto{\pgfqpoint{4.939818in}{0.961453in}}%
\pgfpathlineto{\pgfqpoint{4.894882in}{0.963187in}}%
\pgfpathlineto{\pgfqpoint{4.819179in}{0.967651in}}%
\pgfpathlineto{\pgfqpoint{4.764821in}{0.971059in}}%
\pgfpathclose%
\pgfusepath{fill}%
\end{pgfscope}%
\begin{pgfscope}%
\pgfpathrectangle{\pgfqpoint{3.525000in}{0.100000in}}{\pgfqpoint{2.862985in}{1.832778in}}%
\pgfusepath{clip}%
\pgfsetbuttcap%
\pgfsetmiterjoin%
\definecolor{currentfill}{rgb}{0.620919,0.098934,0.614257}%
\pgfsetfillcolor{currentfill}%
\pgfsetlinewidth{0.000000pt}%
\definecolor{currentstroke}{rgb}{0.000000,0.000000,0.000000}%
\pgfsetstrokecolor{currentstroke}%
\pgfsetstrokeopacity{0.000000}%
\pgfsetdash{}{0pt}%
\pgfpathmoveto{\pgfqpoint{5.786621in}{0.981538in}}%
\pgfpathlineto{\pgfqpoint{5.709707in}{0.970734in}}%
\pgfpathlineto{\pgfqpoint{5.686414in}{0.968427in}}%
\pgfpathlineto{\pgfqpoint{5.693132in}{0.973246in}}%
\pgfpathlineto{\pgfqpoint{5.712084in}{0.980665in}}%
\pgfpathlineto{\pgfqpoint{5.714691in}{0.988030in}}%
\pgfpathlineto{\pgfqpoint{5.723788in}{0.991945in}}%
\pgfpathlineto{\pgfqpoint{5.723683in}{0.997053in}}%
\pgfpathlineto{\pgfqpoint{5.729865in}{1.006463in}}%
\pgfpathlineto{\pgfqpoint{5.746661in}{1.018314in}}%
\pgfpathlineto{\pgfqpoint{5.763724in}{1.037893in}}%
\pgfpathlineto{\pgfqpoint{5.762893in}{1.033679in}}%
\pgfpathlineto{\pgfqpoint{5.771890in}{1.023505in}}%
\pgfpathlineto{\pgfqpoint{5.786560in}{1.020647in}}%
\pgfpathlineto{\pgfqpoint{5.795094in}{1.030341in}}%
\pgfpathlineto{\pgfqpoint{5.802730in}{1.025236in}}%
\pgfpathlineto{\pgfqpoint{5.816393in}{1.031228in}}%
\pgfpathlineto{\pgfqpoint{5.818347in}{1.037079in}}%
\pgfpathlineto{\pgfqpoint{5.826826in}{1.038395in}}%
\pgfpathlineto{\pgfqpoint{5.835633in}{1.045308in}}%
\pgfpathlineto{\pgfqpoint{5.843396in}{1.047589in}}%
\pgfpathlineto{\pgfqpoint{5.843643in}{1.059800in}}%
\pgfpathlineto{\pgfqpoint{5.848669in}{1.072991in}}%
\pgfpathlineto{\pgfqpoint{5.855300in}{1.081251in}}%
\pgfpathlineto{\pgfqpoint{5.857449in}{1.093654in}}%
\pgfpathlineto{\pgfqpoint{5.866091in}{1.111076in}}%
\pgfpathlineto{\pgfqpoint{5.866399in}{1.121442in}}%
\pgfpathlineto{\pgfqpoint{5.875326in}{1.114688in}}%
\pgfpathlineto{\pgfqpoint{5.884053in}{1.113049in}}%
\pgfpathlineto{\pgfqpoint{5.887449in}{1.117963in}}%
\pgfpathlineto{\pgfqpoint{5.893517in}{1.139377in}}%
\pgfpathlineto{\pgfqpoint{5.901965in}{1.138556in}}%
\pgfpathlineto{\pgfqpoint{5.904203in}{1.146402in}}%
\pgfpathlineto{\pgfqpoint{5.909159in}{1.149941in}}%
\pgfpathlineto{\pgfqpoint{5.919659in}{1.167732in}}%
\pgfpathlineto{\pgfqpoint{5.918134in}{1.172996in}}%
\pgfpathlineto{\pgfqpoint{5.922755in}{1.184435in}}%
\pgfpathlineto{\pgfqpoint{5.947427in}{1.170500in}}%
\pgfpathlineto{\pgfqpoint{5.950158in}{1.182877in}}%
\pgfpathlineto{\pgfqpoint{5.957953in}{1.183150in}}%
\pgfpathlineto{\pgfqpoint{5.964008in}{1.179234in}}%
\pgfpathlineto{\pgfqpoint{5.962796in}{1.172482in}}%
\pgfpathlineto{\pgfqpoint{5.976154in}{1.169411in}}%
\pgfpathlineto{\pgfqpoint{5.983697in}{1.164803in}}%
\pgfpathlineto{\pgfqpoint{5.989161in}{1.156689in}}%
\pgfpathlineto{\pgfqpoint{5.986823in}{1.147270in}}%
\pgfpathlineto{\pgfqpoint{5.982103in}{1.146498in}}%
\pgfpathlineto{\pgfqpoint{5.979365in}{1.132285in}}%
\pgfpathlineto{\pgfqpoint{5.985362in}{1.126728in}}%
\pgfpathlineto{\pgfqpoint{5.993804in}{1.131220in}}%
\pgfpathlineto{\pgfqpoint{6.001637in}{1.121734in}}%
\pgfpathlineto{\pgfqpoint{6.019137in}{1.120029in}}%
\pgfpathlineto{\pgfqpoint{6.022152in}{1.114573in}}%
\pgfpathlineto{\pgfqpoint{6.038348in}{1.109266in}}%
\pgfpathlineto{\pgfqpoint{6.036349in}{1.103004in}}%
\pgfpathlineto{\pgfqpoint{6.038456in}{1.085305in}}%
\pgfpathlineto{\pgfqpoint{6.044993in}{1.078622in}}%
\pgfpathlineto{\pgfqpoint{6.035347in}{1.075802in}}%
\pgfpathlineto{\pgfqpoint{6.036626in}{1.068246in}}%
\pgfpathlineto{\pgfqpoint{6.046916in}{1.061840in}}%
\pgfpathlineto{\pgfqpoint{6.047844in}{1.055521in}}%
\pgfpathlineto{\pgfqpoint{6.042034in}{1.050772in}}%
\pgfpathlineto{\pgfqpoint{6.030312in}{1.062022in}}%
\pgfpathlineto{\pgfqpoint{6.029153in}{1.053814in}}%
\pgfpathlineto{\pgfqpoint{6.038966in}{1.049911in}}%
\pgfpathlineto{\pgfqpoint{6.039948in}{1.045875in}}%
\pgfpathlineto{\pgfqpoint{6.053408in}{1.051209in}}%
\pgfpathlineto{\pgfqpoint{6.063722in}{1.052620in}}%
\pgfpathlineto{\pgfqpoint{6.074287in}{1.031303in}}%
\pgfpathlineto{\pgfqpoint{6.073108in}{1.031073in}}%
\pgfpathlineto{\pgfqpoint{6.068334in}{1.030086in}}%
\pgfpathlineto{\pgfqpoint{6.066927in}{1.029792in}}%
\pgfpathlineto{\pgfqpoint{6.065997in}{1.029610in}}%
\pgfpathlineto{\pgfqpoint{6.003091in}{1.016561in}}%
\pgfpathlineto{\pgfqpoint{5.946826in}{1.005045in}}%
\pgfpathlineto{\pgfqpoint{5.869033in}{0.991644in}}%
\pgfpathlineto{\pgfqpoint{5.802964in}{0.982918in}}%
\pgfpathlineto{\pgfqpoint{5.786621in}{0.981538in}}%
\pgfpathclose%
\pgfusepath{fill}%
\end{pgfscope}%
\begin{pgfscope}%
\pgfpathrectangle{\pgfqpoint{3.525000in}{0.100000in}}{\pgfqpoint{2.862985in}{1.832778in}}%
\pgfusepath{clip}%
\pgfsetbuttcap%
\pgfsetmiterjoin%
\definecolor{currentfill}{rgb}{0.620919,0.098934,0.614257}%
\pgfsetfillcolor{currentfill}%
\pgfsetlinewidth{0.000000pt}%
\definecolor{currentstroke}{rgb}{0.000000,0.000000,0.000000}%
\pgfsetstrokecolor{currentstroke}%
\pgfsetstrokeopacity{0.000000}%
\pgfsetdash{}{0pt}%
\pgfpathmoveto{\pgfqpoint{6.066870in}{1.120469in}}%
\pgfpathlineto{\pgfqpoint{6.079551in}{1.126221in}}%
\pgfpathlineto{\pgfqpoint{6.069424in}{1.096596in}}%
\pgfpathlineto{\pgfqpoint{6.062727in}{1.080908in}}%
\pgfpathlineto{\pgfqpoint{6.057902in}{1.089786in}}%
\pgfpathlineto{\pgfqpoint{6.066482in}{1.111041in}}%
\pgfpathlineto{\pgfqpoint{6.066870in}{1.120469in}}%
\pgfpathclose%
\pgfusepath{fill}%
\end{pgfscope}%
\begin{pgfscope}%
\pgfpathrectangle{\pgfqpoint{3.525000in}{0.100000in}}{\pgfqpoint{2.862985in}{1.832778in}}%
\pgfusepath{clip}%
\pgfsetbuttcap%
\pgfsetmiterjoin%
\definecolor{currentfill}{rgb}{0.692840,0.165141,0.564522}%
\pgfsetfillcolor{currentfill}%
\pgfsetlinewidth{0.000000pt}%
\definecolor{currentstroke}{rgb}{0.000000,0.000000,0.000000}%
\pgfsetstrokecolor{currentstroke}%
\pgfsetstrokeopacity{0.000000}%
\pgfsetdash{}{0pt}%
\pgfpathmoveto{\pgfqpoint{5.398847in}{0.935544in}}%
\pgfpathlineto{\pgfqpoint{5.395785in}{0.935101in}}%
\pgfpathlineto{\pgfqpoint{5.392896in}{0.934898in}}%
\pgfpathlineto{\pgfqpoint{5.394133in}{0.926031in}}%
\pgfpathlineto{\pgfqpoint{5.386682in}{0.917093in}}%
\pgfpathlineto{\pgfqpoint{5.385213in}{0.903107in}}%
\pgfpathlineto{\pgfqpoint{5.351904in}{0.900644in}}%
\pgfpathlineto{\pgfqpoint{5.354808in}{0.907211in}}%
\pgfpathlineto{\pgfqpoint{5.366818in}{0.919204in}}%
\pgfpathlineto{\pgfqpoint{5.367150in}{0.926155in}}%
\pgfpathlineto{\pgfqpoint{5.361833in}{0.932737in}}%
\pgfpathlineto{\pgfqpoint{5.312257in}{0.930116in}}%
\pgfpathlineto{\pgfqpoint{5.225580in}{0.927364in}}%
\pgfpathlineto{\pgfqpoint{5.174856in}{0.926440in}}%
\pgfpathlineto{\pgfqpoint{5.136515in}{0.926094in}}%
\pgfpathlineto{\pgfqpoint{5.136224in}{0.957954in}}%
\pgfpathlineto{\pgfqpoint{5.135919in}{1.050444in}}%
\pgfpathlineto{\pgfqpoint{5.136265in}{1.090665in}}%
\pgfpathlineto{\pgfqpoint{5.125905in}{1.096790in}}%
\pgfpathlineto{\pgfqpoint{5.120797in}{1.107461in}}%
\pgfpathlineto{\pgfqpoint{5.115237in}{1.112582in}}%
\pgfpathlineto{\pgfqpoint{5.114607in}{1.120348in}}%
\pgfpathlineto{\pgfqpoint{5.120414in}{1.132640in}}%
\pgfpathlineto{\pgfqpoint{5.120170in}{1.139746in}}%
\pgfpathlineto{\pgfqpoint{5.112350in}{1.138104in}}%
\pgfpathlineto{\pgfqpoint{5.102367in}{1.146251in}}%
\pgfpathlineto{\pgfqpoint{5.094364in}{1.160551in}}%
\pgfpathlineto{\pgfqpoint{5.087651in}{1.167151in}}%
\pgfpathlineto{\pgfqpoint{5.080637in}{1.183365in}}%
\pgfpathlineto{\pgfqpoint{5.153489in}{1.182229in}}%
\pgfpathlineto{\pgfqpoint{5.225906in}{1.184615in}}%
\pgfpathlineto{\pgfqpoint{5.272339in}{1.187293in}}%
\pgfpathlineto{\pgfqpoint{5.286865in}{1.172908in}}%
\pgfpathlineto{\pgfqpoint{5.284210in}{1.157906in}}%
\pgfpathlineto{\pgfqpoint{5.288083in}{1.141015in}}%
\pgfpathlineto{\pgfqpoint{5.292379in}{1.132498in}}%
\pgfpathlineto{\pgfqpoint{5.306087in}{1.120782in}}%
\pgfpathlineto{\pgfqpoint{5.309086in}{1.115607in}}%
\pgfpathlineto{\pgfqpoint{5.324094in}{1.103945in}}%
\pgfpathlineto{\pgfqpoint{5.326269in}{1.090863in}}%
\pgfpathlineto{\pgfqpoint{5.330579in}{1.081212in}}%
\pgfpathlineto{\pgfqpoint{5.340763in}{1.086834in}}%
\pgfpathlineto{\pgfqpoint{5.349157in}{1.084263in}}%
\pgfpathlineto{\pgfqpoint{5.355495in}{1.077440in}}%
\pgfpathlineto{\pgfqpoint{5.353102in}{1.065365in}}%
\pgfpathlineto{\pgfqpoint{5.345322in}{1.049578in}}%
\pgfpathlineto{\pgfqpoint{5.345442in}{1.041355in}}%
\pgfpathlineto{\pgfqpoint{5.351256in}{1.034614in}}%
\pgfpathlineto{\pgfqpoint{5.363768in}{1.025684in}}%
\pgfpathlineto{\pgfqpoint{5.378174in}{1.018436in}}%
\pgfpathlineto{\pgfqpoint{5.381591in}{1.012632in}}%
\pgfpathlineto{\pgfqpoint{5.388959in}{1.009888in}}%
\pgfpathlineto{\pgfqpoint{5.389252in}{1.002186in}}%
\pgfpathlineto{\pgfqpoint{5.394705in}{0.992402in}}%
\pgfpathlineto{\pgfqpoint{5.390922in}{0.983711in}}%
\pgfpathlineto{\pgfqpoint{5.398328in}{0.969882in}}%
\pgfpathlineto{\pgfqpoint{5.404121in}{0.972628in}}%
\pgfpathlineto{\pgfqpoint{5.411225in}{0.966856in}}%
\pgfpathlineto{\pgfqpoint{5.408837in}{0.941156in}}%
\pgfpathlineto{\pgfqpoint{5.400463in}{0.942786in}}%
\pgfpathlineto{\pgfqpoint{5.398847in}{0.935544in}}%
\pgfpathclose%
\pgfusepath{fill}%
\end{pgfscope}%
\begin{pgfscope}%
\pgfpathrectangle{\pgfqpoint{3.525000in}{0.100000in}}{\pgfqpoint{2.862985in}{1.832778in}}%
\pgfusepath{clip}%
\pgfsetbuttcap%
\pgfsetmiterjoin%
\definecolor{currentfill}{rgb}{0.740143,0.213864,0.524216}%
\pgfsetfillcolor{currentfill}%
\pgfsetlinewidth{0.000000pt}%
\definecolor{currentstroke}{rgb}{0.000000,0.000000,0.000000}%
\pgfsetstrokecolor{currentstroke}%
\pgfsetstrokeopacity{0.000000}%
\pgfsetdash{}{0pt}%
\pgfpathmoveto{\pgfqpoint{4.114569in}{0.941085in}}%
\pgfpathlineto{\pgfqpoint{4.120633in}{0.953995in}}%
\pgfpathlineto{\pgfqpoint{4.118973in}{0.973374in}}%
\pgfpathlineto{\pgfqpoint{4.121877in}{0.978886in}}%
\pgfpathlineto{\pgfqpoint{4.122164in}{1.003024in}}%
\pgfpathlineto{\pgfqpoint{4.124916in}{1.009976in}}%
\pgfpathlineto{\pgfqpoint{4.134613in}{1.011057in}}%
\pgfpathlineto{\pgfqpoint{4.143612in}{1.007971in}}%
\pgfpathlineto{\pgfqpoint{4.147531in}{0.999473in}}%
\pgfpathlineto{\pgfqpoint{4.153025in}{0.999798in}}%
\pgfpathlineto{\pgfqpoint{4.159852in}{1.009534in}}%
\pgfpathlineto{\pgfqpoint{4.169719in}{1.057523in}}%
\pgfpathlineto{\pgfqpoint{4.225868in}{1.045930in}}%
\pgfpathlineto{\pgfqpoint{4.258455in}{1.039504in}}%
\pgfpathlineto{\pgfqpoint{4.344991in}{1.024292in}}%
\pgfpathlineto{\pgfqpoint{4.368932in}{1.019484in}}%
\pgfpathlineto{\pgfqpoint{4.415915in}{1.012202in}}%
\pgfpathlineto{\pgfqpoint{4.406281in}{0.950172in}}%
\pgfpathlineto{\pgfqpoint{4.396262in}{0.885469in}}%
\pgfpathlineto{\pgfqpoint{4.384723in}{0.812676in}}%
\pgfpathlineto{\pgfqpoint{4.371747in}{0.729118in}}%
\pgfpathlineto{\pgfqpoint{4.361264in}{0.660490in}}%
\pgfpathlineto{\pgfqpoint{4.286176in}{0.672414in}}%
\pgfpathlineto{\pgfqpoint{4.253185in}{0.678064in}}%
\pgfpathlineto{\pgfqpoint{4.238447in}{0.686887in}}%
\pgfpathlineto{\pgfqpoint{4.142356in}{0.744864in}}%
\pgfpathlineto{\pgfqpoint{4.070244in}{0.788862in}}%
\pgfpathlineto{\pgfqpoint{4.072646in}{0.796660in}}%
\pgfpathlineto{\pgfqpoint{4.078603in}{0.802117in}}%
\pgfpathlineto{\pgfqpoint{4.084798in}{0.801105in}}%
\pgfpathlineto{\pgfqpoint{4.093785in}{0.806836in}}%
\pgfpathlineto{\pgfqpoint{4.095200in}{0.815071in}}%
\pgfpathlineto{\pgfqpoint{4.084240in}{0.825061in}}%
\pgfpathlineto{\pgfqpoint{4.092035in}{0.844239in}}%
\pgfpathlineto{\pgfqpoint{4.099871in}{0.851614in}}%
\pgfpathlineto{\pgfqpoint{4.103591in}{0.860364in}}%
\pgfpathlineto{\pgfqpoint{4.105882in}{0.876407in}}%
\pgfpathlineto{\pgfqpoint{4.113230in}{0.883673in}}%
\pgfpathlineto{\pgfqpoint{4.128654in}{0.890930in}}%
\pgfpathlineto{\pgfqpoint{4.128313in}{0.897324in}}%
\pgfpathlineto{\pgfqpoint{4.119699in}{0.905259in}}%
\pgfpathlineto{\pgfqpoint{4.118526in}{0.921617in}}%
\pgfpathlineto{\pgfqpoint{4.112595in}{0.933576in}}%
\pgfpathlineto{\pgfqpoint{4.114569in}{0.941085in}}%
\pgfpathclose%
\pgfusepath{fill}%
\end{pgfscope}%
\begin{pgfscope}%
\pgfpathrectangle{\pgfqpoint{3.525000in}{0.100000in}}{\pgfqpoint{2.862985in}{1.832778in}}%
\pgfusepath{clip}%
\pgfsetbuttcap%
\pgfsetmiterjoin%
\definecolor{currentfill}{rgb}{0.748289,0.222711,0.516834}%
\pgfsetfillcolor{currentfill}%
\pgfsetlinewidth{0.000000pt}%
\definecolor{currentstroke}{rgb}{0.000000,0.000000,0.000000}%
\pgfsetstrokecolor{currentstroke}%
\pgfsetstrokeopacity{0.000000}%
\pgfsetdash{}{0pt}%
\pgfpathmoveto{\pgfqpoint{5.144643in}{0.746495in}}%
\pgfpathlineto{\pgfqpoint{5.130724in}{0.750800in}}%
\pgfpathlineto{\pgfqpoint{5.112888in}{0.764459in}}%
\pgfpathlineto{\pgfqpoint{5.104966in}{0.767394in}}%
\pgfpathlineto{\pgfqpoint{5.099932in}{0.761495in}}%
\pgfpathlineto{\pgfqpoint{5.093575in}{0.761197in}}%
\pgfpathlineto{\pgfqpoint{5.085532in}{0.766205in}}%
\pgfpathlineto{\pgfqpoint{5.072911in}{0.759788in}}%
\pgfpathlineto{\pgfqpoint{5.068418in}{0.762948in}}%
\pgfpathlineto{\pgfqpoint{5.055999in}{0.755462in}}%
\pgfpathlineto{\pgfqpoint{5.042897in}{0.756843in}}%
\pgfpathlineto{\pgfqpoint{5.032791in}{0.761708in}}%
\pgfpathlineto{\pgfqpoint{5.025714in}{0.759877in}}%
\pgfpathlineto{\pgfqpoint{5.014386in}{0.766758in}}%
\pgfpathlineto{\pgfqpoint{5.008084in}{0.754622in}}%
\pgfpathlineto{\pgfqpoint{5.001640in}{0.765059in}}%
\pgfpathlineto{\pgfqpoint{4.988911in}{0.761010in}}%
\pgfpathlineto{\pgfqpoint{4.977778in}{0.770922in}}%
\pgfpathlineto{\pgfqpoint{4.968049in}{0.762933in}}%
\pgfpathlineto{\pgfqpoint{4.963177in}{0.770275in}}%
\pgfpathlineto{\pgfqpoint{4.956154in}{0.772659in}}%
\pgfpathlineto{\pgfqpoint{4.951877in}{0.779738in}}%
\pgfpathlineto{\pgfqpoint{4.942692in}{0.781753in}}%
\pgfpathlineto{\pgfqpoint{4.937392in}{0.776425in}}%
\pgfpathlineto{\pgfqpoint{4.928387in}{0.783330in}}%
\pgfpathlineto{\pgfqpoint{4.924189in}{0.781753in}}%
\pgfpathlineto{\pgfqpoint{4.909265in}{0.787349in}}%
\pgfpathlineto{\pgfqpoint{4.899917in}{0.787973in}}%
\pgfpathlineto{\pgfqpoint{4.895733in}{0.799856in}}%
\pgfpathlineto{\pgfqpoint{4.888240in}{0.798376in}}%
\pgfpathlineto{\pgfqpoint{4.879657in}{0.801308in}}%
\pgfpathlineto{\pgfqpoint{4.874005in}{0.799614in}}%
\pgfpathlineto{\pgfqpoint{4.861884in}{0.812959in}}%
\pgfpathlineto{\pgfqpoint{4.858507in}{0.812086in}}%
\pgfpathlineto{\pgfqpoint{4.861574in}{0.866207in}}%
\pgfpathlineto{\pgfqpoint{4.864916in}{0.933195in}}%
\pgfpathlineto{\pgfqpoint{4.810064in}{0.936263in}}%
\pgfpathlineto{\pgfqpoint{4.755936in}{0.940356in}}%
\pgfpathlineto{\pgfqpoint{4.714118in}{0.943987in}}%
\pgfpathlineto{\pgfqpoint{4.716895in}{0.975633in}}%
\pgfpathlineto{\pgfqpoint{4.764821in}{0.971059in}}%
\pgfpathlineto{\pgfqpoint{4.819179in}{0.967651in}}%
\pgfpathlineto{\pgfqpoint{4.894882in}{0.963187in}}%
\pgfpathlineto{\pgfqpoint{4.939818in}{0.961453in}}%
\pgfpathlineto{\pgfqpoint{4.993651in}{0.959554in}}%
\pgfpathlineto{\pgfqpoint{5.066779in}{0.958274in}}%
\pgfpathlineto{\pgfqpoint{5.136224in}{0.957954in}}%
\pgfpathlineto{\pgfqpoint{5.136515in}{0.926094in}}%
\pgfpathlineto{\pgfqpoint{5.140413in}{0.902092in}}%
\pgfpathlineto{\pgfqpoint{5.146468in}{0.857760in}}%
\pgfpathlineto{\pgfqpoint{5.145219in}{0.782052in}}%
\pgfpathlineto{\pgfqpoint{5.144643in}{0.746495in}}%
\pgfpathclose%
\pgfusepath{fill}%
\end{pgfscope}%
\begin{pgfscope}%
\pgfpathrectangle{\pgfqpoint{3.525000in}{0.100000in}}{\pgfqpoint{2.862985in}{1.832778in}}%
\pgfusepath{clip}%
\pgfsetbuttcap%
\pgfsetmiterjoin%
\definecolor{currentfill}{rgb}{0.783383,0.262500,0.483918}%
\pgfsetfillcolor{currentfill}%
\pgfsetlinewidth{0.000000pt}%
\definecolor{currentstroke}{rgb}{0.000000,0.000000,0.000000}%
\pgfsetstrokecolor{currentstroke}%
\pgfsetstrokeopacity{0.000000}%
\pgfsetdash{}{0pt}%
\pgfpathmoveto{\pgfqpoint{5.665820in}{0.864070in}}%
\pgfpathlineto{\pgfqpoint{5.665861in}{0.878097in}}%
\pgfpathlineto{\pgfqpoint{5.678051in}{0.883484in}}%
\pgfpathlineto{\pgfqpoint{5.678562in}{0.892095in}}%
\pgfpathlineto{\pgfqpoint{5.689476in}{0.902614in}}%
\pgfpathlineto{\pgfqpoint{5.703119in}{0.904768in}}%
\pgfpathlineto{\pgfqpoint{5.715113in}{0.916361in}}%
\pgfpathlineto{\pgfqpoint{5.727639in}{0.921543in}}%
\pgfpathlineto{\pgfqpoint{5.737040in}{0.937098in}}%
\pgfpathlineto{\pgfqpoint{5.745602in}{0.935969in}}%
\pgfpathlineto{\pgfqpoint{5.763697in}{0.950142in}}%
\pgfpathlineto{\pgfqpoint{5.773253in}{0.950409in}}%
\pgfpathlineto{\pgfqpoint{5.777288in}{0.961208in}}%
\pgfpathlineto{\pgfqpoint{5.784883in}{0.968727in}}%
\pgfpathlineto{\pgfqpoint{5.786621in}{0.981538in}}%
\pgfpathlineto{\pgfqpoint{5.802964in}{0.982918in}}%
\pgfpathlineto{\pgfqpoint{5.869033in}{0.991644in}}%
\pgfpathlineto{\pgfqpoint{5.946826in}{1.005045in}}%
\pgfpathlineto{\pgfqpoint{6.003091in}{1.016561in}}%
\pgfpathlineto{\pgfqpoint{6.065997in}{1.029610in}}%
\pgfpathlineto{\pgfqpoint{6.083973in}{1.004931in}}%
\pgfpathlineto{\pgfqpoint{6.074236in}{1.006506in}}%
\pgfpathlineto{\pgfqpoint{6.061737in}{1.003461in}}%
\pgfpathlineto{\pgfqpoint{6.049421in}{0.990884in}}%
\pgfpathlineto{\pgfqpoint{6.040570in}{0.991789in}}%
\pgfpathlineto{\pgfqpoint{6.039450in}{0.984319in}}%
\pgfpathlineto{\pgfqpoint{6.055449in}{0.990226in}}%
\pgfpathlineto{\pgfqpoint{6.071548in}{0.992662in}}%
\pgfpathlineto{\pgfqpoint{6.077526in}{0.989421in}}%
\pgfpathlineto{\pgfqpoint{6.085571in}{0.993116in}}%
\pgfpathlineto{\pgfqpoint{6.089722in}{0.990527in}}%
\pgfpathlineto{\pgfqpoint{6.093396in}{0.978199in}}%
\pgfpathlineto{\pgfqpoint{6.085750in}{0.974382in}}%
\pgfpathlineto{\pgfqpoint{6.080502in}{0.959350in}}%
\pgfpathlineto{\pgfqpoint{6.075000in}{0.953497in}}%
\pgfpathlineto{\pgfqpoint{6.058132in}{0.954777in}}%
\pgfpathlineto{\pgfqpoint{6.059039in}{0.963604in}}%
\pgfpathlineto{\pgfqpoint{6.049827in}{0.959748in}}%
\pgfpathlineto{\pgfqpoint{6.047823in}{0.952387in}}%
\pgfpathlineto{\pgfqpoint{6.054137in}{0.944761in}}%
\pgfpathlineto{\pgfqpoint{6.056282in}{0.931819in}}%
\pgfpathlineto{\pgfqpoint{6.045950in}{0.924444in}}%
\pgfpathlineto{\pgfqpoint{6.057125in}{0.921851in}}%
\pgfpathlineto{\pgfqpoint{6.062181in}{0.927273in}}%
\pgfpathlineto{\pgfqpoint{6.073362in}{0.927934in}}%
\pgfpathlineto{\pgfqpoint{6.067616in}{0.916224in}}%
\pgfpathlineto{\pgfqpoint{6.060579in}{0.910582in}}%
\pgfpathlineto{\pgfqpoint{6.039292in}{0.904933in}}%
\pgfpathlineto{\pgfqpoint{6.017436in}{0.885111in}}%
\pgfpathlineto{\pgfqpoint{6.003960in}{0.865580in}}%
\pgfpathlineto{\pgfqpoint{6.003880in}{0.857648in}}%
\pgfpathlineto{\pgfqpoint{5.998497in}{0.846702in}}%
\pgfpathlineto{\pgfqpoint{5.970869in}{0.839501in}}%
\pgfpathlineto{\pgfqpoint{5.905165in}{0.886600in}}%
\pgfpathlineto{\pgfqpoint{5.847304in}{0.878098in}}%
\pgfpathlineto{\pgfqpoint{5.846819in}{0.885949in}}%
\pgfpathlineto{\pgfqpoint{5.838024in}{0.894789in}}%
\pgfpathlineto{\pgfqpoint{5.831366in}{0.896952in}}%
\pgfpathlineto{\pgfqpoint{5.768499in}{0.890420in}}%
\pgfpathlineto{\pgfqpoint{5.754049in}{0.885519in}}%
\pgfpathlineto{\pgfqpoint{5.727958in}{0.872540in}}%
\pgfpathlineto{\pgfqpoint{5.705410in}{0.868948in}}%
\pgfpathlineto{\pgfqpoint{5.665820in}{0.864070in}}%
\pgfpathclose%
\pgfusepath{fill}%
\end{pgfscope}%
\begin{pgfscope}%
\pgfpathrectangle{\pgfqpoint{3.525000in}{0.100000in}}{\pgfqpoint{2.862985in}{1.832778in}}%
\pgfusepath{clip}%
\pgfsetbuttcap%
\pgfsetmiterjoin%
\definecolor{currentfill}{rgb}{0.819651,0.306812,0.448306}%
\pgfsetfillcolor{currentfill}%
\pgfsetlinewidth{0.000000pt}%
\definecolor{currentstroke}{rgb}{0.000000,0.000000,0.000000}%
\pgfsetstrokecolor{currentstroke}%
\pgfsetstrokeopacity{0.000000}%
\pgfsetdash{}{0pt}%
\pgfpathmoveto{\pgfqpoint{5.665820in}{0.864070in}}%
\pgfpathlineto{\pgfqpoint{5.600022in}{0.856843in}}%
\pgfpathlineto{\pgfqpoint{5.539830in}{0.851430in}}%
\pgfpathlineto{\pgfqpoint{5.477965in}{0.847429in}}%
\pgfpathlineto{\pgfqpoint{5.467346in}{0.845890in}}%
\pgfpathlineto{\pgfqpoint{5.425643in}{0.842683in}}%
\pgfpathlineto{\pgfqpoint{5.358848in}{0.838785in}}%
\pgfpathlineto{\pgfqpoint{5.370732in}{0.848649in}}%
\pgfpathlineto{\pgfqpoint{5.368020in}{0.859033in}}%
\pgfpathlineto{\pgfqpoint{5.370647in}{0.871619in}}%
\pgfpathlineto{\pgfqpoint{5.374660in}{0.877545in}}%
\pgfpathlineto{\pgfqpoint{5.374509in}{0.885789in}}%
\pgfpathlineto{\pgfqpoint{5.385195in}{0.890974in}}%
\pgfpathlineto{\pgfqpoint{5.385213in}{0.903107in}}%
\pgfpathlineto{\pgfqpoint{5.386682in}{0.917093in}}%
\pgfpathlineto{\pgfqpoint{5.394133in}{0.926031in}}%
\pgfpathlineto{\pgfqpoint{5.392896in}{0.934898in}}%
\pgfpathlineto{\pgfqpoint{5.395785in}{0.935101in}}%
\pgfpathlineto{\pgfqpoint{5.398847in}{0.935544in}}%
\pgfpathlineto{\pgfqpoint{5.467942in}{0.939936in}}%
\pgfpathlineto{\pgfqpoint{5.465629in}{0.951286in}}%
\pgfpathlineto{\pgfqpoint{5.476748in}{0.949746in}}%
\pgfpathlineto{\pgfqpoint{5.544049in}{0.956404in}}%
\pgfpathlineto{\pgfqpoint{5.598070in}{0.959461in}}%
\pgfpathlineto{\pgfqpoint{5.630991in}{0.962230in}}%
\pgfpathlineto{\pgfqpoint{5.685505in}{0.967399in}}%
\pgfpathlineto{\pgfqpoint{5.686414in}{0.968427in}}%
\pgfpathlineto{\pgfqpoint{5.709707in}{0.970734in}}%
\pgfpathlineto{\pgfqpoint{5.786621in}{0.981538in}}%
\pgfpathlineto{\pgfqpoint{5.784883in}{0.968727in}}%
\pgfpathlineto{\pgfqpoint{5.777288in}{0.961208in}}%
\pgfpathlineto{\pgfqpoint{5.773253in}{0.950409in}}%
\pgfpathlineto{\pgfqpoint{5.763697in}{0.950142in}}%
\pgfpathlineto{\pgfqpoint{5.745602in}{0.935969in}}%
\pgfpathlineto{\pgfqpoint{5.737040in}{0.937098in}}%
\pgfpathlineto{\pgfqpoint{5.727639in}{0.921543in}}%
\pgfpathlineto{\pgfqpoint{5.715113in}{0.916361in}}%
\pgfpathlineto{\pgfqpoint{5.703119in}{0.904768in}}%
\pgfpathlineto{\pgfqpoint{5.689476in}{0.902614in}}%
\pgfpathlineto{\pgfqpoint{5.678562in}{0.892095in}}%
\pgfpathlineto{\pgfqpoint{5.678051in}{0.883484in}}%
\pgfpathlineto{\pgfqpoint{5.665861in}{0.878097in}}%
\pgfpathlineto{\pgfqpoint{5.665820in}{0.864070in}}%
\pgfpathclose%
\pgfusepath{fill}%
\end{pgfscope}%
\begin{pgfscope}%
\pgfpathrectangle{\pgfqpoint{3.525000in}{0.100000in}}{\pgfqpoint{2.862985in}{1.832778in}}%
\pgfusepath{clip}%
\pgfsetbuttcap%
\pgfsetmiterjoin%
\definecolor{currentfill}{rgb}{0.679160,0.151848,0.575189}%
\pgfsetfillcolor{currentfill}%
\pgfsetlinewidth{0.000000pt}%
\definecolor{currentstroke}{rgb}{0.000000,0.000000,0.000000}%
\pgfsetstrokecolor{currentstroke}%
\pgfsetstrokeopacity{0.000000}%
\pgfsetdash{}{0pt}%
\pgfpathmoveto{\pgfqpoint{4.498883in}{0.669271in}}%
\pgfpathlineto{\pgfqpoint{4.494190in}{0.676811in}}%
\pgfpathlineto{\pgfqpoint{4.496132in}{0.683307in}}%
\pgfpathlineto{\pgfqpoint{4.529118in}{0.679215in}}%
\pgfpathlineto{\pgfqpoint{4.590452in}{0.672249in}}%
\pgfpathlineto{\pgfqpoint{4.634801in}{0.667891in}}%
\pgfpathlineto{\pgfqpoint{4.686021in}{0.662722in}}%
\pgfpathlineto{\pgfqpoint{4.688807in}{0.695029in}}%
\pgfpathlineto{\pgfqpoint{4.697520in}{0.777201in}}%
\pgfpathlineto{\pgfqpoint{4.703172in}{0.834767in}}%
\pgfpathlineto{\pgfqpoint{4.708057in}{0.889756in}}%
\pgfpathlineto{\pgfqpoint{4.712612in}{0.944071in}}%
\pgfpathlineto{\pgfqpoint{4.714118in}{0.943987in}}%
\pgfpathlineto{\pgfqpoint{4.755936in}{0.940356in}}%
\pgfpathlineto{\pgfqpoint{4.810064in}{0.936263in}}%
\pgfpathlineto{\pgfqpoint{4.864916in}{0.933195in}}%
\pgfpathlineto{\pgfqpoint{4.861574in}{0.866207in}}%
\pgfpathlineto{\pgfqpoint{4.858507in}{0.812086in}}%
\pgfpathlineto{\pgfqpoint{4.861884in}{0.812959in}}%
\pgfpathlineto{\pgfqpoint{4.874005in}{0.799614in}}%
\pgfpathlineto{\pgfqpoint{4.879657in}{0.801308in}}%
\pgfpathlineto{\pgfqpoint{4.888240in}{0.798376in}}%
\pgfpathlineto{\pgfqpoint{4.895733in}{0.799856in}}%
\pgfpathlineto{\pgfqpoint{4.899917in}{0.787973in}}%
\pgfpathlineto{\pgfqpoint{4.909265in}{0.787349in}}%
\pgfpathlineto{\pgfqpoint{4.924189in}{0.781753in}}%
\pgfpathlineto{\pgfqpoint{4.928387in}{0.783330in}}%
\pgfpathlineto{\pgfqpoint{4.937392in}{0.776425in}}%
\pgfpathlineto{\pgfqpoint{4.942692in}{0.781753in}}%
\pgfpathlineto{\pgfqpoint{4.951877in}{0.779738in}}%
\pgfpathlineto{\pgfqpoint{4.956154in}{0.772659in}}%
\pgfpathlineto{\pgfqpoint{4.963177in}{0.770275in}}%
\pgfpathlineto{\pgfqpoint{4.968049in}{0.762933in}}%
\pgfpathlineto{\pgfqpoint{4.977778in}{0.770922in}}%
\pgfpathlineto{\pgfqpoint{4.988911in}{0.761010in}}%
\pgfpathlineto{\pgfqpoint{5.001640in}{0.765059in}}%
\pgfpathlineto{\pgfqpoint{5.008084in}{0.754622in}}%
\pgfpathlineto{\pgfqpoint{5.014386in}{0.766758in}}%
\pgfpathlineto{\pgfqpoint{5.025714in}{0.759877in}}%
\pgfpathlineto{\pgfqpoint{5.032791in}{0.761708in}}%
\pgfpathlineto{\pgfqpoint{5.042897in}{0.756843in}}%
\pgfpathlineto{\pgfqpoint{5.055999in}{0.755462in}}%
\pgfpathlineto{\pgfqpoint{5.068418in}{0.762948in}}%
\pgfpathlineto{\pgfqpoint{5.072911in}{0.759788in}}%
\pgfpathlineto{\pgfqpoint{5.085532in}{0.766205in}}%
\pgfpathlineto{\pgfqpoint{5.093575in}{0.761197in}}%
\pgfpathlineto{\pgfqpoint{5.099932in}{0.761495in}}%
\pgfpathlineto{\pgfqpoint{5.104966in}{0.767394in}}%
\pgfpathlineto{\pgfqpoint{5.112888in}{0.764459in}}%
\pgfpathlineto{\pgfqpoint{5.130724in}{0.750800in}}%
\pgfpathlineto{\pgfqpoint{5.144643in}{0.746495in}}%
\pgfpathlineto{\pgfqpoint{5.150225in}{0.741221in}}%
\pgfpathlineto{\pgfqpoint{5.157212in}{0.744093in}}%
\pgfpathlineto{\pgfqpoint{5.167799in}{0.741894in}}%
\pgfpathlineto{\pgfqpoint{5.168006in}{0.708326in}}%
\pgfpathlineto{\pgfqpoint{5.168890in}{0.643403in}}%
\pgfpathlineto{\pgfqpoint{5.176243in}{0.637164in}}%
\pgfpathlineto{\pgfqpoint{5.182137in}{0.625668in}}%
\pgfpathlineto{\pgfqpoint{5.180065in}{0.617977in}}%
\pgfpathlineto{\pgfqpoint{5.184693in}{0.614721in}}%
\pgfpathlineto{\pgfqpoint{5.188454in}{0.600538in}}%
\pgfpathlineto{\pgfqpoint{5.195658in}{0.592942in}}%
\pgfpathlineto{\pgfqpoint{5.197272in}{0.576816in}}%
\pgfpathlineto{\pgfqpoint{5.196019in}{0.569988in}}%
\pgfpathlineto{\pgfqpoint{5.186227in}{0.551919in}}%
\pgfpathlineto{\pgfqpoint{5.185101in}{0.538927in}}%
\pgfpathlineto{\pgfqpoint{5.188423in}{0.536215in}}%
\pgfpathlineto{\pgfqpoint{5.188712in}{0.524834in}}%
\pgfpathlineto{\pgfqpoint{5.185335in}{0.517680in}}%
\pgfpathlineto{\pgfqpoint{5.180085in}{0.514985in}}%
\pgfpathlineto{\pgfqpoint{5.174944in}{0.505637in}}%
\pgfpathlineto{\pgfqpoint{5.181496in}{0.496583in}}%
\pgfpathlineto{\pgfqpoint{5.168777in}{0.496405in}}%
\pgfpathlineto{\pgfqpoint{5.134766in}{0.480844in}}%
\pgfpathlineto{\pgfqpoint{5.141259in}{0.490157in}}%
\pgfpathlineto{\pgfqpoint{5.133396in}{0.495178in}}%
\pgfpathlineto{\pgfqpoint{5.131753in}{0.503714in}}%
\pgfpathlineto{\pgfqpoint{5.116406in}{0.488845in}}%
\pgfpathlineto{\pgfqpoint{5.123217in}{0.478682in}}%
\pgfpathlineto{\pgfqpoint{5.113501in}{0.465742in}}%
\pgfpathlineto{\pgfqpoint{5.108273in}{0.466013in}}%
\pgfpathlineto{\pgfqpoint{5.103354in}{0.451920in}}%
\pgfpathlineto{\pgfqpoint{5.087798in}{0.440835in}}%
\pgfpathlineto{\pgfqpoint{5.073265in}{0.436840in}}%
\pgfpathlineto{\pgfqpoint{5.047917in}{0.426447in}}%
\pgfpathlineto{\pgfqpoint{5.045291in}{0.432244in}}%
\pgfpathlineto{\pgfqpoint{5.029501in}{0.426908in}}%
\pgfpathlineto{\pgfqpoint{5.039213in}{0.417822in}}%
\pgfpathlineto{\pgfqpoint{5.023894in}{0.409929in}}%
\pgfpathlineto{\pgfqpoint{5.007371in}{0.398011in}}%
\pgfpathlineto{\pgfqpoint{5.003403in}{0.403551in}}%
\pgfpathlineto{\pgfqpoint{4.989881in}{0.395248in}}%
\pgfpathlineto{\pgfqpoint{5.003135in}{0.392095in}}%
\pgfpathlineto{\pgfqpoint{4.993142in}{0.379039in}}%
\pgfpathlineto{\pgfqpoint{4.988264in}{0.382924in}}%
\pgfpathlineto{\pgfqpoint{4.981713in}{0.376684in}}%
\pgfpathlineto{\pgfqpoint{4.989878in}{0.371233in}}%
\pgfpathlineto{\pgfqpoint{4.985050in}{0.363265in}}%
\pgfpathlineto{\pgfqpoint{4.979117in}{0.344404in}}%
\pgfpathlineto{\pgfqpoint{4.970559in}{0.326238in}}%
\pgfpathlineto{\pgfqpoint{4.974401in}{0.314331in}}%
\pgfpathlineto{\pgfqpoint{4.980932in}{0.286258in}}%
\pgfpathlineto{\pgfqpoint{4.981520in}{0.274896in}}%
\pgfpathlineto{\pgfqpoint{4.991610in}{0.259994in}}%
\pgfpathlineto{\pgfqpoint{4.983829in}{0.260865in}}%
\pgfpathlineto{\pgfqpoint{4.976272in}{0.253347in}}%
\pgfpathlineto{\pgfqpoint{4.967420in}{0.259256in}}%
\pgfpathlineto{\pgfqpoint{4.964243in}{0.265150in}}%
\pgfpathlineto{\pgfqpoint{4.951645in}{0.267892in}}%
\pgfpathlineto{\pgfqpoint{4.932402in}{0.268229in}}%
\pgfpathlineto{\pgfqpoint{4.918224in}{0.279391in}}%
\pgfpathlineto{\pgfqpoint{4.905350in}{0.281256in}}%
\pgfpathlineto{\pgfqpoint{4.897540in}{0.290127in}}%
\pgfpathlineto{\pgfqpoint{4.881183in}{0.293697in}}%
\pgfpathlineto{\pgfqpoint{4.872239in}{0.322221in}}%
\pgfpathlineto{\pgfqpoint{4.863094in}{0.333648in}}%
\pgfpathlineto{\pgfqpoint{4.864648in}{0.344506in}}%
\pgfpathlineto{\pgfqpoint{4.858978in}{0.352440in}}%
\pgfpathlineto{\pgfqpoint{4.862536in}{0.363286in}}%
\pgfpathlineto{\pgfqpoint{4.859601in}{0.371234in}}%
\pgfpathlineto{\pgfqpoint{4.850411in}{0.374834in}}%
\pgfpathlineto{\pgfqpoint{4.841818in}{0.383999in}}%
\pgfpathlineto{\pgfqpoint{4.838687in}{0.396269in}}%
\pgfpathlineto{\pgfqpoint{4.830556in}{0.407420in}}%
\pgfpathlineto{\pgfqpoint{4.819734in}{0.416096in}}%
\pgfpathlineto{\pgfqpoint{4.809986in}{0.440984in}}%
\pgfpathlineto{\pgfqpoint{4.802115in}{0.468203in}}%
\pgfpathlineto{\pgfqpoint{4.795657in}{0.478966in}}%
\pgfpathlineto{\pgfqpoint{4.784460in}{0.488026in}}%
\pgfpathlineto{\pgfqpoint{4.781662in}{0.494603in}}%
\pgfpathlineto{\pgfqpoint{4.771179in}{0.498686in}}%
\pgfpathlineto{\pgfqpoint{4.760831in}{0.516112in}}%
\pgfpathlineto{\pgfqpoint{4.751382in}{0.514774in}}%
\pgfpathlineto{\pgfqpoint{4.741787in}{0.519126in}}%
\pgfpathlineto{\pgfqpoint{4.728186in}{0.518285in}}%
\pgfpathlineto{\pgfqpoint{4.714369in}{0.525467in}}%
\pgfpathlineto{\pgfqpoint{4.710472in}{0.518637in}}%
\pgfpathlineto{\pgfqpoint{4.694325in}{0.518466in}}%
\pgfpathlineto{\pgfqpoint{4.686104in}{0.505517in}}%
\pgfpathlineto{\pgfqpoint{4.678954in}{0.489484in}}%
\pgfpathlineto{\pgfqpoint{4.663750in}{0.472274in}}%
\pgfpathlineto{\pgfqpoint{4.646524in}{0.479827in}}%
\pgfpathlineto{\pgfqpoint{4.644080in}{0.484818in}}%
\pgfpathlineto{\pgfqpoint{4.633611in}{0.488625in}}%
\pgfpathlineto{\pgfqpoint{4.630396in}{0.493833in}}%
\pgfpathlineto{\pgfqpoint{4.616499in}{0.499090in}}%
\pgfpathlineto{\pgfqpoint{4.608725in}{0.509836in}}%
\pgfpathlineto{\pgfqpoint{4.599646in}{0.515021in}}%
\pgfpathlineto{\pgfqpoint{4.591834in}{0.524068in}}%
\pgfpathlineto{\pgfqpoint{4.585749in}{0.539391in}}%
\pgfpathlineto{\pgfqpoint{4.586430in}{0.560329in}}%
\pgfpathlineto{\pgfqpoint{4.579291in}{0.570909in}}%
\pgfpathlineto{\pgfqpoint{4.578469in}{0.582369in}}%
\pgfpathlineto{\pgfqpoint{4.562630in}{0.599527in}}%
\pgfpathlineto{\pgfqpoint{4.553414in}{0.603189in}}%
\pgfpathlineto{\pgfqpoint{4.543634in}{0.619199in}}%
\pgfpathlineto{\pgfqpoint{4.535300in}{0.625560in}}%
\pgfpathlineto{\pgfqpoint{4.524711in}{0.641049in}}%
\pgfpathlineto{\pgfqpoint{4.513876in}{0.647753in}}%
\pgfpathlineto{\pgfqpoint{4.506780in}{0.664929in}}%
\pgfpathlineto{\pgfqpoint{4.498883in}{0.669271in}}%
\pgfpathclose%
\pgfusepath{fill}%
\end{pgfscope}%
\begin{pgfscope}%
\pgfpathrectangle{\pgfqpoint{3.525000in}{0.100000in}}{\pgfqpoint{2.862985in}{1.832778in}}%
\pgfusepath{clip}%
\pgfsetbuttcap%
\pgfsetmiterjoin%
\definecolor{currentfill}{rgb}{0.875376,0.383347,0.389976}%
\pgfsetfillcolor{currentfill}%
\pgfsetlinewidth{0.000000pt}%
\definecolor{currentstroke}{rgb}{0.000000,0.000000,0.000000}%
\pgfsetstrokecolor{currentstroke}%
\pgfsetstrokeopacity{0.000000}%
\pgfsetdash{}{0pt}%
\pgfpathmoveto{\pgfqpoint{4.415915in}{1.012202in}}%
\pgfpathlineto{\pgfqpoint{4.494080in}{1.000850in}}%
\pgfpathlineto{\pgfqpoint{4.524362in}{0.996184in}}%
\pgfpathlineto{\pgfqpoint{4.609722in}{0.986032in}}%
\pgfpathlineto{\pgfqpoint{4.667175in}{0.980039in}}%
\pgfpathlineto{\pgfqpoint{4.716895in}{0.975633in}}%
\pgfpathlineto{\pgfqpoint{4.714118in}{0.943987in}}%
\pgfpathlineto{\pgfqpoint{4.712612in}{0.944071in}}%
\pgfpathlineto{\pgfqpoint{4.708057in}{0.889756in}}%
\pgfpathlineto{\pgfqpoint{4.703172in}{0.834767in}}%
\pgfpathlineto{\pgfqpoint{4.697520in}{0.777201in}}%
\pgfpathlineto{\pgfqpoint{4.688807in}{0.695029in}}%
\pgfpathlineto{\pgfqpoint{4.686021in}{0.662722in}}%
\pgfpathlineto{\pgfqpoint{4.634801in}{0.667891in}}%
\pgfpathlineto{\pgfqpoint{4.590452in}{0.672249in}}%
\pgfpathlineto{\pgfqpoint{4.529118in}{0.679215in}}%
\pgfpathlineto{\pgfqpoint{4.496132in}{0.683307in}}%
\pgfpathlineto{\pgfqpoint{4.494190in}{0.676811in}}%
\pgfpathlineto{\pgfqpoint{4.498883in}{0.669271in}}%
\pgfpathlineto{\pgfqpoint{4.459240in}{0.674419in}}%
\pgfpathlineto{\pgfqpoint{4.410330in}{0.681439in}}%
\pgfpathlineto{\pgfqpoint{4.405878in}{0.653775in}}%
\pgfpathlineto{\pgfqpoint{4.361264in}{0.660490in}}%
\pgfpathlineto{\pgfqpoint{4.371747in}{0.729118in}}%
\pgfpathlineto{\pgfqpoint{4.384723in}{0.812676in}}%
\pgfpathlineto{\pgfqpoint{4.396262in}{0.885469in}}%
\pgfpathlineto{\pgfqpoint{4.406281in}{0.950172in}}%
\pgfpathlineto{\pgfqpoint{4.415915in}{1.012202in}}%
\pgfpathclose%
\pgfusepath{fill}%
\end{pgfscope}%
\begin{pgfscope}%
\pgfpathrectangle{\pgfqpoint{3.525000in}{0.100000in}}{\pgfqpoint{2.862985in}{1.832778in}}%
\pgfusepath{clip}%
\pgfsetbuttcap%
\pgfsetmiterjoin%
\definecolor{currentfill}{rgb}{0.856547,0.356066,0.410322}%
\pgfsetfillcolor{currentfill}%
\pgfsetlinewidth{0.000000pt}%
\definecolor{currentstroke}{rgb}{0.000000,0.000000,0.000000}%
\pgfsetstrokecolor{currentstroke}%
\pgfsetstrokeopacity{0.000000}%
\pgfsetdash{}{0pt}%
\pgfpathmoveto{\pgfqpoint{5.658749in}{0.610119in}}%
\pgfpathlineto{\pgfqpoint{5.584178in}{0.601837in}}%
\pgfpathlineto{\pgfqpoint{5.518446in}{0.596734in}}%
\pgfpathlineto{\pgfqpoint{5.517624in}{0.588678in}}%
\pgfpathlineto{\pgfqpoint{5.523661in}{0.581011in}}%
\pgfpathlineto{\pgfqpoint{5.531031in}{0.576507in}}%
\pgfpathlineto{\pgfqpoint{5.530919in}{0.564647in}}%
\pgfpathlineto{\pgfqpoint{5.528966in}{0.556723in}}%
\pgfpathlineto{\pgfqpoint{5.522462in}{0.551013in}}%
\pgfpathlineto{\pgfqpoint{5.513396in}{0.551618in}}%
\pgfpathlineto{\pgfqpoint{5.504825in}{0.558693in}}%
\pgfpathlineto{\pgfqpoint{5.503297in}{0.571283in}}%
\pgfpathlineto{\pgfqpoint{5.496913in}{0.578606in}}%
\pgfpathlineto{\pgfqpoint{5.492568in}{0.552398in}}%
\pgfpathlineto{\pgfqpoint{5.477803in}{0.554889in}}%
\pgfpathlineto{\pgfqpoint{5.472669in}{0.600717in}}%
\pgfpathlineto{\pgfqpoint{5.467108in}{0.649015in}}%
\pgfpathlineto{\pgfqpoint{5.469173in}{0.718697in}}%
\pgfpathlineto{\pgfqpoint{5.470455in}{0.766513in}}%
\pgfpathlineto{\pgfqpoint{5.473181in}{0.839470in}}%
\pgfpathlineto{\pgfqpoint{5.467346in}{0.845890in}}%
\pgfpathlineto{\pgfqpoint{5.477965in}{0.847429in}}%
\pgfpathlineto{\pgfqpoint{5.539830in}{0.851430in}}%
\pgfpathlineto{\pgfqpoint{5.600022in}{0.856843in}}%
\pgfpathlineto{\pgfqpoint{5.615858in}{0.801373in}}%
\pgfpathlineto{\pgfqpoint{5.626658in}{0.760657in}}%
\pgfpathlineto{\pgfqpoint{5.636339in}{0.726571in}}%
\pgfpathlineto{\pgfqpoint{5.641938in}{0.712845in}}%
\pgfpathlineto{\pgfqpoint{5.650753in}{0.700073in}}%
\pgfpathlineto{\pgfqpoint{5.649330in}{0.693573in}}%
\pgfpathlineto{\pgfqpoint{5.655627in}{0.690406in}}%
\pgfpathlineto{\pgfqpoint{5.648185in}{0.680553in}}%
\pgfpathlineto{\pgfqpoint{5.645672in}{0.663006in}}%
\pgfpathlineto{\pgfqpoint{5.652945in}{0.642483in}}%
\pgfpathlineto{\pgfqpoint{5.651935in}{0.621831in}}%
\pgfpathlineto{\pgfqpoint{5.658749in}{0.610119in}}%
\pgfpathclose%
\pgfusepath{fill}%
\end{pgfscope}%
\begin{pgfscope}%
\pgfpathrectangle{\pgfqpoint{3.525000in}{0.100000in}}{\pgfqpoint{2.862985in}{1.832778in}}%
\pgfusepath{clip}%
\pgfsetbuttcap%
\pgfsetmiterjoin%
\definecolor{currentfill}{rgb}{0.928329,0.472975,0.326067}%
\pgfsetfillcolor{currentfill}%
\pgfsetlinewidth{0.000000pt}%
\definecolor{currentstroke}{rgb}{0.000000,0.000000,0.000000}%
\pgfsetstrokecolor{currentstroke}%
\pgfsetstrokeopacity{0.000000}%
\pgfsetdash{}{0pt}%
\pgfpathmoveto{\pgfqpoint{5.477803in}{0.554889in}}%
\pgfpathlineto{\pgfqpoint{5.462652in}{0.550553in}}%
\pgfpathlineto{\pgfqpoint{5.451358in}{0.553337in}}%
\pgfpathlineto{\pgfqpoint{5.430384in}{0.546666in}}%
\pgfpathlineto{\pgfqpoint{5.414565in}{0.538074in}}%
\pgfpathlineto{\pgfqpoint{5.408075in}{0.553600in}}%
\pgfpathlineto{\pgfqpoint{5.401384in}{0.560107in}}%
\pgfpathlineto{\pgfqpoint{5.398157in}{0.567894in}}%
\pgfpathlineto{\pgfqpoint{5.402888in}{0.588972in}}%
\pgfpathlineto{\pgfqpoint{5.358060in}{0.586211in}}%
\pgfpathlineto{\pgfqpoint{5.299932in}{0.583691in}}%
\pgfpathlineto{\pgfqpoint{5.303390in}{0.588938in}}%
\pgfpathlineto{\pgfqpoint{5.299116in}{0.598847in}}%
\pgfpathlineto{\pgfqpoint{5.303375in}{0.600866in}}%
\pgfpathlineto{\pgfqpoint{5.302430in}{0.610386in}}%
\pgfpathlineto{\pgfqpoint{5.309854in}{0.619619in}}%
\pgfpathlineto{\pgfqpoint{5.313916in}{0.637542in}}%
\pgfpathlineto{\pgfqpoint{5.327503in}{0.649350in}}%
\pgfpathlineto{\pgfqpoint{5.322467in}{0.660818in}}%
\pgfpathlineto{\pgfqpoint{5.332102in}{0.667419in}}%
\pgfpathlineto{\pgfqpoint{5.323791in}{0.679370in}}%
\pgfpathlineto{\pgfqpoint{5.317776in}{0.706727in}}%
\pgfpathlineto{\pgfqpoint{5.320057in}{0.711695in}}%
\pgfpathlineto{\pgfqpoint{5.323659in}{0.721223in}}%
\pgfpathlineto{\pgfqpoint{5.320307in}{0.731223in}}%
\pgfpathlineto{\pgfqpoint{5.321946in}{0.735772in}}%
\pgfpathlineto{\pgfqpoint{5.316048in}{0.752960in}}%
\pgfpathlineto{\pgfqpoint{5.333459in}{0.783821in}}%
\pgfpathlineto{\pgfqpoint{5.334269in}{0.792035in}}%
\pgfpathlineto{\pgfqpoint{5.342656in}{0.798007in}}%
\pgfpathlineto{\pgfqpoint{5.351605in}{0.817725in}}%
\pgfpathlineto{\pgfqpoint{5.351180in}{0.825739in}}%
\pgfpathlineto{\pgfqpoint{5.362326in}{0.833932in}}%
\pgfpathlineto{\pgfqpoint{5.358848in}{0.838785in}}%
\pgfpathlineto{\pgfqpoint{5.425643in}{0.842683in}}%
\pgfpathlineto{\pgfqpoint{5.467346in}{0.845890in}}%
\pgfpathlineto{\pgfqpoint{5.473181in}{0.839470in}}%
\pgfpathlineto{\pgfqpoint{5.470455in}{0.766513in}}%
\pgfpathlineto{\pgfqpoint{5.469173in}{0.718697in}}%
\pgfpathlineto{\pgfqpoint{5.467108in}{0.649015in}}%
\pgfpathlineto{\pgfqpoint{5.472669in}{0.600717in}}%
\pgfpathlineto{\pgfqpoint{5.477803in}{0.554889in}}%
\pgfpathclose%
\pgfusepath{fill}%
\end{pgfscope}%
\begin{pgfscope}%
\pgfpathrectangle{\pgfqpoint{3.525000in}{0.100000in}}{\pgfqpoint{2.862985in}{1.832778in}}%
\pgfusepath{clip}%
\pgfsetbuttcap%
\pgfsetmiterjoin%
\definecolor{currentfill}{rgb}{0.744232,0.218288,0.520524}%
\pgfsetfillcolor{currentfill}%
\pgfsetlinewidth{0.000000pt}%
\definecolor{currentstroke}{rgb}{0.000000,0.000000,0.000000}%
\pgfsetstrokecolor{currentstroke}%
\pgfsetstrokeopacity{0.000000}%
\pgfsetdash{}{0pt}%
\pgfpathmoveto{\pgfqpoint{5.600022in}{0.856843in}}%
\pgfpathlineto{\pgfqpoint{5.665820in}{0.864070in}}%
\pgfpathlineto{\pgfqpoint{5.705410in}{0.868948in}}%
\pgfpathlineto{\pgfqpoint{5.727958in}{0.872540in}}%
\pgfpathlineto{\pgfqpoint{5.718618in}{0.857809in}}%
\pgfpathlineto{\pgfqpoint{5.718623in}{0.850863in}}%
\pgfpathlineto{\pgfqpoint{5.728238in}{0.847117in}}%
\pgfpathlineto{\pgfqpoint{5.734776in}{0.841070in}}%
\pgfpathlineto{\pgfqpoint{5.744657in}{0.840323in}}%
\pgfpathlineto{\pgfqpoint{5.753892in}{0.823299in}}%
\pgfpathlineto{\pgfqpoint{5.763910in}{0.811294in}}%
\pgfpathlineto{\pgfqpoint{5.781702in}{0.801205in}}%
\pgfpathlineto{\pgfqpoint{5.786758in}{0.793157in}}%
\pgfpathlineto{\pgfqpoint{5.802624in}{0.784454in}}%
\pgfpathlineto{\pgfqpoint{5.803817in}{0.778143in}}%
\pgfpathlineto{\pgfqpoint{5.816688in}{0.765593in}}%
\pgfpathlineto{\pgfqpoint{5.826546in}{0.762019in}}%
\pgfpathlineto{\pgfqpoint{5.833528in}{0.750169in}}%
\pgfpathlineto{\pgfqpoint{5.836618in}{0.736849in}}%
\pgfpathlineto{\pgfqpoint{5.846858in}{0.731700in}}%
\pgfpathlineto{\pgfqpoint{5.855102in}{0.717383in}}%
\pgfpathlineto{\pgfqpoint{5.857761in}{0.706853in}}%
\pgfpathlineto{\pgfqpoint{5.869841in}{0.702371in}}%
\pgfpathlineto{\pgfqpoint{5.859704in}{0.682961in}}%
\pgfpathlineto{\pgfqpoint{5.855889in}{0.671491in}}%
\pgfpathlineto{\pgfqpoint{5.858378in}{0.666114in}}%
\pgfpathlineto{\pgfqpoint{5.854032in}{0.657178in}}%
\pgfpathlineto{\pgfqpoint{5.852185in}{0.644824in}}%
\pgfpathlineto{\pgfqpoint{5.844456in}{0.642523in}}%
\pgfpathlineto{\pgfqpoint{5.848537in}{0.631212in}}%
\pgfpathlineto{\pgfqpoint{5.848277in}{0.616862in}}%
\pgfpathlineto{\pgfqpoint{5.834539in}{0.617400in}}%
\pgfpathlineto{\pgfqpoint{5.824969in}{0.620018in}}%
\pgfpathlineto{\pgfqpoint{5.820828in}{0.615151in}}%
\pgfpathlineto{\pgfqpoint{5.823918in}{0.603653in}}%
\pgfpathlineto{\pgfqpoint{5.823246in}{0.590285in}}%
\pgfpathlineto{\pgfqpoint{5.817209in}{0.589262in}}%
\pgfpathlineto{\pgfqpoint{5.812325in}{0.601747in}}%
\pgfpathlineto{\pgfqpoint{5.762593in}{0.598433in}}%
\pgfpathlineto{\pgfqpoint{5.699885in}{0.595011in}}%
\pgfpathlineto{\pgfqpoint{5.668260in}{0.592785in}}%
\pgfpathlineto{\pgfqpoint{5.658749in}{0.610119in}}%
\pgfpathlineto{\pgfqpoint{5.651935in}{0.621831in}}%
\pgfpathlineto{\pgfqpoint{5.652945in}{0.642483in}}%
\pgfpathlineto{\pgfqpoint{5.645672in}{0.663006in}}%
\pgfpathlineto{\pgfqpoint{5.648185in}{0.680553in}}%
\pgfpathlineto{\pgfqpoint{5.655627in}{0.690406in}}%
\pgfpathlineto{\pgfqpoint{5.649330in}{0.693573in}}%
\pgfpathlineto{\pgfqpoint{5.650753in}{0.700073in}}%
\pgfpathlineto{\pgfqpoint{5.641938in}{0.712845in}}%
\pgfpathlineto{\pgfqpoint{5.636339in}{0.726571in}}%
\pgfpathlineto{\pgfqpoint{5.626658in}{0.760657in}}%
\pgfpathlineto{\pgfqpoint{5.615858in}{0.801373in}}%
\pgfpathlineto{\pgfqpoint{5.600022in}{0.856843in}}%
\pgfpathclose%
\pgfusepath{fill}%
\end{pgfscope}%
\begin{pgfscope}%
\pgfpathrectangle{\pgfqpoint{3.525000in}{0.100000in}}{\pgfqpoint{2.862985in}{1.832778in}}%
\pgfusepath{clip}%
\pgfsetbuttcap%
\pgfsetmiterjoin%
\definecolor{currentfill}{rgb}{0.893250,0.411048,0.369768}%
\pgfsetfillcolor{currentfill}%
\pgfsetlinewidth{0.000000pt}%
\definecolor{currentstroke}{rgb}{0.000000,0.000000,0.000000}%
\pgfsetstrokecolor{currentstroke}%
\pgfsetstrokeopacity{0.000000}%
\pgfsetdash{}{0pt}%
\pgfpathmoveto{\pgfqpoint{5.727958in}{0.872540in}}%
\pgfpathlineto{\pgfqpoint{5.754049in}{0.885519in}}%
\pgfpathlineto{\pgfqpoint{5.768499in}{0.890420in}}%
\pgfpathlineto{\pgfqpoint{5.831366in}{0.896952in}}%
\pgfpathlineto{\pgfqpoint{5.838024in}{0.894789in}}%
\pgfpathlineto{\pgfqpoint{5.846819in}{0.885949in}}%
\pgfpathlineto{\pgfqpoint{5.847304in}{0.878098in}}%
\pgfpathlineto{\pgfqpoint{5.905165in}{0.886600in}}%
\pgfpathlineto{\pgfqpoint{5.970869in}{0.839501in}}%
\pgfpathlineto{\pgfqpoint{5.958547in}{0.826658in}}%
\pgfpathlineto{\pgfqpoint{5.941640in}{0.796833in}}%
\pgfpathlineto{\pgfqpoint{5.946453in}{0.790423in}}%
\pgfpathlineto{\pgfqpoint{5.937434in}{0.777981in}}%
\pgfpathlineto{\pgfqpoint{5.928462in}{0.776573in}}%
\pgfpathlineto{\pgfqpoint{5.928429in}{0.769097in}}%
\pgfpathlineto{\pgfqpoint{5.921939in}{0.761274in}}%
\pgfpathlineto{\pgfqpoint{5.913862in}{0.759687in}}%
\pgfpathlineto{\pgfqpoint{5.915625in}{0.752736in}}%
\pgfpathlineto{\pgfqpoint{5.911118in}{0.747400in}}%
\pgfpathlineto{\pgfqpoint{5.900337in}{0.742790in}}%
\pgfpathlineto{\pgfqpoint{5.887196in}{0.732304in}}%
\pgfpathlineto{\pgfqpoint{5.889677in}{0.725520in}}%
\pgfpathlineto{\pgfqpoint{5.881407in}{0.721260in}}%
\pgfpathlineto{\pgfqpoint{5.874432in}{0.725739in}}%
\pgfpathlineto{\pgfqpoint{5.869331in}{0.706267in}}%
\pgfpathlineto{\pgfqpoint{5.857761in}{0.706853in}}%
\pgfpathlineto{\pgfqpoint{5.855102in}{0.717383in}}%
\pgfpathlineto{\pgfqpoint{5.846858in}{0.731700in}}%
\pgfpathlineto{\pgfqpoint{5.836618in}{0.736849in}}%
\pgfpathlineto{\pgfqpoint{5.833528in}{0.750169in}}%
\pgfpathlineto{\pgfqpoint{5.826546in}{0.762019in}}%
\pgfpathlineto{\pgfqpoint{5.816688in}{0.765593in}}%
\pgfpathlineto{\pgfqpoint{5.803817in}{0.778143in}}%
\pgfpathlineto{\pgfqpoint{5.802624in}{0.784454in}}%
\pgfpathlineto{\pgfqpoint{5.786758in}{0.793157in}}%
\pgfpathlineto{\pgfqpoint{5.781702in}{0.801205in}}%
\pgfpathlineto{\pgfqpoint{5.763910in}{0.811294in}}%
\pgfpathlineto{\pgfqpoint{5.753892in}{0.823299in}}%
\pgfpathlineto{\pgfqpoint{5.744657in}{0.840323in}}%
\pgfpathlineto{\pgfqpoint{5.734776in}{0.841070in}}%
\pgfpathlineto{\pgfqpoint{5.728238in}{0.847117in}}%
\pgfpathlineto{\pgfqpoint{5.718623in}{0.850863in}}%
\pgfpathlineto{\pgfqpoint{5.718618in}{0.857809in}}%
\pgfpathlineto{\pgfqpoint{5.727958in}{0.872540in}}%
\pgfpathclose%
\pgfusepath{fill}%
\end{pgfscope}%
\begin{pgfscope}%
\pgfpathrectangle{\pgfqpoint{3.525000in}{0.100000in}}{\pgfqpoint{2.862985in}{1.832778in}}%
\pgfusepath{clip}%
\pgfsetbuttcap%
\pgfsetmiterjoin%
\definecolor{currentfill}{rgb}{0.856547,0.356066,0.410322}%
\pgfsetfillcolor{currentfill}%
\pgfsetlinewidth{0.000000pt}%
\definecolor{currentstroke}{rgb}{0.000000,0.000000,0.000000}%
\pgfsetstrokecolor{currentstroke}%
\pgfsetstrokeopacity{0.000000}%
\pgfsetdash{}{0pt}%
\pgfpathmoveto{\pgfqpoint{5.136515in}{0.926094in}}%
\pgfpathlineto{\pgfqpoint{5.174856in}{0.926440in}}%
\pgfpathlineto{\pgfqpoint{5.225580in}{0.927364in}}%
\pgfpathlineto{\pgfqpoint{5.312257in}{0.930116in}}%
\pgfpathlineto{\pgfqpoint{5.361833in}{0.932737in}}%
\pgfpathlineto{\pgfqpoint{5.367150in}{0.926155in}}%
\pgfpathlineto{\pgfqpoint{5.366818in}{0.919204in}}%
\pgfpathlineto{\pgfqpoint{5.354808in}{0.907211in}}%
\pgfpathlineto{\pgfqpoint{5.351904in}{0.900644in}}%
\pgfpathlineto{\pgfqpoint{5.385213in}{0.903107in}}%
\pgfpathlineto{\pgfqpoint{5.385195in}{0.890974in}}%
\pgfpathlineto{\pgfqpoint{5.374509in}{0.885789in}}%
\pgfpathlineto{\pgfqpoint{5.374660in}{0.877545in}}%
\pgfpathlineto{\pgfqpoint{5.370647in}{0.871619in}}%
\pgfpathlineto{\pgfqpoint{5.368020in}{0.859033in}}%
\pgfpathlineto{\pgfqpoint{5.370732in}{0.848649in}}%
\pgfpathlineto{\pgfqpoint{5.358848in}{0.838785in}}%
\pgfpathlineto{\pgfqpoint{5.362326in}{0.833932in}}%
\pgfpathlineto{\pgfqpoint{5.351180in}{0.825739in}}%
\pgfpathlineto{\pgfqpoint{5.351605in}{0.817725in}}%
\pgfpathlineto{\pgfqpoint{5.342656in}{0.798007in}}%
\pgfpathlineto{\pgfqpoint{5.334269in}{0.792035in}}%
\pgfpathlineto{\pgfqpoint{5.333459in}{0.783821in}}%
\pgfpathlineto{\pgfqpoint{5.316048in}{0.752960in}}%
\pgfpathlineto{\pgfqpoint{5.321946in}{0.735772in}}%
\pgfpathlineto{\pgfqpoint{5.320307in}{0.731223in}}%
\pgfpathlineto{\pgfqpoint{5.323659in}{0.721223in}}%
\pgfpathlineto{\pgfqpoint{5.320057in}{0.711695in}}%
\pgfpathlineto{\pgfqpoint{5.272446in}{0.709729in}}%
\pgfpathlineto{\pgfqpoint{5.210638in}{0.708711in}}%
\pgfpathlineto{\pgfqpoint{5.168006in}{0.708326in}}%
\pgfpathlineto{\pgfqpoint{5.167799in}{0.741894in}}%
\pgfpathlineto{\pgfqpoint{5.157212in}{0.744093in}}%
\pgfpathlineto{\pgfqpoint{5.150225in}{0.741221in}}%
\pgfpathlineto{\pgfqpoint{5.144643in}{0.746495in}}%
\pgfpathlineto{\pgfqpoint{5.145219in}{0.782052in}}%
\pgfpathlineto{\pgfqpoint{5.146468in}{0.857760in}}%
\pgfpathlineto{\pgfqpoint{5.140413in}{0.902092in}}%
\pgfpathlineto{\pgfqpoint{5.136515in}{0.926094in}}%
\pgfpathclose%
\pgfusepath{fill}%
\end{pgfscope}%
\begin{pgfscope}%
\pgfpathrectangle{\pgfqpoint{3.525000in}{0.100000in}}{\pgfqpoint{2.862985in}{1.832778in}}%
\pgfusepath{clip}%
\pgfsetbuttcap%
\pgfsetmiterjoin%
\definecolor{currentfill}{rgb}{0.819651,0.306812,0.448306}%
\pgfsetfillcolor{currentfill}%
\pgfsetlinewidth{0.000000pt}%
\definecolor{currentstroke}{rgb}{0.000000,0.000000,0.000000}%
\pgfsetstrokecolor{currentstroke}%
\pgfsetstrokeopacity{0.000000}%
\pgfsetdash{}{0pt}%
\pgfpathmoveto{\pgfqpoint{5.168006in}{0.708326in}}%
\pgfpathlineto{\pgfqpoint{5.210638in}{0.708711in}}%
\pgfpathlineto{\pgfqpoint{5.272446in}{0.709729in}}%
\pgfpathlineto{\pgfqpoint{5.320057in}{0.711695in}}%
\pgfpathlineto{\pgfqpoint{5.317776in}{0.706727in}}%
\pgfpathlineto{\pgfqpoint{5.323791in}{0.679370in}}%
\pgfpathlineto{\pgfqpoint{5.332102in}{0.667419in}}%
\pgfpathlineto{\pgfqpoint{5.322467in}{0.660818in}}%
\pgfpathlineto{\pgfqpoint{5.327503in}{0.649350in}}%
\pgfpathlineto{\pgfqpoint{5.313916in}{0.637542in}}%
\pgfpathlineto{\pgfqpoint{5.309854in}{0.619619in}}%
\pgfpathlineto{\pgfqpoint{5.302430in}{0.610386in}}%
\pgfpathlineto{\pgfqpoint{5.303375in}{0.600866in}}%
\pgfpathlineto{\pgfqpoint{5.299116in}{0.598847in}}%
\pgfpathlineto{\pgfqpoint{5.303390in}{0.588938in}}%
\pgfpathlineto{\pgfqpoint{5.299932in}{0.583691in}}%
\pgfpathlineto{\pgfqpoint{5.358060in}{0.586211in}}%
\pgfpathlineto{\pgfqpoint{5.402888in}{0.588972in}}%
\pgfpathlineto{\pgfqpoint{5.398157in}{0.567894in}}%
\pgfpathlineto{\pgfqpoint{5.401384in}{0.560107in}}%
\pgfpathlineto{\pgfqpoint{5.408075in}{0.553600in}}%
\pgfpathlineto{\pgfqpoint{5.414565in}{0.538074in}}%
\pgfpathlineto{\pgfqpoint{5.394065in}{0.541651in}}%
\pgfpathlineto{\pgfqpoint{5.386502in}{0.547531in}}%
\pgfpathlineto{\pgfqpoint{5.377492in}{0.547811in}}%
\pgfpathlineto{\pgfqpoint{5.368019in}{0.534934in}}%
\pgfpathlineto{\pgfqpoint{5.369907in}{0.529071in}}%
\pgfpathlineto{\pgfqpoint{5.400169in}{0.525523in}}%
\pgfpathlineto{\pgfqpoint{5.408091in}{0.518783in}}%
\pgfpathlineto{\pgfqpoint{5.420984in}{0.515315in}}%
\pgfpathlineto{\pgfqpoint{5.415353in}{0.507336in}}%
\pgfpathlineto{\pgfqpoint{5.405867in}{0.500376in}}%
\pgfpathlineto{\pgfqpoint{5.430906in}{0.484726in}}%
\pgfpathlineto{\pgfqpoint{5.442569in}{0.482281in}}%
\pgfpathlineto{\pgfqpoint{5.447419in}{0.469488in}}%
\pgfpathlineto{\pgfqpoint{5.436438in}{0.467135in}}%
\pgfpathlineto{\pgfqpoint{5.415799in}{0.483224in}}%
\pgfpathlineto{\pgfqpoint{5.407729in}{0.485461in}}%
\pgfpathlineto{\pgfqpoint{5.403824in}{0.491817in}}%
\pgfpathlineto{\pgfqpoint{5.391854in}{0.489208in}}%
\pgfpathlineto{\pgfqpoint{5.388107in}{0.481018in}}%
\pgfpathlineto{\pgfqpoint{5.390480in}{0.471892in}}%
\pgfpathlineto{\pgfqpoint{5.382422in}{0.466506in}}%
\pgfpathlineto{\pgfqpoint{5.375068in}{0.479774in}}%
\pgfpathlineto{\pgfqpoint{5.362187in}{0.475818in}}%
\pgfpathlineto{\pgfqpoint{5.357360in}{0.467901in}}%
\pgfpathlineto{\pgfqpoint{5.348213in}{0.470153in}}%
\pgfpathlineto{\pgfqpoint{5.342364in}{0.480150in}}%
\pgfpathlineto{\pgfqpoint{5.326847in}{0.483451in}}%
\pgfpathlineto{\pgfqpoint{5.323917in}{0.488653in}}%
\pgfpathlineto{\pgfqpoint{5.307751in}{0.497715in}}%
\pgfpathlineto{\pgfqpoint{5.303738in}{0.505638in}}%
\pgfpathlineto{\pgfqpoint{5.290197in}{0.502400in}}%
\pgfpathlineto{\pgfqpoint{5.291917in}{0.509654in}}%
\pgfpathlineto{\pgfqpoint{5.275086in}{0.502213in}}%
\pgfpathlineto{\pgfqpoint{5.279607in}{0.494492in}}%
\pgfpathlineto{\pgfqpoint{5.266610in}{0.489928in}}%
\pgfpathlineto{\pgfqpoint{5.249395in}{0.492437in}}%
\pgfpathlineto{\pgfqpoint{5.214551in}{0.504372in}}%
\pgfpathlineto{\pgfqpoint{5.187666in}{0.502003in}}%
\pgfpathlineto{\pgfqpoint{5.185335in}{0.517680in}}%
\pgfpathlineto{\pgfqpoint{5.188712in}{0.524834in}}%
\pgfpathlineto{\pgfqpoint{5.188423in}{0.536215in}}%
\pgfpathlineto{\pgfqpoint{5.185101in}{0.538927in}}%
\pgfpathlineto{\pgfqpoint{5.186227in}{0.551919in}}%
\pgfpathlineto{\pgfqpoint{5.196019in}{0.569988in}}%
\pgfpathlineto{\pgfqpoint{5.197272in}{0.576816in}}%
\pgfpathlineto{\pgfqpoint{5.195658in}{0.592942in}}%
\pgfpathlineto{\pgfqpoint{5.188454in}{0.600538in}}%
\pgfpathlineto{\pgfqpoint{5.184693in}{0.614721in}}%
\pgfpathlineto{\pgfqpoint{5.180065in}{0.617977in}}%
\pgfpathlineto{\pgfqpoint{5.182137in}{0.625668in}}%
\pgfpathlineto{\pgfqpoint{5.176243in}{0.637164in}}%
\pgfpathlineto{\pgfqpoint{5.168890in}{0.643403in}}%
\pgfpathlineto{\pgfqpoint{5.168006in}{0.708326in}}%
\pgfpathclose%
\pgfusepath{fill}%
\end{pgfscope}%
\begin{pgfscope}%
\pgfpathrectangle{\pgfqpoint{3.525000in}{0.100000in}}{\pgfqpoint{2.862985in}{1.832778in}}%
\pgfusepath{clip}%
\pgfsetbuttcap%
\pgfsetmiterjoin%
\definecolor{currentfill}{rgb}{0.819651,0.306812,0.448306}%
\pgfsetfillcolor{currentfill}%
\pgfsetlinewidth{0.000000pt}%
\definecolor{currentstroke}{rgb}{0.000000,0.000000,0.000000}%
\pgfsetstrokecolor{currentstroke}%
\pgfsetstrokeopacity{0.000000}%
\pgfsetdash{}{0pt}%
\pgfpathmoveto{\pgfqpoint{5.282025in}{0.493901in}}%
\pgfpathlineto{\pgfqpoint{5.288193in}{0.497570in}}%
\pgfpathlineto{\pgfqpoint{5.295679in}{0.493233in}}%
\pgfpathlineto{\pgfqpoint{5.291504in}{0.487265in}}%
\pgfpathlineto{\pgfqpoint{5.282025in}{0.493901in}}%
\pgfpathclose%
\pgfusepath{fill}%
\end{pgfscope}%
\begin{pgfscope}%
\pgfpathrectangle{\pgfqpoint{3.525000in}{0.100000in}}{\pgfqpoint{2.862985in}{1.832778in}}%
\pgfusepath{clip}%
\pgfsetbuttcap%
\pgfsetmiterjoin%
\definecolor{currentfill}{rgb}{0.798216,0.280197,0.469538}%
\pgfsetfillcolor{currentfill}%
\pgfsetlinewidth{0.000000pt}%
\definecolor{currentstroke}{rgb}{0.000000,0.000000,0.000000}%
\pgfsetstrokecolor{currentstroke}%
\pgfsetstrokeopacity{0.000000}%
\pgfsetdash{}{0pt}%
\pgfpathmoveto{\pgfqpoint{5.530919in}{0.564647in}}%
\pgfpathlineto{\pgfqpoint{5.531031in}{0.576507in}}%
\pgfpathlineto{\pgfqpoint{5.523661in}{0.581011in}}%
\pgfpathlineto{\pgfqpoint{5.517624in}{0.588678in}}%
\pgfpathlineto{\pgfqpoint{5.518446in}{0.596734in}}%
\pgfpathlineto{\pgfqpoint{5.584178in}{0.601837in}}%
\pgfpathlineto{\pgfqpoint{5.658749in}{0.610119in}}%
\pgfpathlineto{\pgfqpoint{5.668260in}{0.592785in}}%
\pgfpathlineto{\pgfqpoint{5.699885in}{0.595011in}}%
\pgfpathlineto{\pgfqpoint{5.762593in}{0.598433in}}%
\pgfpathlineto{\pgfqpoint{5.812325in}{0.601747in}}%
\pgfpathlineto{\pgfqpoint{5.817209in}{0.589262in}}%
\pgfpathlineto{\pgfqpoint{5.823246in}{0.590285in}}%
\pgfpathlineto{\pgfqpoint{5.823918in}{0.603653in}}%
\pgfpathlineto{\pgfqpoint{5.820828in}{0.615151in}}%
\pgfpathlineto{\pgfqpoint{5.824969in}{0.620018in}}%
\pgfpathlineto{\pgfqpoint{5.834539in}{0.617400in}}%
\pgfpathlineto{\pgfqpoint{5.848277in}{0.616862in}}%
\pgfpathlineto{\pgfqpoint{5.850389in}{0.606589in}}%
\pgfpathlineto{\pgfqpoint{5.854615in}{0.600719in}}%
\pgfpathlineto{\pgfqpoint{5.857915in}{0.587877in}}%
\pgfpathlineto{\pgfqpoint{5.868139in}{0.568003in}}%
\pgfpathlineto{\pgfqpoint{5.868190in}{0.562628in}}%
\pgfpathlineto{\pgfqpoint{5.883302in}{0.539291in}}%
\pgfpathlineto{\pgfqpoint{5.884764in}{0.534471in}}%
\pgfpathlineto{\pgfqpoint{5.907434in}{0.501512in}}%
\pgfpathlineto{\pgfqpoint{5.903842in}{0.500979in}}%
\pgfpathlineto{\pgfqpoint{5.913399in}{0.477537in}}%
\pgfpathlineto{\pgfqpoint{5.939653in}{0.436782in}}%
\pgfpathlineto{\pgfqpoint{5.954896in}{0.406818in}}%
\pgfpathlineto{\pgfqpoint{5.959991in}{0.401571in}}%
\pgfpathlineto{\pgfqpoint{5.968722in}{0.382763in}}%
\pgfpathlineto{\pgfqpoint{5.971750in}{0.352646in}}%
\pgfpathlineto{\pgfqpoint{5.972955in}{0.330127in}}%
\pgfpathlineto{\pgfqpoint{5.971500in}{0.315709in}}%
\pgfpathlineto{\pgfqpoint{5.966819in}{0.305401in}}%
\pgfpathlineto{\pgfqpoint{5.968994in}{0.291885in}}%
\pgfpathlineto{\pgfqpoint{5.963954in}{0.281186in}}%
\pgfpathlineto{\pgfqpoint{5.948991in}{0.272413in}}%
\pgfpathlineto{\pgfqpoint{5.939259in}{0.273053in}}%
\pgfpathlineto{\pgfqpoint{5.932940in}{0.268467in}}%
\pgfpathlineto{\pgfqpoint{5.931129in}{0.280712in}}%
\pgfpathlineto{\pgfqpoint{5.920668in}{0.283944in}}%
\pgfpathlineto{\pgfqpoint{5.911290in}{0.301030in}}%
\pgfpathlineto{\pgfqpoint{5.910231in}{0.308810in}}%
\pgfpathlineto{\pgfqpoint{5.893448in}{0.313535in}}%
\pgfpathlineto{\pgfqpoint{5.882633in}{0.312512in}}%
\pgfpathlineto{\pgfqpoint{5.876514in}{0.323801in}}%
\pgfpathlineto{\pgfqpoint{5.869508in}{0.344093in}}%
\pgfpathlineto{\pgfqpoint{5.859786in}{0.348206in}}%
\pgfpathlineto{\pgfqpoint{5.854515in}{0.359834in}}%
\pgfpathlineto{\pgfqpoint{5.841625in}{0.366687in}}%
\pgfpathlineto{\pgfqpoint{5.834174in}{0.375259in}}%
\pgfpathlineto{\pgfqpoint{5.821320in}{0.398564in}}%
\pgfpathlineto{\pgfqpoint{5.819809in}{0.414966in}}%
\pgfpathlineto{\pgfqpoint{5.826813in}{0.425461in}}%
\pgfpathlineto{\pgfqpoint{5.826113in}{0.432754in}}%
\pgfpathlineto{\pgfqpoint{5.818019in}{0.433533in}}%
\pgfpathlineto{\pgfqpoint{5.811256in}{0.438606in}}%
\pgfpathlineto{\pgfqpoint{5.807553in}{0.432412in}}%
\pgfpathlineto{\pgfqpoint{5.814043in}{0.427365in}}%
\pgfpathlineto{\pgfqpoint{5.809319in}{0.418289in}}%
\pgfpathlineto{\pgfqpoint{5.801677in}{0.425822in}}%
\pgfpathlineto{\pgfqpoint{5.802564in}{0.446710in}}%
\pgfpathlineto{\pgfqpoint{5.806250in}{0.463657in}}%
\pgfpathlineto{\pgfqpoint{5.804355in}{0.492736in}}%
\pgfpathlineto{\pgfqpoint{5.796738in}{0.499662in}}%
\pgfpathlineto{\pgfqpoint{5.792907in}{0.508549in}}%
\pgfpathlineto{\pgfqpoint{5.779785in}{0.508352in}}%
\pgfpathlineto{\pgfqpoint{5.766814in}{0.522976in}}%
\pgfpathlineto{\pgfqpoint{5.758117in}{0.527359in}}%
\pgfpathlineto{\pgfqpoint{5.755538in}{0.536620in}}%
\pgfpathlineto{\pgfqpoint{5.747031in}{0.539872in}}%
\pgfpathlineto{\pgfqpoint{5.739981in}{0.550098in}}%
\pgfpathlineto{\pgfqpoint{5.721358in}{0.558444in}}%
\pgfpathlineto{\pgfqpoint{5.706879in}{0.558658in}}%
\pgfpathlineto{\pgfqpoint{5.700576in}{0.555452in}}%
\pgfpathlineto{\pgfqpoint{5.702143in}{0.545445in}}%
\pgfpathlineto{\pgfqpoint{5.695570in}{0.545919in}}%
\pgfpathlineto{\pgfqpoint{5.675353in}{0.531927in}}%
\pgfpathlineto{\pgfqpoint{5.651068in}{0.526381in}}%
\pgfpathlineto{\pgfqpoint{5.650669in}{0.533243in}}%
\pgfpathlineto{\pgfqpoint{5.645287in}{0.539949in}}%
\pgfpathlineto{\pgfqpoint{5.630848in}{0.549208in}}%
\pgfpathlineto{\pgfqpoint{5.610121in}{0.558700in}}%
\pgfpathlineto{\pgfqpoint{5.587608in}{0.563725in}}%
\pgfpathlineto{\pgfqpoint{5.583382in}{0.570558in}}%
\pgfpathlineto{\pgfqpoint{5.575265in}{0.564859in}}%
\pgfpathlineto{\pgfqpoint{5.565491in}{0.563602in}}%
\pgfpathlineto{\pgfqpoint{5.531467in}{0.554637in}}%
\pgfpathlineto{\pgfqpoint{5.530919in}{0.564647in}}%
\pgfpathclose%
\pgfusepath{fill}%
\end{pgfscope}%
\begin{pgfscope}%
\pgfpathrectangle{\pgfqpoint{3.525000in}{0.100000in}}{\pgfqpoint{2.862985in}{1.832778in}}%
\pgfusepath{clip}%
\pgfsetbuttcap%
\pgfsetmiterjoin%
\definecolor{currentfill}{rgb}{0.798216,0.280197,0.469538}%
\pgfsetfillcolor{currentfill}%
\pgfsetlinewidth{0.000000pt}%
\definecolor{currentstroke}{rgb}{0.000000,0.000000,0.000000}%
\pgfsetstrokecolor{currentstroke}%
\pgfsetstrokeopacity{0.000000}%
\pgfsetdash{}{0pt}%
\pgfpathmoveto{\pgfqpoint{5.910737in}{0.501980in}}%
\pgfpathlineto{\pgfqpoint{5.919540in}{0.491740in}}%
\pgfpathlineto{\pgfqpoint{5.909372in}{0.491096in}}%
\pgfpathlineto{\pgfqpoint{5.910737in}{0.501980in}}%
\pgfpathclose%
\pgfusepath{fill}%
\end{pgfscope}%
\begin{pgfscope}%
\pgfpathrectangle{\pgfqpoint{3.525000in}{0.100000in}}{\pgfqpoint{2.862985in}{1.832778in}}%
\pgfusepath{clip}%
\pgfsetbuttcap%
\pgfsetmiterjoin%
\definecolor{currentfill}{rgb}{0.801855,0.284626,0.465971}%
\pgfsetfillcolor{currentfill}%
\pgfsetlinewidth{0.000000pt}%
\definecolor{currentstroke}{rgb}{0.000000,0.000000,0.000000}%
\pgfsetstrokecolor{currentstroke}%
\pgfsetstrokeopacity{0.000000}%
\pgfsetdash{}{0pt}%
\pgfpathmoveto{\pgfqpoint{5.392522in}{1.671190in}}%
\pgfpathlineto{\pgfqpoint{5.387746in}{1.661894in}}%
\pgfpathlineto{\pgfqpoint{5.376351in}{1.656468in}}%
\pgfpathlineto{\pgfqpoint{5.371416in}{1.649166in}}%
\pgfpathlineto{\pgfqpoint{5.365635in}{1.654425in}}%
\pgfpathlineto{\pgfqpoint{5.392522in}{1.671190in}}%
\pgfpathclose%
\pgfusepath{fill}%
\end{pgfscope}%
\begin{pgfscope}%
\pgfpathrectangle{\pgfqpoint{3.525000in}{0.100000in}}{\pgfqpoint{2.862985in}{1.832778in}}%
\pgfusepath{clip}%
\pgfsetbuttcap%
\pgfsetmiterjoin%
\definecolor{currentfill}{rgb}{0.801855,0.284626,0.465971}%
\pgfsetfillcolor{currentfill}%
\pgfsetlinewidth{0.000000pt}%
\definecolor{currentstroke}{rgb}{0.000000,0.000000,0.000000}%
\pgfsetstrokecolor{currentstroke}%
\pgfsetstrokeopacity{0.000000}%
\pgfsetdash{}{0pt}%
\pgfpathmoveto{\pgfqpoint{5.396396in}{1.615244in}}%
\pgfpathlineto{\pgfqpoint{5.408054in}{1.626120in}}%
\pgfpathlineto{\pgfqpoint{5.426037in}{1.628969in}}%
\pgfpathlineto{\pgfqpoint{5.421084in}{1.621405in}}%
\pgfpathlineto{\pgfqpoint{5.408756in}{1.610466in}}%
\pgfpathlineto{\pgfqpoint{5.401549in}{1.596426in}}%
\pgfpathlineto{\pgfqpoint{5.398496in}{1.604041in}}%
\pgfpathlineto{\pgfqpoint{5.393076in}{1.605134in}}%
\pgfpathlineto{\pgfqpoint{5.392532in}{1.612009in}}%
\pgfpathlineto{\pgfqpoint{5.396396in}{1.615244in}}%
\pgfpathclose%
\pgfusepath{fill}%
\end{pgfscope}%
\begin{pgfscope}%
\pgfpathrectangle{\pgfqpoint{3.525000in}{0.100000in}}{\pgfqpoint{2.862985in}{1.832778in}}%
\pgfusepath{clip}%
\pgfsetbuttcap%
\pgfsetmiterjoin%
\definecolor{currentfill}{rgb}{0.801855,0.284626,0.465971}%
\pgfsetfillcolor{currentfill}%
\pgfsetlinewidth{0.000000pt}%
\definecolor{currentstroke}{rgb}{0.000000,0.000000,0.000000}%
\pgfsetstrokecolor{currentstroke}%
\pgfsetstrokeopacity{0.000000}%
\pgfsetdash{}{0pt}%
\pgfpathmoveto{\pgfqpoint{5.442832in}{1.482495in}}%
\pgfpathlineto{\pgfqpoint{5.439727in}{1.485941in}}%
\pgfpathlineto{\pgfqpoint{5.442991in}{1.495641in}}%
\pgfpathlineto{\pgfqpoint{5.433301in}{1.496259in}}%
\pgfpathlineto{\pgfqpoint{5.435870in}{1.504623in}}%
\pgfpathlineto{\pgfqpoint{5.434272in}{1.517943in}}%
\pgfpathlineto{\pgfqpoint{5.425578in}{1.522561in}}%
\pgfpathlineto{\pgfqpoint{5.416480in}{1.531907in}}%
\pgfpathlineto{\pgfqpoint{5.402460in}{1.534643in}}%
\pgfpathlineto{\pgfqpoint{5.388816in}{1.534565in}}%
\pgfpathlineto{\pgfqpoint{5.375397in}{1.541198in}}%
\pgfpathlineto{\pgfqpoint{5.330502in}{1.550866in}}%
\pgfpathlineto{\pgfqpoint{5.325596in}{1.561110in}}%
\pgfpathlineto{\pgfqpoint{5.316845in}{1.564606in}}%
\pgfpathlineto{\pgfqpoint{5.333375in}{1.572449in}}%
\pgfpathlineto{\pgfqpoint{5.342713in}{1.582234in}}%
\pgfpathlineto{\pgfqpoint{5.360106in}{1.584888in}}%
\pgfpathlineto{\pgfqpoint{5.370809in}{1.594855in}}%
\pgfpathlineto{\pgfqpoint{5.376425in}{1.595254in}}%
\pgfpathlineto{\pgfqpoint{5.380710in}{1.602362in}}%
\pgfpathlineto{\pgfqpoint{5.391157in}{1.610783in}}%
\pgfpathlineto{\pgfqpoint{5.392067in}{1.604831in}}%
\pgfpathlineto{\pgfqpoint{5.396774in}{1.603602in}}%
\pgfpathlineto{\pgfqpoint{5.400316in}{1.596504in}}%
\pgfpathlineto{\pgfqpoint{5.400954in}{1.584403in}}%
\pgfpathlineto{\pgfqpoint{5.411599in}{1.591634in}}%
\pgfpathlineto{\pgfqpoint{5.424002in}{1.593145in}}%
\pgfpathlineto{\pgfqpoint{5.434590in}{1.589356in}}%
\pgfpathlineto{\pgfqpoint{5.448949in}{1.569658in}}%
\pgfpathlineto{\pgfqpoint{5.464631in}{1.572806in}}%
\pgfpathlineto{\pgfqpoint{5.471003in}{1.567535in}}%
\pgfpathlineto{\pgfqpoint{5.481252in}{1.567064in}}%
\pgfpathlineto{\pgfqpoint{5.488061in}{1.576518in}}%
\pgfpathlineto{\pgfqpoint{5.500982in}{1.584876in}}%
\pgfpathlineto{\pgfqpoint{5.513400in}{1.587493in}}%
\pgfpathlineto{\pgfqpoint{5.528808in}{1.587766in}}%
\pgfpathlineto{\pgfqpoint{5.540068in}{1.594222in}}%
\pgfpathlineto{\pgfqpoint{5.549258in}{1.591248in}}%
\pgfpathlineto{\pgfqpoint{5.551215in}{1.577594in}}%
\pgfpathlineto{\pgfqpoint{5.570930in}{1.575456in}}%
\pgfpathlineto{\pgfqpoint{5.581638in}{1.581847in}}%
\pgfpathlineto{\pgfqpoint{5.589061in}{1.567449in}}%
\pgfpathlineto{\pgfqpoint{5.603225in}{1.550799in}}%
\pgfpathlineto{\pgfqpoint{5.583401in}{1.550895in}}%
\pgfpathlineto{\pgfqpoint{5.577137in}{1.548837in}}%
\pgfpathlineto{\pgfqpoint{5.568551in}{1.551511in}}%
\pgfpathlineto{\pgfqpoint{5.567974in}{1.539985in}}%
\pgfpathlineto{\pgfqpoint{5.552384in}{1.548993in}}%
\pgfpathlineto{\pgfqpoint{5.532335in}{1.551746in}}%
\pgfpathlineto{\pgfqpoint{5.526816in}{1.542973in}}%
\pgfpathlineto{\pgfqpoint{5.500584in}{1.538707in}}%
\pgfpathlineto{\pgfqpoint{5.497566in}{1.531276in}}%
\pgfpathlineto{\pgfqpoint{5.479312in}{1.529046in}}%
\pgfpathlineto{\pgfqpoint{5.473794in}{1.521492in}}%
\pgfpathlineto{\pgfqpoint{5.464139in}{1.519466in}}%
\pgfpathlineto{\pgfqpoint{5.456427in}{1.501552in}}%
\pgfpathlineto{\pgfqpoint{5.446663in}{1.484203in}}%
\pgfpathlineto{\pgfqpoint{5.442832in}{1.482495in}}%
\pgfpathclose%
\pgfusepath{fill}%
\end{pgfscope}%
\begin{pgfscope}%
\pgfpathrectangle{\pgfqpoint{3.525000in}{0.100000in}}{\pgfqpoint{2.862985in}{1.832778in}}%
\pgfusepath{clip}%
\pgfsetbuttcap%
\pgfsetmiterjoin%
\definecolor{currentfill}{rgb}{0.801855,0.284626,0.465971}%
\pgfsetfillcolor{currentfill}%
\pgfsetlinewidth{0.000000pt}%
\definecolor{currentstroke}{rgb}{0.000000,0.000000,0.000000}%
\pgfsetstrokecolor{currentstroke}%
\pgfsetstrokeopacity{0.000000}%
\pgfsetdash{}{0pt}%
\pgfpathmoveto{\pgfqpoint{5.499029in}{1.274203in}}%
\pgfpathlineto{\pgfqpoint{5.508339in}{1.284002in}}%
\pgfpathlineto{\pgfqpoint{5.512580in}{1.298209in}}%
\pgfpathlineto{\pgfqpoint{5.517623in}{1.306443in}}%
\pgfpathlineto{\pgfqpoint{5.520716in}{1.317649in}}%
\pgfpathlineto{\pgfqpoint{5.521678in}{1.339993in}}%
\pgfpathlineto{\pgfqpoint{5.517021in}{1.361407in}}%
\pgfpathlineto{\pgfqpoint{5.501633in}{1.394197in}}%
\pgfpathlineto{\pgfqpoint{5.505800in}{1.404507in}}%
\pgfpathlineto{\pgfqpoint{5.500372in}{1.418758in}}%
\pgfpathlineto{\pgfqpoint{5.509684in}{1.438466in}}%
\pgfpathlineto{\pgfqpoint{5.508166in}{1.460442in}}%
\pgfpathlineto{\pgfqpoint{5.514653in}{1.463207in}}%
\pgfpathlineto{\pgfqpoint{5.515468in}{1.473664in}}%
\pgfpathlineto{\pgfqpoint{5.526982in}{1.480375in}}%
\pgfpathlineto{\pgfqpoint{5.533494in}{1.479291in}}%
\pgfpathlineto{\pgfqpoint{5.535324in}{1.468077in}}%
\pgfpathlineto{\pgfqpoint{5.540407in}{1.467638in}}%
\pgfpathlineto{\pgfqpoint{5.545044in}{1.483813in}}%
\pgfpathlineto{\pgfqpoint{5.543468in}{1.496393in}}%
\pgfpathlineto{\pgfqpoint{5.546483in}{1.503617in}}%
\pgfpathlineto{\pgfqpoint{5.562785in}{1.511080in}}%
\pgfpathlineto{\pgfqpoint{5.555356in}{1.513735in}}%
\pgfpathlineto{\pgfqpoint{5.552936in}{1.520146in}}%
\pgfpathlineto{\pgfqpoint{5.558288in}{1.531404in}}%
\pgfpathlineto{\pgfqpoint{5.568841in}{1.535292in}}%
\pgfpathlineto{\pgfqpoint{5.581095in}{1.528650in}}%
\pgfpathlineto{\pgfqpoint{5.592648in}{1.528567in}}%
\pgfpathlineto{\pgfqpoint{5.598015in}{1.520808in}}%
\pgfpathlineto{\pgfqpoint{5.606106in}{1.521347in}}%
\pgfpathlineto{\pgfqpoint{5.631083in}{1.510561in}}%
\pgfpathlineto{\pgfqpoint{5.636074in}{1.500120in}}%
\pgfpathlineto{\pgfqpoint{5.630624in}{1.496504in}}%
\pgfpathlineto{\pgfqpoint{5.632303in}{1.488679in}}%
\pgfpathlineto{\pgfqpoint{5.637687in}{1.485198in}}%
\pgfpathlineto{\pgfqpoint{5.640688in}{1.475579in}}%
\pgfpathlineto{\pgfqpoint{5.640268in}{1.452144in}}%
\pgfpathlineto{\pgfqpoint{5.633159in}{1.446546in}}%
\pgfpathlineto{\pgfqpoint{5.631564in}{1.434236in}}%
\pgfpathlineto{\pgfqpoint{5.618377in}{1.422859in}}%
\pgfpathlineto{\pgfqpoint{5.619159in}{1.409089in}}%
\pgfpathlineto{\pgfqpoint{5.630717in}{1.404265in}}%
\pgfpathlineto{\pgfqpoint{5.643756in}{1.421446in}}%
\pgfpathlineto{\pgfqpoint{5.644833in}{1.427678in}}%
\pgfpathlineto{\pgfqpoint{5.661094in}{1.438035in}}%
\pgfpathlineto{\pgfqpoint{5.671435in}{1.433238in}}%
\pgfpathlineto{\pgfqpoint{5.677924in}{1.422395in}}%
\pgfpathlineto{\pgfqpoint{5.688393in}{1.384749in}}%
\pgfpathlineto{\pgfqpoint{5.693940in}{1.372839in}}%
\pgfpathlineto{\pgfqpoint{5.692203in}{1.367915in}}%
\pgfpathlineto{\pgfqpoint{5.692372in}{1.351150in}}%
\pgfpathlineto{\pgfqpoint{5.682281in}{1.352759in}}%
\pgfpathlineto{\pgfqpoint{5.676581in}{1.340231in}}%
\pgfpathlineto{\pgfqpoint{5.675791in}{1.331714in}}%
\pgfpathlineto{\pgfqpoint{5.668169in}{1.326246in}}%
\pgfpathlineto{\pgfqpoint{5.666475in}{1.309636in}}%
\pgfpathlineto{\pgfqpoint{5.655390in}{1.288653in}}%
\pgfpathlineto{\pgfqpoint{5.594752in}{1.279623in}}%
\pgfpathlineto{\pgfqpoint{5.594393in}{1.283592in}}%
\pgfpathlineto{\pgfqpoint{5.553831in}{1.279317in}}%
\pgfpathlineto{\pgfqpoint{5.499029in}{1.274203in}}%
\pgfpathclose%
\pgfusepath{fill}%
\end{pgfscope}%
\begin{pgfscope}%
\pgfpathrectangle{\pgfqpoint{3.525000in}{0.100000in}}{\pgfqpoint{2.862985in}{1.832778in}}%
\pgfusepath{clip}%
\pgfsetbuttcap%
\pgfsetmiterjoin%
\definecolor{currentfill}{rgb}{0.645872,0.120898,0.598867}%
\pgfsetfillcolor{currentfill}%
\pgfsetlinewidth{0.000000pt}%
\definecolor{currentstroke}{rgb}{0.000000,0.000000,0.000000}%
\pgfsetstrokecolor{currentstroke}%
\pgfsetstrokeopacity{0.000000}%
\pgfsetdash{}{0pt}%
\pgfpathmoveto{\pgfqpoint{3.967548in}{0.431818in}}%
\pgfpathlineto{\pgfqpoint{3.966337in}{0.430752in}}%
\pgfpathlineto{\pgfqpoint{3.960781in}{0.432173in}}%
\pgfpathlineto{\pgfqpoint{3.961196in}{0.434969in}}%
\pgfpathlineto{\pgfqpoint{3.964341in}{0.436249in}}%
\pgfpathlineto{\pgfqpoint{3.968824in}{0.442694in}}%
\pgfpathlineto{\pgfqpoint{3.969569in}{0.447051in}}%
\pgfpathlineto{\pgfqpoint{3.972090in}{0.448192in}}%
\pgfpathlineto{\pgfqpoint{3.976190in}{0.448197in}}%
\pgfpathlineto{\pgfqpoint{3.980477in}{0.462333in}}%
\pgfpathlineto{\pgfqpoint{3.982157in}{0.463496in}}%
\pgfpathlineto{\pgfqpoint{3.980651in}{0.466260in}}%
\pgfpathlineto{\pgfqpoint{3.977125in}{0.463464in}}%
\pgfpathlineto{\pgfqpoint{3.966766in}{0.467233in}}%
\pgfpathlineto{\pgfqpoint{3.960583in}{0.473582in}}%
\pgfpathlineto{\pgfqpoint{3.963280in}{0.476479in}}%
\pgfpathlineto{\pgfqpoint{3.962053in}{0.480565in}}%
\pgfpathlineto{\pgfqpoint{3.962809in}{0.485392in}}%
\pgfpathlineto{\pgfqpoint{3.961479in}{0.491035in}}%
\pgfpathlineto{\pgfqpoint{3.961266in}{0.496684in}}%
\pgfpathlineto{\pgfqpoint{3.967217in}{0.495739in}}%
\pgfpathlineto{\pgfqpoint{3.968395in}{0.497404in}}%
\pgfpathlineto{\pgfqpoint{3.971355in}{0.497243in}}%
\pgfpathlineto{\pgfqpoint{3.974764in}{0.491217in}}%
\pgfpathlineto{\pgfqpoint{3.978051in}{0.486806in}}%
\pgfpathlineto{\pgfqpoint{3.975418in}{0.483935in}}%
\pgfpathlineto{\pgfqpoint{3.979970in}{0.482833in}}%
\pgfpathlineto{\pgfqpoint{3.981007in}{0.484241in}}%
\pgfpathlineto{\pgfqpoint{3.978573in}{0.489956in}}%
\pgfpathlineto{\pgfqpoint{3.973447in}{0.494692in}}%
\pgfpathlineto{\pgfqpoint{3.974074in}{0.496676in}}%
\pgfpathlineto{\pgfqpoint{3.969984in}{0.501691in}}%
\pgfpathlineto{\pgfqpoint{3.974161in}{0.502710in}}%
\pgfpathlineto{\pgfqpoint{3.970542in}{0.513448in}}%
\pgfpathlineto{\pgfqpoint{3.973484in}{0.514936in}}%
\pgfpathlineto{\pgfqpoint{3.974202in}{0.517511in}}%
\pgfpathlineto{\pgfqpoint{3.972288in}{0.520569in}}%
\pgfpathlineto{\pgfqpoint{3.976239in}{0.521335in}}%
\pgfpathlineto{\pgfqpoint{3.976948in}{0.523648in}}%
\pgfpathlineto{\pgfqpoint{3.981660in}{0.520176in}}%
\pgfpathlineto{\pgfqpoint{3.982707in}{0.523941in}}%
\pgfpathlineto{\pgfqpoint{3.985183in}{0.525597in}}%
\pgfpathlineto{\pgfqpoint{3.997233in}{0.527901in}}%
\pgfpathlineto{\pgfqpoint{4.000932in}{0.527086in}}%
\pgfpathlineto{\pgfqpoint{4.004031in}{0.531659in}}%
\pgfpathlineto{\pgfqpoint{4.011384in}{0.534658in}}%
\pgfpathlineto{\pgfqpoint{4.017047in}{0.533521in}}%
\pgfpathlineto{\pgfqpoint{4.019631in}{0.529538in}}%
\pgfpathlineto{\pgfqpoint{4.019495in}{0.523212in}}%
\pgfpathlineto{\pgfqpoint{4.022394in}{0.521778in}}%
\pgfpathlineto{\pgfqpoint{4.028306in}{0.521978in}}%
\pgfpathlineto{\pgfqpoint{4.035878in}{0.523624in}}%
\pgfpathlineto{\pgfqpoint{4.035785in}{0.521601in}}%
\pgfpathlineto{\pgfqpoint{4.045357in}{0.515832in}}%
\pgfpathlineto{\pgfqpoint{4.052052in}{0.517379in}}%
\pgfpathlineto{\pgfqpoint{4.054022in}{0.519381in}}%
\pgfpathlineto{\pgfqpoint{4.056133in}{0.524508in}}%
\pgfpathlineto{\pgfqpoint{4.058750in}{0.527792in}}%
\pgfpathlineto{\pgfqpoint{4.058984in}{0.532712in}}%
\pgfpathlineto{\pgfqpoint{4.056587in}{0.534647in}}%
\pgfpathlineto{\pgfqpoint{4.060136in}{0.536396in}}%
\pgfpathlineto{\pgfqpoint{4.066098in}{0.533768in}}%
\pgfpathlineto{\pgfqpoint{4.068040in}{0.535390in}}%
\pgfpathlineto{\pgfqpoint{4.068382in}{0.542352in}}%
\pgfpathlineto{\pgfqpoint{4.064528in}{0.540897in}}%
\pgfpathlineto{\pgfqpoint{4.061855in}{0.543309in}}%
\pgfpathlineto{\pgfqpoint{4.057516in}{0.544812in}}%
\pgfpathlineto{\pgfqpoint{4.050629in}{0.544561in}}%
\pgfpathlineto{\pgfqpoint{4.044030in}{0.546848in}}%
\pgfpathlineto{\pgfqpoint{4.043699in}{0.552545in}}%
\pgfpathlineto{\pgfqpoint{4.037797in}{0.558597in}}%
\pgfpathlineto{\pgfqpoint{4.030313in}{0.561138in}}%
\pgfpathlineto{\pgfqpoint{4.023780in}{0.572853in}}%
\pgfpathlineto{\pgfqpoint{4.024316in}{0.577189in}}%
\pgfpathlineto{\pgfqpoint{4.027833in}{0.579165in}}%
\pgfpathlineto{\pgfqpoint{4.027315in}{0.583269in}}%
\pgfpathlineto{\pgfqpoint{4.028132in}{0.587352in}}%
\pgfpathlineto{\pgfqpoint{4.030915in}{0.584557in}}%
\pgfpathlineto{\pgfqpoint{4.034590in}{0.585723in}}%
\pgfpathlineto{\pgfqpoint{4.034407in}{0.589270in}}%
\pgfpathlineto{\pgfqpoint{4.029311in}{0.596390in}}%
\pgfpathlineto{\pgfqpoint{4.027820in}{0.603734in}}%
\pgfpathlineto{\pgfqpoint{4.032085in}{0.604294in}}%
\pgfpathlineto{\pgfqpoint{4.038468in}{0.603843in}}%
\pgfpathlineto{\pgfqpoint{4.040478in}{0.601387in}}%
\pgfpathlineto{\pgfqpoint{4.046201in}{0.602266in}}%
\pgfpathlineto{\pgfqpoint{4.051520in}{0.601112in}}%
\pgfpathlineto{\pgfqpoint{4.055115in}{0.599083in}}%
\pgfpathlineto{\pgfqpoint{4.057604in}{0.600010in}}%
\pgfpathlineto{\pgfqpoint{4.064212in}{0.597983in}}%
\pgfpathlineto{\pgfqpoint{4.064273in}{0.596187in}}%
\pgfpathlineto{\pgfqpoint{4.069305in}{0.596794in}}%
\pgfpathlineto{\pgfqpoint{4.076053in}{0.592414in}}%
\pgfpathlineto{\pgfqpoint{4.075791in}{0.589506in}}%
\pgfpathlineto{\pgfqpoint{4.072580in}{0.588164in}}%
\pgfpathlineto{\pgfqpoint{4.069167in}{0.584972in}}%
\pgfpathlineto{\pgfqpoint{4.068755in}{0.580588in}}%
\pgfpathlineto{\pgfqpoint{4.076106in}{0.574997in}}%
\pgfpathlineto{\pgfqpoint{4.074930in}{0.572318in}}%
\pgfpathlineto{\pgfqpoint{4.080161in}{0.570260in}}%
\pgfpathlineto{\pgfqpoint{4.081445in}{0.566900in}}%
\pgfpathlineto{\pgfqpoint{4.088083in}{0.569570in}}%
\pgfpathlineto{\pgfqpoint{4.091007in}{0.566071in}}%
\pgfpathlineto{\pgfqpoint{4.092356in}{0.568189in}}%
\pgfpathlineto{\pgfqpoint{4.088383in}{0.575083in}}%
\pgfpathlineto{\pgfqpoint{4.089795in}{0.577148in}}%
\pgfpathlineto{\pgfqpoint{4.090421in}{0.582594in}}%
\pgfpathlineto{\pgfqpoint{4.088838in}{0.584991in}}%
\pgfpathlineto{\pgfqpoint{4.090025in}{0.588211in}}%
\pgfpathlineto{\pgfqpoint{4.093915in}{0.587908in}}%
\pgfpathlineto{\pgfqpoint{4.093110in}{0.582429in}}%
\pgfpathlineto{\pgfqpoint{4.090627in}{0.580602in}}%
\pgfpathlineto{\pgfqpoint{4.090803in}{0.573231in}}%
\pgfpathlineto{\pgfqpoint{4.095258in}{0.572350in}}%
\pgfpathlineto{\pgfqpoint{4.095312in}{0.566870in}}%
\pgfpathlineto{\pgfqpoint{4.100178in}{0.563315in}}%
\pgfpathlineto{\pgfqpoint{4.102415in}{0.565557in}}%
\pgfpathlineto{\pgfqpoint{4.099970in}{0.572380in}}%
\pgfpathlineto{\pgfqpoint{4.097442in}{0.574106in}}%
\pgfpathlineto{\pgfqpoint{4.094955in}{0.572970in}}%
\pgfpathlineto{\pgfqpoint{4.092767in}{0.574420in}}%
\pgfpathlineto{\pgfqpoint{4.093149in}{0.580725in}}%
\pgfpathlineto{\pgfqpoint{4.096800in}{0.583304in}}%
\pgfpathlineto{\pgfqpoint{4.100395in}{0.583066in}}%
\pgfpathlineto{\pgfqpoint{4.098905in}{0.586634in}}%
\pgfpathlineto{\pgfqpoint{4.093518in}{0.589366in}}%
\pgfpathlineto{\pgfqpoint{4.086383in}{0.600253in}}%
\pgfpathlineto{\pgfqpoint{4.089944in}{0.605475in}}%
\pgfpathlineto{\pgfqpoint{4.091926in}{0.610393in}}%
\pgfpathlineto{\pgfqpoint{4.092082in}{0.620269in}}%
\pgfpathlineto{\pgfqpoint{4.091319in}{0.628601in}}%
\pgfpathlineto{\pgfqpoint{4.089142in}{0.634072in}}%
\pgfpathlineto{\pgfqpoint{4.089576in}{0.639541in}}%
\pgfpathlineto{\pgfqpoint{4.095144in}{0.644285in}}%
\pgfpathlineto{\pgfqpoint{4.100681in}{0.649922in}}%
\pgfpathlineto{\pgfqpoint{4.114564in}{0.638405in}}%
\pgfpathlineto{\pgfqpoint{4.123640in}{0.637556in}}%
\pgfpathlineto{\pgfqpoint{4.131132in}{0.640424in}}%
\pgfpathlineto{\pgfqpoint{4.137708in}{0.644201in}}%
\pgfpathlineto{\pgfqpoint{4.139105in}{0.646000in}}%
\pgfpathlineto{\pgfqpoint{4.155340in}{0.650103in}}%
\pgfpathlineto{\pgfqpoint{4.156749in}{0.647098in}}%
\pgfpathlineto{\pgfqpoint{4.162266in}{0.644638in}}%
\pgfpathlineto{\pgfqpoint{4.172793in}{0.645392in}}%
\pgfpathlineto{\pgfqpoint{4.174716in}{0.644720in}}%
\pgfpathlineto{\pgfqpoint{4.179448in}{0.646515in}}%
\pgfpathlineto{\pgfqpoint{4.181350in}{0.643411in}}%
\pgfpathlineto{\pgfqpoint{4.185894in}{0.640575in}}%
\pgfpathlineto{\pgfqpoint{4.188167in}{0.640896in}}%
\pgfpathlineto{\pgfqpoint{4.192566in}{0.637488in}}%
\pgfpathlineto{\pgfqpoint{4.200000in}{0.638115in}}%
\pgfpathlineto{\pgfqpoint{4.207425in}{0.640874in}}%
\pgfpathlineto{\pgfqpoint{4.213489in}{0.632051in}}%
\pgfpathlineto{\pgfqpoint{4.213019in}{0.630071in}}%
\pgfpathlineto{\pgfqpoint{4.206328in}{0.627796in}}%
\pgfpathlineto{\pgfqpoint{4.208115in}{0.624355in}}%
\pgfpathlineto{\pgfqpoint{4.214275in}{0.627976in}}%
\pgfpathlineto{\pgfqpoint{4.216790in}{0.627100in}}%
\pgfpathlineto{\pgfqpoint{4.217860in}{0.622768in}}%
\pgfpathlineto{\pgfqpoint{4.214183in}{0.620936in}}%
\pgfpathlineto{\pgfqpoint{4.216809in}{0.615461in}}%
\pgfpathlineto{\pgfqpoint{4.220623in}{0.616386in}}%
\pgfpathlineto{\pgfqpoint{4.226157in}{0.613560in}}%
\pgfpathlineto{\pgfqpoint{4.231450in}{0.605718in}}%
\pgfpathlineto{\pgfqpoint{4.226793in}{0.603540in}}%
\pgfpathlineto{\pgfqpoint{4.230907in}{0.597761in}}%
\pgfpathlineto{\pgfqpoint{4.227441in}{0.596410in}}%
\pgfpathlineto{\pgfqpoint{4.231753in}{0.590922in}}%
\pgfpathlineto{\pgfqpoint{4.237012in}{0.591506in}}%
\pgfpathlineto{\pgfqpoint{4.239781in}{0.589037in}}%
\pgfpathlineto{\pgfqpoint{4.246785in}{0.585396in}}%
\pgfpathlineto{\pgfqpoint{4.255820in}{0.572470in}}%
\pgfpathlineto{\pgfqpoint{4.255616in}{0.569379in}}%
\pgfpathlineto{\pgfqpoint{4.269166in}{0.559477in}}%
\pgfpathlineto{\pgfqpoint{4.269458in}{0.555975in}}%
\pgfpathlineto{\pgfqpoint{4.273069in}{0.550735in}}%
\pgfpathlineto{\pgfqpoint{4.284163in}{0.548115in}}%
\pgfpathlineto{\pgfqpoint{4.288277in}{0.546252in}}%
\pgfpathlineto{\pgfqpoint{4.292148in}{0.540476in}}%
\pgfpathlineto{\pgfqpoint{4.292646in}{0.535415in}}%
\pgfpathlineto{\pgfqpoint{4.295973in}{0.531291in}}%
\pgfpathlineto{\pgfqpoint{4.296972in}{0.526640in}}%
\pgfpathlineto{\pgfqpoint{4.300090in}{0.524941in}}%
\pgfpathlineto{\pgfqpoint{4.284646in}{0.497067in}}%
\pgfpathlineto{\pgfqpoint{4.252590in}{0.439200in}}%
\pgfpathlineto{\pgfqpoint{4.205557in}{0.354280in}}%
\pgfpathlineto{\pgfqpoint{4.187143in}{0.321025in}}%
\pgfpathlineto{\pgfqpoint{4.191065in}{0.316559in}}%
\pgfpathlineto{\pgfqpoint{4.192793in}{0.317984in}}%
\pgfpathlineto{\pgfqpoint{4.196296in}{0.312785in}}%
\pgfpathlineto{\pgfqpoint{4.201159in}{0.314293in}}%
\pgfpathlineto{\pgfqpoint{4.207614in}{0.311178in}}%
\pgfpathlineto{\pgfqpoint{4.203433in}{0.306310in}}%
\pgfpathlineto{\pgfqpoint{4.206662in}{0.299741in}}%
\pgfpathlineto{\pgfqpoint{4.205995in}{0.296543in}}%
\pgfpathlineto{\pgfqpoint{4.211145in}{0.279677in}}%
\pgfpathlineto{\pgfqpoint{4.208966in}{0.272007in}}%
\pgfpathlineto{\pgfqpoint{4.218805in}{0.273906in}}%
\pgfpathlineto{\pgfqpoint{4.221221in}{0.272662in}}%
\pgfpathlineto{\pgfqpoint{4.223844in}{0.274685in}}%
\pgfpathlineto{\pgfqpoint{4.225709in}{0.278509in}}%
\pgfpathlineto{\pgfqpoint{4.228381in}{0.280701in}}%
\pgfpathlineto{\pgfqpoint{4.239784in}{0.280424in}}%
\pgfpathlineto{\pgfqpoint{4.242361in}{0.273064in}}%
\pgfpathlineto{\pgfqpoint{4.240149in}{0.266603in}}%
\pgfpathlineto{\pgfqpoint{4.242662in}{0.264471in}}%
\pgfpathlineto{\pgfqpoint{4.243983in}{0.260825in}}%
\pgfpathlineto{\pgfqpoint{4.243695in}{0.253967in}}%
\pgfpathlineto{\pgfqpoint{4.247048in}{0.249051in}}%
\pgfpathlineto{\pgfqpoint{4.248886in}{0.241346in}}%
\pgfpathlineto{\pgfqpoint{4.247942in}{0.237184in}}%
\pgfpathlineto{\pgfqpoint{4.249080in}{0.232966in}}%
\pgfpathlineto{\pgfqpoint{4.250722in}{0.208539in}}%
\pgfpathlineto{\pgfqpoint{4.248370in}{0.206615in}}%
\pgfpathlineto{\pgfqpoint{4.251472in}{0.203957in}}%
\pgfpathlineto{\pgfqpoint{4.249034in}{0.200723in}}%
\pgfpathlineto{\pgfqpoint{4.251233in}{0.198087in}}%
\pgfpathlineto{\pgfqpoint{4.249817in}{0.193586in}}%
\pgfpathlineto{\pgfqpoint{4.252984in}{0.192366in}}%
\pgfpathlineto{\pgfqpoint{4.257058in}{0.185522in}}%
\pgfpathlineto{\pgfqpoint{4.260085in}{0.183460in}}%
\pgfpathlineto{\pgfqpoint{4.260877in}{0.180445in}}%
\pgfpathlineto{\pgfqpoint{4.262631in}{0.179139in}}%
\pgfpathlineto{\pgfqpoint{4.262265in}{0.176636in}}%
\pgfpathlineto{\pgfqpoint{4.266046in}{0.174822in}}%
\pgfpathlineto{\pgfqpoint{4.265112in}{0.170006in}}%
\pgfpathlineto{\pgfqpoint{4.261858in}{0.167483in}}%
\pgfpathlineto{\pgfqpoint{4.260380in}{0.162795in}}%
\pgfpathlineto{\pgfqpoint{4.259988in}{0.156518in}}%
\pgfpathlineto{\pgfqpoint{4.258180in}{0.155917in}}%
\pgfpathlineto{\pgfqpoint{4.252271in}{0.150415in}}%
\pgfpathlineto{\pgfqpoint{4.247019in}{0.148962in}}%
\pgfpathlineto{\pgfqpoint{4.244513in}{0.151181in}}%
\pgfpathlineto{\pgfqpoint{4.246935in}{0.157071in}}%
\pgfpathlineto{\pgfqpoint{4.246195in}{0.159962in}}%
\pgfpathlineto{\pgfqpoint{4.249665in}{0.161426in}}%
\pgfpathlineto{\pgfqpoint{4.252909in}{0.169947in}}%
\pgfpathlineto{\pgfqpoint{4.251539in}{0.177225in}}%
\pgfpathlineto{\pgfqpoint{4.249346in}{0.178943in}}%
\pgfpathlineto{\pgfqpoint{4.242012in}{0.178387in}}%
\pgfpathlineto{\pgfqpoint{4.243490in}{0.176486in}}%
\pgfpathlineto{\pgfqpoint{4.238181in}{0.171028in}}%
\pgfpathlineto{\pgfqpoint{4.236471in}{0.174063in}}%
\pgfpathlineto{\pgfqpoint{4.236916in}{0.178027in}}%
\pgfpathlineto{\pgfqpoint{4.239963in}{0.177954in}}%
\pgfpathlineto{\pgfqpoint{4.242677in}{0.180866in}}%
\pgfpathlineto{\pgfqpoint{4.244282in}{0.185118in}}%
\pgfpathlineto{\pgfqpoint{4.250331in}{0.183918in}}%
\pgfpathlineto{\pgfqpoint{4.245488in}{0.186190in}}%
\pgfpathlineto{\pgfqpoint{4.245915in}{0.189233in}}%
\pgfpathlineto{\pgfqpoint{4.243899in}{0.190607in}}%
\pgfpathlineto{\pgfqpoint{4.243074in}{0.194752in}}%
\pgfpathlineto{\pgfqpoint{4.244505in}{0.196911in}}%
\pgfpathlineto{\pgfqpoint{4.241093in}{0.203617in}}%
\pgfpathlineto{\pgfqpoint{4.241381in}{0.207883in}}%
\pgfpathlineto{\pgfqpoint{4.237383in}{0.211682in}}%
\pgfpathlineto{\pgfqpoint{4.235220in}{0.214756in}}%
\pgfpathlineto{\pgfqpoint{4.238207in}{0.217685in}}%
\pgfpathlineto{\pgfqpoint{4.239313in}{0.222063in}}%
\pgfpathlineto{\pgfqpoint{4.238326in}{0.224937in}}%
\pgfpathlineto{\pgfqpoint{4.239481in}{0.228366in}}%
\pgfpathlineto{\pgfqpoint{4.238009in}{0.239543in}}%
\pgfpathlineto{\pgfqpoint{4.235807in}{0.244614in}}%
\pgfpathlineto{\pgfqpoint{4.233288in}{0.246519in}}%
\pgfpathlineto{\pgfqpoint{4.234360in}{0.253116in}}%
\pgfpathlineto{\pgfqpoint{4.233604in}{0.258268in}}%
\pgfpathlineto{\pgfqpoint{4.234974in}{0.261572in}}%
\pgfpathlineto{\pgfqpoint{4.232140in}{0.262035in}}%
\pgfpathlineto{\pgfqpoint{4.231380in}{0.254009in}}%
\pgfpathlineto{\pgfqpoint{4.228196in}{0.245055in}}%
\pgfpathlineto{\pgfqpoint{4.225639in}{0.247094in}}%
\pgfpathlineto{\pgfqpoint{4.225297in}{0.250037in}}%
\pgfpathlineto{\pgfqpoint{4.222168in}{0.254057in}}%
\pgfpathlineto{\pgfqpoint{4.223758in}{0.256688in}}%
\pgfpathlineto{\pgfqpoint{4.223264in}{0.262613in}}%
\pgfpathlineto{\pgfqpoint{4.219215in}{0.260530in}}%
\pgfpathlineto{\pgfqpoint{4.220009in}{0.257316in}}%
\pgfpathlineto{\pgfqpoint{4.219247in}{0.253236in}}%
\pgfpathlineto{\pgfqpoint{4.214702in}{0.253244in}}%
\pgfpathlineto{\pgfqpoint{4.212665in}{0.255773in}}%
\pgfpathlineto{\pgfqpoint{4.209590in}{0.256255in}}%
\pgfpathlineto{\pgfqpoint{4.207521in}{0.258978in}}%
\pgfpathlineto{\pgfqpoint{4.203961in}{0.266454in}}%
\pgfpathlineto{\pgfqpoint{4.202807in}{0.272259in}}%
\pgfpathlineto{\pgfqpoint{4.203546in}{0.274726in}}%
\pgfpathlineto{\pgfqpoint{4.201822in}{0.280160in}}%
\pgfpathlineto{\pgfqpoint{4.198967in}{0.283305in}}%
\pgfpathlineto{\pgfqpoint{4.194029in}{0.291344in}}%
\pgfpathlineto{\pgfqpoint{4.190447in}{0.298326in}}%
\pgfpathlineto{\pgfqpoint{4.194102in}{0.298431in}}%
\pgfpathlineto{\pgfqpoint{4.196252in}{0.299897in}}%
\pgfpathlineto{\pgfqpoint{4.196675in}{0.304435in}}%
\pgfpathlineto{\pgfqpoint{4.186334in}{0.304983in}}%
\pgfpathlineto{\pgfqpoint{4.177539in}{0.314807in}}%
\pgfpathlineto{\pgfqpoint{4.179852in}{0.317392in}}%
\pgfpathlineto{\pgfqpoint{4.175599in}{0.317833in}}%
\pgfpathlineto{\pgfqpoint{4.167152in}{0.326967in}}%
\pgfpathlineto{\pgfqpoint{4.153770in}{0.331547in}}%
\pgfpathlineto{\pgfqpoint{4.151968in}{0.336915in}}%
\pgfpathlineto{\pgfqpoint{4.149191in}{0.339711in}}%
\pgfpathlineto{\pgfqpoint{4.149065in}{0.342336in}}%
\pgfpathlineto{\pgfqpoint{4.146976in}{0.344286in}}%
\pgfpathlineto{\pgfqpoint{4.149717in}{0.347544in}}%
\pgfpathlineto{\pgfqpoint{4.143465in}{0.347858in}}%
\pgfpathlineto{\pgfqpoint{4.141876in}{0.352075in}}%
\pgfpathlineto{\pgfqpoint{4.139030in}{0.353870in}}%
\pgfpathlineto{\pgfqpoint{4.144729in}{0.356037in}}%
\pgfpathlineto{\pgfqpoint{4.141765in}{0.359902in}}%
\pgfpathlineto{\pgfqpoint{4.136311in}{0.361524in}}%
\pgfpathlineto{\pgfqpoint{4.137160in}{0.368663in}}%
\pgfpathlineto{\pgfqpoint{4.135596in}{0.371190in}}%
\pgfpathlineto{\pgfqpoint{4.132780in}{0.372970in}}%
\pgfpathlineto{\pgfqpoint{4.130372in}{0.371185in}}%
\pgfpathlineto{\pgfqpoint{4.128950in}{0.372864in}}%
\pgfpathlineto{\pgfqpoint{4.124800in}{0.373345in}}%
\pgfpathlineto{\pgfqpoint{4.125134in}{0.379823in}}%
\pgfpathlineto{\pgfqpoint{4.119119in}{0.375448in}}%
\pgfpathlineto{\pgfqpoint{4.119246in}{0.366430in}}%
\pgfpathlineto{\pgfqpoint{4.112788in}{0.364702in}}%
\pgfpathlineto{\pgfqpoint{4.108572in}{0.361101in}}%
\pgfpathlineto{\pgfqpoint{4.105182in}{0.360342in}}%
\pgfpathlineto{\pgfqpoint{4.101514in}{0.363971in}}%
\pgfpathlineto{\pgfqpoint{4.098248in}{0.363490in}}%
\pgfpathlineto{\pgfqpoint{4.098854in}{0.366805in}}%
\pgfpathlineto{\pgfqpoint{4.094211in}{0.365020in}}%
\pgfpathlineto{\pgfqpoint{4.089909in}{0.364820in}}%
\pgfpathlineto{\pgfqpoint{4.084874in}{0.363010in}}%
\pgfpathlineto{\pgfqpoint{4.074656in}{0.364028in}}%
\pgfpathlineto{\pgfqpoint{4.071376in}{0.363358in}}%
\pgfpathlineto{\pgfqpoint{4.069009in}{0.366339in}}%
\pgfpathlineto{\pgfqpoint{4.064640in}{0.366637in}}%
\pgfpathlineto{\pgfqpoint{4.063731in}{0.370500in}}%
\pgfpathlineto{\pgfqpoint{4.066335in}{0.372604in}}%
\pgfpathlineto{\pgfqpoint{4.069390in}{0.372406in}}%
\pgfpathlineto{\pgfqpoint{4.071940in}{0.369537in}}%
\pgfpathlineto{\pgfqpoint{4.073619in}{0.373944in}}%
\pgfpathlineto{\pgfqpoint{4.071276in}{0.378827in}}%
\pgfpathlineto{\pgfqpoint{4.076579in}{0.383125in}}%
\pgfpathlineto{\pgfqpoint{4.081893in}{0.384709in}}%
\pgfpathlineto{\pgfqpoint{4.085593in}{0.387292in}}%
\pgfpathlineto{\pgfqpoint{4.087990in}{0.390076in}}%
\pgfpathlineto{\pgfqpoint{4.089327in}{0.394638in}}%
\pgfpathlineto{\pgfqpoint{4.093552in}{0.393396in}}%
\pgfpathlineto{\pgfqpoint{4.102714in}{0.394034in}}%
\pgfpathlineto{\pgfqpoint{4.104680in}{0.389032in}}%
\pgfpathlineto{\pgfqpoint{4.107894in}{0.388814in}}%
\pgfpathlineto{\pgfqpoint{4.113742in}{0.381156in}}%
\pgfpathlineto{\pgfqpoint{4.113753in}{0.383947in}}%
\pgfpathlineto{\pgfqpoint{4.111710in}{0.384855in}}%
\pgfpathlineto{\pgfqpoint{4.107971in}{0.393996in}}%
\pgfpathlineto{\pgfqpoint{4.109790in}{0.395238in}}%
\pgfpathlineto{\pgfqpoint{4.104797in}{0.400007in}}%
\pgfpathlineto{\pgfqpoint{4.100334in}{0.400936in}}%
\pgfpathlineto{\pgfqpoint{4.096081in}{0.399287in}}%
\pgfpathlineto{\pgfqpoint{4.088798in}{0.400571in}}%
\pgfpathlineto{\pgfqpoint{4.085344in}{0.398261in}}%
\pgfpathlineto{\pgfqpoint{4.077478in}{0.396476in}}%
\pgfpathlineto{\pgfqpoint{4.076759in}{0.393907in}}%
\pgfpathlineto{\pgfqpoint{4.070716in}{0.393207in}}%
\pgfpathlineto{\pgfqpoint{4.069165in}{0.390270in}}%
\pgfpathlineto{\pgfqpoint{4.065709in}{0.388119in}}%
\pgfpathlineto{\pgfqpoint{4.061533in}{0.388489in}}%
\pgfpathlineto{\pgfqpoint{4.059448in}{0.386408in}}%
\pgfpathlineto{\pgfqpoint{4.056818in}{0.386506in}}%
\pgfpathlineto{\pgfqpoint{4.054304in}{0.388378in}}%
\pgfpathlineto{\pgfqpoint{4.049909in}{0.388044in}}%
\pgfpathlineto{\pgfqpoint{4.048874in}{0.386318in}}%
\pgfpathlineto{\pgfqpoint{4.044258in}{0.388090in}}%
\pgfpathlineto{\pgfqpoint{4.040339in}{0.383518in}}%
\pgfpathlineto{\pgfqpoint{4.044097in}{0.379095in}}%
\pgfpathlineto{\pgfqpoint{4.045543in}{0.375457in}}%
\pgfpathlineto{\pgfqpoint{4.044210in}{0.372376in}}%
\pgfpathlineto{\pgfqpoint{4.038870in}{0.370250in}}%
\pgfpathlineto{\pgfqpoint{4.035776in}{0.371980in}}%
\pgfpathlineto{\pgfqpoint{4.031882in}{0.371043in}}%
\pgfpathlineto{\pgfqpoint{4.031413in}{0.368343in}}%
\pgfpathlineto{\pgfqpoint{4.026673in}{0.369298in}}%
\pgfpathlineto{\pgfqpoint{4.027722in}{0.366600in}}%
\pgfpathlineto{\pgfqpoint{4.020731in}{0.366031in}}%
\pgfpathlineto{\pgfqpoint{4.016205in}{0.369456in}}%
\pgfpathlineto{\pgfqpoint{4.013787in}{0.367278in}}%
\pgfpathlineto{\pgfqpoint{4.007354in}{0.369597in}}%
\pgfpathlineto{\pgfqpoint{4.005752in}{0.367248in}}%
\pgfpathlineto{\pgfqpoint{4.002563in}{0.366203in}}%
\pgfpathlineto{\pgfqpoint{4.000005in}{0.368957in}}%
\pgfpathlineto{\pgfqpoint{3.991257in}{0.368272in}}%
\pgfpathlineto{\pgfqpoint{3.991266in}{0.364211in}}%
\pgfpathlineto{\pgfqpoint{3.987752in}{0.363136in}}%
\pgfpathlineto{\pgfqpoint{3.985018in}{0.365224in}}%
\pgfpathlineto{\pgfqpoint{3.981404in}{0.364064in}}%
\pgfpathlineto{\pgfqpoint{3.976727in}{0.363787in}}%
\pgfpathlineto{\pgfqpoint{3.969759in}{0.366724in}}%
\pgfpathlineto{\pgfqpoint{3.965769in}{0.363696in}}%
\pgfpathlineto{\pgfqpoint{3.960063in}{0.367743in}}%
\pgfpathlineto{\pgfqpoint{3.958197in}{0.366209in}}%
\pgfpathlineto{\pgfqpoint{3.956887in}{0.362246in}}%
\pgfpathlineto{\pgfqpoint{3.953218in}{0.359922in}}%
\pgfpathlineto{\pgfqpoint{3.947236in}{0.362582in}}%
\pgfpathlineto{\pgfqpoint{3.937326in}{0.369153in}}%
\pgfpathlineto{\pgfqpoint{3.934438in}{0.369737in}}%
\pgfpathlineto{\pgfqpoint{3.932808in}{0.368368in}}%
\pgfpathlineto{\pgfqpoint{3.928449in}{0.369299in}}%
\pgfpathlineto{\pgfqpoint{3.926322in}{0.371344in}}%
\pgfpathlineto{\pgfqpoint{3.921969in}{0.372446in}}%
\pgfpathlineto{\pgfqpoint{3.916023in}{0.372510in}}%
\pgfpathlineto{\pgfqpoint{3.913662in}{0.374950in}}%
\pgfpathlineto{\pgfqpoint{3.918187in}{0.377474in}}%
\pgfpathlineto{\pgfqpoint{3.917131in}{0.379966in}}%
\pgfpathlineto{\pgfqpoint{3.914022in}{0.379406in}}%
\pgfpathlineto{\pgfqpoint{3.912470in}{0.377447in}}%
\pgfpathlineto{\pgfqpoint{3.905192in}{0.374639in}}%
\pgfpathlineto{\pgfqpoint{3.901512in}{0.375594in}}%
\pgfpathlineto{\pgfqpoint{3.899964in}{0.381297in}}%
\pgfpathlineto{\pgfqpoint{3.895373in}{0.378845in}}%
\pgfpathlineto{\pgfqpoint{3.894039in}{0.386073in}}%
\pgfpathlineto{\pgfqpoint{3.891864in}{0.384210in}}%
\pgfpathlineto{\pgfqpoint{3.892228in}{0.381429in}}%
\pgfpathlineto{\pgfqpoint{3.887226in}{0.382163in}}%
\pgfpathlineto{\pgfqpoint{3.892429in}{0.387010in}}%
\pgfpathlineto{\pgfqpoint{3.897702in}{0.384130in}}%
\pgfpathlineto{\pgfqpoint{3.898833in}{0.385576in}}%
\pgfpathlineto{\pgfqpoint{3.904380in}{0.383843in}}%
\pgfpathlineto{\pgfqpoint{3.904033in}{0.385816in}}%
\pgfpathlineto{\pgfqpoint{3.919168in}{0.386469in}}%
\pgfpathlineto{\pgfqpoint{3.926551in}{0.383391in}}%
\pgfpathlineto{\pgfqpoint{3.929408in}{0.380355in}}%
\pgfpathlineto{\pgfqpoint{3.928561in}{0.377705in}}%
\pgfpathlineto{\pgfqpoint{3.931429in}{0.376702in}}%
\pgfpathlineto{\pgfqpoint{3.938558in}{0.380885in}}%
\pgfpathlineto{\pgfqpoint{3.947779in}{0.381153in}}%
\pgfpathlineto{\pgfqpoint{3.956155in}{0.378631in}}%
\pgfpathlineto{\pgfqpoint{3.960899in}{0.379054in}}%
\pgfpathlineto{\pgfqpoint{3.962665in}{0.375142in}}%
\pgfpathlineto{\pgfqpoint{3.965318in}{0.379455in}}%
\pgfpathlineto{\pgfqpoint{3.966950in}{0.380389in}}%
\pgfpathlineto{\pgfqpoint{3.974498in}{0.382264in}}%
\pgfpathlineto{\pgfqpoint{3.977304in}{0.381118in}}%
\pgfpathlineto{\pgfqpoint{3.980338in}{0.382479in}}%
\pgfpathlineto{\pgfqpoint{3.983958in}{0.381767in}}%
\pgfpathlineto{\pgfqpoint{3.984812in}{0.383352in}}%
\pgfpathlineto{\pgfqpoint{3.992363in}{0.390625in}}%
\pgfpathlineto{\pgfqpoint{3.995561in}{0.391703in}}%
\pgfpathlineto{\pgfqpoint{3.997451in}{0.395629in}}%
\pgfpathlineto{\pgfqpoint{4.007481in}{0.398668in}}%
\pgfpathlineto{\pgfqpoint{4.008737in}{0.400987in}}%
\pgfpathlineto{\pgfqpoint{3.994522in}{0.404566in}}%
\pgfpathlineto{\pgfqpoint{3.995081in}{0.408415in}}%
\pgfpathlineto{\pgfqpoint{3.993619in}{0.413499in}}%
\pgfpathlineto{\pgfqpoint{3.990309in}{0.413077in}}%
\pgfpathlineto{\pgfqpoint{3.988057in}{0.406479in}}%
\pgfpathlineto{\pgfqpoint{3.985443in}{0.405936in}}%
\pgfpathlineto{\pgfqpoint{3.983742in}{0.408071in}}%
\pgfpathlineto{\pgfqpoint{3.985726in}{0.417413in}}%
\pgfpathlineto{\pgfqpoint{3.983750in}{0.421328in}}%
\pgfpathlineto{\pgfqpoint{3.979922in}{0.425943in}}%
\pgfpathlineto{\pgfqpoint{3.981762in}{0.429465in}}%
\pgfpathlineto{\pgfqpoint{3.973683in}{0.429591in}}%
\pgfpathlineto{\pgfqpoint{3.973867in}{0.430958in}}%
\pgfpathlineto{\pgfqpoint{3.969046in}{0.432871in}}%
\pgfpathlineto{\pgfqpoint{3.967548in}{0.431818in}}%
\pgfpathclose%
\pgfusepath{fill}%
\end{pgfscope}%
\begin{pgfscope}%
\pgfpathrectangle{\pgfqpoint{3.525000in}{0.100000in}}{\pgfqpoint{2.862985in}{1.832778in}}%
\pgfusepath{clip}%
\pgfsetbuttcap%
\pgfsetmiterjoin%
\definecolor{currentfill}{rgb}{0.645872,0.120898,0.598867}%
\pgfsetfillcolor{currentfill}%
\pgfsetlinewidth{0.000000pt}%
\definecolor{currentstroke}{rgb}{0.000000,0.000000,0.000000}%
\pgfsetstrokecolor{currentstroke}%
\pgfsetstrokeopacity{0.000000}%
\pgfsetdash{}{0pt}%
\pgfpathmoveto{\pgfqpoint{3.953131in}{0.499862in}}%
\pgfpathlineto{\pgfqpoint{3.952132in}{0.498106in}}%
\pgfpathlineto{\pgfqpoint{3.954775in}{0.494368in}}%
\pgfpathlineto{\pgfqpoint{3.950597in}{0.490871in}}%
\pgfpathlineto{\pgfqpoint{3.950713in}{0.487863in}}%
\pgfpathlineto{\pgfqpoint{3.948824in}{0.487153in}}%
\pgfpathlineto{\pgfqpoint{3.942331in}{0.491639in}}%
\pgfpathlineto{\pgfqpoint{3.937703in}{0.501458in}}%
\pgfpathlineto{\pgfqpoint{3.937393in}{0.504295in}}%
\pgfpathlineto{\pgfqpoint{3.938583in}{0.507691in}}%
\pgfpathlineto{\pgfqpoint{3.943598in}{0.502496in}}%
\pgfpathlineto{\pgfqpoint{3.945132in}{0.503498in}}%
\pgfpathlineto{\pgfqpoint{3.949386in}{0.502545in}}%
\pgfpathlineto{\pgfqpoint{3.953131in}{0.499862in}}%
\pgfpathclose%
\pgfusepath{fill}%
\end{pgfscope}%
\begin{pgfscope}%
\pgfpathrectangle{\pgfqpoint{3.525000in}{0.100000in}}{\pgfqpoint{2.862985in}{1.832778in}}%
\pgfusepath{clip}%
\pgfsetbuttcap%
\pgfsetmiterjoin%
\definecolor{currentfill}{rgb}{0.645872,0.120898,0.598867}%
\pgfsetfillcolor{currentfill}%
\pgfsetlinewidth{0.000000pt}%
\definecolor{currentstroke}{rgb}{0.000000,0.000000,0.000000}%
\pgfsetstrokecolor{currentstroke}%
\pgfsetstrokeopacity{0.000000}%
\pgfsetdash{}{0pt}%
\pgfpathmoveto{\pgfqpoint{3.876126in}{0.386109in}}%
\pgfpathlineto{\pgfqpoint{3.871562in}{0.385368in}}%
\pgfpathlineto{\pgfqpoint{3.868260in}{0.386895in}}%
\pgfpathlineto{\pgfqpoint{3.866816in}{0.389380in}}%
\pgfpathlineto{\pgfqpoint{3.868453in}{0.392945in}}%
\pgfpathlineto{\pgfqpoint{3.872096in}{0.392002in}}%
\pgfpathlineto{\pgfqpoint{3.878086in}{0.394117in}}%
\pgfpathlineto{\pgfqpoint{3.880184in}{0.391196in}}%
\pgfpathlineto{\pgfqpoint{3.882217in}{0.391735in}}%
\pgfpathlineto{\pgfqpoint{3.887141in}{0.389917in}}%
\pgfpathlineto{\pgfqpoint{3.889268in}{0.387608in}}%
\pgfpathlineto{\pgfqpoint{3.886672in}{0.381567in}}%
\pgfpathlineto{\pgfqpoint{3.884090in}{0.380286in}}%
\pgfpathlineto{\pgfqpoint{3.881769in}{0.380976in}}%
\pgfpathlineto{\pgfqpoint{3.876126in}{0.386109in}}%
\pgfpathclose%
\pgfusepath{fill}%
\end{pgfscope}%
\begin{pgfscope}%
\pgfpathrectangle{\pgfqpoint{3.525000in}{0.100000in}}{\pgfqpoint{2.862985in}{1.832778in}}%
\pgfusepath{clip}%
\pgfsetbuttcap%
\pgfsetmiterjoin%
\definecolor{currentfill}{rgb}{0.645872,0.120898,0.598867}%
\pgfsetfillcolor{currentfill}%
\pgfsetlinewidth{0.000000pt}%
\definecolor{currentstroke}{rgb}{0.000000,0.000000,0.000000}%
\pgfsetstrokecolor{currentstroke}%
\pgfsetstrokeopacity{0.000000}%
\pgfsetdash{}{0pt}%
\pgfpathmoveto{\pgfqpoint{3.835323in}{0.391598in}}%
\pgfpathlineto{\pgfqpoint{3.828857in}{0.394682in}}%
\pgfpathlineto{\pgfqpoint{3.822873in}{0.395600in}}%
\pgfpathlineto{\pgfqpoint{3.820987in}{0.396913in}}%
\pgfpathlineto{\pgfqpoint{3.823135in}{0.399499in}}%
\pgfpathlineto{\pgfqpoint{3.826946in}{0.396269in}}%
\pgfpathlineto{\pgfqpoint{3.830142in}{0.396714in}}%
\pgfpathlineto{\pgfqpoint{3.835341in}{0.398951in}}%
\pgfpathlineto{\pgfqpoint{3.835761in}{0.401832in}}%
\pgfpathlineto{\pgfqpoint{3.842146in}{0.400406in}}%
\pgfpathlineto{\pgfqpoint{3.840651in}{0.396069in}}%
\pgfpathlineto{\pgfqpoint{3.844719in}{0.395751in}}%
\pgfpathlineto{\pgfqpoint{3.841306in}{0.390592in}}%
\pgfpathlineto{\pgfqpoint{3.838001in}{0.392264in}}%
\pgfpathlineto{\pgfqpoint{3.835323in}{0.391598in}}%
\pgfpathclose%
\pgfusepath{fill}%
\end{pgfscope}%
\begin{pgfscope}%
\pgfpathrectangle{\pgfqpoint{3.525000in}{0.100000in}}{\pgfqpoint{2.862985in}{1.832778in}}%
\pgfusepath{clip}%
\pgfsetbuttcap%
\pgfsetmiterjoin%
\definecolor{currentfill}{rgb}{0.645872,0.120898,0.598867}%
\pgfsetfillcolor{currentfill}%
\pgfsetlinewidth{0.000000pt}%
\definecolor{currentstroke}{rgb}{0.000000,0.000000,0.000000}%
\pgfsetstrokecolor{currentstroke}%
\pgfsetstrokeopacity{0.000000}%
\pgfsetdash{}{0pt}%
\pgfpathmoveto{\pgfqpoint{3.823962in}{0.402261in}}%
\pgfpathlineto{\pgfqpoint{3.821586in}{0.400969in}}%
\pgfpathlineto{\pgfqpoint{3.815110in}{0.402955in}}%
\pgfpathlineto{\pgfqpoint{3.810271in}{0.401805in}}%
\pgfpathlineto{\pgfqpoint{3.807813in}{0.405073in}}%
\pgfpathlineto{\pgfqpoint{3.808885in}{0.406569in}}%
\pgfpathlineto{\pgfqpoint{3.812552in}{0.406819in}}%
\pgfpathlineto{\pgfqpoint{3.815357in}{0.405880in}}%
\pgfpathlineto{\pgfqpoint{3.818414in}{0.407367in}}%
\pgfpathlineto{\pgfqpoint{3.823100in}{0.405378in}}%
\pgfpathlineto{\pgfqpoint{3.823962in}{0.402261in}}%
\pgfpathclose%
\pgfusepath{fill}%
\end{pgfscope}%
\begin{pgfscope}%
\pgfpathrectangle{\pgfqpoint{3.525000in}{0.100000in}}{\pgfqpoint{2.862985in}{1.832778in}}%
\pgfusepath{clip}%
\pgfsetbuttcap%
\pgfsetmiterjoin%
\definecolor{currentfill}{rgb}{0.645872,0.120898,0.598867}%
\pgfsetfillcolor{currentfill}%
\pgfsetlinewidth{0.000000pt}%
\definecolor{currentstroke}{rgb}{0.000000,0.000000,0.000000}%
\pgfsetstrokecolor{currentstroke}%
\pgfsetstrokeopacity{0.000000}%
\pgfsetdash{}{0pt}%
\pgfpathmoveto{\pgfqpoint{3.747628in}{0.450357in}}%
\pgfpathlineto{\pgfqpoint{3.747456in}{0.447153in}}%
\pgfpathlineto{\pgfqpoint{3.742729in}{0.445139in}}%
\pgfpathlineto{\pgfqpoint{3.738412in}{0.450180in}}%
\pgfpathlineto{\pgfqpoint{3.743253in}{0.451860in}}%
\pgfpathlineto{\pgfqpoint{3.747628in}{0.450357in}}%
\pgfpathclose%
\pgfusepath{fill}%
\end{pgfscope}%
\begin{pgfscope}%
\pgfpathrectangle{\pgfqpoint{3.525000in}{0.100000in}}{\pgfqpoint{2.862985in}{1.832778in}}%
\pgfusepath{clip}%
\pgfsetbuttcap%
\pgfsetmiterjoin%
\definecolor{currentfill}{rgb}{0.645872,0.120898,0.598867}%
\pgfsetfillcolor{currentfill}%
\pgfsetlinewidth{0.000000pt}%
\definecolor{currentstroke}{rgb}{0.000000,0.000000,0.000000}%
\pgfsetstrokecolor{currentstroke}%
\pgfsetstrokeopacity{0.000000}%
\pgfsetdash{}{0pt}%
\pgfpathmoveto{\pgfqpoint{3.708638in}{0.468371in}}%
\pgfpathlineto{\pgfqpoint{3.709642in}{0.469764in}}%
\pgfpathlineto{\pgfqpoint{3.715440in}{0.469524in}}%
\pgfpathlineto{\pgfqpoint{3.716458in}{0.471597in}}%
\pgfpathlineto{\pgfqpoint{3.719002in}{0.469788in}}%
\pgfpathlineto{\pgfqpoint{3.717533in}{0.466308in}}%
\pgfpathlineto{\pgfqpoint{3.715372in}{0.466014in}}%
\pgfpathlineto{\pgfqpoint{3.710207in}{0.467106in}}%
\pgfpathlineto{\pgfqpoint{3.708638in}{0.468371in}}%
\pgfpathclose%
\pgfusepath{fill}%
\end{pgfscope}%
\begin{pgfscope}%
\pgfpathrectangle{\pgfqpoint{3.525000in}{0.100000in}}{\pgfqpoint{2.862985in}{1.832778in}}%
\pgfusepath{clip}%
\pgfsetbuttcap%
\pgfsetmiterjoin%
\definecolor{currentfill}{rgb}{0.645872,0.120898,0.598867}%
\pgfsetfillcolor{currentfill}%
\pgfsetlinewidth{0.000000pt}%
\definecolor{currentstroke}{rgb}{0.000000,0.000000,0.000000}%
\pgfsetstrokecolor{currentstroke}%
\pgfsetstrokeopacity{0.000000}%
\pgfsetdash{}{0pt}%
\pgfpathmoveto{\pgfqpoint{3.700415in}{0.479156in}}%
\pgfpathlineto{\pgfqpoint{3.700063in}{0.482294in}}%
\pgfpathlineto{\pgfqpoint{3.703088in}{0.482062in}}%
\pgfpathlineto{\pgfqpoint{3.702878in}{0.486470in}}%
\pgfpathlineto{\pgfqpoint{3.705921in}{0.484091in}}%
\pgfpathlineto{\pgfqpoint{3.704868in}{0.480653in}}%
\pgfpathlineto{\pgfqpoint{3.700415in}{0.479156in}}%
\pgfpathclose%
\pgfusepath{fill}%
\end{pgfscope}%
\begin{pgfscope}%
\pgfpathrectangle{\pgfqpoint{3.525000in}{0.100000in}}{\pgfqpoint{2.862985in}{1.832778in}}%
\pgfusepath{clip}%
\pgfsetbuttcap%
\pgfsetmiterjoin%
\definecolor{currentfill}{rgb}{0.645872,0.120898,0.598867}%
\pgfsetfillcolor{currentfill}%
\pgfsetlinewidth{0.000000pt}%
\definecolor{currentstroke}{rgb}{0.000000,0.000000,0.000000}%
\pgfsetstrokecolor{currentstroke}%
\pgfsetstrokeopacity{0.000000}%
\pgfsetdash{}{0pt}%
\pgfpathmoveto{\pgfqpoint{3.966980in}{0.591016in}}%
\pgfpathlineto{\pgfqpoint{3.962154in}{0.592058in}}%
\pgfpathlineto{\pgfqpoint{3.961118in}{0.595283in}}%
\pgfpathlineto{\pgfqpoint{3.962776in}{0.598507in}}%
\pgfpathlineto{\pgfqpoint{3.966491in}{0.600190in}}%
\pgfpathlineto{\pgfqpoint{3.970744in}{0.591916in}}%
\pgfpathlineto{\pgfqpoint{3.974601in}{0.591548in}}%
\pgfpathlineto{\pgfqpoint{3.977479in}{0.588995in}}%
\pgfpathlineto{\pgfqpoint{3.977422in}{0.585482in}}%
\pgfpathlineto{\pgfqpoint{3.975318in}{0.583797in}}%
\pgfpathlineto{\pgfqpoint{3.978740in}{0.573408in}}%
\pgfpathlineto{\pgfqpoint{3.982484in}{0.569595in}}%
\pgfpathlineto{\pgfqpoint{3.978656in}{0.568385in}}%
\pgfpathlineto{\pgfqpoint{3.975557in}{0.572663in}}%
\pgfpathlineto{\pgfqpoint{3.968823in}{0.571338in}}%
\pgfpathlineto{\pgfqpoint{3.970918in}{0.575572in}}%
\pgfpathlineto{\pgfqpoint{3.969943in}{0.577258in}}%
\pgfpathlineto{\pgfqpoint{3.970514in}{0.581671in}}%
\pgfpathlineto{\pgfqpoint{3.969180in}{0.589648in}}%
\pgfpathlineto{\pgfqpoint{3.966980in}{0.591016in}}%
\pgfpathclose%
\pgfusepath{fill}%
\end{pgfscope}%
\begin{pgfscope}%
\pgfpathrectangle{\pgfqpoint{3.525000in}{0.100000in}}{\pgfqpoint{2.862985in}{1.832778in}}%
\pgfusepath{clip}%
\pgfsetbuttcap%
\pgfsetmiterjoin%
\definecolor{currentfill}{rgb}{0.645872,0.120898,0.598867}%
\pgfsetfillcolor{currentfill}%
\pgfsetlinewidth{0.000000pt}%
\definecolor{currentstroke}{rgb}{0.000000,0.000000,0.000000}%
\pgfsetstrokecolor{currentstroke}%
\pgfsetstrokeopacity{0.000000}%
\pgfsetdash{}{0pt}%
\pgfpathmoveto{\pgfqpoint{4.136519in}{0.353334in}}%
\pgfpathlineto{\pgfqpoint{4.129523in}{0.353412in}}%
\pgfpathlineto{\pgfqpoint{4.130062in}{0.356924in}}%
\pgfpathlineto{\pgfqpoint{4.132862in}{0.358180in}}%
\pgfpathlineto{\pgfqpoint{4.136519in}{0.353334in}}%
\pgfpathclose%
\pgfusepath{fill}%
\end{pgfscope}%
\begin{pgfscope}%
\pgfpathrectangle{\pgfqpoint{3.525000in}{0.100000in}}{\pgfqpoint{2.862985in}{1.832778in}}%
\pgfusepath{clip}%
\pgfsetbuttcap%
\pgfsetmiterjoin%
\definecolor{currentfill}{rgb}{0.645872,0.120898,0.598867}%
\pgfsetfillcolor{currentfill}%
\pgfsetlinewidth{0.000000pt}%
\definecolor{currentstroke}{rgb}{0.000000,0.000000,0.000000}%
\pgfsetstrokecolor{currentstroke}%
\pgfsetstrokeopacity{0.000000}%
\pgfsetdash{}{0pt}%
\pgfpathmoveto{\pgfqpoint{4.126812in}{0.356514in}}%
\pgfpathlineto{\pgfqpoint{4.121800in}{0.355301in}}%
\pgfpathlineto{\pgfqpoint{4.117038in}{0.353224in}}%
\pgfpathlineto{\pgfqpoint{4.115567in}{0.355671in}}%
\pgfpathlineto{\pgfqpoint{4.123757in}{0.357328in}}%
\pgfpathlineto{\pgfqpoint{4.126812in}{0.356514in}}%
\pgfpathclose%
\pgfusepath{fill}%
\end{pgfscope}%
\begin{pgfscope}%
\pgfpathrectangle{\pgfqpoint{3.525000in}{0.100000in}}{\pgfqpoint{2.862985in}{1.832778in}}%
\pgfusepath{clip}%
\pgfsetbuttcap%
\pgfsetmiterjoin%
\definecolor{currentfill}{rgb}{0.645872,0.120898,0.598867}%
\pgfsetfillcolor{currentfill}%
\pgfsetlinewidth{0.000000pt}%
\definecolor{currentstroke}{rgb}{0.000000,0.000000,0.000000}%
\pgfsetstrokecolor{currentstroke}%
\pgfsetstrokeopacity{0.000000}%
\pgfsetdash{}{0pt}%
\pgfpathmoveto{\pgfqpoint{4.049990in}{0.354028in}}%
\pgfpathlineto{\pgfqpoint{4.050182in}{0.351599in}}%
\pgfpathlineto{\pgfqpoint{4.047294in}{0.350090in}}%
\pgfpathlineto{\pgfqpoint{4.039235in}{0.353552in}}%
\pgfpathlineto{\pgfqpoint{4.037986in}{0.352159in}}%
\pgfpathlineto{\pgfqpoint{4.036036in}{0.358747in}}%
\pgfpathlineto{\pgfqpoint{4.040074in}{0.359664in}}%
\pgfpathlineto{\pgfqpoint{4.041940in}{0.356957in}}%
\pgfpathlineto{\pgfqpoint{4.045992in}{0.360345in}}%
\pgfpathlineto{\pgfqpoint{4.049990in}{0.354028in}}%
\pgfpathclose%
\pgfusepath{fill}%
\end{pgfscope}%
\begin{pgfscope}%
\pgfpathrectangle{\pgfqpoint{3.525000in}{0.100000in}}{\pgfqpoint{2.862985in}{1.832778in}}%
\pgfusepath{clip}%
\pgfsetbuttcap%
\pgfsetmiterjoin%
\definecolor{currentfill}{rgb}{0.645872,0.120898,0.598867}%
\pgfsetfillcolor{currentfill}%
\pgfsetlinewidth{0.000000pt}%
\definecolor{currentstroke}{rgb}{0.000000,0.000000,0.000000}%
\pgfsetstrokecolor{currentstroke}%
\pgfsetstrokeopacity{0.000000}%
\pgfsetdash{}{0pt}%
\pgfpathmoveto{\pgfqpoint{4.231774in}{0.219412in}}%
\pgfpathlineto{\pgfqpoint{4.225017in}{0.216896in}}%
\pgfpathlineto{\pgfqpoint{4.221437in}{0.216745in}}%
\pgfpathlineto{\pgfqpoint{4.224350in}{0.222082in}}%
\pgfpathlineto{\pgfqpoint{4.226840in}{0.224073in}}%
\pgfpathlineto{\pgfqpoint{4.226760in}{0.229600in}}%
\pgfpathlineto{\pgfqpoint{4.229465in}{0.239578in}}%
\pgfpathlineto{\pgfqpoint{4.231957in}{0.241473in}}%
\pgfpathlineto{\pgfqpoint{4.237710in}{0.238729in}}%
\pgfpathlineto{\pgfqpoint{4.236870in}{0.237158in}}%
\pgfpathlineto{\pgfqpoint{4.237306in}{0.229604in}}%
\pgfpathlineto{\pgfqpoint{4.235440in}{0.227566in}}%
\pgfpathlineto{\pgfqpoint{4.234640in}{0.222201in}}%
\pgfpathlineto{\pgfqpoint{4.231774in}{0.219412in}}%
\pgfpathclose%
\pgfusepath{fill}%
\end{pgfscope}%
\begin{pgfscope}%
\pgfpathrectangle{\pgfqpoint{3.525000in}{0.100000in}}{\pgfqpoint{2.862985in}{1.832778in}}%
\pgfusepath{clip}%
\pgfsetbuttcap%
\pgfsetmiterjoin%
\definecolor{currentfill}{rgb}{0.645872,0.120898,0.598867}%
\pgfsetfillcolor{currentfill}%
\pgfsetlinewidth{0.000000pt}%
\definecolor{currentstroke}{rgb}{0.000000,0.000000,0.000000}%
\pgfsetstrokecolor{currentstroke}%
\pgfsetstrokeopacity{0.000000}%
\pgfsetdash{}{0pt}%
\pgfpathmoveto{\pgfqpoint{4.216809in}{0.244146in}}%
\pgfpathlineto{\pgfqpoint{4.220768in}{0.237310in}}%
\pgfpathlineto{\pgfqpoint{4.220034in}{0.235680in}}%
\pgfpathlineto{\pgfqpoint{4.225329in}{0.234022in}}%
\pgfpathlineto{\pgfqpoint{4.223809in}{0.227738in}}%
\pgfpathlineto{\pgfqpoint{4.221503in}{0.228029in}}%
\pgfpathlineto{\pgfqpoint{4.215445in}{0.238982in}}%
\pgfpathlineto{\pgfqpoint{4.216307in}{0.234052in}}%
\pgfpathlineto{\pgfqpoint{4.215275in}{0.231213in}}%
\pgfpathlineto{\pgfqpoint{4.212709in}{0.229950in}}%
\pgfpathlineto{\pgfqpoint{4.211251in}{0.231311in}}%
\pgfpathlineto{\pgfqpoint{4.210437in}{0.237623in}}%
\pgfpathlineto{\pgfqpoint{4.210917in}{0.241717in}}%
\pgfpathlineto{\pgfqpoint{4.209208in}{0.243504in}}%
\pgfpathlineto{\pgfqpoint{4.213722in}{0.248675in}}%
\pgfpathlineto{\pgfqpoint{4.216684in}{0.250209in}}%
\pgfpathlineto{\pgfqpoint{4.217670in}{0.248193in}}%
\pgfpathlineto{\pgfqpoint{4.220760in}{0.249753in}}%
\pgfpathlineto{\pgfqpoint{4.222999in}{0.245516in}}%
\pgfpathlineto{\pgfqpoint{4.227840in}{0.240126in}}%
\pgfpathlineto{\pgfqpoint{4.227179in}{0.237684in}}%
\pgfpathlineto{\pgfqpoint{4.224491in}{0.235024in}}%
\pgfpathlineto{\pgfqpoint{4.222005in}{0.236278in}}%
\pgfpathlineto{\pgfqpoint{4.216809in}{0.244146in}}%
\pgfpathclose%
\pgfusepath{fill}%
\end{pgfscope}%
\begin{pgfscope}%
\pgfpathrectangle{\pgfqpoint{3.525000in}{0.100000in}}{\pgfqpoint{2.862985in}{1.832778in}}%
\pgfusepath{clip}%
\pgfsetbuttcap%
\pgfsetmiterjoin%
\definecolor{currentfill}{rgb}{0.645872,0.120898,0.598867}%
\pgfsetfillcolor{currentfill}%
\pgfsetlinewidth{0.000000pt}%
\definecolor{currentstroke}{rgb}{0.000000,0.000000,0.000000}%
\pgfsetstrokecolor{currentstroke}%
\pgfsetstrokeopacity{0.000000}%
\pgfsetdash{}{0pt}%
\pgfpathmoveto{\pgfqpoint{4.020353in}{0.341071in}}%
\pgfpathlineto{\pgfqpoint{4.015449in}{0.340433in}}%
\pgfpathlineto{\pgfqpoint{4.008542in}{0.337778in}}%
\pgfpathlineto{\pgfqpoint{4.006931in}{0.339221in}}%
\pgfpathlineto{\pgfqpoint{4.014146in}{0.340851in}}%
\pgfpathlineto{\pgfqpoint{4.012212in}{0.342273in}}%
\pgfpathlineto{\pgfqpoint{4.007739in}{0.343315in}}%
\pgfpathlineto{\pgfqpoint{4.006544in}{0.346649in}}%
\pgfpathlineto{\pgfqpoint{4.008795in}{0.349521in}}%
\pgfpathlineto{\pgfqpoint{4.010690in}{0.356112in}}%
\pgfpathlineto{\pgfqpoint{4.017614in}{0.357340in}}%
\pgfpathlineto{\pgfqpoint{4.020966in}{0.354492in}}%
\pgfpathlineto{\pgfqpoint{4.024115in}{0.357248in}}%
\pgfpathlineto{\pgfqpoint{4.028064in}{0.356572in}}%
\pgfpathlineto{\pgfqpoint{4.030902in}{0.351482in}}%
\pgfpathlineto{\pgfqpoint{4.033305in}{0.356006in}}%
\pgfpathlineto{\pgfqpoint{4.040126in}{0.345423in}}%
\pgfpathlineto{\pgfqpoint{4.037662in}{0.344577in}}%
\pgfpathlineto{\pgfqpoint{4.036150in}{0.342283in}}%
\pgfpathlineto{\pgfqpoint{4.039007in}{0.340190in}}%
\pgfpathlineto{\pgfqpoint{4.034576in}{0.338227in}}%
\pgfpathlineto{\pgfqpoint{4.031364in}{0.339308in}}%
\pgfpathlineto{\pgfqpoint{4.025296in}{0.342890in}}%
\pgfpathlineto{\pgfqpoint{4.020353in}{0.341071in}}%
\pgfpathclose%
\pgfusepath{fill}%
\end{pgfscope}%
\begin{pgfscope}%
\pgfpathrectangle{\pgfqpoint{3.525000in}{0.100000in}}{\pgfqpoint{2.862985in}{1.832778in}}%
\pgfusepath{clip}%
\pgfsetbuttcap%
\pgfsetmiterjoin%
\definecolor{currentfill}{rgb}{0.645872,0.120898,0.598867}%
\pgfsetfillcolor{currentfill}%
\pgfsetlinewidth{0.000000pt}%
\definecolor{currentstroke}{rgb}{0.000000,0.000000,0.000000}%
\pgfsetstrokecolor{currentstroke}%
\pgfsetstrokeopacity{0.000000}%
\pgfsetdash{}{0pt}%
\pgfpathmoveto{\pgfqpoint{4.212691in}{0.199220in}}%
\pgfpathlineto{\pgfqpoint{4.211719in}{0.208994in}}%
\pgfpathlineto{\pgfqpoint{4.214340in}{0.212583in}}%
\pgfpathlineto{\pgfqpoint{4.211538in}{0.212832in}}%
\pgfpathlineto{\pgfqpoint{4.212562in}{0.219812in}}%
\pgfpathlineto{\pgfqpoint{4.213770in}{0.223955in}}%
\pgfpathlineto{\pgfqpoint{4.212720in}{0.229522in}}%
\pgfpathlineto{\pgfqpoint{4.215287in}{0.230600in}}%
\pgfpathlineto{\pgfqpoint{4.216124in}{0.232039in}}%
\pgfpathlineto{\pgfqpoint{4.218653in}{0.231493in}}%
\pgfpathlineto{\pgfqpoint{4.220842in}{0.227033in}}%
\pgfpathlineto{\pgfqpoint{4.221249in}{0.222947in}}%
\pgfpathlineto{\pgfqpoint{4.218571in}{0.210709in}}%
\pgfpathlineto{\pgfqpoint{4.212691in}{0.199220in}}%
\pgfpathclose%
\pgfusepath{fill}%
\end{pgfscope}%
\begin{pgfscope}%
\pgfpathrectangle{\pgfqpoint{3.525000in}{0.100000in}}{\pgfqpoint{2.862985in}{1.832778in}}%
\pgfusepath{clip}%
\pgfsetbuttcap%
\pgfsetmiterjoin%
\definecolor{currentfill}{rgb}{0.645872,0.120898,0.598867}%
\pgfsetfillcolor{currentfill}%
\pgfsetlinewidth{0.000000pt}%
\definecolor{currentstroke}{rgb}{0.000000,0.000000,0.000000}%
\pgfsetstrokecolor{currentstroke}%
\pgfsetstrokeopacity{0.000000}%
\pgfsetdash{}{0pt}%
\pgfpathmoveto{\pgfqpoint{4.212110in}{0.229056in}}%
\pgfpathlineto{\pgfqpoint{4.211975in}{0.224290in}}%
\pgfpathlineto{\pgfqpoint{4.209778in}{0.222127in}}%
\pgfpathlineto{\pgfqpoint{4.207238in}{0.222926in}}%
\pgfpathlineto{\pgfqpoint{4.210448in}{0.229848in}}%
\pgfpathlineto{\pgfqpoint{4.212110in}{0.229056in}}%
\pgfpathclose%
\pgfusepath{fill}%
\end{pgfscope}%
\begin{pgfscope}%
\pgfpathrectangle{\pgfqpoint{3.525000in}{0.100000in}}{\pgfqpoint{2.862985in}{1.832778in}}%
\pgfusepath{clip}%
\pgfsetbuttcap%
\pgfsetmiterjoin%
\definecolor{currentfill}{rgb}{0.645872,0.120898,0.598867}%
\pgfsetfillcolor{currentfill}%
\pgfsetlinewidth{0.000000pt}%
\definecolor{currentstroke}{rgb}{0.000000,0.000000,0.000000}%
\pgfsetstrokecolor{currentstroke}%
\pgfsetstrokeopacity{0.000000}%
\pgfsetdash{}{0pt}%
\pgfpathmoveto{\pgfqpoint{4.239076in}{0.207898in}}%
\pgfpathlineto{\pgfqpoint{4.237202in}{0.199923in}}%
\pgfpathlineto{\pgfqpoint{4.232835in}{0.197402in}}%
\pgfpathlineto{\pgfqpoint{4.227107in}{0.199724in}}%
\pgfpathlineto{\pgfqpoint{4.228685in}{0.202824in}}%
\pgfpathlineto{\pgfqpoint{4.228653in}{0.204970in}}%
\pgfpathlineto{\pgfqpoint{4.230783in}{0.208616in}}%
\pgfpathlineto{\pgfqpoint{4.229286in}{0.207882in}}%
\pgfpathlineto{\pgfqpoint{4.230377in}{0.209663in}}%
\pgfpathlineto{\pgfqpoint{4.228304in}{0.213709in}}%
\pgfpathlineto{\pgfqpoint{4.230657in}{0.214455in}}%
\pgfpathlineto{\pgfqpoint{4.236104in}{0.209721in}}%
\pgfpathlineto{\pgfqpoint{4.239076in}{0.207898in}}%
\pgfpathclose%
\pgfusepath{fill}%
\end{pgfscope}%
\begin{pgfscope}%
\pgfpathrectangle{\pgfqpoint{3.525000in}{0.100000in}}{\pgfqpoint{2.862985in}{1.832778in}}%
\pgfusepath{clip}%
\pgfsetbuttcap%
\pgfsetmiterjoin%
\definecolor{currentfill}{rgb}{0.645872,0.120898,0.598867}%
\pgfsetfillcolor{currentfill}%
\pgfsetlinewidth{0.000000pt}%
\definecolor{currentstroke}{rgb}{0.000000,0.000000,0.000000}%
\pgfsetstrokecolor{currentstroke}%
\pgfsetstrokeopacity{0.000000}%
\pgfsetdash{}{0pt}%
\pgfpathmoveto{\pgfqpoint{4.219962in}{0.193796in}}%
\pgfpathlineto{\pgfqpoint{4.217492in}{0.192836in}}%
\pgfpathlineto{\pgfqpoint{4.218902in}{0.201252in}}%
\pgfpathlineto{\pgfqpoint{4.221866in}{0.203216in}}%
\pgfpathlineto{\pgfqpoint{4.221144in}{0.209521in}}%
\pgfpathlineto{\pgfqpoint{4.222348in}{0.212317in}}%
\pgfpathlineto{\pgfqpoint{4.224731in}{0.213365in}}%
\pgfpathlineto{\pgfqpoint{4.227232in}{0.210622in}}%
\pgfpathlineto{\pgfqpoint{4.229379in}{0.206901in}}%
\pgfpathlineto{\pgfqpoint{4.222478in}{0.199610in}}%
\pgfpathlineto{\pgfqpoint{4.219962in}{0.193796in}}%
\pgfpathclose%
\pgfusepath{fill}%
\end{pgfscope}%
\begin{pgfscope}%
\pgfpathrectangle{\pgfqpoint{3.525000in}{0.100000in}}{\pgfqpoint{2.862985in}{1.832778in}}%
\pgfusepath{clip}%
\pgfsetbuttcap%
\pgfsetmiterjoin%
\definecolor{currentfill}{rgb}{0.645872,0.120898,0.598867}%
\pgfsetfillcolor{currentfill}%
\pgfsetlinewidth{0.000000pt}%
\definecolor{currentstroke}{rgb}{0.000000,0.000000,0.000000}%
\pgfsetstrokecolor{currentstroke}%
\pgfsetstrokeopacity{0.000000}%
\pgfsetdash{}{0pt}%
\pgfpathmoveto{\pgfqpoint{4.240172in}{0.202568in}}%
\pgfpathlineto{\pgfqpoint{4.241382in}{0.196760in}}%
\pgfpathlineto{\pgfqpoint{4.235974in}{0.197425in}}%
\pgfpathlineto{\pgfqpoint{4.240172in}{0.202568in}}%
\pgfpathclose%
\pgfusepath{fill}%
\end{pgfscope}%
\begin{pgfscope}%
\pgfpathrectangle{\pgfqpoint{3.525000in}{0.100000in}}{\pgfqpoint{2.862985in}{1.832778in}}%
\pgfusepath{clip}%
\pgfsetbuttcap%
\pgfsetmiterjoin%
\definecolor{currentfill}{rgb}{0.645872,0.120898,0.598867}%
\pgfsetfillcolor{currentfill}%
\pgfsetlinewidth{0.000000pt}%
\definecolor{currentstroke}{rgb}{0.000000,0.000000,0.000000}%
\pgfsetstrokecolor{currentstroke}%
\pgfsetstrokeopacity{0.000000}%
\pgfsetdash{}{0pt}%
\pgfpathmoveto{\pgfqpoint{4.242752in}{0.193694in}}%
\pgfpathlineto{\pgfqpoint{4.243638in}{0.190111in}}%
\pgfpathlineto{\pgfqpoint{4.245335in}{0.188965in}}%
\pgfpathlineto{\pgfqpoint{4.245026in}{0.185593in}}%
\pgfpathlineto{\pgfqpoint{4.242656in}{0.184459in}}%
\pgfpathlineto{\pgfqpoint{4.240732in}{0.189493in}}%
\pgfpathlineto{\pgfqpoint{4.242752in}{0.193694in}}%
\pgfpathclose%
\pgfusepath{fill}%
\end{pgfscope}%
\begin{pgfscope}%
\pgfpathrectangle{\pgfqpoint{3.525000in}{0.100000in}}{\pgfqpoint{2.862985in}{1.832778in}}%
\pgfusepath{clip}%
\pgfsetbuttcap%
\pgfsetmiterjoin%
\definecolor{currentfill}{rgb}{0.645872,0.120898,0.598867}%
\pgfsetfillcolor{currentfill}%
\pgfsetlinewidth{0.000000pt}%
\definecolor{currentstroke}{rgb}{0.000000,0.000000,0.000000}%
\pgfsetstrokecolor{currentstroke}%
\pgfsetstrokeopacity{0.000000}%
\pgfsetdash{}{0pt}%
\pgfpathmoveto{\pgfqpoint{4.238239in}{0.195068in}}%
\pgfpathlineto{\pgfqpoint{4.236768in}{0.190907in}}%
\pgfpathlineto{\pgfqpoint{4.234484in}{0.191060in}}%
\pgfpathlineto{\pgfqpoint{4.233043in}{0.194492in}}%
\pgfpathlineto{\pgfqpoint{4.235390in}{0.196117in}}%
\pgfpathlineto{\pgfqpoint{4.238239in}{0.195068in}}%
\pgfpathclose%
\pgfusepath{fill}%
\end{pgfscope}%
\begin{pgfscope}%
\pgfpathrectangle{\pgfqpoint{3.525000in}{0.100000in}}{\pgfqpoint{2.862985in}{1.832778in}}%
\pgfusepath{clip}%
\pgfsetbuttcap%
\pgfsetmiterjoin%
\definecolor{currentfill}{rgb}{0.645872,0.120898,0.598867}%
\pgfsetfillcolor{currentfill}%
\pgfsetlinewidth{0.000000pt}%
\definecolor{currentstroke}{rgb}{0.000000,0.000000,0.000000}%
\pgfsetstrokecolor{currentstroke}%
\pgfsetstrokeopacity{0.000000}%
\pgfsetdash{}{0pt}%
\pgfpathmoveto{\pgfqpoint{4.232076in}{0.173613in}}%
\pgfpathlineto{\pgfqpoint{4.234114in}{0.171384in}}%
\pgfpathlineto{\pgfqpoint{4.234563in}{0.166686in}}%
\pgfpathlineto{\pgfqpoint{4.235636in}{0.165321in}}%
\pgfpathlineto{\pgfqpoint{4.233884in}{0.161070in}}%
\pgfpathlineto{\pgfqpoint{4.231578in}{0.158676in}}%
\pgfpathlineto{\pgfqpoint{4.230352in}{0.154236in}}%
\pgfpathlineto{\pgfqpoint{4.227832in}{0.161156in}}%
\pgfpathlineto{\pgfqpoint{4.227414in}{0.167439in}}%
\pgfpathlineto{\pgfqpoint{4.226004in}{0.170567in}}%
\pgfpathlineto{\pgfqpoint{4.221147in}{0.173102in}}%
\pgfpathlineto{\pgfqpoint{4.225703in}{0.178390in}}%
\pgfpathlineto{\pgfqpoint{4.222666in}{0.180650in}}%
\pgfpathlineto{\pgfqpoint{4.225341in}{0.182902in}}%
\pgfpathlineto{\pgfqpoint{4.228893in}{0.191549in}}%
\pgfpathlineto{\pgfqpoint{4.225212in}{0.194597in}}%
\pgfpathlineto{\pgfqpoint{4.226630in}{0.197524in}}%
\pgfpathlineto{\pgfqpoint{4.231388in}{0.194928in}}%
\pgfpathlineto{\pgfqpoint{4.231870in}{0.190831in}}%
\pgfpathlineto{\pgfqpoint{4.230070in}{0.188498in}}%
\pgfpathlineto{\pgfqpoint{4.233638in}{0.185341in}}%
\pgfpathlineto{\pgfqpoint{4.234761in}{0.180424in}}%
\pgfpathlineto{\pgfqpoint{4.234609in}{0.173528in}}%
\pgfpathlineto{\pgfqpoint{4.232076in}{0.173613in}}%
\pgfpathclose%
\pgfusepath{fill}%
\end{pgfscope}%
\begin{pgfscope}%
\pgfpathrectangle{\pgfqpoint{3.525000in}{0.100000in}}{\pgfqpoint{2.862985in}{1.832778in}}%
\pgfusepath{clip}%
\pgfsetbuttcap%
\pgfsetmiterjoin%
\definecolor{currentfill}{rgb}{0.645872,0.120898,0.598867}%
\pgfsetfillcolor{currentfill}%
\pgfsetlinewidth{0.000000pt}%
\definecolor{currentstroke}{rgb}{0.000000,0.000000,0.000000}%
\pgfsetstrokecolor{currentstroke}%
\pgfsetstrokeopacity{0.000000}%
\pgfsetdash{}{0pt}%
\pgfpathmoveto{\pgfqpoint{4.240789in}{0.190924in}}%
\pgfpathlineto{\pgfqpoint{4.239679in}{0.188532in}}%
\pgfpathlineto{\pgfqpoint{4.240722in}{0.184648in}}%
\pgfpathlineto{\pgfqpoint{4.238053in}{0.184321in}}%
\pgfpathlineto{\pgfqpoint{4.235255in}{0.186119in}}%
\pgfpathlineto{\pgfqpoint{4.236163in}{0.190080in}}%
\pgfpathlineto{\pgfqpoint{4.240789in}{0.190924in}}%
\pgfpathclose%
\pgfusepath{fill}%
\end{pgfscope}%
\begin{pgfscope}%
\pgfpathrectangle{\pgfqpoint{3.525000in}{0.100000in}}{\pgfqpoint{2.862985in}{1.832778in}}%
\pgfusepath{clip}%
\pgfsetbuttcap%
\pgfsetmiterjoin%
\definecolor{currentfill}{rgb}{0.645872,0.120898,0.598867}%
\pgfsetfillcolor{currentfill}%
\pgfsetlinewidth{0.000000pt}%
\definecolor{currentstroke}{rgb}{0.000000,0.000000,0.000000}%
\pgfsetstrokecolor{currentstroke}%
\pgfsetstrokeopacity{0.000000}%
\pgfsetdash{}{0pt}%
\pgfpathmoveto{\pgfqpoint{4.228351in}{0.191195in}}%
\pgfpathlineto{\pgfqpoint{4.222461in}{0.188553in}}%
\pgfpathlineto{\pgfqpoint{4.220459in}{0.189503in}}%
\pgfpathlineto{\pgfqpoint{4.224476in}{0.192989in}}%
\pgfpathlineto{\pgfqpoint{4.228351in}{0.191195in}}%
\pgfpathclose%
\pgfusepath{fill}%
\end{pgfscope}%
\begin{pgfscope}%
\pgfpathrectangle{\pgfqpoint{3.525000in}{0.100000in}}{\pgfqpoint{2.862985in}{1.832778in}}%
\pgfusepath{clip}%
\pgfsetbuttcap%
\pgfsetmiterjoin%
\definecolor{currentfill}{rgb}{0.645872,0.120898,0.598867}%
\pgfsetfillcolor{currentfill}%
\pgfsetlinewidth{0.000000pt}%
\definecolor{currentstroke}{rgb}{0.000000,0.000000,0.000000}%
\pgfsetstrokecolor{currentstroke}%
\pgfsetstrokeopacity{0.000000}%
\pgfsetdash{}{0pt}%
\pgfpathmoveto{\pgfqpoint{4.240854in}{0.165422in}}%
\pgfpathlineto{\pgfqpoint{4.239449in}{0.169400in}}%
\pgfpathlineto{\pgfqpoint{4.242419in}{0.170378in}}%
\pgfpathlineto{\pgfqpoint{4.243319in}{0.174668in}}%
\pgfpathlineto{\pgfqpoint{4.246538in}{0.174673in}}%
\pgfpathlineto{\pgfqpoint{4.246524in}{0.177660in}}%
\pgfpathlineto{\pgfqpoint{4.250997in}{0.176953in}}%
\pgfpathlineto{\pgfqpoint{4.251760in}{0.168856in}}%
\pgfpathlineto{\pgfqpoint{4.249165in}{0.163694in}}%
\pgfpathlineto{\pgfqpoint{4.245315in}{0.160469in}}%
\pgfpathlineto{\pgfqpoint{4.240854in}{0.165422in}}%
\pgfpathclose%
\pgfusepath{fill}%
\end{pgfscope}%
\begin{pgfscope}%
\pgfpathrectangle{\pgfqpoint{3.525000in}{0.100000in}}{\pgfqpoint{2.862985in}{1.832778in}}%
\pgfusepath{clip}%
\pgfsetbuttcap%
\pgfsetmiterjoin%
\definecolor{currentfill}{rgb}{0.645872,0.120898,0.598867}%
\pgfsetfillcolor{currentfill}%
\pgfsetlinewidth{0.000000pt}%
\definecolor{currentstroke}{rgb}{0.000000,0.000000,0.000000}%
\pgfsetstrokecolor{currentstroke}%
\pgfsetstrokeopacity{0.000000}%
\pgfsetdash{}{0pt}%
\pgfpathmoveto{\pgfqpoint{4.239114in}{0.168482in}}%
\pgfpathlineto{\pgfqpoint{4.240383in}{0.164715in}}%
\pgfpathlineto{\pgfqpoint{4.237255in}{0.164257in}}%
\pgfpathlineto{\pgfqpoint{4.239114in}{0.168482in}}%
\pgfpathclose%
\pgfusepath{fill}%
\end{pgfscope}%
\begin{pgfscope}%
\pgfpathrectangle{\pgfqpoint{3.525000in}{0.100000in}}{\pgfqpoint{2.862985in}{1.832778in}}%
\pgfusepath{clip}%
\pgfsetbuttcap%
\pgfsetmiterjoin%
\definecolor{currentfill}{rgb}{0.645872,0.120898,0.598867}%
\pgfsetfillcolor{currentfill}%
\pgfsetlinewidth{0.000000pt}%
\definecolor{currentstroke}{rgb}{0.000000,0.000000,0.000000}%
\pgfsetstrokecolor{currentstroke}%
\pgfsetstrokeopacity{0.000000}%
\pgfsetdash{}{0pt}%
\pgfpathmoveto{\pgfqpoint{4.241751in}{0.162942in}}%
\pgfpathlineto{\pgfqpoint{4.241315in}{0.158141in}}%
\pgfpathlineto{\pgfqpoint{4.238015in}{0.158538in}}%
\pgfpathlineto{\pgfqpoint{4.239083in}{0.160858in}}%
\pgfpathlineto{\pgfqpoint{4.241751in}{0.162942in}}%
\pgfpathclose%
\pgfusepath{fill}%
\end{pgfscope}%
\begin{pgfscope}%
\definecolor{textcolor}{rgb}{0.000000,0.000000,0.000000}%
\pgfsetstrokecolor{textcolor}%
\pgfsetfillcolor{textcolor}%
\pgftext[x=4.727454in,y=1.896122in,left,base]{\color{textcolor}\setmainfont{Lato}\rmfamily\fontsize{9.000000}{10.800000}\selectfont November 2022}%
\end{pgfscope}%
\begin{pgfscope}%
\pgfpathrectangle{\pgfqpoint{2.893540in}{0.169444in}}{\pgfqpoint{0.858895in}{1.282944in}}%
\pgfusepath{clip}%
\pgfsetbuttcap%
\pgfsetmiterjoin%
\definecolor{currentfill}{rgb}{0.991209,0.790537,0.149377}%
\pgfsetfillcolor{currentfill}%
\pgfsetlinewidth{0.000000pt}%
\definecolor{currentstroke}{rgb}{0.000000,0.000000,0.000000}%
\pgfsetstrokecolor{currentstroke}%
\pgfsetstrokeopacity{0.000000}%
\pgfsetdash{}{0pt}%
\pgfpathmoveto{\pgfqpoint{3.142620in}{0.314167in}}%
\pgfpathlineto{\pgfqpoint{3.108264in}{0.314167in}}%
\pgfpathlineto{\pgfqpoint{3.108264in}{0.350822in}}%
\pgfpathlineto{\pgfqpoint{3.142620in}{0.350822in}}%
\pgfpathlineto{\pgfqpoint{3.142620in}{0.314167in}}%
\pgfpathclose%
\pgfusepath{fill}%
\end{pgfscope}%
\begin{pgfscope}%
\pgfpathrectangle{\pgfqpoint{2.893540in}{0.169444in}}{\pgfqpoint{0.858895in}{1.282944in}}%
\pgfusepath{clip}%
\pgfsetbuttcap%
\pgfsetmiterjoin%
\definecolor{currentfill}{rgb}{0.993851,0.759304,0.159092}%
\pgfsetfillcolor{currentfill}%
\pgfsetlinewidth{0.000000pt}%
\definecolor{currentstroke}{rgb}{0.000000,0.000000,0.000000}%
\pgfsetstrokecolor{currentstroke}%
\pgfsetstrokeopacity{0.000000}%
\pgfsetdash{}{0pt}%
\pgfpathmoveto{\pgfqpoint{3.142620in}{0.339240in}}%
\pgfpathlineto{\pgfqpoint{3.108264in}{0.339240in}}%
\pgfpathlineto{\pgfqpoint{3.108264in}{0.375896in}}%
\pgfpathlineto{\pgfqpoint{3.142620in}{0.375896in}}%
\pgfpathlineto{\pgfqpoint{3.142620in}{0.339240in}}%
\pgfpathclose%
\pgfusepath{fill}%
\end{pgfscope}%
\begin{pgfscope}%
\pgfpathrectangle{\pgfqpoint{2.893540in}{0.169444in}}{\pgfqpoint{0.858895in}{1.282944in}}%
\pgfusepath{clip}%
\pgfsetbuttcap%
\pgfsetmiterjoin%
\definecolor{currentfill}{rgb}{0.994553,0.728728,0.171622}%
\pgfsetfillcolor{currentfill}%
\pgfsetlinewidth{0.000000pt}%
\definecolor{currentstroke}{rgb}{0.000000,0.000000,0.000000}%
\pgfsetstrokecolor{currentstroke}%
\pgfsetstrokeopacity{0.000000}%
\pgfsetdash{}{0pt}%
\pgfpathmoveto{\pgfqpoint{3.142620in}{0.364313in}}%
\pgfpathlineto{\pgfqpoint{3.108264in}{0.364313in}}%
\pgfpathlineto{\pgfqpoint{3.108264in}{0.400969in}}%
\pgfpathlineto{\pgfqpoint{3.142620in}{0.400969in}}%
\pgfpathlineto{\pgfqpoint{3.142620in}{0.364313in}}%
\pgfpathclose%
\pgfusepath{fill}%
\end{pgfscope}%
\begin{pgfscope}%
\pgfpathrectangle{\pgfqpoint{2.893540in}{0.169444in}}{\pgfqpoint{0.858895in}{1.282944in}}%
\pgfusepath{clip}%
\pgfsetbuttcap%
\pgfsetmiterjoin%
\definecolor{currentfill}{rgb}{0.993456,0.698810,0.186041}%
\pgfsetfillcolor{currentfill}%
\pgfsetlinewidth{0.000000pt}%
\definecolor{currentstroke}{rgb}{0.000000,0.000000,0.000000}%
\pgfsetstrokecolor{currentstroke}%
\pgfsetstrokeopacity{0.000000}%
\pgfsetdash{}{0pt}%
\pgfpathmoveto{\pgfqpoint{3.142620in}{0.389386in}}%
\pgfpathlineto{\pgfqpoint{3.108264in}{0.389386in}}%
\pgfpathlineto{\pgfqpoint{3.108264in}{0.426042in}}%
\pgfpathlineto{\pgfqpoint{3.142620in}{0.426042in}}%
\pgfpathlineto{\pgfqpoint{3.142620in}{0.389386in}}%
\pgfpathclose%
\pgfusepath{fill}%
\end{pgfscope}%
\begin{pgfscope}%
\pgfpathrectangle{\pgfqpoint{2.893540in}{0.169444in}}{\pgfqpoint{0.858895in}{1.282944in}}%
\pgfusepath{clip}%
\pgfsetbuttcap%
\pgfsetmiterjoin%
\definecolor{currentfill}{rgb}{0.990681,0.669558,0.201642}%
\pgfsetfillcolor{currentfill}%
\pgfsetlinewidth{0.000000pt}%
\definecolor{currentstroke}{rgb}{0.000000,0.000000,0.000000}%
\pgfsetstrokecolor{currentstroke}%
\pgfsetstrokeopacity{0.000000}%
\pgfsetdash{}{0pt}%
\pgfpathmoveto{\pgfqpoint{3.142620in}{0.414460in}}%
\pgfpathlineto{\pgfqpoint{3.108264in}{0.414460in}}%
\pgfpathlineto{\pgfqpoint{3.108264in}{0.451115in}}%
\pgfpathlineto{\pgfqpoint{3.142620in}{0.451115in}}%
\pgfpathlineto{\pgfqpoint{3.142620in}{0.414460in}}%
\pgfpathclose%
\pgfusepath{fill}%
\end{pgfscope}%
\begin{pgfscope}%
\pgfpathrectangle{\pgfqpoint{2.893540in}{0.169444in}}{\pgfqpoint{0.858895in}{1.282944in}}%
\pgfusepath{clip}%
\pgfsetbuttcap%
\pgfsetmiterjoin%
\definecolor{currentfill}{rgb}{0.986345,0.640969,0.217948}%
\pgfsetfillcolor{currentfill}%
\pgfsetlinewidth{0.000000pt}%
\definecolor{currentstroke}{rgb}{0.000000,0.000000,0.000000}%
\pgfsetstrokecolor{currentstroke}%
\pgfsetstrokeopacity{0.000000}%
\pgfsetdash{}{0pt}%
\pgfpathmoveto{\pgfqpoint{3.142620in}{0.439533in}}%
\pgfpathlineto{\pgfqpoint{3.108264in}{0.439533in}}%
\pgfpathlineto{\pgfqpoint{3.108264in}{0.476188in}}%
\pgfpathlineto{\pgfqpoint{3.142620in}{0.476188in}}%
\pgfpathlineto{\pgfqpoint{3.142620in}{0.439533in}}%
\pgfpathclose%
\pgfusepath{fill}%
\end{pgfscope}%
\begin{pgfscope}%
\pgfpathrectangle{\pgfqpoint{2.893540in}{0.169444in}}{\pgfqpoint{0.858895in}{1.282944in}}%
\pgfusepath{clip}%
\pgfsetbuttcap%
\pgfsetmiterjoin%
\definecolor{currentfill}{rgb}{0.980556,0.613039,0.234646}%
\pgfsetfillcolor{currentfill}%
\pgfsetlinewidth{0.000000pt}%
\definecolor{currentstroke}{rgb}{0.000000,0.000000,0.000000}%
\pgfsetstrokecolor{currentstroke}%
\pgfsetstrokeopacity{0.000000}%
\pgfsetdash{}{0pt}%
\pgfpathmoveto{\pgfqpoint{3.142620in}{0.464606in}}%
\pgfpathlineto{\pgfqpoint{3.108264in}{0.464606in}}%
\pgfpathlineto{\pgfqpoint{3.108264in}{0.501261in}}%
\pgfpathlineto{\pgfqpoint{3.142620in}{0.501261in}}%
\pgfpathlineto{\pgfqpoint{3.142620in}{0.464606in}}%
\pgfpathclose%
\pgfusepath{fill}%
\end{pgfscope}%
\begin{pgfscope}%
\pgfpathrectangle{\pgfqpoint{2.893540in}{0.169444in}}{\pgfqpoint{0.858895in}{1.282944in}}%
\pgfusepath{clip}%
\pgfsetbuttcap%
\pgfsetmiterjoin%
\definecolor{currentfill}{rgb}{0.973416,0.585761,0.251540}%
\pgfsetfillcolor{currentfill}%
\pgfsetlinewidth{0.000000pt}%
\definecolor{currentstroke}{rgb}{0.000000,0.000000,0.000000}%
\pgfsetstrokecolor{currentstroke}%
\pgfsetstrokeopacity{0.000000}%
\pgfsetdash{}{0pt}%
\pgfpathmoveto{\pgfqpoint{3.142620in}{0.489679in}}%
\pgfpathlineto{\pgfqpoint{3.108264in}{0.489679in}}%
\pgfpathlineto{\pgfqpoint{3.108264in}{0.526335in}}%
\pgfpathlineto{\pgfqpoint{3.142620in}{0.526335in}}%
\pgfpathlineto{\pgfqpoint{3.142620in}{0.489679in}}%
\pgfpathclose%
\pgfusepath{fill}%
\end{pgfscope}%
\begin{pgfscope}%
\pgfpathrectangle{\pgfqpoint{2.893540in}{0.169444in}}{\pgfqpoint{0.858895in}{1.282944in}}%
\pgfusepath{clip}%
\pgfsetbuttcap%
\pgfsetmiterjoin%
\definecolor{currentfill}{rgb}{0.965024,0.559118,0.268513}%
\pgfsetfillcolor{currentfill}%
\pgfsetlinewidth{0.000000pt}%
\definecolor{currentstroke}{rgb}{0.000000,0.000000,0.000000}%
\pgfsetstrokecolor{currentstroke}%
\pgfsetstrokeopacity{0.000000}%
\pgfsetdash{}{0pt}%
\pgfpathmoveto{\pgfqpoint{3.142620in}{0.514752in}}%
\pgfpathlineto{\pgfqpoint{3.108264in}{0.514752in}}%
\pgfpathlineto{\pgfqpoint{3.108264in}{0.551408in}}%
\pgfpathlineto{\pgfqpoint{3.142620in}{0.551408in}}%
\pgfpathlineto{\pgfqpoint{3.142620in}{0.514752in}}%
\pgfpathclose%
\pgfusepath{fill}%
\end{pgfscope}%
\begin{pgfscope}%
\pgfpathrectangle{\pgfqpoint{2.893540in}{0.169444in}}{\pgfqpoint{0.858895in}{1.282944in}}%
\pgfusepath{clip}%
\pgfsetbuttcap%
\pgfsetmiterjoin%
\definecolor{currentfill}{rgb}{0.955470,0.533093,0.285490}%
\pgfsetfillcolor{currentfill}%
\pgfsetlinewidth{0.000000pt}%
\definecolor{currentstroke}{rgb}{0.000000,0.000000,0.000000}%
\pgfsetstrokecolor{currentstroke}%
\pgfsetstrokeopacity{0.000000}%
\pgfsetdash{}{0pt}%
\pgfpathmoveto{\pgfqpoint{3.142620in}{0.539825in}}%
\pgfpathlineto{\pgfqpoint{3.108264in}{0.539825in}}%
\pgfpathlineto{\pgfqpoint{3.108264in}{0.576481in}}%
\pgfpathlineto{\pgfqpoint{3.142620in}{0.576481in}}%
\pgfpathlineto{\pgfqpoint{3.142620in}{0.539825in}}%
\pgfpathclose%
\pgfusepath{fill}%
\end{pgfscope}%
\begin{pgfscope}%
\pgfpathrectangle{\pgfqpoint{2.893540in}{0.169444in}}{\pgfqpoint{0.858895in}{1.282944in}}%
\pgfusepath{clip}%
\pgfsetbuttcap%
\pgfsetmiterjoin%
\definecolor{currentfill}{rgb}{0.944844,0.507658,0.302433}%
\pgfsetfillcolor{currentfill}%
\pgfsetlinewidth{0.000000pt}%
\definecolor{currentstroke}{rgb}{0.000000,0.000000,0.000000}%
\pgfsetstrokecolor{currentstroke}%
\pgfsetstrokeopacity{0.000000}%
\pgfsetdash{}{0pt}%
\pgfpathmoveto{\pgfqpoint{3.142620in}{0.564899in}}%
\pgfpathlineto{\pgfqpoint{3.108264in}{0.564899in}}%
\pgfpathlineto{\pgfqpoint{3.108264in}{0.601554in}}%
\pgfpathlineto{\pgfqpoint{3.142620in}{0.601554in}}%
\pgfpathlineto{\pgfqpoint{3.142620in}{0.564899in}}%
\pgfpathclose%
\pgfusepath{fill}%
\end{pgfscope}%
\begin{pgfscope}%
\pgfpathrectangle{\pgfqpoint{2.893540in}{0.169444in}}{\pgfqpoint{0.858895in}{1.282944in}}%
\pgfusepath{clip}%
\pgfsetbuttcap%
\pgfsetmiterjoin%
\definecolor{currentfill}{rgb}{0.933232,0.482780,0.319325}%
\pgfsetfillcolor{currentfill}%
\pgfsetlinewidth{0.000000pt}%
\definecolor{currentstroke}{rgb}{0.000000,0.000000,0.000000}%
\pgfsetstrokecolor{currentstroke}%
\pgfsetstrokeopacity{0.000000}%
\pgfsetdash{}{0pt}%
\pgfpathmoveto{\pgfqpoint{3.142620in}{0.589972in}}%
\pgfpathlineto{\pgfqpoint{3.108264in}{0.589972in}}%
\pgfpathlineto{\pgfqpoint{3.108264in}{0.626627in}}%
\pgfpathlineto{\pgfqpoint{3.142620in}{0.626627in}}%
\pgfpathlineto{\pgfqpoint{3.142620in}{0.589972in}}%
\pgfpathclose%
\pgfusepath{fill}%
\end{pgfscope}%
\begin{pgfscope}%
\pgfpathrectangle{\pgfqpoint{2.893540in}{0.169444in}}{\pgfqpoint{0.858895in}{1.282944in}}%
\pgfusepath{clip}%
\pgfsetbuttcap%
\pgfsetmiterjoin%
\definecolor{currentfill}{rgb}{0.920714,0.458417,0.336166}%
\pgfsetfillcolor{currentfill}%
\pgfsetlinewidth{0.000000pt}%
\definecolor{currentstroke}{rgb}{0.000000,0.000000,0.000000}%
\pgfsetstrokecolor{currentstroke}%
\pgfsetstrokeopacity{0.000000}%
\pgfsetdash{}{0pt}%
\pgfpathmoveto{\pgfqpoint{3.142620in}{0.615045in}}%
\pgfpathlineto{\pgfqpoint{3.108264in}{0.615045in}}%
\pgfpathlineto{\pgfqpoint{3.108264in}{0.651701in}}%
\pgfpathlineto{\pgfqpoint{3.142620in}{0.651701in}}%
\pgfpathlineto{\pgfqpoint{3.142620in}{0.615045in}}%
\pgfpathclose%
\pgfusepath{fill}%
\end{pgfscope}%
\begin{pgfscope}%
\pgfpathrectangle{\pgfqpoint{2.893540in}{0.169444in}}{\pgfqpoint{0.858895in}{1.282944in}}%
\pgfusepath{clip}%
\pgfsetbuttcap%
\pgfsetmiterjoin%
\definecolor{currentfill}{rgb}{0.907365,0.434524,0.352970}%
\pgfsetfillcolor{currentfill}%
\pgfsetlinewidth{0.000000pt}%
\definecolor{currentstroke}{rgb}{0.000000,0.000000,0.000000}%
\pgfsetstrokecolor{currentstroke}%
\pgfsetstrokeopacity{0.000000}%
\pgfsetdash{}{0pt}%
\pgfpathmoveto{\pgfqpoint{3.142620in}{0.640118in}}%
\pgfpathlineto{\pgfqpoint{3.108264in}{0.640118in}}%
\pgfpathlineto{\pgfqpoint{3.108264in}{0.676774in}}%
\pgfpathlineto{\pgfqpoint{3.142620in}{0.676774in}}%
\pgfpathlineto{\pgfqpoint{3.142620in}{0.640118in}}%
\pgfpathclose%
\pgfusepath{fill}%
\end{pgfscope}%
\begin{pgfscope}%
\pgfpathrectangle{\pgfqpoint{2.893540in}{0.169444in}}{\pgfqpoint{0.858895in}{1.282944in}}%
\pgfusepath{clip}%
\pgfsetbuttcap%
\pgfsetmiterjoin%
\definecolor{currentfill}{rgb}{0.893250,0.411048,0.369768}%
\pgfsetfillcolor{currentfill}%
\pgfsetlinewidth{0.000000pt}%
\definecolor{currentstroke}{rgb}{0.000000,0.000000,0.000000}%
\pgfsetstrokecolor{currentstroke}%
\pgfsetstrokeopacity{0.000000}%
\pgfsetdash{}{0pt}%
\pgfpathmoveto{\pgfqpoint{3.142620in}{0.665191in}}%
\pgfpathlineto{\pgfqpoint{3.108264in}{0.665191in}}%
\pgfpathlineto{\pgfqpoint{3.108264in}{0.701847in}}%
\pgfpathlineto{\pgfqpoint{3.142620in}{0.701847in}}%
\pgfpathlineto{\pgfqpoint{3.142620in}{0.665191in}}%
\pgfpathclose%
\pgfusepath{fill}%
\end{pgfscope}%
\begin{pgfscope}%
\pgfpathrectangle{\pgfqpoint{2.893540in}{0.169444in}}{\pgfqpoint{0.858895in}{1.282944in}}%
\pgfusepath{clip}%
\pgfsetbuttcap%
\pgfsetmiterjoin%
\definecolor{currentfill}{rgb}{0.878423,0.387932,0.386600}%
\pgfsetfillcolor{currentfill}%
\pgfsetlinewidth{0.000000pt}%
\definecolor{currentstroke}{rgb}{0.000000,0.000000,0.000000}%
\pgfsetstrokecolor{currentstroke}%
\pgfsetstrokeopacity{0.000000}%
\pgfsetdash{}{0pt}%
\pgfpathmoveto{\pgfqpoint{3.142620in}{0.690265in}}%
\pgfpathlineto{\pgfqpoint{3.108264in}{0.690265in}}%
\pgfpathlineto{\pgfqpoint{3.108264in}{0.726920in}}%
\pgfpathlineto{\pgfqpoint{3.142620in}{0.726920in}}%
\pgfpathlineto{\pgfqpoint{3.142620in}{0.690265in}}%
\pgfpathclose%
\pgfusepath{fill}%
\end{pgfscope}%
\begin{pgfscope}%
\pgfpathrectangle{\pgfqpoint{2.893540in}{0.169444in}}{\pgfqpoint{0.858895in}{1.282944in}}%
\pgfusepath{clip}%
\pgfsetbuttcap%
\pgfsetmiterjoin%
\definecolor{currentfill}{rgb}{0.862927,0.365119,0.403519}%
\pgfsetfillcolor{currentfill}%
\pgfsetlinewidth{0.000000pt}%
\definecolor{currentstroke}{rgb}{0.000000,0.000000,0.000000}%
\pgfsetstrokecolor{currentstroke}%
\pgfsetstrokeopacity{0.000000}%
\pgfsetdash{}{0pt}%
\pgfpathmoveto{\pgfqpoint{3.142620in}{0.715338in}}%
\pgfpathlineto{\pgfqpoint{3.108264in}{0.715338in}}%
\pgfpathlineto{\pgfqpoint{3.108264in}{0.751993in}}%
\pgfpathlineto{\pgfqpoint{3.142620in}{0.751993in}}%
\pgfpathlineto{\pgfqpoint{3.142620in}{0.715338in}}%
\pgfpathclose%
\pgfusepath{fill}%
\end{pgfscope}%
\begin{pgfscope}%
\pgfpathrectangle{\pgfqpoint{2.893540in}{0.169444in}}{\pgfqpoint{0.858895in}{1.282944in}}%
\pgfusepath{clip}%
\pgfsetbuttcap%
\pgfsetmiterjoin%
\definecolor{currentfill}{rgb}{0.846788,0.342551,0.420579}%
\pgfsetfillcolor{currentfill}%
\pgfsetlinewidth{0.000000pt}%
\definecolor{currentstroke}{rgb}{0.000000,0.000000,0.000000}%
\pgfsetstrokecolor{currentstroke}%
\pgfsetstrokeopacity{0.000000}%
\pgfsetdash{}{0pt}%
\pgfpathmoveto{\pgfqpoint{3.142620in}{0.740411in}}%
\pgfpathlineto{\pgfqpoint{3.108264in}{0.740411in}}%
\pgfpathlineto{\pgfqpoint{3.108264in}{0.777066in}}%
\pgfpathlineto{\pgfqpoint{3.142620in}{0.777066in}}%
\pgfpathlineto{\pgfqpoint{3.142620in}{0.740411in}}%
\pgfpathclose%
\pgfusepath{fill}%
\end{pgfscope}%
\begin{pgfscope}%
\pgfpathrectangle{\pgfqpoint{2.893540in}{0.169444in}}{\pgfqpoint{0.858895in}{1.282944in}}%
\pgfusepath{clip}%
\pgfsetbuttcap%
\pgfsetmiterjoin%
\definecolor{currentfill}{rgb}{0.830018,0.320172,0.437836}%
\pgfsetfillcolor{currentfill}%
\pgfsetlinewidth{0.000000pt}%
\definecolor{currentstroke}{rgb}{0.000000,0.000000,0.000000}%
\pgfsetstrokecolor{currentstroke}%
\pgfsetstrokeopacity{0.000000}%
\pgfsetdash{}{0pt}%
\pgfpathmoveto{\pgfqpoint{3.142620in}{0.765484in}}%
\pgfpathlineto{\pgfqpoint{3.108264in}{0.765484in}}%
\pgfpathlineto{\pgfqpoint{3.108264in}{0.802140in}}%
\pgfpathlineto{\pgfqpoint{3.142620in}{0.802140in}}%
\pgfpathlineto{\pgfqpoint{3.142620in}{0.765484in}}%
\pgfpathclose%
\pgfusepath{fill}%
\end{pgfscope}%
\begin{pgfscope}%
\pgfpathrectangle{\pgfqpoint{2.893540in}{0.169444in}}{\pgfqpoint{0.858895in}{1.282944in}}%
\pgfusepath{clip}%
\pgfsetbuttcap%
\pgfsetmiterjoin%
\definecolor{currentfill}{rgb}{0.812612,0.297928,0.455338}%
\pgfsetfillcolor{currentfill}%
\pgfsetlinewidth{0.000000pt}%
\definecolor{currentstroke}{rgb}{0.000000,0.000000,0.000000}%
\pgfsetstrokecolor{currentstroke}%
\pgfsetstrokeopacity{0.000000}%
\pgfsetdash{}{0pt}%
\pgfpathmoveto{\pgfqpoint{3.142620in}{0.790557in}}%
\pgfpathlineto{\pgfqpoint{3.108264in}{0.790557in}}%
\pgfpathlineto{\pgfqpoint{3.108264in}{0.827213in}}%
\pgfpathlineto{\pgfqpoint{3.142620in}{0.827213in}}%
\pgfpathlineto{\pgfqpoint{3.142620in}{0.790557in}}%
\pgfpathclose%
\pgfusepath{fill}%
\end{pgfscope}%
\begin{pgfscope}%
\pgfpathrectangle{\pgfqpoint{2.893540in}{0.169444in}}{\pgfqpoint{0.858895in}{1.282944in}}%
\pgfusepath{clip}%
\pgfsetbuttcap%
\pgfsetmiterjoin%
\definecolor{currentfill}{rgb}{0.794549,0.275770,0.473117}%
\pgfsetfillcolor{currentfill}%
\pgfsetlinewidth{0.000000pt}%
\definecolor{currentstroke}{rgb}{0.000000,0.000000,0.000000}%
\pgfsetstrokecolor{currentstroke}%
\pgfsetstrokeopacity{0.000000}%
\pgfsetdash{}{0pt}%
\pgfpathmoveto{\pgfqpoint{3.142620in}{0.815630in}}%
\pgfpathlineto{\pgfqpoint{3.108264in}{0.815630in}}%
\pgfpathlineto{\pgfqpoint{3.108264in}{0.852286in}}%
\pgfpathlineto{\pgfqpoint{3.142620in}{0.852286in}}%
\pgfpathlineto{\pgfqpoint{3.142620in}{0.815630in}}%
\pgfpathclose%
\pgfusepath{fill}%
\end{pgfscope}%
\begin{pgfscope}%
\pgfpathrectangle{\pgfqpoint{2.893540in}{0.169444in}}{\pgfqpoint{0.858895in}{1.282944in}}%
\pgfusepath{clip}%
\pgfsetbuttcap%
\pgfsetmiterjoin%
\definecolor{currentfill}{rgb}{0.775796,0.253658,0.491171}%
\pgfsetfillcolor{currentfill}%
\pgfsetlinewidth{0.000000pt}%
\definecolor{currentstroke}{rgb}{0.000000,0.000000,0.000000}%
\pgfsetstrokecolor{currentstroke}%
\pgfsetstrokeopacity{0.000000}%
\pgfsetdash{}{0pt}%
\pgfpathmoveto{\pgfqpoint{3.142620in}{0.840704in}}%
\pgfpathlineto{\pgfqpoint{3.108264in}{0.840704in}}%
\pgfpathlineto{\pgfqpoint{3.108264in}{0.877359in}}%
\pgfpathlineto{\pgfqpoint{3.142620in}{0.877359in}}%
\pgfpathlineto{\pgfqpoint{3.142620in}{0.840704in}}%
\pgfpathclose%
\pgfusepath{fill}%
\end{pgfscope}%
\begin{pgfscope}%
\pgfpathrectangle{\pgfqpoint{2.893540in}{0.169444in}}{\pgfqpoint{0.858895in}{1.282944in}}%
\pgfusepath{clip}%
\pgfsetbuttcap%
\pgfsetmiterjoin%
\definecolor{currentfill}{rgb}{0.756304,0.231555,0.509468}%
\pgfsetfillcolor{currentfill}%
\pgfsetlinewidth{0.000000pt}%
\definecolor{currentstroke}{rgb}{0.000000,0.000000,0.000000}%
\pgfsetstrokecolor{currentstroke}%
\pgfsetstrokeopacity{0.000000}%
\pgfsetdash{}{0pt}%
\pgfpathmoveto{\pgfqpoint{3.142620in}{0.865777in}}%
\pgfpathlineto{\pgfqpoint{3.108264in}{0.865777in}}%
\pgfpathlineto{\pgfqpoint{3.108264in}{0.902432in}}%
\pgfpathlineto{\pgfqpoint{3.142620in}{0.902432in}}%
\pgfpathlineto{\pgfqpoint{3.142620in}{0.865777in}}%
\pgfpathclose%
\pgfusepath{fill}%
\end{pgfscope}%
\begin{pgfscope}%
\pgfpathrectangle{\pgfqpoint{2.893540in}{0.169444in}}{\pgfqpoint{0.858895in}{1.282944in}}%
\pgfusepath{clip}%
\pgfsetbuttcap%
\pgfsetmiterjoin%
\definecolor{currentfill}{rgb}{0.736019,0.209439,0.527908}%
\pgfsetfillcolor{currentfill}%
\pgfsetlinewidth{0.000000pt}%
\definecolor{currentstroke}{rgb}{0.000000,0.000000,0.000000}%
\pgfsetstrokecolor{currentstroke}%
\pgfsetstrokeopacity{0.000000}%
\pgfsetdash{}{0pt}%
\pgfpathmoveto{\pgfqpoint{3.142620in}{0.890850in}}%
\pgfpathlineto{\pgfqpoint{3.108264in}{0.890850in}}%
\pgfpathlineto{\pgfqpoint{3.108264in}{0.927506in}}%
\pgfpathlineto{\pgfqpoint{3.142620in}{0.927506in}}%
\pgfpathlineto{\pgfqpoint{3.142620in}{0.890850in}}%
\pgfpathclose%
\pgfusepath{fill}%
\end{pgfscope}%
\begin{pgfscope}%
\pgfpathrectangle{\pgfqpoint{2.893540in}{0.169444in}}{\pgfqpoint{0.858895in}{1.282944in}}%
\pgfusepath{clip}%
\pgfsetbuttcap%
\pgfsetmiterjoin%
\definecolor{currentfill}{rgb}{0.714883,0.187299,0.546338}%
\pgfsetfillcolor{currentfill}%
\pgfsetlinewidth{0.000000pt}%
\definecolor{currentstroke}{rgb}{0.000000,0.000000,0.000000}%
\pgfsetstrokecolor{currentstroke}%
\pgfsetstrokeopacity{0.000000}%
\pgfsetdash{}{0pt}%
\pgfpathmoveto{\pgfqpoint{3.142620in}{0.915923in}}%
\pgfpathlineto{\pgfqpoint{3.108264in}{0.915923in}}%
\pgfpathlineto{\pgfqpoint{3.108264in}{0.952579in}}%
\pgfpathlineto{\pgfqpoint{3.142620in}{0.952579in}}%
\pgfpathlineto{\pgfqpoint{3.142620in}{0.915923in}}%
\pgfpathclose%
\pgfusepath{fill}%
\end{pgfscope}%
\begin{pgfscope}%
\pgfpathrectangle{\pgfqpoint{2.893540in}{0.169444in}}{\pgfqpoint{0.858895in}{1.282944in}}%
\pgfusepath{clip}%
\pgfsetbuttcap%
\pgfsetmiterjoin%
\definecolor{currentfill}{rgb}{0.692840,0.165141,0.564522}%
\pgfsetfillcolor{currentfill}%
\pgfsetlinewidth{0.000000pt}%
\definecolor{currentstroke}{rgb}{0.000000,0.000000,0.000000}%
\pgfsetstrokecolor{currentstroke}%
\pgfsetstrokeopacity{0.000000}%
\pgfsetdash{}{0pt}%
\pgfpathmoveto{\pgfqpoint{3.142620in}{0.940996in}}%
\pgfpathlineto{\pgfqpoint{3.108264in}{0.940996in}}%
\pgfpathlineto{\pgfqpoint{3.108264in}{0.977652in}}%
\pgfpathlineto{\pgfqpoint{3.142620in}{0.977652in}}%
\pgfpathlineto{\pgfqpoint{3.142620in}{0.940996in}}%
\pgfpathclose%
\pgfusepath{fill}%
\end{pgfscope}%
\begin{pgfscope}%
\pgfpathrectangle{\pgfqpoint{2.893540in}{0.169444in}}{\pgfqpoint{0.858895in}{1.282944in}}%
\pgfusepath{clip}%
\pgfsetbuttcap%
\pgfsetmiterjoin%
\definecolor{currentfill}{rgb}{0.669845,0.142992,0.582154}%
\pgfsetfillcolor{currentfill}%
\pgfsetlinewidth{0.000000pt}%
\definecolor{currentstroke}{rgb}{0.000000,0.000000,0.000000}%
\pgfsetstrokecolor{currentstroke}%
\pgfsetstrokeopacity{0.000000}%
\pgfsetdash{}{0pt}%
\pgfpathmoveto{\pgfqpoint{3.142620in}{0.966070in}}%
\pgfpathlineto{\pgfqpoint{3.108264in}{0.966070in}}%
\pgfpathlineto{\pgfqpoint{3.108264in}{1.002725in}}%
\pgfpathlineto{\pgfqpoint{3.142620in}{1.002725in}}%
\pgfpathlineto{\pgfqpoint{3.142620in}{0.966070in}}%
\pgfpathclose%
\pgfusepath{fill}%
\end{pgfscope}%
\begin{pgfscope}%
\pgfpathrectangle{\pgfqpoint{2.893540in}{0.169444in}}{\pgfqpoint{0.858895in}{1.282944in}}%
\pgfusepath{clip}%
\pgfsetbuttcap%
\pgfsetmiterjoin%
\definecolor{currentfill}{rgb}{0.645872,0.120898,0.598867}%
\pgfsetfillcolor{currentfill}%
\pgfsetlinewidth{0.000000pt}%
\definecolor{currentstroke}{rgb}{0.000000,0.000000,0.000000}%
\pgfsetstrokecolor{currentstroke}%
\pgfsetstrokeopacity{0.000000}%
\pgfsetdash{}{0pt}%
\pgfpathmoveto{\pgfqpoint{3.142620in}{0.991143in}}%
\pgfpathlineto{\pgfqpoint{3.108264in}{0.991143in}}%
\pgfpathlineto{\pgfqpoint{3.108264in}{1.027798in}}%
\pgfpathlineto{\pgfqpoint{3.142620in}{1.027798in}}%
\pgfpathlineto{\pgfqpoint{3.142620in}{0.991143in}}%
\pgfpathclose%
\pgfusepath{fill}%
\end{pgfscope}%
\begin{pgfscope}%
\pgfpathrectangle{\pgfqpoint{2.893540in}{0.169444in}}{\pgfqpoint{0.858895in}{1.282944in}}%
\pgfusepath{clip}%
\pgfsetbuttcap%
\pgfsetmiterjoin%
\definecolor{currentfill}{rgb}{0.620919,0.098934,0.614257}%
\pgfsetfillcolor{currentfill}%
\pgfsetlinewidth{0.000000pt}%
\definecolor{currentstroke}{rgb}{0.000000,0.000000,0.000000}%
\pgfsetstrokecolor{currentstroke}%
\pgfsetstrokeopacity{0.000000}%
\pgfsetdash{}{0pt}%
\pgfpathmoveto{\pgfqpoint{3.142620in}{1.016216in}}%
\pgfpathlineto{\pgfqpoint{3.108264in}{1.016216in}}%
\pgfpathlineto{\pgfqpoint{3.108264in}{1.052871in}}%
\pgfpathlineto{\pgfqpoint{3.142620in}{1.052871in}}%
\pgfpathlineto{\pgfqpoint{3.142620in}{1.016216in}}%
\pgfpathclose%
\pgfusepath{fill}%
\end{pgfscope}%
\begin{pgfscope}%
\pgfpathrectangle{\pgfqpoint{2.893540in}{0.169444in}}{\pgfqpoint{0.858895in}{1.282944in}}%
\pgfusepath{clip}%
\pgfsetbuttcap%
\pgfsetmiterjoin%
\definecolor{currentfill}{rgb}{0.595011,0.077190,0.627917}%
\pgfsetfillcolor{currentfill}%
\pgfsetlinewidth{0.000000pt}%
\definecolor{currentstroke}{rgb}{0.000000,0.000000,0.000000}%
\pgfsetstrokecolor{currentstroke}%
\pgfsetstrokeopacity{0.000000}%
\pgfsetdash{}{0pt}%
\pgfpathmoveto{\pgfqpoint{3.142620in}{1.041289in}}%
\pgfpathlineto{\pgfqpoint{3.108264in}{1.041289in}}%
\pgfpathlineto{\pgfqpoint{3.108264in}{1.077945in}}%
\pgfpathlineto{\pgfqpoint{3.142620in}{1.077945in}}%
\pgfpathlineto{\pgfqpoint{3.142620in}{1.041289in}}%
\pgfpathclose%
\pgfusepath{fill}%
\end{pgfscope}%
\begin{pgfscope}%
\pgfpathrectangle{\pgfqpoint{2.893540in}{0.169444in}}{\pgfqpoint{0.858895in}{1.282944in}}%
\pgfusepath{clip}%
\pgfsetbuttcap%
\pgfsetmiterjoin%
\definecolor{currentfill}{rgb}{0.568201,0.055778,0.639477}%
\pgfsetfillcolor{currentfill}%
\pgfsetlinewidth{0.000000pt}%
\definecolor{currentstroke}{rgb}{0.000000,0.000000,0.000000}%
\pgfsetstrokecolor{currentstroke}%
\pgfsetstrokeopacity{0.000000}%
\pgfsetdash{}{0pt}%
\pgfpathmoveto{\pgfqpoint{3.142620in}{1.066362in}}%
\pgfpathlineto{\pgfqpoint{3.108264in}{1.066362in}}%
\pgfpathlineto{\pgfqpoint{3.108264in}{1.103018in}}%
\pgfpathlineto{\pgfqpoint{3.142620in}{1.103018in}}%
\pgfpathlineto{\pgfqpoint{3.142620in}{1.066362in}}%
\pgfpathclose%
\pgfusepath{fill}%
\end{pgfscope}%
\begin{pgfscope}%
\pgfpathrectangle{\pgfqpoint{2.893540in}{0.169444in}}{\pgfqpoint{0.858895in}{1.282944in}}%
\pgfusepath{clip}%
\pgfsetbuttcap%
\pgfsetmiterjoin%
\definecolor{currentfill}{rgb}{0.540570,0.034950,0.648640}%
\pgfsetfillcolor{currentfill}%
\pgfsetlinewidth{0.000000pt}%
\definecolor{currentstroke}{rgb}{0.000000,0.000000,0.000000}%
\pgfsetstrokecolor{currentstroke}%
\pgfsetstrokeopacity{0.000000}%
\pgfsetdash{}{0pt}%
\pgfpathmoveto{\pgfqpoint{3.142620in}{1.091435in}}%
\pgfpathlineto{\pgfqpoint{3.108264in}{1.091435in}}%
\pgfpathlineto{\pgfqpoint{3.108264in}{1.128091in}}%
\pgfpathlineto{\pgfqpoint{3.142620in}{1.128091in}}%
\pgfpathlineto{\pgfqpoint{3.142620in}{1.091435in}}%
\pgfpathclose%
\pgfusepath{fill}%
\end{pgfscope}%
\begin{pgfscope}%
\pgfpathrectangle{\pgfqpoint{2.893540in}{0.169444in}}{\pgfqpoint{0.858895in}{1.282944in}}%
\pgfusepath{clip}%
\pgfsetbuttcap%
\pgfsetmiterjoin%
\definecolor{currentfill}{rgb}{0.512206,0.018833,0.655209}%
\pgfsetfillcolor{currentfill}%
\pgfsetlinewidth{0.000000pt}%
\definecolor{currentstroke}{rgb}{0.000000,0.000000,0.000000}%
\pgfsetstrokecolor{currentstroke}%
\pgfsetstrokeopacity{0.000000}%
\pgfsetdash{}{0pt}%
\pgfpathmoveto{\pgfqpoint{3.142620in}{1.116509in}}%
\pgfpathlineto{\pgfqpoint{3.108264in}{1.116509in}}%
\pgfpathlineto{\pgfqpoint{3.108264in}{1.153164in}}%
\pgfpathlineto{\pgfqpoint{3.142620in}{1.153164in}}%
\pgfpathlineto{\pgfqpoint{3.142620in}{1.116509in}}%
\pgfpathclose%
\pgfusepath{fill}%
\end{pgfscope}%
\begin{pgfscope}%
\pgfpathrectangle{\pgfqpoint{2.893540in}{0.169444in}}{\pgfqpoint{0.858895in}{1.282944in}}%
\pgfusepath{clip}%
\pgfsetbuttcap%
\pgfsetmiterjoin%
\definecolor{currentfill}{rgb}{0.483210,0.008460,0.659095}%
\pgfsetfillcolor{currentfill}%
\pgfsetlinewidth{0.000000pt}%
\definecolor{currentstroke}{rgb}{0.000000,0.000000,0.000000}%
\pgfsetstrokecolor{currentstroke}%
\pgfsetstrokeopacity{0.000000}%
\pgfsetdash{}{0pt}%
\pgfpathmoveto{\pgfqpoint{3.142620in}{1.141582in}}%
\pgfpathlineto{\pgfqpoint{3.108264in}{1.141582in}}%
\pgfpathlineto{\pgfqpoint{3.108264in}{1.178237in}}%
\pgfpathlineto{\pgfqpoint{3.142620in}{1.178237in}}%
\pgfpathlineto{\pgfqpoint{3.142620in}{1.141582in}}%
\pgfpathclose%
\pgfusepath{fill}%
\end{pgfscope}%
\begin{pgfscope}%
\pgfpathrectangle{\pgfqpoint{2.893540in}{0.169444in}}{\pgfqpoint{0.858895in}{1.282944in}}%
\pgfusepath{clip}%
\pgfsetbuttcap%
\pgfsetmiterjoin%
\definecolor{currentfill}{rgb}{0.453677,0.002755,0.660310}%
\pgfsetfillcolor{currentfill}%
\pgfsetlinewidth{0.000000pt}%
\definecolor{currentstroke}{rgb}{0.000000,0.000000,0.000000}%
\pgfsetstrokecolor{currentstroke}%
\pgfsetstrokeopacity{0.000000}%
\pgfsetdash{}{0pt}%
\pgfpathmoveto{\pgfqpoint{3.142620in}{1.166655in}}%
\pgfpathlineto{\pgfqpoint{3.108264in}{1.166655in}}%
\pgfpathlineto{\pgfqpoint{3.108264in}{1.203310in}}%
\pgfpathlineto{\pgfqpoint{3.142620in}{1.203310in}}%
\pgfpathlineto{\pgfqpoint{3.142620in}{1.166655in}}%
\pgfpathclose%
\pgfusepath{fill}%
\end{pgfscope}%
\begin{pgfscope}%
\pgfpathrectangle{\pgfqpoint{2.893540in}{0.169444in}}{\pgfqpoint{0.858895in}{1.282944in}}%
\pgfusepath{clip}%
\pgfsetbuttcap%
\pgfsetmiterjoin%
\definecolor{currentfill}{rgb}{0.423689,0.000646,0.658956}%
\pgfsetfillcolor{currentfill}%
\pgfsetlinewidth{0.000000pt}%
\definecolor{currentstroke}{rgb}{0.000000,0.000000,0.000000}%
\pgfsetstrokecolor{currentstroke}%
\pgfsetstrokeopacity{0.000000}%
\pgfsetdash{}{0pt}%
\pgfpathmoveto{\pgfqpoint{3.142620in}{1.191728in}}%
\pgfpathlineto{\pgfqpoint{3.108264in}{1.191728in}}%
\pgfpathlineto{\pgfqpoint{3.108264in}{1.228384in}}%
\pgfpathlineto{\pgfqpoint{3.142620in}{1.228384in}}%
\pgfpathlineto{\pgfqpoint{3.142620in}{1.191728in}}%
\pgfpathclose%
\pgfusepath{fill}%
\end{pgfscope}%
\begin{pgfscope}%
\pgfpathrectangle{\pgfqpoint{2.893540in}{0.169444in}}{\pgfqpoint{0.858895in}{1.282944in}}%
\pgfusepath{clip}%
\pgfsetbuttcap%
\pgfsetmiterjoin%
\definecolor{currentfill}{rgb}{0.393304,0.001114,0.655199}%
\pgfsetfillcolor{currentfill}%
\pgfsetlinewidth{0.000000pt}%
\definecolor{currentstroke}{rgb}{0.000000,0.000000,0.000000}%
\pgfsetstrokecolor{currentstroke}%
\pgfsetstrokeopacity{0.000000}%
\pgfsetdash{}{0pt}%
\pgfpathmoveto{\pgfqpoint{3.142620in}{1.216801in}}%
\pgfpathlineto{\pgfqpoint{3.108264in}{1.216801in}}%
\pgfpathlineto{\pgfqpoint{3.108264in}{1.253457in}}%
\pgfpathlineto{\pgfqpoint{3.142620in}{1.253457in}}%
\pgfpathlineto{\pgfqpoint{3.142620in}{1.216801in}}%
\pgfpathclose%
\pgfusepath{fill}%
\end{pgfscope}%
\begin{pgfscope}%
\pgfpathrectangle{\pgfqpoint{2.893540in}{0.169444in}}{\pgfqpoint{0.858895in}{1.282944in}}%
\pgfusepath{clip}%
\pgfsetbuttcap%
\pgfsetmiterjoin%
\definecolor{currentfill}{rgb}{0.362553,0.003243,0.649245}%
\pgfsetfillcolor{currentfill}%
\pgfsetlinewidth{0.000000pt}%
\definecolor{currentstroke}{rgb}{0.000000,0.000000,0.000000}%
\pgfsetstrokecolor{currentstroke}%
\pgfsetstrokeopacity{0.000000}%
\pgfsetdash{}{0pt}%
\pgfpathmoveto{\pgfqpoint{3.142620in}{1.241874in}}%
\pgfpathlineto{\pgfqpoint{3.108264in}{1.241874in}}%
\pgfpathlineto{\pgfqpoint{3.108264in}{1.278530in}}%
\pgfpathlineto{\pgfqpoint{3.142620in}{1.278530in}}%
\pgfpathlineto{\pgfqpoint{3.142620in}{1.241874in}}%
\pgfpathclose%
\pgfusepath{fill}%
\end{pgfscope}%
\begin{pgfscope}%
\pgfpathrectangle{\pgfqpoint{2.893540in}{0.169444in}}{\pgfqpoint{0.858895in}{1.282944in}}%
\pgfusepath{clip}%
\pgfsetbuttcap%
\pgfsetmiterjoin%
\definecolor{currentfill}{rgb}{0.331426,0.006261,0.641316}%
\pgfsetfillcolor{currentfill}%
\pgfsetlinewidth{0.000000pt}%
\definecolor{currentstroke}{rgb}{0.000000,0.000000,0.000000}%
\pgfsetstrokecolor{currentstroke}%
\pgfsetstrokeopacity{0.000000}%
\pgfsetdash{}{0pt}%
\pgfpathmoveto{\pgfqpoint{3.142620in}{1.266948in}}%
\pgfpathlineto{\pgfqpoint{3.108264in}{1.266948in}}%
\pgfpathlineto{\pgfqpoint{3.108264in}{1.303603in}}%
\pgfpathlineto{\pgfqpoint{3.142620in}{1.303603in}}%
\pgfpathlineto{\pgfqpoint{3.142620in}{1.266948in}}%
\pgfpathclose%
\pgfusepath{fill}%
\end{pgfscope}%
\begin{pgfscope}%
\pgfpathrectangle{\pgfqpoint{2.893540in}{0.169444in}}{\pgfqpoint{0.858895in}{1.282944in}}%
\pgfusepath{clip}%
\pgfsetbuttcap%
\pgfsetmiterjoin%
\definecolor{currentfill}{rgb}{0.293478,0.010213,0.629490}%
\pgfsetfillcolor{currentfill}%
\pgfsetlinewidth{0.000000pt}%
\definecolor{currentstroke}{rgb}{0.000000,0.000000,0.000000}%
\pgfsetstrokecolor{currentstroke}%
\pgfsetstrokeopacity{0.000000}%
\pgfsetdash{}{0pt}%
\pgfpathmoveto{\pgfqpoint{3.142620in}{1.292021in}}%
\pgfpathlineto{\pgfqpoint{3.108264in}{1.292021in}}%
\pgfpathlineto{\pgfqpoint{3.108264in}{1.328676in}}%
\pgfpathlineto{\pgfqpoint{3.142620in}{1.328676in}}%
\pgfpathlineto{\pgfqpoint{3.142620in}{1.292021in}}%
\pgfpathclose%
\pgfusepath{fill}%
\end{pgfscope}%
\begin{pgfscope}%
\pgfpathrectangle{\pgfqpoint{2.893540in}{0.169444in}}{\pgfqpoint{0.858895in}{1.282944in}}%
\pgfusepath{clip}%
\pgfsetbuttcap%
\pgfsetmiterjoin%
\definecolor{currentfill}{rgb}{0.933232,0.482780,0.319325}%
\pgfsetfillcolor{currentfill}%
\pgfsetlinewidth{0.000000pt}%
\definecolor{currentstroke}{rgb}{0.000000,0.000000,0.000000}%
\pgfsetstrokecolor{currentstroke}%
\pgfsetstrokeopacity{0.000000}%
\pgfsetdash{}{0pt}%
\pgfpathmoveto{\pgfqpoint{3.030964in}{0.589178in}}%
\pgfpathlineto{\pgfqpoint{3.065319in}{0.589178in}}%
\pgfpathlineto{\pgfqpoint{3.065319in}{0.628056in}}%
\pgfpathlineto{\pgfqpoint{3.030964in}{0.628056in}}%
\pgfpathlineto{\pgfqpoint{3.030964in}{0.589178in}}%
\pgfpathclose%
\pgfusepath{fill}%
\end{pgfscope}%
\begin{pgfscope}%
\pgfpathrectangle{\pgfqpoint{2.893540in}{0.169444in}}{\pgfqpoint{0.858895in}{1.282944in}}%
\pgfusepath{clip}%
\pgfsetbuttcap%
\pgfsetmiterjoin%
\definecolor{currentfill}{rgb}{0.881443,0.392529,0.383229}%
\pgfsetfillcolor{currentfill}%
\pgfsetlinewidth{0.000000pt}%
\definecolor{currentstroke}{rgb}{0.000000,0.000000,0.000000}%
\pgfsetstrokecolor{currentstroke}%
\pgfsetstrokeopacity{0.000000}%
\pgfsetdash{}{0pt}%
\pgfpathmoveto{\pgfqpoint{3.048142in}{0.680862in}}%
\pgfpathlineto{\pgfqpoint{3.065319in}{0.680862in}}%
\pgfpathlineto{\pgfqpoint{3.065319in}{0.719739in}}%
\pgfpathlineto{\pgfqpoint{3.048142in}{0.719739in}}%
\pgfpathlineto{\pgfqpoint{3.048142in}{0.680862in}}%
\pgfpathclose%
\pgfusepath{fill}%
\end{pgfscope}%
\begin{pgfscope}%
\pgfpathrectangle{\pgfqpoint{2.893540in}{0.169444in}}{\pgfqpoint{0.858895in}{1.282944in}}%
\pgfusepath{clip}%
\pgfsetbuttcap%
\pgfsetmiterjoin%
\definecolor{currentfill}{rgb}{0.862927,0.365119,0.403519}%
\pgfsetfillcolor{currentfill}%
\pgfsetlinewidth{0.000000pt}%
\definecolor{currentstroke}{rgb}{0.000000,0.000000,0.000000}%
\pgfsetstrokecolor{currentstroke}%
\pgfsetstrokeopacity{0.000000}%
\pgfsetdash{}{0pt}%
\pgfpathmoveto{\pgfqpoint{3.013786in}{0.711410in}}%
\pgfpathlineto{\pgfqpoint{3.065319in}{0.711410in}}%
\pgfpathlineto{\pgfqpoint{3.065319in}{0.750287in}}%
\pgfpathlineto{\pgfqpoint{3.013786in}{0.750287in}}%
\pgfpathlineto{\pgfqpoint{3.013786in}{0.711410in}}%
\pgfpathclose%
\pgfusepath{fill}%
\end{pgfscope}%
\begin{pgfscope}%
\pgfpathrectangle{\pgfqpoint{2.893540in}{0.169444in}}{\pgfqpoint{0.858895in}{1.282944in}}%
\pgfusepath{clip}%
\pgfsetbuttcap%
\pgfsetmiterjoin%
\definecolor{currentfill}{rgb}{0.843484,0.338062,0.424013}%
\pgfsetfillcolor{currentfill}%
\pgfsetlinewidth{0.000000pt}%
\definecolor{currentstroke}{rgb}{0.000000,0.000000,0.000000}%
\pgfsetstrokecolor{currentstroke}%
\pgfsetstrokeopacity{0.000000}%
\pgfsetdash{}{0pt}%
\pgfpathmoveto{\pgfqpoint{3.048142in}{0.741998in}}%
\pgfpathlineto{\pgfqpoint{3.065319in}{0.741998in}}%
\pgfpathlineto{\pgfqpoint{3.065319in}{0.780875in}}%
\pgfpathlineto{\pgfqpoint{3.048142in}{0.780875in}}%
\pgfpathlineto{\pgfqpoint{3.048142in}{0.741998in}}%
\pgfpathclose%
\pgfusepath{fill}%
\end{pgfscope}%
\begin{pgfscope}%
\pgfpathrectangle{\pgfqpoint{2.893540in}{0.169444in}}{\pgfqpoint{0.858895in}{1.282944in}}%
\pgfusepath{clip}%
\pgfsetbuttcap%
\pgfsetmiterjoin%
\definecolor{currentfill}{rgb}{0.823132,0.311261,0.444806}%
\pgfsetfillcolor{currentfill}%
\pgfsetlinewidth{0.000000pt}%
\definecolor{currentstroke}{rgb}{0.000000,0.000000,0.000000}%
\pgfsetstrokecolor{currentstroke}%
\pgfsetstrokeopacity{0.000000}%
\pgfsetdash{}{0pt}%
\pgfpathmoveto{\pgfqpoint{3.013786in}{0.772546in}}%
\pgfpathlineto{\pgfqpoint{3.065319in}{0.772546in}}%
\pgfpathlineto{\pgfqpoint{3.065319in}{0.811423in}}%
\pgfpathlineto{\pgfqpoint{3.013786in}{0.811423in}}%
\pgfpathlineto{\pgfqpoint{3.013786in}{0.772546in}}%
\pgfpathclose%
\pgfusepath{fill}%
\end{pgfscope}%
\begin{pgfscope}%
\pgfpathrectangle{\pgfqpoint{2.893540in}{0.169444in}}{\pgfqpoint{0.858895in}{1.282944in}}%
\pgfusepath{clip}%
\pgfsetbuttcap%
\pgfsetmiterjoin%
\definecolor{currentfill}{rgb}{0.801855,0.284626,0.465971}%
\pgfsetfillcolor{currentfill}%
\pgfsetlinewidth{0.000000pt}%
\definecolor{currentstroke}{rgb}{0.000000,0.000000,0.000000}%
\pgfsetstrokecolor{currentstroke}%
\pgfsetstrokeopacity{0.000000}%
\pgfsetdash{}{0pt}%
\pgfpathmoveto{\pgfqpoint{3.030964in}{0.803094in}}%
\pgfpathlineto{\pgfqpoint{3.065319in}{0.803094in}}%
\pgfpathlineto{\pgfqpoint{3.065319in}{0.841971in}}%
\pgfpathlineto{\pgfqpoint{3.030964in}{0.841971in}}%
\pgfpathlineto{\pgfqpoint{3.030964in}{0.803094in}}%
\pgfpathclose%
\pgfusepath{fill}%
\end{pgfscope}%
\begin{pgfscope}%
\pgfpathrectangle{\pgfqpoint{2.893540in}{0.169444in}}{\pgfqpoint{0.858895in}{1.282944in}}%
\pgfusepath{clip}%
\pgfsetbuttcap%
\pgfsetmiterjoin%
\definecolor{currentfill}{rgb}{0.779604,0.258078,0.487539}%
\pgfsetfillcolor{currentfill}%
\pgfsetlinewidth{0.000000pt}%
\definecolor{currentstroke}{rgb}{0.000000,0.000000,0.000000}%
\pgfsetstrokecolor{currentstroke}%
\pgfsetstrokeopacity{0.000000}%
\pgfsetdash{}{0pt}%
\pgfpathmoveto{\pgfqpoint{3.030964in}{0.833642in}}%
\pgfpathlineto{\pgfqpoint{3.065319in}{0.833642in}}%
\pgfpathlineto{\pgfqpoint{3.065319in}{0.872519in}}%
\pgfpathlineto{\pgfqpoint{3.030964in}{0.872519in}}%
\pgfpathlineto{\pgfqpoint{3.030964in}{0.833642in}}%
\pgfpathclose%
\pgfusepath{fill}%
\end{pgfscope}%
\begin{pgfscope}%
\pgfpathrectangle{\pgfqpoint{2.893540in}{0.169444in}}{\pgfqpoint{0.858895in}{1.282944in}}%
\pgfusepath{clip}%
\pgfsetbuttcap%
\pgfsetmiterjoin%
\definecolor{currentfill}{rgb}{0.756304,0.231555,0.509468}%
\pgfsetfillcolor{currentfill}%
\pgfsetlinewidth{0.000000pt}%
\definecolor{currentstroke}{rgb}{0.000000,0.000000,0.000000}%
\pgfsetstrokecolor{currentstroke}%
\pgfsetstrokeopacity{0.000000}%
\pgfsetdash{}{0pt}%
\pgfpathmoveto{\pgfqpoint{3.013786in}{0.864190in}}%
\pgfpathlineto{\pgfqpoint{3.065319in}{0.864190in}}%
\pgfpathlineto{\pgfqpoint{3.065319in}{0.903067in}}%
\pgfpathlineto{\pgfqpoint{3.013786in}{0.903067in}}%
\pgfpathlineto{\pgfqpoint{3.013786in}{0.864190in}}%
\pgfpathclose%
\pgfusepath{fill}%
\end{pgfscope}%
\begin{pgfscope}%
\pgfpathrectangle{\pgfqpoint{2.893540in}{0.169444in}}{\pgfqpoint{0.858895in}{1.282944in}}%
\pgfusepath{clip}%
\pgfsetbuttcap%
\pgfsetmiterjoin%
\definecolor{currentfill}{rgb}{0.731862,0.205013,0.531601}%
\pgfsetfillcolor{currentfill}%
\pgfsetlinewidth{0.000000pt}%
\definecolor{currentstroke}{rgb}{0.000000,0.000000,0.000000}%
\pgfsetstrokecolor{currentstroke}%
\pgfsetstrokeopacity{0.000000}%
\pgfsetdash{}{0pt}%
\pgfpathmoveto{\pgfqpoint{2.962252in}{0.894777in}}%
\pgfpathlineto{\pgfqpoint{3.065319in}{0.894777in}}%
\pgfpathlineto{\pgfqpoint{3.065319in}{0.933655in}}%
\pgfpathlineto{\pgfqpoint{2.962252in}{0.933655in}}%
\pgfpathlineto{\pgfqpoint{2.962252in}{0.894777in}}%
\pgfpathclose%
\pgfusepath{fill}%
\end{pgfscope}%
\begin{pgfscope}%
\pgfpathrectangle{\pgfqpoint{2.893540in}{0.169444in}}{\pgfqpoint{0.858895in}{1.282944in}}%
\pgfusepath{clip}%
\pgfsetbuttcap%
\pgfsetmiterjoin%
\definecolor{currentfill}{rgb}{0.706178,0.178437,0.553657}%
\pgfsetfillcolor{currentfill}%
\pgfsetlinewidth{0.000000pt}%
\definecolor{currentstroke}{rgb}{0.000000,0.000000,0.000000}%
\pgfsetstrokecolor{currentstroke}%
\pgfsetstrokeopacity{0.000000}%
\pgfsetdash{}{0pt}%
\pgfpathmoveto{\pgfqpoint{2.979430in}{0.925326in}}%
\pgfpathlineto{\pgfqpoint{3.065319in}{0.925326in}}%
\pgfpathlineto{\pgfqpoint{3.065319in}{0.964203in}}%
\pgfpathlineto{\pgfqpoint{2.979430in}{0.964203in}}%
\pgfpathlineto{\pgfqpoint{2.979430in}{0.925326in}}%
\pgfpathclose%
\pgfusepath{fill}%
\end{pgfscope}%
\begin{pgfscope}%
\pgfpathrectangle{\pgfqpoint{2.893540in}{0.169444in}}{\pgfqpoint{0.858895in}{1.282944in}}%
\pgfusepath{clip}%
\pgfsetbuttcap%
\pgfsetmiterjoin%
\definecolor{currentfill}{rgb}{0.679160,0.151848,0.575189}%
\pgfsetfillcolor{currentfill}%
\pgfsetlinewidth{0.000000pt}%
\definecolor{currentstroke}{rgb}{0.000000,0.000000,0.000000}%
\pgfsetstrokecolor{currentstroke}%
\pgfsetstrokeopacity{0.000000}%
\pgfsetdash{}{0pt}%
\pgfpathmoveto{\pgfqpoint{2.996608in}{0.955874in}}%
\pgfpathlineto{\pgfqpoint{3.065319in}{0.955874in}}%
\pgfpathlineto{\pgfqpoint{3.065319in}{0.994751in}}%
\pgfpathlineto{\pgfqpoint{2.996608in}{0.994751in}}%
\pgfpathlineto{\pgfqpoint{2.996608in}{0.955874in}}%
\pgfpathclose%
\pgfusepath{fill}%
\end{pgfscope}%
\begin{pgfscope}%
\pgfpathrectangle{\pgfqpoint{2.893540in}{0.169444in}}{\pgfqpoint{0.858895in}{1.282944in}}%
\pgfusepath{clip}%
\pgfsetbuttcap%
\pgfsetmiterjoin%
\definecolor{currentfill}{rgb}{0.645872,0.120898,0.598867}%
\pgfsetfillcolor{currentfill}%
\pgfsetlinewidth{0.000000pt}%
\definecolor{currentstroke}{rgb}{0.000000,0.000000,0.000000}%
\pgfsetstrokecolor{currentstroke}%
\pgfsetstrokeopacity{0.000000}%
\pgfsetdash{}{0pt}%
\pgfpathmoveto{\pgfqpoint{2.996608in}{0.986461in}}%
\pgfpathlineto{\pgfqpoint{3.065319in}{0.986461in}}%
\pgfpathlineto{\pgfqpoint{3.065319in}{1.025338in}}%
\pgfpathlineto{\pgfqpoint{2.996608in}{1.025338in}}%
\pgfpathlineto{\pgfqpoint{2.996608in}{0.986461in}}%
\pgfpathclose%
\pgfusepath{fill}%
\end{pgfscope}%
\begin{pgfscope}%
\pgfpathrectangle{\pgfqpoint{2.893540in}{0.169444in}}{\pgfqpoint{0.858895in}{1.282944in}}%
\pgfusepath{clip}%
\pgfsetbuttcap%
\pgfsetmiterjoin%
\definecolor{currentfill}{rgb}{0.615812,0.094564,0.617140}%
\pgfsetfillcolor{currentfill}%
\pgfsetlinewidth{0.000000pt}%
\definecolor{currentstroke}{rgb}{0.000000,0.000000,0.000000}%
\pgfsetstrokecolor{currentstroke}%
\pgfsetstrokeopacity{0.000000}%
\pgfsetdash{}{0pt}%
\pgfpathmoveto{\pgfqpoint{3.013786in}{1.017009in}}%
\pgfpathlineto{\pgfqpoint{3.065319in}{1.017009in}}%
\pgfpathlineto{\pgfqpoint{3.065319in}{1.055886in}}%
\pgfpathlineto{\pgfqpoint{3.013786in}{1.055886in}}%
\pgfpathlineto{\pgfqpoint{3.013786in}{1.017009in}}%
\pgfpathclose%
\pgfusepath{fill}%
\end{pgfscope}%
\begin{pgfscope}%
\pgfpathrectangle{\pgfqpoint{2.893540in}{0.169444in}}{\pgfqpoint{0.858895in}{1.282944in}}%
\pgfusepath{clip}%
\pgfsetbuttcap%
\pgfsetmiterjoin%
\definecolor{currentfill}{rgb}{0.584391,0.068579,0.632812}%
\pgfsetfillcolor{currentfill}%
\pgfsetlinewidth{0.000000pt}%
\definecolor{currentstroke}{rgb}{0.000000,0.000000,0.000000}%
\pgfsetstrokecolor{currentstroke}%
\pgfsetstrokeopacity{0.000000}%
\pgfsetdash{}{0pt}%
\pgfpathmoveto{\pgfqpoint{3.030964in}{1.047557in}}%
\pgfpathlineto{\pgfqpoint{3.065319in}{1.047557in}}%
\pgfpathlineto{\pgfqpoint{3.065319in}{1.086434in}}%
\pgfpathlineto{\pgfqpoint{3.030964in}{1.086434in}}%
\pgfpathlineto{\pgfqpoint{3.030964in}{1.047557in}}%
\pgfpathclose%
\pgfusepath{fill}%
\end{pgfscope}%
\begin{pgfscope}%
\pgfpathrectangle{\pgfqpoint{2.893540in}{0.169444in}}{\pgfqpoint{0.858895in}{1.282944in}}%
\pgfusepath{clip}%
\pgfsetbuttcap%
\pgfsetmiterjoin%
\definecolor{currentfill}{rgb}{0.551715,0.043136,0.645277}%
\pgfsetfillcolor{currentfill}%
\pgfsetlinewidth{0.000000pt}%
\definecolor{currentstroke}{rgb}{0.000000,0.000000,0.000000}%
\pgfsetstrokecolor{currentstroke}%
\pgfsetstrokeopacity{0.000000}%
\pgfsetdash{}{0pt}%
\pgfpathmoveto{\pgfqpoint{3.013786in}{1.078105in}}%
\pgfpathlineto{\pgfqpoint{3.065319in}{1.078105in}}%
\pgfpathlineto{\pgfqpoint{3.065319in}{1.116983in}}%
\pgfpathlineto{\pgfqpoint{3.013786in}{1.116983in}}%
\pgfpathlineto{\pgfqpoint{3.013786in}{1.078105in}}%
\pgfpathclose%
\pgfusepath{fill}%
\end{pgfscope}%
\begin{pgfscope}%
\pgfpathrectangle{\pgfqpoint{2.893540in}{0.169444in}}{\pgfqpoint{0.858895in}{1.282944in}}%
\pgfusepath{clip}%
\pgfsetbuttcap%
\pgfsetmiterjoin%
\definecolor{currentfill}{rgb}{0.483210,0.008460,0.659095}%
\pgfsetfillcolor{currentfill}%
\pgfsetlinewidth{0.000000pt}%
\definecolor{currentstroke}{rgb}{0.000000,0.000000,0.000000}%
\pgfsetstrokecolor{currentstroke}%
\pgfsetstrokeopacity{0.000000}%
\pgfsetdash{}{0pt}%
\pgfpathmoveto{\pgfqpoint{3.013786in}{1.139241in}}%
\pgfpathlineto{\pgfqpoint{3.065319in}{1.139241in}}%
\pgfpathlineto{\pgfqpoint{3.065319in}{1.178118in}}%
\pgfpathlineto{\pgfqpoint{3.013786in}{1.178118in}}%
\pgfpathlineto{\pgfqpoint{3.013786in}{1.139241in}}%
\pgfpathclose%
\pgfusepath{fill}%
\end{pgfscope}%
\begin{pgfscope}%
\pgfpathrectangle{\pgfqpoint{2.893540in}{0.169444in}}{\pgfqpoint{0.858895in}{1.282944in}}%
\pgfusepath{clip}%
\pgfsetbuttcap%
\pgfsetmiterjoin%
\definecolor{currentfill}{rgb}{0.447714,0.002080,0.660240}%
\pgfsetfillcolor{currentfill}%
\pgfsetlinewidth{0.000000pt}%
\definecolor{currentstroke}{rgb}{0.000000,0.000000,0.000000}%
\pgfsetstrokecolor{currentstroke}%
\pgfsetstrokeopacity{0.000000}%
\pgfsetdash{}{0pt}%
\pgfpathmoveto{\pgfqpoint{3.030964in}{1.169789in}}%
\pgfpathlineto{\pgfqpoint{3.065319in}{1.169789in}}%
\pgfpathlineto{\pgfqpoint{3.065319in}{1.208666in}}%
\pgfpathlineto{\pgfqpoint{3.030964in}{1.208666in}}%
\pgfpathlineto{\pgfqpoint{3.030964in}{1.169789in}}%
\pgfpathclose%
\pgfusepath{fill}%
\end{pgfscope}%
\begin{pgfscope}%
\pgfpathrectangle{\pgfqpoint{2.893540in}{0.169444in}}{\pgfqpoint{0.858895in}{1.282944in}}%
\pgfusepath{clip}%
\pgfsetbuttcap%
\pgfsetmiterjoin%
\definecolor{currentfill}{rgb}{0.411580,0.000577,0.657730}%
\pgfsetfillcolor{currentfill}%
\pgfsetlinewidth{0.000000pt}%
\definecolor{currentstroke}{rgb}{0.000000,0.000000,0.000000}%
\pgfsetstrokecolor{currentstroke}%
\pgfsetstrokeopacity{0.000000}%
\pgfsetdash{}{0pt}%
\pgfpathmoveto{\pgfqpoint{3.048142in}{1.200337in}}%
\pgfpathlineto{\pgfqpoint{3.065319in}{1.200337in}}%
\pgfpathlineto{\pgfqpoint{3.065319in}{1.239214in}}%
\pgfpathlineto{\pgfqpoint{3.048142in}{1.239214in}}%
\pgfpathlineto{\pgfqpoint{3.048142in}{1.200337in}}%
\pgfpathclose%
\pgfusepath{fill}%
\end{pgfscope}%
\begin{pgfscope}%
\pgfpathrectangle{\pgfqpoint{2.893540in}{0.169444in}}{\pgfqpoint{0.858895in}{1.282944in}}%
\pgfusepath{clip}%
\pgfsetbuttcap%
\pgfsetmiterjoin%
\definecolor{currentfill}{rgb}{0.374897,0.002245,0.651876}%
\pgfsetfillcolor{currentfill}%
\pgfsetlinewidth{0.000000pt}%
\definecolor{currentstroke}{rgb}{0.000000,0.000000,0.000000}%
\pgfsetstrokecolor{currentstroke}%
\pgfsetstrokeopacity{0.000000}%
\pgfsetdash{}{0pt}%
\pgfpathmoveto{\pgfqpoint{3.048142in}{1.230925in}}%
\pgfpathlineto{\pgfqpoint{3.065319in}{1.230925in}}%
\pgfpathlineto{\pgfqpoint{3.065319in}{1.269802in}}%
\pgfpathlineto{\pgfqpoint{3.048142in}{1.269802in}}%
\pgfpathlineto{\pgfqpoint{3.048142in}{1.230925in}}%
\pgfpathclose%
\pgfusepath{fill}%
\end{pgfscope}%
\begin{pgfscope}%
\pgfpathrectangle{\pgfqpoint{2.893540in}{0.169444in}}{\pgfqpoint{0.858895in}{1.282944in}}%
\pgfusepath{clip}%
\pgfsetbuttcap%
\pgfsetmiterjoin%
\definecolor{currentfill}{rgb}{0.991209,0.790537,0.149377}%
\pgfsetfillcolor{currentfill}%
\pgfsetlinewidth{0.000000pt}%
\definecolor{currentstroke}{rgb}{0.000000,0.000000,0.000000}%
\pgfsetstrokecolor{currentstroke}%
\pgfsetstrokeopacity{0.000000}%
\pgfsetdash{}{0pt}%
\pgfpathmoveto{\pgfqpoint{3.374522in}{0.314167in}}%
\pgfpathlineto{\pgfqpoint{3.340166in}{0.314167in}}%
\pgfpathlineto{\pgfqpoint{3.340166in}{0.350822in}}%
\pgfpathlineto{\pgfqpoint{3.374522in}{0.350822in}}%
\pgfpathlineto{\pgfqpoint{3.374522in}{0.314167in}}%
\pgfpathclose%
\pgfusepath{fill}%
\end{pgfscope}%
\begin{pgfscope}%
\pgfpathrectangle{\pgfqpoint{2.893540in}{0.169444in}}{\pgfqpoint{0.858895in}{1.282944in}}%
\pgfusepath{clip}%
\pgfsetbuttcap%
\pgfsetmiterjoin%
\definecolor{currentfill}{rgb}{0.993851,0.759304,0.159092}%
\pgfsetfillcolor{currentfill}%
\pgfsetlinewidth{0.000000pt}%
\definecolor{currentstroke}{rgb}{0.000000,0.000000,0.000000}%
\pgfsetstrokecolor{currentstroke}%
\pgfsetstrokeopacity{0.000000}%
\pgfsetdash{}{0pt}%
\pgfpathmoveto{\pgfqpoint{3.374522in}{0.339240in}}%
\pgfpathlineto{\pgfqpoint{3.340166in}{0.339240in}}%
\pgfpathlineto{\pgfqpoint{3.340166in}{0.375896in}}%
\pgfpathlineto{\pgfqpoint{3.374522in}{0.375896in}}%
\pgfpathlineto{\pgfqpoint{3.374522in}{0.339240in}}%
\pgfpathclose%
\pgfusepath{fill}%
\end{pgfscope}%
\begin{pgfscope}%
\pgfpathrectangle{\pgfqpoint{2.893540in}{0.169444in}}{\pgfqpoint{0.858895in}{1.282944in}}%
\pgfusepath{clip}%
\pgfsetbuttcap%
\pgfsetmiterjoin%
\definecolor{currentfill}{rgb}{0.994553,0.728728,0.171622}%
\pgfsetfillcolor{currentfill}%
\pgfsetlinewidth{0.000000pt}%
\definecolor{currentstroke}{rgb}{0.000000,0.000000,0.000000}%
\pgfsetstrokecolor{currentstroke}%
\pgfsetstrokeopacity{0.000000}%
\pgfsetdash{}{0pt}%
\pgfpathmoveto{\pgfqpoint{3.374522in}{0.364313in}}%
\pgfpathlineto{\pgfqpoint{3.340166in}{0.364313in}}%
\pgfpathlineto{\pgfqpoint{3.340166in}{0.400969in}}%
\pgfpathlineto{\pgfqpoint{3.374522in}{0.400969in}}%
\pgfpathlineto{\pgfqpoint{3.374522in}{0.364313in}}%
\pgfpathclose%
\pgfusepath{fill}%
\end{pgfscope}%
\begin{pgfscope}%
\pgfpathrectangle{\pgfqpoint{2.893540in}{0.169444in}}{\pgfqpoint{0.858895in}{1.282944in}}%
\pgfusepath{clip}%
\pgfsetbuttcap%
\pgfsetmiterjoin%
\definecolor{currentfill}{rgb}{0.993456,0.698810,0.186041}%
\pgfsetfillcolor{currentfill}%
\pgfsetlinewidth{0.000000pt}%
\definecolor{currentstroke}{rgb}{0.000000,0.000000,0.000000}%
\pgfsetstrokecolor{currentstroke}%
\pgfsetstrokeopacity{0.000000}%
\pgfsetdash{}{0pt}%
\pgfpathmoveto{\pgfqpoint{3.374522in}{0.389386in}}%
\pgfpathlineto{\pgfqpoint{3.340166in}{0.389386in}}%
\pgfpathlineto{\pgfqpoint{3.340166in}{0.426042in}}%
\pgfpathlineto{\pgfqpoint{3.374522in}{0.426042in}}%
\pgfpathlineto{\pgfqpoint{3.374522in}{0.389386in}}%
\pgfpathclose%
\pgfusepath{fill}%
\end{pgfscope}%
\begin{pgfscope}%
\pgfpathrectangle{\pgfqpoint{2.893540in}{0.169444in}}{\pgfqpoint{0.858895in}{1.282944in}}%
\pgfusepath{clip}%
\pgfsetbuttcap%
\pgfsetmiterjoin%
\definecolor{currentfill}{rgb}{0.990681,0.669558,0.201642}%
\pgfsetfillcolor{currentfill}%
\pgfsetlinewidth{0.000000pt}%
\definecolor{currentstroke}{rgb}{0.000000,0.000000,0.000000}%
\pgfsetstrokecolor{currentstroke}%
\pgfsetstrokeopacity{0.000000}%
\pgfsetdash{}{0pt}%
\pgfpathmoveto{\pgfqpoint{3.374522in}{0.414460in}}%
\pgfpathlineto{\pgfqpoint{3.340166in}{0.414460in}}%
\pgfpathlineto{\pgfqpoint{3.340166in}{0.451115in}}%
\pgfpathlineto{\pgfqpoint{3.374522in}{0.451115in}}%
\pgfpathlineto{\pgfqpoint{3.374522in}{0.414460in}}%
\pgfpathclose%
\pgfusepath{fill}%
\end{pgfscope}%
\begin{pgfscope}%
\pgfpathrectangle{\pgfqpoint{2.893540in}{0.169444in}}{\pgfqpoint{0.858895in}{1.282944in}}%
\pgfusepath{clip}%
\pgfsetbuttcap%
\pgfsetmiterjoin%
\definecolor{currentfill}{rgb}{0.986345,0.640969,0.217948}%
\pgfsetfillcolor{currentfill}%
\pgfsetlinewidth{0.000000pt}%
\definecolor{currentstroke}{rgb}{0.000000,0.000000,0.000000}%
\pgfsetstrokecolor{currentstroke}%
\pgfsetstrokeopacity{0.000000}%
\pgfsetdash{}{0pt}%
\pgfpathmoveto{\pgfqpoint{3.374522in}{0.439533in}}%
\pgfpathlineto{\pgfqpoint{3.340166in}{0.439533in}}%
\pgfpathlineto{\pgfqpoint{3.340166in}{0.476188in}}%
\pgfpathlineto{\pgfqpoint{3.374522in}{0.476188in}}%
\pgfpathlineto{\pgfqpoint{3.374522in}{0.439533in}}%
\pgfpathclose%
\pgfusepath{fill}%
\end{pgfscope}%
\begin{pgfscope}%
\pgfpathrectangle{\pgfqpoint{2.893540in}{0.169444in}}{\pgfqpoint{0.858895in}{1.282944in}}%
\pgfusepath{clip}%
\pgfsetbuttcap%
\pgfsetmiterjoin%
\definecolor{currentfill}{rgb}{0.980556,0.613039,0.234646}%
\pgfsetfillcolor{currentfill}%
\pgfsetlinewidth{0.000000pt}%
\definecolor{currentstroke}{rgb}{0.000000,0.000000,0.000000}%
\pgfsetstrokecolor{currentstroke}%
\pgfsetstrokeopacity{0.000000}%
\pgfsetdash{}{0pt}%
\pgfpathmoveto{\pgfqpoint{3.374522in}{0.464606in}}%
\pgfpathlineto{\pgfqpoint{3.340166in}{0.464606in}}%
\pgfpathlineto{\pgfqpoint{3.340166in}{0.501261in}}%
\pgfpathlineto{\pgfqpoint{3.374522in}{0.501261in}}%
\pgfpathlineto{\pgfqpoint{3.374522in}{0.464606in}}%
\pgfpathclose%
\pgfusepath{fill}%
\end{pgfscope}%
\begin{pgfscope}%
\pgfpathrectangle{\pgfqpoint{2.893540in}{0.169444in}}{\pgfqpoint{0.858895in}{1.282944in}}%
\pgfusepath{clip}%
\pgfsetbuttcap%
\pgfsetmiterjoin%
\definecolor{currentfill}{rgb}{0.973416,0.585761,0.251540}%
\pgfsetfillcolor{currentfill}%
\pgfsetlinewidth{0.000000pt}%
\definecolor{currentstroke}{rgb}{0.000000,0.000000,0.000000}%
\pgfsetstrokecolor{currentstroke}%
\pgfsetstrokeopacity{0.000000}%
\pgfsetdash{}{0pt}%
\pgfpathmoveto{\pgfqpoint{3.374522in}{0.489679in}}%
\pgfpathlineto{\pgfqpoint{3.340166in}{0.489679in}}%
\pgfpathlineto{\pgfqpoint{3.340166in}{0.526335in}}%
\pgfpathlineto{\pgfqpoint{3.374522in}{0.526335in}}%
\pgfpathlineto{\pgfqpoint{3.374522in}{0.489679in}}%
\pgfpathclose%
\pgfusepath{fill}%
\end{pgfscope}%
\begin{pgfscope}%
\pgfpathrectangle{\pgfqpoint{2.893540in}{0.169444in}}{\pgfqpoint{0.858895in}{1.282944in}}%
\pgfusepath{clip}%
\pgfsetbuttcap%
\pgfsetmiterjoin%
\definecolor{currentfill}{rgb}{0.965024,0.559118,0.268513}%
\pgfsetfillcolor{currentfill}%
\pgfsetlinewidth{0.000000pt}%
\definecolor{currentstroke}{rgb}{0.000000,0.000000,0.000000}%
\pgfsetstrokecolor{currentstroke}%
\pgfsetstrokeopacity{0.000000}%
\pgfsetdash{}{0pt}%
\pgfpathmoveto{\pgfqpoint{3.374522in}{0.514752in}}%
\pgfpathlineto{\pgfqpoint{3.340166in}{0.514752in}}%
\pgfpathlineto{\pgfqpoint{3.340166in}{0.551408in}}%
\pgfpathlineto{\pgfqpoint{3.374522in}{0.551408in}}%
\pgfpathlineto{\pgfqpoint{3.374522in}{0.514752in}}%
\pgfpathclose%
\pgfusepath{fill}%
\end{pgfscope}%
\begin{pgfscope}%
\pgfpathrectangle{\pgfqpoint{2.893540in}{0.169444in}}{\pgfqpoint{0.858895in}{1.282944in}}%
\pgfusepath{clip}%
\pgfsetbuttcap%
\pgfsetmiterjoin%
\definecolor{currentfill}{rgb}{0.955470,0.533093,0.285490}%
\pgfsetfillcolor{currentfill}%
\pgfsetlinewidth{0.000000pt}%
\definecolor{currentstroke}{rgb}{0.000000,0.000000,0.000000}%
\pgfsetstrokecolor{currentstroke}%
\pgfsetstrokeopacity{0.000000}%
\pgfsetdash{}{0pt}%
\pgfpathmoveto{\pgfqpoint{3.374522in}{0.539825in}}%
\pgfpathlineto{\pgfqpoint{3.340166in}{0.539825in}}%
\pgfpathlineto{\pgfqpoint{3.340166in}{0.576481in}}%
\pgfpathlineto{\pgfqpoint{3.374522in}{0.576481in}}%
\pgfpathlineto{\pgfqpoint{3.374522in}{0.539825in}}%
\pgfpathclose%
\pgfusepath{fill}%
\end{pgfscope}%
\begin{pgfscope}%
\pgfpathrectangle{\pgfqpoint{2.893540in}{0.169444in}}{\pgfqpoint{0.858895in}{1.282944in}}%
\pgfusepath{clip}%
\pgfsetbuttcap%
\pgfsetmiterjoin%
\definecolor{currentfill}{rgb}{0.944844,0.507658,0.302433}%
\pgfsetfillcolor{currentfill}%
\pgfsetlinewidth{0.000000pt}%
\definecolor{currentstroke}{rgb}{0.000000,0.000000,0.000000}%
\pgfsetstrokecolor{currentstroke}%
\pgfsetstrokeopacity{0.000000}%
\pgfsetdash{}{0pt}%
\pgfpathmoveto{\pgfqpoint{3.374522in}{0.564899in}}%
\pgfpathlineto{\pgfqpoint{3.340166in}{0.564899in}}%
\pgfpathlineto{\pgfqpoint{3.340166in}{0.601554in}}%
\pgfpathlineto{\pgfqpoint{3.374522in}{0.601554in}}%
\pgfpathlineto{\pgfqpoint{3.374522in}{0.564899in}}%
\pgfpathclose%
\pgfusepath{fill}%
\end{pgfscope}%
\begin{pgfscope}%
\pgfpathrectangle{\pgfqpoint{2.893540in}{0.169444in}}{\pgfqpoint{0.858895in}{1.282944in}}%
\pgfusepath{clip}%
\pgfsetbuttcap%
\pgfsetmiterjoin%
\definecolor{currentfill}{rgb}{0.933232,0.482780,0.319325}%
\pgfsetfillcolor{currentfill}%
\pgfsetlinewidth{0.000000pt}%
\definecolor{currentstroke}{rgb}{0.000000,0.000000,0.000000}%
\pgfsetstrokecolor{currentstroke}%
\pgfsetstrokeopacity{0.000000}%
\pgfsetdash{}{0pt}%
\pgfpathmoveto{\pgfqpoint{3.374522in}{0.589972in}}%
\pgfpathlineto{\pgfqpoint{3.340166in}{0.589972in}}%
\pgfpathlineto{\pgfqpoint{3.340166in}{0.626627in}}%
\pgfpathlineto{\pgfqpoint{3.374522in}{0.626627in}}%
\pgfpathlineto{\pgfqpoint{3.374522in}{0.589972in}}%
\pgfpathclose%
\pgfusepath{fill}%
\end{pgfscope}%
\begin{pgfscope}%
\pgfpathrectangle{\pgfqpoint{2.893540in}{0.169444in}}{\pgfqpoint{0.858895in}{1.282944in}}%
\pgfusepath{clip}%
\pgfsetbuttcap%
\pgfsetmiterjoin%
\definecolor{currentfill}{rgb}{0.920714,0.458417,0.336166}%
\pgfsetfillcolor{currentfill}%
\pgfsetlinewidth{0.000000pt}%
\definecolor{currentstroke}{rgb}{0.000000,0.000000,0.000000}%
\pgfsetstrokecolor{currentstroke}%
\pgfsetstrokeopacity{0.000000}%
\pgfsetdash{}{0pt}%
\pgfpathmoveto{\pgfqpoint{3.374522in}{0.615045in}}%
\pgfpathlineto{\pgfqpoint{3.340166in}{0.615045in}}%
\pgfpathlineto{\pgfqpoint{3.340166in}{0.651701in}}%
\pgfpathlineto{\pgfqpoint{3.374522in}{0.651701in}}%
\pgfpathlineto{\pgfqpoint{3.374522in}{0.615045in}}%
\pgfpathclose%
\pgfusepath{fill}%
\end{pgfscope}%
\begin{pgfscope}%
\pgfpathrectangle{\pgfqpoint{2.893540in}{0.169444in}}{\pgfqpoint{0.858895in}{1.282944in}}%
\pgfusepath{clip}%
\pgfsetbuttcap%
\pgfsetmiterjoin%
\definecolor{currentfill}{rgb}{0.907365,0.434524,0.352970}%
\pgfsetfillcolor{currentfill}%
\pgfsetlinewidth{0.000000pt}%
\definecolor{currentstroke}{rgb}{0.000000,0.000000,0.000000}%
\pgfsetstrokecolor{currentstroke}%
\pgfsetstrokeopacity{0.000000}%
\pgfsetdash{}{0pt}%
\pgfpathmoveto{\pgfqpoint{3.374522in}{0.640118in}}%
\pgfpathlineto{\pgfqpoint{3.340166in}{0.640118in}}%
\pgfpathlineto{\pgfqpoint{3.340166in}{0.676774in}}%
\pgfpathlineto{\pgfqpoint{3.374522in}{0.676774in}}%
\pgfpathlineto{\pgfqpoint{3.374522in}{0.640118in}}%
\pgfpathclose%
\pgfusepath{fill}%
\end{pgfscope}%
\begin{pgfscope}%
\pgfpathrectangle{\pgfqpoint{2.893540in}{0.169444in}}{\pgfqpoint{0.858895in}{1.282944in}}%
\pgfusepath{clip}%
\pgfsetbuttcap%
\pgfsetmiterjoin%
\definecolor{currentfill}{rgb}{0.893250,0.411048,0.369768}%
\pgfsetfillcolor{currentfill}%
\pgfsetlinewidth{0.000000pt}%
\definecolor{currentstroke}{rgb}{0.000000,0.000000,0.000000}%
\pgfsetstrokecolor{currentstroke}%
\pgfsetstrokeopacity{0.000000}%
\pgfsetdash{}{0pt}%
\pgfpathmoveto{\pgfqpoint{3.374522in}{0.665191in}}%
\pgfpathlineto{\pgfqpoint{3.340166in}{0.665191in}}%
\pgfpathlineto{\pgfqpoint{3.340166in}{0.701847in}}%
\pgfpathlineto{\pgfqpoint{3.374522in}{0.701847in}}%
\pgfpathlineto{\pgfqpoint{3.374522in}{0.665191in}}%
\pgfpathclose%
\pgfusepath{fill}%
\end{pgfscope}%
\begin{pgfscope}%
\pgfpathrectangle{\pgfqpoint{2.893540in}{0.169444in}}{\pgfqpoint{0.858895in}{1.282944in}}%
\pgfusepath{clip}%
\pgfsetbuttcap%
\pgfsetmiterjoin%
\definecolor{currentfill}{rgb}{0.878423,0.387932,0.386600}%
\pgfsetfillcolor{currentfill}%
\pgfsetlinewidth{0.000000pt}%
\definecolor{currentstroke}{rgb}{0.000000,0.000000,0.000000}%
\pgfsetstrokecolor{currentstroke}%
\pgfsetstrokeopacity{0.000000}%
\pgfsetdash{}{0pt}%
\pgfpathmoveto{\pgfqpoint{3.374522in}{0.690265in}}%
\pgfpathlineto{\pgfqpoint{3.340166in}{0.690265in}}%
\pgfpathlineto{\pgfqpoint{3.340166in}{0.726920in}}%
\pgfpathlineto{\pgfqpoint{3.374522in}{0.726920in}}%
\pgfpathlineto{\pgfqpoint{3.374522in}{0.690265in}}%
\pgfpathclose%
\pgfusepath{fill}%
\end{pgfscope}%
\begin{pgfscope}%
\pgfpathrectangle{\pgfqpoint{2.893540in}{0.169444in}}{\pgfqpoint{0.858895in}{1.282944in}}%
\pgfusepath{clip}%
\pgfsetbuttcap%
\pgfsetmiterjoin%
\definecolor{currentfill}{rgb}{0.862927,0.365119,0.403519}%
\pgfsetfillcolor{currentfill}%
\pgfsetlinewidth{0.000000pt}%
\definecolor{currentstroke}{rgb}{0.000000,0.000000,0.000000}%
\pgfsetstrokecolor{currentstroke}%
\pgfsetstrokeopacity{0.000000}%
\pgfsetdash{}{0pt}%
\pgfpathmoveto{\pgfqpoint{3.374522in}{0.715338in}}%
\pgfpathlineto{\pgfqpoint{3.340166in}{0.715338in}}%
\pgfpathlineto{\pgfqpoint{3.340166in}{0.751993in}}%
\pgfpathlineto{\pgfqpoint{3.374522in}{0.751993in}}%
\pgfpathlineto{\pgfqpoint{3.374522in}{0.715338in}}%
\pgfpathclose%
\pgfusepath{fill}%
\end{pgfscope}%
\begin{pgfscope}%
\pgfpathrectangle{\pgfqpoint{2.893540in}{0.169444in}}{\pgfqpoint{0.858895in}{1.282944in}}%
\pgfusepath{clip}%
\pgfsetbuttcap%
\pgfsetmiterjoin%
\definecolor{currentfill}{rgb}{0.846788,0.342551,0.420579}%
\pgfsetfillcolor{currentfill}%
\pgfsetlinewidth{0.000000pt}%
\definecolor{currentstroke}{rgb}{0.000000,0.000000,0.000000}%
\pgfsetstrokecolor{currentstroke}%
\pgfsetstrokeopacity{0.000000}%
\pgfsetdash{}{0pt}%
\pgfpathmoveto{\pgfqpoint{3.374522in}{0.740411in}}%
\pgfpathlineto{\pgfqpoint{3.340166in}{0.740411in}}%
\pgfpathlineto{\pgfqpoint{3.340166in}{0.777066in}}%
\pgfpathlineto{\pgfqpoint{3.374522in}{0.777066in}}%
\pgfpathlineto{\pgfqpoint{3.374522in}{0.740411in}}%
\pgfpathclose%
\pgfusepath{fill}%
\end{pgfscope}%
\begin{pgfscope}%
\pgfpathrectangle{\pgfqpoint{2.893540in}{0.169444in}}{\pgfqpoint{0.858895in}{1.282944in}}%
\pgfusepath{clip}%
\pgfsetbuttcap%
\pgfsetmiterjoin%
\definecolor{currentfill}{rgb}{0.830018,0.320172,0.437836}%
\pgfsetfillcolor{currentfill}%
\pgfsetlinewidth{0.000000pt}%
\definecolor{currentstroke}{rgb}{0.000000,0.000000,0.000000}%
\pgfsetstrokecolor{currentstroke}%
\pgfsetstrokeopacity{0.000000}%
\pgfsetdash{}{0pt}%
\pgfpathmoveto{\pgfqpoint{3.374522in}{0.765484in}}%
\pgfpathlineto{\pgfqpoint{3.340166in}{0.765484in}}%
\pgfpathlineto{\pgfqpoint{3.340166in}{0.802140in}}%
\pgfpathlineto{\pgfqpoint{3.374522in}{0.802140in}}%
\pgfpathlineto{\pgfqpoint{3.374522in}{0.765484in}}%
\pgfpathclose%
\pgfusepath{fill}%
\end{pgfscope}%
\begin{pgfscope}%
\pgfpathrectangle{\pgfqpoint{2.893540in}{0.169444in}}{\pgfqpoint{0.858895in}{1.282944in}}%
\pgfusepath{clip}%
\pgfsetbuttcap%
\pgfsetmiterjoin%
\definecolor{currentfill}{rgb}{0.812612,0.297928,0.455338}%
\pgfsetfillcolor{currentfill}%
\pgfsetlinewidth{0.000000pt}%
\definecolor{currentstroke}{rgb}{0.000000,0.000000,0.000000}%
\pgfsetstrokecolor{currentstroke}%
\pgfsetstrokeopacity{0.000000}%
\pgfsetdash{}{0pt}%
\pgfpathmoveto{\pgfqpoint{3.374522in}{0.790557in}}%
\pgfpathlineto{\pgfqpoint{3.340166in}{0.790557in}}%
\pgfpathlineto{\pgfqpoint{3.340166in}{0.827213in}}%
\pgfpathlineto{\pgfqpoint{3.374522in}{0.827213in}}%
\pgfpathlineto{\pgfqpoint{3.374522in}{0.790557in}}%
\pgfpathclose%
\pgfusepath{fill}%
\end{pgfscope}%
\begin{pgfscope}%
\pgfpathrectangle{\pgfqpoint{2.893540in}{0.169444in}}{\pgfqpoint{0.858895in}{1.282944in}}%
\pgfusepath{clip}%
\pgfsetbuttcap%
\pgfsetmiterjoin%
\definecolor{currentfill}{rgb}{0.794549,0.275770,0.473117}%
\pgfsetfillcolor{currentfill}%
\pgfsetlinewidth{0.000000pt}%
\definecolor{currentstroke}{rgb}{0.000000,0.000000,0.000000}%
\pgfsetstrokecolor{currentstroke}%
\pgfsetstrokeopacity{0.000000}%
\pgfsetdash{}{0pt}%
\pgfpathmoveto{\pgfqpoint{3.374522in}{0.815630in}}%
\pgfpathlineto{\pgfqpoint{3.340166in}{0.815630in}}%
\pgfpathlineto{\pgfqpoint{3.340166in}{0.852286in}}%
\pgfpathlineto{\pgfqpoint{3.374522in}{0.852286in}}%
\pgfpathlineto{\pgfqpoint{3.374522in}{0.815630in}}%
\pgfpathclose%
\pgfusepath{fill}%
\end{pgfscope}%
\begin{pgfscope}%
\pgfpathrectangle{\pgfqpoint{2.893540in}{0.169444in}}{\pgfqpoint{0.858895in}{1.282944in}}%
\pgfusepath{clip}%
\pgfsetbuttcap%
\pgfsetmiterjoin%
\definecolor{currentfill}{rgb}{0.775796,0.253658,0.491171}%
\pgfsetfillcolor{currentfill}%
\pgfsetlinewidth{0.000000pt}%
\definecolor{currentstroke}{rgb}{0.000000,0.000000,0.000000}%
\pgfsetstrokecolor{currentstroke}%
\pgfsetstrokeopacity{0.000000}%
\pgfsetdash{}{0pt}%
\pgfpathmoveto{\pgfqpoint{3.374522in}{0.840704in}}%
\pgfpathlineto{\pgfqpoint{3.340166in}{0.840704in}}%
\pgfpathlineto{\pgfqpoint{3.340166in}{0.877359in}}%
\pgfpathlineto{\pgfqpoint{3.374522in}{0.877359in}}%
\pgfpathlineto{\pgfqpoint{3.374522in}{0.840704in}}%
\pgfpathclose%
\pgfusepath{fill}%
\end{pgfscope}%
\begin{pgfscope}%
\pgfpathrectangle{\pgfqpoint{2.893540in}{0.169444in}}{\pgfqpoint{0.858895in}{1.282944in}}%
\pgfusepath{clip}%
\pgfsetbuttcap%
\pgfsetmiterjoin%
\definecolor{currentfill}{rgb}{0.756304,0.231555,0.509468}%
\pgfsetfillcolor{currentfill}%
\pgfsetlinewidth{0.000000pt}%
\definecolor{currentstroke}{rgb}{0.000000,0.000000,0.000000}%
\pgfsetstrokecolor{currentstroke}%
\pgfsetstrokeopacity{0.000000}%
\pgfsetdash{}{0pt}%
\pgfpathmoveto{\pgfqpoint{3.374522in}{0.865777in}}%
\pgfpathlineto{\pgfqpoint{3.340166in}{0.865777in}}%
\pgfpathlineto{\pgfqpoint{3.340166in}{0.902432in}}%
\pgfpathlineto{\pgfqpoint{3.374522in}{0.902432in}}%
\pgfpathlineto{\pgfqpoint{3.374522in}{0.865777in}}%
\pgfpathclose%
\pgfusepath{fill}%
\end{pgfscope}%
\begin{pgfscope}%
\pgfpathrectangle{\pgfqpoint{2.893540in}{0.169444in}}{\pgfqpoint{0.858895in}{1.282944in}}%
\pgfusepath{clip}%
\pgfsetbuttcap%
\pgfsetmiterjoin%
\definecolor{currentfill}{rgb}{0.736019,0.209439,0.527908}%
\pgfsetfillcolor{currentfill}%
\pgfsetlinewidth{0.000000pt}%
\definecolor{currentstroke}{rgb}{0.000000,0.000000,0.000000}%
\pgfsetstrokecolor{currentstroke}%
\pgfsetstrokeopacity{0.000000}%
\pgfsetdash{}{0pt}%
\pgfpathmoveto{\pgfqpoint{3.374522in}{0.890850in}}%
\pgfpathlineto{\pgfqpoint{3.340166in}{0.890850in}}%
\pgfpathlineto{\pgfqpoint{3.340166in}{0.927506in}}%
\pgfpathlineto{\pgfqpoint{3.374522in}{0.927506in}}%
\pgfpathlineto{\pgfqpoint{3.374522in}{0.890850in}}%
\pgfpathclose%
\pgfusepath{fill}%
\end{pgfscope}%
\begin{pgfscope}%
\pgfpathrectangle{\pgfqpoint{2.893540in}{0.169444in}}{\pgfqpoint{0.858895in}{1.282944in}}%
\pgfusepath{clip}%
\pgfsetbuttcap%
\pgfsetmiterjoin%
\definecolor{currentfill}{rgb}{0.714883,0.187299,0.546338}%
\pgfsetfillcolor{currentfill}%
\pgfsetlinewidth{0.000000pt}%
\definecolor{currentstroke}{rgb}{0.000000,0.000000,0.000000}%
\pgfsetstrokecolor{currentstroke}%
\pgfsetstrokeopacity{0.000000}%
\pgfsetdash{}{0pt}%
\pgfpathmoveto{\pgfqpoint{3.374522in}{0.915923in}}%
\pgfpathlineto{\pgfqpoint{3.340166in}{0.915923in}}%
\pgfpathlineto{\pgfqpoint{3.340166in}{0.952579in}}%
\pgfpathlineto{\pgfqpoint{3.374522in}{0.952579in}}%
\pgfpathlineto{\pgfqpoint{3.374522in}{0.915923in}}%
\pgfpathclose%
\pgfusepath{fill}%
\end{pgfscope}%
\begin{pgfscope}%
\pgfpathrectangle{\pgfqpoint{2.893540in}{0.169444in}}{\pgfqpoint{0.858895in}{1.282944in}}%
\pgfusepath{clip}%
\pgfsetbuttcap%
\pgfsetmiterjoin%
\definecolor{currentfill}{rgb}{0.692840,0.165141,0.564522}%
\pgfsetfillcolor{currentfill}%
\pgfsetlinewidth{0.000000pt}%
\definecolor{currentstroke}{rgb}{0.000000,0.000000,0.000000}%
\pgfsetstrokecolor{currentstroke}%
\pgfsetstrokeopacity{0.000000}%
\pgfsetdash{}{0pt}%
\pgfpathmoveto{\pgfqpoint{3.374522in}{0.940996in}}%
\pgfpathlineto{\pgfqpoint{3.340166in}{0.940996in}}%
\pgfpathlineto{\pgfqpoint{3.340166in}{0.977652in}}%
\pgfpathlineto{\pgfqpoint{3.374522in}{0.977652in}}%
\pgfpathlineto{\pgfqpoint{3.374522in}{0.940996in}}%
\pgfpathclose%
\pgfusepath{fill}%
\end{pgfscope}%
\begin{pgfscope}%
\pgfpathrectangle{\pgfqpoint{2.893540in}{0.169444in}}{\pgfqpoint{0.858895in}{1.282944in}}%
\pgfusepath{clip}%
\pgfsetbuttcap%
\pgfsetmiterjoin%
\definecolor{currentfill}{rgb}{0.669845,0.142992,0.582154}%
\pgfsetfillcolor{currentfill}%
\pgfsetlinewidth{0.000000pt}%
\definecolor{currentstroke}{rgb}{0.000000,0.000000,0.000000}%
\pgfsetstrokecolor{currentstroke}%
\pgfsetstrokeopacity{0.000000}%
\pgfsetdash{}{0pt}%
\pgfpathmoveto{\pgfqpoint{3.374522in}{0.966070in}}%
\pgfpathlineto{\pgfqpoint{3.340166in}{0.966070in}}%
\pgfpathlineto{\pgfqpoint{3.340166in}{1.002725in}}%
\pgfpathlineto{\pgfqpoint{3.374522in}{1.002725in}}%
\pgfpathlineto{\pgfqpoint{3.374522in}{0.966070in}}%
\pgfpathclose%
\pgfusepath{fill}%
\end{pgfscope}%
\begin{pgfscope}%
\pgfpathrectangle{\pgfqpoint{2.893540in}{0.169444in}}{\pgfqpoint{0.858895in}{1.282944in}}%
\pgfusepath{clip}%
\pgfsetbuttcap%
\pgfsetmiterjoin%
\definecolor{currentfill}{rgb}{0.645872,0.120898,0.598867}%
\pgfsetfillcolor{currentfill}%
\pgfsetlinewidth{0.000000pt}%
\definecolor{currentstroke}{rgb}{0.000000,0.000000,0.000000}%
\pgfsetstrokecolor{currentstroke}%
\pgfsetstrokeopacity{0.000000}%
\pgfsetdash{}{0pt}%
\pgfpathmoveto{\pgfqpoint{3.374522in}{0.991143in}}%
\pgfpathlineto{\pgfqpoint{3.340166in}{0.991143in}}%
\pgfpathlineto{\pgfqpoint{3.340166in}{1.027798in}}%
\pgfpathlineto{\pgfqpoint{3.374522in}{1.027798in}}%
\pgfpathlineto{\pgfqpoint{3.374522in}{0.991143in}}%
\pgfpathclose%
\pgfusepath{fill}%
\end{pgfscope}%
\begin{pgfscope}%
\pgfpathrectangle{\pgfqpoint{2.893540in}{0.169444in}}{\pgfqpoint{0.858895in}{1.282944in}}%
\pgfusepath{clip}%
\pgfsetbuttcap%
\pgfsetmiterjoin%
\definecolor{currentfill}{rgb}{0.620919,0.098934,0.614257}%
\pgfsetfillcolor{currentfill}%
\pgfsetlinewidth{0.000000pt}%
\definecolor{currentstroke}{rgb}{0.000000,0.000000,0.000000}%
\pgfsetstrokecolor{currentstroke}%
\pgfsetstrokeopacity{0.000000}%
\pgfsetdash{}{0pt}%
\pgfpathmoveto{\pgfqpoint{3.374522in}{1.016216in}}%
\pgfpathlineto{\pgfqpoint{3.340166in}{1.016216in}}%
\pgfpathlineto{\pgfqpoint{3.340166in}{1.052871in}}%
\pgfpathlineto{\pgfqpoint{3.374522in}{1.052871in}}%
\pgfpathlineto{\pgfqpoint{3.374522in}{1.016216in}}%
\pgfpathclose%
\pgfusepath{fill}%
\end{pgfscope}%
\begin{pgfscope}%
\pgfpathrectangle{\pgfqpoint{2.893540in}{0.169444in}}{\pgfqpoint{0.858895in}{1.282944in}}%
\pgfusepath{clip}%
\pgfsetbuttcap%
\pgfsetmiterjoin%
\definecolor{currentfill}{rgb}{0.595011,0.077190,0.627917}%
\pgfsetfillcolor{currentfill}%
\pgfsetlinewidth{0.000000pt}%
\definecolor{currentstroke}{rgb}{0.000000,0.000000,0.000000}%
\pgfsetstrokecolor{currentstroke}%
\pgfsetstrokeopacity{0.000000}%
\pgfsetdash{}{0pt}%
\pgfpathmoveto{\pgfqpoint{3.374522in}{1.041289in}}%
\pgfpathlineto{\pgfqpoint{3.340166in}{1.041289in}}%
\pgfpathlineto{\pgfqpoint{3.340166in}{1.077945in}}%
\pgfpathlineto{\pgfqpoint{3.374522in}{1.077945in}}%
\pgfpathlineto{\pgfqpoint{3.374522in}{1.041289in}}%
\pgfpathclose%
\pgfusepath{fill}%
\end{pgfscope}%
\begin{pgfscope}%
\pgfpathrectangle{\pgfqpoint{2.893540in}{0.169444in}}{\pgfqpoint{0.858895in}{1.282944in}}%
\pgfusepath{clip}%
\pgfsetbuttcap%
\pgfsetmiterjoin%
\definecolor{currentfill}{rgb}{0.568201,0.055778,0.639477}%
\pgfsetfillcolor{currentfill}%
\pgfsetlinewidth{0.000000pt}%
\definecolor{currentstroke}{rgb}{0.000000,0.000000,0.000000}%
\pgfsetstrokecolor{currentstroke}%
\pgfsetstrokeopacity{0.000000}%
\pgfsetdash{}{0pt}%
\pgfpathmoveto{\pgfqpoint{3.374522in}{1.066362in}}%
\pgfpathlineto{\pgfqpoint{3.340166in}{1.066362in}}%
\pgfpathlineto{\pgfqpoint{3.340166in}{1.103018in}}%
\pgfpathlineto{\pgfqpoint{3.374522in}{1.103018in}}%
\pgfpathlineto{\pgfqpoint{3.374522in}{1.066362in}}%
\pgfpathclose%
\pgfusepath{fill}%
\end{pgfscope}%
\begin{pgfscope}%
\pgfpathrectangle{\pgfqpoint{2.893540in}{0.169444in}}{\pgfqpoint{0.858895in}{1.282944in}}%
\pgfusepath{clip}%
\pgfsetbuttcap%
\pgfsetmiterjoin%
\definecolor{currentfill}{rgb}{0.540570,0.034950,0.648640}%
\pgfsetfillcolor{currentfill}%
\pgfsetlinewidth{0.000000pt}%
\definecolor{currentstroke}{rgb}{0.000000,0.000000,0.000000}%
\pgfsetstrokecolor{currentstroke}%
\pgfsetstrokeopacity{0.000000}%
\pgfsetdash{}{0pt}%
\pgfpathmoveto{\pgfqpoint{3.374522in}{1.091435in}}%
\pgfpathlineto{\pgfqpoint{3.340166in}{1.091435in}}%
\pgfpathlineto{\pgfqpoint{3.340166in}{1.128091in}}%
\pgfpathlineto{\pgfqpoint{3.374522in}{1.128091in}}%
\pgfpathlineto{\pgfqpoint{3.374522in}{1.091435in}}%
\pgfpathclose%
\pgfusepath{fill}%
\end{pgfscope}%
\begin{pgfscope}%
\pgfpathrectangle{\pgfqpoint{2.893540in}{0.169444in}}{\pgfqpoint{0.858895in}{1.282944in}}%
\pgfusepath{clip}%
\pgfsetbuttcap%
\pgfsetmiterjoin%
\definecolor{currentfill}{rgb}{0.512206,0.018833,0.655209}%
\pgfsetfillcolor{currentfill}%
\pgfsetlinewidth{0.000000pt}%
\definecolor{currentstroke}{rgb}{0.000000,0.000000,0.000000}%
\pgfsetstrokecolor{currentstroke}%
\pgfsetstrokeopacity{0.000000}%
\pgfsetdash{}{0pt}%
\pgfpathmoveto{\pgfqpoint{3.374522in}{1.116509in}}%
\pgfpathlineto{\pgfqpoint{3.340166in}{1.116509in}}%
\pgfpathlineto{\pgfqpoint{3.340166in}{1.153164in}}%
\pgfpathlineto{\pgfqpoint{3.374522in}{1.153164in}}%
\pgfpathlineto{\pgfqpoint{3.374522in}{1.116509in}}%
\pgfpathclose%
\pgfusepath{fill}%
\end{pgfscope}%
\begin{pgfscope}%
\pgfpathrectangle{\pgfqpoint{2.893540in}{0.169444in}}{\pgfqpoint{0.858895in}{1.282944in}}%
\pgfusepath{clip}%
\pgfsetbuttcap%
\pgfsetmiterjoin%
\definecolor{currentfill}{rgb}{0.483210,0.008460,0.659095}%
\pgfsetfillcolor{currentfill}%
\pgfsetlinewidth{0.000000pt}%
\definecolor{currentstroke}{rgb}{0.000000,0.000000,0.000000}%
\pgfsetstrokecolor{currentstroke}%
\pgfsetstrokeopacity{0.000000}%
\pgfsetdash{}{0pt}%
\pgfpathmoveto{\pgfqpoint{3.374522in}{1.141582in}}%
\pgfpathlineto{\pgfqpoint{3.340166in}{1.141582in}}%
\pgfpathlineto{\pgfqpoint{3.340166in}{1.178237in}}%
\pgfpathlineto{\pgfqpoint{3.374522in}{1.178237in}}%
\pgfpathlineto{\pgfqpoint{3.374522in}{1.141582in}}%
\pgfpathclose%
\pgfusepath{fill}%
\end{pgfscope}%
\begin{pgfscope}%
\pgfpathrectangle{\pgfqpoint{2.893540in}{0.169444in}}{\pgfqpoint{0.858895in}{1.282944in}}%
\pgfusepath{clip}%
\pgfsetbuttcap%
\pgfsetmiterjoin%
\definecolor{currentfill}{rgb}{0.453677,0.002755,0.660310}%
\pgfsetfillcolor{currentfill}%
\pgfsetlinewidth{0.000000pt}%
\definecolor{currentstroke}{rgb}{0.000000,0.000000,0.000000}%
\pgfsetstrokecolor{currentstroke}%
\pgfsetstrokeopacity{0.000000}%
\pgfsetdash{}{0pt}%
\pgfpathmoveto{\pgfqpoint{3.374522in}{1.166655in}}%
\pgfpathlineto{\pgfqpoint{3.340166in}{1.166655in}}%
\pgfpathlineto{\pgfqpoint{3.340166in}{1.203310in}}%
\pgfpathlineto{\pgfqpoint{3.374522in}{1.203310in}}%
\pgfpathlineto{\pgfqpoint{3.374522in}{1.166655in}}%
\pgfpathclose%
\pgfusepath{fill}%
\end{pgfscope}%
\begin{pgfscope}%
\pgfpathrectangle{\pgfqpoint{2.893540in}{0.169444in}}{\pgfqpoint{0.858895in}{1.282944in}}%
\pgfusepath{clip}%
\pgfsetbuttcap%
\pgfsetmiterjoin%
\definecolor{currentfill}{rgb}{0.423689,0.000646,0.658956}%
\pgfsetfillcolor{currentfill}%
\pgfsetlinewidth{0.000000pt}%
\definecolor{currentstroke}{rgb}{0.000000,0.000000,0.000000}%
\pgfsetstrokecolor{currentstroke}%
\pgfsetstrokeopacity{0.000000}%
\pgfsetdash{}{0pt}%
\pgfpathmoveto{\pgfqpoint{3.374522in}{1.191728in}}%
\pgfpathlineto{\pgfqpoint{3.340166in}{1.191728in}}%
\pgfpathlineto{\pgfqpoint{3.340166in}{1.228384in}}%
\pgfpathlineto{\pgfqpoint{3.374522in}{1.228384in}}%
\pgfpathlineto{\pgfqpoint{3.374522in}{1.191728in}}%
\pgfpathclose%
\pgfusepath{fill}%
\end{pgfscope}%
\begin{pgfscope}%
\pgfpathrectangle{\pgfqpoint{2.893540in}{0.169444in}}{\pgfqpoint{0.858895in}{1.282944in}}%
\pgfusepath{clip}%
\pgfsetbuttcap%
\pgfsetmiterjoin%
\definecolor{currentfill}{rgb}{0.393304,0.001114,0.655199}%
\pgfsetfillcolor{currentfill}%
\pgfsetlinewidth{0.000000pt}%
\definecolor{currentstroke}{rgb}{0.000000,0.000000,0.000000}%
\pgfsetstrokecolor{currentstroke}%
\pgfsetstrokeopacity{0.000000}%
\pgfsetdash{}{0pt}%
\pgfpathmoveto{\pgfqpoint{3.374522in}{1.216801in}}%
\pgfpathlineto{\pgfqpoint{3.340166in}{1.216801in}}%
\pgfpathlineto{\pgfqpoint{3.340166in}{1.253457in}}%
\pgfpathlineto{\pgfqpoint{3.374522in}{1.253457in}}%
\pgfpathlineto{\pgfqpoint{3.374522in}{1.216801in}}%
\pgfpathclose%
\pgfusepath{fill}%
\end{pgfscope}%
\begin{pgfscope}%
\pgfpathrectangle{\pgfqpoint{2.893540in}{0.169444in}}{\pgfqpoint{0.858895in}{1.282944in}}%
\pgfusepath{clip}%
\pgfsetbuttcap%
\pgfsetmiterjoin%
\definecolor{currentfill}{rgb}{0.362553,0.003243,0.649245}%
\pgfsetfillcolor{currentfill}%
\pgfsetlinewidth{0.000000pt}%
\definecolor{currentstroke}{rgb}{0.000000,0.000000,0.000000}%
\pgfsetstrokecolor{currentstroke}%
\pgfsetstrokeopacity{0.000000}%
\pgfsetdash{}{0pt}%
\pgfpathmoveto{\pgfqpoint{3.374522in}{1.241874in}}%
\pgfpathlineto{\pgfqpoint{3.340166in}{1.241874in}}%
\pgfpathlineto{\pgfqpoint{3.340166in}{1.278530in}}%
\pgfpathlineto{\pgfqpoint{3.374522in}{1.278530in}}%
\pgfpathlineto{\pgfqpoint{3.374522in}{1.241874in}}%
\pgfpathclose%
\pgfusepath{fill}%
\end{pgfscope}%
\begin{pgfscope}%
\pgfpathrectangle{\pgfqpoint{2.893540in}{0.169444in}}{\pgfqpoint{0.858895in}{1.282944in}}%
\pgfusepath{clip}%
\pgfsetbuttcap%
\pgfsetmiterjoin%
\definecolor{currentfill}{rgb}{0.331426,0.006261,0.641316}%
\pgfsetfillcolor{currentfill}%
\pgfsetlinewidth{0.000000pt}%
\definecolor{currentstroke}{rgb}{0.000000,0.000000,0.000000}%
\pgfsetstrokecolor{currentstroke}%
\pgfsetstrokeopacity{0.000000}%
\pgfsetdash{}{0pt}%
\pgfpathmoveto{\pgfqpoint{3.374522in}{1.266948in}}%
\pgfpathlineto{\pgfqpoint{3.340166in}{1.266948in}}%
\pgfpathlineto{\pgfqpoint{3.340166in}{1.303603in}}%
\pgfpathlineto{\pgfqpoint{3.374522in}{1.303603in}}%
\pgfpathlineto{\pgfqpoint{3.374522in}{1.266948in}}%
\pgfpathclose%
\pgfusepath{fill}%
\end{pgfscope}%
\begin{pgfscope}%
\pgfpathrectangle{\pgfqpoint{2.893540in}{0.169444in}}{\pgfqpoint{0.858895in}{1.282944in}}%
\pgfusepath{clip}%
\pgfsetbuttcap%
\pgfsetmiterjoin%
\definecolor{currentfill}{rgb}{0.293478,0.010213,0.629490}%
\pgfsetfillcolor{currentfill}%
\pgfsetlinewidth{0.000000pt}%
\definecolor{currentstroke}{rgb}{0.000000,0.000000,0.000000}%
\pgfsetstrokecolor{currentstroke}%
\pgfsetstrokeopacity{0.000000}%
\pgfsetdash{}{0pt}%
\pgfpathmoveto{\pgfqpoint{3.374522in}{1.292021in}}%
\pgfpathlineto{\pgfqpoint{3.340166in}{1.292021in}}%
\pgfpathlineto{\pgfqpoint{3.340166in}{1.328676in}}%
\pgfpathlineto{\pgfqpoint{3.374522in}{1.328676in}}%
\pgfpathlineto{\pgfqpoint{3.374522in}{1.292021in}}%
\pgfpathclose%
\pgfusepath{fill}%
\end{pgfscope}%
\begin{pgfscope}%
\pgfpathrectangle{\pgfqpoint{2.893540in}{0.169444in}}{\pgfqpoint{0.858895in}{1.282944in}}%
\pgfusepath{clip}%
\pgfsetbuttcap%
\pgfsetmiterjoin%
\definecolor{currentfill}{rgb}{0.933232,0.482780,0.319325}%
\pgfsetfillcolor{currentfill}%
\pgfsetlinewidth{0.000000pt}%
\definecolor{currentstroke}{rgb}{0.000000,0.000000,0.000000}%
\pgfsetstrokecolor{currentstroke}%
\pgfsetstrokeopacity{0.000000}%
\pgfsetdash{}{0pt}%
\pgfpathmoveto{\pgfqpoint{3.408878in}{0.589178in}}%
\pgfpathlineto{\pgfqpoint{3.443233in}{0.589178in}}%
\pgfpathlineto{\pgfqpoint{3.443233in}{0.628056in}}%
\pgfpathlineto{\pgfqpoint{3.408878in}{0.628056in}}%
\pgfpathlineto{\pgfqpoint{3.408878in}{0.589178in}}%
\pgfpathclose%
\pgfusepath{fill}%
\end{pgfscope}%
\begin{pgfscope}%
\pgfpathrectangle{\pgfqpoint{2.893540in}{0.169444in}}{\pgfqpoint{0.858895in}{1.282944in}}%
\pgfusepath{clip}%
\pgfsetbuttcap%
\pgfsetmiterjoin%
\definecolor{currentfill}{rgb}{0.901807,0.425087,0.359688}%
\pgfsetfillcolor{currentfill}%
\pgfsetlinewidth{0.000000pt}%
\definecolor{currentstroke}{rgb}{0.000000,0.000000,0.000000}%
\pgfsetstrokecolor{currentstroke}%
\pgfsetstrokeopacity{0.000000}%
\pgfsetdash{}{0pt}%
\pgfpathmoveto{\pgfqpoint{3.408878in}{0.650314in}}%
\pgfpathlineto{\pgfqpoint{3.426056in}{0.650314in}}%
\pgfpathlineto{\pgfqpoint{3.426056in}{0.689191in}}%
\pgfpathlineto{\pgfqpoint{3.408878in}{0.689191in}}%
\pgfpathlineto{\pgfqpoint{3.408878in}{0.650314in}}%
\pgfpathclose%
\pgfusepath{fill}%
\end{pgfscope}%
\begin{pgfscope}%
\pgfpathrectangle{\pgfqpoint{2.893540in}{0.169444in}}{\pgfqpoint{0.858895in}{1.282944in}}%
\pgfusepath{clip}%
\pgfsetbuttcap%
\pgfsetmiterjoin%
\definecolor{currentfill}{rgb}{0.881443,0.392529,0.383229}%
\pgfsetfillcolor{currentfill}%
\pgfsetlinewidth{0.000000pt}%
\definecolor{currentstroke}{rgb}{0.000000,0.000000,0.000000}%
\pgfsetstrokecolor{currentstroke}%
\pgfsetstrokeopacity{0.000000}%
\pgfsetdash{}{0pt}%
\pgfpathmoveto{\pgfqpoint{3.408878in}{0.680862in}}%
\pgfpathlineto{\pgfqpoint{3.426056in}{0.680862in}}%
\pgfpathlineto{\pgfqpoint{3.426056in}{0.719739in}}%
\pgfpathlineto{\pgfqpoint{3.408878in}{0.719739in}}%
\pgfpathlineto{\pgfqpoint{3.408878in}{0.680862in}}%
\pgfpathclose%
\pgfusepath{fill}%
\end{pgfscope}%
\begin{pgfscope}%
\pgfpathrectangle{\pgfqpoint{2.893540in}{0.169444in}}{\pgfqpoint{0.858895in}{1.282944in}}%
\pgfusepath{clip}%
\pgfsetbuttcap%
\pgfsetmiterjoin%
\definecolor{currentfill}{rgb}{0.862927,0.365119,0.403519}%
\pgfsetfillcolor{currentfill}%
\pgfsetlinewidth{0.000000pt}%
\definecolor{currentstroke}{rgb}{0.000000,0.000000,0.000000}%
\pgfsetstrokecolor{currentstroke}%
\pgfsetstrokeopacity{0.000000}%
\pgfsetdash{}{0pt}%
\pgfpathmoveto{\pgfqpoint{3.408878in}{0.711410in}}%
\pgfpathlineto{\pgfqpoint{3.477589in}{0.711410in}}%
\pgfpathlineto{\pgfqpoint{3.477589in}{0.750287in}}%
\pgfpathlineto{\pgfqpoint{3.408878in}{0.750287in}}%
\pgfpathlineto{\pgfqpoint{3.408878in}{0.711410in}}%
\pgfpathclose%
\pgfusepath{fill}%
\end{pgfscope}%
\begin{pgfscope}%
\pgfpathrectangle{\pgfqpoint{2.893540in}{0.169444in}}{\pgfqpoint{0.858895in}{1.282944in}}%
\pgfusepath{clip}%
\pgfsetbuttcap%
\pgfsetmiterjoin%
\definecolor{currentfill}{rgb}{0.823132,0.311261,0.444806}%
\pgfsetfillcolor{currentfill}%
\pgfsetlinewidth{0.000000pt}%
\definecolor{currentstroke}{rgb}{0.000000,0.000000,0.000000}%
\pgfsetstrokecolor{currentstroke}%
\pgfsetstrokeopacity{0.000000}%
\pgfsetdash{}{0pt}%
\pgfpathmoveto{\pgfqpoint{3.408878in}{0.772546in}}%
\pgfpathlineto{\pgfqpoint{3.477589in}{0.772546in}}%
\pgfpathlineto{\pgfqpoint{3.477589in}{0.811423in}}%
\pgfpathlineto{\pgfqpoint{3.408878in}{0.811423in}}%
\pgfpathlineto{\pgfqpoint{3.408878in}{0.772546in}}%
\pgfpathclose%
\pgfusepath{fill}%
\end{pgfscope}%
\begin{pgfscope}%
\pgfpathrectangle{\pgfqpoint{2.893540in}{0.169444in}}{\pgfqpoint{0.858895in}{1.282944in}}%
\pgfusepath{clip}%
\pgfsetbuttcap%
\pgfsetmiterjoin%
\definecolor{currentfill}{rgb}{0.801855,0.284626,0.465971}%
\pgfsetfillcolor{currentfill}%
\pgfsetlinewidth{0.000000pt}%
\definecolor{currentstroke}{rgb}{0.000000,0.000000,0.000000}%
\pgfsetstrokecolor{currentstroke}%
\pgfsetstrokeopacity{0.000000}%
\pgfsetdash{}{0pt}%
\pgfpathmoveto{\pgfqpoint{3.408878in}{0.803094in}}%
\pgfpathlineto{\pgfqpoint{3.494767in}{0.803094in}}%
\pgfpathlineto{\pgfqpoint{3.494767in}{0.841971in}}%
\pgfpathlineto{\pgfqpoint{3.408878in}{0.841971in}}%
\pgfpathlineto{\pgfqpoint{3.408878in}{0.803094in}}%
\pgfpathclose%
\pgfusepath{fill}%
\end{pgfscope}%
\begin{pgfscope}%
\pgfpathrectangle{\pgfqpoint{2.893540in}{0.169444in}}{\pgfqpoint{0.858895in}{1.282944in}}%
\pgfusepath{clip}%
\pgfsetbuttcap%
\pgfsetmiterjoin%
\definecolor{currentfill}{rgb}{0.779604,0.258078,0.487539}%
\pgfsetfillcolor{currentfill}%
\pgfsetlinewidth{0.000000pt}%
\definecolor{currentstroke}{rgb}{0.000000,0.000000,0.000000}%
\pgfsetstrokecolor{currentstroke}%
\pgfsetstrokeopacity{0.000000}%
\pgfsetdash{}{0pt}%
\pgfpathmoveto{\pgfqpoint{3.408878in}{0.833642in}}%
\pgfpathlineto{\pgfqpoint{3.426056in}{0.833642in}}%
\pgfpathlineto{\pgfqpoint{3.426056in}{0.872519in}}%
\pgfpathlineto{\pgfqpoint{3.408878in}{0.872519in}}%
\pgfpathlineto{\pgfqpoint{3.408878in}{0.833642in}}%
\pgfpathclose%
\pgfusepath{fill}%
\end{pgfscope}%
\begin{pgfscope}%
\pgfpathrectangle{\pgfqpoint{2.893540in}{0.169444in}}{\pgfqpoint{0.858895in}{1.282944in}}%
\pgfusepath{clip}%
\pgfsetbuttcap%
\pgfsetmiterjoin%
\definecolor{currentfill}{rgb}{0.756304,0.231555,0.509468}%
\pgfsetfillcolor{currentfill}%
\pgfsetlinewidth{0.000000pt}%
\definecolor{currentstroke}{rgb}{0.000000,0.000000,0.000000}%
\pgfsetstrokecolor{currentstroke}%
\pgfsetstrokeopacity{0.000000}%
\pgfsetdash{}{0pt}%
\pgfpathmoveto{\pgfqpoint{3.408878in}{0.864190in}}%
\pgfpathlineto{\pgfqpoint{3.494767in}{0.864190in}}%
\pgfpathlineto{\pgfqpoint{3.494767in}{0.903067in}}%
\pgfpathlineto{\pgfqpoint{3.408878in}{0.903067in}}%
\pgfpathlineto{\pgfqpoint{3.408878in}{0.864190in}}%
\pgfpathclose%
\pgfusepath{fill}%
\end{pgfscope}%
\begin{pgfscope}%
\pgfpathrectangle{\pgfqpoint{2.893540in}{0.169444in}}{\pgfqpoint{0.858895in}{1.282944in}}%
\pgfusepath{clip}%
\pgfsetbuttcap%
\pgfsetmiterjoin%
\definecolor{currentfill}{rgb}{0.731862,0.205013,0.531601}%
\pgfsetfillcolor{currentfill}%
\pgfsetlinewidth{0.000000pt}%
\definecolor{currentstroke}{rgb}{0.000000,0.000000,0.000000}%
\pgfsetstrokecolor{currentstroke}%
\pgfsetstrokeopacity{0.000000}%
\pgfsetdash{}{0pt}%
\pgfpathmoveto{\pgfqpoint{3.408878in}{0.894777in}}%
\pgfpathlineto{\pgfqpoint{3.426056in}{0.894777in}}%
\pgfpathlineto{\pgfqpoint{3.426056in}{0.933655in}}%
\pgfpathlineto{\pgfqpoint{3.408878in}{0.933655in}}%
\pgfpathlineto{\pgfqpoint{3.408878in}{0.894777in}}%
\pgfpathclose%
\pgfusepath{fill}%
\end{pgfscope}%
\begin{pgfscope}%
\pgfpathrectangle{\pgfqpoint{2.893540in}{0.169444in}}{\pgfqpoint{0.858895in}{1.282944in}}%
\pgfusepath{clip}%
\pgfsetbuttcap%
\pgfsetmiterjoin%
\definecolor{currentfill}{rgb}{0.706178,0.178437,0.553657}%
\pgfsetfillcolor{currentfill}%
\pgfsetlinewidth{0.000000pt}%
\definecolor{currentstroke}{rgb}{0.000000,0.000000,0.000000}%
\pgfsetstrokecolor{currentstroke}%
\pgfsetstrokeopacity{0.000000}%
\pgfsetdash{}{0pt}%
\pgfpathmoveto{\pgfqpoint{3.408878in}{0.925326in}}%
\pgfpathlineto{\pgfqpoint{3.529123in}{0.925326in}}%
\pgfpathlineto{\pgfqpoint{3.529123in}{0.964203in}}%
\pgfpathlineto{\pgfqpoint{3.408878in}{0.964203in}}%
\pgfpathlineto{\pgfqpoint{3.408878in}{0.925326in}}%
\pgfpathclose%
\pgfusepath{fill}%
\end{pgfscope}%
\begin{pgfscope}%
\pgfpathrectangle{\pgfqpoint{2.893540in}{0.169444in}}{\pgfqpoint{0.858895in}{1.282944in}}%
\pgfusepath{clip}%
\pgfsetbuttcap%
\pgfsetmiterjoin%
\definecolor{currentfill}{rgb}{0.679160,0.151848,0.575189}%
\pgfsetfillcolor{currentfill}%
\pgfsetlinewidth{0.000000pt}%
\definecolor{currentstroke}{rgb}{0.000000,0.000000,0.000000}%
\pgfsetstrokecolor{currentstroke}%
\pgfsetstrokeopacity{0.000000}%
\pgfsetdash{}{0pt}%
\pgfpathmoveto{\pgfqpoint{3.408878in}{0.955874in}}%
\pgfpathlineto{\pgfqpoint{3.477589in}{0.955874in}}%
\pgfpathlineto{\pgfqpoint{3.477589in}{0.994751in}}%
\pgfpathlineto{\pgfqpoint{3.408878in}{0.994751in}}%
\pgfpathlineto{\pgfqpoint{3.408878in}{0.955874in}}%
\pgfpathclose%
\pgfusepath{fill}%
\end{pgfscope}%
\begin{pgfscope}%
\pgfpathrectangle{\pgfqpoint{2.893540in}{0.169444in}}{\pgfqpoint{0.858895in}{1.282944in}}%
\pgfusepath{clip}%
\pgfsetbuttcap%
\pgfsetmiterjoin%
\definecolor{currentfill}{rgb}{0.645872,0.120898,0.598867}%
\pgfsetfillcolor{currentfill}%
\pgfsetlinewidth{0.000000pt}%
\definecolor{currentstroke}{rgb}{0.000000,0.000000,0.000000}%
\pgfsetstrokecolor{currentstroke}%
\pgfsetstrokeopacity{0.000000}%
\pgfsetdash{}{0pt}%
\pgfpathmoveto{\pgfqpoint{3.408878in}{0.986461in}}%
\pgfpathlineto{\pgfqpoint{3.511945in}{0.986461in}}%
\pgfpathlineto{\pgfqpoint{3.511945in}{1.025338in}}%
\pgfpathlineto{\pgfqpoint{3.408878in}{1.025338in}}%
\pgfpathlineto{\pgfqpoint{3.408878in}{0.986461in}}%
\pgfpathclose%
\pgfusepath{fill}%
\end{pgfscope}%
\begin{pgfscope}%
\pgfpathrectangle{\pgfqpoint{2.893540in}{0.169444in}}{\pgfqpoint{0.858895in}{1.282944in}}%
\pgfusepath{clip}%
\pgfsetbuttcap%
\pgfsetmiterjoin%
\definecolor{currentfill}{rgb}{0.584391,0.068579,0.632812}%
\pgfsetfillcolor{currentfill}%
\pgfsetlinewidth{0.000000pt}%
\definecolor{currentstroke}{rgb}{0.000000,0.000000,0.000000}%
\pgfsetstrokecolor{currentstroke}%
\pgfsetstrokeopacity{0.000000}%
\pgfsetdash{}{0pt}%
\pgfpathmoveto{\pgfqpoint{3.408878in}{1.047557in}}%
\pgfpathlineto{\pgfqpoint{3.426056in}{1.047557in}}%
\pgfpathlineto{\pgfqpoint{3.426056in}{1.086434in}}%
\pgfpathlineto{\pgfqpoint{3.408878in}{1.086434in}}%
\pgfpathlineto{\pgfqpoint{3.408878in}{1.047557in}}%
\pgfpathclose%
\pgfusepath{fill}%
\end{pgfscope}%
\begin{pgfscope}%
\pgfpathrectangle{\pgfqpoint{2.893540in}{0.169444in}}{\pgfqpoint{0.858895in}{1.282944in}}%
\pgfusepath{clip}%
\pgfsetbuttcap%
\pgfsetmiterjoin%
\definecolor{currentfill}{rgb}{0.551715,0.043136,0.645277}%
\pgfsetfillcolor{currentfill}%
\pgfsetlinewidth{0.000000pt}%
\definecolor{currentstroke}{rgb}{0.000000,0.000000,0.000000}%
\pgfsetstrokecolor{currentstroke}%
\pgfsetstrokeopacity{0.000000}%
\pgfsetdash{}{0pt}%
\pgfpathmoveto{\pgfqpoint{3.408878in}{1.078105in}}%
\pgfpathlineto{\pgfqpoint{3.426056in}{1.078105in}}%
\pgfpathlineto{\pgfqpoint{3.426056in}{1.116983in}}%
\pgfpathlineto{\pgfqpoint{3.408878in}{1.116983in}}%
\pgfpathlineto{\pgfqpoint{3.408878in}{1.078105in}}%
\pgfpathclose%
\pgfusepath{fill}%
\end{pgfscope}%
\begin{pgfscope}%
\pgfpathrectangle{\pgfqpoint{2.893540in}{0.169444in}}{\pgfqpoint{0.858895in}{1.282944in}}%
\pgfusepath{clip}%
\pgfsetbuttcap%
\pgfsetmiterjoin%
\definecolor{currentfill}{rgb}{0.483210,0.008460,0.659095}%
\pgfsetfillcolor{currentfill}%
\pgfsetlinewidth{0.000000pt}%
\definecolor{currentstroke}{rgb}{0.000000,0.000000,0.000000}%
\pgfsetstrokecolor{currentstroke}%
\pgfsetstrokeopacity{0.000000}%
\pgfsetdash{}{0pt}%
\pgfpathmoveto{\pgfqpoint{3.408878in}{1.139241in}}%
\pgfpathlineto{\pgfqpoint{3.477589in}{1.139241in}}%
\pgfpathlineto{\pgfqpoint{3.477589in}{1.178118in}}%
\pgfpathlineto{\pgfqpoint{3.408878in}{1.178118in}}%
\pgfpathlineto{\pgfqpoint{3.408878in}{1.139241in}}%
\pgfpathclose%
\pgfusepath{fill}%
\end{pgfscope}%
\begin{pgfscope}%
\pgfpathrectangle{\pgfqpoint{2.893540in}{0.169444in}}{\pgfqpoint{0.858895in}{1.282944in}}%
\pgfusepath{clip}%
\pgfsetbuttcap%
\pgfsetmiterjoin%
\definecolor{currentfill}{rgb}{0.447714,0.002080,0.660240}%
\pgfsetfillcolor{currentfill}%
\pgfsetlinewidth{0.000000pt}%
\definecolor{currentstroke}{rgb}{0.000000,0.000000,0.000000}%
\pgfsetstrokecolor{currentstroke}%
\pgfsetstrokeopacity{0.000000}%
\pgfsetdash{}{0pt}%
\pgfpathmoveto{\pgfqpoint{3.408878in}{1.169789in}}%
\pgfpathlineto{\pgfqpoint{3.460411in}{1.169789in}}%
\pgfpathlineto{\pgfqpoint{3.460411in}{1.208666in}}%
\pgfpathlineto{\pgfqpoint{3.408878in}{1.208666in}}%
\pgfpathlineto{\pgfqpoint{3.408878in}{1.169789in}}%
\pgfpathclose%
\pgfusepath{fill}%
\end{pgfscope}%
\begin{pgfscope}%
\pgfpathrectangle{\pgfqpoint{2.893540in}{0.169444in}}{\pgfqpoint{0.858895in}{1.282944in}}%
\pgfusepath{clip}%
\pgfsetbuttcap%
\pgfsetmiterjoin%
\definecolor{currentfill}{rgb}{0.411580,0.000577,0.657730}%
\pgfsetfillcolor{currentfill}%
\pgfsetlinewidth{0.000000pt}%
\definecolor{currentstroke}{rgb}{0.000000,0.000000,0.000000}%
\pgfsetstrokecolor{currentstroke}%
\pgfsetstrokeopacity{0.000000}%
\pgfsetdash{}{0pt}%
\pgfpathmoveto{\pgfqpoint{3.408878in}{1.200337in}}%
\pgfpathlineto{\pgfqpoint{3.426056in}{1.200337in}}%
\pgfpathlineto{\pgfqpoint{3.426056in}{1.239214in}}%
\pgfpathlineto{\pgfqpoint{3.408878in}{1.239214in}}%
\pgfpathlineto{\pgfqpoint{3.408878in}{1.200337in}}%
\pgfpathclose%
\pgfusepath{fill}%
\end{pgfscope}%
\begin{pgfscope}%
\definecolor{textcolor}{rgb}{0.000000,0.000000,0.000000}%
\pgfsetstrokecolor{textcolor}%
\pgfsetfillcolor{textcolor}%
\pgftext[x=3.237099in,y=0.314167in,,]{\color{textcolor}\setmainfont{Lato}\rmfamily\fontsize{8.000000}{9.600000}\selectfont 45}%
\end{pgfscope}%
\begin{pgfscope}%
\definecolor{textcolor}{rgb}{0.000000,0.000000,0.000000}%
\pgfsetstrokecolor{textcolor}%
\pgfsetfillcolor{textcolor}%
\pgftext[x=3.237099in,y=0.509738in,,]{\color{textcolor}\setmainfont{Lato}\rmfamily\fontsize{8.000000}{9.600000}\selectfont 50}%
\end{pgfscope}%
\begin{pgfscope}%
\definecolor{textcolor}{rgb}{0.000000,0.000000,0.000000}%
\pgfsetstrokecolor{textcolor}%
\pgfsetfillcolor{textcolor}%
\pgftext[x=3.237099in,y=0.900879in,,]{\color{textcolor}\setmainfont{Lato}\rmfamily\fontsize{8.000000}{9.600000}\selectfont 60}%
\end{pgfscope}%
\begin{pgfscope}%
\definecolor{textcolor}{rgb}{0.000000,0.000000,0.000000}%
\pgfsetstrokecolor{textcolor}%
\pgfsetfillcolor{textcolor}%
\pgftext[x=3.237099in,y=1.292021in,,]{\color{textcolor}\setmainfont{Lato}\rmfamily\fontsize{8.000000}{9.600000}\selectfont 70}%
\end{pgfscope}%
\end{pgfpicture}%
\makeatother%
\endgroup%
 \\
\footnotesize{Source: Bureau of Labor Statistics}
\newpage
\begin{minipage}{0.76\textwidth} 
\small A tight local labor market will employ those ages 25 to 54 at a very high rate, barring any local labor supply constraints, for example availability or cost of child care or high rates of disability. \input{text/state_pa_emp_t3.txt}

\input{text/state_pa_emp.txt}
\end{minipage}

\normalsize \textbf{Age 25--54 Employment Rate by State}\\
\footnotesize{\textit{employed share of age 25--54 population, percent, not seasonally adjusted}}\\
\vspace{-3mm}
\hspace{-9mm} %% Creator: Matplotlib, PGF backend
%%
%% To include the figure in your LaTeX document, write
%%   \input{<filename>.pgf}
%%
%% Make sure the required packages are loaded in your preamble
%%   \usepackage{pgf}
%%
%% Also ensure that all the required font packages are loaded; for instance,
%% the lmodern package is sometimes necessary when using math font.
%%   \usepackage{lmodern}
%%
%% Figures using additional raster images can only be included by \input if
%% they are in the same directory as the main LaTeX file. For loading figures
%% from other directories you can use the `import` package
%%   \usepackage{import}
%%
%% and then include the figures with
%%   \import{<path to file>}{<filename>.pgf}
%%
%% Matplotlib used the following preamble
%%   
%%   \usepackage{fontspec}
%%   \setmainfont{DejaVuSerif.ttf}[Path=\detokenize{/home/brian/miniconda3/lib/python3.8/site-packages/matplotlib/mpl-data/fonts/ttf/}]
%%   \setsansfont{DejaVuSans.ttf}[Path=\detokenize{/home/brian/miniconda3/lib/python3.8/site-packages/matplotlib/mpl-data/fonts/ttf/}]
%%   \setmonofont{DejaVuSansMono.ttf}[Path=\detokenize{/home/brian/miniconda3/lib/python3.8/site-packages/matplotlib/mpl-data/fonts/ttf/}]
%%   \makeatletter\@ifpackageloaded{underscore}{}{\usepackage[strings]{underscore}}\makeatother
%%
\begingroup%
\makeatletter%
\begin{pgfpicture}%
\pgfpathrectangle{\pgfpointorigin}{\pgfqpoint{6.487985in}{2.069857in}}%
\pgfusepath{use as bounding box, clip}%
\begin{pgfscope}%
\pgfsetbuttcap%
\pgfsetmiterjoin%
\definecolor{currentfill}{rgb}{1.000000,1.000000,1.000000}%
\pgfsetfillcolor{currentfill}%
\pgfsetlinewidth{0.000000pt}%
\definecolor{currentstroke}{rgb}{1.000000,1.000000,1.000000}%
\pgfsetstrokecolor{currentstroke}%
\pgfsetdash{}{0pt}%
\pgfpathmoveto{\pgfqpoint{0.000000in}{0.000000in}}%
\pgfpathlineto{\pgfqpoint{6.487985in}{0.000000in}}%
\pgfpathlineto{\pgfqpoint{6.487985in}{2.069857in}}%
\pgfpathlineto{\pgfqpoint{0.000000in}{2.069857in}}%
\pgfpathlineto{\pgfqpoint{0.000000in}{0.000000in}}%
\pgfpathclose%
\pgfusepath{fill}%
\end{pgfscope}%
\begin{pgfscope}%
\pgfpathrectangle{\pgfqpoint{0.100000in}{0.100000in}}{\pgfqpoint{2.862985in}{1.832778in}}%
\pgfusepath{clip}%
\pgfsetbuttcap%
\pgfsetmiterjoin%
\definecolor{currentfill}{rgb}{0.084583,0.772703,0.653210}%
\pgfsetfillcolor{currentfill}%
\pgfsetlinewidth{0.000000pt}%
\definecolor{currentstroke}{rgb}{0.000000,0.000000,0.000000}%
\pgfsetstrokecolor{currentstroke}%
\pgfsetstrokeopacity{0.000000}%
\pgfsetdash{}{0pt}%
\pgfpathmoveto{\pgfqpoint{1.055609in}{0.529755in}}%
\pgfpathlineto{\pgfqpoint{1.047380in}{0.529889in}}%
\pgfpathlineto{\pgfqpoint{1.041780in}{0.536088in}}%
\pgfpathlineto{\pgfqpoint{1.037972in}{0.552550in}}%
\pgfpathlineto{\pgfqpoint{1.046822in}{0.554612in}}%
\pgfpathlineto{\pgfqpoint{1.055813in}{0.552394in}}%
\pgfpathlineto{\pgfqpoint{1.064963in}{0.545909in}}%
\pgfpathlineto{\pgfqpoint{1.064959in}{0.536940in}}%
\pgfpathlineto{\pgfqpoint{1.055609in}{0.529755in}}%
\pgfpathclose%
\pgfusepath{fill}%
\end{pgfscope}%
\begin{pgfscope}%
\pgfpathrectangle{\pgfqpoint{0.100000in}{0.100000in}}{\pgfqpoint{2.862985in}{1.832778in}}%
\pgfusepath{clip}%
\pgfsetbuttcap%
\pgfsetmiterjoin%
\definecolor{currentfill}{rgb}{0.084583,0.772703,0.653210}%
\pgfsetfillcolor{currentfill}%
\pgfsetlinewidth{0.000000pt}%
\definecolor{currentstroke}{rgb}{0.000000,0.000000,0.000000}%
\pgfsetstrokecolor{currentstroke}%
\pgfsetstrokeopacity{0.000000}%
\pgfsetdash{}{0pt}%
\pgfpathmoveto{\pgfqpoint{1.105694in}{0.432023in}}%
\pgfpathlineto{\pgfqpoint{1.096488in}{0.435528in}}%
\pgfpathlineto{\pgfqpoint{1.096446in}{0.442404in}}%
\pgfpathlineto{\pgfqpoint{1.085326in}{0.448219in}}%
\pgfpathlineto{\pgfqpoint{1.086680in}{0.462860in}}%
\pgfpathlineto{\pgfqpoint{1.089901in}{0.469769in}}%
\pgfpathlineto{\pgfqpoint{1.096701in}{0.463410in}}%
\pgfpathlineto{\pgfqpoint{1.107760in}{0.464472in}}%
\pgfpathlineto{\pgfqpoint{1.104967in}{0.444236in}}%
\pgfpathlineto{\pgfqpoint{1.105694in}{0.432023in}}%
\pgfpathclose%
\pgfusepath{fill}%
\end{pgfscope}%
\begin{pgfscope}%
\pgfpathrectangle{\pgfqpoint{0.100000in}{0.100000in}}{\pgfqpoint{2.862985in}{1.832778in}}%
\pgfusepath{clip}%
\pgfsetbuttcap%
\pgfsetmiterjoin%
\definecolor{currentfill}{rgb}{0.084583,0.772703,0.653210}%
\pgfsetfillcolor{currentfill}%
\pgfsetlinewidth{0.000000pt}%
\definecolor{currentstroke}{rgb}{0.000000,0.000000,0.000000}%
\pgfsetstrokecolor{currentstroke}%
\pgfsetstrokeopacity{0.000000}%
\pgfsetdash{}{0pt}%
\pgfpathmoveto{\pgfqpoint{1.144333in}{0.386887in}}%
\pgfpathlineto{\pgfqpoint{1.132651in}{0.388876in}}%
\pgfpathlineto{\pgfqpoint{1.126191in}{0.398877in}}%
\pgfpathlineto{\pgfqpoint{1.125279in}{0.407100in}}%
\pgfpathlineto{\pgfqpoint{1.144333in}{0.386887in}}%
\pgfpathclose%
\pgfusepath{fill}%
\end{pgfscope}%
\begin{pgfscope}%
\pgfpathrectangle{\pgfqpoint{0.100000in}{0.100000in}}{\pgfqpoint{2.862985in}{1.832778in}}%
\pgfusepath{clip}%
\pgfsetbuttcap%
\pgfsetmiterjoin%
\definecolor{currentfill}{rgb}{0.084583,0.772703,0.653210}%
\pgfsetfillcolor{currentfill}%
\pgfsetlinewidth{0.000000pt}%
\definecolor{currentstroke}{rgb}{0.000000,0.000000,0.000000}%
\pgfsetstrokecolor{currentstroke}%
\pgfsetstrokeopacity{0.000000}%
\pgfsetdash{}{0pt}%
\pgfpathmoveto{\pgfqpoint{1.120394in}{0.388631in}}%
\pgfpathlineto{\pgfqpoint{1.126519in}{0.383659in}}%
\pgfpathlineto{\pgfqpoint{1.127848in}{0.375110in}}%
\pgfpathlineto{\pgfqpoint{1.115519in}{0.375583in}}%
\pgfpathlineto{\pgfqpoint{1.120394in}{0.388631in}}%
\pgfpathclose%
\pgfusepath{fill}%
\end{pgfscope}%
\begin{pgfscope}%
\pgfpathrectangle{\pgfqpoint{0.100000in}{0.100000in}}{\pgfqpoint{2.862985in}{1.832778in}}%
\pgfusepath{clip}%
\pgfsetbuttcap%
\pgfsetmiterjoin%
\definecolor{currentfill}{rgb}{0.084583,0.772703,0.653210}%
\pgfsetfillcolor{currentfill}%
\pgfsetlinewidth{0.000000pt}%
\definecolor{currentstroke}{rgb}{0.000000,0.000000,0.000000}%
\pgfsetstrokecolor{currentstroke}%
\pgfsetstrokeopacity{0.000000}%
\pgfsetdash{}{0pt}%
\pgfpathmoveto{\pgfqpoint{1.148521in}{0.340880in}}%
\pgfpathlineto{\pgfqpoint{1.136217in}{0.345303in}}%
\pgfpathlineto{\pgfqpoint{1.139744in}{0.356026in}}%
\pgfpathlineto{\pgfqpoint{1.134586in}{0.366952in}}%
\pgfpathlineto{\pgfqpoint{1.136206in}{0.374463in}}%
\pgfpathlineto{\pgfqpoint{1.146107in}{0.377326in}}%
\pgfpathlineto{\pgfqpoint{1.145916in}{0.365292in}}%
\pgfpathlineto{\pgfqpoint{1.157627in}{0.358618in}}%
\pgfpathlineto{\pgfqpoint{1.164326in}{0.342480in}}%
\pgfpathlineto{\pgfqpoint{1.156861in}{0.336266in}}%
\pgfpathlineto{\pgfqpoint{1.148521in}{0.340880in}}%
\pgfpathclose%
\pgfusepath{fill}%
\end{pgfscope}%
\begin{pgfscope}%
\pgfpathrectangle{\pgfqpoint{0.100000in}{0.100000in}}{\pgfqpoint{2.862985in}{1.832778in}}%
\pgfusepath{clip}%
\pgfsetbuttcap%
\pgfsetmiterjoin%
\definecolor{currentfill}{rgb}{0.084583,0.772703,0.653210}%
\pgfsetfillcolor{currentfill}%
\pgfsetlinewidth{0.000000pt}%
\definecolor{currentstroke}{rgb}{0.000000,0.000000,0.000000}%
\pgfsetstrokecolor{currentstroke}%
\pgfsetstrokeopacity{0.000000}%
\pgfsetdash{}{0pt}%
\pgfpathmoveto{\pgfqpoint{1.104084in}{0.231768in}}%
\pgfpathlineto{\pgfqpoint{1.098222in}{0.243877in}}%
\pgfpathlineto{\pgfqpoint{1.100818in}{0.252022in}}%
\pgfpathlineto{\pgfqpoint{1.112315in}{0.263296in}}%
\pgfpathlineto{\pgfqpoint{1.113404in}{0.271938in}}%
\pgfpathlineto{\pgfqpoint{1.120928in}{0.291072in}}%
\pgfpathlineto{\pgfqpoint{1.140718in}{0.294341in}}%
\pgfpathlineto{\pgfqpoint{1.144506in}{0.305541in}}%
\pgfpathlineto{\pgfqpoint{1.153232in}{0.301341in}}%
\pgfpathlineto{\pgfqpoint{1.170961in}{0.269908in}}%
\pgfpathlineto{\pgfqpoint{1.171185in}{0.259645in}}%
\pgfpathlineto{\pgfqpoint{1.166173in}{0.250763in}}%
\pgfpathlineto{\pgfqpoint{1.172532in}{0.230369in}}%
\pgfpathlineto{\pgfqpoint{1.168089in}{0.227081in}}%
\pgfpathlineto{\pgfqpoint{1.152955in}{0.228037in}}%
\pgfpathlineto{\pgfqpoint{1.134756in}{0.234119in}}%
\pgfpathlineto{\pgfqpoint{1.119460in}{0.236631in}}%
\pgfpathlineto{\pgfqpoint{1.115375in}{0.233063in}}%
\pgfpathlineto{\pgfqpoint{1.104084in}{0.231768in}}%
\pgfpathclose%
\pgfusepath{fill}%
\end{pgfscope}%
\begin{pgfscope}%
\pgfpathrectangle{\pgfqpoint{0.100000in}{0.100000in}}{\pgfqpoint{2.862985in}{1.832778in}}%
\pgfusepath{clip}%
\pgfsetbuttcap%
\pgfsetmiterjoin%
\definecolor{currentfill}{rgb}{0.096886,0.776148,0.602768}%
\pgfsetfillcolor{currentfill}%
\pgfsetlinewidth{0.000000pt}%
\definecolor{currentstroke}{rgb}{0.000000,0.000000,0.000000}%
\pgfsetstrokecolor{currentstroke}%
\pgfsetstrokeopacity{0.000000}%
\pgfsetdash{}{0pt}%
\pgfpathmoveto{\pgfqpoint{0.780273in}{1.824988in}}%
\pgfpathlineto{\pgfqpoint{0.765426in}{1.766159in}}%
\pgfpathlineto{\pgfqpoint{0.754869in}{1.724007in}}%
\pgfpathlineto{\pgfqpoint{0.742317in}{1.673299in}}%
\pgfpathlineto{\pgfqpoint{0.739852in}{1.660801in}}%
\pgfpathlineto{\pgfqpoint{0.741462in}{1.649338in}}%
\pgfpathlineto{\pgfqpoint{0.739329in}{1.638779in}}%
\pgfpathlineto{\pgfqpoint{0.651709in}{1.661643in}}%
\pgfpathlineto{\pgfqpoint{0.643794in}{1.658987in}}%
\pgfpathlineto{\pgfqpoint{0.637021in}{1.661292in}}%
\pgfpathlineto{\pgfqpoint{0.605660in}{1.662045in}}%
\pgfpathlineto{\pgfqpoint{0.595053in}{1.659055in}}%
\pgfpathlineto{\pgfqpoint{0.584520in}{1.660065in}}%
\pgfpathlineto{\pgfqpoint{0.580055in}{1.664746in}}%
\pgfpathlineto{\pgfqpoint{0.551979in}{1.663725in}}%
\pgfpathlineto{\pgfqpoint{0.547545in}{1.671374in}}%
\pgfpathlineto{\pgfqpoint{0.539127in}{1.675404in}}%
\pgfpathlineto{\pgfqpoint{0.526880in}{1.677860in}}%
\pgfpathlineto{\pgfqpoint{0.505740in}{1.674264in}}%
\pgfpathlineto{\pgfqpoint{0.497937in}{1.677787in}}%
\pgfpathlineto{\pgfqpoint{0.485905in}{1.687156in}}%
\pgfpathlineto{\pgfqpoint{0.489522in}{1.705519in}}%
\pgfpathlineto{\pgfqpoint{0.488254in}{1.714897in}}%
\pgfpathlineto{\pgfqpoint{0.478718in}{1.724852in}}%
\pgfpathlineto{\pgfqpoint{0.472610in}{1.724234in}}%
\pgfpathlineto{\pgfqpoint{0.468229in}{1.734304in}}%
\pgfpathlineto{\pgfqpoint{0.457849in}{1.738374in}}%
\pgfpathlineto{\pgfqpoint{0.450305in}{1.737856in}}%
\pgfpathlineto{\pgfqpoint{0.447864in}{1.748198in}}%
\pgfpathlineto{\pgfqpoint{0.455380in}{1.747098in}}%
\pgfpathlineto{\pgfqpoint{0.454783in}{1.761472in}}%
\pgfpathlineto{\pgfqpoint{0.451507in}{1.770008in}}%
\pgfpathlineto{\pgfqpoint{0.452951in}{1.780932in}}%
\pgfpathlineto{\pgfqpoint{0.456884in}{1.789149in}}%
\pgfpathlineto{\pgfqpoint{0.456662in}{1.804741in}}%
\pgfpathlineto{\pgfqpoint{0.454563in}{1.810398in}}%
\pgfpathlineto{\pgfqpoint{0.458188in}{1.828548in}}%
\pgfpathlineto{\pgfqpoint{0.457142in}{1.840246in}}%
\pgfpathlineto{\pgfqpoint{0.453517in}{1.845854in}}%
\pgfpathlineto{\pgfqpoint{0.454070in}{1.864192in}}%
\pgfpathlineto{\pgfqpoint{0.459342in}{1.877678in}}%
\pgfpathlineto{\pgfqpoint{0.465077in}{1.874378in}}%
\pgfpathlineto{\pgfqpoint{0.484064in}{1.854764in}}%
\pgfpathlineto{\pgfqpoint{0.507071in}{1.844031in}}%
\pgfpathlineto{\pgfqpoint{0.518852in}{1.842761in}}%
\pgfpathlineto{\pgfqpoint{0.529951in}{1.838191in}}%
\pgfpathlineto{\pgfqpoint{0.533101in}{1.822814in}}%
\pgfpathlineto{\pgfqpoint{0.523375in}{1.819901in}}%
\pgfpathlineto{\pgfqpoint{0.516929in}{1.807782in}}%
\pgfpathlineto{\pgfqpoint{0.524497in}{1.808455in}}%
\pgfpathlineto{\pgfqpoint{0.527549in}{1.813895in}}%
\pgfpathlineto{\pgfqpoint{0.538221in}{1.820697in}}%
\pgfpathlineto{\pgfqpoint{0.537670in}{1.810596in}}%
\pgfpathlineto{\pgfqpoint{0.530544in}{1.808967in}}%
\pgfpathlineto{\pgfqpoint{0.531727in}{1.795983in}}%
\pgfpathlineto{\pgfqpoint{0.523050in}{1.782976in}}%
\pgfpathlineto{\pgfqpoint{0.517496in}{1.789161in}}%
\pgfpathlineto{\pgfqpoint{0.500991in}{1.785798in}}%
\pgfpathlineto{\pgfqpoint{0.500123in}{1.778232in}}%
\pgfpathlineto{\pgfqpoint{0.505813in}{1.773622in}}%
\pgfpathlineto{\pgfqpoint{0.514493in}{1.773169in}}%
\pgfpathlineto{\pgfqpoint{0.526495in}{1.783019in}}%
\pgfpathlineto{\pgfqpoint{0.535994in}{1.783837in}}%
\pgfpathlineto{\pgfqpoint{0.536358in}{1.794741in}}%
\pgfpathlineto{\pgfqpoint{0.541254in}{1.810756in}}%
\pgfpathlineto{\pgfqpoint{0.552856in}{1.822666in}}%
\pgfpathlineto{\pgfqpoint{0.548988in}{1.831891in}}%
\pgfpathlineto{\pgfqpoint{0.551630in}{1.841823in}}%
\pgfpathlineto{\pgfqpoint{0.546483in}{1.851722in}}%
\pgfpathlineto{\pgfqpoint{0.553818in}{1.864381in}}%
\pgfpathlineto{\pgfqpoint{0.554917in}{1.871996in}}%
\pgfpathlineto{\pgfqpoint{0.548548in}{1.877001in}}%
\pgfpathlineto{\pgfqpoint{0.549609in}{1.889804in}}%
\pgfpathlineto{\pgfqpoint{0.625918in}{1.866620in}}%
\pgfpathlineto{\pgfqpoint{0.706952in}{1.843896in}}%
\pgfpathlineto{\pgfqpoint{0.780273in}{1.824988in}}%
\pgfpathclose%
\pgfusepath{fill}%
\end{pgfscope}%
\begin{pgfscope}%
\pgfpathrectangle{\pgfqpoint{0.100000in}{0.100000in}}{\pgfqpoint{2.862985in}{1.832778in}}%
\pgfusepath{clip}%
\pgfsetbuttcap%
\pgfsetmiterjoin%
\definecolor{currentfill}{rgb}{0.096886,0.776148,0.602768}%
\pgfsetfillcolor{currentfill}%
\pgfsetlinewidth{0.000000pt}%
\definecolor{currentstroke}{rgb}{0.000000,0.000000,0.000000}%
\pgfsetstrokecolor{currentstroke}%
\pgfsetstrokeopacity{0.000000}%
\pgfsetdash{}{0pt}%
\pgfpathmoveto{\pgfqpoint{0.537056in}{1.845512in}}%
\pgfpathlineto{\pgfqpoint{0.540813in}{1.831026in}}%
\pgfpathlineto{\pgfqpoint{0.546706in}{1.826214in}}%
\pgfpathlineto{\pgfqpoint{0.542004in}{1.820089in}}%
\pgfpathlineto{\pgfqpoint{0.537421in}{1.829064in}}%
\pgfpathlineto{\pgfqpoint{0.537056in}{1.845512in}}%
\pgfpathclose%
\pgfusepath{fill}%
\end{pgfscope}%
\begin{pgfscope}%
\pgfpathrectangle{\pgfqpoint{0.100000in}{0.100000in}}{\pgfqpoint{2.862985in}{1.832778in}}%
\pgfusepath{clip}%
\pgfsetbuttcap%
\pgfsetmiterjoin%
\definecolor{currentfill}{rgb}{0.179931,0.799400,0.262284}%
\pgfsetfillcolor{currentfill}%
\pgfsetlinewidth{0.000000pt}%
\definecolor{currentstroke}{rgb}{0.000000,0.000000,0.000000}%
\pgfsetstrokecolor{currentstroke}%
\pgfsetstrokeopacity{0.000000}%
\pgfsetdash{}{0pt}%
\pgfpathmoveto{\pgfqpoint{0.819624in}{1.815480in}}%
\pgfpathlineto{\pgfqpoint{0.901194in}{1.797175in}}%
\pgfpathlineto{\pgfqpoint{0.978016in}{1.781647in}}%
\pgfpathlineto{\pgfqpoint{1.037121in}{1.770817in}}%
\pgfpathlineto{\pgfqpoint{1.088650in}{1.762160in}}%
\pgfpathlineto{\pgfqpoint{1.140293in}{1.754214in}}%
\pgfpathlineto{\pgfqpoint{1.184272in}{1.748018in}}%
\pgfpathlineto{\pgfqpoint{1.228320in}{1.742337in}}%
\pgfpathlineto{\pgfqpoint{1.272432in}{1.737171in}}%
\pgfpathlineto{\pgfqpoint{1.314001in}{1.732781in}}%
\pgfpathlineto{\pgfqpoint{1.308284in}{1.669428in}}%
\pgfpathlineto{\pgfqpoint{1.299775in}{1.583815in}}%
\pgfpathlineto{\pgfqpoint{1.295312in}{1.539783in}}%
\pgfpathlineto{\pgfqpoint{1.288893in}{1.480441in}}%
\pgfpathlineto{\pgfqpoint{1.243453in}{1.485403in}}%
\pgfpathlineto{\pgfqpoint{1.191429in}{1.491308in}}%
\pgfpathlineto{\pgfqpoint{1.119194in}{1.501195in}}%
\pgfpathlineto{\pgfqpoint{1.086931in}{1.505797in}}%
\pgfpathlineto{\pgfqpoint{1.007445in}{1.518274in}}%
\pgfpathlineto{\pgfqpoint{0.980083in}{1.523280in}}%
\pgfpathlineto{\pgfqpoint{0.974380in}{1.490854in}}%
\pgfpathlineto{\pgfqpoint{0.971262in}{1.493168in}}%
\pgfpathlineto{\pgfqpoint{0.965403in}{1.508767in}}%
\pgfpathlineto{\pgfqpoint{0.957871in}{1.507806in}}%
\pgfpathlineto{\pgfqpoint{0.952480in}{1.499920in}}%
\pgfpathlineto{\pgfqpoint{0.941293in}{1.498866in}}%
\pgfpathlineto{\pgfqpoint{0.939045in}{1.502638in}}%
\pgfpathlineto{\pgfqpoint{0.928566in}{1.502181in}}%
\pgfpathlineto{\pgfqpoint{0.923297in}{1.505835in}}%
\pgfpathlineto{\pgfqpoint{0.915931in}{1.500183in}}%
\pgfpathlineto{\pgfqpoint{0.898038in}{1.505269in}}%
\pgfpathlineto{\pgfqpoint{0.890230in}{1.502538in}}%
\pgfpathlineto{\pgfqpoint{0.886029in}{1.514658in}}%
\pgfpathlineto{\pgfqpoint{0.885833in}{1.526190in}}%
\pgfpathlineto{\pgfqpoint{0.873499in}{1.534484in}}%
\pgfpathlineto{\pgfqpoint{0.875510in}{1.546436in}}%
\pgfpathlineto{\pgfqpoint{0.866728in}{1.566180in}}%
\pgfpathlineto{\pgfqpoint{0.867319in}{1.582964in}}%
\pgfpathlineto{\pgfqpoint{0.859772in}{1.591178in}}%
\pgfpathlineto{\pgfqpoint{0.852690in}{1.583911in}}%
\pgfpathlineto{\pgfqpoint{0.843050in}{1.579980in}}%
\pgfpathlineto{\pgfqpoint{0.835712in}{1.588082in}}%
\pgfpathlineto{\pgfqpoint{0.836905in}{1.601019in}}%
\pgfpathlineto{\pgfqpoint{0.845705in}{1.606116in}}%
\pgfpathlineto{\pgfqpoint{0.843366in}{1.614312in}}%
\pgfpathlineto{\pgfqpoint{0.858663in}{1.653985in}}%
\pgfpathlineto{\pgfqpoint{0.846601in}{1.654979in}}%
\pgfpathlineto{\pgfqpoint{0.845361in}{1.662098in}}%
\pgfpathlineto{\pgfqpoint{0.836575in}{1.668238in}}%
\pgfpathlineto{\pgfqpoint{0.837137in}{1.675101in}}%
\pgfpathlineto{\pgfqpoint{0.832502in}{1.680433in}}%
\pgfpathlineto{\pgfqpoint{0.824957in}{1.699791in}}%
\pgfpathlineto{\pgfqpoint{0.818063in}{1.703739in}}%
\pgfpathlineto{\pgfqpoint{0.813126in}{1.720446in}}%
\pgfpathlineto{\pgfqpoint{0.814291in}{1.731551in}}%
\pgfpathlineto{\pgfqpoint{0.805080in}{1.752004in}}%
\pgfpathlineto{\pgfqpoint{0.819624in}{1.815480in}}%
\pgfpathclose%
\pgfusepath{fill}%
\end{pgfscope}%
\begin{pgfscope}%
\pgfpathrectangle{\pgfqpoint{0.100000in}{0.100000in}}{\pgfqpoint{2.862985in}{1.832778in}}%
\pgfusepath{clip}%
\pgfsetbuttcap%
\pgfsetmiterjoin%
\definecolor{currentfill}{rgb}{0.063053,0.766674,0.741484}%
\pgfsetfillcolor{currentfill}%
\pgfsetlinewidth{0.000000pt}%
\definecolor{currentstroke}{rgb}{0.000000,0.000000,0.000000}%
\pgfsetstrokecolor{currentstroke}%
\pgfsetstrokeopacity{0.000000}%
\pgfsetdash{}{0pt}%
\pgfpathmoveto{\pgfqpoint{2.788801in}{1.489012in}}%
\pgfpathlineto{\pgfqpoint{2.787168in}{1.495954in}}%
\pgfpathlineto{\pgfqpoint{2.778094in}{1.502035in}}%
\pgfpathlineto{\pgfqpoint{2.754117in}{1.580468in}}%
\pgfpathlineto{\pgfqpoint{2.741192in}{1.618318in}}%
\pgfpathlineto{\pgfqpoint{2.751583in}{1.629222in}}%
\pgfpathlineto{\pgfqpoint{2.761683in}{1.655908in}}%
\pgfpathlineto{\pgfqpoint{2.766720in}{1.664437in}}%
\pgfpathlineto{\pgfqpoint{2.763152in}{1.668018in}}%
\pgfpathlineto{\pgfqpoint{2.761957in}{1.691707in}}%
\pgfpathlineto{\pgfqpoint{2.766463in}{1.698969in}}%
\pgfpathlineto{\pgfqpoint{2.764488in}{1.715849in}}%
\pgfpathlineto{\pgfqpoint{2.782474in}{1.771228in}}%
\pgfpathlineto{\pgfqpoint{2.790578in}{1.771530in}}%
\pgfpathlineto{\pgfqpoint{2.793942in}{1.761638in}}%
\pgfpathlineto{\pgfqpoint{2.801144in}{1.758798in}}%
\pgfpathlineto{\pgfqpoint{2.814749in}{1.770433in}}%
\pgfpathlineto{\pgfqpoint{2.825417in}{1.777306in}}%
\pgfpathlineto{\pgfqpoint{2.848922in}{1.765204in}}%
\pgfpathlineto{\pgfqpoint{2.870157in}{1.698008in}}%
\pgfpathlineto{\pgfqpoint{2.874194in}{1.681476in}}%
\pgfpathlineto{\pgfqpoint{2.883495in}{1.679537in}}%
\pgfpathlineto{\pgfqpoint{2.896110in}{1.668309in}}%
\pgfpathlineto{\pgfqpoint{2.895396in}{1.661767in}}%
\pgfpathlineto{\pgfqpoint{2.903989in}{1.654013in}}%
\pgfpathlineto{\pgfqpoint{2.910970in}{1.658462in}}%
\pgfpathlineto{\pgfqpoint{2.925577in}{1.641426in}}%
\pgfpathlineto{\pgfqpoint{2.919049in}{1.627792in}}%
\pgfpathlineto{\pgfqpoint{2.910294in}{1.627522in}}%
\pgfpathlineto{\pgfqpoint{2.903271in}{1.615351in}}%
\pgfpathlineto{\pgfqpoint{2.894762in}{1.613625in}}%
\pgfpathlineto{\pgfqpoint{2.888516in}{1.607145in}}%
\pgfpathlineto{\pgfqpoint{2.869860in}{1.600195in}}%
\pgfpathlineto{\pgfqpoint{2.858568in}{1.588995in}}%
\pgfpathlineto{\pgfqpoint{2.852832in}{1.597034in}}%
\pgfpathlineto{\pgfqpoint{2.847620in}{1.591272in}}%
\pgfpathlineto{\pgfqpoint{2.849047in}{1.567979in}}%
\pgfpathlineto{\pgfqpoint{2.844920in}{1.558753in}}%
\pgfpathlineto{\pgfqpoint{2.835918in}{1.561273in}}%
\pgfpathlineto{\pgfqpoint{2.834465in}{1.551814in}}%
\pgfpathlineto{\pgfqpoint{2.830576in}{1.547924in}}%
\pgfpathlineto{\pgfqpoint{2.822791in}{1.548872in}}%
\pgfpathlineto{\pgfqpoint{2.823281in}{1.540028in}}%
\pgfpathlineto{\pgfqpoint{2.811482in}{1.542540in}}%
\pgfpathlineto{\pgfqpoint{2.805038in}{1.530287in}}%
\pgfpathlineto{\pgfqpoint{2.807472in}{1.523875in}}%
\pgfpathlineto{\pgfqpoint{2.803639in}{1.513212in}}%
\pgfpathlineto{\pgfqpoint{2.797604in}{1.505369in}}%
\pgfpathlineto{\pgfqpoint{2.796086in}{1.488995in}}%
\pgfpathlineto{\pgfqpoint{2.788801in}{1.489012in}}%
\pgfpathclose%
\pgfusepath{fill}%
\end{pgfscope}%
\begin{pgfscope}%
\pgfpathrectangle{\pgfqpoint{0.100000in}{0.100000in}}{\pgfqpoint{2.862985in}{1.832778in}}%
\pgfusepath{clip}%
\pgfsetbuttcap%
\pgfsetmiterjoin%
\definecolor{currentfill}{rgb}{0.063053,0.766674,0.741484}%
\pgfsetfillcolor{currentfill}%
\pgfsetlinewidth{0.000000pt}%
\definecolor{currentstroke}{rgb}{0.000000,0.000000,0.000000}%
\pgfsetstrokecolor{currentstroke}%
\pgfsetstrokeopacity{0.000000}%
\pgfsetdash{}{0pt}%
\pgfpathmoveto{\pgfqpoint{2.873211in}{1.595412in}}%
\pgfpathlineto{\pgfqpoint{2.878538in}{1.600991in}}%
\pgfpathlineto{\pgfqpoint{2.883609in}{1.595747in}}%
\pgfpathlineto{\pgfqpoint{2.874500in}{1.588798in}}%
\pgfpathlineto{\pgfqpoint{2.873211in}{1.595412in}}%
\pgfpathclose%
\pgfusepath{fill}%
\end{pgfscope}%
\begin{pgfscope}%
\pgfpathrectangle{\pgfqpoint{0.100000in}{0.100000in}}{\pgfqpoint{2.862985in}{1.832778in}}%
\pgfusepath{clip}%
\pgfsetbuttcap%
\pgfsetmiterjoin%
\definecolor{currentfill}{rgb}{0.104575,0.611765,0.104575}%
\pgfsetfillcolor{currentfill}%
\pgfsetlinewidth{0.000000pt}%
\definecolor{currentstroke}{rgb}{0.000000,0.000000,0.000000}%
\pgfsetstrokecolor{currentstroke}%
\pgfsetstrokeopacity{0.000000}%
\pgfsetdash{}{0pt}%
\pgfpathmoveto{\pgfqpoint{1.295312in}{1.539783in}}%
\pgfpathlineto{\pgfqpoint{1.299775in}{1.583815in}}%
\pgfpathlineto{\pgfqpoint{1.308284in}{1.669428in}}%
\pgfpathlineto{\pgfqpoint{1.314001in}{1.732781in}}%
\pgfpathlineto{\pgfqpoint{1.360821in}{1.728389in}}%
\pgfpathlineto{\pgfqpoint{1.420718in}{1.723620in}}%
\pgfpathlineto{\pgfqpoint{1.475466in}{1.720093in}}%
\pgfpathlineto{\pgfqpoint{1.525039in}{1.717585in}}%
\pgfpathlineto{\pgfqpoint{1.599019in}{1.715047in}}%
\pgfpathlineto{\pgfqpoint{1.604068in}{1.694197in}}%
\pgfpathlineto{\pgfqpoint{1.602261in}{1.684232in}}%
\pgfpathlineto{\pgfqpoint{1.601666in}{1.663751in}}%
\pgfpathlineto{\pgfqpoint{1.605101in}{1.648413in}}%
\pgfpathlineto{\pgfqpoint{1.612970in}{1.625784in}}%
\pgfpathlineto{\pgfqpoint{1.612952in}{1.597294in}}%
\pgfpathlineto{\pgfqpoint{1.614408in}{1.564197in}}%
\pgfpathlineto{\pgfqpoint{1.616408in}{1.555277in}}%
\pgfpathlineto{\pgfqpoint{1.622257in}{1.545485in}}%
\pgfpathlineto{\pgfqpoint{1.624186in}{1.530226in}}%
\pgfpathlineto{\pgfqpoint{1.623353in}{1.520039in}}%
\pgfpathlineto{\pgfqpoint{1.561333in}{1.521383in}}%
\pgfpathlineto{\pgfqpoint{1.516215in}{1.523670in}}%
\pgfpathlineto{\pgfqpoint{1.450074in}{1.527191in}}%
\pgfpathlineto{\pgfqpoint{1.384849in}{1.531824in}}%
\pgfpathlineto{\pgfqpoint{1.341407in}{1.535329in}}%
\pgfpathlineto{\pgfqpoint{1.295312in}{1.539783in}}%
\pgfpathclose%
\pgfusepath{fill}%
\end{pgfscope}%
\begin{pgfscope}%
\pgfpathrectangle{\pgfqpoint{0.100000in}{0.100000in}}{\pgfqpoint{2.862985in}{1.832778in}}%
\pgfusepath{clip}%
\pgfsetbuttcap%
\pgfsetmiterjoin%
\definecolor{currentfill}{rgb}{0.079969,0.560092,0.079969}%
\pgfsetfillcolor{currentfill}%
\pgfsetlinewidth{0.000000pt}%
\definecolor{currentstroke}{rgb}{0.000000,0.000000,0.000000}%
\pgfsetstrokecolor{currentstroke}%
\pgfsetstrokeopacity{0.000000}%
\pgfsetdash{}{0pt}%
\pgfpathmoveto{\pgfqpoint{1.276565in}{1.355506in}}%
\pgfpathlineto{\pgfqpoint{1.283725in}{1.429317in}}%
\pgfpathlineto{\pgfqpoint{1.288893in}{1.480441in}}%
\pgfpathlineto{\pgfqpoint{1.295312in}{1.539783in}}%
\pgfpathlineto{\pgfqpoint{1.341407in}{1.535329in}}%
\pgfpathlineto{\pgfqpoint{1.384849in}{1.531824in}}%
\pgfpathlineto{\pgfqpoint{1.450074in}{1.527191in}}%
\pgfpathlineto{\pgfqpoint{1.516215in}{1.523670in}}%
\pgfpathlineto{\pgfqpoint{1.561333in}{1.521383in}}%
\pgfpathlineto{\pgfqpoint{1.623353in}{1.520039in}}%
\pgfpathlineto{\pgfqpoint{1.619154in}{1.507786in}}%
\pgfpathlineto{\pgfqpoint{1.610772in}{1.498167in}}%
\pgfpathlineto{\pgfqpoint{1.617191in}{1.487092in}}%
\pgfpathlineto{\pgfqpoint{1.624268in}{1.484725in}}%
\pgfpathlineto{\pgfqpoint{1.627626in}{1.478363in}}%
\pgfpathlineto{\pgfqpoint{1.626855in}{1.431914in}}%
\pgfpathlineto{\pgfqpoint{1.625565in}{1.366544in}}%
\pgfpathlineto{\pgfqpoint{1.620791in}{1.349375in}}%
\pgfpathlineto{\pgfqpoint{1.625061in}{1.339878in}}%
\pgfpathlineto{\pgfqpoint{1.616489in}{1.319470in}}%
\pgfpathlineto{\pgfqpoint{1.625519in}{1.303018in}}%
\pgfpathlineto{\pgfqpoint{1.617847in}{1.304277in}}%
\pgfpathlineto{\pgfqpoint{1.612612in}{1.314513in}}%
\pgfpathlineto{\pgfqpoint{1.593930in}{1.321529in}}%
\pgfpathlineto{\pgfqpoint{1.582148in}{1.327701in}}%
\pgfpathlineto{\pgfqpoint{1.562377in}{1.328214in}}%
\pgfpathlineto{\pgfqpoint{1.555515in}{1.322591in}}%
\pgfpathlineto{\pgfqpoint{1.533137in}{1.333663in}}%
\pgfpathlineto{\pgfqpoint{1.531420in}{1.337164in}}%
\pgfpathlineto{\pgfqpoint{1.453322in}{1.340860in}}%
\pgfpathlineto{\pgfqpoint{1.405874in}{1.343580in}}%
\pgfpathlineto{\pgfqpoint{1.366682in}{1.346649in}}%
\pgfpathlineto{\pgfqpoint{1.301922in}{1.352768in}}%
\pgfpathlineto{\pgfqpoint{1.276565in}{1.355506in}}%
\pgfpathclose%
\pgfusepath{fill}%
\end{pgfscope}%
\begin{pgfscope}%
\pgfpathrectangle{\pgfqpoint{0.100000in}{0.100000in}}{\pgfqpoint{2.862985in}{1.832778in}}%
\pgfusepath{clip}%
\pgfsetbuttcap%
\pgfsetmiterjoin%
\definecolor{currentfill}{rgb}{0.192234,0.802845,0.211842}%
\pgfsetfillcolor{currentfill}%
\pgfsetlinewidth{0.000000pt}%
\definecolor{currentstroke}{rgb}{0.000000,0.000000,0.000000}%
\pgfsetstrokecolor{currentstroke}%
\pgfsetstrokeopacity{0.000000}%
\pgfsetdash{}{0pt}%
\pgfpathmoveto{\pgfqpoint{1.264283in}{1.230524in}}%
\pgfpathlineto{\pgfqpoint{1.222647in}{1.234326in}}%
\pgfpathlineto{\pgfqpoint{1.131892in}{1.245527in}}%
\pgfpathlineto{\pgfqpoint{1.082522in}{1.252613in}}%
\pgfpathlineto{\pgfqpoint{1.029571in}{1.260214in}}%
\pgfpathlineto{\pgfqpoint{0.984967in}{1.267352in}}%
\pgfpathlineto{\pgfqpoint{0.936009in}{1.275734in}}%
\pgfpathlineto{\pgfqpoint{0.947141in}{1.337483in}}%
\pgfpathlineto{\pgfqpoint{0.958416in}{1.400800in}}%
\pgfpathlineto{\pgfqpoint{0.974380in}{1.490854in}}%
\pgfpathlineto{\pgfqpoint{0.980083in}{1.523280in}}%
\pgfpathlineto{\pgfqpoint{1.007445in}{1.518274in}}%
\pgfpathlineto{\pgfqpoint{1.086931in}{1.505797in}}%
\pgfpathlineto{\pgfqpoint{1.119194in}{1.501195in}}%
\pgfpathlineto{\pgfqpoint{1.191429in}{1.491308in}}%
\pgfpathlineto{\pgfqpoint{1.243453in}{1.485403in}}%
\pgfpathlineto{\pgfqpoint{1.288893in}{1.480441in}}%
\pgfpathlineto{\pgfqpoint{1.283725in}{1.429317in}}%
\pgfpathlineto{\pgfqpoint{1.276565in}{1.355506in}}%
\pgfpathlineto{\pgfqpoint{1.270419in}{1.292777in}}%
\pgfpathlineto{\pgfqpoint{1.264283in}{1.230524in}}%
\pgfpathclose%
\pgfusepath{fill}%
\end{pgfscope}%
\begin{pgfscope}%
\pgfpathrectangle{\pgfqpoint{0.100000in}{0.100000in}}{\pgfqpoint{2.862985in}{1.832778in}}%
\pgfusepath{clip}%
\pgfsetbuttcap%
\pgfsetmiterjoin%
\definecolor{currentfill}{rgb}{0.150711,0.708651,0.150711}%
\pgfsetfillcolor{currentfill}%
\pgfsetlinewidth{0.000000pt}%
\definecolor{currentstroke}{rgb}{0.000000,0.000000,0.000000}%
\pgfsetstrokecolor{currentstroke}%
\pgfsetstrokeopacity{0.000000}%
\pgfsetdash{}{0pt}%
\pgfpathmoveto{\pgfqpoint{2.025520in}{1.315700in}}%
\pgfpathlineto{\pgfqpoint{1.972758in}{1.311952in}}%
\pgfpathlineto{\pgfqpoint{1.894114in}{1.308600in}}%
\pgfpathlineto{\pgfqpoint{1.891125in}{1.316544in}}%
\pgfpathlineto{\pgfqpoint{1.873686in}{1.322483in}}%
\pgfpathlineto{\pgfqpoint{1.869839in}{1.333702in}}%
\pgfpathlineto{\pgfqpoint{1.868228in}{1.347582in}}%
\pgfpathlineto{\pgfqpoint{1.872158in}{1.354693in}}%
\pgfpathlineto{\pgfqpoint{1.865951in}{1.361517in}}%
\pgfpathlineto{\pgfqpoint{1.864454in}{1.369657in}}%
\pgfpathlineto{\pgfqpoint{1.862453in}{1.387663in}}%
\pgfpathlineto{\pgfqpoint{1.856505in}{1.397444in}}%
\pgfpathlineto{\pgfqpoint{1.845940in}{1.402933in}}%
\pgfpathlineto{\pgfqpoint{1.834444in}{1.411992in}}%
\pgfpathlineto{\pgfqpoint{1.828501in}{1.422710in}}%
\pgfpathlineto{\pgfqpoint{1.817838in}{1.427026in}}%
\pgfpathlineto{\pgfqpoint{1.811570in}{1.434051in}}%
\pgfpathlineto{\pgfqpoint{1.803976in}{1.435243in}}%
\pgfpathlineto{\pgfqpoint{1.790419in}{1.445673in}}%
\pgfpathlineto{\pgfqpoint{1.792621in}{1.457675in}}%
\pgfpathlineto{\pgfqpoint{1.792197in}{1.480500in}}%
\pgfpathlineto{\pgfqpoint{1.796378in}{1.486783in}}%
\pgfpathlineto{\pgfqpoint{1.792621in}{1.496270in}}%
\pgfpathlineto{\pgfqpoint{1.786004in}{1.498105in}}%
\pgfpathlineto{\pgfqpoint{1.786550in}{1.506437in}}%
\pgfpathlineto{\pgfqpoint{1.794761in}{1.519600in}}%
\pgfpathlineto{\pgfqpoint{1.811020in}{1.530009in}}%
\pgfpathlineto{\pgfqpoint{1.810007in}{1.567025in}}%
\pgfpathlineto{\pgfqpoint{1.818143in}{1.572586in}}%
\pgfpathlineto{\pgfqpoint{1.825842in}{1.568882in}}%
\pgfpathlineto{\pgfqpoint{1.841526in}{1.574304in}}%
\pgfpathlineto{\pgfqpoint{1.871034in}{1.587932in}}%
\pgfpathlineto{\pgfqpoint{1.874878in}{1.583713in}}%
\pgfpathlineto{\pgfqpoint{1.869289in}{1.564600in}}%
\pgfpathlineto{\pgfqpoint{1.877604in}{1.568794in}}%
\pgfpathlineto{\pgfqpoint{1.891845in}{1.564606in}}%
\pgfpathlineto{\pgfqpoint{1.900596in}{1.561110in}}%
\pgfpathlineto{\pgfqpoint{1.905502in}{1.550866in}}%
\pgfpathlineto{\pgfqpoint{1.950397in}{1.541198in}}%
\pgfpathlineto{\pgfqpoint{1.963816in}{1.534565in}}%
\pgfpathlineto{\pgfqpoint{1.977460in}{1.534643in}}%
\pgfpathlineto{\pgfqpoint{1.991480in}{1.531907in}}%
\pgfpathlineto{\pgfqpoint{2.000578in}{1.522561in}}%
\pgfpathlineto{\pgfqpoint{2.009272in}{1.517943in}}%
\pgfpathlineto{\pgfqpoint{2.010870in}{1.504623in}}%
\pgfpathlineto{\pgfqpoint{2.008301in}{1.496259in}}%
\pgfpathlineto{\pgfqpoint{2.017991in}{1.495641in}}%
\pgfpathlineto{\pgfqpoint{2.014727in}{1.485941in}}%
\pgfpathlineto{\pgfqpoint{2.017832in}{1.482495in}}%
\pgfpathlineto{\pgfqpoint{2.020919in}{1.473339in}}%
\pgfpathlineto{\pgfqpoint{2.011476in}{1.468491in}}%
\pgfpathlineto{\pgfqpoint{2.005993in}{1.454984in}}%
\pgfpathlineto{\pgfqpoint{2.004273in}{1.445433in}}%
\pgfpathlineto{\pgfqpoint{2.009533in}{1.443798in}}%
\pgfpathlineto{\pgfqpoint{2.016268in}{1.450960in}}%
\pgfpathlineto{\pgfqpoint{2.021994in}{1.463372in}}%
\pgfpathlineto{\pgfqpoint{2.029741in}{1.467683in}}%
\pgfpathlineto{\pgfqpoint{2.035563in}{1.462075in}}%
\pgfpathlineto{\pgfqpoint{2.029832in}{1.445104in}}%
\pgfpathlineto{\pgfqpoint{2.028034in}{1.431928in}}%
\pgfpathlineto{\pgfqpoint{2.029728in}{1.422457in}}%
\pgfpathlineto{\pgfqpoint{2.024406in}{1.417099in}}%
\pgfpathlineto{\pgfqpoint{2.021748in}{1.404007in}}%
\pgfpathlineto{\pgfqpoint{2.023920in}{1.390247in}}%
\pgfpathlineto{\pgfqpoint{2.017640in}{1.369897in}}%
\pgfpathlineto{\pgfqpoint{2.017772in}{1.359728in}}%
\pgfpathlineto{\pgfqpoint{2.022734in}{1.337693in}}%
\pgfpathlineto{\pgfqpoint{2.025951in}{1.333913in}}%
\pgfpathlineto{\pgfqpoint{2.025520in}{1.315700in}}%
\pgfpathclose%
\pgfusepath{fill}%
\end{pgfscope}%
\begin{pgfscope}%
\pgfpathrectangle{\pgfqpoint{0.100000in}{0.100000in}}{\pgfqpoint{2.862985in}{1.832778in}}%
\pgfusepath{clip}%
\pgfsetbuttcap%
\pgfsetmiterjoin%
\definecolor{currentfill}{rgb}{0.150711,0.708651,0.150711}%
\pgfsetfillcolor{currentfill}%
\pgfsetlinewidth{0.000000pt}%
\definecolor{currentstroke}{rgb}{0.000000,0.000000,0.000000}%
\pgfsetstrokecolor{currentstroke}%
\pgfsetstrokeopacity{0.000000}%
\pgfsetdash{}{0pt}%
\pgfpathmoveto{\pgfqpoint{2.045302in}{1.494291in}}%
\pgfpathlineto{\pgfqpoint{2.046480in}{1.485502in}}%
\pgfpathlineto{\pgfqpoint{2.035690in}{1.462343in}}%
\pgfpathlineto{\pgfqpoint{2.030896in}{1.469051in}}%
\pgfpathlineto{\pgfqpoint{2.045302in}{1.494291in}}%
\pgfpathclose%
\pgfusepath{fill}%
\end{pgfscope}%
\begin{pgfscope}%
\pgfpathrectangle{\pgfqpoint{0.100000in}{0.100000in}}{\pgfqpoint{2.862985in}{1.832778in}}%
\pgfusepath{clip}%
\pgfsetbuttcap%
\pgfsetmiterjoin%
\definecolor{currentfill}{rgb}{0.127643,0.784760,0.476663}%
\pgfsetfillcolor{currentfill}%
\pgfsetlinewidth{0.000000pt}%
\definecolor{currentstroke}{rgb}{0.000000,0.000000,0.000000}%
\pgfsetstrokecolor{currentstroke}%
\pgfsetstrokeopacity{0.000000}%
\pgfsetdash{}{0pt}%
\pgfpathmoveto{\pgfqpoint{0.739329in}{1.638779in}}%
\pgfpathlineto{\pgfqpoint{0.741462in}{1.649338in}}%
\pgfpathlineto{\pgfqpoint{0.739852in}{1.660801in}}%
\pgfpathlineto{\pgfqpoint{0.742317in}{1.673299in}}%
\pgfpathlineto{\pgfqpoint{0.754869in}{1.724007in}}%
\pgfpathlineto{\pgfqpoint{0.765426in}{1.766159in}}%
\pgfpathlineto{\pgfqpoint{0.780273in}{1.824988in}}%
\pgfpathlineto{\pgfqpoint{0.819624in}{1.815480in}}%
\pgfpathlineto{\pgfqpoint{0.805080in}{1.752004in}}%
\pgfpathlineto{\pgfqpoint{0.814291in}{1.731551in}}%
\pgfpathlineto{\pgfqpoint{0.813126in}{1.720446in}}%
\pgfpathlineto{\pgfqpoint{0.818063in}{1.703739in}}%
\pgfpathlineto{\pgfqpoint{0.824957in}{1.699791in}}%
\pgfpathlineto{\pgfqpoint{0.832502in}{1.680433in}}%
\pgfpathlineto{\pgfqpoint{0.837137in}{1.675101in}}%
\pgfpathlineto{\pgfqpoint{0.836575in}{1.668238in}}%
\pgfpathlineto{\pgfqpoint{0.845361in}{1.662098in}}%
\pgfpathlineto{\pgfqpoint{0.846601in}{1.654979in}}%
\pgfpathlineto{\pgfqpoint{0.858663in}{1.653985in}}%
\pgfpathlineto{\pgfqpoint{0.843366in}{1.614312in}}%
\pgfpathlineto{\pgfqpoint{0.845705in}{1.606116in}}%
\pgfpathlineto{\pgfqpoint{0.836905in}{1.601019in}}%
\pgfpathlineto{\pgfqpoint{0.835712in}{1.588082in}}%
\pgfpathlineto{\pgfqpoint{0.843050in}{1.579980in}}%
\pgfpathlineto{\pgfqpoint{0.852690in}{1.583911in}}%
\pgfpathlineto{\pgfqpoint{0.859772in}{1.591178in}}%
\pgfpathlineto{\pgfqpoint{0.867319in}{1.582964in}}%
\pgfpathlineto{\pgfqpoint{0.866728in}{1.566180in}}%
\pgfpathlineto{\pgfqpoint{0.875510in}{1.546436in}}%
\pgfpathlineto{\pgfqpoint{0.873499in}{1.534484in}}%
\pgfpathlineto{\pgfqpoint{0.885833in}{1.526190in}}%
\pgfpathlineto{\pgfqpoint{0.886029in}{1.514658in}}%
\pgfpathlineto{\pgfqpoint{0.890230in}{1.502538in}}%
\pgfpathlineto{\pgfqpoint{0.898038in}{1.505269in}}%
\pgfpathlineto{\pgfqpoint{0.915931in}{1.500183in}}%
\pgfpathlineto{\pgfqpoint{0.923297in}{1.505835in}}%
\pgfpathlineto{\pgfqpoint{0.928566in}{1.502181in}}%
\pgfpathlineto{\pgfqpoint{0.939045in}{1.502638in}}%
\pgfpathlineto{\pgfqpoint{0.941293in}{1.498866in}}%
\pgfpathlineto{\pgfqpoint{0.952480in}{1.499920in}}%
\pgfpathlineto{\pgfqpoint{0.957871in}{1.507806in}}%
\pgfpathlineto{\pgfqpoint{0.965403in}{1.508767in}}%
\pgfpathlineto{\pgfqpoint{0.971262in}{1.493168in}}%
\pgfpathlineto{\pgfqpoint{0.974380in}{1.490854in}}%
\pgfpathlineto{\pgfqpoint{0.958416in}{1.400800in}}%
\pgfpathlineto{\pgfqpoint{0.947141in}{1.337483in}}%
\pgfpathlineto{\pgfqpoint{0.858212in}{1.354650in}}%
\pgfpathlineto{\pgfqpoint{0.810177in}{1.364205in}}%
\pgfpathlineto{\pgfqpoint{0.765246in}{1.374085in}}%
\pgfpathlineto{\pgfqpoint{0.723250in}{1.383593in}}%
\pgfpathlineto{\pgfqpoint{0.674647in}{1.395314in}}%
\pgfpathlineto{\pgfqpoint{0.699717in}{1.498158in}}%
\pgfpathlineto{\pgfqpoint{0.701069in}{1.505674in}}%
\pgfpathlineto{\pgfqpoint{0.712217in}{1.525368in}}%
\pgfpathlineto{\pgfqpoint{0.700673in}{1.537205in}}%
\pgfpathlineto{\pgfqpoint{0.703070in}{1.548828in}}%
\pgfpathlineto{\pgfqpoint{0.707464in}{1.551738in}}%
\pgfpathlineto{\pgfqpoint{0.715313in}{1.563715in}}%
\pgfpathlineto{\pgfqpoint{0.726752in}{1.572002in}}%
\pgfpathlineto{\pgfqpoint{0.727377in}{1.578133in}}%
\pgfpathlineto{\pgfqpoint{0.734282in}{1.584259in}}%
\pgfpathlineto{\pgfqpoint{0.740015in}{1.595725in}}%
\pgfpathlineto{\pgfqpoint{0.751774in}{1.607856in}}%
\pgfpathlineto{\pgfqpoint{0.750939in}{1.619838in}}%
\pgfpathlineto{\pgfqpoint{0.742899in}{1.626494in}}%
\pgfpathlineto{\pgfqpoint{0.739329in}{1.638779in}}%
\pgfpathclose%
\pgfusepath{fill}%
\end{pgfscope}%
\begin{pgfscope}%
\pgfpathrectangle{\pgfqpoint{0.100000in}{0.100000in}}{\pgfqpoint{2.862985in}{1.832778in}}%
\pgfusepath{clip}%
\pgfsetbuttcap%
\pgfsetmiterjoin%
\definecolor{currentfill}{rgb}{0.193772,0.799077,0.193772}%
\pgfsetfillcolor{currentfill}%
\pgfsetlinewidth{0.000000pt}%
\definecolor{currentstroke}{rgb}{0.000000,0.000000,0.000000}%
\pgfsetstrokecolor{currentstroke}%
\pgfsetstrokeopacity{0.000000}%
\pgfsetdash{}{0pt}%
\pgfpathmoveto{\pgfqpoint{2.686015in}{1.438080in}}%
\pgfpathlineto{\pgfqpoint{2.682320in}{1.449744in}}%
\pgfpathlineto{\pgfqpoint{2.675463in}{1.485127in}}%
\pgfpathlineto{\pgfqpoint{2.666598in}{1.496013in}}%
\pgfpathlineto{\pgfqpoint{2.658822in}{1.515590in}}%
\pgfpathlineto{\pgfqpoint{2.661375in}{1.530104in}}%
\pgfpathlineto{\pgfqpoint{2.659327in}{1.540801in}}%
\pgfpathlineto{\pgfqpoint{2.652683in}{1.551308in}}%
\pgfpathlineto{\pgfqpoint{2.652410in}{1.562882in}}%
\pgfpathlineto{\pgfqpoint{2.648552in}{1.575365in}}%
\pgfpathlineto{\pgfqpoint{2.683077in}{1.583862in}}%
\pgfpathlineto{\pgfqpoint{2.727922in}{1.595966in}}%
\pgfpathlineto{\pgfqpoint{2.729711in}{1.589025in}}%
\pgfpathlineto{\pgfqpoint{2.726856in}{1.577975in}}%
\pgfpathlineto{\pgfqpoint{2.733559in}{1.569133in}}%
\pgfpathlineto{\pgfqpoint{2.729999in}{1.557930in}}%
\pgfpathlineto{\pgfqpoint{2.715852in}{1.543846in}}%
\pgfpathlineto{\pgfqpoint{2.719741in}{1.533267in}}%
\pgfpathlineto{\pgfqpoint{2.716859in}{1.510906in}}%
\pgfpathlineto{\pgfqpoint{2.712824in}{1.498298in}}%
\pgfpathlineto{\pgfqpoint{2.717972in}{1.462444in}}%
\pgfpathlineto{\pgfqpoint{2.717278in}{1.449830in}}%
\pgfpathlineto{\pgfqpoint{2.722258in}{1.445782in}}%
\pgfpathlineto{\pgfqpoint{2.686015in}{1.438080in}}%
\pgfpathclose%
\pgfusepath{fill}%
\end{pgfscope}%
\begin{pgfscope}%
\pgfpathrectangle{\pgfqpoint{0.100000in}{0.100000in}}{\pgfqpoint{2.862985in}{1.832778in}}%
\pgfusepath{clip}%
\pgfsetbuttcap%
\pgfsetmiterjoin%
\definecolor{currentfill}{rgb}{0.079969,0.560092,0.079969}%
\pgfsetfillcolor{currentfill}%
\pgfsetlinewidth{0.000000pt}%
\definecolor{currentstroke}{rgb}{0.000000,0.000000,0.000000}%
\pgfsetstrokecolor{currentstroke}%
\pgfsetstrokeopacity{0.000000}%
\pgfsetdash{}{0pt}%
\pgfpathmoveto{\pgfqpoint{1.625565in}{1.366544in}}%
\pgfpathlineto{\pgfqpoint{1.626855in}{1.431914in}}%
\pgfpathlineto{\pgfqpoint{1.627626in}{1.478363in}}%
\pgfpathlineto{\pgfqpoint{1.624268in}{1.484725in}}%
\pgfpathlineto{\pgfqpoint{1.617191in}{1.487092in}}%
\pgfpathlineto{\pgfqpoint{1.610772in}{1.498167in}}%
\pgfpathlineto{\pgfqpoint{1.619154in}{1.507786in}}%
\pgfpathlineto{\pgfqpoint{1.623353in}{1.520039in}}%
\pgfpathlineto{\pgfqpoint{1.624186in}{1.530226in}}%
\pgfpathlineto{\pgfqpoint{1.622257in}{1.545485in}}%
\pgfpathlineto{\pgfqpoint{1.616408in}{1.555277in}}%
\pgfpathlineto{\pgfqpoint{1.614408in}{1.564197in}}%
\pgfpathlineto{\pgfqpoint{1.612952in}{1.597294in}}%
\pgfpathlineto{\pgfqpoint{1.612970in}{1.625784in}}%
\pgfpathlineto{\pgfqpoint{1.605101in}{1.648413in}}%
\pgfpathlineto{\pgfqpoint{1.601666in}{1.663751in}}%
\pgfpathlineto{\pgfqpoint{1.602261in}{1.684232in}}%
\pgfpathlineto{\pgfqpoint{1.604068in}{1.694197in}}%
\pgfpathlineto{\pgfqpoint{1.599019in}{1.715047in}}%
\pgfpathlineto{\pgfqpoint{1.633393in}{1.714359in}}%
\pgfpathlineto{\pgfqpoint{1.685604in}{1.713911in}}%
\pgfpathlineto{\pgfqpoint{1.685890in}{1.737609in}}%
\pgfpathlineto{\pgfqpoint{1.699180in}{1.735000in}}%
\pgfpathlineto{\pgfqpoint{1.705549in}{1.706104in}}%
\pgfpathlineto{\pgfqpoint{1.710245in}{1.695714in}}%
\pgfpathlineto{\pgfqpoint{1.721922in}{1.695407in}}%
\pgfpathlineto{\pgfqpoint{1.724536in}{1.691882in}}%
\pgfpathlineto{\pgfqpoint{1.740825in}{1.690320in}}%
\pgfpathlineto{\pgfqpoint{1.743562in}{1.683156in}}%
\pgfpathlineto{\pgfqpoint{1.754788in}{1.684767in}}%
\pgfpathlineto{\pgfqpoint{1.763507in}{1.691467in}}%
\pgfpathlineto{\pgfqpoint{1.778542in}{1.691216in}}%
\pgfpathlineto{\pgfqpoint{1.787845in}{1.685828in}}%
\pgfpathlineto{\pgfqpoint{1.788915in}{1.680774in}}%
\pgfpathlineto{\pgfqpoint{1.797767in}{1.679717in}}%
\pgfpathlineto{\pgfqpoint{1.803544in}{1.665932in}}%
\pgfpathlineto{\pgfqpoint{1.807272in}{1.674408in}}%
\pgfpathlineto{\pgfqpoint{1.817441in}{1.674929in}}%
\pgfpathlineto{\pgfqpoint{1.820013in}{1.668328in}}%
\pgfpathlineto{\pgfqpoint{1.831455in}{1.665315in}}%
\pgfpathlineto{\pgfqpoint{1.837765in}{1.655808in}}%
\pgfpathlineto{\pgfqpoint{1.851759in}{1.658760in}}%
\pgfpathlineto{\pgfqpoint{1.867158in}{1.670427in}}%
\pgfpathlineto{\pgfqpoint{1.872772in}{1.660142in}}%
\pgfpathlineto{\pgfqpoint{1.898039in}{1.662957in}}%
\pgfpathlineto{\pgfqpoint{1.908840in}{1.655898in}}%
\pgfpathlineto{\pgfqpoint{1.915126in}{1.658418in}}%
\pgfpathlineto{\pgfqpoint{1.920167in}{1.654444in}}%
\pgfpathlineto{\pgfqpoint{1.905213in}{1.645016in}}%
\pgfpathlineto{\pgfqpoint{1.883866in}{1.636615in}}%
\pgfpathlineto{\pgfqpoint{1.862725in}{1.619824in}}%
\pgfpathlineto{\pgfqpoint{1.844396in}{1.597742in}}%
\pgfpathlineto{\pgfqpoint{1.830521in}{1.584687in}}%
\pgfpathlineto{\pgfqpoint{1.818367in}{1.575713in}}%
\pgfpathlineto{\pgfqpoint{1.810007in}{1.567025in}}%
\pgfpathlineto{\pgfqpoint{1.811020in}{1.530009in}}%
\pgfpathlineto{\pgfqpoint{1.794761in}{1.519600in}}%
\pgfpathlineto{\pgfqpoint{1.786550in}{1.506437in}}%
\pgfpathlineto{\pgfqpoint{1.786004in}{1.498105in}}%
\pgfpathlineto{\pgfqpoint{1.792621in}{1.496270in}}%
\pgfpathlineto{\pgfqpoint{1.796378in}{1.486783in}}%
\pgfpathlineto{\pgfqpoint{1.792197in}{1.480500in}}%
\pgfpathlineto{\pgfqpoint{1.792621in}{1.457675in}}%
\pgfpathlineto{\pgfqpoint{1.790419in}{1.445673in}}%
\pgfpathlineto{\pgfqpoint{1.803976in}{1.435243in}}%
\pgfpathlineto{\pgfqpoint{1.811570in}{1.434051in}}%
\pgfpathlineto{\pgfqpoint{1.817838in}{1.427026in}}%
\pgfpathlineto{\pgfqpoint{1.828501in}{1.422710in}}%
\pgfpathlineto{\pgfqpoint{1.834444in}{1.411992in}}%
\pgfpathlineto{\pgfqpoint{1.845940in}{1.402933in}}%
\pgfpathlineto{\pgfqpoint{1.856505in}{1.397444in}}%
\pgfpathlineto{\pgfqpoint{1.862453in}{1.387663in}}%
\pgfpathlineto{\pgfqpoint{1.864454in}{1.369657in}}%
\pgfpathlineto{\pgfqpoint{1.808388in}{1.367621in}}%
\pgfpathlineto{\pgfqpoint{1.760597in}{1.366622in}}%
\pgfpathlineto{\pgfqpoint{1.717055in}{1.365982in}}%
\pgfpathlineto{\pgfqpoint{1.670992in}{1.366052in}}%
\pgfpathlineto{\pgfqpoint{1.625565in}{1.366544in}}%
\pgfpathclose%
\pgfusepath{fill}%
\end{pgfscope}%
\begin{pgfscope}%
\pgfpathrectangle{\pgfqpoint{0.100000in}{0.100000in}}{\pgfqpoint{2.862985in}{1.832778in}}%
\pgfusepath{clip}%
\pgfsetbuttcap%
\pgfsetmiterjoin%
\definecolor{currentfill}{rgb}{0.183007,0.800261,0.249673}%
\pgfsetfillcolor{currentfill}%
\pgfsetlinewidth{0.000000pt}%
\definecolor{currentstroke}{rgb}{0.000000,0.000000,0.000000}%
\pgfsetstrokecolor{currentstroke}%
\pgfsetstrokeopacity{0.000000}%
\pgfsetdash{}{0pt}%
\pgfpathmoveto{\pgfqpoint{0.352592in}{1.488751in}}%
\pgfpathlineto{\pgfqpoint{0.348152in}{1.496921in}}%
\pgfpathlineto{\pgfqpoint{0.350987in}{1.517893in}}%
\pgfpathlineto{\pgfqpoint{0.356459in}{1.528853in}}%
\pgfpathlineto{\pgfqpoint{0.353724in}{1.543679in}}%
\pgfpathlineto{\pgfqpoint{0.359432in}{1.549914in}}%
\pgfpathlineto{\pgfqpoint{0.369847in}{1.566720in}}%
\pgfpathlineto{\pgfqpoint{0.378677in}{1.576864in}}%
\pgfpathlineto{\pgfqpoint{0.391585in}{1.598959in}}%
\pgfpathlineto{\pgfqpoint{0.401485in}{1.623028in}}%
\pgfpathlineto{\pgfqpoint{0.411995in}{1.645601in}}%
\pgfpathlineto{\pgfqpoint{0.414119in}{1.655029in}}%
\pgfpathlineto{\pgfqpoint{0.428617in}{1.682000in}}%
\pgfpathlineto{\pgfqpoint{0.431409in}{1.693845in}}%
\pgfpathlineto{\pgfqpoint{0.437582in}{1.706284in}}%
\pgfpathlineto{\pgfqpoint{0.439788in}{1.721458in}}%
\pgfpathlineto{\pgfqpoint{0.451573in}{1.728857in}}%
\pgfpathlineto{\pgfqpoint{0.465546in}{1.732587in}}%
\pgfpathlineto{\pgfqpoint{0.472610in}{1.724234in}}%
\pgfpathlineto{\pgfqpoint{0.478718in}{1.724852in}}%
\pgfpathlineto{\pgfqpoint{0.488254in}{1.714897in}}%
\pgfpathlineto{\pgfqpoint{0.489522in}{1.705519in}}%
\pgfpathlineto{\pgfqpoint{0.485905in}{1.687156in}}%
\pgfpathlineto{\pgfqpoint{0.497937in}{1.677787in}}%
\pgfpathlineto{\pgfqpoint{0.505740in}{1.674264in}}%
\pgfpathlineto{\pgfqpoint{0.526880in}{1.677860in}}%
\pgfpathlineto{\pgfqpoint{0.539127in}{1.675404in}}%
\pgfpathlineto{\pgfqpoint{0.547545in}{1.671374in}}%
\pgfpathlineto{\pgfqpoint{0.551979in}{1.663725in}}%
\pgfpathlineto{\pgfqpoint{0.580055in}{1.664746in}}%
\pgfpathlineto{\pgfqpoint{0.584520in}{1.660065in}}%
\pgfpathlineto{\pgfqpoint{0.595053in}{1.659055in}}%
\pgfpathlineto{\pgfqpoint{0.605660in}{1.662045in}}%
\pgfpathlineto{\pgfqpoint{0.637021in}{1.661292in}}%
\pgfpathlineto{\pgfqpoint{0.643794in}{1.658987in}}%
\pgfpathlineto{\pgfqpoint{0.651709in}{1.661643in}}%
\pgfpathlineto{\pgfqpoint{0.739329in}{1.638779in}}%
\pgfpathlineto{\pgfqpoint{0.742899in}{1.626494in}}%
\pgfpathlineto{\pgfqpoint{0.750939in}{1.619838in}}%
\pgfpathlineto{\pgfqpoint{0.751774in}{1.607856in}}%
\pgfpathlineto{\pgfqpoint{0.740015in}{1.595725in}}%
\pgfpathlineto{\pgfqpoint{0.734282in}{1.584259in}}%
\pgfpathlineto{\pgfqpoint{0.727377in}{1.578133in}}%
\pgfpathlineto{\pgfqpoint{0.726752in}{1.572002in}}%
\pgfpathlineto{\pgfqpoint{0.715313in}{1.563715in}}%
\pgfpathlineto{\pgfqpoint{0.707464in}{1.551738in}}%
\pgfpathlineto{\pgfqpoint{0.703070in}{1.548828in}}%
\pgfpathlineto{\pgfqpoint{0.700673in}{1.537205in}}%
\pgfpathlineto{\pgfqpoint{0.712217in}{1.525368in}}%
\pgfpathlineto{\pgfqpoint{0.701069in}{1.505674in}}%
\pgfpathlineto{\pgfqpoint{0.699717in}{1.498158in}}%
\pgfpathlineto{\pgfqpoint{0.674647in}{1.395314in}}%
\pgfpathlineto{\pgfqpoint{0.621916in}{1.408826in}}%
\pgfpathlineto{\pgfqpoint{0.571046in}{1.421943in}}%
\pgfpathlineto{\pgfqpoint{0.540358in}{1.430483in}}%
\pgfpathlineto{\pgfqpoint{0.500927in}{1.441715in}}%
\pgfpathlineto{\pgfqpoint{0.438032in}{1.461511in}}%
\pgfpathlineto{\pgfqpoint{0.369661in}{1.482798in}}%
\pgfpathlineto{\pgfqpoint{0.352592in}{1.488751in}}%
\pgfpathclose%
\pgfusepath{fill}%
\end{pgfscope}%
\begin{pgfscope}%
\pgfpathrectangle{\pgfqpoint{0.100000in}{0.100000in}}{\pgfqpoint{2.862985in}{1.832778in}}%
\pgfusepath{clip}%
\pgfsetbuttcap%
\pgfsetmiterjoin%
\definecolor{currentfill}{rgb}{0.175317,0.760323,0.175317}%
\pgfsetfillcolor{currentfill}%
\pgfsetlinewidth{0.000000pt}%
\definecolor{currentstroke}{rgb}{0.000000,0.000000,0.000000}%
\pgfsetstrokecolor{currentstroke}%
\pgfsetstrokeopacity{0.000000}%
\pgfsetdash{}{0pt}%
\pgfpathmoveto{\pgfqpoint{2.722258in}{1.445782in}}%
\pgfpathlineto{\pgfqpoint{2.717278in}{1.449830in}}%
\pgfpathlineto{\pgfqpoint{2.717972in}{1.462444in}}%
\pgfpathlineto{\pgfqpoint{2.712824in}{1.498298in}}%
\pgfpathlineto{\pgfqpoint{2.716859in}{1.510906in}}%
\pgfpathlineto{\pgfqpoint{2.719741in}{1.533267in}}%
\pgfpathlineto{\pgfqpoint{2.715852in}{1.543846in}}%
\pgfpathlineto{\pgfqpoint{2.729999in}{1.557930in}}%
\pgfpathlineto{\pgfqpoint{2.733559in}{1.569133in}}%
\pgfpathlineto{\pgfqpoint{2.726856in}{1.577975in}}%
\pgfpathlineto{\pgfqpoint{2.729711in}{1.589025in}}%
\pgfpathlineto{\pgfqpoint{2.727922in}{1.595966in}}%
\pgfpathlineto{\pgfqpoint{2.729459in}{1.610832in}}%
\pgfpathlineto{\pgfqpoint{2.732332in}{1.615421in}}%
\pgfpathlineto{\pgfqpoint{2.741192in}{1.618318in}}%
\pgfpathlineto{\pgfqpoint{2.754117in}{1.580468in}}%
\pgfpathlineto{\pgfqpoint{2.778094in}{1.502035in}}%
\pgfpathlineto{\pgfqpoint{2.787168in}{1.495954in}}%
\pgfpathlineto{\pgfqpoint{2.788801in}{1.489012in}}%
\pgfpathlineto{\pgfqpoint{2.793590in}{1.486207in}}%
\pgfpathlineto{\pgfqpoint{2.793226in}{1.473617in}}%
\pgfpathlineto{\pgfqpoint{2.788151in}{1.473415in}}%
\pgfpathlineto{\pgfqpoint{2.777871in}{1.465574in}}%
\pgfpathlineto{\pgfqpoint{2.774906in}{1.457709in}}%
\pgfpathlineto{\pgfqpoint{2.747390in}{1.450901in}}%
\pgfpathlineto{\pgfqpoint{2.722258in}{1.445782in}}%
\pgfpathclose%
\pgfusepath{fill}%
\end{pgfscope}%
\begin{pgfscope}%
\pgfpathrectangle{\pgfqpoint{0.100000in}{0.100000in}}{\pgfqpoint{2.862985in}{1.832778in}}%
\pgfusepath{clip}%
\pgfsetbuttcap%
\pgfsetmiterjoin%
\definecolor{currentfill}{rgb}{0.073818,0.547174,0.073818}%
\pgfsetfillcolor{currentfill}%
\pgfsetlinewidth{0.000000pt}%
\definecolor{currentstroke}{rgb}{0.000000,0.000000,0.000000}%
\pgfsetstrokecolor{currentstroke}%
\pgfsetstrokeopacity{0.000000}%
\pgfsetdash{}{0pt}%
\pgfpathmoveto{\pgfqpoint{1.861865in}{1.172908in}}%
\pgfpathlineto{\pgfqpoint{1.847339in}{1.187293in}}%
\pgfpathlineto{\pgfqpoint{1.800906in}{1.184615in}}%
\pgfpathlineto{\pgfqpoint{1.728489in}{1.182229in}}%
\pgfpathlineto{\pgfqpoint{1.655637in}{1.183365in}}%
\pgfpathlineto{\pgfqpoint{1.650523in}{1.192282in}}%
\pgfpathlineto{\pgfqpoint{1.652611in}{1.201038in}}%
\pgfpathlineto{\pgfqpoint{1.651620in}{1.216037in}}%
\pgfpathlineto{\pgfqpoint{1.647770in}{1.230565in}}%
\pgfpathlineto{\pgfqpoint{1.648217in}{1.238237in}}%
\pgfpathlineto{\pgfqpoint{1.640033in}{1.246905in}}%
\pgfpathlineto{\pgfqpoint{1.641814in}{1.258967in}}%
\pgfpathlineto{\pgfqpoint{1.629260in}{1.282791in}}%
\pgfpathlineto{\pgfqpoint{1.625519in}{1.303018in}}%
\pgfpathlineto{\pgfqpoint{1.616489in}{1.319470in}}%
\pgfpathlineto{\pgfqpoint{1.625061in}{1.339878in}}%
\pgfpathlineto{\pgfqpoint{1.620791in}{1.349375in}}%
\pgfpathlineto{\pgfqpoint{1.625565in}{1.366544in}}%
\pgfpathlineto{\pgfqpoint{1.670992in}{1.366052in}}%
\pgfpathlineto{\pgfqpoint{1.717055in}{1.365982in}}%
\pgfpathlineto{\pgfqpoint{1.760597in}{1.366622in}}%
\pgfpathlineto{\pgfqpoint{1.808388in}{1.367621in}}%
\pgfpathlineto{\pgfqpoint{1.864454in}{1.369657in}}%
\pgfpathlineto{\pgfqpoint{1.865951in}{1.361517in}}%
\pgfpathlineto{\pgfqpoint{1.872158in}{1.354693in}}%
\pgfpathlineto{\pgfqpoint{1.868228in}{1.347582in}}%
\pgfpathlineto{\pgfqpoint{1.869839in}{1.333702in}}%
\pgfpathlineto{\pgfqpoint{1.873686in}{1.322483in}}%
\pgfpathlineto{\pgfqpoint{1.891125in}{1.316544in}}%
\pgfpathlineto{\pgfqpoint{1.894114in}{1.308600in}}%
\pgfpathlineto{\pgfqpoint{1.903684in}{1.299681in}}%
\pgfpathlineto{\pgfqpoint{1.907588in}{1.290454in}}%
\pgfpathlineto{\pgfqpoint{1.917285in}{1.284265in}}%
\pgfpathlineto{\pgfqpoint{1.918802in}{1.276813in}}%
\pgfpathlineto{\pgfqpoint{1.916914in}{1.265532in}}%
\pgfpathlineto{\pgfqpoint{1.911979in}{1.262151in}}%
\pgfpathlineto{\pgfqpoint{1.910490in}{1.251413in}}%
\pgfpathlineto{\pgfqpoint{1.896306in}{1.242884in}}%
\pgfpathlineto{\pgfqpoint{1.877818in}{1.238212in}}%
\pgfpathlineto{\pgfqpoint{1.876124in}{1.227469in}}%
\pgfpathlineto{\pgfqpoint{1.883256in}{1.219794in}}%
\pgfpathlineto{\pgfqpoint{1.883547in}{1.210143in}}%
\pgfpathlineto{\pgfqpoint{1.877793in}{1.202558in}}%
\pgfpathlineto{\pgfqpoint{1.874773in}{1.191288in}}%
\pgfpathlineto{\pgfqpoint{1.864771in}{1.187558in}}%
\pgfpathlineto{\pgfqpoint{1.865408in}{1.174999in}}%
\pgfpathlineto{\pgfqpoint{1.861865in}{1.172908in}}%
\pgfpathclose%
\pgfusepath{fill}%
\end{pgfscope}%
\begin{pgfscope}%
\pgfpathrectangle{\pgfqpoint{0.100000in}{0.100000in}}{\pgfqpoint{2.862985in}{1.832778in}}%
\pgfusepath{clip}%
\pgfsetbuttcap%
\pgfsetmiterjoin%
\definecolor{currentfill}{rgb}{0.158401,0.793372,0.350557}%
\pgfsetfillcolor{currentfill}%
\pgfsetlinewidth{0.000000pt}%
\definecolor{currentstroke}{rgb}{0.000000,0.000000,0.000000}%
\pgfsetstrokecolor{currentstroke}%
\pgfsetstrokeopacity{0.000000}%
\pgfsetdash{}{0pt}%
\pgfpathmoveto{\pgfqpoint{2.763120in}{1.409488in}}%
\pgfpathlineto{\pgfqpoint{2.749362in}{1.407318in}}%
\pgfpathlineto{\pgfqpoint{2.686166in}{1.392957in}}%
\pgfpathlineto{\pgfqpoint{2.685061in}{1.394631in}}%
\pgfpathlineto{\pgfqpoint{2.686015in}{1.438080in}}%
\pgfpathlineto{\pgfqpoint{2.722258in}{1.445782in}}%
\pgfpathlineto{\pgfqpoint{2.747390in}{1.450901in}}%
\pgfpathlineto{\pgfqpoint{2.774906in}{1.457709in}}%
\pgfpathlineto{\pgfqpoint{2.777871in}{1.465574in}}%
\pgfpathlineto{\pgfqpoint{2.788151in}{1.473415in}}%
\pgfpathlineto{\pgfqpoint{2.793226in}{1.473617in}}%
\pgfpathlineto{\pgfqpoint{2.799896in}{1.462179in}}%
\pgfpathlineto{\pgfqpoint{2.793843in}{1.445482in}}%
\pgfpathlineto{\pgfqpoint{2.792955in}{1.435691in}}%
\pgfpathlineto{\pgfqpoint{2.805207in}{1.436611in}}%
\pgfpathlineto{\pgfqpoint{2.810763in}{1.431914in}}%
\pgfpathlineto{\pgfqpoint{2.823224in}{1.412717in}}%
\pgfpathlineto{\pgfqpoint{2.829417in}{1.410382in}}%
\pgfpathlineto{\pgfqpoint{2.839790in}{1.411248in}}%
\pgfpathlineto{\pgfqpoint{2.847009in}{1.417772in}}%
\pgfpathlineto{\pgfqpoint{2.851808in}{1.411950in}}%
\pgfpathlineto{\pgfqpoint{2.821631in}{1.396025in}}%
\pgfpathlineto{\pgfqpoint{2.820688in}{1.407452in}}%
\pgfpathlineto{\pgfqpoint{2.806980in}{1.389693in}}%
\pgfpathlineto{\pgfqpoint{2.802160in}{1.386633in}}%
\pgfpathlineto{\pgfqpoint{2.795438in}{1.396881in}}%
\pgfpathlineto{\pgfqpoint{2.793598in}{1.398285in}}%
\pgfpathlineto{\pgfqpoint{2.787340in}{1.401600in}}%
\pgfpathlineto{\pgfqpoint{2.781882in}{1.415045in}}%
\pgfpathlineto{\pgfqpoint{2.763120in}{1.409488in}}%
\pgfpathclose%
\pgfusepath{fill}%
\end{pgfscope}%
\begin{pgfscope}%
\pgfpathrectangle{\pgfqpoint{0.100000in}{0.100000in}}{\pgfqpoint{2.862985in}{1.832778in}}%
\pgfusepath{clip}%
\pgfsetbuttcap%
\pgfsetmiterjoin%
\definecolor{currentfill}{rgb}{0.129181,0.663437,0.129181}%
\pgfsetfillcolor{currentfill}%
\pgfsetlinewidth{0.000000pt}%
\definecolor{currentstroke}{rgb}{0.000000,0.000000,0.000000}%
\pgfsetstrokecolor{currentstroke}%
\pgfsetstrokeopacity{0.000000}%
\pgfsetdash{}{0pt}%
\pgfpathmoveto{\pgfqpoint{1.353739in}{1.159333in}}%
\pgfpathlineto{\pgfqpoint{1.358777in}{1.221808in}}%
\pgfpathlineto{\pgfqpoint{1.330243in}{1.224124in}}%
\pgfpathlineto{\pgfqpoint{1.264283in}{1.230524in}}%
\pgfpathlineto{\pgfqpoint{1.270419in}{1.292777in}}%
\pgfpathlineto{\pgfqpoint{1.276565in}{1.355506in}}%
\pgfpathlineto{\pgfqpoint{1.301922in}{1.352768in}}%
\pgfpathlineto{\pgfqpoint{1.366682in}{1.346649in}}%
\pgfpathlineto{\pgfqpoint{1.405874in}{1.343580in}}%
\pgfpathlineto{\pgfqpoint{1.453322in}{1.340860in}}%
\pgfpathlineto{\pgfqpoint{1.531420in}{1.337164in}}%
\pgfpathlineto{\pgfqpoint{1.533137in}{1.333663in}}%
\pgfpathlineto{\pgfqpoint{1.555515in}{1.322591in}}%
\pgfpathlineto{\pgfqpoint{1.562377in}{1.328214in}}%
\pgfpathlineto{\pgfqpoint{1.582148in}{1.327701in}}%
\pgfpathlineto{\pgfqpoint{1.593930in}{1.321529in}}%
\pgfpathlineto{\pgfqpoint{1.612612in}{1.314513in}}%
\pgfpathlineto{\pgfqpoint{1.617847in}{1.304277in}}%
\pgfpathlineto{\pgfqpoint{1.625519in}{1.303018in}}%
\pgfpathlineto{\pgfqpoint{1.629260in}{1.282791in}}%
\pgfpathlineto{\pgfqpoint{1.641814in}{1.258967in}}%
\pgfpathlineto{\pgfqpoint{1.640033in}{1.246905in}}%
\pgfpathlineto{\pgfqpoint{1.648217in}{1.238237in}}%
\pgfpathlineto{\pgfqpoint{1.647770in}{1.230565in}}%
\pgfpathlineto{\pgfqpoint{1.651620in}{1.216037in}}%
\pgfpathlineto{\pgfqpoint{1.652611in}{1.201038in}}%
\pgfpathlineto{\pgfqpoint{1.650523in}{1.192282in}}%
\pgfpathlineto{\pgfqpoint{1.655637in}{1.183365in}}%
\pgfpathlineto{\pgfqpoint{1.662651in}{1.167151in}}%
\pgfpathlineto{\pgfqpoint{1.669364in}{1.160551in}}%
\pgfpathlineto{\pgfqpoint{1.677367in}{1.146251in}}%
\pgfpathlineto{\pgfqpoint{1.654686in}{1.146016in}}%
\pgfpathlineto{\pgfqpoint{1.605647in}{1.146774in}}%
\pgfpathlineto{\pgfqpoint{1.551466in}{1.148434in}}%
\pgfpathlineto{\pgfqpoint{1.496966in}{1.150526in}}%
\pgfpathlineto{\pgfqpoint{1.416835in}{1.154903in}}%
\pgfpathlineto{\pgfqpoint{1.353739in}{1.159333in}}%
\pgfpathclose%
\pgfusepath{fill}%
\end{pgfscope}%
\begin{pgfscope}%
\pgfpathrectangle{\pgfqpoint{0.100000in}{0.100000in}}{\pgfqpoint{2.862985in}{1.832778in}}%
\pgfusepath{clip}%
\pgfsetbuttcap%
\pgfsetmiterjoin%
\definecolor{currentfill}{rgb}{0.066128,0.767536,0.728874}%
\pgfsetfillcolor{currentfill}%
\pgfsetlinewidth{0.000000pt}%
\definecolor{currentstroke}{rgb}{0.000000,0.000000,0.000000}%
\pgfsetstrokecolor{currentstroke}%
\pgfsetstrokeopacity{0.000000}%
\pgfsetdash{}{0pt}%
\pgfpathmoveto{\pgfqpoint{2.397084in}{1.347781in}}%
\pgfpathlineto{\pgfqpoint{2.422175in}{1.371705in}}%
\pgfpathlineto{\pgfqpoint{2.425294in}{1.380209in}}%
\pgfpathlineto{\pgfqpoint{2.432646in}{1.387490in}}%
\pgfpathlineto{\pgfqpoint{2.427129in}{1.398090in}}%
\pgfpathlineto{\pgfqpoint{2.420212in}{1.404290in}}%
\pgfpathlineto{\pgfqpoint{2.418213in}{1.415277in}}%
\pgfpathlineto{\pgfqpoint{2.443959in}{1.426558in}}%
\pgfpathlineto{\pgfqpoint{2.465274in}{1.430126in}}%
\pgfpathlineto{\pgfqpoint{2.476728in}{1.430365in}}%
\pgfpathlineto{\pgfqpoint{2.485453in}{1.426054in}}%
\pgfpathlineto{\pgfqpoint{2.493958in}{1.429903in}}%
\pgfpathlineto{\pgfqpoint{2.514702in}{1.434217in}}%
\pgfpathlineto{\pgfqpoint{2.521870in}{1.439785in}}%
\pgfpathlineto{\pgfqpoint{2.532498in}{1.452111in}}%
\pgfpathlineto{\pgfqpoint{2.542167in}{1.457558in}}%
\pgfpathlineto{\pgfqpoint{2.542849in}{1.462782in}}%
\pgfpathlineto{\pgfqpoint{2.537774in}{1.474701in}}%
\pgfpathlineto{\pgfqpoint{2.541443in}{1.481714in}}%
\pgfpathlineto{\pgfqpoint{2.536503in}{1.489256in}}%
\pgfpathlineto{\pgfqpoint{2.528939in}{1.489782in}}%
\pgfpathlineto{\pgfqpoint{2.547870in}{1.512563in}}%
\pgfpathlineto{\pgfqpoint{2.550119in}{1.521245in}}%
\pgfpathlineto{\pgfqpoint{2.565018in}{1.543386in}}%
\pgfpathlineto{\pgfqpoint{2.578850in}{1.555368in}}%
\pgfpathlineto{\pgfqpoint{2.588347in}{1.560371in}}%
\pgfpathlineto{\pgfqpoint{2.619425in}{1.567413in}}%
\pgfpathlineto{\pgfqpoint{2.648552in}{1.575365in}}%
\pgfpathlineto{\pgfqpoint{2.652410in}{1.562882in}}%
\pgfpathlineto{\pgfqpoint{2.652683in}{1.551308in}}%
\pgfpathlineto{\pgfqpoint{2.659327in}{1.540801in}}%
\pgfpathlineto{\pgfqpoint{2.661375in}{1.530104in}}%
\pgfpathlineto{\pgfqpoint{2.658822in}{1.515590in}}%
\pgfpathlineto{\pgfqpoint{2.666598in}{1.496013in}}%
\pgfpathlineto{\pgfqpoint{2.675463in}{1.485127in}}%
\pgfpathlineto{\pgfqpoint{2.682320in}{1.449744in}}%
\pgfpathlineto{\pgfqpoint{2.686015in}{1.438080in}}%
\pgfpathlineto{\pgfqpoint{2.685061in}{1.394631in}}%
\pgfpathlineto{\pgfqpoint{2.686166in}{1.392957in}}%
\pgfpathlineto{\pgfqpoint{2.694238in}{1.346227in}}%
\pgfpathlineto{\pgfqpoint{2.698765in}{1.341969in}}%
\pgfpathlineto{\pgfqpoint{2.689028in}{1.332508in}}%
\pgfpathlineto{\pgfqpoint{2.693832in}{1.327078in}}%
\pgfpathlineto{\pgfqpoint{2.689609in}{1.318865in}}%
\pgfpathlineto{\pgfqpoint{2.689645in}{1.315368in}}%
\pgfpathlineto{\pgfqpoint{2.684366in}{1.312213in}}%
\pgfpathlineto{\pgfqpoint{2.681795in}{1.305239in}}%
\pgfpathlineto{\pgfqpoint{2.682610in}{1.324417in}}%
\pgfpathlineto{\pgfqpoint{2.666246in}{1.328639in}}%
\pgfpathlineto{\pgfqpoint{2.640655in}{1.337368in}}%
\pgfpathlineto{\pgfqpoint{2.637035in}{1.341668in}}%
\pgfpathlineto{\pgfqpoint{2.626342in}{1.342698in}}%
\pgfpathlineto{\pgfqpoint{2.620190in}{1.349098in}}%
\pgfpathlineto{\pgfqpoint{2.616871in}{1.361844in}}%
\pgfpathlineto{\pgfqpoint{2.608151in}{1.363429in}}%
\pgfpathlineto{\pgfqpoint{2.602319in}{1.370116in}}%
\pgfpathlineto{\pgfqpoint{2.528076in}{1.355148in}}%
\pgfpathlineto{\pgfqpoint{2.492594in}{1.347810in}}%
\pgfpathlineto{\pgfqpoint{2.438721in}{1.337977in}}%
\pgfpathlineto{\pgfqpoint{2.399926in}{1.331432in}}%
\pgfpathlineto{\pgfqpoint{2.397084in}{1.347781in}}%
\pgfpathclose%
\pgfusepath{fill}%
\end{pgfscope}%
\begin{pgfscope}%
\pgfpathrectangle{\pgfqpoint{0.100000in}{0.100000in}}{\pgfqpoint{2.862985in}{1.832778in}}%
\pgfusepath{clip}%
\pgfsetbuttcap%
\pgfsetmiterjoin%
\definecolor{currentfill}{rgb}{0.066128,0.767536,0.728874}%
\pgfsetfillcolor{currentfill}%
\pgfsetlinewidth{0.000000pt}%
\definecolor{currentstroke}{rgb}{0.000000,0.000000,0.000000}%
\pgfsetstrokecolor{currentstroke}%
\pgfsetstrokeopacity{0.000000}%
\pgfsetdash{}{0pt}%
\pgfpathmoveto{\pgfqpoint{2.695139in}{1.301344in}}%
\pgfpathlineto{\pgfqpoint{2.683657in}{1.297767in}}%
\pgfpathlineto{\pgfqpoint{2.681726in}{1.301055in}}%
\pgfpathlineto{\pgfqpoint{2.685411in}{1.312094in}}%
\pgfpathlineto{\pgfqpoint{2.692237in}{1.313181in}}%
\pgfpathlineto{\pgfqpoint{2.691626in}{1.316656in}}%
\pgfpathlineto{\pgfqpoint{2.697754in}{1.321872in}}%
\pgfpathlineto{\pgfqpoint{2.715472in}{1.326028in}}%
\pgfpathlineto{\pgfqpoint{2.723360in}{1.332319in}}%
\pgfpathlineto{\pgfqpoint{2.741079in}{1.337561in}}%
\pgfpathlineto{\pgfqpoint{2.749161in}{1.335643in}}%
\pgfpathlineto{\pgfqpoint{2.755959in}{1.344095in}}%
\pgfpathlineto{\pgfqpoint{2.766234in}{1.345211in}}%
\pgfpathlineto{\pgfqpoint{2.748704in}{1.328725in}}%
\pgfpathlineto{\pgfqpoint{2.695139in}{1.301344in}}%
\pgfpathclose%
\pgfusepath{fill}%
\end{pgfscope}%
\begin{pgfscope}%
\pgfpathrectangle{\pgfqpoint{0.100000in}{0.100000in}}{\pgfqpoint{2.862985in}{1.832778in}}%
\pgfusepath{clip}%
\pgfsetbuttcap%
\pgfsetmiterjoin%
\definecolor{currentfill}{rgb}{0.192234,0.802845,0.211842}%
\pgfsetfillcolor{currentfill}%
\pgfsetlinewidth{0.000000pt}%
\definecolor{currentstroke}{rgb}{0.000000,0.000000,0.000000}%
\pgfsetstrokecolor{currentstroke}%
\pgfsetstrokeopacity{0.000000}%
\pgfsetdash{}{0pt}%
\pgfpathmoveto{\pgfqpoint{2.437178in}{1.192485in}}%
\pgfpathlineto{\pgfqpoint{2.387526in}{1.184275in}}%
\pgfpathlineto{\pgfqpoint{2.378519in}{1.241018in}}%
\pgfpathlineto{\pgfqpoint{2.365146in}{1.324661in}}%
\pgfpathlineto{\pgfqpoint{2.397084in}{1.347781in}}%
\pgfpathlineto{\pgfqpoint{2.399926in}{1.331432in}}%
\pgfpathlineto{\pgfqpoint{2.438721in}{1.337977in}}%
\pgfpathlineto{\pgfqpoint{2.492594in}{1.347810in}}%
\pgfpathlineto{\pgfqpoint{2.528076in}{1.355148in}}%
\pgfpathlineto{\pgfqpoint{2.602319in}{1.370116in}}%
\pgfpathlineto{\pgfqpoint{2.608151in}{1.363429in}}%
\pgfpathlineto{\pgfqpoint{2.616871in}{1.361844in}}%
\pgfpathlineto{\pgfqpoint{2.620190in}{1.349098in}}%
\pgfpathlineto{\pgfqpoint{2.626342in}{1.342698in}}%
\pgfpathlineto{\pgfqpoint{2.637035in}{1.341668in}}%
\pgfpathlineto{\pgfqpoint{2.640655in}{1.337368in}}%
\pgfpathlineto{\pgfqpoint{2.636981in}{1.334051in}}%
\pgfpathlineto{\pgfqpoint{2.633670in}{1.322314in}}%
\pgfpathlineto{\pgfqpoint{2.625913in}{1.309125in}}%
\pgfpathlineto{\pgfqpoint{2.631120in}{1.303390in}}%
\pgfpathlineto{\pgfqpoint{2.626859in}{1.297245in}}%
\pgfpathlineto{\pgfqpoint{2.629263in}{1.283769in}}%
\pgfpathlineto{\pgfqpoint{2.636936in}{1.276695in}}%
\pgfpathlineto{\pgfqpoint{2.655144in}{1.265202in}}%
\pgfpathlineto{\pgfqpoint{2.640481in}{1.248986in}}%
\pgfpathlineto{\pgfqpoint{2.640275in}{1.242831in}}%
\pgfpathlineto{\pgfqpoint{2.628338in}{1.234894in}}%
\pgfpathlineto{\pgfqpoint{2.615138in}{1.233404in}}%
\pgfpathlineto{\pgfqpoint{2.611873in}{1.226505in}}%
\pgfpathlineto{\pgfqpoint{2.575158in}{1.218561in}}%
\pgfpathlineto{\pgfqpoint{2.532316in}{1.209933in}}%
\pgfpathlineto{\pgfqpoint{2.502868in}{1.204659in}}%
\pgfpathlineto{\pgfqpoint{2.437178in}{1.192485in}}%
\pgfpathclose%
\pgfusepath{fill}%
\end{pgfscope}%
\begin{pgfscope}%
\pgfpathrectangle{\pgfqpoint{0.100000in}{0.100000in}}{\pgfqpoint{2.862985in}{1.832778in}}%
\pgfusepath{clip}%
\pgfsetbuttcap%
\pgfsetmiterjoin%
\definecolor{currentfill}{rgb}{0.163014,0.734487,0.163014}%
\pgfsetfillcolor{currentfill}%
\pgfsetlinewidth{0.000000pt}%
\definecolor{currentstroke}{rgb}{0.000000,0.000000,0.000000}%
\pgfsetstrokecolor{currentstroke}%
\pgfsetstrokeopacity{0.000000}%
\pgfsetdash{}{0pt}%
\pgfpathmoveto{\pgfqpoint{2.686166in}{1.392957in}}%
\pgfpathlineto{\pgfqpoint{2.749362in}{1.407318in}}%
\pgfpathlineto{\pgfqpoint{2.763120in}{1.409488in}}%
\pgfpathlineto{\pgfqpoint{2.772226in}{1.373686in}}%
\pgfpathlineto{\pgfqpoint{2.770785in}{1.367275in}}%
\pgfpathlineto{\pgfqpoint{2.741533in}{1.355951in}}%
\pgfpathlineto{\pgfqpoint{2.724071in}{1.351992in}}%
\pgfpathlineto{\pgfqpoint{2.716651in}{1.343114in}}%
\pgfpathlineto{\pgfqpoint{2.693832in}{1.327078in}}%
\pgfpathlineto{\pgfqpoint{2.689028in}{1.332508in}}%
\pgfpathlineto{\pgfqpoint{2.698765in}{1.341969in}}%
\pgfpathlineto{\pgfqpoint{2.694238in}{1.346227in}}%
\pgfpathlineto{\pgfqpoint{2.686166in}{1.392957in}}%
\pgfpathclose%
\pgfusepath{fill}%
\end{pgfscope}%
\begin{pgfscope}%
\pgfpathrectangle{\pgfqpoint{0.100000in}{0.100000in}}{\pgfqpoint{2.862985in}{1.832778in}}%
\pgfusepath{clip}%
\pgfsetbuttcap%
\pgfsetmiterjoin%
\definecolor{currentfill}{rgb}{0.164552,0.795094,0.325336}%
\pgfsetfillcolor{currentfill}%
\pgfsetlinewidth{0.000000pt}%
\definecolor{currentstroke}{rgb}{0.000000,0.000000,0.000000}%
\pgfsetstrokecolor{currentstroke}%
\pgfsetstrokeopacity{0.000000}%
\pgfsetdash{}{0pt}%
\pgfpathmoveto{\pgfqpoint{2.763120in}{1.409488in}}%
\pgfpathlineto{\pgfqpoint{2.781882in}{1.415045in}}%
\pgfpathlineto{\pgfqpoint{2.787340in}{1.401600in}}%
\pgfpathlineto{\pgfqpoint{2.793598in}{1.398285in}}%
\pgfpathlineto{\pgfqpoint{2.785875in}{1.392622in}}%
\pgfpathlineto{\pgfqpoint{2.786849in}{1.375988in}}%
\pgfpathlineto{\pgfqpoint{2.770785in}{1.367275in}}%
\pgfpathlineto{\pgfqpoint{2.772226in}{1.373686in}}%
\pgfpathlineto{\pgfqpoint{2.763120in}{1.409488in}}%
\pgfpathclose%
\pgfusepath{fill}%
\end{pgfscope}%
\begin{pgfscope}%
\pgfpathrectangle{\pgfqpoint{0.100000in}{0.100000in}}{\pgfqpoint{2.862985in}{1.832778in}}%
\pgfusepath{clip}%
\pgfsetbuttcap%
\pgfsetmiterjoin%
\definecolor{currentfill}{rgb}{0.146098,0.789927,0.401000}%
\pgfsetfillcolor{currentfill}%
\pgfsetlinewidth{0.000000pt}%
\definecolor{currentstroke}{rgb}{0.000000,0.000000,0.000000}%
\pgfsetstrokecolor{currentstroke}%
\pgfsetstrokeopacity{0.000000}%
\pgfsetdash{}{0pt}%
\pgfpathmoveto{\pgfqpoint{2.626223in}{1.229031in}}%
\pgfpathlineto{\pgfqpoint{2.628338in}{1.234894in}}%
\pgfpathlineto{\pgfqpoint{2.640275in}{1.242831in}}%
\pgfpathlineto{\pgfqpoint{2.640481in}{1.248986in}}%
\pgfpathlineto{\pgfqpoint{2.655144in}{1.265202in}}%
\pgfpathlineto{\pgfqpoint{2.636936in}{1.276695in}}%
\pgfpathlineto{\pgfqpoint{2.629263in}{1.283769in}}%
\pgfpathlineto{\pgfqpoint{2.626859in}{1.297245in}}%
\pgfpathlineto{\pgfqpoint{2.631120in}{1.303390in}}%
\pgfpathlineto{\pgfqpoint{2.625913in}{1.309125in}}%
\pgfpathlineto{\pgfqpoint{2.633670in}{1.322314in}}%
\pgfpathlineto{\pgfqpoint{2.636981in}{1.334051in}}%
\pgfpathlineto{\pgfqpoint{2.640655in}{1.337368in}}%
\pgfpathlineto{\pgfqpoint{2.666246in}{1.328639in}}%
\pgfpathlineto{\pgfqpoint{2.682610in}{1.324417in}}%
\pgfpathlineto{\pgfqpoint{2.681795in}{1.305239in}}%
\pgfpathlineto{\pgfqpoint{2.671854in}{1.290696in}}%
\pgfpathlineto{\pgfqpoint{2.680045in}{1.288556in}}%
\pgfpathlineto{\pgfqpoint{2.688550in}{1.282315in}}%
\pgfpathlineto{\pgfqpoint{2.689051in}{1.265245in}}%
\pgfpathlineto{\pgfqpoint{2.686477in}{1.253161in}}%
\pgfpathlineto{\pgfqpoint{2.688190in}{1.243235in}}%
\pgfpathlineto{\pgfqpoint{2.673407in}{1.209740in}}%
\pgfpathlineto{\pgfqpoint{2.668094in}{1.194100in}}%
\pgfpathlineto{\pgfqpoint{2.660676in}{1.201703in}}%
\pgfpathlineto{\pgfqpoint{2.650838in}{1.200407in}}%
\pgfpathlineto{\pgfqpoint{2.626218in}{1.214632in}}%
\pgfpathlineto{\pgfqpoint{2.623699in}{1.222249in}}%
\pgfpathlineto{\pgfqpoint{2.626223in}{1.229031in}}%
\pgfpathclose%
\pgfusepath{fill}%
\end{pgfscope}%
\begin{pgfscope}%
\pgfpathrectangle{\pgfqpoint{0.100000in}{0.100000in}}{\pgfqpoint{2.862985in}{1.832778in}}%
\pgfusepath{clip}%
\pgfsetbuttcap%
\pgfsetmiterjoin%
\definecolor{currentfill}{rgb}{0.152249,0.791649,0.375779}%
\pgfsetfillcolor{currentfill}%
\pgfsetlinewidth{0.000000pt}%
\definecolor{currentstroke}{rgb}{0.000000,0.000000,0.000000}%
\pgfsetstrokecolor{currentstroke}%
\pgfsetstrokeopacity{0.000000}%
\pgfsetdash{}{0pt}%
\pgfpathmoveto{\pgfqpoint{2.037013in}{1.021868in}}%
\pgfpathlineto{\pgfqpoint{2.035862in}{1.037407in}}%
\pgfpathlineto{\pgfqpoint{2.040336in}{1.045116in}}%
\pgfpathlineto{\pgfqpoint{2.037387in}{1.048893in}}%
\pgfpathlineto{\pgfqpoint{2.048140in}{1.061261in}}%
\pgfpathlineto{\pgfqpoint{2.047524in}{1.064241in}}%
\pgfpathlineto{\pgfqpoint{2.057971in}{1.085818in}}%
\pgfpathlineto{\pgfqpoint{2.055859in}{1.096224in}}%
\pgfpathlineto{\pgfqpoint{2.048283in}{1.107118in}}%
\pgfpathlineto{\pgfqpoint{2.053542in}{1.120373in}}%
\pgfpathlineto{\pgfqpoint{2.046726in}{1.207597in}}%
\pgfpathlineto{\pgfqpoint{2.041818in}{1.268787in}}%
\pgfpathlineto{\pgfqpoint{2.048598in}{1.263717in}}%
\pgfpathlineto{\pgfqpoint{2.056163in}{1.263854in}}%
\pgfpathlineto{\pgfqpoint{2.074029in}{1.274203in}}%
\pgfpathlineto{\pgfqpoint{2.128831in}{1.279317in}}%
\pgfpathlineto{\pgfqpoint{2.169393in}{1.283592in}}%
\pgfpathlineto{\pgfqpoint{2.169752in}{1.279623in}}%
\pgfpathlineto{\pgfqpoint{2.179011in}{1.195770in}}%
\pgfpathlineto{\pgfqpoint{2.186981in}{1.117748in}}%
\pgfpathlineto{\pgfqpoint{2.183545in}{1.114085in}}%
\pgfpathlineto{\pgfqpoint{2.188822in}{1.104993in}}%
\pgfpathlineto{\pgfqpoint{2.188795in}{1.098441in}}%
\pgfpathlineto{\pgfqpoint{2.181265in}{1.096795in}}%
\pgfpathlineto{\pgfqpoint{2.172839in}{1.090478in}}%
\pgfpathlineto{\pgfqpoint{2.167134in}{1.092965in}}%
\pgfpathlineto{\pgfqpoint{2.158595in}{1.088922in}}%
\pgfpathlineto{\pgfqpoint{2.161237in}{1.080801in}}%
\pgfpathlineto{\pgfqpoint{2.152460in}{1.072642in}}%
\pgfpathlineto{\pgfqpoint{2.150037in}{1.063205in}}%
\pgfpathlineto{\pgfqpoint{2.144002in}{1.061656in}}%
\pgfpathlineto{\pgfqpoint{2.139498in}{1.054509in}}%
\pgfpathlineto{\pgfqpoint{2.140088in}{1.047311in}}%
\pgfpathlineto{\pgfqpoint{2.134790in}{1.042250in}}%
\pgfpathlineto{\pgfqpoint{2.126817in}{1.043031in}}%
\pgfpathlineto{\pgfqpoint{2.120757in}{1.050797in}}%
\pgfpathlineto{\pgfqpoint{2.110489in}{1.043318in}}%
\pgfpathlineto{\pgfqpoint{2.111209in}{1.036783in}}%
\pgfpathlineto{\pgfqpoint{2.099789in}{1.032965in}}%
\pgfpathlineto{\pgfqpoint{2.096919in}{1.037771in}}%
\pgfpathlineto{\pgfqpoint{2.085801in}{1.032321in}}%
\pgfpathlineto{\pgfqpoint{2.081739in}{1.024520in}}%
\pgfpathlineto{\pgfqpoint{2.068349in}{1.032520in}}%
\pgfpathlineto{\pgfqpoint{2.047127in}{1.027220in}}%
\pgfpathlineto{\pgfqpoint{2.037013in}{1.021868in}}%
\pgfpathclose%
\pgfusepath{fill}%
\end{pgfscope}%
\begin{pgfscope}%
\pgfpathrectangle{\pgfqpoint{0.100000in}{0.100000in}}{\pgfqpoint{2.862985in}{1.832778in}}%
\pgfusepath{clip}%
\pgfsetbuttcap%
\pgfsetmiterjoin%
\definecolor{currentfill}{rgb}{0.072280,0.769258,0.703652}%
\pgfsetfillcolor{currentfill}%
\pgfsetlinewidth{0.000000pt}%
\definecolor{currentstroke}{rgb}{0.000000,0.000000,0.000000}%
\pgfsetstrokecolor{currentstroke}%
\pgfsetstrokeopacity{0.000000}%
\pgfsetdash{}{0pt}%
\pgfpathmoveto{\pgfqpoint{0.540358in}{1.430483in}}%
\pgfpathlineto{\pgfqpoint{0.571046in}{1.421943in}}%
\pgfpathlineto{\pgfqpoint{0.621916in}{1.408826in}}%
\pgfpathlineto{\pgfqpoint{0.674647in}{1.395314in}}%
\pgfpathlineto{\pgfqpoint{0.723250in}{1.383593in}}%
\pgfpathlineto{\pgfqpoint{0.765246in}{1.374085in}}%
\pgfpathlineto{\pgfqpoint{0.810177in}{1.364205in}}%
\pgfpathlineto{\pgfqpoint{0.797195in}{1.302942in}}%
\pgfpathlineto{\pgfqpoint{0.785630in}{1.248537in}}%
\pgfpathlineto{\pgfqpoint{0.766658in}{1.160754in}}%
\pgfpathlineto{\pgfqpoint{0.752414in}{1.094498in}}%
\pgfpathlineto{\pgfqpoint{0.744719in}{1.057523in}}%
\pgfpathlineto{\pgfqpoint{0.734852in}{1.009534in}}%
\pgfpathlineto{\pgfqpoint{0.728025in}{0.999798in}}%
\pgfpathlineto{\pgfqpoint{0.722531in}{0.999473in}}%
\pgfpathlineto{\pgfqpoint{0.718612in}{1.007971in}}%
\pgfpathlineto{\pgfqpoint{0.709613in}{1.011057in}}%
\pgfpathlineto{\pgfqpoint{0.699916in}{1.009976in}}%
\pgfpathlineto{\pgfqpoint{0.697164in}{1.003024in}}%
\pgfpathlineto{\pgfqpoint{0.696877in}{0.978886in}}%
\pgfpathlineto{\pgfqpoint{0.693973in}{0.973374in}}%
\pgfpathlineto{\pgfqpoint{0.695633in}{0.953995in}}%
\pgfpathlineto{\pgfqpoint{0.689569in}{0.941085in}}%
\pgfpathlineto{\pgfqpoint{0.640397in}{1.016802in}}%
\pgfpathlineto{\pgfqpoint{0.591971in}{1.090379in}}%
\pgfpathlineto{\pgfqpoint{0.566763in}{1.128960in}}%
\pgfpathlineto{\pgfqpoint{0.545735in}{1.162262in}}%
\pgfpathlineto{\pgfqpoint{0.519238in}{1.203456in}}%
\pgfpathlineto{\pgfqpoint{0.489244in}{1.249665in}}%
\pgfpathlineto{\pgfqpoint{0.501576in}{1.293524in}}%
\pgfpathlineto{\pgfqpoint{0.526391in}{1.381486in}}%
\pgfpathlineto{\pgfqpoint{0.540358in}{1.430483in}}%
\pgfpathclose%
\pgfusepath{fill}%
\end{pgfscope}%
\begin{pgfscope}%
\pgfpathrectangle{\pgfqpoint{0.100000in}{0.100000in}}{\pgfqpoint{2.862985in}{1.832778in}}%
\pgfusepath{clip}%
\pgfsetbuttcap%
\pgfsetmiterjoin%
\definecolor{currentfill}{rgb}{0.139946,0.788205,0.426221}%
\pgfsetfillcolor{currentfill}%
\pgfsetlinewidth{0.000000pt}%
\definecolor{currentstroke}{rgb}{0.000000,0.000000,0.000000}%
\pgfsetstrokecolor{currentstroke}%
\pgfsetstrokeopacity{0.000000}%
\pgfsetdash{}{0pt}%
\pgfpathmoveto{\pgfqpoint{0.810177in}{1.364205in}}%
\pgfpathlineto{\pgfqpoint{0.858212in}{1.354650in}}%
\pgfpathlineto{\pgfqpoint{0.947141in}{1.337483in}}%
\pgfpathlineto{\pgfqpoint{0.936009in}{1.275734in}}%
\pgfpathlineto{\pgfqpoint{0.984967in}{1.267352in}}%
\pgfpathlineto{\pgfqpoint{1.029571in}{1.260214in}}%
\pgfpathlineto{\pgfqpoint{1.021820in}{1.211396in}}%
\pgfpathlineto{\pgfqpoint{1.013608in}{1.158753in}}%
\pgfpathlineto{\pgfqpoint{1.002613in}{1.089632in}}%
\pgfpathlineto{\pgfqpoint{1.002328in}{1.083837in}}%
\pgfpathlineto{\pgfqpoint{0.990915in}{1.012202in}}%
\pgfpathlineto{\pgfqpoint{0.943932in}{1.019484in}}%
\pgfpathlineto{\pgfqpoint{0.919991in}{1.024292in}}%
\pgfpathlineto{\pgfqpoint{0.833455in}{1.039504in}}%
\pgfpathlineto{\pgfqpoint{0.800868in}{1.045930in}}%
\pgfpathlineto{\pgfqpoint{0.744719in}{1.057523in}}%
\pgfpathlineto{\pgfqpoint{0.752414in}{1.094498in}}%
\pgfpathlineto{\pgfqpoint{0.766658in}{1.160754in}}%
\pgfpathlineto{\pgfqpoint{0.785630in}{1.248537in}}%
\pgfpathlineto{\pgfqpoint{0.797195in}{1.302942in}}%
\pgfpathlineto{\pgfqpoint{0.810177in}{1.364205in}}%
\pgfpathclose%
\pgfusepath{fill}%
\end{pgfscope}%
\begin{pgfscope}%
\pgfpathrectangle{\pgfqpoint{0.100000in}{0.100000in}}{\pgfqpoint{2.862985in}{1.832778in}}%
\pgfusepath{clip}%
\pgfsetbuttcap%
\pgfsetmiterjoin%
\definecolor{currentfill}{rgb}{0.032295,0.758062,0.867589}%
\pgfsetfillcolor{currentfill}%
\pgfsetlinewidth{0.000000pt}%
\definecolor{currentstroke}{rgb}{0.000000,0.000000,0.000000}%
\pgfsetstrokecolor{currentstroke}%
\pgfsetstrokeopacity{0.000000}%
\pgfsetdash{}{0pt}%
\pgfpathmoveto{\pgfqpoint{0.352592in}{1.488751in}}%
\pgfpathlineto{\pgfqpoint{0.369661in}{1.482798in}}%
\pgfpathlineto{\pgfqpoint{0.438032in}{1.461511in}}%
\pgfpathlineto{\pgfqpoint{0.500927in}{1.441715in}}%
\pgfpathlineto{\pgfqpoint{0.540358in}{1.430483in}}%
\pgfpathlineto{\pgfqpoint{0.526391in}{1.381486in}}%
\pgfpathlineto{\pgfqpoint{0.501576in}{1.293524in}}%
\pgfpathlineto{\pgfqpoint{0.489244in}{1.249665in}}%
\pgfpathlineto{\pgfqpoint{0.519238in}{1.203456in}}%
\pgfpathlineto{\pgfqpoint{0.545735in}{1.162262in}}%
\pgfpathlineto{\pgfqpoint{0.566763in}{1.128960in}}%
\pgfpathlineto{\pgfqpoint{0.591971in}{1.090379in}}%
\pgfpathlineto{\pgfqpoint{0.640397in}{1.016802in}}%
\pgfpathlineto{\pgfqpoint{0.689569in}{0.941085in}}%
\pgfpathlineto{\pgfqpoint{0.687595in}{0.933576in}}%
\pgfpathlineto{\pgfqpoint{0.693526in}{0.921617in}}%
\pgfpathlineto{\pgfqpoint{0.694699in}{0.905259in}}%
\pgfpathlineto{\pgfqpoint{0.703313in}{0.897324in}}%
\pgfpathlineto{\pgfqpoint{0.703654in}{0.890930in}}%
\pgfpathlineto{\pgfqpoint{0.688230in}{0.883673in}}%
\pgfpathlineto{\pgfqpoint{0.680882in}{0.876407in}}%
\pgfpathlineto{\pgfqpoint{0.678591in}{0.860364in}}%
\pgfpathlineto{\pgfqpoint{0.674871in}{0.851614in}}%
\pgfpathlineto{\pgfqpoint{0.667035in}{0.844239in}}%
\pgfpathlineto{\pgfqpoint{0.659240in}{0.825061in}}%
\pgfpathlineto{\pgfqpoint{0.670200in}{0.815071in}}%
\pgfpathlineto{\pgfqpoint{0.668785in}{0.806836in}}%
\pgfpathlineto{\pgfqpoint{0.659798in}{0.801105in}}%
\pgfpathlineto{\pgfqpoint{0.653603in}{0.802117in}}%
\pgfpathlineto{\pgfqpoint{0.580470in}{0.812266in}}%
\pgfpathlineto{\pgfqpoint{0.526454in}{0.819928in}}%
\pgfpathlineto{\pgfqpoint{0.528816in}{0.828649in}}%
\pgfpathlineto{\pgfqpoint{0.525300in}{0.843128in}}%
\pgfpathlineto{\pgfqpoint{0.524945in}{0.857733in}}%
\pgfpathlineto{\pgfqpoint{0.522647in}{0.866279in}}%
\pgfpathlineto{\pgfqpoint{0.515568in}{0.878498in}}%
\pgfpathlineto{\pgfqpoint{0.495206in}{0.906691in}}%
\pgfpathlineto{\pgfqpoint{0.488537in}{0.910145in}}%
\pgfpathlineto{\pgfqpoint{0.479944in}{0.910090in}}%
\pgfpathlineto{\pgfqpoint{0.481908in}{0.918984in}}%
\pgfpathlineto{\pgfqpoint{0.477839in}{0.930106in}}%
\pgfpathlineto{\pgfqpoint{0.457853in}{0.935633in}}%
\pgfpathlineto{\pgfqpoint{0.445671in}{0.945869in}}%
\pgfpathlineto{\pgfqpoint{0.444661in}{0.952132in}}%
\pgfpathlineto{\pgfqpoint{0.430633in}{0.967637in}}%
\pgfpathlineto{\pgfqpoint{0.417280in}{0.970604in}}%
\pgfpathlineto{\pgfqpoint{0.404908in}{0.978484in}}%
\pgfpathlineto{\pgfqpoint{0.388637in}{0.981210in}}%
\pgfpathlineto{\pgfqpoint{0.381685in}{0.991719in}}%
\pgfpathlineto{\pgfqpoint{0.388287in}{1.008354in}}%
\pgfpathlineto{\pgfqpoint{0.386268in}{1.012095in}}%
\pgfpathlineto{\pgfqpoint{0.391766in}{1.025962in}}%
\pgfpathlineto{\pgfqpoint{0.381986in}{1.033342in}}%
\pgfpathlineto{\pgfqpoint{0.385159in}{1.046721in}}%
\pgfpathlineto{\pgfqpoint{0.379935in}{1.050141in}}%
\pgfpathlineto{\pgfqpoint{0.375424in}{1.062787in}}%
\pgfpathlineto{\pgfqpoint{0.369981in}{1.066644in}}%
\pgfpathlineto{\pgfqpoint{0.369572in}{1.075778in}}%
\pgfpathlineto{\pgfqpoint{0.365307in}{1.082219in}}%
\pgfpathlineto{\pgfqpoint{0.358881in}{1.103911in}}%
\pgfpathlineto{\pgfqpoint{0.351850in}{1.114301in}}%
\pgfpathlineto{\pgfqpoint{0.353378in}{1.131939in}}%
\pgfpathlineto{\pgfqpoint{0.361647in}{1.133714in}}%
\pgfpathlineto{\pgfqpoint{0.367031in}{1.143280in}}%
\pgfpathlineto{\pgfqpoint{0.363778in}{1.153656in}}%
\pgfpathlineto{\pgfqpoint{0.354984in}{1.155388in}}%
\pgfpathlineto{\pgfqpoint{0.350596in}{1.160241in}}%
\pgfpathlineto{\pgfqpoint{0.343491in}{1.178092in}}%
\pgfpathlineto{\pgfqpoint{0.346812in}{1.184503in}}%
\pgfpathlineto{\pgfqpoint{0.344457in}{1.196440in}}%
\pgfpathlineto{\pgfqpoint{0.349672in}{1.211900in}}%
\pgfpathlineto{\pgfqpoint{0.355782in}{1.206207in}}%
\pgfpathlineto{\pgfqpoint{0.353009in}{1.199485in}}%
\pgfpathlineto{\pgfqpoint{0.363597in}{1.188830in}}%
\pgfpathlineto{\pgfqpoint{0.362921in}{1.204652in}}%
\pgfpathlineto{\pgfqpoint{0.358384in}{1.208896in}}%
\pgfpathlineto{\pgfqpoint{0.363570in}{1.222808in}}%
\pgfpathlineto{\pgfqpoint{0.383185in}{1.220593in}}%
\pgfpathlineto{\pgfqpoint{0.380545in}{1.225766in}}%
\pgfpathlineto{\pgfqpoint{0.367599in}{1.225257in}}%
\pgfpathlineto{\pgfqpoint{0.361441in}{1.233093in}}%
\pgfpathlineto{\pgfqpoint{0.353689in}{1.226135in}}%
\pgfpathlineto{\pgfqpoint{0.349576in}{1.214504in}}%
\pgfpathlineto{\pgfqpoint{0.338597in}{1.230154in}}%
\pgfpathlineto{\pgfqpoint{0.334369in}{1.233003in}}%
\pgfpathlineto{\pgfqpoint{0.335980in}{1.250036in}}%
\pgfpathlineto{\pgfqpoint{0.332634in}{1.260081in}}%
\pgfpathlineto{\pgfqpoint{0.326568in}{1.269517in}}%
\pgfpathlineto{\pgfqpoint{0.314145in}{1.298432in}}%
\pgfpathlineto{\pgfqpoint{0.318207in}{1.304826in}}%
\pgfpathlineto{\pgfqpoint{0.318159in}{1.325033in}}%
\pgfpathlineto{\pgfqpoint{0.324864in}{1.336303in}}%
\pgfpathlineto{\pgfqpoint{0.326401in}{1.353918in}}%
\pgfpathlineto{\pgfqpoint{0.320079in}{1.374099in}}%
\pgfpathlineto{\pgfqpoint{0.311692in}{1.386945in}}%
\pgfpathlineto{\pgfqpoint{0.313188in}{1.398536in}}%
\pgfpathlineto{\pgfqpoint{0.336740in}{1.426591in}}%
\pgfpathlineto{\pgfqpoint{0.337910in}{1.436161in}}%
\pgfpathlineto{\pgfqpoint{0.348523in}{1.454416in}}%
\pgfpathlineto{\pgfqpoint{0.349995in}{1.471719in}}%
\pgfpathlineto{\pgfqpoint{0.346583in}{1.476134in}}%
\pgfpathlineto{\pgfqpoint{0.352592in}{1.488751in}}%
\pgfpathclose%
\pgfusepath{fill}%
\end{pgfscope}%
\begin{pgfscope}%
\pgfpathrectangle{\pgfqpoint{0.100000in}{0.100000in}}{\pgfqpoint{2.862985in}{1.832778in}}%
\pgfusepath{clip}%
\pgfsetbuttcap%
\pgfsetmiterjoin%
\definecolor{currentfill}{rgb}{0.139946,0.788205,0.426221}%
\pgfsetfillcolor{currentfill}%
\pgfsetlinewidth{0.000000pt}%
\definecolor{currentstroke}{rgb}{0.000000,0.000000,0.000000}%
\pgfsetstrokecolor{currentstroke}%
\pgfsetstrokeopacity{0.000000}%
\pgfsetdash{}{0pt}%
\pgfpathmoveto{\pgfqpoint{2.186981in}{1.117748in}}%
\pgfpathlineto{\pgfqpoint{2.179011in}{1.195770in}}%
\pgfpathlineto{\pgfqpoint{2.169752in}{1.279623in}}%
\pgfpathlineto{\pgfqpoint{2.230390in}{1.288653in}}%
\pgfpathlineto{\pgfqpoint{2.246483in}{1.284465in}}%
\pgfpathlineto{\pgfqpoint{2.254202in}{1.279918in}}%
\pgfpathlineto{\pgfqpoint{2.263874in}{1.281179in}}%
\pgfpathlineto{\pgfqpoint{2.276627in}{1.273651in}}%
\pgfpathlineto{\pgfqpoint{2.300380in}{1.284855in}}%
\pgfpathlineto{\pgfqpoint{2.313489in}{1.285229in}}%
\pgfpathlineto{\pgfqpoint{2.328797in}{1.302307in}}%
\pgfpathlineto{\pgfqpoint{2.344372in}{1.312695in}}%
\pgfpathlineto{\pgfqpoint{2.365146in}{1.324661in}}%
\pgfpathlineto{\pgfqpoint{2.378519in}{1.241018in}}%
\pgfpathlineto{\pgfqpoint{2.372306in}{1.235646in}}%
\pgfpathlineto{\pgfqpoint{2.376317in}{1.230712in}}%
\pgfpathlineto{\pgfqpoint{2.377915in}{1.219896in}}%
\pgfpathlineto{\pgfqpoint{2.374803in}{1.209737in}}%
\pgfpathlineto{\pgfqpoint{2.374391in}{1.194887in}}%
\pgfpathlineto{\pgfqpoint{2.371468in}{1.175560in}}%
\pgfpathlineto{\pgfqpoint{2.356678in}{1.158322in}}%
\pgfpathlineto{\pgfqpoint{2.350452in}{1.154663in}}%
\pgfpathlineto{\pgfqpoint{2.345610in}{1.157813in}}%
\pgfpathlineto{\pgfqpoint{2.333627in}{1.141383in}}%
\pgfpathlineto{\pgfqpoint{2.334750in}{1.125564in}}%
\pgfpathlineto{\pgfqpoint{2.321406in}{1.129368in}}%
\pgfpathlineto{\pgfqpoint{2.315087in}{1.113569in}}%
\pgfpathlineto{\pgfqpoint{2.318281in}{1.102373in}}%
\pgfpathlineto{\pgfqpoint{2.313281in}{1.100719in}}%
\pgfpathlineto{\pgfqpoint{2.312581in}{1.091873in}}%
\pgfpathlineto{\pgfqpoint{2.300308in}{1.088305in}}%
\pgfpathlineto{\pgfqpoint{2.293930in}{1.095447in}}%
\pgfpathlineto{\pgfqpoint{2.285724in}{1.098217in}}%
\pgfpathlineto{\pgfqpoint{2.283796in}{1.105459in}}%
\pgfpathlineto{\pgfqpoint{2.276376in}{1.104185in}}%
\pgfpathlineto{\pgfqpoint{2.271515in}{1.097524in}}%
\pgfpathlineto{\pgfqpoint{2.264561in}{1.095181in}}%
\pgfpathlineto{\pgfqpoint{2.252277in}{1.099895in}}%
\pgfpathlineto{\pgfqpoint{2.245483in}{1.094286in}}%
\pgfpathlineto{\pgfqpoint{2.235827in}{1.100922in}}%
\pgfpathlineto{\pgfqpoint{2.217303in}{1.102959in}}%
\pgfpathlineto{\pgfqpoint{2.211715in}{1.115014in}}%
\pgfpathlineto{\pgfqpoint{2.204636in}{1.120355in}}%
\pgfpathlineto{\pgfqpoint{2.197773in}{1.116936in}}%
\pgfpathlineto{\pgfqpoint{2.186981in}{1.117748in}}%
\pgfpathclose%
\pgfusepath{fill}%
\end{pgfscope}%
\begin{pgfscope}%
\pgfpathrectangle{\pgfqpoint{0.100000in}{0.100000in}}{\pgfqpoint{2.862985in}{1.832778in}}%
\pgfusepath{clip}%
\pgfsetbuttcap%
\pgfsetmiterjoin%
\definecolor{currentfill}{rgb}{0.133795,0.786482,0.451442}%
\pgfsetfillcolor{currentfill}%
\pgfsetlinewidth{0.000000pt}%
\definecolor{currentstroke}{rgb}{0.000000,0.000000,0.000000}%
\pgfsetstrokecolor{currentstroke}%
\pgfsetstrokeopacity{0.000000}%
\pgfsetdash{}{0pt}%
\pgfpathmoveto{\pgfqpoint{1.986225in}{0.966856in}}%
\pgfpathlineto{\pgfqpoint{1.979121in}{0.972628in}}%
\pgfpathlineto{\pgfqpoint{1.973328in}{0.969882in}}%
\pgfpathlineto{\pgfqpoint{1.965922in}{0.983711in}}%
\pgfpathlineto{\pgfqpoint{1.969705in}{0.992402in}}%
\pgfpathlineto{\pgfqpoint{1.964252in}{1.002186in}}%
\pgfpathlineto{\pgfqpoint{1.963959in}{1.009888in}}%
\pgfpathlineto{\pgfqpoint{1.956591in}{1.012632in}}%
\pgfpathlineto{\pgfqpoint{1.953174in}{1.018436in}}%
\pgfpathlineto{\pgfqpoint{1.938768in}{1.025684in}}%
\pgfpathlineto{\pgfqpoint{1.926256in}{1.034614in}}%
\pgfpathlineto{\pgfqpoint{1.920442in}{1.041355in}}%
\pgfpathlineto{\pgfqpoint{1.920322in}{1.049578in}}%
\pgfpathlineto{\pgfqpoint{1.928102in}{1.065365in}}%
\pgfpathlineto{\pgfqpoint{1.930495in}{1.077440in}}%
\pgfpathlineto{\pgfqpoint{1.924157in}{1.084263in}}%
\pgfpathlineto{\pgfqpoint{1.915763in}{1.086834in}}%
\pgfpathlineto{\pgfqpoint{1.905579in}{1.081212in}}%
\pgfpathlineto{\pgfqpoint{1.901269in}{1.090863in}}%
\pgfpathlineto{\pgfqpoint{1.899094in}{1.103945in}}%
\pgfpathlineto{\pgfqpoint{1.884086in}{1.115607in}}%
\pgfpathlineto{\pgfqpoint{1.881087in}{1.120782in}}%
\pgfpathlineto{\pgfqpoint{1.867379in}{1.132498in}}%
\pgfpathlineto{\pgfqpoint{1.863083in}{1.141015in}}%
\pgfpathlineto{\pgfqpoint{1.859210in}{1.157906in}}%
\pgfpathlineto{\pgfqpoint{1.861865in}{1.172908in}}%
\pgfpathlineto{\pgfqpoint{1.865408in}{1.174999in}}%
\pgfpathlineto{\pgfqpoint{1.864771in}{1.187558in}}%
\pgfpathlineto{\pgfqpoint{1.874773in}{1.191288in}}%
\pgfpathlineto{\pgfqpoint{1.877793in}{1.202558in}}%
\pgfpathlineto{\pgfqpoint{1.883547in}{1.210143in}}%
\pgfpathlineto{\pgfqpoint{1.883256in}{1.219794in}}%
\pgfpathlineto{\pgfqpoint{1.876124in}{1.227469in}}%
\pgfpathlineto{\pgfqpoint{1.877818in}{1.238212in}}%
\pgfpathlineto{\pgfqpoint{1.896306in}{1.242884in}}%
\pgfpathlineto{\pgfqpoint{1.910490in}{1.251413in}}%
\pgfpathlineto{\pgfqpoint{1.911979in}{1.262151in}}%
\pgfpathlineto{\pgfqpoint{1.916914in}{1.265532in}}%
\pgfpathlineto{\pgfqpoint{1.918802in}{1.276813in}}%
\pgfpathlineto{\pgfqpoint{1.917285in}{1.284265in}}%
\pgfpathlineto{\pgfqpoint{1.907588in}{1.290454in}}%
\pgfpathlineto{\pgfqpoint{1.903684in}{1.299681in}}%
\pgfpathlineto{\pgfqpoint{1.894114in}{1.308600in}}%
\pgfpathlineto{\pgfqpoint{1.972758in}{1.311952in}}%
\pgfpathlineto{\pgfqpoint{2.025520in}{1.315700in}}%
\pgfpathlineto{\pgfqpoint{2.024554in}{1.304601in}}%
\pgfpathlineto{\pgfqpoint{2.033546in}{1.289293in}}%
\pgfpathlineto{\pgfqpoint{2.037322in}{1.276213in}}%
\pgfpathlineto{\pgfqpoint{2.041818in}{1.268787in}}%
\pgfpathlineto{\pgfqpoint{2.046726in}{1.207597in}}%
\pgfpathlineto{\pgfqpoint{2.053542in}{1.120373in}}%
\pgfpathlineto{\pgfqpoint{2.048283in}{1.107118in}}%
\pgfpathlineto{\pgfqpoint{2.055859in}{1.096224in}}%
\pgfpathlineto{\pgfqpoint{2.057971in}{1.085818in}}%
\pgfpathlineto{\pgfqpoint{2.047524in}{1.064241in}}%
\pgfpathlineto{\pgfqpoint{2.048140in}{1.061261in}}%
\pgfpathlineto{\pgfqpoint{2.037387in}{1.048893in}}%
\pgfpathlineto{\pgfqpoint{2.040336in}{1.045116in}}%
\pgfpathlineto{\pgfqpoint{2.035862in}{1.037407in}}%
\pgfpathlineto{\pgfqpoint{2.037013in}{1.021868in}}%
\pgfpathlineto{\pgfqpoint{2.031580in}{1.012333in}}%
\pgfpathlineto{\pgfqpoint{2.035999in}{1.001067in}}%
\pgfpathlineto{\pgfqpoint{2.017483in}{0.994944in}}%
\pgfpathlineto{\pgfqpoint{2.015777in}{0.988286in}}%
\pgfpathlineto{\pgfqpoint{2.020833in}{0.979847in}}%
\pgfpathlineto{\pgfqpoint{2.018506in}{0.974346in}}%
\pgfpathlineto{\pgfqpoint{1.998652in}{0.981141in}}%
\pgfpathlineto{\pgfqpoint{1.992101in}{0.981825in}}%
\pgfpathlineto{\pgfqpoint{1.983934in}{0.971495in}}%
\pgfpathlineto{\pgfqpoint{1.986225in}{0.966856in}}%
\pgfpathclose%
\pgfusepath{fill}%
\end{pgfscope}%
\begin{pgfscope}%
\pgfpathrectangle{\pgfqpoint{0.100000in}{0.100000in}}{\pgfqpoint{2.862985in}{1.832778in}}%
\pgfusepath{clip}%
\pgfsetbuttcap%
\pgfsetmiterjoin%
\definecolor{currentfill}{rgb}{0.076894,0.553633,0.076894}%
\pgfsetfillcolor{currentfill}%
\pgfsetlinewidth{0.000000pt}%
\definecolor{currentstroke}{rgb}{0.000000,0.000000,0.000000}%
\pgfsetstrokecolor{currentstroke}%
\pgfsetstrokeopacity{0.000000}%
\pgfsetdash{}{0pt}%
\pgfpathmoveto{\pgfqpoint{2.564161in}{1.156689in}}%
\pgfpathlineto{\pgfqpoint{2.558697in}{1.164803in}}%
\pgfpathlineto{\pgfqpoint{2.561780in}{1.169343in}}%
\pgfpathlineto{\pgfqpoint{2.569331in}{1.164244in}}%
\pgfpathlineto{\pgfqpoint{2.564161in}{1.156689in}}%
\pgfpathclose%
\pgfusepath{fill}%
\end{pgfscope}%
\begin{pgfscope}%
\pgfpathrectangle{\pgfqpoint{0.100000in}{0.100000in}}{\pgfqpoint{2.862985in}{1.832778in}}%
\pgfusepath{clip}%
\pgfsetbuttcap%
\pgfsetmiterjoin%
\definecolor{currentfill}{rgb}{0.179931,0.799400,0.262284}%
\pgfsetfillcolor{currentfill}%
\pgfsetlinewidth{0.000000pt}%
\definecolor{currentstroke}{rgb}{0.000000,0.000000,0.000000}%
\pgfsetstrokecolor{currentstroke}%
\pgfsetstrokeopacity{0.000000}%
\pgfsetdash{}{0pt}%
\pgfpathmoveto{\pgfqpoint{2.611873in}{1.226505in}}%
\pgfpathlineto{\pgfqpoint{2.615138in}{1.233404in}}%
\pgfpathlineto{\pgfqpoint{2.628338in}{1.234894in}}%
\pgfpathlineto{\pgfqpoint{2.626223in}{1.229031in}}%
\pgfpathlineto{\pgfqpoint{2.621868in}{1.221542in}}%
\pgfpathlineto{\pgfqpoint{2.624819in}{1.212619in}}%
\pgfpathlineto{\pgfqpoint{2.636465in}{1.201927in}}%
\pgfpathlineto{\pgfqpoint{2.639168in}{1.190668in}}%
\pgfpathlineto{\pgfqpoint{2.652585in}{1.176648in}}%
\pgfpathlineto{\pgfqpoint{2.657849in}{1.177248in}}%
\pgfpathlineto{\pgfqpoint{2.664406in}{1.156201in}}%
\pgfpathlineto{\pgfqpoint{2.663331in}{1.155991in}}%
\pgfpathlineto{\pgfqpoint{2.662135in}{1.155756in}}%
\pgfpathlineto{\pgfqpoint{2.632891in}{1.150158in}}%
\pgfpathlineto{\pgfqpoint{2.617245in}{1.205802in}}%
\pgfpathlineto{\pgfqpoint{2.611873in}{1.226505in}}%
\pgfpathclose%
\pgfusepath{fill}%
\end{pgfscope}%
\begin{pgfscope}%
\pgfpathrectangle{\pgfqpoint{0.100000in}{0.100000in}}{\pgfqpoint{2.862985in}{1.832778in}}%
\pgfusepath{clip}%
\pgfsetbuttcap%
\pgfsetmiterjoin%
\definecolor{currentfill}{rgb}{0.124506,0.772364,0.982314}%
\pgfsetfillcolor{currentfill}%
\pgfsetlinewidth{0.000000pt}%
\definecolor{currentstroke}{rgb}{0.000000,0.000000,0.000000}%
\pgfsetstrokecolor{currentstroke}%
\pgfsetstrokeopacity{0.000000}%
\pgfsetdash{}{0pt}%
\pgfpathmoveto{\pgfqpoint{2.338724in}{1.037893in}}%
\pgfpathlineto{\pgfqpoint{2.329483in}{1.038233in}}%
\pgfpathlineto{\pgfqpoint{2.320968in}{1.044098in}}%
\pgfpathlineto{\pgfqpoint{2.309465in}{1.061926in}}%
\pgfpathlineto{\pgfqpoint{2.299672in}{1.071367in}}%
\pgfpathlineto{\pgfqpoint{2.302231in}{1.078659in}}%
\pgfpathlineto{\pgfqpoint{2.300308in}{1.088305in}}%
\pgfpathlineto{\pgfqpoint{2.312581in}{1.091873in}}%
\pgfpathlineto{\pgfqpoint{2.313281in}{1.100719in}}%
\pgfpathlineto{\pgfqpoint{2.318281in}{1.102373in}}%
\pgfpathlineto{\pgfqpoint{2.315087in}{1.113569in}}%
\pgfpathlineto{\pgfqpoint{2.321406in}{1.129368in}}%
\pgfpathlineto{\pgfqpoint{2.334750in}{1.125564in}}%
\pgfpathlineto{\pgfqpoint{2.333627in}{1.141383in}}%
\pgfpathlineto{\pgfqpoint{2.345610in}{1.157813in}}%
\pgfpathlineto{\pgfqpoint{2.350452in}{1.154663in}}%
\pgfpathlineto{\pgfqpoint{2.356678in}{1.158322in}}%
\pgfpathlineto{\pgfqpoint{2.371468in}{1.175560in}}%
\pgfpathlineto{\pgfqpoint{2.374391in}{1.194887in}}%
\pgfpathlineto{\pgfqpoint{2.374803in}{1.209737in}}%
\pgfpathlineto{\pgfqpoint{2.377915in}{1.219896in}}%
\pgfpathlineto{\pgfqpoint{2.376317in}{1.230712in}}%
\pgfpathlineto{\pgfqpoint{2.372306in}{1.235646in}}%
\pgfpathlineto{\pgfqpoint{2.378519in}{1.241018in}}%
\pgfpathlineto{\pgfqpoint{2.387526in}{1.184275in}}%
\pgfpathlineto{\pgfqpoint{2.437178in}{1.192485in}}%
\pgfpathlineto{\pgfqpoint{2.442325in}{1.160097in}}%
\pgfpathlineto{\pgfqpoint{2.460300in}{1.181474in}}%
\pgfpathlineto{\pgfqpoint{2.464536in}{1.179343in}}%
\pgfpathlineto{\pgfqpoint{2.470767in}{1.191730in}}%
\pgfpathlineto{\pgfqpoint{2.479399in}{1.187846in}}%
\pgfpathlineto{\pgfqpoint{2.486925in}{1.188584in}}%
\pgfpathlineto{\pgfqpoint{2.491891in}{1.197195in}}%
\pgfpathlineto{\pgfqpoint{2.499098in}{1.201985in}}%
\pgfpathlineto{\pgfqpoint{2.509119in}{1.197766in}}%
\pgfpathlineto{\pgfqpoint{2.515712in}{1.199224in}}%
\pgfpathlineto{\pgfqpoint{2.525158in}{1.182877in}}%
\pgfpathlineto{\pgfqpoint{2.522427in}{1.170500in}}%
\pgfpathlineto{\pgfqpoint{2.497755in}{1.184435in}}%
\pgfpathlineto{\pgfqpoint{2.493134in}{1.172996in}}%
\pgfpathlineto{\pgfqpoint{2.494659in}{1.167732in}}%
\pgfpathlineto{\pgfqpoint{2.484159in}{1.149941in}}%
\pgfpathlineto{\pgfqpoint{2.479203in}{1.146402in}}%
\pgfpathlineto{\pgfqpoint{2.476965in}{1.138556in}}%
\pgfpathlineto{\pgfqpoint{2.468517in}{1.139377in}}%
\pgfpathlineto{\pgfqpoint{2.462449in}{1.117963in}}%
\pgfpathlineto{\pgfqpoint{2.459053in}{1.113049in}}%
\pgfpathlineto{\pgfqpoint{2.450326in}{1.114688in}}%
\pgfpathlineto{\pgfqpoint{2.441399in}{1.121442in}}%
\pgfpathlineto{\pgfqpoint{2.441091in}{1.111076in}}%
\pgfpathlineto{\pgfqpoint{2.432449in}{1.093654in}}%
\pgfpathlineto{\pgfqpoint{2.430300in}{1.081251in}}%
\pgfpathlineto{\pgfqpoint{2.423669in}{1.072991in}}%
\pgfpathlineto{\pgfqpoint{2.418643in}{1.059800in}}%
\pgfpathlineto{\pgfqpoint{2.418396in}{1.047589in}}%
\pgfpathlineto{\pgfqpoint{2.410633in}{1.045308in}}%
\pgfpathlineto{\pgfqpoint{2.401826in}{1.038395in}}%
\pgfpathlineto{\pgfqpoint{2.393347in}{1.037079in}}%
\pgfpathlineto{\pgfqpoint{2.391393in}{1.031228in}}%
\pgfpathlineto{\pgfqpoint{2.377730in}{1.025236in}}%
\pgfpathlineto{\pgfqpoint{2.370094in}{1.030341in}}%
\pgfpathlineto{\pgfqpoint{2.361560in}{1.020647in}}%
\pgfpathlineto{\pgfqpoint{2.346890in}{1.023505in}}%
\pgfpathlineto{\pgfqpoint{2.337893in}{1.033679in}}%
\pgfpathlineto{\pgfqpoint{2.338724in}{1.037893in}}%
\pgfpathclose%
\pgfusepath{fill}%
\end{pgfscope}%
\begin{pgfscope}%
\pgfpathrectangle{\pgfqpoint{0.100000in}{0.100000in}}{\pgfqpoint{2.862985in}{1.832778in}}%
\pgfusepath{clip}%
\pgfsetbuttcap%
\pgfsetmiterjoin%
\definecolor{currentfill}{rgb}{0.163014,0.734487,0.163014}%
\pgfsetfillcolor{currentfill}%
\pgfsetlinewidth{0.000000pt}%
\definecolor{currentstroke}{rgb}{0.000000,0.000000,0.000000}%
\pgfsetstrokecolor{currentstroke}%
\pgfsetstrokeopacity{0.000000}%
\pgfsetdash{}{0pt}%
\pgfpathmoveto{\pgfqpoint{2.662135in}{1.155756in}}%
\pgfpathlineto{\pgfqpoint{2.661760in}{1.144322in}}%
\pgfpathlineto{\pgfqpoint{2.657360in}{1.138701in}}%
\pgfpathlineto{\pgfqpoint{2.654551in}{1.126221in}}%
\pgfpathlineto{\pgfqpoint{2.641870in}{1.120469in}}%
\pgfpathlineto{\pgfqpoint{2.631229in}{1.118793in}}%
\pgfpathlineto{\pgfqpoint{2.634330in}{1.127000in}}%
\pgfpathlineto{\pgfqpoint{2.617947in}{1.133877in}}%
\pgfpathlineto{\pgfqpoint{2.604638in}{1.142487in}}%
\pgfpathlineto{\pgfqpoint{2.612878in}{1.155358in}}%
\pgfpathlineto{\pgfqpoint{2.606274in}{1.165363in}}%
\pgfpathlineto{\pgfqpoint{2.607051in}{1.173993in}}%
\pgfpathlineto{\pgfqpoint{2.602276in}{1.176362in}}%
\pgfpathlineto{\pgfqpoint{2.598402in}{1.190382in}}%
\pgfpathlineto{\pgfqpoint{2.607564in}{1.209388in}}%
\pgfpathlineto{\pgfqpoint{2.600681in}{1.212492in}}%
\pgfpathlineto{\pgfqpoint{2.598896in}{1.203115in}}%
\pgfpathlineto{\pgfqpoint{2.589070in}{1.200505in}}%
\pgfpathlineto{\pgfqpoint{2.590103in}{1.169655in}}%
\pgfpathlineto{\pgfqpoint{2.588306in}{1.159726in}}%
\pgfpathlineto{\pgfqpoint{2.593235in}{1.145575in}}%
\pgfpathlineto{\pgfqpoint{2.600818in}{1.138767in}}%
\pgfpathlineto{\pgfqpoint{2.597379in}{1.134492in}}%
\pgfpathlineto{\pgfqpoint{2.605119in}{1.128246in}}%
\pgfpathlineto{\pgfqpoint{2.607917in}{1.118107in}}%
\pgfpathlineto{\pgfqpoint{2.593730in}{1.126503in}}%
\pgfpathlineto{\pgfqpoint{2.584752in}{1.125412in}}%
\pgfpathlineto{\pgfqpoint{2.577780in}{1.134049in}}%
\pgfpathlineto{\pgfqpoint{2.570682in}{1.134890in}}%
\pgfpathlineto{\pgfqpoint{2.560594in}{1.130564in}}%
\pgfpathlineto{\pgfqpoint{2.556681in}{1.135955in}}%
\pgfpathlineto{\pgfqpoint{2.561823in}{1.147270in}}%
\pgfpathlineto{\pgfqpoint{2.564161in}{1.156689in}}%
\pgfpathlineto{\pgfqpoint{2.569331in}{1.164244in}}%
\pgfpathlineto{\pgfqpoint{2.561780in}{1.169343in}}%
\pgfpathlineto{\pgfqpoint{2.558697in}{1.164803in}}%
\pgfpathlineto{\pgfqpoint{2.551154in}{1.169411in}}%
\pgfpathlineto{\pgfqpoint{2.537796in}{1.172482in}}%
\pgfpathlineto{\pgfqpoint{2.539008in}{1.179234in}}%
\pgfpathlineto{\pgfqpoint{2.532953in}{1.183150in}}%
\pgfpathlineto{\pgfqpoint{2.525158in}{1.182877in}}%
\pgfpathlineto{\pgfqpoint{2.515712in}{1.199224in}}%
\pgfpathlineto{\pgfqpoint{2.509119in}{1.197766in}}%
\pgfpathlineto{\pgfqpoint{2.499098in}{1.201985in}}%
\pgfpathlineto{\pgfqpoint{2.491891in}{1.197195in}}%
\pgfpathlineto{\pgfqpoint{2.486925in}{1.188584in}}%
\pgfpathlineto{\pgfqpoint{2.479399in}{1.187846in}}%
\pgfpathlineto{\pgfqpoint{2.470767in}{1.191730in}}%
\pgfpathlineto{\pgfqpoint{2.464536in}{1.179343in}}%
\pgfpathlineto{\pgfqpoint{2.460300in}{1.181474in}}%
\pgfpathlineto{\pgfqpoint{2.442325in}{1.160097in}}%
\pgfpathlineto{\pgfqpoint{2.437178in}{1.192485in}}%
\pgfpathlineto{\pgfqpoint{2.502868in}{1.204659in}}%
\pgfpathlineto{\pgfqpoint{2.532316in}{1.209933in}}%
\pgfpathlineto{\pgfqpoint{2.575158in}{1.218561in}}%
\pgfpathlineto{\pgfqpoint{2.611873in}{1.226505in}}%
\pgfpathlineto{\pgfqpoint{2.617245in}{1.205802in}}%
\pgfpathlineto{\pgfqpoint{2.632891in}{1.150158in}}%
\pgfpathlineto{\pgfqpoint{2.662135in}{1.155756in}}%
\pgfpathclose%
\pgfusepath{fill}%
\end{pgfscope}%
\begin{pgfscope}%
\pgfpathrectangle{\pgfqpoint{0.100000in}{0.100000in}}{\pgfqpoint{2.862985in}{1.832778in}}%
\pgfusepath{clip}%
\pgfsetbuttcap%
\pgfsetmiterjoin%
\definecolor{currentfill}{rgb}{0.153787,0.715110,0.153787}%
\pgfsetfillcolor{currentfill}%
\pgfsetlinewidth{0.000000pt}%
\definecolor{currentstroke}{rgb}{0.000000,0.000000,0.000000}%
\pgfsetstrokecolor{currentstroke}%
\pgfsetstrokeopacity{0.000000}%
\pgfsetdash{}{0pt}%
\pgfpathmoveto{\pgfqpoint{1.339821in}{0.971059in}}%
\pgfpathlineto{\pgfqpoint{1.291895in}{0.975633in}}%
\pgfpathlineto{\pgfqpoint{1.242175in}{0.980039in}}%
\pgfpathlineto{\pgfqpoint{1.184722in}{0.986032in}}%
\pgfpathlineto{\pgfqpoint{1.099362in}{0.996184in}}%
\pgfpathlineto{\pgfqpoint{1.069080in}{1.000850in}}%
\pgfpathlineto{\pgfqpoint{0.990915in}{1.012202in}}%
\pgfpathlineto{\pgfqpoint{1.002328in}{1.083837in}}%
\pgfpathlineto{\pgfqpoint{1.002613in}{1.089632in}}%
\pgfpathlineto{\pgfqpoint{1.013608in}{1.158753in}}%
\pgfpathlineto{\pgfqpoint{1.021820in}{1.211396in}}%
\pgfpathlineto{\pgfqpoint{1.029571in}{1.260214in}}%
\pgfpathlineto{\pgfqpoint{1.082522in}{1.252613in}}%
\pgfpathlineto{\pgfqpoint{1.131892in}{1.245527in}}%
\pgfpathlineto{\pgfqpoint{1.222647in}{1.234326in}}%
\pgfpathlineto{\pgfqpoint{1.264283in}{1.230524in}}%
\pgfpathlineto{\pgfqpoint{1.330243in}{1.224124in}}%
\pgfpathlineto{\pgfqpoint{1.358777in}{1.221808in}}%
\pgfpathlineto{\pgfqpoint{1.353739in}{1.159333in}}%
\pgfpathlineto{\pgfqpoint{1.349188in}{1.099180in}}%
\pgfpathlineto{\pgfqpoint{1.345525in}{1.050201in}}%
\pgfpathlineto{\pgfqpoint{1.339821in}{0.971059in}}%
\pgfpathclose%
\pgfusepath{fill}%
\end{pgfscope}%
\begin{pgfscope}%
\pgfpathrectangle{\pgfqpoint{0.100000in}{0.100000in}}{\pgfqpoint{2.862985in}{1.832778in}}%
\pgfusepath{clip}%
\pgfsetbuttcap%
\pgfsetmiterjoin%
\definecolor{currentfill}{rgb}{0.146159,0.776424,0.979239}%
\pgfsetfillcolor{currentfill}%
\pgfsetlinewidth{0.000000pt}%
\definecolor{currentstroke}{rgb}{0.000000,0.000000,0.000000}%
\pgfsetstrokecolor{currentstroke}%
\pgfsetstrokeopacity{0.000000}%
\pgfsetdash{}{0pt}%
\pgfpathmoveto{\pgfqpoint{1.973847in}{0.935544in}}%
\pgfpathlineto{\pgfqpoint{1.975463in}{0.942786in}}%
\pgfpathlineto{\pgfqpoint{1.983837in}{0.941156in}}%
\pgfpathlineto{\pgfqpoint{1.986225in}{0.966856in}}%
\pgfpathlineto{\pgfqpoint{1.983934in}{0.971495in}}%
\pgfpathlineto{\pgfqpoint{1.992101in}{0.981825in}}%
\pgfpathlineto{\pgfqpoint{1.998652in}{0.981141in}}%
\pgfpathlineto{\pgfqpoint{2.018506in}{0.974346in}}%
\pgfpathlineto{\pgfqpoint{2.020833in}{0.979847in}}%
\pgfpathlineto{\pgfqpoint{2.015777in}{0.988286in}}%
\pgfpathlineto{\pgfqpoint{2.017483in}{0.994944in}}%
\pgfpathlineto{\pgfqpoint{2.035999in}{1.001067in}}%
\pgfpathlineto{\pgfqpoint{2.031580in}{1.012333in}}%
\pgfpathlineto{\pgfqpoint{2.037013in}{1.021868in}}%
\pgfpathlineto{\pgfqpoint{2.047127in}{1.027220in}}%
\pgfpathlineto{\pgfqpoint{2.068349in}{1.032520in}}%
\pgfpathlineto{\pgfqpoint{2.081739in}{1.024520in}}%
\pgfpathlineto{\pgfqpoint{2.085801in}{1.032321in}}%
\pgfpathlineto{\pgfqpoint{2.096919in}{1.037771in}}%
\pgfpathlineto{\pgfqpoint{2.099789in}{1.032965in}}%
\pgfpathlineto{\pgfqpoint{2.111209in}{1.036783in}}%
\pgfpathlineto{\pgfqpoint{2.110489in}{1.043318in}}%
\pgfpathlineto{\pgfqpoint{2.120757in}{1.050797in}}%
\pgfpathlineto{\pgfqpoint{2.126817in}{1.043031in}}%
\pgfpathlineto{\pgfqpoint{2.134790in}{1.042250in}}%
\pgfpathlineto{\pgfqpoint{2.140088in}{1.047311in}}%
\pgfpathlineto{\pgfqpoint{2.139498in}{1.054509in}}%
\pgfpathlineto{\pgfqpoint{2.144002in}{1.061656in}}%
\pgfpathlineto{\pgfqpoint{2.150037in}{1.063205in}}%
\pgfpathlineto{\pgfqpoint{2.152460in}{1.072642in}}%
\pgfpathlineto{\pgfqpoint{2.161237in}{1.080801in}}%
\pgfpathlineto{\pgfqpoint{2.158595in}{1.088922in}}%
\pgfpathlineto{\pgfqpoint{2.167134in}{1.092965in}}%
\pgfpathlineto{\pgfqpoint{2.172839in}{1.090478in}}%
\pgfpathlineto{\pgfqpoint{2.181265in}{1.096795in}}%
\pgfpathlineto{\pgfqpoint{2.188795in}{1.098441in}}%
\pgfpathlineto{\pgfqpoint{2.188822in}{1.104993in}}%
\pgfpathlineto{\pgfqpoint{2.183545in}{1.114085in}}%
\pgfpathlineto{\pgfqpoint{2.186981in}{1.117748in}}%
\pgfpathlineto{\pgfqpoint{2.197773in}{1.116936in}}%
\pgfpathlineto{\pgfqpoint{2.204636in}{1.120355in}}%
\pgfpathlineto{\pgfqpoint{2.211715in}{1.115014in}}%
\pgfpathlineto{\pgfqpoint{2.217303in}{1.102959in}}%
\pgfpathlineto{\pgfqpoint{2.235827in}{1.100922in}}%
\pgfpathlineto{\pgfqpoint{2.245483in}{1.094286in}}%
\pgfpathlineto{\pgfqpoint{2.252277in}{1.099895in}}%
\pgfpathlineto{\pgfqpoint{2.264561in}{1.095181in}}%
\pgfpathlineto{\pgfqpoint{2.271515in}{1.097524in}}%
\pgfpathlineto{\pgfqpoint{2.276376in}{1.104185in}}%
\pgfpathlineto{\pgfqpoint{2.283796in}{1.105459in}}%
\pgfpathlineto{\pgfqpoint{2.285724in}{1.098217in}}%
\pgfpathlineto{\pgfqpoint{2.293930in}{1.095447in}}%
\pgfpathlineto{\pgfqpoint{2.300308in}{1.088305in}}%
\pgfpathlineto{\pgfqpoint{2.302231in}{1.078659in}}%
\pgfpathlineto{\pgfqpoint{2.299672in}{1.071367in}}%
\pgfpathlineto{\pgfqpoint{2.309465in}{1.061926in}}%
\pgfpathlineto{\pgfqpoint{2.320968in}{1.044098in}}%
\pgfpathlineto{\pgfqpoint{2.329483in}{1.038233in}}%
\pgfpathlineto{\pgfqpoint{2.338724in}{1.037893in}}%
\pgfpathlineto{\pgfqpoint{2.321661in}{1.018314in}}%
\pgfpathlineto{\pgfqpoint{2.304865in}{1.006463in}}%
\pgfpathlineto{\pgfqpoint{2.298683in}{0.997053in}}%
\pgfpathlineto{\pgfqpoint{2.298788in}{0.991945in}}%
\pgfpathlineto{\pgfqpoint{2.289691in}{0.988030in}}%
\pgfpathlineto{\pgfqpoint{2.287084in}{0.980665in}}%
\pgfpathlineto{\pgfqpoint{2.268132in}{0.973246in}}%
\pgfpathlineto{\pgfqpoint{2.261414in}{0.968427in}}%
\pgfpathlineto{\pgfqpoint{2.260505in}{0.967399in}}%
\pgfpathlineto{\pgfqpoint{2.205991in}{0.962230in}}%
\pgfpathlineto{\pgfqpoint{2.173070in}{0.959461in}}%
\pgfpathlineto{\pgfqpoint{2.119049in}{0.956404in}}%
\pgfpathlineto{\pgfqpoint{2.051748in}{0.949746in}}%
\pgfpathlineto{\pgfqpoint{2.040629in}{0.951286in}}%
\pgfpathlineto{\pgfqpoint{2.042942in}{0.939936in}}%
\pgfpathlineto{\pgfqpoint{1.973847in}{0.935544in}}%
\pgfpathclose%
\pgfusepath{fill}%
\end{pgfscope}%
\begin{pgfscope}%
\pgfpathrectangle{\pgfqpoint{0.100000in}{0.100000in}}{\pgfqpoint{2.862985in}{1.832778in}}%
\pgfusepath{clip}%
\pgfsetbuttcap%
\pgfsetmiterjoin%
\definecolor{currentfill}{rgb}{0.163014,0.734487,0.163014}%
\pgfsetfillcolor{currentfill}%
\pgfsetlinewidth{0.000000pt}%
\definecolor{currentstroke}{rgb}{0.000000,0.000000,0.000000}%
\pgfsetstrokecolor{currentstroke}%
\pgfsetstrokeopacity{0.000000}%
\pgfsetdash{}{0pt}%
\pgfpathmoveto{\pgfqpoint{1.339821in}{0.971059in}}%
\pgfpathlineto{\pgfqpoint{1.345525in}{1.050201in}}%
\pgfpathlineto{\pgfqpoint{1.349188in}{1.099180in}}%
\pgfpathlineto{\pgfqpoint{1.353739in}{1.159333in}}%
\pgfpathlineto{\pgfqpoint{1.416835in}{1.154903in}}%
\pgfpathlineto{\pgfqpoint{1.496966in}{1.150526in}}%
\pgfpathlineto{\pgfqpoint{1.551466in}{1.148434in}}%
\pgfpathlineto{\pgfqpoint{1.605647in}{1.146774in}}%
\pgfpathlineto{\pgfqpoint{1.654686in}{1.146016in}}%
\pgfpathlineto{\pgfqpoint{1.677367in}{1.146251in}}%
\pgfpathlineto{\pgfqpoint{1.687350in}{1.138104in}}%
\pgfpathlineto{\pgfqpoint{1.695170in}{1.139746in}}%
\pgfpathlineto{\pgfqpoint{1.695414in}{1.132640in}}%
\pgfpathlineto{\pgfqpoint{1.689607in}{1.120348in}}%
\pgfpathlineto{\pgfqpoint{1.690237in}{1.112582in}}%
\pgfpathlineto{\pgfqpoint{1.695797in}{1.107461in}}%
\pgfpathlineto{\pgfqpoint{1.700905in}{1.096790in}}%
\pgfpathlineto{\pgfqpoint{1.711265in}{1.090665in}}%
\pgfpathlineto{\pgfqpoint{1.710919in}{1.050444in}}%
\pgfpathlineto{\pgfqpoint{1.711224in}{0.957954in}}%
\pgfpathlineto{\pgfqpoint{1.641779in}{0.958274in}}%
\pgfpathlineto{\pgfqpoint{1.568651in}{0.959554in}}%
\pgfpathlineto{\pgfqpoint{1.514818in}{0.961453in}}%
\pgfpathlineto{\pgfqpoint{1.469882in}{0.963187in}}%
\pgfpathlineto{\pgfqpoint{1.394179in}{0.967651in}}%
\pgfpathlineto{\pgfqpoint{1.339821in}{0.971059in}}%
\pgfpathclose%
\pgfusepath{fill}%
\end{pgfscope}%
\begin{pgfscope}%
\pgfpathrectangle{\pgfqpoint{0.100000in}{0.100000in}}{\pgfqpoint{2.862985in}{1.832778in}}%
\pgfusepath{clip}%
\pgfsetbuttcap%
\pgfsetmiterjoin%
\definecolor{currentfill}{rgb}{0.166090,0.740946,0.166090}%
\pgfsetfillcolor{currentfill}%
\pgfsetlinewidth{0.000000pt}%
\definecolor{currentstroke}{rgb}{0.000000,0.000000,0.000000}%
\pgfsetstrokecolor{currentstroke}%
\pgfsetstrokeopacity{0.000000}%
\pgfsetdash{}{0pt}%
\pgfpathmoveto{\pgfqpoint{2.361621in}{0.981538in}}%
\pgfpathlineto{\pgfqpoint{2.284707in}{0.970734in}}%
\pgfpathlineto{\pgfqpoint{2.261414in}{0.968427in}}%
\pgfpathlineto{\pgfqpoint{2.268132in}{0.973246in}}%
\pgfpathlineto{\pgfqpoint{2.287084in}{0.980665in}}%
\pgfpathlineto{\pgfqpoint{2.289691in}{0.988030in}}%
\pgfpathlineto{\pgfqpoint{2.298788in}{0.991945in}}%
\pgfpathlineto{\pgfqpoint{2.298683in}{0.997053in}}%
\pgfpathlineto{\pgfqpoint{2.304865in}{1.006463in}}%
\pgfpathlineto{\pgfqpoint{2.321661in}{1.018314in}}%
\pgfpathlineto{\pgfqpoint{2.338724in}{1.037893in}}%
\pgfpathlineto{\pgfqpoint{2.337893in}{1.033679in}}%
\pgfpathlineto{\pgfqpoint{2.346890in}{1.023505in}}%
\pgfpathlineto{\pgfqpoint{2.361560in}{1.020647in}}%
\pgfpathlineto{\pgfqpoint{2.370094in}{1.030341in}}%
\pgfpathlineto{\pgfqpoint{2.377730in}{1.025236in}}%
\pgfpathlineto{\pgfqpoint{2.391393in}{1.031228in}}%
\pgfpathlineto{\pgfqpoint{2.393347in}{1.037079in}}%
\pgfpathlineto{\pgfqpoint{2.401826in}{1.038395in}}%
\pgfpathlineto{\pgfqpoint{2.410633in}{1.045308in}}%
\pgfpathlineto{\pgfqpoint{2.418396in}{1.047589in}}%
\pgfpathlineto{\pgfqpoint{2.418643in}{1.059800in}}%
\pgfpathlineto{\pgfqpoint{2.423669in}{1.072991in}}%
\pgfpathlineto{\pgfqpoint{2.430300in}{1.081251in}}%
\pgfpathlineto{\pgfqpoint{2.432449in}{1.093654in}}%
\pgfpathlineto{\pgfqpoint{2.441091in}{1.111076in}}%
\pgfpathlineto{\pgfqpoint{2.441399in}{1.121442in}}%
\pgfpathlineto{\pgfqpoint{2.450326in}{1.114688in}}%
\pgfpathlineto{\pgfqpoint{2.459053in}{1.113049in}}%
\pgfpathlineto{\pgfqpoint{2.462449in}{1.117963in}}%
\pgfpathlineto{\pgfqpoint{2.468517in}{1.139377in}}%
\pgfpathlineto{\pgfqpoint{2.476965in}{1.138556in}}%
\pgfpathlineto{\pgfqpoint{2.479203in}{1.146402in}}%
\pgfpathlineto{\pgfqpoint{2.484159in}{1.149941in}}%
\pgfpathlineto{\pgfqpoint{2.494659in}{1.167732in}}%
\pgfpathlineto{\pgfqpoint{2.493134in}{1.172996in}}%
\pgfpathlineto{\pgfqpoint{2.497755in}{1.184435in}}%
\pgfpathlineto{\pgfqpoint{2.522427in}{1.170500in}}%
\pgfpathlineto{\pgfqpoint{2.525158in}{1.182877in}}%
\pgfpathlineto{\pgfqpoint{2.532953in}{1.183150in}}%
\pgfpathlineto{\pgfqpoint{2.539008in}{1.179234in}}%
\pgfpathlineto{\pgfqpoint{2.537796in}{1.172482in}}%
\pgfpathlineto{\pgfqpoint{2.551154in}{1.169411in}}%
\pgfpathlineto{\pgfqpoint{2.558697in}{1.164803in}}%
\pgfpathlineto{\pgfqpoint{2.564161in}{1.156689in}}%
\pgfpathlineto{\pgfqpoint{2.561823in}{1.147270in}}%
\pgfpathlineto{\pgfqpoint{2.557103in}{1.146498in}}%
\pgfpathlineto{\pgfqpoint{2.554365in}{1.132285in}}%
\pgfpathlineto{\pgfqpoint{2.560362in}{1.126728in}}%
\pgfpathlineto{\pgfqpoint{2.568804in}{1.131220in}}%
\pgfpathlineto{\pgfqpoint{2.576637in}{1.121734in}}%
\pgfpathlineto{\pgfqpoint{2.594137in}{1.120029in}}%
\pgfpathlineto{\pgfqpoint{2.597152in}{1.114573in}}%
\pgfpathlineto{\pgfqpoint{2.613348in}{1.109266in}}%
\pgfpathlineto{\pgfqpoint{2.611349in}{1.103004in}}%
\pgfpathlineto{\pgfqpoint{2.613456in}{1.085305in}}%
\pgfpathlineto{\pgfqpoint{2.619993in}{1.078622in}}%
\pgfpathlineto{\pgfqpoint{2.610347in}{1.075802in}}%
\pgfpathlineto{\pgfqpoint{2.611626in}{1.068246in}}%
\pgfpathlineto{\pgfqpoint{2.621916in}{1.061840in}}%
\pgfpathlineto{\pgfqpoint{2.622844in}{1.055521in}}%
\pgfpathlineto{\pgfqpoint{2.617034in}{1.050772in}}%
\pgfpathlineto{\pgfqpoint{2.605312in}{1.062022in}}%
\pgfpathlineto{\pgfqpoint{2.604153in}{1.053814in}}%
\pgfpathlineto{\pgfqpoint{2.613966in}{1.049911in}}%
\pgfpathlineto{\pgfqpoint{2.614948in}{1.045875in}}%
\pgfpathlineto{\pgfqpoint{2.628408in}{1.051209in}}%
\pgfpathlineto{\pgfqpoint{2.638722in}{1.052620in}}%
\pgfpathlineto{\pgfqpoint{2.649287in}{1.031303in}}%
\pgfpathlineto{\pgfqpoint{2.648108in}{1.031073in}}%
\pgfpathlineto{\pgfqpoint{2.643334in}{1.030086in}}%
\pgfpathlineto{\pgfqpoint{2.641927in}{1.029792in}}%
\pgfpathlineto{\pgfqpoint{2.640997in}{1.029610in}}%
\pgfpathlineto{\pgfqpoint{2.578091in}{1.016561in}}%
\pgfpathlineto{\pgfqpoint{2.521826in}{1.005045in}}%
\pgfpathlineto{\pgfqpoint{2.444033in}{0.991644in}}%
\pgfpathlineto{\pgfqpoint{2.377964in}{0.982918in}}%
\pgfpathlineto{\pgfqpoint{2.361621in}{0.981538in}}%
\pgfpathclose%
\pgfusepath{fill}%
\end{pgfscope}%
\begin{pgfscope}%
\pgfpathrectangle{\pgfqpoint{0.100000in}{0.100000in}}{\pgfqpoint{2.862985in}{1.832778in}}%
\pgfusepath{clip}%
\pgfsetbuttcap%
\pgfsetmiterjoin%
\definecolor{currentfill}{rgb}{0.166090,0.740946,0.166090}%
\pgfsetfillcolor{currentfill}%
\pgfsetlinewidth{0.000000pt}%
\definecolor{currentstroke}{rgb}{0.000000,0.000000,0.000000}%
\pgfsetstrokecolor{currentstroke}%
\pgfsetstrokeopacity{0.000000}%
\pgfsetdash{}{0pt}%
\pgfpathmoveto{\pgfqpoint{2.641870in}{1.120469in}}%
\pgfpathlineto{\pgfqpoint{2.654551in}{1.126221in}}%
\pgfpathlineto{\pgfqpoint{2.644424in}{1.096596in}}%
\pgfpathlineto{\pgfqpoint{2.637727in}{1.080908in}}%
\pgfpathlineto{\pgfqpoint{2.632902in}{1.089786in}}%
\pgfpathlineto{\pgfqpoint{2.641482in}{1.111041in}}%
\pgfpathlineto{\pgfqpoint{2.641870in}{1.120469in}}%
\pgfpathclose%
\pgfusepath{fill}%
\end{pgfscope}%
\begin{pgfscope}%
\pgfpathrectangle{\pgfqpoint{0.100000in}{0.100000in}}{\pgfqpoint{2.862985in}{1.832778in}}%
\pgfusepath{clip}%
\pgfsetbuttcap%
\pgfsetmiterjoin%
\definecolor{currentfill}{rgb}{0.167628,0.795955,0.312726}%
\pgfsetfillcolor{currentfill}%
\pgfsetlinewidth{0.000000pt}%
\definecolor{currentstroke}{rgb}{0.000000,0.000000,0.000000}%
\pgfsetstrokecolor{currentstroke}%
\pgfsetstrokeopacity{0.000000}%
\pgfsetdash{}{0pt}%
\pgfpathmoveto{\pgfqpoint{1.973847in}{0.935544in}}%
\pgfpathlineto{\pgfqpoint{1.970785in}{0.935101in}}%
\pgfpathlineto{\pgfqpoint{1.967896in}{0.934898in}}%
\pgfpathlineto{\pgfqpoint{1.969133in}{0.926031in}}%
\pgfpathlineto{\pgfqpoint{1.961682in}{0.917093in}}%
\pgfpathlineto{\pgfqpoint{1.960213in}{0.903107in}}%
\pgfpathlineto{\pgfqpoint{1.926904in}{0.900644in}}%
\pgfpathlineto{\pgfqpoint{1.929808in}{0.907211in}}%
\pgfpathlineto{\pgfqpoint{1.941818in}{0.919204in}}%
\pgfpathlineto{\pgfqpoint{1.942150in}{0.926155in}}%
\pgfpathlineto{\pgfqpoint{1.936833in}{0.932737in}}%
\pgfpathlineto{\pgfqpoint{1.887257in}{0.930116in}}%
\pgfpathlineto{\pgfqpoint{1.800580in}{0.927364in}}%
\pgfpathlineto{\pgfqpoint{1.749856in}{0.926440in}}%
\pgfpathlineto{\pgfqpoint{1.711515in}{0.926094in}}%
\pgfpathlineto{\pgfqpoint{1.711224in}{0.957954in}}%
\pgfpathlineto{\pgfqpoint{1.710919in}{1.050444in}}%
\pgfpathlineto{\pgfqpoint{1.711265in}{1.090665in}}%
\pgfpathlineto{\pgfqpoint{1.700905in}{1.096790in}}%
\pgfpathlineto{\pgfqpoint{1.695797in}{1.107461in}}%
\pgfpathlineto{\pgfqpoint{1.690237in}{1.112582in}}%
\pgfpathlineto{\pgfqpoint{1.689607in}{1.120348in}}%
\pgfpathlineto{\pgfqpoint{1.695414in}{1.132640in}}%
\pgfpathlineto{\pgfqpoint{1.695170in}{1.139746in}}%
\pgfpathlineto{\pgfqpoint{1.687350in}{1.138104in}}%
\pgfpathlineto{\pgfqpoint{1.677367in}{1.146251in}}%
\pgfpathlineto{\pgfqpoint{1.669364in}{1.160551in}}%
\pgfpathlineto{\pgfqpoint{1.662651in}{1.167151in}}%
\pgfpathlineto{\pgfqpoint{1.655637in}{1.183365in}}%
\pgfpathlineto{\pgfqpoint{1.728489in}{1.182229in}}%
\pgfpathlineto{\pgfqpoint{1.800906in}{1.184615in}}%
\pgfpathlineto{\pgfqpoint{1.847339in}{1.187293in}}%
\pgfpathlineto{\pgfqpoint{1.861865in}{1.172908in}}%
\pgfpathlineto{\pgfqpoint{1.859210in}{1.157906in}}%
\pgfpathlineto{\pgfqpoint{1.863083in}{1.141015in}}%
\pgfpathlineto{\pgfqpoint{1.867379in}{1.132498in}}%
\pgfpathlineto{\pgfqpoint{1.881087in}{1.120782in}}%
\pgfpathlineto{\pgfqpoint{1.884086in}{1.115607in}}%
\pgfpathlineto{\pgfqpoint{1.899094in}{1.103945in}}%
\pgfpathlineto{\pgfqpoint{1.901269in}{1.090863in}}%
\pgfpathlineto{\pgfqpoint{1.905579in}{1.081212in}}%
\pgfpathlineto{\pgfqpoint{1.915763in}{1.086834in}}%
\pgfpathlineto{\pgfqpoint{1.924157in}{1.084263in}}%
\pgfpathlineto{\pgfqpoint{1.930495in}{1.077440in}}%
\pgfpathlineto{\pgfqpoint{1.928102in}{1.065365in}}%
\pgfpathlineto{\pgfqpoint{1.920322in}{1.049578in}}%
\pgfpathlineto{\pgfqpoint{1.920442in}{1.041355in}}%
\pgfpathlineto{\pgfqpoint{1.926256in}{1.034614in}}%
\pgfpathlineto{\pgfqpoint{1.938768in}{1.025684in}}%
\pgfpathlineto{\pgfqpoint{1.953174in}{1.018436in}}%
\pgfpathlineto{\pgfqpoint{1.956591in}{1.012632in}}%
\pgfpathlineto{\pgfqpoint{1.963959in}{1.009888in}}%
\pgfpathlineto{\pgfqpoint{1.964252in}{1.002186in}}%
\pgfpathlineto{\pgfqpoint{1.969705in}{0.992402in}}%
\pgfpathlineto{\pgfqpoint{1.965922in}{0.983711in}}%
\pgfpathlineto{\pgfqpoint{1.973328in}{0.969882in}}%
\pgfpathlineto{\pgfqpoint{1.979121in}{0.972628in}}%
\pgfpathlineto{\pgfqpoint{1.986225in}{0.966856in}}%
\pgfpathlineto{\pgfqpoint{1.983837in}{0.941156in}}%
\pgfpathlineto{\pgfqpoint{1.975463in}{0.942786in}}%
\pgfpathlineto{\pgfqpoint{1.973847in}{0.935544in}}%
\pgfpathclose%
\pgfusepath{fill}%
\end{pgfscope}%
\begin{pgfscope}%
\pgfpathrectangle{\pgfqpoint{0.100000in}{0.100000in}}{\pgfqpoint{2.862985in}{1.832778in}}%
\pgfusepath{clip}%
\pgfsetbuttcap%
\pgfsetmiterjoin%
\definecolor{currentfill}{rgb}{0.183007,0.800261,0.249673}%
\pgfsetfillcolor{currentfill}%
\pgfsetlinewidth{0.000000pt}%
\definecolor{currentstroke}{rgb}{0.000000,0.000000,0.000000}%
\pgfsetstrokecolor{currentstroke}%
\pgfsetstrokeopacity{0.000000}%
\pgfsetdash{}{0pt}%
\pgfpathmoveto{\pgfqpoint{0.689569in}{0.941085in}}%
\pgfpathlineto{\pgfqpoint{0.695633in}{0.953995in}}%
\pgfpathlineto{\pgfqpoint{0.693973in}{0.973374in}}%
\pgfpathlineto{\pgfqpoint{0.696877in}{0.978886in}}%
\pgfpathlineto{\pgfqpoint{0.697164in}{1.003024in}}%
\pgfpathlineto{\pgfqpoint{0.699916in}{1.009976in}}%
\pgfpathlineto{\pgfqpoint{0.709613in}{1.011057in}}%
\pgfpathlineto{\pgfqpoint{0.718612in}{1.007971in}}%
\pgfpathlineto{\pgfqpoint{0.722531in}{0.999473in}}%
\pgfpathlineto{\pgfqpoint{0.728025in}{0.999798in}}%
\pgfpathlineto{\pgfqpoint{0.734852in}{1.009534in}}%
\pgfpathlineto{\pgfqpoint{0.744719in}{1.057523in}}%
\pgfpathlineto{\pgfqpoint{0.800868in}{1.045930in}}%
\pgfpathlineto{\pgfqpoint{0.833455in}{1.039504in}}%
\pgfpathlineto{\pgfqpoint{0.919991in}{1.024292in}}%
\pgfpathlineto{\pgfqpoint{0.943932in}{1.019484in}}%
\pgfpathlineto{\pgfqpoint{0.990915in}{1.012202in}}%
\pgfpathlineto{\pgfqpoint{0.981281in}{0.950172in}}%
\pgfpathlineto{\pgfqpoint{0.971262in}{0.885469in}}%
\pgfpathlineto{\pgfqpoint{0.959723in}{0.812676in}}%
\pgfpathlineto{\pgfqpoint{0.946747in}{0.729118in}}%
\pgfpathlineto{\pgfqpoint{0.936264in}{0.660490in}}%
\pgfpathlineto{\pgfqpoint{0.861176in}{0.672414in}}%
\pgfpathlineto{\pgfqpoint{0.828185in}{0.678064in}}%
\pgfpathlineto{\pgfqpoint{0.813447in}{0.686887in}}%
\pgfpathlineto{\pgfqpoint{0.717356in}{0.744864in}}%
\pgfpathlineto{\pgfqpoint{0.645244in}{0.788862in}}%
\pgfpathlineto{\pgfqpoint{0.647646in}{0.796660in}}%
\pgfpathlineto{\pgfqpoint{0.653603in}{0.802117in}}%
\pgfpathlineto{\pgfqpoint{0.659798in}{0.801105in}}%
\pgfpathlineto{\pgfqpoint{0.668785in}{0.806836in}}%
\pgfpathlineto{\pgfqpoint{0.670200in}{0.815071in}}%
\pgfpathlineto{\pgfqpoint{0.659240in}{0.825061in}}%
\pgfpathlineto{\pgfqpoint{0.667035in}{0.844239in}}%
\pgfpathlineto{\pgfqpoint{0.674871in}{0.851614in}}%
\pgfpathlineto{\pgfqpoint{0.678591in}{0.860364in}}%
\pgfpathlineto{\pgfqpoint{0.680882in}{0.876407in}}%
\pgfpathlineto{\pgfqpoint{0.688230in}{0.883673in}}%
\pgfpathlineto{\pgfqpoint{0.703654in}{0.890930in}}%
\pgfpathlineto{\pgfqpoint{0.703313in}{0.897324in}}%
\pgfpathlineto{\pgfqpoint{0.694699in}{0.905259in}}%
\pgfpathlineto{\pgfqpoint{0.693526in}{0.921617in}}%
\pgfpathlineto{\pgfqpoint{0.687595in}{0.933576in}}%
\pgfpathlineto{\pgfqpoint{0.689569in}{0.941085in}}%
\pgfpathclose%
\pgfusepath{fill}%
\end{pgfscope}%
\begin{pgfscope}%
\pgfpathrectangle{\pgfqpoint{0.100000in}{0.100000in}}{\pgfqpoint{2.862985in}{1.832778in}}%
\pgfusepath{clip}%
\pgfsetbuttcap%
\pgfsetmiterjoin%
\definecolor{currentfill}{rgb}{0.124567,0.783899,0.489273}%
\pgfsetfillcolor{currentfill}%
\pgfsetlinewidth{0.000000pt}%
\definecolor{currentstroke}{rgb}{0.000000,0.000000,0.000000}%
\pgfsetstrokecolor{currentstroke}%
\pgfsetstrokeopacity{0.000000}%
\pgfsetdash{}{0pt}%
\pgfpathmoveto{\pgfqpoint{1.719643in}{0.746495in}}%
\pgfpathlineto{\pgfqpoint{1.705724in}{0.750800in}}%
\pgfpathlineto{\pgfqpoint{1.687888in}{0.764459in}}%
\pgfpathlineto{\pgfqpoint{1.679966in}{0.767394in}}%
\pgfpathlineto{\pgfqpoint{1.674932in}{0.761495in}}%
\pgfpathlineto{\pgfqpoint{1.668575in}{0.761197in}}%
\pgfpathlineto{\pgfqpoint{1.660532in}{0.766205in}}%
\pgfpathlineto{\pgfqpoint{1.647911in}{0.759788in}}%
\pgfpathlineto{\pgfqpoint{1.643418in}{0.762948in}}%
\pgfpathlineto{\pgfqpoint{1.630999in}{0.755462in}}%
\pgfpathlineto{\pgfqpoint{1.617897in}{0.756843in}}%
\pgfpathlineto{\pgfqpoint{1.607791in}{0.761708in}}%
\pgfpathlineto{\pgfqpoint{1.600714in}{0.759877in}}%
\pgfpathlineto{\pgfqpoint{1.589386in}{0.766758in}}%
\pgfpathlineto{\pgfqpoint{1.583084in}{0.754622in}}%
\pgfpathlineto{\pgfqpoint{1.576640in}{0.765059in}}%
\pgfpathlineto{\pgfqpoint{1.563911in}{0.761010in}}%
\pgfpathlineto{\pgfqpoint{1.552778in}{0.770922in}}%
\pgfpathlineto{\pgfqpoint{1.543049in}{0.762933in}}%
\pgfpathlineto{\pgfqpoint{1.538177in}{0.770275in}}%
\pgfpathlineto{\pgfqpoint{1.531154in}{0.772659in}}%
\pgfpathlineto{\pgfqpoint{1.526877in}{0.779738in}}%
\pgfpathlineto{\pgfqpoint{1.517692in}{0.781753in}}%
\pgfpathlineto{\pgfqpoint{1.512392in}{0.776425in}}%
\pgfpathlineto{\pgfqpoint{1.503387in}{0.783330in}}%
\pgfpathlineto{\pgfqpoint{1.499189in}{0.781753in}}%
\pgfpathlineto{\pgfqpoint{1.484265in}{0.787349in}}%
\pgfpathlineto{\pgfqpoint{1.474917in}{0.787973in}}%
\pgfpathlineto{\pgfqpoint{1.470733in}{0.799856in}}%
\pgfpathlineto{\pgfqpoint{1.463240in}{0.798376in}}%
\pgfpathlineto{\pgfqpoint{1.454657in}{0.801308in}}%
\pgfpathlineto{\pgfqpoint{1.449005in}{0.799614in}}%
\pgfpathlineto{\pgfqpoint{1.436884in}{0.812959in}}%
\pgfpathlineto{\pgfqpoint{1.433507in}{0.812086in}}%
\pgfpathlineto{\pgfqpoint{1.436574in}{0.866207in}}%
\pgfpathlineto{\pgfqpoint{1.439916in}{0.933195in}}%
\pgfpathlineto{\pgfqpoint{1.385064in}{0.936263in}}%
\pgfpathlineto{\pgfqpoint{1.330936in}{0.940356in}}%
\pgfpathlineto{\pgfqpoint{1.289118in}{0.943987in}}%
\pgfpathlineto{\pgfqpoint{1.291895in}{0.975633in}}%
\pgfpathlineto{\pgfqpoint{1.339821in}{0.971059in}}%
\pgfpathlineto{\pgfqpoint{1.394179in}{0.967651in}}%
\pgfpathlineto{\pgfqpoint{1.469882in}{0.963187in}}%
\pgfpathlineto{\pgfqpoint{1.514818in}{0.961453in}}%
\pgfpathlineto{\pgfqpoint{1.568651in}{0.959554in}}%
\pgfpathlineto{\pgfqpoint{1.641779in}{0.958274in}}%
\pgfpathlineto{\pgfqpoint{1.711224in}{0.957954in}}%
\pgfpathlineto{\pgfqpoint{1.711515in}{0.926094in}}%
\pgfpathlineto{\pgfqpoint{1.715413in}{0.902092in}}%
\pgfpathlineto{\pgfqpoint{1.721468in}{0.857760in}}%
\pgfpathlineto{\pgfqpoint{1.720219in}{0.782052in}}%
\pgfpathlineto{\pgfqpoint{1.719643in}{0.746495in}}%
\pgfpathclose%
\pgfusepath{fill}%
\end{pgfscope}%
\begin{pgfscope}%
\pgfpathrectangle{\pgfqpoint{0.100000in}{0.100000in}}{\pgfqpoint{2.862985in}{1.832778in}}%
\pgfusepath{clip}%
\pgfsetbuttcap%
\pgfsetmiterjoin%
\definecolor{currentfill}{rgb}{0.047674,0.762368,0.804537}%
\pgfsetfillcolor{currentfill}%
\pgfsetlinewidth{0.000000pt}%
\definecolor{currentstroke}{rgb}{0.000000,0.000000,0.000000}%
\pgfsetstrokecolor{currentstroke}%
\pgfsetstrokeopacity{0.000000}%
\pgfsetdash{}{0pt}%
\pgfpathmoveto{\pgfqpoint{2.240820in}{0.864070in}}%
\pgfpathlineto{\pgfqpoint{2.240861in}{0.878097in}}%
\pgfpathlineto{\pgfqpoint{2.253051in}{0.883484in}}%
\pgfpathlineto{\pgfqpoint{2.253562in}{0.892095in}}%
\pgfpathlineto{\pgfqpoint{2.264476in}{0.902614in}}%
\pgfpathlineto{\pgfqpoint{2.278119in}{0.904768in}}%
\pgfpathlineto{\pgfqpoint{2.290113in}{0.916361in}}%
\pgfpathlineto{\pgfqpoint{2.302639in}{0.921543in}}%
\pgfpathlineto{\pgfqpoint{2.312040in}{0.937098in}}%
\pgfpathlineto{\pgfqpoint{2.320602in}{0.935969in}}%
\pgfpathlineto{\pgfqpoint{2.338697in}{0.950142in}}%
\pgfpathlineto{\pgfqpoint{2.348253in}{0.950409in}}%
\pgfpathlineto{\pgfqpoint{2.352288in}{0.961208in}}%
\pgfpathlineto{\pgfqpoint{2.359883in}{0.968727in}}%
\pgfpathlineto{\pgfqpoint{2.361621in}{0.981538in}}%
\pgfpathlineto{\pgfqpoint{2.377964in}{0.982918in}}%
\pgfpathlineto{\pgfqpoint{2.444033in}{0.991644in}}%
\pgfpathlineto{\pgfqpoint{2.521826in}{1.005045in}}%
\pgfpathlineto{\pgfqpoint{2.578091in}{1.016561in}}%
\pgfpathlineto{\pgfqpoint{2.640997in}{1.029610in}}%
\pgfpathlineto{\pgfqpoint{2.658973in}{1.004931in}}%
\pgfpathlineto{\pgfqpoint{2.649236in}{1.006506in}}%
\pgfpathlineto{\pgfqpoint{2.636737in}{1.003461in}}%
\pgfpathlineto{\pgfqpoint{2.624421in}{0.990884in}}%
\pgfpathlineto{\pgfqpoint{2.615570in}{0.991789in}}%
\pgfpathlineto{\pgfqpoint{2.614450in}{0.984319in}}%
\pgfpathlineto{\pgfqpoint{2.630449in}{0.990226in}}%
\pgfpathlineto{\pgfqpoint{2.646548in}{0.992662in}}%
\pgfpathlineto{\pgfqpoint{2.652526in}{0.989421in}}%
\pgfpathlineto{\pgfqpoint{2.660571in}{0.993116in}}%
\pgfpathlineto{\pgfqpoint{2.664722in}{0.990527in}}%
\pgfpathlineto{\pgfqpoint{2.668396in}{0.978199in}}%
\pgfpathlineto{\pgfqpoint{2.660750in}{0.974382in}}%
\pgfpathlineto{\pgfqpoint{2.655502in}{0.959350in}}%
\pgfpathlineto{\pgfqpoint{2.650000in}{0.953497in}}%
\pgfpathlineto{\pgfqpoint{2.633132in}{0.954777in}}%
\pgfpathlineto{\pgfqpoint{2.634039in}{0.963604in}}%
\pgfpathlineto{\pgfqpoint{2.624827in}{0.959748in}}%
\pgfpathlineto{\pgfqpoint{2.622823in}{0.952387in}}%
\pgfpathlineto{\pgfqpoint{2.629137in}{0.944761in}}%
\pgfpathlineto{\pgfqpoint{2.631282in}{0.931819in}}%
\pgfpathlineto{\pgfqpoint{2.620950in}{0.924444in}}%
\pgfpathlineto{\pgfqpoint{2.632125in}{0.921851in}}%
\pgfpathlineto{\pgfqpoint{2.637181in}{0.927273in}}%
\pgfpathlineto{\pgfqpoint{2.648362in}{0.927934in}}%
\pgfpathlineto{\pgfqpoint{2.642616in}{0.916224in}}%
\pgfpathlineto{\pgfqpoint{2.635579in}{0.910582in}}%
\pgfpathlineto{\pgfqpoint{2.614292in}{0.904933in}}%
\pgfpathlineto{\pgfqpoint{2.592436in}{0.885111in}}%
\pgfpathlineto{\pgfqpoint{2.578960in}{0.865580in}}%
\pgfpathlineto{\pgfqpoint{2.578880in}{0.857648in}}%
\pgfpathlineto{\pgfqpoint{2.573497in}{0.846702in}}%
\pgfpathlineto{\pgfqpoint{2.545869in}{0.839501in}}%
\pgfpathlineto{\pgfqpoint{2.480165in}{0.886600in}}%
\pgfpathlineto{\pgfqpoint{2.422304in}{0.878098in}}%
\pgfpathlineto{\pgfqpoint{2.421819in}{0.885949in}}%
\pgfpathlineto{\pgfqpoint{2.413024in}{0.894789in}}%
\pgfpathlineto{\pgfqpoint{2.406366in}{0.896952in}}%
\pgfpathlineto{\pgfqpoint{2.343499in}{0.890420in}}%
\pgfpathlineto{\pgfqpoint{2.329049in}{0.885519in}}%
\pgfpathlineto{\pgfqpoint{2.302958in}{0.872540in}}%
\pgfpathlineto{\pgfqpoint{2.280410in}{0.868948in}}%
\pgfpathlineto{\pgfqpoint{2.240820in}{0.864070in}}%
\pgfpathclose%
\pgfusepath{fill}%
\end{pgfscope}%
\begin{pgfscope}%
\pgfpathrectangle{\pgfqpoint{0.100000in}{0.100000in}}{\pgfqpoint{2.862985in}{1.832778in}}%
\pgfusepath{clip}%
\pgfsetbuttcap%
\pgfsetmiterjoin%
\definecolor{currentfill}{rgb}{0.059977,0.765813,0.754095}%
\pgfsetfillcolor{currentfill}%
\pgfsetlinewidth{0.000000pt}%
\definecolor{currentstroke}{rgb}{0.000000,0.000000,0.000000}%
\pgfsetstrokecolor{currentstroke}%
\pgfsetstrokeopacity{0.000000}%
\pgfsetdash{}{0pt}%
\pgfpathmoveto{\pgfqpoint{2.240820in}{0.864070in}}%
\pgfpathlineto{\pgfqpoint{2.175022in}{0.856843in}}%
\pgfpathlineto{\pgfqpoint{2.114830in}{0.851430in}}%
\pgfpathlineto{\pgfqpoint{2.052965in}{0.847429in}}%
\pgfpathlineto{\pgfqpoint{2.042346in}{0.845890in}}%
\pgfpathlineto{\pgfqpoint{2.000643in}{0.842683in}}%
\pgfpathlineto{\pgfqpoint{1.933848in}{0.838785in}}%
\pgfpathlineto{\pgfqpoint{1.945732in}{0.848649in}}%
\pgfpathlineto{\pgfqpoint{1.943020in}{0.859033in}}%
\pgfpathlineto{\pgfqpoint{1.945647in}{0.871619in}}%
\pgfpathlineto{\pgfqpoint{1.949660in}{0.877545in}}%
\pgfpathlineto{\pgfqpoint{1.949509in}{0.885789in}}%
\pgfpathlineto{\pgfqpoint{1.960195in}{0.890974in}}%
\pgfpathlineto{\pgfqpoint{1.960213in}{0.903107in}}%
\pgfpathlineto{\pgfqpoint{1.961682in}{0.917093in}}%
\pgfpathlineto{\pgfqpoint{1.969133in}{0.926031in}}%
\pgfpathlineto{\pgfqpoint{1.967896in}{0.934898in}}%
\pgfpathlineto{\pgfqpoint{1.970785in}{0.935101in}}%
\pgfpathlineto{\pgfqpoint{1.973847in}{0.935544in}}%
\pgfpathlineto{\pgfqpoint{2.042942in}{0.939936in}}%
\pgfpathlineto{\pgfqpoint{2.040629in}{0.951286in}}%
\pgfpathlineto{\pgfqpoint{2.051748in}{0.949746in}}%
\pgfpathlineto{\pgfqpoint{2.119049in}{0.956404in}}%
\pgfpathlineto{\pgfqpoint{2.173070in}{0.959461in}}%
\pgfpathlineto{\pgfqpoint{2.205991in}{0.962230in}}%
\pgfpathlineto{\pgfqpoint{2.260505in}{0.967399in}}%
\pgfpathlineto{\pgfqpoint{2.261414in}{0.968427in}}%
\pgfpathlineto{\pgfqpoint{2.284707in}{0.970734in}}%
\pgfpathlineto{\pgfqpoint{2.361621in}{0.981538in}}%
\pgfpathlineto{\pgfqpoint{2.359883in}{0.968727in}}%
\pgfpathlineto{\pgfqpoint{2.352288in}{0.961208in}}%
\pgfpathlineto{\pgfqpoint{2.348253in}{0.950409in}}%
\pgfpathlineto{\pgfqpoint{2.338697in}{0.950142in}}%
\pgfpathlineto{\pgfqpoint{2.320602in}{0.935969in}}%
\pgfpathlineto{\pgfqpoint{2.312040in}{0.937098in}}%
\pgfpathlineto{\pgfqpoint{2.302639in}{0.921543in}}%
\pgfpathlineto{\pgfqpoint{2.290113in}{0.916361in}}%
\pgfpathlineto{\pgfqpoint{2.278119in}{0.904768in}}%
\pgfpathlineto{\pgfqpoint{2.264476in}{0.902614in}}%
\pgfpathlineto{\pgfqpoint{2.253562in}{0.892095in}}%
\pgfpathlineto{\pgfqpoint{2.253051in}{0.883484in}}%
\pgfpathlineto{\pgfqpoint{2.240861in}{0.878097in}}%
\pgfpathlineto{\pgfqpoint{2.240820in}{0.864070in}}%
\pgfpathclose%
\pgfusepath{fill}%
\end{pgfscope}%
\begin{pgfscope}%
\pgfpathrectangle{\pgfqpoint{0.100000in}{0.100000in}}{\pgfqpoint{2.862985in}{1.832778in}}%
\pgfusepath{clip}%
\pgfsetbuttcap%
\pgfsetmiterjoin%
\definecolor{currentfill}{rgb}{0.112265,0.780454,0.539715}%
\pgfsetfillcolor{currentfill}%
\pgfsetlinewidth{0.000000pt}%
\definecolor{currentstroke}{rgb}{0.000000,0.000000,0.000000}%
\pgfsetstrokecolor{currentstroke}%
\pgfsetstrokeopacity{0.000000}%
\pgfsetdash{}{0pt}%
\pgfpathmoveto{\pgfqpoint{1.073883in}{0.669271in}}%
\pgfpathlineto{\pgfqpoint{1.069190in}{0.676811in}}%
\pgfpathlineto{\pgfqpoint{1.071132in}{0.683307in}}%
\pgfpathlineto{\pgfqpoint{1.104118in}{0.679215in}}%
\pgfpathlineto{\pgfqpoint{1.165452in}{0.672249in}}%
\pgfpathlineto{\pgfqpoint{1.209801in}{0.667891in}}%
\pgfpathlineto{\pgfqpoint{1.261021in}{0.662722in}}%
\pgfpathlineto{\pgfqpoint{1.263807in}{0.695029in}}%
\pgfpathlineto{\pgfqpoint{1.272520in}{0.777201in}}%
\pgfpathlineto{\pgfqpoint{1.278172in}{0.834767in}}%
\pgfpathlineto{\pgfqpoint{1.283057in}{0.889756in}}%
\pgfpathlineto{\pgfqpoint{1.287612in}{0.944071in}}%
\pgfpathlineto{\pgfqpoint{1.289118in}{0.943987in}}%
\pgfpathlineto{\pgfqpoint{1.330936in}{0.940356in}}%
\pgfpathlineto{\pgfqpoint{1.385064in}{0.936263in}}%
\pgfpathlineto{\pgfqpoint{1.439916in}{0.933195in}}%
\pgfpathlineto{\pgfqpoint{1.436574in}{0.866207in}}%
\pgfpathlineto{\pgfqpoint{1.433507in}{0.812086in}}%
\pgfpathlineto{\pgfqpoint{1.436884in}{0.812959in}}%
\pgfpathlineto{\pgfqpoint{1.449005in}{0.799614in}}%
\pgfpathlineto{\pgfqpoint{1.454657in}{0.801308in}}%
\pgfpathlineto{\pgfqpoint{1.463240in}{0.798376in}}%
\pgfpathlineto{\pgfqpoint{1.470733in}{0.799856in}}%
\pgfpathlineto{\pgfqpoint{1.474917in}{0.787973in}}%
\pgfpathlineto{\pgfqpoint{1.484265in}{0.787349in}}%
\pgfpathlineto{\pgfqpoint{1.499189in}{0.781753in}}%
\pgfpathlineto{\pgfqpoint{1.503387in}{0.783330in}}%
\pgfpathlineto{\pgfqpoint{1.512392in}{0.776425in}}%
\pgfpathlineto{\pgfqpoint{1.517692in}{0.781753in}}%
\pgfpathlineto{\pgfqpoint{1.526877in}{0.779738in}}%
\pgfpathlineto{\pgfqpoint{1.531154in}{0.772659in}}%
\pgfpathlineto{\pgfqpoint{1.538177in}{0.770275in}}%
\pgfpathlineto{\pgfqpoint{1.543049in}{0.762933in}}%
\pgfpathlineto{\pgfqpoint{1.552778in}{0.770922in}}%
\pgfpathlineto{\pgfqpoint{1.563911in}{0.761010in}}%
\pgfpathlineto{\pgfqpoint{1.576640in}{0.765059in}}%
\pgfpathlineto{\pgfqpoint{1.583084in}{0.754622in}}%
\pgfpathlineto{\pgfqpoint{1.589386in}{0.766758in}}%
\pgfpathlineto{\pgfqpoint{1.600714in}{0.759877in}}%
\pgfpathlineto{\pgfqpoint{1.607791in}{0.761708in}}%
\pgfpathlineto{\pgfqpoint{1.617897in}{0.756843in}}%
\pgfpathlineto{\pgfqpoint{1.630999in}{0.755462in}}%
\pgfpathlineto{\pgfqpoint{1.643418in}{0.762948in}}%
\pgfpathlineto{\pgfqpoint{1.647911in}{0.759788in}}%
\pgfpathlineto{\pgfqpoint{1.660532in}{0.766205in}}%
\pgfpathlineto{\pgfqpoint{1.668575in}{0.761197in}}%
\pgfpathlineto{\pgfqpoint{1.674932in}{0.761495in}}%
\pgfpathlineto{\pgfqpoint{1.679966in}{0.767394in}}%
\pgfpathlineto{\pgfqpoint{1.687888in}{0.764459in}}%
\pgfpathlineto{\pgfqpoint{1.705724in}{0.750800in}}%
\pgfpathlineto{\pgfqpoint{1.719643in}{0.746495in}}%
\pgfpathlineto{\pgfqpoint{1.725225in}{0.741221in}}%
\pgfpathlineto{\pgfqpoint{1.732212in}{0.744093in}}%
\pgfpathlineto{\pgfqpoint{1.742799in}{0.741894in}}%
\pgfpathlineto{\pgfqpoint{1.743006in}{0.708326in}}%
\pgfpathlineto{\pgfqpoint{1.743890in}{0.643403in}}%
\pgfpathlineto{\pgfqpoint{1.751243in}{0.637164in}}%
\pgfpathlineto{\pgfqpoint{1.757137in}{0.625668in}}%
\pgfpathlineto{\pgfqpoint{1.755065in}{0.617977in}}%
\pgfpathlineto{\pgfqpoint{1.759693in}{0.614721in}}%
\pgfpathlineto{\pgfqpoint{1.763454in}{0.600538in}}%
\pgfpathlineto{\pgfqpoint{1.770658in}{0.592942in}}%
\pgfpathlineto{\pgfqpoint{1.772272in}{0.576816in}}%
\pgfpathlineto{\pgfqpoint{1.771019in}{0.569988in}}%
\pgfpathlineto{\pgfqpoint{1.761227in}{0.551919in}}%
\pgfpathlineto{\pgfqpoint{1.760101in}{0.538927in}}%
\pgfpathlineto{\pgfqpoint{1.763423in}{0.536215in}}%
\pgfpathlineto{\pgfqpoint{1.763712in}{0.524834in}}%
\pgfpathlineto{\pgfqpoint{1.760335in}{0.517680in}}%
\pgfpathlineto{\pgfqpoint{1.755085in}{0.514985in}}%
\pgfpathlineto{\pgfqpoint{1.749944in}{0.505637in}}%
\pgfpathlineto{\pgfqpoint{1.756496in}{0.496583in}}%
\pgfpathlineto{\pgfqpoint{1.743777in}{0.496405in}}%
\pgfpathlineto{\pgfqpoint{1.709766in}{0.480844in}}%
\pgfpathlineto{\pgfqpoint{1.716259in}{0.490157in}}%
\pgfpathlineto{\pgfqpoint{1.708396in}{0.495178in}}%
\pgfpathlineto{\pgfqpoint{1.706753in}{0.503714in}}%
\pgfpathlineto{\pgfqpoint{1.691406in}{0.488845in}}%
\pgfpathlineto{\pgfqpoint{1.698217in}{0.478682in}}%
\pgfpathlineto{\pgfqpoint{1.688501in}{0.465742in}}%
\pgfpathlineto{\pgfqpoint{1.683273in}{0.466013in}}%
\pgfpathlineto{\pgfqpoint{1.678354in}{0.451920in}}%
\pgfpathlineto{\pgfqpoint{1.662798in}{0.440835in}}%
\pgfpathlineto{\pgfqpoint{1.648265in}{0.436840in}}%
\pgfpathlineto{\pgfqpoint{1.622917in}{0.426447in}}%
\pgfpathlineto{\pgfqpoint{1.620291in}{0.432244in}}%
\pgfpathlineto{\pgfqpoint{1.604501in}{0.426908in}}%
\pgfpathlineto{\pgfqpoint{1.614213in}{0.417822in}}%
\pgfpathlineto{\pgfqpoint{1.598894in}{0.409929in}}%
\pgfpathlineto{\pgfqpoint{1.582371in}{0.398011in}}%
\pgfpathlineto{\pgfqpoint{1.578403in}{0.403551in}}%
\pgfpathlineto{\pgfqpoint{1.564881in}{0.395248in}}%
\pgfpathlineto{\pgfqpoint{1.578135in}{0.392095in}}%
\pgfpathlineto{\pgfqpoint{1.568142in}{0.379039in}}%
\pgfpathlineto{\pgfqpoint{1.563264in}{0.382924in}}%
\pgfpathlineto{\pgfqpoint{1.556713in}{0.376684in}}%
\pgfpathlineto{\pgfqpoint{1.564878in}{0.371233in}}%
\pgfpathlineto{\pgfqpoint{1.560050in}{0.363265in}}%
\pgfpathlineto{\pgfqpoint{1.554117in}{0.344404in}}%
\pgfpathlineto{\pgfqpoint{1.545559in}{0.326238in}}%
\pgfpathlineto{\pgfqpoint{1.549401in}{0.314331in}}%
\pgfpathlineto{\pgfqpoint{1.555932in}{0.286258in}}%
\pgfpathlineto{\pgfqpoint{1.556520in}{0.274896in}}%
\pgfpathlineto{\pgfqpoint{1.566610in}{0.259994in}}%
\pgfpathlineto{\pgfqpoint{1.558829in}{0.260865in}}%
\pgfpathlineto{\pgfqpoint{1.551272in}{0.253347in}}%
\pgfpathlineto{\pgfqpoint{1.542420in}{0.259256in}}%
\pgfpathlineto{\pgfqpoint{1.539243in}{0.265150in}}%
\pgfpathlineto{\pgfqpoint{1.526645in}{0.267892in}}%
\pgfpathlineto{\pgfqpoint{1.507402in}{0.268229in}}%
\pgfpathlineto{\pgfqpoint{1.493224in}{0.279391in}}%
\pgfpathlineto{\pgfqpoint{1.480350in}{0.281256in}}%
\pgfpathlineto{\pgfqpoint{1.472540in}{0.290127in}}%
\pgfpathlineto{\pgfqpoint{1.456183in}{0.293697in}}%
\pgfpathlineto{\pgfqpoint{1.447239in}{0.322221in}}%
\pgfpathlineto{\pgfqpoint{1.438094in}{0.333648in}}%
\pgfpathlineto{\pgfqpoint{1.439648in}{0.344506in}}%
\pgfpathlineto{\pgfqpoint{1.433978in}{0.352440in}}%
\pgfpathlineto{\pgfqpoint{1.437536in}{0.363286in}}%
\pgfpathlineto{\pgfqpoint{1.434601in}{0.371234in}}%
\pgfpathlineto{\pgfqpoint{1.425411in}{0.374834in}}%
\pgfpathlineto{\pgfqpoint{1.416818in}{0.383999in}}%
\pgfpathlineto{\pgfqpoint{1.413687in}{0.396269in}}%
\pgfpathlineto{\pgfqpoint{1.405556in}{0.407420in}}%
\pgfpathlineto{\pgfqpoint{1.394734in}{0.416096in}}%
\pgfpathlineto{\pgfqpoint{1.384986in}{0.440984in}}%
\pgfpathlineto{\pgfqpoint{1.377115in}{0.468203in}}%
\pgfpathlineto{\pgfqpoint{1.370657in}{0.478966in}}%
\pgfpathlineto{\pgfqpoint{1.359460in}{0.488026in}}%
\pgfpathlineto{\pgfqpoint{1.356662in}{0.494603in}}%
\pgfpathlineto{\pgfqpoint{1.346179in}{0.498686in}}%
\pgfpathlineto{\pgfqpoint{1.335831in}{0.516112in}}%
\pgfpathlineto{\pgfqpoint{1.326382in}{0.514774in}}%
\pgfpathlineto{\pgfqpoint{1.316787in}{0.519126in}}%
\pgfpathlineto{\pgfqpoint{1.303186in}{0.518285in}}%
\pgfpathlineto{\pgfqpoint{1.289369in}{0.525467in}}%
\pgfpathlineto{\pgfqpoint{1.285472in}{0.518637in}}%
\pgfpathlineto{\pgfqpoint{1.269325in}{0.518466in}}%
\pgfpathlineto{\pgfqpoint{1.261104in}{0.505517in}}%
\pgfpathlineto{\pgfqpoint{1.253954in}{0.489484in}}%
\pgfpathlineto{\pgfqpoint{1.238750in}{0.472274in}}%
\pgfpathlineto{\pgfqpoint{1.221524in}{0.479827in}}%
\pgfpathlineto{\pgfqpoint{1.219080in}{0.484818in}}%
\pgfpathlineto{\pgfqpoint{1.208611in}{0.488625in}}%
\pgfpathlineto{\pgfqpoint{1.205396in}{0.493833in}}%
\pgfpathlineto{\pgfqpoint{1.191499in}{0.499090in}}%
\pgfpathlineto{\pgfqpoint{1.183725in}{0.509836in}}%
\pgfpathlineto{\pgfqpoint{1.174646in}{0.515021in}}%
\pgfpathlineto{\pgfqpoint{1.166834in}{0.524068in}}%
\pgfpathlineto{\pgfqpoint{1.160749in}{0.539391in}}%
\pgfpathlineto{\pgfqpoint{1.161430in}{0.560329in}}%
\pgfpathlineto{\pgfqpoint{1.154291in}{0.570909in}}%
\pgfpathlineto{\pgfqpoint{1.153469in}{0.582369in}}%
\pgfpathlineto{\pgfqpoint{1.137630in}{0.599527in}}%
\pgfpathlineto{\pgfqpoint{1.128414in}{0.603189in}}%
\pgfpathlineto{\pgfqpoint{1.118634in}{0.619199in}}%
\pgfpathlineto{\pgfqpoint{1.110300in}{0.625560in}}%
\pgfpathlineto{\pgfqpoint{1.099711in}{0.641049in}}%
\pgfpathlineto{\pgfqpoint{1.088876in}{0.647753in}}%
\pgfpathlineto{\pgfqpoint{1.081780in}{0.664929in}}%
\pgfpathlineto{\pgfqpoint{1.073883in}{0.669271in}}%
\pgfpathclose%
\pgfusepath{fill}%
\end{pgfscope}%
\begin{pgfscope}%
\pgfpathrectangle{\pgfqpoint{0.100000in}{0.100000in}}{\pgfqpoint{2.862985in}{1.832778in}}%
\pgfusepath{clip}%
\pgfsetbuttcap%
\pgfsetmiterjoin%
\definecolor{currentfill}{rgb}{0.005413,0.750035,0.999231}%
\pgfsetfillcolor{currentfill}%
\pgfsetlinewidth{0.000000pt}%
\definecolor{currentstroke}{rgb}{0.000000,0.000000,0.000000}%
\pgfsetstrokecolor{currentstroke}%
\pgfsetstrokeopacity{0.000000}%
\pgfsetdash{}{0pt}%
\pgfpathmoveto{\pgfqpoint{0.990915in}{1.012202in}}%
\pgfpathlineto{\pgfqpoint{1.069080in}{1.000850in}}%
\pgfpathlineto{\pgfqpoint{1.099362in}{0.996184in}}%
\pgfpathlineto{\pgfqpoint{1.184722in}{0.986032in}}%
\pgfpathlineto{\pgfqpoint{1.242175in}{0.980039in}}%
\pgfpathlineto{\pgfqpoint{1.291895in}{0.975633in}}%
\pgfpathlineto{\pgfqpoint{1.289118in}{0.943987in}}%
\pgfpathlineto{\pgfqpoint{1.287612in}{0.944071in}}%
\pgfpathlineto{\pgfqpoint{1.283057in}{0.889756in}}%
\pgfpathlineto{\pgfqpoint{1.278172in}{0.834767in}}%
\pgfpathlineto{\pgfqpoint{1.272520in}{0.777201in}}%
\pgfpathlineto{\pgfqpoint{1.263807in}{0.695029in}}%
\pgfpathlineto{\pgfqpoint{1.261021in}{0.662722in}}%
\pgfpathlineto{\pgfqpoint{1.209801in}{0.667891in}}%
\pgfpathlineto{\pgfqpoint{1.165452in}{0.672249in}}%
\pgfpathlineto{\pgfqpoint{1.104118in}{0.679215in}}%
\pgfpathlineto{\pgfqpoint{1.071132in}{0.683307in}}%
\pgfpathlineto{\pgfqpoint{1.069190in}{0.676811in}}%
\pgfpathlineto{\pgfqpoint{1.073883in}{0.669271in}}%
\pgfpathlineto{\pgfqpoint{1.034240in}{0.674419in}}%
\pgfpathlineto{\pgfqpoint{0.985330in}{0.681439in}}%
\pgfpathlineto{\pgfqpoint{0.980878in}{0.653775in}}%
\pgfpathlineto{\pgfqpoint{0.936264in}{0.660490in}}%
\pgfpathlineto{\pgfqpoint{0.946747in}{0.729118in}}%
\pgfpathlineto{\pgfqpoint{0.959723in}{0.812676in}}%
\pgfpathlineto{\pgfqpoint{0.971262in}{0.885469in}}%
\pgfpathlineto{\pgfqpoint{0.981281in}{0.950172in}}%
\pgfpathlineto{\pgfqpoint{0.990915in}{1.012202in}}%
\pgfpathclose%
\pgfusepath{fill}%
\end{pgfscope}%
\begin{pgfscope}%
\pgfpathrectangle{\pgfqpoint{0.100000in}{0.100000in}}{\pgfqpoint{2.862985in}{1.832778in}}%
\pgfusepath{clip}%
\pgfsetbuttcap%
\pgfsetmiterjoin%
\definecolor{currentfill}{rgb}{0.023068,0.755479,0.905421}%
\pgfsetfillcolor{currentfill}%
\pgfsetlinewidth{0.000000pt}%
\definecolor{currentstroke}{rgb}{0.000000,0.000000,0.000000}%
\pgfsetstrokecolor{currentstroke}%
\pgfsetstrokeopacity{0.000000}%
\pgfsetdash{}{0pt}%
\pgfpathmoveto{\pgfqpoint{2.233749in}{0.610119in}}%
\pgfpathlineto{\pgfqpoint{2.159178in}{0.601837in}}%
\pgfpathlineto{\pgfqpoint{2.093446in}{0.596734in}}%
\pgfpathlineto{\pgfqpoint{2.092624in}{0.588678in}}%
\pgfpathlineto{\pgfqpoint{2.098661in}{0.581011in}}%
\pgfpathlineto{\pgfqpoint{2.106031in}{0.576507in}}%
\pgfpathlineto{\pgfqpoint{2.105919in}{0.564647in}}%
\pgfpathlineto{\pgfqpoint{2.103966in}{0.556723in}}%
\pgfpathlineto{\pgfqpoint{2.097462in}{0.551013in}}%
\pgfpathlineto{\pgfqpoint{2.088396in}{0.551618in}}%
\pgfpathlineto{\pgfqpoint{2.079825in}{0.558693in}}%
\pgfpathlineto{\pgfqpoint{2.078297in}{0.571283in}}%
\pgfpathlineto{\pgfqpoint{2.071913in}{0.578606in}}%
\pgfpathlineto{\pgfqpoint{2.067568in}{0.552398in}}%
\pgfpathlineto{\pgfqpoint{2.052803in}{0.554889in}}%
\pgfpathlineto{\pgfqpoint{2.047669in}{0.600717in}}%
\pgfpathlineto{\pgfqpoint{2.042108in}{0.649015in}}%
\pgfpathlineto{\pgfqpoint{2.044173in}{0.718697in}}%
\pgfpathlineto{\pgfqpoint{2.045455in}{0.766513in}}%
\pgfpathlineto{\pgfqpoint{2.048181in}{0.839470in}}%
\pgfpathlineto{\pgfqpoint{2.042346in}{0.845890in}}%
\pgfpathlineto{\pgfqpoint{2.052965in}{0.847429in}}%
\pgfpathlineto{\pgfqpoint{2.114830in}{0.851430in}}%
\pgfpathlineto{\pgfqpoint{2.175022in}{0.856843in}}%
\pgfpathlineto{\pgfqpoint{2.190858in}{0.801373in}}%
\pgfpathlineto{\pgfqpoint{2.201658in}{0.760657in}}%
\pgfpathlineto{\pgfqpoint{2.211339in}{0.726571in}}%
\pgfpathlineto{\pgfqpoint{2.216938in}{0.712845in}}%
\pgfpathlineto{\pgfqpoint{2.225753in}{0.700073in}}%
\pgfpathlineto{\pgfqpoint{2.224330in}{0.693573in}}%
\pgfpathlineto{\pgfqpoint{2.230627in}{0.690406in}}%
\pgfpathlineto{\pgfqpoint{2.223185in}{0.680553in}}%
\pgfpathlineto{\pgfqpoint{2.220672in}{0.663006in}}%
\pgfpathlineto{\pgfqpoint{2.227945in}{0.642483in}}%
\pgfpathlineto{\pgfqpoint{2.226935in}{0.621831in}}%
\pgfpathlineto{\pgfqpoint{2.233749in}{0.610119in}}%
\pgfpathclose%
\pgfusepath{fill}%
\end{pgfscope}%
\begin{pgfscope}%
\pgfpathrectangle{\pgfqpoint{0.100000in}{0.100000in}}{\pgfqpoint{2.862985in}{1.832778in}}%
\pgfusepath{clip}%
\pgfsetbuttcap%
\pgfsetmiterjoin%
\definecolor{currentfill}{rgb}{0.200292,0.786574,0.971549}%
\pgfsetfillcolor{currentfill}%
\pgfsetlinewidth{0.000000pt}%
\definecolor{currentstroke}{rgb}{0.000000,0.000000,0.000000}%
\pgfsetstrokecolor{currentstroke}%
\pgfsetstrokeopacity{0.000000}%
\pgfsetdash{}{0pt}%
\pgfpathmoveto{\pgfqpoint{2.052803in}{0.554889in}}%
\pgfpathlineto{\pgfqpoint{2.037652in}{0.550553in}}%
\pgfpathlineto{\pgfqpoint{2.026358in}{0.553337in}}%
\pgfpathlineto{\pgfqpoint{2.005384in}{0.546666in}}%
\pgfpathlineto{\pgfqpoint{1.989565in}{0.538074in}}%
\pgfpathlineto{\pgfqpoint{1.983075in}{0.553600in}}%
\pgfpathlineto{\pgfqpoint{1.976384in}{0.560107in}}%
\pgfpathlineto{\pgfqpoint{1.973157in}{0.567894in}}%
\pgfpathlineto{\pgfqpoint{1.977888in}{0.588972in}}%
\pgfpathlineto{\pgfqpoint{1.933060in}{0.586211in}}%
\pgfpathlineto{\pgfqpoint{1.874932in}{0.583691in}}%
\pgfpathlineto{\pgfqpoint{1.878390in}{0.588938in}}%
\pgfpathlineto{\pgfqpoint{1.874116in}{0.598847in}}%
\pgfpathlineto{\pgfqpoint{1.878375in}{0.600866in}}%
\pgfpathlineto{\pgfqpoint{1.877430in}{0.610386in}}%
\pgfpathlineto{\pgfqpoint{1.884854in}{0.619619in}}%
\pgfpathlineto{\pgfqpoint{1.888916in}{0.637542in}}%
\pgfpathlineto{\pgfqpoint{1.902503in}{0.649350in}}%
\pgfpathlineto{\pgfqpoint{1.897467in}{0.660818in}}%
\pgfpathlineto{\pgfqpoint{1.907102in}{0.667419in}}%
\pgfpathlineto{\pgfqpoint{1.898791in}{0.679370in}}%
\pgfpathlineto{\pgfqpoint{1.892776in}{0.706727in}}%
\pgfpathlineto{\pgfqpoint{1.895057in}{0.711695in}}%
\pgfpathlineto{\pgfqpoint{1.898659in}{0.721223in}}%
\pgfpathlineto{\pgfqpoint{1.895307in}{0.731223in}}%
\pgfpathlineto{\pgfqpoint{1.896946in}{0.735772in}}%
\pgfpathlineto{\pgfqpoint{1.891048in}{0.752960in}}%
\pgfpathlineto{\pgfqpoint{1.908459in}{0.783821in}}%
\pgfpathlineto{\pgfqpoint{1.909269in}{0.792035in}}%
\pgfpathlineto{\pgfqpoint{1.917656in}{0.798007in}}%
\pgfpathlineto{\pgfqpoint{1.926605in}{0.817725in}}%
\pgfpathlineto{\pgfqpoint{1.926180in}{0.825739in}}%
\pgfpathlineto{\pgfqpoint{1.937326in}{0.833932in}}%
\pgfpathlineto{\pgfqpoint{1.933848in}{0.838785in}}%
\pgfpathlineto{\pgfqpoint{2.000643in}{0.842683in}}%
\pgfpathlineto{\pgfqpoint{2.042346in}{0.845890in}}%
\pgfpathlineto{\pgfqpoint{2.048181in}{0.839470in}}%
\pgfpathlineto{\pgfqpoint{2.045455in}{0.766513in}}%
\pgfpathlineto{\pgfqpoint{2.044173in}{0.718697in}}%
\pgfpathlineto{\pgfqpoint{2.042108in}{0.649015in}}%
\pgfpathlineto{\pgfqpoint{2.047669in}{0.600717in}}%
\pgfpathlineto{\pgfqpoint{2.052803in}{0.554889in}}%
\pgfpathclose%
\pgfusepath{fill}%
\end{pgfscope}%
\begin{pgfscope}%
\pgfpathrectangle{\pgfqpoint{0.100000in}{0.100000in}}{\pgfqpoint{2.862985in}{1.832778in}}%
\pgfusepath{clip}%
\pgfsetbuttcap%
\pgfsetmiterjoin%
\definecolor{currentfill}{rgb}{0.118416,0.782176,0.514494}%
\pgfsetfillcolor{currentfill}%
\pgfsetlinewidth{0.000000pt}%
\definecolor{currentstroke}{rgb}{0.000000,0.000000,0.000000}%
\pgfsetstrokecolor{currentstroke}%
\pgfsetstrokeopacity{0.000000}%
\pgfsetdash{}{0pt}%
\pgfpathmoveto{\pgfqpoint{2.175022in}{0.856843in}}%
\pgfpathlineto{\pgfqpoint{2.240820in}{0.864070in}}%
\pgfpathlineto{\pgfqpoint{2.280410in}{0.868948in}}%
\pgfpathlineto{\pgfqpoint{2.302958in}{0.872540in}}%
\pgfpathlineto{\pgfqpoint{2.293618in}{0.857809in}}%
\pgfpathlineto{\pgfqpoint{2.293623in}{0.850863in}}%
\pgfpathlineto{\pgfqpoint{2.303238in}{0.847117in}}%
\pgfpathlineto{\pgfqpoint{2.309776in}{0.841070in}}%
\pgfpathlineto{\pgfqpoint{2.319657in}{0.840323in}}%
\pgfpathlineto{\pgfqpoint{2.328892in}{0.823299in}}%
\pgfpathlineto{\pgfqpoint{2.338910in}{0.811294in}}%
\pgfpathlineto{\pgfqpoint{2.356702in}{0.801205in}}%
\pgfpathlineto{\pgfqpoint{2.361758in}{0.793157in}}%
\pgfpathlineto{\pgfqpoint{2.377624in}{0.784454in}}%
\pgfpathlineto{\pgfqpoint{2.378817in}{0.778143in}}%
\pgfpathlineto{\pgfqpoint{2.391688in}{0.765593in}}%
\pgfpathlineto{\pgfqpoint{2.401546in}{0.762019in}}%
\pgfpathlineto{\pgfqpoint{2.408528in}{0.750169in}}%
\pgfpathlineto{\pgfqpoint{2.411618in}{0.736849in}}%
\pgfpathlineto{\pgfqpoint{2.421858in}{0.731700in}}%
\pgfpathlineto{\pgfqpoint{2.430102in}{0.717383in}}%
\pgfpathlineto{\pgfqpoint{2.432761in}{0.706853in}}%
\pgfpathlineto{\pgfqpoint{2.444841in}{0.702371in}}%
\pgfpathlineto{\pgfqpoint{2.434704in}{0.682961in}}%
\pgfpathlineto{\pgfqpoint{2.430889in}{0.671491in}}%
\pgfpathlineto{\pgfqpoint{2.433378in}{0.666114in}}%
\pgfpathlineto{\pgfqpoint{2.429032in}{0.657178in}}%
\pgfpathlineto{\pgfqpoint{2.427185in}{0.644824in}}%
\pgfpathlineto{\pgfqpoint{2.419456in}{0.642523in}}%
\pgfpathlineto{\pgfqpoint{2.423537in}{0.631212in}}%
\pgfpathlineto{\pgfqpoint{2.423277in}{0.616862in}}%
\pgfpathlineto{\pgfqpoint{2.409539in}{0.617400in}}%
\pgfpathlineto{\pgfqpoint{2.399969in}{0.620018in}}%
\pgfpathlineto{\pgfqpoint{2.395828in}{0.615151in}}%
\pgfpathlineto{\pgfqpoint{2.398918in}{0.603653in}}%
\pgfpathlineto{\pgfqpoint{2.398246in}{0.590285in}}%
\pgfpathlineto{\pgfqpoint{2.392209in}{0.589262in}}%
\pgfpathlineto{\pgfqpoint{2.387325in}{0.601747in}}%
\pgfpathlineto{\pgfqpoint{2.337593in}{0.598433in}}%
\pgfpathlineto{\pgfqpoint{2.274885in}{0.595011in}}%
\pgfpathlineto{\pgfqpoint{2.243260in}{0.592785in}}%
\pgfpathlineto{\pgfqpoint{2.233749in}{0.610119in}}%
\pgfpathlineto{\pgfqpoint{2.226935in}{0.621831in}}%
\pgfpathlineto{\pgfqpoint{2.227945in}{0.642483in}}%
\pgfpathlineto{\pgfqpoint{2.220672in}{0.663006in}}%
\pgfpathlineto{\pgfqpoint{2.223185in}{0.680553in}}%
\pgfpathlineto{\pgfqpoint{2.230627in}{0.690406in}}%
\pgfpathlineto{\pgfqpoint{2.224330in}{0.693573in}}%
\pgfpathlineto{\pgfqpoint{2.225753in}{0.700073in}}%
\pgfpathlineto{\pgfqpoint{2.216938in}{0.712845in}}%
\pgfpathlineto{\pgfqpoint{2.211339in}{0.726571in}}%
\pgfpathlineto{\pgfqpoint{2.201658in}{0.760657in}}%
\pgfpathlineto{\pgfqpoint{2.190858in}{0.801373in}}%
\pgfpathlineto{\pgfqpoint{2.175022in}{0.856843in}}%
\pgfpathclose%
\pgfusepath{fill}%
\end{pgfscope}%
\begin{pgfscope}%
\pgfpathrectangle{\pgfqpoint{0.100000in}{0.100000in}}{\pgfqpoint{2.862985in}{1.832778in}}%
\pgfusepath{clip}%
\pgfsetbuttcap%
\pgfsetmiterjoin%
\definecolor{currentfill}{rgb}{0.112265,0.780454,0.539715}%
\pgfsetfillcolor{currentfill}%
\pgfsetlinewidth{0.000000pt}%
\definecolor{currentstroke}{rgb}{0.000000,0.000000,0.000000}%
\pgfsetstrokecolor{currentstroke}%
\pgfsetstrokeopacity{0.000000}%
\pgfsetdash{}{0pt}%
\pgfpathmoveto{\pgfqpoint{2.302958in}{0.872540in}}%
\pgfpathlineto{\pgfqpoint{2.329049in}{0.885519in}}%
\pgfpathlineto{\pgfqpoint{2.343499in}{0.890420in}}%
\pgfpathlineto{\pgfqpoint{2.406366in}{0.896952in}}%
\pgfpathlineto{\pgfqpoint{2.413024in}{0.894789in}}%
\pgfpathlineto{\pgfqpoint{2.421819in}{0.885949in}}%
\pgfpathlineto{\pgfqpoint{2.422304in}{0.878098in}}%
\pgfpathlineto{\pgfqpoint{2.480165in}{0.886600in}}%
\pgfpathlineto{\pgfqpoint{2.545869in}{0.839501in}}%
\pgfpathlineto{\pgfqpoint{2.533547in}{0.826658in}}%
\pgfpathlineto{\pgfqpoint{2.516640in}{0.796833in}}%
\pgfpathlineto{\pgfqpoint{2.521453in}{0.790423in}}%
\pgfpathlineto{\pgfqpoint{2.512434in}{0.777981in}}%
\pgfpathlineto{\pgfqpoint{2.503462in}{0.776573in}}%
\pgfpathlineto{\pgfqpoint{2.503429in}{0.769097in}}%
\pgfpathlineto{\pgfqpoint{2.496939in}{0.761274in}}%
\pgfpathlineto{\pgfqpoint{2.488862in}{0.759687in}}%
\pgfpathlineto{\pgfqpoint{2.490625in}{0.752736in}}%
\pgfpathlineto{\pgfqpoint{2.486118in}{0.747400in}}%
\pgfpathlineto{\pgfqpoint{2.475337in}{0.742790in}}%
\pgfpathlineto{\pgfqpoint{2.462196in}{0.732304in}}%
\pgfpathlineto{\pgfqpoint{2.464677in}{0.725520in}}%
\pgfpathlineto{\pgfqpoint{2.456407in}{0.721260in}}%
\pgfpathlineto{\pgfqpoint{2.449432in}{0.725739in}}%
\pgfpathlineto{\pgfqpoint{2.444331in}{0.706267in}}%
\pgfpathlineto{\pgfqpoint{2.432761in}{0.706853in}}%
\pgfpathlineto{\pgfqpoint{2.430102in}{0.717383in}}%
\pgfpathlineto{\pgfqpoint{2.421858in}{0.731700in}}%
\pgfpathlineto{\pgfqpoint{2.411618in}{0.736849in}}%
\pgfpathlineto{\pgfqpoint{2.408528in}{0.750169in}}%
\pgfpathlineto{\pgfqpoint{2.401546in}{0.762019in}}%
\pgfpathlineto{\pgfqpoint{2.391688in}{0.765593in}}%
\pgfpathlineto{\pgfqpoint{2.378817in}{0.778143in}}%
\pgfpathlineto{\pgfqpoint{2.377624in}{0.784454in}}%
\pgfpathlineto{\pgfqpoint{2.361758in}{0.793157in}}%
\pgfpathlineto{\pgfqpoint{2.356702in}{0.801205in}}%
\pgfpathlineto{\pgfqpoint{2.338910in}{0.811294in}}%
\pgfpathlineto{\pgfqpoint{2.328892in}{0.823299in}}%
\pgfpathlineto{\pgfqpoint{2.319657in}{0.840323in}}%
\pgfpathlineto{\pgfqpoint{2.309776in}{0.841070in}}%
\pgfpathlineto{\pgfqpoint{2.303238in}{0.847117in}}%
\pgfpathlineto{\pgfqpoint{2.293623in}{0.850863in}}%
\pgfpathlineto{\pgfqpoint{2.293618in}{0.857809in}}%
\pgfpathlineto{\pgfqpoint{2.302958in}{0.872540in}}%
\pgfpathclose%
\pgfusepath{fill}%
\end{pgfscope}%
\begin{pgfscope}%
\pgfpathrectangle{\pgfqpoint{0.100000in}{0.100000in}}{\pgfqpoint{2.862985in}{1.832778in}}%
\pgfusepath{clip}%
\pgfsetbuttcap%
\pgfsetmiterjoin%
\definecolor{currentfill}{rgb}{0.070373,0.762215,0.990004}%
\pgfsetfillcolor{currentfill}%
\pgfsetlinewidth{0.000000pt}%
\definecolor{currentstroke}{rgb}{0.000000,0.000000,0.000000}%
\pgfsetstrokecolor{currentstroke}%
\pgfsetstrokeopacity{0.000000}%
\pgfsetdash{}{0pt}%
\pgfpathmoveto{\pgfqpoint{1.711515in}{0.926094in}}%
\pgfpathlineto{\pgfqpoint{1.749856in}{0.926440in}}%
\pgfpathlineto{\pgfqpoint{1.800580in}{0.927364in}}%
\pgfpathlineto{\pgfqpoint{1.887257in}{0.930116in}}%
\pgfpathlineto{\pgfqpoint{1.936833in}{0.932737in}}%
\pgfpathlineto{\pgfqpoint{1.942150in}{0.926155in}}%
\pgfpathlineto{\pgfqpoint{1.941818in}{0.919204in}}%
\pgfpathlineto{\pgfqpoint{1.929808in}{0.907211in}}%
\pgfpathlineto{\pgfqpoint{1.926904in}{0.900644in}}%
\pgfpathlineto{\pgfqpoint{1.960213in}{0.903107in}}%
\pgfpathlineto{\pgfqpoint{1.960195in}{0.890974in}}%
\pgfpathlineto{\pgfqpoint{1.949509in}{0.885789in}}%
\pgfpathlineto{\pgfqpoint{1.949660in}{0.877545in}}%
\pgfpathlineto{\pgfqpoint{1.945647in}{0.871619in}}%
\pgfpathlineto{\pgfqpoint{1.943020in}{0.859033in}}%
\pgfpathlineto{\pgfqpoint{1.945732in}{0.848649in}}%
\pgfpathlineto{\pgfqpoint{1.933848in}{0.838785in}}%
\pgfpathlineto{\pgfqpoint{1.937326in}{0.833932in}}%
\pgfpathlineto{\pgfqpoint{1.926180in}{0.825739in}}%
\pgfpathlineto{\pgfqpoint{1.926605in}{0.817725in}}%
\pgfpathlineto{\pgfqpoint{1.917656in}{0.798007in}}%
\pgfpathlineto{\pgfqpoint{1.909269in}{0.792035in}}%
\pgfpathlineto{\pgfqpoint{1.908459in}{0.783821in}}%
\pgfpathlineto{\pgfqpoint{1.891048in}{0.752960in}}%
\pgfpathlineto{\pgfqpoint{1.896946in}{0.735772in}}%
\pgfpathlineto{\pgfqpoint{1.895307in}{0.731223in}}%
\pgfpathlineto{\pgfqpoint{1.898659in}{0.721223in}}%
\pgfpathlineto{\pgfqpoint{1.895057in}{0.711695in}}%
\pgfpathlineto{\pgfqpoint{1.847446in}{0.709729in}}%
\pgfpathlineto{\pgfqpoint{1.785638in}{0.708711in}}%
\pgfpathlineto{\pgfqpoint{1.743006in}{0.708326in}}%
\pgfpathlineto{\pgfqpoint{1.742799in}{0.741894in}}%
\pgfpathlineto{\pgfqpoint{1.732212in}{0.744093in}}%
\pgfpathlineto{\pgfqpoint{1.725225in}{0.741221in}}%
\pgfpathlineto{\pgfqpoint{1.719643in}{0.746495in}}%
\pgfpathlineto{\pgfqpoint{1.720219in}{0.782052in}}%
\pgfpathlineto{\pgfqpoint{1.721468in}{0.857760in}}%
\pgfpathlineto{\pgfqpoint{1.715413in}{0.902092in}}%
\pgfpathlineto{\pgfqpoint{1.711515in}{0.926094in}}%
\pgfpathclose%
\pgfusepath{fill}%
\end{pgfscope}%
\begin{pgfscope}%
\pgfpathrectangle{\pgfqpoint{0.100000in}{0.100000in}}{\pgfqpoint{2.862985in}{1.832778in}}%
\pgfusepath{clip}%
\pgfsetbuttcap%
\pgfsetmiterjoin%
\definecolor{currentfill}{rgb}{0.167812,0.780484,0.976163}%
\pgfsetfillcolor{currentfill}%
\pgfsetlinewidth{0.000000pt}%
\definecolor{currentstroke}{rgb}{0.000000,0.000000,0.000000}%
\pgfsetstrokecolor{currentstroke}%
\pgfsetstrokeopacity{0.000000}%
\pgfsetdash{}{0pt}%
\pgfpathmoveto{\pgfqpoint{1.743006in}{0.708326in}}%
\pgfpathlineto{\pgfqpoint{1.785638in}{0.708711in}}%
\pgfpathlineto{\pgfqpoint{1.847446in}{0.709729in}}%
\pgfpathlineto{\pgfqpoint{1.895057in}{0.711695in}}%
\pgfpathlineto{\pgfqpoint{1.892776in}{0.706727in}}%
\pgfpathlineto{\pgfqpoint{1.898791in}{0.679370in}}%
\pgfpathlineto{\pgfqpoint{1.907102in}{0.667419in}}%
\pgfpathlineto{\pgfqpoint{1.897467in}{0.660818in}}%
\pgfpathlineto{\pgfqpoint{1.902503in}{0.649350in}}%
\pgfpathlineto{\pgfqpoint{1.888916in}{0.637542in}}%
\pgfpathlineto{\pgfqpoint{1.884854in}{0.619619in}}%
\pgfpathlineto{\pgfqpoint{1.877430in}{0.610386in}}%
\pgfpathlineto{\pgfqpoint{1.878375in}{0.600866in}}%
\pgfpathlineto{\pgfqpoint{1.874116in}{0.598847in}}%
\pgfpathlineto{\pgfqpoint{1.878390in}{0.588938in}}%
\pgfpathlineto{\pgfqpoint{1.874932in}{0.583691in}}%
\pgfpathlineto{\pgfqpoint{1.933060in}{0.586211in}}%
\pgfpathlineto{\pgfqpoint{1.977888in}{0.588972in}}%
\pgfpathlineto{\pgfqpoint{1.973157in}{0.567894in}}%
\pgfpathlineto{\pgfqpoint{1.976384in}{0.560107in}}%
\pgfpathlineto{\pgfqpoint{1.983075in}{0.553600in}}%
\pgfpathlineto{\pgfqpoint{1.989565in}{0.538074in}}%
\pgfpathlineto{\pgfqpoint{1.969065in}{0.541651in}}%
\pgfpathlineto{\pgfqpoint{1.961502in}{0.547531in}}%
\pgfpathlineto{\pgfqpoint{1.952492in}{0.547811in}}%
\pgfpathlineto{\pgfqpoint{1.943019in}{0.534934in}}%
\pgfpathlineto{\pgfqpoint{1.944907in}{0.529071in}}%
\pgfpathlineto{\pgfqpoint{1.975169in}{0.525523in}}%
\pgfpathlineto{\pgfqpoint{1.983091in}{0.518783in}}%
\pgfpathlineto{\pgfqpoint{1.995984in}{0.515315in}}%
\pgfpathlineto{\pgfqpoint{1.990353in}{0.507336in}}%
\pgfpathlineto{\pgfqpoint{1.980867in}{0.500376in}}%
\pgfpathlineto{\pgfqpoint{2.005906in}{0.484726in}}%
\pgfpathlineto{\pgfqpoint{2.017569in}{0.482281in}}%
\pgfpathlineto{\pgfqpoint{2.022419in}{0.469488in}}%
\pgfpathlineto{\pgfqpoint{2.011438in}{0.467135in}}%
\pgfpathlineto{\pgfqpoint{1.990799in}{0.483224in}}%
\pgfpathlineto{\pgfqpoint{1.982729in}{0.485461in}}%
\pgfpathlineto{\pgfqpoint{1.978824in}{0.491817in}}%
\pgfpathlineto{\pgfqpoint{1.966854in}{0.489208in}}%
\pgfpathlineto{\pgfqpoint{1.963107in}{0.481018in}}%
\pgfpathlineto{\pgfqpoint{1.965480in}{0.471892in}}%
\pgfpathlineto{\pgfqpoint{1.957422in}{0.466506in}}%
\pgfpathlineto{\pgfqpoint{1.950068in}{0.479774in}}%
\pgfpathlineto{\pgfqpoint{1.937187in}{0.475818in}}%
\pgfpathlineto{\pgfqpoint{1.932360in}{0.467901in}}%
\pgfpathlineto{\pgfqpoint{1.923213in}{0.470153in}}%
\pgfpathlineto{\pgfqpoint{1.917364in}{0.480150in}}%
\pgfpathlineto{\pgfqpoint{1.901847in}{0.483451in}}%
\pgfpathlineto{\pgfqpoint{1.898917in}{0.488653in}}%
\pgfpathlineto{\pgfqpoint{1.882751in}{0.497715in}}%
\pgfpathlineto{\pgfqpoint{1.878738in}{0.505638in}}%
\pgfpathlineto{\pgfqpoint{1.865197in}{0.502400in}}%
\pgfpathlineto{\pgfqpoint{1.866917in}{0.509654in}}%
\pgfpathlineto{\pgfqpoint{1.850086in}{0.502213in}}%
\pgfpathlineto{\pgfqpoint{1.854607in}{0.494492in}}%
\pgfpathlineto{\pgfqpoint{1.841610in}{0.489928in}}%
\pgfpathlineto{\pgfqpoint{1.824395in}{0.492437in}}%
\pgfpathlineto{\pgfqpoint{1.789551in}{0.504372in}}%
\pgfpathlineto{\pgfqpoint{1.762666in}{0.502003in}}%
\pgfpathlineto{\pgfqpoint{1.760335in}{0.517680in}}%
\pgfpathlineto{\pgfqpoint{1.763712in}{0.524834in}}%
\pgfpathlineto{\pgfqpoint{1.763423in}{0.536215in}}%
\pgfpathlineto{\pgfqpoint{1.760101in}{0.538927in}}%
\pgfpathlineto{\pgfqpoint{1.761227in}{0.551919in}}%
\pgfpathlineto{\pgfqpoint{1.771019in}{0.569988in}}%
\pgfpathlineto{\pgfqpoint{1.772272in}{0.576816in}}%
\pgfpathlineto{\pgfqpoint{1.770658in}{0.592942in}}%
\pgfpathlineto{\pgfqpoint{1.763454in}{0.600538in}}%
\pgfpathlineto{\pgfqpoint{1.759693in}{0.614721in}}%
\pgfpathlineto{\pgfqpoint{1.755065in}{0.617977in}}%
\pgfpathlineto{\pgfqpoint{1.757137in}{0.625668in}}%
\pgfpathlineto{\pgfqpoint{1.751243in}{0.637164in}}%
\pgfpathlineto{\pgfqpoint{1.743890in}{0.643403in}}%
\pgfpathlineto{\pgfqpoint{1.743006in}{0.708326in}}%
\pgfpathclose%
\pgfusepath{fill}%
\end{pgfscope}%
\begin{pgfscope}%
\pgfpathrectangle{\pgfqpoint{0.100000in}{0.100000in}}{\pgfqpoint{2.862985in}{1.832778in}}%
\pgfusepath{clip}%
\pgfsetbuttcap%
\pgfsetmiterjoin%
\definecolor{currentfill}{rgb}{0.167812,0.780484,0.976163}%
\pgfsetfillcolor{currentfill}%
\pgfsetlinewidth{0.000000pt}%
\definecolor{currentstroke}{rgb}{0.000000,0.000000,0.000000}%
\pgfsetstrokecolor{currentstroke}%
\pgfsetstrokeopacity{0.000000}%
\pgfsetdash{}{0pt}%
\pgfpathmoveto{\pgfqpoint{1.857025in}{0.493901in}}%
\pgfpathlineto{\pgfqpoint{1.863193in}{0.497570in}}%
\pgfpathlineto{\pgfqpoint{1.870679in}{0.493233in}}%
\pgfpathlineto{\pgfqpoint{1.866504in}{0.487265in}}%
\pgfpathlineto{\pgfqpoint{1.857025in}{0.493901in}}%
\pgfpathclose%
\pgfusepath{fill}%
\end{pgfscope}%
\begin{pgfscope}%
\pgfpathrectangle{\pgfqpoint{0.100000in}{0.100000in}}{\pgfqpoint{2.862985in}{1.832778in}}%
\pgfusepath{clip}%
\pgfsetbuttcap%
\pgfsetmiterjoin%
\definecolor{currentfill}{rgb}{0.075356,0.770119,0.691042}%
\pgfsetfillcolor{currentfill}%
\pgfsetlinewidth{0.000000pt}%
\definecolor{currentstroke}{rgb}{0.000000,0.000000,0.000000}%
\pgfsetstrokecolor{currentstroke}%
\pgfsetstrokeopacity{0.000000}%
\pgfsetdash{}{0pt}%
\pgfpathmoveto{\pgfqpoint{2.105919in}{0.564647in}}%
\pgfpathlineto{\pgfqpoint{2.106031in}{0.576507in}}%
\pgfpathlineto{\pgfqpoint{2.098661in}{0.581011in}}%
\pgfpathlineto{\pgfqpoint{2.092624in}{0.588678in}}%
\pgfpathlineto{\pgfqpoint{2.093446in}{0.596734in}}%
\pgfpathlineto{\pgfqpoint{2.159178in}{0.601837in}}%
\pgfpathlineto{\pgfqpoint{2.233749in}{0.610119in}}%
\pgfpathlineto{\pgfqpoint{2.243260in}{0.592785in}}%
\pgfpathlineto{\pgfqpoint{2.274885in}{0.595011in}}%
\pgfpathlineto{\pgfqpoint{2.337593in}{0.598433in}}%
\pgfpathlineto{\pgfqpoint{2.387325in}{0.601747in}}%
\pgfpathlineto{\pgfqpoint{2.392209in}{0.589262in}}%
\pgfpathlineto{\pgfqpoint{2.398246in}{0.590285in}}%
\pgfpathlineto{\pgfqpoint{2.398918in}{0.603653in}}%
\pgfpathlineto{\pgfqpoint{2.395828in}{0.615151in}}%
\pgfpathlineto{\pgfqpoint{2.399969in}{0.620018in}}%
\pgfpathlineto{\pgfqpoint{2.409539in}{0.617400in}}%
\pgfpathlineto{\pgfqpoint{2.423277in}{0.616862in}}%
\pgfpathlineto{\pgfqpoint{2.425389in}{0.606589in}}%
\pgfpathlineto{\pgfqpoint{2.429615in}{0.600719in}}%
\pgfpathlineto{\pgfqpoint{2.432915in}{0.587877in}}%
\pgfpathlineto{\pgfqpoint{2.443139in}{0.568003in}}%
\pgfpathlineto{\pgfqpoint{2.443190in}{0.562628in}}%
\pgfpathlineto{\pgfqpoint{2.458302in}{0.539291in}}%
\pgfpathlineto{\pgfqpoint{2.459764in}{0.534471in}}%
\pgfpathlineto{\pgfqpoint{2.482434in}{0.501512in}}%
\pgfpathlineto{\pgfqpoint{2.478842in}{0.500979in}}%
\pgfpathlineto{\pgfqpoint{2.488399in}{0.477537in}}%
\pgfpathlineto{\pgfqpoint{2.514653in}{0.436782in}}%
\pgfpathlineto{\pgfqpoint{2.529896in}{0.406818in}}%
\pgfpathlineto{\pgfqpoint{2.534991in}{0.401571in}}%
\pgfpathlineto{\pgfqpoint{2.543722in}{0.382763in}}%
\pgfpathlineto{\pgfqpoint{2.546750in}{0.352646in}}%
\pgfpathlineto{\pgfqpoint{2.547955in}{0.330127in}}%
\pgfpathlineto{\pgfqpoint{2.546500in}{0.315709in}}%
\pgfpathlineto{\pgfqpoint{2.541819in}{0.305401in}}%
\pgfpathlineto{\pgfqpoint{2.543994in}{0.291885in}}%
\pgfpathlineto{\pgfqpoint{2.538954in}{0.281186in}}%
\pgfpathlineto{\pgfqpoint{2.523991in}{0.272413in}}%
\pgfpathlineto{\pgfqpoint{2.514259in}{0.273053in}}%
\pgfpathlineto{\pgfqpoint{2.507940in}{0.268467in}}%
\pgfpathlineto{\pgfqpoint{2.506129in}{0.280712in}}%
\pgfpathlineto{\pgfqpoint{2.495668in}{0.283944in}}%
\pgfpathlineto{\pgfqpoint{2.486290in}{0.301030in}}%
\pgfpathlineto{\pgfqpoint{2.485231in}{0.308810in}}%
\pgfpathlineto{\pgfqpoint{2.468448in}{0.313535in}}%
\pgfpathlineto{\pgfqpoint{2.457633in}{0.312512in}}%
\pgfpathlineto{\pgfqpoint{2.451514in}{0.323801in}}%
\pgfpathlineto{\pgfqpoint{2.444508in}{0.344093in}}%
\pgfpathlineto{\pgfqpoint{2.434786in}{0.348206in}}%
\pgfpathlineto{\pgfqpoint{2.429515in}{0.359834in}}%
\pgfpathlineto{\pgfqpoint{2.416625in}{0.366687in}}%
\pgfpathlineto{\pgfqpoint{2.409174in}{0.375259in}}%
\pgfpathlineto{\pgfqpoint{2.396320in}{0.398564in}}%
\pgfpathlineto{\pgfqpoint{2.394809in}{0.414966in}}%
\pgfpathlineto{\pgfqpoint{2.401813in}{0.425461in}}%
\pgfpathlineto{\pgfqpoint{2.401113in}{0.432754in}}%
\pgfpathlineto{\pgfqpoint{2.393019in}{0.433533in}}%
\pgfpathlineto{\pgfqpoint{2.386256in}{0.438606in}}%
\pgfpathlineto{\pgfqpoint{2.382553in}{0.432412in}}%
\pgfpathlineto{\pgfqpoint{2.389043in}{0.427365in}}%
\pgfpathlineto{\pgfqpoint{2.384319in}{0.418289in}}%
\pgfpathlineto{\pgfqpoint{2.376677in}{0.425822in}}%
\pgfpathlineto{\pgfqpoint{2.377564in}{0.446710in}}%
\pgfpathlineto{\pgfqpoint{2.381250in}{0.463657in}}%
\pgfpathlineto{\pgfqpoint{2.379355in}{0.492736in}}%
\pgfpathlineto{\pgfqpoint{2.371738in}{0.499662in}}%
\pgfpathlineto{\pgfqpoint{2.367907in}{0.508549in}}%
\pgfpathlineto{\pgfqpoint{2.354785in}{0.508352in}}%
\pgfpathlineto{\pgfqpoint{2.341814in}{0.522976in}}%
\pgfpathlineto{\pgfqpoint{2.333117in}{0.527359in}}%
\pgfpathlineto{\pgfqpoint{2.330538in}{0.536620in}}%
\pgfpathlineto{\pgfqpoint{2.322031in}{0.539872in}}%
\pgfpathlineto{\pgfqpoint{2.314981in}{0.550098in}}%
\pgfpathlineto{\pgfqpoint{2.296358in}{0.558444in}}%
\pgfpathlineto{\pgfqpoint{2.281879in}{0.558658in}}%
\pgfpathlineto{\pgfqpoint{2.275576in}{0.555452in}}%
\pgfpathlineto{\pgfqpoint{2.277143in}{0.545445in}}%
\pgfpathlineto{\pgfqpoint{2.270570in}{0.545919in}}%
\pgfpathlineto{\pgfqpoint{2.250353in}{0.531927in}}%
\pgfpathlineto{\pgfqpoint{2.226068in}{0.526381in}}%
\pgfpathlineto{\pgfqpoint{2.225669in}{0.533243in}}%
\pgfpathlineto{\pgfqpoint{2.220287in}{0.539949in}}%
\pgfpathlineto{\pgfqpoint{2.205848in}{0.549208in}}%
\pgfpathlineto{\pgfqpoint{2.185121in}{0.558700in}}%
\pgfpathlineto{\pgfqpoint{2.162608in}{0.563725in}}%
\pgfpathlineto{\pgfqpoint{2.158382in}{0.570558in}}%
\pgfpathlineto{\pgfqpoint{2.150265in}{0.564859in}}%
\pgfpathlineto{\pgfqpoint{2.140491in}{0.563602in}}%
\pgfpathlineto{\pgfqpoint{2.106467in}{0.554637in}}%
\pgfpathlineto{\pgfqpoint{2.105919in}{0.564647in}}%
\pgfpathclose%
\pgfusepath{fill}%
\end{pgfscope}%
\begin{pgfscope}%
\pgfpathrectangle{\pgfqpoint{0.100000in}{0.100000in}}{\pgfqpoint{2.862985in}{1.832778in}}%
\pgfusepath{clip}%
\pgfsetbuttcap%
\pgfsetmiterjoin%
\definecolor{currentfill}{rgb}{0.075356,0.770119,0.691042}%
\pgfsetfillcolor{currentfill}%
\pgfsetlinewidth{0.000000pt}%
\definecolor{currentstroke}{rgb}{0.000000,0.000000,0.000000}%
\pgfsetstrokecolor{currentstroke}%
\pgfsetstrokeopacity{0.000000}%
\pgfsetdash{}{0pt}%
\pgfpathmoveto{\pgfqpoint{2.485737in}{0.501980in}}%
\pgfpathlineto{\pgfqpoint{2.494540in}{0.491740in}}%
\pgfpathlineto{\pgfqpoint{2.484372in}{0.491096in}}%
\pgfpathlineto{\pgfqpoint{2.485737in}{0.501980in}}%
\pgfpathclose%
\pgfusepath{fill}%
\end{pgfscope}%
\begin{pgfscope}%
\pgfpathrectangle{\pgfqpoint{0.100000in}{0.100000in}}{\pgfqpoint{2.862985in}{1.832778in}}%
\pgfusepath{clip}%
\pgfsetbuttcap%
\pgfsetmiterjoin%
\definecolor{currentfill}{rgb}{0.146098,0.789927,0.401000}%
\pgfsetfillcolor{currentfill}%
\pgfsetlinewidth{0.000000pt}%
\definecolor{currentstroke}{rgb}{0.000000,0.000000,0.000000}%
\pgfsetstrokecolor{currentstroke}%
\pgfsetstrokeopacity{0.000000}%
\pgfsetdash{}{0pt}%
\pgfpathmoveto{\pgfqpoint{1.967522in}{1.671190in}}%
\pgfpathlineto{\pgfqpoint{1.962746in}{1.661894in}}%
\pgfpathlineto{\pgfqpoint{1.951351in}{1.656468in}}%
\pgfpathlineto{\pgfqpoint{1.946416in}{1.649166in}}%
\pgfpathlineto{\pgfqpoint{1.940635in}{1.654425in}}%
\pgfpathlineto{\pgfqpoint{1.967522in}{1.671190in}}%
\pgfpathclose%
\pgfusepath{fill}%
\end{pgfscope}%
\begin{pgfscope}%
\pgfpathrectangle{\pgfqpoint{0.100000in}{0.100000in}}{\pgfqpoint{2.862985in}{1.832778in}}%
\pgfusepath{clip}%
\pgfsetbuttcap%
\pgfsetmiterjoin%
\definecolor{currentfill}{rgb}{0.146098,0.789927,0.401000}%
\pgfsetfillcolor{currentfill}%
\pgfsetlinewidth{0.000000pt}%
\definecolor{currentstroke}{rgb}{0.000000,0.000000,0.000000}%
\pgfsetstrokecolor{currentstroke}%
\pgfsetstrokeopacity{0.000000}%
\pgfsetdash{}{0pt}%
\pgfpathmoveto{\pgfqpoint{1.971396in}{1.615244in}}%
\pgfpathlineto{\pgfqpoint{1.983054in}{1.626120in}}%
\pgfpathlineto{\pgfqpoint{2.001037in}{1.628969in}}%
\pgfpathlineto{\pgfqpoint{1.996084in}{1.621405in}}%
\pgfpathlineto{\pgfqpoint{1.983756in}{1.610466in}}%
\pgfpathlineto{\pgfqpoint{1.976549in}{1.596426in}}%
\pgfpathlineto{\pgfqpoint{1.973496in}{1.604041in}}%
\pgfpathlineto{\pgfqpoint{1.968076in}{1.605134in}}%
\pgfpathlineto{\pgfqpoint{1.967532in}{1.612009in}}%
\pgfpathlineto{\pgfqpoint{1.971396in}{1.615244in}}%
\pgfpathclose%
\pgfusepath{fill}%
\end{pgfscope}%
\begin{pgfscope}%
\pgfpathrectangle{\pgfqpoint{0.100000in}{0.100000in}}{\pgfqpoint{2.862985in}{1.832778in}}%
\pgfusepath{clip}%
\pgfsetbuttcap%
\pgfsetmiterjoin%
\definecolor{currentfill}{rgb}{0.146098,0.789927,0.401000}%
\pgfsetfillcolor{currentfill}%
\pgfsetlinewidth{0.000000pt}%
\definecolor{currentstroke}{rgb}{0.000000,0.000000,0.000000}%
\pgfsetstrokecolor{currentstroke}%
\pgfsetstrokeopacity{0.000000}%
\pgfsetdash{}{0pt}%
\pgfpathmoveto{\pgfqpoint{2.017832in}{1.482495in}}%
\pgfpathlineto{\pgfqpoint{2.014727in}{1.485941in}}%
\pgfpathlineto{\pgfqpoint{2.017991in}{1.495641in}}%
\pgfpathlineto{\pgfqpoint{2.008301in}{1.496259in}}%
\pgfpathlineto{\pgfqpoint{2.010870in}{1.504623in}}%
\pgfpathlineto{\pgfqpoint{2.009272in}{1.517943in}}%
\pgfpathlineto{\pgfqpoint{2.000578in}{1.522561in}}%
\pgfpathlineto{\pgfqpoint{1.991480in}{1.531907in}}%
\pgfpathlineto{\pgfqpoint{1.977460in}{1.534643in}}%
\pgfpathlineto{\pgfqpoint{1.963816in}{1.534565in}}%
\pgfpathlineto{\pgfqpoint{1.950397in}{1.541198in}}%
\pgfpathlineto{\pgfqpoint{1.905502in}{1.550866in}}%
\pgfpathlineto{\pgfqpoint{1.900596in}{1.561110in}}%
\pgfpathlineto{\pgfqpoint{1.891845in}{1.564606in}}%
\pgfpathlineto{\pgfqpoint{1.908375in}{1.572449in}}%
\pgfpathlineto{\pgfqpoint{1.917713in}{1.582234in}}%
\pgfpathlineto{\pgfqpoint{1.935106in}{1.584888in}}%
\pgfpathlineto{\pgfqpoint{1.945809in}{1.594855in}}%
\pgfpathlineto{\pgfqpoint{1.951425in}{1.595254in}}%
\pgfpathlineto{\pgfqpoint{1.955710in}{1.602362in}}%
\pgfpathlineto{\pgfqpoint{1.966157in}{1.610783in}}%
\pgfpathlineto{\pgfqpoint{1.967067in}{1.604831in}}%
\pgfpathlineto{\pgfqpoint{1.971774in}{1.603602in}}%
\pgfpathlineto{\pgfqpoint{1.975316in}{1.596504in}}%
\pgfpathlineto{\pgfqpoint{1.975954in}{1.584403in}}%
\pgfpathlineto{\pgfqpoint{1.986599in}{1.591634in}}%
\pgfpathlineto{\pgfqpoint{1.999002in}{1.593145in}}%
\pgfpathlineto{\pgfqpoint{2.009590in}{1.589356in}}%
\pgfpathlineto{\pgfqpoint{2.023949in}{1.569658in}}%
\pgfpathlineto{\pgfqpoint{2.039631in}{1.572806in}}%
\pgfpathlineto{\pgfqpoint{2.046003in}{1.567535in}}%
\pgfpathlineto{\pgfqpoint{2.056252in}{1.567064in}}%
\pgfpathlineto{\pgfqpoint{2.063061in}{1.576518in}}%
\pgfpathlineto{\pgfqpoint{2.075982in}{1.584876in}}%
\pgfpathlineto{\pgfqpoint{2.088400in}{1.587493in}}%
\pgfpathlineto{\pgfqpoint{2.103808in}{1.587766in}}%
\pgfpathlineto{\pgfqpoint{2.115068in}{1.594222in}}%
\pgfpathlineto{\pgfqpoint{2.124258in}{1.591248in}}%
\pgfpathlineto{\pgfqpoint{2.126215in}{1.577594in}}%
\pgfpathlineto{\pgfqpoint{2.145930in}{1.575456in}}%
\pgfpathlineto{\pgfqpoint{2.156638in}{1.581847in}}%
\pgfpathlineto{\pgfqpoint{2.164061in}{1.567449in}}%
\pgfpathlineto{\pgfqpoint{2.178225in}{1.550799in}}%
\pgfpathlineto{\pgfqpoint{2.158401in}{1.550895in}}%
\pgfpathlineto{\pgfqpoint{2.152137in}{1.548837in}}%
\pgfpathlineto{\pgfqpoint{2.143551in}{1.551511in}}%
\pgfpathlineto{\pgfqpoint{2.142974in}{1.539985in}}%
\pgfpathlineto{\pgfqpoint{2.127384in}{1.548993in}}%
\pgfpathlineto{\pgfqpoint{2.107335in}{1.551746in}}%
\pgfpathlineto{\pgfqpoint{2.101816in}{1.542973in}}%
\pgfpathlineto{\pgfqpoint{2.075584in}{1.538707in}}%
\pgfpathlineto{\pgfqpoint{2.072566in}{1.531276in}}%
\pgfpathlineto{\pgfqpoint{2.054312in}{1.529046in}}%
\pgfpathlineto{\pgfqpoint{2.048794in}{1.521492in}}%
\pgfpathlineto{\pgfqpoint{2.039139in}{1.519466in}}%
\pgfpathlineto{\pgfqpoint{2.031427in}{1.501552in}}%
\pgfpathlineto{\pgfqpoint{2.021663in}{1.484203in}}%
\pgfpathlineto{\pgfqpoint{2.017832in}{1.482495in}}%
\pgfpathclose%
\pgfusepath{fill}%
\end{pgfscope}%
\begin{pgfscope}%
\pgfpathrectangle{\pgfqpoint{0.100000in}{0.100000in}}{\pgfqpoint{2.862985in}{1.832778in}}%
\pgfusepath{clip}%
\pgfsetbuttcap%
\pgfsetmiterjoin%
\definecolor{currentfill}{rgb}{0.146098,0.789927,0.401000}%
\pgfsetfillcolor{currentfill}%
\pgfsetlinewidth{0.000000pt}%
\definecolor{currentstroke}{rgb}{0.000000,0.000000,0.000000}%
\pgfsetstrokecolor{currentstroke}%
\pgfsetstrokeopacity{0.000000}%
\pgfsetdash{}{0pt}%
\pgfpathmoveto{\pgfqpoint{2.074029in}{1.274203in}}%
\pgfpathlineto{\pgfqpoint{2.083339in}{1.284002in}}%
\pgfpathlineto{\pgfqpoint{2.087580in}{1.298209in}}%
\pgfpathlineto{\pgfqpoint{2.092623in}{1.306443in}}%
\pgfpathlineto{\pgfqpoint{2.095716in}{1.317649in}}%
\pgfpathlineto{\pgfqpoint{2.096678in}{1.339993in}}%
\pgfpathlineto{\pgfqpoint{2.092021in}{1.361407in}}%
\pgfpathlineto{\pgfqpoint{2.076633in}{1.394197in}}%
\pgfpathlineto{\pgfqpoint{2.080800in}{1.404507in}}%
\pgfpathlineto{\pgfqpoint{2.075372in}{1.418758in}}%
\pgfpathlineto{\pgfqpoint{2.084684in}{1.438466in}}%
\pgfpathlineto{\pgfqpoint{2.083166in}{1.460442in}}%
\pgfpathlineto{\pgfqpoint{2.089653in}{1.463207in}}%
\pgfpathlineto{\pgfqpoint{2.090468in}{1.473664in}}%
\pgfpathlineto{\pgfqpoint{2.101982in}{1.480375in}}%
\pgfpathlineto{\pgfqpoint{2.108494in}{1.479291in}}%
\pgfpathlineto{\pgfqpoint{2.110324in}{1.468077in}}%
\pgfpathlineto{\pgfqpoint{2.115407in}{1.467638in}}%
\pgfpathlineto{\pgfqpoint{2.120044in}{1.483813in}}%
\pgfpathlineto{\pgfqpoint{2.118468in}{1.496393in}}%
\pgfpathlineto{\pgfqpoint{2.121483in}{1.503617in}}%
\pgfpathlineto{\pgfqpoint{2.137785in}{1.511080in}}%
\pgfpathlineto{\pgfqpoint{2.130356in}{1.513735in}}%
\pgfpathlineto{\pgfqpoint{2.127936in}{1.520146in}}%
\pgfpathlineto{\pgfqpoint{2.133288in}{1.531404in}}%
\pgfpathlineto{\pgfqpoint{2.143841in}{1.535292in}}%
\pgfpathlineto{\pgfqpoint{2.156095in}{1.528650in}}%
\pgfpathlineto{\pgfqpoint{2.167648in}{1.528567in}}%
\pgfpathlineto{\pgfqpoint{2.173015in}{1.520808in}}%
\pgfpathlineto{\pgfqpoint{2.181106in}{1.521347in}}%
\pgfpathlineto{\pgfqpoint{2.206083in}{1.510561in}}%
\pgfpathlineto{\pgfqpoint{2.211074in}{1.500120in}}%
\pgfpathlineto{\pgfqpoint{2.205624in}{1.496504in}}%
\pgfpathlineto{\pgfqpoint{2.207303in}{1.488679in}}%
\pgfpathlineto{\pgfqpoint{2.212687in}{1.485198in}}%
\pgfpathlineto{\pgfqpoint{2.215688in}{1.475579in}}%
\pgfpathlineto{\pgfqpoint{2.215268in}{1.452144in}}%
\pgfpathlineto{\pgfqpoint{2.208159in}{1.446546in}}%
\pgfpathlineto{\pgfqpoint{2.206564in}{1.434236in}}%
\pgfpathlineto{\pgfqpoint{2.193377in}{1.422859in}}%
\pgfpathlineto{\pgfqpoint{2.194159in}{1.409089in}}%
\pgfpathlineto{\pgfqpoint{2.205717in}{1.404265in}}%
\pgfpathlineto{\pgfqpoint{2.218756in}{1.421446in}}%
\pgfpathlineto{\pgfqpoint{2.219833in}{1.427678in}}%
\pgfpathlineto{\pgfqpoint{2.236094in}{1.438035in}}%
\pgfpathlineto{\pgfqpoint{2.246435in}{1.433238in}}%
\pgfpathlineto{\pgfqpoint{2.252924in}{1.422395in}}%
\pgfpathlineto{\pgfqpoint{2.263393in}{1.384749in}}%
\pgfpathlineto{\pgfqpoint{2.268940in}{1.372839in}}%
\pgfpathlineto{\pgfqpoint{2.267203in}{1.367915in}}%
\pgfpathlineto{\pgfqpoint{2.267372in}{1.351150in}}%
\pgfpathlineto{\pgfqpoint{2.257281in}{1.352759in}}%
\pgfpathlineto{\pgfqpoint{2.251581in}{1.340231in}}%
\pgfpathlineto{\pgfqpoint{2.250791in}{1.331714in}}%
\pgfpathlineto{\pgfqpoint{2.243169in}{1.326246in}}%
\pgfpathlineto{\pgfqpoint{2.241475in}{1.309636in}}%
\pgfpathlineto{\pgfqpoint{2.230390in}{1.288653in}}%
\pgfpathlineto{\pgfqpoint{2.169752in}{1.279623in}}%
\pgfpathlineto{\pgfqpoint{2.169393in}{1.283592in}}%
\pgfpathlineto{\pgfqpoint{2.128831in}{1.279317in}}%
\pgfpathlineto{\pgfqpoint{2.074029in}{1.274203in}}%
\pgfpathclose%
\pgfusepath{fill}%
\end{pgfscope}%
\begin{pgfscope}%
\pgfpathrectangle{\pgfqpoint{0.100000in}{0.100000in}}{\pgfqpoint{2.862985in}{1.832778in}}%
\pgfusepath{clip}%
\pgfsetbuttcap%
\pgfsetmiterjoin%
\definecolor{currentfill}{rgb}{0.041522,0.760646,0.829758}%
\pgfsetfillcolor{currentfill}%
\pgfsetlinewidth{0.000000pt}%
\definecolor{currentstroke}{rgb}{0.000000,0.000000,0.000000}%
\pgfsetstrokecolor{currentstroke}%
\pgfsetstrokeopacity{0.000000}%
\pgfsetdash{}{0pt}%
\pgfpathmoveto{\pgfqpoint{0.542548in}{0.431818in}}%
\pgfpathlineto{\pgfqpoint{0.541337in}{0.430752in}}%
\pgfpathlineto{\pgfqpoint{0.535781in}{0.432173in}}%
\pgfpathlineto{\pgfqpoint{0.536196in}{0.434969in}}%
\pgfpathlineto{\pgfqpoint{0.539341in}{0.436249in}}%
\pgfpathlineto{\pgfqpoint{0.543824in}{0.442694in}}%
\pgfpathlineto{\pgfqpoint{0.544569in}{0.447051in}}%
\pgfpathlineto{\pgfqpoint{0.547090in}{0.448192in}}%
\pgfpathlineto{\pgfqpoint{0.551190in}{0.448197in}}%
\pgfpathlineto{\pgfqpoint{0.555477in}{0.462333in}}%
\pgfpathlineto{\pgfqpoint{0.557157in}{0.463496in}}%
\pgfpathlineto{\pgfqpoint{0.555651in}{0.466260in}}%
\pgfpathlineto{\pgfqpoint{0.552125in}{0.463464in}}%
\pgfpathlineto{\pgfqpoint{0.541766in}{0.467233in}}%
\pgfpathlineto{\pgfqpoint{0.535583in}{0.473582in}}%
\pgfpathlineto{\pgfqpoint{0.538280in}{0.476479in}}%
\pgfpathlineto{\pgfqpoint{0.537053in}{0.480565in}}%
\pgfpathlineto{\pgfqpoint{0.537809in}{0.485392in}}%
\pgfpathlineto{\pgfqpoint{0.536479in}{0.491035in}}%
\pgfpathlineto{\pgfqpoint{0.536266in}{0.496684in}}%
\pgfpathlineto{\pgfqpoint{0.542217in}{0.495739in}}%
\pgfpathlineto{\pgfqpoint{0.543395in}{0.497404in}}%
\pgfpathlineto{\pgfqpoint{0.546355in}{0.497243in}}%
\pgfpathlineto{\pgfqpoint{0.549764in}{0.491217in}}%
\pgfpathlineto{\pgfqpoint{0.553051in}{0.486806in}}%
\pgfpathlineto{\pgfqpoint{0.550418in}{0.483935in}}%
\pgfpathlineto{\pgfqpoint{0.554970in}{0.482833in}}%
\pgfpathlineto{\pgfqpoint{0.556007in}{0.484241in}}%
\pgfpathlineto{\pgfqpoint{0.553573in}{0.489956in}}%
\pgfpathlineto{\pgfqpoint{0.548447in}{0.494692in}}%
\pgfpathlineto{\pgfqpoint{0.549074in}{0.496676in}}%
\pgfpathlineto{\pgfqpoint{0.544984in}{0.501691in}}%
\pgfpathlineto{\pgfqpoint{0.549161in}{0.502710in}}%
\pgfpathlineto{\pgfqpoint{0.545542in}{0.513448in}}%
\pgfpathlineto{\pgfqpoint{0.548484in}{0.514936in}}%
\pgfpathlineto{\pgfqpoint{0.549202in}{0.517511in}}%
\pgfpathlineto{\pgfqpoint{0.547288in}{0.520569in}}%
\pgfpathlineto{\pgfqpoint{0.551239in}{0.521335in}}%
\pgfpathlineto{\pgfqpoint{0.551948in}{0.523648in}}%
\pgfpathlineto{\pgfqpoint{0.556660in}{0.520176in}}%
\pgfpathlineto{\pgfqpoint{0.557707in}{0.523941in}}%
\pgfpathlineto{\pgfqpoint{0.560183in}{0.525597in}}%
\pgfpathlineto{\pgfqpoint{0.572233in}{0.527901in}}%
\pgfpathlineto{\pgfqpoint{0.575932in}{0.527086in}}%
\pgfpathlineto{\pgfqpoint{0.579031in}{0.531659in}}%
\pgfpathlineto{\pgfqpoint{0.586384in}{0.534658in}}%
\pgfpathlineto{\pgfqpoint{0.592047in}{0.533521in}}%
\pgfpathlineto{\pgfqpoint{0.594631in}{0.529538in}}%
\pgfpathlineto{\pgfqpoint{0.594495in}{0.523212in}}%
\pgfpathlineto{\pgfqpoint{0.597394in}{0.521778in}}%
\pgfpathlineto{\pgfqpoint{0.603306in}{0.521978in}}%
\pgfpathlineto{\pgfqpoint{0.610878in}{0.523624in}}%
\pgfpathlineto{\pgfqpoint{0.610785in}{0.521601in}}%
\pgfpathlineto{\pgfqpoint{0.620357in}{0.515832in}}%
\pgfpathlineto{\pgfqpoint{0.627052in}{0.517379in}}%
\pgfpathlineto{\pgfqpoint{0.629022in}{0.519381in}}%
\pgfpathlineto{\pgfqpoint{0.631133in}{0.524508in}}%
\pgfpathlineto{\pgfqpoint{0.633750in}{0.527792in}}%
\pgfpathlineto{\pgfqpoint{0.633984in}{0.532712in}}%
\pgfpathlineto{\pgfqpoint{0.631587in}{0.534647in}}%
\pgfpathlineto{\pgfqpoint{0.635136in}{0.536396in}}%
\pgfpathlineto{\pgfqpoint{0.641098in}{0.533768in}}%
\pgfpathlineto{\pgfqpoint{0.643040in}{0.535390in}}%
\pgfpathlineto{\pgfqpoint{0.643382in}{0.542352in}}%
\pgfpathlineto{\pgfqpoint{0.639528in}{0.540897in}}%
\pgfpathlineto{\pgfqpoint{0.636855in}{0.543309in}}%
\pgfpathlineto{\pgfqpoint{0.632516in}{0.544812in}}%
\pgfpathlineto{\pgfqpoint{0.625629in}{0.544561in}}%
\pgfpathlineto{\pgfqpoint{0.619030in}{0.546848in}}%
\pgfpathlineto{\pgfqpoint{0.618699in}{0.552545in}}%
\pgfpathlineto{\pgfqpoint{0.612797in}{0.558597in}}%
\pgfpathlineto{\pgfqpoint{0.605313in}{0.561138in}}%
\pgfpathlineto{\pgfqpoint{0.598780in}{0.572853in}}%
\pgfpathlineto{\pgfqpoint{0.599316in}{0.577189in}}%
\pgfpathlineto{\pgfqpoint{0.602833in}{0.579165in}}%
\pgfpathlineto{\pgfqpoint{0.602315in}{0.583269in}}%
\pgfpathlineto{\pgfqpoint{0.603132in}{0.587352in}}%
\pgfpathlineto{\pgfqpoint{0.605915in}{0.584557in}}%
\pgfpathlineto{\pgfqpoint{0.609590in}{0.585723in}}%
\pgfpathlineto{\pgfqpoint{0.609407in}{0.589270in}}%
\pgfpathlineto{\pgfqpoint{0.604311in}{0.596390in}}%
\pgfpathlineto{\pgfqpoint{0.602820in}{0.603734in}}%
\pgfpathlineto{\pgfqpoint{0.607085in}{0.604294in}}%
\pgfpathlineto{\pgfqpoint{0.613468in}{0.603843in}}%
\pgfpathlineto{\pgfqpoint{0.615478in}{0.601387in}}%
\pgfpathlineto{\pgfqpoint{0.621201in}{0.602266in}}%
\pgfpathlineto{\pgfqpoint{0.626520in}{0.601112in}}%
\pgfpathlineto{\pgfqpoint{0.630115in}{0.599083in}}%
\pgfpathlineto{\pgfqpoint{0.632604in}{0.600010in}}%
\pgfpathlineto{\pgfqpoint{0.639212in}{0.597983in}}%
\pgfpathlineto{\pgfqpoint{0.639273in}{0.596187in}}%
\pgfpathlineto{\pgfqpoint{0.644305in}{0.596794in}}%
\pgfpathlineto{\pgfqpoint{0.651053in}{0.592414in}}%
\pgfpathlineto{\pgfqpoint{0.650791in}{0.589506in}}%
\pgfpathlineto{\pgfqpoint{0.647580in}{0.588164in}}%
\pgfpathlineto{\pgfqpoint{0.644167in}{0.584972in}}%
\pgfpathlineto{\pgfqpoint{0.643755in}{0.580588in}}%
\pgfpathlineto{\pgfqpoint{0.651106in}{0.574997in}}%
\pgfpathlineto{\pgfqpoint{0.649930in}{0.572318in}}%
\pgfpathlineto{\pgfqpoint{0.655161in}{0.570260in}}%
\pgfpathlineto{\pgfqpoint{0.656445in}{0.566900in}}%
\pgfpathlineto{\pgfqpoint{0.663083in}{0.569570in}}%
\pgfpathlineto{\pgfqpoint{0.666007in}{0.566071in}}%
\pgfpathlineto{\pgfqpoint{0.667356in}{0.568189in}}%
\pgfpathlineto{\pgfqpoint{0.663383in}{0.575083in}}%
\pgfpathlineto{\pgfqpoint{0.664795in}{0.577148in}}%
\pgfpathlineto{\pgfqpoint{0.665421in}{0.582594in}}%
\pgfpathlineto{\pgfqpoint{0.663838in}{0.584991in}}%
\pgfpathlineto{\pgfqpoint{0.665025in}{0.588211in}}%
\pgfpathlineto{\pgfqpoint{0.668915in}{0.587908in}}%
\pgfpathlineto{\pgfqpoint{0.668110in}{0.582429in}}%
\pgfpathlineto{\pgfqpoint{0.665627in}{0.580602in}}%
\pgfpathlineto{\pgfqpoint{0.665803in}{0.573231in}}%
\pgfpathlineto{\pgfqpoint{0.670258in}{0.572350in}}%
\pgfpathlineto{\pgfqpoint{0.670312in}{0.566870in}}%
\pgfpathlineto{\pgfqpoint{0.675178in}{0.563315in}}%
\pgfpathlineto{\pgfqpoint{0.677415in}{0.565557in}}%
\pgfpathlineto{\pgfqpoint{0.674970in}{0.572380in}}%
\pgfpathlineto{\pgfqpoint{0.672442in}{0.574106in}}%
\pgfpathlineto{\pgfqpoint{0.669955in}{0.572970in}}%
\pgfpathlineto{\pgfqpoint{0.667767in}{0.574420in}}%
\pgfpathlineto{\pgfqpoint{0.668149in}{0.580725in}}%
\pgfpathlineto{\pgfqpoint{0.671800in}{0.583304in}}%
\pgfpathlineto{\pgfqpoint{0.675395in}{0.583066in}}%
\pgfpathlineto{\pgfqpoint{0.673905in}{0.586634in}}%
\pgfpathlineto{\pgfqpoint{0.668518in}{0.589366in}}%
\pgfpathlineto{\pgfqpoint{0.661383in}{0.600253in}}%
\pgfpathlineto{\pgfqpoint{0.664944in}{0.605475in}}%
\pgfpathlineto{\pgfqpoint{0.666926in}{0.610393in}}%
\pgfpathlineto{\pgfqpoint{0.667082in}{0.620269in}}%
\pgfpathlineto{\pgfqpoint{0.666319in}{0.628601in}}%
\pgfpathlineto{\pgfqpoint{0.664142in}{0.634072in}}%
\pgfpathlineto{\pgfqpoint{0.664576in}{0.639541in}}%
\pgfpathlineto{\pgfqpoint{0.670144in}{0.644285in}}%
\pgfpathlineto{\pgfqpoint{0.675681in}{0.649922in}}%
\pgfpathlineto{\pgfqpoint{0.689564in}{0.638405in}}%
\pgfpathlineto{\pgfqpoint{0.698640in}{0.637556in}}%
\pgfpathlineto{\pgfqpoint{0.706132in}{0.640424in}}%
\pgfpathlineto{\pgfqpoint{0.712708in}{0.644201in}}%
\pgfpathlineto{\pgfqpoint{0.714105in}{0.646000in}}%
\pgfpathlineto{\pgfqpoint{0.730340in}{0.650103in}}%
\pgfpathlineto{\pgfqpoint{0.731749in}{0.647098in}}%
\pgfpathlineto{\pgfqpoint{0.737266in}{0.644638in}}%
\pgfpathlineto{\pgfqpoint{0.747793in}{0.645392in}}%
\pgfpathlineto{\pgfqpoint{0.749716in}{0.644720in}}%
\pgfpathlineto{\pgfqpoint{0.754448in}{0.646515in}}%
\pgfpathlineto{\pgfqpoint{0.756350in}{0.643411in}}%
\pgfpathlineto{\pgfqpoint{0.760894in}{0.640575in}}%
\pgfpathlineto{\pgfqpoint{0.763167in}{0.640896in}}%
\pgfpathlineto{\pgfqpoint{0.767566in}{0.637488in}}%
\pgfpathlineto{\pgfqpoint{0.775000in}{0.638115in}}%
\pgfpathlineto{\pgfqpoint{0.782425in}{0.640874in}}%
\pgfpathlineto{\pgfqpoint{0.788489in}{0.632051in}}%
\pgfpathlineto{\pgfqpoint{0.788019in}{0.630071in}}%
\pgfpathlineto{\pgfqpoint{0.781328in}{0.627796in}}%
\pgfpathlineto{\pgfqpoint{0.783115in}{0.624355in}}%
\pgfpathlineto{\pgfqpoint{0.789275in}{0.627976in}}%
\pgfpathlineto{\pgfqpoint{0.791790in}{0.627100in}}%
\pgfpathlineto{\pgfqpoint{0.792860in}{0.622768in}}%
\pgfpathlineto{\pgfqpoint{0.789183in}{0.620936in}}%
\pgfpathlineto{\pgfqpoint{0.791809in}{0.615461in}}%
\pgfpathlineto{\pgfqpoint{0.795623in}{0.616386in}}%
\pgfpathlineto{\pgfqpoint{0.801157in}{0.613560in}}%
\pgfpathlineto{\pgfqpoint{0.806450in}{0.605718in}}%
\pgfpathlineto{\pgfqpoint{0.801793in}{0.603540in}}%
\pgfpathlineto{\pgfqpoint{0.805907in}{0.597761in}}%
\pgfpathlineto{\pgfqpoint{0.802441in}{0.596410in}}%
\pgfpathlineto{\pgfqpoint{0.806753in}{0.590922in}}%
\pgfpathlineto{\pgfqpoint{0.812012in}{0.591506in}}%
\pgfpathlineto{\pgfqpoint{0.814781in}{0.589037in}}%
\pgfpathlineto{\pgfqpoint{0.821785in}{0.585396in}}%
\pgfpathlineto{\pgfqpoint{0.830820in}{0.572470in}}%
\pgfpathlineto{\pgfqpoint{0.830616in}{0.569379in}}%
\pgfpathlineto{\pgfqpoint{0.844166in}{0.559477in}}%
\pgfpathlineto{\pgfqpoint{0.844458in}{0.555975in}}%
\pgfpathlineto{\pgfqpoint{0.848069in}{0.550735in}}%
\pgfpathlineto{\pgfqpoint{0.859163in}{0.548115in}}%
\pgfpathlineto{\pgfqpoint{0.863277in}{0.546252in}}%
\pgfpathlineto{\pgfqpoint{0.867148in}{0.540476in}}%
\pgfpathlineto{\pgfqpoint{0.867646in}{0.535415in}}%
\pgfpathlineto{\pgfqpoint{0.870973in}{0.531291in}}%
\pgfpathlineto{\pgfqpoint{0.871972in}{0.526640in}}%
\pgfpathlineto{\pgfqpoint{0.875090in}{0.524941in}}%
\pgfpathlineto{\pgfqpoint{0.859646in}{0.497067in}}%
\pgfpathlineto{\pgfqpoint{0.827590in}{0.439200in}}%
\pgfpathlineto{\pgfqpoint{0.780557in}{0.354280in}}%
\pgfpathlineto{\pgfqpoint{0.762143in}{0.321025in}}%
\pgfpathlineto{\pgfqpoint{0.766065in}{0.316559in}}%
\pgfpathlineto{\pgfqpoint{0.767793in}{0.317984in}}%
\pgfpathlineto{\pgfqpoint{0.771296in}{0.312785in}}%
\pgfpathlineto{\pgfqpoint{0.776159in}{0.314293in}}%
\pgfpathlineto{\pgfqpoint{0.782614in}{0.311178in}}%
\pgfpathlineto{\pgfqpoint{0.778433in}{0.306310in}}%
\pgfpathlineto{\pgfqpoint{0.781662in}{0.299741in}}%
\pgfpathlineto{\pgfqpoint{0.780995in}{0.296543in}}%
\pgfpathlineto{\pgfqpoint{0.786145in}{0.279677in}}%
\pgfpathlineto{\pgfqpoint{0.783966in}{0.272007in}}%
\pgfpathlineto{\pgfqpoint{0.793805in}{0.273906in}}%
\pgfpathlineto{\pgfqpoint{0.796221in}{0.272662in}}%
\pgfpathlineto{\pgfqpoint{0.798844in}{0.274685in}}%
\pgfpathlineto{\pgfqpoint{0.800709in}{0.278509in}}%
\pgfpathlineto{\pgfqpoint{0.803381in}{0.280701in}}%
\pgfpathlineto{\pgfqpoint{0.814784in}{0.280424in}}%
\pgfpathlineto{\pgfqpoint{0.817361in}{0.273064in}}%
\pgfpathlineto{\pgfqpoint{0.815149in}{0.266603in}}%
\pgfpathlineto{\pgfqpoint{0.817662in}{0.264471in}}%
\pgfpathlineto{\pgfqpoint{0.818983in}{0.260825in}}%
\pgfpathlineto{\pgfqpoint{0.818695in}{0.253967in}}%
\pgfpathlineto{\pgfqpoint{0.822048in}{0.249051in}}%
\pgfpathlineto{\pgfqpoint{0.823886in}{0.241346in}}%
\pgfpathlineto{\pgfqpoint{0.822942in}{0.237184in}}%
\pgfpathlineto{\pgfqpoint{0.824080in}{0.232966in}}%
\pgfpathlineto{\pgfqpoint{0.825722in}{0.208539in}}%
\pgfpathlineto{\pgfqpoint{0.823370in}{0.206615in}}%
\pgfpathlineto{\pgfqpoint{0.826472in}{0.203957in}}%
\pgfpathlineto{\pgfqpoint{0.824034in}{0.200723in}}%
\pgfpathlineto{\pgfqpoint{0.826233in}{0.198087in}}%
\pgfpathlineto{\pgfqpoint{0.824817in}{0.193586in}}%
\pgfpathlineto{\pgfqpoint{0.827984in}{0.192366in}}%
\pgfpathlineto{\pgfqpoint{0.832058in}{0.185522in}}%
\pgfpathlineto{\pgfqpoint{0.835085in}{0.183460in}}%
\pgfpathlineto{\pgfqpoint{0.835877in}{0.180445in}}%
\pgfpathlineto{\pgfqpoint{0.837631in}{0.179139in}}%
\pgfpathlineto{\pgfqpoint{0.837265in}{0.176636in}}%
\pgfpathlineto{\pgfqpoint{0.841046in}{0.174822in}}%
\pgfpathlineto{\pgfqpoint{0.840112in}{0.170006in}}%
\pgfpathlineto{\pgfqpoint{0.836858in}{0.167483in}}%
\pgfpathlineto{\pgfqpoint{0.835380in}{0.162795in}}%
\pgfpathlineto{\pgfqpoint{0.834988in}{0.156518in}}%
\pgfpathlineto{\pgfqpoint{0.833180in}{0.155917in}}%
\pgfpathlineto{\pgfqpoint{0.827271in}{0.150415in}}%
\pgfpathlineto{\pgfqpoint{0.822019in}{0.148962in}}%
\pgfpathlineto{\pgfqpoint{0.819513in}{0.151181in}}%
\pgfpathlineto{\pgfqpoint{0.821935in}{0.157071in}}%
\pgfpathlineto{\pgfqpoint{0.821195in}{0.159962in}}%
\pgfpathlineto{\pgfqpoint{0.824665in}{0.161426in}}%
\pgfpathlineto{\pgfqpoint{0.827909in}{0.169947in}}%
\pgfpathlineto{\pgfqpoint{0.826539in}{0.177225in}}%
\pgfpathlineto{\pgfqpoint{0.824346in}{0.178943in}}%
\pgfpathlineto{\pgfqpoint{0.817012in}{0.178387in}}%
\pgfpathlineto{\pgfqpoint{0.818490in}{0.176486in}}%
\pgfpathlineto{\pgfqpoint{0.813181in}{0.171028in}}%
\pgfpathlineto{\pgfqpoint{0.811471in}{0.174063in}}%
\pgfpathlineto{\pgfqpoint{0.811916in}{0.178027in}}%
\pgfpathlineto{\pgfqpoint{0.814963in}{0.177954in}}%
\pgfpathlineto{\pgfqpoint{0.817677in}{0.180866in}}%
\pgfpathlineto{\pgfqpoint{0.819282in}{0.185118in}}%
\pgfpathlineto{\pgfqpoint{0.825331in}{0.183918in}}%
\pgfpathlineto{\pgfqpoint{0.820488in}{0.186190in}}%
\pgfpathlineto{\pgfqpoint{0.820915in}{0.189233in}}%
\pgfpathlineto{\pgfqpoint{0.818899in}{0.190607in}}%
\pgfpathlineto{\pgfqpoint{0.818074in}{0.194752in}}%
\pgfpathlineto{\pgfqpoint{0.819505in}{0.196911in}}%
\pgfpathlineto{\pgfqpoint{0.816093in}{0.203617in}}%
\pgfpathlineto{\pgfqpoint{0.816381in}{0.207883in}}%
\pgfpathlineto{\pgfqpoint{0.812383in}{0.211682in}}%
\pgfpathlineto{\pgfqpoint{0.810220in}{0.214756in}}%
\pgfpathlineto{\pgfqpoint{0.813207in}{0.217685in}}%
\pgfpathlineto{\pgfqpoint{0.814313in}{0.222063in}}%
\pgfpathlineto{\pgfqpoint{0.813326in}{0.224937in}}%
\pgfpathlineto{\pgfqpoint{0.814481in}{0.228366in}}%
\pgfpathlineto{\pgfqpoint{0.813009in}{0.239543in}}%
\pgfpathlineto{\pgfqpoint{0.810807in}{0.244614in}}%
\pgfpathlineto{\pgfqpoint{0.808288in}{0.246519in}}%
\pgfpathlineto{\pgfqpoint{0.809360in}{0.253116in}}%
\pgfpathlineto{\pgfqpoint{0.808604in}{0.258268in}}%
\pgfpathlineto{\pgfqpoint{0.809974in}{0.261572in}}%
\pgfpathlineto{\pgfqpoint{0.807140in}{0.262035in}}%
\pgfpathlineto{\pgfqpoint{0.806380in}{0.254009in}}%
\pgfpathlineto{\pgfqpoint{0.803196in}{0.245055in}}%
\pgfpathlineto{\pgfqpoint{0.800639in}{0.247094in}}%
\pgfpathlineto{\pgfqpoint{0.800297in}{0.250037in}}%
\pgfpathlineto{\pgfqpoint{0.797168in}{0.254057in}}%
\pgfpathlineto{\pgfqpoint{0.798758in}{0.256688in}}%
\pgfpathlineto{\pgfqpoint{0.798264in}{0.262613in}}%
\pgfpathlineto{\pgfqpoint{0.794215in}{0.260530in}}%
\pgfpathlineto{\pgfqpoint{0.795009in}{0.257316in}}%
\pgfpathlineto{\pgfqpoint{0.794247in}{0.253236in}}%
\pgfpathlineto{\pgfqpoint{0.789702in}{0.253244in}}%
\pgfpathlineto{\pgfqpoint{0.787665in}{0.255773in}}%
\pgfpathlineto{\pgfqpoint{0.784590in}{0.256255in}}%
\pgfpathlineto{\pgfqpoint{0.782521in}{0.258978in}}%
\pgfpathlineto{\pgfqpoint{0.778961in}{0.266454in}}%
\pgfpathlineto{\pgfqpoint{0.777807in}{0.272259in}}%
\pgfpathlineto{\pgfqpoint{0.778546in}{0.274726in}}%
\pgfpathlineto{\pgfqpoint{0.776822in}{0.280160in}}%
\pgfpathlineto{\pgfqpoint{0.773967in}{0.283305in}}%
\pgfpathlineto{\pgfqpoint{0.769029in}{0.291344in}}%
\pgfpathlineto{\pgfqpoint{0.765447in}{0.298326in}}%
\pgfpathlineto{\pgfqpoint{0.769102in}{0.298431in}}%
\pgfpathlineto{\pgfqpoint{0.771252in}{0.299897in}}%
\pgfpathlineto{\pgfqpoint{0.771675in}{0.304435in}}%
\pgfpathlineto{\pgfqpoint{0.761334in}{0.304983in}}%
\pgfpathlineto{\pgfqpoint{0.752539in}{0.314807in}}%
\pgfpathlineto{\pgfqpoint{0.754852in}{0.317392in}}%
\pgfpathlineto{\pgfqpoint{0.750599in}{0.317833in}}%
\pgfpathlineto{\pgfqpoint{0.742152in}{0.326967in}}%
\pgfpathlineto{\pgfqpoint{0.728770in}{0.331547in}}%
\pgfpathlineto{\pgfqpoint{0.726968in}{0.336915in}}%
\pgfpathlineto{\pgfqpoint{0.724191in}{0.339711in}}%
\pgfpathlineto{\pgfqpoint{0.724065in}{0.342336in}}%
\pgfpathlineto{\pgfqpoint{0.721976in}{0.344286in}}%
\pgfpathlineto{\pgfqpoint{0.724717in}{0.347544in}}%
\pgfpathlineto{\pgfqpoint{0.718465in}{0.347858in}}%
\pgfpathlineto{\pgfqpoint{0.716876in}{0.352075in}}%
\pgfpathlineto{\pgfqpoint{0.714030in}{0.353870in}}%
\pgfpathlineto{\pgfqpoint{0.719729in}{0.356037in}}%
\pgfpathlineto{\pgfqpoint{0.716765in}{0.359902in}}%
\pgfpathlineto{\pgfqpoint{0.711311in}{0.361524in}}%
\pgfpathlineto{\pgfqpoint{0.712160in}{0.368663in}}%
\pgfpathlineto{\pgfqpoint{0.710596in}{0.371190in}}%
\pgfpathlineto{\pgfqpoint{0.707780in}{0.372970in}}%
\pgfpathlineto{\pgfqpoint{0.705372in}{0.371185in}}%
\pgfpathlineto{\pgfqpoint{0.703950in}{0.372864in}}%
\pgfpathlineto{\pgfqpoint{0.699800in}{0.373345in}}%
\pgfpathlineto{\pgfqpoint{0.700134in}{0.379823in}}%
\pgfpathlineto{\pgfqpoint{0.694119in}{0.375448in}}%
\pgfpathlineto{\pgfqpoint{0.694246in}{0.366430in}}%
\pgfpathlineto{\pgfqpoint{0.687788in}{0.364702in}}%
\pgfpathlineto{\pgfqpoint{0.683572in}{0.361101in}}%
\pgfpathlineto{\pgfqpoint{0.680182in}{0.360342in}}%
\pgfpathlineto{\pgfqpoint{0.676514in}{0.363971in}}%
\pgfpathlineto{\pgfqpoint{0.673248in}{0.363490in}}%
\pgfpathlineto{\pgfqpoint{0.673854in}{0.366805in}}%
\pgfpathlineto{\pgfqpoint{0.669211in}{0.365020in}}%
\pgfpathlineto{\pgfqpoint{0.664909in}{0.364820in}}%
\pgfpathlineto{\pgfqpoint{0.659874in}{0.363010in}}%
\pgfpathlineto{\pgfqpoint{0.649656in}{0.364028in}}%
\pgfpathlineto{\pgfqpoint{0.646376in}{0.363358in}}%
\pgfpathlineto{\pgfqpoint{0.644009in}{0.366339in}}%
\pgfpathlineto{\pgfqpoint{0.639640in}{0.366637in}}%
\pgfpathlineto{\pgfqpoint{0.638731in}{0.370500in}}%
\pgfpathlineto{\pgfqpoint{0.641335in}{0.372604in}}%
\pgfpathlineto{\pgfqpoint{0.644390in}{0.372406in}}%
\pgfpathlineto{\pgfqpoint{0.646940in}{0.369537in}}%
\pgfpathlineto{\pgfqpoint{0.648619in}{0.373944in}}%
\pgfpathlineto{\pgfqpoint{0.646276in}{0.378827in}}%
\pgfpathlineto{\pgfqpoint{0.651579in}{0.383125in}}%
\pgfpathlineto{\pgfqpoint{0.656893in}{0.384709in}}%
\pgfpathlineto{\pgfqpoint{0.660593in}{0.387292in}}%
\pgfpathlineto{\pgfqpoint{0.662990in}{0.390076in}}%
\pgfpathlineto{\pgfqpoint{0.664327in}{0.394638in}}%
\pgfpathlineto{\pgfqpoint{0.668552in}{0.393396in}}%
\pgfpathlineto{\pgfqpoint{0.677714in}{0.394034in}}%
\pgfpathlineto{\pgfqpoint{0.679680in}{0.389032in}}%
\pgfpathlineto{\pgfqpoint{0.682894in}{0.388814in}}%
\pgfpathlineto{\pgfqpoint{0.688742in}{0.381156in}}%
\pgfpathlineto{\pgfqpoint{0.688753in}{0.383947in}}%
\pgfpathlineto{\pgfqpoint{0.686710in}{0.384855in}}%
\pgfpathlineto{\pgfqpoint{0.682971in}{0.393996in}}%
\pgfpathlineto{\pgfqpoint{0.684790in}{0.395238in}}%
\pgfpathlineto{\pgfqpoint{0.679797in}{0.400007in}}%
\pgfpathlineto{\pgfqpoint{0.675334in}{0.400936in}}%
\pgfpathlineto{\pgfqpoint{0.671081in}{0.399287in}}%
\pgfpathlineto{\pgfqpoint{0.663798in}{0.400571in}}%
\pgfpathlineto{\pgfqpoint{0.660344in}{0.398261in}}%
\pgfpathlineto{\pgfqpoint{0.652478in}{0.396476in}}%
\pgfpathlineto{\pgfqpoint{0.651759in}{0.393907in}}%
\pgfpathlineto{\pgfqpoint{0.645716in}{0.393207in}}%
\pgfpathlineto{\pgfqpoint{0.644165in}{0.390270in}}%
\pgfpathlineto{\pgfqpoint{0.640709in}{0.388119in}}%
\pgfpathlineto{\pgfqpoint{0.636533in}{0.388489in}}%
\pgfpathlineto{\pgfqpoint{0.634448in}{0.386408in}}%
\pgfpathlineto{\pgfqpoint{0.631818in}{0.386506in}}%
\pgfpathlineto{\pgfqpoint{0.629304in}{0.388378in}}%
\pgfpathlineto{\pgfqpoint{0.624909in}{0.388044in}}%
\pgfpathlineto{\pgfqpoint{0.623874in}{0.386318in}}%
\pgfpathlineto{\pgfqpoint{0.619258in}{0.388090in}}%
\pgfpathlineto{\pgfqpoint{0.615339in}{0.383518in}}%
\pgfpathlineto{\pgfqpoint{0.619097in}{0.379095in}}%
\pgfpathlineto{\pgfqpoint{0.620543in}{0.375457in}}%
\pgfpathlineto{\pgfqpoint{0.619210in}{0.372376in}}%
\pgfpathlineto{\pgfqpoint{0.613870in}{0.370250in}}%
\pgfpathlineto{\pgfqpoint{0.610776in}{0.371980in}}%
\pgfpathlineto{\pgfqpoint{0.606882in}{0.371043in}}%
\pgfpathlineto{\pgfqpoint{0.606413in}{0.368343in}}%
\pgfpathlineto{\pgfqpoint{0.601673in}{0.369298in}}%
\pgfpathlineto{\pgfqpoint{0.602722in}{0.366600in}}%
\pgfpathlineto{\pgfqpoint{0.595731in}{0.366031in}}%
\pgfpathlineto{\pgfqpoint{0.591205in}{0.369456in}}%
\pgfpathlineto{\pgfqpoint{0.588787in}{0.367278in}}%
\pgfpathlineto{\pgfqpoint{0.582354in}{0.369597in}}%
\pgfpathlineto{\pgfqpoint{0.580752in}{0.367248in}}%
\pgfpathlineto{\pgfqpoint{0.577563in}{0.366203in}}%
\pgfpathlineto{\pgfqpoint{0.575005in}{0.368957in}}%
\pgfpathlineto{\pgfqpoint{0.566257in}{0.368272in}}%
\pgfpathlineto{\pgfqpoint{0.566266in}{0.364211in}}%
\pgfpathlineto{\pgfqpoint{0.562752in}{0.363136in}}%
\pgfpathlineto{\pgfqpoint{0.560018in}{0.365224in}}%
\pgfpathlineto{\pgfqpoint{0.556404in}{0.364064in}}%
\pgfpathlineto{\pgfqpoint{0.551727in}{0.363787in}}%
\pgfpathlineto{\pgfqpoint{0.544759in}{0.366724in}}%
\pgfpathlineto{\pgfqpoint{0.540769in}{0.363696in}}%
\pgfpathlineto{\pgfqpoint{0.535063in}{0.367743in}}%
\pgfpathlineto{\pgfqpoint{0.533197in}{0.366209in}}%
\pgfpathlineto{\pgfqpoint{0.531887in}{0.362246in}}%
\pgfpathlineto{\pgfqpoint{0.528218in}{0.359922in}}%
\pgfpathlineto{\pgfqpoint{0.522236in}{0.362582in}}%
\pgfpathlineto{\pgfqpoint{0.512326in}{0.369153in}}%
\pgfpathlineto{\pgfqpoint{0.509438in}{0.369737in}}%
\pgfpathlineto{\pgfqpoint{0.507808in}{0.368368in}}%
\pgfpathlineto{\pgfqpoint{0.503449in}{0.369299in}}%
\pgfpathlineto{\pgfqpoint{0.501322in}{0.371344in}}%
\pgfpathlineto{\pgfqpoint{0.496969in}{0.372446in}}%
\pgfpathlineto{\pgfqpoint{0.491023in}{0.372510in}}%
\pgfpathlineto{\pgfqpoint{0.488662in}{0.374950in}}%
\pgfpathlineto{\pgfqpoint{0.493187in}{0.377474in}}%
\pgfpathlineto{\pgfqpoint{0.492131in}{0.379966in}}%
\pgfpathlineto{\pgfqpoint{0.489022in}{0.379406in}}%
\pgfpathlineto{\pgfqpoint{0.487470in}{0.377447in}}%
\pgfpathlineto{\pgfqpoint{0.480192in}{0.374639in}}%
\pgfpathlineto{\pgfqpoint{0.476512in}{0.375594in}}%
\pgfpathlineto{\pgfqpoint{0.474964in}{0.381297in}}%
\pgfpathlineto{\pgfqpoint{0.470373in}{0.378845in}}%
\pgfpathlineto{\pgfqpoint{0.469039in}{0.386073in}}%
\pgfpathlineto{\pgfqpoint{0.466864in}{0.384210in}}%
\pgfpathlineto{\pgfqpoint{0.467228in}{0.381429in}}%
\pgfpathlineto{\pgfqpoint{0.462226in}{0.382163in}}%
\pgfpathlineto{\pgfqpoint{0.467429in}{0.387010in}}%
\pgfpathlineto{\pgfqpoint{0.472702in}{0.384130in}}%
\pgfpathlineto{\pgfqpoint{0.473833in}{0.385576in}}%
\pgfpathlineto{\pgfqpoint{0.479380in}{0.383843in}}%
\pgfpathlineto{\pgfqpoint{0.479033in}{0.385816in}}%
\pgfpathlineto{\pgfqpoint{0.494168in}{0.386469in}}%
\pgfpathlineto{\pgfqpoint{0.501551in}{0.383391in}}%
\pgfpathlineto{\pgfqpoint{0.504408in}{0.380355in}}%
\pgfpathlineto{\pgfqpoint{0.503561in}{0.377705in}}%
\pgfpathlineto{\pgfqpoint{0.506429in}{0.376702in}}%
\pgfpathlineto{\pgfqpoint{0.513558in}{0.380885in}}%
\pgfpathlineto{\pgfqpoint{0.522779in}{0.381153in}}%
\pgfpathlineto{\pgfqpoint{0.531155in}{0.378631in}}%
\pgfpathlineto{\pgfqpoint{0.535899in}{0.379054in}}%
\pgfpathlineto{\pgfqpoint{0.537665in}{0.375142in}}%
\pgfpathlineto{\pgfqpoint{0.540318in}{0.379455in}}%
\pgfpathlineto{\pgfqpoint{0.541950in}{0.380389in}}%
\pgfpathlineto{\pgfqpoint{0.549498in}{0.382264in}}%
\pgfpathlineto{\pgfqpoint{0.552304in}{0.381118in}}%
\pgfpathlineto{\pgfqpoint{0.555338in}{0.382479in}}%
\pgfpathlineto{\pgfqpoint{0.558958in}{0.381767in}}%
\pgfpathlineto{\pgfqpoint{0.559812in}{0.383352in}}%
\pgfpathlineto{\pgfqpoint{0.567363in}{0.390625in}}%
\pgfpathlineto{\pgfqpoint{0.570561in}{0.391703in}}%
\pgfpathlineto{\pgfqpoint{0.572451in}{0.395629in}}%
\pgfpathlineto{\pgfqpoint{0.582481in}{0.398668in}}%
\pgfpathlineto{\pgfqpoint{0.583737in}{0.400987in}}%
\pgfpathlineto{\pgfqpoint{0.569522in}{0.404566in}}%
\pgfpathlineto{\pgfqpoint{0.570081in}{0.408415in}}%
\pgfpathlineto{\pgfqpoint{0.568619in}{0.413499in}}%
\pgfpathlineto{\pgfqpoint{0.565309in}{0.413077in}}%
\pgfpathlineto{\pgfqpoint{0.563057in}{0.406479in}}%
\pgfpathlineto{\pgfqpoint{0.560443in}{0.405936in}}%
\pgfpathlineto{\pgfqpoint{0.558742in}{0.408071in}}%
\pgfpathlineto{\pgfqpoint{0.560726in}{0.417413in}}%
\pgfpathlineto{\pgfqpoint{0.558750in}{0.421328in}}%
\pgfpathlineto{\pgfqpoint{0.554922in}{0.425943in}}%
\pgfpathlineto{\pgfqpoint{0.556762in}{0.429465in}}%
\pgfpathlineto{\pgfqpoint{0.548683in}{0.429591in}}%
\pgfpathlineto{\pgfqpoint{0.548867in}{0.430958in}}%
\pgfpathlineto{\pgfqpoint{0.544046in}{0.432871in}}%
\pgfpathlineto{\pgfqpoint{0.542548in}{0.431818in}}%
\pgfpathclose%
\pgfusepath{fill}%
\end{pgfscope}%
\begin{pgfscope}%
\pgfpathrectangle{\pgfqpoint{0.100000in}{0.100000in}}{\pgfqpoint{2.862985in}{1.832778in}}%
\pgfusepath{clip}%
\pgfsetbuttcap%
\pgfsetmiterjoin%
\definecolor{currentfill}{rgb}{0.041522,0.760646,0.829758}%
\pgfsetfillcolor{currentfill}%
\pgfsetlinewidth{0.000000pt}%
\definecolor{currentstroke}{rgb}{0.000000,0.000000,0.000000}%
\pgfsetstrokecolor{currentstroke}%
\pgfsetstrokeopacity{0.000000}%
\pgfsetdash{}{0pt}%
\pgfpathmoveto{\pgfqpoint{0.528131in}{0.499862in}}%
\pgfpathlineto{\pgfqpoint{0.527132in}{0.498106in}}%
\pgfpathlineto{\pgfqpoint{0.529775in}{0.494368in}}%
\pgfpathlineto{\pgfqpoint{0.525597in}{0.490871in}}%
\pgfpathlineto{\pgfqpoint{0.525713in}{0.487863in}}%
\pgfpathlineto{\pgfqpoint{0.523824in}{0.487153in}}%
\pgfpathlineto{\pgfqpoint{0.517331in}{0.491639in}}%
\pgfpathlineto{\pgfqpoint{0.512703in}{0.501458in}}%
\pgfpathlineto{\pgfqpoint{0.512393in}{0.504295in}}%
\pgfpathlineto{\pgfqpoint{0.513583in}{0.507691in}}%
\pgfpathlineto{\pgfqpoint{0.518598in}{0.502496in}}%
\pgfpathlineto{\pgfqpoint{0.520132in}{0.503498in}}%
\pgfpathlineto{\pgfqpoint{0.524386in}{0.502545in}}%
\pgfpathlineto{\pgfqpoint{0.528131in}{0.499862in}}%
\pgfpathclose%
\pgfusepath{fill}%
\end{pgfscope}%
\begin{pgfscope}%
\pgfpathrectangle{\pgfqpoint{0.100000in}{0.100000in}}{\pgfqpoint{2.862985in}{1.832778in}}%
\pgfusepath{clip}%
\pgfsetbuttcap%
\pgfsetmiterjoin%
\definecolor{currentfill}{rgb}{0.041522,0.760646,0.829758}%
\pgfsetfillcolor{currentfill}%
\pgfsetlinewidth{0.000000pt}%
\definecolor{currentstroke}{rgb}{0.000000,0.000000,0.000000}%
\pgfsetstrokecolor{currentstroke}%
\pgfsetstrokeopacity{0.000000}%
\pgfsetdash{}{0pt}%
\pgfpathmoveto{\pgfqpoint{0.451126in}{0.386109in}}%
\pgfpathlineto{\pgfqpoint{0.446562in}{0.385368in}}%
\pgfpathlineto{\pgfqpoint{0.443260in}{0.386895in}}%
\pgfpathlineto{\pgfqpoint{0.441816in}{0.389380in}}%
\pgfpathlineto{\pgfqpoint{0.443453in}{0.392945in}}%
\pgfpathlineto{\pgfqpoint{0.447096in}{0.392002in}}%
\pgfpathlineto{\pgfqpoint{0.453086in}{0.394117in}}%
\pgfpathlineto{\pgfqpoint{0.455184in}{0.391196in}}%
\pgfpathlineto{\pgfqpoint{0.457217in}{0.391735in}}%
\pgfpathlineto{\pgfqpoint{0.462141in}{0.389917in}}%
\pgfpathlineto{\pgfqpoint{0.464268in}{0.387608in}}%
\pgfpathlineto{\pgfqpoint{0.461672in}{0.381567in}}%
\pgfpathlineto{\pgfqpoint{0.459090in}{0.380286in}}%
\pgfpathlineto{\pgfqpoint{0.456769in}{0.380976in}}%
\pgfpathlineto{\pgfqpoint{0.451126in}{0.386109in}}%
\pgfpathclose%
\pgfusepath{fill}%
\end{pgfscope}%
\begin{pgfscope}%
\pgfpathrectangle{\pgfqpoint{0.100000in}{0.100000in}}{\pgfqpoint{2.862985in}{1.832778in}}%
\pgfusepath{clip}%
\pgfsetbuttcap%
\pgfsetmiterjoin%
\definecolor{currentfill}{rgb}{0.041522,0.760646,0.829758}%
\pgfsetfillcolor{currentfill}%
\pgfsetlinewidth{0.000000pt}%
\definecolor{currentstroke}{rgb}{0.000000,0.000000,0.000000}%
\pgfsetstrokecolor{currentstroke}%
\pgfsetstrokeopacity{0.000000}%
\pgfsetdash{}{0pt}%
\pgfpathmoveto{\pgfqpoint{0.410323in}{0.391598in}}%
\pgfpathlineto{\pgfqpoint{0.403857in}{0.394682in}}%
\pgfpathlineto{\pgfqpoint{0.397873in}{0.395600in}}%
\pgfpathlineto{\pgfqpoint{0.395987in}{0.396913in}}%
\pgfpathlineto{\pgfqpoint{0.398135in}{0.399499in}}%
\pgfpathlineto{\pgfqpoint{0.401946in}{0.396269in}}%
\pgfpathlineto{\pgfqpoint{0.405142in}{0.396714in}}%
\pgfpathlineto{\pgfqpoint{0.410341in}{0.398951in}}%
\pgfpathlineto{\pgfqpoint{0.410761in}{0.401832in}}%
\pgfpathlineto{\pgfqpoint{0.417146in}{0.400406in}}%
\pgfpathlineto{\pgfqpoint{0.415651in}{0.396069in}}%
\pgfpathlineto{\pgfqpoint{0.419719in}{0.395751in}}%
\pgfpathlineto{\pgfqpoint{0.416306in}{0.390592in}}%
\pgfpathlineto{\pgfqpoint{0.413001in}{0.392264in}}%
\pgfpathlineto{\pgfqpoint{0.410323in}{0.391598in}}%
\pgfpathclose%
\pgfusepath{fill}%
\end{pgfscope}%
\begin{pgfscope}%
\pgfpathrectangle{\pgfqpoint{0.100000in}{0.100000in}}{\pgfqpoint{2.862985in}{1.832778in}}%
\pgfusepath{clip}%
\pgfsetbuttcap%
\pgfsetmiterjoin%
\definecolor{currentfill}{rgb}{0.041522,0.760646,0.829758}%
\pgfsetfillcolor{currentfill}%
\pgfsetlinewidth{0.000000pt}%
\definecolor{currentstroke}{rgb}{0.000000,0.000000,0.000000}%
\pgfsetstrokecolor{currentstroke}%
\pgfsetstrokeopacity{0.000000}%
\pgfsetdash{}{0pt}%
\pgfpathmoveto{\pgfqpoint{0.398962in}{0.402261in}}%
\pgfpathlineto{\pgfqpoint{0.396586in}{0.400969in}}%
\pgfpathlineto{\pgfqpoint{0.390110in}{0.402955in}}%
\pgfpathlineto{\pgfqpoint{0.385271in}{0.401805in}}%
\pgfpathlineto{\pgfqpoint{0.382813in}{0.405073in}}%
\pgfpathlineto{\pgfqpoint{0.383885in}{0.406569in}}%
\pgfpathlineto{\pgfqpoint{0.387552in}{0.406819in}}%
\pgfpathlineto{\pgfqpoint{0.390357in}{0.405880in}}%
\pgfpathlineto{\pgfqpoint{0.393414in}{0.407367in}}%
\pgfpathlineto{\pgfqpoint{0.398100in}{0.405378in}}%
\pgfpathlineto{\pgfqpoint{0.398962in}{0.402261in}}%
\pgfpathclose%
\pgfusepath{fill}%
\end{pgfscope}%
\begin{pgfscope}%
\pgfpathrectangle{\pgfqpoint{0.100000in}{0.100000in}}{\pgfqpoint{2.862985in}{1.832778in}}%
\pgfusepath{clip}%
\pgfsetbuttcap%
\pgfsetmiterjoin%
\definecolor{currentfill}{rgb}{0.041522,0.760646,0.829758}%
\pgfsetfillcolor{currentfill}%
\pgfsetlinewidth{0.000000pt}%
\definecolor{currentstroke}{rgb}{0.000000,0.000000,0.000000}%
\pgfsetstrokecolor{currentstroke}%
\pgfsetstrokeopacity{0.000000}%
\pgfsetdash{}{0pt}%
\pgfpathmoveto{\pgfqpoint{0.322628in}{0.450357in}}%
\pgfpathlineto{\pgfqpoint{0.322456in}{0.447153in}}%
\pgfpathlineto{\pgfqpoint{0.317729in}{0.445139in}}%
\pgfpathlineto{\pgfqpoint{0.313412in}{0.450180in}}%
\pgfpathlineto{\pgfqpoint{0.318253in}{0.451860in}}%
\pgfpathlineto{\pgfqpoint{0.322628in}{0.450357in}}%
\pgfpathclose%
\pgfusepath{fill}%
\end{pgfscope}%
\begin{pgfscope}%
\pgfpathrectangle{\pgfqpoint{0.100000in}{0.100000in}}{\pgfqpoint{2.862985in}{1.832778in}}%
\pgfusepath{clip}%
\pgfsetbuttcap%
\pgfsetmiterjoin%
\definecolor{currentfill}{rgb}{0.041522,0.760646,0.829758}%
\pgfsetfillcolor{currentfill}%
\pgfsetlinewidth{0.000000pt}%
\definecolor{currentstroke}{rgb}{0.000000,0.000000,0.000000}%
\pgfsetstrokecolor{currentstroke}%
\pgfsetstrokeopacity{0.000000}%
\pgfsetdash{}{0pt}%
\pgfpathmoveto{\pgfqpoint{0.283638in}{0.468371in}}%
\pgfpathlineto{\pgfqpoint{0.284642in}{0.469764in}}%
\pgfpathlineto{\pgfqpoint{0.290440in}{0.469524in}}%
\pgfpathlineto{\pgfqpoint{0.291458in}{0.471597in}}%
\pgfpathlineto{\pgfqpoint{0.294002in}{0.469788in}}%
\pgfpathlineto{\pgfqpoint{0.292533in}{0.466308in}}%
\pgfpathlineto{\pgfqpoint{0.290372in}{0.466014in}}%
\pgfpathlineto{\pgfqpoint{0.285207in}{0.467106in}}%
\pgfpathlineto{\pgfqpoint{0.283638in}{0.468371in}}%
\pgfpathclose%
\pgfusepath{fill}%
\end{pgfscope}%
\begin{pgfscope}%
\pgfpathrectangle{\pgfqpoint{0.100000in}{0.100000in}}{\pgfqpoint{2.862985in}{1.832778in}}%
\pgfusepath{clip}%
\pgfsetbuttcap%
\pgfsetmiterjoin%
\definecolor{currentfill}{rgb}{0.041522,0.760646,0.829758}%
\pgfsetfillcolor{currentfill}%
\pgfsetlinewidth{0.000000pt}%
\definecolor{currentstroke}{rgb}{0.000000,0.000000,0.000000}%
\pgfsetstrokecolor{currentstroke}%
\pgfsetstrokeopacity{0.000000}%
\pgfsetdash{}{0pt}%
\pgfpathmoveto{\pgfqpoint{0.275415in}{0.479156in}}%
\pgfpathlineto{\pgfqpoint{0.275063in}{0.482294in}}%
\pgfpathlineto{\pgfqpoint{0.278088in}{0.482062in}}%
\pgfpathlineto{\pgfqpoint{0.277878in}{0.486470in}}%
\pgfpathlineto{\pgfqpoint{0.280921in}{0.484091in}}%
\pgfpathlineto{\pgfqpoint{0.279868in}{0.480653in}}%
\pgfpathlineto{\pgfqpoint{0.275415in}{0.479156in}}%
\pgfpathclose%
\pgfusepath{fill}%
\end{pgfscope}%
\begin{pgfscope}%
\pgfpathrectangle{\pgfqpoint{0.100000in}{0.100000in}}{\pgfqpoint{2.862985in}{1.832778in}}%
\pgfusepath{clip}%
\pgfsetbuttcap%
\pgfsetmiterjoin%
\definecolor{currentfill}{rgb}{0.041522,0.760646,0.829758}%
\pgfsetfillcolor{currentfill}%
\pgfsetlinewidth{0.000000pt}%
\definecolor{currentstroke}{rgb}{0.000000,0.000000,0.000000}%
\pgfsetstrokecolor{currentstroke}%
\pgfsetstrokeopacity{0.000000}%
\pgfsetdash{}{0pt}%
\pgfpathmoveto{\pgfqpoint{0.541980in}{0.591016in}}%
\pgfpathlineto{\pgfqpoint{0.537154in}{0.592058in}}%
\pgfpathlineto{\pgfqpoint{0.536118in}{0.595283in}}%
\pgfpathlineto{\pgfqpoint{0.537776in}{0.598507in}}%
\pgfpathlineto{\pgfqpoint{0.541491in}{0.600190in}}%
\pgfpathlineto{\pgfqpoint{0.545744in}{0.591916in}}%
\pgfpathlineto{\pgfqpoint{0.549601in}{0.591548in}}%
\pgfpathlineto{\pgfqpoint{0.552479in}{0.588995in}}%
\pgfpathlineto{\pgfqpoint{0.552422in}{0.585482in}}%
\pgfpathlineto{\pgfqpoint{0.550318in}{0.583797in}}%
\pgfpathlineto{\pgfqpoint{0.553740in}{0.573408in}}%
\pgfpathlineto{\pgfqpoint{0.557484in}{0.569595in}}%
\pgfpathlineto{\pgfqpoint{0.553656in}{0.568385in}}%
\pgfpathlineto{\pgfqpoint{0.550557in}{0.572663in}}%
\pgfpathlineto{\pgfqpoint{0.543823in}{0.571338in}}%
\pgfpathlineto{\pgfqpoint{0.545918in}{0.575572in}}%
\pgfpathlineto{\pgfqpoint{0.544943in}{0.577258in}}%
\pgfpathlineto{\pgfqpoint{0.545514in}{0.581671in}}%
\pgfpathlineto{\pgfqpoint{0.544180in}{0.589648in}}%
\pgfpathlineto{\pgfqpoint{0.541980in}{0.591016in}}%
\pgfpathclose%
\pgfusepath{fill}%
\end{pgfscope}%
\begin{pgfscope}%
\pgfpathrectangle{\pgfqpoint{0.100000in}{0.100000in}}{\pgfqpoint{2.862985in}{1.832778in}}%
\pgfusepath{clip}%
\pgfsetbuttcap%
\pgfsetmiterjoin%
\definecolor{currentfill}{rgb}{0.041522,0.760646,0.829758}%
\pgfsetfillcolor{currentfill}%
\pgfsetlinewidth{0.000000pt}%
\definecolor{currentstroke}{rgb}{0.000000,0.000000,0.000000}%
\pgfsetstrokecolor{currentstroke}%
\pgfsetstrokeopacity{0.000000}%
\pgfsetdash{}{0pt}%
\pgfpathmoveto{\pgfqpoint{0.711519in}{0.353334in}}%
\pgfpathlineto{\pgfqpoint{0.704523in}{0.353412in}}%
\pgfpathlineto{\pgfqpoint{0.705062in}{0.356924in}}%
\pgfpathlineto{\pgfqpoint{0.707862in}{0.358180in}}%
\pgfpathlineto{\pgfqpoint{0.711519in}{0.353334in}}%
\pgfpathclose%
\pgfusepath{fill}%
\end{pgfscope}%
\begin{pgfscope}%
\pgfpathrectangle{\pgfqpoint{0.100000in}{0.100000in}}{\pgfqpoint{2.862985in}{1.832778in}}%
\pgfusepath{clip}%
\pgfsetbuttcap%
\pgfsetmiterjoin%
\definecolor{currentfill}{rgb}{0.041522,0.760646,0.829758}%
\pgfsetfillcolor{currentfill}%
\pgfsetlinewidth{0.000000pt}%
\definecolor{currentstroke}{rgb}{0.000000,0.000000,0.000000}%
\pgfsetstrokecolor{currentstroke}%
\pgfsetstrokeopacity{0.000000}%
\pgfsetdash{}{0pt}%
\pgfpathmoveto{\pgfqpoint{0.701812in}{0.356514in}}%
\pgfpathlineto{\pgfqpoint{0.696800in}{0.355301in}}%
\pgfpathlineto{\pgfqpoint{0.692038in}{0.353224in}}%
\pgfpathlineto{\pgfqpoint{0.690567in}{0.355671in}}%
\pgfpathlineto{\pgfqpoint{0.698757in}{0.357328in}}%
\pgfpathlineto{\pgfqpoint{0.701812in}{0.356514in}}%
\pgfpathclose%
\pgfusepath{fill}%
\end{pgfscope}%
\begin{pgfscope}%
\pgfpathrectangle{\pgfqpoint{0.100000in}{0.100000in}}{\pgfqpoint{2.862985in}{1.832778in}}%
\pgfusepath{clip}%
\pgfsetbuttcap%
\pgfsetmiterjoin%
\definecolor{currentfill}{rgb}{0.041522,0.760646,0.829758}%
\pgfsetfillcolor{currentfill}%
\pgfsetlinewidth{0.000000pt}%
\definecolor{currentstroke}{rgb}{0.000000,0.000000,0.000000}%
\pgfsetstrokecolor{currentstroke}%
\pgfsetstrokeopacity{0.000000}%
\pgfsetdash{}{0pt}%
\pgfpathmoveto{\pgfqpoint{0.624990in}{0.354028in}}%
\pgfpathlineto{\pgfqpoint{0.625182in}{0.351599in}}%
\pgfpathlineto{\pgfqpoint{0.622294in}{0.350090in}}%
\pgfpathlineto{\pgfqpoint{0.614235in}{0.353552in}}%
\pgfpathlineto{\pgfqpoint{0.612986in}{0.352159in}}%
\pgfpathlineto{\pgfqpoint{0.611036in}{0.358747in}}%
\pgfpathlineto{\pgfqpoint{0.615074in}{0.359664in}}%
\pgfpathlineto{\pgfqpoint{0.616940in}{0.356957in}}%
\pgfpathlineto{\pgfqpoint{0.620992in}{0.360345in}}%
\pgfpathlineto{\pgfqpoint{0.624990in}{0.354028in}}%
\pgfpathclose%
\pgfusepath{fill}%
\end{pgfscope}%
\begin{pgfscope}%
\pgfpathrectangle{\pgfqpoint{0.100000in}{0.100000in}}{\pgfqpoint{2.862985in}{1.832778in}}%
\pgfusepath{clip}%
\pgfsetbuttcap%
\pgfsetmiterjoin%
\definecolor{currentfill}{rgb}{0.041522,0.760646,0.829758}%
\pgfsetfillcolor{currentfill}%
\pgfsetlinewidth{0.000000pt}%
\definecolor{currentstroke}{rgb}{0.000000,0.000000,0.000000}%
\pgfsetstrokecolor{currentstroke}%
\pgfsetstrokeopacity{0.000000}%
\pgfsetdash{}{0pt}%
\pgfpathmoveto{\pgfqpoint{0.806774in}{0.219412in}}%
\pgfpathlineto{\pgfqpoint{0.800017in}{0.216896in}}%
\pgfpathlineto{\pgfqpoint{0.796437in}{0.216745in}}%
\pgfpathlineto{\pgfqpoint{0.799350in}{0.222082in}}%
\pgfpathlineto{\pgfqpoint{0.801840in}{0.224073in}}%
\pgfpathlineto{\pgfqpoint{0.801760in}{0.229600in}}%
\pgfpathlineto{\pgfqpoint{0.804465in}{0.239578in}}%
\pgfpathlineto{\pgfqpoint{0.806957in}{0.241473in}}%
\pgfpathlineto{\pgfqpoint{0.812710in}{0.238729in}}%
\pgfpathlineto{\pgfqpoint{0.811870in}{0.237158in}}%
\pgfpathlineto{\pgfqpoint{0.812306in}{0.229604in}}%
\pgfpathlineto{\pgfqpoint{0.810440in}{0.227566in}}%
\pgfpathlineto{\pgfqpoint{0.809640in}{0.222201in}}%
\pgfpathlineto{\pgfqpoint{0.806774in}{0.219412in}}%
\pgfpathclose%
\pgfusepath{fill}%
\end{pgfscope}%
\begin{pgfscope}%
\pgfpathrectangle{\pgfqpoint{0.100000in}{0.100000in}}{\pgfqpoint{2.862985in}{1.832778in}}%
\pgfusepath{clip}%
\pgfsetbuttcap%
\pgfsetmiterjoin%
\definecolor{currentfill}{rgb}{0.041522,0.760646,0.829758}%
\pgfsetfillcolor{currentfill}%
\pgfsetlinewidth{0.000000pt}%
\definecolor{currentstroke}{rgb}{0.000000,0.000000,0.000000}%
\pgfsetstrokecolor{currentstroke}%
\pgfsetstrokeopacity{0.000000}%
\pgfsetdash{}{0pt}%
\pgfpathmoveto{\pgfqpoint{0.791809in}{0.244146in}}%
\pgfpathlineto{\pgfqpoint{0.795768in}{0.237310in}}%
\pgfpathlineto{\pgfqpoint{0.795034in}{0.235680in}}%
\pgfpathlineto{\pgfqpoint{0.800329in}{0.234022in}}%
\pgfpathlineto{\pgfqpoint{0.798809in}{0.227738in}}%
\pgfpathlineto{\pgfqpoint{0.796503in}{0.228029in}}%
\pgfpathlineto{\pgfqpoint{0.790445in}{0.238982in}}%
\pgfpathlineto{\pgfqpoint{0.791307in}{0.234052in}}%
\pgfpathlineto{\pgfqpoint{0.790275in}{0.231213in}}%
\pgfpathlineto{\pgfqpoint{0.787709in}{0.229950in}}%
\pgfpathlineto{\pgfqpoint{0.786251in}{0.231311in}}%
\pgfpathlineto{\pgfqpoint{0.785437in}{0.237623in}}%
\pgfpathlineto{\pgfqpoint{0.785917in}{0.241717in}}%
\pgfpathlineto{\pgfqpoint{0.784208in}{0.243504in}}%
\pgfpathlineto{\pgfqpoint{0.788722in}{0.248675in}}%
\pgfpathlineto{\pgfqpoint{0.791684in}{0.250209in}}%
\pgfpathlineto{\pgfqpoint{0.792670in}{0.248193in}}%
\pgfpathlineto{\pgfqpoint{0.795760in}{0.249753in}}%
\pgfpathlineto{\pgfqpoint{0.797999in}{0.245516in}}%
\pgfpathlineto{\pgfqpoint{0.802840in}{0.240126in}}%
\pgfpathlineto{\pgfqpoint{0.802179in}{0.237684in}}%
\pgfpathlineto{\pgfqpoint{0.799491in}{0.235024in}}%
\pgfpathlineto{\pgfqpoint{0.797005in}{0.236278in}}%
\pgfpathlineto{\pgfqpoint{0.791809in}{0.244146in}}%
\pgfpathclose%
\pgfusepath{fill}%
\end{pgfscope}%
\begin{pgfscope}%
\pgfpathrectangle{\pgfqpoint{0.100000in}{0.100000in}}{\pgfqpoint{2.862985in}{1.832778in}}%
\pgfusepath{clip}%
\pgfsetbuttcap%
\pgfsetmiterjoin%
\definecolor{currentfill}{rgb}{0.041522,0.760646,0.829758}%
\pgfsetfillcolor{currentfill}%
\pgfsetlinewidth{0.000000pt}%
\definecolor{currentstroke}{rgb}{0.000000,0.000000,0.000000}%
\pgfsetstrokecolor{currentstroke}%
\pgfsetstrokeopacity{0.000000}%
\pgfsetdash{}{0pt}%
\pgfpathmoveto{\pgfqpoint{0.595353in}{0.341071in}}%
\pgfpathlineto{\pgfqpoint{0.590449in}{0.340433in}}%
\pgfpathlineto{\pgfqpoint{0.583542in}{0.337778in}}%
\pgfpathlineto{\pgfqpoint{0.581931in}{0.339221in}}%
\pgfpathlineto{\pgfqpoint{0.589146in}{0.340851in}}%
\pgfpathlineto{\pgfqpoint{0.587212in}{0.342273in}}%
\pgfpathlineto{\pgfqpoint{0.582739in}{0.343315in}}%
\pgfpathlineto{\pgfqpoint{0.581544in}{0.346649in}}%
\pgfpathlineto{\pgfqpoint{0.583795in}{0.349521in}}%
\pgfpathlineto{\pgfqpoint{0.585690in}{0.356112in}}%
\pgfpathlineto{\pgfqpoint{0.592614in}{0.357340in}}%
\pgfpathlineto{\pgfqpoint{0.595966in}{0.354492in}}%
\pgfpathlineto{\pgfqpoint{0.599115in}{0.357248in}}%
\pgfpathlineto{\pgfqpoint{0.603064in}{0.356572in}}%
\pgfpathlineto{\pgfqpoint{0.605902in}{0.351482in}}%
\pgfpathlineto{\pgfqpoint{0.608305in}{0.356006in}}%
\pgfpathlineto{\pgfqpoint{0.615126in}{0.345423in}}%
\pgfpathlineto{\pgfqpoint{0.612662in}{0.344577in}}%
\pgfpathlineto{\pgfqpoint{0.611150in}{0.342283in}}%
\pgfpathlineto{\pgfqpoint{0.614007in}{0.340190in}}%
\pgfpathlineto{\pgfqpoint{0.609576in}{0.338227in}}%
\pgfpathlineto{\pgfqpoint{0.606364in}{0.339308in}}%
\pgfpathlineto{\pgfqpoint{0.600296in}{0.342890in}}%
\pgfpathlineto{\pgfqpoint{0.595353in}{0.341071in}}%
\pgfpathclose%
\pgfusepath{fill}%
\end{pgfscope}%
\begin{pgfscope}%
\pgfpathrectangle{\pgfqpoint{0.100000in}{0.100000in}}{\pgfqpoint{2.862985in}{1.832778in}}%
\pgfusepath{clip}%
\pgfsetbuttcap%
\pgfsetmiterjoin%
\definecolor{currentfill}{rgb}{0.041522,0.760646,0.829758}%
\pgfsetfillcolor{currentfill}%
\pgfsetlinewidth{0.000000pt}%
\definecolor{currentstroke}{rgb}{0.000000,0.000000,0.000000}%
\pgfsetstrokecolor{currentstroke}%
\pgfsetstrokeopacity{0.000000}%
\pgfsetdash{}{0pt}%
\pgfpathmoveto{\pgfqpoint{0.787691in}{0.199220in}}%
\pgfpathlineto{\pgfqpoint{0.786719in}{0.208994in}}%
\pgfpathlineto{\pgfqpoint{0.789340in}{0.212583in}}%
\pgfpathlineto{\pgfqpoint{0.786538in}{0.212832in}}%
\pgfpathlineto{\pgfqpoint{0.787562in}{0.219812in}}%
\pgfpathlineto{\pgfqpoint{0.788770in}{0.223955in}}%
\pgfpathlineto{\pgfqpoint{0.787720in}{0.229522in}}%
\pgfpathlineto{\pgfqpoint{0.790287in}{0.230600in}}%
\pgfpathlineto{\pgfqpoint{0.791124in}{0.232039in}}%
\pgfpathlineto{\pgfqpoint{0.793653in}{0.231493in}}%
\pgfpathlineto{\pgfqpoint{0.795842in}{0.227033in}}%
\pgfpathlineto{\pgfqpoint{0.796249in}{0.222947in}}%
\pgfpathlineto{\pgfqpoint{0.793571in}{0.210709in}}%
\pgfpathlineto{\pgfqpoint{0.787691in}{0.199220in}}%
\pgfpathclose%
\pgfusepath{fill}%
\end{pgfscope}%
\begin{pgfscope}%
\pgfpathrectangle{\pgfqpoint{0.100000in}{0.100000in}}{\pgfqpoint{2.862985in}{1.832778in}}%
\pgfusepath{clip}%
\pgfsetbuttcap%
\pgfsetmiterjoin%
\definecolor{currentfill}{rgb}{0.041522,0.760646,0.829758}%
\pgfsetfillcolor{currentfill}%
\pgfsetlinewidth{0.000000pt}%
\definecolor{currentstroke}{rgb}{0.000000,0.000000,0.000000}%
\pgfsetstrokecolor{currentstroke}%
\pgfsetstrokeopacity{0.000000}%
\pgfsetdash{}{0pt}%
\pgfpathmoveto{\pgfqpoint{0.787110in}{0.229056in}}%
\pgfpathlineto{\pgfqpoint{0.786975in}{0.224290in}}%
\pgfpathlineto{\pgfqpoint{0.784778in}{0.222127in}}%
\pgfpathlineto{\pgfqpoint{0.782238in}{0.222926in}}%
\pgfpathlineto{\pgfqpoint{0.785448in}{0.229848in}}%
\pgfpathlineto{\pgfqpoint{0.787110in}{0.229056in}}%
\pgfpathclose%
\pgfusepath{fill}%
\end{pgfscope}%
\begin{pgfscope}%
\pgfpathrectangle{\pgfqpoint{0.100000in}{0.100000in}}{\pgfqpoint{2.862985in}{1.832778in}}%
\pgfusepath{clip}%
\pgfsetbuttcap%
\pgfsetmiterjoin%
\definecolor{currentfill}{rgb}{0.041522,0.760646,0.829758}%
\pgfsetfillcolor{currentfill}%
\pgfsetlinewidth{0.000000pt}%
\definecolor{currentstroke}{rgb}{0.000000,0.000000,0.000000}%
\pgfsetstrokecolor{currentstroke}%
\pgfsetstrokeopacity{0.000000}%
\pgfsetdash{}{0pt}%
\pgfpathmoveto{\pgfqpoint{0.814076in}{0.207898in}}%
\pgfpathlineto{\pgfqpoint{0.812202in}{0.199923in}}%
\pgfpathlineto{\pgfqpoint{0.807835in}{0.197402in}}%
\pgfpathlineto{\pgfqpoint{0.802107in}{0.199724in}}%
\pgfpathlineto{\pgfqpoint{0.803685in}{0.202824in}}%
\pgfpathlineto{\pgfqpoint{0.803653in}{0.204970in}}%
\pgfpathlineto{\pgfqpoint{0.805783in}{0.208616in}}%
\pgfpathlineto{\pgfqpoint{0.804286in}{0.207882in}}%
\pgfpathlineto{\pgfqpoint{0.805377in}{0.209663in}}%
\pgfpathlineto{\pgfqpoint{0.803304in}{0.213709in}}%
\pgfpathlineto{\pgfqpoint{0.805657in}{0.214455in}}%
\pgfpathlineto{\pgfqpoint{0.811104in}{0.209721in}}%
\pgfpathlineto{\pgfqpoint{0.814076in}{0.207898in}}%
\pgfpathclose%
\pgfusepath{fill}%
\end{pgfscope}%
\begin{pgfscope}%
\pgfpathrectangle{\pgfqpoint{0.100000in}{0.100000in}}{\pgfqpoint{2.862985in}{1.832778in}}%
\pgfusepath{clip}%
\pgfsetbuttcap%
\pgfsetmiterjoin%
\definecolor{currentfill}{rgb}{0.041522,0.760646,0.829758}%
\pgfsetfillcolor{currentfill}%
\pgfsetlinewidth{0.000000pt}%
\definecolor{currentstroke}{rgb}{0.000000,0.000000,0.000000}%
\pgfsetstrokecolor{currentstroke}%
\pgfsetstrokeopacity{0.000000}%
\pgfsetdash{}{0pt}%
\pgfpathmoveto{\pgfqpoint{0.794962in}{0.193796in}}%
\pgfpathlineto{\pgfqpoint{0.792492in}{0.192836in}}%
\pgfpathlineto{\pgfqpoint{0.793902in}{0.201252in}}%
\pgfpathlineto{\pgfqpoint{0.796866in}{0.203216in}}%
\pgfpathlineto{\pgfqpoint{0.796144in}{0.209521in}}%
\pgfpathlineto{\pgfqpoint{0.797348in}{0.212317in}}%
\pgfpathlineto{\pgfqpoint{0.799731in}{0.213365in}}%
\pgfpathlineto{\pgfqpoint{0.802232in}{0.210622in}}%
\pgfpathlineto{\pgfqpoint{0.804379in}{0.206901in}}%
\pgfpathlineto{\pgfqpoint{0.797478in}{0.199610in}}%
\pgfpathlineto{\pgfqpoint{0.794962in}{0.193796in}}%
\pgfpathclose%
\pgfusepath{fill}%
\end{pgfscope}%
\begin{pgfscope}%
\pgfpathrectangle{\pgfqpoint{0.100000in}{0.100000in}}{\pgfqpoint{2.862985in}{1.832778in}}%
\pgfusepath{clip}%
\pgfsetbuttcap%
\pgfsetmiterjoin%
\definecolor{currentfill}{rgb}{0.041522,0.760646,0.829758}%
\pgfsetfillcolor{currentfill}%
\pgfsetlinewidth{0.000000pt}%
\definecolor{currentstroke}{rgb}{0.000000,0.000000,0.000000}%
\pgfsetstrokecolor{currentstroke}%
\pgfsetstrokeopacity{0.000000}%
\pgfsetdash{}{0pt}%
\pgfpathmoveto{\pgfqpoint{0.815172in}{0.202568in}}%
\pgfpathlineto{\pgfqpoint{0.816382in}{0.196760in}}%
\pgfpathlineto{\pgfqpoint{0.810974in}{0.197425in}}%
\pgfpathlineto{\pgfqpoint{0.815172in}{0.202568in}}%
\pgfpathclose%
\pgfusepath{fill}%
\end{pgfscope}%
\begin{pgfscope}%
\pgfpathrectangle{\pgfqpoint{0.100000in}{0.100000in}}{\pgfqpoint{2.862985in}{1.832778in}}%
\pgfusepath{clip}%
\pgfsetbuttcap%
\pgfsetmiterjoin%
\definecolor{currentfill}{rgb}{0.041522,0.760646,0.829758}%
\pgfsetfillcolor{currentfill}%
\pgfsetlinewidth{0.000000pt}%
\definecolor{currentstroke}{rgb}{0.000000,0.000000,0.000000}%
\pgfsetstrokecolor{currentstroke}%
\pgfsetstrokeopacity{0.000000}%
\pgfsetdash{}{0pt}%
\pgfpathmoveto{\pgfqpoint{0.817752in}{0.193694in}}%
\pgfpathlineto{\pgfqpoint{0.818638in}{0.190111in}}%
\pgfpathlineto{\pgfqpoint{0.820335in}{0.188965in}}%
\pgfpathlineto{\pgfqpoint{0.820026in}{0.185593in}}%
\pgfpathlineto{\pgfqpoint{0.817656in}{0.184459in}}%
\pgfpathlineto{\pgfqpoint{0.815732in}{0.189493in}}%
\pgfpathlineto{\pgfqpoint{0.817752in}{0.193694in}}%
\pgfpathclose%
\pgfusepath{fill}%
\end{pgfscope}%
\begin{pgfscope}%
\pgfpathrectangle{\pgfqpoint{0.100000in}{0.100000in}}{\pgfqpoint{2.862985in}{1.832778in}}%
\pgfusepath{clip}%
\pgfsetbuttcap%
\pgfsetmiterjoin%
\definecolor{currentfill}{rgb}{0.041522,0.760646,0.829758}%
\pgfsetfillcolor{currentfill}%
\pgfsetlinewidth{0.000000pt}%
\definecolor{currentstroke}{rgb}{0.000000,0.000000,0.000000}%
\pgfsetstrokecolor{currentstroke}%
\pgfsetstrokeopacity{0.000000}%
\pgfsetdash{}{0pt}%
\pgfpathmoveto{\pgfqpoint{0.813239in}{0.195068in}}%
\pgfpathlineto{\pgfqpoint{0.811768in}{0.190907in}}%
\pgfpathlineto{\pgfqpoint{0.809484in}{0.191060in}}%
\pgfpathlineto{\pgfqpoint{0.808043in}{0.194492in}}%
\pgfpathlineto{\pgfqpoint{0.810390in}{0.196117in}}%
\pgfpathlineto{\pgfqpoint{0.813239in}{0.195068in}}%
\pgfpathclose%
\pgfusepath{fill}%
\end{pgfscope}%
\begin{pgfscope}%
\pgfpathrectangle{\pgfqpoint{0.100000in}{0.100000in}}{\pgfqpoint{2.862985in}{1.832778in}}%
\pgfusepath{clip}%
\pgfsetbuttcap%
\pgfsetmiterjoin%
\definecolor{currentfill}{rgb}{0.041522,0.760646,0.829758}%
\pgfsetfillcolor{currentfill}%
\pgfsetlinewidth{0.000000pt}%
\definecolor{currentstroke}{rgb}{0.000000,0.000000,0.000000}%
\pgfsetstrokecolor{currentstroke}%
\pgfsetstrokeopacity{0.000000}%
\pgfsetdash{}{0pt}%
\pgfpathmoveto{\pgfqpoint{0.807076in}{0.173613in}}%
\pgfpathlineto{\pgfqpoint{0.809114in}{0.171384in}}%
\pgfpathlineto{\pgfqpoint{0.809563in}{0.166686in}}%
\pgfpathlineto{\pgfqpoint{0.810636in}{0.165321in}}%
\pgfpathlineto{\pgfqpoint{0.808884in}{0.161070in}}%
\pgfpathlineto{\pgfqpoint{0.806578in}{0.158676in}}%
\pgfpathlineto{\pgfqpoint{0.805352in}{0.154236in}}%
\pgfpathlineto{\pgfqpoint{0.802832in}{0.161156in}}%
\pgfpathlineto{\pgfqpoint{0.802414in}{0.167439in}}%
\pgfpathlineto{\pgfqpoint{0.801004in}{0.170567in}}%
\pgfpathlineto{\pgfqpoint{0.796147in}{0.173102in}}%
\pgfpathlineto{\pgfqpoint{0.800703in}{0.178390in}}%
\pgfpathlineto{\pgfqpoint{0.797666in}{0.180650in}}%
\pgfpathlineto{\pgfqpoint{0.800341in}{0.182902in}}%
\pgfpathlineto{\pgfqpoint{0.803893in}{0.191549in}}%
\pgfpathlineto{\pgfqpoint{0.800212in}{0.194597in}}%
\pgfpathlineto{\pgfqpoint{0.801630in}{0.197524in}}%
\pgfpathlineto{\pgfqpoint{0.806388in}{0.194928in}}%
\pgfpathlineto{\pgfqpoint{0.806870in}{0.190831in}}%
\pgfpathlineto{\pgfqpoint{0.805070in}{0.188498in}}%
\pgfpathlineto{\pgfqpoint{0.808638in}{0.185341in}}%
\pgfpathlineto{\pgfqpoint{0.809761in}{0.180424in}}%
\pgfpathlineto{\pgfqpoint{0.809609in}{0.173528in}}%
\pgfpathlineto{\pgfqpoint{0.807076in}{0.173613in}}%
\pgfpathclose%
\pgfusepath{fill}%
\end{pgfscope}%
\begin{pgfscope}%
\pgfpathrectangle{\pgfqpoint{0.100000in}{0.100000in}}{\pgfqpoint{2.862985in}{1.832778in}}%
\pgfusepath{clip}%
\pgfsetbuttcap%
\pgfsetmiterjoin%
\definecolor{currentfill}{rgb}{0.041522,0.760646,0.829758}%
\pgfsetfillcolor{currentfill}%
\pgfsetlinewidth{0.000000pt}%
\definecolor{currentstroke}{rgb}{0.000000,0.000000,0.000000}%
\pgfsetstrokecolor{currentstroke}%
\pgfsetstrokeopacity{0.000000}%
\pgfsetdash{}{0pt}%
\pgfpathmoveto{\pgfqpoint{0.815789in}{0.190924in}}%
\pgfpathlineto{\pgfqpoint{0.814679in}{0.188532in}}%
\pgfpathlineto{\pgfqpoint{0.815722in}{0.184648in}}%
\pgfpathlineto{\pgfqpoint{0.813053in}{0.184321in}}%
\pgfpathlineto{\pgfqpoint{0.810255in}{0.186119in}}%
\pgfpathlineto{\pgfqpoint{0.811163in}{0.190080in}}%
\pgfpathlineto{\pgfqpoint{0.815789in}{0.190924in}}%
\pgfpathclose%
\pgfusepath{fill}%
\end{pgfscope}%
\begin{pgfscope}%
\pgfpathrectangle{\pgfqpoint{0.100000in}{0.100000in}}{\pgfqpoint{2.862985in}{1.832778in}}%
\pgfusepath{clip}%
\pgfsetbuttcap%
\pgfsetmiterjoin%
\definecolor{currentfill}{rgb}{0.041522,0.760646,0.829758}%
\pgfsetfillcolor{currentfill}%
\pgfsetlinewidth{0.000000pt}%
\definecolor{currentstroke}{rgb}{0.000000,0.000000,0.000000}%
\pgfsetstrokecolor{currentstroke}%
\pgfsetstrokeopacity{0.000000}%
\pgfsetdash{}{0pt}%
\pgfpathmoveto{\pgfqpoint{0.803351in}{0.191195in}}%
\pgfpathlineto{\pgfqpoint{0.797461in}{0.188553in}}%
\pgfpathlineto{\pgfqpoint{0.795459in}{0.189503in}}%
\pgfpathlineto{\pgfqpoint{0.799476in}{0.192989in}}%
\pgfpathlineto{\pgfqpoint{0.803351in}{0.191195in}}%
\pgfpathclose%
\pgfusepath{fill}%
\end{pgfscope}%
\begin{pgfscope}%
\pgfpathrectangle{\pgfqpoint{0.100000in}{0.100000in}}{\pgfqpoint{2.862985in}{1.832778in}}%
\pgfusepath{clip}%
\pgfsetbuttcap%
\pgfsetmiterjoin%
\definecolor{currentfill}{rgb}{0.041522,0.760646,0.829758}%
\pgfsetfillcolor{currentfill}%
\pgfsetlinewidth{0.000000pt}%
\definecolor{currentstroke}{rgb}{0.000000,0.000000,0.000000}%
\pgfsetstrokecolor{currentstroke}%
\pgfsetstrokeopacity{0.000000}%
\pgfsetdash{}{0pt}%
\pgfpathmoveto{\pgfqpoint{0.815854in}{0.165422in}}%
\pgfpathlineto{\pgfqpoint{0.814449in}{0.169400in}}%
\pgfpathlineto{\pgfqpoint{0.817419in}{0.170378in}}%
\pgfpathlineto{\pgfqpoint{0.818319in}{0.174668in}}%
\pgfpathlineto{\pgfqpoint{0.821538in}{0.174673in}}%
\pgfpathlineto{\pgfqpoint{0.821524in}{0.177660in}}%
\pgfpathlineto{\pgfqpoint{0.825997in}{0.176953in}}%
\pgfpathlineto{\pgfqpoint{0.826760in}{0.168856in}}%
\pgfpathlineto{\pgfqpoint{0.824165in}{0.163694in}}%
\pgfpathlineto{\pgfqpoint{0.820315in}{0.160469in}}%
\pgfpathlineto{\pgfqpoint{0.815854in}{0.165422in}}%
\pgfpathclose%
\pgfusepath{fill}%
\end{pgfscope}%
\begin{pgfscope}%
\pgfpathrectangle{\pgfqpoint{0.100000in}{0.100000in}}{\pgfqpoint{2.862985in}{1.832778in}}%
\pgfusepath{clip}%
\pgfsetbuttcap%
\pgfsetmiterjoin%
\definecolor{currentfill}{rgb}{0.041522,0.760646,0.829758}%
\pgfsetfillcolor{currentfill}%
\pgfsetlinewidth{0.000000pt}%
\definecolor{currentstroke}{rgb}{0.000000,0.000000,0.000000}%
\pgfsetstrokecolor{currentstroke}%
\pgfsetstrokeopacity{0.000000}%
\pgfsetdash{}{0pt}%
\pgfpathmoveto{\pgfqpoint{0.814114in}{0.168482in}}%
\pgfpathlineto{\pgfqpoint{0.815383in}{0.164715in}}%
\pgfpathlineto{\pgfqpoint{0.812255in}{0.164257in}}%
\pgfpathlineto{\pgfqpoint{0.814114in}{0.168482in}}%
\pgfpathclose%
\pgfusepath{fill}%
\end{pgfscope}%
\begin{pgfscope}%
\pgfpathrectangle{\pgfqpoint{0.100000in}{0.100000in}}{\pgfqpoint{2.862985in}{1.832778in}}%
\pgfusepath{clip}%
\pgfsetbuttcap%
\pgfsetmiterjoin%
\definecolor{currentfill}{rgb}{0.041522,0.760646,0.829758}%
\pgfsetfillcolor{currentfill}%
\pgfsetlinewidth{0.000000pt}%
\definecolor{currentstroke}{rgb}{0.000000,0.000000,0.000000}%
\pgfsetstrokecolor{currentstroke}%
\pgfsetstrokeopacity{0.000000}%
\pgfsetdash{}{0pt}%
\pgfpathmoveto{\pgfqpoint{0.816751in}{0.162942in}}%
\pgfpathlineto{\pgfqpoint{0.816315in}{0.158141in}}%
\pgfpathlineto{\pgfqpoint{0.813015in}{0.158538in}}%
\pgfpathlineto{\pgfqpoint{0.814083in}{0.160858in}}%
\pgfpathlineto{\pgfqpoint{0.816751in}{0.162942in}}%
\pgfpathclose%
\pgfusepath{fill}%
\end{pgfscope}%
\begin{pgfscope}%
\definecolor{textcolor}{rgb}{0.000000,0.000000,0.000000}%
\pgfsetstrokecolor{textcolor}%
\pgfsetfillcolor{textcolor}%
\pgftext[x=1.302454in,y=1.877794in,left,base]{\color{textcolor}\setmainfont{Lato}\rmfamily\fontsize{9.000000}{10.800000}\selectfont December 2023}%
\end{pgfscope}%
\begin{pgfscope}%
\pgfpathrectangle{\pgfqpoint{3.525000in}{0.100000in}}{\pgfqpoint{2.862985in}{1.832778in}}%
\pgfusepath{clip}%
\pgfsetbuttcap%
\pgfsetmiterjoin%
\definecolor{currentfill}{rgb}{0.099962,0.777009,0.590158}%
\pgfsetfillcolor{currentfill}%
\pgfsetlinewidth{0.000000pt}%
\definecolor{currentstroke}{rgb}{0.000000,0.000000,0.000000}%
\pgfsetstrokecolor{currentstroke}%
\pgfsetstrokeopacity{0.000000}%
\pgfsetdash{}{0pt}%
\pgfpathmoveto{\pgfqpoint{4.480609in}{0.529755in}}%
\pgfpathlineto{\pgfqpoint{4.472380in}{0.529889in}}%
\pgfpathlineto{\pgfqpoint{4.466780in}{0.536088in}}%
\pgfpathlineto{\pgfqpoint{4.462972in}{0.552550in}}%
\pgfpathlineto{\pgfqpoint{4.471822in}{0.554612in}}%
\pgfpathlineto{\pgfqpoint{4.480813in}{0.552394in}}%
\pgfpathlineto{\pgfqpoint{4.489963in}{0.545909in}}%
\pgfpathlineto{\pgfqpoint{4.489959in}{0.536940in}}%
\pgfpathlineto{\pgfqpoint{4.480609in}{0.529755in}}%
\pgfpathclose%
\pgfusepath{fill}%
\end{pgfscope}%
\begin{pgfscope}%
\pgfpathrectangle{\pgfqpoint{3.525000in}{0.100000in}}{\pgfqpoint{2.862985in}{1.832778in}}%
\pgfusepath{clip}%
\pgfsetbuttcap%
\pgfsetmiterjoin%
\definecolor{currentfill}{rgb}{0.099962,0.777009,0.590158}%
\pgfsetfillcolor{currentfill}%
\pgfsetlinewidth{0.000000pt}%
\definecolor{currentstroke}{rgb}{0.000000,0.000000,0.000000}%
\pgfsetstrokecolor{currentstroke}%
\pgfsetstrokeopacity{0.000000}%
\pgfsetdash{}{0pt}%
\pgfpathmoveto{\pgfqpoint{4.530694in}{0.432023in}}%
\pgfpathlineto{\pgfqpoint{4.521488in}{0.435528in}}%
\pgfpathlineto{\pgfqpoint{4.521446in}{0.442404in}}%
\pgfpathlineto{\pgfqpoint{4.510326in}{0.448219in}}%
\pgfpathlineto{\pgfqpoint{4.511680in}{0.462860in}}%
\pgfpathlineto{\pgfqpoint{4.514901in}{0.469769in}}%
\pgfpathlineto{\pgfqpoint{4.521701in}{0.463410in}}%
\pgfpathlineto{\pgfqpoint{4.532760in}{0.464472in}}%
\pgfpathlineto{\pgfqpoint{4.529967in}{0.444236in}}%
\pgfpathlineto{\pgfqpoint{4.530694in}{0.432023in}}%
\pgfpathclose%
\pgfusepath{fill}%
\end{pgfscope}%
\begin{pgfscope}%
\pgfpathrectangle{\pgfqpoint{3.525000in}{0.100000in}}{\pgfqpoint{2.862985in}{1.832778in}}%
\pgfusepath{clip}%
\pgfsetbuttcap%
\pgfsetmiterjoin%
\definecolor{currentfill}{rgb}{0.099962,0.777009,0.590158}%
\pgfsetfillcolor{currentfill}%
\pgfsetlinewidth{0.000000pt}%
\definecolor{currentstroke}{rgb}{0.000000,0.000000,0.000000}%
\pgfsetstrokecolor{currentstroke}%
\pgfsetstrokeopacity{0.000000}%
\pgfsetdash{}{0pt}%
\pgfpathmoveto{\pgfqpoint{4.569333in}{0.386887in}}%
\pgfpathlineto{\pgfqpoint{4.557651in}{0.388876in}}%
\pgfpathlineto{\pgfqpoint{4.551191in}{0.398877in}}%
\pgfpathlineto{\pgfqpoint{4.550279in}{0.407100in}}%
\pgfpathlineto{\pgfqpoint{4.569333in}{0.386887in}}%
\pgfpathclose%
\pgfusepath{fill}%
\end{pgfscope}%
\begin{pgfscope}%
\pgfpathrectangle{\pgfqpoint{3.525000in}{0.100000in}}{\pgfqpoint{2.862985in}{1.832778in}}%
\pgfusepath{clip}%
\pgfsetbuttcap%
\pgfsetmiterjoin%
\definecolor{currentfill}{rgb}{0.099962,0.777009,0.590158}%
\pgfsetfillcolor{currentfill}%
\pgfsetlinewidth{0.000000pt}%
\definecolor{currentstroke}{rgb}{0.000000,0.000000,0.000000}%
\pgfsetstrokecolor{currentstroke}%
\pgfsetstrokeopacity{0.000000}%
\pgfsetdash{}{0pt}%
\pgfpathmoveto{\pgfqpoint{4.545394in}{0.388631in}}%
\pgfpathlineto{\pgfqpoint{4.551519in}{0.383659in}}%
\pgfpathlineto{\pgfqpoint{4.552848in}{0.375110in}}%
\pgfpathlineto{\pgfqpoint{4.540519in}{0.375583in}}%
\pgfpathlineto{\pgfqpoint{4.545394in}{0.388631in}}%
\pgfpathclose%
\pgfusepath{fill}%
\end{pgfscope}%
\begin{pgfscope}%
\pgfpathrectangle{\pgfqpoint{3.525000in}{0.100000in}}{\pgfqpoint{2.862985in}{1.832778in}}%
\pgfusepath{clip}%
\pgfsetbuttcap%
\pgfsetmiterjoin%
\definecolor{currentfill}{rgb}{0.099962,0.777009,0.590158}%
\pgfsetfillcolor{currentfill}%
\pgfsetlinewidth{0.000000pt}%
\definecolor{currentstroke}{rgb}{0.000000,0.000000,0.000000}%
\pgfsetstrokecolor{currentstroke}%
\pgfsetstrokeopacity{0.000000}%
\pgfsetdash{}{0pt}%
\pgfpathmoveto{\pgfqpoint{4.573521in}{0.340880in}}%
\pgfpathlineto{\pgfqpoint{4.561217in}{0.345303in}}%
\pgfpathlineto{\pgfqpoint{4.564744in}{0.356026in}}%
\pgfpathlineto{\pgfqpoint{4.559586in}{0.366952in}}%
\pgfpathlineto{\pgfqpoint{4.561206in}{0.374463in}}%
\pgfpathlineto{\pgfqpoint{4.571107in}{0.377326in}}%
\pgfpathlineto{\pgfqpoint{4.570916in}{0.365292in}}%
\pgfpathlineto{\pgfqpoint{4.582627in}{0.358618in}}%
\pgfpathlineto{\pgfqpoint{4.589326in}{0.342480in}}%
\pgfpathlineto{\pgfqpoint{4.581861in}{0.336266in}}%
\pgfpathlineto{\pgfqpoint{4.573521in}{0.340880in}}%
\pgfpathclose%
\pgfusepath{fill}%
\end{pgfscope}%
\begin{pgfscope}%
\pgfpathrectangle{\pgfqpoint{3.525000in}{0.100000in}}{\pgfqpoint{2.862985in}{1.832778in}}%
\pgfusepath{clip}%
\pgfsetbuttcap%
\pgfsetmiterjoin%
\definecolor{currentfill}{rgb}{0.099962,0.777009,0.590158}%
\pgfsetfillcolor{currentfill}%
\pgfsetlinewidth{0.000000pt}%
\definecolor{currentstroke}{rgb}{0.000000,0.000000,0.000000}%
\pgfsetstrokecolor{currentstroke}%
\pgfsetstrokeopacity{0.000000}%
\pgfsetdash{}{0pt}%
\pgfpathmoveto{\pgfqpoint{4.529084in}{0.231768in}}%
\pgfpathlineto{\pgfqpoint{4.523222in}{0.243877in}}%
\pgfpathlineto{\pgfqpoint{4.525818in}{0.252022in}}%
\pgfpathlineto{\pgfqpoint{4.537315in}{0.263296in}}%
\pgfpathlineto{\pgfqpoint{4.538404in}{0.271938in}}%
\pgfpathlineto{\pgfqpoint{4.545928in}{0.291072in}}%
\pgfpathlineto{\pgfqpoint{4.565718in}{0.294341in}}%
\pgfpathlineto{\pgfqpoint{4.569506in}{0.305541in}}%
\pgfpathlineto{\pgfqpoint{4.578232in}{0.301341in}}%
\pgfpathlineto{\pgfqpoint{4.595961in}{0.269908in}}%
\pgfpathlineto{\pgfqpoint{4.596185in}{0.259645in}}%
\pgfpathlineto{\pgfqpoint{4.591173in}{0.250763in}}%
\pgfpathlineto{\pgfqpoint{4.597532in}{0.230369in}}%
\pgfpathlineto{\pgfqpoint{4.593089in}{0.227081in}}%
\pgfpathlineto{\pgfqpoint{4.577955in}{0.228037in}}%
\pgfpathlineto{\pgfqpoint{4.559756in}{0.234119in}}%
\pgfpathlineto{\pgfqpoint{4.544460in}{0.236631in}}%
\pgfpathlineto{\pgfqpoint{4.540375in}{0.233063in}}%
\pgfpathlineto{\pgfqpoint{4.529084in}{0.231768in}}%
\pgfpathclose%
\pgfusepath{fill}%
\end{pgfscope}%
\begin{pgfscope}%
\pgfpathrectangle{\pgfqpoint{3.525000in}{0.100000in}}{\pgfqpoint{2.862985in}{1.832778in}}%
\pgfusepath{clip}%
\pgfsetbuttcap%
\pgfsetmiterjoin%
\definecolor{currentfill}{rgb}{0.155325,0.792511,0.363168}%
\pgfsetfillcolor{currentfill}%
\pgfsetlinewidth{0.000000pt}%
\definecolor{currentstroke}{rgb}{0.000000,0.000000,0.000000}%
\pgfsetstrokecolor{currentstroke}%
\pgfsetstrokeopacity{0.000000}%
\pgfsetdash{}{0pt}%
\pgfpathmoveto{\pgfqpoint{4.205273in}{1.824988in}}%
\pgfpathlineto{\pgfqpoint{4.190426in}{1.766159in}}%
\pgfpathlineto{\pgfqpoint{4.179869in}{1.724007in}}%
\pgfpathlineto{\pgfqpoint{4.167317in}{1.673299in}}%
\pgfpathlineto{\pgfqpoint{4.164852in}{1.660801in}}%
\pgfpathlineto{\pgfqpoint{4.166462in}{1.649338in}}%
\pgfpathlineto{\pgfqpoint{4.164329in}{1.638779in}}%
\pgfpathlineto{\pgfqpoint{4.076709in}{1.661643in}}%
\pgfpathlineto{\pgfqpoint{4.068794in}{1.658987in}}%
\pgfpathlineto{\pgfqpoint{4.062021in}{1.661292in}}%
\pgfpathlineto{\pgfqpoint{4.030660in}{1.662045in}}%
\pgfpathlineto{\pgfqpoint{4.020053in}{1.659055in}}%
\pgfpathlineto{\pgfqpoint{4.009520in}{1.660065in}}%
\pgfpathlineto{\pgfqpoint{4.005055in}{1.664746in}}%
\pgfpathlineto{\pgfqpoint{3.976979in}{1.663725in}}%
\pgfpathlineto{\pgfqpoint{3.972545in}{1.671374in}}%
\pgfpathlineto{\pgfqpoint{3.964127in}{1.675404in}}%
\pgfpathlineto{\pgfqpoint{3.951880in}{1.677860in}}%
\pgfpathlineto{\pgfqpoint{3.930740in}{1.674264in}}%
\pgfpathlineto{\pgfqpoint{3.922937in}{1.677787in}}%
\pgfpathlineto{\pgfqpoint{3.910905in}{1.687156in}}%
\pgfpathlineto{\pgfqpoint{3.914522in}{1.705519in}}%
\pgfpathlineto{\pgfqpoint{3.913254in}{1.714897in}}%
\pgfpathlineto{\pgfqpoint{3.903718in}{1.724852in}}%
\pgfpathlineto{\pgfqpoint{3.897610in}{1.724234in}}%
\pgfpathlineto{\pgfqpoint{3.893229in}{1.734304in}}%
\pgfpathlineto{\pgfqpoint{3.882849in}{1.738374in}}%
\pgfpathlineto{\pgfqpoint{3.875305in}{1.737856in}}%
\pgfpathlineto{\pgfqpoint{3.872864in}{1.748198in}}%
\pgfpathlineto{\pgfqpoint{3.880380in}{1.747098in}}%
\pgfpathlineto{\pgfqpoint{3.879783in}{1.761472in}}%
\pgfpathlineto{\pgfqpoint{3.876507in}{1.770008in}}%
\pgfpathlineto{\pgfqpoint{3.877951in}{1.780932in}}%
\pgfpathlineto{\pgfqpoint{3.881884in}{1.789149in}}%
\pgfpathlineto{\pgfqpoint{3.881662in}{1.804741in}}%
\pgfpathlineto{\pgfqpoint{3.879563in}{1.810398in}}%
\pgfpathlineto{\pgfqpoint{3.883188in}{1.828548in}}%
\pgfpathlineto{\pgfqpoint{3.882142in}{1.840246in}}%
\pgfpathlineto{\pgfqpoint{3.878517in}{1.845854in}}%
\pgfpathlineto{\pgfqpoint{3.879070in}{1.864192in}}%
\pgfpathlineto{\pgfqpoint{3.884342in}{1.877678in}}%
\pgfpathlineto{\pgfqpoint{3.890077in}{1.874378in}}%
\pgfpathlineto{\pgfqpoint{3.909064in}{1.854764in}}%
\pgfpathlineto{\pgfqpoint{3.932071in}{1.844031in}}%
\pgfpathlineto{\pgfqpoint{3.943852in}{1.842761in}}%
\pgfpathlineto{\pgfqpoint{3.954951in}{1.838191in}}%
\pgfpathlineto{\pgfqpoint{3.958101in}{1.822814in}}%
\pgfpathlineto{\pgfqpoint{3.948375in}{1.819901in}}%
\pgfpathlineto{\pgfqpoint{3.941929in}{1.807782in}}%
\pgfpathlineto{\pgfqpoint{3.949497in}{1.808455in}}%
\pgfpathlineto{\pgfqpoint{3.952549in}{1.813895in}}%
\pgfpathlineto{\pgfqpoint{3.963221in}{1.820697in}}%
\pgfpathlineto{\pgfqpoint{3.962670in}{1.810596in}}%
\pgfpathlineto{\pgfqpoint{3.955544in}{1.808967in}}%
\pgfpathlineto{\pgfqpoint{3.956727in}{1.795983in}}%
\pgfpathlineto{\pgfqpoint{3.948050in}{1.782976in}}%
\pgfpathlineto{\pgfqpoint{3.942496in}{1.789161in}}%
\pgfpathlineto{\pgfqpoint{3.925991in}{1.785798in}}%
\pgfpathlineto{\pgfqpoint{3.925123in}{1.778232in}}%
\pgfpathlineto{\pgfqpoint{3.930813in}{1.773622in}}%
\pgfpathlineto{\pgfqpoint{3.939493in}{1.773169in}}%
\pgfpathlineto{\pgfqpoint{3.951495in}{1.783019in}}%
\pgfpathlineto{\pgfqpoint{3.960994in}{1.783837in}}%
\pgfpathlineto{\pgfqpoint{3.961358in}{1.794741in}}%
\pgfpathlineto{\pgfqpoint{3.966254in}{1.810756in}}%
\pgfpathlineto{\pgfqpoint{3.977856in}{1.822666in}}%
\pgfpathlineto{\pgfqpoint{3.973988in}{1.831891in}}%
\pgfpathlineto{\pgfqpoint{3.976630in}{1.841823in}}%
\pgfpathlineto{\pgfqpoint{3.971483in}{1.851722in}}%
\pgfpathlineto{\pgfqpoint{3.978818in}{1.864381in}}%
\pgfpathlineto{\pgfqpoint{3.979917in}{1.871996in}}%
\pgfpathlineto{\pgfqpoint{3.973548in}{1.877001in}}%
\pgfpathlineto{\pgfqpoint{3.974609in}{1.889804in}}%
\pgfpathlineto{\pgfqpoint{4.050918in}{1.866620in}}%
\pgfpathlineto{\pgfqpoint{4.131952in}{1.843896in}}%
\pgfpathlineto{\pgfqpoint{4.205273in}{1.824988in}}%
\pgfpathclose%
\pgfusepath{fill}%
\end{pgfscope}%
\begin{pgfscope}%
\pgfpathrectangle{\pgfqpoint{3.525000in}{0.100000in}}{\pgfqpoint{2.862985in}{1.832778in}}%
\pgfusepath{clip}%
\pgfsetbuttcap%
\pgfsetmiterjoin%
\definecolor{currentfill}{rgb}{0.155325,0.792511,0.363168}%
\pgfsetfillcolor{currentfill}%
\pgfsetlinewidth{0.000000pt}%
\definecolor{currentstroke}{rgb}{0.000000,0.000000,0.000000}%
\pgfsetstrokecolor{currentstroke}%
\pgfsetstrokeopacity{0.000000}%
\pgfsetdash{}{0pt}%
\pgfpathmoveto{\pgfqpoint{3.962056in}{1.845512in}}%
\pgfpathlineto{\pgfqpoint{3.965813in}{1.831026in}}%
\pgfpathlineto{\pgfqpoint{3.971706in}{1.826214in}}%
\pgfpathlineto{\pgfqpoint{3.967004in}{1.820089in}}%
\pgfpathlineto{\pgfqpoint{3.962421in}{1.829064in}}%
\pgfpathlineto{\pgfqpoint{3.962056in}{1.845512in}}%
\pgfpathclose%
\pgfusepath{fill}%
\end{pgfscope}%
\begin{pgfscope}%
\pgfpathrectangle{\pgfqpoint{3.525000in}{0.100000in}}{\pgfqpoint{2.862985in}{1.832778in}}%
\pgfusepath{clip}%
\pgfsetbuttcap%
\pgfsetmiterjoin%
\definecolor{currentfill}{rgb}{0.170704,0.796817,0.300115}%
\pgfsetfillcolor{currentfill}%
\pgfsetlinewidth{0.000000pt}%
\definecolor{currentstroke}{rgb}{0.000000,0.000000,0.000000}%
\pgfsetstrokecolor{currentstroke}%
\pgfsetstrokeopacity{0.000000}%
\pgfsetdash{}{0pt}%
\pgfpathmoveto{\pgfqpoint{4.244624in}{1.815480in}}%
\pgfpathlineto{\pgfqpoint{4.326194in}{1.797175in}}%
\pgfpathlineto{\pgfqpoint{4.403016in}{1.781647in}}%
\pgfpathlineto{\pgfqpoint{4.462121in}{1.770817in}}%
\pgfpathlineto{\pgfqpoint{4.513650in}{1.762160in}}%
\pgfpathlineto{\pgfqpoint{4.565293in}{1.754214in}}%
\pgfpathlineto{\pgfqpoint{4.609272in}{1.748018in}}%
\pgfpathlineto{\pgfqpoint{4.653320in}{1.742337in}}%
\pgfpathlineto{\pgfqpoint{4.697432in}{1.737171in}}%
\pgfpathlineto{\pgfqpoint{4.739001in}{1.732781in}}%
\pgfpathlineto{\pgfqpoint{4.733284in}{1.669428in}}%
\pgfpathlineto{\pgfqpoint{4.724775in}{1.583815in}}%
\pgfpathlineto{\pgfqpoint{4.720312in}{1.539783in}}%
\pgfpathlineto{\pgfqpoint{4.713893in}{1.480441in}}%
\pgfpathlineto{\pgfqpoint{4.668453in}{1.485403in}}%
\pgfpathlineto{\pgfqpoint{4.616429in}{1.491308in}}%
\pgfpathlineto{\pgfqpoint{4.544194in}{1.501195in}}%
\pgfpathlineto{\pgfqpoint{4.511931in}{1.505797in}}%
\pgfpathlineto{\pgfqpoint{4.432445in}{1.518274in}}%
\pgfpathlineto{\pgfqpoint{4.405083in}{1.523280in}}%
\pgfpathlineto{\pgfqpoint{4.399380in}{1.490854in}}%
\pgfpathlineto{\pgfqpoint{4.396262in}{1.493168in}}%
\pgfpathlineto{\pgfqpoint{4.390403in}{1.508767in}}%
\pgfpathlineto{\pgfqpoint{4.382871in}{1.507806in}}%
\pgfpathlineto{\pgfqpoint{4.377480in}{1.499920in}}%
\pgfpathlineto{\pgfqpoint{4.366293in}{1.498866in}}%
\pgfpathlineto{\pgfqpoint{4.364045in}{1.502638in}}%
\pgfpathlineto{\pgfqpoint{4.353566in}{1.502181in}}%
\pgfpathlineto{\pgfqpoint{4.348297in}{1.505835in}}%
\pgfpathlineto{\pgfqpoint{4.340931in}{1.500183in}}%
\pgfpathlineto{\pgfqpoint{4.323038in}{1.505269in}}%
\pgfpathlineto{\pgfqpoint{4.315230in}{1.502538in}}%
\pgfpathlineto{\pgfqpoint{4.311029in}{1.514658in}}%
\pgfpathlineto{\pgfqpoint{4.310833in}{1.526190in}}%
\pgfpathlineto{\pgfqpoint{4.298499in}{1.534484in}}%
\pgfpathlineto{\pgfqpoint{4.300510in}{1.546436in}}%
\pgfpathlineto{\pgfqpoint{4.291728in}{1.566180in}}%
\pgfpathlineto{\pgfqpoint{4.292319in}{1.582964in}}%
\pgfpathlineto{\pgfqpoint{4.284772in}{1.591178in}}%
\pgfpathlineto{\pgfqpoint{4.277690in}{1.583911in}}%
\pgfpathlineto{\pgfqpoint{4.268050in}{1.579980in}}%
\pgfpathlineto{\pgfqpoint{4.260712in}{1.588082in}}%
\pgfpathlineto{\pgfqpoint{4.261905in}{1.601019in}}%
\pgfpathlineto{\pgfqpoint{4.270705in}{1.606116in}}%
\pgfpathlineto{\pgfqpoint{4.268366in}{1.614312in}}%
\pgfpathlineto{\pgfqpoint{4.283663in}{1.653985in}}%
\pgfpathlineto{\pgfqpoint{4.271601in}{1.654979in}}%
\pgfpathlineto{\pgfqpoint{4.270361in}{1.662098in}}%
\pgfpathlineto{\pgfqpoint{4.261575in}{1.668238in}}%
\pgfpathlineto{\pgfqpoint{4.262137in}{1.675101in}}%
\pgfpathlineto{\pgfqpoint{4.257502in}{1.680433in}}%
\pgfpathlineto{\pgfqpoint{4.249957in}{1.699791in}}%
\pgfpathlineto{\pgfqpoint{4.243063in}{1.703739in}}%
\pgfpathlineto{\pgfqpoint{4.238126in}{1.720446in}}%
\pgfpathlineto{\pgfqpoint{4.239291in}{1.731551in}}%
\pgfpathlineto{\pgfqpoint{4.230080in}{1.752004in}}%
\pgfpathlineto{\pgfqpoint{4.244624in}{1.815480in}}%
\pgfpathclose%
\pgfusepath{fill}%
\end{pgfscope}%
\begin{pgfscope}%
\pgfpathrectangle{\pgfqpoint{3.525000in}{0.100000in}}{\pgfqpoint{2.862985in}{1.832778in}}%
\pgfusepath{clip}%
\pgfsetbuttcap%
\pgfsetmiterjoin%
\definecolor{currentfill}{rgb}{0.161476,0.794233,0.337947}%
\pgfsetfillcolor{currentfill}%
\pgfsetlinewidth{0.000000pt}%
\definecolor{currentstroke}{rgb}{0.000000,0.000000,0.000000}%
\pgfsetstrokecolor{currentstroke}%
\pgfsetstrokeopacity{0.000000}%
\pgfsetdash{}{0pt}%
\pgfpathmoveto{\pgfqpoint{6.213801in}{1.489012in}}%
\pgfpathlineto{\pgfqpoint{6.212168in}{1.495954in}}%
\pgfpathlineto{\pgfqpoint{6.203094in}{1.502035in}}%
\pgfpathlineto{\pgfqpoint{6.179117in}{1.580468in}}%
\pgfpathlineto{\pgfqpoint{6.166192in}{1.618318in}}%
\pgfpathlineto{\pgfqpoint{6.176583in}{1.629222in}}%
\pgfpathlineto{\pgfqpoint{6.186683in}{1.655908in}}%
\pgfpathlineto{\pgfqpoint{6.191720in}{1.664437in}}%
\pgfpathlineto{\pgfqpoint{6.188152in}{1.668018in}}%
\pgfpathlineto{\pgfqpoint{6.186957in}{1.691707in}}%
\pgfpathlineto{\pgfqpoint{6.191463in}{1.698969in}}%
\pgfpathlineto{\pgfqpoint{6.189488in}{1.715849in}}%
\pgfpathlineto{\pgfqpoint{6.207474in}{1.771228in}}%
\pgfpathlineto{\pgfqpoint{6.215578in}{1.771530in}}%
\pgfpathlineto{\pgfqpoint{6.218942in}{1.761638in}}%
\pgfpathlineto{\pgfqpoint{6.226144in}{1.758798in}}%
\pgfpathlineto{\pgfqpoint{6.239749in}{1.770433in}}%
\pgfpathlineto{\pgfqpoint{6.250417in}{1.777306in}}%
\pgfpathlineto{\pgfqpoint{6.273922in}{1.765204in}}%
\pgfpathlineto{\pgfqpoint{6.295157in}{1.698008in}}%
\pgfpathlineto{\pgfqpoint{6.299194in}{1.681476in}}%
\pgfpathlineto{\pgfqpoint{6.308495in}{1.679537in}}%
\pgfpathlineto{\pgfqpoint{6.321110in}{1.668309in}}%
\pgfpathlineto{\pgfqpoint{6.320396in}{1.661767in}}%
\pgfpathlineto{\pgfqpoint{6.328989in}{1.654013in}}%
\pgfpathlineto{\pgfqpoint{6.335970in}{1.658462in}}%
\pgfpathlineto{\pgfqpoint{6.350577in}{1.641426in}}%
\pgfpathlineto{\pgfqpoint{6.344049in}{1.627792in}}%
\pgfpathlineto{\pgfqpoint{6.335294in}{1.627522in}}%
\pgfpathlineto{\pgfqpoint{6.328271in}{1.615351in}}%
\pgfpathlineto{\pgfqpoint{6.319762in}{1.613625in}}%
\pgfpathlineto{\pgfqpoint{6.313516in}{1.607145in}}%
\pgfpathlineto{\pgfqpoint{6.294860in}{1.600195in}}%
\pgfpathlineto{\pgfqpoint{6.283568in}{1.588995in}}%
\pgfpathlineto{\pgfqpoint{6.277832in}{1.597034in}}%
\pgfpathlineto{\pgfqpoint{6.272620in}{1.591272in}}%
\pgfpathlineto{\pgfqpoint{6.274047in}{1.567979in}}%
\pgfpathlineto{\pgfqpoint{6.269920in}{1.558753in}}%
\pgfpathlineto{\pgfqpoint{6.260918in}{1.561273in}}%
\pgfpathlineto{\pgfqpoint{6.259465in}{1.551814in}}%
\pgfpathlineto{\pgfqpoint{6.255576in}{1.547924in}}%
\pgfpathlineto{\pgfqpoint{6.247791in}{1.548872in}}%
\pgfpathlineto{\pgfqpoint{6.248281in}{1.540028in}}%
\pgfpathlineto{\pgfqpoint{6.236482in}{1.542540in}}%
\pgfpathlineto{\pgfqpoint{6.230038in}{1.530287in}}%
\pgfpathlineto{\pgfqpoint{6.232472in}{1.523875in}}%
\pgfpathlineto{\pgfqpoint{6.228639in}{1.513212in}}%
\pgfpathlineto{\pgfqpoint{6.222604in}{1.505369in}}%
\pgfpathlineto{\pgfqpoint{6.221086in}{1.488995in}}%
\pgfpathlineto{\pgfqpoint{6.213801in}{1.489012in}}%
\pgfpathclose%
\pgfusepath{fill}%
\end{pgfscope}%
\begin{pgfscope}%
\pgfpathrectangle{\pgfqpoint{3.525000in}{0.100000in}}{\pgfqpoint{2.862985in}{1.832778in}}%
\pgfusepath{clip}%
\pgfsetbuttcap%
\pgfsetmiterjoin%
\definecolor{currentfill}{rgb}{0.161476,0.794233,0.337947}%
\pgfsetfillcolor{currentfill}%
\pgfsetlinewidth{0.000000pt}%
\definecolor{currentstroke}{rgb}{0.000000,0.000000,0.000000}%
\pgfsetstrokecolor{currentstroke}%
\pgfsetstrokeopacity{0.000000}%
\pgfsetdash{}{0pt}%
\pgfpathmoveto{\pgfqpoint{6.298211in}{1.595412in}}%
\pgfpathlineto{\pgfqpoint{6.303538in}{1.600991in}}%
\pgfpathlineto{\pgfqpoint{6.308609in}{1.595747in}}%
\pgfpathlineto{\pgfqpoint{6.299500in}{1.588798in}}%
\pgfpathlineto{\pgfqpoint{6.298211in}{1.595412in}}%
\pgfpathclose%
\pgfusepath{fill}%
\end{pgfscope}%
\begin{pgfscope}%
\pgfpathrectangle{\pgfqpoint{3.525000in}{0.100000in}}{\pgfqpoint{2.862985in}{1.832778in}}%
\pgfusepath{clip}%
\pgfsetbuttcap%
\pgfsetmiterjoin%
\definecolor{currentfill}{rgb}{0.147636,0.702191,0.147636}%
\pgfsetfillcolor{currentfill}%
\pgfsetlinewidth{0.000000pt}%
\definecolor{currentstroke}{rgb}{0.000000,0.000000,0.000000}%
\pgfsetstrokecolor{currentstroke}%
\pgfsetstrokeopacity{0.000000}%
\pgfsetdash{}{0pt}%
\pgfpathmoveto{\pgfqpoint{4.720312in}{1.539783in}}%
\pgfpathlineto{\pgfqpoint{4.724775in}{1.583815in}}%
\pgfpathlineto{\pgfqpoint{4.733284in}{1.669428in}}%
\pgfpathlineto{\pgfqpoint{4.739001in}{1.732781in}}%
\pgfpathlineto{\pgfqpoint{4.785821in}{1.728389in}}%
\pgfpathlineto{\pgfqpoint{4.845718in}{1.723620in}}%
\pgfpathlineto{\pgfqpoint{4.900466in}{1.720093in}}%
\pgfpathlineto{\pgfqpoint{4.950039in}{1.717585in}}%
\pgfpathlineto{\pgfqpoint{5.024019in}{1.715047in}}%
\pgfpathlineto{\pgfqpoint{5.029068in}{1.694197in}}%
\pgfpathlineto{\pgfqpoint{5.027261in}{1.684232in}}%
\pgfpathlineto{\pgfqpoint{5.026666in}{1.663751in}}%
\pgfpathlineto{\pgfqpoint{5.030101in}{1.648413in}}%
\pgfpathlineto{\pgfqpoint{5.037970in}{1.625784in}}%
\pgfpathlineto{\pgfqpoint{5.037952in}{1.597294in}}%
\pgfpathlineto{\pgfqpoint{5.039408in}{1.564197in}}%
\pgfpathlineto{\pgfqpoint{5.041408in}{1.555277in}}%
\pgfpathlineto{\pgfqpoint{5.047257in}{1.545485in}}%
\pgfpathlineto{\pgfqpoint{5.049186in}{1.530226in}}%
\pgfpathlineto{\pgfqpoint{5.048353in}{1.520039in}}%
\pgfpathlineto{\pgfqpoint{4.986333in}{1.521383in}}%
\pgfpathlineto{\pgfqpoint{4.941215in}{1.523670in}}%
\pgfpathlineto{\pgfqpoint{4.875074in}{1.527191in}}%
\pgfpathlineto{\pgfqpoint{4.809849in}{1.531824in}}%
\pgfpathlineto{\pgfqpoint{4.766407in}{1.535329in}}%
\pgfpathlineto{\pgfqpoint{4.720312in}{1.539783in}}%
\pgfpathclose%
\pgfusepath{fill}%
\end{pgfscope}%
\begin{pgfscope}%
\pgfpathrectangle{\pgfqpoint{3.525000in}{0.100000in}}{\pgfqpoint{2.862985in}{1.832778in}}%
\pgfusepath{clip}%
\pgfsetbuttcap%
\pgfsetmiterjoin%
\definecolor{currentfill}{rgb}{0.107651,0.618224,0.107651}%
\pgfsetfillcolor{currentfill}%
\pgfsetlinewidth{0.000000pt}%
\definecolor{currentstroke}{rgb}{0.000000,0.000000,0.000000}%
\pgfsetstrokecolor{currentstroke}%
\pgfsetstrokeopacity{0.000000}%
\pgfsetdash{}{0pt}%
\pgfpathmoveto{\pgfqpoint{4.701565in}{1.355506in}}%
\pgfpathlineto{\pgfqpoint{4.708725in}{1.429317in}}%
\pgfpathlineto{\pgfqpoint{4.713893in}{1.480441in}}%
\pgfpathlineto{\pgfqpoint{4.720312in}{1.539783in}}%
\pgfpathlineto{\pgfqpoint{4.766407in}{1.535329in}}%
\pgfpathlineto{\pgfqpoint{4.809849in}{1.531824in}}%
\pgfpathlineto{\pgfqpoint{4.875074in}{1.527191in}}%
\pgfpathlineto{\pgfqpoint{4.941215in}{1.523670in}}%
\pgfpathlineto{\pgfqpoint{4.986333in}{1.521383in}}%
\pgfpathlineto{\pgfqpoint{5.048353in}{1.520039in}}%
\pgfpathlineto{\pgfqpoint{5.044154in}{1.507786in}}%
\pgfpathlineto{\pgfqpoint{5.035772in}{1.498167in}}%
\pgfpathlineto{\pgfqpoint{5.042191in}{1.487092in}}%
\pgfpathlineto{\pgfqpoint{5.049268in}{1.484725in}}%
\pgfpathlineto{\pgfqpoint{5.052626in}{1.478363in}}%
\pgfpathlineto{\pgfqpoint{5.051855in}{1.431914in}}%
\pgfpathlineto{\pgfqpoint{5.050565in}{1.366544in}}%
\pgfpathlineto{\pgfqpoint{5.045791in}{1.349375in}}%
\pgfpathlineto{\pgfqpoint{5.050061in}{1.339878in}}%
\pgfpathlineto{\pgfqpoint{5.041489in}{1.319470in}}%
\pgfpathlineto{\pgfqpoint{5.050519in}{1.303018in}}%
\pgfpathlineto{\pgfqpoint{5.042847in}{1.304277in}}%
\pgfpathlineto{\pgfqpoint{5.037612in}{1.314513in}}%
\pgfpathlineto{\pgfqpoint{5.018930in}{1.321529in}}%
\pgfpathlineto{\pgfqpoint{5.007148in}{1.327701in}}%
\pgfpathlineto{\pgfqpoint{4.987377in}{1.328214in}}%
\pgfpathlineto{\pgfqpoint{4.980515in}{1.322591in}}%
\pgfpathlineto{\pgfqpoint{4.958137in}{1.333663in}}%
\pgfpathlineto{\pgfqpoint{4.956420in}{1.337164in}}%
\pgfpathlineto{\pgfqpoint{4.878322in}{1.340860in}}%
\pgfpathlineto{\pgfqpoint{4.830874in}{1.343580in}}%
\pgfpathlineto{\pgfqpoint{4.791682in}{1.346649in}}%
\pgfpathlineto{\pgfqpoint{4.726922in}{1.352768in}}%
\pgfpathlineto{\pgfqpoint{4.701565in}{1.355506in}}%
\pgfpathclose%
\pgfusepath{fill}%
\end{pgfscope}%
\begin{pgfscope}%
\pgfpathrectangle{\pgfqpoint{3.525000in}{0.100000in}}{\pgfqpoint{2.862985in}{1.832778in}}%
\pgfusepath{clip}%
\pgfsetbuttcap%
\pgfsetmiterjoin%
\definecolor{currentfill}{rgb}{0.169166,0.747405,0.169166}%
\pgfsetfillcolor{currentfill}%
\pgfsetlinewidth{0.000000pt}%
\definecolor{currentstroke}{rgb}{0.000000,0.000000,0.000000}%
\pgfsetstrokecolor{currentstroke}%
\pgfsetstrokeopacity{0.000000}%
\pgfsetdash{}{0pt}%
\pgfpathmoveto{\pgfqpoint{4.689283in}{1.230524in}}%
\pgfpathlineto{\pgfqpoint{4.647647in}{1.234326in}}%
\pgfpathlineto{\pgfqpoint{4.556892in}{1.245527in}}%
\pgfpathlineto{\pgfqpoint{4.507522in}{1.252613in}}%
\pgfpathlineto{\pgfqpoint{4.454571in}{1.260214in}}%
\pgfpathlineto{\pgfqpoint{4.409967in}{1.267352in}}%
\pgfpathlineto{\pgfqpoint{4.361009in}{1.275734in}}%
\pgfpathlineto{\pgfqpoint{4.372141in}{1.337483in}}%
\pgfpathlineto{\pgfqpoint{4.383416in}{1.400800in}}%
\pgfpathlineto{\pgfqpoint{4.399380in}{1.490854in}}%
\pgfpathlineto{\pgfqpoint{4.405083in}{1.523280in}}%
\pgfpathlineto{\pgfqpoint{4.432445in}{1.518274in}}%
\pgfpathlineto{\pgfqpoint{4.511931in}{1.505797in}}%
\pgfpathlineto{\pgfqpoint{4.544194in}{1.501195in}}%
\pgfpathlineto{\pgfqpoint{4.616429in}{1.491308in}}%
\pgfpathlineto{\pgfqpoint{4.668453in}{1.485403in}}%
\pgfpathlineto{\pgfqpoint{4.713893in}{1.480441in}}%
\pgfpathlineto{\pgfqpoint{4.708725in}{1.429317in}}%
\pgfpathlineto{\pgfqpoint{4.701565in}{1.355506in}}%
\pgfpathlineto{\pgfqpoint{4.695419in}{1.292777in}}%
\pgfpathlineto{\pgfqpoint{4.689283in}{1.230524in}}%
\pgfpathclose%
\pgfusepath{fill}%
\end{pgfscope}%
\begin{pgfscope}%
\pgfpathrectangle{\pgfqpoint{3.525000in}{0.100000in}}{\pgfqpoint{2.862985in}{1.832778in}}%
\pgfusepath{clip}%
\pgfsetbuttcap%
\pgfsetmiterjoin%
\definecolor{currentfill}{rgb}{0.195309,0.803706,0.199231}%
\pgfsetfillcolor{currentfill}%
\pgfsetlinewidth{0.000000pt}%
\definecolor{currentstroke}{rgb}{0.000000,0.000000,0.000000}%
\pgfsetstrokecolor{currentstroke}%
\pgfsetstrokeopacity{0.000000}%
\pgfsetdash{}{0pt}%
\pgfpathmoveto{\pgfqpoint{5.450520in}{1.315700in}}%
\pgfpathlineto{\pgfqpoint{5.397758in}{1.311952in}}%
\pgfpathlineto{\pgfqpoint{5.319114in}{1.308600in}}%
\pgfpathlineto{\pgfqpoint{5.316125in}{1.316544in}}%
\pgfpathlineto{\pgfqpoint{5.298686in}{1.322483in}}%
\pgfpathlineto{\pgfqpoint{5.294839in}{1.333702in}}%
\pgfpathlineto{\pgfqpoint{5.293228in}{1.347582in}}%
\pgfpathlineto{\pgfqpoint{5.297158in}{1.354693in}}%
\pgfpathlineto{\pgfqpoint{5.290951in}{1.361517in}}%
\pgfpathlineto{\pgfqpoint{5.289454in}{1.369657in}}%
\pgfpathlineto{\pgfqpoint{5.287453in}{1.387663in}}%
\pgfpathlineto{\pgfqpoint{5.281505in}{1.397444in}}%
\pgfpathlineto{\pgfqpoint{5.270940in}{1.402933in}}%
\pgfpathlineto{\pgfqpoint{5.259444in}{1.411992in}}%
\pgfpathlineto{\pgfqpoint{5.253501in}{1.422710in}}%
\pgfpathlineto{\pgfqpoint{5.242838in}{1.427026in}}%
\pgfpathlineto{\pgfqpoint{5.236570in}{1.434051in}}%
\pgfpathlineto{\pgfqpoint{5.228976in}{1.435243in}}%
\pgfpathlineto{\pgfqpoint{5.215419in}{1.445673in}}%
\pgfpathlineto{\pgfqpoint{5.217621in}{1.457675in}}%
\pgfpathlineto{\pgfqpoint{5.217197in}{1.480500in}}%
\pgfpathlineto{\pgfqpoint{5.221378in}{1.486783in}}%
\pgfpathlineto{\pgfqpoint{5.217621in}{1.496270in}}%
\pgfpathlineto{\pgfqpoint{5.211004in}{1.498105in}}%
\pgfpathlineto{\pgfqpoint{5.211550in}{1.506437in}}%
\pgfpathlineto{\pgfqpoint{5.219761in}{1.519600in}}%
\pgfpathlineto{\pgfqpoint{5.236020in}{1.530009in}}%
\pgfpathlineto{\pgfqpoint{5.235007in}{1.567025in}}%
\pgfpathlineto{\pgfqpoint{5.243143in}{1.572586in}}%
\pgfpathlineto{\pgfqpoint{5.250842in}{1.568882in}}%
\pgfpathlineto{\pgfqpoint{5.266526in}{1.574304in}}%
\pgfpathlineto{\pgfqpoint{5.296034in}{1.587932in}}%
\pgfpathlineto{\pgfqpoint{5.299878in}{1.583713in}}%
\pgfpathlineto{\pgfqpoint{5.294289in}{1.564600in}}%
\pgfpathlineto{\pgfqpoint{5.302604in}{1.568794in}}%
\pgfpathlineto{\pgfqpoint{5.316845in}{1.564606in}}%
\pgfpathlineto{\pgfqpoint{5.325596in}{1.561110in}}%
\pgfpathlineto{\pgfqpoint{5.330502in}{1.550866in}}%
\pgfpathlineto{\pgfqpoint{5.375397in}{1.541198in}}%
\pgfpathlineto{\pgfqpoint{5.388816in}{1.534565in}}%
\pgfpathlineto{\pgfqpoint{5.402460in}{1.534643in}}%
\pgfpathlineto{\pgfqpoint{5.416480in}{1.531907in}}%
\pgfpathlineto{\pgfqpoint{5.425578in}{1.522561in}}%
\pgfpathlineto{\pgfqpoint{5.434272in}{1.517943in}}%
\pgfpathlineto{\pgfqpoint{5.435870in}{1.504623in}}%
\pgfpathlineto{\pgfqpoint{5.433301in}{1.496259in}}%
\pgfpathlineto{\pgfqpoint{5.442991in}{1.495641in}}%
\pgfpathlineto{\pgfqpoint{5.439727in}{1.485941in}}%
\pgfpathlineto{\pgfqpoint{5.442832in}{1.482495in}}%
\pgfpathlineto{\pgfqpoint{5.445919in}{1.473339in}}%
\pgfpathlineto{\pgfqpoint{5.436476in}{1.468491in}}%
\pgfpathlineto{\pgfqpoint{5.430993in}{1.454984in}}%
\pgfpathlineto{\pgfqpoint{5.429273in}{1.445433in}}%
\pgfpathlineto{\pgfqpoint{5.434533in}{1.443798in}}%
\pgfpathlineto{\pgfqpoint{5.441268in}{1.450960in}}%
\pgfpathlineto{\pgfqpoint{5.446994in}{1.463372in}}%
\pgfpathlineto{\pgfqpoint{5.454741in}{1.467683in}}%
\pgfpathlineto{\pgfqpoint{5.460563in}{1.462075in}}%
\pgfpathlineto{\pgfqpoint{5.454832in}{1.445104in}}%
\pgfpathlineto{\pgfqpoint{5.453034in}{1.431928in}}%
\pgfpathlineto{\pgfqpoint{5.454728in}{1.422457in}}%
\pgfpathlineto{\pgfqpoint{5.449406in}{1.417099in}}%
\pgfpathlineto{\pgfqpoint{5.446748in}{1.404007in}}%
\pgfpathlineto{\pgfqpoint{5.448920in}{1.390247in}}%
\pgfpathlineto{\pgfqpoint{5.442640in}{1.369897in}}%
\pgfpathlineto{\pgfqpoint{5.442772in}{1.359728in}}%
\pgfpathlineto{\pgfqpoint{5.447734in}{1.337693in}}%
\pgfpathlineto{\pgfqpoint{5.450951in}{1.333913in}}%
\pgfpathlineto{\pgfqpoint{5.450520in}{1.315700in}}%
\pgfpathclose%
\pgfusepath{fill}%
\end{pgfscope}%
\begin{pgfscope}%
\pgfpathrectangle{\pgfqpoint{3.525000in}{0.100000in}}{\pgfqpoint{2.862985in}{1.832778in}}%
\pgfusepath{clip}%
\pgfsetbuttcap%
\pgfsetmiterjoin%
\definecolor{currentfill}{rgb}{0.195309,0.803706,0.199231}%
\pgfsetfillcolor{currentfill}%
\pgfsetlinewidth{0.000000pt}%
\definecolor{currentstroke}{rgb}{0.000000,0.000000,0.000000}%
\pgfsetstrokecolor{currentstroke}%
\pgfsetstrokeopacity{0.000000}%
\pgfsetdash{}{0pt}%
\pgfpathmoveto{\pgfqpoint{5.470302in}{1.494291in}}%
\pgfpathlineto{\pgfqpoint{5.471480in}{1.485502in}}%
\pgfpathlineto{\pgfqpoint{5.460690in}{1.462343in}}%
\pgfpathlineto{\pgfqpoint{5.455896in}{1.469051in}}%
\pgfpathlineto{\pgfqpoint{5.470302in}{1.494291in}}%
\pgfpathclose%
\pgfusepath{fill}%
\end{pgfscope}%
\begin{pgfscope}%
\pgfpathrectangle{\pgfqpoint{3.525000in}{0.100000in}}{\pgfqpoint{2.862985in}{1.832778in}}%
\pgfusepath{clip}%
\pgfsetbuttcap%
\pgfsetmiterjoin%
\definecolor{currentfill}{rgb}{0.152249,0.791649,0.375779}%
\pgfsetfillcolor{currentfill}%
\pgfsetlinewidth{0.000000pt}%
\definecolor{currentstroke}{rgb}{0.000000,0.000000,0.000000}%
\pgfsetstrokecolor{currentstroke}%
\pgfsetstrokeopacity{0.000000}%
\pgfsetdash{}{0pt}%
\pgfpathmoveto{\pgfqpoint{4.164329in}{1.638779in}}%
\pgfpathlineto{\pgfqpoint{4.166462in}{1.649338in}}%
\pgfpathlineto{\pgfqpoint{4.164852in}{1.660801in}}%
\pgfpathlineto{\pgfqpoint{4.167317in}{1.673299in}}%
\pgfpathlineto{\pgfqpoint{4.179869in}{1.724007in}}%
\pgfpathlineto{\pgfqpoint{4.190426in}{1.766159in}}%
\pgfpathlineto{\pgfqpoint{4.205273in}{1.824988in}}%
\pgfpathlineto{\pgfqpoint{4.244624in}{1.815480in}}%
\pgfpathlineto{\pgfqpoint{4.230080in}{1.752004in}}%
\pgfpathlineto{\pgfqpoint{4.239291in}{1.731551in}}%
\pgfpathlineto{\pgfqpoint{4.238126in}{1.720446in}}%
\pgfpathlineto{\pgfqpoint{4.243063in}{1.703739in}}%
\pgfpathlineto{\pgfqpoint{4.249957in}{1.699791in}}%
\pgfpathlineto{\pgfqpoint{4.257502in}{1.680433in}}%
\pgfpathlineto{\pgfqpoint{4.262137in}{1.675101in}}%
\pgfpathlineto{\pgfqpoint{4.261575in}{1.668238in}}%
\pgfpathlineto{\pgfqpoint{4.270361in}{1.662098in}}%
\pgfpathlineto{\pgfqpoint{4.271601in}{1.654979in}}%
\pgfpathlineto{\pgfqpoint{4.283663in}{1.653985in}}%
\pgfpathlineto{\pgfqpoint{4.268366in}{1.614312in}}%
\pgfpathlineto{\pgfqpoint{4.270705in}{1.606116in}}%
\pgfpathlineto{\pgfqpoint{4.261905in}{1.601019in}}%
\pgfpathlineto{\pgfqpoint{4.260712in}{1.588082in}}%
\pgfpathlineto{\pgfqpoint{4.268050in}{1.579980in}}%
\pgfpathlineto{\pgfqpoint{4.277690in}{1.583911in}}%
\pgfpathlineto{\pgfqpoint{4.284772in}{1.591178in}}%
\pgfpathlineto{\pgfqpoint{4.292319in}{1.582964in}}%
\pgfpathlineto{\pgfqpoint{4.291728in}{1.566180in}}%
\pgfpathlineto{\pgfqpoint{4.300510in}{1.546436in}}%
\pgfpathlineto{\pgfqpoint{4.298499in}{1.534484in}}%
\pgfpathlineto{\pgfqpoint{4.310833in}{1.526190in}}%
\pgfpathlineto{\pgfqpoint{4.311029in}{1.514658in}}%
\pgfpathlineto{\pgfqpoint{4.315230in}{1.502538in}}%
\pgfpathlineto{\pgfqpoint{4.323038in}{1.505269in}}%
\pgfpathlineto{\pgfqpoint{4.340931in}{1.500183in}}%
\pgfpathlineto{\pgfqpoint{4.348297in}{1.505835in}}%
\pgfpathlineto{\pgfqpoint{4.353566in}{1.502181in}}%
\pgfpathlineto{\pgfqpoint{4.364045in}{1.502638in}}%
\pgfpathlineto{\pgfqpoint{4.366293in}{1.498866in}}%
\pgfpathlineto{\pgfqpoint{4.377480in}{1.499920in}}%
\pgfpathlineto{\pgfqpoint{4.382871in}{1.507806in}}%
\pgfpathlineto{\pgfqpoint{4.390403in}{1.508767in}}%
\pgfpathlineto{\pgfqpoint{4.396262in}{1.493168in}}%
\pgfpathlineto{\pgfqpoint{4.399380in}{1.490854in}}%
\pgfpathlineto{\pgfqpoint{4.383416in}{1.400800in}}%
\pgfpathlineto{\pgfqpoint{4.372141in}{1.337483in}}%
\pgfpathlineto{\pgfqpoint{4.283212in}{1.354650in}}%
\pgfpathlineto{\pgfqpoint{4.235177in}{1.364205in}}%
\pgfpathlineto{\pgfqpoint{4.190246in}{1.374085in}}%
\pgfpathlineto{\pgfqpoint{4.148250in}{1.383593in}}%
\pgfpathlineto{\pgfqpoint{4.099647in}{1.395314in}}%
\pgfpathlineto{\pgfqpoint{4.124717in}{1.498158in}}%
\pgfpathlineto{\pgfqpoint{4.126069in}{1.505674in}}%
\pgfpathlineto{\pgfqpoint{4.137217in}{1.525368in}}%
\pgfpathlineto{\pgfqpoint{4.125673in}{1.537205in}}%
\pgfpathlineto{\pgfqpoint{4.128070in}{1.548828in}}%
\pgfpathlineto{\pgfqpoint{4.132464in}{1.551738in}}%
\pgfpathlineto{\pgfqpoint{4.140313in}{1.563715in}}%
\pgfpathlineto{\pgfqpoint{4.151752in}{1.572002in}}%
\pgfpathlineto{\pgfqpoint{4.152377in}{1.578133in}}%
\pgfpathlineto{\pgfqpoint{4.159282in}{1.584259in}}%
\pgfpathlineto{\pgfqpoint{4.165015in}{1.595725in}}%
\pgfpathlineto{\pgfqpoint{4.176774in}{1.607856in}}%
\pgfpathlineto{\pgfqpoint{4.175939in}{1.619838in}}%
\pgfpathlineto{\pgfqpoint{4.167899in}{1.626494in}}%
\pgfpathlineto{\pgfqpoint{4.164329in}{1.638779in}}%
\pgfpathclose%
\pgfusepath{fill}%
\end{pgfscope}%
\begin{pgfscope}%
\pgfpathrectangle{\pgfqpoint{3.525000in}{0.100000in}}{\pgfqpoint{2.862985in}{1.832778in}}%
\pgfusepath{clip}%
\pgfsetbuttcap%
\pgfsetmiterjoin%
\definecolor{currentfill}{rgb}{0.156863,0.721569,0.156863}%
\pgfsetfillcolor{currentfill}%
\pgfsetlinewidth{0.000000pt}%
\definecolor{currentstroke}{rgb}{0.000000,0.000000,0.000000}%
\pgfsetstrokecolor{currentstroke}%
\pgfsetstrokeopacity{0.000000}%
\pgfsetdash{}{0pt}%
\pgfpathmoveto{\pgfqpoint{6.111015in}{1.438080in}}%
\pgfpathlineto{\pgfqpoint{6.107320in}{1.449744in}}%
\pgfpathlineto{\pgfqpoint{6.100463in}{1.485127in}}%
\pgfpathlineto{\pgfqpoint{6.091598in}{1.496013in}}%
\pgfpathlineto{\pgfqpoint{6.083822in}{1.515590in}}%
\pgfpathlineto{\pgfqpoint{6.086375in}{1.530104in}}%
\pgfpathlineto{\pgfqpoint{6.084327in}{1.540801in}}%
\pgfpathlineto{\pgfqpoint{6.077683in}{1.551308in}}%
\pgfpathlineto{\pgfqpoint{6.077410in}{1.562882in}}%
\pgfpathlineto{\pgfqpoint{6.073552in}{1.575365in}}%
\pgfpathlineto{\pgfqpoint{6.108077in}{1.583862in}}%
\pgfpathlineto{\pgfqpoint{6.152922in}{1.595966in}}%
\pgfpathlineto{\pgfqpoint{6.154711in}{1.589025in}}%
\pgfpathlineto{\pgfqpoint{6.151856in}{1.577975in}}%
\pgfpathlineto{\pgfqpoint{6.158559in}{1.569133in}}%
\pgfpathlineto{\pgfqpoint{6.154999in}{1.557930in}}%
\pgfpathlineto{\pgfqpoint{6.140852in}{1.543846in}}%
\pgfpathlineto{\pgfqpoint{6.144741in}{1.533267in}}%
\pgfpathlineto{\pgfqpoint{6.141859in}{1.510906in}}%
\pgfpathlineto{\pgfqpoint{6.137824in}{1.498298in}}%
\pgfpathlineto{\pgfqpoint{6.142972in}{1.462444in}}%
\pgfpathlineto{\pgfqpoint{6.142278in}{1.449830in}}%
\pgfpathlineto{\pgfqpoint{6.147258in}{1.445782in}}%
\pgfpathlineto{\pgfqpoint{6.111015in}{1.438080in}}%
\pgfpathclose%
\pgfusepath{fill}%
\end{pgfscope}%
\begin{pgfscope}%
\pgfpathrectangle{\pgfqpoint{3.525000in}{0.100000in}}{\pgfqpoint{2.862985in}{1.832778in}}%
\pgfusepath{clip}%
\pgfsetbuttcap%
\pgfsetmiterjoin%
\definecolor{currentfill}{rgb}{0.132257,0.669896,0.132257}%
\pgfsetfillcolor{currentfill}%
\pgfsetlinewidth{0.000000pt}%
\definecolor{currentstroke}{rgb}{0.000000,0.000000,0.000000}%
\pgfsetstrokecolor{currentstroke}%
\pgfsetstrokeopacity{0.000000}%
\pgfsetdash{}{0pt}%
\pgfpathmoveto{\pgfqpoint{5.050565in}{1.366544in}}%
\pgfpathlineto{\pgfqpoint{5.051855in}{1.431914in}}%
\pgfpathlineto{\pgfqpoint{5.052626in}{1.478363in}}%
\pgfpathlineto{\pgfqpoint{5.049268in}{1.484725in}}%
\pgfpathlineto{\pgfqpoint{5.042191in}{1.487092in}}%
\pgfpathlineto{\pgfqpoint{5.035772in}{1.498167in}}%
\pgfpathlineto{\pgfqpoint{5.044154in}{1.507786in}}%
\pgfpathlineto{\pgfqpoint{5.048353in}{1.520039in}}%
\pgfpathlineto{\pgfqpoint{5.049186in}{1.530226in}}%
\pgfpathlineto{\pgfqpoint{5.047257in}{1.545485in}}%
\pgfpathlineto{\pgfqpoint{5.041408in}{1.555277in}}%
\pgfpathlineto{\pgfqpoint{5.039408in}{1.564197in}}%
\pgfpathlineto{\pgfqpoint{5.037952in}{1.597294in}}%
\pgfpathlineto{\pgfqpoint{5.037970in}{1.625784in}}%
\pgfpathlineto{\pgfqpoint{5.030101in}{1.648413in}}%
\pgfpathlineto{\pgfqpoint{5.026666in}{1.663751in}}%
\pgfpathlineto{\pgfqpoint{5.027261in}{1.684232in}}%
\pgfpathlineto{\pgfqpoint{5.029068in}{1.694197in}}%
\pgfpathlineto{\pgfqpoint{5.024019in}{1.715047in}}%
\pgfpathlineto{\pgfqpoint{5.058393in}{1.714359in}}%
\pgfpathlineto{\pgfqpoint{5.110604in}{1.713911in}}%
\pgfpathlineto{\pgfqpoint{5.110890in}{1.737609in}}%
\pgfpathlineto{\pgfqpoint{5.124180in}{1.735000in}}%
\pgfpathlineto{\pgfqpoint{5.130549in}{1.706104in}}%
\pgfpathlineto{\pgfqpoint{5.135245in}{1.695714in}}%
\pgfpathlineto{\pgfqpoint{5.146922in}{1.695407in}}%
\pgfpathlineto{\pgfqpoint{5.149536in}{1.691882in}}%
\pgfpathlineto{\pgfqpoint{5.165825in}{1.690320in}}%
\pgfpathlineto{\pgfqpoint{5.168562in}{1.683156in}}%
\pgfpathlineto{\pgfqpoint{5.179788in}{1.684767in}}%
\pgfpathlineto{\pgfqpoint{5.188507in}{1.691467in}}%
\pgfpathlineto{\pgfqpoint{5.203542in}{1.691216in}}%
\pgfpathlineto{\pgfqpoint{5.212845in}{1.685828in}}%
\pgfpathlineto{\pgfqpoint{5.213915in}{1.680774in}}%
\pgfpathlineto{\pgfqpoint{5.222767in}{1.679717in}}%
\pgfpathlineto{\pgfqpoint{5.228544in}{1.665932in}}%
\pgfpathlineto{\pgfqpoint{5.232272in}{1.674408in}}%
\pgfpathlineto{\pgfqpoint{5.242441in}{1.674929in}}%
\pgfpathlineto{\pgfqpoint{5.245013in}{1.668328in}}%
\pgfpathlineto{\pgfqpoint{5.256455in}{1.665315in}}%
\pgfpathlineto{\pgfqpoint{5.262765in}{1.655808in}}%
\pgfpathlineto{\pgfqpoint{5.276759in}{1.658760in}}%
\pgfpathlineto{\pgfqpoint{5.292158in}{1.670427in}}%
\pgfpathlineto{\pgfqpoint{5.297772in}{1.660142in}}%
\pgfpathlineto{\pgfqpoint{5.323039in}{1.662957in}}%
\pgfpathlineto{\pgfqpoint{5.333840in}{1.655898in}}%
\pgfpathlineto{\pgfqpoint{5.340126in}{1.658418in}}%
\pgfpathlineto{\pgfqpoint{5.345167in}{1.654444in}}%
\pgfpathlineto{\pgfqpoint{5.330213in}{1.645016in}}%
\pgfpathlineto{\pgfqpoint{5.308866in}{1.636615in}}%
\pgfpathlineto{\pgfqpoint{5.287725in}{1.619824in}}%
\pgfpathlineto{\pgfqpoint{5.269396in}{1.597742in}}%
\pgfpathlineto{\pgfqpoint{5.255521in}{1.584687in}}%
\pgfpathlineto{\pgfqpoint{5.243367in}{1.575713in}}%
\pgfpathlineto{\pgfqpoint{5.235007in}{1.567025in}}%
\pgfpathlineto{\pgfqpoint{5.236020in}{1.530009in}}%
\pgfpathlineto{\pgfqpoint{5.219761in}{1.519600in}}%
\pgfpathlineto{\pgfqpoint{5.211550in}{1.506437in}}%
\pgfpathlineto{\pgfqpoint{5.211004in}{1.498105in}}%
\pgfpathlineto{\pgfqpoint{5.217621in}{1.496270in}}%
\pgfpathlineto{\pgfqpoint{5.221378in}{1.486783in}}%
\pgfpathlineto{\pgfqpoint{5.217197in}{1.480500in}}%
\pgfpathlineto{\pgfqpoint{5.217621in}{1.457675in}}%
\pgfpathlineto{\pgfqpoint{5.215419in}{1.445673in}}%
\pgfpathlineto{\pgfqpoint{5.228976in}{1.435243in}}%
\pgfpathlineto{\pgfqpoint{5.236570in}{1.434051in}}%
\pgfpathlineto{\pgfqpoint{5.242838in}{1.427026in}}%
\pgfpathlineto{\pgfqpoint{5.253501in}{1.422710in}}%
\pgfpathlineto{\pgfqpoint{5.259444in}{1.411992in}}%
\pgfpathlineto{\pgfqpoint{5.270940in}{1.402933in}}%
\pgfpathlineto{\pgfqpoint{5.281505in}{1.397444in}}%
\pgfpathlineto{\pgfqpoint{5.287453in}{1.387663in}}%
\pgfpathlineto{\pgfqpoint{5.289454in}{1.369657in}}%
\pgfpathlineto{\pgfqpoint{5.233388in}{1.367621in}}%
\pgfpathlineto{\pgfqpoint{5.185597in}{1.366622in}}%
\pgfpathlineto{\pgfqpoint{5.142055in}{1.365982in}}%
\pgfpathlineto{\pgfqpoint{5.095992in}{1.366052in}}%
\pgfpathlineto{\pgfqpoint{5.050565in}{1.366544in}}%
\pgfpathclose%
\pgfusepath{fill}%
\end{pgfscope}%
\begin{pgfscope}%
\pgfpathrectangle{\pgfqpoint{3.525000in}{0.100000in}}{\pgfqpoint{2.862985in}{1.832778in}}%
\pgfusepath{clip}%
\pgfsetbuttcap%
\pgfsetmiterjoin%
\definecolor{currentfill}{rgb}{0.106113,0.778731,0.564937}%
\pgfsetfillcolor{currentfill}%
\pgfsetlinewidth{0.000000pt}%
\definecolor{currentstroke}{rgb}{0.000000,0.000000,0.000000}%
\pgfsetstrokecolor{currentstroke}%
\pgfsetstrokeopacity{0.000000}%
\pgfsetdash{}{0pt}%
\pgfpathmoveto{\pgfqpoint{3.777592in}{1.488751in}}%
\pgfpathlineto{\pgfqpoint{3.773152in}{1.496921in}}%
\pgfpathlineto{\pgfqpoint{3.775987in}{1.517893in}}%
\pgfpathlineto{\pgfqpoint{3.781459in}{1.528853in}}%
\pgfpathlineto{\pgfqpoint{3.778724in}{1.543679in}}%
\pgfpathlineto{\pgfqpoint{3.784432in}{1.549914in}}%
\pgfpathlineto{\pgfqpoint{3.794847in}{1.566720in}}%
\pgfpathlineto{\pgfqpoint{3.803677in}{1.576864in}}%
\pgfpathlineto{\pgfqpoint{3.816585in}{1.598959in}}%
\pgfpathlineto{\pgfqpoint{3.826485in}{1.623028in}}%
\pgfpathlineto{\pgfqpoint{3.836995in}{1.645601in}}%
\pgfpathlineto{\pgfqpoint{3.839119in}{1.655029in}}%
\pgfpathlineto{\pgfqpoint{3.853617in}{1.682000in}}%
\pgfpathlineto{\pgfqpoint{3.856409in}{1.693845in}}%
\pgfpathlineto{\pgfqpoint{3.862582in}{1.706284in}}%
\pgfpathlineto{\pgfqpoint{3.864788in}{1.721458in}}%
\pgfpathlineto{\pgfqpoint{3.876573in}{1.728857in}}%
\pgfpathlineto{\pgfqpoint{3.890546in}{1.732587in}}%
\pgfpathlineto{\pgfqpoint{3.897610in}{1.724234in}}%
\pgfpathlineto{\pgfqpoint{3.903718in}{1.724852in}}%
\pgfpathlineto{\pgfqpoint{3.913254in}{1.714897in}}%
\pgfpathlineto{\pgfqpoint{3.914522in}{1.705519in}}%
\pgfpathlineto{\pgfqpoint{3.910905in}{1.687156in}}%
\pgfpathlineto{\pgfqpoint{3.922937in}{1.677787in}}%
\pgfpathlineto{\pgfqpoint{3.930740in}{1.674264in}}%
\pgfpathlineto{\pgfqpoint{3.951880in}{1.677860in}}%
\pgfpathlineto{\pgfqpoint{3.964127in}{1.675404in}}%
\pgfpathlineto{\pgfqpoint{3.972545in}{1.671374in}}%
\pgfpathlineto{\pgfqpoint{3.976979in}{1.663725in}}%
\pgfpathlineto{\pgfqpoint{4.005055in}{1.664746in}}%
\pgfpathlineto{\pgfqpoint{4.009520in}{1.660065in}}%
\pgfpathlineto{\pgfqpoint{4.020053in}{1.659055in}}%
\pgfpathlineto{\pgfqpoint{4.030660in}{1.662045in}}%
\pgfpathlineto{\pgfqpoint{4.062021in}{1.661292in}}%
\pgfpathlineto{\pgfqpoint{4.068794in}{1.658987in}}%
\pgfpathlineto{\pgfqpoint{4.076709in}{1.661643in}}%
\pgfpathlineto{\pgfqpoint{4.164329in}{1.638779in}}%
\pgfpathlineto{\pgfqpoint{4.167899in}{1.626494in}}%
\pgfpathlineto{\pgfqpoint{4.175939in}{1.619838in}}%
\pgfpathlineto{\pgfqpoint{4.176774in}{1.607856in}}%
\pgfpathlineto{\pgfqpoint{4.165015in}{1.595725in}}%
\pgfpathlineto{\pgfqpoint{4.159282in}{1.584259in}}%
\pgfpathlineto{\pgfqpoint{4.152377in}{1.578133in}}%
\pgfpathlineto{\pgfqpoint{4.151752in}{1.572002in}}%
\pgfpathlineto{\pgfqpoint{4.140313in}{1.563715in}}%
\pgfpathlineto{\pgfqpoint{4.132464in}{1.551738in}}%
\pgfpathlineto{\pgfqpoint{4.128070in}{1.548828in}}%
\pgfpathlineto{\pgfqpoint{4.125673in}{1.537205in}}%
\pgfpathlineto{\pgfqpoint{4.137217in}{1.525368in}}%
\pgfpathlineto{\pgfqpoint{4.126069in}{1.505674in}}%
\pgfpathlineto{\pgfqpoint{4.124717in}{1.498158in}}%
\pgfpathlineto{\pgfqpoint{4.099647in}{1.395314in}}%
\pgfpathlineto{\pgfqpoint{4.046916in}{1.408826in}}%
\pgfpathlineto{\pgfqpoint{3.996046in}{1.421943in}}%
\pgfpathlineto{\pgfqpoint{3.965358in}{1.430483in}}%
\pgfpathlineto{\pgfqpoint{3.925927in}{1.441715in}}%
\pgfpathlineto{\pgfqpoint{3.863032in}{1.461511in}}%
\pgfpathlineto{\pgfqpoint{3.794661in}{1.482798in}}%
\pgfpathlineto{\pgfqpoint{3.777592in}{1.488751in}}%
\pgfpathclose%
\pgfusepath{fill}%
\end{pgfscope}%
\begin{pgfscope}%
\pgfpathrectangle{\pgfqpoint{3.525000in}{0.100000in}}{\pgfqpoint{2.862985in}{1.832778in}}%
\pgfusepath{clip}%
\pgfsetbuttcap%
\pgfsetmiterjoin%
\definecolor{currentfill}{rgb}{0.123030,0.650519,0.123030}%
\pgfsetfillcolor{currentfill}%
\pgfsetlinewidth{0.000000pt}%
\definecolor{currentstroke}{rgb}{0.000000,0.000000,0.000000}%
\pgfsetstrokecolor{currentstroke}%
\pgfsetstrokeopacity{0.000000}%
\pgfsetdash{}{0pt}%
\pgfpathmoveto{\pgfqpoint{6.147258in}{1.445782in}}%
\pgfpathlineto{\pgfqpoint{6.142278in}{1.449830in}}%
\pgfpathlineto{\pgfqpoint{6.142972in}{1.462444in}}%
\pgfpathlineto{\pgfqpoint{6.137824in}{1.498298in}}%
\pgfpathlineto{\pgfqpoint{6.141859in}{1.510906in}}%
\pgfpathlineto{\pgfqpoint{6.144741in}{1.533267in}}%
\pgfpathlineto{\pgfqpoint{6.140852in}{1.543846in}}%
\pgfpathlineto{\pgfqpoint{6.154999in}{1.557930in}}%
\pgfpathlineto{\pgfqpoint{6.158559in}{1.569133in}}%
\pgfpathlineto{\pgfqpoint{6.151856in}{1.577975in}}%
\pgfpathlineto{\pgfqpoint{6.154711in}{1.589025in}}%
\pgfpathlineto{\pgfqpoint{6.152922in}{1.595966in}}%
\pgfpathlineto{\pgfqpoint{6.154459in}{1.610832in}}%
\pgfpathlineto{\pgfqpoint{6.157332in}{1.615421in}}%
\pgfpathlineto{\pgfqpoint{6.166192in}{1.618318in}}%
\pgfpathlineto{\pgfqpoint{6.179117in}{1.580468in}}%
\pgfpathlineto{\pgfqpoint{6.203094in}{1.502035in}}%
\pgfpathlineto{\pgfqpoint{6.212168in}{1.495954in}}%
\pgfpathlineto{\pgfqpoint{6.213801in}{1.489012in}}%
\pgfpathlineto{\pgfqpoint{6.218590in}{1.486207in}}%
\pgfpathlineto{\pgfqpoint{6.218226in}{1.473617in}}%
\pgfpathlineto{\pgfqpoint{6.213151in}{1.473415in}}%
\pgfpathlineto{\pgfqpoint{6.202871in}{1.465574in}}%
\pgfpathlineto{\pgfqpoint{6.199906in}{1.457709in}}%
\pgfpathlineto{\pgfqpoint{6.172390in}{1.450901in}}%
\pgfpathlineto{\pgfqpoint{6.147258in}{1.445782in}}%
\pgfpathclose%
\pgfusepath{fill}%
\end{pgfscope}%
\begin{pgfscope}%
\pgfpathrectangle{\pgfqpoint{3.525000in}{0.100000in}}{\pgfqpoint{2.862985in}{1.832778in}}%
\pgfusepath{clip}%
\pgfsetbuttcap%
\pgfsetmiterjoin%
\definecolor{currentfill}{rgb}{0.126105,0.656978,0.126105}%
\pgfsetfillcolor{currentfill}%
\pgfsetlinewidth{0.000000pt}%
\definecolor{currentstroke}{rgb}{0.000000,0.000000,0.000000}%
\pgfsetstrokecolor{currentstroke}%
\pgfsetstrokeopacity{0.000000}%
\pgfsetdash{}{0pt}%
\pgfpathmoveto{\pgfqpoint{5.286865in}{1.172908in}}%
\pgfpathlineto{\pgfqpoint{5.272339in}{1.187293in}}%
\pgfpathlineto{\pgfqpoint{5.225906in}{1.184615in}}%
\pgfpathlineto{\pgfqpoint{5.153489in}{1.182229in}}%
\pgfpathlineto{\pgfqpoint{5.080637in}{1.183365in}}%
\pgfpathlineto{\pgfqpoint{5.075523in}{1.192282in}}%
\pgfpathlineto{\pgfqpoint{5.077611in}{1.201038in}}%
\pgfpathlineto{\pgfqpoint{5.076620in}{1.216037in}}%
\pgfpathlineto{\pgfqpoint{5.072770in}{1.230565in}}%
\pgfpathlineto{\pgfqpoint{5.073217in}{1.238237in}}%
\pgfpathlineto{\pgfqpoint{5.065033in}{1.246905in}}%
\pgfpathlineto{\pgfqpoint{5.066814in}{1.258967in}}%
\pgfpathlineto{\pgfqpoint{5.054260in}{1.282791in}}%
\pgfpathlineto{\pgfqpoint{5.050519in}{1.303018in}}%
\pgfpathlineto{\pgfqpoint{5.041489in}{1.319470in}}%
\pgfpathlineto{\pgfqpoint{5.050061in}{1.339878in}}%
\pgfpathlineto{\pgfqpoint{5.045791in}{1.349375in}}%
\pgfpathlineto{\pgfqpoint{5.050565in}{1.366544in}}%
\pgfpathlineto{\pgfqpoint{5.095992in}{1.366052in}}%
\pgfpathlineto{\pgfqpoint{5.142055in}{1.365982in}}%
\pgfpathlineto{\pgfqpoint{5.185597in}{1.366622in}}%
\pgfpathlineto{\pgfqpoint{5.233388in}{1.367621in}}%
\pgfpathlineto{\pgfqpoint{5.289454in}{1.369657in}}%
\pgfpathlineto{\pgfqpoint{5.290951in}{1.361517in}}%
\pgfpathlineto{\pgfqpoint{5.297158in}{1.354693in}}%
\pgfpathlineto{\pgfqpoint{5.293228in}{1.347582in}}%
\pgfpathlineto{\pgfqpoint{5.294839in}{1.333702in}}%
\pgfpathlineto{\pgfqpoint{5.298686in}{1.322483in}}%
\pgfpathlineto{\pgfqpoint{5.316125in}{1.316544in}}%
\pgfpathlineto{\pgfqpoint{5.319114in}{1.308600in}}%
\pgfpathlineto{\pgfqpoint{5.328684in}{1.299681in}}%
\pgfpathlineto{\pgfqpoint{5.332588in}{1.290454in}}%
\pgfpathlineto{\pgfqpoint{5.342285in}{1.284265in}}%
\pgfpathlineto{\pgfqpoint{5.343802in}{1.276813in}}%
\pgfpathlineto{\pgfqpoint{5.341914in}{1.265532in}}%
\pgfpathlineto{\pgfqpoint{5.336979in}{1.262151in}}%
\pgfpathlineto{\pgfqpoint{5.335490in}{1.251413in}}%
\pgfpathlineto{\pgfqpoint{5.321306in}{1.242884in}}%
\pgfpathlineto{\pgfqpoint{5.302818in}{1.238212in}}%
\pgfpathlineto{\pgfqpoint{5.301124in}{1.227469in}}%
\pgfpathlineto{\pgfqpoint{5.308256in}{1.219794in}}%
\pgfpathlineto{\pgfqpoint{5.308547in}{1.210143in}}%
\pgfpathlineto{\pgfqpoint{5.302793in}{1.202558in}}%
\pgfpathlineto{\pgfqpoint{5.299773in}{1.191288in}}%
\pgfpathlineto{\pgfqpoint{5.289771in}{1.187558in}}%
\pgfpathlineto{\pgfqpoint{5.290408in}{1.174999in}}%
\pgfpathlineto{\pgfqpoint{5.286865in}{1.172908in}}%
\pgfpathclose%
\pgfusepath{fill}%
\end{pgfscope}%
\begin{pgfscope}%
\pgfpathrectangle{\pgfqpoint{3.525000in}{0.100000in}}{\pgfqpoint{2.862985in}{1.832778in}}%
\pgfusepath{clip}%
\pgfsetbuttcap%
\pgfsetmiterjoin%
\definecolor{currentfill}{rgb}{0.173779,0.797678,0.287505}%
\pgfsetfillcolor{currentfill}%
\pgfsetlinewidth{0.000000pt}%
\definecolor{currentstroke}{rgb}{0.000000,0.000000,0.000000}%
\pgfsetstrokecolor{currentstroke}%
\pgfsetstrokeopacity{0.000000}%
\pgfsetdash{}{0pt}%
\pgfpathmoveto{\pgfqpoint{6.188120in}{1.409488in}}%
\pgfpathlineto{\pgfqpoint{6.174362in}{1.407318in}}%
\pgfpathlineto{\pgfqpoint{6.111166in}{1.392957in}}%
\pgfpathlineto{\pgfqpoint{6.110061in}{1.394631in}}%
\pgfpathlineto{\pgfqpoint{6.111015in}{1.438080in}}%
\pgfpathlineto{\pgfqpoint{6.147258in}{1.445782in}}%
\pgfpathlineto{\pgfqpoint{6.172390in}{1.450901in}}%
\pgfpathlineto{\pgfqpoint{6.199906in}{1.457709in}}%
\pgfpathlineto{\pgfqpoint{6.202871in}{1.465574in}}%
\pgfpathlineto{\pgfqpoint{6.213151in}{1.473415in}}%
\pgfpathlineto{\pgfqpoint{6.218226in}{1.473617in}}%
\pgfpathlineto{\pgfqpoint{6.224896in}{1.462179in}}%
\pgfpathlineto{\pgfqpoint{6.218843in}{1.445482in}}%
\pgfpathlineto{\pgfqpoint{6.217955in}{1.435691in}}%
\pgfpathlineto{\pgfqpoint{6.230207in}{1.436611in}}%
\pgfpathlineto{\pgfqpoint{6.235763in}{1.431914in}}%
\pgfpathlineto{\pgfqpoint{6.248224in}{1.412717in}}%
\pgfpathlineto{\pgfqpoint{6.254417in}{1.410382in}}%
\pgfpathlineto{\pgfqpoint{6.264790in}{1.411248in}}%
\pgfpathlineto{\pgfqpoint{6.272009in}{1.417772in}}%
\pgfpathlineto{\pgfqpoint{6.276808in}{1.411950in}}%
\pgfpathlineto{\pgfqpoint{6.246631in}{1.396025in}}%
\pgfpathlineto{\pgfqpoint{6.245688in}{1.407452in}}%
\pgfpathlineto{\pgfqpoint{6.231980in}{1.389693in}}%
\pgfpathlineto{\pgfqpoint{6.227160in}{1.386633in}}%
\pgfpathlineto{\pgfqpoint{6.220438in}{1.396881in}}%
\pgfpathlineto{\pgfqpoint{6.218598in}{1.398285in}}%
\pgfpathlineto{\pgfqpoint{6.212340in}{1.401600in}}%
\pgfpathlineto{\pgfqpoint{6.206882in}{1.415045in}}%
\pgfpathlineto{\pgfqpoint{6.188120in}{1.409488in}}%
\pgfpathclose%
\pgfusepath{fill}%
\end{pgfscope}%
\begin{pgfscope}%
\pgfpathrectangle{\pgfqpoint{3.525000in}{0.100000in}}{\pgfqpoint{2.862985in}{1.832778in}}%
\pgfusepath{clip}%
\pgfsetbuttcap%
\pgfsetmiterjoin%
\definecolor{currentfill}{rgb}{0.067666,0.534256,0.067666}%
\pgfsetfillcolor{currentfill}%
\pgfsetlinewidth{0.000000pt}%
\definecolor{currentstroke}{rgb}{0.000000,0.000000,0.000000}%
\pgfsetstrokecolor{currentstroke}%
\pgfsetstrokeopacity{0.000000}%
\pgfsetdash{}{0pt}%
\pgfpathmoveto{\pgfqpoint{4.778739in}{1.159333in}}%
\pgfpathlineto{\pgfqpoint{4.783777in}{1.221808in}}%
\pgfpathlineto{\pgfqpoint{4.755243in}{1.224124in}}%
\pgfpathlineto{\pgfqpoint{4.689283in}{1.230524in}}%
\pgfpathlineto{\pgfqpoint{4.695419in}{1.292777in}}%
\pgfpathlineto{\pgfqpoint{4.701565in}{1.355506in}}%
\pgfpathlineto{\pgfqpoint{4.726922in}{1.352768in}}%
\pgfpathlineto{\pgfqpoint{4.791682in}{1.346649in}}%
\pgfpathlineto{\pgfqpoint{4.830874in}{1.343580in}}%
\pgfpathlineto{\pgfqpoint{4.878322in}{1.340860in}}%
\pgfpathlineto{\pgfqpoint{4.956420in}{1.337164in}}%
\pgfpathlineto{\pgfqpoint{4.958137in}{1.333663in}}%
\pgfpathlineto{\pgfqpoint{4.980515in}{1.322591in}}%
\pgfpathlineto{\pgfqpoint{4.987377in}{1.328214in}}%
\pgfpathlineto{\pgfqpoint{5.007148in}{1.327701in}}%
\pgfpathlineto{\pgfqpoint{5.018930in}{1.321529in}}%
\pgfpathlineto{\pgfqpoint{5.037612in}{1.314513in}}%
\pgfpathlineto{\pgfqpoint{5.042847in}{1.304277in}}%
\pgfpathlineto{\pgfqpoint{5.050519in}{1.303018in}}%
\pgfpathlineto{\pgfqpoint{5.054260in}{1.282791in}}%
\pgfpathlineto{\pgfqpoint{5.066814in}{1.258967in}}%
\pgfpathlineto{\pgfqpoint{5.065033in}{1.246905in}}%
\pgfpathlineto{\pgfqpoint{5.073217in}{1.238237in}}%
\pgfpathlineto{\pgfqpoint{5.072770in}{1.230565in}}%
\pgfpathlineto{\pgfqpoint{5.076620in}{1.216037in}}%
\pgfpathlineto{\pgfqpoint{5.077611in}{1.201038in}}%
\pgfpathlineto{\pgfqpoint{5.075523in}{1.192282in}}%
\pgfpathlineto{\pgfqpoint{5.080637in}{1.183365in}}%
\pgfpathlineto{\pgfqpoint{5.087651in}{1.167151in}}%
\pgfpathlineto{\pgfqpoint{5.094364in}{1.160551in}}%
\pgfpathlineto{\pgfqpoint{5.102367in}{1.146251in}}%
\pgfpathlineto{\pgfqpoint{5.079686in}{1.146016in}}%
\pgfpathlineto{\pgfqpoint{5.030647in}{1.146774in}}%
\pgfpathlineto{\pgfqpoint{4.976466in}{1.148434in}}%
\pgfpathlineto{\pgfqpoint{4.921966in}{1.150526in}}%
\pgfpathlineto{\pgfqpoint{4.841835in}{1.154903in}}%
\pgfpathlineto{\pgfqpoint{4.778739in}{1.159333in}}%
\pgfpathclose%
\pgfusepath{fill}%
\end{pgfscope}%
\begin{pgfscope}%
\pgfpathrectangle{\pgfqpoint{3.525000in}{0.100000in}}{\pgfqpoint{2.862985in}{1.832778in}}%
\pgfusepath{clip}%
\pgfsetbuttcap%
\pgfsetmiterjoin%
\definecolor{currentfill}{rgb}{0.059977,0.765813,0.754095}%
\pgfsetfillcolor{currentfill}%
\pgfsetlinewidth{0.000000pt}%
\definecolor{currentstroke}{rgb}{0.000000,0.000000,0.000000}%
\pgfsetstrokecolor{currentstroke}%
\pgfsetstrokeopacity{0.000000}%
\pgfsetdash{}{0pt}%
\pgfpathmoveto{\pgfqpoint{5.822084in}{1.347781in}}%
\pgfpathlineto{\pgfqpoint{5.847175in}{1.371705in}}%
\pgfpathlineto{\pgfqpoint{5.850294in}{1.380209in}}%
\pgfpathlineto{\pgfqpoint{5.857646in}{1.387490in}}%
\pgfpathlineto{\pgfqpoint{5.852129in}{1.398090in}}%
\pgfpathlineto{\pgfqpoint{5.845212in}{1.404290in}}%
\pgfpathlineto{\pgfqpoint{5.843213in}{1.415277in}}%
\pgfpathlineto{\pgfqpoint{5.868959in}{1.426558in}}%
\pgfpathlineto{\pgfqpoint{5.890274in}{1.430126in}}%
\pgfpathlineto{\pgfqpoint{5.901728in}{1.430365in}}%
\pgfpathlineto{\pgfqpoint{5.910453in}{1.426054in}}%
\pgfpathlineto{\pgfqpoint{5.918958in}{1.429903in}}%
\pgfpathlineto{\pgfqpoint{5.939702in}{1.434217in}}%
\pgfpathlineto{\pgfqpoint{5.946870in}{1.439785in}}%
\pgfpathlineto{\pgfqpoint{5.957498in}{1.452111in}}%
\pgfpathlineto{\pgfqpoint{5.967167in}{1.457558in}}%
\pgfpathlineto{\pgfqpoint{5.967849in}{1.462782in}}%
\pgfpathlineto{\pgfqpoint{5.962774in}{1.474701in}}%
\pgfpathlineto{\pgfqpoint{5.966443in}{1.481714in}}%
\pgfpathlineto{\pgfqpoint{5.961503in}{1.489256in}}%
\pgfpathlineto{\pgfqpoint{5.953939in}{1.489782in}}%
\pgfpathlineto{\pgfqpoint{5.972870in}{1.512563in}}%
\pgfpathlineto{\pgfqpoint{5.975119in}{1.521245in}}%
\pgfpathlineto{\pgfqpoint{5.990018in}{1.543386in}}%
\pgfpathlineto{\pgfqpoint{6.003850in}{1.555368in}}%
\pgfpathlineto{\pgfqpoint{6.013347in}{1.560371in}}%
\pgfpathlineto{\pgfqpoint{6.044425in}{1.567413in}}%
\pgfpathlineto{\pgfqpoint{6.073552in}{1.575365in}}%
\pgfpathlineto{\pgfqpoint{6.077410in}{1.562882in}}%
\pgfpathlineto{\pgfqpoint{6.077683in}{1.551308in}}%
\pgfpathlineto{\pgfqpoint{6.084327in}{1.540801in}}%
\pgfpathlineto{\pgfqpoint{6.086375in}{1.530104in}}%
\pgfpathlineto{\pgfqpoint{6.083822in}{1.515590in}}%
\pgfpathlineto{\pgfqpoint{6.091598in}{1.496013in}}%
\pgfpathlineto{\pgfqpoint{6.100463in}{1.485127in}}%
\pgfpathlineto{\pgfqpoint{6.107320in}{1.449744in}}%
\pgfpathlineto{\pgfqpoint{6.111015in}{1.438080in}}%
\pgfpathlineto{\pgfqpoint{6.110061in}{1.394631in}}%
\pgfpathlineto{\pgfqpoint{6.111166in}{1.392957in}}%
\pgfpathlineto{\pgfqpoint{6.119238in}{1.346227in}}%
\pgfpathlineto{\pgfqpoint{6.123765in}{1.341969in}}%
\pgfpathlineto{\pgfqpoint{6.114028in}{1.332508in}}%
\pgfpathlineto{\pgfqpoint{6.118832in}{1.327078in}}%
\pgfpathlineto{\pgfqpoint{6.114609in}{1.318865in}}%
\pgfpathlineto{\pgfqpoint{6.114645in}{1.315368in}}%
\pgfpathlineto{\pgfqpoint{6.109366in}{1.312213in}}%
\pgfpathlineto{\pgfqpoint{6.106795in}{1.305239in}}%
\pgfpathlineto{\pgfqpoint{6.107610in}{1.324417in}}%
\pgfpathlineto{\pgfqpoint{6.091246in}{1.328639in}}%
\pgfpathlineto{\pgfqpoint{6.065655in}{1.337368in}}%
\pgfpathlineto{\pgfqpoint{6.062035in}{1.341668in}}%
\pgfpathlineto{\pgfqpoint{6.051342in}{1.342698in}}%
\pgfpathlineto{\pgfqpoint{6.045190in}{1.349098in}}%
\pgfpathlineto{\pgfqpoint{6.041871in}{1.361844in}}%
\pgfpathlineto{\pgfqpoint{6.033151in}{1.363429in}}%
\pgfpathlineto{\pgfqpoint{6.027319in}{1.370116in}}%
\pgfpathlineto{\pgfqpoint{5.953076in}{1.355148in}}%
\pgfpathlineto{\pgfqpoint{5.917594in}{1.347810in}}%
\pgfpathlineto{\pgfqpoint{5.863721in}{1.337977in}}%
\pgfpathlineto{\pgfqpoint{5.824926in}{1.331432in}}%
\pgfpathlineto{\pgfqpoint{5.822084in}{1.347781in}}%
\pgfpathclose%
\pgfusepath{fill}%
\end{pgfscope}%
\begin{pgfscope}%
\pgfpathrectangle{\pgfqpoint{3.525000in}{0.100000in}}{\pgfqpoint{2.862985in}{1.832778in}}%
\pgfusepath{clip}%
\pgfsetbuttcap%
\pgfsetmiterjoin%
\definecolor{currentfill}{rgb}{0.059977,0.765813,0.754095}%
\pgfsetfillcolor{currentfill}%
\pgfsetlinewidth{0.000000pt}%
\definecolor{currentstroke}{rgb}{0.000000,0.000000,0.000000}%
\pgfsetstrokecolor{currentstroke}%
\pgfsetstrokeopacity{0.000000}%
\pgfsetdash{}{0pt}%
\pgfpathmoveto{\pgfqpoint{6.120139in}{1.301344in}}%
\pgfpathlineto{\pgfqpoint{6.108657in}{1.297767in}}%
\pgfpathlineto{\pgfqpoint{6.106726in}{1.301055in}}%
\pgfpathlineto{\pgfqpoint{6.110411in}{1.312094in}}%
\pgfpathlineto{\pgfqpoint{6.117237in}{1.313181in}}%
\pgfpathlineto{\pgfqpoint{6.116626in}{1.316656in}}%
\pgfpathlineto{\pgfqpoint{6.122754in}{1.321872in}}%
\pgfpathlineto{\pgfqpoint{6.140472in}{1.326028in}}%
\pgfpathlineto{\pgfqpoint{6.148360in}{1.332319in}}%
\pgfpathlineto{\pgfqpoint{6.166079in}{1.337561in}}%
\pgfpathlineto{\pgfqpoint{6.174161in}{1.335643in}}%
\pgfpathlineto{\pgfqpoint{6.180959in}{1.344095in}}%
\pgfpathlineto{\pgfqpoint{6.191234in}{1.345211in}}%
\pgfpathlineto{\pgfqpoint{6.173704in}{1.328725in}}%
\pgfpathlineto{\pgfqpoint{6.120139in}{1.301344in}}%
\pgfpathclose%
\pgfusepath{fill}%
\end{pgfscope}%
\begin{pgfscope}%
\pgfpathrectangle{\pgfqpoint{3.525000in}{0.100000in}}{\pgfqpoint{2.862985in}{1.832778in}}%
\pgfusepath{clip}%
\pgfsetbuttcap%
\pgfsetmiterjoin%
\definecolor{currentfill}{rgb}{0.109189,0.779592,0.552326}%
\pgfsetfillcolor{currentfill}%
\pgfsetlinewidth{0.000000pt}%
\definecolor{currentstroke}{rgb}{0.000000,0.000000,0.000000}%
\pgfsetstrokecolor{currentstroke}%
\pgfsetstrokeopacity{0.000000}%
\pgfsetdash{}{0pt}%
\pgfpathmoveto{\pgfqpoint{5.862178in}{1.192485in}}%
\pgfpathlineto{\pgfqpoint{5.812526in}{1.184275in}}%
\pgfpathlineto{\pgfqpoint{5.803519in}{1.241018in}}%
\pgfpathlineto{\pgfqpoint{5.790146in}{1.324661in}}%
\pgfpathlineto{\pgfqpoint{5.822084in}{1.347781in}}%
\pgfpathlineto{\pgfqpoint{5.824926in}{1.331432in}}%
\pgfpathlineto{\pgfqpoint{5.863721in}{1.337977in}}%
\pgfpathlineto{\pgfqpoint{5.917594in}{1.347810in}}%
\pgfpathlineto{\pgfqpoint{5.953076in}{1.355148in}}%
\pgfpathlineto{\pgfqpoint{6.027319in}{1.370116in}}%
\pgfpathlineto{\pgfqpoint{6.033151in}{1.363429in}}%
\pgfpathlineto{\pgfqpoint{6.041871in}{1.361844in}}%
\pgfpathlineto{\pgfqpoint{6.045190in}{1.349098in}}%
\pgfpathlineto{\pgfqpoint{6.051342in}{1.342698in}}%
\pgfpathlineto{\pgfqpoint{6.062035in}{1.341668in}}%
\pgfpathlineto{\pgfqpoint{6.065655in}{1.337368in}}%
\pgfpathlineto{\pgfqpoint{6.061981in}{1.334051in}}%
\pgfpathlineto{\pgfqpoint{6.058670in}{1.322314in}}%
\pgfpathlineto{\pgfqpoint{6.050913in}{1.309125in}}%
\pgfpathlineto{\pgfqpoint{6.056120in}{1.303390in}}%
\pgfpathlineto{\pgfqpoint{6.051859in}{1.297245in}}%
\pgfpathlineto{\pgfqpoint{6.054263in}{1.283769in}}%
\pgfpathlineto{\pgfqpoint{6.061936in}{1.276695in}}%
\pgfpathlineto{\pgfqpoint{6.080144in}{1.265202in}}%
\pgfpathlineto{\pgfqpoint{6.065481in}{1.248986in}}%
\pgfpathlineto{\pgfqpoint{6.065275in}{1.242831in}}%
\pgfpathlineto{\pgfqpoint{6.053338in}{1.234894in}}%
\pgfpathlineto{\pgfqpoint{6.040138in}{1.233404in}}%
\pgfpathlineto{\pgfqpoint{6.036873in}{1.226505in}}%
\pgfpathlineto{\pgfqpoint{6.000158in}{1.218561in}}%
\pgfpathlineto{\pgfqpoint{5.957316in}{1.209933in}}%
\pgfpathlineto{\pgfqpoint{5.927868in}{1.204659in}}%
\pgfpathlineto{\pgfqpoint{5.862178in}{1.192485in}}%
\pgfpathclose%
\pgfusepath{fill}%
\end{pgfscope}%
\begin{pgfscope}%
\pgfpathrectangle{\pgfqpoint{3.525000in}{0.100000in}}{\pgfqpoint{2.862985in}{1.832778in}}%
\pgfusepath{clip}%
\pgfsetbuttcap%
\pgfsetmiterjoin%
\definecolor{currentfill}{rgb}{0.167628,0.795955,0.312726}%
\pgfsetfillcolor{currentfill}%
\pgfsetlinewidth{0.000000pt}%
\definecolor{currentstroke}{rgb}{0.000000,0.000000,0.000000}%
\pgfsetstrokecolor{currentstroke}%
\pgfsetstrokeopacity{0.000000}%
\pgfsetdash{}{0pt}%
\pgfpathmoveto{\pgfqpoint{6.111166in}{1.392957in}}%
\pgfpathlineto{\pgfqpoint{6.174362in}{1.407318in}}%
\pgfpathlineto{\pgfqpoint{6.188120in}{1.409488in}}%
\pgfpathlineto{\pgfqpoint{6.197226in}{1.373686in}}%
\pgfpathlineto{\pgfqpoint{6.195785in}{1.367275in}}%
\pgfpathlineto{\pgfqpoint{6.166533in}{1.355951in}}%
\pgfpathlineto{\pgfqpoint{6.149071in}{1.351992in}}%
\pgfpathlineto{\pgfqpoint{6.141651in}{1.343114in}}%
\pgfpathlineto{\pgfqpoint{6.118832in}{1.327078in}}%
\pgfpathlineto{\pgfqpoint{6.114028in}{1.332508in}}%
\pgfpathlineto{\pgfqpoint{6.123765in}{1.341969in}}%
\pgfpathlineto{\pgfqpoint{6.119238in}{1.346227in}}%
\pgfpathlineto{\pgfqpoint{6.111166in}{1.392957in}}%
\pgfpathclose%
\pgfusepath{fill}%
\end{pgfscope}%
\begin{pgfscope}%
\pgfpathrectangle{\pgfqpoint{3.525000in}{0.100000in}}{\pgfqpoint{2.862985in}{1.832778in}}%
\pgfusepath{clip}%
\pgfsetbuttcap%
\pgfsetmiterjoin%
\definecolor{currentfill}{rgb}{0.172241,0.753864,0.172241}%
\pgfsetfillcolor{currentfill}%
\pgfsetlinewidth{0.000000pt}%
\definecolor{currentstroke}{rgb}{0.000000,0.000000,0.000000}%
\pgfsetstrokecolor{currentstroke}%
\pgfsetstrokeopacity{0.000000}%
\pgfsetdash{}{0pt}%
\pgfpathmoveto{\pgfqpoint{6.188120in}{1.409488in}}%
\pgfpathlineto{\pgfqpoint{6.206882in}{1.415045in}}%
\pgfpathlineto{\pgfqpoint{6.212340in}{1.401600in}}%
\pgfpathlineto{\pgfqpoint{6.218598in}{1.398285in}}%
\pgfpathlineto{\pgfqpoint{6.210875in}{1.392622in}}%
\pgfpathlineto{\pgfqpoint{6.211849in}{1.375988in}}%
\pgfpathlineto{\pgfqpoint{6.195785in}{1.367275in}}%
\pgfpathlineto{\pgfqpoint{6.197226in}{1.373686in}}%
\pgfpathlineto{\pgfqpoint{6.188120in}{1.409488in}}%
\pgfpathclose%
\pgfusepath{fill}%
\end{pgfscope}%
\begin{pgfscope}%
\pgfpathrectangle{\pgfqpoint{3.525000in}{0.100000in}}{\pgfqpoint{2.862985in}{1.832778in}}%
\pgfusepath{clip}%
\pgfsetbuttcap%
\pgfsetmiterjoin%
\definecolor{currentfill}{rgb}{0.170704,0.796817,0.300115}%
\pgfsetfillcolor{currentfill}%
\pgfsetlinewidth{0.000000pt}%
\definecolor{currentstroke}{rgb}{0.000000,0.000000,0.000000}%
\pgfsetstrokecolor{currentstroke}%
\pgfsetstrokeopacity{0.000000}%
\pgfsetdash{}{0pt}%
\pgfpathmoveto{\pgfqpoint{6.051223in}{1.229031in}}%
\pgfpathlineto{\pgfqpoint{6.053338in}{1.234894in}}%
\pgfpathlineto{\pgfqpoint{6.065275in}{1.242831in}}%
\pgfpathlineto{\pgfqpoint{6.065481in}{1.248986in}}%
\pgfpathlineto{\pgfqpoint{6.080144in}{1.265202in}}%
\pgfpathlineto{\pgfqpoint{6.061936in}{1.276695in}}%
\pgfpathlineto{\pgfqpoint{6.054263in}{1.283769in}}%
\pgfpathlineto{\pgfqpoint{6.051859in}{1.297245in}}%
\pgfpathlineto{\pgfqpoint{6.056120in}{1.303390in}}%
\pgfpathlineto{\pgfqpoint{6.050913in}{1.309125in}}%
\pgfpathlineto{\pgfqpoint{6.058670in}{1.322314in}}%
\pgfpathlineto{\pgfqpoint{6.061981in}{1.334051in}}%
\pgfpathlineto{\pgfqpoint{6.065655in}{1.337368in}}%
\pgfpathlineto{\pgfqpoint{6.091246in}{1.328639in}}%
\pgfpathlineto{\pgfqpoint{6.107610in}{1.324417in}}%
\pgfpathlineto{\pgfqpoint{6.106795in}{1.305239in}}%
\pgfpathlineto{\pgfqpoint{6.096854in}{1.290696in}}%
\pgfpathlineto{\pgfqpoint{6.105045in}{1.288556in}}%
\pgfpathlineto{\pgfqpoint{6.113550in}{1.282315in}}%
\pgfpathlineto{\pgfqpoint{6.114051in}{1.265245in}}%
\pgfpathlineto{\pgfqpoint{6.111477in}{1.253161in}}%
\pgfpathlineto{\pgfqpoint{6.113190in}{1.243235in}}%
\pgfpathlineto{\pgfqpoint{6.098407in}{1.209740in}}%
\pgfpathlineto{\pgfqpoint{6.093094in}{1.194100in}}%
\pgfpathlineto{\pgfqpoint{6.085676in}{1.201703in}}%
\pgfpathlineto{\pgfqpoint{6.075838in}{1.200407in}}%
\pgfpathlineto{\pgfqpoint{6.051218in}{1.214632in}}%
\pgfpathlineto{\pgfqpoint{6.048699in}{1.222249in}}%
\pgfpathlineto{\pgfqpoint{6.051223in}{1.229031in}}%
\pgfpathclose%
\pgfusepath{fill}%
\end{pgfscope}%
\begin{pgfscope}%
\pgfpathrectangle{\pgfqpoint{3.525000in}{0.100000in}}{\pgfqpoint{2.862985in}{1.832778in}}%
\pgfusepath{clip}%
\pgfsetbuttcap%
\pgfsetmiterjoin%
\definecolor{currentfill}{rgb}{0.127643,0.784760,0.476663}%
\pgfsetfillcolor{currentfill}%
\pgfsetlinewidth{0.000000pt}%
\definecolor{currentstroke}{rgb}{0.000000,0.000000,0.000000}%
\pgfsetstrokecolor{currentstroke}%
\pgfsetstrokeopacity{0.000000}%
\pgfsetdash{}{0pt}%
\pgfpathmoveto{\pgfqpoint{5.462013in}{1.021868in}}%
\pgfpathlineto{\pgfqpoint{5.460862in}{1.037407in}}%
\pgfpathlineto{\pgfqpoint{5.465336in}{1.045116in}}%
\pgfpathlineto{\pgfqpoint{5.462387in}{1.048893in}}%
\pgfpathlineto{\pgfqpoint{5.473140in}{1.061261in}}%
\pgfpathlineto{\pgfqpoint{5.472524in}{1.064241in}}%
\pgfpathlineto{\pgfqpoint{5.482971in}{1.085818in}}%
\pgfpathlineto{\pgfqpoint{5.480859in}{1.096224in}}%
\pgfpathlineto{\pgfqpoint{5.473283in}{1.107118in}}%
\pgfpathlineto{\pgfqpoint{5.478542in}{1.120373in}}%
\pgfpathlineto{\pgfqpoint{5.471726in}{1.207597in}}%
\pgfpathlineto{\pgfqpoint{5.466818in}{1.268787in}}%
\pgfpathlineto{\pgfqpoint{5.473598in}{1.263717in}}%
\pgfpathlineto{\pgfqpoint{5.481163in}{1.263854in}}%
\pgfpathlineto{\pgfqpoint{5.499029in}{1.274203in}}%
\pgfpathlineto{\pgfqpoint{5.553831in}{1.279317in}}%
\pgfpathlineto{\pgfqpoint{5.594393in}{1.283592in}}%
\pgfpathlineto{\pgfqpoint{5.594752in}{1.279623in}}%
\pgfpathlineto{\pgfqpoint{5.604011in}{1.195770in}}%
\pgfpathlineto{\pgfqpoint{5.611981in}{1.117748in}}%
\pgfpathlineto{\pgfqpoint{5.608545in}{1.114085in}}%
\pgfpathlineto{\pgfqpoint{5.613822in}{1.104993in}}%
\pgfpathlineto{\pgfqpoint{5.613795in}{1.098441in}}%
\pgfpathlineto{\pgfqpoint{5.606265in}{1.096795in}}%
\pgfpathlineto{\pgfqpoint{5.597839in}{1.090478in}}%
\pgfpathlineto{\pgfqpoint{5.592134in}{1.092965in}}%
\pgfpathlineto{\pgfqpoint{5.583595in}{1.088922in}}%
\pgfpathlineto{\pgfqpoint{5.586237in}{1.080801in}}%
\pgfpathlineto{\pgfqpoint{5.577460in}{1.072642in}}%
\pgfpathlineto{\pgfqpoint{5.575037in}{1.063205in}}%
\pgfpathlineto{\pgfqpoint{5.569002in}{1.061656in}}%
\pgfpathlineto{\pgfqpoint{5.564498in}{1.054509in}}%
\pgfpathlineto{\pgfqpoint{5.565088in}{1.047311in}}%
\pgfpathlineto{\pgfqpoint{5.559790in}{1.042250in}}%
\pgfpathlineto{\pgfqpoint{5.551817in}{1.043031in}}%
\pgfpathlineto{\pgfqpoint{5.545757in}{1.050797in}}%
\pgfpathlineto{\pgfqpoint{5.535489in}{1.043318in}}%
\pgfpathlineto{\pgfqpoint{5.536209in}{1.036783in}}%
\pgfpathlineto{\pgfqpoint{5.524789in}{1.032965in}}%
\pgfpathlineto{\pgfqpoint{5.521919in}{1.037771in}}%
\pgfpathlineto{\pgfqpoint{5.510801in}{1.032321in}}%
\pgfpathlineto{\pgfqpoint{5.506739in}{1.024520in}}%
\pgfpathlineto{\pgfqpoint{5.493349in}{1.032520in}}%
\pgfpathlineto{\pgfqpoint{5.472127in}{1.027220in}}%
\pgfpathlineto{\pgfqpoint{5.462013in}{1.021868in}}%
\pgfpathclose%
\pgfusepath{fill}%
\end{pgfscope}%
\begin{pgfscope}%
\pgfpathrectangle{\pgfqpoint{3.525000in}{0.100000in}}{\pgfqpoint{2.862985in}{1.832778in}}%
\pgfusepath{clip}%
\pgfsetbuttcap%
\pgfsetmiterjoin%
\definecolor{currentfill}{rgb}{0.139946,0.788205,0.426221}%
\pgfsetfillcolor{currentfill}%
\pgfsetlinewidth{0.000000pt}%
\definecolor{currentstroke}{rgb}{0.000000,0.000000,0.000000}%
\pgfsetstrokecolor{currentstroke}%
\pgfsetstrokeopacity{0.000000}%
\pgfsetdash{}{0pt}%
\pgfpathmoveto{\pgfqpoint{3.965358in}{1.430483in}}%
\pgfpathlineto{\pgfqpoint{3.996046in}{1.421943in}}%
\pgfpathlineto{\pgfqpoint{4.046916in}{1.408826in}}%
\pgfpathlineto{\pgfqpoint{4.099647in}{1.395314in}}%
\pgfpathlineto{\pgfqpoint{4.148250in}{1.383593in}}%
\pgfpathlineto{\pgfqpoint{4.190246in}{1.374085in}}%
\pgfpathlineto{\pgfqpoint{4.235177in}{1.364205in}}%
\pgfpathlineto{\pgfqpoint{4.222195in}{1.302942in}}%
\pgfpathlineto{\pgfqpoint{4.210630in}{1.248537in}}%
\pgfpathlineto{\pgfqpoint{4.191658in}{1.160754in}}%
\pgfpathlineto{\pgfqpoint{4.177414in}{1.094498in}}%
\pgfpathlineto{\pgfqpoint{4.169719in}{1.057523in}}%
\pgfpathlineto{\pgfqpoint{4.159852in}{1.009534in}}%
\pgfpathlineto{\pgfqpoint{4.153025in}{0.999798in}}%
\pgfpathlineto{\pgfqpoint{4.147531in}{0.999473in}}%
\pgfpathlineto{\pgfqpoint{4.143612in}{1.007971in}}%
\pgfpathlineto{\pgfqpoint{4.134613in}{1.011057in}}%
\pgfpathlineto{\pgfqpoint{4.124916in}{1.009976in}}%
\pgfpathlineto{\pgfqpoint{4.122164in}{1.003024in}}%
\pgfpathlineto{\pgfqpoint{4.121877in}{0.978886in}}%
\pgfpathlineto{\pgfqpoint{4.118973in}{0.973374in}}%
\pgfpathlineto{\pgfqpoint{4.120633in}{0.953995in}}%
\pgfpathlineto{\pgfqpoint{4.114569in}{0.941085in}}%
\pgfpathlineto{\pgfqpoint{4.065397in}{1.016802in}}%
\pgfpathlineto{\pgfqpoint{4.016971in}{1.090379in}}%
\pgfpathlineto{\pgfqpoint{3.991763in}{1.128960in}}%
\pgfpathlineto{\pgfqpoint{3.970735in}{1.162262in}}%
\pgfpathlineto{\pgfqpoint{3.944238in}{1.203456in}}%
\pgfpathlineto{\pgfqpoint{3.914244in}{1.249665in}}%
\pgfpathlineto{\pgfqpoint{3.926576in}{1.293524in}}%
\pgfpathlineto{\pgfqpoint{3.951391in}{1.381486in}}%
\pgfpathlineto{\pgfqpoint{3.965358in}{1.430483in}}%
\pgfpathclose%
\pgfusepath{fill}%
\end{pgfscope}%
\begin{pgfscope}%
\pgfpathrectangle{\pgfqpoint{3.525000in}{0.100000in}}{\pgfqpoint{2.862985in}{1.832778in}}%
\pgfusepath{clip}%
\pgfsetbuttcap%
\pgfsetmiterjoin%
\definecolor{currentfill}{rgb}{0.161476,0.794233,0.337947}%
\pgfsetfillcolor{currentfill}%
\pgfsetlinewidth{0.000000pt}%
\definecolor{currentstroke}{rgb}{0.000000,0.000000,0.000000}%
\pgfsetstrokecolor{currentstroke}%
\pgfsetstrokeopacity{0.000000}%
\pgfsetdash{}{0pt}%
\pgfpathmoveto{\pgfqpoint{4.235177in}{1.364205in}}%
\pgfpathlineto{\pgfqpoint{4.283212in}{1.354650in}}%
\pgfpathlineto{\pgfqpoint{4.372141in}{1.337483in}}%
\pgfpathlineto{\pgfqpoint{4.361009in}{1.275734in}}%
\pgfpathlineto{\pgfqpoint{4.409967in}{1.267352in}}%
\pgfpathlineto{\pgfqpoint{4.454571in}{1.260214in}}%
\pgfpathlineto{\pgfqpoint{4.446820in}{1.211396in}}%
\pgfpathlineto{\pgfqpoint{4.438608in}{1.158753in}}%
\pgfpathlineto{\pgfqpoint{4.427613in}{1.089632in}}%
\pgfpathlineto{\pgfqpoint{4.427328in}{1.083837in}}%
\pgfpathlineto{\pgfqpoint{4.415915in}{1.012202in}}%
\pgfpathlineto{\pgfqpoint{4.368932in}{1.019484in}}%
\pgfpathlineto{\pgfqpoint{4.344991in}{1.024292in}}%
\pgfpathlineto{\pgfqpoint{4.258455in}{1.039504in}}%
\pgfpathlineto{\pgfqpoint{4.225868in}{1.045930in}}%
\pgfpathlineto{\pgfqpoint{4.169719in}{1.057523in}}%
\pgfpathlineto{\pgfqpoint{4.177414in}{1.094498in}}%
\pgfpathlineto{\pgfqpoint{4.191658in}{1.160754in}}%
\pgfpathlineto{\pgfqpoint{4.210630in}{1.248537in}}%
\pgfpathlineto{\pgfqpoint{4.222195in}{1.302942in}}%
\pgfpathlineto{\pgfqpoint{4.235177in}{1.364205in}}%
\pgfpathclose%
\pgfusepath{fill}%
\end{pgfscope}%
\begin{pgfscope}%
\pgfpathrectangle{\pgfqpoint{3.525000in}{0.100000in}}{\pgfqpoint{2.862985in}{1.832778in}}%
\pgfusepath{clip}%
\pgfsetbuttcap%
\pgfsetmiterjoin%
\definecolor{currentfill}{rgb}{0.038447,0.759785,0.842368}%
\pgfsetfillcolor{currentfill}%
\pgfsetlinewidth{0.000000pt}%
\definecolor{currentstroke}{rgb}{0.000000,0.000000,0.000000}%
\pgfsetstrokecolor{currentstroke}%
\pgfsetstrokeopacity{0.000000}%
\pgfsetdash{}{0pt}%
\pgfpathmoveto{\pgfqpoint{3.777592in}{1.488751in}}%
\pgfpathlineto{\pgfqpoint{3.794661in}{1.482798in}}%
\pgfpathlineto{\pgfqpoint{3.863032in}{1.461511in}}%
\pgfpathlineto{\pgfqpoint{3.925927in}{1.441715in}}%
\pgfpathlineto{\pgfqpoint{3.965358in}{1.430483in}}%
\pgfpathlineto{\pgfqpoint{3.951391in}{1.381486in}}%
\pgfpathlineto{\pgfqpoint{3.926576in}{1.293524in}}%
\pgfpathlineto{\pgfqpoint{3.914244in}{1.249665in}}%
\pgfpathlineto{\pgfqpoint{3.944238in}{1.203456in}}%
\pgfpathlineto{\pgfqpoint{3.970735in}{1.162262in}}%
\pgfpathlineto{\pgfqpoint{3.991763in}{1.128960in}}%
\pgfpathlineto{\pgfqpoint{4.016971in}{1.090379in}}%
\pgfpathlineto{\pgfqpoint{4.065397in}{1.016802in}}%
\pgfpathlineto{\pgfqpoint{4.114569in}{0.941085in}}%
\pgfpathlineto{\pgfqpoint{4.112595in}{0.933576in}}%
\pgfpathlineto{\pgfqpoint{4.118526in}{0.921617in}}%
\pgfpathlineto{\pgfqpoint{4.119699in}{0.905259in}}%
\pgfpathlineto{\pgfqpoint{4.128313in}{0.897324in}}%
\pgfpathlineto{\pgfqpoint{4.128654in}{0.890930in}}%
\pgfpathlineto{\pgfqpoint{4.113230in}{0.883673in}}%
\pgfpathlineto{\pgfqpoint{4.105882in}{0.876407in}}%
\pgfpathlineto{\pgfqpoint{4.103591in}{0.860364in}}%
\pgfpathlineto{\pgfqpoint{4.099871in}{0.851614in}}%
\pgfpathlineto{\pgfqpoint{4.092035in}{0.844239in}}%
\pgfpathlineto{\pgfqpoint{4.084240in}{0.825061in}}%
\pgfpathlineto{\pgfqpoint{4.095200in}{0.815071in}}%
\pgfpathlineto{\pgfqpoint{4.093785in}{0.806836in}}%
\pgfpathlineto{\pgfqpoint{4.084798in}{0.801105in}}%
\pgfpathlineto{\pgfqpoint{4.078603in}{0.802117in}}%
\pgfpathlineto{\pgfqpoint{4.005470in}{0.812266in}}%
\pgfpathlineto{\pgfqpoint{3.951454in}{0.819928in}}%
\pgfpathlineto{\pgfqpoint{3.953816in}{0.828649in}}%
\pgfpathlineto{\pgfqpoint{3.950300in}{0.843128in}}%
\pgfpathlineto{\pgfqpoint{3.949945in}{0.857733in}}%
\pgfpathlineto{\pgfqpoint{3.947647in}{0.866279in}}%
\pgfpathlineto{\pgfqpoint{3.940568in}{0.878498in}}%
\pgfpathlineto{\pgfqpoint{3.920206in}{0.906691in}}%
\pgfpathlineto{\pgfqpoint{3.913537in}{0.910145in}}%
\pgfpathlineto{\pgfqpoint{3.904944in}{0.910090in}}%
\pgfpathlineto{\pgfqpoint{3.906908in}{0.918984in}}%
\pgfpathlineto{\pgfqpoint{3.902839in}{0.930106in}}%
\pgfpathlineto{\pgfqpoint{3.882853in}{0.935633in}}%
\pgfpathlineto{\pgfqpoint{3.870671in}{0.945869in}}%
\pgfpathlineto{\pgfqpoint{3.869661in}{0.952132in}}%
\pgfpathlineto{\pgfqpoint{3.855633in}{0.967637in}}%
\pgfpathlineto{\pgfqpoint{3.842280in}{0.970604in}}%
\pgfpathlineto{\pgfqpoint{3.829908in}{0.978484in}}%
\pgfpathlineto{\pgfqpoint{3.813637in}{0.981210in}}%
\pgfpathlineto{\pgfqpoint{3.806685in}{0.991719in}}%
\pgfpathlineto{\pgfqpoint{3.813287in}{1.008354in}}%
\pgfpathlineto{\pgfqpoint{3.811268in}{1.012095in}}%
\pgfpathlineto{\pgfqpoint{3.816766in}{1.025962in}}%
\pgfpathlineto{\pgfqpoint{3.806986in}{1.033342in}}%
\pgfpathlineto{\pgfqpoint{3.810159in}{1.046721in}}%
\pgfpathlineto{\pgfqpoint{3.804935in}{1.050141in}}%
\pgfpathlineto{\pgfqpoint{3.800424in}{1.062787in}}%
\pgfpathlineto{\pgfqpoint{3.794981in}{1.066644in}}%
\pgfpathlineto{\pgfqpoint{3.794572in}{1.075778in}}%
\pgfpathlineto{\pgfqpoint{3.790307in}{1.082219in}}%
\pgfpathlineto{\pgfqpoint{3.783881in}{1.103911in}}%
\pgfpathlineto{\pgfqpoint{3.776850in}{1.114301in}}%
\pgfpathlineto{\pgfqpoint{3.778378in}{1.131939in}}%
\pgfpathlineto{\pgfqpoint{3.786647in}{1.133714in}}%
\pgfpathlineto{\pgfqpoint{3.792031in}{1.143280in}}%
\pgfpathlineto{\pgfqpoint{3.788778in}{1.153656in}}%
\pgfpathlineto{\pgfqpoint{3.779984in}{1.155388in}}%
\pgfpathlineto{\pgfqpoint{3.775596in}{1.160241in}}%
\pgfpathlineto{\pgfqpoint{3.768491in}{1.178092in}}%
\pgfpathlineto{\pgfqpoint{3.771812in}{1.184503in}}%
\pgfpathlineto{\pgfqpoint{3.769457in}{1.196440in}}%
\pgfpathlineto{\pgfqpoint{3.774672in}{1.211900in}}%
\pgfpathlineto{\pgfqpoint{3.780782in}{1.206207in}}%
\pgfpathlineto{\pgfqpoint{3.778009in}{1.199485in}}%
\pgfpathlineto{\pgfqpoint{3.788597in}{1.188830in}}%
\pgfpathlineto{\pgfqpoint{3.787921in}{1.204652in}}%
\pgfpathlineto{\pgfqpoint{3.783384in}{1.208896in}}%
\pgfpathlineto{\pgfqpoint{3.788570in}{1.222808in}}%
\pgfpathlineto{\pgfqpoint{3.808185in}{1.220593in}}%
\pgfpathlineto{\pgfqpoint{3.805545in}{1.225766in}}%
\pgfpathlineto{\pgfqpoint{3.792599in}{1.225257in}}%
\pgfpathlineto{\pgfqpoint{3.786441in}{1.233093in}}%
\pgfpathlineto{\pgfqpoint{3.778689in}{1.226135in}}%
\pgfpathlineto{\pgfqpoint{3.774576in}{1.214504in}}%
\pgfpathlineto{\pgfqpoint{3.763597in}{1.230154in}}%
\pgfpathlineto{\pgfqpoint{3.759369in}{1.233003in}}%
\pgfpathlineto{\pgfqpoint{3.760980in}{1.250036in}}%
\pgfpathlineto{\pgfqpoint{3.757634in}{1.260081in}}%
\pgfpathlineto{\pgfqpoint{3.751568in}{1.269517in}}%
\pgfpathlineto{\pgfqpoint{3.739145in}{1.298432in}}%
\pgfpathlineto{\pgfqpoint{3.743207in}{1.304826in}}%
\pgfpathlineto{\pgfqpoint{3.743159in}{1.325033in}}%
\pgfpathlineto{\pgfqpoint{3.749864in}{1.336303in}}%
\pgfpathlineto{\pgfqpoint{3.751401in}{1.353918in}}%
\pgfpathlineto{\pgfqpoint{3.745079in}{1.374099in}}%
\pgfpathlineto{\pgfqpoint{3.736692in}{1.386945in}}%
\pgfpathlineto{\pgfqpoint{3.738188in}{1.398536in}}%
\pgfpathlineto{\pgfqpoint{3.761740in}{1.426591in}}%
\pgfpathlineto{\pgfqpoint{3.762910in}{1.436161in}}%
\pgfpathlineto{\pgfqpoint{3.773523in}{1.454416in}}%
\pgfpathlineto{\pgfqpoint{3.774995in}{1.471719in}}%
\pgfpathlineto{\pgfqpoint{3.771583in}{1.476134in}}%
\pgfpathlineto{\pgfqpoint{3.777592in}{1.488751in}}%
\pgfpathclose%
\pgfusepath{fill}%
\end{pgfscope}%
\begin{pgfscope}%
\pgfpathrectangle{\pgfqpoint{3.525000in}{0.100000in}}{\pgfqpoint{2.862985in}{1.832778in}}%
\pgfusepath{clip}%
\pgfsetbuttcap%
\pgfsetmiterjoin%
\definecolor{currentfill}{rgb}{0.127643,0.784760,0.476663}%
\pgfsetfillcolor{currentfill}%
\pgfsetlinewidth{0.000000pt}%
\definecolor{currentstroke}{rgb}{0.000000,0.000000,0.000000}%
\pgfsetstrokecolor{currentstroke}%
\pgfsetstrokeopacity{0.000000}%
\pgfsetdash{}{0pt}%
\pgfpathmoveto{\pgfqpoint{5.611981in}{1.117748in}}%
\pgfpathlineto{\pgfqpoint{5.604011in}{1.195770in}}%
\pgfpathlineto{\pgfqpoint{5.594752in}{1.279623in}}%
\pgfpathlineto{\pgfqpoint{5.655390in}{1.288653in}}%
\pgfpathlineto{\pgfqpoint{5.671483in}{1.284465in}}%
\pgfpathlineto{\pgfqpoint{5.679202in}{1.279918in}}%
\pgfpathlineto{\pgfqpoint{5.688874in}{1.281179in}}%
\pgfpathlineto{\pgfqpoint{5.701627in}{1.273651in}}%
\pgfpathlineto{\pgfqpoint{5.725380in}{1.284855in}}%
\pgfpathlineto{\pgfqpoint{5.738489in}{1.285229in}}%
\pgfpathlineto{\pgfqpoint{5.753797in}{1.302307in}}%
\pgfpathlineto{\pgfqpoint{5.769372in}{1.312695in}}%
\pgfpathlineto{\pgfqpoint{5.790146in}{1.324661in}}%
\pgfpathlineto{\pgfqpoint{5.803519in}{1.241018in}}%
\pgfpathlineto{\pgfqpoint{5.797306in}{1.235646in}}%
\pgfpathlineto{\pgfqpoint{5.801317in}{1.230712in}}%
\pgfpathlineto{\pgfqpoint{5.802915in}{1.219896in}}%
\pgfpathlineto{\pgfqpoint{5.799803in}{1.209737in}}%
\pgfpathlineto{\pgfqpoint{5.799391in}{1.194887in}}%
\pgfpathlineto{\pgfqpoint{5.796468in}{1.175560in}}%
\pgfpathlineto{\pgfqpoint{5.781678in}{1.158322in}}%
\pgfpathlineto{\pgfqpoint{5.775452in}{1.154663in}}%
\pgfpathlineto{\pgfqpoint{5.770610in}{1.157813in}}%
\pgfpathlineto{\pgfqpoint{5.758627in}{1.141383in}}%
\pgfpathlineto{\pgfqpoint{5.759750in}{1.125564in}}%
\pgfpathlineto{\pgfqpoint{5.746406in}{1.129368in}}%
\pgfpathlineto{\pgfqpoint{5.740087in}{1.113569in}}%
\pgfpathlineto{\pgfqpoint{5.743281in}{1.102373in}}%
\pgfpathlineto{\pgfqpoint{5.738281in}{1.100719in}}%
\pgfpathlineto{\pgfqpoint{5.737581in}{1.091873in}}%
\pgfpathlineto{\pgfqpoint{5.725308in}{1.088305in}}%
\pgfpathlineto{\pgfqpoint{5.718930in}{1.095447in}}%
\pgfpathlineto{\pgfqpoint{5.710724in}{1.098217in}}%
\pgfpathlineto{\pgfqpoint{5.708796in}{1.105459in}}%
\pgfpathlineto{\pgfqpoint{5.701376in}{1.104185in}}%
\pgfpathlineto{\pgfqpoint{5.696515in}{1.097524in}}%
\pgfpathlineto{\pgfqpoint{5.689561in}{1.095181in}}%
\pgfpathlineto{\pgfqpoint{5.677277in}{1.099895in}}%
\pgfpathlineto{\pgfqpoint{5.670483in}{1.094286in}}%
\pgfpathlineto{\pgfqpoint{5.660827in}{1.100922in}}%
\pgfpathlineto{\pgfqpoint{5.642303in}{1.102959in}}%
\pgfpathlineto{\pgfqpoint{5.636715in}{1.115014in}}%
\pgfpathlineto{\pgfqpoint{5.629636in}{1.120355in}}%
\pgfpathlineto{\pgfqpoint{5.622773in}{1.116936in}}%
\pgfpathlineto{\pgfqpoint{5.611981in}{1.117748in}}%
\pgfpathclose%
\pgfusepath{fill}%
\end{pgfscope}%
\begin{pgfscope}%
\pgfpathrectangle{\pgfqpoint{3.525000in}{0.100000in}}{\pgfqpoint{2.862985in}{1.832778in}}%
\pgfusepath{clip}%
\pgfsetbuttcap%
\pgfsetmiterjoin%
\definecolor{currentfill}{rgb}{0.090734,0.774425,0.627989}%
\pgfsetfillcolor{currentfill}%
\pgfsetlinewidth{0.000000pt}%
\definecolor{currentstroke}{rgb}{0.000000,0.000000,0.000000}%
\pgfsetstrokecolor{currentstroke}%
\pgfsetstrokeopacity{0.000000}%
\pgfsetdash{}{0pt}%
\pgfpathmoveto{\pgfqpoint{5.411225in}{0.966856in}}%
\pgfpathlineto{\pgfqpoint{5.404121in}{0.972628in}}%
\pgfpathlineto{\pgfqpoint{5.398328in}{0.969882in}}%
\pgfpathlineto{\pgfqpoint{5.390922in}{0.983711in}}%
\pgfpathlineto{\pgfqpoint{5.394705in}{0.992402in}}%
\pgfpathlineto{\pgfqpoint{5.389252in}{1.002186in}}%
\pgfpathlineto{\pgfqpoint{5.388959in}{1.009888in}}%
\pgfpathlineto{\pgfqpoint{5.381591in}{1.012632in}}%
\pgfpathlineto{\pgfqpoint{5.378174in}{1.018436in}}%
\pgfpathlineto{\pgfqpoint{5.363768in}{1.025684in}}%
\pgfpathlineto{\pgfqpoint{5.351256in}{1.034614in}}%
\pgfpathlineto{\pgfqpoint{5.345442in}{1.041355in}}%
\pgfpathlineto{\pgfqpoint{5.345322in}{1.049578in}}%
\pgfpathlineto{\pgfqpoint{5.353102in}{1.065365in}}%
\pgfpathlineto{\pgfqpoint{5.355495in}{1.077440in}}%
\pgfpathlineto{\pgfqpoint{5.349157in}{1.084263in}}%
\pgfpathlineto{\pgfqpoint{5.340763in}{1.086834in}}%
\pgfpathlineto{\pgfqpoint{5.330579in}{1.081212in}}%
\pgfpathlineto{\pgfqpoint{5.326269in}{1.090863in}}%
\pgfpathlineto{\pgfqpoint{5.324094in}{1.103945in}}%
\pgfpathlineto{\pgfqpoint{5.309086in}{1.115607in}}%
\pgfpathlineto{\pgfqpoint{5.306087in}{1.120782in}}%
\pgfpathlineto{\pgfqpoint{5.292379in}{1.132498in}}%
\pgfpathlineto{\pgfqpoint{5.288083in}{1.141015in}}%
\pgfpathlineto{\pgfqpoint{5.284210in}{1.157906in}}%
\pgfpathlineto{\pgfqpoint{5.286865in}{1.172908in}}%
\pgfpathlineto{\pgfqpoint{5.290408in}{1.174999in}}%
\pgfpathlineto{\pgfqpoint{5.289771in}{1.187558in}}%
\pgfpathlineto{\pgfqpoint{5.299773in}{1.191288in}}%
\pgfpathlineto{\pgfqpoint{5.302793in}{1.202558in}}%
\pgfpathlineto{\pgfqpoint{5.308547in}{1.210143in}}%
\pgfpathlineto{\pgfqpoint{5.308256in}{1.219794in}}%
\pgfpathlineto{\pgfqpoint{5.301124in}{1.227469in}}%
\pgfpathlineto{\pgfqpoint{5.302818in}{1.238212in}}%
\pgfpathlineto{\pgfqpoint{5.321306in}{1.242884in}}%
\pgfpathlineto{\pgfqpoint{5.335490in}{1.251413in}}%
\pgfpathlineto{\pgfqpoint{5.336979in}{1.262151in}}%
\pgfpathlineto{\pgfqpoint{5.341914in}{1.265532in}}%
\pgfpathlineto{\pgfqpoint{5.343802in}{1.276813in}}%
\pgfpathlineto{\pgfqpoint{5.342285in}{1.284265in}}%
\pgfpathlineto{\pgfqpoint{5.332588in}{1.290454in}}%
\pgfpathlineto{\pgfqpoint{5.328684in}{1.299681in}}%
\pgfpathlineto{\pgfqpoint{5.319114in}{1.308600in}}%
\pgfpathlineto{\pgfqpoint{5.397758in}{1.311952in}}%
\pgfpathlineto{\pgfqpoint{5.450520in}{1.315700in}}%
\pgfpathlineto{\pgfqpoint{5.449554in}{1.304601in}}%
\pgfpathlineto{\pgfqpoint{5.458546in}{1.289293in}}%
\pgfpathlineto{\pgfqpoint{5.462322in}{1.276213in}}%
\pgfpathlineto{\pgfqpoint{5.466818in}{1.268787in}}%
\pgfpathlineto{\pgfqpoint{5.471726in}{1.207597in}}%
\pgfpathlineto{\pgfqpoint{5.478542in}{1.120373in}}%
\pgfpathlineto{\pgfqpoint{5.473283in}{1.107118in}}%
\pgfpathlineto{\pgfqpoint{5.480859in}{1.096224in}}%
\pgfpathlineto{\pgfqpoint{5.482971in}{1.085818in}}%
\pgfpathlineto{\pgfqpoint{5.472524in}{1.064241in}}%
\pgfpathlineto{\pgfqpoint{5.473140in}{1.061261in}}%
\pgfpathlineto{\pgfqpoint{5.462387in}{1.048893in}}%
\pgfpathlineto{\pgfqpoint{5.465336in}{1.045116in}}%
\pgfpathlineto{\pgfqpoint{5.460862in}{1.037407in}}%
\pgfpathlineto{\pgfqpoint{5.462013in}{1.021868in}}%
\pgfpathlineto{\pgfqpoint{5.456580in}{1.012333in}}%
\pgfpathlineto{\pgfqpoint{5.460999in}{1.001067in}}%
\pgfpathlineto{\pgfqpoint{5.442483in}{0.994944in}}%
\pgfpathlineto{\pgfqpoint{5.440777in}{0.988286in}}%
\pgfpathlineto{\pgfqpoint{5.445833in}{0.979847in}}%
\pgfpathlineto{\pgfqpoint{5.443506in}{0.974346in}}%
\pgfpathlineto{\pgfqpoint{5.423652in}{0.981141in}}%
\pgfpathlineto{\pgfqpoint{5.417101in}{0.981825in}}%
\pgfpathlineto{\pgfqpoint{5.408934in}{0.971495in}}%
\pgfpathlineto{\pgfqpoint{5.411225in}{0.966856in}}%
\pgfpathclose%
\pgfusepath{fill}%
\end{pgfscope}%
\begin{pgfscope}%
\pgfpathrectangle{\pgfqpoint{3.525000in}{0.100000in}}{\pgfqpoint{2.862985in}{1.832778in}}%
\pgfusepath{clip}%
\pgfsetbuttcap%
\pgfsetmiterjoin%
\definecolor{currentfill}{rgb}{0.135333,0.676355,0.135333}%
\pgfsetfillcolor{currentfill}%
\pgfsetlinewidth{0.000000pt}%
\definecolor{currentstroke}{rgb}{0.000000,0.000000,0.000000}%
\pgfsetstrokecolor{currentstroke}%
\pgfsetstrokeopacity{0.000000}%
\pgfsetdash{}{0pt}%
\pgfpathmoveto{\pgfqpoint{5.989161in}{1.156689in}}%
\pgfpathlineto{\pgfqpoint{5.983697in}{1.164803in}}%
\pgfpathlineto{\pgfqpoint{5.986780in}{1.169343in}}%
\pgfpathlineto{\pgfqpoint{5.994331in}{1.164244in}}%
\pgfpathlineto{\pgfqpoint{5.989161in}{1.156689in}}%
\pgfpathclose%
\pgfusepath{fill}%
\end{pgfscope}%
\begin{pgfscope}%
\pgfpathrectangle{\pgfqpoint{3.525000in}{0.100000in}}{\pgfqpoint{2.862985in}{1.832778in}}%
\pgfusepath{clip}%
\pgfsetbuttcap%
\pgfsetmiterjoin%
\definecolor{currentfill}{rgb}{0.155325,0.792511,0.363168}%
\pgfsetfillcolor{currentfill}%
\pgfsetlinewidth{0.000000pt}%
\definecolor{currentstroke}{rgb}{0.000000,0.000000,0.000000}%
\pgfsetstrokecolor{currentstroke}%
\pgfsetstrokeopacity{0.000000}%
\pgfsetdash{}{0pt}%
\pgfpathmoveto{\pgfqpoint{6.036873in}{1.226505in}}%
\pgfpathlineto{\pgfqpoint{6.040138in}{1.233404in}}%
\pgfpathlineto{\pgfqpoint{6.053338in}{1.234894in}}%
\pgfpathlineto{\pgfqpoint{6.051223in}{1.229031in}}%
\pgfpathlineto{\pgfqpoint{6.046868in}{1.221542in}}%
\pgfpathlineto{\pgfqpoint{6.049819in}{1.212619in}}%
\pgfpathlineto{\pgfqpoint{6.061465in}{1.201927in}}%
\pgfpathlineto{\pgfqpoint{6.064168in}{1.190668in}}%
\pgfpathlineto{\pgfqpoint{6.077585in}{1.176648in}}%
\pgfpathlineto{\pgfqpoint{6.082849in}{1.177248in}}%
\pgfpathlineto{\pgfqpoint{6.089406in}{1.156201in}}%
\pgfpathlineto{\pgfqpoint{6.088331in}{1.155991in}}%
\pgfpathlineto{\pgfqpoint{6.087135in}{1.155756in}}%
\pgfpathlineto{\pgfqpoint{6.057891in}{1.150158in}}%
\pgfpathlineto{\pgfqpoint{6.042245in}{1.205802in}}%
\pgfpathlineto{\pgfqpoint{6.036873in}{1.226505in}}%
\pgfpathclose%
\pgfusepath{fill}%
\end{pgfscope}%
\begin{pgfscope}%
\pgfpathrectangle{\pgfqpoint{3.525000in}{0.100000in}}{\pgfqpoint{2.862985in}{1.832778in}}%
\pgfusepath{clip}%
\pgfsetbuttcap%
\pgfsetmiterjoin%
\definecolor{currentfill}{rgb}{0.135333,0.774394,0.980777}%
\pgfsetfillcolor{currentfill}%
\pgfsetlinewidth{0.000000pt}%
\definecolor{currentstroke}{rgb}{0.000000,0.000000,0.000000}%
\pgfsetstrokecolor{currentstroke}%
\pgfsetstrokeopacity{0.000000}%
\pgfsetdash{}{0pt}%
\pgfpathmoveto{\pgfqpoint{5.763724in}{1.037893in}}%
\pgfpathlineto{\pgfqpoint{5.754483in}{1.038233in}}%
\pgfpathlineto{\pgfqpoint{5.745968in}{1.044098in}}%
\pgfpathlineto{\pgfqpoint{5.734465in}{1.061926in}}%
\pgfpathlineto{\pgfqpoint{5.724672in}{1.071367in}}%
\pgfpathlineto{\pgfqpoint{5.727231in}{1.078659in}}%
\pgfpathlineto{\pgfqpoint{5.725308in}{1.088305in}}%
\pgfpathlineto{\pgfqpoint{5.737581in}{1.091873in}}%
\pgfpathlineto{\pgfqpoint{5.738281in}{1.100719in}}%
\pgfpathlineto{\pgfqpoint{5.743281in}{1.102373in}}%
\pgfpathlineto{\pgfqpoint{5.740087in}{1.113569in}}%
\pgfpathlineto{\pgfqpoint{5.746406in}{1.129368in}}%
\pgfpathlineto{\pgfqpoint{5.759750in}{1.125564in}}%
\pgfpathlineto{\pgfqpoint{5.758627in}{1.141383in}}%
\pgfpathlineto{\pgfqpoint{5.770610in}{1.157813in}}%
\pgfpathlineto{\pgfqpoint{5.775452in}{1.154663in}}%
\pgfpathlineto{\pgfqpoint{5.781678in}{1.158322in}}%
\pgfpathlineto{\pgfqpoint{5.796468in}{1.175560in}}%
\pgfpathlineto{\pgfqpoint{5.799391in}{1.194887in}}%
\pgfpathlineto{\pgfqpoint{5.799803in}{1.209737in}}%
\pgfpathlineto{\pgfqpoint{5.802915in}{1.219896in}}%
\pgfpathlineto{\pgfqpoint{5.801317in}{1.230712in}}%
\pgfpathlineto{\pgfqpoint{5.797306in}{1.235646in}}%
\pgfpathlineto{\pgfqpoint{5.803519in}{1.241018in}}%
\pgfpathlineto{\pgfqpoint{5.812526in}{1.184275in}}%
\pgfpathlineto{\pgfqpoint{5.862178in}{1.192485in}}%
\pgfpathlineto{\pgfqpoint{5.867325in}{1.160097in}}%
\pgfpathlineto{\pgfqpoint{5.885300in}{1.181474in}}%
\pgfpathlineto{\pgfqpoint{5.889536in}{1.179343in}}%
\pgfpathlineto{\pgfqpoint{5.895767in}{1.191730in}}%
\pgfpathlineto{\pgfqpoint{5.904399in}{1.187846in}}%
\pgfpathlineto{\pgfqpoint{5.911925in}{1.188584in}}%
\pgfpathlineto{\pgfqpoint{5.916891in}{1.197195in}}%
\pgfpathlineto{\pgfqpoint{5.924098in}{1.201985in}}%
\pgfpathlineto{\pgfqpoint{5.934119in}{1.197766in}}%
\pgfpathlineto{\pgfqpoint{5.940712in}{1.199224in}}%
\pgfpathlineto{\pgfqpoint{5.950158in}{1.182877in}}%
\pgfpathlineto{\pgfqpoint{5.947427in}{1.170500in}}%
\pgfpathlineto{\pgfqpoint{5.922755in}{1.184435in}}%
\pgfpathlineto{\pgfqpoint{5.918134in}{1.172996in}}%
\pgfpathlineto{\pgfqpoint{5.919659in}{1.167732in}}%
\pgfpathlineto{\pgfqpoint{5.909159in}{1.149941in}}%
\pgfpathlineto{\pgfqpoint{5.904203in}{1.146402in}}%
\pgfpathlineto{\pgfqpoint{5.901965in}{1.138556in}}%
\pgfpathlineto{\pgfqpoint{5.893517in}{1.139377in}}%
\pgfpathlineto{\pgfqpoint{5.887449in}{1.117963in}}%
\pgfpathlineto{\pgfqpoint{5.884053in}{1.113049in}}%
\pgfpathlineto{\pgfqpoint{5.875326in}{1.114688in}}%
\pgfpathlineto{\pgfqpoint{5.866399in}{1.121442in}}%
\pgfpathlineto{\pgfqpoint{5.866091in}{1.111076in}}%
\pgfpathlineto{\pgfqpoint{5.857449in}{1.093654in}}%
\pgfpathlineto{\pgfqpoint{5.855300in}{1.081251in}}%
\pgfpathlineto{\pgfqpoint{5.848669in}{1.072991in}}%
\pgfpathlineto{\pgfqpoint{5.843643in}{1.059800in}}%
\pgfpathlineto{\pgfqpoint{5.843396in}{1.047589in}}%
\pgfpathlineto{\pgfqpoint{5.835633in}{1.045308in}}%
\pgfpathlineto{\pgfqpoint{5.826826in}{1.038395in}}%
\pgfpathlineto{\pgfqpoint{5.818347in}{1.037079in}}%
\pgfpathlineto{\pgfqpoint{5.816393in}{1.031228in}}%
\pgfpathlineto{\pgfqpoint{5.802730in}{1.025236in}}%
\pgfpathlineto{\pgfqpoint{5.795094in}{1.030341in}}%
\pgfpathlineto{\pgfqpoint{5.786560in}{1.020647in}}%
\pgfpathlineto{\pgfqpoint{5.771890in}{1.023505in}}%
\pgfpathlineto{\pgfqpoint{5.762893in}{1.033679in}}%
\pgfpathlineto{\pgfqpoint{5.763724in}{1.037893in}}%
\pgfpathclose%
\pgfusepath{fill}%
\end{pgfscope}%
\begin{pgfscope}%
\pgfpathrectangle{\pgfqpoint{3.525000in}{0.100000in}}{\pgfqpoint{2.862985in}{1.832778in}}%
\pgfusepath{clip}%
\pgfsetbuttcap%
\pgfsetmiterjoin%
\definecolor{currentfill}{rgb}{0.163014,0.734487,0.163014}%
\pgfsetfillcolor{currentfill}%
\pgfsetlinewidth{0.000000pt}%
\definecolor{currentstroke}{rgb}{0.000000,0.000000,0.000000}%
\pgfsetstrokecolor{currentstroke}%
\pgfsetstrokeopacity{0.000000}%
\pgfsetdash{}{0pt}%
\pgfpathmoveto{\pgfqpoint{6.087135in}{1.155756in}}%
\pgfpathlineto{\pgfqpoint{6.086760in}{1.144322in}}%
\pgfpathlineto{\pgfqpoint{6.082360in}{1.138701in}}%
\pgfpathlineto{\pgfqpoint{6.079551in}{1.126221in}}%
\pgfpathlineto{\pgfqpoint{6.066870in}{1.120469in}}%
\pgfpathlineto{\pgfqpoint{6.056229in}{1.118793in}}%
\pgfpathlineto{\pgfqpoint{6.059330in}{1.127000in}}%
\pgfpathlineto{\pgfqpoint{6.042947in}{1.133877in}}%
\pgfpathlineto{\pgfqpoint{6.029638in}{1.142487in}}%
\pgfpathlineto{\pgfqpoint{6.037878in}{1.155358in}}%
\pgfpathlineto{\pgfqpoint{6.031274in}{1.165363in}}%
\pgfpathlineto{\pgfqpoint{6.032051in}{1.173993in}}%
\pgfpathlineto{\pgfqpoint{6.027276in}{1.176362in}}%
\pgfpathlineto{\pgfqpoint{6.023402in}{1.190382in}}%
\pgfpathlineto{\pgfqpoint{6.032564in}{1.209388in}}%
\pgfpathlineto{\pgfqpoint{6.025681in}{1.212492in}}%
\pgfpathlineto{\pgfqpoint{6.023896in}{1.203115in}}%
\pgfpathlineto{\pgfqpoint{6.014070in}{1.200505in}}%
\pgfpathlineto{\pgfqpoint{6.015103in}{1.169655in}}%
\pgfpathlineto{\pgfqpoint{6.013306in}{1.159726in}}%
\pgfpathlineto{\pgfqpoint{6.018235in}{1.145575in}}%
\pgfpathlineto{\pgfqpoint{6.025818in}{1.138767in}}%
\pgfpathlineto{\pgfqpoint{6.022379in}{1.134492in}}%
\pgfpathlineto{\pgfqpoint{6.030119in}{1.128246in}}%
\pgfpathlineto{\pgfqpoint{6.032917in}{1.118107in}}%
\pgfpathlineto{\pgfqpoint{6.018730in}{1.126503in}}%
\pgfpathlineto{\pgfqpoint{6.009752in}{1.125412in}}%
\pgfpathlineto{\pgfqpoint{6.002780in}{1.134049in}}%
\pgfpathlineto{\pgfqpoint{5.995682in}{1.134890in}}%
\pgfpathlineto{\pgfqpoint{5.985594in}{1.130564in}}%
\pgfpathlineto{\pgfqpoint{5.981681in}{1.135955in}}%
\pgfpathlineto{\pgfqpoint{5.986823in}{1.147270in}}%
\pgfpathlineto{\pgfqpoint{5.989161in}{1.156689in}}%
\pgfpathlineto{\pgfqpoint{5.994331in}{1.164244in}}%
\pgfpathlineto{\pgfqpoint{5.986780in}{1.169343in}}%
\pgfpathlineto{\pgfqpoint{5.983697in}{1.164803in}}%
\pgfpathlineto{\pgfqpoint{5.976154in}{1.169411in}}%
\pgfpathlineto{\pgfqpoint{5.962796in}{1.172482in}}%
\pgfpathlineto{\pgfqpoint{5.964008in}{1.179234in}}%
\pgfpathlineto{\pgfqpoint{5.957953in}{1.183150in}}%
\pgfpathlineto{\pgfqpoint{5.950158in}{1.182877in}}%
\pgfpathlineto{\pgfqpoint{5.940712in}{1.199224in}}%
\pgfpathlineto{\pgfqpoint{5.934119in}{1.197766in}}%
\pgfpathlineto{\pgfqpoint{5.924098in}{1.201985in}}%
\pgfpathlineto{\pgfqpoint{5.916891in}{1.197195in}}%
\pgfpathlineto{\pgfqpoint{5.911925in}{1.188584in}}%
\pgfpathlineto{\pgfqpoint{5.904399in}{1.187846in}}%
\pgfpathlineto{\pgfqpoint{5.895767in}{1.191730in}}%
\pgfpathlineto{\pgfqpoint{5.889536in}{1.179343in}}%
\pgfpathlineto{\pgfqpoint{5.885300in}{1.181474in}}%
\pgfpathlineto{\pgfqpoint{5.867325in}{1.160097in}}%
\pgfpathlineto{\pgfqpoint{5.862178in}{1.192485in}}%
\pgfpathlineto{\pgfqpoint{5.927868in}{1.204659in}}%
\pgfpathlineto{\pgfqpoint{5.957316in}{1.209933in}}%
\pgfpathlineto{\pgfqpoint{6.000158in}{1.218561in}}%
\pgfpathlineto{\pgfqpoint{6.036873in}{1.226505in}}%
\pgfpathlineto{\pgfqpoint{6.042245in}{1.205802in}}%
\pgfpathlineto{\pgfqpoint{6.057891in}{1.150158in}}%
\pgfpathlineto{\pgfqpoint{6.087135in}{1.155756in}}%
\pgfpathclose%
\pgfusepath{fill}%
\end{pgfscope}%
\begin{pgfscope}%
\pgfpathrectangle{\pgfqpoint{3.525000in}{0.100000in}}{\pgfqpoint{2.862985in}{1.832778in}}%
\pgfusepath{clip}%
\pgfsetbuttcap%
\pgfsetmiterjoin%
\definecolor{currentfill}{rgb}{0.192234,0.802845,0.211842}%
\pgfsetfillcolor{currentfill}%
\pgfsetlinewidth{0.000000pt}%
\definecolor{currentstroke}{rgb}{0.000000,0.000000,0.000000}%
\pgfsetstrokecolor{currentstroke}%
\pgfsetstrokeopacity{0.000000}%
\pgfsetdash{}{0pt}%
\pgfpathmoveto{\pgfqpoint{4.764821in}{0.971059in}}%
\pgfpathlineto{\pgfqpoint{4.716895in}{0.975633in}}%
\pgfpathlineto{\pgfqpoint{4.667175in}{0.980039in}}%
\pgfpathlineto{\pgfqpoint{4.609722in}{0.986032in}}%
\pgfpathlineto{\pgfqpoint{4.524362in}{0.996184in}}%
\pgfpathlineto{\pgfqpoint{4.494080in}{1.000850in}}%
\pgfpathlineto{\pgfqpoint{4.415915in}{1.012202in}}%
\pgfpathlineto{\pgfqpoint{4.427328in}{1.083837in}}%
\pgfpathlineto{\pgfqpoint{4.427613in}{1.089632in}}%
\pgfpathlineto{\pgfqpoint{4.438608in}{1.158753in}}%
\pgfpathlineto{\pgfqpoint{4.446820in}{1.211396in}}%
\pgfpathlineto{\pgfqpoint{4.454571in}{1.260214in}}%
\pgfpathlineto{\pgfqpoint{4.507522in}{1.252613in}}%
\pgfpathlineto{\pgfqpoint{4.556892in}{1.245527in}}%
\pgfpathlineto{\pgfqpoint{4.647647in}{1.234326in}}%
\pgfpathlineto{\pgfqpoint{4.689283in}{1.230524in}}%
\pgfpathlineto{\pgfqpoint{4.755243in}{1.224124in}}%
\pgfpathlineto{\pgfqpoint{4.783777in}{1.221808in}}%
\pgfpathlineto{\pgfqpoint{4.778739in}{1.159333in}}%
\pgfpathlineto{\pgfqpoint{4.774188in}{1.099180in}}%
\pgfpathlineto{\pgfqpoint{4.770525in}{1.050201in}}%
\pgfpathlineto{\pgfqpoint{4.764821in}{0.971059in}}%
\pgfpathclose%
\pgfusepath{fill}%
\end{pgfscope}%
\begin{pgfscope}%
\pgfpathrectangle{\pgfqpoint{3.525000in}{0.100000in}}{\pgfqpoint{2.862985in}{1.832778in}}%
\pgfusepath{clip}%
\pgfsetbuttcap%
\pgfsetmiterjoin%
\definecolor{currentfill}{rgb}{0.063053,0.766674,0.741484}%
\pgfsetfillcolor{currentfill}%
\pgfsetlinewidth{0.000000pt}%
\definecolor{currentstroke}{rgb}{0.000000,0.000000,0.000000}%
\pgfsetstrokecolor{currentstroke}%
\pgfsetstrokeopacity{0.000000}%
\pgfsetdash{}{0pt}%
\pgfpathmoveto{\pgfqpoint{5.398847in}{0.935544in}}%
\pgfpathlineto{\pgfqpoint{5.400463in}{0.942786in}}%
\pgfpathlineto{\pgfqpoint{5.408837in}{0.941156in}}%
\pgfpathlineto{\pgfqpoint{5.411225in}{0.966856in}}%
\pgfpathlineto{\pgfqpoint{5.408934in}{0.971495in}}%
\pgfpathlineto{\pgfqpoint{5.417101in}{0.981825in}}%
\pgfpathlineto{\pgfqpoint{5.423652in}{0.981141in}}%
\pgfpathlineto{\pgfqpoint{5.443506in}{0.974346in}}%
\pgfpathlineto{\pgfqpoint{5.445833in}{0.979847in}}%
\pgfpathlineto{\pgfqpoint{5.440777in}{0.988286in}}%
\pgfpathlineto{\pgfqpoint{5.442483in}{0.994944in}}%
\pgfpathlineto{\pgfqpoint{5.460999in}{1.001067in}}%
\pgfpathlineto{\pgfqpoint{5.456580in}{1.012333in}}%
\pgfpathlineto{\pgfqpoint{5.462013in}{1.021868in}}%
\pgfpathlineto{\pgfqpoint{5.472127in}{1.027220in}}%
\pgfpathlineto{\pgfqpoint{5.493349in}{1.032520in}}%
\pgfpathlineto{\pgfqpoint{5.506739in}{1.024520in}}%
\pgfpathlineto{\pgfqpoint{5.510801in}{1.032321in}}%
\pgfpathlineto{\pgfqpoint{5.521919in}{1.037771in}}%
\pgfpathlineto{\pgfqpoint{5.524789in}{1.032965in}}%
\pgfpathlineto{\pgfqpoint{5.536209in}{1.036783in}}%
\pgfpathlineto{\pgfqpoint{5.535489in}{1.043318in}}%
\pgfpathlineto{\pgfqpoint{5.545757in}{1.050797in}}%
\pgfpathlineto{\pgfqpoint{5.551817in}{1.043031in}}%
\pgfpathlineto{\pgfqpoint{5.559790in}{1.042250in}}%
\pgfpathlineto{\pgfqpoint{5.565088in}{1.047311in}}%
\pgfpathlineto{\pgfqpoint{5.564498in}{1.054509in}}%
\pgfpathlineto{\pgfqpoint{5.569002in}{1.061656in}}%
\pgfpathlineto{\pgfqpoint{5.575037in}{1.063205in}}%
\pgfpathlineto{\pgfqpoint{5.577460in}{1.072642in}}%
\pgfpathlineto{\pgfqpoint{5.586237in}{1.080801in}}%
\pgfpathlineto{\pgfqpoint{5.583595in}{1.088922in}}%
\pgfpathlineto{\pgfqpoint{5.592134in}{1.092965in}}%
\pgfpathlineto{\pgfqpoint{5.597839in}{1.090478in}}%
\pgfpathlineto{\pgfqpoint{5.606265in}{1.096795in}}%
\pgfpathlineto{\pgfqpoint{5.613795in}{1.098441in}}%
\pgfpathlineto{\pgfqpoint{5.613822in}{1.104993in}}%
\pgfpathlineto{\pgfqpoint{5.608545in}{1.114085in}}%
\pgfpathlineto{\pgfqpoint{5.611981in}{1.117748in}}%
\pgfpathlineto{\pgfqpoint{5.622773in}{1.116936in}}%
\pgfpathlineto{\pgfqpoint{5.629636in}{1.120355in}}%
\pgfpathlineto{\pgfqpoint{5.636715in}{1.115014in}}%
\pgfpathlineto{\pgfqpoint{5.642303in}{1.102959in}}%
\pgfpathlineto{\pgfqpoint{5.660827in}{1.100922in}}%
\pgfpathlineto{\pgfqpoint{5.670483in}{1.094286in}}%
\pgfpathlineto{\pgfqpoint{5.677277in}{1.099895in}}%
\pgfpathlineto{\pgfqpoint{5.689561in}{1.095181in}}%
\pgfpathlineto{\pgfqpoint{5.696515in}{1.097524in}}%
\pgfpathlineto{\pgfqpoint{5.701376in}{1.104185in}}%
\pgfpathlineto{\pgfqpoint{5.708796in}{1.105459in}}%
\pgfpathlineto{\pgfqpoint{5.710724in}{1.098217in}}%
\pgfpathlineto{\pgfqpoint{5.718930in}{1.095447in}}%
\pgfpathlineto{\pgfqpoint{5.725308in}{1.088305in}}%
\pgfpathlineto{\pgfqpoint{5.727231in}{1.078659in}}%
\pgfpathlineto{\pgfqpoint{5.724672in}{1.071367in}}%
\pgfpathlineto{\pgfqpoint{5.734465in}{1.061926in}}%
\pgfpathlineto{\pgfqpoint{5.745968in}{1.044098in}}%
\pgfpathlineto{\pgfqpoint{5.754483in}{1.038233in}}%
\pgfpathlineto{\pgfqpoint{5.763724in}{1.037893in}}%
\pgfpathlineto{\pgfqpoint{5.746661in}{1.018314in}}%
\pgfpathlineto{\pgfqpoint{5.729865in}{1.006463in}}%
\pgfpathlineto{\pgfqpoint{5.723683in}{0.997053in}}%
\pgfpathlineto{\pgfqpoint{5.723788in}{0.991945in}}%
\pgfpathlineto{\pgfqpoint{5.714691in}{0.988030in}}%
\pgfpathlineto{\pgfqpoint{5.712084in}{0.980665in}}%
\pgfpathlineto{\pgfqpoint{5.693132in}{0.973246in}}%
\pgfpathlineto{\pgfqpoint{5.686414in}{0.968427in}}%
\pgfpathlineto{\pgfqpoint{5.685505in}{0.967399in}}%
\pgfpathlineto{\pgfqpoint{5.630991in}{0.962230in}}%
\pgfpathlineto{\pgfqpoint{5.598070in}{0.959461in}}%
\pgfpathlineto{\pgfqpoint{5.544049in}{0.956404in}}%
\pgfpathlineto{\pgfqpoint{5.476748in}{0.949746in}}%
\pgfpathlineto{\pgfqpoint{5.465629in}{0.951286in}}%
\pgfpathlineto{\pgfqpoint{5.467942in}{0.939936in}}%
\pgfpathlineto{\pgfqpoint{5.398847in}{0.935544in}}%
\pgfpathclose%
\pgfusepath{fill}%
\end{pgfscope}%
\begin{pgfscope}%
\pgfpathrectangle{\pgfqpoint{3.525000in}{0.100000in}}{\pgfqpoint{2.862985in}{1.832778in}}%
\pgfusepath{clip}%
\pgfsetbuttcap%
\pgfsetmiterjoin%
\definecolor{currentfill}{rgb}{0.169166,0.747405,0.169166}%
\pgfsetfillcolor{currentfill}%
\pgfsetlinewidth{0.000000pt}%
\definecolor{currentstroke}{rgb}{0.000000,0.000000,0.000000}%
\pgfsetstrokecolor{currentstroke}%
\pgfsetstrokeopacity{0.000000}%
\pgfsetdash{}{0pt}%
\pgfpathmoveto{\pgfqpoint{4.764821in}{0.971059in}}%
\pgfpathlineto{\pgfqpoint{4.770525in}{1.050201in}}%
\pgfpathlineto{\pgfqpoint{4.774188in}{1.099180in}}%
\pgfpathlineto{\pgfqpoint{4.778739in}{1.159333in}}%
\pgfpathlineto{\pgfqpoint{4.841835in}{1.154903in}}%
\pgfpathlineto{\pgfqpoint{4.921966in}{1.150526in}}%
\pgfpathlineto{\pgfqpoint{4.976466in}{1.148434in}}%
\pgfpathlineto{\pgfqpoint{5.030647in}{1.146774in}}%
\pgfpathlineto{\pgfqpoint{5.079686in}{1.146016in}}%
\pgfpathlineto{\pgfqpoint{5.102367in}{1.146251in}}%
\pgfpathlineto{\pgfqpoint{5.112350in}{1.138104in}}%
\pgfpathlineto{\pgfqpoint{5.120170in}{1.139746in}}%
\pgfpathlineto{\pgfqpoint{5.120414in}{1.132640in}}%
\pgfpathlineto{\pgfqpoint{5.114607in}{1.120348in}}%
\pgfpathlineto{\pgfqpoint{5.115237in}{1.112582in}}%
\pgfpathlineto{\pgfqpoint{5.120797in}{1.107461in}}%
\pgfpathlineto{\pgfqpoint{5.125905in}{1.096790in}}%
\pgfpathlineto{\pgfqpoint{5.136265in}{1.090665in}}%
\pgfpathlineto{\pgfqpoint{5.135919in}{1.050444in}}%
\pgfpathlineto{\pgfqpoint{5.136224in}{0.957954in}}%
\pgfpathlineto{\pgfqpoint{5.066779in}{0.958274in}}%
\pgfpathlineto{\pgfqpoint{4.993651in}{0.959554in}}%
\pgfpathlineto{\pgfqpoint{4.939818in}{0.961453in}}%
\pgfpathlineto{\pgfqpoint{4.894882in}{0.963187in}}%
\pgfpathlineto{\pgfqpoint{4.819179in}{0.967651in}}%
\pgfpathlineto{\pgfqpoint{4.764821in}{0.971059in}}%
\pgfpathclose%
\pgfusepath{fill}%
\end{pgfscope}%
\begin{pgfscope}%
\pgfpathrectangle{\pgfqpoint{3.525000in}{0.100000in}}{\pgfqpoint{2.862985in}{1.832778in}}%
\pgfusepath{clip}%
\pgfsetbuttcap%
\pgfsetmiterjoin%
\definecolor{currentfill}{rgb}{0.176855,0.798539,0.274894}%
\pgfsetfillcolor{currentfill}%
\pgfsetlinewidth{0.000000pt}%
\definecolor{currentstroke}{rgb}{0.000000,0.000000,0.000000}%
\pgfsetstrokecolor{currentstroke}%
\pgfsetstrokeopacity{0.000000}%
\pgfsetdash{}{0pt}%
\pgfpathmoveto{\pgfqpoint{5.786621in}{0.981538in}}%
\pgfpathlineto{\pgfqpoint{5.709707in}{0.970734in}}%
\pgfpathlineto{\pgfqpoint{5.686414in}{0.968427in}}%
\pgfpathlineto{\pgfqpoint{5.693132in}{0.973246in}}%
\pgfpathlineto{\pgfqpoint{5.712084in}{0.980665in}}%
\pgfpathlineto{\pgfqpoint{5.714691in}{0.988030in}}%
\pgfpathlineto{\pgfqpoint{5.723788in}{0.991945in}}%
\pgfpathlineto{\pgfqpoint{5.723683in}{0.997053in}}%
\pgfpathlineto{\pgfqpoint{5.729865in}{1.006463in}}%
\pgfpathlineto{\pgfqpoint{5.746661in}{1.018314in}}%
\pgfpathlineto{\pgfqpoint{5.763724in}{1.037893in}}%
\pgfpathlineto{\pgfqpoint{5.762893in}{1.033679in}}%
\pgfpathlineto{\pgfqpoint{5.771890in}{1.023505in}}%
\pgfpathlineto{\pgfqpoint{5.786560in}{1.020647in}}%
\pgfpathlineto{\pgfqpoint{5.795094in}{1.030341in}}%
\pgfpathlineto{\pgfqpoint{5.802730in}{1.025236in}}%
\pgfpathlineto{\pgfqpoint{5.816393in}{1.031228in}}%
\pgfpathlineto{\pgfqpoint{5.818347in}{1.037079in}}%
\pgfpathlineto{\pgfqpoint{5.826826in}{1.038395in}}%
\pgfpathlineto{\pgfqpoint{5.835633in}{1.045308in}}%
\pgfpathlineto{\pgfqpoint{5.843396in}{1.047589in}}%
\pgfpathlineto{\pgfqpoint{5.843643in}{1.059800in}}%
\pgfpathlineto{\pgfqpoint{5.848669in}{1.072991in}}%
\pgfpathlineto{\pgfqpoint{5.855300in}{1.081251in}}%
\pgfpathlineto{\pgfqpoint{5.857449in}{1.093654in}}%
\pgfpathlineto{\pgfqpoint{5.866091in}{1.111076in}}%
\pgfpathlineto{\pgfqpoint{5.866399in}{1.121442in}}%
\pgfpathlineto{\pgfqpoint{5.875326in}{1.114688in}}%
\pgfpathlineto{\pgfqpoint{5.884053in}{1.113049in}}%
\pgfpathlineto{\pgfqpoint{5.887449in}{1.117963in}}%
\pgfpathlineto{\pgfqpoint{5.893517in}{1.139377in}}%
\pgfpathlineto{\pgfqpoint{5.901965in}{1.138556in}}%
\pgfpathlineto{\pgfqpoint{5.904203in}{1.146402in}}%
\pgfpathlineto{\pgfqpoint{5.909159in}{1.149941in}}%
\pgfpathlineto{\pgfqpoint{5.919659in}{1.167732in}}%
\pgfpathlineto{\pgfqpoint{5.918134in}{1.172996in}}%
\pgfpathlineto{\pgfqpoint{5.922755in}{1.184435in}}%
\pgfpathlineto{\pgfqpoint{5.947427in}{1.170500in}}%
\pgfpathlineto{\pgfqpoint{5.950158in}{1.182877in}}%
\pgfpathlineto{\pgfqpoint{5.957953in}{1.183150in}}%
\pgfpathlineto{\pgfqpoint{5.964008in}{1.179234in}}%
\pgfpathlineto{\pgfqpoint{5.962796in}{1.172482in}}%
\pgfpathlineto{\pgfqpoint{5.976154in}{1.169411in}}%
\pgfpathlineto{\pgfqpoint{5.983697in}{1.164803in}}%
\pgfpathlineto{\pgfqpoint{5.989161in}{1.156689in}}%
\pgfpathlineto{\pgfqpoint{5.986823in}{1.147270in}}%
\pgfpathlineto{\pgfqpoint{5.982103in}{1.146498in}}%
\pgfpathlineto{\pgfqpoint{5.979365in}{1.132285in}}%
\pgfpathlineto{\pgfqpoint{5.985362in}{1.126728in}}%
\pgfpathlineto{\pgfqpoint{5.993804in}{1.131220in}}%
\pgfpathlineto{\pgfqpoint{6.001637in}{1.121734in}}%
\pgfpathlineto{\pgfqpoint{6.019137in}{1.120029in}}%
\pgfpathlineto{\pgfqpoint{6.022152in}{1.114573in}}%
\pgfpathlineto{\pgfqpoint{6.038348in}{1.109266in}}%
\pgfpathlineto{\pgfqpoint{6.036349in}{1.103004in}}%
\pgfpathlineto{\pgfqpoint{6.038456in}{1.085305in}}%
\pgfpathlineto{\pgfqpoint{6.044993in}{1.078622in}}%
\pgfpathlineto{\pgfqpoint{6.035347in}{1.075802in}}%
\pgfpathlineto{\pgfqpoint{6.036626in}{1.068246in}}%
\pgfpathlineto{\pgfqpoint{6.046916in}{1.061840in}}%
\pgfpathlineto{\pgfqpoint{6.047844in}{1.055521in}}%
\pgfpathlineto{\pgfqpoint{6.042034in}{1.050772in}}%
\pgfpathlineto{\pgfqpoint{6.030312in}{1.062022in}}%
\pgfpathlineto{\pgfqpoint{6.029153in}{1.053814in}}%
\pgfpathlineto{\pgfqpoint{6.038966in}{1.049911in}}%
\pgfpathlineto{\pgfqpoint{6.039948in}{1.045875in}}%
\pgfpathlineto{\pgfqpoint{6.053408in}{1.051209in}}%
\pgfpathlineto{\pgfqpoint{6.063722in}{1.052620in}}%
\pgfpathlineto{\pgfqpoint{6.074287in}{1.031303in}}%
\pgfpathlineto{\pgfqpoint{6.073108in}{1.031073in}}%
\pgfpathlineto{\pgfqpoint{6.068334in}{1.030086in}}%
\pgfpathlineto{\pgfqpoint{6.066927in}{1.029792in}}%
\pgfpathlineto{\pgfqpoint{6.065997in}{1.029610in}}%
\pgfpathlineto{\pgfqpoint{6.003091in}{1.016561in}}%
\pgfpathlineto{\pgfqpoint{5.946826in}{1.005045in}}%
\pgfpathlineto{\pgfqpoint{5.869033in}{0.991644in}}%
\pgfpathlineto{\pgfqpoint{5.802964in}{0.982918in}}%
\pgfpathlineto{\pgfqpoint{5.786621in}{0.981538in}}%
\pgfpathclose%
\pgfusepath{fill}%
\end{pgfscope}%
\begin{pgfscope}%
\pgfpathrectangle{\pgfqpoint{3.525000in}{0.100000in}}{\pgfqpoint{2.862985in}{1.832778in}}%
\pgfusepath{clip}%
\pgfsetbuttcap%
\pgfsetmiterjoin%
\definecolor{currentfill}{rgb}{0.176855,0.798539,0.274894}%
\pgfsetfillcolor{currentfill}%
\pgfsetlinewidth{0.000000pt}%
\definecolor{currentstroke}{rgb}{0.000000,0.000000,0.000000}%
\pgfsetstrokecolor{currentstroke}%
\pgfsetstrokeopacity{0.000000}%
\pgfsetdash{}{0pt}%
\pgfpathmoveto{\pgfqpoint{6.066870in}{1.120469in}}%
\pgfpathlineto{\pgfqpoint{6.079551in}{1.126221in}}%
\pgfpathlineto{\pgfqpoint{6.069424in}{1.096596in}}%
\pgfpathlineto{\pgfqpoint{6.062727in}{1.080908in}}%
\pgfpathlineto{\pgfqpoint{6.057902in}{1.089786in}}%
\pgfpathlineto{\pgfqpoint{6.066482in}{1.111041in}}%
\pgfpathlineto{\pgfqpoint{6.066870in}{1.120469in}}%
\pgfpathclose%
\pgfusepath{fill}%
\end{pgfscope}%
\begin{pgfscope}%
\pgfpathrectangle{\pgfqpoint{3.525000in}{0.100000in}}{\pgfqpoint{2.862985in}{1.832778in}}%
\pgfusepath{clip}%
\pgfsetbuttcap%
\pgfsetmiterjoin%
\definecolor{currentfill}{rgb}{0.155325,0.792511,0.363168}%
\pgfsetfillcolor{currentfill}%
\pgfsetlinewidth{0.000000pt}%
\definecolor{currentstroke}{rgb}{0.000000,0.000000,0.000000}%
\pgfsetstrokecolor{currentstroke}%
\pgfsetstrokeopacity{0.000000}%
\pgfsetdash{}{0pt}%
\pgfpathmoveto{\pgfqpoint{5.398847in}{0.935544in}}%
\pgfpathlineto{\pgfqpoint{5.395785in}{0.935101in}}%
\pgfpathlineto{\pgfqpoint{5.392896in}{0.934898in}}%
\pgfpathlineto{\pgfqpoint{5.394133in}{0.926031in}}%
\pgfpathlineto{\pgfqpoint{5.386682in}{0.917093in}}%
\pgfpathlineto{\pgfqpoint{5.385213in}{0.903107in}}%
\pgfpathlineto{\pgfqpoint{5.351904in}{0.900644in}}%
\pgfpathlineto{\pgfqpoint{5.354808in}{0.907211in}}%
\pgfpathlineto{\pgfqpoint{5.366818in}{0.919204in}}%
\pgfpathlineto{\pgfqpoint{5.367150in}{0.926155in}}%
\pgfpathlineto{\pgfqpoint{5.361833in}{0.932737in}}%
\pgfpathlineto{\pgfqpoint{5.312257in}{0.930116in}}%
\pgfpathlineto{\pgfqpoint{5.225580in}{0.927364in}}%
\pgfpathlineto{\pgfqpoint{5.174856in}{0.926440in}}%
\pgfpathlineto{\pgfqpoint{5.136515in}{0.926094in}}%
\pgfpathlineto{\pgfqpoint{5.136224in}{0.957954in}}%
\pgfpathlineto{\pgfqpoint{5.135919in}{1.050444in}}%
\pgfpathlineto{\pgfqpoint{5.136265in}{1.090665in}}%
\pgfpathlineto{\pgfqpoint{5.125905in}{1.096790in}}%
\pgfpathlineto{\pgfqpoint{5.120797in}{1.107461in}}%
\pgfpathlineto{\pgfqpoint{5.115237in}{1.112582in}}%
\pgfpathlineto{\pgfqpoint{5.114607in}{1.120348in}}%
\pgfpathlineto{\pgfqpoint{5.120414in}{1.132640in}}%
\pgfpathlineto{\pgfqpoint{5.120170in}{1.139746in}}%
\pgfpathlineto{\pgfqpoint{5.112350in}{1.138104in}}%
\pgfpathlineto{\pgfqpoint{5.102367in}{1.146251in}}%
\pgfpathlineto{\pgfqpoint{5.094364in}{1.160551in}}%
\pgfpathlineto{\pgfqpoint{5.087651in}{1.167151in}}%
\pgfpathlineto{\pgfqpoint{5.080637in}{1.183365in}}%
\pgfpathlineto{\pgfqpoint{5.153489in}{1.182229in}}%
\pgfpathlineto{\pgfqpoint{5.225906in}{1.184615in}}%
\pgfpathlineto{\pgfqpoint{5.272339in}{1.187293in}}%
\pgfpathlineto{\pgfqpoint{5.286865in}{1.172908in}}%
\pgfpathlineto{\pgfqpoint{5.284210in}{1.157906in}}%
\pgfpathlineto{\pgfqpoint{5.288083in}{1.141015in}}%
\pgfpathlineto{\pgfqpoint{5.292379in}{1.132498in}}%
\pgfpathlineto{\pgfqpoint{5.306087in}{1.120782in}}%
\pgfpathlineto{\pgfqpoint{5.309086in}{1.115607in}}%
\pgfpathlineto{\pgfqpoint{5.324094in}{1.103945in}}%
\pgfpathlineto{\pgfqpoint{5.326269in}{1.090863in}}%
\pgfpathlineto{\pgfqpoint{5.330579in}{1.081212in}}%
\pgfpathlineto{\pgfqpoint{5.340763in}{1.086834in}}%
\pgfpathlineto{\pgfqpoint{5.349157in}{1.084263in}}%
\pgfpathlineto{\pgfqpoint{5.355495in}{1.077440in}}%
\pgfpathlineto{\pgfqpoint{5.353102in}{1.065365in}}%
\pgfpathlineto{\pgfqpoint{5.345322in}{1.049578in}}%
\pgfpathlineto{\pgfqpoint{5.345442in}{1.041355in}}%
\pgfpathlineto{\pgfqpoint{5.351256in}{1.034614in}}%
\pgfpathlineto{\pgfqpoint{5.363768in}{1.025684in}}%
\pgfpathlineto{\pgfqpoint{5.378174in}{1.018436in}}%
\pgfpathlineto{\pgfqpoint{5.381591in}{1.012632in}}%
\pgfpathlineto{\pgfqpoint{5.388959in}{1.009888in}}%
\pgfpathlineto{\pgfqpoint{5.389252in}{1.002186in}}%
\pgfpathlineto{\pgfqpoint{5.394705in}{0.992402in}}%
\pgfpathlineto{\pgfqpoint{5.390922in}{0.983711in}}%
\pgfpathlineto{\pgfqpoint{5.398328in}{0.969882in}}%
\pgfpathlineto{\pgfqpoint{5.404121in}{0.972628in}}%
\pgfpathlineto{\pgfqpoint{5.411225in}{0.966856in}}%
\pgfpathlineto{\pgfqpoint{5.408837in}{0.941156in}}%
\pgfpathlineto{\pgfqpoint{5.400463in}{0.942786in}}%
\pgfpathlineto{\pgfqpoint{5.398847in}{0.935544in}}%
\pgfpathclose%
\pgfusepath{fill}%
\end{pgfscope}%
\begin{pgfscope}%
\pgfpathrectangle{\pgfqpoint{3.525000in}{0.100000in}}{\pgfqpoint{2.862985in}{1.832778in}}%
\pgfusepath{clip}%
\pgfsetbuttcap%
\pgfsetmiterjoin%
\definecolor{currentfill}{rgb}{0.099962,0.777009,0.590158}%
\pgfsetfillcolor{currentfill}%
\pgfsetlinewidth{0.000000pt}%
\definecolor{currentstroke}{rgb}{0.000000,0.000000,0.000000}%
\pgfsetstrokecolor{currentstroke}%
\pgfsetstrokeopacity{0.000000}%
\pgfsetdash{}{0pt}%
\pgfpathmoveto{\pgfqpoint{4.114569in}{0.941085in}}%
\pgfpathlineto{\pgfqpoint{4.120633in}{0.953995in}}%
\pgfpathlineto{\pgfqpoint{4.118973in}{0.973374in}}%
\pgfpathlineto{\pgfqpoint{4.121877in}{0.978886in}}%
\pgfpathlineto{\pgfqpoint{4.122164in}{1.003024in}}%
\pgfpathlineto{\pgfqpoint{4.124916in}{1.009976in}}%
\pgfpathlineto{\pgfqpoint{4.134613in}{1.011057in}}%
\pgfpathlineto{\pgfqpoint{4.143612in}{1.007971in}}%
\pgfpathlineto{\pgfqpoint{4.147531in}{0.999473in}}%
\pgfpathlineto{\pgfqpoint{4.153025in}{0.999798in}}%
\pgfpathlineto{\pgfqpoint{4.159852in}{1.009534in}}%
\pgfpathlineto{\pgfqpoint{4.169719in}{1.057523in}}%
\pgfpathlineto{\pgfqpoint{4.225868in}{1.045930in}}%
\pgfpathlineto{\pgfqpoint{4.258455in}{1.039504in}}%
\pgfpathlineto{\pgfqpoint{4.344991in}{1.024292in}}%
\pgfpathlineto{\pgfqpoint{4.368932in}{1.019484in}}%
\pgfpathlineto{\pgfqpoint{4.415915in}{1.012202in}}%
\pgfpathlineto{\pgfqpoint{4.406281in}{0.950172in}}%
\pgfpathlineto{\pgfqpoint{4.396262in}{0.885469in}}%
\pgfpathlineto{\pgfqpoint{4.384723in}{0.812676in}}%
\pgfpathlineto{\pgfqpoint{4.371747in}{0.729118in}}%
\pgfpathlineto{\pgfqpoint{4.361264in}{0.660490in}}%
\pgfpathlineto{\pgfqpoint{4.286176in}{0.672414in}}%
\pgfpathlineto{\pgfqpoint{4.253185in}{0.678064in}}%
\pgfpathlineto{\pgfqpoint{4.238447in}{0.686887in}}%
\pgfpathlineto{\pgfqpoint{4.142356in}{0.744864in}}%
\pgfpathlineto{\pgfqpoint{4.070244in}{0.788862in}}%
\pgfpathlineto{\pgfqpoint{4.072646in}{0.796660in}}%
\pgfpathlineto{\pgfqpoint{4.078603in}{0.802117in}}%
\pgfpathlineto{\pgfqpoint{4.084798in}{0.801105in}}%
\pgfpathlineto{\pgfqpoint{4.093785in}{0.806836in}}%
\pgfpathlineto{\pgfqpoint{4.095200in}{0.815071in}}%
\pgfpathlineto{\pgfqpoint{4.084240in}{0.825061in}}%
\pgfpathlineto{\pgfqpoint{4.092035in}{0.844239in}}%
\pgfpathlineto{\pgfqpoint{4.099871in}{0.851614in}}%
\pgfpathlineto{\pgfqpoint{4.103591in}{0.860364in}}%
\pgfpathlineto{\pgfqpoint{4.105882in}{0.876407in}}%
\pgfpathlineto{\pgfqpoint{4.113230in}{0.883673in}}%
\pgfpathlineto{\pgfqpoint{4.128654in}{0.890930in}}%
\pgfpathlineto{\pgfqpoint{4.128313in}{0.897324in}}%
\pgfpathlineto{\pgfqpoint{4.119699in}{0.905259in}}%
\pgfpathlineto{\pgfqpoint{4.118526in}{0.921617in}}%
\pgfpathlineto{\pgfqpoint{4.112595in}{0.933576in}}%
\pgfpathlineto{\pgfqpoint{4.114569in}{0.941085in}}%
\pgfpathclose%
\pgfusepath{fill}%
\end{pgfscope}%
\begin{pgfscope}%
\pgfpathrectangle{\pgfqpoint{3.525000in}{0.100000in}}{\pgfqpoint{2.862985in}{1.832778in}}%
\pgfusepath{clip}%
\pgfsetbuttcap%
\pgfsetmiterjoin%
\definecolor{currentfill}{rgb}{0.016240,0.752065,0.997693}%
\pgfsetfillcolor{currentfill}%
\pgfsetlinewidth{0.000000pt}%
\definecolor{currentstroke}{rgb}{0.000000,0.000000,0.000000}%
\pgfsetstrokecolor{currentstroke}%
\pgfsetstrokeopacity{0.000000}%
\pgfsetdash{}{0pt}%
\pgfpathmoveto{\pgfqpoint{5.144643in}{0.746495in}}%
\pgfpathlineto{\pgfqpoint{5.130724in}{0.750800in}}%
\pgfpathlineto{\pgfqpoint{5.112888in}{0.764459in}}%
\pgfpathlineto{\pgfqpoint{5.104966in}{0.767394in}}%
\pgfpathlineto{\pgfqpoint{5.099932in}{0.761495in}}%
\pgfpathlineto{\pgfqpoint{5.093575in}{0.761197in}}%
\pgfpathlineto{\pgfqpoint{5.085532in}{0.766205in}}%
\pgfpathlineto{\pgfqpoint{5.072911in}{0.759788in}}%
\pgfpathlineto{\pgfqpoint{5.068418in}{0.762948in}}%
\pgfpathlineto{\pgfqpoint{5.055999in}{0.755462in}}%
\pgfpathlineto{\pgfqpoint{5.042897in}{0.756843in}}%
\pgfpathlineto{\pgfqpoint{5.032791in}{0.761708in}}%
\pgfpathlineto{\pgfqpoint{5.025714in}{0.759877in}}%
\pgfpathlineto{\pgfqpoint{5.014386in}{0.766758in}}%
\pgfpathlineto{\pgfqpoint{5.008084in}{0.754622in}}%
\pgfpathlineto{\pgfqpoint{5.001640in}{0.765059in}}%
\pgfpathlineto{\pgfqpoint{4.988911in}{0.761010in}}%
\pgfpathlineto{\pgfqpoint{4.977778in}{0.770922in}}%
\pgfpathlineto{\pgfqpoint{4.968049in}{0.762933in}}%
\pgfpathlineto{\pgfqpoint{4.963177in}{0.770275in}}%
\pgfpathlineto{\pgfqpoint{4.956154in}{0.772659in}}%
\pgfpathlineto{\pgfqpoint{4.951877in}{0.779738in}}%
\pgfpathlineto{\pgfqpoint{4.942692in}{0.781753in}}%
\pgfpathlineto{\pgfqpoint{4.937392in}{0.776425in}}%
\pgfpathlineto{\pgfqpoint{4.928387in}{0.783330in}}%
\pgfpathlineto{\pgfqpoint{4.924189in}{0.781753in}}%
\pgfpathlineto{\pgfqpoint{4.909265in}{0.787349in}}%
\pgfpathlineto{\pgfqpoint{4.899917in}{0.787973in}}%
\pgfpathlineto{\pgfqpoint{4.895733in}{0.799856in}}%
\pgfpathlineto{\pgfqpoint{4.888240in}{0.798376in}}%
\pgfpathlineto{\pgfqpoint{4.879657in}{0.801308in}}%
\pgfpathlineto{\pgfqpoint{4.874005in}{0.799614in}}%
\pgfpathlineto{\pgfqpoint{4.861884in}{0.812959in}}%
\pgfpathlineto{\pgfqpoint{4.858507in}{0.812086in}}%
\pgfpathlineto{\pgfqpoint{4.861574in}{0.866207in}}%
\pgfpathlineto{\pgfqpoint{4.864916in}{0.933195in}}%
\pgfpathlineto{\pgfqpoint{4.810064in}{0.936263in}}%
\pgfpathlineto{\pgfqpoint{4.755936in}{0.940356in}}%
\pgfpathlineto{\pgfqpoint{4.714118in}{0.943987in}}%
\pgfpathlineto{\pgfqpoint{4.716895in}{0.975633in}}%
\pgfpathlineto{\pgfqpoint{4.764821in}{0.971059in}}%
\pgfpathlineto{\pgfqpoint{4.819179in}{0.967651in}}%
\pgfpathlineto{\pgfqpoint{4.894882in}{0.963187in}}%
\pgfpathlineto{\pgfqpoint{4.939818in}{0.961453in}}%
\pgfpathlineto{\pgfqpoint{4.993651in}{0.959554in}}%
\pgfpathlineto{\pgfqpoint{5.066779in}{0.958274in}}%
\pgfpathlineto{\pgfqpoint{5.136224in}{0.957954in}}%
\pgfpathlineto{\pgfqpoint{5.136515in}{0.926094in}}%
\pgfpathlineto{\pgfqpoint{5.140413in}{0.902092in}}%
\pgfpathlineto{\pgfqpoint{5.146468in}{0.857760in}}%
\pgfpathlineto{\pgfqpoint{5.145219in}{0.782052in}}%
\pgfpathlineto{\pgfqpoint{5.144643in}{0.746495in}}%
\pgfpathclose%
\pgfusepath{fill}%
\end{pgfscope}%
\begin{pgfscope}%
\pgfpathrectangle{\pgfqpoint{3.525000in}{0.100000in}}{\pgfqpoint{2.862985in}{1.832778in}}%
\pgfusepath{clip}%
\pgfsetbuttcap%
\pgfsetmiterjoin%
\definecolor{currentfill}{rgb}{0.099962,0.777009,0.590158}%
\pgfsetfillcolor{currentfill}%
\pgfsetlinewidth{0.000000pt}%
\definecolor{currentstroke}{rgb}{0.000000,0.000000,0.000000}%
\pgfsetstrokecolor{currentstroke}%
\pgfsetstrokeopacity{0.000000}%
\pgfsetdash{}{0pt}%
\pgfpathmoveto{\pgfqpoint{5.665820in}{0.864070in}}%
\pgfpathlineto{\pgfqpoint{5.665861in}{0.878097in}}%
\pgfpathlineto{\pgfqpoint{5.678051in}{0.883484in}}%
\pgfpathlineto{\pgfqpoint{5.678562in}{0.892095in}}%
\pgfpathlineto{\pgfqpoint{5.689476in}{0.902614in}}%
\pgfpathlineto{\pgfqpoint{5.703119in}{0.904768in}}%
\pgfpathlineto{\pgfqpoint{5.715113in}{0.916361in}}%
\pgfpathlineto{\pgfqpoint{5.727639in}{0.921543in}}%
\pgfpathlineto{\pgfqpoint{5.737040in}{0.937098in}}%
\pgfpathlineto{\pgfqpoint{5.745602in}{0.935969in}}%
\pgfpathlineto{\pgfqpoint{5.763697in}{0.950142in}}%
\pgfpathlineto{\pgfqpoint{5.773253in}{0.950409in}}%
\pgfpathlineto{\pgfqpoint{5.777288in}{0.961208in}}%
\pgfpathlineto{\pgfqpoint{5.784883in}{0.968727in}}%
\pgfpathlineto{\pgfqpoint{5.786621in}{0.981538in}}%
\pgfpathlineto{\pgfqpoint{5.802964in}{0.982918in}}%
\pgfpathlineto{\pgfqpoint{5.869033in}{0.991644in}}%
\pgfpathlineto{\pgfqpoint{5.946826in}{1.005045in}}%
\pgfpathlineto{\pgfqpoint{6.003091in}{1.016561in}}%
\pgfpathlineto{\pgfqpoint{6.065997in}{1.029610in}}%
\pgfpathlineto{\pgfqpoint{6.083973in}{1.004931in}}%
\pgfpathlineto{\pgfqpoint{6.074236in}{1.006506in}}%
\pgfpathlineto{\pgfqpoint{6.061737in}{1.003461in}}%
\pgfpathlineto{\pgfqpoint{6.049421in}{0.990884in}}%
\pgfpathlineto{\pgfqpoint{6.040570in}{0.991789in}}%
\pgfpathlineto{\pgfqpoint{6.039450in}{0.984319in}}%
\pgfpathlineto{\pgfqpoint{6.055449in}{0.990226in}}%
\pgfpathlineto{\pgfqpoint{6.071548in}{0.992662in}}%
\pgfpathlineto{\pgfqpoint{6.077526in}{0.989421in}}%
\pgfpathlineto{\pgfqpoint{6.085571in}{0.993116in}}%
\pgfpathlineto{\pgfqpoint{6.089722in}{0.990527in}}%
\pgfpathlineto{\pgfqpoint{6.093396in}{0.978199in}}%
\pgfpathlineto{\pgfqpoint{6.085750in}{0.974382in}}%
\pgfpathlineto{\pgfqpoint{6.080502in}{0.959350in}}%
\pgfpathlineto{\pgfqpoint{6.075000in}{0.953497in}}%
\pgfpathlineto{\pgfqpoint{6.058132in}{0.954777in}}%
\pgfpathlineto{\pgfqpoint{6.059039in}{0.963604in}}%
\pgfpathlineto{\pgfqpoint{6.049827in}{0.959748in}}%
\pgfpathlineto{\pgfqpoint{6.047823in}{0.952387in}}%
\pgfpathlineto{\pgfqpoint{6.054137in}{0.944761in}}%
\pgfpathlineto{\pgfqpoint{6.056282in}{0.931819in}}%
\pgfpathlineto{\pgfqpoint{6.045950in}{0.924444in}}%
\pgfpathlineto{\pgfqpoint{6.057125in}{0.921851in}}%
\pgfpathlineto{\pgfqpoint{6.062181in}{0.927273in}}%
\pgfpathlineto{\pgfqpoint{6.073362in}{0.927934in}}%
\pgfpathlineto{\pgfqpoint{6.067616in}{0.916224in}}%
\pgfpathlineto{\pgfqpoint{6.060579in}{0.910582in}}%
\pgfpathlineto{\pgfqpoint{6.039292in}{0.904933in}}%
\pgfpathlineto{\pgfqpoint{6.017436in}{0.885111in}}%
\pgfpathlineto{\pgfqpoint{6.003960in}{0.865580in}}%
\pgfpathlineto{\pgfqpoint{6.003880in}{0.857648in}}%
\pgfpathlineto{\pgfqpoint{5.998497in}{0.846702in}}%
\pgfpathlineto{\pgfqpoint{5.970869in}{0.839501in}}%
\pgfpathlineto{\pgfqpoint{5.905165in}{0.886600in}}%
\pgfpathlineto{\pgfqpoint{5.847304in}{0.878098in}}%
\pgfpathlineto{\pgfqpoint{5.846819in}{0.885949in}}%
\pgfpathlineto{\pgfqpoint{5.838024in}{0.894789in}}%
\pgfpathlineto{\pgfqpoint{5.831366in}{0.896952in}}%
\pgfpathlineto{\pgfqpoint{5.768499in}{0.890420in}}%
\pgfpathlineto{\pgfqpoint{5.754049in}{0.885519in}}%
\pgfpathlineto{\pgfqpoint{5.727958in}{0.872540in}}%
\pgfpathlineto{\pgfqpoint{5.705410in}{0.868948in}}%
\pgfpathlineto{\pgfqpoint{5.665820in}{0.864070in}}%
\pgfpathclose%
\pgfusepath{fill}%
\end{pgfscope}%
\begin{pgfscope}%
\pgfpathrectangle{\pgfqpoint{3.525000in}{0.100000in}}{\pgfqpoint{2.862985in}{1.832778in}}%
\pgfusepath{clip}%
\pgfsetbuttcap%
\pgfsetmiterjoin%
\definecolor{currentfill}{rgb}{0.029220,0.757201,0.880200}%
\pgfsetfillcolor{currentfill}%
\pgfsetlinewidth{0.000000pt}%
\definecolor{currentstroke}{rgb}{0.000000,0.000000,0.000000}%
\pgfsetstrokecolor{currentstroke}%
\pgfsetstrokeopacity{0.000000}%
\pgfsetdash{}{0pt}%
\pgfpathmoveto{\pgfqpoint{5.665820in}{0.864070in}}%
\pgfpathlineto{\pgfqpoint{5.600022in}{0.856843in}}%
\pgfpathlineto{\pgfqpoint{5.539830in}{0.851430in}}%
\pgfpathlineto{\pgfqpoint{5.477965in}{0.847429in}}%
\pgfpathlineto{\pgfqpoint{5.467346in}{0.845890in}}%
\pgfpathlineto{\pgfqpoint{5.425643in}{0.842683in}}%
\pgfpathlineto{\pgfqpoint{5.358848in}{0.838785in}}%
\pgfpathlineto{\pgfqpoint{5.370732in}{0.848649in}}%
\pgfpathlineto{\pgfqpoint{5.368020in}{0.859033in}}%
\pgfpathlineto{\pgfqpoint{5.370647in}{0.871619in}}%
\pgfpathlineto{\pgfqpoint{5.374660in}{0.877545in}}%
\pgfpathlineto{\pgfqpoint{5.374509in}{0.885789in}}%
\pgfpathlineto{\pgfqpoint{5.385195in}{0.890974in}}%
\pgfpathlineto{\pgfqpoint{5.385213in}{0.903107in}}%
\pgfpathlineto{\pgfqpoint{5.386682in}{0.917093in}}%
\pgfpathlineto{\pgfqpoint{5.394133in}{0.926031in}}%
\pgfpathlineto{\pgfqpoint{5.392896in}{0.934898in}}%
\pgfpathlineto{\pgfqpoint{5.395785in}{0.935101in}}%
\pgfpathlineto{\pgfqpoint{5.398847in}{0.935544in}}%
\pgfpathlineto{\pgfqpoint{5.467942in}{0.939936in}}%
\pgfpathlineto{\pgfqpoint{5.465629in}{0.951286in}}%
\pgfpathlineto{\pgfqpoint{5.476748in}{0.949746in}}%
\pgfpathlineto{\pgfqpoint{5.544049in}{0.956404in}}%
\pgfpathlineto{\pgfqpoint{5.598070in}{0.959461in}}%
\pgfpathlineto{\pgfqpoint{5.630991in}{0.962230in}}%
\pgfpathlineto{\pgfqpoint{5.685505in}{0.967399in}}%
\pgfpathlineto{\pgfqpoint{5.686414in}{0.968427in}}%
\pgfpathlineto{\pgfqpoint{5.709707in}{0.970734in}}%
\pgfpathlineto{\pgfqpoint{5.786621in}{0.981538in}}%
\pgfpathlineto{\pgfqpoint{5.784883in}{0.968727in}}%
\pgfpathlineto{\pgfqpoint{5.777288in}{0.961208in}}%
\pgfpathlineto{\pgfqpoint{5.773253in}{0.950409in}}%
\pgfpathlineto{\pgfqpoint{5.763697in}{0.950142in}}%
\pgfpathlineto{\pgfqpoint{5.745602in}{0.935969in}}%
\pgfpathlineto{\pgfqpoint{5.737040in}{0.937098in}}%
\pgfpathlineto{\pgfqpoint{5.727639in}{0.921543in}}%
\pgfpathlineto{\pgfqpoint{5.715113in}{0.916361in}}%
\pgfpathlineto{\pgfqpoint{5.703119in}{0.904768in}}%
\pgfpathlineto{\pgfqpoint{5.689476in}{0.902614in}}%
\pgfpathlineto{\pgfqpoint{5.678562in}{0.892095in}}%
\pgfpathlineto{\pgfqpoint{5.678051in}{0.883484in}}%
\pgfpathlineto{\pgfqpoint{5.665861in}{0.878097in}}%
\pgfpathlineto{\pgfqpoint{5.665820in}{0.864070in}}%
\pgfpathclose%
\pgfusepath{fill}%
\end{pgfscope}%
\begin{pgfscope}%
\pgfpathrectangle{\pgfqpoint{3.525000in}{0.100000in}}{\pgfqpoint{2.862985in}{1.832778in}}%
\pgfusepath{clip}%
\pgfsetbuttcap%
\pgfsetmiterjoin%
\definecolor{currentfill}{rgb}{0.087659,0.773564,0.640600}%
\pgfsetfillcolor{currentfill}%
\pgfsetlinewidth{0.000000pt}%
\definecolor{currentstroke}{rgb}{0.000000,0.000000,0.000000}%
\pgfsetstrokecolor{currentstroke}%
\pgfsetstrokeopacity{0.000000}%
\pgfsetdash{}{0pt}%
\pgfpathmoveto{\pgfqpoint{4.498883in}{0.669271in}}%
\pgfpathlineto{\pgfqpoint{4.494190in}{0.676811in}}%
\pgfpathlineto{\pgfqpoint{4.496132in}{0.683307in}}%
\pgfpathlineto{\pgfqpoint{4.529118in}{0.679215in}}%
\pgfpathlineto{\pgfqpoint{4.590452in}{0.672249in}}%
\pgfpathlineto{\pgfqpoint{4.634801in}{0.667891in}}%
\pgfpathlineto{\pgfqpoint{4.686021in}{0.662722in}}%
\pgfpathlineto{\pgfqpoint{4.688807in}{0.695029in}}%
\pgfpathlineto{\pgfqpoint{4.697520in}{0.777201in}}%
\pgfpathlineto{\pgfqpoint{4.703172in}{0.834767in}}%
\pgfpathlineto{\pgfqpoint{4.708057in}{0.889756in}}%
\pgfpathlineto{\pgfqpoint{4.712612in}{0.944071in}}%
\pgfpathlineto{\pgfqpoint{4.714118in}{0.943987in}}%
\pgfpathlineto{\pgfqpoint{4.755936in}{0.940356in}}%
\pgfpathlineto{\pgfqpoint{4.810064in}{0.936263in}}%
\pgfpathlineto{\pgfqpoint{4.864916in}{0.933195in}}%
\pgfpathlineto{\pgfqpoint{4.861574in}{0.866207in}}%
\pgfpathlineto{\pgfqpoint{4.858507in}{0.812086in}}%
\pgfpathlineto{\pgfqpoint{4.861884in}{0.812959in}}%
\pgfpathlineto{\pgfqpoint{4.874005in}{0.799614in}}%
\pgfpathlineto{\pgfqpoint{4.879657in}{0.801308in}}%
\pgfpathlineto{\pgfqpoint{4.888240in}{0.798376in}}%
\pgfpathlineto{\pgfqpoint{4.895733in}{0.799856in}}%
\pgfpathlineto{\pgfqpoint{4.899917in}{0.787973in}}%
\pgfpathlineto{\pgfqpoint{4.909265in}{0.787349in}}%
\pgfpathlineto{\pgfqpoint{4.924189in}{0.781753in}}%
\pgfpathlineto{\pgfqpoint{4.928387in}{0.783330in}}%
\pgfpathlineto{\pgfqpoint{4.937392in}{0.776425in}}%
\pgfpathlineto{\pgfqpoint{4.942692in}{0.781753in}}%
\pgfpathlineto{\pgfqpoint{4.951877in}{0.779738in}}%
\pgfpathlineto{\pgfqpoint{4.956154in}{0.772659in}}%
\pgfpathlineto{\pgfqpoint{4.963177in}{0.770275in}}%
\pgfpathlineto{\pgfqpoint{4.968049in}{0.762933in}}%
\pgfpathlineto{\pgfqpoint{4.977778in}{0.770922in}}%
\pgfpathlineto{\pgfqpoint{4.988911in}{0.761010in}}%
\pgfpathlineto{\pgfqpoint{5.001640in}{0.765059in}}%
\pgfpathlineto{\pgfqpoint{5.008084in}{0.754622in}}%
\pgfpathlineto{\pgfqpoint{5.014386in}{0.766758in}}%
\pgfpathlineto{\pgfqpoint{5.025714in}{0.759877in}}%
\pgfpathlineto{\pgfqpoint{5.032791in}{0.761708in}}%
\pgfpathlineto{\pgfqpoint{5.042897in}{0.756843in}}%
\pgfpathlineto{\pgfqpoint{5.055999in}{0.755462in}}%
\pgfpathlineto{\pgfqpoint{5.068418in}{0.762948in}}%
\pgfpathlineto{\pgfqpoint{5.072911in}{0.759788in}}%
\pgfpathlineto{\pgfqpoint{5.085532in}{0.766205in}}%
\pgfpathlineto{\pgfqpoint{5.093575in}{0.761197in}}%
\pgfpathlineto{\pgfqpoint{5.099932in}{0.761495in}}%
\pgfpathlineto{\pgfqpoint{5.104966in}{0.767394in}}%
\pgfpathlineto{\pgfqpoint{5.112888in}{0.764459in}}%
\pgfpathlineto{\pgfqpoint{5.130724in}{0.750800in}}%
\pgfpathlineto{\pgfqpoint{5.144643in}{0.746495in}}%
\pgfpathlineto{\pgfqpoint{5.150225in}{0.741221in}}%
\pgfpathlineto{\pgfqpoint{5.157212in}{0.744093in}}%
\pgfpathlineto{\pgfqpoint{5.167799in}{0.741894in}}%
\pgfpathlineto{\pgfqpoint{5.168006in}{0.708326in}}%
\pgfpathlineto{\pgfqpoint{5.168890in}{0.643403in}}%
\pgfpathlineto{\pgfqpoint{5.176243in}{0.637164in}}%
\pgfpathlineto{\pgfqpoint{5.182137in}{0.625668in}}%
\pgfpathlineto{\pgfqpoint{5.180065in}{0.617977in}}%
\pgfpathlineto{\pgfqpoint{5.184693in}{0.614721in}}%
\pgfpathlineto{\pgfqpoint{5.188454in}{0.600538in}}%
\pgfpathlineto{\pgfqpoint{5.195658in}{0.592942in}}%
\pgfpathlineto{\pgfqpoint{5.197272in}{0.576816in}}%
\pgfpathlineto{\pgfqpoint{5.196019in}{0.569988in}}%
\pgfpathlineto{\pgfqpoint{5.186227in}{0.551919in}}%
\pgfpathlineto{\pgfqpoint{5.185101in}{0.538927in}}%
\pgfpathlineto{\pgfqpoint{5.188423in}{0.536215in}}%
\pgfpathlineto{\pgfqpoint{5.188712in}{0.524834in}}%
\pgfpathlineto{\pgfqpoint{5.185335in}{0.517680in}}%
\pgfpathlineto{\pgfqpoint{5.180085in}{0.514985in}}%
\pgfpathlineto{\pgfqpoint{5.174944in}{0.505637in}}%
\pgfpathlineto{\pgfqpoint{5.181496in}{0.496583in}}%
\pgfpathlineto{\pgfqpoint{5.168777in}{0.496405in}}%
\pgfpathlineto{\pgfqpoint{5.134766in}{0.480844in}}%
\pgfpathlineto{\pgfqpoint{5.141259in}{0.490157in}}%
\pgfpathlineto{\pgfqpoint{5.133396in}{0.495178in}}%
\pgfpathlineto{\pgfqpoint{5.131753in}{0.503714in}}%
\pgfpathlineto{\pgfqpoint{5.116406in}{0.488845in}}%
\pgfpathlineto{\pgfqpoint{5.123217in}{0.478682in}}%
\pgfpathlineto{\pgfqpoint{5.113501in}{0.465742in}}%
\pgfpathlineto{\pgfqpoint{5.108273in}{0.466013in}}%
\pgfpathlineto{\pgfqpoint{5.103354in}{0.451920in}}%
\pgfpathlineto{\pgfqpoint{5.087798in}{0.440835in}}%
\pgfpathlineto{\pgfqpoint{5.073265in}{0.436840in}}%
\pgfpathlineto{\pgfqpoint{5.047917in}{0.426447in}}%
\pgfpathlineto{\pgfqpoint{5.045291in}{0.432244in}}%
\pgfpathlineto{\pgfqpoint{5.029501in}{0.426908in}}%
\pgfpathlineto{\pgfqpoint{5.039213in}{0.417822in}}%
\pgfpathlineto{\pgfqpoint{5.023894in}{0.409929in}}%
\pgfpathlineto{\pgfqpoint{5.007371in}{0.398011in}}%
\pgfpathlineto{\pgfqpoint{5.003403in}{0.403551in}}%
\pgfpathlineto{\pgfqpoint{4.989881in}{0.395248in}}%
\pgfpathlineto{\pgfqpoint{5.003135in}{0.392095in}}%
\pgfpathlineto{\pgfqpoint{4.993142in}{0.379039in}}%
\pgfpathlineto{\pgfqpoint{4.988264in}{0.382924in}}%
\pgfpathlineto{\pgfqpoint{4.981713in}{0.376684in}}%
\pgfpathlineto{\pgfqpoint{4.989878in}{0.371233in}}%
\pgfpathlineto{\pgfqpoint{4.985050in}{0.363265in}}%
\pgfpathlineto{\pgfqpoint{4.979117in}{0.344404in}}%
\pgfpathlineto{\pgfqpoint{4.970559in}{0.326238in}}%
\pgfpathlineto{\pgfqpoint{4.974401in}{0.314331in}}%
\pgfpathlineto{\pgfqpoint{4.980932in}{0.286258in}}%
\pgfpathlineto{\pgfqpoint{4.981520in}{0.274896in}}%
\pgfpathlineto{\pgfqpoint{4.991610in}{0.259994in}}%
\pgfpathlineto{\pgfqpoint{4.983829in}{0.260865in}}%
\pgfpathlineto{\pgfqpoint{4.976272in}{0.253347in}}%
\pgfpathlineto{\pgfqpoint{4.967420in}{0.259256in}}%
\pgfpathlineto{\pgfqpoint{4.964243in}{0.265150in}}%
\pgfpathlineto{\pgfqpoint{4.951645in}{0.267892in}}%
\pgfpathlineto{\pgfqpoint{4.932402in}{0.268229in}}%
\pgfpathlineto{\pgfqpoint{4.918224in}{0.279391in}}%
\pgfpathlineto{\pgfqpoint{4.905350in}{0.281256in}}%
\pgfpathlineto{\pgfqpoint{4.897540in}{0.290127in}}%
\pgfpathlineto{\pgfqpoint{4.881183in}{0.293697in}}%
\pgfpathlineto{\pgfqpoint{4.872239in}{0.322221in}}%
\pgfpathlineto{\pgfqpoint{4.863094in}{0.333648in}}%
\pgfpathlineto{\pgfqpoint{4.864648in}{0.344506in}}%
\pgfpathlineto{\pgfqpoint{4.858978in}{0.352440in}}%
\pgfpathlineto{\pgfqpoint{4.862536in}{0.363286in}}%
\pgfpathlineto{\pgfqpoint{4.859601in}{0.371234in}}%
\pgfpathlineto{\pgfqpoint{4.850411in}{0.374834in}}%
\pgfpathlineto{\pgfqpoint{4.841818in}{0.383999in}}%
\pgfpathlineto{\pgfqpoint{4.838687in}{0.396269in}}%
\pgfpathlineto{\pgfqpoint{4.830556in}{0.407420in}}%
\pgfpathlineto{\pgfqpoint{4.819734in}{0.416096in}}%
\pgfpathlineto{\pgfqpoint{4.809986in}{0.440984in}}%
\pgfpathlineto{\pgfqpoint{4.802115in}{0.468203in}}%
\pgfpathlineto{\pgfqpoint{4.795657in}{0.478966in}}%
\pgfpathlineto{\pgfqpoint{4.784460in}{0.488026in}}%
\pgfpathlineto{\pgfqpoint{4.781662in}{0.494603in}}%
\pgfpathlineto{\pgfqpoint{4.771179in}{0.498686in}}%
\pgfpathlineto{\pgfqpoint{4.760831in}{0.516112in}}%
\pgfpathlineto{\pgfqpoint{4.751382in}{0.514774in}}%
\pgfpathlineto{\pgfqpoint{4.741787in}{0.519126in}}%
\pgfpathlineto{\pgfqpoint{4.728186in}{0.518285in}}%
\pgfpathlineto{\pgfqpoint{4.714369in}{0.525467in}}%
\pgfpathlineto{\pgfqpoint{4.710472in}{0.518637in}}%
\pgfpathlineto{\pgfqpoint{4.694325in}{0.518466in}}%
\pgfpathlineto{\pgfqpoint{4.686104in}{0.505517in}}%
\pgfpathlineto{\pgfqpoint{4.678954in}{0.489484in}}%
\pgfpathlineto{\pgfqpoint{4.663750in}{0.472274in}}%
\pgfpathlineto{\pgfqpoint{4.646524in}{0.479827in}}%
\pgfpathlineto{\pgfqpoint{4.644080in}{0.484818in}}%
\pgfpathlineto{\pgfqpoint{4.633611in}{0.488625in}}%
\pgfpathlineto{\pgfqpoint{4.630396in}{0.493833in}}%
\pgfpathlineto{\pgfqpoint{4.616499in}{0.499090in}}%
\pgfpathlineto{\pgfqpoint{4.608725in}{0.509836in}}%
\pgfpathlineto{\pgfqpoint{4.599646in}{0.515021in}}%
\pgfpathlineto{\pgfqpoint{4.591834in}{0.524068in}}%
\pgfpathlineto{\pgfqpoint{4.585749in}{0.539391in}}%
\pgfpathlineto{\pgfqpoint{4.586430in}{0.560329in}}%
\pgfpathlineto{\pgfqpoint{4.579291in}{0.570909in}}%
\pgfpathlineto{\pgfqpoint{4.578469in}{0.582369in}}%
\pgfpathlineto{\pgfqpoint{4.562630in}{0.599527in}}%
\pgfpathlineto{\pgfqpoint{4.553414in}{0.603189in}}%
\pgfpathlineto{\pgfqpoint{4.543634in}{0.619199in}}%
\pgfpathlineto{\pgfqpoint{4.535300in}{0.625560in}}%
\pgfpathlineto{\pgfqpoint{4.524711in}{0.641049in}}%
\pgfpathlineto{\pgfqpoint{4.513876in}{0.647753in}}%
\pgfpathlineto{\pgfqpoint{4.506780in}{0.664929in}}%
\pgfpathlineto{\pgfqpoint{4.498883in}{0.669271in}}%
\pgfpathclose%
\pgfusepath{fill}%
\end{pgfscope}%
\begin{pgfscope}%
\pgfpathrectangle{\pgfqpoint{3.525000in}{0.100000in}}{\pgfqpoint{2.862985in}{1.832778in}}%
\pgfusepath{clip}%
\pgfsetbuttcap%
\pgfsetmiterjoin%
\definecolor{currentfill}{rgb}{0.276078,0.800784,0.960784}%
\pgfsetfillcolor{currentfill}%
\pgfsetlinewidth{0.000000pt}%
\definecolor{currentstroke}{rgb}{0.000000,0.000000,0.000000}%
\pgfsetstrokecolor{currentstroke}%
\pgfsetstrokeopacity{0.000000}%
\pgfsetdash{}{0pt}%
\pgfpathmoveto{\pgfqpoint{4.415915in}{1.012202in}}%
\pgfpathlineto{\pgfqpoint{4.494080in}{1.000850in}}%
\pgfpathlineto{\pgfqpoint{4.524362in}{0.996184in}}%
\pgfpathlineto{\pgfqpoint{4.609722in}{0.986032in}}%
\pgfpathlineto{\pgfqpoint{4.667175in}{0.980039in}}%
\pgfpathlineto{\pgfqpoint{4.716895in}{0.975633in}}%
\pgfpathlineto{\pgfqpoint{4.714118in}{0.943987in}}%
\pgfpathlineto{\pgfqpoint{4.712612in}{0.944071in}}%
\pgfpathlineto{\pgfqpoint{4.708057in}{0.889756in}}%
\pgfpathlineto{\pgfqpoint{4.703172in}{0.834767in}}%
\pgfpathlineto{\pgfqpoint{4.697520in}{0.777201in}}%
\pgfpathlineto{\pgfqpoint{4.688807in}{0.695029in}}%
\pgfpathlineto{\pgfqpoint{4.686021in}{0.662722in}}%
\pgfpathlineto{\pgfqpoint{4.634801in}{0.667891in}}%
\pgfpathlineto{\pgfqpoint{4.590452in}{0.672249in}}%
\pgfpathlineto{\pgfqpoint{4.529118in}{0.679215in}}%
\pgfpathlineto{\pgfqpoint{4.496132in}{0.683307in}}%
\pgfpathlineto{\pgfqpoint{4.494190in}{0.676811in}}%
\pgfpathlineto{\pgfqpoint{4.498883in}{0.669271in}}%
\pgfpathlineto{\pgfqpoint{4.459240in}{0.674419in}}%
\pgfpathlineto{\pgfqpoint{4.410330in}{0.681439in}}%
\pgfpathlineto{\pgfqpoint{4.405878in}{0.653775in}}%
\pgfpathlineto{\pgfqpoint{4.361264in}{0.660490in}}%
\pgfpathlineto{\pgfqpoint{4.371747in}{0.729118in}}%
\pgfpathlineto{\pgfqpoint{4.384723in}{0.812676in}}%
\pgfpathlineto{\pgfqpoint{4.396262in}{0.885469in}}%
\pgfpathlineto{\pgfqpoint{4.406281in}{0.950172in}}%
\pgfpathlineto{\pgfqpoint{4.415915in}{1.012202in}}%
\pgfpathclose%
\pgfusepath{fill}%
\end{pgfscope}%
\begin{pgfscope}%
\pgfpathrectangle{\pgfqpoint{3.525000in}{0.100000in}}{\pgfqpoint{2.862985in}{1.832778in}}%
\pgfusepath{clip}%
\pgfsetbuttcap%
\pgfsetmiterjoin%
\definecolor{currentfill}{rgb}{0.135333,0.774394,0.980777}%
\pgfsetfillcolor{currentfill}%
\pgfsetlinewidth{0.000000pt}%
\definecolor{currentstroke}{rgb}{0.000000,0.000000,0.000000}%
\pgfsetstrokecolor{currentstroke}%
\pgfsetstrokeopacity{0.000000}%
\pgfsetdash{}{0pt}%
\pgfpathmoveto{\pgfqpoint{5.658749in}{0.610119in}}%
\pgfpathlineto{\pgfqpoint{5.584178in}{0.601837in}}%
\pgfpathlineto{\pgfqpoint{5.518446in}{0.596734in}}%
\pgfpathlineto{\pgfqpoint{5.517624in}{0.588678in}}%
\pgfpathlineto{\pgfqpoint{5.523661in}{0.581011in}}%
\pgfpathlineto{\pgfqpoint{5.531031in}{0.576507in}}%
\pgfpathlineto{\pgfqpoint{5.530919in}{0.564647in}}%
\pgfpathlineto{\pgfqpoint{5.528966in}{0.556723in}}%
\pgfpathlineto{\pgfqpoint{5.522462in}{0.551013in}}%
\pgfpathlineto{\pgfqpoint{5.513396in}{0.551618in}}%
\pgfpathlineto{\pgfqpoint{5.504825in}{0.558693in}}%
\pgfpathlineto{\pgfqpoint{5.503297in}{0.571283in}}%
\pgfpathlineto{\pgfqpoint{5.496913in}{0.578606in}}%
\pgfpathlineto{\pgfqpoint{5.492568in}{0.552398in}}%
\pgfpathlineto{\pgfqpoint{5.477803in}{0.554889in}}%
\pgfpathlineto{\pgfqpoint{5.472669in}{0.600717in}}%
\pgfpathlineto{\pgfqpoint{5.467108in}{0.649015in}}%
\pgfpathlineto{\pgfqpoint{5.469173in}{0.718697in}}%
\pgfpathlineto{\pgfqpoint{5.470455in}{0.766513in}}%
\pgfpathlineto{\pgfqpoint{5.473181in}{0.839470in}}%
\pgfpathlineto{\pgfqpoint{5.467346in}{0.845890in}}%
\pgfpathlineto{\pgfqpoint{5.477965in}{0.847429in}}%
\pgfpathlineto{\pgfqpoint{5.539830in}{0.851430in}}%
\pgfpathlineto{\pgfqpoint{5.600022in}{0.856843in}}%
\pgfpathlineto{\pgfqpoint{5.615858in}{0.801373in}}%
\pgfpathlineto{\pgfqpoint{5.626658in}{0.760657in}}%
\pgfpathlineto{\pgfqpoint{5.636339in}{0.726571in}}%
\pgfpathlineto{\pgfqpoint{5.641938in}{0.712845in}}%
\pgfpathlineto{\pgfqpoint{5.650753in}{0.700073in}}%
\pgfpathlineto{\pgfqpoint{5.649330in}{0.693573in}}%
\pgfpathlineto{\pgfqpoint{5.655627in}{0.690406in}}%
\pgfpathlineto{\pgfqpoint{5.648185in}{0.680553in}}%
\pgfpathlineto{\pgfqpoint{5.645672in}{0.663006in}}%
\pgfpathlineto{\pgfqpoint{5.652945in}{0.642483in}}%
\pgfpathlineto{\pgfqpoint{5.651935in}{0.621831in}}%
\pgfpathlineto{\pgfqpoint{5.658749in}{0.610119in}}%
\pgfpathclose%
\pgfusepath{fill}%
\end{pgfscope}%
\begin{pgfscope}%
\pgfpathrectangle{\pgfqpoint{3.525000in}{0.100000in}}{\pgfqpoint{2.862985in}{1.832778in}}%
\pgfusepath{clip}%
\pgfsetbuttcap%
\pgfsetmiterjoin%
\definecolor{currentfill}{rgb}{0.373518,0.819054,0.946943}%
\pgfsetfillcolor{currentfill}%
\pgfsetlinewidth{0.000000pt}%
\definecolor{currentstroke}{rgb}{0.000000,0.000000,0.000000}%
\pgfsetstrokecolor{currentstroke}%
\pgfsetstrokeopacity{0.000000}%
\pgfsetdash{}{0pt}%
\pgfpathmoveto{\pgfqpoint{5.477803in}{0.554889in}}%
\pgfpathlineto{\pgfqpoint{5.462652in}{0.550553in}}%
\pgfpathlineto{\pgfqpoint{5.451358in}{0.553337in}}%
\pgfpathlineto{\pgfqpoint{5.430384in}{0.546666in}}%
\pgfpathlineto{\pgfqpoint{5.414565in}{0.538074in}}%
\pgfpathlineto{\pgfqpoint{5.408075in}{0.553600in}}%
\pgfpathlineto{\pgfqpoint{5.401384in}{0.560107in}}%
\pgfpathlineto{\pgfqpoint{5.398157in}{0.567894in}}%
\pgfpathlineto{\pgfqpoint{5.402888in}{0.588972in}}%
\pgfpathlineto{\pgfqpoint{5.358060in}{0.586211in}}%
\pgfpathlineto{\pgfqpoint{5.299932in}{0.583691in}}%
\pgfpathlineto{\pgfqpoint{5.303390in}{0.588938in}}%
\pgfpathlineto{\pgfqpoint{5.299116in}{0.598847in}}%
\pgfpathlineto{\pgfqpoint{5.303375in}{0.600866in}}%
\pgfpathlineto{\pgfqpoint{5.302430in}{0.610386in}}%
\pgfpathlineto{\pgfqpoint{5.309854in}{0.619619in}}%
\pgfpathlineto{\pgfqpoint{5.313916in}{0.637542in}}%
\pgfpathlineto{\pgfqpoint{5.327503in}{0.649350in}}%
\pgfpathlineto{\pgfqpoint{5.322467in}{0.660818in}}%
\pgfpathlineto{\pgfqpoint{5.332102in}{0.667419in}}%
\pgfpathlineto{\pgfqpoint{5.323791in}{0.679370in}}%
\pgfpathlineto{\pgfqpoint{5.317776in}{0.706727in}}%
\pgfpathlineto{\pgfqpoint{5.320057in}{0.711695in}}%
\pgfpathlineto{\pgfqpoint{5.323659in}{0.721223in}}%
\pgfpathlineto{\pgfqpoint{5.320307in}{0.731223in}}%
\pgfpathlineto{\pgfqpoint{5.321946in}{0.735772in}}%
\pgfpathlineto{\pgfqpoint{5.316048in}{0.752960in}}%
\pgfpathlineto{\pgfqpoint{5.333459in}{0.783821in}}%
\pgfpathlineto{\pgfqpoint{5.334269in}{0.792035in}}%
\pgfpathlineto{\pgfqpoint{5.342656in}{0.798007in}}%
\pgfpathlineto{\pgfqpoint{5.351605in}{0.817725in}}%
\pgfpathlineto{\pgfqpoint{5.351180in}{0.825739in}}%
\pgfpathlineto{\pgfqpoint{5.362326in}{0.833932in}}%
\pgfpathlineto{\pgfqpoint{5.358848in}{0.838785in}}%
\pgfpathlineto{\pgfqpoint{5.425643in}{0.842683in}}%
\pgfpathlineto{\pgfqpoint{5.467346in}{0.845890in}}%
\pgfpathlineto{\pgfqpoint{5.473181in}{0.839470in}}%
\pgfpathlineto{\pgfqpoint{5.470455in}{0.766513in}}%
\pgfpathlineto{\pgfqpoint{5.469173in}{0.718697in}}%
\pgfpathlineto{\pgfqpoint{5.467108in}{0.649015in}}%
\pgfpathlineto{\pgfqpoint{5.472669in}{0.600717in}}%
\pgfpathlineto{\pgfqpoint{5.477803in}{0.554889in}}%
\pgfpathclose%
\pgfusepath{fill}%
\end{pgfscope}%
\begin{pgfscope}%
\pgfpathrectangle{\pgfqpoint{3.525000in}{0.100000in}}{\pgfqpoint{2.862985in}{1.832778in}}%
\pgfusepath{clip}%
\pgfsetbuttcap%
\pgfsetmiterjoin%
\definecolor{currentfill}{rgb}{0.050750,0.763230,0.791926}%
\pgfsetfillcolor{currentfill}%
\pgfsetlinewidth{0.000000pt}%
\definecolor{currentstroke}{rgb}{0.000000,0.000000,0.000000}%
\pgfsetstrokecolor{currentstroke}%
\pgfsetstrokeopacity{0.000000}%
\pgfsetdash{}{0pt}%
\pgfpathmoveto{\pgfqpoint{5.600022in}{0.856843in}}%
\pgfpathlineto{\pgfqpoint{5.665820in}{0.864070in}}%
\pgfpathlineto{\pgfqpoint{5.705410in}{0.868948in}}%
\pgfpathlineto{\pgfqpoint{5.727958in}{0.872540in}}%
\pgfpathlineto{\pgfqpoint{5.718618in}{0.857809in}}%
\pgfpathlineto{\pgfqpoint{5.718623in}{0.850863in}}%
\pgfpathlineto{\pgfqpoint{5.728238in}{0.847117in}}%
\pgfpathlineto{\pgfqpoint{5.734776in}{0.841070in}}%
\pgfpathlineto{\pgfqpoint{5.744657in}{0.840323in}}%
\pgfpathlineto{\pgfqpoint{5.753892in}{0.823299in}}%
\pgfpathlineto{\pgfqpoint{5.763910in}{0.811294in}}%
\pgfpathlineto{\pgfqpoint{5.781702in}{0.801205in}}%
\pgfpathlineto{\pgfqpoint{5.786758in}{0.793157in}}%
\pgfpathlineto{\pgfqpoint{5.802624in}{0.784454in}}%
\pgfpathlineto{\pgfqpoint{5.803817in}{0.778143in}}%
\pgfpathlineto{\pgfqpoint{5.816688in}{0.765593in}}%
\pgfpathlineto{\pgfqpoint{5.826546in}{0.762019in}}%
\pgfpathlineto{\pgfqpoint{5.833528in}{0.750169in}}%
\pgfpathlineto{\pgfqpoint{5.836618in}{0.736849in}}%
\pgfpathlineto{\pgfqpoint{5.846858in}{0.731700in}}%
\pgfpathlineto{\pgfqpoint{5.855102in}{0.717383in}}%
\pgfpathlineto{\pgfqpoint{5.857761in}{0.706853in}}%
\pgfpathlineto{\pgfqpoint{5.869841in}{0.702371in}}%
\pgfpathlineto{\pgfqpoint{5.859704in}{0.682961in}}%
\pgfpathlineto{\pgfqpoint{5.855889in}{0.671491in}}%
\pgfpathlineto{\pgfqpoint{5.858378in}{0.666114in}}%
\pgfpathlineto{\pgfqpoint{5.854032in}{0.657178in}}%
\pgfpathlineto{\pgfqpoint{5.852185in}{0.644824in}}%
\pgfpathlineto{\pgfqpoint{5.844456in}{0.642523in}}%
\pgfpathlineto{\pgfqpoint{5.848537in}{0.631212in}}%
\pgfpathlineto{\pgfqpoint{5.848277in}{0.616862in}}%
\pgfpathlineto{\pgfqpoint{5.834539in}{0.617400in}}%
\pgfpathlineto{\pgfqpoint{5.824969in}{0.620018in}}%
\pgfpathlineto{\pgfqpoint{5.820828in}{0.615151in}}%
\pgfpathlineto{\pgfqpoint{5.823918in}{0.603653in}}%
\pgfpathlineto{\pgfqpoint{5.823246in}{0.590285in}}%
\pgfpathlineto{\pgfqpoint{5.817209in}{0.589262in}}%
\pgfpathlineto{\pgfqpoint{5.812325in}{0.601747in}}%
\pgfpathlineto{\pgfqpoint{5.762593in}{0.598433in}}%
\pgfpathlineto{\pgfqpoint{5.699885in}{0.595011in}}%
\pgfpathlineto{\pgfqpoint{5.668260in}{0.592785in}}%
\pgfpathlineto{\pgfqpoint{5.658749in}{0.610119in}}%
\pgfpathlineto{\pgfqpoint{5.651935in}{0.621831in}}%
\pgfpathlineto{\pgfqpoint{5.652945in}{0.642483in}}%
\pgfpathlineto{\pgfqpoint{5.645672in}{0.663006in}}%
\pgfpathlineto{\pgfqpoint{5.648185in}{0.680553in}}%
\pgfpathlineto{\pgfqpoint{5.655627in}{0.690406in}}%
\pgfpathlineto{\pgfqpoint{5.649330in}{0.693573in}}%
\pgfpathlineto{\pgfqpoint{5.650753in}{0.700073in}}%
\pgfpathlineto{\pgfqpoint{5.641938in}{0.712845in}}%
\pgfpathlineto{\pgfqpoint{5.636339in}{0.726571in}}%
\pgfpathlineto{\pgfqpoint{5.626658in}{0.760657in}}%
\pgfpathlineto{\pgfqpoint{5.615858in}{0.801373in}}%
\pgfpathlineto{\pgfqpoint{5.600022in}{0.856843in}}%
\pgfpathclose%
\pgfusepath{fill}%
\end{pgfscope}%
\begin{pgfscope}%
\pgfpathrectangle{\pgfqpoint{3.525000in}{0.100000in}}{\pgfqpoint{2.862985in}{1.832778in}}%
\pgfusepath{clip}%
\pgfsetbuttcap%
\pgfsetmiterjoin%
\definecolor{currentfill}{rgb}{0.050750,0.763230,0.791926}%
\pgfsetfillcolor{currentfill}%
\pgfsetlinewidth{0.000000pt}%
\definecolor{currentstroke}{rgb}{0.000000,0.000000,0.000000}%
\pgfsetstrokecolor{currentstroke}%
\pgfsetstrokeopacity{0.000000}%
\pgfsetdash{}{0pt}%
\pgfpathmoveto{\pgfqpoint{5.727958in}{0.872540in}}%
\pgfpathlineto{\pgfqpoint{5.754049in}{0.885519in}}%
\pgfpathlineto{\pgfqpoint{5.768499in}{0.890420in}}%
\pgfpathlineto{\pgfqpoint{5.831366in}{0.896952in}}%
\pgfpathlineto{\pgfqpoint{5.838024in}{0.894789in}}%
\pgfpathlineto{\pgfqpoint{5.846819in}{0.885949in}}%
\pgfpathlineto{\pgfqpoint{5.847304in}{0.878098in}}%
\pgfpathlineto{\pgfqpoint{5.905165in}{0.886600in}}%
\pgfpathlineto{\pgfqpoint{5.970869in}{0.839501in}}%
\pgfpathlineto{\pgfqpoint{5.958547in}{0.826658in}}%
\pgfpathlineto{\pgfqpoint{5.941640in}{0.796833in}}%
\pgfpathlineto{\pgfqpoint{5.946453in}{0.790423in}}%
\pgfpathlineto{\pgfqpoint{5.937434in}{0.777981in}}%
\pgfpathlineto{\pgfqpoint{5.928462in}{0.776573in}}%
\pgfpathlineto{\pgfqpoint{5.928429in}{0.769097in}}%
\pgfpathlineto{\pgfqpoint{5.921939in}{0.761274in}}%
\pgfpathlineto{\pgfqpoint{5.913862in}{0.759687in}}%
\pgfpathlineto{\pgfqpoint{5.915625in}{0.752736in}}%
\pgfpathlineto{\pgfqpoint{5.911118in}{0.747400in}}%
\pgfpathlineto{\pgfqpoint{5.900337in}{0.742790in}}%
\pgfpathlineto{\pgfqpoint{5.887196in}{0.732304in}}%
\pgfpathlineto{\pgfqpoint{5.889677in}{0.725520in}}%
\pgfpathlineto{\pgfqpoint{5.881407in}{0.721260in}}%
\pgfpathlineto{\pgfqpoint{5.874432in}{0.725739in}}%
\pgfpathlineto{\pgfqpoint{5.869331in}{0.706267in}}%
\pgfpathlineto{\pgfqpoint{5.857761in}{0.706853in}}%
\pgfpathlineto{\pgfqpoint{5.855102in}{0.717383in}}%
\pgfpathlineto{\pgfqpoint{5.846858in}{0.731700in}}%
\pgfpathlineto{\pgfqpoint{5.836618in}{0.736849in}}%
\pgfpathlineto{\pgfqpoint{5.833528in}{0.750169in}}%
\pgfpathlineto{\pgfqpoint{5.826546in}{0.762019in}}%
\pgfpathlineto{\pgfqpoint{5.816688in}{0.765593in}}%
\pgfpathlineto{\pgfqpoint{5.803817in}{0.778143in}}%
\pgfpathlineto{\pgfqpoint{5.802624in}{0.784454in}}%
\pgfpathlineto{\pgfqpoint{5.786758in}{0.793157in}}%
\pgfpathlineto{\pgfqpoint{5.781702in}{0.801205in}}%
\pgfpathlineto{\pgfqpoint{5.763910in}{0.811294in}}%
\pgfpathlineto{\pgfqpoint{5.753892in}{0.823299in}}%
\pgfpathlineto{\pgfqpoint{5.744657in}{0.840323in}}%
\pgfpathlineto{\pgfqpoint{5.734776in}{0.841070in}}%
\pgfpathlineto{\pgfqpoint{5.728238in}{0.847117in}}%
\pgfpathlineto{\pgfqpoint{5.718623in}{0.850863in}}%
\pgfpathlineto{\pgfqpoint{5.718618in}{0.857809in}}%
\pgfpathlineto{\pgfqpoint{5.727958in}{0.872540in}}%
\pgfpathclose%
\pgfusepath{fill}%
\end{pgfscope}%
\begin{pgfscope}%
\pgfpathrectangle{\pgfqpoint{3.525000in}{0.100000in}}{\pgfqpoint{2.862985in}{1.832778in}}%
\pgfusepath{clip}%
\pgfsetbuttcap%
\pgfsetmiterjoin%
\definecolor{currentfill}{rgb}{0.027067,0.754095,0.996155}%
\pgfsetfillcolor{currentfill}%
\pgfsetlinewidth{0.000000pt}%
\definecolor{currentstroke}{rgb}{0.000000,0.000000,0.000000}%
\pgfsetstrokecolor{currentstroke}%
\pgfsetstrokeopacity{0.000000}%
\pgfsetdash{}{0pt}%
\pgfpathmoveto{\pgfqpoint{5.136515in}{0.926094in}}%
\pgfpathlineto{\pgfqpoint{5.174856in}{0.926440in}}%
\pgfpathlineto{\pgfqpoint{5.225580in}{0.927364in}}%
\pgfpathlineto{\pgfqpoint{5.312257in}{0.930116in}}%
\pgfpathlineto{\pgfqpoint{5.361833in}{0.932737in}}%
\pgfpathlineto{\pgfqpoint{5.367150in}{0.926155in}}%
\pgfpathlineto{\pgfqpoint{5.366818in}{0.919204in}}%
\pgfpathlineto{\pgfqpoint{5.354808in}{0.907211in}}%
\pgfpathlineto{\pgfqpoint{5.351904in}{0.900644in}}%
\pgfpathlineto{\pgfqpoint{5.385213in}{0.903107in}}%
\pgfpathlineto{\pgfqpoint{5.385195in}{0.890974in}}%
\pgfpathlineto{\pgfqpoint{5.374509in}{0.885789in}}%
\pgfpathlineto{\pgfqpoint{5.374660in}{0.877545in}}%
\pgfpathlineto{\pgfqpoint{5.370647in}{0.871619in}}%
\pgfpathlineto{\pgfqpoint{5.368020in}{0.859033in}}%
\pgfpathlineto{\pgfqpoint{5.370732in}{0.848649in}}%
\pgfpathlineto{\pgfqpoint{5.358848in}{0.838785in}}%
\pgfpathlineto{\pgfqpoint{5.362326in}{0.833932in}}%
\pgfpathlineto{\pgfqpoint{5.351180in}{0.825739in}}%
\pgfpathlineto{\pgfqpoint{5.351605in}{0.817725in}}%
\pgfpathlineto{\pgfqpoint{5.342656in}{0.798007in}}%
\pgfpathlineto{\pgfqpoint{5.334269in}{0.792035in}}%
\pgfpathlineto{\pgfqpoint{5.333459in}{0.783821in}}%
\pgfpathlineto{\pgfqpoint{5.316048in}{0.752960in}}%
\pgfpathlineto{\pgfqpoint{5.321946in}{0.735772in}}%
\pgfpathlineto{\pgfqpoint{5.320307in}{0.731223in}}%
\pgfpathlineto{\pgfqpoint{5.323659in}{0.721223in}}%
\pgfpathlineto{\pgfqpoint{5.320057in}{0.711695in}}%
\pgfpathlineto{\pgfqpoint{5.272446in}{0.709729in}}%
\pgfpathlineto{\pgfqpoint{5.210638in}{0.708711in}}%
\pgfpathlineto{\pgfqpoint{5.168006in}{0.708326in}}%
\pgfpathlineto{\pgfqpoint{5.167799in}{0.741894in}}%
\pgfpathlineto{\pgfqpoint{5.157212in}{0.744093in}}%
\pgfpathlineto{\pgfqpoint{5.150225in}{0.741221in}}%
\pgfpathlineto{\pgfqpoint{5.144643in}{0.746495in}}%
\pgfpathlineto{\pgfqpoint{5.145219in}{0.782052in}}%
\pgfpathlineto{\pgfqpoint{5.146468in}{0.857760in}}%
\pgfpathlineto{\pgfqpoint{5.140413in}{0.902092in}}%
\pgfpathlineto{\pgfqpoint{5.136515in}{0.926094in}}%
\pgfpathclose%
\pgfusepath{fill}%
\end{pgfscope}%
\begin{pgfscope}%
\pgfpathrectangle{\pgfqpoint{3.525000in}{0.100000in}}{\pgfqpoint{2.862985in}{1.832778in}}%
\pgfusepath{clip}%
\pgfsetbuttcap%
\pgfsetmiterjoin%
\definecolor{currentfill}{rgb}{0.135333,0.774394,0.980777}%
\pgfsetfillcolor{currentfill}%
\pgfsetlinewidth{0.000000pt}%
\definecolor{currentstroke}{rgb}{0.000000,0.000000,0.000000}%
\pgfsetstrokecolor{currentstroke}%
\pgfsetstrokeopacity{0.000000}%
\pgfsetdash{}{0pt}%
\pgfpathmoveto{\pgfqpoint{5.168006in}{0.708326in}}%
\pgfpathlineto{\pgfqpoint{5.210638in}{0.708711in}}%
\pgfpathlineto{\pgfqpoint{5.272446in}{0.709729in}}%
\pgfpathlineto{\pgfqpoint{5.320057in}{0.711695in}}%
\pgfpathlineto{\pgfqpoint{5.317776in}{0.706727in}}%
\pgfpathlineto{\pgfqpoint{5.323791in}{0.679370in}}%
\pgfpathlineto{\pgfqpoint{5.332102in}{0.667419in}}%
\pgfpathlineto{\pgfqpoint{5.322467in}{0.660818in}}%
\pgfpathlineto{\pgfqpoint{5.327503in}{0.649350in}}%
\pgfpathlineto{\pgfqpoint{5.313916in}{0.637542in}}%
\pgfpathlineto{\pgfqpoint{5.309854in}{0.619619in}}%
\pgfpathlineto{\pgfqpoint{5.302430in}{0.610386in}}%
\pgfpathlineto{\pgfqpoint{5.303375in}{0.600866in}}%
\pgfpathlineto{\pgfqpoint{5.299116in}{0.598847in}}%
\pgfpathlineto{\pgfqpoint{5.303390in}{0.588938in}}%
\pgfpathlineto{\pgfqpoint{5.299932in}{0.583691in}}%
\pgfpathlineto{\pgfqpoint{5.358060in}{0.586211in}}%
\pgfpathlineto{\pgfqpoint{5.402888in}{0.588972in}}%
\pgfpathlineto{\pgfqpoint{5.398157in}{0.567894in}}%
\pgfpathlineto{\pgfqpoint{5.401384in}{0.560107in}}%
\pgfpathlineto{\pgfqpoint{5.408075in}{0.553600in}}%
\pgfpathlineto{\pgfqpoint{5.414565in}{0.538074in}}%
\pgfpathlineto{\pgfqpoint{5.394065in}{0.541651in}}%
\pgfpathlineto{\pgfqpoint{5.386502in}{0.547531in}}%
\pgfpathlineto{\pgfqpoint{5.377492in}{0.547811in}}%
\pgfpathlineto{\pgfqpoint{5.368019in}{0.534934in}}%
\pgfpathlineto{\pgfqpoint{5.369907in}{0.529071in}}%
\pgfpathlineto{\pgfqpoint{5.400169in}{0.525523in}}%
\pgfpathlineto{\pgfqpoint{5.408091in}{0.518783in}}%
\pgfpathlineto{\pgfqpoint{5.420984in}{0.515315in}}%
\pgfpathlineto{\pgfqpoint{5.415353in}{0.507336in}}%
\pgfpathlineto{\pgfqpoint{5.405867in}{0.500376in}}%
\pgfpathlineto{\pgfqpoint{5.430906in}{0.484726in}}%
\pgfpathlineto{\pgfqpoint{5.442569in}{0.482281in}}%
\pgfpathlineto{\pgfqpoint{5.447419in}{0.469488in}}%
\pgfpathlineto{\pgfqpoint{5.436438in}{0.467135in}}%
\pgfpathlineto{\pgfqpoint{5.415799in}{0.483224in}}%
\pgfpathlineto{\pgfqpoint{5.407729in}{0.485461in}}%
\pgfpathlineto{\pgfqpoint{5.403824in}{0.491817in}}%
\pgfpathlineto{\pgfqpoint{5.391854in}{0.489208in}}%
\pgfpathlineto{\pgfqpoint{5.388107in}{0.481018in}}%
\pgfpathlineto{\pgfqpoint{5.390480in}{0.471892in}}%
\pgfpathlineto{\pgfqpoint{5.382422in}{0.466506in}}%
\pgfpathlineto{\pgfqpoint{5.375068in}{0.479774in}}%
\pgfpathlineto{\pgfqpoint{5.362187in}{0.475818in}}%
\pgfpathlineto{\pgfqpoint{5.357360in}{0.467901in}}%
\pgfpathlineto{\pgfqpoint{5.348213in}{0.470153in}}%
\pgfpathlineto{\pgfqpoint{5.342364in}{0.480150in}}%
\pgfpathlineto{\pgfqpoint{5.326847in}{0.483451in}}%
\pgfpathlineto{\pgfqpoint{5.323917in}{0.488653in}}%
\pgfpathlineto{\pgfqpoint{5.307751in}{0.497715in}}%
\pgfpathlineto{\pgfqpoint{5.303738in}{0.505638in}}%
\pgfpathlineto{\pgfqpoint{5.290197in}{0.502400in}}%
\pgfpathlineto{\pgfqpoint{5.291917in}{0.509654in}}%
\pgfpathlineto{\pgfqpoint{5.275086in}{0.502213in}}%
\pgfpathlineto{\pgfqpoint{5.279607in}{0.494492in}}%
\pgfpathlineto{\pgfqpoint{5.266610in}{0.489928in}}%
\pgfpathlineto{\pgfqpoint{5.249395in}{0.492437in}}%
\pgfpathlineto{\pgfqpoint{5.214551in}{0.504372in}}%
\pgfpathlineto{\pgfqpoint{5.187666in}{0.502003in}}%
\pgfpathlineto{\pgfqpoint{5.185335in}{0.517680in}}%
\pgfpathlineto{\pgfqpoint{5.188712in}{0.524834in}}%
\pgfpathlineto{\pgfqpoint{5.188423in}{0.536215in}}%
\pgfpathlineto{\pgfqpoint{5.185101in}{0.538927in}}%
\pgfpathlineto{\pgfqpoint{5.186227in}{0.551919in}}%
\pgfpathlineto{\pgfqpoint{5.196019in}{0.569988in}}%
\pgfpathlineto{\pgfqpoint{5.197272in}{0.576816in}}%
\pgfpathlineto{\pgfqpoint{5.195658in}{0.592942in}}%
\pgfpathlineto{\pgfqpoint{5.188454in}{0.600538in}}%
\pgfpathlineto{\pgfqpoint{5.184693in}{0.614721in}}%
\pgfpathlineto{\pgfqpoint{5.180065in}{0.617977in}}%
\pgfpathlineto{\pgfqpoint{5.182137in}{0.625668in}}%
\pgfpathlineto{\pgfqpoint{5.176243in}{0.637164in}}%
\pgfpathlineto{\pgfqpoint{5.168890in}{0.643403in}}%
\pgfpathlineto{\pgfqpoint{5.168006in}{0.708326in}}%
\pgfpathclose%
\pgfusepath{fill}%
\end{pgfscope}%
\begin{pgfscope}%
\pgfpathrectangle{\pgfqpoint{3.525000in}{0.100000in}}{\pgfqpoint{2.862985in}{1.832778in}}%
\pgfusepath{clip}%
\pgfsetbuttcap%
\pgfsetmiterjoin%
\definecolor{currentfill}{rgb}{0.135333,0.774394,0.980777}%
\pgfsetfillcolor{currentfill}%
\pgfsetlinewidth{0.000000pt}%
\definecolor{currentstroke}{rgb}{0.000000,0.000000,0.000000}%
\pgfsetstrokecolor{currentstroke}%
\pgfsetstrokeopacity{0.000000}%
\pgfsetdash{}{0pt}%
\pgfpathmoveto{\pgfqpoint{5.282025in}{0.493901in}}%
\pgfpathlineto{\pgfqpoint{5.288193in}{0.497570in}}%
\pgfpathlineto{\pgfqpoint{5.295679in}{0.493233in}}%
\pgfpathlineto{\pgfqpoint{5.291504in}{0.487265in}}%
\pgfpathlineto{\pgfqpoint{5.282025in}{0.493901in}}%
\pgfpathclose%
\pgfusepath{fill}%
\end{pgfscope}%
\begin{pgfscope}%
\pgfpathrectangle{\pgfqpoint{3.525000in}{0.100000in}}{\pgfqpoint{2.862985in}{1.832778in}}%
\pgfusepath{clip}%
\pgfsetbuttcap%
\pgfsetmiterjoin%
\definecolor{currentfill}{rgb}{0.096886,0.776148,0.602768}%
\pgfsetfillcolor{currentfill}%
\pgfsetlinewidth{0.000000pt}%
\definecolor{currentstroke}{rgb}{0.000000,0.000000,0.000000}%
\pgfsetstrokecolor{currentstroke}%
\pgfsetstrokeopacity{0.000000}%
\pgfsetdash{}{0pt}%
\pgfpathmoveto{\pgfqpoint{5.530919in}{0.564647in}}%
\pgfpathlineto{\pgfqpoint{5.531031in}{0.576507in}}%
\pgfpathlineto{\pgfqpoint{5.523661in}{0.581011in}}%
\pgfpathlineto{\pgfqpoint{5.517624in}{0.588678in}}%
\pgfpathlineto{\pgfqpoint{5.518446in}{0.596734in}}%
\pgfpathlineto{\pgfqpoint{5.584178in}{0.601837in}}%
\pgfpathlineto{\pgfqpoint{5.658749in}{0.610119in}}%
\pgfpathlineto{\pgfqpoint{5.668260in}{0.592785in}}%
\pgfpathlineto{\pgfqpoint{5.699885in}{0.595011in}}%
\pgfpathlineto{\pgfqpoint{5.762593in}{0.598433in}}%
\pgfpathlineto{\pgfqpoint{5.812325in}{0.601747in}}%
\pgfpathlineto{\pgfqpoint{5.817209in}{0.589262in}}%
\pgfpathlineto{\pgfqpoint{5.823246in}{0.590285in}}%
\pgfpathlineto{\pgfqpoint{5.823918in}{0.603653in}}%
\pgfpathlineto{\pgfqpoint{5.820828in}{0.615151in}}%
\pgfpathlineto{\pgfqpoint{5.824969in}{0.620018in}}%
\pgfpathlineto{\pgfqpoint{5.834539in}{0.617400in}}%
\pgfpathlineto{\pgfqpoint{5.848277in}{0.616862in}}%
\pgfpathlineto{\pgfqpoint{5.850389in}{0.606589in}}%
\pgfpathlineto{\pgfqpoint{5.854615in}{0.600719in}}%
\pgfpathlineto{\pgfqpoint{5.857915in}{0.587877in}}%
\pgfpathlineto{\pgfqpoint{5.868139in}{0.568003in}}%
\pgfpathlineto{\pgfqpoint{5.868190in}{0.562628in}}%
\pgfpathlineto{\pgfqpoint{5.883302in}{0.539291in}}%
\pgfpathlineto{\pgfqpoint{5.884764in}{0.534471in}}%
\pgfpathlineto{\pgfqpoint{5.907434in}{0.501512in}}%
\pgfpathlineto{\pgfqpoint{5.903842in}{0.500979in}}%
\pgfpathlineto{\pgfqpoint{5.913399in}{0.477537in}}%
\pgfpathlineto{\pgfqpoint{5.939653in}{0.436782in}}%
\pgfpathlineto{\pgfqpoint{5.954896in}{0.406818in}}%
\pgfpathlineto{\pgfqpoint{5.959991in}{0.401571in}}%
\pgfpathlineto{\pgfqpoint{5.968722in}{0.382763in}}%
\pgfpathlineto{\pgfqpoint{5.971750in}{0.352646in}}%
\pgfpathlineto{\pgfqpoint{5.972955in}{0.330127in}}%
\pgfpathlineto{\pgfqpoint{5.971500in}{0.315709in}}%
\pgfpathlineto{\pgfqpoint{5.966819in}{0.305401in}}%
\pgfpathlineto{\pgfqpoint{5.968994in}{0.291885in}}%
\pgfpathlineto{\pgfqpoint{5.963954in}{0.281186in}}%
\pgfpathlineto{\pgfqpoint{5.948991in}{0.272413in}}%
\pgfpathlineto{\pgfqpoint{5.939259in}{0.273053in}}%
\pgfpathlineto{\pgfqpoint{5.932940in}{0.268467in}}%
\pgfpathlineto{\pgfqpoint{5.931129in}{0.280712in}}%
\pgfpathlineto{\pgfqpoint{5.920668in}{0.283944in}}%
\pgfpathlineto{\pgfqpoint{5.911290in}{0.301030in}}%
\pgfpathlineto{\pgfqpoint{5.910231in}{0.308810in}}%
\pgfpathlineto{\pgfqpoint{5.893448in}{0.313535in}}%
\pgfpathlineto{\pgfqpoint{5.882633in}{0.312512in}}%
\pgfpathlineto{\pgfqpoint{5.876514in}{0.323801in}}%
\pgfpathlineto{\pgfqpoint{5.869508in}{0.344093in}}%
\pgfpathlineto{\pgfqpoint{5.859786in}{0.348206in}}%
\pgfpathlineto{\pgfqpoint{5.854515in}{0.359834in}}%
\pgfpathlineto{\pgfqpoint{5.841625in}{0.366687in}}%
\pgfpathlineto{\pgfqpoint{5.834174in}{0.375259in}}%
\pgfpathlineto{\pgfqpoint{5.821320in}{0.398564in}}%
\pgfpathlineto{\pgfqpoint{5.819809in}{0.414966in}}%
\pgfpathlineto{\pgfqpoint{5.826813in}{0.425461in}}%
\pgfpathlineto{\pgfqpoint{5.826113in}{0.432754in}}%
\pgfpathlineto{\pgfqpoint{5.818019in}{0.433533in}}%
\pgfpathlineto{\pgfqpoint{5.811256in}{0.438606in}}%
\pgfpathlineto{\pgfqpoint{5.807553in}{0.432412in}}%
\pgfpathlineto{\pgfqpoint{5.814043in}{0.427365in}}%
\pgfpathlineto{\pgfqpoint{5.809319in}{0.418289in}}%
\pgfpathlineto{\pgfqpoint{5.801677in}{0.425822in}}%
\pgfpathlineto{\pgfqpoint{5.802564in}{0.446710in}}%
\pgfpathlineto{\pgfqpoint{5.806250in}{0.463657in}}%
\pgfpathlineto{\pgfqpoint{5.804355in}{0.492736in}}%
\pgfpathlineto{\pgfqpoint{5.796738in}{0.499662in}}%
\pgfpathlineto{\pgfqpoint{5.792907in}{0.508549in}}%
\pgfpathlineto{\pgfqpoint{5.779785in}{0.508352in}}%
\pgfpathlineto{\pgfqpoint{5.766814in}{0.522976in}}%
\pgfpathlineto{\pgfqpoint{5.758117in}{0.527359in}}%
\pgfpathlineto{\pgfqpoint{5.755538in}{0.536620in}}%
\pgfpathlineto{\pgfqpoint{5.747031in}{0.539872in}}%
\pgfpathlineto{\pgfqpoint{5.739981in}{0.550098in}}%
\pgfpathlineto{\pgfqpoint{5.721358in}{0.558444in}}%
\pgfpathlineto{\pgfqpoint{5.706879in}{0.558658in}}%
\pgfpathlineto{\pgfqpoint{5.700576in}{0.555452in}}%
\pgfpathlineto{\pgfqpoint{5.702143in}{0.545445in}}%
\pgfpathlineto{\pgfqpoint{5.695570in}{0.545919in}}%
\pgfpathlineto{\pgfqpoint{5.675353in}{0.531927in}}%
\pgfpathlineto{\pgfqpoint{5.651068in}{0.526381in}}%
\pgfpathlineto{\pgfqpoint{5.650669in}{0.533243in}}%
\pgfpathlineto{\pgfqpoint{5.645287in}{0.539949in}}%
\pgfpathlineto{\pgfqpoint{5.630848in}{0.549208in}}%
\pgfpathlineto{\pgfqpoint{5.610121in}{0.558700in}}%
\pgfpathlineto{\pgfqpoint{5.587608in}{0.563725in}}%
\pgfpathlineto{\pgfqpoint{5.583382in}{0.570558in}}%
\pgfpathlineto{\pgfqpoint{5.575265in}{0.564859in}}%
\pgfpathlineto{\pgfqpoint{5.565491in}{0.563602in}}%
\pgfpathlineto{\pgfqpoint{5.531467in}{0.554637in}}%
\pgfpathlineto{\pgfqpoint{5.530919in}{0.564647in}}%
\pgfpathclose%
\pgfusepath{fill}%
\end{pgfscope}%
\begin{pgfscope}%
\pgfpathrectangle{\pgfqpoint{3.525000in}{0.100000in}}{\pgfqpoint{2.862985in}{1.832778in}}%
\pgfusepath{clip}%
\pgfsetbuttcap%
\pgfsetmiterjoin%
\definecolor{currentfill}{rgb}{0.096886,0.776148,0.602768}%
\pgfsetfillcolor{currentfill}%
\pgfsetlinewidth{0.000000pt}%
\definecolor{currentstroke}{rgb}{0.000000,0.000000,0.000000}%
\pgfsetstrokecolor{currentstroke}%
\pgfsetstrokeopacity{0.000000}%
\pgfsetdash{}{0pt}%
\pgfpathmoveto{\pgfqpoint{5.910737in}{0.501980in}}%
\pgfpathlineto{\pgfqpoint{5.919540in}{0.491740in}}%
\pgfpathlineto{\pgfqpoint{5.909372in}{0.491096in}}%
\pgfpathlineto{\pgfqpoint{5.910737in}{0.501980in}}%
\pgfpathclose%
\pgfusepath{fill}%
\end{pgfscope}%
\begin{pgfscope}%
\pgfpathrectangle{\pgfqpoint{3.525000in}{0.100000in}}{\pgfqpoint{2.862985in}{1.832778in}}%
\pgfusepath{clip}%
\pgfsetbuttcap%
\pgfsetmiterjoin%
\definecolor{currentfill}{rgb}{0.093810,0.775286,0.615379}%
\pgfsetfillcolor{currentfill}%
\pgfsetlinewidth{0.000000pt}%
\definecolor{currentstroke}{rgb}{0.000000,0.000000,0.000000}%
\pgfsetstrokecolor{currentstroke}%
\pgfsetstrokeopacity{0.000000}%
\pgfsetdash{}{0pt}%
\pgfpathmoveto{\pgfqpoint{5.392522in}{1.671190in}}%
\pgfpathlineto{\pgfqpoint{5.387746in}{1.661894in}}%
\pgfpathlineto{\pgfqpoint{5.376351in}{1.656468in}}%
\pgfpathlineto{\pgfqpoint{5.371416in}{1.649166in}}%
\pgfpathlineto{\pgfqpoint{5.365635in}{1.654425in}}%
\pgfpathlineto{\pgfqpoint{5.392522in}{1.671190in}}%
\pgfpathclose%
\pgfusepath{fill}%
\end{pgfscope}%
\begin{pgfscope}%
\pgfpathrectangle{\pgfqpoint{3.525000in}{0.100000in}}{\pgfqpoint{2.862985in}{1.832778in}}%
\pgfusepath{clip}%
\pgfsetbuttcap%
\pgfsetmiterjoin%
\definecolor{currentfill}{rgb}{0.093810,0.775286,0.615379}%
\pgfsetfillcolor{currentfill}%
\pgfsetlinewidth{0.000000pt}%
\definecolor{currentstroke}{rgb}{0.000000,0.000000,0.000000}%
\pgfsetstrokecolor{currentstroke}%
\pgfsetstrokeopacity{0.000000}%
\pgfsetdash{}{0pt}%
\pgfpathmoveto{\pgfqpoint{5.396396in}{1.615244in}}%
\pgfpathlineto{\pgfqpoint{5.408054in}{1.626120in}}%
\pgfpathlineto{\pgfqpoint{5.426037in}{1.628969in}}%
\pgfpathlineto{\pgfqpoint{5.421084in}{1.621405in}}%
\pgfpathlineto{\pgfqpoint{5.408756in}{1.610466in}}%
\pgfpathlineto{\pgfqpoint{5.401549in}{1.596426in}}%
\pgfpathlineto{\pgfqpoint{5.398496in}{1.604041in}}%
\pgfpathlineto{\pgfqpoint{5.393076in}{1.605134in}}%
\pgfpathlineto{\pgfqpoint{5.392532in}{1.612009in}}%
\pgfpathlineto{\pgfqpoint{5.396396in}{1.615244in}}%
\pgfpathclose%
\pgfusepath{fill}%
\end{pgfscope}%
\begin{pgfscope}%
\pgfpathrectangle{\pgfqpoint{3.525000in}{0.100000in}}{\pgfqpoint{2.862985in}{1.832778in}}%
\pgfusepath{clip}%
\pgfsetbuttcap%
\pgfsetmiterjoin%
\definecolor{currentfill}{rgb}{0.093810,0.775286,0.615379}%
\pgfsetfillcolor{currentfill}%
\pgfsetlinewidth{0.000000pt}%
\definecolor{currentstroke}{rgb}{0.000000,0.000000,0.000000}%
\pgfsetstrokecolor{currentstroke}%
\pgfsetstrokeopacity{0.000000}%
\pgfsetdash{}{0pt}%
\pgfpathmoveto{\pgfqpoint{5.442832in}{1.482495in}}%
\pgfpathlineto{\pgfqpoint{5.439727in}{1.485941in}}%
\pgfpathlineto{\pgfqpoint{5.442991in}{1.495641in}}%
\pgfpathlineto{\pgfqpoint{5.433301in}{1.496259in}}%
\pgfpathlineto{\pgfqpoint{5.435870in}{1.504623in}}%
\pgfpathlineto{\pgfqpoint{5.434272in}{1.517943in}}%
\pgfpathlineto{\pgfqpoint{5.425578in}{1.522561in}}%
\pgfpathlineto{\pgfqpoint{5.416480in}{1.531907in}}%
\pgfpathlineto{\pgfqpoint{5.402460in}{1.534643in}}%
\pgfpathlineto{\pgfqpoint{5.388816in}{1.534565in}}%
\pgfpathlineto{\pgfqpoint{5.375397in}{1.541198in}}%
\pgfpathlineto{\pgfqpoint{5.330502in}{1.550866in}}%
\pgfpathlineto{\pgfqpoint{5.325596in}{1.561110in}}%
\pgfpathlineto{\pgfqpoint{5.316845in}{1.564606in}}%
\pgfpathlineto{\pgfqpoint{5.333375in}{1.572449in}}%
\pgfpathlineto{\pgfqpoint{5.342713in}{1.582234in}}%
\pgfpathlineto{\pgfqpoint{5.360106in}{1.584888in}}%
\pgfpathlineto{\pgfqpoint{5.370809in}{1.594855in}}%
\pgfpathlineto{\pgfqpoint{5.376425in}{1.595254in}}%
\pgfpathlineto{\pgfqpoint{5.380710in}{1.602362in}}%
\pgfpathlineto{\pgfqpoint{5.391157in}{1.610783in}}%
\pgfpathlineto{\pgfqpoint{5.392067in}{1.604831in}}%
\pgfpathlineto{\pgfqpoint{5.396774in}{1.603602in}}%
\pgfpathlineto{\pgfqpoint{5.400316in}{1.596504in}}%
\pgfpathlineto{\pgfqpoint{5.400954in}{1.584403in}}%
\pgfpathlineto{\pgfqpoint{5.411599in}{1.591634in}}%
\pgfpathlineto{\pgfqpoint{5.424002in}{1.593145in}}%
\pgfpathlineto{\pgfqpoint{5.434590in}{1.589356in}}%
\pgfpathlineto{\pgfqpoint{5.448949in}{1.569658in}}%
\pgfpathlineto{\pgfqpoint{5.464631in}{1.572806in}}%
\pgfpathlineto{\pgfqpoint{5.471003in}{1.567535in}}%
\pgfpathlineto{\pgfqpoint{5.481252in}{1.567064in}}%
\pgfpathlineto{\pgfqpoint{5.488061in}{1.576518in}}%
\pgfpathlineto{\pgfqpoint{5.500982in}{1.584876in}}%
\pgfpathlineto{\pgfqpoint{5.513400in}{1.587493in}}%
\pgfpathlineto{\pgfqpoint{5.528808in}{1.587766in}}%
\pgfpathlineto{\pgfqpoint{5.540068in}{1.594222in}}%
\pgfpathlineto{\pgfqpoint{5.549258in}{1.591248in}}%
\pgfpathlineto{\pgfqpoint{5.551215in}{1.577594in}}%
\pgfpathlineto{\pgfqpoint{5.570930in}{1.575456in}}%
\pgfpathlineto{\pgfqpoint{5.581638in}{1.581847in}}%
\pgfpathlineto{\pgfqpoint{5.589061in}{1.567449in}}%
\pgfpathlineto{\pgfqpoint{5.603225in}{1.550799in}}%
\pgfpathlineto{\pgfqpoint{5.583401in}{1.550895in}}%
\pgfpathlineto{\pgfqpoint{5.577137in}{1.548837in}}%
\pgfpathlineto{\pgfqpoint{5.568551in}{1.551511in}}%
\pgfpathlineto{\pgfqpoint{5.567974in}{1.539985in}}%
\pgfpathlineto{\pgfqpoint{5.552384in}{1.548993in}}%
\pgfpathlineto{\pgfqpoint{5.532335in}{1.551746in}}%
\pgfpathlineto{\pgfqpoint{5.526816in}{1.542973in}}%
\pgfpathlineto{\pgfqpoint{5.500584in}{1.538707in}}%
\pgfpathlineto{\pgfqpoint{5.497566in}{1.531276in}}%
\pgfpathlineto{\pgfqpoint{5.479312in}{1.529046in}}%
\pgfpathlineto{\pgfqpoint{5.473794in}{1.521492in}}%
\pgfpathlineto{\pgfqpoint{5.464139in}{1.519466in}}%
\pgfpathlineto{\pgfqpoint{5.456427in}{1.501552in}}%
\pgfpathlineto{\pgfqpoint{5.446663in}{1.484203in}}%
\pgfpathlineto{\pgfqpoint{5.442832in}{1.482495in}}%
\pgfpathclose%
\pgfusepath{fill}%
\end{pgfscope}%
\begin{pgfscope}%
\pgfpathrectangle{\pgfqpoint{3.525000in}{0.100000in}}{\pgfqpoint{2.862985in}{1.832778in}}%
\pgfusepath{clip}%
\pgfsetbuttcap%
\pgfsetmiterjoin%
\definecolor{currentfill}{rgb}{0.093810,0.775286,0.615379}%
\pgfsetfillcolor{currentfill}%
\pgfsetlinewidth{0.000000pt}%
\definecolor{currentstroke}{rgb}{0.000000,0.000000,0.000000}%
\pgfsetstrokecolor{currentstroke}%
\pgfsetstrokeopacity{0.000000}%
\pgfsetdash{}{0pt}%
\pgfpathmoveto{\pgfqpoint{5.499029in}{1.274203in}}%
\pgfpathlineto{\pgfqpoint{5.508339in}{1.284002in}}%
\pgfpathlineto{\pgfqpoint{5.512580in}{1.298209in}}%
\pgfpathlineto{\pgfqpoint{5.517623in}{1.306443in}}%
\pgfpathlineto{\pgfqpoint{5.520716in}{1.317649in}}%
\pgfpathlineto{\pgfqpoint{5.521678in}{1.339993in}}%
\pgfpathlineto{\pgfqpoint{5.517021in}{1.361407in}}%
\pgfpathlineto{\pgfqpoint{5.501633in}{1.394197in}}%
\pgfpathlineto{\pgfqpoint{5.505800in}{1.404507in}}%
\pgfpathlineto{\pgfqpoint{5.500372in}{1.418758in}}%
\pgfpathlineto{\pgfqpoint{5.509684in}{1.438466in}}%
\pgfpathlineto{\pgfqpoint{5.508166in}{1.460442in}}%
\pgfpathlineto{\pgfqpoint{5.514653in}{1.463207in}}%
\pgfpathlineto{\pgfqpoint{5.515468in}{1.473664in}}%
\pgfpathlineto{\pgfqpoint{5.526982in}{1.480375in}}%
\pgfpathlineto{\pgfqpoint{5.533494in}{1.479291in}}%
\pgfpathlineto{\pgfqpoint{5.535324in}{1.468077in}}%
\pgfpathlineto{\pgfqpoint{5.540407in}{1.467638in}}%
\pgfpathlineto{\pgfqpoint{5.545044in}{1.483813in}}%
\pgfpathlineto{\pgfqpoint{5.543468in}{1.496393in}}%
\pgfpathlineto{\pgfqpoint{5.546483in}{1.503617in}}%
\pgfpathlineto{\pgfqpoint{5.562785in}{1.511080in}}%
\pgfpathlineto{\pgfqpoint{5.555356in}{1.513735in}}%
\pgfpathlineto{\pgfqpoint{5.552936in}{1.520146in}}%
\pgfpathlineto{\pgfqpoint{5.558288in}{1.531404in}}%
\pgfpathlineto{\pgfqpoint{5.568841in}{1.535292in}}%
\pgfpathlineto{\pgfqpoint{5.581095in}{1.528650in}}%
\pgfpathlineto{\pgfqpoint{5.592648in}{1.528567in}}%
\pgfpathlineto{\pgfqpoint{5.598015in}{1.520808in}}%
\pgfpathlineto{\pgfqpoint{5.606106in}{1.521347in}}%
\pgfpathlineto{\pgfqpoint{5.631083in}{1.510561in}}%
\pgfpathlineto{\pgfqpoint{5.636074in}{1.500120in}}%
\pgfpathlineto{\pgfqpoint{5.630624in}{1.496504in}}%
\pgfpathlineto{\pgfqpoint{5.632303in}{1.488679in}}%
\pgfpathlineto{\pgfqpoint{5.637687in}{1.485198in}}%
\pgfpathlineto{\pgfqpoint{5.640688in}{1.475579in}}%
\pgfpathlineto{\pgfqpoint{5.640268in}{1.452144in}}%
\pgfpathlineto{\pgfqpoint{5.633159in}{1.446546in}}%
\pgfpathlineto{\pgfqpoint{5.631564in}{1.434236in}}%
\pgfpathlineto{\pgfqpoint{5.618377in}{1.422859in}}%
\pgfpathlineto{\pgfqpoint{5.619159in}{1.409089in}}%
\pgfpathlineto{\pgfqpoint{5.630717in}{1.404265in}}%
\pgfpathlineto{\pgfqpoint{5.643756in}{1.421446in}}%
\pgfpathlineto{\pgfqpoint{5.644833in}{1.427678in}}%
\pgfpathlineto{\pgfqpoint{5.661094in}{1.438035in}}%
\pgfpathlineto{\pgfqpoint{5.671435in}{1.433238in}}%
\pgfpathlineto{\pgfqpoint{5.677924in}{1.422395in}}%
\pgfpathlineto{\pgfqpoint{5.688393in}{1.384749in}}%
\pgfpathlineto{\pgfqpoint{5.693940in}{1.372839in}}%
\pgfpathlineto{\pgfqpoint{5.692203in}{1.367915in}}%
\pgfpathlineto{\pgfqpoint{5.692372in}{1.351150in}}%
\pgfpathlineto{\pgfqpoint{5.682281in}{1.352759in}}%
\pgfpathlineto{\pgfqpoint{5.676581in}{1.340231in}}%
\pgfpathlineto{\pgfqpoint{5.675791in}{1.331714in}}%
\pgfpathlineto{\pgfqpoint{5.668169in}{1.326246in}}%
\pgfpathlineto{\pgfqpoint{5.666475in}{1.309636in}}%
\pgfpathlineto{\pgfqpoint{5.655390in}{1.288653in}}%
\pgfpathlineto{\pgfqpoint{5.594752in}{1.279623in}}%
\pgfpathlineto{\pgfqpoint{5.594393in}{1.283592in}}%
\pgfpathlineto{\pgfqpoint{5.553831in}{1.279317in}}%
\pgfpathlineto{\pgfqpoint{5.499029in}{1.274203in}}%
\pgfpathclose%
\pgfusepath{fill}%
\end{pgfscope}%
\begin{pgfscope}%
\pgfpathrectangle{\pgfqpoint{3.525000in}{0.100000in}}{\pgfqpoint{2.862985in}{1.832778in}}%
\pgfusepath{clip}%
\pgfsetbuttcap%
\pgfsetmiterjoin%
\definecolor{currentfill}{rgb}{0.081200,0.764245,0.988466}%
\pgfsetfillcolor{currentfill}%
\pgfsetlinewidth{0.000000pt}%
\definecolor{currentstroke}{rgb}{0.000000,0.000000,0.000000}%
\pgfsetstrokecolor{currentstroke}%
\pgfsetstrokeopacity{0.000000}%
\pgfsetdash{}{0pt}%
\pgfpathmoveto{\pgfqpoint{3.967548in}{0.431818in}}%
\pgfpathlineto{\pgfqpoint{3.966337in}{0.430752in}}%
\pgfpathlineto{\pgfqpoint{3.960781in}{0.432173in}}%
\pgfpathlineto{\pgfqpoint{3.961196in}{0.434969in}}%
\pgfpathlineto{\pgfqpoint{3.964341in}{0.436249in}}%
\pgfpathlineto{\pgfqpoint{3.968824in}{0.442694in}}%
\pgfpathlineto{\pgfqpoint{3.969569in}{0.447051in}}%
\pgfpathlineto{\pgfqpoint{3.972090in}{0.448192in}}%
\pgfpathlineto{\pgfqpoint{3.976190in}{0.448197in}}%
\pgfpathlineto{\pgfqpoint{3.980477in}{0.462333in}}%
\pgfpathlineto{\pgfqpoint{3.982157in}{0.463496in}}%
\pgfpathlineto{\pgfqpoint{3.980651in}{0.466260in}}%
\pgfpathlineto{\pgfqpoint{3.977125in}{0.463464in}}%
\pgfpathlineto{\pgfqpoint{3.966766in}{0.467233in}}%
\pgfpathlineto{\pgfqpoint{3.960583in}{0.473582in}}%
\pgfpathlineto{\pgfqpoint{3.963280in}{0.476479in}}%
\pgfpathlineto{\pgfqpoint{3.962053in}{0.480565in}}%
\pgfpathlineto{\pgfqpoint{3.962809in}{0.485392in}}%
\pgfpathlineto{\pgfqpoint{3.961479in}{0.491035in}}%
\pgfpathlineto{\pgfqpoint{3.961266in}{0.496684in}}%
\pgfpathlineto{\pgfqpoint{3.967217in}{0.495739in}}%
\pgfpathlineto{\pgfqpoint{3.968395in}{0.497404in}}%
\pgfpathlineto{\pgfqpoint{3.971355in}{0.497243in}}%
\pgfpathlineto{\pgfqpoint{3.974764in}{0.491217in}}%
\pgfpathlineto{\pgfqpoint{3.978051in}{0.486806in}}%
\pgfpathlineto{\pgfqpoint{3.975418in}{0.483935in}}%
\pgfpathlineto{\pgfqpoint{3.979970in}{0.482833in}}%
\pgfpathlineto{\pgfqpoint{3.981007in}{0.484241in}}%
\pgfpathlineto{\pgfqpoint{3.978573in}{0.489956in}}%
\pgfpathlineto{\pgfqpoint{3.973447in}{0.494692in}}%
\pgfpathlineto{\pgfqpoint{3.974074in}{0.496676in}}%
\pgfpathlineto{\pgfqpoint{3.969984in}{0.501691in}}%
\pgfpathlineto{\pgfqpoint{3.974161in}{0.502710in}}%
\pgfpathlineto{\pgfqpoint{3.970542in}{0.513448in}}%
\pgfpathlineto{\pgfqpoint{3.973484in}{0.514936in}}%
\pgfpathlineto{\pgfqpoint{3.974202in}{0.517511in}}%
\pgfpathlineto{\pgfqpoint{3.972288in}{0.520569in}}%
\pgfpathlineto{\pgfqpoint{3.976239in}{0.521335in}}%
\pgfpathlineto{\pgfqpoint{3.976948in}{0.523648in}}%
\pgfpathlineto{\pgfqpoint{3.981660in}{0.520176in}}%
\pgfpathlineto{\pgfqpoint{3.982707in}{0.523941in}}%
\pgfpathlineto{\pgfqpoint{3.985183in}{0.525597in}}%
\pgfpathlineto{\pgfqpoint{3.997233in}{0.527901in}}%
\pgfpathlineto{\pgfqpoint{4.000932in}{0.527086in}}%
\pgfpathlineto{\pgfqpoint{4.004031in}{0.531659in}}%
\pgfpathlineto{\pgfqpoint{4.011384in}{0.534658in}}%
\pgfpathlineto{\pgfqpoint{4.017047in}{0.533521in}}%
\pgfpathlineto{\pgfqpoint{4.019631in}{0.529538in}}%
\pgfpathlineto{\pgfqpoint{4.019495in}{0.523212in}}%
\pgfpathlineto{\pgfqpoint{4.022394in}{0.521778in}}%
\pgfpathlineto{\pgfqpoint{4.028306in}{0.521978in}}%
\pgfpathlineto{\pgfqpoint{4.035878in}{0.523624in}}%
\pgfpathlineto{\pgfqpoint{4.035785in}{0.521601in}}%
\pgfpathlineto{\pgfqpoint{4.045357in}{0.515832in}}%
\pgfpathlineto{\pgfqpoint{4.052052in}{0.517379in}}%
\pgfpathlineto{\pgfqpoint{4.054022in}{0.519381in}}%
\pgfpathlineto{\pgfqpoint{4.056133in}{0.524508in}}%
\pgfpathlineto{\pgfqpoint{4.058750in}{0.527792in}}%
\pgfpathlineto{\pgfqpoint{4.058984in}{0.532712in}}%
\pgfpathlineto{\pgfqpoint{4.056587in}{0.534647in}}%
\pgfpathlineto{\pgfqpoint{4.060136in}{0.536396in}}%
\pgfpathlineto{\pgfqpoint{4.066098in}{0.533768in}}%
\pgfpathlineto{\pgfqpoint{4.068040in}{0.535390in}}%
\pgfpathlineto{\pgfqpoint{4.068382in}{0.542352in}}%
\pgfpathlineto{\pgfqpoint{4.064528in}{0.540897in}}%
\pgfpathlineto{\pgfqpoint{4.061855in}{0.543309in}}%
\pgfpathlineto{\pgfqpoint{4.057516in}{0.544812in}}%
\pgfpathlineto{\pgfqpoint{4.050629in}{0.544561in}}%
\pgfpathlineto{\pgfqpoint{4.044030in}{0.546848in}}%
\pgfpathlineto{\pgfqpoint{4.043699in}{0.552545in}}%
\pgfpathlineto{\pgfqpoint{4.037797in}{0.558597in}}%
\pgfpathlineto{\pgfqpoint{4.030313in}{0.561138in}}%
\pgfpathlineto{\pgfqpoint{4.023780in}{0.572853in}}%
\pgfpathlineto{\pgfqpoint{4.024316in}{0.577189in}}%
\pgfpathlineto{\pgfqpoint{4.027833in}{0.579165in}}%
\pgfpathlineto{\pgfqpoint{4.027315in}{0.583269in}}%
\pgfpathlineto{\pgfqpoint{4.028132in}{0.587352in}}%
\pgfpathlineto{\pgfqpoint{4.030915in}{0.584557in}}%
\pgfpathlineto{\pgfqpoint{4.034590in}{0.585723in}}%
\pgfpathlineto{\pgfqpoint{4.034407in}{0.589270in}}%
\pgfpathlineto{\pgfqpoint{4.029311in}{0.596390in}}%
\pgfpathlineto{\pgfqpoint{4.027820in}{0.603734in}}%
\pgfpathlineto{\pgfqpoint{4.032085in}{0.604294in}}%
\pgfpathlineto{\pgfqpoint{4.038468in}{0.603843in}}%
\pgfpathlineto{\pgfqpoint{4.040478in}{0.601387in}}%
\pgfpathlineto{\pgfqpoint{4.046201in}{0.602266in}}%
\pgfpathlineto{\pgfqpoint{4.051520in}{0.601112in}}%
\pgfpathlineto{\pgfqpoint{4.055115in}{0.599083in}}%
\pgfpathlineto{\pgfqpoint{4.057604in}{0.600010in}}%
\pgfpathlineto{\pgfqpoint{4.064212in}{0.597983in}}%
\pgfpathlineto{\pgfqpoint{4.064273in}{0.596187in}}%
\pgfpathlineto{\pgfqpoint{4.069305in}{0.596794in}}%
\pgfpathlineto{\pgfqpoint{4.076053in}{0.592414in}}%
\pgfpathlineto{\pgfqpoint{4.075791in}{0.589506in}}%
\pgfpathlineto{\pgfqpoint{4.072580in}{0.588164in}}%
\pgfpathlineto{\pgfqpoint{4.069167in}{0.584972in}}%
\pgfpathlineto{\pgfqpoint{4.068755in}{0.580588in}}%
\pgfpathlineto{\pgfqpoint{4.076106in}{0.574997in}}%
\pgfpathlineto{\pgfqpoint{4.074930in}{0.572318in}}%
\pgfpathlineto{\pgfqpoint{4.080161in}{0.570260in}}%
\pgfpathlineto{\pgfqpoint{4.081445in}{0.566900in}}%
\pgfpathlineto{\pgfqpoint{4.088083in}{0.569570in}}%
\pgfpathlineto{\pgfqpoint{4.091007in}{0.566071in}}%
\pgfpathlineto{\pgfqpoint{4.092356in}{0.568189in}}%
\pgfpathlineto{\pgfqpoint{4.088383in}{0.575083in}}%
\pgfpathlineto{\pgfqpoint{4.089795in}{0.577148in}}%
\pgfpathlineto{\pgfqpoint{4.090421in}{0.582594in}}%
\pgfpathlineto{\pgfqpoint{4.088838in}{0.584991in}}%
\pgfpathlineto{\pgfqpoint{4.090025in}{0.588211in}}%
\pgfpathlineto{\pgfqpoint{4.093915in}{0.587908in}}%
\pgfpathlineto{\pgfqpoint{4.093110in}{0.582429in}}%
\pgfpathlineto{\pgfqpoint{4.090627in}{0.580602in}}%
\pgfpathlineto{\pgfqpoint{4.090803in}{0.573231in}}%
\pgfpathlineto{\pgfqpoint{4.095258in}{0.572350in}}%
\pgfpathlineto{\pgfqpoint{4.095312in}{0.566870in}}%
\pgfpathlineto{\pgfqpoint{4.100178in}{0.563315in}}%
\pgfpathlineto{\pgfqpoint{4.102415in}{0.565557in}}%
\pgfpathlineto{\pgfqpoint{4.099970in}{0.572380in}}%
\pgfpathlineto{\pgfqpoint{4.097442in}{0.574106in}}%
\pgfpathlineto{\pgfqpoint{4.094955in}{0.572970in}}%
\pgfpathlineto{\pgfqpoint{4.092767in}{0.574420in}}%
\pgfpathlineto{\pgfqpoint{4.093149in}{0.580725in}}%
\pgfpathlineto{\pgfqpoint{4.096800in}{0.583304in}}%
\pgfpathlineto{\pgfqpoint{4.100395in}{0.583066in}}%
\pgfpathlineto{\pgfqpoint{4.098905in}{0.586634in}}%
\pgfpathlineto{\pgfqpoint{4.093518in}{0.589366in}}%
\pgfpathlineto{\pgfqpoint{4.086383in}{0.600253in}}%
\pgfpathlineto{\pgfqpoint{4.089944in}{0.605475in}}%
\pgfpathlineto{\pgfqpoint{4.091926in}{0.610393in}}%
\pgfpathlineto{\pgfqpoint{4.092082in}{0.620269in}}%
\pgfpathlineto{\pgfqpoint{4.091319in}{0.628601in}}%
\pgfpathlineto{\pgfqpoint{4.089142in}{0.634072in}}%
\pgfpathlineto{\pgfqpoint{4.089576in}{0.639541in}}%
\pgfpathlineto{\pgfqpoint{4.095144in}{0.644285in}}%
\pgfpathlineto{\pgfqpoint{4.100681in}{0.649922in}}%
\pgfpathlineto{\pgfqpoint{4.114564in}{0.638405in}}%
\pgfpathlineto{\pgfqpoint{4.123640in}{0.637556in}}%
\pgfpathlineto{\pgfqpoint{4.131132in}{0.640424in}}%
\pgfpathlineto{\pgfqpoint{4.137708in}{0.644201in}}%
\pgfpathlineto{\pgfqpoint{4.139105in}{0.646000in}}%
\pgfpathlineto{\pgfqpoint{4.155340in}{0.650103in}}%
\pgfpathlineto{\pgfqpoint{4.156749in}{0.647098in}}%
\pgfpathlineto{\pgfqpoint{4.162266in}{0.644638in}}%
\pgfpathlineto{\pgfqpoint{4.172793in}{0.645392in}}%
\pgfpathlineto{\pgfqpoint{4.174716in}{0.644720in}}%
\pgfpathlineto{\pgfqpoint{4.179448in}{0.646515in}}%
\pgfpathlineto{\pgfqpoint{4.181350in}{0.643411in}}%
\pgfpathlineto{\pgfqpoint{4.185894in}{0.640575in}}%
\pgfpathlineto{\pgfqpoint{4.188167in}{0.640896in}}%
\pgfpathlineto{\pgfqpoint{4.192566in}{0.637488in}}%
\pgfpathlineto{\pgfqpoint{4.200000in}{0.638115in}}%
\pgfpathlineto{\pgfqpoint{4.207425in}{0.640874in}}%
\pgfpathlineto{\pgfqpoint{4.213489in}{0.632051in}}%
\pgfpathlineto{\pgfqpoint{4.213019in}{0.630071in}}%
\pgfpathlineto{\pgfqpoint{4.206328in}{0.627796in}}%
\pgfpathlineto{\pgfqpoint{4.208115in}{0.624355in}}%
\pgfpathlineto{\pgfqpoint{4.214275in}{0.627976in}}%
\pgfpathlineto{\pgfqpoint{4.216790in}{0.627100in}}%
\pgfpathlineto{\pgfqpoint{4.217860in}{0.622768in}}%
\pgfpathlineto{\pgfqpoint{4.214183in}{0.620936in}}%
\pgfpathlineto{\pgfqpoint{4.216809in}{0.615461in}}%
\pgfpathlineto{\pgfqpoint{4.220623in}{0.616386in}}%
\pgfpathlineto{\pgfqpoint{4.226157in}{0.613560in}}%
\pgfpathlineto{\pgfqpoint{4.231450in}{0.605718in}}%
\pgfpathlineto{\pgfqpoint{4.226793in}{0.603540in}}%
\pgfpathlineto{\pgfqpoint{4.230907in}{0.597761in}}%
\pgfpathlineto{\pgfqpoint{4.227441in}{0.596410in}}%
\pgfpathlineto{\pgfqpoint{4.231753in}{0.590922in}}%
\pgfpathlineto{\pgfqpoint{4.237012in}{0.591506in}}%
\pgfpathlineto{\pgfqpoint{4.239781in}{0.589037in}}%
\pgfpathlineto{\pgfqpoint{4.246785in}{0.585396in}}%
\pgfpathlineto{\pgfqpoint{4.255820in}{0.572470in}}%
\pgfpathlineto{\pgfqpoint{4.255616in}{0.569379in}}%
\pgfpathlineto{\pgfqpoint{4.269166in}{0.559477in}}%
\pgfpathlineto{\pgfqpoint{4.269458in}{0.555975in}}%
\pgfpathlineto{\pgfqpoint{4.273069in}{0.550735in}}%
\pgfpathlineto{\pgfqpoint{4.284163in}{0.548115in}}%
\pgfpathlineto{\pgfqpoint{4.288277in}{0.546252in}}%
\pgfpathlineto{\pgfqpoint{4.292148in}{0.540476in}}%
\pgfpathlineto{\pgfqpoint{4.292646in}{0.535415in}}%
\pgfpathlineto{\pgfqpoint{4.295973in}{0.531291in}}%
\pgfpathlineto{\pgfqpoint{4.296972in}{0.526640in}}%
\pgfpathlineto{\pgfqpoint{4.300090in}{0.524941in}}%
\pgfpathlineto{\pgfqpoint{4.284646in}{0.497067in}}%
\pgfpathlineto{\pgfqpoint{4.252590in}{0.439200in}}%
\pgfpathlineto{\pgfqpoint{4.205557in}{0.354280in}}%
\pgfpathlineto{\pgfqpoint{4.187143in}{0.321025in}}%
\pgfpathlineto{\pgfqpoint{4.191065in}{0.316559in}}%
\pgfpathlineto{\pgfqpoint{4.192793in}{0.317984in}}%
\pgfpathlineto{\pgfqpoint{4.196296in}{0.312785in}}%
\pgfpathlineto{\pgfqpoint{4.201159in}{0.314293in}}%
\pgfpathlineto{\pgfqpoint{4.207614in}{0.311178in}}%
\pgfpathlineto{\pgfqpoint{4.203433in}{0.306310in}}%
\pgfpathlineto{\pgfqpoint{4.206662in}{0.299741in}}%
\pgfpathlineto{\pgfqpoint{4.205995in}{0.296543in}}%
\pgfpathlineto{\pgfqpoint{4.211145in}{0.279677in}}%
\pgfpathlineto{\pgfqpoint{4.208966in}{0.272007in}}%
\pgfpathlineto{\pgfqpoint{4.218805in}{0.273906in}}%
\pgfpathlineto{\pgfqpoint{4.221221in}{0.272662in}}%
\pgfpathlineto{\pgfqpoint{4.223844in}{0.274685in}}%
\pgfpathlineto{\pgfqpoint{4.225709in}{0.278509in}}%
\pgfpathlineto{\pgfqpoint{4.228381in}{0.280701in}}%
\pgfpathlineto{\pgfqpoint{4.239784in}{0.280424in}}%
\pgfpathlineto{\pgfqpoint{4.242361in}{0.273064in}}%
\pgfpathlineto{\pgfqpoint{4.240149in}{0.266603in}}%
\pgfpathlineto{\pgfqpoint{4.242662in}{0.264471in}}%
\pgfpathlineto{\pgfqpoint{4.243983in}{0.260825in}}%
\pgfpathlineto{\pgfqpoint{4.243695in}{0.253967in}}%
\pgfpathlineto{\pgfqpoint{4.247048in}{0.249051in}}%
\pgfpathlineto{\pgfqpoint{4.248886in}{0.241346in}}%
\pgfpathlineto{\pgfqpoint{4.247942in}{0.237184in}}%
\pgfpathlineto{\pgfqpoint{4.249080in}{0.232966in}}%
\pgfpathlineto{\pgfqpoint{4.250722in}{0.208539in}}%
\pgfpathlineto{\pgfqpoint{4.248370in}{0.206615in}}%
\pgfpathlineto{\pgfqpoint{4.251472in}{0.203957in}}%
\pgfpathlineto{\pgfqpoint{4.249034in}{0.200723in}}%
\pgfpathlineto{\pgfqpoint{4.251233in}{0.198087in}}%
\pgfpathlineto{\pgfqpoint{4.249817in}{0.193586in}}%
\pgfpathlineto{\pgfqpoint{4.252984in}{0.192366in}}%
\pgfpathlineto{\pgfqpoint{4.257058in}{0.185522in}}%
\pgfpathlineto{\pgfqpoint{4.260085in}{0.183460in}}%
\pgfpathlineto{\pgfqpoint{4.260877in}{0.180445in}}%
\pgfpathlineto{\pgfqpoint{4.262631in}{0.179139in}}%
\pgfpathlineto{\pgfqpoint{4.262265in}{0.176636in}}%
\pgfpathlineto{\pgfqpoint{4.266046in}{0.174822in}}%
\pgfpathlineto{\pgfqpoint{4.265112in}{0.170006in}}%
\pgfpathlineto{\pgfqpoint{4.261858in}{0.167483in}}%
\pgfpathlineto{\pgfqpoint{4.260380in}{0.162795in}}%
\pgfpathlineto{\pgfqpoint{4.259988in}{0.156518in}}%
\pgfpathlineto{\pgfqpoint{4.258180in}{0.155917in}}%
\pgfpathlineto{\pgfqpoint{4.252271in}{0.150415in}}%
\pgfpathlineto{\pgfqpoint{4.247019in}{0.148962in}}%
\pgfpathlineto{\pgfqpoint{4.244513in}{0.151181in}}%
\pgfpathlineto{\pgfqpoint{4.246935in}{0.157071in}}%
\pgfpathlineto{\pgfqpoint{4.246195in}{0.159962in}}%
\pgfpathlineto{\pgfqpoint{4.249665in}{0.161426in}}%
\pgfpathlineto{\pgfqpoint{4.252909in}{0.169947in}}%
\pgfpathlineto{\pgfqpoint{4.251539in}{0.177225in}}%
\pgfpathlineto{\pgfqpoint{4.249346in}{0.178943in}}%
\pgfpathlineto{\pgfqpoint{4.242012in}{0.178387in}}%
\pgfpathlineto{\pgfqpoint{4.243490in}{0.176486in}}%
\pgfpathlineto{\pgfqpoint{4.238181in}{0.171028in}}%
\pgfpathlineto{\pgfqpoint{4.236471in}{0.174063in}}%
\pgfpathlineto{\pgfqpoint{4.236916in}{0.178027in}}%
\pgfpathlineto{\pgfqpoint{4.239963in}{0.177954in}}%
\pgfpathlineto{\pgfqpoint{4.242677in}{0.180866in}}%
\pgfpathlineto{\pgfqpoint{4.244282in}{0.185118in}}%
\pgfpathlineto{\pgfqpoint{4.250331in}{0.183918in}}%
\pgfpathlineto{\pgfqpoint{4.245488in}{0.186190in}}%
\pgfpathlineto{\pgfqpoint{4.245915in}{0.189233in}}%
\pgfpathlineto{\pgfqpoint{4.243899in}{0.190607in}}%
\pgfpathlineto{\pgfqpoint{4.243074in}{0.194752in}}%
\pgfpathlineto{\pgfqpoint{4.244505in}{0.196911in}}%
\pgfpathlineto{\pgfqpoint{4.241093in}{0.203617in}}%
\pgfpathlineto{\pgfqpoint{4.241381in}{0.207883in}}%
\pgfpathlineto{\pgfqpoint{4.237383in}{0.211682in}}%
\pgfpathlineto{\pgfqpoint{4.235220in}{0.214756in}}%
\pgfpathlineto{\pgfqpoint{4.238207in}{0.217685in}}%
\pgfpathlineto{\pgfqpoint{4.239313in}{0.222063in}}%
\pgfpathlineto{\pgfqpoint{4.238326in}{0.224937in}}%
\pgfpathlineto{\pgfqpoint{4.239481in}{0.228366in}}%
\pgfpathlineto{\pgfqpoint{4.238009in}{0.239543in}}%
\pgfpathlineto{\pgfqpoint{4.235807in}{0.244614in}}%
\pgfpathlineto{\pgfqpoint{4.233288in}{0.246519in}}%
\pgfpathlineto{\pgfqpoint{4.234360in}{0.253116in}}%
\pgfpathlineto{\pgfqpoint{4.233604in}{0.258268in}}%
\pgfpathlineto{\pgfqpoint{4.234974in}{0.261572in}}%
\pgfpathlineto{\pgfqpoint{4.232140in}{0.262035in}}%
\pgfpathlineto{\pgfqpoint{4.231380in}{0.254009in}}%
\pgfpathlineto{\pgfqpoint{4.228196in}{0.245055in}}%
\pgfpathlineto{\pgfqpoint{4.225639in}{0.247094in}}%
\pgfpathlineto{\pgfqpoint{4.225297in}{0.250037in}}%
\pgfpathlineto{\pgfqpoint{4.222168in}{0.254057in}}%
\pgfpathlineto{\pgfqpoint{4.223758in}{0.256688in}}%
\pgfpathlineto{\pgfqpoint{4.223264in}{0.262613in}}%
\pgfpathlineto{\pgfqpoint{4.219215in}{0.260530in}}%
\pgfpathlineto{\pgfqpoint{4.220009in}{0.257316in}}%
\pgfpathlineto{\pgfqpoint{4.219247in}{0.253236in}}%
\pgfpathlineto{\pgfqpoint{4.214702in}{0.253244in}}%
\pgfpathlineto{\pgfqpoint{4.212665in}{0.255773in}}%
\pgfpathlineto{\pgfqpoint{4.209590in}{0.256255in}}%
\pgfpathlineto{\pgfqpoint{4.207521in}{0.258978in}}%
\pgfpathlineto{\pgfqpoint{4.203961in}{0.266454in}}%
\pgfpathlineto{\pgfqpoint{4.202807in}{0.272259in}}%
\pgfpathlineto{\pgfqpoint{4.203546in}{0.274726in}}%
\pgfpathlineto{\pgfqpoint{4.201822in}{0.280160in}}%
\pgfpathlineto{\pgfqpoint{4.198967in}{0.283305in}}%
\pgfpathlineto{\pgfqpoint{4.194029in}{0.291344in}}%
\pgfpathlineto{\pgfqpoint{4.190447in}{0.298326in}}%
\pgfpathlineto{\pgfqpoint{4.194102in}{0.298431in}}%
\pgfpathlineto{\pgfqpoint{4.196252in}{0.299897in}}%
\pgfpathlineto{\pgfqpoint{4.196675in}{0.304435in}}%
\pgfpathlineto{\pgfqpoint{4.186334in}{0.304983in}}%
\pgfpathlineto{\pgfqpoint{4.177539in}{0.314807in}}%
\pgfpathlineto{\pgfqpoint{4.179852in}{0.317392in}}%
\pgfpathlineto{\pgfqpoint{4.175599in}{0.317833in}}%
\pgfpathlineto{\pgfqpoint{4.167152in}{0.326967in}}%
\pgfpathlineto{\pgfqpoint{4.153770in}{0.331547in}}%
\pgfpathlineto{\pgfqpoint{4.151968in}{0.336915in}}%
\pgfpathlineto{\pgfqpoint{4.149191in}{0.339711in}}%
\pgfpathlineto{\pgfqpoint{4.149065in}{0.342336in}}%
\pgfpathlineto{\pgfqpoint{4.146976in}{0.344286in}}%
\pgfpathlineto{\pgfqpoint{4.149717in}{0.347544in}}%
\pgfpathlineto{\pgfqpoint{4.143465in}{0.347858in}}%
\pgfpathlineto{\pgfqpoint{4.141876in}{0.352075in}}%
\pgfpathlineto{\pgfqpoint{4.139030in}{0.353870in}}%
\pgfpathlineto{\pgfqpoint{4.144729in}{0.356037in}}%
\pgfpathlineto{\pgfqpoint{4.141765in}{0.359902in}}%
\pgfpathlineto{\pgfqpoint{4.136311in}{0.361524in}}%
\pgfpathlineto{\pgfqpoint{4.137160in}{0.368663in}}%
\pgfpathlineto{\pgfqpoint{4.135596in}{0.371190in}}%
\pgfpathlineto{\pgfqpoint{4.132780in}{0.372970in}}%
\pgfpathlineto{\pgfqpoint{4.130372in}{0.371185in}}%
\pgfpathlineto{\pgfqpoint{4.128950in}{0.372864in}}%
\pgfpathlineto{\pgfqpoint{4.124800in}{0.373345in}}%
\pgfpathlineto{\pgfqpoint{4.125134in}{0.379823in}}%
\pgfpathlineto{\pgfqpoint{4.119119in}{0.375448in}}%
\pgfpathlineto{\pgfqpoint{4.119246in}{0.366430in}}%
\pgfpathlineto{\pgfqpoint{4.112788in}{0.364702in}}%
\pgfpathlineto{\pgfqpoint{4.108572in}{0.361101in}}%
\pgfpathlineto{\pgfqpoint{4.105182in}{0.360342in}}%
\pgfpathlineto{\pgfqpoint{4.101514in}{0.363971in}}%
\pgfpathlineto{\pgfqpoint{4.098248in}{0.363490in}}%
\pgfpathlineto{\pgfqpoint{4.098854in}{0.366805in}}%
\pgfpathlineto{\pgfqpoint{4.094211in}{0.365020in}}%
\pgfpathlineto{\pgfqpoint{4.089909in}{0.364820in}}%
\pgfpathlineto{\pgfqpoint{4.084874in}{0.363010in}}%
\pgfpathlineto{\pgfqpoint{4.074656in}{0.364028in}}%
\pgfpathlineto{\pgfqpoint{4.071376in}{0.363358in}}%
\pgfpathlineto{\pgfqpoint{4.069009in}{0.366339in}}%
\pgfpathlineto{\pgfqpoint{4.064640in}{0.366637in}}%
\pgfpathlineto{\pgfqpoint{4.063731in}{0.370500in}}%
\pgfpathlineto{\pgfqpoint{4.066335in}{0.372604in}}%
\pgfpathlineto{\pgfqpoint{4.069390in}{0.372406in}}%
\pgfpathlineto{\pgfqpoint{4.071940in}{0.369537in}}%
\pgfpathlineto{\pgfqpoint{4.073619in}{0.373944in}}%
\pgfpathlineto{\pgfqpoint{4.071276in}{0.378827in}}%
\pgfpathlineto{\pgfqpoint{4.076579in}{0.383125in}}%
\pgfpathlineto{\pgfqpoint{4.081893in}{0.384709in}}%
\pgfpathlineto{\pgfqpoint{4.085593in}{0.387292in}}%
\pgfpathlineto{\pgfqpoint{4.087990in}{0.390076in}}%
\pgfpathlineto{\pgfqpoint{4.089327in}{0.394638in}}%
\pgfpathlineto{\pgfqpoint{4.093552in}{0.393396in}}%
\pgfpathlineto{\pgfqpoint{4.102714in}{0.394034in}}%
\pgfpathlineto{\pgfqpoint{4.104680in}{0.389032in}}%
\pgfpathlineto{\pgfqpoint{4.107894in}{0.388814in}}%
\pgfpathlineto{\pgfqpoint{4.113742in}{0.381156in}}%
\pgfpathlineto{\pgfqpoint{4.113753in}{0.383947in}}%
\pgfpathlineto{\pgfqpoint{4.111710in}{0.384855in}}%
\pgfpathlineto{\pgfqpoint{4.107971in}{0.393996in}}%
\pgfpathlineto{\pgfqpoint{4.109790in}{0.395238in}}%
\pgfpathlineto{\pgfqpoint{4.104797in}{0.400007in}}%
\pgfpathlineto{\pgfqpoint{4.100334in}{0.400936in}}%
\pgfpathlineto{\pgfqpoint{4.096081in}{0.399287in}}%
\pgfpathlineto{\pgfqpoint{4.088798in}{0.400571in}}%
\pgfpathlineto{\pgfqpoint{4.085344in}{0.398261in}}%
\pgfpathlineto{\pgfqpoint{4.077478in}{0.396476in}}%
\pgfpathlineto{\pgfqpoint{4.076759in}{0.393907in}}%
\pgfpathlineto{\pgfqpoint{4.070716in}{0.393207in}}%
\pgfpathlineto{\pgfqpoint{4.069165in}{0.390270in}}%
\pgfpathlineto{\pgfqpoint{4.065709in}{0.388119in}}%
\pgfpathlineto{\pgfqpoint{4.061533in}{0.388489in}}%
\pgfpathlineto{\pgfqpoint{4.059448in}{0.386408in}}%
\pgfpathlineto{\pgfqpoint{4.056818in}{0.386506in}}%
\pgfpathlineto{\pgfqpoint{4.054304in}{0.388378in}}%
\pgfpathlineto{\pgfqpoint{4.049909in}{0.388044in}}%
\pgfpathlineto{\pgfqpoint{4.048874in}{0.386318in}}%
\pgfpathlineto{\pgfqpoint{4.044258in}{0.388090in}}%
\pgfpathlineto{\pgfqpoint{4.040339in}{0.383518in}}%
\pgfpathlineto{\pgfqpoint{4.044097in}{0.379095in}}%
\pgfpathlineto{\pgfqpoint{4.045543in}{0.375457in}}%
\pgfpathlineto{\pgfqpoint{4.044210in}{0.372376in}}%
\pgfpathlineto{\pgfqpoint{4.038870in}{0.370250in}}%
\pgfpathlineto{\pgfqpoint{4.035776in}{0.371980in}}%
\pgfpathlineto{\pgfqpoint{4.031882in}{0.371043in}}%
\pgfpathlineto{\pgfqpoint{4.031413in}{0.368343in}}%
\pgfpathlineto{\pgfqpoint{4.026673in}{0.369298in}}%
\pgfpathlineto{\pgfqpoint{4.027722in}{0.366600in}}%
\pgfpathlineto{\pgfqpoint{4.020731in}{0.366031in}}%
\pgfpathlineto{\pgfqpoint{4.016205in}{0.369456in}}%
\pgfpathlineto{\pgfqpoint{4.013787in}{0.367278in}}%
\pgfpathlineto{\pgfqpoint{4.007354in}{0.369597in}}%
\pgfpathlineto{\pgfqpoint{4.005752in}{0.367248in}}%
\pgfpathlineto{\pgfqpoint{4.002563in}{0.366203in}}%
\pgfpathlineto{\pgfqpoint{4.000005in}{0.368957in}}%
\pgfpathlineto{\pgfqpoint{3.991257in}{0.368272in}}%
\pgfpathlineto{\pgfqpoint{3.991266in}{0.364211in}}%
\pgfpathlineto{\pgfqpoint{3.987752in}{0.363136in}}%
\pgfpathlineto{\pgfqpoint{3.985018in}{0.365224in}}%
\pgfpathlineto{\pgfqpoint{3.981404in}{0.364064in}}%
\pgfpathlineto{\pgfqpoint{3.976727in}{0.363787in}}%
\pgfpathlineto{\pgfqpoint{3.969759in}{0.366724in}}%
\pgfpathlineto{\pgfqpoint{3.965769in}{0.363696in}}%
\pgfpathlineto{\pgfqpoint{3.960063in}{0.367743in}}%
\pgfpathlineto{\pgfqpoint{3.958197in}{0.366209in}}%
\pgfpathlineto{\pgfqpoint{3.956887in}{0.362246in}}%
\pgfpathlineto{\pgfqpoint{3.953218in}{0.359922in}}%
\pgfpathlineto{\pgfqpoint{3.947236in}{0.362582in}}%
\pgfpathlineto{\pgfqpoint{3.937326in}{0.369153in}}%
\pgfpathlineto{\pgfqpoint{3.934438in}{0.369737in}}%
\pgfpathlineto{\pgfqpoint{3.932808in}{0.368368in}}%
\pgfpathlineto{\pgfqpoint{3.928449in}{0.369299in}}%
\pgfpathlineto{\pgfqpoint{3.926322in}{0.371344in}}%
\pgfpathlineto{\pgfqpoint{3.921969in}{0.372446in}}%
\pgfpathlineto{\pgfqpoint{3.916023in}{0.372510in}}%
\pgfpathlineto{\pgfqpoint{3.913662in}{0.374950in}}%
\pgfpathlineto{\pgfqpoint{3.918187in}{0.377474in}}%
\pgfpathlineto{\pgfqpoint{3.917131in}{0.379966in}}%
\pgfpathlineto{\pgfqpoint{3.914022in}{0.379406in}}%
\pgfpathlineto{\pgfqpoint{3.912470in}{0.377447in}}%
\pgfpathlineto{\pgfqpoint{3.905192in}{0.374639in}}%
\pgfpathlineto{\pgfqpoint{3.901512in}{0.375594in}}%
\pgfpathlineto{\pgfqpoint{3.899964in}{0.381297in}}%
\pgfpathlineto{\pgfqpoint{3.895373in}{0.378845in}}%
\pgfpathlineto{\pgfqpoint{3.894039in}{0.386073in}}%
\pgfpathlineto{\pgfqpoint{3.891864in}{0.384210in}}%
\pgfpathlineto{\pgfqpoint{3.892228in}{0.381429in}}%
\pgfpathlineto{\pgfqpoint{3.887226in}{0.382163in}}%
\pgfpathlineto{\pgfqpoint{3.892429in}{0.387010in}}%
\pgfpathlineto{\pgfqpoint{3.897702in}{0.384130in}}%
\pgfpathlineto{\pgfqpoint{3.898833in}{0.385576in}}%
\pgfpathlineto{\pgfqpoint{3.904380in}{0.383843in}}%
\pgfpathlineto{\pgfqpoint{3.904033in}{0.385816in}}%
\pgfpathlineto{\pgfqpoint{3.919168in}{0.386469in}}%
\pgfpathlineto{\pgfqpoint{3.926551in}{0.383391in}}%
\pgfpathlineto{\pgfqpoint{3.929408in}{0.380355in}}%
\pgfpathlineto{\pgfqpoint{3.928561in}{0.377705in}}%
\pgfpathlineto{\pgfqpoint{3.931429in}{0.376702in}}%
\pgfpathlineto{\pgfqpoint{3.938558in}{0.380885in}}%
\pgfpathlineto{\pgfqpoint{3.947779in}{0.381153in}}%
\pgfpathlineto{\pgfqpoint{3.956155in}{0.378631in}}%
\pgfpathlineto{\pgfqpoint{3.960899in}{0.379054in}}%
\pgfpathlineto{\pgfqpoint{3.962665in}{0.375142in}}%
\pgfpathlineto{\pgfqpoint{3.965318in}{0.379455in}}%
\pgfpathlineto{\pgfqpoint{3.966950in}{0.380389in}}%
\pgfpathlineto{\pgfqpoint{3.974498in}{0.382264in}}%
\pgfpathlineto{\pgfqpoint{3.977304in}{0.381118in}}%
\pgfpathlineto{\pgfqpoint{3.980338in}{0.382479in}}%
\pgfpathlineto{\pgfqpoint{3.983958in}{0.381767in}}%
\pgfpathlineto{\pgfqpoint{3.984812in}{0.383352in}}%
\pgfpathlineto{\pgfqpoint{3.992363in}{0.390625in}}%
\pgfpathlineto{\pgfqpoint{3.995561in}{0.391703in}}%
\pgfpathlineto{\pgfqpoint{3.997451in}{0.395629in}}%
\pgfpathlineto{\pgfqpoint{4.007481in}{0.398668in}}%
\pgfpathlineto{\pgfqpoint{4.008737in}{0.400987in}}%
\pgfpathlineto{\pgfqpoint{3.994522in}{0.404566in}}%
\pgfpathlineto{\pgfqpoint{3.995081in}{0.408415in}}%
\pgfpathlineto{\pgfqpoint{3.993619in}{0.413499in}}%
\pgfpathlineto{\pgfqpoint{3.990309in}{0.413077in}}%
\pgfpathlineto{\pgfqpoint{3.988057in}{0.406479in}}%
\pgfpathlineto{\pgfqpoint{3.985443in}{0.405936in}}%
\pgfpathlineto{\pgfqpoint{3.983742in}{0.408071in}}%
\pgfpathlineto{\pgfqpoint{3.985726in}{0.417413in}}%
\pgfpathlineto{\pgfqpoint{3.983750in}{0.421328in}}%
\pgfpathlineto{\pgfqpoint{3.979922in}{0.425943in}}%
\pgfpathlineto{\pgfqpoint{3.981762in}{0.429465in}}%
\pgfpathlineto{\pgfqpoint{3.973683in}{0.429591in}}%
\pgfpathlineto{\pgfqpoint{3.973867in}{0.430958in}}%
\pgfpathlineto{\pgfqpoint{3.969046in}{0.432871in}}%
\pgfpathlineto{\pgfqpoint{3.967548in}{0.431818in}}%
\pgfpathclose%
\pgfusepath{fill}%
\end{pgfscope}%
\begin{pgfscope}%
\pgfpathrectangle{\pgfqpoint{3.525000in}{0.100000in}}{\pgfqpoint{2.862985in}{1.832778in}}%
\pgfusepath{clip}%
\pgfsetbuttcap%
\pgfsetmiterjoin%
\definecolor{currentfill}{rgb}{0.081200,0.764245,0.988466}%
\pgfsetfillcolor{currentfill}%
\pgfsetlinewidth{0.000000pt}%
\definecolor{currentstroke}{rgb}{0.000000,0.000000,0.000000}%
\pgfsetstrokecolor{currentstroke}%
\pgfsetstrokeopacity{0.000000}%
\pgfsetdash{}{0pt}%
\pgfpathmoveto{\pgfqpoint{3.953131in}{0.499862in}}%
\pgfpathlineto{\pgfqpoint{3.952132in}{0.498106in}}%
\pgfpathlineto{\pgfqpoint{3.954775in}{0.494368in}}%
\pgfpathlineto{\pgfqpoint{3.950597in}{0.490871in}}%
\pgfpathlineto{\pgfqpoint{3.950713in}{0.487863in}}%
\pgfpathlineto{\pgfqpoint{3.948824in}{0.487153in}}%
\pgfpathlineto{\pgfqpoint{3.942331in}{0.491639in}}%
\pgfpathlineto{\pgfqpoint{3.937703in}{0.501458in}}%
\pgfpathlineto{\pgfqpoint{3.937393in}{0.504295in}}%
\pgfpathlineto{\pgfqpoint{3.938583in}{0.507691in}}%
\pgfpathlineto{\pgfqpoint{3.943598in}{0.502496in}}%
\pgfpathlineto{\pgfqpoint{3.945132in}{0.503498in}}%
\pgfpathlineto{\pgfqpoint{3.949386in}{0.502545in}}%
\pgfpathlineto{\pgfqpoint{3.953131in}{0.499862in}}%
\pgfpathclose%
\pgfusepath{fill}%
\end{pgfscope}%
\begin{pgfscope}%
\pgfpathrectangle{\pgfqpoint{3.525000in}{0.100000in}}{\pgfqpoint{2.862985in}{1.832778in}}%
\pgfusepath{clip}%
\pgfsetbuttcap%
\pgfsetmiterjoin%
\definecolor{currentfill}{rgb}{0.081200,0.764245,0.988466}%
\pgfsetfillcolor{currentfill}%
\pgfsetlinewidth{0.000000pt}%
\definecolor{currentstroke}{rgb}{0.000000,0.000000,0.000000}%
\pgfsetstrokecolor{currentstroke}%
\pgfsetstrokeopacity{0.000000}%
\pgfsetdash{}{0pt}%
\pgfpathmoveto{\pgfqpoint{3.876126in}{0.386109in}}%
\pgfpathlineto{\pgfqpoint{3.871562in}{0.385368in}}%
\pgfpathlineto{\pgfqpoint{3.868260in}{0.386895in}}%
\pgfpathlineto{\pgfqpoint{3.866816in}{0.389380in}}%
\pgfpathlineto{\pgfqpoint{3.868453in}{0.392945in}}%
\pgfpathlineto{\pgfqpoint{3.872096in}{0.392002in}}%
\pgfpathlineto{\pgfqpoint{3.878086in}{0.394117in}}%
\pgfpathlineto{\pgfqpoint{3.880184in}{0.391196in}}%
\pgfpathlineto{\pgfqpoint{3.882217in}{0.391735in}}%
\pgfpathlineto{\pgfqpoint{3.887141in}{0.389917in}}%
\pgfpathlineto{\pgfqpoint{3.889268in}{0.387608in}}%
\pgfpathlineto{\pgfqpoint{3.886672in}{0.381567in}}%
\pgfpathlineto{\pgfqpoint{3.884090in}{0.380286in}}%
\pgfpathlineto{\pgfqpoint{3.881769in}{0.380976in}}%
\pgfpathlineto{\pgfqpoint{3.876126in}{0.386109in}}%
\pgfpathclose%
\pgfusepath{fill}%
\end{pgfscope}%
\begin{pgfscope}%
\pgfpathrectangle{\pgfqpoint{3.525000in}{0.100000in}}{\pgfqpoint{2.862985in}{1.832778in}}%
\pgfusepath{clip}%
\pgfsetbuttcap%
\pgfsetmiterjoin%
\definecolor{currentfill}{rgb}{0.081200,0.764245,0.988466}%
\pgfsetfillcolor{currentfill}%
\pgfsetlinewidth{0.000000pt}%
\definecolor{currentstroke}{rgb}{0.000000,0.000000,0.000000}%
\pgfsetstrokecolor{currentstroke}%
\pgfsetstrokeopacity{0.000000}%
\pgfsetdash{}{0pt}%
\pgfpathmoveto{\pgfqpoint{3.835323in}{0.391598in}}%
\pgfpathlineto{\pgfqpoint{3.828857in}{0.394682in}}%
\pgfpathlineto{\pgfqpoint{3.822873in}{0.395600in}}%
\pgfpathlineto{\pgfqpoint{3.820987in}{0.396913in}}%
\pgfpathlineto{\pgfqpoint{3.823135in}{0.399499in}}%
\pgfpathlineto{\pgfqpoint{3.826946in}{0.396269in}}%
\pgfpathlineto{\pgfqpoint{3.830142in}{0.396714in}}%
\pgfpathlineto{\pgfqpoint{3.835341in}{0.398951in}}%
\pgfpathlineto{\pgfqpoint{3.835761in}{0.401832in}}%
\pgfpathlineto{\pgfqpoint{3.842146in}{0.400406in}}%
\pgfpathlineto{\pgfqpoint{3.840651in}{0.396069in}}%
\pgfpathlineto{\pgfqpoint{3.844719in}{0.395751in}}%
\pgfpathlineto{\pgfqpoint{3.841306in}{0.390592in}}%
\pgfpathlineto{\pgfqpoint{3.838001in}{0.392264in}}%
\pgfpathlineto{\pgfqpoint{3.835323in}{0.391598in}}%
\pgfpathclose%
\pgfusepath{fill}%
\end{pgfscope}%
\begin{pgfscope}%
\pgfpathrectangle{\pgfqpoint{3.525000in}{0.100000in}}{\pgfqpoint{2.862985in}{1.832778in}}%
\pgfusepath{clip}%
\pgfsetbuttcap%
\pgfsetmiterjoin%
\definecolor{currentfill}{rgb}{0.081200,0.764245,0.988466}%
\pgfsetfillcolor{currentfill}%
\pgfsetlinewidth{0.000000pt}%
\definecolor{currentstroke}{rgb}{0.000000,0.000000,0.000000}%
\pgfsetstrokecolor{currentstroke}%
\pgfsetstrokeopacity{0.000000}%
\pgfsetdash{}{0pt}%
\pgfpathmoveto{\pgfqpoint{3.823962in}{0.402261in}}%
\pgfpathlineto{\pgfqpoint{3.821586in}{0.400969in}}%
\pgfpathlineto{\pgfqpoint{3.815110in}{0.402955in}}%
\pgfpathlineto{\pgfqpoint{3.810271in}{0.401805in}}%
\pgfpathlineto{\pgfqpoint{3.807813in}{0.405073in}}%
\pgfpathlineto{\pgfqpoint{3.808885in}{0.406569in}}%
\pgfpathlineto{\pgfqpoint{3.812552in}{0.406819in}}%
\pgfpathlineto{\pgfqpoint{3.815357in}{0.405880in}}%
\pgfpathlineto{\pgfqpoint{3.818414in}{0.407367in}}%
\pgfpathlineto{\pgfqpoint{3.823100in}{0.405378in}}%
\pgfpathlineto{\pgfqpoint{3.823962in}{0.402261in}}%
\pgfpathclose%
\pgfusepath{fill}%
\end{pgfscope}%
\begin{pgfscope}%
\pgfpathrectangle{\pgfqpoint{3.525000in}{0.100000in}}{\pgfqpoint{2.862985in}{1.832778in}}%
\pgfusepath{clip}%
\pgfsetbuttcap%
\pgfsetmiterjoin%
\definecolor{currentfill}{rgb}{0.081200,0.764245,0.988466}%
\pgfsetfillcolor{currentfill}%
\pgfsetlinewidth{0.000000pt}%
\definecolor{currentstroke}{rgb}{0.000000,0.000000,0.000000}%
\pgfsetstrokecolor{currentstroke}%
\pgfsetstrokeopacity{0.000000}%
\pgfsetdash{}{0pt}%
\pgfpathmoveto{\pgfqpoint{3.747628in}{0.450357in}}%
\pgfpathlineto{\pgfqpoint{3.747456in}{0.447153in}}%
\pgfpathlineto{\pgfqpoint{3.742729in}{0.445139in}}%
\pgfpathlineto{\pgfqpoint{3.738412in}{0.450180in}}%
\pgfpathlineto{\pgfqpoint{3.743253in}{0.451860in}}%
\pgfpathlineto{\pgfqpoint{3.747628in}{0.450357in}}%
\pgfpathclose%
\pgfusepath{fill}%
\end{pgfscope}%
\begin{pgfscope}%
\pgfpathrectangle{\pgfqpoint{3.525000in}{0.100000in}}{\pgfqpoint{2.862985in}{1.832778in}}%
\pgfusepath{clip}%
\pgfsetbuttcap%
\pgfsetmiterjoin%
\definecolor{currentfill}{rgb}{0.081200,0.764245,0.988466}%
\pgfsetfillcolor{currentfill}%
\pgfsetlinewidth{0.000000pt}%
\definecolor{currentstroke}{rgb}{0.000000,0.000000,0.000000}%
\pgfsetstrokecolor{currentstroke}%
\pgfsetstrokeopacity{0.000000}%
\pgfsetdash{}{0pt}%
\pgfpathmoveto{\pgfqpoint{3.708638in}{0.468371in}}%
\pgfpathlineto{\pgfqpoint{3.709642in}{0.469764in}}%
\pgfpathlineto{\pgfqpoint{3.715440in}{0.469524in}}%
\pgfpathlineto{\pgfqpoint{3.716458in}{0.471597in}}%
\pgfpathlineto{\pgfqpoint{3.719002in}{0.469788in}}%
\pgfpathlineto{\pgfqpoint{3.717533in}{0.466308in}}%
\pgfpathlineto{\pgfqpoint{3.715372in}{0.466014in}}%
\pgfpathlineto{\pgfqpoint{3.710207in}{0.467106in}}%
\pgfpathlineto{\pgfqpoint{3.708638in}{0.468371in}}%
\pgfpathclose%
\pgfusepath{fill}%
\end{pgfscope}%
\begin{pgfscope}%
\pgfpathrectangle{\pgfqpoint{3.525000in}{0.100000in}}{\pgfqpoint{2.862985in}{1.832778in}}%
\pgfusepath{clip}%
\pgfsetbuttcap%
\pgfsetmiterjoin%
\definecolor{currentfill}{rgb}{0.081200,0.764245,0.988466}%
\pgfsetfillcolor{currentfill}%
\pgfsetlinewidth{0.000000pt}%
\definecolor{currentstroke}{rgb}{0.000000,0.000000,0.000000}%
\pgfsetstrokecolor{currentstroke}%
\pgfsetstrokeopacity{0.000000}%
\pgfsetdash{}{0pt}%
\pgfpathmoveto{\pgfqpoint{3.700415in}{0.479156in}}%
\pgfpathlineto{\pgfqpoint{3.700063in}{0.482294in}}%
\pgfpathlineto{\pgfqpoint{3.703088in}{0.482062in}}%
\pgfpathlineto{\pgfqpoint{3.702878in}{0.486470in}}%
\pgfpathlineto{\pgfqpoint{3.705921in}{0.484091in}}%
\pgfpathlineto{\pgfqpoint{3.704868in}{0.480653in}}%
\pgfpathlineto{\pgfqpoint{3.700415in}{0.479156in}}%
\pgfpathclose%
\pgfusepath{fill}%
\end{pgfscope}%
\begin{pgfscope}%
\pgfpathrectangle{\pgfqpoint{3.525000in}{0.100000in}}{\pgfqpoint{2.862985in}{1.832778in}}%
\pgfusepath{clip}%
\pgfsetbuttcap%
\pgfsetmiterjoin%
\definecolor{currentfill}{rgb}{0.081200,0.764245,0.988466}%
\pgfsetfillcolor{currentfill}%
\pgfsetlinewidth{0.000000pt}%
\definecolor{currentstroke}{rgb}{0.000000,0.000000,0.000000}%
\pgfsetstrokecolor{currentstroke}%
\pgfsetstrokeopacity{0.000000}%
\pgfsetdash{}{0pt}%
\pgfpathmoveto{\pgfqpoint{3.966980in}{0.591016in}}%
\pgfpathlineto{\pgfqpoint{3.962154in}{0.592058in}}%
\pgfpathlineto{\pgfqpoint{3.961118in}{0.595283in}}%
\pgfpathlineto{\pgfqpoint{3.962776in}{0.598507in}}%
\pgfpathlineto{\pgfqpoint{3.966491in}{0.600190in}}%
\pgfpathlineto{\pgfqpoint{3.970744in}{0.591916in}}%
\pgfpathlineto{\pgfqpoint{3.974601in}{0.591548in}}%
\pgfpathlineto{\pgfqpoint{3.977479in}{0.588995in}}%
\pgfpathlineto{\pgfqpoint{3.977422in}{0.585482in}}%
\pgfpathlineto{\pgfqpoint{3.975318in}{0.583797in}}%
\pgfpathlineto{\pgfqpoint{3.978740in}{0.573408in}}%
\pgfpathlineto{\pgfqpoint{3.982484in}{0.569595in}}%
\pgfpathlineto{\pgfqpoint{3.978656in}{0.568385in}}%
\pgfpathlineto{\pgfqpoint{3.975557in}{0.572663in}}%
\pgfpathlineto{\pgfqpoint{3.968823in}{0.571338in}}%
\pgfpathlineto{\pgfqpoint{3.970918in}{0.575572in}}%
\pgfpathlineto{\pgfqpoint{3.969943in}{0.577258in}}%
\pgfpathlineto{\pgfqpoint{3.970514in}{0.581671in}}%
\pgfpathlineto{\pgfqpoint{3.969180in}{0.589648in}}%
\pgfpathlineto{\pgfqpoint{3.966980in}{0.591016in}}%
\pgfpathclose%
\pgfusepath{fill}%
\end{pgfscope}%
\begin{pgfscope}%
\pgfpathrectangle{\pgfqpoint{3.525000in}{0.100000in}}{\pgfqpoint{2.862985in}{1.832778in}}%
\pgfusepath{clip}%
\pgfsetbuttcap%
\pgfsetmiterjoin%
\definecolor{currentfill}{rgb}{0.081200,0.764245,0.988466}%
\pgfsetfillcolor{currentfill}%
\pgfsetlinewidth{0.000000pt}%
\definecolor{currentstroke}{rgb}{0.000000,0.000000,0.000000}%
\pgfsetstrokecolor{currentstroke}%
\pgfsetstrokeopacity{0.000000}%
\pgfsetdash{}{0pt}%
\pgfpathmoveto{\pgfqpoint{4.136519in}{0.353334in}}%
\pgfpathlineto{\pgfqpoint{4.129523in}{0.353412in}}%
\pgfpathlineto{\pgfqpoint{4.130062in}{0.356924in}}%
\pgfpathlineto{\pgfqpoint{4.132862in}{0.358180in}}%
\pgfpathlineto{\pgfqpoint{4.136519in}{0.353334in}}%
\pgfpathclose%
\pgfusepath{fill}%
\end{pgfscope}%
\begin{pgfscope}%
\pgfpathrectangle{\pgfqpoint{3.525000in}{0.100000in}}{\pgfqpoint{2.862985in}{1.832778in}}%
\pgfusepath{clip}%
\pgfsetbuttcap%
\pgfsetmiterjoin%
\definecolor{currentfill}{rgb}{0.081200,0.764245,0.988466}%
\pgfsetfillcolor{currentfill}%
\pgfsetlinewidth{0.000000pt}%
\definecolor{currentstroke}{rgb}{0.000000,0.000000,0.000000}%
\pgfsetstrokecolor{currentstroke}%
\pgfsetstrokeopacity{0.000000}%
\pgfsetdash{}{0pt}%
\pgfpathmoveto{\pgfqpoint{4.126812in}{0.356514in}}%
\pgfpathlineto{\pgfqpoint{4.121800in}{0.355301in}}%
\pgfpathlineto{\pgfqpoint{4.117038in}{0.353224in}}%
\pgfpathlineto{\pgfqpoint{4.115567in}{0.355671in}}%
\pgfpathlineto{\pgfqpoint{4.123757in}{0.357328in}}%
\pgfpathlineto{\pgfqpoint{4.126812in}{0.356514in}}%
\pgfpathclose%
\pgfusepath{fill}%
\end{pgfscope}%
\begin{pgfscope}%
\pgfpathrectangle{\pgfqpoint{3.525000in}{0.100000in}}{\pgfqpoint{2.862985in}{1.832778in}}%
\pgfusepath{clip}%
\pgfsetbuttcap%
\pgfsetmiterjoin%
\definecolor{currentfill}{rgb}{0.081200,0.764245,0.988466}%
\pgfsetfillcolor{currentfill}%
\pgfsetlinewidth{0.000000pt}%
\definecolor{currentstroke}{rgb}{0.000000,0.000000,0.000000}%
\pgfsetstrokecolor{currentstroke}%
\pgfsetstrokeopacity{0.000000}%
\pgfsetdash{}{0pt}%
\pgfpathmoveto{\pgfqpoint{4.049990in}{0.354028in}}%
\pgfpathlineto{\pgfqpoint{4.050182in}{0.351599in}}%
\pgfpathlineto{\pgfqpoint{4.047294in}{0.350090in}}%
\pgfpathlineto{\pgfqpoint{4.039235in}{0.353552in}}%
\pgfpathlineto{\pgfqpoint{4.037986in}{0.352159in}}%
\pgfpathlineto{\pgfqpoint{4.036036in}{0.358747in}}%
\pgfpathlineto{\pgfqpoint{4.040074in}{0.359664in}}%
\pgfpathlineto{\pgfqpoint{4.041940in}{0.356957in}}%
\pgfpathlineto{\pgfqpoint{4.045992in}{0.360345in}}%
\pgfpathlineto{\pgfqpoint{4.049990in}{0.354028in}}%
\pgfpathclose%
\pgfusepath{fill}%
\end{pgfscope}%
\begin{pgfscope}%
\pgfpathrectangle{\pgfqpoint{3.525000in}{0.100000in}}{\pgfqpoint{2.862985in}{1.832778in}}%
\pgfusepath{clip}%
\pgfsetbuttcap%
\pgfsetmiterjoin%
\definecolor{currentfill}{rgb}{0.081200,0.764245,0.988466}%
\pgfsetfillcolor{currentfill}%
\pgfsetlinewidth{0.000000pt}%
\definecolor{currentstroke}{rgb}{0.000000,0.000000,0.000000}%
\pgfsetstrokecolor{currentstroke}%
\pgfsetstrokeopacity{0.000000}%
\pgfsetdash{}{0pt}%
\pgfpathmoveto{\pgfqpoint{4.231774in}{0.219412in}}%
\pgfpathlineto{\pgfqpoint{4.225017in}{0.216896in}}%
\pgfpathlineto{\pgfqpoint{4.221437in}{0.216745in}}%
\pgfpathlineto{\pgfqpoint{4.224350in}{0.222082in}}%
\pgfpathlineto{\pgfqpoint{4.226840in}{0.224073in}}%
\pgfpathlineto{\pgfqpoint{4.226760in}{0.229600in}}%
\pgfpathlineto{\pgfqpoint{4.229465in}{0.239578in}}%
\pgfpathlineto{\pgfqpoint{4.231957in}{0.241473in}}%
\pgfpathlineto{\pgfqpoint{4.237710in}{0.238729in}}%
\pgfpathlineto{\pgfqpoint{4.236870in}{0.237158in}}%
\pgfpathlineto{\pgfqpoint{4.237306in}{0.229604in}}%
\pgfpathlineto{\pgfqpoint{4.235440in}{0.227566in}}%
\pgfpathlineto{\pgfqpoint{4.234640in}{0.222201in}}%
\pgfpathlineto{\pgfqpoint{4.231774in}{0.219412in}}%
\pgfpathclose%
\pgfusepath{fill}%
\end{pgfscope}%
\begin{pgfscope}%
\pgfpathrectangle{\pgfqpoint{3.525000in}{0.100000in}}{\pgfqpoint{2.862985in}{1.832778in}}%
\pgfusepath{clip}%
\pgfsetbuttcap%
\pgfsetmiterjoin%
\definecolor{currentfill}{rgb}{0.081200,0.764245,0.988466}%
\pgfsetfillcolor{currentfill}%
\pgfsetlinewidth{0.000000pt}%
\definecolor{currentstroke}{rgb}{0.000000,0.000000,0.000000}%
\pgfsetstrokecolor{currentstroke}%
\pgfsetstrokeopacity{0.000000}%
\pgfsetdash{}{0pt}%
\pgfpathmoveto{\pgfqpoint{4.216809in}{0.244146in}}%
\pgfpathlineto{\pgfqpoint{4.220768in}{0.237310in}}%
\pgfpathlineto{\pgfqpoint{4.220034in}{0.235680in}}%
\pgfpathlineto{\pgfqpoint{4.225329in}{0.234022in}}%
\pgfpathlineto{\pgfqpoint{4.223809in}{0.227738in}}%
\pgfpathlineto{\pgfqpoint{4.221503in}{0.228029in}}%
\pgfpathlineto{\pgfqpoint{4.215445in}{0.238982in}}%
\pgfpathlineto{\pgfqpoint{4.216307in}{0.234052in}}%
\pgfpathlineto{\pgfqpoint{4.215275in}{0.231213in}}%
\pgfpathlineto{\pgfqpoint{4.212709in}{0.229950in}}%
\pgfpathlineto{\pgfqpoint{4.211251in}{0.231311in}}%
\pgfpathlineto{\pgfqpoint{4.210437in}{0.237623in}}%
\pgfpathlineto{\pgfqpoint{4.210917in}{0.241717in}}%
\pgfpathlineto{\pgfqpoint{4.209208in}{0.243504in}}%
\pgfpathlineto{\pgfqpoint{4.213722in}{0.248675in}}%
\pgfpathlineto{\pgfqpoint{4.216684in}{0.250209in}}%
\pgfpathlineto{\pgfqpoint{4.217670in}{0.248193in}}%
\pgfpathlineto{\pgfqpoint{4.220760in}{0.249753in}}%
\pgfpathlineto{\pgfqpoint{4.222999in}{0.245516in}}%
\pgfpathlineto{\pgfqpoint{4.227840in}{0.240126in}}%
\pgfpathlineto{\pgfqpoint{4.227179in}{0.237684in}}%
\pgfpathlineto{\pgfqpoint{4.224491in}{0.235024in}}%
\pgfpathlineto{\pgfqpoint{4.222005in}{0.236278in}}%
\pgfpathlineto{\pgfqpoint{4.216809in}{0.244146in}}%
\pgfpathclose%
\pgfusepath{fill}%
\end{pgfscope}%
\begin{pgfscope}%
\pgfpathrectangle{\pgfqpoint{3.525000in}{0.100000in}}{\pgfqpoint{2.862985in}{1.832778in}}%
\pgfusepath{clip}%
\pgfsetbuttcap%
\pgfsetmiterjoin%
\definecolor{currentfill}{rgb}{0.081200,0.764245,0.988466}%
\pgfsetfillcolor{currentfill}%
\pgfsetlinewidth{0.000000pt}%
\definecolor{currentstroke}{rgb}{0.000000,0.000000,0.000000}%
\pgfsetstrokecolor{currentstroke}%
\pgfsetstrokeopacity{0.000000}%
\pgfsetdash{}{0pt}%
\pgfpathmoveto{\pgfqpoint{4.020353in}{0.341071in}}%
\pgfpathlineto{\pgfqpoint{4.015449in}{0.340433in}}%
\pgfpathlineto{\pgfqpoint{4.008542in}{0.337778in}}%
\pgfpathlineto{\pgfqpoint{4.006931in}{0.339221in}}%
\pgfpathlineto{\pgfqpoint{4.014146in}{0.340851in}}%
\pgfpathlineto{\pgfqpoint{4.012212in}{0.342273in}}%
\pgfpathlineto{\pgfqpoint{4.007739in}{0.343315in}}%
\pgfpathlineto{\pgfqpoint{4.006544in}{0.346649in}}%
\pgfpathlineto{\pgfqpoint{4.008795in}{0.349521in}}%
\pgfpathlineto{\pgfqpoint{4.010690in}{0.356112in}}%
\pgfpathlineto{\pgfqpoint{4.017614in}{0.357340in}}%
\pgfpathlineto{\pgfqpoint{4.020966in}{0.354492in}}%
\pgfpathlineto{\pgfqpoint{4.024115in}{0.357248in}}%
\pgfpathlineto{\pgfqpoint{4.028064in}{0.356572in}}%
\pgfpathlineto{\pgfqpoint{4.030902in}{0.351482in}}%
\pgfpathlineto{\pgfqpoint{4.033305in}{0.356006in}}%
\pgfpathlineto{\pgfqpoint{4.040126in}{0.345423in}}%
\pgfpathlineto{\pgfqpoint{4.037662in}{0.344577in}}%
\pgfpathlineto{\pgfqpoint{4.036150in}{0.342283in}}%
\pgfpathlineto{\pgfqpoint{4.039007in}{0.340190in}}%
\pgfpathlineto{\pgfqpoint{4.034576in}{0.338227in}}%
\pgfpathlineto{\pgfqpoint{4.031364in}{0.339308in}}%
\pgfpathlineto{\pgfqpoint{4.025296in}{0.342890in}}%
\pgfpathlineto{\pgfqpoint{4.020353in}{0.341071in}}%
\pgfpathclose%
\pgfusepath{fill}%
\end{pgfscope}%
\begin{pgfscope}%
\pgfpathrectangle{\pgfqpoint{3.525000in}{0.100000in}}{\pgfqpoint{2.862985in}{1.832778in}}%
\pgfusepath{clip}%
\pgfsetbuttcap%
\pgfsetmiterjoin%
\definecolor{currentfill}{rgb}{0.081200,0.764245,0.988466}%
\pgfsetfillcolor{currentfill}%
\pgfsetlinewidth{0.000000pt}%
\definecolor{currentstroke}{rgb}{0.000000,0.000000,0.000000}%
\pgfsetstrokecolor{currentstroke}%
\pgfsetstrokeopacity{0.000000}%
\pgfsetdash{}{0pt}%
\pgfpathmoveto{\pgfqpoint{4.212691in}{0.199220in}}%
\pgfpathlineto{\pgfqpoint{4.211719in}{0.208994in}}%
\pgfpathlineto{\pgfqpoint{4.214340in}{0.212583in}}%
\pgfpathlineto{\pgfqpoint{4.211538in}{0.212832in}}%
\pgfpathlineto{\pgfqpoint{4.212562in}{0.219812in}}%
\pgfpathlineto{\pgfqpoint{4.213770in}{0.223955in}}%
\pgfpathlineto{\pgfqpoint{4.212720in}{0.229522in}}%
\pgfpathlineto{\pgfqpoint{4.215287in}{0.230600in}}%
\pgfpathlineto{\pgfqpoint{4.216124in}{0.232039in}}%
\pgfpathlineto{\pgfqpoint{4.218653in}{0.231493in}}%
\pgfpathlineto{\pgfqpoint{4.220842in}{0.227033in}}%
\pgfpathlineto{\pgfqpoint{4.221249in}{0.222947in}}%
\pgfpathlineto{\pgfqpoint{4.218571in}{0.210709in}}%
\pgfpathlineto{\pgfqpoint{4.212691in}{0.199220in}}%
\pgfpathclose%
\pgfusepath{fill}%
\end{pgfscope}%
\begin{pgfscope}%
\pgfpathrectangle{\pgfqpoint{3.525000in}{0.100000in}}{\pgfqpoint{2.862985in}{1.832778in}}%
\pgfusepath{clip}%
\pgfsetbuttcap%
\pgfsetmiterjoin%
\definecolor{currentfill}{rgb}{0.081200,0.764245,0.988466}%
\pgfsetfillcolor{currentfill}%
\pgfsetlinewidth{0.000000pt}%
\definecolor{currentstroke}{rgb}{0.000000,0.000000,0.000000}%
\pgfsetstrokecolor{currentstroke}%
\pgfsetstrokeopacity{0.000000}%
\pgfsetdash{}{0pt}%
\pgfpathmoveto{\pgfqpoint{4.212110in}{0.229056in}}%
\pgfpathlineto{\pgfqpoint{4.211975in}{0.224290in}}%
\pgfpathlineto{\pgfqpoint{4.209778in}{0.222127in}}%
\pgfpathlineto{\pgfqpoint{4.207238in}{0.222926in}}%
\pgfpathlineto{\pgfqpoint{4.210448in}{0.229848in}}%
\pgfpathlineto{\pgfqpoint{4.212110in}{0.229056in}}%
\pgfpathclose%
\pgfusepath{fill}%
\end{pgfscope}%
\begin{pgfscope}%
\pgfpathrectangle{\pgfqpoint{3.525000in}{0.100000in}}{\pgfqpoint{2.862985in}{1.832778in}}%
\pgfusepath{clip}%
\pgfsetbuttcap%
\pgfsetmiterjoin%
\definecolor{currentfill}{rgb}{0.081200,0.764245,0.988466}%
\pgfsetfillcolor{currentfill}%
\pgfsetlinewidth{0.000000pt}%
\definecolor{currentstroke}{rgb}{0.000000,0.000000,0.000000}%
\pgfsetstrokecolor{currentstroke}%
\pgfsetstrokeopacity{0.000000}%
\pgfsetdash{}{0pt}%
\pgfpathmoveto{\pgfqpoint{4.239076in}{0.207898in}}%
\pgfpathlineto{\pgfqpoint{4.237202in}{0.199923in}}%
\pgfpathlineto{\pgfqpoint{4.232835in}{0.197402in}}%
\pgfpathlineto{\pgfqpoint{4.227107in}{0.199724in}}%
\pgfpathlineto{\pgfqpoint{4.228685in}{0.202824in}}%
\pgfpathlineto{\pgfqpoint{4.228653in}{0.204970in}}%
\pgfpathlineto{\pgfqpoint{4.230783in}{0.208616in}}%
\pgfpathlineto{\pgfqpoint{4.229286in}{0.207882in}}%
\pgfpathlineto{\pgfqpoint{4.230377in}{0.209663in}}%
\pgfpathlineto{\pgfqpoint{4.228304in}{0.213709in}}%
\pgfpathlineto{\pgfqpoint{4.230657in}{0.214455in}}%
\pgfpathlineto{\pgfqpoint{4.236104in}{0.209721in}}%
\pgfpathlineto{\pgfqpoint{4.239076in}{0.207898in}}%
\pgfpathclose%
\pgfusepath{fill}%
\end{pgfscope}%
\begin{pgfscope}%
\pgfpathrectangle{\pgfqpoint{3.525000in}{0.100000in}}{\pgfqpoint{2.862985in}{1.832778in}}%
\pgfusepath{clip}%
\pgfsetbuttcap%
\pgfsetmiterjoin%
\definecolor{currentfill}{rgb}{0.081200,0.764245,0.988466}%
\pgfsetfillcolor{currentfill}%
\pgfsetlinewidth{0.000000pt}%
\definecolor{currentstroke}{rgb}{0.000000,0.000000,0.000000}%
\pgfsetstrokecolor{currentstroke}%
\pgfsetstrokeopacity{0.000000}%
\pgfsetdash{}{0pt}%
\pgfpathmoveto{\pgfqpoint{4.219962in}{0.193796in}}%
\pgfpathlineto{\pgfqpoint{4.217492in}{0.192836in}}%
\pgfpathlineto{\pgfqpoint{4.218902in}{0.201252in}}%
\pgfpathlineto{\pgfqpoint{4.221866in}{0.203216in}}%
\pgfpathlineto{\pgfqpoint{4.221144in}{0.209521in}}%
\pgfpathlineto{\pgfqpoint{4.222348in}{0.212317in}}%
\pgfpathlineto{\pgfqpoint{4.224731in}{0.213365in}}%
\pgfpathlineto{\pgfqpoint{4.227232in}{0.210622in}}%
\pgfpathlineto{\pgfqpoint{4.229379in}{0.206901in}}%
\pgfpathlineto{\pgfqpoint{4.222478in}{0.199610in}}%
\pgfpathlineto{\pgfqpoint{4.219962in}{0.193796in}}%
\pgfpathclose%
\pgfusepath{fill}%
\end{pgfscope}%
\begin{pgfscope}%
\pgfpathrectangle{\pgfqpoint{3.525000in}{0.100000in}}{\pgfqpoint{2.862985in}{1.832778in}}%
\pgfusepath{clip}%
\pgfsetbuttcap%
\pgfsetmiterjoin%
\definecolor{currentfill}{rgb}{0.081200,0.764245,0.988466}%
\pgfsetfillcolor{currentfill}%
\pgfsetlinewidth{0.000000pt}%
\definecolor{currentstroke}{rgb}{0.000000,0.000000,0.000000}%
\pgfsetstrokecolor{currentstroke}%
\pgfsetstrokeopacity{0.000000}%
\pgfsetdash{}{0pt}%
\pgfpathmoveto{\pgfqpoint{4.240172in}{0.202568in}}%
\pgfpathlineto{\pgfqpoint{4.241382in}{0.196760in}}%
\pgfpathlineto{\pgfqpoint{4.235974in}{0.197425in}}%
\pgfpathlineto{\pgfqpoint{4.240172in}{0.202568in}}%
\pgfpathclose%
\pgfusepath{fill}%
\end{pgfscope}%
\begin{pgfscope}%
\pgfpathrectangle{\pgfqpoint{3.525000in}{0.100000in}}{\pgfqpoint{2.862985in}{1.832778in}}%
\pgfusepath{clip}%
\pgfsetbuttcap%
\pgfsetmiterjoin%
\definecolor{currentfill}{rgb}{0.081200,0.764245,0.988466}%
\pgfsetfillcolor{currentfill}%
\pgfsetlinewidth{0.000000pt}%
\definecolor{currentstroke}{rgb}{0.000000,0.000000,0.000000}%
\pgfsetstrokecolor{currentstroke}%
\pgfsetstrokeopacity{0.000000}%
\pgfsetdash{}{0pt}%
\pgfpathmoveto{\pgfqpoint{4.242752in}{0.193694in}}%
\pgfpathlineto{\pgfqpoint{4.243638in}{0.190111in}}%
\pgfpathlineto{\pgfqpoint{4.245335in}{0.188965in}}%
\pgfpathlineto{\pgfqpoint{4.245026in}{0.185593in}}%
\pgfpathlineto{\pgfqpoint{4.242656in}{0.184459in}}%
\pgfpathlineto{\pgfqpoint{4.240732in}{0.189493in}}%
\pgfpathlineto{\pgfqpoint{4.242752in}{0.193694in}}%
\pgfpathclose%
\pgfusepath{fill}%
\end{pgfscope}%
\begin{pgfscope}%
\pgfpathrectangle{\pgfqpoint{3.525000in}{0.100000in}}{\pgfqpoint{2.862985in}{1.832778in}}%
\pgfusepath{clip}%
\pgfsetbuttcap%
\pgfsetmiterjoin%
\definecolor{currentfill}{rgb}{0.081200,0.764245,0.988466}%
\pgfsetfillcolor{currentfill}%
\pgfsetlinewidth{0.000000pt}%
\definecolor{currentstroke}{rgb}{0.000000,0.000000,0.000000}%
\pgfsetstrokecolor{currentstroke}%
\pgfsetstrokeopacity{0.000000}%
\pgfsetdash{}{0pt}%
\pgfpathmoveto{\pgfqpoint{4.238239in}{0.195068in}}%
\pgfpathlineto{\pgfqpoint{4.236768in}{0.190907in}}%
\pgfpathlineto{\pgfqpoint{4.234484in}{0.191060in}}%
\pgfpathlineto{\pgfqpoint{4.233043in}{0.194492in}}%
\pgfpathlineto{\pgfqpoint{4.235390in}{0.196117in}}%
\pgfpathlineto{\pgfqpoint{4.238239in}{0.195068in}}%
\pgfpathclose%
\pgfusepath{fill}%
\end{pgfscope}%
\begin{pgfscope}%
\pgfpathrectangle{\pgfqpoint{3.525000in}{0.100000in}}{\pgfqpoint{2.862985in}{1.832778in}}%
\pgfusepath{clip}%
\pgfsetbuttcap%
\pgfsetmiterjoin%
\definecolor{currentfill}{rgb}{0.081200,0.764245,0.988466}%
\pgfsetfillcolor{currentfill}%
\pgfsetlinewidth{0.000000pt}%
\definecolor{currentstroke}{rgb}{0.000000,0.000000,0.000000}%
\pgfsetstrokecolor{currentstroke}%
\pgfsetstrokeopacity{0.000000}%
\pgfsetdash{}{0pt}%
\pgfpathmoveto{\pgfqpoint{4.232076in}{0.173613in}}%
\pgfpathlineto{\pgfqpoint{4.234114in}{0.171384in}}%
\pgfpathlineto{\pgfqpoint{4.234563in}{0.166686in}}%
\pgfpathlineto{\pgfqpoint{4.235636in}{0.165321in}}%
\pgfpathlineto{\pgfqpoint{4.233884in}{0.161070in}}%
\pgfpathlineto{\pgfqpoint{4.231578in}{0.158676in}}%
\pgfpathlineto{\pgfqpoint{4.230352in}{0.154236in}}%
\pgfpathlineto{\pgfqpoint{4.227832in}{0.161156in}}%
\pgfpathlineto{\pgfqpoint{4.227414in}{0.167439in}}%
\pgfpathlineto{\pgfqpoint{4.226004in}{0.170567in}}%
\pgfpathlineto{\pgfqpoint{4.221147in}{0.173102in}}%
\pgfpathlineto{\pgfqpoint{4.225703in}{0.178390in}}%
\pgfpathlineto{\pgfqpoint{4.222666in}{0.180650in}}%
\pgfpathlineto{\pgfqpoint{4.225341in}{0.182902in}}%
\pgfpathlineto{\pgfqpoint{4.228893in}{0.191549in}}%
\pgfpathlineto{\pgfqpoint{4.225212in}{0.194597in}}%
\pgfpathlineto{\pgfqpoint{4.226630in}{0.197524in}}%
\pgfpathlineto{\pgfqpoint{4.231388in}{0.194928in}}%
\pgfpathlineto{\pgfqpoint{4.231870in}{0.190831in}}%
\pgfpathlineto{\pgfqpoint{4.230070in}{0.188498in}}%
\pgfpathlineto{\pgfqpoint{4.233638in}{0.185341in}}%
\pgfpathlineto{\pgfqpoint{4.234761in}{0.180424in}}%
\pgfpathlineto{\pgfqpoint{4.234609in}{0.173528in}}%
\pgfpathlineto{\pgfqpoint{4.232076in}{0.173613in}}%
\pgfpathclose%
\pgfusepath{fill}%
\end{pgfscope}%
\begin{pgfscope}%
\pgfpathrectangle{\pgfqpoint{3.525000in}{0.100000in}}{\pgfqpoint{2.862985in}{1.832778in}}%
\pgfusepath{clip}%
\pgfsetbuttcap%
\pgfsetmiterjoin%
\definecolor{currentfill}{rgb}{0.081200,0.764245,0.988466}%
\pgfsetfillcolor{currentfill}%
\pgfsetlinewidth{0.000000pt}%
\definecolor{currentstroke}{rgb}{0.000000,0.000000,0.000000}%
\pgfsetstrokecolor{currentstroke}%
\pgfsetstrokeopacity{0.000000}%
\pgfsetdash{}{0pt}%
\pgfpathmoveto{\pgfqpoint{4.240789in}{0.190924in}}%
\pgfpathlineto{\pgfqpoint{4.239679in}{0.188532in}}%
\pgfpathlineto{\pgfqpoint{4.240722in}{0.184648in}}%
\pgfpathlineto{\pgfqpoint{4.238053in}{0.184321in}}%
\pgfpathlineto{\pgfqpoint{4.235255in}{0.186119in}}%
\pgfpathlineto{\pgfqpoint{4.236163in}{0.190080in}}%
\pgfpathlineto{\pgfqpoint{4.240789in}{0.190924in}}%
\pgfpathclose%
\pgfusepath{fill}%
\end{pgfscope}%
\begin{pgfscope}%
\pgfpathrectangle{\pgfqpoint{3.525000in}{0.100000in}}{\pgfqpoint{2.862985in}{1.832778in}}%
\pgfusepath{clip}%
\pgfsetbuttcap%
\pgfsetmiterjoin%
\definecolor{currentfill}{rgb}{0.081200,0.764245,0.988466}%
\pgfsetfillcolor{currentfill}%
\pgfsetlinewidth{0.000000pt}%
\definecolor{currentstroke}{rgb}{0.000000,0.000000,0.000000}%
\pgfsetstrokecolor{currentstroke}%
\pgfsetstrokeopacity{0.000000}%
\pgfsetdash{}{0pt}%
\pgfpathmoveto{\pgfqpoint{4.228351in}{0.191195in}}%
\pgfpathlineto{\pgfqpoint{4.222461in}{0.188553in}}%
\pgfpathlineto{\pgfqpoint{4.220459in}{0.189503in}}%
\pgfpathlineto{\pgfqpoint{4.224476in}{0.192989in}}%
\pgfpathlineto{\pgfqpoint{4.228351in}{0.191195in}}%
\pgfpathclose%
\pgfusepath{fill}%
\end{pgfscope}%
\begin{pgfscope}%
\pgfpathrectangle{\pgfqpoint{3.525000in}{0.100000in}}{\pgfqpoint{2.862985in}{1.832778in}}%
\pgfusepath{clip}%
\pgfsetbuttcap%
\pgfsetmiterjoin%
\definecolor{currentfill}{rgb}{0.081200,0.764245,0.988466}%
\pgfsetfillcolor{currentfill}%
\pgfsetlinewidth{0.000000pt}%
\definecolor{currentstroke}{rgb}{0.000000,0.000000,0.000000}%
\pgfsetstrokecolor{currentstroke}%
\pgfsetstrokeopacity{0.000000}%
\pgfsetdash{}{0pt}%
\pgfpathmoveto{\pgfqpoint{4.240854in}{0.165422in}}%
\pgfpathlineto{\pgfqpoint{4.239449in}{0.169400in}}%
\pgfpathlineto{\pgfqpoint{4.242419in}{0.170378in}}%
\pgfpathlineto{\pgfqpoint{4.243319in}{0.174668in}}%
\pgfpathlineto{\pgfqpoint{4.246538in}{0.174673in}}%
\pgfpathlineto{\pgfqpoint{4.246524in}{0.177660in}}%
\pgfpathlineto{\pgfqpoint{4.250997in}{0.176953in}}%
\pgfpathlineto{\pgfqpoint{4.251760in}{0.168856in}}%
\pgfpathlineto{\pgfqpoint{4.249165in}{0.163694in}}%
\pgfpathlineto{\pgfqpoint{4.245315in}{0.160469in}}%
\pgfpathlineto{\pgfqpoint{4.240854in}{0.165422in}}%
\pgfpathclose%
\pgfusepath{fill}%
\end{pgfscope}%
\begin{pgfscope}%
\pgfpathrectangle{\pgfqpoint{3.525000in}{0.100000in}}{\pgfqpoint{2.862985in}{1.832778in}}%
\pgfusepath{clip}%
\pgfsetbuttcap%
\pgfsetmiterjoin%
\definecolor{currentfill}{rgb}{0.081200,0.764245,0.988466}%
\pgfsetfillcolor{currentfill}%
\pgfsetlinewidth{0.000000pt}%
\definecolor{currentstroke}{rgb}{0.000000,0.000000,0.000000}%
\pgfsetstrokecolor{currentstroke}%
\pgfsetstrokeopacity{0.000000}%
\pgfsetdash{}{0pt}%
\pgfpathmoveto{\pgfqpoint{4.239114in}{0.168482in}}%
\pgfpathlineto{\pgfqpoint{4.240383in}{0.164715in}}%
\pgfpathlineto{\pgfqpoint{4.237255in}{0.164257in}}%
\pgfpathlineto{\pgfqpoint{4.239114in}{0.168482in}}%
\pgfpathclose%
\pgfusepath{fill}%
\end{pgfscope}%
\begin{pgfscope}%
\pgfpathrectangle{\pgfqpoint{3.525000in}{0.100000in}}{\pgfqpoint{2.862985in}{1.832778in}}%
\pgfusepath{clip}%
\pgfsetbuttcap%
\pgfsetmiterjoin%
\definecolor{currentfill}{rgb}{0.081200,0.764245,0.988466}%
\pgfsetfillcolor{currentfill}%
\pgfsetlinewidth{0.000000pt}%
\definecolor{currentstroke}{rgb}{0.000000,0.000000,0.000000}%
\pgfsetstrokecolor{currentstroke}%
\pgfsetstrokeopacity{0.000000}%
\pgfsetdash{}{0pt}%
\pgfpathmoveto{\pgfqpoint{4.241751in}{0.162942in}}%
\pgfpathlineto{\pgfqpoint{4.241315in}{0.158141in}}%
\pgfpathlineto{\pgfqpoint{4.238015in}{0.158538in}}%
\pgfpathlineto{\pgfqpoint{4.239083in}{0.160858in}}%
\pgfpathlineto{\pgfqpoint{4.241751in}{0.162942in}}%
\pgfpathclose%
\pgfusepath{fill}%
\end{pgfscope}%
\begin{pgfscope}%
\definecolor{textcolor}{rgb}{0.000000,0.000000,0.000000}%
\pgfsetstrokecolor{textcolor}%
\pgfsetfillcolor{textcolor}%
\pgftext[x=4.727454in,y=1.877794in,left,base]{\color{textcolor}\setmainfont{Lato}\rmfamily\fontsize{9.000000}{10.800000}\selectfont December 2019}%
\end{pgfscope}%
\begin{pgfscope}%
\pgfpathrectangle{\pgfqpoint{2.893540in}{0.169444in}}{\pgfqpoint{0.858895in}{1.282944in}}%
\pgfusepath{clip}%
\pgfsetbuttcap%
\pgfsetmiterjoin%
\definecolor{currentfill}{rgb}{1.000000,1.000000,1.000000}%
\pgfsetfillcolor{currentfill}%
\pgfsetlinewidth{0.000000pt}%
\definecolor{currentstroke}{rgb}{0.000000,0.000000,0.000000}%
\pgfsetstrokecolor{currentstroke}%
\pgfsetstrokeopacity{0.000000}%
\pgfsetdash{}{0pt}%
\pgfpathmoveto{\pgfqpoint{3.142620in}{0.151754in}}%
\pgfpathlineto{\pgfqpoint{3.108264in}{0.151754in}}%
\pgfpathlineto{\pgfqpoint{3.108264in}{0.188409in}}%
\pgfpathlineto{\pgfqpoint{3.142620in}{0.188409in}}%
\pgfpathlineto{\pgfqpoint{3.142620in}{0.151754in}}%
\pgfpathclose%
\pgfusepath{fill}%
\end{pgfscope}%
\begin{pgfscope}%
\pgfpathrectangle{\pgfqpoint{2.893540in}{0.169444in}}{\pgfqpoint{0.858895in}{1.282944in}}%
\pgfusepath{clip}%
\pgfsetbuttcap%
\pgfsetmiterjoin%
\definecolor{currentfill}{rgb}{0.985421,0.994279,0.995386}%
\pgfsetfillcolor{currentfill}%
\pgfsetlinewidth{0.000000pt}%
\definecolor{currentstroke}{rgb}{0.000000,0.000000,0.000000}%
\pgfsetstrokecolor{currentstroke}%
\pgfsetstrokeopacity{0.000000}%
\pgfsetdash{}{0pt}%
\pgfpathmoveto{\pgfqpoint{3.142620in}{0.186881in}}%
\pgfpathlineto{\pgfqpoint{3.108264in}{0.186881in}}%
\pgfpathlineto{\pgfqpoint{3.108264in}{0.223536in}}%
\pgfpathlineto{\pgfqpoint{3.142620in}{0.223536in}}%
\pgfpathlineto{\pgfqpoint{3.142620in}{0.186881in}}%
\pgfpathclose%
\pgfusepath{fill}%
\end{pgfscope}%
\begin{pgfscope}%
\pgfpathrectangle{\pgfqpoint{2.893540in}{0.169444in}}{\pgfqpoint{0.858895in}{1.282944in}}%
\pgfusepath{clip}%
\pgfsetbuttcap%
\pgfsetmiterjoin%
\definecolor{currentfill}{rgb}{0.951403,0.980930,0.984621}%
\pgfsetfillcolor{currentfill}%
\pgfsetlinewidth{0.000000pt}%
\definecolor{currentstroke}{rgb}{0.000000,0.000000,0.000000}%
\pgfsetstrokecolor{currentstroke}%
\pgfsetstrokeopacity{0.000000}%
\pgfsetdash{}{0pt}%
\pgfpathmoveto{\pgfqpoint{3.142620in}{0.222007in}}%
\pgfpathlineto{\pgfqpoint{3.108264in}{0.222007in}}%
\pgfpathlineto{\pgfqpoint{3.108264in}{0.258663in}}%
\pgfpathlineto{\pgfqpoint{3.142620in}{0.258663in}}%
\pgfpathlineto{\pgfqpoint{3.142620in}{0.222007in}}%
\pgfpathclose%
\pgfusepath{fill}%
\end{pgfscope}%
\begin{pgfscope}%
\pgfpathrectangle{\pgfqpoint{2.893540in}{0.169444in}}{\pgfqpoint{0.858895in}{1.282944in}}%
\pgfusepath{clip}%
\pgfsetbuttcap%
\pgfsetmiterjoin%
\definecolor{currentfill}{rgb}{0.917386,0.967582,0.973856}%
\pgfsetfillcolor{currentfill}%
\pgfsetlinewidth{0.000000pt}%
\definecolor{currentstroke}{rgb}{0.000000,0.000000,0.000000}%
\pgfsetstrokecolor{currentstroke}%
\pgfsetstrokeopacity{0.000000}%
\pgfsetdash{}{0pt}%
\pgfpathmoveto{\pgfqpoint{3.142620in}{0.257134in}}%
\pgfpathlineto{\pgfqpoint{3.108264in}{0.257134in}}%
\pgfpathlineto{\pgfqpoint{3.108264in}{0.293789in}}%
\pgfpathlineto{\pgfqpoint{3.142620in}{0.293789in}}%
\pgfpathlineto{\pgfqpoint{3.142620in}{0.257134in}}%
\pgfpathclose%
\pgfusepath{fill}%
\end{pgfscope}%
\begin{pgfscope}%
\pgfpathrectangle{\pgfqpoint{2.893540in}{0.169444in}}{\pgfqpoint{0.858895in}{1.282944in}}%
\pgfusepath{clip}%
\pgfsetbuttcap%
\pgfsetmiterjoin%
\definecolor{currentfill}{rgb}{0.883368,0.954233,0.963091}%
\pgfsetfillcolor{currentfill}%
\pgfsetlinewidth{0.000000pt}%
\definecolor{currentstroke}{rgb}{0.000000,0.000000,0.000000}%
\pgfsetstrokecolor{currentstroke}%
\pgfsetstrokeopacity{0.000000}%
\pgfsetdash{}{0pt}%
\pgfpathmoveto{\pgfqpoint{3.142620in}{0.292260in}}%
\pgfpathlineto{\pgfqpoint{3.108264in}{0.292260in}}%
\pgfpathlineto{\pgfqpoint{3.108264in}{0.328916in}}%
\pgfpathlineto{\pgfqpoint{3.142620in}{0.328916in}}%
\pgfpathlineto{\pgfqpoint{3.142620in}{0.292260in}}%
\pgfpathclose%
\pgfusepath{fill}%
\end{pgfscope}%
\begin{pgfscope}%
\pgfpathrectangle{\pgfqpoint{2.893540in}{0.169444in}}{\pgfqpoint{0.858895in}{1.282944in}}%
\pgfusepath{clip}%
\pgfsetbuttcap%
\pgfsetmiterjoin%
\definecolor{currentfill}{rgb}{0.849350,0.940884,0.952326}%
\pgfsetfillcolor{currentfill}%
\pgfsetlinewidth{0.000000pt}%
\definecolor{currentstroke}{rgb}{0.000000,0.000000,0.000000}%
\pgfsetstrokecolor{currentstroke}%
\pgfsetstrokeopacity{0.000000}%
\pgfsetdash{}{0pt}%
\pgfpathmoveto{\pgfqpoint{3.142620in}{0.327387in}}%
\pgfpathlineto{\pgfqpoint{3.108264in}{0.327387in}}%
\pgfpathlineto{\pgfqpoint{3.108264in}{0.364042in}}%
\pgfpathlineto{\pgfqpoint{3.142620in}{0.364042in}}%
\pgfpathlineto{\pgfqpoint{3.142620in}{0.327387in}}%
\pgfpathclose%
\pgfusepath{fill}%
\end{pgfscope}%
\begin{pgfscope}%
\pgfpathrectangle{\pgfqpoint{2.893540in}{0.169444in}}{\pgfqpoint{0.858895in}{1.282944in}}%
\pgfusepath{clip}%
\pgfsetbuttcap%
\pgfsetmiterjoin%
\definecolor{currentfill}{rgb}{0.815333,0.927536,0.941561}%
\pgfsetfillcolor{currentfill}%
\pgfsetlinewidth{0.000000pt}%
\definecolor{currentstroke}{rgb}{0.000000,0.000000,0.000000}%
\pgfsetstrokecolor{currentstroke}%
\pgfsetstrokeopacity{0.000000}%
\pgfsetdash{}{0pt}%
\pgfpathmoveto{\pgfqpoint{3.142620in}{0.362514in}}%
\pgfpathlineto{\pgfqpoint{3.108264in}{0.362514in}}%
\pgfpathlineto{\pgfqpoint{3.108264in}{0.399169in}}%
\pgfpathlineto{\pgfqpoint{3.142620in}{0.399169in}}%
\pgfpathlineto{\pgfqpoint{3.142620in}{0.362514in}}%
\pgfpathclose%
\pgfusepath{fill}%
\end{pgfscope}%
\begin{pgfscope}%
\pgfpathrectangle{\pgfqpoint{2.893540in}{0.169444in}}{\pgfqpoint{0.858895in}{1.282944in}}%
\pgfusepath{clip}%
\pgfsetbuttcap%
\pgfsetmiterjoin%
\definecolor{currentfill}{rgb}{0.781315,0.914187,0.930796}%
\pgfsetfillcolor{currentfill}%
\pgfsetlinewidth{0.000000pt}%
\definecolor{currentstroke}{rgb}{0.000000,0.000000,0.000000}%
\pgfsetstrokecolor{currentstroke}%
\pgfsetstrokeopacity{0.000000}%
\pgfsetdash{}{0pt}%
\pgfpathmoveto{\pgfqpoint{3.142620in}{0.397640in}}%
\pgfpathlineto{\pgfqpoint{3.108264in}{0.397640in}}%
\pgfpathlineto{\pgfqpoint{3.108264in}{0.434296in}}%
\pgfpathlineto{\pgfqpoint{3.142620in}{0.434296in}}%
\pgfpathlineto{\pgfqpoint{3.142620in}{0.397640in}}%
\pgfpathclose%
\pgfusepath{fill}%
\end{pgfscope}%
\begin{pgfscope}%
\pgfpathrectangle{\pgfqpoint{2.893540in}{0.169444in}}{\pgfqpoint{0.858895in}{1.282944in}}%
\pgfusepath{clip}%
\pgfsetbuttcap%
\pgfsetmiterjoin%
\definecolor{currentfill}{rgb}{0.747297,0.900838,0.920031}%
\pgfsetfillcolor{currentfill}%
\pgfsetlinewidth{0.000000pt}%
\definecolor{currentstroke}{rgb}{0.000000,0.000000,0.000000}%
\pgfsetstrokecolor{currentstroke}%
\pgfsetstrokeopacity{0.000000}%
\pgfsetdash{}{0pt}%
\pgfpathmoveto{\pgfqpoint{3.142620in}{0.432767in}}%
\pgfpathlineto{\pgfqpoint{3.108264in}{0.432767in}}%
\pgfpathlineto{\pgfqpoint{3.108264in}{0.469422in}}%
\pgfpathlineto{\pgfqpoint{3.142620in}{0.469422in}}%
\pgfpathlineto{\pgfqpoint{3.142620in}{0.432767in}}%
\pgfpathclose%
\pgfusepath{fill}%
\end{pgfscope}%
\begin{pgfscope}%
\pgfpathrectangle{\pgfqpoint{2.893540in}{0.169444in}}{\pgfqpoint{0.858895in}{1.282944in}}%
\pgfusepath{clip}%
\pgfsetbuttcap%
\pgfsetmiterjoin%
\definecolor{currentfill}{rgb}{0.713280,0.887489,0.909266}%
\pgfsetfillcolor{currentfill}%
\pgfsetlinewidth{0.000000pt}%
\definecolor{currentstroke}{rgb}{0.000000,0.000000,0.000000}%
\pgfsetstrokecolor{currentstroke}%
\pgfsetstrokeopacity{0.000000}%
\pgfsetdash{}{0pt}%
\pgfpathmoveto{\pgfqpoint{3.142620in}{0.467893in}}%
\pgfpathlineto{\pgfqpoint{3.108264in}{0.467893in}}%
\pgfpathlineto{\pgfqpoint{3.108264in}{0.504549in}}%
\pgfpathlineto{\pgfqpoint{3.142620in}{0.504549in}}%
\pgfpathlineto{\pgfqpoint{3.142620in}{0.467893in}}%
\pgfpathclose%
\pgfusepath{fill}%
\end{pgfscope}%
\begin{pgfscope}%
\pgfpathrectangle{\pgfqpoint{2.893540in}{0.169444in}}{\pgfqpoint{0.858895in}{1.282944in}}%
\pgfusepath{clip}%
\pgfsetbuttcap%
\pgfsetmiterjoin%
\definecolor{currentfill}{rgb}{0.665836,0.873864,0.905421}%
\pgfsetfillcolor{currentfill}%
\pgfsetlinewidth{0.000000pt}%
\definecolor{currentstroke}{rgb}{0.000000,0.000000,0.000000}%
\pgfsetstrokecolor{currentstroke}%
\pgfsetstrokeopacity{0.000000}%
\pgfsetdash{}{0pt}%
\pgfpathmoveto{\pgfqpoint{3.142620in}{0.503020in}}%
\pgfpathlineto{\pgfqpoint{3.108264in}{0.503020in}}%
\pgfpathlineto{\pgfqpoint{3.108264in}{0.539675in}}%
\pgfpathlineto{\pgfqpoint{3.142620in}{0.539675in}}%
\pgfpathlineto{\pgfqpoint{3.142620in}{0.503020in}}%
\pgfpathclose%
\pgfusepath{fill}%
\end{pgfscope}%
\begin{pgfscope}%
\pgfpathrectangle{\pgfqpoint{2.893540in}{0.169444in}}{\pgfqpoint{0.858895in}{1.282944in}}%
\pgfusepath{clip}%
\pgfsetbuttcap%
\pgfsetmiterjoin%
\definecolor{currentfill}{rgb}{0.590050,0.859654,0.916186}%
\pgfsetfillcolor{currentfill}%
\pgfsetlinewidth{0.000000pt}%
\definecolor{currentstroke}{rgb}{0.000000,0.000000,0.000000}%
\pgfsetstrokecolor{currentstroke}%
\pgfsetstrokeopacity{0.000000}%
\pgfsetdash{}{0pt}%
\pgfpathmoveto{\pgfqpoint{3.142620in}{0.538146in}}%
\pgfpathlineto{\pgfqpoint{3.108264in}{0.538146in}}%
\pgfpathlineto{\pgfqpoint{3.108264in}{0.574802in}}%
\pgfpathlineto{\pgfqpoint{3.142620in}{0.574802in}}%
\pgfpathlineto{\pgfqpoint{3.142620in}{0.538146in}}%
\pgfpathclose%
\pgfusepath{fill}%
\end{pgfscope}%
\begin{pgfscope}%
\pgfpathrectangle{\pgfqpoint{2.893540in}{0.169444in}}{\pgfqpoint{0.858895in}{1.282944in}}%
\pgfusepath{clip}%
\pgfsetbuttcap%
\pgfsetmiterjoin%
\definecolor{currentfill}{rgb}{0.514264,0.845444,0.926951}%
\pgfsetfillcolor{currentfill}%
\pgfsetlinewidth{0.000000pt}%
\definecolor{currentstroke}{rgb}{0.000000,0.000000,0.000000}%
\pgfsetstrokecolor{currentstroke}%
\pgfsetstrokeopacity{0.000000}%
\pgfsetdash{}{0pt}%
\pgfpathmoveto{\pgfqpoint{3.142620in}{0.573273in}}%
\pgfpathlineto{\pgfqpoint{3.108264in}{0.573273in}}%
\pgfpathlineto{\pgfqpoint{3.108264in}{0.609929in}}%
\pgfpathlineto{\pgfqpoint{3.142620in}{0.609929in}}%
\pgfpathlineto{\pgfqpoint{3.142620in}{0.573273in}}%
\pgfpathclose%
\pgfusepath{fill}%
\end{pgfscope}%
\begin{pgfscope}%
\pgfpathrectangle{\pgfqpoint{2.893540in}{0.169444in}}{\pgfqpoint{0.858895in}{1.282944in}}%
\pgfusepath{clip}%
\pgfsetbuttcap%
\pgfsetmiterjoin%
\definecolor{currentfill}{rgb}{0.438478,0.831234,0.937716}%
\pgfsetfillcolor{currentfill}%
\pgfsetlinewidth{0.000000pt}%
\definecolor{currentstroke}{rgb}{0.000000,0.000000,0.000000}%
\pgfsetstrokecolor{currentstroke}%
\pgfsetstrokeopacity{0.000000}%
\pgfsetdash{}{0pt}%
\pgfpathmoveto{\pgfqpoint{3.142620in}{0.608400in}}%
\pgfpathlineto{\pgfqpoint{3.108264in}{0.608400in}}%
\pgfpathlineto{\pgfqpoint{3.108264in}{0.645055in}}%
\pgfpathlineto{\pgfqpoint{3.142620in}{0.645055in}}%
\pgfpathlineto{\pgfqpoint{3.142620in}{0.608400in}}%
\pgfpathclose%
\pgfusepath{fill}%
\end{pgfscope}%
\begin{pgfscope}%
\pgfpathrectangle{\pgfqpoint{2.893540in}{0.169444in}}{\pgfqpoint{0.858895in}{1.282944in}}%
\pgfusepath{clip}%
\pgfsetbuttcap%
\pgfsetmiterjoin%
\definecolor{currentfill}{rgb}{0.362691,0.817024,0.948481}%
\pgfsetfillcolor{currentfill}%
\pgfsetlinewidth{0.000000pt}%
\definecolor{currentstroke}{rgb}{0.000000,0.000000,0.000000}%
\pgfsetstrokecolor{currentstroke}%
\pgfsetstrokeopacity{0.000000}%
\pgfsetdash{}{0pt}%
\pgfpathmoveto{\pgfqpoint{3.142620in}{0.643526in}}%
\pgfpathlineto{\pgfqpoint{3.108264in}{0.643526in}}%
\pgfpathlineto{\pgfqpoint{3.108264in}{0.680182in}}%
\pgfpathlineto{\pgfqpoint{3.142620in}{0.680182in}}%
\pgfpathlineto{\pgfqpoint{3.142620in}{0.643526in}}%
\pgfpathclose%
\pgfusepath{fill}%
\end{pgfscope}%
\begin{pgfscope}%
\pgfpathrectangle{\pgfqpoint{2.893540in}{0.169444in}}{\pgfqpoint{0.858895in}{1.282944in}}%
\pgfusepath{clip}%
\pgfsetbuttcap%
\pgfsetmiterjoin%
\definecolor{currentfill}{rgb}{0.286905,0.802814,0.959246}%
\pgfsetfillcolor{currentfill}%
\pgfsetlinewidth{0.000000pt}%
\definecolor{currentstroke}{rgb}{0.000000,0.000000,0.000000}%
\pgfsetstrokecolor{currentstroke}%
\pgfsetstrokeopacity{0.000000}%
\pgfsetdash{}{0pt}%
\pgfpathmoveto{\pgfqpoint{3.142620in}{0.678653in}}%
\pgfpathlineto{\pgfqpoint{3.108264in}{0.678653in}}%
\pgfpathlineto{\pgfqpoint{3.108264in}{0.715308in}}%
\pgfpathlineto{\pgfqpoint{3.142620in}{0.715308in}}%
\pgfpathlineto{\pgfqpoint{3.142620in}{0.678653in}}%
\pgfpathclose%
\pgfusepath{fill}%
\end{pgfscope}%
\begin{pgfscope}%
\pgfpathrectangle{\pgfqpoint{2.893540in}{0.169444in}}{\pgfqpoint{0.858895in}{1.282944in}}%
\pgfusepath{clip}%
\pgfsetbuttcap%
\pgfsetmiterjoin%
\definecolor{currentfill}{rgb}{0.211119,0.788604,0.970012}%
\pgfsetfillcolor{currentfill}%
\pgfsetlinewidth{0.000000pt}%
\definecolor{currentstroke}{rgb}{0.000000,0.000000,0.000000}%
\pgfsetstrokecolor{currentstroke}%
\pgfsetstrokeopacity{0.000000}%
\pgfsetdash{}{0pt}%
\pgfpathmoveto{\pgfqpoint{3.142620in}{0.713779in}}%
\pgfpathlineto{\pgfqpoint{3.108264in}{0.713779in}}%
\pgfpathlineto{\pgfqpoint{3.108264in}{0.750435in}}%
\pgfpathlineto{\pgfqpoint{3.142620in}{0.750435in}}%
\pgfpathlineto{\pgfqpoint{3.142620in}{0.713779in}}%
\pgfpathclose%
\pgfusepath{fill}%
\end{pgfscope}%
\begin{pgfscope}%
\pgfpathrectangle{\pgfqpoint{2.893540in}{0.169444in}}{\pgfqpoint{0.858895in}{1.282944in}}%
\pgfusepath{clip}%
\pgfsetbuttcap%
\pgfsetmiterjoin%
\definecolor{currentfill}{rgb}{0.135333,0.774394,0.980777}%
\pgfsetfillcolor{currentfill}%
\pgfsetlinewidth{0.000000pt}%
\definecolor{currentstroke}{rgb}{0.000000,0.000000,0.000000}%
\pgfsetstrokecolor{currentstroke}%
\pgfsetstrokeopacity{0.000000}%
\pgfsetdash{}{0pt}%
\pgfpathmoveto{\pgfqpoint{3.142620in}{0.748906in}}%
\pgfpathlineto{\pgfqpoint{3.108264in}{0.748906in}}%
\pgfpathlineto{\pgfqpoint{3.108264in}{0.785562in}}%
\pgfpathlineto{\pgfqpoint{3.142620in}{0.785562in}}%
\pgfpathlineto{\pgfqpoint{3.142620in}{0.748906in}}%
\pgfpathclose%
\pgfusepath{fill}%
\end{pgfscope}%
\begin{pgfscope}%
\pgfpathrectangle{\pgfqpoint{2.893540in}{0.169444in}}{\pgfqpoint{0.858895in}{1.282944in}}%
\pgfusepath{clip}%
\pgfsetbuttcap%
\pgfsetmiterjoin%
\definecolor{currentfill}{rgb}{0.059546,0.760185,0.991542}%
\pgfsetfillcolor{currentfill}%
\pgfsetlinewidth{0.000000pt}%
\definecolor{currentstroke}{rgb}{0.000000,0.000000,0.000000}%
\pgfsetstrokecolor{currentstroke}%
\pgfsetstrokeopacity{0.000000}%
\pgfsetdash{}{0pt}%
\pgfpathmoveto{\pgfqpoint{3.142620in}{0.784033in}}%
\pgfpathlineto{\pgfqpoint{3.108264in}{0.784033in}}%
\pgfpathlineto{\pgfqpoint{3.108264in}{0.820688in}}%
\pgfpathlineto{\pgfqpoint{3.142620in}{0.820688in}}%
\pgfpathlineto{\pgfqpoint{3.142620in}{0.784033in}}%
\pgfpathclose%
\pgfusepath{fill}%
\end{pgfscope}%
\begin{pgfscope}%
\pgfpathrectangle{\pgfqpoint{2.893540in}{0.169444in}}{\pgfqpoint{0.858895in}{1.282944in}}%
\pgfusepath{clip}%
\pgfsetbuttcap%
\pgfsetmiterjoin%
\definecolor{currentfill}{rgb}{0.004614,0.750311,0.981084}%
\pgfsetfillcolor{currentfill}%
\pgfsetlinewidth{0.000000pt}%
\definecolor{currentstroke}{rgb}{0.000000,0.000000,0.000000}%
\pgfsetstrokecolor{currentstroke}%
\pgfsetstrokeopacity{0.000000}%
\pgfsetdash{}{0pt}%
\pgfpathmoveto{\pgfqpoint{3.142620in}{0.819159in}}%
\pgfpathlineto{\pgfqpoint{3.108264in}{0.819159in}}%
\pgfpathlineto{\pgfqpoint{3.108264in}{0.855815in}}%
\pgfpathlineto{\pgfqpoint{3.142620in}{0.855815in}}%
\pgfpathlineto{\pgfqpoint{3.142620in}{0.819159in}}%
\pgfpathclose%
\pgfusepath{fill}%
\end{pgfscope}%
\begin{pgfscope}%
\pgfpathrectangle{\pgfqpoint{2.893540in}{0.169444in}}{\pgfqpoint{0.858895in}{1.282944in}}%
\pgfusepath{clip}%
\pgfsetbuttcap%
\pgfsetmiterjoin%
\definecolor{currentfill}{rgb}{0.026144,0.756340,0.892810}%
\pgfsetfillcolor{currentfill}%
\pgfsetlinewidth{0.000000pt}%
\definecolor{currentstroke}{rgb}{0.000000,0.000000,0.000000}%
\pgfsetstrokecolor{currentstroke}%
\pgfsetstrokeopacity{0.000000}%
\pgfsetdash{}{0pt}%
\pgfpathmoveto{\pgfqpoint{3.142620in}{0.854286in}}%
\pgfpathlineto{\pgfqpoint{3.108264in}{0.854286in}}%
\pgfpathlineto{\pgfqpoint{3.108264in}{0.890941in}}%
\pgfpathlineto{\pgfqpoint{3.142620in}{0.890941in}}%
\pgfpathlineto{\pgfqpoint{3.142620in}{0.854286in}}%
\pgfpathclose%
\pgfusepath{fill}%
\end{pgfscope}%
\begin{pgfscope}%
\pgfpathrectangle{\pgfqpoint{2.893540in}{0.169444in}}{\pgfqpoint{0.858895in}{1.282944in}}%
\pgfusepath{clip}%
\pgfsetbuttcap%
\pgfsetmiterjoin%
\definecolor{currentfill}{rgb}{0.047674,0.762368,0.804537}%
\pgfsetfillcolor{currentfill}%
\pgfsetlinewidth{0.000000pt}%
\definecolor{currentstroke}{rgb}{0.000000,0.000000,0.000000}%
\pgfsetstrokecolor{currentstroke}%
\pgfsetstrokeopacity{0.000000}%
\pgfsetdash{}{0pt}%
\pgfpathmoveto{\pgfqpoint{3.142620in}{0.889412in}}%
\pgfpathlineto{\pgfqpoint{3.108264in}{0.889412in}}%
\pgfpathlineto{\pgfqpoint{3.108264in}{0.926068in}}%
\pgfpathlineto{\pgfqpoint{3.142620in}{0.926068in}}%
\pgfpathlineto{\pgfqpoint{3.142620in}{0.889412in}}%
\pgfpathclose%
\pgfusepath{fill}%
\end{pgfscope}%
\begin{pgfscope}%
\pgfpathrectangle{\pgfqpoint{2.893540in}{0.169444in}}{\pgfqpoint{0.858895in}{1.282944in}}%
\pgfusepath{clip}%
\pgfsetbuttcap%
\pgfsetmiterjoin%
\definecolor{currentfill}{rgb}{0.069204,0.768397,0.716263}%
\pgfsetfillcolor{currentfill}%
\pgfsetlinewidth{0.000000pt}%
\definecolor{currentstroke}{rgb}{0.000000,0.000000,0.000000}%
\pgfsetstrokecolor{currentstroke}%
\pgfsetstrokeopacity{0.000000}%
\pgfsetdash{}{0pt}%
\pgfpathmoveto{\pgfqpoint{3.142620in}{0.924539in}}%
\pgfpathlineto{\pgfqpoint{3.108264in}{0.924539in}}%
\pgfpathlineto{\pgfqpoint{3.108264in}{0.961195in}}%
\pgfpathlineto{\pgfqpoint{3.142620in}{0.961195in}}%
\pgfpathlineto{\pgfqpoint{3.142620in}{0.924539in}}%
\pgfpathclose%
\pgfusepath{fill}%
\end{pgfscope}%
\begin{pgfscope}%
\pgfpathrectangle{\pgfqpoint{2.893540in}{0.169444in}}{\pgfqpoint{0.858895in}{1.282944in}}%
\pgfusepath{clip}%
\pgfsetbuttcap%
\pgfsetmiterjoin%
\definecolor{currentfill}{rgb}{0.090734,0.774425,0.627989}%
\pgfsetfillcolor{currentfill}%
\pgfsetlinewidth{0.000000pt}%
\definecolor{currentstroke}{rgb}{0.000000,0.000000,0.000000}%
\pgfsetstrokecolor{currentstroke}%
\pgfsetstrokeopacity{0.000000}%
\pgfsetdash{}{0pt}%
\pgfpathmoveto{\pgfqpoint{3.142620in}{0.959666in}}%
\pgfpathlineto{\pgfqpoint{3.108264in}{0.959666in}}%
\pgfpathlineto{\pgfqpoint{3.108264in}{0.996321in}}%
\pgfpathlineto{\pgfqpoint{3.142620in}{0.996321in}}%
\pgfpathlineto{\pgfqpoint{3.142620in}{0.959666in}}%
\pgfpathclose%
\pgfusepath{fill}%
\end{pgfscope}%
\begin{pgfscope}%
\pgfpathrectangle{\pgfqpoint{2.893540in}{0.169444in}}{\pgfqpoint{0.858895in}{1.282944in}}%
\pgfusepath{clip}%
\pgfsetbuttcap%
\pgfsetmiterjoin%
\definecolor{currentfill}{rgb}{0.112265,0.780454,0.539715}%
\pgfsetfillcolor{currentfill}%
\pgfsetlinewidth{0.000000pt}%
\definecolor{currentstroke}{rgb}{0.000000,0.000000,0.000000}%
\pgfsetstrokecolor{currentstroke}%
\pgfsetstrokeopacity{0.000000}%
\pgfsetdash{}{0pt}%
\pgfpathmoveto{\pgfqpoint{3.142620in}{0.994792in}}%
\pgfpathlineto{\pgfqpoint{3.108264in}{0.994792in}}%
\pgfpathlineto{\pgfqpoint{3.108264in}{1.031448in}}%
\pgfpathlineto{\pgfqpoint{3.142620in}{1.031448in}}%
\pgfpathlineto{\pgfqpoint{3.142620in}{0.994792in}}%
\pgfpathclose%
\pgfusepath{fill}%
\end{pgfscope}%
\begin{pgfscope}%
\pgfpathrectangle{\pgfqpoint{2.893540in}{0.169444in}}{\pgfqpoint{0.858895in}{1.282944in}}%
\pgfusepath{clip}%
\pgfsetbuttcap%
\pgfsetmiterjoin%
\definecolor{currentfill}{rgb}{0.133795,0.786482,0.451442}%
\pgfsetfillcolor{currentfill}%
\pgfsetlinewidth{0.000000pt}%
\definecolor{currentstroke}{rgb}{0.000000,0.000000,0.000000}%
\pgfsetstrokecolor{currentstroke}%
\pgfsetstrokeopacity{0.000000}%
\pgfsetdash{}{0pt}%
\pgfpathmoveto{\pgfqpoint{3.142620in}{1.029919in}}%
\pgfpathlineto{\pgfqpoint{3.108264in}{1.029919in}}%
\pgfpathlineto{\pgfqpoint{3.108264in}{1.066574in}}%
\pgfpathlineto{\pgfqpoint{3.142620in}{1.066574in}}%
\pgfpathlineto{\pgfqpoint{3.142620in}{1.029919in}}%
\pgfpathclose%
\pgfusepath{fill}%
\end{pgfscope}%
\begin{pgfscope}%
\pgfpathrectangle{\pgfqpoint{2.893540in}{0.169444in}}{\pgfqpoint{0.858895in}{1.282944in}}%
\pgfusepath{clip}%
\pgfsetbuttcap%
\pgfsetmiterjoin%
\definecolor{currentfill}{rgb}{0.155325,0.792511,0.363168}%
\pgfsetfillcolor{currentfill}%
\pgfsetlinewidth{0.000000pt}%
\definecolor{currentstroke}{rgb}{0.000000,0.000000,0.000000}%
\pgfsetstrokecolor{currentstroke}%
\pgfsetstrokeopacity{0.000000}%
\pgfsetdash{}{0pt}%
\pgfpathmoveto{\pgfqpoint{3.142620in}{1.065045in}}%
\pgfpathlineto{\pgfqpoint{3.108264in}{1.065045in}}%
\pgfpathlineto{\pgfqpoint{3.108264in}{1.101701in}}%
\pgfpathlineto{\pgfqpoint{3.142620in}{1.101701in}}%
\pgfpathlineto{\pgfqpoint{3.142620in}{1.065045in}}%
\pgfpathclose%
\pgfusepath{fill}%
\end{pgfscope}%
\begin{pgfscope}%
\pgfpathrectangle{\pgfqpoint{2.893540in}{0.169444in}}{\pgfqpoint{0.858895in}{1.282944in}}%
\pgfusepath{clip}%
\pgfsetbuttcap%
\pgfsetmiterjoin%
\definecolor{currentfill}{rgb}{0.176855,0.798539,0.274894}%
\pgfsetfillcolor{currentfill}%
\pgfsetlinewidth{0.000000pt}%
\definecolor{currentstroke}{rgb}{0.000000,0.000000,0.000000}%
\pgfsetstrokecolor{currentstroke}%
\pgfsetstrokeopacity{0.000000}%
\pgfsetdash{}{0pt}%
\pgfpathmoveto{\pgfqpoint{3.142620in}{1.100172in}}%
\pgfpathlineto{\pgfqpoint{3.108264in}{1.100172in}}%
\pgfpathlineto{\pgfqpoint{3.108264in}{1.136828in}}%
\pgfpathlineto{\pgfqpoint{3.142620in}{1.136828in}}%
\pgfpathlineto{\pgfqpoint{3.142620in}{1.100172in}}%
\pgfpathclose%
\pgfusepath{fill}%
\end{pgfscope}%
\begin{pgfscope}%
\pgfpathrectangle{\pgfqpoint{2.893540in}{0.169444in}}{\pgfqpoint{0.858895in}{1.282944in}}%
\pgfusepath{clip}%
\pgfsetbuttcap%
\pgfsetmiterjoin%
\definecolor{currentfill}{rgb}{0.193772,0.799077,0.193772}%
\pgfsetfillcolor{currentfill}%
\pgfsetlinewidth{0.000000pt}%
\definecolor{currentstroke}{rgb}{0.000000,0.000000,0.000000}%
\pgfsetstrokecolor{currentstroke}%
\pgfsetstrokeopacity{0.000000}%
\pgfsetdash{}{0pt}%
\pgfpathmoveto{\pgfqpoint{3.142620in}{1.135299in}}%
\pgfpathlineto{\pgfqpoint{3.108264in}{1.135299in}}%
\pgfpathlineto{\pgfqpoint{3.108264in}{1.171954in}}%
\pgfpathlineto{\pgfqpoint{3.142620in}{1.171954in}}%
\pgfpathlineto{\pgfqpoint{3.142620in}{1.135299in}}%
\pgfpathclose%
\pgfusepath{fill}%
\end{pgfscope}%
\begin{pgfscope}%
\pgfpathrectangle{\pgfqpoint{2.893540in}{0.169444in}}{\pgfqpoint{0.858895in}{1.282944in}}%
\pgfusepath{clip}%
\pgfsetbuttcap%
\pgfsetmiterjoin%
\definecolor{currentfill}{rgb}{0.172241,0.753864,0.172241}%
\pgfsetfillcolor{currentfill}%
\pgfsetlinewidth{0.000000pt}%
\definecolor{currentstroke}{rgb}{0.000000,0.000000,0.000000}%
\pgfsetstrokecolor{currentstroke}%
\pgfsetstrokeopacity{0.000000}%
\pgfsetdash{}{0pt}%
\pgfpathmoveto{\pgfqpoint{3.142620in}{1.170425in}}%
\pgfpathlineto{\pgfqpoint{3.108264in}{1.170425in}}%
\pgfpathlineto{\pgfqpoint{3.108264in}{1.207081in}}%
\pgfpathlineto{\pgfqpoint{3.142620in}{1.207081in}}%
\pgfpathlineto{\pgfqpoint{3.142620in}{1.170425in}}%
\pgfpathclose%
\pgfusepath{fill}%
\end{pgfscope}%
\begin{pgfscope}%
\pgfpathrectangle{\pgfqpoint{2.893540in}{0.169444in}}{\pgfqpoint{0.858895in}{1.282944in}}%
\pgfusepath{clip}%
\pgfsetbuttcap%
\pgfsetmiterjoin%
\definecolor{currentfill}{rgb}{0.150711,0.708651,0.150711}%
\pgfsetfillcolor{currentfill}%
\pgfsetlinewidth{0.000000pt}%
\definecolor{currentstroke}{rgb}{0.000000,0.000000,0.000000}%
\pgfsetstrokecolor{currentstroke}%
\pgfsetstrokeopacity{0.000000}%
\pgfsetdash{}{0pt}%
\pgfpathmoveto{\pgfqpoint{3.142620in}{1.205552in}}%
\pgfpathlineto{\pgfqpoint{3.108264in}{1.205552in}}%
\pgfpathlineto{\pgfqpoint{3.108264in}{1.242207in}}%
\pgfpathlineto{\pgfqpoint{3.142620in}{1.242207in}}%
\pgfpathlineto{\pgfqpoint{3.142620in}{1.205552in}}%
\pgfpathclose%
\pgfusepath{fill}%
\end{pgfscope}%
\begin{pgfscope}%
\pgfpathrectangle{\pgfqpoint{2.893540in}{0.169444in}}{\pgfqpoint{0.858895in}{1.282944in}}%
\pgfusepath{clip}%
\pgfsetbuttcap%
\pgfsetmiterjoin%
\definecolor{currentfill}{rgb}{0.129181,0.663437,0.129181}%
\pgfsetfillcolor{currentfill}%
\pgfsetlinewidth{0.000000pt}%
\definecolor{currentstroke}{rgb}{0.000000,0.000000,0.000000}%
\pgfsetstrokecolor{currentstroke}%
\pgfsetstrokeopacity{0.000000}%
\pgfsetdash{}{0pt}%
\pgfpathmoveto{\pgfqpoint{3.142620in}{1.240678in}}%
\pgfpathlineto{\pgfqpoint{3.108264in}{1.240678in}}%
\pgfpathlineto{\pgfqpoint{3.108264in}{1.277334in}}%
\pgfpathlineto{\pgfqpoint{3.142620in}{1.277334in}}%
\pgfpathlineto{\pgfqpoint{3.142620in}{1.240678in}}%
\pgfpathclose%
\pgfusepath{fill}%
\end{pgfscope}%
\begin{pgfscope}%
\pgfpathrectangle{\pgfqpoint{2.893540in}{0.169444in}}{\pgfqpoint{0.858895in}{1.282944in}}%
\pgfusepath{clip}%
\pgfsetbuttcap%
\pgfsetmiterjoin%
\definecolor{currentfill}{rgb}{0.107651,0.618224,0.107651}%
\pgfsetfillcolor{currentfill}%
\pgfsetlinewidth{0.000000pt}%
\definecolor{currentstroke}{rgb}{0.000000,0.000000,0.000000}%
\pgfsetstrokecolor{currentstroke}%
\pgfsetstrokeopacity{0.000000}%
\pgfsetdash{}{0pt}%
\pgfpathmoveto{\pgfqpoint{3.142620in}{1.275805in}}%
\pgfpathlineto{\pgfqpoint{3.108264in}{1.275805in}}%
\pgfpathlineto{\pgfqpoint{3.108264in}{1.312460in}}%
\pgfpathlineto{\pgfqpoint{3.142620in}{1.312460in}}%
\pgfpathlineto{\pgfqpoint{3.142620in}{1.275805in}}%
\pgfpathclose%
\pgfusepath{fill}%
\end{pgfscope}%
\begin{pgfscope}%
\pgfpathrectangle{\pgfqpoint{2.893540in}{0.169444in}}{\pgfqpoint{0.858895in}{1.282944in}}%
\pgfusepath{clip}%
\pgfsetbuttcap%
\pgfsetmiterjoin%
\definecolor{currentfill}{rgb}{0.086121,0.573010,0.086121}%
\pgfsetfillcolor{currentfill}%
\pgfsetlinewidth{0.000000pt}%
\definecolor{currentstroke}{rgb}{0.000000,0.000000,0.000000}%
\pgfsetstrokecolor{currentstroke}%
\pgfsetstrokeopacity{0.000000}%
\pgfsetdash{}{0pt}%
\pgfpathmoveto{\pgfqpoint{3.142620in}{1.310932in}}%
\pgfpathlineto{\pgfqpoint{3.108264in}{1.310932in}}%
\pgfpathlineto{\pgfqpoint{3.108264in}{1.347587in}}%
\pgfpathlineto{\pgfqpoint{3.142620in}{1.347587in}}%
\pgfpathlineto{\pgfqpoint{3.142620in}{1.310932in}}%
\pgfpathclose%
\pgfusepath{fill}%
\end{pgfscope}%
\begin{pgfscope}%
\pgfpathrectangle{\pgfqpoint{2.893540in}{0.169444in}}{\pgfqpoint{0.858895in}{1.282944in}}%
\pgfusepath{clip}%
\pgfsetbuttcap%
\pgfsetmiterjoin%
\definecolor{currentfill}{rgb}{0.064591,0.527797,0.064591}%
\pgfsetfillcolor{currentfill}%
\pgfsetlinewidth{0.000000pt}%
\definecolor{currentstroke}{rgb}{0.000000,0.000000,0.000000}%
\pgfsetstrokecolor{currentstroke}%
\pgfsetstrokeopacity{0.000000}%
\pgfsetdash{}{0pt}%
\pgfpathmoveto{\pgfqpoint{3.142620in}{1.346058in}}%
\pgfpathlineto{\pgfqpoint{3.108264in}{1.346058in}}%
\pgfpathlineto{\pgfqpoint{3.108264in}{1.382714in}}%
\pgfpathlineto{\pgfqpoint{3.142620in}{1.382714in}}%
\pgfpathlineto{\pgfqpoint{3.142620in}{1.346058in}}%
\pgfpathclose%
\pgfusepath{fill}%
\end{pgfscope}%
\begin{pgfscope}%
\pgfpathrectangle{\pgfqpoint{2.893540in}{0.169444in}}{\pgfqpoint{0.858895in}{1.282944in}}%
\pgfusepath{clip}%
\pgfsetbuttcap%
\pgfsetmiterjoin%
\definecolor{currentfill}{rgb}{0.043060,0.482584,0.043060}%
\pgfsetfillcolor{currentfill}%
\pgfsetlinewidth{0.000000pt}%
\definecolor{currentstroke}{rgb}{0.000000,0.000000,0.000000}%
\pgfsetstrokecolor{currentstroke}%
\pgfsetstrokeopacity{0.000000}%
\pgfsetdash{}{0pt}%
\pgfpathmoveto{\pgfqpoint{3.142620in}{1.381185in}}%
\pgfpathlineto{\pgfqpoint{3.108264in}{1.381185in}}%
\pgfpathlineto{\pgfqpoint{3.108264in}{1.417840in}}%
\pgfpathlineto{\pgfqpoint{3.142620in}{1.417840in}}%
\pgfpathlineto{\pgfqpoint{3.142620in}{1.381185in}}%
\pgfpathclose%
\pgfusepath{fill}%
\end{pgfscope}%
\begin{pgfscope}%
\pgfpathrectangle{\pgfqpoint{2.893540in}{0.169444in}}{\pgfqpoint{0.858895in}{1.282944in}}%
\pgfusepath{clip}%
\pgfsetbuttcap%
\pgfsetmiterjoin%
\definecolor{currentfill}{rgb}{0.021530,0.437370,0.021530}%
\pgfsetfillcolor{currentfill}%
\pgfsetlinewidth{0.000000pt}%
\definecolor{currentstroke}{rgb}{0.000000,0.000000,0.000000}%
\pgfsetstrokecolor{currentstroke}%
\pgfsetstrokeopacity{0.000000}%
\pgfsetdash{}{0pt}%
\pgfpathmoveto{\pgfqpoint{3.142620in}{1.416311in}}%
\pgfpathlineto{\pgfqpoint{3.108264in}{1.416311in}}%
\pgfpathlineto{\pgfqpoint{3.108264in}{1.452967in}}%
\pgfpathlineto{\pgfqpoint{3.142620in}{1.452967in}}%
\pgfpathlineto{\pgfqpoint{3.142620in}{1.416311in}}%
\pgfpathclose%
\pgfusepath{fill}%
\end{pgfscope}%
\begin{pgfscope}%
\pgfpathrectangle{\pgfqpoint{2.893540in}{0.169444in}}{\pgfqpoint{0.858895in}{1.282944in}}%
\pgfusepath{clip}%
\pgfsetbuttcap%
\pgfsetmiterjoin%
\definecolor{currentfill}{rgb}{0.000000,0.392157,0.000000}%
\pgfsetfillcolor{currentfill}%
\pgfsetlinewidth{0.000000pt}%
\definecolor{currentstroke}{rgb}{0.000000,0.000000,0.000000}%
\pgfsetstrokecolor{currentstroke}%
\pgfsetstrokeopacity{0.000000}%
\pgfsetdash{}{0pt}%
\pgfpathmoveto{\pgfqpoint{3.142620in}{1.451438in}}%
\pgfpathlineto{\pgfqpoint{3.108264in}{1.451438in}}%
\pgfpathlineto{\pgfqpoint{3.108264in}{1.488093in}}%
\pgfpathlineto{\pgfqpoint{3.142620in}{1.488093in}}%
\pgfpathlineto{\pgfqpoint{3.142620in}{1.451438in}}%
\pgfpathclose%
\pgfusepath{fill}%
\end{pgfscope}%
\begin{pgfscope}%
\pgfpathrectangle{\pgfqpoint{2.893540in}{0.169444in}}{\pgfqpoint{0.858895in}{1.282944in}}%
\pgfusepath{clip}%
\pgfsetbuttcap%
\pgfsetmiterjoin%
\definecolor{currentfill}{rgb}{0.000000,0.392157,0.000000}%
\pgfsetfillcolor{currentfill}%
\pgfsetlinewidth{0.000000pt}%
\definecolor{currentstroke}{rgb}{0.000000,0.000000,0.000000}%
\pgfsetstrokecolor{currentstroke}%
\pgfsetstrokeopacity{0.000000}%
\pgfsetdash{}{0pt}%
\pgfpathmoveto{\pgfqpoint{3.142620in}{1.486564in}}%
\pgfpathlineto{\pgfqpoint{3.108264in}{1.486564in}}%
\pgfpathlineto{\pgfqpoint{3.108264in}{1.523220in}}%
\pgfpathlineto{\pgfqpoint{3.142620in}{1.523220in}}%
\pgfpathlineto{\pgfqpoint{3.142620in}{1.486564in}}%
\pgfpathclose%
\pgfusepath{fill}%
\end{pgfscope}%
\begin{pgfscope}%
\pgfpathrectangle{\pgfqpoint{2.893540in}{0.169444in}}{\pgfqpoint{0.858895in}{1.282944in}}%
\pgfusepath{clip}%
\pgfsetbuttcap%
\pgfsetmiterjoin%
\definecolor{currentfill}{rgb}{0.000000,0.392157,0.000000}%
\pgfsetfillcolor{currentfill}%
\pgfsetlinewidth{0.000000pt}%
\definecolor{currentstroke}{rgb}{0.000000,0.000000,0.000000}%
\pgfsetstrokecolor{currentstroke}%
\pgfsetstrokeopacity{0.000000}%
\pgfsetdash{}{0pt}%
\pgfpathmoveto{\pgfqpoint{3.142620in}{1.521691in}}%
\pgfpathlineto{\pgfqpoint{3.108264in}{1.521691in}}%
\pgfpathlineto{\pgfqpoint{3.108264in}{1.558347in}}%
\pgfpathlineto{\pgfqpoint{3.142620in}{1.558347in}}%
\pgfpathlineto{\pgfqpoint{3.142620in}{1.521691in}}%
\pgfpathclose%
\pgfusepath{fill}%
\end{pgfscope}%
\begin{pgfscope}%
\pgfpathrectangle{\pgfqpoint{2.893540in}{0.169444in}}{\pgfqpoint{0.858895in}{1.282944in}}%
\pgfusepath{clip}%
\pgfsetbuttcap%
\pgfsetmiterjoin%
\definecolor{currentfill}{rgb}{0.221945,0.790634,0.968474}%
\pgfsetfillcolor{currentfill}%
\pgfsetlinewidth{0.000000pt}%
\definecolor{currentstroke}{rgb}{0.000000,0.000000,0.000000}%
\pgfsetstrokecolor{currentstroke}%
\pgfsetstrokeopacity{0.000000}%
\pgfsetdash{}{0pt}%
\pgfpathmoveto{\pgfqpoint{3.013786in}{0.708314in}}%
\pgfpathlineto{\pgfqpoint{3.065319in}{0.708314in}}%
\pgfpathlineto{\pgfqpoint{3.065319in}{0.747191in}}%
\pgfpathlineto{\pgfqpoint{3.013786in}{0.747191in}}%
\pgfpathlineto{\pgfqpoint{3.013786in}{0.708314in}}%
\pgfpathclose%
\pgfusepath{fill}%
\end{pgfscope}%
\begin{pgfscope}%
\pgfpathrectangle{\pgfqpoint{2.893540in}{0.169444in}}{\pgfqpoint{0.858895in}{1.282944in}}%
\pgfusepath{clip}%
\pgfsetbuttcap%
\pgfsetmiterjoin%
\definecolor{currentfill}{rgb}{0.124506,0.772364,0.982314}%
\pgfsetfillcolor{currentfill}%
\pgfsetlinewidth{0.000000pt}%
\definecolor{currentstroke}{rgb}{0.000000,0.000000,0.000000}%
\pgfsetstrokecolor{currentstroke}%
\pgfsetstrokeopacity{0.000000}%
\pgfsetdash{}{0pt}%
\pgfpathmoveto{\pgfqpoint{3.030964in}{0.751101in}}%
\pgfpathlineto{\pgfqpoint{3.065319in}{0.751101in}}%
\pgfpathlineto{\pgfqpoint{3.065319in}{0.789979in}}%
\pgfpathlineto{\pgfqpoint{3.030964in}{0.789979in}}%
\pgfpathlineto{\pgfqpoint{3.030964in}{0.751101in}}%
\pgfpathclose%
\pgfusepath{fill}%
\end{pgfscope}%
\begin{pgfscope}%
\pgfpathrectangle{\pgfqpoint{2.893540in}{0.169444in}}{\pgfqpoint{0.858895in}{1.282944in}}%
\pgfusepath{clip}%
\pgfsetbuttcap%
\pgfsetmiterjoin%
\definecolor{currentfill}{rgb}{0.037893,0.756125,0.994617}%
\pgfsetfillcolor{currentfill}%
\pgfsetlinewidth{0.000000pt}%
\definecolor{currentstroke}{rgb}{0.000000,0.000000,0.000000}%
\pgfsetstrokecolor{currentstroke}%
\pgfsetstrokeopacity{0.000000}%
\pgfsetdash{}{0pt}%
\pgfpathmoveto{\pgfqpoint{3.048142in}{0.793889in}}%
\pgfpathlineto{\pgfqpoint{3.065319in}{0.793889in}}%
\pgfpathlineto{\pgfqpoint{3.065319in}{0.832766in}}%
\pgfpathlineto{\pgfqpoint{3.048142in}{0.832766in}}%
\pgfpathlineto{\pgfqpoint{3.048142in}{0.793889in}}%
\pgfpathclose%
\pgfusepath{fill}%
\end{pgfscope}%
\begin{pgfscope}%
\pgfpathrectangle{\pgfqpoint{2.893540in}{0.169444in}}{\pgfqpoint{0.858895in}{1.282944in}}%
\pgfusepath{clip}%
\pgfsetbuttcap%
\pgfsetmiterjoin%
\definecolor{currentfill}{rgb}{0.016917,0.753756,0.930642}%
\pgfsetfillcolor{currentfill}%
\pgfsetlinewidth{0.000000pt}%
\definecolor{currentstroke}{rgb}{0.000000,0.000000,0.000000}%
\pgfsetstrokecolor{currentstroke}%
\pgfsetstrokeopacity{0.000000}%
\pgfsetdash{}{0pt}%
\pgfpathmoveto{\pgfqpoint{3.013786in}{0.836723in}}%
\pgfpathlineto{\pgfqpoint{3.065319in}{0.836723in}}%
\pgfpathlineto{\pgfqpoint{3.065319in}{0.875600in}}%
\pgfpathlineto{\pgfqpoint{3.013786in}{0.875600in}}%
\pgfpathlineto{\pgfqpoint{3.013786in}{0.836723in}}%
\pgfpathclose%
\pgfusepath{fill}%
\end{pgfscope}%
\begin{pgfscope}%
\pgfpathrectangle{\pgfqpoint{2.893540in}{0.169444in}}{\pgfqpoint{0.858895in}{1.282944in}}%
\pgfusepath{clip}%
\pgfsetbuttcap%
\pgfsetmiterjoin%
\definecolor{currentfill}{rgb}{0.041522,0.760646,0.829758}%
\pgfsetfillcolor{currentfill}%
\pgfsetlinewidth{0.000000pt}%
\definecolor{currentstroke}{rgb}{0.000000,0.000000,0.000000}%
\pgfsetstrokecolor{currentstroke}%
\pgfsetstrokeopacity{0.000000}%
\pgfsetdash{}{0pt}%
\pgfpathmoveto{\pgfqpoint{2.996608in}{0.879556in}}%
\pgfpathlineto{\pgfqpoint{3.065319in}{0.879556in}}%
\pgfpathlineto{\pgfqpoint{3.065319in}{0.918433in}}%
\pgfpathlineto{\pgfqpoint{2.996608in}{0.918433in}}%
\pgfpathlineto{\pgfqpoint{2.996608in}{0.879556in}}%
\pgfpathclose%
\pgfusepath{fill}%
\end{pgfscope}%
\begin{pgfscope}%
\pgfpathrectangle{\pgfqpoint{2.893540in}{0.169444in}}{\pgfqpoint{0.858895in}{1.282944in}}%
\pgfusepath{clip}%
\pgfsetbuttcap%
\pgfsetmiterjoin%
\definecolor{currentfill}{rgb}{0.069204,0.768397,0.716263}%
\pgfsetfillcolor{currentfill}%
\pgfsetlinewidth{0.000000pt}%
\definecolor{currentstroke}{rgb}{0.000000,0.000000,0.000000}%
\pgfsetstrokecolor{currentstroke}%
\pgfsetstrokeopacity{0.000000}%
\pgfsetdash{}{0pt}%
\pgfpathmoveto{\pgfqpoint{3.013786in}{0.922344in}}%
\pgfpathlineto{\pgfqpoint{3.065319in}{0.922344in}}%
\pgfpathlineto{\pgfqpoint{3.065319in}{0.961221in}}%
\pgfpathlineto{\pgfqpoint{3.013786in}{0.961221in}}%
\pgfpathlineto{\pgfqpoint{3.013786in}{0.922344in}}%
\pgfpathclose%
\pgfusepath{fill}%
\end{pgfscope}%
\begin{pgfscope}%
\pgfpathrectangle{\pgfqpoint{2.893540in}{0.169444in}}{\pgfqpoint{0.858895in}{1.282944in}}%
\pgfusepath{clip}%
\pgfsetbuttcap%
\pgfsetmiterjoin%
\definecolor{currentfill}{rgb}{0.093810,0.775286,0.615379}%
\pgfsetfillcolor{currentfill}%
\pgfsetlinewidth{0.000000pt}%
\definecolor{currentstroke}{rgb}{0.000000,0.000000,0.000000}%
\pgfsetstrokecolor{currentstroke}%
\pgfsetstrokeopacity{0.000000}%
\pgfsetdash{}{0pt}%
\pgfpathmoveto{\pgfqpoint{2.996608in}{0.965131in}}%
\pgfpathlineto{\pgfqpoint{3.065319in}{0.965131in}}%
\pgfpathlineto{\pgfqpoint{3.065319in}{1.004008in}}%
\pgfpathlineto{\pgfqpoint{2.996608in}{1.004008in}}%
\pgfpathlineto{\pgfqpoint{2.996608in}{0.965131in}}%
\pgfpathclose%
\pgfusepath{fill}%
\end{pgfscope}%
\begin{pgfscope}%
\pgfpathrectangle{\pgfqpoint{2.893540in}{0.169444in}}{\pgfqpoint{0.858895in}{1.282944in}}%
\pgfusepath{clip}%
\pgfsetbuttcap%
\pgfsetmiterjoin%
\definecolor{currentfill}{rgb}{0.121492,0.783037,0.501884}%
\pgfsetfillcolor{currentfill}%
\pgfsetlinewidth{0.000000pt}%
\definecolor{currentstroke}{rgb}{0.000000,0.000000,0.000000}%
\pgfsetstrokecolor{currentstroke}%
\pgfsetstrokeopacity{0.000000}%
\pgfsetdash{}{0pt}%
\pgfpathmoveto{\pgfqpoint{2.962252in}{1.007965in}}%
\pgfpathlineto{\pgfqpoint{3.065319in}{1.007965in}}%
\pgfpathlineto{\pgfqpoint{3.065319in}{1.046842in}}%
\pgfpathlineto{\pgfqpoint{2.962252in}{1.046842in}}%
\pgfpathlineto{\pgfqpoint{2.962252in}{1.007965in}}%
\pgfpathclose%
\pgfusepath{fill}%
\end{pgfscope}%
\begin{pgfscope}%
\pgfpathrectangle{\pgfqpoint{2.893540in}{0.169444in}}{\pgfqpoint{0.858895in}{1.282944in}}%
\pgfusepath{clip}%
\pgfsetbuttcap%
\pgfsetmiterjoin%
\definecolor{currentfill}{rgb}{0.146098,0.789927,0.401000}%
\pgfsetfillcolor{currentfill}%
\pgfsetlinewidth{0.000000pt}%
\definecolor{currentstroke}{rgb}{0.000000,0.000000,0.000000}%
\pgfsetstrokecolor{currentstroke}%
\pgfsetstrokeopacity{0.000000}%
\pgfsetdash{}{0pt}%
\pgfpathmoveto{\pgfqpoint{2.979430in}{1.050798in}}%
\pgfpathlineto{\pgfqpoint{3.065319in}{1.050798in}}%
\pgfpathlineto{\pgfqpoint{3.065319in}{1.089675in}}%
\pgfpathlineto{\pgfqpoint{2.979430in}{1.089675in}}%
\pgfpathlineto{\pgfqpoint{2.979430in}{1.050798in}}%
\pgfpathclose%
\pgfusepath{fill}%
\end{pgfscope}%
\begin{pgfscope}%
\pgfpathrectangle{\pgfqpoint{2.893540in}{0.169444in}}{\pgfqpoint{0.858895in}{1.282944in}}%
\pgfusepath{clip}%
\pgfsetbuttcap%
\pgfsetmiterjoin%
\definecolor{currentfill}{rgb}{0.173779,0.797678,0.287505}%
\pgfsetfillcolor{currentfill}%
\pgfsetlinewidth{0.000000pt}%
\definecolor{currentstroke}{rgb}{0.000000,0.000000,0.000000}%
\pgfsetstrokecolor{currentstroke}%
\pgfsetstrokeopacity{0.000000}%
\pgfsetdash{}{0pt}%
\pgfpathmoveto{\pgfqpoint{2.945074in}{1.093586in}}%
\pgfpathlineto{\pgfqpoint{3.065319in}{1.093586in}}%
\pgfpathlineto{\pgfqpoint{3.065319in}{1.132463in}}%
\pgfpathlineto{\pgfqpoint{2.945074in}{1.132463in}}%
\pgfpathlineto{\pgfqpoint{2.945074in}{1.093586in}}%
\pgfpathclose%
\pgfusepath{fill}%
\end{pgfscope}%
\begin{pgfscope}%
\pgfpathrectangle{\pgfqpoint{2.893540in}{0.169444in}}{\pgfqpoint{0.858895in}{1.282944in}}%
\pgfusepath{clip}%
\pgfsetbuttcap%
\pgfsetmiterjoin%
\definecolor{currentfill}{rgb}{0.193772,0.799077,0.193772}%
\pgfsetfillcolor{currentfill}%
\pgfsetlinewidth{0.000000pt}%
\definecolor{currentstroke}{rgb}{0.000000,0.000000,0.000000}%
\pgfsetstrokecolor{currentstroke}%
\pgfsetstrokeopacity{0.000000}%
\pgfsetdash{}{0pt}%
\pgfpathmoveto{\pgfqpoint{3.048142in}{1.136373in}}%
\pgfpathlineto{\pgfqpoint{3.065319in}{1.136373in}}%
\pgfpathlineto{\pgfqpoint{3.065319in}{1.175251in}}%
\pgfpathlineto{\pgfqpoint{3.048142in}{1.175251in}}%
\pgfpathlineto{\pgfqpoint{3.048142in}{1.136373in}}%
\pgfpathclose%
\pgfusepath{fill}%
\end{pgfscope}%
\begin{pgfscope}%
\pgfpathrectangle{\pgfqpoint{2.893540in}{0.169444in}}{\pgfqpoint{0.858895in}{1.282944in}}%
\pgfusepath{clip}%
\pgfsetbuttcap%
\pgfsetmiterjoin%
\definecolor{currentfill}{rgb}{0.166090,0.740946,0.166090}%
\pgfsetfillcolor{currentfill}%
\pgfsetlinewidth{0.000000pt}%
\definecolor{currentstroke}{rgb}{0.000000,0.000000,0.000000}%
\pgfsetstrokecolor{currentstroke}%
\pgfsetstrokeopacity{0.000000}%
\pgfsetdash{}{0pt}%
\pgfpathmoveto{\pgfqpoint{2.962252in}{1.179207in}}%
\pgfpathlineto{\pgfqpoint{3.065319in}{1.179207in}}%
\pgfpathlineto{\pgfqpoint{3.065319in}{1.218084in}}%
\pgfpathlineto{\pgfqpoint{2.962252in}{1.218084in}}%
\pgfpathlineto{\pgfqpoint{2.962252in}{1.179207in}}%
\pgfpathclose%
\pgfusepath{fill}%
\end{pgfscope}%
\begin{pgfscope}%
\pgfpathrectangle{\pgfqpoint{2.893540in}{0.169444in}}{\pgfqpoint{0.858895in}{1.282944in}}%
\pgfusepath{clip}%
\pgfsetbuttcap%
\pgfsetmiterjoin%
\definecolor{currentfill}{rgb}{0.138408,0.682814,0.138408}%
\pgfsetfillcolor{currentfill}%
\pgfsetlinewidth{0.000000pt}%
\definecolor{currentstroke}{rgb}{0.000000,0.000000,0.000000}%
\pgfsetstrokecolor{currentstroke}%
\pgfsetstrokeopacity{0.000000}%
\pgfsetdash{}{0pt}%
\pgfpathmoveto{\pgfqpoint{3.048142in}{1.222040in}}%
\pgfpathlineto{\pgfqpoint{3.065319in}{1.222040in}}%
\pgfpathlineto{\pgfqpoint{3.065319in}{1.260917in}}%
\pgfpathlineto{\pgfqpoint{3.048142in}{1.260917in}}%
\pgfpathlineto{\pgfqpoint{3.048142in}{1.222040in}}%
\pgfpathclose%
\pgfusepath{fill}%
\end{pgfscope}%
\begin{pgfscope}%
\pgfpathrectangle{\pgfqpoint{2.893540in}{0.169444in}}{\pgfqpoint{0.858895in}{1.282944in}}%
\pgfusepath{clip}%
\pgfsetbuttcap%
\pgfsetmiterjoin%
\definecolor{currentfill}{rgb}{0.113802,0.631142,0.113802}%
\pgfsetfillcolor{currentfill}%
\pgfsetlinewidth{0.000000pt}%
\definecolor{currentstroke}{rgb}{0.000000,0.000000,0.000000}%
\pgfsetstrokecolor{currentstroke}%
\pgfsetstrokeopacity{0.000000}%
\pgfsetdash{}{0pt}%
\pgfpathmoveto{\pgfqpoint{3.048142in}{1.264828in}}%
\pgfpathlineto{\pgfqpoint{3.065319in}{1.264828in}}%
\pgfpathlineto{\pgfqpoint{3.065319in}{1.303705in}}%
\pgfpathlineto{\pgfqpoint{3.048142in}{1.303705in}}%
\pgfpathlineto{\pgfqpoint{3.048142in}{1.264828in}}%
\pgfpathclose%
\pgfusepath{fill}%
\end{pgfscope}%
\begin{pgfscope}%
\pgfpathrectangle{\pgfqpoint{2.893540in}{0.169444in}}{\pgfqpoint{0.858895in}{1.282944in}}%
\pgfusepath{clip}%
\pgfsetbuttcap%
\pgfsetmiterjoin%
\definecolor{currentfill}{rgb}{0.086121,0.573010,0.086121}%
\pgfsetfillcolor{currentfill}%
\pgfsetlinewidth{0.000000pt}%
\definecolor{currentstroke}{rgb}{0.000000,0.000000,0.000000}%
\pgfsetstrokecolor{currentstroke}%
\pgfsetstrokeopacity{0.000000}%
\pgfsetdash{}{0pt}%
\pgfpathmoveto{\pgfqpoint{2.996608in}{1.307616in}}%
\pgfpathlineto{\pgfqpoint{3.065319in}{1.307616in}}%
\pgfpathlineto{\pgfqpoint{3.065319in}{1.346493in}}%
\pgfpathlineto{\pgfqpoint{2.996608in}{1.346493in}}%
\pgfpathlineto{\pgfqpoint{2.996608in}{1.307616in}}%
\pgfpathclose%
\pgfusepath{fill}%
\end{pgfscope}%
\begin{pgfscope}%
\pgfpathrectangle{\pgfqpoint{2.893540in}{0.169444in}}{\pgfqpoint{0.858895in}{1.282944in}}%
\pgfusepath{clip}%
\pgfsetbuttcap%
\pgfsetmiterjoin%
\definecolor{currentfill}{rgb}{1.000000,1.000000,1.000000}%
\pgfsetfillcolor{currentfill}%
\pgfsetlinewidth{0.000000pt}%
\definecolor{currentstroke}{rgb}{0.000000,0.000000,0.000000}%
\pgfsetstrokecolor{currentstroke}%
\pgfsetstrokeopacity{0.000000}%
\pgfsetdash{}{0pt}%
\pgfpathmoveto{\pgfqpoint{3.374522in}{0.151754in}}%
\pgfpathlineto{\pgfqpoint{3.340166in}{0.151754in}}%
\pgfpathlineto{\pgfqpoint{3.340166in}{0.188409in}}%
\pgfpathlineto{\pgfqpoint{3.374522in}{0.188409in}}%
\pgfpathlineto{\pgfqpoint{3.374522in}{0.151754in}}%
\pgfpathclose%
\pgfusepath{fill}%
\end{pgfscope}%
\begin{pgfscope}%
\pgfpathrectangle{\pgfqpoint{2.893540in}{0.169444in}}{\pgfqpoint{0.858895in}{1.282944in}}%
\pgfusepath{clip}%
\pgfsetbuttcap%
\pgfsetmiterjoin%
\definecolor{currentfill}{rgb}{0.985421,0.994279,0.995386}%
\pgfsetfillcolor{currentfill}%
\pgfsetlinewidth{0.000000pt}%
\definecolor{currentstroke}{rgb}{0.000000,0.000000,0.000000}%
\pgfsetstrokecolor{currentstroke}%
\pgfsetstrokeopacity{0.000000}%
\pgfsetdash{}{0pt}%
\pgfpathmoveto{\pgfqpoint{3.374522in}{0.186881in}}%
\pgfpathlineto{\pgfqpoint{3.340166in}{0.186881in}}%
\pgfpathlineto{\pgfqpoint{3.340166in}{0.223536in}}%
\pgfpathlineto{\pgfqpoint{3.374522in}{0.223536in}}%
\pgfpathlineto{\pgfqpoint{3.374522in}{0.186881in}}%
\pgfpathclose%
\pgfusepath{fill}%
\end{pgfscope}%
\begin{pgfscope}%
\pgfpathrectangle{\pgfqpoint{2.893540in}{0.169444in}}{\pgfqpoint{0.858895in}{1.282944in}}%
\pgfusepath{clip}%
\pgfsetbuttcap%
\pgfsetmiterjoin%
\definecolor{currentfill}{rgb}{0.951403,0.980930,0.984621}%
\pgfsetfillcolor{currentfill}%
\pgfsetlinewidth{0.000000pt}%
\definecolor{currentstroke}{rgb}{0.000000,0.000000,0.000000}%
\pgfsetstrokecolor{currentstroke}%
\pgfsetstrokeopacity{0.000000}%
\pgfsetdash{}{0pt}%
\pgfpathmoveto{\pgfqpoint{3.374522in}{0.222007in}}%
\pgfpathlineto{\pgfqpoint{3.340166in}{0.222007in}}%
\pgfpathlineto{\pgfqpoint{3.340166in}{0.258663in}}%
\pgfpathlineto{\pgfqpoint{3.374522in}{0.258663in}}%
\pgfpathlineto{\pgfqpoint{3.374522in}{0.222007in}}%
\pgfpathclose%
\pgfusepath{fill}%
\end{pgfscope}%
\begin{pgfscope}%
\pgfpathrectangle{\pgfqpoint{2.893540in}{0.169444in}}{\pgfqpoint{0.858895in}{1.282944in}}%
\pgfusepath{clip}%
\pgfsetbuttcap%
\pgfsetmiterjoin%
\definecolor{currentfill}{rgb}{0.917386,0.967582,0.973856}%
\pgfsetfillcolor{currentfill}%
\pgfsetlinewidth{0.000000pt}%
\definecolor{currentstroke}{rgb}{0.000000,0.000000,0.000000}%
\pgfsetstrokecolor{currentstroke}%
\pgfsetstrokeopacity{0.000000}%
\pgfsetdash{}{0pt}%
\pgfpathmoveto{\pgfqpoint{3.374522in}{0.257134in}}%
\pgfpathlineto{\pgfqpoint{3.340166in}{0.257134in}}%
\pgfpathlineto{\pgfqpoint{3.340166in}{0.293789in}}%
\pgfpathlineto{\pgfqpoint{3.374522in}{0.293789in}}%
\pgfpathlineto{\pgfqpoint{3.374522in}{0.257134in}}%
\pgfpathclose%
\pgfusepath{fill}%
\end{pgfscope}%
\begin{pgfscope}%
\pgfpathrectangle{\pgfqpoint{2.893540in}{0.169444in}}{\pgfqpoint{0.858895in}{1.282944in}}%
\pgfusepath{clip}%
\pgfsetbuttcap%
\pgfsetmiterjoin%
\definecolor{currentfill}{rgb}{0.883368,0.954233,0.963091}%
\pgfsetfillcolor{currentfill}%
\pgfsetlinewidth{0.000000pt}%
\definecolor{currentstroke}{rgb}{0.000000,0.000000,0.000000}%
\pgfsetstrokecolor{currentstroke}%
\pgfsetstrokeopacity{0.000000}%
\pgfsetdash{}{0pt}%
\pgfpathmoveto{\pgfqpoint{3.374522in}{0.292260in}}%
\pgfpathlineto{\pgfqpoint{3.340166in}{0.292260in}}%
\pgfpathlineto{\pgfqpoint{3.340166in}{0.328916in}}%
\pgfpathlineto{\pgfqpoint{3.374522in}{0.328916in}}%
\pgfpathlineto{\pgfqpoint{3.374522in}{0.292260in}}%
\pgfpathclose%
\pgfusepath{fill}%
\end{pgfscope}%
\begin{pgfscope}%
\pgfpathrectangle{\pgfqpoint{2.893540in}{0.169444in}}{\pgfqpoint{0.858895in}{1.282944in}}%
\pgfusepath{clip}%
\pgfsetbuttcap%
\pgfsetmiterjoin%
\definecolor{currentfill}{rgb}{0.849350,0.940884,0.952326}%
\pgfsetfillcolor{currentfill}%
\pgfsetlinewidth{0.000000pt}%
\definecolor{currentstroke}{rgb}{0.000000,0.000000,0.000000}%
\pgfsetstrokecolor{currentstroke}%
\pgfsetstrokeopacity{0.000000}%
\pgfsetdash{}{0pt}%
\pgfpathmoveto{\pgfqpoint{3.374522in}{0.327387in}}%
\pgfpathlineto{\pgfqpoint{3.340166in}{0.327387in}}%
\pgfpathlineto{\pgfqpoint{3.340166in}{0.364042in}}%
\pgfpathlineto{\pgfqpoint{3.374522in}{0.364042in}}%
\pgfpathlineto{\pgfqpoint{3.374522in}{0.327387in}}%
\pgfpathclose%
\pgfusepath{fill}%
\end{pgfscope}%
\begin{pgfscope}%
\pgfpathrectangle{\pgfqpoint{2.893540in}{0.169444in}}{\pgfqpoint{0.858895in}{1.282944in}}%
\pgfusepath{clip}%
\pgfsetbuttcap%
\pgfsetmiterjoin%
\definecolor{currentfill}{rgb}{0.815333,0.927536,0.941561}%
\pgfsetfillcolor{currentfill}%
\pgfsetlinewidth{0.000000pt}%
\definecolor{currentstroke}{rgb}{0.000000,0.000000,0.000000}%
\pgfsetstrokecolor{currentstroke}%
\pgfsetstrokeopacity{0.000000}%
\pgfsetdash{}{0pt}%
\pgfpathmoveto{\pgfqpoint{3.374522in}{0.362514in}}%
\pgfpathlineto{\pgfqpoint{3.340166in}{0.362514in}}%
\pgfpathlineto{\pgfqpoint{3.340166in}{0.399169in}}%
\pgfpathlineto{\pgfqpoint{3.374522in}{0.399169in}}%
\pgfpathlineto{\pgfqpoint{3.374522in}{0.362514in}}%
\pgfpathclose%
\pgfusepath{fill}%
\end{pgfscope}%
\begin{pgfscope}%
\pgfpathrectangle{\pgfqpoint{2.893540in}{0.169444in}}{\pgfqpoint{0.858895in}{1.282944in}}%
\pgfusepath{clip}%
\pgfsetbuttcap%
\pgfsetmiterjoin%
\definecolor{currentfill}{rgb}{0.781315,0.914187,0.930796}%
\pgfsetfillcolor{currentfill}%
\pgfsetlinewidth{0.000000pt}%
\definecolor{currentstroke}{rgb}{0.000000,0.000000,0.000000}%
\pgfsetstrokecolor{currentstroke}%
\pgfsetstrokeopacity{0.000000}%
\pgfsetdash{}{0pt}%
\pgfpathmoveto{\pgfqpoint{3.374522in}{0.397640in}}%
\pgfpathlineto{\pgfqpoint{3.340166in}{0.397640in}}%
\pgfpathlineto{\pgfqpoint{3.340166in}{0.434296in}}%
\pgfpathlineto{\pgfqpoint{3.374522in}{0.434296in}}%
\pgfpathlineto{\pgfqpoint{3.374522in}{0.397640in}}%
\pgfpathclose%
\pgfusepath{fill}%
\end{pgfscope}%
\begin{pgfscope}%
\pgfpathrectangle{\pgfqpoint{2.893540in}{0.169444in}}{\pgfqpoint{0.858895in}{1.282944in}}%
\pgfusepath{clip}%
\pgfsetbuttcap%
\pgfsetmiterjoin%
\definecolor{currentfill}{rgb}{0.747297,0.900838,0.920031}%
\pgfsetfillcolor{currentfill}%
\pgfsetlinewidth{0.000000pt}%
\definecolor{currentstroke}{rgb}{0.000000,0.000000,0.000000}%
\pgfsetstrokecolor{currentstroke}%
\pgfsetstrokeopacity{0.000000}%
\pgfsetdash{}{0pt}%
\pgfpathmoveto{\pgfqpoint{3.374522in}{0.432767in}}%
\pgfpathlineto{\pgfqpoint{3.340166in}{0.432767in}}%
\pgfpathlineto{\pgfqpoint{3.340166in}{0.469422in}}%
\pgfpathlineto{\pgfqpoint{3.374522in}{0.469422in}}%
\pgfpathlineto{\pgfqpoint{3.374522in}{0.432767in}}%
\pgfpathclose%
\pgfusepath{fill}%
\end{pgfscope}%
\begin{pgfscope}%
\pgfpathrectangle{\pgfqpoint{2.893540in}{0.169444in}}{\pgfqpoint{0.858895in}{1.282944in}}%
\pgfusepath{clip}%
\pgfsetbuttcap%
\pgfsetmiterjoin%
\definecolor{currentfill}{rgb}{0.713280,0.887489,0.909266}%
\pgfsetfillcolor{currentfill}%
\pgfsetlinewidth{0.000000pt}%
\definecolor{currentstroke}{rgb}{0.000000,0.000000,0.000000}%
\pgfsetstrokecolor{currentstroke}%
\pgfsetstrokeopacity{0.000000}%
\pgfsetdash{}{0pt}%
\pgfpathmoveto{\pgfqpoint{3.374522in}{0.467893in}}%
\pgfpathlineto{\pgfqpoint{3.340166in}{0.467893in}}%
\pgfpathlineto{\pgfqpoint{3.340166in}{0.504549in}}%
\pgfpathlineto{\pgfqpoint{3.374522in}{0.504549in}}%
\pgfpathlineto{\pgfqpoint{3.374522in}{0.467893in}}%
\pgfpathclose%
\pgfusepath{fill}%
\end{pgfscope}%
\begin{pgfscope}%
\pgfpathrectangle{\pgfqpoint{2.893540in}{0.169444in}}{\pgfqpoint{0.858895in}{1.282944in}}%
\pgfusepath{clip}%
\pgfsetbuttcap%
\pgfsetmiterjoin%
\definecolor{currentfill}{rgb}{0.665836,0.873864,0.905421}%
\pgfsetfillcolor{currentfill}%
\pgfsetlinewidth{0.000000pt}%
\definecolor{currentstroke}{rgb}{0.000000,0.000000,0.000000}%
\pgfsetstrokecolor{currentstroke}%
\pgfsetstrokeopacity{0.000000}%
\pgfsetdash{}{0pt}%
\pgfpathmoveto{\pgfqpoint{3.374522in}{0.503020in}}%
\pgfpathlineto{\pgfqpoint{3.340166in}{0.503020in}}%
\pgfpathlineto{\pgfqpoint{3.340166in}{0.539675in}}%
\pgfpathlineto{\pgfqpoint{3.374522in}{0.539675in}}%
\pgfpathlineto{\pgfqpoint{3.374522in}{0.503020in}}%
\pgfpathclose%
\pgfusepath{fill}%
\end{pgfscope}%
\begin{pgfscope}%
\pgfpathrectangle{\pgfqpoint{2.893540in}{0.169444in}}{\pgfqpoint{0.858895in}{1.282944in}}%
\pgfusepath{clip}%
\pgfsetbuttcap%
\pgfsetmiterjoin%
\definecolor{currentfill}{rgb}{0.590050,0.859654,0.916186}%
\pgfsetfillcolor{currentfill}%
\pgfsetlinewidth{0.000000pt}%
\definecolor{currentstroke}{rgb}{0.000000,0.000000,0.000000}%
\pgfsetstrokecolor{currentstroke}%
\pgfsetstrokeopacity{0.000000}%
\pgfsetdash{}{0pt}%
\pgfpathmoveto{\pgfqpoint{3.374522in}{0.538146in}}%
\pgfpathlineto{\pgfqpoint{3.340166in}{0.538146in}}%
\pgfpathlineto{\pgfqpoint{3.340166in}{0.574802in}}%
\pgfpathlineto{\pgfqpoint{3.374522in}{0.574802in}}%
\pgfpathlineto{\pgfqpoint{3.374522in}{0.538146in}}%
\pgfpathclose%
\pgfusepath{fill}%
\end{pgfscope}%
\begin{pgfscope}%
\pgfpathrectangle{\pgfqpoint{2.893540in}{0.169444in}}{\pgfqpoint{0.858895in}{1.282944in}}%
\pgfusepath{clip}%
\pgfsetbuttcap%
\pgfsetmiterjoin%
\definecolor{currentfill}{rgb}{0.514264,0.845444,0.926951}%
\pgfsetfillcolor{currentfill}%
\pgfsetlinewidth{0.000000pt}%
\definecolor{currentstroke}{rgb}{0.000000,0.000000,0.000000}%
\pgfsetstrokecolor{currentstroke}%
\pgfsetstrokeopacity{0.000000}%
\pgfsetdash{}{0pt}%
\pgfpathmoveto{\pgfqpoint{3.374522in}{0.573273in}}%
\pgfpathlineto{\pgfqpoint{3.340166in}{0.573273in}}%
\pgfpathlineto{\pgfqpoint{3.340166in}{0.609929in}}%
\pgfpathlineto{\pgfqpoint{3.374522in}{0.609929in}}%
\pgfpathlineto{\pgfqpoint{3.374522in}{0.573273in}}%
\pgfpathclose%
\pgfusepath{fill}%
\end{pgfscope}%
\begin{pgfscope}%
\pgfpathrectangle{\pgfqpoint{2.893540in}{0.169444in}}{\pgfqpoint{0.858895in}{1.282944in}}%
\pgfusepath{clip}%
\pgfsetbuttcap%
\pgfsetmiterjoin%
\definecolor{currentfill}{rgb}{0.438478,0.831234,0.937716}%
\pgfsetfillcolor{currentfill}%
\pgfsetlinewidth{0.000000pt}%
\definecolor{currentstroke}{rgb}{0.000000,0.000000,0.000000}%
\pgfsetstrokecolor{currentstroke}%
\pgfsetstrokeopacity{0.000000}%
\pgfsetdash{}{0pt}%
\pgfpathmoveto{\pgfqpoint{3.374522in}{0.608400in}}%
\pgfpathlineto{\pgfqpoint{3.340166in}{0.608400in}}%
\pgfpathlineto{\pgfqpoint{3.340166in}{0.645055in}}%
\pgfpathlineto{\pgfqpoint{3.374522in}{0.645055in}}%
\pgfpathlineto{\pgfqpoint{3.374522in}{0.608400in}}%
\pgfpathclose%
\pgfusepath{fill}%
\end{pgfscope}%
\begin{pgfscope}%
\pgfpathrectangle{\pgfqpoint{2.893540in}{0.169444in}}{\pgfqpoint{0.858895in}{1.282944in}}%
\pgfusepath{clip}%
\pgfsetbuttcap%
\pgfsetmiterjoin%
\definecolor{currentfill}{rgb}{0.362691,0.817024,0.948481}%
\pgfsetfillcolor{currentfill}%
\pgfsetlinewidth{0.000000pt}%
\definecolor{currentstroke}{rgb}{0.000000,0.000000,0.000000}%
\pgfsetstrokecolor{currentstroke}%
\pgfsetstrokeopacity{0.000000}%
\pgfsetdash{}{0pt}%
\pgfpathmoveto{\pgfqpoint{3.374522in}{0.643526in}}%
\pgfpathlineto{\pgfqpoint{3.340166in}{0.643526in}}%
\pgfpathlineto{\pgfqpoint{3.340166in}{0.680182in}}%
\pgfpathlineto{\pgfqpoint{3.374522in}{0.680182in}}%
\pgfpathlineto{\pgfqpoint{3.374522in}{0.643526in}}%
\pgfpathclose%
\pgfusepath{fill}%
\end{pgfscope}%
\begin{pgfscope}%
\pgfpathrectangle{\pgfqpoint{2.893540in}{0.169444in}}{\pgfqpoint{0.858895in}{1.282944in}}%
\pgfusepath{clip}%
\pgfsetbuttcap%
\pgfsetmiterjoin%
\definecolor{currentfill}{rgb}{0.286905,0.802814,0.959246}%
\pgfsetfillcolor{currentfill}%
\pgfsetlinewidth{0.000000pt}%
\definecolor{currentstroke}{rgb}{0.000000,0.000000,0.000000}%
\pgfsetstrokecolor{currentstroke}%
\pgfsetstrokeopacity{0.000000}%
\pgfsetdash{}{0pt}%
\pgfpathmoveto{\pgfqpoint{3.374522in}{0.678653in}}%
\pgfpathlineto{\pgfqpoint{3.340166in}{0.678653in}}%
\pgfpathlineto{\pgfqpoint{3.340166in}{0.715308in}}%
\pgfpathlineto{\pgfqpoint{3.374522in}{0.715308in}}%
\pgfpathlineto{\pgfqpoint{3.374522in}{0.678653in}}%
\pgfpathclose%
\pgfusepath{fill}%
\end{pgfscope}%
\begin{pgfscope}%
\pgfpathrectangle{\pgfqpoint{2.893540in}{0.169444in}}{\pgfqpoint{0.858895in}{1.282944in}}%
\pgfusepath{clip}%
\pgfsetbuttcap%
\pgfsetmiterjoin%
\definecolor{currentfill}{rgb}{0.211119,0.788604,0.970012}%
\pgfsetfillcolor{currentfill}%
\pgfsetlinewidth{0.000000pt}%
\definecolor{currentstroke}{rgb}{0.000000,0.000000,0.000000}%
\pgfsetstrokecolor{currentstroke}%
\pgfsetstrokeopacity{0.000000}%
\pgfsetdash{}{0pt}%
\pgfpathmoveto{\pgfqpoint{3.374522in}{0.713779in}}%
\pgfpathlineto{\pgfqpoint{3.340166in}{0.713779in}}%
\pgfpathlineto{\pgfqpoint{3.340166in}{0.750435in}}%
\pgfpathlineto{\pgfqpoint{3.374522in}{0.750435in}}%
\pgfpathlineto{\pgfqpoint{3.374522in}{0.713779in}}%
\pgfpathclose%
\pgfusepath{fill}%
\end{pgfscope}%
\begin{pgfscope}%
\pgfpathrectangle{\pgfqpoint{2.893540in}{0.169444in}}{\pgfqpoint{0.858895in}{1.282944in}}%
\pgfusepath{clip}%
\pgfsetbuttcap%
\pgfsetmiterjoin%
\definecolor{currentfill}{rgb}{0.135333,0.774394,0.980777}%
\pgfsetfillcolor{currentfill}%
\pgfsetlinewidth{0.000000pt}%
\definecolor{currentstroke}{rgb}{0.000000,0.000000,0.000000}%
\pgfsetstrokecolor{currentstroke}%
\pgfsetstrokeopacity{0.000000}%
\pgfsetdash{}{0pt}%
\pgfpathmoveto{\pgfqpoint{3.374522in}{0.748906in}}%
\pgfpathlineto{\pgfqpoint{3.340166in}{0.748906in}}%
\pgfpathlineto{\pgfqpoint{3.340166in}{0.785562in}}%
\pgfpathlineto{\pgfqpoint{3.374522in}{0.785562in}}%
\pgfpathlineto{\pgfqpoint{3.374522in}{0.748906in}}%
\pgfpathclose%
\pgfusepath{fill}%
\end{pgfscope}%
\begin{pgfscope}%
\pgfpathrectangle{\pgfqpoint{2.893540in}{0.169444in}}{\pgfqpoint{0.858895in}{1.282944in}}%
\pgfusepath{clip}%
\pgfsetbuttcap%
\pgfsetmiterjoin%
\definecolor{currentfill}{rgb}{0.059546,0.760185,0.991542}%
\pgfsetfillcolor{currentfill}%
\pgfsetlinewidth{0.000000pt}%
\definecolor{currentstroke}{rgb}{0.000000,0.000000,0.000000}%
\pgfsetstrokecolor{currentstroke}%
\pgfsetstrokeopacity{0.000000}%
\pgfsetdash{}{0pt}%
\pgfpathmoveto{\pgfqpoint{3.374522in}{0.784033in}}%
\pgfpathlineto{\pgfqpoint{3.340166in}{0.784033in}}%
\pgfpathlineto{\pgfqpoint{3.340166in}{0.820688in}}%
\pgfpathlineto{\pgfqpoint{3.374522in}{0.820688in}}%
\pgfpathlineto{\pgfqpoint{3.374522in}{0.784033in}}%
\pgfpathclose%
\pgfusepath{fill}%
\end{pgfscope}%
\begin{pgfscope}%
\pgfpathrectangle{\pgfqpoint{2.893540in}{0.169444in}}{\pgfqpoint{0.858895in}{1.282944in}}%
\pgfusepath{clip}%
\pgfsetbuttcap%
\pgfsetmiterjoin%
\definecolor{currentfill}{rgb}{0.004614,0.750311,0.981084}%
\pgfsetfillcolor{currentfill}%
\pgfsetlinewidth{0.000000pt}%
\definecolor{currentstroke}{rgb}{0.000000,0.000000,0.000000}%
\pgfsetstrokecolor{currentstroke}%
\pgfsetstrokeopacity{0.000000}%
\pgfsetdash{}{0pt}%
\pgfpathmoveto{\pgfqpoint{3.374522in}{0.819159in}}%
\pgfpathlineto{\pgfqpoint{3.340166in}{0.819159in}}%
\pgfpathlineto{\pgfqpoint{3.340166in}{0.855815in}}%
\pgfpathlineto{\pgfqpoint{3.374522in}{0.855815in}}%
\pgfpathlineto{\pgfqpoint{3.374522in}{0.819159in}}%
\pgfpathclose%
\pgfusepath{fill}%
\end{pgfscope}%
\begin{pgfscope}%
\pgfpathrectangle{\pgfqpoint{2.893540in}{0.169444in}}{\pgfqpoint{0.858895in}{1.282944in}}%
\pgfusepath{clip}%
\pgfsetbuttcap%
\pgfsetmiterjoin%
\definecolor{currentfill}{rgb}{0.026144,0.756340,0.892810}%
\pgfsetfillcolor{currentfill}%
\pgfsetlinewidth{0.000000pt}%
\definecolor{currentstroke}{rgb}{0.000000,0.000000,0.000000}%
\pgfsetstrokecolor{currentstroke}%
\pgfsetstrokeopacity{0.000000}%
\pgfsetdash{}{0pt}%
\pgfpathmoveto{\pgfqpoint{3.374522in}{0.854286in}}%
\pgfpathlineto{\pgfqpoint{3.340166in}{0.854286in}}%
\pgfpathlineto{\pgfqpoint{3.340166in}{0.890941in}}%
\pgfpathlineto{\pgfqpoint{3.374522in}{0.890941in}}%
\pgfpathlineto{\pgfqpoint{3.374522in}{0.854286in}}%
\pgfpathclose%
\pgfusepath{fill}%
\end{pgfscope}%
\begin{pgfscope}%
\pgfpathrectangle{\pgfqpoint{2.893540in}{0.169444in}}{\pgfqpoint{0.858895in}{1.282944in}}%
\pgfusepath{clip}%
\pgfsetbuttcap%
\pgfsetmiterjoin%
\definecolor{currentfill}{rgb}{0.047674,0.762368,0.804537}%
\pgfsetfillcolor{currentfill}%
\pgfsetlinewidth{0.000000pt}%
\definecolor{currentstroke}{rgb}{0.000000,0.000000,0.000000}%
\pgfsetstrokecolor{currentstroke}%
\pgfsetstrokeopacity{0.000000}%
\pgfsetdash{}{0pt}%
\pgfpathmoveto{\pgfqpoint{3.374522in}{0.889412in}}%
\pgfpathlineto{\pgfqpoint{3.340166in}{0.889412in}}%
\pgfpathlineto{\pgfqpoint{3.340166in}{0.926068in}}%
\pgfpathlineto{\pgfqpoint{3.374522in}{0.926068in}}%
\pgfpathlineto{\pgfqpoint{3.374522in}{0.889412in}}%
\pgfpathclose%
\pgfusepath{fill}%
\end{pgfscope}%
\begin{pgfscope}%
\pgfpathrectangle{\pgfqpoint{2.893540in}{0.169444in}}{\pgfqpoint{0.858895in}{1.282944in}}%
\pgfusepath{clip}%
\pgfsetbuttcap%
\pgfsetmiterjoin%
\definecolor{currentfill}{rgb}{0.069204,0.768397,0.716263}%
\pgfsetfillcolor{currentfill}%
\pgfsetlinewidth{0.000000pt}%
\definecolor{currentstroke}{rgb}{0.000000,0.000000,0.000000}%
\pgfsetstrokecolor{currentstroke}%
\pgfsetstrokeopacity{0.000000}%
\pgfsetdash{}{0pt}%
\pgfpathmoveto{\pgfqpoint{3.374522in}{0.924539in}}%
\pgfpathlineto{\pgfqpoint{3.340166in}{0.924539in}}%
\pgfpathlineto{\pgfqpoint{3.340166in}{0.961195in}}%
\pgfpathlineto{\pgfqpoint{3.374522in}{0.961195in}}%
\pgfpathlineto{\pgfqpoint{3.374522in}{0.924539in}}%
\pgfpathclose%
\pgfusepath{fill}%
\end{pgfscope}%
\begin{pgfscope}%
\pgfpathrectangle{\pgfqpoint{2.893540in}{0.169444in}}{\pgfqpoint{0.858895in}{1.282944in}}%
\pgfusepath{clip}%
\pgfsetbuttcap%
\pgfsetmiterjoin%
\definecolor{currentfill}{rgb}{0.090734,0.774425,0.627989}%
\pgfsetfillcolor{currentfill}%
\pgfsetlinewidth{0.000000pt}%
\definecolor{currentstroke}{rgb}{0.000000,0.000000,0.000000}%
\pgfsetstrokecolor{currentstroke}%
\pgfsetstrokeopacity{0.000000}%
\pgfsetdash{}{0pt}%
\pgfpathmoveto{\pgfqpoint{3.374522in}{0.959666in}}%
\pgfpathlineto{\pgfqpoint{3.340166in}{0.959666in}}%
\pgfpathlineto{\pgfqpoint{3.340166in}{0.996321in}}%
\pgfpathlineto{\pgfqpoint{3.374522in}{0.996321in}}%
\pgfpathlineto{\pgfqpoint{3.374522in}{0.959666in}}%
\pgfpathclose%
\pgfusepath{fill}%
\end{pgfscope}%
\begin{pgfscope}%
\pgfpathrectangle{\pgfqpoint{2.893540in}{0.169444in}}{\pgfqpoint{0.858895in}{1.282944in}}%
\pgfusepath{clip}%
\pgfsetbuttcap%
\pgfsetmiterjoin%
\definecolor{currentfill}{rgb}{0.112265,0.780454,0.539715}%
\pgfsetfillcolor{currentfill}%
\pgfsetlinewidth{0.000000pt}%
\definecolor{currentstroke}{rgb}{0.000000,0.000000,0.000000}%
\pgfsetstrokecolor{currentstroke}%
\pgfsetstrokeopacity{0.000000}%
\pgfsetdash{}{0pt}%
\pgfpathmoveto{\pgfqpoint{3.374522in}{0.994792in}}%
\pgfpathlineto{\pgfqpoint{3.340166in}{0.994792in}}%
\pgfpathlineto{\pgfqpoint{3.340166in}{1.031448in}}%
\pgfpathlineto{\pgfqpoint{3.374522in}{1.031448in}}%
\pgfpathlineto{\pgfqpoint{3.374522in}{0.994792in}}%
\pgfpathclose%
\pgfusepath{fill}%
\end{pgfscope}%
\begin{pgfscope}%
\pgfpathrectangle{\pgfqpoint{2.893540in}{0.169444in}}{\pgfqpoint{0.858895in}{1.282944in}}%
\pgfusepath{clip}%
\pgfsetbuttcap%
\pgfsetmiterjoin%
\definecolor{currentfill}{rgb}{0.133795,0.786482,0.451442}%
\pgfsetfillcolor{currentfill}%
\pgfsetlinewidth{0.000000pt}%
\definecolor{currentstroke}{rgb}{0.000000,0.000000,0.000000}%
\pgfsetstrokecolor{currentstroke}%
\pgfsetstrokeopacity{0.000000}%
\pgfsetdash{}{0pt}%
\pgfpathmoveto{\pgfqpoint{3.374522in}{1.029919in}}%
\pgfpathlineto{\pgfqpoint{3.340166in}{1.029919in}}%
\pgfpathlineto{\pgfqpoint{3.340166in}{1.066574in}}%
\pgfpathlineto{\pgfqpoint{3.374522in}{1.066574in}}%
\pgfpathlineto{\pgfqpoint{3.374522in}{1.029919in}}%
\pgfpathclose%
\pgfusepath{fill}%
\end{pgfscope}%
\begin{pgfscope}%
\pgfpathrectangle{\pgfqpoint{2.893540in}{0.169444in}}{\pgfqpoint{0.858895in}{1.282944in}}%
\pgfusepath{clip}%
\pgfsetbuttcap%
\pgfsetmiterjoin%
\definecolor{currentfill}{rgb}{0.155325,0.792511,0.363168}%
\pgfsetfillcolor{currentfill}%
\pgfsetlinewidth{0.000000pt}%
\definecolor{currentstroke}{rgb}{0.000000,0.000000,0.000000}%
\pgfsetstrokecolor{currentstroke}%
\pgfsetstrokeopacity{0.000000}%
\pgfsetdash{}{0pt}%
\pgfpathmoveto{\pgfqpoint{3.374522in}{1.065045in}}%
\pgfpathlineto{\pgfqpoint{3.340166in}{1.065045in}}%
\pgfpathlineto{\pgfqpoint{3.340166in}{1.101701in}}%
\pgfpathlineto{\pgfqpoint{3.374522in}{1.101701in}}%
\pgfpathlineto{\pgfqpoint{3.374522in}{1.065045in}}%
\pgfpathclose%
\pgfusepath{fill}%
\end{pgfscope}%
\begin{pgfscope}%
\pgfpathrectangle{\pgfqpoint{2.893540in}{0.169444in}}{\pgfqpoint{0.858895in}{1.282944in}}%
\pgfusepath{clip}%
\pgfsetbuttcap%
\pgfsetmiterjoin%
\definecolor{currentfill}{rgb}{0.176855,0.798539,0.274894}%
\pgfsetfillcolor{currentfill}%
\pgfsetlinewidth{0.000000pt}%
\definecolor{currentstroke}{rgb}{0.000000,0.000000,0.000000}%
\pgfsetstrokecolor{currentstroke}%
\pgfsetstrokeopacity{0.000000}%
\pgfsetdash{}{0pt}%
\pgfpathmoveto{\pgfqpoint{3.374522in}{1.100172in}}%
\pgfpathlineto{\pgfqpoint{3.340166in}{1.100172in}}%
\pgfpathlineto{\pgfqpoint{3.340166in}{1.136828in}}%
\pgfpathlineto{\pgfqpoint{3.374522in}{1.136828in}}%
\pgfpathlineto{\pgfqpoint{3.374522in}{1.100172in}}%
\pgfpathclose%
\pgfusepath{fill}%
\end{pgfscope}%
\begin{pgfscope}%
\pgfpathrectangle{\pgfqpoint{2.893540in}{0.169444in}}{\pgfqpoint{0.858895in}{1.282944in}}%
\pgfusepath{clip}%
\pgfsetbuttcap%
\pgfsetmiterjoin%
\definecolor{currentfill}{rgb}{0.193772,0.799077,0.193772}%
\pgfsetfillcolor{currentfill}%
\pgfsetlinewidth{0.000000pt}%
\definecolor{currentstroke}{rgb}{0.000000,0.000000,0.000000}%
\pgfsetstrokecolor{currentstroke}%
\pgfsetstrokeopacity{0.000000}%
\pgfsetdash{}{0pt}%
\pgfpathmoveto{\pgfqpoint{3.374522in}{1.135299in}}%
\pgfpathlineto{\pgfqpoint{3.340166in}{1.135299in}}%
\pgfpathlineto{\pgfqpoint{3.340166in}{1.171954in}}%
\pgfpathlineto{\pgfqpoint{3.374522in}{1.171954in}}%
\pgfpathlineto{\pgfqpoint{3.374522in}{1.135299in}}%
\pgfpathclose%
\pgfusepath{fill}%
\end{pgfscope}%
\begin{pgfscope}%
\pgfpathrectangle{\pgfqpoint{2.893540in}{0.169444in}}{\pgfqpoint{0.858895in}{1.282944in}}%
\pgfusepath{clip}%
\pgfsetbuttcap%
\pgfsetmiterjoin%
\definecolor{currentfill}{rgb}{0.172241,0.753864,0.172241}%
\pgfsetfillcolor{currentfill}%
\pgfsetlinewidth{0.000000pt}%
\definecolor{currentstroke}{rgb}{0.000000,0.000000,0.000000}%
\pgfsetstrokecolor{currentstroke}%
\pgfsetstrokeopacity{0.000000}%
\pgfsetdash{}{0pt}%
\pgfpathmoveto{\pgfqpoint{3.374522in}{1.170425in}}%
\pgfpathlineto{\pgfqpoint{3.340166in}{1.170425in}}%
\pgfpathlineto{\pgfqpoint{3.340166in}{1.207081in}}%
\pgfpathlineto{\pgfqpoint{3.374522in}{1.207081in}}%
\pgfpathlineto{\pgfqpoint{3.374522in}{1.170425in}}%
\pgfpathclose%
\pgfusepath{fill}%
\end{pgfscope}%
\begin{pgfscope}%
\pgfpathrectangle{\pgfqpoint{2.893540in}{0.169444in}}{\pgfqpoint{0.858895in}{1.282944in}}%
\pgfusepath{clip}%
\pgfsetbuttcap%
\pgfsetmiterjoin%
\definecolor{currentfill}{rgb}{0.150711,0.708651,0.150711}%
\pgfsetfillcolor{currentfill}%
\pgfsetlinewidth{0.000000pt}%
\definecolor{currentstroke}{rgb}{0.000000,0.000000,0.000000}%
\pgfsetstrokecolor{currentstroke}%
\pgfsetstrokeopacity{0.000000}%
\pgfsetdash{}{0pt}%
\pgfpathmoveto{\pgfqpoint{3.374522in}{1.205552in}}%
\pgfpathlineto{\pgfqpoint{3.340166in}{1.205552in}}%
\pgfpathlineto{\pgfqpoint{3.340166in}{1.242207in}}%
\pgfpathlineto{\pgfqpoint{3.374522in}{1.242207in}}%
\pgfpathlineto{\pgfqpoint{3.374522in}{1.205552in}}%
\pgfpathclose%
\pgfusepath{fill}%
\end{pgfscope}%
\begin{pgfscope}%
\pgfpathrectangle{\pgfqpoint{2.893540in}{0.169444in}}{\pgfqpoint{0.858895in}{1.282944in}}%
\pgfusepath{clip}%
\pgfsetbuttcap%
\pgfsetmiterjoin%
\definecolor{currentfill}{rgb}{0.129181,0.663437,0.129181}%
\pgfsetfillcolor{currentfill}%
\pgfsetlinewidth{0.000000pt}%
\definecolor{currentstroke}{rgb}{0.000000,0.000000,0.000000}%
\pgfsetstrokecolor{currentstroke}%
\pgfsetstrokeopacity{0.000000}%
\pgfsetdash{}{0pt}%
\pgfpathmoveto{\pgfqpoint{3.374522in}{1.240678in}}%
\pgfpathlineto{\pgfqpoint{3.340166in}{1.240678in}}%
\pgfpathlineto{\pgfqpoint{3.340166in}{1.277334in}}%
\pgfpathlineto{\pgfqpoint{3.374522in}{1.277334in}}%
\pgfpathlineto{\pgfqpoint{3.374522in}{1.240678in}}%
\pgfpathclose%
\pgfusepath{fill}%
\end{pgfscope}%
\begin{pgfscope}%
\pgfpathrectangle{\pgfqpoint{2.893540in}{0.169444in}}{\pgfqpoint{0.858895in}{1.282944in}}%
\pgfusepath{clip}%
\pgfsetbuttcap%
\pgfsetmiterjoin%
\definecolor{currentfill}{rgb}{0.107651,0.618224,0.107651}%
\pgfsetfillcolor{currentfill}%
\pgfsetlinewidth{0.000000pt}%
\definecolor{currentstroke}{rgb}{0.000000,0.000000,0.000000}%
\pgfsetstrokecolor{currentstroke}%
\pgfsetstrokeopacity{0.000000}%
\pgfsetdash{}{0pt}%
\pgfpathmoveto{\pgfqpoint{3.374522in}{1.275805in}}%
\pgfpathlineto{\pgfqpoint{3.340166in}{1.275805in}}%
\pgfpathlineto{\pgfqpoint{3.340166in}{1.312460in}}%
\pgfpathlineto{\pgfqpoint{3.374522in}{1.312460in}}%
\pgfpathlineto{\pgfqpoint{3.374522in}{1.275805in}}%
\pgfpathclose%
\pgfusepath{fill}%
\end{pgfscope}%
\begin{pgfscope}%
\pgfpathrectangle{\pgfqpoint{2.893540in}{0.169444in}}{\pgfqpoint{0.858895in}{1.282944in}}%
\pgfusepath{clip}%
\pgfsetbuttcap%
\pgfsetmiterjoin%
\definecolor{currentfill}{rgb}{0.086121,0.573010,0.086121}%
\pgfsetfillcolor{currentfill}%
\pgfsetlinewidth{0.000000pt}%
\definecolor{currentstroke}{rgb}{0.000000,0.000000,0.000000}%
\pgfsetstrokecolor{currentstroke}%
\pgfsetstrokeopacity{0.000000}%
\pgfsetdash{}{0pt}%
\pgfpathmoveto{\pgfqpoint{3.374522in}{1.310932in}}%
\pgfpathlineto{\pgfqpoint{3.340166in}{1.310932in}}%
\pgfpathlineto{\pgfqpoint{3.340166in}{1.347587in}}%
\pgfpathlineto{\pgfqpoint{3.374522in}{1.347587in}}%
\pgfpathlineto{\pgfqpoint{3.374522in}{1.310932in}}%
\pgfpathclose%
\pgfusepath{fill}%
\end{pgfscope}%
\begin{pgfscope}%
\pgfpathrectangle{\pgfqpoint{2.893540in}{0.169444in}}{\pgfqpoint{0.858895in}{1.282944in}}%
\pgfusepath{clip}%
\pgfsetbuttcap%
\pgfsetmiterjoin%
\definecolor{currentfill}{rgb}{0.064591,0.527797,0.064591}%
\pgfsetfillcolor{currentfill}%
\pgfsetlinewidth{0.000000pt}%
\definecolor{currentstroke}{rgb}{0.000000,0.000000,0.000000}%
\pgfsetstrokecolor{currentstroke}%
\pgfsetstrokeopacity{0.000000}%
\pgfsetdash{}{0pt}%
\pgfpathmoveto{\pgfqpoint{3.374522in}{1.346058in}}%
\pgfpathlineto{\pgfqpoint{3.340166in}{1.346058in}}%
\pgfpathlineto{\pgfqpoint{3.340166in}{1.382714in}}%
\pgfpathlineto{\pgfqpoint{3.374522in}{1.382714in}}%
\pgfpathlineto{\pgfqpoint{3.374522in}{1.346058in}}%
\pgfpathclose%
\pgfusepath{fill}%
\end{pgfscope}%
\begin{pgfscope}%
\pgfpathrectangle{\pgfqpoint{2.893540in}{0.169444in}}{\pgfqpoint{0.858895in}{1.282944in}}%
\pgfusepath{clip}%
\pgfsetbuttcap%
\pgfsetmiterjoin%
\definecolor{currentfill}{rgb}{0.043060,0.482584,0.043060}%
\pgfsetfillcolor{currentfill}%
\pgfsetlinewidth{0.000000pt}%
\definecolor{currentstroke}{rgb}{0.000000,0.000000,0.000000}%
\pgfsetstrokecolor{currentstroke}%
\pgfsetstrokeopacity{0.000000}%
\pgfsetdash{}{0pt}%
\pgfpathmoveto{\pgfqpoint{3.374522in}{1.381185in}}%
\pgfpathlineto{\pgfqpoint{3.340166in}{1.381185in}}%
\pgfpathlineto{\pgfqpoint{3.340166in}{1.417840in}}%
\pgfpathlineto{\pgfqpoint{3.374522in}{1.417840in}}%
\pgfpathlineto{\pgfqpoint{3.374522in}{1.381185in}}%
\pgfpathclose%
\pgfusepath{fill}%
\end{pgfscope}%
\begin{pgfscope}%
\pgfpathrectangle{\pgfqpoint{2.893540in}{0.169444in}}{\pgfqpoint{0.858895in}{1.282944in}}%
\pgfusepath{clip}%
\pgfsetbuttcap%
\pgfsetmiterjoin%
\definecolor{currentfill}{rgb}{0.021530,0.437370,0.021530}%
\pgfsetfillcolor{currentfill}%
\pgfsetlinewidth{0.000000pt}%
\definecolor{currentstroke}{rgb}{0.000000,0.000000,0.000000}%
\pgfsetstrokecolor{currentstroke}%
\pgfsetstrokeopacity{0.000000}%
\pgfsetdash{}{0pt}%
\pgfpathmoveto{\pgfqpoint{3.374522in}{1.416311in}}%
\pgfpathlineto{\pgfqpoint{3.340166in}{1.416311in}}%
\pgfpathlineto{\pgfqpoint{3.340166in}{1.452967in}}%
\pgfpathlineto{\pgfqpoint{3.374522in}{1.452967in}}%
\pgfpathlineto{\pgfqpoint{3.374522in}{1.416311in}}%
\pgfpathclose%
\pgfusepath{fill}%
\end{pgfscope}%
\begin{pgfscope}%
\pgfpathrectangle{\pgfqpoint{2.893540in}{0.169444in}}{\pgfqpoint{0.858895in}{1.282944in}}%
\pgfusepath{clip}%
\pgfsetbuttcap%
\pgfsetmiterjoin%
\definecolor{currentfill}{rgb}{0.000000,0.392157,0.000000}%
\pgfsetfillcolor{currentfill}%
\pgfsetlinewidth{0.000000pt}%
\definecolor{currentstroke}{rgb}{0.000000,0.000000,0.000000}%
\pgfsetstrokecolor{currentstroke}%
\pgfsetstrokeopacity{0.000000}%
\pgfsetdash{}{0pt}%
\pgfpathmoveto{\pgfqpoint{3.374522in}{1.451438in}}%
\pgfpathlineto{\pgfqpoint{3.340166in}{1.451438in}}%
\pgfpathlineto{\pgfqpoint{3.340166in}{1.488093in}}%
\pgfpathlineto{\pgfqpoint{3.374522in}{1.488093in}}%
\pgfpathlineto{\pgfqpoint{3.374522in}{1.451438in}}%
\pgfpathclose%
\pgfusepath{fill}%
\end{pgfscope}%
\begin{pgfscope}%
\pgfpathrectangle{\pgfqpoint{2.893540in}{0.169444in}}{\pgfqpoint{0.858895in}{1.282944in}}%
\pgfusepath{clip}%
\pgfsetbuttcap%
\pgfsetmiterjoin%
\definecolor{currentfill}{rgb}{0.000000,0.392157,0.000000}%
\pgfsetfillcolor{currentfill}%
\pgfsetlinewidth{0.000000pt}%
\definecolor{currentstroke}{rgb}{0.000000,0.000000,0.000000}%
\pgfsetstrokecolor{currentstroke}%
\pgfsetstrokeopacity{0.000000}%
\pgfsetdash{}{0pt}%
\pgfpathmoveto{\pgfqpoint{3.374522in}{1.486564in}}%
\pgfpathlineto{\pgfqpoint{3.340166in}{1.486564in}}%
\pgfpathlineto{\pgfqpoint{3.340166in}{1.523220in}}%
\pgfpathlineto{\pgfqpoint{3.374522in}{1.523220in}}%
\pgfpathlineto{\pgfqpoint{3.374522in}{1.486564in}}%
\pgfpathclose%
\pgfusepath{fill}%
\end{pgfscope}%
\begin{pgfscope}%
\pgfpathrectangle{\pgfqpoint{2.893540in}{0.169444in}}{\pgfqpoint{0.858895in}{1.282944in}}%
\pgfusepath{clip}%
\pgfsetbuttcap%
\pgfsetmiterjoin%
\definecolor{currentfill}{rgb}{0.000000,0.392157,0.000000}%
\pgfsetfillcolor{currentfill}%
\pgfsetlinewidth{0.000000pt}%
\definecolor{currentstroke}{rgb}{0.000000,0.000000,0.000000}%
\pgfsetstrokecolor{currentstroke}%
\pgfsetstrokeopacity{0.000000}%
\pgfsetdash{}{0pt}%
\pgfpathmoveto{\pgfqpoint{3.374522in}{1.521691in}}%
\pgfpathlineto{\pgfqpoint{3.340166in}{1.521691in}}%
\pgfpathlineto{\pgfqpoint{3.340166in}{1.558347in}}%
\pgfpathlineto{\pgfqpoint{3.374522in}{1.558347in}}%
\pgfpathlineto{\pgfqpoint{3.374522in}{1.521691in}}%
\pgfpathclose%
\pgfusepath{fill}%
\end{pgfscope}%
\begin{pgfscope}%
\pgfpathrectangle{\pgfqpoint{2.893540in}{0.169444in}}{\pgfqpoint{0.858895in}{1.282944in}}%
\pgfusepath{clip}%
\pgfsetbuttcap%
\pgfsetmiterjoin%
\definecolor{currentfill}{rgb}{0.405998,0.825144,0.942330}%
\pgfsetfillcolor{currentfill}%
\pgfsetlinewidth{0.000000pt}%
\definecolor{currentstroke}{rgb}{0.000000,0.000000,0.000000}%
\pgfsetstrokecolor{currentstroke}%
\pgfsetstrokeopacity{0.000000}%
\pgfsetdash{}{0pt}%
\pgfpathmoveto{\pgfqpoint{3.408878in}{0.622647in}}%
\pgfpathlineto{\pgfqpoint{3.426056in}{0.622647in}}%
\pgfpathlineto{\pgfqpoint{3.426056in}{0.661524in}}%
\pgfpathlineto{\pgfqpoint{3.408878in}{0.661524in}}%
\pgfpathlineto{\pgfqpoint{3.408878in}{0.622647in}}%
\pgfpathclose%
\pgfusepath{fill}%
\end{pgfscope}%
\begin{pgfscope}%
\pgfpathrectangle{\pgfqpoint{2.893540in}{0.169444in}}{\pgfqpoint{0.858895in}{1.282944in}}%
\pgfusepath{clip}%
\pgfsetbuttcap%
\pgfsetmiterjoin%
\definecolor{currentfill}{rgb}{0.319385,0.808904,0.954633}%
\pgfsetfillcolor{currentfill}%
\pgfsetlinewidth{0.000000pt}%
\definecolor{currentstroke}{rgb}{0.000000,0.000000,0.000000}%
\pgfsetstrokecolor{currentstroke}%
\pgfsetstrokeopacity{0.000000}%
\pgfsetdash{}{0pt}%
\pgfpathmoveto{\pgfqpoint{3.408878in}{0.665480in}}%
\pgfpathlineto{\pgfqpoint{3.426056in}{0.665480in}}%
\pgfpathlineto{\pgfqpoint{3.426056in}{0.704357in}}%
\pgfpathlineto{\pgfqpoint{3.408878in}{0.704357in}}%
\pgfpathlineto{\pgfqpoint{3.408878in}{0.665480in}}%
\pgfpathclose%
\pgfusepath{fill}%
\end{pgfscope}%
\begin{pgfscope}%
\pgfpathrectangle{\pgfqpoint{2.893540in}{0.169444in}}{\pgfqpoint{0.858895in}{1.282944in}}%
\pgfusepath{clip}%
\pgfsetbuttcap%
\pgfsetmiterjoin%
\definecolor{currentfill}{rgb}{0.221945,0.790634,0.968474}%
\pgfsetfillcolor{currentfill}%
\pgfsetlinewidth{0.000000pt}%
\definecolor{currentstroke}{rgb}{0.000000,0.000000,0.000000}%
\pgfsetstrokecolor{currentstroke}%
\pgfsetstrokeopacity{0.000000}%
\pgfsetdash{}{0pt}%
\pgfpathmoveto{\pgfqpoint{3.408878in}{0.708314in}}%
\pgfpathlineto{\pgfqpoint{3.460411in}{0.708314in}}%
\pgfpathlineto{\pgfqpoint{3.460411in}{0.747191in}}%
\pgfpathlineto{\pgfqpoint{3.408878in}{0.747191in}}%
\pgfpathlineto{\pgfqpoint{3.408878in}{0.708314in}}%
\pgfpathclose%
\pgfusepath{fill}%
\end{pgfscope}%
\begin{pgfscope}%
\pgfpathrectangle{\pgfqpoint{2.893540in}{0.169444in}}{\pgfqpoint{0.858895in}{1.282944in}}%
\pgfusepath{clip}%
\pgfsetbuttcap%
\pgfsetmiterjoin%
\definecolor{currentfill}{rgb}{0.124506,0.772364,0.982314}%
\pgfsetfillcolor{currentfill}%
\pgfsetlinewidth{0.000000pt}%
\definecolor{currentstroke}{rgb}{0.000000,0.000000,0.000000}%
\pgfsetstrokecolor{currentstroke}%
\pgfsetstrokeopacity{0.000000}%
\pgfsetdash{}{0pt}%
\pgfpathmoveto{\pgfqpoint{3.408878in}{0.751101in}}%
\pgfpathlineto{\pgfqpoint{3.426056in}{0.751101in}}%
\pgfpathlineto{\pgfqpoint{3.426056in}{0.789979in}}%
\pgfpathlineto{\pgfqpoint{3.408878in}{0.789979in}}%
\pgfpathlineto{\pgfqpoint{3.408878in}{0.751101in}}%
\pgfpathclose%
\pgfusepath{fill}%
\end{pgfscope}%
\begin{pgfscope}%
\pgfpathrectangle{\pgfqpoint{2.893540in}{0.169444in}}{\pgfqpoint{0.858895in}{1.282944in}}%
\pgfusepath{clip}%
\pgfsetbuttcap%
\pgfsetmiterjoin%
\definecolor{currentfill}{rgb}{0.037893,0.756125,0.994617}%
\pgfsetfillcolor{currentfill}%
\pgfsetlinewidth{0.000000pt}%
\definecolor{currentstroke}{rgb}{0.000000,0.000000,0.000000}%
\pgfsetstrokecolor{currentstroke}%
\pgfsetstrokeopacity{0.000000}%
\pgfsetdash{}{0pt}%
\pgfpathmoveto{\pgfqpoint{3.408878in}{0.793889in}}%
\pgfpathlineto{\pgfqpoint{3.443233in}{0.793889in}}%
\pgfpathlineto{\pgfqpoint{3.443233in}{0.832766in}}%
\pgfpathlineto{\pgfqpoint{3.408878in}{0.832766in}}%
\pgfpathlineto{\pgfqpoint{3.408878in}{0.793889in}}%
\pgfpathclose%
\pgfusepath{fill}%
\end{pgfscope}%
\begin{pgfscope}%
\pgfpathrectangle{\pgfqpoint{2.893540in}{0.169444in}}{\pgfqpoint{0.858895in}{1.282944in}}%
\pgfusepath{clip}%
\pgfsetbuttcap%
\pgfsetmiterjoin%
\definecolor{currentfill}{rgb}{0.016917,0.753756,0.930642}%
\pgfsetfillcolor{currentfill}%
\pgfsetlinewidth{0.000000pt}%
\definecolor{currentstroke}{rgb}{0.000000,0.000000,0.000000}%
\pgfsetstrokecolor{currentstroke}%
\pgfsetstrokeopacity{0.000000}%
\pgfsetdash{}{0pt}%
\pgfpathmoveto{\pgfqpoint{3.408878in}{0.836723in}}%
\pgfpathlineto{\pgfqpoint{3.443233in}{0.836723in}}%
\pgfpathlineto{\pgfqpoint{3.443233in}{0.875600in}}%
\pgfpathlineto{\pgfqpoint{3.408878in}{0.875600in}}%
\pgfpathlineto{\pgfqpoint{3.408878in}{0.836723in}}%
\pgfpathclose%
\pgfusepath{fill}%
\end{pgfscope}%
\begin{pgfscope}%
\pgfpathrectangle{\pgfqpoint{2.893540in}{0.169444in}}{\pgfqpoint{0.858895in}{1.282944in}}%
\pgfusepath{clip}%
\pgfsetbuttcap%
\pgfsetmiterjoin%
\definecolor{currentfill}{rgb}{0.041522,0.760646,0.829758}%
\pgfsetfillcolor{currentfill}%
\pgfsetlinewidth{0.000000pt}%
\definecolor{currentstroke}{rgb}{0.000000,0.000000,0.000000}%
\pgfsetstrokecolor{currentstroke}%
\pgfsetstrokeopacity{0.000000}%
\pgfsetdash{}{0pt}%
\pgfpathmoveto{\pgfqpoint{3.408878in}{0.879556in}}%
\pgfpathlineto{\pgfqpoint{3.477589in}{0.879556in}}%
\pgfpathlineto{\pgfqpoint{3.477589in}{0.918433in}}%
\pgfpathlineto{\pgfqpoint{3.408878in}{0.918433in}}%
\pgfpathlineto{\pgfqpoint{3.408878in}{0.879556in}}%
\pgfpathclose%
\pgfusepath{fill}%
\end{pgfscope}%
\begin{pgfscope}%
\pgfpathrectangle{\pgfqpoint{2.893540in}{0.169444in}}{\pgfqpoint{0.858895in}{1.282944in}}%
\pgfusepath{clip}%
\pgfsetbuttcap%
\pgfsetmiterjoin%
\definecolor{currentfill}{rgb}{0.069204,0.768397,0.716263}%
\pgfsetfillcolor{currentfill}%
\pgfsetlinewidth{0.000000pt}%
\definecolor{currentstroke}{rgb}{0.000000,0.000000,0.000000}%
\pgfsetstrokecolor{currentstroke}%
\pgfsetstrokeopacity{0.000000}%
\pgfsetdash{}{0pt}%
\pgfpathmoveto{\pgfqpoint{3.408878in}{0.922344in}}%
\pgfpathlineto{\pgfqpoint{3.460411in}{0.922344in}}%
\pgfpathlineto{\pgfqpoint{3.460411in}{0.961221in}}%
\pgfpathlineto{\pgfqpoint{3.408878in}{0.961221in}}%
\pgfpathlineto{\pgfqpoint{3.408878in}{0.922344in}}%
\pgfpathclose%
\pgfusepath{fill}%
\end{pgfscope}%
\begin{pgfscope}%
\pgfpathrectangle{\pgfqpoint{2.893540in}{0.169444in}}{\pgfqpoint{0.858895in}{1.282944in}}%
\pgfusepath{clip}%
\pgfsetbuttcap%
\pgfsetmiterjoin%
\definecolor{currentfill}{rgb}{0.093810,0.775286,0.615379}%
\pgfsetfillcolor{currentfill}%
\pgfsetlinewidth{0.000000pt}%
\definecolor{currentstroke}{rgb}{0.000000,0.000000,0.000000}%
\pgfsetstrokecolor{currentstroke}%
\pgfsetstrokeopacity{0.000000}%
\pgfsetdash{}{0pt}%
\pgfpathmoveto{\pgfqpoint{3.408878in}{0.965131in}}%
\pgfpathlineto{\pgfqpoint{3.511945in}{0.965131in}}%
\pgfpathlineto{\pgfqpoint{3.511945in}{1.004008in}}%
\pgfpathlineto{\pgfqpoint{3.408878in}{1.004008in}}%
\pgfpathlineto{\pgfqpoint{3.408878in}{0.965131in}}%
\pgfpathclose%
\pgfusepath{fill}%
\end{pgfscope}%
\begin{pgfscope}%
\pgfpathrectangle{\pgfqpoint{2.893540in}{0.169444in}}{\pgfqpoint{0.858895in}{1.282944in}}%
\pgfusepath{clip}%
\pgfsetbuttcap%
\pgfsetmiterjoin%
\definecolor{currentfill}{rgb}{0.121492,0.783037,0.501884}%
\pgfsetfillcolor{currentfill}%
\pgfsetlinewidth{0.000000pt}%
\definecolor{currentstroke}{rgb}{0.000000,0.000000,0.000000}%
\pgfsetstrokecolor{currentstroke}%
\pgfsetstrokeopacity{0.000000}%
\pgfsetdash{}{0pt}%
\pgfpathmoveto{\pgfqpoint{3.408878in}{1.007965in}}%
\pgfpathlineto{\pgfqpoint{3.460411in}{1.007965in}}%
\pgfpathlineto{\pgfqpoint{3.460411in}{1.046842in}}%
\pgfpathlineto{\pgfqpoint{3.408878in}{1.046842in}}%
\pgfpathlineto{\pgfqpoint{3.408878in}{1.007965in}}%
\pgfpathclose%
\pgfusepath{fill}%
\end{pgfscope}%
\begin{pgfscope}%
\pgfpathrectangle{\pgfqpoint{2.893540in}{0.169444in}}{\pgfqpoint{0.858895in}{1.282944in}}%
\pgfusepath{clip}%
\pgfsetbuttcap%
\pgfsetmiterjoin%
\definecolor{currentfill}{rgb}{0.146098,0.789927,0.401000}%
\pgfsetfillcolor{currentfill}%
\pgfsetlinewidth{0.000000pt}%
\definecolor{currentstroke}{rgb}{0.000000,0.000000,0.000000}%
\pgfsetstrokecolor{currentstroke}%
\pgfsetstrokeopacity{0.000000}%
\pgfsetdash{}{0pt}%
\pgfpathmoveto{\pgfqpoint{3.408878in}{1.050798in}}%
\pgfpathlineto{\pgfqpoint{3.580657in}{1.050798in}}%
\pgfpathlineto{\pgfqpoint{3.580657in}{1.089675in}}%
\pgfpathlineto{\pgfqpoint{3.408878in}{1.089675in}}%
\pgfpathlineto{\pgfqpoint{3.408878in}{1.050798in}}%
\pgfpathclose%
\pgfusepath{fill}%
\end{pgfscope}%
\begin{pgfscope}%
\pgfpathrectangle{\pgfqpoint{2.893540in}{0.169444in}}{\pgfqpoint{0.858895in}{1.282944in}}%
\pgfusepath{clip}%
\pgfsetbuttcap%
\pgfsetmiterjoin%
\definecolor{currentfill}{rgb}{0.173779,0.797678,0.287505}%
\pgfsetfillcolor{currentfill}%
\pgfsetlinewidth{0.000000pt}%
\definecolor{currentstroke}{rgb}{0.000000,0.000000,0.000000}%
\pgfsetstrokecolor{currentstroke}%
\pgfsetstrokeopacity{0.000000}%
\pgfsetdash{}{0pt}%
\pgfpathmoveto{\pgfqpoint{3.408878in}{1.093586in}}%
\pgfpathlineto{\pgfqpoint{3.460411in}{1.093586in}}%
\pgfpathlineto{\pgfqpoint{3.460411in}{1.132463in}}%
\pgfpathlineto{\pgfqpoint{3.408878in}{1.132463in}}%
\pgfpathlineto{\pgfqpoint{3.408878in}{1.093586in}}%
\pgfpathclose%
\pgfusepath{fill}%
\end{pgfscope}%
\begin{pgfscope}%
\pgfpathrectangle{\pgfqpoint{2.893540in}{0.169444in}}{\pgfqpoint{0.858895in}{1.282944in}}%
\pgfusepath{clip}%
\pgfsetbuttcap%
\pgfsetmiterjoin%
\definecolor{currentfill}{rgb}{0.193772,0.799077,0.193772}%
\pgfsetfillcolor{currentfill}%
\pgfsetlinewidth{0.000000pt}%
\definecolor{currentstroke}{rgb}{0.000000,0.000000,0.000000}%
\pgfsetstrokecolor{currentstroke}%
\pgfsetstrokeopacity{0.000000}%
\pgfsetdash{}{0pt}%
\pgfpathmoveto{\pgfqpoint{3.408878in}{1.136373in}}%
\pgfpathlineto{\pgfqpoint{3.460411in}{1.136373in}}%
\pgfpathlineto{\pgfqpoint{3.460411in}{1.175251in}}%
\pgfpathlineto{\pgfqpoint{3.408878in}{1.175251in}}%
\pgfpathlineto{\pgfqpoint{3.408878in}{1.136373in}}%
\pgfpathclose%
\pgfusepath{fill}%
\end{pgfscope}%
\begin{pgfscope}%
\pgfpathrectangle{\pgfqpoint{2.893540in}{0.169444in}}{\pgfqpoint{0.858895in}{1.282944in}}%
\pgfusepath{clip}%
\pgfsetbuttcap%
\pgfsetmiterjoin%
\definecolor{currentfill}{rgb}{0.166090,0.740946,0.166090}%
\pgfsetfillcolor{currentfill}%
\pgfsetlinewidth{0.000000pt}%
\definecolor{currentstroke}{rgb}{0.000000,0.000000,0.000000}%
\pgfsetstrokecolor{currentstroke}%
\pgfsetstrokeopacity{0.000000}%
\pgfsetdash{}{0pt}%
\pgfpathmoveto{\pgfqpoint{3.408878in}{1.179207in}}%
\pgfpathlineto{\pgfqpoint{3.460411in}{1.179207in}}%
\pgfpathlineto{\pgfqpoint{3.460411in}{1.218084in}}%
\pgfpathlineto{\pgfqpoint{3.408878in}{1.218084in}}%
\pgfpathlineto{\pgfqpoint{3.408878in}{1.179207in}}%
\pgfpathclose%
\pgfusepath{fill}%
\end{pgfscope}%
\begin{pgfscope}%
\pgfpathrectangle{\pgfqpoint{2.893540in}{0.169444in}}{\pgfqpoint{0.858895in}{1.282944in}}%
\pgfusepath{clip}%
\pgfsetbuttcap%
\pgfsetmiterjoin%
\definecolor{currentfill}{rgb}{0.138408,0.682814,0.138408}%
\pgfsetfillcolor{currentfill}%
\pgfsetlinewidth{0.000000pt}%
\definecolor{currentstroke}{rgb}{0.000000,0.000000,0.000000}%
\pgfsetstrokecolor{currentstroke}%
\pgfsetstrokeopacity{0.000000}%
\pgfsetdash{}{0pt}%
\pgfpathmoveto{\pgfqpoint{3.408878in}{1.222040in}}%
\pgfpathlineto{\pgfqpoint{3.477589in}{1.222040in}}%
\pgfpathlineto{\pgfqpoint{3.477589in}{1.260917in}}%
\pgfpathlineto{\pgfqpoint{3.408878in}{1.260917in}}%
\pgfpathlineto{\pgfqpoint{3.408878in}{1.222040in}}%
\pgfpathclose%
\pgfusepath{fill}%
\end{pgfscope}%
\begin{pgfscope}%
\pgfpathrectangle{\pgfqpoint{2.893540in}{0.169444in}}{\pgfqpoint{0.858895in}{1.282944in}}%
\pgfusepath{clip}%
\pgfsetbuttcap%
\pgfsetmiterjoin%
\definecolor{currentfill}{rgb}{0.113802,0.631142,0.113802}%
\pgfsetfillcolor{currentfill}%
\pgfsetlinewidth{0.000000pt}%
\definecolor{currentstroke}{rgb}{0.000000,0.000000,0.000000}%
\pgfsetstrokecolor{currentstroke}%
\pgfsetstrokeopacity{0.000000}%
\pgfsetdash{}{0pt}%
\pgfpathmoveto{\pgfqpoint{3.408878in}{1.264828in}}%
\pgfpathlineto{\pgfqpoint{3.426056in}{1.264828in}}%
\pgfpathlineto{\pgfqpoint{3.426056in}{1.303705in}}%
\pgfpathlineto{\pgfqpoint{3.408878in}{1.303705in}}%
\pgfpathlineto{\pgfqpoint{3.408878in}{1.264828in}}%
\pgfpathclose%
\pgfusepath{fill}%
\end{pgfscope}%
\begin{pgfscope}%
\pgfpathrectangle{\pgfqpoint{2.893540in}{0.169444in}}{\pgfqpoint{0.858895in}{1.282944in}}%
\pgfusepath{clip}%
\pgfsetbuttcap%
\pgfsetmiterjoin%
\definecolor{currentfill}{rgb}{0.086121,0.573010,0.086121}%
\pgfsetfillcolor{currentfill}%
\pgfsetlinewidth{0.000000pt}%
\definecolor{currentstroke}{rgb}{0.000000,0.000000,0.000000}%
\pgfsetstrokecolor{currentstroke}%
\pgfsetstrokeopacity{0.000000}%
\pgfsetdash{}{0pt}%
\pgfpathmoveto{\pgfqpoint{3.408878in}{1.307616in}}%
\pgfpathlineto{\pgfqpoint{3.426056in}{1.307616in}}%
\pgfpathlineto{\pgfqpoint{3.426056in}{1.346493in}}%
\pgfpathlineto{\pgfqpoint{3.408878in}{1.346493in}}%
\pgfpathlineto{\pgfqpoint{3.408878in}{1.307616in}}%
\pgfpathclose%
\pgfusepath{fill}%
\end{pgfscope}%
\begin{pgfscope}%
\definecolor{textcolor}{rgb}{0.000000,0.000000,0.000000}%
\pgfsetstrokecolor{textcolor}%
\pgfsetfillcolor{textcolor}%
\pgftext[x=3.237099in,y=0.288748in,,]{\color{textcolor}\setmainfont{Lato}\rmfamily\fontsize{8.000000}{9.600000}\selectfont 65}%
\end{pgfscope}%
\begin{pgfscope}%
\definecolor{textcolor}{rgb}{0.000000,0.000000,0.000000}%
\pgfsetstrokecolor{textcolor}%
\pgfsetfillcolor{textcolor}%
\pgftext[x=3.237099in,y=0.517071in,,]{\color{textcolor}\setmainfont{Lato}\rmfamily\fontsize{8.000000}{9.600000}\selectfont 70}%
\end{pgfscope}%
\begin{pgfscope}%
\definecolor{textcolor}{rgb}{0.000000,0.000000,0.000000}%
\pgfsetstrokecolor{textcolor}%
\pgfsetfillcolor{textcolor}%
\pgftext[x=3.237099in,y=0.745393in,,]{\color{textcolor}\setmainfont{Lato}\rmfamily\fontsize{8.000000}{9.600000}\selectfont 75}%
\end{pgfscope}%
\begin{pgfscope}%
\definecolor{textcolor}{rgb}{0.000000,0.000000,0.000000}%
\pgfsetstrokecolor{textcolor}%
\pgfsetfillcolor{textcolor}%
\pgftext[x=3.237099in,y=0.973716in,,]{\color{textcolor}\setmainfont{Lato}\rmfamily\fontsize{8.000000}{9.600000}\selectfont 80}%
\end{pgfscope}%
\begin{pgfscope}%
\definecolor{textcolor}{rgb}{0.000000,0.000000,0.000000}%
\pgfsetstrokecolor{textcolor}%
\pgfsetfillcolor{textcolor}%
\pgftext[x=3.237099in,y=1.202039in,,]{\color{textcolor}\setmainfont{Lato}\rmfamily\fontsize{8.000000}{9.600000}\selectfont 85}%
\end{pgfscope}%
\begin{pgfscope}%
\definecolor{textcolor}{rgb}{0.000000,0.000000,0.000000}%
\pgfsetstrokecolor{textcolor}%
\pgfsetfillcolor{textcolor}%
\pgftext[x=3.237099in,y=1.430362in,,]{\color{textcolor}\setmainfont{Lato}\rmfamily\fontsize{8.000000}{9.600000}\selectfont 90}%
\end{pgfscope}%
\end{pgfpicture}%
\makeatother%
\endgroup%
 \\
\footnotesize{Source: Author's Calculations from CPS}
\vspace{2mm}

\normalsize \textbf{Employment Rate by State}\\
\footnotesize{\textit{employed share of age 25--54 population, percent, not seasonally adjusted}}\\
\vspace{-2mm}
\begin{minipage}{0.23\textwidth}
\hspace*{-3mm} \begin{tikzpicture}
	\begin{axis}[clip=false, stack negative=separate, ymajorgrids=false,
            yticklabels from table={\paepop}{name},
			y axis line style={opacity=0}, x axis line style={opacity=0}, 
			yticklabel style={black, font=\scriptsize, 
			align=left, anchor=east}, nodes near coords align={above},    
			nodes near coords style={/pgf/number format/.cd,fixed zerofill, 
				precision=1, assume math mode},  
            ytick=data, tickwidth=0pt, xmin=57, xmax=92,
            bar width=1.2ex, reverse legend, xmajorgrids, grid style={ultra thin, black!20},
            enlarge y limits={0.02}, enlarge x limits={0.02},
            legend cell align={left}, legend plot pos=left,
            width=4.6cm, height=10.7cm, legend style={legend columns=3, anchor=west, 
            at={(-0.1, 1.05)}},
            \barylab{0.5cm}{1.0ex}]
	\addplot[xbar stacked, bar shift=0pt, fill=black!3!white, draw=none] 
		table [y expr=-\coordindex, x index=1] {\paepop};
	\addplot[xbar stacked, bar shift=0pt, fill=cyan!45, draw=cyan!45] 
		table [y expr=-\coordindex, x index=2] {\paepop};	
	\addplot[xbar stacked, bar shift=0pt, fill=green!65!black, draw=cyan!45] 
		table [y expr=-\coordindex, x index=3] {\paepop};
	\addplot[xbar stacked, bar shift=0pt, fill=cyan!45, draw=cyan!45] 
		table [y expr=-\coordindex, x index=4] {\paepop};
	\addplot[xbar stacked, bar shift=0pt, fill=black!3!white, draw=none] 
		table [y expr=-\coordindex, x index=5] {\paepop};	        
	\addplot[only marks, mark=diamond*, mark options={fill=white, 
		draw=magenta, scale=1.6}] 
		table [y expr=-\coordindex, x index=6] {\paepop};
	\addplot [nodes near coords={\pgfmathprintnumber{\pgfkeysvalueof{/data point/x}}}, only marks, mark=none] table[y=y, x index=6] {\paepop};
	\legend{, 1989-- range, April 2020-- range,, ,\input{text/cps_lfs_date.txt}, };
	\end{axis}
\end{tikzpicture}\end{minipage} \hspace{1.5mm}
\begin{minipage}{0.25\textwidth}
\vspace{2.5mm}
\begin{tikzpicture}
	\begin{axis}[clip=false, stack negative=separate, ymajorgrids=false,
            yticklabels from table={\paepopt}{name},
			y axis line style={opacity=0}, x axis line style={opacity=0}, 
			yticklabel style={black, font=\scriptsize, 
			align=left, anchor=east}, nodes near coords align={above},    
			nodes near coords style={/pgf/number format/.cd,fixed zerofill, 
				precision=1, assume math mode},  
            ytick=data, tickwidth=0pt, xmin=57, xmax=92,
            bar width=1.2ex, xmajorgrids, grid style={ultra thin, black!20},
            enlarge y limits={0.02}, enlarge x limits={0.02},
            width=4.6cm, height=10.7cm, 
            \barylab{0.5cm}{1.0ex}]
	\addplot[xbar stacked, bar shift=0pt, fill=black!3!white, draw=none] 
		table [y expr=-\coordindex, x index=1] {\paepopt};
	\addplot[xbar stacked, bar shift=0pt, fill=cyan!45, draw=cyan!45] 
		table [y expr=-\coordindex, x index=2] {\paepopt};	
	\addplot[xbar stacked, bar shift=0pt, fill=green!65!black, draw=cyan!45] 
		table [y expr=-\coordindex, x index=3] {\paepopt};
	\addplot[xbar stacked, bar shift=0pt, fill=cyan!45, draw=cyan!45] 
		table [y expr=-\coordindex, x index=4] {\paepopt};
	\addplot[xbar stacked, bar shift=0pt, fill=black!3!white, draw=none] 
		table [y expr=-\coordindex, x index=5] {\paepopt};	        
	\addplot[only marks, mark=diamond*, mark options={fill=white, 
		draw=magenta, scale=1.6}] 
		table [y expr=-\coordindex, x index=6] {\paepopt};
	\addplot [nodes near coords={\pgfmathprintnumber{\pgfkeysvalueof{/data point/x}}}, only marks, mark=none] table[y=y, x index=6] {\paepopt};
	\end{axis}
\end{tikzpicture}\end{minipage}\hspace{2mm}
\begin{minipage}{0.25\textwidth}
\vspace{2.5mm}
\begin{tikzpicture}
	\begin{axis}[clip=false, stack negative=separate, ymajorgrids=false,
            yticklabels from table={\paepoptt}{name},
			y axis line style={opacity=0}, x axis line style={opacity=0}, 
			yticklabel style={black, font=\scriptsize, 
			align=left, anchor=east}, nodes near coords align={above},    
			nodes near coords style={/pgf/number format/.cd,fixed zerofill, 
				precision=1, assume math mode},  
            ytick=data, tickwidth=0pt, xmin=57, xmax=92,
            bar width=1.2ex, xmajorgrids, grid style={ultra thin, black!20},
            enlarge y limits={0.02}, enlarge x limits={0.02},
            width=4.6cm, height=10.7cm, 
            \barylab{0.5cm}{1.0ex}]
	\addplot[xbar stacked, bar shift=0pt, fill=black!3!white, draw=none] 
		table [y expr=-\coordindex, x index=1] {\paepoptt};
	\addplot[xbar stacked, bar shift=0pt, fill=cyan!45, draw=cyan!45] 
		table [y expr=-\coordindex, x index=2] {\paepoptt};	
	\addplot[xbar stacked, bar shift=0pt, fill=green!65!black, draw=cyan!45] 
		table [y expr=-\coordindex, x index=3] {\paepoptt};
	\addplot[xbar stacked, bar shift=0pt, fill=cyan!45, draw=cyan!45] 
		table [y expr=-\coordindex, x index=4] {\paepoptt};
	\addplot[xbar stacked, bar shift=0pt, fill=black!3!white, draw=none] 
		table [y expr=-\coordindex, x index=5] {\paepoptt};	        
	\addplot[only marks, mark=diamond*, mark options={fill=white, 
		draw=magenta, scale=1.6}] 
		table [y expr=-\coordindex, x index=6] {\paepoptt};
	\addplot [nodes near coords={\pgfmathprintnumber{\pgfkeysvalueof{/data point/x}}}, 
		only marks, mark=none] table[y=y, x index=6] {\paepoptt};
	\end{axis}
\end{tikzpicture}\end{minipage}
\vspace{1mm}

\footnotesize{Source: Author's Calculations from CPS} \hspace{67mm} \tbllink{state_pa_epop_full.csv}
\newpage
\subsubsection*{Payroll Employment}
\begin{minipage}{0.36\textwidth}
\small The \href{https://www.bls.gov/ces/}{establishment survey} from the monthly jobs report asks nonfarm businesses about their payrolls. \input{text/nfp_tot.txt} 

The establishment survey is particularly useful for identifying how many jobs were added or lost in each industry each month. \input{text/nfp_basic_text.txt}
\end{minipage} \hspace{6mm}
\begin{minipage}{0.35\textwidth}
\normalsize \textbf{Total Nonfarm Payroll Growth}\\
\footnotesize{\textit{one-month change, in millions}}\\
\hspace*{-3mm} \begin{tikzpicture}
	\begin{axis}[\bbar{y}{0}, \ltdateaxisticks width=6.6cm, height=8.8cm,
		minor xtick={}, xmin={2018-12-29}, yticklabel style={text width=1.4em},
		enlarge y limits={0.02}, xticklabel={`\short{\year}}]
    \rbar
	\ctsbar{blue!60!purple}{date}{ALL}{data/nfp.csv}{3.0pt}
	\absnode{{2020-05-20}}{-15.5}{May `20\\-20.7\\ million}
	\end{axis}
\end{tikzpicture}\\
\footnotesize{Source: Bureau of Labor Statistics} \hfill \tbllink{nfp.csv}
\end{minipage}
\vspace{3mm}

\normalsize \textbf{Nonfarm Payrolls by Industry Group}\\
\footnotesize{\textit{seasonally adjusted, thousands \hspace{8mm}levels \hspace{16mm} monthly change \hspace{25mm} cumulative}\\
\vspace{-5mm}

\hspace*{-3mm} \noindent \rowcolors{1}{}{black!5} \setlength{\tabcolsep}{3.0pt} \color{black!90}
		{\renewcommand{\arraystretch}{1.57}
		 \begin{tabular}{p{37mm} R{12.6mm} R{12.6mm} | R{9.0mm} R{9.0mm} R{9.0mm} 
		   R{9.0mm} | R{12.0mm} R{12.0mm}}
			 \input data/nfp.tex \hline
		\end{tabular}
		}	\\
		
\vspace{-2mm}
\footnotesize{Source: Bureau of Labor Statistics} \hspace{9.5cm} \tbllink{ces_data.csv}
\newpage
\begin{minipage}{0.76\textwidth} 
\small Employment rates vary over time, but also by age, gender, and education. Over the three months ending \input{text/cps_lfs_date.txt}\unskip, the employment rate for most age and education subgroups is lower than it was before the pandemic (the three months ending \input{text/cps_lfs_dateprepan.txt}\unskip).
\vspace{2mm}

\textbf{Employment Rates} \small \hspace{10mm} \cbox{blue} \input{text/cps_lfs_date.txt} \hspace{4mm} \cbox{cyan!50} \input{text/cps_lfs_dateprepan.txt}\\
\footnotesize{\textit{employed share of age group, percent}}\\
\vspace{2mm}

\centerline{\small Total, Any Education}
\hspace{22mm} Men \hspace{54mm} Women\\
\hspace*{-4mm} \begin{tikzpicture}
  	\begin{axis}[\barplotnogrid axis y line=left, \barylab{1.0cm}{1.5ex}
    	width=6.1cm, bar width=1.6ex, height=5.2cm,
    	enlarge y limits={abs=0.35cm}, enlarge x limits=0.05, \bbar{x}{0},
		yticklabels from table={\empt}{name}, xmin=2, xmax=100,
		yticklabel style={font=\footnotesize}, nodes near coords, 
		nodes near coords style={/pgf/number format/.cd, fixed zerofill,
			precision=1, assume math mode}]
  	\addplot[fill=cyan!50, draw=none] 
  		table [y expr=-\coordindex, x index=1] {\empt};
  	\addplot[fill=blue, draw=none] 
  		table [y expr=-\coordindex, x index=2] {\empt};
  	\end{axis}
\end{tikzpicture} \hfill
\begin{tikzpicture}
  	\begin{axis}[\barplotnogrid axis y line=left, \barylab{1.0cm}{1.5ex}
    	width=6.1cm, bar width=1.6ex, height=5.2cm,
    	enlarge y limits={abs=0.35cm}, enlarge x limits=0.05, \bbar{x}{0},
		yticklabels from table={\empt}{name}, xmin=2, xmax=100,
		yticklabel style={font=\footnotesize}, nodes near coords, 
		nodes near coords style={/pgf/number format/.cd, fixed zerofill,
			precision=1, assume math mode}]
  	\addplot[fill=cyan!50, draw=none] 
  		table [y expr=-\coordindex, x index=3] {\empt};
  	\addplot[fill=blue, draw=none] 
  		table [y expr=-\coordindex, x index=4] {\empt};
  	\end{axis}
\end{tikzpicture}\\
\vspace{1mm}

\centerline{\small High School Degree or Less}
\hspace{22mm} Men \hspace{54mm} Women\\
\hspace*{-4mm} \begin{tikzpicture}
  	\begin{axis}[\barplotnogrid axis y line=left, \barylab{1.0cm}{1.5ex}
    	width=6.1cm, bar width=1.6ex, height=4.6cm,
    	enlarge y limits={abs=0.35cm}, enlarge x limits=0.05, \bbar{x}{0},
		yticklabels from table={\emp}{name}, xmin=2, xmax=100,
		yticklabel style={font=\footnotesize}, nodes near coords, 
		nodes near coords style={/pgf/number format/.cd, fixed zerofill,
			precision=1, assume math mode}]
  	\addplot[fill=cyan!50, draw=none] 
  		table [y expr=-\coordindex, x index=1] {\emp};
  	\addplot[fill=blue, draw=none] 
  		table [y expr=-\coordindex, x index=2] {\emp};
  	\end{axis}
\end{tikzpicture} \hfill
\begin{tikzpicture}
  	\begin{axis}[\barplotnogrid axis y line=left, \barylab{1.0cm}{1.5ex}
    	width=6.1cm, bar width=1.6ex, height=4.6cm,
    	enlarge y limits={abs=0.35cm}, enlarge x limits=0.05, \bbar{x}{0},
		yticklabels from table={\emp}{name}, xmin=2, xmax=100,
		yticklabel style={font=\footnotesize}, nodes near coords, 
		nodes near coords style={/pgf/number format/.cd, fixed zerofill,
			precision=1, assume math mode}]
  	\addplot[fill=cyan!50, draw=none] 
  		table [y expr=-\coordindex, x index=3] {\emp};
  	\addplot[fill=blue, draw=none] 
  		table [y expr=-\coordindex, x index=4] {\emp};
  	\end{axis}
\end{tikzpicture}\\
\vspace{1mm}

\centerline{\small Some College}
\hspace{22mm} Men \hspace{54mm} Women\\
\noindent \hspace*{-4mm} \begin{tikzpicture}
  	\begin{axis}[\barplotnogrid axis y line=left, \barylab{1.0cm}{1.5ex}
    	width=6.1cm, bar width=1.6ex, height=4.6cm,
    	enlarge y limits={abs=0.35cm}, enlarge x limits=0.05, \bbar{x}{0},
		yticklabels from table={\emp}{name}, xmin=2, xmax=100,
		yticklabel style={font=\footnotesize}, nodes near coords, 
		nodes near coords style={/pgf/number format/.cd, fixed zerofill,
			precision=1, assume math mode}]
  	\addplot[fill=cyan!50, draw=none] 
  		table [y expr=-\coordindex, x index=5] {\emp};
  	\addplot[fill=blue, draw=none] 
  		table [y expr=-\coordindex, x index=6] {\emp};
  	\end{axis}
\end{tikzpicture} \hfill
\begin{tikzpicture}
  	\begin{axis}[\barplotnogrid axis y line=left, \barylab{1.0cm}{1.5ex}
    	width=6.1cm, bar width=1.6ex, height=4.6cm,
    	enlarge y limits={abs=0.35cm}, enlarge x limits=0.05, \bbar{x}{0},
		yticklabels from table={\emp}{name}, xmin=2, xmax=100,
		yticklabel style={font=\footnotesize}, nodes near coords, 
		nodes near coords style={/pgf/number format/.cd, fixed zerofill,
			precision=1, assume math mode}]
  	\addplot[fill=cyan!50, draw=none] 
  		table [y expr=-\coordindex, x index=7] {\emp};
  	\addplot[fill=blue, draw=none] 
  		table [y expr=-\coordindex, x index=8] {\emp};
  	\end{axis}
\end{tikzpicture}
\vspace{1mm}

\centerline{\small Bachelor or Advanced Degree}
\hspace{22mm} Men \hspace{54mm} Women \\
\hspace*{-4mm} \begin{tikzpicture}
  	\begin{axis}[\barplotnogrid axis y line=left, \barylab{1.0cm}{1.5ex}
    	width=6.1cm, bar width=1.6ex, height=4.6cm,
    	enlarge y limits={abs=0.35cm}, enlarge x limits=0.05, \bbar{x}{0},
		yticklabels from table={\emp}{name}, xmin=2, xmax=100,
		yticklabel style={font=\footnotesize}, nodes near coords, 
		nodes near coords style={/pgf/number format/.cd, fixed zerofill,
			precision=1, assume math mode}]
  	\addplot[fill=cyan!50, draw=none] 
  		table [y expr=-\coordindex, x index=9] {\emp};
  	\addplot[fill=blue, draw=none] 
  		table [y expr=-\coordindex, x index=10] {\emp};
  	\end{axis}
\end{tikzpicture} \hfill
\begin{tikzpicture}
  	\begin{axis}[\barplotnogrid axis y line=left, \barylab{1.0cm}{1.5ex}
    	width=6.1cm, bar width=1.6ex, height=4.6cm,
    	enlarge y limits={abs=0.35cm}, enlarge x limits=0.05, \bbar{x}{0},
		yticklabels from table={\emp}{name}, xmin=2, xmax=100,
		yticklabel style={font=\footnotesize}, nodes near coords, 
		nodes near coords style={/pgf/number format/.cd, fixed zerofill,
			precision=1, assume math mode}]
  	\addplot[fill=cyan!50, draw=none] 
  		table [y expr=-\coordindex, x index=11] {\emp};
  	\addplot[fill=blue, draw=none] 
  		table [y expr=-\coordindex, x index=12] {\emp};
  	\end{axis}
\end{tikzpicture}
\vspace{-1mm}

\footnotesize{Source: Author's Calculations from CPS} \hfill \tbllink{empgroups.csv}
\end{minipage}
\newpage
\subsubsection*{Changes in Employment by Occupation}
\begin{minipage}{0.76\textwidth} 
\small The COVID-19 pandemic changed the occupational composition of the US workforce. Business closures reduced the share of the US population that is employed in personal care and related services, administrative support, and sales. Many workers trained in these occupations lost or changed jobs and, additionally, some full-time workers became part-time. While employment increased in some occupations, such as healthcare support, the share of people employed full-time decreased in most occupational groups during the past three months, which end \input{text/cps_lfs_date.txt}\unskip, compared to the same three months prior to the pandemic (ending \input{text/cps_lfs_dateprepan.txt}\unskip).
\vspace{2mm}

\normalsize \textbf{Change in Occupational Employment, \input{text/cps_lfs_date.txt}}\\
\footnotesize{\textit{change in share of population employed, latest three months vs pre-pandemic, percentage points}}\\
\hspace*{-9mm}  \begin{tikzpicture}
   \begin{axis}[clip=false, xbar stacked, \bbar{x}{0},
            stack negative=separate, ymajorgrids=false,
            yticklabels from table={\occ}{name},
			y axis line style = { opacity = 0 },   
		    x axis line style = { opacity = 0 }, 
			yticklabel style={black, align=left, anchor=east},            
            ytick=data, tickwidth = 0pt,
            legend style={at={(-0.3, 1.05)},
                anchor=north,legend columns=-1, draw=none},
            legend cell align=left,
            extra x ticks={0}, extra x tick style={
                grid style={black}, xticklabel=\empty,},
            enlarge y limits={0.02}, enlarge x limits={0.03},
            width=8.8cm, height=12cm,
            \barylab{5.0cm}{1.5ex}]
            \addplot[fill=green!85!blue, draw opacity=0] table 
            	[y expr=\coordindex, x=PT, col sep=semicolon] {\occ};
            \addplot[fill=blue!60!cyan, draw opacity=0] table 
            	[y expr=\coordindex, x=FT, col sep=semicolon] {\occ};
            \legend{Part-time, Full-time}
        \end{axis}
    \end{tikzpicture}\\
\footnotesize{Source: Author's Calculation from CPS} \hfill \tbllink{occs.csv}
\end{minipage}
\newpage
\begin{minipage}{0.76\textwidth} 
\subsection*{Unemployment}
\hypertarget{labu}{}
\vspace{-0.5mm}
\small The headline unemployment rate, also known as the U3 unemployment rate, measures people who do not have a job but are looking for one or are on temporary layoff, as a share of the labor force (those employed and unemployed). & Dec  `23 & Nov  `23 & Oct  `23 & Sep  `23 & Dec  `22 & Dec  `21 & GFC  peak & Date  of  peak \\  Under-utilization  Rate  (U6) &7.1&7.0&7.2&7.0&6.5&7.3&17.2& Dec  `09 \\  Unemployment  Rate  (U3) &3.7&3.7&3.8&3.8&3.5&3.9&10.0& Oct  `09 \\  \textit{by  race/ethnicity:} &&&&&&&&\\  \hspace{2mm}  White &3.5&3.3&3.5&3.4&3.0&3.3&9.2& Oct  `09 \\  \hspace{2mm}  Black &5.2&5.8&5.8&5.7&5.7&6.9&16.8& Mar  `10 \\  \hspace{2mm}  Hispanic &5.0&4.6&4.8&4.6&4.2&4.6&13.0& Aug  `09 \\  \hspace{2mm}  Asian &3.1&3.5&3.1&2.9&2.4&3.8&8.4& Dec  `09 \\ 
\end{minipage}

\begin{minipage}{0.29\textwidth}
\small BLS also \href{https://www.bls.gov/news.release/empsit.t15.htm}{report} a broader measure of unemployment, known as U6 or labor under-utilization. Labor under-utilization includes unemployed people counted in U3 as well as people who have given up looking for work and people who are part-time but would like to work full-time. \input{text/unemp2.txt}

\end{minipage}\hspace{6mm}
\begin{minipage}{0.425\textwidth}
\normalsize \textbf{Unemployment Measures}\\
\footnotesize{\textit{share of labor force, percent}}
\vspace{2.2cm}

\hspace{4mm} \begin{tikzpicture}[overlay]
	\begin{axis}[\bbar{y}{0}, \shdateaxisticks ytick={0, 5, 10, 15, 20}, 
		yticklabel style={text width=1.0em}, 
		enlarge y limits={0.1}, ymin=1.6, ymax=18.4, xmin={1989-12-01},
		xticklabel={`\short{\year}}, height=4.2cm, width=7.1cm, clip=false]
	\rbars
	\stdline{blue}{date}{U6}{data/unemp2.csv}
	\stdline{blue!50!cyan}{date}{Total}{data/unemp2.csv}
	\absnode{{2010-01-01}}{7.6}{\footnotesize \color{blue!50!cyan}U3}
	\absnode{{1995-01-01}}{14}{\footnotesize \color{blue}U6}
	\input{text/u3_node.txt}
	\input{text/u6_node.txt}
	\end{axis}
\end{tikzpicture}

\footnotesize{Source: BLS} \hfill \tbllink{unemp2.csv}
\end{minipage}

\begin{minipage}{0.76\textwidth}
\small \input{text/unemp3.txt}
\vspace{1mm}

\normalsize \textbf{Unemployment Rate}\\
\footnotesize{\textit{unemployed share of labor force, percent}}
\vspace{2.8cm}

\hspace*{4mm} \begin{tikzpicture}[overlay]
	\begin{axis}[\bbar{y}{0}, \dateaxisticks ytick={0, 5, 10, 15, 20}, 
		enlarge y limits={0.1}, yticklabel style={text width=1.0em},
		xticklabel={`\short{\year}}, clip=false, height=4.8cm, ymin=1.4, ymax=17,
		legend style={fill=white, at={(0.52, 1.01)}}]
	\rbars
	\stdline{green!50!teal!60!black}{date}{Black}{data/unemp.csv}
	\stdline{cyan!60!white}{date}{Hispanic}{data/unemp.csv}
	\stdline{blue!80!black}{date}{White}{data/unemp.csv}
	\legend{Black, Hispanic, White};
	\end{axis}
\end{tikzpicture}

\footnotesize{Source: Bureau of Labor Statistics} \hfill \tbllink{unemp.csv}
\vspace{1mm}

\normalsize \textbf{Unemployment Measures}\\
\footnotesize{\textit{seasonally adjusted, percent}}\\
\rowcolors{1}{}{black!5} \setlength{\tabcolsep}{3.1pt} \color{black!90}
	{\renewcommand{\arraystretch}{1.55} \hspace*{-1mm}
	\begin{tabular}{p{32mm} R{7.7mm} R{7.7mm} R{7.7mm} R{7.7mm}
		   R{7.7mm} R{7.7mm} R{7.7mm} R{12.8mm} R{11.2mm}}
			 \input data/unemp1.tex \hline
		\end{tabular}}
\vspace{-2mm}

\footnotesize{Source: Bureau of Labor Statistics} \hspace{7.35cm} \tbllink{unemp.csv}
\end{minipage}
\newpage
\begin{minipage}{0.76\textwidth}
\small Current Population Survey microdata are used to calculate unemployment by race and ethnicity over the latest six months, on average, and for the same six months before the COVID-19 pandemic. The groups used to produce these estimates separate those with a non-Hispanic ethnicity by race: white alone, black alone, Asian or Hawaiian/Pacific Islander alone, Native American/American Indian or Alaskan Native alone, or more than one race, from those with a Hispanic ethnicity and any race.  
\end{minipage}
\vspace{1.5mm}

\begin{minipage}{0.28\textwidth} 
\normalsize \textbf{Unemployment Rate}\\
\footnotesize{\textit{unemployed share of labor force}}\\
\footnotesize{\textit{age 18 to 64}}\\
\footnotesize{\textit{not seasonally adjusted}}
\end{minipage}
\begin{minipage}{0.17\textwidth} 
\footnotesize \ 6 mo. ending:\\
\small \cbox{orange!75!yellow} \input{text/cps_lfs_datesh.txt}\\
\small \cbox{blue!40!magenta} \input{text/cps_lfs_dateshprepan.txt}
\end{minipage}\hspace{1mm}
\begin{minipage}{0.28\textwidth} 
\normalsize \textbf{Share of Unemployed}\\
\footnotesize{\textit{group share of unemployed}}\\
\footnotesize{\textit{age 18 to 64}}\\
\footnotesize{\textit{not seasonally adjusted}}
\end{minipage}
\vspace{-1.5mm}

\begin{minipage}{0.37\textwidth} 
\hspace*{-3mm} \begin{tikzpicture}
  	\begin{axis}[\barplotnogrid axis y line=left, \barylab{2.1cm}{1.5ex}
    	width=4.3cm, bar width=2.4ex, height=6.9cm,
    	enlarge y limits={abs=0.42cm}, enlarge x limits=0.03, \bbar{x}{0},
		yticklabels from table={\ungrp}{name}, xmin=0.2, 
		yticklabel style={font=\footnotesize}, nodes near coords, 
		nodes near coords style={/pgf/number format/.cd, fixed zerofill,
			precision=1, assume math mode}]
  	\addplot[fill=blue!40!magenta, draw=none] 
  		table [y expr=-\coordindex, x index=2] {\ungrp};
  	\addplot[fill=orange!75!yellow, draw=none] 
  		table [y expr=-\coordindex, x index=1] {\ungrp};
  	\end{axis}
\end{tikzpicture}
\end{minipage} \hspace{2mm}
\begin{minipage}{0.37\textwidth} 
\begin{tikzpicture}
  	\begin{axis}[\barplotnogrid axis y line=left, \barylab{2.1cm}{1.5ex}
    	width=4.3cm, bar width=1.7ex, height=6.9cm,
    	enlarge y limits={abs=0.42cm}, enlarge x limits=0.03, \bbar{x}{0},
		yticklabels from table={\ungrpsh}{name}, xmin=0.2, 
		yticklabel style={font=\footnotesize}, clip=false, nodes near coords, 
		nodes near coords style={/pgf/number format/.cd, fixed zerofill,
			precision=1, assume math mode}]
  	\addplot[fill=green!80!blue, draw=none] 
  		table [y expr=-\coordindex, x index=3] {\ungrpsh};
  	\addplot[fill=blue!40!magenta, draw=none] 
  		table [y expr=-\coordindex, x index=2] {\ungrpsh};
  	\addplot[fill=orange!75!yellow, draw=none] 
  		table [y expr=-\coordindex, x index=1] {\ungrpsh};
  	\stdnode{2.2cm}{0.8cm}{{\color{green!75!blue} \faSquare}:  \footnotesize Share of\\ \footnotesize Labor Force}
  	\end{axis}
\end{tikzpicture}
\end{minipage}
\vspace{-1.5mm}

\begin{minipage}{0.76\textwidth} 
\footnotesize{Source: Author's calculations from CPS} \hfill \tbllink{unempgroups3.csv} \ \tbllink{unempgroups.csv} \ \tbllink{unempgroups2.csv}
\vspace{3mm}

\normalsize \textbf{Unemployment Rate}
\vspace{0.5mm}

\centerline{\small Total, any race or ethnicity}
\vspace{-2mm}

Total \hspace{34mm} Men \hspace{32mm} Women\\
\hspace*{-4mm} \begin{tikzpicture}
  	\begin{axis}[\barplotnogrid axis y line=left, \barylab{1.0cm}{1.5ex}
    	width=4.0cm, bar width=1.2ex, height=5.6cm,
    	enlarge y limits={abs=0.3cm}, enlarge x limits=0.03, \bbar{x}{0},
		yticklabels from table={\unemptt}{name}, xmin=0.1, xmax=19,
		yticklabel style={font=\footnotesize}, nodes near coords, 
		nodes near coords style={/pgf/number format/.cd, fixed zerofill,
			precision=1, assume math mode}]
  	\addplot[fill=blue!40!magenta, draw=none] 
  		table [y expr=-\coordindex, x index=5] {\unemptt};
  	\addplot[fill=orange!75!yellow, draw=none] 
  		table [y expr=-\coordindex, x index=6] {\unemptt};
  	\end{axis}
\end{tikzpicture} \hfill
\begin{tikzpicture}
  	\begin{axis}[\barplotnogrid axis y line=left, \barylab{1.0cm}{1.5ex}
    	width=4.0cm, bar width=1.2ex, height=5.6cm,
    	enlarge y limits={abs=0.3cm}, enlarge x limits=0.03, \bbar{x}{0},
		yticklabels from table={\unemptt}{name}, xmin=0.1, xmax=19,
		yticklabel style={font=\footnotesize}, nodes near coords, 
		nodes near coords style={/pgf/number format/.cd, fixed zerofill,
			precision=1, assume math mode}]
  	\addplot[fill=blue!40!magenta, draw=none] 
  		table [y expr=-\coordindex, x index=1] {\unemptt};
  	\addplot[fill=orange!75!yellow, draw=none] 
  		table [y expr=-\coordindex, x index=2] {\unemptt};
  	\end{axis}
\end{tikzpicture}\hfill
\begin{tikzpicture}
  	\begin{axis}[\barplotnogrid axis y line=left, \barylab{1.0cm}{1.5ex}
    	width=4.0cm, bar width=1.2ex, height=5.6cm,
    	enlarge y limits={abs=0.3cm}, enlarge x limits=0.03, \bbar{x}{0},
		yticklabels from table={\unemptt}{name}, xmin=0.1, xmax=19,
		yticklabel style={font=\footnotesize}, nodes near coords, 
		nodes near coords style={/pgf/number format/.cd, fixed zerofill,
			precision=1, assume math mode}]
  	\addplot[fill=blue!40!magenta, draw=none] 
  		table [y expr=-\coordindex, x index=3] {\unemptt};
  	\addplot[fill=orange!75!yellow, draw=none] 
  		table [y expr=-\coordindex, x index=4] {\unemptt};
  	\end{axis}
\end{tikzpicture}
\vspace{3mm}

\centerline{\small White, non-Hispanic, alone}
\vspace{-2mm}

Total \hspace{34mm} Men \hspace{32mm} Women\\
\hspace*{-4mm} \begin{tikzpicture}
  	\begin{axis}[\barplotnogrid axis y line=left, \barylab{1.0cm}{1.5ex}
    	width=4.0cm, bar width=1.2ex, height=5.6cm,
    	enlarge y limits={abs=0.3cm}, enlarge x limits=0.03, \bbar{x}{0},
		yticklabels from table={\unemp}{name}, xmin=0.1, xmax=19,
		yticklabel style={font=\footnotesize}, nodes near coords, 
		nodes near coords style={/pgf/number format/.cd, fixed zerofill,
			precision=1, assume math mode}]
  	\addplot[fill=blue!40!magenta, draw=none] 
  		table [y expr=-\coordindex, x index=5] {\unemp};
  	\addplot[fill=orange!75!yellow, draw=none] 
  		table [y expr=-\coordindex, x index=6] {\unemp};
  	\end{axis}
\end{tikzpicture} \hfill
\begin{tikzpicture}
  	\begin{axis}[\barplotnogrid axis y line=left, \barylab{1.0cm}{1.5ex}
    	width=4.0cm, bar width=1.2ex, height=5.6cm,
    	enlarge y limits={abs=0.3cm}, enlarge x limits=0.03, \bbar{x}{0},
		yticklabels from table={\unemp}{name}, xmin=0.1, xmax=19,
		yticklabel style={font=\footnotesize}, nodes near coords, 
		nodes near coords style={/pgf/number format/.cd, fixed zerofill,
			precision=1, assume math mode}]
  	\addplot[fill=blue!40!magenta, draw=none] 
  		table [y expr=-\coordindex, x index=1] {\unemp};
  	\addplot[fill=orange!75!yellow, draw=none] 
  		table [y expr=-\coordindex, x index=2] {\unemp};
  	\end{axis}
\end{tikzpicture}\hfill
\begin{tikzpicture}
  	\begin{axis}[\barplotnogrid axis y line=left, \barylab{1.0cm}{1.5ex}
    	width=4.0cm, bar width=1.2ex, height=5.6cm,
    	enlarge y limits={abs=0.3cm}, enlarge x limits=0.03, \bbar{x}{0},
		yticklabels from table={\unemp}{name}, xmin=0.1, xmax=19,
		yticklabel style={font=\footnotesize}, nodes near coords, 
		nodes near coords style={/pgf/number format/.cd, fixed zerofill,
			precision=1, assume math mode}]
  	\addplot[fill=blue!40!magenta, draw=none] 
  		table [y expr=-\coordindex, x index=3] {\unemp};
  	\addplot[fill=orange!75!yellow, draw=none] 
  		table [y expr=-\coordindex, x index=4] {\unemp};
  	\end{axis}
\end{tikzpicture}
\end{minipage}
\newpage
\begin{minipage}{0.76\textwidth} 
\normalsize \textbf{Unemployment Rate, Continued}
\vspace{1mm}

\centerline{\small Black, non-Hispanic, alone}
\vspace{-2mm}

Total \hspace{34mm} Men \hspace{32mm} Women\\
\hspace*{-4mm} \begin{tikzpicture}
  	\begin{axis}[\barplotnogrid axis y line=left, \barylab{1.0cm}{1.5ex}
    	width=4.0cm, bar width=1.2ex, height=5.6cm,
    	enlarge y limits={abs=0.3cm}, enlarge x limits=0.03, \bbar{x}{0},
		yticklabels from table={\unemp}{name}, xmin=0.1, xmax=19,
		yticklabel style={font=\footnotesize}, nodes near coords, 
		nodes near coords style={/pgf/number format/.cd, fixed zerofill,
			precision=1, assume math mode}]
  	\addplot[fill=blue!40!magenta, draw=none] 
  		table [y expr=-\coordindex, x index=11] {\unemp};
  	\addplot[fill=orange!75!yellow, draw=none] 
  		table [y expr=-\coordindex, x index=12] {\unemp};
  	\end{axis}
\end{tikzpicture} \hfill
\begin{tikzpicture}
  	\begin{axis}[\barplotnogrid axis y line=left, \barylab{1.0cm}{1.5ex}
    	width=4.0cm, bar width=1.2ex, height=5.6cm,
    	enlarge y limits={abs=0.3cm}, enlarge x limits=0.03, \bbar{x}{0},
		yticklabels from table={\unemp}{name}, xmin=0.1, xmax=19,
		yticklabel style={font=\footnotesize}, nodes near coords, 
		nodes near coords style={/pgf/number format/.cd, fixed zerofill,
			precision=1, assume math mode}]
  	\addplot[fill=blue!40!magenta, draw=none] 
  		table [y expr=-\coordindex, x index=7] {\unemp};
  	\addplot[fill=orange!75!yellow, draw=none] 
  		table [y expr=-\coordindex, x index=8] {\unemp};
  	\end{axis}
\end{tikzpicture}\hfill
\begin{tikzpicture}
  	\begin{axis}[\barplotnogrid axis y line=left, \barylab{1.0cm}{1.5ex}
    	width=4.0cm, bar width=1.2ex, height=5.6cm,
    	enlarge y limits={abs=0.3cm}, enlarge x limits=0.03, \bbar{x}{0},
		yticklabels from table={\unemp}{name}, xmin=0.1, xmax=19,
		yticklabel style={font=\footnotesize}, nodes near coords, 
		nodes near coords style={/pgf/number format/.cd, fixed zerofill,
			precision=1, assume math mode}]
  	\addplot[fill=blue!40!magenta, draw=none] 
  		table [y expr=-\coordindex, x index=9] {\unemp};
  	\addplot[fill=orange!75!yellow, draw=none] 
  		table [y expr=-\coordindex, x index=10] {\unemp};
  	\end{axis}
\end{tikzpicture}
\vspace{3mm}

\centerline{\small Hispanic, any race}
\vspace{-2mm}

Total \hspace{34mm} Men \hspace{32mm} Women\\
\hspace*{-4mm} \begin{tikzpicture}
  	\begin{axis}[\barplotnogrid axis y line=left, \barylab{1.0cm}{1.5ex}
    	width=4.0cm, bar width=1.2ex, height=5.6cm,
    	enlarge y limits={abs=0.3cm}, enlarge x limits=0.03, \bbar{x}{0},
		yticklabels from table={\unemp}{name}, xmin=0.1, xmax=19,
		yticklabel style={font=\footnotesize}, nodes near coords, 
		nodes near coords style={/pgf/number format/.cd, fixed zerofill,
			precision=1, assume math mode}]
  	\addplot[fill=blue!40!magenta, draw=none] 
  		table [y expr=-\coordindex, x index=17] {\unemp};
  	\addplot[fill=orange!75!yellow, draw=none] 
  		table [y expr=-\coordindex, x index=18] {\unemp};
  	\end{axis}
\end{tikzpicture} \hfill
\begin{tikzpicture}
  	\begin{axis}[\barplotnogrid axis y line=left, \barylab{1.0cm}{1.5ex}
    	width=4.0cm, bar width=1.2ex, height=5.6cm,
    	enlarge y limits={abs=0.3cm}, enlarge x limits=0.03, \bbar{x}{0},
		yticklabels from table={\unemp}{name}, xmin=0.1, xmax=19,
		yticklabel style={font=\footnotesize}, nodes near coords, 
		nodes near coords style={/pgf/number format/.cd, fixed zerofill,
			precision=1, assume math mode}]
  	\addplot[fill=blue!40!magenta, draw=none] 
  		table [y expr=-\coordindex, x index=13] {\unemp};
  	\addplot[fill=orange!75!yellow, draw=none] 
  		table [y expr=-\coordindex, x index=14] {\unemp};
  	\end{axis}
\end{tikzpicture}\hfill
\begin{tikzpicture}
  	\begin{axis}[\barplotnogrid axis y line=left, \barylab{1.0cm}{1.5ex}
    	width=4.0cm, bar width=1.2ex, height=5.6cm,
    	enlarge y limits={abs=0.3cm}, enlarge x limits=0.03, \bbar{x}{0},
		yticklabels from table={\unemp}{name}, xmin=0.1, xmax=19,
		yticklabel style={font=\footnotesize}, nodes near coords, 
		nodes near coords style={/pgf/number format/.cd, fixed zerofill,
			precision=1, assume math mode}]
  	\addplot[fill=blue!40!magenta, draw=none] 
  		table [y expr=-\coordindex, x index=15] {\unemp};
  	\addplot[fill=orange!75!yellow, draw=none] 
  		table [y expr=-\coordindex, x index=16] {\unemp};
  	\end{axis}
\end{tikzpicture}
\vspace{3mm}

\centerline{\small Asian or Hawaiian/Pacific Islander, non-Hispanic, alone}
\vspace{-2mm}

Total \hspace{34mm} Men \hspace{32mm} Women\\
\hspace*{-4mm} \begin{tikzpicture}
  	\begin{axis}[\barplotnogrid axis y line=left, \barylab{1.0cm}{1.5ex}
    	width=4.0cm, bar width=1.2ex, height=5.6cm,
    	enlarge y limits={abs=0.3cm}, enlarge x limits=0.03, \bbar{x}{0},
		yticklabels from table={\unemp}{name}, xmin=0.1, xmax=19,
		yticklabel style={font=\footnotesize}, nodes near coords, 
		nodes near coords style={/pgf/number format/.cd, fixed zerofill,
			precision=1, assume math mode}]
  	\addplot[fill=blue!40!magenta, draw=none] 
  		table [y expr=-\coordindex, x index=23] {\unemp};
  	\addplot[fill=orange!75!yellow, draw=none] 
  		table [y expr=-\coordindex, x index=24] {\unemp};
  	\end{axis}
\end{tikzpicture} \hfill
\begin{tikzpicture}
  	\begin{axis}[\barplotnogrid axis y line=left, \barylab{1.0cm}{1.5ex}
    	width=4.0cm, bar width=1.2ex, height=5.6cm,
    	enlarge y limits={abs=0.3cm}, enlarge x limits=0.03, \bbar{x}{0},
		yticklabels from table={\unemp}{name}, xmin=0.1, xmax=19,
		yticklabel style={font=\footnotesize}, nodes near coords, 
		nodes near coords style={/pgf/number format/.cd, fixed zerofill,
			precision=1, assume math mode}]
  	\addplot[fill=blue!40!magenta, draw=none] 
  		table [y expr=-\coordindex, x index=19] {\unemp};
  	\addplot[fill=orange!75!yellow, draw=none] 
  		table [y expr=-\coordindex, x index=20] {\unemp};
  	\end{axis}
\end{tikzpicture}\hfill
\begin{tikzpicture}
  	\begin{axis}[\barplotnogrid axis y line=left, \barylab{1.0cm}{1.5ex}
    	width=4.0cm, bar width=1.2ex, height=5.6cm,
    	enlarge y limits={abs=0.3cm}, enlarge x limits=0.03, \bbar{x}{0},
		yticklabels from table={\unemp}{name}, xmin=0.1, xmax=19,
		yticklabel style={font=\footnotesize}, nodes near coords, 
		nodes near coords style={/pgf/number format/.cd, fixed zerofill,
			precision=1, assume math mode}]
  	\addplot[fill=blue!40!magenta, draw=none] 
  		table [y expr=-\coordindex, x index=21] {\unemp};
  	\addplot[fill=orange!75!yellow, draw=none] 
  		table [y expr=-\coordindex, x index=22] {\unemp};
  	\end{axis}
\end{tikzpicture}
\vspace{3mm}

\centerline{\small Native American/American Indian or Alaskan Native, non-Hispanic, alone}
\vspace{-2mm}

Total \hspace{34mm} Men \hspace{32mm} Women\\
\hspace*{-4mm} \begin{tikzpicture}
  	\begin{axis}[\barplotnogrid axis y line=left, \barylab{1.0cm}{1.5ex}
    	width=4.0cm, bar width=1.2ex, height=2.7cm,
    	enlarge y limits={abs=0.3cm}, enlarge x limits=0.03, \bbar{x}{0},
		yticklabels from table={\unempt}{name}, xmin=0.1, xmax=19,
		yticklabel style={font=\footnotesize}, nodes near coords, 
		nodes near coords style={/pgf/number format/.cd, fixed zerofill,
			precision=1, assume math mode}]
  	\addplot[fill=blue!40!magenta, draw=none] 
  		table [y expr=-\coordindex, x index=5] {\unempt};
  	\addplot[fill=orange!75!yellow, draw=none] 
  		table [y expr=-\coordindex, x index=6] {\unempt};
  	\end{axis}
\end{tikzpicture} \hfill
\begin{tikzpicture}
  	\begin{axis}[\barplotnogrid axis y line=left, \barylab{1.0cm}{1.5ex}
    	width=4.0cm, bar width=1.2ex, height=2.7cm,
    	enlarge y limits={abs=0.3cm}, enlarge x limits=0.03, \bbar{x}{0},
		yticklabels from table={\unempt}{name}, xmin=0.1, xmax=19,
		yticklabel style={font=\footnotesize}, nodes near coords, 
		nodes near coords style={/pgf/number format/.cd, fixed zerofill,
			precision=1, assume math mode}]
  	\addplot[fill=blue!40!magenta, draw=none] 
  		table [y expr=-\coordindex, x index=1] {\unempt};
  	\addplot[fill=orange!75!yellow, draw=none] 
  		table [y expr=-\coordindex, x index=2] {\unempt};
  	\end{axis}
\end{tikzpicture}\hfill
\begin{tikzpicture}
  	\begin{axis}[\barplotnogrid axis y line=left, \barylab{1.0cm}{1.5ex}
    	width=4.0cm, bar width=1.2ex, height=2.7cm,
    	enlarge y limits={abs=0.3cm}, enlarge x limits=0.03, \bbar{x}{0},
		yticklabels from table={\unemp}{name}, xmin=0.1, xmax=19,
		yticklabel style={font=\footnotesize}, nodes near coords, 
		nodes near coords style={/pgf/number format/.cd, fixed zerofill,
			precision=1, assume math mode}]
  	\addplot[fill=blue!40!magenta, draw=none] 
  		table [y expr=-\coordindex, x index=3] {\unempt};
  	\addplot[fill=orange!75!yellow, draw=none] 
  		table [y expr=-\coordindex, x index=4] {\unempt};
  	\end{axis}
\end{tikzpicture}
\vspace{3mm}

\centerline{\small More than one race, non-Hispanic}
\vspace{-2mm}

Total \hspace{34mm} Men \hspace{32mm} Women\\
\hspace*{-4mm} \begin{tikzpicture}
  	\begin{axis}[\barplotnogrid axis y line=left, \barylab{1.0cm}{1.5ex}
    	width=4.0cm, bar width=1.2ex, height=2.7cm,
    	enlarge y limits={abs=0.3cm}, enlarge x limits=0.03, \bbar{x}{0},
		yticklabels from table={\unempt}{name}, xmin=0.1, xmax=19,
		yticklabel style={font=\footnotesize}, nodes near coords, 
		nodes near coords style={/pgf/number format/.cd, fixed zerofill,
			precision=1, assume math mode}]
  	\addplot[fill=blue!40!magenta, draw=none] 
  		table [y expr=-\coordindex, x index=11] {\unempt};
  	\addplot[fill=orange!75!yellow, draw=none] 
  		table [y expr=-\coordindex, x index=12] {\unempt};
  	\end{axis}
\end{tikzpicture} \hfill
\begin{tikzpicture}
  	\begin{axis}[\barplotnogrid axis y line=left, \barylab{1.0cm}{1.5ex}
    	width=4.0cm, bar width=1.2ex, height=2.7cm,
    	enlarge y limits={abs=0.3cm}, enlarge x limits=0.03, \bbar{x}{0},
		yticklabels from table={\unempt}{name}, xmin=0.1, xmax=19,
		yticklabel style={font=\footnotesize}, nodes near coords, 
		nodes near coords style={/pgf/number format/.cd, fixed zerofill,
			precision=1, assume math mode}]
  	\addplot[fill=blue!40!magenta, draw=none] 
  		table [y expr=-\coordindex, x index=7] {\unempt};
  	\addplot[fill=orange!75!yellow, draw=none] 
  		table [y expr=-\coordindex, x index=8] {\unempt};
  	\end{axis}
\end{tikzpicture}\hfill
\begin{tikzpicture}
  	\begin{axis}[\barplotnogrid axis y line=left, \barylab{1.0cm}{1.5ex}
    	width=4.0cm, bar width=1.2ex, height=2.7cm,
    	enlarge y limits={abs=0.3cm}, enlarge x limits=0.03, \bbar{x}{0},
		yticklabels from table={\unempt}{name}, xmin=0.1, xmax=19,
		yticklabel style={font=\footnotesize}, nodes near coords, 
		nodes near coords style={/pgf/number format/.cd, fixed zerofill,
			precision=1, assume math mode}]
  	\addplot[fill=blue!40!magenta, draw=none] 
  		table [y expr=-\coordindex, x index=9] {\unempt};
  	\addplot[fill=orange!75!yellow, draw=none] 
  		table [y expr=-\coordindex, x index=10] {\unempt};
  	\end{axis}
\end{tikzpicture}
\vspace{-1mm}

\footnotesize{Source: Author's calculations from CPS} \hfill \tbllink{unemp_grp.csv}
\end{minipage}
\newpage
\begin{minipage}{0.76\textwidth} 
\subsubsection*{Reasons for Unemployment}
\vspace{-0.5mm}
\small \input{text/unemp_rsn.txt}

People's reasons for unemployment reflect economic conditions. In an economic downturn, job losses increase the unemployment rate and workers who lose jobs are a large share of the unemployed. A downturn also makes it harder for young people to find their first job. In contrast, an economic boom reduces job losses and improves job-finding. 

Some categories of unemployment make up a bigger share of the total during an economic boom. An economic boom can entice people to re-enter the job market, and encourage workers to quit and look for a new job. 
\vspace{1mm}
%\input{text/unemp_reason.txt}

\normalsize \textbf{Unemployment by Reason}\\
\footnotesize{\textit{share of labor force, percent, annual average \hspace{36mm} recent data}}
\vspace{3.6cm}

\hspace{3mm} \begin{tikzpicture}[overlay]
	\begin{groupplot}[group style={group size=2 by 1, horizontal sep=30pt}]
	\nextgroupplot[\bbar{y}{0}, \shdateaxisticks ytick={0, 2, 4, 6, 8, 10, 12, 14, 16}, 
		width=9.0cm, height=5.6cm, enlarge x limits=0.01, 
		ymin=0.3, ymax=9.6, legend cell align={left}, xmin={1989-01-01}, xmax={2023-01-01},
		xticklabel={\short{\year}}, clip=false, yticklabel style={text width=1.2em}, 
		legend style={fill=white, legend columns=2, at={(0.56, 1.01)}}]
	\rbars
	\input{text/unemp_rsn_ltbar.txt}
	\ctsbar{red!75!orange!70!white}{date}{Job Loser}{data/unemp_reason.csv}{5.4pt}
	\ctsbar{purple!80!red!85!black}{date}{New entrant}{data/unemp_reason.csv}{5.4pt}
	\ctsbar{blue!70!teal!70!gray!62!white}{date}{Re-entrant}{data/unemp_reason.csv}{5.4pt}
	\ctsbar{blue!72!black}{date}{Job Leaver}{data/unemp_reason.csv}{5.4pt}
	\legend{,Job Loser, New Entrant, Re-entrant, Job Leaver};
	\nextgroupplot[\bbar{y}{0}, width=4.0cm, height=5.6cm, yticklabel style={text width=0.8em}, 
		axis line style={draw=none}, xtick=data, clip=false,
			xticklabel style={black!80, text width=1.7em, align=center}, 
	 	enlarge y limits={0.06}, enlarge x limits={0.12}, ymin=0.1,
	 	xticklabels from table={\unrsn}{date},
	 	nodes near coords style={/pgf/number format/.cd, fixed zerofill,
			precision=1, assume math mode}]
  	\addplot[ybar, bar shift=0pt, fill=none, draw=none, nodes near coords] 
  		table [x expr=\coordindex, y index=7] {\unrsn};
	\addplot[ybar stacked, bar width=9.6pt, draw opacity=0, fill=red!75!orange!70!white] table [x expr=\coordindex, y=Job Loser]{\unrsn};
	\addplot[ybar stacked, bar width=9.6pt, draw opacity=0, fill=purple!80!red!85!black] table [x expr=\coordindex, y=New entrant]{\unrsn};
	\addplot[ybar stacked, bar width=9.6pt, draw opacity=0, fill=blue!70!teal!70!gray!62!white] table [x expr=\coordindex, y=Re-entrant]{\unrsn};
	\addplot[ybar stacked, bar width=9.6pt, draw opacity=0, fill=blue!72!black] table [x expr=\coordindex, y=Job Leaver]{\unrsn};
	\input{text/unemp_rsn_ltval.txt}
	\draw [dashed] (axis cs:{0.5},\pgfkeysvalueof{/pgfplots/ymin}) -- 
		(axis cs:{0.5}, \pgfkeysvalueof{/pgfplots/ymax});
	\end{groupplot}
\end{tikzpicture}

\footnotesize{Source: Bureau of Labor Statistics} \hfill \tbllink{unemp_reason.csv}
\vspace{3mm}

\small Many job losses are temporary, particularly during the COVID-19 recession. Other job separations are permanent. \input{text/unemp_rsn2.txt}
\vspace{1mm}

\normalsize \textbf{Unemployment by Reason}\\
\footnotesize{\textit{share of labor force, percent}}\\
\rowcolors{1}{}{black!5} \setlength{\tabcolsep}{1.2pt} \color{black!90}
	{\renewcommand{\arraystretch}{1.5}
		\begin{tabular}{p{36mm} R{9mm} R{9mm} R{9mm} R{9mm} R{9mm} 
		   R{9mm} R{9mm} R{9mm} R{9mm}}
			 \input data/unempreason_table.tex \hline
		\end{tabular}}\vspace{-1mm}
		
\footnotesize{Source: Bureau of Labor Statistics; Author}\\
\footnotesize{* During the COVID-19 shutdowns some unemployed were incorrectly counted as employed but not at work.}
\end{minipage}
\newpage
\begin{minipage}{0.76\textwidth} 
\subsubsection*{Duration of Unemployment}
\small US unemployment benefits are available for a relatively short duration, compared with other advanced countries. Therefore, the long-term unemployed risk running out of unemployment benefits, causing a sharp reduction in income. Additionally, long periods of unemployment can make re-entering the labor market more challenging. 

\input{text/ltu.txt}
\vspace{1mm}

\normalsize \textbf{Long-term Unemployed}\\
\footnotesize{\textit{unemployed 27+ weeks, share of population \hspace{11mm} unemployed 15+ weeks, share of population}}\\
\hspace*{-2mm} \begin{tikzpicture}
	\begin{axis}[\bbar{y}{0}, \shdateaxisticks ytick={0, 1, 2, 3, 4, 5}, 
		enlarge y limits={0.01}, ymin=0, ymax=4.2,
		xticklabel={`\short{\year}}, height=5.0cm, width=6.0cm, clip=false]
	\rbars
	\stdline{blue}{date}{LT}{data/ltu.csv}
	\input{text/ltu_node.txt}
	\end{axis}
\end{tikzpicture} \hfill \begin{tikzpicture}
	\begin{axis}[\bbar{y}{0}, \shdateaxisticks ytick={0, 1, 2, 3, 4, 5}, 
		enlarge y limits={0.01}, ymin=0, ymax=4.2,
		xticklabel={`\short{\year}}, height=5.0cm, width=6.0cm, clip=false]
	\rbars
	\stdline{violet!90!black}{date}{MT}{data/ltu.csv}
	\input{text/ltu_node2.txt}
	\end{axis}
\end{tikzpicture}\\
\footnotesize{Source: Bureau of Labor Statistics} \hfill \tbllink{ltu.csv}
\vspace{2mm}

\small \input{text/unempdur.txt}
\vspace{1mm}

\normalsize \textbf{Duration of Unemployment}\\
\footnotesize{\textit{in weeks}}\\
\hspace*{-2mm} \begin{tikzpicture}
	\begin{axis}[\bbar{y}{0}, \dateaxisticks ytick={0, 10, 20, 30, 40, 50}, 
		enlarge y limits={0.1}, ymin=3.5, 
		yticklabel style={text width=1.2em},
		xticklabel={`\short{\year}}, height=5.2cm, width=12.0cm, clip=false]
	\rbars
	\stdline{blue!60!cyan}{date}{Mean}{data/unempdur.csv}
	\stdline{green!75!blue}{date}{Median}{data/unempdur.csv}
	\absnode{{2015-04-01}}{17}{\color{green!75!blue}Median}
	\absnode{{2014-04-01}}{39}{\color{blue!60!cyan}Average}
	\input{text/unempdur_nodes.txt}
	\end{axis}
\end{tikzpicture}\\
\footnotesize{Source: Bureau of Labor Statistics} \hfill \tbllink{unempdur.csv}
\end{minipage}
\newpage
\begin{minipage}{0.76\textwidth}
\subsubsection*{Unemployment by Metro Area}
\vspace{-1mm}

\small The Bureau of Labor Statistics \href{https://www.bls.gov/lau/}{produce} local area estimates of unemployment, including the \textbf{unemployment rate for metro areas}. The following map shows changes since 2019 in metro area unemployment rates. An increase in the unemployment rate is shown by a blue circle and a decrease is shown by a light green circle; circle size is the magnitude of the change. 

\input{text/msa_unemp_ch_n.txt}
\end{minipage}
\vspace{1mm}

\begin{minipage}{0.82\textwidth}
\normalsize \textbf{Change in Unemployment Rate by Metro Area}\\
\footnotesize{\textit{from \input{text/unemp_map_date.txt}\unskip, percentage points}}
\vspace{-4mm}

\hspace*{-8mm} %% Creator: Matplotlib, PGF backend
%%
%% To include the figure in your LaTeX document, write
%%   \input{<filename>.pgf}
%%
%% Make sure the required packages are loaded in your preamble
%%   \usepackage{pgf}
%%
%% Also ensure that all the required font packages are loaded; for instance,
%% the lmodern package is sometimes necessary when using math font.
%%   \usepackage{lmodern}
%%
%% Figures using additional raster images can only be included by \input if
%% they are in the same directory as the main LaTeX file. For loading figures
%% from other directories you can use the `import` package
%%   \usepackage{import}
%%
%% and then include the figures with
%%   \import{<path to file>}{<filename>.pgf}
%%
%% Matplotlib used the following preamble
%%   
%%   \usepackage{fontspec}
%%   \setmainfont{DejaVuSerif.ttf}[Path=\detokenize{/home/brian/miniconda3/lib/python3.8/site-packages/matplotlib/mpl-data/fonts/ttf/}]
%%   \setsansfont{DejaVuSans.ttf}[Path=\detokenize{/home/brian/miniconda3/lib/python3.8/site-packages/matplotlib/mpl-data/fonts/ttf/}]
%%   \setmonofont{DejaVuSansMono.ttf}[Path=\detokenize{/home/brian/miniconda3/lib/python3.8/site-packages/matplotlib/mpl-data/fonts/ttf/}]
%%   \makeatletter\@ifpackageloaded{underscore}{}{\usepackage[strings]{underscore}}\makeatother
%%
\begingroup%
\makeatletter%
\begin{pgfpicture}%
\pgfpathrectangle{\pgfpointorigin}{\pgfqpoint{5.237500in}{6.020271in}}%
\pgfusepath{use as bounding box, clip}%
\begin{pgfscope}%
\pgfsetbuttcap%
\pgfsetmiterjoin%
\definecolor{currentfill}{rgb}{1.000000,1.000000,1.000000}%
\pgfsetfillcolor{currentfill}%
\pgfsetlinewidth{0.000000pt}%
\definecolor{currentstroke}{rgb}{1.000000,1.000000,1.000000}%
\pgfsetstrokecolor{currentstroke}%
\pgfsetdash{}{0pt}%
\pgfpathmoveto{\pgfqpoint{0.000000in}{0.000000in}}%
\pgfpathlineto{\pgfqpoint{5.237500in}{0.000000in}}%
\pgfpathlineto{\pgfqpoint{5.237500in}{6.020271in}}%
\pgfpathlineto{\pgfqpoint{0.000000in}{6.020271in}}%
\pgfpathlineto{\pgfqpoint{0.000000in}{0.000000in}}%
\pgfpathclose%
\pgfusepath{fill}%
\end{pgfscope}%
\begin{pgfscope}%
\pgfpathrectangle{\pgfqpoint{0.100000in}{2.493063in}}{\pgfqpoint{5.037500in}{3.427208in}}%
\pgfusepath{clip}%
\pgfsetbuttcap%
\pgfsetmiterjoin%
\definecolor{currentfill}{rgb}{1.000000,1.000000,1.000000}%
\pgfsetfillcolor{currentfill}%
\pgfsetlinewidth{0.501875pt}%
\definecolor{currentstroke}{rgb}{0.827451,0.827451,0.827451}%
\pgfsetstrokecolor{currentstroke}%
\pgfsetdash{}{0pt}%
\pgfpathmoveto{\pgfqpoint{1.635246in}{3.344091in}}%
\pgfpathlineto{\pgfqpoint{1.620137in}{3.344339in}}%
\pgfpathlineto{\pgfqpoint{1.609854in}{3.355720in}}%
\pgfpathlineto{\pgfqpoint{1.602861in}{3.385947in}}%
\pgfpathlineto{\pgfqpoint{1.619112in}{3.389734in}}%
\pgfpathlineto{\pgfqpoint{1.635621in}{3.385661in}}%
\pgfpathlineto{\pgfqpoint{1.652422in}{3.373753in}}%
\pgfpathlineto{\pgfqpoint{1.652414in}{3.357285in}}%
\pgfpathlineto{\pgfqpoint{1.635246in}{3.344091in}}%
\pgfpathclose%
\pgfusepath{stroke,fill}%
\end{pgfscope}%
\begin{pgfscope}%
\pgfpathrectangle{\pgfqpoint{0.100000in}{2.493063in}}{\pgfqpoint{5.037500in}{3.427208in}}%
\pgfusepath{clip}%
\pgfsetbuttcap%
\pgfsetmiterjoin%
\definecolor{currentfill}{rgb}{1.000000,1.000000,1.000000}%
\pgfsetfillcolor{currentfill}%
\pgfsetlinewidth{0.501875pt}%
\definecolor{currentstroke}{rgb}{0.827451,0.827451,0.827451}%
\pgfsetstrokecolor{currentstroke}%
\pgfsetdash{}{0pt}%
\pgfpathmoveto{\pgfqpoint{1.727210in}{3.164640in}}%
\pgfpathlineto{\pgfqpoint{1.710308in}{3.171076in}}%
\pgfpathlineto{\pgfqpoint{1.710230in}{3.183701in}}%
\pgfpathlineto{\pgfqpoint{1.689811in}{3.194378in}}%
\pgfpathlineto{\pgfqpoint{1.692298in}{3.221261in}}%
\pgfpathlineto{\pgfqpoint{1.698212in}{3.233947in}}%
\pgfpathlineto{\pgfqpoint{1.710698in}{3.222271in}}%
\pgfpathlineto{\pgfqpoint{1.731004in}{3.224222in}}%
\pgfpathlineto{\pgfqpoint{1.725876in}{3.187065in}}%
\pgfpathlineto{\pgfqpoint{1.727210in}{3.164640in}}%
\pgfpathclose%
\pgfusepath{stroke,fill}%
\end{pgfscope}%
\begin{pgfscope}%
\pgfpathrectangle{\pgfqpoint{0.100000in}{2.493063in}}{\pgfqpoint{5.037500in}{3.427208in}}%
\pgfusepath{clip}%
\pgfsetbuttcap%
\pgfsetmiterjoin%
\definecolor{currentfill}{rgb}{1.000000,1.000000,1.000000}%
\pgfsetfillcolor{currentfill}%
\pgfsetlinewidth{0.501875pt}%
\definecolor{currentstroke}{rgb}{0.827451,0.827451,0.827451}%
\pgfsetstrokecolor{currentstroke}%
\pgfsetdash{}{0pt}%
\pgfpathmoveto{\pgfqpoint{1.798159in}{3.081762in}}%
\pgfpathlineto{\pgfqpoint{1.776708in}{3.085415in}}%
\pgfpathlineto{\pgfqpoint{1.764847in}{3.103779in}}%
\pgfpathlineto{\pgfqpoint{1.763172in}{3.118878in}}%
\pgfpathlineto{\pgfqpoint{1.798159in}{3.081762in}}%
\pgfpathclose%
\pgfusepath{stroke,fill}%
\end{pgfscope}%
\begin{pgfscope}%
\pgfpathrectangle{\pgfqpoint{0.100000in}{2.493063in}}{\pgfqpoint{5.037500in}{3.427208in}}%
\pgfusepath{clip}%
\pgfsetbuttcap%
\pgfsetmiterjoin%
\definecolor{currentfill}{rgb}{1.000000,1.000000,1.000000}%
\pgfsetfillcolor{currentfill}%
\pgfsetlinewidth{0.501875pt}%
\definecolor{currentstroke}{rgb}{0.827451,0.827451,0.827451}%
\pgfsetstrokecolor{currentstroke}%
\pgfsetdash{}{0pt}%
\pgfpathmoveto{\pgfqpoint{1.754202in}{3.084965in}}%
\pgfpathlineto{\pgfqpoint{1.765448in}{3.075835in}}%
\pgfpathlineto{\pgfqpoint{1.767890in}{3.060139in}}%
\pgfpathlineto{\pgfqpoint{1.745251in}{3.061006in}}%
\pgfpathlineto{\pgfqpoint{1.754202in}{3.084965in}}%
\pgfpathclose%
\pgfusepath{stroke,fill}%
\end{pgfscope}%
\begin{pgfscope}%
\pgfpathrectangle{\pgfqpoint{0.100000in}{2.493063in}}{\pgfqpoint{5.037500in}{3.427208in}}%
\pgfusepath{clip}%
\pgfsetbuttcap%
\pgfsetmiterjoin%
\definecolor{currentfill}{rgb}{1.000000,1.000000,1.000000}%
\pgfsetfillcolor{currentfill}%
\pgfsetlinewidth{0.501875pt}%
\definecolor{currentstroke}{rgb}{0.827451,0.827451,0.827451}%
\pgfsetstrokecolor{currentstroke}%
\pgfsetdash{}{0pt}%
\pgfpathmoveto{\pgfqpoint{1.805849in}{2.997287in}}%
\pgfpathlineto{\pgfqpoint{1.783257in}{3.005408in}}%
\pgfpathlineto{\pgfqpoint{1.789731in}{3.025097in}}%
\pgfpathlineto{\pgfqpoint{1.780262in}{3.045159in}}%
\pgfpathlineto{\pgfqpoint{1.783236in}{3.058951in}}%
\pgfpathlineto{\pgfqpoint{1.801416in}{3.064207in}}%
\pgfpathlineto{\pgfqpoint{1.801066in}{3.042112in}}%
\pgfpathlineto{\pgfqpoint{1.822569in}{3.029856in}}%
\pgfpathlineto{\pgfqpoint{1.834868in}{3.000225in}}%
\pgfpathlineto{\pgfqpoint{1.821162in}{2.988814in}}%
\pgfpathlineto{\pgfqpoint{1.805849in}{2.997287in}}%
\pgfpathclose%
\pgfusepath{stroke,fill}%
\end{pgfscope}%
\begin{pgfscope}%
\pgfpathrectangle{\pgfqpoint{0.100000in}{2.493063in}}{\pgfqpoint{5.037500in}{3.427208in}}%
\pgfusepath{clip}%
\pgfsetbuttcap%
\pgfsetmiterjoin%
\definecolor{currentfill}{rgb}{1.000000,1.000000,1.000000}%
\pgfsetfillcolor{currentfill}%
\pgfsetlinewidth{0.501875pt}%
\definecolor{currentstroke}{rgb}{0.827451,0.827451,0.827451}%
\pgfsetstrokecolor{currentstroke}%
\pgfsetdash{}{0pt}%
\pgfpathmoveto{\pgfqpoint{1.724255in}{2.796939in}}%
\pgfpathlineto{\pgfqpoint{1.713491in}{2.819173in}}%
\pgfpathlineto{\pgfqpoint{1.718257in}{2.834128in}}%
\pgfpathlineto{\pgfqpoint{1.739369in}{2.854829in}}%
\pgfpathlineto{\pgfqpoint{1.741367in}{2.870698in}}%
\pgfpathlineto{\pgfqpoint{1.755183in}{2.905830in}}%
\pgfpathlineto{\pgfqpoint{1.791521in}{2.911833in}}%
\pgfpathlineto{\pgfqpoint{1.798476in}{2.932398in}}%
\pgfpathlineto{\pgfqpoint{1.814497in}{2.924687in}}%
\pgfpathlineto{\pgfqpoint{1.847051in}{2.866970in}}%
\pgfpathlineto{\pgfqpoint{1.847462in}{2.848126in}}%
\pgfpathlineto{\pgfqpoint{1.838260in}{2.831817in}}%
\pgfpathlineto{\pgfqpoint{1.849936in}{2.794370in}}%
\pgfpathlineto{\pgfqpoint{1.841778in}{2.788333in}}%
\pgfpathlineto{\pgfqpoint{1.813989in}{2.790088in}}%
\pgfpathlineto{\pgfqpoint{1.780572in}{2.801257in}}%
\pgfpathlineto{\pgfqpoint{1.752487in}{2.805868in}}%
\pgfpathlineto{\pgfqpoint{1.744986in}{2.799316in}}%
\pgfpathlineto{\pgfqpoint{1.724255in}{2.796939in}}%
\pgfpathclose%
\pgfusepath{stroke,fill}%
\end{pgfscope}%
\begin{pgfscope}%
\pgfpathrectangle{\pgfqpoint{0.100000in}{2.493063in}}{\pgfqpoint{5.037500in}{3.427208in}}%
\pgfusepath{clip}%
\pgfsetbuttcap%
\pgfsetmiterjoin%
\definecolor{currentfill}{rgb}{1.000000,1.000000,1.000000}%
\pgfsetfillcolor{currentfill}%
\pgfsetlinewidth{0.501875pt}%
\definecolor{currentstroke}{rgb}{0.827451,0.827451,0.827451}%
\pgfsetstrokecolor{currentstroke}%
\pgfsetdash{}{0pt}%
\pgfpathmoveto{\pgfqpoint{0.951624in}{3.067308in}}%
\pgfpathlineto{\pgfqpoint{0.949465in}{3.065407in}}%
\pgfpathlineto{\pgfqpoint{0.939556in}{3.067940in}}%
\pgfpathlineto{\pgfqpoint{0.940296in}{3.072929in}}%
\pgfpathlineto{\pgfqpoint{0.945905in}{3.075212in}}%
\pgfpathlineto{\pgfqpoint{0.953902in}{3.086708in}}%
\pgfpathlineto{\pgfqpoint{0.955231in}{3.094478in}}%
\pgfpathlineto{\pgfqpoint{0.959726in}{3.096515in}}%
\pgfpathlineto{\pgfqpoint{0.967041in}{3.096522in}}%
\pgfpathlineto{\pgfqpoint{0.974687in}{3.121738in}}%
\pgfpathlineto{\pgfqpoint{0.977683in}{3.123812in}}%
\pgfpathlineto{\pgfqpoint{0.974996in}{3.128741in}}%
\pgfpathlineto{\pgfqpoint{0.968708in}{3.123755in}}%
\pgfpathlineto{\pgfqpoint{0.950231in}{3.130477in}}%
\pgfpathlineto{\pgfqpoint{0.939203in}{3.141802in}}%
\pgfpathlineto{\pgfqpoint{0.944011in}{3.146969in}}%
\pgfpathlineto{\pgfqpoint{0.941824in}{3.154257in}}%
\pgfpathlineto{\pgfqpoint{0.943172in}{3.162867in}}%
\pgfpathlineto{\pgfqpoint{0.940800in}{3.172932in}}%
\pgfpathlineto{\pgfqpoint{0.940421in}{3.183009in}}%
\pgfpathlineto{\pgfqpoint{0.951035in}{3.181323in}}%
\pgfpathlineto{\pgfqpoint{0.953137in}{3.184293in}}%
\pgfpathlineto{\pgfqpoint{0.958416in}{3.184005in}}%
\pgfpathlineto{\pgfqpoint{0.964496in}{3.173258in}}%
\pgfpathlineto{\pgfqpoint{0.970359in}{3.165389in}}%
\pgfpathlineto{\pgfqpoint{0.965663in}{3.160268in}}%
\pgfpathlineto{\pgfqpoint{0.973783in}{3.158303in}}%
\pgfpathlineto{\pgfqpoint{0.975632in}{3.160815in}}%
\pgfpathlineto{\pgfqpoint{0.971290in}{3.171008in}}%
\pgfpathlineto{\pgfqpoint{0.962148in}{3.179455in}}%
\pgfpathlineto{\pgfqpoint{0.963266in}{3.182995in}}%
\pgfpathlineto{\pgfqpoint{0.955970in}{3.191941in}}%
\pgfpathlineto{\pgfqpoint{0.963421in}{3.193757in}}%
\pgfpathlineto{\pgfqpoint{0.956965in}{3.212911in}}%
\pgfpathlineto{\pgfqpoint{0.962214in}{3.215566in}}%
\pgfpathlineto{\pgfqpoint{0.963494in}{3.220158in}}%
\pgfpathlineto{\pgfqpoint{0.960079in}{3.225613in}}%
\pgfpathlineto{\pgfqpoint{0.967127in}{3.226979in}}%
\pgfpathlineto{\pgfqpoint{0.968391in}{3.231105in}}%
\pgfpathlineto{\pgfqpoint{0.976797in}{3.224912in}}%
\pgfpathlineto{\pgfqpoint{0.978664in}{3.231627in}}%
\pgfpathlineto{\pgfqpoint{0.983081in}{3.234581in}}%
\pgfpathlineto{\pgfqpoint{1.004574in}{3.238691in}}%
\pgfpathlineto{\pgfqpoint{1.011172in}{3.237238in}}%
\pgfpathlineto{\pgfqpoint{1.016700in}{3.245394in}}%
\pgfpathlineto{\pgfqpoint{1.029815in}{3.250744in}}%
\pgfpathlineto{\pgfqpoint{1.039917in}{3.248715in}}%
\pgfpathlineto{\pgfqpoint{1.044526in}{3.241612in}}%
\pgfpathlineto{\pgfqpoint{1.044284in}{3.230327in}}%
\pgfpathlineto{\pgfqpoint{1.049454in}{3.227769in}}%
\pgfpathlineto{\pgfqpoint{1.059999in}{3.228126in}}%
\pgfpathlineto{\pgfqpoint{1.073506in}{3.231062in}}%
\pgfpathlineto{\pgfqpoint{1.073340in}{3.227453in}}%
\pgfpathlineto{\pgfqpoint{1.090413in}{3.217164in}}%
\pgfpathlineto{\pgfqpoint{1.102354in}{3.219923in}}%
\pgfpathlineto{\pgfqpoint{1.105870in}{3.223494in}}%
\pgfpathlineto{\pgfqpoint{1.109634in}{3.232639in}}%
\pgfpathlineto{\pgfqpoint{1.114303in}{3.238496in}}%
\pgfpathlineto{\pgfqpoint{1.114719in}{3.247272in}}%
\pgfpathlineto{\pgfqpoint{1.110445in}{3.250724in}}%
\pgfpathlineto{\pgfqpoint{1.116774in}{3.253843in}}%
\pgfpathlineto{\pgfqpoint{1.127410in}{3.249156in}}%
\pgfpathlineto{\pgfqpoint{1.130874in}{3.252050in}}%
\pgfpathlineto{\pgfqpoint{1.131483in}{3.264466in}}%
\pgfpathlineto{\pgfqpoint{1.124609in}{3.261871in}}%
\pgfpathlineto{\pgfqpoint{1.119841in}{3.266174in}}%
\pgfpathlineto{\pgfqpoint{1.112102in}{3.268855in}}%
\pgfpathlineto{\pgfqpoint{1.099816in}{3.268408in}}%
\pgfpathlineto{\pgfqpoint{1.088046in}{3.272487in}}%
\pgfpathlineto{\pgfqpoint{1.087455in}{3.282649in}}%
\pgfpathlineto{\pgfqpoint{1.076928in}{3.293444in}}%
\pgfpathlineto{\pgfqpoint{1.063579in}{3.297975in}}%
\pgfpathlineto{\pgfqpoint{1.051926in}{3.318872in}}%
\pgfpathlineto{\pgfqpoint{1.052882in}{3.326605in}}%
\pgfpathlineto{\pgfqpoint{1.059156in}{3.330131in}}%
\pgfpathlineto{\pgfqpoint{1.058232in}{3.337452in}}%
\pgfpathlineto{\pgfqpoint{1.059689in}{3.344734in}}%
\pgfpathlineto{\pgfqpoint{1.064654in}{3.339749in}}%
\pgfpathlineto{\pgfqpoint{1.071209in}{3.341829in}}%
\pgfpathlineto{\pgfqpoint{1.070881in}{3.348154in}}%
\pgfpathlineto{\pgfqpoint{1.061792in}{3.360854in}}%
\pgfpathlineto{\pgfqpoint{1.059132in}{3.373954in}}%
\pgfpathlineto{\pgfqpoint{1.066739in}{3.374953in}}%
\pgfpathlineto{\pgfqpoint{1.078126in}{3.374149in}}%
\pgfpathlineto{\pgfqpoint{1.081711in}{3.369767in}}%
\pgfpathlineto{\pgfqpoint{1.091918in}{3.371335in}}%
\pgfpathlineto{\pgfqpoint{1.101406in}{3.369277in}}%
\pgfpathlineto{\pgfqpoint{1.107819in}{3.365658in}}%
\pgfpathlineto{\pgfqpoint{1.112258in}{3.367312in}}%
\pgfpathlineto{\pgfqpoint{1.124045in}{3.363696in}}%
\pgfpathlineto{\pgfqpoint{1.124153in}{3.360493in}}%
\pgfpathlineto{\pgfqpoint{1.133129in}{3.361575in}}%
\pgfpathlineto{\pgfqpoint{1.145165in}{3.353762in}}%
\pgfpathlineto{\pgfqpoint{1.144698in}{3.348577in}}%
\pgfpathlineto{\pgfqpoint{1.138971in}{3.346182in}}%
\pgfpathlineto{\pgfqpoint{1.132884in}{3.340488in}}%
\pgfpathlineto{\pgfqpoint{1.132148in}{3.332669in}}%
\pgfpathlineto{\pgfqpoint{1.145260in}{3.322696in}}%
\pgfpathlineto{\pgfqpoint{1.143163in}{3.317918in}}%
\pgfpathlineto{\pgfqpoint{1.152493in}{3.314247in}}%
\pgfpathlineto{\pgfqpoint{1.154784in}{3.308253in}}%
\pgfpathlineto{\pgfqpoint{1.166624in}{3.313016in}}%
\pgfpathlineto{\pgfqpoint{1.171839in}{3.306774in}}%
\pgfpathlineto{\pgfqpoint{1.174245in}{3.310552in}}%
\pgfpathlineto{\pgfqpoint{1.167159in}{3.322850in}}%
\pgfpathlineto{\pgfqpoint{1.169677in}{3.326534in}}%
\pgfpathlineto{\pgfqpoint{1.170793in}{3.336248in}}%
\pgfpathlineto{\pgfqpoint{1.167971in}{3.340522in}}%
\pgfpathlineto{\pgfqpoint{1.170088in}{3.346265in}}%
\pgfpathlineto{\pgfqpoint{1.177026in}{3.345725in}}%
\pgfpathlineto{\pgfqpoint{1.175590in}{3.335953in}}%
\pgfpathlineto{\pgfqpoint{1.171161in}{3.332693in}}%
\pgfpathlineto{\pgfqpoint{1.171475in}{3.319546in}}%
\pgfpathlineto{\pgfqpoint{1.179422in}{3.317975in}}%
\pgfpathlineto{\pgfqpoint{1.179519in}{3.308201in}}%
\pgfpathlineto{\pgfqpoint{1.188198in}{3.301859in}}%
\pgfpathlineto{\pgfqpoint{1.192188in}{3.305858in}}%
\pgfpathlineto{\pgfqpoint{1.187827in}{3.318028in}}%
\pgfpathlineto{\pgfqpoint{1.183318in}{3.321107in}}%
\pgfpathlineto{\pgfqpoint{1.178882in}{3.319081in}}%
\pgfpathlineto{\pgfqpoint{1.174978in}{3.321667in}}%
\pgfpathlineto{\pgfqpoint{1.175659in}{3.332913in}}%
\pgfpathlineto{\pgfqpoint{1.182173in}{3.337512in}}%
\pgfpathlineto{\pgfqpoint{1.188584in}{3.337089in}}%
\pgfpathlineto{\pgfqpoint{1.185928in}{3.343452in}}%
\pgfpathlineto{\pgfqpoint{1.176318in}{3.348326in}}%
\pgfpathlineto{\pgfqpoint{1.163592in}{3.367746in}}%
\pgfpathlineto{\pgfqpoint{1.169943in}{3.377060in}}%
\pgfpathlineto{\pgfqpoint{1.173478in}{3.385832in}}%
\pgfpathlineto{\pgfqpoint{1.173756in}{3.403448in}}%
\pgfpathlineto{\pgfqpoint{1.172395in}{3.418309in}}%
\pgfpathlineto{\pgfqpoint{1.168512in}{3.428067in}}%
\pgfpathlineto{\pgfqpoint{1.169286in}{3.437824in}}%
\pgfpathlineto{\pgfqpoint{1.179218in}{3.446286in}}%
\pgfpathlineto{\pgfqpoint{1.189095in}{3.456341in}}%
\pgfpathlineto{\pgfqpoint{1.213857in}{3.435796in}}%
\pgfpathlineto{\pgfqpoint{1.230047in}{3.434282in}}%
\pgfpathlineto{\pgfqpoint{1.243410in}{3.439398in}}%
\pgfpathlineto{\pgfqpoint{1.255140in}{3.446135in}}%
\pgfpathlineto{\pgfqpoint{1.257633in}{3.449344in}}%
\pgfpathlineto{\pgfqpoint{1.286589in}{3.456664in}}%
\pgfpathlineto{\pgfqpoint{1.289104in}{3.451302in}}%
\pgfpathlineto{\pgfqpoint{1.298944in}{3.446914in}}%
\pgfpathlineto{\pgfqpoint{1.317722in}{3.448261in}}%
\pgfpathlineto{\pgfqpoint{1.321152in}{3.447062in}}%
\pgfpathlineto{\pgfqpoint{1.329592in}{3.450264in}}%
\pgfpathlineto{\pgfqpoint{1.332985in}{3.444727in}}%
\pgfpathlineto{\pgfqpoint{1.341089in}{3.439668in}}%
\pgfpathlineto{\pgfqpoint{1.345143in}{3.440240in}}%
\pgfpathlineto{\pgfqpoint{1.352990in}{3.434162in}}%
\pgfpathlineto{\pgfqpoint{1.366250in}{3.435279in}}%
\pgfpathlineto{\pgfqpoint{1.379495in}{3.440200in}}%
\pgfpathlineto{\pgfqpoint{1.390311in}{3.424464in}}%
\pgfpathlineto{\pgfqpoint{1.389472in}{3.420932in}}%
\pgfpathlineto{\pgfqpoint{1.377537in}{3.416874in}}%
\pgfpathlineto{\pgfqpoint{1.380726in}{3.410737in}}%
\pgfpathlineto{\pgfqpoint{1.391713in}{3.417196in}}%
\pgfpathlineto{\pgfqpoint{1.396199in}{3.415631in}}%
\pgfpathlineto{\pgfqpoint{1.398107in}{3.407905in}}%
\pgfpathlineto{\pgfqpoint{1.391548in}{3.404638in}}%
\pgfpathlineto{\pgfqpoint{1.396233in}{3.394872in}}%
\pgfpathlineto{\pgfqpoint{1.403035in}{3.396522in}}%
\pgfpathlineto{\pgfqpoint{1.412907in}{3.391482in}}%
\pgfpathlineto{\pgfqpoint{1.422347in}{3.377492in}}%
\pgfpathlineto{\pgfqpoint{1.414040in}{3.373608in}}%
\pgfpathlineto{\pgfqpoint{1.421379in}{3.363300in}}%
\pgfpathlineto{\pgfqpoint{1.415196in}{3.360891in}}%
\pgfpathlineto{\pgfqpoint{1.422888in}{3.351101in}}%
\pgfpathlineto{\pgfqpoint{1.432268in}{3.352144in}}%
\pgfpathlineto{\pgfqpoint{1.437208in}{3.347739in}}%
\pgfpathlineto{\pgfqpoint{1.449701in}{3.341244in}}%
\pgfpathlineto{\pgfqpoint{1.465817in}{3.318189in}}%
\pgfpathlineto{\pgfqpoint{1.465453in}{3.312675in}}%
\pgfpathlineto{\pgfqpoint{1.489622in}{3.295013in}}%
\pgfpathlineto{\pgfqpoint{1.490143in}{3.288766in}}%
\pgfpathlineto{\pgfqpoint{1.496584in}{3.279419in}}%
\pgfpathlineto{\pgfqpoint{1.516372in}{3.274747in}}%
\pgfpathlineto{\pgfqpoint{1.523710in}{3.271424in}}%
\pgfpathlineto{\pgfqpoint{1.530615in}{3.261120in}}%
\pgfpathlineto{\pgfqpoint{1.531504in}{3.252094in}}%
\pgfpathlineto{\pgfqpoint{1.537437in}{3.244738in}}%
\pgfpathlineto{\pgfqpoint{1.539220in}{3.236441in}}%
\pgfpathlineto{\pgfqpoint{1.544780in}{3.233411in}}%
\pgfpathlineto{\pgfqpoint{1.517233in}{3.183693in}}%
\pgfpathlineto{\pgfqpoint{1.460054in}{3.080474in}}%
\pgfpathlineto{\pgfqpoint{1.376162in}{2.929003in}}%
\pgfpathlineto{\pgfqpoint{1.343317in}{2.869686in}}%
\pgfpathlineto{\pgfqpoint{1.350313in}{2.861720in}}%
\pgfpathlineto{\pgfqpoint{1.353395in}{2.864262in}}%
\pgfpathlineto{\pgfqpoint{1.359643in}{2.854988in}}%
\pgfpathlineto{\pgfqpoint{1.368318in}{2.857678in}}%
\pgfpathlineto{\pgfqpoint{1.379832in}{2.852122in}}%
\pgfpathlineto{\pgfqpoint{1.372374in}{2.843439in}}%
\pgfpathlineto{\pgfqpoint{1.378134in}{2.831722in}}%
\pgfpathlineto{\pgfqpoint{1.376943in}{2.826017in}}%
\pgfpathlineto{\pgfqpoint{1.386129in}{2.795933in}}%
\pgfpathlineto{\pgfqpoint{1.382242in}{2.782252in}}%
\pgfpathlineto{\pgfqpoint{1.399793in}{2.785640in}}%
\pgfpathlineto{\pgfqpoint{1.404102in}{2.783420in}}%
\pgfpathlineto{\pgfqpoint{1.408780in}{2.787029in}}%
\pgfpathlineto{\pgfqpoint{1.412108in}{2.793850in}}%
\pgfpathlineto{\pgfqpoint{1.416874in}{2.797760in}}%
\pgfpathlineto{\pgfqpoint{1.437213in}{2.797265in}}%
\pgfpathlineto{\pgfqpoint{1.441809in}{2.784137in}}%
\pgfpathlineto{\pgfqpoint{1.437865in}{2.772613in}}%
\pgfpathlineto{\pgfqpoint{1.442347in}{2.768811in}}%
\pgfpathlineto{\pgfqpoint{1.444703in}{2.762307in}}%
\pgfpathlineto{\pgfqpoint{1.444189in}{2.750075in}}%
\pgfpathlineto{\pgfqpoint{1.450169in}{2.741305in}}%
\pgfpathlineto{\pgfqpoint{1.453448in}{2.727561in}}%
\pgfpathlineto{\pgfqpoint{1.451765in}{2.720138in}}%
\pgfpathlineto{\pgfqpoint{1.453795in}{2.712614in}}%
\pgfpathlineto{\pgfqpoint{1.456722in}{2.669045in}}%
\pgfpathlineto{\pgfqpoint{1.452528in}{2.665613in}}%
\pgfpathlineto{\pgfqpoint{1.458062in}{2.660871in}}%
\pgfpathlineto{\pgfqpoint{1.453713in}{2.655103in}}%
\pgfpathlineto{\pgfqpoint{1.457635in}{2.650400in}}%
\pgfpathlineto{\pgfqpoint{1.455109in}{2.642374in}}%
\pgfpathlineto{\pgfqpoint{1.460758in}{2.640197in}}%
\pgfpathlineto{\pgfqpoint{1.468024in}{2.627989in}}%
\pgfpathlineto{\pgfqpoint{1.473424in}{2.624311in}}%
\pgfpathlineto{\pgfqpoint{1.474837in}{2.618933in}}%
\pgfpathlineto{\pgfqpoint{1.477965in}{2.616604in}}%
\pgfpathlineto{\pgfqpoint{1.477313in}{2.612139in}}%
\pgfpathlineto{\pgfqpoint{1.484056in}{2.608904in}}%
\pgfpathlineto{\pgfqpoint{1.482391in}{2.600313in}}%
\pgfpathlineto{\pgfqpoint{1.476587in}{2.595813in}}%
\pgfpathlineto{\pgfqpoint{1.473950in}{2.587451in}}%
\pgfpathlineto{\pgfqpoint{1.473250in}{2.576254in}}%
\pgfpathlineto{\pgfqpoint{1.470025in}{2.575183in}}%
\pgfpathlineto{\pgfqpoint{1.459486in}{2.565368in}}%
\pgfpathlineto{\pgfqpoint{1.450119in}{2.562777in}}%
\pgfpathlineto{\pgfqpoint{1.445648in}{2.566735in}}%
\pgfpathlineto{\pgfqpoint{1.449968in}{2.577240in}}%
\pgfpathlineto{\pgfqpoint{1.448648in}{2.582397in}}%
\pgfpathlineto{\pgfqpoint{1.454837in}{2.585010in}}%
\pgfpathlineto{\pgfqpoint{1.460625in}{2.600208in}}%
\pgfpathlineto{\pgfqpoint{1.458180in}{2.613190in}}%
\pgfpathlineto{\pgfqpoint{1.454269in}{2.616255in}}%
\pgfpathlineto{\pgfqpoint{1.441186in}{2.615262in}}%
\pgfpathlineto{\pgfqpoint{1.443823in}{2.611871in}}%
\pgfpathlineto{\pgfqpoint{1.434354in}{2.602137in}}%
\pgfpathlineto{\pgfqpoint{1.431303in}{2.607549in}}%
\pgfpathlineto{\pgfqpoint{1.432097in}{2.614620in}}%
\pgfpathlineto{\pgfqpoint{1.437533in}{2.614490in}}%
\pgfpathlineto{\pgfqpoint{1.442372in}{2.619684in}}%
\pgfpathlineto{\pgfqpoint{1.445236in}{2.627269in}}%
\pgfpathlineto{\pgfqpoint{1.456026in}{2.625128in}}%
\pgfpathlineto{\pgfqpoint{1.447388in}{2.629180in}}%
\pgfpathlineto{\pgfqpoint{1.448148in}{2.634607in}}%
\pgfpathlineto{\pgfqpoint{1.444553in}{2.637058in}}%
\pgfpathlineto{\pgfqpoint{1.443082in}{2.644452in}}%
\pgfpathlineto{\pgfqpoint{1.445633in}{2.648303in}}%
\pgfpathlineto{\pgfqpoint{1.439548in}{2.660265in}}%
\pgfpathlineto{\pgfqpoint{1.440061in}{2.667875in}}%
\pgfpathlineto{\pgfqpoint{1.432930in}{2.674651in}}%
\pgfpathlineto{\pgfqpoint{1.429073in}{2.680133in}}%
\pgfpathlineto{\pgfqpoint{1.434400in}{2.685358in}}%
\pgfpathlineto{\pgfqpoint{1.436374in}{2.693167in}}%
\pgfpathlineto{\pgfqpoint{1.434612in}{2.698294in}}%
\pgfpathlineto{\pgfqpoint{1.436673in}{2.704409in}}%
\pgfpathlineto{\pgfqpoint{1.434046in}{2.724346in}}%
\pgfpathlineto{\pgfqpoint{1.430119in}{2.733391in}}%
\pgfpathlineto{\pgfqpoint{1.425627in}{2.736790in}}%
\pgfpathlineto{\pgfqpoint{1.427539in}{2.748556in}}%
\pgfpathlineto{\pgfqpoint{1.426190in}{2.757745in}}%
\pgfpathlineto{\pgfqpoint{1.428633in}{2.763640in}}%
\pgfpathlineto{\pgfqpoint{1.423578in}{2.764465in}}%
\pgfpathlineto{\pgfqpoint{1.422223in}{2.750149in}}%
\pgfpathlineto{\pgfqpoint{1.416543in}{2.734178in}}%
\pgfpathlineto{\pgfqpoint{1.411982in}{2.737815in}}%
\pgfpathlineto{\pgfqpoint{1.411372in}{2.743064in}}%
\pgfpathlineto{\pgfqpoint{1.405792in}{2.750234in}}%
\pgfpathlineto{\pgfqpoint{1.408627in}{2.754929in}}%
\pgfpathlineto{\pgfqpoint{1.407746in}{2.765497in}}%
\pgfpathlineto{\pgfqpoint{1.400524in}{2.761780in}}%
\pgfpathlineto{\pgfqpoint{1.401941in}{2.756048in}}%
\pgfpathlineto{\pgfqpoint{1.400581in}{2.748771in}}%
\pgfpathlineto{\pgfqpoint{1.392474in}{2.748784in}}%
\pgfpathlineto{\pgfqpoint{1.388840in}{2.753295in}}%
\pgfpathlineto{\pgfqpoint{1.383356in}{2.754156in}}%
\pgfpathlineto{\pgfqpoint{1.379665in}{2.759012in}}%
\pgfpathlineto{\pgfqpoint{1.373315in}{2.772347in}}%
\pgfpathlineto{\pgfqpoint{1.371257in}{2.782701in}}%
\pgfpathlineto{\pgfqpoint{1.372574in}{2.787102in}}%
\pgfpathlineto{\pgfqpoint{1.369501in}{2.796794in}}%
\pgfpathlineto{\pgfqpoint{1.364408in}{2.802405in}}%
\pgfpathlineto{\pgfqpoint{1.355600in}{2.816743in}}%
\pgfpathlineto{\pgfqpoint{1.349211in}{2.829198in}}%
\pgfpathlineto{\pgfqpoint{1.355730in}{2.829385in}}%
\pgfpathlineto{\pgfqpoint{1.359566in}{2.832000in}}%
\pgfpathlineto{\pgfqpoint{1.360320in}{2.840094in}}%
\pgfpathlineto{\pgfqpoint{1.341875in}{2.841071in}}%
\pgfpathlineto{\pgfqpoint{1.326187in}{2.858596in}}%
\pgfpathlineto{\pgfqpoint{1.330312in}{2.863205in}}%
\pgfpathlineto{\pgfqpoint{1.322726in}{2.863992in}}%
\pgfpathlineto{\pgfqpoint{1.307659in}{2.880285in}}%
\pgfpathlineto{\pgfqpoint{1.283789in}{2.888454in}}%
\pgfpathlineto{\pgfqpoint{1.280575in}{2.898029in}}%
\pgfpathlineto{\pgfqpoint{1.275622in}{2.903016in}}%
\pgfpathlineto{\pgfqpoint{1.275397in}{2.907699in}}%
\pgfpathlineto{\pgfqpoint{1.271671in}{2.911176in}}%
\pgfpathlineto{\pgfqpoint{1.276561in}{2.916988in}}%
\pgfpathlineto{\pgfqpoint{1.265409in}{2.917549in}}%
\pgfpathlineto{\pgfqpoint{1.262574in}{2.925070in}}%
\pgfpathlineto{\pgfqpoint{1.257499in}{2.928272in}}%
\pgfpathlineto{\pgfqpoint{1.267663in}{2.932137in}}%
\pgfpathlineto{\pgfqpoint{1.262377in}{2.939031in}}%
\pgfpathlineto{\pgfqpoint{1.252649in}{2.941924in}}%
\pgfpathlineto{\pgfqpoint{1.254162in}{2.954658in}}%
\pgfpathlineto{\pgfqpoint{1.251373in}{2.959166in}}%
\pgfpathlineto{\pgfqpoint{1.246350in}{2.962341in}}%
\pgfpathlineto{\pgfqpoint{1.242055in}{2.959156in}}%
\pgfpathlineto{\pgfqpoint{1.239517in}{2.962151in}}%
\pgfpathlineto{\pgfqpoint{1.232116in}{2.963009in}}%
\pgfpathlineto{\pgfqpoint{1.232712in}{2.974564in}}%
\pgfpathlineto{\pgfqpoint{1.221982in}{2.966760in}}%
\pgfpathlineto{\pgfqpoint{1.222208in}{2.950676in}}%
\pgfpathlineto{\pgfqpoint{1.210691in}{2.947592in}}%
\pgfpathlineto{\pgfqpoint{1.203169in}{2.941169in}}%
\pgfpathlineto{\pgfqpoint{1.197123in}{2.939816in}}%
\pgfpathlineto{\pgfqpoint{1.190581in}{2.946289in}}%
\pgfpathlineto{\pgfqpoint{1.184755in}{2.945431in}}%
\pgfpathlineto{\pgfqpoint{1.185835in}{2.951344in}}%
\pgfpathlineto{\pgfqpoint{1.177554in}{2.948160in}}%
\pgfpathlineto{\pgfqpoint{1.169881in}{2.947802in}}%
\pgfpathlineto{\pgfqpoint{1.160900in}{2.944574in}}%
\pgfpathlineto{\pgfqpoint{1.142675in}{2.946391in}}%
\pgfpathlineto{\pgfqpoint{1.136823in}{2.945196in}}%
\pgfpathlineto{\pgfqpoint{1.132601in}{2.950513in}}%
\pgfpathlineto{\pgfqpoint{1.124808in}{2.951044in}}%
\pgfpathlineto{\pgfqpoint{1.123187in}{2.957934in}}%
\pgfpathlineto{\pgfqpoint{1.127832in}{2.961688in}}%
\pgfpathlineto{\pgfqpoint{1.133281in}{2.961335in}}%
\pgfpathlineto{\pgfqpoint{1.137829in}{2.956217in}}%
\pgfpathlineto{\pgfqpoint{1.140824in}{2.964078in}}%
\pgfpathlineto{\pgfqpoint{1.136645in}{2.972787in}}%
\pgfpathlineto{\pgfqpoint{1.146104in}{2.980453in}}%
\pgfpathlineto{\pgfqpoint{1.155583in}{2.983278in}}%
\pgfpathlineto{\pgfqpoint{1.162183in}{2.987887in}}%
\pgfpathlineto{\pgfqpoint{1.166457in}{2.992851in}}%
\pgfpathlineto{\pgfqpoint{1.168843in}{3.000989in}}%
\pgfpathlineto{\pgfqpoint{1.176380in}{2.998774in}}%
\pgfpathlineto{\pgfqpoint{1.192721in}{2.999912in}}%
\pgfpathlineto{\pgfqpoint{1.196229in}{2.990989in}}%
\pgfpathlineto{\pgfqpoint{1.201960in}{2.990600in}}%
\pgfpathlineto{\pgfqpoint{1.212392in}{2.976942in}}%
\pgfpathlineto{\pgfqpoint{1.212411in}{2.981921in}}%
\pgfpathlineto{\pgfqpoint{1.208767in}{2.983540in}}%
\pgfpathlineto{\pgfqpoint{1.202098in}{2.999845in}}%
\pgfpathlineto{\pgfqpoint{1.205343in}{3.002060in}}%
\pgfpathlineto{\pgfqpoint{1.196436in}{3.010567in}}%
\pgfpathlineto{\pgfqpoint{1.188475in}{3.012223in}}%
\pgfpathlineto{\pgfqpoint{1.180890in}{3.009282in}}%
\pgfpathlineto{\pgfqpoint{1.167898in}{3.011573in}}%
\pgfpathlineto{\pgfqpoint{1.161737in}{3.007452in}}%
\pgfpathlineto{\pgfqpoint{1.147708in}{3.004268in}}%
\pgfpathlineto{\pgfqpoint{1.146425in}{2.999686in}}%
\pgfpathlineto{\pgfqpoint{1.135646in}{2.998437in}}%
\pgfpathlineto{\pgfqpoint{1.132880in}{2.993199in}}%
\pgfpathlineto{\pgfqpoint{1.126714in}{2.989361in}}%
\pgfpathlineto{\pgfqpoint{1.119266in}{2.990021in}}%
\pgfpathlineto{\pgfqpoint{1.115548in}{2.986310in}}%
\pgfpathlineto{\pgfqpoint{1.110856in}{2.986485in}}%
\pgfpathlineto{\pgfqpoint{1.106372in}{2.989824in}}%
\pgfpathlineto{\pgfqpoint{1.098533in}{2.989228in}}%
\pgfpathlineto{\pgfqpoint{1.096687in}{2.986150in}}%
\pgfpathlineto{\pgfqpoint{1.088453in}{2.989309in}}%
\pgfpathlineto{\pgfqpoint{1.081462in}{2.981155in}}%
\pgfpathlineto{\pgfqpoint{1.088166in}{2.973265in}}%
\pgfpathlineto{\pgfqpoint{1.090745in}{2.966777in}}%
\pgfpathlineto{\pgfqpoint{1.088367in}{2.961281in}}%
\pgfpathlineto{\pgfqpoint{1.078842in}{2.957488in}}%
\pgfpathlineto{\pgfqpoint{1.073324in}{2.960575in}}%
\pgfpathlineto{\pgfqpoint{1.066378in}{2.958904in}}%
\pgfpathlineto{\pgfqpoint{1.065541in}{2.954087in}}%
\pgfpathlineto{\pgfqpoint{1.057086in}{2.955791in}}%
\pgfpathlineto{\pgfqpoint{1.058958in}{2.950978in}}%
\pgfpathlineto{\pgfqpoint{1.046487in}{2.949963in}}%
\pgfpathlineto{\pgfqpoint{1.038414in}{2.956072in}}%
\pgfpathlineto{\pgfqpoint{1.034102in}{2.952187in}}%
\pgfpathlineto{\pgfqpoint{1.022628in}{2.956323in}}%
\pgfpathlineto{\pgfqpoint{1.019770in}{2.952134in}}%
\pgfpathlineto{\pgfqpoint{1.014082in}{2.950270in}}%
\pgfpathlineto{\pgfqpoint{1.009518in}{2.955182in}}%
\pgfpathlineto{\pgfqpoint{0.993915in}{2.953961in}}%
\pgfpathlineto{\pgfqpoint{0.993931in}{2.946717in}}%
\pgfpathlineto{\pgfqpoint{0.987663in}{2.944800in}}%
\pgfpathlineto{\pgfqpoint{0.982786in}{2.948523in}}%
\pgfpathlineto{\pgfqpoint{0.976340in}{2.946454in}}%
\pgfpathlineto{\pgfqpoint{0.967998in}{2.945961in}}%
\pgfpathlineto{\pgfqpoint{0.955569in}{2.951199in}}%
\pgfpathlineto{\pgfqpoint{0.948453in}{2.945798in}}%
\pgfpathlineto{\pgfqpoint{0.938273in}{2.953016in}}%
\pgfpathlineto{\pgfqpoint{0.934945in}{2.950281in}}%
\pgfpathlineto{\pgfqpoint{0.932608in}{2.943211in}}%
\pgfpathlineto{\pgfqpoint{0.926065in}{2.939066in}}%
\pgfpathlineto{\pgfqpoint{0.915395in}{2.943811in}}%
\pgfpathlineto{\pgfqpoint{0.897718in}{2.955532in}}%
\pgfpathlineto{\pgfqpoint{0.892566in}{2.956574in}}%
\pgfpathlineto{\pgfqpoint{0.889659in}{2.954131in}}%
\pgfpathlineto{\pgfqpoint{0.881884in}{2.955792in}}%
\pgfpathlineto{\pgfqpoint{0.878091in}{2.959440in}}%
\pgfpathlineto{\pgfqpoint{0.870326in}{2.961406in}}%
\pgfpathlineto{\pgfqpoint{0.859720in}{2.961520in}}%
\pgfpathlineto{\pgfqpoint{0.855509in}{2.965871in}}%
\pgfpathlineto{\pgfqpoint{0.863579in}{2.970374in}}%
\pgfpathlineto{\pgfqpoint{0.861696in}{2.974819in}}%
\pgfpathlineto{\pgfqpoint{0.856152in}{2.973821in}}%
\pgfpathlineto{\pgfqpoint{0.853382in}{2.970327in}}%
\pgfpathlineto{\pgfqpoint{0.840401in}{2.965316in}}%
\pgfpathlineto{\pgfqpoint{0.833837in}{2.967021in}}%
\pgfpathlineto{\pgfqpoint{0.831075in}{2.977193in}}%
\pgfpathlineto{\pgfqpoint{0.822887in}{2.972819in}}%
\pgfpathlineto{\pgfqpoint{0.820507in}{2.985712in}}%
\pgfpathlineto{\pgfqpoint{0.816627in}{2.982389in}}%
\pgfpathlineto{\pgfqpoint{0.817276in}{2.977428in}}%
\pgfpathlineto{\pgfqpoint{0.808354in}{2.978738in}}%
\pgfpathlineto{\pgfqpoint{0.817635in}{2.987383in}}%
\pgfpathlineto{\pgfqpoint{0.827041in}{2.982246in}}%
\pgfpathlineto{\pgfqpoint{0.829057in}{2.984825in}}%
\pgfpathlineto{\pgfqpoint{0.838952in}{2.981734in}}%
\pgfpathlineto{\pgfqpoint{0.838334in}{2.985253in}}%
\pgfpathlineto{\pgfqpoint{0.865329in}{2.986418in}}%
\pgfpathlineto{\pgfqpoint{0.878499in}{2.980929in}}%
\pgfpathlineto{\pgfqpoint{0.883595in}{2.975513in}}%
\pgfpathlineto{\pgfqpoint{0.882083in}{2.970785in}}%
\pgfpathlineto{\pgfqpoint{0.887200in}{2.968996in}}%
\pgfpathlineto{\pgfqpoint{0.899915in}{2.976458in}}%
\pgfpathlineto{\pgfqpoint{0.916364in}{2.976936in}}%
\pgfpathlineto{\pgfqpoint{0.931303in}{2.972438in}}%
\pgfpathlineto{\pgfqpoint{0.939764in}{2.973193in}}%
\pgfpathlineto{\pgfqpoint{0.942916in}{2.966215in}}%
\pgfpathlineto{\pgfqpoint{0.947648in}{2.973908in}}%
\pgfpathlineto{\pgfqpoint{0.950558in}{2.975574in}}%
\pgfpathlineto{\pgfqpoint{0.964023in}{2.978918in}}%
\pgfpathlineto{\pgfqpoint{0.969028in}{2.976874in}}%
\pgfpathlineto{\pgfqpoint{0.974439in}{2.979301in}}%
\pgfpathlineto{\pgfqpoint{0.980895in}{2.978032in}}%
\pgfpathlineto{\pgfqpoint{0.982418in}{2.980859in}}%
\pgfpathlineto{\pgfqpoint{0.995888in}{2.993831in}}%
\pgfpathlineto{\pgfqpoint{1.001592in}{2.995755in}}%
\pgfpathlineto{\pgfqpoint{1.004963in}{3.002756in}}%
\pgfpathlineto{\pgfqpoint{1.022854in}{3.008179in}}%
\pgfpathlineto{\pgfqpoint{1.025094in}{3.012315in}}%
\pgfpathlineto{\pgfqpoint{0.999738in}{3.018698in}}%
\pgfpathlineto{\pgfqpoint{1.000735in}{3.025564in}}%
\pgfpathlineto{\pgfqpoint{0.998128in}{3.034632in}}%
\pgfpathlineto{\pgfqpoint{0.992224in}{3.033878in}}%
\pgfpathlineto{\pgfqpoint{0.988207in}{3.022111in}}%
\pgfpathlineto{\pgfqpoint{0.983545in}{3.021141in}}%
\pgfpathlineto{\pgfqpoint{0.980510in}{3.024950in}}%
\pgfpathlineto{\pgfqpoint{0.984050in}{3.041614in}}%
\pgfpathlineto{\pgfqpoint{0.980525in}{3.048597in}}%
\pgfpathlineto{\pgfqpoint{0.973696in}{3.056829in}}%
\pgfpathlineto{\pgfqpoint{0.976978in}{3.063110in}}%
\pgfpathlineto{\pgfqpoint{0.962567in}{3.063336in}}%
\pgfpathlineto{\pgfqpoint{0.962896in}{3.065773in}}%
\pgfpathlineto{\pgfqpoint{0.954297in}{3.069187in}}%
\pgfpathlineto{\pgfqpoint{0.951624in}{3.067308in}}%
\pgfpathclose%
\pgfusepath{stroke,fill}%
\end{pgfscope}%
\begin{pgfscope}%
\pgfpathrectangle{\pgfqpoint{0.100000in}{2.493063in}}{\pgfqpoint{5.037500in}{3.427208in}}%
\pgfusepath{clip}%
\pgfsetbuttcap%
\pgfsetmiterjoin%
\definecolor{currentfill}{rgb}{1.000000,1.000000,1.000000}%
\pgfsetfillcolor{currentfill}%
\pgfsetlinewidth{0.501875pt}%
\definecolor{currentstroke}{rgb}{0.827451,0.827451,0.827451}%
\pgfsetstrokecolor{currentstroke}%
\pgfsetdash{}{0pt}%
\pgfpathmoveto{\pgfqpoint{0.925909in}{3.188677in}}%
\pgfpathlineto{\pgfqpoint{0.924127in}{3.185546in}}%
\pgfpathlineto{\pgfqpoint{0.928841in}{3.178878in}}%
\pgfpathlineto{\pgfqpoint{0.921389in}{3.172640in}}%
\pgfpathlineto{\pgfqpoint{0.921597in}{3.167275in}}%
\pgfpathlineto{\pgfqpoint{0.918227in}{3.166009in}}%
\pgfpathlineto{\pgfqpoint{0.906646in}{3.174010in}}%
\pgfpathlineto{\pgfqpoint{0.898390in}{3.191524in}}%
\pgfpathlineto{\pgfqpoint{0.897838in}{3.196585in}}%
\pgfpathlineto{\pgfqpoint{0.899960in}{3.202642in}}%
\pgfpathlineto{\pgfqpoint{0.908905in}{3.193376in}}%
\pgfpathlineto{\pgfqpoint{0.911642in}{3.195164in}}%
\pgfpathlineto{\pgfqpoint{0.919230in}{3.193463in}}%
\pgfpathlineto{\pgfqpoint{0.925909in}{3.188677in}}%
\pgfpathclose%
\pgfusepath{stroke,fill}%
\end{pgfscope}%
\begin{pgfscope}%
\pgfpathrectangle{\pgfqpoint{0.100000in}{2.493063in}}{\pgfqpoint{5.037500in}{3.427208in}}%
\pgfusepath{clip}%
\pgfsetbuttcap%
\pgfsetmiterjoin%
\definecolor{currentfill}{rgb}{1.000000,1.000000,1.000000}%
\pgfsetfillcolor{currentfill}%
\pgfsetlinewidth{0.501875pt}%
\definecolor{currentstroke}{rgb}{0.827451,0.827451,0.827451}%
\pgfsetstrokecolor{currentstroke}%
\pgfsetdash{}{0pt}%
\pgfpathmoveto{\pgfqpoint{0.788555in}{2.985777in}}%
\pgfpathlineto{\pgfqpoint{0.780416in}{2.984455in}}%
\pgfpathlineto{\pgfqpoint{0.774525in}{2.987178in}}%
\pgfpathlineto{\pgfqpoint{0.771950in}{2.991611in}}%
\pgfpathlineto{\pgfqpoint{0.774868in}{2.997970in}}%
\pgfpathlineto{\pgfqpoint{0.781368in}{2.996288in}}%
\pgfpathlineto{\pgfqpoint{0.792051in}{3.000061in}}%
\pgfpathlineto{\pgfqpoint{0.795793in}{2.994849in}}%
\pgfpathlineto{\pgfqpoint{0.799421in}{2.995812in}}%
\pgfpathlineto{\pgfqpoint{0.808203in}{2.992569in}}%
\pgfpathlineto{\pgfqpoint{0.811998in}{2.988449in}}%
\pgfpathlineto{\pgfqpoint{0.807367in}{2.977674in}}%
\pgfpathlineto{\pgfqpoint{0.802761in}{2.975390in}}%
\pgfpathlineto{\pgfqpoint{0.798621in}{2.976621in}}%
\pgfpathlineto{\pgfqpoint{0.788555in}{2.985777in}}%
\pgfpathclose%
\pgfusepath{stroke,fill}%
\end{pgfscope}%
\begin{pgfscope}%
\pgfpathrectangle{\pgfqpoint{0.100000in}{2.493063in}}{\pgfqpoint{5.037500in}{3.427208in}}%
\pgfusepath{clip}%
\pgfsetbuttcap%
\pgfsetmiterjoin%
\definecolor{currentfill}{rgb}{1.000000,1.000000,1.000000}%
\pgfsetfillcolor{currentfill}%
\pgfsetlinewidth{0.501875pt}%
\definecolor{currentstroke}{rgb}{0.827451,0.827451,0.827451}%
\pgfsetstrokecolor{currentstroke}%
\pgfsetdash{}{0pt}%
\pgfpathmoveto{\pgfqpoint{0.715775in}{2.995567in}}%
\pgfpathlineto{\pgfqpoint{0.704241in}{3.001069in}}%
\pgfpathlineto{\pgfqpoint{0.693568in}{3.002705in}}%
\pgfpathlineto{\pgfqpoint{0.690204in}{3.005048in}}%
\pgfpathlineto{\pgfqpoint{0.694035in}{3.009660in}}%
\pgfpathlineto{\pgfqpoint{0.700833in}{3.003899in}}%
\pgfpathlineto{\pgfqpoint{0.706535in}{3.004693in}}%
\pgfpathlineto{\pgfqpoint{0.715807in}{3.008683in}}%
\pgfpathlineto{\pgfqpoint{0.716557in}{3.013822in}}%
\pgfpathlineto{\pgfqpoint{0.727946in}{3.011278in}}%
\pgfpathlineto{\pgfqpoint{0.725279in}{3.003541in}}%
\pgfpathlineto{\pgfqpoint{0.732535in}{3.002975in}}%
\pgfpathlineto{\pgfqpoint{0.726447in}{2.993772in}}%
\pgfpathlineto{\pgfqpoint{0.720552in}{2.996755in}}%
\pgfpathlineto{\pgfqpoint{0.715775in}{2.995567in}}%
\pgfpathclose%
\pgfusepath{stroke,fill}%
\end{pgfscope}%
\begin{pgfscope}%
\pgfpathrectangle{\pgfqpoint{0.100000in}{2.493063in}}{\pgfqpoint{5.037500in}{3.427208in}}%
\pgfusepath{clip}%
\pgfsetbuttcap%
\pgfsetmiterjoin%
\definecolor{currentfill}{rgb}{1.000000,1.000000,1.000000}%
\pgfsetfillcolor{currentfill}%
\pgfsetlinewidth{0.501875pt}%
\definecolor{currentstroke}{rgb}{0.827451,0.827451,0.827451}%
\pgfsetstrokecolor{currentstroke}%
\pgfsetdash{}{0pt}%
\pgfpathmoveto{\pgfqpoint{0.695511in}{3.014587in}}%
\pgfpathlineto{\pgfqpoint{0.691272in}{3.012282in}}%
\pgfpathlineto{\pgfqpoint{0.679721in}{3.015825in}}%
\pgfpathlineto{\pgfqpoint{0.671090in}{3.013773in}}%
\pgfpathlineto{\pgfqpoint{0.666705in}{3.019602in}}%
\pgfpathlineto{\pgfqpoint{0.668618in}{3.022271in}}%
\pgfpathlineto{\pgfqpoint{0.675159in}{3.022717in}}%
\pgfpathlineto{\pgfqpoint{0.680162in}{3.021042in}}%
\pgfpathlineto{\pgfqpoint{0.685614in}{3.023693in}}%
\pgfpathlineto{\pgfqpoint{0.693972in}{3.020147in}}%
\pgfpathlineto{\pgfqpoint{0.695511in}{3.014587in}}%
\pgfpathclose%
\pgfusepath{stroke,fill}%
\end{pgfscope}%
\begin{pgfscope}%
\pgfpathrectangle{\pgfqpoint{0.100000in}{2.493063in}}{\pgfqpoint{5.037500in}{3.427208in}}%
\pgfusepath{clip}%
\pgfsetbuttcap%
\pgfsetmiterjoin%
\definecolor{currentfill}{rgb}{1.000000,1.000000,1.000000}%
\pgfsetfillcolor{currentfill}%
\pgfsetlinewidth{0.501875pt}%
\definecolor{currentstroke}{rgb}{0.827451,0.827451,0.827451}%
\pgfsetstrokecolor{currentstroke}%
\pgfsetdash{}{0pt}%
\pgfpathmoveto{\pgfqpoint{0.559354in}{3.100376in}}%
\pgfpathlineto{\pgfqpoint{0.559047in}{3.094660in}}%
\pgfpathlineto{\pgfqpoint{0.550615in}{3.091068in}}%
\pgfpathlineto{\pgfqpoint{0.542915in}{3.100060in}}%
\pgfpathlineto{\pgfqpoint{0.551550in}{3.103056in}}%
\pgfpathlineto{\pgfqpoint{0.559354in}{3.100376in}}%
\pgfpathclose%
\pgfusepath{stroke,fill}%
\end{pgfscope}%
\begin{pgfscope}%
\pgfpathrectangle{\pgfqpoint{0.100000in}{2.493063in}}{\pgfqpoint{5.037500in}{3.427208in}}%
\pgfusepath{clip}%
\pgfsetbuttcap%
\pgfsetmiterjoin%
\definecolor{currentfill}{rgb}{1.000000,1.000000,1.000000}%
\pgfsetfillcolor{currentfill}%
\pgfsetlinewidth{0.501875pt}%
\definecolor{currentstroke}{rgb}{0.827451,0.827451,0.827451}%
\pgfsetstrokecolor{currentstroke}%
\pgfsetdash{}{0pt}%
\pgfpathmoveto{\pgfqpoint{0.489807in}{3.132507in}}%
\pgfpathlineto{\pgfqpoint{0.491597in}{3.134992in}}%
\pgfpathlineto{\pgfqpoint{0.501941in}{3.134564in}}%
\pgfpathlineto{\pgfqpoint{0.503755in}{3.138261in}}%
\pgfpathlineto{\pgfqpoint{0.508293in}{3.135035in}}%
\pgfpathlineto{\pgfqpoint{0.505673in}{3.128828in}}%
\pgfpathlineto{\pgfqpoint{0.501818in}{3.128302in}}%
\pgfpathlineto{\pgfqpoint{0.492605in}{3.130250in}}%
\pgfpathlineto{\pgfqpoint{0.489807in}{3.132507in}}%
\pgfpathclose%
\pgfusepath{stroke,fill}%
\end{pgfscope}%
\begin{pgfscope}%
\pgfpathrectangle{\pgfqpoint{0.100000in}{2.493063in}}{\pgfqpoint{5.037500in}{3.427208in}}%
\pgfusepath{clip}%
\pgfsetbuttcap%
\pgfsetmiterjoin%
\definecolor{currentfill}{rgb}{1.000000,1.000000,1.000000}%
\pgfsetfillcolor{currentfill}%
\pgfsetlinewidth{0.501875pt}%
\definecolor{currentstroke}{rgb}{0.827451,0.827451,0.827451}%
\pgfsetstrokecolor{currentstroke}%
\pgfsetdash{}{0pt}%
\pgfpathmoveto{\pgfqpoint{0.475139in}{3.151744in}}%
\pgfpathlineto{\pgfqpoint{0.474512in}{3.157342in}}%
\pgfpathlineto{\pgfqpoint{0.479908in}{3.156928in}}%
\pgfpathlineto{\pgfqpoint{0.479533in}{3.164790in}}%
\pgfpathlineto{\pgfqpoint{0.484961in}{3.160547in}}%
\pgfpathlineto{\pgfqpoint{0.483083in}{3.154414in}}%
\pgfpathlineto{\pgfqpoint{0.475139in}{3.151744in}}%
\pgfpathclose%
\pgfusepath{stroke,fill}%
\end{pgfscope}%
\begin{pgfscope}%
\pgfpathrectangle{\pgfqpoint{0.100000in}{2.493063in}}{\pgfqpoint{5.037500in}{3.427208in}}%
\pgfusepath{clip}%
\pgfsetbuttcap%
\pgfsetmiterjoin%
\definecolor{currentfill}{rgb}{1.000000,1.000000,1.000000}%
\pgfsetfillcolor{currentfill}%
\pgfsetlinewidth{0.501875pt}%
\definecolor{currentstroke}{rgb}{0.827451,0.827451,0.827451}%
\pgfsetstrokecolor{currentstroke}%
\pgfsetdash{}{0pt}%
\pgfpathmoveto{\pgfqpoint{0.950612in}{3.351270in}}%
\pgfpathlineto{\pgfqpoint{0.942005in}{3.353128in}}%
\pgfpathlineto{\pgfqpoint{0.940155in}{3.358880in}}%
\pgfpathlineto{\pgfqpoint{0.943114in}{3.364630in}}%
\pgfpathlineto{\pgfqpoint{0.949739in}{3.367633in}}%
\pgfpathlineto{\pgfqpoint{0.957325in}{3.352874in}}%
\pgfpathlineto{\pgfqpoint{0.964205in}{3.352219in}}%
\pgfpathlineto{\pgfqpoint{0.969338in}{3.347664in}}%
\pgfpathlineto{\pgfqpoint{0.969238in}{3.341398in}}%
\pgfpathlineto{\pgfqpoint{0.965484in}{3.338394in}}%
\pgfpathlineto{\pgfqpoint{0.971588in}{3.319861in}}%
\pgfpathlineto{\pgfqpoint{0.978267in}{3.313060in}}%
\pgfpathlineto{\pgfqpoint{0.971438in}{3.310902in}}%
\pgfpathlineto{\pgfqpoint{0.965912in}{3.318533in}}%
\pgfpathlineto{\pgfqpoint{0.953900in}{3.316170in}}%
\pgfpathlineto{\pgfqpoint{0.957635in}{3.323721in}}%
\pgfpathlineto{\pgfqpoint{0.955896in}{3.326729in}}%
\pgfpathlineto{\pgfqpoint{0.956916in}{3.334600in}}%
\pgfpathlineto{\pgfqpoint{0.954536in}{3.348829in}}%
\pgfpathlineto{\pgfqpoint{0.950612in}{3.351270in}}%
\pgfpathclose%
\pgfusepath{stroke,fill}%
\end{pgfscope}%
\begin{pgfscope}%
\pgfpathrectangle{\pgfqpoint{0.100000in}{2.493063in}}{\pgfqpoint{5.037500in}{3.427208in}}%
\pgfusepath{clip}%
\pgfsetbuttcap%
\pgfsetmiterjoin%
\definecolor{currentfill}{rgb}{1.000000,1.000000,1.000000}%
\pgfsetfillcolor{currentfill}%
\pgfsetlinewidth{0.501875pt}%
\definecolor{currentstroke}{rgb}{0.827451,0.827451,0.827451}%
\pgfsetstrokecolor{currentstroke}%
\pgfsetdash{}{0pt}%
\pgfpathmoveto{\pgfqpoint{1.253018in}{2.927315in}}%
\pgfpathlineto{\pgfqpoint{1.240541in}{2.927455in}}%
\pgfpathlineto{\pgfqpoint{1.241502in}{2.933719in}}%
\pgfpathlineto{\pgfqpoint{1.246496in}{2.935959in}}%
\pgfpathlineto{\pgfqpoint{1.253018in}{2.927315in}}%
\pgfpathclose%
\pgfusepath{stroke,fill}%
\end{pgfscope}%
\begin{pgfscope}%
\pgfpathrectangle{\pgfqpoint{0.100000in}{2.493063in}}{\pgfqpoint{5.037500in}{3.427208in}}%
\pgfusepath{clip}%
\pgfsetbuttcap%
\pgfsetmiterjoin%
\definecolor{currentfill}{rgb}{1.000000,1.000000,1.000000}%
\pgfsetfillcolor{currentfill}%
\pgfsetlinewidth{0.501875pt}%
\definecolor{currentstroke}{rgb}{0.827451,0.827451,0.827451}%
\pgfsetstrokecolor{currentstroke}%
\pgfsetdash{}{0pt}%
\pgfpathmoveto{\pgfqpoint{1.235704in}{2.932987in}}%
\pgfpathlineto{\pgfqpoint{1.226764in}{2.930825in}}%
\pgfpathlineto{\pgfqpoint{1.218270in}{2.927120in}}%
\pgfpathlineto{\pgfqpoint{1.215647in}{2.931485in}}%
\pgfpathlineto{\pgfqpoint{1.230255in}{2.934439in}}%
\pgfpathlineto{\pgfqpoint{1.235704in}{2.932987in}}%
\pgfpathclose%
\pgfusepath{stroke,fill}%
\end{pgfscope}%
\begin{pgfscope}%
\pgfpathrectangle{\pgfqpoint{0.100000in}{2.493063in}}{\pgfqpoint{5.037500in}{3.427208in}}%
\pgfusepath{clip}%
\pgfsetbuttcap%
\pgfsetmiterjoin%
\definecolor{currentfill}{rgb}{1.000000,1.000000,1.000000}%
\pgfsetfillcolor{currentfill}%
\pgfsetlinewidth{0.501875pt}%
\definecolor{currentstroke}{rgb}{0.827451,0.827451,0.827451}%
\pgfsetstrokecolor{currentstroke}%
\pgfsetdash{}{0pt}%
\pgfpathmoveto{\pgfqpoint{1.098678in}{2.928553in}}%
\pgfpathlineto{\pgfqpoint{1.099020in}{2.924222in}}%
\pgfpathlineto{\pgfqpoint{1.093868in}{2.921529in}}%
\pgfpathlineto{\pgfqpoint{1.079494in}{2.927704in}}%
\pgfpathlineto{\pgfqpoint{1.077265in}{2.925220in}}%
\pgfpathlineto{\pgfqpoint{1.073788in}{2.936971in}}%
\pgfpathlineto{\pgfqpoint{1.080989in}{2.938606in}}%
\pgfpathlineto{\pgfqpoint{1.084318in}{2.933777in}}%
\pgfpathlineto{\pgfqpoint{1.091547in}{2.939822in}}%
\pgfpathlineto{\pgfqpoint{1.098678in}{2.928553in}}%
\pgfpathclose%
\pgfusepath{stroke,fill}%
\end{pgfscope}%
\begin{pgfscope}%
\pgfpathrectangle{\pgfqpoint{0.100000in}{2.493063in}}{\pgfqpoint{5.037500in}{3.427208in}}%
\pgfusepath{clip}%
\pgfsetbuttcap%
\pgfsetmiterjoin%
\definecolor{currentfill}{rgb}{1.000000,1.000000,1.000000}%
\pgfsetfillcolor{currentfill}%
\pgfsetlinewidth{0.501875pt}%
\definecolor{currentstroke}{rgb}{0.827451,0.827451,0.827451}%
\pgfsetstrokecolor{currentstroke}%
\pgfsetdash{}{0pt}%
\pgfpathmoveto{\pgfqpoint{1.422926in}{2.688438in}}%
\pgfpathlineto{\pgfqpoint{1.410872in}{2.683950in}}%
\pgfpathlineto{\pgfqpoint{1.404488in}{2.683681in}}%
\pgfpathlineto{\pgfqpoint{1.409684in}{2.693202in}}%
\pgfpathlineto{\pgfqpoint{1.414124in}{2.696752in}}%
\pgfpathlineto{\pgfqpoint{1.413981in}{2.706611in}}%
\pgfpathlineto{\pgfqpoint{1.418807in}{2.724409in}}%
\pgfpathlineto{\pgfqpoint{1.423252in}{2.727789in}}%
\pgfpathlineto{\pgfqpoint{1.433513in}{2.722895in}}%
\pgfpathlineto{\pgfqpoint{1.432015in}{2.720091in}}%
\pgfpathlineto{\pgfqpoint{1.432792in}{2.706619in}}%
\pgfpathlineto{\pgfqpoint{1.429465in}{2.702982in}}%
\pgfpathlineto{\pgfqpoint{1.428037in}{2.693413in}}%
\pgfpathlineto{\pgfqpoint{1.422926in}{2.688438in}}%
\pgfpathclose%
\pgfusepath{stroke,fill}%
\end{pgfscope}%
\begin{pgfscope}%
\pgfpathrectangle{\pgfqpoint{0.100000in}{2.493063in}}{\pgfqpoint{5.037500in}{3.427208in}}%
\pgfusepath{clip}%
\pgfsetbuttcap%
\pgfsetmiterjoin%
\definecolor{currentfill}{rgb}{1.000000,1.000000,1.000000}%
\pgfsetfillcolor{currentfill}%
\pgfsetlinewidth{0.501875pt}%
\definecolor{currentstroke}{rgb}{0.827451,0.827451,0.827451}%
\pgfsetstrokecolor{currentstroke}%
\pgfsetdash{}{0pt}%
\pgfpathmoveto{\pgfqpoint{1.396233in}{2.732557in}}%
\pgfpathlineto{\pgfqpoint{1.403294in}{2.720364in}}%
\pgfpathlineto{\pgfqpoint{1.401985in}{2.717455in}}%
\pgfpathlineto{\pgfqpoint{1.411430in}{2.714498in}}%
\pgfpathlineto{\pgfqpoint{1.408718in}{2.703290in}}%
\pgfpathlineto{\pgfqpoint{1.404606in}{2.703808in}}%
\pgfpathlineto{\pgfqpoint{1.393800in}{2.723346in}}%
\pgfpathlineto{\pgfqpoint{1.395336in}{2.714553in}}%
\pgfpathlineto{\pgfqpoint{1.393496in}{2.709488in}}%
\pgfpathlineto{\pgfqpoint{1.388920in}{2.707235in}}%
\pgfpathlineto{\pgfqpoint{1.386318in}{2.709663in}}%
\pgfpathlineto{\pgfqpoint{1.384866in}{2.720922in}}%
\pgfpathlineto{\pgfqpoint{1.385723in}{2.728224in}}%
\pgfpathlineto{\pgfqpoint{1.382675in}{2.731411in}}%
\pgfpathlineto{\pgfqpoint{1.390726in}{2.740634in}}%
\pgfpathlineto{\pgfqpoint{1.396009in}{2.743371in}}%
\pgfpathlineto{\pgfqpoint{1.397768in}{2.739775in}}%
\pgfpathlineto{\pgfqpoint{1.403280in}{2.742558in}}%
\pgfpathlineto{\pgfqpoint{1.407273in}{2.735001in}}%
\pgfpathlineto{\pgfqpoint{1.415909in}{2.725386in}}%
\pgfpathlineto{\pgfqpoint{1.414730in}{2.721031in}}%
\pgfpathlineto{\pgfqpoint{1.409935in}{2.716286in}}%
\pgfpathlineto{\pgfqpoint{1.405501in}{2.718522in}}%
\pgfpathlineto{\pgfqpoint{1.396233in}{2.732557in}}%
\pgfpathclose%
\pgfusepath{stroke,fill}%
\end{pgfscope}%
\begin{pgfscope}%
\pgfpathrectangle{\pgfqpoint{0.100000in}{2.493063in}}{\pgfqpoint{5.037500in}{3.427208in}}%
\pgfusepath{clip}%
\pgfsetbuttcap%
\pgfsetmiterjoin%
\definecolor{currentfill}{rgb}{1.000000,1.000000,1.000000}%
\pgfsetfillcolor{currentfill}%
\pgfsetlinewidth{0.501875pt}%
\definecolor{currentstroke}{rgb}{0.827451,0.827451,0.827451}%
\pgfsetstrokecolor{currentstroke}%
\pgfsetdash{}{0pt}%
\pgfpathmoveto{\pgfqpoint{1.045813in}{2.905442in}}%
\pgfpathlineto{\pgfqpoint{1.037066in}{2.904304in}}%
\pgfpathlineto{\pgfqpoint{1.024746in}{2.899568in}}%
\pgfpathlineto{\pgfqpoint{1.021873in}{2.902141in}}%
\pgfpathlineto{\pgfqpoint{1.034741in}{2.905049in}}%
\pgfpathlineto{\pgfqpoint{1.031292in}{2.907587in}}%
\pgfpathlineto{\pgfqpoint{1.023313in}{2.909445in}}%
\pgfpathlineto{\pgfqpoint{1.021183in}{2.915391in}}%
\pgfpathlineto{\pgfqpoint{1.025198in}{2.920514in}}%
\pgfpathlineto{\pgfqpoint{1.028578in}{2.932270in}}%
\pgfpathlineto{\pgfqpoint{1.040928in}{2.934460in}}%
\pgfpathlineto{\pgfqpoint{1.046907in}{2.929381in}}%
\pgfpathlineto{\pgfqpoint{1.052523in}{2.934297in}}%
\pgfpathlineto{\pgfqpoint{1.059567in}{2.933091in}}%
\pgfpathlineto{\pgfqpoint{1.064629in}{2.924012in}}%
\pgfpathlineto{\pgfqpoint{1.068916in}{2.932081in}}%
\pgfpathlineto{\pgfqpoint{1.081083in}{2.913205in}}%
\pgfpathlineto{\pgfqpoint{1.076688in}{2.911696in}}%
\pgfpathlineto{\pgfqpoint{1.073990in}{2.907604in}}%
\pgfpathlineto{\pgfqpoint{1.079087in}{2.903870in}}%
\pgfpathlineto{\pgfqpoint{1.071184in}{2.900369in}}%
\pgfpathlineto{\pgfqpoint{1.065454in}{2.902297in}}%
\pgfpathlineto{\pgfqpoint{1.054630in}{2.908686in}}%
\pgfpathlineto{\pgfqpoint{1.045813in}{2.905442in}}%
\pgfpathclose%
\pgfusepath{stroke,fill}%
\end{pgfscope}%
\begin{pgfscope}%
\pgfpathrectangle{\pgfqpoint{0.100000in}{2.493063in}}{\pgfqpoint{5.037500in}{3.427208in}}%
\pgfusepath{clip}%
\pgfsetbuttcap%
\pgfsetmiterjoin%
\definecolor{currentfill}{rgb}{1.000000,1.000000,1.000000}%
\pgfsetfillcolor{currentfill}%
\pgfsetlinewidth{0.501875pt}%
\definecolor{currentstroke}{rgb}{0.827451,0.827451,0.827451}%
\pgfsetstrokecolor{currentstroke}%
\pgfsetdash{}{0pt}%
\pgfpathmoveto{\pgfqpoint{1.388888in}{2.652423in}}%
\pgfpathlineto{\pgfqpoint{1.387153in}{2.669856in}}%
\pgfpathlineto{\pgfqpoint{1.391828in}{2.676257in}}%
\pgfpathlineto{\pgfqpoint{1.386830in}{2.676702in}}%
\pgfpathlineto{\pgfqpoint{1.388657in}{2.689151in}}%
\pgfpathlineto{\pgfqpoint{1.390811in}{2.696542in}}%
\pgfpathlineto{\pgfqpoint{1.388939in}{2.706471in}}%
\pgfpathlineto{\pgfqpoint{1.393517in}{2.708394in}}%
\pgfpathlineto{\pgfqpoint{1.395011in}{2.710962in}}%
\pgfpathlineto{\pgfqpoint{1.399522in}{2.709988in}}%
\pgfpathlineto{\pgfqpoint{1.403426in}{2.702032in}}%
\pgfpathlineto{\pgfqpoint{1.404151in}{2.694744in}}%
\pgfpathlineto{\pgfqpoint{1.399375in}{2.672915in}}%
\pgfpathlineto{\pgfqpoint{1.388888in}{2.652423in}}%
\pgfpathclose%
\pgfusepath{stroke,fill}%
\end{pgfscope}%
\begin{pgfscope}%
\pgfpathrectangle{\pgfqpoint{0.100000in}{2.493063in}}{\pgfqpoint{5.037500in}{3.427208in}}%
\pgfusepath{clip}%
\pgfsetbuttcap%
\pgfsetmiterjoin%
\definecolor{currentfill}{rgb}{1.000000,1.000000,1.000000}%
\pgfsetfillcolor{currentfill}%
\pgfsetlinewidth{0.501875pt}%
\definecolor{currentstroke}{rgb}{0.827451,0.827451,0.827451}%
\pgfsetstrokecolor{currentstroke}%
\pgfsetdash{}{0pt}%
\pgfpathmoveto{\pgfqpoint{1.387850in}{2.705640in}}%
\pgfpathlineto{\pgfqpoint{1.387611in}{2.697140in}}%
\pgfpathlineto{\pgfqpoint{1.383691in}{2.693281in}}%
\pgfpathlineto{\pgfqpoint{1.379160in}{2.694707in}}%
\pgfpathlineto{\pgfqpoint{1.384887in}{2.707054in}}%
\pgfpathlineto{\pgfqpoint{1.387850in}{2.705640in}}%
\pgfpathclose%
\pgfusepath{stroke,fill}%
\end{pgfscope}%
\begin{pgfscope}%
\pgfpathrectangle{\pgfqpoint{0.100000in}{2.493063in}}{\pgfqpoint{5.037500in}{3.427208in}}%
\pgfusepath{clip}%
\pgfsetbuttcap%
\pgfsetmiterjoin%
\definecolor{currentfill}{rgb}{1.000000,1.000000,1.000000}%
\pgfsetfillcolor{currentfill}%
\pgfsetlinewidth{0.501875pt}%
\definecolor{currentstroke}{rgb}{0.827451,0.827451,0.827451}%
\pgfsetstrokecolor{currentstroke}%
\pgfsetdash{}{0pt}%
\pgfpathmoveto{\pgfqpoint{1.435950in}{2.667901in}}%
\pgfpathlineto{\pgfqpoint{1.432607in}{2.653676in}}%
\pgfpathlineto{\pgfqpoint{1.424819in}{2.649179in}}%
\pgfpathlineto{\pgfqpoint{1.414600in}{2.653322in}}%
\pgfpathlineto{\pgfqpoint{1.417416in}{2.658850in}}%
\pgfpathlineto{\pgfqpoint{1.417359in}{2.662678in}}%
\pgfpathlineto{\pgfqpoint{1.421158in}{2.669182in}}%
\pgfpathlineto{\pgfqpoint{1.418488in}{2.667873in}}%
\pgfpathlineto{\pgfqpoint{1.420434in}{2.671049in}}%
\pgfpathlineto{\pgfqpoint{1.416736in}{2.678265in}}%
\pgfpathlineto{\pgfqpoint{1.420933in}{2.679596in}}%
\pgfpathlineto{\pgfqpoint{1.430648in}{2.671152in}}%
\pgfpathlineto{\pgfqpoint{1.435950in}{2.667901in}}%
\pgfpathclose%
\pgfusepath{stroke,fill}%
\end{pgfscope}%
\begin{pgfscope}%
\pgfpathrectangle{\pgfqpoint{0.100000in}{2.493063in}}{\pgfqpoint{5.037500in}{3.427208in}}%
\pgfusepath{clip}%
\pgfsetbuttcap%
\pgfsetmiterjoin%
\definecolor{currentfill}{rgb}{1.000000,1.000000,1.000000}%
\pgfsetfillcolor{currentfill}%
\pgfsetlinewidth{0.501875pt}%
\definecolor{currentstroke}{rgb}{0.827451,0.827451,0.827451}%
\pgfsetstrokecolor{currentstroke}%
\pgfsetdash{}{0pt}%
\pgfpathmoveto{\pgfqpoint{1.401856in}{2.642748in}}%
\pgfpathlineto{\pgfqpoint{1.397450in}{2.641035in}}%
\pgfpathlineto{\pgfqpoint{1.399966in}{2.656046in}}%
\pgfpathlineto{\pgfqpoint{1.405253in}{2.659549in}}%
\pgfpathlineto{\pgfqpoint{1.403964in}{2.670797in}}%
\pgfpathlineto{\pgfqpoint{1.406113in}{2.675784in}}%
\pgfpathlineto{\pgfqpoint{1.410363in}{2.677653in}}%
\pgfpathlineto{\pgfqpoint{1.414824in}{2.672761in}}%
\pgfpathlineto{\pgfqpoint{1.418653in}{2.666122in}}%
\pgfpathlineto{\pgfqpoint{1.406344in}{2.653117in}}%
\pgfpathlineto{\pgfqpoint{1.401856in}{2.642748in}}%
\pgfpathclose%
\pgfusepath{stroke,fill}%
\end{pgfscope}%
\begin{pgfscope}%
\pgfpathrectangle{\pgfqpoint{0.100000in}{2.493063in}}{\pgfqpoint{5.037500in}{3.427208in}}%
\pgfusepath{clip}%
\pgfsetbuttcap%
\pgfsetmiterjoin%
\definecolor{currentfill}{rgb}{1.000000,1.000000,1.000000}%
\pgfsetfillcolor{currentfill}%
\pgfsetlinewidth{0.501875pt}%
\definecolor{currentstroke}{rgb}{0.827451,0.827451,0.827451}%
\pgfsetstrokecolor{currentstroke}%
\pgfsetdash{}{0pt}%
\pgfpathmoveto{\pgfqpoint{1.437904in}{2.658393in}}%
\pgfpathlineto{\pgfqpoint{1.440063in}{2.648035in}}%
\pgfpathlineto{\pgfqpoint{1.430417in}{2.649220in}}%
\pgfpathlineto{\pgfqpoint{1.437904in}{2.658393in}}%
\pgfpathclose%
\pgfusepath{stroke,fill}%
\end{pgfscope}%
\begin{pgfscope}%
\pgfpathrectangle{\pgfqpoint{0.100000in}{2.493063in}}{\pgfqpoint{5.037500in}{3.427208in}}%
\pgfusepath{clip}%
\pgfsetbuttcap%
\pgfsetmiterjoin%
\definecolor{currentfill}{rgb}{1.000000,1.000000,1.000000}%
\pgfsetfillcolor{currentfill}%
\pgfsetlinewidth{0.501875pt}%
\definecolor{currentstroke}{rgb}{0.827451,0.827451,0.827451}%
\pgfsetstrokecolor{currentstroke}%
\pgfsetdash{}{0pt}%
\pgfpathmoveto{\pgfqpoint{1.442508in}{2.642566in}}%
\pgfpathlineto{\pgfqpoint{1.444088in}{2.636175in}}%
\pgfpathlineto{\pgfqpoint{1.447114in}{2.634131in}}%
\pgfpathlineto{\pgfqpoint{1.446563in}{2.628116in}}%
\pgfpathlineto{\pgfqpoint{1.442336in}{2.626093in}}%
\pgfpathlineto{\pgfqpoint{1.438904in}{2.635072in}}%
\pgfpathlineto{\pgfqpoint{1.442508in}{2.642566in}}%
\pgfpathclose%
\pgfusepath{stroke,fill}%
\end{pgfscope}%
\begin{pgfscope}%
\pgfpathrectangle{\pgfqpoint{0.100000in}{2.493063in}}{\pgfqpoint{5.037500in}{3.427208in}}%
\pgfusepath{clip}%
\pgfsetbuttcap%
\pgfsetmiterjoin%
\definecolor{currentfill}{rgb}{1.000000,1.000000,1.000000}%
\pgfsetfillcolor{currentfill}%
\pgfsetlinewidth{0.501875pt}%
\definecolor{currentstroke}{rgb}{0.827451,0.827451,0.827451}%
\pgfsetstrokecolor{currentstroke}%
\pgfsetdash{}{0pt}%
\pgfpathmoveto{\pgfqpoint{1.434457in}{2.645017in}}%
\pgfpathlineto{\pgfqpoint{1.431834in}{2.637594in}}%
\pgfpathlineto{\pgfqpoint{1.427759in}{2.637867in}}%
\pgfpathlineto{\pgfqpoint{1.425189in}{2.643988in}}%
\pgfpathlineto{\pgfqpoint{1.429376in}{2.646887in}}%
\pgfpathlineto{\pgfqpoint{1.434457in}{2.645017in}}%
\pgfpathclose%
\pgfusepath{stroke,fill}%
\end{pgfscope}%
\begin{pgfscope}%
\pgfpathrectangle{\pgfqpoint{0.100000in}{2.493063in}}{\pgfqpoint{5.037500in}{3.427208in}}%
\pgfusepath{clip}%
\pgfsetbuttcap%
\pgfsetmiterjoin%
\definecolor{currentfill}{rgb}{1.000000,1.000000,1.000000}%
\pgfsetfillcolor{currentfill}%
\pgfsetlinewidth{0.501875pt}%
\definecolor{currentstroke}{rgb}{0.827451,0.827451,0.827451}%
\pgfsetstrokecolor{currentstroke}%
\pgfsetdash{}{0pt}%
\pgfpathmoveto{\pgfqpoint{1.423464in}{2.606748in}}%
\pgfpathlineto{\pgfqpoint{1.427100in}{2.602771in}}%
\pgfpathlineto{\pgfqpoint{1.427901in}{2.594391in}}%
\pgfpathlineto{\pgfqpoint{1.429814in}{2.591956in}}%
\pgfpathlineto{\pgfqpoint{1.426689in}{2.584374in}}%
\pgfpathlineto{\pgfqpoint{1.422577in}{2.580104in}}%
\pgfpathlineto{\pgfqpoint{1.420388in}{2.572184in}}%
\pgfpathlineto{\pgfqpoint{1.415893in}{2.584527in}}%
\pgfpathlineto{\pgfqpoint{1.415148in}{2.595734in}}%
\pgfpathlineto{\pgfqpoint{1.412633in}{2.601313in}}%
\pgfpathlineto{\pgfqpoint{1.403970in}{2.605836in}}%
\pgfpathlineto{\pgfqpoint{1.412096in}{2.615268in}}%
\pgfpathlineto{\pgfqpoint{1.406679in}{2.619299in}}%
\pgfpathlineto{\pgfqpoint{1.411451in}{2.623316in}}%
\pgfpathlineto{\pgfqpoint{1.417787in}{2.638739in}}%
\pgfpathlineto{\pgfqpoint{1.411220in}{2.644177in}}%
\pgfpathlineto{\pgfqpoint{1.413750in}{2.649397in}}%
\pgfpathlineto{\pgfqpoint{1.422238in}{2.644766in}}%
\pgfpathlineto{\pgfqpoint{1.423098in}{2.637458in}}%
\pgfpathlineto{\pgfqpoint{1.419886in}{2.633296in}}%
\pgfpathlineto{\pgfqpoint{1.426250in}{2.627666in}}%
\pgfpathlineto{\pgfqpoint{1.428254in}{2.618896in}}%
\pgfpathlineto{\pgfqpoint{1.427982in}{2.606595in}}%
\pgfpathlineto{\pgfqpoint{1.423464in}{2.606748in}}%
\pgfpathclose%
\pgfusepath{stroke,fill}%
\end{pgfscope}%
\begin{pgfscope}%
\pgfpathrectangle{\pgfqpoint{0.100000in}{2.493063in}}{\pgfqpoint{5.037500in}{3.427208in}}%
\pgfusepath{clip}%
\pgfsetbuttcap%
\pgfsetmiterjoin%
\definecolor{currentfill}{rgb}{1.000000,1.000000,1.000000}%
\pgfsetfillcolor{currentfill}%
\pgfsetlinewidth{0.501875pt}%
\definecolor{currentstroke}{rgb}{0.827451,0.827451,0.827451}%
\pgfsetstrokecolor{currentstroke}%
\pgfsetdash{}{0pt}%
\pgfpathmoveto{\pgfqpoint{1.439005in}{2.637624in}}%
\pgfpathlineto{\pgfqpoint{1.437026in}{2.633358in}}%
\pgfpathlineto{\pgfqpoint{1.438885in}{2.626429in}}%
\pgfpathlineto{\pgfqpoint{1.434125in}{2.625847in}}%
\pgfpathlineto{\pgfqpoint{1.429134in}{2.629054in}}%
\pgfpathlineto{\pgfqpoint{1.430754in}{2.636120in}}%
\pgfpathlineto{\pgfqpoint{1.439005in}{2.637624in}}%
\pgfpathclose%
\pgfusepath{stroke,fill}%
\end{pgfscope}%
\begin{pgfscope}%
\pgfpathrectangle{\pgfqpoint{0.100000in}{2.493063in}}{\pgfqpoint{5.037500in}{3.427208in}}%
\pgfusepath{clip}%
\pgfsetbuttcap%
\pgfsetmiterjoin%
\definecolor{currentfill}{rgb}{1.000000,1.000000,1.000000}%
\pgfsetfillcolor{currentfill}%
\pgfsetlinewidth{0.501875pt}%
\definecolor{currentstroke}{rgb}{0.827451,0.827451,0.827451}%
\pgfsetstrokecolor{currentstroke}%
\pgfsetdash{}{0pt}%
\pgfpathmoveto{\pgfqpoint{1.416821in}{2.638107in}}%
\pgfpathlineto{\pgfqpoint{1.406314in}{2.633395in}}%
\pgfpathlineto{\pgfqpoint{1.402743in}{2.635090in}}%
\pgfpathlineto{\pgfqpoint{1.409908in}{2.641307in}}%
\pgfpathlineto{\pgfqpoint{1.416821in}{2.638107in}}%
\pgfpathclose%
\pgfusepath{stroke,fill}%
\end{pgfscope}%
\begin{pgfscope}%
\pgfpathrectangle{\pgfqpoint{0.100000in}{2.493063in}}{\pgfqpoint{5.037500in}{3.427208in}}%
\pgfusepath{clip}%
\pgfsetbuttcap%
\pgfsetmiterjoin%
\definecolor{currentfill}{rgb}{1.000000,1.000000,1.000000}%
\pgfsetfillcolor{currentfill}%
\pgfsetlinewidth{0.501875pt}%
\definecolor{currentstroke}{rgb}{0.827451,0.827451,0.827451}%
\pgfsetstrokecolor{currentstroke}%
\pgfsetdash{}{0pt}%
\pgfpathmoveto{\pgfqpoint{1.439121in}{2.592137in}}%
\pgfpathlineto{\pgfqpoint{1.436616in}{2.599232in}}%
\pgfpathlineto{\pgfqpoint{1.441914in}{2.600977in}}%
\pgfpathlineto{\pgfqpoint{1.443518in}{2.608629in}}%
\pgfpathlineto{\pgfqpoint{1.449259in}{2.608637in}}%
\pgfpathlineto{\pgfqpoint{1.449235in}{2.613965in}}%
\pgfpathlineto{\pgfqpoint{1.457214in}{2.612705in}}%
\pgfpathlineto{\pgfqpoint{1.458574in}{2.598262in}}%
\pgfpathlineto{\pgfqpoint{1.453946in}{2.589054in}}%
\pgfpathlineto{\pgfqpoint{1.447079in}{2.583302in}}%
\pgfpathlineto{\pgfqpoint{1.439121in}{2.592137in}}%
\pgfpathclose%
\pgfusepath{stroke,fill}%
\end{pgfscope}%
\begin{pgfscope}%
\pgfpathrectangle{\pgfqpoint{0.100000in}{2.493063in}}{\pgfqpoint{5.037500in}{3.427208in}}%
\pgfusepath{clip}%
\pgfsetbuttcap%
\pgfsetmiterjoin%
\definecolor{currentfill}{rgb}{1.000000,1.000000,1.000000}%
\pgfsetfillcolor{currentfill}%
\pgfsetlinewidth{0.501875pt}%
\definecolor{currentstroke}{rgb}{0.827451,0.827451,0.827451}%
\pgfsetstrokecolor{currentstroke}%
\pgfsetdash{}{0pt}%
\pgfpathmoveto{\pgfqpoint{1.436019in}{2.597594in}}%
\pgfpathlineto{\pgfqpoint{1.438281in}{2.590876in}}%
\pgfpathlineto{\pgfqpoint{1.432702in}{2.590059in}}%
\pgfpathlineto{\pgfqpoint{1.436019in}{2.597594in}}%
\pgfpathclose%
\pgfusepath{stroke,fill}%
\end{pgfscope}%
\begin{pgfscope}%
\pgfpathrectangle{\pgfqpoint{0.100000in}{2.493063in}}{\pgfqpoint{5.037500in}{3.427208in}}%
\pgfusepath{clip}%
\pgfsetbuttcap%
\pgfsetmiterjoin%
\definecolor{currentfill}{rgb}{1.000000,1.000000,1.000000}%
\pgfsetfillcolor{currentfill}%
\pgfsetlinewidth{0.501875pt}%
\definecolor{currentstroke}{rgb}{0.827451,0.827451,0.827451}%
\pgfsetstrokecolor{currentstroke}%
\pgfsetdash{}{0pt}%
\pgfpathmoveto{\pgfqpoint{1.440722in}{2.587713in}}%
\pgfpathlineto{\pgfqpoint{1.439944in}{2.579149in}}%
\pgfpathlineto{\pgfqpoint{1.434058in}{2.579858in}}%
\pgfpathlineto{\pgfqpoint{1.435962in}{2.583996in}}%
\pgfpathlineto{\pgfqpoint{1.440722in}{2.587713in}}%
\pgfpathclose%
\pgfusepath{stroke,fill}%
\end{pgfscope}%
\begin{pgfscope}%
\pgfpathrectangle{\pgfqpoint{0.100000in}{2.493063in}}{\pgfqpoint{5.037500in}{3.427208in}}%
\pgfusepath{clip}%
\pgfsetbuttcap%
\pgfsetroundjoin%
\pgfsetlinewidth{0.501875pt}%
\definecolor{currentstroke}{rgb}{0.827451,0.827451,0.827451}%
\pgfsetstrokecolor{currentstroke}%
\pgfsetdash{}{0pt}%
\pgfusepath{stroke}%
\end{pgfscope}%
\begin{pgfscope}%
\pgfpathrectangle{\pgfqpoint{0.100000in}{2.493063in}}{\pgfqpoint{5.037500in}{3.427208in}}%
\pgfusepath{clip}%
\pgfsetbuttcap%
\pgfsetroundjoin%
\pgfsetlinewidth{0.501875pt}%
\definecolor{currentstroke}{rgb}{0.827451,0.827451,0.827451}%
\pgfsetstrokecolor{currentstroke}%
\pgfsetdash{}{0pt}%
\pgfusepath{stroke}%
\end{pgfscope}%
\begin{pgfscope}%
\pgfpathrectangle{\pgfqpoint{0.100000in}{2.493063in}}{\pgfqpoint{5.037500in}{3.427208in}}%
\pgfusepath{clip}%
\pgfsetbuttcap%
\pgfsetroundjoin%
\pgfsetlinewidth{0.501875pt}%
\definecolor{currentstroke}{rgb}{0.827451,0.827451,0.827451}%
\pgfsetstrokecolor{currentstroke}%
\pgfsetdash{}{0pt}%
\pgfusepath{stroke}%
\end{pgfscope}%
\begin{pgfscope}%
\pgfpathrectangle{\pgfqpoint{0.100000in}{2.493063in}}{\pgfqpoint{5.037500in}{3.427208in}}%
\pgfusepath{clip}%
\pgfsetbuttcap%
\pgfsetroundjoin%
\pgfsetlinewidth{0.501875pt}%
\definecolor{currentstroke}{rgb}{0.827451,0.827451,0.827451}%
\pgfsetstrokecolor{currentstroke}%
\pgfsetdash{}{0pt}%
\pgfusepath{stroke}%
\end{pgfscope}%
\begin{pgfscope}%
\pgfpathrectangle{\pgfqpoint{0.100000in}{2.493063in}}{\pgfqpoint{5.037500in}{3.427208in}}%
\pgfusepath{clip}%
\pgfsetbuttcap%
\pgfsetroundjoin%
\pgfsetlinewidth{0.501875pt}%
\definecolor{currentstroke}{rgb}{0.827451,0.827451,0.827451}%
\pgfsetstrokecolor{currentstroke}%
\pgfsetdash{}{0pt}%
\pgfusepath{stroke}%
\end{pgfscope}%
\begin{pgfscope}%
\pgfpathrectangle{\pgfqpoint{0.100000in}{2.493063in}}{\pgfqpoint{5.037500in}{3.427208in}}%
\pgfusepath{clip}%
\pgfsetbuttcap%
\pgfsetroundjoin%
\pgfsetlinewidth{0.501875pt}%
\definecolor{currentstroke}{rgb}{0.827451,0.827451,0.827451}%
\pgfsetstrokecolor{currentstroke}%
\pgfsetdash{}{0pt}%
\pgfusepath{stroke}%
\end{pgfscope}%
\begin{pgfscope}%
\pgfpathrectangle{\pgfqpoint{0.100000in}{2.493063in}}{\pgfqpoint{5.037500in}{3.427208in}}%
\pgfusepath{clip}%
\pgfsetbuttcap%
\pgfsetroundjoin%
\pgfsetlinewidth{0.501875pt}%
\definecolor{currentstroke}{rgb}{0.827451,0.827451,0.827451}%
\pgfsetstrokecolor{currentstroke}%
\pgfsetdash{}{0pt}%
\pgfpathmoveto{\pgfqpoint{1.129685in}{5.722352in}}%
\pgfpathlineto{\pgfqpoint{1.102423in}{5.614332in}}%
\pgfpathlineto{\pgfqpoint{1.083039in}{5.536934in}}%
\pgfpathlineto{\pgfqpoint{1.059991in}{5.443826in}}%
\pgfpathlineto{\pgfqpoint{1.055466in}{5.420878in}}%
\pgfpathlineto{\pgfqpoint{1.058421in}{5.399829in}}%
\pgfpathlineto{\pgfqpoint{1.054504in}{5.380440in}}%
\pgfpathlineto{\pgfqpoint{0.893620in}{5.422423in}}%
\pgfpathlineto{\pgfqpoint{0.879086in}{5.417547in}}%
\pgfpathlineto{\pgfqpoint{0.866649in}{5.421778in}}%
\pgfpathlineto{\pgfqpoint{0.809066in}{5.423162in}}%
\pgfpathlineto{\pgfqpoint{0.789590in}{5.417671in}}%
\pgfpathlineto{\pgfqpoint{0.770250in}{5.419526in}}%
\pgfpathlineto{\pgfqpoint{0.762052in}{5.428120in}}%
\pgfpathlineto{\pgfqpoint{0.710499in}{5.426247in}}%
\pgfpathlineto{\pgfqpoint{0.702357in}{5.440292in}}%
\pgfpathlineto{\pgfqpoint{0.686901in}{5.447691in}}%
\pgfpathlineto{\pgfqpoint{0.664413in}{5.452200in}}%
\pgfpathlineto{\pgfqpoint{0.625598in}{5.445598in}}%
\pgfpathlineto{\pgfqpoint{0.611269in}{5.452066in}}%
\pgfpathlineto{\pgfqpoint{0.589177in}{5.469269in}}%
\pgfpathlineto{\pgfqpoint{0.595818in}{5.502986in}}%
\pgfpathlineto{\pgfqpoint{0.593490in}{5.520205in}}%
\pgfpathlineto{\pgfqpoint{0.575981in}{5.538485in}}%
\pgfpathlineto{\pgfqpoint{0.564765in}{5.537351in}}%
\pgfpathlineto{\pgfqpoint{0.556720in}{5.555841in}}%
\pgfpathlineto{\pgfqpoint{0.537662in}{5.563315in}}%
\pgfpathlineto{\pgfqpoint{0.523809in}{5.562363in}}%
\pgfpathlineto{\pgfqpoint{0.519327in}{5.581353in}}%
\pgfpathlineto{\pgfqpoint{0.533128in}{5.579332in}}%
\pgfpathlineto{\pgfqpoint{0.532031in}{5.605726in}}%
\pgfpathlineto{\pgfqpoint{0.526017in}{5.621398in}}%
\pgfpathlineto{\pgfqpoint{0.528668in}{5.641458in}}%
\pgfpathlineto{\pgfqpoint{0.535889in}{5.656545in}}%
\pgfpathlineto{\pgfqpoint{0.535482in}{5.685175in}}%
\pgfpathlineto{\pgfqpoint{0.531627in}{5.695562in}}%
\pgfpathlineto{\pgfqpoint{0.538284in}{5.728888in}}%
\pgfpathlineto{\pgfqpoint{0.536362in}{5.750367in}}%
\pgfpathlineto{\pgfqpoint{0.529707in}{5.760665in}}%
\pgfpathlineto{\pgfqpoint{0.530722in}{5.794337in}}%
\pgfpathlineto{\pgfqpoint{0.540403in}{5.819099in}}%
\pgfpathlineto{\pgfqpoint{0.550932in}{5.813039in}}%
\pgfpathlineto{\pgfqpoint{0.585796in}{5.777025in}}%
\pgfpathlineto{\pgfqpoint{0.628042in}{5.757317in}}%
\pgfpathlineto{\pgfqpoint{0.649673in}{5.754985in}}%
\pgfpathlineto{\pgfqpoint{0.670052in}{5.746594in}}%
\pgfpathlineto{\pgfqpoint{0.675837in}{5.718359in}}%
\pgfpathlineto{\pgfqpoint{0.657977in}{5.713010in}}%
\pgfpathlineto{\pgfqpoint{0.646142in}{5.690758in}}%
\pgfpathlineto{\pgfqpoint{0.660038in}{5.691994in}}%
\pgfpathlineto{\pgfqpoint{0.665641in}{5.701983in}}%
\pgfpathlineto{\pgfqpoint{0.685238in}{5.714473in}}%
\pgfpathlineto{\pgfqpoint{0.684226in}{5.695924in}}%
\pgfpathlineto{\pgfqpoint{0.671141in}{5.692934in}}%
\pgfpathlineto{\pgfqpoint{0.673314in}{5.669093in}}%
\pgfpathlineto{\pgfqpoint{0.657380in}{5.645210in}}%
\pgfpathlineto{\pgfqpoint{0.647182in}{5.656567in}}%
\pgfpathlineto{\pgfqpoint{0.616877in}{5.650392in}}%
\pgfpathlineto{\pgfqpoint{0.615283in}{5.636499in}}%
\pgfpathlineto{\pgfqpoint{0.625730in}{5.628035in}}%
\pgfpathlineto{\pgfqpoint{0.641668in}{5.627203in}}%
\pgfpathlineto{\pgfqpoint{0.663706in}{5.645290in}}%
\pgfpathlineto{\pgfqpoint{0.681148in}{5.646792in}}%
\pgfpathlineto{\pgfqpoint{0.681817in}{5.666812in}}%
\pgfpathlineto{\pgfqpoint{0.690807in}{5.696218in}}%
\pgfpathlineto{\pgfqpoint{0.712110in}{5.718088in}}%
\pgfpathlineto{\pgfqpoint{0.705007in}{5.735027in}}%
\pgfpathlineto{\pgfqpoint{0.709859in}{5.753264in}}%
\pgfpathlineto{\pgfqpoint{0.700407in}{5.771439in}}%
\pgfpathlineto{\pgfqpoint{0.713877in}{5.794682in}}%
\pgfpathlineto{\pgfqpoint{0.715895in}{5.808665in}}%
\pgfpathlineto{\pgfqpoint{0.704200in}{5.817856in}}%
\pgfpathlineto{\pgfqpoint{0.706147in}{5.841364in}}%
\pgfpathlineto{\pgfqpoint{0.846264in}{5.798794in}}%
\pgfpathlineto{\pgfqpoint{0.995055in}{5.757070in}}%
\pgfpathlineto{\pgfqpoint{1.129685in}{5.722352in}}%
\pgfusepath{stroke}%
\end{pgfscope}%
\begin{pgfscope}%
\pgfpathrectangle{\pgfqpoint{0.100000in}{2.493063in}}{\pgfqpoint{5.037500in}{3.427208in}}%
\pgfusepath{clip}%
\pgfsetbuttcap%
\pgfsetroundjoin%
\pgfsetlinewidth{0.501875pt}%
\definecolor{currentstroke}{rgb}{0.827451,0.827451,0.827451}%
\pgfsetstrokecolor{currentstroke}%
\pgfsetdash{}{0pt}%
\pgfpathmoveto{\pgfqpoint{0.683098in}{5.760037in}}%
\pgfpathlineto{\pgfqpoint{0.689997in}{5.733438in}}%
\pgfpathlineto{\pgfqpoint{0.700816in}{5.724603in}}%
\pgfpathlineto{\pgfqpoint{0.692183in}{5.713355in}}%
\pgfpathlineto{\pgfqpoint{0.683769in}{5.729836in}}%
\pgfpathlineto{\pgfqpoint{0.683098in}{5.760037in}}%
\pgfusepath{stroke}%
\end{pgfscope}%
\begin{pgfscope}%
\pgfpathrectangle{\pgfqpoint{0.100000in}{2.493063in}}{\pgfqpoint{5.037500in}{3.427208in}}%
\pgfusepath{clip}%
\pgfsetbuttcap%
\pgfsetroundjoin%
\pgfsetlinewidth{0.501875pt}%
\definecolor{currentstroke}{rgb}{0.827451,0.827451,0.827451}%
\pgfsetstrokecolor{currentstroke}%
\pgfsetdash{}{0pt}%
\pgfpathmoveto{\pgfqpoint{1.201939in}{5.704893in}}%
\pgfpathlineto{\pgfqpoint{1.351716in}{5.671281in}}%
\pgfpathlineto{\pgfqpoint{1.492772in}{5.642770in}}%
\pgfpathlineto{\pgfqpoint{1.601299in}{5.622885in}}%
\pgfpathlineto{\pgfqpoint{1.695915in}{5.606989in}}%
\pgfpathlineto{\pgfqpoint{1.790740in}{5.592398in}}%
\pgfpathlineto{\pgfqpoint{1.871493in}{5.581022in}}%
\pgfpathlineto{\pgfqpoint{1.952373in}{5.570591in}}%
\pgfpathlineto{\pgfqpoint{2.033369in}{5.561105in}}%
\pgfpathlineto{\pgfqpoint{2.109697in}{5.553044in}}%
\pgfpathlineto{\pgfqpoint{2.099200in}{5.436717in}}%
\pgfpathlineto{\pgfqpoint{2.083575in}{5.279518in}}%
\pgfpathlineto{\pgfqpoint{2.075380in}{5.198668in}}%
\pgfpathlineto{\pgfqpoint{2.063595in}{5.089707in}}%
\pgfpathlineto{\pgfqpoint{1.980159in}{5.098817in}}%
\pgfpathlineto{\pgfqpoint{1.884635in}{5.109661in}}%
\pgfpathlineto{\pgfqpoint{1.751998in}{5.127814in}}%
\pgfpathlineto{\pgfqpoint{1.692759in}{5.136264in}}%
\pgfpathlineto{\pgfqpoint{1.546810in}{5.159174in}}%
\pgfpathlineto{\pgfqpoint{1.496569in}{5.168365in}}%
\pgfpathlineto{\pgfqpoint{1.486097in}{5.108827in}}%
\pgfpathlineto{\pgfqpoint{1.480372in}{5.113075in}}%
\pgfpathlineto{\pgfqpoint{1.469613in}{5.141717in}}%
\pgfpathlineto{\pgfqpoint{1.455784in}{5.139953in}}%
\pgfpathlineto{\pgfqpoint{1.445885in}{5.125474in}}%
\pgfpathlineto{\pgfqpoint{1.425343in}{5.123537in}}%
\pgfpathlineto{\pgfqpoint{1.421216in}{5.130465in}}%
\pgfpathlineto{\pgfqpoint{1.401974in}{5.129626in}}%
\pgfpathlineto{\pgfqpoint{1.392301in}{5.136334in}}%
\pgfpathlineto{\pgfqpoint{1.378775in}{5.125956in}}%
\pgfpathlineto{\pgfqpoint{1.345921in}{5.135295in}}%
\pgfpathlineto{\pgfqpoint{1.331583in}{5.130280in}}%
\pgfpathlineto{\pgfqpoint{1.323870in}{5.152534in}}%
\pgfpathlineto{\pgfqpoint{1.323509in}{5.173709in}}%
\pgfpathlineto{\pgfqpoint{1.300862in}{5.188938in}}%
\pgfpathlineto{\pgfqpoint{1.304555in}{5.210885in}}%
\pgfpathlineto{\pgfqpoint{1.288430in}{5.247138in}}%
\pgfpathlineto{\pgfqpoint{1.289516in}{5.277955in}}%
\pgfpathlineto{\pgfqpoint{1.275657in}{5.293038in}}%
\pgfpathlineto{\pgfqpoint{1.262653in}{5.279694in}}%
\pgfpathlineto{\pgfqpoint{1.244953in}{5.272477in}}%
\pgfpathlineto{\pgfqpoint{1.231480in}{5.287352in}}%
\pgfpathlineto{\pgfqpoint{1.233671in}{5.311107in}}%
\pgfpathlineto{\pgfqpoint{1.249829in}{5.320466in}}%
\pgfpathlineto{\pgfqpoint{1.245534in}{5.335516in}}%
\pgfpathlineto{\pgfqpoint{1.273621in}{5.408361in}}%
\pgfpathlineto{\pgfqpoint{1.251474in}{5.410187in}}%
\pgfpathlineto{\pgfqpoint{1.249197in}{5.423259in}}%
\pgfpathlineto{\pgfqpoint{1.233064in}{5.434532in}}%
\pgfpathlineto{\pgfqpoint{1.234095in}{5.447134in}}%
\pgfpathlineto{\pgfqpoint{1.225585in}{5.456924in}}%
\pgfpathlineto{\pgfqpoint{1.211732in}{5.492469in}}%
\pgfpathlineto{\pgfqpoint{1.199072in}{5.499718in}}%
\pgfpathlineto{\pgfqpoint{1.190007in}{5.530394in}}%
\pgfpathlineto{\pgfqpoint{1.192146in}{5.550785in}}%
\pgfpathlineto{\pgfqpoint{1.175235in}{5.588340in}}%
\pgfpathlineto{\pgfqpoint{1.201939in}{5.704893in}}%
\pgfusepath{stroke}%
\end{pgfscope}%
\begin{pgfscope}%
\pgfpathrectangle{\pgfqpoint{0.100000in}{2.493063in}}{\pgfqpoint{5.037500in}{3.427208in}}%
\pgfusepath{clip}%
\pgfsetbuttcap%
\pgfsetroundjoin%
\pgfsetlinewidth{0.501875pt}%
\definecolor{currentstroke}{rgb}{0.827451,0.827451,0.827451}%
\pgfsetstrokecolor{currentstroke}%
\pgfsetdash{}{0pt}%
\pgfpathmoveto{\pgfqpoint{4.817670in}{5.105444in}}%
\pgfpathlineto{\pgfqpoint{4.814672in}{5.118191in}}%
\pgfpathlineto{\pgfqpoint{4.798011in}{5.129357in}}%
\pgfpathlineto{\pgfqpoint{4.753984in}{5.273372in}}%
\pgfpathlineto{\pgfqpoint{4.730253in}{5.342870in}}%
\pgfpathlineto{\pgfqpoint{4.749331in}{5.362892in}}%
\pgfpathlineto{\pgfqpoint{4.767877in}{5.411893in}}%
\pgfpathlineto{\pgfqpoint{4.777126in}{5.427553in}}%
\pgfpathlineto{\pgfqpoint{4.770574in}{5.434129in}}%
\pgfpathlineto{\pgfqpoint{4.768379in}{5.477626in}}%
\pgfpathlineto{\pgfqpoint{4.776654in}{5.490960in}}%
\pgfpathlineto{\pgfqpoint{4.773027in}{5.521953in}}%
\pgfpathlineto{\pgfqpoint{4.806052in}{5.623639in}}%
\pgfpathlineto{\pgfqpoint{4.820933in}{5.624194in}}%
\pgfpathlineto{\pgfqpoint{4.827110in}{5.606031in}}%
\pgfpathlineto{\pgfqpoint{4.840334in}{5.600815in}}%
\pgfpathlineto{\pgfqpoint{4.865315in}{5.622180in}}%
\pgfpathlineto{\pgfqpoint{4.884903in}{5.634799in}}%
\pgfpathlineto{\pgfqpoint{4.928062in}{5.612577in}}%
\pgfpathlineto{\pgfqpoint{4.967053in}{5.489195in}}%
\pgfpathlineto{\pgfqpoint{4.974465in}{5.458840in}}%
\pgfpathlineto{\pgfqpoint{4.991544in}{5.455280in}}%
\pgfpathlineto{\pgfqpoint{5.014707in}{5.434663in}}%
\pgfpathlineto{\pgfqpoint{5.013396in}{5.422651in}}%
\pgfpathlineto{\pgfqpoint{5.029175in}{5.408413in}}%
\pgfpathlineto{\pgfqpoint{5.041992in}{5.416582in}}%
\pgfpathlineto{\pgfqpoint{5.068814in}{5.385301in}}%
\pgfpathlineto{\pgfqpoint{5.056826in}{5.360267in}}%
\pgfpathlineto{\pgfqpoint{5.040752in}{5.359772in}}%
\pgfpathlineto{\pgfqpoint{5.027855in}{5.337424in}}%
\pgfpathlineto{\pgfqpoint{5.012232in}{5.334255in}}%
\pgfpathlineto{\pgfqpoint{5.000763in}{5.322355in}}%
\pgfpathlineto{\pgfqpoint{4.966507in}{5.309594in}}%
\pgfpathlineto{\pgfqpoint{4.945773in}{5.289029in}}%
\pgfpathlineto{\pgfqpoint{4.935241in}{5.303790in}}%
\pgfpathlineto{\pgfqpoint{4.925671in}{5.293211in}}%
\pgfpathlineto{\pgfqpoint{4.928292in}{5.250441in}}%
\pgfpathlineto{\pgfqpoint{4.920714in}{5.233500in}}%
\pgfpathlineto{\pgfqpoint{4.904185in}{5.238127in}}%
\pgfpathlineto{\pgfqpoint{4.901518in}{5.220759in}}%
\pgfpathlineto{\pgfqpoint{4.894375in}{5.213617in}}%
\pgfpathlineto{\pgfqpoint{4.880081in}{5.215357in}}%
\pgfpathlineto{\pgfqpoint{4.880982in}{5.199118in}}%
\pgfpathlineto{\pgfqpoint{4.859316in}{5.203730in}}%
\pgfpathlineto{\pgfqpoint{4.847484in}{5.181231in}}%
\pgfpathlineto{\pgfqpoint{4.851952in}{5.169459in}}%
\pgfpathlineto{\pgfqpoint{4.844914in}{5.149879in}}%
\pgfpathlineto{\pgfqpoint{4.833835in}{5.135478in}}%
\pgfpathlineto{\pgfqpoint{4.831046in}{5.105414in}}%
\pgfpathlineto{\pgfqpoint{4.817670in}{5.105444in}}%
\pgfusepath{stroke}%
\end{pgfscope}%
\begin{pgfscope}%
\pgfpathrectangle{\pgfqpoint{0.100000in}{2.493063in}}{\pgfqpoint{5.037500in}{3.427208in}}%
\pgfusepath{clip}%
\pgfsetbuttcap%
\pgfsetroundjoin%
\pgfsetlinewidth{0.501875pt}%
\definecolor{currentstroke}{rgb}{0.827451,0.827451,0.827451}%
\pgfsetstrokecolor{currentstroke}%
\pgfsetdash{}{0pt}%
\pgfpathmoveto{\pgfqpoint{4.972660in}{5.300812in}}%
\pgfpathlineto{\pgfqpoint{4.982442in}{5.311057in}}%
\pgfpathlineto{\pgfqpoint{4.991754in}{5.301427in}}%
\pgfpathlineto{\pgfqpoint{4.975027in}{5.288668in}}%
\pgfpathlineto{\pgfqpoint{4.972660in}{5.300812in}}%
\pgfusepath{stroke}%
\end{pgfscope}%
\begin{pgfscope}%
\pgfpathrectangle{\pgfqpoint{0.100000in}{2.493063in}}{\pgfqpoint{5.037500in}{3.427208in}}%
\pgfusepath{clip}%
\pgfsetbuttcap%
\pgfsetroundjoin%
\pgfsetlinewidth{0.501875pt}%
\definecolor{currentstroke}{rgb}{0.827451,0.827451,0.827451}%
\pgfsetstrokecolor{currentstroke}%
\pgfsetdash{}{0pt}%
\pgfpathmoveto{\pgfqpoint{2.075380in}{5.198668in}}%
\pgfpathlineto{\pgfqpoint{2.083575in}{5.279518in}}%
\pgfpathlineto{\pgfqpoint{2.099200in}{5.436717in}}%
\pgfpathlineto{\pgfqpoint{2.109697in}{5.553044in}}%
\pgfpathlineto{\pgfqpoint{2.195665in}{5.544979in}}%
\pgfpathlineto{\pgfqpoint{2.305646in}{5.536223in}}%
\pgfpathlineto{\pgfqpoint{2.406173in}{5.529747in}}%
\pgfpathlineto{\pgfqpoint{2.497197in}{5.525142in}}%
\pgfpathlineto{\pgfqpoint{2.633036in}{5.520482in}}%
\pgfpathlineto{\pgfqpoint{2.642307in}{5.482197in}}%
\pgfpathlineto{\pgfqpoint{2.638988in}{5.463901in}}%
\pgfpathlineto{\pgfqpoint{2.637896in}{5.426294in}}%
\pgfpathlineto{\pgfqpoint{2.644203in}{5.398130in}}%
\pgfpathlineto{\pgfqpoint{2.658653in}{5.356579in}}%
\pgfpathlineto{\pgfqpoint{2.658619in}{5.304267in}}%
\pgfpathlineto{\pgfqpoint{2.661293in}{5.243497in}}%
\pgfpathlineto{\pgfqpoint{2.664964in}{5.227119in}}%
\pgfpathlineto{\pgfqpoint{2.675704in}{5.209137in}}%
\pgfpathlineto{\pgfqpoint{2.679246in}{5.181121in}}%
\pgfpathlineto{\pgfqpoint{2.677717in}{5.162416in}}%
\pgfpathlineto{\pgfqpoint{2.563838in}{5.164883in}}%
\pgfpathlineto{\pgfqpoint{2.480994in}{5.169082in}}%
\pgfpathlineto{\pgfqpoint{2.359549in}{5.175547in}}%
\pgfpathlineto{\pgfqpoint{2.239785in}{5.184054in}}%
\pgfpathlineto{\pgfqpoint{2.160017in}{5.190489in}}%
\pgfpathlineto{\pgfqpoint{2.075380in}{5.198668in}}%
\pgfusepath{stroke}%
\end{pgfscope}%
\begin{pgfscope}%
\pgfpathrectangle{\pgfqpoint{0.100000in}{2.493063in}}{\pgfqpoint{5.037500in}{3.427208in}}%
\pgfusepath{clip}%
\pgfsetbuttcap%
\pgfsetroundjoin%
\pgfsetlinewidth{0.501875pt}%
\definecolor{currentstroke}{rgb}{0.827451,0.827451,0.827451}%
\pgfsetstrokecolor{currentstroke}%
\pgfsetdash{}{0pt}%
\pgfpathmoveto{\pgfqpoint{2.040958in}{4.860305in}}%
\pgfpathlineto{\pgfqpoint{2.054105in}{4.995834in}}%
\pgfpathlineto{\pgfqpoint{2.063595in}{5.089707in}}%
\pgfpathlineto{\pgfqpoint{2.075380in}{5.198668in}}%
\pgfpathlineto{\pgfqpoint{2.160017in}{5.190489in}}%
\pgfpathlineto{\pgfqpoint{2.239785in}{5.184054in}}%
\pgfpathlineto{\pgfqpoint{2.359549in}{5.175547in}}%
\pgfpathlineto{\pgfqpoint{2.480994in}{5.169082in}}%
\pgfpathlineto{\pgfqpoint{2.563838in}{5.164883in}}%
\pgfpathlineto{\pgfqpoint{2.677717in}{5.162416in}}%
\pgfpathlineto{\pgfqpoint{2.670006in}{5.139917in}}%
\pgfpathlineto{\pgfqpoint{2.654617in}{5.122254in}}%
\pgfpathlineto{\pgfqpoint{2.666402in}{5.101918in}}%
\pgfpathlineto{\pgfqpoint{2.679397in}{5.097573in}}%
\pgfpathlineto{\pgfqpoint{2.685563in}{5.085892in}}%
\pgfpathlineto{\pgfqpoint{2.684148in}{5.000604in}}%
\pgfpathlineto{\pgfqpoint{2.681778in}{4.880574in}}%
\pgfpathlineto{\pgfqpoint{2.673013in}{4.849048in}}%
\pgfpathlineto{\pgfqpoint{2.680853in}{4.831610in}}%
\pgfpathlineto{\pgfqpoint{2.665113in}{4.794137in}}%
\pgfpathlineto{\pgfqpoint{2.681695in}{4.763929in}}%
\pgfpathlineto{\pgfqpoint{2.667607in}{4.766241in}}%
\pgfpathlineto{\pgfqpoint{2.657995in}{4.785035in}}%
\pgfpathlineto{\pgfqpoint{2.623692in}{4.797918in}}%
\pgfpathlineto{\pgfqpoint{2.602058in}{4.809251in}}%
\pgfpathlineto{\pgfqpoint{2.565756in}{4.810193in}}%
\pgfpathlineto{\pgfqpoint{2.553155in}{4.799869in}}%
\pgfpathlineto{\pgfqpoint{2.512066in}{4.820198in}}%
\pgfpathlineto{\pgfqpoint{2.508914in}{4.826627in}}%
\pgfpathlineto{\pgfqpoint{2.365512in}{4.833414in}}%
\pgfpathlineto{\pgfqpoint{2.278390in}{4.838408in}}%
\pgfpathlineto{\pgfqpoint{2.206426in}{4.844043in}}%
\pgfpathlineto{\pgfqpoint{2.087517in}{4.855278in}}%
\pgfpathlineto{\pgfqpoint{2.040958in}{4.860305in}}%
\pgfusepath{stroke}%
\end{pgfscope}%
\begin{pgfscope}%
\pgfpathrectangle{\pgfqpoint{0.100000in}{2.493063in}}{\pgfqpoint{5.037500in}{3.427208in}}%
\pgfusepath{clip}%
\pgfsetbuttcap%
\pgfsetroundjoin%
\pgfsetlinewidth{0.501875pt}%
\definecolor{currentstroke}{rgb}{0.827451,0.827451,0.827451}%
\pgfsetstrokecolor{currentstroke}%
\pgfsetdash{}{0pt}%
\pgfpathmoveto{\pgfqpoint{2.018405in}{4.630818in}}%
\pgfpathlineto{\pgfqpoint{1.941955in}{4.637799in}}%
\pgfpathlineto{\pgfqpoint{1.775315in}{4.658367in}}%
\pgfpathlineto{\pgfqpoint{1.684663in}{4.671377in}}%
\pgfpathlineto{\pgfqpoint{1.587436in}{4.685334in}}%
\pgfpathlineto{\pgfqpoint{1.505537in}{4.698440in}}%
\pgfpathlineto{\pgfqpoint{1.415642in}{4.713831in}}%
\pgfpathlineto{\pgfqpoint{1.436081in}{4.827212in}}%
\pgfpathlineto{\pgfqpoint{1.456784in}{4.943472in}}%
\pgfpathlineto{\pgfqpoint{1.486097in}{5.108827in}}%
\pgfpathlineto{\pgfqpoint{1.496569in}{5.168365in}}%
\pgfpathlineto{\pgfqpoint{1.546810in}{5.159174in}}%
\pgfpathlineto{\pgfqpoint{1.692759in}{5.136264in}}%
\pgfpathlineto{\pgfqpoint{1.751998in}{5.127814in}}%
\pgfpathlineto{\pgfqpoint{1.884635in}{5.109661in}}%
\pgfpathlineto{\pgfqpoint{1.980159in}{5.098817in}}%
\pgfpathlineto{\pgfqpoint{2.063595in}{5.089707in}}%
\pgfpathlineto{\pgfqpoint{2.054105in}{4.995834in}}%
\pgfpathlineto{\pgfqpoint{2.040958in}{4.860305in}}%
\pgfpathlineto{\pgfqpoint{2.029673in}{4.745125in}}%
\pgfpathlineto{\pgfqpoint{2.018405in}{4.630818in}}%
\pgfusepath{stroke}%
\end{pgfscope}%
\begin{pgfscope}%
\pgfpathrectangle{\pgfqpoint{0.100000in}{2.493063in}}{\pgfqpoint{5.037500in}{3.427208in}}%
\pgfusepath{clip}%
\pgfsetbuttcap%
\pgfsetroundjoin%
\pgfsetlinewidth{0.501875pt}%
\definecolor{currentstroke}{rgb}{0.827451,0.827451,0.827451}%
\pgfsetstrokecolor{currentstroke}%
\pgfsetdash{}{0pt}%
\pgfpathmoveto{\pgfqpoint{3.416161in}{4.787215in}}%
\pgfpathlineto{\pgfqpoint{3.319282in}{4.780333in}}%
\pgfpathlineto{\pgfqpoint{3.174878in}{4.774179in}}%
\pgfpathlineto{\pgfqpoint{3.169391in}{4.788766in}}%
\pgfpathlineto{\pgfqpoint{3.137369in}{4.799670in}}%
\pgfpathlineto{\pgfqpoint{3.130306in}{4.820271in}}%
\pgfpathlineto{\pgfqpoint{3.127348in}{4.845756in}}%
\pgfpathlineto{\pgfqpoint{3.134564in}{4.858814in}}%
\pgfpathlineto{\pgfqpoint{3.123167in}{4.871344in}}%
\pgfpathlineto{\pgfqpoint{3.120418in}{4.886289in}}%
\pgfpathlineto{\pgfqpoint{3.116744in}{4.919351in}}%
\pgfpathlineto{\pgfqpoint{3.105822in}{4.937310in}}%
\pgfpathlineto{\pgfqpoint{3.086424in}{4.947390in}}%
\pgfpathlineto{\pgfqpoint{3.065315in}{4.964023in}}%
\pgfpathlineto{\pgfqpoint{3.054403in}{4.983703in}}%
\pgfpathlineto{\pgfqpoint{3.034824in}{4.991628in}}%
\pgfpathlineto{\pgfqpoint{3.023314in}{5.004526in}}%
\pgfpathlineto{\pgfqpoint{3.009371in}{5.006715in}}%
\pgfpathlineto{\pgfqpoint{2.984477in}{5.025866in}}%
\pgfpathlineto{\pgfqpoint{2.988521in}{5.047905in}}%
\pgfpathlineto{\pgfqpoint{2.987741in}{5.089815in}}%
\pgfpathlineto{\pgfqpoint{2.995419in}{5.101352in}}%
\pgfpathlineto{\pgfqpoint{2.988521in}{5.118771in}}%
\pgfpathlineto{\pgfqpoint{2.976370in}{5.122141in}}%
\pgfpathlineto{\pgfqpoint{2.977374in}{5.137440in}}%
\pgfpathlineto{\pgfqpoint{2.992450in}{5.161609in}}%
\pgfpathlineto{\pgfqpoint{3.022305in}{5.180721in}}%
\pgfpathlineto{\pgfqpoint{3.020445in}{5.248690in}}%
\pgfpathlineto{\pgfqpoint{3.035383in}{5.258900in}}%
\pgfpathlineto{\pgfqpoint{3.049520in}{5.252098in}}%
\pgfpathlineto{\pgfqpoint{3.078318in}{5.262055in}}%
\pgfpathlineto{\pgfqpoint{3.132500in}{5.287078in}}%
\pgfpathlineto{\pgfqpoint{3.139558in}{5.279330in}}%
\pgfpathlineto{\pgfqpoint{3.129296in}{5.244236in}}%
\pgfpathlineto{\pgfqpoint{3.144563in}{5.251938in}}%
\pgfpathlineto{\pgfqpoint{3.170712in}{5.244247in}}%
\pgfpathlineto{\pgfqpoint{3.186780in}{5.237828in}}%
\pgfpathlineto{\pgfqpoint{3.195789in}{5.219018in}}%
\pgfpathlineto{\pgfqpoint{3.278223in}{5.201266in}}%
\pgfpathlineto{\pgfqpoint{3.302863in}{5.189087in}}%
\pgfpathlineto{\pgfqpoint{3.327916in}{5.189231in}}%
\pgfpathlineto{\pgfqpoint{3.353658in}{5.184206in}}%
\pgfpathlineto{\pgfqpoint{3.370364in}{5.167046in}}%
\pgfpathlineto{\pgfqpoint{3.386328in}{5.158566in}}%
\pgfpathlineto{\pgfqpoint{3.389261in}{5.134109in}}%
\pgfpathlineto{\pgfqpoint{3.384545in}{5.118752in}}%
\pgfpathlineto{\pgfqpoint{3.402337in}{5.117617in}}%
\pgfpathlineto{\pgfqpoint{3.396343in}{5.099806in}}%
\pgfpathlineto{\pgfqpoint{3.402044in}{5.093479in}}%
\pgfpathlineto{\pgfqpoint{3.407714in}{5.076666in}}%
\pgfpathlineto{\pgfqpoint{3.390375in}{5.067764in}}%
\pgfpathlineto{\pgfqpoint{3.380307in}{5.042963in}}%
\pgfpathlineto{\pgfqpoint{3.377148in}{5.025425in}}%
\pgfpathlineto{\pgfqpoint{3.386807in}{5.022424in}}%
\pgfpathlineto{\pgfqpoint{3.399173in}{5.035574in}}%
\pgfpathlineto{\pgfqpoint{3.409687in}{5.058365in}}%
\pgfpathlineto{\pgfqpoint{3.423912in}{5.066280in}}%
\pgfpathlineto{\pgfqpoint{3.434602in}{5.055984in}}%
\pgfpathlineto{\pgfqpoint{3.424080in}{5.024821in}}%
\pgfpathlineto{\pgfqpoint{3.420778in}{5.000629in}}%
\pgfpathlineto{\pgfqpoint{3.423887in}{4.983239in}}%
\pgfpathlineto{\pgfqpoint{3.414116in}{4.973401in}}%
\pgfpathlineto{\pgfqpoint{3.409236in}{4.949361in}}%
\pgfpathlineto{\pgfqpoint{3.413223in}{4.924095in}}%
\pgfpathlineto{\pgfqpoint{3.401692in}{4.886729in}}%
\pgfpathlineto{\pgfqpoint{3.401934in}{4.868058in}}%
\pgfpathlineto{\pgfqpoint{3.411046in}{4.827597in}}%
\pgfpathlineto{\pgfqpoint{3.416952in}{4.820656in}}%
\pgfpathlineto{\pgfqpoint{3.416161in}{4.787215in}}%
\pgfusepath{stroke}%
\end{pgfscope}%
\begin{pgfscope}%
\pgfpathrectangle{\pgfqpoint{0.100000in}{2.493063in}}{\pgfqpoint{5.037500in}{3.427208in}}%
\pgfusepath{clip}%
\pgfsetbuttcap%
\pgfsetroundjoin%
\pgfsetlinewidth{0.501875pt}%
\definecolor{currentstroke}{rgb}{0.827451,0.827451,0.827451}%
\pgfsetstrokecolor{currentstroke}%
\pgfsetdash{}{0pt}%
\pgfpathmoveto{\pgfqpoint{3.452484in}{5.115138in}}%
\pgfpathlineto{\pgfqpoint{3.454646in}{5.098999in}}%
\pgfpathlineto{\pgfqpoint{3.434836in}{5.056475in}}%
\pgfpathlineto{\pgfqpoint{3.426032in}{5.068793in}}%
\pgfpathlineto{\pgfqpoint{3.452484in}{5.115138in}}%
\pgfusepath{stroke}%
\end{pgfscope}%
\begin{pgfscope}%
\pgfpathrectangle{\pgfqpoint{0.100000in}{2.493063in}}{\pgfqpoint{5.037500in}{3.427208in}}%
\pgfusepath{clip}%
\pgfsetbuttcap%
\pgfsetroundjoin%
\pgfsetlinewidth{0.501875pt}%
\definecolor{currentstroke}{rgb}{0.827451,0.827451,0.827451}%
\pgfsetstrokecolor{currentstroke}%
\pgfsetdash{}{0pt}%
\pgfpathmoveto{\pgfqpoint{1.054504in}{5.380440in}}%
\pgfpathlineto{\pgfqpoint{1.058421in}{5.399829in}}%
\pgfpathlineto{\pgfqpoint{1.055466in}{5.420878in}}%
\pgfpathlineto{\pgfqpoint{1.059991in}{5.443826in}}%
\pgfpathlineto{\pgfqpoint{1.083039in}{5.536934in}}%
\pgfpathlineto{\pgfqpoint{1.102423in}{5.614332in}}%
\pgfpathlineto{\pgfqpoint{1.129685in}{5.722352in}}%
\pgfpathlineto{\pgfqpoint{1.201939in}{5.704893in}}%
\pgfpathlineto{\pgfqpoint{1.175235in}{5.588340in}}%
\pgfpathlineto{\pgfqpoint{1.192146in}{5.550785in}}%
\pgfpathlineto{\pgfqpoint{1.190007in}{5.530394in}}%
\pgfpathlineto{\pgfqpoint{1.199072in}{5.499718in}}%
\pgfpathlineto{\pgfqpoint{1.211732in}{5.492469in}}%
\pgfpathlineto{\pgfqpoint{1.225585in}{5.456924in}}%
\pgfpathlineto{\pgfqpoint{1.234095in}{5.447134in}}%
\pgfpathlineto{\pgfqpoint{1.233064in}{5.434532in}}%
\pgfpathlineto{\pgfqpoint{1.249197in}{5.423259in}}%
\pgfpathlineto{\pgfqpoint{1.251474in}{5.410187in}}%
\pgfpathlineto{\pgfqpoint{1.273621in}{5.408361in}}%
\pgfpathlineto{\pgfqpoint{1.245534in}{5.335516in}}%
\pgfpathlineto{\pgfqpoint{1.249829in}{5.320466in}}%
\pgfpathlineto{\pgfqpoint{1.233671in}{5.311107in}}%
\pgfpathlineto{\pgfqpoint{1.231480in}{5.287352in}}%
\pgfpathlineto{\pgfqpoint{1.244953in}{5.272477in}}%
\pgfpathlineto{\pgfqpoint{1.262653in}{5.279694in}}%
\pgfpathlineto{\pgfqpoint{1.275657in}{5.293038in}}%
\pgfpathlineto{\pgfqpoint{1.289516in}{5.277955in}}%
\pgfpathlineto{\pgfqpoint{1.288430in}{5.247138in}}%
\pgfpathlineto{\pgfqpoint{1.304555in}{5.210885in}}%
\pgfpathlineto{\pgfqpoint{1.300862in}{5.188938in}}%
\pgfpathlineto{\pgfqpoint{1.323509in}{5.173709in}}%
\pgfpathlineto{\pgfqpoint{1.323870in}{5.152534in}}%
\pgfpathlineto{\pgfqpoint{1.331583in}{5.130280in}}%
\pgfpathlineto{\pgfqpoint{1.345921in}{5.135295in}}%
\pgfpathlineto{\pgfqpoint{1.378775in}{5.125956in}}%
\pgfpathlineto{\pgfqpoint{1.392301in}{5.136334in}}%
\pgfpathlineto{\pgfqpoint{1.401974in}{5.129626in}}%
\pgfpathlineto{\pgfqpoint{1.421216in}{5.130465in}}%
\pgfpathlineto{\pgfqpoint{1.425343in}{5.123537in}}%
\pgfpathlineto{\pgfqpoint{1.445885in}{5.125474in}}%
\pgfpathlineto{\pgfqpoint{1.455784in}{5.139953in}}%
\pgfpathlineto{\pgfqpoint{1.469613in}{5.141717in}}%
\pgfpathlineto{\pgfqpoint{1.480372in}{5.113075in}}%
\pgfpathlineto{\pgfqpoint{1.486097in}{5.108827in}}%
\pgfpathlineto{\pgfqpoint{1.456784in}{4.943472in}}%
\pgfpathlineto{\pgfqpoint{1.436081in}{4.827212in}}%
\pgfpathlineto{\pgfqpoint{1.272793in}{4.858734in}}%
\pgfpathlineto{\pgfqpoint{1.184593in}{4.876279in}}%
\pgfpathlineto{\pgfqpoint{1.102093in}{4.894420in}}%
\pgfpathlineto{\pgfqpoint{1.024981in}{4.911878in}}%
\pgfpathlineto{\pgfqpoint{0.935738in}{4.933400in}}%
\pgfpathlineto{\pgfqpoint{0.981770in}{5.122238in}}%
\pgfpathlineto{\pgfqpoint{0.984252in}{5.136039in}}%
\pgfpathlineto{\pgfqpoint{1.004722in}{5.172200in}}%
\pgfpathlineto{\pgfqpoint{0.983526in}{5.193934in}}%
\pgfpathlineto{\pgfqpoint{0.987927in}{5.215277in}}%
\pgfpathlineto{\pgfqpoint{0.995996in}{5.220619in}}%
\pgfpathlineto{\pgfqpoint{1.010406in}{5.242611in}}%
\pgfpathlineto{\pgfqpoint{1.031412in}{5.257827in}}%
\pgfpathlineto{\pgfqpoint{1.032559in}{5.269084in}}%
\pgfpathlineto{\pgfqpoint{1.045237in}{5.280333in}}%
\pgfpathlineto{\pgfqpoint{1.055763in}{5.301387in}}%
\pgfpathlineto{\pgfqpoint{1.077356in}{5.323662in}}%
\pgfpathlineto{\pgfqpoint{1.075823in}{5.345662in}}%
\pgfpathlineto{\pgfqpoint{1.061060in}{5.357883in}}%
\pgfpathlineto{\pgfqpoint{1.054504in}{5.380440in}}%
\pgfusepath{stroke}%
\end{pgfscope}%
\begin{pgfscope}%
\pgfpathrectangle{\pgfqpoint{0.100000in}{2.493063in}}{\pgfqpoint{5.037500in}{3.427208in}}%
\pgfusepath{clip}%
\pgfsetbuttcap%
\pgfsetroundjoin%
\pgfsetlinewidth{0.501875pt}%
\definecolor{currentstroke}{rgb}{0.827451,0.827451,0.827451}%
\pgfsetstrokecolor{currentstroke}%
\pgfsetdash{}{0pt}%
\pgfpathmoveto{\pgfqpoint{4.628939in}{5.011925in}}%
\pgfpathlineto{\pgfqpoint{4.622154in}{5.033342in}}%
\pgfpathlineto{\pgfqpoint{4.609564in}{5.098312in}}%
\pgfpathlineto{\pgfqpoint{4.593285in}{5.118300in}}%
\pgfpathlineto{\pgfqpoint{4.579007in}{5.154245in}}%
\pgfpathlineto{\pgfqpoint{4.583695in}{5.180896in}}%
\pgfpathlineto{\pgfqpoint{4.579935in}{5.200538in}}%
\pgfpathlineto{\pgfqpoint{4.567736in}{5.219830in}}%
\pgfpathlineto{\pgfqpoint{4.567234in}{5.241082in}}%
\pgfpathlineto{\pgfqpoint{4.560150in}{5.264002in}}%
\pgfpathlineto{\pgfqpoint{4.623544in}{5.279604in}}%
\pgfpathlineto{\pgfqpoint{4.705886in}{5.301830in}}%
\pgfpathlineto{\pgfqpoint{4.709172in}{5.289085in}}%
\pgfpathlineto{\pgfqpoint{4.703929in}{5.268795in}}%
\pgfpathlineto{\pgfqpoint{4.716237in}{5.252560in}}%
\pgfpathlineto{\pgfqpoint{4.709700in}{5.231988in}}%
\pgfpathlineto{\pgfqpoint{4.683724in}{5.206129in}}%
\pgfpathlineto{\pgfqpoint{4.690864in}{5.186703in}}%
\pgfpathlineto{\pgfqpoint{4.685572in}{5.145644in}}%
\pgfpathlineto{\pgfqpoint{4.678164in}{5.122494in}}%
\pgfpathlineto{\pgfqpoint{4.687616in}{5.056662in}}%
\pgfpathlineto{\pgfqpoint{4.686342in}{5.033500in}}%
\pgfpathlineto{\pgfqpoint{4.695486in}{5.026066in}}%
\pgfpathlineto{\pgfqpoint{4.628939in}{5.011925in}}%
\pgfusepath{stroke}%
\end{pgfscope}%
\begin{pgfscope}%
\pgfpathrectangle{\pgfqpoint{0.100000in}{2.493063in}}{\pgfqpoint{5.037500in}{3.427208in}}%
\pgfusepath{clip}%
\pgfsetbuttcap%
\pgfsetroundjoin%
\pgfsetlinewidth{0.501875pt}%
\definecolor{currentstroke}{rgb}{0.827451,0.827451,0.827451}%
\pgfsetstrokecolor{currentstroke}%
\pgfsetdash{}{0pt}%
\pgfpathmoveto{\pgfqpoint{2.681778in}{4.880574in}}%
\pgfpathlineto{\pgfqpoint{2.684148in}{5.000604in}}%
\pgfpathlineto{\pgfqpoint{2.685563in}{5.085892in}}%
\pgfpathlineto{\pgfqpoint{2.679397in}{5.097573in}}%
\pgfpathlineto{\pgfqpoint{2.666402in}{5.101918in}}%
\pgfpathlineto{\pgfqpoint{2.654617in}{5.122254in}}%
\pgfpathlineto{\pgfqpoint{2.670006in}{5.139917in}}%
\pgfpathlineto{\pgfqpoint{2.677717in}{5.162416in}}%
\pgfpathlineto{\pgfqpoint{2.679246in}{5.181121in}}%
\pgfpathlineto{\pgfqpoint{2.675704in}{5.209137in}}%
\pgfpathlineto{\pgfqpoint{2.664964in}{5.227119in}}%
\pgfpathlineto{\pgfqpoint{2.661293in}{5.243497in}}%
\pgfpathlineto{\pgfqpoint{2.658619in}{5.304267in}}%
\pgfpathlineto{\pgfqpoint{2.658653in}{5.356579in}}%
\pgfpathlineto{\pgfqpoint{2.644203in}{5.398130in}}%
\pgfpathlineto{\pgfqpoint{2.637896in}{5.426294in}}%
\pgfpathlineto{\pgfqpoint{2.638988in}{5.463901in}}%
\pgfpathlineto{\pgfqpoint{2.642307in}{5.482197in}}%
\pgfpathlineto{\pgfqpoint{2.633036in}{5.520482in}}%
\pgfpathlineto{\pgfqpoint{2.696152in}{5.519219in}}%
\pgfpathlineto{\pgfqpoint{2.792021in}{5.518395in}}%
\pgfpathlineto{\pgfqpoint{2.792545in}{5.561909in}}%
\pgfpathlineto{\pgfqpoint{2.816948in}{5.557118in}}%
\pgfpathlineto{\pgfqpoint{2.828643in}{5.504060in}}%
\pgfpathlineto{\pgfqpoint{2.837265in}{5.484983in}}%
\pgfpathlineto{\pgfqpoint{2.858707in}{5.484420in}}%
\pgfpathlineto{\pgfqpoint{2.863506in}{5.477946in}}%
\pgfpathlineto{\pgfqpoint{2.893414in}{5.475079in}}%
\pgfpathlineto{\pgfqpoint{2.898441in}{5.461925in}}%
\pgfpathlineto{\pgfqpoint{2.919053in}{5.464882in}}%
\pgfpathlineto{\pgfqpoint{2.935062in}{5.477185in}}%
\pgfpathlineto{\pgfqpoint{2.962670in}{5.476725in}}%
\pgfpathlineto{\pgfqpoint{2.979752in}{5.466831in}}%
\pgfpathlineto{\pgfqpoint{2.981716in}{5.457552in}}%
\pgfpathlineto{\pgfqpoint{2.997970in}{5.455609in}}%
\pgfpathlineto{\pgfqpoint{3.008578in}{5.430299in}}%
\pgfpathlineto{\pgfqpoint{3.015422in}{5.445861in}}%
\pgfpathlineto{\pgfqpoint{3.034094in}{5.446818in}}%
\pgfpathlineto{\pgfqpoint{3.038817in}{5.434698in}}%
\pgfpathlineto{\pgfqpoint{3.059826in}{5.429166in}}%
\pgfpathlineto{\pgfqpoint{3.071412in}{5.411709in}}%
\pgfpathlineto{\pgfqpoint{3.097108in}{5.417129in}}%
\pgfpathlineto{\pgfqpoint{3.125382in}{5.438552in}}%
\pgfpathlineto{\pgfqpoint{3.135691in}{5.419667in}}%
\pgfpathlineto{\pgfqpoint{3.182085in}{5.424836in}}%
\pgfpathlineto{\pgfqpoint{3.201917in}{5.411874in}}%
\pgfpathlineto{\pgfqpoint{3.213460in}{5.416502in}}%
\pgfpathlineto{\pgfqpoint{3.222716in}{5.409205in}}%
\pgfpathlineto{\pgfqpoint{3.195257in}{5.391893in}}%
\pgfpathlineto{\pgfqpoint{3.156062in}{5.376468in}}%
\pgfpathlineto{\pgfqpoint{3.117242in}{5.345636in}}%
\pgfpathlineto{\pgfqpoint{3.083587in}{5.305090in}}%
\pgfpathlineto{\pgfqpoint{3.058111in}{5.281119in}}%
\pgfpathlineto{\pgfqpoint{3.035795in}{5.264642in}}%
\pgfpathlineto{\pgfqpoint{3.020445in}{5.248690in}}%
\pgfpathlineto{\pgfqpoint{3.022305in}{5.180721in}}%
\pgfpathlineto{\pgfqpoint{2.992450in}{5.161609in}}%
\pgfpathlineto{\pgfqpoint{2.977374in}{5.137440in}}%
\pgfpathlineto{\pgfqpoint{2.976370in}{5.122141in}}%
\pgfpathlineto{\pgfqpoint{2.988521in}{5.118771in}}%
\pgfpathlineto{\pgfqpoint{2.995419in}{5.101352in}}%
\pgfpathlineto{\pgfqpoint{2.987741in}{5.089815in}}%
\pgfpathlineto{\pgfqpoint{2.988521in}{5.047905in}}%
\pgfpathlineto{\pgfqpoint{2.984477in}{5.025866in}}%
\pgfpathlineto{\pgfqpoint{3.009371in}{5.006715in}}%
\pgfpathlineto{\pgfqpoint{3.023314in}{5.004526in}}%
\pgfpathlineto{\pgfqpoint{3.034824in}{4.991628in}}%
\pgfpathlineto{\pgfqpoint{3.054403in}{4.983703in}}%
\pgfpathlineto{\pgfqpoint{3.065315in}{4.964023in}}%
\pgfpathlineto{\pgfqpoint{3.086424in}{4.947390in}}%
\pgfpathlineto{\pgfqpoint{3.105822in}{4.937310in}}%
\pgfpathlineto{\pgfqpoint{3.116744in}{4.919351in}}%
\pgfpathlineto{\pgfqpoint{3.120418in}{4.886289in}}%
\pgfpathlineto{\pgfqpoint{3.017472in}{4.882551in}}%
\pgfpathlineto{\pgfqpoint{2.929720in}{4.880716in}}%
\pgfpathlineto{\pgfqpoint{2.849769in}{4.879541in}}%
\pgfpathlineto{\pgfqpoint{2.765190in}{4.879670in}}%
\pgfpathlineto{\pgfqpoint{2.681778in}{4.880574in}}%
\pgfusepath{stroke}%
\end{pgfscope}%
\begin{pgfscope}%
\pgfpathrectangle{\pgfqpoint{0.100000in}{2.493063in}}{\pgfqpoint{5.037500in}{3.427208in}}%
\pgfusepath{clip}%
\pgfsetbuttcap%
\pgfsetroundjoin%
\pgfsetlinewidth{0.501875pt}%
\definecolor{currentstroke}{rgb}{0.827451,0.827451,0.827451}%
\pgfsetstrokecolor{currentstroke}%
\pgfsetdash{}{0pt}%
\pgfpathmoveto{\pgfqpoint{0.344393in}{5.104965in}}%
\pgfpathlineto{\pgfqpoint{0.336240in}{5.119967in}}%
\pgfpathlineto{\pgfqpoint{0.341446in}{5.158475in}}%
\pgfpathlineto{\pgfqpoint{0.351493in}{5.178599in}}%
\pgfpathlineto{\pgfqpoint{0.346470in}{5.205823in}}%
\pgfpathlineto{\pgfqpoint{0.356952in}{5.217270in}}%
\pgfpathlineto{\pgfqpoint{0.376075in}{5.248128in}}%
\pgfpathlineto{\pgfqpoint{0.392289in}{5.266755in}}%
\pgfpathlineto{\pgfqpoint{0.415990in}{5.307325in}}%
\pgfpathlineto{\pgfqpoint{0.434168in}{5.351519in}}%
\pgfpathlineto{\pgfqpoint{0.453465in}{5.392967in}}%
\pgfpathlineto{\pgfqpoint{0.457366in}{5.410279in}}%
\pgfpathlineto{\pgfqpoint{0.483987in}{5.459801in}}%
\pgfpathlineto{\pgfqpoint{0.489113in}{5.481551in}}%
\pgfpathlineto{\pgfqpoint{0.500447in}{5.504392in}}%
\pgfpathlineto{\pgfqpoint{0.504499in}{5.532254in}}%
\pgfpathlineto{\pgfqpoint{0.526137in}{5.545840in}}%
\pgfpathlineto{\pgfqpoint{0.551794in}{5.552687in}}%
\pgfpathlineto{\pgfqpoint{0.564765in}{5.537351in}}%
\pgfpathlineto{\pgfqpoint{0.575981in}{5.538485in}}%
\pgfpathlineto{\pgfqpoint{0.593490in}{5.520205in}}%
\pgfpathlineto{\pgfqpoint{0.595818in}{5.502986in}}%
\pgfpathlineto{\pgfqpoint{0.589177in}{5.469269in}}%
\pgfpathlineto{\pgfqpoint{0.611269in}{5.452066in}}%
\pgfpathlineto{\pgfqpoint{0.625598in}{5.445598in}}%
\pgfpathlineto{\pgfqpoint{0.664413in}{5.452200in}}%
\pgfpathlineto{\pgfqpoint{0.686901in}{5.447691in}}%
\pgfpathlineto{\pgfqpoint{0.702357in}{5.440292in}}%
\pgfpathlineto{\pgfqpoint{0.710499in}{5.426247in}}%
\pgfpathlineto{\pgfqpoint{0.762052in}{5.428120in}}%
\pgfpathlineto{\pgfqpoint{0.770250in}{5.419526in}}%
\pgfpathlineto{\pgfqpoint{0.789590in}{5.417671in}}%
\pgfpathlineto{\pgfqpoint{0.809066in}{5.423162in}}%
\pgfpathlineto{\pgfqpoint{0.866649in}{5.421778in}}%
\pgfpathlineto{\pgfqpoint{0.879086in}{5.417547in}}%
\pgfpathlineto{\pgfqpoint{0.893620in}{5.422423in}}%
\pgfpathlineto{\pgfqpoint{1.054504in}{5.380440in}}%
\pgfpathlineto{\pgfqpoint{1.061060in}{5.357883in}}%
\pgfpathlineto{\pgfqpoint{1.075823in}{5.345662in}}%
\pgfpathlineto{\pgfqpoint{1.077356in}{5.323662in}}%
\pgfpathlineto{\pgfqpoint{1.055763in}{5.301387in}}%
\pgfpathlineto{\pgfqpoint{1.045237in}{5.280333in}}%
\pgfpathlineto{\pgfqpoint{1.032559in}{5.269084in}}%
\pgfpathlineto{\pgfqpoint{1.031412in}{5.257827in}}%
\pgfpathlineto{\pgfqpoint{1.010406in}{5.242611in}}%
\pgfpathlineto{\pgfqpoint{0.995996in}{5.220619in}}%
\pgfpathlineto{\pgfqpoint{0.987927in}{5.215277in}}%
\pgfpathlineto{\pgfqpoint{0.983526in}{5.193934in}}%
\pgfpathlineto{\pgfqpoint{1.004722in}{5.172200in}}%
\pgfpathlineto{\pgfqpoint{0.984252in}{5.136039in}}%
\pgfpathlineto{\pgfqpoint{0.981770in}{5.122238in}}%
\pgfpathlineto{\pgfqpoint{0.935738in}{4.933400in}}%
\pgfpathlineto{\pgfqpoint{0.838915in}{4.958210in}}%
\pgfpathlineto{\pgfqpoint{0.745509in}{4.982295in}}%
\pgfpathlineto{\pgfqpoint{0.689161in}{4.997975in}}%
\pgfpathlineto{\pgfqpoint{0.616760in}{5.018599in}}%
\pgfpathlineto{\pgfqpoint{0.501273in}{5.054949in}}%
\pgfpathlineto{\pgfqpoint{0.375733in}{5.094034in}}%
\pgfpathlineto{\pgfqpoint{0.344393in}{5.104965in}}%
\pgfusepath{stroke}%
\end{pgfscope}%
\begin{pgfscope}%
\pgfpathrectangle{\pgfqpoint{0.100000in}{2.493063in}}{\pgfqpoint{5.037500in}{3.427208in}}%
\pgfusepath{clip}%
\pgfsetbuttcap%
\pgfsetroundjoin%
\pgfsetlinewidth{0.501875pt}%
\definecolor{currentstroke}{rgb}{0.827451,0.827451,0.827451}%
\pgfsetstrokecolor{currentstroke}%
\pgfsetdash{}{0pt}%
\pgfpathmoveto{\pgfqpoint{4.695486in}{5.026066in}}%
\pgfpathlineto{\pgfqpoint{4.686342in}{5.033500in}}%
\pgfpathlineto{\pgfqpoint{4.687616in}{5.056662in}}%
\pgfpathlineto{\pgfqpoint{4.678164in}{5.122494in}}%
\pgfpathlineto{\pgfqpoint{4.685572in}{5.145644in}}%
\pgfpathlineto{\pgfqpoint{4.690864in}{5.186703in}}%
\pgfpathlineto{\pgfqpoint{4.683724in}{5.206129in}}%
\pgfpathlineto{\pgfqpoint{4.709700in}{5.231988in}}%
\pgfpathlineto{\pgfqpoint{4.716237in}{5.252560in}}%
\pgfpathlineto{\pgfqpoint{4.703929in}{5.268795in}}%
\pgfpathlineto{\pgfqpoint{4.709172in}{5.289085in}}%
\pgfpathlineto{\pgfqpoint{4.705886in}{5.301830in}}%
\pgfpathlineto{\pgfqpoint{4.708708in}{5.329125in}}%
\pgfpathlineto{\pgfqpoint{4.713985in}{5.337552in}}%
\pgfpathlineto{\pgfqpoint{4.730253in}{5.342870in}}%
\pgfpathlineto{\pgfqpoint{4.753984in}{5.273372in}}%
\pgfpathlineto{\pgfqpoint{4.798011in}{5.129357in}}%
\pgfpathlineto{\pgfqpoint{4.814672in}{5.118191in}}%
\pgfpathlineto{\pgfqpoint{4.817670in}{5.105444in}}%
\pgfpathlineto{\pgfqpoint{4.826463in}{5.100294in}}%
\pgfpathlineto{\pgfqpoint{4.825794in}{5.077176in}}%
\pgfpathlineto{\pgfqpoint{4.816477in}{5.076805in}}%
\pgfpathlineto{\pgfqpoint{4.797601in}{5.062409in}}%
\pgfpathlineto{\pgfqpoint{4.792157in}{5.047966in}}%
\pgfpathlineto{\pgfqpoint{4.741634in}{5.035466in}}%
\pgfpathlineto{\pgfqpoint{4.695486in}{5.026066in}}%
\pgfusepath{stroke}%
\end{pgfscope}%
\begin{pgfscope}%
\pgfpathrectangle{\pgfqpoint{0.100000in}{2.493063in}}{\pgfqpoint{5.037500in}{3.427208in}}%
\pgfusepath{clip}%
\pgfsetbuttcap%
\pgfsetroundjoin%
\pgfsetlinewidth{0.501875pt}%
\definecolor{currentstroke}{rgb}{0.827451,0.827451,0.827451}%
\pgfsetstrokecolor{currentstroke}%
\pgfsetdash{}{0pt}%
\pgfpathmoveto{\pgfqpoint{3.115663in}{4.525025in}}%
\pgfpathlineto{\pgfqpoint{3.088992in}{4.551439in}}%
\pgfpathlineto{\pgfqpoint{3.003734in}{4.546521in}}%
\pgfpathlineto{\pgfqpoint{2.870765in}{4.542140in}}%
\pgfpathlineto{\pgfqpoint{2.736996in}{4.544227in}}%
\pgfpathlineto{\pgfqpoint{2.727607in}{4.560599in}}%
\pgfpathlineto{\pgfqpoint{2.731440in}{4.576676in}}%
\pgfpathlineto{\pgfqpoint{2.729620in}{4.604217in}}%
\pgfpathlineto{\pgfqpoint{2.722550in}{4.630893in}}%
\pgfpathlineto{\pgfqpoint{2.723372in}{4.644980in}}%
\pgfpathlineto{\pgfqpoint{2.708345in}{4.660897in}}%
\pgfpathlineto{\pgfqpoint{2.711614in}{4.683044in}}%
\pgfpathlineto{\pgfqpoint{2.688563in}{4.726790in}}%
\pgfpathlineto{\pgfqpoint{2.681695in}{4.763929in}}%
\pgfpathlineto{\pgfqpoint{2.665113in}{4.794137in}}%
\pgfpathlineto{\pgfqpoint{2.680853in}{4.831610in}}%
\pgfpathlineto{\pgfqpoint{2.673013in}{4.849048in}}%
\pgfpathlineto{\pgfqpoint{2.681778in}{4.880574in}}%
\pgfpathlineto{\pgfqpoint{2.765190in}{4.879670in}}%
\pgfpathlineto{\pgfqpoint{2.849769in}{4.879541in}}%
\pgfpathlineto{\pgfqpoint{2.929720in}{4.880716in}}%
\pgfpathlineto{\pgfqpoint{3.017472in}{4.882551in}}%
\pgfpathlineto{\pgfqpoint{3.120418in}{4.886289in}}%
\pgfpathlineto{\pgfqpoint{3.123167in}{4.871344in}}%
\pgfpathlineto{\pgfqpoint{3.134564in}{4.858814in}}%
\pgfpathlineto{\pgfqpoint{3.127348in}{4.845756in}}%
\pgfpathlineto{\pgfqpoint{3.130306in}{4.820271in}}%
\pgfpathlineto{\pgfqpoint{3.137369in}{4.799670in}}%
\pgfpathlineto{\pgfqpoint{3.169391in}{4.788766in}}%
\pgfpathlineto{\pgfqpoint{3.174878in}{4.774179in}}%
\pgfpathlineto{\pgfqpoint{3.192451in}{4.757802in}}%
\pgfpathlineto{\pgfqpoint{3.199619in}{4.740860in}}%
\pgfpathlineto{\pgfqpoint{3.217424in}{4.729495in}}%
\pgfpathlineto{\pgfqpoint{3.220210in}{4.715813in}}%
\pgfpathlineto{\pgfqpoint{3.216743in}{4.695099in}}%
\pgfpathlineto{\pgfqpoint{3.207681in}{4.688891in}}%
\pgfpathlineto{\pgfqpoint{3.204948in}{4.669174in}}%
\pgfpathlineto{\pgfqpoint{3.178904in}{4.653514in}}%
\pgfpathlineto{\pgfqpoint{3.144957in}{4.644935in}}%
\pgfpathlineto{\pgfqpoint{3.141846in}{4.625209in}}%
\pgfpathlineto{\pgfqpoint{3.154942in}{4.611116in}}%
\pgfpathlineto{\pgfqpoint{3.155476in}{4.593396in}}%
\pgfpathlineto{\pgfqpoint{3.144911in}{4.579468in}}%
\pgfpathlineto{\pgfqpoint{3.139365in}{4.558775in}}%
\pgfpathlineto{\pgfqpoint{3.120999in}{4.551925in}}%
\pgfpathlineto{\pgfqpoint{3.122170in}{4.528866in}}%
\pgfpathlineto{\pgfqpoint{3.115663in}{4.525025in}}%
\pgfusepath{stroke}%
\end{pgfscope}%
\begin{pgfscope}%
\pgfpathrectangle{\pgfqpoint{0.100000in}{2.493063in}}{\pgfqpoint{5.037500in}{3.427208in}}%
\pgfusepath{clip}%
\pgfsetbuttcap%
\pgfsetroundjoin%
\pgfsetlinewidth{0.501875pt}%
\definecolor{currentstroke}{rgb}{0.827451,0.827451,0.827451}%
\pgfsetstrokecolor{currentstroke}%
\pgfsetdash{}{0pt}%
\pgfpathmoveto{\pgfqpoint{4.770516in}{4.959426in}}%
\pgfpathlineto{\pgfqpoint{4.745254in}{4.955440in}}%
\pgfpathlineto{\pgfqpoint{4.629215in}{4.929072in}}%
\pgfpathlineto{\pgfqpoint{4.627187in}{4.932145in}}%
\pgfpathlineto{\pgfqpoint{4.628939in}{5.011925in}}%
\pgfpathlineto{\pgfqpoint{4.695486in}{5.026066in}}%
\pgfpathlineto{\pgfqpoint{4.741634in}{5.035466in}}%
\pgfpathlineto{\pgfqpoint{4.792157in}{5.047966in}}%
\pgfpathlineto{\pgfqpoint{4.797601in}{5.062409in}}%
\pgfpathlineto{\pgfqpoint{4.816477in}{5.076805in}}%
\pgfpathlineto{\pgfqpoint{4.825794in}{5.077176in}}%
\pgfpathlineto{\pgfqpoint{4.838043in}{5.056174in}}%
\pgfpathlineto{\pgfqpoint{4.826929in}{5.025516in}}%
\pgfpathlineto{\pgfqpoint{4.825298in}{5.007539in}}%
\pgfpathlineto{\pgfqpoint{4.847794in}{5.009227in}}%
\pgfpathlineto{\pgfqpoint{4.857997in}{5.000604in}}%
\pgfpathlineto{\pgfqpoint{4.880876in}{4.965354in}}%
\pgfpathlineto{\pgfqpoint{4.892247in}{4.961067in}}%
\pgfpathlineto{\pgfqpoint{4.911295in}{4.962657in}}%
\pgfpathlineto{\pgfqpoint{4.924550in}{4.974636in}}%
\pgfpathlineto{\pgfqpoint{4.933361in}{4.963945in}}%
\pgfpathlineto{\pgfqpoint{4.877952in}{4.934706in}}%
\pgfpathlineto{\pgfqpoint{4.876220in}{4.955687in}}%
\pgfpathlineto{\pgfqpoint{4.851050in}{4.923078in}}%
\pgfpathlineto{\pgfqpoint{4.842200in}{4.917460in}}%
\pgfpathlineto{\pgfqpoint{4.829857in}{4.936276in}}%
\pgfpathlineto{\pgfqpoint{4.826478in}{4.938855in}}%
\pgfpathlineto{\pgfqpoint{4.814987in}{4.944941in}}%
\pgfpathlineto{\pgfqpoint{4.804966in}{4.969628in}}%
\pgfpathlineto{\pgfqpoint{4.770516in}{4.959426in}}%
\pgfusepath{stroke}%
\end{pgfscope}%
\begin{pgfscope}%
\pgfpathrectangle{\pgfqpoint{0.100000in}{2.493063in}}{\pgfqpoint{5.037500in}{3.427208in}}%
\pgfusepath{clip}%
\pgfsetbuttcap%
\pgfsetroundjoin%
\pgfsetlinewidth{0.501875pt}%
\definecolor{currentstroke}{rgb}{0.827451,0.827451,0.827451}%
\pgfsetstrokecolor{currentstroke}%
\pgfsetdash{}{0pt}%
\pgfpathmoveto{\pgfqpoint{2.182662in}{4.500100in}}%
\pgfpathlineto{\pgfqpoint{2.191913in}{4.614814in}}%
\pgfpathlineto{\pgfqpoint{2.139519in}{4.619067in}}%
\pgfpathlineto{\pgfqpoint{2.018405in}{4.630818in}}%
\pgfpathlineto{\pgfqpoint{2.029673in}{4.745125in}}%
\pgfpathlineto{\pgfqpoint{2.040958in}{4.860305in}}%
\pgfpathlineto{\pgfqpoint{2.087517in}{4.855278in}}%
\pgfpathlineto{\pgfqpoint{2.206426in}{4.844043in}}%
\pgfpathlineto{\pgfqpoint{2.278390in}{4.838408in}}%
\pgfpathlineto{\pgfqpoint{2.365512in}{4.833414in}}%
\pgfpathlineto{\pgfqpoint{2.508914in}{4.826627in}}%
\pgfpathlineto{\pgfqpoint{2.512066in}{4.820198in}}%
\pgfpathlineto{\pgfqpoint{2.553155in}{4.799869in}}%
\pgfpathlineto{\pgfqpoint{2.565756in}{4.810193in}}%
\pgfpathlineto{\pgfqpoint{2.602058in}{4.809251in}}%
\pgfpathlineto{\pgfqpoint{2.623692in}{4.797918in}}%
\pgfpathlineto{\pgfqpoint{2.657995in}{4.785035in}}%
\pgfpathlineto{\pgfqpoint{2.667607in}{4.766241in}}%
\pgfpathlineto{\pgfqpoint{2.681695in}{4.763929in}}%
\pgfpathlineto{\pgfqpoint{2.688563in}{4.726790in}}%
\pgfpathlineto{\pgfqpoint{2.711614in}{4.683044in}}%
\pgfpathlineto{\pgfqpoint{2.708345in}{4.660897in}}%
\pgfpathlineto{\pgfqpoint{2.723372in}{4.644980in}}%
\pgfpathlineto{\pgfqpoint{2.722550in}{4.630893in}}%
\pgfpathlineto{\pgfqpoint{2.729620in}{4.604217in}}%
\pgfpathlineto{\pgfqpoint{2.731440in}{4.576676in}}%
\pgfpathlineto{\pgfqpoint{2.727607in}{4.560599in}}%
\pgfpathlineto{\pgfqpoint{2.736996in}{4.544227in}}%
\pgfpathlineto{\pgfqpoint{2.749875in}{4.514454in}}%
\pgfpathlineto{\pgfqpoint{2.762201in}{4.502337in}}%
\pgfpathlineto{\pgfqpoint{2.776895in}{4.476079in}}%
\pgfpathlineto{\pgfqpoint{2.735250in}{4.475647in}}%
\pgfpathlineto{\pgfqpoint{2.645206in}{4.477040in}}%
\pgfpathlineto{\pgfqpoint{2.545721in}{4.480088in}}%
\pgfpathlineto{\pgfqpoint{2.445651in}{4.483928in}}%
\pgfpathlineto{\pgfqpoint{2.298517in}{4.491965in}}%
\pgfpathlineto{\pgfqpoint{2.182662in}{4.500100in}}%
\pgfusepath{stroke}%
\end{pgfscope}%
\begin{pgfscope}%
\pgfpathrectangle{\pgfqpoint{0.100000in}{2.493063in}}{\pgfqpoint{5.037500in}{3.427208in}}%
\pgfusepath{clip}%
\pgfsetbuttcap%
\pgfsetroundjoin%
\pgfsetlinewidth{0.501875pt}%
\definecolor{currentstroke}{rgb}{0.827451,0.827451,0.827451}%
\pgfsetstrokecolor{currentstroke}%
\pgfsetdash{}{0pt}%
\pgfpathmoveto{\pgfqpoint{4.098414in}{4.846121in}}%
\pgfpathlineto{\pgfqpoint{4.144485in}{4.890050in}}%
\pgfpathlineto{\pgfqpoint{4.150212in}{4.905665in}}%
\pgfpathlineto{\pgfqpoint{4.163712in}{4.919033in}}%
\pgfpathlineto{\pgfqpoint{4.153582in}{4.938496in}}%
\pgfpathlineto{\pgfqpoint{4.140881in}{4.949880in}}%
\pgfpathlineto{\pgfqpoint{4.137210in}{4.970054in}}%
\pgfpathlineto{\pgfqpoint{4.184483in}{4.990768in}}%
\pgfpathlineto{\pgfqpoint{4.223622in}{4.997320in}}%
\pgfpathlineto{\pgfqpoint{4.244653in}{4.997758in}}%
\pgfpathlineto{\pgfqpoint{4.260673in}{4.989843in}}%
\pgfpathlineto{\pgfqpoint{4.276290in}{4.996911in}}%
\pgfpathlineto{\pgfqpoint{4.314380in}{5.004832in}}%
\pgfpathlineto{\pgfqpoint{4.327541in}{5.015056in}}%
\pgfpathlineto{\pgfqpoint{4.347056in}{5.037689in}}%
\pgfpathlineto{\pgfqpoint{4.364809in}{5.047690in}}%
\pgfpathlineto{\pgfqpoint{4.366061in}{5.057283in}}%
\pgfpathlineto{\pgfqpoint{4.356744in}{5.079166in}}%
\pgfpathlineto{\pgfqpoint{4.363480in}{5.092045in}}%
\pgfpathlineto{\pgfqpoint{4.354410in}{5.105892in}}%
\pgfpathlineto{\pgfqpoint{4.340522in}{5.106859in}}%
\pgfpathlineto{\pgfqpoint{4.375282in}{5.148688in}}%
\pgfpathlineto{\pgfqpoint{4.379411in}{5.164629in}}%
\pgfpathlineto{\pgfqpoint{4.406769in}{5.205284in}}%
\pgfpathlineto{\pgfqpoint{4.432165in}{5.227284in}}%
\pgfpathlineto{\pgfqpoint{4.449604in}{5.236471in}}%
\pgfpathlineto{\pgfqpoint{4.506667in}{5.249401in}}%
\pgfpathlineto{\pgfqpoint{4.560150in}{5.264002in}}%
\pgfpathlineto{\pgfqpoint{4.567234in}{5.241082in}}%
\pgfpathlineto{\pgfqpoint{4.567736in}{5.219830in}}%
\pgfpathlineto{\pgfqpoint{4.579935in}{5.200538in}}%
\pgfpathlineto{\pgfqpoint{4.583695in}{5.180896in}}%
\pgfpathlineto{\pgfqpoint{4.579007in}{5.154245in}}%
\pgfpathlineto{\pgfqpoint{4.593285in}{5.118300in}}%
\pgfpathlineto{\pgfqpoint{4.609564in}{5.098312in}}%
\pgfpathlineto{\pgfqpoint{4.622154in}{5.033342in}}%
\pgfpathlineto{\pgfqpoint{4.628939in}{5.011925in}}%
\pgfpathlineto{\pgfqpoint{4.627187in}{4.932145in}}%
\pgfpathlineto{\pgfqpoint{4.629215in}{4.929072in}}%
\pgfpathlineto{\pgfqpoint{4.644037in}{4.843268in}}%
\pgfpathlineto{\pgfqpoint{4.652350in}{4.835449in}}%
\pgfpathlineto{\pgfqpoint{4.634470in}{4.818078in}}%
\pgfpathlineto{\pgfqpoint{4.643292in}{4.808107in}}%
\pgfpathlineto{\pgfqpoint{4.635538in}{4.793027in}}%
\pgfpathlineto{\pgfqpoint{4.635604in}{4.786606in}}%
\pgfpathlineto{\pgfqpoint{4.625910in}{4.780813in}}%
\pgfpathlineto{\pgfqpoint{4.621189in}{4.768007in}}%
\pgfpathlineto{\pgfqpoint{4.622686in}{4.803221in}}%
\pgfpathlineto{\pgfqpoint{4.592639in}{4.810974in}}%
\pgfpathlineto{\pgfqpoint{4.545650in}{4.827002in}}%
\pgfpathlineto{\pgfqpoint{4.539004in}{4.834897in}}%
\pgfpathlineto{\pgfqpoint{4.519368in}{4.836789in}}%
\pgfpathlineto{\pgfqpoint{4.508073in}{4.848540in}}%
\pgfpathlineto{\pgfqpoint{4.501979in}{4.871943in}}%
\pgfpathlineto{\pgfqpoint{4.485967in}{4.874854in}}%
\pgfpathlineto{\pgfqpoint{4.475258in}{4.887131in}}%
\pgfpathlineto{\pgfqpoint{4.338937in}{4.859648in}}%
\pgfpathlineto{\pgfqpoint{4.273785in}{4.846174in}}%
\pgfpathlineto{\pgfqpoint{4.174866in}{4.828119in}}%
\pgfpathlineto{\pgfqpoint{4.103631in}{4.816102in}}%
\pgfpathlineto{\pgfqpoint{4.098414in}{4.846121in}}%
\pgfusepath{stroke}%
\end{pgfscope}%
\begin{pgfscope}%
\pgfpathrectangle{\pgfqpoint{0.100000in}{2.493063in}}{\pgfqpoint{5.037500in}{3.427208in}}%
\pgfusepath{clip}%
\pgfsetbuttcap%
\pgfsetroundjoin%
\pgfsetlinewidth{0.501875pt}%
\definecolor{currentstroke}{rgb}{0.827451,0.827451,0.827451}%
\pgfsetstrokecolor{currentstroke}%
\pgfsetdash{}{0pt}%
\pgfpathmoveto{\pgfqpoint{4.645692in}{4.760855in}}%
\pgfpathlineto{\pgfqpoint{4.624609in}{4.754287in}}%
\pgfpathlineto{\pgfqpoint{4.621062in}{4.760325in}}%
\pgfpathlineto{\pgfqpoint{4.627829in}{4.780594in}}%
\pgfpathlineto{\pgfqpoint{4.640362in}{4.782589in}}%
\pgfpathlineto{\pgfqpoint{4.639240in}{4.788971in}}%
\pgfpathlineto{\pgfqpoint{4.650493in}{4.798549in}}%
\pgfpathlineto{\pgfqpoint{4.683026in}{4.806179in}}%
\pgfpathlineto{\pgfqpoint{4.697509in}{4.817730in}}%
\pgfpathlineto{\pgfqpoint{4.730045in}{4.827355in}}%
\pgfpathlineto{\pgfqpoint{4.744884in}{4.823833in}}%
\pgfpathlineto{\pgfqpoint{4.757367in}{4.839352in}}%
\pgfpathlineto{\pgfqpoint{4.776234in}{4.841403in}}%
\pgfpathlineto{\pgfqpoint{4.744046in}{4.811132in}}%
\pgfpathlineto{\pgfqpoint{4.645692in}{4.760855in}}%
\pgfusepath{stroke}%
\end{pgfscope}%
\begin{pgfscope}%
\pgfpathrectangle{\pgfqpoint{0.100000in}{2.493063in}}{\pgfqpoint{5.037500in}{3.427208in}}%
\pgfusepath{clip}%
\pgfsetbuttcap%
\pgfsetroundjoin%
\pgfsetlinewidth{0.501875pt}%
\definecolor{currentstroke}{rgb}{0.827451,0.827451,0.827451}%
\pgfsetstrokecolor{currentstroke}%
\pgfsetdash{}{0pt}%
\pgfpathmoveto{\pgfqpoint{4.172033in}{4.560971in}}%
\pgfpathlineto{\pgfqpoint{4.080863in}{4.545897in}}%
\pgfpathlineto{\pgfqpoint{4.064326in}{4.650087in}}%
\pgfpathlineto{\pgfqpoint{4.039771in}{4.803670in}}%
\pgfpathlineto{\pgfqpoint{4.098414in}{4.846121in}}%
\pgfpathlineto{\pgfqpoint{4.103631in}{4.816102in}}%
\pgfpathlineto{\pgfqpoint{4.174866in}{4.828119in}}%
\pgfpathlineto{\pgfqpoint{4.273785in}{4.846174in}}%
\pgfpathlineto{\pgfqpoint{4.338937in}{4.859648in}}%
\pgfpathlineto{\pgfqpoint{4.475258in}{4.887131in}}%
\pgfpathlineto{\pgfqpoint{4.485967in}{4.874854in}}%
\pgfpathlineto{\pgfqpoint{4.501979in}{4.871943in}}%
\pgfpathlineto{\pgfqpoint{4.508073in}{4.848540in}}%
\pgfpathlineto{\pgfqpoint{4.519368in}{4.836789in}}%
\pgfpathlineto{\pgfqpoint{4.539004in}{4.834897in}}%
\pgfpathlineto{\pgfqpoint{4.545650in}{4.827002in}}%
\pgfpathlineto{\pgfqpoint{4.538903in}{4.820911in}}%
\pgfpathlineto{\pgfqpoint{4.532824in}{4.799359in}}%
\pgfpathlineto{\pgfqpoint{4.518581in}{4.775143in}}%
\pgfpathlineto{\pgfqpoint{4.528142in}{4.764611in}}%
\pgfpathlineto{\pgfqpoint{4.520318in}{4.753328in}}%
\pgfpathlineto{\pgfqpoint{4.524732in}{4.728584in}}%
\pgfpathlineto{\pgfqpoint{4.538821in}{4.715595in}}%
\pgfpathlineto{\pgfqpoint{4.572254in}{4.694493in}}%
\pgfpathlineto{\pgfqpoint{4.545330in}{4.664718in}}%
\pgfpathlineto{\pgfqpoint{4.544953in}{4.653417in}}%
\pgfpathlineto{\pgfqpoint{4.523034in}{4.638842in}}%
\pgfpathlineto{\pgfqpoint{4.498797in}{4.636107in}}%
\pgfpathlineto{\pgfqpoint{4.492802in}{4.623439in}}%
\pgfpathlineto{\pgfqpoint{4.425387in}{4.608852in}}%
\pgfpathlineto{\pgfqpoint{4.346722in}{4.593009in}}%
\pgfpathlineto{\pgfqpoint{4.292651in}{4.583327in}}%
\pgfpathlineto{\pgfqpoint{4.172033in}{4.560971in}}%
\pgfusepath{stroke}%
\end{pgfscope}%
\begin{pgfscope}%
\pgfpathrectangle{\pgfqpoint{0.100000in}{2.493063in}}{\pgfqpoint{5.037500in}{3.427208in}}%
\pgfusepath{clip}%
\pgfsetbuttcap%
\pgfsetroundjoin%
\pgfsetlinewidth{0.501875pt}%
\definecolor{currentstroke}{rgb}{0.827451,0.827451,0.827451}%
\pgfsetstrokecolor{currentstroke}%
\pgfsetdash{}{0pt}%
\pgfpathmoveto{\pgfqpoint{4.629215in}{4.929072in}}%
\pgfpathlineto{\pgfqpoint{4.745254in}{4.955440in}}%
\pgfpathlineto{\pgfqpoint{4.770516in}{4.959426in}}%
\pgfpathlineto{\pgfqpoint{4.787236in}{4.893688in}}%
\pgfpathlineto{\pgfqpoint{4.784590in}{4.881916in}}%
\pgfpathlineto{\pgfqpoint{4.730879in}{4.861122in}}%
\pgfpathlineto{\pgfqpoint{4.698815in}{4.853854in}}%
\pgfpathlineto{\pgfqpoint{4.685191in}{4.837551in}}%
\pgfpathlineto{\pgfqpoint{4.643292in}{4.808107in}}%
\pgfpathlineto{\pgfqpoint{4.634470in}{4.818078in}}%
\pgfpathlineto{\pgfqpoint{4.652350in}{4.835449in}}%
\pgfpathlineto{\pgfqpoint{4.644037in}{4.843268in}}%
\pgfpathlineto{\pgfqpoint{4.629215in}{4.929072in}}%
\pgfusepath{stroke}%
\end{pgfscope}%
\begin{pgfscope}%
\pgfpathrectangle{\pgfqpoint{0.100000in}{2.493063in}}{\pgfqpoint{5.037500in}{3.427208in}}%
\pgfusepath{clip}%
\pgfsetbuttcap%
\pgfsetroundjoin%
\pgfsetlinewidth{0.501875pt}%
\definecolor{currentstroke}{rgb}{0.827451,0.827451,0.827451}%
\pgfsetstrokecolor{currentstroke}%
\pgfsetdash{}{0pt}%
\pgfpathmoveto{\pgfqpoint{4.770516in}{4.959426in}}%
\pgfpathlineto{\pgfqpoint{4.804966in}{4.969628in}}%
\pgfpathlineto{\pgfqpoint{4.814987in}{4.944941in}}%
\pgfpathlineto{\pgfqpoint{4.826478in}{4.938855in}}%
\pgfpathlineto{\pgfqpoint{4.812298in}{4.928456in}}%
\pgfpathlineto{\pgfqpoint{4.814086in}{4.897915in}}%
\pgfpathlineto{\pgfqpoint{4.784590in}{4.881916in}}%
\pgfpathlineto{\pgfqpoint{4.787236in}{4.893688in}}%
\pgfpathlineto{\pgfqpoint{4.770516in}{4.959426in}}%
\pgfusepath{stroke}%
\end{pgfscope}%
\begin{pgfscope}%
\pgfpathrectangle{\pgfqpoint{0.100000in}{2.493063in}}{\pgfqpoint{5.037500in}{3.427208in}}%
\pgfusepath{clip}%
\pgfsetbuttcap%
\pgfsetroundjoin%
\pgfsetlinewidth{0.501875pt}%
\definecolor{currentstroke}{rgb}{0.827451,0.827451,0.827451}%
\pgfsetstrokecolor{currentstroke}%
\pgfsetdash{}{0pt}%
\pgfpathmoveto{\pgfqpoint{4.519150in}{4.628078in}}%
\pgfpathlineto{\pgfqpoint{4.523034in}{4.638842in}}%
\pgfpathlineto{\pgfqpoint{4.544953in}{4.653417in}}%
\pgfpathlineto{\pgfqpoint{4.545330in}{4.664718in}}%
\pgfpathlineto{\pgfqpoint{4.572254in}{4.694493in}}%
\pgfpathlineto{\pgfqpoint{4.538821in}{4.715595in}}%
\pgfpathlineto{\pgfqpoint{4.524732in}{4.728584in}}%
\pgfpathlineto{\pgfqpoint{4.520318in}{4.753328in}}%
\pgfpathlineto{\pgfqpoint{4.528142in}{4.764611in}}%
\pgfpathlineto{\pgfqpoint{4.518581in}{4.775143in}}%
\pgfpathlineto{\pgfqpoint{4.532824in}{4.799359in}}%
\pgfpathlineto{\pgfqpoint{4.538903in}{4.820911in}}%
\pgfpathlineto{\pgfqpoint{4.545650in}{4.827002in}}%
\pgfpathlineto{\pgfqpoint{4.592639in}{4.810974in}}%
\pgfpathlineto{\pgfqpoint{4.622686in}{4.803221in}}%
\pgfpathlineto{\pgfqpoint{4.621189in}{4.768007in}}%
\pgfpathlineto{\pgfqpoint{4.602937in}{4.741303in}}%
\pgfpathlineto{\pgfqpoint{4.617977in}{4.737374in}}%
\pgfpathlineto{\pgfqpoint{4.633592in}{4.725915in}}%
\pgfpathlineto{\pgfqpoint{4.634513in}{4.694571in}}%
\pgfpathlineto{\pgfqpoint{4.629787in}{4.672383in}}%
\pgfpathlineto{\pgfqpoint{4.632932in}{4.654157in}}%
\pgfpathlineto{\pgfqpoint{4.605787in}{4.592656in}}%
\pgfpathlineto{\pgfqpoint{4.596032in}{4.563939in}}%
\pgfpathlineto{\pgfqpoint{4.582411in}{4.577898in}}%
\pgfpathlineto{\pgfqpoint{4.564348in}{4.575519in}}%
\pgfpathlineto{\pgfqpoint{4.519141in}{4.601637in}}%
\pgfpathlineto{\pgfqpoint{4.514516in}{4.615624in}}%
\pgfpathlineto{\pgfqpoint{4.519150in}{4.628078in}}%
\pgfusepath{stroke}%
\end{pgfscope}%
\begin{pgfscope}%
\pgfpathrectangle{\pgfqpoint{0.100000in}{2.493063in}}{\pgfqpoint{5.037500in}{3.427208in}}%
\pgfusepath{clip}%
\pgfsetbuttcap%
\pgfsetroundjoin%
\pgfsetlinewidth{0.501875pt}%
\definecolor{currentstroke}{rgb}{0.827451,0.827451,0.827451}%
\pgfsetstrokecolor{currentstroke}%
\pgfsetdash{}{0pt}%
\pgfpathmoveto{\pgfqpoint{3.437265in}{4.247692in}}%
\pgfpathlineto{\pgfqpoint{3.435151in}{4.276223in}}%
\pgfpathlineto{\pgfqpoint{3.443366in}{4.290378in}}%
\pgfpathlineto{\pgfqpoint{3.437951in}{4.297314in}}%
\pgfpathlineto{\pgfqpoint{3.457695in}{4.320023in}}%
\pgfpathlineto{\pgfqpoint{3.456564in}{4.325496in}}%
\pgfpathlineto{\pgfqpoint{3.475746in}{4.365114in}}%
\pgfpathlineto{\pgfqpoint{3.471869in}{4.384221in}}%
\pgfpathlineto{\pgfqpoint{3.457959in}{4.404225in}}%
\pgfpathlineto{\pgfqpoint{3.467615in}{4.428563in}}%
\pgfpathlineto{\pgfqpoint{3.455099in}{4.588720in}}%
\pgfpathlineto{\pgfqpoint{3.446087in}{4.701075in}}%
\pgfpathlineto{\pgfqpoint{3.458537in}{4.691765in}}%
\pgfpathlineto{\pgfqpoint{3.472426in}{4.692017in}}%
\pgfpathlineto{\pgfqpoint{3.505233in}{4.711021in}}%
\pgfpathlineto{\pgfqpoint{3.605857in}{4.720409in}}%
\pgfpathlineto{\pgfqpoint{3.680336in}{4.728259in}}%
\pgfpathlineto{\pgfqpoint{3.680994in}{4.720971in}}%
\pgfpathlineto{\pgfqpoint{3.697996in}{4.567004in}}%
\pgfpathlineto{\pgfqpoint{3.712630in}{4.423743in}}%
\pgfpathlineto{\pgfqpoint{3.706322in}{4.417017in}}%
\pgfpathlineto{\pgfqpoint{3.716011in}{4.400323in}}%
\pgfpathlineto{\pgfqpoint{3.715962in}{4.388292in}}%
\pgfpathlineto{\pgfqpoint{3.702135in}{4.385269in}}%
\pgfpathlineto{\pgfqpoint{3.686663in}{4.373671in}}%
\pgfpathlineto{\pgfqpoint{3.676188in}{4.378237in}}%
\pgfpathlineto{\pgfqpoint{3.660509in}{4.370813in}}%
\pgfpathlineto{\pgfqpoint{3.665360in}{4.355902in}}%
\pgfpathlineto{\pgfqpoint{3.649244in}{4.340922in}}%
\pgfpathlineto{\pgfqpoint{3.644795in}{4.323592in}}%
\pgfpathlineto{\pgfqpoint{3.633714in}{4.320750in}}%
\pgfpathlineto{\pgfqpoint{3.625445in}{4.307626in}}%
\pgfpathlineto{\pgfqpoint{3.626527in}{4.294409in}}%
\pgfpathlineto{\pgfqpoint{3.616798in}{4.285115in}}%
\pgfpathlineto{\pgfqpoint{3.602160in}{4.286550in}}%
\pgfpathlineto{\pgfqpoint{3.591032in}{4.300810in}}%
\pgfpathlineto{\pgfqpoint{3.572178in}{4.287078in}}%
\pgfpathlineto{\pgfqpoint{3.573500in}{4.275078in}}%
\pgfpathlineto{\pgfqpoint{3.552531in}{4.268067in}}%
\pgfpathlineto{\pgfqpoint{3.547262in}{4.276892in}}%
\pgfpathlineto{\pgfqpoint{3.526847in}{4.266886in}}%
\pgfpathlineto{\pgfqpoint{3.519389in}{4.252561in}}%
\pgfpathlineto{\pgfqpoint{3.494803in}{4.267250in}}%
\pgfpathlineto{\pgfqpoint{3.455836in}{4.257519in}}%
\pgfpathlineto{\pgfqpoint{3.437265in}{4.247692in}}%
\pgfusepath{stroke}%
\end{pgfscope}%
\begin{pgfscope}%
\pgfpathrectangle{\pgfqpoint{0.100000in}{2.493063in}}{\pgfqpoint{5.037500in}{3.427208in}}%
\pgfusepath{clip}%
\pgfsetbuttcap%
\pgfsetroundjoin%
\pgfsetlinewidth{0.501875pt}%
\definecolor{currentstroke}{rgb}{0.827451,0.827451,0.827451}%
\pgfsetstrokecolor{currentstroke}%
\pgfsetdash{}{0pt}%
\pgfpathmoveto{\pgfqpoint{0.689161in}{4.997975in}}%
\pgfpathlineto{\pgfqpoint{0.745509in}{4.982295in}}%
\pgfpathlineto{\pgfqpoint{0.838915in}{4.958210in}}%
\pgfpathlineto{\pgfqpoint{0.935738in}{4.933400in}}%
\pgfpathlineto{\pgfqpoint{1.024981in}{4.911878in}}%
\pgfpathlineto{\pgfqpoint{1.102093in}{4.894420in}}%
\pgfpathlineto{\pgfqpoint{1.184593in}{4.876279in}}%
\pgfpathlineto{\pgfqpoint{1.160756in}{4.763789in}}%
\pgfpathlineto{\pgfqpoint{1.139521in}{4.663892in}}%
\pgfpathlineto{\pgfqpoint{1.104685in}{4.502710in}}%
\pgfpathlineto{\pgfqpoint{1.078531in}{4.381052in}}%
\pgfpathlineto{\pgfqpoint{1.064402in}{4.313160in}}%
\pgfpathlineto{\pgfqpoint{1.046284in}{4.225043in}}%
\pgfpathlineto{\pgfqpoint{1.033749in}{4.207168in}}%
\pgfpathlineto{\pgfqpoint{1.023660in}{4.206570in}}%
\pgfpathlineto{\pgfqpoint{1.016464in}{4.222175in}}%
\pgfpathlineto{\pgfqpoint{0.999941in}{4.227841in}}%
\pgfpathlineto{\pgfqpoint{0.982136in}{4.225855in}}%
\pgfpathlineto{\pgfqpoint{0.977083in}{4.213092in}}%
\pgfpathlineto{\pgfqpoint{0.976556in}{4.168769in}}%
\pgfpathlineto{\pgfqpoint{0.971224in}{4.158649in}}%
\pgfpathlineto{\pgfqpoint{0.974271in}{4.123066in}}%
\pgfpathlineto{\pgfqpoint{0.963137in}{4.099360in}}%
\pgfpathlineto{\pgfqpoint{0.872849in}{4.238390in}}%
\pgfpathlineto{\pgfqpoint{0.783931in}{4.373488in}}%
\pgfpathlineto{\pgfqpoint{0.737645in}{4.444331in}}%
\pgfpathlineto{\pgfqpoint{0.699033in}{4.505478in}}%
\pgfpathlineto{\pgfqpoint{0.650381in}{4.581117in}}%
\pgfpathlineto{\pgfqpoint{0.595308in}{4.665963in}}%
\pgfpathlineto{\pgfqpoint{0.617950in}{4.746497in}}%
\pgfpathlineto{\pgfqpoint{0.663516in}{4.908009in}}%
\pgfpathlineto{\pgfqpoint{0.689161in}{4.997975in}}%
\pgfusepath{stroke}%
\end{pgfscope}%
\begin{pgfscope}%
\pgfpathrectangle{\pgfqpoint{0.100000in}{2.493063in}}{\pgfqpoint{5.037500in}{3.427208in}}%
\pgfusepath{clip}%
\pgfsetbuttcap%
\pgfsetroundjoin%
\pgfsetlinewidth{0.501875pt}%
\definecolor{currentstroke}{rgb}{0.827451,0.827451,0.827451}%
\pgfsetstrokecolor{currentstroke}%
\pgfsetdash{}{0pt}%
\pgfpathmoveto{\pgfqpoint{1.184593in}{4.876279in}}%
\pgfpathlineto{\pgfqpoint{1.272793in}{4.858734in}}%
\pgfpathlineto{\pgfqpoint{1.436081in}{4.827212in}}%
\pgfpathlineto{\pgfqpoint{1.415642in}{4.713831in}}%
\pgfpathlineto{\pgfqpoint{1.505537in}{4.698440in}}%
\pgfpathlineto{\pgfqpoint{1.587436in}{4.685334in}}%
\pgfpathlineto{\pgfqpoint{1.573205in}{4.595697in}}%
\pgfpathlineto{\pgfqpoint{1.558125in}{4.499035in}}%
\pgfpathlineto{\pgfqpoint{1.537936in}{4.372117in}}%
\pgfpathlineto{\pgfqpoint{1.537414in}{4.361477in}}%
\pgfpathlineto{\pgfqpoint{1.516457in}{4.229943in}}%
\pgfpathlineto{\pgfqpoint{1.430188in}{4.243315in}}%
\pgfpathlineto{\pgfqpoint{1.386229in}{4.252143in}}%
\pgfpathlineto{\pgfqpoint{1.227336in}{4.280074in}}%
\pgfpathlineto{\pgfqpoint{1.167500in}{4.291874in}}%
\pgfpathlineto{\pgfqpoint{1.064402in}{4.313160in}}%
\pgfpathlineto{\pgfqpoint{1.078531in}{4.381052in}}%
\pgfpathlineto{\pgfqpoint{1.104685in}{4.502710in}}%
\pgfpathlineto{\pgfqpoint{1.139521in}{4.663892in}}%
\pgfpathlineto{\pgfqpoint{1.160756in}{4.763789in}}%
\pgfpathlineto{\pgfqpoint{1.184593in}{4.876279in}}%
\pgfusepath{stroke}%
\end{pgfscope}%
\begin{pgfscope}%
\pgfpathrectangle{\pgfqpoint{0.100000in}{2.493063in}}{\pgfqpoint{5.037500in}{3.427208in}}%
\pgfusepath{clip}%
\pgfsetbuttcap%
\pgfsetroundjoin%
\pgfsetlinewidth{0.501875pt}%
\definecolor{currentstroke}{rgb}{0.827451,0.827451,0.827451}%
\pgfsetstrokecolor{currentstroke}%
\pgfsetdash{}{0pt}%
\pgfpathmoveto{\pgfqpoint{0.344393in}{5.104965in}}%
\pgfpathlineto{\pgfqpoint{0.375733in}{5.094034in}}%
\pgfpathlineto{\pgfqpoint{0.501273in}{5.054949in}}%
\pgfpathlineto{\pgfqpoint{0.616760in}{5.018599in}}%
\pgfpathlineto{\pgfqpoint{0.689161in}{4.997975in}}%
\pgfpathlineto{\pgfqpoint{0.663516in}{4.908009in}}%
\pgfpathlineto{\pgfqpoint{0.617950in}{4.746497in}}%
\pgfpathlineto{\pgfqpoint{0.595308in}{4.665963in}}%
\pgfpathlineto{\pgfqpoint{0.650381in}{4.581117in}}%
\pgfpathlineto{\pgfqpoint{0.699033in}{4.505478in}}%
\pgfpathlineto{\pgfqpoint{0.737645in}{4.444331in}}%
\pgfpathlineto{\pgfqpoint{0.783931in}{4.373488in}}%
\pgfpathlineto{\pgfqpoint{0.872849in}{4.238390in}}%
\pgfpathlineto{\pgfqpoint{0.963137in}{4.099360in}}%
\pgfpathlineto{\pgfqpoint{0.959513in}{4.085573in}}%
\pgfpathlineto{\pgfqpoint{0.970403in}{4.063615in}}%
\pgfpathlineto{\pgfqpoint{0.972557in}{4.033579in}}%
\pgfpathlineto{\pgfqpoint{0.988373in}{4.019008in}}%
\pgfpathlineto{\pgfqpoint{0.988999in}{4.007268in}}%
\pgfpathlineto{\pgfqpoint{0.960679in}{3.993943in}}%
\pgfpathlineto{\pgfqpoint{0.947187in}{3.980602in}}%
\pgfpathlineto{\pgfqpoint{0.942979in}{3.951143in}}%
\pgfpathlineto{\pgfqpoint{0.936150in}{3.935078in}}%
\pgfpathlineto{\pgfqpoint{0.921760in}{3.921535in}}%
\pgfpathlineto{\pgfqpoint{0.907448in}{3.886322in}}%
\pgfpathlineto{\pgfqpoint{0.927572in}{3.867979in}}%
\pgfpathlineto{\pgfqpoint{0.924974in}{3.852858in}}%
\pgfpathlineto{\pgfqpoint{0.908473in}{3.842335in}}%
\pgfpathlineto{\pgfqpoint{0.897097in}{3.844192in}}%
\pgfpathlineto{\pgfqpoint{0.762813in}{3.862827in}}%
\pgfpathlineto{\pgfqpoint{0.663632in}{3.876898in}}%
\pgfpathlineto{\pgfqpoint{0.667969in}{3.892911in}}%
\pgfpathlineto{\pgfqpoint{0.661512in}{3.919496in}}%
\pgfpathlineto{\pgfqpoint{0.660861in}{3.946312in}}%
\pgfpathlineto{\pgfqpoint{0.656642in}{3.962005in}}%
\pgfpathlineto{\pgfqpoint{0.643643in}{3.984441in}}%
\pgfpathlineto{\pgfqpoint{0.606255in}{4.036207in}}%
\pgfpathlineto{\pgfqpoint{0.594010in}{4.042550in}}%
\pgfpathlineto{\pgfqpoint{0.578232in}{4.042450in}}%
\pgfpathlineto{\pgfqpoint{0.581837in}{4.058780in}}%
\pgfpathlineto{\pgfqpoint{0.574366in}{4.079201in}}%
\pgfpathlineto{\pgfqpoint{0.537669in}{4.089351in}}%
\pgfpathlineto{\pgfqpoint{0.515300in}{4.108145in}}%
\pgfpathlineto{\pgfqpoint{0.513445in}{4.119644in}}%
\pgfpathlineto{\pgfqpoint{0.487688in}{4.148115in}}%
\pgfpathlineto{\pgfqpoint{0.463170in}{4.153562in}}%
\pgfpathlineto{\pgfqpoint{0.440454in}{4.168031in}}%
\pgfpathlineto{\pgfqpoint{0.410577in}{4.173037in}}%
\pgfpathlineto{\pgfqpoint{0.397812in}{4.192333in}}%
\pgfpathlineto{\pgfqpoint{0.409935in}{4.222878in}}%
\pgfpathlineto{\pgfqpoint{0.406227in}{4.229746in}}%
\pgfpathlineto{\pgfqpoint{0.416322in}{4.255209in}}%
\pgfpathlineto{\pgfqpoint{0.398364in}{4.268760in}}%
\pgfpathlineto{\pgfqpoint{0.404191in}{4.293326in}}%
\pgfpathlineto{\pgfqpoint{0.394598in}{4.299606in}}%
\pgfpathlineto{\pgfqpoint{0.386316in}{4.322826in}}%
\pgfpathlineto{\pgfqpoint{0.376322in}{4.329907in}}%
\pgfpathlineto{\pgfqpoint{0.375570in}{4.346679in}}%
\pgfpathlineto{\pgfqpoint{0.367740in}{4.358505in}}%
\pgfpathlineto{\pgfqpoint{0.355940in}{4.398337in}}%
\pgfpathlineto{\pgfqpoint{0.343030in}{4.417414in}}%
\pgfpathlineto{\pgfqpoint{0.345836in}{4.449800in}}%
\pgfpathlineto{\pgfqpoint{0.361018in}{4.453059in}}%
\pgfpathlineto{\pgfqpoint{0.370904in}{4.470624in}}%
\pgfpathlineto{\pgfqpoint{0.364931in}{4.489676in}}%
\pgfpathlineto{\pgfqpoint{0.348784in}{4.492856in}}%
\pgfpathlineto{\pgfqpoint{0.340727in}{4.501767in}}%
\pgfpathlineto{\pgfqpoint{0.327682in}{4.534544in}}%
\pgfpathlineto{\pgfqpoint{0.333779in}{4.546316in}}%
\pgfpathlineto{\pgfqpoint{0.329454in}{4.568234in}}%
\pgfpathlineto{\pgfqpoint{0.339031in}{4.596622in}}%
\pgfpathlineto{\pgfqpoint{0.350249in}{4.586168in}}%
\pgfpathlineto{\pgfqpoint{0.345158in}{4.573826in}}%
\pgfpathlineto{\pgfqpoint{0.364600in}{4.554261in}}%
\pgfpathlineto{\pgfqpoint{0.363359in}{4.583314in}}%
\pgfpathlineto{\pgfqpoint{0.355028in}{4.591106in}}%
\pgfpathlineto{\pgfqpoint{0.364550in}{4.616651in}}%
\pgfpathlineto{\pgfqpoint{0.400566in}{4.612583in}}%
\pgfpathlineto{\pgfqpoint{0.395719in}{4.622082in}}%
\pgfpathlineto{\pgfqpoint{0.371948in}{4.621147in}}%
\pgfpathlineto{\pgfqpoint{0.360641in}{4.635535in}}%
\pgfpathlineto{\pgfqpoint{0.346407in}{4.622758in}}%
\pgfpathlineto{\pgfqpoint{0.338854in}{4.601403in}}%
\pgfpathlineto{\pgfqpoint{0.318696in}{4.630139in}}%
\pgfpathlineto{\pgfqpoint{0.310932in}{4.635371in}}%
\pgfpathlineto{\pgfqpoint{0.313889in}{4.666646in}}%
\pgfpathlineto{\pgfqpoint{0.307746in}{4.685089in}}%
\pgfpathlineto{\pgfqpoint{0.296607in}{4.702417in}}%
\pgfpathlineto{\pgfqpoint{0.273798in}{4.755509in}}%
\pgfpathlineto{\pgfqpoint{0.281255in}{4.767249in}}%
\pgfpathlineto{\pgfqpoint{0.281168in}{4.804352in}}%
\pgfpathlineto{\pgfqpoint{0.293478in}{4.825047in}}%
\pgfpathlineto{\pgfqpoint{0.296301in}{4.857389in}}%
\pgfpathlineto{\pgfqpoint{0.284693in}{4.894445in}}%
\pgfpathlineto{\pgfqpoint{0.269293in}{4.918032in}}%
\pgfpathlineto{\pgfqpoint{0.272040in}{4.939316in}}%
\pgfpathlineto{\pgfqpoint{0.315286in}{4.990828in}}%
\pgfpathlineto{\pgfqpoint{0.317434in}{5.008400in}}%
\pgfpathlineto{\pgfqpoint{0.336921in}{5.041920in}}%
\pgfpathlineto{\pgfqpoint{0.339624in}{5.073692in}}%
\pgfpathlineto{\pgfqpoint{0.333359in}{5.081799in}}%
\pgfpathlineto{\pgfqpoint{0.344393in}{5.104965in}}%
\pgfusepath{stroke}%
\end{pgfscope}%
\begin{pgfscope}%
\pgfpathrectangle{\pgfqpoint{0.100000in}{2.493063in}}{\pgfqpoint{5.037500in}{3.427208in}}%
\pgfusepath{clip}%
\pgfsetbuttcap%
\pgfsetroundjoin%
\pgfsetlinewidth{0.501875pt}%
\definecolor{currentstroke}{rgb}{0.827451,0.827451,0.827451}%
\pgfsetstrokecolor{currentstroke}%
\pgfsetdash{}{0pt}%
\pgfpathmoveto{\pgfqpoint{3.712630in}{4.423743in}}%
\pgfpathlineto{\pgfqpoint{3.697996in}{4.567004in}}%
\pgfpathlineto{\pgfqpoint{3.680994in}{4.720971in}}%
\pgfpathlineto{\pgfqpoint{3.792336in}{4.737553in}}%
\pgfpathlineto{\pgfqpoint{3.821886in}{4.729862in}}%
\pgfpathlineto{\pgfqpoint{3.836059in}{4.721513in}}%
\pgfpathlineto{\pgfqpoint{3.853819in}{4.723829in}}%
\pgfpathlineto{\pgfqpoint{3.877234in}{4.710006in}}%
\pgfpathlineto{\pgfqpoint{3.920850in}{4.730579in}}%
\pgfpathlineto{\pgfqpoint{3.944919in}{4.731266in}}%
\pgfpathlineto{\pgfqpoint{3.973027in}{4.762624in}}%
\pgfpathlineto{\pgfqpoint{4.001625in}{4.781697in}}%
\pgfpathlineto{\pgfqpoint{4.039771in}{4.803670in}}%
\pgfpathlineto{\pgfqpoint{4.064326in}{4.650087in}}%
\pgfpathlineto{\pgfqpoint{4.052917in}{4.640224in}}%
\pgfpathlineto{\pgfqpoint{4.060283in}{4.631164in}}%
\pgfpathlineto{\pgfqpoint{4.063217in}{4.611303in}}%
\pgfpathlineto{\pgfqpoint{4.057501in}{4.592649in}}%
\pgfpathlineto{\pgfqpoint{4.056746in}{4.565383in}}%
\pgfpathlineto{\pgfqpoint{4.051378in}{4.529895in}}%
\pgfpathlineto{\pgfqpoint{4.024221in}{4.498243in}}%
\pgfpathlineto{\pgfqpoint{4.012789in}{4.491525in}}%
\pgfpathlineto{\pgfqpoint{4.003899in}{4.497309in}}%
\pgfpathlineto{\pgfqpoint{3.981897in}{4.467141in}}%
\pgfpathlineto{\pgfqpoint{3.983958in}{4.438094in}}%
\pgfpathlineto{\pgfqpoint{3.959456in}{4.445080in}}%
\pgfpathlineto{\pgfqpoint{3.947853in}{4.416069in}}%
\pgfpathlineto{\pgfqpoint{3.953718in}{4.395512in}}%
\pgfpathlineto{\pgfqpoint{3.944537in}{4.392475in}}%
\pgfpathlineto{\pgfqpoint{3.943252in}{4.376232in}}%
\pgfpathlineto{\pgfqpoint{3.920717in}{4.369680in}}%
\pgfpathlineto{\pgfqpoint{3.909006in}{4.382795in}}%
\pgfpathlineto{\pgfqpoint{3.893938in}{4.387880in}}%
\pgfpathlineto{\pgfqpoint{3.890398in}{4.401178in}}%
\pgfpathlineto{\pgfqpoint{3.876773in}{4.398839in}}%
\pgfpathlineto{\pgfqpoint{3.867848in}{4.386608in}}%
\pgfpathlineto{\pgfqpoint{3.855079in}{4.382307in}}%
\pgfpathlineto{\pgfqpoint{3.832524in}{4.390962in}}%
\pgfpathlineto{\pgfqpoint{3.820049in}{4.380664in}}%
\pgfpathlineto{\pgfqpoint{3.802320in}{4.392848in}}%
\pgfpathlineto{\pgfqpoint{3.768306in}{4.396588in}}%
\pgfpathlineto{\pgfqpoint{3.758047in}{4.418723in}}%
\pgfpathlineto{\pgfqpoint{3.745048in}{4.428530in}}%
\pgfpathlineto{\pgfqpoint{3.732447in}{4.422252in}}%
\pgfpathlineto{\pgfqpoint{3.712630in}{4.423743in}}%
\pgfusepath{stroke}%
\end{pgfscope}%
\begin{pgfscope}%
\pgfpathrectangle{\pgfqpoint{0.100000in}{2.493063in}}{\pgfqpoint{5.037500in}{3.427208in}}%
\pgfusepath{clip}%
\pgfsetbuttcap%
\pgfsetroundjoin%
\pgfsetlinewidth{0.501875pt}%
\definecolor{currentstroke}{rgb}{0.827451,0.827451,0.827451}%
\pgfsetstrokecolor{currentstroke}%
\pgfsetdash{}{0pt}%
\pgfpathmoveto{\pgfqpoint{3.344010in}{4.146681in}}%
\pgfpathlineto{\pgfqpoint{3.330964in}{4.157278in}}%
\pgfpathlineto{\pgfqpoint{3.320328in}{4.152236in}}%
\pgfpathlineto{\pgfqpoint{3.306729in}{4.177628in}}%
\pgfpathlineto{\pgfqpoint{3.313676in}{4.193588in}}%
\pgfpathlineto{\pgfqpoint{3.303664in}{4.211552in}}%
\pgfpathlineto{\pgfqpoint{3.303126in}{4.225694in}}%
\pgfpathlineto{\pgfqpoint{3.289597in}{4.230734in}}%
\pgfpathlineto{\pgfqpoint{3.283323in}{4.241390in}}%
\pgfpathlineto{\pgfqpoint{3.256870in}{4.254698in}}%
\pgfpathlineto{\pgfqpoint{3.233897in}{4.271095in}}%
\pgfpathlineto{\pgfqpoint{3.223222in}{4.283474in}}%
\pgfpathlineto{\pgfqpoint{3.223001in}{4.298571in}}%
\pgfpathlineto{\pgfqpoint{3.237285in}{4.327559in}}%
\pgfpathlineto{\pgfqpoint{3.241679in}{4.349731in}}%
\pgfpathlineto{\pgfqpoint{3.230042in}{4.362258in}}%
\pgfpathlineto{\pgfqpoint{3.214629in}{4.366979in}}%
\pgfpathlineto{\pgfqpoint{3.195930in}{4.356657in}}%
\pgfpathlineto{\pgfqpoint{3.188016in}{4.374377in}}%
\pgfpathlineto{\pgfqpoint{3.184022in}{4.398399in}}%
\pgfpathlineto{\pgfqpoint{3.156465in}{4.419812in}}%
\pgfpathlineto{\pgfqpoint{3.150959in}{4.429315in}}%
\pgfpathlineto{\pgfqpoint{3.125789in}{4.450826in}}%
\pgfpathlineto{\pgfqpoint{3.117900in}{4.466465in}}%
\pgfpathlineto{\pgfqpoint{3.110789in}{4.497479in}}%
\pgfpathlineto{\pgfqpoint{3.115663in}{4.525025in}}%
\pgfpathlineto{\pgfqpoint{3.122170in}{4.528866in}}%
\pgfpathlineto{\pgfqpoint{3.120999in}{4.551925in}}%
\pgfpathlineto{\pgfqpoint{3.139365in}{4.558775in}}%
\pgfpathlineto{\pgfqpoint{3.144911in}{4.579468in}}%
\pgfpathlineto{\pgfqpoint{3.155476in}{4.593396in}}%
\pgfpathlineto{\pgfqpoint{3.154942in}{4.611116in}}%
\pgfpathlineto{\pgfqpoint{3.141846in}{4.625209in}}%
\pgfpathlineto{\pgfqpoint{3.144957in}{4.644935in}}%
\pgfpathlineto{\pgfqpoint{3.178904in}{4.653514in}}%
\pgfpathlineto{\pgfqpoint{3.204948in}{4.669174in}}%
\pgfpathlineto{\pgfqpoint{3.207681in}{4.688891in}}%
\pgfpathlineto{\pgfqpoint{3.216743in}{4.695099in}}%
\pgfpathlineto{\pgfqpoint{3.220210in}{4.715813in}}%
\pgfpathlineto{\pgfqpoint{3.217424in}{4.729495in}}%
\pgfpathlineto{\pgfqpoint{3.199619in}{4.740860in}}%
\pgfpathlineto{\pgfqpoint{3.192451in}{4.757802in}}%
\pgfpathlineto{\pgfqpoint{3.174878in}{4.774179in}}%
\pgfpathlineto{\pgfqpoint{3.319282in}{4.780333in}}%
\pgfpathlineto{\pgfqpoint{3.416161in}{4.787215in}}%
\pgfpathlineto{\pgfqpoint{3.414389in}{4.766836in}}%
\pgfpathlineto{\pgfqpoint{3.430899in}{4.738727in}}%
\pgfpathlineto{\pgfqpoint{3.437831in}{4.714711in}}%
\pgfpathlineto{\pgfqpoint{3.446087in}{4.701075in}}%
\pgfpathlineto{\pgfqpoint{3.455099in}{4.588720in}}%
\pgfpathlineto{\pgfqpoint{3.467615in}{4.428563in}}%
\pgfpathlineto{\pgfqpoint{3.457959in}{4.404225in}}%
\pgfpathlineto{\pgfqpoint{3.471869in}{4.384221in}}%
\pgfpathlineto{\pgfqpoint{3.475746in}{4.365114in}}%
\pgfpathlineto{\pgfqpoint{3.456564in}{4.325496in}}%
\pgfpathlineto{\pgfqpoint{3.457695in}{4.320023in}}%
\pgfpathlineto{\pgfqpoint{3.437951in}{4.297314in}}%
\pgfpathlineto{\pgfqpoint{3.443366in}{4.290378in}}%
\pgfpathlineto{\pgfqpoint{3.435151in}{4.276223in}}%
\pgfpathlineto{\pgfqpoint{3.437265in}{4.247692in}}%
\pgfpathlineto{\pgfqpoint{3.427288in}{4.230184in}}%
\pgfpathlineto{\pgfqpoint{3.435402in}{4.209497in}}%
\pgfpathlineto{\pgfqpoint{3.401404in}{4.198255in}}%
\pgfpathlineto{\pgfqpoint{3.398271in}{4.186030in}}%
\pgfpathlineto{\pgfqpoint{3.407555in}{4.170534in}}%
\pgfpathlineto{\pgfqpoint{3.403283in}{4.160434in}}%
\pgfpathlineto{\pgfqpoint{3.366827in}{4.172910in}}%
\pgfpathlineto{\pgfqpoint{3.354799in}{4.174166in}}%
\pgfpathlineto{\pgfqpoint{3.339802in}{4.155199in}}%
\pgfpathlineto{\pgfqpoint{3.344010in}{4.146681in}}%
\pgfusepath{stroke}%
\end{pgfscope}%
\begin{pgfscope}%
\pgfpathrectangle{\pgfqpoint{0.100000in}{2.493063in}}{\pgfqpoint{5.037500in}{3.427208in}}%
\pgfusepath{clip}%
\pgfsetbuttcap%
\pgfsetroundjoin%
\pgfsetlinewidth{0.501875pt}%
\definecolor{currentstroke}{rgb}{0.827451,0.827451,0.827451}%
\pgfsetstrokecolor{currentstroke}%
\pgfsetdash{}{0pt}%
\pgfpathmoveto{\pgfqpoint{4.405194in}{4.495246in}}%
\pgfpathlineto{\pgfqpoint{4.395161in}{4.510143in}}%
\pgfpathlineto{\pgfqpoint{4.400822in}{4.518480in}}%
\pgfpathlineto{\pgfqpoint{4.414687in}{4.509117in}}%
\pgfpathlineto{\pgfqpoint{4.405194in}{4.495246in}}%
\pgfusepath{stroke}%
\end{pgfscope}%
\begin{pgfscope}%
\pgfpathrectangle{\pgfqpoint{0.100000in}{2.493063in}}{\pgfqpoint{5.037500in}{3.427208in}}%
\pgfusepath{clip}%
\pgfsetbuttcap%
\pgfsetroundjoin%
\pgfsetlinewidth{0.501875pt}%
\definecolor{currentstroke}{rgb}{0.827451,0.827451,0.827451}%
\pgfsetstrokecolor{currentstroke}%
\pgfsetdash{}{0pt}%
\pgfpathmoveto{\pgfqpoint{4.492802in}{4.623439in}}%
\pgfpathlineto{\pgfqpoint{4.498797in}{4.636107in}}%
\pgfpathlineto{\pgfqpoint{4.523034in}{4.638842in}}%
\pgfpathlineto{\pgfqpoint{4.519150in}{4.628078in}}%
\pgfpathlineto{\pgfqpoint{4.511154in}{4.614325in}}%
\pgfpathlineto{\pgfqpoint{4.516572in}{4.597941in}}%
\pgfpathlineto{\pgfqpoint{4.537956in}{4.578309in}}%
\pgfpathlineto{\pgfqpoint{4.542919in}{4.557636in}}%
\pgfpathlineto{\pgfqpoint{4.567556in}{4.531893in}}%
\pgfpathlineto{\pgfqpoint{4.577221in}{4.532995in}}%
\pgfpathlineto{\pgfqpoint{4.589261in}{4.494349in}}%
\pgfpathlineto{\pgfqpoint{4.587286in}{4.493964in}}%
\pgfpathlineto{\pgfqpoint{4.585091in}{4.493533in}}%
\pgfpathlineto{\pgfqpoint{4.531394in}{4.483253in}}%
\pgfpathlineto{\pgfqpoint{4.502665in}{4.585425in}}%
\pgfpathlineto{\pgfqpoint{4.492802in}{4.623439in}}%
\pgfusepath{stroke}%
\end{pgfscope}%
\begin{pgfscope}%
\pgfpathrectangle{\pgfqpoint{0.100000in}{2.493063in}}{\pgfqpoint{5.037500in}{3.427208in}}%
\pgfusepath{clip}%
\pgfsetbuttcap%
\pgfsetroundjoin%
\pgfsetlinewidth{0.501875pt}%
\definecolor{currentstroke}{rgb}{0.827451,0.827451,0.827451}%
\pgfsetstrokecolor{currentstroke}%
\pgfsetdash{}{0pt}%
\pgfpathmoveto{\pgfqpoint{3.991255in}{4.277116in}}%
\pgfpathlineto{\pgfqpoint{3.974287in}{4.277740in}}%
\pgfpathlineto{\pgfqpoint{3.958652in}{4.288510in}}%
\pgfpathlineto{\pgfqpoint{3.937530in}{4.321244in}}%
\pgfpathlineto{\pgfqpoint{3.919550in}{4.338580in}}%
\pgfpathlineto{\pgfqpoint{3.924249in}{4.351969in}}%
\pgfpathlineto{\pgfqpoint{3.920717in}{4.369680in}}%
\pgfpathlineto{\pgfqpoint{3.943252in}{4.376232in}}%
\pgfpathlineto{\pgfqpoint{3.944537in}{4.392475in}}%
\pgfpathlineto{\pgfqpoint{3.953718in}{4.395512in}}%
\pgfpathlineto{\pgfqpoint{3.947853in}{4.416069in}}%
\pgfpathlineto{\pgfqpoint{3.959456in}{4.445080in}}%
\pgfpathlineto{\pgfqpoint{3.983958in}{4.438094in}}%
\pgfpathlineto{\pgfqpoint{3.981897in}{4.467141in}}%
\pgfpathlineto{\pgfqpoint{4.003899in}{4.497309in}}%
\pgfpathlineto{\pgfqpoint{4.012789in}{4.491525in}}%
\pgfpathlineto{\pgfqpoint{4.024221in}{4.498243in}}%
\pgfpathlineto{\pgfqpoint{4.051378in}{4.529895in}}%
\pgfpathlineto{\pgfqpoint{4.056746in}{4.565383in}}%
\pgfpathlineto{\pgfqpoint{4.057501in}{4.592649in}}%
\pgfpathlineto{\pgfqpoint{4.063217in}{4.611303in}}%
\pgfpathlineto{\pgfqpoint{4.060283in}{4.631164in}}%
\pgfpathlineto{\pgfqpoint{4.052917in}{4.640224in}}%
\pgfpathlineto{\pgfqpoint{4.064326in}{4.650087in}}%
\pgfpathlineto{\pgfqpoint{4.080863in}{4.545897in}}%
\pgfpathlineto{\pgfqpoint{4.172033in}{4.560971in}}%
\pgfpathlineto{\pgfqpoint{4.181483in}{4.501503in}}%
\pgfpathlineto{\pgfqpoint{4.214489in}{4.540755in}}%
\pgfpathlineto{\pgfqpoint{4.222266in}{4.536841in}}%
\pgfpathlineto{\pgfqpoint{4.233707in}{4.559585in}}%
\pgfpathlineto{\pgfqpoint{4.249557in}{4.552454in}}%
\pgfpathlineto{\pgfqpoint{4.263377in}{4.553810in}}%
\pgfpathlineto{\pgfqpoint{4.272494in}{4.569620in}}%
\pgfpathlineto{\pgfqpoint{4.285729in}{4.578416in}}%
\pgfpathlineto{\pgfqpoint{4.304128in}{4.570670in}}%
\pgfpathlineto{\pgfqpoint{4.316235in}{4.573347in}}%
\pgfpathlineto{\pgfqpoint{4.333579in}{4.543330in}}%
\pgfpathlineto{\pgfqpoint{4.328563in}{4.520605in}}%
\pgfpathlineto{\pgfqpoint{4.283263in}{4.546191in}}%
\pgfpathlineto{\pgfqpoint{4.274778in}{4.525187in}}%
\pgfpathlineto{\pgfqpoint{4.277577in}{4.515523in}}%
\pgfpathlineto{\pgfqpoint{4.258297in}{4.482854in}}%
\pgfpathlineto{\pgfqpoint{4.249198in}{4.476356in}}%
\pgfpathlineto{\pgfqpoint{4.245088in}{4.461950in}}%
\pgfpathlineto{\pgfqpoint{4.229575in}{4.463457in}}%
\pgfpathlineto{\pgfqpoint{4.218435in}{4.424138in}}%
\pgfpathlineto{\pgfqpoint{4.212199in}{4.415115in}}%
\pgfpathlineto{\pgfqpoint{4.196175in}{4.418124in}}%
\pgfpathlineto{\pgfqpoint{4.179784in}{4.430525in}}%
\pgfpathlineto{\pgfqpoint{4.179217in}{4.411491in}}%
\pgfpathlineto{\pgfqpoint{4.163350in}{4.379502in}}%
\pgfpathlineto{\pgfqpoint{4.159403in}{4.356729in}}%
\pgfpathlineto{\pgfqpoint{4.147229in}{4.341562in}}%
\pgfpathlineto{\pgfqpoint{4.137999in}{4.317341in}}%
\pgfpathlineto{\pgfqpoint{4.137545in}{4.294919in}}%
\pgfpathlineto{\pgfqpoint{4.123291in}{4.290732in}}%
\pgfpathlineto{\pgfqpoint{4.107121in}{4.278039in}}%
\pgfpathlineto{\pgfqpoint{4.091551in}{4.275621in}}%
\pgfpathlineto{\pgfqpoint{4.087963in}{4.264878in}}%
\pgfpathlineto{\pgfqpoint{4.062876in}{4.253876in}}%
\pgfpathlineto{\pgfqpoint{4.048856in}{4.263250in}}%
\pgfpathlineto{\pgfqpoint{4.033186in}{4.245450in}}%
\pgfpathlineto{\pgfqpoint{4.006250in}{4.250697in}}%
\pgfpathlineto{\pgfqpoint{3.989729in}{4.269379in}}%
\pgfpathlineto{\pgfqpoint{3.991255in}{4.277116in}}%
\pgfusepath{stroke}%
\end{pgfscope}%
\begin{pgfscope}%
\pgfpathrectangle{\pgfqpoint{0.100000in}{2.493063in}}{\pgfqpoint{5.037500in}{3.427208in}}%
\pgfusepath{clip}%
\pgfsetbuttcap%
\pgfsetroundjoin%
\pgfsetlinewidth{0.501875pt}%
\definecolor{currentstroke}{rgb}{0.827451,0.827451,0.827451}%
\pgfsetstrokecolor{currentstroke}%
\pgfsetdash{}{0pt}%
\pgfpathmoveto{\pgfqpoint{4.585091in}{4.493533in}}%
\pgfpathlineto{\pgfqpoint{4.584403in}{4.472537in}}%
\pgfpathlineto{\pgfqpoint{4.576322in}{4.462215in}}%
\pgfpathlineto{\pgfqpoint{4.571165in}{4.439301in}}%
\pgfpathlineto{\pgfqpoint{4.547880in}{4.428739in}}%
\pgfpathlineto{\pgfqpoint{4.528343in}{4.425661in}}%
\pgfpathlineto{\pgfqpoint{4.534036in}{4.440731in}}%
\pgfpathlineto{\pgfqpoint{4.503955in}{4.453358in}}%
\pgfpathlineto{\pgfqpoint{4.479516in}{4.469169in}}%
\pgfpathlineto{\pgfqpoint{4.494646in}{4.492802in}}%
\pgfpathlineto{\pgfqpoint{4.482520in}{4.511173in}}%
\pgfpathlineto{\pgfqpoint{4.483948in}{4.527018in}}%
\pgfpathlineto{\pgfqpoint{4.475180in}{4.531368in}}%
\pgfpathlineto{\pgfqpoint{4.468066in}{4.557111in}}%
\pgfpathlineto{\pgfqpoint{4.484890in}{4.592009in}}%
\pgfpathlineto{\pgfqpoint{4.472250in}{4.597708in}}%
\pgfpathlineto{\pgfqpoint{4.468974in}{4.580491in}}%
\pgfpathlineto{\pgfqpoint{4.450931in}{4.575698in}}%
\pgfpathlineto{\pgfqpoint{4.452828in}{4.519052in}}%
\pgfpathlineto{\pgfqpoint{4.449529in}{4.500822in}}%
\pgfpathlineto{\pgfqpoint{4.458579in}{4.474839in}}%
\pgfpathlineto{\pgfqpoint{4.472502in}{4.462338in}}%
\pgfpathlineto{\pgfqpoint{4.466188in}{4.454488in}}%
\pgfpathlineto{\pgfqpoint{4.480400in}{4.443018in}}%
\pgfpathlineto{\pgfqpoint{4.485537in}{4.424401in}}%
\pgfpathlineto{\pgfqpoint{4.459488in}{4.439818in}}%
\pgfpathlineto{\pgfqpoint{4.443002in}{4.437815in}}%
\pgfpathlineto{\pgfqpoint{4.430202in}{4.453674in}}%
\pgfpathlineto{\pgfqpoint{4.417168in}{4.455219in}}%
\pgfpathlineto{\pgfqpoint{4.398644in}{4.447275in}}%
\pgfpathlineto{\pgfqpoint{4.391460in}{4.457174in}}%
\pgfpathlineto{\pgfqpoint{4.400902in}{4.477950in}}%
\pgfpathlineto{\pgfqpoint{4.405194in}{4.495246in}}%
\pgfpathlineto{\pgfqpoint{4.414687in}{4.509117in}}%
\pgfpathlineto{\pgfqpoint{4.400822in}{4.518480in}}%
\pgfpathlineto{\pgfqpoint{4.395161in}{4.510143in}}%
\pgfpathlineto{\pgfqpoint{4.381310in}{4.518606in}}%
\pgfpathlineto{\pgfqpoint{4.356785in}{4.524244in}}%
\pgfpathlineto{\pgfqpoint{4.359009in}{4.536642in}}%
\pgfpathlineto{\pgfqpoint{4.347892in}{4.543833in}}%
\pgfpathlineto{\pgfqpoint{4.333579in}{4.543330in}}%
\pgfpathlineto{\pgfqpoint{4.316235in}{4.573347in}}%
\pgfpathlineto{\pgfqpoint{4.304128in}{4.570670in}}%
\pgfpathlineto{\pgfqpoint{4.285729in}{4.578416in}}%
\pgfpathlineto{\pgfqpoint{4.272494in}{4.569620in}}%
\pgfpathlineto{\pgfqpoint{4.263377in}{4.553810in}}%
\pgfpathlineto{\pgfqpoint{4.249557in}{4.552454in}}%
\pgfpathlineto{\pgfqpoint{4.233707in}{4.559585in}}%
\pgfpathlineto{\pgfqpoint{4.222266in}{4.536841in}}%
\pgfpathlineto{\pgfqpoint{4.214489in}{4.540755in}}%
\pgfpathlineto{\pgfqpoint{4.181483in}{4.501503in}}%
\pgfpathlineto{\pgfqpoint{4.172033in}{4.560971in}}%
\pgfpathlineto{\pgfqpoint{4.292651in}{4.583327in}}%
\pgfpathlineto{\pgfqpoint{4.346722in}{4.593009in}}%
\pgfpathlineto{\pgfqpoint{4.425387in}{4.608852in}}%
\pgfpathlineto{\pgfqpoint{4.492802in}{4.623439in}}%
\pgfpathlineto{\pgfqpoint{4.502665in}{4.585425in}}%
\pgfpathlineto{\pgfqpoint{4.531394in}{4.483253in}}%
\pgfpathlineto{\pgfqpoint{4.585091in}{4.493533in}}%
\pgfusepath{stroke}%
\end{pgfscope}%
\begin{pgfscope}%
\pgfpathrectangle{\pgfqpoint{0.100000in}{2.493063in}}{\pgfqpoint{5.037500in}{3.427208in}}%
\pgfusepath{clip}%
\pgfsetbuttcap%
\pgfsetroundjoin%
\pgfsetlinewidth{0.501875pt}%
\definecolor{currentstroke}{rgb}{0.827451,0.827451,0.827451}%
\pgfsetstrokecolor{currentstroke}%
\pgfsetdash{}{0pt}%
\pgfpathmoveto{\pgfqpoint{2.157107in}{4.154398in}}%
\pgfpathlineto{\pgfqpoint{2.069106in}{4.162796in}}%
\pgfpathlineto{\pgfqpoint{1.977813in}{4.170887in}}%
\pgfpathlineto{\pgfqpoint{1.872319in}{4.181891in}}%
\pgfpathlineto{\pgfqpoint{1.715584in}{4.200532in}}%
\pgfpathlineto{\pgfqpoint{1.659981in}{4.209100in}}%
\pgfpathlineto{\pgfqpoint{1.516457in}{4.229943in}}%
\pgfpathlineto{\pgfqpoint{1.537414in}{4.361477in}}%
\pgfpathlineto{\pgfqpoint{1.537936in}{4.372117in}}%
\pgfpathlineto{\pgfqpoint{1.558125in}{4.499035in}}%
\pgfpathlineto{\pgfqpoint{1.573205in}{4.595697in}}%
\pgfpathlineto{\pgfqpoint{1.587436in}{4.685334in}}%
\pgfpathlineto{\pgfqpoint{1.684663in}{4.671377in}}%
\pgfpathlineto{\pgfqpoint{1.775315in}{4.658367in}}%
\pgfpathlineto{\pgfqpoint{1.941955in}{4.637799in}}%
\pgfpathlineto{\pgfqpoint{2.018405in}{4.630818in}}%
\pgfpathlineto{\pgfqpoint{2.139519in}{4.619067in}}%
\pgfpathlineto{\pgfqpoint{2.191913in}{4.614814in}}%
\pgfpathlineto{\pgfqpoint{2.182662in}{4.500100in}}%
\pgfpathlineto{\pgfqpoint{2.174305in}{4.389649in}}%
\pgfpathlineto{\pgfqpoint{2.167579in}{4.299715in}}%
\pgfpathlineto{\pgfqpoint{2.157107in}{4.154398in}}%
\pgfusepath{stroke}%
\end{pgfscope}%
\begin{pgfscope}%
\pgfpathrectangle{\pgfqpoint{0.100000in}{2.493063in}}{\pgfqpoint{5.037500in}{3.427208in}}%
\pgfusepath{clip}%
\pgfsetbuttcap%
\pgfsetroundjoin%
\pgfsetlinewidth{0.501875pt}%
\definecolor{currentstroke}{rgb}{0.827451,0.827451,0.827451}%
\pgfsetstrokecolor{currentstroke}%
\pgfsetdash{}{0pt}%
\pgfpathmoveto{\pgfqpoint{3.321282in}{4.089187in}}%
\pgfpathlineto{\pgfqpoint{3.324248in}{4.102485in}}%
\pgfpathlineto{\pgfqpoint{3.339624in}{4.099492in}}%
\pgfpathlineto{\pgfqpoint{3.344010in}{4.146681in}}%
\pgfpathlineto{\pgfqpoint{3.339802in}{4.155199in}}%
\pgfpathlineto{\pgfqpoint{3.354799in}{4.174166in}}%
\pgfpathlineto{\pgfqpoint{3.366827in}{4.172910in}}%
\pgfpathlineto{\pgfqpoint{3.403283in}{4.160434in}}%
\pgfpathlineto{\pgfqpoint{3.407555in}{4.170534in}}%
\pgfpathlineto{\pgfqpoint{3.398271in}{4.186030in}}%
\pgfpathlineto{\pgfqpoint{3.401404in}{4.198255in}}%
\pgfpathlineto{\pgfqpoint{3.435402in}{4.209497in}}%
\pgfpathlineto{\pgfqpoint{3.427288in}{4.230184in}}%
\pgfpathlineto{\pgfqpoint{3.437265in}{4.247692in}}%
\pgfpathlineto{\pgfqpoint{3.455836in}{4.257519in}}%
\pgfpathlineto{\pgfqpoint{3.494803in}{4.267250in}}%
\pgfpathlineto{\pgfqpoint{3.519389in}{4.252561in}}%
\pgfpathlineto{\pgfqpoint{3.526847in}{4.266886in}}%
\pgfpathlineto{\pgfqpoint{3.547262in}{4.276892in}}%
\pgfpathlineto{\pgfqpoint{3.552531in}{4.268067in}}%
\pgfpathlineto{\pgfqpoint{3.573500in}{4.275078in}}%
\pgfpathlineto{\pgfqpoint{3.572178in}{4.287078in}}%
\pgfpathlineto{\pgfqpoint{3.591032in}{4.300810in}}%
\pgfpathlineto{\pgfqpoint{3.602160in}{4.286550in}}%
\pgfpathlineto{\pgfqpoint{3.616798in}{4.285115in}}%
\pgfpathlineto{\pgfqpoint{3.626527in}{4.294409in}}%
\pgfpathlineto{\pgfqpoint{3.625445in}{4.307626in}}%
\pgfpathlineto{\pgfqpoint{3.633714in}{4.320750in}}%
\pgfpathlineto{\pgfqpoint{3.644795in}{4.323592in}}%
\pgfpathlineto{\pgfqpoint{3.649244in}{4.340922in}}%
\pgfpathlineto{\pgfqpoint{3.665360in}{4.355902in}}%
\pgfpathlineto{\pgfqpoint{3.660509in}{4.370813in}}%
\pgfpathlineto{\pgfqpoint{3.676188in}{4.378237in}}%
\pgfpathlineto{\pgfqpoint{3.686663in}{4.373671in}}%
\pgfpathlineto{\pgfqpoint{3.702135in}{4.385269in}}%
\pgfpathlineto{\pgfqpoint{3.715962in}{4.388292in}}%
\pgfpathlineto{\pgfqpoint{3.716011in}{4.400323in}}%
\pgfpathlineto{\pgfqpoint{3.706322in}{4.417017in}}%
\pgfpathlineto{\pgfqpoint{3.712630in}{4.423743in}}%
\pgfpathlineto{\pgfqpoint{3.732447in}{4.422252in}}%
\pgfpathlineto{\pgfqpoint{3.745048in}{4.428530in}}%
\pgfpathlineto{\pgfqpoint{3.758047in}{4.418723in}}%
\pgfpathlineto{\pgfqpoint{3.768306in}{4.396588in}}%
\pgfpathlineto{\pgfqpoint{3.802320in}{4.392848in}}%
\pgfpathlineto{\pgfqpoint{3.820049in}{4.380664in}}%
\pgfpathlineto{\pgfqpoint{3.832524in}{4.390962in}}%
\pgfpathlineto{\pgfqpoint{3.855079in}{4.382307in}}%
\pgfpathlineto{\pgfqpoint{3.867848in}{4.386608in}}%
\pgfpathlineto{\pgfqpoint{3.876773in}{4.398839in}}%
\pgfpathlineto{\pgfqpoint{3.890398in}{4.401178in}}%
\pgfpathlineto{\pgfqpoint{3.893938in}{4.387880in}}%
\pgfpathlineto{\pgfqpoint{3.909006in}{4.382795in}}%
\pgfpathlineto{\pgfqpoint{3.920717in}{4.369680in}}%
\pgfpathlineto{\pgfqpoint{3.924249in}{4.351969in}}%
\pgfpathlineto{\pgfqpoint{3.919550in}{4.338580in}}%
\pgfpathlineto{\pgfqpoint{3.937530in}{4.321244in}}%
\pgfpathlineto{\pgfqpoint{3.958652in}{4.288510in}}%
\pgfpathlineto{\pgfqpoint{3.974287in}{4.277740in}}%
\pgfpathlineto{\pgfqpoint{3.991255in}{4.277116in}}%
\pgfpathlineto{\pgfqpoint{3.959924in}{4.241166in}}%
\pgfpathlineto{\pgfqpoint{3.929084in}{4.219406in}}%
\pgfpathlineto{\pgfqpoint{3.917734in}{4.202126in}}%
\pgfpathlineto{\pgfqpoint{3.917926in}{4.192748in}}%
\pgfpathlineto{\pgfqpoint{3.901222in}{4.185560in}}%
\pgfpathlineto{\pgfqpoint{3.896435in}{4.172036in}}%
\pgfpathlineto{\pgfqpoint{3.861637in}{4.158413in}}%
\pgfpathlineto{\pgfqpoint{3.849302in}{4.149566in}}%
\pgfpathlineto{\pgfqpoint{3.847633in}{4.147678in}}%
\pgfpathlineto{\pgfqpoint{3.747536in}{4.138186in}}%
\pgfpathlineto{\pgfqpoint{3.687088in}{4.133103in}}%
\pgfpathlineto{\pgfqpoint{3.587896in}{4.127489in}}%
\pgfpathlineto{\pgfqpoint{3.464320in}{4.115264in}}%
\pgfpathlineto{\pgfqpoint{3.443904in}{4.118092in}}%
\pgfpathlineto{\pgfqpoint{3.448152in}{4.097250in}}%
\pgfpathlineto{\pgfqpoint{3.321282in}{4.089187in}}%
\pgfusepath{stroke}%
\end{pgfscope}%
\begin{pgfscope}%
\pgfpathrectangle{\pgfqpoint{0.100000in}{2.493063in}}{\pgfqpoint{5.037500in}{3.427208in}}%
\pgfusepath{clip}%
\pgfsetbuttcap%
\pgfsetroundjoin%
\pgfsetlinewidth{0.501875pt}%
\definecolor{currentstroke}{rgb}{0.827451,0.827451,0.827451}%
\pgfsetstrokecolor{currentstroke}%
\pgfsetdash{}{0pt}%
\pgfpathmoveto{\pgfqpoint{2.157107in}{4.154398in}}%
\pgfpathlineto{\pgfqpoint{2.167579in}{4.299715in}}%
\pgfpathlineto{\pgfqpoint{2.174305in}{4.389649in}}%
\pgfpathlineto{\pgfqpoint{2.182662in}{4.500100in}}%
\pgfpathlineto{\pgfqpoint{2.298517in}{4.491965in}}%
\pgfpathlineto{\pgfqpoint{2.445651in}{4.483928in}}%
\pgfpathlineto{\pgfqpoint{2.545721in}{4.480088in}}%
\pgfpathlineto{\pgfqpoint{2.645206in}{4.477040in}}%
\pgfpathlineto{\pgfqpoint{2.735250in}{4.475647in}}%
\pgfpathlineto{\pgfqpoint{2.776895in}{4.476079in}}%
\pgfpathlineto{\pgfqpoint{2.795226in}{4.461120in}}%
\pgfpathlineto{\pgfqpoint{2.809584in}{4.464135in}}%
\pgfpathlineto{\pgfqpoint{2.810034in}{4.451087in}}%
\pgfpathlineto{\pgfqpoint{2.799370in}{4.428517in}}%
\pgfpathlineto{\pgfqpoint{2.800527in}{4.414257in}}%
\pgfpathlineto{\pgfqpoint{2.810737in}{4.404854in}}%
\pgfpathlineto{\pgfqpoint{2.820114in}{4.385262in}}%
\pgfpathlineto{\pgfqpoint{2.839138in}{4.374013in}}%
\pgfpathlineto{\pgfqpoint{2.838503in}{4.300162in}}%
\pgfpathlineto{\pgfqpoint{2.839063in}{4.130335in}}%
\pgfpathlineto{\pgfqpoint{2.711550in}{4.130923in}}%
\pgfpathlineto{\pgfqpoint{2.577276in}{4.133273in}}%
\pgfpathlineto{\pgfqpoint{2.478429in}{4.136760in}}%
\pgfpathlineto{\pgfqpoint{2.395919in}{4.139944in}}%
\pgfpathlineto{\pgfqpoint{2.256916in}{4.148140in}}%
\pgfpathlineto{\pgfqpoint{2.157107in}{4.154398in}}%
\pgfusepath{stroke}%
\end{pgfscope}%
\begin{pgfscope}%
\pgfpathrectangle{\pgfqpoint{0.100000in}{2.493063in}}{\pgfqpoint{5.037500in}{3.427208in}}%
\pgfusepath{clip}%
\pgfsetbuttcap%
\pgfsetroundjoin%
\pgfsetlinewidth{0.501875pt}%
\definecolor{currentstroke}{rgb}{0.827451,0.827451,0.827451}%
\pgfsetstrokecolor{currentstroke}%
\pgfsetdash{}{0pt}%
\pgfpathmoveto{\pgfqpoint{4.033298in}{4.173639in}}%
\pgfpathlineto{\pgfqpoint{3.892070in}{4.153801in}}%
\pgfpathlineto{\pgfqpoint{3.849302in}{4.149566in}}%
\pgfpathlineto{\pgfqpoint{3.861637in}{4.158413in}}%
\pgfpathlineto{\pgfqpoint{3.896435in}{4.172036in}}%
\pgfpathlineto{\pgfqpoint{3.901222in}{4.185560in}}%
\pgfpathlineto{\pgfqpoint{3.917926in}{4.192748in}}%
\pgfpathlineto{\pgfqpoint{3.917734in}{4.202126in}}%
\pgfpathlineto{\pgfqpoint{3.929084in}{4.219406in}}%
\pgfpathlineto{\pgfqpoint{3.959924in}{4.241166in}}%
\pgfpathlineto{\pgfqpoint{3.991255in}{4.277116in}}%
\pgfpathlineto{\pgfqpoint{3.989729in}{4.269379in}}%
\pgfpathlineto{\pgfqpoint{4.006250in}{4.250697in}}%
\pgfpathlineto{\pgfqpoint{4.033186in}{4.245450in}}%
\pgfpathlineto{\pgfqpoint{4.048856in}{4.263250in}}%
\pgfpathlineto{\pgfqpoint{4.062876in}{4.253876in}}%
\pgfpathlineto{\pgfqpoint{4.087963in}{4.264878in}}%
\pgfpathlineto{\pgfqpoint{4.091551in}{4.275621in}}%
\pgfpathlineto{\pgfqpoint{4.107121in}{4.278039in}}%
\pgfpathlineto{\pgfqpoint{4.123291in}{4.290732in}}%
\pgfpathlineto{\pgfqpoint{4.137545in}{4.294919in}}%
\pgfpathlineto{\pgfqpoint{4.137999in}{4.317341in}}%
\pgfpathlineto{\pgfqpoint{4.147229in}{4.341562in}}%
\pgfpathlineto{\pgfqpoint{4.159403in}{4.356729in}}%
\pgfpathlineto{\pgfqpoint{4.163350in}{4.379502in}}%
\pgfpathlineto{\pgfqpoint{4.179217in}{4.411491in}}%
\pgfpathlineto{\pgfqpoint{4.179784in}{4.430525in}}%
\pgfpathlineto{\pgfqpoint{4.196175in}{4.418124in}}%
\pgfpathlineto{\pgfqpoint{4.212199in}{4.415115in}}%
\pgfpathlineto{\pgfqpoint{4.218435in}{4.424138in}}%
\pgfpathlineto{\pgfqpoint{4.229575in}{4.463457in}}%
\pgfpathlineto{\pgfqpoint{4.245088in}{4.461950in}}%
\pgfpathlineto{\pgfqpoint{4.249198in}{4.476356in}}%
\pgfpathlineto{\pgfqpoint{4.258297in}{4.482854in}}%
\pgfpathlineto{\pgfqpoint{4.277577in}{4.515523in}}%
\pgfpathlineto{\pgfqpoint{4.274778in}{4.525187in}}%
\pgfpathlineto{\pgfqpoint{4.283263in}{4.546191in}}%
\pgfpathlineto{\pgfqpoint{4.328563in}{4.520605in}}%
\pgfpathlineto{\pgfqpoint{4.333579in}{4.543330in}}%
\pgfpathlineto{\pgfqpoint{4.347892in}{4.543833in}}%
\pgfpathlineto{\pgfqpoint{4.359009in}{4.536642in}}%
\pgfpathlineto{\pgfqpoint{4.356785in}{4.524244in}}%
\pgfpathlineto{\pgfqpoint{4.381310in}{4.518606in}}%
\pgfpathlineto{\pgfqpoint{4.395161in}{4.510143in}}%
\pgfpathlineto{\pgfqpoint{4.405194in}{4.495246in}}%
\pgfpathlineto{\pgfqpoint{4.400902in}{4.477950in}}%
\pgfpathlineto{\pgfqpoint{4.392234in}{4.476533in}}%
\pgfpathlineto{\pgfqpoint{4.387208in}{4.450435in}}%
\pgfpathlineto{\pgfqpoint{4.398219in}{4.440231in}}%
\pgfpathlineto{\pgfqpoint{4.413719in}{4.448479in}}%
\pgfpathlineto{\pgfqpoint{4.428102in}{4.431062in}}%
\pgfpathlineto{\pgfqpoint{4.460236in}{4.427931in}}%
\pgfpathlineto{\pgfqpoint{4.465772in}{4.417912in}}%
\pgfpathlineto{\pgfqpoint{4.495510in}{4.408168in}}%
\pgfpathlineto{\pgfqpoint{4.491840in}{4.396670in}}%
\pgfpathlineto{\pgfqpoint{4.495709in}{4.364173in}}%
\pgfpathlineto{\pgfqpoint{4.507712in}{4.351902in}}%
\pgfpathlineto{\pgfqpoint{4.490000in}{4.346722in}}%
\pgfpathlineto{\pgfqpoint{4.492348in}{4.332849in}}%
\pgfpathlineto{\pgfqpoint{4.511242in}{4.321087in}}%
\pgfpathlineto{\pgfqpoint{4.512946in}{4.309484in}}%
\pgfpathlineto{\pgfqpoint{4.502278in}{4.300765in}}%
\pgfpathlineto{\pgfqpoint{4.480755in}{4.321421in}}%
\pgfpathlineto{\pgfqpoint{4.478627in}{4.306350in}}%
\pgfpathlineto{\pgfqpoint{4.496645in}{4.299183in}}%
\pgfpathlineto{\pgfqpoint{4.498448in}{4.291772in}}%
\pgfpathlineto{\pgfqpoint{4.523162in}{4.301567in}}%
\pgfpathlineto{\pgfqpoint{4.542100in}{4.304158in}}%
\pgfpathlineto{\pgfqpoint{4.561500in}{4.265016in}}%
\pgfpathlineto{\pgfqpoint{4.559336in}{4.264593in}}%
\pgfpathlineto{\pgfqpoint{4.550569in}{4.262781in}}%
\pgfpathlineto{\pgfqpoint{4.547986in}{4.262241in}}%
\pgfpathlineto{\pgfqpoint{4.546278in}{4.261907in}}%
\pgfpathlineto{\pgfqpoint{4.430773in}{4.237948in}}%
\pgfpathlineto{\pgfqpoint{4.327461in}{4.216802in}}%
\pgfpathlineto{\pgfqpoint{4.184620in}{4.192195in}}%
\pgfpathlineto{\pgfqpoint{4.063306in}{4.176172in}}%
\pgfpathlineto{\pgfqpoint{4.033298in}{4.173639in}}%
\pgfusepath{stroke}%
\end{pgfscope}%
\begin{pgfscope}%
\pgfpathrectangle{\pgfqpoint{0.100000in}{2.493063in}}{\pgfqpoint{5.037500in}{3.427208in}}%
\pgfusepath{clip}%
\pgfsetbuttcap%
\pgfsetroundjoin%
\pgfsetlinewidth{0.501875pt}%
\definecolor{currentstroke}{rgb}{0.827451,0.827451,0.827451}%
\pgfsetstrokecolor{currentstroke}%
\pgfsetdash{}{0pt}%
\pgfpathmoveto{\pgfqpoint{4.547880in}{4.428739in}}%
\pgfpathlineto{\pgfqpoint{4.571165in}{4.439301in}}%
\pgfpathlineto{\pgfqpoint{4.552570in}{4.384904in}}%
\pgfpathlineto{\pgfqpoint{4.540273in}{4.356098in}}%
\pgfpathlineto{\pgfqpoint{4.531414in}{4.372400in}}%
\pgfpathlineto{\pgfqpoint{4.547168in}{4.411428in}}%
\pgfpathlineto{\pgfqpoint{4.547880in}{4.428739in}}%
\pgfusepath{stroke}%
\end{pgfscope}%
\begin{pgfscope}%
\pgfpathrectangle{\pgfqpoint{0.100000in}{2.493063in}}{\pgfqpoint{5.037500in}{3.427208in}}%
\pgfusepath{clip}%
\pgfsetbuttcap%
\pgfsetroundjoin%
\pgfsetlinewidth{0.501875pt}%
\definecolor{currentstroke}{rgb}{0.827451,0.827451,0.827451}%
\pgfsetstrokecolor{currentstroke}%
\pgfsetdash{}{0pt}%
\pgfpathmoveto{\pgfqpoint{3.321282in}{4.089187in}}%
\pgfpathlineto{\pgfqpoint{3.315658in}{4.088373in}}%
\pgfpathlineto{\pgfqpoint{3.310355in}{4.088001in}}%
\pgfpathlineto{\pgfqpoint{3.312625in}{4.071720in}}%
\pgfpathlineto{\pgfqpoint{3.298944in}{4.055308in}}%
\pgfpathlineto{\pgfqpoint{3.296248in}{4.029628in}}%
\pgfpathlineto{\pgfqpoint{3.235087in}{4.025104in}}%
\pgfpathlineto{\pgfqpoint{3.240418in}{4.037162in}}%
\pgfpathlineto{\pgfqpoint{3.262470in}{4.059183in}}%
\pgfpathlineto{\pgfqpoint{3.263080in}{4.071948in}}%
\pgfpathlineto{\pgfqpoint{3.253318in}{4.084033in}}%
\pgfpathlineto{\pgfqpoint{3.162287in}{4.079219in}}%
\pgfpathlineto{\pgfqpoint{3.003135in}{4.074167in}}%
\pgfpathlineto{\pgfqpoint{2.909997in}{4.072469in}}%
\pgfpathlineto{\pgfqpoint{2.839596in}{4.071836in}}%
\pgfpathlineto{\pgfqpoint{2.839063in}{4.130335in}}%
\pgfpathlineto{\pgfqpoint{2.838503in}{4.300162in}}%
\pgfpathlineto{\pgfqpoint{2.839138in}{4.374013in}}%
\pgfpathlineto{\pgfqpoint{2.820114in}{4.385262in}}%
\pgfpathlineto{\pgfqpoint{2.810737in}{4.404854in}}%
\pgfpathlineto{\pgfqpoint{2.800527in}{4.414257in}}%
\pgfpathlineto{\pgfqpoint{2.799370in}{4.428517in}}%
\pgfpathlineto{\pgfqpoint{2.810034in}{4.451087in}}%
\pgfpathlineto{\pgfqpoint{2.809584in}{4.464135in}}%
\pgfpathlineto{\pgfqpoint{2.795226in}{4.461120in}}%
\pgfpathlineto{\pgfqpoint{2.776895in}{4.476079in}}%
\pgfpathlineto{\pgfqpoint{2.762201in}{4.502337in}}%
\pgfpathlineto{\pgfqpoint{2.749875in}{4.514454in}}%
\pgfpathlineto{\pgfqpoint{2.736996in}{4.544227in}}%
\pgfpathlineto{\pgfqpoint{2.870765in}{4.542140in}}%
\pgfpathlineto{\pgfqpoint{3.003734in}{4.546521in}}%
\pgfpathlineto{\pgfqpoint{3.088992in}{4.551439in}}%
\pgfpathlineto{\pgfqpoint{3.115663in}{4.525025in}}%
\pgfpathlineto{\pgfqpoint{3.110789in}{4.497479in}}%
\pgfpathlineto{\pgfqpoint{3.117900in}{4.466465in}}%
\pgfpathlineto{\pgfqpoint{3.125789in}{4.450826in}}%
\pgfpathlineto{\pgfqpoint{3.150959in}{4.429315in}}%
\pgfpathlineto{\pgfqpoint{3.156465in}{4.419812in}}%
\pgfpathlineto{\pgfqpoint{3.184022in}{4.398399in}}%
\pgfpathlineto{\pgfqpoint{3.188016in}{4.374377in}}%
\pgfpathlineto{\pgfqpoint{3.195930in}{4.356657in}}%
\pgfpathlineto{\pgfqpoint{3.214629in}{4.366979in}}%
\pgfpathlineto{\pgfqpoint{3.230042in}{4.362258in}}%
\pgfpathlineto{\pgfqpoint{3.241679in}{4.349731in}}%
\pgfpathlineto{\pgfqpoint{3.237285in}{4.327559in}}%
\pgfpathlineto{\pgfqpoint{3.223001in}{4.298571in}}%
\pgfpathlineto{\pgfqpoint{3.223222in}{4.283474in}}%
\pgfpathlineto{\pgfqpoint{3.233897in}{4.271095in}}%
\pgfpathlineto{\pgfqpoint{3.256870in}{4.254698in}}%
\pgfpathlineto{\pgfqpoint{3.283323in}{4.241390in}}%
\pgfpathlineto{\pgfqpoint{3.289597in}{4.230734in}}%
\pgfpathlineto{\pgfqpoint{3.303126in}{4.225694in}}%
\pgfpathlineto{\pgfqpoint{3.303664in}{4.211552in}}%
\pgfpathlineto{\pgfqpoint{3.313676in}{4.193588in}}%
\pgfpathlineto{\pgfqpoint{3.306729in}{4.177628in}}%
\pgfpathlineto{\pgfqpoint{3.320328in}{4.152236in}}%
\pgfpathlineto{\pgfqpoint{3.330964in}{4.157278in}}%
\pgfpathlineto{\pgfqpoint{3.344010in}{4.146681in}}%
\pgfpathlineto{\pgfqpoint{3.339624in}{4.099492in}}%
\pgfpathlineto{\pgfqpoint{3.324248in}{4.102485in}}%
\pgfpathlineto{\pgfqpoint{3.321282in}{4.089187in}}%
\pgfusepath{stroke}%
\end{pgfscope}%
\begin{pgfscope}%
\pgfpathrectangle{\pgfqpoint{0.100000in}{2.493063in}}{\pgfqpoint{5.037500in}{3.427208in}}%
\pgfusepath{clip}%
\pgfsetbuttcap%
\pgfsetroundjoin%
\pgfsetlinewidth{0.501875pt}%
\definecolor{currentstroke}{rgb}{0.827451,0.827451,0.827451}%
\pgfsetstrokecolor{currentstroke}%
\pgfsetdash{}{0pt}%
\pgfpathmoveto{\pgfqpoint{0.963137in}{4.099360in}}%
\pgfpathlineto{\pgfqpoint{0.974271in}{4.123066in}}%
\pgfpathlineto{\pgfqpoint{0.971224in}{4.158649in}}%
\pgfpathlineto{\pgfqpoint{0.976556in}{4.168769in}}%
\pgfpathlineto{\pgfqpoint{0.977083in}{4.213092in}}%
\pgfpathlineto{\pgfqpoint{0.982136in}{4.225855in}}%
\pgfpathlineto{\pgfqpoint{0.999941in}{4.227841in}}%
\pgfpathlineto{\pgfqpoint{1.016464in}{4.222175in}}%
\pgfpathlineto{\pgfqpoint{1.023660in}{4.206570in}}%
\pgfpathlineto{\pgfqpoint{1.033749in}{4.207168in}}%
\pgfpathlineto{\pgfqpoint{1.046284in}{4.225043in}}%
\pgfpathlineto{\pgfqpoint{1.064402in}{4.313160in}}%
\pgfpathlineto{\pgfqpoint{1.167500in}{4.291874in}}%
\pgfpathlineto{\pgfqpoint{1.227336in}{4.280074in}}%
\pgfpathlineto{\pgfqpoint{1.386229in}{4.252143in}}%
\pgfpathlineto{\pgfqpoint{1.430188in}{4.243315in}}%
\pgfpathlineto{\pgfqpoint{1.516457in}{4.229943in}}%
\pgfpathlineto{\pgfqpoint{1.498768in}{4.116047in}}%
\pgfpathlineto{\pgfqpoint{1.480371in}{3.997241in}}%
\pgfpathlineto{\pgfqpoint{1.459184in}{3.863582in}}%
\pgfpathlineto{\pgfqpoint{1.435357in}{3.710155in}}%
\pgfpathlineto{\pgfqpoint{1.416109in}{3.584143in}}%
\pgfpathlineto{\pgfqpoint{1.278236in}{3.606036in}}%
\pgfpathlineto{\pgfqpoint{1.217659in}{3.616411in}}%
\pgfpathlineto{\pgfqpoint{1.190598in}{3.632612in}}%
\pgfpathlineto{\pgfqpoint{1.014159in}{3.739067in}}%
\pgfpathlineto{\pgfqpoint{0.881749in}{3.819855in}}%
\pgfpathlineto{\pgfqpoint{0.886159in}{3.834174in}}%
\pgfpathlineto{\pgfqpoint{0.897097in}{3.844192in}}%
\pgfpathlineto{\pgfqpoint{0.908473in}{3.842335in}}%
\pgfpathlineto{\pgfqpoint{0.924974in}{3.852858in}}%
\pgfpathlineto{\pgfqpoint{0.927572in}{3.867979in}}%
\pgfpathlineto{\pgfqpoint{0.907448in}{3.886322in}}%
\pgfpathlineto{\pgfqpoint{0.921760in}{3.921535in}}%
\pgfpathlineto{\pgfqpoint{0.936150in}{3.935078in}}%
\pgfpathlineto{\pgfqpoint{0.942979in}{3.951143in}}%
\pgfpathlineto{\pgfqpoint{0.947187in}{3.980602in}}%
\pgfpathlineto{\pgfqpoint{0.960679in}{3.993943in}}%
\pgfpathlineto{\pgfqpoint{0.988999in}{4.007268in}}%
\pgfpathlineto{\pgfqpoint{0.988373in}{4.019008in}}%
\pgfpathlineto{\pgfqpoint{0.972557in}{4.033579in}}%
\pgfpathlineto{\pgfqpoint{0.970403in}{4.063615in}}%
\pgfpathlineto{\pgfqpoint{0.959513in}{4.085573in}}%
\pgfpathlineto{\pgfqpoint{0.963137in}{4.099360in}}%
\pgfusepath{stroke}%
\end{pgfscope}%
\begin{pgfscope}%
\pgfpathrectangle{\pgfqpoint{0.100000in}{2.493063in}}{\pgfqpoint{5.037500in}{3.427208in}}%
\pgfusepath{clip}%
\pgfsetbuttcap%
\pgfsetroundjoin%
\pgfsetlinewidth{0.501875pt}%
\definecolor{currentstroke}{rgb}{0.827451,0.827451,0.827451}%
\pgfsetstrokecolor{currentstroke}%
\pgfsetdash{}{0pt}%
\pgfpathmoveto{\pgfqpoint{2.854521in}{3.742061in}}%
\pgfpathlineto{\pgfqpoint{2.828963in}{3.749966in}}%
\pgfpathlineto{\pgfqpoint{2.796214in}{3.775046in}}%
\pgfpathlineto{\pgfqpoint{2.781668in}{3.780436in}}%
\pgfpathlineto{\pgfqpoint{2.772424in}{3.769604in}}%
\pgfpathlineto{\pgfqpoint{2.760751in}{3.769057in}}%
\pgfpathlineto{\pgfqpoint{2.745983in}{3.778253in}}%
\pgfpathlineto{\pgfqpoint{2.722809in}{3.766470in}}%
\pgfpathlineto{\pgfqpoint{2.714559in}{3.772273in}}%
\pgfpathlineto{\pgfqpoint{2.691756in}{3.758526in}}%
\pgfpathlineto{\pgfqpoint{2.667699in}{3.761062in}}%
\pgfpathlineto{\pgfqpoint{2.649143in}{3.769996in}}%
\pgfpathlineto{\pgfqpoint{2.636148in}{3.766633in}}%
\pgfpathlineto{\pgfqpoint{2.615348in}{3.779267in}}%
\pgfpathlineto{\pgfqpoint{2.603777in}{3.756985in}}%
\pgfpathlineto{\pgfqpoint{2.591944in}{3.776148in}}%
\pgfpathlineto{\pgfqpoint{2.568572in}{3.768713in}}%
\pgfpathlineto{\pgfqpoint{2.548129in}{3.786915in}}%
\pgfpathlineto{\pgfqpoint{2.530266in}{3.772244in}}%
\pgfpathlineto{\pgfqpoint{2.521320in}{3.785726in}}%
\pgfpathlineto{\pgfqpoint{2.508425in}{3.790104in}}%
\pgfpathlineto{\pgfqpoint{2.500572in}{3.803102in}}%
\pgfpathlineto{\pgfqpoint{2.483707in}{3.806801in}}%
\pgfpathlineto{\pgfqpoint{2.473975in}{3.797018in}}%
\pgfpathlineto{\pgfqpoint{2.457440in}{3.809697in}}%
\pgfpathlineto{\pgfqpoint{2.449732in}{3.806802in}}%
\pgfpathlineto{\pgfqpoint{2.422328in}{3.817076in}}%
\pgfpathlineto{\pgfqpoint{2.405164in}{3.818222in}}%
\pgfpathlineto{\pgfqpoint{2.397482in}{3.840041in}}%
\pgfpathlineto{\pgfqpoint{2.383723in}{3.837323in}}%
\pgfpathlineto{\pgfqpoint{2.367963in}{3.842707in}}%
\pgfpathlineto{\pgfqpoint{2.357586in}{3.839596in}}%
\pgfpathlineto{\pgfqpoint{2.335330in}{3.864101in}}%
\pgfpathlineto{\pgfqpoint{2.329128in}{3.862498in}}%
\pgfpathlineto{\pgfqpoint{2.334761in}{3.961873in}}%
\pgfpathlineto{\pgfqpoint{2.340896in}{4.084874in}}%
\pgfpathlineto{\pgfqpoint{2.240180in}{4.090507in}}%
\pgfpathlineto{\pgfqpoint{2.140792in}{4.098021in}}%
\pgfpathlineto{\pgfqpoint{2.064008in}{4.104689in}}%
\pgfpathlineto{\pgfqpoint{2.069106in}{4.162796in}}%
\pgfpathlineto{\pgfqpoint{2.157107in}{4.154398in}}%
\pgfpathlineto{\pgfqpoint{2.256916in}{4.148140in}}%
\pgfpathlineto{\pgfqpoint{2.395919in}{4.139944in}}%
\pgfpathlineto{\pgfqpoint{2.478429in}{4.136760in}}%
\pgfpathlineto{\pgfqpoint{2.577276in}{4.133273in}}%
\pgfpathlineto{\pgfqpoint{2.711550in}{4.130923in}}%
\pgfpathlineto{\pgfqpoint{2.839063in}{4.130335in}}%
\pgfpathlineto{\pgfqpoint{2.839596in}{4.071836in}}%
\pgfpathlineto{\pgfqpoint{2.846754in}{4.027764in}}%
\pgfpathlineto{\pgfqpoint{2.857873in}{3.946362in}}%
\pgfpathlineto{\pgfqpoint{2.855579in}{3.807350in}}%
\pgfpathlineto{\pgfqpoint{2.854521in}{3.742061in}}%
\pgfusepath{stroke}%
\end{pgfscope}%
\begin{pgfscope}%
\pgfpathrectangle{\pgfqpoint{0.100000in}{2.493063in}}{\pgfqpoint{5.037500in}{3.427208in}}%
\pgfusepath{clip}%
\pgfsetbuttcap%
\pgfsetroundjoin%
\pgfsetlinewidth{0.501875pt}%
\definecolor{currentstroke}{rgb}{0.827451,0.827451,0.827451}%
\pgfsetstrokecolor{currentstroke}%
\pgfsetdash{}{0pt}%
\pgfpathmoveto{\pgfqpoint{3.811487in}{3.957948in}}%
\pgfpathlineto{\pgfqpoint{3.811562in}{3.983705in}}%
\pgfpathlineto{\pgfqpoint{3.833945in}{3.993596in}}%
\pgfpathlineto{\pgfqpoint{3.834884in}{4.009407in}}%
\pgfpathlineto{\pgfqpoint{3.854924in}{4.028722in}}%
\pgfpathlineto{\pgfqpoint{3.879975in}{4.032678in}}%
\pgfpathlineto{\pgfqpoint{3.901997in}{4.053964in}}%
\pgfpathlineto{\pgfqpoint{3.924997in}{4.063478in}}%
\pgfpathlineto{\pgfqpoint{3.942259in}{4.092040in}}%
\pgfpathlineto{\pgfqpoint{3.957980in}{4.089966in}}%
\pgfpathlineto{\pgfqpoint{3.991206in}{4.115991in}}%
\pgfpathlineto{\pgfqpoint{4.008752in}{4.116482in}}%
\pgfpathlineto{\pgfqpoint{4.016161in}{4.136311in}}%
\pgfpathlineto{\pgfqpoint{4.030107in}{4.150117in}}%
\pgfpathlineto{\pgfqpoint{4.033298in}{4.173639in}}%
\pgfpathlineto{\pgfqpoint{4.063306in}{4.176172in}}%
\pgfpathlineto{\pgfqpoint{4.184620in}{4.192195in}}%
\pgfpathlineto{\pgfqpoint{4.327461in}{4.216802in}}%
\pgfpathlineto{\pgfqpoint{4.430773in}{4.237948in}}%
\pgfpathlineto{\pgfqpoint{4.546278in}{4.261907in}}%
\pgfpathlineto{\pgfqpoint{4.579284in}{4.216592in}}%
\pgfpathlineto{\pgfqpoint{4.561406in}{4.219485in}}%
\pgfpathlineto{\pgfqpoint{4.538456in}{4.213894in}}%
\pgfpathlineto{\pgfqpoint{4.515842in}{4.190800in}}%
\pgfpathlineto{\pgfqpoint{4.499590in}{4.192461in}}%
\pgfpathlineto{\pgfqpoint{4.497534in}{4.178745in}}%
\pgfpathlineto{\pgfqpoint{4.526910in}{4.189592in}}%
\pgfpathlineto{\pgfqpoint{4.556471in}{4.194065in}}%
\pgfpathlineto{\pgfqpoint{4.567447in}{4.188113in}}%
\pgfpathlineto{\pgfqpoint{4.582218in}{4.194897in}}%
\pgfpathlineto{\pgfqpoint{4.589841in}{4.190144in}}%
\pgfpathlineto{\pgfqpoint{4.596587in}{4.167509in}}%
\pgfpathlineto{\pgfqpoint{4.582547in}{4.160499in}}%
\pgfpathlineto{\pgfqpoint{4.572911in}{4.132898in}}%
\pgfpathlineto{\pgfqpoint{4.562808in}{4.122151in}}%
\pgfpathlineto{\pgfqpoint{4.531837in}{4.124502in}}%
\pgfpathlineto{\pgfqpoint{4.533501in}{4.140709in}}%
\pgfpathlineto{\pgfqpoint{4.516587in}{4.133630in}}%
\pgfpathlineto{\pgfqpoint{4.512908in}{4.120112in}}%
\pgfpathlineto{\pgfqpoint{4.524502in}{4.106111in}}%
\pgfpathlineto{\pgfqpoint{4.528440in}{4.082346in}}%
\pgfpathlineto{\pgfqpoint{4.509469in}{4.068806in}}%
\pgfpathlineto{\pgfqpoint{4.529986in}{4.064043in}}%
\pgfpathlineto{\pgfqpoint{4.539271in}{4.074000in}}%
\pgfpathlineto{\pgfqpoint{4.559800in}{4.075213in}}%
\pgfpathlineto{\pgfqpoint{4.549250in}{4.053712in}}%
\pgfpathlineto{\pgfqpoint{4.536329in}{4.043352in}}%
\pgfpathlineto{\pgfqpoint{4.497243in}{4.032980in}}%
\pgfpathlineto{\pgfqpoint{4.457112in}{3.996583in}}%
\pgfpathlineto{\pgfqpoint{4.432367in}{3.960721in}}%
\pgfpathlineto{\pgfqpoint{4.432221in}{3.946156in}}%
\pgfpathlineto{\pgfqpoint{4.422337in}{3.926058in}}%
\pgfpathlineto{\pgfqpoint{4.371607in}{3.912835in}}%
\pgfpathlineto{\pgfqpoint{4.250965in}{3.999317in}}%
\pgfpathlineto{\pgfqpoint{4.144721in}{3.983707in}}%
\pgfpathlineto{\pgfqpoint{4.143831in}{3.998121in}}%
\pgfpathlineto{\pgfqpoint{4.127682in}{4.014353in}}%
\pgfpathlineto{\pgfqpoint{4.115457in}{4.018326in}}%
\pgfpathlineto{\pgfqpoint{4.000022in}{4.006332in}}%
\pgfpathlineto{\pgfqpoint{3.973490in}{3.997333in}}%
\pgfpathlineto{\pgfqpoint{3.925582in}{3.973502in}}%
\pgfpathlineto{\pgfqpoint{3.884181in}{3.966905in}}%
\pgfpathlineto{\pgfqpoint{3.811487in}{3.957948in}}%
\pgfusepath{stroke}%
\end{pgfscope}%
\begin{pgfscope}%
\pgfpathrectangle{\pgfqpoint{0.100000in}{2.493063in}}{\pgfqpoint{5.037500in}{3.427208in}}%
\pgfusepath{clip}%
\pgfsetbuttcap%
\pgfsetroundjoin%
\pgfsetlinewidth{0.501875pt}%
\definecolor{currentstroke}{rgb}{0.827451,0.827451,0.827451}%
\pgfsetstrokecolor{currentstroke}%
\pgfsetdash{}{0pt}%
\pgfpathmoveto{\pgfqpoint{3.811487in}{3.957948in}}%
\pgfpathlineto{\pgfqpoint{3.690671in}{3.944679in}}%
\pgfpathlineto{\pgfqpoint{3.580149in}{3.934740in}}%
\pgfpathlineto{\pgfqpoint{3.466555in}{3.927394in}}%
\pgfpathlineto{\pgfqpoint{3.447056in}{3.924568in}}%
\pgfpathlineto{\pgfqpoint{3.370483in}{3.918678in}}%
\pgfpathlineto{\pgfqpoint{3.247837in}{3.911521in}}%
\pgfpathlineto{\pgfqpoint{3.269658in}{3.929634in}}%
\pgfpathlineto{\pgfqpoint{3.264679in}{3.948699in}}%
\pgfpathlineto{\pgfqpoint{3.269502in}{3.971811in}}%
\pgfpathlineto{\pgfqpoint{3.276871in}{3.982691in}}%
\pgfpathlineto{\pgfqpoint{3.276593in}{3.997829in}}%
\pgfpathlineto{\pgfqpoint{3.296214in}{4.007350in}}%
\pgfpathlineto{\pgfqpoint{3.296248in}{4.029628in}}%
\pgfpathlineto{\pgfqpoint{3.298944in}{4.055308in}}%
\pgfpathlineto{\pgfqpoint{3.312625in}{4.071720in}}%
\pgfpathlineto{\pgfqpoint{3.310355in}{4.088001in}}%
\pgfpathlineto{\pgfqpoint{3.315658in}{4.088373in}}%
\pgfpathlineto{\pgfqpoint{3.321282in}{4.089187in}}%
\pgfpathlineto{\pgfqpoint{3.448152in}{4.097250in}}%
\pgfpathlineto{\pgfqpoint{3.443904in}{4.118092in}}%
\pgfpathlineto{\pgfqpoint{3.464320in}{4.115264in}}%
\pgfpathlineto{\pgfqpoint{3.587896in}{4.127489in}}%
\pgfpathlineto{\pgfqpoint{3.687088in}{4.133103in}}%
\pgfpathlineto{\pgfqpoint{3.747536in}{4.138186in}}%
\pgfpathlineto{\pgfqpoint{3.847633in}{4.147678in}}%
\pgfpathlineto{\pgfqpoint{3.849302in}{4.149566in}}%
\pgfpathlineto{\pgfqpoint{3.892070in}{4.153801in}}%
\pgfpathlineto{\pgfqpoint{4.033298in}{4.173639in}}%
\pgfpathlineto{\pgfqpoint{4.030107in}{4.150117in}}%
\pgfpathlineto{\pgfqpoint{4.016161in}{4.136311in}}%
\pgfpathlineto{\pgfqpoint{4.008752in}{4.116482in}}%
\pgfpathlineto{\pgfqpoint{3.991206in}{4.115991in}}%
\pgfpathlineto{\pgfqpoint{3.957980in}{4.089966in}}%
\pgfpathlineto{\pgfqpoint{3.942259in}{4.092040in}}%
\pgfpathlineto{\pgfqpoint{3.924997in}{4.063478in}}%
\pgfpathlineto{\pgfqpoint{3.901997in}{4.053964in}}%
\pgfpathlineto{\pgfqpoint{3.879975in}{4.032678in}}%
\pgfpathlineto{\pgfqpoint{3.854924in}{4.028722in}}%
\pgfpathlineto{\pgfqpoint{3.834884in}{4.009407in}}%
\pgfpathlineto{\pgfqpoint{3.833945in}{3.993596in}}%
\pgfpathlineto{\pgfqpoint{3.811562in}{3.983705in}}%
\pgfpathlineto{\pgfqpoint{3.811487in}{3.957948in}}%
\pgfusepath{stroke}%
\end{pgfscope}%
\begin{pgfscope}%
\pgfpathrectangle{\pgfqpoint{0.100000in}{2.493063in}}{\pgfqpoint{5.037500in}{3.427208in}}%
\pgfusepath{clip}%
\pgfsetbuttcap%
\pgfsetroundjoin%
\pgfsetlinewidth{0.501875pt}%
\definecolor{currentstroke}{rgb}{0.827451,0.827451,0.827451}%
\pgfsetstrokecolor{currentstroke}%
\pgfsetdash{}{0pt}%
\pgfpathmoveto{\pgfqpoint{1.668801in}{3.600265in}}%
\pgfpathlineto{\pgfqpoint{1.660184in}{3.614111in}}%
\pgfpathlineto{\pgfqpoint{1.663749in}{3.626038in}}%
\pgfpathlineto{\pgfqpoint{1.724316in}{3.618525in}}%
\pgfpathlineto{\pgfqpoint{1.836936in}{3.605734in}}%
\pgfpathlineto{\pgfqpoint{2.012416in}{3.588241in}}%
\pgfpathlineto{\pgfqpoint{2.017532in}{3.647562in}}%
\pgfpathlineto{\pgfqpoint{2.043908in}{3.904143in}}%
\pgfpathlineto{\pgfqpoint{2.061241in}{4.104844in}}%
\pgfpathlineto{\pgfqpoint{2.064008in}{4.104689in}}%
\pgfpathlineto{\pgfqpoint{2.140792in}{4.098021in}}%
\pgfpathlineto{\pgfqpoint{2.240180in}{4.090507in}}%
\pgfpathlineto{\pgfqpoint{2.340896in}{4.084874in}}%
\pgfpathlineto{\pgfqpoint{2.329128in}{3.862498in}}%
\pgfpathlineto{\pgfqpoint{2.335330in}{3.864101in}}%
\pgfpathlineto{\pgfqpoint{2.357586in}{3.839596in}}%
\pgfpathlineto{\pgfqpoint{2.367963in}{3.842707in}}%
\pgfpathlineto{\pgfqpoint{2.383723in}{3.837323in}}%
\pgfpathlineto{\pgfqpoint{2.397482in}{3.840041in}}%
\pgfpathlineto{\pgfqpoint{2.405164in}{3.818222in}}%
\pgfpathlineto{\pgfqpoint{2.422328in}{3.817076in}}%
\pgfpathlineto{\pgfqpoint{2.449732in}{3.806802in}}%
\pgfpathlineto{\pgfqpoint{2.457440in}{3.809697in}}%
\pgfpathlineto{\pgfqpoint{2.473975in}{3.797018in}}%
\pgfpathlineto{\pgfqpoint{2.483707in}{3.806801in}}%
\pgfpathlineto{\pgfqpoint{2.500572in}{3.803102in}}%
\pgfpathlineto{\pgfqpoint{2.508425in}{3.790104in}}%
\pgfpathlineto{\pgfqpoint{2.521320in}{3.785726in}}%
\pgfpathlineto{\pgfqpoint{2.530266in}{3.772244in}}%
\pgfpathlineto{\pgfqpoint{2.548129in}{3.786915in}}%
\pgfpathlineto{\pgfqpoint{2.568572in}{3.768713in}}%
\pgfpathlineto{\pgfqpoint{2.591944in}{3.776148in}}%
\pgfpathlineto{\pgfqpoint{2.603777in}{3.756985in}}%
\pgfpathlineto{\pgfqpoint{2.615348in}{3.779267in}}%
\pgfpathlineto{\pgfqpoint{2.636148in}{3.766633in}}%
\pgfpathlineto{\pgfqpoint{2.649143in}{3.769996in}}%
\pgfpathlineto{\pgfqpoint{2.667699in}{3.761062in}}%
\pgfpathlineto{\pgfqpoint{2.691756in}{3.758526in}}%
\pgfpathlineto{\pgfqpoint{2.714559in}{3.772273in}}%
\pgfpathlineto{\pgfqpoint{2.722809in}{3.766470in}}%
\pgfpathlineto{\pgfqpoint{2.745983in}{3.778253in}}%
\pgfpathlineto{\pgfqpoint{2.760751in}{3.769057in}}%
\pgfpathlineto{\pgfqpoint{2.772424in}{3.769604in}}%
\pgfpathlineto{\pgfqpoint{2.781668in}{3.780436in}}%
\pgfpathlineto{\pgfqpoint{2.796214in}{3.775046in}}%
\pgfpathlineto{\pgfqpoint{2.828963in}{3.749966in}}%
\pgfpathlineto{\pgfqpoint{2.854521in}{3.742061in}}%
\pgfpathlineto{\pgfqpoint{2.864771in}{3.732378in}}%
\pgfpathlineto{\pgfqpoint{2.877599in}{3.737651in}}%
\pgfpathlineto{\pgfqpoint{2.897040in}{3.733613in}}%
\pgfpathlineto{\pgfqpoint{2.899043in}{3.552768in}}%
\pgfpathlineto{\pgfqpoint{2.912544in}{3.541313in}}%
\pgfpathlineto{\pgfqpoint{2.923366in}{3.520204in}}%
\pgfpathlineto{\pgfqpoint{2.919562in}{3.506082in}}%
\pgfpathlineto{\pgfqpoint{2.928060in}{3.500103in}}%
\pgfpathlineto{\pgfqpoint{2.934966in}{3.474061in}}%
\pgfpathlineto{\pgfqpoint{2.948193in}{3.460114in}}%
\pgfpathlineto{\pgfqpoint{2.951156in}{3.430504in}}%
\pgfpathlineto{\pgfqpoint{2.948856in}{3.417967in}}%
\pgfpathlineto{\pgfqpoint{2.930876in}{3.384788in}}%
\pgfpathlineto{\pgfqpoint{2.928809in}{3.360933in}}%
\pgfpathlineto{\pgfqpoint{2.934908in}{3.355954in}}%
\pgfpathlineto{\pgfqpoint{2.935438in}{3.335055in}}%
\pgfpathlineto{\pgfqpoint{2.929239in}{3.321920in}}%
\pgfpathlineto{\pgfqpoint{2.919598in}{3.316972in}}%
\pgfpathlineto{\pgfqpoint{2.910158in}{3.299808in}}%
\pgfpathlineto{\pgfqpoint{2.922190in}{3.283183in}}%
\pgfpathlineto{\pgfqpoint{2.898834in}{3.282856in}}%
\pgfpathlineto{\pgfqpoint{2.836385in}{3.254284in}}%
\pgfpathlineto{\pgfqpoint{2.848308in}{3.271383in}}%
\pgfpathlineto{\pgfqpoint{2.833870in}{3.280602in}}%
\pgfpathlineto{\pgfqpoint{2.830853in}{3.296277in}}%
\pgfpathlineto{\pgfqpoint{2.802674in}{3.268975in}}%
\pgfpathlineto{\pgfqpoint{2.815180in}{3.250314in}}%
\pgfpathlineto{\pgfqpoint{2.797340in}{3.226553in}}%
\pgfpathlineto{\pgfqpoint{2.787739in}{3.227050in}}%
\pgfpathlineto{\pgfqpoint{2.778707in}{3.201174in}}%
\pgfpathlineto{\pgfqpoint{2.750144in}{3.180819in}}%
\pgfpathlineto{\pgfqpoint{2.723459in}{3.173484in}}%
\pgfpathlineto{\pgfqpoint{2.676917in}{3.154401in}}%
\pgfpathlineto{\pgfqpoint{2.672094in}{3.165046in}}%
\pgfpathlineto{\pgfqpoint{2.643101in}{3.155248in}}%
\pgfpathlineto{\pgfqpoint{2.660934in}{3.138565in}}%
\pgfpathlineto{\pgfqpoint{2.632807in}{3.124071in}}%
\pgfpathlineto{\pgfqpoint{2.602468in}{3.102188in}}%
\pgfpathlineto{\pgfqpoint{2.595182in}{3.112361in}}%
\pgfpathlineto{\pgfqpoint{2.570354in}{3.097114in}}%
\pgfpathlineto{\pgfqpoint{2.594689in}{3.091325in}}%
\pgfpathlineto{\pgfqpoint{2.576341in}{3.067353in}}%
\pgfpathlineto{\pgfqpoint{2.567385in}{3.074487in}}%
\pgfpathlineto{\pgfqpoint{2.555355in}{3.063028in}}%
\pgfpathlineto{\pgfqpoint{2.570347in}{3.053019in}}%
\pgfpathlineto{\pgfqpoint{2.561483in}{3.038390in}}%
\pgfpathlineto{\pgfqpoint{2.550589in}{3.003757in}}%
\pgfpathlineto{\pgfqpoint{2.534875in}{2.970402in}}%
\pgfpathlineto{\pgfqpoint{2.541929in}{2.948538in}}%
\pgfpathlineto{\pgfqpoint{2.553920in}{2.896991in}}%
\pgfpathlineto{\pgfqpoint{2.555001in}{2.876129in}}%
\pgfpathlineto{\pgfqpoint{2.573528in}{2.848767in}}%
\pgfpathlineto{\pgfqpoint{2.559241in}{2.850366in}}%
\pgfpathlineto{\pgfqpoint{2.545364in}{2.836562in}}%
\pgfpathlineto{\pgfqpoint{2.529111in}{2.847411in}}%
\pgfpathlineto{\pgfqpoint{2.523277in}{2.858234in}}%
\pgfpathlineto{\pgfqpoint{2.500146in}{2.863269in}}%
\pgfpathlineto{\pgfqpoint{2.464813in}{2.863886in}}%
\pgfpathlineto{\pgfqpoint{2.438779in}{2.884383in}}%
\pgfpathlineto{\pgfqpoint{2.415141in}{2.887806in}}%
\pgfpathlineto{\pgfqpoint{2.400800in}{2.904095in}}%
\pgfpathlineto{\pgfqpoint{2.370766in}{2.910651in}}%
\pgfpathlineto{\pgfqpoint{2.354344in}{2.963026in}}%
\pgfpathlineto{\pgfqpoint{2.337552in}{2.984007in}}%
\pgfpathlineto{\pgfqpoint{2.340405in}{3.003943in}}%
\pgfpathlineto{\pgfqpoint{2.329993in}{3.018513in}}%
\pgfpathlineto{\pgfqpoint{2.336527in}{3.038428in}}%
\pgfpathlineto{\pgfqpoint{2.331137in}{3.053021in}}%
\pgfpathlineto{\pgfqpoint{2.314263in}{3.059632in}}%
\pgfpathlineto{\pgfqpoint{2.298486in}{3.076460in}}%
\pgfpathlineto{\pgfqpoint{2.292736in}{3.098989in}}%
\pgfpathlineto{\pgfqpoint{2.277807in}{3.119464in}}%
\pgfpathlineto{\pgfqpoint{2.257935in}{3.135395in}}%
\pgfpathlineto{\pgfqpoint{2.240037in}{3.181094in}}%
\pgfpathlineto{\pgfqpoint{2.225584in}{3.231073in}}%
\pgfpathlineto{\pgfqpoint{2.213725in}{3.250834in}}%
\pgfpathlineto{\pgfqpoint{2.193167in}{3.267471in}}%
\pgfpathlineto{\pgfqpoint{2.188029in}{3.279547in}}%
\pgfpathlineto{\pgfqpoint{2.168780in}{3.287045in}}%
\pgfpathlineto{\pgfqpoint{2.149781in}{3.319041in}}%
\pgfpathlineto{\pgfqpoint{2.132430in}{3.316585in}}%
\pgfpathlineto{\pgfqpoint{2.114812in}{3.324576in}}%
\pgfpathlineto{\pgfqpoint{2.089838in}{3.323031in}}%
\pgfpathlineto{\pgfqpoint{2.064467in}{3.336219in}}%
\pgfpathlineto{\pgfqpoint{2.057313in}{3.323677in}}%
\pgfpathlineto{\pgfqpoint{2.027664in}{3.323363in}}%
\pgfpathlineto{\pgfqpoint{2.012569in}{3.299586in}}%
\pgfpathlineto{\pgfqpoint{1.999441in}{3.270148in}}%
\pgfpathlineto{\pgfqpoint{1.971522in}{3.238547in}}%
\pgfpathlineto{\pgfqpoint{1.939894in}{3.252415in}}%
\pgfpathlineto{\pgfqpoint{1.935405in}{3.261580in}}%
\pgfpathlineto{\pgfqpoint{1.916184in}{3.268571in}}%
\pgfpathlineto{\pgfqpoint{1.910279in}{3.278133in}}%
\pgfpathlineto{\pgfqpoint{1.884762in}{3.287785in}}%
\pgfpathlineto{\pgfqpoint{1.870487in}{3.307517in}}%
\pgfpathlineto{\pgfqpoint{1.853817in}{3.317037in}}%
\pgfpathlineto{\pgfqpoint{1.839474in}{3.333650in}}%
\pgfpathlineto{\pgfqpoint{1.828300in}{3.361784in}}%
\pgfpathlineto{\pgfqpoint{1.829550in}{3.400230in}}%
\pgfpathlineto{\pgfqpoint{1.816442in}{3.419658in}}%
\pgfpathlineto{\pgfqpoint{1.814934in}{3.440700in}}%
\pgfpathlineto{\pgfqpoint{1.785850in}{3.472204in}}%
\pgfpathlineto{\pgfqpoint{1.768928in}{3.478929in}}%
\pgfpathlineto{\pgfqpoint{1.750970in}{3.508326in}}%
\pgfpathlineto{\pgfqpoint{1.735668in}{3.520005in}}%
\pgfpathlineto{\pgfqpoint{1.716224in}{3.548446in}}%
\pgfpathlineto{\pgfqpoint{1.696330in}{3.560754in}}%
\pgfpathlineto{\pgfqpoint{1.683301in}{3.592292in}}%
\pgfpathlineto{\pgfqpoint{1.668801in}{3.600265in}}%
\pgfpathlineto{\pgfqpoint{1.668801in}{3.600265in}}%
\pgfusepath{stroke}%
\end{pgfscope}%
\begin{pgfscope}%
\pgfpathrectangle{\pgfqpoint{0.100000in}{2.493063in}}{\pgfqpoint{5.037500in}{3.427208in}}%
\pgfusepath{clip}%
\pgfsetbuttcap%
\pgfsetroundjoin%
\pgfsetlinewidth{0.501875pt}%
\definecolor{currentstroke}{rgb}{0.827451,0.827451,0.827451}%
\pgfsetstrokecolor{currentstroke}%
\pgfsetdash{}{0pt}%
\pgfpathmoveto{\pgfqpoint{1.516457in}{4.229943in}}%
\pgfpathlineto{\pgfqpoint{1.659981in}{4.209100in}}%
\pgfpathlineto{\pgfqpoint{1.715584in}{4.200532in}}%
\pgfpathlineto{\pgfqpoint{1.872319in}{4.181891in}}%
\pgfpathlineto{\pgfqpoint{1.977813in}{4.170887in}}%
\pgfpathlineto{\pgfqpoint{2.069106in}{4.162796in}}%
\pgfpathlineto{\pgfqpoint{2.064008in}{4.104689in}}%
\pgfpathlineto{\pgfqpoint{2.061241in}{4.104844in}}%
\pgfpathlineto{\pgfqpoint{2.052879in}{4.005113in}}%
\pgfpathlineto{\pgfqpoint{2.043908in}{3.904143in}}%
\pgfpathlineto{\pgfqpoint{2.033530in}{3.798444in}}%
\pgfpathlineto{\pgfqpoint{2.017532in}{3.647562in}}%
\pgfpathlineto{\pgfqpoint{2.012416in}{3.588241in}}%
\pgfpathlineto{\pgfqpoint{1.918368in}{3.597731in}}%
\pgfpathlineto{\pgfqpoint{1.836936in}{3.605734in}}%
\pgfpathlineto{\pgfqpoint{1.724316in}{3.618525in}}%
\pgfpathlineto{\pgfqpoint{1.663749in}{3.626038in}}%
\pgfpathlineto{\pgfqpoint{1.660184in}{3.614111in}}%
\pgfpathlineto{\pgfqpoint{1.668801in}{3.600265in}}%
\pgfpathlineto{\pgfqpoint{1.596010in}{3.609718in}}%
\pgfpathlineto{\pgfqpoint{1.506202in}{3.622608in}}%
\pgfpathlineto{\pgfqpoint{1.498028in}{3.571813in}}%
\pgfpathlineto{\pgfqpoint{1.416109in}{3.584143in}}%
\pgfpathlineto{\pgfqpoint{1.435357in}{3.710155in}}%
\pgfpathlineto{\pgfqpoint{1.459184in}{3.863582in}}%
\pgfpathlineto{\pgfqpoint{1.480371in}{3.997241in}}%
\pgfpathlineto{\pgfqpoint{1.498768in}{4.116047in}}%
\pgfpathlineto{\pgfqpoint{1.516457in}{4.229943in}}%
\pgfusepath{stroke}%
\end{pgfscope}%
\begin{pgfscope}%
\pgfpathrectangle{\pgfqpoint{0.100000in}{2.493063in}}{\pgfqpoint{5.037500in}{3.427208in}}%
\pgfusepath{clip}%
\pgfsetbuttcap%
\pgfsetroundjoin%
\pgfsetlinewidth{0.501875pt}%
\definecolor{currentstroke}{rgb}{0.827451,0.827451,0.827451}%
\pgfsetstrokecolor{currentstroke}%
\pgfsetdash{}{0pt}%
\pgfpathmoveto{\pgfqpoint{3.798503in}{3.491653in}}%
\pgfpathlineto{\pgfqpoint{3.661580in}{3.476445in}}%
\pgfpathlineto{\pgfqpoint{3.540884in}{3.467076in}}%
\pgfpathlineto{\pgfqpoint{3.539374in}{3.452285in}}%
\pgfpathlineto{\pgfqpoint{3.550461in}{3.438206in}}%
\pgfpathlineto{\pgfqpoint{3.563992in}{3.429936in}}%
\pgfpathlineto{\pgfqpoint{3.563787in}{3.408159in}}%
\pgfpathlineto{\pgfqpoint{3.560202in}{3.393610in}}%
\pgfpathlineto{\pgfqpoint{3.548259in}{3.383125in}}%
\pgfpathlineto{\pgfqpoint{3.531613in}{3.384235in}}%
\pgfpathlineto{\pgfqpoint{3.515874in}{3.397227in}}%
\pgfpathlineto{\pgfqpoint{3.513069in}{3.420345in}}%
\pgfpathlineto{\pgfqpoint{3.501347in}{3.433791in}}%
\pgfpathlineto{\pgfqpoint{3.493369in}{3.385669in}}%
\pgfpathlineto{\pgfqpoint{3.466258in}{3.390242in}}%
\pgfpathlineto{\pgfqpoint{3.456830in}{3.474390in}}%
\pgfpathlineto{\pgfqpoint{3.446620in}{3.563072in}}%
\pgfpathlineto{\pgfqpoint{3.450412in}{3.691019in}}%
\pgfpathlineto{\pgfqpoint{3.452766in}{3.778819in}}%
\pgfpathlineto{\pgfqpoint{3.457771in}{3.912779in}}%
\pgfpathlineto{\pgfqpoint{3.447056in}{3.924568in}}%
\pgfpathlineto{\pgfqpoint{3.466555in}{3.927394in}}%
\pgfpathlineto{\pgfqpoint{3.580149in}{3.934740in}}%
\pgfpathlineto{\pgfqpoint{3.690671in}{3.944679in}}%
\pgfpathlineto{\pgfqpoint{3.719749in}{3.842827in}}%
\pgfpathlineto{\pgfqpoint{3.739580in}{3.768066in}}%
\pgfpathlineto{\pgfqpoint{3.757355in}{3.705478in}}%
\pgfpathlineto{\pgfqpoint{3.767636in}{3.680276in}}%
\pgfpathlineto{\pgfqpoint{3.783823in}{3.656822in}}%
\pgfpathlineto{\pgfqpoint{3.781209in}{3.644887in}}%
\pgfpathlineto{\pgfqpoint{3.792772in}{3.639073in}}%
\pgfpathlineto{\pgfqpoint{3.779107in}{3.620981in}}%
\pgfpathlineto{\pgfqpoint{3.774492in}{3.588762in}}%
\pgfpathlineto{\pgfqpoint{3.787846in}{3.551078in}}%
\pgfpathlineto{\pgfqpoint{3.785992in}{3.513159in}}%
\pgfpathlineto{\pgfqpoint{3.798503in}{3.491653in}}%
\pgfusepath{stroke}%
\end{pgfscope}%
\begin{pgfscope}%
\pgfpathrectangle{\pgfqpoint{0.100000in}{2.493063in}}{\pgfqpoint{5.037500in}{3.427208in}}%
\pgfusepath{clip}%
\pgfsetbuttcap%
\pgfsetroundjoin%
\pgfsetlinewidth{0.501875pt}%
\definecolor{currentstroke}{rgb}{0.827451,0.827451,0.827451}%
\pgfsetstrokecolor{currentstroke}%
\pgfsetdash{}{0pt}%
\pgfpathmoveto{\pgfqpoint{3.466258in}{3.390242in}}%
\pgfpathlineto{\pgfqpoint{3.438437in}{3.382281in}}%
\pgfpathlineto{\pgfqpoint{3.417699in}{3.387393in}}%
\pgfpathlineto{\pgfqpoint{3.379188in}{3.375143in}}%
\pgfpathlineto{\pgfqpoint{3.350143in}{3.359367in}}%
\pgfpathlineto{\pgfqpoint{3.338225in}{3.387874in}}%
\pgfpathlineto{\pgfqpoint{3.325940in}{3.399823in}}%
\pgfpathlineto{\pgfqpoint{3.320014in}{3.414122in}}%
\pgfpathlineto{\pgfqpoint{3.328701in}{3.452823in}}%
\pgfpathlineto{\pgfqpoint{3.246390in}{3.447755in}}%
\pgfpathlineto{\pgfqpoint{3.139657in}{3.443127in}}%
\pgfpathlineto{\pgfqpoint{3.146007in}{3.452761in}}%
\pgfpathlineto{\pgfqpoint{3.138158in}{3.470955in}}%
\pgfpathlineto{\pgfqpoint{3.145979in}{3.474663in}}%
\pgfpathlineto{\pgfqpoint{3.144244in}{3.492144in}}%
\pgfpathlineto{\pgfqpoint{3.157876in}{3.509096in}}%
\pgfpathlineto{\pgfqpoint{3.165334in}{3.542006in}}%
\pgfpathlineto{\pgfqpoint{3.190283in}{3.563688in}}%
\pgfpathlineto{\pgfqpoint{3.181036in}{3.584745in}}%
\pgfpathlineto{\pgfqpoint{3.198727in}{3.596866in}}%
\pgfpathlineto{\pgfqpoint{3.183467in}{3.618809in}}%
\pgfpathlineto{\pgfqpoint{3.172421in}{3.669042in}}%
\pgfpathlineto{\pgfqpoint{3.176610in}{3.678163in}}%
\pgfpathlineto{\pgfqpoint{3.183224in}{3.695659in}}%
\pgfpathlineto{\pgfqpoint{3.177068in}{3.714020in}}%
\pgfpathlineto{\pgfqpoint{3.180079in}{3.722373in}}%
\pgfpathlineto{\pgfqpoint{3.169249in}{3.753932in}}%
\pgfpathlineto{\pgfqpoint{3.201219in}{3.810598in}}%
\pgfpathlineto{\pgfqpoint{3.202705in}{3.825680in}}%
\pgfpathlineto{\pgfqpoint{3.218105in}{3.836647in}}%
\pgfpathlineto{\pgfqpoint{3.234537in}{3.872851in}}%
\pgfpathlineto{\pgfqpoint{3.233756in}{3.887568in}}%
\pgfpathlineto{\pgfqpoint{3.254223in}{3.902611in}}%
\pgfpathlineto{\pgfqpoint{3.247837in}{3.911521in}}%
\pgfpathlineto{\pgfqpoint{3.370483in}{3.918678in}}%
\pgfpathlineto{\pgfqpoint{3.447056in}{3.924568in}}%
\pgfpathlineto{\pgfqpoint{3.457771in}{3.912779in}}%
\pgfpathlineto{\pgfqpoint{3.452766in}{3.778819in}}%
\pgfpathlineto{\pgfqpoint{3.450412in}{3.691019in}}%
\pgfpathlineto{\pgfqpoint{3.446620in}{3.563072in}}%
\pgfpathlineto{\pgfqpoint{3.456830in}{3.474390in}}%
\pgfpathlineto{\pgfqpoint{3.466258in}{3.390242in}}%
\pgfusepath{stroke}%
\end{pgfscope}%
\begin{pgfscope}%
\pgfpathrectangle{\pgfqpoint{0.100000in}{2.493063in}}{\pgfqpoint{5.037500in}{3.427208in}}%
\pgfusepath{clip}%
\pgfsetbuttcap%
\pgfsetroundjoin%
\pgfsetlinewidth{0.501875pt}%
\definecolor{currentstroke}{rgb}{0.827451,0.827451,0.827451}%
\pgfsetstrokecolor{currentstroke}%
\pgfsetdash{}{0pt}%
\pgfpathmoveto{\pgfqpoint{3.690671in}{3.944679in}}%
\pgfpathlineto{\pgfqpoint{3.811487in}{3.957948in}}%
\pgfpathlineto{\pgfqpoint{3.884181in}{3.966905in}}%
\pgfpathlineto{\pgfqpoint{3.925582in}{3.973502in}}%
\pgfpathlineto{\pgfqpoint{3.908432in}{3.946452in}}%
\pgfpathlineto{\pgfqpoint{3.908443in}{3.933698in}}%
\pgfpathlineto{\pgfqpoint{3.926096in}{3.926821in}}%
\pgfpathlineto{\pgfqpoint{3.938102in}{3.915717in}}%
\pgfpathlineto{\pgfqpoint{3.956245in}{3.914345in}}%
\pgfpathlineto{\pgfqpoint{3.973202in}{3.883087in}}%
\pgfpathlineto{\pgfqpoint{3.991596in}{3.861044in}}%
\pgfpathlineto{\pgfqpoint{4.024265in}{3.842518in}}%
\pgfpathlineto{\pgfqpoint{4.033550in}{3.827742in}}%
\pgfpathlineto{\pgfqpoint{4.062682in}{3.811761in}}%
\pgfpathlineto{\pgfqpoint{4.064872in}{3.800172in}}%
\pgfpathlineto{\pgfqpoint{4.088505in}{3.777128in}}%
\pgfpathlineto{\pgfqpoint{4.106607in}{3.770567in}}%
\pgfpathlineto{\pgfqpoint{4.119426in}{3.748808in}}%
\pgfpathlineto{\pgfqpoint{4.125101in}{3.724350in}}%
\pgfpathlineto{\pgfqpoint{4.143903in}{3.714896in}}%
\pgfpathlineto{\pgfqpoint{4.159040in}{3.688607in}}%
\pgfpathlineto{\pgfqpoint{4.163922in}{3.669272in}}%
\pgfpathlineto{\pgfqpoint{4.186104in}{3.661043in}}%
\pgfpathlineto{\pgfqpoint{4.167490in}{3.625404in}}%
\pgfpathlineto{\pgfqpoint{4.160485in}{3.604342in}}%
\pgfpathlineto{\pgfqpoint{4.165056in}{3.594470in}}%
\pgfpathlineto{\pgfqpoint{4.157076in}{3.578062in}}%
\pgfpathlineto{\pgfqpoint{4.153683in}{3.555378in}}%
\pgfpathlineto{\pgfqpoint{4.139492in}{3.551153in}}%
\pgfpathlineto{\pgfqpoint{4.146986in}{3.530384in}}%
\pgfpathlineto{\pgfqpoint{4.146508in}{3.504034in}}%
\pgfpathlineto{\pgfqpoint{4.121284in}{3.505022in}}%
\pgfpathlineto{\pgfqpoint{4.103711in}{3.509829in}}%
\pgfpathlineto{\pgfqpoint{4.096108in}{3.500892in}}%
\pgfpathlineto{\pgfqpoint{4.101781in}{3.479780in}}%
\pgfpathlineto{\pgfqpoint{4.100548in}{3.455235in}}%
\pgfpathlineto{\pgfqpoint{4.089462in}{3.453357in}}%
\pgfpathlineto{\pgfqpoint{4.080494in}{3.476281in}}%
\pgfpathlineto{\pgfqpoint{3.989179in}{3.470195in}}%
\pgfpathlineto{\pgfqpoint{3.874037in}{3.463913in}}%
\pgfpathlineto{\pgfqpoint{3.815967in}{3.459825in}}%
\pgfpathlineto{\pgfqpoint{3.798503in}{3.491653in}}%
\pgfpathlineto{\pgfqpoint{3.785992in}{3.513159in}}%
\pgfpathlineto{\pgfqpoint{3.787846in}{3.551078in}}%
\pgfpathlineto{\pgfqpoint{3.774492in}{3.588762in}}%
\pgfpathlineto{\pgfqpoint{3.779107in}{3.620981in}}%
\pgfpathlineto{\pgfqpoint{3.792772in}{3.639073in}}%
\pgfpathlineto{\pgfqpoint{3.781209in}{3.644887in}}%
\pgfpathlineto{\pgfqpoint{3.783823in}{3.656822in}}%
\pgfpathlineto{\pgfqpoint{3.767636in}{3.680276in}}%
\pgfpathlineto{\pgfqpoint{3.757355in}{3.705478in}}%
\pgfpathlineto{\pgfqpoint{3.739580in}{3.768066in}}%
\pgfpathlineto{\pgfqpoint{3.719749in}{3.842827in}}%
\pgfpathlineto{\pgfqpoint{3.690671in}{3.944679in}}%
\pgfusepath{stroke}%
\end{pgfscope}%
\begin{pgfscope}%
\pgfpathrectangle{\pgfqpoint{0.100000in}{2.493063in}}{\pgfqpoint{5.037500in}{3.427208in}}%
\pgfusepath{clip}%
\pgfsetbuttcap%
\pgfsetroundjoin%
\pgfsetlinewidth{0.501875pt}%
\definecolor{currentstroke}{rgb}{0.827451,0.827451,0.827451}%
\pgfsetstrokecolor{currentstroke}%
\pgfsetdash{}{0pt}%
\pgfpathmoveto{\pgfqpoint{3.925582in}{3.973502in}}%
\pgfpathlineto{\pgfqpoint{3.973490in}{3.997333in}}%
\pgfpathlineto{\pgfqpoint{4.000022in}{4.006332in}}%
\pgfpathlineto{\pgfqpoint{4.115457in}{4.018326in}}%
\pgfpathlineto{\pgfqpoint{4.127682in}{4.014353in}}%
\pgfpathlineto{\pgfqpoint{4.143831in}{3.998121in}}%
\pgfpathlineto{\pgfqpoint{4.144721in}{3.983707in}}%
\pgfpathlineto{\pgfqpoint{4.250965in}{3.999317in}}%
\pgfpathlineto{\pgfqpoint{4.371607in}{3.912835in}}%
\pgfpathlineto{\pgfqpoint{4.348983in}{3.889254in}}%
\pgfpathlineto{\pgfqpoint{4.317938in}{3.834490in}}%
\pgfpathlineto{\pgfqpoint{4.326775in}{3.822720in}}%
\pgfpathlineto{\pgfqpoint{4.310214in}{3.799876in}}%
\pgfpathlineto{\pgfqpoint{4.293741in}{3.797290in}}%
\pgfpathlineto{\pgfqpoint{4.293680in}{3.783562in}}%
\pgfpathlineto{\pgfqpoint{4.281764in}{3.769198in}}%
\pgfpathlineto{\pgfqpoint{4.266934in}{3.766285in}}%
\pgfpathlineto{\pgfqpoint{4.270170in}{3.753521in}}%
\pgfpathlineto{\pgfqpoint{4.261894in}{3.743723in}}%
\pgfpathlineto{\pgfqpoint{4.242099in}{3.735259in}}%
\pgfpathlineto{\pgfqpoint{4.217970in}{3.716005in}}%
\pgfpathlineto{\pgfqpoint{4.222525in}{3.703548in}}%
\pgfpathlineto{\pgfqpoint{4.207340in}{3.695727in}}%
\pgfpathlineto{\pgfqpoint{4.194533in}{3.703950in}}%
\pgfpathlineto{\pgfqpoint{4.185167in}{3.668197in}}%
\pgfpathlineto{\pgfqpoint{4.163922in}{3.669272in}}%
\pgfpathlineto{\pgfqpoint{4.159040in}{3.688607in}}%
\pgfpathlineto{\pgfqpoint{4.143903in}{3.714896in}}%
\pgfpathlineto{\pgfqpoint{4.125101in}{3.724350in}}%
\pgfpathlineto{\pgfqpoint{4.119426in}{3.748808in}}%
\pgfpathlineto{\pgfqpoint{4.106607in}{3.770567in}}%
\pgfpathlineto{\pgfqpoint{4.088505in}{3.777128in}}%
\pgfpathlineto{\pgfqpoint{4.064872in}{3.800172in}}%
\pgfpathlineto{\pgfqpoint{4.062682in}{3.811761in}}%
\pgfpathlineto{\pgfqpoint{4.033550in}{3.827742in}}%
\pgfpathlineto{\pgfqpoint{4.024265in}{3.842518in}}%
\pgfpathlineto{\pgfqpoint{3.991596in}{3.861044in}}%
\pgfpathlineto{\pgfqpoint{3.973202in}{3.883087in}}%
\pgfpathlineto{\pgfqpoint{3.956245in}{3.914345in}}%
\pgfpathlineto{\pgfqpoint{3.938102in}{3.915717in}}%
\pgfpathlineto{\pgfqpoint{3.926096in}{3.926821in}}%
\pgfpathlineto{\pgfqpoint{3.908443in}{3.933698in}}%
\pgfpathlineto{\pgfqpoint{3.908432in}{3.946452in}}%
\pgfpathlineto{\pgfqpoint{3.925582in}{3.973502in}}%
\pgfusepath{stroke}%
\end{pgfscope}%
\begin{pgfscope}%
\pgfpathrectangle{\pgfqpoint{0.100000in}{2.493063in}}{\pgfqpoint{5.037500in}{3.427208in}}%
\pgfusepath{clip}%
\pgfsetbuttcap%
\pgfsetroundjoin%
\pgfsetlinewidth{0.501875pt}%
\definecolor{currentstroke}{rgb}{0.827451,0.827451,0.827451}%
\pgfsetstrokecolor{currentstroke}%
\pgfsetdash{}{0pt}%
\pgfpathmoveto{\pgfqpoint{2.839596in}{4.071836in}}%
\pgfpathlineto{\pgfqpoint{2.909997in}{4.072469in}}%
\pgfpathlineto{\pgfqpoint{3.003135in}{4.074167in}}%
\pgfpathlineto{\pgfqpoint{3.162287in}{4.079219in}}%
\pgfpathlineto{\pgfqpoint{3.253318in}{4.084033in}}%
\pgfpathlineto{\pgfqpoint{3.263080in}{4.071948in}}%
\pgfpathlineto{\pgfqpoint{3.262470in}{4.059183in}}%
\pgfpathlineto{\pgfqpoint{3.240418in}{4.037162in}}%
\pgfpathlineto{\pgfqpoint{3.235087in}{4.025104in}}%
\pgfpathlineto{\pgfqpoint{3.296248in}{4.029628in}}%
\pgfpathlineto{\pgfqpoint{3.296214in}{4.007350in}}%
\pgfpathlineto{\pgfqpoint{3.276593in}{3.997829in}}%
\pgfpathlineto{\pgfqpoint{3.276871in}{3.982691in}}%
\pgfpathlineto{\pgfqpoint{3.269502in}{3.971811in}}%
\pgfpathlineto{\pgfqpoint{3.264679in}{3.948699in}}%
\pgfpathlineto{\pgfqpoint{3.269658in}{3.929634in}}%
\pgfpathlineto{\pgfqpoint{3.247837in}{3.911521in}}%
\pgfpathlineto{\pgfqpoint{3.254223in}{3.902611in}}%
\pgfpathlineto{\pgfqpoint{3.233756in}{3.887568in}}%
\pgfpathlineto{\pgfqpoint{3.234537in}{3.872851in}}%
\pgfpathlineto{\pgfqpoint{3.218105in}{3.836647in}}%
\pgfpathlineto{\pgfqpoint{3.202705in}{3.825680in}}%
\pgfpathlineto{\pgfqpoint{3.201219in}{3.810598in}}%
\pgfpathlineto{\pgfqpoint{3.169249in}{3.753932in}}%
\pgfpathlineto{\pgfqpoint{3.180079in}{3.722373in}}%
\pgfpathlineto{\pgfqpoint{3.177068in}{3.714020in}}%
\pgfpathlineto{\pgfqpoint{3.183224in}{3.695659in}}%
\pgfpathlineto{\pgfqpoint{3.176610in}{3.678163in}}%
\pgfpathlineto{\pgfqpoint{3.089189in}{3.674554in}}%
\pgfpathlineto{\pgfqpoint{2.975699in}{3.672684in}}%
\pgfpathlineto{\pgfqpoint{2.897420in}{3.671978in}}%
\pgfpathlineto{\pgfqpoint{2.897040in}{3.733613in}}%
\pgfpathlineto{\pgfqpoint{2.877599in}{3.737651in}}%
\pgfpathlineto{\pgfqpoint{2.864771in}{3.732378in}}%
\pgfpathlineto{\pgfqpoint{2.854521in}{3.742061in}}%
\pgfpathlineto{\pgfqpoint{2.855579in}{3.807350in}}%
\pgfpathlineto{\pgfqpoint{2.857873in}{3.946362in}}%
\pgfpathlineto{\pgfqpoint{2.846754in}{4.027764in}}%
\pgfpathlineto{\pgfqpoint{2.839596in}{4.071836in}}%
\pgfusepath{stroke}%
\end{pgfscope}%
\begin{pgfscope}%
\pgfpathrectangle{\pgfqpoint{0.100000in}{2.493063in}}{\pgfqpoint{5.037500in}{3.427208in}}%
\pgfusepath{clip}%
\pgfsetbuttcap%
\pgfsetroundjoin%
\pgfsetlinewidth{0.501875pt}%
\definecolor{currentstroke}{rgb}{0.827451,0.827451,0.827451}%
\pgfsetstrokecolor{currentstroke}%
\pgfsetdash{}{0pt}%
\pgfpathmoveto{\pgfqpoint{2.897420in}{3.671978in}}%
\pgfpathlineto{\pgfqpoint{2.975699in}{3.672684in}}%
\pgfpathlineto{\pgfqpoint{3.089189in}{3.674554in}}%
\pgfpathlineto{\pgfqpoint{3.176610in}{3.678163in}}%
\pgfpathlineto{\pgfqpoint{3.172421in}{3.669042in}}%
\pgfpathlineto{\pgfqpoint{3.183467in}{3.618809in}}%
\pgfpathlineto{\pgfqpoint{3.198727in}{3.596866in}}%
\pgfpathlineto{\pgfqpoint{3.181036in}{3.584745in}}%
\pgfpathlineto{\pgfqpoint{3.190283in}{3.563688in}}%
\pgfpathlineto{\pgfqpoint{3.165334in}{3.542006in}}%
\pgfpathlineto{\pgfqpoint{3.157876in}{3.509096in}}%
\pgfpathlineto{\pgfqpoint{3.144244in}{3.492144in}}%
\pgfpathlineto{\pgfqpoint{3.145979in}{3.474663in}}%
\pgfpathlineto{\pgfqpoint{3.138158in}{3.470955in}}%
\pgfpathlineto{\pgfqpoint{3.146007in}{3.452761in}}%
\pgfpathlineto{\pgfqpoint{3.139657in}{3.443127in}}%
\pgfpathlineto{\pgfqpoint{3.246390in}{3.447755in}}%
\pgfpathlineto{\pgfqpoint{3.328701in}{3.452823in}}%
\pgfpathlineto{\pgfqpoint{3.320014in}{3.414122in}}%
\pgfpathlineto{\pgfqpoint{3.325940in}{3.399823in}}%
\pgfpathlineto{\pgfqpoint{3.338225in}{3.387874in}}%
\pgfpathlineto{\pgfqpoint{3.350143in}{3.359367in}}%
\pgfpathlineto{\pgfqpoint{3.312500in}{3.365934in}}%
\pgfpathlineto{\pgfqpoint{3.298615in}{3.376732in}}%
\pgfpathlineto{\pgfqpoint{3.282070in}{3.377246in}}%
\pgfpathlineto{\pgfqpoint{3.264677in}{3.353602in}}%
\pgfpathlineto{\pgfqpoint{3.268143in}{3.342835in}}%
\pgfpathlineto{\pgfqpoint{3.323709in}{3.336321in}}%
\pgfpathlineto{\pgfqpoint{3.338256in}{3.323945in}}%
\pgfpathlineto{\pgfqpoint{3.361928in}{3.317578in}}%
\pgfpathlineto{\pgfqpoint{3.351589in}{3.302926in}}%
\pgfpathlineto{\pgfqpoint{3.334171in}{3.290148in}}%
\pgfpathlineto{\pgfqpoint{3.380146in}{3.261411in}}%
\pgfpathlineto{\pgfqpoint{3.401562in}{3.256921in}}%
\pgfpathlineto{\pgfqpoint{3.410468in}{3.233432in}}%
\pgfpathlineto{\pgfqpoint{3.390304in}{3.229111in}}%
\pgfpathlineto{\pgfqpoint{3.352407in}{3.258654in}}%
\pgfpathlineto{\pgfqpoint{3.337590in}{3.262760in}}%
\pgfpathlineto{\pgfqpoint{3.330420in}{3.274431in}}%
\pgfpathlineto{\pgfqpoint{3.308442in}{3.269641in}}%
\pgfpathlineto{\pgfqpoint{3.301561in}{3.254602in}}%
\pgfpathlineto{\pgfqpoint{3.305919in}{3.237845in}}%
\pgfpathlineto{\pgfqpoint{3.291123in}{3.227957in}}%
\pgfpathlineto{\pgfqpoint{3.277619in}{3.252318in}}%
\pgfpathlineto{\pgfqpoint{3.253967in}{3.245054in}}%
\pgfpathlineto{\pgfqpoint{3.245105in}{3.230519in}}%
\pgfpathlineto{\pgfqpoint{3.228308in}{3.234653in}}%
\pgfpathlineto{\pgfqpoint{3.217569in}{3.253009in}}%
\pgfpathlineto{\pgfqpoint{3.189078in}{3.259071in}}%
\pgfpathlineto{\pgfqpoint{3.183698in}{3.268622in}}%
\pgfpathlineto{\pgfqpoint{3.154014in}{3.285261in}}%
\pgfpathlineto{\pgfqpoint{3.146645in}{3.299809in}}%
\pgfpathlineto{\pgfqpoint{3.121782in}{3.293864in}}%
\pgfpathlineto{\pgfqpoint{3.124940in}{3.307183in}}%
\pgfpathlineto{\pgfqpoint{3.094036in}{3.293520in}}%
\pgfpathlineto{\pgfqpoint{3.102338in}{3.279344in}}%
\pgfpathlineto{\pgfqpoint{3.078472in}{3.270963in}}%
\pgfpathlineto{\pgfqpoint{3.046862in}{3.275569in}}%
\pgfpathlineto{\pgfqpoint{2.982884in}{3.297484in}}%
\pgfpathlineto{\pgfqpoint{2.933518in}{3.293134in}}%
\pgfpathlineto{\pgfqpoint{2.929239in}{3.321920in}}%
\pgfpathlineto{\pgfqpoint{2.935438in}{3.335055in}}%
\pgfpathlineto{\pgfqpoint{2.934908in}{3.355954in}}%
\pgfpathlineto{\pgfqpoint{2.928809in}{3.360933in}}%
\pgfpathlineto{\pgfqpoint{2.930876in}{3.384788in}}%
\pgfpathlineto{\pgfqpoint{2.948856in}{3.417967in}}%
\pgfpathlineto{\pgfqpoint{2.951156in}{3.430504in}}%
\pgfpathlineto{\pgfqpoint{2.948193in}{3.460114in}}%
\pgfpathlineto{\pgfqpoint{2.934966in}{3.474061in}}%
\pgfpathlineto{\pgfqpoint{2.928060in}{3.500103in}}%
\pgfpathlineto{\pgfqpoint{2.919562in}{3.506082in}}%
\pgfpathlineto{\pgfqpoint{2.923366in}{3.520204in}}%
\pgfpathlineto{\pgfqpoint{2.912544in}{3.541313in}}%
\pgfpathlineto{\pgfqpoint{2.899043in}{3.552768in}}%
\pgfpathlineto{\pgfqpoint{2.897420in}{3.671978in}}%
\pgfusepath{stroke}%
\end{pgfscope}%
\begin{pgfscope}%
\pgfpathrectangle{\pgfqpoint{0.100000in}{2.493063in}}{\pgfqpoint{5.037500in}{3.427208in}}%
\pgfusepath{clip}%
\pgfsetbuttcap%
\pgfsetroundjoin%
\pgfsetlinewidth{0.501875pt}%
\definecolor{currentstroke}{rgb}{0.827451,0.827451,0.827451}%
\pgfsetstrokecolor{currentstroke}%
\pgfsetdash{}{0pt}%
\pgfpathmoveto{\pgfqpoint{3.106777in}{3.278259in}}%
\pgfpathlineto{\pgfqpoint{3.118102in}{3.284994in}}%
\pgfpathlineto{\pgfqpoint{3.131847in}{3.277031in}}%
\pgfpathlineto{\pgfqpoint{3.124181in}{3.266073in}}%
\pgfpathlineto{\pgfqpoint{3.106777in}{3.278259in}}%
\pgfusepath{stroke}%
\end{pgfscope}%
\begin{pgfscope}%
\pgfpathrectangle{\pgfqpoint{0.100000in}{2.493063in}}{\pgfqpoint{5.037500in}{3.427208in}}%
\pgfusepath{clip}%
\pgfsetbuttcap%
\pgfsetroundjoin%
\pgfsetlinewidth{0.501875pt}%
\definecolor{currentstroke}{rgb}{0.827451,0.827451,0.827451}%
\pgfsetstrokecolor{currentstroke}%
\pgfsetdash{}{0pt}%
\pgfpathmoveto{\pgfqpoint{3.563787in}{3.408159in}}%
\pgfpathlineto{\pgfqpoint{3.563992in}{3.429936in}}%
\pgfpathlineto{\pgfqpoint{3.550461in}{3.438206in}}%
\pgfpathlineto{\pgfqpoint{3.539374in}{3.452285in}}%
\pgfpathlineto{\pgfqpoint{3.540884in}{3.467076in}}%
\pgfpathlineto{\pgfqpoint{3.661580in}{3.476445in}}%
\pgfpathlineto{\pgfqpoint{3.798503in}{3.491653in}}%
\pgfpathlineto{\pgfqpoint{3.815967in}{3.459825in}}%
\pgfpathlineto{\pgfqpoint{3.874037in}{3.463913in}}%
\pgfpathlineto{\pgfqpoint{3.989179in}{3.470195in}}%
\pgfpathlineto{\pgfqpoint{4.080494in}{3.476281in}}%
\pgfpathlineto{\pgfqpoint{4.089462in}{3.453357in}}%
\pgfpathlineto{\pgfqpoint{4.100548in}{3.455235in}}%
\pgfpathlineto{\pgfqpoint{4.101781in}{3.479780in}}%
\pgfpathlineto{\pgfqpoint{4.096108in}{3.500892in}}%
\pgfpathlineto{\pgfqpoint{4.103711in}{3.509829in}}%
\pgfpathlineto{\pgfqpoint{4.121284in}{3.505022in}}%
\pgfpathlineto{\pgfqpoint{4.146508in}{3.504034in}}%
\pgfpathlineto{\pgfqpoint{4.150386in}{3.485172in}}%
\pgfpathlineto{\pgfqpoint{4.158145in}{3.474394in}}%
\pgfpathlineto{\pgfqpoint{4.164205in}{3.450814in}}%
\pgfpathlineto{\pgfqpoint{4.182979in}{3.414321in}}%
\pgfpathlineto{\pgfqpoint{4.183073in}{3.404452in}}%
\pgfpathlineto{\pgfqpoint{4.210820in}{3.361602in}}%
\pgfpathlineto{\pgfqpoint{4.213504in}{3.352750in}}%
\pgfpathlineto{\pgfqpoint{4.255129in}{3.292232in}}%
\pgfpathlineto{\pgfqpoint{4.248535in}{3.291255in}}%
\pgfpathlineto{\pgfqpoint{4.266083in}{3.248210in}}%
\pgfpathlineto{\pgfqpoint{4.314290in}{3.173379in}}%
\pgfpathlineto{\pgfqpoint{4.342279in}{3.118358in}}%
\pgfpathlineto{\pgfqpoint{4.351634in}{3.108725in}}%
\pgfpathlineto{\pgfqpoint{4.367664in}{3.074190in}}%
\pgfpathlineto{\pgfqpoint{4.373225in}{3.018891in}}%
\pgfpathlineto{\pgfqpoint{4.375438in}{2.977542in}}%
\pgfpathlineto{\pgfqpoint{4.372765in}{2.951067in}}%
\pgfpathlineto{\pgfqpoint{4.364170in}{2.932141in}}%
\pgfpathlineto{\pgfqpoint{4.368164in}{2.907323in}}%
\pgfpathlineto{\pgfqpoint{4.358910in}{2.887678in}}%
\pgfpathlineto{\pgfqpoint{4.331435in}{2.871570in}}%
\pgfpathlineto{\pgfqpoint{4.313566in}{2.872746in}}%
\pgfpathlineto{\pgfqpoint{4.301962in}{2.864325in}}%
\pgfpathlineto{\pgfqpoint{4.298638in}{2.886808in}}%
\pgfpathlineto{\pgfqpoint{4.279430in}{2.892743in}}%
\pgfpathlineto{\pgfqpoint{4.262210in}{2.924115in}}%
\pgfpathlineto{\pgfqpoint{4.260267in}{2.938401in}}%
\pgfpathlineto{\pgfqpoint{4.229450in}{2.947076in}}%
\pgfpathlineto{\pgfqpoint{4.209592in}{2.945198in}}%
\pgfpathlineto{\pgfqpoint{4.198356in}{2.965927in}}%
\pgfpathlineto{\pgfqpoint{4.185492in}{3.003185in}}%
\pgfpathlineto{\pgfqpoint{4.167640in}{3.010738in}}%
\pgfpathlineto{\pgfqpoint{4.157962in}{3.032089in}}%
\pgfpathlineto{\pgfqpoint{4.134293in}{3.044672in}}%
\pgfpathlineto{\pgfqpoint{4.120612in}{3.060412in}}%
\pgfpathlineto{\pgfqpoint{4.097012in}{3.103203in}}%
\pgfpathlineto{\pgfqpoint{4.094236in}{3.133321in}}%
\pgfpathlineto{\pgfqpoint{4.107097in}{3.152591in}}%
\pgfpathlineto{\pgfqpoint{4.105811in}{3.165982in}}%
\pgfpathlineto{\pgfqpoint{4.090949in}{3.167413in}}%
\pgfpathlineto{\pgfqpoint{4.078532in}{3.176726in}}%
\pgfpathlineto{\pgfqpoint{4.071732in}{3.165353in}}%
\pgfpathlineto{\pgfqpoint{4.083649in}{3.156086in}}%
\pgfpathlineto{\pgfqpoint{4.074976in}{3.139422in}}%
\pgfpathlineto{\pgfqpoint{4.060944in}{3.153253in}}%
\pgfpathlineto{\pgfqpoint{4.062571in}{3.191608in}}%
\pgfpathlineto{\pgfqpoint{4.069340in}{3.222725in}}%
\pgfpathlineto{\pgfqpoint{4.065861in}{3.276118in}}%
\pgfpathlineto{\pgfqpoint{4.051874in}{3.288836in}}%
\pgfpathlineto{\pgfqpoint{4.044840in}{3.305155in}}%
\pgfpathlineto{\pgfqpoint{4.020746in}{3.304791in}}%
\pgfpathlineto{\pgfqpoint{3.996928in}{3.331644in}}%
\pgfpathlineto{\pgfqpoint{3.980959in}{3.339692in}}%
\pgfpathlineto{\pgfqpoint{3.976224in}{3.356697in}}%
\pgfpathlineto{\pgfqpoint{3.960604in}{3.362669in}}%
\pgfpathlineto{\pgfqpoint{3.947660in}{3.381445in}}%
\pgfpathlineto{\pgfqpoint{3.913464in}{3.396770in}}%
\pgfpathlineto{\pgfqpoint{3.886879in}{3.397162in}}%
\pgfpathlineto{\pgfqpoint{3.875306in}{3.391276in}}%
\pgfpathlineto{\pgfqpoint{3.878182in}{3.372901in}}%
\pgfpathlineto{\pgfqpoint{3.866114in}{3.373771in}}%
\pgfpathlineto{\pgfqpoint{3.828991in}{3.348080in}}%
\pgfpathlineto{\pgfqpoint{3.784400in}{3.337897in}}%
\pgfpathlineto{\pgfqpoint{3.783668in}{3.350496in}}%
\pgfpathlineto{\pgfqpoint{3.773786in}{3.362810in}}%
\pgfpathlineto{\pgfqpoint{3.747272in}{3.379810in}}%
\pgfpathlineto{\pgfqpoint{3.709215in}{3.397240in}}%
\pgfpathlineto{\pgfqpoint{3.667877in}{3.406466in}}%
\pgfpathlineto{\pgfqpoint{3.660118in}{3.419013in}}%
\pgfpathlineto{\pgfqpoint{3.645213in}{3.408549in}}%
\pgfpathlineto{\pgfqpoint{3.627266in}{3.406240in}}%
\pgfpathlineto{\pgfqpoint{3.564793in}{3.389780in}}%
\pgfpathlineto{\pgfqpoint{3.563787in}{3.408159in}}%
\pgfusepath{stroke}%
\end{pgfscope}%
\begin{pgfscope}%
\pgfpathrectangle{\pgfqpoint{0.100000in}{2.493063in}}{\pgfqpoint{5.037500in}{3.427208in}}%
\pgfusepath{clip}%
\pgfsetbuttcap%
\pgfsetroundjoin%
\pgfsetlinewidth{0.501875pt}%
\definecolor{currentstroke}{rgb}{0.827451,0.827451,0.827451}%
\pgfsetstrokecolor{currentstroke}%
\pgfsetdash{}{0pt}%
\pgfpathmoveto{\pgfqpoint{4.261195in}{3.293093in}}%
\pgfpathlineto{\pgfqpoint{4.277358in}{3.274291in}}%
\pgfpathlineto{\pgfqpoint{4.258688in}{3.273108in}}%
\pgfpathlineto{\pgfqpoint{4.261195in}{3.293093in}}%
\pgfusepath{stroke}%
\end{pgfscope}%
\begin{pgfscope}%
\pgfpathrectangle{\pgfqpoint{0.100000in}{2.493063in}}{\pgfqpoint{5.037500in}{3.427208in}}%
\pgfusepath{clip}%
\pgfsetbuttcap%
\pgfsetroundjoin%
\pgfsetlinewidth{0.501875pt}%
\definecolor{currentstroke}{rgb}{0.827451,0.827451,0.827451}%
\pgfsetstrokecolor{currentstroke}%
\pgfsetdash{}{0pt}%
\pgfpathmoveto{\pgfqpoint{3.309667in}{5.439952in}}%
\pgfpathlineto{\pgfqpoint{3.300898in}{5.422883in}}%
\pgfpathlineto{\pgfqpoint{3.279976in}{5.412921in}}%
\pgfpathlineto{\pgfqpoint{3.270914in}{5.399513in}}%
\pgfpathlineto{\pgfqpoint{3.260299in}{5.409170in}}%
\pgfpathlineto{\pgfqpoint{3.309667in}{5.439952in}}%
\pgfusepath{stroke}%
\end{pgfscope}%
\begin{pgfscope}%
\pgfpathrectangle{\pgfqpoint{0.100000in}{2.493063in}}{\pgfqpoint{5.037500in}{3.427208in}}%
\pgfusepath{clip}%
\pgfsetbuttcap%
\pgfsetroundjoin%
\pgfsetlinewidth{0.501875pt}%
\definecolor{currentstroke}{rgb}{0.827451,0.827451,0.827451}%
\pgfsetstrokecolor{currentstroke}%
\pgfsetdash{}{0pt}%
\pgfpathmoveto{\pgfqpoint{3.316781in}{5.337228in}}%
\pgfpathlineto{\pgfqpoint{3.338188in}{5.357198in}}%
\pgfpathlineto{\pgfqpoint{3.371207in}{5.362428in}}%
\pgfpathlineto{\pgfqpoint{3.362113in}{5.348539in}}%
\pgfpathlineto{\pgfqpoint{3.339476in}{5.328453in}}%
\pgfpathlineto{\pgfqpoint{3.326242in}{5.302673in}}%
\pgfpathlineto{\pgfqpoint{3.320637in}{5.316656in}}%
\pgfpathlineto{\pgfqpoint{3.310685in}{5.318663in}}%
\pgfpathlineto{\pgfqpoint{3.309685in}{5.331287in}}%
\pgfpathlineto{\pgfqpoint{3.316781in}{5.337228in}}%
\pgfusepath{stroke}%
\end{pgfscope}%
\begin{pgfscope}%
\pgfpathrectangle{\pgfqpoint{0.100000in}{2.493063in}}{\pgfqpoint{5.037500in}{3.427208in}}%
\pgfusepath{clip}%
\pgfsetbuttcap%
\pgfsetroundjoin%
\pgfsetlinewidth{0.501875pt}%
\definecolor{currentstroke}{rgb}{0.827451,0.827451,0.827451}%
\pgfsetstrokecolor{currentstroke}%
\pgfsetdash{}{0pt}%
\pgfpathmoveto{\pgfqpoint{3.402044in}{5.093479in}}%
\pgfpathlineto{\pgfqpoint{3.396343in}{5.099806in}}%
\pgfpathlineto{\pgfqpoint{3.402337in}{5.117617in}}%
\pgfpathlineto{\pgfqpoint{3.384545in}{5.118752in}}%
\pgfpathlineto{\pgfqpoint{3.389261in}{5.134109in}}%
\pgfpathlineto{\pgfqpoint{3.386328in}{5.158566in}}%
\pgfpathlineto{\pgfqpoint{3.370364in}{5.167046in}}%
\pgfpathlineto{\pgfqpoint{3.353658in}{5.184206in}}%
\pgfpathlineto{\pgfqpoint{3.327916in}{5.189231in}}%
\pgfpathlineto{\pgfqpoint{3.302863in}{5.189087in}}%
\pgfpathlineto{\pgfqpoint{3.278223in}{5.201266in}}%
\pgfpathlineto{\pgfqpoint{3.195789in}{5.219018in}}%
\pgfpathlineto{\pgfqpoint{3.186780in}{5.237828in}}%
\pgfpathlineto{\pgfqpoint{3.170712in}{5.244247in}}%
\pgfpathlineto{\pgfqpoint{3.201064in}{5.258648in}}%
\pgfpathlineto{\pgfqpoint{3.218211in}{5.276616in}}%
\pgfpathlineto{\pgfqpoint{3.250147in}{5.281489in}}%
\pgfpathlineto{\pgfqpoint{3.269800in}{5.299789in}}%
\pgfpathlineto{\pgfqpoint{3.280110in}{5.300522in}}%
\pgfpathlineto{\pgfqpoint{3.287979in}{5.313574in}}%
\pgfpathlineto{\pgfqpoint{3.307161in}{5.329035in}}%
\pgfpathlineto{\pgfqpoint{3.308833in}{5.318107in}}%
\pgfpathlineto{\pgfqpoint{3.317476in}{5.315850in}}%
\pgfpathlineto{\pgfqpoint{3.323979in}{5.302817in}}%
\pgfpathlineto{\pgfqpoint{3.325151in}{5.280597in}}%
\pgfpathlineto{\pgfqpoint{3.344696in}{5.293876in}}%
\pgfpathlineto{\pgfqpoint{3.367471in}{5.296649in}}%
\pgfpathlineto{\pgfqpoint{3.386912in}{5.289693in}}%
\pgfpathlineto{\pgfqpoint{3.413276in}{5.253524in}}%
\pgfpathlineto{\pgfqpoint{3.442071in}{5.259304in}}%
\pgfpathlineto{\pgfqpoint{3.453772in}{5.249624in}}%
\pgfpathlineto{\pgfqpoint{3.472590in}{5.248760in}}%
\pgfpathlineto{\pgfqpoint{3.485093in}{5.266120in}}%
\pgfpathlineto{\pgfqpoint{3.508817in}{5.281466in}}%
\pgfpathlineto{\pgfqpoint{3.531619in}{5.286272in}}%
\pgfpathlineto{\pgfqpoint{3.559912in}{5.286773in}}%
\pgfpathlineto{\pgfqpoint{3.580587in}{5.298627in}}%
\pgfpathlineto{\pgfqpoint{3.597460in}{5.293167in}}%
\pgfpathlineto{\pgfqpoint{3.601054in}{5.268095in}}%
\pgfpathlineto{\pgfqpoint{3.637253in}{5.264169in}}%
\pgfpathlineto{\pgfqpoint{3.656916in}{5.275905in}}%
\pgfpathlineto{\pgfqpoint{3.670546in}{5.249467in}}%
\pgfpathlineto{\pgfqpoint{3.696553in}{5.218895in}}%
\pgfpathlineto{\pgfqpoint{3.660152in}{5.219072in}}%
\pgfpathlineto{\pgfqpoint{3.648651in}{5.215292in}}%
\pgfpathlineto{\pgfqpoint{3.632886in}{5.220204in}}%
\pgfpathlineto{\pgfqpoint{3.631825in}{5.199039in}}%
\pgfpathlineto{\pgfqpoint{3.603200in}{5.215579in}}%
\pgfpathlineto{\pgfqpoint{3.566388in}{5.220633in}}%
\pgfpathlineto{\pgfqpoint{3.556253in}{5.204526in}}%
\pgfpathlineto{\pgfqpoint{3.508088in}{5.196693in}}%
\pgfpathlineto{\pgfqpoint{3.502546in}{5.183048in}}%
\pgfpathlineto{\pgfqpoint{3.469029in}{5.178953in}}%
\pgfpathlineto{\pgfqpoint{3.458896in}{5.165082in}}%
\pgfpathlineto{\pgfqpoint{3.441168in}{5.161363in}}%
\pgfpathlineto{\pgfqpoint{3.427007in}{5.128470in}}%
\pgfpathlineto{\pgfqpoint{3.409079in}{5.096614in}}%
\pgfpathlineto{\pgfqpoint{3.402044in}{5.093479in}}%
\pgfusepath{stroke}%
\end{pgfscope}%
\begin{pgfscope}%
\pgfpathrectangle{\pgfqpoint{0.100000in}{2.493063in}}{\pgfqpoint{5.037500in}{3.427208in}}%
\pgfusepath{clip}%
\pgfsetbuttcap%
\pgfsetroundjoin%
\pgfsetlinewidth{0.501875pt}%
\definecolor{currentstroke}{rgb}{0.827451,0.827451,0.827451}%
\pgfsetstrokecolor{currentstroke}%
\pgfsetdash{}{0pt}%
\pgfpathmoveto{\pgfqpoint{3.505233in}{4.711021in}}%
\pgfpathlineto{\pgfqpoint{3.522326in}{4.729013in}}%
\pgfpathlineto{\pgfqpoint{3.530113in}{4.755098in}}%
\pgfpathlineto{\pgfqpoint{3.539374in}{4.770218in}}%
\pgfpathlineto{\pgfqpoint{3.545053in}{4.790794in}}%
\pgfpathlineto{\pgfqpoint{3.546819in}{4.831821in}}%
\pgfpathlineto{\pgfqpoint{3.538268in}{4.871140in}}%
\pgfpathlineto{\pgfqpoint{3.510013in}{4.931348in}}%
\pgfpathlineto{\pgfqpoint{3.517665in}{4.950279in}}%
\pgfpathlineto{\pgfqpoint{3.507697in}{4.976446in}}%
\pgfpathlineto{\pgfqpoint{3.524796in}{5.012634in}}%
\pgfpathlineto{\pgfqpoint{3.522010in}{5.052984in}}%
\pgfpathlineto{\pgfqpoint{3.533920in}{5.058063in}}%
\pgfpathlineto{\pgfqpoint{3.535416in}{5.077262in}}%
\pgfpathlineto{\pgfqpoint{3.556558in}{5.089585in}}%
\pgfpathlineto{\pgfqpoint{3.568515in}{5.087596in}}%
\pgfpathlineto{\pgfqpoint{3.571876in}{5.067005in}}%
\pgfpathlineto{\pgfqpoint{3.581208in}{5.066198in}}%
\pgfpathlineto{\pgfqpoint{3.589723in}{5.095898in}}%
\pgfpathlineto{\pgfqpoint{3.586829in}{5.118997in}}%
\pgfpathlineto{\pgfqpoint{3.592365in}{5.132261in}}%
\pgfpathlineto{\pgfqpoint{3.622298in}{5.145964in}}%
\pgfpathlineto{\pgfqpoint{3.608658in}{5.150841in}}%
\pgfpathlineto{\pgfqpoint{3.604215in}{5.162612in}}%
\pgfpathlineto{\pgfqpoint{3.614041in}{5.183283in}}%
\pgfpathlineto{\pgfqpoint{3.633417in}{5.190422in}}%
\pgfpathlineto{\pgfqpoint{3.655918in}{5.178226in}}%
\pgfpathlineto{\pgfqpoint{3.677133in}{5.178074in}}%
\pgfpathlineto{\pgfqpoint{3.686985in}{5.163827in}}%
\pgfpathlineto{\pgfqpoint{3.701842in}{5.164817in}}%
\pgfpathlineto{\pgfqpoint{3.747705in}{5.145012in}}%
\pgfpathlineto{\pgfqpoint{3.756869in}{5.125840in}}%
\pgfpathlineto{\pgfqpoint{3.746861in}{5.119200in}}%
\pgfpathlineto{\pgfqpoint{3.749945in}{5.104832in}}%
\pgfpathlineto{\pgfqpoint{3.759831in}{5.098442in}}%
\pgfpathlineto{\pgfqpoint{3.765341in}{5.080779in}}%
\pgfpathlineto{\pgfqpoint{3.764570in}{5.037748in}}%
\pgfpathlineto{\pgfqpoint{3.751517in}{5.027469in}}%
\pgfpathlineto{\pgfqpoint{3.748587in}{5.004866in}}%
\pgfpathlineto{\pgfqpoint{3.724374in}{4.983977in}}%
\pgfpathlineto{\pgfqpoint{3.725811in}{4.958692in}}%
\pgfpathlineto{\pgfqpoint{3.747033in}{4.949835in}}%
\pgfpathlineto{\pgfqpoint{3.770974in}{4.981382in}}%
\pgfpathlineto{\pgfqpoint{3.772953in}{4.992824in}}%
\pgfpathlineto{\pgfqpoint{3.802811in}{5.011842in}}%
\pgfpathlineto{\pgfqpoint{3.821798in}{5.003033in}}%
\pgfpathlineto{\pgfqpoint{3.833713in}{4.983124in}}%
\pgfpathlineto{\pgfqpoint{3.852934in}{4.914001in}}%
\pgfpathlineto{\pgfqpoint{3.863121in}{4.892132in}}%
\pgfpathlineto{\pgfqpoint{3.859932in}{4.883091in}}%
\pgfpathlineto{\pgfqpoint{3.860241in}{4.852307in}}%
\pgfpathlineto{\pgfqpoint{3.841712in}{4.855262in}}%
\pgfpathlineto{\pgfqpoint{3.831247in}{4.832258in}}%
\pgfpathlineto{\pgfqpoint{3.829796in}{4.816619in}}%
\pgfpathlineto{\pgfqpoint{3.815800in}{4.806580in}}%
\pgfpathlineto{\pgfqpoint{3.812690in}{4.776081in}}%
\pgfpathlineto{\pgfqpoint{3.792336in}{4.737553in}}%
\pgfpathlineto{\pgfqpoint{3.680994in}{4.720971in}}%
\pgfpathlineto{\pgfqpoint{3.680336in}{4.728259in}}%
\pgfpathlineto{\pgfqpoint{3.605857in}{4.720409in}}%
\pgfpathlineto{\pgfqpoint{3.505233in}{4.711021in}}%
\pgfusepath{stroke}%
\end{pgfscope}%
\begin{pgfscope}%
\pgfpathrectangle{\pgfqpoint{0.100000in}{2.493063in}}{\pgfqpoint{5.037500in}{3.427208in}}%
\pgfusepath{clip}%
\pgfsetbuttcap%
\pgfsetroundjoin%
\pgfsetlinewidth{0.501875pt}%
\definecolor{currentstroke}{rgb}{0.827451,0.827451,0.827451}%
\pgfsetstrokecolor{currentstroke}%
\pgfsetdash{}{0pt}%
\pgfpathmoveto{\pgfqpoint{0.000000in}{0.000000in}}%
\pgfusepath{stroke}%
\end{pgfscope}%
\begin{pgfscope}%
\pgfpathrectangle{\pgfqpoint{0.100000in}{2.493063in}}{\pgfqpoint{5.037500in}{3.427208in}}%
\pgfusepath{clip}%
\pgfsetbuttcap%
\pgfsetroundjoin%
\pgfsetlinewidth{0.501875pt}%
\definecolor{currentstroke}{rgb}{0.827451,0.827451,0.827451}%
\pgfsetstrokecolor{currentstroke}%
\pgfsetdash{}{0pt}%
\pgfusepath{stroke}%
\end{pgfscope}%
\begin{pgfscope}%
\pgfpathrectangle{\pgfqpoint{0.100000in}{2.493063in}}{\pgfqpoint{5.037500in}{3.427208in}}%
\pgfusepath{clip}%
\pgfsetbuttcap%
\pgfsetroundjoin%
\pgfsetlinewidth{0.501875pt}%
\definecolor{currentstroke}{rgb}{0.827451,0.827451,0.827451}%
\pgfsetstrokecolor{currentstroke}%
\pgfsetdash{}{0pt}%
\pgfusepath{stroke}%
\end{pgfscope}%
\begin{pgfscope}%
\pgfpathrectangle{\pgfqpoint{0.100000in}{2.493063in}}{\pgfqpoint{5.037500in}{3.427208in}}%
\pgfusepath{clip}%
\pgfsetbuttcap%
\pgfsetroundjoin%
\pgfsetlinewidth{0.501875pt}%
\definecolor{currentstroke}{rgb}{0.827451,0.827451,0.827451}%
\pgfsetstrokecolor{currentstroke}%
\pgfsetdash{}{0pt}%
\pgfusepath{stroke}%
\end{pgfscope}%
\begin{pgfscope}%
\pgfpathrectangle{\pgfqpoint{0.100000in}{2.493063in}}{\pgfqpoint{5.037500in}{3.427208in}}%
\pgfusepath{clip}%
\pgfsetbuttcap%
\pgfsetroundjoin%
\pgfsetlinewidth{0.501875pt}%
\definecolor{currentstroke}{rgb}{0.827451,0.827451,0.827451}%
\pgfsetstrokecolor{currentstroke}%
\pgfsetdash{}{0pt}%
\pgfusepath{stroke}%
\end{pgfscope}%
\begin{pgfscope}%
\pgfpathrectangle{\pgfqpoint{0.100000in}{2.493063in}}{\pgfqpoint{5.037500in}{3.427208in}}%
\pgfusepath{clip}%
\pgfsetbuttcap%
\pgfsetroundjoin%
\pgfsetlinewidth{0.501875pt}%
\definecolor{currentstroke}{rgb}{0.827451,0.827451,0.827451}%
\pgfsetstrokecolor{currentstroke}%
\pgfsetdash{}{0pt}%
\pgfusepath{stroke}%
\end{pgfscope}%
\begin{pgfscope}%
\pgfpathrectangle{\pgfqpoint{0.100000in}{2.493063in}}{\pgfqpoint{5.037500in}{3.427208in}}%
\pgfusepath{clip}%
\pgfsetbuttcap%
\pgfsetroundjoin%
\pgfsetlinewidth{0.501875pt}%
\definecolor{currentstroke}{rgb}{0.827451,0.827451,0.827451}%
\pgfsetstrokecolor{currentstroke}%
\pgfsetdash{}{0pt}%
\pgfusepath{stroke}%
\end{pgfscope}%
\begin{pgfscope}%
\pgfpathrectangle{\pgfqpoint{0.100000in}{2.493063in}}{\pgfqpoint{5.037500in}{3.427208in}}%
\pgfusepath{clip}%
\pgfsetbuttcap%
\pgfsetroundjoin%
\pgfsetlinewidth{0.501875pt}%
\definecolor{currentstroke}{rgb}{0.827451,0.827451,0.827451}%
\pgfsetstrokecolor{currentstroke}%
\pgfsetdash{}{0pt}%
\pgfusepath{stroke}%
\end{pgfscope}%
\begin{pgfscope}%
\pgfpathrectangle{\pgfqpoint{0.100000in}{2.493063in}}{\pgfqpoint{5.037500in}{3.427208in}}%
\pgfusepath{clip}%
\pgfsetbuttcap%
\pgfsetroundjoin%
\pgfsetlinewidth{0.501875pt}%
\definecolor{currentstroke}{rgb}{0.827451,0.827451,0.827451}%
\pgfsetstrokecolor{currentstroke}%
\pgfsetdash{}{0pt}%
\pgfusepath{stroke}%
\end{pgfscope}%
\begin{pgfscope}%
\pgfpathrectangle{\pgfqpoint{0.100000in}{2.493063in}}{\pgfqpoint{5.037500in}{3.427208in}}%
\pgfusepath{clip}%
\pgfsetbuttcap%
\pgfsetroundjoin%
\pgfsetlinewidth{0.501875pt}%
\definecolor{currentstroke}{rgb}{0.827451,0.827451,0.827451}%
\pgfsetstrokecolor{currentstroke}%
\pgfsetdash{}{0pt}%
\pgfusepath{stroke}%
\end{pgfscope}%
\begin{pgfscope}%
\pgfpathrectangle{\pgfqpoint{0.100000in}{2.493063in}}{\pgfqpoint{5.037500in}{3.427208in}}%
\pgfusepath{clip}%
\pgfsetbuttcap%
\pgfsetroundjoin%
\pgfsetlinewidth{0.501875pt}%
\definecolor{currentstroke}{rgb}{0.827451,0.827451,0.827451}%
\pgfsetstrokecolor{currentstroke}%
\pgfsetdash{}{0pt}%
\pgfusepath{stroke}%
\end{pgfscope}%
\begin{pgfscope}%
\pgfpathrectangle{\pgfqpoint{0.100000in}{2.493063in}}{\pgfqpoint{5.037500in}{3.427208in}}%
\pgfusepath{clip}%
\pgfsetbuttcap%
\pgfsetroundjoin%
\pgfsetlinewidth{0.501875pt}%
\definecolor{currentstroke}{rgb}{0.827451,0.827451,0.827451}%
\pgfsetstrokecolor{currentstroke}%
\pgfsetdash{}{0pt}%
\pgfusepath{stroke}%
\end{pgfscope}%
\begin{pgfscope}%
\pgfpathrectangle{\pgfqpoint{0.100000in}{2.493063in}}{\pgfqpoint{5.037500in}{3.427208in}}%
\pgfusepath{clip}%
\pgfsetbuttcap%
\pgfsetroundjoin%
\pgfsetlinewidth{0.501875pt}%
\definecolor{currentstroke}{rgb}{0.827451,0.827451,0.827451}%
\pgfsetstrokecolor{currentstroke}%
\pgfsetdash{}{0pt}%
\pgfusepath{stroke}%
\end{pgfscope}%
\begin{pgfscope}%
\pgfpathrectangle{\pgfqpoint{0.100000in}{2.493063in}}{\pgfqpoint{5.037500in}{3.427208in}}%
\pgfusepath{clip}%
\pgfsetbuttcap%
\pgfsetroundjoin%
\pgfsetlinewidth{0.501875pt}%
\definecolor{currentstroke}{rgb}{0.827451,0.827451,0.827451}%
\pgfsetstrokecolor{currentstroke}%
\pgfsetdash{}{0pt}%
\pgfusepath{stroke}%
\end{pgfscope}%
\begin{pgfscope}%
\pgfpathrectangle{\pgfqpoint{0.100000in}{2.493063in}}{\pgfqpoint{5.037500in}{3.427208in}}%
\pgfusepath{clip}%
\pgfsetbuttcap%
\pgfsetroundjoin%
\pgfsetlinewidth{0.501875pt}%
\definecolor{currentstroke}{rgb}{0.827451,0.827451,0.827451}%
\pgfsetstrokecolor{currentstroke}%
\pgfsetdash{}{0pt}%
\pgfusepath{stroke}%
\end{pgfscope}%
\begin{pgfscope}%
\pgfpathrectangle{\pgfqpoint{0.100000in}{2.493063in}}{\pgfqpoint{5.037500in}{3.427208in}}%
\pgfusepath{clip}%
\pgfsetbuttcap%
\pgfsetroundjoin%
\pgfsetlinewidth{0.501875pt}%
\definecolor{currentstroke}{rgb}{0.827451,0.827451,0.827451}%
\pgfsetstrokecolor{currentstroke}%
\pgfsetdash{}{0pt}%
\pgfusepath{stroke}%
\end{pgfscope}%
\begin{pgfscope}%
\pgfpathrectangle{\pgfqpoint{0.100000in}{2.493063in}}{\pgfqpoint{5.037500in}{3.427208in}}%
\pgfusepath{clip}%
\pgfsetbuttcap%
\pgfsetroundjoin%
\pgfsetlinewidth{0.501875pt}%
\definecolor{currentstroke}{rgb}{0.827451,0.827451,0.827451}%
\pgfsetstrokecolor{currentstroke}%
\pgfsetdash{}{0pt}%
\pgfusepath{stroke}%
\end{pgfscope}%
\begin{pgfscope}%
\pgfpathrectangle{\pgfqpoint{0.100000in}{2.493063in}}{\pgfqpoint{5.037500in}{3.427208in}}%
\pgfusepath{clip}%
\pgfsetbuttcap%
\pgfsetroundjoin%
\pgfsetlinewidth{0.501875pt}%
\definecolor{currentstroke}{rgb}{0.827451,0.827451,0.827451}%
\pgfsetstrokecolor{currentstroke}%
\pgfsetdash{}{0pt}%
\pgfusepath{stroke}%
\end{pgfscope}%
\begin{pgfscope}%
\pgfpathrectangle{\pgfqpoint{0.100000in}{2.493063in}}{\pgfqpoint{5.037500in}{3.427208in}}%
\pgfusepath{clip}%
\pgfsetbuttcap%
\pgfsetroundjoin%
\pgfsetlinewidth{0.501875pt}%
\definecolor{currentstroke}{rgb}{0.827451,0.827451,0.827451}%
\pgfsetstrokecolor{currentstroke}%
\pgfsetdash{}{0pt}%
\pgfusepath{stroke}%
\end{pgfscope}%
\begin{pgfscope}%
\pgfpathrectangle{\pgfqpoint{0.100000in}{2.493063in}}{\pgfqpoint{5.037500in}{3.427208in}}%
\pgfusepath{clip}%
\pgfsetbuttcap%
\pgfsetroundjoin%
\pgfsetlinewidth{0.501875pt}%
\definecolor{currentstroke}{rgb}{0.827451,0.827451,0.827451}%
\pgfsetstrokecolor{currentstroke}%
\pgfsetdash{}{0pt}%
\pgfusepath{stroke}%
\end{pgfscope}%
\begin{pgfscope}%
\pgfpathrectangle{\pgfqpoint{0.100000in}{2.493063in}}{\pgfqpoint{5.037500in}{3.427208in}}%
\pgfusepath{clip}%
\pgfsetbuttcap%
\pgfsetroundjoin%
\pgfsetlinewidth{0.501875pt}%
\definecolor{currentstroke}{rgb}{0.827451,0.827451,0.827451}%
\pgfsetstrokecolor{currentstroke}%
\pgfsetdash{}{0pt}%
\pgfusepath{stroke}%
\end{pgfscope}%
\begin{pgfscope}%
\pgfpathrectangle{\pgfqpoint{0.100000in}{2.493063in}}{\pgfqpoint{5.037500in}{3.427208in}}%
\pgfusepath{clip}%
\pgfsetbuttcap%
\pgfsetroundjoin%
\pgfsetlinewidth{0.501875pt}%
\definecolor{currentstroke}{rgb}{0.827451,0.827451,0.827451}%
\pgfsetstrokecolor{currentstroke}%
\pgfsetdash{}{0pt}%
\pgfusepath{stroke}%
\end{pgfscope}%
\begin{pgfscope}%
\pgfpathrectangle{\pgfqpoint{0.100000in}{2.493063in}}{\pgfqpoint{5.037500in}{3.427208in}}%
\pgfusepath{clip}%
\pgfsetbuttcap%
\pgfsetroundjoin%
\pgfsetlinewidth{0.501875pt}%
\definecolor{currentstroke}{rgb}{0.827451,0.827451,0.827451}%
\pgfsetstrokecolor{currentstroke}%
\pgfsetdash{}{0pt}%
\pgfusepath{stroke}%
\end{pgfscope}%
\begin{pgfscope}%
\pgfpathrectangle{\pgfqpoint{0.100000in}{2.493063in}}{\pgfqpoint{5.037500in}{3.427208in}}%
\pgfusepath{clip}%
\pgfsetbuttcap%
\pgfsetroundjoin%
\pgfsetlinewidth{0.501875pt}%
\definecolor{currentstroke}{rgb}{0.827451,0.827451,0.827451}%
\pgfsetstrokecolor{currentstroke}%
\pgfsetdash{}{0pt}%
\pgfusepath{stroke}%
\end{pgfscope}%
\begin{pgfscope}%
\pgfpathrectangle{\pgfqpoint{0.100000in}{2.493063in}}{\pgfqpoint{5.037500in}{3.427208in}}%
\pgfusepath{clip}%
\pgfsetbuttcap%
\pgfsetroundjoin%
\pgfsetlinewidth{0.501875pt}%
\definecolor{currentstroke}{rgb}{0.827451,0.827451,0.827451}%
\pgfsetstrokecolor{currentstroke}%
\pgfsetdash{}{0pt}%
\pgfusepath{stroke}%
\end{pgfscope}%
\begin{pgfscope}%
\pgfpathrectangle{\pgfqpoint{0.100000in}{2.493063in}}{\pgfqpoint{5.037500in}{3.427208in}}%
\pgfusepath{clip}%
\pgfsetbuttcap%
\pgfsetroundjoin%
\pgfsetlinewidth{0.501875pt}%
\definecolor{currentstroke}{rgb}{0.827451,0.827451,0.827451}%
\pgfsetstrokecolor{currentstroke}%
\pgfsetdash{}{0pt}%
\pgfusepath{stroke}%
\end{pgfscope}%
\begin{pgfscope}%
\pgfpathrectangle{\pgfqpoint{0.100000in}{2.493063in}}{\pgfqpoint{5.037500in}{3.427208in}}%
\pgfusepath{clip}%
\pgfsetbuttcap%
\pgfsetroundjoin%
\pgfsetlinewidth{0.501875pt}%
\definecolor{currentstroke}{rgb}{0.827451,0.827451,0.827451}%
\pgfsetstrokecolor{currentstroke}%
\pgfsetdash{}{0pt}%
\pgfusepath{stroke}%
\end{pgfscope}%
\begin{pgfscope}%
\pgfpathrectangle{\pgfqpoint{0.100000in}{2.493063in}}{\pgfqpoint{5.037500in}{3.427208in}}%
\pgfusepath{clip}%
\pgfsetbuttcap%
\pgfsetroundjoin%
\pgfsetlinewidth{0.501875pt}%
\definecolor{currentstroke}{rgb}{0.827451,0.827451,0.827451}%
\pgfsetstrokecolor{currentstroke}%
\pgfsetdash{}{0pt}%
\pgfusepath{stroke}%
\end{pgfscope}%
\begin{pgfscope}%
\pgfpathrectangle{\pgfqpoint{0.100000in}{2.493063in}}{\pgfqpoint{5.037500in}{3.427208in}}%
\pgfusepath{clip}%
\pgfsetrectcap%
\pgfsetroundjoin%
\pgfsetlinewidth{1.505625pt}%
\definecolor{currentstroke}{rgb}{0.678431,1.000000,0.184314}%
\pgfsetstrokecolor{currentstroke}%
\pgfsetstrokeopacity{0.500000}%
\pgfsetdash{}{0pt}%
\pgfpathmoveto{\pgfqpoint{3.684003in}{3.793619in}}%
\pgfusepath{stroke}%
\end{pgfscope}%
\begin{pgfscope}%
\pgfpathrectangle{\pgfqpoint{0.100000in}{2.493063in}}{\pgfqpoint{5.037500in}{3.427208in}}%
\pgfusepath{clip}%
\pgfsetbuttcap%
\pgfsetroundjoin%
\definecolor{currentfill}{rgb}{0.678431,1.000000,0.184314}%
\pgfsetfillcolor{currentfill}%
\pgfsetfillopacity{0.500000}%
\pgfsetlinewidth{0.250937pt}%
\definecolor{currentstroke}{rgb}{0.000000,0.000000,0.000000}%
\pgfsetstrokecolor{currentstroke}%
\pgfsetstrokeopacity{0.500000}%
\pgfsetdash{}{0pt}%
\pgfsys@defobject{currentmarker}{\pgfqpoint{-0.008333in}{-0.008333in}}{\pgfqpoint{0.008333in}{0.008333in}}{%
\pgfpathmoveto{\pgfqpoint{0.000000in}{-0.008333in}}%
\pgfpathcurveto{\pgfqpoint{0.002210in}{-0.008333in}}{\pgfqpoint{0.004330in}{-0.007455in}}{\pgfqpoint{0.005893in}{-0.005893in}}%
\pgfpathcurveto{\pgfqpoint{0.007455in}{-0.004330in}}{\pgfqpoint{0.008333in}{-0.002210in}}{\pgfqpoint{0.008333in}{0.000000in}}%
\pgfpathcurveto{\pgfqpoint{0.008333in}{0.002210in}}{\pgfqpoint{0.007455in}{0.004330in}}{\pgfqpoint{0.005893in}{0.005893in}}%
\pgfpathcurveto{\pgfqpoint{0.004330in}{0.007455in}}{\pgfqpoint{0.002210in}{0.008333in}}{\pgfqpoint{0.000000in}{0.008333in}}%
\pgfpathcurveto{\pgfqpoint{-0.002210in}{0.008333in}}{\pgfqpoint{-0.004330in}{0.007455in}}{\pgfqpoint{-0.005893in}{0.005893in}}%
\pgfpathcurveto{\pgfqpoint{-0.007455in}{0.004330in}}{\pgfqpoint{-0.008333in}{0.002210in}}{\pgfqpoint{-0.008333in}{0.000000in}}%
\pgfpathcurveto{\pgfqpoint{-0.008333in}{-0.002210in}}{\pgfqpoint{-0.007455in}{-0.004330in}}{\pgfqpoint{-0.005893in}{-0.005893in}}%
\pgfpathcurveto{\pgfqpoint{-0.004330in}{-0.007455in}}{\pgfqpoint{-0.002210in}{-0.008333in}}{\pgfqpoint{0.000000in}{-0.008333in}}%
\pgfpathlineto{\pgfqpoint{0.000000in}{-0.008333in}}%
\pgfpathclose%
\pgfusepath{stroke,fill}%
}%
\begin{pgfscope}%
\pgfsys@transformshift{3.684003in}{3.793619in}%
\pgfsys@useobject{currentmarker}{}%
\end{pgfscope}%
\end{pgfscope}%
\begin{pgfscope}%
\pgfpathrectangle{\pgfqpoint{0.100000in}{2.493063in}}{\pgfqpoint{5.037500in}{3.427208in}}%
\pgfusepath{clip}%
\pgfsetrectcap%
\pgfsetroundjoin%
\pgfsetlinewidth{1.505625pt}%
\definecolor{currentstroke}{rgb}{0.501961,0.501961,0.501961}%
\pgfsetstrokecolor{currentstroke}%
\pgfsetstrokeopacity{0.500000}%
\pgfsetdash{}{0pt}%
\pgfpathmoveto{\pgfqpoint{3.728741in}{3.677179in}}%
\pgfusepath{stroke}%
\end{pgfscope}%
\begin{pgfscope}%
\pgfpathrectangle{\pgfqpoint{0.100000in}{2.493063in}}{\pgfqpoint{5.037500in}{3.427208in}}%
\pgfusepath{clip}%
\pgfsetbuttcap%
\pgfsetroundjoin%
\definecolor{currentfill}{rgb}{0.501961,0.501961,0.501961}%
\pgfsetfillcolor{currentfill}%
\pgfsetfillopacity{0.500000}%
\pgfsetlinewidth{0.250937pt}%
\definecolor{currentstroke}{rgb}{0.000000,0.000000,0.000000}%
\pgfsetstrokecolor{currentstroke}%
\pgfsetstrokeopacity{0.500000}%
\pgfsetdash{}{0pt}%
\pgfsys@defobject{currentmarker}{\pgfqpoint{-0.013889in}{-0.013889in}}{\pgfqpoint{0.013889in}{0.013889in}}{%
\pgfpathmoveto{\pgfqpoint{0.000000in}{-0.013889in}}%
\pgfpathcurveto{\pgfqpoint{0.003683in}{-0.013889in}}{\pgfqpoint{0.007216in}{-0.012425in}}{\pgfqpoint{0.009821in}{-0.009821in}}%
\pgfpathcurveto{\pgfqpoint{0.012425in}{-0.007216in}}{\pgfqpoint{0.013889in}{-0.003683in}}{\pgfqpoint{0.013889in}{0.000000in}}%
\pgfpathcurveto{\pgfqpoint{0.013889in}{0.003683in}}{\pgfqpoint{0.012425in}{0.007216in}}{\pgfqpoint{0.009821in}{0.009821in}}%
\pgfpathcurveto{\pgfqpoint{0.007216in}{0.012425in}}{\pgfqpoint{0.003683in}{0.013889in}}{\pgfqpoint{0.000000in}{0.013889in}}%
\pgfpathcurveto{\pgfqpoint{-0.003683in}{0.013889in}}{\pgfqpoint{-0.007216in}{0.012425in}}{\pgfqpoint{-0.009821in}{0.009821in}}%
\pgfpathcurveto{\pgfqpoint{-0.012425in}{0.007216in}}{\pgfqpoint{-0.013889in}{0.003683in}}{\pgfqpoint{-0.013889in}{0.000000in}}%
\pgfpathcurveto{\pgfqpoint{-0.013889in}{-0.003683in}}{\pgfqpoint{-0.012425in}{-0.007216in}}{\pgfqpoint{-0.009821in}{-0.009821in}}%
\pgfpathcurveto{\pgfqpoint{-0.007216in}{-0.012425in}}{\pgfqpoint{-0.003683in}{-0.013889in}}{\pgfqpoint{0.000000in}{-0.013889in}}%
\pgfpathlineto{\pgfqpoint{0.000000in}{-0.013889in}}%
\pgfpathclose%
\pgfusepath{stroke,fill}%
}%
\begin{pgfscope}%
\pgfsys@transformshift{3.728741in}{3.677179in}%
\pgfsys@useobject{currentmarker}{}%
\end{pgfscope}%
\end{pgfscope}%
\begin{pgfscope}%
\pgfpathrectangle{\pgfqpoint{0.100000in}{2.493063in}}{\pgfqpoint{5.037500in}{3.427208in}}%
\pgfusepath{clip}%
\pgfsetrectcap%
\pgfsetroundjoin%
\pgfsetlinewidth{1.505625pt}%
\definecolor{currentstroke}{rgb}{0.501961,0.501961,0.501961}%
\pgfsetstrokecolor{currentstroke}%
\pgfsetstrokeopacity{0.500000}%
\pgfsetdash{}{0pt}%
\pgfpathmoveto{\pgfqpoint{3.593628in}{3.764664in}}%
\pgfusepath{stroke}%
\end{pgfscope}%
\begin{pgfscope}%
\pgfpathrectangle{\pgfqpoint{0.100000in}{2.493063in}}{\pgfqpoint{5.037500in}{3.427208in}}%
\pgfusepath{clip}%
\pgfsetbuttcap%
\pgfsetroundjoin%
\definecolor{currentfill}{rgb}{0.501961,0.501961,0.501961}%
\pgfsetfillcolor{currentfill}%
\pgfsetfillopacity{0.500000}%
\pgfsetlinewidth{0.250937pt}%
\definecolor{currentstroke}{rgb}{0.000000,0.000000,0.000000}%
\pgfsetstrokecolor{currentstroke}%
\pgfsetstrokeopacity{0.500000}%
\pgfsetdash{}{0pt}%
\pgfsys@defobject{currentmarker}{\pgfqpoint{-0.013889in}{-0.013889in}}{\pgfqpoint{0.013889in}{0.013889in}}{%
\pgfpathmoveto{\pgfqpoint{0.000000in}{-0.013889in}}%
\pgfpathcurveto{\pgfqpoint{0.003683in}{-0.013889in}}{\pgfqpoint{0.007216in}{-0.012425in}}{\pgfqpoint{0.009821in}{-0.009821in}}%
\pgfpathcurveto{\pgfqpoint{0.012425in}{-0.007216in}}{\pgfqpoint{0.013889in}{-0.003683in}}{\pgfqpoint{0.013889in}{0.000000in}}%
\pgfpathcurveto{\pgfqpoint{0.013889in}{0.003683in}}{\pgfqpoint{0.012425in}{0.007216in}}{\pgfqpoint{0.009821in}{0.009821in}}%
\pgfpathcurveto{\pgfqpoint{0.007216in}{0.012425in}}{\pgfqpoint{0.003683in}{0.013889in}}{\pgfqpoint{0.000000in}{0.013889in}}%
\pgfpathcurveto{\pgfqpoint{-0.003683in}{0.013889in}}{\pgfqpoint{-0.007216in}{0.012425in}}{\pgfqpoint{-0.009821in}{0.009821in}}%
\pgfpathcurveto{\pgfqpoint{-0.012425in}{0.007216in}}{\pgfqpoint{-0.013889in}{0.003683in}}{\pgfqpoint{-0.013889in}{0.000000in}}%
\pgfpathcurveto{\pgfqpoint{-0.013889in}{-0.003683in}}{\pgfqpoint{-0.012425in}{-0.007216in}}{\pgfqpoint{-0.009821in}{-0.009821in}}%
\pgfpathcurveto{\pgfqpoint{-0.007216in}{-0.012425in}}{\pgfqpoint{-0.003683in}{-0.013889in}}{\pgfqpoint{0.000000in}{-0.013889in}}%
\pgfpathlineto{\pgfqpoint{0.000000in}{-0.013889in}}%
\pgfpathclose%
\pgfusepath{stroke,fill}%
}%
\begin{pgfscope}%
\pgfsys@transformshift{3.593628in}{3.764664in}%
\pgfsys@useobject{currentmarker}{}%
\end{pgfscope}%
\end{pgfscope}%
\begin{pgfscope}%
\pgfpathrectangle{\pgfqpoint{0.100000in}{2.493063in}}{\pgfqpoint{5.037500in}{3.427208in}}%
\pgfusepath{clip}%
\pgfsetrectcap%
\pgfsetroundjoin%
\pgfsetlinewidth{1.505625pt}%
\definecolor{currentstroke}{rgb}{0.678431,1.000000,0.184314}%
\pgfsetstrokecolor{currentstroke}%
\pgfsetstrokeopacity{0.500000}%
\pgfsetdash{}{0pt}%
\pgfpathmoveto{\pgfqpoint{3.519290in}{3.415074in}}%
\pgfusepath{stroke}%
\end{pgfscope}%
\begin{pgfscope}%
\pgfpathrectangle{\pgfqpoint{0.100000in}{2.493063in}}{\pgfqpoint{5.037500in}{3.427208in}}%
\pgfusepath{clip}%
\pgfsetbuttcap%
\pgfsetroundjoin%
\definecolor{currentfill}{rgb}{0.678431,1.000000,0.184314}%
\pgfsetfillcolor{currentfill}%
\pgfsetfillopacity{0.500000}%
\pgfsetlinewidth{0.250937pt}%
\definecolor{currentstroke}{rgb}{0.000000,0.000000,0.000000}%
\pgfsetstrokecolor{currentstroke}%
\pgfsetstrokeopacity{0.500000}%
\pgfsetdash{}{0pt}%
\pgfsys@defobject{currentmarker}{\pgfqpoint{-0.005556in}{-0.005556in}}{\pgfqpoint{0.005556in}{0.005556in}}{%
\pgfpathmoveto{\pgfqpoint{0.000000in}{-0.005556in}}%
\pgfpathcurveto{\pgfqpoint{0.001473in}{-0.005556in}}{\pgfqpoint{0.002887in}{-0.004970in}}{\pgfqpoint{0.003928in}{-0.003928in}}%
\pgfpathcurveto{\pgfqpoint{0.004970in}{-0.002887in}}{\pgfqpoint{0.005556in}{-0.001473in}}{\pgfqpoint{0.005556in}{0.000000in}}%
\pgfpathcurveto{\pgfqpoint{0.005556in}{0.001473in}}{\pgfqpoint{0.004970in}{0.002887in}}{\pgfqpoint{0.003928in}{0.003928in}}%
\pgfpathcurveto{\pgfqpoint{0.002887in}{0.004970in}}{\pgfqpoint{0.001473in}{0.005556in}}{\pgfqpoint{0.000000in}{0.005556in}}%
\pgfpathcurveto{\pgfqpoint{-0.001473in}{0.005556in}}{\pgfqpoint{-0.002887in}{0.004970in}}{\pgfqpoint{-0.003928in}{0.003928in}}%
\pgfpathcurveto{\pgfqpoint{-0.004970in}{0.002887in}}{\pgfqpoint{-0.005556in}{0.001473in}}{\pgfqpoint{-0.005556in}{0.000000in}}%
\pgfpathcurveto{\pgfqpoint{-0.005556in}{-0.001473in}}{\pgfqpoint{-0.004970in}{-0.002887in}}{\pgfqpoint{-0.003928in}{-0.003928in}}%
\pgfpathcurveto{\pgfqpoint{-0.002887in}{-0.004970in}}{\pgfqpoint{-0.001473in}{-0.005556in}}{\pgfqpoint{0.000000in}{-0.005556in}}%
\pgfpathlineto{\pgfqpoint{0.000000in}{-0.005556in}}%
\pgfpathclose%
\pgfusepath{stroke,fill}%
}%
\begin{pgfscope}%
\pgfsys@transformshift{3.519290in}{3.415074in}%
\pgfsys@useobject{currentmarker}{}%
\end{pgfscope}%
\end{pgfscope}%
\begin{pgfscope}%
\pgfpathrectangle{\pgfqpoint{0.100000in}{2.493063in}}{\pgfqpoint{5.037500in}{3.427208in}}%
\pgfusepath{clip}%
\pgfsetrectcap%
\pgfsetroundjoin%
\pgfsetlinewidth{1.505625pt}%
\definecolor{currentstroke}{rgb}{0.678431,1.000000,0.184314}%
\pgfsetstrokecolor{currentstroke}%
\pgfsetstrokeopacity{0.500000}%
\pgfsetdash{}{0pt}%
\pgfpathmoveto{\pgfqpoint{3.557956in}{3.889517in}}%
\pgfusepath{stroke}%
\end{pgfscope}%
\begin{pgfscope}%
\pgfpathrectangle{\pgfqpoint{0.100000in}{2.493063in}}{\pgfqpoint{5.037500in}{3.427208in}}%
\pgfusepath{clip}%
\pgfsetbuttcap%
\pgfsetroundjoin%
\definecolor{currentfill}{rgb}{0.678431,1.000000,0.184314}%
\pgfsetfillcolor{currentfill}%
\pgfsetfillopacity{0.500000}%
\pgfsetlinewidth{0.250937pt}%
\definecolor{currentstroke}{rgb}{0.000000,0.000000,0.000000}%
\pgfsetstrokecolor{currentstroke}%
\pgfsetstrokeopacity{0.500000}%
\pgfsetdash{}{0pt}%
\pgfsys@defobject{currentmarker}{\pgfqpoint{-0.011111in}{-0.011111in}}{\pgfqpoint{0.011111in}{0.011111in}}{%
\pgfpathmoveto{\pgfqpoint{0.000000in}{-0.011111in}}%
\pgfpathcurveto{\pgfqpoint{0.002947in}{-0.011111in}}{\pgfqpoint{0.005773in}{-0.009940in}}{\pgfqpoint{0.007857in}{-0.007857in}}%
\pgfpathcurveto{\pgfqpoint{0.009940in}{-0.005773in}}{\pgfqpoint{0.011111in}{-0.002947in}}{\pgfqpoint{0.011111in}{0.000000in}}%
\pgfpathcurveto{\pgfqpoint{0.011111in}{0.002947in}}{\pgfqpoint{0.009940in}{0.005773in}}{\pgfqpoint{0.007857in}{0.007857in}}%
\pgfpathcurveto{\pgfqpoint{0.005773in}{0.009940in}}{\pgfqpoint{0.002947in}{0.011111in}}{\pgfqpoint{0.000000in}{0.011111in}}%
\pgfpathcurveto{\pgfqpoint{-0.002947in}{0.011111in}}{\pgfqpoint{-0.005773in}{0.009940in}}{\pgfqpoint{-0.007857in}{0.007857in}}%
\pgfpathcurveto{\pgfqpoint{-0.009940in}{0.005773in}}{\pgfqpoint{-0.011111in}{0.002947in}}{\pgfqpoint{-0.011111in}{0.000000in}}%
\pgfpathcurveto{\pgfqpoint{-0.011111in}{-0.002947in}}{\pgfqpoint{-0.009940in}{-0.005773in}}{\pgfqpoint{-0.007857in}{-0.007857in}}%
\pgfpathcurveto{\pgfqpoint{-0.005773in}{-0.009940in}}{\pgfqpoint{-0.002947in}{-0.011111in}}{\pgfqpoint{0.000000in}{-0.011111in}}%
\pgfpathlineto{\pgfqpoint{0.000000in}{-0.011111in}}%
\pgfpathclose%
\pgfusepath{stroke,fill}%
}%
\begin{pgfscope}%
\pgfsys@transformshift{3.557956in}{3.889517in}%
\pgfsys@useobject{currentmarker}{}%
\end{pgfscope}%
\end{pgfscope}%
\begin{pgfscope}%
\pgfpathrectangle{\pgfqpoint{0.100000in}{2.493063in}}{\pgfqpoint{5.037500in}{3.427208in}}%
\pgfusepath{clip}%
\pgfsetrectcap%
\pgfsetroundjoin%
\pgfsetlinewidth{1.505625pt}%
\definecolor{currentstroke}{rgb}{0.678431,1.000000,0.184314}%
\pgfsetstrokecolor{currentstroke}%
\pgfsetstrokeopacity{0.500000}%
\pgfsetdash{}{0pt}%
\pgfpathmoveto{\pgfqpoint{3.757252in}{3.513137in}}%
\pgfusepath{stroke}%
\end{pgfscope}%
\begin{pgfscope}%
\pgfpathrectangle{\pgfqpoint{0.100000in}{2.493063in}}{\pgfqpoint{5.037500in}{3.427208in}}%
\pgfusepath{clip}%
\pgfsetbuttcap%
\pgfsetroundjoin%
\definecolor{currentfill}{rgb}{0.678431,1.000000,0.184314}%
\pgfsetfillcolor{currentfill}%
\pgfsetfillopacity{0.500000}%
\pgfsetlinewidth{0.250937pt}%
\definecolor{currentstroke}{rgb}{0.000000,0.000000,0.000000}%
\pgfsetstrokecolor{currentstroke}%
\pgfsetstrokeopacity{0.500000}%
\pgfsetdash{}{0pt}%
\pgfsys@defobject{currentmarker}{\pgfqpoint{-0.005556in}{-0.005556in}}{\pgfqpoint{0.005556in}{0.005556in}}{%
\pgfpathmoveto{\pgfqpoint{0.000000in}{-0.005556in}}%
\pgfpathcurveto{\pgfqpoint{0.001473in}{-0.005556in}}{\pgfqpoint{0.002887in}{-0.004970in}}{\pgfqpoint{0.003928in}{-0.003928in}}%
\pgfpathcurveto{\pgfqpoint{0.004970in}{-0.002887in}}{\pgfqpoint{0.005556in}{-0.001473in}}{\pgfqpoint{0.005556in}{0.000000in}}%
\pgfpathcurveto{\pgfqpoint{0.005556in}{0.001473in}}{\pgfqpoint{0.004970in}{0.002887in}}{\pgfqpoint{0.003928in}{0.003928in}}%
\pgfpathcurveto{\pgfqpoint{0.002887in}{0.004970in}}{\pgfqpoint{0.001473in}{0.005556in}}{\pgfqpoint{0.000000in}{0.005556in}}%
\pgfpathcurveto{\pgfqpoint{-0.001473in}{0.005556in}}{\pgfqpoint{-0.002887in}{0.004970in}}{\pgfqpoint{-0.003928in}{0.003928in}}%
\pgfpathcurveto{\pgfqpoint{-0.004970in}{0.002887in}}{\pgfqpoint{-0.005556in}{0.001473in}}{\pgfqpoint{-0.005556in}{0.000000in}}%
\pgfpathcurveto{\pgfqpoint{-0.005556in}{-0.001473in}}{\pgfqpoint{-0.004970in}{-0.002887in}}{\pgfqpoint{-0.003928in}{-0.003928in}}%
\pgfpathcurveto{\pgfqpoint{-0.002887in}{-0.004970in}}{\pgfqpoint{-0.001473in}{-0.005556in}}{\pgfqpoint{0.000000in}{-0.005556in}}%
\pgfpathlineto{\pgfqpoint{0.000000in}{-0.005556in}}%
\pgfpathclose%
\pgfusepath{stroke,fill}%
}%
\begin{pgfscope}%
\pgfsys@transformshift{3.757252in}{3.513137in}%
\pgfsys@useobject{currentmarker}{}%
\end{pgfscope}%
\end{pgfscope}%
\begin{pgfscope}%
\pgfpathrectangle{\pgfqpoint{0.100000in}{2.493063in}}{\pgfqpoint{5.037500in}{3.427208in}}%
\pgfusepath{clip}%
\pgfsetrectcap%
\pgfsetroundjoin%
\pgfsetlinewidth{1.505625pt}%
\definecolor{currentstroke}{rgb}{0.678431,1.000000,0.184314}%
\pgfsetstrokecolor{currentstroke}%
\pgfsetstrokeopacity{0.500000}%
\pgfsetdash{}{0pt}%
\pgfpathmoveto{\pgfqpoint{3.497670in}{3.904813in}}%
\pgfusepath{stroke}%
\end{pgfscope}%
\begin{pgfscope}%
\pgfpathrectangle{\pgfqpoint{0.100000in}{2.493063in}}{\pgfqpoint{5.037500in}{3.427208in}}%
\pgfusepath{clip}%
\pgfsetbuttcap%
\pgfsetroundjoin%
\definecolor{currentfill}{rgb}{0.678431,1.000000,0.184314}%
\pgfsetfillcolor{currentfill}%
\pgfsetfillopacity{0.500000}%
\pgfsetlinewidth{0.250937pt}%
\definecolor{currentstroke}{rgb}{0.000000,0.000000,0.000000}%
\pgfsetstrokecolor{currentstroke}%
\pgfsetstrokeopacity{0.500000}%
\pgfsetdash{}{0pt}%
\pgfsys@defobject{currentmarker}{\pgfqpoint{-0.011111in}{-0.011111in}}{\pgfqpoint{0.011111in}{0.011111in}}{%
\pgfpathmoveto{\pgfqpoint{0.000000in}{-0.011111in}}%
\pgfpathcurveto{\pgfqpoint{0.002947in}{-0.011111in}}{\pgfqpoint{0.005773in}{-0.009940in}}{\pgfqpoint{0.007857in}{-0.007857in}}%
\pgfpathcurveto{\pgfqpoint{0.009940in}{-0.005773in}}{\pgfqpoint{0.011111in}{-0.002947in}}{\pgfqpoint{0.011111in}{0.000000in}}%
\pgfpathcurveto{\pgfqpoint{0.011111in}{0.002947in}}{\pgfqpoint{0.009940in}{0.005773in}}{\pgfqpoint{0.007857in}{0.007857in}}%
\pgfpathcurveto{\pgfqpoint{0.005773in}{0.009940in}}{\pgfqpoint{0.002947in}{0.011111in}}{\pgfqpoint{0.000000in}{0.011111in}}%
\pgfpathcurveto{\pgfqpoint{-0.002947in}{0.011111in}}{\pgfqpoint{-0.005773in}{0.009940in}}{\pgfqpoint{-0.007857in}{0.007857in}}%
\pgfpathcurveto{\pgfqpoint{-0.009940in}{0.005773in}}{\pgfqpoint{-0.011111in}{0.002947in}}{\pgfqpoint{-0.011111in}{0.000000in}}%
\pgfpathcurveto{\pgfqpoint{-0.011111in}{-0.002947in}}{\pgfqpoint{-0.009940in}{-0.005773in}}{\pgfqpoint{-0.007857in}{-0.007857in}}%
\pgfpathcurveto{\pgfqpoint{-0.005773in}{-0.009940in}}{\pgfqpoint{-0.002947in}{-0.011111in}}{\pgfqpoint{0.000000in}{-0.011111in}}%
\pgfpathlineto{\pgfqpoint{0.000000in}{-0.011111in}}%
\pgfpathclose%
\pgfusepath{stroke,fill}%
}%
\begin{pgfscope}%
\pgfsys@transformshift{3.497670in}{3.904813in}%
\pgfsys@useobject{currentmarker}{}%
\end{pgfscope}%
\end{pgfscope}%
\begin{pgfscope}%
\pgfpathrectangle{\pgfqpoint{0.100000in}{2.493063in}}{\pgfqpoint{5.037500in}{3.427208in}}%
\pgfusepath{clip}%
\pgfsetrectcap%
\pgfsetroundjoin%
\pgfsetlinewidth{1.505625pt}%
\definecolor{currentstroke}{rgb}{0.678431,1.000000,0.184314}%
\pgfsetstrokecolor{currentstroke}%
\pgfsetstrokeopacity{0.500000}%
\pgfsetdash{}{0pt}%
\pgfpathmoveto{\pgfqpoint{3.664388in}{3.828720in}}%
\pgfusepath{stroke}%
\end{pgfscope}%
\begin{pgfscope}%
\pgfpathrectangle{\pgfqpoint{0.100000in}{2.493063in}}{\pgfqpoint{5.037500in}{3.427208in}}%
\pgfusepath{clip}%
\pgfsetbuttcap%
\pgfsetroundjoin%
\definecolor{currentfill}{rgb}{0.678431,1.000000,0.184314}%
\pgfsetfillcolor{currentfill}%
\pgfsetfillopacity{0.500000}%
\pgfsetlinewidth{0.250937pt}%
\definecolor{currentstroke}{rgb}{0.000000,0.000000,0.000000}%
\pgfsetstrokecolor{currentstroke}%
\pgfsetstrokeopacity{0.500000}%
\pgfsetdash{}{0pt}%
\pgfsys@defobject{currentmarker}{\pgfqpoint{-0.005556in}{-0.005556in}}{\pgfqpoint{0.005556in}{0.005556in}}{%
\pgfpathmoveto{\pgfqpoint{0.000000in}{-0.005556in}}%
\pgfpathcurveto{\pgfqpoint{0.001473in}{-0.005556in}}{\pgfqpoint{0.002887in}{-0.004970in}}{\pgfqpoint{0.003928in}{-0.003928in}}%
\pgfpathcurveto{\pgfqpoint{0.004970in}{-0.002887in}}{\pgfqpoint{0.005556in}{-0.001473in}}{\pgfqpoint{0.005556in}{0.000000in}}%
\pgfpathcurveto{\pgfqpoint{0.005556in}{0.001473in}}{\pgfqpoint{0.004970in}{0.002887in}}{\pgfqpoint{0.003928in}{0.003928in}}%
\pgfpathcurveto{\pgfqpoint{0.002887in}{0.004970in}}{\pgfqpoint{0.001473in}{0.005556in}}{\pgfqpoint{0.000000in}{0.005556in}}%
\pgfpathcurveto{\pgfqpoint{-0.001473in}{0.005556in}}{\pgfqpoint{-0.002887in}{0.004970in}}{\pgfqpoint{-0.003928in}{0.003928in}}%
\pgfpathcurveto{\pgfqpoint{-0.004970in}{0.002887in}}{\pgfqpoint{-0.005556in}{0.001473in}}{\pgfqpoint{-0.005556in}{0.000000in}}%
\pgfpathcurveto{\pgfqpoint{-0.005556in}{-0.001473in}}{\pgfqpoint{-0.004970in}{-0.002887in}}{\pgfqpoint{-0.003928in}{-0.003928in}}%
\pgfpathcurveto{\pgfqpoint{-0.002887in}{-0.004970in}}{\pgfqpoint{-0.001473in}{-0.005556in}}{\pgfqpoint{0.000000in}{-0.005556in}}%
\pgfpathlineto{\pgfqpoint{0.000000in}{-0.005556in}}%
\pgfpathclose%
\pgfusepath{stroke,fill}%
}%
\begin{pgfscope}%
\pgfsys@transformshift{3.664388in}{3.828720in}%
\pgfsys@useobject{currentmarker}{}%
\end{pgfscope}%
\end{pgfscope}%
\begin{pgfscope}%
\pgfpathrectangle{\pgfqpoint{0.100000in}{2.493063in}}{\pgfqpoint{5.037500in}{3.427208in}}%
\pgfusepath{clip}%
\pgfsetrectcap%
\pgfsetroundjoin%
\pgfsetlinewidth{1.505625pt}%
\definecolor{currentstroke}{rgb}{0.501961,0.501961,0.501961}%
\pgfsetstrokecolor{currentstroke}%
\pgfsetstrokeopacity{0.500000}%
\pgfsetdash{}{0pt}%
\pgfpathmoveto{\pgfqpoint{3.601503in}{3.905717in}}%
\pgfusepath{stroke}%
\end{pgfscope}%
\begin{pgfscope}%
\pgfpathrectangle{\pgfqpoint{0.100000in}{2.493063in}}{\pgfqpoint{5.037500in}{3.427208in}}%
\pgfusepath{clip}%
\pgfsetbuttcap%
\pgfsetroundjoin%
\definecolor{currentfill}{rgb}{0.501961,0.501961,0.501961}%
\pgfsetfillcolor{currentfill}%
\pgfsetfillopacity{0.500000}%
\pgfsetlinewidth{0.250937pt}%
\definecolor{currentstroke}{rgb}{0.000000,0.000000,0.000000}%
\pgfsetstrokecolor{currentstroke}%
\pgfsetstrokeopacity{0.500000}%
\pgfsetdash{}{0pt}%
\pgfsys@defobject{currentmarker}{\pgfqpoint{-0.013889in}{-0.013889in}}{\pgfqpoint{0.013889in}{0.013889in}}{%
\pgfpathmoveto{\pgfqpoint{0.000000in}{-0.013889in}}%
\pgfpathcurveto{\pgfqpoint{0.003683in}{-0.013889in}}{\pgfqpoint{0.007216in}{-0.012425in}}{\pgfqpoint{0.009821in}{-0.009821in}}%
\pgfpathcurveto{\pgfqpoint{0.012425in}{-0.007216in}}{\pgfqpoint{0.013889in}{-0.003683in}}{\pgfqpoint{0.013889in}{0.000000in}}%
\pgfpathcurveto{\pgfqpoint{0.013889in}{0.003683in}}{\pgfqpoint{0.012425in}{0.007216in}}{\pgfqpoint{0.009821in}{0.009821in}}%
\pgfpathcurveto{\pgfqpoint{0.007216in}{0.012425in}}{\pgfqpoint{0.003683in}{0.013889in}}{\pgfqpoint{0.000000in}{0.013889in}}%
\pgfpathcurveto{\pgfqpoint{-0.003683in}{0.013889in}}{\pgfqpoint{-0.007216in}{0.012425in}}{\pgfqpoint{-0.009821in}{0.009821in}}%
\pgfpathcurveto{\pgfqpoint{-0.012425in}{0.007216in}}{\pgfqpoint{-0.013889in}{0.003683in}}{\pgfqpoint{-0.013889in}{0.000000in}}%
\pgfpathcurveto{\pgfqpoint{-0.013889in}{-0.003683in}}{\pgfqpoint{-0.012425in}{-0.007216in}}{\pgfqpoint{-0.009821in}{-0.009821in}}%
\pgfpathcurveto{\pgfqpoint{-0.007216in}{-0.012425in}}{\pgfqpoint{-0.003683in}{-0.013889in}}{\pgfqpoint{0.000000in}{-0.013889in}}%
\pgfpathlineto{\pgfqpoint{0.000000in}{-0.013889in}}%
\pgfpathclose%
\pgfusepath{stroke,fill}%
}%
\begin{pgfscope}%
\pgfsys@transformshift{3.601503in}{3.905717in}%
\pgfsys@useobject{currentmarker}{}%
\end{pgfscope}%
\end{pgfscope}%
\begin{pgfscope}%
\pgfpathrectangle{\pgfqpoint{0.100000in}{2.493063in}}{\pgfqpoint{5.037500in}{3.427208in}}%
\pgfusepath{clip}%
\pgfsetrectcap%
\pgfsetroundjoin%
\pgfsetlinewidth{1.505625pt}%
\definecolor{currentstroke}{rgb}{0.501961,0.501961,0.501961}%
\pgfsetstrokecolor{currentstroke}%
\pgfsetstrokeopacity{0.500000}%
\pgfsetdash{}{0pt}%
\pgfpathmoveto{\pgfqpoint{3.499471in}{3.427832in}}%
\pgfusepath{stroke}%
\end{pgfscope}%
\begin{pgfscope}%
\pgfpathrectangle{\pgfqpoint{0.100000in}{2.493063in}}{\pgfqpoint{5.037500in}{3.427208in}}%
\pgfusepath{clip}%
\pgfsetbuttcap%
\pgfsetroundjoin%
\definecolor{currentfill}{rgb}{0.501961,0.501961,0.501961}%
\pgfsetfillcolor{currentfill}%
\pgfsetfillopacity{0.500000}%
\pgfsetlinewidth{0.250937pt}%
\definecolor{currentstroke}{rgb}{0.000000,0.000000,0.000000}%
\pgfsetstrokecolor{currentstroke}%
\pgfsetstrokeopacity{0.500000}%
\pgfsetdash{}{0pt}%
\pgfsys@defobject{currentmarker}{\pgfqpoint{-0.013889in}{-0.013889in}}{\pgfqpoint{0.013889in}{0.013889in}}{%
\pgfpathmoveto{\pgfqpoint{0.000000in}{-0.013889in}}%
\pgfpathcurveto{\pgfqpoint{0.003683in}{-0.013889in}}{\pgfqpoint{0.007216in}{-0.012425in}}{\pgfqpoint{0.009821in}{-0.009821in}}%
\pgfpathcurveto{\pgfqpoint{0.012425in}{-0.007216in}}{\pgfqpoint{0.013889in}{-0.003683in}}{\pgfqpoint{0.013889in}{0.000000in}}%
\pgfpathcurveto{\pgfqpoint{0.013889in}{0.003683in}}{\pgfqpoint{0.012425in}{0.007216in}}{\pgfqpoint{0.009821in}{0.009821in}}%
\pgfpathcurveto{\pgfqpoint{0.007216in}{0.012425in}}{\pgfqpoint{0.003683in}{0.013889in}}{\pgfqpoint{0.000000in}{0.013889in}}%
\pgfpathcurveto{\pgfqpoint{-0.003683in}{0.013889in}}{\pgfqpoint{-0.007216in}{0.012425in}}{\pgfqpoint{-0.009821in}{0.009821in}}%
\pgfpathcurveto{\pgfqpoint{-0.012425in}{0.007216in}}{\pgfqpoint{-0.013889in}{0.003683in}}{\pgfqpoint{-0.013889in}{0.000000in}}%
\pgfpathcurveto{\pgfqpoint{-0.013889in}{-0.003683in}}{\pgfqpoint{-0.012425in}{-0.007216in}}{\pgfqpoint{-0.009821in}{-0.009821in}}%
\pgfpathcurveto{\pgfqpoint{-0.007216in}{-0.012425in}}{\pgfqpoint{-0.003683in}{-0.013889in}}{\pgfqpoint{0.000000in}{-0.013889in}}%
\pgfpathlineto{\pgfqpoint{0.000000in}{-0.013889in}}%
\pgfpathclose%
\pgfusepath{stroke,fill}%
}%
\begin{pgfscope}%
\pgfsys@transformshift{3.499471in}{3.427832in}%
\pgfsys@useobject{currentmarker}{}%
\end{pgfscope}%
\end{pgfscope}%
\begin{pgfscope}%
\pgfpathrectangle{\pgfqpoint{0.100000in}{2.493063in}}{\pgfqpoint{5.037500in}{3.427208in}}%
\pgfusepath{clip}%
\pgfsetrectcap%
\pgfsetroundjoin%
\pgfsetlinewidth{1.505625pt}%
\definecolor{currentstroke}{rgb}{0.501961,0.501961,0.501961}%
\pgfsetstrokecolor{currentstroke}%
\pgfsetstrokeopacity{0.500000}%
\pgfsetdash{}{0pt}%
\pgfpathmoveto{\pgfqpoint{3.654580in}{3.636172in}}%
\pgfusepath{stroke}%
\end{pgfscope}%
\begin{pgfscope}%
\pgfpathrectangle{\pgfqpoint{0.100000in}{2.493063in}}{\pgfqpoint{5.037500in}{3.427208in}}%
\pgfusepath{clip}%
\pgfsetbuttcap%
\pgfsetroundjoin%
\definecolor{currentfill}{rgb}{0.501961,0.501961,0.501961}%
\pgfsetfillcolor{currentfill}%
\pgfsetfillopacity{0.500000}%
\pgfsetlinewidth{0.250937pt}%
\definecolor{currentstroke}{rgb}{0.000000,0.000000,0.000000}%
\pgfsetstrokecolor{currentstroke}%
\pgfsetstrokeopacity{0.500000}%
\pgfsetdash{}{0pt}%
\pgfsys@defobject{currentmarker}{\pgfqpoint{-0.013889in}{-0.013889in}}{\pgfqpoint{0.013889in}{0.013889in}}{%
\pgfpathmoveto{\pgfqpoint{0.000000in}{-0.013889in}}%
\pgfpathcurveto{\pgfqpoint{0.003683in}{-0.013889in}}{\pgfqpoint{0.007216in}{-0.012425in}}{\pgfqpoint{0.009821in}{-0.009821in}}%
\pgfpathcurveto{\pgfqpoint{0.012425in}{-0.007216in}}{\pgfqpoint{0.013889in}{-0.003683in}}{\pgfqpoint{0.013889in}{0.000000in}}%
\pgfpathcurveto{\pgfqpoint{0.013889in}{0.003683in}}{\pgfqpoint{0.012425in}{0.007216in}}{\pgfqpoint{0.009821in}{0.009821in}}%
\pgfpathcurveto{\pgfqpoint{0.007216in}{0.012425in}}{\pgfqpoint{0.003683in}{0.013889in}}{\pgfqpoint{0.000000in}{0.013889in}}%
\pgfpathcurveto{\pgfqpoint{-0.003683in}{0.013889in}}{\pgfqpoint{-0.007216in}{0.012425in}}{\pgfqpoint{-0.009821in}{0.009821in}}%
\pgfpathcurveto{\pgfqpoint{-0.012425in}{0.007216in}}{\pgfqpoint{-0.013889in}{0.003683in}}{\pgfqpoint{-0.013889in}{0.000000in}}%
\pgfpathcurveto{\pgfqpoint{-0.013889in}{-0.003683in}}{\pgfqpoint{-0.012425in}{-0.007216in}}{\pgfqpoint{-0.009821in}{-0.009821in}}%
\pgfpathcurveto{\pgfqpoint{-0.007216in}{-0.012425in}}{\pgfqpoint{-0.003683in}{-0.013889in}}{\pgfqpoint{0.000000in}{-0.013889in}}%
\pgfpathlineto{\pgfqpoint{0.000000in}{-0.013889in}}%
\pgfpathclose%
\pgfusepath{stroke,fill}%
}%
\begin{pgfscope}%
\pgfsys@transformshift{3.654580in}{3.636172in}%
\pgfsys@useobject{currentmarker}{}%
\end{pgfscope}%
\end{pgfscope}%
\begin{pgfscope}%
\pgfpathrectangle{\pgfqpoint{0.100000in}{2.493063in}}{\pgfqpoint{5.037500in}{3.427208in}}%
\pgfusepath{clip}%
\pgfsetrectcap%
\pgfsetroundjoin%
\pgfsetlinewidth{1.505625pt}%
\definecolor{currentstroke}{rgb}{0.678431,1.000000,0.184314}%
\pgfsetstrokecolor{currentstroke}%
\pgfsetstrokeopacity{0.500000}%
\pgfsetdash{}{0pt}%
\pgfpathmoveto{\pgfqpoint{3.530356in}{3.729700in}}%
\pgfusepath{stroke}%
\end{pgfscope}%
\begin{pgfscope}%
\pgfpathrectangle{\pgfqpoint{0.100000in}{2.493063in}}{\pgfqpoint{5.037500in}{3.427208in}}%
\pgfusepath{clip}%
\pgfsetbuttcap%
\pgfsetroundjoin%
\definecolor{currentfill}{rgb}{0.678431,1.000000,0.184314}%
\pgfsetfillcolor{currentfill}%
\pgfsetfillopacity{0.500000}%
\pgfsetlinewidth{0.250937pt}%
\definecolor{currentstroke}{rgb}{0.000000,0.000000,0.000000}%
\pgfsetstrokecolor{currentstroke}%
\pgfsetstrokeopacity{0.500000}%
\pgfsetdash{}{0pt}%
\pgfsys@defobject{currentmarker}{\pgfqpoint{-0.005556in}{-0.005556in}}{\pgfqpoint{0.005556in}{0.005556in}}{%
\pgfpathmoveto{\pgfqpoint{0.000000in}{-0.005556in}}%
\pgfpathcurveto{\pgfqpoint{0.001473in}{-0.005556in}}{\pgfqpoint{0.002887in}{-0.004970in}}{\pgfqpoint{0.003928in}{-0.003928in}}%
\pgfpathcurveto{\pgfqpoint{0.004970in}{-0.002887in}}{\pgfqpoint{0.005556in}{-0.001473in}}{\pgfqpoint{0.005556in}{0.000000in}}%
\pgfpathcurveto{\pgfqpoint{0.005556in}{0.001473in}}{\pgfqpoint{0.004970in}{0.002887in}}{\pgfqpoint{0.003928in}{0.003928in}}%
\pgfpathcurveto{\pgfqpoint{0.002887in}{0.004970in}}{\pgfqpoint{0.001473in}{0.005556in}}{\pgfqpoint{0.000000in}{0.005556in}}%
\pgfpathcurveto{\pgfqpoint{-0.001473in}{0.005556in}}{\pgfqpoint{-0.002887in}{0.004970in}}{\pgfqpoint{-0.003928in}{0.003928in}}%
\pgfpathcurveto{\pgfqpoint{-0.004970in}{0.002887in}}{\pgfqpoint{-0.005556in}{0.001473in}}{\pgfqpoint{-0.005556in}{0.000000in}}%
\pgfpathcurveto{\pgfqpoint{-0.005556in}{-0.001473in}}{\pgfqpoint{-0.004970in}{-0.002887in}}{\pgfqpoint{-0.003928in}{-0.003928in}}%
\pgfpathcurveto{\pgfqpoint{-0.002887in}{-0.004970in}}{\pgfqpoint{-0.001473in}{-0.005556in}}{\pgfqpoint{0.000000in}{-0.005556in}}%
\pgfpathlineto{\pgfqpoint{0.000000in}{-0.005556in}}%
\pgfpathclose%
\pgfusepath{stroke,fill}%
}%
\begin{pgfscope}%
\pgfsys@transformshift{3.530356in}{3.729700in}%
\pgfsys@useobject{currentmarker}{}%
\end{pgfscope}%
\end{pgfscope}%
\begin{pgfscope}%
\pgfpathrectangle{\pgfqpoint{0.100000in}{2.493063in}}{\pgfqpoint{5.037500in}{3.427208in}}%
\pgfusepath{clip}%
\pgfsetrectcap%
\pgfsetroundjoin%
\pgfsetlinewidth{1.505625pt}%
\definecolor{currentstroke}{rgb}{0.678431,1.000000,0.184314}%
\pgfsetstrokecolor{currentstroke}%
\pgfsetstrokeopacity{0.500000}%
\pgfsetdash{}{0pt}%
\pgfpathmoveto{\pgfqpoint{1.206521in}{3.000030in}}%
\pgfusepath{stroke}%
\end{pgfscope}%
\begin{pgfscope}%
\pgfpathrectangle{\pgfqpoint{0.100000in}{2.493063in}}{\pgfqpoint{5.037500in}{3.427208in}}%
\pgfusepath{clip}%
\pgfsetbuttcap%
\pgfsetroundjoin%
\definecolor{currentfill}{rgb}{0.678431,1.000000,0.184314}%
\pgfsetfillcolor{currentfill}%
\pgfsetfillopacity{0.500000}%
\pgfsetlinewidth{0.250937pt}%
\definecolor{currentstroke}{rgb}{0.000000,0.000000,0.000000}%
\pgfsetstrokecolor{currentstroke}%
\pgfsetstrokeopacity{0.500000}%
\pgfsetdash{}{0pt}%
\pgfsys@defobject{currentmarker}{\pgfqpoint{-0.022222in}{-0.022222in}}{\pgfqpoint{0.022222in}{0.022222in}}{%
\pgfpathmoveto{\pgfqpoint{0.000000in}{-0.022222in}}%
\pgfpathcurveto{\pgfqpoint{0.005893in}{-0.022222in}}{\pgfqpoint{0.011546in}{-0.019881in}}{\pgfqpoint{0.015713in}{-0.015713in}}%
\pgfpathcurveto{\pgfqpoint{0.019881in}{-0.011546in}}{\pgfqpoint{0.022222in}{-0.005893in}}{\pgfqpoint{0.022222in}{0.000000in}}%
\pgfpathcurveto{\pgfqpoint{0.022222in}{0.005893in}}{\pgfqpoint{0.019881in}{0.011546in}}{\pgfqpoint{0.015713in}{0.015713in}}%
\pgfpathcurveto{\pgfqpoint{0.011546in}{0.019881in}}{\pgfqpoint{0.005893in}{0.022222in}}{\pgfqpoint{0.000000in}{0.022222in}}%
\pgfpathcurveto{\pgfqpoint{-0.005893in}{0.022222in}}{\pgfqpoint{-0.011546in}{0.019881in}}{\pgfqpoint{-0.015713in}{0.015713in}}%
\pgfpathcurveto{\pgfqpoint{-0.019881in}{0.011546in}}{\pgfqpoint{-0.022222in}{0.005893in}}{\pgfqpoint{-0.022222in}{0.000000in}}%
\pgfpathcurveto{\pgfqpoint{-0.022222in}{-0.005893in}}{\pgfqpoint{-0.019881in}{-0.011546in}}{\pgfqpoint{-0.015713in}{-0.015713in}}%
\pgfpathcurveto{\pgfqpoint{-0.011546in}{-0.019881in}}{\pgfqpoint{-0.005893in}{-0.022222in}}{\pgfqpoint{0.000000in}{-0.022222in}}%
\pgfpathlineto{\pgfqpoint{0.000000in}{-0.022222in}}%
\pgfpathclose%
\pgfusepath{stroke,fill}%
}%
\begin{pgfscope}%
\pgfsys@transformshift{1.206521in}{3.000030in}%
\pgfsys@useobject{currentmarker}{}%
\end{pgfscope}%
\end{pgfscope}%
\begin{pgfscope}%
\pgfpathrectangle{\pgfqpoint{0.100000in}{2.493063in}}{\pgfqpoint{5.037500in}{3.427208in}}%
\pgfusepath{clip}%
\pgfsetrectcap%
\pgfsetroundjoin%
\pgfsetlinewidth{1.505625pt}%
\definecolor{currentstroke}{rgb}{0.678431,1.000000,0.184314}%
\pgfsetstrokecolor{currentstroke}%
\pgfsetstrokeopacity{0.500000}%
\pgfsetdash{}{0pt}%
\pgfpathmoveto{\pgfqpoint{1.330071in}{3.106783in}}%
\pgfusepath{stroke}%
\end{pgfscope}%
\begin{pgfscope}%
\pgfpathrectangle{\pgfqpoint{0.100000in}{2.493063in}}{\pgfqpoint{5.037500in}{3.427208in}}%
\pgfusepath{clip}%
\pgfsetbuttcap%
\pgfsetroundjoin%
\definecolor{currentfill}{rgb}{0.678431,1.000000,0.184314}%
\pgfsetfillcolor{currentfill}%
\pgfsetfillopacity{0.500000}%
\pgfsetlinewidth{0.250937pt}%
\definecolor{currentstroke}{rgb}{0.000000,0.000000,0.000000}%
\pgfsetstrokecolor{currentstroke}%
\pgfsetstrokeopacity{0.500000}%
\pgfsetdash{}{0pt}%
\pgfsys@defobject{currentmarker}{\pgfqpoint{-0.013889in}{-0.013889in}}{\pgfqpoint{0.013889in}{0.013889in}}{%
\pgfpathmoveto{\pgfqpoint{0.000000in}{-0.013889in}}%
\pgfpathcurveto{\pgfqpoint{0.003683in}{-0.013889in}}{\pgfqpoint{0.007216in}{-0.012425in}}{\pgfqpoint{0.009821in}{-0.009821in}}%
\pgfpathcurveto{\pgfqpoint{0.012425in}{-0.007216in}}{\pgfqpoint{0.013889in}{-0.003683in}}{\pgfqpoint{0.013889in}{0.000000in}}%
\pgfpathcurveto{\pgfqpoint{0.013889in}{0.003683in}}{\pgfqpoint{0.012425in}{0.007216in}}{\pgfqpoint{0.009821in}{0.009821in}}%
\pgfpathcurveto{\pgfqpoint{0.007216in}{0.012425in}}{\pgfqpoint{0.003683in}{0.013889in}}{\pgfqpoint{0.000000in}{0.013889in}}%
\pgfpathcurveto{\pgfqpoint{-0.003683in}{0.013889in}}{\pgfqpoint{-0.007216in}{0.012425in}}{\pgfqpoint{-0.009821in}{0.009821in}}%
\pgfpathcurveto{\pgfqpoint{-0.012425in}{0.007216in}}{\pgfqpoint{-0.013889in}{0.003683in}}{\pgfqpoint{-0.013889in}{0.000000in}}%
\pgfpathcurveto{\pgfqpoint{-0.013889in}{-0.003683in}}{\pgfqpoint{-0.012425in}{-0.007216in}}{\pgfqpoint{-0.009821in}{-0.009821in}}%
\pgfpathcurveto{\pgfqpoint{-0.007216in}{-0.012425in}}{\pgfqpoint{-0.003683in}{-0.013889in}}{\pgfqpoint{0.000000in}{-0.013889in}}%
\pgfpathlineto{\pgfqpoint{0.000000in}{-0.013889in}}%
\pgfpathclose%
\pgfusepath{stroke,fill}%
}%
\begin{pgfscope}%
\pgfsys@transformshift{1.330071in}{3.106783in}%
\pgfsys@useobject{currentmarker}{}%
\end{pgfscope}%
\end{pgfscope}%
\begin{pgfscope}%
\pgfpathrectangle{\pgfqpoint{0.100000in}{2.493063in}}{\pgfqpoint{5.037500in}{3.427208in}}%
\pgfusepath{clip}%
\pgfsetrectcap%
\pgfsetroundjoin%
\pgfsetlinewidth{1.505625pt}%
\definecolor{currentstroke}{rgb}{0.678431,1.000000,0.184314}%
\pgfsetstrokecolor{currentstroke}%
\pgfsetstrokeopacity{0.500000}%
\pgfsetdash{}{0pt}%
\pgfpathmoveto{\pgfqpoint{1.242579in}{4.066452in}}%
\pgfusepath{stroke}%
\end{pgfscope}%
\begin{pgfscope}%
\pgfpathrectangle{\pgfqpoint{0.100000in}{2.493063in}}{\pgfqpoint{5.037500in}{3.427208in}}%
\pgfusepath{clip}%
\pgfsetbuttcap%
\pgfsetroundjoin%
\definecolor{currentfill}{rgb}{0.678431,1.000000,0.184314}%
\pgfsetfillcolor{currentfill}%
\pgfsetfillopacity{0.500000}%
\pgfsetlinewidth{0.250937pt}%
\definecolor{currentstroke}{rgb}{0.000000,0.000000,0.000000}%
\pgfsetstrokecolor{currentstroke}%
\pgfsetstrokeopacity{0.500000}%
\pgfsetdash{}{0pt}%
\pgfsys@defobject{currentmarker}{\pgfqpoint{-0.022222in}{-0.022222in}}{\pgfqpoint{0.022222in}{0.022222in}}{%
\pgfpathmoveto{\pgfqpoint{0.000000in}{-0.022222in}}%
\pgfpathcurveto{\pgfqpoint{0.005893in}{-0.022222in}}{\pgfqpoint{0.011546in}{-0.019881in}}{\pgfqpoint{0.015713in}{-0.015713in}}%
\pgfpathcurveto{\pgfqpoint{0.019881in}{-0.011546in}}{\pgfqpoint{0.022222in}{-0.005893in}}{\pgfqpoint{0.022222in}{0.000000in}}%
\pgfpathcurveto{\pgfqpoint{0.022222in}{0.005893in}}{\pgfqpoint{0.019881in}{0.011546in}}{\pgfqpoint{0.015713in}{0.015713in}}%
\pgfpathcurveto{\pgfqpoint{0.011546in}{0.019881in}}{\pgfqpoint{0.005893in}{0.022222in}}{\pgfqpoint{0.000000in}{0.022222in}}%
\pgfpathcurveto{\pgfqpoint{-0.005893in}{0.022222in}}{\pgfqpoint{-0.011546in}{0.019881in}}{\pgfqpoint{-0.015713in}{0.015713in}}%
\pgfpathcurveto{\pgfqpoint{-0.019881in}{0.011546in}}{\pgfqpoint{-0.022222in}{0.005893in}}{\pgfqpoint{-0.022222in}{0.000000in}}%
\pgfpathcurveto{\pgfqpoint{-0.022222in}{-0.005893in}}{\pgfqpoint{-0.019881in}{-0.011546in}}{\pgfqpoint{-0.015713in}{-0.015713in}}%
\pgfpathcurveto{\pgfqpoint{-0.011546in}{-0.019881in}}{\pgfqpoint{-0.005893in}{-0.022222in}}{\pgfqpoint{0.000000in}{-0.022222in}}%
\pgfpathlineto{\pgfqpoint{0.000000in}{-0.022222in}}%
\pgfpathclose%
\pgfusepath{stroke,fill}%
}%
\begin{pgfscope}%
\pgfsys@transformshift{1.242579in}{4.066452in}%
\pgfsys@useobject{currentmarker}{}%
\end{pgfscope}%
\end{pgfscope}%
\begin{pgfscope}%
\pgfpathrectangle{\pgfqpoint{0.100000in}{2.493063in}}{\pgfqpoint{5.037500in}{3.427208in}}%
\pgfusepath{clip}%
\pgfsetrectcap%
\pgfsetroundjoin%
\pgfsetlinewidth{1.505625pt}%
\definecolor{currentstroke}{rgb}{0.678431,1.000000,0.184314}%
\pgfsetstrokecolor{currentstroke}%
\pgfsetstrokeopacity{0.500000}%
\pgfsetdash{}{0pt}%
\pgfpathmoveto{\pgfqpoint{0.975939in}{4.034572in}}%
\pgfusepath{stroke}%
\end{pgfscope}%
\begin{pgfscope}%
\pgfpathrectangle{\pgfqpoint{0.100000in}{2.493063in}}{\pgfqpoint{5.037500in}{3.427208in}}%
\pgfusepath{clip}%
\pgfsetbuttcap%
\pgfsetroundjoin%
\definecolor{currentfill}{rgb}{0.678431,1.000000,0.184314}%
\pgfsetfillcolor{currentfill}%
\pgfsetfillopacity{0.500000}%
\pgfsetlinewidth{0.250937pt}%
\definecolor{currentstroke}{rgb}{0.000000,0.000000,0.000000}%
\pgfsetstrokecolor{currentstroke}%
\pgfsetstrokeopacity{0.500000}%
\pgfsetdash{}{0pt}%
\pgfsys@defobject{currentmarker}{\pgfqpoint{-0.027778in}{-0.027778in}}{\pgfqpoint{0.027778in}{0.027778in}}{%
\pgfpathmoveto{\pgfqpoint{0.000000in}{-0.027778in}}%
\pgfpathcurveto{\pgfqpoint{0.007367in}{-0.027778in}}{\pgfqpoint{0.014433in}{-0.024851in}}{\pgfqpoint{0.019642in}{-0.019642in}}%
\pgfpathcurveto{\pgfqpoint{0.024851in}{-0.014433in}}{\pgfqpoint{0.027778in}{-0.007367in}}{\pgfqpoint{0.027778in}{0.000000in}}%
\pgfpathcurveto{\pgfqpoint{0.027778in}{0.007367in}}{\pgfqpoint{0.024851in}{0.014433in}}{\pgfqpoint{0.019642in}{0.019642in}}%
\pgfpathcurveto{\pgfqpoint{0.014433in}{0.024851in}}{\pgfqpoint{0.007367in}{0.027778in}}{\pgfqpoint{0.000000in}{0.027778in}}%
\pgfpathcurveto{\pgfqpoint{-0.007367in}{0.027778in}}{\pgfqpoint{-0.014433in}{0.024851in}}{\pgfqpoint{-0.019642in}{0.019642in}}%
\pgfpathcurveto{\pgfqpoint{-0.024851in}{0.014433in}}{\pgfqpoint{-0.027778in}{0.007367in}}{\pgfqpoint{-0.027778in}{0.000000in}}%
\pgfpathcurveto{\pgfqpoint{-0.027778in}{-0.007367in}}{\pgfqpoint{-0.024851in}{-0.014433in}}{\pgfqpoint{-0.019642in}{-0.019642in}}%
\pgfpathcurveto{\pgfqpoint{-0.014433in}{-0.024851in}}{\pgfqpoint{-0.007367in}{-0.027778in}}{\pgfqpoint{0.000000in}{-0.027778in}}%
\pgfpathlineto{\pgfqpoint{0.000000in}{-0.027778in}}%
\pgfpathclose%
\pgfusepath{stroke,fill}%
}%
\begin{pgfscope}%
\pgfsys@transformshift{0.975939in}{4.034572in}%
\pgfsys@useobject{currentmarker}{}%
\end{pgfscope}%
\end{pgfscope}%
\begin{pgfscope}%
\pgfpathrectangle{\pgfqpoint{0.100000in}{2.493063in}}{\pgfqpoint{5.037500in}{3.427208in}}%
\pgfusepath{clip}%
\pgfsetrectcap%
\pgfsetroundjoin%
\pgfsetlinewidth{1.505625pt}%
\definecolor{currentstroke}{rgb}{0.678431,1.000000,0.184314}%
\pgfsetstrokecolor{currentstroke}%
\pgfsetstrokeopacity{0.500000}%
\pgfsetdash{}{0pt}%
\pgfpathmoveto{\pgfqpoint{1.165310in}{3.875127in}}%
\pgfusepath{stroke}%
\end{pgfscope}%
\begin{pgfscope}%
\pgfpathrectangle{\pgfqpoint{0.100000in}{2.493063in}}{\pgfqpoint{5.037500in}{3.427208in}}%
\pgfusepath{clip}%
\pgfsetbuttcap%
\pgfsetroundjoin%
\definecolor{currentfill}{rgb}{0.678431,1.000000,0.184314}%
\pgfsetfillcolor{currentfill}%
\pgfsetfillopacity{0.500000}%
\pgfsetlinewidth{0.250937pt}%
\definecolor{currentstroke}{rgb}{0.000000,0.000000,0.000000}%
\pgfsetstrokecolor{currentstroke}%
\pgfsetstrokeopacity{0.500000}%
\pgfsetdash{}{0pt}%
\pgfsys@defobject{currentmarker}{\pgfqpoint{-0.008333in}{-0.008333in}}{\pgfqpoint{0.008333in}{0.008333in}}{%
\pgfpathmoveto{\pgfqpoint{0.000000in}{-0.008333in}}%
\pgfpathcurveto{\pgfqpoint{0.002210in}{-0.008333in}}{\pgfqpoint{0.004330in}{-0.007455in}}{\pgfqpoint{0.005893in}{-0.005893in}}%
\pgfpathcurveto{\pgfqpoint{0.007455in}{-0.004330in}}{\pgfqpoint{0.008333in}{-0.002210in}}{\pgfqpoint{0.008333in}{0.000000in}}%
\pgfpathcurveto{\pgfqpoint{0.008333in}{0.002210in}}{\pgfqpoint{0.007455in}{0.004330in}}{\pgfqpoint{0.005893in}{0.005893in}}%
\pgfpathcurveto{\pgfqpoint{0.004330in}{0.007455in}}{\pgfqpoint{0.002210in}{0.008333in}}{\pgfqpoint{0.000000in}{0.008333in}}%
\pgfpathcurveto{\pgfqpoint{-0.002210in}{0.008333in}}{\pgfqpoint{-0.004330in}{0.007455in}}{\pgfqpoint{-0.005893in}{0.005893in}}%
\pgfpathcurveto{\pgfqpoint{-0.007455in}{0.004330in}}{\pgfqpoint{-0.008333in}{0.002210in}}{\pgfqpoint{-0.008333in}{0.000000in}}%
\pgfpathcurveto{\pgfqpoint{-0.008333in}{-0.002210in}}{\pgfqpoint{-0.007455in}{-0.004330in}}{\pgfqpoint{-0.005893in}{-0.005893in}}%
\pgfpathcurveto{\pgfqpoint{-0.004330in}{-0.007455in}}{\pgfqpoint{-0.002210in}{-0.008333in}}{\pgfqpoint{0.000000in}{-0.008333in}}%
\pgfpathlineto{\pgfqpoint{0.000000in}{-0.008333in}}%
\pgfpathclose%
\pgfusepath{stroke,fill}%
}%
\begin{pgfscope}%
\pgfsys@transformshift{1.165310in}{3.875127in}%
\pgfsys@useobject{currentmarker}{}%
\end{pgfscope}%
\end{pgfscope}%
\begin{pgfscope}%
\pgfpathrectangle{\pgfqpoint{0.100000in}{2.493063in}}{\pgfqpoint{5.037500in}{3.427208in}}%
\pgfusepath{clip}%
\pgfsetrectcap%
\pgfsetroundjoin%
\pgfsetlinewidth{1.505625pt}%
\definecolor{currentstroke}{rgb}{0.678431,1.000000,0.184314}%
\pgfsetstrokecolor{currentstroke}%
\pgfsetstrokeopacity{0.500000}%
\pgfsetdash{}{0pt}%
\pgfpathmoveto{\pgfqpoint{1.152245in}{4.006148in}}%
\pgfusepath{stroke}%
\end{pgfscope}%
\begin{pgfscope}%
\pgfpathrectangle{\pgfqpoint{0.100000in}{2.493063in}}{\pgfqpoint{5.037500in}{3.427208in}}%
\pgfusepath{clip}%
\pgfsetbuttcap%
\pgfsetroundjoin%
\definecolor{currentfill}{rgb}{0.678431,1.000000,0.184314}%
\pgfsetfillcolor{currentfill}%
\pgfsetfillopacity{0.500000}%
\pgfsetlinewidth{0.250937pt}%
\definecolor{currentstroke}{rgb}{0.000000,0.000000,0.000000}%
\pgfsetstrokecolor{currentstroke}%
\pgfsetstrokeopacity{0.500000}%
\pgfsetdash{}{0pt}%
\pgfsys@defobject{currentmarker}{\pgfqpoint{-0.025000in}{-0.025000in}}{\pgfqpoint{0.025000in}{0.025000in}}{%
\pgfpathmoveto{\pgfqpoint{0.000000in}{-0.025000in}}%
\pgfpathcurveto{\pgfqpoint{0.006630in}{-0.025000in}}{\pgfqpoint{0.012989in}{-0.022366in}}{\pgfqpoint{0.017678in}{-0.017678in}}%
\pgfpathcurveto{\pgfqpoint{0.022366in}{-0.012989in}}{\pgfqpoint{0.025000in}{-0.006630in}}{\pgfqpoint{0.025000in}{0.000000in}}%
\pgfpathcurveto{\pgfqpoint{0.025000in}{0.006630in}}{\pgfqpoint{0.022366in}{0.012989in}}{\pgfqpoint{0.017678in}{0.017678in}}%
\pgfpathcurveto{\pgfqpoint{0.012989in}{0.022366in}}{\pgfqpoint{0.006630in}{0.025000in}}{\pgfqpoint{0.000000in}{0.025000in}}%
\pgfpathcurveto{\pgfqpoint{-0.006630in}{0.025000in}}{\pgfqpoint{-0.012989in}{0.022366in}}{\pgfqpoint{-0.017678in}{0.017678in}}%
\pgfpathcurveto{\pgfqpoint{-0.022366in}{0.012989in}}{\pgfqpoint{-0.025000in}{0.006630in}}{\pgfqpoint{-0.025000in}{0.000000in}}%
\pgfpathcurveto{\pgfqpoint{-0.025000in}{-0.006630in}}{\pgfqpoint{-0.022366in}{-0.012989in}}{\pgfqpoint{-0.017678in}{-0.017678in}}%
\pgfpathcurveto{\pgfqpoint{-0.012989in}{-0.022366in}}{\pgfqpoint{-0.006630in}{-0.025000in}}{\pgfqpoint{0.000000in}{-0.025000in}}%
\pgfpathlineto{\pgfqpoint{0.000000in}{-0.025000in}}%
\pgfpathclose%
\pgfusepath{stroke,fill}%
}%
\begin{pgfscope}%
\pgfsys@transformshift{1.152245in}{4.006148in}%
\pgfsys@useobject{currentmarker}{}%
\end{pgfscope}%
\end{pgfscope}%
\begin{pgfscope}%
\pgfpathrectangle{\pgfqpoint{0.100000in}{2.493063in}}{\pgfqpoint{5.037500in}{3.427208in}}%
\pgfusepath{clip}%
\pgfsetrectcap%
\pgfsetroundjoin%
\pgfsetlinewidth{1.505625pt}%
\definecolor{currentstroke}{rgb}{0.678431,1.000000,0.184314}%
\pgfsetstrokecolor{currentstroke}%
\pgfsetstrokeopacity{0.500000}%
\pgfsetdash{}{0pt}%
\pgfpathmoveto{\pgfqpoint{1.297020in}{3.628381in}}%
\pgfusepath{stroke}%
\end{pgfscope}%
\begin{pgfscope}%
\pgfpathrectangle{\pgfqpoint{0.100000in}{2.493063in}}{\pgfqpoint{5.037500in}{3.427208in}}%
\pgfusepath{clip}%
\pgfsetbuttcap%
\pgfsetroundjoin%
\definecolor{currentfill}{rgb}{0.678431,1.000000,0.184314}%
\pgfsetfillcolor{currentfill}%
\pgfsetfillopacity{0.500000}%
\pgfsetlinewidth{0.250937pt}%
\definecolor{currentstroke}{rgb}{0.000000,0.000000,0.000000}%
\pgfsetstrokecolor{currentstroke}%
\pgfsetstrokeopacity{0.500000}%
\pgfsetdash{}{0pt}%
\pgfsys@defobject{currentmarker}{\pgfqpoint{-0.019444in}{-0.019444in}}{\pgfqpoint{0.019444in}{0.019444in}}{%
\pgfpathmoveto{\pgfqpoint{0.000000in}{-0.019444in}}%
\pgfpathcurveto{\pgfqpoint{0.005157in}{-0.019444in}}{\pgfqpoint{0.010103in}{-0.017396in}}{\pgfqpoint{0.013749in}{-0.013749in}}%
\pgfpathcurveto{\pgfqpoint{0.017396in}{-0.010103in}}{\pgfqpoint{0.019444in}{-0.005157in}}{\pgfqpoint{0.019444in}{0.000000in}}%
\pgfpathcurveto{\pgfqpoint{0.019444in}{0.005157in}}{\pgfqpoint{0.017396in}{0.010103in}}{\pgfqpoint{0.013749in}{0.013749in}}%
\pgfpathcurveto{\pgfqpoint{0.010103in}{0.017396in}}{\pgfqpoint{0.005157in}{0.019444in}}{\pgfqpoint{0.000000in}{0.019444in}}%
\pgfpathcurveto{\pgfqpoint{-0.005157in}{0.019444in}}{\pgfqpoint{-0.010103in}{0.017396in}}{\pgfqpoint{-0.013749in}{0.013749in}}%
\pgfpathcurveto{\pgfqpoint{-0.017396in}{0.010103in}}{\pgfqpoint{-0.019444in}{0.005157in}}{\pgfqpoint{-0.019444in}{0.000000in}}%
\pgfpathcurveto{\pgfqpoint{-0.019444in}{-0.005157in}}{\pgfqpoint{-0.017396in}{-0.010103in}}{\pgfqpoint{-0.013749in}{-0.013749in}}%
\pgfpathcurveto{\pgfqpoint{-0.010103in}{-0.017396in}}{\pgfqpoint{-0.005157in}{-0.019444in}}{\pgfqpoint{0.000000in}{-0.019444in}}%
\pgfpathlineto{\pgfqpoint{0.000000in}{-0.019444in}}%
\pgfpathclose%
\pgfusepath{stroke,fill}%
}%
\begin{pgfscope}%
\pgfsys@transformshift{1.297020in}{3.628381in}%
\pgfsys@useobject{currentmarker}{}%
\end{pgfscope}%
\end{pgfscope}%
\begin{pgfscope}%
\pgfpathrectangle{\pgfqpoint{0.100000in}{2.493063in}}{\pgfqpoint{5.037500in}{3.427208in}}%
\pgfusepath{clip}%
\pgfsetrectcap%
\pgfsetroundjoin%
\pgfsetlinewidth{1.505625pt}%
\definecolor{currentstroke}{rgb}{0.678431,1.000000,0.184314}%
\pgfsetstrokecolor{currentstroke}%
\pgfsetstrokeopacity{0.500000}%
\pgfsetdash{}{0pt}%
\pgfpathmoveto{\pgfqpoint{1.249803in}{3.715149in}}%
\pgfusepath{stroke}%
\end{pgfscope}%
\begin{pgfscope}%
\pgfpathrectangle{\pgfqpoint{0.100000in}{2.493063in}}{\pgfqpoint{5.037500in}{3.427208in}}%
\pgfusepath{clip}%
\pgfsetbuttcap%
\pgfsetroundjoin%
\definecolor{currentfill}{rgb}{0.678431,1.000000,0.184314}%
\pgfsetfillcolor{currentfill}%
\pgfsetfillopacity{0.500000}%
\pgfsetlinewidth{0.250937pt}%
\definecolor{currentstroke}{rgb}{0.000000,0.000000,0.000000}%
\pgfsetstrokecolor{currentstroke}%
\pgfsetstrokeopacity{0.500000}%
\pgfsetdash{}{0pt}%
\pgfsys@defobject{currentmarker}{\pgfqpoint{-0.011111in}{-0.011111in}}{\pgfqpoint{0.011111in}{0.011111in}}{%
\pgfpathmoveto{\pgfqpoint{0.000000in}{-0.011111in}}%
\pgfpathcurveto{\pgfqpoint{0.002947in}{-0.011111in}}{\pgfqpoint{0.005773in}{-0.009940in}}{\pgfqpoint{0.007857in}{-0.007857in}}%
\pgfpathcurveto{\pgfqpoint{0.009940in}{-0.005773in}}{\pgfqpoint{0.011111in}{-0.002947in}}{\pgfqpoint{0.011111in}{0.000000in}}%
\pgfpathcurveto{\pgfqpoint{0.011111in}{0.002947in}}{\pgfqpoint{0.009940in}{0.005773in}}{\pgfqpoint{0.007857in}{0.007857in}}%
\pgfpathcurveto{\pgfqpoint{0.005773in}{0.009940in}}{\pgfqpoint{0.002947in}{0.011111in}}{\pgfqpoint{0.000000in}{0.011111in}}%
\pgfpathcurveto{\pgfqpoint{-0.002947in}{0.011111in}}{\pgfqpoint{-0.005773in}{0.009940in}}{\pgfqpoint{-0.007857in}{0.007857in}}%
\pgfpathcurveto{\pgfqpoint{-0.009940in}{0.005773in}}{\pgfqpoint{-0.011111in}{0.002947in}}{\pgfqpoint{-0.011111in}{0.000000in}}%
\pgfpathcurveto{\pgfqpoint{-0.011111in}{-0.002947in}}{\pgfqpoint{-0.009940in}{-0.005773in}}{\pgfqpoint{-0.007857in}{-0.007857in}}%
\pgfpathcurveto{\pgfqpoint{-0.005773in}{-0.009940in}}{\pgfqpoint{-0.002947in}{-0.011111in}}{\pgfqpoint{0.000000in}{-0.011111in}}%
\pgfpathlineto{\pgfqpoint{0.000000in}{-0.011111in}}%
\pgfpathclose%
\pgfusepath{stroke,fill}%
}%
\begin{pgfscope}%
\pgfsys@transformshift{1.249803in}{3.715149in}%
\pgfsys@useobject{currentmarker}{}%
\end{pgfscope}%
\end{pgfscope}%
\begin{pgfscope}%
\pgfpathrectangle{\pgfqpoint{0.100000in}{2.493063in}}{\pgfqpoint{5.037500in}{3.427208in}}%
\pgfusepath{clip}%
\pgfsetrectcap%
\pgfsetroundjoin%
\pgfsetlinewidth{1.505625pt}%
\definecolor{currentstroke}{rgb}{0.678431,1.000000,0.184314}%
\pgfsetstrokecolor{currentstroke}%
\pgfsetstrokeopacity{0.500000}%
\pgfsetdash{}{0pt}%
\pgfpathmoveto{\pgfqpoint{0.919145in}{3.832754in}}%
\pgfusepath{stroke}%
\end{pgfscope}%
\begin{pgfscope}%
\pgfpathrectangle{\pgfqpoint{0.100000in}{2.493063in}}{\pgfqpoint{5.037500in}{3.427208in}}%
\pgfusepath{clip}%
\pgfsetbuttcap%
\pgfsetroundjoin%
\definecolor{currentfill}{rgb}{0.678431,1.000000,0.184314}%
\pgfsetfillcolor{currentfill}%
\pgfsetfillopacity{0.500000}%
\pgfsetlinewidth{0.250937pt}%
\definecolor{currentstroke}{rgb}{0.000000,0.000000,0.000000}%
\pgfsetstrokecolor{currentstroke}%
\pgfsetstrokeopacity{0.500000}%
\pgfsetdash{}{0pt}%
\pgfsys@defobject{currentmarker}{\pgfqpoint{-0.038889in}{-0.038889in}}{\pgfqpoint{0.038889in}{0.038889in}}{%
\pgfpathmoveto{\pgfqpoint{0.000000in}{-0.038889in}}%
\pgfpathcurveto{\pgfqpoint{0.010313in}{-0.038889in}}{\pgfqpoint{0.020206in}{-0.034791in}}{\pgfqpoint{0.027499in}{-0.027499in}}%
\pgfpathcurveto{\pgfqpoint{0.034791in}{-0.020206in}}{\pgfqpoint{0.038889in}{-0.010313in}}{\pgfqpoint{0.038889in}{0.000000in}}%
\pgfpathcurveto{\pgfqpoint{0.038889in}{0.010313in}}{\pgfqpoint{0.034791in}{0.020206in}}{\pgfqpoint{0.027499in}{0.027499in}}%
\pgfpathcurveto{\pgfqpoint{0.020206in}{0.034791in}}{\pgfqpoint{0.010313in}{0.038889in}}{\pgfqpoint{0.000000in}{0.038889in}}%
\pgfpathcurveto{\pgfqpoint{-0.010313in}{0.038889in}}{\pgfqpoint{-0.020206in}{0.034791in}}{\pgfqpoint{-0.027499in}{0.027499in}}%
\pgfpathcurveto{\pgfqpoint{-0.034791in}{0.020206in}}{\pgfqpoint{-0.038889in}{0.010313in}}{\pgfqpoint{-0.038889in}{0.000000in}}%
\pgfpathcurveto{\pgfqpoint{-0.038889in}{-0.010313in}}{\pgfqpoint{-0.034791in}{-0.020206in}}{\pgfqpoint{-0.027499in}{-0.027499in}}%
\pgfpathcurveto{\pgfqpoint{-0.020206in}{-0.034791in}}{\pgfqpoint{-0.010313in}{-0.038889in}}{\pgfqpoint{0.000000in}{-0.038889in}}%
\pgfpathlineto{\pgfqpoint{0.000000in}{-0.038889in}}%
\pgfpathclose%
\pgfusepath{stroke,fill}%
}%
\begin{pgfscope}%
\pgfsys@transformshift{0.919145in}{3.832754in}%
\pgfsys@useobject{currentmarker}{}%
\end{pgfscope}%
\end{pgfscope}%
\begin{pgfscope}%
\pgfpathrectangle{\pgfqpoint{0.100000in}{2.493063in}}{\pgfqpoint{5.037500in}{3.427208in}}%
\pgfusepath{clip}%
\pgfsetrectcap%
\pgfsetroundjoin%
\pgfsetlinewidth{1.505625pt}%
\definecolor{currentstroke}{rgb}{0.501961,0.501961,0.501961}%
\pgfsetstrokecolor{currentstroke}%
\pgfsetstrokeopacity{0.500000}%
\pgfsetdash{}{0pt}%
\pgfpathmoveto{\pgfqpoint{2.882663in}{4.022918in}}%
\pgfusepath{stroke}%
\end{pgfscope}%
\begin{pgfscope}%
\pgfpathrectangle{\pgfqpoint{0.100000in}{2.493063in}}{\pgfqpoint{5.037500in}{3.427208in}}%
\pgfusepath{clip}%
\pgfsetbuttcap%
\pgfsetroundjoin%
\definecolor{currentfill}{rgb}{0.501961,0.501961,0.501961}%
\pgfsetfillcolor{currentfill}%
\pgfsetfillopacity{0.500000}%
\pgfsetlinewidth{0.250937pt}%
\definecolor{currentstroke}{rgb}{0.000000,0.000000,0.000000}%
\pgfsetstrokecolor{currentstroke}%
\pgfsetstrokeopacity{0.500000}%
\pgfsetdash{}{0pt}%
\pgfsys@defobject{currentmarker}{\pgfqpoint{-0.013889in}{-0.013889in}}{\pgfqpoint{0.013889in}{0.013889in}}{%
\pgfpathmoveto{\pgfqpoint{0.000000in}{-0.013889in}}%
\pgfpathcurveto{\pgfqpoint{0.003683in}{-0.013889in}}{\pgfqpoint{0.007216in}{-0.012425in}}{\pgfqpoint{0.009821in}{-0.009821in}}%
\pgfpathcurveto{\pgfqpoint{0.012425in}{-0.007216in}}{\pgfqpoint{0.013889in}{-0.003683in}}{\pgfqpoint{0.013889in}{0.000000in}}%
\pgfpathcurveto{\pgfqpoint{0.013889in}{0.003683in}}{\pgfqpoint{0.012425in}{0.007216in}}{\pgfqpoint{0.009821in}{0.009821in}}%
\pgfpathcurveto{\pgfqpoint{0.007216in}{0.012425in}}{\pgfqpoint{0.003683in}{0.013889in}}{\pgfqpoint{0.000000in}{0.013889in}}%
\pgfpathcurveto{\pgfqpoint{-0.003683in}{0.013889in}}{\pgfqpoint{-0.007216in}{0.012425in}}{\pgfqpoint{-0.009821in}{0.009821in}}%
\pgfpathcurveto{\pgfqpoint{-0.012425in}{0.007216in}}{\pgfqpoint{-0.013889in}{0.003683in}}{\pgfqpoint{-0.013889in}{0.000000in}}%
\pgfpathcurveto{\pgfqpoint{-0.013889in}{-0.003683in}}{\pgfqpoint{-0.012425in}{-0.007216in}}{\pgfqpoint{-0.009821in}{-0.009821in}}%
\pgfpathcurveto{\pgfqpoint{-0.007216in}{-0.012425in}}{\pgfqpoint{-0.003683in}{-0.013889in}}{\pgfqpoint{0.000000in}{-0.013889in}}%
\pgfpathlineto{\pgfqpoint{0.000000in}{-0.013889in}}%
\pgfpathclose%
\pgfusepath{stroke,fill}%
}%
\begin{pgfscope}%
\pgfsys@transformshift{2.882663in}{4.022918in}%
\pgfsys@useobject{currentmarker}{}%
\end{pgfscope}%
\end{pgfscope}%
\begin{pgfscope}%
\pgfpathrectangle{\pgfqpoint{0.100000in}{2.493063in}}{\pgfqpoint{5.037500in}{3.427208in}}%
\pgfusepath{clip}%
\pgfsetrectcap%
\pgfsetroundjoin%
\pgfsetlinewidth{1.505625pt}%
\definecolor{currentstroke}{rgb}{0.000000,0.000000,1.000000}%
\pgfsetstrokecolor{currentstroke}%
\pgfsetstrokeopacity{0.500000}%
\pgfsetdash{}{0pt}%
\pgfpathmoveto{\pgfqpoint{2.858226in}{3.944827in}}%
\pgfusepath{stroke}%
\end{pgfscope}%
\begin{pgfscope}%
\pgfpathrectangle{\pgfqpoint{0.100000in}{2.493063in}}{\pgfqpoint{5.037500in}{3.427208in}}%
\pgfusepath{clip}%
\pgfsetbuttcap%
\pgfsetroundjoin%
\definecolor{currentfill}{rgb}{0.000000,0.000000,1.000000}%
\pgfsetfillcolor{currentfill}%
\pgfsetfillopacity{0.500000}%
\pgfsetlinewidth{0.250937pt}%
\definecolor{currentstroke}{rgb}{0.000000,0.000000,0.000000}%
\pgfsetstrokecolor{currentstroke}%
\pgfsetstrokeopacity{0.500000}%
\pgfsetdash{}{0pt}%
\pgfsys@defobject{currentmarker}{\pgfqpoint{-0.011111in}{-0.011111in}}{\pgfqpoint{0.011111in}{0.011111in}}{%
\pgfpathmoveto{\pgfqpoint{0.000000in}{-0.011111in}}%
\pgfpathcurveto{\pgfqpoint{0.002947in}{-0.011111in}}{\pgfqpoint{0.005773in}{-0.009940in}}{\pgfqpoint{0.007857in}{-0.007857in}}%
\pgfpathcurveto{\pgfqpoint{0.009940in}{-0.005773in}}{\pgfqpoint{0.011111in}{-0.002947in}}{\pgfqpoint{0.011111in}{0.000000in}}%
\pgfpathcurveto{\pgfqpoint{0.011111in}{0.002947in}}{\pgfqpoint{0.009940in}{0.005773in}}{\pgfqpoint{0.007857in}{0.007857in}}%
\pgfpathcurveto{\pgfqpoint{0.005773in}{0.009940in}}{\pgfqpoint{0.002947in}{0.011111in}}{\pgfqpoint{0.000000in}{0.011111in}}%
\pgfpathcurveto{\pgfqpoint{-0.002947in}{0.011111in}}{\pgfqpoint{-0.005773in}{0.009940in}}{\pgfqpoint{-0.007857in}{0.007857in}}%
\pgfpathcurveto{\pgfqpoint{-0.009940in}{0.005773in}}{\pgfqpoint{-0.011111in}{0.002947in}}{\pgfqpoint{-0.011111in}{0.000000in}}%
\pgfpathcurveto{\pgfqpoint{-0.011111in}{-0.002947in}}{\pgfqpoint{-0.009940in}{-0.005773in}}{\pgfqpoint{-0.007857in}{-0.007857in}}%
\pgfpathcurveto{\pgfqpoint{-0.005773in}{-0.009940in}}{\pgfqpoint{-0.002947in}{-0.011111in}}{\pgfqpoint{0.000000in}{-0.011111in}}%
\pgfpathlineto{\pgfqpoint{0.000000in}{-0.011111in}}%
\pgfpathclose%
\pgfusepath{stroke,fill}%
}%
\begin{pgfscope}%
\pgfsys@transformshift{2.858226in}{3.944827in}%
\pgfsys@useobject{currentmarker}{}%
\end{pgfscope}%
\end{pgfscope}%
\begin{pgfscope}%
\pgfpathrectangle{\pgfqpoint{0.100000in}{2.493063in}}{\pgfqpoint{5.037500in}{3.427208in}}%
\pgfusepath{clip}%
\pgfsetrectcap%
\pgfsetroundjoin%
\pgfsetlinewidth{1.505625pt}%
\definecolor{currentstroke}{rgb}{0.501961,0.501961,0.501961}%
\pgfsetstrokecolor{currentstroke}%
\pgfsetstrokeopacity{0.500000}%
\pgfsetdash{}{0pt}%
\pgfpathmoveto{\pgfqpoint{2.990171in}{3.838232in}}%
\pgfusepath{stroke}%
\end{pgfscope}%
\begin{pgfscope}%
\pgfpathrectangle{\pgfqpoint{0.100000in}{2.493063in}}{\pgfqpoint{5.037500in}{3.427208in}}%
\pgfusepath{clip}%
\pgfsetbuttcap%
\pgfsetroundjoin%
\definecolor{currentfill}{rgb}{0.501961,0.501961,0.501961}%
\pgfsetfillcolor{currentfill}%
\pgfsetfillopacity{0.500000}%
\pgfsetlinewidth{0.250937pt}%
\definecolor{currentstroke}{rgb}{0.000000,0.000000,0.000000}%
\pgfsetstrokecolor{currentstroke}%
\pgfsetstrokeopacity{0.500000}%
\pgfsetdash{}{0pt}%
\pgfsys@defobject{currentmarker}{\pgfqpoint{-0.013889in}{-0.013889in}}{\pgfqpoint{0.013889in}{0.013889in}}{%
\pgfpathmoveto{\pgfqpoint{0.000000in}{-0.013889in}}%
\pgfpathcurveto{\pgfqpoint{0.003683in}{-0.013889in}}{\pgfqpoint{0.007216in}{-0.012425in}}{\pgfqpoint{0.009821in}{-0.009821in}}%
\pgfpathcurveto{\pgfqpoint{0.012425in}{-0.007216in}}{\pgfqpoint{0.013889in}{-0.003683in}}{\pgfqpoint{0.013889in}{0.000000in}}%
\pgfpathcurveto{\pgfqpoint{0.013889in}{0.003683in}}{\pgfqpoint{0.012425in}{0.007216in}}{\pgfqpoint{0.009821in}{0.009821in}}%
\pgfpathcurveto{\pgfqpoint{0.007216in}{0.012425in}}{\pgfqpoint{0.003683in}{0.013889in}}{\pgfqpoint{0.000000in}{0.013889in}}%
\pgfpathcurveto{\pgfqpoint{-0.003683in}{0.013889in}}{\pgfqpoint{-0.007216in}{0.012425in}}{\pgfqpoint{-0.009821in}{0.009821in}}%
\pgfpathcurveto{\pgfqpoint{-0.012425in}{0.007216in}}{\pgfqpoint{-0.013889in}{0.003683in}}{\pgfqpoint{-0.013889in}{0.000000in}}%
\pgfpathcurveto{\pgfqpoint{-0.013889in}{-0.003683in}}{\pgfqpoint{-0.012425in}{-0.007216in}}{\pgfqpoint{-0.009821in}{-0.009821in}}%
\pgfpathcurveto{\pgfqpoint{-0.007216in}{-0.012425in}}{\pgfqpoint{-0.003683in}{-0.013889in}}{\pgfqpoint{0.000000in}{-0.013889in}}%
\pgfpathlineto{\pgfqpoint{0.000000in}{-0.013889in}}%
\pgfpathclose%
\pgfusepath{stroke,fill}%
}%
\begin{pgfscope}%
\pgfsys@transformshift{2.990171in}{3.838232in}%
\pgfsys@useobject{currentmarker}{}%
\end{pgfscope}%
\end{pgfscope}%
\begin{pgfscope}%
\pgfpathrectangle{\pgfqpoint{0.100000in}{2.493063in}}{\pgfqpoint{5.037500in}{3.427208in}}%
\pgfusepath{clip}%
\pgfsetrectcap%
\pgfsetroundjoin%
\pgfsetlinewidth{1.505625pt}%
\definecolor{currentstroke}{rgb}{0.000000,0.000000,1.000000}%
\pgfsetstrokecolor{currentstroke}%
\pgfsetstrokeopacity{0.500000}%
\pgfsetdash{}{0pt}%
\pgfpathmoveto{\pgfqpoint{3.205657in}{4.006711in}}%
\pgfusepath{stroke}%
\end{pgfscope}%
\begin{pgfscope}%
\pgfpathrectangle{\pgfqpoint{0.100000in}{2.493063in}}{\pgfqpoint{5.037500in}{3.427208in}}%
\pgfusepath{clip}%
\pgfsetbuttcap%
\pgfsetroundjoin%
\definecolor{currentfill}{rgb}{0.000000,0.000000,1.000000}%
\pgfsetfillcolor{currentfill}%
\pgfsetfillopacity{0.500000}%
\pgfsetlinewidth{0.250937pt}%
\definecolor{currentstroke}{rgb}{0.000000,0.000000,0.000000}%
\pgfsetstrokecolor{currentstroke}%
\pgfsetstrokeopacity{0.500000}%
\pgfsetdash{}{0pt}%
\pgfsys@defobject{currentmarker}{\pgfqpoint{-0.011111in}{-0.011111in}}{\pgfqpoint{0.011111in}{0.011111in}}{%
\pgfpathmoveto{\pgfqpoint{0.000000in}{-0.011111in}}%
\pgfpathcurveto{\pgfqpoint{0.002947in}{-0.011111in}}{\pgfqpoint{0.005773in}{-0.009940in}}{\pgfqpoint{0.007857in}{-0.007857in}}%
\pgfpathcurveto{\pgfqpoint{0.009940in}{-0.005773in}}{\pgfqpoint{0.011111in}{-0.002947in}}{\pgfqpoint{0.011111in}{0.000000in}}%
\pgfpathcurveto{\pgfqpoint{0.011111in}{0.002947in}}{\pgfqpoint{0.009940in}{0.005773in}}{\pgfqpoint{0.007857in}{0.007857in}}%
\pgfpathcurveto{\pgfqpoint{0.005773in}{0.009940in}}{\pgfqpoint{0.002947in}{0.011111in}}{\pgfqpoint{0.000000in}{0.011111in}}%
\pgfpathcurveto{\pgfqpoint{-0.002947in}{0.011111in}}{\pgfqpoint{-0.005773in}{0.009940in}}{\pgfqpoint{-0.007857in}{0.007857in}}%
\pgfpathcurveto{\pgfqpoint{-0.009940in}{0.005773in}}{\pgfqpoint{-0.011111in}{0.002947in}}{\pgfqpoint{-0.011111in}{0.000000in}}%
\pgfpathcurveto{\pgfqpoint{-0.011111in}{-0.002947in}}{\pgfqpoint{-0.009940in}{-0.005773in}}{\pgfqpoint{-0.007857in}{-0.007857in}}%
\pgfpathcurveto{\pgfqpoint{-0.005773in}{-0.009940in}}{\pgfqpoint{-0.002947in}{-0.011111in}}{\pgfqpoint{0.000000in}{-0.011111in}}%
\pgfpathlineto{\pgfqpoint{0.000000in}{-0.011111in}}%
\pgfpathclose%
\pgfusepath{stroke,fill}%
}%
\begin{pgfscope}%
\pgfsys@transformshift{3.205657in}{4.006711in}%
\pgfsys@useobject{currentmarker}{}%
\end{pgfscope}%
\end{pgfscope}%
\begin{pgfscope}%
\pgfpathrectangle{\pgfqpoint{0.100000in}{2.493063in}}{\pgfqpoint{5.037500in}{3.427208in}}%
\pgfusepath{clip}%
\pgfsetrectcap%
\pgfsetroundjoin%
\pgfsetlinewidth{1.505625pt}%
\definecolor{currentstroke}{rgb}{0.501961,0.501961,0.501961}%
\pgfsetstrokecolor{currentstroke}%
\pgfsetstrokeopacity{0.500000}%
\pgfsetdash{}{0pt}%
\pgfpathmoveto{\pgfqpoint{3.061256in}{3.874532in}}%
\pgfusepath{stroke}%
\end{pgfscope}%
\begin{pgfscope}%
\pgfpathrectangle{\pgfqpoint{0.100000in}{2.493063in}}{\pgfqpoint{5.037500in}{3.427208in}}%
\pgfusepath{clip}%
\pgfsetbuttcap%
\pgfsetroundjoin%
\definecolor{currentfill}{rgb}{0.501961,0.501961,0.501961}%
\pgfsetfillcolor{currentfill}%
\pgfsetfillopacity{0.500000}%
\pgfsetlinewidth{0.250937pt}%
\definecolor{currentstroke}{rgb}{0.000000,0.000000,0.000000}%
\pgfsetstrokecolor{currentstroke}%
\pgfsetstrokeopacity{0.500000}%
\pgfsetdash{}{0pt}%
\pgfsys@defobject{currentmarker}{\pgfqpoint{-0.013889in}{-0.013889in}}{\pgfqpoint{0.013889in}{0.013889in}}{%
\pgfpathmoveto{\pgfqpoint{0.000000in}{-0.013889in}}%
\pgfpathcurveto{\pgfqpoint{0.003683in}{-0.013889in}}{\pgfqpoint{0.007216in}{-0.012425in}}{\pgfqpoint{0.009821in}{-0.009821in}}%
\pgfpathcurveto{\pgfqpoint{0.012425in}{-0.007216in}}{\pgfqpoint{0.013889in}{-0.003683in}}{\pgfqpoint{0.013889in}{0.000000in}}%
\pgfpathcurveto{\pgfqpoint{0.013889in}{0.003683in}}{\pgfqpoint{0.012425in}{0.007216in}}{\pgfqpoint{0.009821in}{0.009821in}}%
\pgfpathcurveto{\pgfqpoint{0.007216in}{0.012425in}}{\pgfqpoint{0.003683in}{0.013889in}}{\pgfqpoint{0.000000in}{0.013889in}}%
\pgfpathcurveto{\pgfqpoint{-0.003683in}{0.013889in}}{\pgfqpoint{-0.007216in}{0.012425in}}{\pgfqpoint{-0.009821in}{0.009821in}}%
\pgfpathcurveto{\pgfqpoint{-0.012425in}{0.007216in}}{\pgfqpoint{-0.013889in}{0.003683in}}{\pgfqpoint{-0.013889in}{0.000000in}}%
\pgfpathcurveto{\pgfqpoint{-0.013889in}{-0.003683in}}{\pgfqpoint{-0.012425in}{-0.007216in}}{\pgfqpoint{-0.009821in}{-0.009821in}}%
\pgfpathcurveto{\pgfqpoint{-0.007216in}{-0.012425in}}{\pgfqpoint{-0.003683in}{-0.013889in}}{\pgfqpoint{0.000000in}{-0.013889in}}%
\pgfpathlineto{\pgfqpoint{0.000000in}{-0.013889in}}%
\pgfpathclose%
\pgfusepath{stroke,fill}%
}%
\begin{pgfscope}%
\pgfsys@transformshift{3.061256in}{3.874532in}%
\pgfsys@useobject{currentmarker}{}%
\end{pgfscope}%
\end{pgfscope}%
\begin{pgfscope}%
\pgfpathrectangle{\pgfqpoint{0.100000in}{2.493063in}}{\pgfqpoint{5.037500in}{3.427208in}}%
\pgfusepath{clip}%
\pgfsetrectcap%
\pgfsetroundjoin%
\pgfsetlinewidth{1.505625pt}%
\definecolor{currentstroke}{rgb}{0.678431,1.000000,0.184314}%
\pgfsetstrokecolor{currentstroke}%
\pgfsetstrokeopacity{0.500000}%
\pgfsetdash{}{0pt}%
\pgfpathmoveto{\pgfqpoint{3.090389in}{3.815577in}}%
\pgfusepath{stroke}%
\end{pgfscope}%
\begin{pgfscope}%
\pgfpathrectangle{\pgfqpoint{0.100000in}{2.493063in}}{\pgfqpoint{5.037500in}{3.427208in}}%
\pgfusepath{clip}%
\pgfsetbuttcap%
\pgfsetroundjoin%
\definecolor{currentfill}{rgb}{0.678431,1.000000,0.184314}%
\pgfsetfillcolor{currentfill}%
\pgfsetfillopacity{0.500000}%
\pgfsetlinewidth{0.250937pt}%
\definecolor{currentstroke}{rgb}{0.000000,0.000000,0.000000}%
\pgfsetstrokecolor{currentstroke}%
\pgfsetstrokeopacity{0.500000}%
\pgfsetdash{}{0pt}%
\pgfsys@defobject{currentmarker}{\pgfqpoint{-0.011111in}{-0.011111in}}{\pgfqpoint{0.011111in}{0.011111in}}{%
\pgfpathmoveto{\pgfqpoint{0.000000in}{-0.011111in}}%
\pgfpathcurveto{\pgfqpoint{0.002947in}{-0.011111in}}{\pgfqpoint{0.005773in}{-0.009940in}}{\pgfqpoint{0.007857in}{-0.007857in}}%
\pgfpathcurveto{\pgfqpoint{0.009940in}{-0.005773in}}{\pgfqpoint{0.011111in}{-0.002947in}}{\pgfqpoint{0.011111in}{0.000000in}}%
\pgfpathcurveto{\pgfqpoint{0.011111in}{0.002947in}}{\pgfqpoint{0.009940in}{0.005773in}}{\pgfqpoint{0.007857in}{0.007857in}}%
\pgfpathcurveto{\pgfqpoint{0.005773in}{0.009940in}}{\pgfqpoint{0.002947in}{0.011111in}}{\pgfqpoint{0.000000in}{0.011111in}}%
\pgfpathcurveto{\pgfqpoint{-0.002947in}{0.011111in}}{\pgfqpoint{-0.005773in}{0.009940in}}{\pgfqpoint{-0.007857in}{0.007857in}}%
\pgfpathcurveto{\pgfqpoint{-0.009940in}{0.005773in}}{\pgfqpoint{-0.011111in}{0.002947in}}{\pgfqpoint{-0.011111in}{0.000000in}}%
\pgfpathcurveto{\pgfqpoint{-0.011111in}{-0.002947in}}{\pgfqpoint{-0.009940in}{-0.005773in}}{\pgfqpoint{-0.007857in}{-0.007857in}}%
\pgfpathcurveto{\pgfqpoint{-0.005773in}{-0.009940in}}{\pgfqpoint{-0.002947in}{-0.011111in}}{\pgfqpoint{0.000000in}{-0.011111in}}%
\pgfpathlineto{\pgfqpoint{0.000000in}{-0.011111in}}%
\pgfpathclose%
\pgfusepath{stroke,fill}%
}%
\begin{pgfscope}%
\pgfsys@transformshift{3.090389in}{3.815577in}%
\pgfsys@useobject{currentmarker}{}%
\end{pgfscope}%
\end{pgfscope}%
\begin{pgfscope}%
\pgfpathrectangle{\pgfqpoint{0.100000in}{2.493063in}}{\pgfqpoint{5.037500in}{3.427208in}}%
\pgfusepath{clip}%
\pgfsetrectcap%
\pgfsetroundjoin%
\pgfsetlinewidth{1.505625pt}%
\definecolor{currentstroke}{rgb}{0.000000,0.000000,1.000000}%
\pgfsetstrokecolor{currentstroke}%
\pgfsetstrokeopacity{0.500000}%
\pgfsetdash{}{0pt}%
\pgfpathmoveto{\pgfqpoint{0.570500in}{4.240029in}}%
\pgfusepath{stroke}%
\end{pgfscope}%
\begin{pgfscope}%
\pgfpathrectangle{\pgfqpoint{0.100000in}{2.493063in}}{\pgfqpoint{5.037500in}{3.427208in}}%
\pgfusepath{clip}%
\pgfsetbuttcap%
\pgfsetroundjoin%
\definecolor{currentfill}{rgb}{0.000000,0.000000,1.000000}%
\pgfsetfillcolor{currentfill}%
\pgfsetfillopacity{0.500000}%
\pgfsetlinewidth{0.250937pt}%
\definecolor{currentstroke}{rgb}{0.000000,0.000000,0.000000}%
\pgfsetstrokecolor{currentstroke}%
\pgfsetstrokeopacity{0.500000}%
\pgfsetdash{}{0pt}%
\pgfsys@defobject{currentmarker}{\pgfqpoint{-0.030556in}{-0.030556in}}{\pgfqpoint{0.030556in}{0.030556in}}{%
\pgfpathmoveto{\pgfqpoint{0.000000in}{-0.030556in}}%
\pgfpathcurveto{\pgfqpoint{0.008103in}{-0.030556in}}{\pgfqpoint{0.015876in}{-0.027336in}}{\pgfqpoint{0.021606in}{-0.021606in}}%
\pgfpathcurveto{\pgfqpoint{0.027336in}{-0.015876in}}{\pgfqpoint{0.030556in}{-0.008103in}}{\pgfqpoint{0.030556in}{0.000000in}}%
\pgfpathcurveto{\pgfqpoint{0.030556in}{0.008103in}}{\pgfqpoint{0.027336in}{0.015876in}}{\pgfqpoint{0.021606in}{0.021606in}}%
\pgfpathcurveto{\pgfqpoint{0.015876in}{0.027336in}}{\pgfqpoint{0.008103in}{0.030556in}}{\pgfqpoint{0.000000in}{0.030556in}}%
\pgfpathcurveto{\pgfqpoint{-0.008103in}{0.030556in}}{\pgfqpoint{-0.015876in}{0.027336in}}{\pgfqpoint{-0.021606in}{0.021606in}}%
\pgfpathcurveto{\pgfqpoint{-0.027336in}{0.015876in}}{\pgfqpoint{-0.030556in}{0.008103in}}{\pgfqpoint{-0.030556in}{0.000000in}}%
\pgfpathcurveto{\pgfqpoint{-0.030556in}{-0.008103in}}{\pgfqpoint{-0.027336in}{-0.015876in}}{\pgfqpoint{-0.021606in}{-0.021606in}}%
\pgfpathcurveto{\pgfqpoint{-0.015876in}{-0.027336in}}{\pgfqpoint{-0.008103in}{-0.030556in}}{\pgfqpoint{0.000000in}{-0.030556in}}%
\pgfpathlineto{\pgfqpoint{0.000000in}{-0.030556in}}%
\pgfpathclose%
\pgfusepath{stroke,fill}%
}%
\begin{pgfscope}%
\pgfsys@transformshift{0.570500in}{4.240029in}%
\pgfsys@useobject{currentmarker}{}%
\end{pgfscope}%
\end{pgfscope}%
\begin{pgfscope}%
\pgfpathrectangle{\pgfqpoint{0.100000in}{2.493063in}}{\pgfqpoint{5.037500in}{3.427208in}}%
\pgfusepath{clip}%
\pgfsetrectcap%
\pgfsetroundjoin%
\pgfsetlinewidth{1.505625pt}%
\definecolor{currentstroke}{rgb}{0.000000,0.000000,1.000000}%
\pgfsetstrokecolor{currentstroke}%
\pgfsetstrokeopacity{0.500000}%
\pgfsetdash{}{0pt}%
\pgfpathmoveto{\pgfqpoint{0.461714in}{4.793248in}}%
\pgfusepath{stroke}%
\end{pgfscope}%
\begin{pgfscope}%
\pgfpathrectangle{\pgfqpoint{0.100000in}{2.493063in}}{\pgfqpoint{5.037500in}{3.427208in}}%
\pgfusepath{clip}%
\pgfsetbuttcap%
\pgfsetroundjoin%
\definecolor{currentfill}{rgb}{0.000000,0.000000,1.000000}%
\pgfsetfillcolor{currentfill}%
\pgfsetfillopacity{0.500000}%
\pgfsetlinewidth{0.250937pt}%
\definecolor{currentstroke}{rgb}{0.000000,0.000000,0.000000}%
\pgfsetstrokecolor{currentstroke}%
\pgfsetstrokeopacity{0.500000}%
\pgfsetdash{}{0pt}%
\pgfsys@defobject{currentmarker}{\pgfqpoint{-0.022222in}{-0.022222in}}{\pgfqpoint{0.022222in}{0.022222in}}{%
\pgfpathmoveto{\pgfqpoint{0.000000in}{-0.022222in}}%
\pgfpathcurveto{\pgfqpoint{0.005893in}{-0.022222in}}{\pgfqpoint{0.011546in}{-0.019881in}}{\pgfqpoint{0.015713in}{-0.015713in}}%
\pgfpathcurveto{\pgfqpoint{0.019881in}{-0.011546in}}{\pgfqpoint{0.022222in}{-0.005893in}}{\pgfqpoint{0.022222in}{0.000000in}}%
\pgfpathcurveto{\pgfqpoint{0.022222in}{0.005893in}}{\pgfqpoint{0.019881in}{0.011546in}}{\pgfqpoint{0.015713in}{0.015713in}}%
\pgfpathcurveto{\pgfqpoint{0.011546in}{0.019881in}}{\pgfqpoint{0.005893in}{0.022222in}}{\pgfqpoint{0.000000in}{0.022222in}}%
\pgfpathcurveto{\pgfqpoint{-0.005893in}{0.022222in}}{\pgfqpoint{-0.011546in}{0.019881in}}{\pgfqpoint{-0.015713in}{0.015713in}}%
\pgfpathcurveto{\pgfqpoint{-0.019881in}{0.011546in}}{\pgfqpoint{-0.022222in}{0.005893in}}{\pgfqpoint{-0.022222in}{0.000000in}}%
\pgfpathcurveto{\pgfqpoint{-0.022222in}{-0.005893in}}{\pgfqpoint{-0.019881in}{-0.011546in}}{\pgfqpoint{-0.015713in}{-0.015713in}}%
\pgfpathcurveto{\pgfqpoint{-0.011546in}{-0.019881in}}{\pgfqpoint{-0.005893in}{-0.022222in}}{\pgfqpoint{0.000000in}{-0.022222in}}%
\pgfpathlineto{\pgfqpoint{0.000000in}{-0.022222in}}%
\pgfpathclose%
\pgfusepath{stroke,fill}%
}%
\begin{pgfscope}%
\pgfsys@transformshift{0.461714in}{4.793248in}%
\pgfsys@useobject{currentmarker}{}%
\end{pgfscope}%
\end{pgfscope}%
\begin{pgfscope}%
\pgfpathrectangle{\pgfqpoint{0.100000in}{2.493063in}}{\pgfqpoint{5.037500in}{3.427208in}}%
\pgfusepath{clip}%
\pgfsetrectcap%
\pgfsetroundjoin%
\pgfsetlinewidth{1.505625pt}%
\definecolor{currentstroke}{rgb}{0.678431,1.000000,0.184314}%
\pgfsetstrokecolor{currentstroke}%
\pgfsetstrokeopacity{0.500000}%
\pgfsetdash{}{0pt}%
\pgfpathmoveto{\pgfqpoint{0.818988in}{3.870241in}}%
\pgfusepath{stroke}%
\end{pgfscope}%
\begin{pgfscope}%
\pgfpathrectangle{\pgfqpoint{0.100000in}{2.493063in}}{\pgfqpoint{5.037500in}{3.427208in}}%
\pgfusepath{clip}%
\pgfsetbuttcap%
\pgfsetroundjoin%
\definecolor{currentfill}{rgb}{0.678431,1.000000,0.184314}%
\pgfsetfillcolor{currentfill}%
\pgfsetfillopacity{0.500000}%
\pgfsetlinewidth{0.250937pt}%
\definecolor{currentstroke}{rgb}{0.000000,0.000000,0.000000}%
\pgfsetstrokecolor{currentstroke}%
\pgfsetstrokeopacity{0.500000}%
\pgfsetdash{}{0pt}%
\pgfsys@defobject{currentmarker}{\pgfqpoint{-0.091667in}{-0.091667in}}{\pgfqpoint{0.091667in}{0.091667in}}{%
\pgfpathmoveto{\pgfqpoint{0.000000in}{-0.091667in}}%
\pgfpathcurveto{\pgfqpoint{0.024310in}{-0.091667in}}{\pgfqpoint{0.047628in}{-0.082008in}}{\pgfqpoint{0.064818in}{-0.064818in}}%
\pgfpathcurveto{\pgfqpoint{0.082008in}{-0.047628in}}{\pgfqpoint{0.091667in}{-0.024310in}}{\pgfqpoint{0.091667in}{0.000000in}}%
\pgfpathcurveto{\pgfqpoint{0.091667in}{0.024310in}}{\pgfqpoint{0.082008in}{0.047628in}}{\pgfqpoint{0.064818in}{0.064818in}}%
\pgfpathcurveto{\pgfqpoint{0.047628in}{0.082008in}}{\pgfqpoint{0.024310in}{0.091667in}}{\pgfqpoint{0.000000in}{0.091667in}}%
\pgfpathcurveto{\pgfqpoint{-0.024310in}{0.091667in}}{\pgfqpoint{-0.047628in}{0.082008in}}{\pgfqpoint{-0.064818in}{0.064818in}}%
\pgfpathcurveto{\pgfqpoint{-0.082008in}{0.047628in}}{\pgfqpoint{-0.091667in}{0.024310in}}{\pgfqpoint{-0.091667in}{0.000000in}}%
\pgfpathcurveto{\pgfqpoint{-0.091667in}{-0.024310in}}{\pgfqpoint{-0.082008in}{-0.047628in}}{\pgfqpoint{-0.064818in}{-0.064818in}}%
\pgfpathcurveto{\pgfqpoint{-0.047628in}{-0.082008in}}{\pgfqpoint{-0.024310in}{-0.091667in}}{\pgfqpoint{0.000000in}{-0.091667in}}%
\pgfpathlineto{\pgfqpoint{0.000000in}{-0.091667in}}%
\pgfpathclose%
\pgfusepath{stroke,fill}%
}%
\begin{pgfscope}%
\pgfsys@transformshift{0.818988in}{3.870241in}%
\pgfsys@useobject{currentmarker}{}%
\end{pgfscope}%
\end{pgfscope}%
\begin{pgfscope}%
\pgfpathrectangle{\pgfqpoint{0.100000in}{2.493063in}}{\pgfqpoint{5.037500in}{3.427208in}}%
\pgfusepath{clip}%
\pgfsetrectcap%
\pgfsetroundjoin%
\pgfsetlinewidth{1.505625pt}%
\definecolor{currentstroke}{rgb}{0.000000,0.000000,1.000000}%
\pgfsetstrokecolor{currentstroke}%
\pgfsetstrokeopacity{0.500000}%
\pgfsetdash{}{0pt}%
\pgfpathmoveto{\pgfqpoint{0.549937in}{4.407788in}}%
\pgfusepath{stroke}%
\end{pgfscope}%
\begin{pgfscope}%
\pgfpathrectangle{\pgfqpoint{0.100000in}{2.493063in}}{\pgfqpoint{5.037500in}{3.427208in}}%
\pgfusepath{clip}%
\pgfsetbuttcap%
\pgfsetroundjoin%
\definecolor{currentfill}{rgb}{0.000000,0.000000,1.000000}%
\pgfsetfillcolor{currentfill}%
\pgfsetfillopacity{0.500000}%
\pgfsetlinewidth{0.250937pt}%
\definecolor{currentstroke}{rgb}{0.000000,0.000000,0.000000}%
\pgfsetstrokecolor{currentstroke}%
\pgfsetstrokeopacity{0.500000}%
\pgfsetdash{}{0pt}%
\pgfsys@defobject{currentmarker}{\pgfqpoint{-0.019444in}{-0.019444in}}{\pgfqpoint{0.019444in}{0.019444in}}{%
\pgfpathmoveto{\pgfqpoint{0.000000in}{-0.019444in}}%
\pgfpathcurveto{\pgfqpoint{0.005157in}{-0.019444in}}{\pgfqpoint{0.010103in}{-0.017396in}}{\pgfqpoint{0.013749in}{-0.013749in}}%
\pgfpathcurveto{\pgfqpoint{0.017396in}{-0.010103in}}{\pgfqpoint{0.019444in}{-0.005157in}}{\pgfqpoint{0.019444in}{0.000000in}}%
\pgfpathcurveto{\pgfqpoint{0.019444in}{0.005157in}}{\pgfqpoint{0.017396in}{0.010103in}}{\pgfqpoint{0.013749in}{0.013749in}}%
\pgfpathcurveto{\pgfqpoint{0.010103in}{0.017396in}}{\pgfqpoint{0.005157in}{0.019444in}}{\pgfqpoint{0.000000in}{0.019444in}}%
\pgfpathcurveto{\pgfqpoint{-0.005157in}{0.019444in}}{\pgfqpoint{-0.010103in}{0.017396in}}{\pgfqpoint{-0.013749in}{0.013749in}}%
\pgfpathcurveto{\pgfqpoint{-0.017396in}{0.010103in}}{\pgfqpoint{-0.019444in}{0.005157in}}{\pgfqpoint{-0.019444in}{0.000000in}}%
\pgfpathcurveto{\pgfqpoint{-0.019444in}{-0.005157in}}{\pgfqpoint{-0.017396in}{-0.010103in}}{\pgfqpoint{-0.013749in}{-0.013749in}}%
\pgfpathcurveto{\pgfqpoint{-0.010103in}{-0.017396in}}{\pgfqpoint{-0.005157in}{-0.019444in}}{\pgfqpoint{0.000000in}{-0.019444in}}%
\pgfpathlineto{\pgfqpoint{0.000000in}{-0.019444in}}%
\pgfpathclose%
\pgfusepath{stroke,fill}%
}%
\begin{pgfscope}%
\pgfsys@transformshift{0.549937in}{4.407788in}%
\pgfsys@useobject{currentmarker}{}%
\end{pgfscope}%
\end{pgfscope}%
\begin{pgfscope}%
\pgfpathrectangle{\pgfqpoint{0.100000in}{2.493063in}}{\pgfqpoint{5.037500in}{3.427208in}}%
\pgfusepath{clip}%
\pgfsetrectcap%
\pgfsetroundjoin%
\pgfsetlinewidth{1.505625pt}%
\definecolor{currentstroke}{rgb}{0.000000,0.000000,1.000000}%
\pgfsetstrokecolor{currentstroke}%
\pgfsetstrokeopacity{0.500000}%
\pgfsetdash{}{0pt}%
\pgfpathmoveto{\pgfqpoint{0.543143in}{4.361573in}}%
\pgfusepath{stroke}%
\end{pgfscope}%
\begin{pgfscope}%
\pgfpathrectangle{\pgfqpoint{0.100000in}{2.493063in}}{\pgfqpoint{5.037500in}{3.427208in}}%
\pgfusepath{clip}%
\pgfsetbuttcap%
\pgfsetroundjoin%
\definecolor{currentfill}{rgb}{0.000000,0.000000,1.000000}%
\pgfsetfillcolor{currentfill}%
\pgfsetfillopacity{0.500000}%
\pgfsetlinewidth{0.250937pt}%
\definecolor{currentstroke}{rgb}{0.000000,0.000000,0.000000}%
\pgfsetstrokecolor{currentstroke}%
\pgfsetstrokeopacity{0.500000}%
\pgfsetdash{}{0pt}%
\pgfsys@defobject{currentmarker}{\pgfqpoint{-0.013889in}{-0.013889in}}{\pgfqpoint{0.013889in}{0.013889in}}{%
\pgfpathmoveto{\pgfqpoint{0.000000in}{-0.013889in}}%
\pgfpathcurveto{\pgfqpoint{0.003683in}{-0.013889in}}{\pgfqpoint{0.007216in}{-0.012425in}}{\pgfqpoint{0.009821in}{-0.009821in}}%
\pgfpathcurveto{\pgfqpoint{0.012425in}{-0.007216in}}{\pgfqpoint{0.013889in}{-0.003683in}}{\pgfqpoint{0.013889in}{0.000000in}}%
\pgfpathcurveto{\pgfqpoint{0.013889in}{0.003683in}}{\pgfqpoint{0.012425in}{0.007216in}}{\pgfqpoint{0.009821in}{0.009821in}}%
\pgfpathcurveto{\pgfqpoint{0.007216in}{0.012425in}}{\pgfqpoint{0.003683in}{0.013889in}}{\pgfqpoint{0.000000in}{0.013889in}}%
\pgfpathcurveto{\pgfqpoint{-0.003683in}{0.013889in}}{\pgfqpoint{-0.007216in}{0.012425in}}{\pgfqpoint{-0.009821in}{0.009821in}}%
\pgfpathcurveto{\pgfqpoint{-0.012425in}{0.007216in}}{\pgfqpoint{-0.013889in}{0.003683in}}{\pgfqpoint{-0.013889in}{0.000000in}}%
\pgfpathcurveto{\pgfqpoint{-0.013889in}{-0.003683in}}{\pgfqpoint{-0.012425in}{-0.007216in}}{\pgfqpoint{-0.009821in}{-0.009821in}}%
\pgfpathcurveto{\pgfqpoint{-0.007216in}{-0.012425in}}{\pgfqpoint{-0.003683in}{-0.013889in}}{\pgfqpoint{0.000000in}{-0.013889in}}%
\pgfpathlineto{\pgfqpoint{0.000000in}{-0.013889in}}%
\pgfpathclose%
\pgfusepath{stroke,fill}%
}%
\begin{pgfscope}%
\pgfsys@transformshift{0.543143in}{4.361573in}%
\pgfsys@useobject{currentmarker}{}%
\end{pgfscope}%
\end{pgfscope}%
\begin{pgfscope}%
\pgfpathrectangle{\pgfqpoint{0.100000in}{2.493063in}}{\pgfqpoint{5.037500in}{3.427208in}}%
\pgfusepath{clip}%
\pgfsetrectcap%
\pgfsetroundjoin%
\pgfsetlinewidth{1.505625pt}%
\definecolor{currentstroke}{rgb}{0.000000,0.000000,1.000000}%
\pgfsetstrokecolor{currentstroke}%
\pgfsetstrokeopacity{0.500000}%
\pgfsetdash{}{0pt}%
\pgfpathmoveto{\pgfqpoint{0.602428in}{4.073927in}}%
\pgfusepath{stroke}%
\end{pgfscope}%
\begin{pgfscope}%
\pgfpathrectangle{\pgfqpoint{0.100000in}{2.493063in}}{\pgfqpoint{5.037500in}{3.427208in}}%
\pgfusepath{clip}%
\pgfsetbuttcap%
\pgfsetroundjoin%
\definecolor{currentfill}{rgb}{0.000000,0.000000,1.000000}%
\pgfsetfillcolor{currentfill}%
\pgfsetfillopacity{0.500000}%
\pgfsetlinewidth{0.250937pt}%
\definecolor{currentstroke}{rgb}{0.000000,0.000000,0.000000}%
\pgfsetstrokecolor{currentstroke}%
\pgfsetstrokeopacity{0.500000}%
\pgfsetdash{}{0pt}%
\pgfsys@defobject{currentmarker}{\pgfqpoint{-0.022222in}{-0.022222in}}{\pgfqpoint{0.022222in}{0.022222in}}{%
\pgfpathmoveto{\pgfqpoint{0.000000in}{-0.022222in}}%
\pgfpathcurveto{\pgfqpoint{0.005893in}{-0.022222in}}{\pgfqpoint{0.011546in}{-0.019881in}}{\pgfqpoint{0.015713in}{-0.015713in}}%
\pgfpathcurveto{\pgfqpoint{0.019881in}{-0.011546in}}{\pgfqpoint{0.022222in}{-0.005893in}}{\pgfqpoint{0.022222in}{0.000000in}}%
\pgfpathcurveto{\pgfqpoint{0.022222in}{0.005893in}}{\pgfqpoint{0.019881in}{0.011546in}}{\pgfqpoint{0.015713in}{0.015713in}}%
\pgfpathcurveto{\pgfqpoint{0.011546in}{0.019881in}}{\pgfqpoint{0.005893in}{0.022222in}}{\pgfqpoint{0.000000in}{0.022222in}}%
\pgfpathcurveto{\pgfqpoint{-0.005893in}{0.022222in}}{\pgfqpoint{-0.011546in}{0.019881in}}{\pgfqpoint{-0.015713in}{0.015713in}}%
\pgfpathcurveto{\pgfqpoint{-0.019881in}{0.011546in}}{\pgfqpoint{-0.022222in}{0.005893in}}{\pgfqpoint{-0.022222in}{0.000000in}}%
\pgfpathcurveto{\pgfqpoint{-0.022222in}{-0.005893in}}{\pgfqpoint{-0.019881in}{-0.011546in}}{\pgfqpoint{-0.015713in}{-0.015713in}}%
\pgfpathcurveto{\pgfqpoint{-0.011546in}{-0.019881in}}{\pgfqpoint{-0.005893in}{-0.022222in}}{\pgfqpoint{0.000000in}{-0.022222in}}%
\pgfpathlineto{\pgfqpoint{0.000000in}{-0.022222in}}%
\pgfpathclose%
\pgfusepath{stroke,fill}%
}%
\begin{pgfscope}%
\pgfsys@transformshift{0.602428in}{4.073927in}%
\pgfsys@useobject{currentmarker}{}%
\end{pgfscope}%
\end{pgfscope}%
\begin{pgfscope}%
\pgfpathrectangle{\pgfqpoint{0.100000in}{2.493063in}}{\pgfqpoint{5.037500in}{3.427208in}}%
\pgfusepath{clip}%
\pgfsetrectcap%
\pgfsetroundjoin%
\pgfsetlinewidth{1.505625pt}%
\definecolor{currentstroke}{rgb}{0.000000,0.000000,1.000000}%
\pgfsetstrokecolor{currentstroke}%
\pgfsetstrokeopacity{0.500000}%
\pgfsetdash{}{0pt}%
\pgfpathmoveto{\pgfqpoint{0.557875in}{4.458461in}}%
\pgfusepath{stroke}%
\end{pgfscope}%
\begin{pgfscope}%
\pgfpathrectangle{\pgfqpoint{0.100000in}{2.493063in}}{\pgfqpoint{5.037500in}{3.427208in}}%
\pgfusepath{clip}%
\pgfsetbuttcap%
\pgfsetroundjoin%
\definecolor{currentfill}{rgb}{0.000000,0.000000,1.000000}%
\pgfsetfillcolor{currentfill}%
\pgfsetfillopacity{0.500000}%
\pgfsetlinewidth{0.250937pt}%
\definecolor{currentstroke}{rgb}{0.000000,0.000000,0.000000}%
\pgfsetstrokecolor{currentstroke}%
\pgfsetstrokeopacity{0.500000}%
\pgfsetdash{}{0pt}%
\pgfsys@defobject{currentmarker}{\pgfqpoint{-0.025000in}{-0.025000in}}{\pgfqpoint{0.025000in}{0.025000in}}{%
\pgfpathmoveto{\pgfqpoint{0.000000in}{-0.025000in}}%
\pgfpathcurveto{\pgfqpoint{0.006630in}{-0.025000in}}{\pgfqpoint{0.012989in}{-0.022366in}}{\pgfqpoint{0.017678in}{-0.017678in}}%
\pgfpathcurveto{\pgfqpoint{0.022366in}{-0.012989in}}{\pgfqpoint{0.025000in}{-0.006630in}}{\pgfqpoint{0.025000in}{0.000000in}}%
\pgfpathcurveto{\pgfqpoint{0.025000in}{0.006630in}}{\pgfqpoint{0.022366in}{0.012989in}}{\pgfqpoint{0.017678in}{0.017678in}}%
\pgfpathcurveto{\pgfqpoint{0.012989in}{0.022366in}}{\pgfqpoint{0.006630in}{0.025000in}}{\pgfqpoint{0.000000in}{0.025000in}}%
\pgfpathcurveto{\pgfqpoint{-0.006630in}{0.025000in}}{\pgfqpoint{-0.012989in}{0.022366in}}{\pgfqpoint{-0.017678in}{0.017678in}}%
\pgfpathcurveto{\pgfqpoint{-0.022366in}{0.012989in}}{\pgfqpoint{-0.025000in}{0.006630in}}{\pgfqpoint{-0.025000in}{0.000000in}}%
\pgfpathcurveto{\pgfqpoint{-0.025000in}{-0.006630in}}{\pgfqpoint{-0.022366in}{-0.012989in}}{\pgfqpoint{-0.017678in}{-0.017678in}}%
\pgfpathcurveto{\pgfqpoint{-0.012989in}{-0.022366in}}{\pgfqpoint{-0.006630in}{-0.025000in}}{\pgfqpoint{0.000000in}{-0.025000in}}%
\pgfpathlineto{\pgfqpoint{0.000000in}{-0.025000in}}%
\pgfpathclose%
\pgfusepath{stroke,fill}%
}%
\begin{pgfscope}%
\pgfsys@transformshift{0.557875in}{4.458461in}%
\pgfsys@useobject{currentmarker}{}%
\end{pgfscope}%
\end{pgfscope}%
\begin{pgfscope}%
\pgfpathrectangle{\pgfqpoint{0.100000in}{2.493063in}}{\pgfqpoint{5.037500in}{3.427208in}}%
\pgfusepath{clip}%
\pgfsetrectcap%
\pgfsetroundjoin%
\pgfsetlinewidth{1.505625pt}%
\definecolor{currentstroke}{rgb}{0.000000,0.000000,1.000000}%
\pgfsetstrokecolor{currentstroke}%
\pgfsetstrokeopacity{0.500000}%
\pgfsetdash{}{0pt}%
\pgfpathmoveto{\pgfqpoint{0.499312in}{4.490808in}}%
\pgfusepath{stroke}%
\end{pgfscope}%
\begin{pgfscope}%
\pgfpathrectangle{\pgfqpoint{0.100000in}{2.493063in}}{\pgfqpoint{5.037500in}{3.427208in}}%
\pgfusepath{clip}%
\pgfsetbuttcap%
\pgfsetroundjoin%
\definecolor{currentfill}{rgb}{0.000000,0.000000,1.000000}%
\pgfsetfillcolor{currentfill}%
\pgfsetfillopacity{0.500000}%
\pgfsetlinewidth{0.250937pt}%
\definecolor{currentstroke}{rgb}{0.000000,0.000000,0.000000}%
\pgfsetstrokecolor{currentstroke}%
\pgfsetstrokeopacity{0.500000}%
\pgfsetdash{}{0pt}%
\pgfsys@defobject{currentmarker}{\pgfqpoint{-0.044444in}{-0.044444in}}{\pgfqpoint{0.044444in}{0.044444in}}{%
\pgfpathmoveto{\pgfqpoint{0.000000in}{-0.044444in}}%
\pgfpathcurveto{\pgfqpoint{0.011787in}{-0.044444in}}{\pgfqpoint{0.023092in}{-0.039761in}}{\pgfqpoint{0.031427in}{-0.031427in}}%
\pgfpathcurveto{\pgfqpoint{0.039761in}{-0.023092in}}{\pgfqpoint{0.044444in}{-0.011787in}}{\pgfqpoint{0.044444in}{0.000000in}}%
\pgfpathcurveto{\pgfqpoint{0.044444in}{0.011787in}}{\pgfqpoint{0.039761in}{0.023092in}}{\pgfqpoint{0.031427in}{0.031427in}}%
\pgfpathcurveto{\pgfqpoint{0.023092in}{0.039761in}}{\pgfqpoint{0.011787in}{0.044444in}}{\pgfqpoint{0.000000in}{0.044444in}}%
\pgfpathcurveto{\pgfqpoint{-0.011787in}{0.044444in}}{\pgfqpoint{-0.023092in}{0.039761in}}{\pgfqpoint{-0.031427in}{0.031427in}}%
\pgfpathcurveto{\pgfqpoint{-0.039761in}{0.023092in}}{\pgfqpoint{-0.044444in}{0.011787in}}{\pgfqpoint{-0.044444in}{0.000000in}}%
\pgfpathcurveto{\pgfqpoint{-0.044444in}{-0.011787in}}{\pgfqpoint{-0.039761in}{-0.023092in}}{\pgfqpoint{-0.031427in}{-0.031427in}}%
\pgfpathcurveto{\pgfqpoint{-0.023092in}{-0.039761in}}{\pgfqpoint{-0.011787in}{-0.044444in}}{\pgfqpoint{0.000000in}{-0.044444in}}%
\pgfpathlineto{\pgfqpoint{0.000000in}{-0.044444in}}%
\pgfpathclose%
\pgfusepath{stroke,fill}%
}%
\begin{pgfscope}%
\pgfsys@transformshift{0.499312in}{4.490808in}%
\pgfsys@useobject{currentmarker}{}%
\end{pgfscope}%
\end{pgfscope}%
\begin{pgfscope}%
\pgfpathrectangle{\pgfqpoint{0.100000in}{2.493063in}}{\pgfqpoint{5.037500in}{3.427208in}}%
\pgfusepath{clip}%
\pgfsetrectcap%
\pgfsetroundjoin%
\pgfsetlinewidth{1.505625pt}%
\definecolor{currentstroke}{rgb}{0.000000,0.000000,1.000000}%
\pgfsetstrokecolor{currentstroke}%
\pgfsetstrokeopacity{0.500000}%
\pgfsetdash{}{0pt}%
\pgfpathmoveto{\pgfqpoint{0.465142in}{4.541082in}}%
\pgfusepath{stroke}%
\end{pgfscope}%
\begin{pgfscope}%
\pgfpathrectangle{\pgfqpoint{0.100000in}{2.493063in}}{\pgfqpoint{5.037500in}{3.427208in}}%
\pgfusepath{clip}%
\pgfsetbuttcap%
\pgfsetroundjoin%
\definecolor{currentfill}{rgb}{0.000000,0.000000,1.000000}%
\pgfsetfillcolor{currentfill}%
\pgfsetfillopacity{0.500000}%
\pgfsetlinewidth{0.250937pt}%
\definecolor{currentstroke}{rgb}{0.000000,0.000000,0.000000}%
\pgfsetstrokecolor{currentstroke}%
\pgfsetstrokeopacity{0.500000}%
\pgfsetdash{}{0pt}%
\pgfsys@defobject{currentmarker}{\pgfqpoint{-0.025000in}{-0.025000in}}{\pgfqpoint{0.025000in}{0.025000in}}{%
\pgfpathmoveto{\pgfqpoint{0.000000in}{-0.025000in}}%
\pgfpathcurveto{\pgfqpoint{0.006630in}{-0.025000in}}{\pgfqpoint{0.012989in}{-0.022366in}}{\pgfqpoint{0.017678in}{-0.017678in}}%
\pgfpathcurveto{\pgfqpoint{0.022366in}{-0.012989in}}{\pgfqpoint{0.025000in}{-0.006630in}}{\pgfqpoint{0.025000in}{0.000000in}}%
\pgfpathcurveto{\pgfqpoint{0.025000in}{0.006630in}}{\pgfqpoint{0.022366in}{0.012989in}}{\pgfqpoint{0.017678in}{0.017678in}}%
\pgfpathcurveto{\pgfqpoint{0.012989in}{0.022366in}}{\pgfqpoint{0.006630in}{0.025000in}}{\pgfqpoint{0.000000in}{0.025000in}}%
\pgfpathcurveto{\pgfqpoint{-0.006630in}{0.025000in}}{\pgfqpoint{-0.012989in}{0.022366in}}{\pgfqpoint{-0.017678in}{0.017678in}}%
\pgfpathcurveto{\pgfqpoint{-0.022366in}{0.012989in}}{\pgfqpoint{-0.025000in}{0.006630in}}{\pgfqpoint{-0.025000in}{0.000000in}}%
\pgfpathcurveto{\pgfqpoint{-0.025000in}{-0.006630in}}{\pgfqpoint{-0.022366in}{-0.012989in}}{\pgfqpoint{-0.017678in}{-0.017678in}}%
\pgfpathcurveto{\pgfqpoint{-0.012989in}{-0.022366in}}{\pgfqpoint{-0.006630in}{-0.025000in}}{\pgfqpoint{0.000000in}{-0.025000in}}%
\pgfpathlineto{\pgfqpoint{0.000000in}{-0.025000in}}%
\pgfpathclose%
\pgfusepath{stroke,fill}%
}%
\begin{pgfscope}%
\pgfsys@transformshift{0.465142in}{4.541082in}%
\pgfsys@useobject{currentmarker}{}%
\end{pgfscope}%
\end{pgfscope}%
\begin{pgfscope}%
\pgfpathrectangle{\pgfqpoint{0.100000in}{2.493063in}}{\pgfqpoint{5.037500in}{3.427208in}}%
\pgfusepath{clip}%
\pgfsetrectcap%
\pgfsetroundjoin%
\pgfsetlinewidth{1.505625pt}%
\definecolor{currentstroke}{rgb}{0.000000,0.000000,1.000000}%
\pgfsetstrokecolor{currentstroke}%
\pgfsetstrokeopacity{0.500000}%
\pgfsetdash{}{0pt}%
\pgfpathmoveto{\pgfqpoint{0.375033in}{4.647443in}}%
\pgfusepath{stroke}%
\end{pgfscope}%
\begin{pgfscope}%
\pgfpathrectangle{\pgfqpoint{0.100000in}{2.493063in}}{\pgfqpoint{5.037500in}{3.427208in}}%
\pgfusepath{clip}%
\pgfsetbuttcap%
\pgfsetroundjoin%
\definecolor{currentfill}{rgb}{0.000000,0.000000,1.000000}%
\pgfsetfillcolor{currentfill}%
\pgfsetfillopacity{0.500000}%
\pgfsetlinewidth{0.250937pt}%
\definecolor{currentstroke}{rgb}{0.000000,0.000000,0.000000}%
\pgfsetstrokecolor{currentstroke}%
\pgfsetstrokeopacity{0.500000}%
\pgfsetdash{}{0pt}%
\pgfsys@defobject{currentmarker}{\pgfqpoint{-0.025000in}{-0.025000in}}{\pgfqpoint{0.025000in}{0.025000in}}{%
\pgfpathmoveto{\pgfqpoint{0.000000in}{-0.025000in}}%
\pgfpathcurveto{\pgfqpoint{0.006630in}{-0.025000in}}{\pgfqpoint{0.012989in}{-0.022366in}}{\pgfqpoint{0.017678in}{-0.017678in}}%
\pgfpathcurveto{\pgfqpoint{0.022366in}{-0.012989in}}{\pgfqpoint{0.025000in}{-0.006630in}}{\pgfqpoint{0.025000in}{0.000000in}}%
\pgfpathcurveto{\pgfqpoint{0.025000in}{0.006630in}}{\pgfqpoint{0.022366in}{0.012989in}}{\pgfqpoint{0.017678in}{0.017678in}}%
\pgfpathcurveto{\pgfqpoint{0.012989in}{0.022366in}}{\pgfqpoint{0.006630in}{0.025000in}}{\pgfqpoint{0.000000in}{0.025000in}}%
\pgfpathcurveto{\pgfqpoint{-0.006630in}{0.025000in}}{\pgfqpoint{-0.012989in}{0.022366in}}{\pgfqpoint{-0.017678in}{0.017678in}}%
\pgfpathcurveto{\pgfqpoint{-0.022366in}{0.012989in}}{\pgfqpoint{-0.025000in}{0.006630in}}{\pgfqpoint{-0.025000in}{0.000000in}}%
\pgfpathcurveto{\pgfqpoint{-0.025000in}{-0.006630in}}{\pgfqpoint{-0.022366in}{-0.012989in}}{\pgfqpoint{-0.017678in}{-0.017678in}}%
\pgfpathcurveto{\pgfqpoint{-0.012989in}{-0.022366in}}{\pgfqpoint{-0.006630in}{-0.025000in}}{\pgfqpoint{0.000000in}{-0.025000in}}%
\pgfpathlineto{\pgfqpoint{0.000000in}{-0.025000in}}%
\pgfpathclose%
\pgfusepath{stroke,fill}%
}%
\begin{pgfscope}%
\pgfsys@transformshift{0.375033in}{4.647443in}%
\pgfsys@useobject{currentmarker}{}%
\end{pgfscope}%
\end{pgfscope}%
\begin{pgfscope}%
\pgfpathrectangle{\pgfqpoint{0.100000in}{2.493063in}}{\pgfqpoint{5.037500in}{3.427208in}}%
\pgfusepath{clip}%
\pgfsetrectcap%
\pgfsetroundjoin%
\pgfsetlinewidth{1.505625pt}%
\definecolor{currentstroke}{rgb}{0.000000,0.000000,1.000000}%
\pgfsetstrokecolor{currentstroke}%
\pgfsetstrokeopacity{0.500000}%
\pgfsetdash{}{0pt}%
\pgfpathmoveto{\pgfqpoint{0.520176in}{4.112959in}}%
\pgfusepath{stroke}%
\end{pgfscope}%
\begin{pgfscope}%
\pgfpathrectangle{\pgfqpoint{0.100000in}{2.493063in}}{\pgfqpoint{5.037500in}{3.427208in}}%
\pgfusepath{clip}%
\pgfsetbuttcap%
\pgfsetroundjoin%
\definecolor{currentfill}{rgb}{0.000000,0.000000,1.000000}%
\pgfsetfillcolor{currentfill}%
\pgfsetfillopacity{0.500000}%
\pgfsetlinewidth{0.250937pt}%
\definecolor{currentstroke}{rgb}{0.000000,0.000000,0.000000}%
\pgfsetstrokecolor{currentstroke}%
\pgfsetstrokeopacity{0.500000}%
\pgfsetdash{}{0pt}%
\pgfsys@defobject{currentmarker}{\pgfqpoint{-0.030556in}{-0.030556in}}{\pgfqpoint{0.030556in}{0.030556in}}{%
\pgfpathmoveto{\pgfqpoint{0.000000in}{-0.030556in}}%
\pgfpathcurveto{\pgfqpoint{0.008103in}{-0.030556in}}{\pgfqpoint{0.015876in}{-0.027336in}}{\pgfqpoint{0.021606in}{-0.021606in}}%
\pgfpathcurveto{\pgfqpoint{0.027336in}{-0.015876in}}{\pgfqpoint{0.030556in}{-0.008103in}}{\pgfqpoint{0.030556in}{0.000000in}}%
\pgfpathcurveto{\pgfqpoint{0.030556in}{0.008103in}}{\pgfqpoint{0.027336in}{0.015876in}}{\pgfqpoint{0.021606in}{0.021606in}}%
\pgfpathcurveto{\pgfqpoint{0.015876in}{0.027336in}}{\pgfqpoint{0.008103in}{0.030556in}}{\pgfqpoint{0.000000in}{0.030556in}}%
\pgfpathcurveto{\pgfqpoint{-0.008103in}{0.030556in}}{\pgfqpoint{-0.015876in}{0.027336in}}{\pgfqpoint{-0.021606in}{0.021606in}}%
\pgfpathcurveto{\pgfqpoint{-0.027336in}{0.015876in}}{\pgfqpoint{-0.030556in}{0.008103in}}{\pgfqpoint{-0.030556in}{0.000000in}}%
\pgfpathcurveto{\pgfqpoint{-0.030556in}{-0.008103in}}{\pgfqpoint{-0.027336in}{-0.015876in}}{\pgfqpoint{-0.021606in}{-0.021606in}}%
\pgfpathcurveto{\pgfqpoint{-0.015876in}{-0.027336in}}{\pgfqpoint{-0.008103in}{-0.030556in}}{\pgfqpoint{0.000000in}{-0.030556in}}%
\pgfpathlineto{\pgfqpoint{0.000000in}{-0.030556in}}%
\pgfpathclose%
\pgfusepath{stroke,fill}%
}%
\begin{pgfscope}%
\pgfsys@transformshift{0.520176in}{4.112959in}%
\pgfsys@useobject{currentmarker}{}%
\end{pgfscope}%
\end{pgfscope}%
\begin{pgfscope}%
\pgfpathrectangle{\pgfqpoint{0.100000in}{2.493063in}}{\pgfqpoint{5.037500in}{3.427208in}}%
\pgfusepath{clip}%
\pgfsetrectcap%
\pgfsetroundjoin%
\pgfsetlinewidth{1.505625pt}%
\definecolor{currentstroke}{rgb}{0.000000,0.000000,1.000000}%
\pgfsetstrokecolor{currentstroke}%
\pgfsetstrokeopacity{0.500000}%
\pgfsetdash{}{0pt}%
\pgfpathmoveto{\pgfqpoint{0.444132in}{4.902021in}}%
\pgfusepath{stroke}%
\end{pgfscope}%
\begin{pgfscope}%
\pgfpathrectangle{\pgfqpoint{0.100000in}{2.493063in}}{\pgfqpoint{5.037500in}{3.427208in}}%
\pgfusepath{clip}%
\pgfsetbuttcap%
\pgfsetroundjoin%
\definecolor{currentfill}{rgb}{0.000000,0.000000,1.000000}%
\pgfsetfillcolor{currentfill}%
\pgfsetfillopacity{0.500000}%
\pgfsetlinewidth{0.250937pt}%
\definecolor{currentstroke}{rgb}{0.000000,0.000000,0.000000}%
\pgfsetstrokecolor{currentstroke}%
\pgfsetstrokeopacity{0.500000}%
\pgfsetdash{}{0pt}%
\pgfsys@defobject{currentmarker}{\pgfqpoint{-0.027778in}{-0.027778in}}{\pgfqpoint{0.027778in}{0.027778in}}{%
\pgfpathmoveto{\pgfqpoint{0.000000in}{-0.027778in}}%
\pgfpathcurveto{\pgfqpoint{0.007367in}{-0.027778in}}{\pgfqpoint{0.014433in}{-0.024851in}}{\pgfqpoint{0.019642in}{-0.019642in}}%
\pgfpathcurveto{\pgfqpoint{0.024851in}{-0.014433in}}{\pgfqpoint{0.027778in}{-0.007367in}}{\pgfqpoint{0.027778in}{0.000000in}}%
\pgfpathcurveto{\pgfqpoint{0.027778in}{0.007367in}}{\pgfqpoint{0.024851in}{0.014433in}}{\pgfqpoint{0.019642in}{0.019642in}}%
\pgfpathcurveto{\pgfqpoint{0.014433in}{0.024851in}}{\pgfqpoint{0.007367in}{0.027778in}}{\pgfqpoint{0.000000in}{0.027778in}}%
\pgfpathcurveto{\pgfqpoint{-0.007367in}{0.027778in}}{\pgfqpoint{-0.014433in}{0.024851in}}{\pgfqpoint{-0.019642in}{0.019642in}}%
\pgfpathcurveto{\pgfqpoint{-0.024851in}{0.014433in}}{\pgfqpoint{-0.027778in}{0.007367in}}{\pgfqpoint{-0.027778in}{0.000000in}}%
\pgfpathcurveto{\pgfqpoint{-0.027778in}{-0.007367in}}{\pgfqpoint{-0.024851in}{-0.014433in}}{\pgfqpoint{-0.019642in}{-0.019642in}}%
\pgfpathcurveto{\pgfqpoint{-0.014433in}{-0.024851in}}{\pgfqpoint{-0.007367in}{-0.027778in}}{\pgfqpoint{0.000000in}{-0.027778in}}%
\pgfpathlineto{\pgfqpoint{0.000000in}{-0.027778in}}%
\pgfpathclose%
\pgfusepath{stroke,fill}%
}%
\begin{pgfscope}%
\pgfsys@transformshift{0.444132in}{4.902021in}%
\pgfsys@useobject{currentmarker}{}%
\end{pgfscope}%
\end{pgfscope}%
\begin{pgfscope}%
\pgfpathrectangle{\pgfqpoint{0.100000in}{2.493063in}}{\pgfqpoint{5.037500in}{3.427208in}}%
\pgfusepath{clip}%
\pgfsetrectcap%
\pgfsetroundjoin%
\pgfsetlinewidth{1.505625pt}%
\definecolor{currentstroke}{rgb}{0.000000,0.000000,1.000000}%
\pgfsetstrokecolor{currentstroke}%
\pgfsetstrokeopacity{0.500000}%
\pgfsetdash{}{0pt}%
\pgfpathmoveto{\pgfqpoint{0.678250in}{4.042355in}}%
\pgfusepath{stroke}%
\end{pgfscope}%
\begin{pgfscope}%
\pgfpathrectangle{\pgfqpoint{0.100000in}{2.493063in}}{\pgfqpoint{5.037500in}{3.427208in}}%
\pgfusepath{clip}%
\pgfsetbuttcap%
\pgfsetroundjoin%
\definecolor{currentfill}{rgb}{0.000000,0.000000,1.000000}%
\pgfsetfillcolor{currentfill}%
\pgfsetfillopacity{0.500000}%
\pgfsetlinewidth{0.250937pt}%
\definecolor{currentstroke}{rgb}{0.000000,0.000000,0.000000}%
\pgfsetstrokecolor{currentstroke}%
\pgfsetstrokeopacity{0.500000}%
\pgfsetdash{}{0pt}%
\pgfsys@defobject{currentmarker}{\pgfqpoint{-0.036111in}{-0.036111in}}{\pgfqpoint{0.036111in}{0.036111in}}{%
\pgfpathmoveto{\pgfqpoint{0.000000in}{-0.036111in}}%
\pgfpathcurveto{\pgfqpoint{0.009577in}{-0.036111in}}{\pgfqpoint{0.018763in}{-0.032306in}}{\pgfqpoint{0.025534in}{-0.025534in}}%
\pgfpathcurveto{\pgfqpoint{0.032306in}{-0.018763in}}{\pgfqpoint{0.036111in}{-0.009577in}}{\pgfqpoint{0.036111in}{0.000000in}}%
\pgfpathcurveto{\pgfqpoint{0.036111in}{0.009577in}}{\pgfqpoint{0.032306in}{0.018763in}}{\pgfqpoint{0.025534in}{0.025534in}}%
\pgfpathcurveto{\pgfqpoint{0.018763in}{0.032306in}}{\pgfqpoint{0.009577in}{0.036111in}}{\pgfqpoint{0.000000in}{0.036111in}}%
\pgfpathcurveto{\pgfqpoint{-0.009577in}{0.036111in}}{\pgfqpoint{-0.018763in}{0.032306in}}{\pgfqpoint{-0.025534in}{0.025534in}}%
\pgfpathcurveto{\pgfqpoint{-0.032306in}{0.018763in}}{\pgfqpoint{-0.036111in}{0.009577in}}{\pgfqpoint{-0.036111in}{0.000000in}}%
\pgfpathcurveto{\pgfqpoint{-0.036111in}{-0.009577in}}{\pgfqpoint{-0.032306in}{-0.018763in}}{\pgfqpoint{-0.025534in}{-0.025534in}}%
\pgfpathcurveto{\pgfqpoint{-0.018763in}{-0.032306in}}{\pgfqpoint{-0.009577in}{-0.036111in}}{\pgfqpoint{0.000000in}{-0.036111in}}%
\pgfpathlineto{\pgfqpoint{0.000000in}{-0.036111in}}%
\pgfpathclose%
\pgfusepath{stroke,fill}%
}%
\begin{pgfscope}%
\pgfsys@transformshift{0.678250in}{4.042355in}%
\pgfsys@useobject{currentmarker}{}%
\end{pgfscope}%
\end{pgfscope}%
\begin{pgfscope}%
\pgfpathrectangle{\pgfqpoint{0.100000in}{2.493063in}}{\pgfqpoint{5.037500in}{3.427208in}}%
\pgfusepath{clip}%
\pgfsetrectcap%
\pgfsetroundjoin%
\pgfsetlinewidth{1.505625pt}%
\definecolor{currentstroke}{rgb}{0.000000,0.000000,1.000000}%
\pgfsetstrokecolor{currentstroke}%
\pgfsetstrokeopacity{0.500000}%
\pgfsetdash{}{0pt}%
\pgfpathmoveto{\pgfqpoint{0.452839in}{4.657966in}}%
\pgfusepath{stroke}%
\end{pgfscope}%
\begin{pgfscope}%
\pgfpathrectangle{\pgfqpoint{0.100000in}{2.493063in}}{\pgfqpoint{5.037500in}{3.427208in}}%
\pgfusepath{clip}%
\pgfsetbuttcap%
\pgfsetroundjoin%
\definecolor{currentfill}{rgb}{0.000000,0.000000,1.000000}%
\pgfsetfillcolor{currentfill}%
\pgfsetfillopacity{0.500000}%
\pgfsetlinewidth{0.250937pt}%
\definecolor{currentstroke}{rgb}{0.000000,0.000000,0.000000}%
\pgfsetstrokecolor{currentstroke}%
\pgfsetstrokeopacity{0.500000}%
\pgfsetdash{}{0pt}%
\pgfsys@defobject{currentmarker}{\pgfqpoint{-0.030556in}{-0.030556in}}{\pgfqpoint{0.030556in}{0.030556in}}{%
\pgfpathmoveto{\pgfqpoint{0.000000in}{-0.030556in}}%
\pgfpathcurveto{\pgfqpoint{0.008103in}{-0.030556in}}{\pgfqpoint{0.015876in}{-0.027336in}}{\pgfqpoint{0.021606in}{-0.021606in}}%
\pgfpathcurveto{\pgfqpoint{0.027336in}{-0.015876in}}{\pgfqpoint{0.030556in}{-0.008103in}}{\pgfqpoint{0.030556in}{0.000000in}}%
\pgfpathcurveto{\pgfqpoint{0.030556in}{0.008103in}}{\pgfqpoint{0.027336in}{0.015876in}}{\pgfqpoint{0.021606in}{0.021606in}}%
\pgfpathcurveto{\pgfqpoint{0.015876in}{0.027336in}}{\pgfqpoint{0.008103in}{0.030556in}}{\pgfqpoint{0.000000in}{0.030556in}}%
\pgfpathcurveto{\pgfqpoint{-0.008103in}{0.030556in}}{\pgfqpoint{-0.015876in}{0.027336in}}{\pgfqpoint{-0.021606in}{0.021606in}}%
\pgfpathcurveto{\pgfqpoint{-0.027336in}{0.015876in}}{\pgfqpoint{-0.030556in}{0.008103in}}{\pgfqpoint{-0.030556in}{0.000000in}}%
\pgfpathcurveto{\pgfqpoint{-0.030556in}{-0.008103in}}{\pgfqpoint{-0.027336in}{-0.015876in}}{\pgfqpoint{-0.021606in}{-0.021606in}}%
\pgfpathcurveto{\pgfqpoint{-0.015876in}{-0.027336in}}{\pgfqpoint{-0.008103in}{-0.030556in}}{\pgfqpoint{0.000000in}{-0.030556in}}%
\pgfpathlineto{\pgfqpoint{0.000000in}{-0.030556in}}%
\pgfpathclose%
\pgfusepath{stroke,fill}%
}%
\begin{pgfscope}%
\pgfsys@transformshift{0.452839in}{4.657966in}%
\pgfsys@useobject{currentmarker}{}%
\end{pgfscope}%
\end{pgfscope}%
\begin{pgfscope}%
\pgfpathrectangle{\pgfqpoint{0.100000in}{2.493063in}}{\pgfqpoint{5.037500in}{3.427208in}}%
\pgfusepath{clip}%
\pgfsetrectcap%
\pgfsetroundjoin%
\pgfsetlinewidth{1.505625pt}%
\definecolor{currentstroke}{rgb}{0.000000,0.000000,1.000000}%
\pgfsetstrokecolor{currentstroke}%
\pgfsetstrokeopacity{0.500000}%
\pgfsetdash{}{0pt}%
\pgfpathmoveto{\pgfqpoint{0.375477in}{4.451832in}}%
\pgfusepath{stroke}%
\end{pgfscope}%
\begin{pgfscope}%
\pgfpathrectangle{\pgfqpoint{0.100000in}{2.493063in}}{\pgfqpoint{5.037500in}{3.427208in}}%
\pgfusepath{clip}%
\pgfsetbuttcap%
\pgfsetroundjoin%
\definecolor{currentfill}{rgb}{0.000000,0.000000,1.000000}%
\pgfsetfillcolor{currentfill}%
\pgfsetfillopacity{0.500000}%
\pgfsetlinewidth{0.250937pt}%
\definecolor{currentstroke}{rgb}{0.000000,0.000000,0.000000}%
\pgfsetstrokecolor{currentstroke}%
\pgfsetstrokeopacity{0.500000}%
\pgfsetdash{}{0pt}%
\pgfsys@defobject{currentmarker}{\pgfqpoint{-0.019444in}{-0.019444in}}{\pgfqpoint{0.019444in}{0.019444in}}{%
\pgfpathmoveto{\pgfqpoint{0.000000in}{-0.019444in}}%
\pgfpathcurveto{\pgfqpoint{0.005157in}{-0.019444in}}{\pgfqpoint{0.010103in}{-0.017396in}}{\pgfqpoint{0.013749in}{-0.013749in}}%
\pgfpathcurveto{\pgfqpoint{0.017396in}{-0.010103in}}{\pgfqpoint{0.019444in}{-0.005157in}}{\pgfqpoint{0.019444in}{0.000000in}}%
\pgfpathcurveto{\pgfqpoint{0.019444in}{0.005157in}}{\pgfqpoint{0.017396in}{0.010103in}}{\pgfqpoint{0.013749in}{0.013749in}}%
\pgfpathcurveto{\pgfqpoint{0.010103in}{0.017396in}}{\pgfqpoint{0.005157in}{0.019444in}}{\pgfqpoint{0.000000in}{0.019444in}}%
\pgfpathcurveto{\pgfqpoint{-0.005157in}{0.019444in}}{\pgfqpoint{-0.010103in}{0.017396in}}{\pgfqpoint{-0.013749in}{0.013749in}}%
\pgfpathcurveto{\pgfqpoint{-0.017396in}{0.010103in}}{\pgfqpoint{-0.019444in}{0.005157in}}{\pgfqpoint{-0.019444in}{0.000000in}}%
\pgfpathcurveto{\pgfqpoint{-0.019444in}{-0.005157in}}{\pgfqpoint{-0.017396in}{-0.010103in}}{\pgfqpoint{-0.013749in}{-0.013749in}}%
\pgfpathcurveto{\pgfqpoint{-0.010103in}{-0.017396in}}{\pgfqpoint{-0.005157in}{-0.019444in}}{\pgfqpoint{0.000000in}{-0.019444in}}%
\pgfpathlineto{\pgfqpoint{0.000000in}{-0.019444in}}%
\pgfpathclose%
\pgfusepath{stroke,fill}%
}%
\begin{pgfscope}%
\pgfsys@transformshift{0.375477in}{4.451832in}%
\pgfsys@useobject{currentmarker}{}%
\end{pgfscope}%
\end{pgfscope}%
\begin{pgfscope}%
\pgfpathrectangle{\pgfqpoint{0.100000in}{2.493063in}}{\pgfqpoint{5.037500in}{3.427208in}}%
\pgfusepath{clip}%
\pgfsetrectcap%
\pgfsetroundjoin%
\pgfsetlinewidth{1.505625pt}%
\definecolor{currentstroke}{rgb}{0.000000,0.000000,1.000000}%
\pgfsetstrokecolor{currentstroke}%
\pgfsetstrokeopacity{0.500000}%
\pgfsetdash{}{0pt}%
\pgfpathmoveto{\pgfqpoint{0.665426in}{3.898137in}}%
\pgfusepath{stroke}%
\end{pgfscope}%
\begin{pgfscope}%
\pgfpathrectangle{\pgfqpoint{0.100000in}{2.493063in}}{\pgfqpoint{5.037500in}{3.427208in}}%
\pgfusepath{clip}%
\pgfsetbuttcap%
\pgfsetroundjoin%
\definecolor{currentfill}{rgb}{0.000000,0.000000,1.000000}%
\pgfsetfillcolor{currentfill}%
\pgfsetfillopacity{0.500000}%
\pgfsetlinewidth{0.250937pt}%
\definecolor{currentstroke}{rgb}{0.000000,0.000000,0.000000}%
\pgfsetstrokecolor{currentstroke}%
\pgfsetstrokeopacity{0.500000}%
\pgfsetdash{}{0pt}%
\pgfsys@defobject{currentmarker}{\pgfqpoint{-0.030556in}{-0.030556in}}{\pgfqpoint{0.030556in}{0.030556in}}{%
\pgfpathmoveto{\pgfqpoint{0.000000in}{-0.030556in}}%
\pgfpathcurveto{\pgfqpoint{0.008103in}{-0.030556in}}{\pgfqpoint{0.015876in}{-0.027336in}}{\pgfqpoint{0.021606in}{-0.021606in}}%
\pgfpathcurveto{\pgfqpoint{0.027336in}{-0.015876in}}{\pgfqpoint{0.030556in}{-0.008103in}}{\pgfqpoint{0.030556in}{0.000000in}}%
\pgfpathcurveto{\pgfqpoint{0.030556in}{0.008103in}}{\pgfqpoint{0.027336in}{0.015876in}}{\pgfqpoint{0.021606in}{0.021606in}}%
\pgfpathcurveto{\pgfqpoint{0.015876in}{0.027336in}}{\pgfqpoint{0.008103in}{0.030556in}}{\pgfqpoint{0.000000in}{0.030556in}}%
\pgfpathcurveto{\pgfqpoint{-0.008103in}{0.030556in}}{\pgfqpoint{-0.015876in}{0.027336in}}{\pgfqpoint{-0.021606in}{0.021606in}}%
\pgfpathcurveto{\pgfqpoint{-0.027336in}{0.015876in}}{\pgfqpoint{-0.030556in}{0.008103in}}{\pgfqpoint{-0.030556in}{0.000000in}}%
\pgfpathcurveto{\pgfqpoint{-0.030556in}{-0.008103in}}{\pgfqpoint{-0.027336in}{-0.015876in}}{\pgfqpoint{-0.021606in}{-0.021606in}}%
\pgfpathcurveto{\pgfqpoint{-0.015876in}{-0.027336in}}{\pgfqpoint{-0.008103in}{-0.030556in}}{\pgfqpoint{0.000000in}{-0.030556in}}%
\pgfpathlineto{\pgfqpoint{0.000000in}{-0.030556in}}%
\pgfpathclose%
\pgfusepath{stroke,fill}%
}%
\begin{pgfscope}%
\pgfsys@transformshift{0.665426in}{3.898137in}%
\pgfsys@useobject{currentmarker}{}%
\end{pgfscope}%
\end{pgfscope}%
\begin{pgfscope}%
\pgfpathrectangle{\pgfqpoint{0.100000in}{2.493063in}}{\pgfqpoint{5.037500in}{3.427208in}}%
\pgfusepath{clip}%
\pgfsetrectcap%
\pgfsetroundjoin%
\pgfsetlinewidth{1.505625pt}%
\definecolor{currentstroke}{rgb}{0.000000,0.000000,1.000000}%
\pgfsetstrokecolor{currentstroke}%
\pgfsetstrokeopacity{0.500000}%
\pgfsetdash{}{0pt}%
\pgfpathmoveto{\pgfqpoint{0.345769in}{4.594080in}}%
\pgfusepath{stroke}%
\end{pgfscope}%
\begin{pgfscope}%
\pgfpathrectangle{\pgfqpoint{0.100000in}{2.493063in}}{\pgfqpoint{5.037500in}{3.427208in}}%
\pgfusepath{clip}%
\pgfsetbuttcap%
\pgfsetroundjoin%
\definecolor{currentfill}{rgb}{0.000000,0.000000,1.000000}%
\pgfsetfillcolor{currentfill}%
\pgfsetfillopacity{0.500000}%
\pgfsetlinewidth{0.250937pt}%
\definecolor{currentstroke}{rgb}{0.000000,0.000000,0.000000}%
\pgfsetstrokecolor{currentstroke}%
\pgfsetstrokeopacity{0.500000}%
\pgfsetdash{}{0pt}%
\pgfsys@defobject{currentmarker}{\pgfqpoint{-0.041667in}{-0.041667in}}{\pgfqpoint{0.041667in}{0.041667in}}{%
\pgfpathmoveto{\pgfqpoint{0.000000in}{-0.041667in}}%
\pgfpathcurveto{\pgfqpoint{0.011050in}{-0.041667in}}{\pgfqpoint{0.021649in}{-0.037276in}}{\pgfqpoint{0.029463in}{-0.029463in}}%
\pgfpathcurveto{\pgfqpoint{0.037276in}{-0.021649in}}{\pgfqpoint{0.041667in}{-0.011050in}}{\pgfqpoint{0.041667in}{0.000000in}}%
\pgfpathcurveto{\pgfqpoint{0.041667in}{0.011050in}}{\pgfqpoint{0.037276in}{0.021649in}}{\pgfqpoint{0.029463in}{0.029463in}}%
\pgfpathcurveto{\pgfqpoint{0.021649in}{0.037276in}}{\pgfqpoint{0.011050in}{0.041667in}}{\pgfqpoint{0.000000in}{0.041667in}}%
\pgfpathcurveto{\pgfqpoint{-0.011050in}{0.041667in}}{\pgfqpoint{-0.021649in}{0.037276in}}{\pgfqpoint{-0.029463in}{0.029463in}}%
\pgfpathcurveto{\pgfqpoint{-0.037276in}{0.021649in}}{\pgfqpoint{-0.041667in}{0.011050in}}{\pgfqpoint{-0.041667in}{0.000000in}}%
\pgfpathcurveto{\pgfqpoint{-0.041667in}{-0.011050in}}{\pgfqpoint{-0.037276in}{-0.021649in}}{\pgfqpoint{-0.029463in}{-0.029463in}}%
\pgfpathcurveto{\pgfqpoint{-0.021649in}{-0.037276in}}{\pgfqpoint{-0.011050in}{-0.041667in}}{\pgfqpoint{0.000000in}{-0.041667in}}%
\pgfpathlineto{\pgfqpoint{0.000000in}{-0.041667in}}%
\pgfpathclose%
\pgfusepath{stroke,fill}%
}%
\begin{pgfscope}%
\pgfsys@transformshift{0.345769in}{4.594080in}%
\pgfsys@useobject{currentmarker}{}%
\end{pgfscope}%
\end{pgfscope}%
\begin{pgfscope}%
\pgfpathrectangle{\pgfqpoint{0.100000in}{2.493063in}}{\pgfqpoint{5.037500in}{3.427208in}}%
\pgfusepath{clip}%
\pgfsetrectcap%
\pgfsetroundjoin%
\pgfsetlinewidth{1.505625pt}%
\definecolor{currentstroke}{rgb}{0.000000,0.000000,1.000000}%
\pgfsetstrokecolor{currentstroke}%
\pgfsetstrokeopacity{0.500000}%
\pgfsetdash{}{0pt}%
\pgfpathmoveto{\pgfqpoint{0.376878in}{4.531092in}}%
\pgfusepath{stroke}%
\end{pgfscope}%
\begin{pgfscope}%
\pgfpathrectangle{\pgfqpoint{0.100000in}{2.493063in}}{\pgfqpoint{5.037500in}{3.427208in}}%
\pgfusepath{clip}%
\pgfsetbuttcap%
\pgfsetroundjoin%
\definecolor{currentfill}{rgb}{0.000000,0.000000,1.000000}%
\pgfsetfillcolor{currentfill}%
\pgfsetfillopacity{0.500000}%
\pgfsetlinewidth{0.250937pt}%
\definecolor{currentstroke}{rgb}{0.000000,0.000000,0.000000}%
\pgfsetstrokecolor{currentstroke}%
\pgfsetstrokeopacity{0.500000}%
\pgfsetdash{}{0pt}%
\pgfsys@defobject{currentmarker}{\pgfqpoint{-0.038889in}{-0.038889in}}{\pgfqpoint{0.038889in}{0.038889in}}{%
\pgfpathmoveto{\pgfqpoint{0.000000in}{-0.038889in}}%
\pgfpathcurveto{\pgfqpoint{0.010313in}{-0.038889in}}{\pgfqpoint{0.020206in}{-0.034791in}}{\pgfqpoint{0.027499in}{-0.027499in}}%
\pgfpathcurveto{\pgfqpoint{0.034791in}{-0.020206in}}{\pgfqpoint{0.038889in}{-0.010313in}}{\pgfqpoint{0.038889in}{0.000000in}}%
\pgfpathcurveto{\pgfqpoint{0.038889in}{0.010313in}}{\pgfqpoint{0.034791in}{0.020206in}}{\pgfqpoint{0.027499in}{0.027499in}}%
\pgfpathcurveto{\pgfqpoint{0.020206in}{0.034791in}}{\pgfqpoint{0.010313in}{0.038889in}}{\pgfqpoint{0.000000in}{0.038889in}}%
\pgfpathcurveto{\pgfqpoint{-0.010313in}{0.038889in}}{\pgfqpoint{-0.020206in}{0.034791in}}{\pgfqpoint{-0.027499in}{0.027499in}}%
\pgfpathcurveto{\pgfqpoint{-0.034791in}{0.020206in}}{\pgfqpoint{-0.038889in}{0.010313in}}{\pgfqpoint{-0.038889in}{0.000000in}}%
\pgfpathcurveto{\pgfqpoint{-0.038889in}{-0.010313in}}{\pgfqpoint{-0.034791in}{-0.020206in}}{\pgfqpoint{-0.027499in}{-0.027499in}}%
\pgfpathcurveto{\pgfqpoint{-0.020206in}{-0.034791in}}{\pgfqpoint{-0.010313in}{-0.038889in}}{\pgfqpoint{0.000000in}{-0.038889in}}%
\pgfpathlineto{\pgfqpoint{0.000000in}{-0.038889in}}%
\pgfpathclose%
\pgfusepath{stroke,fill}%
}%
\begin{pgfscope}%
\pgfsys@transformshift{0.376878in}{4.531092in}%
\pgfsys@useobject{currentmarker}{}%
\end{pgfscope}%
\end{pgfscope}%
\begin{pgfscope}%
\pgfpathrectangle{\pgfqpoint{0.100000in}{2.493063in}}{\pgfqpoint{5.037500in}{3.427208in}}%
\pgfusepath{clip}%
\pgfsetrectcap%
\pgfsetroundjoin%
\pgfsetlinewidth{1.505625pt}%
\definecolor{currentstroke}{rgb}{0.000000,0.000000,1.000000}%
\pgfsetstrokecolor{currentstroke}%
\pgfsetstrokeopacity{0.500000}%
\pgfsetdash{}{0pt}%
\pgfpathmoveto{\pgfqpoint{0.418989in}{4.271586in}}%
\pgfusepath{stroke}%
\end{pgfscope}%
\begin{pgfscope}%
\pgfpathrectangle{\pgfqpoint{0.100000in}{2.493063in}}{\pgfqpoint{5.037500in}{3.427208in}}%
\pgfusepath{clip}%
\pgfsetbuttcap%
\pgfsetroundjoin%
\definecolor{currentfill}{rgb}{0.000000,0.000000,1.000000}%
\pgfsetfillcolor{currentfill}%
\pgfsetfillopacity{0.500000}%
\pgfsetlinewidth{0.250937pt}%
\definecolor{currentstroke}{rgb}{0.000000,0.000000,0.000000}%
\pgfsetstrokecolor{currentstroke}%
\pgfsetstrokeopacity{0.500000}%
\pgfsetdash{}{0pt}%
\pgfsys@defobject{currentmarker}{\pgfqpoint{-0.022222in}{-0.022222in}}{\pgfqpoint{0.022222in}{0.022222in}}{%
\pgfpathmoveto{\pgfqpoint{0.000000in}{-0.022222in}}%
\pgfpathcurveto{\pgfqpoint{0.005893in}{-0.022222in}}{\pgfqpoint{0.011546in}{-0.019881in}}{\pgfqpoint{0.015713in}{-0.015713in}}%
\pgfpathcurveto{\pgfqpoint{0.019881in}{-0.011546in}}{\pgfqpoint{0.022222in}{-0.005893in}}{\pgfqpoint{0.022222in}{0.000000in}}%
\pgfpathcurveto{\pgfqpoint{0.022222in}{0.005893in}}{\pgfqpoint{0.019881in}{0.011546in}}{\pgfqpoint{0.015713in}{0.015713in}}%
\pgfpathcurveto{\pgfqpoint{0.011546in}{0.019881in}}{\pgfqpoint{0.005893in}{0.022222in}}{\pgfqpoint{0.000000in}{0.022222in}}%
\pgfpathcurveto{\pgfqpoint{-0.005893in}{0.022222in}}{\pgfqpoint{-0.011546in}{0.019881in}}{\pgfqpoint{-0.015713in}{0.015713in}}%
\pgfpathcurveto{\pgfqpoint{-0.019881in}{0.011546in}}{\pgfqpoint{-0.022222in}{0.005893in}}{\pgfqpoint{-0.022222in}{0.000000in}}%
\pgfpathcurveto{\pgfqpoint{-0.022222in}{-0.005893in}}{\pgfqpoint{-0.019881in}{-0.011546in}}{\pgfqpoint{-0.015713in}{-0.015713in}}%
\pgfpathcurveto{\pgfqpoint{-0.011546in}{-0.019881in}}{\pgfqpoint{-0.005893in}{-0.022222in}}{\pgfqpoint{0.000000in}{-0.022222in}}%
\pgfpathlineto{\pgfqpoint{0.000000in}{-0.022222in}}%
\pgfpathclose%
\pgfusepath{stroke,fill}%
}%
\begin{pgfscope}%
\pgfsys@transformshift{0.418989in}{4.271586in}%
\pgfsys@useobject{currentmarker}{}%
\end{pgfscope}%
\end{pgfscope}%
\begin{pgfscope}%
\pgfpathrectangle{\pgfqpoint{0.100000in}{2.493063in}}{\pgfqpoint{5.037500in}{3.427208in}}%
\pgfusepath{clip}%
\pgfsetrectcap%
\pgfsetroundjoin%
\pgfsetlinewidth{1.505625pt}%
\definecolor{currentstroke}{rgb}{0.000000,0.000000,1.000000}%
\pgfsetstrokecolor{currentstroke}%
\pgfsetstrokeopacity{0.500000}%
\pgfsetdash{}{0pt}%
\pgfpathmoveto{\pgfqpoint{0.352770in}{4.494777in}}%
\pgfusepath{stroke}%
\end{pgfscope}%
\begin{pgfscope}%
\pgfpathrectangle{\pgfqpoint{0.100000in}{2.493063in}}{\pgfqpoint{5.037500in}{3.427208in}}%
\pgfusepath{clip}%
\pgfsetbuttcap%
\pgfsetroundjoin%
\definecolor{currentfill}{rgb}{0.000000,0.000000,1.000000}%
\pgfsetfillcolor{currentfill}%
\pgfsetfillopacity{0.500000}%
\pgfsetlinewidth{0.250937pt}%
\definecolor{currentstroke}{rgb}{0.000000,0.000000,0.000000}%
\pgfsetstrokecolor{currentstroke}%
\pgfsetstrokeopacity{0.500000}%
\pgfsetdash{}{0pt}%
\pgfsys@defobject{currentmarker}{\pgfqpoint{-0.036111in}{-0.036111in}}{\pgfqpoint{0.036111in}{0.036111in}}{%
\pgfpathmoveto{\pgfqpoint{0.000000in}{-0.036111in}}%
\pgfpathcurveto{\pgfqpoint{0.009577in}{-0.036111in}}{\pgfqpoint{0.018763in}{-0.032306in}}{\pgfqpoint{0.025534in}{-0.025534in}}%
\pgfpathcurveto{\pgfqpoint{0.032306in}{-0.018763in}}{\pgfqpoint{0.036111in}{-0.009577in}}{\pgfqpoint{0.036111in}{0.000000in}}%
\pgfpathcurveto{\pgfqpoint{0.036111in}{0.009577in}}{\pgfqpoint{0.032306in}{0.018763in}}{\pgfqpoint{0.025534in}{0.025534in}}%
\pgfpathcurveto{\pgfqpoint{0.018763in}{0.032306in}}{\pgfqpoint{0.009577in}{0.036111in}}{\pgfqpoint{0.000000in}{0.036111in}}%
\pgfpathcurveto{\pgfqpoint{-0.009577in}{0.036111in}}{\pgfqpoint{-0.018763in}{0.032306in}}{\pgfqpoint{-0.025534in}{0.025534in}}%
\pgfpathcurveto{\pgfqpoint{-0.032306in}{0.018763in}}{\pgfqpoint{-0.036111in}{0.009577in}}{\pgfqpoint{-0.036111in}{0.000000in}}%
\pgfpathcurveto{\pgfqpoint{-0.036111in}{-0.009577in}}{\pgfqpoint{-0.032306in}{-0.018763in}}{\pgfqpoint{-0.025534in}{-0.025534in}}%
\pgfpathcurveto{\pgfqpoint{-0.018763in}{-0.032306in}}{\pgfqpoint{-0.009577in}{-0.036111in}}{\pgfqpoint{0.000000in}{-0.036111in}}%
\pgfpathlineto{\pgfqpoint{0.000000in}{-0.036111in}}%
\pgfpathclose%
\pgfusepath{stroke,fill}%
}%
\begin{pgfscope}%
\pgfsys@transformshift{0.352770in}{4.494777in}%
\pgfsys@useobject{currentmarker}{}%
\end{pgfscope}%
\end{pgfscope}%
\begin{pgfscope}%
\pgfpathrectangle{\pgfqpoint{0.100000in}{2.493063in}}{\pgfqpoint{5.037500in}{3.427208in}}%
\pgfusepath{clip}%
\pgfsetrectcap%
\pgfsetroundjoin%
\pgfsetlinewidth{1.505625pt}%
\definecolor{currentstroke}{rgb}{0.000000,0.000000,1.000000}%
\pgfsetstrokecolor{currentstroke}%
\pgfsetstrokeopacity{0.500000}%
\pgfsetdash{}{0pt}%
\pgfpathmoveto{\pgfqpoint{0.428730in}{4.229174in}}%
\pgfusepath{stroke}%
\end{pgfscope}%
\begin{pgfscope}%
\pgfpathrectangle{\pgfqpoint{0.100000in}{2.493063in}}{\pgfqpoint{5.037500in}{3.427208in}}%
\pgfusepath{clip}%
\pgfsetbuttcap%
\pgfsetroundjoin%
\definecolor{currentfill}{rgb}{0.000000,0.000000,1.000000}%
\pgfsetfillcolor{currentfill}%
\pgfsetfillopacity{0.500000}%
\pgfsetlinewidth{0.250937pt}%
\definecolor{currentstroke}{rgb}{0.000000,0.000000,0.000000}%
\pgfsetstrokecolor{currentstroke}%
\pgfsetstrokeopacity{0.500000}%
\pgfsetdash{}{0pt}%
\pgfsys@defobject{currentmarker}{\pgfqpoint{-0.019444in}{-0.019444in}}{\pgfqpoint{0.019444in}{0.019444in}}{%
\pgfpathmoveto{\pgfqpoint{0.000000in}{-0.019444in}}%
\pgfpathcurveto{\pgfqpoint{0.005157in}{-0.019444in}}{\pgfqpoint{0.010103in}{-0.017396in}}{\pgfqpoint{0.013749in}{-0.013749in}}%
\pgfpathcurveto{\pgfqpoint{0.017396in}{-0.010103in}}{\pgfqpoint{0.019444in}{-0.005157in}}{\pgfqpoint{0.019444in}{0.000000in}}%
\pgfpathcurveto{\pgfqpoint{0.019444in}{0.005157in}}{\pgfqpoint{0.017396in}{0.010103in}}{\pgfqpoint{0.013749in}{0.013749in}}%
\pgfpathcurveto{\pgfqpoint{0.010103in}{0.017396in}}{\pgfqpoint{0.005157in}{0.019444in}}{\pgfqpoint{0.000000in}{0.019444in}}%
\pgfpathcurveto{\pgfqpoint{-0.005157in}{0.019444in}}{\pgfqpoint{-0.010103in}{0.017396in}}{\pgfqpoint{-0.013749in}{0.013749in}}%
\pgfpathcurveto{\pgfqpoint{-0.017396in}{0.010103in}}{\pgfqpoint{-0.019444in}{0.005157in}}{\pgfqpoint{-0.019444in}{0.000000in}}%
\pgfpathcurveto{\pgfqpoint{-0.019444in}{-0.005157in}}{\pgfqpoint{-0.017396in}{-0.010103in}}{\pgfqpoint{-0.013749in}{-0.013749in}}%
\pgfpathcurveto{\pgfqpoint{-0.010103in}{-0.017396in}}{\pgfqpoint{-0.005157in}{-0.019444in}}{\pgfqpoint{0.000000in}{-0.019444in}}%
\pgfpathlineto{\pgfqpoint{0.000000in}{-0.019444in}}%
\pgfpathclose%
\pgfusepath{stroke,fill}%
}%
\begin{pgfscope}%
\pgfsys@transformshift{0.428730in}{4.229174in}%
\pgfsys@useobject{currentmarker}{}%
\end{pgfscope}%
\end{pgfscope}%
\begin{pgfscope}%
\pgfpathrectangle{\pgfqpoint{0.100000in}{2.493063in}}{\pgfqpoint{5.037500in}{3.427208in}}%
\pgfusepath{clip}%
\pgfsetrectcap%
\pgfsetroundjoin%
\pgfsetlinewidth{1.505625pt}%
\definecolor{currentstroke}{rgb}{0.000000,0.000000,1.000000}%
\pgfsetstrokecolor{currentstroke}%
\pgfsetstrokeopacity{0.500000}%
\pgfsetdash{}{0pt}%
\pgfpathmoveto{\pgfqpoint{0.342987in}{4.674739in}}%
\pgfusepath{stroke}%
\end{pgfscope}%
\begin{pgfscope}%
\pgfpathrectangle{\pgfqpoint{0.100000in}{2.493063in}}{\pgfqpoint{5.037500in}{3.427208in}}%
\pgfusepath{clip}%
\pgfsetbuttcap%
\pgfsetroundjoin%
\definecolor{currentfill}{rgb}{0.000000,0.000000,1.000000}%
\pgfsetfillcolor{currentfill}%
\pgfsetfillopacity{0.500000}%
\pgfsetlinewidth{0.250937pt}%
\definecolor{currentstroke}{rgb}{0.000000,0.000000,0.000000}%
\pgfsetstrokecolor{currentstroke}%
\pgfsetstrokeopacity{0.500000}%
\pgfsetdash{}{0pt}%
\pgfsys@defobject{currentmarker}{\pgfqpoint{-0.033333in}{-0.033333in}}{\pgfqpoint{0.033333in}{0.033333in}}{%
\pgfpathmoveto{\pgfqpoint{0.000000in}{-0.033333in}}%
\pgfpathcurveto{\pgfqpoint{0.008840in}{-0.033333in}}{\pgfqpoint{0.017319in}{-0.029821in}}{\pgfqpoint{0.023570in}{-0.023570in}}%
\pgfpathcurveto{\pgfqpoint{0.029821in}{-0.017319in}}{\pgfqpoint{0.033333in}{-0.008840in}}{\pgfqpoint{0.033333in}{0.000000in}}%
\pgfpathcurveto{\pgfqpoint{0.033333in}{0.008840in}}{\pgfqpoint{0.029821in}{0.017319in}}{\pgfqpoint{0.023570in}{0.023570in}}%
\pgfpathcurveto{\pgfqpoint{0.017319in}{0.029821in}}{\pgfqpoint{0.008840in}{0.033333in}}{\pgfqpoint{0.000000in}{0.033333in}}%
\pgfpathcurveto{\pgfqpoint{-0.008840in}{0.033333in}}{\pgfqpoint{-0.017319in}{0.029821in}}{\pgfqpoint{-0.023570in}{0.023570in}}%
\pgfpathcurveto{\pgfqpoint{-0.029821in}{0.017319in}}{\pgfqpoint{-0.033333in}{0.008840in}}{\pgfqpoint{-0.033333in}{0.000000in}}%
\pgfpathcurveto{\pgfqpoint{-0.033333in}{-0.008840in}}{\pgfqpoint{-0.029821in}{-0.017319in}}{\pgfqpoint{-0.023570in}{-0.023570in}}%
\pgfpathcurveto{\pgfqpoint{-0.017319in}{-0.029821in}}{\pgfqpoint{-0.008840in}{-0.033333in}}{\pgfqpoint{0.000000in}{-0.033333in}}%
\pgfpathlineto{\pgfqpoint{0.000000in}{-0.033333in}}%
\pgfpathclose%
\pgfusepath{stroke,fill}%
}%
\begin{pgfscope}%
\pgfsys@transformshift{0.342987in}{4.674739in}%
\pgfsys@useobject{currentmarker}{}%
\end{pgfscope}%
\end{pgfscope}%
\begin{pgfscope}%
\pgfpathrectangle{\pgfqpoint{0.100000in}{2.493063in}}{\pgfqpoint{5.037500in}{3.427208in}}%
\pgfusepath{clip}%
\pgfsetrectcap%
\pgfsetroundjoin%
\pgfsetlinewidth{1.505625pt}%
\definecolor{currentstroke}{rgb}{0.000000,0.000000,1.000000}%
\pgfsetstrokecolor{currentstroke}%
\pgfsetstrokeopacity{0.500000}%
\pgfsetdash{}{0pt}%
\pgfpathmoveto{\pgfqpoint{0.449905in}{4.583858in}}%
\pgfusepath{stroke}%
\end{pgfscope}%
\begin{pgfscope}%
\pgfpathrectangle{\pgfqpoint{0.100000in}{2.493063in}}{\pgfqpoint{5.037500in}{3.427208in}}%
\pgfusepath{clip}%
\pgfsetbuttcap%
\pgfsetroundjoin%
\definecolor{currentfill}{rgb}{0.000000,0.000000,1.000000}%
\pgfsetfillcolor{currentfill}%
\pgfsetfillopacity{0.500000}%
\pgfsetlinewidth{0.250937pt}%
\definecolor{currentstroke}{rgb}{0.000000,0.000000,0.000000}%
\pgfsetstrokecolor{currentstroke}%
\pgfsetstrokeopacity{0.500000}%
\pgfsetdash{}{0pt}%
\pgfsys@defobject{currentmarker}{\pgfqpoint{-0.022222in}{-0.022222in}}{\pgfqpoint{0.022222in}{0.022222in}}{%
\pgfpathmoveto{\pgfqpoint{0.000000in}{-0.022222in}}%
\pgfpathcurveto{\pgfqpoint{0.005893in}{-0.022222in}}{\pgfqpoint{0.011546in}{-0.019881in}}{\pgfqpoint{0.015713in}{-0.015713in}}%
\pgfpathcurveto{\pgfqpoint{0.019881in}{-0.011546in}}{\pgfqpoint{0.022222in}{-0.005893in}}{\pgfqpoint{0.022222in}{0.000000in}}%
\pgfpathcurveto{\pgfqpoint{0.022222in}{0.005893in}}{\pgfqpoint{0.019881in}{0.011546in}}{\pgfqpoint{0.015713in}{0.015713in}}%
\pgfpathcurveto{\pgfqpoint{0.011546in}{0.019881in}}{\pgfqpoint{0.005893in}{0.022222in}}{\pgfqpoint{0.000000in}{0.022222in}}%
\pgfpathcurveto{\pgfqpoint{-0.005893in}{0.022222in}}{\pgfqpoint{-0.011546in}{0.019881in}}{\pgfqpoint{-0.015713in}{0.015713in}}%
\pgfpathcurveto{\pgfqpoint{-0.019881in}{0.011546in}}{\pgfqpoint{-0.022222in}{0.005893in}}{\pgfqpoint{-0.022222in}{0.000000in}}%
\pgfpathcurveto{\pgfqpoint{-0.022222in}{-0.005893in}}{\pgfqpoint{-0.019881in}{-0.011546in}}{\pgfqpoint{-0.015713in}{-0.015713in}}%
\pgfpathcurveto{\pgfqpoint{-0.011546in}{-0.019881in}}{\pgfqpoint{-0.005893in}{-0.022222in}}{\pgfqpoint{0.000000in}{-0.022222in}}%
\pgfpathlineto{\pgfqpoint{0.000000in}{-0.022222in}}%
\pgfpathclose%
\pgfusepath{stroke,fill}%
}%
\begin{pgfscope}%
\pgfsys@transformshift{0.449905in}{4.583858in}%
\pgfsys@useobject{currentmarker}{}%
\end{pgfscope}%
\end{pgfscope}%
\begin{pgfscope}%
\pgfpathrectangle{\pgfqpoint{0.100000in}{2.493063in}}{\pgfqpoint{5.037500in}{3.427208in}}%
\pgfusepath{clip}%
\pgfsetrectcap%
\pgfsetroundjoin%
\pgfsetlinewidth{1.505625pt}%
\definecolor{currentstroke}{rgb}{0.000000,0.000000,1.000000}%
\pgfsetstrokecolor{currentstroke}%
\pgfsetstrokeopacity{0.500000}%
\pgfsetdash{}{0pt}%
\pgfpathmoveto{\pgfqpoint{0.370959in}{4.625432in}}%
\pgfusepath{stroke}%
\end{pgfscope}%
\begin{pgfscope}%
\pgfpathrectangle{\pgfqpoint{0.100000in}{2.493063in}}{\pgfqpoint{5.037500in}{3.427208in}}%
\pgfusepath{clip}%
\pgfsetbuttcap%
\pgfsetroundjoin%
\definecolor{currentfill}{rgb}{0.000000,0.000000,1.000000}%
\pgfsetfillcolor{currentfill}%
\pgfsetfillopacity{0.500000}%
\pgfsetlinewidth{0.250937pt}%
\definecolor{currentstroke}{rgb}{0.000000,0.000000,0.000000}%
\pgfsetstrokecolor{currentstroke}%
\pgfsetstrokeopacity{0.500000}%
\pgfsetdash{}{0pt}%
\pgfsys@defobject{currentmarker}{\pgfqpoint{-0.036111in}{-0.036111in}}{\pgfqpoint{0.036111in}{0.036111in}}{%
\pgfpathmoveto{\pgfqpoint{0.000000in}{-0.036111in}}%
\pgfpathcurveto{\pgfqpoint{0.009577in}{-0.036111in}}{\pgfqpoint{0.018763in}{-0.032306in}}{\pgfqpoint{0.025534in}{-0.025534in}}%
\pgfpathcurveto{\pgfqpoint{0.032306in}{-0.018763in}}{\pgfqpoint{0.036111in}{-0.009577in}}{\pgfqpoint{0.036111in}{0.000000in}}%
\pgfpathcurveto{\pgfqpoint{0.036111in}{0.009577in}}{\pgfqpoint{0.032306in}{0.018763in}}{\pgfqpoint{0.025534in}{0.025534in}}%
\pgfpathcurveto{\pgfqpoint{0.018763in}{0.032306in}}{\pgfqpoint{0.009577in}{0.036111in}}{\pgfqpoint{0.000000in}{0.036111in}}%
\pgfpathcurveto{\pgfqpoint{-0.009577in}{0.036111in}}{\pgfqpoint{-0.018763in}{0.032306in}}{\pgfqpoint{-0.025534in}{0.025534in}}%
\pgfpathcurveto{\pgfqpoint{-0.032306in}{0.018763in}}{\pgfqpoint{-0.036111in}{0.009577in}}{\pgfqpoint{-0.036111in}{0.000000in}}%
\pgfpathcurveto{\pgfqpoint{-0.036111in}{-0.009577in}}{\pgfqpoint{-0.032306in}{-0.018763in}}{\pgfqpoint{-0.025534in}{-0.025534in}}%
\pgfpathcurveto{\pgfqpoint{-0.018763in}{-0.032306in}}{\pgfqpoint{-0.009577in}{-0.036111in}}{\pgfqpoint{0.000000in}{-0.036111in}}%
\pgfpathlineto{\pgfqpoint{0.000000in}{-0.036111in}}%
\pgfpathclose%
\pgfusepath{stroke,fill}%
}%
\begin{pgfscope}%
\pgfsys@transformshift{0.370959in}{4.625432in}%
\pgfsys@useobject{currentmarker}{}%
\end{pgfscope}%
\end{pgfscope}%
\begin{pgfscope}%
\pgfpathrectangle{\pgfqpoint{0.100000in}{2.493063in}}{\pgfqpoint{5.037500in}{3.427208in}}%
\pgfusepath{clip}%
\pgfsetrectcap%
\pgfsetroundjoin%
\pgfsetlinewidth{1.505625pt}%
\definecolor{currentstroke}{rgb}{0.000000,0.000000,1.000000}%
\pgfsetstrokecolor{currentstroke}%
\pgfsetstrokeopacity{0.500000}%
\pgfsetdash{}{0pt}%
\pgfpathmoveto{\pgfqpoint{0.574887in}{4.353145in}}%
\pgfusepath{stroke}%
\end{pgfscope}%
\begin{pgfscope}%
\pgfpathrectangle{\pgfqpoint{0.100000in}{2.493063in}}{\pgfqpoint{5.037500in}{3.427208in}}%
\pgfusepath{clip}%
\pgfsetbuttcap%
\pgfsetroundjoin%
\definecolor{currentfill}{rgb}{0.000000,0.000000,1.000000}%
\pgfsetfillcolor{currentfill}%
\pgfsetfillopacity{0.500000}%
\pgfsetlinewidth{0.250937pt}%
\definecolor{currentstroke}{rgb}{0.000000,0.000000,0.000000}%
\pgfsetstrokecolor{currentstroke}%
\pgfsetstrokeopacity{0.500000}%
\pgfsetdash{}{0pt}%
\pgfsys@defobject{currentmarker}{\pgfqpoint{-0.030556in}{-0.030556in}}{\pgfqpoint{0.030556in}{0.030556in}}{%
\pgfpathmoveto{\pgfqpoint{0.000000in}{-0.030556in}}%
\pgfpathcurveto{\pgfqpoint{0.008103in}{-0.030556in}}{\pgfqpoint{0.015876in}{-0.027336in}}{\pgfqpoint{0.021606in}{-0.021606in}}%
\pgfpathcurveto{\pgfqpoint{0.027336in}{-0.015876in}}{\pgfqpoint{0.030556in}{-0.008103in}}{\pgfqpoint{0.030556in}{0.000000in}}%
\pgfpathcurveto{\pgfqpoint{0.030556in}{0.008103in}}{\pgfqpoint{0.027336in}{0.015876in}}{\pgfqpoint{0.021606in}{0.021606in}}%
\pgfpathcurveto{\pgfqpoint{0.015876in}{0.027336in}}{\pgfqpoint{0.008103in}{0.030556in}}{\pgfqpoint{0.000000in}{0.030556in}}%
\pgfpathcurveto{\pgfqpoint{-0.008103in}{0.030556in}}{\pgfqpoint{-0.015876in}{0.027336in}}{\pgfqpoint{-0.021606in}{0.021606in}}%
\pgfpathcurveto{\pgfqpoint{-0.027336in}{0.015876in}}{\pgfqpoint{-0.030556in}{0.008103in}}{\pgfqpoint{-0.030556in}{0.000000in}}%
\pgfpathcurveto{\pgfqpoint{-0.030556in}{-0.008103in}}{\pgfqpoint{-0.027336in}{-0.015876in}}{\pgfqpoint{-0.021606in}{-0.021606in}}%
\pgfpathcurveto{\pgfqpoint{-0.015876in}{-0.027336in}}{\pgfqpoint{-0.008103in}{-0.030556in}}{\pgfqpoint{0.000000in}{-0.030556in}}%
\pgfpathlineto{\pgfqpoint{0.000000in}{-0.030556in}}%
\pgfpathclose%
\pgfusepath{stroke,fill}%
}%
\begin{pgfscope}%
\pgfsys@transformshift{0.574887in}{4.353145in}%
\pgfsys@useobject{currentmarker}{}%
\end{pgfscope}%
\end{pgfscope}%
\begin{pgfscope}%
\pgfpathrectangle{\pgfqpoint{0.100000in}{2.493063in}}{\pgfqpoint{5.037500in}{3.427208in}}%
\pgfusepath{clip}%
\pgfsetrectcap%
\pgfsetroundjoin%
\pgfsetlinewidth{1.505625pt}%
\definecolor{currentstroke}{rgb}{0.000000,0.000000,1.000000}%
\pgfsetstrokecolor{currentstroke}%
\pgfsetstrokeopacity{0.500000}%
\pgfsetdash{}{0pt}%
\pgfpathmoveto{\pgfqpoint{0.460859in}{4.722756in}}%
\pgfusepath{stroke}%
\end{pgfscope}%
\begin{pgfscope}%
\pgfpathrectangle{\pgfqpoint{0.100000in}{2.493063in}}{\pgfqpoint{5.037500in}{3.427208in}}%
\pgfusepath{clip}%
\pgfsetbuttcap%
\pgfsetroundjoin%
\definecolor{currentfill}{rgb}{0.000000,0.000000,1.000000}%
\pgfsetfillcolor{currentfill}%
\pgfsetfillopacity{0.500000}%
\pgfsetlinewidth{0.250937pt}%
\definecolor{currentstroke}{rgb}{0.000000,0.000000,0.000000}%
\pgfsetstrokecolor{currentstroke}%
\pgfsetstrokeopacity{0.500000}%
\pgfsetdash{}{0pt}%
\pgfsys@defobject{currentmarker}{\pgfqpoint{-0.022222in}{-0.022222in}}{\pgfqpoint{0.022222in}{0.022222in}}{%
\pgfpathmoveto{\pgfqpoint{0.000000in}{-0.022222in}}%
\pgfpathcurveto{\pgfqpoint{0.005893in}{-0.022222in}}{\pgfqpoint{0.011546in}{-0.019881in}}{\pgfqpoint{0.015713in}{-0.015713in}}%
\pgfpathcurveto{\pgfqpoint{0.019881in}{-0.011546in}}{\pgfqpoint{0.022222in}{-0.005893in}}{\pgfqpoint{0.022222in}{0.000000in}}%
\pgfpathcurveto{\pgfqpoint{0.022222in}{0.005893in}}{\pgfqpoint{0.019881in}{0.011546in}}{\pgfqpoint{0.015713in}{0.015713in}}%
\pgfpathcurveto{\pgfqpoint{0.011546in}{0.019881in}}{\pgfqpoint{0.005893in}{0.022222in}}{\pgfqpoint{0.000000in}{0.022222in}}%
\pgfpathcurveto{\pgfqpoint{-0.005893in}{0.022222in}}{\pgfqpoint{-0.011546in}{0.019881in}}{\pgfqpoint{-0.015713in}{0.015713in}}%
\pgfpathcurveto{\pgfqpoint{-0.019881in}{0.011546in}}{\pgfqpoint{-0.022222in}{0.005893in}}{\pgfqpoint{-0.022222in}{0.000000in}}%
\pgfpathcurveto{\pgfqpoint{-0.022222in}{-0.005893in}}{\pgfqpoint{-0.019881in}{-0.011546in}}{\pgfqpoint{-0.015713in}{-0.015713in}}%
\pgfpathcurveto{\pgfqpoint{-0.011546in}{-0.019881in}}{\pgfqpoint{-0.005893in}{-0.022222in}}{\pgfqpoint{0.000000in}{-0.022222in}}%
\pgfpathlineto{\pgfqpoint{0.000000in}{-0.022222in}}%
\pgfpathclose%
\pgfusepath{stroke,fill}%
}%
\begin{pgfscope}%
\pgfsys@transformshift{0.460859in}{4.722756in}%
\pgfsys@useobject{currentmarker}{}%
\end{pgfscope}%
\end{pgfscope}%
\begin{pgfscope}%
\pgfpathrectangle{\pgfqpoint{0.100000in}{2.493063in}}{\pgfqpoint{5.037500in}{3.427208in}}%
\pgfusepath{clip}%
\pgfsetrectcap%
\pgfsetroundjoin%
\pgfsetlinewidth{1.505625pt}%
\definecolor{currentstroke}{rgb}{0.000000,0.000000,1.000000}%
\pgfsetstrokecolor{currentstroke}%
\pgfsetstrokeopacity{0.500000}%
\pgfsetdash{}{0pt}%
\pgfpathmoveto{\pgfqpoint{1.900192in}{4.528961in}}%
\pgfusepath{stroke}%
\end{pgfscope}%
\begin{pgfscope}%
\pgfpathrectangle{\pgfqpoint{0.100000in}{2.493063in}}{\pgfqpoint{5.037500in}{3.427208in}}%
\pgfusepath{clip}%
\pgfsetbuttcap%
\pgfsetroundjoin%
\definecolor{currentfill}{rgb}{0.000000,0.000000,1.000000}%
\pgfsetfillcolor{currentfill}%
\pgfsetfillopacity{0.500000}%
\pgfsetlinewidth{0.250937pt}%
\definecolor{currentstroke}{rgb}{0.000000,0.000000,0.000000}%
\pgfsetstrokecolor{currentstroke}%
\pgfsetstrokeopacity{0.500000}%
\pgfsetdash{}{0pt}%
\pgfsys@defobject{currentmarker}{\pgfqpoint{-0.019444in}{-0.019444in}}{\pgfqpoint{0.019444in}{0.019444in}}{%
\pgfpathmoveto{\pgfqpoint{0.000000in}{-0.019444in}}%
\pgfpathcurveto{\pgfqpoint{0.005157in}{-0.019444in}}{\pgfqpoint{0.010103in}{-0.017396in}}{\pgfqpoint{0.013749in}{-0.013749in}}%
\pgfpathcurveto{\pgfqpoint{0.017396in}{-0.010103in}}{\pgfqpoint{0.019444in}{-0.005157in}}{\pgfqpoint{0.019444in}{0.000000in}}%
\pgfpathcurveto{\pgfqpoint{0.019444in}{0.005157in}}{\pgfqpoint{0.017396in}{0.010103in}}{\pgfqpoint{0.013749in}{0.013749in}}%
\pgfpathcurveto{\pgfqpoint{0.010103in}{0.017396in}}{\pgfqpoint{0.005157in}{0.019444in}}{\pgfqpoint{0.000000in}{0.019444in}}%
\pgfpathcurveto{\pgfqpoint{-0.005157in}{0.019444in}}{\pgfqpoint{-0.010103in}{0.017396in}}{\pgfqpoint{-0.013749in}{0.013749in}}%
\pgfpathcurveto{\pgfqpoint{-0.017396in}{0.010103in}}{\pgfqpoint{-0.019444in}{0.005157in}}{\pgfqpoint{-0.019444in}{0.000000in}}%
\pgfpathcurveto{\pgfqpoint{-0.019444in}{-0.005157in}}{\pgfqpoint{-0.017396in}{-0.010103in}}{\pgfqpoint{-0.013749in}{-0.013749in}}%
\pgfpathcurveto{\pgfqpoint{-0.010103in}{-0.017396in}}{\pgfqpoint{-0.005157in}{-0.019444in}}{\pgfqpoint{0.000000in}{-0.019444in}}%
\pgfpathlineto{\pgfqpoint{0.000000in}{-0.019444in}}%
\pgfpathclose%
\pgfusepath{stroke,fill}%
}%
\begin{pgfscope}%
\pgfsys@transformshift{1.900192in}{4.528961in}%
\pgfsys@useobject{currentmarker}{}%
\end{pgfscope}%
\end{pgfscope}%
\begin{pgfscope}%
\pgfpathrectangle{\pgfqpoint{0.100000in}{2.493063in}}{\pgfqpoint{5.037500in}{3.427208in}}%
\pgfusepath{clip}%
\pgfsetrectcap%
\pgfsetroundjoin%
\pgfsetlinewidth{1.505625pt}%
\definecolor{currentstroke}{rgb}{0.000000,0.000000,1.000000}%
\pgfsetstrokecolor{currentstroke}%
\pgfsetstrokeopacity{0.500000}%
\pgfsetdash{}{0pt}%
\pgfpathmoveto{\pgfqpoint{1.924544in}{4.389425in}}%
\pgfusepath{stroke}%
\end{pgfscope}%
\begin{pgfscope}%
\pgfpathrectangle{\pgfqpoint{0.100000in}{2.493063in}}{\pgfqpoint{5.037500in}{3.427208in}}%
\pgfusepath{clip}%
\pgfsetbuttcap%
\pgfsetroundjoin%
\definecolor{currentfill}{rgb}{0.000000,0.000000,1.000000}%
\pgfsetfillcolor{currentfill}%
\pgfsetfillopacity{0.500000}%
\pgfsetlinewidth{0.250937pt}%
\definecolor{currentstroke}{rgb}{0.000000,0.000000,0.000000}%
\pgfsetstrokecolor{currentstroke}%
\pgfsetstrokeopacity{0.500000}%
\pgfsetdash{}{0pt}%
\pgfsys@defobject{currentmarker}{\pgfqpoint{-0.008333in}{-0.008333in}}{\pgfqpoint{0.008333in}{0.008333in}}{%
\pgfpathmoveto{\pgfqpoint{0.000000in}{-0.008333in}}%
\pgfpathcurveto{\pgfqpoint{0.002210in}{-0.008333in}}{\pgfqpoint{0.004330in}{-0.007455in}}{\pgfqpoint{0.005893in}{-0.005893in}}%
\pgfpathcurveto{\pgfqpoint{0.007455in}{-0.004330in}}{\pgfqpoint{0.008333in}{-0.002210in}}{\pgfqpoint{0.008333in}{0.000000in}}%
\pgfpathcurveto{\pgfqpoint{0.008333in}{0.002210in}}{\pgfqpoint{0.007455in}{0.004330in}}{\pgfqpoint{0.005893in}{0.005893in}}%
\pgfpathcurveto{\pgfqpoint{0.004330in}{0.007455in}}{\pgfqpoint{0.002210in}{0.008333in}}{\pgfqpoint{0.000000in}{0.008333in}}%
\pgfpathcurveto{\pgfqpoint{-0.002210in}{0.008333in}}{\pgfqpoint{-0.004330in}{0.007455in}}{\pgfqpoint{-0.005893in}{0.005893in}}%
\pgfpathcurveto{\pgfqpoint{-0.007455in}{0.004330in}}{\pgfqpoint{-0.008333in}{0.002210in}}{\pgfqpoint{-0.008333in}{0.000000in}}%
\pgfpathcurveto{\pgfqpoint{-0.008333in}{-0.002210in}}{\pgfqpoint{-0.007455in}{-0.004330in}}{\pgfqpoint{-0.005893in}{-0.005893in}}%
\pgfpathcurveto{\pgfqpoint{-0.004330in}{-0.007455in}}{\pgfqpoint{-0.002210in}{-0.008333in}}{\pgfqpoint{0.000000in}{-0.008333in}}%
\pgfpathlineto{\pgfqpoint{0.000000in}{-0.008333in}}%
\pgfpathclose%
\pgfusepath{stroke,fill}%
}%
\begin{pgfscope}%
\pgfsys@transformshift{1.924544in}{4.389425in}%
\pgfsys@useobject{currentmarker}{}%
\end{pgfscope}%
\end{pgfscope}%
\begin{pgfscope}%
\pgfpathrectangle{\pgfqpoint{0.100000in}{2.493063in}}{\pgfqpoint{5.037500in}{3.427208in}}%
\pgfusepath{clip}%
\pgfsetrectcap%
\pgfsetroundjoin%
\pgfsetlinewidth{1.505625pt}%
\definecolor{currentstroke}{rgb}{0.000000,0.000000,1.000000}%
\pgfsetstrokecolor{currentstroke}%
\pgfsetstrokeopacity{0.500000}%
\pgfsetdash{}{0pt}%
\pgfpathmoveto{\pgfqpoint{1.921764in}{4.494584in}}%
\pgfusepath{stroke}%
\end{pgfscope}%
\begin{pgfscope}%
\pgfpathrectangle{\pgfqpoint{0.100000in}{2.493063in}}{\pgfqpoint{5.037500in}{3.427208in}}%
\pgfusepath{clip}%
\pgfsetbuttcap%
\pgfsetroundjoin%
\definecolor{currentfill}{rgb}{0.000000,0.000000,1.000000}%
\pgfsetfillcolor{currentfill}%
\pgfsetfillopacity{0.500000}%
\pgfsetlinewidth{0.250937pt}%
\definecolor{currentstroke}{rgb}{0.000000,0.000000,0.000000}%
\pgfsetstrokecolor{currentstroke}%
\pgfsetstrokeopacity{0.500000}%
\pgfsetdash{}{0pt}%
\pgfsys@defobject{currentmarker}{\pgfqpoint{-0.022222in}{-0.022222in}}{\pgfqpoint{0.022222in}{0.022222in}}{%
\pgfpathmoveto{\pgfqpoint{0.000000in}{-0.022222in}}%
\pgfpathcurveto{\pgfqpoint{0.005893in}{-0.022222in}}{\pgfqpoint{0.011546in}{-0.019881in}}{\pgfqpoint{0.015713in}{-0.015713in}}%
\pgfpathcurveto{\pgfqpoint{0.019881in}{-0.011546in}}{\pgfqpoint{0.022222in}{-0.005893in}}{\pgfqpoint{0.022222in}{0.000000in}}%
\pgfpathcurveto{\pgfqpoint{0.022222in}{0.005893in}}{\pgfqpoint{0.019881in}{0.011546in}}{\pgfqpoint{0.015713in}{0.015713in}}%
\pgfpathcurveto{\pgfqpoint{0.011546in}{0.019881in}}{\pgfqpoint{0.005893in}{0.022222in}}{\pgfqpoint{0.000000in}{0.022222in}}%
\pgfpathcurveto{\pgfqpoint{-0.005893in}{0.022222in}}{\pgfqpoint{-0.011546in}{0.019881in}}{\pgfqpoint{-0.015713in}{0.015713in}}%
\pgfpathcurveto{\pgfqpoint{-0.019881in}{0.011546in}}{\pgfqpoint{-0.022222in}{0.005893in}}{\pgfqpoint{-0.022222in}{0.000000in}}%
\pgfpathcurveto{\pgfqpoint{-0.022222in}{-0.005893in}}{\pgfqpoint{-0.019881in}{-0.011546in}}{\pgfqpoint{-0.015713in}{-0.015713in}}%
\pgfpathcurveto{\pgfqpoint{-0.011546in}{-0.019881in}}{\pgfqpoint{-0.005893in}{-0.022222in}}{\pgfqpoint{0.000000in}{-0.022222in}}%
\pgfpathlineto{\pgfqpoint{0.000000in}{-0.022222in}}%
\pgfpathclose%
\pgfusepath{stroke,fill}%
}%
\begin{pgfscope}%
\pgfsys@transformshift{1.921764in}{4.494584in}%
\pgfsys@useobject{currentmarker}{}%
\end{pgfscope}%
\end{pgfscope}%
\begin{pgfscope}%
\pgfpathrectangle{\pgfqpoint{0.100000in}{2.493063in}}{\pgfqpoint{5.037500in}{3.427208in}}%
\pgfusepath{clip}%
\pgfsetrectcap%
\pgfsetroundjoin%
\pgfsetlinewidth{1.505625pt}%
\definecolor{currentstroke}{rgb}{0.000000,0.000000,1.000000}%
\pgfsetstrokecolor{currentstroke}%
\pgfsetstrokeopacity{0.500000}%
\pgfsetdash{}{0pt}%
\pgfpathmoveto{\pgfqpoint{1.924872in}{4.588293in}}%
\pgfusepath{stroke}%
\end{pgfscope}%
\begin{pgfscope}%
\pgfpathrectangle{\pgfqpoint{0.100000in}{2.493063in}}{\pgfqpoint{5.037500in}{3.427208in}}%
\pgfusepath{clip}%
\pgfsetbuttcap%
\pgfsetroundjoin%
\definecolor{currentfill}{rgb}{0.000000,0.000000,1.000000}%
\pgfsetfillcolor{currentfill}%
\pgfsetfillopacity{0.500000}%
\pgfsetlinewidth{0.250937pt}%
\definecolor{currentstroke}{rgb}{0.000000,0.000000,0.000000}%
\pgfsetstrokecolor{currentstroke}%
\pgfsetstrokeopacity{0.500000}%
\pgfsetdash{}{0pt}%
\pgfsys@defobject{currentmarker}{\pgfqpoint{-0.016667in}{-0.016667in}}{\pgfqpoint{0.016667in}{0.016667in}}{%
\pgfpathmoveto{\pgfqpoint{0.000000in}{-0.016667in}}%
\pgfpathcurveto{\pgfqpoint{0.004420in}{-0.016667in}}{\pgfqpoint{0.008660in}{-0.014911in}}{\pgfqpoint{0.011785in}{-0.011785in}}%
\pgfpathcurveto{\pgfqpoint{0.014911in}{-0.008660in}}{\pgfqpoint{0.016667in}{-0.004420in}}{\pgfqpoint{0.016667in}{0.000000in}}%
\pgfpathcurveto{\pgfqpoint{0.016667in}{0.004420in}}{\pgfqpoint{0.014911in}{0.008660in}}{\pgfqpoint{0.011785in}{0.011785in}}%
\pgfpathcurveto{\pgfqpoint{0.008660in}{0.014911in}}{\pgfqpoint{0.004420in}{0.016667in}}{\pgfqpoint{0.000000in}{0.016667in}}%
\pgfpathcurveto{\pgfqpoint{-0.004420in}{0.016667in}}{\pgfqpoint{-0.008660in}{0.014911in}}{\pgfqpoint{-0.011785in}{0.011785in}}%
\pgfpathcurveto{\pgfqpoint{-0.014911in}{0.008660in}}{\pgfqpoint{-0.016667in}{0.004420in}}{\pgfqpoint{-0.016667in}{0.000000in}}%
\pgfpathcurveto{\pgfqpoint{-0.016667in}{-0.004420in}}{\pgfqpoint{-0.014911in}{-0.008660in}}{\pgfqpoint{-0.011785in}{-0.011785in}}%
\pgfpathcurveto{\pgfqpoint{-0.008660in}{-0.014911in}}{\pgfqpoint{-0.004420in}{-0.016667in}}{\pgfqpoint{0.000000in}{-0.016667in}}%
\pgfpathlineto{\pgfqpoint{0.000000in}{-0.016667in}}%
\pgfpathclose%
\pgfusepath{stroke,fill}%
}%
\begin{pgfscope}%
\pgfsys@transformshift{1.924872in}{4.588293in}%
\pgfsys@useobject{currentmarker}{}%
\end{pgfscope}%
\end{pgfscope}%
\begin{pgfscope}%
\pgfpathrectangle{\pgfqpoint{0.100000in}{2.493063in}}{\pgfqpoint{5.037500in}{3.427208in}}%
\pgfusepath{clip}%
\pgfsetrectcap%
\pgfsetroundjoin%
\pgfsetlinewidth{1.505625pt}%
\definecolor{currentstroke}{rgb}{0.000000,0.000000,1.000000}%
\pgfsetstrokecolor{currentstroke}%
\pgfsetstrokeopacity{0.500000}%
\pgfsetdash{}{0pt}%
\pgfpathmoveto{\pgfqpoint{1.597107in}{4.458471in}}%
\pgfusepath{stroke}%
\end{pgfscope}%
\begin{pgfscope}%
\pgfpathrectangle{\pgfqpoint{0.100000in}{2.493063in}}{\pgfqpoint{5.037500in}{3.427208in}}%
\pgfusepath{clip}%
\pgfsetbuttcap%
\pgfsetroundjoin%
\definecolor{currentfill}{rgb}{0.000000,0.000000,1.000000}%
\pgfsetfillcolor{currentfill}%
\pgfsetfillopacity{0.500000}%
\pgfsetlinewidth{0.250937pt}%
\definecolor{currentstroke}{rgb}{0.000000,0.000000,0.000000}%
\pgfsetstrokecolor{currentstroke}%
\pgfsetstrokeopacity{0.500000}%
\pgfsetdash{}{0pt}%
\pgfsys@defobject{currentmarker}{\pgfqpoint{-0.005556in}{-0.005556in}}{\pgfqpoint{0.005556in}{0.005556in}}{%
\pgfpathmoveto{\pgfqpoint{0.000000in}{-0.005556in}}%
\pgfpathcurveto{\pgfqpoint{0.001473in}{-0.005556in}}{\pgfqpoint{0.002887in}{-0.004970in}}{\pgfqpoint{0.003928in}{-0.003928in}}%
\pgfpathcurveto{\pgfqpoint{0.004970in}{-0.002887in}}{\pgfqpoint{0.005556in}{-0.001473in}}{\pgfqpoint{0.005556in}{0.000000in}}%
\pgfpathcurveto{\pgfqpoint{0.005556in}{0.001473in}}{\pgfqpoint{0.004970in}{0.002887in}}{\pgfqpoint{0.003928in}{0.003928in}}%
\pgfpathcurveto{\pgfqpoint{0.002887in}{0.004970in}}{\pgfqpoint{0.001473in}{0.005556in}}{\pgfqpoint{0.000000in}{0.005556in}}%
\pgfpathcurveto{\pgfqpoint{-0.001473in}{0.005556in}}{\pgfqpoint{-0.002887in}{0.004970in}}{\pgfqpoint{-0.003928in}{0.003928in}}%
\pgfpathcurveto{\pgfqpoint{-0.004970in}{0.002887in}}{\pgfqpoint{-0.005556in}{0.001473in}}{\pgfqpoint{-0.005556in}{0.000000in}}%
\pgfpathcurveto{\pgfqpoint{-0.005556in}{-0.001473in}}{\pgfqpoint{-0.004970in}{-0.002887in}}{\pgfqpoint{-0.003928in}{-0.003928in}}%
\pgfpathcurveto{\pgfqpoint{-0.002887in}{-0.004970in}}{\pgfqpoint{-0.001473in}{-0.005556in}}{\pgfqpoint{0.000000in}{-0.005556in}}%
\pgfpathlineto{\pgfqpoint{0.000000in}{-0.005556in}}%
\pgfpathclose%
\pgfusepath{stroke,fill}%
}%
\begin{pgfscope}%
\pgfsys@transformshift{1.597107in}{4.458471in}%
\pgfsys@useobject{currentmarker}{}%
\end{pgfscope}%
\end{pgfscope}%
\begin{pgfscope}%
\pgfpathrectangle{\pgfqpoint{0.100000in}{2.493063in}}{\pgfqpoint{5.037500in}{3.427208in}}%
\pgfusepath{clip}%
\pgfsetrectcap%
\pgfsetroundjoin%
\pgfsetlinewidth{1.505625pt}%
\definecolor{currentstroke}{rgb}{0.000000,0.000000,1.000000}%
\pgfsetstrokecolor{currentstroke}%
\pgfsetstrokeopacity{0.500000}%
\pgfsetdash{}{0pt}%
\pgfpathmoveto{\pgfqpoint{1.954425in}{4.570288in}}%
\pgfusepath{stroke}%
\end{pgfscope}%
\begin{pgfscope}%
\pgfpathrectangle{\pgfqpoint{0.100000in}{2.493063in}}{\pgfqpoint{5.037500in}{3.427208in}}%
\pgfusepath{clip}%
\pgfsetbuttcap%
\pgfsetroundjoin%
\definecolor{currentfill}{rgb}{0.000000,0.000000,1.000000}%
\pgfsetfillcolor{currentfill}%
\pgfsetfillopacity{0.500000}%
\pgfsetlinewidth{0.250937pt}%
\definecolor{currentstroke}{rgb}{0.000000,0.000000,0.000000}%
\pgfsetstrokecolor{currentstroke}%
\pgfsetstrokeopacity{0.500000}%
\pgfsetdash{}{0pt}%
\pgfsys@defobject{currentmarker}{\pgfqpoint{-0.025000in}{-0.025000in}}{\pgfqpoint{0.025000in}{0.025000in}}{%
\pgfpathmoveto{\pgfqpoint{0.000000in}{-0.025000in}}%
\pgfpathcurveto{\pgfqpoint{0.006630in}{-0.025000in}}{\pgfqpoint{0.012989in}{-0.022366in}}{\pgfqpoint{0.017678in}{-0.017678in}}%
\pgfpathcurveto{\pgfqpoint{0.022366in}{-0.012989in}}{\pgfqpoint{0.025000in}{-0.006630in}}{\pgfqpoint{0.025000in}{0.000000in}}%
\pgfpathcurveto{\pgfqpoint{0.025000in}{0.006630in}}{\pgfqpoint{0.022366in}{0.012989in}}{\pgfqpoint{0.017678in}{0.017678in}}%
\pgfpathcurveto{\pgfqpoint{0.012989in}{0.022366in}}{\pgfqpoint{0.006630in}{0.025000in}}{\pgfqpoint{0.000000in}{0.025000in}}%
\pgfpathcurveto{\pgfqpoint{-0.006630in}{0.025000in}}{\pgfqpoint{-0.012989in}{0.022366in}}{\pgfqpoint{-0.017678in}{0.017678in}}%
\pgfpathcurveto{\pgfqpoint{-0.022366in}{0.012989in}}{\pgfqpoint{-0.025000in}{0.006630in}}{\pgfqpoint{-0.025000in}{0.000000in}}%
\pgfpathcurveto{\pgfqpoint{-0.025000in}{-0.006630in}}{\pgfqpoint{-0.022366in}{-0.012989in}}{\pgfqpoint{-0.017678in}{-0.017678in}}%
\pgfpathcurveto{\pgfqpoint{-0.012989in}{-0.022366in}}{\pgfqpoint{-0.006630in}{-0.025000in}}{\pgfqpoint{0.000000in}{-0.025000in}}%
\pgfpathlineto{\pgfqpoint{0.000000in}{-0.025000in}}%
\pgfpathclose%
\pgfusepath{stroke,fill}%
}%
\begin{pgfscope}%
\pgfsys@transformshift{1.954425in}{4.570288in}%
\pgfsys@useobject{currentmarker}{}%
\end{pgfscope}%
\end{pgfscope}%
\begin{pgfscope}%
\pgfpathrectangle{\pgfqpoint{0.100000in}{2.493063in}}{\pgfqpoint{5.037500in}{3.427208in}}%
\pgfusepath{clip}%
\pgfsetrectcap%
\pgfsetroundjoin%
\pgfsetlinewidth{1.505625pt}%
\definecolor{currentstroke}{rgb}{0.000000,0.000000,1.000000}%
\pgfsetstrokecolor{currentstroke}%
\pgfsetstrokeopacity{0.500000}%
\pgfsetdash{}{0pt}%
\pgfpathmoveto{\pgfqpoint{1.936783in}{4.320996in}}%
\pgfusepath{stroke}%
\end{pgfscope}%
\begin{pgfscope}%
\pgfpathrectangle{\pgfqpoint{0.100000in}{2.493063in}}{\pgfqpoint{5.037500in}{3.427208in}}%
\pgfusepath{clip}%
\pgfsetbuttcap%
\pgfsetroundjoin%
\definecolor{currentfill}{rgb}{0.000000,0.000000,1.000000}%
\pgfsetfillcolor{currentfill}%
\pgfsetfillopacity{0.500000}%
\pgfsetlinewidth{0.250937pt}%
\definecolor{currentstroke}{rgb}{0.000000,0.000000,0.000000}%
\pgfsetstrokecolor{currentstroke}%
\pgfsetstrokeopacity{0.500000}%
\pgfsetdash{}{0pt}%
\pgfsys@defobject{currentmarker}{\pgfqpoint{-0.013889in}{-0.013889in}}{\pgfqpoint{0.013889in}{0.013889in}}{%
\pgfpathmoveto{\pgfqpoint{0.000000in}{-0.013889in}}%
\pgfpathcurveto{\pgfqpoint{0.003683in}{-0.013889in}}{\pgfqpoint{0.007216in}{-0.012425in}}{\pgfqpoint{0.009821in}{-0.009821in}}%
\pgfpathcurveto{\pgfqpoint{0.012425in}{-0.007216in}}{\pgfqpoint{0.013889in}{-0.003683in}}{\pgfqpoint{0.013889in}{0.000000in}}%
\pgfpathcurveto{\pgfqpoint{0.013889in}{0.003683in}}{\pgfqpoint{0.012425in}{0.007216in}}{\pgfqpoint{0.009821in}{0.009821in}}%
\pgfpathcurveto{\pgfqpoint{0.007216in}{0.012425in}}{\pgfqpoint{0.003683in}{0.013889in}}{\pgfqpoint{0.000000in}{0.013889in}}%
\pgfpathcurveto{\pgfqpoint{-0.003683in}{0.013889in}}{\pgfqpoint{-0.007216in}{0.012425in}}{\pgfqpoint{-0.009821in}{0.009821in}}%
\pgfpathcurveto{\pgfqpoint{-0.012425in}{0.007216in}}{\pgfqpoint{-0.013889in}{0.003683in}}{\pgfqpoint{-0.013889in}{0.000000in}}%
\pgfpathcurveto{\pgfqpoint{-0.013889in}{-0.003683in}}{\pgfqpoint{-0.012425in}{-0.007216in}}{\pgfqpoint{-0.009821in}{-0.009821in}}%
\pgfpathcurveto{\pgfqpoint{-0.007216in}{-0.012425in}}{\pgfqpoint{-0.003683in}{-0.013889in}}{\pgfqpoint{0.000000in}{-0.013889in}}%
\pgfpathlineto{\pgfqpoint{0.000000in}{-0.013889in}}%
\pgfpathclose%
\pgfusepath{stroke,fill}%
}%
\begin{pgfscope}%
\pgfsys@transformshift{1.936783in}{4.320996in}%
\pgfsys@useobject{currentmarker}{}%
\end{pgfscope}%
\end{pgfscope}%
\begin{pgfscope}%
\pgfpathrectangle{\pgfqpoint{0.100000in}{2.493063in}}{\pgfqpoint{5.037500in}{3.427208in}}%
\pgfusepath{clip}%
\pgfsetrectcap%
\pgfsetroundjoin%
\pgfsetlinewidth{1.505625pt}%
\definecolor{currentstroke}{rgb}{0.000000,0.000000,1.000000}%
\pgfsetstrokecolor{currentstroke}%
\pgfsetstrokeopacity{0.500000}%
\pgfsetdash{}{0pt}%
\pgfpathmoveto{\pgfqpoint{4.676566in}{4.836149in}}%
\pgfusepath{stroke}%
\end{pgfscope}%
\begin{pgfscope}%
\pgfpathrectangle{\pgfqpoint{0.100000in}{2.493063in}}{\pgfqpoint{5.037500in}{3.427208in}}%
\pgfusepath{clip}%
\pgfsetbuttcap%
\pgfsetroundjoin%
\definecolor{currentfill}{rgb}{0.000000,0.000000,1.000000}%
\pgfsetfillcolor{currentfill}%
\pgfsetfillopacity{0.500000}%
\pgfsetlinewidth{0.250937pt}%
\definecolor{currentstroke}{rgb}{0.000000,0.000000,0.000000}%
\pgfsetstrokecolor{currentstroke}%
\pgfsetstrokeopacity{0.500000}%
\pgfsetdash{}{0pt}%
\pgfsys@defobject{currentmarker}{\pgfqpoint{-0.016667in}{-0.016667in}}{\pgfqpoint{0.016667in}{0.016667in}}{%
\pgfpathmoveto{\pgfqpoint{0.000000in}{-0.016667in}}%
\pgfpathcurveto{\pgfqpoint{0.004420in}{-0.016667in}}{\pgfqpoint{0.008660in}{-0.014911in}}{\pgfqpoint{0.011785in}{-0.011785in}}%
\pgfpathcurveto{\pgfqpoint{0.014911in}{-0.008660in}}{\pgfqpoint{0.016667in}{-0.004420in}}{\pgfqpoint{0.016667in}{0.000000in}}%
\pgfpathcurveto{\pgfqpoint{0.016667in}{0.004420in}}{\pgfqpoint{0.014911in}{0.008660in}}{\pgfqpoint{0.011785in}{0.011785in}}%
\pgfpathcurveto{\pgfqpoint{0.008660in}{0.014911in}}{\pgfqpoint{0.004420in}{0.016667in}}{\pgfqpoint{0.000000in}{0.016667in}}%
\pgfpathcurveto{\pgfqpoint{-0.004420in}{0.016667in}}{\pgfqpoint{-0.008660in}{0.014911in}}{\pgfqpoint{-0.011785in}{0.011785in}}%
\pgfpathcurveto{\pgfqpoint{-0.014911in}{0.008660in}}{\pgfqpoint{-0.016667in}{0.004420in}}{\pgfqpoint{-0.016667in}{0.000000in}}%
\pgfpathcurveto{\pgfqpoint{-0.016667in}{-0.004420in}}{\pgfqpoint{-0.014911in}{-0.008660in}}{\pgfqpoint{-0.011785in}{-0.011785in}}%
\pgfpathcurveto{\pgfqpoint{-0.008660in}{-0.014911in}}{\pgfqpoint{-0.004420in}{-0.016667in}}{\pgfqpoint{0.000000in}{-0.016667in}}%
\pgfpathlineto{\pgfqpoint{0.000000in}{-0.016667in}}%
\pgfpathclose%
\pgfusepath{stroke,fill}%
}%
\begin{pgfscope}%
\pgfsys@transformshift{4.676566in}{4.836149in}%
\pgfsys@useobject{currentmarker}{}%
\end{pgfscope}%
\end{pgfscope}%
\begin{pgfscope}%
\pgfpathrectangle{\pgfqpoint{0.100000in}{2.493063in}}{\pgfqpoint{5.037500in}{3.427208in}}%
\pgfusepath{clip}%
\pgfsetrectcap%
\pgfsetroundjoin%
\pgfsetlinewidth{1.505625pt}%
\definecolor{currentstroke}{rgb}{0.000000,0.000000,1.000000}%
\pgfsetstrokecolor{currentstroke}%
\pgfsetstrokeopacity{0.500000}%
\pgfsetdash{}{0pt}%
\pgfpathmoveto{\pgfqpoint{4.649399in}{4.856556in}}%
\pgfusepath{stroke}%
\end{pgfscope}%
\begin{pgfscope}%
\pgfpathrectangle{\pgfqpoint{0.100000in}{2.493063in}}{\pgfqpoint{5.037500in}{3.427208in}}%
\pgfusepath{clip}%
\pgfsetbuttcap%
\pgfsetroundjoin%
\definecolor{currentfill}{rgb}{0.000000,0.000000,1.000000}%
\pgfsetfillcolor{currentfill}%
\pgfsetfillopacity{0.500000}%
\pgfsetlinewidth{0.250937pt}%
\definecolor{currentstroke}{rgb}{0.000000,0.000000,0.000000}%
\pgfsetstrokecolor{currentstroke}%
\pgfsetstrokeopacity{0.500000}%
\pgfsetdash{}{0pt}%
\pgfsys@defobject{currentmarker}{\pgfqpoint{-0.016667in}{-0.016667in}}{\pgfqpoint{0.016667in}{0.016667in}}{%
\pgfpathmoveto{\pgfqpoint{0.000000in}{-0.016667in}}%
\pgfpathcurveto{\pgfqpoint{0.004420in}{-0.016667in}}{\pgfqpoint{0.008660in}{-0.014911in}}{\pgfqpoint{0.011785in}{-0.011785in}}%
\pgfpathcurveto{\pgfqpoint{0.014911in}{-0.008660in}}{\pgfqpoint{0.016667in}{-0.004420in}}{\pgfqpoint{0.016667in}{0.000000in}}%
\pgfpathcurveto{\pgfqpoint{0.016667in}{0.004420in}}{\pgfqpoint{0.014911in}{0.008660in}}{\pgfqpoint{0.011785in}{0.011785in}}%
\pgfpathcurveto{\pgfqpoint{0.008660in}{0.014911in}}{\pgfqpoint{0.004420in}{0.016667in}}{\pgfqpoint{0.000000in}{0.016667in}}%
\pgfpathcurveto{\pgfqpoint{-0.004420in}{0.016667in}}{\pgfqpoint{-0.008660in}{0.014911in}}{\pgfqpoint{-0.011785in}{0.011785in}}%
\pgfpathcurveto{\pgfqpoint{-0.014911in}{0.008660in}}{\pgfqpoint{-0.016667in}{0.004420in}}{\pgfqpoint{-0.016667in}{0.000000in}}%
\pgfpathcurveto{\pgfqpoint{-0.016667in}{-0.004420in}}{\pgfqpoint{-0.014911in}{-0.008660in}}{\pgfqpoint{-0.011785in}{-0.011785in}}%
\pgfpathcurveto{\pgfqpoint{-0.008660in}{-0.014911in}}{\pgfqpoint{-0.004420in}{-0.016667in}}{\pgfqpoint{0.000000in}{-0.016667in}}%
\pgfpathlineto{\pgfqpoint{0.000000in}{-0.016667in}}%
\pgfpathclose%
\pgfusepath{stroke,fill}%
}%
\begin{pgfscope}%
\pgfsys@transformshift{4.649399in}{4.856556in}%
\pgfsys@useobject{currentmarker}{}%
\end{pgfscope}%
\end{pgfscope}%
\begin{pgfscope}%
\pgfpathrectangle{\pgfqpoint{0.100000in}{2.493063in}}{\pgfqpoint{5.037500in}{3.427208in}}%
\pgfusepath{clip}%
\pgfsetrectcap%
\pgfsetroundjoin%
\pgfsetlinewidth{1.505625pt}%
\definecolor{currentstroke}{rgb}{0.000000,0.000000,1.000000}%
\pgfsetstrokecolor{currentstroke}%
\pgfsetstrokeopacity{0.500000}%
\pgfsetdash{}{0pt}%
\pgfpathmoveto{\pgfqpoint{4.703787in}{4.913704in}}%
\pgfusepath{stroke}%
\end{pgfscope}%
\begin{pgfscope}%
\pgfpathrectangle{\pgfqpoint{0.100000in}{2.493063in}}{\pgfqpoint{5.037500in}{3.427208in}}%
\pgfusepath{clip}%
\pgfsetbuttcap%
\pgfsetroundjoin%
\definecolor{currentfill}{rgb}{0.000000,0.000000,1.000000}%
\pgfsetfillcolor{currentfill}%
\pgfsetfillopacity{0.500000}%
\pgfsetlinewidth{0.250937pt}%
\definecolor{currentstroke}{rgb}{0.000000,0.000000,0.000000}%
\pgfsetstrokecolor{currentstroke}%
\pgfsetstrokeopacity{0.500000}%
\pgfsetdash{}{0pt}%
\pgfsys@defobject{currentmarker}{\pgfqpoint{-0.008333in}{-0.008333in}}{\pgfqpoint{0.008333in}{0.008333in}}{%
\pgfpathmoveto{\pgfqpoint{0.000000in}{-0.008333in}}%
\pgfpathcurveto{\pgfqpoint{0.002210in}{-0.008333in}}{\pgfqpoint{0.004330in}{-0.007455in}}{\pgfqpoint{0.005893in}{-0.005893in}}%
\pgfpathcurveto{\pgfqpoint{0.007455in}{-0.004330in}}{\pgfqpoint{0.008333in}{-0.002210in}}{\pgfqpoint{0.008333in}{0.000000in}}%
\pgfpathcurveto{\pgfqpoint{0.008333in}{0.002210in}}{\pgfqpoint{0.007455in}{0.004330in}}{\pgfqpoint{0.005893in}{0.005893in}}%
\pgfpathcurveto{\pgfqpoint{0.004330in}{0.007455in}}{\pgfqpoint{0.002210in}{0.008333in}}{\pgfqpoint{0.000000in}{0.008333in}}%
\pgfpathcurveto{\pgfqpoint{-0.002210in}{0.008333in}}{\pgfqpoint{-0.004330in}{0.007455in}}{\pgfqpoint{-0.005893in}{0.005893in}}%
\pgfpathcurveto{\pgfqpoint{-0.007455in}{0.004330in}}{\pgfqpoint{-0.008333in}{0.002210in}}{\pgfqpoint{-0.008333in}{0.000000in}}%
\pgfpathcurveto{\pgfqpoint{-0.008333in}{-0.002210in}}{\pgfqpoint{-0.007455in}{-0.004330in}}{\pgfqpoint{-0.005893in}{-0.005893in}}%
\pgfpathcurveto{\pgfqpoint{-0.004330in}{-0.007455in}}{\pgfqpoint{-0.002210in}{-0.008333in}}{\pgfqpoint{0.000000in}{-0.008333in}}%
\pgfpathlineto{\pgfqpoint{0.000000in}{-0.008333in}}%
\pgfpathclose%
\pgfusepath{stroke,fill}%
}%
\begin{pgfscope}%
\pgfsys@transformshift{4.703787in}{4.913704in}%
\pgfsys@useobject{currentmarker}{}%
\end{pgfscope}%
\end{pgfscope}%
\begin{pgfscope}%
\pgfpathrectangle{\pgfqpoint{0.100000in}{2.493063in}}{\pgfqpoint{5.037500in}{3.427208in}}%
\pgfusepath{clip}%
\pgfsetrectcap%
\pgfsetroundjoin%
\pgfsetlinewidth{1.505625pt}%
\definecolor{currentstroke}{rgb}{0.000000,0.000000,1.000000}%
\pgfsetstrokecolor{currentstroke}%
\pgfsetstrokeopacity{0.500000}%
\pgfsetdash{}{0pt}%
\pgfpathmoveto{\pgfqpoint{4.696218in}{4.857739in}}%
\pgfusepath{stroke}%
\end{pgfscope}%
\begin{pgfscope}%
\pgfpathrectangle{\pgfqpoint{0.100000in}{2.493063in}}{\pgfqpoint{5.037500in}{3.427208in}}%
\pgfusepath{clip}%
\pgfsetbuttcap%
\pgfsetroundjoin%
\definecolor{currentfill}{rgb}{0.000000,0.000000,1.000000}%
\pgfsetfillcolor{currentfill}%
\pgfsetfillopacity{0.500000}%
\pgfsetlinewidth{0.250937pt}%
\definecolor{currentstroke}{rgb}{0.000000,0.000000,0.000000}%
\pgfsetstrokecolor{currentstroke}%
\pgfsetstrokeopacity{0.500000}%
\pgfsetdash{}{0pt}%
\pgfsys@defobject{currentmarker}{\pgfqpoint{-0.008333in}{-0.008333in}}{\pgfqpoint{0.008333in}{0.008333in}}{%
\pgfpathmoveto{\pgfqpoint{0.000000in}{-0.008333in}}%
\pgfpathcurveto{\pgfqpoint{0.002210in}{-0.008333in}}{\pgfqpoint{0.004330in}{-0.007455in}}{\pgfqpoint{0.005893in}{-0.005893in}}%
\pgfpathcurveto{\pgfqpoint{0.007455in}{-0.004330in}}{\pgfqpoint{0.008333in}{-0.002210in}}{\pgfqpoint{0.008333in}{0.000000in}}%
\pgfpathcurveto{\pgfqpoint{0.008333in}{0.002210in}}{\pgfqpoint{0.007455in}{0.004330in}}{\pgfqpoint{0.005893in}{0.005893in}}%
\pgfpathcurveto{\pgfqpoint{0.004330in}{0.007455in}}{\pgfqpoint{0.002210in}{0.008333in}}{\pgfqpoint{0.000000in}{0.008333in}}%
\pgfpathcurveto{\pgfqpoint{-0.002210in}{0.008333in}}{\pgfqpoint{-0.004330in}{0.007455in}}{\pgfqpoint{-0.005893in}{0.005893in}}%
\pgfpathcurveto{\pgfqpoint{-0.007455in}{0.004330in}}{\pgfqpoint{-0.008333in}{0.002210in}}{\pgfqpoint{-0.008333in}{0.000000in}}%
\pgfpathcurveto{\pgfqpoint{-0.008333in}{-0.002210in}}{\pgfqpoint{-0.007455in}{-0.004330in}}{\pgfqpoint{-0.005893in}{-0.005893in}}%
\pgfpathcurveto{\pgfqpoint{-0.004330in}{-0.007455in}}{\pgfqpoint{-0.002210in}{-0.008333in}}{\pgfqpoint{0.000000in}{-0.008333in}}%
\pgfpathlineto{\pgfqpoint{0.000000in}{-0.008333in}}%
\pgfpathclose%
\pgfusepath{stroke,fill}%
}%
\begin{pgfscope}%
\pgfsys@transformshift{4.696218in}{4.857739in}%
\pgfsys@useobject{currentmarker}{}%
\end{pgfscope}%
\end{pgfscope}%
\begin{pgfscope}%
\pgfpathrectangle{\pgfqpoint{0.100000in}{2.493063in}}{\pgfqpoint{5.037500in}{3.427208in}}%
\pgfusepath{clip}%
\pgfsetrectcap%
\pgfsetroundjoin%
\pgfsetlinewidth{1.505625pt}%
\definecolor{currentstroke}{rgb}{0.501961,0.501961,0.501961}%
\pgfsetstrokecolor{currentstroke}%
\pgfsetstrokeopacity{0.500000}%
\pgfsetdash{}{0pt}%
\pgfpathmoveto{\pgfqpoint{4.761265in}{4.899862in}}%
\pgfusepath{stroke}%
\end{pgfscope}%
\begin{pgfscope}%
\pgfpathrectangle{\pgfqpoint{0.100000in}{2.493063in}}{\pgfqpoint{5.037500in}{3.427208in}}%
\pgfusepath{clip}%
\pgfsetbuttcap%
\pgfsetroundjoin%
\definecolor{currentfill}{rgb}{0.501961,0.501961,0.501961}%
\pgfsetfillcolor{currentfill}%
\pgfsetfillopacity{0.500000}%
\pgfsetlinewidth{0.250937pt}%
\definecolor{currentstroke}{rgb}{0.000000,0.000000,0.000000}%
\pgfsetstrokecolor{currentstroke}%
\pgfsetstrokeopacity{0.500000}%
\pgfsetdash{}{0pt}%
\pgfsys@defobject{currentmarker}{\pgfqpoint{-0.013889in}{-0.013889in}}{\pgfqpoint{0.013889in}{0.013889in}}{%
\pgfpathmoveto{\pgfqpoint{0.000000in}{-0.013889in}}%
\pgfpathcurveto{\pgfqpoint{0.003683in}{-0.013889in}}{\pgfqpoint{0.007216in}{-0.012425in}}{\pgfqpoint{0.009821in}{-0.009821in}}%
\pgfpathcurveto{\pgfqpoint{0.012425in}{-0.007216in}}{\pgfqpoint{0.013889in}{-0.003683in}}{\pgfqpoint{0.013889in}{0.000000in}}%
\pgfpathcurveto{\pgfqpoint{0.013889in}{0.003683in}}{\pgfqpoint{0.012425in}{0.007216in}}{\pgfqpoint{0.009821in}{0.009821in}}%
\pgfpathcurveto{\pgfqpoint{0.007216in}{0.012425in}}{\pgfqpoint{0.003683in}{0.013889in}}{\pgfqpoint{0.000000in}{0.013889in}}%
\pgfpathcurveto{\pgfqpoint{-0.003683in}{0.013889in}}{\pgfqpoint{-0.007216in}{0.012425in}}{\pgfqpoint{-0.009821in}{0.009821in}}%
\pgfpathcurveto{\pgfqpoint{-0.012425in}{0.007216in}}{\pgfqpoint{-0.013889in}{0.003683in}}{\pgfqpoint{-0.013889in}{0.000000in}}%
\pgfpathcurveto{\pgfqpoint{-0.013889in}{-0.003683in}}{\pgfqpoint{-0.012425in}{-0.007216in}}{\pgfqpoint{-0.009821in}{-0.009821in}}%
\pgfpathcurveto{\pgfqpoint{-0.007216in}{-0.012425in}}{\pgfqpoint{-0.003683in}{-0.013889in}}{\pgfqpoint{0.000000in}{-0.013889in}}%
\pgfpathlineto{\pgfqpoint{0.000000in}{-0.013889in}}%
\pgfpathclose%
\pgfusepath{stroke,fill}%
}%
\begin{pgfscope}%
\pgfsys@transformshift{4.761265in}{4.899862in}%
\pgfsys@useobject{currentmarker}{}%
\end{pgfscope}%
\end{pgfscope}%
\begin{pgfscope}%
\pgfpathrectangle{\pgfqpoint{0.100000in}{2.493063in}}{\pgfqpoint{5.037500in}{3.427208in}}%
\pgfusepath{clip}%
\pgfsetrectcap%
\pgfsetroundjoin%
\pgfsetlinewidth{1.505625pt}%
\definecolor{currentstroke}{rgb}{0.000000,0.000000,1.000000}%
\pgfsetstrokecolor{currentstroke}%
\pgfsetstrokeopacity{0.500000}%
\pgfsetdash{}{0pt}%
\pgfpathmoveto{\pgfqpoint{4.678703in}{4.883082in}}%
\pgfusepath{stroke}%
\end{pgfscope}%
\begin{pgfscope}%
\pgfpathrectangle{\pgfqpoint{0.100000in}{2.493063in}}{\pgfqpoint{5.037500in}{3.427208in}}%
\pgfusepath{clip}%
\pgfsetbuttcap%
\pgfsetroundjoin%
\definecolor{currentfill}{rgb}{0.000000,0.000000,1.000000}%
\pgfsetfillcolor{currentfill}%
\pgfsetfillopacity{0.500000}%
\pgfsetlinewidth{0.250937pt}%
\definecolor{currentstroke}{rgb}{0.000000,0.000000,0.000000}%
\pgfsetstrokecolor{currentstroke}%
\pgfsetstrokeopacity{0.500000}%
\pgfsetdash{}{0pt}%
\pgfsys@defobject{currentmarker}{\pgfqpoint{-0.008333in}{-0.008333in}}{\pgfqpoint{0.008333in}{0.008333in}}{%
\pgfpathmoveto{\pgfqpoint{0.000000in}{-0.008333in}}%
\pgfpathcurveto{\pgfqpoint{0.002210in}{-0.008333in}}{\pgfqpoint{0.004330in}{-0.007455in}}{\pgfqpoint{0.005893in}{-0.005893in}}%
\pgfpathcurveto{\pgfqpoint{0.007455in}{-0.004330in}}{\pgfqpoint{0.008333in}{-0.002210in}}{\pgfqpoint{0.008333in}{0.000000in}}%
\pgfpathcurveto{\pgfqpoint{0.008333in}{0.002210in}}{\pgfqpoint{0.007455in}{0.004330in}}{\pgfqpoint{0.005893in}{0.005893in}}%
\pgfpathcurveto{\pgfqpoint{0.004330in}{0.007455in}}{\pgfqpoint{0.002210in}{0.008333in}}{\pgfqpoint{0.000000in}{0.008333in}}%
\pgfpathcurveto{\pgfqpoint{-0.002210in}{0.008333in}}{\pgfqpoint{-0.004330in}{0.007455in}}{\pgfqpoint{-0.005893in}{0.005893in}}%
\pgfpathcurveto{\pgfqpoint{-0.007455in}{0.004330in}}{\pgfqpoint{-0.008333in}{0.002210in}}{\pgfqpoint{-0.008333in}{0.000000in}}%
\pgfpathcurveto{\pgfqpoint{-0.008333in}{-0.002210in}}{\pgfqpoint{-0.007455in}{-0.004330in}}{\pgfqpoint{-0.005893in}{-0.005893in}}%
\pgfpathcurveto{\pgfqpoint{-0.004330in}{-0.007455in}}{\pgfqpoint{-0.002210in}{-0.008333in}}{\pgfqpoint{0.000000in}{-0.008333in}}%
\pgfpathlineto{\pgfqpoint{0.000000in}{-0.008333in}}%
\pgfpathclose%
\pgfusepath{stroke,fill}%
}%
\begin{pgfscope}%
\pgfsys@transformshift{4.678703in}{4.883082in}%
\pgfsys@useobject{currentmarker}{}%
\end{pgfscope}%
\end{pgfscope}%
\begin{pgfscope}%
\pgfpathrectangle{\pgfqpoint{0.100000in}{2.493063in}}{\pgfqpoint{5.037500in}{3.427208in}}%
\pgfusepath{clip}%
\pgfsetrectcap%
\pgfsetroundjoin%
\pgfsetlinewidth{1.505625pt}%
\definecolor{currentstroke}{rgb}{0.000000,0.000000,1.000000}%
\pgfsetstrokecolor{currentstroke}%
\pgfsetstrokeopacity{0.500000}%
\pgfsetdash{}{0pt}%
\pgfpathmoveto{\pgfqpoint{4.529746in}{4.564780in}}%
\pgfusepath{stroke}%
\end{pgfscope}%
\begin{pgfscope}%
\pgfpathrectangle{\pgfqpoint{0.100000in}{2.493063in}}{\pgfqpoint{5.037500in}{3.427208in}}%
\pgfusepath{clip}%
\pgfsetbuttcap%
\pgfsetroundjoin%
\definecolor{currentfill}{rgb}{0.000000,0.000000,1.000000}%
\pgfsetfillcolor{currentfill}%
\pgfsetfillopacity{0.500000}%
\pgfsetlinewidth{0.250937pt}%
\definecolor{currentstroke}{rgb}{0.000000,0.000000,0.000000}%
\pgfsetstrokecolor{currentstroke}%
\pgfsetstrokeopacity{0.500000}%
\pgfsetdash{}{0pt}%
\pgfsys@defobject{currentmarker}{\pgfqpoint{-0.022222in}{-0.022222in}}{\pgfqpoint{0.022222in}{0.022222in}}{%
\pgfpathmoveto{\pgfqpoint{0.000000in}{-0.022222in}}%
\pgfpathcurveto{\pgfqpoint{0.005893in}{-0.022222in}}{\pgfqpoint{0.011546in}{-0.019881in}}{\pgfqpoint{0.015713in}{-0.015713in}}%
\pgfpathcurveto{\pgfqpoint{0.019881in}{-0.011546in}}{\pgfqpoint{0.022222in}{-0.005893in}}{\pgfqpoint{0.022222in}{0.000000in}}%
\pgfpathcurveto{\pgfqpoint{0.022222in}{0.005893in}}{\pgfqpoint{0.019881in}{0.011546in}}{\pgfqpoint{0.015713in}{0.015713in}}%
\pgfpathcurveto{\pgfqpoint{0.011546in}{0.019881in}}{\pgfqpoint{0.005893in}{0.022222in}}{\pgfqpoint{0.000000in}{0.022222in}}%
\pgfpathcurveto{\pgfqpoint{-0.005893in}{0.022222in}}{\pgfqpoint{-0.011546in}{0.019881in}}{\pgfqpoint{-0.015713in}{0.015713in}}%
\pgfpathcurveto{\pgfqpoint{-0.019881in}{0.011546in}}{\pgfqpoint{-0.022222in}{0.005893in}}{\pgfqpoint{-0.022222in}{0.000000in}}%
\pgfpathcurveto{\pgfqpoint{-0.022222in}{-0.005893in}}{\pgfqpoint{-0.019881in}{-0.011546in}}{\pgfqpoint{-0.015713in}{-0.015713in}}%
\pgfpathcurveto{\pgfqpoint{-0.011546in}{-0.019881in}}{\pgfqpoint{-0.005893in}{-0.022222in}}{\pgfqpoint{0.000000in}{-0.022222in}}%
\pgfpathlineto{\pgfqpoint{0.000000in}{-0.022222in}}%
\pgfpathclose%
\pgfusepath{stroke,fill}%
}%
\begin{pgfscope}%
\pgfsys@transformshift{4.529746in}{4.564780in}%
\pgfsys@useobject{currentmarker}{}%
\end{pgfscope}%
\end{pgfscope}%
\begin{pgfscope}%
\pgfpathrectangle{\pgfqpoint{0.100000in}{2.493063in}}{\pgfqpoint{5.037500in}{3.427208in}}%
\pgfusepath{clip}%
\pgfsetrectcap%
\pgfsetroundjoin%
\pgfsetlinewidth{1.505625pt}%
\definecolor{currentstroke}{rgb}{0.678431,1.000000,0.184314}%
\pgfsetstrokecolor{currentstroke}%
\pgfsetstrokeopacity{0.500000}%
\pgfsetdash{}{0pt}%
\pgfpathmoveto{\pgfqpoint{4.542388in}{4.474206in}}%
\pgfusepath{stroke}%
\end{pgfscope}%
\begin{pgfscope}%
\pgfpathrectangle{\pgfqpoint{0.100000in}{2.493063in}}{\pgfqpoint{5.037500in}{3.427208in}}%
\pgfusepath{clip}%
\pgfsetbuttcap%
\pgfsetroundjoin%
\definecolor{currentfill}{rgb}{0.678431,1.000000,0.184314}%
\pgfsetfillcolor{currentfill}%
\pgfsetfillopacity{0.500000}%
\pgfsetlinewidth{0.250937pt}%
\definecolor{currentstroke}{rgb}{0.000000,0.000000,0.000000}%
\pgfsetstrokecolor{currentstroke}%
\pgfsetstrokeopacity{0.500000}%
\pgfsetdash{}{0pt}%
\pgfsys@defobject{currentmarker}{\pgfqpoint{-0.027778in}{-0.027778in}}{\pgfqpoint{0.027778in}{0.027778in}}{%
\pgfpathmoveto{\pgfqpoint{0.000000in}{-0.027778in}}%
\pgfpathcurveto{\pgfqpoint{0.007367in}{-0.027778in}}{\pgfqpoint{0.014433in}{-0.024851in}}{\pgfqpoint{0.019642in}{-0.019642in}}%
\pgfpathcurveto{\pgfqpoint{0.024851in}{-0.014433in}}{\pgfqpoint{0.027778in}{-0.007367in}}{\pgfqpoint{0.027778in}{0.000000in}}%
\pgfpathcurveto{\pgfqpoint{0.027778in}{0.007367in}}{\pgfqpoint{0.024851in}{0.014433in}}{\pgfqpoint{0.019642in}{0.019642in}}%
\pgfpathcurveto{\pgfqpoint{0.014433in}{0.024851in}}{\pgfqpoint{0.007367in}{0.027778in}}{\pgfqpoint{0.000000in}{0.027778in}}%
\pgfpathcurveto{\pgfqpoint{-0.007367in}{0.027778in}}{\pgfqpoint{-0.014433in}{0.024851in}}{\pgfqpoint{-0.019642in}{0.019642in}}%
\pgfpathcurveto{\pgfqpoint{-0.024851in}{0.014433in}}{\pgfqpoint{-0.027778in}{0.007367in}}{\pgfqpoint{-0.027778in}{0.000000in}}%
\pgfpathcurveto{\pgfqpoint{-0.027778in}{-0.007367in}}{\pgfqpoint{-0.024851in}{-0.014433in}}{\pgfqpoint{-0.019642in}{-0.019642in}}%
\pgfpathcurveto{\pgfqpoint{-0.014433in}{-0.024851in}}{\pgfqpoint{-0.007367in}{-0.027778in}}{\pgfqpoint{0.000000in}{-0.027778in}}%
\pgfpathlineto{\pgfqpoint{0.000000in}{-0.027778in}}%
\pgfpathclose%
\pgfusepath{stroke,fill}%
}%
\begin{pgfscope}%
\pgfsys@transformshift{4.542388in}{4.474206in}%
\pgfsys@useobject{currentmarker}{}%
\end{pgfscope}%
\end{pgfscope}%
\begin{pgfscope}%
\pgfpathrectangle{\pgfqpoint{0.100000in}{2.493063in}}{\pgfqpoint{5.037500in}{3.427208in}}%
\pgfusepath{clip}%
\pgfsetrectcap%
\pgfsetroundjoin%
\pgfsetlinewidth{1.505625pt}%
\definecolor{currentstroke}{rgb}{0.678431,1.000000,0.184314}%
\pgfsetstrokecolor{currentstroke}%
\pgfsetstrokeopacity{0.500000}%
\pgfsetdash{}{0pt}%
\pgfpathmoveto{\pgfqpoint{4.403529in}{4.507445in}}%
\pgfusepath{stroke}%
\end{pgfscope}%
\begin{pgfscope}%
\pgfpathrectangle{\pgfqpoint{0.100000in}{2.493063in}}{\pgfqpoint{5.037500in}{3.427208in}}%
\pgfusepath{clip}%
\pgfsetbuttcap%
\pgfsetroundjoin%
\definecolor{currentfill}{rgb}{0.678431,1.000000,0.184314}%
\pgfsetfillcolor{currentfill}%
\pgfsetfillopacity{0.500000}%
\pgfsetlinewidth{0.250937pt}%
\definecolor{currentstroke}{rgb}{0.000000,0.000000,0.000000}%
\pgfsetstrokecolor{currentstroke}%
\pgfsetstrokeopacity{0.500000}%
\pgfsetdash{}{0pt}%
\pgfsys@defobject{currentmarker}{\pgfqpoint{-0.005556in}{-0.005556in}}{\pgfqpoint{0.005556in}{0.005556in}}{%
\pgfpathmoveto{\pgfqpoint{0.000000in}{-0.005556in}}%
\pgfpathcurveto{\pgfqpoint{0.001473in}{-0.005556in}}{\pgfqpoint{0.002887in}{-0.004970in}}{\pgfqpoint{0.003928in}{-0.003928in}}%
\pgfpathcurveto{\pgfqpoint{0.004970in}{-0.002887in}}{\pgfqpoint{0.005556in}{-0.001473in}}{\pgfqpoint{0.005556in}{0.000000in}}%
\pgfpathcurveto{\pgfqpoint{0.005556in}{0.001473in}}{\pgfqpoint{0.004970in}{0.002887in}}{\pgfqpoint{0.003928in}{0.003928in}}%
\pgfpathcurveto{\pgfqpoint{0.002887in}{0.004970in}}{\pgfqpoint{0.001473in}{0.005556in}}{\pgfqpoint{0.000000in}{0.005556in}}%
\pgfpathcurveto{\pgfqpoint{-0.001473in}{0.005556in}}{\pgfqpoint{-0.002887in}{0.004970in}}{\pgfqpoint{-0.003928in}{0.003928in}}%
\pgfpathcurveto{\pgfqpoint{-0.004970in}{0.002887in}}{\pgfqpoint{-0.005556in}{0.001473in}}{\pgfqpoint{-0.005556in}{0.000000in}}%
\pgfpathcurveto{\pgfqpoint{-0.005556in}{-0.001473in}}{\pgfqpoint{-0.004970in}{-0.002887in}}{\pgfqpoint{-0.003928in}{-0.003928in}}%
\pgfpathcurveto{\pgfqpoint{-0.002887in}{-0.004970in}}{\pgfqpoint{-0.001473in}{-0.005556in}}{\pgfqpoint{0.000000in}{-0.005556in}}%
\pgfpathlineto{\pgfqpoint{0.000000in}{-0.005556in}}%
\pgfpathclose%
\pgfusepath{stroke,fill}%
}%
\begin{pgfscope}%
\pgfsys@transformshift{4.403529in}{4.507445in}%
\pgfsys@useobject{currentmarker}{}%
\end{pgfscope}%
\end{pgfscope}%
\begin{pgfscope}%
\pgfpathrectangle{\pgfqpoint{0.100000in}{2.493063in}}{\pgfqpoint{5.037500in}{3.427208in}}%
\pgfusepath{clip}%
\pgfsetrectcap%
\pgfsetroundjoin%
\pgfsetlinewidth{1.505625pt}%
\definecolor{currentstroke}{rgb}{0.000000,0.000000,1.000000}%
\pgfsetstrokecolor{currentstroke}%
\pgfsetstrokeopacity{0.500000}%
\pgfsetdash{}{0pt}%
\pgfpathmoveto{\pgfqpoint{4.170525in}{3.020820in}}%
\pgfusepath{stroke}%
\end{pgfscope}%
\begin{pgfscope}%
\pgfpathrectangle{\pgfqpoint{0.100000in}{2.493063in}}{\pgfqpoint{5.037500in}{3.427208in}}%
\pgfusepath{clip}%
\pgfsetbuttcap%
\pgfsetroundjoin%
\definecolor{currentfill}{rgb}{0.000000,0.000000,1.000000}%
\pgfsetfillcolor{currentfill}%
\pgfsetfillopacity{0.500000}%
\pgfsetlinewidth{0.250937pt}%
\definecolor{currentstroke}{rgb}{0.000000,0.000000,0.000000}%
\pgfsetstrokecolor{currentstroke}%
\pgfsetstrokeopacity{0.500000}%
\pgfsetdash{}{0pt}%
\pgfsys@defobject{currentmarker}{\pgfqpoint{-0.005556in}{-0.005556in}}{\pgfqpoint{0.005556in}{0.005556in}}{%
\pgfpathmoveto{\pgfqpoint{0.000000in}{-0.005556in}}%
\pgfpathcurveto{\pgfqpoint{0.001473in}{-0.005556in}}{\pgfqpoint{0.002887in}{-0.004970in}}{\pgfqpoint{0.003928in}{-0.003928in}}%
\pgfpathcurveto{\pgfqpoint{0.004970in}{-0.002887in}}{\pgfqpoint{0.005556in}{-0.001473in}}{\pgfqpoint{0.005556in}{0.000000in}}%
\pgfpathcurveto{\pgfqpoint{0.005556in}{0.001473in}}{\pgfqpoint{0.004970in}{0.002887in}}{\pgfqpoint{0.003928in}{0.003928in}}%
\pgfpathcurveto{\pgfqpoint{0.002887in}{0.004970in}}{\pgfqpoint{0.001473in}{0.005556in}}{\pgfqpoint{0.000000in}{0.005556in}}%
\pgfpathcurveto{\pgfqpoint{-0.001473in}{0.005556in}}{\pgfqpoint{-0.002887in}{0.004970in}}{\pgfqpoint{-0.003928in}{0.003928in}}%
\pgfpathcurveto{\pgfqpoint{-0.004970in}{0.002887in}}{\pgfqpoint{-0.005556in}{0.001473in}}{\pgfqpoint{-0.005556in}{0.000000in}}%
\pgfpathcurveto{\pgfqpoint{-0.005556in}{-0.001473in}}{\pgfqpoint{-0.004970in}{-0.002887in}}{\pgfqpoint{-0.003928in}{-0.003928in}}%
\pgfpathcurveto{\pgfqpoint{-0.002887in}{-0.004970in}}{\pgfqpoint{-0.001473in}{-0.005556in}}{\pgfqpoint{0.000000in}{-0.005556in}}%
\pgfpathlineto{\pgfqpoint{0.000000in}{-0.005556in}}%
\pgfpathclose%
\pgfusepath{stroke,fill}%
}%
\begin{pgfscope}%
\pgfsys@transformshift{4.170525in}{3.020820in}%
\pgfsys@useobject{currentmarker}{}%
\end{pgfscope}%
\end{pgfscope}%
\begin{pgfscope}%
\pgfpathrectangle{\pgfqpoint{0.100000in}{2.493063in}}{\pgfqpoint{5.037500in}{3.427208in}}%
\pgfusepath{clip}%
\pgfsetrectcap%
\pgfsetroundjoin%
\pgfsetlinewidth{1.505625pt}%
\definecolor{currentstroke}{rgb}{0.501961,0.501961,0.501961}%
\pgfsetstrokecolor{currentstroke}%
\pgfsetstrokeopacity{0.500000}%
\pgfsetdash{}{0pt}%
\pgfpathmoveto{\pgfqpoint{3.645536in}{3.448139in}}%
\pgfusepath{stroke}%
\end{pgfscope}%
\begin{pgfscope}%
\pgfpathrectangle{\pgfqpoint{0.100000in}{2.493063in}}{\pgfqpoint{5.037500in}{3.427208in}}%
\pgfusepath{clip}%
\pgfsetbuttcap%
\pgfsetroundjoin%
\definecolor{currentfill}{rgb}{0.501961,0.501961,0.501961}%
\pgfsetfillcolor{currentfill}%
\pgfsetfillopacity{0.500000}%
\pgfsetlinewidth{0.250937pt}%
\definecolor{currentstroke}{rgb}{0.000000,0.000000,0.000000}%
\pgfsetstrokecolor{currentstroke}%
\pgfsetstrokeopacity{0.500000}%
\pgfsetdash{}{0pt}%
\pgfsys@defobject{currentmarker}{\pgfqpoint{-0.013889in}{-0.013889in}}{\pgfqpoint{0.013889in}{0.013889in}}{%
\pgfpathmoveto{\pgfqpoint{0.000000in}{-0.013889in}}%
\pgfpathcurveto{\pgfqpoint{0.003683in}{-0.013889in}}{\pgfqpoint{0.007216in}{-0.012425in}}{\pgfqpoint{0.009821in}{-0.009821in}}%
\pgfpathcurveto{\pgfqpoint{0.012425in}{-0.007216in}}{\pgfqpoint{0.013889in}{-0.003683in}}{\pgfqpoint{0.013889in}{0.000000in}}%
\pgfpathcurveto{\pgfqpoint{0.013889in}{0.003683in}}{\pgfqpoint{0.012425in}{0.007216in}}{\pgfqpoint{0.009821in}{0.009821in}}%
\pgfpathcurveto{\pgfqpoint{0.007216in}{0.012425in}}{\pgfqpoint{0.003683in}{0.013889in}}{\pgfqpoint{0.000000in}{0.013889in}}%
\pgfpathcurveto{\pgfqpoint{-0.003683in}{0.013889in}}{\pgfqpoint{-0.007216in}{0.012425in}}{\pgfqpoint{-0.009821in}{0.009821in}}%
\pgfpathcurveto{\pgfqpoint{-0.012425in}{0.007216in}}{\pgfqpoint{-0.013889in}{0.003683in}}{\pgfqpoint{-0.013889in}{0.000000in}}%
\pgfpathcurveto{\pgfqpoint{-0.013889in}{-0.003683in}}{\pgfqpoint{-0.012425in}{-0.007216in}}{\pgfqpoint{-0.009821in}{-0.009821in}}%
\pgfpathcurveto{\pgfqpoint{-0.007216in}{-0.012425in}}{\pgfqpoint{-0.003683in}{-0.013889in}}{\pgfqpoint{0.000000in}{-0.013889in}}%
\pgfpathlineto{\pgfqpoint{0.000000in}{-0.013889in}}%
\pgfpathclose%
\pgfusepath{stroke,fill}%
}%
\begin{pgfscope}%
\pgfsys@transformshift{3.645536in}{3.448139in}%
\pgfsys@useobject{currentmarker}{}%
\end{pgfscope}%
\end{pgfscope}%
\begin{pgfscope}%
\pgfpathrectangle{\pgfqpoint{0.100000in}{2.493063in}}{\pgfqpoint{5.037500in}{3.427208in}}%
\pgfusepath{clip}%
\pgfsetrectcap%
\pgfsetroundjoin%
\pgfsetlinewidth{1.505625pt}%
\definecolor{currentstroke}{rgb}{0.501961,0.501961,0.501961}%
\pgfsetstrokecolor{currentstroke}%
\pgfsetstrokeopacity{0.500000}%
\pgfsetdash{}{0pt}%
\pgfpathmoveto{\pgfqpoint{3.644545in}{3.406395in}}%
\pgfusepath{stroke}%
\end{pgfscope}%
\begin{pgfscope}%
\pgfpathrectangle{\pgfqpoint{0.100000in}{2.493063in}}{\pgfqpoint{5.037500in}{3.427208in}}%
\pgfusepath{clip}%
\pgfsetbuttcap%
\pgfsetroundjoin%
\definecolor{currentfill}{rgb}{0.501961,0.501961,0.501961}%
\pgfsetfillcolor{currentfill}%
\pgfsetfillopacity{0.500000}%
\pgfsetlinewidth{0.250937pt}%
\definecolor{currentstroke}{rgb}{0.000000,0.000000,0.000000}%
\pgfsetstrokecolor{currentstroke}%
\pgfsetstrokeopacity{0.500000}%
\pgfsetdash{}{0pt}%
\pgfsys@defobject{currentmarker}{\pgfqpoint{-0.013889in}{-0.013889in}}{\pgfqpoint{0.013889in}{0.013889in}}{%
\pgfpathmoveto{\pgfqpoint{0.000000in}{-0.013889in}}%
\pgfpathcurveto{\pgfqpoint{0.003683in}{-0.013889in}}{\pgfqpoint{0.007216in}{-0.012425in}}{\pgfqpoint{0.009821in}{-0.009821in}}%
\pgfpathcurveto{\pgfqpoint{0.012425in}{-0.007216in}}{\pgfqpoint{0.013889in}{-0.003683in}}{\pgfqpoint{0.013889in}{0.000000in}}%
\pgfpathcurveto{\pgfqpoint{0.013889in}{0.003683in}}{\pgfqpoint{0.012425in}{0.007216in}}{\pgfqpoint{0.009821in}{0.009821in}}%
\pgfpathcurveto{\pgfqpoint{0.007216in}{0.012425in}}{\pgfqpoint{0.003683in}{0.013889in}}{\pgfqpoint{0.000000in}{0.013889in}}%
\pgfpathcurveto{\pgfqpoint{-0.003683in}{0.013889in}}{\pgfqpoint{-0.007216in}{0.012425in}}{\pgfqpoint{-0.009821in}{0.009821in}}%
\pgfpathcurveto{\pgfqpoint{-0.012425in}{0.007216in}}{\pgfqpoint{-0.013889in}{0.003683in}}{\pgfqpoint{-0.013889in}{0.000000in}}%
\pgfpathcurveto{\pgfqpoint{-0.013889in}{-0.003683in}}{\pgfqpoint{-0.012425in}{-0.007216in}}{\pgfqpoint{-0.009821in}{-0.009821in}}%
\pgfpathcurveto{\pgfqpoint{-0.007216in}{-0.012425in}}{\pgfqpoint{-0.003683in}{-0.013889in}}{\pgfqpoint{0.000000in}{-0.013889in}}%
\pgfpathlineto{\pgfqpoint{0.000000in}{-0.013889in}}%
\pgfpathclose%
\pgfusepath{stroke,fill}%
}%
\begin{pgfscope}%
\pgfsys@transformshift{3.644545in}{3.406395in}%
\pgfsys@useobject{currentmarker}{}%
\end{pgfscope}%
\end{pgfscope}%
\begin{pgfscope}%
\pgfpathrectangle{\pgfqpoint{0.100000in}{2.493063in}}{\pgfqpoint{5.037500in}{3.427208in}}%
\pgfusepath{clip}%
\pgfsetrectcap%
\pgfsetroundjoin%
\pgfsetlinewidth{1.505625pt}%
\definecolor{currentstroke}{rgb}{0.501961,0.501961,0.501961}%
\pgfsetstrokecolor{currentstroke}%
\pgfsetstrokeopacity{0.500000}%
\pgfsetdash{}{0pt}%
\pgfpathmoveto{\pgfqpoint{4.204682in}{3.298061in}}%
\pgfusepath{stroke}%
\end{pgfscope}%
\begin{pgfscope}%
\pgfpathrectangle{\pgfqpoint{0.100000in}{2.493063in}}{\pgfqpoint{5.037500in}{3.427208in}}%
\pgfusepath{clip}%
\pgfsetbuttcap%
\pgfsetroundjoin%
\definecolor{currentfill}{rgb}{0.501961,0.501961,0.501961}%
\pgfsetfillcolor{currentfill}%
\pgfsetfillopacity{0.500000}%
\pgfsetlinewidth{0.250937pt}%
\definecolor{currentstroke}{rgb}{0.000000,0.000000,0.000000}%
\pgfsetstrokecolor{currentstroke}%
\pgfsetstrokeopacity{0.500000}%
\pgfsetdash{}{0pt}%
\pgfsys@defobject{currentmarker}{\pgfqpoint{-0.013889in}{-0.013889in}}{\pgfqpoint{0.013889in}{0.013889in}}{%
\pgfpathmoveto{\pgfqpoint{0.000000in}{-0.013889in}}%
\pgfpathcurveto{\pgfqpoint{0.003683in}{-0.013889in}}{\pgfqpoint{0.007216in}{-0.012425in}}{\pgfqpoint{0.009821in}{-0.009821in}}%
\pgfpathcurveto{\pgfqpoint{0.012425in}{-0.007216in}}{\pgfqpoint{0.013889in}{-0.003683in}}{\pgfqpoint{0.013889in}{0.000000in}}%
\pgfpathcurveto{\pgfqpoint{0.013889in}{0.003683in}}{\pgfqpoint{0.012425in}{0.007216in}}{\pgfqpoint{0.009821in}{0.009821in}}%
\pgfpathcurveto{\pgfqpoint{0.007216in}{0.012425in}}{\pgfqpoint{0.003683in}{0.013889in}}{\pgfqpoint{0.000000in}{0.013889in}}%
\pgfpathcurveto{\pgfqpoint{-0.003683in}{0.013889in}}{\pgfqpoint{-0.007216in}{0.012425in}}{\pgfqpoint{-0.009821in}{0.009821in}}%
\pgfpathcurveto{\pgfqpoint{-0.012425in}{0.007216in}}{\pgfqpoint{-0.013889in}{0.003683in}}{\pgfqpoint{-0.013889in}{0.000000in}}%
\pgfpathcurveto{\pgfqpoint{-0.013889in}{-0.003683in}}{\pgfqpoint{-0.012425in}{-0.007216in}}{\pgfqpoint{-0.009821in}{-0.009821in}}%
\pgfpathcurveto{\pgfqpoint{-0.007216in}{-0.012425in}}{\pgfqpoint{-0.003683in}{-0.013889in}}{\pgfqpoint{0.000000in}{-0.013889in}}%
\pgfpathlineto{\pgfqpoint{0.000000in}{-0.013889in}}%
\pgfpathclose%
\pgfusepath{stroke,fill}%
}%
\begin{pgfscope}%
\pgfsys@transformshift{4.204682in}{3.298061in}%
\pgfsys@useobject{currentmarker}{}%
\end{pgfscope}%
\end{pgfscope}%
\begin{pgfscope}%
\pgfpathrectangle{\pgfqpoint{0.100000in}{2.493063in}}{\pgfqpoint{5.037500in}{3.427208in}}%
\pgfusepath{clip}%
\pgfsetrectcap%
\pgfsetroundjoin%
\pgfsetlinewidth{1.505625pt}%
\definecolor{currentstroke}{rgb}{0.000000,0.000000,1.000000}%
\pgfsetstrokecolor{currentstroke}%
\pgfsetstrokeopacity{0.500000}%
\pgfsetdash{}{0pt}%
\pgfpathmoveto{\pgfqpoint{4.085003in}{3.369230in}}%
\pgfusepath{stroke}%
\end{pgfscope}%
\begin{pgfscope}%
\pgfpathrectangle{\pgfqpoint{0.100000in}{2.493063in}}{\pgfqpoint{5.037500in}{3.427208in}}%
\pgfusepath{clip}%
\pgfsetbuttcap%
\pgfsetroundjoin%
\definecolor{currentfill}{rgb}{0.000000,0.000000,1.000000}%
\pgfsetfillcolor{currentfill}%
\pgfsetfillopacity{0.500000}%
\pgfsetlinewidth{0.250937pt}%
\definecolor{currentstroke}{rgb}{0.000000,0.000000,0.000000}%
\pgfsetstrokecolor{currentstroke}%
\pgfsetstrokeopacity{0.500000}%
\pgfsetdash{}{0pt}%
\pgfsys@defobject{currentmarker}{\pgfqpoint{-0.013889in}{-0.013889in}}{\pgfqpoint{0.013889in}{0.013889in}}{%
\pgfpathmoveto{\pgfqpoint{0.000000in}{-0.013889in}}%
\pgfpathcurveto{\pgfqpoint{0.003683in}{-0.013889in}}{\pgfqpoint{0.007216in}{-0.012425in}}{\pgfqpoint{0.009821in}{-0.009821in}}%
\pgfpathcurveto{\pgfqpoint{0.012425in}{-0.007216in}}{\pgfqpoint{0.013889in}{-0.003683in}}{\pgfqpoint{0.013889in}{0.000000in}}%
\pgfpathcurveto{\pgfqpoint{0.013889in}{0.003683in}}{\pgfqpoint{0.012425in}{0.007216in}}{\pgfqpoint{0.009821in}{0.009821in}}%
\pgfpathcurveto{\pgfqpoint{0.007216in}{0.012425in}}{\pgfqpoint{0.003683in}{0.013889in}}{\pgfqpoint{0.000000in}{0.013889in}}%
\pgfpathcurveto{\pgfqpoint{-0.003683in}{0.013889in}}{\pgfqpoint{-0.007216in}{0.012425in}}{\pgfqpoint{-0.009821in}{0.009821in}}%
\pgfpathcurveto{\pgfqpoint{-0.012425in}{0.007216in}}{\pgfqpoint{-0.013889in}{0.003683in}}{\pgfqpoint{-0.013889in}{0.000000in}}%
\pgfpathcurveto{\pgfqpoint{-0.013889in}{-0.003683in}}{\pgfqpoint{-0.012425in}{-0.007216in}}{\pgfqpoint{-0.009821in}{-0.009821in}}%
\pgfpathcurveto{\pgfqpoint{-0.007216in}{-0.012425in}}{\pgfqpoint{-0.003683in}{-0.013889in}}{\pgfqpoint{0.000000in}{-0.013889in}}%
\pgfpathlineto{\pgfqpoint{0.000000in}{-0.013889in}}%
\pgfpathclose%
\pgfusepath{stroke,fill}%
}%
\begin{pgfscope}%
\pgfsys@transformshift{4.085003in}{3.369230in}%
\pgfsys@useobject{currentmarker}{}%
\end{pgfscope}%
\end{pgfscope}%
\begin{pgfscope}%
\pgfpathrectangle{\pgfqpoint{0.100000in}{2.493063in}}{\pgfqpoint{5.037500in}{3.427208in}}%
\pgfusepath{clip}%
\pgfsetrectcap%
\pgfsetroundjoin%
\pgfsetlinewidth{1.505625pt}%
\definecolor{currentstroke}{rgb}{0.501961,0.501961,0.501961}%
\pgfsetstrokecolor{currentstroke}%
\pgfsetstrokeopacity{0.500000}%
\pgfsetdash{}{0pt}%
\pgfpathmoveto{\pgfqpoint{4.073308in}{3.267567in}}%
\pgfusepath{stroke}%
\end{pgfscope}%
\begin{pgfscope}%
\pgfpathrectangle{\pgfqpoint{0.100000in}{2.493063in}}{\pgfqpoint{5.037500in}{3.427208in}}%
\pgfusepath{clip}%
\pgfsetbuttcap%
\pgfsetroundjoin%
\definecolor{currentfill}{rgb}{0.501961,0.501961,0.501961}%
\pgfsetfillcolor{currentfill}%
\pgfsetfillopacity{0.500000}%
\pgfsetlinewidth{0.250937pt}%
\definecolor{currentstroke}{rgb}{0.000000,0.000000,0.000000}%
\pgfsetstrokecolor{currentstroke}%
\pgfsetstrokeopacity{0.500000}%
\pgfsetdash{}{0pt}%
\pgfsys@defobject{currentmarker}{\pgfqpoint{-0.013889in}{-0.013889in}}{\pgfqpoint{0.013889in}{0.013889in}}{%
\pgfpathmoveto{\pgfqpoint{0.000000in}{-0.013889in}}%
\pgfpathcurveto{\pgfqpoint{0.003683in}{-0.013889in}}{\pgfqpoint{0.007216in}{-0.012425in}}{\pgfqpoint{0.009821in}{-0.009821in}}%
\pgfpathcurveto{\pgfqpoint{0.012425in}{-0.007216in}}{\pgfqpoint{0.013889in}{-0.003683in}}{\pgfqpoint{0.013889in}{0.000000in}}%
\pgfpathcurveto{\pgfqpoint{0.013889in}{0.003683in}}{\pgfqpoint{0.012425in}{0.007216in}}{\pgfqpoint{0.009821in}{0.009821in}}%
\pgfpathcurveto{\pgfqpoint{0.007216in}{0.012425in}}{\pgfqpoint{0.003683in}{0.013889in}}{\pgfqpoint{0.000000in}{0.013889in}}%
\pgfpathcurveto{\pgfqpoint{-0.003683in}{0.013889in}}{\pgfqpoint{-0.007216in}{0.012425in}}{\pgfqpoint{-0.009821in}{0.009821in}}%
\pgfpathcurveto{\pgfqpoint{-0.012425in}{0.007216in}}{\pgfqpoint{-0.013889in}{0.003683in}}{\pgfqpoint{-0.013889in}{0.000000in}}%
\pgfpathcurveto{\pgfqpoint{-0.013889in}{-0.003683in}}{\pgfqpoint{-0.012425in}{-0.007216in}}{\pgfqpoint{-0.009821in}{-0.009821in}}%
\pgfpathcurveto{\pgfqpoint{-0.007216in}{-0.012425in}}{\pgfqpoint{-0.003683in}{-0.013889in}}{\pgfqpoint{0.000000in}{-0.013889in}}%
\pgfpathlineto{\pgfqpoint{0.000000in}{-0.013889in}}%
\pgfpathclose%
\pgfusepath{stroke,fill}%
}%
\begin{pgfscope}%
\pgfsys@transformshift{4.073308in}{3.267567in}%
\pgfsys@useobject{currentmarker}{}%
\end{pgfscope}%
\end{pgfscope}%
\begin{pgfscope}%
\pgfpathrectangle{\pgfqpoint{0.100000in}{2.493063in}}{\pgfqpoint{5.037500in}{3.427208in}}%
\pgfusepath{clip}%
\pgfsetrectcap%
\pgfsetroundjoin%
\pgfsetlinewidth{1.505625pt}%
\definecolor{currentstroke}{rgb}{0.501961,0.501961,0.501961}%
\pgfsetstrokecolor{currentstroke}%
\pgfsetstrokeopacity{0.500000}%
\pgfsetdash{}{0pt}%
\pgfpathmoveto{\pgfqpoint{4.140568in}{3.457345in}}%
\pgfusepath{stroke}%
\end{pgfscope}%
\begin{pgfscope}%
\pgfpathrectangle{\pgfqpoint{0.100000in}{2.493063in}}{\pgfqpoint{5.037500in}{3.427208in}}%
\pgfusepath{clip}%
\pgfsetbuttcap%
\pgfsetroundjoin%
\definecolor{currentfill}{rgb}{0.501961,0.501961,0.501961}%
\pgfsetfillcolor{currentfill}%
\pgfsetfillopacity{0.500000}%
\pgfsetlinewidth{0.250937pt}%
\definecolor{currentstroke}{rgb}{0.000000,0.000000,0.000000}%
\pgfsetstrokecolor{currentstroke}%
\pgfsetstrokeopacity{0.500000}%
\pgfsetdash{}{0pt}%
\pgfsys@defobject{currentmarker}{\pgfqpoint{-0.013889in}{-0.013889in}}{\pgfqpoint{0.013889in}{0.013889in}}{%
\pgfpathmoveto{\pgfqpoint{0.000000in}{-0.013889in}}%
\pgfpathcurveto{\pgfqpoint{0.003683in}{-0.013889in}}{\pgfqpoint{0.007216in}{-0.012425in}}{\pgfqpoint{0.009821in}{-0.009821in}}%
\pgfpathcurveto{\pgfqpoint{0.012425in}{-0.007216in}}{\pgfqpoint{0.013889in}{-0.003683in}}{\pgfqpoint{0.013889in}{0.000000in}}%
\pgfpathcurveto{\pgfqpoint{0.013889in}{0.003683in}}{\pgfqpoint{0.012425in}{0.007216in}}{\pgfqpoint{0.009821in}{0.009821in}}%
\pgfpathcurveto{\pgfqpoint{0.007216in}{0.012425in}}{\pgfqpoint{0.003683in}{0.013889in}}{\pgfqpoint{0.000000in}{0.013889in}}%
\pgfpathcurveto{\pgfqpoint{-0.003683in}{0.013889in}}{\pgfqpoint{-0.007216in}{0.012425in}}{\pgfqpoint{-0.009821in}{0.009821in}}%
\pgfpathcurveto{\pgfqpoint{-0.012425in}{0.007216in}}{\pgfqpoint{-0.013889in}{0.003683in}}{\pgfqpoint{-0.013889in}{0.000000in}}%
\pgfpathcurveto{\pgfqpoint{-0.013889in}{-0.003683in}}{\pgfqpoint{-0.012425in}{-0.007216in}}{\pgfqpoint{-0.009821in}{-0.009821in}}%
\pgfpathcurveto{\pgfqpoint{-0.007216in}{-0.012425in}}{\pgfqpoint{-0.003683in}{-0.013889in}}{\pgfqpoint{0.000000in}{-0.013889in}}%
\pgfpathlineto{\pgfqpoint{0.000000in}{-0.013889in}}%
\pgfpathclose%
\pgfusepath{stroke,fill}%
}%
\begin{pgfscope}%
\pgfsys@transformshift{4.140568in}{3.457345in}%
\pgfsys@useobject{currentmarker}{}%
\end{pgfscope}%
\end{pgfscope}%
\begin{pgfscope}%
\pgfpathrectangle{\pgfqpoint{0.100000in}{2.493063in}}{\pgfqpoint{5.037500in}{3.427208in}}%
\pgfusepath{clip}%
\pgfsetrectcap%
\pgfsetroundjoin%
\pgfsetlinewidth{1.505625pt}%
\definecolor{currentstroke}{rgb}{0.000000,0.000000,1.000000}%
\pgfsetstrokecolor{currentstroke}%
\pgfsetstrokeopacity{0.500000}%
\pgfsetdash{}{0pt}%
\pgfpathmoveto{\pgfqpoint{4.149644in}{3.187988in}}%
\pgfusepath{stroke}%
\end{pgfscope}%
\begin{pgfscope}%
\pgfpathrectangle{\pgfqpoint{0.100000in}{2.493063in}}{\pgfqpoint{5.037500in}{3.427208in}}%
\pgfusepath{clip}%
\pgfsetbuttcap%
\pgfsetroundjoin%
\definecolor{currentfill}{rgb}{0.000000,0.000000,1.000000}%
\pgfsetfillcolor{currentfill}%
\pgfsetfillopacity{0.500000}%
\pgfsetlinewidth{0.250937pt}%
\definecolor{currentstroke}{rgb}{0.000000,0.000000,0.000000}%
\pgfsetstrokecolor{currentstroke}%
\pgfsetstrokeopacity{0.500000}%
\pgfsetdash{}{0pt}%
\pgfsys@defobject{currentmarker}{\pgfqpoint{-0.005556in}{-0.005556in}}{\pgfqpoint{0.005556in}{0.005556in}}{%
\pgfpathmoveto{\pgfqpoint{0.000000in}{-0.005556in}}%
\pgfpathcurveto{\pgfqpoint{0.001473in}{-0.005556in}}{\pgfqpoint{0.002887in}{-0.004970in}}{\pgfqpoint{0.003928in}{-0.003928in}}%
\pgfpathcurveto{\pgfqpoint{0.004970in}{-0.002887in}}{\pgfqpoint{0.005556in}{-0.001473in}}{\pgfqpoint{0.005556in}{0.000000in}}%
\pgfpathcurveto{\pgfqpoint{0.005556in}{0.001473in}}{\pgfqpoint{0.004970in}{0.002887in}}{\pgfqpoint{0.003928in}{0.003928in}}%
\pgfpathcurveto{\pgfqpoint{0.002887in}{0.004970in}}{\pgfqpoint{0.001473in}{0.005556in}}{\pgfqpoint{0.000000in}{0.005556in}}%
\pgfpathcurveto{\pgfqpoint{-0.001473in}{0.005556in}}{\pgfqpoint{-0.002887in}{0.004970in}}{\pgfqpoint{-0.003928in}{0.003928in}}%
\pgfpathcurveto{\pgfqpoint{-0.004970in}{0.002887in}}{\pgfqpoint{-0.005556in}{0.001473in}}{\pgfqpoint{-0.005556in}{0.000000in}}%
\pgfpathcurveto{\pgfqpoint{-0.005556in}{-0.001473in}}{\pgfqpoint{-0.004970in}{-0.002887in}}{\pgfqpoint{-0.003928in}{-0.003928in}}%
\pgfpathcurveto{\pgfqpoint{-0.002887in}{-0.004970in}}{\pgfqpoint{-0.001473in}{-0.005556in}}{\pgfqpoint{0.000000in}{-0.005556in}}%
\pgfpathlineto{\pgfqpoint{0.000000in}{-0.005556in}}%
\pgfpathclose%
\pgfusepath{stroke,fill}%
}%
\begin{pgfscope}%
\pgfsys@transformshift{4.149644in}{3.187988in}%
\pgfsys@useobject{currentmarker}{}%
\end{pgfscope}%
\end{pgfscope}%
\begin{pgfscope}%
\pgfpathrectangle{\pgfqpoint{0.100000in}{2.493063in}}{\pgfqpoint{5.037500in}{3.427208in}}%
\pgfusepath{clip}%
\pgfsetrectcap%
\pgfsetroundjoin%
\pgfsetlinewidth{1.505625pt}%
\definecolor{currentstroke}{rgb}{0.678431,1.000000,0.184314}%
\pgfsetstrokecolor{currentstroke}%
\pgfsetstrokeopacity{0.500000}%
\pgfsetdash{}{0pt}%
\pgfpathmoveto{\pgfqpoint{4.372433in}{2.952717in}}%
\pgfusepath{stroke}%
\end{pgfscope}%
\begin{pgfscope}%
\pgfpathrectangle{\pgfqpoint{0.100000in}{2.493063in}}{\pgfqpoint{5.037500in}{3.427208in}}%
\pgfusepath{clip}%
\pgfsetbuttcap%
\pgfsetroundjoin%
\definecolor{currentfill}{rgb}{0.678431,1.000000,0.184314}%
\pgfsetfillcolor{currentfill}%
\pgfsetfillopacity{0.500000}%
\pgfsetlinewidth{0.250937pt}%
\definecolor{currentstroke}{rgb}{0.000000,0.000000,0.000000}%
\pgfsetstrokecolor{currentstroke}%
\pgfsetstrokeopacity{0.500000}%
\pgfsetdash{}{0pt}%
\pgfsys@defobject{currentmarker}{\pgfqpoint{-0.011111in}{-0.011111in}}{\pgfqpoint{0.011111in}{0.011111in}}{%
\pgfpathmoveto{\pgfqpoint{0.000000in}{-0.011111in}}%
\pgfpathcurveto{\pgfqpoint{0.002947in}{-0.011111in}}{\pgfqpoint{0.005773in}{-0.009940in}}{\pgfqpoint{0.007857in}{-0.007857in}}%
\pgfpathcurveto{\pgfqpoint{0.009940in}{-0.005773in}}{\pgfqpoint{0.011111in}{-0.002947in}}{\pgfqpoint{0.011111in}{0.000000in}}%
\pgfpathcurveto{\pgfqpoint{0.011111in}{0.002947in}}{\pgfqpoint{0.009940in}{0.005773in}}{\pgfqpoint{0.007857in}{0.007857in}}%
\pgfpathcurveto{\pgfqpoint{0.005773in}{0.009940in}}{\pgfqpoint{0.002947in}{0.011111in}}{\pgfqpoint{0.000000in}{0.011111in}}%
\pgfpathcurveto{\pgfqpoint{-0.002947in}{0.011111in}}{\pgfqpoint{-0.005773in}{0.009940in}}{\pgfqpoint{-0.007857in}{0.007857in}}%
\pgfpathcurveto{\pgfqpoint{-0.009940in}{0.005773in}}{\pgfqpoint{-0.011111in}{0.002947in}}{\pgfqpoint{-0.011111in}{0.000000in}}%
\pgfpathcurveto{\pgfqpoint{-0.011111in}{-0.002947in}}{\pgfqpoint{-0.009940in}{-0.005773in}}{\pgfqpoint{-0.007857in}{-0.007857in}}%
\pgfpathcurveto{\pgfqpoint{-0.005773in}{-0.009940in}}{\pgfqpoint{-0.002947in}{-0.011111in}}{\pgfqpoint{0.000000in}{-0.011111in}}%
\pgfpathlineto{\pgfqpoint{0.000000in}{-0.011111in}}%
\pgfpathclose%
\pgfusepath{stroke,fill}%
}%
\begin{pgfscope}%
\pgfsys@transformshift{4.372433in}{2.952717in}%
\pgfsys@useobject{currentmarker}{}%
\end{pgfscope}%
\end{pgfscope}%
\begin{pgfscope}%
\pgfpathrectangle{\pgfqpoint{0.100000in}{2.493063in}}{\pgfqpoint{5.037500in}{3.427208in}}%
\pgfusepath{clip}%
\pgfsetrectcap%
\pgfsetroundjoin%
\pgfsetlinewidth{1.505625pt}%
\definecolor{currentstroke}{rgb}{0.501961,0.501961,0.501961}%
\pgfsetstrokecolor{currentstroke}%
\pgfsetstrokeopacity{0.500000}%
\pgfsetdash{}{0pt}%
\pgfpathmoveto{\pgfqpoint{4.197819in}{2.969530in}}%
\pgfusepath{stroke}%
\end{pgfscope}%
\begin{pgfscope}%
\pgfpathrectangle{\pgfqpoint{0.100000in}{2.493063in}}{\pgfqpoint{5.037500in}{3.427208in}}%
\pgfusepath{clip}%
\pgfsetbuttcap%
\pgfsetroundjoin%
\definecolor{currentfill}{rgb}{0.501961,0.501961,0.501961}%
\pgfsetfillcolor{currentfill}%
\pgfsetfillopacity{0.500000}%
\pgfsetlinewidth{0.250937pt}%
\definecolor{currentstroke}{rgb}{0.000000,0.000000,0.000000}%
\pgfsetstrokecolor{currentstroke}%
\pgfsetstrokeopacity{0.500000}%
\pgfsetdash{}{0pt}%
\pgfsys@defobject{currentmarker}{\pgfqpoint{-0.013889in}{-0.013889in}}{\pgfqpoint{0.013889in}{0.013889in}}{%
\pgfpathmoveto{\pgfqpoint{0.000000in}{-0.013889in}}%
\pgfpathcurveto{\pgfqpoint{0.003683in}{-0.013889in}}{\pgfqpoint{0.007216in}{-0.012425in}}{\pgfqpoint{0.009821in}{-0.009821in}}%
\pgfpathcurveto{\pgfqpoint{0.012425in}{-0.007216in}}{\pgfqpoint{0.013889in}{-0.003683in}}{\pgfqpoint{0.013889in}{0.000000in}}%
\pgfpathcurveto{\pgfqpoint{0.013889in}{0.003683in}}{\pgfqpoint{0.012425in}{0.007216in}}{\pgfqpoint{0.009821in}{0.009821in}}%
\pgfpathcurveto{\pgfqpoint{0.007216in}{0.012425in}}{\pgfqpoint{0.003683in}{0.013889in}}{\pgfqpoint{0.000000in}{0.013889in}}%
\pgfpathcurveto{\pgfqpoint{-0.003683in}{0.013889in}}{\pgfqpoint{-0.007216in}{0.012425in}}{\pgfqpoint{-0.009821in}{0.009821in}}%
\pgfpathcurveto{\pgfqpoint{-0.012425in}{0.007216in}}{\pgfqpoint{-0.013889in}{0.003683in}}{\pgfqpoint{-0.013889in}{0.000000in}}%
\pgfpathcurveto{\pgfqpoint{-0.013889in}{-0.003683in}}{\pgfqpoint{-0.012425in}{-0.007216in}}{\pgfqpoint{-0.009821in}{-0.009821in}}%
\pgfpathcurveto{\pgfqpoint{-0.007216in}{-0.012425in}}{\pgfqpoint{-0.003683in}{-0.013889in}}{\pgfqpoint{0.000000in}{-0.013889in}}%
\pgfpathlineto{\pgfqpoint{0.000000in}{-0.013889in}}%
\pgfpathclose%
\pgfusepath{stroke,fill}%
}%
\begin{pgfscope}%
\pgfsys@transformshift{4.197819in}{2.969530in}%
\pgfsys@useobject{currentmarker}{}%
\end{pgfscope}%
\end{pgfscope}%
\begin{pgfscope}%
\pgfpathrectangle{\pgfqpoint{0.100000in}{2.493063in}}{\pgfqpoint{5.037500in}{3.427208in}}%
\pgfusepath{clip}%
\pgfsetrectcap%
\pgfsetroundjoin%
\pgfsetlinewidth{1.505625pt}%
\definecolor{currentstroke}{rgb}{0.000000,0.000000,1.000000}%
\pgfsetstrokecolor{currentstroke}%
\pgfsetstrokeopacity{0.500000}%
\pgfsetdash{}{0pt}%
\pgfpathmoveto{\pgfqpoint{4.136698in}{3.068100in}}%
\pgfusepath{stroke}%
\end{pgfscope}%
\begin{pgfscope}%
\pgfpathrectangle{\pgfqpoint{0.100000in}{2.493063in}}{\pgfqpoint{5.037500in}{3.427208in}}%
\pgfusepath{clip}%
\pgfsetbuttcap%
\pgfsetroundjoin%
\definecolor{currentfill}{rgb}{0.000000,0.000000,1.000000}%
\pgfsetfillcolor{currentfill}%
\pgfsetfillopacity{0.500000}%
\pgfsetlinewidth{0.250937pt}%
\definecolor{currentstroke}{rgb}{0.000000,0.000000,0.000000}%
\pgfsetstrokecolor{currentstroke}%
\pgfsetstrokeopacity{0.500000}%
\pgfsetdash{}{0pt}%
\pgfsys@defobject{currentmarker}{\pgfqpoint{-0.005556in}{-0.005556in}}{\pgfqpoint{0.005556in}{0.005556in}}{%
\pgfpathmoveto{\pgfqpoint{0.000000in}{-0.005556in}}%
\pgfpathcurveto{\pgfqpoint{0.001473in}{-0.005556in}}{\pgfqpoint{0.002887in}{-0.004970in}}{\pgfqpoint{0.003928in}{-0.003928in}}%
\pgfpathcurveto{\pgfqpoint{0.004970in}{-0.002887in}}{\pgfqpoint{0.005556in}{-0.001473in}}{\pgfqpoint{0.005556in}{0.000000in}}%
\pgfpathcurveto{\pgfqpoint{0.005556in}{0.001473in}}{\pgfqpoint{0.004970in}{0.002887in}}{\pgfqpoint{0.003928in}{0.003928in}}%
\pgfpathcurveto{\pgfqpoint{0.002887in}{0.004970in}}{\pgfqpoint{0.001473in}{0.005556in}}{\pgfqpoint{0.000000in}{0.005556in}}%
\pgfpathcurveto{\pgfqpoint{-0.001473in}{0.005556in}}{\pgfqpoint{-0.002887in}{0.004970in}}{\pgfqpoint{-0.003928in}{0.003928in}}%
\pgfpathcurveto{\pgfqpoint{-0.004970in}{0.002887in}}{\pgfqpoint{-0.005556in}{0.001473in}}{\pgfqpoint{-0.005556in}{0.000000in}}%
\pgfpathcurveto{\pgfqpoint{-0.005556in}{-0.001473in}}{\pgfqpoint{-0.004970in}{-0.002887in}}{\pgfqpoint{-0.003928in}{-0.003928in}}%
\pgfpathcurveto{\pgfqpoint{-0.002887in}{-0.004970in}}{\pgfqpoint{-0.001473in}{-0.005556in}}{\pgfqpoint{0.000000in}{-0.005556in}}%
\pgfpathlineto{\pgfqpoint{0.000000in}{-0.005556in}}%
\pgfpathclose%
\pgfusepath{stroke,fill}%
}%
\begin{pgfscope}%
\pgfsys@transformshift{4.136698in}{3.068100in}%
\pgfsys@useobject{currentmarker}{}%
\end{pgfscope}%
\end{pgfscope}%
\begin{pgfscope}%
\pgfpathrectangle{\pgfqpoint{0.100000in}{2.493063in}}{\pgfqpoint{5.037500in}{3.427208in}}%
\pgfusepath{clip}%
\pgfsetrectcap%
\pgfsetroundjoin%
\pgfsetlinewidth{1.505625pt}%
\definecolor{currentstroke}{rgb}{0.501961,0.501961,0.501961}%
\pgfsetstrokecolor{currentstroke}%
\pgfsetstrokeopacity{0.500000}%
\pgfsetdash{}{0pt}%
\pgfpathmoveto{\pgfqpoint{4.111216in}{3.318146in}}%
\pgfusepath{stroke}%
\end{pgfscope}%
\begin{pgfscope}%
\pgfpathrectangle{\pgfqpoint{0.100000in}{2.493063in}}{\pgfqpoint{5.037500in}{3.427208in}}%
\pgfusepath{clip}%
\pgfsetbuttcap%
\pgfsetroundjoin%
\definecolor{currentfill}{rgb}{0.501961,0.501961,0.501961}%
\pgfsetfillcolor{currentfill}%
\pgfsetfillopacity{0.500000}%
\pgfsetlinewidth{0.250937pt}%
\definecolor{currentstroke}{rgb}{0.000000,0.000000,0.000000}%
\pgfsetstrokecolor{currentstroke}%
\pgfsetstrokeopacity{0.500000}%
\pgfsetdash{}{0pt}%
\pgfsys@defobject{currentmarker}{\pgfqpoint{-0.013889in}{-0.013889in}}{\pgfqpoint{0.013889in}{0.013889in}}{%
\pgfpathmoveto{\pgfqpoint{0.000000in}{-0.013889in}}%
\pgfpathcurveto{\pgfqpoint{0.003683in}{-0.013889in}}{\pgfqpoint{0.007216in}{-0.012425in}}{\pgfqpoint{0.009821in}{-0.009821in}}%
\pgfpathcurveto{\pgfqpoint{0.012425in}{-0.007216in}}{\pgfqpoint{0.013889in}{-0.003683in}}{\pgfqpoint{0.013889in}{0.000000in}}%
\pgfpathcurveto{\pgfqpoint{0.013889in}{0.003683in}}{\pgfqpoint{0.012425in}{0.007216in}}{\pgfqpoint{0.009821in}{0.009821in}}%
\pgfpathcurveto{\pgfqpoint{0.007216in}{0.012425in}}{\pgfqpoint{0.003683in}{0.013889in}}{\pgfqpoint{0.000000in}{0.013889in}}%
\pgfpathcurveto{\pgfqpoint{-0.003683in}{0.013889in}}{\pgfqpoint{-0.007216in}{0.012425in}}{\pgfqpoint{-0.009821in}{0.009821in}}%
\pgfpathcurveto{\pgfqpoint{-0.012425in}{0.007216in}}{\pgfqpoint{-0.013889in}{0.003683in}}{\pgfqpoint{-0.013889in}{0.000000in}}%
\pgfpathcurveto{\pgfqpoint{-0.013889in}{-0.003683in}}{\pgfqpoint{-0.012425in}{-0.007216in}}{\pgfqpoint{-0.009821in}{-0.009821in}}%
\pgfpathcurveto{\pgfqpoint{-0.007216in}{-0.012425in}}{\pgfqpoint{-0.003683in}{-0.013889in}}{\pgfqpoint{0.000000in}{-0.013889in}}%
\pgfpathlineto{\pgfqpoint{0.000000in}{-0.013889in}}%
\pgfpathclose%
\pgfusepath{stroke,fill}%
}%
\begin{pgfscope}%
\pgfsys@transformshift{4.111216in}{3.318146in}%
\pgfsys@useobject{currentmarker}{}%
\end{pgfscope}%
\end{pgfscope}%
\begin{pgfscope}%
\pgfpathrectangle{\pgfqpoint{0.100000in}{2.493063in}}{\pgfqpoint{5.037500in}{3.427208in}}%
\pgfusepath{clip}%
\pgfsetrectcap%
\pgfsetroundjoin%
\pgfsetlinewidth{1.505625pt}%
\definecolor{currentstroke}{rgb}{0.000000,0.000000,1.000000}%
\pgfsetstrokecolor{currentstroke}%
\pgfsetstrokeopacity{0.500000}%
\pgfsetdash{}{0pt}%
\pgfpathmoveto{\pgfqpoint{4.199267in}{3.254835in}}%
\pgfusepath{stroke}%
\end{pgfscope}%
\begin{pgfscope}%
\pgfpathrectangle{\pgfqpoint{0.100000in}{2.493063in}}{\pgfqpoint{5.037500in}{3.427208in}}%
\pgfusepath{clip}%
\pgfsetbuttcap%
\pgfsetroundjoin%
\definecolor{currentfill}{rgb}{0.000000,0.000000,1.000000}%
\pgfsetfillcolor{currentfill}%
\pgfsetfillopacity{0.500000}%
\pgfsetlinewidth{0.250937pt}%
\definecolor{currentstroke}{rgb}{0.000000,0.000000,0.000000}%
\pgfsetstrokecolor{currentstroke}%
\pgfsetstrokeopacity{0.500000}%
\pgfsetdash{}{0pt}%
\pgfsys@defobject{currentmarker}{\pgfqpoint{-0.005556in}{-0.005556in}}{\pgfqpoint{0.005556in}{0.005556in}}{%
\pgfpathmoveto{\pgfqpoint{0.000000in}{-0.005556in}}%
\pgfpathcurveto{\pgfqpoint{0.001473in}{-0.005556in}}{\pgfqpoint{0.002887in}{-0.004970in}}{\pgfqpoint{0.003928in}{-0.003928in}}%
\pgfpathcurveto{\pgfqpoint{0.004970in}{-0.002887in}}{\pgfqpoint{0.005556in}{-0.001473in}}{\pgfqpoint{0.005556in}{0.000000in}}%
\pgfpathcurveto{\pgfqpoint{0.005556in}{0.001473in}}{\pgfqpoint{0.004970in}{0.002887in}}{\pgfqpoint{0.003928in}{0.003928in}}%
\pgfpathcurveto{\pgfqpoint{0.002887in}{0.004970in}}{\pgfqpoint{0.001473in}{0.005556in}}{\pgfqpoint{0.000000in}{0.005556in}}%
\pgfpathcurveto{\pgfqpoint{-0.001473in}{0.005556in}}{\pgfqpoint{-0.002887in}{0.004970in}}{\pgfqpoint{-0.003928in}{0.003928in}}%
\pgfpathcurveto{\pgfqpoint{-0.004970in}{0.002887in}}{\pgfqpoint{-0.005556in}{0.001473in}}{\pgfqpoint{-0.005556in}{0.000000in}}%
\pgfpathcurveto{\pgfqpoint{-0.005556in}{-0.001473in}}{\pgfqpoint{-0.004970in}{-0.002887in}}{\pgfqpoint{-0.003928in}{-0.003928in}}%
\pgfpathcurveto{\pgfqpoint{-0.002887in}{-0.004970in}}{\pgfqpoint{-0.001473in}{-0.005556in}}{\pgfqpoint{0.000000in}{-0.005556in}}%
\pgfpathlineto{\pgfqpoint{0.000000in}{-0.005556in}}%
\pgfpathclose%
\pgfusepath{stroke,fill}%
}%
\begin{pgfscope}%
\pgfsys@transformshift{4.199267in}{3.254835in}%
\pgfsys@useobject{currentmarker}{}%
\end{pgfscope}%
\end{pgfscope}%
\begin{pgfscope}%
\pgfpathrectangle{\pgfqpoint{0.100000in}{2.493063in}}{\pgfqpoint{5.037500in}{3.427208in}}%
\pgfusepath{clip}%
\pgfsetrectcap%
\pgfsetroundjoin%
\pgfsetlinewidth{1.505625pt}%
\definecolor{currentstroke}{rgb}{0.000000,0.000000,1.000000}%
\pgfsetstrokecolor{currentstroke}%
\pgfsetstrokeopacity{0.500000}%
\pgfsetdash{}{0pt}%
\pgfpathmoveto{\pgfqpoint{4.283427in}{3.207572in}}%
\pgfusepath{stroke}%
\end{pgfscope}%
\begin{pgfscope}%
\pgfpathrectangle{\pgfqpoint{0.100000in}{2.493063in}}{\pgfqpoint{5.037500in}{3.427208in}}%
\pgfusepath{clip}%
\pgfsetbuttcap%
\pgfsetroundjoin%
\definecolor{currentfill}{rgb}{0.000000,0.000000,1.000000}%
\pgfsetfillcolor{currentfill}%
\pgfsetfillopacity{0.500000}%
\pgfsetlinewidth{0.250937pt}%
\definecolor{currentstroke}{rgb}{0.000000,0.000000,0.000000}%
\pgfsetstrokecolor{currentstroke}%
\pgfsetstrokeopacity{0.500000}%
\pgfsetdash{}{0pt}%
\pgfsys@defobject{currentmarker}{\pgfqpoint{-0.005556in}{-0.005556in}}{\pgfqpoint{0.005556in}{0.005556in}}{%
\pgfpathmoveto{\pgfqpoint{0.000000in}{-0.005556in}}%
\pgfpathcurveto{\pgfqpoint{0.001473in}{-0.005556in}}{\pgfqpoint{0.002887in}{-0.004970in}}{\pgfqpoint{0.003928in}{-0.003928in}}%
\pgfpathcurveto{\pgfqpoint{0.004970in}{-0.002887in}}{\pgfqpoint{0.005556in}{-0.001473in}}{\pgfqpoint{0.005556in}{0.000000in}}%
\pgfpathcurveto{\pgfqpoint{0.005556in}{0.001473in}}{\pgfqpoint{0.004970in}{0.002887in}}{\pgfqpoint{0.003928in}{0.003928in}}%
\pgfpathcurveto{\pgfqpoint{0.002887in}{0.004970in}}{\pgfqpoint{0.001473in}{0.005556in}}{\pgfqpoint{0.000000in}{0.005556in}}%
\pgfpathcurveto{\pgfqpoint{-0.001473in}{0.005556in}}{\pgfqpoint{-0.002887in}{0.004970in}}{\pgfqpoint{-0.003928in}{0.003928in}}%
\pgfpathcurveto{\pgfqpoint{-0.004970in}{0.002887in}}{\pgfqpoint{-0.005556in}{0.001473in}}{\pgfqpoint{-0.005556in}{0.000000in}}%
\pgfpathcurveto{\pgfqpoint{-0.005556in}{-0.001473in}}{\pgfqpoint{-0.004970in}{-0.002887in}}{\pgfqpoint{-0.003928in}{-0.003928in}}%
\pgfpathcurveto{\pgfqpoint{-0.002887in}{-0.004970in}}{\pgfqpoint{-0.001473in}{-0.005556in}}{\pgfqpoint{0.000000in}{-0.005556in}}%
\pgfpathlineto{\pgfqpoint{0.000000in}{-0.005556in}}%
\pgfpathclose%
\pgfusepath{stroke,fill}%
}%
\begin{pgfscope}%
\pgfsys@transformshift{4.283427in}{3.207572in}%
\pgfsys@useobject{currentmarker}{}%
\end{pgfscope}%
\end{pgfscope}%
\begin{pgfscope}%
\pgfpathrectangle{\pgfqpoint{0.100000in}{2.493063in}}{\pgfqpoint{5.037500in}{3.427208in}}%
\pgfusepath{clip}%
\pgfsetrectcap%
\pgfsetroundjoin%
\pgfsetlinewidth{1.505625pt}%
\definecolor{currentstroke}{rgb}{0.501961,0.501961,0.501961}%
\pgfsetstrokecolor{currentstroke}%
\pgfsetstrokeopacity{0.500000}%
\pgfsetdash{}{0pt}%
\pgfpathmoveto{\pgfqpoint{3.743850in}{3.387362in}}%
\pgfusepath{stroke}%
\end{pgfscope}%
\begin{pgfscope}%
\pgfpathrectangle{\pgfqpoint{0.100000in}{2.493063in}}{\pgfqpoint{5.037500in}{3.427208in}}%
\pgfusepath{clip}%
\pgfsetbuttcap%
\pgfsetroundjoin%
\definecolor{currentfill}{rgb}{0.501961,0.501961,0.501961}%
\pgfsetfillcolor{currentfill}%
\pgfsetfillopacity{0.500000}%
\pgfsetlinewidth{0.250937pt}%
\definecolor{currentstroke}{rgb}{0.000000,0.000000,0.000000}%
\pgfsetstrokecolor{currentstroke}%
\pgfsetstrokeopacity{0.500000}%
\pgfsetdash{}{0pt}%
\pgfsys@defobject{currentmarker}{\pgfqpoint{-0.013889in}{-0.013889in}}{\pgfqpoint{0.013889in}{0.013889in}}{%
\pgfpathmoveto{\pgfqpoint{0.000000in}{-0.013889in}}%
\pgfpathcurveto{\pgfqpoint{0.003683in}{-0.013889in}}{\pgfqpoint{0.007216in}{-0.012425in}}{\pgfqpoint{0.009821in}{-0.009821in}}%
\pgfpathcurveto{\pgfqpoint{0.012425in}{-0.007216in}}{\pgfqpoint{0.013889in}{-0.003683in}}{\pgfqpoint{0.013889in}{0.000000in}}%
\pgfpathcurveto{\pgfqpoint{0.013889in}{0.003683in}}{\pgfqpoint{0.012425in}{0.007216in}}{\pgfqpoint{0.009821in}{0.009821in}}%
\pgfpathcurveto{\pgfqpoint{0.007216in}{0.012425in}}{\pgfqpoint{0.003683in}{0.013889in}}{\pgfqpoint{0.000000in}{0.013889in}}%
\pgfpathcurveto{\pgfqpoint{-0.003683in}{0.013889in}}{\pgfqpoint{-0.007216in}{0.012425in}}{\pgfqpoint{-0.009821in}{0.009821in}}%
\pgfpathcurveto{\pgfqpoint{-0.012425in}{0.007216in}}{\pgfqpoint{-0.013889in}{0.003683in}}{\pgfqpoint{-0.013889in}{0.000000in}}%
\pgfpathcurveto{\pgfqpoint{-0.013889in}{-0.003683in}}{\pgfqpoint{-0.012425in}{-0.007216in}}{\pgfqpoint{-0.009821in}{-0.009821in}}%
\pgfpathcurveto{\pgfqpoint{-0.007216in}{-0.012425in}}{\pgfqpoint{-0.003683in}{-0.013889in}}{\pgfqpoint{0.000000in}{-0.013889in}}%
\pgfpathlineto{\pgfqpoint{0.000000in}{-0.013889in}}%
\pgfpathclose%
\pgfusepath{stroke,fill}%
}%
\begin{pgfscope}%
\pgfsys@transformshift{3.743850in}{3.387362in}%
\pgfsys@useobject{currentmarker}{}%
\end{pgfscope}%
\end{pgfscope}%
\begin{pgfscope}%
\pgfpathrectangle{\pgfqpoint{0.100000in}{2.493063in}}{\pgfqpoint{5.037500in}{3.427208in}}%
\pgfusepath{clip}%
\pgfsetrectcap%
\pgfsetroundjoin%
\pgfsetlinewidth{1.505625pt}%
\definecolor{currentstroke}{rgb}{0.000000,0.000000,1.000000}%
\pgfsetstrokecolor{currentstroke}%
\pgfsetstrokeopacity{0.500000}%
\pgfsetdash{}{0pt}%
\pgfpathmoveto{\pgfqpoint{3.584571in}{3.402865in}}%
\pgfusepath{stroke}%
\end{pgfscope}%
\begin{pgfscope}%
\pgfpathrectangle{\pgfqpoint{0.100000in}{2.493063in}}{\pgfqpoint{5.037500in}{3.427208in}}%
\pgfusepath{clip}%
\pgfsetbuttcap%
\pgfsetroundjoin%
\definecolor{currentfill}{rgb}{0.000000,0.000000,1.000000}%
\pgfsetfillcolor{currentfill}%
\pgfsetfillopacity{0.500000}%
\pgfsetlinewidth{0.250937pt}%
\definecolor{currentstroke}{rgb}{0.000000,0.000000,0.000000}%
\pgfsetstrokecolor{currentstroke}%
\pgfsetstrokeopacity{0.500000}%
\pgfsetdash{}{0pt}%
\pgfsys@defobject{currentmarker}{\pgfqpoint{-0.005556in}{-0.005556in}}{\pgfqpoint{0.005556in}{0.005556in}}{%
\pgfpathmoveto{\pgfqpoint{0.000000in}{-0.005556in}}%
\pgfpathcurveto{\pgfqpoint{0.001473in}{-0.005556in}}{\pgfqpoint{0.002887in}{-0.004970in}}{\pgfqpoint{0.003928in}{-0.003928in}}%
\pgfpathcurveto{\pgfqpoint{0.004970in}{-0.002887in}}{\pgfqpoint{0.005556in}{-0.001473in}}{\pgfqpoint{0.005556in}{0.000000in}}%
\pgfpathcurveto{\pgfqpoint{0.005556in}{0.001473in}}{\pgfqpoint{0.004970in}{0.002887in}}{\pgfqpoint{0.003928in}{0.003928in}}%
\pgfpathcurveto{\pgfqpoint{0.002887in}{0.004970in}}{\pgfqpoint{0.001473in}{0.005556in}}{\pgfqpoint{0.000000in}{0.005556in}}%
\pgfpathcurveto{\pgfqpoint{-0.001473in}{0.005556in}}{\pgfqpoint{-0.002887in}{0.004970in}}{\pgfqpoint{-0.003928in}{0.003928in}}%
\pgfpathcurveto{\pgfqpoint{-0.004970in}{0.002887in}}{\pgfqpoint{-0.005556in}{0.001473in}}{\pgfqpoint{-0.005556in}{0.000000in}}%
\pgfpathcurveto{\pgfqpoint{-0.005556in}{-0.001473in}}{\pgfqpoint{-0.004970in}{-0.002887in}}{\pgfqpoint{-0.003928in}{-0.003928in}}%
\pgfpathcurveto{\pgfqpoint{-0.002887in}{-0.004970in}}{\pgfqpoint{-0.001473in}{-0.005556in}}{\pgfqpoint{0.000000in}{-0.005556in}}%
\pgfpathlineto{\pgfqpoint{0.000000in}{-0.005556in}}%
\pgfpathclose%
\pgfusepath{stroke,fill}%
}%
\begin{pgfscope}%
\pgfsys@transformshift{3.584571in}{3.402865in}%
\pgfsys@useobject{currentmarker}{}%
\end{pgfscope}%
\end{pgfscope}%
\begin{pgfscope}%
\pgfpathrectangle{\pgfqpoint{0.100000in}{2.493063in}}{\pgfqpoint{5.037500in}{3.427208in}}%
\pgfusepath{clip}%
\pgfsetrectcap%
\pgfsetroundjoin%
\pgfsetlinewidth{1.505625pt}%
\definecolor{currentstroke}{rgb}{0.501961,0.501961,0.501961}%
\pgfsetstrokecolor{currentstroke}%
\pgfsetstrokeopacity{0.500000}%
\pgfsetdash{}{0pt}%
\pgfpathmoveto{\pgfqpoint{4.327246in}{3.126750in}}%
\pgfusepath{stroke}%
\end{pgfscope}%
\begin{pgfscope}%
\pgfpathrectangle{\pgfqpoint{0.100000in}{2.493063in}}{\pgfqpoint{5.037500in}{3.427208in}}%
\pgfusepath{clip}%
\pgfsetbuttcap%
\pgfsetroundjoin%
\definecolor{currentfill}{rgb}{0.501961,0.501961,0.501961}%
\pgfsetfillcolor{currentfill}%
\pgfsetfillopacity{0.500000}%
\pgfsetlinewidth{0.250937pt}%
\definecolor{currentstroke}{rgb}{0.000000,0.000000,0.000000}%
\pgfsetstrokecolor{currentstroke}%
\pgfsetstrokeopacity{0.500000}%
\pgfsetdash{}{0pt}%
\pgfsys@defobject{currentmarker}{\pgfqpoint{-0.013889in}{-0.013889in}}{\pgfqpoint{0.013889in}{0.013889in}}{%
\pgfpathmoveto{\pgfqpoint{0.000000in}{-0.013889in}}%
\pgfpathcurveto{\pgfqpoint{0.003683in}{-0.013889in}}{\pgfqpoint{0.007216in}{-0.012425in}}{\pgfqpoint{0.009821in}{-0.009821in}}%
\pgfpathcurveto{\pgfqpoint{0.012425in}{-0.007216in}}{\pgfqpoint{0.013889in}{-0.003683in}}{\pgfqpoint{0.013889in}{0.000000in}}%
\pgfpathcurveto{\pgfqpoint{0.013889in}{0.003683in}}{\pgfqpoint{0.012425in}{0.007216in}}{\pgfqpoint{0.009821in}{0.009821in}}%
\pgfpathcurveto{\pgfqpoint{0.007216in}{0.012425in}}{\pgfqpoint{0.003683in}{0.013889in}}{\pgfqpoint{0.000000in}{0.013889in}}%
\pgfpathcurveto{\pgfqpoint{-0.003683in}{0.013889in}}{\pgfqpoint{-0.007216in}{0.012425in}}{\pgfqpoint{-0.009821in}{0.009821in}}%
\pgfpathcurveto{\pgfqpoint{-0.012425in}{0.007216in}}{\pgfqpoint{-0.013889in}{0.003683in}}{\pgfqpoint{-0.013889in}{0.000000in}}%
\pgfpathcurveto{\pgfqpoint{-0.013889in}{-0.003683in}}{\pgfqpoint{-0.012425in}{-0.007216in}}{\pgfqpoint{-0.009821in}{-0.009821in}}%
\pgfpathcurveto{\pgfqpoint{-0.007216in}{-0.012425in}}{\pgfqpoint{-0.003683in}{-0.013889in}}{\pgfqpoint{0.000000in}{-0.013889in}}%
\pgfpathlineto{\pgfqpoint{0.000000in}{-0.013889in}}%
\pgfpathclose%
\pgfusepath{stroke,fill}%
}%
\begin{pgfscope}%
\pgfsys@transformshift{4.327246in}{3.126750in}%
\pgfsys@useobject{currentmarker}{}%
\end{pgfscope}%
\end{pgfscope}%
\begin{pgfscope}%
\pgfpathrectangle{\pgfqpoint{0.100000in}{2.493063in}}{\pgfqpoint{5.037500in}{3.427208in}}%
\pgfusepath{clip}%
\pgfsetrectcap%
\pgfsetroundjoin%
\pgfsetlinewidth{1.505625pt}%
\definecolor{currentstroke}{rgb}{0.501961,0.501961,0.501961}%
\pgfsetstrokecolor{currentstroke}%
\pgfsetstrokeopacity{0.500000}%
\pgfsetdash{}{0pt}%
\pgfpathmoveto{\pgfqpoint{4.158371in}{3.057590in}}%
\pgfusepath{stroke}%
\end{pgfscope}%
\begin{pgfscope}%
\pgfpathrectangle{\pgfqpoint{0.100000in}{2.493063in}}{\pgfqpoint{5.037500in}{3.427208in}}%
\pgfusepath{clip}%
\pgfsetbuttcap%
\pgfsetroundjoin%
\definecolor{currentfill}{rgb}{0.501961,0.501961,0.501961}%
\pgfsetfillcolor{currentfill}%
\pgfsetfillopacity{0.500000}%
\pgfsetlinewidth{0.250937pt}%
\definecolor{currentstroke}{rgb}{0.000000,0.000000,0.000000}%
\pgfsetstrokecolor{currentstroke}%
\pgfsetstrokeopacity{0.500000}%
\pgfsetdash{}{0pt}%
\pgfsys@defobject{currentmarker}{\pgfqpoint{-0.013889in}{-0.013889in}}{\pgfqpoint{0.013889in}{0.013889in}}{%
\pgfpathmoveto{\pgfqpoint{0.000000in}{-0.013889in}}%
\pgfpathcurveto{\pgfqpoint{0.003683in}{-0.013889in}}{\pgfqpoint{0.007216in}{-0.012425in}}{\pgfqpoint{0.009821in}{-0.009821in}}%
\pgfpathcurveto{\pgfqpoint{0.012425in}{-0.007216in}}{\pgfqpoint{0.013889in}{-0.003683in}}{\pgfqpoint{0.013889in}{0.000000in}}%
\pgfpathcurveto{\pgfqpoint{0.013889in}{0.003683in}}{\pgfqpoint{0.012425in}{0.007216in}}{\pgfqpoint{0.009821in}{0.009821in}}%
\pgfpathcurveto{\pgfqpoint{0.007216in}{0.012425in}}{\pgfqpoint{0.003683in}{0.013889in}}{\pgfqpoint{0.000000in}{0.013889in}}%
\pgfpathcurveto{\pgfqpoint{-0.003683in}{0.013889in}}{\pgfqpoint{-0.007216in}{0.012425in}}{\pgfqpoint{-0.009821in}{0.009821in}}%
\pgfpathcurveto{\pgfqpoint{-0.012425in}{0.007216in}}{\pgfqpoint{-0.013889in}{0.003683in}}{\pgfqpoint{-0.013889in}{0.000000in}}%
\pgfpathcurveto{\pgfqpoint{-0.013889in}{-0.003683in}}{\pgfqpoint{-0.012425in}{-0.007216in}}{\pgfqpoint{-0.009821in}{-0.009821in}}%
\pgfpathcurveto{\pgfqpoint{-0.007216in}{-0.012425in}}{\pgfqpoint{-0.003683in}{-0.013889in}}{\pgfqpoint{0.000000in}{-0.013889in}}%
\pgfpathlineto{\pgfqpoint{0.000000in}{-0.013889in}}%
\pgfpathclose%
\pgfusepath{stroke,fill}%
}%
\begin{pgfscope}%
\pgfsys@transformshift{4.158371in}{3.057590in}%
\pgfsys@useobject{currentmarker}{}%
\end{pgfscope}%
\end{pgfscope}%
\begin{pgfscope}%
\pgfpathrectangle{\pgfqpoint{0.100000in}{2.493063in}}{\pgfqpoint{5.037500in}{3.427208in}}%
\pgfusepath{clip}%
\pgfsetrectcap%
\pgfsetroundjoin%
\pgfsetlinewidth{1.505625pt}%
\definecolor{currentstroke}{rgb}{0.501961,0.501961,0.501961}%
\pgfsetstrokecolor{currentstroke}%
\pgfsetstrokeopacity{0.500000}%
\pgfsetdash{}{0pt}%
\pgfpathmoveto{\pgfqpoint{4.305082in}{3.185300in}}%
\pgfusepath{stroke}%
\end{pgfscope}%
\begin{pgfscope}%
\pgfpathrectangle{\pgfqpoint{0.100000in}{2.493063in}}{\pgfqpoint{5.037500in}{3.427208in}}%
\pgfusepath{clip}%
\pgfsetbuttcap%
\pgfsetroundjoin%
\definecolor{currentfill}{rgb}{0.501961,0.501961,0.501961}%
\pgfsetfillcolor{currentfill}%
\pgfsetfillopacity{0.500000}%
\pgfsetlinewidth{0.250937pt}%
\definecolor{currentstroke}{rgb}{0.000000,0.000000,0.000000}%
\pgfsetstrokecolor{currentstroke}%
\pgfsetstrokeopacity{0.500000}%
\pgfsetdash{}{0pt}%
\pgfsys@defobject{currentmarker}{\pgfqpoint{-0.013889in}{-0.013889in}}{\pgfqpoint{0.013889in}{0.013889in}}{%
\pgfpathmoveto{\pgfqpoint{0.000000in}{-0.013889in}}%
\pgfpathcurveto{\pgfqpoint{0.003683in}{-0.013889in}}{\pgfqpoint{0.007216in}{-0.012425in}}{\pgfqpoint{0.009821in}{-0.009821in}}%
\pgfpathcurveto{\pgfqpoint{0.012425in}{-0.007216in}}{\pgfqpoint{0.013889in}{-0.003683in}}{\pgfqpoint{0.013889in}{0.000000in}}%
\pgfpathcurveto{\pgfqpoint{0.013889in}{0.003683in}}{\pgfqpoint{0.012425in}{0.007216in}}{\pgfqpoint{0.009821in}{0.009821in}}%
\pgfpathcurveto{\pgfqpoint{0.007216in}{0.012425in}}{\pgfqpoint{0.003683in}{0.013889in}}{\pgfqpoint{0.000000in}{0.013889in}}%
\pgfpathcurveto{\pgfqpoint{-0.003683in}{0.013889in}}{\pgfqpoint{-0.007216in}{0.012425in}}{\pgfqpoint{-0.009821in}{0.009821in}}%
\pgfpathcurveto{\pgfqpoint{-0.012425in}{0.007216in}}{\pgfqpoint{-0.013889in}{0.003683in}}{\pgfqpoint{-0.013889in}{0.000000in}}%
\pgfpathcurveto{\pgfqpoint{-0.013889in}{-0.003683in}}{\pgfqpoint{-0.012425in}{-0.007216in}}{\pgfqpoint{-0.009821in}{-0.009821in}}%
\pgfpathcurveto{\pgfqpoint{-0.007216in}{-0.012425in}}{\pgfqpoint{-0.003683in}{-0.013889in}}{\pgfqpoint{0.000000in}{-0.013889in}}%
\pgfpathlineto{\pgfqpoint{0.000000in}{-0.013889in}}%
\pgfpathclose%
\pgfusepath{stroke,fill}%
}%
\begin{pgfscope}%
\pgfsys@transformshift{4.305082in}{3.185300in}%
\pgfsys@useobject{currentmarker}{}%
\end{pgfscope}%
\end{pgfscope}%
\begin{pgfscope}%
\pgfpathrectangle{\pgfqpoint{0.100000in}{2.493063in}}{\pgfqpoint{5.037500in}{3.427208in}}%
\pgfusepath{clip}%
\pgfsetrectcap%
\pgfsetroundjoin%
\pgfsetlinewidth{1.505625pt}%
\definecolor{currentstroke}{rgb}{0.501961,0.501961,0.501961}%
\pgfsetstrokecolor{currentstroke}%
\pgfsetstrokeopacity{0.500000}%
\pgfsetdash{}{0pt}%
\pgfpathmoveto{\pgfqpoint{4.211198in}{3.132586in}}%
\pgfusepath{stroke}%
\end{pgfscope}%
\begin{pgfscope}%
\pgfpathrectangle{\pgfqpoint{0.100000in}{2.493063in}}{\pgfqpoint{5.037500in}{3.427208in}}%
\pgfusepath{clip}%
\pgfsetbuttcap%
\pgfsetroundjoin%
\definecolor{currentfill}{rgb}{0.501961,0.501961,0.501961}%
\pgfsetfillcolor{currentfill}%
\pgfsetfillopacity{0.500000}%
\pgfsetlinewidth{0.250937pt}%
\definecolor{currentstroke}{rgb}{0.000000,0.000000,0.000000}%
\pgfsetstrokecolor{currentstroke}%
\pgfsetstrokeopacity{0.500000}%
\pgfsetdash{}{0pt}%
\pgfsys@defobject{currentmarker}{\pgfqpoint{-0.013889in}{-0.013889in}}{\pgfqpoint{0.013889in}{0.013889in}}{%
\pgfpathmoveto{\pgfqpoint{0.000000in}{-0.013889in}}%
\pgfpathcurveto{\pgfqpoint{0.003683in}{-0.013889in}}{\pgfqpoint{0.007216in}{-0.012425in}}{\pgfqpoint{0.009821in}{-0.009821in}}%
\pgfpathcurveto{\pgfqpoint{0.012425in}{-0.007216in}}{\pgfqpoint{0.013889in}{-0.003683in}}{\pgfqpoint{0.013889in}{0.000000in}}%
\pgfpathcurveto{\pgfqpoint{0.013889in}{0.003683in}}{\pgfqpoint{0.012425in}{0.007216in}}{\pgfqpoint{0.009821in}{0.009821in}}%
\pgfpathcurveto{\pgfqpoint{0.007216in}{0.012425in}}{\pgfqpoint{0.003683in}{0.013889in}}{\pgfqpoint{0.000000in}{0.013889in}}%
\pgfpathcurveto{\pgfqpoint{-0.003683in}{0.013889in}}{\pgfqpoint{-0.007216in}{0.012425in}}{\pgfqpoint{-0.009821in}{0.009821in}}%
\pgfpathcurveto{\pgfqpoint{-0.012425in}{0.007216in}}{\pgfqpoint{-0.013889in}{0.003683in}}{\pgfqpoint{-0.013889in}{0.000000in}}%
\pgfpathcurveto{\pgfqpoint{-0.013889in}{-0.003683in}}{\pgfqpoint{-0.012425in}{-0.007216in}}{\pgfqpoint{-0.009821in}{-0.009821in}}%
\pgfpathcurveto{\pgfqpoint{-0.007216in}{-0.012425in}}{\pgfqpoint{-0.003683in}{-0.013889in}}{\pgfqpoint{0.000000in}{-0.013889in}}%
\pgfpathlineto{\pgfqpoint{0.000000in}{-0.013889in}}%
\pgfpathclose%
\pgfusepath{stroke,fill}%
}%
\begin{pgfscope}%
\pgfsys@transformshift{4.211198in}{3.132586in}%
\pgfsys@useobject{currentmarker}{}%
\end{pgfscope}%
\end{pgfscope}%
\begin{pgfscope}%
\pgfpathrectangle{\pgfqpoint{0.100000in}{2.493063in}}{\pgfqpoint{5.037500in}{3.427208in}}%
\pgfusepath{clip}%
\pgfsetrectcap%
\pgfsetroundjoin%
\pgfsetlinewidth{1.505625pt}%
\definecolor{currentstroke}{rgb}{0.000000,0.000000,1.000000}%
\pgfsetstrokecolor{currentstroke}%
\pgfsetstrokeopacity{0.500000}%
\pgfsetdash{}{0pt}%
\pgfpathmoveto{\pgfqpoint{3.877566in}{3.434757in}}%
\pgfusepath{stroke}%
\end{pgfscope}%
\begin{pgfscope}%
\pgfpathrectangle{\pgfqpoint{0.100000in}{2.493063in}}{\pgfqpoint{5.037500in}{3.427208in}}%
\pgfusepath{clip}%
\pgfsetbuttcap%
\pgfsetroundjoin%
\definecolor{currentfill}{rgb}{0.000000,0.000000,1.000000}%
\pgfsetfillcolor{currentfill}%
\pgfsetfillopacity{0.500000}%
\pgfsetlinewidth{0.250937pt}%
\definecolor{currentstroke}{rgb}{0.000000,0.000000,0.000000}%
\pgfsetstrokecolor{currentstroke}%
\pgfsetstrokeopacity{0.500000}%
\pgfsetdash{}{0pt}%
\pgfsys@defobject{currentmarker}{\pgfqpoint{-0.008333in}{-0.008333in}}{\pgfqpoint{0.008333in}{0.008333in}}{%
\pgfpathmoveto{\pgfqpoint{0.000000in}{-0.008333in}}%
\pgfpathcurveto{\pgfqpoint{0.002210in}{-0.008333in}}{\pgfqpoint{0.004330in}{-0.007455in}}{\pgfqpoint{0.005893in}{-0.005893in}}%
\pgfpathcurveto{\pgfqpoint{0.007455in}{-0.004330in}}{\pgfqpoint{0.008333in}{-0.002210in}}{\pgfqpoint{0.008333in}{0.000000in}}%
\pgfpathcurveto{\pgfqpoint{0.008333in}{0.002210in}}{\pgfqpoint{0.007455in}{0.004330in}}{\pgfqpoint{0.005893in}{0.005893in}}%
\pgfpathcurveto{\pgfqpoint{0.004330in}{0.007455in}}{\pgfqpoint{0.002210in}{0.008333in}}{\pgfqpoint{0.000000in}{0.008333in}}%
\pgfpathcurveto{\pgfqpoint{-0.002210in}{0.008333in}}{\pgfqpoint{-0.004330in}{0.007455in}}{\pgfqpoint{-0.005893in}{0.005893in}}%
\pgfpathcurveto{\pgfqpoint{-0.007455in}{0.004330in}}{\pgfqpoint{-0.008333in}{0.002210in}}{\pgfqpoint{-0.008333in}{0.000000in}}%
\pgfpathcurveto{\pgfqpoint{-0.008333in}{-0.002210in}}{\pgfqpoint{-0.007455in}{-0.004330in}}{\pgfqpoint{-0.005893in}{-0.005893in}}%
\pgfpathcurveto{\pgfqpoint{-0.004330in}{-0.007455in}}{\pgfqpoint{-0.002210in}{-0.008333in}}{\pgfqpoint{0.000000in}{-0.008333in}}%
\pgfpathlineto{\pgfqpoint{0.000000in}{-0.008333in}}%
\pgfpathclose%
\pgfusepath{stroke,fill}%
}%
\begin{pgfscope}%
\pgfsys@transformshift{3.877566in}{3.434757in}%
\pgfsys@useobject{currentmarker}{}%
\end{pgfscope}%
\end{pgfscope}%
\begin{pgfscope}%
\pgfpathrectangle{\pgfqpoint{0.100000in}{2.493063in}}{\pgfqpoint{5.037500in}{3.427208in}}%
\pgfusepath{clip}%
\pgfsetrectcap%
\pgfsetroundjoin%
\pgfsetlinewidth{1.505625pt}%
\definecolor{currentstroke}{rgb}{0.501961,0.501961,0.501961}%
\pgfsetstrokecolor{currentstroke}%
\pgfsetstrokeopacity{0.500000}%
\pgfsetdash{}{0pt}%
\pgfpathmoveto{\pgfqpoint{4.099020in}{3.169954in}}%
\pgfusepath{stroke}%
\end{pgfscope}%
\begin{pgfscope}%
\pgfpathrectangle{\pgfqpoint{0.100000in}{2.493063in}}{\pgfqpoint{5.037500in}{3.427208in}}%
\pgfusepath{clip}%
\pgfsetbuttcap%
\pgfsetroundjoin%
\definecolor{currentfill}{rgb}{0.501961,0.501961,0.501961}%
\pgfsetfillcolor{currentfill}%
\pgfsetfillopacity{0.500000}%
\pgfsetlinewidth{0.250937pt}%
\definecolor{currentstroke}{rgb}{0.000000,0.000000,0.000000}%
\pgfsetstrokecolor{currentstroke}%
\pgfsetstrokeopacity{0.500000}%
\pgfsetdash{}{0pt}%
\pgfsys@defobject{currentmarker}{\pgfqpoint{-0.013889in}{-0.013889in}}{\pgfqpoint{0.013889in}{0.013889in}}{%
\pgfpathmoveto{\pgfqpoint{0.000000in}{-0.013889in}}%
\pgfpathcurveto{\pgfqpoint{0.003683in}{-0.013889in}}{\pgfqpoint{0.007216in}{-0.012425in}}{\pgfqpoint{0.009821in}{-0.009821in}}%
\pgfpathcurveto{\pgfqpoint{0.012425in}{-0.007216in}}{\pgfqpoint{0.013889in}{-0.003683in}}{\pgfqpoint{0.013889in}{0.000000in}}%
\pgfpathcurveto{\pgfqpoint{0.013889in}{0.003683in}}{\pgfqpoint{0.012425in}{0.007216in}}{\pgfqpoint{0.009821in}{0.009821in}}%
\pgfpathcurveto{\pgfqpoint{0.007216in}{0.012425in}}{\pgfqpoint{0.003683in}{0.013889in}}{\pgfqpoint{0.000000in}{0.013889in}}%
\pgfpathcurveto{\pgfqpoint{-0.003683in}{0.013889in}}{\pgfqpoint{-0.007216in}{0.012425in}}{\pgfqpoint{-0.009821in}{0.009821in}}%
\pgfpathcurveto{\pgfqpoint{-0.012425in}{0.007216in}}{\pgfqpoint{-0.013889in}{0.003683in}}{\pgfqpoint{-0.013889in}{0.000000in}}%
\pgfpathcurveto{\pgfqpoint{-0.013889in}{-0.003683in}}{\pgfqpoint{-0.012425in}{-0.007216in}}{\pgfqpoint{-0.009821in}{-0.009821in}}%
\pgfpathcurveto{\pgfqpoint{-0.007216in}{-0.012425in}}{\pgfqpoint{-0.003683in}{-0.013889in}}{\pgfqpoint{0.000000in}{-0.013889in}}%
\pgfpathlineto{\pgfqpoint{0.000000in}{-0.013889in}}%
\pgfpathclose%
\pgfusepath{stroke,fill}%
}%
\begin{pgfscope}%
\pgfsys@transformshift{4.099020in}{3.169954in}%
\pgfsys@useobject{currentmarker}{}%
\end{pgfscope}%
\end{pgfscope}%
\begin{pgfscope}%
\pgfpathrectangle{\pgfqpoint{0.100000in}{2.493063in}}{\pgfqpoint{5.037500in}{3.427208in}}%
\pgfusepath{clip}%
\pgfsetrectcap%
\pgfsetroundjoin%
\pgfsetlinewidth{1.505625pt}%
\definecolor{currentstroke}{rgb}{0.501961,0.501961,0.501961}%
\pgfsetstrokecolor{currentstroke}%
\pgfsetstrokeopacity{0.500000}%
\pgfsetdash{}{0pt}%
\pgfpathmoveto{\pgfqpoint{4.134085in}{3.293055in}}%
\pgfusepath{stroke}%
\end{pgfscope}%
\begin{pgfscope}%
\pgfpathrectangle{\pgfqpoint{0.100000in}{2.493063in}}{\pgfqpoint{5.037500in}{3.427208in}}%
\pgfusepath{clip}%
\pgfsetbuttcap%
\pgfsetroundjoin%
\definecolor{currentfill}{rgb}{0.501961,0.501961,0.501961}%
\pgfsetfillcolor{currentfill}%
\pgfsetfillopacity{0.500000}%
\pgfsetlinewidth{0.250937pt}%
\definecolor{currentstroke}{rgb}{0.000000,0.000000,0.000000}%
\pgfsetstrokecolor{currentstroke}%
\pgfsetstrokeopacity{0.500000}%
\pgfsetdash{}{0pt}%
\pgfsys@defobject{currentmarker}{\pgfqpoint{-0.013889in}{-0.013889in}}{\pgfqpoint{0.013889in}{0.013889in}}{%
\pgfpathmoveto{\pgfqpoint{0.000000in}{-0.013889in}}%
\pgfpathcurveto{\pgfqpoint{0.003683in}{-0.013889in}}{\pgfqpoint{0.007216in}{-0.012425in}}{\pgfqpoint{0.009821in}{-0.009821in}}%
\pgfpathcurveto{\pgfqpoint{0.012425in}{-0.007216in}}{\pgfqpoint{0.013889in}{-0.003683in}}{\pgfqpoint{0.013889in}{0.000000in}}%
\pgfpathcurveto{\pgfqpoint{0.013889in}{0.003683in}}{\pgfqpoint{0.012425in}{0.007216in}}{\pgfqpoint{0.009821in}{0.009821in}}%
\pgfpathcurveto{\pgfqpoint{0.007216in}{0.012425in}}{\pgfqpoint{0.003683in}{0.013889in}}{\pgfqpoint{0.000000in}{0.013889in}}%
\pgfpathcurveto{\pgfqpoint{-0.003683in}{0.013889in}}{\pgfqpoint{-0.007216in}{0.012425in}}{\pgfqpoint{-0.009821in}{0.009821in}}%
\pgfpathcurveto{\pgfqpoint{-0.012425in}{0.007216in}}{\pgfqpoint{-0.013889in}{0.003683in}}{\pgfqpoint{-0.013889in}{0.000000in}}%
\pgfpathcurveto{\pgfqpoint{-0.013889in}{-0.003683in}}{\pgfqpoint{-0.012425in}{-0.007216in}}{\pgfqpoint{-0.009821in}{-0.009821in}}%
\pgfpathcurveto{\pgfqpoint{-0.007216in}{-0.012425in}}{\pgfqpoint{-0.003683in}{-0.013889in}}{\pgfqpoint{0.000000in}{-0.013889in}}%
\pgfpathlineto{\pgfqpoint{0.000000in}{-0.013889in}}%
\pgfpathclose%
\pgfusepath{stroke,fill}%
}%
\begin{pgfscope}%
\pgfsys@transformshift{4.134085in}{3.293055in}%
\pgfsys@useobject{currentmarker}{}%
\end{pgfscope}%
\end{pgfscope}%
\begin{pgfscope}%
\pgfpathrectangle{\pgfqpoint{0.100000in}{2.493063in}}{\pgfqpoint{5.037500in}{3.427208in}}%
\pgfusepath{clip}%
\pgfsetrectcap%
\pgfsetroundjoin%
\pgfsetlinewidth{1.505625pt}%
\definecolor{currentstroke}{rgb}{0.678431,1.000000,0.184314}%
\pgfsetstrokecolor{currentstroke}%
\pgfsetstrokeopacity{0.500000}%
\pgfsetdash{}{0pt}%
\pgfpathmoveto{\pgfqpoint{3.874269in}{3.567855in}}%
\pgfusepath{stroke}%
\end{pgfscope}%
\begin{pgfscope}%
\pgfpathrectangle{\pgfqpoint{0.100000in}{2.493063in}}{\pgfqpoint{5.037500in}{3.427208in}}%
\pgfusepath{clip}%
\pgfsetbuttcap%
\pgfsetroundjoin%
\definecolor{currentfill}{rgb}{0.678431,1.000000,0.184314}%
\pgfsetfillcolor{currentfill}%
\pgfsetfillopacity{0.500000}%
\pgfsetlinewidth{0.250937pt}%
\definecolor{currentstroke}{rgb}{0.000000,0.000000,0.000000}%
\pgfsetstrokecolor{currentstroke}%
\pgfsetstrokeopacity{0.500000}%
\pgfsetdash{}{0pt}%
\pgfsys@defobject{currentmarker}{\pgfqpoint{-0.008333in}{-0.008333in}}{\pgfqpoint{0.008333in}{0.008333in}}{%
\pgfpathmoveto{\pgfqpoint{0.000000in}{-0.008333in}}%
\pgfpathcurveto{\pgfqpoint{0.002210in}{-0.008333in}}{\pgfqpoint{0.004330in}{-0.007455in}}{\pgfqpoint{0.005893in}{-0.005893in}}%
\pgfpathcurveto{\pgfqpoint{0.007455in}{-0.004330in}}{\pgfqpoint{0.008333in}{-0.002210in}}{\pgfqpoint{0.008333in}{0.000000in}}%
\pgfpathcurveto{\pgfqpoint{0.008333in}{0.002210in}}{\pgfqpoint{0.007455in}{0.004330in}}{\pgfqpoint{0.005893in}{0.005893in}}%
\pgfpathcurveto{\pgfqpoint{0.004330in}{0.007455in}}{\pgfqpoint{0.002210in}{0.008333in}}{\pgfqpoint{0.000000in}{0.008333in}}%
\pgfpathcurveto{\pgfqpoint{-0.002210in}{0.008333in}}{\pgfqpoint{-0.004330in}{0.007455in}}{\pgfqpoint{-0.005893in}{0.005893in}}%
\pgfpathcurveto{\pgfqpoint{-0.007455in}{0.004330in}}{\pgfqpoint{-0.008333in}{0.002210in}}{\pgfqpoint{-0.008333in}{0.000000in}}%
\pgfpathcurveto{\pgfqpoint{-0.008333in}{-0.002210in}}{\pgfqpoint{-0.007455in}{-0.004330in}}{\pgfqpoint{-0.005893in}{-0.005893in}}%
\pgfpathcurveto{\pgfqpoint{-0.004330in}{-0.007455in}}{\pgfqpoint{-0.002210in}{-0.008333in}}{\pgfqpoint{0.000000in}{-0.008333in}}%
\pgfpathlineto{\pgfqpoint{0.000000in}{-0.008333in}}%
\pgfpathclose%
\pgfusepath{stroke,fill}%
}%
\begin{pgfscope}%
\pgfsys@transformshift{3.874269in}{3.567855in}%
\pgfsys@useobject{currentmarker}{}%
\end{pgfscope}%
\end{pgfscope}%
\begin{pgfscope}%
\pgfpathrectangle{\pgfqpoint{0.100000in}{2.493063in}}{\pgfqpoint{5.037500in}{3.427208in}}%
\pgfusepath{clip}%
\pgfsetrectcap%
\pgfsetroundjoin%
\pgfsetlinewidth{1.505625pt}%
\definecolor{currentstroke}{rgb}{0.501961,0.501961,0.501961}%
\pgfsetstrokecolor{currentstroke}%
\pgfsetstrokeopacity{0.500000}%
\pgfsetdash{}{0pt}%
\pgfpathmoveto{\pgfqpoint{3.914523in}{3.850303in}}%
\pgfusepath{stroke}%
\end{pgfscope}%
\begin{pgfscope}%
\pgfpathrectangle{\pgfqpoint{0.100000in}{2.493063in}}{\pgfqpoint{5.037500in}{3.427208in}}%
\pgfusepath{clip}%
\pgfsetbuttcap%
\pgfsetroundjoin%
\definecolor{currentfill}{rgb}{0.501961,0.501961,0.501961}%
\pgfsetfillcolor{currentfill}%
\pgfsetfillopacity{0.500000}%
\pgfsetlinewidth{0.250937pt}%
\definecolor{currentstroke}{rgb}{0.000000,0.000000,0.000000}%
\pgfsetstrokecolor{currentstroke}%
\pgfsetstrokeopacity{0.500000}%
\pgfsetdash{}{0pt}%
\pgfsys@defobject{currentmarker}{\pgfqpoint{-0.013889in}{-0.013889in}}{\pgfqpoint{0.013889in}{0.013889in}}{%
\pgfpathmoveto{\pgfqpoint{0.000000in}{-0.013889in}}%
\pgfpathcurveto{\pgfqpoint{0.003683in}{-0.013889in}}{\pgfqpoint{0.007216in}{-0.012425in}}{\pgfqpoint{0.009821in}{-0.009821in}}%
\pgfpathcurveto{\pgfqpoint{0.012425in}{-0.007216in}}{\pgfqpoint{0.013889in}{-0.003683in}}{\pgfqpoint{0.013889in}{0.000000in}}%
\pgfpathcurveto{\pgfqpoint{0.013889in}{0.003683in}}{\pgfqpoint{0.012425in}{0.007216in}}{\pgfqpoint{0.009821in}{0.009821in}}%
\pgfpathcurveto{\pgfqpoint{0.007216in}{0.012425in}}{\pgfqpoint{0.003683in}{0.013889in}}{\pgfqpoint{0.000000in}{0.013889in}}%
\pgfpathcurveto{\pgfqpoint{-0.003683in}{0.013889in}}{\pgfqpoint{-0.007216in}{0.012425in}}{\pgfqpoint{-0.009821in}{0.009821in}}%
\pgfpathcurveto{\pgfqpoint{-0.012425in}{0.007216in}}{\pgfqpoint{-0.013889in}{0.003683in}}{\pgfqpoint{-0.013889in}{0.000000in}}%
\pgfpathcurveto{\pgfqpoint{-0.013889in}{-0.003683in}}{\pgfqpoint{-0.012425in}{-0.007216in}}{\pgfqpoint{-0.009821in}{-0.009821in}}%
\pgfpathcurveto{\pgfqpoint{-0.007216in}{-0.012425in}}{\pgfqpoint{-0.003683in}{-0.013889in}}{\pgfqpoint{0.000000in}{-0.013889in}}%
\pgfpathlineto{\pgfqpoint{0.000000in}{-0.013889in}}%
\pgfpathclose%
\pgfusepath{stroke,fill}%
}%
\begin{pgfscope}%
\pgfsys@transformshift{3.914523in}{3.850303in}%
\pgfsys@useobject{currentmarker}{}%
\end{pgfscope}%
\end{pgfscope}%
\begin{pgfscope}%
\pgfpathrectangle{\pgfqpoint{0.100000in}{2.493063in}}{\pgfqpoint{5.037500in}{3.427208in}}%
\pgfusepath{clip}%
\pgfsetrectcap%
\pgfsetroundjoin%
\pgfsetlinewidth{1.505625pt}%
\definecolor{currentstroke}{rgb}{0.000000,0.000000,1.000000}%
\pgfsetstrokecolor{currentstroke}%
\pgfsetstrokeopacity{0.500000}%
\pgfsetdash{}{0pt}%
\pgfpathmoveto{\pgfqpoint{3.821936in}{3.814931in}}%
\pgfusepath{stroke}%
\end{pgfscope}%
\begin{pgfscope}%
\pgfpathrectangle{\pgfqpoint{0.100000in}{2.493063in}}{\pgfqpoint{5.037500in}{3.427208in}}%
\pgfusepath{clip}%
\pgfsetbuttcap%
\pgfsetroundjoin%
\definecolor{currentfill}{rgb}{0.000000,0.000000,1.000000}%
\pgfsetfillcolor{currentfill}%
\pgfsetfillopacity{0.500000}%
\pgfsetlinewidth{0.250937pt}%
\definecolor{currentstroke}{rgb}{0.000000,0.000000,0.000000}%
\pgfsetstrokecolor{currentstroke}%
\pgfsetstrokeopacity{0.500000}%
\pgfsetdash{}{0pt}%
\pgfsys@defobject{currentmarker}{\pgfqpoint{-0.005556in}{-0.005556in}}{\pgfqpoint{0.005556in}{0.005556in}}{%
\pgfpathmoveto{\pgfqpoint{0.000000in}{-0.005556in}}%
\pgfpathcurveto{\pgfqpoint{0.001473in}{-0.005556in}}{\pgfqpoint{0.002887in}{-0.004970in}}{\pgfqpoint{0.003928in}{-0.003928in}}%
\pgfpathcurveto{\pgfqpoint{0.004970in}{-0.002887in}}{\pgfqpoint{0.005556in}{-0.001473in}}{\pgfqpoint{0.005556in}{0.000000in}}%
\pgfpathcurveto{\pgfqpoint{0.005556in}{0.001473in}}{\pgfqpoint{0.004970in}{0.002887in}}{\pgfqpoint{0.003928in}{0.003928in}}%
\pgfpathcurveto{\pgfqpoint{0.002887in}{0.004970in}}{\pgfqpoint{0.001473in}{0.005556in}}{\pgfqpoint{0.000000in}{0.005556in}}%
\pgfpathcurveto{\pgfqpoint{-0.001473in}{0.005556in}}{\pgfqpoint{-0.002887in}{0.004970in}}{\pgfqpoint{-0.003928in}{0.003928in}}%
\pgfpathcurveto{\pgfqpoint{-0.004970in}{0.002887in}}{\pgfqpoint{-0.005556in}{0.001473in}}{\pgfqpoint{-0.005556in}{0.000000in}}%
\pgfpathcurveto{\pgfqpoint{-0.005556in}{-0.001473in}}{\pgfqpoint{-0.004970in}{-0.002887in}}{\pgfqpoint{-0.003928in}{-0.003928in}}%
\pgfpathcurveto{\pgfqpoint{-0.002887in}{-0.004970in}}{\pgfqpoint{-0.001473in}{-0.005556in}}{\pgfqpoint{0.000000in}{-0.005556in}}%
\pgfpathlineto{\pgfqpoint{0.000000in}{-0.005556in}}%
\pgfpathclose%
\pgfusepath{stroke,fill}%
}%
\begin{pgfscope}%
\pgfsys@transformshift{3.821936in}{3.814931in}%
\pgfsys@useobject{currentmarker}{}%
\end{pgfscope}%
\end{pgfscope}%
\begin{pgfscope}%
\pgfpathrectangle{\pgfqpoint{0.100000in}{2.493063in}}{\pgfqpoint{5.037500in}{3.427208in}}%
\pgfusepath{clip}%
\pgfsetrectcap%
\pgfsetroundjoin%
\pgfsetlinewidth{1.505625pt}%
\definecolor{currentstroke}{rgb}{0.000000,0.000000,1.000000}%
\pgfsetstrokecolor{currentstroke}%
\pgfsetstrokeopacity{0.500000}%
\pgfsetdash{}{0pt}%
\pgfpathmoveto{\pgfqpoint{4.056797in}{3.813205in}}%
\pgfusepath{stroke}%
\end{pgfscope}%
\begin{pgfscope}%
\pgfpathrectangle{\pgfqpoint{0.100000in}{2.493063in}}{\pgfqpoint{5.037500in}{3.427208in}}%
\pgfusepath{clip}%
\pgfsetbuttcap%
\pgfsetroundjoin%
\definecolor{currentfill}{rgb}{0.000000,0.000000,1.000000}%
\pgfsetfillcolor{currentfill}%
\pgfsetfillopacity{0.500000}%
\pgfsetlinewidth{0.250937pt}%
\definecolor{currentstroke}{rgb}{0.000000,0.000000,0.000000}%
\pgfsetstrokecolor{currentstroke}%
\pgfsetstrokeopacity{0.500000}%
\pgfsetdash{}{0pt}%
\pgfsys@defobject{currentmarker}{\pgfqpoint{-0.008333in}{-0.008333in}}{\pgfqpoint{0.008333in}{0.008333in}}{%
\pgfpathmoveto{\pgfqpoint{0.000000in}{-0.008333in}}%
\pgfpathcurveto{\pgfqpoint{0.002210in}{-0.008333in}}{\pgfqpoint{0.004330in}{-0.007455in}}{\pgfqpoint{0.005893in}{-0.005893in}}%
\pgfpathcurveto{\pgfqpoint{0.007455in}{-0.004330in}}{\pgfqpoint{0.008333in}{-0.002210in}}{\pgfqpoint{0.008333in}{0.000000in}}%
\pgfpathcurveto{\pgfqpoint{0.008333in}{0.002210in}}{\pgfqpoint{0.007455in}{0.004330in}}{\pgfqpoint{0.005893in}{0.005893in}}%
\pgfpathcurveto{\pgfqpoint{0.004330in}{0.007455in}}{\pgfqpoint{0.002210in}{0.008333in}}{\pgfqpoint{0.000000in}{0.008333in}}%
\pgfpathcurveto{\pgfqpoint{-0.002210in}{0.008333in}}{\pgfqpoint{-0.004330in}{0.007455in}}{\pgfqpoint{-0.005893in}{0.005893in}}%
\pgfpathcurveto{\pgfqpoint{-0.007455in}{0.004330in}}{\pgfqpoint{-0.008333in}{0.002210in}}{\pgfqpoint{-0.008333in}{0.000000in}}%
\pgfpathcurveto{\pgfqpoint{-0.008333in}{-0.002210in}}{\pgfqpoint{-0.007455in}{-0.004330in}}{\pgfqpoint{-0.005893in}{-0.005893in}}%
\pgfpathcurveto{\pgfqpoint{-0.004330in}{-0.007455in}}{\pgfqpoint{-0.002210in}{-0.008333in}}{\pgfqpoint{0.000000in}{-0.008333in}}%
\pgfpathlineto{\pgfqpoint{0.000000in}{-0.008333in}}%
\pgfpathclose%
\pgfusepath{stroke,fill}%
}%
\begin{pgfscope}%
\pgfsys@transformshift{4.056797in}{3.813205in}%
\pgfsys@useobject{currentmarker}{}%
\end{pgfscope}%
\end{pgfscope}%
\begin{pgfscope}%
\pgfpathrectangle{\pgfqpoint{0.100000in}{2.493063in}}{\pgfqpoint{5.037500in}{3.427208in}}%
\pgfusepath{clip}%
\pgfsetrectcap%
\pgfsetroundjoin%
\pgfsetlinewidth{1.505625pt}%
\definecolor{currentstroke}{rgb}{0.678431,1.000000,0.184314}%
\pgfsetstrokecolor{currentstroke}%
\pgfsetstrokeopacity{0.500000}%
\pgfsetdash{}{0pt}%
\pgfpathmoveto{\pgfqpoint{4.142799in}{3.553859in}}%
\pgfusepath{stroke}%
\end{pgfscope}%
\begin{pgfscope}%
\pgfpathrectangle{\pgfqpoint{0.100000in}{2.493063in}}{\pgfqpoint{5.037500in}{3.427208in}}%
\pgfusepath{clip}%
\pgfsetbuttcap%
\pgfsetroundjoin%
\definecolor{currentfill}{rgb}{0.678431,1.000000,0.184314}%
\pgfsetfillcolor{currentfill}%
\pgfsetfillopacity{0.500000}%
\pgfsetlinewidth{0.250937pt}%
\definecolor{currentstroke}{rgb}{0.000000,0.000000,0.000000}%
\pgfsetstrokecolor{currentstroke}%
\pgfsetstrokeopacity{0.500000}%
\pgfsetdash{}{0pt}%
\pgfsys@defobject{currentmarker}{\pgfqpoint{-0.011111in}{-0.011111in}}{\pgfqpoint{0.011111in}{0.011111in}}{%
\pgfpathmoveto{\pgfqpoint{0.000000in}{-0.011111in}}%
\pgfpathcurveto{\pgfqpoint{0.002947in}{-0.011111in}}{\pgfqpoint{0.005773in}{-0.009940in}}{\pgfqpoint{0.007857in}{-0.007857in}}%
\pgfpathcurveto{\pgfqpoint{0.009940in}{-0.005773in}}{\pgfqpoint{0.011111in}{-0.002947in}}{\pgfqpoint{0.011111in}{0.000000in}}%
\pgfpathcurveto{\pgfqpoint{0.011111in}{0.002947in}}{\pgfqpoint{0.009940in}{0.005773in}}{\pgfqpoint{0.007857in}{0.007857in}}%
\pgfpathcurveto{\pgfqpoint{0.005773in}{0.009940in}}{\pgfqpoint{0.002947in}{0.011111in}}{\pgfqpoint{0.000000in}{0.011111in}}%
\pgfpathcurveto{\pgfqpoint{-0.002947in}{0.011111in}}{\pgfqpoint{-0.005773in}{0.009940in}}{\pgfqpoint{-0.007857in}{0.007857in}}%
\pgfpathcurveto{\pgfqpoint{-0.009940in}{0.005773in}}{\pgfqpoint{-0.011111in}{0.002947in}}{\pgfqpoint{-0.011111in}{0.000000in}}%
\pgfpathcurveto{\pgfqpoint{-0.011111in}{-0.002947in}}{\pgfqpoint{-0.009940in}{-0.005773in}}{\pgfqpoint{-0.007857in}{-0.007857in}}%
\pgfpathcurveto{\pgfqpoint{-0.005773in}{-0.009940in}}{\pgfqpoint{-0.002947in}{-0.011111in}}{\pgfqpoint{0.000000in}{-0.011111in}}%
\pgfpathlineto{\pgfqpoint{0.000000in}{-0.011111in}}%
\pgfpathclose%
\pgfusepath{stroke,fill}%
}%
\begin{pgfscope}%
\pgfsys@transformshift{4.142799in}{3.553859in}%
\pgfsys@useobject{currentmarker}{}%
\end{pgfscope}%
\end{pgfscope}%
\begin{pgfscope}%
\pgfpathrectangle{\pgfqpoint{0.100000in}{2.493063in}}{\pgfqpoint{5.037500in}{3.427208in}}%
\pgfusepath{clip}%
\pgfsetrectcap%
\pgfsetroundjoin%
\pgfsetlinewidth{1.505625pt}%
\definecolor{currentstroke}{rgb}{0.678431,1.000000,0.184314}%
\pgfsetstrokecolor{currentstroke}%
\pgfsetstrokeopacity{0.500000}%
\pgfsetdash{}{0pt}%
\pgfpathmoveto{\pgfqpoint{3.781250in}{3.660214in}}%
\pgfusepath{stroke}%
\end{pgfscope}%
\begin{pgfscope}%
\pgfpathrectangle{\pgfqpoint{0.100000in}{2.493063in}}{\pgfqpoint{5.037500in}{3.427208in}}%
\pgfusepath{clip}%
\pgfsetbuttcap%
\pgfsetroundjoin%
\definecolor{currentfill}{rgb}{0.678431,1.000000,0.184314}%
\pgfsetfillcolor{currentfill}%
\pgfsetfillopacity{0.500000}%
\pgfsetlinewidth{0.250937pt}%
\definecolor{currentstroke}{rgb}{0.000000,0.000000,0.000000}%
\pgfsetstrokecolor{currentstroke}%
\pgfsetstrokeopacity{0.500000}%
\pgfsetdash{}{0pt}%
\pgfsys@defobject{currentmarker}{\pgfqpoint{-0.005556in}{-0.005556in}}{\pgfqpoint{0.005556in}{0.005556in}}{%
\pgfpathmoveto{\pgfqpoint{0.000000in}{-0.005556in}}%
\pgfpathcurveto{\pgfqpoint{0.001473in}{-0.005556in}}{\pgfqpoint{0.002887in}{-0.004970in}}{\pgfqpoint{0.003928in}{-0.003928in}}%
\pgfpathcurveto{\pgfqpoint{0.004970in}{-0.002887in}}{\pgfqpoint{0.005556in}{-0.001473in}}{\pgfqpoint{0.005556in}{0.000000in}}%
\pgfpathcurveto{\pgfqpoint{0.005556in}{0.001473in}}{\pgfqpoint{0.004970in}{0.002887in}}{\pgfqpoint{0.003928in}{0.003928in}}%
\pgfpathcurveto{\pgfqpoint{0.002887in}{0.004970in}}{\pgfqpoint{0.001473in}{0.005556in}}{\pgfqpoint{0.000000in}{0.005556in}}%
\pgfpathcurveto{\pgfqpoint{-0.001473in}{0.005556in}}{\pgfqpoint{-0.002887in}{0.004970in}}{\pgfqpoint{-0.003928in}{0.003928in}}%
\pgfpathcurveto{\pgfqpoint{-0.004970in}{0.002887in}}{\pgfqpoint{-0.005556in}{0.001473in}}{\pgfqpoint{-0.005556in}{0.000000in}}%
\pgfpathcurveto{\pgfqpoint{-0.005556in}{-0.001473in}}{\pgfqpoint{-0.004970in}{-0.002887in}}{\pgfqpoint{-0.003928in}{-0.003928in}}%
\pgfpathcurveto{\pgfqpoint{-0.002887in}{-0.004970in}}{\pgfqpoint{-0.001473in}{-0.005556in}}{\pgfqpoint{0.000000in}{-0.005556in}}%
\pgfpathlineto{\pgfqpoint{0.000000in}{-0.005556in}}%
\pgfpathclose%
\pgfusepath{stroke,fill}%
}%
\begin{pgfscope}%
\pgfsys@transformshift{3.781250in}{3.660214in}%
\pgfsys@useobject{currentmarker}{}%
\end{pgfscope}%
\end{pgfscope}%
\begin{pgfscope}%
\pgfpathrectangle{\pgfqpoint{0.100000in}{2.493063in}}{\pgfqpoint{5.037500in}{3.427208in}}%
\pgfusepath{clip}%
\pgfsetrectcap%
\pgfsetroundjoin%
\pgfsetlinewidth{1.505625pt}%
\definecolor{currentstroke}{rgb}{0.678431,1.000000,0.184314}%
\pgfsetstrokecolor{currentstroke}%
\pgfsetstrokeopacity{0.500000}%
\pgfsetdash{}{0pt}%
\pgfpathmoveto{\pgfqpoint{3.763715in}{3.921616in}}%
\pgfusepath{stroke}%
\end{pgfscope}%
\begin{pgfscope}%
\pgfpathrectangle{\pgfqpoint{0.100000in}{2.493063in}}{\pgfqpoint{5.037500in}{3.427208in}}%
\pgfusepath{clip}%
\pgfsetbuttcap%
\pgfsetroundjoin%
\definecolor{currentfill}{rgb}{0.678431,1.000000,0.184314}%
\pgfsetfillcolor{currentfill}%
\pgfsetfillopacity{0.500000}%
\pgfsetlinewidth{0.250937pt}%
\definecolor{currentstroke}{rgb}{0.000000,0.000000,0.000000}%
\pgfsetstrokecolor{currentstroke}%
\pgfsetstrokeopacity{0.500000}%
\pgfsetdash{}{0pt}%
\pgfsys@defobject{currentmarker}{\pgfqpoint{-0.022222in}{-0.022222in}}{\pgfqpoint{0.022222in}{0.022222in}}{%
\pgfpathmoveto{\pgfqpoint{0.000000in}{-0.022222in}}%
\pgfpathcurveto{\pgfqpoint{0.005893in}{-0.022222in}}{\pgfqpoint{0.011546in}{-0.019881in}}{\pgfqpoint{0.015713in}{-0.015713in}}%
\pgfpathcurveto{\pgfqpoint{0.019881in}{-0.011546in}}{\pgfqpoint{0.022222in}{-0.005893in}}{\pgfqpoint{0.022222in}{0.000000in}}%
\pgfpathcurveto{\pgfqpoint{0.022222in}{0.005893in}}{\pgfqpoint{0.019881in}{0.011546in}}{\pgfqpoint{0.015713in}{0.015713in}}%
\pgfpathcurveto{\pgfqpoint{0.011546in}{0.019881in}}{\pgfqpoint{0.005893in}{0.022222in}}{\pgfqpoint{0.000000in}{0.022222in}}%
\pgfpathcurveto{\pgfqpoint{-0.005893in}{0.022222in}}{\pgfqpoint{-0.011546in}{0.019881in}}{\pgfqpoint{-0.015713in}{0.015713in}}%
\pgfpathcurveto{\pgfqpoint{-0.019881in}{0.011546in}}{\pgfqpoint{-0.022222in}{0.005893in}}{\pgfqpoint{-0.022222in}{0.000000in}}%
\pgfpathcurveto{\pgfqpoint{-0.022222in}{-0.005893in}}{\pgfqpoint{-0.019881in}{-0.011546in}}{\pgfqpoint{-0.015713in}{-0.015713in}}%
\pgfpathcurveto{\pgfqpoint{-0.011546in}{-0.019881in}}{\pgfqpoint{-0.005893in}{-0.022222in}}{\pgfqpoint{0.000000in}{-0.022222in}}%
\pgfpathlineto{\pgfqpoint{0.000000in}{-0.022222in}}%
\pgfpathclose%
\pgfusepath{stroke,fill}%
}%
\begin{pgfscope}%
\pgfsys@transformshift{3.763715in}{3.921616in}%
\pgfsys@useobject{currentmarker}{}%
\end{pgfscope}%
\end{pgfscope}%
\begin{pgfscope}%
\pgfpathrectangle{\pgfqpoint{0.100000in}{2.493063in}}{\pgfqpoint{5.037500in}{3.427208in}}%
\pgfusepath{clip}%
\pgfsetrectcap%
\pgfsetroundjoin%
\pgfsetlinewidth{1.505625pt}%
\definecolor{currentstroke}{rgb}{0.678431,1.000000,0.184314}%
\pgfsetstrokecolor{currentstroke}%
\pgfsetstrokeopacity{0.500000}%
\pgfsetdash{}{0pt}%
\pgfpathmoveto{\pgfqpoint{3.868196in}{3.884393in}}%
\pgfusepath{stroke}%
\end{pgfscope}%
\begin{pgfscope}%
\pgfpathrectangle{\pgfqpoint{0.100000in}{2.493063in}}{\pgfqpoint{5.037500in}{3.427208in}}%
\pgfusepath{clip}%
\pgfsetbuttcap%
\pgfsetroundjoin%
\definecolor{currentfill}{rgb}{0.678431,1.000000,0.184314}%
\pgfsetfillcolor{currentfill}%
\pgfsetfillopacity{0.500000}%
\pgfsetlinewidth{0.250937pt}%
\definecolor{currentstroke}{rgb}{0.000000,0.000000,0.000000}%
\pgfsetstrokecolor{currentstroke}%
\pgfsetstrokeopacity{0.500000}%
\pgfsetdash{}{0pt}%
\pgfsys@defobject{currentmarker}{\pgfqpoint{-0.008333in}{-0.008333in}}{\pgfqpoint{0.008333in}{0.008333in}}{%
\pgfpathmoveto{\pgfqpoint{0.000000in}{-0.008333in}}%
\pgfpathcurveto{\pgfqpoint{0.002210in}{-0.008333in}}{\pgfqpoint{0.004330in}{-0.007455in}}{\pgfqpoint{0.005893in}{-0.005893in}}%
\pgfpathcurveto{\pgfqpoint{0.007455in}{-0.004330in}}{\pgfqpoint{0.008333in}{-0.002210in}}{\pgfqpoint{0.008333in}{0.000000in}}%
\pgfpathcurveto{\pgfqpoint{0.008333in}{0.002210in}}{\pgfqpoint{0.007455in}{0.004330in}}{\pgfqpoint{0.005893in}{0.005893in}}%
\pgfpathcurveto{\pgfqpoint{0.004330in}{0.007455in}}{\pgfqpoint{0.002210in}{0.008333in}}{\pgfqpoint{0.000000in}{0.008333in}}%
\pgfpathcurveto{\pgfqpoint{-0.002210in}{0.008333in}}{\pgfqpoint{-0.004330in}{0.007455in}}{\pgfqpoint{-0.005893in}{0.005893in}}%
\pgfpathcurveto{\pgfqpoint{-0.007455in}{0.004330in}}{\pgfqpoint{-0.008333in}{0.002210in}}{\pgfqpoint{-0.008333in}{0.000000in}}%
\pgfpathcurveto{\pgfqpoint{-0.008333in}{-0.002210in}}{\pgfqpoint{-0.007455in}{-0.004330in}}{\pgfqpoint{-0.005893in}{-0.005893in}}%
\pgfpathcurveto{\pgfqpoint{-0.004330in}{-0.007455in}}{\pgfqpoint{-0.002210in}{-0.008333in}}{\pgfqpoint{0.000000in}{-0.008333in}}%
\pgfpathlineto{\pgfqpoint{0.000000in}{-0.008333in}}%
\pgfpathclose%
\pgfusepath{stroke,fill}%
}%
\begin{pgfscope}%
\pgfsys@transformshift{3.868196in}{3.884393in}%
\pgfsys@useobject{currentmarker}{}%
\end{pgfscope}%
\end{pgfscope}%
\begin{pgfscope}%
\pgfpathrectangle{\pgfqpoint{0.100000in}{2.493063in}}{\pgfqpoint{5.037500in}{3.427208in}}%
\pgfusepath{clip}%
\pgfsetrectcap%
\pgfsetroundjoin%
\pgfsetlinewidth{1.505625pt}%
\definecolor{currentstroke}{rgb}{0.678431,1.000000,0.184314}%
\pgfsetstrokecolor{currentstroke}%
\pgfsetstrokeopacity{0.500000}%
\pgfsetdash{}{0pt}%
\pgfpathmoveto{\pgfqpoint{4.120605in}{3.632479in}}%
\pgfusepath{stroke}%
\end{pgfscope}%
\begin{pgfscope}%
\pgfpathrectangle{\pgfqpoint{0.100000in}{2.493063in}}{\pgfqpoint{5.037500in}{3.427208in}}%
\pgfusepath{clip}%
\pgfsetbuttcap%
\pgfsetroundjoin%
\definecolor{currentfill}{rgb}{0.678431,1.000000,0.184314}%
\pgfsetfillcolor{currentfill}%
\pgfsetfillopacity{0.500000}%
\pgfsetlinewidth{0.250937pt}%
\definecolor{currentstroke}{rgb}{0.000000,0.000000,0.000000}%
\pgfsetstrokecolor{currentstroke}%
\pgfsetstrokeopacity{0.500000}%
\pgfsetdash{}{0pt}%
\pgfsys@defobject{currentmarker}{\pgfqpoint{-0.011111in}{-0.011111in}}{\pgfqpoint{0.011111in}{0.011111in}}{%
\pgfpathmoveto{\pgfqpoint{0.000000in}{-0.011111in}}%
\pgfpathcurveto{\pgfqpoint{0.002947in}{-0.011111in}}{\pgfqpoint{0.005773in}{-0.009940in}}{\pgfqpoint{0.007857in}{-0.007857in}}%
\pgfpathcurveto{\pgfqpoint{0.009940in}{-0.005773in}}{\pgfqpoint{0.011111in}{-0.002947in}}{\pgfqpoint{0.011111in}{0.000000in}}%
\pgfpathcurveto{\pgfqpoint{0.011111in}{0.002947in}}{\pgfqpoint{0.009940in}{0.005773in}}{\pgfqpoint{0.007857in}{0.007857in}}%
\pgfpathcurveto{\pgfqpoint{0.005773in}{0.009940in}}{\pgfqpoint{0.002947in}{0.011111in}}{\pgfqpoint{0.000000in}{0.011111in}}%
\pgfpathcurveto{\pgfqpoint{-0.002947in}{0.011111in}}{\pgfqpoint{-0.005773in}{0.009940in}}{\pgfqpoint{-0.007857in}{0.007857in}}%
\pgfpathcurveto{\pgfqpoint{-0.009940in}{0.005773in}}{\pgfqpoint{-0.011111in}{0.002947in}}{\pgfqpoint{-0.011111in}{0.000000in}}%
\pgfpathcurveto{\pgfqpoint{-0.011111in}{-0.002947in}}{\pgfqpoint{-0.009940in}{-0.005773in}}{\pgfqpoint{-0.007857in}{-0.007857in}}%
\pgfpathcurveto{\pgfqpoint{-0.005773in}{-0.009940in}}{\pgfqpoint{-0.002947in}{-0.011111in}}{\pgfqpoint{0.000000in}{-0.011111in}}%
\pgfpathlineto{\pgfqpoint{0.000000in}{-0.011111in}}%
\pgfpathclose%
\pgfusepath{stroke,fill}%
}%
\begin{pgfscope}%
\pgfsys@transformshift{4.120605in}{3.632479in}%
\pgfsys@useobject{currentmarker}{}%
\end{pgfscope}%
\end{pgfscope}%
\begin{pgfscope}%
\pgfpathrectangle{\pgfqpoint{0.100000in}{2.493063in}}{\pgfqpoint{5.037500in}{3.427208in}}%
\pgfusepath{clip}%
\pgfsetrectcap%
\pgfsetroundjoin%
\pgfsetlinewidth{1.505625pt}%
\definecolor{currentstroke}{rgb}{0.501961,0.501961,0.501961}%
\pgfsetstrokecolor{currentstroke}%
\pgfsetstrokeopacity{0.500000}%
\pgfsetdash{}{0pt}%
\pgfpathmoveto{\pgfqpoint{3.907380in}{3.719417in}}%
\pgfusepath{stroke}%
\end{pgfscope}%
\begin{pgfscope}%
\pgfpathrectangle{\pgfqpoint{0.100000in}{2.493063in}}{\pgfqpoint{5.037500in}{3.427208in}}%
\pgfusepath{clip}%
\pgfsetbuttcap%
\pgfsetroundjoin%
\definecolor{currentfill}{rgb}{0.501961,0.501961,0.501961}%
\pgfsetfillcolor{currentfill}%
\pgfsetfillopacity{0.500000}%
\pgfsetlinewidth{0.250937pt}%
\definecolor{currentstroke}{rgb}{0.000000,0.000000,0.000000}%
\pgfsetstrokecolor{currentstroke}%
\pgfsetstrokeopacity{0.500000}%
\pgfsetdash{}{0pt}%
\pgfsys@defobject{currentmarker}{\pgfqpoint{-0.013889in}{-0.013889in}}{\pgfqpoint{0.013889in}{0.013889in}}{%
\pgfpathmoveto{\pgfqpoint{0.000000in}{-0.013889in}}%
\pgfpathcurveto{\pgfqpoint{0.003683in}{-0.013889in}}{\pgfqpoint{0.007216in}{-0.012425in}}{\pgfqpoint{0.009821in}{-0.009821in}}%
\pgfpathcurveto{\pgfqpoint{0.012425in}{-0.007216in}}{\pgfqpoint{0.013889in}{-0.003683in}}{\pgfqpoint{0.013889in}{0.000000in}}%
\pgfpathcurveto{\pgfqpoint{0.013889in}{0.003683in}}{\pgfqpoint{0.012425in}{0.007216in}}{\pgfqpoint{0.009821in}{0.009821in}}%
\pgfpathcurveto{\pgfqpoint{0.007216in}{0.012425in}}{\pgfqpoint{0.003683in}{0.013889in}}{\pgfqpoint{0.000000in}{0.013889in}}%
\pgfpathcurveto{\pgfqpoint{-0.003683in}{0.013889in}}{\pgfqpoint{-0.007216in}{0.012425in}}{\pgfqpoint{-0.009821in}{0.009821in}}%
\pgfpathcurveto{\pgfqpoint{-0.012425in}{0.007216in}}{\pgfqpoint{-0.013889in}{0.003683in}}{\pgfqpoint{-0.013889in}{0.000000in}}%
\pgfpathcurveto{\pgfqpoint{-0.013889in}{-0.003683in}}{\pgfqpoint{-0.012425in}{-0.007216in}}{\pgfqpoint{-0.009821in}{-0.009821in}}%
\pgfpathcurveto{\pgfqpoint{-0.007216in}{-0.012425in}}{\pgfqpoint{-0.003683in}{-0.013889in}}{\pgfqpoint{0.000000in}{-0.013889in}}%
\pgfpathlineto{\pgfqpoint{0.000000in}{-0.013889in}}%
\pgfpathclose%
\pgfusepath{stroke,fill}%
}%
\begin{pgfscope}%
\pgfsys@transformshift{3.907380in}{3.719417in}%
\pgfsys@useobject{currentmarker}{}%
\end{pgfscope}%
\end{pgfscope}%
\begin{pgfscope}%
\pgfpathrectangle{\pgfqpoint{0.100000in}{2.493063in}}{\pgfqpoint{5.037500in}{3.427208in}}%
\pgfusepath{clip}%
\pgfsetrectcap%
\pgfsetroundjoin%
\pgfsetlinewidth{1.505625pt}%
\definecolor{currentstroke}{rgb}{0.678431,1.000000,0.184314}%
\pgfsetstrokecolor{currentstroke}%
\pgfsetstrokeopacity{0.500000}%
\pgfsetdash{}{0pt}%
\pgfpathmoveto{\pgfqpoint{3.741599in}{3.864942in}}%
\pgfusepath{stroke}%
\end{pgfscope}%
\begin{pgfscope}%
\pgfpathrectangle{\pgfqpoint{0.100000in}{2.493063in}}{\pgfqpoint{5.037500in}{3.427208in}}%
\pgfusepath{clip}%
\pgfsetbuttcap%
\pgfsetroundjoin%
\definecolor{currentfill}{rgb}{0.678431,1.000000,0.184314}%
\pgfsetfillcolor{currentfill}%
\pgfsetfillopacity{0.500000}%
\pgfsetlinewidth{0.250937pt}%
\definecolor{currentstroke}{rgb}{0.000000,0.000000,0.000000}%
\pgfsetstrokecolor{currentstroke}%
\pgfsetstrokeopacity{0.500000}%
\pgfsetdash{}{0pt}%
\pgfsys@defobject{currentmarker}{\pgfqpoint{-0.022222in}{-0.022222in}}{\pgfqpoint{0.022222in}{0.022222in}}{%
\pgfpathmoveto{\pgfqpoint{0.000000in}{-0.022222in}}%
\pgfpathcurveto{\pgfqpoint{0.005893in}{-0.022222in}}{\pgfqpoint{0.011546in}{-0.019881in}}{\pgfqpoint{0.015713in}{-0.015713in}}%
\pgfpathcurveto{\pgfqpoint{0.019881in}{-0.011546in}}{\pgfqpoint{0.022222in}{-0.005893in}}{\pgfqpoint{0.022222in}{0.000000in}}%
\pgfpathcurveto{\pgfqpoint{0.022222in}{0.005893in}}{\pgfqpoint{0.019881in}{0.011546in}}{\pgfqpoint{0.015713in}{0.015713in}}%
\pgfpathcurveto{\pgfqpoint{0.011546in}{0.019881in}}{\pgfqpoint{0.005893in}{0.022222in}}{\pgfqpoint{0.000000in}{0.022222in}}%
\pgfpathcurveto{\pgfqpoint{-0.005893in}{0.022222in}}{\pgfqpoint{-0.011546in}{0.019881in}}{\pgfqpoint{-0.015713in}{0.015713in}}%
\pgfpathcurveto{\pgfqpoint{-0.019881in}{0.011546in}}{\pgfqpoint{-0.022222in}{0.005893in}}{\pgfqpoint{-0.022222in}{0.000000in}}%
\pgfpathcurveto{\pgfqpoint{-0.022222in}{-0.005893in}}{\pgfqpoint{-0.019881in}{-0.011546in}}{\pgfqpoint{-0.015713in}{-0.015713in}}%
\pgfpathcurveto{\pgfqpoint{-0.011546in}{-0.019881in}}{\pgfqpoint{-0.005893in}{-0.022222in}}{\pgfqpoint{0.000000in}{-0.022222in}}%
\pgfpathlineto{\pgfqpoint{0.000000in}{-0.022222in}}%
\pgfpathclose%
\pgfusepath{stroke,fill}%
}%
\begin{pgfscope}%
\pgfsys@transformshift{3.741599in}{3.864942in}%
\pgfsys@useobject{currentmarker}{}%
\end{pgfscope}%
\end{pgfscope}%
\begin{pgfscope}%
\pgfpathrectangle{\pgfqpoint{0.100000in}{2.493063in}}{\pgfqpoint{5.037500in}{3.427208in}}%
\pgfusepath{clip}%
\pgfsetrectcap%
\pgfsetroundjoin%
\pgfsetlinewidth{1.505625pt}%
\definecolor{currentstroke}{rgb}{0.501961,0.501961,0.501961}%
\pgfsetstrokecolor{currentstroke}%
\pgfsetstrokeopacity{0.500000}%
\pgfsetdash{}{0pt}%
\pgfpathmoveto{\pgfqpoint{4.164798in}{3.666837in}}%
\pgfusepath{stroke}%
\end{pgfscope}%
\begin{pgfscope}%
\pgfpathrectangle{\pgfqpoint{0.100000in}{2.493063in}}{\pgfqpoint{5.037500in}{3.427208in}}%
\pgfusepath{clip}%
\pgfsetbuttcap%
\pgfsetroundjoin%
\definecolor{currentfill}{rgb}{0.501961,0.501961,0.501961}%
\pgfsetfillcolor{currentfill}%
\pgfsetfillopacity{0.500000}%
\pgfsetlinewidth{0.250937pt}%
\definecolor{currentstroke}{rgb}{0.000000,0.000000,0.000000}%
\pgfsetstrokecolor{currentstroke}%
\pgfsetstrokeopacity{0.500000}%
\pgfsetdash{}{0pt}%
\pgfsys@defobject{currentmarker}{\pgfqpoint{-0.013889in}{-0.013889in}}{\pgfqpoint{0.013889in}{0.013889in}}{%
\pgfpathmoveto{\pgfqpoint{0.000000in}{-0.013889in}}%
\pgfpathcurveto{\pgfqpoint{0.003683in}{-0.013889in}}{\pgfqpoint{0.007216in}{-0.012425in}}{\pgfqpoint{0.009821in}{-0.009821in}}%
\pgfpathcurveto{\pgfqpoint{0.012425in}{-0.007216in}}{\pgfqpoint{0.013889in}{-0.003683in}}{\pgfqpoint{0.013889in}{0.000000in}}%
\pgfpathcurveto{\pgfqpoint{0.013889in}{0.003683in}}{\pgfqpoint{0.012425in}{0.007216in}}{\pgfqpoint{0.009821in}{0.009821in}}%
\pgfpathcurveto{\pgfqpoint{0.007216in}{0.012425in}}{\pgfqpoint{0.003683in}{0.013889in}}{\pgfqpoint{0.000000in}{0.013889in}}%
\pgfpathcurveto{\pgfqpoint{-0.003683in}{0.013889in}}{\pgfqpoint{-0.007216in}{0.012425in}}{\pgfqpoint{-0.009821in}{0.009821in}}%
\pgfpathcurveto{\pgfqpoint{-0.012425in}{0.007216in}}{\pgfqpoint{-0.013889in}{0.003683in}}{\pgfqpoint{-0.013889in}{0.000000in}}%
\pgfpathcurveto{\pgfqpoint{-0.013889in}{-0.003683in}}{\pgfqpoint{-0.012425in}{-0.007216in}}{\pgfqpoint{-0.009821in}{-0.009821in}}%
\pgfpathcurveto{\pgfqpoint{-0.007216in}{-0.012425in}}{\pgfqpoint{-0.003683in}{-0.013889in}}{\pgfqpoint{0.000000in}{-0.013889in}}%
\pgfpathlineto{\pgfqpoint{0.000000in}{-0.013889in}}%
\pgfpathclose%
\pgfusepath{stroke,fill}%
}%
\begin{pgfscope}%
\pgfsys@transformshift{4.164798in}{3.666837in}%
\pgfsys@useobject{currentmarker}{}%
\end{pgfscope}%
\end{pgfscope}%
\begin{pgfscope}%
\pgfpathrectangle{\pgfqpoint{0.100000in}{2.493063in}}{\pgfqpoint{5.037500in}{3.427208in}}%
\pgfusepath{clip}%
\pgfsetrectcap%
\pgfsetroundjoin%
\pgfsetlinewidth{1.505625pt}%
\definecolor{currentstroke}{rgb}{0.501961,0.501961,0.501961}%
\pgfsetstrokecolor{currentstroke}%
\pgfsetstrokeopacity{0.500000}%
\pgfsetdash{}{0pt}%
\pgfpathmoveto{\pgfqpoint{3.971514in}{3.492687in}}%
\pgfusepath{stroke}%
\end{pgfscope}%
\begin{pgfscope}%
\pgfpathrectangle{\pgfqpoint{0.100000in}{2.493063in}}{\pgfqpoint{5.037500in}{3.427208in}}%
\pgfusepath{clip}%
\pgfsetbuttcap%
\pgfsetroundjoin%
\definecolor{currentfill}{rgb}{0.501961,0.501961,0.501961}%
\pgfsetfillcolor{currentfill}%
\pgfsetfillopacity{0.500000}%
\pgfsetlinewidth{0.250937pt}%
\definecolor{currentstroke}{rgb}{0.000000,0.000000,0.000000}%
\pgfsetstrokecolor{currentstroke}%
\pgfsetstrokeopacity{0.500000}%
\pgfsetdash{}{0pt}%
\pgfsys@defobject{currentmarker}{\pgfqpoint{-0.013889in}{-0.013889in}}{\pgfqpoint{0.013889in}{0.013889in}}{%
\pgfpathmoveto{\pgfqpoint{0.000000in}{-0.013889in}}%
\pgfpathcurveto{\pgfqpoint{0.003683in}{-0.013889in}}{\pgfqpoint{0.007216in}{-0.012425in}}{\pgfqpoint{0.009821in}{-0.009821in}}%
\pgfpathcurveto{\pgfqpoint{0.012425in}{-0.007216in}}{\pgfqpoint{0.013889in}{-0.003683in}}{\pgfqpoint{0.013889in}{0.000000in}}%
\pgfpathcurveto{\pgfqpoint{0.013889in}{0.003683in}}{\pgfqpoint{0.012425in}{0.007216in}}{\pgfqpoint{0.009821in}{0.009821in}}%
\pgfpathcurveto{\pgfqpoint{0.007216in}{0.012425in}}{\pgfqpoint{0.003683in}{0.013889in}}{\pgfqpoint{0.000000in}{0.013889in}}%
\pgfpathcurveto{\pgfqpoint{-0.003683in}{0.013889in}}{\pgfqpoint{-0.007216in}{0.012425in}}{\pgfqpoint{-0.009821in}{0.009821in}}%
\pgfpathcurveto{\pgfqpoint{-0.012425in}{0.007216in}}{\pgfqpoint{-0.013889in}{0.003683in}}{\pgfqpoint{-0.013889in}{0.000000in}}%
\pgfpathcurveto{\pgfqpoint{-0.013889in}{-0.003683in}}{\pgfqpoint{-0.012425in}{-0.007216in}}{\pgfqpoint{-0.009821in}{-0.009821in}}%
\pgfpathcurveto{\pgfqpoint{-0.007216in}{-0.012425in}}{\pgfqpoint{-0.003683in}{-0.013889in}}{\pgfqpoint{0.000000in}{-0.013889in}}%
\pgfpathlineto{\pgfqpoint{0.000000in}{-0.013889in}}%
\pgfpathclose%
\pgfusepath{stroke,fill}%
}%
\begin{pgfscope}%
\pgfsys@transformshift{3.971514in}{3.492687in}%
\pgfsys@useobject{currentmarker}{}%
\end{pgfscope}%
\end{pgfscope}%
\begin{pgfscope}%
\pgfpathrectangle{\pgfqpoint{0.100000in}{2.493063in}}{\pgfqpoint{5.037500in}{3.427208in}}%
\pgfusepath{clip}%
\pgfsetrectcap%
\pgfsetroundjoin%
\pgfsetlinewidth{1.505625pt}%
\definecolor{currentstroke}{rgb}{0.678431,1.000000,0.184314}%
\pgfsetstrokecolor{currentstroke}%
\pgfsetstrokeopacity{0.500000}%
\pgfsetdash{}{0pt}%
\pgfpathmoveto{\pgfqpoint{3.911534in}{3.691641in}}%
\pgfusepath{stroke}%
\end{pgfscope}%
\begin{pgfscope}%
\pgfpathrectangle{\pgfqpoint{0.100000in}{2.493063in}}{\pgfqpoint{5.037500in}{3.427208in}}%
\pgfusepath{clip}%
\pgfsetbuttcap%
\pgfsetroundjoin%
\definecolor{currentfill}{rgb}{0.678431,1.000000,0.184314}%
\pgfsetfillcolor{currentfill}%
\pgfsetfillopacity{0.500000}%
\pgfsetlinewidth{0.250937pt}%
\definecolor{currentstroke}{rgb}{0.000000,0.000000,0.000000}%
\pgfsetstrokecolor{currentstroke}%
\pgfsetstrokeopacity{0.500000}%
\pgfsetdash{}{0pt}%
\pgfsys@defobject{currentmarker}{\pgfqpoint{-0.005556in}{-0.005556in}}{\pgfqpoint{0.005556in}{0.005556in}}{%
\pgfpathmoveto{\pgfqpoint{0.000000in}{-0.005556in}}%
\pgfpathcurveto{\pgfqpoint{0.001473in}{-0.005556in}}{\pgfqpoint{0.002887in}{-0.004970in}}{\pgfqpoint{0.003928in}{-0.003928in}}%
\pgfpathcurveto{\pgfqpoint{0.004970in}{-0.002887in}}{\pgfqpoint{0.005556in}{-0.001473in}}{\pgfqpoint{0.005556in}{0.000000in}}%
\pgfpathcurveto{\pgfqpoint{0.005556in}{0.001473in}}{\pgfqpoint{0.004970in}{0.002887in}}{\pgfqpoint{0.003928in}{0.003928in}}%
\pgfpathcurveto{\pgfqpoint{0.002887in}{0.004970in}}{\pgfqpoint{0.001473in}{0.005556in}}{\pgfqpoint{0.000000in}{0.005556in}}%
\pgfpathcurveto{\pgfqpoint{-0.001473in}{0.005556in}}{\pgfqpoint{-0.002887in}{0.004970in}}{\pgfqpoint{-0.003928in}{0.003928in}}%
\pgfpathcurveto{\pgfqpoint{-0.004970in}{0.002887in}}{\pgfqpoint{-0.005556in}{0.001473in}}{\pgfqpoint{-0.005556in}{0.000000in}}%
\pgfpathcurveto{\pgfqpoint{-0.005556in}{-0.001473in}}{\pgfqpoint{-0.004970in}{-0.002887in}}{\pgfqpoint{-0.003928in}{-0.003928in}}%
\pgfpathcurveto{\pgfqpoint{-0.002887in}{-0.004970in}}{\pgfqpoint{-0.001473in}{-0.005556in}}{\pgfqpoint{0.000000in}{-0.005556in}}%
\pgfpathlineto{\pgfqpoint{0.000000in}{-0.005556in}}%
\pgfpathclose%
\pgfusepath{stroke,fill}%
}%
\begin{pgfscope}%
\pgfsys@transformshift{3.911534in}{3.691641in}%
\pgfsys@useobject{currentmarker}{}%
\end{pgfscope}%
\end{pgfscope}%
\begin{pgfscope}%
\pgfpathrectangle{\pgfqpoint{0.100000in}{2.493063in}}{\pgfqpoint{5.037500in}{3.427208in}}%
\pgfusepath{clip}%
\pgfsetrectcap%
\pgfsetroundjoin%
\pgfsetlinewidth{1.505625pt}%
\definecolor{currentstroke}{rgb}{0.000000,0.000000,1.000000}%
\pgfsetstrokecolor{currentstroke}%
\pgfsetstrokeopacity{0.500000}%
\pgfsetdash{}{0pt}%
\pgfpathmoveto{\pgfqpoint{1.800904in}{3.041134in}}%
\pgfusepath{stroke}%
\end{pgfscope}%
\begin{pgfscope}%
\pgfpathrectangle{\pgfqpoint{0.100000in}{2.493063in}}{\pgfqpoint{5.037500in}{3.427208in}}%
\pgfusepath{clip}%
\pgfsetbuttcap%
\pgfsetroundjoin%
\definecolor{currentfill}{rgb}{0.000000,0.000000,1.000000}%
\pgfsetfillcolor{currentfill}%
\pgfsetfillopacity{0.500000}%
\pgfsetlinewidth{0.250937pt}%
\definecolor{currentstroke}{rgb}{0.000000,0.000000,0.000000}%
\pgfsetstrokecolor{currentstroke}%
\pgfsetstrokeopacity{0.500000}%
\pgfsetdash{}{0pt}%
\pgfsys@defobject{currentmarker}{\pgfqpoint{-0.111111in}{-0.111111in}}{\pgfqpoint{0.111111in}{0.111111in}}{%
\pgfpathmoveto{\pgfqpoint{0.000000in}{-0.111111in}}%
\pgfpathcurveto{\pgfqpoint{0.029467in}{-0.111111in}}{\pgfqpoint{0.057731in}{-0.099404in}}{\pgfqpoint{0.078567in}{-0.078567in}}%
\pgfpathcurveto{\pgfqpoint{0.099404in}{-0.057731in}}{\pgfqpoint{0.111111in}{-0.029467in}}{\pgfqpoint{0.111111in}{0.000000in}}%
\pgfpathcurveto{\pgfqpoint{0.111111in}{0.029467in}}{\pgfqpoint{0.099404in}{0.057731in}}{\pgfqpoint{0.078567in}{0.078567in}}%
\pgfpathcurveto{\pgfqpoint{0.057731in}{0.099404in}}{\pgfqpoint{0.029467in}{0.111111in}}{\pgfqpoint{0.000000in}{0.111111in}}%
\pgfpathcurveto{\pgfqpoint{-0.029467in}{0.111111in}}{\pgfqpoint{-0.057731in}{0.099404in}}{\pgfqpoint{-0.078567in}{0.078567in}}%
\pgfpathcurveto{\pgfqpoint{-0.099404in}{0.057731in}}{\pgfqpoint{-0.111111in}{0.029467in}}{\pgfqpoint{-0.111111in}{0.000000in}}%
\pgfpathcurveto{\pgfqpoint{-0.111111in}{-0.029467in}}{\pgfqpoint{-0.099404in}{-0.057731in}}{\pgfqpoint{-0.078567in}{-0.078567in}}%
\pgfpathcurveto{\pgfqpoint{-0.057731in}{-0.099404in}}{\pgfqpoint{-0.029467in}{-0.111111in}}{\pgfqpoint{0.000000in}{-0.111111in}}%
\pgfpathlineto{\pgfqpoint{0.000000in}{-0.111111in}}%
\pgfpathclose%
\pgfusepath{stroke,fill}%
}%
\begin{pgfscope}%
\pgfsys@transformshift{1.800904in}{3.041134in}%
\pgfsys@useobject{currentmarker}{}%
\end{pgfscope}%
\end{pgfscope}%
\begin{pgfscope}%
\pgfpathrectangle{\pgfqpoint{0.100000in}{2.493063in}}{\pgfqpoint{5.037500in}{3.427208in}}%
\pgfusepath{clip}%
\pgfsetrectcap%
\pgfsetroundjoin%
\pgfsetlinewidth{1.505625pt}%
\definecolor{currentstroke}{rgb}{0.000000,0.000000,1.000000}%
\pgfsetstrokecolor{currentstroke}%
\pgfsetstrokeopacity{0.500000}%
\pgfsetdash{}{0pt}%
\pgfpathmoveto{\pgfqpoint{1.714465in}{3.179522in}}%
\pgfusepath{stroke}%
\end{pgfscope}%
\begin{pgfscope}%
\pgfpathrectangle{\pgfqpoint{0.100000in}{2.493063in}}{\pgfqpoint{5.037500in}{3.427208in}}%
\pgfusepath{clip}%
\pgfsetbuttcap%
\pgfsetroundjoin%
\definecolor{currentfill}{rgb}{0.000000,0.000000,1.000000}%
\pgfsetfillcolor{currentfill}%
\pgfsetfillopacity{0.500000}%
\pgfsetlinewidth{0.250937pt}%
\definecolor{currentstroke}{rgb}{0.000000,0.000000,0.000000}%
\pgfsetstrokecolor{currentstroke}%
\pgfsetstrokeopacity{0.500000}%
\pgfsetdash{}{0pt}%
\pgfsys@defobject{currentmarker}{\pgfqpoint{-0.008333in}{-0.008333in}}{\pgfqpoint{0.008333in}{0.008333in}}{%
\pgfpathmoveto{\pgfqpoint{0.000000in}{-0.008333in}}%
\pgfpathcurveto{\pgfqpoint{0.002210in}{-0.008333in}}{\pgfqpoint{0.004330in}{-0.007455in}}{\pgfqpoint{0.005893in}{-0.005893in}}%
\pgfpathcurveto{\pgfqpoint{0.007455in}{-0.004330in}}{\pgfqpoint{0.008333in}{-0.002210in}}{\pgfqpoint{0.008333in}{0.000000in}}%
\pgfpathcurveto{\pgfqpoint{0.008333in}{0.002210in}}{\pgfqpoint{0.007455in}{0.004330in}}{\pgfqpoint{0.005893in}{0.005893in}}%
\pgfpathcurveto{\pgfqpoint{0.004330in}{0.007455in}}{\pgfqpoint{0.002210in}{0.008333in}}{\pgfqpoint{0.000000in}{0.008333in}}%
\pgfpathcurveto{\pgfqpoint{-0.002210in}{0.008333in}}{\pgfqpoint{-0.004330in}{0.007455in}}{\pgfqpoint{-0.005893in}{0.005893in}}%
\pgfpathcurveto{\pgfqpoint{-0.007455in}{0.004330in}}{\pgfqpoint{-0.008333in}{0.002210in}}{\pgfqpoint{-0.008333in}{0.000000in}}%
\pgfpathcurveto{\pgfqpoint{-0.008333in}{-0.002210in}}{\pgfqpoint{-0.007455in}{-0.004330in}}{\pgfqpoint{-0.005893in}{-0.005893in}}%
\pgfpathcurveto{\pgfqpoint{-0.004330in}{-0.007455in}}{\pgfqpoint{-0.002210in}{-0.008333in}}{\pgfqpoint{0.000000in}{-0.008333in}}%
\pgfpathlineto{\pgfqpoint{0.000000in}{-0.008333in}}%
\pgfpathclose%
\pgfusepath{stroke,fill}%
}%
\begin{pgfscope}%
\pgfsys@transformshift{1.714465in}{3.179522in}%
\pgfsys@useobject{currentmarker}{}%
\end{pgfscope}%
\end{pgfscope}%
\begin{pgfscope}%
\pgfpathrectangle{\pgfqpoint{0.100000in}{2.493063in}}{\pgfqpoint{5.037500in}{3.427208in}}%
\pgfusepath{clip}%
\pgfsetrectcap%
\pgfsetroundjoin%
\pgfsetlinewidth{1.505625pt}%
\definecolor{currentstroke}{rgb}{0.000000,0.000000,1.000000}%
\pgfsetstrokecolor{currentstroke}%
\pgfsetstrokeopacity{0.500000}%
\pgfsetdash{}{0pt}%
\pgfpathmoveto{\pgfqpoint{1.043203in}{5.097798in}}%
\pgfusepath{stroke}%
\end{pgfscope}%
\begin{pgfscope}%
\pgfpathrectangle{\pgfqpoint{0.100000in}{2.493063in}}{\pgfqpoint{5.037500in}{3.427208in}}%
\pgfusepath{clip}%
\pgfsetbuttcap%
\pgfsetroundjoin%
\definecolor{currentfill}{rgb}{0.000000,0.000000,1.000000}%
\pgfsetfillcolor{currentfill}%
\pgfsetfillopacity{0.500000}%
\pgfsetlinewidth{0.250937pt}%
\definecolor{currentstroke}{rgb}{0.000000,0.000000,0.000000}%
\pgfsetstrokecolor{currentstroke}%
\pgfsetstrokeopacity{0.500000}%
\pgfsetdash{}{0pt}%
\pgfsys@defobject{currentmarker}{\pgfqpoint{-0.013889in}{-0.013889in}}{\pgfqpoint{0.013889in}{0.013889in}}{%
\pgfpathmoveto{\pgfqpoint{0.000000in}{-0.013889in}}%
\pgfpathcurveto{\pgfqpoint{0.003683in}{-0.013889in}}{\pgfqpoint{0.007216in}{-0.012425in}}{\pgfqpoint{0.009821in}{-0.009821in}}%
\pgfpathcurveto{\pgfqpoint{0.012425in}{-0.007216in}}{\pgfqpoint{0.013889in}{-0.003683in}}{\pgfqpoint{0.013889in}{0.000000in}}%
\pgfpathcurveto{\pgfqpoint{0.013889in}{0.003683in}}{\pgfqpoint{0.012425in}{0.007216in}}{\pgfqpoint{0.009821in}{0.009821in}}%
\pgfpathcurveto{\pgfqpoint{0.007216in}{0.012425in}}{\pgfqpoint{0.003683in}{0.013889in}}{\pgfqpoint{0.000000in}{0.013889in}}%
\pgfpathcurveto{\pgfqpoint{-0.003683in}{0.013889in}}{\pgfqpoint{-0.007216in}{0.012425in}}{\pgfqpoint{-0.009821in}{0.009821in}}%
\pgfpathcurveto{\pgfqpoint{-0.012425in}{0.007216in}}{\pgfqpoint{-0.013889in}{0.003683in}}{\pgfqpoint{-0.013889in}{0.000000in}}%
\pgfpathcurveto{\pgfqpoint{-0.013889in}{-0.003683in}}{\pgfqpoint{-0.012425in}{-0.007216in}}{\pgfqpoint{-0.009821in}{-0.009821in}}%
\pgfpathcurveto{\pgfqpoint{-0.007216in}{-0.012425in}}{\pgfqpoint{-0.003683in}{-0.013889in}}{\pgfqpoint{0.000000in}{-0.013889in}}%
\pgfpathlineto{\pgfqpoint{0.000000in}{-0.013889in}}%
\pgfpathclose%
\pgfusepath{stroke,fill}%
}%
\begin{pgfscope}%
\pgfsys@transformshift{1.043203in}{5.097798in}%
\pgfsys@useobject{currentmarker}{}%
\end{pgfscope}%
\end{pgfscope}%
\begin{pgfscope}%
\pgfpathrectangle{\pgfqpoint{0.100000in}{2.493063in}}{\pgfqpoint{5.037500in}{3.427208in}}%
\pgfusepath{clip}%
\pgfsetrectcap%
\pgfsetroundjoin%
\pgfsetlinewidth{1.505625pt}%
\definecolor{currentstroke}{rgb}{0.000000,0.000000,1.000000}%
\pgfsetstrokecolor{currentstroke}%
\pgfsetstrokeopacity{0.500000}%
\pgfsetdash{}{0pt}%
\pgfpathmoveto{\pgfqpoint{1.111792in}{5.567791in}}%
\pgfusepath{stroke}%
\end{pgfscope}%
\begin{pgfscope}%
\pgfpathrectangle{\pgfqpoint{0.100000in}{2.493063in}}{\pgfqpoint{5.037500in}{3.427208in}}%
\pgfusepath{clip}%
\pgfsetbuttcap%
\pgfsetroundjoin%
\definecolor{currentfill}{rgb}{0.000000,0.000000,1.000000}%
\pgfsetfillcolor{currentfill}%
\pgfsetfillopacity{0.500000}%
\pgfsetlinewidth{0.250937pt}%
\definecolor{currentstroke}{rgb}{0.000000,0.000000,0.000000}%
\pgfsetstrokecolor{currentstroke}%
\pgfsetstrokeopacity{0.500000}%
\pgfsetdash{}{0pt}%
\pgfsys@defobject{currentmarker}{\pgfqpoint{-0.016667in}{-0.016667in}}{\pgfqpoint{0.016667in}{0.016667in}}{%
\pgfpathmoveto{\pgfqpoint{0.000000in}{-0.016667in}}%
\pgfpathcurveto{\pgfqpoint{0.004420in}{-0.016667in}}{\pgfqpoint{0.008660in}{-0.014911in}}{\pgfqpoint{0.011785in}{-0.011785in}}%
\pgfpathcurveto{\pgfqpoint{0.014911in}{-0.008660in}}{\pgfqpoint{0.016667in}{-0.004420in}}{\pgfqpoint{0.016667in}{0.000000in}}%
\pgfpathcurveto{\pgfqpoint{0.016667in}{0.004420in}}{\pgfqpoint{0.014911in}{0.008660in}}{\pgfqpoint{0.011785in}{0.011785in}}%
\pgfpathcurveto{\pgfqpoint{0.008660in}{0.014911in}}{\pgfqpoint{0.004420in}{0.016667in}}{\pgfqpoint{0.000000in}{0.016667in}}%
\pgfpathcurveto{\pgfqpoint{-0.004420in}{0.016667in}}{\pgfqpoint{-0.008660in}{0.014911in}}{\pgfqpoint{-0.011785in}{0.011785in}}%
\pgfpathcurveto{\pgfqpoint{-0.014911in}{0.008660in}}{\pgfqpoint{-0.016667in}{0.004420in}}{\pgfqpoint{-0.016667in}{0.000000in}}%
\pgfpathcurveto{\pgfqpoint{-0.016667in}{-0.004420in}}{\pgfqpoint{-0.014911in}{-0.008660in}}{\pgfqpoint{-0.011785in}{-0.011785in}}%
\pgfpathcurveto{\pgfqpoint{-0.008660in}{-0.014911in}}{\pgfqpoint{-0.004420in}{-0.016667in}}{\pgfqpoint{0.000000in}{-0.016667in}}%
\pgfpathlineto{\pgfqpoint{0.000000in}{-0.016667in}}%
\pgfpathclose%
\pgfusepath{stroke,fill}%
}%
\begin{pgfscope}%
\pgfsys@transformshift{1.111792in}{5.567791in}%
\pgfsys@useobject{currentmarker}{}%
\end{pgfscope}%
\end{pgfscope}%
\begin{pgfscope}%
\pgfpathrectangle{\pgfqpoint{0.100000in}{2.493063in}}{\pgfqpoint{5.037500in}{3.427208in}}%
\pgfusepath{clip}%
\pgfsetrectcap%
\pgfsetroundjoin%
\pgfsetlinewidth{1.505625pt}%
\definecolor{currentstroke}{rgb}{0.000000,0.000000,1.000000}%
\pgfsetstrokecolor{currentstroke}%
\pgfsetstrokeopacity{0.500000}%
\pgfsetdash{}{0pt}%
\pgfpathmoveto{\pgfqpoint{1.384741in}{5.012340in}}%
\pgfusepath{stroke}%
\end{pgfscope}%
\begin{pgfscope}%
\pgfpathrectangle{\pgfqpoint{0.100000in}{2.493063in}}{\pgfqpoint{5.037500in}{3.427208in}}%
\pgfusepath{clip}%
\pgfsetbuttcap%
\pgfsetroundjoin%
\definecolor{currentfill}{rgb}{0.000000,0.000000,1.000000}%
\pgfsetfillcolor{currentfill}%
\pgfsetfillopacity{0.500000}%
\pgfsetlinewidth{0.250937pt}%
\definecolor{currentstroke}{rgb}{0.000000,0.000000,0.000000}%
\pgfsetstrokecolor{currentstroke}%
\pgfsetstrokeopacity{0.500000}%
\pgfsetdash{}{0pt}%
\pgfsys@defobject{currentmarker}{\pgfqpoint{-0.011111in}{-0.011111in}}{\pgfqpoint{0.011111in}{0.011111in}}{%
\pgfpathmoveto{\pgfqpoint{0.000000in}{-0.011111in}}%
\pgfpathcurveto{\pgfqpoint{0.002947in}{-0.011111in}}{\pgfqpoint{0.005773in}{-0.009940in}}{\pgfqpoint{0.007857in}{-0.007857in}}%
\pgfpathcurveto{\pgfqpoint{0.009940in}{-0.005773in}}{\pgfqpoint{0.011111in}{-0.002947in}}{\pgfqpoint{0.011111in}{0.000000in}}%
\pgfpathcurveto{\pgfqpoint{0.011111in}{0.002947in}}{\pgfqpoint{0.009940in}{0.005773in}}{\pgfqpoint{0.007857in}{0.007857in}}%
\pgfpathcurveto{\pgfqpoint{0.005773in}{0.009940in}}{\pgfqpoint{0.002947in}{0.011111in}}{\pgfqpoint{0.000000in}{0.011111in}}%
\pgfpathcurveto{\pgfqpoint{-0.002947in}{0.011111in}}{\pgfqpoint{-0.005773in}{0.009940in}}{\pgfqpoint{-0.007857in}{0.007857in}}%
\pgfpathcurveto{\pgfqpoint{-0.009940in}{0.005773in}}{\pgfqpoint{-0.011111in}{0.002947in}}{\pgfqpoint{-0.011111in}{0.000000in}}%
\pgfpathcurveto{\pgfqpoint{-0.011111in}{-0.002947in}}{\pgfqpoint{-0.009940in}{-0.005773in}}{\pgfqpoint{-0.007857in}{-0.007857in}}%
\pgfpathcurveto{\pgfqpoint{-0.005773in}{-0.009940in}}{\pgfqpoint{-0.002947in}{-0.011111in}}{\pgfqpoint{0.000000in}{-0.011111in}}%
\pgfpathlineto{\pgfqpoint{0.000000in}{-0.011111in}}%
\pgfpathclose%
\pgfusepath{stroke,fill}%
}%
\begin{pgfscope}%
\pgfsys@transformshift{1.384741in}{5.012340in}%
\pgfsys@useobject{currentmarker}{}%
\end{pgfscope}%
\end{pgfscope}%
\begin{pgfscope}%
\pgfpathrectangle{\pgfqpoint{0.100000in}{2.493063in}}{\pgfqpoint{5.037500in}{3.427208in}}%
\pgfusepath{clip}%
\pgfsetrectcap%
\pgfsetroundjoin%
\pgfsetlinewidth{1.505625pt}%
\definecolor{currentstroke}{rgb}{0.501961,0.501961,0.501961}%
\pgfsetstrokecolor{currentstroke}%
\pgfsetstrokeopacity{0.500000}%
\pgfsetdash{}{0pt}%
\pgfpathmoveto{\pgfqpoint{1.058529in}{5.429824in}}%
\pgfusepath{stroke}%
\end{pgfscope}%
\begin{pgfscope}%
\pgfpathrectangle{\pgfqpoint{0.100000in}{2.493063in}}{\pgfqpoint{5.037500in}{3.427208in}}%
\pgfusepath{clip}%
\pgfsetbuttcap%
\pgfsetroundjoin%
\definecolor{currentfill}{rgb}{0.501961,0.501961,0.501961}%
\pgfsetfillcolor{currentfill}%
\pgfsetfillopacity{0.500000}%
\pgfsetlinewidth{0.250937pt}%
\definecolor{currentstroke}{rgb}{0.000000,0.000000,0.000000}%
\pgfsetstrokecolor{currentstroke}%
\pgfsetstrokeopacity{0.500000}%
\pgfsetdash{}{0pt}%
\pgfsys@defobject{currentmarker}{\pgfqpoint{-0.013889in}{-0.013889in}}{\pgfqpoint{0.013889in}{0.013889in}}{%
\pgfpathmoveto{\pgfqpoint{0.000000in}{-0.013889in}}%
\pgfpathcurveto{\pgfqpoint{0.003683in}{-0.013889in}}{\pgfqpoint{0.007216in}{-0.012425in}}{\pgfqpoint{0.009821in}{-0.009821in}}%
\pgfpathcurveto{\pgfqpoint{0.012425in}{-0.007216in}}{\pgfqpoint{0.013889in}{-0.003683in}}{\pgfqpoint{0.013889in}{0.000000in}}%
\pgfpathcurveto{\pgfqpoint{0.013889in}{0.003683in}}{\pgfqpoint{0.012425in}{0.007216in}}{\pgfqpoint{0.009821in}{0.009821in}}%
\pgfpathcurveto{\pgfqpoint{0.007216in}{0.012425in}}{\pgfqpoint{0.003683in}{0.013889in}}{\pgfqpoint{0.000000in}{0.013889in}}%
\pgfpathcurveto{\pgfqpoint{-0.003683in}{0.013889in}}{\pgfqpoint{-0.007216in}{0.012425in}}{\pgfqpoint{-0.009821in}{0.009821in}}%
\pgfpathcurveto{\pgfqpoint{-0.012425in}{0.007216in}}{\pgfqpoint{-0.013889in}{0.003683in}}{\pgfqpoint{-0.013889in}{0.000000in}}%
\pgfpathcurveto{\pgfqpoint{-0.013889in}{-0.003683in}}{\pgfqpoint{-0.012425in}{-0.007216in}}{\pgfqpoint{-0.009821in}{-0.009821in}}%
\pgfpathcurveto{\pgfqpoint{-0.007216in}{-0.012425in}}{\pgfqpoint{-0.003683in}{-0.013889in}}{\pgfqpoint{0.000000in}{-0.013889in}}%
\pgfpathlineto{\pgfqpoint{0.000000in}{-0.013889in}}%
\pgfpathclose%
\pgfusepath{stroke,fill}%
}%
\begin{pgfscope}%
\pgfsys@transformshift{1.058529in}{5.429824in}%
\pgfsys@useobject{currentmarker}{}%
\end{pgfscope}%
\end{pgfscope}%
\begin{pgfscope}%
\pgfpathrectangle{\pgfqpoint{0.100000in}{2.493063in}}{\pgfqpoint{5.037500in}{3.427208in}}%
\pgfusepath{clip}%
\pgfsetrectcap%
\pgfsetroundjoin%
\pgfsetlinewidth{1.505625pt}%
\definecolor{currentstroke}{rgb}{0.000000,0.000000,1.000000}%
\pgfsetstrokecolor{currentstroke}%
\pgfsetstrokeopacity{0.500000}%
\pgfsetdash{}{0pt}%
\pgfpathmoveto{\pgfqpoint{1.338071in}{4.947867in}}%
\pgfusepath{stroke}%
\end{pgfscope}%
\begin{pgfscope}%
\pgfpathrectangle{\pgfqpoint{0.100000in}{2.493063in}}{\pgfqpoint{5.037500in}{3.427208in}}%
\pgfusepath{clip}%
\pgfsetbuttcap%
\pgfsetroundjoin%
\definecolor{currentfill}{rgb}{0.000000,0.000000,1.000000}%
\pgfsetfillcolor{currentfill}%
\pgfsetfillopacity{0.500000}%
\pgfsetlinewidth{0.250937pt}%
\definecolor{currentstroke}{rgb}{0.000000,0.000000,0.000000}%
\pgfsetstrokecolor{currentstroke}%
\pgfsetstrokeopacity{0.500000}%
\pgfsetdash{}{0pt}%
\pgfsys@defobject{currentmarker}{\pgfqpoint{-0.013889in}{-0.013889in}}{\pgfqpoint{0.013889in}{0.013889in}}{%
\pgfpathmoveto{\pgfqpoint{0.000000in}{-0.013889in}}%
\pgfpathcurveto{\pgfqpoint{0.003683in}{-0.013889in}}{\pgfqpoint{0.007216in}{-0.012425in}}{\pgfqpoint{0.009821in}{-0.009821in}}%
\pgfpathcurveto{\pgfqpoint{0.012425in}{-0.007216in}}{\pgfqpoint{0.013889in}{-0.003683in}}{\pgfqpoint{0.013889in}{0.000000in}}%
\pgfpathcurveto{\pgfqpoint{0.013889in}{0.003683in}}{\pgfqpoint{0.012425in}{0.007216in}}{\pgfqpoint{0.009821in}{0.009821in}}%
\pgfpathcurveto{\pgfqpoint{0.007216in}{0.012425in}}{\pgfqpoint{0.003683in}{0.013889in}}{\pgfqpoint{0.000000in}{0.013889in}}%
\pgfpathcurveto{\pgfqpoint{-0.003683in}{0.013889in}}{\pgfqpoint{-0.007216in}{0.012425in}}{\pgfqpoint{-0.009821in}{0.009821in}}%
\pgfpathcurveto{\pgfqpoint{-0.012425in}{0.007216in}}{\pgfqpoint{-0.013889in}{0.003683in}}{\pgfqpoint{-0.013889in}{0.000000in}}%
\pgfpathcurveto{\pgfqpoint{-0.013889in}{-0.003683in}}{\pgfqpoint{-0.012425in}{-0.007216in}}{\pgfqpoint{-0.009821in}{-0.009821in}}%
\pgfpathcurveto{\pgfqpoint{-0.007216in}{-0.012425in}}{\pgfqpoint{-0.003683in}{-0.013889in}}{\pgfqpoint{0.000000in}{-0.013889in}}%
\pgfpathlineto{\pgfqpoint{0.000000in}{-0.013889in}}%
\pgfpathclose%
\pgfusepath{stroke,fill}%
}%
\begin{pgfscope}%
\pgfsys@transformshift{1.338071in}{4.947867in}%
\pgfsys@useobject{currentmarker}{}%
\end{pgfscope}%
\end{pgfscope}%
\begin{pgfscope}%
\pgfpathrectangle{\pgfqpoint{0.100000in}{2.493063in}}{\pgfqpoint{5.037500in}{3.427208in}}%
\pgfusepath{clip}%
\pgfsetrectcap%
\pgfsetroundjoin%
\pgfsetlinewidth{1.505625pt}%
\definecolor{currentstroke}{rgb}{0.000000,0.000000,1.000000}%
\pgfsetstrokecolor{currentstroke}%
\pgfsetstrokeopacity{0.500000}%
\pgfsetdash{}{0pt}%
\pgfpathmoveto{\pgfqpoint{3.329976in}{4.547971in}}%
\pgfusepath{stroke}%
\end{pgfscope}%
\begin{pgfscope}%
\pgfpathrectangle{\pgfqpoint{0.100000in}{2.493063in}}{\pgfqpoint{5.037500in}{3.427208in}}%
\pgfusepath{clip}%
\pgfsetbuttcap%
\pgfsetroundjoin%
\definecolor{currentfill}{rgb}{0.000000,0.000000,1.000000}%
\pgfsetfillcolor{currentfill}%
\pgfsetfillopacity{0.500000}%
\pgfsetlinewidth{0.250937pt}%
\definecolor{currentstroke}{rgb}{0.000000,0.000000,0.000000}%
\pgfsetstrokecolor{currentstroke}%
\pgfsetstrokeopacity{0.500000}%
\pgfsetdash{}{0pt}%
\pgfsys@defobject{currentmarker}{\pgfqpoint{-0.013889in}{-0.013889in}}{\pgfqpoint{0.013889in}{0.013889in}}{%
\pgfpathmoveto{\pgfqpoint{0.000000in}{-0.013889in}}%
\pgfpathcurveto{\pgfqpoint{0.003683in}{-0.013889in}}{\pgfqpoint{0.007216in}{-0.012425in}}{\pgfqpoint{0.009821in}{-0.009821in}}%
\pgfpathcurveto{\pgfqpoint{0.012425in}{-0.007216in}}{\pgfqpoint{0.013889in}{-0.003683in}}{\pgfqpoint{0.013889in}{0.000000in}}%
\pgfpathcurveto{\pgfqpoint{0.013889in}{0.003683in}}{\pgfqpoint{0.012425in}{0.007216in}}{\pgfqpoint{0.009821in}{0.009821in}}%
\pgfpathcurveto{\pgfqpoint{0.007216in}{0.012425in}}{\pgfqpoint{0.003683in}{0.013889in}}{\pgfqpoint{0.000000in}{0.013889in}}%
\pgfpathcurveto{\pgfqpoint{-0.003683in}{0.013889in}}{\pgfqpoint{-0.007216in}{0.012425in}}{\pgfqpoint{-0.009821in}{0.009821in}}%
\pgfpathcurveto{\pgfqpoint{-0.012425in}{0.007216in}}{\pgfqpoint{-0.013889in}{0.003683in}}{\pgfqpoint{-0.013889in}{0.000000in}}%
\pgfpathcurveto{\pgfqpoint{-0.013889in}{-0.003683in}}{\pgfqpoint{-0.012425in}{-0.007216in}}{\pgfqpoint{-0.009821in}{-0.009821in}}%
\pgfpathcurveto{\pgfqpoint{-0.007216in}{-0.012425in}}{\pgfqpoint{-0.003683in}{-0.013889in}}{\pgfqpoint{0.000000in}{-0.013889in}}%
\pgfpathlineto{\pgfqpoint{0.000000in}{-0.013889in}}%
\pgfpathclose%
\pgfusepath{stroke,fill}%
}%
\begin{pgfscope}%
\pgfsys@transformshift{3.329976in}{4.547971in}%
\pgfsys@useobject{currentmarker}{}%
\end{pgfscope}%
\end{pgfscope}%
\begin{pgfscope}%
\pgfpathrectangle{\pgfqpoint{0.100000in}{2.493063in}}{\pgfqpoint{5.037500in}{3.427208in}}%
\pgfusepath{clip}%
\pgfsetrectcap%
\pgfsetroundjoin%
\pgfsetlinewidth{1.505625pt}%
\definecolor{currentstroke}{rgb}{0.000000,0.000000,1.000000}%
\pgfsetstrokecolor{currentstroke}%
\pgfsetstrokeopacity{0.500000}%
\pgfsetdash{}{0pt}%
\pgfpathmoveto{\pgfqpoint{3.330628in}{4.231165in}}%
\pgfusepath{stroke}%
\end{pgfscope}%
\begin{pgfscope}%
\pgfpathrectangle{\pgfqpoint{0.100000in}{2.493063in}}{\pgfqpoint{5.037500in}{3.427208in}}%
\pgfusepath{clip}%
\pgfsetbuttcap%
\pgfsetroundjoin%
\definecolor{currentfill}{rgb}{0.000000,0.000000,1.000000}%
\pgfsetfillcolor{currentfill}%
\pgfsetfillopacity{0.500000}%
\pgfsetlinewidth{0.250937pt}%
\definecolor{currentstroke}{rgb}{0.000000,0.000000,0.000000}%
\pgfsetstrokecolor{currentstroke}%
\pgfsetstrokeopacity{0.500000}%
\pgfsetdash{}{0pt}%
\pgfsys@defobject{currentmarker}{\pgfqpoint{-0.016667in}{-0.016667in}}{\pgfqpoint{0.016667in}{0.016667in}}{%
\pgfpathmoveto{\pgfqpoint{0.000000in}{-0.016667in}}%
\pgfpathcurveto{\pgfqpoint{0.004420in}{-0.016667in}}{\pgfqpoint{0.008660in}{-0.014911in}}{\pgfqpoint{0.011785in}{-0.011785in}}%
\pgfpathcurveto{\pgfqpoint{0.014911in}{-0.008660in}}{\pgfqpoint{0.016667in}{-0.004420in}}{\pgfqpoint{0.016667in}{0.000000in}}%
\pgfpathcurveto{\pgfqpoint{0.016667in}{0.004420in}}{\pgfqpoint{0.014911in}{0.008660in}}{\pgfqpoint{0.011785in}{0.011785in}}%
\pgfpathcurveto{\pgfqpoint{0.008660in}{0.014911in}}{\pgfqpoint{0.004420in}{0.016667in}}{\pgfqpoint{0.000000in}{0.016667in}}%
\pgfpathcurveto{\pgfqpoint{-0.004420in}{0.016667in}}{\pgfqpoint{-0.008660in}{0.014911in}}{\pgfqpoint{-0.011785in}{0.011785in}}%
\pgfpathcurveto{\pgfqpoint{-0.014911in}{0.008660in}}{\pgfqpoint{-0.016667in}{0.004420in}}{\pgfqpoint{-0.016667in}{0.000000in}}%
\pgfpathcurveto{\pgfqpoint{-0.016667in}{-0.004420in}}{\pgfqpoint{-0.014911in}{-0.008660in}}{\pgfqpoint{-0.011785in}{-0.011785in}}%
\pgfpathcurveto{\pgfqpoint{-0.008660in}{-0.014911in}}{\pgfqpoint{-0.004420in}{-0.016667in}}{\pgfqpoint{0.000000in}{-0.016667in}}%
\pgfpathlineto{\pgfqpoint{0.000000in}{-0.016667in}}%
\pgfpathclose%
\pgfusepath{stroke,fill}%
}%
\begin{pgfscope}%
\pgfsys@transformshift{3.330628in}{4.231165in}%
\pgfsys@useobject{currentmarker}{}%
\end{pgfscope}%
\end{pgfscope}%
\begin{pgfscope}%
\pgfpathrectangle{\pgfqpoint{0.100000in}{2.493063in}}{\pgfqpoint{5.037500in}{3.427208in}}%
\pgfusepath{clip}%
\pgfsetrectcap%
\pgfsetroundjoin%
\pgfsetlinewidth{1.505625pt}%
\definecolor{currentstroke}{rgb}{0.000000,0.000000,1.000000}%
\pgfsetstrokecolor{currentstroke}%
\pgfsetstrokeopacity{0.500000}%
\pgfsetdash{}{0pt}%
\pgfpathmoveto{\pgfqpoint{3.398646in}{4.511619in}}%
\pgfusepath{stroke}%
\end{pgfscope}%
\begin{pgfscope}%
\pgfpathrectangle{\pgfqpoint{0.100000in}{2.493063in}}{\pgfqpoint{5.037500in}{3.427208in}}%
\pgfusepath{clip}%
\pgfsetbuttcap%
\pgfsetroundjoin%
\definecolor{currentfill}{rgb}{0.000000,0.000000,1.000000}%
\pgfsetfillcolor{currentfill}%
\pgfsetfillopacity{0.500000}%
\pgfsetlinewidth{0.250937pt}%
\definecolor{currentstroke}{rgb}{0.000000,0.000000,0.000000}%
\pgfsetstrokecolor{currentstroke}%
\pgfsetstrokeopacity{0.500000}%
\pgfsetdash{}{0pt}%
\pgfsys@defobject{currentmarker}{\pgfqpoint{-0.019444in}{-0.019444in}}{\pgfqpoint{0.019444in}{0.019444in}}{%
\pgfpathmoveto{\pgfqpoint{0.000000in}{-0.019444in}}%
\pgfpathcurveto{\pgfqpoint{0.005157in}{-0.019444in}}{\pgfqpoint{0.010103in}{-0.017396in}}{\pgfqpoint{0.013749in}{-0.013749in}}%
\pgfpathcurveto{\pgfqpoint{0.017396in}{-0.010103in}}{\pgfqpoint{0.019444in}{-0.005157in}}{\pgfqpoint{0.019444in}{0.000000in}}%
\pgfpathcurveto{\pgfqpoint{0.019444in}{0.005157in}}{\pgfqpoint{0.017396in}{0.010103in}}{\pgfqpoint{0.013749in}{0.013749in}}%
\pgfpathcurveto{\pgfqpoint{0.010103in}{0.017396in}}{\pgfqpoint{0.005157in}{0.019444in}}{\pgfqpoint{0.000000in}{0.019444in}}%
\pgfpathcurveto{\pgfqpoint{-0.005157in}{0.019444in}}{\pgfqpoint{-0.010103in}{0.017396in}}{\pgfqpoint{-0.013749in}{0.013749in}}%
\pgfpathcurveto{\pgfqpoint{-0.017396in}{0.010103in}}{\pgfqpoint{-0.019444in}{0.005157in}}{\pgfqpoint{-0.019444in}{0.000000in}}%
\pgfpathcurveto{\pgfqpoint{-0.019444in}{-0.005157in}}{\pgfqpoint{-0.017396in}{-0.010103in}}{\pgfqpoint{-0.013749in}{-0.013749in}}%
\pgfpathcurveto{\pgfqpoint{-0.010103in}{-0.017396in}}{\pgfqpoint{-0.005157in}{-0.019444in}}{\pgfqpoint{0.000000in}{-0.019444in}}%
\pgfpathlineto{\pgfqpoint{0.000000in}{-0.019444in}}%
\pgfpathclose%
\pgfusepath{stroke,fill}%
}%
\begin{pgfscope}%
\pgfsys@transformshift{3.398646in}{4.511619in}%
\pgfsys@useobject{currentmarker}{}%
\end{pgfscope}%
\end{pgfscope}%
\begin{pgfscope}%
\pgfpathrectangle{\pgfqpoint{0.100000in}{2.493063in}}{\pgfqpoint{5.037500in}{3.427208in}}%
\pgfusepath{clip}%
\pgfsetrectcap%
\pgfsetroundjoin%
\pgfsetlinewidth{1.505625pt}%
\definecolor{currentstroke}{rgb}{0.000000,0.000000,1.000000}%
\pgfsetstrokecolor{currentstroke}%
\pgfsetstrokeopacity{0.500000}%
\pgfsetdash{}{0pt}%
\pgfpathmoveto{\pgfqpoint{3.436563in}{4.717868in}}%
\pgfusepath{stroke}%
\end{pgfscope}%
\begin{pgfscope}%
\pgfpathrectangle{\pgfqpoint{0.100000in}{2.493063in}}{\pgfqpoint{5.037500in}{3.427208in}}%
\pgfusepath{clip}%
\pgfsetbuttcap%
\pgfsetroundjoin%
\definecolor{currentfill}{rgb}{0.000000,0.000000,1.000000}%
\pgfsetfillcolor{currentfill}%
\pgfsetfillopacity{0.500000}%
\pgfsetlinewidth{0.250937pt}%
\definecolor{currentstroke}{rgb}{0.000000,0.000000,0.000000}%
\pgfsetstrokecolor{currentstroke}%
\pgfsetstrokeopacity{0.500000}%
\pgfsetdash{}{0pt}%
\pgfsys@defobject{currentmarker}{\pgfqpoint{-0.030556in}{-0.030556in}}{\pgfqpoint{0.030556in}{0.030556in}}{%
\pgfpathmoveto{\pgfqpoint{0.000000in}{-0.030556in}}%
\pgfpathcurveto{\pgfqpoint{0.008103in}{-0.030556in}}{\pgfqpoint{0.015876in}{-0.027336in}}{\pgfqpoint{0.021606in}{-0.021606in}}%
\pgfpathcurveto{\pgfqpoint{0.027336in}{-0.015876in}}{\pgfqpoint{0.030556in}{-0.008103in}}{\pgfqpoint{0.030556in}{0.000000in}}%
\pgfpathcurveto{\pgfqpoint{0.030556in}{0.008103in}}{\pgfqpoint{0.027336in}{0.015876in}}{\pgfqpoint{0.021606in}{0.021606in}}%
\pgfpathcurveto{\pgfqpoint{0.015876in}{0.027336in}}{\pgfqpoint{0.008103in}{0.030556in}}{\pgfqpoint{0.000000in}{0.030556in}}%
\pgfpathcurveto{\pgfqpoint{-0.008103in}{0.030556in}}{\pgfqpoint{-0.015876in}{0.027336in}}{\pgfqpoint{-0.021606in}{0.021606in}}%
\pgfpathcurveto{\pgfqpoint{-0.027336in}{0.015876in}}{\pgfqpoint{-0.030556in}{0.008103in}}{\pgfqpoint{-0.030556in}{0.000000in}}%
\pgfpathcurveto{\pgfqpoint{-0.030556in}{-0.008103in}}{\pgfqpoint{-0.027336in}{-0.015876in}}{\pgfqpoint{-0.021606in}{-0.021606in}}%
\pgfpathcurveto{\pgfqpoint{-0.015876in}{-0.027336in}}{\pgfqpoint{-0.008103in}{-0.030556in}}{\pgfqpoint{0.000000in}{-0.030556in}}%
\pgfpathlineto{\pgfqpoint{0.000000in}{-0.030556in}}%
\pgfpathclose%
\pgfusepath{stroke,fill}%
}%
\begin{pgfscope}%
\pgfsys@transformshift{3.436563in}{4.717868in}%
\pgfsys@useobject{currentmarker}{}%
\end{pgfscope}%
\end{pgfscope}%
\begin{pgfscope}%
\pgfpathrectangle{\pgfqpoint{0.100000in}{2.493063in}}{\pgfqpoint{5.037500in}{3.427208in}}%
\pgfusepath{clip}%
\pgfsetrectcap%
\pgfsetroundjoin%
\pgfsetlinewidth{1.505625pt}%
\definecolor{currentstroke}{rgb}{0.000000,0.000000,1.000000}%
\pgfsetstrokecolor{currentstroke}%
\pgfsetstrokeopacity{0.500000}%
\pgfsetdash{}{0pt}%
\pgfpathmoveto{\pgfqpoint{3.452382in}{4.516821in}}%
\pgfusepath{stroke}%
\end{pgfscope}%
\begin{pgfscope}%
\pgfpathrectangle{\pgfqpoint{0.100000in}{2.493063in}}{\pgfqpoint{5.037500in}{3.427208in}}%
\pgfusepath{clip}%
\pgfsetbuttcap%
\pgfsetroundjoin%
\definecolor{currentfill}{rgb}{0.000000,0.000000,1.000000}%
\pgfsetfillcolor{currentfill}%
\pgfsetfillopacity{0.500000}%
\pgfsetlinewidth{0.250937pt}%
\definecolor{currentstroke}{rgb}{0.000000,0.000000,0.000000}%
\pgfsetstrokecolor{currentstroke}%
\pgfsetstrokeopacity{0.500000}%
\pgfsetdash{}{0pt}%
\pgfsys@defobject{currentmarker}{\pgfqpoint{-0.016667in}{-0.016667in}}{\pgfqpoint{0.016667in}{0.016667in}}{%
\pgfpathmoveto{\pgfqpoint{0.000000in}{-0.016667in}}%
\pgfpathcurveto{\pgfqpoint{0.004420in}{-0.016667in}}{\pgfqpoint{0.008660in}{-0.014911in}}{\pgfqpoint{0.011785in}{-0.011785in}}%
\pgfpathcurveto{\pgfqpoint{0.014911in}{-0.008660in}}{\pgfqpoint{0.016667in}{-0.004420in}}{\pgfqpoint{0.016667in}{0.000000in}}%
\pgfpathcurveto{\pgfqpoint{0.016667in}{0.004420in}}{\pgfqpoint{0.014911in}{0.008660in}}{\pgfqpoint{0.011785in}{0.011785in}}%
\pgfpathcurveto{\pgfqpoint{0.008660in}{0.014911in}}{\pgfqpoint{0.004420in}{0.016667in}}{\pgfqpoint{0.000000in}{0.016667in}}%
\pgfpathcurveto{\pgfqpoint{-0.004420in}{0.016667in}}{\pgfqpoint{-0.008660in}{0.014911in}}{\pgfqpoint{-0.011785in}{0.011785in}}%
\pgfpathcurveto{\pgfqpoint{-0.014911in}{0.008660in}}{\pgfqpoint{-0.016667in}{0.004420in}}{\pgfqpoint{-0.016667in}{0.000000in}}%
\pgfpathcurveto{\pgfqpoint{-0.016667in}{-0.004420in}}{\pgfqpoint{-0.014911in}{-0.008660in}}{\pgfqpoint{-0.011785in}{-0.011785in}}%
\pgfpathcurveto{\pgfqpoint{-0.008660in}{-0.014911in}}{\pgfqpoint{-0.004420in}{-0.016667in}}{\pgfqpoint{0.000000in}{-0.016667in}}%
\pgfpathlineto{\pgfqpoint{0.000000in}{-0.016667in}}%
\pgfpathclose%
\pgfusepath{stroke,fill}%
}%
\begin{pgfscope}%
\pgfsys@transformshift{3.452382in}{4.516821in}%
\pgfsys@useobject{currentmarker}{}%
\end{pgfscope}%
\end{pgfscope}%
\begin{pgfscope}%
\pgfpathrectangle{\pgfqpoint{0.100000in}{2.493063in}}{\pgfqpoint{5.037500in}{3.427208in}}%
\pgfusepath{clip}%
\pgfsetrectcap%
\pgfsetroundjoin%
\pgfsetlinewidth{1.505625pt}%
\definecolor{currentstroke}{rgb}{0.501961,0.501961,0.501961}%
\pgfsetstrokecolor{currentstroke}%
\pgfsetstrokeopacity{0.500000}%
\pgfsetdash{}{0pt}%
\pgfpathmoveto{\pgfqpoint{3.185573in}{4.660908in}}%
\pgfusepath{stroke}%
\end{pgfscope}%
\begin{pgfscope}%
\pgfpathrectangle{\pgfqpoint{0.100000in}{2.493063in}}{\pgfqpoint{5.037500in}{3.427208in}}%
\pgfusepath{clip}%
\pgfsetbuttcap%
\pgfsetroundjoin%
\definecolor{currentfill}{rgb}{0.501961,0.501961,0.501961}%
\pgfsetfillcolor{currentfill}%
\pgfsetfillopacity{0.500000}%
\pgfsetlinewidth{0.250937pt}%
\definecolor{currentstroke}{rgb}{0.000000,0.000000,0.000000}%
\pgfsetstrokecolor{currentstroke}%
\pgfsetstrokeopacity{0.500000}%
\pgfsetdash{}{0pt}%
\pgfsys@defobject{currentmarker}{\pgfqpoint{-0.013889in}{-0.013889in}}{\pgfqpoint{0.013889in}{0.013889in}}{%
\pgfpathmoveto{\pgfqpoint{0.000000in}{-0.013889in}}%
\pgfpathcurveto{\pgfqpoint{0.003683in}{-0.013889in}}{\pgfqpoint{0.007216in}{-0.012425in}}{\pgfqpoint{0.009821in}{-0.009821in}}%
\pgfpathcurveto{\pgfqpoint{0.012425in}{-0.007216in}}{\pgfqpoint{0.013889in}{-0.003683in}}{\pgfqpoint{0.013889in}{0.000000in}}%
\pgfpathcurveto{\pgfqpoint{0.013889in}{0.003683in}}{\pgfqpoint{0.012425in}{0.007216in}}{\pgfqpoint{0.009821in}{0.009821in}}%
\pgfpathcurveto{\pgfqpoint{0.007216in}{0.012425in}}{\pgfqpoint{0.003683in}{0.013889in}}{\pgfqpoint{0.000000in}{0.013889in}}%
\pgfpathcurveto{\pgfqpoint{-0.003683in}{0.013889in}}{\pgfqpoint{-0.007216in}{0.012425in}}{\pgfqpoint{-0.009821in}{0.009821in}}%
\pgfpathcurveto{\pgfqpoint{-0.012425in}{0.007216in}}{\pgfqpoint{-0.013889in}{0.003683in}}{\pgfqpoint{-0.013889in}{0.000000in}}%
\pgfpathcurveto{\pgfqpoint{-0.013889in}{-0.003683in}}{\pgfqpoint{-0.012425in}{-0.007216in}}{\pgfqpoint{-0.009821in}{-0.009821in}}%
\pgfpathcurveto{\pgfqpoint{-0.007216in}{-0.012425in}}{\pgfqpoint{-0.003683in}{-0.013889in}}{\pgfqpoint{0.000000in}{-0.013889in}}%
\pgfpathlineto{\pgfqpoint{0.000000in}{-0.013889in}}%
\pgfpathclose%
\pgfusepath{stroke,fill}%
}%
\begin{pgfscope}%
\pgfsys@transformshift{3.185573in}{4.660908in}%
\pgfsys@useobject{currentmarker}{}%
\end{pgfscope}%
\end{pgfscope}%
\begin{pgfscope}%
\pgfpathrectangle{\pgfqpoint{0.100000in}{2.493063in}}{\pgfqpoint{5.037500in}{3.427208in}}%
\pgfusepath{clip}%
\pgfsetrectcap%
\pgfsetroundjoin%
\pgfsetlinewidth{1.505625pt}%
\definecolor{currentstroke}{rgb}{0.000000,0.000000,1.000000}%
\pgfsetstrokecolor{currentstroke}%
\pgfsetstrokeopacity{0.500000}%
\pgfsetdash{}{0pt}%
\pgfpathmoveto{\pgfqpoint{3.338431in}{4.476084in}}%
\pgfusepath{stroke}%
\end{pgfscope}%
\begin{pgfscope}%
\pgfpathrectangle{\pgfqpoint{0.100000in}{2.493063in}}{\pgfqpoint{5.037500in}{3.427208in}}%
\pgfusepath{clip}%
\pgfsetbuttcap%
\pgfsetroundjoin%
\definecolor{currentfill}{rgb}{0.000000,0.000000,1.000000}%
\pgfsetfillcolor{currentfill}%
\pgfsetfillopacity{0.500000}%
\pgfsetlinewidth{0.250937pt}%
\definecolor{currentstroke}{rgb}{0.000000,0.000000,0.000000}%
\pgfsetstrokecolor{currentstroke}%
\pgfsetstrokeopacity{0.500000}%
\pgfsetdash{}{0pt}%
\pgfsys@defobject{currentmarker}{\pgfqpoint{-0.022222in}{-0.022222in}}{\pgfqpoint{0.022222in}{0.022222in}}{%
\pgfpathmoveto{\pgfqpoint{0.000000in}{-0.022222in}}%
\pgfpathcurveto{\pgfqpoint{0.005893in}{-0.022222in}}{\pgfqpoint{0.011546in}{-0.019881in}}{\pgfqpoint{0.015713in}{-0.015713in}}%
\pgfpathcurveto{\pgfqpoint{0.019881in}{-0.011546in}}{\pgfqpoint{0.022222in}{-0.005893in}}{\pgfqpoint{0.022222in}{0.000000in}}%
\pgfpathcurveto{\pgfqpoint{0.022222in}{0.005893in}}{\pgfqpoint{0.019881in}{0.011546in}}{\pgfqpoint{0.015713in}{0.015713in}}%
\pgfpathcurveto{\pgfqpoint{0.011546in}{0.019881in}}{\pgfqpoint{0.005893in}{0.022222in}}{\pgfqpoint{0.000000in}{0.022222in}}%
\pgfpathcurveto{\pgfqpoint{-0.005893in}{0.022222in}}{\pgfqpoint{-0.011546in}{0.019881in}}{\pgfqpoint{-0.015713in}{0.015713in}}%
\pgfpathcurveto{\pgfqpoint{-0.019881in}{0.011546in}}{\pgfqpoint{-0.022222in}{0.005893in}}{\pgfqpoint{-0.022222in}{0.000000in}}%
\pgfpathcurveto{\pgfqpoint{-0.022222in}{-0.005893in}}{\pgfqpoint{-0.019881in}{-0.011546in}}{\pgfqpoint{-0.015713in}{-0.015713in}}%
\pgfpathcurveto{\pgfqpoint{-0.011546in}{-0.019881in}}{\pgfqpoint{-0.005893in}{-0.022222in}}{\pgfqpoint{0.000000in}{-0.022222in}}%
\pgfpathlineto{\pgfqpoint{0.000000in}{-0.022222in}}%
\pgfpathclose%
\pgfusepath{stroke,fill}%
}%
\begin{pgfscope}%
\pgfsys@transformshift{3.338431in}{4.476084in}%
\pgfsys@useobject{currentmarker}{}%
\end{pgfscope}%
\end{pgfscope}%
\begin{pgfscope}%
\pgfpathrectangle{\pgfqpoint{0.100000in}{2.493063in}}{\pgfqpoint{5.037500in}{3.427208in}}%
\pgfusepath{clip}%
\pgfsetrectcap%
\pgfsetroundjoin%
\pgfsetlinewidth{1.505625pt}%
\definecolor{currentstroke}{rgb}{0.000000,0.000000,1.000000}%
\pgfsetstrokecolor{currentstroke}%
\pgfsetstrokeopacity{0.500000}%
\pgfsetdash{}{0pt}%
\pgfpathmoveto{\pgfqpoint{3.423139in}{4.629440in}}%
\pgfusepath{stroke}%
\end{pgfscope}%
\begin{pgfscope}%
\pgfpathrectangle{\pgfqpoint{0.100000in}{2.493063in}}{\pgfqpoint{5.037500in}{3.427208in}}%
\pgfusepath{clip}%
\pgfsetbuttcap%
\pgfsetroundjoin%
\definecolor{currentfill}{rgb}{0.000000,0.000000,1.000000}%
\pgfsetfillcolor{currentfill}%
\pgfsetfillopacity{0.500000}%
\pgfsetlinewidth{0.250937pt}%
\definecolor{currentstroke}{rgb}{0.000000,0.000000,0.000000}%
\pgfsetstrokecolor{currentstroke}%
\pgfsetstrokeopacity{0.500000}%
\pgfsetdash{}{0pt}%
\pgfsys@defobject{currentmarker}{\pgfqpoint{-0.013889in}{-0.013889in}}{\pgfqpoint{0.013889in}{0.013889in}}{%
\pgfpathmoveto{\pgfqpoint{0.000000in}{-0.013889in}}%
\pgfpathcurveto{\pgfqpoint{0.003683in}{-0.013889in}}{\pgfqpoint{0.007216in}{-0.012425in}}{\pgfqpoint{0.009821in}{-0.009821in}}%
\pgfpathcurveto{\pgfqpoint{0.012425in}{-0.007216in}}{\pgfqpoint{0.013889in}{-0.003683in}}{\pgfqpoint{0.013889in}{0.000000in}}%
\pgfpathcurveto{\pgfqpoint{0.013889in}{0.003683in}}{\pgfqpoint{0.012425in}{0.007216in}}{\pgfqpoint{0.009821in}{0.009821in}}%
\pgfpathcurveto{\pgfqpoint{0.007216in}{0.012425in}}{\pgfqpoint{0.003683in}{0.013889in}}{\pgfqpoint{0.000000in}{0.013889in}}%
\pgfpathcurveto{\pgfqpoint{-0.003683in}{0.013889in}}{\pgfqpoint{-0.007216in}{0.012425in}}{\pgfqpoint{-0.009821in}{0.009821in}}%
\pgfpathcurveto{\pgfqpoint{-0.012425in}{0.007216in}}{\pgfqpoint{-0.013889in}{0.003683in}}{\pgfqpoint{-0.013889in}{0.000000in}}%
\pgfpathcurveto{\pgfqpoint{-0.013889in}{-0.003683in}}{\pgfqpoint{-0.012425in}{-0.007216in}}{\pgfqpoint{-0.009821in}{-0.009821in}}%
\pgfpathcurveto{\pgfqpoint{-0.007216in}{-0.012425in}}{\pgfqpoint{-0.003683in}{-0.013889in}}{\pgfqpoint{0.000000in}{-0.013889in}}%
\pgfpathlineto{\pgfqpoint{0.000000in}{-0.013889in}}%
\pgfpathclose%
\pgfusepath{stroke,fill}%
}%
\begin{pgfscope}%
\pgfsys@transformshift{3.423139in}{4.629440in}%
\pgfsys@useobject{currentmarker}{}%
\end{pgfscope}%
\end{pgfscope}%
\begin{pgfscope}%
\pgfpathrectangle{\pgfqpoint{0.100000in}{2.493063in}}{\pgfqpoint{5.037500in}{3.427208in}}%
\pgfusepath{clip}%
\pgfsetrectcap%
\pgfsetroundjoin%
\pgfsetlinewidth{1.505625pt}%
\definecolor{currentstroke}{rgb}{0.000000,0.000000,1.000000}%
\pgfsetstrokecolor{currentstroke}%
\pgfsetstrokeopacity{0.500000}%
\pgfsetdash{}{0pt}%
\pgfpathmoveto{\pgfqpoint{3.276435in}{4.570082in}}%
\pgfusepath{stroke}%
\end{pgfscope}%
\begin{pgfscope}%
\pgfpathrectangle{\pgfqpoint{0.100000in}{2.493063in}}{\pgfqpoint{5.037500in}{3.427208in}}%
\pgfusepath{clip}%
\pgfsetbuttcap%
\pgfsetroundjoin%
\definecolor{currentfill}{rgb}{0.000000,0.000000,1.000000}%
\pgfsetfillcolor{currentfill}%
\pgfsetfillopacity{0.500000}%
\pgfsetlinewidth{0.250937pt}%
\definecolor{currentstroke}{rgb}{0.000000,0.000000,0.000000}%
\pgfsetstrokecolor{currentstroke}%
\pgfsetstrokeopacity{0.500000}%
\pgfsetdash{}{0pt}%
\pgfsys@defobject{currentmarker}{\pgfqpoint{-0.011111in}{-0.011111in}}{\pgfqpoint{0.011111in}{0.011111in}}{%
\pgfpathmoveto{\pgfqpoint{0.000000in}{-0.011111in}}%
\pgfpathcurveto{\pgfqpoint{0.002947in}{-0.011111in}}{\pgfqpoint{0.005773in}{-0.009940in}}{\pgfqpoint{0.007857in}{-0.007857in}}%
\pgfpathcurveto{\pgfqpoint{0.009940in}{-0.005773in}}{\pgfqpoint{0.011111in}{-0.002947in}}{\pgfqpoint{0.011111in}{0.000000in}}%
\pgfpathcurveto{\pgfqpoint{0.011111in}{0.002947in}}{\pgfqpoint{0.009940in}{0.005773in}}{\pgfqpoint{0.007857in}{0.007857in}}%
\pgfpathcurveto{\pgfqpoint{0.005773in}{0.009940in}}{\pgfqpoint{0.002947in}{0.011111in}}{\pgfqpoint{0.000000in}{0.011111in}}%
\pgfpathcurveto{\pgfqpoint{-0.002947in}{0.011111in}}{\pgfqpoint{-0.005773in}{0.009940in}}{\pgfqpoint{-0.007857in}{0.007857in}}%
\pgfpathcurveto{\pgfqpoint{-0.009940in}{0.005773in}}{\pgfqpoint{-0.011111in}{0.002947in}}{\pgfqpoint{-0.011111in}{0.000000in}}%
\pgfpathcurveto{\pgfqpoint{-0.011111in}{-0.002947in}}{\pgfqpoint{-0.009940in}{-0.005773in}}{\pgfqpoint{-0.007857in}{-0.007857in}}%
\pgfpathcurveto{\pgfqpoint{-0.005773in}{-0.009940in}}{\pgfqpoint{-0.002947in}{-0.011111in}}{\pgfqpoint{0.000000in}{-0.011111in}}%
\pgfpathlineto{\pgfqpoint{0.000000in}{-0.011111in}}%
\pgfpathclose%
\pgfusepath{stroke,fill}%
}%
\begin{pgfscope}%
\pgfsys@transformshift{3.276435in}{4.570082in}%
\pgfsys@useobject{currentmarker}{}%
\end{pgfscope}%
\end{pgfscope}%
\begin{pgfscope}%
\pgfpathrectangle{\pgfqpoint{0.100000in}{2.493063in}}{\pgfqpoint{5.037500in}{3.427208in}}%
\pgfusepath{clip}%
\pgfsetrectcap%
\pgfsetroundjoin%
\pgfsetlinewidth{1.505625pt}%
\definecolor{currentstroke}{rgb}{0.000000,0.000000,1.000000}%
\pgfsetstrokecolor{currentstroke}%
\pgfsetstrokeopacity{0.500000}%
\pgfsetdash{}{0pt}%
\pgfpathmoveto{\pgfqpoint{3.307789in}{4.754204in}}%
\pgfusepath{stroke}%
\end{pgfscope}%
\begin{pgfscope}%
\pgfpathrectangle{\pgfqpoint{0.100000in}{2.493063in}}{\pgfqpoint{5.037500in}{3.427208in}}%
\pgfusepath{clip}%
\pgfsetbuttcap%
\pgfsetroundjoin%
\definecolor{currentfill}{rgb}{0.000000,0.000000,1.000000}%
\pgfsetfillcolor{currentfill}%
\pgfsetfillopacity{0.500000}%
\pgfsetlinewidth{0.250937pt}%
\definecolor{currentstroke}{rgb}{0.000000,0.000000,0.000000}%
\pgfsetstrokecolor{currentstroke}%
\pgfsetstrokeopacity{0.500000}%
\pgfsetdash{}{0pt}%
\pgfsys@defobject{currentmarker}{\pgfqpoint{-0.013889in}{-0.013889in}}{\pgfqpoint{0.013889in}{0.013889in}}{%
\pgfpathmoveto{\pgfqpoint{0.000000in}{-0.013889in}}%
\pgfpathcurveto{\pgfqpoint{0.003683in}{-0.013889in}}{\pgfqpoint{0.007216in}{-0.012425in}}{\pgfqpoint{0.009821in}{-0.009821in}}%
\pgfpathcurveto{\pgfqpoint{0.012425in}{-0.007216in}}{\pgfqpoint{0.013889in}{-0.003683in}}{\pgfqpoint{0.013889in}{0.000000in}}%
\pgfpathcurveto{\pgfqpoint{0.013889in}{0.003683in}}{\pgfqpoint{0.012425in}{0.007216in}}{\pgfqpoint{0.009821in}{0.009821in}}%
\pgfpathcurveto{\pgfqpoint{0.007216in}{0.012425in}}{\pgfqpoint{0.003683in}{0.013889in}}{\pgfqpoint{0.000000in}{0.013889in}}%
\pgfpathcurveto{\pgfqpoint{-0.003683in}{0.013889in}}{\pgfqpoint{-0.007216in}{0.012425in}}{\pgfqpoint{-0.009821in}{0.009821in}}%
\pgfpathcurveto{\pgfqpoint{-0.012425in}{0.007216in}}{\pgfqpoint{-0.013889in}{0.003683in}}{\pgfqpoint{-0.013889in}{0.000000in}}%
\pgfpathcurveto{\pgfqpoint{-0.013889in}{-0.003683in}}{\pgfqpoint{-0.012425in}{-0.007216in}}{\pgfqpoint{-0.009821in}{-0.009821in}}%
\pgfpathcurveto{\pgfqpoint{-0.007216in}{-0.012425in}}{\pgfqpoint{-0.003683in}{-0.013889in}}{\pgfqpoint{0.000000in}{-0.013889in}}%
\pgfpathlineto{\pgfqpoint{0.000000in}{-0.013889in}}%
\pgfpathclose%
\pgfusepath{stroke,fill}%
}%
\begin{pgfscope}%
\pgfsys@transformshift{3.307789in}{4.754204in}%
\pgfsys@useobject{currentmarker}{}%
\end{pgfscope}%
\end{pgfscope}%
\begin{pgfscope}%
\pgfpathrectangle{\pgfqpoint{0.100000in}{2.493063in}}{\pgfqpoint{5.037500in}{3.427208in}}%
\pgfusepath{clip}%
\pgfsetrectcap%
\pgfsetroundjoin%
\pgfsetlinewidth{1.505625pt}%
\definecolor{currentstroke}{rgb}{0.000000,0.000000,1.000000}%
\pgfsetstrokecolor{currentstroke}%
\pgfsetstrokeopacity{0.500000}%
\pgfsetdash{}{0pt}%
\pgfpathmoveto{\pgfqpoint{3.277687in}{4.466909in}}%
\pgfusepath{stroke}%
\end{pgfscope}%
\begin{pgfscope}%
\pgfpathrectangle{\pgfqpoint{0.100000in}{2.493063in}}{\pgfqpoint{5.037500in}{3.427208in}}%
\pgfusepath{clip}%
\pgfsetbuttcap%
\pgfsetroundjoin%
\definecolor{currentfill}{rgb}{0.000000,0.000000,1.000000}%
\pgfsetfillcolor{currentfill}%
\pgfsetfillopacity{0.500000}%
\pgfsetlinewidth{0.250937pt}%
\definecolor{currentstroke}{rgb}{0.000000,0.000000,0.000000}%
\pgfsetstrokecolor{currentstroke}%
\pgfsetstrokeopacity{0.500000}%
\pgfsetdash{}{0pt}%
\pgfsys@defobject{currentmarker}{\pgfqpoint{-0.011111in}{-0.011111in}}{\pgfqpoint{0.011111in}{0.011111in}}{%
\pgfpathmoveto{\pgfqpoint{0.000000in}{-0.011111in}}%
\pgfpathcurveto{\pgfqpoint{0.002947in}{-0.011111in}}{\pgfqpoint{0.005773in}{-0.009940in}}{\pgfqpoint{0.007857in}{-0.007857in}}%
\pgfpathcurveto{\pgfqpoint{0.009940in}{-0.005773in}}{\pgfqpoint{0.011111in}{-0.002947in}}{\pgfqpoint{0.011111in}{0.000000in}}%
\pgfpathcurveto{\pgfqpoint{0.011111in}{0.002947in}}{\pgfqpoint{0.009940in}{0.005773in}}{\pgfqpoint{0.007857in}{0.007857in}}%
\pgfpathcurveto{\pgfqpoint{0.005773in}{0.009940in}}{\pgfqpoint{0.002947in}{0.011111in}}{\pgfqpoint{0.000000in}{0.011111in}}%
\pgfpathcurveto{\pgfqpoint{-0.002947in}{0.011111in}}{\pgfqpoint{-0.005773in}{0.009940in}}{\pgfqpoint{-0.007857in}{0.007857in}}%
\pgfpathcurveto{\pgfqpoint{-0.009940in}{0.005773in}}{\pgfqpoint{-0.011111in}{0.002947in}}{\pgfqpoint{-0.011111in}{0.000000in}}%
\pgfpathcurveto{\pgfqpoint{-0.011111in}{-0.002947in}}{\pgfqpoint{-0.009940in}{-0.005773in}}{\pgfqpoint{-0.007857in}{-0.007857in}}%
\pgfpathcurveto{\pgfqpoint{-0.005773in}{-0.009940in}}{\pgfqpoint{-0.002947in}{-0.011111in}}{\pgfqpoint{0.000000in}{-0.011111in}}%
\pgfpathlineto{\pgfqpoint{0.000000in}{-0.011111in}}%
\pgfpathclose%
\pgfusepath{stroke,fill}%
}%
\begin{pgfscope}%
\pgfsys@transformshift{3.277687in}{4.466909in}%
\pgfsys@useobject{currentmarker}{}%
\end{pgfscope}%
\end{pgfscope}%
\begin{pgfscope}%
\pgfpathrectangle{\pgfqpoint{0.100000in}{2.493063in}}{\pgfqpoint{5.037500in}{3.427208in}}%
\pgfusepath{clip}%
\pgfsetrectcap%
\pgfsetroundjoin%
\pgfsetlinewidth{1.505625pt}%
\definecolor{currentstroke}{rgb}{0.000000,0.000000,1.000000}%
\pgfsetstrokecolor{currentstroke}%
\pgfsetstrokeopacity{0.500000}%
\pgfsetdash{}{0pt}%
\pgfpathmoveto{\pgfqpoint{3.558808in}{4.415360in}}%
\pgfusepath{stroke}%
\end{pgfscope}%
\begin{pgfscope}%
\pgfpathrectangle{\pgfqpoint{0.100000in}{2.493063in}}{\pgfqpoint{5.037500in}{3.427208in}}%
\pgfusepath{clip}%
\pgfsetbuttcap%
\pgfsetroundjoin%
\definecolor{currentfill}{rgb}{0.000000,0.000000,1.000000}%
\pgfsetfillcolor{currentfill}%
\pgfsetfillopacity{0.500000}%
\pgfsetlinewidth{0.250937pt}%
\definecolor{currentstroke}{rgb}{0.000000,0.000000,0.000000}%
\pgfsetstrokecolor{currentstroke}%
\pgfsetstrokeopacity{0.500000}%
\pgfsetdash{}{0pt}%
\pgfsys@defobject{currentmarker}{\pgfqpoint{-0.011111in}{-0.011111in}}{\pgfqpoint{0.011111in}{0.011111in}}{%
\pgfpathmoveto{\pgfqpoint{0.000000in}{-0.011111in}}%
\pgfpathcurveto{\pgfqpoint{0.002947in}{-0.011111in}}{\pgfqpoint{0.005773in}{-0.009940in}}{\pgfqpoint{0.007857in}{-0.007857in}}%
\pgfpathcurveto{\pgfqpoint{0.009940in}{-0.005773in}}{\pgfqpoint{0.011111in}{-0.002947in}}{\pgfqpoint{0.011111in}{0.000000in}}%
\pgfpathcurveto{\pgfqpoint{0.011111in}{0.002947in}}{\pgfqpoint{0.009940in}{0.005773in}}{\pgfqpoint{0.007857in}{0.007857in}}%
\pgfpathcurveto{\pgfqpoint{0.005773in}{0.009940in}}{\pgfqpoint{0.002947in}{0.011111in}}{\pgfqpoint{0.000000in}{0.011111in}}%
\pgfpathcurveto{\pgfqpoint{-0.002947in}{0.011111in}}{\pgfqpoint{-0.005773in}{0.009940in}}{\pgfqpoint{-0.007857in}{0.007857in}}%
\pgfpathcurveto{\pgfqpoint{-0.009940in}{0.005773in}}{\pgfqpoint{-0.011111in}{0.002947in}}{\pgfqpoint{-0.011111in}{0.000000in}}%
\pgfpathcurveto{\pgfqpoint{-0.011111in}{-0.002947in}}{\pgfqpoint{-0.009940in}{-0.005773in}}{\pgfqpoint{-0.007857in}{-0.007857in}}%
\pgfpathcurveto{\pgfqpoint{-0.005773in}{-0.009940in}}{\pgfqpoint{-0.002947in}{-0.011111in}}{\pgfqpoint{0.000000in}{-0.011111in}}%
\pgfpathlineto{\pgfqpoint{0.000000in}{-0.011111in}}%
\pgfpathclose%
\pgfusepath{stroke,fill}%
}%
\begin{pgfscope}%
\pgfsys@transformshift{3.558808in}{4.415360in}%
\pgfsys@useobject{currentmarker}{}%
\end{pgfscope}%
\end{pgfscope}%
\begin{pgfscope}%
\pgfpathrectangle{\pgfqpoint{0.100000in}{2.493063in}}{\pgfqpoint{5.037500in}{3.427208in}}%
\pgfusepath{clip}%
\pgfsetrectcap%
\pgfsetroundjoin%
\pgfsetlinewidth{1.505625pt}%
\definecolor{currentstroke}{rgb}{0.000000,0.000000,1.000000}%
\pgfsetstrokecolor{currentstroke}%
\pgfsetstrokeopacity{0.500000}%
\pgfsetdash{}{0pt}%
\pgfpathmoveto{\pgfqpoint{3.612876in}{4.424575in}}%
\pgfusepath{stroke}%
\end{pgfscope}%
\begin{pgfscope}%
\pgfpathrectangle{\pgfqpoint{0.100000in}{2.493063in}}{\pgfqpoint{5.037500in}{3.427208in}}%
\pgfusepath{clip}%
\pgfsetbuttcap%
\pgfsetroundjoin%
\definecolor{currentfill}{rgb}{0.000000,0.000000,1.000000}%
\pgfsetfillcolor{currentfill}%
\pgfsetfillopacity{0.500000}%
\pgfsetlinewidth{0.250937pt}%
\definecolor{currentstroke}{rgb}{0.000000,0.000000,0.000000}%
\pgfsetstrokecolor{currentstroke}%
\pgfsetstrokeopacity{0.500000}%
\pgfsetdash{}{0pt}%
\pgfsys@defobject{currentmarker}{\pgfqpoint{-0.019444in}{-0.019444in}}{\pgfqpoint{0.019444in}{0.019444in}}{%
\pgfpathmoveto{\pgfqpoint{0.000000in}{-0.019444in}}%
\pgfpathcurveto{\pgfqpoint{0.005157in}{-0.019444in}}{\pgfqpoint{0.010103in}{-0.017396in}}{\pgfqpoint{0.013749in}{-0.013749in}}%
\pgfpathcurveto{\pgfqpoint{0.017396in}{-0.010103in}}{\pgfqpoint{0.019444in}{-0.005157in}}{\pgfqpoint{0.019444in}{0.000000in}}%
\pgfpathcurveto{\pgfqpoint{0.019444in}{0.005157in}}{\pgfqpoint{0.017396in}{0.010103in}}{\pgfqpoint{0.013749in}{0.013749in}}%
\pgfpathcurveto{\pgfqpoint{0.010103in}{0.017396in}}{\pgfqpoint{0.005157in}{0.019444in}}{\pgfqpoint{0.000000in}{0.019444in}}%
\pgfpathcurveto{\pgfqpoint{-0.005157in}{0.019444in}}{\pgfqpoint{-0.010103in}{0.017396in}}{\pgfqpoint{-0.013749in}{0.013749in}}%
\pgfpathcurveto{\pgfqpoint{-0.017396in}{0.010103in}}{\pgfqpoint{-0.019444in}{0.005157in}}{\pgfqpoint{-0.019444in}{0.000000in}}%
\pgfpathcurveto{\pgfqpoint{-0.019444in}{-0.005157in}}{\pgfqpoint{-0.017396in}{-0.010103in}}{\pgfqpoint{-0.013749in}{-0.013749in}}%
\pgfpathcurveto{\pgfqpoint{-0.010103in}{-0.017396in}}{\pgfqpoint{-0.005157in}{-0.019444in}}{\pgfqpoint{0.000000in}{-0.019444in}}%
\pgfpathlineto{\pgfqpoint{0.000000in}{-0.019444in}}%
\pgfpathclose%
\pgfusepath{stroke,fill}%
}%
\begin{pgfscope}%
\pgfsys@transformshift{3.612876in}{4.424575in}%
\pgfsys@useobject{currentmarker}{}%
\end{pgfscope}%
\end{pgfscope}%
\begin{pgfscope}%
\pgfpathrectangle{\pgfqpoint{0.100000in}{2.493063in}}{\pgfqpoint{5.037500in}{3.427208in}}%
\pgfusepath{clip}%
\pgfsetrectcap%
\pgfsetroundjoin%
\pgfsetlinewidth{1.505625pt}%
\definecolor{currentstroke}{rgb}{0.000000,0.000000,1.000000}%
\pgfsetstrokecolor{currentstroke}%
\pgfsetstrokeopacity{0.500000}%
\pgfsetdash{}{0pt}%
\pgfpathmoveto{\pgfqpoint{3.579642in}{4.708409in}}%
\pgfusepath{stroke}%
\end{pgfscope}%
\begin{pgfscope}%
\pgfpathrectangle{\pgfqpoint{0.100000in}{2.493063in}}{\pgfqpoint{5.037500in}{3.427208in}}%
\pgfusepath{clip}%
\pgfsetbuttcap%
\pgfsetroundjoin%
\definecolor{currentfill}{rgb}{0.000000,0.000000,1.000000}%
\pgfsetfillcolor{currentfill}%
\pgfsetfillopacity{0.500000}%
\pgfsetlinewidth{0.250937pt}%
\definecolor{currentstroke}{rgb}{0.000000,0.000000,0.000000}%
\pgfsetstrokecolor{currentstroke}%
\pgfsetstrokeopacity{0.500000}%
\pgfsetdash{}{0pt}%
\pgfsys@defobject{currentmarker}{\pgfqpoint{-0.019444in}{-0.019444in}}{\pgfqpoint{0.019444in}{0.019444in}}{%
\pgfpathmoveto{\pgfqpoint{0.000000in}{-0.019444in}}%
\pgfpathcurveto{\pgfqpoint{0.005157in}{-0.019444in}}{\pgfqpoint{0.010103in}{-0.017396in}}{\pgfqpoint{0.013749in}{-0.013749in}}%
\pgfpathcurveto{\pgfqpoint{0.017396in}{-0.010103in}}{\pgfqpoint{0.019444in}{-0.005157in}}{\pgfqpoint{0.019444in}{0.000000in}}%
\pgfpathcurveto{\pgfqpoint{0.019444in}{0.005157in}}{\pgfqpoint{0.017396in}{0.010103in}}{\pgfqpoint{0.013749in}{0.013749in}}%
\pgfpathcurveto{\pgfqpoint{0.010103in}{0.017396in}}{\pgfqpoint{0.005157in}{0.019444in}}{\pgfqpoint{0.000000in}{0.019444in}}%
\pgfpathcurveto{\pgfqpoint{-0.005157in}{0.019444in}}{\pgfqpoint{-0.010103in}{0.017396in}}{\pgfqpoint{-0.013749in}{0.013749in}}%
\pgfpathcurveto{\pgfqpoint{-0.017396in}{0.010103in}}{\pgfqpoint{-0.019444in}{0.005157in}}{\pgfqpoint{-0.019444in}{0.000000in}}%
\pgfpathcurveto{\pgfqpoint{-0.019444in}{-0.005157in}}{\pgfqpoint{-0.017396in}{-0.010103in}}{\pgfqpoint{-0.013749in}{-0.013749in}}%
\pgfpathcurveto{\pgfqpoint{-0.010103in}{-0.017396in}}{\pgfqpoint{-0.005157in}{-0.019444in}}{\pgfqpoint{0.000000in}{-0.019444in}}%
\pgfpathlineto{\pgfqpoint{0.000000in}{-0.019444in}}%
\pgfpathclose%
\pgfusepath{stroke,fill}%
}%
\begin{pgfscope}%
\pgfsys@transformshift{3.579642in}{4.708409in}%
\pgfsys@useobject{currentmarker}{}%
\end{pgfscope}%
\end{pgfscope}%
\begin{pgfscope}%
\pgfpathrectangle{\pgfqpoint{0.100000in}{2.493063in}}{\pgfqpoint{5.037500in}{3.427208in}}%
\pgfusepath{clip}%
\pgfsetrectcap%
\pgfsetroundjoin%
\pgfsetlinewidth{1.505625pt}%
\definecolor{currentstroke}{rgb}{0.000000,0.000000,1.000000}%
\pgfsetstrokecolor{currentstroke}%
\pgfsetstrokeopacity{0.500000}%
\pgfsetdash{}{0pt}%
\pgfpathmoveto{\pgfqpoint{3.479193in}{4.270553in}}%
\pgfusepath{stroke}%
\end{pgfscope}%
\begin{pgfscope}%
\pgfpathrectangle{\pgfqpoint{0.100000in}{2.493063in}}{\pgfqpoint{5.037500in}{3.427208in}}%
\pgfusepath{clip}%
\pgfsetbuttcap%
\pgfsetroundjoin%
\definecolor{currentfill}{rgb}{0.000000,0.000000,1.000000}%
\pgfsetfillcolor{currentfill}%
\pgfsetfillopacity{0.500000}%
\pgfsetlinewidth{0.250937pt}%
\definecolor{currentstroke}{rgb}{0.000000,0.000000,0.000000}%
\pgfsetstrokecolor{currentstroke}%
\pgfsetstrokeopacity{0.500000}%
\pgfsetdash{}{0pt}%
\pgfsys@defobject{currentmarker}{\pgfqpoint{-0.008333in}{-0.008333in}}{\pgfqpoint{0.008333in}{0.008333in}}{%
\pgfpathmoveto{\pgfqpoint{0.000000in}{-0.008333in}}%
\pgfpathcurveto{\pgfqpoint{0.002210in}{-0.008333in}}{\pgfqpoint{0.004330in}{-0.007455in}}{\pgfqpoint{0.005893in}{-0.005893in}}%
\pgfpathcurveto{\pgfqpoint{0.007455in}{-0.004330in}}{\pgfqpoint{0.008333in}{-0.002210in}}{\pgfqpoint{0.008333in}{0.000000in}}%
\pgfpathcurveto{\pgfqpoint{0.008333in}{0.002210in}}{\pgfqpoint{0.007455in}{0.004330in}}{\pgfqpoint{0.005893in}{0.005893in}}%
\pgfpathcurveto{\pgfqpoint{0.004330in}{0.007455in}}{\pgfqpoint{0.002210in}{0.008333in}}{\pgfqpoint{0.000000in}{0.008333in}}%
\pgfpathcurveto{\pgfqpoint{-0.002210in}{0.008333in}}{\pgfqpoint{-0.004330in}{0.007455in}}{\pgfqpoint{-0.005893in}{0.005893in}}%
\pgfpathcurveto{\pgfqpoint{-0.007455in}{0.004330in}}{\pgfqpoint{-0.008333in}{0.002210in}}{\pgfqpoint{-0.008333in}{0.000000in}}%
\pgfpathcurveto{\pgfqpoint{-0.008333in}{-0.002210in}}{\pgfqpoint{-0.007455in}{-0.004330in}}{\pgfqpoint{-0.005893in}{-0.005893in}}%
\pgfpathcurveto{\pgfqpoint{-0.004330in}{-0.007455in}}{\pgfqpoint{-0.002210in}{-0.008333in}}{\pgfqpoint{0.000000in}{-0.008333in}}%
\pgfpathlineto{\pgfqpoint{0.000000in}{-0.008333in}}%
\pgfpathclose%
\pgfusepath{stroke,fill}%
}%
\begin{pgfscope}%
\pgfsys@transformshift{3.479193in}{4.270553in}%
\pgfsys@useobject{currentmarker}{}%
\end{pgfscope}%
\end{pgfscope}%
\begin{pgfscope}%
\pgfpathrectangle{\pgfqpoint{0.100000in}{2.493063in}}{\pgfqpoint{5.037500in}{3.427208in}}%
\pgfusepath{clip}%
\pgfsetrectcap%
\pgfsetroundjoin%
\pgfsetlinewidth{1.505625pt}%
\definecolor{currentstroke}{rgb}{0.501961,0.501961,0.501961}%
\pgfsetstrokecolor{currentstroke}%
\pgfsetstrokeopacity{0.500000}%
\pgfsetdash{}{0pt}%
\pgfpathmoveto{\pgfqpoint{3.658913in}{4.646924in}}%
\pgfusepath{stroke}%
\end{pgfscope}%
\begin{pgfscope}%
\pgfpathrectangle{\pgfqpoint{0.100000in}{2.493063in}}{\pgfqpoint{5.037500in}{3.427208in}}%
\pgfusepath{clip}%
\pgfsetbuttcap%
\pgfsetroundjoin%
\definecolor{currentfill}{rgb}{0.501961,0.501961,0.501961}%
\pgfsetfillcolor{currentfill}%
\pgfsetfillopacity{0.500000}%
\pgfsetlinewidth{0.250937pt}%
\definecolor{currentstroke}{rgb}{0.000000,0.000000,0.000000}%
\pgfsetstrokecolor{currentstroke}%
\pgfsetstrokeopacity{0.500000}%
\pgfsetdash{}{0pt}%
\pgfsys@defobject{currentmarker}{\pgfqpoint{-0.013889in}{-0.013889in}}{\pgfqpoint{0.013889in}{0.013889in}}{%
\pgfpathmoveto{\pgfqpoint{0.000000in}{-0.013889in}}%
\pgfpathcurveto{\pgfqpoint{0.003683in}{-0.013889in}}{\pgfqpoint{0.007216in}{-0.012425in}}{\pgfqpoint{0.009821in}{-0.009821in}}%
\pgfpathcurveto{\pgfqpoint{0.012425in}{-0.007216in}}{\pgfqpoint{0.013889in}{-0.003683in}}{\pgfqpoint{0.013889in}{0.000000in}}%
\pgfpathcurveto{\pgfqpoint{0.013889in}{0.003683in}}{\pgfqpoint{0.012425in}{0.007216in}}{\pgfqpoint{0.009821in}{0.009821in}}%
\pgfpathcurveto{\pgfqpoint{0.007216in}{0.012425in}}{\pgfqpoint{0.003683in}{0.013889in}}{\pgfqpoint{0.000000in}{0.013889in}}%
\pgfpathcurveto{\pgfqpoint{-0.003683in}{0.013889in}}{\pgfqpoint{-0.007216in}{0.012425in}}{\pgfqpoint{-0.009821in}{0.009821in}}%
\pgfpathcurveto{\pgfqpoint{-0.012425in}{0.007216in}}{\pgfqpoint{-0.013889in}{0.003683in}}{\pgfqpoint{-0.013889in}{0.000000in}}%
\pgfpathcurveto{\pgfqpoint{-0.013889in}{-0.003683in}}{\pgfqpoint{-0.012425in}{-0.007216in}}{\pgfqpoint{-0.009821in}{-0.009821in}}%
\pgfpathcurveto{\pgfqpoint{-0.007216in}{-0.012425in}}{\pgfqpoint{-0.003683in}{-0.013889in}}{\pgfqpoint{0.000000in}{-0.013889in}}%
\pgfpathlineto{\pgfqpoint{0.000000in}{-0.013889in}}%
\pgfpathclose%
\pgfusepath{stroke,fill}%
}%
\begin{pgfscope}%
\pgfsys@transformshift{3.658913in}{4.646924in}%
\pgfsys@useobject{currentmarker}{}%
\end{pgfscope}%
\end{pgfscope}%
\begin{pgfscope}%
\pgfpathrectangle{\pgfqpoint{0.100000in}{2.493063in}}{\pgfqpoint{5.037500in}{3.427208in}}%
\pgfusepath{clip}%
\pgfsetrectcap%
\pgfsetroundjoin%
\pgfsetlinewidth{1.505625pt}%
\definecolor{currentstroke}{rgb}{0.000000,0.000000,1.000000}%
\pgfsetstrokecolor{currentstroke}%
\pgfsetstrokeopacity{0.500000}%
\pgfsetdash{}{0pt}%
\pgfpathmoveto{\pgfqpoint{3.623626in}{4.530288in}}%
\pgfusepath{stroke}%
\end{pgfscope}%
\begin{pgfscope}%
\pgfpathrectangle{\pgfqpoint{0.100000in}{2.493063in}}{\pgfqpoint{5.037500in}{3.427208in}}%
\pgfusepath{clip}%
\pgfsetbuttcap%
\pgfsetroundjoin%
\definecolor{currentfill}{rgb}{0.000000,0.000000,1.000000}%
\pgfsetfillcolor{currentfill}%
\pgfsetfillopacity{0.500000}%
\pgfsetlinewidth{0.250937pt}%
\definecolor{currentstroke}{rgb}{0.000000,0.000000,0.000000}%
\pgfsetstrokecolor{currentstroke}%
\pgfsetstrokeopacity{0.500000}%
\pgfsetdash{}{0pt}%
\pgfsys@defobject{currentmarker}{\pgfqpoint{-0.008333in}{-0.008333in}}{\pgfqpoint{0.008333in}{0.008333in}}{%
\pgfpathmoveto{\pgfqpoint{0.000000in}{-0.008333in}}%
\pgfpathcurveto{\pgfqpoint{0.002210in}{-0.008333in}}{\pgfqpoint{0.004330in}{-0.007455in}}{\pgfqpoint{0.005893in}{-0.005893in}}%
\pgfpathcurveto{\pgfqpoint{0.007455in}{-0.004330in}}{\pgfqpoint{0.008333in}{-0.002210in}}{\pgfqpoint{0.008333in}{0.000000in}}%
\pgfpathcurveto{\pgfqpoint{0.008333in}{0.002210in}}{\pgfqpoint{0.007455in}{0.004330in}}{\pgfqpoint{0.005893in}{0.005893in}}%
\pgfpathcurveto{\pgfqpoint{0.004330in}{0.007455in}}{\pgfqpoint{0.002210in}{0.008333in}}{\pgfqpoint{0.000000in}{0.008333in}}%
\pgfpathcurveto{\pgfqpoint{-0.002210in}{0.008333in}}{\pgfqpoint{-0.004330in}{0.007455in}}{\pgfqpoint{-0.005893in}{0.005893in}}%
\pgfpathcurveto{\pgfqpoint{-0.007455in}{0.004330in}}{\pgfqpoint{-0.008333in}{0.002210in}}{\pgfqpoint{-0.008333in}{0.000000in}}%
\pgfpathcurveto{\pgfqpoint{-0.008333in}{-0.002210in}}{\pgfqpoint{-0.007455in}{-0.004330in}}{\pgfqpoint{-0.005893in}{-0.005893in}}%
\pgfpathcurveto{\pgfqpoint{-0.004330in}{-0.007455in}}{\pgfqpoint{-0.002210in}{-0.008333in}}{\pgfqpoint{0.000000in}{-0.008333in}}%
\pgfpathlineto{\pgfqpoint{0.000000in}{-0.008333in}}%
\pgfpathclose%
\pgfusepath{stroke,fill}%
}%
\begin{pgfscope}%
\pgfsys@transformshift{3.623626in}{4.530288in}%
\pgfsys@useobject{currentmarker}{}%
\end{pgfscope}%
\end{pgfscope}%
\begin{pgfscope}%
\pgfpathrectangle{\pgfqpoint{0.100000in}{2.493063in}}{\pgfqpoint{5.037500in}{3.427208in}}%
\pgfusepath{clip}%
\pgfsetrectcap%
\pgfsetroundjoin%
\pgfsetlinewidth{1.505625pt}%
\definecolor{currentstroke}{rgb}{0.000000,0.000000,1.000000}%
\pgfsetstrokecolor{currentstroke}%
\pgfsetstrokeopacity{0.500000}%
\pgfsetdash{}{0pt}%
\pgfpathmoveto{\pgfqpoint{3.585492in}{4.487463in}}%
\pgfusepath{stroke}%
\end{pgfscope}%
\begin{pgfscope}%
\pgfpathrectangle{\pgfqpoint{0.100000in}{2.493063in}}{\pgfqpoint{5.037500in}{3.427208in}}%
\pgfusepath{clip}%
\pgfsetbuttcap%
\pgfsetroundjoin%
\definecolor{currentfill}{rgb}{0.000000,0.000000,1.000000}%
\pgfsetfillcolor{currentfill}%
\pgfsetfillopacity{0.500000}%
\pgfsetlinewidth{0.250937pt}%
\definecolor{currentstroke}{rgb}{0.000000,0.000000,0.000000}%
\pgfsetstrokecolor{currentstroke}%
\pgfsetstrokeopacity{0.500000}%
\pgfsetdash{}{0pt}%
\pgfsys@defobject{currentmarker}{\pgfqpoint{-0.008333in}{-0.008333in}}{\pgfqpoint{0.008333in}{0.008333in}}{%
\pgfpathmoveto{\pgfqpoint{0.000000in}{-0.008333in}}%
\pgfpathcurveto{\pgfqpoint{0.002210in}{-0.008333in}}{\pgfqpoint{0.004330in}{-0.007455in}}{\pgfqpoint{0.005893in}{-0.005893in}}%
\pgfpathcurveto{\pgfqpoint{0.007455in}{-0.004330in}}{\pgfqpoint{0.008333in}{-0.002210in}}{\pgfqpoint{0.008333in}{0.000000in}}%
\pgfpathcurveto{\pgfqpoint{0.008333in}{0.002210in}}{\pgfqpoint{0.007455in}{0.004330in}}{\pgfqpoint{0.005893in}{0.005893in}}%
\pgfpathcurveto{\pgfqpoint{0.004330in}{0.007455in}}{\pgfqpoint{0.002210in}{0.008333in}}{\pgfqpoint{0.000000in}{0.008333in}}%
\pgfpathcurveto{\pgfqpoint{-0.002210in}{0.008333in}}{\pgfqpoint{-0.004330in}{0.007455in}}{\pgfqpoint{-0.005893in}{0.005893in}}%
\pgfpathcurveto{\pgfqpoint{-0.007455in}{0.004330in}}{\pgfqpoint{-0.008333in}{0.002210in}}{\pgfqpoint{-0.008333in}{0.000000in}}%
\pgfpathcurveto{\pgfqpoint{-0.008333in}{-0.002210in}}{\pgfqpoint{-0.007455in}{-0.004330in}}{\pgfqpoint{-0.005893in}{-0.005893in}}%
\pgfpathcurveto{\pgfqpoint{-0.004330in}{-0.007455in}}{\pgfqpoint{-0.002210in}{-0.008333in}}{\pgfqpoint{0.000000in}{-0.008333in}}%
\pgfpathlineto{\pgfqpoint{0.000000in}{-0.008333in}}%
\pgfpathclose%
\pgfusepath{stroke,fill}%
}%
\begin{pgfscope}%
\pgfsys@transformshift{3.585492in}{4.487463in}%
\pgfsys@useobject{currentmarker}{}%
\end{pgfscope}%
\end{pgfscope}%
\begin{pgfscope}%
\pgfpathrectangle{\pgfqpoint{0.100000in}{2.493063in}}{\pgfqpoint{5.037500in}{3.427208in}}%
\pgfusepath{clip}%
\pgfsetrectcap%
\pgfsetroundjoin%
\pgfsetlinewidth{1.505625pt}%
\definecolor{currentstroke}{rgb}{0.000000,0.000000,1.000000}%
\pgfsetstrokecolor{currentstroke}%
\pgfsetstrokeopacity{0.500000}%
\pgfsetdash{}{0pt}%
\pgfpathmoveto{\pgfqpoint{3.579613in}{4.569985in}}%
\pgfusepath{stroke}%
\end{pgfscope}%
\begin{pgfscope}%
\pgfpathrectangle{\pgfqpoint{0.100000in}{2.493063in}}{\pgfqpoint{5.037500in}{3.427208in}}%
\pgfusepath{clip}%
\pgfsetbuttcap%
\pgfsetroundjoin%
\definecolor{currentfill}{rgb}{0.000000,0.000000,1.000000}%
\pgfsetfillcolor{currentfill}%
\pgfsetfillopacity{0.500000}%
\pgfsetlinewidth{0.250937pt}%
\definecolor{currentstroke}{rgb}{0.000000,0.000000,0.000000}%
\pgfsetstrokecolor{currentstroke}%
\pgfsetstrokeopacity{0.500000}%
\pgfsetdash{}{0pt}%
\pgfsys@defobject{currentmarker}{\pgfqpoint{-0.086111in}{-0.086111in}}{\pgfqpoint{0.086111in}{0.086111in}}{%
\pgfpathmoveto{\pgfqpoint{0.000000in}{-0.086111in}}%
\pgfpathcurveto{\pgfqpoint{0.022837in}{-0.086111in}}{\pgfqpoint{0.044742in}{-0.077038in}}{\pgfqpoint{0.060890in}{-0.060890in}}%
\pgfpathcurveto{\pgfqpoint{0.077038in}{-0.044742in}}{\pgfqpoint{0.086111in}{-0.022837in}}{\pgfqpoint{0.086111in}{0.000000in}}%
\pgfpathcurveto{\pgfqpoint{0.086111in}{0.022837in}}{\pgfqpoint{0.077038in}{0.044742in}}{\pgfqpoint{0.060890in}{0.060890in}}%
\pgfpathcurveto{\pgfqpoint{0.044742in}{0.077038in}}{\pgfqpoint{0.022837in}{0.086111in}}{\pgfqpoint{0.000000in}{0.086111in}}%
\pgfpathcurveto{\pgfqpoint{-0.022837in}{0.086111in}}{\pgfqpoint{-0.044742in}{0.077038in}}{\pgfqpoint{-0.060890in}{0.060890in}}%
\pgfpathcurveto{\pgfqpoint{-0.077038in}{0.044742in}}{\pgfqpoint{-0.086111in}{0.022837in}}{\pgfqpoint{-0.086111in}{0.000000in}}%
\pgfpathcurveto{\pgfqpoint{-0.086111in}{-0.022837in}}{\pgfqpoint{-0.077038in}{-0.044742in}}{\pgfqpoint{-0.060890in}{-0.060890in}}%
\pgfpathcurveto{\pgfqpoint{-0.044742in}{-0.077038in}}{\pgfqpoint{-0.022837in}{-0.086111in}}{\pgfqpoint{0.000000in}{-0.086111in}}%
\pgfpathlineto{\pgfqpoint{0.000000in}{-0.086111in}}%
\pgfpathclose%
\pgfusepath{stroke,fill}%
}%
\begin{pgfscope}%
\pgfsys@transformshift{3.579613in}{4.569985in}%
\pgfsys@useobject{currentmarker}{}%
\end{pgfscope}%
\end{pgfscope}%
\begin{pgfscope}%
\pgfpathrectangle{\pgfqpoint{0.100000in}{2.493063in}}{\pgfqpoint{5.037500in}{3.427208in}}%
\pgfusepath{clip}%
\pgfsetrectcap%
\pgfsetroundjoin%
\pgfsetlinewidth{1.505625pt}%
\definecolor{currentstroke}{rgb}{0.000000,0.000000,1.000000}%
\pgfsetstrokecolor{currentstroke}%
\pgfsetstrokeopacity{0.500000}%
\pgfsetdash{}{0pt}%
\pgfpathmoveto{\pgfqpoint{3.515644in}{4.555920in}}%
\pgfusepath{stroke}%
\end{pgfscope}%
\begin{pgfscope}%
\pgfpathrectangle{\pgfqpoint{0.100000in}{2.493063in}}{\pgfqpoint{5.037500in}{3.427208in}}%
\pgfusepath{clip}%
\pgfsetbuttcap%
\pgfsetroundjoin%
\definecolor{currentfill}{rgb}{0.000000,0.000000,1.000000}%
\pgfsetfillcolor{currentfill}%
\pgfsetfillopacity{0.500000}%
\pgfsetlinewidth{0.250937pt}%
\definecolor{currentstroke}{rgb}{0.000000,0.000000,0.000000}%
\pgfsetstrokecolor{currentstroke}%
\pgfsetstrokeopacity{0.500000}%
\pgfsetdash{}{0pt}%
\pgfsys@defobject{currentmarker}{\pgfqpoint{-0.008333in}{-0.008333in}}{\pgfqpoint{0.008333in}{0.008333in}}{%
\pgfpathmoveto{\pgfqpoint{0.000000in}{-0.008333in}}%
\pgfpathcurveto{\pgfqpoint{0.002210in}{-0.008333in}}{\pgfqpoint{0.004330in}{-0.007455in}}{\pgfqpoint{0.005893in}{-0.005893in}}%
\pgfpathcurveto{\pgfqpoint{0.007455in}{-0.004330in}}{\pgfqpoint{0.008333in}{-0.002210in}}{\pgfqpoint{0.008333in}{0.000000in}}%
\pgfpathcurveto{\pgfqpoint{0.008333in}{0.002210in}}{\pgfqpoint{0.007455in}{0.004330in}}{\pgfqpoint{0.005893in}{0.005893in}}%
\pgfpathcurveto{\pgfqpoint{0.004330in}{0.007455in}}{\pgfqpoint{0.002210in}{0.008333in}}{\pgfqpoint{0.000000in}{0.008333in}}%
\pgfpathcurveto{\pgfqpoint{-0.002210in}{0.008333in}}{\pgfqpoint{-0.004330in}{0.007455in}}{\pgfqpoint{-0.005893in}{0.005893in}}%
\pgfpathcurveto{\pgfqpoint{-0.007455in}{0.004330in}}{\pgfqpoint{-0.008333in}{0.002210in}}{\pgfqpoint{-0.008333in}{0.000000in}}%
\pgfpathcurveto{\pgfqpoint{-0.008333in}{-0.002210in}}{\pgfqpoint{-0.007455in}{-0.004330in}}{\pgfqpoint{-0.005893in}{-0.005893in}}%
\pgfpathcurveto{\pgfqpoint{-0.004330in}{-0.007455in}}{\pgfqpoint{-0.002210in}{-0.008333in}}{\pgfqpoint{0.000000in}{-0.008333in}}%
\pgfpathlineto{\pgfqpoint{0.000000in}{-0.008333in}}%
\pgfpathclose%
\pgfusepath{stroke,fill}%
}%
\begin{pgfscope}%
\pgfsys@transformshift{3.515644in}{4.555920in}%
\pgfsys@useobject{currentmarker}{}%
\end{pgfscope}%
\end{pgfscope}%
\begin{pgfscope}%
\pgfpathrectangle{\pgfqpoint{0.100000in}{2.493063in}}{\pgfqpoint{5.037500in}{3.427208in}}%
\pgfusepath{clip}%
\pgfsetrectcap%
\pgfsetroundjoin%
\pgfsetlinewidth{1.505625pt}%
\definecolor{currentstroke}{rgb}{0.000000,0.000000,1.000000}%
\pgfsetstrokecolor{currentstroke}%
\pgfsetstrokeopacity{0.500000}%
\pgfsetdash{}{0pt}%
\pgfpathmoveto{\pgfqpoint{3.500672in}{4.703903in}}%
\pgfusepath{stroke}%
\end{pgfscope}%
\begin{pgfscope}%
\pgfpathrectangle{\pgfqpoint{0.100000in}{2.493063in}}{\pgfqpoint{5.037500in}{3.427208in}}%
\pgfusepath{clip}%
\pgfsetbuttcap%
\pgfsetroundjoin%
\definecolor{currentfill}{rgb}{0.000000,0.000000,1.000000}%
\pgfsetfillcolor{currentfill}%
\pgfsetfillopacity{0.500000}%
\pgfsetlinewidth{0.250937pt}%
\definecolor{currentstroke}{rgb}{0.000000,0.000000,0.000000}%
\pgfsetstrokecolor{currentstroke}%
\pgfsetstrokeopacity{0.500000}%
\pgfsetdash{}{0pt}%
\pgfsys@defobject{currentmarker}{\pgfqpoint{-0.005556in}{-0.005556in}}{\pgfqpoint{0.005556in}{0.005556in}}{%
\pgfpathmoveto{\pgfqpoint{0.000000in}{-0.005556in}}%
\pgfpathcurveto{\pgfqpoint{0.001473in}{-0.005556in}}{\pgfqpoint{0.002887in}{-0.004970in}}{\pgfqpoint{0.003928in}{-0.003928in}}%
\pgfpathcurveto{\pgfqpoint{0.004970in}{-0.002887in}}{\pgfqpoint{0.005556in}{-0.001473in}}{\pgfqpoint{0.005556in}{0.000000in}}%
\pgfpathcurveto{\pgfqpoint{0.005556in}{0.001473in}}{\pgfqpoint{0.004970in}{0.002887in}}{\pgfqpoint{0.003928in}{0.003928in}}%
\pgfpathcurveto{\pgfqpoint{0.002887in}{0.004970in}}{\pgfqpoint{0.001473in}{0.005556in}}{\pgfqpoint{0.000000in}{0.005556in}}%
\pgfpathcurveto{\pgfqpoint{-0.001473in}{0.005556in}}{\pgfqpoint{-0.002887in}{0.004970in}}{\pgfqpoint{-0.003928in}{0.003928in}}%
\pgfpathcurveto{\pgfqpoint{-0.004970in}{0.002887in}}{\pgfqpoint{-0.005556in}{0.001473in}}{\pgfqpoint{-0.005556in}{0.000000in}}%
\pgfpathcurveto{\pgfqpoint{-0.005556in}{-0.001473in}}{\pgfqpoint{-0.004970in}{-0.002887in}}{\pgfqpoint{-0.003928in}{-0.003928in}}%
\pgfpathcurveto{\pgfqpoint{-0.002887in}{-0.004970in}}{\pgfqpoint{-0.001473in}{-0.005556in}}{\pgfqpoint{0.000000in}{-0.005556in}}%
\pgfpathlineto{\pgfqpoint{0.000000in}{-0.005556in}}%
\pgfpathclose%
\pgfusepath{stroke,fill}%
}%
\begin{pgfscope}%
\pgfsys@transformshift{3.500672in}{4.703903in}%
\pgfsys@useobject{currentmarker}{}%
\end{pgfscope}%
\end{pgfscope}%
\begin{pgfscope}%
\pgfpathrectangle{\pgfqpoint{0.100000in}{2.493063in}}{\pgfqpoint{5.037500in}{3.427208in}}%
\pgfusepath{clip}%
\pgfsetrectcap%
\pgfsetroundjoin%
\pgfsetlinewidth{1.505625pt}%
\definecolor{currentstroke}{rgb}{0.000000,0.000000,1.000000}%
\pgfsetstrokecolor{currentstroke}%
\pgfsetstrokeopacity{0.500000}%
\pgfsetdash{}{0pt}%
\pgfpathmoveto{\pgfqpoint{3.648286in}{4.543096in}}%
\pgfusepath{stroke}%
\end{pgfscope}%
\begin{pgfscope}%
\pgfpathrectangle{\pgfqpoint{0.100000in}{2.493063in}}{\pgfqpoint{5.037500in}{3.427208in}}%
\pgfusepath{clip}%
\pgfsetbuttcap%
\pgfsetroundjoin%
\definecolor{currentfill}{rgb}{0.000000,0.000000,1.000000}%
\pgfsetfillcolor{currentfill}%
\pgfsetfillopacity{0.500000}%
\pgfsetlinewidth{0.250937pt}%
\definecolor{currentstroke}{rgb}{0.000000,0.000000,0.000000}%
\pgfsetstrokecolor{currentstroke}%
\pgfsetstrokeopacity{0.500000}%
\pgfsetdash{}{0pt}%
\pgfsys@defobject{currentmarker}{\pgfqpoint{-0.008333in}{-0.008333in}}{\pgfqpoint{0.008333in}{0.008333in}}{%
\pgfpathmoveto{\pgfqpoint{0.000000in}{-0.008333in}}%
\pgfpathcurveto{\pgfqpoint{0.002210in}{-0.008333in}}{\pgfqpoint{0.004330in}{-0.007455in}}{\pgfqpoint{0.005893in}{-0.005893in}}%
\pgfpathcurveto{\pgfqpoint{0.007455in}{-0.004330in}}{\pgfqpoint{0.008333in}{-0.002210in}}{\pgfqpoint{0.008333in}{0.000000in}}%
\pgfpathcurveto{\pgfqpoint{0.008333in}{0.002210in}}{\pgfqpoint{0.007455in}{0.004330in}}{\pgfqpoint{0.005893in}{0.005893in}}%
\pgfpathcurveto{\pgfqpoint{0.004330in}{0.007455in}}{\pgfqpoint{0.002210in}{0.008333in}}{\pgfqpoint{0.000000in}{0.008333in}}%
\pgfpathcurveto{\pgfqpoint{-0.002210in}{0.008333in}}{\pgfqpoint{-0.004330in}{0.007455in}}{\pgfqpoint{-0.005893in}{0.005893in}}%
\pgfpathcurveto{\pgfqpoint{-0.007455in}{0.004330in}}{\pgfqpoint{-0.008333in}{0.002210in}}{\pgfqpoint{-0.008333in}{0.000000in}}%
\pgfpathcurveto{\pgfqpoint{-0.008333in}{-0.002210in}}{\pgfqpoint{-0.007455in}{-0.004330in}}{\pgfqpoint{-0.005893in}{-0.005893in}}%
\pgfpathcurveto{\pgfqpoint{-0.004330in}{-0.007455in}}{\pgfqpoint{-0.002210in}{-0.008333in}}{\pgfqpoint{0.000000in}{-0.008333in}}%
\pgfpathlineto{\pgfqpoint{0.000000in}{-0.008333in}}%
\pgfpathclose%
\pgfusepath{stroke,fill}%
}%
\begin{pgfscope}%
\pgfsys@transformshift{3.648286in}{4.543096in}%
\pgfsys@useobject{currentmarker}{}%
\end{pgfscope}%
\end{pgfscope}%
\begin{pgfscope}%
\pgfpathrectangle{\pgfqpoint{0.100000in}{2.493063in}}{\pgfqpoint{5.037500in}{3.427208in}}%
\pgfusepath{clip}%
\pgfsetrectcap%
\pgfsetroundjoin%
\pgfsetlinewidth{1.505625pt}%
\definecolor{currentstroke}{rgb}{0.000000,0.000000,1.000000}%
\pgfsetstrokecolor{currentstroke}%
\pgfsetstrokeopacity{0.500000}%
\pgfsetdash{}{0pt}%
\pgfpathmoveto{\pgfqpoint{3.556209in}{4.706271in}}%
\pgfusepath{stroke}%
\end{pgfscope}%
\begin{pgfscope}%
\pgfpathrectangle{\pgfqpoint{0.100000in}{2.493063in}}{\pgfqpoint{5.037500in}{3.427208in}}%
\pgfusepath{clip}%
\pgfsetbuttcap%
\pgfsetroundjoin%
\definecolor{currentfill}{rgb}{0.000000,0.000000,1.000000}%
\pgfsetfillcolor{currentfill}%
\pgfsetfillopacity{0.500000}%
\pgfsetlinewidth{0.250937pt}%
\definecolor{currentstroke}{rgb}{0.000000,0.000000,0.000000}%
\pgfsetstrokecolor{currentstroke}%
\pgfsetstrokeopacity{0.500000}%
\pgfsetdash{}{0pt}%
\pgfsys@defobject{currentmarker}{\pgfqpoint{-0.013889in}{-0.013889in}}{\pgfqpoint{0.013889in}{0.013889in}}{%
\pgfpathmoveto{\pgfqpoint{0.000000in}{-0.013889in}}%
\pgfpathcurveto{\pgfqpoint{0.003683in}{-0.013889in}}{\pgfqpoint{0.007216in}{-0.012425in}}{\pgfqpoint{0.009821in}{-0.009821in}}%
\pgfpathcurveto{\pgfqpoint{0.012425in}{-0.007216in}}{\pgfqpoint{0.013889in}{-0.003683in}}{\pgfqpoint{0.013889in}{0.000000in}}%
\pgfpathcurveto{\pgfqpoint{0.013889in}{0.003683in}}{\pgfqpoint{0.012425in}{0.007216in}}{\pgfqpoint{0.009821in}{0.009821in}}%
\pgfpathcurveto{\pgfqpoint{0.007216in}{0.012425in}}{\pgfqpoint{0.003683in}{0.013889in}}{\pgfqpoint{0.000000in}{0.013889in}}%
\pgfpathcurveto{\pgfqpoint{-0.003683in}{0.013889in}}{\pgfqpoint{-0.007216in}{0.012425in}}{\pgfqpoint{-0.009821in}{0.009821in}}%
\pgfpathcurveto{\pgfqpoint{-0.012425in}{0.007216in}}{\pgfqpoint{-0.013889in}{0.003683in}}{\pgfqpoint{-0.013889in}{0.000000in}}%
\pgfpathcurveto{\pgfqpoint{-0.013889in}{-0.003683in}}{\pgfqpoint{-0.012425in}{-0.007216in}}{\pgfqpoint{-0.009821in}{-0.009821in}}%
\pgfpathcurveto{\pgfqpoint{-0.007216in}{-0.012425in}}{\pgfqpoint{-0.003683in}{-0.013889in}}{\pgfqpoint{0.000000in}{-0.013889in}}%
\pgfpathlineto{\pgfqpoint{0.000000in}{-0.013889in}}%
\pgfpathclose%
\pgfusepath{stroke,fill}%
}%
\begin{pgfscope}%
\pgfsys@transformshift{3.556209in}{4.706271in}%
\pgfsys@useobject{currentmarker}{}%
\end{pgfscope}%
\end{pgfscope}%
\begin{pgfscope}%
\pgfpathrectangle{\pgfqpoint{0.100000in}{2.493063in}}{\pgfqpoint{5.037500in}{3.427208in}}%
\pgfusepath{clip}%
\pgfsetrectcap%
\pgfsetroundjoin%
\pgfsetlinewidth{1.505625pt}%
\definecolor{currentstroke}{rgb}{0.678431,1.000000,0.184314}%
\pgfsetstrokecolor{currentstroke}%
\pgfsetstrokeopacity{0.500000}%
\pgfsetdash{}{0pt}%
\pgfpathmoveto{\pgfqpoint{3.477710in}{4.442820in}}%
\pgfusepath{stroke}%
\end{pgfscope}%
\begin{pgfscope}%
\pgfpathrectangle{\pgfqpoint{0.100000in}{2.493063in}}{\pgfqpoint{5.037500in}{3.427208in}}%
\pgfusepath{clip}%
\pgfsetbuttcap%
\pgfsetroundjoin%
\definecolor{currentfill}{rgb}{0.678431,1.000000,0.184314}%
\pgfsetfillcolor{currentfill}%
\pgfsetfillopacity{0.500000}%
\pgfsetlinewidth{0.250937pt}%
\definecolor{currentstroke}{rgb}{0.000000,0.000000,0.000000}%
\pgfsetstrokecolor{currentstroke}%
\pgfsetstrokeopacity{0.500000}%
\pgfsetdash{}{0pt}%
\pgfsys@defobject{currentmarker}{\pgfqpoint{-0.005556in}{-0.005556in}}{\pgfqpoint{0.005556in}{0.005556in}}{%
\pgfpathmoveto{\pgfqpoint{0.000000in}{-0.005556in}}%
\pgfpathcurveto{\pgfqpoint{0.001473in}{-0.005556in}}{\pgfqpoint{0.002887in}{-0.004970in}}{\pgfqpoint{0.003928in}{-0.003928in}}%
\pgfpathcurveto{\pgfqpoint{0.004970in}{-0.002887in}}{\pgfqpoint{0.005556in}{-0.001473in}}{\pgfqpoint{0.005556in}{0.000000in}}%
\pgfpathcurveto{\pgfqpoint{0.005556in}{0.001473in}}{\pgfqpoint{0.004970in}{0.002887in}}{\pgfqpoint{0.003928in}{0.003928in}}%
\pgfpathcurveto{\pgfqpoint{0.002887in}{0.004970in}}{\pgfqpoint{0.001473in}{0.005556in}}{\pgfqpoint{0.000000in}{0.005556in}}%
\pgfpathcurveto{\pgfqpoint{-0.001473in}{0.005556in}}{\pgfqpoint{-0.002887in}{0.004970in}}{\pgfqpoint{-0.003928in}{0.003928in}}%
\pgfpathcurveto{\pgfqpoint{-0.004970in}{0.002887in}}{\pgfqpoint{-0.005556in}{0.001473in}}{\pgfqpoint{-0.005556in}{0.000000in}}%
\pgfpathcurveto{\pgfqpoint{-0.005556in}{-0.001473in}}{\pgfqpoint{-0.004970in}{-0.002887in}}{\pgfqpoint{-0.003928in}{-0.003928in}}%
\pgfpathcurveto{\pgfqpoint{-0.002887in}{-0.004970in}}{\pgfqpoint{-0.001473in}{-0.005556in}}{\pgfqpoint{0.000000in}{-0.005556in}}%
\pgfpathlineto{\pgfqpoint{0.000000in}{-0.005556in}}%
\pgfpathclose%
\pgfusepath{stroke,fill}%
}%
\begin{pgfscope}%
\pgfsys@transformshift{3.477710in}{4.442820in}%
\pgfsys@useobject{currentmarker}{}%
\end{pgfscope}%
\end{pgfscope}%
\begin{pgfscope}%
\pgfpathrectangle{\pgfqpoint{0.100000in}{2.493063in}}{\pgfqpoint{5.037500in}{3.427208in}}%
\pgfusepath{clip}%
\pgfsetrectcap%
\pgfsetroundjoin%
\pgfsetlinewidth{1.505625pt}%
\definecolor{currentstroke}{rgb}{0.000000,0.000000,1.000000}%
\pgfsetstrokecolor{currentstroke}%
\pgfsetstrokeopacity{0.500000}%
\pgfsetdash{}{0pt}%
\pgfpathmoveto{\pgfqpoint{2.922572in}{4.710525in}}%
\pgfusepath{stroke}%
\end{pgfscope}%
\begin{pgfscope}%
\pgfpathrectangle{\pgfqpoint{0.100000in}{2.493063in}}{\pgfqpoint{5.037500in}{3.427208in}}%
\pgfusepath{clip}%
\pgfsetbuttcap%
\pgfsetroundjoin%
\definecolor{currentfill}{rgb}{0.000000,0.000000,1.000000}%
\pgfsetfillcolor{currentfill}%
\pgfsetfillopacity{0.500000}%
\pgfsetlinewidth{0.250937pt}%
\definecolor{currentstroke}{rgb}{0.000000,0.000000,0.000000}%
\pgfsetstrokecolor{currentstroke}%
\pgfsetstrokeopacity{0.500000}%
\pgfsetdash{}{0pt}%
\pgfsys@defobject{currentmarker}{\pgfqpoint{-0.016667in}{-0.016667in}}{\pgfqpoint{0.016667in}{0.016667in}}{%
\pgfpathmoveto{\pgfqpoint{0.000000in}{-0.016667in}}%
\pgfpathcurveto{\pgfqpoint{0.004420in}{-0.016667in}}{\pgfqpoint{0.008660in}{-0.014911in}}{\pgfqpoint{0.011785in}{-0.011785in}}%
\pgfpathcurveto{\pgfqpoint{0.014911in}{-0.008660in}}{\pgfqpoint{0.016667in}{-0.004420in}}{\pgfqpoint{0.016667in}{0.000000in}}%
\pgfpathcurveto{\pgfqpoint{0.016667in}{0.004420in}}{\pgfqpoint{0.014911in}{0.008660in}}{\pgfqpoint{0.011785in}{0.011785in}}%
\pgfpathcurveto{\pgfqpoint{0.008660in}{0.014911in}}{\pgfqpoint{0.004420in}{0.016667in}}{\pgfqpoint{0.000000in}{0.016667in}}%
\pgfpathcurveto{\pgfqpoint{-0.004420in}{0.016667in}}{\pgfqpoint{-0.008660in}{0.014911in}}{\pgfqpoint{-0.011785in}{0.011785in}}%
\pgfpathcurveto{\pgfqpoint{-0.014911in}{0.008660in}}{\pgfqpoint{-0.016667in}{0.004420in}}{\pgfqpoint{-0.016667in}{0.000000in}}%
\pgfpathcurveto{\pgfqpoint{-0.016667in}{-0.004420in}}{\pgfqpoint{-0.014911in}{-0.008660in}}{\pgfqpoint{-0.011785in}{-0.011785in}}%
\pgfpathcurveto{\pgfqpoint{-0.008660in}{-0.014911in}}{\pgfqpoint{-0.004420in}{-0.016667in}}{\pgfqpoint{0.000000in}{-0.016667in}}%
\pgfpathlineto{\pgfqpoint{0.000000in}{-0.016667in}}%
\pgfpathclose%
\pgfusepath{stroke,fill}%
}%
\begin{pgfscope}%
\pgfsys@transformshift{2.922572in}{4.710525in}%
\pgfsys@useobject{currentmarker}{}%
\end{pgfscope}%
\end{pgfscope}%
\begin{pgfscope}%
\pgfpathrectangle{\pgfqpoint{0.100000in}{2.493063in}}{\pgfqpoint{5.037500in}{3.427208in}}%
\pgfusepath{clip}%
\pgfsetrectcap%
\pgfsetroundjoin%
\pgfsetlinewidth{1.505625pt}%
\definecolor{currentstroke}{rgb}{0.000000,0.000000,1.000000}%
\pgfsetstrokecolor{currentstroke}%
\pgfsetstrokeopacity{0.500000}%
\pgfsetdash{}{0pt}%
\pgfpathmoveto{\pgfqpoint{3.089457in}{4.708984in}}%
\pgfusepath{stroke}%
\end{pgfscope}%
\begin{pgfscope}%
\pgfpathrectangle{\pgfqpoint{0.100000in}{2.493063in}}{\pgfqpoint{5.037500in}{3.427208in}}%
\pgfusepath{clip}%
\pgfsetbuttcap%
\pgfsetroundjoin%
\definecolor{currentfill}{rgb}{0.000000,0.000000,1.000000}%
\pgfsetfillcolor{currentfill}%
\pgfsetfillopacity{0.500000}%
\pgfsetlinewidth{0.250937pt}%
\definecolor{currentstroke}{rgb}{0.000000,0.000000,0.000000}%
\pgfsetstrokecolor{currentstroke}%
\pgfsetstrokeopacity{0.500000}%
\pgfsetdash{}{0pt}%
\pgfsys@defobject{currentmarker}{\pgfqpoint{-0.019444in}{-0.019444in}}{\pgfqpoint{0.019444in}{0.019444in}}{%
\pgfpathmoveto{\pgfqpoint{0.000000in}{-0.019444in}}%
\pgfpathcurveto{\pgfqpoint{0.005157in}{-0.019444in}}{\pgfqpoint{0.010103in}{-0.017396in}}{\pgfqpoint{0.013749in}{-0.013749in}}%
\pgfpathcurveto{\pgfqpoint{0.017396in}{-0.010103in}}{\pgfqpoint{0.019444in}{-0.005157in}}{\pgfqpoint{0.019444in}{0.000000in}}%
\pgfpathcurveto{\pgfqpoint{0.019444in}{0.005157in}}{\pgfqpoint{0.017396in}{0.010103in}}{\pgfqpoint{0.013749in}{0.013749in}}%
\pgfpathcurveto{\pgfqpoint{0.010103in}{0.017396in}}{\pgfqpoint{0.005157in}{0.019444in}}{\pgfqpoint{0.000000in}{0.019444in}}%
\pgfpathcurveto{\pgfqpoint{-0.005157in}{0.019444in}}{\pgfqpoint{-0.010103in}{0.017396in}}{\pgfqpoint{-0.013749in}{0.013749in}}%
\pgfpathcurveto{\pgfqpoint{-0.017396in}{0.010103in}}{\pgfqpoint{-0.019444in}{0.005157in}}{\pgfqpoint{-0.019444in}{0.000000in}}%
\pgfpathcurveto{\pgfqpoint{-0.019444in}{-0.005157in}}{\pgfqpoint{-0.017396in}{-0.010103in}}{\pgfqpoint{-0.013749in}{-0.013749in}}%
\pgfpathcurveto{\pgfqpoint{-0.010103in}{-0.017396in}}{\pgfqpoint{-0.005157in}{-0.019444in}}{\pgfqpoint{0.000000in}{-0.019444in}}%
\pgfpathlineto{\pgfqpoint{0.000000in}{-0.019444in}}%
\pgfpathclose%
\pgfusepath{stroke,fill}%
}%
\begin{pgfscope}%
\pgfsys@transformshift{3.089457in}{4.708984in}%
\pgfsys@useobject{currentmarker}{}%
\end{pgfscope}%
\end{pgfscope}%
\begin{pgfscope}%
\pgfpathrectangle{\pgfqpoint{0.100000in}{2.493063in}}{\pgfqpoint{5.037500in}{3.427208in}}%
\pgfusepath{clip}%
\pgfsetrectcap%
\pgfsetroundjoin%
\pgfsetlinewidth{1.505625pt}%
\definecolor{currentstroke}{rgb}{0.000000,0.000000,1.000000}%
\pgfsetstrokecolor{currentstroke}%
\pgfsetstrokeopacity{0.500000}%
\pgfsetdash{}{0pt}%
\pgfpathmoveto{\pgfqpoint{2.924486in}{4.660320in}}%
\pgfusepath{stroke}%
\end{pgfscope}%
\begin{pgfscope}%
\pgfpathrectangle{\pgfqpoint{0.100000in}{2.493063in}}{\pgfqpoint{5.037500in}{3.427208in}}%
\pgfusepath{clip}%
\pgfsetbuttcap%
\pgfsetroundjoin%
\definecolor{currentfill}{rgb}{0.000000,0.000000,1.000000}%
\pgfsetfillcolor{currentfill}%
\pgfsetfillopacity{0.500000}%
\pgfsetlinewidth{0.250937pt}%
\definecolor{currentstroke}{rgb}{0.000000,0.000000,0.000000}%
\pgfsetstrokecolor{currentstroke}%
\pgfsetstrokeopacity{0.500000}%
\pgfsetdash{}{0pt}%
\pgfsys@defobject{currentmarker}{\pgfqpoint{-0.013889in}{-0.013889in}}{\pgfqpoint{0.013889in}{0.013889in}}{%
\pgfpathmoveto{\pgfqpoint{0.000000in}{-0.013889in}}%
\pgfpathcurveto{\pgfqpoint{0.003683in}{-0.013889in}}{\pgfqpoint{0.007216in}{-0.012425in}}{\pgfqpoint{0.009821in}{-0.009821in}}%
\pgfpathcurveto{\pgfqpoint{0.012425in}{-0.007216in}}{\pgfqpoint{0.013889in}{-0.003683in}}{\pgfqpoint{0.013889in}{0.000000in}}%
\pgfpathcurveto{\pgfqpoint{0.013889in}{0.003683in}}{\pgfqpoint{0.012425in}{0.007216in}}{\pgfqpoint{0.009821in}{0.009821in}}%
\pgfpathcurveto{\pgfqpoint{0.007216in}{0.012425in}}{\pgfqpoint{0.003683in}{0.013889in}}{\pgfqpoint{0.000000in}{0.013889in}}%
\pgfpathcurveto{\pgfqpoint{-0.003683in}{0.013889in}}{\pgfqpoint{-0.007216in}{0.012425in}}{\pgfqpoint{-0.009821in}{0.009821in}}%
\pgfpathcurveto{\pgfqpoint{-0.012425in}{0.007216in}}{\pgfqpoint{-0.013889in}{0.003683in}}{\pgfqpoint{-0.013889in}{0.000000in}}%
\pgfpathcurveto{\pgfqpoint{-0.013889in}{-0.003683in}}{\pgfqpoint{-0.012425in}{-0.007216in}}{\pgfqpoint{-0.009821in}{-0.009821in}}%
\pgfpathcurveto{\pgfqpoint{-0.007216in}{-0.012425in}}{\pgfqpoint{-0.003683in}{-0.013889in}}{\pgfqpoint{0.000000in}{-0.013889in}}%
\pgfpathlineto{\pgfqpoint{0.000000in}{-0.013889in}}%
\pgfpathclose%
\pgfusepath{stroke,fill}%
}%
\begin{pgfscope}%
\pgfsys@transformshift{2.924486in}{4.660320in}%
\pgfsys@useobject{currentmarker}{}%
\end{pgfscope}%
\end{pgfscope}%
\begin{pgfscope}%
\pgfpathrectangle{\pgfqpoint{0.100000in}{2.493063in}}{\pgfqpoint{5.037500in}{3.427208in}}%
\pgfusepath{clip}%
\pgfsetrectcap%
\pgfsetroundjoin%
\pgfsetlinewidth{1.505625pt}%
\definecolor{currentstroke}{rgb}{0.000000,0.000000,1.000000}%
\pgfsetstrokecolor{currentstroke}%
\pgfsetstrokeopacity{0.500000}%
\pgfsetdash{}{0pt}%
\pgfpathmoveto{\pgfqpoint{3.172689in}{4.773064in}}%
\pgfusepath{stroke}%
\end{pgfscope}%
\begin{pgfscope}%
\pgfpathrectangle{\pgfqpoint{0.100000in}{2.493063in}}{\pgfqpoint{5.037500in}{3.427208in}}%
\pgfusepath{clip}%
\pgfsetbuttcap%
\pgfsetroundjoin%
\definecolor{currentfill}{rgb}{0.000000,0.000000,1.000000}%
\pgfsetfillcolor{currentfill}%
\pgfsetfillopacity{0.500000}%
\pgfsetlinewidth{0.250937pt}%
\definecolor{currentstroke}{rgb}{0.000000,0.000000,0.000000}%
\pgfsetstrokecolor{currentstroke}%
\pgfsetstrokeopacity{0.500000}%
\pgfsetdash{}{0pt}%
\pgfsys@defobject{currentmarker}{\pgfqpoint{-0.013889in}{-0.013889in}}{\pgfqpoint{0.013889in}{0.013889in}}{%
\pgfpathmoveto{\pgfqpoint{0.000000in}{-0.013889in}}%
\pgfpathcurveto{\pgfqpoint{0.003683in}{-0.013889in}}{\pgfqpoint{0.007216in}{-0.012425in}}{\pgfqpoint{0.009821in}{-0.009821in}}%
\pgfpathcurveto{\pgfqpoint{0.012425in}{-0.007216in}}{\pgfqpoint{0.013889in}{-0.003683in}}{\pgfqpoint{0.013889in}{0.000000in}}%
\pgfpathcurveto{\pgfqpoint{0.013889in}{0.003683in}}{\pgfqpoint{0.012425in}{0.007216in}}{\pgfqpoint{0.009821in}{0.009821in}}%
\pgfpathcurveto{\pgfqpoint{0.007216in}{0.012425in}}{\pgfqpoint{0.003683in}{0.013889in}}{\pgfqpoint{0.000000in}{0.013889in}}%
\pgfpathcurveto{\pgfqpoint{-0.003683in}{0.013889in}}{\pgfqpoint{-0.007216in}{0.012425in}}{\pgfqpoint{-0.009821in}{0.009821in}}%
\pgfpathcurveto{\pgfqpoint{-0.012425in}{0.007216in}}{\pgfqpoint{-0.013889in}{0.003683in}}{\pgfqpoint{-0.013889in}{0.000000in}}%
\pgfpathcurveto{\pgfqpoint{-0.013889in}{-0.003683in}}{\pgfqpoint{-0.012425in}{-0.007216in}}{\pgfqpoint{-0.009821in}{-0.009821in}}%
\pgfpathcurveto{\pgfqpoint{-0.007216in}{-0.012425in}}{\pgfqpoint{-0.003683in}{-0.013889in}}{\pgfqpoint{0.000000in}{-0.013889in}}%
\pgfpathlineto{\pgfqpoint{0.000000in}{-0.013889in}}%
\pgfpathclose%
\pgfusepath{stroke,fill}%
}%
\begin{pgfscope}%
\pgfsys@transformshift{3.172689in}{4.773064in}%
\pgfsys@useobject{currentmarker}{}%
\end{pgfscope}%
\end{pgfscope}%
\begin{pgfscope}%
\pgfpathrectangle{\pgfqpoint{0.100000in}{2.493063in}}{\pgfqpoint{5.037500in}{3.427208in}}%
\pgfusepath{clip}%
\pgfsetrectcap%
\pgfsetroundjoin%
\pgfsetlinewidth{1.505625pt}%
\definecolor{currentstroke}{rgb}{0.000000,0.000000,1.000000}%
\pgfsetstrokecolor{currentstroke}%
\pgfsetstrokeopacity{0.500000}%
\pgfsetdash{}{0pt}%
\pgfpathmoveto{\pgfqpoint{3.102876in}{4.673191in}}%
\pgfusepath{stroke}%
\end{pgfscope}%
\begin{pgfscope}%
\pgfpathrectangle{\pgfqpoint{0.100000in}{2.493063in}}{\pgfqpoint{5.037500in}{3.427208in}}%
\pgfusepath{clip}%
\pgfsetbuttcap%
\pgfsetroundjoin%
\definecolor{currentfill}{rgb}{0.000000,0.000000,1.000000}%
\pgfsetfillcolor{currentfill}%
\pgfsetfillopacity{0.500000}%
\pgfsetlinewidth{0.250937pt}%
\definecolor{currentstroke}{rgb}{0.000000,0.000000,0.000000}%
\pgfsetstrokecolor{currentstroke}%
\pgfsetstrokeopacity{0.500000}%
\pgfsetdash{}{0pt}%
\pgfsys@defobject{currentmarker}{\pgfqpoint{-0.016667in}{-0.016667in}}{\pgfqpoint{0.016667in}{0.016667in}}{%
\pgfpathmoveto{\pgfqpoint{0.000000in}{-0.016667in}}%
\pgfpathcurveto{\pgfqpoint{0.004420in}{-0.016667in}}{\pgfqpoint{0.008660in}{-0.014911in}}{\pgfqpoint{0.011785in}{-0.011785in}}%
\pgfpathcurveto{\pgfqpoint{0.014911in}{-0.008660in}}{\pgfqpoint{0.016667in}{-0.004420in}}{\pgfqpoint{0.016667in}{0.000000in}}%
\pgfpathcurveto{\pgfqpoint{0.016667in}{0.004420in}}{\pgfqpoint{0.014911in}{0.008660in}}{\pgfqpoint{0.011785in}{0.011785in}}%
\pgfpathcurveto{\pgfqpoint{0.008660in}{0.014911in}}{\pgfqpoint{0.004420in}{0.016667in}}{\pgfqpoint{0.000000in}{0.016667in}}%
\pgfpathcurveto{\pgfqpoint{-0.004420in}{0.016667in}}{\pgfqpoint{-0.008660in}{0.014911in}}{\pgfqpoint{-0.011785in}{0.011785in}}%
\pgfpathcurveto{\pgfqpoint{-0.014911in}{0.008660in}}{\pgfqpoint{-0.016667in}{0.004420in}}{\pgfqpoint{-0.016667in}{0.000000in}}%
\pgfpathcurveto{\pgfqpoint{-0.016667in}{-0.004420in}}{\pgfqpoint{-0.014911in}{-0.008660in}}{\pgfqpoint{-0.011785in}{-0.011785in}}%
\pgfpathcurveto{\pgfqpoint{-0.008660in}{-0.014911in}}{\pgfqpoint{-0.004420in}{-0.016667in}}{\pgfqpoint{0.000000in}{-0.016667in}}%
\pgfpathlineto{\pgfqpoint{0.000000in}{-0.016667in}}%
\pgfpathclose%
\pgfusepath{stroke,fill}%
}%
\begin{pgfscope}%
\pgfsys@transformshift{3.102876in}{4.673191in}%
\pgfsys@useobject{currentmarker}{}%
\end{pgfscope}%
\end{pgfscope}%
\begin{pgfscope}%
\pgfpathrectangle{\pgfqpoint{0.100000in}{2.493063in}}{\pgfqpoint{5.037500in}{3.427208in}}%
\pgfusepath{clip}%
\pgfsetrectcap%
\pgfsetroundjoin%
\pgfsetlinewidth{1.505625pt}%
\definecolor{currentstroke}{rgb}{0.501961,0.501961,0.501961}%
\pgfsetstrokecolor{currentstroke}%
\pgfsetstrokeopacity{0.500000}%
\pgfsetdash{}{0pt}%
\pgfpathmoveto{\pgfqpoint{2.684559in}{4.764742in}}%
\pgfusepath{stroke}%
\end{pgfscope}%
\begin{pgfscope}%
\pgfpathrectangle{\pgfqpoint{0.100000in}{2.493063in}}{\pgfqpoint{5.037500in}{3.427208in}}%
\pgfusepath{clip}%
\pgfsetbuttcap%
\pgfsetroundjoin%
\definecolor{currentfill}{rgb}{0.501961,0.501961,0.501961}%
\pgfsetfillcolor{currentfill}%
\pgfsetfillopacity{0.500000}%
\pgfsetlinewidth{0.250937pt}%
\definecolor{currentstroke}{rgb}{0.000000,0.000000,0.000000}%
\pgfsetstrokecolor{currentstroke}%
\pgfsetstrokeopacity{0.500000}%
\pgfsetdash{}{0pt}%
\pgfsys@defobject{currentmarker}{\pgfqpoint{-0.013889in}{-0.013889in}}{\pgfqpoint{0.013889in}{0.013889in}}{%
\pgfpathmoveto{\pgfqpoint{0.000000in}{-0.013889in}}%
\pgfpathcurveto{\pgfqpoint{0.003683in}{-0.013889in}}{\pgfqpoint{0.007216in}{-0.012425in}}{\pgfqpoint{0.009821in}{-0.009821in}}%
\pgfpathcurveto{\pgfqpoint{0.012425in}{-0.007216in}}{\pgfqpoint{0.013889in}{-0.003683in}}{\pgfqpoint{0.013889in}{0.000000in}}%
\pgfpathcurveto{\pgfqpoint{0.013889in}{0.003683in}}{\pgfqpoint{0.012425in}{0.007216in}}{\pgfqpoint{0.009821in}{0.009821in}}%
\pgfpathcurveto{\pgfqpoint{0.007216in}{0.012425in}}{\pgfqpoint{0.003683in}{0.013889in}}{\pgfqpoint{0.000000in}{0.013889in}}%
\pgfpathcurveto{\pgfqpoint{-0.003683in}{0.013889in}}{\pgfqpoint{-0.007216in}{0.012425in}}{\pgfqpoint{-0.009821in}{0.009821in}}%
\pgfpathcurveto{\pgfqpoint{-0.012425in}{0.007216in}}{\pgfqpoint{-0.013889in}{0.003683in}}{\pgfqpoint{-0.013889in}{0.000000in}}%
\pgfpathcurveto{\pgfqpoint{-0.013889in}{-0.003683in}}{\pgfqpoint{-0.012425in}{-0.007216in}}{\pgfqpoint{-0.009821in}{-0.009821in}}%
\pgfpathcurveto{\pgfqpoint{-0.007216in}{-0.012425in}}{\pgfqpoint{-0.003683in}{-0.013889in}}{\pgfqpoint{0.000000in}{-0.013889in}}%
\pgfpathlineto{\pgfqpoint{0.000000in}{-0.013889in}}%
\pgfpathclose%
\pgfusepath{stroke,fill}%
}%
\begin{pgfscope}%
\pgfsys@transformshift{2.684559in}{4.764742in}%
\pgfsys@useobject{currentmarker}{}%
\end{pgfscope}%
\end{pgfscope}%
\begin{pgfscope}%
\pgfpathrectangle{\pgfqpoint{0.100000in}{2.493063in}}{\pgfqpoint{5.037500in}{3.427208in}}%
\pgfusepath{clip}%
\pgfsetrectcap%
\pgfsetroundjoin%
\pgfsetlinewidth{1.505625pt}%
\definecolor{currentstroke}{rgb}{0.000000,0.000000,1.000000}%
\pgfsetstrokecolor{currentstroke}%
\pgfsetstrokeopacity{0.500000}%
\pgfsetdash{}{0pt}%
\pgfpathmoveto{\pgfqpoint{3.030598in}{4.767321in}}%
\pgfusepath{stroke}%
\end{pgfscope}%
\begin{pgfscope}%
\pgfpathrectangle{\pgfqpoint{0.100000in}{2.493063in}}{\pgfqpoint{5.037500in}{3.427208in}}%
\pgfusepath{clip}%
\pgfsetbuttcap%
\pgfsetroundjoin%
\definecolor{currentfill}{rgb}{0.000000,0.000000,1.000000}%
\pgfsetfillcolor{currentfill}%
\pgfsetfillopacity{0.500000}%
\pgfsetlinewidth{0.250937pt}%
\definecolor{currentstroke}{rgb}{0.000000,0.000000,0.000000}%
\pgfsetstrokecolor{currentstroke}%
\pgfsetstrokeopacity{0.500000}%
\pgfsetdash{}{0pt}%
\pgfsys@defobject{currentmarker}{\pgfqpoint{-0.008333in}{-0.008333in}}{\pgfqpoint{0.008333in}{0.008333in}}{%
\pgfpathmoveto{\pgfqpoint{0.000000in}{-0.008333in}}%
\pgfpathcurveto{\pgfqpoint{0.002210in}{-0.008333in}}{\pgfqpoint{0.004330in}{-0.007455in}}{\pgfqpoint{0.005893in}{-0.005893in}}%
\pgfpathcurveto{\pgfqpoint{0.007455in}{-0.004330in}}{\pgfqpoint{0.008333in}{-0.002210in}}{\pgfqpoint{0.008333in}{0.000000in}}%
\pgfpathcurveto{\pgfqpoint{0.008333in}{0.002210in}}{\pgfqpoint{0.007455in}{0.004330in}}{\pgfqpoint{0.005893in}{0.005893in}}%
\pgfpathcurveto{\pgfqpoint{0.004330in}{0.007455in}}{\pgfqpoint{0.002210in}{0.008333in}}{\pgfqpoint{0.000000in}{0.008333in}}%
\pgfpathcurveto{\pgfqpoint{-0.002210in}{0.008333in}}{\pgfqpoint{-0.004330in}{0.007455in}}{\pgfqpoint{-0.005893in}{0.005893in}}%
\pgfpathcurveto{\pgfqpoint{-0.007455in}{0.004330in}}{\pgfqpoint{-0.008333in}{0.002210in}}{\pgfqpoint{-0.008333in}{0.000000in}}%
\pgfpathcurveto{\pgfqpoint{-0.008333in}{-0.002210in}}{\pgfqpoint{-0.007455in}{-0.004330in}}{\pgfqpoint{-0.005893in}{-0.005893in}}%
\pgfpathcurveto{\pgfqpoint{-0.004330in}{-0.007455in}}{\pgfqpoint{-0.002210in}{-0.008333in}}{\pgfqpoint{0.000000in}{-0.008333in}}%
\pgfpathlineto{\pgfqpoint{0.000000in}{-0.008333in}}%
\pgfpathclose%
\pgfusepath{stroke,fill}%
}%
\begin{pgfscope}%
\pgfsys@transformshift{3.030598in}{4.767321in}%
\pgfsys@useobject{currentmarker}{}%
\end{pgfscope}%
\end{pgfscope}%
\begin{pgfscope}%
\pgfpathrectangle{\pgfqpoint{0.100000in}{2.493063in}}{\pgfqpoint{5.037500in}{3.427208in}}%
\pgfusepath{clip}%
\pgfsetrectcap%
\pgfsetroundjoin%
\pgfsetlinewidth{1.505625pt}%
\definecolor{currentstroke}{rgb}{0.678431,1.000000,0.184314}%
\pgfsetstrokecolor{currentstroke}%
\pgfsetstrokeopacity{0.500000}%
\pgfsetdash{}{0pt}%
\pgfpathmoveto{\pgfqpoint{2.783006in}{4.357745in}}%
\pgfusepath{stroke}%
\end{pgfscope}%
\begin{pgfscope}%
\pgfpathrectangle{\pgfqpoint{0.100000in}{2.493063in}}{\pgfqpoint{5.037500in}{3.427208in}}%
\pgfusepath{clip}%
\pgfsetbuttcap%
\pgfsetroundjoin%
\definecolor{currentfill}{rgb}{0.678431,1.000000,0.184314}%
\pgfsetfillcolor{currentfill}%
\pgfsetfillopacity{0.500000}%
\pgfsetlinewidth{0.250937pt}%
\definecolor{currentstroke}{rgb}{0.000000,0.000000,0.000000}%
\pgfsetstrokecolor{currentstroke}%
\pgfsetstrokeopacity{0.500000}%
\pgfsetdash{}{0pt}%
\pgfsys@defobject{currentmarker}{\pgfqpoint{-0.025000in}{-0.025000in}}{\pgfqpoint{0.025000in}{0.025000in}}{%
\pgfpathmoveto{\pgfqpoint{0.000000in}{-0.025000in}}%
\pgfpathcurveto{\pgfqpoint{0.006630in}{-0.025000in}}{\pgfqpoint{0.012989in}{-0.022366in}}{\pgfqpoint{0.017678in}{-0.017678in}}%
\pgfpathcurveto{\pgfqpoint{0.022366in}{-0.012989in}}{\pgfqpoint{0.025000in}{-0.006630in}}{\pgfqpoint{0.025000in}{0.000000in}}%
\pgfpathcurveto{\pgfqpoint{0.025000in}{0.006630in}}{\pgfqpoint{0.022366in}{0.012989in}}{\pgfqpoint{0.017678in}{0.017678in}}%
\pgfpathcurveto{\pgfqpoint{0.012989in}{0.022366in}}{\pgfqpoint{0.006630in}{0.025000in}}{\pgfqpoint{0.000000in}{0.025000in}}%
\pgfpathcurveto{\pgfqpoint{-0.006630in}{0.025000in}}{\pgfqpoint{-0.012989in}{0.022366in}}{\pgfqpoint{-0.017678in}{0.017678in}}%
\pgfpathcurveto{\pgfqpoint{-0.022366in}{0.012989in}}{\pgfqpoint{-0.025000in}{0.006630in}}{\pgfqpoint{-0.025000in}{0.000000in}}%
\pgfpathcurveto{\pgfqpoint{-0.025000in}{-0.006630in}}{\pgfqpoint{-0.022366in}{-0.012989in}}{\pgfqpoint{-0.017678in}{-0.017678in}}%
\pgfpathcurveto{\pgfqpoint{-0.012989in}{-0.022366in}}{\pgfqpoint{-0.006630in}{-0.025000in}}{\pgfqpoint{0.000000in}{-0.025000in}}%
\pgfpathlineto{\pgfqpoint{0.000000in}{-0.025000in}}%
\pgfpathclose%
\pgfusepath{stroke,fill}%
}%
\begin{pgfscope}%
\pgfsys@transformshift{2.783006in}{4.357745in}%
\pgfsys@useobject{currentmarker}{}%
\end{pgfscope}%
\end{pgfscope}%
\begin{pgfscope}%
\pgfpathrectangle{\pgfqpoint{0.100000in}{2.493063in}}{\pgfqpoint{5.037500in}{3.427208in}}%
\pgfusepath{clip}%
\pgfsetrectcap%
\pgfsetroundjoin%
\pgfsetlinewidth{1.505625pt}%
\definecolor{currentstroke}{rgb}{0.678431,1.000000,0.184314}%
\pgfsetstrokecolor{currentstroke}%
\pgfsetstrokeopacity{0.500000}%
\pgfsetdash{}{0pt}%
\pgfpathmoveto{\pgfqpoint{2.663858in}{4.383304in}}%
\pgfusepath{stroke}%
\end{pgfscope}%
\begin{pgfscope}%
\pgfpathrectangle{\pgfqpoint{0.100000in}{2.493063in}}{\pgfqpoint{5.037500in}{3.427208in}}%
\pgfusepath{clip}%
\pgfsetbuttcap%
\pgfsetroundjoin%
\definecolor{currentfill}{rgb}{0.678431,1.000000,0.184314}%
\pgfsetfillcolor{currentfill}%
\pgfsetfillopacity{0.500000}%
\pgfsetlinewidth{0.250937pt}%
\definecolor{currentstroke}{rgb}{0.000000,0.000000,0.000000}%
\pgfsetstrokecolor{currentstroke}%
\pgfsetstrokeopacity{0.500000}%
\pgfsetdash{}{0pt}%
\pgfsys@defobject{currentmarker}{\pgfqpoint{-0.022222in}{-0.022222in}}{\pgfqpoint{0.022222in}{0.022222in}}{%
\pgfpathmoveto{\pgfqpoint{0.000000in}{-0.022222in}}%
\pgfpathcurveto{\pgfqpoint{0.005893in}{-0.022222in}}{\pgfqpoint{0.011546in}{-0.019881in}}{\pgfqpoint{0.015713in}{-0.015713in}}%
\pgfpathcurveto{\pgfqpoint{0.019881in}{-0.011546in}}{\pgfqpoint{0.022222in}{-0.005893in}}{\pgfqpoint{0.022222in}{0.000000in}}%
\pgfpathcurveto{\pgfqpoint{0.022222in}{0.005893in}}{\pgfqpoint{0.019881in}{0.011546in}}{\pgfqpoint{0.015713in}{0.015713in}}%
\pgfpathcurveto{\pgfqpoint{0.011546in}{0.019881in}}{\pgfqpoint{0.005893in}{0.022222in}}{\pgfqpoint{0.000000in}{0.022222in}}%
\pgfpathcurveto{\pgfqpoint{-0.005893in}{0.022222in}}{\pgfqpoint{-0.011546in}{0.019881in}}{\pgfqpoint{-0.015713in}{0.015713in}}%
\pgfpathcurveto{\pgfqpoint{-0.019881in}{0.011546in}}{\pgfqpoint{-0.022222in}{0.005893in}}{\pgfqpoint{-0.022222in}{0.000000in}}%
\pgfpathcurveto{\pgfqpoint{-0.022222in}{-0.005893in}}{\pgfqpoint{-0.019881in}{-0.011546in}}{\pgfqpoint{-0.015713in}{-0.015713in}}%
\pgfpathcurveto{\pgfqpoint{-0.011546in}{-0.019881in}}{\pgfqpoint{-0.005893in}{-0.022222in}}{\pgfqpoint{0.000000in}{-0.022222in}}%
\pgfpathlineto{\pgfqpoint{0.000000in}{-0.022222in}}%
\pgfpathclose%
\pgfusepath{stroke,fill}%
}%
\begin{pgfscope}%
\pgfsys@transformshift{2.663858in}{4.383304in}%
\pgfsys@useobject{currentmarker}{}%
\end{pgfscope}%
\end{pgfscope}%
\begin{pgfscope}%
\pgfpathrectangle{\pgfqpoint{0.100000in}{2.493063in}}{\pgfqpoint{5.037500in}{3.427208in}}%
\pgfusepath{clip}%
\pgfsetrectcap%
\pgfsetroundjoin%
\pgfsetlinewidth{1.505625pt}%
\definecolor{currentstroke}{rgb}{0.678431,1.000000,0.184314}%
\pgfsetstrokecolor{currentstroke}%
\pgfsetstrokeopacity{0.500000}%
\pgfsetdash{}{0pt}%
\pgfpathmoveto{\pgfqpoint{2.743539in}{4.366819in}}%
\pgfusepath{stroke}%
\end{pgfscope}%
\begin{pgfscope}%
\pgfpathrectangle{\pgfqpoint{0.100000in}{2.493063in}}{\pgfqpoint{5.037500in}{3.427208in}}%
\pgfusepath{clip}%
\pgfsetbuttcap%
\pgfsetroundjoin%
\definecolor{currentfill}{rgb}{0.678431,1.000000,0.184314}%
\pgfsetfillcolor{currentfill}%
\pgfsetfillopacity{0.500000}%
\pgfsetlinewidth{0.250937pt}%
\definecolor{currentstroke}{rgb}{0.000000,0.000000,0.000000}%
\pgfsetstrokecolor{currentstroke}%
\pgfsetstrokeopacity{0.500000}%
\pgfsetdash{}{0pt}%
\pgfsys@defobject{currentmarker}{\pgfqpoint{-0.022222in}{-0.022222in}}{\pgfqpoint{0.022222in}{0.022222in}}{%
\pgfpathmoveto{\pgfqpoint{0.000000in}{-0.022222in}}%
\pgfpathcurveto{\pgfqpoint{0.005893in}{-0.022222in}}{\pgfqpoint{0.011546in}{-0.019881in}}{\pgfqpoint{0.015713in}{-0.015713in}}%
\pgfpathcurveto{\pgfqpoint{0.019881in}{-0.011546in}}{\pgfqpoint{0.022222in}{-0.005893in}}{\pgfqpoint{0.022222in}{0.000000in}}%
\pgfpathcurveto{\pgfqpoint{0.022222in}{0.005893in}}{\pgfqpoint{0.019881in}{0.011546in}}{\pgfqpoint{0.015713in}{0.015713in}}%
\pgfpathcurveto{\pgfqpoint{0.011546in}{0.019881in}}{\pgfqpoint{0.005893in}{0.022222in}}{\pgfqpoint{0.000000in}{0.022222in}}%
\pgfpathcurveto{\pgfqpoint{-0.005893in}{0.022222in}}{\pgfqpoint{-0.011546in}{0.019881in}}{\pgfqpoint{-0.015713in}{0.015713in}}%
\pgfpathcurveto{\pgfqpoint{-0.019881in}{0.011546in}}{\pgfqpoint{-0.022222in}{0.005893in}}{\pgfqpoint{-0.022222in}{0.000000in}}%
\pgfpathcurveto{\pgfqpoint{-0.022222in}{-0.005893in}}{\pgfqpoint{-0.019881in}{-0.011546in}}{\pgfqpoint{-0.015713in}{-0.015713in}}%
\pgfpathcurveto{\pgfqpoint{-0.011546in}{-0.019881in}}{\pgfqpoint{-0.005893in}{-0.022222in}}{\pgfqpoint{0.000000in}{-0.022222in}}%
\pgfpathlineto{\pgfqpoint{0.000000in}{-0.022222in}}%
\pgfpathclose%
\pgfusepath{stroke,fill}%
}%
\begin{pgfscope}%
\pgfsys@transformshift{2.743539in}{4.366819in}%
\pgfsys@useobject{currentmarker}{}%
\end{pgfscope}%
\end{pgfscope}%
\begin{pgfscope}%
\pgfpathrectangle{\pgfqpoint{0.100000in}{2.493063in}}{\pgfqpoint{5.037500in}{3.427208in}}%
\pgfusepath{clip}%
\pgfsetrectcap%
\pgfsetroundjoin%
\pgfsetlinewidth{1.505625pt}%
\definecolor{currentstroke}{rgb}{0.678431,1.000000,0.184314}%
\pgfsetstrokecolor{currentstroke}%
\pgfsetstrokeopacity{0.500000}%
\pgfsetdash{}{0pt}%
\pgfpathmoveto{\pgfqpoint{2.591100in}{4.213088in}}%
\pgfusepath{stroke}%
\end{pgfscope}%
\begin{pgfscope}%
\pgfpathrectangle{\pgfqpoint{0.100000in}{2.493063in}}{\pgfqpoint{5.037500in}{3.427208in}}%
\pgfusepath{clip}%
\pgfsetbuttcap%
\pgfsetroundjoin%
\definecolor{currentfill}{rgb}{0.678431,1.000000,0.184314}%
\pgfsetfillcolor{currentfill}%
\pgfsetfillopacity{0.500000}%
\pgfsetlinewidth{0.250937pt}%
\definecolor{currentstroke}{rgb}{0.000000,0.000000,0.000000}%
\pgfsetstrokecolor{currentstroke}%
\pgfsetstrokeopacity{0.500000}%
\pgfsetdash{}{0pt}%
\pgfsys@defobject{currentmarker}{\pgfqpoint{-0.016667in}{-0.016667in}}{\pgfqpoint{0.016667in}{0.016667in}}{%
\pgfpathmoveto{\pgfqpoint{0.000000in}{-0.016667in}}%
\pgfpathcurveto{\pgfqpoint{0.004420in}{-0.016667in}}{\pgfqpoint{0.008660in}{-0.014911in}}{\pgfqpoint{0.011785in}{-0.011785in}}%
\pgfpathcurveto{\pgfqpoint{0.014911in}{-0.008660in}}{\pgfqpoint{0.016667in}{-0.004420in}}{\pgfqpoint{0.016667in}{0.000000in}}%
\pgfpathcurveto{\pgfqpoint{0.016667in}{0.004420in}}{\pgfqpoint{0.014911in}{0.008660in}}{\pgfqpoint{0.011785in}{0.011785in}}%
\pgfpathcurveto{\pgfqpoint{0.008660in}{0.014911in}}{\pgfqpoint{0.004420in}{0.016667in}}{\pgfqpoint{0.000000in}{0.016667in}}%
\pgfpathcurveto{\pgfqpoint{-0.004420in}{0.016667in}}{\pgfqpoint{-0.008660in}{0.014911in}}{\pgfqpoint{-0.011785in}{0.011785in}}%
\pgfpathcurveto{\pgfqpoint{-0.014911in}{0.008660in}}{\pgfqpoint{-0.016667in}{0.004420in}}{\pgfqpoint{-0.016667in}{0.000000in}}%
\pgfpathcurveto{\pgfqpoint{-0.016667in}{-0.004420in}}{\pgfqpoint{-0.014911in}{-0.008660in}}{\pgfqpoint{-0.011785in}{-0.011785in}}%
\pgfpathcurveto{\pgfqpoint{-0.008660in}{-0.014911in}}{\pgfqpoint{-0.004420in}{-0.016667in}}{\pgfqpoint{0.000000in}{-0.016667in}}%
\pgfpathlineto{\pgfqpoint{0.000000in}{-0.016667in}}%
\pgfpathclose%
\pgfusepath{stroke,fill}%
}%
\begin{pgfscope}%
\pgfsys@transformshift{2.591100in}{4.213088in}%
\pgfsys@useobject{currentmarker}{}%
\end{pgfscope}%
\end{pgfscope}%
\begin{pgfscope}%
\pgfpathrectangle{\pgfqpoint{0.100000in}{2.493063in}}{\pgfqpoint{5.037500in}{3.427208in}}%
\pgfusepath{clip}%
\pgfsetrectcap%
\pgfsetroundjoin%
\pgfsetlinewidth{1.505625pt}%
\definecolor{currentstroke}{rgb}{0.000000,0.000000,1.000000}%
\pgfsetstrokecolor{currentstroke}%
\pgfsetstrokeopacity{0.500000}%
\pgfsetdash{}{0pt}%
\pgfpathmoveto{\pgfqpoint{3.590438in}{4.166531in}}%
\pgfusepath{stroke}%
\end{pgfscope}%
\begin{pgfscope}%
\pgfpathrectangle{\pgfqpoint{0.100000in}{2.493063in}}{\pgfqpoint{5.037500in}{3.427208in}}%
\pgfusepath{clip}%
\pgfsetbuttcap%
\pgfsetroundjoin%
\definecolor{currentfill}{rgb}{0.000000,0.000000,1.000000}%
\pgfsetfillcolor{currentfill}%
\pgfsetfillopacity{0.500000}%
\pgfsetlinewidth{0.250937pt}%
\definecolor{currentstroke}{rgb}{0.000000,0.000000,0.000000}%
\pgfsetstrokecolor{currentstroke}%
\pgfsetstrokeopacity{0.500000}%
\pgfsetdash{}{0pt}%
\pgfsys@defobject{currentmarker}{\pgfqpoint{-0.016667in}{-0.016667in}}{\pgfqpoint{0.016667in}{0.016667in}}{%
\pgfpathmoveto{\pgfqpoint{0.000000in}{-0.016667in}}%
\pgfpathcurveto{\pgfqpoint{0.004420in}{-0.016667in}}{\pgfqpoint{0.008660in}{-0.014911in}}{\pgfqpoint{0.011785in}{-0.011785in}}%
\pgfpathcurveto{\pgfqpoint{0.014911in}{-0.008660in}}{\pgfqpoint{0.016667in}{-0.004420in}}{\pgfqpoint{0.016667in}{0.000000in}}%
\pgfpathcurveto{\pgfqpoint{0.016667in}{0.004420in}}{\pgfqpoint{0.014911in}{0.008660in}}{\pgfqpoint{0.011785in}{0.011785in}}%
\pgfpathcurveto{\pgfqpoint{0.008660in}{0.014911in}}{\pgfqpoint{0.004420in}{0.016667in}}{\pgfqpoint{0.000000in}{0.016667in}}%
\pgfpathcurveto{\pgfqpoint{-0.004420in}{0.016667in}}{\pgfqpoint{-0.008660in}{0.014911in}}{\pgfqpoint{-0.011785in}{0.011785in}}%
\pgfpathcurveto{\pgfqpoint{-0.014911in}{0.008660in}}{\pgfqpoint{-0.016667in}{0.004420in}}{\pgfqpoint{-0.016667in}{0.000000in}}%
\pgfpathcurveto{\pgfqpoint{-0.016667in}{-0.004420in}}{\pgfqpoint{-0.014911in}{-0.008660in}}{\pgfqpoint{-0.011785in}{-0.011785in}}%
\pgfpathcurveto{\pgfqpoint{-0.008660in}{-0.014911in}}{\pgfqpoint{-0.004420in}{-0.016667in}}{\pgfqpoint{0.000000in}{-0.016667in}}%
\pgfpathlineto{\pgfqpoint{0.000000in}{-0.016667in}}%
\pgfpathclose%
\pgfusepath{stroke,fill}%
}%
\begin{pgfscope}%
\pgfsys@transformshift{3.590438in}{4.166531in}%
\pgfsys@useobject{currentmarker}{}%
\end{pgfscope}%
\end{pgfscope}%
\begin{pgfscope}%
\pgfpathrectangle{\pgfqpoint{0.100000in}{2.493063in}}{\pgfqpoint{5.037500in}{3.427208in}}%
\pgfusepath{clip}%
\pgfsetrectcap%
\pgfsetroundjoin%
\pgfsetlinewidth{1.505625pt}%
\definecolor{currentstroke}{rgb}{0.000000,0.000000,1.000000}%
\pgfsetstrokecolor{currentstroke}%
\pgfsetstrokeopacity{0.500000}%
\pgfsetdash{}{0pt}%
\pgfpathmoveto{\pgfqpoint{3.635835in}{4.252410in}}%
\pgfusepath{stroke}%
\end{pgfscope}%
\begin{pgfscope}%
\pgfpathrectangle{\pgfqpoint{0.100000in}{2.493063in}}{\pgfqpoint{5.037500in}{3.427208in}}%
\pgfusepath{clip}%
\pgfsetbuttcap%
\pgfsetroundjoin%
\definecolor{currentfill}{rgb}{0.000000,0.000000,1.000000}%
\pgfsetfillcolor{currentfill}%
\pgfsetfillopacity{0.500000}%
\pgfsetlinewidth{0.250937pt}%
\definecolor{currentstroke}{rgb}{0.000000,0.000000,0.000000}%
\pgfsetstrokecolor{currentstroke}%
\pgfsetstrokeopacity{0.500000}%
\pgfsetdash{}{0pt}%
\pgfsys@defobject{currentmarker}{\pgfqpoint{-0.016667in}{-0.016667in}}{\pgfqpoint{0.016667in}{0.016667in}}{%
\pgfpathmoveto{\pgfqpoint{0.000000in}{-0.016667in}}%
\pgfpathcurveto{\pgfqpoint{0.004420in}{-0.016667in}}{\pgfqpoint{0.008660in}{-0.014911in}}{\pgfqpoint{0.011785in}{-0.011785in}}%
\pgfpathcurveto{\pgfqpoint{0.014911in}{-0.008660in}}{\pgfqpoint{0.016667in}{-0.004420in}}{\pgfqpoint{0.016667in}{0.000000in}}%
\pgfpathcurveto{\pgfqpoint{0.016667in}{0.004420in}}{\pgfqpoint{0.014911in}{0.008660in}}{\pgfqpoint{0.011785in}{0.011785in}}%
\pgfpathcurveto{\pgfqpoint{0.008660in}{0.014911in}}{\pgfqpoint{0.004420in}{0.016667in}}{\pgfqpoint{0.000000in}{0.016667in}}%
\pgfpathcurveto{\pgfqpoint{-0.004420in}{0.016667in}}{\pgfqpoint{-0.008660in}{0.014911in}}{\pgfqpoint{-0.011785in}{0.011785in}}%
\pgfpathcurveto{\pgfqpoint{-0.014911in}{0.008660in}}{\pgfqpoint{-0.016667in}{0.004420in}}{\pgfqpoint{-0.016667in}{0.000000in}}%
\pgfpathcurveto{\pgfqpoint{-0.016667in}{-0.004420in}}{\pgfqpoint{-0.014911in}{-0.008660in}}{\pgfqpoint{-0.011785in}{-0.011785in}}%
\pgfpathcurveto{\pgfqpoint{-0.008660in}{-0.014911in}}{\pgfqpoint{-0.004420in}{-0.016667in}}{\pgfqpoint{0.000000in}{-0.016667in}}%
\pgfpathlineto{\pgfqpoint{0.000000in}{-0.016667in}}%
\pgfpathclose%
\pgfusepath{stroke,fill}%
}%
\begin{pgfscope}%
\pgfsys@transformshift{3.635835in}{4.252410in}%
\pgfsys@useobject{currentmarker}{}%
\end{pgfscope}%
\end{pgfscope}%
\begin{pgfscope}%
\pgfpathrectangle{\pgfqpoint{0.100000in}{2.493063in}}{\pgfqpoint{5.037500in}{3.427208in}}%
\pgfusepath{clip}%
\pgfsetrectcap%
\pgfsetroundjoin%
\pgfsetlinewidth{1.505625pt}%
\definecolor{currentstroke}{rgb}{0.000000,0.000000,1.000000}%
\pgfsetstrokecolor{currentstroke}%
\pgfsetstrokeopacity{0.500000}%
\pgfsetdash{}{0pt}%
\pgfpathmoveto{\pgfqpoint{3.754570in}{4.306123in}}%
\pgfusepath{stroke}%
\end{pgfscope}%
\begin{pgfscope}%
\pgfpathrectangle{\pgfqpoint{0.100000in}{2.493063in}}{\pgfqpoint{5.037500in}{3.427208in}}%
\pgfusepath{clip}%
\pgfsetbuttcap%
\pgfsetroundjoin%
\definecolor{currentfill}{rgb}{0.000000,0.000000,1.000000}%
\pgfsetfillcolor{currentfill}%
\pgfsetfillopacity{0.500000}%
\pgfsetlinewidth{0.250937pt}%
\definecolor{currentstroke}{rgb}{0.000000,0.000000,0.000000}%
\pgfsetstrokecolor{currentstroke}%
\pgfsetstrokeopacity{0.500000}%
\pgfsetdash{}{0pt}%
\pgfsys@defobject{currentmarker}{\pgfqpoint{-0.013889in}{-0.013889in}}{\pgfqpoint{0.013889in}{0.013889in}}{%
\pgfpathmoveto{\pgfqpoint{0.000000in}{-0.013889in}}%
\pgfpathcurveto{\pgfqpoint{0.003683in}{-0.013889in}}{\pgfqpoint{0.007216in}{-0.012425in}}{\pgfqpoint{0.009821in}{-0.009821in}}%
\pgfpathcurveto{\pgfqpoint{0.012425in}{-0.007216in}}{\pgfqpoint{0.013889in}{-0.003683in}}{\pgfqpoint{0.013889in}{0.000000in}}%
\pgfpathcurveto{\pgfqpoint{0.013889in}{0.003683in}}{\pgfqpoint{0.012425in}{0.007216in}}{\pgfqpoint{0.009821in}{0.009821in}}%
\pgfpathcurveto{\pgfqpoint{0.007216in}{0.012425in}}{\pgfqpoint{0.003683in}{0.013889in}}{\pgfqpoint{0.000000in}{0.013889in}}%
\pgfpathcurveto{\pgfqpoint{-0.003683in}{0.013889in}}{\pgfqpoint{-0.007216in}{0.012425in}}{\pgfqpoint{-0.009821in}{0.009821in}}%
\pgfpathcurveto{\pgfqpoint{-0.012425in}{0.007216in}}{\pgfqpoint{-0.013889in}{0.003683in}}{\pgfqpoint{-0.013889in}{0.000000in}}%
\pgfpathcurveto{\pgfqpoint{-0.013889in}{-0.003683in}}{\pgfqpoint{-0.012425in}{-0.007216in}}{\pgfqpoint{-0.009821in}{-0.009821in}}%
\pgfpathcurveto{\pgfqpoint{-0.007216in}{-0.012425in}}{\pgfqpoint{-0.003683in}{-0.013889in}}{\pgfqpoint{0.000000in}{-0.013889in}}%
\pgfpathlineto{\pgfqpoint{0.000000in}{-0.013889in}}%
\pgfpathclose%
\pgfusepath{stroke,fill}%
}%
\begin{pgfscope}%
\pgfsys@transformshift{3.754570in}{4.306123in}%
\pgfsys@useobject{currentmarker}{}%
\end{pgfscope}%
\end{pgfscope}%
\begin{pgfscope}%
\pgfpathrectangle{\pgfqpoint{0.100000in}{2.493063in}}{\pgfqpoint{5.037500in}{3.427208in}}%
\pgfusepath{clip}%
\pgfsetrectcap%
\pgfsetroundjoin%
\pgfsetlinewidth{1.505625pt}%
\definecolor{currentstroke}{rgb}{0.000000,0.000000,1.000000}%
\pgfsetstrokecolor{currentstroke}%
\pgfsetstrokeopacity{0.500000}%
\pgfsetdash{}{0pt}%
\pgfpathmoveto{\pgfqpoint{3.638327in}{4.317522in}}%
\pgfusepath{stroke}%
\end{pgfscope}%
\begin{pgfscope}%
\pgfpathrectangle{\pgfqpoint{0.100000in}{2.493063in}}{\pgfqpoint{5.037500in}{3.427208in}}%
\pgfusepath{clip}%
\pgfsetbuttcap%
\pgfsetroundjoin%
\definecolor{currentfill}{rgb}{0.000000,0.000000,1.000000}%
\pgfsetfillcolor{currentfill}%
\pgfsetfillopacity{0.500000}%
\pgfsetlinewidth{0.250937pt}%
\definecolor{currentstroke}{rgb}{0.000000,0.000000,0.000000}%
\pgfsetstrokecolor{currentstroke}%
\pgfsetstrokeopacity{0.500000}%
\pgfsetdash{}{0pt}%
\pgfsys@defobject{currentmarker}{\pgfqpoint{-0.013889in}{-0.013889in}}{\pgfqpoint{0.013889in}{0.013889in}}{%
\pgfpathmoveto{\pgfqpoint{0.000000in}{-0.013889in}}%
\pgfpathcurveto{\pgfqpoint{0.003683in}{-0.013889in}}{\pgfqpoint{0.007216in}{-0.012425in}}{\pgfqpoint{0.009821in}{-0.009821in}}%
\pgfpathcurveto{\pgfqpoint{0.012425in}{-0.007216in}}{\pgfqpoint{0.013889in}{-0.003683in}}{\pgfqpoint{0.013889in}{0.000000in}}%
\pgfpathcurveto{\pgfqpoint{0.013889in}{0.003683in}}{\pgfqpoint{0.012425in}{0.007216in}}{\pgfqpoint{0.009821in}{0.009821in}}%
\pgfpathcurveto{\pgfqpoint{0.007216in}{0.012425in}}{\pgfqpoint{0.003683in}{0.013889in}}{\pgfqpoint{0.000000in}{0.013889in}}%
\pgfpathcurveto{\pgfqpoint{-0.003683in}{0.013889in}}{\pgfqpoint{-0.007216in}{0.012425in}}{\pgfqpoint{-0.009821in}{0.009821in}}%
\pgfpathcurveto{\pgfqpoint{-0.012425in}{0.007216in}}{\pgfqpoint{-0.013889in}{0.003683in}}{\pgfqpoint{-0.013889in}{0.000000in}}%
\pgfpathcurveto{\pgfqpoint{-0.013889in}{-0.003683in}}{\pgfqpoint{-0.012425in}{-0.007216in}}{\pgfqpoint{-0.009821in}{-0.009821in}}%
\pgfpathcurveto{\pgfqpoint{-0.007216in}{-0.012425in}}{\pgfqpoint{-0.003683in}{-0.013889in}}{\pgfqpoint{0.000000in}{-0.013889in}}%
\pgfpathlineto{\pgfqpoint{0.000000in}{-0.013889in}}%
\pgfpathclose%
\pgfusepath{stroke,fill}%
}%
\begin{pgfscope}%
\pgfsys@transformshift{3.638327in}{4.317522in}%
\pgfsys@useobject{currentmarker}{}%
\end{pgfscope}%
\end{pgfscope}%
\begin{pgfscope}%
\pgfpathrectangle{\pgfqpoint{0.100000in}{2.493063in}}{\pgfqpoint{5.037500in}{3.427208in}}%
\pgfusepath{clip}%
\pgfsetrectcap%
\pgfsetroundjoin%
\pgfsetlinewidth{1.505625pt}%
\definecolor{currentstroke}{rgb}{0.000000,0.000000,1.000000}%
\pgfsetstrokecolor{currentstroke}%
\pgfsetstrokeopacity{0.500000}%
\pgfsetdash{}{0pt}%
\pgfpathmoveto{\pgfqpoint{3.521226in}{4.250926in}}%
\pgfusepath{stroke}%
\end{pgfscope}%
\begin{pgfscope}%
\pgfpathrectangle{\pgfqpoint{0.100000in}{2.493063in}}{\pgfqpoint{5.037500in}{3.427208in}}%
\pgfusepath{clip}%
\pgfsetbuttcap%
\pgfsetroundjoin%
\definecolor{currentfill}{rgb}{0.000000,0.000000,1.000000}%
\pgfsetfillcolor{currentfill}%
\pgfsetfillopacity{0.500000}%
\pgfsetlinewidth{0.250937pt}%
\definecolor{currentstroke}{rgb}{0.000000,0.000000,0.000000}%
\pgfsetstrokecolor{currentstroke}%
\pgfsetstrokeopacity{0.500000}%
\pgfsetdash{}{0pt}%
\pgfsys@defobject{currentmarker}{\pgfqpoint{-0.025000in}{-0.025000in}}{\pgfqpoint{0.025000in}{0.025000in}}{%
\pgfpathmoveto{\pgfqpoint{0.000000in}{-0.025000in}}%
\pgfpathcurveto{\pgfqpoint{0.006630in}{-0.025000in}}{\pgfqpoint{0.012989in}{-0.022366in}}{\pgfqpoint{0.017678in}{-0.017678in}}%
\pgfpathcurveto{\pgfqpoint{0.022366in}{-0.012989in}}{\pgfqpoint{0.025000in}{-0.006630in}}{\pgfqpoint{0.025000in}{0.000000in}}%
\pgfpathcurveto{\pgfqpoint{0.025000in}{0.006630in}}{\pgfqpoint{0.022366in}{0.012989in}}{\pgfqpoint{0.017678in}{0.017678in}}%
\pgfpathcurveto{\pgfqpoint{0.012989in}{0.022366in}}{\pgfqpoint{0.006630in}{0.025000in}}{\pgfqpoint{0.000000in}{0.025000in}}%
\pgfpathcurveto{\pgfqpoint{-0.006630in}{0.025000in}}{\pgfqpoint{-0.012989in}{0.022366in}}{\pgfqpoint{-0.017678in}{0.017678in}}%
\pgfpathcurveto{\pgfqpoint{-0.022366in}{0.012989in}}{\pgfqpoint{-0.025000in}{0.006630in}}{\pgfqpoint{-0.025000in}{0.000000in}}%
\pgfpathcurveto{\pgfqpoint{-0.025000in}{-0.006630in}}{\pgfqpoint{-0.022366in}{-0.012989in}}{\pgfqpoint{-0.017678in}{-0.017678in}}%
\pgfpathcurveto{\pgfqpoint{-0.012989in}{-0.022366in}}{\pgfqpoint{-0.006630in}{-0.025000in}}{\pgfqpoint{0.000000in}{-0.025000in}}%
\pgfpathlineto{\pgfqpoint{0.000000in}{-0.025000in}}%
\pgfpathclose%
\pgfusepath{stroke,fill}%
}%
\begin{pgfscope}%
\pgfsys@transformshift{3.521226in}{4.250926in}%
\pgfsys@useobject{currentmarker}{}%
\end{pgfscope}%
\end{pgfscope}%
\begin{pgfscope}%
\pgfpathrectangle{\pgfqpoint{0.100000in}{2.493063in}}{\pgfqpoint{5.037500in}{3.427208in}}%
\pgfusepath{clip}%
\pgfsetrectcap%
\pgfsetroundjoin%
\pgfsetlinewidth{1.505625pt}%
\definecolor{currentstroke}{rgb}{0.678431,1.000000,0.184314}%
\pgfsetstrokecolor{currentstroke}%
\pgfsetstrokeopacity{0.500000}%
\pgfsetdash{}{0pt}%
\pgfpathmoveto{\pgfqpoint{3.058052in}{3.463452in}}%
\pgfusepath{stroke}%
\end{pgfscope}%
\begin{pgfscope}%
\pgfpathrectangle{\pgfqpoint{0.100000in}{2.493063in}}{\pgfqpoint{5.037500in}{3.427208in}}%
\pgfusepath{clip}%
\pgfsetbuttcap%
\pgfsetroundjoin%
\definecolor{currentfill}{rgb}{0.678431,1.000000,0.184314}%
\pgfsetfillcolor{currentfill}%
\pgfsetfillopacity{0.500000}%
\pgfsetlinewidth{0.250937pt}%
\definecolor{currentstroke}{rgb}{0.000000,0.000000,0.000000}%
\pgfsetstrokecolor{currentstroke}%
\pgfsetstrokeopacity{0.500000}%
\pgfsetdash{}{0pt}%
\pgfsys@defobject{currentmarker}{\pgfqpoint{-0.041667in}{-0.041667in}}{\pgfqpoint{0.041667in}{0.041667in}}{%
\pgfpathmoveto{\pgfqpoint{0.000000in}{-0.041667in}}%
\pgfpathcurveto{\pgfqpoint{0.011050in}{-0.041667in}}{\pgfqpoint{0.021649in}{-0.037276in}}{\pgfqpoint{0.029463in}{-0.029463in}}%
\pgfpathcurveto{\pgfqpoint{0.037276in}{-0.021649in}}{\pgfqpoint{0.041667in}{-0.011050in}}{\pgfqpoint{0.041667in}{0.000000in}}%
\pgfpathcurveto{\pgfqpoint{0.041667in}{0.011050in}}{\pgfqpoint{0.037276in}{0.021649in}}{\pgfqpoint{0.029463in}{0.029463in}}%
\pgfpathcurveto{\pgfqpoint{0.021649in}{0.037276in}}{\pgfqpoint{0.011050in}{0.041667in}}{\pgfqpoint{0.000000in}{0.041667in}}%
\pgfpathcurveto{\pgfqpoint{-0.011050in}{0.041667in}}{\pgfqpoint{-0.021649in}{0.037276in}}{\pgfqpoint{-0.029463in}{0.029463in}}%
\pgfpathcurveto{\pgfqpoint{-0.037276in}{0.021649in}}{\pgfqpoint{-0.041667in}{0.011050in}}{\pgfqpoint{-0.041667in}{0.000000in}}%
\pgfpathcurveto{\pgfqpoint{-0.041667in}{-0.011050in}}{\pgfqpoint{-0.037276in}{-0.021649in}}{\pgfqpoint{-0.029463in}{-0.029463in}}%
\pgfpathcurveto{\pgfqpoint{-0.021649in}{-0.037276in}}{\pgfqpoint{-0.011050in}{-0.041667in}}{\pgfqpoint{0.000000in}{-0.041667in}}%
\pgfpathlineto{\pgfqpoint{0.000000in}{-0.041667in}}%
\pgfpathclose%
\pgfusepath{stroke,fill}%
}%
\begin{pgfscope}%
\pgfsys@transformshift{3.058052in}{3.463452in}%
\pgfsys@useobject{currentmarker}{}%
\end{pgfscope}%
\end{pgfscope}%
\begin{pgfscope}%
\pgfpathrectangle{\pgfqpoint{0.100000in}{2.493063in}}{\pgfqpoint{5.037500in}{3.427208in}}%
\pgfusepath{clip}%
\pgfsetrectcap%
\pgfsetroundjoin%
\pgfsetlinewidth{1.505625pt}%
\definecolor{currentstroke}{rgb}{0.678431,1.000000,0.184314}%
\pgfsetstrokecolor{currentstroke}%
\pgfsetstrokeopacity{0.500000}%
\pgfsetdash{}{0pt}%
\pgfpathmoveto{\pgfqpoint{3.189939in}{3.381015in}}%
\pgfusepath{stroke}%
\end{pgfscope}%
\begin{pgfscope}%
\pgfpathrectangle{\pgfqpoint{0.100000in}{2.493063in}}{\pgfqpoint{5.037500in}{3.427208in}}%
\pgfusepath{clip}%
\pgfsetbuttcap%
\pgfsetroundjoin%
\definecolor{currentfill}{rgb}{0.678431,1.000000,0.184314}%
\pgfsetfillcolor{currentfill}%
\pgfsetfillopacity{0.500000}%
\pgfsetlinewidth{0.250937pt}%
\definecolor{currentstroke}{rgb}{0.000000,0.000000,0.000000}%
\pgfsetstrokecolor{currentstroke}%
\pgfsetstrokeopacity{0.500000}%
\pgfsetdash{}{0pt}%
\pgfsys@defobject{currentmarker}{\pgfqpoint{-0.030556in}{-0.030556in}}{\pgfqpoint{0.030556in}{0.030556in}}{%
\pgfpathmoveto{\pgfqpoint{0.000000in}{-0.030556in}}%
\pgfpathcurveto{\pgfqpoint{0.008103in}{-0.030556in}}{\pgfqpoint{0.015876in}{-0.027336in}}{\pgfqpoint{0.021606in}{-0.021606in}}%
\pgfpathcurveto{\pgfqpoint{0.027336in}{-0.015876in}}{\pgfqpoint{0.030556in}{-0.008103in}}{\pgfqpoint{0.030556in}{0.000000in}}%
\pgfpathcurveto{\pgfqpoint{0.030556in}{0.008103in}}{\pgfqpoint{0.027336in}{0.015876in}}{\pgfqpoint{0.021606in}{0.021606in}}%
\pgfpathcurveto{\pgfqpoint{0.015876in}{0.027336in}}{\pgfqpoint{0.008103in}{0.030556in}}{\pgfqpoint{0.000000in}{0.030556in}}%
\pgfpathcurveto{\pgfqpoint{-0.008103in}{0.030556in}}{\pgfqpoint{-0.015876in}{0.027336in}}{\pgfqpoint{-0.021606in}{0.021606in}}%
\pgfpathcurveto{\pgfqpoint{-0.027336in}{0.015876in}}{\pgfqpoint{-0.030556in}{0.008103in}}{\pgfqpoint{-0.030556in}{0.000000in}}%
\pgfpathcurveto{\pgfqpoint{-0.030556in}{-0.008103in}}{\pgfqpoint{-0.027336in}{-0.015876in}}{\pgfqpoint{-0.021606in}{-0.021606in}}%
\pgfpathcurveto{\pgfqpoint{-0.015876in}{-0.027336in}}{\pgfqpoint{-0.008103in}{-0.030556in}}{\pgfqpoint{0.000000in}{-0.030556in}}%
\pgfpathlineto{\pgfqpoint{0.000000in}{-0.030556in}}%
\pgfpathclose%
\pgfusepath{stroke,fill}%
}%
\begin{pgfscope}%
\pgfsys@transformshift{3.189939in}{3.381015in}%
\pgfsys@useobject{currentmarker}{}%
\end{pgfscope}%
\end{pgfscope}%
\begin{pgfscope}%
\pgfpathrectangle{\pgfqpoint{0.100000in}{2.493063in}}{\pgfqpoint{5.037500in}{3.427208in}}%
\pgfusepath{clip}%
\pgfsetrectcap%
\pgfsetroundjoin%
\pgfsetlinewidth{1.505625pt}%
\definecolor{currentstroke}{rgb}{0.678431,1.000000,0.184314}%
\pgfsetstrokecolor{currentstroke}%
\pgfsetstrokeopacity{0.500000}%
\pgfsetdash{}{0pt}%
\pgfpathmoveto{\pgfqpoint{3.449980in}{4.685174in}}%
\pgfusepath{stroke}%
\end{pgfscope}%
\begin{pgfscope}%
\pgfpathrectangle{\pgfqpoint{0.100000in}{2.493063in}}{\pgfqpoint{5.037500in}{3.427208in}}%
\pgfusepath{clip}%
\pgfsetbuttcap%
\pgfsetroundjoin%
\definecolor{currentfill}{rgb}{0.678431,1.000000,0.184314}%
\pgfsetfillcolor{currentfill}%
\pgfsetfillopacity{0.500000}%
\pgfsetlinewidth{0.250937pt}%
\definecolor{currentstroke}{rgb}{0.000000,0.000000,0.000000}%
\pgfsetstrokecolor{currentstroke}%
\pgfsetstrokeopacity{0.500000}%
\pgfsetdash{}{0pt}%
\pgfsys@defobject{currentmarker}{\pgfqpoint{-0.030556in}{-0.030556in}}{\pgfqpoint{0.030556in}{0.030556in}}{%
\pgfpathmoveto{\pgfqpoint{0.000000in}{-0.030556in}}%
\pgfpathcurveto{\pgfqpoint{0.008103in}{-0.030556in}}{\pgfqpoint{0.015876in}{-0.027336in}}{\pgfqpoint{0.021606in}{-0.021606in}}%
\pgfpathcurveto{\pgfqpoint{0.027336in}{-0.015876in}}{\pgfqpoint{0.030556in}{-0.008103in}}{\pgfqpoint{0.030556in}{0.000000in}}%
\pgfpathcurveto{\pgfqpoint{0.030556in}{0.008103in}}{\pgfqpoint{0.027336in}{0.015876in}}{\pgfqpoint{0.021606in}{0.021606in}}%
\pgfpathcurveto{\pgfqpoint{0.015876in}{0.027336in}}{\pgfqpoint{0.008103in}{0.030556in}}{\pgfqpoint{0.000000in}{0.030556in}}%
\pgfpathcurveto{\pgfqpoint{-0.008103in}{0.030556in}}{\pgfqpoint{-0.015876in}{0.027336in}}{\pgfqpoint{-0.021606in}{0.021606in}}%
\pgfpathcurveto{\pgfqpoint{-0.027336in}{0.015876in}}{\pgfqpoint{-0.030556in}{0.008103in}}{\pgfqpoint{-0.030556in}{0.000000in}}%
\pgfpathcurveto{\pgfqpoint{-0.030556in}{-0.008103in}}{\pgfqpoint{-0.027336in}{-0.015876in}}{\pgfqpoint{-0.021606in}{-0.021606in}}%
\pgfpathcurveto{\pgfqpoint{-0.015876in}{-0.027336in}}{\pgfqpoint{-0.008103in}{-0.030556in}}{\pgfqpoint{0.000000in}{-0.030556in}}%
\pgfpathlineto{\pgfqpoint{0.000000in}{-0.030556in}}%
\pgfpathclose%
\pgfusepath{stroke,fill}%
}%
\begin{pgfscope}%
\pgfsys@transformshift{3.449980in}{4.685174in}%
\pgfsys@useobject{currentmarker}{}%
\end{pgfscope}%
\end{pgfscope}%
\begin{pgfscope}%
\pgfpathrectangle{\pgfqpoint{0.100000in}{2.493063in}}{\pgfqpoint{5.037500in}{3.427208in}}%
\pgfusepath{clip}%
\pgfsetrectcap%
\pgfsetroundjoin%
\pgfsetlinewidth{1.505625pt}%
\definecolor{currentstroke}{rgb}{0.678431,1.000000,0.184314}%
\pgfsetstrokecolor{currentstroke}%
\pgfsetstrokeopacity{0.500000}%
\pgfsetdash{}{0pt}%
\pgfpathmoveto{\pgfqpoint{3.238245in}{3.283433in}}%
\pgfusepath{stroke}%
\end{pgfscope}%
\begin{pgfscope}%
\pgfpathrectangle{\pgfqpoint{0.100000in}{2.493063in}}{\pgfqpoint{5.037500in}{3.427208in}}%
\pgfusepath{clip}%
\pgfsetbuttcap%
\pgfsetroundjoin%
\definecolor{currentfill}{rgb}{0.678431,1.000000,0.184314}%
\pgfsetfillcolor{currentfill}%
\pgfsetfillopacity{0.500000}%
\pgfsetlinewidth{0.250937pt}%
\definecolor{currentstroke}{rgb}{0.000000,0.000000,0.000000}%
\pgfsetstrokecolor{currentstroke}%
\pgfsetstrokeopacity{0.500000}%
\pgfsetdash{}{0pt}%
\pgfsys@defobject{currentmarker}{\pgfqpoint{-0.030556in}{-0.030556in}}{\pgfqpoint{0.030556in}{0.030556in}}{%
\pgfpathmoveto{\pgfqpoint{0.000000in}{-0.030556in}}%
\pgfpathcurveto{\pgfqpoint{0.008103in}{-0.030556in}}{\pgfqpoint{0.015876in}{-0.027336in}}{\pgfqpoint{0.021606in}{-0.021606in}}%
\pgfpathcurveto{\pgfqpoint{0.027336in}{-0.015876in}}{\pgfqpoint{0.030556in}{-0.008103in}}{\pgfqpoint{0.030556in}{0.000000in}}%
\pgfpathcurveto{\pgfqpoint{0.030556in}{0.008103in}}{\pgfqpoint{0.027336in}{0.015876in}}{\pgfqpoint{0.021606in}{0.021606in}}%
\pgfpathcurveto{\pgfqpoint{0.015876in}{0.027336in}}{\pgfqpoint{0.008103in}{0.030556in}}{\pgfqpoint{0.000000in}{0.030556in}}%
\pgfpathcurveto{\pgfqpoint{-0.008103in}{0.030556in}}{\pgfqpoint{-0.015876in}{0.027336in}}{\pgfqpoint{-0.021606in}{0.021606in}}%
\pgfpathcurveto{\pgfqpoint{-0.027336in}{0.015876in}}{\pgfqpoint{-0.030556in}{0.008103in}}{\pgfqpoint{-0.030556in}{0.000000in}}%
\pgfpathcurveto{\pgfqpoint{-0.030556in}{-0.008103in}}{\pgfqpoint{-0.027336in}{-0.015876in}}{\pgfqpoint{-0.021606in}{-0.021606in}}%
\pgfpathcurveto{\pgfqpoint{-0.015876in}{-0.027336in}}{\pgfqpoint{-0.008103in}{-0.030556in}}{\pgfqpoint{0.000000in}{-0.030556in}}%
\pgfpathlineto{\pgfqpoint{0.000000in}{-0.030556in}}%
\pgfpathclose%
\pgfusepath{stroke,fill}%
}%
\begin{pgfscope}%
\pgfsys@transformshift{3.238245in}{3.283433in}%
\pgfsys@useobject{currentmarker}{}%
\end{pgfscope}%
\end{pgfscope}%
\begin{pgfscope}%
\pgfpathrectangle{\pgfqpoint{0.100000in}{2.493063in}}{\pgfqpoint{5.037500in}{3.427208in}}%
\pgfusepath{clip}%
\pgfsetrectcap%
\pgfsetroundjoin%
\pgfsetlinewidth{1.505625pt}%
\definecolor{currentstroke}{rgb}{0.678431,1.000000,0.184314}%
\pgfsetstrokecolor{currentstroke}%
\pgfsetstrokeopacity{0.500000}%
\pgfsetdash{}{0pt}%
\pgfpathmoveto{\pgfqpoint{3.104026in}{3.351375in}}%
\pgfusepath{stroke}%
\end{pgfscope}%
\begin{pgfscope}%
\pgfpathrectangle{\pgfqpoint{0.100000in}{2.493063in}}{\pgfqpoint{5.037500in}{3.427208in}}%
\pgfusepath{clip}%
\pgfsetbuttcap%
\pgfsetroundjoin%
\definecolor{currentfill}{rgb}{0.678431,1.000000,0.184314}%
\pgfsetfillcolor{currentfill}%
\pgfsetfillopacity{0.500000}%
\pgfsetlinewidth{0.250937pt}%
\definecolor{currentstroke}{rgb}{0.000000,0.000000,0.000000}%
\pgfsetstrokecolor{currentstroke}%
\pgfsetstrokeopacity{0.500000}%
\pgfsetdash{}{0pt}%
\pgfsys@defobject{currentmarker}{\pgfqpoint{-0.033333in}{-0.033333in}}{\pgfqpoint{0.033333in}{0.033333in}}{%
\pgfpathmoveto{\pgfqpoint{0.000000in}{-0.033333in}}%
\pgfpathcurveto{\pgfqpoint{0.008840in}{-0.033333in}}{\pgfqpoint{0.017319in}{-0.029821in}}{\pgfqpoint{0.023570in}{-0.023570in}}%
\pgfpathcurveto{\pgfqpoint{0.029821in}{-0.017319in}}{\pgfqpoint{0.033333in}{-0.008840in}}{\pgfqpoint{0.033333in}{0.000000in}}%
\pgfpathcurveto{\pgfqpoint{0.033333in}{0.008840in}}{\pgfqpoint{0.029821in}{0.017319in}}{\pgfqpoint{0.023570in}{0.023570in}}%
\pgfpathcurveto{\pgfqpoint{0.017319in}{0.029821in}}{\pgfqpoint{0.008840in}{0.033333in}}{\pgfqpoint{0.000000in}{0.033333in}}%
\pgfpathcurveto{\pgfqpoint{-0.008840in}{0.033333in}}{\pgfqpoint{-0.017319in}{0.029821in}}{\pgfqpoint{-0.023570in}{0.023570in}}%
\pgfpathcurveto{\pgfqpoint{-0.029821in}{0.017319in}}{\pgfqpoint{-0.033333in}{0.008840in}}{\pgfqpoint{-0.033333in}{0.000000in}}%
\pgfpathcurveto{\pgfqpoint{-0.033333in}{-0.008840in}}{\pgfqpoint{-0.029821in}{-0.017319in}}{\pgfqpoint{-0.023570in}{-0.023570in}}%
\pgfpathcurveto{\pgfqpoint{-0.017319in}{-0.029821in}}{\pgfqpoint{-0.008840in}{-0.033333in}}{\pgfqpoint{0.000000in}{-0.033333in}}%
\pgfpathlineto{\pgfqpoint{0.000000in}{-0.033333in}}%
\pgfpathclose%
\pgfusepath{stroke,fill}%
}%
\begin{pgfscope}%
\pgfsys@transformshift{3.104026in}{3.351375in}%
\pgfsys@useobject{currentmarker}{}%
\end{pgfscope}%
\end{pgfscope}%
\begin{pgfscope}%
\pgfpathrectangle{\pgfqpoint{0.100000in}{2.493063in}}{\pgfqpoint{5.037500in}{3.427208in}}%
\pgfusepath{clip}%
\pgfsetrectcap%
\pgfsetroundjoin%
\pgfsetlinewidth{1.505625pt}%
\definecolor{currentstroke}{rgb}{0.678431,1.000000,0.184314}%
\pgfsetstrokecolor{currentstroke}%
\pgfsetstrokeopacity{0.500000}%
\pgfsetdash{}{0pt}%
\pgfpathmoveto{\pgfqpoint{2.983531in}{3.348508in}}%
\pgfusepath{stroke}%
\end{pgfscope}%
\begin{pgfscope}%
\pgfpathrectangle{\pgfqpoint{0.100000in}{2.493063in}}{\pgfqpoint{5.037500in}{3.427208in}}%
\pgfusepath{clip}%
\pgfsetbuttcap%
\pgfsetroundjoin%
\definecolor{currentfill}{rgb}{0.678431,1.000000,0.184314}%
\pgfsetfillcolor{currentfill}%
\pgfsetfillopacity{0.500000}%
\pgfsetlinewidth{0.250937pt}%
\definecolor{currentstroke}{rgb}{0.000000,0.000000,0.000000}%
\pgfsetstrokecolor{currentstroke}%
\pgfsetstrokeopacity{0.500000}%
\pgfsetdash{}{0pt}%
\pgfsys@defobject{currentmarker}{\pgfqpoint{-0.022222in}{-0.022222in}}{\pgfqpoint{0.022222in}{0.022222in}}{%
\pgfpathmoveto{\pgfqpoint{0.000000in}{-0.022222in}}%
\pgfpathcurveto{\pgfqpoint{0.005893in}{-0.022222in}}{\pgfqpoint{0.011546in}{-0.019881in}}{\pgfqpoint{0.015713in}{-0.015713in}}%
\pgfpathcurveto{\pgfqpoint{0.019881in}{-0.011546in}}{\pgfqpoint{0.022222in}{-0.005893in}}{\pgfqpoint{0.022222in}{0.000000in}}%
\pgfpathcurveto{\pgfqpoint{0.022222in}{0.005893in}}{\pgfqpoint{0.019881in}{0.011546in}}{\pgfqpoint{0.015713in}{0.015713in}}%
\pgfpathcurveto{\pgfqpoint{0.011546in}{0.019881in}}{\pgfqpoint{0.005893in}{0.022222in}}{\pgfqpoint{0.000000in}{0.022222in}}%
\pgfpathcurveto{\pgfqpoint{-0.005893in}{0.022222in}}{\pgfqpoint{-0.011546in}{0.019881in}}{\pgfqpoint{-0.015713in}{0.015713in}}%
\pgfpathcurveto{\pgfqpoint{-0.019881in}{0.011546in}}{\pgfqpoint{-0.022222in}{0.005893in}}{\pgfqpoint{-0.022222in}{0.000000in}}%
\pgfpathcurveto{\pgfqpoint{-0.022222in}{-0.005893in}}{\pgfqpoint{-0.019881in}{-0.011546in}}{\pgfqpoint{-0.015713in}{-0.015713in}}%
\pgfpathcurveto{\pgfqpoint{-0.011546in}{-0.019881in}}{\pgfqpoint{-0.005893in}{-0.022222in}}{\pgfqpoint{0.000000in}{-0.022222in}}%
\pgfpathlineto{\pgfqpoint{0.000000in}{-0.022222in}}%
\pgfpathclose%
\pgfusepath{stroke,fill}%
}%
\begin{pgfscope}%
\pgfsys@transformshift{2.983531in}{3.348508in}%
\pgfsys@useobject{currentmarker}{}%
\end{pgfscope}%
\end{pgfscope}%
\begin{pgfscope}%
\pgfpathrectangle{\pgfqpoint{0.100000in}{2.493063in}}{\pgfqpoint{5.037500in}{3.427208in}}%
\pgfusepath{clip}%
\pgfsetrectcap%
\pgfsetroundjoin%
\pgfsetlinewidth{1.505625pt}%
\definecolor{currentstroke}{rgb}{0.678431,1.000000,0.184314}%
\pgfsetstrokecolor{currentstroke}%
\pgfsetstrokeopacity{0.500000}%
\pgfsetdash{}{0pt}%
\pgfpathmoveto{\pgfqpoint{3.239760in}{4.288703in}}%
\pgfusepath{stroke}%
\end{pgfscope}%
\begin{pgfscope}%
\pgfpathrectangle{\pgfqpoint{0.100000in}{2.493063in}}{\pgfqpoint{5.037500in}{3.427208in}}%
\pgfusepath{clip}%
\pgfsetbuttcap%
\pgfsetroundjoin%
\definecolor{currentfill}{rgb}{0.678431,1.000000,0.184314}%
\pgfsetfillcolor{currentfill}%
\pgfsetfillopacity{0.500000}%
\pgfsetlinewidth{0.250937pt}%
\definecolor{currentstroke}{rgb}{0.000000,0.000000,0.000000}%
\pgfsetstrokecolor{currentstroke}%
\pgfsetstrokeopacity{0.500000}%
\pgfsetdash{}{0pt}%
\pgfsys@defobject{currentmarker}{\pgfqpoint{-0.041667in}{-0.041667in}}{\pgfqpoint{0.041667in}{0.041667in}}{%
\pgfpathmoveto{\pgfqpoint{0.000000in}{-0.041667in}}%
\pgfpathcurveto{\pgfqpoint{0.011050in}{-0.041667in}}{\pgfqpoint{0.021649in}{-0.037276in}}{\pgfqpoint{0.029463in}{-0.029463in}}%
\pgfpathcurveto{\pgfqpoint{0.037276in}{-0.021649in}}{\pgfqpoint{0.041667in}{-0.011050in}}{\pgfqpoint{0.041667in}{0.000000in}}%
\pgfpathcurveto{\pgfqpoint{0.041667in}{0.011050in}}{\pgfqpoint{0.037276in}{0.021649in}}{\pgfqpoint{0.029463in}{0.029463in}}%
\pgfpathcurveto{\pgfqpoint{0.021649in}{0.037276in}}{\pgfqpoint{0.011050in}{0.041667in}}{\pgfqpoint{0.000000in}{0.041667in}}%
\pgfpathcurveto{\pgfqpoint{-0.011050in}{0.041667in}}{\pgfqpoint{-0.021649in}{0.037276in}}{\pgfqpoint{-0.029463in}{0.029463in}}%
\pgfpathcurveto{\pgfqpoint{-0.037276in}{0.021649in}}{\pgfqpoint{-0.041667in}{0.011050in}}{\pgfqpoint{-0.041667in}{0.000000in}}%
\pgfpathcurveto{\pgfqpoint{-0.041667in}{-0.011050in}}{\pgfqpoint{-0.037276in}{-0.021649in}}{\pgfqpoint{-0.029463in}{-0.029463in}}%
\pgfpathcurveto{\pgfqpoint{-0.021649in}{-0.037276in}}{\pgfqpoint{-0.011050in}{-0.041667in}}{\pgfqpoint{0.000000in}{-0.041667in}}%
\pgfpathlineto{\pgfqpoint{0.000000in}{-0.041667in}}%
\pgfpathclose%
\pgfusepath{stroke,fill}%
}%
\begin{pgfscope}%
\pgfsys@transformshift{3.239760in}{4.288703in}%
\pgfsys@useobject{currentmarker}{}%
\end{pgfscope}%
\end{pgfscope}%
\begin{pgfscope}%
\pgfpathrectangle{\pgfqpoint{0.100000in}{2.493063in}}{\pgfqpoint{5.037500in}{3.427208in}}%
\pgfusepath{clip}%
\pgfsetrectcap%
\pgfsetroundjoin%
\pgfsetlinewidth{1.505625pt}%
\definecolor{currentstroke}{rgb}{0.678431,1.000000,0.184314}%
\pgfsetstrokecolor{currentstroke}%
\pgfsetstrokeopacity{0.500000}%
\pgfsetdash{}{0pt}%
\pgfpathmoveto{\pgfqpoint{3.301808in}{3.328013in}}%
\pgfusepath{stroke}%
\end{pgfscope}%
\begin{pgfscope}%
\pgfpathrectangle{\pgfqpoint{0.100000in}{2.493063in}}{\pgfqpoint{5.037500in}{3.427208in}}%
\pgfusepath{clip}%
\pgfsetbuttcap%
\pgfsetroundjoin%
\definecolor{currentfill}{rgb}{0.678431,1.000000,0.184314}%
\pgfsetfillcolor{currentfill}%
\pgfsetfillopacity{0.500000}%
\pgfsetlinewidth{0.250937pt}%
\definecolor{currentstroke}{rgb}{0.000000,0.000000,0.000000}%
\pgfsetstrokecolor{currentstroke}%
\pgfsetstrokeopacity{0.500000}%
\pgfsetdash{}{0pt}%
\pgfsys@defobject{currentmarker}{\pgfqpoint{-0.016667in}{-0.016667in}}{\pgfqpoint{0.016667in}{0.016667in}}{%
\pgfpathmoveto{\pgfqpoint{0.000000in}{-0.016667in}}%
\pgfpathcurveto{\pgfqpoint{0.004420in}{-0.016667in}}{\pgfqpoint{0.008660in}{-0.014911in}}{\pgfqpoint{0.011785in}{-0.011785in}}%
\pgfpathcurveto{\pgfqpoint{0.014911in}{-0.008660in}}{\pgfqpoint{0.016667in}{-0.004420in}}{\pgfqpoint{0.016667in}{0.000000in}}%
\pgfpathcurveto{\pgfqpoint{0.016667in}{0.004420in}}{\pgfqpoint{0.014911in}{0.008660in}}{\pgfqpoint{0.011785in}{0.011785in}}%
\pgfpathcurveto{\pgfqpoint{0.008660in}{0.014911in}}{\pgfqpoint{0.004420in}{0.016667in}}{\pgfqpoint{0.000000in}{0.016667in}}%
\pgfpathcurveto{\pgfqpoint{-0.004420in}{0.016667in}}{\pgfqpoint{-0.008660in}{0.014911in}}{\pgfqpoint{-0.011785in}{0.011785in}}%
\pgfpathcurveto{\pgfqpoint{-0.014911in}{0.008660in}}{\pgfqpoint{-0.016667in}{0.004420in}}{\pgfqpoint{-0.016667in}{0.000000in}}%
\pgfpathcurveto{\pgfqpoint{-0.016667in}{-0.004420in}}{\pgfqpoint{-0.014911in}{-0.008660in}}{\pgfqpoint{-0.011785in}{-0.011785in}}%
\pgfpathcurveto{\pgfqpoint{-0.008660in}{-0.014911in}}{\pgfqpoint{-0.004420in}{-0.016667in}}{\pgfqpoint{0.000000in}{-0.016667in}}%
\pgfpathlineto{\pgfqpoint{0.000000in}{-0.016667in}}%
\pgfpathclose%
\pgfusepath{stroke,fill}%
}%
\begin{pgfscope}%
\pgfsys@transformshift{3.301808in}{3.328013in}%
\pgfsys@useobject{currentmarker}{}%
\end{pgfscope}%
\end{pgfscope}%
\begin{pgfscope}%
\pgfpathrectangle{\pgfqpoint{0.100000in}{2.493063in}}{\pgfqpoint{5.037500in}{3.427208in}}%
\pgfusepath{clip}%
\pgfsetrectcap%
\pgfsetroundjoin%
\pgfsetlinewidth{1.505625pt}%
\definecolor{currentstroke}{rgb}{0.678431,1.000000,0.184314}%
\pgfsetstrokecolor{currentstroke}%
\pgfsetstrokeopacity{0.500000}%
\pgfsetdash{}{0pt}%
\pgfpathmoveto{\pgfqpoint{2.924780in}{3.614245in}}%
\pgfusepath{stroke}%
\end{pgfscope}%
\begin{pgfscope}%
\pgfpathrectangle{\pgfqpoint{0.100000in}{2.493063in}}{\pgfqpoint{5.037500in}{3.427208in}}%
\pgfusepath{clip}%
\pgfsetbuttcap%
\pgfsetroundjoin%
\definecolor{currentfill}{rgb}{0.678431,1.000000,0.184314}%
\pgfsetfillcolor{currentfill}%
\pgfsetfillopacity{0.500000}%
\pgfsetlinewidth{0.250937pt}%
\definecolor{currentstroke}{rgb}{0.000000,0.000000,0.000000}%
\pgfsetstrokecolor{currentstroke}%
\pgfsetstrokeopacity{0.500000}%
\pgfsetdash{}{0pt}%
\pgfsys@defobject{currentmarker}{\pgfqpoint{-0.033333in}{-0.033333in}}{\pgfqpoint{0.033333in}{0.033333in}}{%
\pgfpathmoveto{\pgfqpoint{0.000000in}{-0.033333in}}%
\pgfpathcurveto{\pgfqpoint{0.008840in}{-0.033333in}}{\pgfqpoint{0.017319in}{-0.029821in}}{\pgfqpoint{0.023570in}{-0.023570in}}%
\pgfpathcurveto{\pgfqpoint{0.029821in}{-0.017319in}}{\pgfqpoint{0.033333in}{-0.008840in}}{\pgfqpoint{0.033333in}{0.000000in}}%
\pgfpathcurveto{\pgfqpoint{0.033333in}{0.008840in}}{\pgfqpoint{0.029821in}{0.017319in}}{\pgfqpoint{0.023570in}{0.023570in}}%
\pgfpathcurveto{\pgfqpoint{0.017319in}{0.029821in}}{\pgfqpoint{0.008840in}{0.033333in}}{\pgfqpoint{0.000000in}{0.033333in}}%
\pgfpathcurveto{\pgfqpoint{-0.008840in}{0.033333in}}{\pgfqpoint{-0.017319in}{0.029821in}}{\pgfqpoint{-0.023570in}{0.023570in}}%
\pgfpathcurveto{\pgfqpoint{-0.029821in}{0.017319in}}{\pgfqpoint{-0.033333in}{0.008840in}}{\pgfqpoint{-0.033333in}{0.000000in}}%
\pgfpathcurveto{\pgfqpoint{-0.033333in}{-0.008840in}}{\pgfqpoint{-0.029821in}{-0.017319in}}{\pgfqpoint{-0.023570in}{-0.023570in}}%
\pgfpathcurveto{\pgfqpoint{-0.017319in}{-0.029821in}}{\pgfqpoint{-0.008840in}{-0.033333in}}{\pgfqpoint{0.000000in}{-0.033333in}}%
\pgfpathlineto{\pgfqpoint{0.000000in}{-0.033333in}}%
\pgfpathclose%
\pgfusepath{stroke,fill}%
}%
\begin{pgfscope}%
\pgfsys@transformshift{2.924780in}{3.614245in}%
\pgfsys@useobject{currentmarker}{}%
\end{pgfscope}%
\end{pgfscope}%
\begin{pgfscope}%
\pgfpathrectangle{\pgfqpoint{0.100000in}{2.493063in}}{\pgfqpoint{5.037500in}{3.427208in}}%
\pgfusepath{clip}%
\pgfsetrectcap%
\pgfsetroundjoin%
\pgfsetlinewidth{1.505625pt}%
\definecolor{currentstroke}{rgb}{0.000000,0.000000,1.000000}%
\pgfsetstrokecolor{currentstroke}%
\pgfsetstrokeopacity{0.500000}%
\pgfsetdash{}{0pt}%
\pgfpathmoveto{\pgfqpoint{4.927826in}{5.338593in}}%
\pgfusepath{stroke}%
\end{pgfscope}%
\begin{pgfscope}%
\pgfpathrectangle{\pgfqpoint{0.100000in}{2.493063in}}{\pgfqpoint{5.037500in}{3.427208in}}%
\pgfusepath{clip}%
\pgfsetbuttcap%
\pgfsetroundjoin%
\definecolor{currentfill}{rgb}{0.000000,0.000000,1.000000}%
\pgfsetfillcolor{currentfill}%
\pgfsetfillopacity{0.500000}%
\pgfsetlinewidth{0.250937pt}%
\definecolor{currentstroke}{rgb}{0.000000,0.000000,0.000000}%
\pgfsetstrokecolor{currentstroke}%
\pgfsetstrokeopacity{0.500000}%
\pgfsetdash{}{0pt}%
\pgfsys@defobject{currentmarker}{\pgfqpoint{-0.008333in}{-0.008333in}}{\pgfqpoint{0.008333in}{0.008333in}}{%
\pgfpathmoveto{\pgfqpoint{0.000000in}{-0.008333in}}%
\pgfpathcurveto{\pgfqpoint{0.002210in}{-0.008333in}}{\pgfqpoint{0.004330in}{-0.007455in}}{\pgfqpoint{0.005893in}{-0.005893in}}%
\pgfpathcurveto{\pgfqpoint{0.007455in}{-0.004330in}}{\pgfqpoint{0.008333in}{-0.002210in}}{\pgfqpoint{0.008333in}{0.000000in}}%
\pgfpathcurveto{\pgfqpoint{0.008333in}{0.002210in}}{\pgfqpoint{0.007455in}{0.004330in}}{\pgfqpoint{0.005893in}{0.005893in}}%
\pgfpathcurveto{\pgfqpoint{0.004330in}{0.007455in}}{\pgfqpoint{0.002210in}{0.008333in}}{\pgfqpoint{0.000000in}{0.008333in}}%
\pgfpathcurveto{\pgfqpoint{-0.002210in}{0.008333in}}{\pgfqpoint{-0.004330in}{0.007455in}}{\pgfqpoint{-0.005893in}{0.005893in}}%
\pgfpathcurveto{\pgfqpoint{-0.007455in}{0.004330in}}{\pgfqpoint{-0.008333in}{0.002210in}}{\pgfqpoint{-0.008333in}{0.000000in}}%
\pgfpathcurveto{\pgfqpoint{-0.008333in}{-0.002210in}}{\pgfqpoint{-0.007455in}{-0.004330in}}{\pgfqpoint{-0.005893in}{-0.005893in}}%
\pgfpathcurveto{\pgfqpoint{-0.004330in}{-0.007455in}}{\pgfqpoint{-0.002210in}{-0.008333in}}{\pgfqpoint{0.000000in}{-0.008333in}}%
\pgfpathlineto{\pgfqpoint{0.000000in}{-0.008333in}}%
\pgfpathclose%
\pgfusepath{stroke,fill}%
}%
\begin{pgfscope}%
\pgfsys@transformshift{4.927826in}{5.338593in}%
\pgfsys@useobject{currentmarker}{}%
\end{pgfscope}%
\end{pgfscope}%
\begin{pgfscope}%
\pgfpathrectangle{\pgfqpoint{0.100000in}{2.493063in}}{\pgfqpoint{5.037500in}{3.427208in}}%
\pgfusepath{clip}%
\pgfsetrectcap%
\pgfsetroundjoin%
\pgfsetlinewidth{1.505625pt}%
\definecolor{currentstroke}{rgb}{0.000000,0.000000,1.000000}%
\pgfsetstrokecolor{currentstroke}%
\pgfsetstrokeopacity{0.500000}%
\pgfsetdash{}{0pt}%
\pgfpathmoveto{\pgfqpoint{4.836271in}{5.227841in}}%
\pgfusepath{stroke}%
\end{pgfscope}%
\begin{pgfscope}%
\pgfpathrectangle{\pgfqpoint{0.100000in}{2.493063in}}{\pgfqpoint{5.037500in}{3.427208in}}%
\pgfusepath{clip}%
\pgfsetbuttcap%
\pgfsetroundjoin%
\definecolor{currentfill}{rgb}{0.000000,0.000000,1.000000}%
\pgfsetfillcolor{currentfill}%
\pgfsetfillopacity{0.500000}%
\pgfsetlinewidth{0.250937pt}%
\definecolor{currentstroke}{rgb}{0.000000,0.000000,0.000000}%
\pgfsetstrokecolor{currentstroke}%
\pgfsetstrokeopacity{0.500000}%
\pgfsetdash{}{0pt}%
\pgfsys@defobject{currentmarker}{\pgfqpoint{-0.011111in}{-0.011111in}}{\pgfqpoint{0.011111in}{0.011111in}}{%
\pgfpathmoveto{\pgfqpoint{0.000000in}{-0.011111in}}%
\pgfpathcurveto{\pgfqpoint{0.002947in}{-0.011111in}}{\pgfqpoint{0.005773in}{-0.009940in}}{\pgfqpoint{0.007857in}{-0.007857in}}%
\pgfpathcurveto{\pgfqpoint{0.009940in}{-0.005773in}}{\pgfqpoint{0.011111in}{-0.002947in}}{\pgfqpoint{0.011111in}{0.000000in}}%
\pgfpathcurveto{\pgfqpoint{0.011111in}{0.002947in}}{\pgfqpoint{0.009940in}{0.005773in}}{\pgfqpoint{0.007857in}{0.007857in}}%
\pgfpathcurveto{\pgfqpoint{0.005773in}{0.009940in}}{\pgfqpoint{0.002947in}{0.011111in}}{\pgfqpoint{0.000000in}{0.011111in}}%
\pgfpathcurveto{\pgfqpoint{-0.002947in}{0.011111in}}{\pgfqpoint{-0.005773in}{0.009940in}}{\pgfqpoint{-0.007857in}{0.007857in}}%
\pgfpathcurveto{\pgfqpoint{-0.009940in}{0.005773in}}{\pgfqpoint{-0.011111in}{0.002947in}}{\pgfqpoint{-0.011111in}{0.000000in}}%
\pgfpathcurveto{\pgfqpoint{-0.011111in}{-0.002947in}}{\pgfqpoint{-0.009940in}{-0.005773in}}{\pgfqpoint{-0.007857in}{-0.007857in}}%
\pgfpathcurveto{\pgfqpoint{-0.005773in}{-0.009940in}}{\pgfqpoint{-0.002947in}{-0.011111in}}{\pgfqpoint{0.000000in}{-0.011111in}}%
\pgfpathlineto{\pgfqpoint{0.000000in}{-0.011111in}}%
\pgfpathclose%
\pgfusepath{stroke,fill}%
}%
\begin{pgfscope}%
\pgfsys@transformshift{4.836271in}{5.227841in}%
\pgfsys@useobject{currentmarker}{}%
\end{pgfscope}%
\end{pgfscope}%
\begin{pgfscope}%
\pgfpathrectangle{\pgfqpoint{0.100000in}{2.493063in}}{\pgfqpoint{5.037500in}{3.427208in}}%
\pgfusepath{clip}%
\pgfsetrectcap%
\pgfsetroundjoin%
\pgfsetlinewidth{1.505625pt}%
\definecolor{currentstroke}{rgb}{0.000000,0.000000,1.000000}%
\pgfsetstrokecolor{currentstroke}%
\pgfsetstrokeopacity{0.500000}%
\pgfsetdash{}{0pt}%
\pgfpathmoveto{\pgfqpoint{4.846720in}{5.178047in}}%
\pgfusepath{stroke}%
\end{pgfscope}%
\begin{pgfscope}%
\pgfpathrectangle{\pgfqpoint{0.100000in}{2.493063in}}{\pgfqpoint{5.037500in}{3.427208in}}%
\pgfusepath{clip}%
\pgfsetbuttcap%
\pgfsetroundjoin%
\definecolor{currentfill}{rgb}{0.000000,0.000000,1.000000}%
\pgfsetfillcolor{currentfill}%
\pgfsetfillopacity{0.500000}%
\pgfsetlinewidth{0.250937pt}%
\definecolor{currentstroke}{rgb}{0.000000,0.000000,0.000000}%
\pgfsetstrokecolor{currentstroke}%
\pgfsetstrokeopacity{0.500000}%
\pgfsetdash{}{0pt}%
\pgfsys@defobject{currentmarker}{\pgfqpoint{-0.011111in}{-0.011111in}}{\pgfqpoint{0.011111in}{0.011111in}}{%
\pgfpathmoveto{\pgfqpoint{0.000000in}{-0.011111in}}%
\pgfpathcurveto{\pgfqpoint{0.002947in}{-0.011111in}}{\pgfqpoint{0.005773in}{-0.009940in}}{\pgfqpoint{0.007857in}{-0.007857in}}%
\pgfpathcurveto{\pgfqpoint{0.009940in}{-0.005773in}}{\pgfqpoint{0.011111in}{-0.002947in}}{\pgfqpoint{0.011111in}{0.000000in}}%
\pgfpathcurveto{\pgfqpoint{0.011111in}{0.002947in}}{\pgfqpoint{0.009940in}{0.005773in}}{\pgfqpoint{0.007857in}{0.007857in}}%
\pgfpathcurveto{\pgfqpoint{0.005773in}{0.009940in}}{\pgfqpoint{0.002947in}{0.011111in}}{\pgfqpoint{0.000000in}{0.011111in}}%
\pgfpathcurveto{\pgfqpoint{-0.002947in}{0.011111in}}{\pgfqpoint{-0.005773in}{0.009940in}}{\pgfqpoint{-0.007857in}{0.007857in}}%
\pgfpathcurveto{\pgfqpoint{-0.009940in}{0.005773in}}{\pgfqpoint{-0.011111in}{0.002947in}}{\pgfqpoint{-0.011111in}{0.000000in}}%
\pgfpathcurveto{\pgfqpoint{-0.011111in}{-0.002947in}}{\pgfqpoint{-0.009940in}{-0.005773in}}{\pgfqpoint{-0.007857in}{-0.007857in}}%
\pgfpathcurveto{\pgfqpoint{-0.005773in}{-0.009940in}}{\pgfqpoint{-0.002947in}{-0.011111in}}{\pgfqpoint{0.000000in}{-0.011111in}}%
\pgfpathlineto{\pgfqpoint{0.000000in}{-0.011111in}}%
\pgfpathclose%
\pgfusepath{stroke,fill}%
}%
\begin{pgfscope}%
\pgfsys@transformshift{4.846720in}{5.178047in}%
\pgfsys@useobject{currentmarker}{}%
\end{pgfscope}%
\end{pgfscope}%
\begin{pgfscope}%
\pgfpathrectangle{\pgfqpoint{0.100000in}{2.493063in}}{\pgfqpoint{5.037500in}{3.427208in}}%
\pgfusepath{clip}%
\pgfsetrectcap%
\pgfsetroundjoin%
\pgfsetlinewidth{1.505625pt}%
\definecolor{currentstroke}{rgb}{0.678431,1.000000,0.184314}%
\pgfsetstrokecolor{currentstroke}%
\pgfsetstrokeopacity{0.500000}%
\pgfsetdash{}{0pt}%
\pgfpathmoveto{\pgfqpoint{4.431773in}{4.559688in}}%
\pgfusepath{stroke}%
\end{pgfscope}%
\begin{pgfscope}%
\pgfpathrectangle{\pgfqpoint{0.100000in}{2.493063in}}{\pgfqpoint{5.037500in}{3.427208in}}%
\pgfusepath{clip}%
\pgfsetbuttcap%
\pgfsetroundjoin%
\definecolor{currentfill}{rgb}{0.678431,1.000000,0.184314}%
\pgfsetfillcolor{currentfill}%
\pgfsetfillopacity{0.500000}%
\pgfsetlinewidth{0.250937pt}%
\definecolor{currentstroke}{rgb}{0.000000,0.000000,0.000000}%
\pgfsetstrokecolor{currentstroke}%
\pgfsetstrokeopacity{0.500000}%
\pgfsetdash{}{0pt}%
\pgfsys@defobject{currentmarker}{\pgfqpoint{-0.033333in}{-0.033333in}}{\pgfqpoint{0.033333in}{0.033333in}}{%
\pgfpathmoveto{\pgfqpoint{0.000000in}{-0.033333in}}%
\pgfpathcurveto{\pgfqpoint{0.008840in}{-0.033333in}}{\pgfqpoint{0.017319in}{-0.029821in}}{\pgfqpoint{0.023570in}{-0.023570in}}%
\pgfpathcurveto{\pgfqpoint{0.029821in}{-0.017319in}}{\pgfqpoint{0.033333in}{-0.008840in}}{\pgfqpoint{0.033333in}{0.000000in}}%
\pgfpathcurveto{\pgfqpoint{0.033333in}{0.008840in}}{\pgfqpoint{0.029821in}{0.017319in}}{\pgfqpoint{0.023570in}{0.023570in}}%
\pgfpathcurveto{\pgfqpoint{0.017319in}{0.029821in}}{\pgfqpoint{0.008840in}{0.033333in}}{\pgfqpoint{0.000000in}{0.033333in}}%
\pgfpathcurveto{\pgfqpoint{-0.008840in}{0.033333in}}{\pgfqpoint{-0.017319in}{0.029821in}}{\pgfqpoint{-0.023570in}{0.023570in}}%
\pgfpathcurveto{\pgfqpoint{-0.029821in}{0.017319in}}{\pgfqpoint{-0.033333in}{0.008840in}}{\pgfqpoint{-0.033333in}{0.000000in}}%
\pgfpathcurveto{\pgfqpoint{-0.033333in}{-0.008840in}}{\pgfqpoint{-0.029821in}{-0.017319in}}{\pgfqpoint{-0.023570in}{-0.023570in}}%
\pgfpathcurveto{\pgfqpoint{-0.017319in}{-0.029821in}}{\pgfqpoint{-0.008840in}{-0.033333in}}{\pgfqpoint{0.000000in}{-0.033333in}}%
\pgfpathlineto{\pgfqpoint{0.000000in}{-0.033333in}}%
\pgfpathclose%
\pgfusepath{stroke,fill}%
}%
\begin{pgfscope}%
\pgfsys@transformshift{4.431773in}{4.559688in}%
\pgfsys@useobject{currentmarker}{}%
\end{pgfscope}%
\end{pgfscope}%
\begin{pgfscope}%
\pgfpathrectangle{\pgfqpoint{0.100000in}{2.493063in}}{\pgfqpoint{5.037500in}{3.427208in}}%
\pgfusepath{clip}%
\pgfsetrectcap%
\pgfsetroundjoin%
\pgfsetlinewidth{1.505625pt}%
\definecolor{currentstroke}{rgb}{0.678431,1.000000,0.184314}%
\pgfsetstrokecolor{currentstroke}%
\pgfsetstrokeopacity{0.500000}%
\pgfsetdash{}{0pt}%
\pgfpathmoveto{\pgfqpoint{4.463839in}{4.449843in}}%
\pgfusepath{stroke}%
\end{pgfscope}%
\begin{pgfscope}%
\pgfpathrectangle{\pgfqpoint{0.100000in}{2.493063in}}{\pgfqpoint{5.037500in}{3.427208in}}%
\pgfusepath{clip}%
\pgfsetbuttcap%
\pgfsetroundjoin%
\definecolor{currentfill}{rgb}{0.678431,1.000000,0.184314}%
\pgfsetfillcolor{currentfill}%
\pgfsetfillopacity{0.500000}%
\pgfsetlinewidth{0.250937pt}%
\definecolor{currentstroke}{rgb}{0.000000,0.000000,0.000000}%
\pgfsetstrokecolor{currentstroke}%
\pgfsetstrokeopacity{0.500000}%
\pgfsetdash{}{0pt}%
\pgfsys@defobject{currentmarker}{\pgfqpoint{-0.027778in}{-0.027778in}}{\pgfqpoint{0.027778in}{0.027778in}}{%
\pgfpathmoveto{\pgfqpoint{0.000000in}{-0.027778in}}%
\pgfpathcurveto{\pgfqpoint{0.007367in}{-0.027778in}}{\pgfqpoint{0.014433in}{-0.024851in}}{\pgfqpoint{0.019642in}{-0.019642in}}%
\pgfpathcurveto{\pgfqpoint{0.024851in}{-0.014433in}}{\pgfqpoint{0.027778in}{-0.007367in}}{\pgfqpoint{0.027778in}{0.000000in}}%
\pgfpathcurveto{\pgfqpoint{0.027778in}{0.007367in}}{\pgfqpoint{0.024851in}{0.014433in}}{\pgfqpoint{0.019642in}{0.019642in}}%
\pgfpathcurveto{\pgfqpoint{0.014433in}{0.024851in}}{\pgfqpoint{0.007367in}{0.027778in}}{\pgfqpoint{0.000000in}{0.027778in}}%
\pgfpathcurveto{\pgfqpoint{-0.007367in}{0.027778in}}{\pgfqpoint{-0.014433in}{0.024851in}}{\pgfqpoint{-0.019642in}{0.019642in}}%
\pgfpathcurveto{\pgfqpoint{-0.024851in}{0.014433in}}{\pgfqpoint{-0.027778in}{0.007367in}}{\pgfqpoint{-0.027778in}{0.000000in}}%
\pgfpathcurveto{\pgfqpoint{-0.027778in}{-0.007367in}}{\pgfqpoint{-0.024851in}{-0.014433in}}{\pgfqpoint{-0.019642in}{-0.019642in}}%
\pgfpathcurveto{\pgfqpoint{-0.014433in}{-0.024851in}}{\pgfqpoint{-0.007367in}{-0.027778in}}{\pgfqpoint{0.000000in}{-0.027778in}}%
\pgfpathlineto{\pgfqpoint{0.000000in}{-0.027778in}}%
\pgfpathclose%
\pgfusepath{stroke,fill}%
}%
\begin{pgfscope}%
\pgfsys@transformshift{4.463839in}{4.449843in}%
\pgfsys@useobject{currentmarker}{}%
\end{pgfscope}%
\end{pgfscope}%
\begin{pgfscope}%
\pgfpathrectangle{\pgfqpoint{0.100000in}{2.493063in}}{\pgfqpoint{5.037500in}{3.427208in}}%
\pgfusepath{clip}%
\pgfsetrectcap%
\pgfsetroundjoin%
\pgfsetlinewidth{1.505625pt}%
\definecolor{currentstroke}{rgb}{0.678431,1.000000,0.184314}%
\pgfsetstrokecolor{currentstroke}%
\pgfsetstrokeopacity{0.500000}%
\pgfsetdash{}{0pt}%
\pgfpathmoveto{\pgfqpoint{4.236084in}{4.564394in}}%
\pgfusepath{stroke}%
\end{pgfscope}%
\begin{pgfscope}%
\pgfpathrectangle{\pgfqpoint{0.100000in}{2.493063in}}{\pgfqpoint{5.037500in}{3.427208in}}%
\pgfusepath{clip}%
\pgfsetbuttcap%
\pgfsetroundjoin%
\definecolor{currentfill}{rgb}{0.678431,1.000000,0.184314}%
\pgfsetfillcolor{currentfill}%
\pgfsetfillopacity{0.500000}%
\pgfsetlinewidth{0.250937pt}%
\definecolor{currentstroke}{rgb}{0.000000,0.000000,0.000000}%
\pgfsetstrokecolor{currentstroke}%
\pgfsetstrokeopacity{0.500000}%
\pgfsetdash{}{0pt}%
\pgfsys@defobject{currentmarker}{\pgfqpoint{-0.055556in}{-0.055556in}}{\pgfqpoint{0.055556in}{0.055556in}}{%
\pgfpathmoveto{\pgfqpoint{0.000000in}{-0.055556in}}%
\pgfpathcurveto{\pgfqpoint{0.014734in}{-0.055556in}}{\pgfqpoint{0.028866in}{-0.049702in}}{\pgfqpoint{0.039284in}{-0.039284in}}%
\pgfpathcurveto{\pgfqpoint{0.049702in}{-0.028866in}}{\pgfqpoint{0.055556in}{-0.014734in}}{\pgfqpoint{0.055556in}{0.000000in}}%
\pgfpathcurveto{\pgfqpoint{0.055556in}{0.014734in}}{\pgfqpoint{0.049702in}{0.028866in}}{\pgfqpoint{0.039284in}{0.039284in}}%
\pgfpathcurveto{\pgfqpoint{0.028866in}{0.049702in}}{\pgfqpoint{0.014734in}{0.055556in}}{\pgfqpoint{0.000000in}{0.055556in}}%
\pgfpathcurveto{\pgfqpoint{-0.014734in}{0.055556in}}{\pgfqpoint{-0.028866in}{0.049702in}}{\pgfqpoint{-0.039284in}{0.039284in}}%
\pgfpathcurveto{\pgfqpoint{-0.049702in}{0.028866in}}{\pgfqpoint{-0.055556in}{0.014734in}}{\pgfqpoint{-0.055556in}{0.000000in}}%
\pgfpathcurveto{\pgfqpoint{-0.055556in}{-0.014734in}}{\pgfqpoint{-0.049702in}{-0.028866in}}{\pgfqpoint{-0.039284in}{-0.039284in}}%
\pgfpathcurveto{\pgfqpoint{-0.028866in}{-0.049702in}}{\pgfqpoint{-0.014734in}{-0.055556in}}{\pgfqpoint{0.000000in}{-0.055556in}}%
\pgfpathlineto{\pgfqpoint{0.000000in}{-0.055556in}}%
\pgfpathclose%
\pgfusepath{stroke,fill}%
}%
\begin{pgfscope}%
\pgfsys@transformshift{4.236084in}{4.564394in}%
\pgfsys@useobject{currentmarker}{}%
\end{pgfscope}%
\end{pgfscope}%
\begin{pgfscope}%
\pgfpathrectangle{\pgfqpoint{0.100000in}{2.493063in}}{\pgfqpoint{5.037500in}{3.427208in}}%
\pgfusepath{clip}%
\pgfsetrectcap%
\pgfsetroundjoin%
\pgfsetlinewidth{1.505625pt}%
\definecolor{currentstroke}{rgb}{0.678431,1.000000,0.184314}%
\pgfsetstrokecolor{currentstroke}%
\pgfsetstrokeopacity{0.500000}%
\pgfsetdash{}{0pt}%
\pgfpathmoveto{\pgfqpoint{4.327155in}{4.580101in}}%
\pgfusepath{stroke}%
\end{pgfscope}%
\begin{pgfscope}%
\pgfpathrectangle{\pgfqpoint{0.100000in}{2.493063in}}{\pgfqpoint{5.037500in}{3.427208in}}%
\pgfusepath{clip}%
\pgfsetbuttcap%
\pgfsetroundjoin%
\definecolor{currentfill}{rgb}{0.678431,1.000000,0.184314}%
\pgfsetfillcolor{currentfill}%
\pgfsetfillopacity{0.500000}%
\pgfsetlinewidth{0.250937pt}%
\definecolor{currentstroke}{rgb}{0.000000,0.000000,0.000000}%
\pgfsetstrokecolor{currentstroke}%
\pgfsetstrokeopacity{0.500000}%
\pgfsetdash{}{0pt}%
\pgfsys@defobject{currentmarker}{\pgfqpoint{-0.013889in}{-0.013889in}}{\pgfqpoint{0.013889in}{0.013889in}}{%
\pgfpathmoveto{\pgfqpoint{0.000000in}{-0.013889in}}%
\pgfpathcurveto{\pgfqpoint{0.003683in}{-0.013889in}}{\pgfqpoint{0.007216in}{-0.012425in}}{\pgfqpoint{0.009821in}{-0.009821in}}%
\pgfpathcurveto{\pgfqpoint{0.012425in}{-0.007216in}}{\pgfqpoint{0.013889in}{-0.003683in}}{\pgfqpoint{0.013889in}{0.000000in}}%
\pgfpathcurveto{\pgfqpoint{0.013889in}{0.003683in}}{\pgfqpoint{0.012425in}{0.007216in}}{\pgfqpoint{0.009821in}{0.009821in}}%
\pgfpathcurveto{\pgfqpoint{0.007216in}{0.012425in}}{\pgfqpoint{0.003683in}{0.013889in}}{\pgfqpoint{0.000000in}{0.013889in}}%
\pgfpathcurveto{\pgfqpoint{-0.003683in}{0.013889in}}{\pgfqpoint{-0.007216in}{0.012425in}}{\pgfqpoint{-0.009821in}{0.009821in}}%
\pgfpathcurveto{\pgfqpoint{-0.012425in}{0.007216in}}{\pgfqpoint{-0.013889in}{0.003683in}}{\pgfqpoint{-0.013889in}{0.000000in}}%
\pgfpathcurveto{\pgfqpoint{-0.013889in}{-0.003683in}}{\pgfqpoint{-0.012425in}{-0.007216in}}{\pgfqpoint{-0.009821in}{-0.009821in}}%
\pgfpathcurveto{\pgfqpoint{-0.007216in}{-0.012425in}}{\pgfqpoint{-0.003683in}{-0.013889in}}{\pgfqpoint{0.000000in}{-0.013889in}}%
\pgfpathlineto{\pgfqpoint{0.000000in}{-0.013889in}}%
\pgfpathclose%
\pgfusepath{stroke,fill}%
}%
\begin{pgfscope}%
\pgfsys@transformshift{4.327155in}{4.580101in}%
\pgfsys@useobject{currentmarker}{}%
\end{pgfscope}%
\end{pgfscope}%
\begin{pgfscope}%
\pgfpathrectangle{\pgfqpoint{0.100000in}{2.493063in}}{\pgfqpoint{5.037500in}{3.427208in}}%
\pgfusepath{clip}%
\pgfsetrectcap%
\pgfsetroundjoin%
\pgfsetlinewidth{1.505625pt}%
\definecolor{currentstroke}{rgb}{0.501961,0.501961,0.501961}%
\pgfsetstrokecolor{currentstroke}%
\pgfsetstrokeopacity{0.500000}%
\pgfsetdash{}{0pt}%
\pgfpathmoveto{\pgfqpoint{4.906699in}{4.954236in}}%
\pgfusepath{stroke}%
\end{pgfscope}%
\begin{pgfscope}%
\pgfpathrectangle{\pgfqpoint{0.100000in}{2.493063in}}{\pgfqpoint{5.037500in}{3.427208in}}%
\pgfusepath{clip}%
\pgfsetbuttcap%
\pgfsetroundjoin%
\definecolor{currentfill}{rgb}{0.501961,0.501961,0.501961}%
\pgfsetfillcolor{currentfill}%
\pgfsetfillopacity{0.500000}%
\pgfsetlinewidth{0.250937pt}%
\definecolor{currentstroke}{rgb}{0.000000,0.000000,0.000000}%
\pgfsetstrokecolor{currentstroke}%
\pgfsetstrokeopacity{0.500000}%
\pgfsetdash{}{0pt}%
\pgfsys@defobject{currentmarker}{\pgfqpoint{-0.013889in}{-0.013889in}}{\pgfqpoint{0.013889in}{0.013889in}}{%
\pgfpathmoveto{\pgfqpoint{0.000000in}{-0.013889in}}%
\pgfpathcurveto{\pgfqpoint{0.003683in}{-0.013889in}}{\pgfqpoint{0.007216in}{-0.012425in}}{\pgfqpoint{0.009821in}{-0.009821in}}%
\pgfpathcurveto{\pgfqpoint{0.012425in}{-0.007216in}}{\pgfqpoint{0.013889in}{-0.003683in}}{\pgfqpoint{0.013889in}{0.000000in}}%
\pgfpathcurveto{\pgfqpoint{0.013889in}{0.003683in}}{\pgfqpoint{0.012425in}{0.007216in}}{\pgfqpoint{0.009821in}{0.009821in}}%
\pgfpathcurveto{\pgfqpoint{0.007216in}{0.012425in}}{\pgfqpoint{0.003683in}{0.013889in}}{\pgfqpoint{0.000000in}{0.013889in}}%
\pgfpathcurveto{\pgfqpoint{-0.003683in}{0.013889in}}{\pgfqpoint{-0.007216in}{0.012425in}}{\pgfqpoint{-0.009821in}{0.009821in}}%
\pgfpathcurveto{\pgfqpoint{-0.012425in}{0.007216in}}{\pgfqpoint{-0.013889in}{0.003683in}}{\pgfqpoint{-0.013889in}{0.000000in}}%
\pgfpathcurveto{\pgfqpoint{-0.013889in}{-0.003683in}}{\pgfqpoint{-0.012425in}{-0.007216in}}{\pgfqpoint{-0.009821in}{-0.009821in}}%
\pgfpathcurveto{\pgfqpoint{-0.007216in}{-0.012425in}}{\pgfqpoint{-0.003683in}{-0.013889in}}{\pgfqpoint{0.000000in}{-0.013889in}}%
\pgfpathlineto{\pgfqpoint{0.000000in}{-0.013889in}}%
\pgfpathclose%
\pgfusepath{stroke,fill}%
}%
\begin{pgfscope}%
\pgfsys@transformshift{4.906699in}{4.954236in}%
\pgfsys@useobject{currentmarker}{}%
\end{pgfscope}%
\end{pgfscope}%
\begin{pgfscope}%
\pgfpathrectangle{\pgfqpoint{0.100000in}{2.493063in}}{\pgfqpoint{5.037500in}{3.427208in}}%
\pgfusepath{clip}%
\pgfsetrectcap%
\pgfsetroundjoin%
\pgfsetlinewidth{1.505625pt}%
\definecolor{currentstroke}{rgb}{0.000000,0.000000,1.000000}%
\pgfsetstrokecolor{currentstroke}%
\pgfsetstrokeopacity{0.500000}%
\pgfsetdash{}{0pt}%
\pgfpathmoveto{\pgfqpoint{4.820957in}{5.015309in}}%
\pgfusepath{stroke}%
\end{pgfscope}%
\begin{pgfscope}%
\pgfpathrectangle{\pgfqpoint{0.100000in}{2.493063in}}{\pgfqpoint{5.037500in}{3.427208in}}%
\pgfusepath{clip}%
\pgfsetbuttcap%
\pgfsetroundjoin%
\definecolor{currentfill}{rgb}{0.000000,0.000000,1.000000}%
\pgfsetfillcolor{currentfill}%
\pgfsetfillopacity{0.500000}%
\pgfsetlinewidth{0.250937pt}%
\definecolor{currentstroke}{rgb}{0.000000,0.000000,0.000000}%
\pgfsetstrokecolor{currentstroke}%
\pgfsetstrokeopacity{0.500000}%
\pgfsetdash{}{0pt}%
\pgfsys@defobject{currentmarker}{\pgfqpoint{-0.013889in}{-0.013889in}}{\pgfqpoint{0.013889in}{0.013889in}}{%
\pgfpathmoveto{\pgfqpoint{0.000000in}{-0.013889in}}%
\pgfpathcurveto{\pgfqpoint{0.003683in}{-0.013889in}}{\pgfqpoint{0.007216in}{-0.012425in}}{\pgfqpoint{0.009821in}{-0.009821in}}%
\pgfpathcurveto{\pgfqpoint{0.012425in}{-0.007216in}}{\pgfqpoint{0.013889in}{-0.003683in}}{\pgfqpoint{0.013889in}{0.000000in}}%
\pgfpathcurveto{\pgfqpoint{0.013889in}{0.003683in}}{\pgfqpoint{0.012425in}{0.007216in}}{\pgfqpoint{0.009821in}{0.009821in}}%
\pgfpathcurveto{\pgfqpoint{0.007216in}{0.012425in}}{\pgfqpoint{0.003683in}{0.013889in}}{\pgfqpoint{0.000000in}{0.013889in}}%
\pgfpathcurveto{\pgfqpoint{-0.003683in}{0.013889in}}{\pgfqpoint{-0.007216in}{0.012425in}}{\pgfqpoint{-0.009821in}{0.009821in}}%
\pgfpathcurveto{\pgfqpoint{-0.012425in}{0.007216in}}{\pgfqpoint{-0.013889in}{0.003683in}}{\pgfqpoint{-0.013889in}{0.000000in}}%
\pgfpathcurveto{\pgfqpoint{-0.013889in}{-0.003683in}}{\pgfqpoint{-0.012425in}{-0.007216in}}{\pgfqpoint{-0.009821in}{-0.009821in}}%
\pgfpathcurveto{\pgfqpoint{-0.007216in}{-0.012425in}}{\pgfqpoint{-0.003683in}{-0.013889in}}{\pgfqpoint{0.000000in}{-0.013889in}}%
\pgfpathlineto{\pgfqpoint{0.000000in}{-0.013889in}}%
\pgfpathclose%
\pgfusepath{stroke,fill}%
}%
\begin{pgfscope}%
\pgfsys@transformshift{4.820957in}{5.015309in}%
\pgfsys@useobject{currentmarker}{}%
\end{pgfscope}%
\end{pgfscope}%
\begin{pgfscope}%
\pgfpathrectangle{\pgfqpoint{0.100000in}{2.493063in}}{\pgfqpoint{5.037500in}{3.427208in}}%
\pgfusepath{clip}%
\pgfsetrectcap%
\pgfsetroundjoin%
\pgfsetlinewidth{1.505625pt}%
\definecolor{currentstroke}{rgb}{0.000000,0.000000,1.000000}%
\pgfsetstrokecolor{currentstroke}%
\pgfsetstrokeopacity{0.500000}%
\pgfsetdash{}{0pt}%
\pgfpathmoveto{\pgfqpoint{4.758476in}{5.018361in}}%
\pgfusepath{stroke}%
\end{pgfscope}%
\begin{pgfscope}%
\pgfpathrectangle{\pgfqpoint{0.100000in}{2.493063in}}{\pgfqpoint{5.037500in}{3.427208in}}%
\pgfusepath{clip}%
\pgfsetbuttcap%
\pgfsetroundjoin%
\definecolor{currentfill}{rgb}{0.000000,0.000000,1.000000}%
\pgfsetfillcolor{currentfill}%
\pgfsetfillopacity{0.500000}%
\pgfsetlinewidth{0.250937pt}%
\definecolor{currentstroke}{rgb}{0.000000,0.000000,0.000000}%
\pgfsetstrokecolor{currentstroke}%
\pgfsetstrokeopacity{0.500000}%
\pgfsetdash{}{0pt}%
\pgfsys@defobject{currentmarker}{\pgfqpoint{-0.005556in}{-0.005556in}}{\pgfqpoint{0.005556in}{0.005556in}}{%
\pgfpathmoveto{\pgfqpoint{0.000000in}{-0.005556in}}%
\pgfpathcurveto{\pgfqpoint{0.001473in}{-0.005556in}}{\pgfqpoint{0.002887in}{-0.004970in}}{\pgfqpoint{0.003928in}{-0.003928in}}%
\pgfpathcurveto{\pgfqpoint{0.004970in}{-0.002887in}}{\pgfqpoint{0.005556in}{-0.001473in}}{\pgfqpoint{0.005556in}{0.000000in}}%
\pgfpathcurveto{\pgfqpoint{0.005556in}{0.001473in}}{\pgfqpoint{0.004970in}{0.002887in}}{\pgfqpoint{0.003928in}{0.003928in}}%
\pgfpathcurveto{\pgfqpoint{0.002887in}{0.004970in}}{\pgfqpoint{0.001473in}{0.005556in}}{\pgfqpoint{0.000000in}{0.005556in}}%
\pgfpathcurveto{\pgfqpoint{-0.001473in}{0.005556in}}{\pgfqpoint{-0.002887in}{0.004970in}}{\pgfqpoint{-0.003928in}{0.003928in}}%
\pgfpathcurveto{\pgfqpoint{-0.004970in}{0.002887in}}{\pgfqpoint{-0.005556in}{0.001473in}}{\pgfqpoint{-0.005556in}{0.000000in}}%
\pgfpathcurveto{\pgfqpoint{-0.005556in}{-0.001473in}}{\pgfqpoint{-0.004970in}{-0.002887in}}{\pgfqpoint{-0.003928in}{-0.003928in}}%
\pgfpathcurveto{\pgfqpoint{-0.002887in}{-0.004970in}}{\pgfqpoint{-0.001473in}{-0.005556in}}{\pgfqpoint{0.000000in}{-0.005556in}}%
\pgfpathlineto{\pgfqpoint{0.000000in}{-0.005556in}}%
\pgfpathclose%
\pgfusepath{stroke,fill}%
}%
\begin{pgfscope}%
\pgfsys@transformshift{4.758476in}{5.018361in}%
\pgfsys@useobject{currentmarker}{}%
\end{pgfscope}%
\end{pgfscope}%
\begin{pgfscope}%
\pgfpathrectangle{\pgfqpoint{0.100000in}{2.493063in}}{\pgfqpoint{5.037500in}{3.427208in}}%
\pgfusepath{clip}%
\pgfsetrectcap%
\pgfsetroundjoin%
\pgfsetlinewidth{1.505625pt}%
\definecolor{currentstroke}{rgb}{0.501961,0.501961,0.501961}%
\pgfsetstrokecolor{currentstroke}%
\pgfsetstrokeopacity{0.500000}%
\pgfsetdash{}{0pt}%
\pgfpathmoveto{\pgfqpoint{4.853127in}{4.937502in}}%
\pgfusepath{stroke}%
\end{pgfscope}%
\begin{pgfscope}%
\pgfpathrectangle{\pgfqpoint{0.100000in}{2.493063in}}{\pgfqpoint{5.037500in}{3.427208in}}%
\pgfusepath{clip}%
\pgfsetbuttcap%
\pgfsetroundjoin%
\definecolor{currentfill}{rgb}{0.501961,0.501961,0.501961}%
\pgfsetfillcolor{currentfill}%
\pgfsetfillopacity{0.500000}%
\pgfsetlinewidth{0.250937pt}%
\definecolor{currentstroke}{rgb}{0.000000,0.000000,0.000000}%
\pgfsetstrokecolor{currentstroke}%
\pgfsetstrokeopacity{0.500000}%
\pgfsetdash{}{0pt}%
\pgfsys@defobject{currentmarker}{\pgfqpoint{-0.013889in}{-0.013889in}}{\pgfqpoint{0.013889in}{0.013889in}}{%
\pgfpathmoveto{\pgfqpoint{0.000000in}{-0.013889in}}%
\pgfpathcurveto{\pgfqpoint{0.003683in}{-0.013889in}}{\pgfqpoint{0.007216in}{-0.012425in}}{\pgfqpoint{0.009821in}{-0.009821in}}%
\pgfpathcurveto{\pgfqpoint{0.012425in}{-0.007216in}}{\pgfqpoint{0.013889in}{-0.003683in}}{\pgfqpoint{0.013889in}{0.000000in}}%
\pgfpathcurveto{\pgfqpoint{0.013889in}{0.003683in}}{\pgfqpoint{0.012425in}{0.007216in}}{\pgfqpoint{0.009821in}{0.009821in}}%
\pgfpathcurveto{\pgfqpoint{0.007216in}{0.012425in}}{\pgfqpoint{0.003683in}{0.013889in}}{\pgfqpoint{0.000000in}{0.013889in}}%
\pgfpathcurveto{\pgfqpoint{-0.003683in}{0.013889in}}{\pgfqpoint{-0.007216in}{0.012425in}}{\pgfqpoint{-0.009821in}{0.009821in}}%
\pgfpathcurveto{\pgfqpoint{-0.012425in}{0.007216in}}{\pgfqpoint{-0.013889in}{0.003683in}}{\pgfqpoint{-0.013889in}{0.000000in}}%
\pgfpathcurveto{\pgfqpoint{-0.013889in}{-0.003683in}}{\pgfqpoint{-0.012425in}{-0.007216in}}{\pgfqpoint{-0.009821in}{-0.009821in}}%
\pgfpathcurveto{\pgfqpoint{-0.007216in}{-0.012425in}}{\pgfqpoint{-0.003683in}{-0.013889in}}{\pgfqpoint{0.000000in}{-0.013889in}}%
\pgfpathlineto{\pgfqpoint{0.000000in}{-0.013889in}}%
\pgfpathclose%
\pgfusepath{stroke,fill}%
}%
\begin{pgfscope}%
\pgfsys@transformshift{4.853127in}{4.937502in}%
\pgfsys@useobject{currentmarker}{}%
\end{pgfscope}%
\end{pgfscope}%
\begin{pgfscope}%
\pgfpathrectangle{\pgfqpoint{0.100000in}{2.493063in}}{\pgfqpoint{5.037500in}{3.427208in}}%
\pgfusepath{clip}%
\pgfsetrectcap%
\pgfsetroundjoin%
\pgfsetlinewidth{1.505625pt}%
\definecolor{currentstroke}{rgb}{0.501961,0.501961,0.501961}%
\pgfsetstrokecolor{currentstroke}%
\pgfsetstrokeopacity{0.500000}%
\pgfsetdash{}{0pt}%
\pgfpathmoveto{\pgfqpoint{4.638079in}{4.979009in}}%
\pgfusepath{stroke}%
\end{pgfscope}%
\begin{pgfscope}%
\pgfpathrectangle{\pgfqpoint{0.100000in}{2.493063in}}{\pgfqpoint{5.037500in}{3.427208in}}%
\pgfusepath{clip}%
\pgfsetbuttcap%
\pgfsetroundjoin%
\definecolor{currentfill}{rgb}{0.501961,0.501961,0.501961}%
\pgfsetfillcolor{currentfill}%
\pgfsetfillopacity{0.500000}%
\pgfsetlinewidth{0.250937pt}%
\definecolor{currentstroke}{rgb}{0.000000,0.000000,0.000000}%
\pgfsetstrokecolor{currentstroke}%
\pgfsetstrokeopacity{0.500000}%
\pgfsetdash{}{0pt}%
\pgfsys@defobject{currentmarker}{\pgfqpoint{-0.013889in}{-0.013889in}}{\pgfqpoint{0.013889in}{0.013889in}}{%
\pgfpathmoveto{\pgfqpoint{0.000000in}{-0.013889in}}%
\pgfpathcurveto{\pgfqpoint{0.003683in}{-0.013889in}}{\pgfqpoint{0.007216in}{-0.012425in}}{\pgfqpoint{0.009821in}{-0.009821in}}%
\pgfpathcurveto{\pgfqpoint{0.012425in}{-0.007216in}}{\pgfqpoint{0.013889in}{-0.003683in}}{\pgfqpoint{0.013889in}{0.000000in}}%
\pgfpathcurveto{\pgfqpoint{0.013889in}{0.003683in}}{\pgfqpoint{0.012425in}{0.007216in}}{\pgfqpoint{0.009821in}{0.009821in}}%
\pgfpathcurveto{\pgfqpoint{0.007216in}{0.012425in}}{\pgfqpoint{0.003683in}{0.013889in}}{\pgfqpoint{0.000000in}{0.013889in}}%
\pgfpathcurveto{\pgfqpoint{-0.003683in}{0.013889in}}{\pgfqpoint{-0.007216in}{0.012425in}}{\pgfqpoint{-0.009821in}{0.009821in}}%
\pgfpathcurveto{\pgfqpoint{-0.012425in}{0.007216in}}{\pgfqpoint{-0.013889in}{0.003683in}}{\pgfqpoint{-0.013889in}{0.000000in}}%
\pgfpathcurveto{\pgfqpoint{-0.013889in}{-0.003683in}}{\pgfqpoint{-0.012425in}{-0.007216in}}{\pgfqpoint{-0.009821in}{-0.009821in}}%
\pgfpathcurveto{\pgfqpoint{-0.007216in}{-0.012425in}}{\pgfqpoint{-0.003683in}{-0.013889in}}{\pgfqpoint{0.000000in}{-0.013889in}}%
\pgfpathlineto{\pgfqpoint{0.000000in}{-0.013889in}}%
\pgfpathclose%
\pgfusepath{stroke,fill}%
}%
\begin{pgfscope}%
\pgfsys@transformshift{4.638079in}{4.979009in}%
\pgfsys@useobject{currentmarker}{}%
\end{pgfscope}%
\end{pgfscope}%
\begin{pgfscope}%
\pgfpathrectangle{\pgfqpoint{0.100000in}{2.493063in}}{\pgfqpoint{5.037500in}{3.427208in}}%
\pgfusepath{clip}%
\pgfsetrectcap%
\pgfsetroundjoin%
\pgfsetlinewidth{1.505625pt}%
\definecolor{currentstroke}{rgb}{0.501961,0.501961,0.501961}%
\pgfsetstrokecolor{currentstroke}%
\pgfsetstrokeopacity{0.500000}%
\pgfsetdash{}{0pt}%
\pgfpathmoveto{\pgfqpoint{4.702051in}{4.953445in}}%
\pgfusepath{stroke}%
\end{pgfscope}%
\begin{pgfscope}%
\pgfpathrectangle{\pgfqpoint{0.100000in}{2.493063in}}{\pgfqpoint{5.037500in}{3.427208in}}%
\pgfusepath{clip}%
\pgfsetbuttcap%
\pgfsetroundjoin%
\definecolor{currentfill}{rgb}{0.501961,0.501961,0.501961}%
\pgfsetfillcolor{currentfill}%
\pgfsetfillopacity{0.500000}%
\pgfsetlinewidth{0.250937pt}%
\definecolor{currentstroke}{rgb}{0.000000,0.000000,0.000000}%
\pgfsetstrokecolor{currentstroke}%
\pgfsetstrokeopacity{0.500000}%
\pgfsetdash{}{0pt}%
\pgfsys@defobject{currentmarker}{\pgfqpoint{-0.013889in}{-0.013889in}}{\pgfqpoint{0.013889in}{0.013889in}}{%
\pgfpathmoveto{\pgfqpoint{0.000000in}{-0.013889in}}%
\pgfpathcurveto{\pgfqpoint{0.003683in}{-0.013889in}}{\pgfqpoint{0.007216in}{-0.012425in}}{\pgfqpoint{0.009821in}{-0.009821in}}%
\pgfpathcurveto{\pgfqpoint{0.012425in}{-0.007216in}}{\pgfqpoint{0.013889in}{-0.003683in}}{\pgfqpoint{0.013889in}{0.000000in}}%
\pgfpathcurveto{\pgfqpoint{0.013889in}{0.003683in}}{\pgfqpoint{0.012425in}{0.007216in}}{\pgfqpoint{0.009821in}{0.009821in}}%
\pgfpathcurveto{\pgfqpoint{0.007216in}{0.012425in}}{\pgfqpoint{0.003683in}{0.013889in}}{\pgfqpoint{0.000000in}{0.013889in}}%
\pgfpathcurveto{\pgfqpoint{-0.003683in}{0.013889in}}{\pgfqpoint{-0.007216in}{0.012425in}}{\pgfqpoint{-0.009821in}{0.009821in}}%
\pgfpathcurveto{\pgfqpoint{-0.012425in}{0.007216in}}{\pgfqpoint{-0.013889in}{0.003683in}}{\pgfqpoint{-0.013889in}{0.000000in}}%
\pgfpathcurveto{\pgfqpoint{-0.013889in}{-0.003683in}}{\pgfqpoint{-0.012425in}{-0.007216in}}{\pgfqpoint{-0.009821in}{-0.009821in}}%
\pgfpathcurveto{\pgfqpoint{-0.007216in}{-0.012425in}}{\pgfqpoint{-0.003683in}{-0.013889in}}{\pgfqpoint{0.000000in}{-0.013889in}}%
\pgfpathlineto{\pgfqpoint{0.000000in}{-0.013889in}}%
\pgfpathclose%
\pgfusepath{stroke,fill}%
}%
\begin{pgfscope}%
\pgfsys@transformshift{4.702051in}{4.953445in}%
\pgfsys@useobject{currentmarker}{}%
\end{pgfscope}%
\end{pgfscope}%
\begin{pgfscope}%
\pgfpathrectangle{\pgfqpoint{0.100000in}{2.493063in}}{\pgfqpoint{5.037500in}{3.427208in}}%
\pgfusepath{clip}%
\pgfsetrectcap%
\pgfsetroundjoin%
\pgfsetlinewidth{1.505625pt}%
\definecolor{currentstroke}{rgb}{0.501961,0.501961,0.501961}%
\pgfsetstrokecolor{currentstroke}%
\pgfsetstrokeopacity{0.500000}%
\pgfsetdash{}{0pt}%
\pgfpathmoveto{\pgfqpoint{4.762625in}{4.988152in}}%
\pgfusepath{stroke}%
\end{pgfscope}%
\begin{pgfscope}%
\pgfpathrectangle{\pgfqpoint{0.100000in}{2.493063in}}{\pgfqpoint{5.037500in}{3.427208in}}%
\pgfusepath{clip}%
\pgfsetbuttcap%
\pgfsetroundjoin%
\definecolor{currentfill}{rgb}{0.501961,0.501961,0.501961}%
\pgfsetfillcolor{currentfill}%
\pgfsetfillopacity{0.500000}%
\pgfsetlinewidth{0.250937pt}%
\definecolor{currentstroke}{rgb}{0.000000,0.000000,0.000000}%
\pgfsetstrokecolor{currentstroke}%
\pgfsetstrokeopacity{0.500000}%
\pgfsetdash{}{0pt}%
\pgfsys@defobject{currentmarker}{\pgfqpoint{-0.013889in}{-0.013889in}}{\pgfqpoint{0.013889in}{0.013889in}}{%
\pgfpathmoveto{\pgfqpoint{0.000000in}{-0.013889in}}%
\pgfpathcurveto{\pgfqpoint{0.003683in}{-0.013889in}}{\pgfqpoint{0.007216in}{-0.012425in}}{\pgfqpoint{0.009821in}{-0.009821in}}%
\pgfpathcurveto{\pgfqpoint{0.012425in}{-0.007216in}}{\pgfqpoint{0.013889in}{-0.003683in}}{\pgfqpoint{0.013889in}{0.000000in}}%
\pgfpathcurveto{\pgfqpoint{0.013889in}{0.003683in}}{\pgfqpoint{0.012425in}{0.007216in}}{\pgfqpoint{0.009821in}{0.009821in}}%
\pgfpathcurveto{\pgfqpoint{0.007216in}{0.012425in}}{\pgfqpoint{0.003683in}{0.013889in}}{\pgfqpoint{0.000000in}{0.013889in}}%
\pgfpathcurveto{\pgfqpoint{-0.003683in}{0.013889in}}{\pgfqpoint{-0.007216in}{0.012425in}}{\pgfqpoint{-0.009821in}{0.009821in}}%
\pgfpathcurveto{\pgfqpoint{-0.012425in}{0.007216in}}{\pgfqpoint{-0.013889in}{0.003683in}}{\pgfqpoint{-0.013889in}{0.000000in}}%
\pgfpathcurveto{\pgfqpoint{-0.013889in}{-0.003683in}}{\pgfqpoint{-0.012425in}{-0.007216in}}{\pgfqpoint{-0.009821in}{-0.009821in}}%
\pgfpathcurveto{\pgfqpoint{-0.007216in}{-0.012425in}}{\pgfqpoint{-0.003683in}{-0.013889in}}{\pgfqpoint{0.000000in}{-0.013889in}}%
\pgfpathlineto{\pgfqpoint{0.000000in}{-0.013889in}}%
\pgfpathclose%
\pgfusepath{stroke,fill}%
}%
\begin{pgfscope}%
\pgfsys@transformshift{4.762625in}{4.988152in}%
\pgfsys@useobject{currentmarker}{}%
\end{pgfscope}%
\end{pgfscope}%
\begin{pgfscope}%
\pgfpathrectangle{\pgfqpoint{0.100000in}{2.493063in}}{\pgfqpoint{5.037500in}{3.427208in}}%
\pgfusepath{clip}%
\pgfsetrectcap%
\pgfsetroundjoin%
\pgfsetlinewidth{1.505625pt}%
\definecolor{currentstroke}{rgb}{0.000000,0.000000,1.000000}%
\pgfsetstrokecolor{currentstroke}%
\pgfsetstrokeopacity{0.500000}%
\pgfsetdash{}{0pt}%
\pgfpathmoveto{\pgfqpoint{3.763378in}{4.797009in}}%
\pgfusepath{stroke}%
\end{pgfscope}%
\begin{pgfscope}%
\pgfpathrectangle{\pgfqpoint{0.100000in}{2.493063in}}{\pgfqpoint{5.037500in}{3.427208in}}%
\pgfusepath{clip}%
\pgfsetbuttcap%
\pgfsetroundjoin%
\definecolor{currentfill}{rgb}{0.000000,0.000000,1.000000}%
\pgfsetfillcolor{currentfill}%
\pgfsetfillopacity{0.500000}%
\pgfsetlinewidth{0.250937pt}%
\definecolor{currentstroke}{rgb}{0.000000,0.000000,0.000000}%
\pgfsetstrokecolor{currentstroke}%
\pgfsetstrokeopacity{0.500000}%
\pgfsetdash{}{0pt}%
\pgfsys@defobject{currentmarker}{\pgfqpoint{-0.013889in}{-0.013889in}}{\pgfqpoint{0.013889in}{0.013889in}}{%
\pgfpathmoveto{\pgfqpoint{0.000000in}{-0.013889in}}%
\pgfpathcurveto{\pgfqpoint{0.003683in}{-0.013889in}}{\pgfqpoint{0.007216in}{-0.012425in}}{\pgfqpoint{0.009821in}{-0.009821in}}%
\pgfpathcurveto{\pgfqpoint{0.012425in}{-0.007216in}}{\pgfqpoint{0.013889in}{-0.003683in}}{\pgfqpoint{0.013889in}{0.000000in}}%
\pgfpathcurveto{\pgfqpoint{0.013889in}{0.003683in}}{\pgfqpoint{0.012425in}{0.007216in}}{\pgfqpoint{0.009821in}{0.009821in}}%
\pgfpathcurveto{\pgfqpoint{0.007216in}{0.012425in}}{\pgfqpoint{0.003683in}{0.013889in}}{\pgfqpoint{0.000000in}{0.013889in}}%
\pgfpathcurveto{\pgfqpoint{-0.003683in}{0.013889in}}{\pgfqpoint{-0.007216in}{0.012425in}}{\pgfqpoint{-0.009821in}{0.009821in}}%
\pgfpathcurveto{\pgfqpoint{-0.012425in}{0.007216in}}{\pgfqpoint{-0.013889in}{0.003683in}}{\pgfqpoint{-0.013889in}{0.000000in}}%
\pgfpathcurveto{\pgfqpoint{-0.013889in}{-0.003683in}}{\pgfqpoint{-0.012425in}{-0.007216in}}{\pgfqpoint{-0.009821in}{-0.009821in}}%
\pgfpathcurveto{\pgfqpoint{-0.007216in}{-0.012425in}}{\pgfqpoint{-0.003683in}{-0.013889in}}{\pgfqpoint{0.000000in}{-0.013889in}}%
\pgfpathlineto{\pgfqpoint{0.000000in}{-0.013889in}}%
\pgfpathclose%
\pgfusepath{stroke,fill}%
}%
\begin{pgfscope}%
\pgfsys@transformshift{3.763378in}{4.797009in}%
\pgfsys@useobject{currentmarker}{}%
\end{pgfscope}%
\end{pgfscope}%
\begin{pgfscope}%
\pgfpathrectangle{\pgfqpoint{0.100000in}{2.493063in}}{\pgfqpoint{5.037500in}{3.427208in}}%
\pgfusepath{clip}%
\pgfsetrectcap%
\pgfsetroundjoin%
\pgfsetlinewidth{1.505625pt}%
\definecolor{currentstroke}{rgb}{0.000000,0.000000,1.000000}%
\pgfsetstrokecolor{currentstroke}%
\pgfsetstrokeopacity{0.500000}%
\pgfsetdash{}{0pt}%
\pgfpathmoveto{\pgfqpoint{3.639727in}{4.788532in}}%
\pgfusepath{stroke}%
\end{pgfscope}%
\begin{pgfscope}%
\pgfpathrectangle{\pgfqpoint{0.100000in}{2.493063in}}{\pgfqpoint{5.037500in}{3.427208in}}%
\pgfusepath{clip}%
\pgfsetbuttcap%
\pgfsetroundjoin%
\definecolor{currentfill}{rgb}{0.000000,0.000000,1.000000}%
\pgfsetfillcolor{currentfill}%
\pgfsetfillopacity{0.500000}%
\pgfsetlinewidth{0.250937pt}%
\definecolor{currentstroke}{rgb}{0.000000,0.000000,0.000000}%
\pgfsetstrokecolor{currentstroke}%
\pgfsetstrokeopacity{0.500000}%
\pgfsetdash{}{0pt}%
\pgfsys@defobject{currentmarker}{\pgfqpoint{-0.008333in}{-0.008333in}}{\pgfqpoint{0.008333in}{0.008333in}}{%
\pgfpathmoveto{\pgfqpoint{0.000000in}{-0.008333in}}%
\pgfpathcurveto{\pgfqpoint{0.002210in}{-0.008333in}}{\pgfqpoint{0.004330in}{-0.007455in}}{\pgfqpoint{0.005893in}{-0.005893in}}%
\pgfpathcurveto{\pgfqpoint{0.007455in}{-0.004330in}}{\pgfqpoint{0.008333in}{-0.002210in}}{\pgfqpoint{0.008333in}{0.000000in}}%
\pgfpathcurveto{\pgfqpoint{0.008333in}{0.002210in}}{\pgfqpoint{0.007455in}{0.004330in}}{\pgfqpoint{0.005893in}{0.005893in}}%
\pgfpathcurveto{\pgfqpoint{0.004330in}{0.007455in}}{\pgfqpoint{0.002210in}{0.008333in}}{\pgfqpoint{0.000000in}{0.008333in}}%
\pgfpathcurveto{\pgfqpoint{-0.002210in}{0.008333in}}{\pgfqpoint{-0.004330in}{0.007455in}}{\pgfqpoint{-0.005893in}{0.005893in}}%
\pgfpathcurveto{\pgfqpoint{-0.007455in}{0.004330in}}{\pgfqpoint{-0.008333in}{0.002210in}}{\pgfqpoint{-0.008333in}{0.000000in}}%
\pgfpathcurveto{\pgfqpoint{-0.008333in}{-0.002210in}}{\pgfqpoint{-0.007455in}{-0.004330in}}{\pgfqpoint{-0.005893in}{-0.005893in}}%
\pgfpathcurveto{\pgfqpoint{-0.004330in}{-0.007455in}}{\pgfqpoint{-0.002210in}{-0.008333in}}{\pgfqpoint{0.000000in}{-0.008333in}}%
\pgfpathlineto{\pgfqpoint{0.000000in}{-0.008333in}}%
\pgfpathclose%
\pgfusepath{stroke,fill}%
}%
\begin{pgfscope}%
\pgfsys@transformshift{3.639727in}{4.788532in}%
\pgfsys@useobject{currentmarker}{}%
\end{pgfscope}%
\end{pgfscope}%
\begin{pgfscope}%
\pgfpathrectangle{\pgfqpoint{0.100000in}{2.493063in}}{\pgfqpoint{5.037500in}{3.427208in}}%
\pgfusepath{clip}%
\pgfsetrectcap%
\pgfsetroundjoin%
\pgfsetlinewidth{1.505625pt}%
\definecolor{currentstroke}{rgb}{0.501961,0.501961,0.501961}%
\pgfsetstrokecolor{currentstroke}%
\pgfsetstrokeopacity{0.500000}%
\pgfsetdash{}{0pt}%
\pgfpathmoveto{\pgfqpoint{3.731354in}{4.947293in}}%
\pgfusepath{stroke}%
\end{pgfscope}%
\begin{pgfscope}%
\pgfpathrectangle{\pgfqpoint{0.100000in}{2.493063in}}{\pgfqpoint{5.037500in}{3.427208in}}%
\pgfusepath{clip}%
\pgfsetbuttcap%
\pgfsetroundjoin%
\definecolor{currentfill}{rgb}{0.501961,0.501961,0.501961}%
\pgfsetfillcolor{currentfill}%
\pgfsetfillopacity{0.500000}%
\pgfsetlinewidth{0.250937pt}%
\definecolor{currentstroke}{rgb}{0.000000,0.000000,0.000000}%
\pgfsetstrokecolor{currentstroke}%
\pgfsetstrokeopacity{0.500000}%
\pgfsetdash{}{0pt}%
\pgfsys@defobject{currentmarker}{\pgfqpoint{-0.013889in}{-0.013889in}}{\pgfqpoint{0.013889in}{0.013889in}}{%
\pgfpathmoveto{\pgfqpoint{0.000000in}{-0.013889in}}%
\pgfpathcurveto{\pgfqpoint{0.003683in}{-0.013889in}}{\pgfqpoint{0.007216in}{-0.012425in}}{\pgfqpoint{0.009821in}{-0.009821in}}%
\pgfpathcurveto{\pgfqpoint{0.012425in}{-0.007216in}}{\pgfqpoint{0.013889in}{-0.003683in}}{\pgfqpoint{0.013889in}{0.000000in}}%
\pgfpathcurveto{\pgfqpoint{0.013889in}{0.003683in}}{\pgfqpoint{0.012425in}{0.007216in}}{\pgfqpoint{0.009821in}{0.009821in}}%
\pgfpathcurveto{\pgfqpoint{0.007216in}{0.012425in}}{\pgfqpoint{0.003683in}{0.013889in}}{\pgfqpoint{0.000000in}{0.013889in}}%
\pgfpathcurveto{\pgfqpoint{-0.003683in}{0.013889in}}{\pgfqpoint{-0.007216in}{0.012425in}}{\pgfqpoint{-0.009821in}{0.009821in}}%
\pgfpathcurveto{\pgfqpoint{-0.012425in}{0.007216in}}{\pgfqpoint{-0.013889in}{0.003683in}}{\pgfqpoint{-0.013889in}{0.000000in}}%
\pgfpathcurveto{\pgfqpoint{-0.013889in}{-0.003683in}}{\pgfqpoint{-0.012425in}{-0.007216in}}{\pgfqpoint{-0.009821in}{-0.009821in}}%
\pgfpathcurveto{\pgfqpoint{-0.007216in}{-0.012425in}}{\pgfqpoint{-0.003683in}{-0.013889in}}{\pgfqpoint{0.000000in}{-0.013889in}}%
\pgfpathlineto{\pgfqpoint{0.000000in}{-0.013889in}}%
\pgfpathclose%
\pgfusepath{stroke,fill}%
}%
\begin{pgfscope}%
\pgfsys@transformshift{3.731354in}{4.947293in}%
\pgfsys@useobject{currentmarker}{}%
\end{pgfscope}%
\end{pgfscope}%
\begin{pgfscope}%
\pgfpathrectangle{\pgfqpoint{0.100000in}{2.493063in}}{\pgfqpoint{5.037500in}{3.427208in}}%
\pgfusepath{clip}%
\pgfsetrectcap%
\pgfsetroundjoin%
\pgfsetlinewidth{1.505625pt}%
\definecolor{currentstroke}{rgb}{0.000000,0.000000,1.000000}%
\pgfsetstrokecolor{currentstroke}%
\pgfsetstrokeopacity{0.500000}%
\pgfsetdash{}{0pt}%
\pgfpathmoveto{\pgfqpoint{3.819255in}{4.813545in}}%
\pgfusepath{stroke}%
\end{pgfscope}%
\begin{pgfscope}%
\pgfpathrectangle{\pgfqpoint{0.100000in}{2.493063in}}{\pgfqpoint{5.037500in}{3.427208in}}%
\pgfusepath{clip}%
\pgfsetbuttcap%
\pgfsetroundjoin%
\definecolor{currentfill}{rgb}{0.000000,0.000000,1.000000}%
\pgfsetfillcolor{currentfill}%
\pgfsetfillopacity{0.500000}%
\pgfsetlinewidth{0.250937pt}%
\definecolor{currentstroke}{rgb}{0.000000,0.000000,0.000000}%
\pgfsetstrokecolor{currentstroke}%
\pgfsetstrokeopacity{0.500000}%
\pgfsetdash{}{0pt}%
\pgfsys@defobject{currentmarker}{\pgfqpoint{-0.013889in}{-0.013889in}}{\pgfqpoint{0.013889in}{0.013889in}}{%
\pgfpathmoveto{\pgfqpoint{0.000000in}{-0.013889in}}%
\pgfpathcurveto{\pgfqpoint{0.003683in}{-0.013889in}}{\pgfqpoint{0.007216in}{-0.012425in}}{\pgfqpoint{0.009821in}{-0.009821in}}%
\pgfpathcurveto{\pgfqpoint{0.012425in}{-0.007216in}}{\pgfqpoint{0.013889in}{-0.003683in}}{\pgfqpoint{0.013889in}{0.000000in}}%
\pgfpathcurveto{\pgfqpoint{0.013889in}{0.003683in}}{\pgfqpoint{0.012425in}{0.007216in}}{\pgfqpoint{0.009821in}{0.009821in}}%
\pgfpathcurveto{\pgfqpoint{0.007216in}{0.012425in}}{\pgfqpoint{0.003683in}{0.013889in}}{\pgfqpoint{0.000000in}{0.013889in}}%
\pgfpathcurveto{\pgfqpoint{-0.003683in}{0.013889in}}{\pgfqpoint{-0.007216in}{0.012425in}}{\pgfqpoint{-0.009821in}{0.009821in}}%
\pgfpathcurveto{\pgfqpoint{-0.012425in}{0.007216in}}{\pgfqpoint{-0.013889in}{0.003683in}}{\pgfqpoint{-0.013889in}{0.000000in}}%
\pgfpathcurveto{\pgfqpoint{-0.013889in}{-0.003683in}}{\pgfqpoint{-0.012425in}{-0.007216in}}{\pgfqpoint{-0.009821in}{-0.009821in}}%
\pgfpathcurveto{\pgfqpoint{-0.007216in}{-0.012425in}}{\pgfqpoint{-0.003683in}{-0.013889in}}{\pgfqpoint{0.000000in}{-0.013889in}}%
\pgfpathlineto{\pgfqpoint{0.000000in}{-0.013889in}}%
\pgfpathclose%
\pgfusepath{stroke,fill}%
}%
\begin{pgfscope}%
\pgfsys@transformshift{3.819255in}{4.813545in}%
\pgfsys@useobject{currentmarker}{}%
\end{pgfscope}%
\end{pgfscope}%
\begin{pgfscope}%
\pgfpathrectangle{\pgfqpoint{0.100000in}{2.493063in}}{\pgfqpoint{5.037500in}{3.427208in}}%
\pgfusepath{clip}%
\pgfsetrectcap%
\pgfsetroundjoin%
\pgfsetlinewidth{1.505625pt}%
\definecolor{currentstroke}{rgb}{0.000000,0.000000,1.000000}%
\pgfsetstrokecolor{currentstroke}%
\pgfsetstrokeopacity{0.500000}%
\pgfsetdash{}{0pt}%
\pgfpathmoveto{\pgfqpoint{3.753328in}{4.882843in}}%
\pgfusepath{stroke}%
\end{pgfscope}%
\begin{pgfscope}%
\pgfpathrectangle{\pgfqpoint{0.100000in}{2.493063in}}{\pgfqpoint{5.037500in}{3.427208in}}%
\pgfusepath{clip}%
\pgfsetbuttcap%
\pgfsetroundjoin%
\definecolor{currentfill}{rgb}{0.000000,0.000000,1.000000}%
\pgfsetfillcolor{currentfill}%
\pgfsetfillopacity{0.500000}%
\pgfsetlinewidth{0.250937pt}%
\definecolor{currentstroke}{rgb}{0.000000,0.000000,0.000000}%
\pgfsetstrokecolor{currentstroke}%
\pgfsetstrokeopacity{0.500000}%
\pgfsetdash{}{0pt}%
\pgfsys@defobject{currentmarker}{\pgfqpoint{-0.008333in}{-0.008333in}}{\pgfqpoint{0.008333in}{0.008333in}}{%
\pgfpathmoveto{\pgfqpoint{0.000000in}{-0.008333in}}%
\pgfpathcurveto{\pgfqpoint{0.002210in}{-0.008333in}}{\pgfqpoint{0.004330in}{-0.007455in}}{\pgfqpoint{0.005893in}{-0.005893in}}%
\pgfpathcurveto{\pgfqpoint{0.007455in}{-0.004330in}}{\pgfqpoint{0.008333in}{-0.002210in}}{\pgfqpoint{0.008333in}{0.000000in}}%
\pgfpathcurveto{\pgfqpoint{0.008333in}{0.002210in}}{\pgfqpoint{0.007455in}{0.004330in}}{\pgfqpoint{0.005893in}{0.005893in}}%
\pgfpathcurveto{\pgfqpoint{0.004330in}{0.007455in}}{\pgfqpoint{0.002210in}{0.008333in}}{\pgfqpoint{0.000000in}{0.008333in}}%
\pgfpathcurveto{\pgfqpoint{-0.002210in}{0.008333in}}{\pgfqpoint{-0.004330in}{0.007455in}}{\pgfqpoint{-0.005893in}{0.005893in}}%
\pgfpathcurveto{\pgfqpoint{-0.007455in}{0.004330in}}{\pgfqpoint{-0.008333in}{0.002210in}}{\pgfqpoint{-0.008333in}{0.000000in}}%
\pgfpathcurveto{\pgfqpoint{-0.008333in}{-0.002210in}}{\pgfqpoint{-0.007455in}{-0.004330in}}{\pgfqpoint{-0.005893in}{-0.005893in}}%
\pgfpathcurveto{\pgfqpoint{-0.004330in}{-0.007455in}}{\pgfqpoint{-0.002210in}{-0.008333in}}{\pgfqpoint{0.000000in}{-0.008333in}}%
\pgfpathlineto{\pgfqpoint{0.000000in}{-0.008333in}}%
\pgfpathclose%
\pgfusepath{stroke,fill}%
}%
\begin{pgfscope}%
\pgfsys@transformshift{3.753328in}{4.882843in}%
\pgfsys@useobject{currentmarker}{}%
\end{pgfscope}%
\end{pgfscope}%
\begin{pgfscope}%
\pgfpathrectangle{\pgfqpoint{0.100000in}{2.493063in}}{\pgfqpoint{5.037500in}{3.427208in}}%
\pgfusepath{clip}%
\pgfsetrectcap%
\pgfsetroundjoin%
\pgfsetlinewidth{1.505625pt}%
\definecolor{currentstroke}{rgb}{0.000000,0.000000,1.000000}%
\pgfsetstrokecolor{currentstroke}%
\pgfsetstrokeopacity{0.500000}%
\pgfsetdash{}{0pt}%
\pgfpathmoveto{\pgfqpoint{3.590943in}{4.858086in}}%
\pgfusepath{stroke}%
\end{pgfscope}%
\begin{pgfscope}%
\pgfpathrectangle{\pgfqpoint{0.100000in}{2.493063in}}{\pgfqpoint{5.037500in}{3.427208in}}%
\pgfusepath{clip}%
\pgfsetbuttcap%
\pgfsetroundjoin%
\definecolor{currentfill}{rgb}{0.000000,0.000000,1.000000}%
\pgfsetfillcolor{currentfill}%
\pgfsetfillopacity{0.500000}%
\pgfsetlinewidth{0.250937pt}%
\definecolor{currentstroke}{rgb}{0.000000,0.000000,0.000000}%
\pgfsetstrokecolor{currentstroke}%
\pgfsetstrokeopacity{0.500000}%
\pgfsetdash{}{0pt}%
\pgfsys@defobject{currentmarker}{\pgfqpoint{-0.011111in}{-0.011111in}}{\pgfqpoint{0.011111in}{0.011111in}}{%
\pgfpathmoveto{\pgfqpoint{0.000000in}{-0.011111in}}%
\pgfpathcurveto{\pgfqpoint{0.002947in}{-0.011111in}}{\pgfqpoint{0.005773in}{-0.009940in}}{\pgfqpoint{0.007857in}{-0.007857in}}%
\pgfpathcurveto{\pgfqpoint{0.009940in}{-0.005773in}}{\pgfqpoint{0.011111in}{-0.002947in}}{\pgfqpoint{0.011111in}{0.000000in}}%
\pgfpathcurveto{\pgfqpoint{0.011111in}{0.002947in}}{\pgfqpoint{0.009940in}{0.005773in}}{\pgfqpoint{0.007857in}{0.007857in}}%
\pgfpathcurveto{\pgfqpoint{0.005773in}{0.009940in}}{\pgfqpoint{0.002947in}{0.011111in}}{\pgfqpoint{0.000000in}{0.011111in}}%
\pgfpathcurveto{\pgfqpoint{-0.002947in}{0.011111in}}{\pgfqpoint{-0.005773in}{0.009940in}}{\pgfqpoint{-0.007857in}{0.007857in}}%
\pgfpathcurveto{\pgfqpoint{-0.009940in}{0.005773in}}{\pgfqpoint{-0.011111in}{0.002947in}}{\pgfqpoint{-0.011111in}{0.000000in}}%
\pgfpathcurveto{\pgfqpoint{-0.011111in}{-0.002947in}}{\pgfqpoint{-0.009940in}{-0.005773in}}{\pgfqpoint{-0.007857in}{-0.007857in}}%
\pgfpathcurveto{\pgfqpoint{-0.005773in}{-0.009940in}}{\pgfqpoint{-0.002947in}{-0.011111in}}{\pgfqpoint{0.000000in}{-0.011111in}}%
\pgfpathlineto{\pgfqpoint{0.000000in}{-0.011111in}}%
\pgfpathclose%
\pgfusepath{stroke,fill}%
}%
\begin{pgfscope}%
\pgfsys@transformshift{3.590943in}{4.858086in}%
\pgfsys@useobject{currentmarker}{}%
\end{pgfscope}%
\end{pgfscope}%
\begin{pgfscope}%
\pgfpathrectangle{\pgfqpoint{0.100000in}{2.493063in}}{\pgfqpoint{5.037500in}{3.427208in}}%
\pgfusepath{clip}%
\pgfsetrectcap%
\pgfsetroundjoin%
\pgfsetlinewidth{1.505625pt}%
\definecolor{currentstroke}{rgb}{0.000000,0.000000,1.000000}%
\pgfsetstrokecolor{currentstroke}%
\pgfsetstrokeopacity{0.500000}%
\pgfsetdash{}{0pt}%
\pgfpathmoveto{\pgfqpoint{3.706510in}{4.787536in}}%
\pgfusepath{stroke}%
\end{pgfscope}%
\begin{pgfscope}%
\pgfpathrectangle{\pgfqpoint{0.100000in}{2.493063in}}{\pgfqpoint{5.037500in}{3.427208in}}%
\pgfusepath{clip}%
\pgfsetbuttcap%
\pgfsetroundjoin%
\definecolor{currentfill}{rgb}{0.000000,0.000000,1.000000}%
\pgfsetfillcolor{currentfill}%
\pgfsetfillopacity{0.500000}%
\pgfsetlinewidth{0.250937pt}%
\definecolor{currentstroke}{rgb}{0.000000,0.000000,0.000000}%
\pgfsetstrokecolor{currentstroke}%
\pgfsetstrokeopacity{0.500000}%
\pgfsetdash{}{0pt}%
\pgfsys@defobject{currentmarker}{\pgfqpoint{-0.008333in}{-0.008333in}}{\pgfqpoint{0.008333in}{0.008333in}}{%
\pgfpathmoveto{\pgfqpoint{0.000000in}{-0.008333in}}%
\pgfpathcurveto{\pgfqpoint{0.002210in}{-0.008333in}}{\pgfqpoint{0.004330in}{-0.007455in}}{\pgfqpoint{0.005893in}{-0.005893in}}%
\pgfpathcurveto{\pgfqpoint{0.007455in}{-0.004330in}}{\pgfqpoint{0.008333in}{-0.002210in}}{\pgfqpoint{0.008333in}{0.000000in}}%
\pgfpathcurveto{\pgfqpoint{0.008333in}{0.002210in}}{\pgfqpoint{0.007455in}{0.004330in}}{\pgfqpoint{0.005893in}{0.005893in}}%
\pgfpathcurveto{\pgfqpoint{0.004330in}{0.007455in}}{\pgfqpoint{0.002210in}{0.008333in}}{\pgfqpoint{0.000000in}{0.008333in}}%
\pgfpathcurveto{\pgfqpoint{-0.002210in}{0.008333in}}{\pgfqpoint{-0.004330in}{0.007455in}}{\pgfqpoint{-0.005893in}{0.005893in}}%
\pgfpathcurveto{\pgfqpoint{-0.007455in}{0.004330in}}{\pgfqpoint{-0.008333in}{0.002210in}}{\pgfqpoint{-0.008333in}{0.000000in}}%
\pgfpathcurveto{\pgfqpoint{-0.008333in}{-0.002210in}}{\pgfqpoint{-0.007455in}{-0.004330in}}{\pgfqpoint{-0.005893in}{-0.005893in}}%
\pgfpathcurveto{\pgfqpoint{-0.004330in}{-0.007455in}}{\pgfqpoint{-0.002210in}{-0.008333in}}{\pgfqpoint{0.000000in}{-0.008333in}}%
\pgfpathlineto{\pgfqpoint{0.000000in}{-0.008333in}}%
\pgfpathclose%
\pgfusepath{stroke,fill}%
}%
\begin{pgfscope}%
\pgfsys@transformshift{3.706510in}{4.787536in}%
\pgfsys@useobject{currentmarker}{}%
\end{pgfscope}%
\end{pgfscope}%
\begin{pgfscope}%
\pgfpathrectangle{\pgfqpoint{0.100000in}{2.493063in}}{\pgfqpoint{5.037500in}{3.427208in}}%
\pgfusepath{clip}%
\pgfsetrectcap%
\pgfsetroundjoin%
\pgfsetlinewidth{1.505625pt}%
\definecolor{currentstroke}{rgb}{0.000000,0.000000,1.000000}%
\pgfsetstrokecolor{currentstroke}%
\pgfsetstrokeopacity{0.500000}%
\pgfsetdash{}{0pt}%
\pgfpathmoveto{\pgfqpoint{3.605727in}{4.781726in}}%
\pgfusepath{stroke}%
\end{pgfscope}%
\begin{pgfscope}%
\pgfpathrectangle{\pgfqpoint{0.100000in}{2.493063in}}{\pgfqpoint{5.037500in}{3.427208in}}%
\pgfusepath{clip}%
\pgfsetbuttcap%
\pgfsetroundjoin%
\definecolor{currentfill}{rgb}{0.000000,0.000000,1.000000}%
\pgfsetfillcolor{currentfill}%
\pgfsetfillopacity{0.500000}%
\pgfsetlinewidth{0.250937pt}%
\definecolor{currentstroke}{rgb}{0.000000,0.000000,0.000000}%
\pgfsetstrokecolor{currentstroke}%
\pgfsetstrokeopacity{0.500000}%
\pgfsetdash{}{0pt}%
\pgfsys@defobject{currentmarker}{\pgfqpoint{-0.011111in}{-0.011111in}}{\pgfqpoint{0.011111in}{0.011111in}}{%
\pgfpathmoveto{\pgfqpoint{0.000000in}{-0.011111in}}%
\pgfpathcurveto{\pgfqpoint{0.002947in}{-0.011111in}}{\pgfqpoint{0.005773in}{-0.009940in}}{\pgfqpoint{0.007857in}{-0.007857in}}%
\pgfpathcurveto{\pgfqpoint{0.009940in}{-0.005773in}}{\pgfqpoint{0.011111in}{-0.002947in}}{\pgfqpoint{0.011111in}{0.000000in}}%
\pgfpathcurveto{\pgfqpoint{0.011111in}{0.002947in}}{\pgfqpoint{0.009940in}{0.005773in}}{\pgfqpoint{0.007857in}{0.007857in}}%
\pgfpathcurveto{\pgfqpoint{0.005773in}{0.009940in}}{\pgfqpoint{0.002947in}{0.011111in}}{\pgfqpoint{0.000000in}{0.011111in}}%
\pgfpathcurveto{\pgfqpoint{-0.002947in}{0.011111in}}{\pgfqpoint{-0.005773in}{0.009940in}}{\pgfqpoint{-0.007857in}{0.007857in}}%
\pgfpathcurveto{\pgfqpoint{-0.009940in}{0.005773in}}{\pgfqpoint{-0.011111in}{0.002947in}}{\pgfqpoint{-0.011111in}{0.000000in}}%
\pgfpathcurveto{\pgfqpoint{-0.011111in}{-0.002947in}}{\pgfqpoint{-0.009940in}{-0.005773in}}{\pgfqpoint{-0.007857in}{-0.007857in}}%
\pgfpathcurveto{\pgfqpoint{-0.005773in}{-0.009940in}}{\pgfqpoint{-0.002947in}{-0.011111in}}{\pgfqpoint{0.000000in}{-0.011111in}}%
\pgfpathlineto{\pgfqpoint{0.000000in}{-0.011111in}}%
\pgfpathclose%
\pgfusepath{stroke,fill}%
}%
\begin{pgfscope}%
\pgfsys@transformshift{3.605727in}{4.781726in}%
\pgfsys@useobject{currentmarker}{}%
\end{pgfscope}%
\end{pgfscope}%
\begin{pgfscope}%
\pgfpathrectangle{\pgfqpoint{0.100000in}{2.493063in}}{\pgfqpoint{5.037500in}{3.427208in}}%
\pgfusepath{clip}%
\pgfsetrectcap%
\pgfsetroundjoin%
\pgfsetlinewidth{1.505625pt}%
\definecolor{currentstroke}{rgb}{0.000000,0.000000,1.000000}%
\pgfsetstrokecolor{currentstroke}%
\pgfsetstrokeopacity{0.500000}%
\pgfsetdash{}{0pt}%
\pgfpathmoveto{\pgfqpoint{3.687384in}{4.841983in}}%
\pgfusepath{stroke}%
\end{pgfscope}%
\begin{pgfscope}%
\pgfpathrectangle{\pgfqpoint{0.100000in}{2.493063in}}{\pgfqpoint{5.037500in}{3.427208in}}%
\pgfusepath{clip}%
\pgfsetbuttcap%
\pgfsetroundjoin%
\definecolor{currentfill}{rgb}{0.000000,0.000000,1.000000}%
\pgfsetfillcolor{currentfill}%
\pgfsetfillopacity{0.500000}%
\pgfsetlinewidth{0.250937pt}%
\definecolor{currentstroke}{rgb}{0.000000,0.000000,0.000000}%
\pgfsetstrokecolor{currentstroke}%
\pgfsetstrokeopacity{0.500000}%
\pgfsetdash{}{0pt}%
\pgfsys@defobject{currentmarker}{\pgfqpoint{-0.013889in}{-0.013889in}}{\pgfqpoint{0.013889in}{0.013889in}}{%
\pgfpathmoveto{\pgfqpoint{0.000000in}{-0.013889in}}%
\pgfpathcurveto{\pgfqpoint{0.003683in}{-0.013889in}}{\pgfqpoint{0.007216in}{-0.012425in}}{\pgfqpoint{0.009821in}{-0.009821in}}%
\pgfpathcurveto{\pgfqpoint{0.012425in}{-0.007216in}}{\pgfqpoint{0.013889in}{-0.003683in}}{\pgfqpoint{0.013889in}{0.000000in}}%
\pgfpathcurveto{\pgfqpoint{0.013889in}{0.003683in}}{\pgfqpoint{0.012425in}{0.007216in}}{\pgfqpoint{0.009821in}{0.009821in}}%
\pgfpathcurveto{\pgfqpoint{0.007216in}{0.012425in}}{\pgfqpoint{0.003683in}{0.013889in}}{\pgfqpoint{0.000000in}{0.013889in}}%
\pgfpathcurveto{\pgfqpoint{-0.003683in}{0.013889in}}{\pgfqpoint{-0.007216in}{0.012425in}}{\pgfqpoint{-0.009821in}{0.009821in}}%
\pgfpathcurveto{\pgfqpoint{-0.012425in}{0.007216in}}{\pgfqpoint{-0.013889in}{0.003683in}}{\pgfqpoint{-0.013889in}{0.000000in}}%
\pgfpathcurveto{\pgfqpoint{-0.013889in}{-0.003683in}}{\pgfqpoint{-0.012425in}{-0.007216in}}{\pgfqpoint{-0.009821in}{-0.009821in}}%
\pgfpathcurveto{\pgfqpoint{-0.007216in}{-0.012425in}}{\pgfqpoint{-0.003683in}{-0.013889in}}{\pgfqpoint{0.000000in}{-0.013889in}}%
\pgfpathlineto{\pgfqpoint{0.000000in}{-0.013889in}}%
\pgfpathclose%
\pgfusepath{stroke,fill}%
}%
\begin{pgfscope}%
\pgfsys@transformshift{3.687384in}{4.841983in}%
\pgfsys@useobject{currentmarker}{}%
\end{pgfscope}%
\end{pgfscope}%
\begin{pgfscope}%
\pgfpathrectangle{\pgfqpoint{0.100000in}{2.493063in}}{\pgfqpoint{5.037500in}{3.427208in}}%
\pgfusepath{clip}%
\pgfsetrectcap%
\pgfsetroundjoin%
\pgfsetlinewidth{1.505625pt}%
\definecolor{currentstroke}{rgb}{0.501961,0.501961,0.501961}%
\pgfsetstrokecolor{currentstroke}%
\pgfsetstrokeopacity{0.500000}%
\pgfsetdash{}{0pt}%
\pgfpathmoveto{\pgfqpoint{3.701307in}{4.946114in}}%
\pgfusepath{stroke}%
\end{pgfscope}%
\begin{pgfscope}%
\pgfpathrectangle{\pgfqpoint{0.100000in}{2.493063in}}{\pgfqpoint{5.037500in}{3.427208in}}%
\pgfusepath{clip}%
\pgfsetbuttcap%
\pgfsetroundjoin%
\definecolor{currentfill}{rgb}{0.501961,0.501961,0.501961}%
\pgfsetfillcolor{currentfill}%
\pgfsetfillopacity{0.500000}%
\pgfsetlinewidth{0.250937pt}%
\definecolor{currentstroke}{rgb}{0.000000,0.000000,0.000000}%
\pgfsetstrokecolor{currentstroke}%
\pgfsetstrokeopacity{0.500000}%
\pgfsetdash{}{0pt}%
\pgfsys@defobject{currentmarker}{\pgfqpoint{-0.013889in}{-0.013889in}}{\pgfqpoint{0.013889in}{0.013889in}}{%
\pgfpathmoveto{\pgfqpoint{0.000000in}{-0.013889in}}%
\pgfpathcurveto{\pgfqpoint{0.003683in}{-0.013889in}}{\pgfqpoint{0.007216in}{-0.012425in}}{\pgfqpoint{0.009821in}{-0.009821in}}%
\pgfpathcurveto{\pgfqpoint{0.012425in}{-0.007216in}}{\pgfqpoint{0.013889in}{-0.003683in}}{\pgfqpoint{0.013889in}{0.000000in}}%
\pgfpathcurveto{\pgfqpoint{0.013889in}{0.003683in}}{\pgfqpoint{0.012425in}{0.007216in}}{\pgfqpoint{0.009821in}{0.009821in}}%
\pgfpathcurveto{\pgfqpoint{0.007216in}{0.012425in}}{\pgfqpoint{0.003683in}{0.013889in}}{\pgfqpoint{0.000000in}{0.013889in}}%
\pgfpathcurveto{\pgfqpoint{-0.003683in}{0.013889in}}{\pgfqpoint{-0.007216in}{0.012425in}}{\pgfqpoint{-0.009821in}{0.009821in}}%
\pgfpathcurveto{\pgfqpoint{-0.012425in}{0.007216in}}{\pgfqpoint{-0.013889in}{0.003683in}}{\pgfqpoint{-0.013889in}{0.000000in}}%
\pgfpathcurveto{\pgfqpoint{-0.013889in}{-0.003683in}}{\pgfqpoint{-0.012425in}{-0.007216in}}{\pgfqpoint{-0.009821in}{-0.009821in}}%
\pgfpathcurveto{\pgfqpoint{-0.007216in}{-0.012425in}}{\pgfqpoint{-0.003683in}{-0.013889in}}{\pgfqpoint{0.000000in}{-0.013889in}}%
\pgfpathlineto{\pgfqpoint{0.000000in}{-0.013889in}}%
\pgfpathclose%
\pgfusepath{stroke,fill}%
}%
\begin{pgfscope}%
\pgfsys@transformshift{3.701307in}{4.946114in}%
\pgfsys@useobject{currentmarker}{}%
\end{pgfscope}%
\end{pgfscope}%
\begin{pgfscope}%
\pgfpathrectangle{\pgfqpoint{0.100000in}{2.493063in}}{\pgfqpoint{5.037500in}{3.427208in}}%
\pgfusepath{clip}%
\pgfsetrectcap%
\pgfsetroundjoin%
\pgfsetlinewidth{1.505625pt}%
\definecolor{currentstroke}{rgb}{0.000000,0.000000,1.000000}%
\pgfsetstrokecolor{currentstroke}%
\pgfsetstrokeopacity{0.500000}%
\pgfsetdash{}{0pt}%
\pgfpathmoveto{\pgfqpoint{3.796884in}{4.760144in}}%
\pgfusepath{stroke}%
\end{pgfscope}%
\begin{pgfscope}%
\pgfpathrectangle{\pgfqpoint{0.100000in}{2.493063in}}{\pgfqpoint{5.037500in}{3.427208in}}%
\pgfusepath{clip}%
\pgfsetbuttcap%
\pgfsetroundjoin%
\definecolor{currentfill}{rgb}{0.000000,0.000000,1.000000}%
\pgfsetfillcolor{currentfill}%
\pgfsetfillopacity{0.500000}%
\pgfsetlinewidth{0.250937pt}%
\definecolor{currentstroke}{rgb}{0.000000,0.000000,0.000000}%
\pgfsetstrokecolor{currentstroke}%
\pgfsetstrokeopacity{0.500000}%
\pgfsetdash{}{0pt}%
\pgfsys@defobject{currentmarker}{\pgfqpoint{-0.025000in}{-0.025000in}}{\pgfqpoint{0.025000in}{0.025000in}}{%
\pgfpathmoveto{\pgfqpoint{0.000000in}{-0.025000in}}%
\pgfpathcurveto{\pgfqpoint{0.006630in}{-0.025000in}}{\pgfqpoint{0.012989in}{-0.022366in}}{\pgfqpoint{0.017678in}{-0.017678in}}%
\pgfpathcurveto{\pgfqpoint{0.022366in}{-0.012989in}}{\pgfqpoint{0.025000in}{-0.006630in}}{\pgfqpoint{0.025000in}{0.000000in}}%
\pgfpathcurveto{\pgfqpoint{0.025000in}{0.006630in}}{\pgfqpoint{0.022366in}{0.012989in}}{\pgfqpoint{0.017678in}{0.017678in}}%
\pgfpathcurveto{\pgfqpoint{0.012989in}{0.022366in}}{\pgfqpoint{0.006630in}{0.025000in}}{\pgfqpoint{0.000000in}{0.025000in}}%
\pgfpathcurveto{\pgfqpoint{-0.006630in}{0.025000in}}{\pgfqpoint{-0.012989in}{0.022366in}}{\pgfqpoint{-0.017678in}{0.017678in}}%
\pgfpathcurveto{\pgfqpoint{-0.022366in}{0.012989in}}{\pgfqpoint{-0.025000in}{0.006630in}}{\pgfqpoint{-0.025000in}{0.000000in}}%
\pgfpathcurveto{\pgfqpoint{-0.025000in}{-0.006630in}}{\pgfqpoint{-0.022366in}{-0.012989in}}{\pgfqpoint{-0.017678in}{-0.017678in}}%
\pgfpathcurveto{\pgfqpoint{-0.012989in}{-0.022366in}}{\pgfqpoint{-0.006630in}{-0.025000in}}{\pgfqpoint{0.000000in}{-0.025000in}}%
\pgfpathlineto{\pgfqpoint{0.000000in}{-0.025000in}}%
\pgfpathclose%
\pgfusepath{stroke,fill}%
}%
\begin{pgfscope}%
\pgfsys@transformshift{3.796884in}{4.760144in}%
\pgfsys@useobject{currentmarker}{}%
\end{pgfscope}%
\end{pgfscope}%
\begin{pgfscope}%
\pgfpathrectangle{\pgfqpoint{0.100000in}{2.493063in}}{\pgfqpoint{5.037500in}{3.427208in}}%
\pgfusepath{clip}%
\pgfsetrectcap%
\pgfsetroundjoin%
\pgfsetlinewidth{1.505625pt}%
\definecolor{currentstroke}{rgb}{0.501961,0.501961,0.501961}%
\pgfsetstrokecolor{currentstroke}%
\pgfsetstrokeopacity{0.500000}%
\pgfsetdash{}{0pt}%
\pgfpathmoveto{\pgfqpoint{3.539146in}{4.884347in}}%
\pgfusepath{stroke}%
\end{pgfscope}%
\begin{pgfscope}%
\pgfpathrectangle{\pgfqpoint{0.100000in}{2.493063in}}{\pgfqpoint{5.037500in}{3.427208in}}%
\pgfusepath{clip}%
\pgfsetbuttcap%
\pgfsetroundjoin%
\definecolor{currentfill}{rgb}{0.501961,0.501961,0.501961}%
\pgfsetfillcolor{currentfill}%
\pgfsetfillopacity{0.500000}%
\pgfsetlinewidth{0.250937pt}%
\definecolor{currentstroke}{rgb}{0.000000,0.000000,0.000000}%
\pgfsetstrokecolor{currentstroke}%
\pgfsetstrokeopacity{0.500000}%
\pgfsetdash{}{0pt}%
\pgfsys@defobject{currentmarker}{\pgfqpoint{-0.013889in}{-0.013889in}}{\pgfqpoint{0.013889in}{0.013889in}}{%
\pgfpathmoveto{\pgfqpoint{0.000000in}{-0.013889in}}%
\pgfpathcurveto{\pgfqpoint{0.003683in}{-0.013889in}}{\pgfqpoint{0.007216in}{-0.012425in}}{\pgfqpoint{0.009821in}{-0.009821in}}%
\pgfpathcurveto{\pgfqpoint{0.012425in}{-0.007216in}}{\pgfqpoint{0.013889in}{-0.003683in}}{\pgfqpoint{0.013889in}{0.000000in}}%
\pgfpathcurveto{\pgfqpoint{0.013889in}{0.003683in}}{\pgfqpoint{0.012425in}{0.007216in}}{\pgfqpoint{0.009821in}{0.009821in}}%
\pgfpathcurveto{\pgfqpoint{0.007216in}{0.012425in}}{\pgfqpoint{0.003683in}{0.013889in}}{\pgfqpoint{0.000000in}{0.013889in}}%
\pgfpathcurveto{\pgfqpoint{-0.003683in}{0.013889in}}{\pgfqpoint{-0.007216in}{0.012425in}}{\pgfqpoint{-0.009821in}{0.009821in}}%
\pgfpathcurveto{\pgfqpoint{-0.012425in}{0.007216in}}{\pgfqpoint{-0.013889in}{0.003683in}}{\pgfqpoint{-0.013889in}{0.000000in}}%
\pgfpathcurveto{\pgfqpoint{-0.013889in}{-0.003683in}}{\pgfqpoint{-0.012425in}{-0.007216in}}{\pgfqpoint{-0.009821in}{-0.009821in}}%
\pgfpathcurveto{\pgfqpoint{-0.007216in}{-0.012425in}}{\pgfqpoint{-0.003683in}{-0.013889in}}{\pgfqpoint{0.000000in}{-0.013889in}}%
\pgfpathlineto{\pgfqpoint{0.000000in}{-0.013889in}}%
\pgfpathclose%
\pgfusepath{stroke,fill}%
}%
\begin{pgfscope}%
\pgfsys@transformshift{3.539146in}{4.884347in}%
\pgfsys@useobject{currentmarker}{}%
\end{pgfscope}%
\end{pgfscope}%
\begin{pgfscope}%
\pgfpathrectangle{\pgfqpoint{0.100000in}{2.493063in}}{\pgfqpoint{5.037500in}{3.427208in}}%
\pgfusepath{clip}%
\pgfsetrectcap%
\pgfsetroundjoin%
\pgfsetlinewidth{1.505625pt}%
\definecolor{currentstroke}{rgb}{0.000000,0.000000,1.000000}%
\pgfsetstrokecolor{currentstroke}%
\pgfsetstrokeopacity{0.500000}%
\pgfsetdash{}{0pt}%
\pgfpathmoveto{\pgfqpoint{3.554230in}{4.723033in}}%
\pgfusepath{stroke}%
\end{pgfscope}%
\begin{pgfscope}%
\pgfpathrectangle{\pgfqpoint{0.100000in}{2.493063in}}{\pgfqpoint{5.037500in}{3.427208in}}%
\pgfusepath{clip}%
\pgfsetbuttcap%
\pgfsetroundjoin%
\definecolor{currentfill}{rgb}{0.000000,0.000000,1.000000}%
\pgfsetfillcolor{currentfill}%
\pgfsetfillopacity{0.500000}%
\pgfsetlinewidth{0.250937pt}%
\definecolor{currentstroke}{rgb}{0.000000,0.000000,0.000000}%
\pgfsetstrokecolor{currentstroke}%
\pgfsetstrokeopacity{0.500000}%
\pgfsetdash{}{0pt}%
\pgfsys@defobject{currentmarker}{\pgfqpoint{-0.011111in}{-0.011111in}}{\pgfqpoint{0.011111in}{0.011111in}}{%
\pgfpathmoveto{\pgfqpoint{0.000000in}{-0.011111in}}%
\pgfpathcurveto{\pgfqpoint{0.002947in}{-0.011111in}}{\pgfqpoint{0.005773in}{-0.009940in}}{\pgfqpoint{0.007857in}{-0.007857in}}%
\pgfpathcurveto{\pgfqpoint{0.009940in}{-0.005773in}}{\pgfqpoint{0.011111in}{-0.002947in}}{\pgfqpoint{0.011111in}{0.000000in}}%
\pgfpathcurveto{\pgfqpoint{0.011111in}{0.002947in}}{\pgfqpoint{0.009940in}{0.005773in}}{\pgfqpoint{0.007857in}{0.007857in}}%
\pgfpathcurveto{\pgfqpoint{0.005773in}{0.009940in}}{\pgfqpoint{0.002947in}{0.011111in}}{\pgfqpoint{0.000000in}{0.011111in}}%
\pgfpathcurveto{\pgfqpoint{-0.002947in}{0.011111in}}{\pgfqpoint{-0.005773in}{0.009940in}}{\pgfqpoint{-0.007857in}{0.007857in}}%
\pgfpathcurveto{\pgfqpoint{-0.009940in}{0.005773in}}{\pgfqpoint{-0.011111in}{0.002947in}}{\pgfqpoint{-0.011111in}{0.000000in}}%
\pgfpathcurveto{\pgfqpoint{-0.011111in}{-0.002947in}}{\pgfqpoint{-0.009940in}{-0.005773in}}{\pgfqpoint{-0.007857in}{-0.007857in}}%
\pgfpathcurveto{\pgfqpoint{-0.005773in}{-0.009940in}}{\pgfqpoint{-0.002947in}{-0.011111in}}{\pgfqpoint{0.000000in}{-0.011111in}}%
\pgfpathlineto{\pgfqpoint{0.000000in}{-0.011111in}}%
\pgfpathclose%
\pgfusepath{stroke,fill}%
}%
\begin{pgfscope}%
\pgfsys@transformshift{3.554230in}{4.723033in}%
\pgfsys@useobject{currentmarker}{}%
\end{pgfscope}%
\end{pgfscope}%
\begin{pgfscope}%
\pgfpathrectangle{\pgfqpoint{0.100000in}{2.493063in}}{\pgfqpoint{5.037500in}{3.427208in}}%
\pgfusepath{clip}%
\pgfsetrectcap%
\pgfsetroundjoin%
\pgfsetlinewidth{1.505625pt}%
\definecolor{currentstroke}{rgb}{0.000000,0.000000,1.000000}%
\pgfsetstrokecolor{currentstroke}%
\pgfsetstrokeopacity{0.500000}%
\pgfsetdash{}{0pt}%
\pgfpathmoveto{\pgfqpoint{3.728805in}{4.926680in}}%
\pgfusepath{stroke}%
\end{pgfscope}%
\begin{pgfscope}%
\pgfpathrectangle{\pgfqpoint{0.100000in}{2.493063in}}{\pgfqpoint{5.037500in}{3.427208in}}%
\pgfusepath{clip}%
\pgfsetbuttcap%
\pgfsetroundjoin%
\definecolor{currentfill}{rgb}{0.000000,0.000000,1.000000}%
\pgfsetfillcolor{currentfill}%
\pgfsetfillopacity{0.500000}%
\pgfsetlinewidth{0.250937pt}%
\definecolor{currentstroke}{rgb}{0.000000,0.000000,0.000000}%
\pgfsetstrokecolor{currentstroke}%
\pgfsetstrokeopacity{0.500000}%
\pgfsetdash{}{0pt}%
\pgfsys@defobject{currentmarker}{\pgfqpoint{-0.008333in}{-0.008333in}}{\pgfqpoint{0.008333in}{0.008333in}}{%
\pgfpathmoveto{\pgfqpoint{0.000000in}{-0.008333in}}%
\pgfpathcurveto{\pgfqpoint{0.002210in}{-0.008333in}}{\pgfqpoint{0.004330in}{-0.007455in}}{\pgfqpoint{0.005893in}{-0.005893in}}%
\pgfpathcurveto{\pgfqpoint{0.007455in}{-0.004330in}}{\pgfqpoint{0.008333in}{-0.002210in}}{\pgfqpoint{0.008333in}{0.000000in}}%
\pgfpathcurveto{\pgfqpoint{0.008333in}{0.002210in}}{\pgfqpoint{0.007455in}{0.004330in}}{\pgfqpoint{0.005893in}{0.005893in}}%
\pgfpathcurveto{\pgfqpoint{0.004330in}{0.007455in}}{\pgfqpoint{0.002210in}{0.008333in}}{\pgfqpoint{0.000000in}{0.008333in}}%
\pgfpathcurveto{\pgfqpoint{-0.002210in}{0.008333in}}{\pgfqpoint{-0.004330in}{0.007455in}}{\pgfqpoint{-0.005893in}{0.005893in}}%
\pgfpathcurveto{\pgfqpoint{-0.007455in}{0.004330in}}{\pgfqpoint{-0.008333in}{0.002210in}}{\pgfqpoint{-0.008333in}{0.000000in}}%
\pgfpathcurveto{\pgfqpoint{-0.008333in}{-0.002210in}}{\pgfqpoint{-0.007455in}{-0.004330in}}{\pgfqpoint{-0.005893in}{-0.005893in}}%
\pgfpathcurveto{\pgfqpoint{-0.004330in}{-0.007455in}}{\pgfqpoint{-0.002210in}{-0.008333in}}{\pgfqpoint{0.000000in}{-0.008333in}}%
\pgfpathlineto{\pgfqpoint{0.000000in}{-0.008333in}}%
\pgfpathclose%
\pgfusepath{stroke,fill}%
}%
\begin{pgfscope}%
\pgfsys@transformshift{3.728805in}{4.926680in}%
\pgfsys@useobject{currentmarker}{}%
\end{pgfscope}%
\end{pgfscope}%
\begin{pgfscope}%
\pgfpathrectangle{\pgfqpoint{0.100000in}{2.493063in}}{\pgfqpoint{5.037500in}{3.427208in}}%
\pgfusepath{clip}%
\pgfsetrectcap%
\pgfsetroundjoin%
\pgfsetlinewidth{1.505625pt}%
\definecolor{currentstroke}{rgb}{0.678431,1.000000,0.184314}%
\pgfsetstrokecolor{currentstroke}%
\pgfsetstrokeopacity{0.500000}%
\pgfsetdash{}{0pt}%
\pgfpathmoveto{\pgfqpoint{3.032951in}{5.262346in}}%
\pgfusepath{stroke}%
\end{pgfscope}%
\begin{pgfscope}%
\pgfpathrectangle{\pgfqpoint{0.100000in}{2.493063in}}{\pgfqpoint{5.037500in}{3.427208in}}%
\pgfusepath{clip}%
\pgfsetbuttcap%
\pgfsetroundjoin%
\definecolor{currentfill}{rgb}{0.678431,1.000000,0.184314}%
\pgfsetfillcolor{currentfill}%
\pgfsetfillopacity{0.500000}%
\pgfsetlinewidth{0.250937pt}%
\definecolor{currentstroke}{rgb}{0.000000,0.000000,0.000000}%
\pgfsetstrokecolor{currentstroke}%
\pgfsetstrokeopacity{0.500000}%
\pgfsetdash{}{0pt}%
\pgfsys@defobject{currentmarker}{\pgfqpoint{-0.036111in}{-0.036111in}}{\pgfqpoint{0.036111in}{0.036111in}}{%
\pgfpathmoveto{\pgfqpoint{0.000000in}{-0.036111in}}%
\pgfpathcurveto{\pgfqpoint{0.009577in}{-0.036111in}}{\pgfqpoint{0.018763in}{-0.032306in}}{\pgfqpoint{0.025534in}{-0.025534in}}%
\pgfpathcurveto{\pgfqpoint{0.032306in}{-0.018763in}}{\pgfqpoint{0.036111in}{-0.009577in}}{\pgfqpoint{0.036111in}{0.000000in}}%
\pgfpathcurveto{\pgfqpoint{0.036111in}{0.009577in}}{\pgfqpoint{0.032306in}{0.018763in}}{\pgfqpoint{0.025534in}{0.025534in}}%
\pgfpathcurveto{\pgfqpoint{0.018763in}{0.032306in}}{\pgfqpoint{0.009577in}{0.036111in}}{\pgfqpoint{0.000000in}{0.036111in}}%
\pgfpathcurveto{\pgfqpoint{-0.009577in}{0.036111in}}{\pgfqpoint{-0.018763in}{0.032306in}}{\pgfqpoint{-0.025534in}{0.025534in}}%
\pgfpathcurveto{\pgfqpoint{-0.032306in}{0.018763in}}{\pgfqpoint{-0.036111in}{0.009577in}}{\pgfqpoint{-0.036111in}{0.000000in}}%
\pgfpathcurveto{\pgfqpoint{-0.036111in}{-0.009577in}}{\pgfqpoint{-0.032306in}{-0.018763in}}{\pgfqpoint{-0.025534in}{-0.025534in}}%
\pgfpathcurveto{\pgfqpoint{-0.018763in}{-0.032306in}}{\pgfqpoint{-0.009577in}{-0.036111in}}{\pgfqpoint{0.000000in}{-0.036111in}}%
\pgfpathlineto{\pgfqpoint{0.000000in}{-0.036111in}}%
\pgfpathclose%
\pgfusepath{stroke,fill}%
}%
\begin{pgfscope}%
\pgfsys@transformshift{3.032951in}{5.262346in}%
\pgfsys@useobject{currentmarker}{}%
\end{pgfscope}%
\end{pgfscope}%
\begin{pgfscope}%
\pgfpathrectangle{\pgfqpoint{0.100000in}{2.493063in}}{\pgfqpoint{5.037500in}{3.427208in}}%
\pgfusepath{clip}%
\pgfsetrectcap%
\pgfsetroundjoin%
\pgfsetlinewidth{1.505625pt}%
\definecolor{currentstroke}{rgb}{0.678431,1.000000,0.184314}%
\pgfsetstrokecolor{currentstroke}%
\pgfsetstrokeopacity{0.500000}%
\pgfsetdash{}{0pt}%
\pgfpathmoveto{\pgfqpoint{2.887116in}{4.956746in}}%
\pgfusepath{stroke}%
\end{pgfscope}%
\begin{pgfscope}%
\pgfpathrectangle{\pgfqpoint{0.100000in}{2.493063in}}{\pgfqpoint{5.037500in}{3.427208in}}%
\pgfusepath{clip}%
\pgfsetbuttcap%
\pgfsetroundjoin%
\definecolor{currentfill}{rgb}{0.678431,1.000000,0.184314}%
\pgfsetfillcolor{currentfill}%
\pgfsetfillopacity{0.500000}%
\pgfsetlinewidth{0.250937pt}%
\definecolor{currentstroke}{rgb}{0.000000,0.000000,0.000000}%
\pgfsetstrokecolor{currentstroke}%
\pgfsetstrokeopacity{0.500000}%
\pgfsetdash{}{0pt}%
\pgfsys@defobject{currentmarker}{\pgfqpoint{-0.016667in}{-0.016667in}}{\pgfqpoint{0.016667in}{0.016667in}}{%
\pgfpathmoveto{\pgfqpoint{0.000000in}{-0.016667in}}%
\pgfpathcurveto{\pgfqpoint{0.004420in}{-0.016667in}}{\pgfqpoint{0.008660in}{-0.014911in}}{\pgfqpoint{0.011785in}{-0.011785in}}%
\pgfpathcurveto{\pgfqpoint{0.014911in}{-0.008660in}}{\pgfqpoint{0.016667in}{-0.004420in}}{\pgfqpoint{0.016667in}{0.000000in}}%
\pgfpathcurveto{\pgfqpoint{0.016667in}{0.004420in}}{\pgfqpoint{0.014911in}{0.008660in}}{\pgfqpoint{0.011785in}{0.011785in}}%
\pgfpathcurveto{\pgfqpoint{0.008660in}{0.014911in}}{\pgfqpoint{0.004420in}{0.016667in}}{\pgfqpoint{0.000000in}{0.016667in}}%
\pgfpathcurveto{\pgfqpoint{-0.004420in}{0.016667in}}{\pgfqpoint{-0.008660in}{0.014911in}}{\pgfqpoint{-0.011785in}{0.011785in}}%
\pgfpathcurveto{\pgfqpoint{-0.014911in}{0.008660in}}{\pgfqpoint{-0.016667in}{0.004420in}}{\pgfqpoint{-0.016667in}{0.000000in}}%
\pgfpathcurveto{\pgfqpoint{-0.016667in}{-0.004420in}}{\pgfqpoint{-0.014911in}{-0.008660in}}{\pgfqpoint{-0.011785in}{-0.011785in}}%
\pgfpathcurveto{\pgfqpoint{-0.008660in}{-0.014911in}}{\pgfqpoint{-0.004420in}{-0.016667in}}{\pgfqpoint{0.000000in}{-0.016667in}}%
\pgfpathlineto{\pgfqpoint{0.000000in}{-0.016667in}}%
\pgfpathclose%
\pgfusepath{stroke,fill}%
}%
\begin{pgfscope}%
\pgfsys@transformshift{2.887116in}{4.956746in}%
\pgfsys@useobject{currentmarker}{}%
\end{pgfscope}%
\end{pgfscope}%
\begin{pgfscope}%
\pgfpathrectangle{\pgfqpoint{0.100000in}{2.493063in}}{\pgfqpoint{5.037500in}{3.427208in}}%
\pgfusepath{clip}%
\pgfsetrectcap%
\pgfsetroundjoin%
\pgfsetlinewidth{1.505625pt}%
\definecolor{currentstroke}{rgb}{0.678431,1.000000,0.184314}%
\pgfsetstrokecolor{currentstroke}%
\pgfsetstrokeopacity{0.500000}%
\pgfsetdash{}{0pt}%
\pgfpathmoveto{\pgfqpoint{2.946176in}{5.051799in}}%
\pgfusepath{stroke}%
\end{pgfscope}%
\begin{pgfscope}%
\pgfpathrectangle{\pgfqpoint{0.100000in}{2.493063in}}{\pgfqpoint{5.037500in}{3.427208in}}%
\pgfusepath{clip}%
\pgfsetbuttcap%
\pgfsetroundjoin%
\definecolor{currentfill}{rgb}{0.678431,1.000000,0.184314}%
\pgfsetfillcolor{currentfill}%
\pgfsetfillopacity{0.500000}%
\pgfsetlinewidth{0.250937pt}%
\definecolor{currentstroke}{rgb}{0.000000,0.000000,0.000000}%
\pgfsetstrokecolor{currentstroke}%
\pgfsetstrokeopacity{0.500000}%
\pgfsetdash{}{0pt}%
\pgfsys@defobject{currentmarker}{\pgfqpoint{-0.019444in}{-0.019444in}}{\pgfqpoint{0.019444in}{0.019444in}}{%
\pgfpathmoveto{\pgfqpoint{0.000000in}{-0.019444in}}%
\pgfpathcurveto{\pgfqpoint{0.005157in}{-0.019444in}}{\pgfqpoint{0.010103in}{-0.017396in}}{\pgfqpoint{0.013749in}{-0.013749in}}%
\pgfpathcurveto{\pgfqpoint{0.017396in}{-0.010103in}}{\pgfqpoint{0.019444in}{-0.005157in}}{\pgfqpoint{0.019444in}{0.000000in}}%
\pgfpathcurveto{\pgfqpoint{0.019444in}{0.005157in}}{\pgfqpoint{0.017396in}{0.010103in}}{\pgfqpoint{0.013749in}{0.013749in}}%
\pgfpathcurveto{\pgfqpoint{0.010103in}{0.017396in}}{\pgfqpoint{0.005157in}{0.019444in}}{\pgfqpoint{0.000000in}{0.019444in}}%
\pgfpathcurveto{\pgfqpoint{-0.005157in}{0.019444in}}{\pgfqpoint{-0.010103in}{0.017396in}}{\pgfqpoint{-0.013749in}{0.013749in}}%
\pgfpathcurveto{\pgfqpoint{-0.017396in}{0.010103in}}{\pgfqpoint{-0.019444in}{0.005157in}}{\pgfqpoint{-0.019444in}{0.000000in}}%
\pgfpathcurveto{\pgfqpoint{-0.019444in}{-0.005157in}}{\pgfqpoint{-0.017396in}{-0.010103in}}{\pgfqpoint{-0.013749in}{-0.013749in}}%
\pgfpathcurveto{\pgfqpoint{-0.010103in}{-0.017396in}}{\pgfqpoint{-0.005157in}{-0.019444in}}{\pgfqpoint{0.000000in}{-0.019444in}}%
\pgfpathlineto{\pgfqpoint{0.000000in}{-0.019444in}}%
\pgfpathclose%
\pgfusepath{stroke,fill}%
}%
\begin{pgfscope}%
\pgfsys@transformshift{2.946176in}{5.051799in}%
\pgfsys@useobject{currentmarker}{}%
\end{pgfscope}%
\end{pgfscope}%
\begin{pgfscope}%
\pgfpathrectangle{\pgfqpoint{0.100000in}{2.493063in}}{\pgfqpoint{5.037500in}{3.427208in}}%
\pgfusepath{clip}%
\pgfsetrectcap%
\pgfsetroundjoin%
\pgfsetlinewidth{1.505625pt}%
\definecolor{currentstroke}{rgb}{0.678431,1.000000,0.184314}%
\pgfsetstrokecolor{currentstroke}%
\pgfsetstrokeopacity{0.500000}%
\pgfsetdash{}{0pt}%
\pgfpathmoveto{\pgfqpoint{3.014958in}{4.943032in}}%
\pgfusepath{stroke}%
\end{pgfscope}%
\begin{pgfscope}%
\pgfpathrectangle{\pgfqpoint{0.100000in}{2.493063in}}{\pgfqpoint{5.037500in}{3.427208in}}%
\pgfusepath{clip}%
\pgfsetbuttcap%
\pgfsetroundjoin%
\definecolor{currentfill}{rgb}{0.678431,1.000000,0.184314}%
\pgfsetfillcolor{currentfill}%
\pgfsetfillopacity{0.500000}%
\pgfsetlinewidth{0.250937pt}%
\definecolor{currentstroke}{rgb}{0.000000,0.000000,0.000000}%
\pgfsetstrokecolor{currentstroke}%
\pgfsetstrokeopacity{0.500000}%
\pgfsetdash{}{0pt}%
\pgfsys@defobject{currentmarker}{\pgfqpoint{-0.025000in}{-0.025000in}}{\pgfqpoint{0.025000in}{0.025000in}}{%
\pgfpathmoveto{\pgfqpoint{0.000000in}{-0.025000in}}%
\pgfpathcurveto{\pgfqpoint{0.006630in}{-0.025000in}}{\pgfqpoint{0.012989in}{-0.022366in}}{\pgfqpoint{0.017678in}{-0.017678in}}%
\pgfpathcurveto{\pgfqpoint{0.022366in}{-0.012989in}}{\pgfqpoint{0.025000in}{-0.006630in}}{\pgfqpoint{0.025000in}{0.000000in}}%
\pgfpathcurveto{\pgfqpoint{0.025000in}{0.006630in}}{\pgfqpoint{0.022366in}{0.012989in}}{\pgfqpoint{0.017678in}{0.017678in}}%
\pgfpathcurveto{\pgfqpoint{0.012989in}{0.022366in}}{\pgfqpoint{0.006630in}{0.025000in}}{\pgfqpoint{0.000000in}{0.025000in}}%
\pgfpathcurveto{\pgfqpoint{-0.006630in}{0.025000in}}{\pgfqpoint{-0.012989in}{0.022366in}}{\pgfqpoint{-0.017678in}{0.017678in}}%
\pgfpathcurveto{\pgfqpoint{-0.022366in}{0.012989in}}{\pgfqpoint{-0.025000in}{0.006630in}}{\pgfqpoint{-0.025000in}{0.000000in}}%
\pgfpathcurveto{\pgfqpoint{-0.025000in}{-0.006630in}}{\pgfqpoint{-0.022366in}{-0.012989in}}{\pgfqpoint{-0.017678in}{-0.017678in}}%
\pgfpathcurveto{\pgfqpoint{-0.012989in}{-0.022366in}}{\pgfqpoint{-0.006630in}{-0.025000in}}{\pgfqpoint{0.000000in}{-0.025000in}}%
\pgfpathlineto{\pgfqpoint{0.000000in}{-0.025000in}}%
\pgfpathclose%
\pgfusepath{stroke,fill}%
}%
\begin{pgfscope}%
\pgfsys@transformshift{3.014958in}{4.943032in}%
\pgfsys@useobject{currentmarker}{}%
\end{pgfscope}%
\end{pgfscope}%
\begin{pgfscope}%
\pgfpathrectangle{\pgfqpoint{0.100000in}{2.493063in}}{\pgfqpoint{5.037500in}{3.427208in}}%
\pgfusepath{clip}%
\pgfsetrectcap%
\pgfsetroundjoin%
\pgfsetlinewidth{1.505625pt}%
\definecolor{currentstroke}{rgb}{0.678431,1.000000,0.184314}%
\pgfsetstrokecolor{currentstroke}%
\pgfsetstrokeopacity{0.500000}%
\pgfsetdash{}{0pt}%
\pgfpathmoveto{\pgfqpoint{2.871956in}{5.118443in}}%
\pgfusepath{stroke}%
\end{pgfscope}%
\begin{pgfscope}%
\pgfpathrectangle{\pgfqpoint{0.100000in}{2.493063in}}{\pgfqpoint{5.037500in}{3.427208in}}%
\pgfusepath{clip}%
\pgfsetbuttcap%
\pgfsetroundjoin%
\definecolor{currentfill}{rgb}{0.678431,1.000000,0.184314}%
\pgfsetfillcolor{currentfill}%
\pgfsetfillopacity{0.500000}%
\pgfsetlinewidth{0.250937pt}%
\definecolor{currentstroke}{rgb}{0.000000,0.000000,0.000000}%
\pgfsetstrokecolor{currentstroke}%
\pgfsetstrokeopacity{0.500000}%
\pgfsetdash{}{0pt}%
\pgfsys@defobject{currentmarker}{\pgfqpoint{-0.022222in}{-0.022222in}}{\pgfqpoint{0.022222in}{0.022222in}}{%
\pgfpathmoveto{\pgfqpoint{0.000000in}{-0.022222in}}%
\pgfpathcurveto{\pgfqpoint{0.005893in}{-0.022222in}}{\pgfqpoint{0.011546in}{-0.019881in}}{\pgfqpoint{0.015713in}{-0.015713in}}%
\pgfpathcurveto{\pgfqpoint{0.019881in}{-0.011546in}}{\pgfqpoint{0.022222in}{-0.005893in}}{\pgfqpoint{0.022222in}{0.000000in}}%
\pgfpathcurveto{\pgfqpoint{0.022222in}{0.005893in}}{\pgfqpoint{0.019881in}{0.011546in}}{\pgfqpoint{0.015713in}{0.015713in}}%
\pgfpathcurveto{\pgfqpoint{0.011546in}{0.019881in}}{\pgfqpoint{0.005893in}{0.022222in}}{\pgfqpoint{0.000000in}{0.022222in}}%
\pgfpathcurveto{\pgfqpoint{-0.005893in}{0.022222in}}{\pgfqpoint{-0.011546in}{0.019881in}}{\pgfqpoint{-0.015713in}{0.015713in}}%
\pgfpathcurveto{\pgfqpoint{-0.019881in}{0.011546in}}{\pgfqpoint{-0.022222in}{0.005893in}}{\pgfqpoint{-0.022222in}{0.000000in}}%
\pgfpathcurveto{\pgfqpoint{-0.022222in}{-0.005893in}}{\pgfqpoint{-0.019881in}{-0.011546in}}{\pgfqpoint{-0.015713in}{-0.015713in}}%
\pgfpathcurveto{\pgfqpoint{-0.011546in}{-0.019881in}}{\pgfqpoint{-0.005893in}{-0.022222in}}{\pgfqpoint{0.000000in}{-0.022222in}}%
\pgfpathlineto{\pgfqpoint{0.000000in}{-0.022222in}}%
\pgfpathclose%
\pgfusepath{stroke,fill}%
}%
\begin{pgfscope}%
\pgfsys@transformshift{2.871956in}{5.118443in}%
\pgfsys@useobject{currentmarker}{}%
\end{pgfscope}%
\end{pgfscope}%
\begin{pgfscope}%
\pgfpathrectangle{\pgfqpoint{0.100000in}{2.493063in}}{\pgfqpoint{5.037500in}{3.427208in}}%
\pgfusepath{clip}%
\pgfsetrectcap%
\pgfsetroundjoin%
\pgfsetlinewidth{1.505625pt}%
\definecolor{currentstroke}{rgb}{0.678431,1.000000,0.184314}%
\pgfsetstrokecolor{currentstroke}%
\pgfsetstrokeopacity{0.500000}%
\pgfsetdash{}{0pt}%
\pgfpathmoveto{\pgfqpoint{3.397181in}{3.382431in}}%
\pgfusepath{stroke}%
\end{pgfscope}%
\begin{pgfscope}%
\pgfpathrectangle{\pgfqpoint{0.100000in}{2.493063in}}{\pgfqpoint{5.037500in}{3.427208in}}%
\pgfusepath{clip}%
\pgfsetbuttcap%
\pgfsetroundjoin%
\definecolor{currentfill}{rgb}{0.678431,1.000000,0.184314}%
\pgfsetfillcolor{currentfill}%
\pgfsetfillopacity{0.500000}%
\pgfsetlinewidth{0.250937pt}%
\definecolor{currentstroke}{rgb}{0.000000,0.000000,0.000000}%
\pgfsetstrokecolor{currentstroke}%
\pgfsetstrokeopacity{0.500000}%
\pgfsetdash{}{0pt}%
\pgfsys@defobject{currentmarker}{\pgfqpoint{-0.075000in}{-0.075000in}}{\pgfqpoint{0.075000in}{0.075000in}}{%
\pgfpathmoveto{\pgfqpoint{0.000000in}{-0.075000in}}%
\pgfpathcurveto{\pgfqpoint{0.019890in}{-0.075000in}}{\pgfqpoint{0.038968in}{-0.067098in}}{\pgfqpoint{0.053033in}{-0.053033in}}%
\pgfpathcurveto{\pgfqpoint{0.067098in}{-0.038968in}}{\pgfqpoint{0.075000in}{-0.019890in}}{\pgfqpoint{0.075000in}{0.000000in}}%
\pgfpathcurveto{\pgfqpoint{0.075000in}{0.019890in}}{\pgfqpoint{0.067098in}{0.038968in}}{\pgfqpoint{0.053033in}{0.053033in}}%
\pgfpathcurveto{\pgfqpoint{0.038968in}{0.067098in}}{\pgfqpoint{0.019890in}{0.075000in}}{\pgfqpoint{0.000000in}{0.075000in}}%
\pgfpathcurveto{\pgfqpoint{-0.019890in}{0.075000in}}{\pgfqpoint{-0.038968in}{0.067098in}}{\pgfqpoint{-0.053033in}{0.053033in}}%
\pgfpathcurveto{\pgfqpoint{-0.067098in}{0.038968in}}{\pgfqpoint{-0.075000in}{0.019890in}}{\pgfqpoint{-0.075000in}{0.000000in}}%
\pgfpathcurveto{\pgfqpoint{-0.075000in}{-0.019890in}}{\pgfqpoint{-0.067098in}{-0.038968in}}{\pgfqpoint{-0.053033in}{-0.053033in}}%
\pgfpathcurveto{\pgfqpoint{-0.038968in}{-0.067098in}}{\pgfqpoint{-0.019890in}{-0.075000in}}{\pgfqpoint{0.000000in}{-0.075000in}}%
\pgfpathlineto{\pgfqpoint{0.000000in}{-0.075000in}}%
\pgfpathclose%
\pgfusepath{stroke,fill}%
}%
\begin{pgfscope}%
\pgfsys@transformshift{3.397181in}{3.382431in}%
\pgfsys@useobject{currentmarker}{}%
\end{pgfscope}%
\end{pgfscope}%
\begin{pgfscope}%
\pgfpathrectangle{\pgfqpoint{0.100000in}{2.493063in}}{\pgfqpoint{5.037500in}{3.427208in}}%
\pgfusepath{clip}%
\pgfsetrectcap%
\pgfsetroundjoin%
\pgfsetlinewidth{1.505625pt}%
\definecolor{currentstroke}{rgb}{0.678431,1.000000,0.184314}%
\pgfsetstrokecolor{currentstroke}%
\pgfsetstrokeopacity{0.500000}%
\pgfsetdash{}{0pt}%
\pgfpathmoveto{\pgfqpoint{3.370369in}{3.492562in}}%
\pgfusepath{stroke}%
\end{pgfscope}%
\begin{pgfscope}%
\pgfpathrectangle{\pgfqpoint{0.100000in}{2.493063in}}{\pgfqpoint{5.037500in}{3.427208in}}%
\pgfusepath{clip}%
\pgfsetbuttcap%
\pgfsetroundjoin%
\definecolor{currentfill}{rgb}{0.678431,1.000000,0.184314}%
\pgfsetfillcolor{currentfill}%
\pgfsetfillopacity{0.500000}%
\pgfsetlinewidth{0.250937pt}%
\definecolor{currentstroke}{rgb}{0.000000,0.000000,0.000000}%
\pgfsetstrokecolor{currentstroke}%
\pgfsetstrokeopacity{0.500000}%
\pgfsetdash{}{0pt}%
\pgfsys@defobject{currentmarker}{\pgfqpoint{-0.063889in}{-0.063889in}}{\pgfqpoint{0.063889in}{0.063889in}}{%
\pgfpathmoveto{\pgfqpoint{0.000000in}{-0.063889in}}%
\pgfpathcurveto{\pgfqpoint{0.016944in}{-0.063889in}}{\pgfqpoint{0.033195in}{-0.057157in}}{\pgfqpoint{0.045176in}{-0.045176in}}%
\pgfpathcurveto{\pgfqpoint{0.057157in}{-0.033195in}}{\pgfqpoint{0.063889in}{-0.016944in}}{\pgfqpoint{0.063889in}{0.000000in}}%
\pgfpathcurveto{\pgfqpoint{0.063889in}{0.016944in}}{\pgfqpoint{0.057157in}{0.033195in}}{\pgfqpoint{0.045176in}{0.045176in}}%
\pgfpathcurveto{\pgfqpoint{0.033195in}{0.057157in}}{\pgfqpoint{0.016944in}{0.063889in}}{\pgfqpoint{0.000000in}{0.063889in}}%
\pgfpathcurveto{\pgfqpoint{-0.016944in}{0.063889in}}{\pgfqpoint{-0.033195in}{0.057157in}}{\pgfqpoint{-0.045176in}{0.045176in}}%
\pgfpathcurveto{\pgfqpoint{-0.057157in}{0.033195in}}{\pgfqpoint{-0.063889in}{0.016944in}}{\pgfqpoint{-0.063889in}{0.000000in}}%
\pgfpathcurveto{\pgfqpoint{-0.063889in}{-0.016944in}}{\pgfqpoint{-0.057157in}{-0.033195in}}{\pgfqpoint{-0.045176in}{-0.045176in}}%
\pgfpathcurveto{\pgfqpoint{-0.033195in}{-0.057157in}}{\pgfqpoint{-0.016944in}{-0.063889in}}{\pgfqpoint{0.000000in}{-0.063889in}}%
\pgfpathlineto{\pgfqpoint{0.000000in}{-0.063889in}}%
\pgfpathclose%
\pgfusepath{stroke,fill}%
}%
\begin{pgfscope}%
\pgfsys@transformshift{3.370369in}{3.492562in}%
\pgfsys@useobject{currentmarker}{}%
\end{pgfscope}%
\end{pgfscope}%
\begin{pgfscope}%
\pgfpathrectangle{\pgfqpoint{0.100000in}{2.493063in}}{\pgfqpoint{5.037500in}{3.427208in}}%
\pgfusepath{clip}%
\pgfsetrectcap%
\pgfsetroundjoin%
\pgfsetlinewidth{1.505625pt}%
\definecolor{currentstroke}{rgb}{0.678431,1.000000,0.184314}%
\pgfsetstrokecolor{currentstroke}%
\pgfsetstrokeopacity{0.500000}%
\pgfsetdash{}{0pt}%
\pgfpathmoveto{\pgfqpoint{3.277318in}{3.613073in}}%
\pgfusepath{stroke}%
\end{pgfscope}%
\begin{pgfscope}%
\pgfpathrectangle{\pgfqpoint{0.100000in}{2.493063in}}{\pgfqpoint{5.037500in}{3.427208in}}%
\pgfusepath{clip}%
\pgfsetbuttcap%
\pgfsetroundjoin%
\definecolor{currentfill}{rgb}{0.678431,1.000000,0.184314}%
\pgfsetfillcolor{currentfill}%
\pgfsetfillopacity{0.500000}%
\pgfsetlinewidth{0.250937pt}%
\definecolor{currentstroke}{rgb}{0.000000,0.000000,0.000000}%
\pgfsetstrokecolor{currentstroke}%
\pgfsetstrokeopacity{0.500000}%
\pgfsetdash{}{0pt}%
\pgfsys@defobject{currentmarker}{\pgfqpoint{-0.066667in}{-0.066667in}}{\pgfqpoint{0.066667in}{0.066667in}}{%
\pgfpathmoveto{\pgfqpoint{0.000000in}{-0.066667in}}%
\pgfpathcurveto{\pgfqpoint{0.017680in}{-0.066667in}}{\pgfqpoint{0.034639in}{-0.059642in}}{\pgfqpoint{0.047140in}{-0.047140in}}%
\pgfpathcurveto{\pgfqpoint{0.059642in}{-0.034639in}}{\pgfqpoint{0.066667in}{-0.017680in}}{\pgfqpoint{0.066667in}{0.000000in}}%
\pgfpathcurveto{\pgfqpoint{0.066667in}{0.017680in}}{\pgfqpoint{0.059642in}{0.034639in}}{\pgfqpoint{0.047140in}{0.047140in}}%
\pgfpathcurveto{\pgfqpoint{0.034639in}{0.059642in}}{\pgfqpoint{0.017680in}{0.066667in}}{\pgfqpoint{0.000000in}{0.066667in}}%
\pgfpathcurveto{\pgfqpoint{-0.017680in}{0.066667in}}{\pgfqpoint{-0.034639in}{0.059642in}}{\pgfqpoint{-0.047140in}{0.047140in}}%
\pgfpathcurveto{\pgfqpoint{-0.059642in}{0.034639in}}{\pgfqpoint{-0.066667in}{0.017680in}}{\pgfqpoint{-0.066667in}{0.000000in}}%
\pgfpathcurveto{\pgfqpoint{-0.066667in}{-0.017680in}}{\pgfqpoint{-0.059642in}{-0.034639in}}{\pgfqpoint{-0.047140in}{-0.047140in}}%
\pgfpathcurveto{\pgfqpoint{-0.034639in}{-0.059642in}}{\pgfqpoint{-0.017680in}{-0.066667in}}{\pgfqpoint{0.000000in}{-0.066667in}}%
\pgfpathlineto{\pgfqpoint{0.000000in}{-0.066667in}}%
\pgfpathclose%
\pgfusepath{stroke,fill}%
}%
\begin{pgfscope}%
\pgfsys@transformshift{3.277318in}{3.613073in}%
\pgfsys@useobject{currentmarker}{}%
\end{pgfscope}%
\end{pgfscope}%
\begin{pgfscope}%
\pgfpathrectangle{\pgfqpoint{0.100000in}{2.493063in}}{\pgfqpoint{5.037500in}{3.427208in}}%
\pgfusepath{clip}%
\pgfsetrectcap%
\pgfsetroundjoin%
\pgfsetlinewidth{1.505625pt}%
\definecolor{currentstroke}{rgb}{0.678431,1.000000,0.184314}%
\pgfsetstrokecolor{currentstroke}%
\pgfsetstrokeopacity{0.500000}%
\pgfsetdash{}{0pt}%
\pgfpathmoveto{\pgfqpoint{3.306143in}{4.180990in}}%
\pgfusepath{stroke}%
\end{pgfscope}%
\begin{pgfscope}%
\pgfpathrectangle{\pgfqpoint{0.100000in}{2.493063in}}{\pgfqpoint{5.037500in}{3.427208in}}%
\pgfusepath{clip}%
\pgfsetbuttcap%
\pgfsetroundjoin%
\definecolor{currentfill}{rgb}{0.678431,1.000000,0.184314}%
\pgfsetfillcolor{currentfill}%
\pgfsetfillopacity{0.500000}%
\pgfsetlinewidth{0.250937pt}%
\definecolor{currentstroke}{rgb}{0.000000,0.000000,0.000000}%
\pgfsetstrokecolor{currentstroke}%
\pgfsetstrokeopacity{0.500000}%
\pgfsetdash{}{0pt}%
\pgfsys@defobject{currentmarker}{\pgfqpoint{-0.008333in}{-0.008333in}}{\pgfqpoint{0.008333in}{0.008333in}}{%
\pgfpathmoveto{\pgfqpoint{0.000000in}{-0.008333in}}%
\pgfpathcurveto{\pgfqpoint{0.002210in}{-0.008333in}}{\pgfqpoint{0.004330in}{-0.007455in}}{\pgfqpoint{0.005893in}{-0.005893in}}%
\pgfpathcurveto{\pgfqpoint{0.007455in}{-0.004330in}}{\pgfqpoint{0.008333in}{-0.002210in}}{\pgfqpoint{0.008333in}{0.000000in}}%
\pgfpathcurveto{\pgfqpoint{0.008333in}{0.002210in}}{\pgfqpoint{0.007455in}{0.004330in}}{\pgfqpoint{0.005893in}{0.005893in}}%
\pgfpathcurveto{\pgfqpoint{0.004330in}{0.007455in}}{\pgfqpoint{0.002210in}{0.008333in}}{\pgfqpoint{0.000000in}{0.008333in}}%
\pgfpathcurveto{\pgfqpoint{-0.002210in}{0.008333in}}{\pgfqpoint{-0.004330in}{0.007455in}}{\pgfqpoint{-0.005893in}{0.005893in}}%
\pgfpathcurveto{\pgfqpoint{-0.007455in}{0.004330in}}{\pgfqpoint{-0.008333in}{0.002210in}}{\pgfqpoint{-0.008333in}{0.000000in}}%
\pgfpathcurveto{\pgfqpoint{-0.008333in}{-0.002210in}}{\pgfqpoint{-0.007455in}{-0.004330in}}{\pgfqpoint{-0.005893in}{-0.005893in}}%
\pgfpathcurveto{\pgfqpoint{-0.004330in}{-0.007455in}}{\pgfqpoint{-0.002210in}{-0.008333in}}{\pgfqpoint{0.000000in}{-0.008333in}}%
\pgfpathlineto{\pgfqpoint{0.000000in}{-0.008333in}}%
\pgfpathclose%
\pgfusepath{stroke,fill}%
}%
\begin{pgfscope}%
\pgfsys@transformshift{3.306143in}{4.180990in}%
\pgfsys@useobject{currentmarker}{}%
\end{pgfscope}%
\end{pgfscope}%
\begin{pgfscope}%
\pgfpathrectangle{\pgfqpoint{0.100000in}{2.493063in}}{\pgfqpoint{5.037500in}{3.427208in}}%
\pgfusepath{clip}%
\pgfsetrectcap%
\pgfsetroundjoin%
\pgfsetlinewidth{1.505625pt}%
\definecolor{currentstroke}{rgb}{0.501961,0.501961,0.501961}%
\pgfsetstrokecolor{currentstroke}%
\pgfsetstrokeopacity{0.500000}%
\pgfsetdash{}{0pt}%
\pgfpathmoveto{\pgfqpoint{3.042850in}{4.358911in}}%
\pgfusepath{stroke}%
\end{pgfscope}%
\begin{pgfscope}%
\pgfpathrectangle{\pgfqpoint{0.100000in}{2.493063in}}{\pgfqpoint{5.037500in}{3.427208in}}%
\pgfusepath{clip}%
\pgfsetbuttcap%
\pgfsetroundjoin%
\definecolor{currentfill}{rgb}{0.501961,0.501961,0.501961}%
\pgfsetfillcolor{currentfill}%
\pgfsetfillopacity{0.500000}%
\pgfsetlinewidth{0.250937pt}%
\definecolor{currentstroke}{rgb}{0.000000,0.000000,0.000000}%
\pgfsetstrokecolor{currentstroke}%
\pgfsetstrokeopacity{0.500000}%
\pgfsetdash{}{0pt}%
\pgfsys@defobject{currentmarker}{\pgfqpoint{-0.013889in}{-0.013889in}}{\pgfqpoint{0.013889in}{0.013889in}}{%
\pgfpathmoveto{\pgfqpoint{0.000000in}{-0.013889in}}%
\pgfpathcurveto{\pgfqpoint{0.003683in}{-0.013889in}}{\pgfqpoint{0.007216in}{-0.012425in}}{\pgfqpoint{0.009821in}{-0.009821in}}%
\pgfpathcurveto{\pgfqpoint{0.012425in}{-0.007216in}}{\pgfqpoint{0.013889in}{-0.003683in}}{\pgfqpoint{0.013889in}{0.000000in}}%
\pgfpathcurveto{\pgfqpoint{0.013889in}{0.003683in}}{\pgfqpoint{0.012425in}{0.007216in}}{\pgfqpoint{0.009821in}{0.009821in}}%
\pgfpathcurveto{\pgfqpoint{0.007216in}{0.012425in}}{\pgfqpoint{0.003683in}{0.013889in}}{\pgfqpoint{0.000000in}{0.013889in}}%
\pgfpathcurveto{\pgfqpoint{-0.003683in}{0.013889in}}{\pgfqpoint{-0.007216in}{0.012425in}}{\pgfqpoint{-0.009821in}{0.009821in}}%
\pgfpathcurveto{\pgfqpoint{-0.012425in}{0.007216in}}{\pgfqpoint{-0.013889in}{0.003683in}}{\pgfqpoint{-0.013889in}{0.000000in}}%
\pgfpathcurveto{\pgfqpoint{-0.013889in}{-0.003683in}}{\pgfqpoint{-0.012425in}{-0.007216in}}{\pgfqpoint{-0.009821in}{-0.009821in}}%
\pgfpathcurveto{\pgfqpoint{-0.007216in}{-0.012425in}}{\pgfqpoint{-0.003683in}{-0.013889in}}{\pgfqpoint{0.000000in}{-0.013889in}}%
\pgfpathlineto{\pgfqpoint{0.000000in}{-0.013889in}}%
\pgfpathclose%
\pgfusepath{stroke,fill}%
}%
\begin{pgfscope}%
\pgfsys@transformshift{3.042850in}{4.358911in}%
\pgfsys@useobject{currentmarker}{}%
\end{pgfscope}%
\end{pgfscope}%
\begin{pgfscope}%
\pgfpathrectangle{\pgfqpoint{0.100000in}{2.493063in}}{\pgfqpoint{5.037500in}{3.427208in}}%
\pgfusepath{clip}%
\pgfsetrectcap%
\pgfsetroundjoin%
\pgfsetlinewidth{1.505625pt}%
\definecolor{currentstroke}{rgb}{0.501961,0.501961,0.501961}%
\pgfsetstrokecolor{currentstroke}%
\pgfsetstrokeopacity{0.500000}%
\pgfsetdash{}{0pt}%
\pgfpathmoveto{\pgfqpoint{3.058630in}{4.316242in}}%
\pgfusepath{stroke}%
\end{pgfscope}%
\begin{pgfscope}%
\pgfpathrectangle{\pgfqpoint{0.100000in}{2.493063in}}{\pgfqpoint{5.037500in}{3.427208in}}%
\pgfusepath{clip}%
\pgfsetbuttcap%
\pgfsetroundjoin%
\definecolor{currentfill}{rgb}{0.501961,0.501961,0.501961}%
\pgfsetfillcolor{currentfill}%
\pgfsetfillopacity{0.500000}%
\pgfsetlinewidth{0.250937pt}%
\definecolor{currentstroke}{rgb}{0.000000,0.000000,0.000000}%
\pgfsetstrokecolor{currentstroke}%
\pgfsetstrokeopacity{0.500000}%
\pgfsetdash{}{0pt}%
\pgfsys@defobject{currentmarker}{\pgfqpoint{-0.013889in}{-0.013889in}}{\pgfqpoint{0.013889in}{0.013889in}}{%
\pgfpathmoveto{\pgfqpoint{0.000000in}{-0.013889in}}%
\pgfpathcurveto{\pgfqpoint{0.003683in}{-0.013889in}}{\pgfqpoint{0.007216in}{-0.012425in}}{\pgfqpoint{0.009821in}{-0.009821in}}%
\pgfpathcurveto{\pgfqpoint{0.012425in}{-0.007216in}}{\pgfqpoint{0.013889in}{-0.003683in}}{\pgfqpoint{0.013889in}{0.000000in}}%
\pgfpathcurveto{\pgfqpoint{0.013889in}{0.003683in}}{\pgfqpoint{0.012425in}{0.007216in}}{\pgfqpoint{0.009821in}{0.009821in}}%
\pgfpathcurveto{\pgfqpoint{0.007216in}{0.012425in}}{\pgfqpoint{0.003683in}{0.013889in}}{\pgfqpoint{0.000000in}{0.013889in}}%
\pgfpathcurveto{\pgfqpoint{-0.003683in}{0.013889in}}{\pgfqpoint{-0.007216in}{0.012425in}}{\pgfqpoint{-0.009821in}{0.009821in}}%
\pgfpathcurveto{\pgfqpoint{-0.012425in}{0.007216in}}{\pgfqpoint{-0.013889in}{0.003683in}}{\pgfqpoint{-0.013889in}{0.000000in}}%
\pgfpathcurveto{\pgfqpoint{-0.013889in}{-0.003683in}}{\pgfqpoint{-0.012425in}{-0.007216in}}{\pgfqpoint{-0.009821in}{-0.009821in}}%
\pgfpathcurveto{\pgfqpoint{-0.007216in}{-0.012425in}}{\pgfqpoint{-0.003683in}{-0.013889in}}{\pgfqpoint{0.000000in}{-0.013889in}}%
\pgfpathlineto{\pgfqpoint{0.000000in}{-0.013889in}}%
\pgfpathclose%
\pgfusepath{stroke,fill}%
}%
\begin{pgfscope}%
\pgfsys@transformshift{3.058630in}{4.316242in}%
\pgfsys@useobject{currentmarker}{}%
\end{pgfscope}%
\end{pgfscope}%
\begin{pgfscope}%
\pgfpathrectangle{\pgfqpoint{0.100000in}{2.493063in}}{\pgfqpoint{5.037500in}{3.427208in}}%
\pgfusepath{clip}%
\pgfsetrectcap%
\pgfsetroundjoin%
\pgfsetlinewidth{1.505625pt}%
\definecolor{currentstroke}{rgb}{0.501961,0.501961,0.501961}%
\pgfsetstrokecolor{currentstroke}%
\pgfsetstrokeopacity{0.500000}%
\pgfsetdash{}{0pt}%
\pgfpathmoveto{\pgfqpoint{2.848864in}{4.140427in}}%
\pgfusepath{stroke}%
\end{pgfscope}%
\begin{pgfscope}%
\pgfpathrectangle{\pgfqpoint{0.100000in}{2.493063in}}{\pgfqpoint{5.037500in}{3.427208in}}%
\pgfusepath{clip}%
\pgfsetbuttcap%
\pgfsetroundjoin%
\definecolor{currentfill}{rgb}{0.501961,0.501961,0.501961}%
\pgfsetfillcolor{currentfill}%
\pgfsetfillopacity{0.500000}%
\pgfsetlinewidth{0.250937pt}%
\definecolor{currentstroke}{rgb}{0.000000,0.000000,0.000000}%
\pgfsetstrokecolor{currentstroke}%
\pgfsetstrokeopacity{0.500000}%
\pgfsetdash{}{0pt}%
\pgfsys@defobject{currentmarker}{\pgfqpoint{-0.013889in}{-0.013889in}}{\pgfqpoint{0.013889in}{0.013889in}}{%
\pgfpathmoveto{\pgfqpoint{0.000000in}{-0.013889in}}%
\pgfpathcurveto{\pgfqpoint{0.003683in}{-0.013889in}}{\pgfqpoint{0.007216in}{-0.012425in}}{\pgfqpoint{0.009821in}{-0.009821in}}%
\pgfpathcurveto{\pgfqpoint{0.012425in}{-0.007216in}}{\pgfqpoint{0.013889in}{-0.003683in}}{\pgfqpoint{0.013889in}{0.000000in}}%
\pgfpathcurveto{\pgfqpoint{0.013889in}{0.003683in}}{\pgfqpoint{0.012425in}{0.007216in}}{\pgfqpoint{0.009821in}{0.009821in}}%
\pgfpathcurveto{\pgfqpoint{0.007216in}{0.012425in}}{\pgfqpoint{0.003683in}{0.013889in}}{\pgfqpoint{0.000000in}{0.013889in}}%
\pgfpathcurveto{\pgfqpoint{-0.003683in}{0.013889in}}{\pgfqpoint{-0.007216in}{0.012425in}}{\pgfqpoint{-0.009821in}{0.009821in}}%
\pgfpathcurveto{\pgfqpoint{-0.012425in}{0.007216in}}{\pgfqpoint{-0.013889in}{0.003683in}}{\pgfqpoint{-0.013889in}{0.000000in}}%
\pgfpathcurveto{\pgfqpoint{-0.013889in}{-0.003683in}}{\pgfqpoint{-0.012425in}{-0.007216in}}{\pgfqpoint{-0.009821in}{-0.009821in}}%
\pgfpathcurveto{\pgfqpoint{-0.007216in}{-0.012425in}}{\pgfqpoint{-0.003683in}{-0.013889in}}{\pgfqpoint{0.000000in}{-0.013889in}}%
\pgfpathlineto{\pgfqpoint{0.000000in}{-0.013889in}}%
\pgfpathclose%
\pgfusepath{stroke,fill}%
}%
\begin{pgfscope}%
\pgfsys@transformshift{2.848864in}{4.140427in}%
\pgfsys@useobject{currentmarker}{}%
\end{pgfscope}%
\end{pgfscope}%
\begin{pgfscope}%
\pgfpathrectangle{\pgfqpoint{0.100000in}{2.493063in}}{\pgfqpoint{5.037500in}{3.427208in}}%
\pgfusepath{clip}%
\pgfsetrectcap%
\pgfsetroundjoin%
\pgfsetlinewidth{1.505625pt}%
\definecolor{currentstroke}{rgb}{0.678431,1.000000,0.184314}%
\pgfsetstrokecolor{currentstroke}%
\pgfsetstrokeopacity{0.500000}%
\pgfsetdash{}{0pt}%
\pgfpathmoveto{\pgfqpoint{2.843184in}{4.370769in}}%
\pgfusepath{stroke}%
\end{pgfscope}%
\begin{pgfscope}%
\pgfpathrectangle{\pgfqpoint{0.100000in}{2.493063in}}{\pgfqpoint{5.037500in}{3.427208in}}%
\pgfusepath{clip}%
\pgfsetbuttcap%
\pgfsetroundjoin%
\definecolor{currentfill}{rgb}{0.678431,1.000000,0.184314}%
\pgfsetfillcolor{currentfill}%
\pgfsetfillopacity{0.500000}%
\pgfsetlinewidth{0.250937pt}%
\definecolor{currentstroke}{rgb}{0.000000,0.000000,0.000000}%
\pgfsetstrokecolor{currentstroke}%
\pgfsetstrokeopacity{0.500000}%
\pgfsetdash{}{0pt}%
\pgfsys@defobject{currentmarker}{\pgfqpoint{-0.008333in}{-0.008333in}}{\pgfqpoint{0.008333in}{0.008333in}}{%
\pgfpathmoveto{\pgfqpoint{0.000000in}{-0.008333in}}%
\pgfpathcurveto{\pgfqpoint{0.002210in}{-0.008333in}}{\pgfqpoint{0.004330in}{-0.007455in}}{\pgfqpoint{0.005893in}{-0.005893in}}%
\pgfpathcurveto{\pgfqpoint{0.007455in}{-0.004330in}}{\pgfqpoint{0.008333in}{-0.002210in}}{\pgfqpoint{0.008333in}{0.000000in}}%
\pgfpathcurveto{\pgfqpoint{0.008333in}{0.002210in}}{\pgfqpoint{0.007455in}{0.004330in}}{\pgfqpoint{0.005893in}{0.005893in}}%
\pgfpathcurveto{\pgfqpoint{0.004330in}{0.007455in}}{\pgfqpoint{0.002210in}{0.008333in}}{\pgfqpoint{0.000000in}{0.008333in}}%
\pgfpathcurveto{\pgfqpoint{-0.002210in}{0.008333in}}{\pgfqpoint{-0.004330in}{0.007455in}}{\pgfqpoint{-0.005893in}{0.005893in}}%
\pgfpathcurveto{\pgfqpoint{-0.007455in}{0.004330in}}{\pgfqpoint{-0.008333in}{0.002210in}}{\pgfqpoint{-0.008333in}{0.000000in}}%
\pgfpathcurveto{\pgfqpoint{-0.008333in}{-0.002210in}}{\pgfqpoint{-0.007455in}{-0.004330in}}{\pgfqpoint{-0.005893in}{-0.005893in}}%
\pgfpathcurveto{\pgfqpoint{-0.004330in}{-0.007455in}}{\pgfqpoint{-0.002210in}{-0.008333in}}{\pgfqpoint{0.000000in}{-0.008333in}}%
\pgfpathlineto{\pgfqpoint{0.000000in}{-0.008333in}}%
\pgfpathclose%
\pgfusepath{stroke,fill}%
}%
\begin{pgfscope}%
\pgfsys@transformshift{2.843184in}{4.370769in}%
\pgfsys@useobject{currentmarker}{}%
\end{pgfscope}%
\end{pgfscope}%
\begin{pgfscope}%
\pgfpathrectangle{\pgfqpoint{0.100000in}{2.493063in}}{\pgfqpoint{5.037500in}{3.427208in}}%
\pgfusepath{clip}%
\pgfsetrectcap%
\pgfsetroundjoin%
\pgfsetlinewidth{1.505625pt}%
\definecolor{currentstroke}{rgb}{0.501961,0.501961,0.501961}%
\pgfsetstrokecolor{currentstroke}%
\pgfsetstrokeopacity{0.500000}%
\pgfsetdash{}{0pt}%
\pgfpathmoveto{\pgfqpoint{2.817703in}{4.449465in}}%
\pgfusepath{stroke}%
\end{pgfscope}%
\begin{pgfscope}%
\pgfpathrectangle{\pgfqpoint{0.100000in}{2.493063in}}{\pgfqpoint{5.037500in}{3.427208in}}%
\pgfusepath{clip}%
\pgfsetbuttcap%
\pgfsetroundjoin%
\definecolor{currentfill}{rgb}{0.501961,0.501961,0.501961}%
\pgfsetfillcolor{currentfill}%
\pgfsetfillopacity{0.500000}%
\pgfsetlinewidth{0.250937pt}%
\definecolor{currentstroke}{rgb}{0.000000,0.000000,0.000000}%
\pgfsetstrokecolor{currentstroke}%
\pgfsetstrokeopacity{0.500000}%
\pgfsetdash{}{0pt}%
\pgfsys@defobject{currentmarker}{\pgfqpoint{-0.013889in}{-0.013889in}}{\pgfqpoint{0.013889in}{0.013889in}}{%
\pgfpathmoveto{\pgfqpoint{0.000000in}{-0.013889in}}%
\pgfpathcurveto{\pgfqpoint{0.003683in}{-0.013889in}}{\pgfqpoint{0.007216in}{-0.012425in}}{\pgfqpoint{0.009821in}{-0.009821in}}%
\pgfpathcurveto{\pgfqpoint{0.012425in}{-0.007216in}}{\pgfqpoint{0.013889in}{-0.003683in}}{\pgfqpoint{0.013889in}{0.000000in}}%
\pgfpathcurveto{\pgfqpoint{0.013889in}{0.003683in}}{\pgfqpoint{0.012425in}{0.007216in}}{\pgfqpoint{0.009821in}{0.009821in}}%
\pgfpathcurveto{\pgfqpoint{0.007216in}{0.012425in}}{\pgfqpoint{0.003683in}{0.013889in}}{\pgfqpoint{0.000000in}{0.013889in}}%
\pgfpathcurveto{\pgfqpoint{-0.003683in}{0.013889in}}{\pgfqpoint{-0.007216in}{0.012425in}}{\pgfqpoint{-0.009821in}{0.009821in}}%
\pgfpathcurveto{\pgfqpoint{-0.012425in}{0.007216in}}{\pgfqpoint{-0.013889in}{0.003683in}}{\pgfqpoint{-0.013889in}{0.000000in}}%
\pgfpathcurveto{\pgfqpoint{-0.013889in}{-0.003683in}}{\pgfqpoint{-0.012425in}{-0.007216in}}{\pgfqpoint{-0.009821in}{-0.009821in}}%
\pgfpathcurveto{\pgfqpoint{-0.007216in}{-0.012425in}}{\pgfqpoint{-0.003683in}{-0.013889in}}{\pgfqpoint{0.000000in}{-0.013889in}}%
\pgfpathlineto{\pgfqpoint{0.000000in}{-0.013889in}}%
\pgfpathclose%
\pgfusepath{stroke,fill}%
}%
\begin{pgfscope}%
\pgfsys@transformshift{2.817703in}{4.449465in}%
\pgfsys@useobject{currentmarker}{}%
\end{pgfscope}%
\end{pgfscope}%
\begin{pgfscope}%
\pgfpathrectangle{\pgfqpoint{0.100000in}{2.493063in}}{\pgfqpoint{5.037500in}{3.427208in}}%
\pgfusepath{clip}%
\pgfsetrectcap%
\pgfsetroundjoin%
\pgfsetlinewidth{1.505625pt}%
\definecolor{currentstroke}{rgb}{0.000000,0.000000,1.000000}%
\pgfsetstrokecolor{currentstroke}%
\pgfsetstrokeopacity{0.500000}%
\pgfsetdash{}{0pt}%
\pgfpathmoveto{\pgfqpoint{3.235918in}{4.329441in}}%
\pgfusepath{stroke}%
\end{pgfscope}%
\begin{pgfscope}%
\pgfpathrectangle{\pgfqpoint{0.100000in}{2.493063in}}{\pgfqpoint{5.037500in}{3.427208in}}%
\pgfusepath{clip}%
\pgfsetbuttcap%
\pgfsetroundjoin%
\definecolor{currentfill}{rgb}{0.000000,0.000000,1.000000}%
\pgfsetfillcolor{currentfill}%
\pgfsetfillopacity{0.500000}%
\pgfsetlinewidth{0.250937pt}%
\definecolor{currentstroke}{rgb}{0.000000,0.000000,0.000000}%
\pgfsetstrokecolor{currentstroke}%
\pgfsetstrokeopacity{0.500000}%
\pgfsetdash{}{0pt}%
\pgfsys@defobject{currentmarker}{\pgfqpoint{-0.005556in}{-0.005556in}}{\pgfqpoint{0.005556in}{0.005556in}}{%
\pgfpathmoveto{\pgfqpoint{0.000000in}{-0.005556in}}%
\pgfpathcurveto{\pgfqpoint{0.001473in}{-0.005556in}}{\pgfqpoint{0.002887in}{-0.004970in}}{\pgfqpoint{0.003928in}{-0.003928in}}%
\pgfpathcurveto{\pgfqpoint{0.004970in}{-0.002887in}}{\pgfqpoint{0.005556in}{-0.001473in}}{\pgfqpoint{0.005556in}{0.000000in}}%
\pgfpathcurveto{\pgfqpoint{0.005556in}{0.001473in}}{\pgfqpoint{0.004970in}{0.002887in}}{\pgfqpoint{0.003928in}{0.003928in}}%
\pgfpathcurveto{\pgfqpoint{0.002887in}{0.004970in}}{\pgfqpoint{0.001473in}{0.005556in}}{\pgfqpoint{0.000000in}{0.005556in}}%
\pgfpathcurveto{\pgfqpoint{-0.001473in}{0.005556in}}{\pgfqpoint{-0.002887in}{0.004970in}}{\pgfqpoint{-0.003928in}{0.003928in}}%
\pgfpathcurveto{\pgfqpoint{-0.004970in}{0.002887in}}{\pgfqpoint{-0.005556in}{0.001473in}}{\pgfqpoint{-0.005556in}{0.000000in}}%
\pgfpathcurveto{\pgfqpoint{-0.005556in}{-0.001473in}}{\pgfqpoint{-0.004970in}{-0.002887in}}{\pgfqpoint{-0.003928in}{-0.003928in}}%
\pgfpathcurveto{\pgfqpoint{-0.002887in}{-0.004970in}}{\pgfqpoint{-0.001473in}{-0.005556in}}{\pgfqpoint{0.000000in}{-0.005556in}}%
\pgfpathlineto{\pgfqpoint{0.000000in}{-0.005556in}}%
\pgfpathclose%
\pgfusepath{stroke,fill}%
}%
\begin{pgfscope}%
\pgfsys@transformshift{3.235918in}{4.329441in}%
\pgfsys@useobject{currentmarker}{}%
\end{pgfscope}%
\end{pgfscope}%
\begin{pgfscope}%
\pgfpathrectangle{\pgfqpoint{0.100000in}{2.493063in}}{\pgfqpoint{5.037500in}{3.427208in}}%
\pgfusepath{clip}%
\pgfsetrectcap%
\pgfsetroundjoin%
\pgfsetlinewidth{1.505625pt}%
\definecolor{currentstroke}{rgb}{0.678431,1.000000,0.184314}%
\pgfsetstrokecolor{currentstroke}%
\pgfsetstrokeopacity{0.500000}%
\pgfsetdash{}{0pt}%
\pgfpathmoveto{\pgfqpoint{2.960251in}{4.156874in}}%
\pgfusepath{stroke}%
\end{pgfscope}%
\begin{pgfscope}%
\pgfpathrectangle{\pgfqpoint{0.100000in}{2.493063in}}{\pgfqpoint{5.037500in}{3.427208in}}%
\pgfusepath{clip}%
\pgfsetbuttcap%
\pgfsetroundjoin%
\definecolor{currentfill}{rgb}{0.678431,1.000000,0.184314}%
\pgfsetfillcolor{currentfill}%
\pgfsetfillopacity{0.500000}%
\pgfsetlinewidth{0.250937pt}%
\definecolor{currentstroke}{rgb}{0.000000,0.000000,0.000000}%
\pgfsetstrokecolor{currentstroke}%
\pgfsetstrokeopacity{0.500000}%
\pgfsetdash{}{0pt}%
\pgfsys@defobject{currentmarker}{\pgfqpoint{-0.005556in}{-0.005556in}}{\pgfqpoint{0.005556in}{0.005556in}}{%
\pgfpathmoveto{\pgfqpoint{0.000000in}{-0.005556in}}%
\pgfpathcurveto{\pgfqpoint{0.001473in}{-0.005556in}}{\pgfqpoint{0.002887in}{-0.004970in}}{\pgfqpoint{0.003928in}{-0.003928in}}%
\pgfpathcurveto{\pgfqpoint{0.004970in}{-0.002887in}}{\pgfqpoint{0.005556in}{-0.001473in}}{\pgfqpoint{0.005556in}{0.000000in}}%
\pgfpathcurveto{\pgfqpoint{0.005556in}{0.001473in}}{\pgfqpoint{0.004970in}{0.002887in}}{\pgfqpoint{0.003928in}{0.003928in}}%
\pgfpathcurveto{\pgfqpoint{0.002887in}{0.004970in}}{\pgfqpoint{0.001473in}{0.005556in}}{\pgfqpoint{0.000000in}{0.005556in}}%
\pgfpathcurveto{\pgfqpoint{-0.001473in}{0.005556in}}{\pgfqpoint{-0.002887in}{0.004970in}}{\pgfqpoint{-0.003928in}{0.003928in}}%
\pgfpathcurveto{\pgfqpoint{-0.004970in}{0.002887in}}{\pgfqpoint{-0.005556in}{0.001473in}}{\pgfqpoint{-0.005556in}{0.000000in}}%
\pgfpathcurveto{\pgfqpoint{-0.005556in}{-0.001473in}}{\pgfqpoint{-0.004970in}{-0.002887in}}{\pgfqpoint{-0.003928in}{-0.003928in}}%
\pgfpathcurveto{\pgfqpoint{-0.002887in}{-0.004970in}}{\pgfqpoint{-0.001473in}{-0.005556in}}{\pgfqpoint{0.000000in}{-0.005556in}}%
\pgfpathlineto{\pgfqpoint{0.000000in}{-0.005556in}}%
\pgfpathclose%
\pgfusepath{stroke,fill}%
}%
\begin{pgfscope}%
\pgfsys@transformshift{2.960251in}{4.156874in}%
\pgfsys@useobject{currentmarker}{}%
\end{pgfscope}%
\end{pgfscope}%
\begin{pgfscope}%
\pgfpathrectangle{\pgfqpoint{0.100000in}{2.493063in}}{\pgfqpoint{5.037500in}{3.427208in}}%
\pgfusepath{clip}%
\pgfsetrectcap%
\pgfsetroundjoin%
\pgfsetlinewidth{1.505625pt}%
\definecolor{currentstroke}{rgb}{0.678431,1.000000,0.184314}%
\pgfsetstrokecolor{currentstroke}%
\pgfsetstrokeopacity{0.500000}%
\pgfsetdash{}{0pt}%
\pgfpathmoveto{\pgfqpoint{1.716761in}{5.225212in}}%
\pgfusepath{stroke}%
\end{pgfscope}%
\begin{pgfscope}%
\pgfpathrectangle{\pgfqpoint{0.100000in}{2.493063in}}{\pgfqpoint{5.037500in}{3.427208in}}%
\pgfusepath{clip}%
\pgfsetbuttcap%
\pgfsetroundjoin%
\definecolor{currentfill}{rgb}{0.678431,1.000000,0.184314}%
\pgfsetfillcolor{currentfill}%
\pgfsetfillopacity{0.500000}%
\pgfsetlinewidth{0.250937pt}%
\definecolor{currentstroke}{rgb}{0.000000,0.000000,0.000000}%
\pgfsetstrokecolor{currentstroke}%
\pgfsetstrokeopacity{0.500000}%
\pgfsetdash{}{0pt}%
\pgfsys@defobject{currentmarker}{\pgfqpoint{-0.008333in}{-0.008333in}}{\pgfqpoint{0.008333in}{0.008333in}}{%
\pgfpathmoveto{\pgfqpoint{0.000000in}{-0.008333in}}%
\pgfpathcurveto{\pgfqpoint{0.002210in}{-0.008333in}}{\pgfqpoint{0.004330in}{-0.007455in}}{\pgfqpoint{0.005893in}{-0.005893in}}%
\pgfpathcurveto{\pgfqpoint{0.007455in}{-0.004330in}}{\pgfqpoint{0.008333in}{-0.002210in}}{\pgfqpoint{0.008333in}{0.000000in}}%
\pgfpathcurveto{\pgfqpoint{0.008333in}{0.002210in}}{\pgfqpoint{0.007455in}{0.004330in}}{\pgfqpoint{0.005893in}{0.005893in}}%
\pgfpathcurveto{\pgfqpoint{0.004330in}{0.007455in}}{\pgfqpoint{0.002210in}{0.008333in}}{\pgfqpoint{0.000000in}{0.008333in}}%
\pgfpathcurveto{\pgfqpoint{-0.002210in}{0.008333in}}{\pgfqpoint{-0.004330in}{0.007455in}}{\pgfqpoint{-0.005893in}{0.005893in}}%
\pgfpathcurveto{\pgfqpoint{-0.007455in}{0.004330in}}{\pgfqpoint{-0.008333in}{0.002210in}}{\pgfqpoint{-0.008333in}{0.000000in}}%
\pgfpathcurveto{\pgfqpoint{-0.008333in}{-0.002210in}}{\pgfqpoint{-0.007455in}{-0.004330in}}{\pgfqpoint{-0.005893in}{-0.005893in}}%
\pgfpathcurveto{\pgfqpoint{-0.004330in}{-0.007455in}}{\pgfqpoint{-0.002210in}{-0.008333in}}{\pgfqpoint{0.000000in}{-0.008333in}}%
\pgfpathlineto{\pgfqpoint{0.000000in}{-0.008333in}}%
\pgfpathclose%
\pgfusepath{stroke,fill}%
}%
\begin{pgfscope}%
\pgfsys@transformshift{1.716761in}{5.225212in}%
\pgfsys@useobject{currentmarker}{}%
\end{pgfscope}%
\end{pgfscope}%
\begin{pgfscope}%
\pgfpathrectangle{\pgfqpoint{0.100000in}{2.493063in}}{\pgfqpoint{5.037500in}{3.427208in}}%
\pgfusepath{clip}%
\pgfsetrectcap%
\pgfsetroundjoin%
\pgfsetlinewidth{1.505625pt}%
\definecolor{currentstroke}{rgb}{0.501961,0.501961,0.501961}%
\pgfsetstrokecolor{currentstroke}%
\pgfsetstrokeopacity{0.500000}%
\pgfsetdash{}{0pt}%
\pgfpathmoveto{\pgfqpoint{1.529298in}{5.458464in}}%
\pgfusepath{stroke}%
\end{pgfscope}%
\begin{pgfscope}%
\pgfpathrectangle{\pgfqpoint{0.100000in}{2.493063in}}{\pgfqpoint{5.037500in}{3.427208in}}%
\pgfusepath{clip}%
\pgfsetbuttcap%
\pgfsetroundjoin%
\definecolor{currentfill}{rgb}{0.501961,0.501961,0.501961}%
\pgfsetfillcolor{currentfill}%
\pgfsetfillopacity{0.500000}%
\pgfsetlinewidth{0.250937pt}%
\definecolor{currentstroke}{rgb}{0.000000,0.000000,0.000000}%
\pgfsetstrokecolor{currentstroke}%
\pgfsetstrokeopacity{0.500000}%
\pgfsetdash{}{0pt}%
\pgfsys@defobject{currentmarker}{\pgfqpoint{-0.013889in}{-0.013889in}}{\pgfqpoint{0.013889in}{0.013889in}}{%
\pgfpathmoveto{\pgfqpoint{0.000000in}{-0.013889in}}%
\pgfpathcurveto{\pgfqpoint{0.003683in}{-0.013889in}}{\pgfqpoint{0.007216in}{-0.012425in}}{\pgfqpoint{0.009821in}{-0.009821in}}%
\pgfpathcurveto{\pgfqpoint{0.012425in}{-0.007216in}}{\pgfqpoint{0.013889in}{-0.003683in}}{\pgfqpoint{0.013889in}{0.000000in}}%
\pgfpathcurveto{\pgfqpoint{0.013889in}{0.003683in}}{\pgfqpoint{0.012425in}{0.007216in}}{\pgfqpoint{0.009821in}{0.009821in}}%
\pgfpathcurveto{\pgfqpoint{0.007216in}{0.012425in}}{\pgfqpoint{0.003683in}{0.013889in}}{\pgfqpoint{0.000000in}{0.013889in}}%
\pgfpathcurveto{\pgfqpoint{-0.003683in}{0.013889in}}{\pgfqpoint{-0.007216in}{0.012425in}}{\pgfqpoint{-0.009821in}{0.009821in}}%
\pgfpathcurveto{\pgfqpoint{-0.012425in}{0.007216in}}{\pgfqpoint{-0.013889in}{0.003683in}}{\pgfqpoint{-0.013889in}{0.000000in}}%
\pgfpathcurveto{\pgfqpoint{-0.013889in}{-0.003683in}}{\pgfqpoint{-0.012425in}{-0.007216in}}{\pgfqpoint{-0.009821in}{-0.009821in}}%
\pgfpathcurveto{\pgfqpoint{-0.007216in}{-0.012425in}}{\pgfqpoint{-0.003683in}{-0.013889in}}{\pgfqpoint{0.000000in}{-0.013889in}}%
\pgfpathlineto{\pgfqpoint{0.000000in}{-0.013889in}}%
\pgfpathclose%
\pgfusepath{stroke,fill}%
}%
\begin{pgfscope}%
\pgfsys@transformshift{1.529298in}{5.458464in}%
\pgfsys@useobject{currentmarker}{}%
\end{pgfscope}%
\end{pgfscope}%
\begin{pgfscope}%
\pgfpathrectangle{\pgfqpoint{0.100000in}{2.493063in}}{\pgfqpoint{5.037500in}{3.427208in}}%
\pgfusepath{clip}%
\pgfsetrectcap%
\pgfsetroundjoin%
\pgfsetlinewidth{1.505625pt}%
\definecolor{currentstroke}{rgb}{0.501961,0.501961,0.501961}%
\pgfsetstrokecolor{currentstroke}%
\pgfsetstrokeopacity{0.500000}%
\pgfsetdash{}{0pt}%
\pgfpathmoveto{\pgfqpoint{1.305298in}{5.427237in}}%
\pgfusepath{stroke}%
\end{pgfscope}%
\begin{pgfscope}%
\pgfpathrectangle{\pgfqpoint{0.100000in}{2.493063in}}{\pgfqpoint{5.037500in}{3.427208in}}%
\pgfusepath{clip}%
\pgfsetbuttcap%
\pgfsetroundjoin%
\definecolor{currentfill}{rgb}{0.501961,0.501961,0.501961}%
\pgfsetfillcolor{currentfill}%
\pgfsetfillopacity{0.500000}%
\pgfsetlinewidth{0.250937pt}%
\definecolor{currentstroke}{rgb}{0.000000,0.000000,0.000000}%
\pgfsetstrokecolor{currentstroke}%
\pgfsetstrokeopacity{0.500000}%
\pgfsetdash{}{0pt}%
\pgfsys@defobject{currentmarker}{\pgfqpoint{-0.013889in}{-0.013889in}}{\pgfqpoint{0.013889in}{0.013889in}}{%
\pgfpathmoveto{\pgfqpoint{0.000000in}{-0.013889in}}%
\pgfpathcurveto{\pgfqpoint{0.003683in}{-0.013889in}}{\pgfqpoint{0.007216in}{-0.012425in}}{\pgfqpoint{0.009821in}{-0.009821in}}%
\pgfpathcurveto{\pgfqpoint{0.012425in}{-0.007216in}}{\pgfqpoint{0.013889in}{-0.003683in}}{\pgfqpoint{0.013889in}{0.000000in}}%
\pgfpathcurveto{\pgfqpoint{0.013889in}{0.003683in}}{\pgfqpoint{0.012425in}{0.007216in}}{\pgfqpoint{0.009821in}{0.009821in}}%
\pgfpathcurveto{\pgfqpoint{0.007216in}{0.012425in}}{\pgfqpoint{0.003683in}{0.013889in}}{\pgfqpoint{0.000000in}{0.013889in}}%
\pgfpathcurveto{\pgfqpoint{-0.003683in}{0.013889in}}{\pgfqpoint{-0.007216in}{0.012425in}}{\pgfqpoint{-0.009821in}{0.009821in}}%
\pgfpathcurveto{\pgfqpoint{-0.012425in}{0.007216in}}{\pgfqpoint{-0.013889in}{0.003683in}}{\pgfqpoint{-0.013889in}{0.000000in}}%
\pgfpathcurveto{\pgfqpoint{-0.013889in}{-0.003683in}}{\pgfqpoint{-0.012425in}{-0.007216in}}{\pgfqpoint{-0.009821in}{-0.009821in}}%
\pgfpathcurveto{\pgfqpoint{-0.007216in}{-0.012425in}}{\pgfqpoint{-0.003683in}{-0.013889in}}{\pgfqpoint{0.000000in}{-0.013889in}}%
\pgfpathlineto{\pgfqpoint{0.000000in}{-0.013889in}}%
\pgfpathclose%
\pgfusepath{stroke,fill}%
}%
\begin{pgfscope}%
\pgfsys@transformshift{1.305298in}{5.427237in}%
\pgfsys@useobject{currentmarker}{}%
\end{pgfscope}%
\end{pgfscope}%
\begin{pgfscope}%
\pgfpathrectangle{\pgfqpoint{0.100000in}{2.493063in}}{\pgfqpoint{5.037500in}{3.427208in}}%
\pgfusepath{clip}%
\pgfsetrectcap%
\pgfsetroundjoin%
\pgfsetlinewidth{1.505625pt}%
\definecolor{currentstroke}{rgb}{0.678431,1.000000,0.184314}%
\pgfsetstrokecolor{currentstroke}%
\pgfsetstrokeopacity{0.500000}%
\pgfsetdash{}{0pt}%
\pgfpathmoveto{\pgfqpoint{2.513569in}{4.587898in}}%
\pgfusepath{stroke}%
\end{pgfscope}%
\begin{pgfscope}%
\pgfpathrectangle{\pgfqpoint{0.100000in}{2.493063in}}{\pgfqpoint{5.037500in}{3.427208in}}%
\pgfusepath{clip}%
\pgfsetbuttcap%
\pgfsetroundjoin%
\definecolor{currentfill}{rgb}{0.678431,1.000000,0.184314}%
\pgfsetfillcolor{currentfill}%
\pgfsetfillopacity{0.500000}%
\pgfsetlinewidth{0.250937pt}%
\definecolor{currentstroke}{rgb}{0.000000,0.000000,0.000000}%
\pgfsetstrokecolor{currentstroke}%
\pgfsetstrokeopacity{0.500000}%
\pgfsetdash{}{0pt}%
\pgfsys@defobject{currentmarker}{\pgfqpoint{-0.027778in}{-0.027778in}}{\pgfqpoint{0.027778in}{0.027778in}}{%
\pgfpathmoveto{\pgfqpoint{0.000000in}{-0.027778in}}%
\pgfpathcurveto{\pgfqpoint{0.007367in}{-0.027778in}}{\pgfqpoint{0.014433in}{-0.024851in}}{\pgfqpoint{0.019642in}{-0.019642in}}%
\pgfpathcurveto{\pgfqpoint{0.024851in}{-0.014433in}}{\pgfqpoint{0.027778in}{-0.007367in}}{\pgfqpoint{0.027778in}{0.000000in}}%
\pgfpathcurveto{\pgfqpoint{0.027778in}{0.007367in}}{\pgfqpoint{0.024851in}{0.014433in}}{\pgfqpoint{0.019642in}{0.019642in}}%
\pgfpathcurveto{\pgfqpoint{0.014433in}{0.024851in}}{\pgfqpoint{0.007367in}{0.027778in}}{\pgfqpoint{0.000000in}{0.027778in}}%
\pgfpathcurveto{\pgfqpoint{-0.007367in}{0.027778in}}{\pgfqpoint{-0.014433in}{0.024851in}}{\pgfqpoint{-0.019642in}{0.019642in}}%
\pgfpathcurveto{\pgfqpoint{-0.024851in}{0.014433in}}{\pgfqpoint{-0.027778in}{0.007367in}}{\pgfqpoint{-0.027778in}{0.000000in}}%
\pgfpathcurveto{\pgfqpoint{-0.027778in}{-0.007367in}}{\pgfqpoint{-0.024851in}{-0.014433in}}{\pgfqpoint{-0.019642in}{-0.019642in}}%
\pgfpathcurveto{\pgfqpoint{-0.014433in}{-0.024851in}}{\pgfqpoint{-0.007367in}{-0.027778in}}{\pgfqpoint{0.000000in}{-0.027778in}}%
\pgfpathlineto{\pgfqpoint{0.000000in}{-0.027778in}}%
\pgfpathclose%
\pgfusepath{stroke,fill}%
}%
\begin{pgfscope}%
\pgfsys@transformshift{2.513569in}{4.587898in}%
\pgfsys@useobject{currentmarker}{}%
\end{pgfscope}%
\end{pgfscope}%
\begin{pgfscope}%
\pgfpathrectangle{\pgfqpoint{0.100000in}{2.493063in}}{\pgfqpoint{5.037500in}{3.427208in}}%
\pgfusepath{clip}%
\pgfsetrectcap%
\pgfsetroundjoin%
\pgfsetlinewidth{1.505625pt}%
\definecolor{currentstroke}{rgb}{0.678431,1.000000,0.184314}%
\pgfsetstrokecolor{currentstroke}%
\pgfsetstrokeopacity{0.500000}%
\pgfsetdash{}{0pt}%
\pgfpathmoveto{\pgfqpoint{2.658694in}{4.569578in}}%
\pgfusepath{stroke}%
\end{pgfscope}%
\begin{pgfscope}%
\pgfpathrectangle{\pgfqpoint{0.100000in}{2.493063in}}{\pgfqpoint{5.037500in}{3.427208in}}%
\pgfusepath{clip}%
\pgfsetbuttcap%
\pgfsetroundjoin%
\definecolor{currentfill}{rgb}{0.678431,1.000000,0.184314}%
\pgfsetfillcolor{currentfill}%
\pgfsetfillopacity{0.500000}%
\pgfsetlinewidth{0.250937pt}%
\definecolor{currentstroke}{rgb}{0.000000,0.000000,0.000000}%
\pgfsetstrokecolor{currentstroke}%
\pgfsetstrokeopacity{0.500000}%
\pgfsetdash{}{0pt}%
\pgfsys@defobject{currentmarker}{\pgfqpoint{-0.019444in}{-0.019444in}}{\pgfqpoint{0.019444in}{0.019444in}}{%
\pgfpathmoveto{\pgfqpoint{0.000000in}{-0.019444in}}%
\pgfpathcurveto{\pgfqpoint{0.005157in}{-0.019444in}}{\pgfqpoint{0.010103in}{-0.017396in}}{\pgfqpoint{0.013749in}{-0.013749in}}%
\pgfpathcurveto{\pgfqpoint{0.017396in}{-0.010103in}}{\pgfqpoint{0.019444in}{-0.005157in}}{\pgfqpoint{0.019444in}{0.000000in}}%
\pgfpathcurveto{\pgfqpoint{0.019444in}{0.005157in}}{\pgfqpoint{0.017396in}{0.010103in}}{\pgfqpoint{0.013749in}{0.013749in}}%
\pgfpathcurveto{\pgfqpoint{0.010103in}{0.017396in}}{\pgfqpoint{0.005157in}{0.019444in}}{\pgfqpoint{0.000000in}{0.019444in}}%
\pgfpathcurveto{\pgfqpoint{-0.005157in}{0.019444in}}{\pgfqpoint{-0.010103in}{0.017396in}}{\pgfqpoint{-0.013749in}{0.013749in}}%
\pgfpathcurveto{\pgfqpoint{-0.017396in}{0.010103in}}{\pgfqpoint{-0.019444in}{0.005157in}}{\pgfqpoint{-0.019444in}{0.000000in}}%
\pgfpathcurveto{\pgfqpoint{-0.019444in}{-0.005157in}}{\pgfqpoint{-0.017396in}{-0.010103in}}{\pgfqpoint{-0.013749in}{-0.013749in}}%
\pgfpathcurveto{\pgfqpoint{-0.010103in}{-0.017396in}}{\pgfqpoint{-0.005157in}{-0.019444in}}{\pgfqpoint{0.000000in}{-0.019444in}}%
\pgfpathlineto{\pgfqpoint{0.000000in}{-0.019444in}}%
\pgfpathclose%
\pgfusepath{stroke,fill}%
}%
\begin{pgfscope}%
\pgfsys@transformshift{2.658694in}{4.569578in}%
\pgfsys@useobject{currentmarker}{}%
\end{pgfscope}%
\end{pgfscope}%
\begin{pgfscope}%
\pgfpathrectangle{\pgfqpoint{0.100000in}{2.493063in}}{\pgfqpoint{5.037500in}{3.427208in}}%
\pgfusepath{clip}%
\pgfsetrectcap%
\pgfsetroundjoin%
\pgfsetlinewidth{1.505625pt}%
\definecolor{currentstroke}{rgb}{0.678431,1.000000,0.184314}%
\pgfsetstrokecolor{currentstroke}%
\pgfsetstrokeopacity{0.500000}%
\pgfsetdash{}{0pt}%
\pgfpathmoveto{\pgfqpoint{2.722888in}{4.621515in}}%
\pgfusepath{stroke}%
\end{pgfscope}%
\begin{pgfscope}%
\pgfpathrectangle{\pgfqpoint{0.100000in}{2.493063in}}{\pgfqpoint{5.037500in}{3.427208in}}%
\pgfusepath{clip}%
\pgfsetbuttcap%
\pgfsetroundjoin%
\definecolor{currentfill}{rgb}{0.678431,1.000000,0.184314}%
\pgfsetfillcolor{currentfill}%
\pgfsetfillopacity{0.500000}%
\pgfsetlinewidth{0.250937pt}%
\definecolor{currentstroke}{rgb}{0.000000,0.000000,0.000000}%
\pgfsetstrokecolor{currentstroke}%
\pgfsetstrokeopacity{0.500000}%
\pgfsetdash{}{0pt}%
\pgfsys@defobject{currentmarker}{\pgfqpoint{-0.013889in}{-0.013889in}}{\pgfqpoint{0.013889in}{0.013889in}}{%
\pgfpathmoveto{\pgfqpoint{0.000000in}{-0.013889in}}%
\pgfpathcurveto{\pgfqpoint{0.003683in}{-0.013889in}}{\pgfqpoint{0.007216in}{-0.012425in}}{\pgfqpoint{0.009821in}{-0.009821in}}%
\pgfpathcurveto{\pgfqpoint{0.012425in}{-0.007216in}}{\pgfqpoint{0.013889in}{-0.003683in}}{\pgfqpoint{0.013889in}{0.000000in}}%
\pgfpathcurveto{\pgfqpoint{0.013889in}{0.003683in}}{\pgfqpoint{0.012425in}{0.007216in}}{\pgfqpoint{0.009821in}{0.009821in}}%
\pgfpathcurveto{\pgfqpoint{0.007216in}{0.012425in}}{\pgfqpoint{0.003683in}{0.013889in}}{\pgfqpoint{0.000000in}{0.013889in}}%
\pgfpathcurveto{\pgfqpoint{-0.003683in}{0.013889in}}{\pgfqpoint{-0.007216in}{0.012425in}}{\pgfqpoint{-0.009821in}{0.009821in}}%
\pgfpathcurveto{\pgfqpoint{-0.012425in}{0.007216in}}{\pgfqpoint{-0.013889in}{0.003683in}}{\pgfqpoint{-0.013889in}{0.000000in}}%
\pgfpathcurveto{\pgfqpoint{-0.013889in}{-0.003683in}}{\pgfqpoint{-0.012425in}{-0.007216in}}{\pgfqpoint{-0.009821in}{-0.009821in}}%
\pgfpathcurveto{\pgfqpoint{-0.007216in}{-0.012425in}}{\pgfqpoint{-0.003683in}{-0.013889in}}{\pgfqpoint{0.000000in}{-0.013889in}}%
\pgfpathlineto{\pgfqpoint{0.000000in}{-0.013889in}}%
\pgfpathclose%
\pgfusepath{stroke,fill}%
}%
\begin{pgfscope}%
\pgfsys@transformshift{2.722888in}{4.621515in}%
\pgfsys@useobject{currentmarker}{}%
\end{pgfscope}%
\end{pgfscope}%
\begin{pgfscope}%
\pgfpathrectangle{\pgfqpoint{0.100000in}{2.493063in}}{\pgfqpoint{5.037500in}{3.427208in}}%
\pgfusepath{clip}%
\pgfsetrectcap%
\pgfsetroundjoin%
\pgfsetlinewidth{1.505625pt}%
\definecolor{currentstroke}{rgb}{0.000000,0.000000,1.000000}%
\pgfsetstrokecolor{currentstroke}%
\pgfsetstrokeopacity{0.500000}%
\pgfsetdash{}{0pt}%
\pgfpathmoveto{\pgfqpoint{0.620199in}{4.679268in}}%
\pgfusepath{stroke}%
\end{pgfscope}%
\begin{pgfscope}%
\pgfpathrectangle{\pgfqpoint{0.100000in}{2.493063in}}{\pgfqpoint{5.037500in}{3.427208in}}%
\pgfusepath{clip}%
\pgfsetbuttcap%
\pgfsetroundjoin%
\definecolor{currentfill}{rgb}{0.000000,0.000000,1.000000}%
\pgfsetfillcolor{currentfill}%
\pgfsetfillopacity{0.500000}%
\pgfsetlinewidth{0.250937pt}%
\definecolor{currentstroke}{rgb}{0.000000,0.000000,0.000000}%
\pgfsetstrokecolor{currentstroke}%
\pgfsetstrokeopacity{0.500000}%
\pgfsetdash{}{0pt}%
\pgfsys@defobject{currentmarker}{\pgfqpoint{-0.011111in}{-0.011111in}}{\pgfqpoint{0.011111in}{0.011111in}}{%
\pgfpathmoveto{\pgfqpoint{0.000000in}{-0.011111in}}%
\pgfpathcurveto{\pgfqpoint{0.002947in}{-0.011111in}}{\pgfqpoint{0.005773in}{-0.009940in}}{\pgfqpoint{0.007857in}{-0.007857in}}%
\pgfpathcurveto{\pgfqpoint{0.009940in}{-0.005773in}}{\pgfqpoint{0.011111in}{-0.002947in}}{\pgfqpoint{0.011111in}{0.000000in}}%
\pgfpathcurveto{\pgfqpoint{0.011111in}{0.002947in}}{\pgfqpoint{0.009940in}{0.005773in}}{\pgfqpoint{0.007857in}{0.007857in}}%
\pgfpathcurveto{\pgfqpoint{0.005773in}{0.009940in}}{\pgfqpoint{0.002947in}{0.011111in}}{\pgfqpoint{0.000000in}{0.011111in}}%
\pgfpathcurveto{\pgfqpoint{-0.002947in}{0.011111in}}{\pgfqpoint{-0.005773in}{0.009940in}}{\pgfqpoint{-0.007857in}{0.007857in}}%
\pgfpathcurveto{\pgfqpoint{-0.009940in}{0.005773in}}{\pgfqpoint{-0.011111in}{0.002947in}}{\pgfqpoint{-0.011111in}{0.000000in}}%
\pgfpathcurveto{\pgfqpoint{-0.011111in}{-0.002947in}}{\pgfqpoint{-0.009940in}{-0.005773in}}{\pgfqpoint{-0.007857in}{-0.007857in}}%
\pgfpathcurveto{\pgfqpoint{-0.005773in}{-0.009940in}}{\pgfqpoint{-0.002947in}{-0.011111in}}{\pgfqpoint{0.000000in}{-0.011111in}}%
\pgfpathlineto{\pgfqpoint{0.000000in}{-0.011111in}}%
\pgfpathclose%
\pgfusepath{stroke,fill}%
}%
\begin{pgfscope}%
\pgfsys@transformshift{0.620199in}{4.679268in}%
\pgfsys@useobject{currentmarker}{}%
\end{pgfscope}%
\end{pgfscope}%
\begin{pgfscope}%
\pgfpathrectangle{\pgfqpoint{0.100000in}{2.493063in}}{\pgfqpoint{5.037500in}{3.427208in}}%
\pgfusepath{clip}%
\pgfsetrectcap%
\pgfsetroundjoin%
\pgfsetlinewidth{1.505625pt}%
\definecolor{currentstroke}{rgb}{0.000000,0.000000,1.000000}%
\pgfsetstrokecolor{currentstroke}%
\pgfsetstrokeopacity{0.500000}%
\pgfsetdash{}{0pt}%
\pgfpathmoveto{\pgfqpoint{0.944001in}{4.241574in}}%
\pgfusepath{stroke}%
\end{pgfscope}%
\begin{pgfscope}%
\pgfpathrectangle{\pgfqpoint{0.100000in}{2.493063in}}{\pgfqpoint{5.037500in}{3.427208in}}%
\pgfusepath{clip}%
\pgfsetbuttcap%
\pgfsetroundjoin%
\definecolor{currentfill}{rgb}{0.000000,0.000000,1.000000}%
\pgfsetfillcolor{currentfill}%
\pgfsetfillopacity{0.500000}%
\pgfsetlinewidth{0.250937pt}%
\definecolor{currentstroke}{rgb}{0.000000,0.000000,0.000000}%
\pgfsetstrokecolor{currentstroke}%
\pgfsetstrokeopacity{0.500000}%
\pgfsetdash{}{0pt}%
\pgfsys@defobject{currentmarker}{\pgfqpoint{-0.030556in}{-0.030556in}}{\pgfqpoint{0.030556in}{0.030556in}}{%
\pgfpathmoveto{\pgfqpoint{0.000000in}{-0.030556in}}%
\pgfpathcurveto{\pgfqpoint{0.008103in}{-0.030556in}}{\pgfqpoint{0.015876in}{-0.027336in}}{\pgfqpoint{0.021606in}{-0.021606in}}%
\pgfpathcurveto{\pgfqpoint{0.027336in}{-0.015876in}}{\pgfqpoint{0.030556in}{-0.008103in}}{\pgfqpoint{0.030556in}{0.000000in}}%
\pgfpathcurveto{\pgfqpoint{0.030556in}{0.008103in}}{\pgfqpoint{0.027336in}{0.015876in}}{\pgfqpoint{0.021606in}{0.021606in}}%
\pgfpathcurveto{\pgfqpoint{0.015876in}{0.027336in}}{\pgfqpoint{0.008103in}{0.030556in}}{\pgfqpoint{0.000000in}{0.030556in}}%
\pgfpathcurveto{\pgfqpoint{-0.008103in}{0.030556in}}{\pgfqpoint{-0.015876in}{0.027336in}}{\pgfqpoint{-0.021606in}{0.021606in}}%
\pgfpathcurveto{\pgfqpoint{-0.027336in}{0.015876in}}{\pgfqpoint{-0.030556in}{0.008103in}}{\pgfqpoint{-0.030556in}{0.000000in}}%
\pgfpathcurveto{\pgfqpoint{-0.030556in}{-0.008103in}}{\pgfqpoint{-0.027336in}{-0.015876in}}{\pgfqpoint{-0.021606in}{-0.021606in}}%
\pgfpathcurveto{\pgfqpoint{-0.015876in}{-0.027336in}}{\pgfqpoint{-0.008103in}{-0.030556in}}{\pgfqpoint{0.000000in}{-0.030556in}}%
\pgfpathlineto{\pgfqpoint{0.000000in}{-0.030556in}}%
\pgfpathclose%
\pgfusepath{stroke,fill}%
}%
\begin{pgfscope}%
\pgfsys@transformshift{0.944001in}{4.241574in}%
\pgfsys@useobject{currentmarker}{}%
\end{pgfscope}%
\end{pgfscope}%
\begin{pgfscope}%
\pgfpathrectangle{\pgfqpoint{0.100000in}{2.493063in}}{\pgfqpoint{5.037500in}{3.427208in}}%
\pgfusepath{clip}%
\pgfsetrectcap%
\pgfsetroundjoin%
\pgfsetlinewidth{1.505625pt}%
\definecolor{currentstroke}{rgb}{0.000000,0.000000,1.000000}%
\pgfsetstrokecolor{currentstroke}%
\pgfsetstrokeopacity{0.500000}%
\pgfsetdash{}{0pt}%
\pgfpathmoveto{\pgfqpoint{0.627569in}{4.720906in}}%
\pgfusepath{stroke}%
\end{pgfscope}%
\begin{pgfscope}%
\pgfpathrectangle{\pgfqpoint{0.100000in}{2.493063in}}{\pgfqpoint{5.037500in}{3.427208in}}%
\pgfusepath{clip}%
\pgfsetbuttcap%
\pgfsetroundjoin%
\definecolor{currentfill}{rgb}{0.000000,0.000000,1.000000}%
\pgfsetfillcolor{currentfill}%
\pgfsetfillopacity{0.500000}%
\pgfsetlinewidth{0.250937pt}%
\definecolor{currentstroke}{rgb}{0.000000,0.000000,0.000000}%
\pgfsetstrokecolor{currentstroke}%
\pgfsetstrokeopacity{0.500000}%
\pgfsetdash{}{0pt}%
\pgfsys@defobject{currentmarker}{\pgfqpoint{-0.016667in}{-0.016667in}}{\pgfqpoint{0.016667in}{0.016667in}}{%
\pgfpathmoveto{\pgfqpoint{0.000000in}{-0.016667in}}%
\pgfpathcurveto{\pgfqpoint{0.004420in}{-0.016667in}}{\pgfqpoint{0.008660in}{-0.014911in}}{\pgfqpoint{0.011785in}{-0.011785in}}%
\pgfpathcurveto{\pgfqpoint{0.014911in}{-0.008660in}}{\pgfqpoint{0.016667in}{-0.004420in}}{\pgfqpoint{0.016667in}{0.000000in}}%
\pgfpathcurveto{\pgfqpoint{0.016667in}{0.004420in}}{\pgfqpoint{0.014911in}{0.008660in}}{\pgfqpoint{0.011785in}{0.011785in}}%
\pgfpathcurveto{\pgfqpoint{0.008660in}{0.014911in}}{\pgfqpoint{0.004420in}{0.016667in}}{\pgfqpoint{0.000000in}{0.016667in}}%
\pgfpathcurveto{\pgfqpoint{-0.004420in}{0.016667in}}{\pgfqpoint{-0.008660in}{0.014911in}}{\pgfqpoint{-0.011785in}{0.011785in}}%
\pgfpathcurveto{\pgfqpoint{-0.014911in}{0.008660in}}{\pgfqpoint{-0.016667in}{0.004420in}}{\pgfqpoint{-0.016667in}{0.000000in}}%
\pgfpathcurveto{\pgfqpoint{-0.016667in}{-0.004420in}}{\pgfqpoint{-0.014911in}{-0.008660in}}{\pgfqpoint{-0.011785in}{-0.011785in}}%
\pgfpathcurveto{\pgfqpoint{-0.008660in}{-0.014911in}}{\pgfqpoint{-0.004420in}{-0.016667in}}{\pgfqpoint{0.000000in}{-0.016667in}}%
\pgfpathlineto{\pgfqpoint{0.000000in}{-0.016667in}}%
\pgfpathclose%
\pgfusepath{stroke,fill}%
}%
\begin{pgfscope}%
\pgfsys@transformshift{0.627569in}{4.720906in}%
\pgfsys@useobject{currentmarker}{}%
\end{pgfscope}%
\end{pgfscope}%
\begin{pgfscope}%
\pgfpathrectangle{\pgfqpoint{0.100000in}{2.493063in}}{\pgfqpoint{5.037500in}{3.427208in}}%
\pgfusepath{clip}%
\pgfsetrectcap%
\pgfsetroundjoin%
\pgfsetlinewidth{1.505625pt}%
\definecolor{currentstroke}{rgb}{0.501961,0.501961,0.501961}%
\pgfsetstrokecolor{currentstroke}%
\pgfsetstrokeopacity{0.500000}%
\pgfsetdash{}{0pt}%
\pgfpathmoveto{\pgfqpoint{4.810896in}{5.112692in}}%
\pgfusepath{stroke}%
\end{pgfscope}%
\begin{pgfscope}%
\pgfpathrectangle{\pgfqpoint{0.100000in}{2.493063in}}{\pgfqpoint{5.037500in}{3.427208in}}%
\pgfusepath{clip}%
\pgfsetbuttcap%
\pgfsetroundjoin%
\definecolor{currentfill}{rgb}{0.501961,0.501961,0.501961}%
\pgfsetfillcolor{currentfill}%
\pgfsetfillopacity{0.500000}%
\pgfsetlinewidth{0.250937pt}%
\definecolor{currentstroke}{rgb}{0.000000,0.000000,0.000000}%
\pgfsetstrokecolor{currentstroke}%
\pgfsetstrokeopacity{0.500000}%
\pgfsetdash{}{0pt}%
\pgfsys@defobject{currentmarker}{\pgfqpoint{-0.013889in}{-0.013889in}}{\pgfqpoint{0.013889in}{0.013889in}}{%
\pgfpathmoveto{\pgfqpoint{0.000000in}{-0.013889in}}%
\pgfpathcurveto{\pgfqpoint{0.003683in}{-0.013889in}}{\pgfqpoint{0.007216in}{-0.012425in}}{\pgfqpoint{0.009821in}{-0.009821in}}%
\pgfpathcurveto{\pgfqpoint{0.012425in}{-0.007216in}}{\pgfqpoint{0.013889in}{-0.003683in}}{\pgfqpoint{0.013889in}{0.000000in}}%
\pgfpathcurveto{\pgfqpoint{0.013889in}{0.003683in}}{\pgfqpoint{0.012425in}{0.007216in}}{\pgfqpoint{0.009821in}{0.009821in}}%
\pgfpathcurveto{\pgfqpoint{0.007216in}{0.012425in}}{\pgfqpoint{0.003683in}{0.013889in}}{\pgfqpoint{0.000000in}{0.013889in}}%
\pgfpathcurveto{\pgfqpoint{-0.003683in}{0.013889in}}{\pgfqpoint{-0.007216in}{0.012425in}}{\pgfqpoint{-0.009821in}{0.009821in}}%
\pgfpathcurveto{\pgfqpoint{-0.012425in}{0.007216in}}{\pgfqpoint{-0.013889in}{0.003683in}}{\pgfqpoint{-0.013889in}{0.000000in}}%
\pgfpathcurveto{\pgfqpoint{-0.013889in}{-0.003683in}}{\pgfqpoint{-0.012425in}{-0.007216in}}{\pgfqpoint{-0.009821in}{-0.009821in}}%
\pgfpathcurveto{\pgfqpoint{-0.007216in}{-0.012425in}}{\pgfqpoint{-0.003683in}{-0.013889in}}{\pgfqpoint{0.000000in}{-0.013889in}}%
\pgfpathlineto{\pgfqpoint{0.000000in}{-0.013889in}}%
\pgfpathclose%
\pgfusepath{stroke,fill}%
}%
\begin{pgfscope}%
\pgfsys@transformshift{4.810896in}{5.112692in}%
\pgfsys@useobject{currentmarker}{}%
\end{pgfscope}%
\end{pgfscope}%
\begin{pgfscope}%
\pgfpathrectangle{\pgfqpoint{0.100000in}{2.493063in}}{\pgfqpoint{5.037500in}{3.427208in}}%
\pgfusepath{clip}%
\pgfsetrectcap%
\pgfsetroundjoin%
\pgfsetlinewidth{1.505625pt}%
\definecolor{currentstroke}{rgb}{0.000000,0.000000,1.000000}%
\pgfsetstrokecolor{currentstroke}%
\pgfsetstrokeopacity{0.500000}%
\pgfsetdash{}{0pt}%
\pgfpathmoveto{\pgfqpoint{4.769641in}{5.077446in}}%
\pgfusepath{stroke}%
\end{pgfscope}%
\begin{pgfscope}%
\pgfpathrectangle{\pgfqpoint{0.100000in}{2.493063in}}{\pgfqpoint{5.037500in}{3.427208in}}%
\pgfusepath{clip}%
\pgfsetbuttcap%
\pgfsetroundjoin%
\definecolor{currentfill}{rgb}{0.000000,0.000000,1.000000}%
\pgfsetfillcolor{currentfill}%
\pgfsetfillopacity{0.500000}%
\pgfsetlinewidth{0.250937pt}%
\definecolor{currentstroke}{rgb}{0.000000,0.000000,0.000000}%
\pgfsetstrokecolor{currentstroke}%
\pgfsetstrokeopacity{0.500000}%
\pgfsetdash{}{0pt}%
\pgfsys@defobject{currentmarker}{\pgfqpoint{-0.008333in}{-0.008333in}}{\pgfqpoint{0.008333in}{0.008333in}}{%
\pgfpathmoveto{\pgfqpoint{0.000000in}{-0.008333in}}%
\pgfpathcurveto{\pgfqpoint{0.002210in}{-0.008333in}}{\pgfqpoint{0.004330in}{-0.007455in}}{\pgfqpoint{0.005893in}{-0.005893in}}%
\pgfpathcurveto{\pgfqpoint{0.007455in}{-0.004330in}}{\pgfqpoint{0.008333in}{-0.002210in}}{\pgfqpoint{0.008333in}{0.000000in}}%
\pgfpathcurveto{\pgfqpoint{0.008333in}{0.002210in}}{\pgfqpoint{0.007455in}{0.004330in}}{\pgfqpoint{0.005893in}{0.005893in}}%
\pgfpathcurveto{\pgfqpoint{0.004330in}{0.007455in}}{\pgfqpoint{0.002210in}{0.008333in}}{\pgfqpoint{0.000000in}{0.008333in}}%
\pgfpathcurveto{\pgfqpoint{-0.002210in}{0.008333in}}{\pgfqpoint{-0.004330in}{0.007455in}}{\pgfqpoint{-0.005893in}{0.005893in}}%
\pgfpathcurveto{\pgfqpoint{-0.007455in}{0.004330in}}{\pgfqpoint{-0.008333in}{0.002210in}}{\pgfqpoint{-0.008333in}{0.000000in}}%
\pgfpathcurveto{\pgfqpoint{-0.008333in}{-0.002210in}}{\pgfqpoint{-0.007455in}{-0.004330in}}{\pgfqpoint{-0.005893in}{-0.005893in}}%
\pgfpathcurveto{\pgfqpoint{-0.004330in}{-0.007455in}}{\pgfqpoint{-0.002210in}{-0.008333in}}{\pgfqpoint{0.000000in}{-0.008333in}}%
\pgfpathlineto{\pgfqpoint{0.000000in}{-0.008333in}}%
\pgfpathclose%
\pgfusepath{stroke,fill}%
}%
\begin{pgfscope}%
\pgfsys@transformshift{4.769641in}{5.077446in}%
\pgfsys@useobject{currentmarker}{}%
\end{pgfscope}%
\end{pgfscope}%
\begin{pgfscope}%
\pgfpathrectangle{\pgfqpoint{0.100000in}{2.493063in}}{\pgfqpoint{5.037500in}{3.427208in}}%
\pgfusepath{clip}%
\pgfsetrectcap%
\pgfsetroundjoin%
\pgfsetlinewidth{1.505625pt}%
\definecolor{currentstroke}{rgb}{0.000000,0.000000,1.000000}%
\pgfsetstrokecolor{currentstroke}%
\pgfsetstrokeopacity{0.500000}%
\pgfsetdash{}{0pt}%
\pgfpathmoveto{\pgfqpoint{4.823756in}{5.101108in}}%
\pgfusepath{stroke}%
\end{pgfscope}%
\begin{pgfscope}%
\pgfpathrectangle{\pgfqpoint{0.100000in}{2.493063in}}{\pgfqpoint{5.037500in}{3.427208in}}%
\pgfusepath{clip}%
\pgfsetbuttcap%
\pgfsetroundjoin%
\definecolor{currentfill}{rgb}{0.000000,0.000000,1.000000}%
\pgfsetfillcolor{currentfill}%
\pgfsetfillopacity{0.500000}%
\pgfsetlinewidth{0.250937pt}%
\definecolor{currentstroke}{rgb}{0.000000,0.000000,0.000000}%
\pgfsetstrokecolor{currentstroke}%
\pgfsetstrokeopacity{0.500000}%
\pgfsetdash{}{0pt}%
\pgfsys@defobject{currentmarker}{\pgfqpoint{-0.008333in}{-0.008333in}}{\pgfqpoint{0.008333in}{0.008333in}}{%
\pgfpathmoveto{\pgfqpoint{0.000000in}{-0.008333in}}%
\pgfpathcurveto{\pgfqpoint{0.002210in}{-0.008333in}}{\pgfqpoint{0.004330in}{-0.007455in}}{\pgfqpoint{0.005893in}{-0.005893in}}%
\pgfpathcurveto{\pgfqpoint{0.007455in}{-0.004330in}}{\pgfqpoint{0.008333in}{-0.002210in}}{\pgfqpoint{0.008333in}{0.000000in}}%
\pgfpathcurveto{\pgfqpoint{0.008333in}{0.002210in}}{\pgfqpoint{0.007455in}{0.004330in}}{\pgfqpoint{0.005893in}{0.005893in}}%
\pgfpathcurveto{\pgfqpoint{0.004330in}{0.007455in}}{\pgfqpoint{0.002210in}{0.008333in}}{\pgfqpoint{0.000000in}{0.008333in}}%
\pgfpathcurveto{\pgfqpoint{-0.002210in}{0.008333in}}{\pgfqpoint{-0.004330in}{0.007455in}}{\pgfqpoint{-0.005893in}{0.005893in}}%
\pgfpathcurveto{\pgfqpoint{-0.007455in}{0.004330in}}{\pgfqpoint{-0.008333in}{0.002210in}}{\pgfqpoint{-0.008333in}{0.000000in}}%
\pgfpathcurveto{\pgfqpoint{-0.008333in}{-0.002210in}}{\pgfqpoint{-0.007455in}{-0.004330in}}{\pgfqpoint{-0.005893in}{-0.005893in}}%
\pgfpathcurveto{\pgfqpoint{-0.004330in}{-0.007455in}}{\pgfqpoint{-0.002210in}{-0.008333in}}{\pgfqpoint{0.000000in}{-0.008333in}}%
\pgfpathlineto{\pgfqpoint{0.000000in}{-0.008333in}}%
\pgfpathclose%
\pgfusepath{stroke,fill}%
}%
\begin{pgfscope}%
\pgfsys@transformshift{4.823756in}{5.101108in}%
\pgfsys@useobject{currentmarker}{}%
\end{pgfscope}%
\end{pgfscope}%
\begin{pgfscope}%
\pgfpathrectangle{\pgfqpoint{0.100000in}{2.493063in}}{\pgfqpoint{5.037500in}{3.427208in}}%
\pgfusepath{clip}%
\pgfsetrectcap%
\pgfsetroundjoin%
\pgfsetlinewidth{1.505625pt}%
\definecolor{currentstroke}{rgb}{0.000000,0.000000,1.000000}%
\pgfsetstrokecolor{currentstroke}%
\pgfsetstrokeopacity{0.500000}%
\pgfsetdash{}{0pt}%
\pgfpathmoveto{\pgfqpoint{4.620381in}{4.609404in}}%
\pgfusepath{stroke}%
\end{pgfscope}%
\begin{pgfscope}%
\pgfpathrectangle{\pgfqpoint{0.100000in}{2.493063in}}{\pgfqpoint{5.037500in}{3.427208in}}%
\pgfusepath{clip}%
\pgfsetbuttcap%
\pgfsetroundjoin%
\definecolor{currentfill}{rgb}{0.000000,0.000000,1.000000}%
\pgfsetfillcolor{currentfill}%
\pgfsetfillopacity{0.500000}%
\pgfsetlinewidth{0.250937pt}%
\definecolor{currentstroke}{rgb}{0.000000,0.000000,0.000000}%
\pgfsetstrokecolor{currentstroke}%
\pgfsetstrokeopacity{0.500000}%
\pgfsetdash{}{0pt}%
\pgfsys@defobject{currentmarker}{\pgfqpoint{-0.022222in}{-0.022222in}}{\pgfqpoint{0.022222in}{0.022222in}}{%
\pgfpathmoveto{\pgfqpoint{0.000000in}{-0.022222in}}%
\pgfpathcurveto{\pgfqpoint{0.005893in}{-0.022222in}}{\pgfqpoint{0.011546in}{-0.019881in}}{\pgfqpoint{0.015713in}{-0.015713in}}%
\pgfpathcurveto{\pgfqpoint{0.019881in}{-0.011546in}}{\pgfqpoint{0.022222in}{-0.005893in}}{\pgfqpoint{0.022222in}{0.000000in}}%
\pgfpathcurveto{\pgfqpoint{0.022222in}{0.005893in}}{\pgfqpoint{0.019881in}{0.011546in}}{\pgfqpoint{0.015713in}{0.015713in}}%
\pgfpathcurveto{\pgfqpoint{0.011546in}{0.019881in}}{\pgfqpoint{0.005893in}{0.022222in}}{\pgfqpoint{0.000000in}{0.022222in}}%
\pgfpathcurveto{\pgfqpoint{-0.005893in}{0.022222in}}{\pgfqpoint{-0.011546in}{0.019881in}}{\pgfqpoint{-0.015713in}{0.015713in}}%
\pgfpathcurveto{\pgfqpoint{-0.019881in}{0.011546in}}{\pgfqpoint{-0.022222in}{0.005893in}}{\pgfqpoint{-0.022222in}{0.000000in}}%
\pgfpathcurveto{\pgfqpoint{-0.022222in}{-0.005893in}}{\pgfqpoint{-0.019881in}{-0.011546in}}{\pgfqpoint{-0.015713in}{-0.015713in}}%
\pgfpathcurveto{\pgfqpoint{-0.011546in}{-0.019881in}}{\pgfqpoint{-0.005893in}{-0.022222in}}{\pgfqpoint{0.000000in}{-0.022222in}}%
\pgfpathlineto{\pgfqpoint{0.000000in}{-0.022222in}}%
\pgfpathclose%
\pgfusepath{stroke,fill}%
}%
\begin{pgfscope}%
\pgfsys@transformshift{4.620381in}{4.609404in}%
\pgfsys@useobject{currentmarker}{}%
\end{pgfscope}%
\end{pgfscope}%
\begin{pgfscope}%
\pgfpathrectangle{\pgfqpoint{0.100000in}{2.493063in}}{\pgfqpoint{5.037500in}{3.427208in}}%
\pgfusepath{clip}%
\pgfsetrectcap%
\pgfsetroundjoin%
\pgfsetlinewidth{1.505625pt}%
\definecolor{currentstroke}{rgb}{0.678431,1.000000,0.184314}%
\pgfsetstrokecolor{currentstroke}%
\pgfsetstrokeopacity{0.500000}%
\pgfsetdash{}{0pt}%
\pgfpathmoveto{\pgfqpoint{4.609445in}{4.596655in}}%
\pgfusepath{stroke}%
\end{pgfscope}%
\begin{pgfscope}%
\pgfpathrectangle{\pgfqpoint{0.100000in}{2.493063in}}{\pgfqpoint{5.037500in}{3.427208in}}%
\pgfusepath{clip}%
\pgfsetbuttcap%
\pgfsetroundjoin%
\definecolor{currentfill}{rgb}{0.678431,1.000000,0.184314}%
\pgfsetfillcolor{currentfill}%
\pgfsetfillopacity{0.500000}%
\pgfsetlinewidth{0.250937pt}%
\definecolor{currentstroke}{rgb}{0.000000,0.000000,0.000000}%
\pgfsetstrokecolor{currentstroke}%
\pgfsetstrokeopacity{0.500000}%
\pgfsetdash{}{0pt}%
\pgfsys@defobject{currentmarker}{\pgfqpoint{-0.016667in}{-0.016667in}}{\pgfqpoint{0.016667in}{0.016667in}}{%
\pgfpathmoveto{\pgfqpoint{0.000000in}{-0.016667in}}%
\pgfpathcurveto{\pgfqpoint{0.004420in}{-0.016667in}}{\pgfqpoint{0.008660in}{-0.014911in}}{\pgfqpoint{0.011785in}{-0.011785in}}%
\pgfpathcurveto{\pgfqpoint{0.014911in}{-0.008660in}}{\pgfqpoint{0.016667in}{-0.004420in}}{\pgfqpoint{0.016667in}{0.000000in}}%
\pgfpathcurveto{\pgfqpoint{0.016667in}{0.004420in}}{\pgfqpoint{0.014911in}{0.008660in}}{\pgfqpoint{0.011785in}{0.011785in}}%
\pgfpathcurveto{\pgfqpoint{0.008660in}{0.014911in}}{\pgfqpoint{0.004420in}{0.016667in}}{\pgfqpoint{0.000000in}{0.016667in}}%
\pgfpathcurveto{\pgfqpoint{-0.004420in}{0.016667in}}{\pgfqpoint{-0.008660in}{0.014911in}}{\pgfqpoint{-0.011785in}{0.011785in}}%
\pgfpathcurveto{\pgfqpoint{-0.014911in}{0.008660in}}{\pgfqpoint{-0.016667in}{0.004420in}}{\pgfqpoint{-0.016667in}{0.000000in}}%
\pgfpathcurveto{\pgfqpoint{-0.016667in}{-0.004420in}}{\pgfqpoint{-0.014911in}{-0.008660in}}{\pgfqpoint{-0.011785in}{-0.011785in}}%
\pgfpathcurveto{\pgfqpoint{-0.008660in}{-0.014911in}}{\pgfqpoint{-0.004420in}{-0.016667in}}{\pgfqpoint{0.000000in}{-0.016667in}}%
\pgfpathlineto{\pgfqpoint{0.000000in}{-0.016667in}}%
\pgfpathclose%
\pgfusepath{stroke,fill}%
}%
\begin{pgfscope}%
\pgfsys@transformshift{4.609445in}{4.596655in}%
\pgfsys@useobject{currentmarker}{}%
\end{pgfscope}%
\end{pgfscope}%
\begin{pgfscope}%
\pgfpathrectangle{\pgfqpoint{0.100000in}{2.493063in}}{\pgfqpoint{5.037500in}{3.427208in}}%
\pgfusepath{clip}%
\pgfsetrectcap%
\pgfsetroundjoin%
\pgfsetlinewidth{1.505625pt}%
\definecolor{currentstroke}{rgb}{0.000000,0.000000,1.000000}%
\pgfsetstrokecolor{currentstroke}%
\pgfsetstrokeopacity{0.500000}%
\pgfsetdash{}{0pt}%
\pgfpathmoveto{\pgfqpoint{4.570899in}{4.698831in}}%
\pgfusepath{stroke}%
\end{pgfscope}%
\begin{pgfscope}%
\pgfpathrectangle{\pgfqpoint{0.100000in}{2.493063in}}{\pgfqpoint{5.037500in}{3.427208in}}%
\pgfusepath{clip}%
\pgfsetbuttcap%
\pgfsetroundjoin%
\definecolor{currentfill}{rgb}{0.000000,0.000000,1.000000}%
\pgfsetfillcolor{currentfill}%
\pgfsetfillopacity{0.500000}%
\pgfsetlinewidth{0.250937pt}%
\definecolor{currentstroke}{rgb}{0.000000,0.000000,0.000000}%
\pgfsetstrokecolor{currentstroke}%
\pgfsetstrokeopacity{0.500000}%
\pgfsetdash{}{0pt}%
\pgfsys@defobject{currentmarker}{\pgfqpoint{-0.022222in}{-0.022222in}}{\pgfqpoint{0.022222in}{0.022222in}}{%
\pgfpathmoveto{\pgfqpoint{0.000000in}{-0.022222in}}%
\pgfpathcurveto{\pgfqpoint{0.005893in}{-0.022222in}}{\pgfqpoint{0.011546in}{-0.019881in}}{\pgfqpoint{0.015713in}{-0.015713in}}%
\pgfpathcurveto{\pgfqpoint{0.019881in}{-0.011546in}}{\pgfqpoint{0.022222in}{-0.005893in}}{\pgfqpoint{0.022222in}{0.000000in}}%
\pgfpathcurveto{\pgfqpoint{0.022222in}{0.005893in}}{\pgfqpoint{0.019881in}{0.011546in}}{\pgfqpoint{0.015713in}{0.015713in}}%
\pgfpathcurveto{\pgfqpoint{0.011546in}{0.019881in}}{\pgfqpoint{0.005893in}{0.022222in}}{\pgfqpoint{0.000000in}{0.022222in}}%
\pgfpathcurveto{\pgfqpoint{-0.005893in}{0.022222in}}{\pgfqpoint{-0.011546in}{0.019881in}}{\pgfqpoint{-0.015713in}{0.015713in}}%
\pgfpathcurveto{\pgfqpoint{-0.019881in}{0.011546in}}{\pgfqpoint{-0.022222in}{0.005893in}}{\pgfqpoint{-0.022222in}{0.000000in}}%
\pgfpathcurveto{\pgfqpoint{-0.022222in}{-0.005893in}}{\pgfqpoint{-0.019881in}{-0.011546in}}{\pgfqpoint{-0.015713in}{-0.015713in}}%
\pgfpathcurveto{\pgfqpoint{-0.011546in}{-0.019881in}}{\pgfqpoint{-0.005893in}{-0.022222in}}{\pgfqpoint{0.000000in}{-0.022222in}}%
\pgfpathlineto{\pgfqpoint{0.000000in}{-0.022222in}}%
\pgfpathclose%
\pgfusepath{stroke,fill}%
}%
\begin{pgfscope}%
\pgfsys@transformshift{4.570899in}{4.698831in}%
\pgfsys@useobject{currentmarker}{}%
\end{pgfscope}%
\end{pgfscope}%
\begin{pgfscope}%
\pgfpathrectangle{\pgfqpoint{0.100000in}{2.493063in}}{\pgfqpoint{5.037500in}{3.427208in}}%
\pgfusepath{clip}%
\pgfsetrectcap%
\pgfsetroundjoin%
\pgfsetlinewidth{1.505625pt}%
\definecolor{currentstroke}{rgb}{0.000000,0.000000,1.000000}%
\pgfsetstrokecolor{currentstroke}%
\pgfsetstrokeopacity{0.500000}%
\pgfsetdash{}{0pt}%
\pgfpathmoveto{\pgfqpoint{4.567371in}{4.610224in}}%
\pgfusepath{stroke}%
\end{pgfscope}%
\begin{pgfscope}%
\pgfpathrectangle{\pgfqpoint{0.100000in}{2.493063in}}{\pgfqpoint{5.037500in}{3.427208in}}%
\pgfusepath{clip}%
\pgfsetbuttcap%
\pgfsetroundjoin%
\definecolor{currentfill}{rgb}{0.000000,0.000000,1.000000}%
\pgfsetfillcolor{currentfill}%
\pgfsetfillopacity{0.500000}%
\pgfsetlinewidth{0.250937pt}%
\definecolor{currentstroke}{rgb}{0.000000,0.000000,0.000000}%
\pgfsetstrokecolor{currentstroke}%
\pgfsetstrokeopacity{0.500000}%
\pgfsetdash{}{0pt}%
\pgfsys@defobject{currentmarker}{\pgfqpoint{-0.027778in}{-0.027778in}}{\pgfqpoint{0.027778in}{0.027778in}}{%
\pgfpathmoveto{\pgfqpoint{0.000000in}{-0.027778in}}%
\pgfpathcurveto{\pgfqpoint{0.007367in}{-0.027778in}}{\pgfqpoint{0.014433in}{-0.024851in}}{\pgfqpoint{0.019642in}{-0.019642in}}%
\pgfpathcurveto{\pgfqpoint{0.024851in}{-0.014433in}}{\pgfqpoint{0.027778in}{-0.007367in}}{\pgfqpoint{0.027778in}{0.000000in}}%
\pgfpathcurveto{\pgfqpoint{0.027778in}{0.007367in}}{\pgfqpoint{0.024851in}{0.014433in}}{\pgfqpoint{0.019642in}{0.019642in}}%
\pgfpathcurveto{\pgfqpoint{0.014433in}{0.024851in}}{\pgfqpoint{0.007367in}{0.027778in}}{\pgfqpoint{0.000000in}{0.027778in}}%
\pgfpathcurveto{\pgfqpoint{-0.007367in}{0.027778in}}{\pgfqpoint{-0.014433in}{0.024851in}}{\pgfqpoint{-0.019642in}{0.019642in}}%
\pgfpathcurveto{\pgfqpoint{-0.024851in}{0.014433in}}{\pgfqpoint{-0.027778in}{0.007367in}}{\pgfqpoint{-0.027778in}{0.000000in}}%
\pgfpathcurveto{\pgfqpoint{-0.027778in}{-0.007367in}}{\pgfqpoint{-0.024851in}{-0.014433in}}{\pgfqpoint{-0.019642in}{-0.019642in}}%
\pgfpathcurveto{\pgfqpoint{-0.014433in}{-0.024851in}}{\pgfqpoint{-0.007367in}{-0.027778in}}{\pgfqpoint{0.000000in}{-0.027778in}}%
\pgfpathlineto{\pgfqpoint{0.000000in}{-0.027778in}}%
\pgfpathclose%
\pgfusepath{stroke,fill}%
}%
\begin{pgfscope}%
\pgfsys@transformshift{4.567371in}{4.610224in}%
\pgfsys@useobject{currentmarker}{}%
\end{pgfscope}%
\end{pgfscope}%
\begin{pgfscope}%
\pgfpathrectangle{\pgfqpoint{0.100000in}{2.493063in}}{\pgfqpoint{5.037500in}{3.427208in}}%
\pgfusepath{clip}%
\pgfsetrectcap%
\pgfsetroundjoin%
\pgfsetlinewidth{1.505625pt}%
\definecolor{currentstroke}{rgb}{0.678431,1.000000,0.184314}%
\pgfsetstrokecolor{currentstroke}%
\pgfsetstrokeopacity{0.500000}%
\pgfsetdash{}{0pt}%
\pgfpathmoveto{\pgfqpoint{1.706681in}{3.980122in}}%
\pgfusepath{stroke}%
\end{pgfscope}%
\begin{pgfscope}%
\pgfpathrectangle{\pgfqpoint{0.100000in}{2.493063in}}{\pgfqpoint{5.037500in}{3.427208in}}%
\pgfusepath{clip}%
\pgfsetbuttcap%
\pgfsetroundjoin%
\definecolor{currentfill}{rgb}{0.678431,1.000000,0.184314}%
\pgfsetfillcolor{currentfill}%
\pgfsetfillopacity{0.500000}%
\pgfsetlinewidth{0.250937pt}%
\definecolor{currentstroke}{rgb}{0.000000,0.000000,0.000000}%
\pgfsetstrokecolor{currentstroke}%
\pgfsetstrokeopacity{0.500000}%
\pgfsetdash{}{0pt}%
\pgfsys@defobject{currentmarker}{\pgfqpoint{-0.019444in}{-0.019444in}}{\pgfqpoint{0.019444in}{0.019444in}}{%
\pgfpathmoveto{\pgfqpoint{0.000000in}{-0.019444in}}%
\pgfpathcurveto{\pgfqpoint{0.005157in}{-0.019444in}}{\pgfqpoint{0.010103in}{-0.017396in}}{\pgfqpoint{0.013749in}{-0.013749in}}%
\pgfpathcurveto{\pgfqpoint{0.017396in}{-0.010103in}}{\pgfqpoint{0.019444in}{-0.005157in}}{\pgfqpoint{0.019444in}{0.000000in}}%
\pgfpathcurveto{\pgfqpoint{0.019444in}{0.005157in}}{\pgfqpoint{0.017396in}{0.010103in}}{\pgfqpoint{0.013749in}{0.013749in}}%
\pgfpathcurveto{\pgfqpoint{0.010103in}{0.017396in}}{\pgfqpoint{0.005157in}{0.019444in}}{\pgfqpoint{0.000000in}{0.019444in}}%
\pgfpathcurveto{\pgfqpoint{-0.005157in}{0.019444in}}{\pgfqpoint{-0.010103in}{0.017396in}}{\pgfqpoint{-0.013749in}{0.013749in}}%
\pgfpathcurveto{\pgfqpoint{-0.017396in}{0.010103in}}{\pgfqpoint{-0.019444in}{0.005157in}}{\pgfqpoint{-0.019444in}{0.000000in}}%
\pgfpathcurveto{\pgfqpoint{-0.019444in}{-0.005157in}}{\pgfqpoint{-0.017396in}{-0.010103in}}{\pgfqpoint{-0.013749in}{-0.013749in}}%
\pgfpathcurveto{\pgfqpoint{-0.010103in}{-0.017396in}}{\pgfqpoint{-0.005157in}{-0.019444in}}{\pgfqpoint{0.000000in}{-0.019444in}}%
\pgfpathlineto{\pgfqpoint{0.000000in}{-0.019444in}}%
\pgfpathclose%
\pgfusepath{stroke,fill}%
}%
\begin{pgfscope}%
\pgfsys@transformshift{1.706681in}{3.980122in}%
\pgfsys@useobject{currentmarker}{}%
\end{pgfscope}%
\end{pgfscope}%
\begin{pgfscope}%
\pgfpathrectangle{\pgfqpoint{0.100000in}{2.493063in}}{\pgfqpoint{5.037500in}{3.427208in}}%
\pgfusepath{clip}%
\pgfsetrectcap%
\pgfsetroundjoin%
\pgfsetlinewidth{1.505625pt}%
\definecolor{currentstroke}{rgb}{0.678431,1.000000,0.184314}%
\pgfsetstrokecolor{currentstroke}%
\pgfsetstrokeopacity{0.500000}%
\pgfsetdash{}{0pt}%
\pgfpathmoveto{\pgfqpoint{1.588407in}{4.188047in}}%
\pgfusepath{stroke}%
\end{pgfscope}%
\begin{pgfscope}%
\pgfpathrectangle{\pgfqpoint{0.100000in}{2.493063in}}{\pgfqpoint{5.037500in}{3.427208in}}%
\pgfusepath{clip}%
\pgfsetbuttcap%
\pgfsetroundjoin%
\definecolor{currentfill}{rgb}{0.678431,1.000000,0.184314}%
\pgfsetfillcolor{currentfill}%
\pgfsetfillopacity{0.500000}%
\pgfsetlinewidth{0.250937pt}%
\definecolor{currentstroke}{rgb}{0.000000,0.000000,0.000000}%
\pgfsetstrokecolor{currentstroke}%
\pgfsetstrokeopacity{0.500000}%
\pgfsetdash{}{0pt}%
\pgfsys@defobject{currentmarker}{\pgfqpoint{-0.044444in}{-0.044444in}}{\pgfqpoint{0.044444in}{0.044444in}}{%
\pgfpathmoveto{\pgfqpoint{0.000000in}{-0.044444in}}%
\pgfpathcurveto{\pgfqpoint{0.011787in}{-0.044444in}}{\pgfqpoint{0.023092in}{-0.039761in}}{\pgfqpoint{0.031427in}{-0.031427in}}%
\pgfpathcurveto{\pgfqpoint{0.039761in}{-0.023092in}}{\pgfqpoint{0.044444in}{-0.011787in}}{\pgfqpoint{0.044444in}{0.000000in}}%
\pgfpathcurveto{\pgfqpoint{0.044444in}{0.011787in}}{\pgfqpoint{0.039761in}{0.023092in}}{\pgfqpoint{0.031427in}{0.031427in}}%
\pgfpathcurveto{\pgfqpoint{0.023092in}{0.039761in}}{\pgfqpoint{0.011787in}{0.044444in}}{\pgfqpoint{0.000000in}{0.044444in}}%
\pgfpathcurveto{\pgfqpoint{-0.011787in}{0.044444in}}{\pgfqpoint{-0.023092in}{0.039761in}}{\pgfqpoint{-0.031427in}{0.031427in}}%
\pgfpathcurveto{\pgfqpoint{-0.039761in}{0.023092in}}{\pgfqpoint{-0.044444in}{0.011787in}}{\pgfqpoint{-0.044444in}{0.000000in}}%
\pgfpathcurveto{\pgfqpoint{-0.044444in}{-0.011787in}}{\pgfqpoint{-0.039761in}{-0.023092in}}{\pgfqpoint{-0.031427in}{-0.031427in}}%
\pgfpathcurveto{\pgfqpoint{-0.023092in}{-0.039761in}}{\pgfqpoint{-0.011787in}{-0.044444in}}{\pgfqpoint{0.000000in}{-0.044444in}}%
\pgfpathlineto{\pgfqpoint{0.000000in}{-0.044444in}}%
\pgfpathclose%
\pgfusepath{stroke,fill}%
}%
\begin{pgfscope}%
\pgfsys@transformshift{1.588407in}{4.188047in}%
\pgfsys@useobject{currentmarker}{}%
\end{pgfscope}%
\end{pgfscope}%
\begin{pgfscope}%
\pgfpathrectangle{\pgfqpoint{0.100000in}{2.493063in}}{\pgfqpoint{5.037500in}{3.427208in}}%
\pgfusepath{clip}%
\pgfsetrectcap%
\pgfsetroundjoin%
\pgfsetlinewidth{1.505625pt}%
\definecolor{currentstroke}{rgb}{0.678431,1.000000,0.184314}%
\pgfsetstrokecolor{currentstroke}%
\pgfsetstrokeopacity{0.500000}%
\pgfsetdash{}{0pt}%
\pgfpathmoveto{\pgfqpoint{1.653203in}{3.663998in}}%
\pgfusepath{stroke}%
\end{pgfscope}%
\begin{pgfscope}%
\pgfpathrectangle{\pgfqpoint{0.100000in}{2.493063in}}{\pgfqpoint{5.037500in}{3.427208in}}%
\pgfusepath{clip}%
\pgfsetbuttcap%
\pgfsetroundjoin%
\definecolor{currentfill}{rgb}{0.678431,1.000000,0.184314}%
\pgfsetfillcolor{currentfill}%
\pgfsetfillopacity{0.500000}%
\pgfsetlinewidth{0.250937pt}%
\definecolor{currentstroke}{rgb}{0.000000,0.000000,0.000000}%
\pgfsetstrokecolor{currentstroke}%
\pgfsetstrokeopacity{0.500000}%
\pgfsetdash{}{0pt}%
\pgfsys@defobject{currentmarker}{\pgfqpoint{-0.033333in}{-0.033333in}}{\pgfqpoint{0.033333in}{0.033333in}}{%
\pgfpathmoveto{\pgfqpoint{0.000000in}{-0.033333in}}%
\pgfpathcurveto{\pgfqpoint{0.008840in}{-0.033333in}}{\pgfqpoint{0.017319in}{-0.029821in}}{\pgfqpoint{0.023570in}{-0.023570in}}%
\pgfpathcurveto{\pgfqpoint{0.029821in}{-0.017319in}}{\pgfqpoint{0.033333in}{-0.008840in}}{\pgfqpoint{0.033333in}{0.000000in}}%
\pgfpathcurveto{\pgfqpoint{0.033333in}{0.008840in}}{\pgfqpoint{0.029821in}{0.017319in}}{\pgfqpoint{0.023570in}{0.023570in}}%
\pgfpathcurveto{\pgfqpoint{0.017319in}{0.029821in}}{\pgfqpoint{0.008840in}{0.033333in}}{\pgfqpoint{0.000000in}{0.033333in}}%
\pgfpathcurveto{\pgfqpoint{-0.008840in}{0.033333in}}{\pgfqpoint{-0.017319in}{0.029821in}}{\pgfqpoint{-0.023570in}{0.023570in}}%
\pgfpathcurveto{\pgfqpoint{-0.029821in}{0.017319in}}{\pgfqpoint{-0.033333in}{0.008840in}}{\pgfqpoint{-0.033333in}{0.000000in}}%
\pgfpathcurveto{\pgfqpoint{-0.033333in}{-0.008840in}}{\pgfqpoint{-0.029821in}{-0.017319in}}{\pgfqpoint{-0.023570in}{-0.023570in}}%
\pgfpathcurveto{\pgfqpoint{-0.017319in}{-0.029821in}}{\pgfqpoint{-0.008840in}{-0.033333in}}{\pgfqpoint{0.000000in}{-0.033333in}}%
\pgfpathlineto{\pgfqpoint{0.000000in}{-0.033333in}}%
\pgfpathclose%
\pgfusepath{stroke,fill}%
}%
\begin{pgfscope}%
\pgfsys@transformshift{1.653203in}{3.663998in}%
\pgfsys@useobject{currentmarker}{}%
\end{pgfscope}%
\end{pgfscope}%
\begin{pgfscope}%
\pgfpathrectangle{\pgfqpoint{0.100000in}{2.493063in}}{\pgfqpoint{5.037500in}{3.427208in}}%
\pgfusepath{clip}%
\pgfsetrectcap%
\pgfsetroundjoin%
\pgfsetlinewidth{1.505625pt}%
\definecolor{currentstroke}{rgb}{0.678431,1.000000,0.184314}%
\pgfsetstrokecolor{currentstroke}%
\pgfsetstrokeopacity{0.500000}%
\pgfsetdash{}{0pt}%
\pgfpathmoveto{\pgfqpoint{1.781876in}{4.040833in}}%
\pgfusepath{stroke}%
\end{pgfscope}%
\begin{pgfscope}%
\pgfpathrectangle{\pgfqpoint{0.100000in}{2.493063in}}{\pgfqpoint{5.037500in}{3.427208in}}%
\pgfusepath{clip}%
\pgfsetbuttcap%
\pgfsetroundjoin%
\definecolor{currentfill}{rgb}{0.678431,1.000000,0.184314}%
\pgfsetfillcolor{currentfill}%
\pgfsetfillopacity{0.500000}%
\pgfsetlinewidth{0.250937pt}%
\definecolor{currentstroke}{rgb}{0.000000,0.000000,0.000000}%
\pgfsetstrokecolor{currentstroke}%
\pgfsetstrokeopacity{0.500000}%
\pgfsetdash{}{0pt}%
\pgfsys@defobject{currentmarker}{\pgfqpoint{-0.008333in}{-0.008333in}}{\pgfqpoint{0.008333in}{0.008333in}}{%
\pgfpathmoveto{\pgfqpoint{0.000000in}{-0.008333in}}%
\pgfpathcurveto{\pgfqpoint{0.002210in}{-0.008333in}}{\pgfqpoint{0.004330in}{-0.007455in}}{\pgfqpoint{0.005893in}{-0.005893in}}%
\pgfpathcurveto{\pgfqpoint{0.007455in}{-0.004330in}}{\pgfqpoint{0.008333in}{-0.002210in}}{\pgfqpoint{0.008333in}{0.000000in}}%
\pgfpathcurveto{\pgfqpoint{0.008333in}{0.002210in}}{\pgfqpoint{0.007455in}{0.004330in}}{\pgfqpoint{0.005893in}{0.005893in}}%
\pgfpathcurveto{\pgfqpoint{0.004330in}{0.007455in}}{\pgfqpoint{0.002210in}{0.008333in}}{\pgfqpoint{0.000000in}{0.008333in}}%
\pgfpathcurveto{\pgfqpoint{-0.002210in}{0.008333in}}{\pgfqpoint{-0.004330in}{0.007455in}}{\pgfqpoint{-0.005893in}{0.005893in}}%
\pgfpathcurveto{\pgfqpoint{-0.007455in}{0.004330in}}{\pgfqpoint{-0.008333in}{0.002210in}}{\pgfqpoint{-0.008333in}{0.000000in}}%
\pgfpathcurveto{\pgfqpoint{-0.008333in}{-0.002210in}}{\pgfqpoint{-0.007455in}{-0.004330in}}{\pgfqpoint{-0.005893in}{-0.005893in}}%
\pgfpathcurveto{\pgfqpoint{-0.004330in}{-0.007455in}}{\pgfqpoint{-0.002210in}{-0.008333in}}{\pgfqpoint{0.000000in}{-0.008333in}}%
\pgfpathlineto{\pgfqpoint{0.000000in}{-0.008333in}}%
\pgfpathclose%
\pgfusepath{stroke,fill}%
}%
\begin{pgfscope}%
\pgfsys@transformshift{1.781876in}{4.040833in}%
\pgfsys@useobject{currentmarker}{}%
\end{pgfscope}%
\end{pgfscope}%
\begin{pgfscope}%
\pgfpathrectangle{\pgfqpoint{0.100000in}{2.493063in}}{\pgfqpoint{5.037500in}{3.427208in}}%
\pgfusepath{clip}%
\pgfsetrectcap%
\pgfsetroundjoin%
\pgfsetlinewidth{1.505625pt}%
\definecolor{currentstroke}{rgb}{0.501961,0.501961,0.501961}%
\pgfsetstrokecolor{currentstroke}%
\pgfsetstrokeopacity{0.500000}%
\pgfsetdash{}{0pt}%
\pgfpathmoveto{\pgfqpoint{4.590541in}{4.991408in}}%
\pgfusepath{stroke}%
\end{pgfscope}%
\begin{pgfscope}%
\pgfpathrectangle{\pgfqpoint{0.100000in}{2.493063in}}{\pgfqpoint{5.037500in}{3.427208in}}%
\pgfusepath{clip}%
\pgfsetbuttcap%
\pgfsetroundjoin%
\definecolor{currentfill}{rgb}{0.501961,0.501961,0.501961}%
\pgfsetfillcolor{currentfill}%
\pgfsetfillopacity{0.500000}%
\pgfsetlinewidth{0.250937pt}%
\definecolor{currentstroke}{rgb}{0.000000,0.000000,0.000000}%
\pgfsetstrokecolor{currentstroke}%
\pgfsetstrokeopacity{0.500000}%
\pgfsetdash{}{0pt}%
\pgfsys@defobject{currentmarker}{\pgfqpoint{-0.013889in}{-0.013889in}}{\pgfqpoint{0.013889in}{0.013889in}}{%
\pgfpathmoveto{\pgfqpoint{0.000000in}{-0.013889in}}%
\pgfpathcurveto{\pgfqpoint{0.003683in}{-0.013889in}}{\pgfqpoint{0.007216in}{-0.012425in}}{\pgfqpoint{0.009821in}{-0.009821in}}%
\pgfpathcurveto{\pgfqpoint{0.012425in}{-0.007216in}}{\pgfqpoint{0.013889in}{-0.003683in}}{\pgfqpoint{0.013889in}{0.000000in}}%
\pgfpathcurveto{\pgfqpoint{0.013889in}{0.003683in}}{\pgfqpoint{0.012425in}{0.007216in}}{\pgfqpoint{0.009821in}{0.009821in}}%
\pgfpathcurveto{\pgfqpoint{0.007216in}{0.012425in}}{\pgfqpoint{0.003683in}{0.013889in}}{\pgfqpoint{0.000000in}{0.013889in}}%
\pgfpathcurveto{\pgfqpoint{-0.003683in}{0.013889in}}{\pgfqpoint{-0.007216in}{0.012425in}}{\pgfqpoint{-0.009821in}{0.009821in}}%
\pgfpathcurveto{\pgfqpoint{-0.012425in}{0.007216in}}{\pgfqpoint{-0.013889in}{0.003683in}}{\pgfqpoint{-0.013889in}{0.000000in}}%
\pgfpathcurveto{\pgfqpoint{-0.013889in}{-0.003683in}}{\pgfqpoint{-0.012425in}{-0.007216in}}{\pgfqpoint{-0.009821in}{-0.009821in}}%
\pgfpathcurveto{\pgfqpoint{-0.007216in}{-0.012425in}}{\pgfqpoint{-0.003683in}{-0.013889in}}{\pgfqpoint{0.000000in}{-0.013889in}}%
\pgfpathlineto{\pgfqpoint{0.000000in}{-0.013889in}}%
\pgfpathclose%
\pgfusepath{stroke,fill}%
}%
\begin{pgfscope}%
\pgfsys@transformshift{4.590541in}{4.991408in}%
\pgfsys@useobject{currentmarker}{}%
\end{pgfscope}%
\end{pgfscope}%
\begin{pgfscope}%
\pgfpathrectangle{\pgfqpoint{0.100000in}{2.493063in}}{\pgfqpoint{5.037500in}{3.427208in}}%
\pgfusepath{clip}%
\pgfsetrectcap%
\pgfsetroundjoin%
\pgfsetlinewidth{1.505625pt}%
\definecolor{currentstroke}{rgb}{0.678431,1.000000,0.184314}%
\pgfsetstrokecolor{currentstroke}%
\pgfsetstrokeopacity{0.500000}%
\pgfsetdash{}{0pt}%
\pgfpathmoveto{\pgfqpoint{4.425126in}{4.888565in}}%
\pgfusepath{stroke}%
\end{pgfscope}%
\begin{pgfscope}%
\pgfpathrectangle{\pgfqpoint{0.100000in}{2.493063in}}{\pgfqpoint{5.037500in}{3.427208in}}%
\pgfusepath{clip}%
\pgfsetbuttcap%
\pgfsetroundjoin%
\definecolor{currentfill}{rgb}{0.678431,1.000000,0.184314}%
\pgfsetfillcolor{currentfill}%
\pgfsetfillopacity{0.500000}%
\pgfsetlinewidth{0.250937pt}%
\definecolor{currentstroke}{rgb}{0.000000,0.000000,0.000000}%
\pgfsetstrokecolor{currentstroke}%
\pgfsetstrokeopacity{0.500000}%
\pgfsetdash{}{0pt}%
\pgfsys@defobject{currentmarker}{\pgfqpoint{-0.005556in}{-0.005556in}}{\pgfqpoint{0.005556in}{0.005556in}}{%
\pgfpathmoveto{\pgfqpoint{0.000000in}{-0.005556in}}%
\pgfpathcurveto{\pgfqpoint{0.001473in}{-0.005556in}}{\pgfqpoint{0.002887in}{-0.004970in}}{\pgfqpoint{0.003928in}{-0.003928in}}%
\pgfpathcurveto{\pgfqpoint{0.004970in}{-0.002887in}}{\pgfqpoint{0.005556in}{-0.001473in}}{\pgfqpoint{0.005556in}{0.000000in}}%
\pgfpathcurveto{\pgfqpoint{0.005556in}{0.001473in}}{\pgfqpoint{0.004970in}{0.002887in}}{\pgfqpoint{0.003928in}{0.003928in}}%
\pgfpathcurveto{\pgfqpoint{0.002887in}{0.004970in}}{\pgfqpoint{0.001473in}{0.005556in}}{\pgfqpoint{0.000000in}{0.005556in}}%
\pgfpathcurveto{\pgfqpoint{-0.001473in}{0.005556in}}{\pgfqpoint{-0.002887in}{0.004970in}}{\pgfqpoint{-0.003928in}{0.003928in}}%
\pgfpathcurveto{\pgfqpoint{-0.004970in}{0.002887in}}{\pgfqpoint{-0.005556in}{0.001473in}}{\pgfqpoint{-0.005556in}{0.000000in}}%
\pgfpathcurveto{\pgfqpoint{-0.005556in}{-0.001473in}}{\pgfqpoint{-0.004970in}{-0.002887in}}{\pgfqpoint{-0.003928in}{-0.003928in}}%
\pgfpathcurveto{\pgfqpoint{-0.002887in}{-0.004970in}}{\pgfqpoint{-0.001473in}{-0.005556in}}{\pgfqpoint{0.000000in}{-0.005556in}}%
\pgfpathlineto{\pgfqpoint{0.000000in}{-0.005556in}}%
\pgfpathclose%
\pgfusepath{stroke,fill}%
}%
\begin{pgfscope}%
\pgfsys@transformshift{4.425126in}{4.888565in}%
\pgfsys@useobject{currentmarker}{}%
\end{pgfscope}%
\end{pgfscope}%
\begin{pgfscope}%
\pgfpathrectangle{\pgfqpoint{0.100000in}{2.493063in}}{\pgfqpoint{5.037500in}{3.427208in}}%
\pgfusepath{clip}%
\pgfsetrectcap%
\pgfsetroundjoin%
\pgfsetlinewidth{1.505625pt}%
\definecolor{currentstroke}{rgb}{0.501961,0.501961,0.501961}%
\pgfsetstrokecolor{currentstroke}%
\pgfsetstrokeopacity{0.500000}%
\pgfsetdash{}{0pt}%
\pgfpathmoveto{\pgfqpoint{4.160198in}{4.929420in}}%
\pgfusepath{stroke}%
\end{pgfscope}%
\begin{pgfscope}%
\pgfpathrectangle{\pgfqpoint{0.100000in}{2.493063in}}{\pgfqpoint{5.037500in}{3.427208in}}%
\pgfusepath{clip}%
\pgfsetbuttcap%
\pgfsetroundjoin%
\definecolor{currentfill}{rgb}{0.501961,0.501961,0.501961}%
\pgfsetfillcolor{currentfill}%
\pgfsetfillopacity{0.500000}%
\pgfsetlinewidth{0.250937pt}%
\definecolor{currentstroke}{rgb}{0.000000,0.000000,0.000000}%
\pgfsetstrokecolor{currentstroke}%
\pgfsetstrokeopacity{0.500000}%
\pgfsetdash{}{0pt}%
\pgfsys@defobject{currentmarker}{\pgfqpoint{-0.013889in}{-0.013889in}}{\pgfqpoint{0.013889in}{0.013889in}}{%
\pgfpathmoveto{\pgfqpoint{0.000000in}{-0.013889in}}%
\pgfpathcurveto{\pgfqpoint{0.003683in}{-0.013889in}}{\pgfqpoint{0.007216in}{-0.012425in}}{\pgfqpoint{0.009821in}{-0.009821in}}%
\pgfpathcurveto{\pgfqpoint{0.012425in}{-0.007216in}}{\pgfqpoint{0.013889in}{-0.003683in}}{\pgfqpoint{0.013889in}{0.000000in}}%
\pgfpathcurveto{\pgfqpoint{0.013889in}{0.003683in}}{\pgfqpoint{0.012425in}{0.007216in}}{\pgfqpoint{0.009821in}{0.009821in}}%
\pgfpathcurveto{\pgfqpoint{0.007216in}{0.012425in}}{\pgfqpoint{0.003683in}{0.013889in}}{\pgfqpoint{0.000000in}{0.013889in}}%
\pgfpathcurveto{\pgfqpoint{-0.003683in}{0.013889in}}{\pgfqpoint{-0.007216in}{0.012425in}}{\pgfqpoint{-0.009821in}{0.009821in}}%
\pgfpathcurveto{\pgfqpoint{-0.012425in}{0.007216in}}{\pgfqpoint{-0.013889in}{0.003683in}}{\pgfqpoint{-0.013889in}{0.000000in}}%
\pgfpathcurveto{\pgfqpoint{-0.013889in}{-0.003683in}}{\pgfqpoint{-0.012425in}{-0.007216in}}{\pgfqpoint{-0.009821in}{-0.009821in}}%
\pgfpathcurveto{\pgfqpoint{-0.007216in}{-0.012425in}}{\pgfqpoint{-0.003683in}{-0.013889in}}{\pgfqpoint{0.000000in}{-0.013889in}}%
\pgfpathlineto{\pgfqpoint{0.000000in}{-0.013889in}}%
\pgfpathclose%
\pgfusepath{stroke,fill}%
}%
\begin{pgfscope}%
\pgfsys@transformshift{4.160198in}{4.929420in}%
\pgfsys@useobject{currentmarker}{}%
\end{pgfscope}%
\end{pgfscope}%
\begin{pgfscope}%
\pgfpathrectangle{\pgfqpoint{0.100000in}{2.493063in}}{\pgfqpoint{5.037500in}{3.427208in}}%
\pgfusepath{clip}%
\pgfsetrectcap%
\pgfsetroundjoin%
\pgfsetlinewidth{1.505625pt}%
\definecolor{currentstroke}{rgb}{0.000000,0.000000,1.000000}%
\pgfsetstrokecolor{currentstroke}%
\pgfsetstrokeopacity{0.500000}%
\pgfsetdash{}{0pt}%
\pgfpathmoveto{\pgfqpoint{4.350454in}{4.872186in}}%
\pgfusepath{stroke}%
\end{pgfscope}%
\begin{pgfscope}%
\pgfpathrectangle{\pgfqpoint{0.100000in}{2.493063in}}{\pgfqpoint{5.037500in}{3.427208in}}%
\pgfusepath{clip}%
\pgfsetbuttcap%
\pgfsetroundjoin%
\definecolor{currentfill}{rgb}{0.000000,0.000000,1.000000}%
\pgfsetfillcolor{currentfill}%
\pgfsetfillopacity{0.500000}%
\pgfsetlinewidth{0.250937pt}%
\definecolor{currentstroke}{rgb}{0.000000,0.000000,0.000000}%
\pgfsetstrokecolor{currentstroke}%
\pgfsetstrokeopacity{0.500000}%
\pgfsetdash{}{0pt}%
\pgfsys@defobject{currentmarker}{\pgfqpoint{-0.011111in}{-0.011111in}}{\pgfqpoint{0.011111in}{0.011111in}}{%
\pgfpathmoveto{\pgfqpoint{0.000000in}{-0.011111in}}%
\pgfpathcurveto{\pgfqpoint{0.002947in}{-0.011111in}}{\pgfqpoint{0.005773in}{-0.009940in}}{\pgfqpoint{0.007857in}{-0.007857in}}%
\pgfpathcurveto{\pgfqpoint{0.009940in}{-0.005773in}}{\pgfqpoint{0.011111in}{-0.002947in}}{\pgfqpoint{0.011111in}{0.000000in}}%
\pgfpathcurveto{\pgfqpoint{0.011111in}{0.002947in}}{\pgfqpoint{0.009940in}{0.005773in}}{\pgfqpoint{0.007857in}{0.007857in}}%
\pgfpathcurveto{\pgfqpoint{0.005773in}{0.009940in}}{\pgfqpoint{0.002947in}{0.011111in}}{\pgfqpoint{0.000000in}{0.011111in}}%
\pgfpathcurveto{\pgfqpoint{-0.002947in}{0.011111in}}{\pgfqpoint{-0.005773in}{0.009940in}}{\pgfqpoint{-0.007857in}{0.007857in}}%
\pgfpathcurveto{\pgfqpoint{-0.009940in}{0.005773in}}{\pgfqpoint{-0.011111in}{0.002947in}}{\pgfqpoint{-0.011111in}{0.000000in}}%
\pgfpathcurveto{\pgfqpoint{-0.011111in}{-0.002947in}}{\pgfqpoint{-0.009940in}{-0.005773in}}{\pgfqpoint{-0.007857in}{-0.007857in}}%
\pgfpathcurveto{\pgfqpoint{-0.005773in}{-0.009940in}}{\pgfqpoint{-0.002947in}{-0.011111in}}{\pgfqpoint{0.000000in}{-0.011111in}}%
\pgfpathlineto{\pgfqpoint{0.000000in}{-0.011111in}}%
\pgfpathclose%
\pgfusepath{stroke,fill}%
}%
\begin{pgfscope}%
\pgfsys@transformshift{4.350454in}{4.872186in}%
\pgfsys@useobject{currentmarker}{}%
\end{pgfscope}%
\end{pgfscope}%
\begin{pgfscope}%
\pgfpathrectangle{\pgfqpoint{0.100000in}{2.493063in}}{\pgfqpoint{5.037500in}{3.427208in}}%
\pgfusepath{clip}%
\pgfsetrectcap%
\pgfsetroundjoin%
\pgfsetlinewidth{1.505625pt}%
\definecolor{currentstroke}{rgb}{0.678431,1.000000,0.184314}%
\pgfsetstrokecolor{currentstroke}%
\pgfsetstrokeopacity{0.500000}%
\pgfsetdash{}{0pt}%
\pgfpathmoveto{\pgfqpoint{4.581959in}{5.067532in}}%
\pgfusepath{stroke}%
\end{pgfscope}%
\begin{pgfscope}%
\pgfpathrectangle{\pgfqpoint{0.100000in}{2.493063in}}{\pgfqpoint{5.037500in}{3.427208in}}%
\pgfusepath{clip}%
\pgfsetbuttcap%
\pgfsetroundjoin%
\definecolor{currentfill}{rgb}{0.678431,1.000000,0.184314}%
\pgfsetfillcolor{currentfill}%
\pgfsetfillopacity{0.500000}%
\pgfsetlinewidth{0.250937pt}%
\definecolor{currentstroke}{rgb}{0.000000,0.000000,0.000000}%
\pgfsetstrokecolor{currentstroke}%
\pgfsetstrokeopacity{0.500000}%
\pgfsetdash{}{0pt}%
\pgfsys@defobject{currentmarker}{\pgfqpoint{-0.011111in}{-0.011111in}}{\pgfqpoint{0.011111in}{0.011111in}}{%
\pgfpathmoveto{\pgfqpoint{0.000000in}{-0.011111in}}%
\pgfpathcurveto{\pgfqpoint{0.002947in}{-0.011111in}}{\pgfqpoint{0.005773in}{-0.009940in}}{\pgfqpoint{0.007857in}{-0.007857in}}%
\pgfpathcurveto{\pgfqpoint{0.009940in}{-0.005773in}}{\pgfqpoint{0.011111in}{-0.002947in}}{\pgfqpoint{0.011111in}{0.000000in}}%
\pgfpathcurveto{\pgfqpoint{0.011111in}{0.002947in}}{\pgfqpoint{0.009940in}{0.005773in}}{\pgfqpoint{0.007857in}{0.007857in}}%
\pgfpathcurveto{\pgfqpoint{0.005773in}{0.009940in}}{\pgfqpoint{0.002947in}{0.011111in}}{\pgfqpoint{0.000000in}{0.011111in}}%
\pgfpathcurveto{\pgfqpoint{-0.002947in}{0.011111in}}{\pgfqpoint{-0.005773in}{0.009940in}}{\pgfqpoint{-0.007857in}{0.007857in}}%
\pgfpathcurveto{\pgfqpoint{-0.009940in}{0.005773in}}{\pgfqpoint{-0.011111in}{0.002947in}}{\pgfqpoint{-0.011111in}{0.000000in}}%
\pgfpathcurveto{\pgfqpoint{-0.011111in}{-0.002947in}}{\pgfqpoint{-0.009940in}{-0.005773in}}{\pgfqpoint{-0.007857in}{-0.007857in}}%
\pgfpathcurveto{\pgfqpoint{-0.005773in}{-0.009940in}}{\pgfqpoint{-0.002947in}{-0.011111in}}{\pgfqpoint{0.000000in}{-0.011111in}}%
\pgfpathlineto{\pgfqpoint{0.000000in}{-0.011111in}}%
\pgfpathclose%
\pgfusepath{stroke,fill}%
}%
\begin{pgfscope}%
\pgfsys@transformshift{4.581959in}{5.067532in}%
\pgfsys@useobject{currentmarker}{}%
\end{pgfscope}%
\end{pgfscope}%
\begin{pgfscope}%
\pgfpathrectangle{\pgfqpoint{0.100000in}{2.493063in}}{\pgfqpoint{5.037500in}{3.427208in}}%
\pgfusepath{clip}%
\pgfsetrectcap%
\pgfsetroundjoin%
\pgfsetlinewidth{1.505625pt}%
\definecolor{currentstroke}{rgb}{0.678431,1.000000,0.184314}%
\pgfsetstrokecolor{currentstroke}%
\pgfsetstrokeopacity{0.500000}%
\pgfsetdash{}{0pt}%
\pgfpathmoveto{\pgfqpoint{4.368363in}{4.917031in}}%
\pgfusepath{stroke}%
\end{pgfscope}%
\begin{pgfscope}%
\pgfpathrectangle{\pgfqpoint{0.100000in}{2.493063in}}{\pgfqpoint{5.037500in}{3.427208in}}%
\pgfusepath{clip}%
\pgfsetbuttcap%
\pgfsetroundjoin%
\definecolor{currentfill}{rgb}{0.678431,1.000000,0.184314}%
\pgfsetfillcolor{currentfill}%
\pgfsetfillopacity{0.500000}%
\pgfsetlinewidth{0.250937pt}%
\definecolor{currentstroke}{rgb}{0.000000,0.000000,0.000000}%
\pgfsetstrokecolor{currentstroke}%
\pgfsetstrokeopacity{0.500000}%
\pgfsetdash{}{0pt}%
\pgfsys@defobject{currentmarker}{\pgfqpoint{-0.008333in}{-0.008333in}}{\pgfqpoint{0.008333in}{0.008333in}}{%
\pgfpathmoveto{\pgfqpoint{0.000000in}{-0.008333in}}%
\pgfpathcurveto{\pgfqpoint{0.002210in}{-0.008333in}}{\pgfqpoint{0.004330in}{-0.007455in}}{\pgfqpoint{0.005893in}{-0.005893in}}%
\pgfpathcurveto{\pgfqpoint{0.007455in}{-0.004330in}}{\pgfqpoint{0.008333in}{-0.002210in}}{\pgfqpoint{0.008333in}{0.000000in}}%
\pgfpathcurveto{\pgfqpoint{0.008333in}{0.002210in}}{\pgfqpoint{0.007455in}{0.004330in}}{\pgfqpoint{0.005893in}{0.005893in}}%
\pgfpathcurveto{\pgfqpoint{0.004330in}{0.007455in}}{\pgfqpoint{0.002210in}{0.008333in}}{\pgfqpoint{0.000000in}{0.008333in}}%
\pgfpathcurveto{\pgfqpoint{-0.002210in}{0.008333in}}{\pgfqpoint{-0.004330in}{0.007455in}}{\pgfqpoint{-0.005893in}{0.005893in}}%
\pgfpathcurveto{\pgfqpoint{-0.007455in}{0.004330in}}{\pgfqpoint{-0.008333in}{0.002210in}}{\pgfqpoint{-0.008333in}{0.000000in}}%
\pgfpathcurveto{\pgfqpoint{-0.008333in}{-0.002210in}}{\pgfqpoint{-0.007455in}{-0.004330in}}{\pgfqpoint{-0.005893in}{-0.005893in}}%
\pgfpathcurveto{\pgfqpoint{-0.004330in}{-0.007455in}}{\pgfqpoint{-0.002210in}{-0.008333in}}{\pgfqpoint{0.000000in}{-0.008333in}}%
\pgfpathlineto{\pgfqpoint{0.000000in}{-0.008333in}}%
\pgfpathclose%
\pgfusepath{stroke,fill}%
}%
\begin{pgfscope}%
\pgfsys@transformshift{4.368363in}{4.917031in}%
\pgfsys@useobject{currentmarker}{}%
\end{pgfscope}%
\end{pgfscope}%
\begin{pgfscope}%
\pgfpathrectangle{\pgfqpoint{0.100000in}{2.493063in}}{\pgfqpoint{5.037500in}{3.427208in}}%
\pgfusepath{clip}%
\pgfsetrectcap%
\pgfsetroundjoin%
\pgfsetlinewidth{1.505625pt}%
\definecolor{currentstroke}{rgb}{0.000000,0.000000,1.000000}%
\pgfsetstrokecolor{currentstroke}%
\pgfsetstrokeopacity{0.500000}%
\pgfsetdash{}{0pt}%
\pgfpathmoveto{\pgfqpoint{4.589191in}{4.905473in}}%
\pgfusepath{stroke}%
\end{pgfscope}%
\begin{pgfscope}%
\pgfpathrectangle{\pgfqpoint{0.100000in}{2.493063in}}{\pgfqpoint{5.037500in}{3.427208in}}%
\pgfusepath{clip}%
\pgfsetbuttcap%
\pgfsetroundjoin%
\definecolor{currentfill}{rgb}{0.000000,0.000000,1.000000}%
\pgfsetfillcolor{currentfill}%
\pgfsetfillopacity{0.500000}%
\pgfsetlinewidth{0.250937pt}%
\definecolor{currentstroke}{rgb}{0.000000,0.000000,0.000000}%
\pgfsetstrokecolor{currentstroke}%
\pgfsetstrokeopacity{0.500000}%
\pgfsetdash{}{0pt}%
\pgfsys@defobject{currentmarker}{\pgfqpoint{-0.005556in}{-0.005556in}}{\pgfqpoint{0.005556in}{0.005556in}}{%
\pgfpathmoveto{\pgfqpoint{0.000000in}{-0.005556in}}%
\pgfpathcurveto{\pgfqpoint{0.001473in}{-0.005556in}}{\pgfqpoint{0.002887in}{-0.004970in}}{\pgfqpoint{0.003928in}{-0.003928in}}%
\pgfpathcurveto{\pgfqpoint{0.004970in}{-0.002887in}}{\pgfqpoint{0.005556in}{-0.001473in}}{\pgfqpoint{0.005556in}{0.000000in}}%
\pgfpathcurveto{\pgfqpoint{0.005556in}{0.001473in}}{\pgfqpoint{0.004970in}{0.002887in}}{\pgfqpoint{0.003928in}{0.003928in}}%
\pgfpathcurveto{\pgfqpoint{0.002887in}{0.004970in}}{\pgfqpoint{0.001473in}{0.005556in}}{\pgfqpoint{0.000000in}{0.005556in}}%
\pgfpathcurveto{\pgfqpoint{-0.001473in}{0.005556in}}{\pgfqpoint{-0.002887in}{0.004970in}}{\pgfqpoint{-0.003928in}{0.003928in}}%
\pgfpathcurveto{\pgfqpoint{-0.004970in}{0.002887in}}{\pgfqpoint{-0.005556in}{0.001473in}}{\pgfqpoint{-0.005556in}{0.000000in}}%
\pgfpathcurveto{\pgfqpoint{-0.005556in}{-0.001473in}}{\pgfqpoint{-0.004970in}{-0.002887in}}{\pgfqpoint{-0.003928in}{-0.003928in}}%
\pgfpathcurveto{\pgfqpoint{-0.002887in}{-0.004970in}}{\pgfqpoint{-0.001473in}{-0.005556in}}{\pgfqpoint{0.000000in}{-0.005556in}}%
\pgfpathlineto{\pgfqpoint{0.000000in}{-0.005556in}}%
\pgfpathclose%
\pgfusepath{stroke,fill}%
}%
\begin{pgfscope}%
\pgfsys@transformshift{4.589191in}{4.905473in}%
\pgfsys@useobject{currentmarker}{}%
\end{pgfscope}%
\end{pgfscope}%
\begin{pgfscope}%
\pgfpathrectangle{\pgfqpoint{0.100000in}{2.493063in}}{\pgfqpoint{5.037500in}{3.427208in}}%
\pgfusepath{clip}%
\pgfsetrectcap%
\pgfsetroundjoin%
\pgfsetlinewidth{1.505625pt}%
\definecolor{currentstroke}{rgb}{0.000000,0.000000,1.000000}%
\pgfsetstrokecolor{currentstroke}%
\pgfsetstrokeopacity{0.500000}%
\pgfsetdash{}{0pt}%
\pgfpathmoveto{\pgfqpoint{4.622153in}{4.771381in}}%
\pgfusepath{stroke}%
\end{pgfscope}%
\begin{pgfscope}%
\pgfpathrectangle{\pgfqpoint{0.100000in}{2.493063in}}{\pgfqpoint{5.037500in}{3.427208in}}%
\pgfusepath{clip}%
\pgfsetbuttcap%
\pgfsetroundjoin%
\definecolor{currentfill}{rgb}{0.000000,0.000000,1.000000}%
\pgfsetfillcolor{currentfill}%
\pgfsetfillopacity{0.500000}%
\pgfsetlinewidth{0.250937pt}%
\definecolor{currentstroke}{rgb}{0.000000,0.000000,0.000000}%
\pgfsetstrokecolor{currentstroke}%
\pgfsetstrokeopacity{0.500000}%
\pgfsetdash{}{0pt}%
\pgfsys@defobject{currentmarker}{\pgfqpoint{-0.022222in}{-0.022222in}}{\pgfqpoint{0.022222in}{0.022222in}}{%
\pgfpathmoveto{\pgfqpoint{0.000000in}{-0.022222in}}%
\pgfpathcurveto{\pgfqpoint{0.005893in}{-0.022222in}}{\pgfqpoint{0.011546in}{-0.019881in}}{\pgfqpoint{0.015713in}{-0.015713in}}%
\pgfpathcurveto{\pgfqpoint{0.019881in}{-0.011546in}}{\pgfqpoint{0.022222in}{-0.005893in}}{\pgfqpoint{0.022222in}{0.000000in}}%
\pgfpathcurveto{\pgfqpoint{0.022222in}{0.005893in}}{\pgfqpoint{0.019881in}{0.011546in}}{\pgfqpoint{0.015713in}{0.015713in}}%
\pgfpathcurveto{\pgfqpoint{0.011546in}{0.019881in}}{\pgfqpoint{0.005893in}{0.022222in}}{\pgfqpoint{0.000000in}{0.022222in}}%
\pgfpathcurveto{\pgfqpoint{-0.005893in}{0.022222in}}{\pgfqpoint{-0.011546in}{0.019881in}}{\pgfqpoint{-0.015713in}{0.015713in}}%
\pgfpathcurveto{\pgfqpoint{-0.019881in}{0.011546in}}{\pgfqpoint{-0.022222in}{0.005893in}}{\pgfqpoint{-0.022222in}{0.000000in}}%
\pgfpathcurveto{\pgfqpoint{-0.022222in}{-0.005893in}}{\pgfqpoint{-0.019881in}{-0.011546in}}{\pgfqpoint{-0.015713in}{-0.015713in}}%
\pgfpathcurveto{\pgfqpoint{-0.011546in}{-0.019881in}}{\pgfqpoint{-0.005893in}{-0.022222in}}{\pgfqpoint{0.000000in}{-0.022222in}}%
\pgfpathlineto{\pgfqpoint{0.000000in}{-0.022222in}}%
\pgfpathclose%
\pgfusepath{stroke,fill}%
}%
\begin{pgfscope}%
\pgfsys@transformshift{4.622153in}{4.771381in}%
\pgfsys@useobject{currentmarker}{}%
\end{pgfscope}%
\end{pgfscope}%
\begin{pgfscope}%
\pgfpathrectangle{\pgfqpoint{0.100000in}{2.493063in}}{\pgfqpoint{5.037500in}{3.427208in}}%
\pgfusepath{clip}%
\pgfsetrectcap%
\pgfsetroundjoin%
\pgfsetlinewidth{1.505625pt}%
\definecolor{currentstroke}{rgb}{0.501961,0.501961,0.501961}%
\pgfsetstrokecolor{currentstroke}%
\pgfsetstrokeopacity{0.500000}%
\pgfsetdash{}{0pt}%
\pgfpathmoveto{\pgfqpoint{4.259261in}{4.979660in}}%
\pgfusepath{stroke}%
\end{pgfscope}%
\begin{pgfscope}%
\pgfpathrectangle{\pgfqpoint{0.100000in}{2.493063in}}{\pgfqpoint{5.037500in}{3.427208in}}%
\pgfusepath{clip}%
\pgfsetbuttcap%
\pgfsetroundjoin%
\definecolor{currentfill}{rgb}{0.501961,0.501961,0.501961}%
\pgfsetfillcolor{currentfill}%
\pgfsetfillopacity{0.500000}%
\pgfsetlinewidth{0.250937pt}%
\definecolor{currentstroke}{rgb}{0.000000,0.000000,0.000000}%
\pgfsetstrokecolor{currentstroke}%
\pgfsetstrokeopacity{0.500000}%
\pgfsetdash{}{0pt}%
\pgfsys@defobject{currentmarker}{\pgfqpoint{-0.013889in}{-0.013889in}}{\pgfqpoint{0.013889in}{0.013889in}}{%
\pgfpathmoveto{\pgfqpoint{0.000000in}{-0.013889in}}%
\pgfpathcurveto{\pgfqpoint{0.003683in}{-0.013889in}}{\pgfqpoint{0.007216in}{-0.012425in}}{\pgfqpoint{0.009821in}{-0.009821in}}%
\pgfpathcurveto{\pgfqpoint{0.012425in}{-0.007216in}}{\pgfqpoint{0.013889in}{-0.003683in}}{\pgfqpoint{0.013889in}{0.000000in}}%
\pgfpathcurveto{\pgfqpoint{0.013889in}{0.003683in}}{\pgfqpoint{0.012425in}{0.007216in}}{\pgfqpoint{0.009821in}{0.009821in}}%
\pgfpathcurveto{\pgfqpoint{0.007216in}{0.012425in}}{\pgfqpoint{0.003683in}{0.013889in}}{\pgfqpoint{0.000000in}{0.013889in}}%
\pgfpathcurveto{\pgfqpoint{-0.003683in}{0.013889in}}{\pgfqpoint{-0.007216in}{0.012425in}}{\pgfqpoint{-0.009821in}{0.009821in}}%
\pgfpathcurveto{\pgfqpoint{-0.012425in}{0.007216in}}{\pgfqpoint{-0.013889in}{0.003683in}}{\pgfqpoint{-0.013889in}{0.000000in}}%
\pgfpathcurveto{\pgfqpoint{-0.013889in}{-0.003683in}}{\pgfqpoint{-0.012425in}{-0.007216in}}{\pgfqpoint{-0.009821in}{-0.009821in}}%
\pgfpathcurveto{\pgfqpoint{-0.007216in}{-0.012425in}}{\pgfqpoint{-0.003683in}{-0.013889in}}{\pgfqpoint{0.000000in}{-0.013889in}}%
\pgfpathlineto{\pgfqpoint{0.000000in}{-0.013889in}}%
\pgfpathclose%
\pgfusepath{stroke,fill}%
}%
\begin{pgfscope}%
\pgfsys@transformshift{4.259261in}{4.979660in}%
\pgfsys@useobject{currentmarker}{}%
\end{pgfscope}%
\end{pgfscope}%
\begin{pgfscope}%
\pgfpathrectangle{\pgfqpoint{0.100000in}{2.493063in}}{\pgfqpoint{5.037500in}{3.427208in}}%
\pgfusepath{clip}%
\pgfsetrectcap%
\pgfsetroundjoin%
\pgfsetlinewidth{1.505625pt}%
\definecolor{currentstroke}{rgb}{0.501961,0.501961,0.501961}%
\pgfsetstrokecolor{currentstroke}%
\pgfsetstrokeopacity{0.500000}%
\pgfsetdash{}{0pt}%
\pgfpathmoveto{\pgfqpoint{4.383016in}{4.991810in}}%
\pgfusepath{stroke}%
\end{pgfscope}%
\begin{pgfscope}%
\pgfpathrectangle{\pgfqpoint{0.100000in}{2.493063in}}{\pgfqpoint{5.037500in}{3.427208in}}%
\pgfusepath{clip}%
\pgfsetbuttcap%
\pgfsetroundjoin%
\definecolor{currentfill}{rgb}{0.501961,0.501961,0.501961}%
\pgfsetfillcolor{currentfill}%
\pgfsetfillopacity{0.500000}%
\pgfsetlinewidth{0.250937pt}%
\definecolor{currentstroke}{rgb}{0.000000,0.000000,0.000000}%
\pgfsetstrokecolor{currentstroke}%
\pgfsetstrokeopacity{0.500000}%
\pgfsetdash{}{0pt}%
\pgfsys@defobject{currentmarker}{\pgfqpoint{-0.013889in}{-0.013889in}}{\pgfqpoint{0.013889in}{0.013889in}}{%
\pgfpathmoveto{\pgfqpoint{0.000000in}{-0.013889in}}%
\pgfpathcurveto{\pgfqpoint{0.003683in}{-0.013889in}}{\pgfqpoint{0.007216in}{-0.012425in}}{\pgfqpoint{0.009821in}{-0.009821in}}%
\pgfpathcurveto{\pgfqpoint{0.012425in}{-0.007216in}}{\pgfqpoint{0.013889in}{-0.003683in}}{\pgfqpoint{0.013889in}{0.000000in}}%
\pgfpathcurveto{\pgfqpoint{0.013889in}{0.003683in}}{\pgfqpoint{0.012425in}{0.007216in}}{\pgfqpoint{0.009821in}{0.009821in}}%
\pgfpathcurveto{\pgfqpoint{0.007216in}{0.012425in}}{\pgfqpoint{0.003683in}{0.013889in}}{\pgfqpoint{0.000000in}{0.013889in}}%
\pgfpathcurveto{\pgfqpoint{-0.003683in}{0.013889in}}{\pgfqpoint{-0.007216in}{0.012425in}}{\pgfqpoint{-0.009821in}{0.009821in}}%
\pgfpathcurveto{\pgfqpoint{-0.012425in}{0.007216in}}{\pgfqpoint{-0.013889in}{0.003683in}}{\pgfqpoint{-0.013889in}{0.000000in}}%
\pgfpathcurveto{\pgfqpoint{-0.013889in}{-0.003683in}}{\pgfqpoint{-0.012425in}{-0.007216in}}{\pgfqpoint{-0.009821in}{-0.009821in}}%
\pgfpathcurveto{\pgfqpoint{-0.007216in}{-0.012425in}}{\pgfqpoint{-0.003683in}{-0.013889in}}{\pgfqpoint{0.000000in}{-0.013889in}}%
\pgfpathlineto{\pgfqpoint{0.000000in}{-0.013889in}}%
\pgfpathclose%
\pgfusepath{stroke,fill}%
}%
\begin{pgfscope}%
\pgfsys@transformshift{4.383016in}{4.991810in}%
\pgfsys@useobject{currentmarker}{}%
\end{pgfscope}%
\end{pgfscope}%
\begin{pgfscope}%
\pgfpathrectangle{\pgfqpoint{0.100000in}{2.493063in}}{\pgfqpoint{5.037500in}{3.427208in}}%
\pgfusepath{clip}%
\pgfsetrectcap%
\pgfsetroundjoin%
\pgfsetlinewidth{1.505625pt}%
\definecolor{currentstroke}{rgb}{0.501961,0.501961,0.501961}%
\pgfsetstrokecolor{currentstroke}%
\pgfsetstrokeopacity{0.500000}%
\pgfsetdash{}{0pt}%
\pgfpathmoveto{\pgfqpoint{4.457128in}{5.014032in}}%
\pgfusepath{stroke}%
\end{pgfscope}%
\begin{pgfscope}%
\pgfpathrectangle{\pgfqpoint{0.100000in}{2.493063in}}{\pgfqpoint{5.037500in}{3.427208in}}%
\pgfusepath{clip}%
\pgfsetbuttcap%
\pgfsetroundjoin%
\definecolor{currentfill}{rgb}{0.501961,0.501961,0.501961}%
\pgfsetfillcolor{currentfill}%
\pgfsetfillopacity{0.500000}%
\pgfsetlinewidth{0.250937pt}%
\definecolor{currentstroke}{rgb}{0.000000,0.000000,0.000000}%
\pgfsetstrokecolor{currentstroke}%
\pgfsetstrokeopacity{0.500000}%
\pgfsetdash{}{0pt}%
\pgfsys@defobject{currentmarker}{\pgfqpoint{-0.013889in}{-0.013889in}}{\pgfqpoint{0.013889in}{0.013889in}}{%
\pgfpathmoveto{\pgfqpoint{0.000000in}{-0.013889in}}%
\pgfpathcurveto{\pgfqpoint{0.003683in}{-0.013889in}}{\pgfqpoint{0.007216in}{-0.012425in}}{\pgfqpoint{0.009821in}{-0.009821in}}%
\pgfpathcurveto{\pgfqpoint{0.012425in}{-0.007216in}}{\pgfqpoint{0.013889in}{-0.003683in}}{\pgfqpoint{0.013889in}{0.000000in}}%
\pgfpathcurveto{\pgfqpoint{0.013889in}{0.003683in}}{\pgfqpoint{0.012425in}{0.007216in}}{\pgfqpoint{0.009821in}{0.009821in}}%
\pgfpathcurveto{\pgfqpoint{0.007216in}{0.012425in}}{\pgfqpoint{0.003683in}{0.013889in}}{\pgfqpoint{0.000000in}{0.013889in}}%
\pgfpathcurveto{\pgfqpoint{-0.003683in}{0.013889in}}{\pgfqpoint{-0.007216in}{0.012425in}}{\pgfqpoint{-0.009821in}{0.009821in}}%
\pgfpathcurveto{\pgfqpoint{-0.012425in}{0.007216in}}{\pgfqpoint{-0.013889in}{0.003683in}}{\pgfqpoint{-0.013889in}{0.000000in}}%
\pgfpathcurveto{\pgfqpoint{-0.013889in}{-0.003683in}}{\pgfqpoint{-0.012425in}{-0.007216in}}{\pgfqpoint{-0.009821in}{-0.009821in}}%
\pgfpathcurveto{\pgfqpoint{-0.007216in}{-0.012425in}}{\pgfqpoint{-0.003683in}{-0.013889in}}{\pgfqpoint{0.000000in}{-0.013889in}}%
\pgfpathlineto{\pgfqpoint{0.000000in}{-0.013889in}}%
\pgfpathclose%
\pgfusepath{stroke,fill}%
}%
\begin{pgfscope}%
\pgfsys@transformshift{4.457128in}{5.014032in}%
\pgfsys@useobject{currentmarker}{}%
\end{pgfscope}%
\end{pgfscope}%
\begin{pgfscope}%
\pgfpathrectangle{\pgfqpoint{0.100000in}{2.493063in}}{\pgfqpoint{5.037500in}{3.427208in}}%
\pgfusepath{clip}%
\pgfsetrectcap%
\pgfsetroundjoin%
\pgfsetlinewidth{1.505625pt}%
\definecolor{currentstroke}{rgb}{0.678431,1.000000,0.184314}%
\pgfsetstrokecolor{currentstroke}%
\pgfsetstrokeopacity{0.500000}%
\pgfsetdash{}{0pt}%
\pgfpathmoveto{\pgfqpoint{4.380229in}{5.100625in}}%
\pgfusepath{stroke}%
\end{pgfscope}%
\begin{pgfscope}%
\pgfpathrectangle{\pgfqpoint{0.100000in}{2.493063in}}{\pgfqpoint{5.037500in}{3.427208in}}%
\pgfusepath{clip}%
\pgfsetbuttcap%
\pgfsetroundjoin%
\definecolor{currentfill}{rgb}{0.678431,1.000000,0.184314}%
\pgfsetfillcolor{currentfill}%
\pgfsetfillopacity{0.500000}%
\pgfsetlinewidth{0.250937pt}%
\definecolor{currentstroke}{rgb}{0.000000,0.000000,0.000000}%
\pgfsetstrokecolor{currentstroke}%
\pgfsetstrokeopacity{0.500000}%
\pgfsetdash{}{0pt}%
\pgfsys@defobject{currentmarker}{\pgfqpoint{-0.030556in}{-0.030556in}}{\pgfqpoint{0.030556in}{0.030556in}}{%
\pgfpathmoveto{\pgfqpoint{0.000000in}{-0.030556in}}%
\pgfpathcurveto{\pgfqpoint{0.008103in}{-0.030556in}}{\pgfqpoint{0.015876in}{-0.027336in}}{\pgfqpoint{0.021606in}{-0.021606in}}%
\pgfpathcurveto{\pgfqpoint{0.027336in}{-0.015876in}}{\pgfqpoint{0.030556in}{-0.008103in}}{\pgfqpoint{0.030556in}{0.000000in}}%
\pgfpathcurveto{\pgfqpoint{0.030556in}{0.008103in}}{\pgfqpoint{0.027336in}{0.015876in}}{\pgfqpoint{0.021606in}{0.021606in}}%
\pgfpathcurveto{\pgfqpoint{0.015876in}{0.027336in}}{\pgfqpoint{0.008103in}{0.030556in}}{\pgfqpoint{0.000000in}{0.030556in}}%
\pgfpathcurveto{\pgfqpoint{-0.008103in}{0.030556in}}{\pgfqpoint{-0.015876in}{0.027336in}}{\pgfqpoint{-0.021606in}{0.021606in}}%
\pgfpathcurveto{\pgfqpoint{-0.027336in}{0.015876in}}{\pgfqpoint{-0.030556in}{0.008103in}}{\pgfqpoint{-0.030556in}{0.000000in}}%
\pgfpathcurveto{\pgfqpoint{-0.030556in}{-0.008103in}}{\pgfqpoint{-0.027336in}{-0.015876in}}{\pgfqpoint{-0.021606in}{-0.021606in}}%
\pgfpathcurveto{\pgfqpoint{-0.015876in}{-0.027336in}}{\pgfqpoint{-0.008103in}{-0.030556in}}{\pgfqpoint{0.000000in}{-0.030556in}}%
\pgfpathlineto{\pgfqpoint{0.000000in}{-0.030556in}}%
\pgfpathclose%
\pgfusepath{stroke,fill}%
}%
\begin{pgfscope}%
\pgfsys@transformshift{4.380229in}{5.100625in}%
\pgfsys@useobject{currentmarker}{}%
\end{pgfscope}%
\end{pgfscope}%
\begin{pgfscope}%
\pgfpathrectangle{\pgfqpoint{0.100000in}{2.493063in}}{\pgfqpoint{5.037500in}{3.427208in}}%
\pgfusepath{clip}%
\pgfsetrectcap%
\pgfsetroundjoin%
\pgfsetlinewidth{1.505625pt}%
\definecolor{currentstroke}{rgb}{0.501961,0.501961,0.501961}%
\pgfsetstrokecolor{currentstroke}%
\pgfsetstrokeopacity{0.500000}%
\pgfsetdash{}{0pt}%
\pgfpathmoveto{\pgfqpoint{3.967863in}{4.049128in}}%
\pgfusepath{stroke}%
\end{pgfscope}%
\begin{pgfscope}%
\pgfpathrectangle{\pgfqpoint{0.100000in}{2.493063in}}{\pgfqpoint{5.037500in}{3.427208in}}%
\pgfusepath{clip}%
\pgfsetbuttcap%
\pgfsetroundjoin%
\definecolor{currentfill}{rgb}{0.501961,0.501961,0.501961}%
\pgfsetfillcolor{currentfill}%
\pgfsetfillopacity{0.500000}%
\pgfsetlinewidth{0.250937pt}%
\definecolor{currentstroke}{rgb}{0.000000,0.000000,0.000000}%
\pgfsetstrokecolor{currentstroke}%
\pgfsetstrokeopacity{0.500000}%
\pgfsetdash{}{0pt}%
\pgfsys@defobject{currentmarker}{\pgfqpoint{-0.013889in}{-0.013889in}}{\pgfqpoint{0.013889in}{0.013889in}}{%
\pgfpathmoveto{\pgfqpoint{0.000000in}{-0.013889in}}%
\pgfpathcurveto{\pgfqpoint{0.003683in}{-0.013889in}}{\pgfqpoint{0.007216in}{-0.012425in}}{\pgfqpoint{0.009821in}{-0.009821in}}%
\pgfpathcurveto{\pgfqpoint{0.012425in}{-0.007216in}}{\pgfqpoint{0.013889in}{-0.003683in}}{\pgfqpoint{0.013889in}{0.000000in}}%
\pgfpathcurveto{\pgfqpoint{0.013889in}{0.003683in}}{\pgfqpoint{0.012425in}{0.007216in}}{\pgfqpoint{0.009821in}{0.009821in}}%
\pgfpathcurveto{\pgfqpoint{0.007216in}{0.012425in}}{\pgfqpoint{0.003683in}{0.013889in}}{\pgfqpoint{0.000000in}{0.013889in}}%
\pgfpathcurveto{\pgfqpoint{-0.003683in}{0.013889in}}{\pgfqpoint{-0.007216in}{0.012425in}}{\pgfqpoint{-0.009821in}{0.009821in}}%
\pgfpathcurveto{\pgfqpoint{-0.012425in}{0.007216in}}{\pgfqpoint{-0.013889in}{0.003683in}}{\pgfqpoint{-0.013889in}{0.000000in}}%
\pgfpathcurveto{\pgfqpoint{-0.013889in}{-0.003683in}}{\pgfqpoint{-0.012425in}{-0.007216in}}{\pgfqpoint{-0.009821in}{-0.009821in}}%
\pgfpathcurveto{\pgfqpoint{-0.007216in}{-0.012425in}}{\pgfqpoint{-0.003683in}{-0.013889in}}{\pgfqpoint{0.000000in}{-0.013889in}}%
\pgfpathlineto{\pgfqpoint{0.000000in}{-0.013889in}}%
\pgfpathclose%
\pgfusepath{stroke,fill}%
}%
\begin{pgfscope}%
\pgfsys@transformshift{3.967863in}{4.049128in}%
\pgfsys@useobject{currentmarker}{}%
\end{pgfscope}%
\end{pgfscope}%
\begin{pgfscope}%
\pgfpathrectangle{\pgfqpoint{0.100000in}{2.493063in}}{\pgfqpoint{5.037500in}{3.427208in}}%
\pgfusepath{clip}%
\pgfsetrectcap%
\pgfsetroundjoin%
\pgfsetlinewidth{1.505625pt}%
\definecolor{currentstroke}{rgb}{0.501961,0.501961,0.501961}%
\pgfsetstrokecolor{currentstroke}%
\pgfsetstrokeopacity{0.500000}%
\pgfsetdash{}{0pt}%
\pgfpathmoveto{\pgfqpoint{4.246950in}{4.150213in}}%
\pgfusepath{stroke}%
\end{pgfscope}%
\begin{pgfscope}%
\pgfpathrectangle{\pgfqpoint{0.100000in}{2.493063in}}{\pgfqpoint{5.037500in}{3.427208in}}%
\pgfusepath{clip}%
\pgfsetbuttcap%
\pgfsetroundjoin%
\definecolor{currentfill}{rgb}{0.501961,0.501961,0.501961}%
\pgfsetfillcolor{currentfill}%
\pgfsetfillopacity{0.500000}%
\pgfsetlinewidth{0.250937pt}%
\definecolor{currentstroke}{rgb}{0.000000,0.000000,0.000000}%
\pgfsetstrokecolor{currentstroke}%
\pgfsetstrokeopacity{0.500000}%
\pgfsetdash{}{0pt}%
\pgfsys@defobject{currentmarker}{\pgfqpoint{-0.013889in}{-0.013889in}}{\pgfqpoint{0.013889in}{0.013889in}}{%
\pgfpathmoveto{\pgfqpoint{0.000000in}{-0.013889in}}%
\pgfpathcurveto{\pgfqpoint{0.003683in}{-0.013889in}}{\pgfqpoint{0.007216in}{-0.012425in}}{\pgfqpoint{0.009821in}{-0.009821in}}%
\pgfpathcurveto{\pgfqpoint{0.012425in}{-0.007216in}}{\pgfqpoint{0.013889in}{-0.003683in}}{\pgfqpoint{0.013889in}{0.000000in}}%
\pgfpathcurveto{\pgfqpoint{0.013889in}{0.003683in}}{\pgfqpoint{0.012425in}{0.007216in}}{\pgfqpoint{0.009821in}{0.009821in}}%
\pgfpathcurveto{\pgfqpoint{0.007216in}{0.012425in}}{\pgfqpoint{0.003683in}{0.013889in}}{\pgfqpoint{0.000000in}{0.013889in}}%
\pgfpathcurveto{\pgfqpoint{-0.003683in}{0.013889in}}{\pgfqpoint{-0.007216in}{0.012425in}}{\pgfqpoint{-0.009821in}{0.009821in}}%
\pgfpathcurveto{\pgfqpoint{-0.012425in}{0.007216in}}{\pgfqpoint{-0.013889in}{0.003683in}}{\pgfqpoint{-0.013889in}{0.000000in}}%
\pgfpathcurveto{\pgfqpoint{-0.013889in}{-0.003683in}}{\pgfqpoint{-0.012425in}{-0.007216in}}{\pgfqpoint{-0.009821in}{-0.009821in}}%
\pgfpathcurveto{\pgfqpoint{-0.007216in}{-0.012425in}}{\pgfqpoint{-0.003683in}{-0.013889in}}{\pgfqpoint{0.000000in}{-0.013889in}}%
\pgfpathlineto{\pgfqpoint{0.000000in}{-0.013889in}}%
\pgfpathclose%
\pgfusepath{stroke,fill}%
}%
\begin{pgfscope}%
\pgfsys@transformshift{4.246950in}{4.150213in}%
\pgfsys@useobject{currentmarker}{}%
\end{pgfscope}%
\end{pgfscope}%
\begin{pgfscope}%
\pgfpathrectangle{\pgfqpoint{0.100000in}{2.493063in}}{\pgfqpoint{5.037500in}{3.427208in}}%
\pgfusepath{clip}%
\pgfsetrectcap%
\pgfsetroundjoin%
\pgfsetlinewidth{1.505625pt}%
\definecolor{currentstroke}{rgb}{0.501961,0.501961,0.501961}%
\pgfsetstrokecolor{currentstroke}%
\pgfsetstrokeopacity{0.500000}%
\pgfsetdash{}{0pt}%
\pgfpathmoveto{\pgfqpoint{4.133311in}{4.029918in}}%
\pgfusepath{stroke}%
\end{pgfscope}%
\begin{pgfscope}%
\pgfpathrectangle{\pgfqpoint{0.100000in}{2.493063in}}{\pgfqpoint{5.037500in}{3.427208in}}%
\pgfusepath{clip}%
\pgfsetbuttcap%
\pgfsetroundjoin%
\definecolor{currentfill}{rgb}{0.501961,0.501961,0.501961}%
\pgfsetfillcolor{currentfill}%
\pgfsetfillopacity{0.500000}%
\pgfsetlinewidth{0.250937pt}%
\definecolor{currentstroke}{rgb}{0.000000,0.000000,0.000000}%
\pgfsetstrokecolor{currentstroke}%
\pgfsetstrokeopacity{0.500000}%
\pgfsetdash{}{0pt}%
\pgfsys@defobject{currentmarker}{\pgfqpoint{-0.013889in}{-0.013889in}}{\pgfqpoint{0.013889in}{0.013889in}}{%
\pgfpathmoveto{\pgfqpoint{0.000000in}{-0.013889in}}%
\pgfpathcurveto{\pgfqpoint{0.003683in}{-0.013889in}}{\pgfqpoint{0.007216in}{-0.012425in}}{\pgfqpoint{0.009821in}{-0.009821in}}%
\pgfpathcurveto{\pgfqpoint{0.012425in}{-0.007216in}}{\pgfqpoint{0.013889in}{-0.003683in}}{\pgfqpoint{0.013889in}{0.000000in}}%
\pgfpathcurveto{\pgfqpoint{0.013889in}{0.003683in}}{\pgfqpoint{0.012425in}{0.007216in}}{\pgfqpoint{0.009821in}{0.009821in}}%
\pgfpathcurveto{\pgfqpoint{0.007216in}{0.012425in}}{\pgfqpoint{0.003683in}{0.013889in}}{\pgfqpoint{0.000000in}{0.013889in}}%
\pgfpathcurveto{\pgfqpoint{-0.003683in}{0.013889in}}{\pgfqpoint{-0.007216in}{0.012425in}}{\pgfqpoint{-0.009821in}{0.009821in}}%
\pgfpathcurveto{\pgfqpoint{-0.012425in}{0.007216in}}{\pgfqpoint{-0.013889in}{0.003683in}}{\pgfqpoint{-0.013889in}{0.000000in}}%
\pgfpathcurveto{\pgfqpoint{-0.013889in}{-0.003683in}}{\pgfqpoint{-0.012425in}{-0.007216in}}{\pgfqpoint{-0.009821in}{-0.009821in}}%
\pgfpathcurveto{\pgfqpoint{-0.007216in}{-0.012425in}}{\pgfqpoint{-0.003683in}{-0.013889in}}{\pgfqpoint{0.000000in}{-0.013889in}}%
\pgfpathlineto{\pgfqpoint{0.000000in}{-0.013889in}}%
\pgfpathclose%
\pgfusepath{stroke,fill}%
}%
\begin{pgfscope}%
\pgfsys@transformshift{4.133311in}{4.029918in}%
\pgfsys@useobject{currentmarker}{}%
\end{pgfscope}%
\end{pgfscope}%
\begin{pgfscope}%
\pgfpathrectangle{\pgfqpoint{0.100000in}{2.493063in}}{\pgfqpoint{5.037500in}{3.427208in}}%
\pgfusepath{clip}%
\pgfsetrectcap%
\pgfsetroundjoin%
\pgfsetlinewidth{1.505625pt}%
\definecolor{currentstroke}{rgb}{0.501961,0.501961,0.501961}%
\pgfsetstrokecolor{currentstroke}%
\pgfsetstrokeopacity{0.500000}%
\pgfsetdash{}{0pt}%
\pgfpathmoveto{\pgfqpoint{4.298489in}{4.147381in}}%
\pgfusepath{stroke}%
\end{pgfscope}%
\begin{pgfscope}%
\pgfpathrectangle{\pgfqpoint{0.100000in}{2.493063in}}{\pgfqpoint{5.037500in}{3.427208in}}%
\pgfusepath{clip}%
\pgfsetbuttcap%
\pgfsetroundjoin%
\definecolor{currentfill}{rgb}{0.501961,0.501961,0.501961}%
\pgfsetfillcolor{currentfill}%
\pgfsetfillopacity{0.500000}%
\pgfsetlinewidth{0.250937pt}%
\definecolor{currentstroke}{rgb}{0.000000,0.000000,0.000000}%
\pgfsetstrokecolor{currentstroke}%
\pgfsetstrokeopacity{0.500000}%
\pgfsetdash{}{0pt}%
\pgfsys@defobject{currentmarker}{\pgfqpoint{-0.013889in}{-0.013889in}}{\pgfqpoint{0.013889in}{0.013889in}}{%
\pgfpathmoveto{\pgfqpoint{0.000000in}{-0.013889in}}%
\pgfpathcurveto{\pgfqpoint{0.003683in}{-0.013889in}}{\pgfqpoint{0.007216in}{-0.012425in}}{\pgfqpoint{0.009821in}{-0.009821in}}%
\pgfpathcurveto{\pgfqpoint{0.012425in}{-0.007216in}}{\pgfqpoint{0.013889in}{-0.003683in}}{\pgfqpoint{0.013889in}{0.000000in}}%
\pgfpathcurveto{\pgfqpoint{0.013889in}{0.003683in}}{\pgfqpoint{0.012425in}{0.007216in}}{\pgfqpoint{0.009821in}{0.009821in}}%
\pgfpathcurveto{\pgfqpoint{0.007216in}{0.012425in}}{\pgfqpoint{0.003683in}{0.013889in}}{\pgfqpoint{0.000000in}{0.013889in}}%
\pgfpathcurveto{\pgfqpoint{-0.003683in}{0.013889in}}{\pgfqpoint{-0.007216in}{0.012425in}}{\pgfqpoint{-0.009821in}{0.009821in}}%
\pgfpathcurveto{\pgfqpoint{-0.012425in}{0.007216in}}{\pgfqpoint{-0.013889in}{0.003683in}}{\pgfqpoint{-0.013889in}{0.000000in}}%
\pgfpathcurveto{\pgfqpoint{-0.013889in}{-0.003683in}}{\pgfqpoint{-0.012425in}{-0.007216in}}{\pgfqpoint{-0.009821in}{-0.009821in}}%
\pgfpathcurveto{\pgfqpoint{-0.007216in}{-0.012425in}}{\pgfqpoint{-0.003683in}{-0.013889in}}{\pgfqpoint{0.000000in}{-0.013889in}}%
\pgfpathlineto{\pgfqpoint{0.000000in}{-0.013889in}}%
\pgfpathclose%
\pgfusepath{stroke,fill}%
}%
\begin{pgfscope}%
\pgfsys@transformshift{4.298489in}{4.147381in}%
\pgfsys@useobject{currentmarker}{}%
\end{pgfscope}%
\end{pgfscope}%
\begin{pgfscope}%
\pgfpathrectangle{\pgfqpoint{0.100000in}{2.493063in}}{\pgfqpoint{5.037500in}{3.427208in}}%
\pgfusepath{clip}%
\pgfsetrectcap%
\pgfsetroundjoin%
\pgfsetlinewidth{1.505625pt}%
\definecolor{currentstroke}{rgb}{0.678431,1.000000,0.184314}%
\pgfsetstrokecolor{currentstroke}%
\pgfsetstrokeopacity{0.500000}%
\pgfsetdash{}{0pt}%
\pgfpathmoveto{\pgfqpoint{4.319477in}{4.040847in}}%
\pgfusepath{stroke}%
\end{pgfscope}%
\begin{pgfscope}%
\pgfpathrectangle{\pgfqpoint{0.100000in}{2.493063in}}{\pgfqpoint{5.037500in}{3.427208in}}%
\pgfusepath{clip}%
\pgfsetbuttcap%
\pgfsetroundjoin%
\definecolor{currentfill}{rgb}{0.678431,1.000000,0.184314}%
\pgfsetfillcolor{currentfill}%
\pgfsetfillopacity{0.500000}%
\pgfsetlinewidth{0.250937pt}%
\definecolor{currentstroke}{rgb}{0.000000,0.000000,0.000000}%
\pgfsetstrokecolor{currentstroke}%
\pgfsetstrokeopacity{0.500000}%
\pgfsetdash{}{0pt}%
\pgfsys@defobject{currentmarker}{\pgfqpoint{-0.005556in}{-0.005556in}}{\pgfqpoint{0.005556in}{0.005556in}}{%
\pgfpathmoveto{\pgfqpoint{0.000000in}{-0.005556in}}%
\pgfpathcurveto{\pgfqpoint{0.001473in}{-0.005556in}}{\pgfqpoint{0.002887in}{-0.004970in}}{\pgfqpoint{0.003928in}{-0.003928in}}%
\pgfpathcurveto{\pgfqpoint{0.004970in}{-0.002887in}}{\pgfqpoint{0.005556in}{-0.001473in}}{\pgfqpoint{0.005556in}{0.000000in}}%
\pgfpathcurveto{\pgfqpoint{0.005556in}{0.001473in}}{\pgfqpoint{0.004970in}{0.002887in}}{\pgfqpoint{0.003928in}{0.003928in}}%
\pgfpathcurveto{\pgfqpoint{0.002887in}{0.004970in}}{\pgfqpoint{0.001473in}{0.005556in}}{\pgfqpoint{0.000000in}{0.005556in}}%
\pgfpathcurveto{\pgfqpoint{-0.001473in}{0.005556in}}{\pgfqpoint{-0.002887in}{0.004970in}}{\pgfqpoint{-0.003928in}{0.003928in}}%
\pgfpathcurveto{\pgfqpoint{-0.004970in}{0.002887in}}{\pgfqpoint{-0.005556in}{0.001473in}}{\pgfqpoint{-0.005556in}{0.000000in}}%
\pgfpathcurveto{\pgfqpoint{-0.005556in}{-0.001473in}}{\pgfqpoint{-0.004970in}{-0.002887in}}{\pgfqpoint{-0.003928in}{-0.003928in}}%
\pgfpathcurveto{\pgfqpoint{-0.002887in}{-0.004970in}}{\pgfqpoint{-0.001473in}{-0.005556in}}{\pgfqpoint{0.000000in}{-0.005556in}}%
\pgfpathlineto{\pgfqpoint{0.000000in}{-0.005556in}}%
\pgfpathclose%
\pgfusepath{stroke,fill}%
}%
\begin{pgfscope}%
\pgfsys@transformshift{4.319477in}{4.040847in}%
\pgfsys@useobject{currentmarker}{}%
\end{pgfscope}%
\end{pgfscope}%
\begin{pgfscope}%
\pgfpathrectangle{\pgfqpoint{0.100000in}{2.493063in}}{\pgfqpoint{5.037500in}{3.427208in}}%
\pgfusepath{clip}%
\pgfsetrectcap%
\pgfsetroundjoin%
\pgfsetlinewidth{1.505625pt}%
\definecolor{currentstroke}{rgb}{0.501961,0.501961,0.501961}%
\pgfsetstrokecolor{currentstroke}%
\pgfsetstrokeopacity{0.500000}%
\pgfsetdash{}{0pt}%
\pgfpathmoveto{\pgfqpoint{4.394677in}{4.093655in}}%
\pgfusepath{stroke}%
\end{pgfscope}%
\begin{pgfscope}%
\pgfpathrectangle{\pgfqpoint{0.100000in}{2.493063in}}{\pgfqpoint{5.037500in}{3.427208in}}%
\pgfusepath{clip}%
\pgfsetbuttcap%
\pgfsetroundjoin%
\definecolor{currentfill}{rgb}{0.501961,0.501961,0.501961}%
\pgfsetfillcolor{currentfill}%
\pgfsetfillopacity{0.500000}%
\pgfsetlinewidth{0.250937pt}%
\definecolor{currentstroke}{rgb}{0.000000,0.000000,0.000000}%
\pgfsetstrokecolor{currentstroke}%
\pgfsetstrokeopacity{0.500000}%
\pgfsetdash{}{0pt}%
\pgfsys@defobject{currentmarker}{\pgfqpoint{-0.013889in}{-0.013889in}}{\pgfqpoint{0.013889in}{0.013889in}}{%
\pgfpathmoveto{\pgfqpoint{0.000000in}{-0.013889in}}%
\pgfpathcurveto{\pgfqpoint{0.003683in}{-0.013889in}}{\pgfqpoint{0.007216in}{-0.012425in}}{\pgfqpoint{0.009821in}{-0.009821in}}%
\pgfpathcurveto{\pgfqpoint{0.012425in}{-0.007216in}}{\pgfqpoint{0.013889in}{-0.003683in}}{\pgfqpoint{0.013889in}{0.000000in}}%
\pgfpathcurveto{\pgfqpoint{0.013889in}{0.003683in}}{\pgfqpoint{0.012425in}{0.007216in}}{\pgfqpoint{0.009821in}{0.009821in}}%
\pgfpathcurveto{\pgfqpoint{0.007216in}{0.012425in}}{\pgfqpoint{0.003683in}{0.013889in}}{\pgfqpoint{0.000000in}{0.013889in}}%
\pgfpathcurveto{\pgfqpoint{-0.003683in}{0.013889in}}{\pgfqpoint{-0.007216in}{0.012425in}}{\pgfqpoint{-0.009821in}{0.009821in}}%
\pgfpathcurveto{\pgfqpoint{-0.012425in}{0.007216in}}{\pgfqpoint{-0.013889in}{0.003683in}}{\pgfqpoint{-0.013889in}{0.000000in}}%
\pgfpathcurveto{\pgfqpoint{-0.013889in}{-0.003683in}}{\pgfqpoint{-0.012425in}{-0.007216in}}{\pgfqpoint{-0.009821in}{-0.009821in}}%
\pgfpathcurveto{\pgfqpoint{-0.007216in}{-0.012425in}}{\pgfqpoint{-0.003683in}{-0.013889in}}{\pgfqpoint{0.000000in}{-0.013889in}}%
\pgfpathlineto{\pgfqpoint{0.000000in}{-0.013889in}}%
\pgfpathclose%
\pgfusepath{stroke,fill}%
}%
\begin{pgfscope}%
\pgfsys@transformshift{4.394677in}{4.093655in}%
\pgfsys@useobject{currentmarker}{}%
\end{pgfscope}%
\end{pgfscope}%
\begin{pgfscope}%
\pgfpathrectangle{\pgfqpoint{0.100000in}{2.493063in}}{\pgfqpoint{5.037500in}{3.427208in}}%
\pgfusepath{clip}%
\pgfsetrectcap%
\pgfsetroundjoin%
\pgfsetlinewidth{1.505625pt}%
\definecolor{currentstroke}{rgb}{0.501961,0.501961,0.501961}%
\pgfsetstrokecolor{currentstroke}%
\pgfsetstrokeopacity{0.500000}%
\pgfsetdash{}{0pt}%
\pgfpathmoveto{\pgfqpoint{4.214866in}{4.142032in}}%
\pgfusepath{stroke}%
\end{pgfscope}%
\begin{pgfscope}%
\pgfpathrectangle{\pgfqpoint{0.100000in}{2.493063in}}{\pgfqpoint{5.037500in}{3.427208in}}%
\pgfusepath{clip}%
\pgfsetbuttcap%
\pgfsetroundjoin%
\definecolor{currentfill}{rgb}{0.501961,0.501961,0.501961}%
\pgfsetfillcolor{currentfill}%
\pgfsetfillopacity{0.500000}%
\pgfsetlinewidth{0.250937pt}%
\definecolor{currentstroke}{rgb}{0.000000,0.000000,0.000000}%
\pgfsetstrokecolor{currentstroke}%
\pgfsetstrokeopacity{0.500000}%
\pgfsetdash{}{0pt}%
\pgfsys@defobject{currentmarker}{\pgfqpoint{-0.013889in}{-0.013889in}}{\pgfqpoint{0.013889in}{0.013889in}}{%
\pgfpathmoveto{\pgfqpoint{0.000000in}{-0.013889in}}%
\pgfpathcurveto{\pgfqpoint{0.003683in}{-0.013889in}}{\pgfqpoint{0.007216in}{-0.012425in}}{\pgfqpoint{0.009821in}{-0.009821in}}%
\pgfpathcurveto{\pgfqpoint{0.012425in}{-0.007216in}}{\pgfqpoint{0.013889in}{-0.003683in}}{\pgfqpoint{0.013889in}{0.000000in}}%
\pgfpathcurveto{\pgfqpoint{0.013889in}{0.003683in}}{\pgfqpoint{0.012425in}{0.007216in}}{\pgfqpoint{0.009821in}{0.009821in}}%
\pgfpathcurveto{\pgfqpoint{0.007216in}{0.012425in}}{\pgfqpoint{0.003683in}{0.013889in}}{\pgfqpoint{0.000000in}{0.013889in}}%
\pgfpathcurveto{\pgfqpoint{-0.003683in}{0.013889in}}{\pgfqpoint{-0.007216in}{0.012425in}}{\pgfqpoint{-0.009821in}{0.009821in}}%
\pgfpathcurveto{\pgfqpoint{-0.012425in}{0.007216in}}{\pgfqpoint{-0.013889in}{0.003683in}}{\pgfqpoint{-0.013889in}{0.000000in}}%
\pgfpathcurveto{\pgfqpoint{-0.013889in}{-0.003683in}}{\pgfqpoint{-0.012425in}{-0.007216in}}{\pgfqpoint{-0.009821in}{-0.009821in}}%
\pgfpathcurveto{\pgfqpoint{-0.007216in}{-0.012425in}}{\pgfqpoint{-0.003683in}{-0.013889in}}{\pgfqpoint{0.000000in}{-0.013889in}}%
\pgfpathlineto{\pgfqpoint{0.000000in}{-0.013889in}}%
\pgfpathclose%
\pgfusepath{stroke,fill}%
}%
\begin{pgfscope}%
\pgfsys@transformshift{4.214866in}{4.142032in}%
\pgfsys@useobject{currentmarker}{}%
\end{pgfscope}%
\end{pgfscope}%
\begin{pgfscope}%
\pgfpathrectangle{\pgfqpoint{0.100000in}{2.493063in}}{\pgfqpoint{5.037500in}{3.427208in}}%
\pgfusepath{clip}%
\pgfsetrectcap%
\pgfsetroundjoin%
\pgfsetlinewidth{1.505625pt}%
\definecolor{currentstroke}{rgb}{0.501961,0.501961,0.501961}%
\pgfsetstrokecolor{currentstroke}%
\pgfsetstrokeopacity{0.500000}%
\pgfsetdash{}{0pt}%
\pgfpathmoveto{\pgfqpoint{4.447469in}{4.130603in}}%
\pgfusepath{stroke}%
\end{pgfscope}%
\begin{pgfscope}%
\pgfpathrectangle{\pgfqpoint{0.100000in}{2.493063in}}{\pgfqpoint{5.037500in}{3.427208in}}%
\pgfusepath{clip}%
\pgfsetbuttcap%
\pgfsetroundjoin%
\definecolor{currentfill}{rgb}{0.501961,0.501961,0.501961}%
\pgfsetfillcolor{currentfill}%
\pgfsetfillopacity{0.500000}%
\pgfsetlinewidth{0.250937pt}%
\definecolor{currentstroke}{rgb}{0.000000,0.000000,0.000000}%
\pgfsetstrokecolor{currentstroke}%
\pgfsetstrokeopacity{0.500000}%
\pgfsetdash{}{0pt}%
\pgfsys@defobject{currentmarker}{\pgfqpoint{-0.013889in}{-0.013889in}}{\pgfqpoint{0.013889in}{0.013889in}}{%
\pgfpathmoveto{\pgfqpoint{0.000000in}{-0.013889in}}%
\pgfpathcurveto{\pgfqpoint{0.003683in}{-0.013889in}}{\pgfqpoint{0.007216in}{-0.012425in}}{\pgfqpoint{0.009821in}{-0.009821in}}%
\pgfpathcurveto{\pgfqpoint{0.012425in}{-0.007216in}}{\pgfqpoint{0.013889in}{-0.003683in}}{\pgfqpoint{0.013889in}{0.000000in}}%
\pgfpathcurveto{\pgfqpoint{0.013889in}{0.003683in}}{\pgfqpoint{0.012425in}{0.007216in}}{\pgfqpoint{0.009821in}{0.009821in}}%
\pgfpathcurveto{\pgfqpoint{0.007216in}{0.012425in}}{\pgfqpoint{0.003683in}{0.013889in}}{\pgfqpoint{0.000000in}{0.013889in}}%
\pgfpathcurveto{\pgfqpoint{-0.003683in}{0.013889in}}{\pgfqpoint{-0.007216in}{0.012425in}}{\pgfqpoint{-0.009821in}{0.009821in}}%
\pgfpathcurveto{\pgfqpoint{-0.012425in}{0.007216in}}{\pgfqpoint{-0.013889in}{0.003683in}}{\pgfqpoint{-0.013889in}{0.000000in}}%
\pgfpathcurveto{\pgfqpoint{-0.013889in}{-0.003683in}}{\pgfqpoint{-0.012425in}{-0.007216in}}{\pgfqpoint{-0.009821in}{-0.009821in}}%
\pgfpathcurveto{\pgfqpoint{-0.007216in}{-0.012425in}}{\pgfqpoint{-0.003683in}{-0.013889in}}{\pgfqpoint{0.000000in}{-0.013889in}}%
\pgfpathlineto{\pgfqpoint{0.000000in}{-0.013889in}}%
\pgfpathclose%
\pgfusepath{stroke,fill}%
}%
\begin{pgfscope}%
\pgfsys@transformshift{4.447469in}{4.130603in}%
\pgfsys@useobject{currentmarker}{}%
\end{pgfscope}%
\end{pgfscope}%
\begin{pgfscope}%
\pgfpathrectangle{\pgfqpoint{0.100000in}{2.493063in}}{\pgfqpoint{5.037500in}{3.427208in}}%
\pgfusepath{clip}%
\pgfsetrectcap%
\pgfsetroundjoin%
\pgfsetlinewidth{1.505625pt}%
\definecolor{currentstroke}{rgb}{0.501961,0.501961,0.501961}%
\pgfsetstrokecolor{currentstroke}%
\pgfsetstrokeopacity{0.500000}%
\pgfsetdash{}{0pt}%
\pgfpathmoveto{\pgfqpoint{4.078163in}{4.080492in}}%
\pgfusepath{stroke}%
\end{pgfscope}%
\begin{pgfscope}%
\pgfpathrectangle{\pgfqpoint{0.100000in}{2.493063in}}{\pgfqpoint{5.037500in}{3.427208in}}%
\pgfusepath{clip}%
\pgfsetbuttcap%
\pgfsetroundjoin%
\definecolor{currentfill}{rgb}{0.501961,0.501961,0.501961}%
\pgfsetfillcolor{currentfill}%
\pgfsetfillopacity{0.500000}%
\pgfsetlinewidth{0.250937pt}%
\definecolor{currentstroke}{rgb}{0.000000,0.000000,0.000000}%
\pgfsetstrokecolor{currentstroke}%
\pgfsetstrokeopacity{0.500000}%
\pgfsetdash{}{0pt}%
\pgfsys@defobject{currentmarker}{\pgfqpoint{-0.013889in}{-0.013889in}}{\pgfqpoint{0.013889in}{0.013889in}}{%
\pgfpathmoveto{\pgfqpoint{0.000000in}{-0.013889in}}%
\pgfpathcurveto{\pgfqpoint{0.003683in}{-0.013889in}}{\pgfqpoint{0.007216in}{-0.012425in}}{\pgfqpoint{0.009821in}{-0.009821in}}%
\pgfpathcurveto{\pgfqpoint{0.012425in}{-0.007216in}}{\pgfqpoint{0.013889in}{-0.003683in}}{\pgfqpoint{0.013889in}{0.000000in}}%
\pgfpathcurveto{\pgfqpoint{0.013889in}{0.003683in}}{\pgfqpoint{0.012425in}{0.007216in}}{\pgfqpoint{0.009821in}{0.009821in}}%
\pgfpathcurveto{\pgfqpoint{0.007216in}{0.012425in}}{\pgfqpoint{0.003683in}{0.013889in}}{\pgfqpoint{0.000000in}{0.013889in}}%
\pgfpathcurveto{\pgfqpoint{-0.003683in}{0.013889in}}{\pgfqpoint{-0.007216in}{0.012425in}}{\pgfqpoint{-0.009821in}{0.009821in}}%
\pgfpathcurveto{\pgfqpoint{-0.012425in}{0.007216in}}{\pgfqpoint{-0.013889in}{0.003683in}}{\pgfqpoint{-0.013889in}{0.000000in}}%
\pgfpathcurveto{\pgfqpoint{-0.013889in}{-0.003683in}}{\pgfqpoint{-0.012425in}{-0.007216in}}{\pgfqpoint{-0.009821in}{-0.009821in}}%
\pgfpathcurveto{\pgfqpoint{-0.007216in}{-0.012425in}}{\pgfqpoint{-0.003683in}{-0.013889in}}{\pgfqpoint{0.000000in}{-0.013889in}}%
\pgfpathlineto{\pgfqpoint{0.000000in}{-0.013889in}}%
\pgfpathclose%
\pgfusepath{stroke,fill}%
}%
\begin{pgfscope}%
\pgfsys@transformshift{4.078163in}{4.080492in}%
\pgfsys@useobject{currentmarker}{}%
\end{pgfscope}%
\end{pgfscope}%
\begin{pgfscope}%
\pgfpathrectangle{\pgfqpoint{0.100000in}{2.493063in}}{\pgfqpoint{5.037500in}{3.427208in}}%
\pgfusepath{clip}%
\pgfsetrectcap%
\pgfsetroundjoin%
\pgfsetlinewidth{1.505625pt}%
\definecolor{currentstroke}{rgb}{0.678431,1.000000,0.184314}%
\pgfsetstrokecolor{currentstroke}%
\pgfsetstrokeopacity{0.500000}%
\pgfsetdash{}{0pt}%
\pgfpathmoveto{\pgfqpoint{4.460629in}{4.032203in}}%
\pgfusepath{stroke}%
\end{pgfscope}%
\begin{pgfscope}%
\pgfpathrectangle{\pgfqpoint{0.100000in}{2.493063in}}{\pgfqpoint{5.037500in}{3.427208in}}%
\pgfusepath{clip}%
\pgfsetbuttcap%
\pgfsetroundjoin%
\definecolor{currentfill}{rgb}{0.678431,1.000000,0.184314}%
\pgfsetfillcolor{currentfill}%
\pgfsetfillopacity{0.500000}%
\pgfsetlinewidth{0.250937pt}%
\definecolor{currentstroke}{rgb}{0.000000,0.000000,0.000000}%
\pgfsetstrokecolor{currentstroke}%
\pgfsetstrokeopacity{0.500000}%
\pgfsetdash{}{0pt}%
\pgfsys@defobject{currentmarker}{\pgfqpoint{-0.005556in}{-0.005556in}}{\pgfqpoint{0.005556in}{0.005556in}}{%
\pgfpathmoveto{\pgfqpoint{0.000000in}{-0.005556in}}%
\pgfpathcurveto{\pgfqpoint{0.001473in}{-0.005556in}}{\pgfqpoint{0.002887in}{-0.004970in}}{\pgfqpoint{0.003928in}{-0.003928in}}%
\pgfpathcurveto{\pgfqpoint{0.004970in}{-0.002887in}}{\pgfqpoint{0.005556in}{-0.001473in}}{\pgfqpoint{0.005556in}{0.000000in}}%
\pgfpathcurveto{\pgfqpoint{0.005556in}{0.001473in}}{\pgfqpoint{0.004970in}{0.002887in}}{\pgfqpoint{0.003928in}{0.003928in}}%
\pgfpathcurveto{\pgfqpoint{0.002887in}{0.004970in}}{\pgfqpoint{0.001473in}{0.005556in}}{\pgfqpoint{0.000000in}{0.005556in}}%
\pgfpathcurveto{\pgfqpoint{-0.001473in}{0.005556in}}{\pgfqpoint{-0.002887in}{0.004970in}}{\pgfqpoint{-0.003928in}{0.003928in}}%
\pgfpathcurveto{\pgfqpoint{-0.004970in}{0.002887in}}{\pgfqpoint{-0.005556in}{0.001473in}}{\pgfqpoint{-0.005556in}{0.000000in}}%
\pgfpathcurveto{\pgfqpoint{-0.005556in}{-0.001473in}}{\pgfqpoint{-0.004970in}{-0.002887in}}{\pgfqpoint{-0.003928in}{-0.003928in}}%
\pgfpathcurveto{\pgfqpoint{-0.002887in}{-0.004970in}}{\pgfqpoint{-0.001473in}{-0.005556in}}{\pgfqpoint{0.000000in}{-0.005556in}}%
\pgfpathlineto{\pgfqpoint{0.000000in}{-0.005556in}}%
\pgfpathclose%
\pgfusepath{stroke,fill}%
}%
\begin{pgfscope}%
\pgfsys@transformshift{4.460629in}{4.032203in}%
\pgfsys@useobject{currentmarker}{}%
\end{pgfscope}%
\end{pgfscope}%
\begin{pgfscope}%
\pgfpathrectangle{\pgfqpoint{0.100000in}{2.493063in}}{\pgfqpoint{5.037500in}{3.427208in}}%
\pgfusepath{clip}%
\pgfsetrectcap%
\pgfsetroundjoin%
\pgfsetlinewidth{1.505625pt}%
\definecolor{currentstroke}{rgb}{0.678431,1.000000,0.184314}%
\pgfsetstrokecolor{currentstroke}%
\pgfsetstrokeopacity{0.500000}%
\pgfsetdash{}{0pt}%
\pgfpathmoveto{\pgfqpoint{4.488531in}{4.079440in}}%
\pgfusepath{stroke}%
\end{pgfscope}%
\begin{pgfscope}%
\pgfpathrectangle{\pgfqpoint{0.100000in}{2.493063in}}{\pgfqpoint{5.037500in}{3.427208in}}%
\pgfusepath{clip}%
\pgfsetbuttcap%
\pgfsetroundjoin%
\definecolor{currentfill}{rgb}{0.678431,1.000000,0.184314}%
\pgfsetfillcolor{currentfill}%
\pgfsetfillopacity{0.500000}%
\pgfsetlinewidth{0.250937pt}%
\definecolor{currentstroke}{rgb}{0.000000,0.000000,0.000000}%
\pgfsetstrokecolor{currentstroke}%
\pgfsetstrokeopacity{0.500000}%
\pgfsetdash{}{0pt}%
\pgfsys@defobject{currentmarker}{\pgfqpoint{-0.008333in}{-0.008333in}}{\pgfqpoint{0.008333in}{0.008333in}}{%
\pgfpathmoveto{\pgfqpoint{0.000000in}{-0.008333in}}%
\pgfpathcurveto{\pgfqpoint{0.002210in}{-0.008333in}}{\pgfqpoint{0.004330in}{-0.007455in}}{\pgfqpoint{0.005893in}{-0.005893in}}%
\pgfpathcurveto{\pgfqpoint{0.007455in}{-0.004330in}}{\pgfqpoint{0.008333in}{-0.002210in}}{\pgfqpoint{0.008333in}{0.000000in}}%
\pgfpathcurveto{\pgfqpoint{0.008333in}{0.002210in}}{\pgfqpoint{0.007455in}{0.004330in}}{\pgfqpoint{0.005893in}{0.005893in}}%
\pgfpathcurveto{\pgfqpoint{0.004330in}{0.007455in}}{\pgfqpoint{0.002210in}{0.008333in}}{\pgfqpoint{0.000000in}{0.008333in}}%
\pgfpathcurveto{\pgfqpoint{-0.002210in}{0.008333in}}{\pgfqpoint{-0.004330in}{0.007455in}}{\pgfqpoint{-0.005893in}{0.005893in}}%
\pgfpathcurveto{\pgfqpoint{-0.007455in}{0.004330in}}{\pgfqpoint{-0.008333in}{0.002210in}}{\pgfqpoint{-0.008333in}{0.000000in}}%
\pgfpathcurveto{\pgfqpoint{-0.008333in}{-0.002210in}}{\pgfqpoint{-0.007455in}{-0.004330in}}{\pgfqpoint{-0.005893in}{-0.005893in}}%
\pgfpathcurveto{\pgfqpoint{-0.004330in}{-0.007455in}}{\pgfqpoint{-0.002210in}{-0.008333in}}{\pgfqpoint{0.000000in}{-0.008333in}}%
\pgfpathlineto{\pgfqpoint{0.000000in}{-0.008333in}}%
\pgfpathclose%
\pgfusepath{stroke,fill}%
}%
\begin{pgfscope}%
\pgfsys@transformshift{4.488531in}{4.079440in}%
\pgfsys@useobject{currentmarker}{}%
\end{pgfscope}%
\end{pgfscope}%
\begin{pgfscope}%
\pgfpathrectangle{\pgfqpoint{0.100000in}{2.493063in}}{\pgfqpoint{5.037500in}{3.427208in}}%
\pgfusepath{clip}%
\pgfsetrectcap%
\pgfsetroundjoin%
\pgfsetlinewidth{1.505625pt}%
\definecolor{currentstroke}{rgb}{0.678431,1.000000,0.184314}%
\pgfsetstrokecolor{currentstroke}%
\pgfsetstrokeopacity{0.500000}%
\pgfsetdash{}{0pt}%
\pgfpathmoveto{\pgfqpoint{4.326728in}{4.127444in}}%
\pgfusepath{stroke}%
\end{pgfscope}%
\begin{pgfscope}%
\pgfpathrectangle{\pgfqpoint{0.100000in}{2.493063in}}{\pgfqpoint{5.037500in}{3.427208in}}%
\pgfusepath{clip}%
\pgfsetbuttcap%
\pgfsetroundjoin%
\definecolor{currentfill}{rgb}{0.678431,1.000000,0.184314}%
\pgfsetfillcolor{currentfill}%
\pgfsetfillopacity{0.500000}%
\pgfsetlinewidth{0.250937pt}%
\definecolor{currentstroke}{rgb}{0.000000,0.000000,0.000000}%
\pgfsetstrokecolor{currentstroke}%
\pgfsetstrokeopacity{0.500000}%
\pgfsetdash{}{0pt}%
\pgfsys@defobject{currentmarker}{\pgfqpoint{-0.005556in}{-0.005556in}}{\pgfqpoint{0.005556in}{0.005556in}}{%
\pgfpathmoveto{\pgfqpoint{0.000000in}{-0.005556in}}%
\pgfpathcurveto{\pgfqpoint{0.001473in}{-0.005556in}}{\pgfqpoint{0.002887in}{-0.004970in}}{\pgfqpoint{0.003928in}{-0.003928in}}%
\pgfpathcurveto{\pgfqpoint{0.004970in}{-0.002887in}}{\pgfqpoint{0.005556in}{-0.001473in}}{\pgfqpoint{0.005556in}{0.000000in}}%
\pgfpathcurveto{\pgfqpoint{0.005556in}{0.001473in}}{\pgfqpoint{0.004970in}{0.002887in}}{\pgfqpoint{0.003928in}{0.003928in}}%
\pgfpathcurveto{\pgfqpoint{0.002887in}{0.004970in}}{\pgfqpoint{0.001473in}{0.005556in}}{\pgfqpoint{0.000000in}{0.005556in}}%
\pgfpathcurveto{\pgfqpoint{-0.001473in}{0.005556in}}{\pgfqpoint{-0.002887in}{0.004970in}}{\pgfqpoint{-0.003928in}{0.003928in}}%
\pgfpathcurveto{\pgfqpoint{-0.004970in}{0.002887in}}{\pgfqpoint{-0.005556in}{0.001473in}}{\pgfqpoint{-0.005556in}{0.000000in}}%
\pgfpathcurveto{\pgfqpoint{-0.005556in}{-0.001473in}}{\pgfqpoint{-0.004970in}{-0.002887in}}{\pgfqpoint{-0.003928in}{-0.003928in}}%
\pgfpathcurveto{\pgfqpoint{-0.002887in}{-0.004970in}}{\pgfqpoint{-0.001473in}{-0.005556in}}{\pgfqpoint{0.000000in}{-0.005556in}}%
\pgfpathlineto{\pgfqpoint{0.000000in}{-0.005556in}}%
\pgfpathclose%
\pgfusepath{stroke,fill}%
}%
\begin{pgfscope}%
\pgfsys@transformshift{4.326728in}{4.127444in}%
\pgfsys@useobject{currentmarker}{}%
\end{pgfscope}%
\end{pgfscope}%
\begin{pgfscope}%
\pgfpathrectangle{\pgfqpoint{0.100000in}{2.493063in}}{\pgfqpoint{5.037500in}{3.427208in}}%
\pgfusepath{clip}%
\pgfsetrectcap%
\pgfsetroundjoin%
\pgfsetlinewidth{1.505625pt}%
\definecolor{currentstroke}{rgb}{0.678431,1.000000,0.184314}%
\pgfsetstrokecolor{currentstroke}%
\pgfsetstrokeopacity{0.500000}%
\pgfsetdash{}{0pt}%
\pgfpathmoveto{\pgfqpoint{4.400786in}{4.160282in}}%
\pgfusepath{stroke}%
\end{pgfscope}%
\begin{pgfscope}%
\pgfpathrectangle{\pgfqpoint{0.100000in}{2.493063in}}{\pgfqpoint{5.037500in}{3.427208in}}%
\pgfusepath{clip}%
\pgfsetbuttcap%
\pgfsetroundjoin%
\definecolor{currentfill}{rgb}{0.678431,1.000000,0.184314}%
\pgfsetfillcolor{currentfill}%
\pgfsetfillopacity{0.500000}%
\pgfsetlinewidth{0.250937pt}%
\definecolor{currentstroke}{rgb}{0.000000,0.000000,0.000000}%
\pgfsetstrokecolor{currentstroke}%
\pgfsetstrokeopacity{0.500000}%
\pgfsetdash{}{0pt}%
\pgfsys@defobject{currentmarker}{\pgfqpoint{-0.005556in}{-0.005556in}}{\pgfqpoint{0.005556in}{0.005556in}}{%
\pgfpathmoveto{\pgfqpoint{0.000000in}{-0.005556in}}%
\pgfpathcurveto{\pgfqpoint{0.001473in}{-0.005556in}}{\pgfqpoint{0.002887in}{-0.004970in}}{\pgfqpoint{0.003928in}{-0.003928in}}%
\pgfpathcurveto{\pgfqpoint{0.004970in}{-0.002887in}}{\pgfqpoint{0.005556in}{-0.001473in}}{\pgfqpoint{0.005556in}{0.000000in}}%
\pgfpathcurveto{\pgfqpoint{0.005556in}{0.001473in}}{\pgfqpoint{0.004970in}{0.002887in}}{\pgfqpoint{0.003928in}{0.003928in}}%
\pgfpathcurveto{\pgfqpoint{0.002887in}{0.004970in}}{\pgfqpoint{0.001473in}{0.005556in}}{\pgfqpoint{0.000000in}{0.005556in}}%
\pgfpathcurveto{\pgfqpoint{-0.001473in}{0.005556in}}{\pgfqpoint{-0.002887in}{0.004970in}}{\pgfqpoint{-0.003928in}{0.003928in}}%
\pgfpathcurveto{\pgfqpoint{-0.004970in}{0.002887in}}{\pgfqpoint{-0.005556in}{0.001473in}}{\pgfqpoint{-0.005556in}{0.000000in}}%
\pgfpathcurveto{\pgfqpoint{-0.005556in}{-0.001473in}}{\pgfqpoint{-0.004970in}{-0.002887in}}{\pgfqpoint{-0.003928in}{-0.003928in}}%
\pgfpathcurveto{\pgfqpoint{-0.002887in}{-0.004970in}}{\pgfqpoint{-0.001473in}{-0.005556in}}{\pgfqpoint{0.000000in}{-0.005556in}}%
\pgfpathlineto{\pgfqpoint{0.000000in}{-0.005556in}}%
\pgfpathclose%
\pgfusepath{stroke,fill}%
}%
\begin{pgfscope}%
\pgfsys@transformshift{4.400786in}{4.160282in}%
\pgfsys@useobject{currentmarker}{}%
\end{pgfscope}%
\end{pgfscope}%
\begin{pgfscope}%
\pgfpathrectangle{\pgfqpoint{0.100000in}{2.493063in}}{\pgfqpoint{5.037500in}{3.427208in}}%
\pgfusepath{clip}%
\pgfsetrectcap%
\pgfsetroundjoin%
\pgfsetlinewidth{1.505625pt}%
\definecolor{currentstroke}{rgb}{0.501961,0.501961,0.501961}%
\pgfsetstrokecolor{currentstroke}%
\pgfsetstrokeopacity{0.500000}%
\pgfsetdash{}{0pt}%
\pgfpathmoveto{\pgfqpoint{4.424090in}{3.963005in}}%
\pgfusepath{stroke}%
\end{pgfscope}%
\begin{pgfscope}%
\pgfpathrectangle{\pgfqpoint{0.100000in}{2.493063in}}{\pgfqpoint{5.037500in}{3.427208in}}%
\pgfusepath{clip}%
\pgfsetbuttcap%
\pgfsetroundjoin%
\definecolor{currentfill}{rgb}{0.501961,0.501961,0.501961}%
\pgfsetfillcolor{currentfill}%
\pgfsetfillopacity{0.500000}%
\pgfsetlinewidth{0.250937pt}%
\definecolor{currentstroke}{rgb}{0.000000,0.000000,0.000000}%
\pgfsetstrokecolor{currentstroke}%
\pgfsetstrokeopacity{0.500000}%
\pgfsetdash{}{0pt}%
\pgfsys@defobject{currentmarker}{\pgfqpoint{-0.013889in}{-0.013889in}}{\pgfqpoint{0.013889in}{0.013889in}}{%
\pgfpathmoveto{\pgfqpoint{0.000000in}{-0.013889in}}%
\pgfpathcurveto{\pgfqpoint{0.003683in}{-0.013889in}}{\pgfqpoint{0.007216in}{-0.012425in}}{\pgfqpoint{0.009821in}{-0.009821in}}%
\pgfpathcurveto{\pgfqpoint{0.012425in}{-0.007216in}}{\pgfqpoint{0.013889in}{-0.003683in}}{\pgfqpoint{0.013889in}{0.000000in}}%
\pgfpathcurveto{\pgfqpoint{0.013889in}{0.003683in}}{\pgfqpoint{0.012425in}{0.007216in}}{\pgfqpoint{0.009821in}{0.009821in}}%
\pgfpathcurveto{\pgfqpoint{0.007216in}{0.012425in}}{\pgfqpoint{0.003683in}{0.013889in}}{\pgfqpoint{0.000000in}{0.013889in}}%
\pgfpathcurveto{\pgfqpoint{-0.003683in}{0.013889in}}{\pgfqpoint{-0.007216in}{0.012425in}}{\pgfqpoint{-0.009821in}{0.009821in}}%
\pgfpathcurveto{\pgfqpoint{-0.012425in}{0.007216in}}{\pgfqpoint{-0.013889in}{0.003683in}}{\pgfqpoint{-0.013889in}{0.000000in}}%
\pgfpathcurveto{\pgfqpoint{-0.013889in}{-0.003683in}}{\pgfqpoint{-0.012425in}{-0.007216in}}{\pgfqpoint{-0.009821in}{-0.009821in}}%
\pgfpathcurveto{\pgfqpoint{-0.007216in}{-0.012425in}}{\pgfqpoint{-0.003683in}{-0.013889in}}{\pgfqpoint{0.000000in}{-0.013889in}}%
\pgfpathlineto{\pgfqpoint{0.000000in}{-0.013889in}}%
\pgfpathclose%
\pgfusepath{stroke,fill}%
}%
\begin{pgfscope}%
\pgfsys@transformshift{4.424090in}{3.963005in}%
\pgfsys@useobject{currentmarker}{}%
\end{pgfscope}%
\end{pgfscope}%
\begin{pgfscope}%
\pgfpathrectangle{\pgfqpoint{0.100000in}{2.493063in}}{\pgfqpoint{5.037500in}{3.427208in}}%
\pgfusepath{clip}%
\pgfsetrectcap%
\pgfsetroundjoin%
\pgfsetlinewidth{1.505625pt}%
\definecolor{currentstroke}{rgb}{0.501961,0.501961,0.501961}%
\pgfsetstrokecolor{currentstroke}%
\pgfsetstrokeopacity{0.500000}%
\pgfsetdash{}{0pt}%
\pgfpathmoveto{\pgfqpoint{4.172799in}{4.138209in}}%
\pgfusepath{stroke}%
\end{pgfscope}%
\begin{pgfscope}%
\pgfpathrectangle{\pgfqpoint{0.100000in}{2.493063in}}{\pgfqpoint{5.037500in}{3.427208in}}%
\pgfusepath{clip}%
\pgfsetbuttcap%
\pgfsetroundjoin%
\definecolor{currentfill}{rgb}{0.501961,0.501961,0.501961}%
\pgfsetfillcolor{currentfill}%
\pgfsetfillopacity{0.500000}%
\pgfsetlinewidth{0.250937pt}%
\definecolor{currentstroke}{rgb}{0.000000,0.000000,0.000000}%
\pgfsetstrokecolor{currentstroke}%
\pgfsetstrokeopacity{0.500000}%
\pgfsetdash{}{0pt}%
\pgfsys@defobject{currentmarker}{\pgfqpoint{-0.013889in}{-0.013889in}}{\pgfqpoint{0.013889in}{0.013889in}}{%
\pgfpathmoveto{\pgfqpoint{0.000000in}{-0.013889in}}%
\pgfpathcurveto{\pgfqpoint{0.003683in}{-0.013889in}}{\pgfqpoint{0.007216in}{-0.012425in}}{\pgfqpoint{0.009821in}{-0.009821in}}%
\pgfpathcurveto{\pgfqpoint{0.012425in}{-0.007216in}}{\pgfqpoint{0.013889in}{-0.003683in}}{\pgfqpoint{0.013889in}{0.000000in}}%
\pgfpathcurveto{\pgfqpoint{0.013889in}{0.003683in}}{\pgfqpoint{0.012425in}{0.007216in}}{\pgfqpoint{0.009821in}{0.009821in}}%
\pgfpathcurveto{\pgfqpoint{0.007216in}{0.012425in}}{\pgfqpoint{0.003683in}{0.013889in}}{\pgfqpoint{0.000000in}{0.013889in}}%
\pgfpathcurveto{\pgfqpoint{-0.003683in}{0.013889in}}{\pgfqpoint{-0.007216in}{0.012425in}}{\pgfqpoint{-0.009821in}{0.009821in}}%
\pgfpathcurveto{\pgfqpoint{-0.012425in}{0.007216in}}{\pgfqpoint{-0.013889in}{0.003683in}}{\pgfqpoint{-0.013889in}{0.000000in}}%
\pgfpathcurveto{\pgfqpoint{-0.013889in}{-0.003683in}}{\pgfqpoint{-0.012425in}{-0.007216in}}{\pgfqpoint{-0.009821in}{-0.009821in}}%
\pgfpathcurveto{\pgfqpoint{-0.007216in}{-0.012425in}}{\pgfqpoint{-0.003683in}{-0.013889in}}{\pgfqpoint{0.000000in}{-0.013889in}}%
\pgfpathlineto{\pgfqpoint{0.000000in}{-0.013889in}}%
\pgfpathclose%
\pgfusepath{stroke,fill}%
}%
\begin{pgfscope}%
\pgfsys@transformshift{4.172799in}{4.138209in}%
\pgfsys@useobject{currentmarker}{}%
\end{pgfscope}%
\end{pgfscope}%
\begin{pgfscope}%
\pgfpathrectangle{\pgfqpoint{0.100000in}{2.493063in}}{\pgfqpoint{5.037500in}{3.427208in}}%
\pgfusepath{clip}%
\pgfsetrectcap%
\pgfsetroundjoin%
\pgfsetlinewidth{1.505625pt}%
\definecolor{currentstroke}{rgb}{0.678431,1.000000,0.184314}%
\pgfsetstrokecolor{currentstroke}%
\pgfsetstrokeopacity{0.500000}%
\pgfsetdash{}{0pt}%
\pgfpathmoveto{\pgfqpoint{2.344274in}{5.277479in}}%
\pgfusepath{stroke}%
\end{pgfscope}%
\begin{pgfscope}%
\pgfpathrectangle{\pgfqpoint{0.100000in}{2.493063in}}{\pgfqpoint{5.037500in}{3.427208in}}%
\pgfusepath{clip}%
\pgfsetbuttcap%
\pgfsetroundjoin%
\definecolor{currentfill}{rgb}{0.678431,1.000000,0.184314}%
\pgfsetfillcolor{currentfill}%
\pgfsetfillopacity{0.500000}%
\pgfsetlinewidth{0.250937pt}%
\definecolor{currentstroke}{rgb}{0.000000,0.000000,0.000000}%
\pgfsetstrokecolor{currentstroke}%
\pgfsetstrokeopacity{0.500000}%
\pgfsetdash{}{0pt}%
\pgfsys@defobject{currentmarker}{\pgfqpoint{-0.011111in}{-0.011111in}}{\pgfqpoint{0.011111in}{0.011111in}}{%
\pgfpathmoveto{\pgfqpoint{0.000000in}{-0.011111in}}%
\pgfpathcurveto{\pgfqpoint{0.002947in}{-0.011111in}}{\pgfqpoint{0.005773in}{-0.009940in}}{\pgfqpoint{0.007857in}{-0.007857in}}%
\pgfpathcurveto{\pgfqpoint{0.009940in}{-0.005773in}}{\pgfqpoint{0.011111in}{-0.002947in}}{\pgfqpoint{0.011111in}{0.000000in}}%
\pgfpathcurveto{\pgfqpoint{0.011111in}{0.002947in}}{\pgfqpoint{0.009940in}{0.005773in}}{\pgfqpoint{0.007857in}{0.007857in}}%
\pgfpathcurveto{\pgfqpoint{0.005773in}{0.009940in}}{\pgfqpoint{0.002947in}{0.011111in}}{\pgfqpoint{0.000000in}{0.011111in}}%
\pgfpathcurveto{\pgfqpoint{-0.002947in}{0.011111in}}{\pgfqpoint{-0.005773in}{0.009940in}}{\pgfqpoint{-0.007857in}{0.007857in}}%
\pgfpathcurveto{\pgfqpoint{-0.009940in}{0.005773in}}{\pgfqpoint{-0.011111in}{0.002947in}}{\pgfqpoint{-0.011111in}{0.000000in}}%
\pgfpathcurveto{\pgfqpoint{-0.011111in}{-0.002947in}}{\pgfqpoint{-0.009940in}{-0.005773in}}{\pgfqpoint{-0.007857in}{-0.007857in}}%
\pgfpathcurveto{\pgfqpoint{-0.005773in}{-0.009940in}}{\pgfqpoint{-0.002947in}{-0.011111in}}{\pgfqpoint{0.000000in}{-0.011111in}}%
\pgfpathlineto{\pgfqpoint{0.000000in}{-0.011111in}}%
\pgfpathclose%
\pgfusepath{stroke,fill}%
}%
\begin{pgfscope}%
\pgfsys@transformshift{2.344274in}{5.277479in}%
\pgfsys@useobject{currentmarker}{}%
\end{pgfscope}%
\end{pgfscope}%
\begin{pgfscope}%
\pgfpathrectangle{\pgfqpoint{0.100000in}{2.493063in}}{\pgfqpoint{5.037500in}{3.427208in}}%
\pgfusepath{clip}%
\pgfsetrectcap%
\pgfsetroundjoin%
\pgfsetlinewidth{1.505625pt}%
\definecolor{currentstroke}{rgb}{0.678431,1.000000,0.184314}%
\pgfsetstrokecolor{currentstroke}%
\pgfsetstrokeopacity{0.500000}%
\pgfsetdash{}{0pt}%
\pgfpathmoveto{\pgfqpoint{2.661894in}{5.272252in}}%
\pgfusepath{stroke}%
\end{pgfscope}%
\begin{pgfscope}%
\pgfpathrectangle{\pgfqpoint{0.100000in}{2.493063in}}{\pgfqpoint{5.037500in}{3.427208in}}%
\pgfusepath{clip}%
\pgfsetbuttcap%
\pgfsetroundjoin%
\definecolor{currentfill}{rgb}{0.678431,1.000000,0.184314}%
\pgfsetfillcolor{currentfill}%
\pgfsetfillopacity{0.500000}%
\pgfsetlinewidth{0.250937pt}%
\definecolor{currentstroke}{rgb}{0.000000,0.000000,0.000000}%
\pgfsetstrokecolor{currentstroke}%
\pgfsetstrokeopacity{0.500000}%
\pgfsetdash{}{0pt}%
\pgfsys@defobject{currentmarker}{\pgfqpoint{-0.008333in}{-0.008333in}}{\pgfqpoint{0.008333in}{0.008333in}}{%
\pgfpathmoveto{\pgfqpoint{0.000000in}{-0.008333in}}%
\pgfpathcurveto{\pgfqpoint{0.002210in}{-0.008333in}}{\pgfqpoint{0.004330in}{-0.007455in}}{\pgfqpoint{0.005893in}{-0.005893in}}%
\pgfpathcurveto{\pgfqpoint{0.007455in}{-0.004330in}}{\pgfqpoint{0.008333in}{-0.002210in}}{\pgfqpoint{0.008333in}{0.000000in}}%
\pgfpathcurveto{\pgfqpoint{0.008333in}{0.002210in}}{\pgfqpoint{0.007455in}{0.004330in}}{\pgfqpoint{0.005893in}{0.005893in}}%
\pgfpathcurveto{\pgfqpoint{0.004330in}{0.007455in}}{\pgfqpoint{0.002210in}{0.008333in}}{\pgfqpoint{0.000000in}{0.008333in}}%
\pgfpathcurveto{\pgfqpoint{-0.002210in}{0.008333in}}{\pgfqpoint{-0.004330in}{0.007455in}}{\pgfqpoint{-0.005893in}{0.005893in}}%
\pgfpathcurveto{\pgfqpoint{-0.007455in}{0.004330in}}{\pgfqpoint{-0.008333in}{0.002210in}}{\pgfqpoint{-0.008333in}{0.000000in}}%
\pgfpathcurveto{\pgfqpoint{-0.008333in}{-0.002210in}}{\pgfqpoint{-0.007455in}{-0.004330in}}{\pgfqpoint{-0.005893in}{-0.005893in}}%
\pgfpathcurveto{\pgfqpoint{-0.004330in}{-0.007455in}}{\pgfqpoint{-0.002210in}{-0.008333in}}{\pgfqpoint{0.000000in}{-0.008333in}}%
\pgfpathlineto{\pgfqpoint{0.000000in}{-0.008333in}}%
\pgfpathclose%
\pgfusepath{stroke,fill}%
}%
\begin{pgfscope}%
\pgfsys@transformshift{2.661894in}{5.272252in}%
\pgfsys@useobject{currentmarker}{}%
\end{pgfscope}%
\end{pgfscope}%
\begin{pgfscope}%
\pgfpathrectangle{\pgfqpoint{0.100000in}{2.493063in}}{\pgfqpoint{5.037500in}{3.427208in}}%
\pgfusepath{clip}%
\pgfsetrectcap%
\pgfsetroundjoin%
\pgfsetlinewidth{1.505625pt}%
\definecolor{currentstroke}{rgb}{0.678431,1.000000,0.184314}%
\pgfsetstrokecolor{currentstroke}%
\pgfsetstrokeopacity{0.500000}%
\pgfsetdash{}{0pt}%
\pgfpathmoveto{\pgfqpoint{2.643208in}{5.392659in}}%
\pgfusepath{stroke}%
\end{pgfscope}%
\begin{pgfscope}%
\pgfpathrectangle{\pgfqpoint{0.100000in}{2.493063in}}{\pgfqpoint{5.037500in}{3.427208in}}%
\pgfusepath{clip}%
\pgfsetbuttcap%
\pgfsetroundjoin%
\definecolor{currentfill}{rgb}{0.678431,1.000000,0.184314}%
\pgfsetfillcolor{currentfill}%
\pgfsetfillopacity{0.500000}%
\pgfsetlinewidth{0.250937pt}%
\definecolor{currentstroke}{rgb}{0.000000,0.000000,0.000000}%
\pgfsetstrokecolor{currentstroke}%
\pgfsetstrokeopacity{0.500000}%
\pgfsetdash{}{0pt}%
\pgfsys@defobject{currentmarker}{\pgfqpoint{-0.013889in}{-0.013889in}}{\pgfqpoint{0.013889in}{0.013889in}}{%
\pgfpathmoveto{\pgfqpoint{0.000000in}{-0.013889in}}%
\pgfpathcurveto{\pgfqpoint{0.003683in}{-0.013889in}}{\pgfqpoint{0.007216in}{-0.012425in}}{\pgfqpoint{0.009821in}{-0.009821in}}%
\pgfpathcurveto{\pgfqpoint{0.012425in}{-0.007216in}}{\pgfqpoint{0.013889in}{-0.003683in}}{\pgfqpoint{0.013889in}{0.000000in}}%
\pgfpathcurveto{\pgfqpoint{0.013889in}{0.003683in}}{\pgfqpoint{0.012425in}{0.007216in}}{\pgfqpoint{0.009821in}{0.009821in}}%
\pgfpathcurveto{\pgfqpoint{0.007216in}{0.012425in}}{\pgfqpoint{0.003683in}{0.013889in}}{\pgfqpoint{0.000000in}{0.013889in}}%
\pgfpathcurveto{\pgfqpoint{-0.003683in}{0.013889in}}{\pgfqpoint{-0.007216in}{0.012425in}}{\pgfqpoint{-0.009821in}{0.009821in}}%
\pgfpathcurveto{\pgfqpoint{-0.012425in}{0.007216in}}{\pgfqpoint{-0.013889in}{0.003683in}}{\pgfqpoint{-0.013889in}{0.000000in}}%
\pgfpathcurveto{\pgfqpoint{-0.013889in}{-0.003683in}}{\pgfqpoint{-0.012425in}{-0.007216in}}{\pgfqpoint{-0.009821in}{-0.009821in}}%
\pgfpathcurveto{\pgfqpoint{-0.007216in}{-0.012425in}}{\pgfqpoint{-0.003683in}{-0.013889in}}{\pgfqpoint{0.000000in}{-0.013889in}}%
\pgfpathlineto{\pgfqpoint{0.000000in}{-0.013889in}}%
\pgfpathclose%
\pgfusepath{stroke,fill}%
}%
\begin{pgfscope}%
\pgfsys@transformshift{2.643208in}{5.392659in}%
\pgfsys@useobject{currentmarker}{}%
\end{pgfscope}%
\end{pgfscope}%
\begin{pgfscope}%
\pgfpathrectangle{\pgfqpoint{0.100000in}{2.493063in}}{\pgfqpoint{5.037500in}{3.427208in}}%
\pgfusepath{clip}%
\pgfsetrectcap%
\pgfsetroundjoin%
\pgfsetlinewidth{1.505625pt}%
\definecolor{currentstroke}{rgb}{0.678431,1.000000,0.184314}%
\pgfsetstrokecolor{currentstroke}%
\pgfsetstrokeopacity{0.500000}%
\pgfsetdash{}{0pt}%
\pgfpathmoveto{\pgfqpoint{3.970658in}{4.687541in}}%
\pgfusepath{stroke}%
\end{pgfscope}%
\begin{pgfscope}%
\pgfpathrectangle{\pgfqpoint{0.100000in}{2.493063in}}{\pgfqpoint{5.037500in}{3.427208in}}%
\pgfusepath{clip}%
\pgfsetbuttcap%
\pgfsetroundjoin%
\definecolor{currentfill}{rgb}{0.678431,1.000000,0.184314}%
\pgfsetfillcolor{currentfill}%
\pgfsetfillopacity{0.500000}%
\pgfsetlinewidth{0.250937pt}%
\definecolor{currentstroke}{rgb}{0.000000,0.000000,0.000000}%
\pgfsetstrokecolor{currentstroke}%
\pgfsetstrokeopacity{0.500000}%
\pgfsetdash{}{0pt}%
\pgfsys@defobject{currentmarker}{\pgfqpoint{-0.019444in}{-0.019444in}}{\pgfqpoint{0.019444in}{0.019444in}}{%
\pgfpathmoveto{\pgfqpoint{0.000000in}{-0.019444in}}%
\pgfpathcurveto{\pgfqpoint{0.005157in}{-0.019444in}}{\pgfqpoint{0.010103in}{-0.017396in}}{\pgfqpoint{0.013749in}{-0.013749in}}%
\pgfpathcurveto{\pgfqpoint{0.017396in}{-0.010103in}}{\pgfqpoint{0.019444in}{-0.005157in}}{\pgfqpoint{0.019444in}{0.000000in}}%
\pgfpathcurveto{\pgfqpoint{0.019444in}{0.005157in}}{\pgfqpoint{0.017396in}{0.010103in}}{\pgfqpoint{0.013749in}{0.013749in}}%
\pgfpathcurveto{\pgfqpoint{0.010103in}{0.017396in}}{\pgfqpoint{0.005157in}{0.019444in}}{\pgfqpoint{0.000000in}{0.019444in}}%
\pgfpathcurveto{\pgfqpoint{-0.005157in}{0.019444in}}{\pgfqpoint{-0.010103in}{0.017396in}}{\pgfqpoint{-0.013749in}{0.013749in}}%
\pgfpathcurveto{\pgfqpoint{-0.017396in}{0.010103in}}{\pgfqpoint{-0.019444in}{0.005157in}}{\pgfqpoint{-0.019444in}{0.000000in}}%
\pgfpathcurveto{\pgfqpoint{-0.019444in}{-0.005157in}}{\pgfqpoint{-0.017396in}{-0.010103in}}{\pgfqpoint{-0.013749in}{-0.013749in}}%
\pgfpathcurveto{\pgfqpoint{-0.010103in}{-0.017396in}}{\pgfqpoint{-0.005157in}{-0.019444in}}{\pgfqpoint{0.000000in}{-0.019444in}}%
\pgfpathlineto{\pgfqpoint{0.000000in}{-0.019444in}}%
\pgfpathclose%
\pgfusepath{stroke,fill}%
}%
\begin{pgfscope}%
\pgfsys@transformshift{3.970658in}{4.687541in}%
\pgfsys@useobject{currentmarker}{}%
\end{pgfscope}%
\end{pgfscope}%
\begin{pgfscope}%
\pgfpathrectangle{\pgfqpoint{0.100000in}{2.493063in}}{\pgfqpoint{5.037500in}{3.427208in}}%
\pgfusepath{clip}%
\pgfsetrectcap%
\pgfsetroundjoin%
\pgfsetlinewidth{1.505625pt}%
\definecolor{currentstroke}{rgb}{0.678431,1.000000,0.184314}%
\pgfsetstrokecolor{currentstroke}%
\pgfsetstrokeopacity{0.500000}%
\pgfsetdash{}{0pt}%
\pgfpathmoveto{\pgfqpoint{3.987573in}{4.656983in}}%
\pgfusepath{stroke}%
\end{pgfscope}%
\begin{pgfscope}%
\pgfpathrectangle{\pgfqpoint{0.100000in}{2.493063in}}{\pgfqpoint{5.037500in}{3.427208in}}%
\pgfusepath{clip}%
\pgfsetbuttcap%
\pgfsetroundjoin%
\definecolor{currentfill}{rgb}{0.678431,1.000000,0.184314}%
\pgfsetfillcolor{currentfill}%
\pgfsetfillopacity{0.500000}%
\pgfsetlinewidth{0.250937pt}%
\definecolor{currentstroke}{rgb}{0.000000,0.000000,0.000000}%
\pgfsetstrokecolor{currentstroke}%
\pgfsetstrokeopacity{0.500000}%
\pgfsetdash{}{0pt}%
\pgfsys@defobject{currentmarker}{\pgfqpoint{-0.030556in}{-0.030556in}}{\pgfqpoint{0.030556in}{0.030556in}}{%
\pgfpathmoveto{\pgfqpoint{0.000000in}{-0.030556in}}%
\pgfpathcurveto{\pgfqpoint{0.008103in}{-0.030556in}}{\pgfqpoint{0.015876in}{-0.027336in}}{\pgfqpoint{0.021606in}{-0.021606in}}%
\pgfpathcurveto{\pgfqpoint{0.027336in}{-0.015876in}}{\pgfqpoint{0.030556in}{-0.008103in}}{\pgfqpoint{0.030556in}{0.000000in}}%
\pgfpathcurveto{\pgfqpoint{0.030556in}{0.008103in}}{\pgfqpoint{0.027336in}{0.015876in}}{\pgfqpoint{0.021606in}{0.021606in}}%
\pgfpathcurveto{\pgfqpoint{0.015876in}{0.027336in}}{\pgfqpoint{0.008103in}{0.030556in}}{\pgfqpoint{0.000000in}{0.030556in}}%
\pgfpathcurveto{\pgfqpoint{-0.008103in}{0.030556in}}{\pgfqpoint{-0.015876in}{0.027336in}}{\pgfqpoint{-0.021606in}{0.021606in}}%
\pgfpathcurveto{\pgfqpoint{-0.027336in}{0.015876in}}{\pgfqpoint{-0.030556in}{0.008103in}}{\pgfqpoint{-0.030556in}{0.000000in}}%
\pgfpathcurveto{\pgfqpoint{-0.030556in}{-0.008103in}}{\pgfqpoint{-0.027336in}{-0.015876in}}{\pgfqpoint{-0.021606in}{-0.021606in}}%
\pgfpathcurveto{\pgfqpoint{-0.015876in}{-0.027336in}}{\pgfqpoint{-0.008103in}{-0.030556in}}{\pgfqpoint{0.000000in}{-0.030556in}}%
\pgfpathlineto{\pgfqpoint{0.000000in}{-0.030556in}}%
\pgfpathclose%
\pgfusepath{stroke,fill}%
}%
\begin{pgfscope}%
\pgfsys@transformshift{3.987573in}{4.656983in}%
\pgfsys@useobject{currentmarker}{}%
\end{pgfscope}%
\end{pgfscope}%
\begin{pgfscope}%
\pgfpathrectangle{\pgfqpoint{0.100000in}{2.493063in}}{\pgfqpoint{5.037500in}{3.427208in}}%
\pgfusepath{clip}%
\pgfsetrectcap%
\pgfsetroundjoin%
\pgfsetlinewidth{1.505625pt}%
\definecolor{currentstroke}{rgb}{0.678431,1.000000,0.184314}%
\pgfsetstrokecolor{currentstroke}%
\pgfsetstrokeopacity{0.500000}%
\pgfsetdash{}{0pt}%
\pgfpathmoveto{\pgfqpoint{3.739192in}{4.426646in}}%
\pgfusepath{stroke}%
\end{pgfscope}%
\begin{pgfscope}%
\pgfpathrectangle{\pgfqpoint{0.100000in}{2.493063in}}{\pgfqpoint{5.037500in}{3.427208in}}%
\pgfusepath{clip}%
\pgfsetbuttcap%
\pgfsetroundjoin%
\definecolor{currentfill}{rgb}{0.678431,1.000000,0.184314}%
\pgfsetfillcolor{currentfill}%
\pgfsetfillopacity{0.500000}%
\pgfsetlinewidth{0.250937pt}%
\definecolor{currentstroke}{rgb}{0.000000,0.000000,0.000000}%
\pgfsetstrokecolor{currentstroke}%
\pgfsetstrokeopacity{0.500000}%
\pgfsetdash{}{0pt}%
\pgfsys@defobject{currentmarker}{\pgfqpoint{-0.008333in}{-0.008333in}}{\pgfqpoint{0.008333in}{0.008333in}}{%
\pgfpathmoveto{\pgfqpoint{0.000000in}{-0.008333in}}%
\pgfpathcurveto{\pgfqpoint{0.002210in}{-0.008333in}}{\pgfqpoint{0.004330in}{-0.007455in}}{\pgfqpoint{0.005893in}{-0.005893in}}%
\pgfpathcurveto{\pgfqpoint{0.007455in}{-0.004330in}}{\pgfqpoint{0.008333in}{-0.002210in}}{\pgfqpoint{0.008333in}{0.000000in}}%
\pgfpathcurveto{\pgfqpoint{0.008333in}{0.002210in}}{\pgfqpoint{0.007455in}{0.004330in}}{\pgfqpoint{0.005893in}{0.005893in}}%
\pgfpathcurveto{\pgfqpoint{0.004330in}{0.007455in}}{\pgfqpoint{0.002210in}{0.008333in}}{\pgfqpoint{0.000000in}{0.008333in}}%
\pgfpathcurveto{\pgfqpoint{-0.002210in}{0.008333in}}{\pgfqpoint{-0.004330in}{0.007455in}}{\pgfqpoint{-0.005893in}{0.005893in}}%
\pgfpathcurveto{\pgfqpoint{-0.007455in}{0.004330in}}{\pgfqpoint{-0.008333in}{0.002210in}}{\pgfqpoint{-0.008333in}{0.000000in}}%
\pgfpathcurveto{\pgfqpoint{-0.008333in}{-0.002210in}}{\pgfqpoint{-0.007455in}{-0.004330in}}{\pgfqpoint{-0.005893in}{-0.005893in}}%
\pgfpathcurveto{\pgfqpoint{-0.004330in}{-0.007455in}}{\pgfqpoint{-0.002210in}{-0.008333in}}{\pgfqpoint{0.000000in}{-0.008333in}}%
\pgfpathlineto{\pgfqpoint{0.000000in}{-0.008333in}}%
\pgfpathclose%
\pgfusepath{stroke,fill}%
}%
\begin{pgfscope}%
\pgfsys@transformshift{3.739192in}{4.426646in}%
\pgfsys@useobject{currentmarker}{}%
\end{pgfscope}%
\end{pgfscope}%
\begin{pgfscope}%
\pgfpathrectangle{\pgfqpoint{0.100000in}{2.493063in}}{\pgfqpoint{5.037500in}{3.427208in}}%
\pgfusepath{clip}%
\pgfsetrectcap%
\pgfsetroundjoin%
\pgfsetlinewidth{1.505625pt}%
\definecolor{currentstroke}{rgb}{0.678431,1.000000,0.184314}%
\pgfsetstrokecolor{currentstroke}%
\pgfsetstrokeopacity{0.500000}%
\pgfsetdash{}{0pt}%
\pgfpathmoveto{\pgfqpoint{3.948532in}{4.733429in}}%
\pgfusepath{stroke}%
\end{pgfscope}%
\begin{pgfscope}%
\pgfpathrectangle{\pgfqpoint{0.100000in}{2.493063in}}{\pgfqpoint{5.037500in}{3.427208in}}%
\pgfusepath{clip}%
\pgfsetbuttcap%
\pgfsetroundjoin%
\definecolor{currentfill}{rgb}{0.678431,1.000000,0.184314}%
\pgfsetfillcolor{currentfill}%
\pgfsetfillopacity{0.500000}%
\pgfsetlinewidth{0.250937pt}%
\definecolor{currentstroke}{rgb}{0.000000,0.000000,0.000000}%
\pgfsetstrokecolor{currentstroke}%
\pgfsetstrokeopacity{0.500000}%
\pgfsetdash{}{0pt}%
\pgfsys@defobject{currentmarker}{\pgfqpoint{-0.019444in}{-0.019444in}}{\pgfqpoint{0.019444in}{0.019444in}}{%
\pgfpathmoveto{\pgfqpoint{0.000000in}{-0.019444in}}%
\pgfpathcurveto{\pgfqpoint{0.005157in}{-0.019444in}}{\pgfqpoint{0.010103in}{-0.017396in}}{\pgfqpoint{0.013749in}{-0.013749in}}%
\pgfpathcurveto{\pgfqpoint{0.017396in}{-0.010103in}}{\pgfqpoint{0.019444in}{-0.005157in}}{\pgfqpoint{0.019444in}{0.000000in}}%
\pgfpathcurveto{\pgfqpoint{0.019444in}{0.005157in}}{\pgfqpoint{0.017396in}{0.010103in}}{\pgfqpoint{0.013749in}{0.013749in}}%
\pgfpathcurveto{\pgfqpoint{0.010103in}{0.017396in}}{\pgfqpoint{0.005157in}{0.019444in}}{\pgfqpoint{0.000000in}{0.019444in}}%
\pgfpathcurveto{\pgfqpoint{-0.005157in}{0.019444in}}{\pgfqpoint{-0.010103in}{0.017396in}}{\pgfqpoint{-0.013749in}{0.013749in}}%
\pgfpathcurveto{\pgfqpoint{-0.017396in}{0.010103in}}{\pgfqpoint{-0.019444in}{0.005157in}}{\pgfqpoint{-0.019444in}{0.000000in}}%
\pgfpathcurveto{\pgfqpoint{-0.019444in}{-0.005157in}}{\pgfqpoint{-0.017396in}{-0.010103in}}{\pgfqpoint{-0.013749in}{-0.013749in}}%
\pgfpathcurveto{\pgfqpoint{-0.010103in}{-0.017396in}}{\pgfqpoint{-0.005157in}{-0.019444in}}{\pgfqpoint{0.000000in}{-0.019444in}}%
\pgfpathlineto{\pgfqpoint{0.000000in}{-0.019444in}}%
\pgfpathclose%
\pgfusepath{stroke,fill}%
}%
\begin{pgfscope}%
\pgfsys@transformshift{3.948532in}{4.733429in}%
\pgfsys@useobject{currentmarker}{}%
\end{pgfscope}%
\end{pgfscope}%
\begin{pgfscope}%
\pgfpathrectangle{\pgfqpoint{0.100000in}{2.493063in}}{\pgfqpoint{5.037500in}{3.427208in}}%
\pgfusepath{clip}%
\pgfsetrectcap%
\pgfsetroundjoin%
\pgfsetlinewidth{1.505625pt}%
\definecolor{currentstroke}{rgb}{0.678431,1.000000,0.184314}%
\pgfsetstrokecolor{currentstroke}%
\pgfsetstrokeopacity{0.500000}%
\pgfsetdash{}{0pt}%
\pgfpathmoveto{\pgfqpoint{3.860194in}{4.541577in}}%
\pgfusepath{stroke}%
\end{pgfscope}%
\begin{pgfscope}%
\pgfpathrectangle{\pgfqpoint{0.100000in}{2.493063in}}{\pgfqpoint{5.037500in}{3.427208in}}%
\pgfusepath{clip}%
\pgfsetbuttcap%
\pgfsetroundjoin%
\definecolor{currentfill}{rgb}{0.678431,1.000000,0.184314}%
\pgfsetfillcolor{currentfill}%
\pgfsetfillopacity{0.500000}%
\pgfsetlinewidth{0.250937pt}%
\definecolor{currentstroke}{rgb}{0.000000,0.000000,0.000000}%
\pgfsetstrokecolor{currentstroke}%
\pgfsetstrokeopacity{0.500000}%
\pgfsetdash{}{0pt}%
\pgfsys@defobject{currentmarker}{\pgfqpoint{-0.013889in}{-0.013889in}}{\pgfqpoint{0.013889in}{0.013889in}}{%
\pgfpathmoveto{\pgfqpoint{0.000000in}{-0.013889in}}%
\pgfpathcurveto{\pgfqpoint{0.003683in}{-0.013889in}}{\pgfqpoint{0.007216in}{-0.012425in}}{\pgfqpoint{0.009821in}{-0.009821in}}%
\pgfpathcurveto{\pgfqpoint{0.012425in}{-0.007216in}}{\pgfqpoint{0.013889in}{-0.003683in}}{\pgfqpoint{0.013889in}{0.000000in}}%
\pgfpathcurveto{\pgfqpoint{0.013889in}{0.003683in}}{\pgfqpoint{0.012425in}{0.007216in}}{\pgfqpoint{0.009821in}{0.009821in}}%
\pgfpathcurveto{\pgfqpoint{0.007216in}{0.012425in}}{\pgfqpoint{0.003683in}{0.013889in}}{\pgfqpoint{0.000000in}{0.013889in}}%
\pgfpathcurveto{\pgfqpoint{-0.003683in}{0.013889in}}{\pgfqpoint{-0.007216in}{0.012425in}}{\pgfqpoint{-0.009821in}{0.009821in}}%
\pgfpathcurveto{\pgfqpoint{-0.012425in}{0.007216in}}{\pgfqpoint{-0.013889in}{0.003683in}}{\pgfqpoint{-0.013889in}{0.000000in}}%
\pgfpathcurveto{\pgfqpoint{-0.013889in}{-0.003683in}}{\pgfqpoint{-0.012425in}{-0.007216in}}{\pgfqpoint{-0.009821in}{-0.009821in}}%
\pgfpathcurveto{\pgfqpoint{-0.007216in}{-0.012425in}}{\pgfqpoint{-0.003683in}{-0.013889in}}{\pgfqpoint{0.000000in}{-0.013889in}}%
\pgfpathlineto{\pgfqpoint{0.000000in}{-0.013889in}}%
\pgfpathclose%
\pgfusepath{stroke,fill}%
}%
\begin{pgfscope}%
\pgfsys@transformshift{3.860194in}{4.541577in}%
\pgfsys@useobject{currentmarker}{}%
\end{pgfscope}%
\end{pgfscope}%
\begin{pgfscope}%
\pgfpathrectangle{\pgfqpoint{0.100000in}{2.493063in}}{\pgfqpoint{5.037500in}{3.427208in}}%
\pgfusepath{clip}%
\pgfsetrectcap%
\pgfsetroundjoin%
\pgfsetlinewidth{1.505625pt}%
\definecolor{currentstroke}{rgb}{0.678431,1.000000,0.184314}%
\pgfsetstrokecolor{currentstroke}%
\pgfsetstrokeopacity{0.500000}%
\pgfsetdash{}{0pt}%
\pgfpathmoveto{\pgfqpoint{3.758682in}{4.505173in}}%
\pgfusepath{stroke}%
\end{pgfscope}%
\begin{pgfscope}%
\pgfpathrectangle{\pgfqpoint{0.100000in}{2.493063in}}{\pgfqpoint{5.037500in}{3.427208in}}%
\pgfusepath{clip}%
\pgfsetbuttcap%
\pgfsetroundjoin%
\definecolor{currentfill}{rgb}{0.678431,1.000000,0.184314}%
\pgfsetfillcolor{currentfill}%
\pgfsetfillopacity{0.500000}%
\pgfsetlinewidth{0.250937pt}%
\definecolor{currentstroke}{rgb}{0.000000,0.000000,0.000000}%
\pgfsetstrokecolor{currentstroke}%
\pgfsetstrokeopacity{0.500000}%
\pgfsetdash{}{0pt}%
\pgfsys@defobject{currentmarker}{\pgfqpoint{-0.013889in}{-0.013889in}}{\pgfqpoint{0.013889in}{0.013889in}}{%
\pgfpathmoveto{\pgfqpoint{0.000000in}{-0.013889in}}%
\pgfpathcurveto{\pgfqpoint{0.003683in}{-0.013889in}}{\pgfqpoint{0.007216in}{-0.012425in}}{\pgfqpoint{0.009821in}{-0.009821in}}%
\pgfpathcurveto{\pgfqpoint{0.012425in}{-0.007216in}}{\pgfqpoint{0.013889in}{-0.003683in}}{\pgfqpoint{0.013889in}{0.000000in}}%
\pgfpathcurveto{\pgfqpoint{0.013889in}{0.003683in}}{\pgfqpoint{0.012425in}{0.007216in}}{\pgfqpoint{0.009821in}{0.009821in}}%
\pgfpathcurveto{\pgfqpoint{0.007216in}{0.012425in}}{\pgfqpoint{0.003683in}{0.013889in}}{\pgfqpoint{0.000000in}{0.013889in}}%
\pgfpathcurveto{\pgfqpoint{-0.003683in}{0.013889in}}{\pgfqpoint{-0.007216in}{0.012425in}}{\pgfqpoint{-0.009821in}{0.009821in}}%
\pgfpathcurveto{\pgfqpoint{-0.012425in}{0.007216in}}{\pgfqpoint{-0.013889in}{0.003683in}}{\pgfqpoint{-0.013889in}{0.000000in}}%
\pgfpathcurveto{\pgfqpoint{-0.013889in}{-0.003683in}}{\pgfqpoint{-0.012425in}{-0.007216in}}{\pgfqpoint{-0.009821in}{-0.009821in}}%
\pgfpathcurveto{\pgfqpoint{-0.007216in}{-0.012425in}}{\pgfqpoint{-0.003683in}{-0.013889in}}{\pgfqpoint{0.000000in}{-0.013889in}}%
\pgfpathlineto{\pgfqpoint{0.000000in}{-0.013889in}}%
\pgfpathclose%
\pgfusepath{stroke,fill}%
}%
\begin{pgfscope}%
\pgfsys@transformshift{3.758682in}{4.505173in}%
\pgfsys@useobject{currentmarker}{}%
\end{pgfscope}%
\end{pgfscope}%
\begin{pgfscope}%
\pgfpathrectangle{\pgfqpoint{0.100000in}{2.493063in}}{\pgfqpoint{5.037500in}{3.427208in}}%
\pgfusepath{clip}%
\pgfsetrectcap%
\pgfsetroundjoin%
\pgfsetlinewidth{1.505625pt}%
\definecolor{currentstroke}{rgb}{0.678431,1.000000,0.184314}%
\pgfsetstrokecolor{currentstroke}%
\pgfsetstrokeopacity{0.500000}%
\pgfsetdash{}{0pt}%
\pgfpathmoveto{\pgfqpoint{3.752726in}{4.618555in}}%
\pgfusepath{stroke}%
\end{pgfscope}%
\begin{pgfscope}%
\pgfpathrectangle{\pgfqpoint{0.100000in}{2.493063in}}{\pgfqpoint{5.037500in}{3.427208in}}%
\pgfusepath{clip}%
\pgfsetbuttcap%
\pgfsetroundjoin%
\definecolor{currentfill}{rgb}{0.678431,1.000000,0.184314}%
\pgfsetfillcolor{currentfill}%
\pgfsetfillopacity{0.500000}%
\pgfsetlinewidth{0.250937pt}%
\definecolor{currentstroke}{rgb}{0.000000,0.000000,0.000000}%
\pgfsetstrokecolor{currentstroke}%
\pgfsetstrokeopacity{0.500000}%
\pgfsetdash{}{0pt}%
\pgfsys@defobject{currentmarker}{\pgfqpoint{-0.019444in}{-0.019444in}}{\pgfqpoint{0.019444in}{0.019444in}}{%
\pgfpathmoveto{\pgfqpoint{0.000000in}{-0.019444in}}%
\pgfpathcurveto{\pgfqpoint{0.005157in}{-0.019444in}}{\pgfqpoint{0.010103in}{-0.017396in}}{\pgfqpoint{0.013749in}{-0.013749in}}%
\pgfpathcurveto{\pgfqpoint{0.017396in}{-0.010103in}}{\pgfqpoint{0.019444in}{-0.005157in}}{\pgfqpoint{0.019444in}{0.000000in}}%
\pgfpathcurveto{\pgfqpoint{0.019444in}{0.005157in}}{\pgfqpoint{0.017396in}{0.010103in}}{\pgfqpoint{0.013749in}{0.013749in}}%
\pgfpathcurveto{\pgfqpoint{0.010103in}{0.017396in}}{\pgfqpoint{0.005157in}{0.019444in}}{\pgfqpoint{0.000000in}{0.019444in}}%
\pgfpathcurveto{\pgfqpoint{-0.005157in}{0.019444in}}{\pgfqpoint{-0.010103in}{0.017396in}}{\pgfqpoint{-0.013749in}{0.013749in}}%
\pgfpathcurveto{\pgfqpoint{-0.017396in}{0.010103in}}{\pgfqpoint{-0.019444in}{0.005157in}}{\pgfqpoint{-0.019444in}{0.000000in}}%
\pgfpathcurveto{\pgfqpoint{-0.019444in}{-0.005157in}}{\pgfqpoint{-0.017396in}{-0.010103in}}{\pgfqpoint{-0.013749in}{-0.013749in}}%
\pgfpathcurveto{\pgfqpoint{-0.010103in}{-0.017396in}}{\pgfqpoint{-0.005157in}{-0.019444in}}{\pgfqpoint{0.000000in}{-0.019444in}}%
\pgfpathlineto{\pgfqpoint{0.000000in}{-0.019444in}}%
\pgfpathclose%
\pgfusepath{stroke,fill}%
}%
\begin{pgfscope}%
\pgfsys@transformshift{3.752726in}{4.618555in}%
\pgfsys@useobject{currentmarker}{}%
\end{pgfscope}%
\end{pgfscope}%
\begin{pgfscope}%
\pgfpathrectangle{\pgfqpoint{0.100000in}{2.493063in}}{\pgfqpoint{5.037500in}{3.427208in}}%
\pgfusepath{clip}%
\pgfsetrectcap%
\pgfsetroundjoin%
\pgfsetlinewidth{1.505625pt}%
\definecolor{currentstroke}{rgb}{0.678431,1.000000,0.184314}%
\pgfsetstrokecolor{currentstroke}%
\pgfsetstrokeopacity{0.500000}%
\pgfsetdash{}{0pt}%
\pgfpathmoveto{\pgfqpoint{3.890081in}{4.638161in}}%
\pgfusepath{stroke}%
\end{pgfscope}%
\begin{pgfscope}%
\pgfpathrectangle{\pgfqpoint{0.100000in}{2.493063in}}{\pgfqpoint{5.037500in}{3.427208in}}%
\pgfusepath{clip}%
\pgfsetbuttcap%
\pgfsetroundjoin%
\definecolor{currentfill}{rgb}{0.678431,1.000000,0.184314}%
\pgfsetfillcolor{currentfill}%
\pgfsetfillopacity{0.500000}%
\pgfsetlinewidth{0.250937pt}%
\definecolor{currentstroke}{rgb}{0.000000,0.000000,0.000000}%
\pgfsetstrokecolor{currentstroke}%
\pgfsetstrokeopacity{0.500000}%
\pgfsetdash{}{0pt}%
\pgfsys@defobject{currentmarker}{\pgfqpoint{-0.025000in}{-0.025000in}}{\pgfqpoint{0.025000in}{0.025000in}}{%
\pgfpathmoveto{\pgfqpoint{0.000000in}{-0.025000in}}%
\pgfpathcurveto{\pgfqpoint{0.006630in}{-0.025000in}}{\pgfqpoint{0.012989in}{-0.022366in}}{\pgfqpoint{0.017678in}{-0.017678in}}%
\pgfpathcurveto{\pgfqpoint{0.022366in}{-0.012989in}}{\pgfqpoint{0.025000in}{-0.006630in}}{\pgfqpoint{0.025000in}{0.000000in}}%
\pgfpathcurveto{\pgfqpoint{0.025000in}{0.006630in}}{\pgfqpoint{0.022366in}{0.012989in}}{\pgfqpoint{0.017678in}{0.017678in}}%
\pgfpathcurveto{\pgfqpoint{0.012989in}{0.022366in}}{\pgfqpoint{0.006630in}{0.025000in}}{\pgfqpoint{0.000000in}{0.025000in}}%
\pgfpathcurveto{\pgfqpoint{-0.006630in}{0.025000in}}{\pgfqpoint{-0.012989in}{0.022366in}}{\pgfqpoint{-0.017678in}{0.017678in}}%
\pgfpathcurveto{\pgfqpoint{-0.022366in}{0.012989in}}{\pgfqpoint{-0.025000in}{0.006630in}}{\pgfqpoint{-0.025000in}{0.000000in}}%
\pgfpathcurveto{\pgfqpoint{-0.025000in}{-0.006630in}}{\pgfqpoint{-0.022366in}{-0.012989in}}{\pgfqpoint{-0.017678in}{-0.017678in}}%
\pgfpathcurveto{\pgfqpoint{-0.012989in}{-0.022366in}}{\pgfqpoint{-0.006630in}{-0.025000in}}{\pgfqpoint{0.000000in}{-0.025000in}}%
\pgfpathlineto{\pgfqpoint{0.000000in}{-0.025000in}}%
\pgfpathclose%
\pgfusepath{stroke,fill}%
}%
\begin{pgfscope}%
\pgfsys@transformshift{3.890081in}{4.638161in}%
\pgfsys@useobject{currentmarker}{}%
\end{pgfscope}%
\end{pgfscope}%
\begin{pgfscope}%
\pgfpathrectangle{\pgfqpoint{0.100000in}{2.493063in}}{\pgfqpoint{5.037500in}{3.427208in}}%
\pgfusepath{clip}%
\pgfsetrectcap%
\pgfsetroundjoin%
\pgfsetlinewidth{1.505625pt}%
\definecolor{currentstroke}{rgb}{0.678431,1.000000,0.184314}%
\pgfsetstrokecolor{currentstroke}%
\pgfsetstrokeopacity{0.500000}%
\pgfsetdash{}{0pt}%
\pgfpathmoveto{\pgfqpoint{3.789983in}{4.528154in}}%
\pgfusepath{stroke}%
\end{pgfscope}%
\begin{pgfscope}%
\pgfpathrectangle{\pgfqpoint{0.100000in}{2.493063in}}{\pgfqpoint{5.037500in}{3.427208in}}%
\pgfusepath{clip}%
\pgfsetbuttcap%
\pgfsetroundjoin%
\definecolor{currentfill}{rgb}{0.678431,1.000000,0.184314}%
\pgfsetfillcolor{currentfill}%
\pgfsetfillopacity{0.500000}%
\pgfsetlinewidth{0.250937pt}%
\definecolor{currentstroke}{rgb}{0.000000,0.000000,0.000000}%
\pgfsetstrokecolor{currentstroke}%
\pgfsetstrokeopacity{0.500000}%
\pgfsetdash{}{0pt}%
\pgfsys@defobject{currentmarker}{\pgfqpoint{-0.027778in}{-0.027778in}}{\pgfqpoint{0.027778in}{0.027778in}}{%
\pgfpathmoveto{\pgfqpoint{0.000000in}{-0.027778in}}%
\pgfpathcurveto{\pgfqpoint{0.007367in}{-0.027778in}}{\pgfqpoint{0.014433in}{-0.024851in}}{\pgfqpoint{0.019642in}{-0.019642in}}%
\pgfpathcurveto{\pgfqpoint{0.024851in}{-0.014433in}}{\pgfqpoint{0.027778in}{-0.007367in}}{\pgfqpoint{0.027778in}{0.000000in}}%
\pgfpathcurveto{\pgfqpoint{0.027778in}{0.007367in}}{\pgfqpoint{0.024851in}{0.014433in}}{\pgfqpoint{0.019642in}{0.019642in}}%
\pgfpathcurveto{\pgfqpoint{0.014433in}{0.024851in}}{\pgfqpoint{0.007367in}{0.027778in}}{\pgfqpoint{0.000000in}{0.027778in}}%
\pgfpathcurveto{\pgfqpoint{-0.007367in}{0.027778in}}{\pgfqpoint{-0.014433in}{0.024851in}}{\pgfqpoint{-0.019642in}{0.019642in}}%
\pgfpathcurveto{\pgfqpoint{-0.024851in}{0.014433in}}{\pgfqpoint{-0.027778in}{0.007367in}}{\pgfqpoint{-0.027778in}{0.000000in}}%
\pgfpathcurveto{\pgfqpoint{-0.027778in}{-0.007367in}}{\pgfqpoint{-0.024851in}{-0.014433in}}{\pgfqpoint{-0.019642in}{-0.019642in}}%
\pgfpathcurveto{\pgfqpoint{-0.014433in}{-0.024851in}}{\pgfqpoint{-0.007367in}{-0.027778in}}{\pgfqpoint{0.000000in}{-0.027778in}}%
\pgfpathlineto{\pgfqpoint{0.000000in}{-0.027778in}}%
\pgfpathclose%
\pgfusepath{stroke,fill}%
}%
\begin{pgfscope}%
\pgfsys@transformshift{3.789983in}{4.528154in}%
\pgfsys@useobject{currentmarker}{}%
\end{pgfscope}%
\end{pgfscope}%
\begin{pgfscope}%
\pgfpathrectangle{\pgfqpoint{0.100000in}{2.493063in}}{\pgfqpoint{5.037500in}{3.427208in}}%
\pgfusepath{clip}%
\pgfsetrectcap%
\pgfsetroundjoin%
\pgfsetlinewidth{1.505625pt}%
\definecolor{currentstroke}{rgb}{0.678431,1.000000,0.184314}%
\pgfsetstrokecolor{currentstroke}%
\pgfsetstrokeopacity{0.500000}%
\pgfsetdash{}{0pt}%
\pgfpathmoveto{\pgfqpoint{3.790448in}{4.731915in}}%
\pgfusepath{stroke}%
\end{pgfscope}%
\begin{pgfscope}%
\pgfpathrectangle{\pgfqpoint{0.100000in}{2.493063in}}{\pgfqpoint{5.037500in}{3.427208in}}%
\pgfusepath{clip}%
\pgfsetbuttcap%
\pgfsetroundjoin%
\definecolor{currentfill}{rgb}{0.678431,1.000000,0.184314}%
\pgfsetfillcolor{currentfill}%
\pgfsetfillopacity{0.500000}%
\pgfsetlinewidth{0.250937pt}%
\definecolor{currentstroke}{rgb}{0.000000,0.000000,0.000000}%
\pgfsetstrokecolor{currentstroke}%
\pgfsetstrokeopacity{0.500000}%
\pgfsetdash{}{0pt}%
\pgfsys@defobject{currentmarker}{\pgfqpoint{-0.008333in}{-0.008333in}}{\pgfqpoint{0.008333in}{0.008333in}}{%
\pgfpathmoveto{\pgfqpoint{0.000000in}{-0.008333in}}%
\pgfpathcurveto{\pgfqpoint{0.002210in}{-0.008333in}}{\pgfqpoint{0.004330in}{-0.007455in}}{\pgfqpoint{0.005893in}{-0.005893in}}%
\pgfpathcurveto{\pgfqpoint{0.007455in}{-0.004330in}}{\pgfqpoint{0.008333in}{-0.002210in}}{\pgfqpoint{0.008333in}{0.000000in}}%
\pgfpathcurveto{\pgfqpoint{0.008333in}{0.002210in}}{\pgfqpoint{0.007455in}{0.004330in}}{\pgfqpoint{0.005893in}{0.005893in}}%
\pgfpathcurveto{\pgfqpoint{0.004330in}{0.007455in}}{\pgfqpoint{0.002210in}{0.008333in}}{\pgfqpoint{0.000000in}{0.008333in}}%
\pgfpathcurveto{\pgfqpoint{-0.002210in}{0.008333in}}{\pgfqpoint{-0.004330in}{0.007455in}}{\pgfqpoint{-0.005893in}{0.005893in}}%
\pgfpathcurveto{\pgfqpoint{-0.007455in}{0.004330in}}{\pgfqpoint{-0.008333in}{0.002210in}}{\pgfqpoint{-0.008333in}{0.000000in}}%
\pgfpathcurveto{\pgfqpoint{-0.008333in}{-0.002210in}}{\pgfqpoint{-0.007455in}{-0.004330in}}{\pgfqpoint{-0.005893in}{-0.005893in}}%
\pgfpathcurveto{\pgfqpoint{-0.004330in}{-0.007455in}}{\pgfqpoint{-0.002210in}{-0.008333in}}{\pgfqpoint{0.000000in}{-0.008333in}}%
\pgfpathlineto{\pgfqpoint{0.000000in}{-0.008333in}}%
\pgfpathclose%
\pgfusepath{stroke,fill}%
}%
\begin{pgfscope}%
\pgfsys@transformshift{3.790448in}{4.731915in}%
\pgfsys@useobject{currentmarker}{}%
\end{pgfscope}%
\end{pgfscope}%
\begin{pgfscope}%
\pgfpathrectangle{\pgfqpoint{0.100000in}{2.493063in}}{\pgfqpoint{5.037500in}{3.427208in}}%
\pgfusepath{clip}%
\pgfsetrectcap%
\pgfsetroundjoin%
\pgfsetlinewidth{1.505625pt}%
\definecolor{currentstroke}{rgb}{0.678431,1.000000,0.184314}%
\pgfsetstrokecolor{currentstroke}%
\pgfsetstrokeopacity{0.500000}%
\pgfsetdash{}{0pt}%
\pgfpathmoveto{\pgfqpoint{4.062468in}{4.624330in}}%
\pgfusepath{stroke}%
\end{pgfscope}%
\begin{pgfscope}%
\pgfpathrectangle{\pgfqpoint{0.100000in}{2.493063in}}{\pgfqpoint{5.037500in}{3.427208in}}%
\pgfusepath{clip}%
\pgfsetbuttcap%
\pgfsetroundjoin%
\definecolor{currentfill}{rgb}{0.678431,1.000000,0.184314}%
\pgfsetfillcolor{currentfill}%
\pgfsetfillopacity{0.500000}%
\pgfsetlinewidth{0.250937pt}%
\definecolor{currentstroke}{rgb}{0.000000,0.000000,0.000000}%
\pgfsetstrokecolor{currentstroke}%
\pgfsetstrokeopacity{0.500000}%
\pgfsetdash{}{0pt}%
\pgfsys@defobject{currentmarker}{\pgfqpoint{-0.013889in}{-0.013889in}}{\pgfqpoint{0.013889in}{0.013889in}}{%
\pgfpathmoveto{\pgfqpoint{0.000000in}{-0.013889in}}%
\pgfpathcurveto{\pgfqpoint{0.003683in}{-0.013889in}}{\pgfqpoint{0.007216in}{-0.012425in}}{\pgfqpoint{0.009821in}{-0.009821in}}%
\pgfpathcurveto{\pgfqpoint{0.012425in}{-0.007216in}}{\pgfqpoint{0.013889in}{-0.003683in}}{\pgfqpoint{0.013889in}{0.000000in}}%
\pgfpathcurveto{\pgfqpoint{0.013889in}{0.003683in}}{\pgfqpoint{0.012425in}{0.007216in}}{\pgfqpoint{0.009821in}{0.009821in}}%
\pgfpathcurveto{\pgfqpoint{0.007216in}{0.012425in}}{\pgfqpoint{0.003683in}{0.013889in}}{\pgfqpoint{0.000000in}{0.013889in}}%
\pgfpathcurveto{\pgfqpoint{-0.003683in}{0.013889in}}{\pgfqpoint{-0.007216in}{0.012425in}}{\pgfqpoint{-0.009821in}{0.009821in}}%
\pgfpathcurveto{\pgfqpoint{-0.012425in}{0.007216in}}{\pgfqpoint{-0.013889in}{0.003683in}}{\pgfqpoint{-0.013889in}{0.000000in}}%
\pgfpathcurveto{\pgfqpoint{-0.013889in}{-0.003683in}}{\pgfqpoint{-0.012425in}{-0.007216in}}{\pgfqpoint{-0.009821in}{-0.009821in}}%
\pgfpathcurveto{\pgfqpoint{-0.007216in}{-0.012425in}}{\pgfqpoint{-0.003683in}{-0.013889in}}{\pgfqpoint{0.000000in}{-0.013889in}}%
\pgfpathlineto{\pgfqpoint{0.000000in}{-0.013889in}}%
\pgfpathclose%
\pgfusepath{stroke,fill}%
}%
\begin{pgfscope}%
\pgfsys@transformshift{4.062468in}{4.624330in}%
\pgfsys@useobject{currentmarker}{}%
\end{pgfscope}%
\end{pgfscope}%
\begin{pgfscope}%
\pgfpathrectangle{\pgfqpoint{0.100000in}{2.493063in}}{\pgfqpoint{5.037500in}{3.427208in}}%
\pgfusepath{clip}%
\pgfsetrectcap%
\pgfsetroundjoin%
\pgfsetlinewidth{1.505625pt}%
\definecolor{currentstroke}{rgb}{0.678431,1.000000,0.184314}%
\pgfsetstrokecolor{currentstroke}%
\pgfsetstrokeopacity{0.500000}%
\pgfsetdash{}{0pt}%
\pgfpathmoveto{\pgfqpoint{4.044939in}{4.700954in}}%
\pgfusepath{stroke}%
\end{pgfscope}%
\begin{pgfscope}%
\pgfpathrectangle{\pgfqpoint{0.100000in}{2.493063in}}{\pgfqpoint{5.037500in}{3.427208in}}%
\pgfusepath{clip}%
\pgfsetbuttcap%
\pgfsetroundjoin%
\definecolor{currentfill}{rgb}{0.678431,1.000000,0.184314}%
\pgfsetfillcolor{currentfill}%
\pgfsetfillopacity{0.500000}%
\pgfsetlinewidth{0.250937pt}%
\definecolor{currentstroke}{rgb}{0.000000,0.000000,0.000000}%
\pgfsetstrokecolor{currentstroke}%
\pgfsetstrokeopacity{0.500000}%
\pgfsetdash{}{0pt}%
\pgfsys@defobject{currentmarker}{\pgfqpoint{-0.044444in}{-0.044444in}}{\pgfqpoint{0.044444in}{0.044444in}}{%
\pgfpathmoveto{\pgfqpoint{0.000000in}{-0.044444in}}%
\pgfpathcurveto{\pgfqpoint{0.011787in}{-0.044444in}}{\pgfqpoint{0.023092in}{-0.039761in}}{\pgfqpoint{0.031427in}{-0.031427in}}%
\pgfpathcurveto{\pgfqpoint{0.039761in}{-0.023092in}}{\pgfqpoint{0.044444in}{-0.011787in}}{\pgfqpoint{0.044444in}{0.000000in}}%
\pgfpathcurveto{\pgfqpoint{0.044444in}{0.011787in}}{\pgfqpoint{0.039761in}{0.023092in}}{\pgfqpoint{0.031427in}{0.031427in}}%
\pgfpathcurveto{\pgfqpoint{0.023092in}{0.039761in}}{\pgfqpoint{0.011787in}{0.044444in}}{\pgfqpoint{0.000000in}{0.044444in}}%
\pgfpathcurveto{\pgfqpoint{-0.011787in}{0.044444in}}{\pgfqpoint{-0.023092in}{0.039761in}}{\pgfqpoint{-0.031427in}{0.031427in}}%
\pgfpathcurveto{\pgfqpoint{-0.039761in}{0.023092in}}{\pgfqpoint{-0.044444in}{0.011787in}}{\pgfqpoint{-0.044444in}{0.000000in}}%
\pgfpathcurveto{\pgfqpoint{-0.044444in}{-0.011787in}}{\pgfqpoint{-0.039761in}{-0.023092in}}{\pgfqpoint{-0.031427in}{-0.031427in}}%
\pgfpathcurveto{\pgfqpoint{-0.023092in}{-0.039761in}}{\pgfqpoint{-0.011787in}{-0.044444in}}{\pgfqpoint{0.000000in}{-0.044444in}}%
\pgfpathlineto{\pgfqpoint{0.000000in}{-0.044444in}}%
\pgfpathclose%
\pgfusepath{stroke,fill}%
}%
\begin{pgfscope}%
\pgfsys@transformshift{4.044939in}{4.700954in}%
\pgfsys@useobject{currentmarker}{}%
\end{pgfscope}%
\end{pgfscope}%
\begin{pgfscope}%
\pgfpathrectangle{\pgfqpoint{0.100000in}{2.493063in}}{\pgfqpoint{5.037500in}{3.427208in}}%
\pgfusepath{clip}%
\pgfsetrectcap%
\pgfsetroundjoin%
\pgfsetlinewidth{1.505625pt}%
\definecolor{currentstroke}{rgb}{0.000000,0.000000,1.000000}%
\pgfsetstrokecolor{currentstroke}%
\pgfsetstrokeopacity{0.500000}%
\pgfsetdash{}{0pt}%
\pgfpathmoveto{\pgfqpoint{2.481932in}{3.860796in}}%
\pgfusepath{stroke}%
\end{pgfscope}%
\begin{pgfscope}%
\pgfpathrectangle{\pgfqpoint{0.100000in}{2.493063in}}{\pgfqpoint{5.037500in}{3.427208in}}%
\pgfusepath{clip}%
\pgfsetbuttcap%
\pgfsetroundjoin%
\definecolor{currentfill}{rgb}{0.000000,0.000000,1.000000}%
\pgfsetfillcolor{currentfill}%
\pgfsetfillopacity{0.500000}%
\pgfsetlinewidth{0.250937pt}%
\definecolor{currentstroke}{rgb}{0.000000,0.000000,0.000000}%
\pgfsetstrokecolor{currentstroke}%
\pgfsetstrokeopacity{0.500000}%
\pgfsetdash{}{0pt}%
\pgfsys@defobject{currentmarker}{\pgfqpoint{-0.016667in}{-0.016667in}}{\pgfqpoint{0.016667in}{0.016667in}}{%
\pgfpathmoveto{\pgfqpoint{0.000000in}{-0.016667in}}%
\pgfpathcurveto{\pgfqpoint{0.004420in}{-0.016667in}}{\pgfqpoint{0.008660in}{-0.014911in}}{\pgfqpoint{0.011785in}{-0.011785in}}%
\pgfpathcurveto{\pgfqpoint{0.014911in}{-0.008660in}}{\pgfqpoint{0.016667in}{-0.004420in}}{\pgfqpoint{0.016667in}{0.000000in}}%
\pgfpathcurveto{\pgfqpoint{0.016667in}{0.004420in}}{\pgfqpoint{0.014911in}{0.008660in}}{\pgfqpoint{0.011785in}{0.011785in}}%
\pgfpathcurveto{\pgfqpoint{0.008660in}{0.014911in}}{\pgfqpoint{0.004420in}{0.016667in}}{\pgfqpoint{0.000000in}{0.016667in}}%
\pgfpathcurveto{\pgfqpoint{-0.004420in}{0.016667in}}{\pgfqpoint{-0.008660in}{0.014911in}}{\pgfqpoint{-0.011785in}{0.011785in}}%
\pgfpathcurveto{\pgfqpoint{-0.014911in}{0.008660in}}{\pgfqpoint{-0.016667in}{0.004420in}}{\pgfqpoint{-0.016667in}{0.000000in}}%
\pgfpathcurveto{\pgfqpoint{-0.016667in}{-0.004420in}}{\pgfqpoint{-0.014911in}{-0.008660in}}{\pgfqpoint{-0.011785in}{-0.011785in}}%
\pgfpathcurveto{\pgfqpoint{-0.008660in}{-0.014911in}}{\pgfqpoint{-0.004420in}{-0.016667in}}{\pgfqpoint{0.000000in}{-0.016667in}}%
\pgfpathlineto{\pgfqpoint{0.000000in}{-0.016667in}}%
\pgfpathclose%
\pgfusepath{stroke,fill}%
}%
\begin{pgfscope}%
\pgfsys@transformshift{2.481932in}{3.860796in}%
\pgfsys@useobject{currentmarker}{}%
\end{pgfscope}%
\end{pgfscope}%
\begin{pgfscope}%
\pgfpathrectangle{\pgfqpoint{0.100000in}{2.493063in}}{\pgfqpoint{5.037500in}{3.427208in}}%
\pgfusepath{clip}%
\pgfsetrectcap%
\pgfsetroundjoin%
\pgfsetlinewidth{1.505625pt}%
\definecolor{currentstroke}{rgb}{0.000000,0.000000,1.000000}%
\pgfsetstrokecolor{currentstroke}%
\pgfsetstrokeopacity{0.500000}%
\pgfsetdash{}{0pt}%
\pgfpathmoveto{\pgfqpoint{2.567661in}{3.957827in}}%
\pgfusepath{stroke}%
\end{pgfscope}%
\begin{pgfscope}%
\pgfpathrectangle{\pgfqpoint{0.100000in}{2.493063in}}{\pgfqpoint{5.037500in}{3.427208in}}%
\pgfusepath{clip}%
\pgfsetbuttcap%
\pgfsetroundjoin%
\definecolor{currentfill}{rgb}{0.000000,0.000000,1.000000}%
\pgfsetfillcolor{currentfill}%
\pgfsetfillopacity{0.500000}%
\pgfsetlinewidth{0.250937pt}%
\definecolor{currentstroke}{rgb}{0.000000,0.000000,0.000000}%
\pgfsetstrokecolor{currentstroke}%
\pgfsetstrokeopacity{0.500000}%
\pgfsetdash{}{0pt}%
\pgfsys@defobject{currentmarker}{\pgfqpoint{-0.011111in}{-0.011111in}}{\pgfqpoint{0.011111in}{0.011111in}}{%
\pgfpathmoveto{\pgfqpoint{0.000000in}{-0.011111in}}%
\pgfpathcurveto{\pgfqpoint{0.002947in}{-0.011111in}}{\pgfqpoint{0.005773in}{-0.009940in}}{\pgfqpoint{0.007857in}{-0.007857in}}%
\pgfpathcurveto{\pgfqpoint{0.009940in}{-0.005773in}}{\pgfqpoint{0.011111in}{-0.002947in}}{\pgfqpoint{0.011111in}{0.000000in}}%
\pgfpathcurveto{\pgfqpoint{0.011111in}{0.002947in}}{\pgfqpoint{0.009940in}{0.005773in}}{\pgfqpoint{0.007857in}{0.007857in}}%
\pgfpathcurveto{\pgfqpoint{0.005773in}{0.009940in}}{\pgfqpoint{0.002947in}{0.011111in}}{\pgfqpoint{0.000000in}{0.011111in}}%
\pgfpathcurveto{\pgfqpoint{-0.002947in}{0.011111in}}{\pgfqpoint{-0.005773in}{0.009940in}}{\pgfqpoint{-0.007857in}{0.007857in}}%
\pgfpathcurveto{\pgfqpoint{-0.009940in}{0.005773in}}{\pgfqpoint{-0.011111in}{0.002947in}}{\pgfqpoint{-0.011111in}{0.000000in}}%
\pgfpathcurveto{\pgfqpoint{-0.011111in}{-0.002947in}}{\pgfqpoint{-0.009940in}{-0.005773in}}{\pgfqpoint{-0.007857in}{-0.007857in}}%
\pgfpathcurveto{\pgfqpoint{-0.005773in}{-0.009940in}}{\pgfqpoint{-0.002947in}{-0.011111in}}{\pgfqpoint{0.000000in}{-0.011111in}}%
\pgfpathlineto{\pgfqpoint{0.000000in}{-0.011111in}}%
\pgfpathclose%
\pgfusepath{stroke,fill}%
}%
\begin{pgfscope}%
\pgfsys@transformshift{2.567661in}{3.957827in}%
\pgfsys@useobject{currentmarker}{}%
\end{pgfscope}%
\end{pgfscope}%
\begin{pgfscope}%
\pgfpathrectangle{\pgfqpoint{0.100000in}{2.493063in}}{\pgfqpoint{5.037500in}{3.427208in}}%
\pgfusepath{clip}%
\pgfsetrectcap%
\pgfsetroundjoin%
\pgfsetlinewidth{1.505625pt}%
\definecolor{currentstroke}{rgb}{0.000000,0.000000,1.000000}%
\pgfsetstrokecolor{currentstroke}%
\pgfsetstrokeopacity{0.500000}%
\pgfsetdash{}{0pt}%
\pgfpathmoveto{\pgfqpoint{2.711698in}{4.033793in}}%
\pgfusepath{stroke}%
\end{pgfscope}%
\begin{pgfscope}%
\pgfpathrectangle{\pgfqpoint{0.100000in}{2.493063in}}{\pgfqpoint{5.037500in}{3.427208in}}%
\pgfusepath{clip}%
\pgfsetbuttcap%
\pgfsetroundjoin%
\definecolor{currentfill}{rgb}{0.000000,0.000000,1.000000}%
\pgfsetfillcolor{currentfill}%
\pgfsetfillopacity{0.500000}%
\pgfsetlinewidth{0.250937pt}%
\definecolor{currentstroke}{rgb}{0.000000,0.000000,0.000000}%
\pgfsetstrokecolor{currentstroke}%
\pgfsetstrokeopacity{0.500000}%
\pgfsetdash{}{0pt}%
\pgfsys@defobject{currentmarker}{\pgfqpoint{-0.008333in}{-0.008333in}}{\pgfqpoint{0.008333in}{0.008333in}}{%
\pgfpathmoveto{\pgfqpoint{0.000000in}{-0.008333in}}%
\pgfpathcurveto{\pgfqpoint{0.002210in}{-0.008333in}}{\pgfqpoint{0.004330in}{-0.007455in}}{\pgfqpoint{0.005893in}{-0.005893in}}%
\pgfpathcurveto{\pgfqpoint{0.007455in}{-0.004330in}}{\pgfqpoint{0.008333in}{-0.002210in}}{\pgfqpoint{0.008333in}{0.000000in}}%
\pgfpathcurveto{\pgfqpoint{0.008333in}{0.002210in}}{\pgfqpoint{0.007455in}{0.004330in}}{\pgfqpoint{0.005893in}{0.005893in}}%
\pgfpathcurveto{\pgfqpoint{0.004330in}{0.007455in}}{\pgfqpoint{0.002210in}{0.008333in}}{\pgfqpoint{0.000000in}{0.008333in}}%
\pgfpathcurveto{\pgfqpoint{-0.002210in}{0.008333in}}{\pgfqpoint{-0.004330in}{0.007455in}}{\pgfqpoint{-0.005893in}{0.005893in}}%
\pgfpathcurveto{\pgfqpoint{-0.007455in}{0.004330in}}{\pgfqpoint{-0.008333in}{0.002210in}}{\pgfqpoint{-0.008333in}{0.000000in}}%
\pgfpathcurveto{\pgfqpoint{-0.008333in}{-0.002210in}}{\pgfqpoint{-0.007455in}{-0.004330in}}{\pgfqpoint{-0.005893in}{-0.005893in}}%
\pgfpathcurveto{\pgfqpoint{-0.004330in}{-0.007455in}}{\pgfqpoint{-0.002210in}{-0.008333in}}{\pgfqpoint{0.000000in}{-0.008333in}}%
\pgfpathlineto{\pgfqpoint{0.000000in}{-0.008333in}}%
\pgfpathclose%
\pgfusepath{stroke,fill}%
}%
\begin{pgfscope}%
\pgfsys@transformshift{2.711698in}{4.033793in}%
\pgfsys@useobject{currentmarker}{}%
\end{pgfscope}%
\end{pgfscope}%
\begin{pgfscope}%
\pgfpathrectangle{\pgfqpoint{0.100000in}{2.493063in}}{\pgfqpoint{5.037500in}{3.427208in}}%
\pgfusepath{clip}%
\pgfsetrectcap%
\pgfsetroundjoin%
\pgfsetlinewidth{1.505625pt}%
\definecolor{currentstroke}{rgb}{0.000000,0.000000,1.000000}%
\pgfsetstrokecolor{currentstroke}%
\pgfsetstrokeopacity{0.500000}%
\pgfsetdash{}{0pt}%
\pgfpathmoveto{\pgfqpoint{0.526793in}{5.365996in}}%
\pgfusepath{stroke}%
\end{pgfscope}%
\begin{pgfscope}%
\pgfpathrectangle{\pgfqpoint{0.100000in}{2.493063in}}{\pgfqpoint{5.037500in}{3.427208in}}%
\pgfusepath{clip}%
\pgfsetbuttcap%
\pgfsetroundjoin%
\definecolor{currentfill}{rgb}{0.000000,0.000000,1.000000}%
\pgfsetfillcolor{currentfill}%
\pgfsetfillopacity{0.500000}%
\pgfsetlinewidth{0.250937pt}%
\definecolor{currentstroke}{rgb}{0.000000,0.000000,0.000000}%
\pgfsetstrokecolor{currentstroke}%
\pgfsetstrokeopacity{0.500000}%
\pgfsetdash{}{0pt}%
\pgfsys@defobject{currentmarker}{\pgfqpoint{-0.011111in}{-0.011111in}}{\pgfqpoint{0.011111in}{0.011111in}}{%
\pgfpathmoveto{\pgfqpoint{0.000000in}{-0.011111in}}%
\pgfpathcurveto{\pgfqpoint{0.002947in}{-0.011111in}}{\pgfqpoint{0.005773in}{-0.009940in}}{\pgfqpoint{0.007857in}{-0.007857in}}%
\pgfpathcurveto{\pgfqpoint{0.009940in}{-0.005773in}}{\pgfqpoint{0.011111in}{-0.002947in}}{\pgfqpoint{0.011111in}{0.000000in}}%
\pgfpathcurveto{\pgfqpoint{0.011111in}{0.002947in}}{\pgfqpoint{0.009940in}{0.005773in}}{\pgfqpoint{0.007857in}{0.007857in}}%
\pgfpathcurveto{\pgfqpoint{0.005773in}{0.009940in}}{\pgfqpoint{0.002947in}{0.011111in}}{\pgfqpoint{0.000000in}{0.011111in}}%
\pgfpathcurveto{\pgfqpoint{-0.002947in}{0.011111in}}{\pgfqpoint{-0.005773in}{0.009940in}}{\pgfqpoint{-0.007857in}{0.007857in}}%
\pgfpathcurveto{\pgfqpoint{-0.009940in}{0.005773in}}{\pgfqpoint{-0.011111in}{0.002947in}}{\pgfqpoint{-0.011111in}{0.000000in}}%
\pgfpathcurveto{\pgfqpoint{-0.011111in}{-0.002947in}}{\pgfqpoint{-0.009940in}{-0.005773in}}{\pgfqpoint{-0.007857in}{-0.007857in}}%
\pgfpathcurveto{\pgfqpoint{-0.005773in}{-0.009940in}}{\pgfqpoint{-0.002947in}{-0.011111in}}{\pgfqpoint{0.000000in}{-0.011111in}}%
\pgfpathlineto{\pgfqpoint{0.000000in}{-0.011111in}}%
\pgfpathclose%
\pgfusepath{stroke,fill}%
}%
\begin{pgfscope}%
\pgfsys@transformshift{0.526793in}{5.365996in}%
\pgfsys@useobject{currentmarker}{}%
\end{pgfscope}%
\end{pgfscope}%
\begin{pgfscope}%
\pgfpathrectangle{\pgfqpoint{0.100000in}{2.493063in}}{\pgfqpoint{5.037500in}{3.427208in}}%
\pgfusepath{clip}%
\pgfsetrectcap%
\pgfsetroundjoin%
\pgfsetlinewidth{1.505625pt}%
\definecolor{currentstroke}{rgb}{0.000000,0.000000,1.000000}%
\pgfsetstrokecolor{currentstroke}%
\pgfsetstrokeopacity{0.500000}%
\pgfsetdash{}{0pt}%
\pgfpathmoveto{\pgfqpoint{0.648561in}{5.258411in}}%
\pgfusepath{stroke}%
\end{pgfscope}%
\begin{pgfscope}%
\pgfpathrectangle{\pgfqpoint{0.100000in}{2.493063in}}{\pgfqpoint{5.037500in}{3.427208in}}%
\pgfusepath{clip}%
\pgfsetbuttcap%
\pgfsetroundjoin%
\definecolor{currentfill}{rgb}{0.000000,0.000000,1.000000}%
\pgfsetfillcolor{currentfill}%
\pgfsetfillopacity{0.500000}%
\pgfsetlinewidth{0.250937pt}%
\definecolor{currentstroke}{rgb}{0.000000,0.000000,0.000000}%
\pgfsetstrokecolor{currentstroke}%
\pgfsetstrokeopacity{0.500000}%
\pgfsetdash{}{0pt}%
\pgfsys@defobject{currentmarker}{\pgfqpoint{-0.016667in}{-0.016667in}}{\pgfqpoint{0.016667in}{0.016667in}}{%
\pgfpathmoveto{\pgfqpoint{0.000000in}{-0.016667in}}%
\pgfpathcurveto{\pgfqpoint{0.004420in}{-0.016667in}}{\pgfqpoint{0.008660in}{-0.014911in}}{\pgfqpoint{0.011785in}{-0.011785in}}%
\pgfpathcurveto{\pgfqpoint{0.014911in}{-0.008660in}}{\pgfqpoint{0.016667in}{-0.004420in}}{\pgfqpoint{0.016667in}{0.000000in}}%
\pgfpathcurveto{\pgfqpoint{0.016667in}{0.004420in}}{\pgfqpoint{0.014911in}{0.008660in}}{\pgfqpoint{0.011785in}{0.011785in}}%
\pgfpathcurveto{\pgfqpoint{0.008660in}{0.014911in}}{\pgfqpoint{0.004420in}{0.016667in}}{\pgfqpoint{0.000000in}{0.016667in}}%
\pgfpathcurveto{\pgfqpoint{-0.004420in}{0.016667in}}{\pgfqpoint{-0.008660in}{0.014911in}}{\pgfqpoint{-0.011785in}{0.011785in}}%
\pgfpathcurveto{\pgfqpoint{-0.014911in}{0.008660in}}{\pgfqpoint{-0.016667in}{0.004420in}}{\pgfqpoint{-0.016667in}{0.000000in}}%
\pgfpathcurveto{\pgfqpoint{-0.016667in}{-0.004420in}}{\pgfqpoint{-0.014911in}{-0.008660in}}{\pgfqpoint{-0.011785in}{-0.011785in}}%
\pgfpathcurveto{\pgfqpoint{-0.008660in}{-0.014911in}}{\pgfqpoint{-0.004420in}{-0.016667in}}{\pgfqpoint{0.000000in}{-0.016667in}}%
\pgfpathlineto{\pgfqpoint{0.000000in}{-0.016667in}}%
\pgfpathclose%
\pgfusepath{stroke,fill}%
}%
\begin{pgfscope}%
\pgfsys@transformshift{0.648561in}{5.258411in}%
\pgfsys@useobject{currentmarker}{}%
\end{pgfscope}%
\end{pgfscope}%
\begin{pgfscope}%
\pgfpathrectangle{\pgfqpoint{0.100000in}{2.493063in}}{\pgfqpoint{5.037500in}{3.427208in}}%
\pgfusepath{clip}%
\pgfsetrectcap%
\pgfsetroundjoin%
\pgfsetlinewidth{1.505625pt}%
\definecolor{currentstroke}{rgb}{0.000000,0.000000,1.000000}%
\pgfsetstrokecolor{currentstroke}%
\pgfsetstrokeopacity{0.500000}%
\pgfsetdash{}{0pt}%
\pgfpathmoveto{\pgfqpoint{0.512006in}{5.361994in}}%
\pgfusepath{stroke}%
\end{pgfscope}%
\begin{pgfscope}%
\pgfpathrectangle{\pgfqpoint{0.100000in}{2.493063in}}{\pgfqpoint{5.037500in}{3.427208in}}%
\pgfusepath{clip}%
\pgfsetbuttcap%
\pgfsetroundjoin%
\definecolor{currentfill}{rgb}{0.000000,0.000000,1.000000}%
\pgfsetfillcolor{currentfill}%
\pgfsetfillopacity{0.500000}%
\pgfsetlinewidth{0.250937pt}%
\definecolor{currentstroke}{rgb}{0.000000,0.000000,0.000000}%
\pgfsetstrokecolor{currentstroke}%
\pgfsetstrokeopacity{0.500000}%
\pgfsetdash{}{0pt}%
\pgfsys@defobject{currentmarker}{\pgfqpoint{-0.013889in}{-0.013889in}}{\pgfqpoint{0.013889in}{0.013889in}}{%
\pgfpathmoveto{\pgfqpoint{0.000000in}{-0.013889in}}%
\pgfpathcurveto{\pgfqpoint{0.003683in}{-0.013889in}}{\pgfqpoint{0.007216in}{-0.012425in}}{\pgfqpoint{0.009821in}{-0.009821in}}%
\pgfpathcurveto{\pgfqpoint{0.012425in}{-0.007216in}}{\pgfqpoint{0.013889in}{-0.003683in}}{\pgfqpoint{0.013889in}{0.000000in}}%
\pgfpathcurveto{\pgfqpoint{0.013889in}{0.003683in}}{\pgfqpoint{0.012425in}{0.007216in}}{\pgfqpoint{0.009821in}{0.009821in}}%
\pgfpathcurveto{\pgfqpoint{0.007216in}{0.012425in}}{\pgfqpoint{0.003683in}{0.013889in}}{\pgfqpoint{0.000000in}{0.013889in}}%
\pgfpathcurveto{\pgfqpoint{-0.003683in}{0.013889in}}{\pgfqpoint{-0.007216in}{0.012425in}}{\pgfqpoint{-0.009821in}{0.009821in}}%
\pgfpathcurveto{\pgfqpoint{-0.012425in}{0.007216in}}{\pgfqpoint{-0.013889in}{0.003683in}}{\pgfqpoint{-0.013889in}{0.000000in}}%
\pgfpathcurveto{\pgfqpoint{-0.013889in}{-0.003683in}}{\pgfqpoint{-0.012425in}{-0.007216in}}{\pgfqpoint{-0.009821in}{-0.009821in}}%
\pgfpathcurveto{\pgfqpoint{-0.007216in}{-0.012425in}}{\pgfqpoint{-0.003683in}{-0.013889in}}{\pgfqpoint{0.000000in}{-0.013889in}}%
\pgfpathlineto{\pgfqpoint{0.000000in}{-0.013889in}}%
\pgfpathclose%
\pgfusepath{stroke,fill}%
}%
\begin{pgfscope}%
\pgfsys@transformshift{0.512006in}{5.361994in}%
\pgfsys@useobject{currentmarker}{}%
\end{pgfscope}%
\end{pgfscope}%
\begin{pgfscope}%
\pgfpathrectangle{\pgfqpoint{0.100000in}{2.493063in}}{\pgfqpoint{5.037500in}{3.427208in}}%
\pgfusepath{clip}%
\pgfsetrectcap%
\pgfsetroundjoin%
\pgfsetlinewidth{1.505625pt}%
\definecolor{currentstroke}{rgb}{0.000000,0.000000,1.000000}%
\pgfsetstrokecolor{currentstroke}%
\pgfsetstrokeopacity{0.500000}%
\pgfsetdash{}{0pt}%
\pgfpathmoveto{\pgfqpoint{0.504276in}{5.308196in}}%
\pgfusepath{stroke}%
\end{pgfscope}%
\begin{pgfscope}%
\pgfpathrectangle{\pgfqpoint{0.100000in}{2.493063in}}{\pgfqpoint{5.037500in}{3.427208in}}%
\pgfusepath{clip}%
\pgfsetbuttcap%
\pgfsetroundjoin%
\definecolor{currentfill}{rgb}{0.000000,0.000000,1.000000}%
\pgfsetfillcolor{currentfill}%
\pgfsetfillopacity{0.500000}%
\pgfsetlinewidth{0.250937pt}%
\definecolor{currentstroke}{rgb}{0.000000,0.000000,0.000000}%
\pgfsetstrokecolor{currentstroke}%
\pgfsetstrokeopacity{0.500000}%
\pgfsetdash{}{0pt}%
\pgfsys@defobject{currentmarker}{\pgfqpoint{-0.019444in}{-0.019444in}}{\pgfqpoint{0.019444in}{0.019444in}}{%
\pgfpathmoveto{\pgfqpoint{0.000000in}{-0.019444in}}%
\pgfpathcurveto{\pgfqpoint{0.005157in}{-0.019444in}}{\pgfqpoint{0.010103in}{-0.017396in}}{\pgfqpoint{0.013749in}{-0.013749in}}%
\pgfpathcurveto{\pgfqpoint{0.017396in}{-0.010103in}}{\pgfqpoint{0.019444in}{-0.005157in}}{\pgfqpoint{0.019444in}{0.000000in}}%
\pgfpathcurveto{\pgfqpoint{0.019444in}{0.005157in}}{\pgfqpoint{0.017396in}{0.010103in}}{\pgfqpoint{0.013749in}{0.013749in}}%
\pgfpathcurveto{\pgfqpoint{0.010103in}{0.017396in}}{\pgfqpoint{0.005157in}{0.019444in}}{\pgfqpoint{0.000000in}{0.019444in}}%
\pgfpathcurveto{\pgfqpoint{-0.005157in}{0.019444in}}{\pgfqpoint{-0.010103in}{0.017396in}}{\pgfqpoint{-0.013749in}{0.013749in}}%
\pgfpathcurveto{\pgfqpoint{-0.017396in}{0.010103in}}{\pgfqpoint{-0.019444in}{0.005157in}}{\pgfqpoint{-0.019444in}{0.000000in}}%
\pgfpathcurveto{\pgfqpoint{-0.019444in}{-0.005157in}}{\pgfqpoint{-0.017396in}{-0.010103in}}{\pgfqpoint{-0.013749in}{-0.013749in}}%
\pgfpathcurveto{\pgfqpoint{-0.010103in}{-0.017396in}}{\pgfqpoint{-0.005157in}{-0.019444in}}{\pgfqpoint{0.000000in}{-0.019444in}}%
\pgfpathlineto{\pgfqpoint{0.000000in}{-0.019444in}}%
\pgfpathclose%
\pgfusepath{stroke,fill}%
}%
\begin{pgfscope}%
\pgfsys@transformshift{0.504276in}{5.308196in}%
\pgfsys@useobject{currentmarker}{}%
\end{pgfscope}%
\end{pgfscope}%
\begin{pgfscope}%
\pgfpathrectangle{\pgfqpoint{0.100000in}{2.493063in}}{\pgfqpoint{5.037500in}{3.427208in}}%
\pgfusepath{clip}%
\pgfsetrectcap%
\pgfsetroundjoin%
\pgfsetlinewidth{1.505625pt}%
\definecolor{currentstroke}{rgb}{0.000000,0.000000,1.000000}%
\pgfsetstrokecolor{currentstroke}%
\pgfsetstrokeopacity{0.500000}%
\pgfsetdash{}{0pt}%
\pgfpathmoveto{\pgfqpoint{0.431607in}{5.130000in}}%
\pgfusepath{stroke}%
\end{pgfscope}%
\begin{pgfscope}%
\pgfpathrectangle{\pgfqpoint{0.100000in}{2.493063in}}{\pgfqpoint{5.037500in}{3.427208in}}%
\pgfusepath{clip}%
\pgfsetbuttcap%
\pgfsetroundjoin%
\definecolor{currentfill}{rgb}{0.000000,0.000000,1.000000}%
\pgfsetfillcolor{currentfill}%
\pgfsetfillopacity{0.500000}%
\pgfsetlinewidth{0.250937pt}%
\definecolor{currentstroke}{rgb}{0.000000,0.000000,0.000000}%
\pgfsetstrokecolor{currentstroke}%
\pgfsetstrokeopacity{0.500000}%
\pgfsetdash{}{0pt}%
\pgfsys@defobject{currentmarker}{\pgfqpoint{-0.022222in}{-0.022222in}}{\pgfqpoint{0.022222in}{0.022222in}}{%
\pgfpathmoveto{\pgfqpoint{0.000000in}{-0.022222in}}%
\pgfpathcurveto{\pgfqpoint{0.005893in}{-0.022222in}}{\pgfqpoint{0.011546in}{-0.019881in}}{\pgfqpoint{0.015713in}{-0.015713in}}%
\pgfpathcurveto{\pgfqpoint{0.019881in}{-0.011546in}}{\pgfqpoint{0.022222in}{-0.005893in}}{\pgfqpoint{0.022222in}{0.000000in}}%
\pgfpathcurveto{\pgfqpoint{0.022222in}{0.005893in}}{\pgfqpoint{0.019881in}{0.011546in}}{\pgfqpoint{0.015713in}{0.015713in}}%
\pgfpathcurveto{\pgfqpoint{0.011546in}{0.019881in}}{\pgfqpoint{0.005893in}{0.022222in}}{\pgfqpoint{0.000000in}{0.022222in}}%
\pgfpathcurveto{\pgfqpoint{-0.005893in}{0.022222in}}{\pgfqpoint{-0.011546in}{0.019881in}}{\pgfqpoint{-0.015713in}{0.015713in}}%
\pgfpathcurveto{\pgfqpoint{-0.019881in}{0.011546in}}{\pgfqpoint{-0.022222in}{0.005893in}}{\pgfqpoint{-0.022222in}{0.000000in}}%
\pgfpathcurveto{\pgfqpoint{-0.022222in}{-0.005893in}}{\pgfqpoint{-0.019881in}{-0.011546in}}{\pgfqpoint{-0.015713in}{-0.015713in}}%
\pgfpathcurveto{\pgfqpoint{-0.011546in}{-0.019881in}}{\pgfqpoint{-0.005893in}{-0.022222in}}{\pgfqpoint{0.000000in}{-0.022222in}}%
\pgfpathlineto{\pgfqpoint{0.000000in}{-0.022222in}}%
\pgfpathclose%
\pgfusepath{stroke,fill}%
}%
\begin{pgfscope}%
\pgfsys@transformshift{0.431607in}{5.130000in}%
\pgfsys@useobject{currentmarker}{}%
\end{pgfscope}%
\end{pgfscope}%
\begin{pgfscope}%
\pgfpathrectangle{\pgfqpoint{0.100000in}{2.493063in}}{\pgfqpoint{5.037500in}{3.427208in}}%
\pgfusepath{clip}%
\pgfsetrectcap%
\pgfsetroundjoin%
\pgfsetlinewidth{1.505625pt}%
\definecolor{currentstroke}{rgb}{0.000000,0.000000,1.000000}%
\pgfsetstrokecolor{currentstroke}%
\pgfsetstrokeopacity{0.500000}%
\pgfsetdash{}{0pt}%
\pgfpathmoveto{\pgfqpoint{0.464601in}{5.105783in}}%
\pgfusepath{stroke}%
\end{pgfscope}%
\begin{pgfscope}%
\pgfpathrectangle{\pgfqpoint{0.100000in}{2.493063in}}{\pgfqpoint{5.037500in}{3.427208in}}%
\pgfusepath{clip}%
\pgfsetbuttcap%
\pgfsetroundjoin%
\definecolor{currentfill}{rgb}{0.000000,0.000000,1.000000}%
\pgfsetfillcolor{currentfill}%
\pgfsetfillopacity{0.500000}%
\pgfsetlinewidth{0.250937pt}%
\definecolor{currentstroke}{rgb}{0.000000,0.000000,0.000000}%
\pgfsetstrokecolor{currentstroke}%
\pgfsetstrokeopacity{0.500000}%
\pgfsetdash{}{0pt}%
\pgfsys@defobject{currentmarker}{\pgfqpoint{-0.022222in}{-0.022222in}}{\pgfqpoint{0.022222in}{0.022222in}}{%
\pgfpathmoveto{\pgfqpoint{0.000000in}{-0.022222in}}%
\pgfpathcurveto{\pgfqpoint{0.005893in}{-0.022222in}}{\pgfqpoint{0.011546in}{-0.019881in}}{\pgfqpoint{0.015713in}{-0.015713in}}%
\pgfpathcurveto{\pgfqpoint{0.019881in}{-0.011546in}}{\pgfqpoint{0.022222in}{-0.005893in}}{\pgfqpoint{0.022222in}{0.000000in}}%
\pgfpathcurveto{\pgfqpoint{0.022222in}{0.005893in}}{\pgfqpoint{0.019881in}{0.011546in}}{\pgfqpoint{0.015713in}{0.015713in}}%
\pgfpathcurveto{\pgfqpoint{0.011546in}{0.019881in}}{\pgfqpoint{0.005893in}{0.022222in}}{\pgfqpoint{0.000000in}{0.022222in}}%
\pgfpathcurveto{\pgfqpoint{-0.005893in}{0.022222in}}{\pgfqpoint{-0.011546in}{0.019881in}}{\pgfqpoint{-0.015713in}{0.015713in}}%
\pgfpathcurveto{\pgfqpoint{-0.019881in}{0.011546in}}{\pgfqpoint{-0.022222in}{0.005893in}}{\pgfqpoint{-0.022222in}{0.000000in}}%
\pgfpathcurveto{\pgfqpoint{-0.022222in}{-0.005893in}}{\pgfqpoint{-0.019881in}{-0.011546in}}{\pgfqpoint{-0.015713in}{-0.015713in}}%
\pgfpathcurveto{\pgfqpoint{-0.011546in}{-0.019881in}}{\pgfqpoint{-0.005893in}{-0.022222in}}{\pgfqpoint{0.000000in}{-0.022222in}}%
\pgfpathlineto{\pgfqpoint{0.000000in}{-0.022222in}}%
\pgfpathclose%
\pgfusepath{stroke,fill}%
}%
\begin{pgfscope}%
\pgfsys@transformshift{0.464601in}{5.105783in}%
\pgfsys@useobject{currentmarker}{}%
\end{pgfscope}%
\end{pgfscope}%
\begin{pgfscope}%
\pgfpathrectangle{\pgfqpoint{0.100000in}{2.493063in}}{\pgfqpoint{5.037500in}{3.427208in}}%
\pgfusepath{clip}%
\pgfsetrectcap%
\pgfsetroundjoin%
\pgfsetlinewidth{1.505625pt}%
\definecolor{currentstroke}{rgb}{0.000000,0.000000,1.000000}%
\pgfsetstrokecolor{currentstroke}%
\pgfsetstrokeopacity{0.500000}%
\pgfsetdash{}{0pt}%
\pgfpathmoveto{\pgfqpoint{0.591363in}{5.452886in}}%
\pgfusepath{stroke}%
\end{pgfscope}%
\begin{pgfscope}%
\pgfpathrectangle{\pgfqpoint{0.100000in}{2.493063in}}{\pgfqpoint{5.037500in}{3.427208in}}%
\pgfusepath{clip}%
\pgfsetbuttcap%
\pgfsetroundjoin%
\definecolor{currentfill}{rgb}{0.000000,0.000000,1.000000}%
\pgfsetfillcolor{currentfill}%
\pgfsetfillopacity{0.500000}%
\pgfsetlinewidth{0.250937pt}%
\definecolor{currentstroke}{rgb}{0.000000,0.000000,0.000000}%
\pgfsetstrokecolor{currentstroke}%
\pgfsetstrokeopacity{0.500000}%
\pgfsetdash{}{0pt}%
\pgfsys@defobject{currentmarker}{\pgfqpoint{-0.013889in}{-0.013889in}}{\pgfqpoint{0.013889in}{0.013889in}}{%
\pgfpathmoveto{\pgfqpoint{0.000000in}{-0.013889in}}%
\pgfpathcurveto{\pgfqpoint{0.003683in}{-0.013889in}}{\pgfqpoint{0.007216in}{-0.012425in}}{\pgfqpoint{0.009821in}{-0.009821in}}%
\pgfpathcurveto{\pgfqpoint{0.012425in}{-0.007216in}}{\pgfqpoint{0.013889in}{-0.003683in}}{\pgfqpoint{0.013889in}{0.000000in}}%
\pgfpathcurveto{\pgfqpoint{0.013889in}{0.003683in}}{\pgfqpoint{0.012425in}{0.007216in}}{\pgfqpoint{0.009821in}{0.009821in}}%
\pgfpathcurveto{\pgfqpoint{0.007216in}{0.012425in}}{\pgfqpoint{0.003683in}{0.013889in}}{\pgfqpoint{0.000000in}{0.013889in}}%
\pgfpathcurveto{\pgfqpoint{-0.003683in}{0.013889in}}{\pgfqpoint{-0.007216in}{0.012425in}}{\pgfqpoint{-0.009821in}{0.009821in}}%
\pgfpathcurveto{\pgfqpoint{-0.012425in}{0.007216in}}{\pgfqpoint{-0.013889in}{0.003683in}}{\pgfqpoint{-0.013889in}{0.000000in}}%
\pgfpathcurveto{\pgfqpoint{-0.013889in}{-0.003683in}}{\pgfqpoint{-0.012425in}{-0.007216in}}{\pgfqpoint{-0.009821in}{-0.009821in}}%
\pgfpathcurveto{\pgfqpoint{-0.007216in}{-0.012425in}}{\pgfqpoint{-0.003683in}{-0.013889in}}{\pgfqpoint{0.000000in}{-0.013889in}}%
\pgfpathlineto{\pgfqpoint{0.000000in}{-0.013889in}}%
\pgfpathclose%
\pgfusepath{stroke,fill}%
}%
\begin{pgfscope}%
\pgfsys@transformshift{0.591363in}{5.452886in}%
\pgfsys@useobject{currentmarker}{}%
\end{pgfscope}%
\end{pgfscope}%
\begin{pgfscope}%
\pgfpathrectangle{\pgfqpoint{0.100000in}{2.493063in}}{\pgfqpoint{5.037500in}{3.427208in}}%
\pgfusepath{clip}%
\pgfsetrectcap%
\pgfsetroundjoin%
\pgfsetlinewidth{1.505625pt}%
\definecolor{currentstroke}{rgb}{0.000000,0.000000,1.000000}%
\pgfsetstrokecolor{currentstroke}%
\pgfsetstrokeopacity{0.500000}%
\pgfsetdash{}{0pt}%
\pgfpathmoveto{\pgfqpoint{0.543139in}{5.397548in}}%
\pgfusepath{stroke}%
\end{pgfscope}%
\begin{pgfscope}%
\pgfpathrectangle{\pgfqpoint{0.100000in}{2.493063in}}{\pgfqpoint{5.037500in}{3.427208in}}%
\pgfusepath{clip}%
\pgfsetbuttcap%
\pgfsetroundjoin%
\definecolor{currentfill}{rgb}{0.000000,0.000000,1.000000}%
\pgfsetfillcolor{currentfill}%
\pgfsetfillopacity{0.500000}%
\pgfsetlinewidth{0.250937pt}%
\definecolor{currentstroke}{rgb}{0.000000,0.000000,0.000000}%
\pgfsetstrokecolor{currentstroke}%
\pgfsetstrokeopacity{0.500000}%
\pgfsetdash{}{0pt}%
\pgfsys@defobject{currentmarker}{\pgfqpoint{-0.013889in}{-0.013889in}}{\pgfqpoint{0.013889in}{0.013889in}}{%
\pgfpathmoveto{\pgfqpoint{0.000000in}{-0.013889in}}%
\pgfpathcurveto{\pgfqpoint{0.003683in}{-0.013889in}}{\pgfqpoint{0.007216in}{-0.012425in}}{\pgfqpoint{0.009821in}{-0.009821in}}%
\pgfpathcurveto{\pgfqpoint{0.012425in}{-0.007216in}}{\pgfqpoint{0.013889in}{-0.003683in}}{\pgfqpoint{0.013889in}{0.000000in}}%
\pgfpathcurveto{\pgfqpoint{0.013889in}{0.003683in}}{\pgfqpoint{0.012425in}{0.007216in}}{\pgfqpoint{0.009821in}{0.009821in}}%
\pgfpathcurveto{\pgfqpoint{0.007216in}{0.012425in}}{\pgfqpoint{0.003683in}{0.013889in}}{\pgfqpoint{0.000000in}{0.013889in}}%
\pgfpathcurveto{\pgfqpoint{-0.003683in}{0.013889in}}{\pgfqpoint{-0.007216in}{0.012425in}}{\pgfqpoint{-0.009821in}{0.009821in}}%
\pgfpathcurveto{\pgfqpoint{-0.012425in}{0.007216in}}{\pgfqpoint{-0.013889in}{0.003683in}}{\pgfqpoint{-0.013889in}{0.000000in}}%
\pgfpathcurveto{\pgfqpoint{-0.013889in}{-0.003683in}}{\pgfqpoint{-0.012425in}{-0.007216in}}{\pgfqpoint{-0.009821in}{-0.009821in}}%
\pgfpathcurveto{\pgfqpoint{-0.007216in}{-0.012425in}}{\pgfqpoint{-0.003683in}{-0.013889in}}{\pgfqpoint{0.000000in}{-0.013889in}}%
\pgfpathlineto{\pgfqpoint{0.000000in}{-0.013889in}}%
\pgfpathclose%
\pgfusepath{stroke,fill}%
}%
\begin{pgfscope}%
\pgfsys@transformshift{0.543139in}{5.397548in}%
\pgfsys@useobject{currentmarker}{}%
\end{pgfscope}%
\end{pgfscope}%
\begin{pgfscope}%
\pgfpathrectangle{\pgfqpoint{0.100000in}{2.493063in}}{\pgfqpoint{5.037500in}{3.427208in}}%
\pgfusepath{clip}%
\pgfsetrectcap%
\pgfsetroundjoin%
\pgfsetlinewidth{1.505625pt}%
\definecolor{currentstroke}{rgb}{0.678431,1.000000,0.184314}%
\pgfsetstrokecolor{currentstroke}%
\pgfsetstrokeopacity{0.500000}%
\pgfsetdash{}{0pt}%
\pgfpathmoveto{\pgfqpoint{4.498915in}{4.728262in}}%
\pgfusepath{stroke}%
\end{pgfscope}%
\begin{pgfscope}%
\pgfpathrectangle{\pgfqpoint{0.100000in}{2.493063in}}{\pgfqpoint{5.037500in}{3.427208in}}%
\pgfusepath{clip}%
\pgfsetbuttcap%
\pgfsetroundjoin%
\definecolor{currentfill}{rgb}{0.678431,1.000000,0.184314}%
\pgfsetfillcolor{currentfill}%
\pgfsetfillopacity{0.500000}%
\pgfsetlinewidth{0.250937pt}%
\definecolor{currentstroke}{rgb}{0.000000,0.000000,0.000000}%
\pgfsetstrokecolor{currentstroke}%
\pgfsetstrokeopacity{0.500000}%
\pgfsetdash{}{0pt}%
\pgfsys@defobject{currentmarker}{\pgfqpoint{-0.013889in}{-0.013889in}}{\pgfqpoint{0.013889in}{0.013889in}}{%
\pgfpathmoveto{\pgfqpoint{0.000000in}{-0.013889in}}%
\pgfpathcurveto{\pgfqpoint{0.003683in}{-0.013889in}}{\pgfqpoint{0.007216in}{-0.012425in}}{\pgfqpoint{0.009821in}{-0.009821in}}%
\pgfpathcurveto{\pgfqpoint{0.012425in}{-0.007216in}}{\pgfqpoint{0.013889in}{-0.003683in}}{\pgfqpoint{0.013889in}{0.000000in}}%
\pgfpathcurveto{\pgfqpoint{0.013889in}{0.003683in}}{\pgfqpoint{0.012425in}{0.007216in}}{\pgfqpoint{0.009821in}{0.009821in}}%
\pgfpathcurveto{\pgfqpoint{0.007216in}{0.012425in}}{\pgfqpoint{0.003683in}{0.013889in}}{\pgfqpoint{0.000000in}{0.013889in}}%
\pgfpathcurveto{\pgfqpoint{-0.003683in}{0.013889in}}{\pgfqpoint{-0.007216in}{0.012425in}}{\pgfqpoint{-0.009821in}{0.009821in}}%
\pgfpathcurveto{\pgfqpoint{-0.012425in}{0.007216in}}{\pgfqpoint{-0.013889in}{0.003683in}}{\pgfqpoint{-0.013889in}{0.000000in}}%
\pgfpathcurveto{\pgfqpoint{-0.013889in}{-0.003683in}}{\pgfqpoint{-0.012425in}{-0.007216in}}{\pgfqpoint{-0.009821in}{-0.009821in}}%
\pgfpathcurveto{\pgfqpoint{-0.007216in}{-0.012425in}}{\pgfqpoint{-0.003683in}{-0.013889in}}{\pgfqpoint{0.000000in}{-0.013889in}}%
\pgfpathlineto{\pgfqpoint{0.000000in}{-0.013889in}}%
\pgfpathclose%
\pgfusepath{stroke,fill}%
}%
\begin{pgfscope}%
\pgfsys@transformshift{4.498915in}{4.728262in}%
\pgfsys@useobject{currentmarker}{}%
\end{pgfscope}%
\end{pgfscope}%
\begin{pgfscope}%
\pgfpathrectangle{\pgfqpoint{0.100000in}{2.493063in}}{\pgfqpoint{5.037500in}{3.427208in}}%
\pgfusepath{clip}%
\pgfsetrectcap%
\pgfsetroundjoin%
\pgfsetlinewidth{1.505625pt}%
\definecolor{currentstroke}{rgb}{0.678431,1.000000,0.184314}%
\pgfsetstrokecolor{currentstroke}%
\pgfsetstrokeopacity{0.500000}%
\pgfsetdash{}{0pt}%
\pgfpathmoveto{\pgfqpoint{4.250037in}{4.668307in}}%
\pgfusepath{stroke}%
\end{pgfscope}%
\begin{pgfscope}%
\pgfpathrectangle{\pgfqpoint{0.100000in}{2.493063in}}{\pgfqpoint{5.037500in}{3.427208in}}%
\pgfusepath{clip}%
\pgfsetbuttcap%
\pgfsetroundjoin%
\definecolor{currentfill}{rgb}{0.678431,1.000000,0.184314}%
\pgfsetfillcolor{currentfill}%
\pgfsetfillopacity{0.500000}%
\pgfsetlinewidth{0.250937pt}%
\definecolor{currentstroke}{rgb}{0.000000,0.000000,0.000000}%
\pgfsetstrokecolor{currentstroke}%
\pgfsetstrokeopacity{0.500000}%
\pgfsetdash{}{0pt}%
\pgfsys@defobject{currentmarker}{\pgfqpoint{-0.041667in}{-0.041667in}}{\pgfqpoint{0.041667in}{0.041667in}}{%
\pgfpathmoveto{\pgfqpoint{0.000000in}{-0.041667in}}%
\pgfpathcurveto{\pgfqpoint{0.011050in}{-0.041667in}}{\pgfqpoint{0.021649in}{-0.037276in}}{\pgfqpoint{0.029463in}{-0.029463in}}%
\pgfpathcurveto{\pgfqpoint{0.037276in}{-0.021649in}}{\pgfqpoint{0.041667in}{-0.011050in}}{\pgfqpoint{0.041667in}{0.000000in}}%
\pgfpathcurveto{\pgfqpoint{0.041667in}{0.011050in}}{\pgfqpoint{0.037276in}{0.021649in}}{\pgfqpoint{0.029463in}{0.029463in}}%
\pgfpathcurveto{\pgfqpoint{0.021649in}{0.037276in}}{\pgfqpoint{0.011050in}{0.041667in}}{\pgfqpoint{0.000000in}{0.041667in}}%
\pgfpathcurveto{\pgfqpoint{-0.011050in}{0.041667in}}{\pgfqpoint{-0.021649in}{0.037276in}}{\pgfqpoint{-0.029463in}{0.029463in}}%
\pgfpathcurveto{\pgfqpoint{-0.037276in}{0.021649in}}{\pgfqpoint{-0.041667in}{0.011050in}}{\pgfqpoint{-0.041667in}{0.000000in}}%
\pgfpathcurveto{\pgfqpoint{-0.041667in}{-0.011050in}}{\pgfqpoint{-0.037276in}{-0.021649in}}{\pgfqpoint{-0.029463in}{-0.029463in}}%
\pgfpathcurveto{\pgfqpoint{-0.021649in}{-0.037276in}}{\pgfqpoint{-0.011050in}{-0.041667in}}{\pgfqpoint{0.000000in}{-0.041667in}}%
\pgfpathlineto{\pgfqpoint{0.000000in}{-0.041667in}}%
\pgfpathclose%
\pgfusepath{stroke,fill}%
}%
\begin{pgfscope}%
\pgfsys@transformshift{4.250037in}{4.668307in}%
\pgfsys@useobject{currentmarker}{}%
\end{pgfscope}%
\end{pgfscope}%
\begin{pgfscope}%
\pgfpathrectangle{\pgfqpoint{0.100000in}{2.493063in}}{\pgfqpoint{5.037500in}{3.427208in}}%
\pgfusepath{clip}%
\pgfsetrectcap%
\pgfsetroundjoin%
\pgfsetlinewidth{1.505625pt}%
\definecolor{currentstroke}{rgb}{0.678431,1.000000,0.184314}%
\pgfsetstrokecolor{currentstroke}%
\pgfsetstrokeopacity{0.500000}%
\pgfsetdash{}{0pt}%
\pgfpathmoveto{\pgfqpoint{4.405482in}{4.755732in}}%
\pgfusepath{stroke}%
\end{pgfscope}%
\begin{pgfscope}%
\pgfpathrectangle{\pgfqpoint{0.100000in}{2.493063in}}{\pgfqpoint{5.037500in}{3.427208in}}%
\pgfusepath{clip}%
\pgfsetbuttcap%
\pgfsetroundjoin%
\definecolor{currentfill}{rgb}{0.678431,1.000000,0.184314}%
\pgfsetfillcolor{currentfill}%
\pgfsetfillopacity{0.500000}%
\pgfsetlinewidth{0.250937pt}%
\definecolor{currentstroke}{rgb}{0.000000,0.000000,0.000000}%
\pgfsetstrokecolor{currentstroke}%
\pgfsetstrokeopacity{0.500000}%
\pgfsetdash{}{0pt}%
\pgfsys@defobject{currentmarker}{\pgfqpoint{-0.025000in}{-0.025000in}}{\pgfqpoint{0.025000in}{0.025000in}}{%
\pgfpathmoveto{\pgfqpoint{0.000000in}{-0.025000in}}%
\pgfpathcurveto{\pgfqpoint{0.006630in}{-0.025000in}}{\pgfqpoint{0.012989in}{-0.022366in}}{\pgfqpoint{0.017678in}{-0.017678in}}%
\pgfpathcurveto{\pgfqpoint{0.022366in}{-0.012989in}}{\pgfqpoint{0.025000in}{-0.006630in}}{\pgfqpoint{0.025000in}{0.000000in}}%
\pgfpathcurveto{\pgfqpoint{0.025000in}{0.006630in}}{\pgfqpoint{0.022366in}{0.012989in}}{\pgfqpoint{0.017678in}{0.017678in}}%
\pgfpathcurveto{\pgfqpoint{0.012989in}{0.022366in}}{\pgfqpoint{0.006630in}{0.025000in}}{\pgfqpoint{0.000000in}{0.025000in}}%
\pgfpathcurveto{\pgfqpoint{-0.006630in}{0.025000in}}{\pgfqpoint{-0.012989in}{0.022366in}}{\pgfqpoint{-0.017678in}{0.017678in}}%
\pgfpathcurveto{\pgfqpoint{-0.022366in}{0.012989in}}{\pgfqpoint{-0.025000in}{0.006630in}}{\pgfqpoint{-0.025000in}{0.000000in}}%
\pgfpathcurveto{\pgfqpoint{-0.025000in}{-0.006630in}}{\pgfqpoint{-0.022366in}{-0.012989in}}{\pgfqpoint{-0.017678in}{-0.017678in}}%
\pgfpathcurveto{\pgfqpoint{-0.012989in}{-0.022366in}}{\pgfqpoint{-0.006630in}{-0.025000in}}{\pgfqpoint{0.000000in}{-0.025000in}}%
\pgfpathlineto{\pgfqpoint{0.000000in}{-0.025000in}}%
\pgfpathclose%
\pgfusepath{stroke,fill}%
}%
\begin{pgfscope}%
\pgfsys@transformshift{4.405482in}{4.755732in}%
\pgfsys@useobject{currentmarker}{}%
\end{pgfscope}%
\end{pgfscope}%
\begin{pgfscope}%
\pgfpathrectangle{\pgfqpoint{0.100000in}{2.493063in}}{\pgfqpoint{5.037500in}{3.427208in}}%
\pgfusepath{clip}%
\pgfsetrectcap%
\pgfsetroundjoin%
\pgfsetlinewidth{1.505625pt}%
\definecolor{currentstroke}{rgb}{0.678431,1.000000,0.184314}%
\pgfsetstrokecolor{currentstroke}%
\pgfsetstrokeopacity{0.500000}%
\pgfsetdash{}{0pt}%
\pgfpathmoveto{\pgfqpoint{4.325892in}{4.614885in}}%
\pgfusepath{stroke}%
\end{pgfscope}%
\begin{pgfscope}%
\pgfpathrectangle{\pgfqpoint{0.100000in}{2.493063in}}{\pgfqpoint{5.037500in}{3.427208in}}%
\pgfusepath{clip}%
\pgfsetbuttcap%
\pgfsetroundjoin%
\definecolor{currentfill}{rgb}{0.678431,1.000000,0.184314}%
\pgfsetfillcolor{currentfill}%
\pgfsetfillopacity{0.500000}%
\pgfsetlinewidth{0.250937pt}%
\definecolor{currentstroke}{rgb}{0.000000,0.000000,0.000000}%
\pgfsetstrokecolor{currentstroke}%
\pgfsetstrokeopacity{0.500000}%
\pgfsetdash{}{0pt}%
\pgfsys@defobject{currentmarker}{\pgfqpoint{-0.030556in}{-0.030556in}}{\pgfqpoint{0.030556in}{0.030556in}}{%
\pgfpathmoveto{\pgfqpoint{0.000000in}{-0.030556in}}%
\pgfpathcurveto{\pgfqpoint{0.008103in}{-0.030556in}}{\pgfqpoint{0.015876in}{-0.027336in}}{\pgfqpoint{0.021606in}{-0.021606in}}%
\pgfpathcurveto{\pgfqpoint{0.027336in}{-0.015876in}}{\pgfqpoint{0.030556in}{-0.008103in}}{\pgfqpoint{0.030556in}{0.000000in}}%
\pgfpathcurveto{\pgfqpoint{0.030556in}{0.008103in}}{\pgfqpoint{0.027336in}{0.015876in}}{\pgfqpoint{0.021606in}{0.021606in}}%
\pgfpathcurveto{\pgfqpoint{0.015876in}{0.027336in}}{\pgfqpoint{0.008103in}{0.030556in}}{\pgfqpoint{0.000000in}{0.030556in}}%
\pgfpathcurveto{\pgfqpoint{-0.008103in}{0.030556in}}{\pgfqpoint{-0.015876in}{0.027336in}}{\pgfqpoint{-0.021606in}{0.021606in}}%
\pgfpathcurveto{\pgfqpoint{-0.027336in}{0.015876in}}{\pgfqpoint{-0.030556in}{0.008103in}}{\pgfqpoint{-0.030556in}{0.000000in}}%
\pgfpathcurveto{\pgfqpoint{-0.030556in}{-0.008103in}}{\pgfqpoint{-0.027336in}{-0.015876in}}{\pgfqpoint{-0.021606in}{-0.021606in}}%
\pgfpathcurveto{\pgfqpoint{-0.015876in}{-0.027336in}}{\pgfqpoint{-0.008103in}{-0.030556in}}{\pgfqpoint{0.000000in}{-0.030556in}}%
\pgfpathlineto{\pgfqpoint{0.000000in}{-0.030556in}}%
\pgfpathclose%
\pgfusepath{stroke,fill}%
}%
\begin{pgfscope}%
\pgfsys@transformshift{4.325892in}{4.614885in}%
\pgfsys@useobject{currentmarker}{}%
\end{pgfscope}%
\end{pgfscope}%
\begin{pgfscope}%
\pgfpathrectangle{\pgfqpoint{0.100000in}{2.493063in}}{\pgfqpoint{5.037500in}{3.427208in}}%
\pgfusepath{clip}%
\pgfsetrectcap%
\pgfsetroundjoin%
\pgfsetlinewidth{1.505625pt}%
\definecolor{currentstroke}{rgb}{0.678431,1.000000,0.184314}%
\pgfsetstrokecolor{currentstroke}%
\pgfsetstrokeopacity{0.500000}%
\pgfsetdash{}{0pt}%
\pgfpathmoveto{\pgfqpoint{4.513781in}{4.778398in}}%
\pgfusepath{stroke}%
\end{pgfscope}%
\begin{pgfscope}%
\pgfpathrectangle{\pgfqpoint{0.100000in}{2.493063in}}{\pgfqpoint{5.037500in}{3.427208in}}%
\pgfusepath{clip}%
\pgfsetbuttcap%
\pgfsetroundjoin%
\definecolor{currentfill}{rgb}{0.678431,1.000000,0.184314}%
\pgfsetfillcolor{currentfill}%
\pgfsetfillopacity{0.500000}%
\pgfsetlinewidth{0.250937pt}%
\definecolor{currentstroke}{rgb}{0.000000,0.000000,0.000000}%
\pgfsetstrokecolor{currentstroke}%
\pgfsetstrokeopacity{0.500000}%
\pgfsetdash{}{0pt}%
\pgfsys@defobject{currentmarker}{\pgfqpoint{-0.030556in}{-0.030556in}}{\pgfqpoint{0.030556in}{0.030556in}}{%
\pgfpathmoveto{\pgfqpoint{0.000000in}{-0.030556in}}%
\pgfpathcurveto{\pgfqpoint{0.008103in}{-0.030556in}}{\pgfqpoint{0.015876in}{-0.027336in}}{\pgfqpoint{0.021606in}{-0.021606in}}%
\pgfpathcurveto{\pgfqpoint{0.027336in}{-0.015876in}}{\pgfqpoint{0.030556in}{-0.008103in}}{\pgfqpoint{0.030556in}{0.000000in}}%
\pgfpathcurveto{\pgfqpoint{0.030556in}{0.008103in}}{\pgfqpoint{0.027336in}{0.015876in}}{\pgfqpoint{0.021606in}{0.021606in}}%
\pgfpathcurveto{\pgfqpoint{0.015876in}{0.027336in}}{\pgfqpoint{0.008103in}{0.030556in}}{\pgfqpoint{0.000000in}{0.030556in}}%
\pgfpathcurveto{\pgfqpoint{-0.008103in}{0.030556in}}{\pgfqpoint{-0.015876in}{0.027336in}}{\pgfqpoint{-0.021606in}{0.021606in}}%
\pgfpathcurveto{\pgfqpoint{-0.027336in}{0.015876in}}{\pgfqpoint{-0.030556in}{0.008103in}}{\pgfqpoint{-0.030556in}{0.000000in}}%
\pgfpathcurveto{\pgfqpoint{-0.030556in}{-0.008103in}}{\pgfqpoint{-0.027336in}{-0.015876in}}{\pgfqpoint{-0.021606in}{-0.021606in}}%
\pgfpathcurveto{\pgfqpoint{-0.015876in}{-0.027336in}}{\pgfqpoint{-0.008103in}{-0.030556in}}{\pgfqpoint{0.000000in}{-0.030556in}}%
\pgfpathlineto{\pgfqpoint{0.000000in}{-0.030556in}}%
\pgfpathclose%
\pgfusepath{stroke,fill}%
}%
\begin{pgfscope}%
\pgfsys@transformshift{4.513781in}{4.778398in}%
\pgfsys@useobject{currentmarker}{}%
\end{pgfscope}%
\end{pgfscope}%
\begin{pgfscope}%
\pgfpathrectangle{\pgfqpoint{0.100000in}{2.493063in}}{\pgfqpoint{5.037500in}{3.427208in}}%
\pgfusepath{clip}%
\pgfsetrectcap%
\pgfsetroundjoin%
\pgfsetlinewidth{1.505625pt}%
\definecolor{currentstroke}{rgb}{0.678431,1.000000,0.184314}%
\pgfsetstrokecolor{currentstroke}%
\pgfsetstrokeopacity{0.500000}%
\pgfsetdash{}{0pt}%
\pgfpathmoveto{\pgfqpoint{4.073915in}{4.825853in}}%
\pgfusepath{stroke}%
\end{pgfscope}%
\begin{pgfscope}%
\pgfpathrectangle{\pgfqpoint{0.100000in}{2.493063in}}{\pgfqpoint{5.037500in}{3.427208in}}%
\pgfusepath{clip}%
\pgfsetbuttcap%
\pgfsetroundjoin%
\definecolor{currentfill}{rgb}{0.678431,1.000000,0.184314}%
\pgfsetfillcolor{currentfill}%
\pgfsetfillopacity{0.500000}%
\pgfsetlinewidth{0.250937pt}%
\definecolor{currentstroke}{rgb}{0.000000,0.000000,0.000000}%
\pgfsetstrokecolor{currentstroke}%
\pgfsetstrokeopacity{0.500000}%
\pgfsetdash{}{0pt}%
\pgfsys@defobject{currentmarker}{\pgfqpoint{-0.027778in}{-0.027778in}}{\pgfqpoint{0.027778in}{0.027778in}}{%
\pgfpathmoveto{\pgfqpoint{0.000000in}{-0.027778in}}%
\pgfpathcurveto{\pgfqpoint{0.007367in}{-0.027778in}}{\pgfqpoint{0.014433in}{-0.024851in}}{\pgfqpoint{0.019642in}{-0.019642in}}%
\pgfpathcurveto{\pgfqpoint{0.024851in}{-0.014433in}}{\pgfqpoint{0.027778in}{-0.007367in}}{\pgfqpoint{0.027778in}{0.000000in}}%
\pgfpathcurveto{\pgfqpoint{0.027778in}{0.007367in}}{\pgfqpoint{0.024851in}{0.014433in}}{\pgfqpoint{0.019642in}{0.019642in}}%
\pgfpathcurveto{\pgfqpoint{0.014433in}{0.024851in}}{\pgfqpoint{0.007367in}{0.027778in}}{\pgfqpoint{0.000000in}{0.027778in}}%
\pgfpathcurveto{\pgfqpoint{-0.007367in}{0.027778in}}{\pgfqpoint{-0.014433in}{0.024851in}}{\pgfqpoint{-0.019642in}{0.019642in}}%
\pgfpathcurveto{\pgfqpoint{-0.024851in}{0.014433in}}{\pgfqpoint{-0.027778in}{0.007367in}}{\pgfqpoint{-0.027778in}{0.000000in}}%
\pgfpathcurveto{\pgfqpoint{-0.027778in}{-0.007367in}}{\pgfqpoint{-0.024851in}{-0.014433in}}{\pgfqpoint{-0.019642in}{-0.019642in}}%
\pgfpathcurveto{\pgfqpoint{-0.014433in}{-0.024851in}}{\pgfqpoint{-0.007367in}{-0.027778in}}{\pgfqpoint{0.000000in}{-0.027778in}}%
\pgfpathlineto{\pgfqpoint{0.000000in}{-0.027778in}}%
\pgfpathclose%
\pgfusepath{stroke,fill}%
}%
\begin{pgfscope}%
\pgfsys@transformshift{4.073915in}{4.825853in}%
\pgfsys@useobject{currentmarker}{}%
\end{pgfscope}%
\end{pgfscope}%
\begin{pgfscope}%
\pgfpathrectangle{\pgfqpoint{0.100000in}{2.493063in}}{\pgfqpoint{5.037500in}{3.427208in}}%
\pgfusepath{clip}%
\pgfsetrectcap%
\pgfsetroundjoin%
\pgfsetlinewidth{1.505625pt}%
\definecolor{currentstroke}{rgb}{0.678431,1.000000,0.184314}%
\pgfsetstrokecolor{currentstroke}%
\pgfsetstrokeopacity{0.500000}%
\pgfsetdash{}{0pt}%
\pgfpathmoveto{\pgfqpoint{4.365499in}{4.609766in}}%
\pgfusepath{stroke}%
\end{pgfscope}%
\begin{pgfscope}%
\pgfpathrectangle{\pgfqpoint{0.100000in}{2.493063in}}{\pgfqpoint{5.037500in}{3.427208in}}%
\pgfusepath{clip}%
\pgfsetbuttcap%
\pgfsetroundjoin%
\definecolor{currentfill}{rgb}{0.678431,1.000000,0.184314}%
\pgfsetfillcolor{currentfill}%
\pgfsetfillopacity{0.500000}%
\pgfsetlinewidth{0.250937pt}%
\definecolor{currentstroke}{rgb}{0.000000,0.000000,0.000000}%
\pgfsetstrokecolor{currentstroke}%
\pgfsetstrokeopacity{0.500000}%
\pgfsetdash{}{0pt}%
\pgfsys@defobject{currentmarker}{\pgfqpoint{-0.022222in}{-0.022222in}}{\pgfqpoint{0.022222in}{0.022222in}}{%
\pgfpathmoveto{\pgfqpoint{0.000000in}{-0.022222in}}%
\pgfpathcurveto{\pgfqpoint{0.005893in}{-0.022222in}}{\pgfqpoint{0.011546in}{-0.019881in}}{\pgfqpoint{0.015713in}{-0.015713in}}%
\pgfpathcurveto{\pgfqpoint{0.019881in}{-0.011546in}}{\pgfqpoint{0.022222in}{-0.005893in}}{\pgfqpoint{0.022222in}{0.000000in}}%
\pgfpathcurveto{\pgfqpoint{0.022222in}{0.005893in}}{\pgfqpoint{0.019881in}{0.011546in}}{\pgfqpoint{0.015713in}{0.015713in}}%
\pgfpathcurveto{\pgfqpoint{0.011546in}{0.019881in}}{\pgfqpoint{0.005893in}{0.022222in}}{\pgfqpoint{0.000000in}{0.022222in}}%
\pgfpathcurveto{\pgfqpoint{-0.005893in}{0.022222in}}{\pgfqpoint{-0.011546in}{0.019881in}}{\pgfqpoint{-0.015713in}{0.015713in}}%
\pgfpathcurveto{\pgfqpoint{-0.019881in}{0.011546in}}{\pgfqpoint{-0.022222in}{0.005893in}}{\pgfqpoint{-0.022222in}{0.000000in}}%
\pgfpathcurveto{\pgfqpoint{-0.022222in}{-0.005893in}}{\pgfqpoint{-0.019881in}{-0.011546in}}{\pgfqpoint{-0.015713in}{-0.015713in}}%
\pgfpathcurveto{\pgfqpoint{-0.011546in}{-0.019881in}}{\pgfqpoint{-0.005893in}{-0.022222in}}{\pgfqpoint{0.000000in}{-0.022222in}}%
\pgfpathlineto{\pgfqpoint{0.000000in}{-0.022222in}}%
\pgfpathclose%
\pgfusepath{stroke,fill}%
}%
\begin{pgfscope}%
\pgfsys@transformshift{4.365499in}{4.609766in}%
\pgfsys@useobject{currentmarker}{}%
\end{pgfscope}%
\end{pgfscope}%
\begin{pgfscope}%
\pgfpathrectangle{\pgfqpoint{0.100000in}{2.493063in}}{\pgfqpoint{5.037500in}{3.427208in}}%
\pgfusepath{clip}%
\pgfsetrectcap%
\pgfsetroundjoin%
\pgfsetlinewidth{1.505625pt}%
\definecolor{currentstroke}{rgb}{0.678431,1.000000,0.184314}%
\pgfsetstrokecolor{currentstroke}%
\pgfsetstrokeopacity{0.500000}%
\pgfsetdash{}{0pt}%
\pgfpathmoveto{\pgfqpoint{4.385486in}{4.664896in}}%
\pgfusepath{stroke}%
\end{pgfscope}%
\begin{pgfscope}%
\pgfpathrectangle{\pgfqpoint{0.100000in}{2.493063in}}{\pgfqpoint{5.037500in}{3.427208in}}%
\pgfusepath{clip}%
\pgfsetbuttcap%
\pgfsetroundjoin%
\definecolor{currentfill}{rgb}{0.678431,1.000000,0.184314}%
\pgfsetfillcolor{currentfill}%
\pgfsetfillopacity{0.500000}%
\pgfsetlinewidth{0.250937pt}%
\definecolor{currentstroke}{rgb}{0.000000,0.000000,0.000000}%
\pgfsetstrokecolor{currentstroke}%
\pgfsetstrokeopacity{0.500000}%
\pgfsetdash{}{0pt}%
\pgfsys@defobject{currentmarker}{\pgfqpoint{-0.025000in}{-0.025000in}}{\pgfqpoint{0.025000in}{0.025000in}}{%
\pgfpathmoveto{\pgfqpoint{0.000000in}{-0.025000in}}%
\pgfpathcurveto{\pgfqpoint{0.006630in}{-0.025000in}}{\pgfqpoint{0.012989in}{-0.022366in}}{\pgfqpoint{0.017678in}{-0.017678in}}%
\pgfpathcurveto{\pgfqpoint{0.022366in}{-0.012989in}}{\pgfqpoint{0.025000in}{-0.006630in}}{\pgfqpoint{0.025000in}{0.000000in}}%
\pgfpathcurveto{\pgfqpoint{0.025000in}{0.006630in}}{\pgfqpoint{0.022366in}{0.012989in}}{\pgfqpoint{0.017678in}{0.017678in}}%
\pgfpathcurveto{\pgfqpoint{0.012989in}{0.022366in}}{\pgfqpoint{0.006630in}{0.025000in}}{\pgfqpoint{0.000000in}{0.025000in}}%
\pgfpathcurveto{\pgfqpoint{-0.006630in}{0.025000in}}{\pgfqpoint{-0.012989in}{0.022366in}}{\pgfqpoint{-0.017678in}{0.017678in}}%
\pgfpathcurveto{\pgfqpoint{-0.022366in}{0.012989in}}{\pgfqpoint{-0.025000in}{0.006630in}}{\pgfqpoint{-0.025000in}{0.000000in}}%
\pgfpathcurveto{\pgfqpoint{-0.025000in}{-0.006630in}}{\pgfqpoint{-0.022366in}{-0.012989in}}{\pgfqpoint{-0.017678in}{-0.017678in}}%
\pgfpathcurveto{\pgfqpoint{-0.012989in}{-0.022366in}}{\pgfqpoint{-0.006630in}{-0.025000in}}{\pgfqpoint{0.000000in}{-0.025000in}}%
\pgfpathlineto{\pgfqpoint{0.000000in}{-0.025000in}}%
\pgfpathclose%
\pgfusepath{stroke,fill}%
}%
\begin{pgfscope}%
\pgfsys@transformshift{4.385486in}{4.664896in}%
\pgfsys@useobject{currentmarker}{}%
\end{pgfscope}%
\end{pgfscope}%
\begin{pgfscope}%
\pgfpathrectangle{\pgfqpoint{0.100000in}{2.493063in}}{\pgfqpoint{5.037500in}{3.427208in}}%
\pgfusepath{clip}%
\pgfsetrectcap%
\pgfsetroundjoin%
\pgfsetlinewidth{1.505625pt}%
\definecolor{currentstroke}{rgb}{0.678431,1.000000,0.184314}%
\pgfsetstrokecolor{currentstroke}%
\pgfsetstrokeopacity{0.500000}%
\pgfsetdash{}{0pt}%
\pgfpathmoveto{\pgfqpoint{4.208507in}{4.638288in}}%
\pgfusepath{stroke}%
\end{pgfscope}%
\begin{pgfscope}%
\pgfpathrectangle{\pgfqpoint{0.100000in}{2.493063in}}{\pgfqpoint{5.037500in}{3.427208in}}%
\pgfusepath{clip}%
\pgfsetbuttcap%
\pgfsetroundjoin%
\definecolor{currentfill}{rgb}{0.678431,1.000000,0.184314}%
\pgfsetfillcolor{currentfill}%
\pgfsetfillopacity{0.500000}%
\pgfsetlinewidth{0.250937pt}%
\definecolor{currentstroke}{rgb}{0.000000,0.000000,0.000000}%
\pgfsetstrokecolor{currentstroke}%
\pgfsetstrokeopacity{0.500000}%
\pgfsetdash{}{0pt}%
\pgfsys@defobject{currentmarker}{\pgfqpoint{-0.036111in}{-0.036111in}}{\pgfqpoint{0.036111in}{0.036111in}}{%
\pgfpathmoveto{\pgfqpoint{0.000000in}{-0.036111in}}%
\pgfpathcurveto{\pgfqpoint{0.009577in}{-0.036111in}}{\pgfqpoint{0.018763in}{-0.032306in}}{\pgfqpoint{0.025534in}{-0.025534in}}%
\pgfpathcurveto{\pgfqpoint{0.032306in}{-0.018763in}}{\pgfqpoint{0.036111in}{-0.009577in}}{\pgfqpoint{0.036111in}{0.000000in}}%
\pgfpathcurveto{\pgfqpoint{0.036111in}{0.009577in}}{\pgfqpoint{0.032306in}{0.018763in}}{\pgfqpoint{0.025534in}{0.025534in}}%
\pgfpathcurveto{\pgfqpoint{0.018763in}{0.032306in}}{\pgfqpoint{0.009577in}{0.036111in}}{\pgfqpoint{0.000000in}{0.036111in}}%
\pgfpathcurveto{\pgfqpoint{-0.009577in}{0.036111in}}{\pgfqpoint{-0.018763in}{0.032306in}}{\pgfqpoint{-0.025534in}{0.025534in}}%
\pgfpathcurveto{\pgfqpoint{-0.032306in}{0.018763in}}{\pgfqpoint{-0.036111in}{0.009577in}}{\pgfqpoint{-0.036111in}{0.000000in}}%
\pgfpathcurveto{\pgfqpoint{-0.036111in}{-0.009577in}}{\pgfqpoint{-0.032306in}{-0.018763in}}{\pgfqpoint{-0.025534in}{-0.025534in}}%
\pgfpathcurveto{\pgfqpoint{-0.018763in}{-0.032306in}}{\pgfqpoint{-0.009577in}{-0.036111in}}{\pgfqpoint{0.000000in}{-0.036111in}}%
\pgfpathlineto{\pgfqpoint{0.000000in}{-0.036111in}}%
\pgfpathclose%
\pgfusepath{stroke,fill}%
}%
\begin{pgfscope}%
\pgfsys@transformshift{4.208507in}{4.638288in}%
\pgfsys@useobject{currentmarker}{}%
\end{pgfscope}%
\end{pgfscope}%
\begin{pgfscope}%
\pgfpathrectangle{\pgfqpoint{0.100000in}{2.493063in}}{\pgfqpoint{5.037500in}{3.427208in}}%
\pgfusepath{clip}%
\pgfsetrectcap%
\pgfsetroundjoin%
\pgfsetlinewidth{1.505625pt}%
\definecolor{currentstroke}{rgb}{0.678431,1.000000,0.184314}%
\pgfsetstrokecolor{currentstroke}%
\pgfsetstrokeopacity{0.500000}%
\pgfsetdash{}{0pt}%
\pgfpathmoveto{\pgfqpoint{4.440822in}{4.649431in}}%
\pgfusepath{stroke}%
\end{pgfscope}%
\begin{pgfscope}%
\pgfpathrectangle{\pgfqpoint{0.100000in}{2.493063in}}{\pgfqpoint{5.037500in}{3.427208in}}%
\pgfusepath{clip}%
\pgfsetbuttcap%
\pgfsetroundjoin%
\definecolor{currentfill}{rgb}{0.678431,1.000000,0.184314}%
\pgfsetfillcolor{currentfill}%
\pgfsetfillopacity{0.500000}%
\pgfsetlinewidth{0.250937pt}%
\definecolor{currentstroke}{rgb}{0.000000,0.000000,0.000000}%
\pgfsetstrokecolor{currentstroke}%
\pgfsetstrokeopacity{0.500000}%
\pgfsetdash{}{0pt}%
\pgfsys@defobject{currentmarker}{\pgfqpoint{-0.022222in}{-0.022222in}}{\pgfqpoint{0.022222in}{0.022222in}}{%
\pgfpathmoveto{\pgfqpoint{0.000000in}{-0.022222in}}%
\pgfpathcurveto{\pgfqpoint{0.005893in}{-0.022222in}}{\pgfqpoint{0.011546in}{-0.019881in}}{\pgfqpoint{0.015713in}{-0.015713in}}%
\pgfpathcurveto{\pgfqpoint{0.019881in}{-0.011546in}}{\pgfqpoint{0.022222in}{-0.005893in}}{\pgfqpoint{0.022222in}{0.000000in}}%
\pgfpathcurveto{\pgfqpoint{0.022222in}{0.005893in}}{\pgfqpoint{0.019881in}{0.011546in}}{\pgfqpoint{0.015713in}{0.015713in}}%
\pgfpathcurveto{\pgfqpoint{0.011546in}{0.019881in}}{\pgfqpoint{0.005893in}{0.022222in}}{\pgfqpoint{0.000000in}{0.022222in}}%
\pgfpathcurveto{\pgfqpoint{-0.005893in}{0.022222in}}{\pgfqpoint{-0.011546in}{0.019881in}}{\pgfqpoint{-0.015713in}{0.015713in}}%
\pgfpathcurveto{\pgfqpoint{-0.019881in}{0.011546in}}{\pgfqpoint{-0.022222in}{0.005893in}}{\pgfqpoint{-0.022222in}{0.000000in}}%
\pgfpathcurveto{\pgfqpoint{-0.022222in}{-0.005893in}}{\pgfqpoint{-0.019881in}{-0.011546in}}{\pgfqpoint{-0.015713in}{-0.015713in}}%
\pgfpathcurveto{\pgfqpoint{-0.011546in}{-0.019881in}}{\pgfqpoint{-0.005893in}{-0.022222in}}{\pgfqpoint{0.000000in}{-0.022222in}}%
\pgfpathlineto{\pgfqpoint{0.000000in}{-0.022222in}}%
\pgfpathclose%
\pgfusepath{stroke,fill}%
}%
\begin{pgfscope}%
\pgfsys@transformshift{4.440822in}{4.649431in}%
\pgfsys@useobject{currentmarker}{}%
\end{pgfscope}%
\end{pgfscope}%
\begin{pgfscope}%
\pgfpathrectangle{\pgfqpoint{0.100000in}{2.493063in}}{\pgfqpoint{5.037500in}{3.427208in}}%
\pgfusepath{clip}%
\pgfsetrectcap%
\pgfsetroundjoin%
\pgfsetlinewidth{1.505625pt}%
\definecolor{currentstroke}{rgb}{0.678431,1.000000,0.184314}%
\pgfsetstrokecolor{currentstroke}%
\pgfsetstrokeopacity{0.500000}%
\pgfsetdash{}{0pt}%
\pgfpathmoveto{\pgfqpoint{4.423020in}{4.682009in}}%
\pgfusepath{stroke}%
\end{pgfscope}%
\begin{pgfscope}%
\pgfpathrectangle{\pgfqpoint{0.100000in}{2.493063in}}{\pgfqpoint{5.037500in}{3.427208in}}%
\pgfusepath{clip}%
\pgfsetbuttcap%
\pgfsetroundjoin%
\definecolor{currentfill}{rgb}{0.678431,1.000000,0.184314}%
\pgfsetfillcolor{currentfill}%
\pgfsetfillopacity{0.500000}%
\pgfsetlinewidth{0.250937pt}%
\definecolor{currentstroke}{rgb}{0.000000,0.000000,0.000000}%
\pgfsetstrokecolor{currentstroke}%
\pgfsetstrokeopacity{0.500000}%
\pgfsetdash{}{0pt}%
\pgfsys@defobject{currentmarker}{\pgfqpoint{-0.036111in}{-0.036111in}}{\pgfqpoint{0.036111in}{0.036111in}}{%
\pgfpathmoveto{\pgfqpoint{0.000000in}{-0.036111in}}%
\pgfpathcurveto{\pgfqpoint{0.009577in}{-0.036111in}}{\pgfqpoint{0.018763in}{-0.032306in}}{\pgfqpoint{0.025534in}{-0.025534in}}%
\pgfpathcurveto{\pgfqpoint{0.032306in}{-0.018763in}}{\pgfqpoint{0.036111in}{-0.009577in}}{\pgfqpoint{0.036111in}{0.000000in}}%
\pgfpathcurveto{\pgfqpoint{0.036111in}{0.009577in}}{\pgfqpoint{0.032306in}{0.018763in}}{\pgfqpoint{0.025534in}{0.025534in}}%
\pgfpathcurveto{\pgfqpoint{0.018763in}{0.032306in}}{\pgfqpoint{0.009577in}{0.036111in}}{\pgfqpoint{0.000000in}{0.036111in}}%
\pgfpathcurveto{\pgfqpoint{-0.009577in}{0.036111in}}{\pgfqpoint{-0.018763in}{0.032306in}}{\pgfqpoint{-0.025534in}{0.025534in}}%
\pgfpathcurveto{\pgfqpoint{-0.032306in}{0.018763in}}{\pgfqpoint{-0.036111in}{0.009577in}}{\pgfqpoint{-0.036111in}{0.000000in}}%
\pgfpathcurveto{\pgfqpoint{-0.036111in}{-0.009577in}}{\pgfqpoint{-0.032306in}{-0.018763in}}{\pgfqpoint{-0.025534in}{-0.025534in}}%
\pgfpathcurveto{\pgfqpoint{-0.018763in}{-0.032306in}}{\pgfqpoint{-0.009577in}{-0.036111in}}{\pgfqpoint{0.000000in}{-0.036111in}}%
\pgfpathlineto{\pgfqpoint{0.000000in}{-0.036111in}}%
\pgfpathclose%
\pgfusepath{stroke,fill}%
}%
\begin{pgfscope}%
\pgfsys@transformshift{4.423020in}{4.682009in}%
\pgfsys@useobject{currentmarker}{}%
\end{pgfscope}%
\end{pgfscope}%
\begin{pgfscope}%
\pgfpathrectangle{\pgfqpoint{0.100000in}{2.493063in}}{\pgfqpoint{5.037500in}{3.427208in}}%
\pgfusepath{clip}%
\pgfsetrectcap%
\pgfsetroundjoin%
\pgfsetlinewidth{1.505625pt}%
\definecolor{currentstroke}{rgb}{0.678431,1.000000,0.184314}%
\pgfsetstrokecolor{currentstroke}%
\pgfsetstrokeopacity{0.500000}%
\pgfsetdash{}{0pt}%
\pgfpathmoveto{\pgfqpoint{4.541403in}{4.660980in}}%
\pgfusepath{stroke}%
\end{pgfscope}%
\begin{pgfscope}%
\pgfpathrectangle{\pgfqpoint{0.100000in}{2.493063in}}{\pgfqpoint{5.037500in}{3.427208in}}%
\pgfusepath{clip}%
\pgfsetbuttcap%
\pgfsetroundjoin%
\definecolor{currentfill}{rgb}{0.678431,1.000000,0.184314}%
\pgfsetfillcolor{currentfill}%
\pgfsetfillopacity{0.500000}%
\pgfsetlinewidth{0.250937pt}%
\definecolor{currentstroke}{rgb}{0.000000,0.000000,0.000000}%
\pgfsetstrokecolor{currentstroke}%
\pgfsetstrokeopacity{0.500000}%
\pgfsetdash{}{0pt}%
\pgfsys@defobject{currentmarker}{\pgfqpoint{-0.011111in}{-0.011111in}}{\pgfqpoint{0.011111in}{0.011111in}}{%
\pgfpathmoveto{\pgfqpoint{0.000000in}{-0.011111in}}%
\pgfpathcurveto{\pgfqpoint{0.002947in}{-0.011111in}}{\pgfqpoint{0.005773in}{-0.009940in}}{\pgfqpoint{0.007857in}{-0.007857in}}%
\pgfpathcurveto{\pgfqpoint{0.009940in}{-0.005773in}}{\pgfqpoint{0.011111in}{-0.002947in}}{\pgfqpoint{0.011111in}{0.000000in}}%
\pgfpathcurveto{\pgfqpoint{0.011111in}{0.002947in}}{\pgfqpoint{0.009940in}{0.005773in}}{\pgfqpoint{0.007857in}{0.007857in}}%
\pgfpathcurveto{\pgfqpoint{0.005773in}{0.009940in}}{\pgfqpoint{0.002947in}{0.011111in}}{\pgfqpoint{0.000000in}{0.011111in}}%
\pgfpathcurveto{\pgfqpoint{-0.002947in}{0.011111in}}{\pgfqpoint{-0.005773in}{0.009940in}}{\pgfqpoint{-0.007857in}{0.007857in}}%
\pgfpathcurveto{\pgfqpoint{-0.009940in}{0.005773in}}{\pgfqpoint{-0.011111in}{0.002947in}}{\pgfqpoint{-0.011111in}{0.000000in}}%
\pgfpathcurveto{\pgfqpoint{-0.011111in}{-0.002947in}}{\pgfqpoint{-0.009940in}{-0.005773in}}{\pgfqpoint{-0.007857in}{-0.007857in}}%
\pgfpathcurveto{\pgfqpoint{-0.005773in}{-0.009940in}}{\pgfqpoint{-0.002947in}{-0.011111in}}{\pgfqpoint{0.000000in}{-0.011111in}}%
\pgfpathlineto{\pgfqpoint{0.000000in}{-0.011111in}}%
\pgfpathclose%
\pgfusepath{stroke,fill}%
}%
\begin{pgfscope}%
\pgfsys@transformshift{4.541403in}{4.660980in}%
\pgfsys@useobject{currentmarker}{}%
\end{pgfscope}%
\end{pgfscope}%
\begin{pgfscope}%
\pgfpathrectangle{\pgfqpoint{0.100000in}{2.493063in}}{\pgfqpoint{5.037500in}{3.427208in}}%
\pgfusepath{clip}%
\pgfsetrectcap%
\pgfsetroundjoin%
\pgfsetlinewidth{1.505625pt}%
\definecolor{currentstroke}{rgb}{0.678431,1.000000,0.184314}%
\pgfsetstrokecolor{currentstroke}%
\pgfsetstrokeopacity{0.500000}%
\pgfsetdash{}{0pt}%
\pgfpathmoveto{\pgfqpoint{4.113913in}{4.635385in}}%
\pgfusepath{stroke}%
\end{pgfscope}%
\begin{pgfscope}%
\pgfpathrectangle{\pgfqpoint{0.100000in}{2.493063in}}{\pgfqpoint{5.037500in}{3.427208in}}%
\pgfusepath{clip}%
\pgfsetbuttcap%
\pgfsetroundjoin%
\definecolor{currentfill}{rgb}{0.678431,1.000000,0.184314}%
\pgfsetfillcolor{currentfill}%
\pgfsetfillopacity{0.500000}%
\pgfsetlinewidth{0.250937pt}%
\definecolor{currentstroke}{rgb}{0.000000,0.000000,0.000000}%
\pgfsetstrokecolor{currentstroke}%
\pgfsetstrokeopacity{0.500000}%
\pgfsetdash{}{0pt}%
\pgfsys@defobject{currentmarker}{\pgfqpoint{-0.027778in}{-0.027778in}}{\pgfqpoint{0.027778in}{0.027778in}}{%
\pgfpathmoveto{\pgfqpoint{0.000000in}{-0.027778in}}%
\pgfpathcurveto{\pgfqpoint{0.007367in}{-0.027778in}}{\pgfqpoint{0.014433in}{-0.024851in}}{\pgfqpoint{0.019642in}{-0.019642in}}%
\pgfpathcurveto{\pgfqpoint{0.024851in}{-0.014433in}}{\pgfqpoint{0.027778in}{-0.007367in}}{\pgfqpoint{0.027778in}{0.000000in}}%
\pgfpathcurveto{\pgfqpoint{0.027778in}{0.007367in}}{\pgfqpoint{0.024851in}{0.014433in}}{\pgfqpoint{0.019642in}{0.019642in}}%
\pgfpathcurveto{\pgfqpoint{0.014433in}{0.024851in}}{\pgfqpoint{0.007367in}{0.027778in}}{\pgfqpoint{0.000000in}{0.027778in}}%
\pgfpathcurveto{\pgfqpoint{-0.007367in}{0.027778in}}{\pgfqpoint{-0.014433in}{0.024851in}}{\pgfqpoint{-0.019642in}{0.019642in}}%
\pgfpathcurveto{\pgfqpoint{-0.024851in}{0.014433in}}{\pgfqpoint{-0.027778in}{0.007367in}}{\pgfqpoint{-0.027778in}{0.000000in}}%
\pgfpathcurveto{\pgfqpoint{-0.027778in}{-0.007367in}}{\pgfqpoint{-0.024851in}{-0.014433in}}{\pgfqpoint{-0.019642in}{-0.019642in}}%
\pgfpathcurveto{\pgfqpoint{-0.014433in}{-0.024851in}}{\pgfqpoint{-0.007367in}{-0.027778in}}{\pgfqpoint{0.000000in}{-0.027778in}}%
\pgfpathlineto{\pgfqpoint{0.000000in}{-0.027778in}}%
\pgfpathclose%
\pgfusepath{stroke,fill}%
}%
\begin{pgfscope}%
\pgfsys@transformshift{4.113913in}{4.635385in}%
\pgfsys@useobject{currentmarker}{}%
\end{pgfscope}%
\end{pgfscope}%
\begin{pgfscope}%
\pgfpathrectangle{\pgfqpoint{0.100000in}{2.493063in}}{\pgfqpoint{5.037500in}{3.427208in}}%
\pgfusepath{clip}%
\pgfsetrectcap%
\pgfsetroundjoin%
\pgfsetlinewidth{1.505625pt}%
\definecolor{currentstroke}{rgb}{0.678431,1.000000,0.184314}%
\pgfsetstrokecolor{currentstroke}%
\pgfsetstrokeopacity{0.500000}%
\pgfsetdash{}{0pt}%
\pgfpathmoveto{\pgfqpoint{4.466278in}{4.689780in}}%
\pgfusepath{stroke}%
\end{pgfscope}%
\begin{pgfscope}%
\pgfpathrectangle{\pgfqpoint{0.100000in}{2.493063in}}{\pgfqpoint{5.037500in}{3.427208in}}%
\pgfusepath{clip}%
\pgfsetbuttcap%
\pgfsetroundjoin%
\definecolor{currentfill}{rgb}{0.678431,1.000000,0.184314}%
\pgfsetfillcolor{currentfill}%
\pgfsetfillopacity{0.500000}%
\pgfsetlinewidth{0.250937pt}%
\definecolor{currentstroke}{rgb}{0.000000,0.000000,0.000000}%
\pgfsetstrokecolor{currentstroke}%
\pgfsetstrokeopacity{0.500000}%
\pgfsetdash{}{0pt}%
\pgfsys@defobject{currentmarker}{\pgfqpoint{-0.025000in}{-0.025000in}}{\pgfqpoint{0.025000in}{0.025000in}}{%
\pgfpathmoveto{\pgfqpoint{0.000000in}{-0.025000in}}%
\pgfpathcurveto{\pgfqpoint{0.006630in}{-0.025000in}}{\pgfqpoint{0.012989in}{-0.022366in}}{\pgfqpoint{0.017678in}{-0.017678in}}%
\pgfpathcurveto{\pgfqpoint{0.022366in}{-0.012989in}}{\pgfqpoint{0.025000in}{-0.006630in}}{\pgfqpoint{0.025000in}{0.000000in}}%
\pgfpathcurveto{\pgfqpoint{0.025000in}{0.006630in}}{\pgfqpoint{0.022366in}{0.012989in}}{\pgfqpoint{0.017678in}{0.017678in}}%
\pgfpathcurveto{\pgfqpoint{0.012989in}{0.022366in}}{\pgfqpoint{0.006630in}{0.025000in}}{\pgfqpoint{0.000000in}{0.025000in}}%
\pgfpathcurveto{\pgfqpoint{-0.006630in}{0.025000in}}{\pgfqpoint{-0.012989in}{0.022366in}}{\pgfqpoint{-0.017678in}{0.017678in}}%
\pgfpathcurveto{\pgfqpoint{-0.022366in}{0.012989in}}{\pgfqpoint{-0.025000in}{0.006630in}}{\pgfqpoint{-0.025000in}{0.000000in}}%
\pgfpathcurveto{\pgfqpoint{-0.025000in}{-0.006630in}}{\pgfqpoint{-0.022366in}{-0.012989in}}{\pgfqpoint{-0.017678in}{-0.017678in}}%
\pgfpathcurveto{\pgfqpoint{-0.012989in}{-0.022366in}}{\pgfqpoint{-0.006630in}{-0.025000in}}{\pgfqpoint{0.000000in}{-0.025000in}}%
\pgfpathlineto{\pgfqpoint{0.000000in}{-0.025000in}}%
\pgfpathclose%
\pgfusepath{stroke,fill}%
}%
\begin{pgfscope}%
\pgfsys@transformshift{4.466278in}{4.689780in}%
\pgfsys@useobject{currentmarker}{}%
\end{pgfscope}%
\end{pgfscope}%
\begin{pgfscope}%
\pgfpathrectangle{\pgfqpoint{0.100000in}{2.493063in}}{\pgfqpoint{5.037500in}{3.427208in}}%
\pgfusepath{clip}%
\pgfsetrectcap%
\pgfsetroundjoin%
\pgfsetlinewidth{1.505625pt}%
\definecolor{currentstroke}{rgb}{0.678431,1.000000,0.184314}%
\pgfsetstrokecolor{currentstroke}%
\pgfsetstrokeopacity{0.500000}%
\pgfsetdash{}{0pt}%
\pgfpathmoveto{\pgfqpoint{4.462981in}{4.815541in}}%
\pgfusepath{stroke}%
\end{pgfscope}%
\begin{pgfscope}%
\pgfpathrectangle{\pgfqpoint{0.100000in}{2.493063in}}{\pgfqpoint{5.037500in}{3.427208in}}%
\pgfusepath{clip}%
\pgfsetbuttcap%
\pgfsetroundjoin%
\definecolor{currentfill}{rgb}{0.678431,1.000000,0.184314}%
\pgfsetfillcolor{currentfill}%
\pgfsetfillopacity{0.500000}%
\pgfsetlinewidth{0.250937pt}%
\definecolor{currentstroke}{rgb}{0.000000,0.000000,0.000000}%
\pgfsetstrokecolor{currentstroke}%
\pgfsetstrokeopacity{0.500000}%
\pgfsetdash{}{0pt}%
\pgfsys@defobject{currentmarker}{\pgfqpoint{-0.036111in}{-0.036111in}}{\pgfqpoint{0.036111in}{0.036111in}}{%
\pgfpathmoveto{\pgfqpoint{0.000000in}{-0.036111in}}%
\pgfpathcurveto{\pgfqpoint{0.009577in}{-0.036111in}}{\pgfqpoint{0.018763in}{-0.032306in}}{\pgfqpoint{0.025534in}{-0.025534in}}%
\pgfpathcurveto{\pgfqpoint{0.032306in}{-0.018763in}}{\pgfqpoint{0.036111in}{-0.009577in}}{\pgfqpoint{0.036111in}{0.000000in}}%
\pgfpathcurveto{\pgfqpoint{0.036111in}{0.009577in}}{\pgfqpoint{0.032306in}{0.018763in}}{\pgfqpoint{0.025534in}{0.025534in}}%
\pgfpathcurveto{\pgfqpoint{0.018763in}{0.032306in}}{\pgfqpoint{0.009577in}{0.036111in}}{\pgfqpoint{0.000000in}{0.036111in}}%
\pgfpathcurveto{\pgfqpoint{-0.009577in}{0.036111in}}{\pgfqpoint{-0.018763in}{0.032306in}}{\pgfqpoint{-0.025534in}{0.025534in}}%
\pgfpathcurveto{\pgfqpoint{-0.032306in}{0.018763in}}{\pgfqpoint{-0.036111in}{0.009577in}}{\pgfqpoint{-0.036111in}{0.000000in}}%
\pgfpathcurveto{\pgfqpoint{-0.036111in}{-0.009577in}}{\pgfqpoint{-0.032306in}{-0.018763in}}{\pgfqpoint{-0.025534in}{-0.025534in}}%
\pgfpathcurveto{\pgfqpoint{-0.018763in}{-0.032306in}}{\pgfqpoint{-0.009577in}{-0.036111in}}{\pgfqpoint{0.000000in}{-0.036111in}}%
\pgfpathlineto{\pgfqpoint{0.000000in}{-0.036111in}}%
\pgfpathclose%
\pgfusepath{stroke,fill}%
}%
\begin{pgfscope}%
\pgfsys@transformshift{4.462981in}{4.815541in}%
\pgfsys@useobject{currentmarker}{}%
\end{pgfscope}%
\end{pgfscope}%
\begin{pgfscope}%
\pgfpathrectangle{\pgfqpoint{0.100000in}{2.493063in}}{\pgfqpoint{5.037500in}{3.427208in}}%
\pgfusepath{clip}%
\pgfsetrectcap%
\pgfsetroundjoin%
\pgfsetlinewidth{1.505625pt}%
\definecolor{currentstroke}{rgb}{0.678431,1.000000,0.184314}%
\pgfsetstrokecolor{currentstroke}%
\pgfsetstrokeopacity{0.500000}%
\pgfsetdash{}{0pt}%
\pgfpathmoveto{\pgfqpoint{4.290048in}{4.708094in}}%
\pgfusepath{stroke}%
\end{pgfscope}%
\begin{pgfscope}%
\pgfpathrectangle{\pgfqpoint{0.100000in}{2.493063in}}{\pgfqpoint{5.037500in}{3.427208in}}%
\pgfusepath{clip}%
\pgfsetbuttcap%
\pgfsetroundjoin%
\definecolor{currentfill}{rgb}{0.678431,1.000000,0.184314}%
\pgfsetfillcolor{currentfill}%
\pgfsetfillopacity{0.500000}%
\pgfsetlinewidth{0.250937pt}%
\definecolor{currentstroke}{rgb}{0.000000,0.000000,0.000000}%
\pgfsetstrokecolor{currentstroke}%
\pgfsetstrokeopacity{0.500000}%
\pgfsetdash{}{0pt}%
\pgfsys@defobject{currentmarker}{\pgfqpoint{-0.025000in}{-0.025000in}}{\pgfqpoint{0.025000in}{0.025000in}}{%
\pgfpathmoveto{\pgfqpoint{0.000000in}{-0.025000in}}%
\pgfpathcurveto{\pgfqpoint{0.006630in}{-0.025000in}}{\pgfqpoint{0.012989in}{-0.022366in}}{\pgfqpoint{0.017678in}{-0.017678in}}%
\pgfpathcurveto{\pgfqpoint{0.022366in}{-0.012989in}}{\pgfqpoint{0.025000in}{-0.006630in}}{\pgfqpoint{0.025000in}{0.000000in}}%
\pgfpathcurveto{\pgfqpoint{0.025000in}{0.006630in}}{\pgfqpoint{0.022366in}{0.012989in}}{\pgfqpoint{0.017678in}{0.017678in}}%
\pgfpathcurveto{\pgfqpoint{0.012989in}{0.022366in}}{\pgfqpoint{0.006630in}{0.025000in}}{\pgfqpoint{0.000000in}{0.025000in}}%
\pgfpathcurveto{\pgfqpoint{-0.006630in}{0.025000in}}{\pgfqpoint{-0.012989in}{0.022366in}}{\pgfqpoint{-0.017678in}{0.017678in}}%
\pgfpathcurveto{\pgfqpoint{-0.022366in}{0.012989in}}{\pgfqpoint{-0.025000in}{0.006630in}}{\pgfqpoint{-0.025000in}{0.000000in}}%
\pgfpathcurveto{\pgfqpoint{-0.025000in}{-0.006630in}}{\pgfqpoint{-0.022366in}{-0.012989in}}{\pgfqpoint{-0.017678in}{-0.017678in}}%
\pgfpathcurveto{\pgfqpoint{-0.012989in}{-0.022366in}}{\pgfqpoint{-0.006630in}{-0.025000in}}{\pgfqpoint{0.000000in}{-0.025000in}}%
\pgfpathlineto{\pgfqpoint{0.000000in}{-0.025000in}}%
\pgfpathclose%
\pgfusepath{stroke,fill}%
}%
\begin{pgfscope}%
\pgfsys@transformshift{4.290048in}{4.708094in}%
\pgfsys@useobject{currentmarker}{}%
\end{pgfscope}%
\end{pgfscope}%
\begin{pgfscope}%
\pgfpathrectangle{\pgfqpoint{0.100000in}{2.493063in}}{\pgfqpoint{5.037500in}{3.427208in}}%
\pgfusepath{clip}%
\pgfsetrectcap%
\pgfsetroundjoin%
\pgfsetlinewidth{1.505625pt}%
\definecolor{currentstroke}{rgb}{0.678431,1.000000,0.184314}%
\pgfsetstrokecolor{currentstroke}%
\pgfsetstrokeopacity{0.500000}%
\pgfsetdash{}{0pt}%
\pgfpathmoveto{\pgfqpoint{4.353156in}{4.773901in}}%
\pgfusepath{stroke}%
\end{pgfscope}%
\begin{pgfscope}%
\pgfpathrectangle{\pgfqpoint{0.100000in}{2.493063in}}{\pgfqpoint{5.037500in}{3.427208in}}%
\pgfusepath{clip}%
\pgfsetbuttcap%
\pgfsetroundjoin%
\definecolor{currentfill}{rgb}{0.678431,1.000000,0.184314}%
\pgfsetfillcolor{currentfill}%
\pgfsetfillopacity{0.500000}%
\pgfsetlinewidth{0.250937pt}%
\definecolor{currentstroke}{rgb}{0.000000,0.000000,0.000000}%
\pgfsetstrokecolor{currentstroke}%
\pgfsetstrokeopacity{0.500000}%
\pgfsetdash{}{0pt}%
\pgfsys@defobject{currentmarker}{\pgfqpoint{-0.036111in}{-0.036111in}}{\pgfqpoint{0.036111in}{0.036111in}}{%
\pgfpathmoveto{\pgfqpoint{0.000000in}{-0.036111in}}%
\pgfpathcurveto{\pgfqpoint{0.009577in}{-0.036111in}}{\pgfqpoint{0.018763in}{-0.032306in}}{\pgfqpoint{0.025534in}{-0.025534in}}%
\pgfpathcurveto{\pgfqpoint{0.032306in}{-0.018763in}}{\pgfqpoint{0.036111in}{-0.009577in}}{\pgfqpoint{0.036111in}{0.000000in}}%
\pgfpathcurveto{\pgfqpoint{0.036111in}{0.009577in}}{\pgfqpoint{0.032306in}{0.018763in}}{\pgfqpoint{0.025534in}{0.025534in}}%
\pgfpathcurveto{\pgfqpoint{0.018763in}{0.032306in}}{\pgfqpoint{0.009577in}{0.036111in}}{\pgfqpoint{0.000000in}{0.036111in}}%
\pgfpathcurveto{\pgfqpoint{-0.009577in}{0.036111in}}{\pgfqpoint{-0.018763in}{0.032306in}}{\pgfqpoint{-0.025534in}{0.025534in}}%
\pgfpathcurveto{\pgfqpoint{-0.032306in}{0.018763in}}{\pgfqpoint{-0.036111in}{0.009577in}}{\pgfqpoint{-0.036111in}{0.000000in}}%
\pgfpathcurveto{\pgfqpoint{-0.036111in}{-0.009577in}}{\pgfqpoint{-0.032306in}{-0.018763in}}{\pgfqpoint{-0.025534in}{-0.025534in}}%
\pgfpathcurveto{\pgfqpoint{-0.018763in}{-0.032306in}}{\pgfqpoint{-0.009577in}{-0.036111in}}{\pgfqpoint{0.000000in}{-0.036111in}}%
\pgfpathlineto{\pgfqpoint{0.000000in}{-0.036111in}}%
\pgfpathclose%
\pgfusepath{stroke,fill}%
}%
\begin{pgfscope}%
\pgfsys@transformshift{4.353156in}{4.773901in}%
\pgfsys@useobject{currentmarker}{}%
\end{pgfscope}%
\end{pgfscope}%
\begin{pgfscope}%
\pgfpathrectangle{\pgfqpoint{0.100000in}{2.493063in}}{\pgfqpoint{5.037500in}{3.427208in}}%
\pgfusepath{clip}%
\pgfsetrectcap%
\pgfsetroundjoin%
\pgfsetlinewidth{1.505625pt}%
\definecolor{currentstroke}{rgb}{0.678431,1.000000,0.184314}%
\pgfsetstrokecolor{currentstroke}%
\pgfsetstrokeopacity{0.500000}%
\pgfsetdash{}{0pt}%
\pgfpathmoveto{\pgfqpoint{4.406146in}{4.633374in}}%
\pgfusepath{stroke}%
\end{pgfscope}%
\begin{pgfscope}%
\pgfpathrectangle{\pgfqpoint{0.100000in}{2.493063in}}{\pgfqpoint{5.037500in}{3.427208in}}%
\pgfusepath{clip}%
\pgfsetbuttcap%
\pgfsetroundjoin%
\definecolor{currentfill}{rgb}{0.678431,1.000000,0.184314}%
\pgfsetfillcolor{currentfill}%
\pgfsetfillopacity{0.500000}%
\pgfsetlinewidth{0.250937pt}%
\definecolor{currentstroke}{rgb}{0.000000,0.000000,0.000000}%
\pgfsetstrokecolor{currentstroke}%
\pgfsetstrokeopacity{0.500000}%
\pgfsetdash{}{0pt}%
\pgfsys@defobject{currentmarker}{\pgfqpoint{-0.022222in}{-0.022222in}}{\pgfqpoint{0.022222in}{0.022222in}}{%
\pgfpathmoveto{\pgfqpoint{0.000000in}{-0.022222in}}%
\pgfpathcurveto{\pgfqpoint{0.005893in}{-0.022222in}}{\pgfqpoint{0.011546in}{-0.019881in}}{\pgfqpoint{0.015713in}{-0.015713in}}%
\pgfpathcurveto{\pgfqpoint{0.019881in}{-0.011546in}}{\pgfqpoint{0.022222in}{-0.005893in}}{\pgfqpoint{0.022222in}{0.000000in}}%
\pgfpathcurveto{\pgfqpoint{0.022222in}{0.005893in}}{\pgfqpoint{0.019881in}{0.011546in}}{\pgfqpoint{0.015713in}{0.015713in}}%
\pgfpathcurveto{\pgfqpoint{0.011546in}{0.019881in}}{\pgfqpoint{0.005893in}{0.022222in}}{\pgfqpoint{0.000000in}{0.022222in}}%
\pgfpathcurveto{\pgfqpoint{-0.005893in}{0.022222in}}{\pgfqpoint{-0.011546in}{0.019881in}}{\pgfqpoint{-0.015713in}{0.015713in}}%
\pgfpathcurveto{\pgfqpoint{-0.019881in}{0.011546in}}{\pgfqpoint{-0.022222in}{0.005893in}}{\pgfqpoint{-0.022222in}{0.000000in}}%
\pgfpathcurveto{\pgfqpoint{-0.022222in}{-0.005893in}}{\pgfqpoint{-0.019881in}{-0.011546in}}{\pgfqpoint{-0.015713in}{-0.015713in}}%
\pgfpathcurveto{\pgfqpoint{-0.011546in}{-0.019881in}}{\pgfqpoint{-0.005893in}{-0.022222in}}{\pgfqpoint{0.000000in}{-0.022222in}}%
\pgfpathlineto{\pgfqpoint{0.000000in}{-0.022222in}}%
\pgfpathclose%
\pgfusepath{stroke,fill}%
}%
\begin{pgfscope}%
\pgfsys@transformshift{4.406146in}{4.633374in}%
\pgfsys@useobject{currentmarker}{}%
\end{pgfscope}%
\end{pgfscope}%
\begin{pgfscope}%
\pgfpathrectangle{\pgfqpoint{0.100000in}{2.493063in}}{\pgfqpoint{5.037500in}{3.427208in}}%
\pgfusepath{clip}%
\pgfsetrectcap%
\pgfsetroundjoin%
\pgfsetlinewidth{1.505625pt}%
\definecolor{currentstroke}{rgb}{0.501961,0.501961,0.501961}%
\pgfsetstrokecolor{currentstroke}%
\pgfsetstrokeopacity{0.500000}%
\pgfsetdash{}{0pt}%
\pgfpathmoveto{\pgfqpoint{4.807749in}{4.947735in}}%
\pgfusepath{stroke}%
\end{pgfscope}%
\begin{pgfscope}%
\pgfpathrectangle{\pgfqpoint{0.100000in}{2.493063in}}{\pgfqpoint{5.037500in}{3.427208in}}%
\pgfusepath{clip}%
\pgfsetbuttcap%
\pgfsetroundjoin%
\definecolor{currentfill}{rgb}{0.501961,0.501961,0.501961}%
\pgfsetfillcolor{currentfill}%
\pgfsetfillopacity{0.500000}%
\pgfsetlinewidth{0.250937pt}%
\definecolor{currentstroke}{rgb}{0.000000,0.000000,0.000000}%
\pgfsetstrokecolor{currentstroke}%
\pgfsetstrokeopacity{0.500000}%
\pgfsetdash{}{0pt}%
\pgfsys@defobject{currentmarker}{\pgfqpoint{-0.013889in}{-0.013889in}}{\pgfqpoint{0.013889in}{0.013889in}}{%
\pgfpathmoveto{\pgfqpoint{0.000000in}{-0.013889in}}%
\pgfpathcurveto{\pgfqpoint{0.003683in}{-0.013889in}}{\pgfqpoint{0.007216in}{-0.012425in}}{\pgfqpoint{0.009821in}{-0.009821in}}%
\pgfpathcurveto{\pgfqpoint{0.012425in}{-0.007216in}}{\pgfqpoint{0.013889in}{-0.003683in}}{\pgfqpoint{0.013889in}{0.000000in}}%
\pgfpathcurveto{\pgfqpoint{0.013889in}{0.003683in}}{\pgfqpoint{0.012425in}{0.007216in}}{\pgfqpoint{0.009821in}{0.009821in}}%
\pgfpathcurveto{\pgfqpoint{0.007216in}{0.012425in}}{\pgfqpoint{0.003683in}{0.013889in}}{\pgfqpoint{0.000000in}{0.013889in}}%
\pgfpathcurveto{\pgfqpoint{-0.003683in}{0.013889in}}{\pgfqpoint{-0.007216in}{0.012425in}}{\pgfqpoint{-0.009821in}{0.009821in}}%
\pgfpathcurveto{\pgfqpoint{-0.012425in}{0.007216in}}{\pgfqpoint{-0.013889in}{0.003683in}}{\pgfqpoint{-0.013889in}{0.000000in}}%
\pgfpathcurveto{\pgfqpoint{-0.013889in}{-0.003683in}}{\pgfqpoint{-0.012425in}{-0.007216in}}{\pgfqpoint{-0.009821in}{-0.009821in}}%
\pgfpathcurveto{\pgfqpoint{-0.007216in}{-0.012425in}}{\pgfqpoint{-0.003683in}{-0.013889in}}{\pgfqpoint{0.000000in}{-0.013889in}}%
\pgfpathlineto{\pgfqpoint{0.000000in}{-0.013889in}}%
\pgfpathclose%
\pgfusepath{stroke,fill}%
}%
\begin{pgfscope}%
\pgfsys@transformshift{4.807749in}{4.947735in}%
\pgfsys@useobject{currentmarker}{}%
\end{pgfscope}%
\end{pgfscope}%
\begin{pgfscope}%
\pgfpathrectangle{\pgfqpoint{0.100000in}{2.493063in}}{\pgfqpoint{5.037500in}{3.427208in}}%
\pgfusepath{clip}%
\pgfsetrectcap%
\pgfsetroundjoin%
\pgfsetlinewidth{1.505625pt}%
\definecolor{currentstroke}{rgb}{0.000000,0.000000,1.000000}%
\pgfsetstrokecolor{currentstroke}%
\pgfsetstrokeopacity{0.500000}%
\pgfsetdash{}{0pt}%
\pgfpathmoveto{\pgfqpoint{4.263845in}{3.765403in}}%
\pgfusepath{stroke}%
\end{pgfscope}%
\begin{pgfscope}%
\pgfpathrectangle{\pgfqpoint{0.100000in}{2.493063in}}{\pgfqpoint{5.037500in}{3.427208in}}%
\pgfusepath{clip}%
\pgfsetbuttcap%
\pgfsetroundjoin%
\definecolor{currentfill}{rgb}{0.000000,0.000000,1.000000}%
\pgfsetfillcolor{currentfill}%
\pgfsetfillopacity{0.500000}%
\pgfsetlinewidth{0.250937pt}%
\definecolor{currentstroke}{rgb}{0.000000,0.000000,0.000000}%
\pgfsetstrokecolor{currentstroke}%
\pgfsetstrokeopacity{0.500000}%
\pgfsetdash{}{0pt}%
\pgfsys@defobject{currentmarker}{\pgfqpoint{-0.013889in}{-0.013889in}}{\pgfqpoint{0.013889in}{0.013889in}}{%
\pgfpathmoveto{\pgfqpoint{0.000000in}{-0.013889in}}%
\pgfpathcurveto{\pgfqpoint{0.003683in}{-0.013889in}}{\pgfqpoint{0.007216in}{-0.012425in}}{\pgfqpoint{0.009821in}{-0.009821in}}%
\pgfpathcurveto{\pgfqpoint{0.012425in}{-0.007216in}}{\pgfqpoint{0.013889in}{-0.003683in}}{\pgfqpoint{0.013889in}{0.000000in}}%
\pgfpathcurveto{\pgfqpoint{0.013889in}{0.003683in}}{\pgfqpoint{0.012425in}{0.007216in}}{\pgfqpoint{0.009821in}{0.009821in}}%
\pgfpathcurveto{\pgfqpoint{0.007216in}{0.012425in}}{\pgfqpoint{0.003683in}{0.013889in}}{\pgfqpoint{0.000000in}{0.013889in}}%
\pgfpathcurveto{\pgfqpoint{-0.003683in}{0.013889in}}{\pgfqpoint{-0.007216in}{0.012425in}}{\pgfqpoint{-0.009821in}{0.009821in}}%
\pgfpathcurveto{\pgfqpoint{-0.012425in}{0.007216in}}{\pgfqpoint{-0.013889in}{0.003683in}}{\pgfqpoint{-0.013889in}{0.000000in}}%
\pgfpathcurveto{\pgfqpoint{-0.013889in}{-0.003683in}}{\pgfqpoint{-0.012425in}{-0.007216in}}{\pgfqpoint{-0.009821in}{-0.009821in}}%
\pgfpathcurveto{\pgfqpoint{-0.007216in}{-0.012425in}}{\pgfqpoint{-0.003683in}{-0.013889in}}{\pgfqpoint{0.000000in}{-0.013889in}}%
\pgfpathlineto{\pgfqpoint{0.000000in}{-0.013889in}}%
\pgfpathclose%
\pgfusepath{stroke,fill}%
}%
\begin{pgfscope}%
\pgfsys@transformshift{4.263845in}{3.765403in}%
\pgfsys@useobject{currentmarker}{}%
\end{pgfscope}%
\end{pgfscope}%
\begin{pgfscope}%
\pgfpathrectangle{\pgfqpoint{0.100000in}{2.493063in}}{\pgfqpoint{5.037500in}{3.427208in}}%
\pgfusepath{clip}%
\pgfsetrectcap%
\pgfsetroundjoin%
\pgfsetlinewidth{1.505625pt}%
\definecolor{currentstroke}{rgb}{0.000000,0.000000,1.000000}%
\pgfsetstrokecolor{currentstroke}%
\pgfsetstrokeopacity{0.500000}%
\pgfsetdash{}{0pt}%
\pgfpathmoveto{\pgfqpoint{4.137182in}{3.887175in}}%
\pgfusepath{stroke}%
\end{pgfscope}%
\begin{pgfscope}%
\pgfpathrectangle{\pgfqpoint{0.100000in}{2.493063in}}{\pgfqpoint{5.037500in}{3.427208in}}%
\pgfusepath{clip}%
\pgfsetbuttcap%
\pgfsetroundjoin%
\definecolor{currentfill}{rgb}{0.000000,0.000000,1.000000}%
\pgfsetfillcolor{currentfill}%
\pgfsetfillopacity{0.500000}%
\pgfsetlinewidth{0.250937pt}%
\definecolor{currentstroke}{rgb}{0.000000,0.000000,0.000000}%
\pgfsetstrokecolor{currentstroke}%
\pgfsetstrokeopacity{0.500000}%
\pgfsetdash{}{0pt}%
\pgfsys@defobject{currentmarker}{\pgfqpoint{-0.016667in}{-0.016667in}}{\pgfqpoint{0.016667in}{0.016667in}}{%
\pgfpathmoveto{\pgfqpoint{0.000000in}{-0.016667in}}%
\pgfpathcurveto{\pgfqpoint{0.004420in}{-0.016667in}}{\pgfqpoint{0.008660in}{-0.014911in}}{\pgfqpoint{0.011785in}{-0.011785in}}%
\pgfpathcurveto{\pgfqpoint{0.014911in}{-0.008660in}}{\pgfqpoint{0.016667in}{-0.004420in}}{\pgfqpoint{0.016667in}{0.000000in}}%
\pgfpathcurveto{\pgfqpoint{0.016667in}{0.004420in}}{\pgfqpoint{0.014911in}{0.008660in}}{\pgfqpoint{0.011785in}{0.011785in}}%
\pgfpathcurveto{\pgfqpoint{0.008660in}{0.014911in}}{\pgfqpoint{0.004420in}{0.016667in}}{\pgfqpoint{0.000000in}{0.016667in}}%
\pgfpathcurveto{\pgfqpoint{-0.004420in}{0.016667in}}{\pgfqpoint{-0.008660in}{0.014911in}}{\pgfqpoint{-0.011785in}{0.011785in}}%
\pgfpathcurveto{\pgfqpoint{-0.014911in}{0.008660in}}{\pgfqpoint{-0.016667in}{0.004420in}}{\pgfqpoint{-0.016667in}{0.000000in}}%
\pgfpathcurveto{\pgfqpoint{-0.016667in}{-0.004420in}}{\pgfqpoint{-0.014911in}{-0.008660in}}{\pgfqpoint{-0.011785in}{-0.011785in}}%
\pgfpathcurveto{\pgfqpoint{-0.008660in}{-0.014911in}}{\pgfqpoint{-0.004420in}{-0.016667in}}{\pgfqpoint{0.000000in}{-0.016667in}}%
\pgfpathlineto{\pgfqpoint{0.000000in}{-0.016667in}}%
\pgfpathclose%
\pgfusepath{stroke,fill}%
}%
\begin{pgfscope}%
\pgfsys@transformshift{4.137182in}{3.887175in}%
\pgfsys@useobject{currentmarker}{}%
\end{pgfscope}%
\end{pgfscope}%
\begin{pgfscope}%
\pgfpathrectangle{\pgfqpoint{0.100000in}{2.493063in}}{\pgfqpoint{5.037500in}{3.427208in}}%
\pgfusepath{clip}%
\pgfsetrectcap%
\pgfsetroundjoin%
\pgfsetlinewidth{1.505625pt}%
\definecolor{currentstroke}{rgb}{0.000000,0.000000,1.000000}%
\pgfsetstrokecolor{currentstroke}%
\pgfsetstrokeopacity{0.500000}%
\pgfsetdash{}{0pt}%
\pgfpathmoveto{\pgfqpoint{4.253237in}{3.929127in}}%
\pgfusepath{stroke}%
\end{pgfscope}%
\begin{pgfscope}%
\pgfpathrectangle{\pgfqpoint{0.100000in}{2.493063in}}{\pgfqpoint{5.037500in}{3.427208in}}%
\pgfusepath{clip}%
\pgfsetbuttcap%
\pgfsetroundjoin%
\definecolor{currentfill}{rgb}{0.000000,0.000000,1.000000}%
\pgfsetfillcolor{currentfill}%
\pgfsetfillopacity{0.500000}%
\pgfsetlinewidth{0.250937pt}%
\definecolor{currentstroke}{rgb}{0.000000,0.000000,0.000000}%
\pgfsetstrokecolor{currentstroke}%
\pgfsetstrokeopacity{0.500000}%
\pgfsetdash{}{0pt}%
\pgfsys@defobject{currentmarker}{\pgfqpoint{-0.013889in}{-0.013889in}}{\pgfqpoint{0.013889in}{0.013889in}}{%
\pgfpathmoveto{\pgfqpoint{0.000000in}{-0.013889in}}%
\pgfpathcurveto{\pgfqpoint{0.003683in}{-0.013889in}}{\pgfqpoint{0.007216in}{-0.012425in}}{\pgfqpoint{0.009821in}{-0.009821in}}%
\pgfpathcurveto{\pgfqpoint{0.012425in}{-0.007216in}}{\pgfqpoint{0.013889in}{-0.003683in}}{\pgfqpoint{0.013889in}{0.000000in}}%
\pgfpathcurveto{\pgfqpoint{0.013889in}{0.003683in}}{\pgfqpoint{0.012425in}{0.007216in}}{\pgfqpoint{0.009821in}{0.009821in}}%
\pgfpathcurveto{\pgfqpoint{0.007216in}{0.012425in}}{\pgfqpoint{0.003683in}{0.013889in}}{\pgfqpoint{0.000000in}{0.013889in}}%
\pgfpathcurveto{\pgfqpoint{-0.003683in}{0.013889in}}{\pgfqpoint{-0.007216in}{0.012425in}}{\pgfqpoint{-0.009821in}{0.009821in}}%
\pgfpathcurveto{\pgfqpoint{-0.012425in}{0.007216in}}{\pgfqpoint{-0.013889in}{0.003683in}}{\pgfqpoint{-0.013889in}{0.000000in}}%
\pgfpathcurveto{\pgfqpoint{-0.013889in}{-0.003683in}}{\pgfqpoint{-0.012425in}{-0.007216in}}{\pgfqpoint{-0.009821in}{-0.009821in}}%
\pgfpathcurveto{\pgfqpoint{-0.007216in}{-0.012425in}}{\pgfqpoint{-0.003683in}{-0.013889in}}{\pgfqpoint{0.000000in}{-0.013889in}}%
\pgfpathlineto{\pgfqpoint{0.000000in}{-0.013889in}}%
\pgfpathclose%
\pgfusepath{stroke,fill}%
}%
\begin{pgfscope}%
\pgfsys@transformshift{4.253237in}{3.929127in}%
\pgfsys@useobject{currentmarker}{}%
\end{pgfscope}%
\end{pgfscope}%
\begin{pgfscope}%
\pgfpathrectangle{\pgfqpoint{0.100000in}{2.493063in}}{\pgfqpoint{5.037500in}{3.427208in}}%
\pgfusepath{clip}%
\pgfsetrectcap%
\pgfsetroundjoin%
\pgfsetlinewidth{1.505625pt}%
\definecolor{currentstroke}{rgb}{0.000000,0.000000,1.000000}%
\pgfsetstrokecolor{currentstroke}%
\pgfsetstrokeopacity{0.500000}%
\pgfsetdash{}{0pt}%
\pgfpathmoveto{\pgfqpoint{3.975994in}{3.922770in}}%
\pgfusepath{stroke}%
\end{pgfscope}%
\begin{pgfscope}%
\pgfpathrectangle{\pgfqpoint{0.100000in}{2.493063in}}{\pgfqpoint{5.037500in}{3.427208in}}%
\pgfusepath{clip}%
\pgfsetbuttcap%
\pgfsetroundjoin%
\definecolor{currentfill}{rgb}{0.000000,0.000000,1.000000}%
\pgfsetfillcolor{currentfill}%
\pgfsetfillopacity{0.500000}%
\pgfsetlinewidth{0.250937pt}%
\definecolor{currentstroke}{rgb}{0.000000,0.000000,0.000000}%
\pgfsetstrokecolor{currentstroke}%
\pgfsetstrokeopacity{0.500000}%
\pgfsetdash{}{0pt}%
\pgfsys@defobject{currentmarker}{\pgfqpoint{-0.016667in}{-0.016667in}}{\pgfqpoint{0.016667in}{0.016667in}}{%
\pgfpathmoveto{\pgfqpoint{0.000000in}{-0.016667in}}%
\pgfpathcurveto{\pgfqpoint{0.004420in}{-0.016667in}}{\pgfqpoint{0.008660in}{-0.014911in}}{\pgfqpoint{0.011785in}{-0.011785in}}%
\pgfpathcurveto{\pgfqpoint{0.014911in}{-0.008660in}}{\pgfqpoint{0.016667in}{-0.004420in}}{\pgfqpoint{0.016667in}{0.000000in}}%
\pgfpathcurveto{\pgfqpoint{0.016667in}{0.004420in}}{\pgfqpoint{0.014911in}{0.008660in}}{\pgfqpoint{0.011785in}{0.011785in}}%
\pgfpathcurveto{\pgfqpoint{0.008660in}{0.014911in}}{\pgfqpoint{0.004420in}{0.016667in}}{\pgfqpoint{0.000000in}{0.016667in}}%
\pgfpathcurveto{\pgfqpoint{-0.004420in}{0.016667in}}{\pgfqpoint{-0.008660in}{0.014911in}}{\pgfqpoint{-0.011785in}{0.011785in}}%
\pgfpathcurveto{\pgfqpoint{-0.014911in}{0.008660in}}{\pgfqpoint{-0.016667in}{0.004420in}}{\pgfqpoint{-0.016667in}{0.000000in}}%
\pgfpathcurveto{\pgfqpoint{-0.016667in}{-0.004420in}}{\pgfqpoint{-0.014911in}{-0.008660in}}{\pgfqpoint{-0.011785in}{-0.011785in}}%
\pgfpathcurveto{\pgfqpoint{-0.008660in}{-0.014911in}}{\pgfqpoint{-0.004420in}{-0.016667in}}{\pgfqpoint{0.000000in}{-0.016667in}}%
\pgfpathlineto{\pgfqpoint{0.000000in}{-0.016667in}}%
\pgfpathclose%
\pgfusepath{stroke,fill}%
}%
\begin{pgfscope}%
\pgfsys@transformshift{3.975994in}{3.922770in}%
\pgfsys@useobject{currentmarker}{}%
\end{pgfscope}%
\end{pgfscope}%
\begin{pgfscope}%
\pgfpathrectangle{\pgfqpoint{0.100000in}{2.493063in}}{\pgfqpoint{5.037500in}{3.427208in}}%
\pgfusepath{clip}%
\pgfsetrectcap%
\pgfsetroundjoin%
\pgfsetlinewidth{1.505625pt}%
\definecolor{currentstroke}{rgb}{0.000000,0.000000,1.000000}%
\pgfsetstrokecolor{currentstroke}%
\pgfsetstrokeopacity{0.500000}%
\pgfsetdash{}{0pt}%
\pgfpathmoveto{\pgfqpoint{3.994269in}{3.965455in}}%
\pgfusepath{stroke}%
\end{pgfscope}%
\begin{pgfscope}%
\pgfpathrectangle{\pgfqpoint{0.100000in}{2.493063in}}{\pgfqpoint{5.037500in}{3.427208in}}%
\pgfusepath{clip}%
\pgfsetbuttcap%
\pgfsetroundjoin%
\definecolor{currentfill}{rgb}{0.000000,0.000000,1.000000}%
\pgfsetfillcolor{currentfill}%
\pgfsetfillopacity{0.500000}%
\pgfsetlinewidth{0.250937pt}%
\definecolor{currentstroke}{rgb}{0.000000,0.000000,0.000000}%
\pgfsetstrokecolor{currentstroke}%
\pgfsetstrokeopacity{0.500000}%
\pgfsetdash{}{0pt}%
\pgfsys@defobject{currentmarker}{\pgfqpoint{-0.016667in}{-0.016667in}}{\pgfqpoint{0.016667in}{0.016667in}}{%
\pgfpathmoveto{\pgfqpoint{0.000000in}{-0.016667in}}%
\pgfpathcurveto{\pgfqpoint{0.004420in}{-0.016667in}}{\pgfqpoint{0.008660in}{-0.014911in}}{\pgfqpoint{0.011785in}{-0.011785in}}%
\pgfpathcurveto{\pgfqpoint{0.014911in}{-0.008660in}}{\pgfqpoint{0.016667in}{-0.004420in}}{\pgfqpoint{0.016667in}{0.000000in}}%
\pgfpathcurveto{\pgfqpoint{0.016667in}{0.004420in}}{\pgfqpoint{0.014911in}{0.008660in}}{\pgfqpoint{0.011785in}{0.011785in}}%
\pgfpathcurveto{\pgfqpoint{0.008660in}{0.014911in}}{\pgfqpoint{0.004420in}{0.016667in}}{\pgfqpoint{0.000000in}{0.016667in}}%
\pgfpathcurveto{\pgfqpoint{-0.004420in}{0.016667in}}{\pgfqpoint{-0.008660in}{0.014911in}}{\pgfqpoint{-0.011785in}{0.011785in}}%
\pgfpathcurveto{\pgfqpoint{-0.014911in}{0.008660in}}{\pgfqpoint{-0.016667in}{0.004420in}}{\pgfqpoint{-0.016667in}{0.000000in}}%
\pgfpathcurveto{\pgfqpoint{-0.016667in}{-0.004420in}}{\pgfqpoint{-0.014911in}{-0.008660in}}{\pgfqpoint{-0.011785in}{-0.011785in}}%
\pgfpathcurveto{\pgfqpoint{-0.008660in}{-0.014911in}}{\pgfqpoint{-0.004420in}{-0.016667in}}{\pgfqpoint{0.000000in}{-0.016667in}}%
\pgfpathlineto{\pgfqpoint{0.000000in}{-0.016667in}}%
\pgfpathclose%
\pgfusepath{stroke,fill}%
}%
\begin{pgfscope}%
\pgfsys@transformshift{3.994269in}{3.965455in}%
\pgfsys@useobject{currentmarker}{}%
\end{pgfscope}%
\end{pgfscope}%
\begin{pgfscope}%
\pgfpathrectangle{\pgfqpoint{0.100000in}{2.493063in}}{\pgfqpoint{5.037500in}{3.427208in}}%
\pgfusepath{clip}%
\pgfsetrectcap%
\pgfsetroundjoin%
\pgfsetlinewidth{1.505625pt}%
\definecolor{currentstroke}{rgb}{0.000000,0.000000,1.000000}%
\pgfsetstrokecolor{currentstroke}%
\pgfsetstrokeopacity{0.500000}%
\pgfsetdash{}{0pt}%
\pgfpathmoveto{\pgfqpoint{4.196108in}{3.687309in}}%
\pgfusepath{stroke}%
\end{pgfscope}%
\begin{pgfscope}%
\pgfpathrectangle{\pgfqpoint{0.100000in}{2.493063in}}{\pgfqpoint{5.037500in}{3.427208in}}%
\pgfusepath{clip}%
\pgfsetbuttcap%
\pgfsetroundjoin%
\definecolor{currentfill}{rgb}{0.000000,0.000000,1.000000}%
\pgfsetfillcolor{currentfill}%
\pgfsetfillopacity{0.500000}%
\pgfsetlinewidth{0.250937pt}%
\definecolor{currentstroke}{rgb}{0.000000,0.000000,0.000000}%
\pgfsetstrokecolor{currentstroke}%
\pgfsetstrokeopacity{0.500000}%
\pgfsetdash{}{0pt}%
\pgfsys@defobject{currentmarker}{\pgfqpoint{-0.019444in}{-0.019444in}}{\pgfqpoint{0.019444in}{0.019444in}}{%
\pgfpathmoveto{\pgfqpoint{0.000000in}{-0.019444in}}%
\pgfpathcurveto{\pgfqpoint{0.005157in}{-0.019444in}}{\pgfqpoint{0.010103in}{-0.017396in}}{\pgfqpoint{0.013749in}{-0.013749in}}%
\pgfpathcurveto{\pgfqpoint{0.017396in}{-0.010103in}}{\pgfqpoint{0.019444in}{-0.005157in}}{\pgfqpoint{0.019444in}{0.000000in}}%
\pgfpathcurveto{\pgfqpoint{0.019444in}{0.005157in}}{\pgfqpoint{0.017396in}{0.010103in}}{\pgfqpoint{0.013749in}{0.013749in}}%
\pgfpathcurveto{\pgfqpoint{0.010103in}{0.017396in}}{\pgfqpoint{0.005157in}{0.019444in}}{\pgfqpoint{0.000000in}{0.019444in}}%
\pgfpathcurveto{\pgfqpoint{-0.005157in}{0.019444in}}{\pgfqpoint{-0.010103in}{0.017396in}}{\pgfqpoint{-0.013749in}{0.013749in}}%
\pgfpathcurveto{\pgfqpoint{-0.017396in}{0.010103in}}{\pgfqpoint{-0.019444in}{0.005157in}}{\pgfqpoint{-0.019444in}{0.000000in}}%
\pgfpathcurveto{\pgfqpoint{-0.019444in}{-0.005157in}}{\pgfqpoint{-0.017396in}{-0.010103in}}{\pgfqpoint{-0.013749in}{-0.013749in}}%
\pgfpathcurveto{\pgfqpoint{-0.010103in}{-0.017396in}}{\pgfqpoint{-0.005157in}{-0.019444in}}{\pgfqpoint{0.000000in}{-0.019444in}}%
\pgfpathlineto{\pgfqpoint{0.000000in}{-0.019444in}}%
\pgfpathclose%
\pgfusepath{stroke,fill}%
}%
\begin{pgfscope}%
\pgfsys@transformshift{4.196108in}{3.687309in}%
\pgfsys@useobject{currentmarker}{}%
\end{pgfscope}%
\end{pgfscope}%
\begin{pgfscope}%
\pgfpathrectangle{\pgfqpoint{0.100000in}{2.493063in}}{\pgfqpoint{5.037500in}{3.427208in}}%
\pgfusepath{clip}%
\pgfsetrectcap%
\pgfsetroundjoin%
\pgfsetlinewidth{1.505625pt}%
\definecolor{currentstroke}{rgb}{0.501961,0.501961,0.501961}%
\pgfsetstrokecolor{currentstroke}%
\pgfsetstrokeopacity{0.500000}%
\pgfsetdash{}{0pt}%
\pgfpathmoveto{\pgfqpoint{4.346076in}{3.886299in}}%
\pgfusepath{stroke}%
\end{pgfscope}%
\begin{pgfscope}%
\pgfpathrectangle{\pgfqpoint{0.100000in}{2.493063in}}{\pgfqpoint{5.037500in}{3.427208in}}%
\pgfusepath{clip}%
\pgfsetbuttcap%
\pgfsetroundjoin%
\definecolor{currentfill}{rgb}{0.501961,0.501961,0.501961}%
\pgfsetfillcolor{currentfill}%
\pgfsetfillopacity{0.500000}%
\pgfsetlinewidth{0.250937pt}%
\definecolor{currentstroke}{rgb}{0.000000,0.000000,0.000000}%
\pgfsetstrokecolor{currentstroke}%
\pgfsetstrokeopacity{0.500000}%
\pgfsetdash{}{0pt}%
\pgfsys@defobject{currentmarker}{\pgfqpoint{-0.013889in}{-0.013889in}}{\pgfqpoint{0.013889in}{0.013889in}}{%
\pgfpathmoveto{\pgfqpoint{0.000000in}{-0.013889in}}%
\pgfpathcurveto{\pgfqpoint{0.003683in}{-0.013889in}}{\pgfqpoint{0.007216in}{-0.012425in}}{\pgfqpoint{0.009821in}{-0.009821in}}%
\pgfpathcurveto{\pgfqpoint{0.012425in}{-0.007216in}}{\pgfqpoint{0.013889in}{-0.003683in}}{\pgfqpoint{0.013889in}{0.000000in}}%
\pgfpathcurveto{\pgfqpoint{0.013889in}{0.003683in}}{\pgfqpoint{0.012425in}{0.007216in}}{\pgfqpoint{0.009821in}{0.009821in}}%
\pgfpathcurveto{\pgfqpoint{0.007216in}{0.012425in}}{\pgfqpoint{0.003683in}{0.013889in}}{\pgfqpoint{0.000000in}{0.013889in}}%
\pgfpathcurveto{\pgfqpoint{-0.003683in}{0.013889in}}{\pgfqpoint{-0.007216in}{0.012425in}}{\pgfqpoint{-0.009821in}{0.009821in}}%
\pgfpathcurveto{\pgfqpoint{-0.012425in}{0.007216in}}{\pgfqpoint{-0.013889in}{0.003683in}}{\pgfqpoint{-0.013889in}{0.000000in}}%
\pgfpathcurveto{\pgfqpoint{-0.013889in}{-0.003683in}}{\pgfqpoint{-0.012425in}{-0.007216in}}{\pgfqpoint{-0.009821in}{-0.009821in}}%
\pgfpathcurveto{\pgfqpoint{-0.007216in}{-0.012425in}}{\pgfqpoint{-0.003683in}{-0.013889in}}{\pgfqpoint{0.000000in}{-0.013889in}}%
\pgfpathlineto{\pgfqpoint{0.000000in}{-0.013889in}}%
\pgfpathclose%
\pgfusepath{stroke,fill}%
}%
\begin{pgfscope}%
\pgfsys@transformshift{4.346076in}{3.886299in}%
\pgfsys@useobject{currentmarker}{}%
\end{pgfscope}%
\end{pgfscope}%
\begin{pgfscope}%
\pgfpathrectangle{\pgfqpoint{0.100000in}{2.493063in}}{\pgfqpoint{5.037500in}{3.427208in}}%
\pgfusepath{clip}%
\pgfsetrectcap%
\pgfsetroundjoin%
\pgfsetlinewidth{1.505625pt}%
\definecolor{currentstroke}{rgb}{0.000000,0.000000,1.000000}%
\pgfsetstrokecolor{currentstroke}%
\pgfsetstrokeopacity{0.500000}%
\pgfsetdash{}{0pt}%
\pgfpathmoveto{\pgfqpoint{4.036398in}{3.982925in}}%
\pgfusepath{stroke}%
\end{pgfscope}%
\begin{pgfscope}%
\pgfpathrectangle{\pgfqpoint{0.100000in}{2.493063in}}{\pgfqpoint{5.037500in}{3.427208in}}%
\pgfusepath{clip}%
\pgfsetbuttcap%
\pgfsetroundjoin%
\definecolor{currentfill}{rgb}{0.000000,0.000000,1.000000}%
\pgfsetfillcolor{currentfill}%
\pgfsetfillopacity{0.500000}%
\pgfsetlinewidth{0.250937pt}%
\definecolor{currentstroke}{rgb}{0.000000,0.000000,0.000000}%
\pgfsetstrokecolor{currentstroke}%
\pgfsetstrokeopacity{0.500000}%
\pgfsetdash{}{0pt}%
\pgfsys@defobject{currentmarker}{\pgfqpoint{-0.019444in}{-0.019444in}}{\pgfqpoint{0.019444in}{0.019444in}}{%
\pgfpathmoveto{\pgfqpoint{0.000000in}{-0.019444in}}%
\pgfpathcurveto{\pgfqpoint{0.005157in}{-0.019444in}}{\pgfqpoint{0.010103in}{-0.017396in}}{\pgfqpoint{0.013749in}{-0.013749in}}%
\pgfpathcurveto{\pgfqpoint{0.017396in}{-0.010103in}}{\pgfqpoint{0.019444in}{-0.005157in}}{\pgfqpoint{0.019444in}{0.000000in}}%
\pgfpathcurveto{\pgfqpoint{0.019444in}{0.005157in}}{\pgfqpoint{0.017396in}{0.010103in}}{\pgfqpoint{0.013749in}{0.013749in}}%
\pgfpathcurveto{\pgfqpoint{0.010103in}{0.017396in}}{\pgfqpoint{0.005157in}{0.019444in}}{\pgfqpoint{0.000000in}{0.019444in}}%
\pgfpathcurveto{\pgfqpoint{-0.005157in}{0.019444in}}{\pgfqpoint{-0.010103in}{0.017396in}}{\pgfqpoint{-0.013749in}{0.013749in}}%
\pgfpathcurveto{\pgfqpoint{-0.017396in}{0.010103in}}{\pgfqpoint{-0.019444in}{0.005157in}}{\pgfqpoint{-0.019444in}{0.000000in}}%
\pgfpathcurveto{\pgfqpoint{-0.019444in}{-0.005157in}}{\pgfqpoint{-0.017396in}{-0.010103in}}{\pgfqpoint{-0.013749in}{-0.013749in}}%
\pgfpathcurveto{\pgfqpoint{-0.010103in}{-0.017396in}}{\pgfqpoint{-0.005157in}{-0.019444in}}{\pgfqpoint{0.000000in}{-0.019444in}}%
\pgfpathlineto{\pgfqpoint{0.000000in}{-0.019444in}}%
\pgfpathclose%
\pgfusepath{stroke,fill}%
}%
\begin{pgfscope}%
\pgfsys@transformshift{4.036398in}{3.982925in}%
\pgfsys@useobject{currentmarker}{}%
\end{pgfscope}%
\end{pgfscope}%
\begin{pgfscope}%
\pgfpathrectangle{\pgfqpoint{0.100000in}{2.493063in}}{\pgfqpoint{5.037500in}{3.427208in}}%
\pgfusepath{clip}%
\pgfsetrectcap%
\pgfsetroundjoin%
\pgfsetlinewidth{1.505625pt}%
\definecolor{currentstroke}{rgb}{0.000000,0.000000,1.000000}%
\pgfsetstrokecolor{currentstroke}%
\pgfsetstrokeopacity{0.500000}%
\pgfsetdash{}{0pt}%
\pgfpathmoveto{\pgfqpoint{4.204247in}{3.888425in}}%
\pgfusepath{stroke}%
\end{pgfscope}%
\begin{pgfscope}%
\pgfpathrectangle{\pgfqpoint{0.100000in}{2.493063in}}{\pgfqpoint{5.037500in}{3.427208in}}%
\pgfusepath{clip}%
\pgfsetbuttcap%
\pgfsetroundjoin%
\definecolor{currentfill}{rgb}{0.000000,0.000000,1.000000}%
\pgfsetfillcolor{currentfill}%
\pgfsetfillopacity{0.500000}%
\pgfsetlinewidth{0.250937pt}%
\definecolor{currentstroke}{rgb}{0.000000,0.000000,0.000000}%
\pgfsetstrokecolor{currentstroke}%
\pgfsetstrokeopacity{0.500000}%
\pgfsetdash{}{0pt}%
\pgfsys@defobject{currentmarker}{\pgfqpoint{-0.019444in}{-0.019444in}}{\pgfqpoint{0.019444in}{0.019444in}}{%
\pgfpathmoveto{\pgfqpoint{0.000000in}{-0.019444in}}%
\pgfpathcurveto{\pgfqpoint{0.005157in}{-0.019444in}}{\pgfqpoint{0.010103in}{-0.017396in}}{\pgfqpoint{0.013749in}{-0.013749in}}%
\pgfpathcurveto{\pgfqpoint{0.017396in}{-0.010103in}}{\pgfqpoint{0.019444in}{-0.005157in}}{\pgfqpoint{0.019444in}{0.000000in}}%
\pgfpathcurveto{\pgfqpoint{0.019444in}{0.005157in}}{\pgfqpoint{0.017396in}{0.010103in}}{\pgfqpoint{0.013749in}{0.013749in}}%
\pgfpathcurveto{\pgfqpoint{0.010103in}{0.017396in}}{\pgfqpoint{0.005157in}{0.019444in}}{\pgfqpoint{0.000000in}{0.019444in}}%
\pgfpathcurveto{\pgfqpoint{-0.005157in}{0.019444in}}{\pgfqpoint{-0.010103in}{0.017396in}}{\pgfqpoint{-0.013749in}{0.013749in}}%
\pgfpathcurveto{\pgfqpoint{-0.017396in}{0.010103in}}{\pgfqpoint{-0.019444in}{0.005157in}}{\pgfqpoint{-0.019444in}{0.000000in}}%
\pgfpathcurveto{\pgfqpoint{-0.019444in}{-0.005157in}}{\pgfqpoint{-0.017396in}{-0.010103in}}{\pgfqpoint{-0.013749in}{-0.013749in}}%
\pgfpathcurveto{\pgfqpoint{-0.010103in}{-0.017396in}}{\pgfqpoint{-0.005157in}{-0.019444in}}{\pgfqpoint{0.000000in}{-0.019444in}}%
\pgfpathlineto{\pgfqpoint{0.000000in}{-0.019444in}}%
\pgfpathclose%
\pgfusepath{stroke,fill}%
}%
\begin{pgfscope}%
\pgfsys@transformshift{4.204247in}{3.888425in}%
\pgfsys@useobject{currentmarker}{}%
\end{pgfscope}%
\end{pgfscope}%
\begin{pgfscope}%
\pgfpathrectangle{\pgfqpoint{0.100000in}{2.493063in}}{\pgfqpoint{5.037500in}{3.427208in}}%
\pgfusepath{clip}%
\pgfsetrectcap%
\pgfsetroundjoin%
\pgfsetlinewidth{1.505625pt}%
\definecolor{currentstroke}{rgb}{0.678431,1.000000,0.184314}%
\pgfsetstrokecolor{currentstroke}%
\pgfsetstrokeopacity{0.500000}%
\pgfsetdash{}{0pt}%
\pgfpathmoveto{\pgfqpoint{2.121809in}{4.977686in}}%
\pgfusepath{stroke}%
\end{pgfscope}%
\begin{pgfscope}%
\pgfpathrectangle{\pgfqpoint{0.100000in}{2.493063in}}{\pgfqpoint{5.037500in}{3.427208in}}%
\pgfusepath{clip}%
\pgfsetbuttcap%
\pgfsetroundjoin%
\definecolor{currentfill}{rgb}{0.678431,1.000000,0.184314}%
\pgfsetfillcolor{currentfill}%
\pgfsetfillopacity{0.500000}%
\pgfsetlinewidth{0.250937pt}%
\definecolor{currentstroke}{rgb}{0.000000,0.000000,0.000000}%
\pgfsetstrokecolor{currentstroke}%
\pgfsetstrokeopacity{0.500000}%
\pgfsetdash{}{0pt}%
\pgfsys@defobject{currentmarker}{\pgfqpoint{-0.025000in}{-0.025000in}}{\pgfqpoint{0.025000in}{0.025000in}}{%
\pgfpathmoveto{\pgfqpoint{0.000000in}{-0.025000in}}%
\pgfpathcurveto{\pgfqpoint{0.006630in}{-0.025000in}}{\pgfqpoint{0.012989in}{-0.022366in}}{\pgfqpoint{0.017678in}{-0.017678in}}%
\pgfpathcurveto{\pgfqpoint{0.022366in}{-0.012989in}}{\pgfqpoint{0.025000in}{-0.006630in}}{\pgfqpoint{0.025000in}{0.000000in}}%
\pgfpathcurveto{\pgfqpoint{0.025000in}{0.006630in}}{\pgfqpoint{0.022366in}{0.012989in}}{\pgfqpoint{0.017678in}{0.017678in}}%
\pgfpathcurveto{\pgfqpoint{0.012989in}{0.022366in}}{\pgfqpoint{0.006630in}{0.025000in}}{\pgfqpoint{0.000000in}{0.025000in}}%
\pgfpathcurveto{\pgfqpoint{-0.006630in}{0.025000in}}{\pgfqpoint{-0.012989in}{0.022366in}}{\pgfqpoint{-0.017678in}{0.017678in}}%
\pgfpathcurveto{\pgfqpoint{-0.022366in}{0.012989in}}{\pgfqpoint{-0.025000in}{0.006630in}}{\pgfqpoint{-0.025000in}{0.000000in}}%
\pgfpathcurveto{\pgfqpoint{-0.025000in}{-0.006630in}}{\pgfqpoint{-0.022366in}{-0.012989in}}{\pgfqpoint{-0.017678in}{-0.017678in}}%
\pgfpathcurveto{\pgfqpoint{-0.012989in}{-0.022366in}}{\pgfqpoint{-0.006630in}{-0.025000in}}{\pgfqpoint{0.000000in}{-0.025000in}}%
\pgfpathlineto{\pgfqpoint{0.000000in}{-0.025000in}}%
\pgfpathclose%
\pgfusepath{stroke,fill}%
}%
\begin{pgfscope}%
\pgfsys@transformshift{2.121809in}{4.977686in}%
\pgfsys@useobject{currentmarker}{}%
\end{pgfscope}%
\end{pgfscope}%
\begin{pgfscope}%
\pgfpathrectangle{\pgfqpoint{0.100000in}{2.493063in}}{\pgfqpoint{5.037500in}{3.427208in}}%
\pgfusepath{clip}%
\pgfsetrectcap%
\pgfsetroundjoin%
\pgfsetlinewidth{1.505625pt}%
\definecolor{currentstroke}{rgb}{0.678431,1.000000,0.184314}%
\pgfsetstrokecolor{currentstroke}%
\pgfsetstrokeopacity{0.500000}%
\pgfsetdash{}{0pt}%
\pgfpathmoveto{\pgfqpoint{2.661793in}{4.886727in}}%
\pgfusepath{stroke}%
\end{pgfscope}%
\begin{pgfscope}%
\pgfpathrectangle{\pgfqpoint{0.100000in}{2.493063in}}{\pgfqpoint{5.037500in}{3.427208in}}%
\pgfusepath{clip}%
\pgfsetbuttcap%
\pgfsetroundjoin%
\definecolor{currentfill}{rgb}{0.678431,1.000000,0.184314}%
\pgfsetfillcolor{currentfill}%
\pgfsetfillopacity{0.500000}%
\pgfsetlinewidth{0.250937pt}%
\definecolor{currentstroke}{rgb}{0.000000,0.000000,0.000000}%
\pgfsetstrokecolor{currentstroke}%
\pgfsetstrokeopacity{0.500000}%
\pgfsetdash{}{0pt}%
\pgfsys@defobject{currentmarker}{\pgfqpoint{-0.016667in}{-0.016667in}}{\pgfqpoint{0.016667in}{0.016667in}}{%
\pgfpathmoveto{\pgfqpoint{0.000000in}{-0.016667in}}%
\pgfpathcurveto{\pgfqpoint{0.004420in}{-0.016667in}}{\pgfqpoint{0.008660in}{-0.014911in}}{\pgfqpoint{0.011785in}{-0.011785in}}%
\pgfpathcurveto{\pgfqpoint{0.014911in}{-0.008660in}}{\pgfqpoint{0.016667in}{-0.004420in}}{\pgfqpoint{0.016667in}{0.000000in}}%
\pgfpathcurveto{\pgfqpoint{0.016667in}{0.004420in}}{\pgfqpoint{0.014911in}{0.008660in}}{\pgfqpoint{0.011785in}{0.011785in}}%
\pgfpathcurveto{\pgfqpoint{0.008660in}{0.014911in}}{\pgfqpoint{0.004420in}{0.016667in}}{\pgfqpoint{0.000000in}{0.016667in}}%
\pgfpathcurveto{\pgfqpoint{-0.004420in}{0.016667in}}{\pgfqpoint{-0.008660in}{0.014911in}}{\pgfqpoint{-0.011785in}{0.011785in}}%
\pgfpathcurveto{\pgfqpoint{-0.014911in}{0.008660in}}{\pgfqpoint{-0.016667in}{0.004420in}}{\pgfqpoint{-0.016667in}{0.000000in}}%
\pgfpathcurveto{\pgfqpoint{-0.016667in}{-0.004420in}}{\pgfqpoint{-0.014911in}{-0.008660in}}{\pgfqpoint{-0.011785in}{-0.011785in}}%
\pgfpathcurveto{\pgfqpoint{-0.008660in}{-0.014911in}}{\pgfqpoint{-0.004420in}{-0.016667in}}{\pgfqpoint{0.000000in}{-0.016667in}}%
\pgfpathlineto{\pgfqpoint{0.000000in}{-0.016667in}}%
\pgfpathclose%
\pgfusepath{stroke,fill}%
}%
\begin{pgfscope}%
\pgfsys@transformshift{2.661793in}{4.886727in}%
\pgfsys@useobject{currentmarker}{}%
\end{pgfscope}%
\end{pgfscope}%
\begin{pgfscope}%
\pgfpathrectangle{\pgfqpoint{0.100000in}{2.493063in}}{\pgfqpoint{5.037500in}{3.427208in}}%
\pgfusepath{clip}%
\pgfsetrectcap%
\pgfsetroundjoin%
\pgfsetlinewidth{1.505625pt}%
\definecolor{currentstroke}{rgb}{0.000000,0.000000,1.000000}%
\pgfsetstrokecolor{currentstroke}%
\pgfsetstrokeopacity{0.500000}%
\pgfsetdash{}{0pt}%
\pgfpathmoveto{\pgfqpoint{3.718137in}{3.954012in}}%
\pgfusepath{stroke}%
\end{pgfscope}%
\begin{pgfscope}%
\pgfpathrectangle{\pgfqpoint{0.100000in}{2.493063in}}{\pgfqpoint{5.037500in}{3.427208in}}%
\pgfusepath{clip}%
\pgfsetbuttcap%
\pgfsetroundjoin%
\definecolor{currentfill}{rgb}{0.000000,0.000000,1.000000}%
\pgfsetfillcolor{currentfill}%
\pgfsetfillopacity{0.500000}%
\pgfsetlinewidth{0.250937pt}%
\definecolor{currentstroke}{rgb}{0.000000,0.000000,0.000000}%
\pgfsetstrokecolor{currentstroke}%
\pgfsetstrokeopacity{0.500000}%
\pgfsetdash{}{0pt}%
\pgfsys@defobject{currentmarker}{\pgfqpoint{-0.005556in}{-0.005556in}}{\pgfqpoint{0.005556in}{0.005556in}}{%
\pgfpathmoveto{\pgfqpoint{0.000000in}{-0.005556in}}%
\pgfpathcurveto{\pgfqpoint{0.001473in}{-0.005556in}}{\pgfqpoint{0.002887in}{-0.004970in}}{\pgfqpoint{0.003928in}{-0.003928in}}%
\pgfpathcurveto{\pgfqpoint{0.004970in}{-0.002887in}}{\pgfqpoint{0.005556in}{-0.001473in}}{\pgfqpoint{0.005556in}{0.000000in}}%
\pgfpathcurveto{\pgfqpoint{0.005556in}{0.001473in}}{\pgfqpoint{0.004970in}{0.002887in}}{\pgfqpoint{0.003928in}{0.003928in}}%
\pgfpathcurveto{\pgfqpoint{0.002887in}{0.004970in}}{\pgfqpoint{0.001473in}{0.005556in}}{\pgfqpoint{0.000000in}{0.005556in}}%
\pgfpathcurveto{\pgfqpoint{-0.001473in}{0.005556in}}{\pgfqpoint{-0.002887in}{0.004970in}}{\pgfqpoint{-0.003928in}{0.003928in}}%
\pgfpathcurveto{\pgfqpoint{-0.004970in}{0.002887in}}{\pgfqpoint{-0.005556in}{0.001473in}}{\pgfqpoint{-0.005556in}{0.000000in}}%
\pgfpathcurveto{\pgfqpoint{-0.005556in}{-0.001473in}}{\pgfqpoint{-0.004970in}{-0.002887in}}{\pgfqpoint{-0.003928in}{-0.003928in}}%
\pgfpathcurveto{\pgfqpoint{-0.002887in}{-0.004970in}}{\pgfqpoint{-0.001473in}{-0.005556in}}{\pgfqpoint{0.000000in}{-0.005556in}}%
\pgfpathlineto{\pgfqpoint{0.000000in}{-0.005556in}}%
\pgfpathclose%
\pgfusepath{stroke,fill}%
}%
\begin{pgfscope}%
\pgfsys@transformshift{3.718137in}{3.954012in}%
\pgfsys@useobject{currentmarker}{}%
\end{pgfscope}%
\end{pgfscope}%
\begin{pgfscope}%
\pgfpathrectangle{\pgfqpoint{0.100000in}{2.493063in}}{\pgfqpoint{5.037500in}{3.427208in}}%
\pgfusepath{clip}%
\pgfsetrectcap%
\pgfsetroundjoin%
\pgfsetlinewidth{1.505625pt}%
\definecolor{currentstroke}{rgb}{0.000000,0.000000,1.000000}%
\pgfsetstrokecolor{currentstroke}%
\pgfsetstrokeopacity{0.500000}%
\pgfsetdash{}{0pt}%
\pgfpathmoveto{\pgfqpoint{3.511014in}{4.105917in}}%
\pgfusepath{stroke}%
\end{pgfscope}%
\begin{pgfscope}%
\pgfpathrectangle{\pgfqpoint{0.100000in}{2.493063in}}{\pgfqpoint{5.037500in}{3.427208in}}%
\pgfusepath{clip}%
\pgfsetbuttcap%
\pgfsetroundjoin%
\definecolor{currentfill}{rgb}{0.000000,0.000000,1.000000}%
\pgfsetfillcolor{currentfill}%
\pgfsetfillopacity{0.500000}%
\pgfsetlinewidth{0.250937pt}%
\definecolor{currentstroke}{rgb}{0.000000,0.000000,0.000000}%
\pgfsetstrokecolor{currentstroke}%
\pgfsetstrokeopacity{0.500000}%
\pgfsetdash{}{0pt}%
\pgfsys@defobject{currentmarker}{\pgfqpoint{-0.011111in}{-0.011111in}}{\pgfqpoint{0.011111in}{0.011111in}}{%
\pgfpathmoveto{\pgfqpoint{0.000000in}{-0.011111in}}%
\pgfpathcurveto{\pgfqpoint{0.002947in}{-0.011111in}}{\pgfqpoint{0.005773in}{-0.009940in}}{\pgfqpoint{0.007857in}{-0.007857in}}%
\pgfpathcurveto{\pgfqpoint{0.009940in}{-0.005773in}}{\pgfqpoint{0.011111in}{-0.002947in}}{\pgfqpoint{0.011111in}{0.000000in}}%
\pgfpathcurveto{\pgfqpoint{0.011111in}{0.002947in}}{\pgfqpoint{0.009940in}{0.005773in}}{\pgfqpoint{0.007857in}{0.007857in}}%
\pgfpathcurveto{\pgfqpoint{0.005773in}{0.009940in}}{\pgfqpoint{0.002947in}{0.011111in}}{\pgfqpoint{0.000000in}{0.011111in}}%
\pgfpathcurveto{\pgfqpoint{-0.002947in}{0.011111in}}{\pgfqpoint{-0.005773in}{0.009940in}}{\pgfqpoint{-0.007857in}{0.007857in}}%
\pgfpathcurveto{\pgfqpoint{-0.009940in}{0.005773in}}{\pgfqpoint{-0.011111in}{0.002947in}}{\pgfqpoint{-0.011111in}{0.000000in}}%
\pgfpathcurveto{\pgfqpoint{-0.011111in}{-0.002947in}}{\pgfqpoint{-0.009940in}{-0.005773in}}{\pgfqpoint{-0.007857in}{-0.007857in}}%
\pgfpathcurveto{\pgfqpoint{-0.005773in}{-0.009940in}}{\pgfqpoint{-0.002947in}{-0.011111in}}{\pgfqpoint{0.000000in}{-0.011111in}}%
\pgfpathlineto{\pgfqpoint{0.000000in}{-0.011111in}}%
\pgfpathclose%
\pgfusepath{stroke,fill}%
}%
\begin{pgfscope}%
\pgfsys@transformshift{3.511014in}{4.105917in}%
\pgfsys@useobject{currentmarker}{}%
\end{pgfscope}%
\end{pgfscope}%
\begin{pgfscope}%
\pgfpathrectangle{\pgfqpoint{0.100000in}{2.493063in}}{\pgfqpoint{5.037500in}{3.427208in}}%
\pgfusepath{clip}%
\pgfsetrectcap%
\pgfsetroundjoin%
\pgfsetlinewidth{1.505625pt}%
\definecolor{currentstroke}{rgb}{0.000000,0.000000,1.000000}%
\pgfsetstrokecolor{currentstroke}%
\pgfsetstrokeopacity{0.500000}%
\pgfsetdash{}{0pt}%
\pgfpathmoveto{\pgfqpoint{3.757359in}{3.971528in}}%
\pgfusepath{stroke}%
\end{pgfscope}%
\begin{pgfscope}%
\pgfpathrectangle{\pgfqpoint{0.100000in}{2.493063in}}{\pgfqpoint{5.037500in}{3.427208in}}%
\pgfusepath{clip}%
\pgfsetbuttcap%
\pgfsetroundjoin%
\definecolor{currentfill}{rgb}{0.000000,0.000000,1.000000}%
\pgfsetfillcolor{currentfill}%
\pgfsetfillopacity{0.500000}%
\pgfsetlinewidth{0.250937pt}%
\definecolor{currentstroke}{rgb}{0.000000,0.000000,0.000000}%
\pgfsetstrokecolor{currentstroke}%
\pgfsetstrokeopacity{0.500000}%
\pgfsetdash{}{0pt}%
\pgfsys@defobject{currentmarker}{\pgfqpoint{-0.011111in}{-0.011111in}}{\pgfqpoint{0.011111in}{0.011111in}}{%
\pgfpathmoveto{\pgfqpoint{0.000000in}{-0.011111in}}%
\pgfpathcurveto{\pgfqpoint{0.002947in}{-0.011111in}}{\pgfqpoint{0.005773in}{-0.009940in}}{\pgfqpoint{0.007857in}{-0.007857in}}%
\pgfpathcurveto{\pgfqpoint{0.009940in}{-0.005773in}}{\pgfqpoint{0.011111in}{-0.002947in}}{\pgfqpoint{0.011111in}{0.000000in}}%
\pgfpathcurveto{\pgfqpoint{0.011111in}{0.002947in}}{\pgfqpoint{0.009940in}{0.005773in}}{\pgfqpoint{0.007857in}{0.007857in}}%
\pgfpathcurveto{\pgfqpoint{0.005773in}{0.009940in}}{\pgfqpoint{0.002947in}{0.011111in}}{\pgfqpoint{0.000000in}{0.011111in}}%
\pgfpathcurveto{\pgfqpoint{-0.002947in}{0.011111in}}{\pgfqpoint{-0.005773in}{0.009940in}}{\pgfqpoint{-0.007857in}{0.007857in}}%
\pgfpathcurveto{\pgfqpoint{-0.009940in}{0.005773in}}{\pgfqpoint{-0.011111in}{0.002947in}}{\pgfqpoint{-0.011111in}{0.000000in}}%
\pgfpathcurveto{\pgfqpoint{-0.011111in}{-0.002947in}}{\pgfqpoint{-0.009940in}{-0.005773in}}{\pgfqpoint{-0.007857in}{-0.007857in}}%
\pgfpathcurveto{\pgfqpoint{-0.005773in}{-0.009940in}}{\pgfqpoint{-0.002947in}{-0.011111in}}{\pgfqpoint{0.000000in}{-0.011111in}}%
\pgfpathlineto{\pgfqpoint{0.000000in}{-0.011111in}}%
\pgfpathclose%
\pgfusepath{stroke,fill}%
}%
\begin{pgfscope}%
\pgfsys@transformshift{3.757359in}{3.971528in}%
\pgfsys@useobject{currentmarker}{}%
\end{pgfscope}%
\end{pgfscope}%
\begin{pgfscope}%
\pgfpathrectangle{\pgfqpoint{0.100000in}{2.493063in}}{\pgfqpoint{5.037500in}{3.427208in}}%
\pgfusepath{clip}%
\pgfsetrectcap%
\pgfsetroundjoin%
\pgfsetlinewidth{1.505625pt}%
\definecolor{currentstroke}{rgb}{0.501961,0.501961,0.501961}%
\pgfsetstrokecolor{currentstroke}%
\pgfsetstrokeopacity{0.500000}%
\pgfsetdash{}{0pt}%
\pgfpathmoveto{\pgfqpoint{3.383445in}{3.990583in}}%
\pgfusepath{stroke}%
\end{pgfscope}%
\begin{pgfscope}%
\pgfpathrectangle{\pgfqpoint{0.100000in}{2.493063in}}{\pgfqpoint{5.037500in}{3.427208in}}%
\pgfusepath{clip}%
\pgfsetbuttcap%
\pgfsetroundjoin%
\definecolor{currentfill}{rgb}{0.501961,0.501961,0.501961}%
\pgfsetfillcolor{currentfill}%
\pgfsetfillopacity{0.500000}%
\pgfsetlinewidth{0.250937pt}%
\definecolor{currentstroke}{rgb}{0.000000,0.000000,0.000000}%
\pgfsetstrokecolor{currentstroke}%
\pgfsetstrokeopacity{0.500000}%
\pgfsetdash{}{0pt}%
\pgfsys@defobject{currentmarker}{\pgfqpoint{-0.013889in}{-0.013889in}}{\pgfqpoint{0.013889in}{0.013889in}}{%
\pgfpathmoveto{\pgfqpoint{0.000000in}{-0.013889in}}%
\pgfpathcurveto{\pgfqpoint{0.003683in}{-0.013889in}}{\pgfqpoint{0.007216in}{-0.012425in}}{\pgfqpoint{0.009821in}{-0.009821in}}%
\pgfpathcurveto{\pgfqpoint{0.012425in}{-0.007216in}}{\pgfqpoint{0.013889in}{-0.003683in}}{\pgfqpoint{0.013889in}{0.000000in}}%
\pgfpathcurveto{\pgfqpoint{0.013889in}{0.003683in}}{\pgfqpoint{0.012425in}{0.007216in}}{\pgfqpoint{0.009821in}{0.009821in}}%
\pgfpathcurveto{\pgfqpoint{0.007216in}{0.012425in}}{\pgfqpoint{0.003683in}{0.013889in}}{\pgfqpoint{0.000000in}{0.013889in}}%
\pgfpathcurveto{\pgfqpoint{-0.003683in}{0.013889in}}{\pgfqpoint{-0.007216in}{0.012425in}}{\pgfqpoint{-0.009821in}{0.009821in}}%
\pgfpathcurveto{\pgfqpoint{-0.012425in}{0.007216in}}{\pgfqpoint{-0.013889in}{0.003683in}}{\pgfqpoint{-0.013889in}{0.000000in}}%
\pgfpathcurveto{\pgfqpoint{-0.013889in}{-0.003683in}}{\pgfqpoint{-0.012425in}{-0.007216in}}{\pgfqpoint{-0.009821in}{-0.009821in}}%
\pgfpathcurveto{\pgfqpoint{-0.007216in}{-0.012425in}}{\pgfqpoint{-0.003683in}{-0.013889in}}{\pgfqpoint{0.000000in}{-0.013889in}}%
\pgfpathlineto{\pgfqpoint{0.000000in}{-0.013889in}}%
\pgfpathclose%
\pgfusepath{stroke,fill}%
}%
\begin{pgfscope}%
\pgfsys@transformshift{3.383445in}{3.990583in}%
\pgfsys@useobject{currentmarker}{}%
\end{pgfscope}%
\end{pgfscope}%
\begin{pgfscope}%
\pgfpathrectangle{\pgfqpoint{0.100000in}{2.493063in}}{\pgfqpoint{5.037500in}{3.427208in}}%
\pgfusepath{clip}%
\pgfsetrectcap%
\pgfsetroundjoin%
\pgfsetlinewidth{1.505625pt}%
\definecolor{currentstroke}{rgb}{0.000000,0.000000,1.000000}%
\pgfsetstrokecolor{currentstroke}%
\pgfsetstrokeopacity{0.500000}%
\pgfsetdash{}{0pt}%
\pgfpathmoveto{\pgfqpoint{3.975097in}{4.133085in}}%
\pgfusepath{stroke}%
\end{pgfscope}%
\begin{pgfscope}%
\pgfpathrectangle{\pgfqpoint{0.100000in}{2.493063in}}{\pgfqpoint{5.037500in}{3.427208in}}%
\pgfusepath{clip}%
\pgfsetbuttcap%
\pgfsetroundjoin%
\definecolor{currentfill}{rgb}{0.000000,0.000000,1.000000}%
\pgfsetfillcolor{currentfill}%
\pgfsetfillopacity{0.500000}%
\pgfsetlinewidth{0.250937pt}%
\definecolor{currentstroke}{rgb}{0.000000,0.000000,0.000000}%
\pgfsetstrokecolor{currentstroke}%
\pgfsetstrokeopacity{0.500000}%
\pgfsetdash{}{0pt}%
\pgfsys@defobject{currentmarker}{\pgfqpoint{-0.005556in}{-0.005556in}}{\pgfqpoint{0.005556in}{0.005556in}}{%
\pgfpathmoveto{\pgfqpoint{0.000000in}{-0.005556in}}%
\pgfpathcurveto{\pgfqpoint{0.001473in}{-0.005556in}}{\pgfqpoint{0.002887in}{-0.004970in}}{\pgfqpoint{0.003928in}{-0.003928in}}%
\pgfpathcurveto{\pgfqpoint{0.004970in}{-0.002887in}}{\pgfqpoint{0.005556in}{-0.001473in}}{\pgfqpoint{0.005556in}{0.000000in}}%
\pgfpathcurveto{\pgfqpoint{0.005556in}{0.001473in}}{\pgfqpoint{0.004970in}{0.002887in}}{\pgfqpoint{0.003928in}{0.003928in}}%
\pgfpathcurveto{\pgfqpoint{0.002887in}{0.004970in}}{\pgfqpoint{0.001473in}{0.005556in}}{\pgfqpoint{0.000000in}{0.005556in}}%
\pgfpathcurveto{\pgfqpoint{-0.001473in}{0.005556in}}{\pgfqpoint{-0.002887in}{0.004970in}}{\pgfqpoint{-0.003928in}{0.003928in}}%
\pgfpathcurveto{\pgfqpoint{-0.004970in}{0.002887in}}{\pgfqpoint{-0.005556in}{0.001473in}}{\pgfqpoint{-0.005556in}{0.000000in}}%
\pgfpathcurveto{\pgfqpoint{-0.005556in}{-0.001473in}}{\pgfqpoint{-0.004970in}{-0.002887in}}{\pgfqpoint{-0.003928in}{-0.003928in}}%
\pgfpathcurveto{\pgfqpoint{-0.002887in}{-0.004970in}}{\pgfqpoint{-0.001473in}{-0.005556in}}{\pgfqpoint{0.000000in}{-0.005556in}}%
\pgfpathlineto{\pgfqpoint{0.000000in}{-0.005556in}}%
\pgfpathclose%
\pgfusepath{stroke,fill}%
}%
\begin{pgfscope}%
\pgfsys@transformshift{3.975097in}{4.133085in}%
\pgfsys@useobject{currentmarker}{}%
\end{pgfscope}%
\end{pgfscope}%
\begin{pgfscope}%
\pgfpathrectangle{\pgfqpoint{0.100000in}{2.493063in}}{\pgfqpoint{5.037500in}{3.427208in}}%
\pgfusepath{clip}%
\pgfsetrectcap%
\pgfsetroundjoin%
\pgfsetlinewidth{1.505625pt}%
\definecolor{currentstroke}{rgb}{0.000000,0.000000,1.000000}%
\pgfsetstrokecolor{currentstroke}%
\pgfsetstrokeopacity{0.500000}%
\pgfsetdash{}{0pt}%
\pgfpathmoveto{\pgfqpoint{3.985695in}{4.167467in}}%
\pgfusepath{stroke}%
\end{pgfscope}%
\begin{pgfscope}%
\pgfpathrectangle{\pgfqpoint{0.100000in}{2.493063in}}{\pgfqpoint{5.037500in}{3.427208in}}%
\pgfusepath{clip}%
\pgfsetbuttcap%
\pgfsetroundjoin%
\definecolor{currentfill}{rgb}{0.000000,0.000000,1.000000}%
\pgfsetfillcolor{currentfill}%
\pgfsetfillopacity{0.500000}%
\pgfsetlinewidth{0.250937pt}%
\definecolor{currentstroke}{rgb}{0.000000,0.000000,0.000000}%
\pgfsetstrokecolor{currentstroke}%
\pgfsetstrokeopacity{0.500000}%
\pgfsetdash{}{0pt}%
\pgfsys@defobject{currentmarker}{\pgfqpoint{-0.005556in}{-0.005556in}}{\pgfqpoint{0.005556in}{0.005556in}}{%
\pgfpathmoveto{\pgfqpoint{0.000000in}{-0.005556in}}%
\pgfpathcurveto{\pgfqpoint{0.001473in}{-0.005556in}}{\pgfqpoint{0.002887in}{-0.004970in}}{\pgfqpoint{0.003928in}{-0.003928in}}%
\pgfpathcurveto{\pgfqpoint{0.004970in}{-0.002887in}}{\pgfqpoint{0.005556in}{-0.001473in}}{\pgfqpoint{0.005556in}{0.000000in}}%
\pgfpathcurveto{\pgfqpoint{0.005556in}{0.001473in}}{\pgfqpoint{0.004970in}{0.002887in}}{\pgfqpoint{0.003928in}{0.003928in}}%
\pgfpathcurveto{\pgfqpoint{0.002887in}{0.004970in}}{\pgfqpoint{0.001473in}{0.005556in}}{\pgfqpoint{0.000000in}{0.005556in}}%
\pgfpathcurveto{\pgfqpoint{-0.001473in}{0.005556in}}{\pgfqpoint{-0.002887in}{0.004970in}}{\pgfqpoint{-0.003928in}{0.003928in}}%
\pgfpathcurveto{\pgfqpoint{-0.004970in}{0.002887in}}{\pgfqpoint{-0.005556in}{0.001473in}}{\pgfqpoint{-0.005556in}{0.000000in}}%
\pgfpathcurveto{\pgfqpoint{-0.005556in}{-0.001473in}}{\pgfqpoint{-0.004970in}{-0.002887in}}{\pgfqpoint{-0.003928in}{-0.003928in}}%
\pgfpathcurveto{\pgfqpoint{-0.002887in}{-0.004970in}}{\pgfqpoint{-0.001473in}{-0.005556in}}{\pgfqpoint{0.000000in}{-0.005556in}}%
\pgfpathlineto{\pgfqpoint{0.000000in}{-0.005556in}}%
\pgfpathclose%
\pgfusepath{stroke,fill}%
}%
\begin{pgfscope}%
\pgfsys@transformshift{3.985695in}{4.167467in}%
\pgfsys@useobject{currentmarker}{}%
\end{pgfscope}%
\end{pgfscope}%
\begin{pgfscope}%
\pgfpathrectangle{\pgfqpoint{0.100000in}{2.493063in}}{\pgfqpoint{5.037500in}{3.427208in}}%
\pgfusepath{clip}%
\pgfsetrectcap%
\pgfsetroundjoin%
\pgfsetlinewidth{1.505625pt}%
\definecolor{currentstroke}{rgb}{0.000000,0.000000,1.000000}%
\pgfsetstrokecolor{currentstroke}%
\pgfsetstrokeopacity{0.500000}%
\pgfsetdash{}{0pt}%
\pgfpathmoveto{\pgfqpoint{3.952422in}{4.157493in}}%
\pgfusepath{stroke}%
\end{pgfscope}%
\begin{pgfscope}%
\pgfpathrectangle{\pgfqpoint{0.100000in}{2.493063in}}{\pgfqpoint{5.037500in}{3.427208in}}%
\pgfusepath{clip}%
\pgfsetbuttcap%
\pgfsetroundjoin%
\definecolor{currentfill}{rgb}{0.000000,0.000000,1.000000}%
\pgfsetfillcolor{currentfill}%
\pgfsetfillopacity{0.500000}%
\pgfsetlinewidth{0.250937pt}%
\definecolor{currentstroke}{rgb}{0.000000,0.000000,0.000000}%
\pgfsetstrokecolor{currentstroke}%
\pgfsetstrokeopacity{0.500000}%
\pgfsetdash{}{0pt}%
\pgfsys@defobject{currentmarker}{\pgfqpoint{-0.005556in}{-0.005556in}}{\pgfqpoint{0.005556in}{0.005556in}}{%
\pgfpathmoveto{\pgfqpoint{0.000000in}{-0.005556in}}%
\pgfpathcurveto{\pgfqpoint{0.001473in}{-0.005556in}}{\pgfqpoint{0.002887in}{-0.004970in}}{\pgfqpoint{0.003928in}{-0.003928in}}%
\pgfpathcurveto{\pgfqpoint{0.004970in}{-0.002887in}}{\pgfqpoint{0.005556in}{-0.001473in}}{\pgfqpoint{0.005556in}{0.000000in}}%
\pgfpathcurveto{\pgfqpoint{0.005556in}{0.001473in}}{\pgfqpoint{0.004970in}{0.002887in}}{\pgfqpoint{0.003928in}{0.003928in}}%
\pgfpathcurveto{\pgfqpoint{0.002887in}{0.004970in}}{\pgfqpoint{0.001473in}{0.005556in}}{\pgfqpoint{0.000000in}{0.005556in}}%
\pgfpathcurveto{\pgfqpoint{-0.001473in}{0.005556in}}{\pgfqpoint{-0.002887in}{0.004970in}}{\pgfqpoint{-0.003928in}{0.003928in}}%
\pgfpathcurveto{\pgfqpoint{-0.004970in}{0.002887in}}{\pgfqpoint{-0.005556in}{0.001473in}}{\pgfqpoint{-0.005556in}{0.000000in}}%
\pgfpathcurveto{\pgfqpoint{-0.005556in}{-0.001473in}}{\pgfqpoint{-0.004970in}{-0.002887in}}{\pgfqpoint{-0.003928in}{-0.003928in}}%
\pgfpathcurveto{\pgfqpoint{-0.002887in}{-0.004970in}}{\pgfqpoint{-0.001473in}{-0.005556in}}{\pgfqpoint{0.000000in}{-0.005556in}}%
\pgfpathlineto{\pgfqpoint{0.000000in}{-0.005556in}}%
\pgfpathclose%
\pgfusepath{stroke,fill}%
}%
\begin{pgfscope}%
\pgfsys@transformshift{3.952422in}{4.157493in}%
\pgfsys@useobject{currentmarker}{}%
\end{pgfscope}%
\end{pgfscope}%
\begin{pgfscope}%
\pgfpathrectangle{\pgfqpoint{0.100000in}{2.493063in}}{\pgfqpoint{5.037500in}{3.427208in}}%
\pgfusepath{clip}%
\pgfsetrectcap%
\pgfsetroundjoin%
\pgfsetlinewidth{1.505625pt}%
\definecolor{currentstroke}{rgb}{0.000000,0.000000,1.000000}%
\pgfsetstrokecolor{currentstroke}%
\pgfsetstrokeopacity{0.500000}%
\pgfsetdash{}{0pt}%
\pgfpathmoveto{\pgfqpoint{3.835675in}{4.073728in}}%
\pgfusepath{stroke}%
\end{pgfscope}%
\begin{pgfscope}%
\pgfpathrectangle{\pgfqpoint{0.100000in}{2.493063in}}{\pgfqpoint{5.037500in}{3.427208in}}%
\pgfusepath{clip}%
\pgfsetbuttcap%
\pgfsetroundjoin%
\definecolor{currentfill}{rgb}{0.000000,0.000000,1.000000}%
\pgfsetfillcolor{currentfill}%
\pgfsetfillopacity{0.500000}%
\pgfsetlinewidth{0.250937pt}%
\definecolor{currentstroke}{rgb}{0.000000,0.000000,0.000000}%
\pgfsetstrokecolor{currentstroke}%
\pgfsetstrokeopacity{0.500000}%
\pgfsetdash{}{0pt}%
\pgfsys@defobject{currentmarker}{\pgfqpoint{-0.005556in}{-0.005556in}}{\pgfqpoint{0.005556in}{0.005556in}}{%
\pgfpathmoveto{\pgfqpoint{0.000000in}{-0.005556in}}%
\pgfpathcurveto{\pgfqpoint{0.001473in}{-0.005556in}}{\pgfqpoint{0.002887in}{-0.004970in}}{\pgfqpoint{0.003928in}{-0.003928in}}%
\pgfpathcurveto{\pgfqpoint{0.004970in}{-0.002887in}}{\pgfqpoint{0.005556in}{-0.001473in}}{\pgfqpoint{0.005556in}{0.000000in}}%
\pgfpathcurveto{\pgfqpoint{0.005556in}{0.001473in}}{\pgfqpoint{0.004970in}{0.002887in}}{\pgfqpoint{0.003928in}{0.003928in}}%
\pgfpathcurveto{\pgfqpoint{0.002887in}{0.004970in}}{\pgfqpoint{0.001473in}{0.005556in}}{\pgfqpoint{0.000000in}{0.005556in}}%
\pgfpathcurveto{\pgfqpoint{-0.001473in}{0.005556in}}{\pgfqpoint{-0.002887in}{0.004970in}}{\pgfqpoint{-0.003928in}{0.003928in}}%
\pgfpathcurveto{\pgfqpoint{-0.004970in}{0.002887in}}{\pgfqpoint{-0.005556in}{0.001473in}}{\pgfqpoint{-0.005556in}{0.000000in}}%
\pgfpathcurveto{\pgfqpoint{-0.005556in}{-0.001473in}}{\pgfqpoint{-0.004970in}{-0.002887in}}{\pgfqpoint{-0.003928in}{-0.003928in}}%
\pgfpathcurveto{\pgfqpoint{-0.002887in}{-0.004970in}}{\pgfqpoint{-0.001473in}{-0.005556in}}{\pgfqpoint{0.000000in}{-0.005556in}}%
\pgfpathlineto{\pgfqpoint{0.000000in}{-0.005556in}}%
\pgfpathclose%
\pgfusepath{stroke,fill}%
}%
\begin{pgfscope}%
\pgfsys@transformshift{3.835675in}{4.073728in}%
\pgfsys@useobject{currentmarker}{}%
\end{pgfscope}%
\end{pgfscope}%
\begin{pgfscope}%
\pgfpathrectangle{\pgfqpoint{0.100000in}{2.493063in}}{\pgfqpoint{5.037500in}{3.427208in}}%
\pgfusepath{clip}%
\pgfsetrectcap%
\pgfsetroundjoin%
\pgfsetlinewidth{1.505625pt}%
\definecolor{currentstroke}{rgb}{0.501961,0.501961,0.501961}%
\pgfsetstrokecolor{currentstroke}%
\pgfsetstrokeopacity{0.500000}%
\pgfsetdash{}{0pt}%
\pgfpathmoveto{\pgfqpoint{3.270873in}{3.929879in}}%
\pgfusepath{stroke}%
\end{pgfscope}%
\begin{pgfscope}%
\pgfpathrectangle{\pgfqpoint{0.100000in}{2.493063in}}{\pgfqpoint{5.037500in}{3.427208in}}%
\pgfusepath{clip}%
\pgfsetbuttcap%
\pgfsetroundjoin%
\definecolor{currentfill}{rgb}{0.501961,0.501961,0.501961}%
\pgfsetfillcolor{currentfill}%
\pgfsetfillopacity{0.500000}%
\pgfsetlinewidth{0.250937pt}%
\definecolor{currentstroke}{rgb}{0.000000,0.000000,0.000000}%
\pgfsetstrokecolor{currentstroke}%
\pgfsetstrokeopacity{0.500000}%
\pgfsetdash{}{0pt}%
\pgfsys@defobject{currentmarker}{\pgfqpoint{-0.013889in}{-0.013889in}}{\pgfqpoint{0.013889in}{0.013889in}}{%
\pgfpathmoveto{\pgfqpoint{0.000000in}{-0.013889in}}%
\pgfpathcurveto{\pgfqpoint{0.003683in}{-0.013889in}}{\pgfqpoint{0.007216in}{-0.012425in}}{\pgfqpoint{0.009821in}{-0.009821in}}%
\pgfpathcurveto{\pgfqpoint{0.012425in}{-0.007216in}}{\pgfqpoint{0.013889in}{-0.003683in}}{\pgfqpoint{0.013889in}{0.000000in}}%
\pgfpathcurveto{\pgfqpoint{0.013889in}{0.003683in}}{\pgfqpoint{0.012425in}{0.007216in}}{\pgfqpoint{0.009821in}{0.009821in}}%
\pgfpathcurveto{\pgfqpoint{0.007216in}{0.012425in}}{\pgfqpoint{0.003683in}{0.013889in}}{\pgfqpoint{0.000000in}{0.013889in}}%
\pgfpathcurveto{\pgfqpoint{-0.003683in}{0.013889in}}{\pgfqpoint{-0.007216in}{0.012425in}}{\pgfqpoint{-0.009821in}{0.009821in}}%
\pgfpathcurveto{\pgfqpoint{-0.012425in}{0.007216in}}{\pgfqpoint{-0.013889in}{0.003683in}}{\pgfqpoint{-0.013889in}{0.000000in}}%
\pgfpathcurveto{\pgfqpoint{-0.013889in}{-0.003683in}}{\pgfqpoint{-0.012425in}{-0.007216in}}{\pgfqpoint{-0.009821in}{-0.009821in}}%
\pgfpathcurveto{\pgfqpoint{-0.007216in}{-0.012425in}}{\pgfqpoint{-0.003683in}{-0.013889in}}{\pgfqpoint{0.000000in}{-0.013889in}}%
\pgfpathlineto{\pgfqpoint{0.000000in}{-0.013889in}}%
\pgfpathclose%
\pgfusepath{stroke,fill}%
}%
\begin{pgfscope}%
\pgfsys@transformshift{3.270873in}{3.929879in}%
\pgfsys@useobject{currentmarker}{}%
\end{pgfscope}%
\end{pgfscope}%
\begin{pgfscope}%
\pgfpathrectangle{\pgfqpoint{0.100000in}{2.493063in}}{\pgfqpoint{5.037500in}{3.427208in}}%
\pgfusepath{clip}%
\pgfsetrectcap%
\pgfsetroundjoin%
\pgfsetlinewidth{1.505625pt}%
\definecolor{currentstroke}{rgb}{0.501961,0.501961,0.501961}%
\pgfsetstrokecolor{currentstroke}%
\pgfsetstrokeopacity{0.500000}%
\pgfsetdash{}{0pt}%
\pgfpathmoveto{\pgfqpoint{3.889904in}{4.110032in}}%
\pgfusepath{stroke}%
\end{pgfscope}%
\begin{pgfscope}%
\pgfpathrectangle{\pgfqpoint{0.100000in}{2.493063in}}{\pgfqpoint{5.037500in}{3.427208in}}%
\pgfusepath{clip}%
\pgfsetbuttcap%
\pgfsetroundjoin%
\definecolor{currentfill}{rgb}{0.501961,0.501961,0.501961}%
\pgfsetfillcolor{currentfill}%
\pgfsetfillopacity{0.500000}%
\pgfsetlinewidth{0.250937pt}%
\definecolor{currentstroke}{rgb}{0.000000,0.000000,0.000000}%
\pgfsetstrokecolor{currentstroke}%
\pgfsetstrokeopacity{0.500000}%
\pgfsetdash{}{0pt}%
\pgfsys@defobject{currentmarker}{\pgfqpoint{-0.013889in}{-0.013889in}}{\pgfqpoint{0.013889in}{0.013889in}}{%
\pgfpathmoveto{\pgfqpoint{0.000000in}{-0.013889in}}%
\pgfpathcurveto{\pgfqpoint{0.003683in}{-0.013889in}}{\pgfqpoint{0.007216in}{-0.012425in}}{\pgfqpoint{0.009821in}{-0.009821in}}%
\pgfpathcurveto{\pgfqpoint{0.012425in}{-0.007216in}}{\pgfqpoint{0.013889in}{-0.003683in}}{\pgfqpoint{0.013889in}{0.000000in}}%
\pgfpathcurveto{\pgfqpoint{0.013889in}{0.003683in}}{\pgfqpoint{0.012425in}{0.007216in}}{\pgfqpoint{0.009821in}{0.009821in}}%
\pgfpathcurveto{\pgfqpoint{0.007216in}{0.012425in}}{\pgfqpoint{0.003683in}{0.013889in}}{\pgfqpoint{0.000000in}{0.013889in}}%
\pgfpathcurveto{\pgfqpoint{-0.003683in}{0.013889in}}{\pgfqpoint{-0.007216in}{0.012425in}}{\pgfqpoint{-0.009821in}{0.009821in}}%
\pgfpathcurveto{\pgfqpoint{-0.012425in}{0.007216in}}{\pgfqpoint{-0.013889in}{0.003683in}}{\pgfqpoint{-0.013889in}{0.000000in}}%
\pgfpathcurveto{\pgfqpoint{-0.013889in}{-0.003683in}}{\pgfqpoint{-0.012425in}{-0.007216in}}{\pgfqpoint{-0.009821in}{-0.009821in}}%
\pgfpathcurveto{\pgfqpoint{-0.007216in}{-0.012425in}}{\pgfqpoint{-0.003683in}{-0.013889in}}{\pgfqpoint{0.000000in}{-0.013889in}}%
\pgfpathlineto{\pgfqpoint{0.000000in}{-0.013889in}}%
\pgfpathclose%
\pgfusepath{stroke,fill}%
}%
\begin{pgfscope}%
\pgfsys@transformshift{3.889904in}{4.110032in}%
\pgfsys@useobject{currentmarker}{}%
\end{pgfscope}%
\end{pgfscope}%
\begin{pgfscope}%
\pgfpathrectangle{\pgfqpoint{0.100000in}{2.493063in}}{\pgfqpoint{5.037500in}{3.427208in}}%
\pgfusepath{clip}%
\pgfsetrectcap%
\pgfsetroundjoin%
\pgfsetlinewidth{1.505625pt}%
\definecolor{currentstroke}{rgb}{0.000000,0.000000,1.000000}%
\pgfsetstrokecolor{currentstroke}%
\pgfsetstrokeopacity{0.500000}%
\pgfsetdash{}{0pt}%
\pgfpathmoveto{\pgfqpoint{3.568755in}{4.068672in}}%
\pgfusepath{stroke}%
\end{pgfscope}%
\begin{pgfscope}%
\pgfpathrectangle{\pgfqpoint{0.100000in}{2.493063in}}{\pgfqpoint{5.037500in}{3.427208in}}%
\pgfusepath{clip}%
\pgfsetbuttcap%
\pgfsetroundjoin%
\definecolor{currentfill}{rgb}{0.000000,0.000000,1.000000}%
\pgfsetfillcolor{currentfill}%
\pgfsetfillopacity{0.500000}%
\pgfsetlinewidth{0.250937pt}%
\definecolor{currentstroke}{rgb}{0.000000,0.000000,0.000000}%
\pgfsetstrokecolor{currentstroke}%
\pgfsetstrokeopacity{0.500000}%
\pgfsetdash{}{0pt}%
\pgfsys@defobject{currentmarker}{\pgfqpoint{-0.008333in}{-0.008333in}}{\pgfqpoint{0.008333in}{0.008333in}}{%
\pgfpathmoveto{\pgfqpoint{0.000000in}{-0.008333in}}%
\pgfpathcurveto{\pgfqpoint{0.002210in}{-0.008333in}}{\pgfqpoint{0.004330in}{-0.007455in}}{\pgfqpoint{0.005893in}{-0.005893in}}%
\pgfpathcurveto{\pgfqpoint{0.007455in}{-0.004330in}}{\pgfqpoint{0.008333in}{-0.002210in}}{\pgfqpoint{0.008333in}{0.000000in}}%
\pgfpathcurveto{\pgfqpoint{0.008333in}{0.002210in}}{\pgfqpoint{0.007455in}{0.004330in}}{\pgfqpoint{0.005893in}{0.005893in}}%
\pgfpathcurveto{\pgfqpoint{0.004330in}{0.007455in}}{\pgfqpoint{0.002210in}{0.008333in}}{\pgfqpoint{0.000000in}{0.008333in}}%
\pgfpathcurveto{\pgfqpoint{-0.002210in}{0.008333in}}{\pgfqpoint{-0.004330in}{0.007455in}}{\pgfqpoint{-0.005893in}{0.005893in}}%
\pgfpathcurveto{\pgfqpoint{-0.007455in}{0.004330in}}{\pgfqpoint{-0.008333in}{0.002210in}}{\pgfqpoint{-0.008333in}{0.000000in}}%
\pgfpathcurveto{\pgfqpoint{-0.008333in}{-0.002210in}}{\pgfqpoint{-0.007455in}{-0.004330in}}{\pgfqpoint{-0.005893in}{-0.005893in}}%
\pgfpathcurveto{\pgfqpoint{-0.004330in}{-0.007455in}}{\pgfqpoint{-0.002210in}{-0.008333in}}{\pgfqpoint{0.000000in}{-0.008333in}}%
\pgfpathlineto{\pgfqpoint{0.000000in}{-0.008333in}}%
\pgfpathclose%
\pgfusepath{stroke,fill}%
}%
\begin{pgfscope}%
\pgfsys@transformshift{3.568755in}{4.068672in}%
\pgfsys@useobject{currentmarker}{}%
\end{pgfscope}%
\end{pgfscope}%
\begin{pgfscope}%
\pgfpathrectangle{\pgfqpoint{0.100000in}{2.493063in}}{\pgfqpoint{5.037500in}{3.427208in}}%
\pgfusepath{clip}%
\pgfsetrectcap%
\pgfsetroundjoin%
\pgfsetlinewidth{1.505625pt}%
\definecolor{currentstroke}{rgb}{0.501961,0.501961,0.501961}%
\pgfsetstrokecolor{currentstroke}%
\pgfsetstrokeopacity{0.500000}%
\pgfsetdash{}{0pt}%
\pgfpathmoveto{\pgfqpoint{2.341376in}{3.616703in}}%
\pgfusepath{stroke}%
\end{pgfscope}%
\begin{pgfscope}%
\pgfpathrectangle{\pgfqpoint{0.100000in}{2.493063in}}{\pgfqpoint{5.037500in}{3.427208in}}%
\pgfusepath{clip}%
\pgfsetbuttcap%
\pgfsetroundjoin%
\definecolor{currentfill}{rgb}{0.501961,0.501961,0.501961}%
\pgfsetfillcolor{currentfill}%
\pgfsetfillopacity{0.500000}%
\pgfsetlinewidth{0.250937pt}%
\definecolor{currentstroke}{rgb}{0.000000,0.000000,0.000000}%
\pgfsetstrokecolor{currentstroke}%
\pgfsetstrokeopacity{0.500000}%
\pgfsetdash{}{0pt}%
\pgfsys@defobject{currentmarker}{\pgfqpoint{-0.013889in}{-0.013889in}}{\pgfqpoint{0.013889in}{0.013889in}}{%
\pgfpathmoveto{\pgfqpoint{0.000000in}{-0.013889in}}%
\pgfpathcurveto{\pgfqpoint{0.003683in}{-0.013889in}}{\pgfqpoint{0.007216in}{-0.012425in}}{\pgfqpoint{0.009821in}{-0.009821in}}%
\pgfpathcurveto{\pgfqpoint{0.012425in}{-0.007216in}}{\pgfqpoint{0.013889in}{-0.003683in}}{\pgfqpoint{0.013889in}{0.000000in}}%
\pgfpathcurveto{\pgfqpoint{0.013889in}{0.003683in}}{\pgfqpoint{0.012425in}{0.007216in}}{\pgfqpoint{0.009821in}{0.009821in}}%
\pgfpathcurveto{\pgfqpoint{0.007216in}{0.012425in}}{\pgfqpoint{0.003683in}{0.013889in}}{\pgfqpoint{0.000000in}{0.013889in}}%
\pgfpathcurveto{\pgfqpoint{-0.003683in}{0.013889in}}{\pgfqpoint{-0.007216in}{0.012425in}}{\pgfqpoint{-0.009821in}{0.009821in}}%
\pgfpathcurveto{\pgfqpoint{-0.012425in}{0.007216in}}{\pgfqpoint{-0.013889in}{0.003683in}}{\pgfqpoint{-0.013889in}{0.000000in}}%
\pgfpathcurveto{\pgfqpoint{-0.013889in}{-0.003683in}}{\pgfqpoint{-0.012425in}{-0.007216in}}{\pgfqpoint{-0.009821in}{-0.009821in}}%
\pgfpathcurveto{\pgfqpoint{-0.007216in}{-0.012425in}}{\pgfqpoint{-0.003683in}{-0.013889in}}{\pgfqpoint{0.000000in}{-0.013889in}}%
\pgfpathlineto{\pgfqpoint{0.000000in}{-0.013889in}}%
\pgfpathclose%
\pgfusepath{stroke,fill}%
}%
\begin{pgfscope}%
\pgfsys@transformshift{2.341376in}{3.616703in}%
\pgfsys@useobject{currentmarker}{}%
\end{pgfscope}%
\end{pgfscope}%
\begin{pgfscope}%
\pgfpathrectangle{\pgfqpoint{0.100000in}{2.493063in}}{\pgfqpoint{5.037500in}{3.427208in}}%
\pgfusepath{clip}%
\pgfsetrectcap%
\pgfsetroundjoin%
\pgfsetlinewidth{1.505625pt}%
\definecolor{currentstroke}{rgb}{0.000000,0.000000,1.000000}%
\pgfsetstrokecolor{currentstroke}%
\pgfsetstrokeopacity{0.500000}%
\pgfsetdash{}{0pt}%
\pgfpathmoveto{\pgfqpoint{2.160335in}{3.948102in}}%
\pgfusepath{stroke}%
\end{pgfscope}%
\begin{pgfscope}%
\pgfpathrectangle{\pgfqpoint{0.100000in}{2.493063in}}{\pgfqpoint{5.037500in}{3.427208in}}%
\pgfusepath{clip}%
\pgfsetbuttcap%
\pgfsetroundjoin%
\definecolor{currentfill}{rgb}{0.000000,0.000000,1.000000}%
\pgfsetfillcolor{currentfill}%
\pgfsetfillopacity{0.500000}%
\pgfsetlinewidth{0.250937pt}%
\definecolor{currentstroke}{rgb}{0.000000,0.000000,0.000000}%
\pgfsetstrokecolor{currentstroke}%
\pgfsetstrokeopacity{0.500000}%
\pgfsetdash{}{0pt}%
\pgfsys@defobject{currentmarker}{\pgfqpoint{-0.005556in}{-0.005556in}}{\pgfqpoint{0.005556in}{0.005556in}}{%
\pgfpathmoveto{\pgfqpoint{0.000000in}{-0.005556in}}%
\pgfpathcurveto{\pgfqpoint{0.001473in}{-0.005556in}}{\pgfqpoint{0.002887in}{-0.004970in}}{\pgfqpoint{0.003928in}{-0.003928in}}%
\pgfpathcurveto{\pgfqpoint{0.004970in}{-0.002887in}}{\pgfqpoint{0.005556in}{-0.001473in}}{\pgfqpoint{0.005556in}{0.000000in}}%
\pgfpathcurveto{\pgfqpoint{0.005556in}{0.001473in}}{\pgfqpoint{0.004970in}{0.002887in}}{\pgfqpoint{0.003928in}{0.003928in}}%
\pgfpathcurveto{\pgfqpoint{0.002887in}{0.004970in}}{\pgfqpoint{0.001473in}{0.005556in}}{\pgfqpoint{0.000000in}{0.005556in}}%
\pgfpathcurveto{\pgfqpoint{-0.001473in}{0.005556in}}{\pgfqpoint{-0.002887in}{0.004970in}}{\pgfqpoint{-0.003928in}{0.003928in}}%
\pgfpathcurveto{\pgfqpoint{-0.004970in}{0.002887in}}{\pgfqpoint{-0.005556in}{0.001473in}}{\pgfqpoint{-0.005556in}{0.000000in}}%
\pgfpathcurveto{\pgfqpoint{-0.005556in}{-0.001473in}}{\pgfqpoint{-0.004970in}{-0.002887in}}{\pgfqpoint{-0.003928in}{-0.003928in}}%
\pgfpathcurveto{\pgfqpoint{-0.002887in}{-0.004970in}}{\pgfqpoint{-0.001473in}{-0.005556in}}{\pgfqpoint{0.000000in}{-0.005556in}}%
\pgfpathlineto{\pgfqpoint{0.000000in}{-0.005556in}}%
\pgfpathclose%
\pgfusepath{stroke,fill}%
}%
\begin{pgfscope}%
\pgfsys@transformshift{2.160335in}{3.948102in}%
\pgfsys@useobject{currentmarker}{}%
\end{pgfscope}%
\end{pgfscope}%
\begin{pgfscope}%
\pgfpathrectangle{\pgfqpoint{0.100000in}{2.493063in}}{\pgfqpoint{5.037500in}{3.427208in}}%
\pgfusepath{clip}%
\pgfsetrectcap%
\pgfsetroundjoin%
\pgfsetlinewidth{1.505625pt}%
\definecolor{currentstroke}{rgb}{0.000000,0.000000,1.000000}%
\pgfsetstrokecolor{currentstroke}%
\pgfsetstrokeopacity{0.500000}%
\pgfsetdash{}{0pt}%
\pgfpathmoveto{\pgfqpoint{2.528183in}{3.356102in}}%
\pgfusepath{stroke}%
\end{pgfscope}%
\begin{pgfscope}%
\pgfpathrectangle{\pgfqpoint{0.100000in}{2.493063in}}{\pgfqpoint{5.037500in}{3.427208in}}%
\pgfusepath{clip}%
\pgfsetbuttcap%
\pgfsetroundjoin%
\definecolor{currentfill}{rgb}{0.000000,0.000000,1.000000}%
\pgfsetfillcolor{currentfill}%
\pgfsetfillopacity{0.500000}%
\pgfsetlinewidth{0.250937pt}%
\definecolor{currentstroke}{rgb}{0.000000,0.000000,0.000000}%
\pgfsetstrokecolor{currentstroke}%
\pgfsetstrokeopacity{0.500000}%
\pgfsetdash{}{0pt}%
\pgfsys@defobject{currentmarker}{\pgfqpoint{-0.013889in}{-0.013889in}}{\pgfqpoint{0.013889in}{0.013889in}}{%
\pgfpathmoveto{\pgfqpoint{0.000000in}{-0.013889in}}%
\pgfpathcurveto{\pgfqpoint{0.003683in}{-0.013889in}}{\pgfqpoint{0.007216in}{-0.012425in}}{\pgfqpoint{0.009821in}{-0.009821in}}%
\pgfpathcurveto{\pgfqpoint{0.012425in}{-0.007216in}}{\pgfqpoint{0.013889in}{-0.003683in}}{\pgfqpoint{0.013889in}{0.000000in}}%
\pgfpathcurveto{\pgfqpoint{0.013889in}{0.003683in}}{\pgfqpoint{0.012425in}{0.007216in}}{\pgfqpoint{0.009821in}{0.009821in}}%
\pgfpathcurveto{\pgfqpoint{0.007216in}{0.012425in}}{\pgfqpoint{0.003683in}{0.013889in}}{\pgfqpoint{0.000000in}{0.013889in}}%
\pgfpathcurveto{\pgfqpoint{-0.003683in}{0.013889in}}{\pgfqpoint{-0.007216in}{0.012425in}}{\pgfqpoint{-0.009821in}{0.009821in}}%
\pgfpathcurveto{\pgfqpoint{-0.012425in}{0.007216in}}{\pgfqpoint{-0.013889in}{0.003683in}}{\pgfqpoint{-0.013889in}{0.000000in}}%
\pgfpathcurveto{\pgfqpoint{-0.013889in}{-0.003683in}}{\pgfqpoint{-0.012425in}{-0.007216in}}{\pgfqpoint{-0.009821in}{-0.009821in}}%
\pgfpathcurveto{\pgfqpoint{-0.007216in}{-0.012425in}}{\pgfqpoint{-0.003683in}{-0.013889in}}{\pgfqpoint{0.000000in}{-0.013889in}}%
\pgfpathlineto{\pgfqpoint{0.000000in}{-0.013889in}}%
\pgfpathclose%
\pgfusepath{stroke,fill}%
}%
\begin{pgfscope}%
\pgfsys@transformshift{2.528183in}{3.356102in}%
\pgfsys@useobject{currentmarker}{}%
\end{pgfscope}%
\end{pgfscope}%
\begin{pgfscope}%
\pgfpathrectangle{\pgfqpoint{0.100000in}{2.493063in}}{\pgfqpoint{5.037500in}{3.427208in}}%
\pgfusepath{clip}%
\pgfsetrectcap%
\pgfsetroundjoin%
\pgfsetlinewidth{1.505625pt}%
\definecolor{currentstroke}{rgb}{0.501961,0.501961,0.501961}%
\pgfsetstrokecolor{currentstroke}%
\pgfsetstrokeopacity{0.500000}%
\pgfsetdash{}{0pt}%
\pgfpathmoveto{\pgfqpoint{2.894684in}{3.330822in}}%
\pgfusepath{stroke}%
\end{pgfscope}%
\begin{pgfscope}%
\pgfpathrectangle{\pgfqpoint{0.100000in}{2.493063in}}{\pgfqpoint{5.037500in}{3.427208in}}%
\pgfusepath{clip}%
\pgfsetbuttcap%
\pgfsetroundjoin%
\definecolor{currentfill}{rgb}{0.501961,0.501961,0.501961}%
\pgfsetfillcolor{currentfill}%
\pgfsetfillopacity{0.500000}%
\pgfsetlinewidth{0.250937pt}%
\definecolor{currentstroke}{rgb}{0.000000,0.000000,0.000000}%
\pgfsetstrokecolor{currentstroke}%
\pgfsetstrokeopacity{0.500000}%
\pgfsetdash{}{0pt}%
\pgfsys@defobject{currentmarker}{\pgfqpoint{-0.013889in}{-0.013889in}}{\pgfqpoint{0.013889in}{0.013889in}}{%
\pgfpathmoveto{\pgfqpoint{0.000000in}{-0.013889in}}%
\pgfpathcurveto{\pgfqpoint{0.003683in}{-0.013889in}}{\pgfqpoint{0.007216in}{-0.012425in}}{\pgfqpoint{0.009821in}{-0.009821in}}%
\pgfpathcurveto{\pgfqpoint{0.012425in}{-0.007216in}}{\pgfqpoint{0.013889in}{-0.003683in}}{\pgfqpoint{0.013889in}{0.000000in}}%
\pgfpathcurveto{\pgfqpoint{0.013889in}{0.003683in}}{\pgfqpoint{0.012425in}{0.007216in}}{\pgfqpoint{0.009821in}{0.009821in}}%
\pgfpathcurveto{\pgfqpoint{0.007216in}{0.012425in}}{\pgfqpoint{0.003683in}{0.013889in}}{\pgfqpoint{0.000000in}{0.013889in}}%
\pgfpathcurveto{\pgfqpoint{-0.003683in}{0.013889in}}{\pgfqpoint{-0.007216in}{0.012425in}}{\pgfqpoint{-0.009821in}{0.009821in}}%
\pgfpathcurveto{\pgfqpoint{-0.012425in}{0.007216in}}{\pgfqpoint{-0.013889in}{0.003683in}}{\pgfqpoint{-0.013889in}{0.000000in}}%
\pgfpathcurveto{\pgfqpoint{-0.013889in}{-0.003683in}}{\pgfqpoint{-0.012425in}{-0.007216in}}{\pgfqpoint{-0.009821in}{-0.009821in}}%
\pgfpathcurveto{\pgfqpoint{-0.007216in}{-0.012425in}}{\pgfqpoint{-0.003683in}{-0.013889in}}{\pgfqpoint{0.000000in}{-0.013889in}}%
\pgfpathlineto{\pgfqpoint{0.000000in}{-0.013889in}}%
\pgfpathclose%
\pgfusepath{stroke,fill}%
}%
\begin{pgfscope}%
\pgfsys@transformshift{2.894684in}{3.330822in}%
\pgfsys@useobject{currentmarker}{}%
\end{pgfscope}%
\end{pgfscope}%
\begin{pgfscope}%
\pgfpathrectangle{\pgfqpoint{0.100000in}{2.493063in}}{\pgfqpoint{5.037500in}{3.427208in}}%
\pgfusepath{clip}%
\pgfsetrectcap%
\pgfsetroundjoin%
\pgfsetlinewidth{1.505625pt}%
\definecolor{currentstroke}{rgb}{0.678431,1.000000,0.184314}%
\pgfsetstrokecolor{currentstroke}%
\pgfsetstrokeopacity{0.500000}%
\pgfsetdash{}{0pt}%
\pgfpathmoveto{\pgfqpoint{2.539796in}{2.844801in}}%
\pgfusepath{stroke}%
\end{pgfscope}%
\begin{pgfscope}%
\pgfpathrectangle{\pgfqpoint{0.100000in}{2.493063in}}{\pgfqpoint{5.037500in}{3.427208in}}%
\pgfusepath{clip}%
\pgfsetbuttcap%
\pgfsetroundjoin%
\definecolor{currentfill}{rgb}{0.678431,1.000000,0.184314}%
\pgfsetfillcolor{currentfill}%
\pgfsetfillopacity{0.500000}%
\pgfsetlinewidth{0.250937pt}%
\definecolor{currentstroke}{rgb}{0.000000,0.000000,0.000000}%
\pgfsetstrokecolor{currentstroke}%
\pgfsetstrokeopacity{0.500000}%
\pgfsetdash{}{0pt}%
\pgfsys@defobject{currentmarker}{\pgfqpoint{-0.016667in}{-0.016667in}}{\pgfqpoint{0.016667in}{0.016667in}}{%
\pgfpathmoveto{\pgfqpoint{0.000000in}{-0.016667in}}%
\pgfpathcurveto{\pgfqpoint{0.004420in}{-0.016667in}}{\pgfqpoint{0.008660in}{-0.014911in}}{\pgfqpoint{0.011785in}{-0.011785in}}%
\pgfpathcurveto{\pgfqpoint{0.014911in}{-0.008660in}}{\pgfqpoint{0.016667in}{-0.004420in}}{\pgfqpoint{0.016667in}{0.000000in}}%
\pgfpathcurveto{\pgfqpoint{0.016667in}{0.004420in}}{\pgfqpoint{0.014911in}{0.008660in}}{\pgfqpoint{0.011785in}{0.011785in}}%
\pgfpathcurveto{\pgfqpoint{0.008660in}{0.014911in}}{\pgfqpoint{0.004420in}{0.016667in}}{\pgfqpoint{0.000000in}{0.016667in}}%
\pgfpathcurveto{\pgfqpoint{-0.004420in}{0.016667in}}{\pgfqpoint{-0.008660in}{0.014911in}}{\pgfqpoint{-0.011785in}{0.011785in}}%
\pgfpathcurveto{\pgfqpoint{-0.014911in}{0.008660in}}{\pgfqpoint{-0.016667in}{0.004420in}}{\pgfqpoint{-0.016667in}{0.000000in}}%
\pgfpathcurveto{\pgfqpoint{-0.016667in}{-0.004420in}}{\pgfqpoint{-0.014911in}{-0.008660in}}{\pgfqpoint{-0.011785in}{-0.011785in}}%
\pgfpathcurveto{\pgfqpoint{-0.008660in}{-0.014911in}}{\pgfqpoint{-0.004420in}{-0.016667in}}{\pgfqpoint{0.000000in}{-0.016667in}}%
\pgfpathlineto{\pgfqpoint{0.000000in}{-0.016667in}}%
\pgfpathclose%
\pgfusepath{stroke,fill}%
}%
\begin{pgfscope}%
\pgfsys@transformshift{2.539796in}{2.844801in}%
\pgfsys@useobject{currentmarker}{}%
\end{pgfscope}%
\end{pgfscope}%
\begin{pgfscope}%
\pgfpathrectangle{\pgfqpoint{0.100000in}{2.493063in}}{\pgfqpoint{5.037500in}{3.427208in}}%
\pgfusepath{clip}%
\pgfsetrectcap%
\pgfsetroundjoin%
\pgfsetlinewidth{1.505625pt}%
\definecolor{currentstroke}{rgb}{0.501961,0.501961,0.501961}%
\pgfsetstrokecolor{currentstroke}%
\pgfsetstrokeopacity{0.500000}%
\pgfsetdash{}{0pt}%
\pgfpathmoveto{\pgfqpoint{2.671238in}{3.394144in}}%
\pgfusepath{stroke}%
\end{pgfscope}%
\begin{pgfscope}%
\pgfpathrectangle{\pgfqpoint{0.100000in}{2.493063in}}{\pgfqpoint{5.037500in}{3.427208in}}%
\pgfusepath{clip}%
\pgfsetbuttcap%
\pgfsetroundjoin%
\definecolor{currentfill}{rgb}{0.501961,0.501961,0.501961}%
\pgfsetfillcolor{currentfill}%
\pgfsetfillopacity{0.500000}%
\pgfsetlinewidth{0.250937pt}%
\definecolor{currentstroke}{rgb}{0.000000,0.000000,0.000000}%
\pgfsetstrokecolor{currentstroke}%
\pgfsetstrokeopacity{0.500000}%
\pgfsetdash{}{0pt}%
\pgfsys@defobject{currentmarker}{\pgfqpoint{-0.013889in}{-0.013889in}}{\pgfqpoint{0.013889in}{0.013889in}}{%
\pgfpathmoveto{\pgfqpoint{0.000000in}{-0.013889in}}%
\pgfpathcurveto{\pgfqpoint{0.003683in}{-0.013889in}}{\pgfqpoint{0.007216in}{-0.012425in}}{\pgfqpoint{0.009821in}{-0.009821in}}%
\pgfpathcurveto{\pgfqpoint{0.012425in}{-0.007216in}}{\pgfqpoint{0.013889in}{-0.003683in}}{\pgfqpoint{0.013889in}{0.000000in}}%
\pgfpathcurveto{\pgfqpoint{0.013889in}{0.003683in}}{\pgfqpoint{0.012425in}{0.007216in}}{\pgfqpoint{0.009821in}{0.009821in}}%
\pgfpathcurveto{\pgfqpoint{0.007216in}{0.012425in}}{\pgfqpoint{0.003683in}{0.013889in}}{\pgfqpoint{0.000000in}{0.013889in}}%
\pgfpathcurveto{\pgfqpoint{-0.003683in}{0.013889in}}{\pgfqpoint{-0.007216in}{0.012425in}}{\pgfqpoint{-0.009821in}{0.009821in}}%
\pgfpathcurveto{\pgfqpoint{-0.012425in}{0.007216in}}{\pgfqpoint{-0.013889in}{0.003683in}}{\pgfqpoint{-0.013889in}{0.000000in}}%
\pgfpathcurveto{\pgfqpoint{-0.013889in}{-0.003683in}}{\pgfqpoint{-0.012425in}{-0.007216in}}{\pgfqpoint{-0.009821in}{-0.009821in}}%
\pgfpathcurveto{\pgfqpoint{-0.007216in}{-0.012425in}}{\pgfqpoint{-0.003683in}{-0.013889in}}{\pgfqpoint{0.000000in}{-0.013889in}}%
\pgfpathlineto{\pgfqpoint{0.000000in}{-0.013889in}}%
\pgfpathclose%
\pgfusepath{stroke,fill}%
}%
\begin{pgfscope}%
\pgfsys@transformshift{2.671238in}{3.394144in}%
\pgfsys@useobject{currentmarker}{}%
\end{pgfscope}%
\end{pgfscope}%
\begin{pgfscope}%
\pgfpathrectangle{\pgfqpoint{0.100000in}{2.493063in}}{\pgfqpoint{5.037500in}{3.427208in}}%
\pgfusepath{clip}%
\pgfsetrectcap%
\pgfsetroundjoin%
\pgfsetlinewidth{1.505625pt}%
\definecolor{currentstroke}{rgb}{0.501961,0.501961,0.501961}%
\pgfsetstrokecolor{currentstroke}%
\pgfsetstrokeopacity{0.500000}%
\pgfsetdash{}{0pt}%
\pgfpathmoveto{\pgfqpoint{2.555575in}{3.066438in}}%
\pgfusepath{stroke}%
\end{pgfscope}%
\begin{pgfscope}%
\pgfpathrectangle{\pgfqpoint{0.100000in}{2.493063in}}{\pgfqpoint{5.037500in}{3.427208in}}%
\pgfusepath{clip}%
\pgfsetbuttcap%
\pgfsetroundjoin%
\definecolor{currentfill}{rgb}{0.501961,0.501961,0.501961}%
\pgfsetfillcolor{currentfill}%
\pgfsetfillopacity{0.500000}%
\pgfsetlinewidth{0.250937pt}%
\definecolor{currentstroke}{rgb}{0.000000,0.000000,0.000000}%
\pgfsetstrokecolor{currentstroke}%
\pgfsetstrokeopacity{0.500000}%
\pgfsetdash{}{0pt}%
\pgfsys@defobject{currentmarker}{\pgfqpoint{-0.013889in}{-0.013889in}}{\pgfqpoint{0.013889in}{0.013889in}}{%
\pgfpathmoveto{\pgfqpoint{0.000000in}{-0.013889in}}%
\pgfpathcurveto{\pgfqpoint{0.003683in}{-0.013889in}}{\pgfqpoint{0.007216in}{-0.012425in}}{\pgfqpoint{0.009821in}{-0.009821in}}%
\pgfpathcurveto{\pgfqpoint{0.012425in}{-0.007216in}}{\pgfqpoint{0.013889in}{-0.003683in}}{\pgfqpoint{0.013889in}{0.000000in}}%
\pgfpathcurveto{\pgfqpoint{0.013889in}{0.003683in}}{\pgfqpoint{0.012425in}{0.007216in}}{\pgfqpoint{0.009821in}{0.009821in}}%
\pgfpathcurveto{\pgfqpoint{0.007216in}{0.012425in}}{\pgfqpoint{0.003683in}{0.013889in}}{\pgfqpoint{0.000000in}{0.013889in}}%
\pgfpathcurveto{\pgfqpoint{-0.003683in}{0.013889in}}{\pgfqpoint{-0.007216in}{0.012425in}}{\pgfqpoint{-0.009821in}{0.009821in}}%
\pgfpathcurveto{\pgfqpoint{-0.012425in}{0.007216in}}{\pgfqpoint{-0.013889in}{0.003683in}}{\pgfqpoint{-0.013889in}{0.000000in}}%
\pgfpathcurveto{\pgfqpoint{-0.013889in}{-0.003683in}}{\pgfqpoint{-0.012425in}{-0.007216in}}{\pgfqpoint{-0.009821in}{-0.009821in}}%
\pgfpathcurveto{\pgfqpoint{-0.007216in}{-0.012425in}}{\pgfqpoint{-0.003683in}{-0.013889in}}{\pgfqpoint{0.000000in}{-0.013889in}}%
\pgfpathlineto{\pgfqpoint{0.000000in}{-0.013889in}}%
\pgfpathclose%
\pgfusepath{stroke,fill}%
}%
\begin{pgfscope}%
\pgfsys@transformshift{2.555575in}{3.066438in}%
\pgfsys@useobject{currentmarker}{}%
\end{pgfscope}%
\end{pgfscope}%
\begin{pgfscope}%
\pgfpathrectangle{\pgfqpoint{0.100000in}{2.493063in}}{\pgfqpoint{5.037500in}{3.427208in}}%
\pgfusepath{clip}%
\pgfsetrectcap%
\pgfsetroundjoin%
\pgfsetlinewidth{1.505625pt}%
\definecolor{currentstroke}{rgb}{0.501961,0.501961,0.501961}%
\pgfsetstrokecolor{currentstroke}%
\pgfsetstrokeopacity{0.500000}%
\pgfsetdash{}{0pt}%
\pgfpathmoveto{\pgfqpoint{2.629147in}{3.644594in}}%
\pgfusepath{stroke}%
\end{pgfscope}%
\begin{pgfscope}%
\pgfpathrectangle{\pgfqpoint{0.100000in}{2.493063in}}{\pgfqpoint{5.037500in}{3.427208in}}%
\pgfusepath{clip}%
\pgfsetbuttcap%
\pgfsetroundjoin%
\definecolor{currentfill}{rgb}{0.501961,0.501961,0.501961}%
\pgfsetfillcolor{currentfill}%
\pgfsetfillopacity{0.500000}%
\pgfsetlinewidth{0.250937pt}%
\definecolor{currentstroke}{rgb}{0.000000,0.000000,0.000000}%
\pgfsetstrokecolor{currentstroke}%
\pgfsetstrokeopacity{0.500000}%
\pgfsetdash{}{0pt}%
\pgfsys@defobject{currentmarker}{\pgfqpoint{-0.013889in}{-0.013889in}}{\pgfqpoint{0.013889in}{0.013889in}}{%
\pgfpathmoveto{\pgfqpoint{0.000000in}{-0.013889in}}%
\pgfpathcurveto{\pgfqpoint{0.003683in}{-0.013889in}}{\pgfqpoint{0.007216in}{-0.012425in}}{\pgfqpoint{0.009821in}{-0.009821in}}%
\pgfpathcurveto{\pgfqpoint{0.012425in}{-0.007216in}}{\pgfqpoint{0.013889in}{-0.003683in}}{\pgfqpoint{0.013889in}{0.000000in}}%
\pgfpathcurveto{\pgfqpoint{0.013889in}{0.003683in}}{\pgfqpoint{0.012425in}{0.007216in}}{\pgfqpoint{0.009821in}{0.009821in}}%
\pgfpathcurveto{\pgfqpoint{0.007216in}{0.012425in}}{\pgfqpoint{0.003683in}{0.013889in}}{\pgfqpoint{0.000000in}{0.013889in}}%
\pgfpathcurveto{\pgfqpoint{-0.003683in}{0.013889in}}{\pgfqpoint{-0.007216in}{0.012425in}}{\pgfqpoint{-0.009821in}{0.009821in}}%
\pgfpathcurveto{\pgfqpoint{-0.012425in}{0.007216in}}{\pgfqpoint{-0.013889in}{0.003683in}}{\pgfqpoint{-0.013889in}{0.000000in}}%
\pgfpathcurveto{\pgfqpoint{-0.013889in}{-0.003683in}}{\pgfqpoint{-0.012425in}{-0.007216in}}{\pgfqpoint{-0.009821in}{-0.009821in}}%
\pgfpathcurveto{\pgfqpoint{-0.007216in}{-0.012425in}}{\pgfqpoint{-0.003683in}{-0.013889in}}{\pgfqpoint{0.000000in}{-0.013889in}}%
\pgfpathlineto{\pgfqpoint{0.000000in}{-0.013889in}}%
\pgfpathclose%
\pgfusepath{stroke,fill}%
}%
\begin{pgfscope}%
\pgfsys@transformshift{2.629147in}{3.644594in}%
\pgfsys@useobject{currentmarker}{}%
\end{pgfscope}%
\end{pgfscope}%
\begin{pgfscope}%
\pgfpathrectangle{\pgfqpoint{0.100000in}{2.493063in}}{\pgfqpoint{5.037500in}{3.427208in}}%
\pgfusepath{clip}%
\pgfsetrectcap%
\pgfsetroundjoin%
\pgfsetlinewidth{1.505625pt}%
\definecolor{currentstroke}{rgb}{0.000000,0.000000,1.000000}%
\pgfsetstrokecolor{currentstroke}%
\pgfsetstrokeopacity{0.500000}%
\pgfsetdash{}{0pt}%
\pgfpathmoveto{\pgfqpoint{1.673527in}{3.596590in}}%
\pgfusepath{stroke}%
\end{pgfscope}%
\begin{pgfscope}%
\pgfpathrectangle{\pgfqpoint{0.100000in}{2.493063in}}{\pgfqpoint{5.037500in}{3.427208in}}%
\pgfusepath{clip}%
\pgfsetbuttcap%
\pgfsetroundjoin%
\definecolor{currentfill}{rgb}{0.000000,0.000000,1.000000}%
\pgfsetfillcolor{currentfill}%
\pgfsetfillopacity{0.500000}%
\pgfsetlinewidth{0.250937pt}%
\definecolor{currentstroke}{rgb}{0.000000,0.000000,0.000000}%
\pgfsetstrokecolor{currentstroke}%
\pgfsetstrokeopacity{0.500000}%
\pgfsetdash{}{0pt}%
\pgfsys@defobject{currentmarker}{\pgfqpoint{-0.011111in}{-0.011111in}}{\pgfqpoint{0.011111in}{0.011111in}}{%
\pgfpathmoveto{\pgfqpoint{0.000000in}{-0.011111in}}%
\pgfpathcurveto{\pgfqpoint{0.002947in}{-0.011111in}}{\pgfqpoint{0.005773in}{-0.009940in}}{\pgfqpoint{0.007857in}{-0.007857in}}%
\pgfpathcurveto{\pgfqpoint{0.009940in}{-0.005773in}}{\pgfqpoint{0.011111in}{-0.002947in}}{\pgfqpoint{0.011111in}{0.000000in}}%
\pgfpathcurveto{\pgfqpoint{0.011111in}{0.002947in}}{\pgfqpoint{0.009940in}{0.005773in}}{\pgfqpoint{0.007857in}{0.007857in}}%
\pgfpathcurveto{\pgfqpoint{0.005773in}{0.009940in}}{\pgfqpoint{0.002947in}{0.011111in}}{\pgfqpoint{0.000000in}{0.011111in}}%
\pgfpathcurveto{\pgfqpoint{-0.002947in}{0.011111in}}{\pgfqpoint{-0.005773in}{0.009940in}}{\pgfqpoint{-0.007857in}{0.007857in}}%
\pgfpathcurveto{\pgfqpoint{-0.009940in}{0.005773in}}{\pgfqpoint{-0.011111in}{0.002947in}}{\pgfqpoint{-0.011111in}{0.000000in}}%
\pgfpathcurveto{\pgfqpoint{-0.011111in}{-0.002947in}}{\pgfqpoint{-0.009940in}{-0.005773in}}{\pgfqpoint{-0.007857in}{-0.007857in}}%
\pgfpathcurveto{\pgfqpoint{-0.005773in}{-0.009940in}}{\pgfqpoint{-0.002947in}{-0.011111in}}{\pgfqpoint{0.000000in}{-0.011111in}}%
\pgfpathlineto{\pgfqpoint{0.000000in}{-0.011111in}}%
\pgfpathclose%
\pgfusepath{stroke,fill}%
}%
\begin{pgfscope}%
\pgfsys@transformshift{1.673527in}{3.596590in}%
\pgfsys@useobject{currentmarker}{}%
\end{pgfscope}%
\end{pgfscope}%
\begin{pgfscope}%
\pgfpathrectangle{\pgfqpoint{0.100000in}{2.493063in}}{\pgfqpoint{5.037500in}{3.427208in}}%
\pgfusepath{clip}%
\pgfsetrectcap%
\pgfsetroundjoin%
\pgfsetlinewidth{1.505625pt}%
\definecolor{currentstroke}{rgb}{0.501961,0.501961,0.501961}%
\pgfsetstrokecolor{currentstroke}%
\pgfsetstrokeopacity{0.500000}%
\pgfsetdash{}{0pt}%
\pgfpathmoveto{\pgfqpoint{2.766900in}{3.292317in}}%
\pgfusepath{stroke}%
\end{pgfscope}%
\begin{pgfscope}%
\pgfpathrectangle{\pgfqpoint{0.100000in}{2.493063in}}{\pgfqpoint{5.037500in}{3.427208in}}%
\pgfusepath{clip}%
\pgfsetbuttcap%
\pgfsetroundjoin%
\definecolor{currentfill}{rgb}{0.501961,0.501961,0.501961}%
\pgfsetfillcolor{currentfill}%
\pgfsetfillopacity{0.500000}%
\pgfsetlinewidth{0.250937pt}%
\definecolor{currentstroke}{rgb}{0.000000,0.000000,0.000000}%
\pgfsetstrokecolor{currentstroke}%
\pgfsetstrokeopacity{0.500000}%
\pgfsetdash{}{0pt}%
\pgfsys@defobject{currentmarker}{\pgfqpoint{-0.013889in}{-0.013889in}}{\pgfqpoint{0.013889in}{0.013889in}}{%
\pgfpathmoveto{\pgfqpoint{0.000000in}{-0.013889in}}%
\pgfpathcurveto{\pgfqpoint{0.003683in}{-0.013889in}}{\pgfqpoint{0.007216in}{-0.012425in}}{\pgfqpoint{0.009821in}{-0.009821in}}%
\pgfpathcurveto{\pgfqpoint{0.012425in}{-0.007216in}}{\pgfqpoint{0.013889in}{-0.003683in}}{\pgfqpoint{0.013889in}{0.000000in}}%
\pgfpathcurveto{\pgfqpoint{0.013889in}{0.003683in}}{\pgfqpoint{0.012425in}{0.007216in}}{\pgfqpoint{0.009821in}{0.009821in}}%
\pgfpathcurveto{\pgfqpoint{0.007216in}{0.012425in}}{\pgfqpoint{0.003683in}{0.013889in}}{\pgfqpoint{0.000000in}{0.013889in}}%
\pgfpathcurveto{\pgfqpoint{-0.003683in}{0.013889in}}{\pgfqpoint{-0.007216in}{0.012425in}}{\pgfqpoint{-0.009821in}{0.009821in}}%
\pgfpathcurveto{\pgfqpoint{-0.012425in}{0.007216in}}{\pgfqpoint{-0.013889in}{0.003683in}}{\pgfqpoint{-0.013889in}{0.000000in}}%
\pgfpathcurveto{\pgfqpoint{-0.013889in}{-0.003683in}}{\pgfqpoint{-0.012425in}{-0.007216in}}{\pgfqpoint{-0.009821in}{-0.009821in}}%
\pgfpathcurveto{\pgfqpoint{-0.007216in}{-0.012425in}}{\pgfqpoint{-0.003683in}{-0.013889in}}{\pgfqpoint{0.000000in}{-0.013889in}}%
\pgfpathlineto{\pgfqpoint{0.000000in}{-0.013889in}}%
\pgfpathclose%
\pgfusepath{stroke,fill}%
}%
\begin{pgfscope}%
\pgfsys@transformshift{2.766900in}{3.292317in}%
\pgfsys@useobject{currentmarker}{}%
\end{pgfscope}%
\end{pgfscope}%
\begin{pgfscope}%
\pgfpathrectangle{\pgfqpoint{0.100000in}{2.493063in}}{\pgfqpoint{5.037500in}{3.427208in}}%
\pgfusepath{clip}%
\pgfsetrectcap%
\pgfsetroundjoin%
\pgfsetlinewidth{1.505625pt}%
\definecolor{currentstroke}{rgb}{0.000000,0.000000,1.000000}%
\pgfsetstrokecolor{currentstroke}%
\pgfsetstrokeopacity{0.500000}%
\pgfsetdash{}{0pt}%
\pgfpathmoveto{\pgfqpoint{2.532736in}{3.454442in}}%
\pgfusepath{stroke}%
\end{pgfscope}%
\begin{pgfscope}%
\pgfpathrectangle{\pgfqpoint{0.100000in}{2.493063in}}{\pgfqpoint{5.037500in}{3.427208in}}%
\pgfusepath{clip}%
\pgfsetbuttcap%
\pgfsetroundjoin%
\definecolor{currentfill}{rgb}{0.000000,0.000000,1.000000}%
\pgfsetfillcolor{currentfill}%
\pgfsetfillopacity{0.500000}%
\pgfsetlinewidth{0.250937pt}%
\definecolor{currentstroke}{rgb}{0.000000,0.000000,0.000000}%
\pgfsetstrokecolor{currentstroke}%
\pgfsetstrokeopacity{0.500000}%
\pgfsetdash{}{0pt}%
\pgfsys@defobject{currentmarker}{\pgfqpoint{-0.008333in}{-0.008333in}}{\pgfqpoint{0.008333in}{0.008333in}}{%
\pgfpathmoveto{\pgfqpoint{0.000000in}{-0.008333in}}%
\pgfpathcurveto{\pgfqpoint{0.002210in}{-0.008333in}}{\pgfqpoint{0.004330in}{-0.007455in}}{\pgfqpoint{0.005893in}{-0.005893in}}%
\pgfpathcurveto{\pgfqpoint{0.007455in}{-0.004330in}}{\pgfqpoint{0.008333in}{-0.002210in}}{\pgfqpoint{0.008333in}{0.000000in}}%
\pgfpathcurveto{\pgfqpoint{0.008333in}{0.002210in}}{\pgfqpoint{0.007455in}{0.004330in}}{\pgfqpoint{0.005893in}{0.005893in}}%
\pgfpathcurveto{\pgfqpoint{0.004330in}{0.007455in}}{\pgfqpoint{0.002210in}{0.008333in}}{\pgfqpoint{0.000000in}{0.008333in}}%
\pgfpathcurveto{\pgfqpoint{-0.002210in}{0.008333in}}{\pgfqpoint{-0.004330in}{0.007455in}}{\pgfqpoint{-0.005893in}{0.005893in}}%
\pgfpathcurveto{\pgfqpoint{-0.007455in}{0.004330in}}{\pgfqpoint{-0.008333in}{0.002210in}}{\pgfqpoint{-0.008333in}{0.000000in}}%
\pgfpathcurveto{\pgfqpoint{-0.008333in}{-0.002210in}}{\pgfqpoint{-0.007455in}{-0.004330in}}{\pgfqpoint{-0.005893in}{-0.005893in}}%
\pgfpathcurveto{\pgfqpoint{-0.004330in}{-0.007455in}}{\pgfqpoint{-0.002210in}{-0.008333in}}{\pgfqpoint{0.000000in}{-0.008333in}}%
\pgfpathlineto{\pgfqpoint{0.000000in}{-0.008333in}}%
\pgfpathclose%
\pgfusepath{stroke,fill}%
}%
\begin{pgfscope}%
\pgfsys@transformshift{2.532736in}{3.454442in}%
\pgfsys@useobject{currentmarker}{}%
\end{pgfscope}%
\end{pgfscope}%
\begin{pgfscope}%
\pgfpathrectangle{\pgfqpoint{0.100000in}{2.493063in}}{\pgfqpoint{5.037500in}{3.427208in}}%
\pgfusepath{clip}%
\pgfsetrectcap%
\pgfsetroundjoin%
\pgfsetlinewidth{1.505625pt}%
\definecolor{currentstroke}{rgb}{0.501961,0.501961,0.501961}%
\pgfsetstrokecolor{currentstroke}%
\pgfsetstrokeopacity{0.500000}%
\pgfsetdash{}{0pt}%
\pgfpathmoveto{\pgfqpoint{2.334733in}{3.040286in}}%
\pgfusepath{stroke}%
\end{pgfscope}%
\begin{pgfscope}%
\pgfpathrectangle{\pgfqpoint{0.100000in}{2.493063in}}{\pgfqpoint{5.037500in}{3.427208in}}%
\pgfusepath{clip}%
\pgfsetbuttcap%
\pgfsetroundjoin%
\definecolor{currentfill}{rgb}{0.501961,0.501961,0.501961}%
\pgfsetfillcolor{currentfill}%
\pgfsetfillopacity{0.500000}%
\pgfsetlinewidth{0.250937pt}%
\definecolor{currentstroke}{rgb}{0.000000,0.000000,0.000000}%
\pgfsetstrokecolor{currentstroke}%
\pgfsetstrokeopacity{0.500000}%
\pgfsetdash{}{0pt}%
\pgfsys@defobject{currentmarker}{\pgfqpoint{-0.013889in}{-0.013889in}}{\pgfqpoint{0.013889in}{0.013889in}}{%
\pgfpathmoveto{\pgfqpoint{0.000000in}{-0.013889in}}%
\pgfpathcurveto{\pgfqpoint{0.003683in}{-0.013889in}}{\pgfqpoint{0.007216in}{-0.012425in}}{\pgfqpoint{0.009821in}{-0.009821in}}%
\pgfpathcurveto{\pgfqpoint{0.012425in}{-0.007216in}}{\pgfqpoint{0.013889in}{-0.003683in}}{\pgfqpoint{0.013889in}{0.000000in}}%
\pgfpathcurveto{\pgfqpoint{0.013889in}{0.003683in}}{\pgfqpoint{0.012425in}{0.007216in}}{\pgfqpoint{0.009821in}{0.009821in}}%
\pgfpathcurveto{\pgfqpoint{0.007216in}{0.012425in}}{\pgfqpoint{0.003683in}{0.013889in}}{\pgfqpoint{0.000000in}{0.013889in}}%
\pgfpathcurveto{\pgfqpoint{-0.003683in}{0.013889in}}{\pgfqpoint{-0.007216in}{0.012425in}}{\pgfqpoint{-0.009821in}{0.009821in}}%
\pgfpathcurveto{\pgfqpoint{-0.012425in}{0.007216in}}{\pgfqpoint{-0.013889in}{0.003683in}}{\pgfqpoint{-0.013889in}{0.000000in}}%
\pgfpathcurveto{\pgfqpoint{-0.013889in}{-0.003683in}}{\pgfqpoint{-0.012425in}{-0.007216in}}{\pgfqpoint{-0.009821in}{-0.009821in}}%
\pgfpathcurveto{\pgfqpoint{-0.007216in}{-0.012425in}}{\pgfqpoint{-0.003683in}{-0.013889in}}{\pgfqpoint{0.000000in}{-0.013889in}}%
\pgfpathlineto{\pgfqpoint{0.000000in}{-0.013889in}}%
\pgfpathclose%
\pgfusepath{stroke,fill}%
}%
\begin{pgfscope}%
\pgfsys@transformshift{2.334733in}{3.040286in}%
\pgfsys@useobject{currentmarker}{}%
\end{pgfscope}%
\end{pgfscope}%
\begin{pgfscope}%
\pgfpathrectangle{\pgfqpoint{0.100000in}{2.493063in}}{\pgfqpoint{5.037500in}{3.427208in}}%
\pgfusepath{clip}%
\pgfsetrectcap%
\pgfsetroundjoin%
\pgfsetlinewidth{1.505625pt}%
\definecolor{currentstroke}{rgb}{0.501961,0.501961,0.501961}%
\pgfsetstrokecolor{currentstroke}%
\pgfsetstrokeopacity{0.500000}%
\pgfsetdash{}{0pt}%
\pgfpathmoveto{\pgfqpoint{2.829493in}{3.610973in}}%
\pgfusepath{stroke}%
\end{pgfscope}%
\begin{pgfscope}%
\pgfpathrectangle{\pgfqpoint{0.100000in}{2.493063in}}{\pgfqpoint{5.037500in}{3.427208in}}%
\pgfusepath{clip}%
\pgfsetbuttcap%
\pgfsetroundjoin%
\definecolor{currentfill}{rgb}{0.501961,0.501961,0.501961}%
\pgfsetfillcolor{currentfill}%
\pgfsetfillopacity{0.500000}%
\pgfsetlinewidth{0.250937pt}%
\definecolor{currentstroke}{rgb}{0.000000,0.000000,0.000000}%
\pgfsetstrokecolor{currentstroke}%
\pgfsetstrokeopacity{0.500000}%
\pgfsetdash{}{0pt}%
\pgfsys@defobject{currentmarker}{\pgfqpoint{-0.013889in}{-0.013889in}}{\pgfqpoint{0.013889in}{0.013889in}}{%
\pgfpathmoveto{\pgfqpoint{0.000000in}{-0.013889in}}%
\pgfpathcurveto{\pgfqpoint{0.003683in}{-0.013889in}}{\pgfqpoint{0.007216in}{-0.012425in}}{\pgfqpoint{0.009821in}{-0.009821in}}%
\pgfpathcurveto{\pgfqpoint{0.012425in}{-0.007216in}}{\pgfqpoint{0.013889in}{-0.003683in}}{\pgfqpoint{0.013889in}{0.000000in}}%
\pgfpathcurveto{\pgfqpoint{0.013889in}{0.003683in}}{\pgfqpoint{0.012425in}{0.007216in}}{\pgfqpoint{0.009821in}{0.009821in}}%
\pgfpathcurveto{\pgfqpoint{0.007216in}{0.012425in}}{\pgfqpoint{0.003683in}{0.013889in}}{\pgfqpoint{0.000000in}{0.013889in}}%
\pgfpathcurveto{\pgfqpoint{-0.003683in}{0.013889in}}{\pgfqpoint{-0.007216in}{0.012425in}}{\pgfqpoint{-0.009821in}{0.009821in}}%
\pgfpathcurveto{\pgfqpoint{-0.012425in}{0.007216in}}{\pgfqpoint{-0.013889in}{0.003683in}}{\pgfqpoint{-0.013889in}{0.000000in}}%
\pgfpathcurveto{\pgfqpoint{-0.013889in}{-0.003683in}}{\pgfqpoint{-0.012425in}{-0.007216in}}{\pgfqpoint{-0.009821in}{-0.009821in}}%
\pgfpathcurveto{\pgfqpoint{-0.007216in}{-0.012425in}}{\pgfqpoint{-0.003683in}{-0.013889in}}{\pgfqpoint{0.000000in}{-0.013889in}}%
\pgfpathlineto{\pgfqpoint{0.000000in}{-0.013889in}}%
\pgfpathclose%
\pgfusepath{stroke,fill}%
}%
\begin{pgfscope}%
\pgfsys@transformshift{2.829493in}{3.610973in}%
\pgfsys@useobject{currentmarker}{}%
\end{pgfscope}%
\end{pgfscope}%
\begin{pgfscope}%
\pgfpathrectangle{\pgfqpoint{0.100000in}{2.493063in}}{\pgfqpoint{5.037500in}{3.427208in}}%
\pgfusepath{clip}%
\pgfsetrectcap%
\pgfsetroundjoin%
\pgfsetlinewidth{1.505625pt}%
\definecolor{currentstroke}{rgb}{0.501961,0.501961,0.501961}%
\pgfsetstrokecolor{currentstroke}%
\pgfsetstrokeopacity{0.500000}%
\pgfsetdash{}{0pt}%
\pgfpathmoveto{\pgfqpoint{2.144143in}{3.760550in}}%
\pgfusepath{stroke}%
\end{pgfscope}%
\begin{pgfscope}%
\pgfpathrectangle{\pgfqpoint{0.100000in}{2.493063in}}{\pgfqpoint{5.037500in}{3.427208in}}%
\pgfusepath{clip}%
\pgfsetbuttcap%
\pgfsetroundjoin%
\definecolor{currentfill}{rgb}{0.501961,0.501961,0.501961}%
\pgfsetfillcolor{currentfill}%
\pgfsetfillopacity{0.500000}%
\pgfsetlinewidth{0.250937pt}%
\definecolor{currentstroke}{rgb}{0.000000,0.000000,0.000000}%
\pgfsetstrokecolor{currentstroke}%
\pgfsetstrokeopacity{0.500000}%
\pgfsetdash{}{0pt}%
\pgfsys@defobject{currentmarker}{\pgfqpoint{-0.013889in}{-0.013889in}}{\pgfqpoint{0.013889in}{0.013889in}}{%
\pgfpathmoveto{\pgfqpoint{0.000000in}{-0.013889in}}%
\pgfpathcurveto{\pgfqpoint{0.003683in}{-0.013889in}}{\pgfqpoint{0.007216in}{-0.012425in}}{\pgfqpoint{0.009821in}{-0.009821in}}%
\pgfpathcurveto{\pgfqpoint{0.012425in}{-0.007216in}}{\pgfqpoint{0.013889in}{-0.003683in}}{\pgfqpoint{0.013889in}{0.000000in}}%
\pgfpathcurveto{\pgfqpoint{0.013889in}{0.003683in}}{\pgfqpoint{0.012425in}{0.007216in}}{\pgfqpoint{0.009821in}{0.009821in}}%
\pgfpathcurveto{\pgfqpoint{0.007216in}{0.012425in}}{\pgfqpoint{0.003683in}{0.013889in}}{\pgfqpoint{0.000000in}{0.013889in}}%
\pgfpathcurveto{\pgfqpoint{-0.003683in}{0.013889in}}{\pgfqpoint{-0.007216in}{0.012425in}}{\pgfqpoint{-0.009821in}{0.009821in}}%
\pgfpathcurveto{\pgfqpoint{-0.012425in}{0.007216in}}{\pgfqpoint{-0.013889in}{0.003683in}}{\pgfqpoint{-0.013889in}{0.000000in}}%
\pgfpathcurveto{\pgfqpoint{-0.013889in}{-0.003683in}}{\pgfqpoint{-0.012425in}{-0.007216in}}{\pgfqpoint{-0.009821in}{-0.009821in}}%
\pgfpathcurveto{\pgfqpoint{-0.007216in}{-0.012425in}}{\pgfqpoint{-0.003683in}{-0.013889in}}{\pgfqpoint{0.000000in}{-0.013889in}}%
\pgfpathlineto{\pgfqpoint{0.000000in}{-0.013889in}}%
\pgfpathclose%
\pgfusepath{stroke,fill}%
}%
\begin{pgfscope}%
\pgfsys@transformshift{2.144143in}{3.760550in}%
\pgfsys@useobject{currentmarker}{}%
\end{pgfscope}%
\end{pgfscope}%
\begin{pgfscope}%
\pgfpathrectangle{\pgfqpoint{0.100000in}{2.493063in}}{\pgfqpoint{5.037500in}{3.427208in}}%
\pgfusepath{clip}%
\pgfsetrectcap%
\pgfsetroundjoin%
\pgfsetlinewidth{1.505625pt}%
\definecolor{currentstroke}{rgb}{0.678431,1.000000,0.184314}%
\pgfsetstrokecolor{currentstroke}%
\pgfsetstrokeopacity{0.500000}%
\pgfsetdash{}{0pt}%
\pgfpathmoveto{\pgfqpoint{2.462352in}{2.881522in}}%
\pgfusepath{stroke}%
\end{pgfscope}%
\begin{pgfscope}%
\pgfpathrectangle{\pgfqpoint{0.100000in}{2.493063in}}{\pgfqpoint{5.037500in}{3.427208in}}%
\pgfusepath{clip}%
\pgfsetbuttcap%
\pgfsetroundjoin%
\definecolor{currentfill}{rgb}{0.678431,1.000000,0.184314}%
\pgfsetfillcolor{currentfill}%
\pgfsetfillopacity{0.500000}%
\pgfsetlinewidth{0.250937pt}%
\definecolor{currentstroke}{rgb}{0.000000,0.000000,0.000000}%
\pgfsetstrokecolor{currentstroke}%
\pgfsetstrokeopacity{0.500000}%
\pgfsetdash{}{0pt}%
\pgfsys@defobject{currentmarker}{\pgfqpoint{-0.022222in}{-0.022222in}}{\pgfqpoint{0.022222in}{0.022222in}}{%
\pgfpathmoveto{\pgfqpoint{0.000000in}{-0.022222in}}%
\pgfpathcurveto{\pgfqpoint{0.005893in}{-0.022222in}}{\pgfqpoint{0.011546in}{-0.019881in}}{\pgfqpoint{0.015713in}{-0.015713in}}%
\pgfpathcurveto{\pgfqpoint{0.019881in}{-0.011546in}}{\pgfqpoint{0.022222in}{-0.005893in}}{\pgfqpoint{0.022222in}{0.000000in}}%
\pgfpathcurveto{\pgfqpoint{0.022222in}{0.005893in}}{\pgfqpoint{0.019881in}{0.011546in}}{\pgfqpoint{0.015713in}{0.015713in}}%
\pgfpathcurveto{\pgfqpoint{0.011546in}{0.019881in}}{\pgfqpoint{0.005893in}{0.022222in}}{\pgfqpoint{0.000000in}{0.022222in}}%
\pgfpathcurveto{\pgfqpoint{-0.005893in}{0.022222in}}{\pgfqpoint{-0.011546in}{0.019881in}}{\pgfqpoint{-0.015713in}{0.015713in}}%
\pgfpathcurveto{\pgfqpoint{-0.019881in}{0.011546in}}{\pgfqpoint{-0.022222in}{0.005893in}}{\pgfqpoint{-0.022222in}{0.000000in}}%
\pgfpathcurveto{\pgfqpoint{-0.022222in}{-0.005893in}}{\pgfqpoint{-0.019881in}{-0.011546in}}{\pgfqpoint{-0.015713in}{-0.015713in}}%
\pgfpathcurveto{\pgfqpoint{-0.011546in}{-0.019881in}}{\pgfqpoint{-0.005893in}{-0.022222in}}{\pgfqpoint{0.000000in}{-0.022222in}}%
\pgfpathlineto{\pgfqpoint{0.000000in}{-0.022222in}}%
\pgfpathclose%
\pgfusepath{stroke,fill}%
}%
\begin{pgfscope}%
\pgfsys@transformshift{2.462352in}{2.881522in}%
\pgfsys@useobject{currentmarker}{}%
\end{pgfscope}%
\end{pgfscope}%
\begin{pgfscope}%
\pgfpathrectangle{\pgfqpoint{0.100000in}{2.493063in}}{\pgfqpoint{5.037500in}{3.427208in}}%
\pgfusepath{clip}%
\pgfsetrectcap%
\pgfsetroundjoin%
\pgfsetlinewidth{1.505625pt}%
\definecolor{currentstroke}{rgb}{0.501961,0.501961,0.501961}%
\pgfsetstrokecolor{currentstroke}%
\pgfsetstrokeopacity{0.500000}%
\pgfsetdash{}{0pt}%
\pgfpathmoveto{\pgfqpoint{2.108493in}{3.579680in}}%
\pgfusepath{stroke}%
\end{pgfscope}%
\begin{pgfscope}%
\pgfpathrectangle{\pgfqpoint{0.100000in}{2.493063in}}{\pgfqpoint{5.037500in}{3.427208in}}%
\pgfusepath{clip}%
\pgfsetbuttcap%
\pgfsetroundjoin%
\definecolor{currentfill}{rgb}{0.501961,0.501961,0.501961}%
\pgfsetfillcolor{currentfill}%
\pgfsetfillopacity{0.500000}%
\pgfsetlinewidth{0.250937pt}%
\definecolor{currentstroke}{rgb}{0.000000,0.000000,0.000000}%
\pgfsetstrokecolor{currentstroke}%
\pgfsetstrokeopacity{0.500000}%
\pgfsetdash{}{0pt}%
\pgfsys@defobject{currentmarker}{\pgfqpoint{-0.013889in}{-0.013889in}}{\pgfqpoint{0.013889in}{0.013889in}}{%
\pgfpathmoveto{\pgfqpoint{0.000000in}{-0.013889in}}%
\pgfpathcurveto{\pgfqpoint{0.003683in}{-0.013889in}}{\pgfqpoint{0.007216in}{-0.012425in}}{\pgfqpoint{0.009821in}{-0.009821in}}%
\pgfpathcurveto{\pgfqpoint{0.012425in}{-0.007216in}}{\pgfqpoint{0.013889in}{-0.003683in}}{\pgfqpoint{0.013889in}{0.000000in}}%
\pgfpathcurveto{\pgfqpoint{0.013889in}{0.003683in}}{\pgfqpoint{0.012425in}{0.007216in}}{\pgfqpoint{0.009821in}{0.009821in}}%
\pgfpathcurveto{\pgfqpoint{0.007216in}{0.012425in}}{\pgfqpoint{0.003683in}{0.013889in}}{\pgfqpoint{0.000000in}{0.013889in}}%
\pgfpathcurveto{\pgfqpoint{-0.003683in}{0.013889in}}{\pgfqpoint{-0.007216in}{0.012425in}}{\pgfqpoint{-0.009821in}{0.009821in}}%
\pgfpathcurveto{\pgfqpoint{-0.012425in}{0.007216in}}{\pgfqpoint{-0.013889in}{0.003683in}}{\pgfqpoint{-0.013889in}{0.000000in}}%
\pgfpathcurveto{\pgfqpoint{-0.013889in}{-0.003683in}}{\pgfqpoint{-0.012425in}{-0.007216in}}{\pgfqpoint{-0.009821in}{-0.009821in}}%
\pgfpathcurveto{\pgfqpoint{-0.007216in}{-0.012425in}}{\pgfqpoint{-0.003683in}{-0.013889in}}{\pgfqpoint{0.000000in}{-0.013889in}}%
\pgfpathlineto{\pgfqpoint{0.000000in}{-0.013889in}}%
\pgfpathclose%
\pgfusepath{stroke,fill}%
}%
\begin{pgfscope}%
\pgfsys@transformshift{2.108493in}{3.579680in}%
\pgfsys@useobject{currentmarker}{}%
\end{pgfscope}%
\end{pgfscope}%
\begin{pgfscope}%
\pgfpathrectangle{\pgfqpoint{0.100000in}{2.493063in}}{\pgfqpoint{5.037500in}{3.427208in}}%
\pgfusepath{clip}%
\pgfsetrectcap%
\pgfsetroundjoin%
\pgfsetlinewidth{1.505625pt}%
\definecolor{currentstroke}{rgb}{0.501961,0.501961,0.501961}%
\pgfsetstrokecolor{currentstroke}%
\pgfsetstrokeopacity{0.500000}%
\pgfsetdash{}{0pt}%
\pgfpathmoveto{\pgfqpoint{2.078653in}{3.564405in}}%
\pgfusepath{stroke}%
\end{pgfscope}%
\begin{pgfscope}%
\pgfpathrectangle{\pgfqpoint{0.100000in}{2.493063in}}{\pgfqpoint{5.037500in}{3.427208in}}%
\pgfusepath{clip}%
\pgfsetbuttcap%
\pgfsetroundjoin%
\definecolor{currentfill}{rgb}{0.501961,0.501961,0.501961}%
\pgfsetfillcolor{currentfill}%
\pgfsetfillopacity{0.500000}%
\pgfsetlinewidth{0.250937pt}%
\definecolor{currentstroke}{rgb}{0.000000,0.000000,0.000000}%
\pgfsetstrokecolor{currentstroke}%
\pgfsetstrokeopacity{0.500000}%
\pgfsetdash{}{0pt}%
\pgfsys@defobject{currentmarker}{\pgfqpoint{-0.013889in}{-0.013889in}}{\pgfqpoint{0.013889in}{0.013889in}}{%
\pgfpathmoveto{\pgfqpoint{0.000000in}{-0.013889in}}%
\pgfpathcurveto{\pgfqpoint{0.003683in}{-0.013889in}}{\pgfqpoint{0.007216in}{-0.012425in}}{\pgfqpoint{0.009821in}{-0.009821in}}%
\pgfpathcurveto{\pgfqpoint{0.012425in}{-0.007216in}}{\pgfqpoint{0.013889in}{-0.003683in}}{\pgfqpoint{0.013889in}{0.000000in}}%
\pgfpathcurveto{\pgfqpoint{0.013889in}{0.003683in}}{\pgfqpoint{0.012425in}{0.007216in}}{\pgfqpoint{0.009821in}{0.009821in}}%
\pgfpathcurveto{\pgfqpoint{0.007216in}{0.012425in}}{\pgfqpoint{0.003683in}{0.013889in}}{\pgfqpoint{0.000000in}{0.013889in}}%
\pgfpathcurveto{\pgfqpoint{-0.003683in}{0.013889in}}{\pgfqpoint{-0.007216in}{0.012425in}}{\pgfqpoint{-0.009821in}{0.009821in}}%
\pgfpathcurveto{\pgfqpoint{-0.012425in}{0.007216in}}{\pgfqpoint{-0.013889in}{0.003683in}}{\pgfqpoint{-0.013889in}{0.000000in}}%
\pgfpathcurveto{\pgfqpoint{-0.013889in}{-0.003683in}}{\pgfqpoint{-0.012425in}{-0.007216in}}{\pgfqpoint{-0.009821in}{-0.009821in}}%
\pgfpathcurveto{\pgfqpoint{-0.007216in}{-0.012425in}}{\pgfqpoint{-0.003683in}{-0.013889in}}{\pgfqpoint{0.000000in}{-0.013889in}}%
\pgfpathlineto{\pgfqpoint{0.000000in}{-0.013889in}}%
\pgfpathclose%
\pgfusepath{stroke,fill}%
}%
\begin{pgfscope}%
\pgfsys@transformshift{2.078653in}{3.564405in}%
\pgfsys@useobject{currentmarker}{}%
\end{pgfscope}%
\end{pgfscope}%
\begin{pgfscope}%
\pgfpathrectangle{\pgfqpoint{0.100000in}{2.493063in}}{\pgfqpoint{5.037500in}{3.427208in}}%
\pgfusepath{clip}%
\pgfsetrectcap%
\pgfsetroundjoin%
\pgfsetlinewidth{1.505625pt}%
\definecolor{currentstroke}{rgb}{0.501961,0.501961,0.501961}%
\pgfsetstrokecolor{currentstroke}%
\pgfsetstrokeopacity{0.500000}%
\pgfsetdash{}{0pt}%
\pgfpathmoveto{\pgfqpoint{2.265570in}{3.506888in}}%
\pgfusepath{stroke}%
\end{pgfscope}%
\begin{pgfscope}%
\pgfpathrectangle{\pgfqpoint{0.100000in}{2.493063in}}{\pgfqpoint{5.037500in}{3.427208in}}%
\pgfusepath{clip}%
\pgfsetbuttcap%
\pgfsetroundjoin%
\definecolor{currentfill}{rgb}{0.501961,0.501961,0.501961}%
\pgfsetfillcolor{currentfill}%
\pgfsetfillopacity{0.500000}%
\pgfsetlinewidth{0.250937pt}%
\definecolor{currentstroke}{rgb}{0.000000,0.000000,0.000000}%
\pgfsetstrokecolor{currentstroke}%
\pgfsetstrokeopacity{0.500000}%
\pgfsetdash{}{0pt}%
\pgfsys@defobject{currentmarker}{\pgfqpoint{-0.013889in}{-0.013889in}}{\pgfqpoint{0.013889in}{0.013889in}}{%
\pgfpathmoveto{\pgfqpoint{0.000000in}{-0.013889in}}%
\pgfpathcurveto{\pgfqpoint{0.003683in}{-0.013889in}}{\pgfqpoint{0.007216in}{-0.012425in}}{\pgfqpoint{0.009821in}{-0.009821in}}%
\pgfpathcurveto{\pgfqpoint{0.012425in}{-0.007216in}}{\pgfqpoint{0.013889in}{-0.003683in}}{\pgfqpoint{0.013889in}{0.000000in}}%
\pgfpathcurveto{\pgfqpoint{0.013889in}{0.003683in}}{\pgfqpoint{0.012425in}{0.007216in}}{\pgfqpoint{0.009821in}{0.009821in}}%
\pgfpathcurveto{\pgfqpoint{0.007216in}{0.012425in}}{\pgfqpoint{0.003683in}{0.013889in}}{\pgfqpoint{0.000000in}{0.013889in}}%
\pgfpathcurveto{\pgfqpoint{-0.003683in}{0.013889in}}{\pgfqpoint{-0.007216in}{0.012425in}}{\pgfqpoint{-0.009821in}{0.009821in}}%
\pgfpathcurveto{\pgfqpoint{-0.012425in}{0.007216in}}{\pgfqpoint{-0.013889in}{0.003683in}}{\pgfqpoint{-0.013889in}{0.000000in}}%
\pgfpathcurveto{\pgfqpoint{-0.013889in}{-0.003683in}}{\pgfqpoint{-0.012425in}{-0.007216in}}{\pgfqpoint{-0.009821in}{-0.009821in}}%
\pgfpathcurveto{\pgfqpoint{-0.007216in}{-0.012425in}}{\pgfqpoint{-0.003683in}{-0.013889in}}{\pgfqpoint{0.000000in}{-0.013889in}}%
\pgfpathlineto{\pgfqpoint{0.000000in}{-0.013889in}}%
\pgfpathclose%
\pgfusepath{stroke,fill}%
}%
\begin{pgfscope}%
\pgfsys@transformshift{2.265570in}{3.506888in}%
\pgfsys@useobject{currentmarker}{}%
\end{pgfscope}%
\end{pgfscope}%
\begin{pgfscope}%
\pgfpathrectangle{\pgfqpoint{0.100000in}{2.493063in}}{\pgfqpoint{5.037500in}{3.427208in}}%
\pgfusepath{clip}%
\pgfsetrectcap%
\pgfsetroundjoin%
\pgfsetlinewidth{1.505625pt}%
\definecolor{currentstroke}{rgb}{0.000000,0.000000,1.000000}%
\pgfsetstrokecolor{currentstroke}%
\pgfsetstrokeopacity{0.500000}%
\pgfsetdash{}{0pt}%
\pgfpathmoveto{\pgfqpoint{2.448846in}{3.260064in}}%
\pgfusepath{stroke}%
\end{pgfscope}%
\begin{pgfscope}%
\pgfpathrectangle{\pgfqpoint{0.100000in}{2.493063in}}{\pgfqpoint{5.037500in}{3.427208in}}%
\pgfusepath{clip}%
\pgfsetbuttcap%
\pgfsetroundjoin%
\definecolor{currentfill}{rgb}{0.000000,0.000000,1.000000}%
\pgfsetfillcolor{currentfill}%
\pgfsetfillopacity{0.500000}%
\pgfsetlinewidth{0.250937pt}%
\definecolor{currentstroke}{rgb}{0.000000,0.000000,0.000000}%
\pgfsetstrokecolor{currentstroke}%
\pgfsetstrokeopacity{0.500000}%
\pgfsetdash{}{0pt}%
\pgfsys@defobject{currentmarker}{\pgfqpoint{-0.005556in}{-0.005556in}}{\pgfqpoint{0.005556in}{0.005556in}}{%
\pgfpathmoveto{\pgfqpoint{0.000000in}{-0.005556in}}%
\pgfpathcurveto{\pgfqpoint{0.001473in}{-0.005556in}}{\pgfqpoint{0.002887in}{-0.004970in}}{\pgfqpoint{0.003928in}{-0.003928in}}%
\pgfpathcurveto{\pgfqpoint{0.004970in}{-0.002887in}}{\pgfqpoint{0.005556in}{-0.001473in}}{\pgfqpoint{0.005556in}{0.000000in}}%
\pgfpathcurveto{\pgfqpoint{0.005556in}{0.001473in}}{\pgfqpoint{0.004970in}{0.002887in}}{\pgfqpoint{0.003928in}{0.003928in}}%
\pgfpathcurveto{\pgfqpoint{0.002887in}{0.004970in}}{\pgfqpoint{0.001473in}{0.005556in}}{\pgfqpoint{0.000000in}{0.005556in}}%
\pgfpathcurveto{\pgfqpoint{-0.001473in}{0.005556in}}{\pgfqpoint{-0.002887in}{0.004970in}}{\pgfqpoint{-0.003928in}{0.003928in}}%
\pgfpathcurveto{\pgfqpoint{-0.004970in}{0.002887in}}{\pgfqpoint{-0.005556in}{0.001473in}}{\pgfqpoint{-0.005556in}{0.000000in}}%
\pgfpathcurveto{\pgfqpoint{-0.005556in}{-0.001473in}}{\pgfqpoint{-0.004970in}{-0.002887in}}{\pgfqpoint{-0.003928in}{-0.003928in}}%
\pgfpathcurveto{\pgfqpoint{-0.002887in}{-0.004970in}}{\pgfqpoint{-0.001473in}{-0.005556in}}{\pgfqpoint{0.000000in}{-0.005556in}}%
\pgfpathlineto{\pgfqpoint{0.000000in}{-0.005556in}}%
\pgfpathclose%
\pgfusepath{stroke,fill}%
}%
\begin{pgfscope}%
\pgfsys@transformshift{2.448846in}{3.260064in}%
\pgfsys@useobject{currentmarker}{}%
\end{pgfscope}%
\end{pgfscope}%
\begin{pgfscope}%
\pgfpathrectangle{\pgfqpoint{0.100000in}{2.493063in}}{\pgfqpoint{5.037500in}{3.427208in}}%
\pgfusepath{clip}%
\pgfsetrectcap%
\pgfsetroundjoin%
\pgfsetlinewidth{1.505625pt}%
\definecolor{currentstroke}{rgb}{0.000000,0.000000,1.000000}%
\pgfsetstrokecolor{currentstroke}%
\pgfsetstrokeopacity{0.500000}%
\pgfsetdash{}{0pt}%
\pgfpathmoveto{\pgfqpoint{2.649125in}{3.743744in}}%
\pgfusepath{stroke}%
\end{pgfscope}%
\begin{pgfscope}%
\pgfpathrectangle{\pgfqpoint{0.100000in}{2.493063in}}{\pgfqpoint{5.037500in}{3.427208in}}%
\pgfusepath{clip}%
\pgfsetbuttcap%
\pgfsetroundjoin%
\definecolor{currentfill}{rgb}{0.000000,0.000000,1.000000}%
\pgfsetfillcolor{currentfill}%
\pgfsetfillopacity{0.500000}%
\pgfsetlinewidth{0.250937pt}%
\definecolor{currentstroke}{rgb}{0.000000,0.000000,0.000000}%
\pgfsetstrokecolor{currentstroke}%
\pgfsetstrokeopacity{0.500000}%
\pgfsetdash{}{0pt}%
\pgfsys@defobject{currentmarker}{\pgfqpoint{-0.008333in}{-0.008333in}}{\pgfqpoint{0.008333in}{0.008333in}}{%
\pgfpathmoveto{\pgfqpoint{0.000000in}{-0.008333in}}%
\pgfpathcurveto{\pgfqpoint{0.002210in}{-0.008333in}}{\pgfqpoint{0.004330in}{-0.007455in}}{\pgfqpoint{0.005893in}{-0.005893in}}%
\pgfpathcurveto{\pgfqpoint{0.007455in}{-0.004330in}}{\pgfqpoint{0.008333in}{-0.002210in}}{\pgfqpoint{0.008333in}{0.000000in}}%
\pgfpathcurveto{\pgfqpoint{0.008333in}{0.002210in}}{\pgfqpoint{0.007455in}{0.004330in}}{\pgfqpoint{0.005893in}{0.005893in}}%
\pgfpathcurveto{\pgfqpoint{0.004330in}{0.007455in}}{\pgfqpoint{0.002210in}{0.008333in}}{\pgfqpoint{0.000000in}{0.008333in}}%
\pgfpathcurveto{\pgfqpoint{-0.002210in}{0.008333in}}{\pgfqpoint{-0.004330in}{0.007455in}}{\pgfqpoint{-0.005893in}{0.005893in}}%
\pgfpathcurveto{\pgfqpoint{-0.007455in}{0.004330in}}{\pgfqpoint{-0.008333in}{0.002210in}}{\pgfqpoint{-0.008333in}{0.000000in}}%
\pgfpathcurveto{\pgfqpoint{-0.008333in}{-0.002210in}}{\pgfqpoint{-0.007455in}{-0.004330in}}{\pgfqpoint{-0.005893in}{-0.005893in}}%
\pgfpathcurveto{\pgfqpoint{-0.004330in}{-0.007455in}}{\pgfqpoint{-0.002210in}{-0.008333in}}{\pgfqpoint{0.000000in}{-0.008333in}}%
\pgfpathlineto{\pgfqpoint{0.000000in}{-0.008333in}}%
\pgfpathclose%
\pgfusepath{stroke,fill}%
}%
\begin{pgfscope}%
\pgfsys@transformshift{2.649125in}{3.743744in}%
\pgfsys@useobject{currentmarker}{}%
\end{pgfscope}%
\end{pgfscope}%
\begin{pgfscope}%
\pgfpathrectangle{\pgfqpoint{0.100000in}{2.493063in}}{\pgfqpoint{5.037500in}{3.427208in}}%
\pgfusepath{clip}%
\pgfsetrectcap%
\pgfsetroundjoin%
\pgfsetlinewidth{1.505625pt}%
\definecolor{currentstroke}{rgb}{0.678431,1.000000,0.184314}%
\pgfsetstrokecolor{currentstroke}%
\pgfsetstrokeopacity{0.500000}%
\pgfsetdash{}{0pt}%
\pgfpathmoveto{\pgfqpoint{2.896083in}{3.718458in}}%
\pgfusepath{stroke}%
\end{pgfscope}%
\begin{pgfscope}%
\pgfpathrectangle{\pgfqpoint{0.100000in}{2.493063in}}{\pgfqpoint{5.037500in}{3.427208in}}%
\pgfusepath{clip}%
\pgfsetbuttcap%
\pgfsetroundjoin%
\definecolor{currentfill}{rgb}{0.678431,1.000000,0.184314}%
\pgfsetfillcolor{currentfill}%
\pgfsetfillopacity{0.500000}%
\pgfsetlinewidth{0.250937pt}%
\definecolor{currentstroke}{rgb}{0.000000,0.000000,0.000000}%
\pgfsetstrokecolor{currentstroke}%
\pgfsetstrokeopacity{0.500000}%
\pgfsetdash{}{0pt}%
\pgfsys@defobject{currentmarker}{\pgfqpoint{-0.011111in}{-0.011111in}}{\pgfqpoint{0.011111in}{0.011111in}}{%
\pgfpathmoveto{\pgfqpoint{0.000000in}{-0.011111in}}%
\pgfpathcurveto{\pgfqpoint{0.002947in}{-0.011111in}}{\pgfqpoint{0.005773in}{-0.009940in}}{\pgfqpoint{0.007857in}{-0.007857in}}%
\pgfpathcurveto{\pgfqpoint{0.009940in}{-0.005773in}}{\pgfqpoint{0.011111in}{-0.002947in}}{\pgfqpoint{0.011111in}{0.000000in}}%
\pgfpathcurveto{\pgfqpoint{0.011111in}{0.002947in}}{\pgfqpoint{0.009940in}{0.005773in}}{\pgfqpoint{0.007857in}{0.007857in}}%
\pgfpathcurveto{\pgfqpoint{0.005773in}{0.009940in}}{\pgfqpoint{0.002947in}{0.011111in}}{\pgfqpoint{0.000000in}{0.011111in}}%
\pgfpathcurveto{\pgfqpoint{-0.002947in}{0.011111in}}{\pgfqpoint{-0.005773in}{0.009940in}}{\pgfqpoint{-0.007857in}{0.007857in}}%
\pgfpathcurveto{\pgfqpoint{-0.009940in}{0.005773in}}{\pgfqpoint{-0.011111in}{0.002947in}}{\pgfqpoint{-0.011111in}{0.000000in}}%
\pgfpathcurveto{\pgfqpoint{-0.011111in}{-0.002947in}}{\pgfqpoint{-0.009940in}{-0.005773in}}{\pgfqpoint{-0.007857in}{-0.007857in}}%
\pgfpathcurveto{\pgfqpoint{-0.005773in}{-0.009940in}}{\pgfqpoint{-0.002947in}{-0.011111in}}{\pgfqpoint{0.000000in}{-0.011111in}}%
\pgfpathlineto{\pgfqpoint{0.000000in}{-0.011111in}}%
\pgfpathclose%
\pgfusepath{stroke,fill}%
}%
\begin{pgfscope}%
\pgfsys@transformshift{2.896083in}{3.718458in}%
\pgfsys@useobject{currentmarker}{}%
\end{pgfscope}%
\end{pgfscope}%
\begin{pgfscope}%
\pgfpathrectangle{\pgfqpoint{0.100000in}{2.493063in}}{\pgfqpoint{5.037500in}{3.427208in}}%
\pgfusepath{clip}%
\pgfsetrectcap%
\pgfsetroundjoin%
\pgfsetlinewidth{1.505625pt}%
\definecolor{currentstroke}{rgb}{0.501961,0.501961,0.501961}%
\pgfsetstrokecolor{currentstroke}%
\pgfsetstrokeopacity{0.500000}%
\pgfsetdash{}{0pt}%
\pgfpathmoveto{\pgfqpoint{2.774645in}{3.593664in}}%
\pgfusepath{stroke}%
\end{pgfscope}%
\begin{pgfscope}%
\pgfpathrectangle{\pgfqpoint{0.100000in}{2.493063in}}{\pgfqpoint{5.037500in}{3.427208in}}%
\pgfusepath{clip}%
\pgfsetbuttcap%
\pgfsetroundjoin%
\definecolor{currentfill}{rgb}{0.501961,0.501961,0.501961}%
\pgfsetfillcolor{currentfill}%
\pgfsetfillopacity{0.500000}%
\pgfsetlinewidth{0.250937pt}%
\definecolor{currentstroke}{rgb}{0.000000,0.000000,0.000000}%
\pgfsetstrokecolor{currentstroke}%
\pgfsetstrokeopacity{0.500000}%
\pgfsetdash{}{0pt}%
\pgfsys@defobject{currentmarker}{\pgfqpoint{-0.013889in}{-0.013889in}}{\pgfqpoint{0.013889in}{0.013889in}}{%
\pgfpathmoveto{\pgfqpoint{0.000000in}{-0.013889in}}%
\pgfpathcurveto{\pgfqpoint{0.003683in}{-0.013889in}}{\pgfqpoint{0.007216in}{-0.012425in}}{\pgfqpoint{0.009821in}{-0.009821in}}%
\pgfpathcurveto{\pgfqpoint{0.012425in}{-0.007216in}}{\pgfqpoint{0.013889in}{-0.003683in}}{\pgfqpoint{0.013889in}{0.000000in}}%
\pgfpathcurveto{\pgfqpoint{0.013889in}{0.003683in}}{\pgfqpoint{0.012425in}{0.007216in}}{\pgfqpoint{0.009821in}{0.009821in}}%
\pgfpathcurveto{\pgfqpoint{0.007216in}{0.012425in}}{\pgfqpoint{0.003683in}{0.013889in}}{\pgfqpoint{0.000000in}{0.013889in}}%
\pgfpathcurveto{\pgfqpoint{-0.003683in}{0.013889in}}{\pgfqpoint{-0.007216in}{0.012425in}}{\pgfqpoint{-0.009821in}{0.009821in}}%
\pgfpathcurveto{\pgfqpoint{-0.012425in}{0.007216in}}{\pgfqpoint{-0.013889in}{0.003683in}}{\pgfqpoint{-0.013889in}{0.000000in}}%
\pgfpathcurveto{\pgfqpoint{-0.013889in}{-0.003683in}}{\pgfqpoint{-0.012425in}{-0.007216in}}{\pgfqpoint{-0.009821in}{-0.009821in}}%
\pgfpathcurveto{\pgfqpoint{-0.007216in}{-0.012425in}}{\pgfqpoint{-0.003683in}{-0.013889in}}{\pgfqpoint{0.000000in}{-0.013889in}}%
\pgfpathlineto{\pgfqpoint{0.000000in}{-0.013889in}}%
\pgfpathclose%
\pgfusepath{stroke,fill}%
}%
\begin{pgfscope}%
\pgfsys@transformshift{2.774645in}{3.593664in}%
\pgfsys@useobject{currentmarker}{}%
\end{pgfscope}%
\end{pgfscope}%
\begin{pgfscope}%
\pgfpathrectangle{\pgfqpoint{0.100000in}{2.493063in}}{\pgfqpoint{5.037500in}{3.427208in}}%
\pgfusepath{clip}%
\pgfsetrectcap%
\pgfsetroundjoin%
\pgfsetlinewidth{1.505625pt}%
\definecolor{currentstroke}{rgb}{0.501961,0.501961,0.501961}%
\pgfsetstrokecolor{currentstroke}%
\pgfsetstrokeopacity{0.500000}%
\pgfsetdash{}{0pt}%
\pgfpathmoveto{\pgfqpoint{2.598833in}{3.183133in}}%
\pgfusepath{stroke}%
\end{pgfscope}%
\begin{pgfscope}%
\pgfpathrectangle{\pgfqpoint{0.100000in}{2.493063in}}{\pgfqpoint{5.037500in}{3.427208in}}%
\pgfusepath{clip}%
\pgfsetbuttcap%
\pgfsetroundjoin%
\definecolor{currentfill}{rgb}{0.501961,0.501961,0.501961}%
\pgfsetfillcolor{currentfill}%
\pgfsetfillopacity{0.500000}%
\pgfsetlinewidth{0.250937pt}%
\definecolor{currentstroke}{rgb}{0.000000,0.000000,0.000000}%
\pgfsetstrokecolor{currentstroke}%
\pgfsetstrokeopacity{0.500000}%
\pgfsetdash{}{0pt}%
\pgfsys@defobject{currentmarker}{\pgfqpoint{-0.013889in}{-0.013889in}}{\pgfqpoint{0.013889in}{0.013889in}}{%
\pgfpathmoveto{\pgfqpoint{0.000000in}{-0.013889in}}%
\pgfpathcurveto{\pgfqpoint{0.003683in}{-0.013889in}}{\pgfqpoint{0.007216in}{-0.012425in}}{\pgfqpoint{0.009821in}{-0.009821in}}%
\pgfpathcurveto{\pgfqpoint{0.012425in}{-0.007216in}}{\pgfqpoint{0.013889in}{-0.003683in}}{\pgfqpoint{0.013889in}{0.000000in}}%
\pgfpathcurveto{\pgfqpoint{0.013889in}{0.003683in}}{\pgfqpoint{0.012425in}{0.007216in}}{\pgfqpoint{0.009821in}{0.009821in}}%
\pgfpathcurveto{\pgfqpoint{0.007216in}{0.012425in}}{\pgfqpoint{0.003683in}{0.013889in}}{\pgfqpoint{0.000000in}{0.013889in}}%
\pgfpathcurveto{\pgfqpoint{-0.003683in}{0.013889in}}{\pgfqpoint{-0.007216in}{0.012425in}}{\pgfqpoint{-0.009821in}{0.009821in}}%
\pgfpathcurveto{\pgfqpoint{-0.012425in}{0.007216in}}{\pgfqpoint{-0.013889in}{0.003683in}}{\pgfqpoint{-0.013889in}{0.000000in}}%
\pgfpathcurveto{\pgfqpoint{-0.013889in}{-0.003683in}}{\pgfqpoint{-0.012425in}{-0.007216in}}{\pgfqpoint{-0.009821in}{-0.009821in}}%
\pgfpathcurveto{\pgfqpoint{-0.007216in}{-0.012425in}}{\pgfqpoint{-0.003683in}{-0.013889in}}{\pgfqpoint{0.000000in}{-0.013889in}}%
\pgfpathlineto{\pgfqpoint{0.000000in}{-0.013889in}}%
\pgfpathclose%
\pgfusepath{stroke,fill}%
}%
\begin{pgfscope}%
\pgfsys@transformshift{2.598833in}{3.183133in}%
\pgfsys@useobject{currentmarker}{}%
\end{pgfscope}%
\end{pgfscope}%
\begin{pgfscope}%
\pgfpathrectangle{\pgfqpoint{0.100000in}{2.493063in}}{\pgfqpoint{5.037500in}{3.427208in}}%
\pgfusepath{clip}%
\pgfsetrectcap%
\pgfsetroundjoin%
\pgfsetlinewidth{1.505625pt}%
\definecolor{currentstroke}{rgb}{0.501961,0.501961,0.501961}%
\pgfsetstrokecolor{currentstroke}%
\pgfsetstrokeopacity{0.500000}%
\pgfsetdash{}{0pt}%
\pgfpathmoveto{\pgfqpoint{2.591727in}{3.503086in}}%
\pgfusepath{stroke}%
\end{pgfscope}%
\begin{pgfscope}%
\pgfpathrectangle{\pgfqpoint{0.100000in}{2.493063in}}{\pgfqpoint{5.037500in}{3.427208in}}%
\pgfusepath{clip}%
\pgfsetbuttcap%
\pgfsetroundjoin%
\definecolor{currentfill}{rgb}{0.501961,0.501961,0.501961}%
\pgfsetfillcolor{currentfill}%
\pgfsetfillopacity{0.500000}%
\pgfsetlinewidth{0.250937pt}%
\definecolor{currentstroke}{rgb}{0.000000,0.000000,0.000000}%
\pgfsetstrokecolor{currentstroke}%
\pgfsetstrokeopacity{0.500000}%
\pgfsetdash{}{0pt}%
\pgfsys@defobject{currentmarker}{\pgfqpoint{-0.013889in}{-0.013889in}}{\pgfqpoint{0.013889in}{0.013889in}}{%
\pgfpathmoveto{\pgfqpoint{0.000000in}{-0.013889in}}%
\pgfpathcurveto{\pgfqpoint{0.003683in}{-0.013889in}}{\pgfqpoint{0.007216in}{-0.012425in}}{\pgfqpoint{0.009821in}{-0.009821in}}%
\pgfpathcurveto{\pgfqpoint{0.012425in}{-0.007216in}}{\pgfqpoint{0.013889in}{-0.003683in}}{\pgfqpoint{0.013889in}{0.000000in}}%
\pgfpathcurveto{\pgfqpoint{0.013889in}{0.003683in}}{\pgfqpoint{0.012425in}{0.007216in}}{\pgfqpoint{0.009821in}{0.009821in}}%
\pgfpathcurveto{\pgfqpoint{0.007216in}{0.012425in}}{\pgfqpoint{0.003683in}{0.013889in}}{\pgfqpoint{0.000000in}{0.013889in}}%
\pgfpathcurveto{\pgfqpoint{-0.003683in}{0.013889in}}{\pgfqpoint{-0.007216in}{0.012425in}}{\pgfqpoint{-0.009821in}{0.009821in}}%
\pgfpathcurveto{\pgfqpoint{-0.012425in}{0.007216in}}{\pgfqpoint{-0.013889in}{0.003683in}}{\pgfqpoint{-0.013889in}{0.000000in}}%
\pgfpathcurveto{\pgfqpoint{-0.013889in}{-0.003683in}}{\pgfqpoint{-0.012425in}{-0.007216in}}{\pgfqpoint{-0.009821in}{-0.009821in}}%
\pgfpathcurveto{\pgfqpoint{-0.007216in}{-0.012425in}}{\pgfqpoint{-0.003683in}{-0.013889in}}{\pgfqpoint{0.000000in}{-0.013889in}}%
\pgfpathlineto{\pgfqpoint{0.000000in}{-0.013889in}}%
\pgfpathclose%
\pgfusepath{stroke,fill}%
}%
\begin{pgfscope}%
\pgfsys@transformshift{2.591727in}{3.503086in}%
\pgfsys@useobject{currentmarker}{}%
\end{pgfscope}%
\end{pgfscope}%
\begin{pgfscope}%
\pgfpathrectangle{\pgfqpoint{0.100000in}{2.493063in}}{\pgfqpoint{5.037500in}{3.427208in}}%
\pgfusepath{clip}%
\pgfsetrectcap%
\pgfsetroundjoin%
\pgfsetlinewidth{1.505625pt}%
\definecolor{currentstroke}{rgb}{0.000000,0.000000,1.000000}%
\pgfsetstrokecolor{currentstroke}%
\pgfsetstrokeopacity{0.500000}%
\pgfsetdash{}{0pt}%
\pgfpathmoveto{\pgfqpoint{2.469056in}{3.780922in}}%
\pgfusepath{stroke}%
\end{pgfscope}%
\begin{pgfscope}%
\pgfpathrectangle{\pgfqpoint{0.100000in}{2.493063in}}{\pgfqpoint{5.037500in}{3.427208in}}%
\pgfusepath{clip}%
\pgfsetbuttcap%
\pgfsetroundjoin%
\definecolor{currentfill}{rgb}{0.000000,0.000000,1.000000}%
\pgfsetfillcolor{currentfill}%
\pgfsetfillopacity{0.500000}%
\pgfsetlinewidth{0.250937pt}%
\definecolor{currentstroke}{rgb}{0.000000,0.000000,0.000000}%
\pgfsetstrokecolor{currentstroke}%
\pgfsetstrokeopacity{0.500000}%
\pgfsetdash{}{0pt}%
\pgfsys@defobject{currentmarker}{\pgfqpoint{-0.011111in}{-0.011111in}}{\pgfqpoint{0.011111in}{0.011111in}}{%
\pgfpathmoveto{\pgfqpoint{0.000000in}{-0.011111in}}%
\pgfpathcurveto{\pgfqpoint{0.002947in}{-0.011111in}}{\pgfqpoint{0.005773in}{-0.009940in}}{\pgfqpoint{0.007857in}{-0.007857in}}%
\pgfpathcurveto{\pgfqpoint{0.009940in}{-0.005773in}}{\pgfqpoint{0.011111in}{-0.002947in}}{\pgfqpoint{0.011111in}{0.000000in}}%
\pgfpathcurveto{\pgfqpoint{0.011111in}{0.002947in}}{\pgfqpoint{0.009940in}{0.005773in}}{\pgfqpoint{0.007857in}{0.007857in}}%
\pgfpathcurveto{\pgfqpoint{0.005773in}{0.009940in}}{\pgfqpoint{0.002947in}{0.011111in}}{\pgfqpoint{0.000000in}{0.011111in}}%
\pgfpathcurveto{\pgfqpoint{-0.002947in}{0.011111in}}{\pgfqpoint{-0.005773in}{0.009940in}}{\pgfqpoint{-0.007857in}{0.007857in}}%
\pgfpathcurveto{\pgfqpoint{-0.009940in}{0.005773in}}{\pgfqpoint{-0.011111in}{0.002947in}}{\pgfqpoint{-0.011111in}{0.000000in}}%
\pgfpathcurveto{\pgfqpoint{-0.011111in}{-0.002947in}}{\pgfqpoint{-0.009940in}{-0.005773in}}{\pgfqpoint{-0.007857in}{-0.007857in}}%
\pgfpathcurveto{\pgfqpoint{-0.005773in}{-0.009940in}}{\pgfqpoint{-0.002947in}{-0.011111in}}{\pgfqpoint{0.000000in}{-0.011111in}}%
\pgfpathlineto{\pgfqpoint{0.000000in}{-0.011111in}}%
\pgfpathclose%
\pgfusepath{stroke,fill}%
}%
\begin{pgfscope}%
\pgfsys@transformshift{2.469056in}{3.780922in}%
\pgfsys@useobject{currentmarker}{}%
\end{pgfscope}%
\end{pgfscope}%
\begin{pgfscope}%
\pgfpathrectangle{\pgfqpoint{0.100000in}{2.493063in}}{\pgfqpoint{5.037500in}{3.427208in}}%
\pgfusepath{clip}%
\pgfsetrectcap%
\pgfsetroundjoin%
\pgfsetlinewidth{1.505625pt}%
\definecolor{currentstroke}{rgb}{0.000000,0.000000,1.000000}%
\pgfsetstrokecolor{currentstroke}%
\pgfsetstrokeopacity{0.500000}%
\pgfsetdash{}{0pt}%
\pgfpathmoveto{\pgfqpoint{1.364191in}{4.809322in}}%
\pgfusepath{stroke}%
\end{pgfscope}%
\begin{pgfscope}%
\pgfpathrectangle{\pgfqpoint{0.100000in}{2.493063in}}{\pgfqpoint{5.037500in}{3.427208in}}%
\pgfusepath{clip}%
\pgfsetbuttcap%
\pgfsetroundjoin%
\definecolor{currentfill}{rgb}{0.000000,0.000000,1.000000}%
\pgfsetfillcolor{currentfill}%
\pgfsetfillopacity{0.500000}%
\pgfsetlinewidth{0.250937pt}%
\definecolor{currentstroke}{rgb}{0.000000,0.000000,0.000000}%
\pgfsetstrokecolor{currentstroke}%
\pgfsetstrokeopacity{0.500000}%
\pgfsetdash{}{0pt}%
\pgfsys@defobject{currentmarker}{\pgfqpoint{-0.013889in}{-0.013889in}}{\pgfqpoint{0.013889in}{0.013889in}}{%
\pgfpathmoveto{\pgfqpoint{0.000000in}{-0.013889in}}%
\pgfpathcurveto{\pgfqpoint{0.003683in}{-0.013889in}}{\pgfqpoint{0.007216in}{-0.012425in}}{\pgfqpoint{0.009821in}{-0.009821in}}%
\pgfpathcurveto{\pgfqpoint{0.012425in}{-0.007216in}}{\pgfqpoint{0.013889in}{-0.003683in}}{\pgfqpoint{0.013889in}{0.000000in}}%
\pgfpathcurveto{\pgfqpoint{0.013889in}{0.003683in}}{\pgfqpoint{0.012425in}{0.007216in}}{\pgfqpoint{0.009821in}{0.009821in}}%
\pgfpathcurveto{\pgfqpoint{0.007216in}{0.012425in}}{\pgfqpoint{0.003683in}{0.013889in}}{\pgfqpoint{0.000000in}{0.013889in}}%
\pgfpathcurveto{\pgfqpoint{-0.003683in}{0.013889in}}{\pgfqpoint{-0.007216in}{0.012425in}}{\pgfqpoint{-0.009821in}{0.009821in}}%
\pgfpathcurveto{\pgfqpoint{-0.012425in}{0.007216in}}{\pgfqpoint{-0.013889in}{0.003683in}}{\pgfqpoint{-0.013889in}{0.000000in}}%
\pgfpathcurveto{\pgfqpoint{-0.013889in}{-0.003683in}}{\pgfqpoint{-0.012425in}{-0.007216in}}{\pgfqpoint{-0.009821in}{-0.009821in}}%
\pgfpathcurveto{\pgfqpoint{-0.007216in}{-0.012425in}}{\pgfqpoint{-0.003683in}{-0.013889in}}{\pgfqpoint{0.000000in}{-0.013889in}}%
\pgfpathlineto{\pgfqpoint{0.000000in}{-0.013889in}}%
\pgfpathclose%
\pgfusepath{stroke,fill}%
}%
\begin{pgfscope}%
\pgfsys@transformshift{1.364191in}{4.809322in}%
\pgfsys@useobject{currentmarker}{}%
\end{pgfscope}%
\end{pgfscope}%
\begin{pgfscope}%
\pgfpathrectangle{\pgfqpoint{0.100000in}{2.493063in}}{\pgfqpoint{5.037500in}{3.427208in}}%
\pgfusepath{clip}%
\pgfsetrectcap%
\pgfsetroundjoin%
\pgfsetlinewidth{1.505625pt}%
\definecolor{currentstroke}{rgb}{0.000000,0.000000,1.000000}%
\pgfsetstrokecolor{currentstroke}%
\pgfsetstrokeopacity{0.500000}%
\pgfsetdash{}{0pt}%
\pgfpathmoveto{\pgfqpoint{1.341563in}{4.753966in}}%
\pgfusepath{stroke}%
\end{pgfscope}%
\begin{pgfscope}%
\pgfpathrectangle{\pgfqpoint{0.100000in}{2.493063in}}{\pgfqpoint{5.037500in}{3.427208in}}%
\pgfusepath{clip}%
\pgfsetbuttcap%
\pgfsetroundjoin%
\definecolor{currentfill}{rgb}{0.000000,0.000000,1.000000}%
\pgfsetfillcolor{currentfill}%
\pgfsetfillopacity{0.500000}%
\pgfsetlinewidth{0.250937pt}%
\definecolor{currentstroke}{rgb}{0.000000,0.000000,0.000000}%
\pgfsetstrokecolor{currentstroke}%
\pgfsetstrokeopacity{0.500000}%
\pgfsetdash{}{0pt}%
\pgfsys@defobject{currentmarker}{\pgfqpoint{-0.011111in}{-0.011111in}}{\pgfqpoint{0.011111in}{0.011111in}}{%
\pgfpathmoveto{\pgfqpoint{0.000000in}{-0.011111in}}%
\pgfpathcurveto{\pgfqpoint{0.002947in}{-0.011111in}}{\pgfqpoint{0.005773in}{-0.009940in}}{\pgfqpoint{0.007857in}{-0.007857in}}%
\pgfpathcurveto{\pgfqpoint{0.009940in}{-0.005773in}}{\pgfqpoint{0.011111in}{-0.002947in}}{\pgfqpoint{0.011111in}{0.000000in}}%
\pgfpathcurveto{\pgfqpoint{0.011111in}{0.002947in}}{\pgfqpoint{0.009940in}{0.005773in}}{\pgfqpoint{0.007857in}{0.007857in}}%
\pgfpathcurveto{\pgfqpoint{0.005773in}{0.009940in}}{\pgfqpoint{0.002947in}{0.011111in}}{\pgfqpoint{0.000000in}{0.011111in}}%
\pgfpathcurveto{\pgfqpoint{-0.002947in}{0.011111in}}{\pgfqpoint{-0.005773in}{0.009940in}}{\pgfqpoint{-0.007857in}{0.007857in}}%
\pgfpathcurveto{\pgfqpoint{-0.009940in}{0.005773in}}{\pgfqpoint{-0.011111in}{0.002947in}}{\pgfqpoint{-0.011111in}{0.000000in}}%
\pgfpathcurveto{\pgfqpoint{-0.011111in}{-0.002947in}}{\pgfqpoint{-0.009940in}{-0.005773in}}{\pgfqpoint{-0.007857in}{-0.007857in}}%
\pgfpathcurveto{\pgfqpoint{-0.005773in}{-0.009940in}}{\pgfqpoint{-0.002947in}{-0.011111in}}{\pgfqpoint{0.000000in}{-0.011111in}}%
\pgfpathlineto{\pgfqpoint{0.000000in}{-0.011111in}}%
\pgfpathclose%
\pgfusepath{stroke,fill}%
}%
\begin{pgfscope}%
\pgfsys@transformshift{1.341563in}{4.753966in}%
\pgfsys@useobject{currentmarker}{}%
\end{pgfscope}%
\end{pgfscope}%
\begin{pgfscope}%
\pgfpathrectangle{\pgfqpoint{0.100000in}{2.493063in}}{\pgfqpoint{5.037500in}{3.427208in}}%
\pgfusepath{clip}%
\pgfsetrectcap%
\pgfsetroundjoin%
\pgfsetlinewidth{1.505625pt}%
\definecolor{currentstroke}{rgb}{0.000000,0.000000,1.000000}%
\pgfsetstrokecolor{currentstroke}%
\pgfsetstrokeopacity{0.500000}%
\pgfsetdash{}{0pt}%
\pgfpathmoveto{\pgfqpoint{1.347653in}{4.636899in}}%
\pgfusepath{stroke}%
\end{pgfscope}%
\begin{pgfscope}%
\pgfpathrectangle{\pgfqpoint{0.100000in}{2.493063in}}{\pgfqpoint{5.037500in}{3.427208in}}%
\pgfusepath{clip}%
\pgfsetbuttcap%
\pgfsetroundjoin%
\definecolor{currentfill}{rgb}{0.000000,0.000000,1.000000}%
\pgfsetfillcolor{currentfill}%
\pgfsetfillopacity{0.500000}%
\pgfsetlinewidth{0.250937pt}%
\definecolor{currentstroke}{rgb}{0.000000,0.000000,0.000000}%
\pgfsetstrokecolor{currentstroke}%
\pgfsetstrokeopacity{0.500000}%
\pgfsetdash{}{0pt}%
\pgfsys@defobject{currentmarker}{\pgfqpoint{-0.019444in}{-0.019444in}}{\pgfqpoint{0.019444in}{0.019444in}}{%
\pgfpathmoveto{\pgfqpoint{0.000000in}{-0.019444in}}%
\pgfpathcurveto{\pgfqpoint{0.005157in}{-0.019444in}}{\pgfqpoint{0.010103in}{-0.017396in}}{\pgfqpoint{0.013749in}{-0.013749in}}%
\pgfpathcurveto{\pgfqpoint{0.017396in}{-0.010103in}}{\pgfqpoint{0.019444in}{-0.005157in}}{\pgfqpoint{0.019444in}{0.000000in}}%
\pgfpathcurveto{\pgfqpoint{0.019444in}{0.005157in}}{\pgfqpoint{0.017396in}{0.010103in}}{\pgfqpoint{0.013749in}{0.013749in}}%
\pgfpathcurveto{\pgfqpoint{0.010103in}{0.017396in}}{\pgfqpoint{0.005157in}{0.019444in}}{\pgfqpoint{0.000000in}{0.019444in}}%
\pgfpathcurveto{\pgfqpoint{-0.005157in}{0.019444in}}{\pgfqpoint{-0.010103in}{0.017396in}}{\pgfqpoint{-0.013749in}{0.013749in}}%
\pgfpathcurveto{\pgfqpoint{-0.017396in}{0.010103in}}{\pgfqpoint{-0.019444in}{0.005157in}}{\pgfqpoint{-0.019444in}{0.000000in}}%
\pgfpathcurveto{\pgfqpoint{-0.019444in}{-0.005157in}}{\pgfqpoint{-0.017396in}{-0.010103in}}{\pgfqpoint{-0.013749in}{-0.013749in}}%
\pgfpathcurveto{\pgfqpoint{-0.010103in}{-0.017396in}}{\pgfqpoint{-0.005157in}{-0.019444in}}{\pgfqpoint{0.000000in}{-0.019444in}}%
\pgfpathlineto{\pgfqpoint{0.000000in}{-0.019444in}}%
\pgfpathclose%
\pgfusepath{stroke,fill}%
}%
\begin{pgfscope}%
\pgfsys@transformshift{1.347653in}{4.636899in}%
\pgfsys@useobject{currentmarker}{}%
\end{pgfscope}%
\end{pgfscope}%
\begin{pgfscope}%
\pgfpathrectangle{\pgfqpoint{0.100000in}{2.493063in}}{\pgfqpoint{5.037500in}{3.427208in}}%
\pgfusepath{clip}%
\pgfsetrectcap%
\pgfsetroundjoin%
\pgfsetlinewidth{1.505625pt}%
\definecolor{currentstroke}{rgb}{0.000000,0.000000,1.000000}%
\pgfsetstrokecolor{currentstroke}%
\pgfsetstrokeopacity{0.500000}%
\pgfsetdash{}{0pt}%
\pgfpathmoveto{\pgfqpoint{1.108330in}{4.316603in}}%
\pgfusepath{stroke}%
\end{pgfscope}%
\begin{pgfscope}%
\pgfpathrectangle{\pgfqpoint{0.100000in}{2.493063in}}{\pgfqpoint{5.037500in}{3.427208in}}%
\pgfusepath{clip}%
\pgfsetbuttcap%
\pgfsetroundjoin%
\definecolor{currentfill}{rgb}{0.000000,0.000000,1.000000}%
\pgfsetfillcolor{currentfill}%
\pgfsetfillopacity{0.500000}%
\pgfsetlinewidth{0.250937pt}%
\definecolor{currentstroke}{rgb}{0.000000,0.000000,0.000000}%
\pgfsetstrokecolor{currentstroke}%
\pgfsetstrokeopacity{0.500000}%
\pgfsetdash{}{0pt}%
\pgfsys@defobject{currentmarker}{\pgfqpoint{-0.008333in}{-0.008333in}}{\pgfqpoint{0.008333in}{0.008333in}}{%
\pgfpathmoveto{\pgfqpoint{0.000000in}{-0.008333in}}%
\pgfpathcurveto{\pgfqpoint{0.002210in}{-0.008333in}}{\pgfqpoint{0.004330in}{-0.007455in}}{\pgfqpoint{0.005893in}{-0.005893in}}%
\pgfpathcurveto{\pgfqpoint{0.007455in}{-0.004330in}}{\pgfqpoint{0.008333in}{-0.002210in}}{\pgfqpoint{0.008333in}{0.000000in}}%
\pgfpathcurveto{\pgfqpoint{0.008333in}{0.002210in}}{\pgfqpoint{0.007455in}{0.004330in}}{\pgfqpoint{0.005893in}{0.005893in}}%
\pgfpathcurveto{\pgfqpoint{0.004330in}{0.007455in}}{\pgfqpoint{0.002210in}{0.008333in}}{\pgfqpoint{0.000000in}{0.008333in}}%
\pgfpathcurveto{\pgfqpoint{-0.002210in}{0.008333in}}{\pgfqpoint{-0.004330in}{0.007455in}}{\pgfqpoint{-0.005893in}{0.005893in}}%
\pgfpathcurveto{\pgfqpoint{-0.007455in}{0.004330in}}{\pgfqpoint{-0.008333in}{0.002210in}}{\pgfqpoint{-0.008333in}{0.000000in}}%
\pgfpathcurveto{\pgfqpoint{-0.008333in}{-0.002210in}}{\pgfqpoint{-0.007455in}{-0.004330in}}{\pgfqpoint{-0.005893in}{-0.005893in}}%
\pgfpathcurveto{\pgfqpoint{-0.004330in}{-0.007455in}}{\pgfqpoint{-0.002210in}{-0.008333in}}{\pgfqpoint{0.000000in}{-0.008333in}}%
\pgfpathlineto{\pgfqpoint{0.000000in}{-0.008333in}}%
\pgfpathclose%
\pgfusepath{stroke,fill}%
}%
\begin{pgfscope}%
\pgfsys@transformshift{1.108330in}{4.316603in}%
\pgfsys@useobject{currentmarker}{}%
\end{pgfscope}%
\end{pgfscope}%
\begin{pgfscope}%
\pgfpathrectangle{\pgfqpoint{0.100000in}{2.493063in}}{\pgfqpoint{5.037500in}{3.427208in}}%
\pgfusepath{clip}%
\pgfsetrectcap%
\pgfsetroundjoin%
\pgfsetlinewidth{1.505625pt}%
\definecolor{currentstroke}{rgb}{0.000000,0.000000,1.000000}%
\pgfsetstrokecolor{currentstroke}%
\pgfsetstrokeopacity{0.500000}%
\pgfsetdash{}{0pt}%
\pgfpathmoveto{\pgfqpoint{1.338959in}{4.700996in}}%
\pgfusepath{stroke}%
\end{pgfscope}%
\begin{pgfscope}%
\pgfpathrectangle{\pgfqpoint{0.100000in}{2.493063in}}{\pgfqpoint{5.037500in}{3.427208in}}%
\pgfusepath{clip}%
\pgfsetbuttcap%
\pgfsetroundjoin%
\definecolor{currentfill}{rgb}{0.000000,0.000000,1.000000}%
\pgfsetfillcolor{currentfill}%
\pgfsetfillopacity{0.500000}%
\pgfsetlinewidth{0.250937pt}%
\definecolor{currentstroke}{rgb}{0.000000,0.000000,0.000000}%
\pgfsetstrokecolor{currentstroke}%
\pgfsetstrokeopacity{0.500000}%
\pgfsetdash{}{0pt}%
\pgfsys@defobject{currentmarker}{\pgfqpoint{-0.013889in}{-0.013889in}}{\pgfqpoint{0.013889in}{0.013889in}}{%
\pgfpathmoveto{\pgfqpoint{0.000000in}{-0.013889in}}%
\pgfpathcurveto{\pgfqpoint{0.003683in}{-0.013889in}}{\pgfqpoint{0.007216in}{-0.012425in}}{\pgfqpoint{0.009821in}{-0.009821in}}%
\pgfpathcurveto{\pgfqpoint{0.012425in}{-0.007216in}}{\pgfqpoint{0.013889in}{-0.003683in}}{\pgfqpoint{0.013889in}{0.000000in}}%
\pgfpathcurveto{\pgfqpoint{0.013889in}{0.003683in}}{\pgfqpoint{0.012425in}{0.007216in}}{\pgfqpoint{0.009821in}{0.009821in}}%
\pgfpathcurveto{\pgfqpoint{0.007216in}{0.012425in}}{\pgfqpoint{0.003683in}{0.013889in}}{\pgfqpoint{0.000000in}{0.013889in}}%
\pgfpathcurveto{\pgfqpoint{-0.003683in}{0.013889in}}{\pgfqpoint{-0.007216in}{0.012425in}}{\pgfqpoint{-0.009821in}{0.009821in}}%
\pgfpathcurveto{\pgfqpoint{-0.012425in}{0.007216in}}{\pgfqpoint{-0.013889in}{0.003683in}}{\pgfqpoint{-0.013889in}{0.000000in}}%
\pgfpathcurveto{\pgfqpoint{-0.013889in}{-0.003683in}}{\pgfqpoint{-0.012425in}{-0.007216in}}{\pgfqpoint{-0.009821in}{-0.009821in}}%
\pgfpathcurveto{\pgfqpoint{-0.007216in}{-0.012425in}}{\pgfqpoint{-0.003683in}{-0.013889in}}{\pgfqpoint{0.000000in}{-0.013889in}}%
\pgfpathlineto{\pgfqpoint{0.000000in}{-0.013889in}}%
\pgfpathclose%
\pgfusepath{stroke,fill}%
}%
\begin{pgfscope}%
\pgfsys@transformshift{1.338959in}{4.700996in}%
\pgfsys@useobject{currentmarker}{}%
\end{pgfscope}%
\end{pgfscope}%
\begin{pgfscope}%
\pgfpathrectangle{\pgfqpoint{0.100000in}{2.493063in}}{\pgfqpoint{5.037500in}{3.427208in}}%
\pgfusepath{clip}%
\pgfsetrectcap%
\pgfsetroundjoin%
\pgfsetlinewidth{1.505625pt}%
\definecolor{currentstroke}{rgb}{0.501961,0.501961,0.501961}%
\pgfsetstrokecolor{currentstroke}%
\pgfsetstrokeopacity{0.500000}%
\pgfsetdash{}{0pt}%
\pgfpathmoveto{\pgfqpoint{4.585415in}{5.206583in}}%
\pgfusepath{stroke}%
\end{pgfscope}%
\begin{pgfscope}%
\pgfpathrectangle{\pgfqpoint{0.100000in}{2.493063in}}{\pgfqpoint{5.037500in}{3.427208in}}%
\pgfusepath{clip}%
\pgfsetbuttcap%
\pgfsetroundjoin%
\definecolor{currentfill}{rgb}{0.501961,0.501961,0.501961}%
\pgfsetfillcolor{currentfill}%
\pgfsetfillopacity{0.500000}%
\pgfsetlinewidth{0.250937pt}%
\definecolor{currentstroke}{rgb}{0.000000,0.000000,0.000000}%
\pgfsetstrokecolor{currentstroke}%
\pgfsetstrokeopacity{0.500000}%
\pgfsetdash{}{0pt}%
\pgfsys@defobject{currentmarker}{\pgfqpoint{-0.013889in}{-0.013889in}}{\pgfqpoint{0.013889in}{0.013889in}}{%
\pgfpathmoveto{\pgfqpoint{0.000000in}{-0.013889in}}%
\pgfpathcurveto{\pgfqpoint{0.003683in}{-0.013889in}}{\pgfqpoint{0.007216in}{-0.012425in}}{\pgfqpoint{0.009821in}{-0.009821in}}%
\pgfpathcurveto{\pgfqpoint{0.012425in}{-0.007216in}}{\pgfqpoint{0.013889in}{-0.003683in}}{\pgfqpoint{0.013889in}{0.000000in}}%
\pgfpathcurveto{\pgfqpoint{0.013889in}{0.003683in}}{\pgfqpoint{0.012425in}{0.007216in}}{\pgfqpoint{0.009821in}{0.009821in}}%
\pgfpathcurveto{\pgfqpoint{0.007216in}{0.012425in}}{\pgfqpoint{0.003683in}{0.013889in}}{\pgfqpoint{0.000000in}{0.013889in}}%
\pgfpathcurveto{\pgfqpoint{-0.003683in}{0.013889in}}{\pgfqpoint{-0.007216in}{0.012425in}}{\pgfqpoint{-0.009821in}{0.009821in}}%
\pgfpathcurveto{\pgfqpoint{-0.012425in}{0.007216in}}{\pgfqpoint{-0.013889in}{0.003683in}}{\pgfqpoint{-0.013889in}{0.000000in}}%
\pgfpathcurveto{\pgfqpoint{-0.013889in}{-0.003683in}}{\pgfqpoint{-0.012425in}{-0.007216in}}{\pgfqpoint{-0.009821in}{-0.009821in}}%
\pgfpathcurveto{\pgfqpoint{-0.007216in}{-0.012425in}}{\pgfqpoint{-0.003683in}{-0.013889in}}{\pgfqpoint{0.000000in}{-0.013889in}}%
\pgfpathlineto{\pgfqpoint{0.000000in}{-0.013889in}}%
\pgfpathclose%
\pgfusepath{stroke,fill}%
}%
\begin{pgfscope}%
\pgfsys@transformshift{4.585415in}{5.206583in}%
\pgfsys@useobject{currentmarker}{}%
\end{pgfscope}%
\end{pgfscope}%
\begin{pgfscope}%
\pgfpathrectangle{\pgfqpoint{0.100000in}{2.493063in}}{\pgfqpoint{5.037500in}{3.427208in}}%
\pgfusepath{clip}%
\pgfsetrectcap%
\pgfsetroundjoin%
\pgfsetlinewidth{1.505625pt}%
\definecolor{currentstroke}{rgb}{0.501961,0.501961,0.501961}%
\pgfsetstrokecolor{currentstroke}%
\pgfsetstrokeopacity{0.500000}%
\pgfsetdash{}{0pt}%
\pgfpathmoveto{\pgfqpoint{4.136391in}{4.264211in}}%
\pgfusepath{stroke}%
\end{pgfscope}%
\begin{pgfscope}%
\pgfpathrectangle{\pgfqpoint{0.100000in}{2.493063in}}{\pgfqpoint{5.037500in}{3.427208in}}%
\pgfusepath{clip}%
\pgfsetbuttcap%
\pgfsetroundjoin%
\definecolor{currentfill}{rgb}{0.501961,0.501961,0.501961}%
\pgfsetfillcolor{currentfill}%
\pgfsetfillopacity{0.500000}%
\pgfsetlinewidth{0.250937pt}%
\definecolor{currentstroke}{rgb}{0.000000,0.000000,0.000000}%
\pgfsetstrokecolor{currentstroke}%
\pgfsetstrokeopacity{0.500000}%
\pgfsetdash{}{0pt}%
\pgfsys@defobject{currentmarker}{\pgfqpoint{-0.013889in}{-0.013889in}}{\pgfqpoint{0.013889in}{0.013889in}}{%
\pgfpathmoveto{\pgfqpoint{0.000000in}{-0.013889in}}%
\pgfpathcurveto{\pgfqpoint{0.003683in}{-0.013889in}}{\pgfqpoint{0.007216in}{-0.012425in}}{\pgfqpoint{0.009821in}{-0.009821in}}%
\pgfpathcurveto{\pgfqpoint{0.012425in}{-0.007216in}}{\pgfqpoint{0.013889in}{-0.003683in}}{\pgfqpoint{0.013889in}{0.000000in}}%
\pgfpathcurveto{\pgfqpoint{0.013889in}{0.003683in}}{\pgfqpoint{0.012425in}{0.007216in}}{\pgfqpoint{0.009821in}{0.009821in}}%
\pgfpathcurveto{\pgfqpoint{0.007216in}{0.012425in}}{\pgfqpoint{0.003683in}{0.013889in}}{\pgfqpoint{0.000000in}{0.013889in}}%
\pgfpathcurveto{\pgfqpoint{-0.003683in}{0.013889in}}{\pgfqpoint{-0.007216in}{0.012425in}}{\pgfqpoint{-0.009821in}{0.009821in}}%
\pgfpathcurveto{\pgfqpoint{-0.012425in}{0.007216in}}{\pgfqpoint{-0.013889in}{0.003683in}}{\pgfqpoint{-0.013889in}{0.000000in}}%
\pgfpathcurveto{\pgfqpoint{-0.013889in}{-0.003683in}}{\pgfqpoint{-0.012425in}{-0.007216in}}{\pgfqpoint{-0.009821in}{-0.009821in}}%
\pgfpathcurveto{\pgfqpoint{-0.007216in}{-0.012425in}}{\pgfqpoint{-0.003683in}{-0.013889in}}{\pgfqpoint{0.000000in}{-0.013889in}}%
\pgfpathlineto{\pgfqpoint{0.000000in}{-0.013889in}}%
\pgfpathclose%
\pgfusepath{stroke,fill}%
}%
\begin{pgfscope}%
\pgfsys@transformshift{4.136391in}{4.264211in}%
\pgfsys@useobject{currentmarker}{}%
\end{pgfscope}%
\end{pgfscope}%
\begin{pgfscope}%
\pgfpathrectangle{\pgfqpoint{0.100000in}{2.493063in}}{\pgfqpoint{5.037500in}{3.427208in}}%
\pgfusepath{clip}%
\pgfsetrectcap%
\pgfsetroundjoin%
\pgfsetlinewidth{1.505625pt}%
\definecolor{currentstroke}{rgb}{0.000000,0.000000,1.000000}%
\pgfsetstrokecolor{currentstroke}%
\pgfsetstrokeopacity{0.500000}%
\pgfsetdash{}{0pt}%
\pgfpathmoveto{\pgfqpoint{4.294817in}{4.385064in}}%
\pgfusepath{stroke}%
\end{pgfscope}%
\begin{pgfscope}%
\pgfpathrectangle{\pgfqpoint{0.100000in}{2.493063in}}{\pgfqpoint{5.037500in}{3.427208in}}%
\pgfusepath{clip}%
\pgfsetbuttcap%
\pgfsetroundjoin%
\definecolor{currentfill}{rgb}{0.000000,0.000000,1.000000}%
\pgfsetfillcolor{currentfill}%
\pgfsetfillopacity{0.500000}%
\pgfsetlinewidth{0.250937pt}%
\definecolor{currentstroke}{rgb}{0.000000,0.000000,0.000000}%
\pgfsetstrokecolor{currentstroke}%
\pgfsetstrokeopacity{0.500000}%
\pgfsetdash{}{0pt}%
\pgfsys@defobject{currentmarker}{\pgfqpoint{-0.011111in}{-0.011111in}}{\pgfqpoint{0.011111in}{0.011111in}}{%
\pgfpathmoveto{\pgfqpoint{0.000000in}{-0.011111in}}%
\pgfpathcurveto{\pgfqpoint{0.002947in}{-0.011111in}}{\pgfqpoint{0.005773in}{-0.009940in}}{\pgfqpoint{0.007857in}{-0.007857in}}%
\pgfpathcurveto{\pgfqpoint{0.009940in}{-0.005773in}}{\pgfqpoint{0.011111in}{-0.002947in}}{\pgfqpoint{0.011111in}{0.000000in}}%
\pgfpathcurveto{\pgfqpoint{0.011111in}{0.002947in}}{\pgfqpoint{0.009940in}{0.005773in}}{\pgfqpoint{0.007857in}{0.007857in}}%
\pgfpathcurveto{\pgfqpoint{0.005773in}{0.009940in}}{\pgfqpoint{0.002947in}{0.011111in}}{\pgfqpoint{0.000000in}{0.011111in}}%
\pgfpathcurveto{\pgfqpoint{-0.002947in}{0.011111in}}{\pgfqpoint{-0.005773in}{0.009940in}}{\pgfqpoint{-0.007857in}{0.007857in}}%
\pgfpathcurveto{\pgfqpoint{-0.009940in}{0.005773in}}{\pgfqpoint{-0.011111in}{0.002947in}}{\pgfqpoint{-0.011111in}{0.000000in}}%
\pgfpathcurveto{\pgfqpoint{-0.011111in}{-0.002947in}}{\pgfqpoint{-0.009940in}{-0.005773in}}{\pgfqpoint{-0.007857in}{-0.007857in}}%
\pgfpathcurveto{\pgfqpoint{-0.005773in}{-0.009940in}}{\pgfqpoint{-0.002947in}{-0.011111in}}{\pgfqpoint{0.000000in}{-0.011111in}}%
\pgfpathlineto{\pgfqpoint{0.000000in}{-0.011111in}}%
\pgfpathclose%
\pgfusepath{stroke,fill}%
}%
\begin{pgfscope}%
\pgfsys@transformshift{4.294817in}{4.385064in}%
\pgfsys@useobject{currentmarker}{}%
\end{pgfscope}%
\end{pgfscope}%
\begin{pgfscope}%
\pgfpathrectangle{\pgfqpoint{0.100000in}{2.493063in}}{\pgfqpoint{5.037500in}{3.427208in}}%
\pgfusepath{clip}%
\pgfsetrectcap%
\pgfsetroundjoin%
\pgfsetlinewidth{1.505625pt}%
\definecolor{currentstroke}{rgb}{0.000000,0.000000,1.000000}%
\pgfsetstrokecolor{currentstroke}%
\pgfsetstrokeopacity{0.500000}%
\pgfsetdash{}{0pt}%
\pgfpathmoveto{\pgfqpoint{4.251269in}{4.426316in}}%
\pgfusepath{stroke}%
\end{pgfscope}%
\begin{pgfscope}%
\pgfpathrectangle{\pgfqpoint{0.100000in}{2.493063in}}{\pgfqpoint{5.037500in}{3.427208in}}%
\pgfusepath{clip}%
\pgfsetbuttcap%
\pgfsetroundjoin%
\definecolor{currentfill}{rgb}{0.000000,0.000000,1.000000}%
\pgfsetfillcolor{currentfill}%
\pgfsetfillopacity{0.500000}%
\pgfsetlinewidth{0.250937pt}%
\definecolor{currentstroke}{rgb}{0.000000,0.000000,0.000000}%
\pgfsetstrokecolor{currentstroke}%
\pgfsetstrokeopacity{0.500000}%
\pgfsetdash{}{0pt}%
\pgfsys@defobject{currentmarker}{\pgfqpoint{-0.013889in}{-0.013889in}}{\pgfqpoint{0.013889in}{0.013889in}}{%
\pgfpathmoveto{\pgfqpoint{0.000000in}{-0.013889in}}%
\pgfpathcurveto{\pgfqpoint{0.003683in}{-0.013889in}}{\pgfqpoint{0.007216in}{-0.012425in}}{\pgfqpoint{0.009821in}{-0.009821in}}%
\pgfpathcurveto{\pgfqpoint{0.012425in}{-0.007216in}}{\pgfqpoint{0.013889in}{-0.003683in}}{\pgfqpoint{0.013889in}{0.000000in}}%
\pgfpathcurveto{\pgfqpoint{0.013889in}{0.003683in}}{\pgfqpoint{0.012425in}{0.007216in}}{\pgfqpoint{0.009821in}{0.009821in}}%
\pgfpathcurveto{\pgfqpoint{0.007216in}{0.012425in}}{\pgfqpoint{0.003683in}{0.013889in}}{\pgfqpoint{0.000000in}{0.013889in}}%
\pgfpathcurveto{\pgfqpoint{-0.003683in}{0.013889in}}{\pgfqpoint{-0.007216in}{0.012425in}}{\pgfqpoint{-0.009821in}{0.009821in}}%
\pgfpathcurveto{\pgfqpoint{-0.012425in}{0.007216in}}{\pgfqpoint{-0.013889in}{0.003683in}}{\pgfqpoint{-0.013889in}{0.000000in}}%
\pgfpathcurveto{\pgfqpoint{-0.013889in}{-0.003683in}}{\pgfqpoint{-0.012425in}{-0.007216in}}{\pgfqpoint{-0.009821in}{-0.009821in}}%
\pgfpathcurveto{\pgfqpoint{-0.007216in}{-0.012425in}}{\pgfqpoint{-0.003683in}{-0.013889in}}{\pgfqpoint{0.000000in}{-0.013889in}}%
\pgfpathlineto{\pgfqpoint{0.000000in}{-0.013889in}}%
\pgfpathclose%
\pgfusepath{stroke,fill}%
}%
\begin{pgfscope}%
\pgfsys@transformshift{4.251269in}{4.426316in}%
\pgfsys@useobject{currentmarker}{}%
\end{pgfscope}%
\end{pgfscope}%
\begin{pgfscope}%
\pgfpathrectangle{\pgfqpoint{0.100000in}{2.493063in}}{\pgfqpoint{5.037500in}{3.427208in}}%
\pgfusepath{clip}%
\pgfsetrectcap%
\pgfsetroundjoin%
\pgfsetlinewidth{1.505625pt}%
\definecolor{currentstroke}{rgb}{0.000000,0.000000,1.000000}%
\pgfsetstrokecolor{currentstroke}%
\pgfsetstrokeopacity{0.500000}%
\pgfsetdash{}{0pt}%
\pgfpathmoveto{\pgfqpoint{4.244154in}{4.302410in}}%
\pgfusepath{stroke}%
\end{pgfscope}%
\begin{pgfscope}%
\pgfpathrectangle{\pgfqpoint{0.100000in}{2.493063in}}{\pgfqpoint{5.037500in}{3.427208in}}%
\pgfusepath{clip}%
\pgfsetbuttcap%
\pgfsetroundjoin%
\definecolor{currentfill}{rgb}{0.000000,0.000000,1.000000}%
\pgfsetfillcolor{currentfill}%
\pgfsetfillopacity{0.500000}%
\pgfsetlinewidth{0.250937pt}%
\definecolor{currentstroke}{rgb}{0.000000,0.000000,0.000000}%
\pgfsetstrokecolor{currentstroke}%
\pgfsetstrokeopacity{0.500000}%
\pgfsetdash{}{0pt}%
\pgfsys@defobject{currentmarker}{\pgfqpoint{-0.013889in}{-0.013889in}}{\pgfqpoint{0.013889in}{0.013889in}}{%
\pgfpathmoveto{\pgfqpoint{0.000000in}{-0.013889in}}%
\pgfpathcurveto{\pgfqpoint{0.003683in}{-0.013889in}}{\pgfqpoint{0.007216in}{-0.012425in}}{\pgfqpoint{0.009821in}{-0.009821in}}%
\pgfpathcurveto{\pgfqpoint{0.012425in}{-0.007216in}}{\pgfqpoint{0.013889in}{-0.003683in}}{\pgfqpoint{0.013889in}{0.000000in}}%
\pgfpathcurveto{\pgfqpoint{0.013889in}{0.003683in}}{\pgfqpoint{0.012425in}{0.007216in}}{\pgfqpoint{0.009821in}{0.009821in}}%
\pgfpathcurveto{\pgfqpoint{0.007216in}{0.012425in}}{\pgfqpoint{0.003683in}{0.013889in}}{\pgfqpoint{0.000000in}{0.013889in}}%
\pgfpathcurveto{\pgfqpoint{-0.003683in}{0.013889in}}{\pgfqpoint{-0.007216in}{0.012425in}}{\pgfqpoint{-0.009821in}{0.009821in}}%
\pgfpathcurveto{\pgfqpoint{-0.012425in}{0.007216in}}{\pgfqpoint{-0.013889in}{0.003683in}}{\pgfqpoint{-0.013889in}{0.000000in}}%
\pgfpathcurveto{\pgfqpoint{-0.013889in}{-0.003683in}}{\pgfqpoint{-0.012425in}{-0.007216in}}{\pgfqpoint{-0.009821in}{-0.009821in}}%
\pgfpathcurveto{\pgfqpoint{-0.007216in}{-0.012425in}}{\pgfqpoint{-0.003683in}{-0.013889in}}{\pgfqpoint{0.000000in}{-0.013889in}}%
\pgfpathlineto{\pgfqpoint{0.000000in}{-0.013889in}}%
\pgfpathclose%
\pgfusepath{stroke,fill}%
}%
\begin{pgfscope}%
\pgfsys@transformshift{4.244154in}{4.302410in}%
\pgfsys@useobject{currentmarker}{}%
\end{pgfscope}%
\end{pgfscope}%
\begin{pgfscope}%
\pgfpathrectangle{\pgfqpoint{0.100000in}{2.493063in}}{\pgfqpoint{5.037500in}{3.427208in}}%
\pgfusepath{clip}%
\pgfsetrectcap%
\pgfsetroundjoin%
\pgfsetlinewidth{1.505625pt}%
\definecolor{currentstroke}{rgb}{0.000000,0.000000,1.000000}%
\pgfsetstrokecolor{currentstroke}%
\pgfsetstrokeopacity{0.500000}%
\pgfsetdash{}{0pt}%
\pgfpathmoveto{\pgfqpoint{4.398579in}{4.347220in}}%
\pgfusepath{stroke}%
\end{pgfscope}%
\begin{pgfscope}%
\pgfpathrectangle{\pgfqpoint{0.100000in}{2.493063in}}{\pgfqpoint{5.037500in}{3.427208in}}%
\pgfusepath{clip}%
\pgfsetbuttcap%
\pgfsetroundjoin%
\definecolor{currentfill}{rgb}{0.000000,0.000000,1.000000}%
\pgfsetfillcolor{currentfill}%
\pgfsetfillopacity{0.500000}%
\pgfsetlinewidth{0.250937pt}%
\definecolor{currentstroke}{rgb}{0.000000,0.000000,0.000000}%
\pgfsetstrokecolor{currentstroke}%
\pgfsetstrokeopacity{0.500000}%
\pgfsetdash{}{0pt}%
\pgfsys@defobject{currentmarker}{\pgfqpoint{-0.011111in}{-0.011111in}}{\pgfqpoint{0.011111in}{0.011111in}}{%
\pgfpathmoveto{\pgfqpoint{0.000000in}{-0.011111in}}%
\pgfpathcurveto{\pgfqpoint{0.002947in}{-0.011111in}}{\pgfqpoint{0.005773in}{-0.009940in}}{\pgfqpoint{0.007857in}{-0.007857in}}%
\pgfpathcurveto{\pgfqpoint{0.009940in}{-0.005773in}}{\pgfqpoint{0.011111in}{-0.002947in}}{\pgfqpoint{0.011111in}{0.000000in}}%
\pgfpathcurveto{\pgfqpoint{0.011111in}{0.002947in}}{\pgfqpoint{0.009940in}{0.005773in}}{\pgfqpoint{0.007857in}{0.007857in}}%
\pgfpathcurveto{\pgfqpoint{0.005773in}{0.009940in}}{\pgfqpoint{0.002947in}{0.011111in}}{\pgfqpoint{0.000000in}{0.011111in}}%
\pgfpathcurveto{\pgfqpoint{-0.002947in}{0.011111in}}{\pgfqpoint{-0.005773in}{0.009940in}}{\pgfqpoint{-0.007857in}{0.007857in}}%
\pgfpathcurveto{\pgfqpoint{-0.009940in}{0.005773in}}{\pgfqpoint{-0.011111in}{0.002947in}}{\pgfqpoint{-0.011111in}{0.000000in}}%
\pgfpathcurveto{\pgfqpoint{-0.011111in}{-0.002947in}}{\pgfqpoint{-0.009940in}{-0.005773in}}{\pgfqpoint{-0.007857in}{-0.007857in}}%
\pgfpathcurveto{\pgfqpoint{-0.005773in}{-0.009940in}}{\pgfqpoint{-0.002947in}{-0.011111in}}{\pgfqpoint{0.000000in}{-0.011111in}}%
\pgfpathlineto{\pgfqpoint{0.000000in}{-0.011111in}}%
\pgfpathclose%
\pgfusepath{stroke,fill}%
}%
\begin{pgfscope}%
\pgfsys@transformshift{4.398579in}{4.347220in}%
\pgfsys@useobject{currentmarker}{}%
\end{pgfscope}%
\end{pgfscope}%
\begin{pgfscope}%
\pgfpathrectangle{\pgfqpoint{0.100000in}{2.493063in}}{\pgfqpoint{5.037500in}{3.427208in}}%
\pgfusepath{clip}%
\pgfsetrectcap%
\pgfsetroundjoin%
\pgfsetlinewidth{1.505625pt}%
\definecolor{currentstroke}{rgb}{0.000000,0.000000,1.000000}%
\pgfsetstrokecolor{currentstroke}%
\pgfsetstrokeopacity{0.500000}%
\pgfsetdash{}{0pt}%
\pgfpathmoveto{\pgfqpoint{4.178350in}{4.275940in}}%
\pgfusepath{stroke}%
\end{pgfscope}%
\begin{pgfscope}%
\pgfpathrectangle{\pgfqpoint{0.100000in}{2.493063in}}{\pgfqpoint{5.037500in}{3.427208in}}%
\pgfusepath{clip}%
\pgfsetbuttcap%
\pgfsetroundjoin%
\definecolor{currentfill}{rgb}{0.000000,0.000000,1.000000}%
\pgfsetfillcolor{currentfill}%
\pgfsetfillopacity{0.500000}%
\pgfsetlinewidth{0.250937pt}%
\definecolor{currentstroke}{rgb}{0.000000,0.000000,0.000000}%
\pgfsetstrokecolor{currentstroke}%
\pgfsetstrokeopacity{0.500000}%
\pgfsetdash{}{0pt}%
\pgfsys@defobject{currentmarker}{\pgfqpoint{-0.008333in}{-0.008333in}}{\pgfqpoint{0.008333in}{0.008333in}}{%
\pgfpathmoveto{\pgfqpoint{0.000000in}{-0.008333in}}%
\pgfpathcurveto{\pgfqpoint{0.002210in}{-0.008333in}}{\pgfqpoint{0.004330in}{-0.007455in}}{\pgfqpoint{0.005893in}{-0.005893in}}%
\pgfpathcurveto{\pgfqpoint{0.007455in}{-0.004330in}}{\pgfqpoint{0.008333in}{-0.002210in}}{\pgfqpoint{0.008333in}{0.000000in}}%
\pgfpathcurveto{\pgfqpoint{0.008333in}{0.002210in}}{\pgfqpoint{0.007455in}{0.004330in}}{\pgfqpoint{0.005893in}{0.005893in}}%
\pgfpathcurveto{\pgfqpoint{0.004330in}{0.007455in}}{\pgfqpoint{0.002210in}{0.008333in}}{\pgfqpoint{0.000000in}{0.008333in}}%
\pgfpathcurveto{\pgfqpoint{-0.002210in}{0.008333in}}{\pgfqpoint{-0.004330in}{0.007455in}}{\pgfqpoint{-0.005893in}{0.005893in}}%
\pgfpathcurveto{\pgfqpoint{-0.007455in}{0.004330in}}{\pgfqpoint{-0.008333in}{0.002210in}}{\pgfqpoint{-0.008333in}{0.000000in}}%
\pgfpathcurveto{\pgfqpoint{-0.008333in}{-0.002210in}}{\pgfqpoint{-0.007455in}{-0.004330in}}{\pgfqpoint{-0.005893in}{-0.005893in}}%
\pgfpathcurveto{\pgfqpoint{-0.004330in}{-0.007455in}}{\pgfqpoint{-0.002210in}{-0.008333in}}{\pgfqpoint{0.000000in}{-0.008333in}}%
\pgfpathlineto{\pgfqpoint{0.000000in}{-0.008333in}}%
\pgfpathclose%
\pgfusepath{stroke,fill}%
}%
\begin{pgfscope}%
\pgfsys@transformshift{4.178350in}{4.275940in}%
\pgfsys@useobject{currentmarker}{}%
\end{pgfscope}%
\end{pgfscope}%
\begin{pgfscope}%
\pgfpathrectangle{\pgfqpoint{0.100000in}{2.493063in}}{\pgfqpoint{5.037500in}{3.427208in}}%
\pgfusepath{clip}%
\pgfsetrectcap%
\pgfsetroundjoin%
\pgfsetlinewidth{1.505625pt}%
\definecolor{currentstroke}{rgb}{0.000000,0.000000,1.000000}%
\pgfsetstrokecolor{currentstroke}%
\pgfsetstrokeopacity{0.500000}%
\pgfsetdash{}{0pt}%
\pgfpathmoveto{\pgfqpoint{4.239205in}{4.389104in}}%
\pgfusepath{stroke}%
\end{pgfscope}%
\begin{pgfscope}%
\pgfpathrectangle{\pgfqpoint{0.100000in}{2.493063in}}{\pgfqpoint{5.037500in}{3.427208in}}%
\pgfusepath{clip}%
\pgfsetbuttcap%
\pgfsetroundjoin%
\definecolor{currentfill}{rgb}{0.000000,0.000000,1.000000}%
\pgfsetfillcolor{currentfill}%
\pgfsetfillopacity{0.500000}%
\pgfsetlinewidth{0.250937pt}%
\definecolor{currentstroke}{rgb}{0.000000,0.000000,0.000000}%
\pgfsetstrokecolor{currentstroke}%
\pgfsetstrokeopacity{0.500000}%
\pgfsetdash{}{0pt}%
\pgfsys@defobject{currentmarker}{\pgfqpoint{-0.008333in}{-0.008333in}}{\pgfqpoint{0.008333in}{0.008333in}}{%
\pgfpathmoveto{\pgfqpoint{0.000000in}{-0.008333in}}%
\pgfpathcurveto{\pgfqpoint{0.002210in}{-0.008333in}}{\pgfqpoint{0.004330in}{-0.007455in}}{\pgfqpoint{0.005893in}{-0.005893in}}%
\pgfpathcurveto{\pgfqpoint{0.007455in}{-0.004330in}}{\pgfqpoint{0.008333in}{-0.002210in}}{\pgfqpoint{0.008333in}{0.000000in}}%
\pgfpathcurveto{\pgfqpoint{0.008333in}{0.002210in}}{\pgfqpoint{0.007455in}{0.004330in}}{\pgfqpoint{0.005893in}{0.005893in}}%
\pgfpathcurveto{\pgfqpoint{0.004330in}{0.007455in}}{\pgfqpoint{0.002210in}{0.008333in}}{\pgfqpoint{0.000000in}{0.008333in}}%
\pgfpathcurveto{\pgfqpoint{-0.002210in}{0.008333in}}{\pgfqpoint{-0.004330in}{0.007455in}}{\pgfqpoint{-0.005893in}{0.005893in}}%
\pgfpathcurveto{\pgfqpoint{-0.007455in}{0.004330in}}{\pgfqpoint{-0.008333in}{0.002210in}}{\pgfqpoint{-0.008333in}{0.000000in}}%
\pgfpathcurveto{\pgfqpoint{-0.008333in}{-0.002210in}}{\pgfqpoint{-0.007455in}{-0.004330in}}{\pgfqpoint{-0.005893in}{-0.005893in}}%
\pgfpathcurveto{\pgfqpoint{-0.004330in}{-0.007455in}}{\pgfqpoint{-0.002210in}{-0.008333in}}{\pgfqpoint{0.000000in}{-0.008333in}}%
\pgfpathlineto{\pgfqpoint{0.000000in}{-0.008333in}}%
\pgfpathclose%
\pgfusepath{stroke,fill}%
}%
\begin{pgfscope}%
\pgfsys@transformshift{4.239205in}{4.389104in}%
\pgfsys@useobject{currentmarker}{}%
\end{pgfscope}%
\end{pgfscope}%
\begin{pgfscope}%
\pgfpathrectangle{\pgfqpoint{0.100000in}{2.493063in}}{\pgfqpoint{5.037500in}{3.427208in}}%
\pgfusepath{clip}%
\pgfsetrectcap%
\pgfsetroundjoin%
\pgfsetlinewidth{1.505625pt}%
\definecolor{currentstroke}{rgb}{0.000000,0.000000,1.000000}%
\pgfsetstrokecolor{currentstroke}%
\pgfsetstrokeopacity{0.500000}%
\pgfsetdash{}{0pt}%
\pgfpathmoveto{\pgfqpoint{4.545369in}{4.296571in}}%
\pgfusepath{stroke}%
\end{pgfscope}%
\begin{pgfscope}%
\pgfpathrectangle{\pgfqpoint{0.100000in}{2.493063in}}{\pgfqpoint{5.037500in}{3.427208in}}%
\pgfusepath{clip}%
\pgfsetbuttcap%
\pgfsetroundjoin%
\definecolor{currentfill}{rgb}{0.000000,0.000000,1.000000}%
\pgfsetfillcolor{currentfill}%
\pgfsetfillopacity{0.500000}%
\pgfsetlinewidth{0.250937pt}%
\definecolor{currentstroke}{rgb}{0.000000,0.000000,0.000000}%
\pgfsetstrokecolor{currentstroke}%
\pgfsetstrokeopacity{0.500000}%
\pgfsetdash{}{0pt}%
\pgfsys@defobject{currentmarker}{\pgfqpoint{-0.011111in}{-0.011111in}}{\pgfqpoint{0.011111in}{0.011111in}}{%
\pgfpathmoveto{\pgfqpoint{0.000000in}{-0.011111in}}%
\pgfpathcurveto{\pgfqpoint{0.002947in}{-0.011111in}}{\pgfqpoint{0.005773in}{-0.009940in}}{\pgfqpoint{0.007857in}{-0.007857in}}%
\pgfpathcurveto{\pgfqpoint{0.009940in}{-0.005773in}}{\pgfqpoint{0.011111in}{-0.002947in}}{\pgfqpoint{0.011111in}{0.000000in}}%
\pgfpathcurveto{\pgfqpoint{0.011111in}{0.002947in}}{\pgfqpoint{0.009940in}{0.005773in}}{\pgfqpoint{0.007857in}{0.007857in}}%
\pgfpathcurveto{\pgfqpoint{0.005773in}{0.009940in}}{\pgfqpoint{0.002947in}{0.011111in}}{\pgfqpoint{0.000000in}{0.011111in}}%
\pgfpathcurveto{\pgfqpoint{-0.002947in}{0.011111in}}{\pgfqpoint{-0.005773in}{0.009940in}}{\pgfqpoint{-0.007857in}{0.007857in}}%
\pgfpathcurveto{\pgfqpoint{-0.009940in}{0.005773in}}{\pgfqpoint{-0.011111in}{0.002947in}}{\pgfqpoint{-0.011111in}{0.000000in}}%
\pgfpathcurveto{\pgfqpoint{-0.011111in}{-0.002947in}}{\pgfqpoint{-0.009940in}{-0.005773in}}{\pgfqpoint{-0.007857in}{-0.007857in}}%
\pgfpathcurveto{\pgfqpoint{-0.005773in}{-0.009940in}}{\pgfqpoint{-0.002947in}{-0.011111in}}{\pgfqpoint{0.000000in}{-0.011111in}}%
\pgfpathlineto{\pgfqpoint{0.000000in}{-0.011111in}}%
\pgfpathclose%
\pgfusepath{stroke,fill}%
}%
\begin{pgfscope}%
\pgfsys@transformshift{4.545369in}{4.296571in}%
\pgfsys@useobject{currentmarker}{}%
\end{pgfscope}%
\end{pgfscope}%
\begin{pgfscope}%
\pgfpathrectangle{\pgfqpoint{0.100000in}{2.493063in}}{\pgfqpoint{5.037500in}{3.427208in}}%
\pgfusepath{clip}%
\pgfsetrectcap%
\pgfsetroundjoin%
\pgfsetlinewidth{1.505625pt}%
\definecolor{currentstroke}{rgb}{0.000000,0.000000,1.000000}%
\pgfsetstrokecolor{currentstroke}%
\pgfsetstrokeopacity{0.500000}%
\pgfsetdash{}{0pt}%
\pgfpathmoveto{\pgfqpoint{4.298077in}{4.521049in}}%
\pgfusepath{stroke}%
\end{pgfscope}%
\begin{pgfscope}%
\pgfpathrectangle{\pgfqpoint{0.100000in}{2.493063in}}{\pgfqpoint{5.037500in}{3.427208in}}%
\pgfusepath{clip}%
\pgfsetbuttcap%
\pgfsetroundjoin%
\definecolor{currentfill}{rgb}{0.000000,0.000000,1.000000}%
\pgfsetfillcolor{currentfill}%
\pgfsetfillopacity{0.500000}%
\pgfsetlinewidth{0.250937pt}%
\definecolor{currentstroke}{rgb}{0.000000,0.000000,0.000000}%
\pgfsetstrokecolor{currentstroke}%
\pgfsetstrokeopacity{0.500000}%
\pgfsetdash{}{0pt}%
\pgfsys@defobject{currentmarker}{\pgfqpoint{-0.005556in}{-0.005556in}}{\pgfqpoint{0.005556in}{0.005556in}}{%
\pgfpathmoveto{\pgfqpoint{0.000000in}{-0.005556in}}%
\pgfpathcurveto{\pgfqpoint{0.001473in}{-0.005556in}}{\pgfqpoint{0.002887in}{-0.004970in}}{\pgfqpoint{0.003928in}{-0.003928in}}%
\pgfpathcurveto{\pgfqpoint{0.004970in}{-0.002887in}}{\pgfqpoint{0.005556in}{-0.001473in}}{\pgfqpoint{0.005556in}{0.000000in}}%
\pgfpathcurveto{\pgfqpoint{0.005556in}{0.001473in}}{\pgfqpoint{0.004970in}{0.002887in}}{\pgfqpoint{0.003928in}{0.003928in}}%
\pgfpathcurveto{\pgfqpoint{0.002887in}{0.004970in}}{\pgfqpoint{0.001473in}{0.005556in}}{\pgfqpoint{0.000000in}{0.005556in}}%
\pgfpathcurveto{\pgfqpoint{-0.001473in}{0.005556in}}{\pgfqpoint{-0.002887in}{0.004970in}}{\pgfqpoint{-0.003928in}{0.003928in}}%
\pgfpathcurveto{\pgfqpoint{-0.004970in}{0.002887in}}{\pgfqpoint{-0.005556in}{0.001473in}}{\pgfqpoint{-0.005556in}{0.000000in}}%
\pgfpathcurveto{\pgfqpoint{-0.005556in}{-0.001473in}}{\pgfqpoint{-0.004970in}{-0.002887in}}{\pgfqpoint{-0.003928in}{-0.003928in}}%
\pgfpathcurveto{\pgfqpoint{-0.002887in}{-0.004970in}}{\pgfqpoint{-0.001473in}{-0.005556in}}{\pgfqpoint{0.000000in}{-0.005556in}}%
\pgfpathlineto{\pgfqpoint{0.000000in}{-0.005556in}}%
\pgfpathclose%
\pgfusepath{stroke,fill}%
}%
\begin{pgfscope}%
\pgfsys@transformshift{4.298077in}{4.521049in}%
\pgfsys@useobject{currentmarker}{}%
\end{pgfscope}%
\end{pgfscope}%
\begin{pgfscope}%
\pgfpathrectangle{\pgfqpoint{0.100000in}{2.493063in}}{\pgfqpoint{5.037500in}{3.427208in}}%
\pgfusepath{clip}%
\pgfsetrectcap%
\pgfsetroundjoin%
\pgfsetlinewidth{1.505625pt}%
\definecolor{currentstroke}{rgb}{0.678431,1.000000,0.184314}%
\pgfsetstrokecolor{currentstroke}%
\pgfsetstrokeopacity{0.500000}%
\pgfsetdash{}{0pt}%
\pgfpathmoveto{\pgfqpoint{0.718567in}{5.807400in}}%
\pgfusepath{stroke}%
\end{pgfscope}%
\begin{pgfscope}%
\pgfpathrectangle{\pgfqpoint{0.100000in}{2.493063in}}{\pgfqpoint{5.037500in}{3.427208in}}%
\pgfusepath{clip}%
\pgfsetbuttcap%
\pgfsetroundjoin%
\definecolor{currentfill}{rgb}{0.678431,1.000000,0.184314}%
\pgfsetfillcolor{currentfill}%
\pgfsetfillopacity{0.500000}%
\pgfsetlinewidth{0.250937pt}%
\definecolor{currentstroke}{rgb}{0.000000,0.000000,0.000000}%
\pgfsetstrokecolor{currentstroke}%
\pgfsetstrokeopacity{0.500000}%
\pgfsetdash{}{0pt}%
\pgfsys@defobject{currentmarker}{\pgfqpoint{-0.011111in}{-0.011111in}}{\pgfqpoint{0.011111in}{0.011111in}}{%
\pgfpathmoveto{\pgfqpoint{0.000000in}{-0.011111in}}%
\pgfpathcurveto{\pgfqpoint{0.002947in}{-0.011111in}}{\pgfqpoint{0.005773in}{-0.009940in}}{\pgfqpoint{0.007857in}{-0.007857in}}%
\pgfpathcurveto{\pgfqpoint{0.009940in}{-0.005773in}}{\pgfqpoint{0.011111in}{-0.002947in}}{\pgfqpoint{0.011111in}{0.000000in}}%
\pgfpathcurveto{\pgfqpoint{0.011111in}{0.002947in}}{\pgfqpoint{0.009940in}{0.005773in}}{\pgfqpoint{0.007857in}{0.007857in}}%
\pgfpathcurveto{\pgfqpoint{0.005773in}{0.009940in}}{\pgfqpoint{0.002947in}{0.011111in}}{\pgfqpoint{0.000000in}{0.011111in}}%
\pgfpathcurveto{\pgfqpoint{-0.002947in}{0.011111in}}{\pgfqpoint{-0.005773in}{0.009940in}}{\pgfqpoint{-0.007857in}{0.007857in}}%
\pgfpathcurveto{\pgfqpoint{-0.009940in}{0.005773in}}{\pgfqpoint{-0.011111in}{0.002947in}}{\pgfqpoint{-0.011111in}{0.000000in}}%
\pgfpathcurveto{\pgfqpoint{-0.011111in}{-0.002947in}}{\pgfqpoint{-0.009940in}{-0.005773in}}{\pgfqpoint{-0.007857in}{-0.007857in}}%
\pgfpathcurveto{\pgfqpoint{-0.005773in}{-0.009940in}}{\pgfqpoint{-0.002947in}{-0.011111in}}{\pgfqpoint{0.000000in}{-0.011111in}}%
\pgfpathlineto{\pgfqpoint{0.000000in}{-0.011111in}}%
\pgfpathclose%
\pgfusepath{stroke,fill}%
}%
\begin{pgfscope}%
\pgfsys@transformshift{0.718567in}{5.807400in}%
\pgfsys@useobject{currentmarker}{}%
\end{pgfscope}%
\end{pgfscope}%
\begin{pgfscope}%
\pgfpathrectangle{\pgfqpoint{0.100000in}{2.493063in}}{\pgfqpoint{5.037500in}{3.427208in}}%
\pgfusepath{clip}%
\pgfsetrectcap%
\pgfsetroundjoin%
\pgfsetlinewidth{1.505625pt}%
\definecolor{currentstroke}{rgb}{0.000000,0.000000,1.000000}%
\pgfsetstrokecolor{currentstroke}%
\pgfsetstrokeopacity{0.500000}%
\pgfsetdash{}{0pt}%
\pgfpathmoveto{\pgfqpoint{0.655467in}{5.677815in}}%
\pgfusepath{stroke}%
\end{pgfscope}%
\begin{pgfscope}%
\pgfpathrectangle{\pgfqpoint{0.100000in}{2.493063in}}{\pgfqpoint{5.037500in}{3.427208in}}%
\pgfusepath{clip}%
\pgfsetbuttcap%
\pgfsetroundjoin%
\definecolor{currentfill}{rgb}{0.000000,0.000000,1.000000}%
\pgfsetfillcolor{currentfill}%
\pgfsetfillopacity{0.500000}%
\pgfsetlinewidth{0.250937pt}%
\definecolor{currentstroke}{rgb}{0.000000,0.000000,0.000000}%
\pgfsetstrokecolor{currentstroke}%
\pgfsetstrokeopacity{0.500000}%
\pgfsetdash{}{0pt}%
\pgfsys@defobject{currentmarker}{\pgfqpoint{-0.005556in}{-0.005556in}}{\pgfqpoint{0.005556in}{0.005556in}}{%
\pgfpathmoveto{\pgfqpoint{0.000000in}{-0.005556in}}%
\pgfpathcurveto{\pgfqpoint{0.001473in}{-0.005556in}}{\pgfqpoint{0.002887in}{-0.004970in}}{\pgfqpoint{0.003928in}{-0.003928in}}%
\pgfpathcurveto{\pgfqpoint{0.004970in}{-0.002887in}}{\pgfqpoint{0.005556in}{-0.001473in}}{\pgfqpoint{0.005556in}{0.000000in}}%
\pgfpathcurveto{\pgfqpoint{0.005556in}{0.001473in}}{\pgfqpoint{0.004970in}{0.002887in}}{\pgfqpoint{0.003928in}{0.003928in}}%
\pgfpathcurveto{\pgfqpoint{0.002887in}{0.004970in}}{\pgfqpoint{0.001473in}{0.005556in}}{\pgfqpoint{0.000000in}{0.005556in}}%
\pgfpathcurveto{\pgfqpoint{-0.001473in}{0.005556in}}{\pgfqpoint{-0.002887in}{0.004970in}}{\pgfqpoint{-0.003928in}{0.003928in}}%
\pgfpathcurveto{\pgfqpoint{-0.004970in}{0.002887in}}{\pgfqpoint{-0.005556in}{0.001473in}}{\pgfqpoint{-0.005556in}{0.000000in}}%
\pgfpathcurveto{\pgfqpoint{-0.005556in}{-0.001473in}}{\pgfqpoint{-0.004970in}{-0.002887in}}{\pgfqpoint{-0.003928in}{-0.003928in}}%
\pgfpathcurveto{\pgfqpoint{-0.002887in}{-0.004970in}}{\pgfqpoint{-0.001473in}{-0.005556in}}{\pgfqpoint{0.000000in}{-0.005556in}}%
\pgfpathlineto{\pgfqpoint{0.000000in}{-0.005556in}}%
\pgfpathclose%
\pgfusepath{stroke,fill}%
}%
\begin{pgfscope}%
\pgfsys@transformshift{0.655467in}{5.677815in}%
\pgfsys@useobject{currentmarker}{}%
\end{pgfscope}%
\end{pgfscope}%
\begin{pgfscope}%
\pgfpathrectangle{\pgfqpoint{0.100000in}{2.493063in}}{\pgfqpoint{5.037500in}{3.427208in}}%
\pgfusepath{clip}%
\pgfsetrectcap%
\pgfsetroundjoin%
\pgfsetlinewidth{1.505625pt}%
\definecolor{currentstroke}{rgb}{0.678431,1.000000,0.184314}%
\pgfsetstrokecolor{currentstroke}%
\pgfsetstrokeopacity{0.500000}%
\pgfsetdash{}{0pt}%
\pgfpathmoveto{\pgfqpoint{0.888019in}{5.449454in}}%
\pgfusepath{stroke}%
\end{pgfscope}%
\begin{pgfscope}%
\pgfpathrectangle{\pgfqpoint{0.100000in}{2.493063in}}{\pgfqpoint{5.037500in}{3.427208in}}%
\pgfusepath{clip}%
\pgfsetbuttcap%
\pgfsetroundjoin%
\definecolor{currentfill}{rgb}{0.678431,1.000000,0.184314}%
\pgfsetfillcolor{currentfill}%
\pgfsetfillopacity{0.500000}%
\pgfsetlinewidth{0.250937pt}%
\definecolor{currentstroke}{rgb}{0.000000,0.000000,0.000000}%
\pgfsetstrokecolor{currentstroke}%
\pgfsetstrokeopacity{0.500000}%
\pgfsetdash{}{0pt}%
\pgfsys@defobject{currentmarker}{\pgfqpoint{-0.011111in}{-0.011111in}}{\pgfqpoint{0.011111in}{0.011111in}}{%
\pgfpathmoveto{\pgfqpoint{0.000000in}{-0.011111in}}%
\pgfpathcurveto{\pgfqpoint{0.002947in}{-0.011111in}}{\pgfqpoint{0.005773in}{-0.009940in}}{\pgfqpoint{0.007857in}{-0.007857in}}%
\pgfpathcurveto{\pgfqpoint{0.009940in}{-0.005773in}}{\pgfqpoint{0.011111in}{-0.002947in}}{\pgfqpoint{0.011111in}{0.000000in}}%
\pgfpathcurveto{\pgfqpoint{0.011111in}{0.002947in}}{\pgfqpoint{0.009940in}{0.005773in}}{\pgfqpoint{0.007857in}{0.007857in}}%
\pgfpathcurveto{\pgfqpoint{0.005773in}{0.009940in}}{\pgfqpoint{0.002947in}{0.011111in}}{\pgfqpoint{0.000000in}{0.011111in}}%
\pgfpathcurveto{\pgfqpoint{-0.002947in}{0.011111in}}{\pgfqpoint{-0.005773in}{0.009940in}}{\pgfqpoint{-0.007857in}{0.007857in}}%
\pgfpathcurveto{\pgfqpoint{-0.009940in}{0.005773in}}{\pgfqpoint{-0.011111in}{0.002947in}}{\pgfqpoint{-0.011111in}{0.000000in}}%
\pgfpathcurveto{\pgfqpoint{-0.011111in}{-0.002947in}}{\pgfqpoint{-0.009940in}{-0.005773in}}{\pgfqpoint{-0.007857in}{-0.007857in}}%
\pgfpathcurveto{\pgfqpoint{-0.005773in}{-0.009940in}}{\pgfqpoint{-0.002947in}{-0.011111in}}{\pgfqpoint{0.000000in}{-0.011111in}}%
\pgfpathlineto{\pgfqpoint{0.000000in}{-0.011111in}}%
\pgfpathclose%
\pgfusepath{stroke,fill}%
}%
\begin{pgfscope}%
\pgfsys@transformshift{0.888019in}{5.449454in}%
\pgfsys@useobject{currentmarker}{}%
\end{pgfscope}%
\end{pgfscope}%
\begin{pgfscope}%
\pgfpathrectangle{\pgfqpoint{0.100000in}{2.493063in}}{\pgfqpoint{5.037500in}{3.427208in}}%
\pgfusepath{clip}%
\pgfsetrectcap%
\pgfsetroundjoin%
\pgfsetlinewidth{1.505625pt}%
\definecolor{currentstroke}{rgb}{0.678431,1.000000,0.184314}%
\pgfsetstrokecolor{currentstroke}%
\pgfsetstrokeopacity{0.500000}%
\pgfsetdash{}{0pt}%
\pgfpathmoveto{\pgfqpoint{0.591430in}{5.528578in}}%
\pgfusepath{stroke}%
\end{pgfscope}%
\begin{pgfscope}%
\pgfpathrectangle{\pgfqpoint{0.100000in}{2.493063in}}{\pgfqpoint{5.037500in}{3.427208in}}%
\pgfusepath{clip}%
\pgfsetbuttcap%
\pgfsetroundjoin%
\definecolor{currentfill}{rgb}{0.678431,1.000000,0.184314}%
\pgfsetfillcolor{currentfill}%
\pgfsetfillopacity{0.500000}%
\pgfsetlinewidth{0.250937pt}%
\definecolor{currentstroke}{rgb}{0.000000,0.000000,0.000000}%
\pgfsetstrokecolor{currentstroke}%
\pgfsetstrokeopacity{0.500000}%
\pgfsetdash{}{0pt}%
\pgfsys@defobject{currentmarker}{\pgfqpoint{-0.019444in}{-0.019444in}}{\pgfqpoint{0.019444in}{0.019444in}}{%
\pgfpathmoveto{\pgfqpoint{0.000000in}{-0.019444in}}%
\pgfpathcurveto{\pgfqpoint{0.005157in}{-0.019444in}}{\pgfqpoint{0.010103in}{-0.017396in}}{\pgfqpoint{0.013749in}{-0.013749in}}%
\pgfpathcurveto{\pgfqpoint{0.017396in}{-0.010103in}}{\pgfqpoint{0.019444in}{-0.005157in}}{\pgfqpoint{0.019444in}{0.000000in}}%
\pgfpathcurveto{\pgfqpoint{0.019444in}{0.005157in}}{\pgfqpoint{0.017396in}{0.010103in}}{\pgfqpoint{0.013749in}{0.013749in}}%
\pgfpathcurveto{\pgfqpoint{0.010103in}{0.017396in}}{\pgfqpoint{0.005157in}{0.019444in}}{\pgfqpoint{0.000000in}{0.019444in}}%
\pgfpathcurveto{\pgfqpoint{-0.005157in}{0.019444in}}{\pgfqpoint{-0.010103in}{0.017396in}}{\pgfqpoint{-0.013749in}{0.013749in}}%
\pgfpathcurveto{\pgfqpoint{-0.017396in}{0.010103in}}{\pgfqpoint{-0.019444in}{0.005157in}}{\pgfqpoint{-0.019444in}{0.000000in}}%
\pgfpathcurveto{\pgfqpoint{-0.019444in}{-0.005157in}}{\pgfqpoint{-0.017396in}{-0.010103in}}{\pgfqpoint{-0.013749in}{-0.013749in}}%
\pgfpathcurveto{\pgfqpoint{-0.010103in}{-0.017396in}}{\pgfqpoint{-0.005157in}{-0.019444in}}{\pgfqpoint{0.000000in}{-0.019444in}}%
\pgfpathlineto{\pgfqpoint{0.000000in}{-0.019444in}}%
\pgfpathclose%
\pgfusepath{stroke,fill}%
}%
\begin{pgfscope}%
\pgfsys@transformshift{0.591430in}{5.528578in}%
\pgfsys@useobject{currentmarker}{}%
\end{pgfscope}%
\end{pgfscope}%
\begin{pgfscope}%
\pgfpathrectangle{\pgfqpoint{0.100000in}{2.493063in}}{\pgfqpoint{5.037500in}{3.427208in}}%
\pgfusepath{clip}%
\pgfsetrectcap%
\pgfsetroundjoin%
\pgfsetlinewidth{1.505625pt}%
\definecolor{currentstroke}{rgb}{0.678431,1.000000,0.184314}%
\pgfsetstrokecolor{currentstroke}%
\pgfsetstrokeopacity{0.500000}%
\pgfsetdash{}{0pt}%
\pgfpathmoveto{\pgfqpoint{0.718215in}{5.766601in}}%
\pgfusepath{stroke}%
\end{pgfscope}%
\begin{pgfscope}%
\pgfpathrectangle{\pgfqpoint{0.100000in}{2.493063in}}{\pgfqpoint{5.037500in}{3.427208in}}%
\pgfusepath{clip}%
\pgfsetbuttcap%
\pgfsetroundjoin%
\definecolor{currentfill}{rgb}{0.678431,1.000000,0.184314}%
\pgfsetfillcolor{currentfill}%
\pgfsetfillopacity{0.500000}%
\pgfsetlinewidth{0.250937pt}%
\definecolor{currentstroke}{rgb}{0.000000,0.000000,0.000000}%
\pgfsetstrokecolor{currentstroke}%
\pgfsetstrokeopacity{0.500000}%
\pgfsetdash{}{0pt}%
\pgfsys@defobject{currentmarker}{\pgfqpoint{-0.008333in}{-0.008333in}}{\pgfqpoint{0.008333in}{0.008333in}}{%
\pgfpathmoveto{\pgfqpoint{0.000000in}{-0.008333in}}%
\pgfpathcurveto{\pgfqpoint{0.002210in}{-0.008333in}}{\pgfqpoint{0.004330in}{-0.007455in}}{\pgfqpoint{0.005893in}{-0.005893in}}%
\pgfpathcurveto{\pgfqpoint{0.007455in}{-0.004330in}}{\pgfqpoint{0.008333in}{-0.002210in}}{\pgfqpoint{0.008333in}{0.000000in}}%
\pgfpathcurveto{\pgfqpoint{0.008333in}{0.002210in}}{\pgfqpoint{0.007455in}{0.004330in}}{\pgfqpoint{0.005893in}{0.005893in}}%
\pgfpathcurveto{\pgfqpoint{0.004330in}{0.007455in}}{\pgfqpoint{0.002210in}{0.008333in}}{\pgfqpoint{0.000000in}{0.008333in}}%
\pgfpathcurveto{\pgfqpoint{-0.002210in}{0.008333in}}{\pgfqpoint{-0.004330in}{0.007455in}}{\pgfqpoint{-0.005893in}{0.005893in}}%
\pgfpathcurveto{\pgfqpoint{-0.007455in}{0.004330in}}{\pgfqpoint{-0.008333in}{0.002210in}}{\pgfqpoint{-0.008333in}{0.000000in}}%
\pgfpathcurveto{\pgfqpoint{-0.008333in}{-0.002210in}}{\pgfqpoint{-0.007455in}{-0.004330in}}{\pgfqpoint{-0.005893in}{-0.005893in}}%
\pgfpathcurveto{\pgfqpoint{-0.004330in}{-0.007455in}}{\pgfqpoint{-0.002210in}{-0.008333in}}{\pgfqpoint{0.000000in}{-0.008333in}}%
\pgfpathlineto{\pgfqpoint{0.000000in}{-0.008333in}}%
\pgfpathclose%
\pgfusepath{stroke,fill}%
}%
\begin{pgfscope}%
\pgfsys@transformshift{0.718215in}{5.766601in}%
\pgfsys@useobject{currentmarker}{}%
\end{pgfscope}%
\end{pgfscope}%
\begin{pgfscope}%
\pgfpathrectangle{\pgfqpoint{0.100000in}{2.493063in}}{\pgfqpoint{5.037500in}{3.427208in}}%
\pgfusepath{clip}%
\pgfsetrectcap%
\pgfsetroundjoin%
\pgfsetlinewidth{1.505625pt}%
\definecolor{currentstroke}{rgb}{0.678431,1.000000,0.184314}%
\pgfsetstrokecolor{currentstroke}%
\pgfsetstrokeopacity{0.500000}%
\pgfsetdash{}{0pt}%
\pgfpathmoveto{\pgfqpoint{0.627763in}{5.627031in}}%
\pgfusepath{stroke}%
\end{pgfscope}%
\begin{pgfscope}%
\pgfpathrectangle{\pgfqpoint{0.100000in}{2.493063in}}{\pgfqpoint{5.037500in}{3.427208in}}%
\pgfusepath{clip}%
\pgfsetbuttcap%
\pgfsetroundjoin%
\definecolor{currentfill}{rgb}{0.678431,1.000000,0.184314}%
\pgfsetfillcolor{currentfill}%
\pgfsetfillopacity{0.500000}%
\pgfsetlinewidth{0.250937pt}%
\definecolor{currentstroke}{rgb}{0.000000,0.000000,0.000000}%
\pgfsetstrokecolor{currentstroke}%
\pgfsetstrokeopacity{0.500000}%
\pgfsetdash{}{0pt}%
\pgfsys@defobject{currentmarker}{\pgfqpoint{-0.011111in}{-0.011111in}}{\pgfqpoint{0.011111in}{0.011111in}}{%
\pgfpathmoveto{\pgfqpoint{0.000000in}{-0.011111in}}%
\pgfpathcurveto{\pgfqpoint{0.002947in}{-0.011111in}}{\pgfqpoint{0.005773in}{-0.009940in}}{\pgfqpoint{0.007857in}{-0.007857in}}%
\pgfpathcurveto{\pgfqpoint{0.009940in}{-0.005773in}}{\pgfqpoint{0.011111in}{-0.002947in}}{\pgfqpoint{0.011111in}{0.000000in}}%
\pgfpathcurveto{\pgfqpoint{0.011111in}{0.002947in}}{\pgfqpoint{0.009940in}{0.005773in}}{\pgfqpoint{0.007857in}{0.007857in}}%
\pgfpathcurveto{\pgfqpoint{0.005773in}{0.009940in}}{\pgfqpoint{0.002947in}{0.011111in}}{\pgfqpoint{0.000000in}{0.011111in}}%
\pgfpathcurveto{\pgfqpoint{-0.002947in}{0.011111in}}{\pgfqpoint{-0.005773in}{0.009940in}}{\pgfqpoint{-0.007857in}{0.007857in}}%
\pgfpathcurveto{\pgfqpoint{-0.009940in}{0.005773in}}{\pgfqpoint{-0.011111in}{0.002947in}}{\pgfqpoint{-0.011111in}{0.000000in}}%
\pgfpathcurveto{\pgfqpoint{-0.011111in}{-0.002947in}}{\pgfqpoint{-0.009940in}{-0.005773in}}{\pgfqpoint{-0.007857in}{-0.007857in}}%
\pgfpathcurveto{\pgfqpoint{-0.005773in}{-0.009940in}}{\pgfqpoint{-0.002947in}{-0.011111in}}{\pgfqpoint{0.000000in}{-0.011111in}}%
\pgfpathlineto{\pgfqpoint{0.000000in}{-0.011111in}}%
\pgfpathclose%
\pgfusepath{stroke,fill}%
}%
\begin{pgfscope}%
\pgfsys@transformshift{0.627763in}{5.627031in}%
\pgfsys@useobject{currentmarker}{}%
\end{pgfscope}%
\end{pgfscope}%
\begin{pgfscope}%
\pgfpathrectangle{\pgfqpoint{0.100000in}{2.493063in}}{\pgfqpoint{5.037500in}{3.427208in}}%
\pgfusepath{clip}%
\pgfsetrectcap%
\pgfsetroundjoin%
\pgfsetlinewidth{1.505625pt}%
\definecolor{currentstroke}{rgb}{0.000000,0.000000,1.000000}%
\pgfsetstrokecolor{currentstroke}%
\pgfsetstrokeopacity{0.500000}%
\pgfsetdash{}{0pt}%
\pgfpathmoveto{\pgfqpoint{0.689747in}{5.675764in}}%
\pgfusepath{stroke}%
\end{pgfscope}%
\begin{pgfscope}%
\pgfpathrectangle{\pgfqpoint{0.100000in}{2.493063in}}{\pgfqpoint{5.037500in}{3.427208in}}%
\pgfusepath{clip}%
\pgfsetbuttcap%
\pgfsetroundjoin%
\definecolor{currentfill}{rgb}{0.000000,0.000000,1.000000}%
\pgfsetfillcolor{currentfill}%
\pgfsetfillopacity{0.500000}%
\pgfsetlinewidth{0.250937pt}%
\definecolor{currentstroke}{rgb}{0.000000,0.000000,0.000000}%
\pgfsetstrokecolor{currentstroke}%
\pgfsetstrokeopacity{0.500000}%
\pgfsetdash{}{0pt}%
\pgfsys@defobject{currentmarker}{\pgfqpoint{-0.038889in}{-0.038889in}}{\pgfqpoint{0.038889in}{0.038889in}}{%
\pgfpathmoveto{\pgfqpoint{0.000000in}{-0.038889in}}%
\pgfpathcurveto{\pgfqpoint{0.010313in}{-0.038889in}}{\pgfqpoint{0.020206in}{-0.034791in}}{\pgfqpoint{0.027499in}{-0.027499in}}%
\pgfpathcurveto{\pgfqpoint{0.034791in}{-0.020206in}}{\pgfqpoint{0.038889in}{-0.010313in}}{\pgfqpoint{0.038889in}{0.000000in}}%
\pgfpathcurveto{\pgfqpoint{0.038889in}{0.010313in}}{\pgfqpoint{0.034791in}{0.020206in}}{\pgfqpoint{0.027499in}{0.027499in}}%
\pgfpathcurveto{\pgfqpoint{0.020206in}{0.034791in}}{\pgfqpoint{0.010313in}{0.038889in}}{\pgfqpoint{0.000000in}{0.038889in}}%
\pgfpathcurveto{\pgfqpoint{-0.010313in}{0.038889in}}{\pgfqpoint{-0.020206in}{0.034791in}}{\pgfqpoint{-0.027499in}{0.027499in}}%
\pgfpathcurveto{\pgfqpoint{-0.034791in}{0.020206in}}{\pgfqpoint{-0.038889in}{0.010313in}}{\pgfqpoint{-0.038889in}{0.000000in}}%
\pgfpathcurveto{\pgfqpoint{-0.038889in}{-0.010313in}}{\pgfqpoint{-0.034791in}{-0.020206in}}{\pgfqpoint{-0.027499in}{-0.027499in}}%
\pgfpathcurveto{\pgfqpoint{-0.020206in}{-0.034791in}}{\pgfqpoint{-0.010313in}{-0.038889in}}{\pgfqpoint{0.000000in}{-0.038889in}}%
\pgfpathlineto{\pgfqpoint{0.000000in}{-0.038889in}}%
\pgfpathclose%
\pgfusepath{stroke,fill}%
}%
\begin{pgfscope}%
\pgfsys@transformshift{0.689747in}{5.675764in}%
\pgfsys@useobject{currentmarker}{}%
\end{pgfscope}%
\end{pgfscope}%
\begin{pgfscope}%
\pgfpathrectangle{\pgfqpoint{0.100000in}{2.493063in}}{\pgfqpoint{5.037500in}{3.427208in}}%
\pgfusepath{clip}%
\pgfsetrectcap%
\pgfsetroundjoin%
\pgfsetlinewidth{1.505625pt}%
\definecolor{currentstroke}{rgb}{0.678431,1.000000,0.184314}%
\pgfsetstrokecolor{currentstroke}%
\pgfsetstrokeopacity{0.500000}%
\pgfsetdash{}{0pt}%
\pgfpathmoveto{\pgfqpoint{1.062204in}{5.577835in}}%
\pgfusepath{stroke}%
\end{pgfscope}%
\begin{pgfscope}%
\pgfpathrectangle{\pgfqpoint{0.100000in}{2.493063in}}{\pgfqpoint{5.037500in}{3.427208in}}%
\pgfusepath{clip}%
\pgfsetbuttcap%
\pgfsetroundjoin%
\definecolor{currentfill}{rgb}{0.678431,1.000000,0.184314}%
\pgfsetfillcolor{currentfill}%
\pgfsetfillopacity{0.500000}%
\pgfsetlinewidth{0.250937pt}%
\definecolor{currentstroke}{rgb}{0.000000,0.000000,0.000000}%
\pgfsetstrokecolor{currentstroke}%
\pgfsetstrokeopacity{0.500000}%
\pgfsetdash{}{0pt}%
\pgfsys@defobject{currentmarker}{\pgfqpoint{-0.016667in}{-0.016667in}}{\pgfqpoint{0.016667in}{0.016667in}}{%
\pgfpathmoveto{\pgfqpoint{0.000000in}{-0.016667in}}%
\pgfpathcurveto{\pgfqpoint{0.004420in}{-0.016667in}}{\pgfqpoint{0.008660in}{-0.014911in}}{\pgfqpoint{0.011785in}{-0.011785in}}%
\pgfpathcurveto{\pgfqpoint{0.014911in}{-0.008660in}}{\pgfqpoint{0.016667in}{-0.004420in}}{\pgfqpoint{0.016667in}{0.000000in}}%
\pgfpathcurveto{\pgfqpoint{0.016667in}{0.004420in}}{\pgfqpoint{0.014911in}{0.008660in}}{\pgfqpoint{0.011785in}{0.011785in}}%
\pgfpathcurveto{\pgfqpoint{0.008660in}{0.014911in}}{\pgfqpoint{0.004420in}{0.016667in}}{\pgfqpoint{0.000000in}{0.016667in}}%
\pgfpathcurveto{\pgfqpoint{-0.004420in}{0.016667in}}{\pgfqpoint{-0.008660in}{0.014911in}}{\pgfqpoint{-0.011785in}{0.011785in}}%
\pgfpathcurveto{\pgfqpoint{-0.014911in}{0.008660in}}{\pgfqpoint{-0.016667in}{0.004420in}}{\pgfqpoint{-0.016667in}{0.000000in}}%
\pgfpathcurveto{\pgfqpoint{-0.016667in}{-0.004420in}}{\pgfqpoint{-0.014911in}{-0.008660in}}{\pgfqpoint{-0.011785in}{-0.011785in}}%
\pgfpathcurveto{\pgfqpoint{-0.008660in}{-0.014911in}}{\pgfqpoint{-0.004420in}{-0.016667in}}{\pgfqpoint{0.000000in}{-0.016667in}}%
\pgfpathlineto{\pgfqpoint{0.000000in}{-0.016667in}}%
\pgfpathclose%
\pgfusepath{stroke,fill}%
}%
\begin{pgfscope}%
\pgfsys@transformshift{1.062204in}{5.577835in}%
\pgfsys@useobject{currentmarker}{}%
\end{pgfscope}%
\end{pgfscope}%
\begin{pgfscope}%
\pgfpathrectangle{\pgfqpoint{0.100000in}{2.493063in}}{\pgfqpoint{5.037500in}{3.427208in}}%
\pgfusepath{clip}%
\pgfsetrectcap%
\pgfsetroundjoin%
\pgfsetlinewidth{1.505625pt}%
\definecolor{currentstroke}{rgb}{0.678431,1.000000,0.184314}%
\pgfsetstrokecolor{currentstroke}%
\pgfsetstrokeopacity{0.500000}%
\pgfsetdash{}{0pt}%
\pgfpathmoveto{\pgfqpoint{0.946329in}{5.416301in}}%
\pgfusepath{stroke}%
\end{pgfscope}%
\begin{pgfscope}%
\pgfpathrectangle{\pgfqpoint{0.100000in}{2.493063in}}{\pgfqpoint{5.037500in}{3.427208in}}%
\pgfusepath{clip}%
\pgfsetbuttcap%
\pgfsetroundjoin%
\definecolor{currentfill}{rgb}{0.678431,1.000000,0.184314}%
\pgfsetfillcolor{currentfill}%
\pgfsetfillopacity{0.500000}%
\pgfsetlinewidth{0.250937pt}%
\definecolor{currentstroke}{rgb}{0.000000,0.000000,0.000000}%
\pgfsetstrokecolor{currentstroke}%
\pgfsetstrokeopacity{0.500000}%
\pgfsetdash{}{0pt}%
\pgfsys@defobject{currentmarker}{\pgfqpoint{-0.008333in}{-0.008333in}}{\pgfqpoint{0.008333in}{0.008333in}}{%
\pgfpathmoveto{\pgfqpoint{0.000000in}{-0.008333in}}%
\pgfpathcurveto{\pgfqpoint{0.002210in}{-0.008333in}}{\pgfqpoint{0.004330in}{-0.007455in}}{\pgfqpoint{0.005893in}{-0.005893in}}%
\pgfpathcurveto{\pgfqpoint{0.007455in}{-0.004330in}}{\pgfqpoint{0.008333in}{-0.002210in}}{\pgfqpoint{0.008333in}{0.000000in}}%
\pgfpathcurveto{\pgfqpoint{0.008333in}{0.002210in}}{\pgfqpoint{0.007455in}{0.004330in}}{\pgfqpoint{0.005893in}{0.005893in}}%
\pgfpathcurveto{\pgfqpoint{0.004330in}{0.007455in}}{\pgfqpoint{0.002210in}{0.008333in}}{\pgfqpoint{0.000000in}{0.008333in}}%
\pgfpathcurveto{\pgfqpoint{-0.002210in}{0.008333in}}{\pgfqpoint{-0.004330in}{0.007455in}}{\pgfqpoint{-0.005893in}{0.005893in}}%
\pgfpathcurveto{\pgfqpoint{-0.007455in}{0.004330in}}{\pgfqpoint{-0.008333in}{0.002210in}}{\pgfqpoint{-0.008333in}{0.000000in}}%
\pgfpathcurveto{\pgfqpoint{-0.008333in}{-0.002210in}}{\pgfqpoint{-0.007455in}{-0.004330in}}{\pgfqpoint{-0.005893in}{-0.005893in}}%
\pgfpathcurveto{\pgfqpoint{-0.004330in}{-0.007455in}}{\pgfqpoint{-0.002210in}{-0.008333in}}{\pgfqpoint{0.000000in}{-0.008333in}}%
\pgfpathlineto{\pgfqpoint{0.000000in}{-0.008333in}}%
\pgfpathclose%
\pgfusepath{stroke,fill}%
}%
\begin{pgfscope}%
\pgfsys@transformshift{0.946329in}{5.416301in}%
\pgfsys@useobject{currentmarker}{}%
\end{pgfscope}%
\end{pgfscope}%
\begin{pgfscope}%
\pgfpathrectangle{\pgfqpoint{0.100000in}{2.493063in}}{\pgfqpoint{5.037500in}{3.427208in}}%
\pgfusepath{clip}%
\pgfsetrectcap%
\pgfsetroundjoin%
\pgfsetlinewidth{1.505625pt}%
\definecolor{currentstroke}{rgb}{0.678431,1.000000,0.184314}%
\pgfsetstrokecolor{currentstroke}%
\pgfsetstrokeopacity{0.500000}%
\pgfsetdash{}{0pt}%
\pgfpathmoveto{\pgfqpoint{0.835996in}{5.610263in}}%
\pgfusepath{stroke}%
\end{pgfscope}%
\begin{pgfscope}%
\pgfpathrectangle{\pgfqpoint{0.100000in}{2.493063in}}{\pgfqpoint{5.037500in}{3.427208in}}%
\pgfusepath{clip}%
\pgfsetbuttcap%
\pgfsetroundjoin%
\definecolor{currentfill}{rgb}{0.678431,1.000000,0.184314}%
\pgfsetfillcolor{currentfill}%
\pgfsetfillopacity{0.500000}%
\pgfsetlinewidth{0.250937pt}%
\definecolor{currentstroke}{rgb}{0.000000,0.000000,0.000000}%
\pgfsetstrokecolor{currentstroke}%
\pgfsetstrokeopacity{0.500000}%
\pgfsetdash{}{0pt}%
\pgfsys@defobject{currentmarker}{\pgfqpoint{-0.013889in}{-0.013889in}}{\pgfqpoint{0.013889in}{0.013889in}}{%
\pgfpathmoveto{\pgfqpoint{0.000000in}{-0.013889in}}%
\pgfpathcurveto{\pgfqpoint{0.003683in}{-0.013889in}}{\pgfqpoint{0.007216in}{-0.012425in}}{\pgfqpoint{0.009821in}{-0.009821in}}%
\pgfpathcurveto{\pgfqpoint{0.012425in}{-0.007216in}}{\pgfqpoint{0.013889in}{-0.003683in}}{\pgfqpoint{0.013889in}{0.000000in}}%
\pgfpathcurveto{\pgfqpoint{0.013889in}{0.003683in}}{\pgfqpoint{0.012425in}{0.007216in}}{\pgfqpoint{0.009821in}{0.009821in}}%
\pgfpathcurveto{\pgfqpoint{0.007216in}{0.012425in}}{\pgfqpoint{0.003683in}{0.013889in}}{\pgfqpoint{0.000000in}{0.013889in}}%
\pgfpathcurveto{\pgfqpoint{-0.003683in}{0.013889in}}{\pgfqpoint{-0.007216in}{0.012425in}}{\pgfqpoint{-0.009821in}{0.009821in}}%
\pgfpathcurveto{\pgfqpoint{-0.012425in}{0.007216in}}{\pgfqpoint{-0.013889in}{0.003683in}}{\pgfqpoint{-0.013889in}{0.000000in}}%
\pgfpathcurveto{\pgfqpoint{-0.013889in}{-0.003683in}}{\pgfqpoint{-0.012425in}{-0.007216in}}{\pgfqpoint{-0.009821in}{-0.009821in}}%
\pgfpathcurveto{\pgfqpoint{-0.007216in}{-0.012425in}}{\pgfqpoint{-0.003683in}{-0.013889in}}{\pgfqpoint{0.000000in}{-0.013889in}}%
\pgfpathlineto{\pgfqpoint{0.000000in}{-0.013889in}}%
\pgfpathclose%
\pgfusepath{stroke,fill}%
}%
\begin{pgfscope}%
\pgfsys@transformshift{0.835996in}{5.610263in}%
\pgfsys@useobject{currentmarker}{}%
\end{pgfscope}%
\end{pgfscope}%
\begin{pgfscope}%
\pgfpathrectangle{\pgfqpoint{0.100000in}{2.493063in}}{\pgfqpoint{5.037500in}{3.427208in}}%
\pgfusepath{clip}%
\pgfsetrectcap%
\pgfsetroundjoin%
\pgfsetlinewidth{1.505625pt}%
\definecolor{currentstroke}{rgb}{0.678431,1.000000,0.184314}%
\pgfsetstrokecolor{currentstroke}%
\pgfsetstrokeopacity{0.500000}%
\pgfsetdash{}{0pt}%
\pgfpathmoveto{\pgfqpoint{0.794333in}{5.522772in}}%
\pgfusepath{stroke}%
\end{pgfscope}%
\begin{pgfscope}%
\pgfpathrectangle{\pgfqpoint{0.100000in}{2.493063in}}{\pgfqpoint{5.037500in}{3.427208in}}%
\pgfusepath{clip}%
\pgfsetbuttcap%
\pgfsetroundjoin%
\definecolor{currentfill}{rgb}{0.678431,1.000000,0.184314}%
\pgfsetfillcolor{currentfill}%
\pgfsetfillopacity{0.500000}%
\pgfsetlinewidth{0.250937pt}%
\definecolor{currentstroke}{rgb}{0.000000,0.000000,0.000000}%
\pgfsetstrokecolor{currentstroke}%
\pgfsetstrokeopacity{0.500000}%
\pgfsetdash{}{0pt}%
\pgfsys@defobject{currentmarker}{\pgfqpoint{-0.025000in}{-0.025000in}}{\pgfqpoint{0.025000in}{0.025000in}}{%
\pgfpathmoveto{\pgfqpoint{0.000000in}{-0.025000in}}%
\pgfpathcurveto{\pgfqpoint{0.006630in}{-0.025000in}}{\pgfqpoint{0.012989in}{-0.022366in}}{\pgfqpoint{0.017678in}{-0.017678in}}%
\pgfpathcurveto{\pgfqpoint{0.022366in}{-0.012989in}}{\pgfqpoint{0.025000in}{-0.006630in}}{\pgfqpoint{0.025000in}{0.000000in}}%
\pgfpathcurveto{\pgfqpoint{0.025000in}{0.006630in}}{\pgfqpoint{0.022366in}{0.012989in}}{\pgfqpoint{0.017678in}{0.017678in}}%
\pgfpathcurveto{\pgfqpoint{0.012989in}{0.022366in}}{\pgfqpoint{0.006630in}{0.025000in}}{\pgfqpoint{0.000000in}{0.025000in}}%
\pgfpathcurveto{\pgfqpoint{-0.006630in}{0.025000in}}{\pgfqpoint{-0.012989in}{0.022366in}}{\pgfqpoint{-0.017678in}{0.017678in}}%
\pgfpathcurveto{\pgfqpoint{-0.022366in}{0.012989in}}{\pgfqpoint{-0.025000in}{0.006630in}}{\pgfqpoint{-0.025000in}{0.000000in}}%
\pgfpathcurveto{\pgfqpoint{-0.025000in}{-0.006630in}}{\pgfqpoint{-0.022366in}{-0.012989in}}{\pgfqpoint{-0.017678in}{-0.017678in}}%
\pgfpathcurveto{\pgfqpoint{-0.012989in}{-0.022366in}}{\pgfqpoint{-0.006630in}{-0.025000in}}{\pgfqpoint{0.000000in}{-0.025000in}}%
\pgfpathlineto{\pgfqpoint{0.000000in}{-0.025000in}}%
\pgfpathclose%
\pgfusepath{stroke,fill}%
}%
\begin{pgfscope}%
\pgfsys@transformshift{0.794333in}{5.522772in}%
\pgfsys@useobject{currentmarker}{}%
\end{pgfscope}%
\end{pgfscope}%
\begin{pgfscope}%
\pgfpathrectangle{\pgfqpoint{0.100000in}{2.493063in}}{\pgfqpoint{5.037500in}{3.427208in}}%
\pgfusepath{clip}%
\pgfsetrectcap%
\pgfsetroundjoin%
\pgfsetlinewidth{1.505625pt}%
\definecolor{currentstroke}{rgb}{0.678431,1.000000,0.184314}%
\pgfsetstrokecolor{currentstroke}%
\pgfsetstrokeopacity{0.500000}%
\pgfsetdash{}{0pt}%
\pgfpathmoveto{\pgfqpoint{4.056648in}{4.315611in}}%
\pgfusepath{stroke}%
\end{pgfscope}%
\begin{pgfscope}%
\pgfpathrectangle{\pgfqpoint{0.100000in}{2.493063in}}{\pgfqpoint{5.037500in}{3.427208in}}%
\pgfusepath{clip}%
\pgfsetbuttcap%
\pgfsetroundjoin%
\definecolor{currentfill}{rgb}{0.678431,1.000000,0.184314}%
\pgfsetfillcolor{currentfill}%
\pgfsetfillopacity{0.500000}%
\pgfsetlinewidth{0.250937pt}%
\definecolor{currentstroke}{rgb}{0.000000,0.000000,0.000000}%
\pgfsetstrokecolor{currentstroke}%
\pgfsetstrokeopacity{0.500000}%
\pgfsetdash{}{0pt}%
\pgfsys@defobject{currentmarker}{\pgfqpoint{-0.030556in}{-0.030556in}}{\pgfqpoint{0.030556in}{0.030556in}}{%
\pgfpathmoveto{\pgfqpoint{0.000000in}{-0.030556in}}%
\pgfpathcurveto{\pgfqpoint{0.008103in}{-0.030556in}}{\pgfqpoint{0.015876in}{-0.027336in}}{\pgfqpoint{0.021606in}{-0.021606in}}%
\pgfpathcurveto{\pgfqpoint{0.027336in}{-0.015876in}}{\pgfqpoint{0.030556in}{-0.008103in}}{\pgfqpoint{0.030556in}{0.000000in}}%
\pgfpathcurveto{\pgfqpoint{0.030556in}{0.008103in}}{\pgfqpoint{0.027336in}{0.015876in}}{\pgfqpoint{0.021606in}{0.021606in}}%
\pgfpathcurveto{\pgfqpoint{0.015876in}{0.027336in}}{\pgfqpoint{0.008103in}{0.030556in}}{\pgfqpoint{0.000000in}{0.030556in}}%
\pgfpathcurveto{\pgfqpoint{-0.008103in}{0.030556in}}{\pgfqpoint{-0.015876in}{0.027336in}}{\pgfqpoint{-0.021606in}{0.021606in}}%
\pgfpathcurveto{\pgfqpoint{-0.027336in}{0.015876in}}{\pgfqpoint{-0.030556in}{0.008103in}}{\pgfqpoint{-0.030556in}{0.000000in}}%
\pgfpathcurveto{\pgfqpoint{-0.030556in}{-0.008103in}}{\pgfqpoint{-0.027336in}{-0.015876in}}{\pgfqpoint{-0.021606in}{-0.021606in}}%
\pgfpathcurveto{\pgfqpoint{-0.015876in}{-0.027336in}}{\pgfqpoint{-0.008103in}{-0.030556in}}{\pgfqpoint{0.000000in}{-0.030556in}}%
\pgfpathlineto{\pgfqpoint{0.000000in}{-0.030556in}}%
\pgfpathclose%
\pgfusepath{stroke,fill}%
}%
\begin{pgfscope}%
\pgfsys@transformshift{4.056648in}{4.315611in}%
\pgfsys@useobject{currentmarker}{}%
\end{pgfscope}%
\end{pgfscope}%
\begin{pgfscope}%
\pgfpathrectangle{\pgfqpoint{0.100000in}{2.493063in}}{\pgfqpoint{5.037500in}{3.427208in}}%
\pgfusepath{clip}%
\pgfsetrectcap%
\pgfsetroundjoin%
\pgfsetlinewidth{1.505625pt}%
\definecolor{currentstroke}{rgb}{0.678431,1.000000,0.184314}%
\pgfsetstrokecolor{currentstroke}%
\pgfsetstrokeopacity{0.500000}%
\pgfsetdash{}{0pt}%
\pgfpathmoveto{\pgfqpoint{4.006916in}{4.374787in}}%
\pgfusepath{stroke}%
\end{pgfscope}%
\begin{pgfscope}%
\pgfpathrectangle{\pgfqpoint{0.100000in}{2.493063in}}{\pgfqpoint{5.037500in}{3.427208in}}%
\pgfusepath{clip}%
\pgfsetbuttcap%
\pgfsetroundjoin%
\definecolor{currentfill}{rgb}{0.678431,1.000000,0.184314}%
\pgfsetfillcolor{currentfill}%
\pgfsetfillopacity{0.500000}%
\pgfsetlinewidth{0.250937pt}%
\definecolor{currentstroke}{rgb}{0.000000,0.000000,0.000000}%
\pgfsetstrokecolor{currentstroke}%
\pgfsetstrokeopacity{0.500000}%
\pgfsetdash{}{0pt}%
\pgfsys@defobject{currentmarker}{\pgfqpoint{-0.022222in}{-0.022222in}}{\pgfqpoint{0.022222in}{0.022222in}}{%
\pgfpathmoveto{\pgfqpoint{0.000000in}{-0.022222in}}%
\pgfpathcurveto{\pgfqpoint{0.005893in}{-0.022222in}}{\pgfqpoint{0.011546in}{-0.019881in}}{\pgfqpoint{0.015713in}{-0.015713in}}%
\pgfpathcurveto{\pgfqpoint{0.019881in}{-0.011546in}}{\pgfqpoint{0.022222in}{-0.005893in}}{\pgfqpoint{0.022222in}{0.000000in}}%
\pgfpathcurveto{\pgfqpoint{0.022222in}{0.005893in}}{\pgfqpoint{0.019881in}{0.011546in}}{\pgfqpoint{0.015713in}{0.015713in}}%
\pgfpathcurveto{\pgfqpoint{0.011546in}{0.019881in}}{\pgfqpoint{0.005893in}{0.022222in}}{\pgfqpoint{0.000000in}{0.022222in}}%
\pgfpathcurveto{\pgfqpoint{-0.005893in}{0.022222in}}{\pgfqpoint{-0.011546in}{0.019881in}}{\pgfqpoint{-0.015713in}{0.015713in}}%
\pgfpathcurveto{\pgfqpoint{-0.019881in}{0.011546in}}{\pgfqpoint{-0.022222in}{0.005893in}}{\pgfqpoint{-0.022222in}{0.000000in}}%
\pgfpathcurveto{\pgfqpoint{-0.022222in}{-0.005893in}}{\pgfqpoint{-0.019881in}{-0.011546in}}{\pgfqpoint{-0.015713in}{-0.015713in}}%
\pgfpathcurveto{\pgfqpoint{-0.011546in}{-0.019881in}}{\pgfqpoint{-0.005893in}{-0.022222in}}{\pgfqpoint{0.000000in}{-0.022222in}}%
\pgfpathlineto{\pgfqpoint{0.000000in}{-0.022222in}}%
\pgfpathclose%
\pgfusepath{stroke,fill}%
}%
\begin{pgfscope}%
\pgfsys@transformshift{4.006916in}{4.374787in}%
\pgfsys@useobject{currentmarker}{}%
\end{pgfscope}%
\end{pgfscope}%
\begin{pgfscope}%
\pgfpathrectangle{\pgfqpoint{0.100000in}{2.493063in}}{\pgfqpoint{5.037500in}{3.427208in}}%
\pgfusepath{clip}%
\pgfsetrectcap%
\pgfsetroundjoin%
\pgfsetlinewidth{1.505625pt}%
\definecolor{currentstroke}{rgb}{0.678431,1.000000,0.184314}%
\pgfsetstrokecolor{currentstroke}%
\pgfsetstrokeopacity{0.500000}%
\pgfsetdash{}{0pt}%
\pgfpathmoveto{\pgfqpoint{3.933253in}{4.372196in}}%
\pgfusepath{stroke}%
\end{pgfscope}%
\begin{pgfscope}%
\pgfpathrectangle{\pgfqpoint{0.100000in}{2.493063in}}{\pgfqpoint{5.037500in}{3.427208in}}%
\pgfusepath{clip}%
\pgfsetbuttcap%
\pgfsetroundjoin%
\definecolor{currentfill}{rgb}{0.678431,1.000000,0.184314}%
\pgfsetfillcolor{currentfill}%
\pgfsetfillopacity{0.500000}%
\pgfsetlinewidth{0.250937pt}%
\definecolor{currentstroke}{rgb}{0.000000,0.000000,0.000000}%
\pgfsetstrokecolor{currentstroke}%
\pgfsetstrokeopacity{0.500000}%
\pgfsetdash{}{0pt}%
\pgfsys@defobject{currentmarker}{\pgfqpoint{-0.013889in}{-0.013889in}}{\pgfqpoint{0.013889in}{0.013889in}}{%
\pgfpathmoveto{\pgfqpoint{0.000000in}{-0.013889in}}%
\pgfpathcurveto{\pgfqpoint{0.003683in}{-0.013889in}}{\pgfqpoint{0.007216in}{-0.012425in}}{\pgfqpoint{0.009821in}{-0.009821in}}%
\pgfpathcurveto{\pgfqpoint{0.012425in}{-0.007216in}}{\pgfqpoint{0.013889in}{-0.003683in}}{\pgfqpoint{0.013889in}{0.000000in}}%
\pgfpathcurveto{\pgfqpoint{0.013889in}{0.003683in}}{\pgfqpoint{0.012425in}{0.007216in}}{\pgfqpoint{0.009821in}{0.009821in}}%
\pgfpathcurveto{\pgfqpoint{0.007216in}{0.012425in}}{\pgfqpoint{0.003683in}{0.013889in}}{\pgfqpoint{0.000000in}{0.013889in}}%
\pgfpathcurveto{\pgfqpoint{-0.003683in}{0.013889in}}{\pgfqpoint{-0.007216in}{0.012425in}}{\pgfqpoint{-0.009821in}{0.009821in}}%
\pgfpathcurveto{\pgfqpoint{-0.012425in}{0.007216in}}{\pgfqpoint{-0.013889in}{0.003683in}}{\pgfqpoint{-0.013889in}{0.000000in}}%
\pgfpathcurveto{\pgfqpoint{-0.013889in}{-0.003683in}}{\pgfqpoint{-0.012425in}{-0.007216in}}{\pgfqpoint{-0.009821in}{-0.009821in}}%
\pgfpathcurveto{\pgfqpoint{-0.007216in}{-0.012425in}}{\pgfqpoint{-0.003683in}{-0.013889in}}{\pgfqpoint{0.000000in}{-0.013889in}}%
\pgfpathlineto{\pgfqpoint{0.000000in}{-0.013889in}}%
\pgfpathclose%
\pgfusepath{stroke,fill}%
}%
\begin{pgfscope}%
\pgfsys@transformshift{3.933253in}{4.372196in}%
\pgfsys@useobject{currentmarker}{}%
\end{pgfscope}%
\end{pgfscope}%
\begin{pgfscope}%
\pgfpathrectangle{\pgfqpoint{0.100000in}{2.493063in}}{\pgfqpoint{5.037500in}{3.427208in}}%
\pgfusepath{clip}%
\pgfsetrectcap%
\pgfsetroundjoin%
\pgfsetlinewidth{1.505625pt}%
\definecolor{currentstroke}{rgb}{0.678431,1.000000,0.184314}%
\pgfsetstrokecolor{currentstroke}%
\pgfsetstrokeopacity{0.500000}%
\pgfsetdash{}{0pt}%
\pgfpathmoveto{\pgfqpoint{4.132275in}{4.543641in}}%
\pgfusepath{stroke}%
\end{pgfscope}%
\begin{pgfscope}%
\pgfpathrectangle{\pgfqpoint{0.100000in}{2.493063in}}{\pgfqpoint{5.037500in}{3.427208in}}%
\pgfusepath{clip}%
\pgfsetbuttcap%
\pgfsetroundjoin%
\definecolor{currentfill}{rgb}{0.678431,1.000000,0.184314}%
\pgfsetfillcolor{currentfill}%
\pgfsetfillopacity{0.500000}%
\pgfsetlinewidth{0.250937pt}%
\definecolor{currentstroke}{rgb}{0.000000,0.000000,0.000000}%
\pgfsetstrokecolor{currentstroke}%
\pgfsetstrokeopacity{0.500000}%
\pgfsetdash{}{0pt}%
\pgfsys@defobject{currentmarker}{\pgfqpoint{-0.011111in}{-0.011111in}}{\pgfqpoint{0.011111in}{0.011111in}}{%
\pgfpathmoveto{\pgfqpoint{0.000000in}{-0.011111in}}%
\pgfpathcurveto{\pgfqpoint{0.002947in}{-0.011111in}}{\pgfqpoint{0.005773in}{-0.009940in}}{\pgfqpoint{0.007857in}{-0.007857in}}%
\pgfpathcurveto{\pgfqpoint{0.009940in}{-0.005773in}}{\pgfqpoint{0.011111in}{-0.002947in}}{\pgfqpoint{0.011111in}{0.000000in}}%
\pgfpathcurveto{\pgfqpoint{0.011111in}{0.002947in}}{\pgfqpoint{0.009940in}{0.005773in}}{\pgfqpoint{0.007857in}{0.007857in}}%
\pgfpathcurveto{\pgfqpoint{0.005773in}{0.009940in}}{\pgfqpoint{0.002947in}{0.011111in}}{\pgfqpoint{0.000000in}{0.011111in}}%
\pgfpathcurveto{\pgfqpoint{-0.002947in}{0.011111in}}{\pgfqpoint{-0.005773in}{0.009940in}}{\pgfqpoint{-0.007857in}{0.007857in}}%
\pgfpathcurveto{\pgfqpoint{-0.009940in}{0.005773in}}{\pgfqpoint{-0.011111in}{0.002947in}}{\pgfqpoint{-0.011111in}{0.000000in}}%
\pgfpathcurveto{\pgfqpoint{-0.011111in}{-0.002947in}}{\pgfqpoint{-0.009940in}{-0.005773in}}{\pgfqpoint{-0.007857in}{-0.007857in}}%
\pgfpathcurveto{\pgfqpoint{-0.005773in}{-0.009940in}}{\pgfqpoint{-0.002947in}{-0.011111in}}{\pgfqpoint{0.000000in}{-0.011111in}}%
\pgfpathlineto{\pgfqpoint{0.000000in}{-0.011111in}}%
\pgfpathclose%
\pgfusepath{stroke,fill}%
}%
\begin{pgfscope}%
\pgfsys@transformshift{4.132275in}{4.543641in}%
\pgfsys@useobject{currentmarker}{}%
\end{pgfscope}%
\end{pgfscope}%
\begin{pgfscope}%
\pgfpathrectangle{\pgfqpoint{0.100000in}{2.493063in}}{\pgfqpoint{5.037500in}{3.427208in}}%
\pgfusepath{clip}%
\pgfsetrectcap%
\pgfsetroundjoin%
\pgfsetlinewidth{1.505625pt}%
\definecolor{currentstroke}{rgb}{0.678431,1.000000,0.184314}%
\pgfsetstrokecolor{currentstroke}%
\pgfsetstrokeopacity{0.500000}%
\pgfsetdash{}{0pt}%
\pgfpathmoveto{\pgfqpoint{3.998684in}{4.480590in}}%
\pgfusepath{stroke}%
\end{pgfscope}%
\begin{pgfscope}%
\pgfpathrectangle{\pgfqpoint{0.100000in}{2.493063in}}{\pgfqpoint{5.037500in}{3.427208in}}%
\pgfusepath{clip}%
\pgfsetbuttcap%
\pgfsetroundjoin%
\definecolor{currentfill}{rgb}{0.678431,1.000000,0.184314}%
\pgfsetfillcolor{currentfill}%
\pgfsetfillopacity{0.500000}%
\pgfsetlinewidth{0.250937pt}%
\definecolor{currentstroke}{rgb}{0.000000,0.000000,0.000000}%
\pgfsetstrokecolor{currentstroke}%
\pgfsetstrokeopacity{0.500000}%
\pgfsetdash{}{0pt}%
\pgfsys@defobject{currentmarker}{\pgfqpoint{-0.030556in}{-0.030556in}}{\pgfqpoint{0.030556in}{0.030556in}}{%
\pgfpathmoveto{\pgfqpoint{0.000000in}{-0.030556in}}%
\pgfpathcurveto{\pgfqpoint{0.008103in}{-0.030556in}}{\pgfqpoint{0.015876in}{-0.027336in}}{\pgfqpoint{0.021606in}{-0.021606in}}%
\pgfpathcurveto{\pgfqpoint{0.027336in}{-0.015876in}}{\pgfqpoint{0.030556in}{-0.008103in}}{\pgfqpoint{0.030556in}{0.000000in}}%
\pgfpathcurveto{\pgfqpoint{0.030556in}{0.008103in}}{\pgfqpoint{0.027336in}{0.015876in}}{\pgfqpoint{0.021606in}{0.021606in}}%
\pgfpathcurveto{\pgfqpoint{0.015876in}{0.027336in}}{\pgfqpoint{0.008103in}{0.030556in}}{\pgfqpoint{0.000000in}{0.030556in}}%
\pgfpathcurveto{\pgfqpoint{-0.008103in}{0.030556in}}{\pgfqpoint{-0.015876in}{0.027336in}}{\pgfqpoint{-0.021606in}{0.021606in}}%
\pgfpathcurveto{\pgfqpoint{-0.027336in}{0.015876in}}{\pgfqpoint{-0.030556in}{0.008103in}}{\pgfqpoint{-0.030556in}{0.000000in}}%
\pgfpathcurveto{\pgfqpoint{-0.030556in}{-0.008103in}}{\pgfqpoint{-0.027336in}{-0.015876in}}{\pgfqpoint{-0.021606in}{-0.021606in}}%
\pgfpathcurveto{\pgfqpoint{-0.015876in}{-0.027336in}}{\pgfqpoint{-0.008103in}{-0.030556in}}{\pgfqpoint{0.000000in}{-0.030556in}}%
\pgfpathlineto{\pgfqpoint{0.000000in}{-0.030556in}}%
\pgfpathclose%
\pgfusepath{stroke,fill}%
}%
\begin{pgfscope}%
\pgfsys@transformshift{3.998684in}{4.480590in}%
\pgfsys@useobject{currentmarker}{}%
\end{pgfscope}%
\end{pgfscope}%
\begin{pgfscope}%
\pgfpathrectangle{\pgfqpoint{0.100000in}{2.493063in}}{\pgfqpoint{5.037500in}{3.427208in}}%
\pgfusepath{clip}%
\pgfsetrectcap%
\pgfsetroundjoin%
\pgfsetlinewidth{1.505625pt}%
\definecolor{currentstroke}{rgb}{0.678431,1.000000,0.184314}%
\pgfsetstrokecolor{currentstroke}%
\pgfsetstrokeopacity{0.500000}%
\pgfsetdash{}{0pt}%
\pgfpathmoveto{\pgfqpoint{4.057487in}{4.582137in}}%
\pgfusepath{stroke}%
\end{pgfscope}%
\begin{pgfscope}%
\pgfpathrectangle{\pgfqpoint{0.100000in}{2.493063in}}{\pgfqpoint{5.037500in}{3.427208in}}%
\pgfusepath{clip}%
\pgfsetbuttcap%
\pgfsetroundjoin%
\definecolor{currentfill}{rgb}{0.678431,1.000000,0.184314}%
\pgfsetfillcolor{currentfill}%
\pgfsetfillopacity{0.500000}%
\pgfsetlinewidth{0.250937pt}%
\definecolor{currentstroke}{rgb}{0.000000,0.000000,0.000000}%
\pgfsetstrokecolor{currentstroke}%
\pgfsetstrokeopacity{0.500000}%
\pgfsetdash{}{0pt}%
\pgfsys@defobject{currentmarker}{\pgfqpoint{-0.050000in}{-0.050000in}}{\pgfqpoint{0.050000in}{0.050000in}}{%
\pgfpathmoveto{\pgfqpoint{0.000000in}{-0.050000in}}%
\pgfpathcurveto{\pgfqpoint{0.013260in}{-0.050000in}}{\pgfqpoint{0.025979in}{-0.044732in}}{\pgfqpoint{0.035355in}{-0.035355in}}%
\pgfpathcurveto{\pgfqpoint{0.044732in}{-0.025979in}}{\pgfqpoint{0.050000in}{-0.013260in}}{\pgfqpoint{0.050000in}{0.000000in}}%
\pgfpathcurveto{\pgfqpoint{0.050000in}{0.013260in}}{\pgfqpoint{0.044732in}{0.025979in}}{\pgfqpoint{0.035355in}{0.035355in}}%
\pgfpathcurveto{\pgfqpoint{0.025979in}{0.044732in}}{\pgfqpoint{0.013260in}{0.050000in}}{\pgfqpoint{0.000000in}{0.050000in}}%
\pgfpathcurveto{\pgfqpoint{-0.013260in}{0.050000in}}{\pgfqpoint{-0.025979in}{0.044732in}}{\pgfqpoint{-0.035355in}{0.035355in}}%
\pgfpathcurveto{\pgfqpoint{-0.044732in}{0.025979in}}{\pgfqpoint{-0.050000in}{0.013260in}}{\pgfqpoint{-0.050000in}{0.000000in}}%
\pgfpathcurveto{\pgfqpoint{-0.050000in}{-0.013260in}}{\pgfqpoint{-0.044732in}{-0.025979in}}{\pgfqpoint{-0.035355in}{-0.035355in}}%
\pgfpathcurveto{\pgfqpoint{-0.025979in}{-0.044732in}}{\pgfqpoint{-0.013260in}{-0.050000in}}{\pgfqpoint{0.000000in}{-0.050000in}}%
\pgfpathlineto{\pgfqpoint{0.000000in}{-0.050000in}}%
\pgfpathclose%
\pgfusepath{stroke,fill}%
}%
\begin{pgfscope}%
\pgfsys@transformshift{4.057487in}{4.582137in}%
\pgfsys@useobject{currentmarker}{}%
\end{pgfscope}%
\end{pgfscope}%
\begin{pgfscope}%
\pgfpathrectangle{\pgfqpoint{0.100000in}{2.493063in}}{\pgfqpoint{5.037500in}{3.427208in}}%
\pgfusepath{clip}%
\pgfsetrectcap%
\pgfsetroundjoin%
\pgfsetlinewidth{1.505625pt}%
\definecolor{currentstroke}{rgb}{0.501961,0.501961,0.501961}%
\pgfsetstrokecolor{currentstroke}%
\pgfsetstrokeopacity{0.500000}%
\pgfsetdash{}{0pt}%
\pgfpathmoveto{\pgfqpoint{3.349641in}{4.987379in}}%
\pgfusepath{stroke}%
\end{pgfscope}%
\begin{pgfscope}%
\pgfpathrectangle{\pgfqpoint{0.100000in}{2.493063in}}{\pgfqpoint{5.037500in}{3.427208in}}%
\pgfusepath{clip}%
\pgfsetbuttcap%
\pgfsetroundjoin%
\definecolor{currentfill}{rgb}{0.501961,0.501961,0.501961}%
\pgfsetfillcolor{currentfill}%
\pgfsetfillopacity{0.500000}%
\pgfsetlinewidth{0.250937pt}%
\definecolor{currentstroke}{rgb}{0.000000,0.000000,0.000000}%
\pgfsetstrokecolor{currentstroke}%
\pgfsetstrokeopacity{0.500000}%
\pgfsetdash{}{0pt}%
\pgfsys@defobject{currentmarker}{\pgfqpoint{-0.013889in}{-0.013889in}}{\pgfqpoint{0.013889in}{0.013889in}}{%
\pgfpathmoveto{\pgfqpoint{0.000000in}{-0.013889in}}%
\pgfpathcurveto{\pgfqpoint{0.003683in}{-0.013889in}}{\pgfqpoint{0.007216in}{-0.012425in}}{\pgfqpoint{0.009821in}{-0.009821in}}%
\pgfpathcurveto{\pgfqpoint{0.012425in}{-0.007216in}}{\pgfqpoint{0.013889in}{-0.003683in}}{\pgfqpoint{0.013889in}{0.000000in}}%
\pgfpathcurveto{\pgfqpoint{0.013889in}{0.003683in}}{\pgfqpoint{0.012425in}{0.007216in}}{\pgfqpoint{0.009821in}{0.009821in}}%
\pgfpathcurveto{\pgfqpoint{0.007216in}{0.012425in}}{\pgfqpoint{0.003683in}{0.013889in}}{\pgfqpoint{0.000000in}{0.013889in}}%
\pgfpathcurveto{\pgfqpoint{-0.003683in}{0.013889in}}{\pgfqpoint{-0.007216in}{0.012425in}}{\pgfqpoint{-0.009821in}{0.009821in}}%
\pgfpathcurveto{\pgfqpoint{-0.012425in}{0.007216in}}{\pgfqpoint{-0.013889in}{0.003683in}}{\pgfqpoint{-0.013889in}{0.000000in}}%
\pgfpathcurveto{\pgfqpoint{-0.013889in}{-0.003683in}}{\pgfqpoint{-0.012425in}{-0.007216in}}{\pgfqpoint{-0.009821in}{-0.009821in}}%
\pgfpathcurveto{\pgfqpoint{-0.007216in}{-0.012425in}}{\pgfqpoint{-0.003683in}{-0.013889in}}{\pgfqpoint{0.000000in}{-0.013889in}}%
\pgfpathlineto{\pgfqpoint{0.000000in}{-0.013889in}}%
\pgfpathclose%
\pgfusepath{stroke,fill}%
}%
\begin{pgfscope}%
\pgfsys@transformshift{3.349641in}{4.987379in}%
\pgfsys@useobject{currentmarker}{}%
\end{pgfscope}%
\end{pgfscope}%
\begin{pgfscope}%
\pgfpathrectangle{\pgfqpoint{0.100000in}{2.493063in}}{\pgfqpoint{5.037500in}{3.427208in}}%
\pgfusepath{clip}%
\pgfsetrectcap%
\pgfsetroundjoin%
\pgfsetlinewidth{1.505625pt}%
\definecolor{currentstroke}{rgb}{0.501961,0.501961,0.501961}%
\pgfsetstrokecolor{currentstroke}%
\pgfsetstrokeopacity{0.500000}%
\pgfsetdash{}{0pt}%
\pgfpathmoveto{\pgfqpoint{3.091567in}{5.036771in}}%
\pgfusepath{stroke}%
\end{pgfscope}%
\begin{pgfscope}%
\pgfpathrectangle{\pgfqpoint{0.100000in}{2.493063in}}{\pgfqpoint{5.037500in}{3.427208in}}%
\pgfusepath{clip}%
\pgfsetbuttcap%
\pgfsetroundjoin%
\definecolor{currentfill}{rgb}{0.501961,0.501961,0.501961}%
\pgfsetfillcolor{currentfill}%
\pgfsetfillopacity{0.500000}%
\pgfsetlinewidth{0.250937pt}%
\definecolor{currentstroke}{rgb}{0.000000,0.000000,0.000000}%
\pgfsetstrokecolor{currentstroke}%
\pgfsetstrokeopacity{0.500000}%
\pgfsetdash{}{0pt}%
\pgfsys@defobject{currentmarker}{\pgfqpoint{-0.013889in}{-0.013889in}}{\pgfqpoint{0.013889in}{0.013889in}}{%
\pgfpathmoveto{\pgfqpoint{0.000000in}{-0.013889in}}%
\pgfpathcurveto{\pgfqpoint{0.003683in}{-0.013889in}}{\pgfqpoint{0.007216in}{-0.012425in}}{\pgfqpoint{0.009821in}{-0.009821in}}%
\pgfpathcurveto{\pgfqpoint{0.012425in}{-0.007216in}}{\pgfqpoint{0.013889in}{-0.003683in}}{\pgfqpoint{0.013889in}{0.000000in}}%
\pgfpathcurveto{\pgfqpoint{0.013889in}{0.003683in}}{\pgfqpoint{0.012425in}{0.007216in}}{\pgfqpoint{0.009821in}{0.009821in}}%
\pgfpathcurveto{\pgfqpoint{0.007216in}{0.012425in}}{\pgfqpoint{0.003683in}{0.013889in}}{\pgfqpoint{0.000000in}{0.013889in}}%
\pgfpathcurveto{\pgfqpoint{-0.003683in}{0.013889in}}{\pgfqpoint{-0.007216in}{0.012425in}}{\pgfqpoint{-0.009821in}{0.009821in}}%
\pgfpathcurveto{\pgfqpoint{-0.012425in}{0.007216in}}{\pgfqpoint{-0.013889in}{0.003683in}}{\pgfqpoint{-0.013889in}{0.000000in}}%
\pgfpathcurveto{\pgfqpoint{-0.013889in}{-0.003683in}}{\pgfqpoint{-0.012425in}{-0.007216in}}{\pgfqpoint{-0.009821in}{-0.009821in}}%
\pgfpathcurveto{\pgfqpoint{-0.007216in}{-0.012425in}}{\pgfqpoint{-0.003683in}{-0.013889in}}{\pgfqpoint{0.000000in}{-0.013889in}}%
\pgfpathlineto{\pgfqpoint{0.000000in}{-0.013889in}}%
\pgfpathclose%
\pgfusepath{stroke,fill}%
}%
\begin{pgfscope}%
\pgfsys@transformshift{3.091567in}{5.036771in}%
\pgfsys@useobject{currentmarker}{}%
\end{pgfscope}%
\end{pgfscope}%
\begin{pgfscope}%
\pgfpathrectangle{\pgfqpoint{0.100000in}{2.493063in}}{\pgfqpoint{5.037500in}{3.427208in}}%
\pgfusepath{clip}%
\pgfsetrectcap%
\pgfsetroundjoin%
\pgfsetlinewidth{1.505625pt}%
\definecolor{currentstroke}{rgb}{0.501961,0.501961,0.501961}%
\pgfsetstrokecolor{currentstroke}%
\pgfsetstrokeopacity{0.500000}%
\pgfsetdash{}{0pt}%
\pgfpathmoveto{\pgfqpoint{3.350467in}{4.931087in}}%
\pgfusepath{stroke}%
\end{pgfscope}%
\begin{pgfscope}%
\pgfpathrectangle{\pgfqpoint{0.100000in}{2.493063in}}{\pgfqpoint{5.037500in}{3.427208in}}%
\pgfusepath{clip}%
\pgfsetbuttcap%
\pgfsetroundjoin%
\definecolor{currentfill}{rgb}{0.501961,0.501961,0.501961}%
\pgfsetfillcolor{currentfill}%
\pgfsetfillopacity{0.500000}%
\pgfsetlinewidth{0.250937pt}%
\definecolor{currentstroke}{rgb}{0.000000,0.000000,0.000000}%
\pgfsetstrokecolor{currentstroke}%
\pgfsetstrokeopacity{0.500000}%
\pgfsetdash{}{0pt}%
\pgfsys@defobject{currentmarker}{\pgfqpoint{-0.013889in}{-0.013889in}}{\pgfqpoint{0.013889in}{0.013889in}}{%
\pgfpathmoveto{\pgfqpoint{0.000000in}{-0.013889in}}%
\pgfpathcurveto{\pgfqpoint{0.003683in}{-0.013889in}}{\pgfqpoint{0.007216in}{-0.012425in}}{\pgfqpoint{0.009821in}{-0.009821in}}%
\pgfpathcurveto{\pgfqpoint{0.012425in}{-0.007216in}}{\pgfqpoint{0.013889in}{-0.003683in}}{\pgfqpoint{0.013889in}{0.000000in}}%
\pgfpathcurveto{\pgfqpoint{0.013889in}{0.003683in}}{\pgfqpoint{0.012425in}{0.007216in}}{\pgfqpoint{0.009821in}{0.009821in}}%
\pgfpathcurveto{\pgfqpoint{0.007216in}{0.012425in}}{\pgfqpoint{0.003683in}{0.013889in}}{\pgfqpoint{0.000000in}{0.013889in}}%
\pgfpathcurveto{\pgfqpoint{-0.003683in}{0.013889in}}{\pgfqpoint{-0.007216in}{0.012425in}}{\pgfqpoint{-0.009821in}{0.009821in}}%
\pgfpathcurveto{\pgfqpoint{-0.012425in}{0.007216in}}{\pgfqpoint{-0.013889in}{0.003683in}}{\pgfqpoint{-0.013889in}{0.000000in}}%
\pgfpathcurveto{\pgfqpoint{-0.013889in}{-0.003683in}}{\pgfqpoint{-0.012425in}{-0.007216in}}{\pgfqpoint{-0.009821in}{-0.009821in}}%
\pgfpathcurveto{\pgfqpoint{-0.007216in}{-0.012425in}}{\pgfqpoint{-0.003683in}{-0.013889in}}{\pgfqpoint{0.000000in}{-0.013889in}}%
\pgfpathlineto{\pgfqpoint{0.000000in}{-0.013889in}}%
\pgfpathclose%
\pgfusepath{stroke,fill}%
}%
\begin{pgfscope}%
\pgfsys@transformshift{3.350467in}{4.931087in}%
\pgfsys@useobject{currentmarker}{}%
\end{pgfscope}%
\end{pgfscope}%
\begin{pgfscope}%
\pgfpathrectangle{\pgfqpoint{0.100000in}{2.493063in}}{\pgfqpoint{5.037500in}{3.427208in}}%
\pgfusepath{clip}%
\pgfsetrectcap%
\pgfsetroundjoin%
\pgfsetlinewidth{1.505625pt}%
\definecolor{currentstroke}{rgb}{0.678431,1.000000,0.184314}%
\pgfsetstrokecolor{currentstroke}%
\pgfsetstrokeopacity{0.500000}%
\pgfsetdash{}{0pt}%
\pgfpathmoveto{\pgfqpoint{3.378797in}{5.020736in}}%
\pgfusepath{stroke}%
\end{pgfscope}%
\begin{pgfscope}%
\pgfpathrectangle{\pgfqpoint{0.100000in}{2.493063in}}{\pgfqpoint{5.037500in}{3.427208in}}%
\pgfusepath{clip}%
\pgfsetbuttcap%
\pgfsetroundjoin%
\definecolor{currentfill}{rgb}{0.678431,1.000000,0.184314}%
\pgfsetfillcolor{currentfill}%
\pgfsetfillopacity{0.500000}%
\pgfsetlinewidth{0.250937pt}%
\definecolor{currentstroke}{rgb}{0.000000,0.000000,0.000000}%
\pgfsetstrokecolor{currentstroke}%
\pgfsetstrokeopacity{0.500000}%
\pgfsetdash{}{0pt}%
\pgfsys@defobject{currentmarker}{\pgfqpoint{-0.005556in}{-0.005556in}}{\pgfqpoint{0.005556in}{0.005556in}}{%
\pgfpathmoveto{\pgfqpoint{0.000000in}{-0.005556in}}%
\pgfpathcurveto{\pgfqpoint{0.001473in}{-0.005556in}}{\pgfqpoint{0.002887in}{-0.004970in}}{\pgfqpoint{0.003928in}{-0.003928in}}%
\pgfpathcurveto{\pgfqpoint{0.004970in}{-0.002887in}}{\pgfqpoint{0.005556in}{-0.001473in}}{\pgfqpoint{0.005556in}{0.000000in}}%
\pgfpathcurveto{\pgfqpoint{0.005556in}{0.001473in}}{\pgfqpoint{0.004970in}{0.002887in}}{\pgfqpoint{0.003928in}{0.003928in}}%
\pgfpathcurveto{\pgfqpoint{0.002887in}{0.004970in}}{\pgfqpoint{0.001473in}{0.005556in}}{\pgfqpoint{0.000000in}{0.005556in}}%
\pgfpathcurveto{\pgfqpoint{-0.001473in}{0.005556in}}{\pgfqpoint{-0.002887in}{0.004970in}}{\pgfqpoint{-0.003928in}{0.003928in}}%
\pgfpathcurveto{\pgfqpoint{-0.004970in}{0.002887in}}{\pgfqpoint{-0.005556in}{0.001473in}}{\pgfqpoint{-0.005556in}{0.000000in}}%
\pgfpathcurveto{\pgfqpoint{-0.005556in}{-0.001473in}}{\pgfqpoint{-0.004970in}{-0.002887in}}{\pgfqpoint{-0.003928in}{-0.003928in}}%
\pgfpathcurveto{\pgfqpoint{-0.002887in}{-0.004970in}}{\pgfqpoint{-0.001473in}{-0.005556in}}{\pgfqpoint{0.000000in}{-0.005556in}}%
\pgfpathlineto{\pgfqpoint{0.000000in}{-0.005556in}}%
\pgfpathclose%
\pgfusepath{stroke,fill}%
}%
\begin{pgfscope}%
\pgfsys@transformshift{3.378797in}{5.020736in}%
\pgfsys@useobject{currentmarker}{}%
\end{pgfscope}%
\end{pgfscope}%
\begin{pgfscope}%
\pgfpathrectangle{\pgfqpoint{0.100000in}{2.493063in}}{\pgfqpoint{5.037500in}{3.427208in}}%
\pgfusepath{clip}%
\pgfsetrectcap%
\pgfsetroundjoin%
\pgfsetlinewidth{1.505625pt}%
\definecolor{currentstroke}{rgb}{0.501961,0.501961,0.501961}%
\pgfsetstrokecolor{currentstroke}%
\pgfsetstrokeopacity{0.500000}%
\pgfsetdash{}{0pt}%
\pgfpathmoveto{\pgfqpoint{3.313193in}{4.805863in}}%
\pgfusepath{stroke}%
\end{pgfscope}%
\begin{pgfscope}%
\pgfpathrectangle{\pgfqpoint{0.100000in}{2.493063in}}{\pgfqpoint{5.037500in}{3.427208in}}%
\pgfusepath{clip}%
\pgfsetbuttcap%
\pgfsetroundjoin%
\definecolor{currentfill}{rgb}{0.501961,0.501961,0.501961}%
\pgfsetfillcolor{currentfill}%
\pgfsetfillopacity{0.500000}%
\pgfsetlinewidth{0.250937pt}%
\definecolor{currentstroke}{rgb}{0.000000,0.000000,0.000000}%
\pgfsetstrokecolor{currentstroke}%
\pgfsetstrokeopacity{0.500000}%
\pgfsetdash{}{0pt}%
\pgfsys@defobject{currentmarker}{\pgfqpoint{-0.013889in}{-0.013889in}}{\pgfqpoint{0.013889in}{0.013889in}}{%
\pgfpathmoveto{\pgfqpoint{0.000000in}{-0.013889in}}%
\pgfpathcurveto{\pgfqpoint{0.003683in}{-0.013889in}}{\pgfqpoint{0.007216in}{-0.012425in}}{\pgfqpoint{0.009821in}{-0.009821in}}%
\pgfpathcurveto{\pgfqpoint{0.012425in}{-0.007216in}}{\pgfqpoint{0.013889in}{-0.003683in}}{\pgfqpoint{0.013889in}{0.000000in}}%
\pgfpathcurveto{\pgfqpoint{0.013889in}{0.003683in}}{\pgfqpoint{0.012425in}{0.007216in}}{\pgfqpoint{0.009821in}{0.009821in}}%
\pgfpathcurveto{\pgfqpoint{0.007216in}{0.012425in}}{\pgfqpoint{0.003683in}{0.013889in}}{\pgfqpoint{0.000000in}{0.013889in}}%
\pgfpathcurveto{\pgfqpoint{-0.003683in}{0.013889in}}{\pgfqpoint{-0.007216in}{0.012425in}}{\pgfqpoint{-0.009821in}{0.009821in}}%
\pgfpathcurveto{\pgfqpoint{-0.012425in}{0.007216in}}{\pgfqpoint{-0.013889in}{0.003683in}}{\pgfqpoint{-0.013889in}{0.000000in}}%
\pgfpathcurveto{\pgfqpoint{-0.013889in}{-0.003683in}}{\pgfqpoint{-0.012425in}{-0.007216in}}{\pgfqpoint{-0.009821in}{-0.009821in}}%
\pgfpathcurveto{\pgfqpoint{-0.007216in}{-0.012425in}}{\pgfqpoint{-0.003683in}{-0.013889in}}{\pgfqpoint{0.000000in}{-0.013889in}}%
\pgfpathlineto{\pgfqpoint{0.000000in}{-0.013889in}}%
\pgfpathclose%
\pgfusepath{stroke,fill}%
}%
\begin{pgfscope}%
\pgfsys@transformshift{3.313193in}{4.805863in}%
\pgfsys@useobject{currentmarker}{}%
\end{pgfscope}%
\end{pgfscope}%
\begin{pgfscope}%
\pgfpathrectangle{\pgfqpoint{0.100000in}{2.493063in}}{\pgfqpoint{5.037500in}{3.427208in}}%
\pgfusepath{clip}%
\pgfsetrectcap%
\pgfsetroundjoin%
\pgfsetlinewidth{1.505625pt}%
\definecolor{currentstroke}{rgb}{0.501961,0.501961,0.501961}%
\pgfsetstrokecolor{currentstroke}%
\pgfsetstrokeopacity{0.500000}%
\pgfsetdash{}{0pt}%
\pgfpathmoveto{\pgfqpoint{3.117641in}{4.920924in}}%
\pgfusepath{stroke}%
\end{pgfscope}%
\begin{pgfscope}%
\pgfpathrectangle{\pgfqpoint{0.100000in}{2.493063in}}{\pgfqpoint{5.037500in}{3.427208in}}%
\pgfusepath{clip}%
\pgfsetbuttcap%
\pgfsetroundjoin%
\definecolor{currentfill}{rgb}{0.501961,0.501961,0.501961}%
\pgfsetfillcolor{currentfill}%
\pgfsetfillopacity{0.500000}%
\pgfsetlinewidth{0.250937pt}%
\definecolor{currentstroke}{rgb}{0.000000,0.000000,0.000000}%
\pgfsetstrokecolor{currentstroke}%
\pgfsetstrokeopacity{0.500000}%
\pgfsetdash{}{0pt}%
\pgfsys@defobject{currentmarker}{\pgfqpoint{-0.013889in}{-0.013889in}}{\pgfqpoint{0.013889in}{0.013889in}}{%
\pgfpathmoveto{\pgfqpoint{0.000000in}{-0.013889in}}%
\pgfpathcurveto{\pgfqpoint{0.003683in}{-0.013889in}}{\pgfqpoint{0.007216in}{-0.012425in}}{\pgfqpoint{0.009821in}{-0.009821in}}%
\pgfpathcurveto{\pgfqpoint{0.012425in}{-0.007216in}}{\pgfqpoint{0.013889in}{-0.003683in}}{\pgfqpoint{0.013889in}{0.000000in}}%
\pgfpathcurveto{\pgfqpoint{0.013889in}{0.003683in}}{\pgfqpoint{0.012425in}{0.007216in}}{\pgfqpoint{0.009821in}{0.009821in}}%
\pgfpathcurveto{\pgfqpoint{0.007216in}{0.012425in}}{\pgfqpoint{0.003683in}{0.013889in}}{\pgfqpoint{0.000000in}{0.013889in}}%
\pgfpathcurveto{\pgfqpoint{-0.003683in}{0.013889in}}{\pgfqpoint{-0.007216in}{0.012425in}}{\pgfqpoint{-0.009821in}{0.009821in}}%
\pgfpathcurveto{\pgfqpoint{-0.012425in}{0.007216in}}{\pgfqpoint{-0.013889in}{0.003683in}}{\pgfqpoint{-0.013889in}{0.000000in}}%
\pgfpathcurveto{\pgfqpoint{-0.013889in}{-0.003683in}}{\pgfqpoint{-0.012425in}{-0.007216in}}{\pgfqpoint{-0.009821in}{-0.009821in}}%
\pgfpathcurveto{\pgfqpoint{-0.007216in}{-0.012425in}}{\pgfqpoint{-0.003683in}{-0.013889in}}{\pgfqpoint{0.000000in}{-0.013889in}}%
\pgfpathlineto{\pgfqpoint{0.000000in}{-0.013889in}}%
\pgfpathclose%
\pgfusepath{stroke,fill}%
}%
\begin{pgfscope}%
\pgfsys@transformshift{3.117641in}{4.920924in}%
\pgfsys@useobject{currentmarker}{}%
\end{pgfscope}%
\end{pgfscope}%
\begin{pgfscope}%
\pgfpathrectangle{\pgfqpoint{0.100000in}{2.493063in}}{\pgfqpoint{5.037500in}{3.427208in}}%
\pgfusepath{clip}%
\pgfsetrectcap%
\pgfsetroundjoin%
\pgfsetlinewidth{1.505625pt}%
\definecolor{currentstroke}{rgb}{0.501961,0.501961,0.501961}%
\pgfsetstrokecolor{currentstroke}%
\pgfsetstrokeopacity{0.500000}%
\pgfsetdash{}{0pt}%
\pgfpathmoveto{\pgfqpoint{3.277383in}{4.845152in}}%
\pgfusepath{stroke}%
\end{pgfscope}%
\begin{pgfscope}%
\pgfpathrectangle{\pgfqpoint{0.100000in}{2.493063in}}{\pgfqpoint{5.037500in}{3.427208in}}%
\pgfusepath{clip}%
\pgfsetbuttcap%
\pgfsetroundjoin%
\definecolor{currentfill}{rgb}{0.501961,0.501961,0.501961}%
\pgfsetfillcolor{currentfill}%
\pgfsetfillopacity{0.500000}%
\pgfsetlinewidth{0.250937pt}%
\definecolor{currentstroke}{rgb}{0.000000,0.000000,0.000000}%
\pgfsetstrokecolor{currentstroke}%
\pgfsetstrokeopacity{0.500000}%
\pgfsetdash{}{0pt}%
\pgfsys@defobject{currentmarker}{\pgfqpoint{-0.013889in}{-0.013889in}}{\pgfqpoint{0.013889in}{0.013889in}}{%
\pgfpathmoveto{\pgfqpoint{0.000000in}{-0.013889in}}%
\pgfpathcurveto{\pgfqpoint{0.003683in}{-0.013889in}}{\pgfqpoint{0.007216in}{-0.012425in}}{\pgfqpoint{0.009821in}{-0.009821in}}%
\pgfpathcurveto{\pgfqpoint{0.012425in}{-0.007216in}}{\pgfqpoint{0.013889in}{-0.003683in}}{\pgfqpoint{0.013889in}{0.000000in}}%
\pgfpathcurveto{\pgfqpoint{0.013889in}{0.003683in}}{\pgfqpoint{0.012425in}{0.007216in}}{\pgfqpoint{0.009821in}{0.009821in}}%
\pgfpathcurveto{\pgfqpoint{0.007216in}{0.012425in}}{\pgfqpoint{0.003683in}{0.013889in}}{\pgfqpoint{0.000000in}{0.013889in}}%
\pgfpathcurveto{\pgfqpoint{-0.003683in}{0.013889in}}{\pgfqpoint{-0.007216in}{0.012425in}}{\pgfqpoint{-0.009821in}{0.009821in}}%
\pgfpathcurveto{\pgfqpoint{-0.012425in}{0.007216in}}{\pgfqpoint{-0.013889in}{0.003683in}}{\pgfqpoint{-0.013889in}{0.000000in}}%
\pgfpathcurveto{\pgfqpoint{-0.013889in}{-0.003683in}}{\pgfqpoint{-0.012425in}{-0.007216in}}{\pgfqpoint{-0.009821in}{-0.009821in}}%
\pgfpathcurveto{\pgfqpoint{-0.007216in}{-0.012425in}}{\pgfqpoint{-0.003683in}{-0.013889in}}{\pgfqpoint{0.000000in}{-0.013889in}}%
\pgfpathlineto{\pgfqpoint{0.000000in}{-0.013889in}}%
\pgfpathclose%
\pgfusepath{stroke,fill}%
}%
\begin{pgfscope}%
\pgfsys@transformshift{3.277383in}{4.845152in}%
\pgfsys@useobject{currentmarker}{}%
\end{pgfscope}%
\end{pgfscope}%
\begin{pgfscope}%
\pgfpathrectangle{\pgfqpoint{0.100000in}{2.493063in}}{\pgfqpoint{5.037500in}{3.427208in}}%
\pgfusepath{clip}%
\pgfsetrectcap%
\pgfsetroundjoin%
\pgfsetlinewidth{1.505625pt}%
\definecolor{currentstroke}{rgb}{0.501961,0.501961,0.501961}%
\pgfsetstrokecolor{currentstroke}%
\pgfsetstrokeopacity{0.500000}%
\pgfsetdash{}{0pt}%
\pgfpathmoveto{\pgfqpoint{3.400651in}{4.849173in}}%
\pgfusepath{stroke}%
\end{pgfscope}%
\begin{pgfscope}%
\pgfpathrectangle{\pgfqpoint{0.100000in}{2.493063in}}{\pgfqpoint{5.037500in}{3.427208in}}%
\pgfusepath{clip}%
\pgfsetbuttcap%
\pgfsetroundjoin%
\definecolor{currentfill}{rgb}{0.501961,0.501961,0.501961}%
\pgfsetfillcolor{currentfill}%
\pgfsetfillopacity{0.500000}%
\pgfsetlinewidth{0.250937pt}%
\definecolor{currentstroke}{rgb}{0.000000,0.000000,0.000000}%
\pgfsetstrokecolor{currentstroke}%
\pgfsetstrokeopacity{0.500000}%
\pgfsetdash{}{0pt}%
\pgfsys@defobject{currentmarker}{\pgfqpoint{-0.013889in}{-0.013889in}}{\pgfqpoint{0.013889in}{0.013889in}}{%
\pgfpathmoveto{\pgfqpoint{0.000000in}{-0.013889in}}%
\pgfpathcurveto{\pgfqpoint{0.003683in}{-0.013889in}}{\pgfqpoint{0.007216in}{-0.012425in}}{\pgfqpoint{0.009821in}{-0.009821in}}%
\pgfpathcurveto{\pgfqpoint{0.012425in}{-0.007216in}}{\pgfqpoint{0.013889in}{-0.003683in}}{\pgfqpoint{0.013889in}{0.000000in}}%
\pgfpathcurveto{\pgfqpoint{0.013889in}{0.003683in}}{\pgfqpoint{0.012425in}{0.007216in}}{\pgfqpoint{0.009821in}{0.009821in}}%
\pgfpathcurveto{\pgfqpoint{0.007216in}{0.012425in}}{\pgfqpoint{0.003683in}{0.013889in}}{\pgfqpoint{0.000000in}{0.013889in}}%
\pgfpathcurveto{\pgfqpoint{-0.003683in}{0.013889in}}{\pgfqpoint{-0.007216in}{0.012425in}}{\pgfqpoint{-0.009821in}{0.009821in}}%
\pgfpathcurveto{\pgfqpoint{-0.012425in}{0.007216in}}{\pgfqpoint{-0.013889in}{0.003683in}}{\pgfqpoint{-0.013889in}{0.000000in}}%
\pgfpathcurveto{\pgfqpoint{-0.013889in}{-0.003683in}}{\pgfqpoint{-0.012425in}{-0.007216in}}{\pgfqpoint{-0.009821in}{-0.009821in}}%
\pgfpathcurveto{\pgfqpoint{-0.007216in}{-0.012425in}}{\pgfqpoint{-0.003683in}{-0.013889in}}{\pgfqpoint{0.000000in}{-0.013889in}}%
\pgfpathlineto{\pgfqpoint{0.000000in}{-0.013889in}}%
\pgfpathclose%
\pgfusepath{stroke,fill}%
}%
\begin{pgfscope}%
\pgfsys@transformshift{3.400651in}{4.849173in}%
\pgfsys@useobject{currentmarker}{}%
\end{pgfscope}%
\end{pgfscope}%
\begin{pgfscope}%
\pgfpathrectangle{\pgfqpoint{0.100000in}{2.493063in}}{\pgfqpoint{5.037500in}{3.427208in}}%
\pgfusepath{clip}%
\pgfsetrectcap%
\pgfsetroundjoin%
\pgfsetlinewidth{1.505625pt}%
\definecolor{currentstroke}{rgb}{0.501961,0.501961,0.501961}%
\pgfsetstrokecolor{currentstroke}%
\pgfsetstrokeopacity{0.500000}%
\pgfsetdash{}{0pt}%
\pgfpathmoveto{\pgfqpoint{3.340539in}{4.958855in}}%
\pgfusepath{stroke}%
\end{pgfscope}%
\begin{pgfscope}%
\pgfpathrectangle{\pgfqpoint{0.100000in}{2.493063in}}{\pgfqpoint{5.037500in}{3.427208in}}%
\pgfusepath{clip}%
\pgfsetbuttcap%
\pgfsetroundjoin%
\definecolor{currentfill}{rgb}{0.501961,0.501961,0.501961}%
\pgfsetfillcolor{currentfill}%
\pgfsetfillopacity{0.500000}%
\pgfsetlinewidth{0.250937pt}%
\definecolor{currentstroke}{rgb}{0.000000,0.000000,0.000000}%
\pgfsetstrokecolor{currentstroke}%
\pgfsetstrokeopacity{0.500000}%
\pgfsetdash{}{0pt}%
\pgfsys@defobject{currentmarker}{\pgfqpoint{-0.013889in}{-0.013889in}}{\pgfqpoint{0.013889in}{0.013889in}}{%
\pgfpathmoveto{\pgfqpoint{0.000000in}{-0.013889in}}%
\pgfpathcurveto{\pgfqpoint{0.003683in}{-0.013889in}}{\pgfqpoint{0.007216in}{-0.012425in}}{\pgfqpoint{0.009821in}{-0.009821in}}%
\pgfpathcurveto{\pgfqpoint{0.012425in}{-0.007216in}}{\pgfqpoint{0.013889in}{-0.003683in}}{\pgfqpoint{0.013889in}{0.000000in}}%
\pgfpathcurveto{\pgfqpoint{0.013889in}{0.003683in}}{\pgfqpoint{0.012425in}{0.007216in}}{\pgfqpoint{0.009821in}{0.009821in}}%
\pgfpathcurveto{\pgfqpoint{0.007216in}{0.012425in}}{\pgfqpoint{0.003683in}{0.013889in}}{\pgfqpoint{0.000000in}{0.013889in}}%
\pgfpathcurveto{\pgfqpoint{-0.003683in}{0.013889in}}{\pgfqpoint{-0.007216in}{0.012425in}}{\pgfqpoint{-0.009821in}{0.009821in}}%
\pgfpathcurveto{\pgfqpoint{-0.012425in}{0.007216in}}{\pgfqpoint{-0.013889in}{0.003683in}}{\pgfqpoint{-0.013889in}{0.000000in}}%
\pgfpathcurveto{\pgfqpoint{-0.013889in}{-0.003683in}}{\pgfqpoint{-0.012425in}{-0.007216in}}{\pgfqpoint{-0.009821in}{-0.009821in}}%
\pgfpathcurveto{\pgfqpoint{-0.007216in}{-0.012425in}}{\pgfqpoint{-0.003683in}{-0.013889in}}{\pgfqpoint{0.000000in}{-0.013889in}}%
\pgfpathlineto{\pgfqpoint{0.000000in}{-0.013889in}}%
\pgfpathclose%
\pgfusepath{stroke,fill}%
}%
\begin{pgfscope}%
\pgfsys@transformshift{3.340539in}{4.958855in}%
\pgfsys@useobject{currentmarker}{}%
\end{pgfscope}%
\end{pgfscope}%
\begin{pgfscope}%
\pgfpathrectangle{\pgfqpoint{0.100000in}{2.493063in}}{\pgfqpoint{5.037500in}{3.427208in}}%
\pgfusepath{clip}%
\pgfsetrectcap%
\pgfsetroundjoin%
\pgfsetlinewidth{1.505625pt}%
\definecolor{currentstroke}{rgb}{0.678431,1.000000,0.184314}%
\pgfsetstrokecolor{currentstroke}%
\pgfsetstrokeopacity{0.500000}%
\pgfsetdash{}{0pt}%
\pgfpathmoveto{\pgfqpoint{3.415252in}{4.814560in}}%
\pgfusepath{stroke}%
\end{pgfscope}%
\begin{pgfscope}%
\pgfpathrectangle{\pgfqpoint{0.100000in}{2.493063in}}{\pgfqpoint{5.037500in}{3.427208in}}%
\pgfusepath{clip}%
\pgfsetbuttcap%
\pgfsetroundjoin%
\definecolor{currentfill}{rgb}{0.678431,1.000000,0.184314}%
\pgfsetfillcolor{currentfill}%
\pgfsetfillopacity{0.500000}%
\pgfsetlinewidth{0.250937pt}%
\definecolor{currentstroke}{rgb}{0.000000,0.000000,0.000000}%
\pgfsetstrokecolor{currentstroke}%
\pgfsetstrokeopacity{0.500000}%
\pgfsetdash{}{0pt}%
\pgfsys@defobject{currentmarker}{\pgfqpoint{-0.008333in}{-0.008333in}}{\pgfqpoint{0.008333in}{0.008333in}}{%
\pgfpathmoveto{\pgfqpoint{0.000000in}{-0.008333in}}%
\pgfpathcurveto{\pgfqpoint{0.002210in}{-0.008333in}}{\pgfqpoint{0.004330in}{-0.007455in}}{\pgfqpoint{0.005893in}{-0.005893in}}%
\pgfpathcurveto{\pgfqpoint{0.007455in}{-0.004330in}}{\pgfqpoint{0.008333in}{-0.002210in}}{\pgfqpoint{0.008333in}{0.000000in}}%
\pgfpathcurveto{\pgfqpoint{0.008333in}{0.002210in}}{\pgfqpoint{0.007455in}{0.004330in}}{\pgfqpoint{0.005893in}{0.005893in}}%
\pgfpathcurveto{\pgfqpoint{0.004330in}{0.007455in}}{\pgfqpoint{0.002210in}{0.008333in}}{\pgfqpoint{0.000000in}{0.008333in}}%
\pgfpathcurveto{\pgfqpoint{-0.002210in}{0.008333in}}{\pgfqpoint{-0.004330in}{0.007455in}}{\pgfqpoint{-0.005893in}{0.005893in}}%
\pgfpathcurveto{\pgfqpoint{-0.007455in}{0.004330in}}{\pgfqpoint{-0.008333in}{0.002210in}}{\pgfqpoint{-0.008333in}{0.000000in}}%
\pgfpathcurveto{\pgfqpoint{-0.008333in}{-0.002210in}}{\pgfqpoint{-0.007455in}{-0.004330in}}{\pgfqpoint{-0.005893in}{-0.005893in}}%
\pgfpathcurveto{\pgfqpoint{-0.004330in}{-0.007455in}}{\pgfqpoint{-0.002210in}{-0.008333in}}{\pgfqpoint{0.000000in}{-0.008333in}}%
\pgfpathlineto{\pgfqpoint{0.000000in}{-0.008333in}}%
\pgfpathclose%
\pgfusepath{stroke,fill}%
}%
\begin{pgfscope}%
\pgfsys@transformshift{3.415252in}{4.814560in}%
\pgfsys@useobject{currentmarker}{}%
\end{pgfscope}%
\end{pgfscope}%
\begin{pgfscope}%
\pgfpathrectangle{\pgfqpoint{0.100000in}{2.493063in}}{\pgfqpoint{5.037500in}{3.427208in}}%
\pgfusepath{clip}%
\pgfsetrectcap%
\pgfsetroundjoin%
\pgfsetlinewidth{1.505625pt}%
\definecolor{currentstroke}{rgb}{0.678431,1.000000,0.184314}%
\pgfsetstrokecolor{currentstroke}%
\pgfsetstrokeopacity{0.500000}%
\pgfsetdash{}{0pt}%
\pgfpathmoveto{\pgfqpoint{3.411522in}{4.932958in}}%
\pgfusepath{stroke}%
\end{pgfscope}%
\begin{pgfscope}%
\pgfpathrectangle{\pgfqpoint{0.100000in}{2.493063in}}{\pgfqpoint{5.037500in}{3.427208in}}%
\pgfusepath{clip}%
\pgfsetbuttcap%
\pgfsetroundjoin%
\definecolor{currentfill}{rgb}{0.678431,1.000000,0.184314}%
\pgfsetfillcolor{currentfill}%
\pgfsetfillopacity{0.500000}%
\pgfsetlinewidth{0.250937pt}%
\definecolor{currentstroke}{rgb}{0.000000,0.000000,0.000000}%
\pgfsetstrokecolor{currentstroke}%
\pgfsetstrokeopacity{0.500000}%
\pgfsetdash{}{0pt}%
\pgfsys@defobject{currentmarker}{\pgfqpoint{-0.005556in}{-0.005556in}}{\pgfqpoint{0.005556in}{0.005556in}}{%
\pgfpathmoveto{\pgfqpoint{0.000000in}{-0.005556in}}%
\pgfpathcurveto{\pgfqpoint{0.001473in}{-0.005556in}}{\pgfqpoint{0.002887in}{-0.004970in}}{\pgfqpoint{0.003928in}{-0.003928in}}%
\pgfpathcurveto{\pgfqpoint{0.004970in}{-0.002887in}}{\pgfqpoint{0.005556in}{-0.001473in}}{\pgfqpoint{0.005556in}{0.000000in}}%
\pgfpathcurveto{\pgfqpoint{0.005556in}{0.001473in}}{\pgfqpoint{0.004970in}{0.002887in}}{\pgfqpoint{0.003928in}{0.003928in}}%
\pgfpathcurveto{\pgfqpoint{0.002887in}{0.004970in}}{\pgfqpoint{0.001473in}{0.005556in}}{\pgfqpoint{0.000000in}{0.005556in}}%
\pgfpathcurveto{\pgfqpoint{-0.001473in}{0.005556in}}{\pgfqpoint{-0.002887in}{0.004970in}}{\pgfqpoint{-0.003928in}{0.003928in}}%
\pgfpathcurveto{\pgfqpoint{-0.004970in}{0.002887in}}{\pgfqpoint{-0.005556in}{0.001473in}}{\pgfqpoint{-0.005556in}{0.000000in}}%
\pgfpathcurveto{\pgfqpoint{-0.005556in}{-0.001473in}}{\pgfqpoint{-0.004970in}{-0.002887in}}{\pgfqpoint{-0.003928in}{-0.003928in}}%
\pgfpathcurveto{\pgfqpoint{-0.002887in}{-0.004970in}}{\pgfqpoint{-0.001473in}{-0.005556in}}{\pgfqpoint{0.000000in}{-0.005556in}}%
\pgfpathlineto{\pgfqpoint{0.000000in}{-0.005556in}}%
\pgfpathclose%
\pgfusepath{stroke,fill}%
}%
\begin{pgfscope}%
\pgfsys@transformshift{3.411522in}{4.932958in}%
\pgfsys@useobject{currentmarker}{}%
\end{pgfscope}%
\end{pgfscope}%
\begin{pgfscope}%
\pgfpathrectangle{\pgfqpoint{0.100000in}{2.493063in}}{\pgfqpoint{5.037500in}{3.427208in}}%
\pgfusepath{clip}%
\pgfsetrectcap%
\pgfsetroundjoin%
\pgfsetlinewidth{1.505625pt}%
\definecolor{currentstroke}{rgb}{0.501961,0.501961,0.501961}%
\pgfsetstrokecolor{currentstroke}%
\pgfsetstrokeopacity{0.500000}%
\pgfsetdash{}{0pt}%
\pgfpathmoveto{\pgfqpoint{3.243806in}{5.061389in}}%
\pgfusepath{stroke}%
\end{pgfscope}%
\begin{pgfscope}%
\pgfpathrectangle{\pgfqpoint{0.100000in}{2.493063in}}{\pgfqpoint{5.037500in}{3.427208in}}%
\pgfusepath{clip}%
\pgfsetbuttcap%
\pgfsetroundjoin%
\definecolor{currentfill}{rgb}{0.501961,0.501961,0.501961}%
\pgfsetfillcolor{currentfill}%
\pgfsetfillopacity{0.500000}%
\pgfsetlinewidth{0.250937pt}%
\definecolor{currentstroke}{rgb}{0.000000,0.000000,0.000000}%
\pgfsetstrokecolor{currentstroke}%
\pgfsetstrokeopacity{0.500000}%
\pgfsetdash{}{0pt}%
\pgfsys@defobject{currentmarker}{\pgfqpoint{-0.013889in}{-0.013889in}}{\pgfqpoint{0.013889in}{0.013889in}}{%
\pgfpathmoveto{\pgfqpoint{0.000000in}{-0.013889in}}%
\pgfpathcurveto{\pgfqpoint{0.003683in}{-0.013889in}}{\pgfqpoint{0.007216in}{-0.012425in}}{\pgfqpoint{0.009821in}{-0.009821in}}%
\pgfpathcurveto{\pgfqpoint{0.012425in}{-0.007216in}}{\pgfqpoint{0.013889in}{-0.003683in}}{\pgfqpoint{0.013889in}{0.000000in}}%
\pgfpathcurveto{\pgfqpoint{0.013889in}{0.003683in}}{\pgfqpoint{0.012425in}{0.007216in}}{\pgfqpoint{0.009821in}{0.009821in}}%
\pgfpathcurveto{\pgfqpoint{0.007216in}{0.012425in}}{\pgfqpoint{0.003683in}{0.013889in}}{\pgfqpoint{0.000000in}{0.013889in}}%
\pgfpathcurveto{\pgfqpoint{-0.003683in}{0.013889in}}{\pgfqpoint{-0.007216in}{0.012425in}}{\pgfqpoint{-0.009821in}{0.009821in}}%
\pgfpathcurveto{\pgfqpoint{-0.012425in}{0.007216in}}{\pgfqpoint{-0.013889in}{0.003683in}}{\pgfqpoint{-0.013889in}{0.000000in}}%
\pgfpathcurveto{\pgfqpoint{-0.013889in}{-0.003683in}}{\pgfqpoint{-0.012425in}{-0.007216in}}{\pgfqpoint{-0.009821in}{-0.009821in}}%
\pgfpathcurveto{\pgfqpoint{-0.007216in}{-0.012425in}}{\pgfqpoint{-0.003683in}{-0.013889in}}{\pgfqpoint{0.000000in}{-0.013889in}}%
\pgfpathlineto{\pgfqpoint{0.000000in}{-0.013889in}}%
\pgfpathclose%
\pgfusepath{stroke,fill}%
}%
\begin{pgfscope}%
\pgfsys@transformshift{3.243806in}{5.061389in}%
\pgfsys@useobject{currentmarker}{}%
\end{pgfscope}%
\end{pgfscope}%
\begin{pgfscope}%
\pgfpathrectangle{\pgfqpoint{0.100000in}{2.493063in}}{\pgfqpoint{5.037500in}{3.427208in}}%
\pgfusepath{clip}%
\pgfsetrectcap%
\pgfsetroundjoin%
\pgfsetlinewidth{1.505625pt}%
\definecolor{currentstroke}{rgb}{0.678431,1.000000,0.184314}%
\pgfsetstrokecolor{currentstroke}%
\pgfsetstrokeopacity{0.500000}%
\pgfsetdash{}{0pt}%
\pgfpathmoveto{\pgfqpoint{1.849699in}{4.865986in}}%
\pgfusepath{stroke}%
\end{pgfscope}%
\begin{pgfscope}%
\pgfpathrectangle{\pgfqpoint{0.100000in}{2.493063in}}{\pgfqpoint{5.037500in}{3.427208in}}%
\pgfusepath{clip}%
\pgfsetbuttcap%
\pgfsetroundjoin%
\definecolor{currentfill}{rgb}{0.678431,1.000000,0.184314}%
\pgfsetfillcolor{currentfill}%
\pgfsetfillopacity{0.500000}%
\pgfsetlinewidth{0.250937pt}%
\definecolor{currentstroke}{rgb}{0.000000,0.000000,0.000000}%
\pgfsetstrokecolor{currentstroke}%
\pgfsetstrokeopacity{0.500000}%
\pgfsetdash{}{0pt}%
\pgfsys@defobject{currentmarker}{\pgfqpoint{-0.025000in}{-0.025000in}}{\pgfqpoint{0.025000in}{0.025000in}}{%
\pgfpathmoveto{\pgfqpoint{0.000000in}{-0.025000in}}%
\pgfpathcurveto{\pgfqpoint{0.006630in}{-0.025000in}}{\pgfqpoint{0.012989in}{-0.022366in}}{\pgfqpoint{0.017678in}{-0.017678in}}%
\pgfpathcurveto{\pgfqpoint{0.022366in}{-0.012989in}}{\pgfqpoint{0.025000in}{-0.006630in}}{\pgfqpoint{0.025000in}{0.000000in}}%
\pgfpathcurveto{\pgfqpoint{0.025000in}{0.006630in}}{\pgfqpoint{0.022366in}{0.012989in}}{\pgfqpoint{0.017678in}{0.017678in}}%
\pgfpathcurveto{\pgfqpoint{0.012989in}{0.022366in}}{\pgfqpoint{0.006630in}{0.025000in}}{\pgfqpoint{0.000000in}{0.025000in}}%
\pgfpathcurveto{\pgfqpoint{-0.006630in}{0.025000in}}{\pgfqpoint{-0.012989in}{0.022366in}}{\pgfqpoint{-0.017678in}{0.017678in}}%
\pgfpathcurveto{\pgfqpoint{-0.022366in}{0.012989in}}{\pgfqpoint{-0.025000in}{0.006630in}}{\pgfqpoint{-0.025000in}{0.000000in}}%
\pgfpathcurveto{\pgfqpoint{-0.025000in}{-0.006630in}}{\pgfqpoint{-0.022366in}{-0.012989in}}{\pgfqpoint{-0.017678in}{-0.017678in}}%
\pgfpathcurveto{\pgfqpoint{-0.012989in}{-0.022366in}}{\pgfqpoint{-0.006630in}{-0.025000in}}{\pgfqpoint{0.000000in}{-0.025000in}}%
\pgfpathlineto{\pgfqpoint{0.000000in}{-0.025000in}}%
\pgfpathclose%
\pgfusepath{stroke,fill}%
}%
\begin{pgfscope}%
\pgfsys@transformshift{1.849699in}{4.865986in}%
\pgfsys@useobject{currentmarker}{}%
\end{pgfscope}%
\end{pgfscope}%
\begin{pgfscope}%
\pgfpathrectangle{\pgfqpoint{0.100000in}{2.493063in}}{\pgfqpoint{5.037500in}{3.427208in}}%
\pgfusepath{clip}%
\pgfsetrectcap%
\pgfsetroundjoin%
\pgfsetlinewidth{1.505625pt}%
\definecolor{currentstroke}{rgb}{0.678431,1.000000,0.184314}%
\pgfsetstrokecolor{currentstroke}%
\pgfsetstrokeopacity{0.500000}%
\pgfsetdash{}{0pt}%
\pgfpathmoveto{\pgfqpoint{1.953642in}{4.653410in}}%
\pgfusepath{stroke}%
\end{pgfscope}%
\begin{pgfscope}%
\pgfpathrectangle{\pgfqpoint{0.100000in}{2.493063in}}{\pgfqpoint{5.037500in}{3.427208in}}%
\pgfusepath{clip}%
\pgfsetbuttcap%
\pgfsetroundjoin%
\definecolor{currentfill}{rgb}{0.678431,1.000000,0.184314}%
\pgfsetfillcolor{currentfill}%
\pgfsetfillopacity{0.500000}%
\pgfsetlinewidth{0.250937pt}%
\definecolor{currentstroke}{rgb}{0.000000,0.000000,0.000000}%
\pgfsetstrokecolor{currentstroke}%
\pgfsetstrokeopacity{0.500000}%
\pgfsetdash{}{0pt}%
\pgfsys@defobject{currentmarker}{\pgfqpoint{-0.011111in}{-0.011111in}}{\pgfqpoint{0.011111in}{0.011111in}}{%
\pgfpathmoveto{\pgfqpoint{0.000000in}{-0.011111in}}%
\pgfpathcurveto{\pgfqpoint{0.002947in}{-0.011111in}}{\pgfqpoint{0.005773in}{-0.009940in}}{\pgfqpoint{0.007857in}{-0.007857in}}%
\pgfpathcurveto{\pgfqpoint{0.009940in}{-0.005773in}}{\pgfqpoint{0.011111in}{-0.002947in}}{\pgfqpoint{0.011111in}{0.000000in}}%
\pgfpathcurveto{\pgfqpoint{0.011111in}{0.002947in}}{\pgfqpoint{0.009940in}{0.005773in}}{\pgfqpoint{0.007857in}{0.007857in}}%
\pgfpathcurveto{\pgfqpoint{0.005773in}{0.009940in}}{\pgfqpoint{0.002947in}{0.011111in}}{\pgfqpoint{0.000000in}{0.011111in}}%
\pgfpathcurveto{\pgfqpoint{-0.002947in}{0.011111in}}{\pgfqpoint{-0.005773in}{0.009940in}}{\pgfqpoint{-0.007857in}{0.007857in}}%
\pgfpathcurveto{\pgfqpoint{-0.009940in}{0.005773in}}{\pgfqpoint{-0.011111in}{0.002947in}}{\pgfqpoint{-0.011111in}{0.000000in}}%
\pgfpathcurveto{\pgfqpoint{-0.011111in}{-0.002947in}}{\pgfqpoint{-0.009940in}{-0.005773in}}{\pgfqpoint{-0.007857in}{-0.007857in}}%
\pgfpathcurveto{\pgfqpoint{-0.005773in}{-0.009940in}}{\pgfqpoint{-0.002947in}{-0.011111in}}{\pgfqpoint{0.000000in}{-0.011111in}}%
\pgfpathlineto{\pgfqpoint{0.000000in}{-0.011111in}}%
\pgfpathclose%
\pgfusepath{stroke,fill}%
}%
\begin{pgfscope}%
\pgfsys@transformshift{1.953642in}{4.653410in}%
\pgfsys@useobject{currentmarker}{}%
\end{pgfscope}%
\end{pgfscope}%
\begin{pgfscope}%
\pgfpathrectangle{\pgfqpoint{0.100000in}{2.493063in}}{\pgfqpoint{5.037500in}{3.427208in}}%
\pgfusepath{clip}%
\pgfsetrectcap%
\pgfsetroundjoin%
\pgfsetlinewidth{1.505625pt}%
\definecolor{currentstroke}{rgb}{0.501961,0.501961,0.501961}%
\pgfsetstrokecolor{currentstroke}%
\pgfsetstrokeopacity{0.500000}%
\pgfsetdash{}{0pt}%
\pgfpathmoveto{\pgfqpoint{2.769875in}{5.748910in}}%
\pgfusepath{stroke}%
\end{pgfscope}%
\begin{pgfscope}%
\pgfpathrectangle{\pgfqpoint{0.100000in}{2.493063in}}{\pgfqpoint{5.037500in}{3.427208in}}%
\pgfusepath{clip}%
\pgfsetbuttcap%
\pgfsetroundjoin%
\definecolor{currentfill}{rgb}{0.501961,0.501961,0.501961}%
\pgfsetfillcolor{currentfill}%
\pgfsetfillopacity{0.500000}%
\pgfsetlinewidth{0.250937pt}%
\definecolor{currentstroke}{rgb}{0.000000,0.000000,0.000000}%
\pgfsetstrokecolor{currentstroke}%
\pgfsetstrokeopacity{0.500000}%
\pgfsetdash{}{0pt}%
\pgfsys@defobject{currentmarker}{\pgfqpoint{-0.013889in}{-0.013889in}}{\pgfqpoint{0.013889in}{0.013889in}}{%
\pgfpathmoveto{\pgfqpoint{0.000000in}{-0.013889in}}%
\pgfpathcurveto{\pgfqpoint{0.003683in}{-0.013889in}}{\pgfqpoint{0.007216in}{-0.012425in}}{\pgfqpoint{0.009821in}{-0.009821in}}%
\pgfpathcurveto{\pgfqpoint{0.012425in}{-0.007216in}}{\pgfqpoint{0.013889in}{-0.003683in}}{\pgfqpoint{0.013889in}{0.000000in}}%
\pgfpathcurveto{\pgfqpoint{0.013889in}{0.003683in}}{\pgfqpoint{0.012425in}{0.007216in}}{\pgfqpoint{0.009821in}{0.009821in}}%
\pgfpathcurveto{\pgfqpoint{0.007216in}{0.012425in}}{\pgfqpoint{0.003683in}{0.013889in}}{\pgfqpoint{0.000000in}{0.013889in}}%
\pgfpathcurveto{\pgfqpoint{-0.003683in}{0.013889in}}{\pgfqpoint{-0.007216in}{0.012425in}}{\pgfqpoint{-0.009821in}{0.009821in}}%
\pgfpathcurveto{\pgfqpoint{-0.012425in}{0.007216in}}{\pgfqpoint{-0.013889in}{0.003683in}}{\pgfqpoint{-0.013889in}{0.000000in}}%
\pgfpathcurveto{\pgfqpoint{-0.013889in}{-0.003683in}}{\pgfqpoint{-0.012425in}{-0.007216in}}{\pgfqpoint{-0.009821in}{-0.009821in}}%
\pgfpathcurveto{\pgfqpoint{-0.007216in}{-0.012425in}}{\pgfqpoint{-0.003683in}{-0.013889in}}{\pgfqpoint{0.000000in}{-0.013889in}}%
\pgfpathlineto{\pgfqpoint{0.000000in}{-0.013889in}}%
\pgfpathclose%
\pgfusepath{stroke,fill}%
}%
\begin{pgfscope}%
\pgfsys@transformshift{2.769875in}{5.748910in}%
\pgfsys@useobject{currentmarker}{}%
\end{pgfscope}%
\end{pgfscope}%
\begin{pgfscope}%
\pgfpathrectangle{\pgfqpoint{0.100000in}{2.493063in}}{\pgfqpoint{5.037500in}{3.427208in}}%
\pgfusepath{clip}%
\pgfsetrectcap%
\pgfsetroundjoin%
\pgfsetlinewidth{1.505625pt}%
\definecolor{currentstroke}{rgb}{0.000000,0.000000,1.000000}%
\pgfsetstrokecolor{currentstroke}%
\pgfsetstrokeopacity{0.500000}%
\pgfsetdash{}{0pt}%
\pgfpathmoveto{\pgfqpoint{3.550688in}{5.748910in}}%
\pgfusepath{stroke}%
\end{pgfscope}%
\begin{pgfscope}%
\pgfpathrectangle{\pgfqpoint{0.100000in}{2.493063in}}{\pgfqpoint{5.037500in}{3.427208in}}%
\pgfusepath{clip}%
\pgfsetbuttcap%
\pgfsetroundjoin%
\definecolor{currentfill}{rgb}{0.000000,0.000000,1.000000}%
\pgfsetfillcolor{currentfill}%
\pgfsetfillopacity{0.500000}%
\pgfsetlinewidth{0.250937pt}%
\definecolor{currentstroke}{rgb}{0.000000,0.000000,0.000000}%
\pgfsetstrokecolor{currentstroke}%
\pgfsetstrokeopacity{0.500000}%
\pgfsetdash{}{0pt}%
\pgfsys@defobject{currentmarker}{\pgfqpoint{-0.083333in}{-0.083333in}}{\pgfqpoint{0.083333in}{0.083333in}}{%
\pgfpathmoveto{\pgfqpoint{0.000000in}{-0.083333in}}%
\pgfpathcurveto{\pgfqpoint{0.022100in}{-0.083333in}}{\pgfqpoint{0.043298in}{-0.074553in}}{\pgfqpoint{0.058926in}{-0.058926in}}%
\pgfpathcurveto{\pgfqpoint{0.074553in}{-0.043298in}}{\pgfqpoint{0.083333in}{-0.022100in}}{\pgfqpoint{0.083333in}{0.000000in}}%
\pgfpathcurveto{\pgfqpoint{0.083333in}{0.022100in}}{\pgfqpoint{0.074553in}{0.043298in}}{\pgfqpoint{0.058926in}{0.058926in}}%
\pgfpathcurveto{\pgfqpoint{0.043298in}{0.074553in}}{\pgfqpoint{0.022100in}{0.083333in}}{\pgfqpoint{0.000000in}{0.083333in}}%
\pgfpathcurveto{\pgfqpoint{-0.022100in}{0.083333in}}{\pgfqpoint{-0.043298in}{0.074553in}}{\pgfqpoint{-0.058926in}{0.058926in}}%
\pgfpathcurveto{\pgfqpoint{-0.074553in}{0.043298in}}{\pgfqpoint{-0.083333in}{0.022100in}}{\pgfqpoint{-0.083333in}{0.000000in}}%
\pgfpathcurveto{\pgfqpoint{-0.083333in}{-0.022100in}}{\pgfqpoint{-0.074553in}{-0.043298in}}{\pgfqpoint{-0.058926in}{-0.058926in}}%
\pgfpathcurveto{\pgfqpoint{-0.043298in}{-0.074553in}}{\pgfqpoint{-0.022100in}{-0.083333in}}{\pgfqpoint{0.000000in}{-0.083333in}}%
\pgfpathlineto{\pgfqpoint{0.000000in}{-0.083333in}}%
\pgfpathclose%
\pgfusepath{stroke,fill}%
}%
\begin{pgfscope}%
\pgfsys@transformshift{3.550688in}{5.748910in}%
\pgfsys@useobject{currentmarker}{}%
\end{pgfscope}%
\end{pgfscope}%
\begin{pgfscope}%
\pgfpathrectangle{\pgfqpoint{0.100000in}{2.493063in}}{\pgfqpoint{5.037500in}{3.427208in}}%
\pgfusepath{clip}%
\pgfsetrectcap%
\pgfsetroundjoin%
\pgfsetlinewidth{1.505625pt}%
\definecolor{currentstroke}{rgb}{0.000000,0.000000,1.000000}%
\pgfsetstrokecolor{currentstroke}%
\pgfsetstrokeopacity{0.500000}%
\pgfsetdash{}{0pt}%
\pgfpathmoveto{\pgfqpoint{3.550688in}{5.577550in}}%
\pgfusepath{stroke}%
\end{pgfscope}%
\begin{pgfscope}%
\pgfpathrectangle{\pgfqpoint{0.100000in}{2.493063in}}{\pgfqpoint{5.037500in}{3.427208in}}%
\pgfusepath{clip}%
\pgfsetbuttcap%
\pgfsetroundjoin%
\definecolor{currentfill}{rgb}{0.000000,0.000000,1.000000}%
\pgfsetfillcolor{currentfill}%
\pgfsetfillopacity{0.500000}%
\pgfsetlinewidth{0.250937pt}%
\definecolor{currentstroke}{rgb}{0.000000,0.000000,0.000000}%
\pgfsetstrokecolor{currentstroke}%
\pgfsetstrokeopacity{0.500000}%
\pgfsetdash{}{0pt}%
\pgfsys@defobject{currentmarker}{\pgfqpoint{-0.055556in}{-0.055556in}}{\pgfqpoint{0.055556in}{0.055556in}}{%
\pgfpathmoveto{\pgfqpoint{0.000000in}{-0.055556in}}%
\pgfpathcurveto{\pgfqpoint{0.014734in}{-0.055556in}}{\pgfqpoint{0.028866in}{-0.049702in}}{\pgfqpoint{0.039284in}{-0.039284in}}%
\pgfpathcurveto{\pgfqpoint{0.049702in}{-0.028866in}}{\pgfqpoint{0.055556in}{-0.014734in}}{\pgfqpoint{0.055556in}{0.000000in}}%
\pgfpathcurveto{\pgfqpoint{0.055556in}{0.014734in}}{\pgfqpoint{0.049702in}{0.028866in}}{\pgfqpoint{0.039284in}{0.039284in}}%
\pgfpathcurveto{\pgfqpoint{0.028866in}{0.049702in}}{\pgfqpoint{0.014734in}{0.055556in}}{\pgfqpoint{0.000000in}{0.055556in}}%
\pgfpathcurveto{\pgfqpoint{-0.014734in}{0.055556in}}{\pgfqpoint{-0.028866in}{0.049702in}}{\pgfqpoint{-0.039284in}{0.039284in}}%
\pgfpathcurveto{\pgfqpoint{-0.049702in}{0.028866in}}{\pgfqpoint{-0.055556in}{0.014734in}}{\pgfqpoint{-0.055556in}{0.000000in}}%
\pgfpathcurveto{\pgfqpoint{-0.055556in}{-0.014734in}}{\pgfqpoint{-0.049702in}{-0.028866in}}{\pgfqpoint{-0.039284in}{-0.039284in}}%
\pgfpathcurveto{\pgfqpoint{-0.028866in}{-0.049702in}}{\pgfqpoint{-0.014734in}{-0.055556in}}{\pgfqpoint{0.000000in}{-0.055556in}}%
\pgfpathlineto{\pgfqpoint{0.000000in}{-0.055556in}}%
\pgfpathclose%
\pgfusepath{stroke,fill}%
}%
\begin{pgfscope}%
\pgfsys@transformshift{3.550688in}{5.577550in}%
\pgfsys@useobject{currentmarker}{}%
\end{pgfscope}%
\end{pgfscope}%
\begin{pgfscope}%
\pgfpathrectangle{\pgfqpoint{0.100000in}{2.493063in}}{\pgfqpoint{5.037500in}{3.427208in}}%
\pgfusepath{clip}%
\pgfsetrectcap%
\pgfsetroundjoin%
\pgfsetlinewidth{1.505625pt}%
\definecolor{currentstroke}{rgb}{0.000000,0.000000,1.000000}%
\pgfsetstrokecolor{currentstroke}%
\pgfsetstrokeopacity{0.500000}%
\pgfsetdash{}{0pt}%
\pgfpathmoveto{\pgfqpoint{3.550688in}{5.440462in}}%
\pgfusepath{stroke}%
\end{pgfscope}%
\begin{pgfscope}%
\pgfpathrectangle{\pgfqpoint{0.100000in}{2.493063in}}{\pgfqpoint{5.037500in}{3.427208in}}%
\pgfusepath{clip}%
\pgfsetbuttcap%
\pgfsetroundjoin%
\definecolor{currentfill}{rgb}{0.000000,0.000000,1.000000}%
\pgfsetfillcolor{currentfill}%
\pgfsetfillopacity{0.500000}%
\pgfsetlinewidth{0.250937pt}%
\definecolor{currentstroke}{rgb}{0.000000,0.000000,0.000000}%
\pgfsetstrokecolor{currentstroke}%
\pgfsetstrokeopacity{0.500000}%
\pgfsetdash{}{0pt}%
\pgfsys@defobject{currentmarker}{\pgfqpoint{-0.027778in}{-0.027778in}}{\pgfqpoint{0.027778in}{0.027778in}}{%
\pgfpathmoveto{\pgfqpoint{0.000000in}{-0.027778in}}%
\pgfpathcurveto{\pgfqpoint{0.007367in}{-0.027778in}}{\pgfqpoint{0.014433in}{-0.024851in}}{\pgfqpoint{0.019642in}{-0.019642in}}%
\pgfpathcurveto{\pgfqpoint{0.024851in}{-0.014433in}}{\pgfqpoint{0.027778in}{-0.007367in}}{\pgfqpoint{0.027778in}{0.000000in}}%
\pgfpathcurveto{\pgfqpoint{0.027778in}{0.007367in}}{\pgfqpoint{0.024851in}{0.014433in}}{\pgfqpoint{0.019642in}{0.019642in}}%
\pgfpathcurveto{\pgfqpoint{0.014433in}{0.024851in}}{\pgfqpoint{0.007367in}{0.027778in}}{\pgfqpoint{0.000000in}{0.027778in}}%
\pgfpathcurveto{\pgfqpoint{-0.007367in}{0.027778in}}{\pgfqpoint{-0.014433in}{0.024851in}}{\pgfqpoint{-0.019642in}{0.019642in}}%
\pgfpathcurveto{\pgfqpoint{-0.024851in}{0.014433in}}{\pgfqpoint{-0.027778in}{0.007367in}}{\pgfqpoint{-0.027778in}{0.000000in}}%
\pgfpathcurveto{\pgfqpoint{-0.027778in}{-0.007367in}}{\pgfqpoint{-0.024851in}{-0.014433in}}{\pgfqpoint{-0.019642in}{-0.019642in}}%
\pgfpathcurveto{\pgfqpoint{-0.014433in}{-0.024851in}}{\pgfqpoint{-0.007367in}{-0.027778in}}{\pgfqpoint{0.000000in}{-0.027778in}}%
\pgfpathlineto{\pgfqpoint{0.000000in}{-0.027778in}}%
\pgfpathclose%
\pgfusepath{stroke,fill}%
}%
\begin{pgfscope}%
\pgfsys@transformshift{3.550688in}{5.440462in}%
\pgfsys@useobject{currentmarker}{}%
\end{pgfscope}%
\end{pgfscope}%
\begin{pgfscope}%
\pgfpathrectangle{\pgfqpoint{0.100000in}{2.493063in}}{\pgfqpoint{5.037500in}{3.427208in}}%
\pgfusepath{clip}%
\pgfsetrectcap%
\pgfsetroundjoin%
\pgfsetlinewidth{1.505625pt}%
\definecolor{currentstroke}{rgb}{0.678431,1.000000,0.184314}%
\pgfsetstrokecolor{currentstroke}%
\pgfsetstrokeopacity{0.500000}%
\pgfsetdash{}{0pt}%
\pgfpathmoveto{\pgfqpoint{4.205563in}{5.748910in}}%
\pgfusepath{stroke}%
\end{pgfscope}%
\begin{pgfscope}%
\pgfpathrectangle{\pgfqpoint{0.100000in}{2.493063in}}{\pgfqpoint{5.037500in}{3.427208in}}%
\pgfusepath{clip}%
\pgfsetbuttcap%
\pgfsetroundjoin%
\definecolor{currentfill}{rgb}{0.678431,1.000000,0.184314}%
\pgfsetfillcolor{currentfill}%
\pgfsetfillopacity{0.500000}%
\pgfsetlinewidth{0.250937pt}%
\definecolor{currentstroke}{rgb}{0.000000,0.000000,0.000000}%
\pgfsetstrokecolor{currentstroke}%
\pgfsetstrokeopacity{0.500000}%
\pgfsetdash{}{0pt}%
\pgfsys@defobject{currentmarker}{\pgfqpoint{-0.083333in}{-0.083333in}}{\pgfqpoint{0.083333in}{0.083333in}}{%
\pgfpathmoveto{\pgfqpoint{0.000000in}{-0.083333in}}%
\pgfpathcurveto{\pgfqpoint{0.022100in}{-0.083333in}}{\pgfqpoint{0.043298in}{-0.074553in}}{\pgfqpoint{0.058926in}{-0.058926in}}%
\pgfpathcurveto{\pgfqpoint{0.074553in}{-0.043298in}}{\pgfqpoint{0.083333in}{-0.022100in}}{\pgfqpoint{0.083333in}{0.000000in}}%
\pgfpathcurveto{\pgfqpoint{0.083333in}{0.022100in}}{\pgfqpoint{0.074553in}{0.043298in}}{\pgfqpoint{0.058926in}{0.058926in}}%
\pgfpathcurveto{\pgfqpoint{0.043298in}{0.074553in}}{\pgfqpoint{0.022100in}{0.083333in}}{\pgfqpoint{0.000000in}{0.083333in}}%
\pgfpathcurveto{\pgfqpoint{-0.022100in}{0.083333in}}{\pgfqpoint{-0.043298in}{0.074553in}}{\pgfqpoint{-0.058926in}{0.058926in}}%
\pgfpathcurveto{\pgfqpoint{-0.074553in}{0.043298in}}{\pgfqpoint{-0.083333in}{0.022100in}}{\pgfqpoint{-0.083333in}{0.000000in}}%
\pgfpathcurveto{\pgfqpoint{-0.083333in}{-0.022100in}}{\pgfqpoint{-0.074553in}{-0.043298in}}{\pgfqpoint{-0.058926in}{-0.058926in}}%
\pgfpathcurveto{\pgfqpoint{-0.043298in}{-0.074553in}}{\pgfqpoint{-0.022100in}{-0.083333in}}{\pgfqpoint{0.000000in}{-0.083333in}}%
\pgfpathlineto{\pgfqpoint{0.000000in}{-0.083333in}}%
\pgfpathclose%
\pgfusepath{stroke,fill}%
}%
\begin{pgfscope}%
\pgfsys@transformshift{4.205563in}{5.748910in}%
\pgfsys@useobject{currentmarker}{}%
\end{pgfscope}%
\end{pgfscope}%
\begin{pgfscope}%
\pgfpathrectangle{\pgfqpoint{0.100000in}{2.493063in}}{\pgfqpoint{5.037500in}{3.427208in}}%
\pgfusepath{clip}%
\pgfsetrectcap%
\pgfsetroundjoin%
\pgfsetlinewidth{1.505625pt}%
\definecolor{currentstroke}{rgb}{0.678431,1.000000,0.184314}%
\pgfsetstrokecolor{currentstroke}%
\pgfsetstrokeopacity{0.500000}%
\pgfsetdash{}{0pt}%
\pgfpathmoveto{\pgfqpoint{4.205563in}{5.577550in}}%
\pgfusepath{stroke}%
\end{pgfscope}%
\begin{pgfscope}%
\pgfpathrectangle{\pgfqpoint{0.100000in}{2.493063in}}{\pgfqpoint{5.037500in}{3.427208in}}%
\pgfusepath{clip}%
\pgfsetbuttcap%
\pgfsetroundjoin%
\definecolor{currentfill}{rgb}{0.678431,1.000000,0.184314}%
\pgfsetfillcolor{currentfill}%
\pgfsetfillopacity{0.500000}%
\pgfsetlinewidth{0.250937pt}%
\definecolor{currentstroke}{rgb}{0.000000,0.000000,0.000000}%
\pgfsetstrokecolor{currentstroke}%
\pgfsetstrokeopacity{0.500000}%
\pgfsetdash{}{0pt}%
\pgfsys@defobject{currentmarker}{\pgfqpoint{-0.055556in}{-0.055556in}}{\pgfqpoint{0.055556in}{0.055556in}}{%
\pgfpathmoveto{\pgfqpoint{0.000000in}{-0.055556in}}%
\pgfpathcurveto{\pgfqpoint{0.014734in}{-0.055556in}}{\pgfqpoint{0.028866in}{-0.049702in}}{\pgfqpoint{0.039284in}{-0.039284in}}%
\pgfpathcurveto{\pgfqpoint{0.049702in}{-0.028866in}}{\pgfqpoint{0.055556in}{-0.014734in}}{\pgfqpoint{0.055556in}{0.000000in}}%
\pgfpathcurveto{\pgfqpoint{0.055556in}{0.014734in}}{\pgfqpoint{0.049702in}{0.028866in}}{\pgfqpoint{0.039284in}{0.039284in}}%
\pgfpathcurveto{\pgfqpoint{0.028866in}{0.049702in}}{\pgfqpoint{0.014734in}{0.055556in}}{\pgfqpoint{0.000000in}{0.055556in}}%
\pgfpathcurveto{\pgfqpoint{-0.014734in}{0.055556in}}{\pgfqpoint{-0.028866in}{0.049702in}}{\pgfqpoint{-0.039284in}{0.039284in}}%
\pgfpathcurveto{\pgfqpoint{-0.049702in}{0.028866in}}{\pgfqpoint{-0.055556in}{0.014734in}}{\pgfqpoint{-0.055556in}{0.000000in}}%
\pgfpathcurveto{\pgfqpoint{-0.055556in}{-0.014734in}}{\pgfqpoint{-0.049702in}{-0.028866in}}{\pgfqpoint{-0.039284in}{-0.039284in}}%
\pgfpathcurveto{\pgfqpoint{-0.028866in}{-0.049702in}}{\pgfqpoint{-0.014734in}{-0.055556in}}{\pgfqpoint{0.000000in}{-0.055556in}}%
\pgfpathlineto{\pgfqpoint{0.000000in}{-0.055556in}}%
\pgfpathclose%
\pgfusepath{stroke,fill}%
}%
\begin{pgfscope}%
\pgfsys@transformshift{4.205563in}{5.577550in}%
\pgfsys@useobject{currentmarker}{}%
\end{pgfscope}%
\end{pgfscope}%
\begin{pgfscope}%
\pgfpathrectangle{\pgfqpoint{0.100000in}{2.493063in}}{\pgfqpoint{5.037500in}{3.427208in}}%
\pgfusepath{clip}%
\pgfsetrectcap%
\pgfsetroundjoin%
\pgfsetlinewidth{1.505625pt}%
\definecolor{currentstroke}{rgb}{0.678431,1.000000,0.184314}%
\pgfsetstrokecolor{currentstroke}%
\pgfsetstrokeopacity{0.500000}%
\pgfsetdash{}{0pt}%
\pgfpathmoveto{\pgfqpoint{4.205563in}{5.440462in}}%
\pgfusepath{stroke}%
\end{pgfscope}%
\begin{pgfscope}%
\pgfpathrectangle{\pgfqpoint{0.100000in}{2.493063in}}{\pgfqpoint{5.037500in}{3.427208in}}%
\pgfusepath{clip}%
\pgfsetbuttcap%
\pgfsetroundjoin%
\definecolor{currentfill}{rgb}{0.678431,1.000000,0.184314}%
\pgfsetfillcolor{currentfill}%
\pgfsetfillopacity{0.500000}%
\pgfsetlinewidth{0.250937pt}%
\definecolor{currentstroke}{rgb}{0.000000,0.000000,0.000000}%
\pgfsetstrokecolor{currentstroke}%
\pgfsetstrokeopacity{0.500000}%
\pgfsetdash{}{0pt}%
\pgfsys@defobject{currentmarker}{\pgfqpoint{-0.027778in}{-0.027778in}}{\pgfqpoint{0.027778in}{0.027778in}}{%
\pgfpathmoveto{\pgfqpoint{0.000000in}{-0.027778in}}%
\pgfpathcurveto{\pgfqpoint{0.007367in}{-0.027778in}}{\pgfqpoint{0.014433in}{-0.024851in}}{\pgfqpoint{0.019642in}{-0.019642in}}%
\pgfpathcurveto{\pgfqpoint{0.024851in}{-0.014433in}}{\pgfqpoint{0.027778in}{-0.007367in}}{\pgfqpoint{0.027778in}{0.000000in}}%
\pgfpathcurveto{\pgfqpoint{0.027778in}{0.007367in}}{\pgfqpoint{0.024851in}{0.014433in}}{\pgfqpoint{0.019642in}{0.019642in}}%
\pgfpathcurveto{\pgfqpoint{0.014433in}{0.024851in}}{\pgfqpoint{0.007367in}{0.027778in}}{\pgfqpoint{0.000000in}{0.027778in}}%
\pgfpathcurveto{\pgfqpoint{-0.007367in}{0.027778in}}{\pgfqpoint{-0.014433in}{0.024851in}}{\pgfqpoint{-0.019642in}{0.019642in}}%
\pgfpathcurveto{\pgfqpoint{-0.024851in}{0.014433in}}{\pgfqpoint{-0.027778in}{0.007367in}}{\pgfqpoint{-0.027778in}{0.000000in}}%
\pgfpathcurveto{\pgfqpoint{-0.027778in}{-0.007367in}}{\pgfqpoint{-0.024851in}{-0.014433in}}{\pgfqpoint{-0.019642in}{-0.019642in}}%
\pgfpathcurveto{\pgfqpoint{-0.014433in}{-0.024851in}}{\pgfqpoint{-0.007367in}{-0.027778in}}{\pgfqpoint{0.000000in}{-0.027778in}}%
\pgfpathlineto{\pgfqpoint{0.000000in}{-0.027778in}}%
\pgfpathclose%
\pgfusepath{stroke,fill}%
}%
\begin{pgfscope}%
\pgfsys@transformshift{4.205563in}{5.440462in}%
\pgfsys@useobject{currentmarker}{}%
\end{pgfscope}%
\end{pgfscope}%
\begin{pgfscope}%
\definecolor{textcolor}{rgb}{0.000000,0.000000,0.000000}%
\pgfsetstrokecolor{textcolor}%
\pgfsetfillcolor{textcolor}%
\pgftext[x=2.895812in,y=5.716352in,left,base]{\color{textcolor}\setmainfont{Lato}\rmfamily\fontsize{7.000000}{8.400000}\selectfont +/- 0.2pp}%
\end{pgfscope}%
\begin{pgfscope}%
\definecolor{textcolor}{rgb}{0.000000,0.000000,0.000000}%
\pgfsetstrokecolor{textcolor}%
\pgfsetfillcolor{textcolor}%
\pgftext[x=3.676625in,y=5.716352in,left,base]{\color{textcolor}\setmainfont{Lato}\rmfamily\fontsize{7.000000}{8.400000}\selectfont +3.0pp}%
\end{pgfscope}%
\begin{pgfscope}%
\definecolor{textcolor}{rgb}{0.000000,0.000000,0.000000}%
\pgfsetstrokecolor{textcolor}%
\pgfsetfillcolor{textcolor}%
\pgftext[x=3.676625in,y=5.563670in,left,base]{\color{textcolor}\setmainfont{Lato}\rmfamily\fontsize{7.000000}{8.400000}\selectfont +2.0pp}%
\end{pgfscope}%
\begin{pgfscope}%
\definecolor{textcolor}{rgb}{0.000000,0.000000,0.000000}%
\pgfsetstrokecolor{textcolor}%
\pgfsetfillcolor{textcolor}%
\pgftext[x=3.676625in,y=5.410988in,left,base]{\color{textcolor}\setmainfont{Lato}\rmfamily\fontsize{7.000000}{8.400000}\selectfont +1.0pp}%
\end{pgfscope}%
\begin{pgfscope}%
\definecolor{textcolor}{rgb}{0.000000,0.000000,0.000000}%
\pgfsetstrokecolor{textcolor}%
\pgfsetfillcolor{textcolor}%
\pgftext[x=4.331500in,y=5.716352in,left,base]{\color{textcolor}\setmainfont{Lato}\rmfamily\fontsize{7.000000}{8.400000}\selectfont -3.0pp}%
\end{pgfscope}%
\begin{pgfscope}%
\definecolor{textcolor}{rgb}{0.000000,0.000000,0.000000}%
\pgfsetstrokecolor{textcolor}%
\pgfsetfillcolor{textcolor}%
\pgftext[x=4.331500in,y=5.563670in,left,base]{\color{textcolor}\setmainfont{Lato}\rmfamily\fontsize{7.000000}{8.400000}\selectfont -2.0pp}%
\end{pgfscope}%
\begin{pgfscope}%
\definecolor{textcolor}{rgb}{0.000000,0.000000,0.000000}%
\pgfsetstrokecolor{textcolor}%
\pgfsetfillcolor{textcolor}%
\pgftext[x=4.331500in,y=5.410988in,left,base]{\color{textcolor}\setmainfont{Lato}\rmfamily\fontsize{7.000000}{8.400000}\selectfont -1.0pp}%
\end{pgfscope}%
\begin{pgfscope}%
\pgfpathrectangle{\pgfqpoint{0.100000in}{0.100000in}}{\pgfqpoint{5.037500in}{2.053333in}}%
\pgfusepath{clip}%
\pgfsetrectcap%
\pgfsetroundjoin%
\pgfsetlinewidth{1.505625pt}%
\definecolor{currentstroke}{rgb}{0.000000,0.000000,1.000000}%
\pgfsetstrokecolor{currentstroke}%
\pgfsetstrokeopacity{0.500000}%
\pgfsetdash{}{0pt}%
\pgfpathmoveto{\pgfqpoint{0.780063in}{2.050667in}}%
\pgfusepath{stroke}%
\end{pgfscope}%
\begin{pgfscope}%
\pgfpathrectangle{\pgfqpoint{0.100000in}{0.100000in}}{\pgfqpoint{5.037500in}{2.053333in}}%
\pgfusepath{clip}%
\pgfsetbuttcap%
\pgfsetroundjoin%
\definecolor{currentfill}{rgb}{0.000000,0.000000,1.000000}%
\pgfsetfillcolor{currentfill}%
\pgfsetfillopacity{0.500000}%
\pgfsetlinewidth{0.250937pt}%
\definecolor{currentstroke}{rgb}{0.000000,0.000000,0.000000}%
\pgfsetstrokecolor{currentstroke}%
\pgfsetstrokeopacity{0.500000}%
\pgfsetdash{}{0pt}%
\pgfsys@defobject{currentmarker}{\pgfqpoint{-0.022222in}{-0.022222in}}{\pgfqpoint{0.022222in}{0.022222in}}{%
\pgfpathmoveto{\pgfqpoint{0.000000in}{-0.022222in}}%
\pgfpathcurveto{\pgfqpoint{0.005893in}{-0.022222in}}{\pgfqpoint{0.011546in}{-0.019881in}}{\pgfqpoint{0.015713in}{-0.015713in}}%
\pgfpathcurveto{\pgfqpoint{0.019881in}{-0.011546in}}{\pgfqpoint{0.022222in}{-0.005893in}}{\pgfqpoint{0.022222in}{0.000000in}}%
\pgfpathcurveto{\pgfqpoint{0.022222in}{0.005893in}}{\pgfqpoint{0.019881in}{0.011546in}}{\pgfqpoint{0.015713in}{0.015713in}}%
\pgfpathcurveto{\pgfqpoint{0.011546in}{0.019881in}}{\pgfqpoint{0.005893in}{0.022222in}}{\pgfqpoint{0.000000in}{0.022222in}}%
\pgfpathcurveto{\pgfqpoint{-0.005893in}{0.022222in}}{\pgfqpoint{-0.011546in}{0.019881in}}{\pgfqpoint{-0.015713in}{0.015713in}}%
\pgfpathcurveto{\pgfqpoint{-0.019881in}{0.011546in}}{\pgfqpoint{-0.022222in}{0.005893in}}{\pgfqpoint{-0.022222in}{0.000000in}}%
\pgfpathcurveto{\pgfqpoint{-0.022222in}{-0.005893in}}{\pgfqpoint{-0.019881in}{-0.011546in}}{\pgfqpoint{-0.015713in}{-0.015713in}}%
\pgfpathcurveto{\pgfqpoint{-0.011546in}{-0.019881in}}{\pgfqpoint{-0.005893in}{-0.022222in}}{\pgfqpoint{0.000000in}{-0.022222in}}%
\pgfpathlineto{\pgfqpoint{0.000000in}{-0.022222in}}%
\pgfpathclose%
\pgfusepath{stroke,fill}%
}%
\begin{pgfscope}%
\pgfsys@transformshift{0.780063in}{2.050667in}%
\pgfsys@useobject{currentmarker}{}%
\end{pgfscope}%
\end{pgfscope}%
\begin{pgfscope}%
\pgfpathrectangle{\pgfqpoint{0.100000in}{0.100000in}}{\pgfqpoint{5.037500in}{2.053333in}}%
\pgfusepath{clip}%
\pgfsetrectcap%
\pgfsetroundjoin%
\pgfsetlinewidth{0.250937pt}%
\definecolor{currentstroke}{rgb}{0.862745,0.862745,0.862745}%
\pgfsetstrokecolor{currentstroke}%
\pgfsetdash{}{0pt}%
\pgfpathmoveto{\pgfqpoint{0.906000in}{1.968533in}}%
\pgfpathlineto{\pgfqpoint{4.180375in}{1.968533in}}%
\pgfusepath{stroke}%
\end{pgfscope}%
\begin{pgfscope}%
\pgfpathrectangle{\pgfqpoint{0.100000in}{0.100000in}}{\pgfqpoint{5.037500in}{2.053333in}}%
\pgfusepath{clip}%
\pgfsetrectcap%
\pgfsetroundjoin%
\pgfsetlinewidth{1.505625pt}%
\definecolor{currentstroke}{rgb}{0.000000,0.000000,1.000000}%
\pgfsetstrokecolor{currentstroke}%
\pgfsetstrokeopacity{0.500000}%
\pgfsetdash{}{0pt}%
\pgfpathmoveto{\pgfqpoint{0.780063in}{1.864000in}}%
\pgfusepath{stroke}%
\end{pgfscope}%
\begin{pgfscope}%
\pgfpathrectangle{\pgfqpoint{0.100000in}{0.100000in}}{\pgfqpoint{5.037500in}{2.053333in}}%
\pgfusepath{clip}%
\pgfsetbuttcap%
\pgfsetroundjoin%
\definecolor{currentfill}{rgb}{0.000000,0.000000,1.000000}%
\pgfsetfillcolor{currentfill}%
\pgfsetfillopacity{0.500000}%
\pgfsetlinewidth{0.250937pt}%
\definecolor{currentstroke}{rgb}{0.000000,0.000000,0.000000}%
\pgfsetstrokecolor{currentstroke}%
\pgfsetstrokeopacity{0.500000}%
\pgfsetdash{}{0pt}%
\pgfsys@defobject{currentmarker}{\pgfqpoint{-0.022222in}{-0.022222in}}{\pgfqpoint{0.022222in}{0.022222in}}{%
\pgfpathmoveto{\pgfqpoint{0.000000in}{-0.022222in}}%
\pgfpathcurveto{\pgfqpoint{0.005893in}{-0.022222in}}{\pgfqpoint{0.011546in}{-0.019881in}}{\pgfqpoint{0.015713in}{-0.015713in}}%
\pgfpathcurveto{\pgfqpoint{0.019881in}{-0.011546in}}{\pgfqpoint{0.022222in}{-0.005893in}}{\pgfqpoint{0.022222in}{0.000000in}}%
\pgfpathcurveto{\pgfqpoint{0.022222in}{0.005893in}}{\pgfqpoint{0.019881in}{0.011546in}}{\pgfqpoint{0.015713in}{0.015713in}}%
\pgfpathcurveto{\pgfqpoint{0.011546in}{0.019881in}}{\pgfqpoint{0.005893in}{0.022222in}}{\pgfqpoint{0.000000in}{0.022222in}}%
\pgfpathcurveto{\pgfqpoint{-0.005893in}{0.022222in}}{\pgfqpoint{-0.011546in}{0.019881in}}{\pgfqpoint{-0.015713in}{0.015713in}}%
\pgfpathcurveto{\pgfqpoint{-0.019881in}{0.011546in}}{\pgfqpoint{-0.022222in}{0.005893in}}{\pgfqpoint{-0.022222in}{0.000000in}}%
\pgfpathcurveto{\pgfqpoint{-0.022222in}{-0.005893in}}{\pgfqpoint{-0.019881in}{-0.011546in}}{\pgfqpoint{-0.015713in}{-0.015713in}}%
\pgfpathcurveto{\pgfqpoint{-0.011546in}{-0.019881in}}{\pgfqpoint{-0.005893in}{-0.022222in}}{\pgfqpoint{0.000000in}{-0.022222in}}%
\pgfpathlineto{\pgfqpoint{0.000000in}{-0.022222in}}%
\pgfpathclose%
\pgfusepath{stroke,fill}%
}%
\begin{pgfscope}%
\pgfsys@transformshift{0.780063in}{1.864000in}%
\pgfsys@useobject{currentmarker}{}%
\end{pgfscope}%
\end{pgfscope}%
\begin{pgfscope}%
\pgfpathrectangle{\pgfqpoint{0.100000in}{0.100000in}}{\pgfqpoint{5.037500in}{2.053333in}}%
\pgfusepath{clip}%
\pgfsetrectcap%
\pgfsetroundjoin%
\pgfsetlinewidth{0.250937pt}%
\definecolor{currentstroke}{rgb}{0.862745,0.862745,0.862745}%
\pgfsetstrokecolor{currentstroke}%
\pgfsetdash{}{0pt}%
\pgfpathmoveto{\pgfqpoint{0.906000in}{1.781867in}}%
\pgfpathlineto{\pgfqpoint{4.180375in}{1.781867in}}%
\pgfusepath{stroke}%
\end{pgfscope}%
\begin{pgfscope}%
\pgfpathrectangle{\pgfqpoint{0.100000in}{0.100000in}}{\pgfqpoint{5.037500in}{2.053333in}}%
\pgfusepath{clip}%
\pgfsetrectcap%
\pgfsetroundjoin%
\pgfsetlinewidth{1.505625pt}%
\definecolor{currentstroke}{rgb}{0.000000,0.000000,1.000000}%
\pgfsetstrokecolor{currentstroke}%
\pgfsetstrokeopacity{0.500000}%
\pgfsetdash{}{0pt}%
\pgfpathmoveto{\pgfqpoint{0.780063in}{1.677333in}}%
\pgfusepath{stroke}%
\end{pgfscope}%
\begin{pgfscope}%
\pgfpathrectangle{\pgfqpoint{0.100000in}{0.100000in}}{\pgfqpoint{5.037500in}{2.053333in}}%
\pgfusepath{clip}%
\pgfsetbuttcap%
\pgfsetroundjoin%
\definecolor{currentfill}{rgb}{0.000000,0.000000,1.000000}%
\pgfsetfillcolor{currentfill}%
\pgfsetfillopacity{0.500000}%
\pgfsetlinewidth{0.250937pt}%
\definecolor{currentstroke}{rgb}{0.000000,0.000000,0.000000}%
\pgfsetstrokecolor{currentstroke}%
\pgfsetstrokeopacity{0.500000}%
\pgfsetdash{}{0pt}%
\pgfsys@defobject{currentmarker}{\pgfqpoint{-0.030556in}{-0.030556in}}{\pgfqpoint{0.030556in}{0.030556in}}{%
\pgfpathmoveto{\pgfqpoint{0.000000in}{-0.030556in}}%
\pgfpathcurveto{\pgfqpoint{0.008103in}{-0.030556in}}{\pgfqpoint{0.015876in}{-0.027336in}}{\pgfqpoint{0.021606in}{-0.021606in}}%
\pgfpathcurveto{\pgfqpoint{0.027336in}{-0.015876in}}{\pgfqpoint{0.030556in}{-0.008103in}}{\pgfqpoint{0.030556in}{0.000000in}}%
\pgfpathcurveto{\pgfqpoint{0.030556in}{0.008103in}}{\pgfqpoint{0.027336in}{0.015876in}}{\pgfqpoint{0.021606in}{0.021606in}}%
\pgfpathcurveto{\pgfqpoint{0.015876in}{0.027336in}}{\pgfqpoint{0.008103in}{0.030556in}}{\pgfqpoint{0.000000in}{0.030556in}}%
\pgfpathcurveto{\pgfqpoint{-0.008103in}{0.030556in}}{\pgfqpoint{-0.015876in}{0.027336in}}{\pgfqpoint{-0.021606in}{0.021606in}}%
\pgfpathcurveto{\pgfqpoint{-0.027336in}{0.015876in}}{\pgfqpoint{-0.030556in}{0.008103in}}{\pgfqpoint{-0.030556in}{0.000000in}}%
\pgfpathcurveto{\pgfqpoint{-0.030556in}{-0.008103in}}{\pgfqpoint{-0.027336in}{-0.015876in}}{\pgfqpoint{-0.021606in}{-0.021606in}}%
\pgfpathcurveto{\pgfqpoint{-0.015876in}{-0.027336in}}{\pgfqpoint{-0.008103in}{-0.030556in}}{\pgfqpoint{0.000000in}{-0.030556in}}%
\pgfpathlineto{\pgfqpoint{0.000000in}{-0.030556in}}%
\pgfpathclose%
\pgfusepath{stroke,fill}%
}%
\begin{pgfscope}%
\pgfsys@transformshift{0.780063in}{1.677333in}%
\pgfsys@useobject{currentmarker}{}%
\end{pgfscope}%
\end{pgfscope}%
\begin{pgfscope}%
\pgfpathrectangle{\pgfqpoint{0.100000in}{0.100000in}}{\pgfqpoint{5.037500in}{2.053333in}}%
\pgfusepath{clip}%
\pgfsetrectcap%
\pgfsetroundjoin%
\pgfsetlinewidth{0.250937pt}%
\definecolor{currentstroke}{rgb}{0.862745,0.862745,0.862745}%
\pgfsetstrokecolor{currentstroke}%
\pgfsetdash{}{0pt}%
\pgfpathmoveto{\pgfqpoint{0.906000in}{1.595200in}}%
\pgfpathlineto{\pgfqpoint{4.180375in}{1.595200in}}%
\pgfusepath{stroke}%
\end{pgfscope}%
\begin{pgfscope}%
\pgfpathrectangle{\pgfqpoint{0.100000in}{0.100000in}}{\pgfqpoint{5.037500in}{2.053333in}}%
\pgfusepath{clip}%
\pgfsetrectcap%
\pgfsetroundjoin%
\pgfsetlinewidth{1.505625pt}%
\definecolor{currentstroke}{rgb}{0.501961,0.501961,0.501961}%
\pgfsetstrokecolor{currentstroke}%
\pgfsetstrokeopacity{0.500000}%
\pgfsetdash{}{0pt}%
\pgfpathmoveto{\pgfqpoint{0.780063in}{1.490667in}}%
\pgfusepath{stroke}%
\end{pgfscope}%
\begin{pgfscope}%
\pgfpathrectangle{\pgfqpoint{0.100000in}{0.100000in}}{\pgfqpoint{5.037500in}{2.053333in}}%
\pgfusepath{clip}%
\pgfsetbuttcap%
\pgfsetroundjoin%
\definecolor{currentfill}{rgb}{0.501961,0.501961,0.501961}%
\pgfsetfillcolor{currentfill}%
\pgfsetfillopacity{0.500000}%
\pgfsetlinewidth{0.250937pt}%
\definecolor{currentstroke}{rgb}{0.000000,0.000000,0.000000}%
\pgfsetstrokecolor{currentstroke}%
\pgfsetstrokeopacity{0.500000}%
\pgfsetdash{}{0pt}%
\pgfsys@defobject{currentmarker}{\pgfqpoint{-0.013889in}{-0.013889in}}{\pgfqpoint{0.013889in}{0.013889in}}{%
\pgfpathmoveto{\pgfqpoint{0.000000in}{-0.013889in}}%
\pgfpathcurveto{\pgfqpoint{0.003683in}{-0.013889in}}{\pgfqpoint{0.007216in}{-0.012425in}}{\pgfqpoint{0.009821in}{-0.009821in}}%
\pgfpathcurveto{\pgfqpoint{0.012425in}{-0.007216in}}{\pgfqpoint{0.013889in}{-0.003683in}}{\pgfqpoint{0.013889in}{0.000000in}}%
\pgfpathcurveto{\pgfqpoint{0.013889in}{0.003683in}}{\pgfqpoint{0.012425in}{0.007216in}}{\pgfqpoint{0.009821in}{0.009821in}}%
\pgfpathcurveto{\pgfqpoint{0.007216in}{0.012425in}}{\pgfqpoint{0.003683in}{0.013889in}}{\pgfqpoint{0.000000in}{0.013889in}}%
\pgfpathcurveto{\pgfqpoint{-0.003683in}{0.013889in}}{\pgfqpoint{-0.007216in}{0.012425in}}{\pgfqpoint{-0.009821in}{0.009821in}}%
\pgfpathcurveto{\pgfqpoint{-0.012425in}{0.007216in}}{\pgfqpoint{-0.013889in}{0.003683in}}{\pgfqpoint{-0.013889in}{0.000000in}}%
\pgfpathcurveto{\pgfqpoint{-0.013889in}{-0.003683in}}{\pgfqpoint{-0.012425in}{-0.007216in}}{\pgfqpoint{-0.009821in}{-0.009821in}}%
\pgfpathcurveto{\pgfqpoint{-0.007216in}{-0.012425in}}{\pgfqpoint{-0.003683in}{-0.013889in}}{\pgfqpoint{0.000000in}{-0.013889in}}%
\pgfpathlineto{\pgfqpoint{0.000000in}{-0.013889in}}%
\pgfpathclose%
\pgfusepath{stroke,fill}%
}%
\begin{pgfscope}%
\pgfsys@transformshift{0.780063in}{1.490667in}%
\pgfsys@useobject{currentmarker}{}%
\end{pgfscope}%
\end{pgfscope}%
\begin{pgfscope}%
\pgfpathrectangle{\pgfqpoint{0.100000in}{0.100000in}}{\pgfqpoint{5.037500in}{2.053333in}}%
\pgfusepath{clip}%
\pgfsetrectcap%
\pgfsetroundjoin%
\pgfsetlinewidth{0.250937pt}%
\definecolor{currentstroke}{rgb}{0.862745,0.862745,0.862745}%
\pgfsetstrokecolor{currentstroke}%
\pgfsetdash{}{0pt}%
\pgfpathmoveto{\pgfqpoint{0.906000in}{1.408533in}}%
\pgfpathlineto{\pgfqpoint{4.180375in}{1.408533in}}%
\pgfusepath{stroke}%
\end{pgfscope}%
\begin{pgfscope}%
\pgfpathrectangle{\pgfqpoint{0.100000in}{0.100000in}}{\pgfqpoint{5.037500in}{2.053333in}}%
\pgfusepath{clip}%
\pgfsetrectcap%
\pgfsetroundjoin%
\pgfsetlinewidth{1.505625pt}%
\definecolor{currentstroke}{rgb}{0.501961,0.501961,0.501961}%
\pgfsetstrokecolor{currentstroke}%
\pgfsetstrokeopacity{0.500000}%
\pgfsetdash{}{0pt}%
\pgfpathmoveto{\pgfqpoint{0.780063in}{1.304000in}}%
\pgfusepath{stroke}%
\end{pgfscope}%
\begin{pgfscope}%
\pgfpathrectangle{\pgfqpoint{0.100000in}{0.100000in}}{\pgfqpoint{5.037500in}{2.053333in}}%
\pgfusepath{clip}%
\pgfsetbuttcap%
\pgfsetroundjoin%
\definecolor{currentfill}{rgb}{0.501961,0.501961,0.501961}%
\pgfsetfillcolor{currentfill}%
\pgfsetfillopacity{0.500000}%
\pgfsetlinewidth{0.250937pt}%
\definecolor{currentstroke}{rgb}{0.000000,0.000000,0.000000}%
\pgfsetstrokecolor{currentstroke}%
\pgfsetstrokeopacity{0.500000}%
\pgfsetdash{}{0pt}%
\pgfsys@defobject{currentmarker}{\pgfqpoint{-0.013889in}{-0.013889in}}{\pgfqpoint{0.013889in}{0.013889in}}{%
\pgfpathmoveto{\pgfqpoint{0.000000in}{-0.013889in}}%
\pgfpathcurveto{\pgfqpoint{0.003683in}{-0.013889in}}{\pgfqpoint{0.007216in}{-0.012425in}}{\pgfqpoint{0.009821in}{-0.009821in}}%
\pgfpathcurveto{\pgfqpoint{0.012425in}{-0.007216in}}{\pgfqpoint{0.013889in}{-0.003683in}}{\pgfqpoint{0.013889in}{0.000000in}}%
\pgfpathcurveto{\pgfqpoint{0.013889in}{0.003683in}}{\pgfqpoint{0.012425in}{0.007216in}}{\pgfqpoint{0.009821in}{0.009821in}}%
\pgfpathcurveto{\pgfqpoint{0.007216in}{0.012425in}}{\pgfqpoint{0.003683in}{0.013889in}}{\pgfqpoint{0.000000in}{0.013889in}}%
\pgfpathcurveto{\pgfqpoint{-0.003683in}{0.013889in}}{\pgfqpoint{-0.007216in}{0.012425in}}{\pgfqpoint{-0.009821in}{0.009821in}}%
\pgfpathcurveto{\pgfqpoint{-0.012425in}{0.007216in}}{\pgfqpoint{-0.013889in}{0.003683in}}{\pgfqpoint{-0.013889in}{0.000000in}}%
\pgfpathcurveto{\pgfqpoint{-0.013889in}{-0.003683in}}{\pgfqpoint{-0.012425in}{-0.007216in}}{\pgfqpoint{-0.009821in}{-0.009821in}}%
\pgfpathcurveto{\pgfqpoint{-0.007216in}{-0.012425in}}{\pgfqpoint{-0.003683in}{-0.013889in}}{\pgfqpoint{0.000000in}{-0.013889in}}%
\pgfpathlineto{\pgfqpoint{0.000000in}{-0.013889in}}%
\pgfpathclose%
\pgfusepath{stroke,fill}%
}%
\begin{pgfscope}%
\pgfsys@transformshift{0.780063in}{1.304000in}%
\pgfsys@useobject{currentmarker}{}%
\end{pgfscope}%
\end{pgfscope}%
\begin{pgfscope}%
\pgfpathrectangle{\pgfqpoint{0.100000in}{0.100000in}}{\pgfqpoint{5.037500in}{2.053333in}}%
\pgfusepath{clip}%
\pgfsetrectcap%
\pgfsetroundjoin%
\pgfsetlinewidth{0.250937pt}%
\definecolor{currentstroke}{rgb}{0.862745,0.862745,0.862745}%
\pgfsetstrokecolor{currentstroke}%
\pgfsetdash{}{0pt}%
\pgfpathmoveto{\pgfqpoint{0.906000in}{1.221867in}}%
\pgfpathlineto{\pgfqpoint{4.180375in}{1.221867in}}%
\pgfusepath{stroke}%
\end{pgfscope}%
\begin{pgfscope}%
\pgfpathrectangle{\pgfqpoint{0.100000in}{0.100000in}}{\pgfqpoint{5.037500in}{2.053333in}}%
\pgfusepath{clip}%
\pgfsetrectcap%
\pgfsetroundjoin%
\pgfsetlinewidth{1.505625pt}%
\definecolor{currentstroke}{rgb}{0.678431,1.000000,0.184314}%
\pgfsetstrokecolor{currentstroke}%
\pgfsetstrokeopacity{0.500000}%
\pgfsetdash{}{0pt}%
\pgfpathmoveto{\pgfqpoint{0.780063in}{1.117333in}}%
\pgfusepath{stroke}%
\end{pgfscope}%
\begin{pgfscope}%
\pgfpathrectangle{\pgfqpoint{0.100000in}{0.100000in}}{\pgfqpoint{5.037500in}{2.053333in}}%
\pgfusepath{clip}%
\pgfsetbuttcap%
\pgfsetroundjoin%
\definecolor{currentfill}{rgb}{0.678431,1.000000,0.184314}%
\pgfsetfillcolor{currentfill}%
\pgfsetfillopacity{0.500000}%
\pgfsetlinewidth{0.250937pt}%
\definecolor{currentstroke}{rgb}{0.000000,0.000000,0.000000}%
\pgfsetstrokecolor{currentstroke}%
\pgfsetstrokeopacity{0.500000}%
\pgfsetdash{}{0pt}%
\pgfsys@defobject{currentmarker}{\pgfqpoint{-0.005556in}{-0.005556in}}{\pgfqpoint{0.005556in}{0.005556in}}{%
\pgfpathmoveto{\pgfqpoint{0.000000in}{-0.005556in}}%
\pgfpathcurveto{\pgfqpoint{0.001473in}{-0.005556in}}{\pgfqpoint{0.002887in}{-0.004970in}}{\pgfqpoint{0.003928in}{-0.003928in}}%
\pgfpathcurveto{\pgfqpoint{0.004970in}{-0.002887in}}{\pgfqpoint{0.005556in}{-0.001473in}}{\pgfqpoint{0.005556in}{0.000000in}}%
\pgfpathcurveto{\pgfqpoint{0.005556in}{0.001473in}}{\pgfqpoint{0.004970in}{0.002887in}}{\pgfqpoint{0.003928in}{0.003928in}}%
\pgfpathcurveto{\pgfqpoint{0.002887in}{0.004970in}}{\pgfqpoint{0.001473in}{0.005556in}}{\pgfqpoint{0.000000in}{0.005556in}}%
\pgfpathcurveto{\pgfqpoint{-0.001473in}{0.005556in}}{\pgfqpoint{-0.002887in}{0.004970in}}{\pgfqpoint{-0.003928in}{0.003928in}}%
\pgfpathcurveto{\pgfqpoint{-0.004970in}{0.002887in}}{\pgfqpoint{-0.005556in}{0.001473in}}{\pgfqpoint{-0.005556in}{0.000000in}}%
\pgfpathcurveto{\pgfqpoint{-0.005556in}{-0.001473in}}{\pgfqpoint{-0.004970in}{-0.002887in}}{\pgfqpoint{-0.003928in}{-0.003928in}}%
\pgfpathcurveto{\pgfqpoint{-0.002887in}{-0.004970in}}{\pgfqpoint{-0.001473in}{-0.005556in}}{\pgfqpoint{0.000000in}{-0.005556in}}%
\pgfpathlineto{\pgfqpoint{0.000000in}{-0.005556in}}%
\pgfpathclose%
\pgfusepath{stroke,fill}%
}%
\begin{pgfscope}%
\pgfsys@transformshift{0.780063in}{1.117333in}%
\pgfsys@useobject{currentmarker}{}%
\end{pgfscope}%
\end{pgfscope}%
\begin{pgfscope}%
\pgfpathrectangle{\pgfqpoint{0.100000in}{0.100000in}}{\pgfqpoint{5.037500in}{2.053333in}}%
\pgfusepath{clip}%
\pgfsetrectcap%
\pgfsetroundjoin%
\pgfsetlinewidth{0.250937pt}%
\definecolor{currentstroke}{rgb}{0.862745,0.862745,0.862745}%
\pgfsetstrokecolor{currentstroke}%
\pgfsetdash{}{0pt}%
\pgfpathmoveto{\pgfqpoint{0.906000in}{1.035200in}}%
\pgfpathlineto{\pgfqpoint{4.180375in}{1.035200in}}%
\pgfusepath{stroke}%
\end{pgfscope}%
\begin{pgfscope}%
\pgfpathrectangle{\pgfqpoint{0.100000in}{0.100000in}}{\pgfqpoint{5.037500in}{2.053333in}}%
\pgfusepath{clip}%
\pgfsetrectcap%
\pgfsetroundjoin%
\pgfsetlinewidth{1.505625pt}%
\definecolor{currentstroke}{rgb}{0.678431,1.000000,0.184314}%
\pgfsetstrokecolor{currentstroke}%
\pgfsetstrokeopacity{0.500000}%
\pgfsetdash{}{0pt}%
\pgfpathmoveto{\pgfqpoint{0.780063in}{0.930667in}}%
\pgfusepath{stroke}%
\end{pgfscope}%
\begin{pgfscope}%
\pgfpathrectangle{\pgfqpoint{0.100000in}{0.100000in}}{\pgfqpoint{5.037500in}{2.053333in}}%
\pgfusepath{clip}%
\pgfsetbuttcap%
\pgfsetroundjoin%
\definecolor{currentfill}{rgb}{0.678431,1.000000,0.184314}%
\pgfsetfillcolor{currentfill}%
\pgfsetfillopacity{0.500000}%
\pgfsetlinewidth{0.250937pt}%
\definecolor{currentstroke}{rgb}{0.000000,0.000000,0.000000}%
\pgfsetstrokecolor{currentstroke}%
\pgfsetstrokeopacity{0.500000}%
\pgfsetdash{}{0pt}%
\pgfsys@defobject{currentmarker}{\pgfqpoint{-0.011111in}{-0.011111in}}{\pgfqpoint{0.011111in}{0.011111in}}{%
\pgfpathmoveto{\pgfqpoint{0.000000in}{-0.011111in}}%
\pgfpathcurveto{\pgfqpoint{0.002947in}{-0.011111in}}{\pgfqpoint{0.005773in}{-0.009940in}}{\pgfqpoint{0.007857in}{-0.007857in}}%
\pgfpathcurveto{\pgfqpoint{0.009940in}{-0.005773in}}{\pgfqpoint{0.011111in}{-0.002947in}}{\pgfqpoint{0.011111in}{0.000000in}}%
\pgfpathcurveto{\pgfqpoint{0.011111in}{0.002947in}}{\pgfqpoint{0.009940in}{0.005773in}}{\pgfqpoint{0.007857in}{0.007857in}}%
\pgfpathcurveto{\pgfqpoint{0.005773in}{0.009940in}}{\pgfqpoint{0.002947in}{0.011111in}}{\pgfqpoint{0.000000in}{0.011111in}}%
\pgfpathcurveto{\pgfqpoint{-0.002947in}{0.011111in}}{\pgfqpoint{-0.005773in}{0.009940in}}{\pgfqpoint{-0.007857in}{0.007857in}}%
\pgfpathcurveto{\pgfqpoint{-0.009940in}{0.005773in}}{\pgfqpoint{-0.011111in}{0.002947in}}{\pgfqpoint{-0.011111in}{0.000000in}}%
\pgfpathcurveto{\pgfqpoint{-0.011111in}{-0.002947in}}{\pgfqpoint{-0.009940in}{-0.005773in}}{\pgfqpoint{-0.007857in}{-0.007857in}}%
\pgfpathcurveto{\pgfqpoint{-0.005773in}{-0.009940in}}{\pgfqpoint{-0.002947in}{-0.011111in}}{\pgfqpoint{0.000000in}{-0.011111in}}%
\pgfpathlineto{\pgfqpoint{0.000000in}{-0.011111in}}%
\pgfpathclose%
\pgfusepath{stroke,fill}%
}%
\begin{pgfscope}%
\pgfsys@transformshift{0.780063in}{0.930667in}%
\pgfsys@useobject{currentmarker}{}%
\end{pgfscope}%
\end{pgfscope}%
\begin{pgfscope}%
\pgfpathrectangle{\pgfqpoint{0.100000in}{0.100000in}}{\pgfqpoint{5.037500in}{2.053333in}}%
\pgfusepath{clip}%
\pgfsetrectcap%
\pgfsetroundjoin%
\pgfsetlinewidth{0.250937pt}%
\definecolor{currentstroke}{rgb}{0.862745,0.862745,0.862745}%
\pgfsetstrokecolor{currentstroke}%
\pgfsetdash{}{0pt}%
\pgfpathmoveto{\pgfqpoint{0.906000in}{0.848533in}}%
\pgfpathlineto{\pgfqpoint{4.180375in}{0.848533in}}%
\pgfusepath{stroke}%
\end{pgfscope}%
\begin{pgfscope}%
\pgfpathrectangle{\pgfqpoint{0.100000in}{0.100000in}}{\pgfqpoint{5.037500in}{2.053333in}}%
\pgfusepath{clip}%
\pgfsetrectcap%
\pgfsetroundjoin%
\pgfsetlinewidth{1.505625pt}%
\definecolor{currentstroke}{rgb}{0.678431,1.000000,0.184314}%
\pgfsetstrokecolor{currentstroke}%
\pgfsetstrokeopacity{0.500000}%
\pgfsetdash{}{0pt}%
\pgfpathmoveto{\pgfqpoint{0.780063in}{0.744000in}}%
\pgfusepath{stroke}%
\end{pgfscope}%
\begin{pgfscope}%
\pgfpathrectangle{\pgfqpoint{0.100000in}{0.100000in}}{\pgfqpoint{5.037500in}{2.053333in}}%
\pgfusepath{clip}%
\pgfsetbuttcap%
\pgfsetroundjoin%
\definecolor{currentfill}{rgb}{0.678431,1.000000,0.184314}%
\pgfsetfillcolor{currentfill}%
\pgfsetfillopacity{0.500000}%
\pgfsetlinewidth{0.250937pt}%
\definecolor{currentstroke}{rgb}{0.000000,0.000000,0.000000}%
\pgfsetstrokecolor{currentstroke}%
\pgfsetstrokeopacity{0.500000}%
\pgfsetdash{}{0pt}%
\pgfsys@defobject{currentmarker}{\pgfqpoint{-0.011111in}{-0.011111in}}{\pgfqpoint{0.011111in}{0.011111in}}{%
\pgfpathmoveto{\pgfqpoint{0.000000in}{-0.011111in}}%
\pgfpathcurveto{\pgfqpoint{0.002947in}{-0.011111in}}{\pgfqpoint{0.005773in}{-0.009940in}}{\pgfqpoint{0.007857in}{-0.007857in}}%
\pgfpathcurveto{\pgfqpoint{0.009940in}{-0.005773in}}{\pgfqpoint{0.011111in}{-0.002947in}}{\pgfqpoint{0.011111in}{0.000000in}}%
\pgfpathcurveto{\pgfqpoint{0.011111in}{0.002947in}}{\pgfqpoint{0.009940in}{0.005773in}}{\pgfqpoint{0.007857in}{0.007857in}}%
\pgfpathcurveto{\pgfqpoint{0.005773in}{0.009940in}}{\pgfqpoint{0.002947in}{0.011111in}}{\pgfqpoint{0.000000in}{0.011111in}}%
\pgfpathcurveto{\pgfqpoint{-0.002947in}{0.011111in}}{\pgfqpoint{-0.005773in}{0.009940in}}{\pgfqpoint{-0.007857in}{0.007857in}}%
\pgfpathcurveto{\pgfqpoint{-0.009940in}{0.005773in}}{\pgfqpoint{-0.011111in}{0.002947in}}{\pgfqpoint{-0.011111in}{0.000000in}}%
\pgfpathcurveto{\pgfqpoint{-0.011111in}{-0.002947in}}{\pgfqpoint{-0.009940in}{-0.005773in}}{\pgfqpoint{-0.007857in}{-0.007857in}}%
\pgfpathcurveto{\pgfqpoint{-0.005773in}{-0.009940in}}{\pgfqpoint{-0.002947in}{-0.011111in}}{\pgfqpoint{0.000000in}{-0.011111in}}%
\pgfpathlineto{\pgfqpoint{0.000000in}{-0.011111in}}%
\pgfpathclose%
\pgfusepath{stroke,fill}%
}%
\begin{pgfscope}%
\pgfsys@transformshift{0.780063in}{0.744000in}%
\pgfsys@useobject{currentmarker}{}%
\end{pgfscope}%
\end{pgfscope}%
\begin{pgfscope}%
\pgfpathrectangle{\pgfqpoint{0.100000in}{0.100000in}}{\pgfqpoint{5.037500in}{2.053333in}}%
\pgfusepath{clip}%
\pgfsetrectcap%
\pgfsetroundjoin%
\pgfsetlinewidth{0.250937pt}%
\definecolor{currentstroke}{rgb}{0.862745,0.862745,0.862745}%
\pgfsetstrokecolor{currentstroke}%
\pgfsetdash{}{0pt}%
\pgfpathmoveto{\pgfqpoint{0.906000in}{0.661867in}}%
\pgfpathlineto{\pgfqpoint{4.180375in}{0.661867in}}%
\pgfusepath{stroke}%
\end{pgfscope}%
\begin{pgfscope}%
\pgfpathrectangle{\pgfqpoint{0.100000in}{0.100000in}}{\pgfqpoint{5.037500in}{2.053333in}}%
\pgfusepath{clip}%
\pgfsetrectcap%
\pgfsetroundjoin%
\pgfsetlinewidth{1.505625pt}%
\definecolor{currentstroke}{rgb}{0.000000,0.000000,1.000000}%
\pgfsetstrokecolor{currentstroke}%
\pgfsetstrokeopacity{0.500000}%
\pgfsetdash{}{0pt}%
\pgfpathmoveto{\pgfqpoint{0.780063in}{0.557333in}}%
\pgfusepath{stroke}%
\end{pgfscope}%
\begin{pgfscope}%
\pgfpathrectangle{\pgfqpoint{0.100000in}{0.100000in}}{\pgfqpoint{5.037500in}{2.053333in}}%
\pgfusepath{clip}%
\pgfsetbuttcap%
\pgfsetroundjoin%
\definecolor{currentfill}{rgb}{0.000000,0.000000,1.000000}%
\pgfsetfillcolor{currentfill}%
\pgfsetfillopacity{0.500000}%
\pgfsetlinewidth{0.250937pt}%
\definecolor{currentstroke}{rgb}{0.000000,0.000000,0.000000}%
\pgfsetstrokecolor{currentstroke}%
\pgfsetstrokeopacity{0.500000}%
\pgfsetdash{}{0pt}%
\pgfsys@defobject{currentmarker}{\pgfqpoint{-0.005556in}{-0.005556in}}{\pgfqpoint{0.005556in}{0.005556in}}{%
\pgfpathmoveto{\pgfqpoint{0.000000in}{-0.005556in}}%
\pgfpathcurveto{\pgfqpoint{0.001473in}{-0.005556in}}{\pgfqpoint{0.002887in}{-0.004970in}}{\pgfqpoint{0.003928in}{-0.003928in}}%
\pgfpathcurveto{\pgfqpoint{0.004970in}{-0.002887in}}{\pgfqpoint{0.005556in}{-0.001473in}}{\pgfqpoint{0.005556in}{0.000000in}}%
\pgfpathcurveto{\pgfqpoint{0.005556in}{0.001473in}}{\pgfqpoint{0.004970in}{0.002887in}}{\pgfqpoint{0.003928in}{0.003928in}}%
\pgfpathcurveto{\pgfqpoint{0.002887in}{0.004970in}}{\pgfqpoint{0.001473in}{0.005556in}}{\pgfqpoint{0.000000in}{0.005556in}}%
\pgfpathcurveto{\pgfqpoint{-0.001473in}{0.005556in}}{\pgfqpoint{-0.002887in}{0.004970in}}{\pgfqpoint{-0.003928in}{0.003928in}}%
\pgfpathcurveto{\pgfqpoint{-0.004970in}{0.002887in}}{\pgfqpoint{-0.005556in}{0.001473in}}{\pgfqpoint{-0.005556in}{0.000000in}}%
\pgfpathcurveto{\pgfqpoint{-0.005556in}{-0.001473in}}{\pgfqpoint{-0.004970in}{-0.002887in}}{\pgfqpoint{-0.003928in}{-0.003928in}}%
\pgfpathcurveto{\pgfqpoint{-0.002887in}{-0.004970in}}{\pgfqpoint{-0.001473in}{-0.005556in}}{\pgfqpoint{0.000000in}{-0.005556in}}%
\pgfpathlineto{\pgfqpoint{0.000000in}{-0.005556in}}%
\pgfpathclose%
\pgfusepath{stroke,fill}%
}%
\begin{pgfscope}%
\pgfsys@transformshift{0.780063in}{0.557333in}%
\pgfsys@useobject{currentmarker}{}%
\end{pgfscope}%
\end{pgfscope}%
\begin{pgfscope}%
\pgfpathrectangle{\pgfqpoint{0.100000in}{0.100000in}}{\pgfqpoint{5.037500in}{2.053333in}}%
\pgfusepath{clip}%
\pgfsetrectcap%
\pgfsetroundjoin%
\pgfsetlinewidth{0.250937pt}%
\definecolor{currentstroke}{rgb}{0.862745,0.862745,0.862745}%
\pgfsetstrokecolor{currentstroke}%
\pgfsetdash{}{0pt}%
\pgfpathmoveto{\pgfqpoint{0.906000in}{0.475200in}}%
\pgfpathlineto{\pgfqpoint{4.180375in}{0.475200in}}%
\pgfusepath{stroke}%
\end{pgfscope}%
\begin{pgfscope}%
\pgfpathrectangle{\pgfqpoint{0.100000in}{0.100000in}}{\pgfqpoint{5.037500in}{2.053333in}}%
\pgfusepath{clip}%
\pgfsetrectcap%
\pgfsetroundjoin%
\pgfsetlinewidth{1.505625pt}%
\definecolor{currentstroke}{rgb}{0.000000,0.000000,1.000000}%
\pgfsetstrokecolor{currentstroke}%
\pgfsetstrokeopacity{0.500000}%
\pgfsetdash{}{0pt}%
\pgfpathmoveto{\pgfqpoint{0.780063in}{0.370667in}}%
\pgfusepath{stroke}%
\end{pgfscope}%
\begin{pgfscope}%
\pgfpathrectangle{\pgfqpoint{0.100000in}{0.100000in}}{\pgfqpoint{5.037500in}{2.053333in}}%
\pgfusepath{clip}%
\pgfsetbuttcap%
\pgfsetroundjoin%
\definecolor{currentfill}{rgb}{0.000000,0.000000,1.000000}%
\pgfsetfillcolor{currentfill}%
\pgfsetfillopacity{0.500000}%
\pgfsetlinewidth{0.250937pt}%
\definecolor{currentstroke}{rgb}{0.000000,0.000000,0.000000}%
\pgfsetstrokecolor{currentstroke}%
\pgfsetstrokeopacity{0.500000}%
\pgfsetdash{}{0pt}%
\pgfsys@defobject{currentmarker}{\pgfqpoint{-0.013889in}{-0.013889in}}{\pgfqpoint{0.013889in}{0.013889in}}{%
\pgfpathmoveto{\pgfqpoint{0.000000in}{-0.013889in}}%
\pgfpathcurveto{\pgfqpoint{0.003683in}{-0.013889in}}{\pgfqpoint{0.007216in}{-0.012425in}}{\pgfqpoint{0.009821in}{-0.009821in}}%
\pgfpathcurveto{\pgfqpoint{0.012425in}{-0.007216in}}{\pgfqpoint{0.013889in}{-0.003683in}}{\pgfqpoint{0.013889in}{0.000000in}}%
\pgfpathcurveto{\pgfqpoint{0.013889in}{0.003683in}}{\pgfqpoint{0.012425in}{0.007216in}}{\pgfqpoint{0.009821in}{0.009821in}}%
\pgfpathcurveto{\pgfqpoint{0.007216in}{0.012425in}}{\pgfqpoint{0.003683in}{0.013889in}}{\pgfqpoint{0.000000in}{0.013889in}}%
\pgfpathcurveto{\pgfqpoint{-0.003683in}{0.013889in}}{\pgfqpoint{-0.007216in}{0.012425in}}{\pgfqpoint{-0.009821in}{0.009821in}}%
\pgfpathcurveto{\pgfqpoint{-0.012425in}{0.007216in}}{\pgfqpoint{-0.013889in}{0.003683in}}{\pgfqpoint{-0.013889in}{0.000000in}}%
\pgfpathcurveto{\pgfqpoint{-0.013889in}{-0.003683in}}{\pgfqpoint{-0.012425in}{-0.007216in}}{\pgfqpoint{-0.009821in}{-0.009821in}}%
\pgfpathcurveto{\pgfqpoint{-0.007216in}{-0.012425in}}{\pgfqpoint{-0.003683in}{-0.013889in}}{\pgfqpoint{0.000000in}{-0.013889in}}%
\pgfpathlineto{\pgfqpoint{0.000000in}{-0.013889in}}%
\pgfpathclose%
\pgfusepath{stroke,fill}%
}%
\begin{pgfscope}%
\pgfsys@transformshift{0.780063in}{0.370667in}%
\pgfsys@useobject{currentmarker}{}%
\end{pgfscope}%
\end{pgfscope}%
\begin{pgfscope}%
\pgfpathrectangle{\pgfqpoint{0.100000in}{0.100000in}}{\pgfqpoint{5.037500in}{2.053333in}}%
\pgfusepath{clip}%
\pgfsetrectcap%
\pgfsetroundjoin%
\pgfsetlinewidth{0.250937pt}%
\definecolor{currentstroke}{rgb}{0.862745,0.862745,0.862745}%
\pgfsetstrokecolor{currentstroke}%
\pgfsetdash{}{0pt}%
\pgfpathmoveto{\pgfqpoint{0.906000in}{0.288533in}}%
\pgfpathlineto{\pgfqpoint{4.180375in}{0.288533in}}%
\pgfusepath{stroke}%
\end{pgfscope}%
\begin{pgfscope}%
\pgfpathrectangle{\pgfqpoint{0.100000in}{0.100000in}}{\pgfqpoint{5.037500in}{2.053333in}}%
\pgfusepath{clip}%
\pgfsetrectcap%
\pgfsetroundjoin%
\pgfsetlinewidth{1.505625pt}%
\definecolor{currentstroke}{rgb}{0.678431,1.000000,0.184314}%
\pgfsetstrokecolor{currentstroke}%
\pgfsetstrokeopacity{0.500000}%
\pgfsetdash{}{0pt}%
\pgfpathmoveto{\pgfqpoint{0.780063in}{0.184000in}}%
\pgfusepath{stroke}%
\end{pgfscope}%
\begin{pgfscope}%
\pgfpathrectangle{\pgfqpoint{0.100000in}{0.100000in}}{\pgfqpoint{5.037500in}{2.053333in}}%
\pgfusepath{clip}%
\pgfsetbuttcap%
\pgfsetroundjoin%
\definecolor{currentfill}{rgb}{0.678431,1.000000,0.184314}%
\pgfsetfillcolor{currentfill}%
\pgfsetfillopacity{0.500000}%
\pgfsetlinewidth{0.250937pt}%
\definecolor{currentstroke}{rgb}{0.000000,0.000000,0.000000}%
\pgfsetstrokecolor{currentstroke}%
\pgfsetstrokeopacity{0.500000}%
\pgfsetdash{}{0pt}%
\pgfsys@defobject{currentmarker}{\pgfqpoint{-0.008333in}{-0.008333in}}{\pgfqpoint{0.008333in}{0.008333in}}{%
\pgfpathmoveto{\pgfqpoint{0.000000in}{-0.008333in}}%
\pgfpathcurveto{\pgfqpoint{0.002210in}{-0.008333in}}{\pgfqpoint{0.004330in}{-0.007455in}}{\pgfqpoint{0.005893in}{-0.005893in}}%
\pgfpathcurveto{\pgfqpoint{0.007455in}{-0.004330in}}{\pgfqpoint{0.008333in}{-0.002210in}}{\pgfqpoint{0.008333in}{0.000000in}}%
\pgfpathcurveto{\pgfqpoint{0.008333in}{0.002210in}}{\pgfqpoint{0.007455in}{0.004330in}}{\pgfqpoint{0.005893in}{0.005893in}}%
\pgfpathcurveto{\pgfqpoint{0.004330in}{0.007455in}}{\pgfqpoint{0.002210in}{0.008333in}}{\pgfqpoint{0.000000in}{0.008333in}}%
\pgfpathcurveto{\pgfqpoint{-0.002210in}{0.008333in}}{\pgfqpoint{-0.004330in}{0.007455in}}{\pgfqpoint{-0.005893in}{0.005893in}}%
\pgfpathcurveto{\pgfqpoint{-0.007455in}{0.004330in}}{\pgfqpoint{-0.008333in}{0.002210in}}{\pgfqpoint{-0.008333in}{0.000000in}}%
\pgfpathcurveto{\pgfqpoint{-0.008333in}{-0.002210in}}{\pgfqpoint{-0.007455in}{-0.004330in}}{\pgfqpoint{-0.005893in}{-0.005893in}}%
\pgfpathcurveto{\pgfqpoint{-0.004330in}{-0.007455in}}{\pgfqpoint{-0.002210in}{-0.008333in}}{\pgfqpoint{0.000000in}{-0.008333in}}%
\pgfpathlineto{\pgfqpoint{0.000000in}{-0.008333in}}%
\pgfpathclose%
\pgfusepath{stroke,fill}%
}%
\begin{pgfscope}%
\pgfsys@transformshift{0.780063in}{0.184000in}%
\pgfsys@useobject{currentmarker}{}%
\end{pgfscope}%
\end{pgfscope}%
\begin{pgfscope}%
\pgfpathrectangle{\pgfqpoint{0.100000in}{0.100000in}}{\pgfqpoint{5.037500in}{2.053333in}}%
\pgfusepath{clip}%
\pgfsetrectcap%
\pgfsetroundjoin%
\pgfsetlinewidth{0.250937pt}%
\definecolor{currentstroke}{rgb}{0.862745,0.862745,0.862745}%
\pgfsetstrokecolor{currentstroke}%
\pgfsetdash{}{0pt}%
\pgfpathmoveto{\pgfqpoint{0.906000in}{0.101867in}}%
\pgfpathlineto{\pgfqpoint{4.180375in}{0.101867in}}%
\pgfusepath{stroke}%
\end{pgfscope}%
\begin{pgfscope}%
\definecolor{textcolor}{rgb}{0.000000,0.000000,0.000000}%
\pgfsetstrokecolor{textcolor}%
\pgfsetfillcolor{textcolor}%
\pgftext[x=2.266125in,y=5.714638in,left,base]{\color{textcolor}\setmainfont{Lato}\rmfamily\fontsize{8.000000}{9.600000}\bfseries\selectfont Legend:}%
\end{pgfscope}%
\begin{pgfscope}%
\definecolor{textcolor}{rgb}{0.000000,0.000000,0.000000}%
\pgfsetstrokecolor{textcolor}%
\pgfsetfillcolor{textcolor}%
\pgftext[x=0.906000in,y=2.009600in,left,base]{\color{textcolor}\setmainfont{Lato}\rmfamily\fontsize{8.000000}{9.600000}\selectfont New York, NY}%
\end{pgfscope}%
\begin{pgfscope}%
\definecolor{textcolor}{rgb}{0.000000,0.000000,0.000000}%
\pgfsetstrokecolor{textcolor}%
\pgfsetfillcolor{textcolor}%
\pgftext[x=1.963875in,y=2.009600in,left,base]{\color{textcolor}\setmainfont{Lato}\rmfamily\fontsize{8.000000}{9.600000}\selectfont 4.2}%
\end{pgfscope}%
\begin{pgfscope}%
\definecolor{textcolor}{rgb}{0.000000,0.000000,0.000000}%
\pgfsetstrokecolor{textcolor}%
\pgfsetfillcolor{textcolor}%
\pgftext[x=2.568375in,y=2.009600in,left,base]{\color{textcolor}\setmainfont{Lato}\rmfamily\fontsize{8.000000}{9.600000}\selectfont 3.4}%
\end{pgfscope}%
\begin{pgfscope}%
\definecolor{textcolor}{rgb}{0.000000,0.000000,0.000000}%
\pgfsetstrokecolor{textcolor}%
\pgfsetfillcolor{textcolor}%
\pgftext[x=0.654125in,y=2.009600in,right,base]{\color{textcolor}\setmainfont{Lato}\rmfamily\fontsize{8.000000}{9.600000}\selectfont +0.8}%
\end{pgfscope}%
\begin{pgfscope}%
\definecolor{textcolor}{rgb}{0.000000,0.000000,0.000000}%
\pgfsetstrokecolor{textcolor}%
\pgfsetfillcolor{textcolor}%
\pgftext[x=3.097312in,y=2.009600in,left,base]{\color{textcolor}\setmainfont{Lato}\rmfamily\fontsize{8.000000}{9.600000}\selectfont 10,348,700}%
\end{pgfscope}%
\begin{pgfscope}%
\definecolor{textcolor}{rgb}{0.000000,0.000000,0.000000}%
\pgfsetstrokecolor{textcolor}%
\pgfsetfillcolor{textcolor}%
\pgftext[x=4.180375in,y=2.009600in,right,base]{\color{textcolor}\setmainfont{Lato}\rmfamily\fontsize{8.000000}{9.600000}\selectfont -0.7}%
\end{pgfscope}%
\begin{pgfscope}%
\definecolor{textcolor}{rgb}{0.000000,0.000000,0.000000}%
\pgfsetstrokecolor{textcolor}%
\pgfsetfillcolor{textcolor}%
\pgftext[x=0.906000in,y=1.822933in,left,base]{\color{textcolor}\setmainfont{Lato}\rmfamily\fontsize{8.000000}{9.600000}\selectfont Los Angeles, CA}%
\end{pgfscope}%
\begin{pgfscope}%
\definecolor{textcolor}{rgb}{0.000000,0.000000,0.000000}%
\pgfsetstrokecolor{textcolor}%
\pgfsetfillcolor{textcolor}%
\pgftext[x=1.963875in,y=1.822933in,left,base]{\color{textcolor}\setmainfont{Lato}\rmfamily\fontsize{8.000000}{9.600000}\selectfont 4.7}%
\end{pgfscope}%
\begin{pgfscope}%
\definecolor{textcolor}{rgb}{0.000000,0.000000,0.000000}%
\pgfsetstrokecolor{textcolor}%
\pgfsetfillcolor{textcolor}%
\pgftext[x=2.568375in,y=1.822933in,left,base]{\color{textcolor}\setmainfont{Lato}\rmfamily\fontsize{8.000000}{9.600000}\selectfont 3.9}%
\end{pgfscope}%
\begin{pgfscope}%
\definecolor{textcolor}{rgb}{0.000000,0.000000,0.000000}%
\pgfsetstrokecolor{textcolor}%
\pgfsetfillcolor{textcolor}%
\pgftext[x=0.654125in,y=1.822933in,right,base]{\color{textcolor}\setmainfont{Lato}\rmfamily\fontsize{8.000000}{9.600000}\selectfont +0.8}%
\end{pgfscope}%
\begin{pgfscope}%
\definecolor{textcolor}{rgb}{0.000000,0.000000,0.000000}%
\pgfsetstrokecolor{textcolor}%
\pgfsetfillcolor{textcolor}%
\pgftext[x=3.097312in,y=1.822933in,left,base]{\color{textcolor}\setmainfont{Lato}\rmfamily\fontsize{8.000000}{9.600000}\selectfont 6,569,200}%
\end{pgfscope}%
\begin{pgfscope}%
\definecolor{textcolor}{rgb}{0.000000,0.000000,0.000000}%
\pgfsetstrokecolor{textcolor}%
\pgfsetfillcolor{textcolor}%
\pgftext[x=4.180375in,y=1.822933in,right,base]{\color{textcolor}\setmainfont{Lato}\rmfamily\fontsize{8.000000}{9.600000}\selectfont -3.5}%
\end{pgfscope}%
\begin{pgfscope}%
\definecolor{textcolor}{rgb}{0.000000,0.000000,0.000000}%
\pgfsetstrokecolor{textcolor}%
\pgfsetfillcolor{textcolor}%
\pgftext[x=0.906000in,y=1.636267in,left,base]{\color{textcolor}\setmainfont{Lato}\rmfamily\fontsize{8.000000}{9.600000}\selectfont Chicago, IL}%
\end{pgfscope}%
\begin{pgfscope}%
\definecolor{textcolor}{rgb}{0.000000,0.000000,0.000000}%
\pgfsetstrokecolor{textcolor}%
\pgfsetfillcolor{textcolor}%
\pgftext[x=1.963875in,y=1.636267in,left,base]{\color{textcolor}\setmainfont{Lato}\rmfamily\fontsize{8.000000}{9.600000}\selectfont 4.3}%
\end{pgfscope}%
\begin{pgfscope}%
\definecolor{textcolor}{rgb}{0.000000,0.000000,0.000000}%
\pgfsetstrokecolor{textcolor}%
\pgfsetfillcolor{textcolor}%
\pgftext[x=2.568375in,y=1.636267in,left,base]{\color{textcolor}\setmainfont{Lato}\rmfamily\fontsize{8.000000}{9.600000}\selectfont 3.2}%
\end{pgfscope}%
\begin{pgfscope}%
\definecolor{textcolor}{rgb}{0.000000,0.000000,0.000000}%
\pgfsetstrokecolor{textcolor}%
\pgfsetfillcolor{textcolor}%
\pgftext[x=0.654125in,y=1.636267in,right,base]{\color{textcolor}\setmainfont{Lato}\rmfamily\fontsize{8.000000}{9.600000}\selectfont +1.1}%
\end{pgfscope}%
\begin{pgfscope}%
\definecolor{textcolor}{rgb}{0.000000,0.000000,0.000000}%
\pgfsetstrokecolor{textcolor}%
\pgfsetfillcolor{textcolor}%
\pgftext[x=3.097312in,y=1.636267in,left,base]{\color{textcolor}\setmainfont{Lato}\rmfamily\fontsize{8.000000}{9.600000}\selectfont 4,960,100}%
\end{pgfscope}%
\begin{pgfscope}%
\definecolor{textcolor}{rgb}{0.000000,0.000000,0.000000}%
\pgfsetstrokecolor{textcolor}%
\pgfsetfillcolor{textcolor}%
\pgftext[x=4.180375in,y=1.636267in,right,base]{\color{textcolor}\setmainfont{Lato}\rmfamily\fontsize{8.000000}{9.600000}\selectfont 0.3}%
\end{pgfscope}%
\begin{pgfscope}%
\definecolor{textcolor}{rgb}{0.000000,0.000000,0.000000}%
\pgfsetstrokecolor{textcolor}%
\pgfsetfillcolor{textcolor}%
\pgftext[x=0.906000in,y=1.449600in,left,base]{\color{textcolor}\setmainfont{Lato}\rmfamily\fontsize{8.000000}{9.600000}\selectfont Dallas, TX}%
\end{pgfscope}%
\begin{pgfscope}%
\definecolor{textcolor}{rgb}{0.000000,0.000000,0.000000}%
\pgfsetstrokecolor{textcolor}%
\pgfsetfillcolor{textcolor}%
\pgftext[x=1.963875in,y=1.449600in,left,base]{\color{textcolor}\setmainfont{Lato}\rmfamily\fontsize{8.000000}{9.600000}\selectfont 3.3}%
\end{pgfscope}%
\begin{pgfscope}%
\definecolor{textcolor}{rgb}{0.000000,0.000000,0.000000}%
\pgfsetstrokecolor{textcolor}%
\pgfsetfillcolor{textcolor}%
\pgftext[x=2.568375in,y=1.449600in,left,base]{\color{textcolor}\setmainfont{Lato}\rmfamily\fontsize{8.000000}{9.600000}\selectfont 3.1}%
\end{pgfscope}%
\begin{pgfscope}%
\definecolor{textcolor}{rgb}{0.000000,0.000000,0.000000}%
\pgfsetstrokecolor{textcolor}%
\pgfsetfillcolor{textcolor}%
\pgftext[x=0.654125in,y=1.449600in,right,base]{\color{textcolor}\setmainfont{Lato}\rmfamily\fontsize{8.000000}{9.600000}\selectfont unch.}%
\end{pgfscope}%
\begin{pgfscope}%
\definecolor{textcolor}{rgb}{0.000000,0.000000,0.000000}%
\pgfsetstrokecolor{textcolor}%
\pgfsetfillcolor{textcolor}%
\pgftext[x=3.097312in,y=1.449600in,left,base]{\color{textcolor}\setmainfont{Lato}\rmfamily\fontsize{8.000000}{9.600000}\selectfont 4,494,200}%
\end{pgfscope}%
\begin{pgfscope}%
\definecolor{textcolor}{rgb}{0.000000,0.000000,0.000000}%
\pgfsetstrokecolor{textcolor}%
\pgfsetfillcolor{textcolor}%
\pgftext[x=4.180375in,y=1.449600in,right,base]{\color{textcolor}\setmainfont{Lato}\rmfamily\fontsize{8.000000}{9.600000}\selectfont 12.6}%
\end{pgfscope}%
\begin{pgfscope}%
\definecolor{textcolor}{rgb}{0.000000,0.000000,0.000000}%
\pgfsetstrokecolor{textcolor}%
\pgfsetfillcolor{textcolor}%
\pgftext[x=0.906000in,y=1.262933in,left,base]{\color{textcolor}\setmainfont{Lato}\rmfamily\fontsize{8.000000}{9.600000}\selectfont Houston, TX}%
\end{pgfscope}%
\begin{pgfscope}%
\definecolor{textcolor}{rgb}{0.000000,0.000000,0.000000}%
\pgfsetstrokecolor{textcolor}%
\pgfsetfillcolor{textcolor}%
\pgftext[x=1.963875in,y=1.262933in,left,base]{\color{textcolor}\setmainfont{Lato}\rmfamily\fontsize{8.000000}{9.600000}\selectfont 3.8}%
\end{pgfscope}%
\begin{pgfscope}%
\definecolor{textcolor}{rgb}{0.000000,0.000000,0.000000}%
\pgfsetstrokecolor{textcolor}%
\pgfsetfillcolor{textcolor}%
\pgftext[x=2.568375in,y=1.262933in,left,base]{\color{textcolor}\setmainfont{Lato}\rmfamily\fontsize{8.000000}{9.600000}\selectfont 3.6}%
\end{pgfscope}%
\begin{pgfscope}%
\definecolor{textcolor}{rgb}{0.000000,0.000000,0.000000}%
\pgfsetstrokecolor{textcolor}%
\pgfsetfillcolor{textcolor}%
\pgftext[x=0.654125in,y=1.262933in,right,base]{\color{textcolor}\setmainfont{Lato}\rmfamily\fontsize{8.000000}{9.600000}\selectfont unch.}%
\end{pgfscope}%
\begin{pgfscope}%
\definecolor{textcolor}{rgb}{0.000000,0.000000,0.000000}%
\pgfsetstrokecolor{textcolor}%
\pgfsetfillcolor{textcolor}%
\pgftext[x=3.097312in,y=1.262933in,left,base]{\color{textcolor}\setmainfont{Lato}\rmfamily\fontsize{8.000000}{9.600000}\selectfont 3,645,100}%
\end{pgfscope}%
\begin{pgfscope}%
\definecolor{textcolor}{rgb}{0.000000,0.000000,0.000000}%
\pgfsetstrokecolor{textcolor}%
\pgfsetfillcolor{textcolor}%
\pgftext[x=4.180375in,y=1.262933in,right,base]{\color{textcolor}\setmainfont{Lato}\rmfamily\fontsize{8.000000}{9.600000}\selectfont 6.6}%
\end{pgfscope}%
\begin{pgfscope}%
\definecolor{textcolor}{rgb}{0.000000,0.000000,0.000000}%
\pgfsetstrokecolor{textcolor}%
\pgfsetfillcolor{textcolor}%
\pgftext[x=0.906000in,y=1.076267in,left,base]{\color{textcolor}\setmainfont{Lato}\rmfamily\fontsize{8.000000}{9.600000}\selectfont Washington, DC}%
\end{pgfscope}%
\begin{pgfscope}%
\definecolor{textcolor}{rgb}{0.000000,0.000000,0.000000}%
\pgfsetstrokecolor{textcolor}%
\pgfsetfillcolor{textcolor}%
\pgftext[x=1.963875in,y=1.076267in,left,base]{\color{textcolor}\setmainfont{Lato}\rmfamily\fontsize{8.000000}{9.600000}\selectfont 2.5}%
\end{pgfscope}%
\begin{pgfscope}%
\definecolor{textcolor}{rgb}{0.000000,0.000000,0.000000}%
\pgfsetstrokecolor{textcolor}%
\pgfsetfillcolor{textcolor}%
\pgftext[x=2.568375in,y=1.076267in,left,base]{\color{textcolor}\setmainfont{Lato}\rmfamily\fontsize{8.000000}{9.600000}\selectfont 2.7}%
\end{pgfscope}%
\begin{pgfscope}%
\definecolor{textcolor}{rgb}{0.000000,0.000000,0.000000}%
\pgfsetstrokecolor{textcolor}%
\pgfsetfillcolor{textcolor}%
\pgftext[x=0.654125in,y=1.076267in,right,base]{\color{textcolor}\setmainfont{Lato}\rmfamily\fontsize{8.000000}{9.600000}\selectfont unch.}%
\end{pgfscope}%
\begin{pgfscope}%
\definecolor{textcolor}{rgb}{0.000000,0.000000,0.000000}%
\pgfsetstrokecolor{textcolor}%
\pgfsetfillcolor{textcolor}%
\pgftext[x=3.097312in,y=1.076267in,left,base]{\color{textcolor}\setmainfont{Lato}\rmfamily\fontsize{8.000000}{9.600000}\selectfont 3,506,200}%
\end{pgfscope}%
\begin{pgfscope}%
\definecolor{textcolor}{rgb}{0.000000,0.000000,0.000000}%
\pgfsetstrokecolor{textcolor}%
\pgfsetfillcolor{textcolor}%
\pgftext[x=4.180375in,y=1.076267in,right,base]{\color{textcolor}\setmainfont{Lato}\rmfamily\fontsize{8.000000}{9.600000}\selectfont -0.1}%
\end{pgfscope}%
\begin{pgfscope}%
\definecolor{textcolor}{rgb}{0.000000,0.000000,0.000000}%
\pgfsetstrokecolor{textcolor}%
\pgfsetfillcolor{textcolor}%
\pgftext[x=0.906000in,y=0.889600in,left,base]{\color{textcolor}\setmainfont{Lato}\rmfamily\fontsize{8.000000}{9.600000}\selectfont Miami, FL}%
\end{pgfscope}%
\begin{pgfscope}%
\definecolor{textcolor}{rgb}{0.000000,0.000000,0.000000}%
\pgfsetstrokecolor{textcolor}%
\pgfsetfillcolor{textcolor}%
\pgftext[x=1.963875in,y=0.889600in,left,base]{\color{textcolor}\setmainfont{Lato}\rmfamily\fontsize{8.000000}{9.600000}\selectfont 2.3}%
\end{pgfscope}%
\begin{pgfscope}%
\definecolor{textcolor}{rgb}{0.000000,0.000000,0.000000}%
\pgfsetstrokecolor{textcolor}%
\pgfsetfillcolor{textcolor}%
\pgftext[x=2.568375in,y=0.889600in,left,base]{\color{textcolor}\setmainfont{Lato}\rmfamily\fontsize{8.000000}{9.600000}\selectfont 2.7}%
\end{pgfscope}%
\begin{pgfscope}%
\definecolor{textcolor}{rgb}{0.000000,0.000000,0.000000}%
\pgfsetstrokecolor{textcolor}%
\pgfsetfillcolor{textcolor}%
\pgftext[x=0.654125in,y=0.889600in,right,base]{\color{textcolor}\setmainfont{Lato}\rmfamily\fontsize{8.000000}{9.600000}\selectfont -0.4}%
\end{pgfscope}%
\begin{pgfscope}%
\definecolor{textcolor}{rgb}{0.000000,0.000000,0.000000}%
\pgfsetstrokecolor{textcolor}%
\pgfsetfillcolor{textcolor}%
\pgftext[x=3.097312in,y=0.889600in,left,base]{\color{textcolor}\setmainfont{Lato}\rmfamily\fontsize{8.000000}{9.600000}\selectfont 3,308,400}%
\end{pgfscope}%
\begin{pgfscope}%
\definecolor{textcolor}{rgb}{0.000000,0.000000,0.000000}%
\pgfsetstrokecolor{textcolor}%
\pgfsetfillcolor{textcolor}%
\pgftext[x=4.180375in,y=0.889600in,right,base]{\color{textcolor}\setmainfont{Lato}\rmfamily\fontsize{8.000000}{9.600000}\selectfont 3.8}%
\end{pgfscope}%
\begin{pgfscope}%
\definecolor{textcolor}{rgb}{0.000000,0.000000,0.000000}%
\pgfsetstrokecolor{textcolor}%
\pgfsetfillcolor{textcolor}%
\pgftext[x=0.906000in,y=0.702933in,left,base]{\color{textcolor}\setmainfont{Lato}\rmfamily\fontsize{8.000000}{9.600000}\selectfont Philadelphia, PA}%
\end{pgfscope}%
\begin{pgfscope}%
\definecolor{textcolor}{rgb}{0.000000,0.000000,0.000000}%
\pgfsetstrokecolor{textcolor}%
\pgfsetfillcolor{textcolor}%
\pgftext[x=1.963875in,y=0.702933in,left,base]{\color{textcolor}\setmainfont{Lato}\rmfamily\fontsize{8.000000}{9.600000}\selectfont 3.3}%
\end{pgfscope}%
\begin{pgfscope}%
\definecolor{textcolor}{rgb}{0.000000,0.000000,0.000000}%
\pgfsetstrokecolor{textcolor}%
\pgfsetfillcolor{textcolor}%
\pgftext[x=2.568375in,y=0.702933in,left,base]{\color{textcolor}\setmainfont{Lato}\rmfamily\fontsize{8.000000}{9.600000}\selectfont 3.7}%
\end{pgfscope}%
\begin{pgfscope}%
\definecolor{textcolor}{rgb}{0.000000,0.000000,0.000000}%
\pgfsetstrokecolor{textcolor}%
\pgfsetfillcolor{textcolor}%
\pgftext[x=0.654125in,y=0.702933in,right,base]{\color{textcolor}\setmainfont{Lato}\rmfamily\fontsize{8.000000}{9.600000}\selectfont -0.4}%
\end{pgfscope}%
\begin{pgfscope}%
\definecolor{textcolor}{rgb}{0.000000,0.000000,0.000000}%
\pgfsetstrokecolor{textcolor}%
\pgfsetfillcolor{textcolor}%
\pgftext[x=3.097312in,y=0.702933in,left,base]{\color{textcolor}\setmainfont{Lato}\rmfamily\fontsize{8.000000}{9.600000}\selectfont 3,290,200}%
\end{pgfscope}%
\begin{pgfscope}%
\definecolor{textcolor}{rgb}{0.000000,0.000000,0.000000}%
\pgfsetstrokecolor{textcolor}%
\pgfsetfillcolor{textcolor}%
\pgftext[x=4.180375in,y=0.702933in,right,base]{\color{textcolor}\setmainfont{Lato}\rmfamily\fontsize{8.000000}{9.600000}\selectfont 2.2}%
\end{pgfscope}%
\begin{pgfscope}%
\definecolor{textcolor}{rgb}{0.000000,0.000000,0.000000}%
\pgfsetstrokecolor{textcolor}%
\pgfsetfillcolor{textcolor}%
\pgftext[x=0.906000in,y=0.516267in,left,base]{\color{textcolor}\setmainfont{Lato}\rmfamily\fontsize{8.000000}{9.600000}\selectfont Atlanta, GA}%
\end{pgfscope}%
\begin{pgfscope}%
\definecolor{textcolor}{rgb}{0.000000,0.000000,0.000000}%
\pgfsetstrokecolor{textcolor}%
\pgfsetfillcolor{textcolor}%
\pgftext[x=1.963875in,y=0.516267in,left,base]{\color{textcolor}\setmainfont{Lato}\rmfamily\fontsize{8.000000}{9.600000}\selectfont 3.1}%
\end{pgfscope}%
\begin{pgfscope}%
\definecolor{textcolor}{rgb}{0.000000,0.000000,0.000000}%
\pgfsetstrokecolor{textcolor}%
\pgfsetfillcolor{textcolor}%
\pgftext[x=2.568375in,y=0.516267in,left,base]{\color{textcolor}\setmainfont{Lato}\rmfamily\fontsize{8.000000}{9.600000}\selectfont 2.9}%
\end{pgfscope}%
\begin{pgfscope}%
\definecolor{textcolor}{rgb}{0.000000,0.000000,0.000000}%
\pgfsetstrokecolor{textcolor}%
\pgfsetfillcolor{textcolor}%
\pgftext[x=0.654125in,y=0.516267in,right,base]{\color{textcolor}\setmainfont{Lato}\rmfamily\fontsize{8.000000}{9.600000}\selectfont unch.}%
\end{pgfscope}%
\begin{pgfscope}%
\definecolor{textcolor}{rgb}{0.000000,0.000000,0.000000}%
\pgfsetstrokecolor{textcolor}%
\pgfsetfillcolor{textcolor}%
\pgftext[x=3.097312in,y=0.516267in,left,base]{\color{textcolor}\setmainfont{Lato}\rmfamily\fontsize{8.000000}{9.600000}\selectfont 3,272,200}%
\end{pgfscope}%
\begin{pgfscope}%
\definecolor{textcolor}{rgb}{0.000000,0.000000,0.000000}%
\pgfsetstrokecolor{textcolor}%
\pgfsetfillcolor{textcolor}%
\pgftext[x=4.180375in,y=0.516267in,right,base]{\color{textcolor}\setmainfont{Lato}\rmfamily\fontsize{8.000000}{9.600000}\selectfont 4.4}%
\end{pgfscope}%
\begin{pgfscope}%
\definecolor{textcolor}{rgb}{0.000000,0.000000,0.000000}%
\pgfsetstrokecolor{textcolor}%
\pgfsetfillcolor{textcolor}%
\pgftext[x=0.906000in,y=0.329600in,left,base]{\color{textcolor}\setmainfont{Lato}\rmfamily\fontsize{8.000000}{9.600000}\selectfont Boston, MA}%
\end{pgfscope}%
\begin{pgfscope}%
\definecolor{textcolor}{rgb}{0.000000,0.000000,0.000000}%
\pgfsetstrokecolor{textcolor}%
\pgfsetfillcolor{textcolor}%
\pgftext[x=1.963875in,y=0.329600in,left,base]{\color{textcolor}\setmainfont{Lato}\rmfamily\fontsize{8.000000}{9.600000}\selectfont 2.7}%
\end{pgfscope}%
\begin{pgfscope}%
\definecolor{textcolor}{rgb}{0.000000,0.000000,0.000000}%
\pgfsetstrokecolor{textcolor}%
\pgfsetfillcolor{textcolor}%
\pgftext[x=2.568375in,y=0.329600in,left,base]{\color{textcolor}\setmainfont{Lato}\rmfamily\fontsize{8.000000}{9.600000}\selectfont 2.2}%
\end{pgfscope}%
\begin{pgfscope}%
\definecolor{textcolor}{rgb}{0.000000,0.000000,0.000000}%
\pgfsetstrokecolor{textcolor}%
\pgfsetfillcolor{textcolor}%
\pgftext[x=0.654125in,y=0.329600in,right,base]{\color{textcolor}\setmainfont{Lato}\rmfamily\fontsize{8.000000}{9.600000}\selectfont +0.5}%
\end{pgfscope}%
\begin{pgfscope}%
\definecolor{textcolor}{rgb}{0.000000,0.000000,0.000000}%
\pgfsetstrokecolor{textcolor}%
\pgfsetfillcolor{textcolor}%
\pgftext[x=3.097312in,y=0.329600in,left,base]{\color{textcolor}\setmainfont{Lato}\rmfamily\fontsize{8.000000}{9.600000}\selectfont 2,805,100}%
\end{pgfscope}%
\begin{pgfscope}%
\definecolor{textcolor}{rgb}{0.000000,0.000000,0.000000}%
\pgfsetstrokecolor{textcolor}%
\pgfsetfillcolor{textcolor}%
\pgftext[x=4.180375in,y=0.329600in,right,base]{\color{textcolor}\setmainfont{Lato}\rmfamily\fontsize{8.000000}{9.600000}\selectfont -1.6}%
\end{pgfscope}%
\begin{pgfscope}%
\definecolor{textcolor}{rgb}{0.000000,0.000000,0.000000}%
\pgfsetstrokecolor{textcolor}%
\pgfsetfillcolor{textcolor}%
\pgftext[x=0.906000in,y=0.142933in,left,base]{\color{textcolor}\setmainfont{Lato}\rmfamily\fontsize{8.000000}{9.600000}\selectfont Phoenix, AZ}%
\end{pgfscope}%
\begin{pgfscope}%
\definecolor{textcolor}{rgb}{0.000000,0.000000,0.000000}%
\pgfsetstrokecolor{textcolor}%
\pgfsetfillcolor{textcolor}%
\pgftext[x=1.963875in,y=0.142933in,left,base]{\color{textcolor}\setmainfont{Lato}\rmfamily\fontsize{8.000000}{9.600000}\selectfont 3.5}%
\end{pgfscope}%
\begin{pgfscope}%
\definecolor{textcolor}{rgb}{0.000000,0.000000,0.000000}%
\pgfsetstrokecolor{textcolor}%
\pgfsetfillcolor{textcolor}%
\pgftext[x=2.568375in,y=0.142933in,left,base]{\color{textcolor}\setmainfont{Lato}\rmfamily\fontsize{8.000000}{9.600000}\selectfont 3.8}%
\end{pgfscope}%
\begin{pgfscope}%
\definecolor{textcolor}{rgb}{0.000000,0.000000,0.000000}%
\pgfsetstrokecolor{textcolor}%
\pgfsetfillcolor{textcolor}%
\pgftext[x=0.654125in,y=0.142933in,right,base]{\color{textcolor}\setmainfont{Lato}\rmfamily\fontsize{8.000000}{9.600000}\selectfont -0.3}%
\end{pgfscope}%
\begin{pgfscope}%
\definecolor{textcolor}{rgb}{0.000000,0.000000,0.000000}%
\pgfsetstrokecolor{textcolor}%
\pgfsetfillcolor{textcolor}%
\pgftext[x=3.097312in,y=0.142933in,left,base]{\color{textcolor}\setmainfont{Lato}\rmfamily\fontsize{8.000000}{9.600000}\selectfont 2,691,000}%
\end{pgfscope}%
\begin{pgfscope}%
\definecolor{textcolor}{rgb}{0.000000,0.000000,0.000000}%
\pgfsetstrokecolor{textcolor}%
\pgfsetfillcolor{textcolor}%
\pgftext[x=4.180375in,y=0.142933in,right,base]{\color{textcolor}\setmainfont{Lato}\rmfamily\fontsize{8.000000}{9.600000}\selectfont 10.5}%
\end{pgfscope}%
\begin{pgfscope}%
\definecolor{textcolor}{rgb}{0.000000,0.000000,0.000000}%
\pgfsetstrokecolor{textcolor}%
\pgfsetfillcolor{textcolor}%
\pgftext[x=0.283182in,y=2.396000in,left,base]{\color{textcolor}\setmainfont{Lato}\rmfamily\fontsize{10.000000}{12.000000}\bfseries\selectfont Largest MSAs:}%
\end{pgfscope}%
\begin{pgfscope}%
\definecolor{textcolor}{rgb}{0.411765,0.411765,0.411765}%
\pgfsetstrokecolor{textcolor}%
\pgfsetfillcolor{textcolor}%
\pgftext[x=0.924318in,y=2.190667in,left,base]{\color{textcolor}\setmainfont{Lato}\rmfamily\fontsize{9.000000}{10.800000}\selectfont Core City \ \ \ \ \ \ \ \ Nov 23  \ \ \ \ \ \  Nov 19 \ \ \ \ \ \ Labor Force \ \ \ \ \ \ Pct Ch*}%
\end{pgfscope}%
\end{pgfpicture}%
\makeatother%
\endgroup%

\vspace{-3mm}

\footnotesize{Source: Bureau of Labor Statistics; \ Full Table: \tbllink{msa_unemp_rate.csv} \\ \**Pct Ch is percent change in labor force from \input{text/unemp_map_date.txt}}
\end{minipage}
\newpage
\begin{minipage}{0.76\textwidth}
\subsubsection*{Unemployment by State}
\small The Bureau of Labor Statistics \href{https://www.bls.gov/lau/}{report} the \textbf{state unemployment rate}--unemployed people as a share of the state labor force--each month, around two weeks after reporting the national unemployment rate. \input{text/state_unemp.txt}
\vspace{1mm}

\normalsize \textbf{Unemployment Rate by State}\\
\footnotesize{\textit{unemployed share of labor force, percent}}\\
\vspace{-2mm}
\hspace{-10mm} %% Creator: Matplotlib, PGF backend
%%
%% To include the figure in your LaTeX document, write
%%   \input{<filename>.pgf}
%%
%% Make sure the required packages are loaded in your preamble
%%   \usepackage{pgf}
%%
%% Also ensure that all the required font packages are loaded; for instance,
%% the lmodern package is sometimes necessary when using math font.
%%   \usepackage{lmodern}
%%
%% Figures using additional raster images can only be included by \input if
%% they are in the same directory as the main LaTeX file. For loading figures
%% from other directories you can use the `import` package
%%   \usepackage{import}
%%
%% and then include the figures with
%%   \import{<path to file>}{<filename>.pgf}
%%
%% Matplotlib used the following preamble
%%   
%%   \usepackage{fontspec}
%%   \setmainfont{DejaVuSerif.ttf}[Path=\detokenize{/home/brian/miniconda3/lib/python3.8/site-packages/matplotlib/mpl-data/fonts/ttf/}]
%%   \setsansfont{DejaVuSans.ttf}[Path=\detokenize{/home/brian/miniconda3/lib/python3.8/site-packages/matplotlib/mpl-data/fonts/ttf/}]
%%   \setmonofont{DejaVuSansMono.ttf}[Path=\detokenize{/home/brian/miniconda3/lib/python3.8/site-packages/matplotlib/mpl-data/fonts/ttf/}]
%%   \makeatletter\@ifpackageloaded{underscore}{}{\usepackage[strings]{underscore}}\makeatother
%%
\begingroup%
\makeatletter%
\begin{pgfpicture}%
\pgfpathrectangle{\pgfpointorigin}{\pgfqpoint{6.487985in}{2.088185in}}%
\pgfusepath{use as bounding box, clip}%
\begin{pgfscope}%
\pgfsetbuttcap%
\pgfsetmiterjoin%
\definecolor{currentfill}{rgb}{1.000000,1.000000,1.000000}%
\pgfsetfillcolor{currentfill}%
\pgfsetlinewidth{0.000000pt}%
\definecolor{currentstroke}{rgb}{1.000000,1.000000,1.000000}%
\pgfsetstrokecolor{currentstroke}%
\pgfsetdash{}{0pt}%
\pgfpathmoveto{\pgfqpoint{0.000000in}{0.000000in}}%
\pgfpathlineto{\pgfqpoint{6.487985in}{0.000000in}}%
\pgfpathlineto{\pgfqpoint{6.487985in}{2.088185in}}%
\pgfpathlineto{\pgfqpoint{0.000000in}{2.088185in}}%
\pgfpathlineto{\pgfqpoint{0.000000in}{0.000000in}}%
\pgfpathclose%
\pgfusepath{fill}%
\end{pgfscope}%
\begin{pgfscope}%
\pgfpathrectangle{\pgfqpoint{0.100000in}{0.100000in}}{\pgfqpoint{2.862985in}{1.832778in}}%
\pgfusepath{clip}%
\pgfsetbuttcap%
\pgfsetmiterjoin%
\definecolor{currentfill}{rgb}{0.524337,0.754325,0.286736}%
\pgfsetfillcolor{currentfill}%
\pgfsetlinewidth{0.000000pt}%
\definecolor{currentstroke}{rgb}{0.000000,0.000000,0.000000}%
\pgfsetstrokecolor{currentstroke}%
\pgfsetstrokeopacity{0.000000}%
\pgfsetdash{}{0pt}%
\pgfpathmoveto{\pgfqpoint{1.055609in}{0.529755in}}%
\pgfpathlineto{\pgfqpoint{1.047380in}{0.529889in}}%
\pgfpathlineto{\pgfqpoint{1.041780in}{0.536088in}}%
\pgfpathlineto{\pgfqpoint{1.037972in}{0.552550in}}%
\pgfpathlineto{\pgfqpoint{1.046822in}{0.554612in}}%
\pgfpathlineto{\pgfqpoint{1.055813in}{0.552394in}}%
\pgfpathlineto{\pgfqpoint{1.064963in}{0.545909in}}%
\pgfpathlineto{\pgfqpoint{1.064959in}{0.536940in}}%
\pgfpathlineto{\pgfqpoint{1.055609in}{0.529755in}}%
\pgfpathclose%
\pgfusepath{fill}%
\end{pgfscope}%
\begin{pgfscope}%
\pgfpathrectangle{\pgfqpoint{0.100000in}{0.100000in}}{\pgfqpoint{2.862985in}{1.832778in}}%
\pgfusepath{clip}%
\pgfsetbuttcap%
\pgfsetmiterjoin%
\definecolor{currentfill}{rgb}{0.524337,0.754325,0.286736}%
\pgfsetfillcolor{currentfill}%
\pgfsetlinewidth{0.000000pt}%
\definecolor{currentstroke}{rgb}{0.000000,0.000000,0.000000}%
\pgfsetstrokecolor{currentstroke}%
\pgfsetstrokeopacity{0.000000}%
\pgfsetdash{}{0pt}%
\pgfpathmoveto{\pgfqpoint{1.105694in}{0.432023in}}%
\pgfpathlineto{\pgfqpoint{1.096488in}{0.435528in}}%
\pgfpathlineto{\pgfqpoint{1.096446in}{0.442404in}}%
\pgfpathlineto{\pgfqpoint{1.085326in}{0.448219in}}%
\pgfpathlineto{\pgfqpoint{1.086680in}{0.462860in}}%
\pgfpathlineto{\pgfqpoint{1.089901in}{0.469769in}}%
\pgfpathlineto{\pgfqpoint{1.096701in}{0.463410in}}%
\pgfpathlineto{\pgfqpoint{1.107760in}{0.464472in}}%
\pgfpathlineto{\pgfqpoint{1.104967in}{0.444236in}}%
\pgfpathlineto{\pgfqpoint{1.105694in}{0.432023in}}%
\pgfpathclose%
\pgfusepath{fill}%
\end{pgfscope}%
\begin{pgfscope}%
\pgfpathrectangle{\pgfqpoint{0.100000in}{0.100000in}}{\pgfqpoint{2.862985in}{1.832778in}}%
\pgfusepath{clip}%
\pgfsetbuttcap%
\pgfsetmiterjoin%
\definecolor{currentfill}{rgb}{0.524337,0.754325,0.286736}%
\pgfsetfillcolor{currentfill}%
\pgfsetlinewidth{0.000000pt}%
\definecolor{currentstroke}{rgb}{0.000000,0.000000,0.000000}%
\pgfsetstrokecolor{currentstroke}%
\pgfsetstrokeopacity{0.000000}%
\pgfsetdash{}{0pt}%
\pgfpathmoveto{\pgfqpoint{1.144333in}{0.386887in}}%
\pgfpathlineto{\pgfqpoint{1.132651in}{0.388876in}}%
\pgfpathlineto{\pgfqpoint{1.126191in}{0.398877in}}%
\pgfpathlineto{\pgfqpoint{1.125279in}{0.407100in}}%
\pgfpathlineto{\pgfqpoint{1.144333in}{0.386887in}}%
\pgfpathclose%
\pgfusepath{fill}%
\end{pgfscope}%
\begin{pgfscope}%
\pgfpathrectangle{\pgfqpoint{0.100000in}{0.100000in}}{\pgfqpoint{2.862985in}{1.832778in}}%
\pgfusepath{clip}%
\pgfsetbuttcap%
\pgfsetmiterjoin%
\definecolor{currentfill}{rgb}{0.524337,0.754325,0.286736}%
\pgfsetfillcolor{currentfill}%
\pgfsetlinewidth{0.000000pt}%
\definecolor{currentstroke}{rgb}{0.000000,0.000000,0.000000}%
\pgfsetstrokecolor{currentstroke}%
\pgfsetstrokeopacity{0.000000}%
\pgfsetdash{}{0pt}%
\pgfpathmoveto{\pgfqpoint{1.120394in}{0.388631in}}%
\pgfpathlineto{\pgfqpoint{1.126519in}{0.383659in}}%
\pgfpathlineto{\pgfqpoint{1.127848in}{0.375110in}}%
\pgfpathlineto{\pgfqpoint{1.115519in}{0.375583in}}%
\pgfpathlineto{\pgfqpoint{1.120394in}{0.388631in}}%
\pgfpathclose%
\pgfusepath{fill}%
\end{pgfscope}%
\begin{pgfscope}%
\pgfpathrectangle{\pgfqpoint{0.100000in}{0.100000in}}{\pgfqpoint{2.862985in}{1.832778in}}%
\pgfusepath{clip}%
\pgfsetbuttcap%
\pgfsetmiterjoin%
\definecolor{currentfill}{rgb}{0.524337,0.754325,0.286736}%
\pgfsetfillcolor{currentfill}%
\pgfsetlinewidth{0.000000pt}%
\definecolor{currentstroke}{rgb}{0.000000,0.000000,0.000000}%
\pgfsetstrokecolor{currentstroke}%
\pgfsetstrokeopacity{0.000000}%
\pgfsetdash{}{0pt}%
\pgfpathmoveto{\pgfqpoint{1.148521in}{0.340880in}}%
\pgfpathlineto{\pgfqpoint{1.136217in}{0.345303in}}%
\pgfpathlineto{\pgfqpoint{1.139744in}{0.356026in}}%
\pgfpathlineto{\pgfqpoint{1.134586in}{0.366952in}}%
\pgfpathlineto{\pgfqpoint{1.136206in}{0.374463in}}%
\pgfpathlineto{\pgfqpoint{1.146107in}{0.377326in}}%
\pgfpathlineto{\pgfqpoint{1.145916in}{0.365292in}}%
\pgfpathlineto{\pgfqpoint{1.157627in}{0.358618in}}%
\pgfpathlineto{\pgfqpoint{1.164326in}{0.342480in}}%
\pgfpathlineto{\pgfqpoint{1.156861in}{0.336266in}}%
\pgfpathlineto{\pgfqpoint{1.148521in}{0.340880in}}%
\pgfpathclose%
\pgfusepath{fill}%
\end{pgfscope}%
\begin{pgfscope}%
\pgfpathrectangle{\pgfqpoint{0.100000in}{0.100000in}}{\pgfqpoint{2.862985in}{1.832778in}}%
\pgfusepath{clip}%
\pgfsetbuttcap%
\pgfsetmiterjoin%
\definecolor{currentfill}{rgb}{0.524337,0.754325,0.286736}%
\pgfsetfillcolor{currentfill}%
\pgfsetlinewidth{0.000000pt}%
\definecolor{currentstroke}{rgb}{0.000000,0.000000,0.000000}%
\pgfsetstrokecolor{currentstroke}%
\pgfsetstrokeopacity{0.000000}%
\pgfsetdash{}{0pt}%
\pgfpathmoveto{\pgfqpoint{1.104084in}{0.231768in}}%
\pgfpathlineto{\pgfqpoint{1.098222in}{0.243877in}}%
\pgfpathlineto{\pgfqpoint{1.100818in}{0.252022in}}%
\pgfpathlineto{\pgfqpoint{1.112315in}{0.263296in}}%
\pgfpathlineto{\pgfqpoint{1.113404in}{0.271938in}}%
\pgfpathlineto{\pgfqpoint{1.120928in}{0.291072in}}%
\pgfpathlineto{\pgfqpoint{1.140718in}{0.294341in}}%
\pgfpathlineto{\pgfqpoint{1.144506in}{0.305541in}}%
\pgfpathlineto{\pgfqpoint{1.153232in}{0.301341in}}%
\pgfpathlineto{\pgfqpoint{1.170961in}{0.269908in}}%
\pgfpathlineto{\pgfqpoint{1.171185in}{0.259645in}}%
\pgfpathlineto{\pgfqpoint{1.166173in}{0.250763in}}%
\pgfpathlineto{\pgfqpoint{1.172532in}{0.230369in}}%
\pgfpathlineto{\pgfqpoint{1.168089in}{0.227081in}}%
\pgfpathlineto{\pgfqpoint{1.152955in}{0.228037in}}%
\pgfpathlineto{\pgfqpoint{1.134756in}{0.234119in}}%
\pgfpathlineto{\pgfqpoint{1.119460in}{0.236631in}}%
\pgfpathlineto{\pgfqpoint{1.115375in}{0.233063in}}%
\pgfpathlineto{\pgfqpoint{1.104084in}{0.231768in}}%
\pgfpathclose%
\pgfusepath{fill}%
\end{pgfscope}%
\begin{pgfscope}%
\pgfpathrectangle{\pgfqpoint{0.100000in}{0.100000in}}{\pgfqpoint{2.862985in}{1.832778in}}%
\pgfusepath{clip}%
\pgfsetbuttcap%
\pgfsetmiterjoin%
\definecolor{currentfill}{rgb}{0.824145,0.926951,0.690504}%
\pgfsetfillcolor{currentfill}%
\pgfsetlinewidth{0.000000pt}%
\definecolor{currentstroke}{rgb}{0.000000,0.000000,0.000000}%
\pgfsetstrokecolor{currentstroke}%
\pgfsetstrokeopacity{0.000000}%
\pgfsetdash{}{0pt}%
\pgfpathmoveto{\pgfqpoint{0.780273in}{1.824988in}}%
\pgfpathlineto{\pgfqpoint{0.765426in}{1.766159in}}%
\pgfpathlineto{\pgfqpoint{0.754869in}{1.724007in}}%
\pgfpathlineto{\pgfqpoint{0.742317in}{1.673299in}}%
\pgfpathlineto{\pgfqpoint{0.739852in}{1.660801in}}%
\pgfpathlineto{\pgfqpoint{0.741462in}{1.649338in}}%
\pgfpathlineto{\pgfqpoint{0.739329in}{1.638779in}}%
\pgfpathlineto{\pgfqpoint{0.651709in}{1.661643in}}%
\pgfpathlineto{\pgfqpoint{0.643794in}{1.658987in}}%
\pgfpathlineto{\pgfqpoint{0.637021in}{1.661292in}}%
\pgfpathlineto{\pgfqpoint{0.605660in}{1.662045in}}%
\pgfpathlineto{\pgfqpoint{0.595053in}{1.659055in}}%
\pgfpathlineto{\pgfqpoint{0.584520in}{1.660065in}}%
\pgfpathlineto{\pgfqpoint{0.580055in}{1.664746in}}%
\pgfpathlineto{\pgfqpoint{0.551979in}{1.663725in}}%
\pgfpathlineto{\pgfqpoint{0.547545in}{1.671374in}}%
\pgfpathlineto{\pgfqpoint{0.539127in}{1.675404in}}%
\pgfpathlineto{\pgfqpoint{0.526880in}{1.677860in}}%
\pgfpathlineto{\pgfqpoint{0.505740in}{1.674264in}}%
\pgfpathlineto{\pgfqpoint{0.497937in}{1.677787in}}%
\pgfpathlineto{\pgfqpoint{0.485905in}{1.687156in}}%
\pgfpathlineto{\pgfqpoint{0.489522in}{1.705519in}}%
\pgfpathlineto{\pgfqpoint{0.488254in}{1.714897in}}%
\pgfpathlineto{\pgfqpoint{0.478718in}{1.724852in}}%
\pgfpathlineto{\pgfqpoint{0.472610in}{1.724234in}}%
\pgfpathlineto{\pgfqpoint{0.468229in}{1.734304in}}%
\pgfpathlineto{\pgfqpoint{0.457849in}{1.738374in}}%
\pgfpathlineto{\pgfqpoint{0.450305in}{1.737856in}}%
\pgfpathlineto{\pgfqpoint{0.447864in}{1.748198in}}%
\pgfpathlineto{\pgfqpoint{0.455380in}{1.747098in}}%
\pgfpathlineto{\pgfqpoint{0.454783in}{1.761472in}}%
\pgfpathlineto{\pgfqpoint{0.451507in}{1.770008in}}%
\pgfpathlineto{\pgfqpoint{0.452951in}{1.780932in}}%
\pgfpathlineto{\pgfqpoint{0.456884in}{1.789149in}}%
\pgfpathlineto{\pgfqpoint{0.456662in}{1.804741in}}%
\pgfpathlineto{\pgfqpoint{0.454563in}{1.810398in}}%
\pgfpathlineto{\pgfqpoint{0.458188in}{1.828548in}}%
\pgfpathlineto{\pgfqpoint{0.457142in}{1.840246in}}%
\pgfpathlineto{\pgfqpoint{0.453517in}{1.845854in}}%
\pgfpathlineto{\pgfqpoint{0.454070in}{1.864192in}}%
\pgfpathlineto{\pgfqpoint{0.459342in}{1.877678in}}%
\pgfpathlineto{\pgfqpoint{0.465077in}{1.874378in}}%
\pgfpathlineto{\pgfqpoint{0.484064in}{1.854764in}}%
\pgfpathlineto{\pgfqpoint{0.507071in}{1.844031in}}%
\pgfpathlineto{\pgfqpoint{0.518852in}{1.842761in}}%
\pgfpathlineto{\pgfqpoint{0.529951in}{1.838191in}}%
\pgfpathlineto{\pgfqpoint{0.533101in}{1.822814in}}%
\pgfpathlineto{\pgfqpoint{0.523375in}{1.819901in}}%
\pgfpathlineto{\pgfqpoint{0.516929in}{1.807782in}}%
\pgfpathlineto{\pgfqpoint{0.524497in}{1.808455in}}%
\pgfpathlineto{\pgfqpoint{0.527549in}{1.813895in}}%
\pgfpathlineto{\pgfqpoint{0.538221in}{1.820697in}}%
\pgfpathlineto{\pgfqpoint{0.537670in}{1.810596in}}%
\pgfpathlineto{\pgfqpoint{0.530544in}{1.808967in}}%
\pgfpathlineto{\pgfqpoint{0.531727in}{1.795983in}}%
\pgfpathlineto{\pgfqpoint{0.523050in}{1.782976in}}%
\pgfpathlineto{\pgfqpoint{0.517496in}{1.789161in}}%
\pgfpathlineto{\pgfqpoint{0.500991in}{1.785798in}}%
\pgfpathlineto{\pgfqpoint{0.500123in}{1.778232in}}%
\pgfpathlineto{\pgfqpoint{0.505813in}{1.773622in}}%
\pgfpathlineto{\pgfqpoint{0.514493in}{1.773169in}}%
\pgfpathlineto{\pgfqpoint{0.526495in}{1.783019in}}%
\pgfpathlineto{\pgfqpoint{0.535994in}{1.783837in}}%
\pgfpathlineto{\pgfqpoint{0.536358in}{1.794741in}}%
\pgfpathlineto{\pgfqpoint{0.541254in}{1.810756in}}%
\pgfpathlineto{\pgfqpoint{0.552856in}{1.822666in}}%
\pgfpathlineto{\pgfqpoint{0.548988in}{1.831891in}}%
\pgfpathlineto{\pgfqpoint{0.551630in}{1.841823in}}%
\pgfpathlineto{\pgfqpoint{0.546483in}{1.851722in}}%
\pgfpathlineto{\pgfqpoint{0.553818in}{1.864381in}}%
\pgfpathlineto{\pgfqpoint{0.554917in}{1.871996in}}%
\pgfpathlineto{\pgfqpoint{0.548548in}{1.877001in}}%
\pgfpathlineto{\pgfqpoint{0.549609in}{1.889804in}}%
\pgfpathlineto{\pgfqpoint{0.625918in}{1.866620in}}%
\pgfpathlineto{\pgfqpoint{0.706952in}{1.843896in}}%
\pgfpathlineto{\pgfqpoint{0.780273in}{1.824988in}}%
\pgfpathclose%
\pgfusepath{fill}%
\end{pgfscope}%
\begin{pgfscope}%
\pgfpathrectangle{\pgfqpoint{0.100000in}{0.100000in}}{\pgfqpoint{2.862985in}{1.832778in}}%
\pgfusepath{clip}%
\pgfsetbuttcap%
\pgfsetmiterjoin%
\definecolor{currentfill}{rgb}{0.824145,0.926951,0.690504}%
\pgfsetfillcolor{currentfill}%
\pgfsetlinewidth{0.000000pt}%
\definecolor{currentstroke}{rgb}{0.000000,0.000000,0.000000}%
\pgfsetstrokecolor{currentstroke}%
\pgfsetstrokeopacity{0.000000}%
\pgfsetdash{}{0pt}%
\pgfpathmoveto{\pgfqpoint{0.537056in}{1.845512in}}%
\pgfpathlineto{\pgfqpoint{0.540813in}{1.831026in}}%
\pgfpathlineto{\pgfqpoint{0.546706in}{1.826214in}}%
\pgfpathlineto{\pgfqpoint{0.542004in}{1.820089in}}%
\pgfpathlineto{\pgfqpoint{0.537421in}{1.829064in}}%
\pgfpathlineto{\pgfqpoint{0.537056in}{1.845512in}}%
\pgfpathclose%
\pgfusepath{fill}%
\end{pgfscope}%
\begin{pgfscope}%
\pgfpathrectangle{\pgfqpoint{0.100000in}{0.100000in}}{\pgfqpoint{2.862985in}{1.832778in}}%
\pgfusepath{clip}%
\pgfsetbuttcap%
\pgfsetmiterjoin%
\definecolor{currentfill}{rgb}{0.541869,0.765705,0.307958}%
\pgfsetfillcolor{currentfill}%
\pgfsetlinewidth{0.000000pt}%
\definecolor{currentstroke}{rgb}{0.000000,0.000000,0.000000}%
\pgfsetstrokecolor{currentstroke}%
\pgfsetstrokeopacity{0.000000}%
\pgfsetdash{}{0pt}%
\pgfpathmoveto{\pgfqpoint{0.819624in}{1.815480in}}%
\pgfpathlineto{\pgfqpoint{0.901194in}{1.797175in}}%
\pgfpathlineto{\pgfqpoint{0.978016in}{1.781647in}}%
\pgfpathlineto{\pgfqpoint{1.037121in}{1.770817in}}%
\pgfpathlineto{\pgfqpoint{1.088650in}{1.762160in}}%
\pgfpathlineto{\pgfqpoint{1.140293in}{1.754214in}}%
\pgfpathlineto{\pgfqpoint{1.184272in}{1.748018in}}%
\pgfpathlineto{\pgfqpoint{1.228320in}{1.742337in}}%
\pgfpathlineto{\pgfqpoint{1.272432in}{1.737171in}}%
\pgfpathlineto{\pgfqpoint{1.314001in}{1.732781in}}%
\pgfpathlineto{\pgfqpoint{1.308284in}{1.669428in}}%
\pgfpathlineto{\pgfqpoint{1.299775in}{1.583815in}}%
\pgfpathlineto{\pgfqpoint{1.295312in}{1.539783in}}%
\pgfpathlineto{\pgfqpoint{1.288893in}{1.480441in}}%
\pgfpathlineto{\pgfqpoint{1.243453in}{1.485403in}}%
\pgfpathlineto{\pgfqpoint{1.191429in}{1.491308in}}%
\pgfpathlineto{\pgfqpoint{1.119194in}{1.501195in}}%
\pgfpathlineto{\pgfqpoint{1.086931in}{1.505797in}}%
\pgfpathlineto{\pgfqpoint{1.007445in}{1.518274in}}%
\pgfpathlineto{\pgfqpoint{0.980083in}{1.523280in}}%
\pgfpathlineto{\pgfqpoint{0.974380in}{1.490854in}}%
\pgfpathlineto{\pgfqpoint{0.971262in}{1.493168in}}%
\pgfpathlineto{\pgfqpoint{0.965403in}{1.508767in}}%
\pgfpathlineto{\pgfqpoint{0.957871in}{1.507806in}}%
\pgfpathlineto{\pgfqpoint{0.952480in}{1.499920in}}%
\pgfpathlineto{\pgfqpoint{0.941293in}{1.498866in}}%
\pgfpathlineto{\pgfqpoint{0.939045in}{1.502638in}}%
\pgfpathlineto{\pgfqpoint{0.928566in}{1.502181in}}%
\pgfpathlineto{\pgfqpoint{0.923297in}{1.505835in}}%
\pgfpathlineto{\pgfqpoint{0.915931in}{1.500183in}}%
\pgfpathlineto{\pgfqpoint{0.898038in}{1.505269in}}%
\pgfpathlineto{\pgfqpoint{0.890230in}{1.502538in}}%
\pgfpathlineto{\pgfqpoint{0.886029in}{1.514658in}}%
\pgfpathlineto{\pgfqpoint{0.885833in}{1.526190in}}%
\pgfpathlineto{\pgfqpoint{0.873499in}{1.534484in}}%
\pgfpathlineto{\pgfqpoint{0.875510in}{1.546436in}}%
\pgfpathlineto{\pgfqpoint{0.866728in}{1.566180in}}%
\pgfpathlineto{\pgfqpoint{0.867319in}{1.582964in}}%
\pgfpathlineto{\pgfqpoint{0.859772in}{1.591178in}}%
\pgfpathlineto{\pgfqpoint{0.852690in}{1.583911in}}%
\pgfpathlineto{\pgfqpoint{0.843050in}{1.579980in}}%
\pgfpathlineto{\pgfqpoint{0.835712in}{1.588082in}}%
\pgfpathlineto{\pgfqpoint{0.836905in}{1.601019in}}%
\pgfpathlineto{\pgfqpoint{0.845705in}{1.606116in}}%
\pgfpathlineto{\pgfqpoint{0.843366in}{1.614312in}}%
\pgfpathlineto{\pgfqpoint{0.858663in}{1.653985in}}%
\pgfpathlineto{\pgfqpoint{0.846601in}{1.654979in}}%
\pgfpathlineto{\pgfqpoint{0.845361in}{1.662098in}}%
\pgfpathlineto{\pgfqpoint{0.836575in}{1.668238in}}%
\pgfpathlineto{\pgfqpoint{0.837137in}{1.675101in}}%
\pgfpathlineto{\pgfqpoint{0.832502in}{1.680433in}}%
\pgfpathlineto{\pgfqpoint{0.824957in}{1.699791in}}%
\pgfpathlineto{\pgfqpoint{0.818063in}{1.703739in}}%
\pgfpathlineto{\pgfqpoint{0.813126in}{1.720446in}}%
\pgfpathlineto{\pgfqpoint{0.814291in}{1.731551in}}%
\pgfpathlineto{\pgfqpoint{0.805080in}{1.752004in}}%
\pgfpathlineto{\pgfqpoint{0.819624in}{1.815480in}}%
\pgfpathclose%
\pgfusepath{fill}%
\end{pgfscope}%
\begin{pgfscope}%
\pgfpathrectangle{\pgfqpoint{0.100000in}{0.100000in}}{\pgfqpoint{2.862985in}{1.832778in}}%
\pgfusepath{clip}%
\pgfsetbuttcap%
\pgfsetmiterjoin%
\definecolor{currentfill}{rgb}{0.594464,0.799846,0.371626}%
\pgfsetfillcolor{currentfill}%
\pgfsetlinewidth{0.000000pt}%
\definecolor{currentstroke}{rgb}{0.000000,0.000000,0.000000}%
\pgfsetstrokecolor{currentstroke}%
\pgfsetstrokeopacity{0.000000}%
\pgfsetdash{}{0pt}%
\pgfpathmoveto{\pgfqpoint{2.788801in}{1.489012in}}%
\pgfpathlineto{\pgfqpoint{2.787168in}{1.495954in}}%
\pgfpathlineto{\pgfqpoint{2.778094in}{1.502035in}}%
\pgfpathlineto{\pgfqpoint{2.754117in}{1.580468in}}%
\pgfpathlineto{\pgfqpoint{2.741192in}{1.618318in}}%
\pgfpathlineto{\pgfqpoint{2.751583in}{1.629222in}}%
\pgfpathlineto{\pgfqpoint{2.761683in}{1.655908in}}%
\pgfpathlineto{\pgfqpoint{2.766720in}{1.664437in}}%
\pgfpathlineto{\pgfqpoint{2.763152in}{1.668018in}}%
\pgfpathlineto{\pgfqpoint{2.761957in}{1.691707in}}%
\pgfpathlineto{\pgfqpoint{2.766463in}{1.698969in}}%
\pgfpathlineto{\pgfqpoint{2.764488in}{1.715849in}}%
\pgfpathlineto{\pgfqpoint{2.782474in}{1.771228in}}%
\pgfpathlineto{\pgfqpoint{2.790578in}{1.771530in}}%
\pgfpathlineto{\pgfqpoint{2.793942in}{1.761638in}}%
\pgfpathlineto{\pgfqpoint{2.801144in}{1.758798in}}%
\pgfpathlineto{\pgfqpoint{2.814749in}{1.770433in}}%
\pgfpathlineto{\pgfqpoint{2.825417in}{1.777306in}}%
\pgfpathlineto{\pgfqpoint{2.848922in}{1.765204in}}%
\pgfpathlineto{\pgfqpoint{2.870157in}{1.698008in}}%
\pgfpathlineto{\pgfqpoint{2.874194in}{1.681476in}}%
\pgfpathlineto{\pgfqpoint{2.883495in}{1.679537in}}%
\pgfpathlineto{\pgfqpoint{2.896110in}{1.668309in}}%
\pgfpathlineto{\pgfqpoint{2.895396in}{1.661767in}}%
\pgfpathlineto{\pgfqpoint{2.903989in}{1.654013in}}%
\pgfpathlineto{\pgfqpoint{2.910970in}{1.658462in}}%
\pgfpathlineto{\pgfqpoint{2.925577in}{1.641426in}}%
\pgfpathlineto{\pgfqpoint{2.919049in}{1.627792in}}%
\pgfpathlineto{\pgfqpoint{2.910294in}{1.627522in}}%
\pgfpathlineto{\pgfqpoint{2.903271in}{1.615351in}}%
\pgfpathlineto{\pgfqpoint{2.894762in}{1.613625in}}%
\pgfpathlineto{\pgfqpoint{2.888516in}{1.607145in}}%
\pgfpathlineto{\pgfqpoint{2.869860in}{1.600195in}}%
\pgfpathlineto{\pgfqpoint{2.858568in}{1.588995in}}%
\pgfpathlineto{\pgfqpoint{2.852832in}{1.597034in}}%
\pgfpathlineto{\pgfqpoint{2.847620in}{1.591272in}}%
\pgfpathlineto{\pgfqpoint{2.849047in}{1.567979in}}%
\pgfpathlineto{\pgfqpoint{2.844920in}{1.558753in}}%
\pgfpathlineto{\pgfqpoint{2.835918in}{1.561273in}}%
\pgfpathlineto{\pgfqpoint{2.834465in}{1.551814in}}%
\pgfpathlineto{\pgfqpoint{2.830576in}{1.547924in}}%
\pgfpathlineto{\pgfqpoint{2.822791in}{1.548872in}}%
\pgfpathlineto{\pgfqpoint{2.823281in}{1.540028in}}%
\pgfpathlineto{\pgfqpoint{2.811482in}{1.542540in}}%
\pgfpathlineto{\pgfqpoint{2.805038in}{1.530287in}}%
\pgfpathlineto{\pgfqpoint{2.807472in}{1.523875in}}%
\pgfpathlineto{\pgfqpoint{2.803639in}{1.513212in}}%
\pgfpathlineto{\pgfqpoint{2.797604in}{1.505369in}}%
\pgfpathlineto{\pgfqpoint{2.796086in}{1.488995in}}%
\pgfpathlineto{\pgfqpoint{2.788801in}{1.489012in}}%
\pgfpathclose%
\pgfusepath{fill}%
\end{pgfscope}%
\begin{pgfscope}%
\pgfpathrectangle{\pgfqpoint{0.100000in}{0.100000in}}{\pgfqpoint{2.862985in}{1.832778in}}%
\pgfusepath{clip}%
\pgfsetbuttcap%
\pgfsetmiterjoin%
\definecolor{currentfill}{rgb}{0.594464,0.799846,0.371626}%
\pgfsetfillcolor{currentfill}%
\pgfsetlinewidth{0.000000pt}%
\definecolor{currentstroke}{rgb}{0.000000,0.000000,0.000000}%
\pgfsetstrokecolor{currentstroke}%
\pgfsetstrokeopacity{0.000000}%
\pgfsetdash{}{0pt}%
\pgfpathmoveto{\pgfqpoint{2.873211in}{1.595412in}}%
\pgfpathlineto{\pgfqpoint{2.878538in}{1.600991in}}%
\pgfpathlineto{\pgfqpoint{2.883609in}{1.595747in}}%
\pgfpathlineto{\pgfqpoint{2.874500in}{1.588798in}}%
\pgfpathlineto{\pgfqpoint{2.873211in}{1.595412in}}%
\pgfpathclose%
\pgfusepath{fill}%
\end{pgfscope}%
\begin{pgfscope}%
\pgfpathrectangle{\pgfqpoint{0.100000in}{0.100000in}}{\pgfqpoint{2.862985in}{1.832778in}}%
\pgfusepath{clip}%
\pgfsetbuttcap%
\pgfsetmiterjoin%
\definecolor{currentfill}{rgb}{0.217224,0.469973,0.111572}%
\pgfsetfillcolor{currentfill}%
\pgfsetlinewidth{0.000000pt}%
\definecolor{currentstroke}{rgb}{0.000000,0.000000,0.000000}%
\pgfsetstrokecolor{currentstroke}%
\pgfsetstrokeopacity{0.000000}%
\pgfsetdash{}{0pt}%
\pgfpathmoveto{\pgfqpoint{1.295312in}{1.539783in}}%
\pgfpathlineto{\pgfqpoint{1.299775in}{1.583815in}}%
\pgfpathlineto{\pgfqpoint{1.308284in}{1.669428in}}%
\pgfpathlineto{\pgfqpoint{1.314001in}{1.732781in}}%
\pgfpathlineto{\pgfqpoint{1.360821in}{1.728389in}}%
\pgfpathlineto{\pgfqpoint{1.420718in}{1.723620in}}%
\pgfpathlineto{\pgfqpoint{1.475466in}{1.720093in}}%
\pgfpathlineto{\pgfqpoint{1.525039in}{1.717585in}}%
\pgfpathlineto{\pgfqpoint{1.599019in}{1.715047in}}%
\pgfpathlineto{\pgfqpoint{1.604068in}{1.694197in}}%
\pgfpathlineto{\pgfqpoint{1.602261in}{1.684232in}}%
\pgfpathlineto{\pgfqpoint{1.601666in}{1.663751in}}%
\pgfpathlineto{\pgfqpoint{1.605101in}{1.648413in}}%
\pgfpathlineto{\pgfqpoint{1.612970in}{1.625784in}}%
\pgfpathlineto{\pgfqpoint{1.612952in}{1.597294in}}%
\pgfpathlineto{\pgfqpoint{1.614408in}{1.564197in}}%
\pgfpathlineto{\pgfqpoint{1.616408in}{1.555277in}}%
\pgfpathlineto{\pgfqpoint{1.622257in}{1.545485in}}%
\pgfpathlineto{\pgfqpoint{1.624186in}{1.530226in}}%
\pgfpathlineto{\pgfqpoint{1.623353in}{1.520039in}}%
\pgfpathlineto{\pgfqpoint{1.561333in}{1.521383in}}%
\pgfpathlineto{\pgfqpoint{1.516215in}{1.523670in}}%
\pgfpathlineto{\pgfqpoint{1.450074in}{1.527191in}}%
\pgfpathlineto{\pgfqpoint{1.384849in}{1.531824in}}%
\pgfpathlineto{\pgfqpoint{1.341407in}{1.535329in}}%
\pgfpathlineto{\pgfqpoint{1.295312in}{1.539783in}}%
\pgfpathclose%
\pgfusepath{fill}%
\end{pgfscope}%
\begin{pgfscope}%
\pgfpathrectangle{\pgfqpoint{0.100000in}{0.100000in}}{\pgfqpoint{2.862985in}{1.832778in}}%
\pgfusepath{clip}%
\pgfsetbuttcap%
\pgfsetmiterjoin%
\definecolor{currentfill}{rgb}{0.263975,0.526567,0.121415}%
\pgfsetfillcolor{currentfill}%
\pgfsetlinewidth{0.000000pt}%
\definecolor{currentstroke}{rgb}{0.000000,0.000000,0.000000}%
\pgfsetstrokecolor{currentstroke}%
\pgfsetstrokeopacity{0.000000}%
\pgfsetdash{}{0pt}%
\pgfpathmoveto{\pgfqpoint{1.276565in}{1.355506in}}%
\pgfpathlineto{\pgfqpoint{1.283725in}{1.429317in}}%
\pgfpathlineto{\pgfqpoint{1.288893in}{1.480441in}}%
\pgfpathlineto{\pgfqpoint{1.295312in}{1.539783in}}%
\pgfpathlineto{\pgfqpoint{1.341407in}{1.535329in}}%
\pgfpathlineto{\pgfqpoint{1.384849in}{1.531824in}}%
\pgfpathlineto{\pgfqpoint{1.450074in}{1.527191in}}%
\pgfpathlineto{\pgfqpoint{1.516215in}{1.523670in}}%
\pgfpathlineto{\pgfqpoint{1.561333in}{1.521383in}}%
\pgfpathlineto{\pgfqpoint{1.623353in}{1.520039in}}%
\pgfpathlineto{\pgfqpoint{1.619154in}{1.507786in}}%
\pgfpathlineto{\pgfqpoint{1.610772in}{1.498167in}}%
\pgfpathlineto{\pgfqpoint{1.617191in}{1.487092in}}%
\pgfpathlineto{\pgfqpoint{1.624268in}{1.484725in}}%
\pgfpathlineto{\pgfqpoint{1.627626in}{1.478363in}}%
\pgfpathlineto{\pgfqpoint{1.626855in}{1.431914in}}%
\pgfpathlineto{\pgfqpoint{1.625565in}{1.366544in}}%
\pgfpathlineto{\pgfqpoint{1.620791in}{1.349375in}}%
\pgfpathlineto{\pgfqpoint{1.625061in}{1.339878in}}%
\pgfpathlineto{\pgfqpoint{1.616489in}{1.319470in}}%
\pgfpathlineto{\pgfqpoint{1.625519in}{1.303018in}}%
\pgfpathlineto{\pgfqpoint{1.617847in}{1.304277in}}%
\pgfpathlineto{\pgfqpoint{1.612612in}{1.314513in}}%
\pgfpathlineto{\pgfqpoint{1.593930in}{1.321529in}}%
\pgfpathlineto{\pgfqpoint{1.582148in}{1.327701in}}%
\pgfpathlineto{\pgfqpoint{1.562377in}{1.328214in}}%
\pgfpathlineto{\pgfqpoint{1.555515in}{1.322591in}}%
\pgfpathlineto{\pgfqpoint{1.533137in}{1.333663in}}%
\pgfpathlineto{\pgfqpoint{1.531420in}{1.337164in}}%
\pgfpathlineto{\pgfqpoint{1.453322in}{1.340860in}}%
\pgfpathlineto{\pgfqpoint{1.405874in}{1.343580in}}%
\pgfpathlineto{\pgfqpoint{1.366682in}{1.346649in}}%
\pgfpathlineto{\pgfqpoint{1.301922in}{1.352768in}}%
\pgfpathlineto{\pgfqpoint{1.276565in}{1.355506in}}%
\pgfpathclose%
\pgfusepath{fill}%
\end{pgfscope}%
\begin{pgfscope}%
\pgfpathrectangle{\pgfqpoint{0.100000in}{0.100000in}}{\pgfqpoint{2.862985in}{1.832778in}}%
\pgfusepath{clip}%
\pgfsetbuttcap%
\pgfsetmiterjoin%
\definecolor{currentfill}{rgb}{0.541869,0.765705,0.307958}%
\pgfsetfillcolor{currentfill}%
\pgfsetlinewidth{0.000000pt}%
\definecolor{currentstroke}{rgb}{0.000000,0.000000,0.000000}%
\pgfsetstrokecolor{currentstroke}%
\pgfsetstrokeopacity{0.000000}%
\pgfsetdash{}{0pt}%
\pgfpathmoveto{\pgfqpoint{1.264283in}{1.230524in}}%
\pgfpathlineto{\pgfqpoint{1.222647in}{1.234326in}}%
\pgfpathlineto{\pgfqpoint{1.131892in}{1.245527in}}%
\pgfpathlineto{\pgfqpoint{1.082522in}{1.252613in}}%
\pgfpathlineto{\pgfqpoint{1.029571in}{1.260214in}}%
\pgfpathlineto{\pgfqpoint{0.984967in}{1.267352in}}%
\pgfpathlineto{\pgfqpoint{0.936009in}{1.275734in}}%
\pgfpathlineto{\pgfqpoint{0.947141in}{1.337483in}}%
\pgfpathlineto{\pgfqpoint{0.958416in}{1.400800in}}%
\pgfpathlineto{\pgfqpoint{0.974380in}{1.490854in}}%
\pgfpathlineto{\pgfqpoint{0.980083in}{1.523280in}}%
\pgfpathlineto{\pgfqpoint{1.007445in}{1.518274in}}%
\pgfpathlineto{\pgfqpoint{1.086931in}{1.505797in}}%
\pgfpathlineto{\pgfqpoint{1.119194in}{1.501195in}}%
\pgfpathlineto{\pgfqpoint{1.191429in}{1.491308in}}%
\pgfpathlineto{\pgfqpoint{1.243453in}{1.485403in}}%
\pgfpathlineto{\pgfqpoint{1.288893in}{1.480441in}}%
\pgfpathlineto{\pgfqpoint{1.283725in}{1.429317in}}%
\pgfpathlineto{\pgfqpoint{1.276565in}{1.355506in}}%
\pgfpathlineto{\pgfqpoint{1.270419in}{1.292777in}}%
\pgfpathlineto{\pgfqpoint{1.264283in}{1.230524in}}%
\pgfpathclose%
\pgfusepath{fill}%
\end{pgfscope}%
\begin{pgfscope}%
\pgfpathrectangle{\pgfqpoint{0.100000in}{0.100000in}}{\pgfqpoint{2.862985in}{1.832778in}}%
\pgfusepath{clip}%
\pgfsetbuttcap%
\pgfsetmiterjoin%
\definecolor{currentfill}{rgb}{0.474971,0.717878,0.240138}%
\pgfsetfillcolor{currentfill}%
\pgfsetlinewidth{0.000000pt}%
\definecolor{currentstroke}{rgb}{0.000000,0.000000,0.000000}%
\pgfsetstrokecolor{currentstroke}%
\pgfsetstrokeopacity{0.000000}%
\pgfsetdash{}{0pt}%
\pgfpathmoveto{\pgfqpoint{2.025520in}{1.315700in}}%
\pgfpathlineto{\pgfqpoint{1.972758in}{1.311952in}}%
\pgfpathlineto{\pgfqpoint{1.894114in}{1.308600in}}%
\pgfpathlineto{\pgfqpoint{1.891125in}{1.316544in}}%
\pgfpathlineto{\pgfqpoint{1.873686in}{1.322483in}}%
\pgfpathlineto{\pgfqpoint{1.869839in}{1.333702in}}%
\pgfpathlineto{\pgfqpoint{1.868228in}{1.347582in}}%
\pgfpathlineto{\pgfqpoint{1.872158in}{1.354693in}}%
\pgfpathlineto{\pgfqpoint{1.865951in}{1.361517in}}%
\pgfpathlineto{\pgfqpoint{1.864454in}{1.369657in}}%
\pgfpathlineto{\pgfqpoint{1.862453in}{1.387663in}}%
\pgfpathlineto{\pgfqpoint{1.856505in}{1.397444in}}%
\pgfpathlineto{\pgfqpoint{1.845940in}{1.402933in}}%
\pgfpathlineto{\pgfqpoint{1.834444in}{1.411992in}}%
\pgfpathlineto{\pgfqpoint{1.828501in}{1.422710in}}%
\pgfpathlineto{\pgfqpoint{1.817838in}{1.427026in}}%
\pgfpathlineto{\pgfqpoint{1.811570in}{1.434051in}}%
\pgfpathlineto{\pgfqpoint{1.803976in}{1.435243in}}%
\pgfpathlineto{\pgfqpoint{1.790419in}{1.445673in}}%
\pgfpathlineto{\pgfqpoint{1.792621in}{1.457675in}}%
\pgfpathlineto{\pgfqpoint{1.792197in}{1.480500in}}%
\pgfpathlineto{\pgfqpoint{1.796378in}{1.486783in}}%
\pgfpathlineto{\pgfqpoint{1.792621in}{1.496270in}}%
\pgfpathlineto{\pgfqpoint{1.786004in}{1.498105in}}%
\pgfpathlineto{\pgfqpoint{1.786550in}{1.506437in}}%
\pgfpathlineto{\pgfqpoint{1.794761in}{1.519600in}}%
\pgfpathlineto{\pgfqpoint{1.811020in}{1.530009in}}%
\pgfpathlineto{\pgfqpoint{1.810007in}{1.567025in}}%
\pgfpathlineto{\pgfqpoint{1.818143in}{1.572586in}}%
\pgfpathlineto{\pgfqpoint{1.825842in}{1.568882in}}%
\pgfpathlineto{\pgfqpoint{1.841526in}{1.574304in}}%
\pgfpathlineto{\pgfqpoint{1.871034in}{1.587932in}}%
\pgfpathlineto{\pgfqpoint{1.874878in}{1.583713in}}%
\pgfpathlineto{\pgfqpoint{1.869289in}{1.564600in}}%
\pgfpathlineto{\pgfqpoint{1.877604in}{1.568794in}}%
\pgfpathlineto{\pgfqpoint{1.891845in}{1.564606in}}%
\pgfpathlineto{\pgfqpoint{1.900596in}{1.561110in}}%
\pgfpathlineto{\pgfqpoint{1.905502in}{1.550866in}}%
\pgfpathlineto{\pgfqpoint{1.950397in}{1.541198in}}%
\pgfpathlineto{\pgfqpoint{1.963816in}{1.534565in}}%
\pgfpathlineto{\pgfqpoint{1.977460in}{1.534643in}}%
\pgfpathlineto{\pgfqpoint{1.991480in}{1.531907in}}%
\pgfpathlineto{\pgfqpoint{2.000578in}{1.522561in}}%
\pgfpathlineto{\pgfqpoint{2.009272in}{1.517943in}}%
\pgfpathlineto{\pgfqpoint{2.010870in}{1.504623in}}%
\pgfpathlineto{\pgfqpoint{2.008301in}{1.496259in}}%
\pgfpathlineto{\pgfqpoint{2.017991in}{1.495641in}}%
\pgfpathlineto{\pgfqpoint{2.014727in}{1.485941in}}%
\pgfpathlineto{\pgfqpoint{2.017832in}{1.482495in}}%
\pgfpathlineto{\pgfqpoint{2.020919in}{1.473339in}}%
\pgfpathlineto{\pgfqpoint{2.011476in}{1.468491in}}%
\pgfpathlineto{\pgfqpoint{2.005993in}{1.454984in}}%
\pgfpathlineto{\pgfqpoint{2.004273in}{1.445433in}}%
\pgfpathlineto{\pgfqpoint{2.009533in}{1.443798in}}%
\pgfpathlineto{\pgfqpoint{2.016268in}{1.450960in}}%
\pgfpathlineto{\pgfqpoint{2.021994in}{1.463372in}}%
\pgfpathlineto{\pgfqpoint{2.029741in}{1.467683in}}%
\pgfpathlineto{\pgfqpoint{2.035563in}{1.462075in}}%
\pgfpathlineto{\pgfqpoint{2.029832in}{1.445104in}}%
\pgfpathlineto{\pgfqpoint{2.028034in}{1.431928in}}%
\pgfpathlineto{\pgfqpoint{2.029728in}{1.422457in}}%
\pgfpathlineto{\pgfqpoint{2.024406in}{1.417099in}}%
\pgfpathlineto{\pgfqpoint{2.021748in}{1.404007in}}%
\pgfpathlineto{\pgfqpoint{2.023920in}{1.390247in}}%
\pgfpathlineto{\pgfqpoint{2.017640in}{1.369897in}}%
\pgfpathlineto{\pgfqpoint{2.017772in}{1.359728in}}%
\pgfpathlineto{\pgfqpoint{2.022734in}{1.337693in}}%
\pgfpathlineto{\pgfqpoint{2.025951in}{1.333913in}}%
\pgfpathlineto{\pgfqpoint{2.025520in}{1.315700in}}%
\pgfpathclose%
\pgfusepath{fill}%
\end{pgfscope}%
\begin{pgfscope}%
\pgfpathrectangle{\pgfqpoint{0.100000in}{0.100000in}}{\pgfqpoint{2.862985in}{1.832778in}}%
\pgfusepath{clip}%
\pgfsetbuttcap%
\pgfsetmiterjoin%
\definecolor{currentfill}{rgb}{0.474971,0.717878,0.240138}%
\pgfsetfillcolor{currentfill}%
\pgfsetlinewidth{0.000000pt}%
\definecolor{currentstroke}{rgb}{0.000000,0.000000,0.000000}%
\pgfsetstrokecolor{currentstroke}%
\pgfsetstrokeopacity{0.000000}%
\pgfsetdash{}{0pt}%
\pgfpathmoveto{\pgfqpoint{2.045302in}{1.494291in}}%
\pgfpathlineto{\pgfqpoint{2.046480in}{1.485502in}}%
\pgfpathlineto{\pgfqpoint{2.035690in}{1.462343in}}%
\pgfpathlineto{\pgfqpoint{2.030896in}{1.469051in}}%
\pgfpathlineto{\pgfqpoint{2.045302in}{1.494291in}}%
\pgfpathclose%
\pgfusepath{fill}%
\end{pgfscope}%
\begin{pgfscope}%
\pgfpathrectangle{\pgfqpoint{0.100000in}{0.100000in}}{\pgfqpoint{2.862985in}{1.832778in}}%
\pgfusepath{clip}%
\pgfsetbuttcap%
\pgfsetmiterjoin%
\definecolor{currentfill}{rgb}{0.568166,0.782776,0.339792}%
\pgfsetfillcolor{currentfill}%
\pgfsetlinewidth{0.000000pt}%
\definecolor{currentstroke}{rgb}{0.000000,0.000000,0.000000}%
\pgfsetstrokecolor{currentstroke}%
\pgfsetstrokeopacity{0.000000}%
\pgfsetdash{}{0pt}%
\pgfpathmoveto{\pgfqpoint{0.739329in}{1.638779in}}%
\pgfpathlineto{\pgfqpoint{0.741462in}{1.649338in}}%
\pgfpathlineto{\pgfqpoint{0.739852in}{1.660801in}}%
\pgfpathlineto{\pgfqpoint{0.742317in}{1.673299in}}%
\pgfpathlineto{\pgfqpoint{0.754869in}{1.724007in}}%
\pgfpathlineto{\pgfqpoint{0.765426in}{1.766159in}}%
\pgfpathlineto{\pgfqpoint{0.780273in}{1.824988in}}%
\pgfpathlineto{\pgfqpoint{0.819624in}{1.815480in}}%
\pgfpathlineto{\pgfqpoint{0.805080in}{1.752004in}}%
\pgfpathlineto{\pgfqpoint{0.814291in}{1.731551in}}%
\pgfpathlineto{\pgfqpoint{0.813126in}{1.720446in}}%
\pgfpathlineto{\pgfqpoint{0.818063in}{1.703739in}}%
\pgfpathlineto{\pgfqpoint{0.824957in}{1.699791in}}%
\pgfpathlineto{\pgfqpoint{0.832502in}{1.680433in}}%
\pgfpathlineto{\pgfqpoint{0.837137in}{1.675101in}}%
\pgfpathlineto{\pgfqpoint{0.836575in}{1.668238in}}%
\pgfpathlineto{\pgfqpoint{0.845361in}{1.662098in}}%
\pgfpathlineto{\pgfqpoint{0.846601in}{1.654979in}}%
\pgfpathlineto{\pgfqpoint{0.858663in}{1.653985in}}%
\pgfpathlineto{\pgfqpoint{0.843366in}{1.614312in}}%
\pgfpathlineto{\pgfqpoint{0.845705in}{1.606116in}}%
\pgfpathlineto{\pgfqpoint{0.836905in}{1.601019in}}%
\pgfpathlineto{\pgfqpoint{0.835712in}{1.588082in}}%
\pgfpathlineto{\pgfqpoint{0.843050in}{1.579980in}}%
\pgfpathlineto{\pgfqpoint{0.852690in}{1.583911in}}%
\pgfpathlineto{\pgfqpoint{0.859772in}{1.591178in}}%
\pgfpathlineto{\pgfqpoint{0.867319in}{1.582964in}}%
\pgfpathlineto{\pgfqpoint{0.866728in}{1.566180in}}%
\pgfpathlineto{\pgfqpoint{0.875510in}{1.546436in}}%
\pgfpathlineto{\pgfqpoint{0.873499in}{1.534484in}}%
\pgfpathlineto{\pgfqpoint{0.885833in}{1.526190in}}%
\pgfpathlineto{\pgfqpoint{0.886029in}{1.514658in}}%
\pgfpathlineto{\pgfqpoint{0.890230in}{1.502538in}}%
\pgfpathlineto{\pgfqpoint{0.898038in}{1.505269in}}%
\pgfpathlineto{\pgfqpoint{0.915931in}{1.500183in}}%
\pgfpathlineto{\pgfqpoint{0.923297in}{1.505835in}}%
\pgfpathlineto{\pgfqpoint{0.928566in}{1.502181in}}%
\pgfpathlineto{\pgfqpoint{0.939045in}{1.502638in}}%
\pgfpathlineto{\pgfqpoint{0.941293in}{1.498866in}}%
\pgfpathlineto{\pgfqpoint{0.952480in}{1.499920in}}%
\pgfpathlineto{\pgfqpoint{0.957871in}{1.507806in}}%
\pgfpathlineto{\pgfqpoint{0.965403in}{1.508767in}}%
\pgfpathlineto{\pgfqpoint{0.971262in}{1.493168in}}%
\pgfpathlineto{\pgfqpoint{0.974380in}{1.490854in}}%
\pgfpathlineto{\pgfqpoint{0.958416in}{1.400800in}}%
\pgfpathlineto{\pgfqpoint{0.947141in}{1.337483in}}%
\pgfpathlineto{\pgfqpoint{0.858212in}{1.354650in}}%
\pgfpathlineto{\pgfqpoint{0.810177in}{1.364205in}}%
\pgfpathlineto{\pgfqpoint{0.765246in}{1.374085in}}%
\pgfpathlineto{\pgfqpoint{0.723250in}{1.383593in}}%
\pgfpathlineto{\pgfqpoint{0.674647in}{1.395314in}}%
\pgfpathlineto{\pgfqpoint{0.699717in}{1.498158in}}%
\pgfpathlineto{\pgfqpoint{0.701069in}{1.505674in}}%
\pgfpathlineto{\pgfqpoint{0.712217in}{1.525368in}}%
\pgfpathlineto{\pgfqpoint{0.700673in}{1.537205in}}%
\pgfpathlineto{\pgfqpoint{0.703070in}{1.548828in}}%
\pgfpathlineto{\pgfqpoint{0.707464in}{1.551738in}}%
\pgfpathlineto{\pgfqpoint{0.715313in}{1.563715in}}%
\pgfpathlineto{\pgfqpoint{0.726752in}{1.572002in}}%
\pgfpathlineto{\pgfqpoint{0.727377in}{1.578133in}}%
\pgfpathlineto{\pgfqpoint{0.734282in}{1.584259in}}%
\pgfpathlineto{\pgfqpoint{0.740015in}{1.595725in}}%
\pgfpathlineto{\pgfqpoint{0.751774in}{1.607856in}}%
\pgfpathlineto{\pgfqpoint{0.750939in}{1.619838in}}%
\pgfpathlineto{\pgfqpoint{0.742899in}{1.626494in}}%
\pgfpathlineto{\pgfqpoint{0.739329in}{1.638779in}}%
\pgfpathclose%
\pgfusepath{fill}%
\end{pgfscope}%
\begin{pgfscope}%
\pgfpathrectangle{\pgfqpoint{0.100000in}{0.100000in}}{\pgfqpoint{2.862985in}{1.832778in}}%
\pgfusepath{clip}%
\pgfsetbuttcap%
\pgfsetmiterjoin%
\definecolor{currentfill}{rgb}{0.321184,0.588697,0.141715}%
\pgfsetfillcolor{currentfill}%
\pgfsetlinewidth{0.000000pt}%
\definecolor{currentstroke}{rgb}{0.000000,0.000000,0.000000}%
\pgfsetstrokecolor{currentstroke}%
\pgfsetstrokeopacity{0.000000}%
\pgfsetdash{}{0pt}%
\pgfpathmoveto{\pgfqpoint{2.686015in}{1.438080in}}%
\pgfpathlineto{\pgfqpoint{2.682320in}{1.449744in}}%
\pgfpathlineto{\pgfqpoint{2.675463in}{1.485127in}}%
\pgfpathlineto{\pgfqpoint{2.666598in}{1.496013in}}%
\pgfpathlineto{\pgfqpoint{2.658822in}{1.515590in}}%
\pgfpathlineto{\pgfqpoint{2.661375in}{1.530104in}}%
\pgfpathlineto{\pgfqpoint{2.659327in}{1.540801in}}%
\pgfpathlineto{\pgfqpoint{2.652683in}{1.551308in}}%
\pgfpathlineto{\pgfqpoint{2.652410in}{1.562882in}}%
\pgfpathlineto{\pgfqpoint{2.648552in}{1.575365in}}%
\pgfpathlineto{\pgfqpoint{2.683077in}{1.583862in}}%
\pgfpathlineto{\pgfqpoint{2.727922in}{1.595966in}}%
\pgfpathlineto{\pgfqpoint{2.729711in}{1.589025in}}%
\pgfpathlineto{\pgfqpoint{2.726856in}{1.577975in}}%
\pgfpathlineto{\pgfqpoint{2.733559in}{1.569133in}}%
\pgfpathlineto{\pgfqpoint{2.729999in}{1.557930in}}%
\pgfpathlineto{\pgfqpoint{2.715852in}{1.543846in}}%
\pgfpathlineto{\pgfqpoint{2.719741in}{1.533267in}}%
\pgfpathlineto{\pgfqpoint{2.716859in}{1.510906in}}%
\pgfpathlineto{\pgfqpoint{2.712824in}{1.498298in}}%
\pgfpathlineto{\pgfqpoint{2.717972in}{1.462444in}}%
\pgfpathlineto{\pgfqpoint{2.717278in}{1.449830in}}%
\pgfpathlineto{\pgfqpoint{2.722258in}{1.445782in}}%
\pgfpathlineto{\pgfqpoint{2.686015in}{1.438080in}}%
\pgfpathclose%
\pgfusepath{fill}%
\end{pgfscope}%
\begin{pgfscope}%
\pgfpathrectangle{\pgfqpoint{0.100000in}{0.100000in}}{\pgfqpoint{2.862985in}{1.832778in}}%
\pgfusepath{clip}%
\pgfsetbuttcap%
\pgfsetmiterjoin%
\definecolor{currentfill}{rgb}{0.299039,0.569012,0.128797}%
\pgfsetfillcolor{currentfill}%
\pgfsetlinewidth{0.000000pt}%
\definecolor{currentstroke}{rgb}{0.000000,0.000000,0.000000}%
\pgfsetstrokecolor{currentstroke}%
\pgfsetstrokeopacity{0.000000}%
\pgfsetdash{}{0pt}%
\pgfpathmoveto{\pgfqpoint{1.625565in}{1.366544in}}%
\pgfpathlineto{\pgfqpoint{1.626855in}{1.431914in}}%
\pgfpathlineto{\pgfqpoint{1.627626in}{1.478363in}}%
\pgfpathlineto{\pgfqpoint{1.624268in}{1.484725in}}%
\pgfpathlineto{\pgfqpoint{1.617191in}{1.487092in}}%
\pgfpathlineto{\pgfqpoint{1.610772in}{1.498167in}}%
\pgfpathlineto{\pgfqpoint{1.619154in}{1.507786in}}%
\pgfpathlineto{\pgfqpoint{1.623353in}{1.520039in}}%
\pgfpathlineto{\pgfqpoint{1.624186in}{1.530226in}}%
\pgfpathlineto{\pgfqpoint{1.622257in}{1.545485in}}%
\pgfpathlineto{\pgfqpoint{1.616408in}{1.555277in}}%
\pgfpathlineto{\pgfqpoint{1.614408in}{1.564197in}}%
\pgfpathlineto{\pgfqpoint{1.612952in}{1.597294in}}%
\pgfpathlineto{\pgfqpoint{1.612970in}{1.625784in}}%
\pgfpathlineto{\pgfqpoint{1.605101in}{1.648413in}}%
\pgfpathlineto{\pgfqpoint{1.601666in}{1.663751in}}%
\pgfpathlineto{\pgfqpoint{1.602261in}{1.684232in}}%
\pgfpathlineto{\pgfqpoint{1.604068in}{1.694197in}}%
\pgfpathlineto{\pgfqpoint{1.599019in}{1.715047in}}%
\pgfpathlineto{\pgfqpoint{1.633393in}{1.714359in}}%
\pgfpathlineto{\pgfqpoint{1.685604in}{1.713911in}}%
\pgfpathlineto{\pgfqpoint{1.685890in}{1.737609in}}%
\pgfpathlineto{\pgfqpoint{1.699180in}{1.735000in}}%
\pgfpathlineto{\pgfqpoint{1.705549in}{1.706104in}}%
\pgfpathlineto{\pgfqpoint{1.710245in}{1.695714in}}%
\pgfpathlineto{\pgfqpoint{1.721922in}{1.695407in}}%
\pgfpathlineto{\pgfqpoint{1.724536in}{1.691882in}}%
\pgfpathlineto{\pgfqpoint{1.740825in}{1.690320in}}%
\pgfpathlineto{\pgfqpoint{1.743562in}{1.683156in}}%
\pgfpathlineto{\pgfqpoint{1.754788in}{1.684767in}}%
\pgfpathlineto{\pgfqpoint{1.763507in}{1.691467in}}%
\pgfpathlineto{\pgfqpoint{1.778542in}{1.691216in}}%
\pgfpathlineto{\pgfqpoint{1.787845in}{1.685828in}}%
\pgfpathlineto{\pgfqpoint{1.788915in}{1.680774in}}%
\pgfpathlineto{\pgfqpoint{1.797767in}{1.679717in}}%
\pgfpathlineto{\pgfqpoint{1.803544in}{1.665932in}}%
\pgfpathlineto{\pgfqpoint{1.807272in}{1.674408in}}%
\pgfpathlineto{\pgfqpoint{1.817441in}{1.674929in}}%
\pgfpathlineto{\pgfqpoint{1.820013in}{1.668328in}}%
\pgfpathlineto{\pgfqpoint{1.831455in}{1.665315in}}%
\pgfpathlineto{\pgfqpoint{1.837765in}{1.655808in}}%
\pgfpathlineto{\pgfqpoint{1.851759in}{1.658760in}}%
\pgfpathlineto{\pgfqpoint{1.867158in}{1.670427in}}%
\pgfpathlineto{\pgfqpoint{1.872772in}{1.660142in}}%
\pgfpathlineto{\pgfqpoint{1.898039in}{1.662957in}}%
\pgfpathlineto{\pgfqpoint{1.908840in}{1.655898in}}%
\pgfpathlineto{\pgfqpoint{1.915126in}{1.658418in}}%
\pgfpathlineto{\pgfqpoint{1.920167in}{1.654444in}}%
\pgfpathlineto{\pgfqpoint{1.905213in}{1.645016in}}%
\pgfpathlineto{\pgfqpoint{1.883866in}{1.636615in}}%
\pgfpathlineto{\pgfqpoint{1.862725in}{1.619824in}}%
\pgfpathlineto{\pgfqpoint{1.844396in}{1.597742in}}%
\pgfpathlineto{\pgfqpoint{1.830521in}{1.584687in}}%
\pgfpathlineto{\pgfqpoint{1.818367in}{1.575713in}}%
\pgfpathlineto{\pgfqpoint{1.810007in}{1.567025in}}%
\pgfpathlineto{\pgfqpoint{1.811020in}{1.530009in}}%
\pgfpathlineto{\pgfqpoint{1.794761in}{1.519600in}}%
\pgfpathlineto{\pgfqpoint{1.786550in}{1.506437in}}%
\pgfpathlineto{\pgfqpoint{1.786004in}{1.498105in}}%
\pgfpathlineto{\pgfqpoint{1.792621in}{1.496270in}}%
\pgfpathlineto{\pgfqpoint{1.796378in}{1.486783in}}%
\pgfpathlineto{\pgfqpoint{1.792197in}{1.480500in}}%
\pgfpathlineto{\pgfqpoint{1.792621in}{1.457675in}}%
\pgfpathlineto{\pgfqpoint{1.790419in}{1.445673in}}%
\pgfpathlineto{\pgfqpoint{1.803976in}{1.435243in}}%
\pgfpathlineto{\pgfqpoint{1.811570in}{1.434051in}}%
\pgfpathlineto{\pgfqpoint{1.817838in}{1.427026in}}%
\pgfpathlineto{\pgfqpoint{1.828501in}{1.422710in}}%
\pgfpathlineto{\pgfqpoint{1.834444in}{1.411992in}}%
\pgfpathlineto{\pgfqpoint{1.845940in}{1.402933in}}%
\pgfpathlineto{\pgfqpoint{1.856505in}{1.397444in}}%
\pgfpathlineto{\pgfqpoint{1.862453in}{1.387663in}}%
\pgfpathlineto{\pgfqpoint{1.864454in}{1.369657in}}%
\pgfpathlineto{\pgfqpoint{1.808388in}{1.367621in}}%
\pgfpathlineto{\pgfqpoint{1.760597in}{1.366622in}}%
\pgfpathlineto{\pgfqpoint{1.717055in}{1.365982in}}%
\pgfpathlineto{\pgfqpoint{1.670992in}{1.366052in}}%
\pgfpathlineto{\pgfqpoint{1.625565in}{1.366544in}}%
\pgfpathclose%
\pgfusepath{fill}%
\end{pgfscope}%
\begin{pgfscope}%
\pgfpathrectangle{\pgfqpoint{0.100000in}{0.100000in}}{\pgfqpoint{2.862985in}{1.832778in}}%
\pgfusepath{clip}%
\pgfsetbuttcap%
\pgfsetmiterjoin%
\definecolor{currentfill}{rgb}{0.647059,0.833987,0.435294}%
\pgfsetfillcolor{currentfill}%
\pgfsetlinewidth{0.000000pt}%
\definecolor{currentstroke}{rgb}{0.000000,0.000000,0.000000}%
\pgfsetstrokecolor{currentstroke}%
\pgfsetstrokeopacity{0.000000}%
\pgfsetdash{}{0pt}%
\pgfpathmoveto{\pgfqpoint{0.352592in}{1.488751in}}%
\pgfpathlineto{\pgfqpoint{0.348152in}{1.496921in}}%
\pgfpathlineto{\pgfqpoint{0.350987in}{1.517893in}}%
\pgfpathlineto{\pgfqpoint{0.356459in}{1.528853in}}%
\pgfpathlineto{\pgfqpoint{0.353724in}{1.543679in}}%
\pgfpathlineto{\pgfqpoint{0.359432in}{1.549914in}}%
\pgfpathlineto{\pgfqpoint{0.369847in}{1.566720in}}%
\pgfpathlineto{\pgfqpoint{0.378677in}{1.576864in}}%
\pgfpathlineto{\pgfqpoint{0.391585in}{1.598959in}}%
\pgfpathlineto{\pgfqpoint{0.401485in}{1.623028in}}%
\pgfpathlineto{\pgfqpoint{0.411995in}{1.645601in}}%
\pgfpathlineto{\pgfqpoint{0.414119in}{1.655029in}}%
\pgfpathlineto{\pgfqpoint{0.428617in}{1.682000in}}%
\pgfpathlineto{\pgfqpoint{0.431409in}{1.693845in}}%
\pgfpathlineto{\pgfqpoint{0.437582in}{1.706284in}}%
\pgfpathlineto{\pgfqpoint{0.439788in}{1.721458in}}%
\pgfpathlineto{\pgfqpoint{0.451573in}{1.728857in}}%
\pgfpathlineto{\pgfqpoint{0.465546in}{1.732587in}}%
\pgfpathlineto{\pgfqpoint{0.472610in}{1.724234in}}%
\pgfpathlineto{\pgfqpoint{0.478718in}{1.724852in}}%
\pgfpathlineto{\pgfqpoint{0.488254in}{1.714897in}}%
\pgfpathlineto{\pgfqpoint{0.489522in}{1.705519in}}%
\pgfpathlineto{\pgfqpoint{0.485905in}{1.687156in}}%
\pgfpathlineto{\pgfqpoint{0.497937in}{1.677787in}}%
\pgfpathlineto{\pgfqpoint{0.505740in}{1.674264in}}%
\pgfpathlineto{\pgfqpoint{0.526880in}{1.677860in}}%
\pgfpathlineto{\pgfqpoint{0.539127in}{1.675404in}}%
\pgfpathlineto{\pgfqpoint{0.547545in}{1.671374in}}%
\pgfpathlineto{\pgfqpoint{0.551979in}{1.663725in}}%
\pgfpathlineto{\pgfqpoint{0.580055in}{1.664746in}}%
\pgfpathlineto{\pgfqpoint{0.584520in}{1.660065in}}%
\pgfpathlineto{\pgfqpoint{0.595053in}{1.659055in}}%
\pgfpathlineto{\pgfqpoint{0.605660in}{1.662045in}}%
\pgfpathlineto{\pgfqpoint{0.637021in}{1.661292in}}%
\pgfpathlineto{\pgfqpoint{0.643794in}{1.658987in}}%
\pgfpathlineto{\pgfqpoint{0.651709in}{1.661643in}}%
\pgfpathlineto{\pgfqpoint{0.739329in}{1.638779in}}%
\pgfpathlineto{\pgfqpoint{0.742899in}{1.626494in}}%
\pgfpathlineto{\pgfqpoint{0.750939in}{1.619838in}}%
\pgfpathlineto{\pgfqpoint{0.751774in}{1.607856in}}%
\pgfpathlineto{\pgfqpoint{0.740015in}{1.595725in}}%
\pgfpathlineto{\pgfqpoint{0.734282in}{1.584259in}}%
\pgfpathlineto{\pgfqpoint{0.727377in}{1.578133in}}%
\pgfpathlineto{\pgfqpoint{0.726752in}{1.572002in}}%
\pgfpathlineto{\pgfqpoint{0.715313in}{1.563715in}}%
\pgfpathlineto{\pgfqpoint{0.707464in}{1.551738in}}%
\pgfpathlineto{\pgfqpoint{0.703070in}{1.548828in}}%
\pgfpathlineto{\pgfqpoint{0.700673in}{1.537205in}}%
\pgfpathlineto{\pgfqpoint{0.712217in}{1.525368in}}%
\pgfpathlineto{\pgfqpoint{0.701069in}{1.505674in}}%
\pgfpathlineto{\pgfqpoint{0.699717in}{1.498158in}}%
\pgfpathlineto{\pgfqpoint{0.674647in}{1.395314in}}%
\pgfpathlineto{\pgfqpoint{0.621916in}{1.408826in}}%
\pgfpathlineto{\pgfqpoint{0.571046in}{1.421943in}}%
\pgfpathlineto{\pgfqpoint{0.540358in}{1.430483in}}%
\pgfpathlineto{\pgfqpoint{0.500927in}{1.441715in}}%
\pgfpathlineto{\pgfqpoint{0.438032in}{1.461511in}}%
\pgfpathlineto{\pgfqpoint{0.369661in}{1.482798in}}%
\pgfpathlineto{\pgfqpoint{0.352592in}{1.488751in}}%
\pgfpathclose%
\pgfusepath{fill}%
\end{pgfscope}%
\begin{pgfscope}%
\pgfpathrectangle{\pgfqpoint{0.100000in}{0.100000in}}{\pgfqpoint{2.862985in}{1.832778in}}%
\pgfusepath{clip}%
\pgfsetbuttcap%
\pgfsetmiterjoin%
\definecolor{currentfill}{rgb}{0.428835,0.679123,0.210611}%
\pgfsetfillcolor{currentfill}%
\pgfsetlinewidth{0.000000pt}%
\definecolor{currentstroke}{rgb}{0.000000,0.000000,0.000000}%
\pgfsetstrokecolor{currentstroke}%
\pgfsetstrokeopacity{0.000000}%
\pgfsetdash{}{0pt}%
\pgfpathmoveto{\pgfqpoint{2.722258in}{1.445782in}}%
\pgfpathlineto{\pgfqpoint{2.717278in}{1.449830in}}%
\pgfpathlineto{\pgfqpoint{2.717972in}{1.462444in}}%
\pgfpathlineto{\pgfqpoint{2.712824in}{1.498298in}}%
\pgfpathlineto{\pgfqpoint{2.716859in}{1.510906in}}%
\pgfpathlineto{\pgfqpoint{2.719741in}{1.533267in}}%
\pgfpathlineto{\pgfqpoint{2.715852in}{1.543846in}}%
\pgfpathlineto{\pgfqpoint{2.729999in}{1.557930in}}%
\pgfpathlineto{\pgfqpoint{2.733559in}{1.569133in}}%
\pgfpathlineto{\pgfqpoint{2.726856in}{1.577975in}}%
\pgfpathlineto{\pgfqpoint{2.729711in}{1.589025in}}%
\pgfpathlineto{\pgfqpoint{2.727922in}{1.595966in}}%
\pgfpathlineto{\pgfqpoint{2.729459in}{1.610832in}}%
\pgfpathlineto{\pgfqpoint{2.732332in}{1.615421in}}%
\pgfpathlineto{\pgfqpoint{2.741192in}{1.618318in}}%
\pgfpathlineto{\pgfqpoint{2.754117in}{1.580468in}}%
\pgfpathlineto{\pgfqpoint{2.778094in}{1.502035in}}%
\pgfpathlineto{\pgfqpoint{2.787168in}{1.495954in}}%
\pgfpathlineto{\pgfqpoint{2.788801in}{1.489012in}}%
\pgfpathlineto{\pgfqpoint{2.793590in}{1.486207in}}%
\pgfpathlineto{\pgfqpoint{2.793226in}{1.473617in}}%
\pgfpathlineto{\pgfqpoint{2.788151in}{1.473415in}}%
\pgfpathlineto{\pgfqpoint{2.777871in}{1.465574in}}%
\pgfpathlineto{\pgfqpoint{2.774906in}{1.457709in}}%
\pgfpathlineto{\pgfqpoint{2.747390in}{1.450901in}}%
\pgfpathlineto{\pgfqpoint{2.722258in}{1.445782in}}%
\pgfpathclose%
\pgfusepath{fill}%
\end{pgfscope}%
\begin{pgfscope}%
\pgfpathrectangle{\pgfqpoint{0.100000in}{0.100000in}}{\pgfqpoint{2.862985in}{1.832778in}}%
\pgfusepath{clip}%
\pgfsetbuttcap%
\pgfsetmiterjoin%
\definecolor{currentfill}{rgb}{0.524337,0.754325,0.286736}%
\pgfsetfillcolor{currentfill}%
\pgfsetlinewidth{0.000000pt}%
\definecolor{currentstroke}{rgb}{0.000000,0.000000,0.000000}%
\pgfsetstrokecolor{currentstroke}%
\pgfsetstrokeopacity{0.000000}%
\pgfsetdash{}{0pt}%
\pgfpathmoveto{\pgfqpoint{1.861865in}{1.172908in}}%
\pgfpathlineto{\pgfqpoint{1.847339in}{1.187293in}}%
\pgfpathlineto{\pgfqpoint{1.800906in}{1.184615in}}%
\pgfpathlineto{\pgfqpoint{1.728489in}{1.182229in}}%
\pgfpathlineto{\pgfqpoint{1.655637in}{1.183365in}}%
\pgfpathlineto{\pgfqpoint{1.650523in}{1.192282in}}%
\pgfpathlineto{\pgfqpoint{1.652611in}{1.201038in}}%
\pgfpathlineto{\pgfqpoint{1.651620in}{1.216037in}}%
\pgfpathlineto{\pgfqpoint{1.647770in}{1.230565in}}%
\pgfpathlineto{\pgfqpoint{1.648217in}{1.238237in}}%
\pgfpathlineto{\pgfqpoint{1.640033in}{1.246905in}}%
\pgfpathlineto{\pgfqpoint{1.641814in}{1.258967in}}%
\pgfpathlineto{\pgfqpoint{1.629260in}{1.282791in}}%
\pgfpathlineto{\pgfqpoint{1.625519in}{1.303018in}}%
\pgfpathlineto{\pgfqpoint{1.616489in}{1.319470in}}%
\pgfpathlineto{\pgfqpoint{1.625061in}{1.339878in}}%
\pgfpathlineto{\pgfqpoint{1.620791in}{1.349375in}}%
\pgfpathlineto{\pgfqpoint{1.625565in}{1.366544in}}%
\pgfpathlineto{\pgfqpoint{1.670992in}{1.366052in}}%
\pgfpathlineto{\pgfqpoint{1.717055in}{1.365982in}}%
\pgfpathlineto{\pgfqpoint{1.760597in}{1.366622in}}%
\pgfpathlineto{\pgfqpoint{1.808388in}{1.367621in}}%
\pgfpathlineto{\pgfqpoint{1.864454in}{1.369657in}}%
\pgfpathlineto{\pgfqpoint{1.865951in}{1.361517in}}%
\pgfpathlineto{\pgfqpoint{1.872158in}{1.354693in}}%
\pgfpathlineto{\pgfqpoint{1.868228in}{1.347582in}}%
\pgfpathlineto{\pgfqpoint{1.869839in}{1.333702in}}%
\pgfpathlineto{\pgfqpoint{1.873686in}{1.322483in}}%
\pgfpathlineto{\pgfqpoint{1.891125in}{1.316544in}}%
\pgfpathlineto{\pgfqpoint{1.894114in}{1.308600in}}%
\pgfpathlineto{\pgfqpoint{1.903684in}{1.299681in}}%
\pgfpathlineto{\pgfqpoint{1.907588in}{1.290454in}}%
\pgfpathlineto{\pgfqpoint{1.917285in}{1.284265in}}%
\pgfpathlineto{\pgfqpoint{1.918802in}{1.276813in}}%
\pgfpathlineto{\pgfqpoint{1.916914in}{1.265532in}}%
\pgfpathlineto{\pgfqpoint{1.911979in}{1.262151in}}%
\pgfpathlineto{\pgfqpoint{1.910490in}{1.251413in}}%
\pgfpathlineto{\pgfqpoint{1.896306in}{1.242884in}}%
\pgfpathlineto{\pgfqpoint{1.877818in}{1.238212in}}%
\pgfpathlineto{\pgfqpoint{1.876124in}{1.227469in}}%
\pgfpathlineto{\pgfqpoint{1.883256in}{1.219794in}}%
\pgfpathlineto{\pgfqpoint{1.883547in}{1.210143in}}%
\pgfpathlineto{\pgfqpoint{1.877793in}{1.202558in}}%
\pgfpathlineto{\pgfqpoint{1.874773in}{1.191288in}}%
\pgfpathlineto{\pgfqpoint{1.864771in}{1.187558in}}%
\pgfpathlineto{\pgfqpoint{1.865408in}{1.174999in}}%
\pgfpathlineto{\pgfqpoint{1.861865in}{1.172908in}}%
\pgfpathclose%
\pgfusepath{fill}%
\end{pgfscope}%
\begin{pgfscope}%
\pgfpathrectangle{\pgfqpoint{0.100000in}{0.100000in}}{\pgfqpoint{2.862985in}{1.832778in}}%
\pgfusepath{clip}%
\pgfsetbuttcap%
\pgfsetmiterjoin%
\definecolor{currentfill}{rgb}{0.498039,0.737255,0.254902}%
\pgfsetfillcolor{currentfill}%
\pgfsetlinewidth{0.000000pt}%
\definecolor{currentstroke}{rgb}{0.000000,0.000000,0.000000}%
\pgfsetstrokecolor{currentstroke}%
\pgfsetstrokeopacity{0.000000}%
\pgfsetdash{}{0pt}%
\pgfpathmoveto{\pgfqpoint{2.763120in}{1.409488in}}%
\pgfpathlineto{\pgfqpoint{2.749362in}{1.407318in}}%
\pgfpathlineto{\pgfqpoint{2.686166in}{1.392957in}}%
\pgfpathlineto{\pgfqpoint{2.685061in}{1.394631in}}%
\pgfpathlineto{\pgfqpoint{2.686015in}{1.438080in}}%
\pgfpathlineto{\pgfqpoint{2.722258in}{1.445782in}}%
\pgfpathlineto{\pgfqpoint{2.747390in}{1.450901in}}%
\pgfpathlineto{\pgfqpoint{2.774906in}{1.457709in}}%
\pgfpathlineto{\pgfqpoint{2.777871in}{1.465574in}}%
\pgfpathlineto{\pgfqpoint{2.788151in}{1.473415in}}%
\pgfpathlineto{\pgfqpoint{2.793226in}{1.473617in}}%
\pgfpathlineto{\pgfqpoint{2.799896in}{1.462179in}}%
\pgfpathlineto{\pgfqpoint{2.793843in}{1.445482in}}%
\pgfpathlineto{\pgfqpoint{2.792955in}{1.435691in}}%
\pgfpathlineto{\pgfqpoint{2.805207in}{1.436611in}}%
\pgfpathlineto{\pgfqpoint{2.810763in}{1.431914in}}%
\pgfpathlineto{\pgfqpoint{2.823224in}{1.412717in}}%
\pgfpathlineto{\pgfqpoint{2.829417in}{1.410382in}}%
\pgfpathlineto{\pgfqpoint{2.839790in}{1.411248in}}%
\pgfpathlineto{\pgfqpoint{2.847009in}{1.417772in}}%
\pgfpathlineto{\pgfqpoint{2.851808in}{1.411950in}}%
\pgfpathlineto{\pgfqpoint{2.821631in}{1.396025in}}%
\pgfpathlineto{\pgfqpoint{2.820688in}{1.407452in}}%
\pgfpathlineto{\pgfqpoint{2.806980in}{1.389693in}}%
\pgfpathlineto{\pgfqpoint{2.802160in}{1.386633in}}%
\pgfpathlineto{\pgfqpoint{2.795438in}{1.396881in}}%
\pgfpathlineto{\pgfqpoint{2.793598in}{1.398285in}}%
\pgfpathlineto{\pgfqpoint{2.787340in}{1.401600in}}%
\pgfpathlineto{\pgfqpoint{2.781882in}{1.415045in}}%
\pgfpathlineto{\pgfqpoint{2.763120in}{1.409488in}}%
\pgfpathclose%
\pgfusepath{fill}%
\end{pgfscope}%
\begin{pgfscope}%
\pgfpathrectangle{\pgfqpoint{0.100000in}{0.100000in}}{\pgfqpoint{2.862985in}{1.832778in}}%
\pgfusepath{clip}%
\pgfsetbuttcap%
\pgfsetmiterjoin%
\definecolor{currentfill}{rgb}{0.321184,0.588697,0.141715}%
\pgfsetfillcolor{currentfill}%
\pgfsetlinewidth{0.000000pt}%
\definecolor{currentstroke}{rgb}{0.000000,0.000000,0.000000}%
\pgfsetstrokecolor{currentstroke}%
\pgfsetstrokeopacity{0.000000}%
\pgfsetdash{}{0pt}%
\pgfpathmoveto{\pgfqpoint{1.353739in}{1.159333in}}%
\pgfpathlineto{\pgfqpoint{1.358777in}{1.221808in}}%
\pgfpathlineto{\pgfqpoint{1.330243in}{1.224124in}}%
\pgfpathlineto{\pgfqpoint{1.264283in}{1.230524in}}%
\pgfpathlineto{\pgfqpoint{1.270419in}{1.292777in}}%
\pgfpathlineto{\pgfqpoint{1.276565in}{1.355506in}}%
\pgfpathlineto{\pgfqpoint{1.301922in}{1.352768in}}%
\pgfpathlineto{\pgfqpoint{1.366682in}{1.346649in}}%
\pgfpathlineto{\pgfqpoint{1.405874in}{1.343580in}}%
\pgfpathlineto{\pgfqpoint{1.453322in}{1.340860in}}%
\pgfpathlineto{\pgfqpoint{1.531420in}{1.337164in}}%
\pgfpathlineto{\pgfqpoint{1.533137in}{1.333663in}}%
\pgfpathlineto{\pgfqpoint{1.555515in}{1.322591in}}%
\pgfpathlineto{\pgfqpoint{1.562377in}{1.328214in}}%
\pgfpathlineto{\pgfqpoint{1.582148in}{1.327701in}}%
\pgfpathlineto{\pgfqpoint{1.593930in}{1.321529in}}%
\pgfpathlineto{\pgfqpoint{1.612612in}{1.314513in}}%
\pgfpathlineto{\pgfqpoint{1.617847in}{1.304277in}}%
\pgfpathlineto{\pgfqpoint{1.625519in}{1.303018in}}%
\pgfpathlineto{\pgfqpoint{1.629260in}{1.282791in}}%
\pgfpathlineto{\pgfqpoint{1.641814in}{1.258967in}}%
\pgfpathlineto{\pgfqpoint{1.640033in}{1.246905in}}%
\pgfpathlineto{\pgfqpoint{1.648217in}{1.238237in}}%
\pgfpathlineto{\pgfqpoint{1.647770in}{1.230565in}}%
\pgfpathlineto{\pgfqpoint{1.651620in}{1.216037in}}%
\pgfpathlineto{\pgfqpoint{1.652611in}{1.201038in}}%
\pgfpathlineto{\pgfqpoint{1.650523in}{1.192282in}}%
\pgfpathlineto{\pgfqpoint{1.655637in}{1.183365in}}%
\pgfpathlineto{\pgfqpoint{1.662651in}{1.167151in}}%
\pgfpathlineto{\pgfqpoint{1.669364in}{1.160551in}}%
\pgfpathlineto{\pgfqpoint{1.677367in}{1.146251in}}%
\pgfpathlineto{\pgfqpoint{1.654686in}{1.146016in}}%
\pgfpathlineto{\pgfqpoint{1.605647in}{1.146774in}}%
\pgfpathlineto{\pgfqpoint{1.551466in}{1.148434in}}%
\pgfpathlineto{\pgfqpoint{1.496966in}{1.150526in}}%
\pgfpathlineto{\pgfqpoint{1.416835in}{1.154903in}}%
\pgfpathlineto{\pgfqpoint{1.353739in}{1.159333in}}%
\pgfpathclose%
\pgfusepath{fill}%
\end{pgfscope}%
\begin{pgfscope}%
\pgfpathrectangle{\pgfqpoint{0.100000in}{0.100000in}}{\pgfqpoint{2.862985in}{1.832778in}}%
\pgfusepath{clip}%
\pgfsetbuttcap%
\pgfsetmiterjoin%
\definecolor{currentfill}{rgb}{0.781699,0.908497,0.622222}%
\pgfsetfillcolor{currentfill}%
\pgfsetlinewidth{0.000000pt}%
\definecolor{currentstroke}{rgb}{0.000000,0.000000,0.000000}%
\pgfsetstrokecolor{currentstroke}%
\pgfsetstrokeopacity{0.000000}%
\pgfsetdash{}{0pt}%
\pgfpathmoveto{\pgfqpoint{2.397084in}{1.347781in}}%
\pgfpathlineto{\pgfqpoint{2.422175in}{1.371705in}}%
\pgfpathlineto{\pgfqpoint{2.425294in}{1.380209in}}%
\pgfpathlineto{\pgfqpoint{2.432646in}{1.387490in}}%
\pgfpathlineto{\pgfqpoint{2.427129in}{1.398090in}}%
\pgfpathlineto{\pgfqpoint{2.420212in}{1.404290in}}%
\pgfpathlineto{\pgfqpoint{2.418213in}{1.415277in}}%
\pgfpathlineto{\pgfqpoint{2.443959in}{1.426558in}}%
\pgfpathlineto{\pgfqpoint{2.465274in}{1.430126in}}%
\pgfpathlineto{\pgfqpoint{2.476728in}{1.430365in}}%
\pgfpathlineto{\pgfqpoint{2.485453in}{1.426054in}}%
\pgfpathlineto{\pgfqpoint{2.493958in}{1.429903in}}%
\pgfpathlineto{\pgfqpoint{2.514702in}{1.434217in}}%
\pgfpathlineto{\pgfqpoint{2.521870in}{1.439785in}}%
\pgfpathlineto{\pgfqpoint{2.532498in}{1.452111in}}%
\pgfpathlineto{\pgfqpoint{2.542167in}{1.457558in}}%
\pgfpathlineto{\pgfqpoint{2.542849in}{1.462782in}}%
\pgfpathlineto{\pgfqpoint{2.537774in}{1.474701in}}%
\pgfpathlineto{\pgfqpoint{2.541443in}{1.481714in}}%
\pgfpathlineto{\pgfqpoint{2.536503in}{1.489256in}}%
\pgfpathlineto{\pgfqpoint{2.528939in}{1.489782in}}%
\pgfpathlineto{\pgfqpoint{2.547870in}{1.512563in}}%
\pgfpathlineto{\pgfqpoint{2.550119in}{1.521245in}}%
\pgfpathlineto{\pgfqpoint{2.565018in}{1.543386in}}%
\pgfpathlineto{\pgfqpoint{2.578850in}{1.555368in}}%
\pgfpathlineto{\pgfqpoint{2.588347in}{1.560371in}}%
\pgfpathlineto{\pgfqpoint{2.619425in}{1.567413in}}%
\pgfpathlineto{\pgfqpoint{2.648552in}{1.575365in}}%
\pgfpathlineto{\pgfqpoint{2.652410in}{1.562882in}}%
\pgfpathlineto{\pgfqpoint{2.652683in}{1.551308in}}%
\pgfpathlineto{\pgfqpoint{2.659327in}{1.540801in}}%
\pgfpathlineto{\pgfqpoint{2.661375in}{1.530104in}}%
\pgfpathlineto{\pgfqpoint{2.658822in}{1.515590in}}%
\pgfpathlineto{\pgfqpoint{2.666598in}{1.496013in}}%
\pgfpathlineto{\pgfqpoint{2.675463in}{1.485127in}}%
\pgfpathlineto{\pgfqpoint{2.682320in}{1.449744in}}%
\pgfpathlineto{\pgfqpoint{2.686015in}{1.438080in}}%
\pgfpathlineto{\pgfqpoint{2.685061in}{1.394631in}}%
\pgfpathlineto{\pgfqpoint{2.686166in}{1.392957in}}%
\pgfpathlineto{\pgfqpoint{2.694238in}{1.346227in}}%
\pgfpathlineto{\pgfqpoint{2.698765in}{1.341969in}}%
\pgfpathlineto{\pgfqpoint{2.689028in}{1.332508in}}%
\pgfpathlineto{\pgfqpoint{2.693832in}{1.327078in}}%
\pgfpathlineto{\pgfqpoint{2.689609in}{1.318865in}}%
\pgfpathlineto{\pgfqpoint{2.689645in}{1.315368in}}%
\pgfpathlineto{\pgfqpoint{2.684366in}{1.312213in}}%
\pgfpathlineto{\pgfqpoint{2.681795in}{1.305239in}}%
\pgfpathlineto{\pgfqpoint{2.682610in}{1.324417in}}%
\pgfpathlineto{\pgfqpoint{2.666246in}{1.328639in}}%
\pgfpathlineto{\pgfqpoint{2.640655in}{1.337368in}}%
\pgfpathlineto{\pgfqpoint{2.637035in}{1.341668in}}%
\pgfpathlineto{\pgfqpoint{2.626342in}{1.342698in}}%
\pgfpathlineto{\pgfqpoint{2.620190in}{1.349098in}}%
\pgfpathlineto{\pgfqpoint{2.616871in}{1.361844in}}%
\pgfpathlineto{\pgfqpoint{2.608151in}{1.363429in}}%
\pgfpathlineto{\pgfqpoint{2.602319in}{1.370116in}}%
\pgfpathlineto{\pgfqpoint{2.528076in}{1.355148in}}%
\pgfpathlineto{\pgfqpoint{2.492594in}{1.347810in}}%
\pgfpathlineto{\pgfqpoint{2.438721in}{1.337977in}}%
\pgfpathlineto{\pgfqpoint{2.399926in}{1.331432in}}%
\pgfpathlineto{\pgfqpoint{2.397084in}{1.347781in}}%
\pgfpathclose%
\pgfusepath{fill}%
\end{pgfscope}%
\begin{pgfscope}%
\pgfpathrectangle{\pgfqpoint{0.100000in}{0.100000in}}{\pgfqpoint{2.862985in}{1.832778in}}%
\pgfusepath{clip}%
\pgfsetbuttcap%
\pgfsetmiterjoin%
\definecolor{currentfill}{rgb}{0.781699,0.908497,0.622222}%
\pgfsetfillcolor{currentfill}%
\pgfsetlinewidth{0.000000pt}%
\definecolor{currentstroke}{rgb}{0.000000,0.000000,0.000000}%
\pgfsetstrokecolor{currentstroke}%
\pgfsetstrokeopacity{0.000000}%
\pgfsetdash{}{0pt}%
\pgfpathmoveto{\pgfqpoint{2.695139in}{1.301344in}}%
\pgfpathlineto{\pgfqpoint{2.683657in}{1.297767in}}%
\pgfpathlineto{\pgfqpoint{2.681726in}{1.301055in}}%
\pgfpathlineto{\pgfqpoint{2.685411in}{1.312094in}}%
\pgfpathlineto{\pgfqpoint{2.692237in}{1.313181in}}%
\pgfpathlineto{\pgfqpoint{2.691626in}{1.316656in}}%
\pgfpathlineto{\pgfqpoint{2.697754in}{1.321872in}}%
\pgfpathlineto{\pgfqpoint{2.715472in}{1.326028in}}%
\pgfpathlineto{\pgfqpoint{2.723360in}{1.332319in}}%
\pgfpathlineto{\pgfqpoint{2.741079in}{1.337561in}}%
\pgfpathlineto{\pgfqpoint{2.749161in}{1.335643in}}%
\pgfpathlineto{\pgfqpoint{2.755959in}{1.344095in}}%
\pgfpathlineto{\pgfqpoint{2.766234in}{1.345211in}}%
\pgfpathlineto{\pgfqpoint{2.748704in}{1.328725in}}%
\pgfpathlineto{\pgfqpoint{2.695139in}{1.301344in}}%
\pgfpathclose%
\pgfusepath{fill}%
\end{pgfscope}%
\begin{pgfscope}%
\pgfpathrectangle{\pgfqpoint{0.100000in}{0.100000in}}{\pgfqpoint{2.862985in}{1.832778in}}%
\pgfusepath{clip}%
\pgfsetbuttcap%
\pgfsetmiterjoin%
\definecolor{currentfill}{rgb}{0.524337,0.754325,0.286736}%
\pgfsetfillcolor{currentfill}%
\pgfsetlinewidth{0.000000pt}%
\definecolor{currentstroke}{rgb}{0.000000,0.000000,0.000000}%
\pgfsetstrokecolor{currentstroke}%
\pgfsetstrokeopacity{0.000000}%
\pgfsetdash{}{0pt}%
\pgfpathmoveto{\pgfqpoint{2.437178in}{1.192485in}}%
\pgfpathlineto{\pgfqpoint{2.387526in}{1.184275in}}%
\pgfpathlineto{\pgfqpoint{2.378519in}{1.241018in}}%
\pgfpathlineto{\pgfqpoint{2.365146in}{1.324661in}}%
\pgfpathlineto{\pgfqpoint{2.397084in}{1.347781in}}%
\pgfpathlineto{\pgfqpoint{2.399926in}{1.331432in}}%
\pgfpathlineto{\pgfqpoint{2.438721in}{1.337977in}}%
\pgfpathlineto{\pgfqpoint{2.492594in}{1.347810in}}%
\pgfpathlineto{\pgfqpoint{2.528076in}{1.355148in}}%
\pgfpathlineto{\pgfqpoint{2.602319in}{1.370116in}}%
\pgfpathlineto{\pgfqpoint{2.608151in}{1.363429in}}%
\pgfpathlineto{\pgfqpoint{2.616871in}{1.361844in}}%
\pgfpathlineto{\pgfqpoint{2.620190in}{1.349098in}}%
\pgfpathlineto{\pgfqpoint{2.626342in}{1.342698in}}%
\pgfpathlineto{\pgfqpoint{2.637035in}{1.341668in}}%
\pgfpathlineto{\pgfqpoint{2.640655in}{1.337368in}}%
\pgfpathlineto{\pgfqpoint{2.636981in}{1.334051in}}%
\pgfpathlineto{\pgfqpoint{2.633670in}{1.322314in}}%
\pgfpathlineto{\pgfqpoint{2.625913in}{1.309125in}}%
\pgfpathlineto{\pgfqpoint{2.631120in}{1.303390in}}%
\pgfpathlineto{\pgfqpoint{2.626859in}{1.297245in}}%
\pgfpathlineto{\pgfqpoint{2.629263in}{1.283769in}}%
\pgfpathlineto{\pgfqpoint{2.636936in}{1.276695in}}%
\pgfpathlineto{\pgfqpoint{2.655144in}{1.265202in}}%
\pgfpathlineto{\pgfqpoint{2.640481in}{1.248986in}}%
\pgfpathlineto{\pgfqpoint{2.640275in}{1.242831in}}%
\pgfpathlineto{\pgfqpoint{2.628338in}{1.234894in}}%
\pgfpathlineto{\pgfqpoint{2.615138in}{1.233404in}}%
\pgfpathlineto{\pgfqpoint{2.611873in}{1.226505in}}%
\pgfpathlineto{\pgfqpoint{2.575158in}{1.218561in}}%
\pgfpathlineto{\pgfqpoint{2.532316in}{1.209933in}}%
\pgfpathlineto{\pgfqpoint{2.502868in}{1.204659in}}%
\pgfpathlineto{\pgfqpoint{2.437178in}{1.192485in}}%
\pgfpathclose%
\pgfusepath{fill}%
\end{pgfscope}%
\begin{pgfscope}%
\pgfpathrectangle{\pgfqpoint{0.100000in}{0.100000in}}{\pgfqpoint{2.862985in}{1.832778in}}%
\pgfusepath{clip}%
\pgfsetbuttcap%
\pgfsetmiterjoin%
\definecolor{currentfill}{rgb}{0.673356,0.851057,0.467128}%
\pgfsetfillcolor{currentfill}%
\pgfsetlinewidth{0.000000pt}%
\definecolor{currentstroke}{rgb}{0.000000,0.000000,0.000000}%
\pgfsetstrokecolor{currentstroke}%
\pgfsetstrokeopacity{0.000000}%
\pgfsetdash{}{0pt}%
\pgfpathmoveto{\pgfqpoint{2.686166in}{1.392957in}}%
\pgfpathlineto{\pgfqpoint{2.749362in}{1.407318in}}%
\pgfpathlineto{\pgfqpoint{2.763120in}{1.409488in}}%
\pgfpathlineto{\pgfqpoint{2.772226in}{1.373686in}}%
\pgfpathlineto{\pgfqpoint{2.770785in}{1.367275in}}%
\pgfpathlineto{\pgfqpoint{2.741533in}{1.355951in}}%
\pgfpathlineto{\pgfqpoint{2.724071in}{1.351992in}}%
\pgfpathlineto{\pgfqpoint{2.716651in}{1.343114in}}%
\pgfpathlineto{\pgfqpoint{2.693832in}{1.327078in}}%
\pgfpathlineto{\pgfqpoint{2.689028in}{1.332508in}}%
\pgfpathlineto{\pgfqpoint{2.698765in}{1.341969in}}%
\pgfpathlineto{\pgfqpoint{2.694238in}{1.346227in}}%
\pgfpathlineto{\pgfqpoint{2.686166in}{1.392957in}}%
\pgfpathclose%
\pgfusepath{fill}%
\end{pgfscope}%
\begin{pgfscope}%
\pgfpathrectangle{\pgfqpoint{0.100000in}{0.100000in}}{\pgfqpoint{2.862985in}{1.832778in}}%
\pgfusepath{clip}%
\pgfsetbuttcap%
\pgfsetmiterjoin%
\definecolor{currentfill}{rgb}{0.594464,0.799846,0.371626}%
\pgfsetfillcolor{currentfill}%
\pgfsetlinewidth{0.000000pt}%
\definecolor{currentstroke}{rgb}{0.000000,0.000000,0.000000}%
\pgfsetstrokecolor{currentstroke}%
\pgfsetstrokeopacity{0.000000}%
\pgfsetdash{}{0pt}%
\pgfpathmoveto{\pgfqpoint{2.763120in}{1.409488in}}%
\pgfpathlineto{\pgfqpoint{2.781882in}{1.415045in}}%
\pgfpathlineto{\pgfqpoint{2.787340in}{1.401600in}}%
\pgfpathlineto{\pgfqpoint{2.793598in}{1.398285in}}%
\pgfpathlineto{\pgfqpoint{2.785875in}{1.392622in}}%
\pgfpathlineto{\pgfqpoint{2.786849in}{1.375988in}}%
\pgfpathlineto{\pgfqpoint{2.770785in}{1.367275in}}%
\pgfpathlineto{\pgfqpoint{2.772226in}{1.373686in}}%
\pgfpathlineto{\pgfqpoint{2.763120in}{1.409488in}}%
\pgfpathclose%
\pgfusepath{fill}%
\end{pgfscope}%
\begin{pgfscope}%
\pgfpathrectangle{\pgfqpoint{0.100000in}{0.100000in}}{\pgfqpoint{2.862985in}{1.832778in}}%
\pgfusepath{clip}%
\pgfsetbuttcap%
\pgfsetmiterjoin%
\definecolor{currentfill}{rgb}{0.838293,0.933103,0.713264}%
\pgfsetfillcolor{currentfill}%
\pgfsetlinewidth{0.000000pt}%
\definecolor{currentstroke}{rgb}{0.000000,0.000000,0.000000}%
\pgfsetstrokecolor{currentstroke}%
\pgfsetstrokeopacity{0.000000}%
\pgfsetdash{}{0pt}%
\pgfpathmoveto{\pgfqpoint{2.626223in}{1.229031in}}%
\pgfpathlineto{\pgfqpoint{2.628338in}{1.234894in}}%
\pgfpathlineto{\pgfqpoint{2.640275in}{1.242831in}}%
\pgfpathlineto{\pgfqpoint{2.640481in}{1.248986in}}%
\pgfpathlineto{\pgfqpoint{2.655144in}{1.265202in}}%
\pgfpathlineto{\pgfqpoint{2.636936in}{1.276695in}}%
\pgfpathlineto{\pgfqpoint{2.629263in}{1.283769in}}%
\pgfpathlineto{\pgfqpoint{2.626859in}{1.297245in}}%
\pgfpathlineto{\pgfqpoint{2.631120in}{1.303390in}}%
\pgfpathlineto{\pgfqpoint{2.625913in}{1.309125in}}%
\pgfpathlineto{\pgfqpoint{2.633670in}{1.322314in}}%
\pgfpathlineto{\pgfqpoint{2.636981in}{1.334051in}}%
\pgfpathlineto{\pgfqpoint{2.640655in}{1.337368in}}%
\pgfpathlineto{\pgfqpoint{2.666246in}{1.328639in}}%
\pgfpathlineto{\pgfqpoint{2.682610in}{1.324417in}}%
\pgfpathlineto{\pgfqpoint{2.681795in}{1.305239in}}%
\pgfpathlineto{\pgfqpoint{2.671854in}{1.290696in}}%
\pgfpathlineto{\pgfqpoint{2.680045in}{1.288556in}}%
\pgfpathlineto{\pgfqpoint{2.688550in}{1.282315in}}%
\pgfpathlineto{\pgfqpoint{2.689051in}{1.265245in}}%
\pgfpathlineto{\pgfqpoint{2.686477in}{1.253161in}}%
\pgfpathlineto{\pgfqpoint{2.688190in}{1.243235in}}%
\pgfpathlineto{\pgfqpoint{2.673407in}{1.209740in}}%
\pgfpathlineto{\pgfqpoint{2.668094in}{1.194100in}}%
\pgfpathlineto{\pgfqpoint{2.660676in}{1.201703in}}%
\pgfpathlineto{\pgfqpoint{2.650838in}{1.200407in}}%
\pgfpathlineto{\pgfqpoint{2.626218in}{1.214632in}}%
\pgfpathlineto{\pgfqpoint{2.623699in}{1.222249in}}%
\pgfpathlineto{\pgfqpoint{2.626223in}{1.229031in}}%
\pgfpathclose%
\pgfusepath{fill}%
\end{pgfscope}%
\begin{pgfscope}%
\pgfpathrectangle{\pgfqpoint{0.100000in}{0.100000in}}{\pgfqpoint{2.862985in}{1.832778in}}%
\pgfusepath{clip}%
\pgfsetbuttcap%
\pgfsetmiterjoin%
\definecolor{currentfill}{rgb}{0.647059,0.833987,0.435294}%
\pgfsetfillcolor{currentfill}%
\pgfsetlinewidth{0.000000pt}%
\definecolor{currentstroke}{rgb}{0.000000,0.000000,0.000000}%
\pgfsetstrokecolor{currentstroke}%
\pgfsetstrokeopacity{0.000000}%
\pgfsetdash{}{0pt}%
\pgfpathmoveto{\pgfqpoint{2.037013in}{1.021868in}}%
\pgfpathlineto{\pgfqpoint{2.035862in}{1.037407in}}%
\pgfpathlineto{\pgfqpoint{2.040336in}{1.045116in}}%
\pgfpathlineto{\pgfqpoint{2.037387in}{1.048893in}}%
\pgfpathlineto{\pgfqpoint{2.048140in}{1.061261in}}%
\pgfpathlineto{\pgfqpoint{2.047524in}{1.064241in}}%
\pgfpathlineto{\pgfqpoint{2.057971in}{1.085818in}}%
\pgfpathlineto{\pgfqpoint{2.055859in}{1.096224in}}%
\pgfpathlineto{\pgfqpoint{2.048283in}{1.107118in}}%
\pgfpathlineto{\pgfqpoint{2.053542in}{1.120373in}}%
\pgfpathlineto{\pgfqpoint{2.046726in}{1.207597in}}%
\pgfpathlineto{\pgfqpoint{2.041818in}{1.268787in}}%
\pgfpathlineto{\pgfqpoint{2.048598in}{1.263717in}}%
\pgfpathlineto{\pgfqpoint{2.056163in}{1.263854in}}%
\pgfpathlineto{\pgfqpoint{2.074029in}{1.274203in}}%
\pgfpathlineto{\pgfqpoint{2.128831in}{1.279317in}}%
\pgfpathlineto{\pgfqpoint{2.169393in}{1.283592in}}%
\pgfpathlineto{\pgfqpoint{2.169752in}{1.279623in}}%
\pgfpathlineto{\pgfqpoint{2.179011in}{1.195770in}}%
\pgfpathlineto{\pgfqpoint{2.186981in}{1.117748in}}%
\pgfpathlineto{\pgfqpoint{2.183545in}{1.114085in}}%
\pgfpathlineto{\pgfqpoint{2.188822in}{1.104993in}}%
\pgfpathlineto{\pgfqpoint{2.188795in}{1.098441in}}%
\pgfpathlineto{\pgfqpoint{2.181265in}{1.096795in}}%
\pgfpathlineto{\pgfqpoint{2.172839in}{1.090478in}}%
\pgfpathlineto{\pgfqpoint{2.167134in}{1.092965in}}%
\pgfpathlineto{\pgfqpoint{2.158595in}{1.088922in}}%
\pgfpathlineto{\pgfqpoint{2.161237in}{1.080801in}}%
\pgfpathlineto{\pgfqpoint{2.152460in}{1.072642in}}%
\pgfpathlineto{\pgfqpoint{2.150037in}{1.063205in}}%
\pgfpathlineto{\pgfqpoint{2.144002in}{1.061656in}}%
\pgfpathlineto{\pgfqpoint{2.139498in}{1.054509in}}%
\pgfpathlineto{\pgfqpoint{2.140088in}{1.047311in}}%
\pgfpathlineto{\pgfqpoint{2.134790in}{1.042250in}}%
\pgfpathlineto{\pgfqpoint{2.126817in}{1.043031in}}%
\pgfpathlineto{\pgfqpoint{2.120757in}{1.050797in}}%
\pgfpathlineto{\pgfqpoint{2.110489in}{1.043318in}}%
\pgfpathlineto{\pgfqpoint{2.111209in}{1.036783in}}%
\pgfpathlineto{\pgfqpoint{2.099789in}{1.032965in}}%
\pgfpathlineto{\pgfqpoint{2.096919in}{1.037771in}}%
\pgfpathlineto{\pgfqpoint{2.085801in}{1.032321in}}%
\pgfpathlineto{\pgfqpoint{2.081739in}{1.024520in}}%
\pgfpathlineto{\pgfqpoint{2.068349in}{1.032520in}}%
\pgfpathlineto{\pgfqpoint{2.047127in}{1.027220in}}%
\pgfpathlineto{\pgfqpoint{2.037013in}{1.021868in}}%
\pgfpathclose%
\pgfusepath{fill}%
\end{pgfscope}%
\begin{pgfscope}%
\pgfpathrectangle{\pgfqpoint{0.100000in}{0.100000in}}{\pgfqpoint{2.862985in}{1.832778in}}%
\pgfusepath{clip}%
\pgfsetbuttcap%
\pgfsetmiterjoin%
\definecolor{currentfill}{rgb}{0.917647,0.962630,0.851672}%
\pgfsetfillcolor{currentfill}%
\pgfsetlinewidth{0.000000pt}%
\definecolor{currentstroke}{rgb}{0.000000,0.000000,0.000000}%
\pgfsetstrokecolor{currentstroke}%
\pgfsetstrokeopacity{0.000000}%
\pgfsetdash{}{0pt}%
\pgfpathmoveto{\pgfqpoint{0.540358in}{1.430483in}}%
\pgfpathlineto{\pgfqpoint{0.571046in}{1.421943in}}%
\pgfpathlineto{\pgfqpoint{0.621916in}{1.408826in}}%
\pgfpathlineto{\pgfqpoint{0.674647in}{1.395314in}}%
\pgfpathlineto{\pgfqpoint{0.723250in}{1.383593in}}%
\pgfpathlineto{\pgfqpoint{0.765246in}{1.374085in}}%
\pgfpathlineto{\pgfqpoint{0.810177in}{1.364205in}}%
\pgfpathlineto{\pgfqpoint{0.797195in}{1.302942in}}%
\pgfpathlineto{\pgfqpoint{0.785630in}{1.248537in}}%
\pgfpathlineto{\pgfqpoint{0.766658in}{1.160754in}}%
\pgfpathlineto{\pgfqpoint{0.752414in}{1.094498in}}%
\pgfpathlineto{\pgfqpoint{0.744719in}{1.057523in}}%
\pgfpathlineto{\pgfqpoint{0.734852in}{1.009534in}}%
\pgfpathlineto{\pgfqpoint{0.728025in}{0.999798in}}%
\pgfpathlineto{\pgfqpoint{0.722531in}{0.999473in}}%
\pgfpathlineto{\pgfqpoint{0.718612in}{1.007971in}}%
\pgfpathlineto{\pgfqpoint{0.709613in}{1.011057in}}%
\pgfpathlineto{\pgfqpoint{0.699916in}{1.009976in}}%
\pgfpathlineto{\pgfqpoint{0.697164in}{1.003024in}}%
\pgfpathlineto{\pgfqpoint{0.696877in}{0.978886in}}%
\pgfpathlineto{\pgfqpoint{0.693973in}{0.973374in}}%
\pgfpathlineto{\pgfqpoint{0.695633in}{0.953995in}}%
\pgfpathlineto{\pgfqpoint{0.689569in}{0.941085in}}%
\pgfpathlineto{\pgfqpoint{0.640397in}{1.016802in}}%
\pgfpathlineto{\pgfqpoint{0.591971in}{1.090379in}}%
\pgfpathlineto{\pgfqpoint{0.566763in}{1.128960in}}%
\pgfpathlineto{\pgfqpoint{0.545735in}{1.162262in}}%
\pgfpathlineto{\pgfqpoint{0.519238in}{1.203456in}}%
\pgfpathlineto{\pgfqpoint{0.489244in}{1.249665in}}%
\pgfpathlineto{\pgfqpoint{0.501576in}{1.293524in}}%
\pgfpathlineto{\pgfqpoint{0.526391in}{1.381486in}}%
\pgfpathlineto{\pgfqpoint{0.540358in}{1.430483in}}%
\pgfpathclose%
\pgfusepath{fill}%
\end{pgfscope}%
\begin{pgfscope}%
\pgfpathrectangle{\pgfqpoint{0.100000in}{0.100000in}}{\pgfqpoint{2.862985in}{1.832778in}}%
\pgfusepath{clip}%
\pgfsetbuttcap%
\pgfsetmiterjoin%
\definecolor{currentfill}{rgb}{0.451903,0.698501,0.225375}%
\pgfsetfillcolor{currentfill}%
\pgfsetlinewidth{0.000000pt}%
\definecolor{currentstroke}{rgb}{0.000000,0.000000,0.000000}%
\pgfsetstrokecolor{currentstroke}%
\pgfsetstrokeopacity{0.000000}%
\pgfsetdash{}{0pt}%
\pgfpathmoveto{\pgfqpoint{0.810177in}{1.364205in}}%
\pgfpathlineto{\pgfqpoint{0.858212in}{1.354650in}}%
\pgfpathlineto{\pgfqpoint{0.947141in}{1.337483in}}%
\pgfpathlineto{\pgfqpoint{0.936009in}{1.275734in}}%
\pgfpathlineto{\pgfqpoint{0.984967in}{1.267352in}}%
\pgfpathlineto{\pgfqpoint{1.029571in}{1.260214in}}%
\pgfpathlineto{\pgfqpoint{1.021820in}{1.211396in}}%
\pgfpathlineto{\pgfqpoint{1.013608in}{1.158753in}}%
\pgfpathlineto{\pgfqpoint{1.002613in}{1.089632in}}%
\pgfpathlineto{\pgfqpoint{1.002328in}{1.083837in}}%
\pgfpathlineto{\pgfqpoint{0.990915in}{1.012202in}}%
\pgfpathlineto{\pgfqpoint{0.943932in}{1.019484in}}%
\pgfpathlineto{\pgfqpoint{0.919991in}{1.024292in}}%
\pgfpathlineto{\pgfqpoint{0.833455in}{1.039504in}}%
\pgfpathlineto{\pgfqpoint{0.800868in}{1.045930in}}%
\pgfpathlineto{\pgfqpoint{0.744719in}{1.057523in}}%
\pgfpathlineto{\pgfqpoint{0.752414in}{1.094498in}}%
\pgfpathlineto{\pgfqpoint{0.766658in}{1.160754in}}%
\pgfpathlineto{\pgfqpoint{0.785630in}{1.248537in}}%
\pgfpathlineto{\pgfqpoint{0.797195in}{1.302942in}}%
\pgfpathlineto{\pgfqpoint{0.810177in}{1.364205in}}%
\pgfpathclose%
\pgfusepath{fill}%
\end{pgfscope}%
\begin{pgfscope}%
\pgfpathrectangle{\pgfqpoint{0.100000in}{0.100000in}}{\pgfqpoint{2.862985in}{1.832778in}}%
\pgfusepath{clip}%
\pgfsetbuttcap%
\pgfsetmiterjoin%
\definecolor{currentfill}{rgb}{0.922876,0.963245,0.863668}%
\pgfsetfillcolor{currentfill}%
\pgfsetlinewidth{0.000000pt}%
\definecolor{currentstroke}{rgb}{0.000000,0.000000,0.000000}%
\pgfsetstrokecolor{currentstroke}%
\pgfsetstrokeopacity{0.000000}%
\pgfsetdash{}{0pt}%
\pgfpathmoveto{\pgfqpoint{0.352592in}{1.488751in}}%
\pgfpathlineto{\pgfqpoint{0.369661in}{1.482798in}}%
\pgfpathlineto{\pgfqpoint{0.438032in}{1.461511in}}%
\pgfpathlineto{\pgfqpoint{0.500927in}{1.441715in}}%
\pgfpathlineto{\pgfqpoint{0.540358in}{1.430483in}}%
\pgfpathlineto{\pgfqpoint{0.526391in}{1.381486in}}%
\pgfpathlineto{\pgfqpoint{0.501576in}{1.293524in}}%
\pgfpathlineto{\pgfqpoint{0.489244in}{1.249665in}}%
\pgfpathlineto{\pgfqpoint{0.519238in}{1.203456in}}%
\pgfpathlineto{\pgfqpoint{0.545735in}{1.162262in}}%
\pgfpathlineto{\pgfqpoint{0.566763in}{1.128960in}}%
\pgfpathlineto{\pgfqpoint{0.591971in}{1.090379in}}%
\pgfpathlineto{\pgfqpoint{0.640397in}{1.016802in}}%
\pgfpathlineto{\pgfqpoint{0.689569in}{0.941085in}}%
\pgfpathlineto{\pgfqpoint{0.687595in}{0.933576in}}%
\pgfpathlineto{\pgfqpoint{0.693526in}{0.921617in}}%
\pgfpathlineto{\pgfqpoint{0.694699in}{0.905259in}}%
\pgfpathlineto{\pgfqpoint{0.703313in}{0.897324in}}%
\pgfpathlineto{\pgfqpoint{0.703654in}{0.890930in}}%
\pgfpathlineto{\pgfqpoint{0.688230in}{0.883673in}}%
\pgfpathlineto{\pgfqpoint{0.680882in}{0.876407in}}%
\pgfpathlineto{\pgfqpoint{0.678591in}{0.860364in}}%
\pgfpathlineto{\pgfqpoint{0.674871in}{0.851614in}}%
\pgfpathlineto{\pgfqpoint{0.667035in}{0.844239in}}%
\pgfpathlineto{\pgfqpoint{0.659240in}{0.825061in}}%
\pgfpathlineto{\pgfqpoint{0.670200in}{0.815071in}}%
\pgfpathlineto{\pgfqpoint{0.668785in}{0.806836in}}%
\pgfpathlineto{\pgfqpoint{0.659798in}{0.801105in}}%
\pgfpathlineto{\pgfqpoint{0.653603in}{0.802117in}}%
\pgfpathlineto{\pgfqpoint{0.580470in}{0.812266in}}%
\pgfpathlineto{\pgfqpoint{0.526454in}{0.819928in}}%
\pgfpathlineto{\pgfqpoint{0.528816in}{0.828649in}}%
\pgfpathlineto{\pgfqpoint{0.525300in}{0.843128in}}%
\pgfpathlineto{\pgfqpoint{0.524945in}{0.857733in}}%
\pgfpathlineto{\pgfqpoint{0.522647in}{0.866279in}}%
\pgfpathlineto{\pgfqpoint{0.515568in}{0.878498in}}%
\pgfpathlineto{\pgfqpoint{0.495206in}{0.906691in}}%
\pgfpathlineto{\pgfqpoint{0.488537in}{0.910145in}}%
\pgfpathlineto{\pgfqpoint{0.479944in}{0.910090in}}%
\pgfpathlineto{\pgfqpoint{0.481908in}{0.918984in}}%
\pgfpathlineto{\pgfqpoint{0.477839in}{0.930106in}}%
\pgfpathlineto{\pgfqpoint{0.457853in}{0.935633in}}%
\pgfpathlineto{\pgfqpoint{0.445671in}{0.945869in}}%
\pgfpathlineto{\pgfqpoint{0.444661in}{0.952132in}}%
\pgfpathlineto{\pgfqpoint{0.430633in}{0.967637in}}%
\pgfpathlineto{\pgfqpoint{0.417280in}{0.970604in}}%
\pgfpathlineto{\pgfqpoint{0.404908in}{0.978484in}}%
\pgfpathlineto{\pgfqpoint{0.388637in}{0.981210in}}%
\pgfpathlineto{\pgfqpoint{0.381685in}{0.991719in}}%
\pgfpathlineto{\pgfqpoint{0.388287in}{1.008354in}}%
\pgfpathlineto{\pgfqpoint{0.386268in}{1.012095in}}%
\pgfpathlineto{\pgfqpoint{0.391766in}{1.025962in}}%
\pgfpathlineto{\pgfqpoint{0.381986in}{1.033342in}}%
\pgfpathlineto{\pgfqpoint{0.385159in}{1.046721in}}%
\pgfpathlineto{\pgfqpoint{0.379935in}{1.050141in}}%
\pgfpathlineto{\pgfqpoint{0.375424in}{1.062787in}}%
\pgfpathlineto{\pgfqpoint{0.369981in}{1.066644in}}%
\pgfpathlineto{\pgfqpoint{0.369572in}{1.075778in}}%
\pgfpathlineto{\pgfqpoint{0.365307in}{1.082219in}}%
\pgfpathlineto{\pgfqpoint{0.358881in}{1.103911in}}%
\pgfpathlineto{\pgfqpoint{0.351850in}{1.114301in}}%
\pgfpathlineto{\pgfqpoint{0.353378in}{1.131939in}}%
\pgfpathlineto{\pgfqpoint{0.361647in}{1.133714in}}%
\pgfpathlineto{\pgfqpoint{0.367031in}{1.143280in}}%
\pgfpathlineto{\pgfqpoint{0.363778in}{1.153656in}}%
\pgfpathlineto{\pgfqpoint{0.354984in}{1.155388in}}%
\pgfpathlineto{\pgfqpoint{0.350596in}{1.160241in}}%
\pgfpathlineto{\pgfqpoint{0.343491in}{1.178092in}}%
\pgfpathlineto{\pgfqpoint{0.346812in}{1.184503in}}%
\pgfpathlineto{\pgfqpoint{0.344457in}{1.196440in}}%
\pgfpathlineto{\pgfqpoint{0.349672in}{1.211900in}}%
\pgfpathlineto{\pgfqpoint{0.355782in}{1.206207in}}%
\pgfpathlineto{\pgfqpoint{0.353009in}{1.199485in}}%
\pgfpathlineto{\pgfqpoint{0.363597in}{1.188830in}}%
\pgfpathlineto{\pgfqpoint{0.362921in}{1.204652in}}%
\pgfpathlineto{\pgfqpoint{0.358384in}{1.208896in}}%
\pgfpathlineto{\pgfqpoint{0.363570in}{1.222808in}}%
\pgfpathlineto{\pgfqpoint{0.383185in}{1.220593in}}%
\pgfpathlineto{\pgfqpoint{0.380545in}{1.225766in}}%
\pgfpathlineto{\pgfqpoint{0.367599in}{1.225257in}}%
\pgfpathlineto{\pgfqpoint{0.361441in}{1.233093in}}%
\pgfpathlineto{\pgfqpoint{0.353689in}{1.226135in}}%
\pgfpathlineto{\pgfqpoint{0.349576in}{1.214504in}}%
\pgfpathlineto{\pgfqpoint{0.338597in}{1.230154in}}%
\pgfpathlineto{\pgfqpoint{0.334369in}{1.233003in}}%
\pgfpathlineto{\pgfqpoint{0.335980in}{1.250036in}}%
\pgfpathlineto{\pgfqpoint{0.332634in}{1.260081in}}%
\pgfpathlineto{\pgfqpoint{0.326568in}{1.269517in}}%
\pgfpathlineto{\pgfqpoint{0.314145in}{1.298432in}}%
\pgfpathlineto{\pgfqpoint{0.318207in}{1.304826in}}%
\pgfpathlineto{\pgfqpoint{0.318159in}{1.325033in}}%
\pgfpathlineto{\pgfqpoint{0.324864in}{1.336303in}}%
\pgfpathlineto{\pgfqpoint{0.326401in}{1.353918in}}%
\pgfpathlineto{\pgfqpoint{0.320079in}{1.374099in}}%
\pgfpathlineto{\pgfqpoint{0.311692in}{1.386945in}}%
\pgfpathlineto{\pgfqpoint{0.313188in}{1.398536in}}%
\pgfpathlineto{\pgfqpoint{0.336740in}{1.426591in}}%
\pgfpathlineto{\pgfqpoint{0.337910in}{1.436161in}}%
\pgfpathlineto{\pgfqpoint{0.348523in}{1.454416in}}%
\pgfpathlineto{\pgfqpoint{0.349995in}{1.471719in}}%
\pgfpathlineto{\pgfqpoint{0.346583in}{1.476134in}}%
\pgfpathlineto{\pgfqpoint{0.352592in}{1.488751in}}%
\pgfpathclose%
\pgfusepath{fill}%
\end{pgfscope}%
\begin{pgfscope}%
\pgfpathrectangle{\pgfqpoint{0.100000in}{0.100000in}}{\pgfqpoint{2.862985in}{1.832778in}}%
\pgfusepath{clip}%
\pgfsetbuttcap%
\pgfsetmiterjoin%
\definecolor{currentfill}{rgb}{0.541869,0.765705,0.307958}%
\pgfsetfillcolor{currentfill}%
\pgfsetlinewidth{0.000000pt}%
\definecolor{currentstroke}{rgb}{0.000000,0.000000,0.000000}%
\pgfsetstrokecolor{currentstroke}%
\pgfsetstrokeopacity{0.000000}%
\pgfsetdash{}{0pt}%
\pgfpathmoveto{\pgfqpoint{2.186981in}{1.117748in}}%
\pgfpathlineto{\pgfqpoint{2.179011in}{1.195770in}}%
\pgfpathlineto{\pgfqpoint{2.169752in}{1.279623in}}%
\pgfpathlineto{\pgfqpoint{2.230390in}{1.288653in}}%
\pgfpathlineto{\pgfqpoint{2.246483in}{1.284465in}}%
\pgfpathlineto{\pgfqpoint{2.254202in}{1.279918in}}%
\pgfpathlineto{\pgfqpoint{2.263874in}{1.281179in}}%
\pgfpathlineto{\pgfqpoint{2.276627in}{1.273651in}}%
\pgfpathlineto{\pgfqpoint{2.300380in}{1.284855in}}%
\pgfpathlineto{\pgfqpoint{2.313489in}{1.285229in}}%
\pgfpathlineto{\pgfqpoint{2.328797in}{1.302307in}}%
\pgfpathlineto{\pgfqpoint{2.344372in}{1.312695in}}%
\pgfpathlineto{\pgfqpoint{2.365146in}{1.324661in}}%
\pgfpathlineto{\pgfqpoint{2.378519in}{1.241018in}}%
\pgfpathlineto{\pgfqpoint{2.372306in}{1.235646in}}%
\pgfpathlineto{\pgfqpoint{2.376317in}{1.230712in}}%
\pgfpathlineto{\pgfqpoint{2.377915in}{1.219896in}}%
\pgfpathlineto{\pgfqpoint{2.374803in}{1.209737in}}%
\pgfpathlineto{\pgfqpoint{2.374391in}{1.194887in}}%
\pgfpathlineto{\pgfqpoint{2.371468in}{1.175560in}}%
\pgfpathlineto{\pgfqpoint{2.356678in}{1.158322in}}%
\pgfpathlineto{\pgfqpoint{2.350452in}{1.154663in}}%
\pgfpathlineto{\pgfqpoint{2.345610in}{1.157813in}}%
\pgfpathlineto{\pgfqpoint{2.333627in}{1.141383in}}%
\pgfpathlineto{\pgfqpoint{2.334750in}{1.125564in}}%
\pgfpathlineto{\pgfqpoint{2.321406in}{1.129368in}}%
\pgfpathlineto{\pgfqpoint{2.315087in}{1.113569in}}%
\pgfpathlineto{\pgfqpoint{2.318281in}{1.102373in}}%
\pgfpathlineto{\pgfqpoint{2.313281in}{1.100719in}}%
\pgfpathlineto{\pgfqpoint{2.312581in}{1.091873in}}%
\pgfpathlineto{\pgfqpoint{2.300308in}{1.088305in}}%
\pgfpathlineto{\pgfqpoint{2.293930in}{1.095447in}}%
\pgfpathlineto{\pgfqpoint{2.285724in}{1.098217in}}%
\pgfpathlineto{\pgfqpoint{2.283796in}{1.105459in}}%
\pgfpathlineto{\pgfqpoint{2.276376in}{1.104185in}}%
\pgfpathlineto{\pgfqpoint{2.271515in}{1.097524in}}%
\pgfpathlineto{\pgfqpoint{2.264561in}{1.095181in}}%
\pgfpathlineto{\pgfqpoint{2.252277in}{1.099895in}}%
\pgfpathlineto{\pgfqpoint{2.245483in}{1.094286in}}%
\pgfpathlineto{\pgfqpoint{2.235827in}{1.100922in}}%
\pgfpathlineto{\pgfqpoint{2.217303in}{1.102959in}}%
\pgfpathlineto{\pgfqpoint{2.211715in}{1.115014in}}%
\pgfpathlineto{\pgfqpoint{2.204636in}{1.120355in}}%
\pgfpathlineto{\pgfqpoint{2.197773in}{1.116936in}}%
\pgfpathlineto{\pgfqpoint{2.186981in}{1.117748in}}%
\pgfpathclose%
\pgfusepath{fill}%
\end{pgfscope}%
\begin{pgfscope}%
\pgfpathrectangle{\pgfqpoint{0.100000in}{0.100000in}}{\pgfqpoint{2.862985in}{1.832778in}}%
\pgfusepath{clip}%
\pgfsetbuttcap%
\pgfsetmiterjoin%
\definecolor{currentfill}{rgb}{0.824145,0.926951,0.690504}%
\pgfsetfillcolor{currentfill}%
\pgfsetlinewidth{0.000000pt}%
\definecolor{currentstroke}{rgb}{0.000000,0.000000,0.000000}%
\pgfsetstrokecolor{currentstroke}%
\pgfsetstrokeopacity{0.000000}%
\pgfsetdash{}{0pt}%
\pgfpathmoveto{\pgfqpoint{1.986225in}{0.966856in}}%
\pgfpathlineto{\pgfqpoint{1.979121in}{0.972628in}}%
\pgfpathlineto{\pgfqpoint{1.973328in}{0.969882in}}%
\pgfpathlineto{\pgfqpoint{1.965922in}{0.983711in}}%
\pgfpathlineto{\pgfqpoint{1.969705in}{0.992402in}}%
\pgfpathlineto{\pgfqpoint{1.964252in}{1.002186in}}%
\pgfpathlineto{\pgfqpoint{1.963959in}{1.009888in}}%
\pgfpathlineto{\pgfqpoint{1.956591in}{1.012632in}}%
\pgfpathlineto{\pgfqpoint{1.953174in}{1.018436in}}%
\pgfpathlineto{\pgfqpoint{1.938768in}{1.025684in}}%
\pgfpathlineto{\pgfqpoint{1.926256in}{1.034614in}}%
\pgfpathlineto{\pgfqpoint{1.920442in}{1.041355in}}%
\pgfpathlineto{\pgfqpoint{1.920322in}{1.049578in}}%
\pgfpathlineto{\pgfqpoint{1.928102in}{1.065365in}}%
\pgfpathlineto{\pgfqpoint{1.930495in}{1.077440in}}%
\pgfpathlineto{\pgfqpoint{1.924157in}{1.084263in}}%
\pgfpathlineto{\pgfqpoint{1.915763in}{1.086834in}}%
\pgfpathlineto{\pgfqpoint{1.905579in}{1.081212in}}%
\pgfpathlineto{\pgfqpoint{1.901269in}{1.090863in}}%
\pgfpathlineto{\pgfqpoint{1.899094in}{1.103945in}}%
\pgfpathlineto{\pgfqpoint{1.884086in}{1.115607in}}%
\pgfpathlineto{\pgfqpoint{1.881087in}{1.120782in}}%
\pgfpathlineto{\pgfqpoint{1.867379in}{1.132498in}}%
\pgfpathlineto{\pgfqpoint{1.863083in}{1.141015in}}%
\pgfpathlineto{\pgfqpoint{1.859210in}{1.157906in}}%
\pgfpathlineto{\pgfqpoint{1.861865in}{1.172908in}}%
\pgfpathlineto{\pgfqpoint{1.865408in}{1.174999in}}%
\pgfpathlineto{\pgfqpoint{1.864771in}{1.187558in}}%
\pgfpathlineto{\pgfqpoint{1.874773in}{1.191288in}}%
\pgfpathlineto{\pgfqpoint{1.877793in}{1.202558in}}%
\pgfpathlineto{\pgfqpoint{1.883547in}{1.210143in}}%
\pgfpathlineto{\pgfqpoint{1.883256in}{1.219794in}}%
\pgfpathlineto{\pgfqpoint{1.876124in}{1.227469in}}%
\pgfpathlineto{\pgfqpoint{1.877818in}{1.238212in}}%
\pgfpathlineto{\pgfqpoint{1.896306in}{1.242884in}}%
\pgfpathlineto{\pgfqpoint{1.910490in}{1.251413in}}%
\pgfpathlineto{\pgfqpoint{1.911979in}{1.262151in}}%
\pgfpathlineto{\pgfqpoint{1.916914in}{1.265532in}}%
\pgfpathlineto{\pgfqpoint{1.918802in}{1.276813in}}%
\pgfpathlineto{\pgfqpoint{1.917285in}{1.284265in}}%
\pgfpathlineto{\pgfqpoint{1.907588in}{1.290454in}}%
\pgfpathlineto{\pgfqpoint{1.903684in}{1.299681in}}%
\pgfpathlineto{\pgfqpoint{1.894114in}{1.308600in}}%
\pgfpathlineto{\pgfqpoint{1.972758in}{1.311952in}}%
\pgfpathlineto{\pgfqpoint{2.025520in}{1.315700in}}%
\pgfpathlineto{\pgfqpoint{2.024554in}{1.304601in}}%
\pgfpathlineto{\pgfqpoint{2.033546in}{1.289293in}}%
\pgfpathlineto{\pgfqpoint{2.037322in}{1.276213in}}%
\pgfpathlineto{\pgfqpoint{2.041818in}{1.268787in}}%
\pgfpathlineto{\pgfqpoint{2.046726in}{1.207597in}}%
\pgfpathlineto{\pgfqpoint{2.053542in}{1.120373in}}%
\pgfpathlineto{\pgfqpoint{2.048283in}{1.107118in}}%
\pgfpathlineto{\pgfqpoint{2.055859in}{1.096224in}}%
\pgfpathlineto{\pgfqpoint{2.057971in}{1.085818in}}%
\pgfpathlineto{\pgfqpoint{2.047524in}{1.064241in}}%
\pgfpathlineto{\pgfqpoint{2.048140in}{1.061261in}}%
\pgfpathlineto{\pgfqpoint{2.037387in}{1.048893in}}%
\pgfpathlineto{\pgfqpoint{2.040336in}{1.045116in}}%
\pgfpathlineto{\pgfqpoint{2.035862in}{1.037407in}}%
\pgfpathlineto{\pgfqpoint{2.037013in}{1.021868in}}%
\pgfpathlineto{\pgfqpoint{2.031580in}{1.012333in}}%
\pgfpathlineto{\pgfqpoint{2.035999in}{1.001067in}}%
\pgfpathlineto{\pgfqpoint{2.017483in}{0.994944in}}%
\pgfpathlineto{\pgfqpoint{2.015777in}{0.988286in}}%
\pgfpathlineto{\pgfqpoint{2.020833in}{0.979847in}}%
\pgfpathlineto{\pgfqpoint{2.018506in}{0.974346in}}%
\pgfpathlineto{\pgfqpoint{1.998652in}{0.981141in}}%
\pgfpathlineto{\pgfqpoint{1.992101in}{0.981825in}}%
\pgfpathlineto{\pgfqpoint{1.983934in}{0.971495in}}%
\pgfpathlineto{\pgfqpoint{1.986225in}{0.966856in}}%
\pgfpathclose%
\pgfusepath{fill}%
\end{pgfscope}%
\begin{pgfscope}%
\pgfpathrectangle{\pgfqpoint{0.100000in}{0.100000in}}{\pgfqpoint{2.862985in}{1.832778in}}%
\pgfusepath{clip}%
\pgfsetbuttcap%
\pgfsetmiterjoin%
\definecolor{currentfill}{rgb}{0.917647,0.962630,0.851672}%
\pgfsetfillcolor{currentfill}%
\pgfsetlinewidth{0.000000pt}%
\definecolor{currentstroke}{rgb}{0.000000,0.000000,0.000000}%
\pgfsetstrokecolor{currentstroke}%
\pgfsetstrokeopacity{0.000000}%
\pgfsetdash{}{0pt}%
\pgfpathmoveto{\pgfqpoint{2.564161in}{1.156689in}}%
\pgfpathlineto{\pgfqpoint{2.558697in}{1.164803in}}%
\pgfpathlineto{\pgfqpoint{2.561780in}{1.169343in}}%
\pgfpathlineto{\pgfqpoint{2.569331in}{1.164244in}}%
\pgfpathlineto{\pgfqpoint{2.564161in}{1.156689in}}%
\pgfpathclose%
\pgfusepath{fill}%
\end{pgfscope}%
\begin{pgfscope}%
\pgfpathrectangle{\pgfqpoint{0.100000in}{0.100000in}}{\pgfqpoint{2.862985in}{1.832778in}}%
\pgfusepath{clip}%
\pgfsetbuttcap%
\pgfsetmiterjoin%
\definecolor{currentfill}{rgb}{0.739254,0.890042,0.553941}%
\pgfsetfillcolor{currentfill}%
\pgfsetlinewidth{0.000000pt}%
\definecolor{currentstroke}{rgb}{0.000000,0.000000,0.000000}%
\pgfsetstrokecolor{currentstroke}%
\pgfsetstrokeopacity{0.000000}%
\pgfsetdash{}{0pt}%
\pgfpathmoveto{\pgfqpoint{2.611873in}{1.226505in}}%
\pgfpathlineto{\pgfqpoint{2.615138in}{1.233404in}}%
\pgfpathlineto{\pgfqpoint{2.628338in}{1.234894in}}%
\pgfpathlineto{\pgfqpoint{2.626223in}{1.229031in}}%
\pgfpathlineto{\pgfqpoint{2.621868in}{1.221542in}}%
\pgfpathlineto{\pgfqpoint{2.624819in}{1.212619in}}%
\pgfpathlineto{\pgfqpoint{2.636465in}{1.201927in}}%
\pgfpathlineto{\pgfqpoint{2.639168in}{1.190668in}}%
\pgfpathlineto{\pgfqpoint{2.652585in}{1.176648in}}%
\pgfpathlineto{\pgfqpoint{2.657849in}{1.177248in}}%
\pgfpathlineto{\pgfqpoint{2.664406in}{1.156201in}}%
\pgfpathlineto{\pgfqpoint{2.663331in}{1.155991in}}%
\pgfpathlineto{\pgfqpoint{2.662135in}{1.155756in}}%
\pgfpathlineto{\pgfqpoint{2.632891in}{1.150158in}}%
\pgfpathlineto{\pgfqpoint{2.617245in}{1.205802in}}%
\pgfpathlineto{\pgfqpoint{2.611873in}{1.226505in}}%
\pgfpathclose%
\pgfusepath{fill}%
\end{pgfscope}%
\begin{pgfscope}%
\pgfpathrectangle{\pgfqpoint{0.100000in}{0.100000in}}{\pgfqpoint{2.862985in}{1.832778in}}%
\pgfusepath{clip}%
\pgfsetbuttcap%
\pgfsetmiterjoin%
\definecolor{currentfill}{rgb}{0.690888,0.862438,0.488351}%
\pgfsetfillcolor{currentfill}%
\pgfsetlinewidth{0.000000pt}%
\definecolor{currentstroke}{rgb}{0.000000,0.000000,0.000000}%
\pgfsetstrokecolor{currentstroke}%
\pgfsetstrokeopacity{0.000000}%
\pgfsetdash{}{0pt}%
\pgfpathmoveto{\pgfqpoint{2.338724in}{1.037893in}}%
\pgfpathlineto{\pgfqpoint{2.329483in}{1.038233in}}%
\pgfpathlineto{\pgfqpoint{2.320968in}{1.044098in}}%
\pgfpathlineto{\pgfqpoint{2.309465in}{1.061926in}}%
\pgfpathlineto{\pgfqpoint{2.299672in}{1.071367in}}%
\pgfpathlineto{\pgfqpoint{2.302231in}{1.078659in}}%
\pgfpathlineto{\pgfqpoint{2.300308in}{1.088305in}}%
\pgfpathlineto{\pgfqpoint{2.312581in}{1.091873in}}%
\pgfpathlineto{\pgfqpoint{2.313281in}{1.100719in}}%
\pgfpathlineto{\pgfqpoint{2.318281in}{1.102373in}}%
\pgfpathlineto{\pgfqpoint{2.315087in}{1.113569in}}%
\pgfpathlineto{\pgfqpoint{2.321406in}{1.129368in}}%
\pgfpathlineto{\pgfqpoint{2.334750in}{1.125564in}}%
\pgfpathlineto{\pgfqpoint{2.333627in}{1.141383in}}%
\pgfpathlineto{\pgfqpoint{2.345610in}{1.157813in}}%
\pgfpathlineto{\pgfqpoint{2.350452in}{1.154663in}}%
\pgfpathlineto{\pgfqpoint{2.356678in}{1.158322in}}%
\pgfpathlineto{\pgfqpoint{2.371468in}{1.175560in}}%
\pgfpathlineto{\pgfqpoint{2.374391in}{1.194887in}}%
\pgfpathlineto{\pgfqpoint{2.374803in}{1.209737in}}%
\pgfpathlineto{\pgfqpoint{2.377915in}{1.219896in}}%
\pgfpathlineto{\pgfqpoint{2.376317in}{1.230712in}}%
\pgfpathlineto{\pgfqpoint{2.372306in}{1.235646in}}%
\pgfpathlineto{\pgfqpoint{2.378519in}{1.241018in}}%
\pgfpathlineto{\pgfqpoint{2.387526in}{1.184275in}}%
\pgfpathlineto{\pgfqpoint{2.437178in}{1.192485in}}%
\pgfpathlineto{\pgfqpoint{2.442325in}{1.160097in}}%
\pgfpathlineto{\pgfqpoint{2.460300in}{1.181474in}}%
\pgfpathlineto{\pgfqpoint{2.464536in}{1.179343in}}%
\pgfpathlineto{\pgfqpoint{2.470767in}{1.191730in}}%
\pgfpathlineto{\pgfqpoint{2.479399in}{1.187846in}}%
\pgfpathlineto{\pgfqpoint{2.486925in}{1.188584in}}%
\pgfpathlineto{\pgfqpoint{2.491891in}{1.197195in}}%
\pgfpathlineto{\pgfqpoint{2.499098in}{1.201985in}}%
\pgfpathlineto{\pgfqpoint{2.509119in}{1.197766in}}%
\pgfpathlineto{\pgfqpoint{2.515712in}{1.199224in}}%
\pgfpathlineto{\pgfqpoint{2.525158in}{1.182877in}}%
\pgfpathlineto{\pgfqpoint{2.522427in}{1.170500in}}%
\pgfpathlineto{\pgfqpoint{2.497755in}{1.184435in}}%
\pgfpathlineto{\pgfqpoint{2.493134in}{1.172996in}}%
\pgfpathlineto{\pgfqpoint{2.494659in}{1.167732in}}%
\pgfpathlineto{\pgfqpoint{2.484159in}{1.149941in}}%
\pgfpathlineto{\pgfqpoint{2.479203in}{1.146402in}}%
\pgfpathlineto{\pgfqpoint{2.476965in}{1.138556in}}%
\pgfpathlineto{\pgfqpoint{2.468517in}{1.139377in}}%
\pgfpathlineto{\pgfqpoint{2.462449in}{1.117963in}}%
\pgfpathlineto{\pgfqpoint{2.459053in}{1.113049in}}%
\pgfpathlineto{\pgfqpoint{2.450326in}{1.114688in}}%
\pgfpathlineto{\pgfqpoint{2.441399in}{1.121442in}}%
\pgfpathlineto{\pgfqpoint{2.441091in}{1.111076in}}%
\pgfpathlineto{\pgfqpoint{2.432449in}{1.093654in}}%
\pgfpathlineto{\pgfqpoint{2.430300in}{1.081251in}}%
\pgfpathlineto{\pgfqpoint{2.423669in}{1.072991in}}%
\pgfpathlineto{\pgfqpoint{2.418643in}{1.059800in}}%
\pgfpathlineto{\pgfqpoint{2.418396in}{1.047589in}}%
\pgfpathlineto{\pgfqpoint{2.410633in}{1.045308in}}%
\pgfpathlineto{\pgfqpoint{2.401826in}{1.038395in}}%
\pgfpathlineto{\pgfqpoint{2.393347in}{1.037079in}}%
\pgfpathlineto{\pgfqpoint{2.391393in}{1.031228in}}%
\pgfpathlineto{\pgfqpoint{2.377730in}{1.025236in}}%
\pgfpathlineto{\pgfqpoint{2.370094in}{1.030341in}}%
\pgfpathlineto{\pgfqpoint{2.361560in}{1.020647in}}%
\pgfpathlineto{\pgfqpoint{2.346890in}{1.023505in}}%
\pgfpathlineto{\pgfqpoint{2.337893in}{1.033679in}}%
\pgfpathlineto{\pgfqpoint{2.338724in}{1.037893in}}%
\pgfpathclose%
\pgfusepath{fill}%
\end{pgfscope}%
\begin{pgfscope}%
\pgfpathrectangle{\pgfqpoint{0.100000in}{0.100000in}}{\pgfqpoint{2.862985in}{1.832778in}}%
\pgfusepath{clip}%
\pgfsetbuttcap%
\pgfsetmiterjoin%
\definecolor{currentfill}{rgb}{0.281507,0.547789,0.125106}%
\pgfsetfillcolor{currentfill}%
\pgfsetlinewidth{0.000000pt}%
\definecolor{currentstroke}{rgb}{0.000000,0.000000,0.000000}%
\pgfsetstrokecolor{currentstroke}%
\pgfsetstrokeopacity{0.000000}%
\pgfsetdash{}{0pt}%
\pgfpathmoveto{\pgfqpoint{2.662135in}{1.155756in}}%
\pgfpathlineto{\pgfqpoint{2.661760in}{1.144322in}}%
\pgfpathlineto{\pgfqpoint{2.657360in}{1.138701in}}%
\pgfpathlineto{\pgfqpoint{2.654551in}{1.126221in}}%
\pgfpathlineto{\pgfqpoint{2.641870in}{1.120469in}}%
\pgfpathlineto{\pgfqpoint{2.631229in}{1.118793in}}%
\pgfpathlineto{\pgfqpoint{2.634330in}{1.127000in}}%
\pgfpathlineto{\pgfqpoint{2.617947in}{1.133877in}}%
\pgfpathlineto{\pgfqpoint{2.604638in}{1.142487in}}%
\pgfpathlineto{\pgfqpoint{2.612878in}{1.155358in}}%
\pgfpathlineto{\pgfqpoint{2.606274in}{1.165363in}}%
\pgfpathlineto{\pgfqpoint{2.607051in}{1.173993in}}%
\pgfpathlineto{\pgfqpoint{2.602276in}{1.176362in}}%
\pgfpathlineto{\pgfqpoint{2.598402in}{1.190382in}}%
\pgfpathlineto{\pgfqpoint{2.607564in}{1.209388in}}%
\pgfpathlineto{\pgfqpoint{2.600681in}{1.212492in}}%
\pgfpathlineto{\pgfqpoint{2.598896in}{1.203115in}}%
\pgfpathlineto{\pgfqpoint{2.589070in}{1.200505in}}%
\pgfpathlineto{\pgfqpoint{2.590103in}{1.169655in}}%
\pgfpathlineto{\pgfqpoint{2.588306in}{1.159726in}}%
\pgfpathlineto{\pgfqpoint{2.593235in}{1.145575in}}%
\pgfpathlineto{\pgfqpoint{2.600818in}{1.138767in}}%
\pgfpathlineto{\pgfqpoint{2.597379in}{1.134492in}}%
\pgfpathlineto{\pgfqpoint{2.605119in}{1.128246in}}%
\pgfpathlineto{\pgfqpoint{2.607917in}{1.118107in}}%
\pgfpathlineto{\pgfqpoint{2.593730in}{1.126503in}}%
\pgfpathlineto{\pgfqpoint{2.584752in}{1.125412in}}%
\pgfpathlineto{\pgfqpoint{2.577780in}{1.134049in}}%
\pgfpathlineto{\pgfqpoint{2.570682in}{1.134890in}}%
\pgfpathlineto{\pgfqpoint{2.560594in}{1.130564in}}%
\pgfpathlineto{\pgfqpoint{2.556681in}{1.135955in}}%
\pgfpathlineto{\pgfqpoint{2.561823in}{1.147270in}}%
\pgfpathlineto{\pgfqpoint{2.564161in}{1.156689in}}%
\pgfpathlineto{\pgfqpoint{2.569331in}{1.164244in}}%
\pgfpathlineto{\pgfqpoint{2.561780in}{1.169343in}}%
\pgfpathlineto{\pgfqpoint{2.558697in}{1.164803in}}%
\pgfpathlineto{\pgfqpoint{2.551154in}{1.169411in}}%
\pgfpathlineto{\pgfqpoint{2.537796in}{1.172482in}}%
\pgfpathlineto{\pgfqpoint{2.539008in}{1.179234in}}%
\pgfpathlineto{\pgfqpoint{2.532953in}{1.183150in}}%
\pgfpathlineto{\pgfqpoint{2.525158in}{1.182877in}}%
\pgfpathlineto{\pgfqpoint{2.515712in}{1.199224in}}%
\pgfpathlineto{\pgfqpoint{2.509119in}{1.197766in}}%
\pgfpathlineto{\pgfqpoint{2.499098in}{1.201985in}}%
\pgfpathlineto{\pgfqpoint{2.491891in}{1.197195in}}%
\pgfpathlineto{\pgfqpoint{2.486925in}{1.188584in}}%
\pgfpathlineto{\pgfqpoint{2.479399in}{1.187846in}}%
\pgfpathlineto{\pgfqpoint{2.470767in}{1.191730in}}%
\pgfpathlineto{\pgfqpoint{2.464536in}{1.179343in}}%
\pgfpathlineto{\pgfqpoint{2.460300in}{1.181474in}}%
\pgfpathlineto{\pgfqpoint{2.442325in}{1.160097in}}%
\pgfpathlineto{\pgfqpoint{2.437178in}{1.192485in}}%
\pgfpathlineto{\pgfqpoint{2.502868in}{1.204659in}}%
\pgfpathlineto{\pgfqpoint{2.532316in}{1.209933in}}%
\pgfpathlineto{\pgfqpoint{2.575158in}{1.218561in}}%
\pgfpathlineto{\pgfqpoint{2.611873in}{1.226505in}}%
\pgfpathlineto{\pgfqpoint{2.617245in}{1.205802in}}%
\pgfpathlineto{\pgfqpoint{2.632891in}{1.150158in}}%
\pgfpathlineto{\pgfqpoint{2.662135in}{1.155756in}}%
\pgfpathclose%
\pgfusepath{fill}%
\end{pgfscope}%
\begin{pgfscope}%
\pgfpathrectangle{\pgfqpoint{0.100000in}{0.100000in}}{\pgfqpoint{2.862985in}{1.832778in}}%
\pgfusepath{clip}%
\pgfsetbuttcap%
\pgfsetmiterjoin%
\definecolor{currentfill}{rgb}{0.594464,0.799846,0.371626}%
\pgfsetfillcolor{currentfill}%
\pgfsetlinewidth{0.000000pt}%
\definecolor{currentstroke}{rgb}{0.000000,0.000000,0.000000}%
\pgfsetstrokecolor{currentstroke}%
\pgfsetstrokeopacity{0.000000}%
\pgfsetdash{}{0pt}%
\pgfpathmoveto{\pgfqpoint{1.339821in}{0.971059in}}%
\pgfpathlineto{\pgfqpoint{1.291895in}{0.975633in}}%
\pgfpathlineto{\pgfqpoint{1.242175in}{0.980039in}}%
\pgfpathlineto{\pgfqpoint{1.184722in}{0.986032in}}%
\pgfpathlineto{\pgfqpoint{1.099362in}{0.996184in}}%
\pgfpathlineto{\pgfqpoint{1.069080in}{1.000850in}}%
\pgfpathlineto{\pgfqpoint{0.990915in}{1.012202in}}%
\pgfpathlineto{\pgfqpoint{1.002328in}{1.083837in}}%
\pgfpathlineto{\pgfqpoint{1.002613in}{1.089632in}}%
\pgfpathlineto{\pgfqpoint{1.013608in}{1.158753in}}%
\pgfpathlineto{\pgfqpoint{1.021820in}{1.211396in}}%
\pgfpathlineto{\pgfqpoint{1.029571in}{1.260214in}}%
\pgfpathlineto{\pgfqpoint{1.082522in}{1.252613in}}%
\pgfpathlineto{\pgfqpoint{1.131892in}{1.245527in}}%
\pgfpathlineto{\pgfqpoint{1.222647in}{1.234326in}}%
\pgfpathlineto{\pgfqpoint{1.264283in}{1.230524in}}%
\pgfpathlineto{\pgfqpoint{1.330243in}{1.224124in}}%
\pgfpathlineto{\pgfqpoint{1.358777in}{1.221808in}}%
\pgfpathlineto{\pgfqpoint{1.353739in}{1.159333in}}%
\pgfpathlineto{\pgfqpoint{1.349188in}{1.099180in}}%
\pgfpathlineto{\pgfqpoint{1.345525in}{1.050201in}}%
\pgfpathlineto{\pgfqpoint{1.339821in}{0.971059in}}%
\pgfpathclose%
\pgfusepath{fill}%
\end{pgfscope}%
\begin{pgfscope}%
\pgfpathrectangle{\pgfqpoint{0.100000in}{0.100000in}}{\pgfqpoint{2.862985in}{1.832778in}}%
\pgfusepath{clip}%
\pgfsetbuttcap%
\pgfsetmiterjoin%
\definecolor{currentfill}{rgb}{0.760477,0.899270,0.588082}%
\pgfsetfillcolor{currentfill}%
\pgfsetlinewidth{0.000000pt}%
\definecolor{currentstroke}{rgb}{0.000000,0.000000,0.000000}%
\pgfsetstrokecolor{currentstroke}%
\pgfsetstrokeopacity{0.000000}%
\pgfsetdash{}{0pt}%
\pgfpathmoveto{\pgfqpoint{1.973847in}{0.935544in}}%
\pgfpathlineto{\pgfqpoint{1.975463in}{0.942786in}}%
\pgfpathlineto{\pgfqpoint{1.983837in}{0.941156in}}%
\pgfpathlineto{\pgfqpoint{1.986225in}{0.966856in}}%
\pgfpathlineto{\pgfqpoint{1.983934in}{0.971495in}}%
\pgfpathlineto{\pgfqpoint{1.992101in}{0.981825in}}%
\pgfpathlineto{\pgfqpoint{1.998652in}{0.981141in}}%
\pgfpathlineto{\pgfqpoint{2.018506in}{0.974346in}}%
\pgfpathlineto{\pgfqpoint{2.020833in}{0.979847in}}%
\pgfpathlineto{\pgfqpoint{2.015777in}{0.988286in}}%
\pgfpathlineto{\pgfqpoint{2.017483in}{0.994944in}}%
\pgfpathlineto{\pgfqpoint{2.035999in}{1.001067in}}%
\pgfpathlineto{\pgfqpoint{2.031580in}{1.012333in}}%
\pgfpathlineto{\pgfqpoint{2.037013in}{1.021868in}}%
\pgfpathlineto{\pgfqpoint{2.047127in}{1.027220in}}%
\pgfpathlineto{\pgfqpoint{2.068349in}{1.032520in}}%
\pgfpathlineto{\pgfqpoint{2.081739in}{1.024520in}}%
\pgfpathlineto{\pgfqpoint{2.085801in}{1.032321in}}%
\pgfpathlineto{\pgfqpoint{2.096919in}{1.037771in}}%
\pgfpathlineto{\pgfqpoint{2.099789in}{1.032965in}}%
\pgfpathlineto{\pgfqpoint{2.111209in}{1.036783in}}%
\pgfpathlineto{\pgfqpoint{2.110489in}{1.043318in}}%
\pgfpathlineto{\pgfqpoint{2.120757in}{1.050797in}}%
\pgfpathlineto{\pgfqpoint{2.126817in}{1.043031in}}%
\pgfpathlineto{\pgfqpoint{2.134790in}{1.042250in}}%
\pgfpathlineto{\pgfqpoint{2.140088in}{1.047311in}}%
\pgfpathlineto{\pgfqpoint{2.139498in}{1.054509in}}%
\pgfpathlineto{\pgfqpoint{2.144002in}{1.061656in}}%
\pgfpathlineto{\pgfqpoint{2.150037in}{1.063205in}}%
\pgfpathlineto{\pgfqpoint{2.152460in}{1.072642in}}%
\pgfpathlineto{\pgfqpoint{2.161237in}{1.080801in}}%
\pgfpathlineto{\pgfqpoint{2.158595in}{1.088922in}}%
\pgfpathlineto{\pgfqpoint{2.167134in}{1.092965in}}%
\pgfpathlineto{\pgfqpoint{2.172839in}{1.090478in}}%
\pgfpathlineto{\pgfqpoint{2.181265in}{1.096795in}}%
\pgfpathlineto{\pgfqpoint{2.188795in}{1.098441in}}%
\pgfpathlineto{\pgfqpoint{2.188822in}{1.104993in}}%
\pgfpathlineto{\pgfqpoint{2.183545in}{1.114085in}}%
\pgfpathlineto{\pgfqpoint{2.186981in}{1.117748in}}%
\pgfpathlineto{\pgfqpoint{2.197773in}{1.116936in}}%
\pgfpathlineto{\pgfqpoint{2.204636in}{1.120355in}}%
\pgfpathlineto{\pgfqpoint{2.211715in}{1.115014in}}%
\pgfpathlineto{\pgfqpoint{2.217303in}{1.102959in}}%
\pgfpathlineto{\pgfqpoint{2.235827in}{1.100922in}}%
\pgfpathlineto{\pgfqpoint{2.245483in}{1.094286in}}%
\pgfpathlineto{\pgfqpoint{2.252277in}{1.099895in}}%
\pgfpathlineto{\pgfqpoint{2.264561in}{1.095181in}}%
\pgfpathlineto{\pgfqpoint{2.271515in}{1.097524in}}%
\pgfpathlineto{\pgfqpoint{2.276376in}{1.104185in}}%
\pgfpathlineto{\pgfqpoint{2.283796in}{1.105459in}}%
\pgfpathlineto{\pgfqpoint{2.285724in}{1.098217in}}%
\pgfpathlineto{\pgfqpoint{2.293930in}{1.095447in}}%
\pgfpathlineto{\pgfqpoint{2.300308in}{1.088305in}}%
\pgfpathlineto{\pgfqpoint{2.302231in}{1.078659in}}%
\pgfpathlineto{\pgfqpoint{2.299672in}{1.071367in}}%
\pgfpathlineto{\pgfqpoint{2.309465in}{1.061926in}}%
\pgfpathlineto{\pgfqpoint{2.320968in}{1.044098in}}%
\pgfpathlineto{\pgfqpoint{2.329483in}{1.038233in}}%
\pgfpathlineto{\pgfqpoint{2.338724in}{1.037893in}}%
\pgfpathlineto{\pgfqpoint{2.321661in}{1.018314in}}%
\pgfpathlineto{\pgfqpoint{2.304865in}{1.006463in}}%
\pgfpathlineto{\pgfqpoint{2.298683in}{0.997053in}}%
\pgfpathlineto{\pgfqpoint{2.298788in}{0.991945in}}%
\pgfpathlineto{\pgfqpoint{2.289691in}{0.988030in}}%
\pgfpathlineto{\pgfqpoint{2.287084in}{0.980665in}}%
\pgfpathlineto{\pgfqpoint{2.268132in}{0.973246in}}%
\pgfpathlineto{\pgfqpoint{2.261414in}{0.968427in}}%
\pgfpathlineto{\pgfqpoint{2.260505in}{0.967399in}}%
\pgfpathlineto{\pgfqpoint{2.205991in}{0.962230in}}%
\pgfpathlineto{\pgfqpoint{2.173070in}{0.959461in}}%
\pgfpathlineto{\pgfqpoint{2.119049in}{0.956404in}}%
\pgfpathlineto{\pgfqpoint{2.051748in}{0.949746in}}%
\pgfpathlineto{\pgfqpoint{2.040629in}{0.951286in}}%
\pgfpathlineto{\pgfqpoint{2.042942in}{0.939936in}}%
\pgfpathlineto{\pgfqpoint{1.973847in}{0.935544in}}%
\pgfpathclose%
\pgfusepath{fill}%
\end{pgfscope}%
\begin{pgfscope}%
\pgfpathrectangle{\pgfqpoint{0.100000in}{0.100000in}}{\pgfqpoint{2.862985in}{1.832778in}}%
\pgfusepath{clip}%
\pgfsetbuttcap%
\pgfsetmiterjoin%
\definecolor{currentfill}{rgb}{0.344252,0.608074,0.156478}%
\pgfsetfillcolor{currentfill}%
\pgfsetlinewidth{0.000000pt}%
\definecolor{currentstroke}{rgb}{0.000000,0.000000,0.000000}%
\pgfsetstrokecolor{currentstroke}%
\pgfsetstrokeopacity{0.000000}%
\pgfsetdash{}{0pt}%
\pgfpathmoveto{\pgfqpoint{1.339821in}{0.971059in}}%
\pgfpathlineto{\pgfqpoint{1.345525in}{1.050201in}}%
\pgfpathlineto{\pgfqpoint{1.349188in}{1.099180in}}%
\pgfpathlineto{\pgfqpoint{1.353739in}{1.159333in}}%
\pgfpathlineto{\pgfqpoint{1.416835in}{1.154903in}}%
\pgfpathlineto{\pgfqpoint{1.496966in}{1.150526in}}%
\pgfpathlineto{\pgfqpoint{1.551466in}{1.148434in}}%
\pgfpathlineto{\pgfqpoint{1.605647in}{1.146774in}}%
\pgfpathlineto{\pgfqpoint{1.654686in}{1.146016in}}%
\pgfpathlineto{\pgfqpoint{1.677367in}{1.146251in}}%
\pgfpathlineto{\pgfqpoint{1.687350in}{1.138104in}}%
\pgfpathlineto{\pgfqpoint{1.695170in}{1.139746in}}%
\pgfpathlineto{\pgfqpoint{1.695414in}{1.132640in}}%
\pgfpathlineto{\pgfqpoint{1.689607in}{1.120348in}}%
\pgfpathlineto{\pgfqpoint{1.690237in}{1.112582in}}%
\pgfpathlineto{\pgfqpoint{1.695797in}{1.107461in}}%
\pgfpathlineto{\pgfqpoint{1.700905in}{1.096790in}}%
\pgfpathlineto{\pgfqpoint{1.711265in}{1.090665in}}%
\pgfpathlineto{\pgfqpoint{1.710919in}{1.050444in}}%
\pgfpathlineto{\pgfqpoint{1.711224in}{0.957954in}}%
\pgfpathlineto{\pgfqpoint{1.641779in}{0.958274in}}%
\pgfpathlineto{\pgfqpoint{1.568651in}{0.959554in}}%
\pgfpathlineto{\pgfqpoint{1.514818in}{0.961453in}}%
\pgfpathlineto{\pgfqpoint{1.469882in}{0.963187in}}%
\pgfpathlineto{\pgfqpoint{1.394179in}{0.967651in}}%
\pgfpathlineto{\pgfqpoint{1.339821in}{0.971059in}}%
\pgfpathclose%
\pgfusepath{fill}%
\end{pgfscope}%
\begin{pgfscope}%
\pgfpathrectangle{\pgfqpoint{0.100000in}{0.100000in}}{\pgfqpoint{2.862985in}{1.832778in}}%
\pgfusepath{clip}%
\pgfsetbuttcap%
\pgfsetmiterjoin%
\definecolor{currentfill}{rgb}{0.524337,0.754325,0.286736}%
\pgfsetfillcolor{currentfill}%
\pgfsetlinewidth{0.000000pt}%
\definecolor{currentstroke}{rgb}{0.000000,0.000000,0.000000}%
\pgfsetstrokecolor{currentstroke}%
\pgfsetstrokeopacity{0.000000}%
\pgfsetdash{}{0pt}%
\pgfpathmoveto{\pgfqpoint{2.361621in}{0.981538in}}%
\pgfpathlineto{\pgfqpoint{2.284707in}{0.970734in}}%
\pgfpathlineto{\pgfqpoint{2.261414in}{0.968427in}}%
\pgfpathlineto{\pgfqpoint{2.268132in}{0.973246in}}%
\pgfpathlineto{\pgfqpoint{2.287084in}{0.980665in}}%
\pgfpathlineto{\pgfqpoint{2.289691in}{0.988030in}}%
\pgfpathlineto{\pgfqpoint{2.298788in}{0.991945in}}%
\pgfpathlineto{\pgfqpoint{2.298683in}{0.997053in}}%
\pgfpathlineto{\pgfqpoint{2.304865in}{1.006463in}}%
\pgfpathlineto{\pgfqpoint{2.321661in}{1.018314in}}%
\pgfpathlineto{\pgfqpoint{2.338724in}{1.037893in}}%
\pgfpathlineto{\pgfqpoint{2.337893in}{1.033679in}}%
\pgfpathlineto{\pgfqpoint{2.346890in}{1.023505in}}%
\pgfpathlineto{\pgfqpoint{2.361560in}{1.020647in}}%
\pgfpathlineto{\pgfqpoint{2.370094in}{1.030341in}}%
\pgfpathlineto{\pgfqpoint{2.377730in}{1.025236in}}%
\pgfpathlineto{\pgfqpoint{2.391393in}{1.031228in}}%
\pgfpathlineto{\pgfqpoint{2.393347in}{1.037079in}}%
\pgfpathlineto{\pgfqpoint{2.401826in}{1.038395in}}%
\pgfpathlineto{\pgfqpoint{2.410633in}{1.045308in}}%
\pgfpathlineto{\pgfqpoint{2.418396in}{1.047589in}}%
\pgfpathlineto{\pgfqpoint{2.418643in}{1.059800in}}%
\pgfpathlineto{\pgfqpoint{2.423669in}{1.072991in}}%
\pgfpathlineto{\pgfqpoint{2.430300in}{1.081251in}}%
\pgfpathlineto{\pgfqpoint{2.432449in}{1.093654in}}%
\pgfpathlineto{\pgfqpoint{2.441091in}{1.111076in}}%
\pgfpathlineto{\pgfqpoint{2.441399in}{1.121442in}}%
\pgfpathlineto{\pgfqpoint{2.450326in}{1.114688in}}%
\pgfpathlineto{\pgfqpoint{2.459053in}{1.113049in}}%
\pgfpathlineto{\pgfqpoint{2.462449in}{1.117963in}}%
\pgfpathlineto{\pgfqpoint{2.468517in}{1.139377in}}%
\pgfpathlineto{\pgfqpoint{2.476965in}{1.138556in}}%
\pgfpathlineto{\pgfqpoint{2.479203in}{1.146402in}}%
\pgfpathlineto{\pgfqpoint{2.484159in}{1.149941in}}%
\pgfpathlineto{\pgfqpoint{2.494659in}{1.167732in}}%
\pgfpathlineto{\pgfqpoint{2.493134in}{1.172996in}}%
\pgfpathlineto{\pgfqpoint{2.497755in}{1.184435in}}%
\pgfpathlineto{\pgfqpoint{2.522427in}{1.170500in}}%
\pgfpathlineto{\pgfqpoint{2.525158in}{1.182877in}}%
\pgfpathlineto{\pgfqpoint{2.532953in}{1.183150in}}%
\pgfpathlineto{\pgfqpoint{2.539008in}{1.179234in}}%
\pgfpathlineto{\pgfqpoint{2.537796in}{1.172482in}}%
\pgfpathlineto{\pgfqpoint{2.551154in}{1.169411in}}%
\pgfpathlineto{\pgfqpoint{2.558697in}{1.164803in}}%
\pgfpathlineto{\pgfqpoint{2.564161in}{1.156689in}}%
\pgfpathlineto{\pgfqpoint{2.561823in}{1.147270in}}%
\pgfpathlineto{\pgfqpoint{2.557103in}{1.146498in}}%
\pgfpathlineto{\pgfqpoint{2.554365in}{1.132285in}}%
\pgfpathlineto{\pgfqpoint{2.560362in}{1.126728in}}%
\pgfpathlineto{\pgfqpoint{2.568804in}{1.131220in}}%
\pgfpathlineto{\pgfqpoint{2.576637in}{1.121734in}}%
\pgfpathlineto{\pgfqpoint{2.594137in}{1.120029in}}%
\pgfpathlineto{\pgfqpoint{2.597152in}{1.114573in}}%
\pgfpathlineto{\pgfqpoint{2.613348in}{1.109266in}}%
\pgfpathlineto{\pgfqpoint{2.611349in}{1.103004in}}%
\pgfpathlineto{\pgfqpoint{2.613456in}{1.085305in}}%
\pgfpathlineto{\pgfqpoint{2.619993in}{1.078622in}}%
\pgfpathlineto{\pgfqpoint{2.610347in}{1.075802in}}%
\pgfpathlineto{\pgfqpoint{2.611626in}{1.068246in}}%
\pgfpathlineto{\pgfqpoint{2.621916in}{1.061840in}}%
\pgfpathlineto{\pgfqpoint{2.622844in}{1.055521in}}%
\pgfpathlineto{\pgfqpoint{2.617034in}{1.050772in}}%
\pgfpathlineto{\pgfqpoint{2.605312in}{1.062022in}}%
\pgfpathlineto{\pgfqpoint{2.604153in}{1.053814in}}%
\pgfpathlineto{\pgfqpoint{2.613966in}{1.049911in}}%
\pgfpathlineto{\pgfqpoint{2.614948in}{1.045875in}}%
\pgfpathlineto{\pgfqpoint{2.628408in}{1.051209in}}%
\pgfpathlineto{\pgfqpoint{2.638722in}{1.052620in}}%
\pgfpathlineto{\pgfqpoint{2.649287in}{1.031303in}}%
\pgfpathlineto{\pgfqpoint{2.648108in}{1.031073in}}%
\pgfpathlineto{\pgfqpoint{2.643334in}{1.030086in}}%
\pgfpathlineto{\pgfqpoint{2.641927in}{1.029792in}}%
\pgfpathlineto{\pgfqpoint{2.640997in}{1.029610in}}%
\pgfpathlineto{\pgfqpoint{2.578091in}{1.016561in}}%
\pgfpathlineto{\pgfqpoint{2.521826in}{1.005045in}}%
\pgfpathlineto{\pgfqpoint{2.444033in}{0.991644in}}%
\pgfpathlineto{\pgfqpoint{2.377964in}{0.982918in}}%
\pgfpathlineto{\pgfqpoint{2.361621in}{0.981538in}}%
\pgfpathclose%
\pgfusepath{fill}%
\end{pgfscope}%
\begin{pgfscope}%
\pgfpathrectangle{\pgfqpoint{0.100000in}{0.100000in}}{\pgfqpoint{2.862985in}{1.832778in}}%
\pgfusepath{clip}%
\pgfsetbuttcap%
\pgfsetmiterjoin%
\definecolor{currentfill}{rgb}{0.524337,0.754325,0.286736}%
\pgfsetfillcolor{currentfill}%
\pgfsetlinewidth{0.000000pt}%
\definecolor{currentstroke}{rgb}{0.000000,0.000000,0.000000}%
\pgfsetstrokecolor{currentstroke}%
\pgfsetstrokeopacity{0.000000}%
\pgfsetdash{}{0pt}%
\pgfpathmoveto{\pgfqpoint{2.641870in}{1.120469in}}%
\pgfpathlineto{\pgfqpoint{2.654551in}{1.126221in}}%
\pgfpathlineto{\pgfqpoint{2.644424in}{1.096596in}}%
\pgfpathlineto{\pgfqpoint{2.637727in}{1.080908in}}%
\pgfpathlineto{\pgfqpoint{2.632902in}{1.089786in}}%
\pgfpathlineto{\pgfqpoint{2.641482in}{1.111041in}}%
\pgfpathlineto{\pgfqpoint{2.641870in}{1.120469in}}%
\pgfpathclose%
\pgfusepath{fill}%
\end{pgfscope}%
\begin{pgfscope}%
\pgfpathrectangle{\pgfqpoint{0.100000in}{0.100000in}}{\pgfqpoint{2.862985in}{1.832778in}}%
\pgfusepath{clip}%
\pgfsetbuttcap%
\pgfsetmiterjoin%
\definecolor{currentfill}{rgb}{0.524337,0.754325,0.286736}%
\pgfsetfillcolor{currentfill}%
\pgfsetlinewidth{0.000000pt}%
\definecolor{currentstroke}{rgb}{0.000000,0.000000,0.000000}%
\pgfsetstrokecolor{currentstroke}%
\pgfsetstrokeopacity{0.000000}%
\pgfsetdash{}{0pt}%
\pgfpathmoveto{\pgfqpoint{1.973847in}{0.935544in}}%
\pgfpathlineto{\pgfqpoint{1.970785in}{0.935101in}}%
\pgfpathlineto{\pgfqpoint{1.967896in}{0.934898in}}%
\pgfpathlineto{\pgfqpoint{1.969133in}{0.926031in}}%
\pgfpathlineto{\pgfqpoint{1.961682in}{0.917093in}}%
\pgfpathlineto{\pgfqpoint{1.960213in}{0.903107in}}%
\pgfpathlineto{\pgfqpoint{1.926904in}{0.900644in}}%
\pgfpathlineto{\pgfqpoint{1.929808in}{0.907211in}}%
\pgfpathlineto{\pgfqpoint{1.941818in}{0.919204in}}%
\pgfpathlineto{\pgfqpoint{1.942150in}{0.926155in}}%
\pgfpathlineto{\pgfqpoint{1.936833in}{0.932737in}}%
\pgfpathlineto{\pgfqpoint{1.887257in}{0.930116in}}%
\pgfpathlineto{\pgfqpoint{1.800580in}{0.927364in}}%
\pgfpathlineto{\pgfqpoint{1.749856in}{0.926440in}}%
\pgfpathlineto{\pgfqpoint{1.711515in}{0.926094in}}%
\pgfpathlineto{\pgfqpoint{1.711224in}{0.957954in}}%
\pgfpathlineto{\pgfqpoint{1.710919in}{1.050444in}}%
\pgfpathlineto{\pgfqpoint{1.711265in}{1.090665in}}%
\pgfpathlineto{\pgfqpoint{1.700905in}{1.096790in}}%
\pgfpathlineto{\pgfqpoint{1.695797in}{1.107461in}}%
\pgfpathlineto{\pgfqpoint{1.690237in}{1.112582in}}%
\pgfpathlineto{\pgfqpoint{1.689607in}{1.120348in}}%
\pgfpathlineto{\pgfqpoint{1.695414in}{1.132640in}}%
\pgfpathlineto{\pgfqpoint{1.695170in}{1.139746in}}%
\pgfpathlineto{\pgfqpoint{1.687350in}{1.138104in}}%
\pgfpathlineto{\pgfqpoint{1.677367in}{1.146251in}}%
\pgfpathlineto{\pgfqpoint{1.669364in}{1.160551in}}%
\pgfpathlineto{\pgfqpoint{1.662651in}{1.167151in}}%
\pgfpathlineto{\pgfqpoint{1.655637in}{1.183365in}}%
\pgfpathlineto{\pgfqpoint{1.728489in}{1.182229in}}%
\pgfpathlineto{\pgfqpoint{1.800906in}{1.184615in}}%
\pgfpathlineto{\pgfqpoint{1.847339in}{1.187293in}}%
\pgfpathlineto{\pgfqpoint{1.861865in}{1.172908in}}%
\pgfpathlineto{\pgfqpoint{1.859210in}{1.157906in}}%
\pgfpathlineto{\pgfqpoint{1.863083in}{1.141015in}}%
\pgfpathlineto{\pgfqpoint{1.867379in}{1.132498in}}%
\pgfpathlineto{\pgfqpoint{1.881087in}{1.120782in}}%
\pgfpathlineto{\pgfqpoint{1.884086in}{1.115607in}}%
\pgfpathlineto{\pgfqpoint{1.899094in}{1.103945in}}%
\pgfpathlineto{\pgfqpoint{1.901269in}{1.090863in}}%
\pgfpathlineto{\pgfqpoint{1.905579in}{1.081212in}}%
\pgfpathlineto{\pgfqpoint{1.915763in}{1.086834in}}%
\pgfpathlineto{\pgfqpoint{1.924157in}{1.084263in}}%
\pgfpathlineto{\pgfqpoint{1.930495in}{1.077440in}}%
\pgfpathlineto{\pgfqpoint{1.928102in}{1.065365in}}%
\pgfpathlineto{\pgfqpoint{1.920322in}{1.049578in}}%
\pgfpathlineto{\pgfqpoint{1.920442in}{1.041355in}}%
\pgfpathlineto{\pgfqpoint{1.926256in}{1.034614in}}%
\pgfpathlineto{\pgfqpoint{1.938768in}{1.025684in}}%
\pgfpathlineto{\pgfqpoint{1.953174in}{1.018436in}}%
\pgfpathlineto{\pgfqpoint{1.956591in}{1.012632in}}%
\pgfpathlineto{\pgfqpoint{1.963959in}{1.009888in}}%
\pgfpathlineto{\pgfqpoint{1.964252in}{1.002186in}}%
\pgfpathlineto{\pgfqpoint{1.969705in}{0.992402in}}%
\pgfpathlineto{\pgfqpoint{1.965922in}{0.983711in}}%
\pgfpathlineto{\pgfqpoint{1.973328in}{0.969882in}}%
\pgfpathlineto{\pgfqpoint{1.979121in}{0.972628in}}%
\pgfpathlineto{\pgfqpoint{1.986225in}{0.966856in}}%
\pgfpathlineto{\pgfqpoint{1.983837in}{0.941156in}}%
\pgfpathlineto{\pgfqpoint{1.975463in}{0.942786in}}%
\pgfpathlineto{\pgfqpoint{1.973847in}{0.935544in}}%
\pgfpathclose%
\pgfusepath{fill}%
\end{pgfscope}%
\begin{pgfscope}%
\pgfpathrectangle{\pgfqpoint{0.100000in}{0.100000in}}{\pgfqpoint{2.862985in}{1.832778in}}%
\pgfusepath{clip}%
\pgfsetbuttcap%
\pgfsetmiterjoin%
\definecolor{currentfill}{rgb}{0.781699,0.908497,0.622222}%
\pgfsetfillcolor{currentfill}%
\pgfsetlinewidth{0.000000pt}%
\definecolor{currentstroke}{rgb}{0.000000,0.000000,0.000000}%
\pgfsetstrokecolor{currentstroke}%
\pgfsetstrokeopacity{0.000000}%
\pgfsetdash{}{0pt}%
\pgfpathmoveto{\pgfqpoint{0.689569in}{0.941085in}}%
\pgfpathlineto{\pgfqpoint{0.695633in}{0.953995in}}%
\pgfpathlineto{\pgfqpoint{0.693973in}{0.973374in}}%
\pgfpathlineto{\pgfqpoint{0.696877in}{0.978886in}}%
\pgfpathlineto{\pgfqpoint{0.697164in}{1.003024in}}%
\pgfpathlineto{\pgfqpoint{0.699916in}{1.009976in}}%
\pgfpathlineto{\pgfqpoint{0.709613in}{1.011057in}}%
\pgfpathlineto{\pgfqpoint{0.718612in}{1.007971in}}%
\pgfpathlineto{\pgfqpoint{0.722531in}{0.999473in}}%
\pgfpathlineto{\pgfqpoint{0.728025in}{0.999798in}}%
\pgfpathlineto{\pgfqpoint{0.734852in}{1.009534in}}%
\pgfpathlineto{\pgfqpoint{0.744719in}{1.057523in}}%
\pgfpathlineto{\pgfqpoint{0.800868in}{1.045930in}}%
\pgfpathlineto{\pgfqpoint{0.833455in}{1.039504in}}%
\pgfpathlineto{\pgfqpoint{0.919991in}{1.024292in}}%
\pgfpathlineto{\pgfqpoint{0.943932in}{1.019484in}}%
\pgfpathlineto{\pgfqpoint{0.990915in}{1.012202in}}%
\pgfpathlineto{\pgfqpoint{0.981281in}{0.950172in}}%
\pgfpathlineto{\pgfqpoint{0.971262in}{0.885469in}}%
\pgfpathlineto{\pgfqpoint{0.959723in}{0.812676in}}%
\pgfpathlineto{\pgfqpoint{0.946747in}{0.729118in}}%
\pgfpathlineto{\pgfqpoint{0.936264in}{0.660490in}}%
\pgfpathlineto{\pgfqpoint{0.861176in}{0.672414in}}%
\pgfpathlineto{\pgfqpoint{0.828185in}{0.678064in}}%
\pgfpathlineto{\pgfqpoint{0.813447in}{0.686887in}}%
\pgfpathlineto{\pgfqpoint{0.717356in}{0.744864in}}%
\pgfpathlineto{\pgfqpoint{0.645244in}{0.788862in}}%
\pgfpathlineto{\pgfqpoint{0.647646in}{0.796660in}}%
\pgfpathlineto{\pgfqpoint{0.653603in}{0.802117in}}%
\pgfpathlineto{\pgfqpoint{0.659798in}{0.801105in}}%
\pgfpathlineto{\pgfqpoint{0.668785in}{0.806836in}}%
\pgfpathlineto{\pgfqpoint{0.670200in}{0.815071in}}%
\pgfpathlineto{\pgfqpoint{0.659240in}{0.825061in}}%
\pgfpathlineto{\pgfqpoint{0.667035in}{0.844239in}}%
\pgfpathlineto{\pgfqpoint{0.674871in}{0.851614in}}%
\pgfpathlineto{\pgfqpoint{0.678591in}{0.860364in}}%
\pgfpathlineto{\pgfqpoint{0.680882in}{0.876407in}}%
\pgfpathlineto{\pgfqpoint{0.688230in}{0.883673in}}%
\pgfpathlineto{\pgfqpoint{0.703654in}{0.890930in}}%
\pgfpathlineto{\pgfqpoint{0.703313in}{0.897324in}}%
\pgfpathlineto{\pgfqpoint{0.694699in}{0.905259in}}%
\pgfpathlineto{\pgfqpoint{0.693526in}{0.921617in}}%
\pgfpathlineto{\pgfqpoint{0.687595in}{0.933576in}}%
\pgfpathlineto{\pgfqpoint{0.689569in}{0.941085in}}%
\pgfpathclose%
\pgfusepath{fill}%
\end{pgfscope}%
\begin{pgfscope}%
\pgfpathrectangle{\pgfqpoint{0.100000in}{0.100000in}}{\pgfqpoint{2.862985in}{1.832778in}}%
\pgfusepath{clip}%
\pgfsetbuttcap%
\pgfsetmiterjoin%
\definecolor{currentfill}{rgb}{0.594464,0.799846,0.371626}%
\pgfsetfillcolor{currentfill}%
\pgfsetlinewidth{0.000000pt}%
\definecolor{currentstroke}{rgb}{0.000000,0.000000,0.000000}%
\pgfsetstrokecolor{currentstroke}%
\pgfsetstrokeopacity{0.000000}%
\pgfsetdash{}{0pt}%
\pgfpathmoveto{\pgfqpoint{1.719643in}{0.746495in}}%
\pgfpathlineto{\pgfqpoint{1.705724in}{0.750800in}}%
\pgfpathlineto{\pgfqpoint{1.687888in}{0.764459in}}%
\pgfpathlineto{\pgfqpoint{1.679966in}{0.767394in}}%
\pgfpathlineto{\pgfqpoint{1.674932in}{0.761495in}}%
\pgfpathlineto{\pgfqpoint{1.668575in}{0.761197in}}%
\pgfpathlineto{\pgfqpoint{1.660532in}{0.766205in}}%
\pgfpathlineto{\pgfqpoint{1.647911in}{0.759788in}}%
\pgfpathlineto{\pgfqpoint{1.643418in}{0.762948in}}%
\pgfpathlineto{\pgfqpoint{1.630999in}{0.755462in}}%
\pgfpathlineto{\pgfqpoint{1.617897in}{0.756843in}}%
\pgfpathlineto{\pgfqpoint{1.607791in}{0.761708in}}%
\pgfpathlineto{\pgfqpoint{1.600714in}{0.759877in}}%
\pgfpathlineto{\pgfqpoint{1.589386in}{0.766758in}}%
\pgfpathlineto{\pgfqpoint{1.583084in}{0.754622in}}%
\pgfpathlineto{\pgfqpoint{1.576640in}{0.765059in}}%
\pgfpathlineto{\pgfqpoint{1.563911in}{0.761010in}}%
\pgfpathlineto{\pgfqpoint{1.552778in}{0.770922in}}%
\pgfpathlineto{\pgfqpoint{1.543049in}{0.762933in}}%
\pgfpathlineto{\pgfqpoint{1.538177in}{0.770275in}}%
\pgfpathlineto{\pgfqpoint{1.531154in}{0.772659in}}%
\pgfpathlineto{\pgfqpoint{1.526877in}{0.779738in}}%
\pgfpathlineto{\pgfqpoint{1.517692in}{0.781753in}}%
\pgfpathlineto{\pgfqpoint{1.512392in}{0.776425in}}%
\pgfpathlineto{\pgfqpoint{1.503387in}{0.783330in}}%
\pgfpathlineto{\pgfqpoint{1.499189in}{0.781753in}}%
\pgfpathlineto{\pgfqpoint{1.484265in}{0.787349in}}%
\pgfpathlineto{\pgfqpoint{1.474917in}{0.787973in}}%
\pgfpathlineto{\pgfqpoint{1.470733in}{0.799856in}}%
\pgfpathlineto{\pgfqpoint{1.463240in}{0.798376in}}%
\pgfpathlineto{\pgfqpoint{1.454657in}{0.801308in}}%
\pgfpathlineto{\pgfqpoint{1.449005in}{0.799614in}}%
\pgfpathlineto{\pgfqpoint{1.436884in}{0.812959in}}%
\pgfpathlineto{\pgfqpoint{1.433507in}{0.812086in}}%
\pgfpathlineto{\pgfqpoint{1.436574in}{0.866207in}}%
\pgfpathlineto{\pgfqpoint{1.439916in}{0.933195in}}%
\pgfpathlineto{\pgfqpoint{1.385064in}{0.936263in}}%
\pgfpathlineto{\pgfqpoint{1.330936in}{0.940356in}}%
\pgfpathlineto{\pgfqpoint{1.289118in}{0.943987in}}%
\pgfpathlineto{\pgfqpoint{1.291895in}{0.975633in}}%
\pgfpathlineto{\pgfqpoint{1.339821in}{0.971059in}}%
\pgfpathlineto{\pgfqpoint{1.394179in}{0.967651in}}%
\pgfpathlineto{\pgfqpoint{1.469882in}{0.963187in}}%
\pgfpathlineto{\pgfqpoint{1.514818in}{0.961453in}}%
\pgfpathlineto{\pgfqpoint{1.568651in}{0.959554in}}%
\pgfpathlineto{\pgfqpoint{1.641779in}{0.958274in}}%
\pgfpathlineto{\pgfqpoint{1.711224in}{0.957954in}}%
\pgfpathlineto{\pgfqpoint{1.711515in}{0.926094in}}%
\pgfpathlineto{\pgfqpoint{1.715413in}{0.902092in}}%
\pgfpathlineto{\pgfqpoint{1.721468in}{0.857760in}}%
\pgfpathlineto{\pgfqpoint{1.720219in}{0.782052in}}%
\pgfpathlineto{\pgfqpoint{1.719643in}{0.746495in}}%
\pgfpathclose%
\pgfusepath{fill}%
\end{pgfscope}%
\begin{pgfscope}%
\pgfpathrectangle{\pgfqpoint{0.100000in}{0.100000in}}{\pgfqpoint{2.862985in}{1.832778in}}%
\pgfusepath{clip}%
\pgfsetbuttcap%
\pgfsetmiterjoin%
\definecolor{currentfill}{rgb}{0.647059,0.833987,0.435294}%
\pgfsetfillcolor{currentfill}%
\pgfsetlinewidth{0.000000pt}%
\definecolor{currentstroke}{rgb}{0.000000,0.000000,0.000000}%
\pgfsetstrokecolor{currentstroke}%
\pgfsetstrokeopacity{0.000000}%
\pgfsetdash{}{0pt}%
\pgfpathmoveto{\pgfqpoint{2.240820in}{0.864070in}}%
\pgfpathlineto{\pgfqpoint{2.240861in}{0.878097in}}%
\pgfpathlineto{\pgfqpoint{2.253051in}{0.883484in}}%
\pgfpathlineto{\pgfqpoint{2.253562in}{0.892095in}}%
\pgfpathlineto{\pgfqpoint{2.264476in}{0.902614in}}%
\pgfpathlineto{\pgfqpoint{2.278119in}{0.904768in}}%
\pgfpathlineto{\pgfqpoint{2.290113in}{0.916361in}}%
\pgfpathlineto{\pgfqpoint{2.302639in}{0.921543in}}%
\pgfpathlineto{\pgfqpoint{2.312040in}{0.937098in}}%
\pgfpathlineto{\pgfqpoint{2.320602in}{0.935969in}}%
\pgfpathlineto{\pgfqpoint{2.338697in}{0.950142in}}%
\pgfpathlineto{\pgfqpoint{2.348253in}{0.950409in}}%
\pgfpathlineto{\pgfqpoint{2.352288in}{0.961208in}}%
\pgfpathlineto{\pgfqpoint{2.359883in}{0.968727in}}%
\pgfpathlineto{\pgfqpoint{2.361621in}{0.981538in}}%
\pgfpathlineto{\pgfqpoint{2.377964in}{0.982918in}}%
\pgfpathlineto{\pgfqpoint{2.444033in}{0.991644in}}%
\pgfpathlineto{\pgfqpoint{2.521826in}{1.005045in}}%
\pgfpathlineto{\pgfqpoint{2.578091in}{1.016561in}}%
\pgfpathlineto{\pgfqpoint{2.640997in}{1.029610in}}%
\pgfpathlineto{\pgfqpoint{2.658973in}{1.004931in}}%
\pgfpathlineto{\pgfqpoint{2.649236in}{1.006506in}}%
\pgfpathlineto{\pgfqpoint{2.636737in}{1.003461in}}%
\pgfpathlineto{\pgfqpoint{2.624421in}{0.990884in}}%
\pgfpathlineto{\pgfqpoint{2.615570in}{0.991789in}}%
\pgfpathlineto{\pgfqpoint{2.614450in}{0.984319in}}%
\pgfpathlineto{\pgfqpoint{2.630449in}{0.990226in}}%
\pgfpathlineto{\pgfqpoint{2.646548in}{0.992662in}}%
\pgfpathlineto{\pgfqpoint{2.652526in}{0.989421in}}%
\pgfpathlineto{\pgfqpoint{2.660571in}{0.993116in}}%
\pgfpathlineto{\pgfqpoint{2.664722in}{0.990527in}}%
\pgfpathlineto{\pgfqpoint{2.668396in}{0.978199in}}%
\pgfpathlineto{\pgfqpoint{2.660750in}{0.974382in}}%
\pgfpathlineto{\pgfqpoint{2.655502in}{0.959350in}}%
\pgfpathlineto{\pgfqpoint{2.650000in}{0.953497in}}%
\pgfpathlineto{\pgfqpoint{2.633132in}{0.954777in}}%
\pgfpathlineto{\pgfqpoint{2.634039in}{0.963604in}}%
\pgfpathlineto{\pgfqpoint{2.624827in}{0.959748in}}%
\pgfpathlineto{\pgfqpoint{2.622823in}{0.952387in}}%
\pgfpathlineto{\pgfqpoint{2.629137in}{0.944761in}}%
\pgfpathlineto{\pgfqpoint{2.631282in}{0.931819in}}%
\pgfpathlineto{\pgfqpoint{2.620950in}{0.924444in}}%
\pgfpathlineto{\pgfqpoint{2.632125in}{0.921851in}}%
\pgfpathlineto{\pgfqpoint{2.637181in}{0.927273in}}%
\pgfpathlineto{\pgfqpoint{2.648362in}{0.927934in}}%
\pgfpathlineto{\pgfqpoint{2.642616in}{0.916224in}}%
\pgfpathlineto{\pgfqpoint{2.635579in}{0.910582in}}%
\pgfpathlineto{\pgfqpoint{2.614292in}{0.904933in}}%
\pgfpathlineto{\pgfqpoint{2.592436in}{0.885111in}}%
\pgfpathlineto{\pgfqpoint{2.578960in}{0.865580in}}%
\pgfpathlineto{\pgfqpoint{2.578880in}{0.857648in}}%
\pgfpathlineto{\pgfqpoint{2.573497in}{0.846702in}}%
\pgfpathlineto{\pgfqpoint{2.545869in}{0.839501in}}%
\pgfpathlineto{\pgfqpoint{2.480165in}{0.886600in}}%
\pgfpathlineto{\pgfqpoint{2.422304in}{0.878098in}}%
\pgfpathlineto{\pgfqpoint{2.421819in}{0.885949in}}%
\pgfpathlineto{\pgfqpoint{2.413024in}{0.894789in}}%
\pgfpathlineto{\pgfqpoint{2.406366in}{0.896952in}}%
\pgfpathlineto{\pgfqpoint{2.343499in}{0.890420in}}%
\pgfpathlineto{\pgfqpoint{2.329049in}{0.885519in}}%
\pgfpathlineto{\pgfqpoint{2.302958in}{0.872540in}}%
\pgfpathlineto{\pgfqpoint{2.280410in}{0.868948in}}%
\pgfpathlineto{\pgfqpoint{2.240820in}{0.864070in}}%
\pgfpathclose%
\pgfusepath{fill}%
\end{pgfscope}%
\begin{pgfscope}%
\pgfpathrectangle{\pgfqpoint{0.100000in}{0.100000in}}{\pgfqpoint{2.862985in}{1.832778in}}%
\pgfusepath{clip}%
\pgfsetbuttcap%
\pgfsetmiterjoin%
\definecolor{currentfill}{rgb}{0.647059,0.833987,0.435294}%
\pgfsetfillcolor{currentfill}%
\pgfsetlinewidth{0.000000pt}%
\definecolor{currentstroke}{rgb}{0.000000,0.000000,0.000000}%
\pgfsetstrokecolor{currentstroke}%
\pgfsetstrokeopacity{0.000000}%
\pgfsetdash{}{0pt}%
\pgfpathmoveto{\pgfqpoint{2.240820in}{0.864070in}}%
\pgfpathlineto{\pgfqpoint{2.175022in}{0.856843in}}%
\pgfpathlineto{\pgfqpoint{2.114830in}{0.851430in}}%
\pgfpathlineto{\pgfqpoint{2.052965in}{0.847429in}}%
\pgfpathlineto{\pgfqpoint{2.042346in}{0.845890in}}%
\pgfpathlineto{\pgfqpoint{2.000643in}{0.842683in}}%
\pgfpathlineto{\pgfqpoint{1.933848in}{0.838785in}}%
\pgfpathlineto{\pgfqpoint{1.945732in}{0.848649in}}%
\pgfpathlineto{\pgfqpoint{1.943020in}{0.859033in}}%
\pgfpathlineto{\pgfqpoint{1.945647in}{0.871619in}}%
\pgfpathlineto{\pgfqpoint{1.949660in}{0.877545in}}%
\pgfpathlineto{\pgfqpoint{1.949509in}{0.885789in}}%
\pgfpathlineto{\pgfqpoint{1.960195in}{0.890974in}}%
\pgfpathlineto{\pgfqpoint{1.960213in}{0.903107in}}%
\pgfpathlineto{\pgfqpoint{1.961682in}{0.917093in}}%
\pgfpathlineto{\pgfqpoint{1.969133in}{0.926031in}}%
\pgfpathlineto{\pgfqpoint{1.967896in}{0.934898in}}%
\pgfpathlineto{\pgfqpoint{1.970785in}{0.935101in}}%
\pgfpathlineto{\pgfqpoint{1.973847in}{0.935544in}}%
\pgfpathlineto{\pgfqpoint{2.042942in}{0.939936in}}%
\pgfpathlineto{\pgfqpoint{2.040629in}{0.951286in}}%
\pgfpathlineto{\pgfqpoint{2.051748in}{0.949746in}}%
\pgfpathlineto{\pgfqpoint{2.119049in}{0.956404in}}%
\pgfpathlineto{\pgfqpoint{2.173070in}{0.959461in}}%
\pgfpathlineto{\pgfqpoint{2.205991in}{0.962230in}}%
\pgfpathlineto{\pgfqpoint{2.260505in}{0.967399in}}%
\pgfpathlineto{\pgfqpoint{2.261414in}{0.968427in}}%
\pgfpathlineto{\pgfqpoint{2.284707in}{0.970734in}}%
\pgfpathlineto{\pgfqpoint{2.361621in}{0.981538in}}%
\pgfpathlineto{\pgfqpoint{2.359883in}{0.968727in}}%
\pgfpathlineto{\pgfqpoint{2.352288in}{0.961208in}}%
\pgfpathlineto{\pgfqpoint{2.348253in}{0.950409in}}%
\pgfpathlineto{\pgfqpoint{2.338697in}{0.950142in}}%
\pgfpathlineto{\pgfqpoint{2.320602in}{0.935969in}}%
\pgfpathlineto{\pgfqpoint{2.312040in}{0.937098in}}%
\pgfpathlineto{\pgfqpoint{2.302639in}{0.921543in}}%
\pgfpathlineto{\pgfqpoint{2.290113in}{0.916361in}}%
\pgfpathlineto{\pgfqpoint{2.278119in}{0.904768in}}%
\pgfpathlineto{\pgfqpoint{2.264476in}{0.902614in}}%
\pgfpathlineto{\pgfqpoint{2.253562in}{0.892095in}}%
\pgfpathlineto{\pgfqpoint{2.253051in}{0.883484in}}%
\pgfpathlineto{\pgfqpoint{2.240861in}{0.878097in}}%
\pgfpathlineto{\pgfqpoint{2.240820in}{0.864070in}}%
\pgfpathclose%
\pgfusepath{fill}%
\end{pgfscope}%
\begin{pgfscope}%
\pgfpathrectangle{\pgfqpoint{0.100000in}{0.100000in}}{\pgfqpoint{2.862985in}{1.832778in}}%
\pgfusepath{clip}%
\pgfsetbuttcap%
\pgfsetmiterjoin%
\definecolor{currentfill}{rgb}{0.673356,0.851057,0.467128}%
\pgfsetfillcolor{currentfill}%
\pgfsetlinewidth{0.000000pt}%
\definecolor{currentstroke}{rgb}{0.000000,0.000000,0.000000}%
\pgfsetstrokecolor{currentstroke}%
\pgfsetstrokeopacity{0.000000}%
\pgfsetdash{}{0pt}%
\pgfpathmoveto{\pgfqpoint{1.073883in}{0.669271in}}%
\pgfpathlineto{\pgfqpoint{1.069190in}{0.676811in}}%
\pgfpathlineto{\pgfqpoint{1.071132in}{0.683307in}}%
\pgfpathlineto{\pgfqpoint{1.104118in}{0.679215in}}%
\pgfpathlineto{\pgfqpoint{1.165452in}{0.672249in}}%
\pgfpathlineto{\pgfqpoint{1.209801in}{0.667891in}}%
\pgfpathlineto{\pgfqpoint{1.261021in}{0.662722in}}%
\pgfpathlineto{\pgfqpoint{1.263807in}{0.695029in}}%
\pgfpathlineto{\pgfqpoint{1.272520in}{0.777201in}}%
\pgfpathlineto{\pgfqpoint{1.278172in}{0.834767in}}%
\pgfpathlineto{\pgfqpoint{1.283057in}{0.889756in}}%
\pgfpathlineto{\pgfqpoint{1.287612in}{0.944071in}}%
\pgfpathlineto{\pgfqpoint{1.289118in}{0.943987in}}%
\pgfpathlineto{\pgfqpoint{1.330936in}{0.940356in}}%
\pgfpathlineto{\pgfqpoint{1.385064in}{0.936263in}}%
\pgfpathlineto{\pgfqpoint{1.439916in}{0.933195in}}%
\pgfpathlineto{\pgfqpoint{1.436574in}{0.866207in}}%
\pgfpathlineto{\pgfqpoint{1.433507in}{0.812086in}}%
\pgfpathlineto{\pgfqpoint{1.436884in}{0.812959in}}%
\pgfpathlineto{\pgfqpoint{1.449005in}{0.799614in}}%
\pgfpathlineto{\pgfqpoint{1.454657in}{0.801308in}}%
\pgfpathlineto{\pgfqpoint{1.463240in}{0.798376in}}%
\pgfpathlineto{\pgfqpoint{1.470733in}{0.799856in}}%
\pgfpathlineto{\pgfqpoint{1.474917in}{0.787973in}}%
\pgfpathlineto{\pgfqpoint{1.484265in}{0.787349in}}%
\pgfpathlineto{\pgfqpoint{1.499189in}{0.781753in}}%
\pgfpathlineto{\pgfqpoint{1.503387in}{0.783330in}}%
\pgfpathlineto{\pgfqpoint{1.512392in}{0.776425in}}%
\pgfpathlineto{\pgfqpoint{1.517692in}{0.781753in}}%
\pgfpathlineto{\pgfqpoint{1.526877in}{0.779738in}}%
\pgfpathlineto{\pgfqpoint{1.531154in}{0.772659in}}%
\pgfpathlineto{\pgfqpoint{1.538177in}{0.770275in}}%
\pgfpathlineto{\pgfqpoint{1.543049in}{0.762933in}}%
\pgfpathlineto{\pgfqpoint{1.552778in}{0.770922in}}%
\pgfpathlineto{\pgfqpoint{1.563911in}{0.761010in}}%
\pgfpathlineto{\pgfqpoint{1.576640in}{0.765059in}}%
\pgfpathlineto{\pgfqpoint{1.583084in}{0.754622in}}%
\pgfpathlineto{\pgfqpoint{1.589386in}{0.766758in}}%
\pgfpathlineto{\pgfqpoint{1.600714in}{0.759877in}}%
\pgfpathlineto{\pgfqpoint{1.607791in}{0.761708in}}%
\pgfpathlineto{\pgfqpoint{1.617897in}{0.756843in}}%
\pgfpathlineto{\pgfqpoint{1.630999in}{0.755462in}}%
\pgfpathlineto{\pgfqpoint{1.643418in}{0.762948in}}%
\pgfpathlineto{\pgfqpoint{1.647911in}{0.759788in}}%
\pgfpathlineto{\pgfqpoint{1.660532in}{0.766205in}}%
\pgfpathlineto{\pgfqpoint{1.668575in}{0.761197in}}%
\pgfpathlineto{\pgfqpoint{1.674932in}{0.761495in}}%
\pgfpathlineto{\pgfqpoint{1.679966in}{0.767394in}}%
\pgfpathlineto{\pgfqpoint{1.687888in}{0.764459in}}%
\pgfpathlineto{\pgfqpoint{1.705724in}{0.750800in}}%
\pgfpathlineto{\pgfqpoint{1.719643in}{0.746495in}}%
\pgfpathlineto{\pgfqpoint{1.725225in}{0.741221in}}%
\pgfpathlineto{\pgfqpoint{1.732212in}{0.744093in}}%
\pgfpathlineto{\pgfqpoint{1.742799in}{0.741894in}}%
\pgfpathlineto{\pgfqpoint{1.743006in}{0.708326in}}%
\pgfpathlineto{\pgfqpoint{1.743890in}{0.643403in}}%
\pgfpathlineto{\pgfqpoint{1.751243in}{0.637164in}}%
\pgfpathlineto{\pgfqpoint{1.757137in}{0.625668in}}%
\pgfpathlineto{\pgfqpoint{1.755065in}{0.617977in}}%
\pgfpathlineto{\pgfqpoint{1.759693in}{0.614721in}}%
\pgfpathlineto{\pgfqpoint{1.763454in}{0.600538in}}%
\pgfpathlineto{\pgfqpoint{1.770658in}{0.592942in}}%
\pgfpathlineto{\pgfqpoint{1.772272in}{0.576816in}}%
\pgfpathlineto{\pgfqpoint{1.771019in}{0.569988in}}%
\pgfpathlineto{\pgfqpoint{1.761227in}{0.551919in}}%
\pgfpathlineto{\pgfqpoint{1.760101in}{0.538927in}}%
\pgfpathlineto{\pgfqpoint{1.763423in}{0.536215in}}%
\pgfpathlineto{\pgfqpoint{1.763712in}{0.524834in}}%
\pgfpathlineto{\pgfqpoint{1.760335in}{0.517680in}}%
\pgfpathlineto{\pgfqpoint{1.755085in}{0.514985in}}%
\pgfpathlineto{\pgfqpoint{1.749944in}{0.505637in}}%
\pgfpathlineto{\pgfqpoint{1.756496in}{0.496583in}}%
\pgfpathlineto{\pgfqpoint{1.743777in}{0.496405in}}%
\pgfpathlineto{\pgfqpoint{1.709766in}{0.480844in}}%
\pgfpathlineto{\pgfqpoint{1.716259in}{0.490157in}}%
\pgfpathlineto{\pgfqpoint{1.708396in}{0.495178in}}%
\pgfpathlineto{\pgfqpoint{1.706753in}{0.503714in}}%
\pgfpathlineto{\pgfqpoint{1.691406in}{0.488845in}}%
\pgfpathlineto{\pgfqpoint{1.698217in}{0.478682in}}%
\pgfpathlineto{\pgfqpoint{1.688501in}{0.465742in}}%
\pgfpathlineto{\pgfqpoint{1.683273in}{0.466013in}}%
\pgfpathlineto{\pgfqpoint{1.678354in}{0.451920in}}%
\pgfpathlineto{\pgfqpoint{1.662798in}{0.440835in}}%
\pgfpathlineto{\pgfqpoint{1.648265in}{0.436840in}}%
\pgfpathlineto{\pgfqpoint{1.622917in}{0.426447in}}%
\pgfpathlineto{\pgfqpoint{1.620291in}{0.432244in}}%
\pgfpathlineto{\pgfqpoint{1.604501in}{0.426908in}}%
\pgfpathlineto{\pgfqpoint{1.614213in}{0.417822in}}%
\pgfpathlineto{\pgfqpoint{1.598894in}{0.409929in}}%
\pgfpathlineto{\pgfqpoint{1.582371in}{0.398011in}}%
\pgfpathlineto{\pgfqpoint{1.578403in}{0.403551in}}%
\pgfpathlineto{\pgfqpoint{1.564881in}{0.395248in}}%
\pgfpathlineto{\pgfqpoint{1.578135in}{0.392095in}}%
\pgfpathlineto{\pgfqpoint{1.568142in}{0.379039in}}%
\pgfpathlineto{\pgfqpoint{1.563264in}{0.382924in}}%
\pgfpathlineto{\pgfqpoint{1.556713in}{0.376684in}}%
\pgfpathlineto{\pgfqpoint{1.564878in}{0.371233in}}%
\pgfpathlineto{\pgfqpoint{1.560050in}{0.363265in}}%
\pgfpathlineto{\pgfqpoint{1.554117in}{0.344404in}}%
\pgfpathlineto{\pgfqpoint{1.545559in}{0.326238in}}%
\pgfpathlineto{\pgfqpoint{1.549401in}{0.314331in}}%
\pgfpathlineto{\pgfqpoint{1.555932in}{0.286258in}}%
\pgfpathlineto{\pgfqpoint{1.556520in}{0.274896in}}%
\pgfpathlineto{\pgfqpoint{1.566610in}{0.259994in}}%
\pgfpathlineto{\pgfqpoint{1.558829in}{0.260865in}}%
\pgfpathlineto{\pgfqpoint{1.551272in}{0.253347in}}%
\pgfpathlineto{\pgfqpoint{1.542420in}{0.259256in}}%
\pgfpathlineto{\pgfqpoint{1.539243in}{0.265150in}}%
\pgfpathlineto{\pgfqpoint{1.526645in}{0.267892in}}%
\pgfpathlineto{\pgfqpoint{1.507402in}{0.268229in}}%
\pgfpathlineto{\pgfqpoint{1.493224in}{0.279391in}}%
\pgfpathlineto{\pgfqpoint{1.480350in}{0.281256in}}%
\pgfpathlineto{\pgfqpoint{1.472540in}{0.290127in}}%
\pgfpathlineto{\pgfqpoint{1.456183in}{0.293697in}}%
\pgfpathlineto{\pgfqpoint{1.447239in}{0.322221in}}%
\pgfpathlineto{\pgfqpoint{1.438094in}{0.333648in}}%
\pgfpathlineto{\pgfqpoint{1.439648in}{0.344506in}}%
\pgfpathlineto{\pgfqpoint{1.433978in}{0.352440in}}%
\pgfpathlineto{\pgfqpoint{1.437536in}{0.363286in}}%
\pgfpathlineto{\pgfqpoint{1.434601in}{0.371234in}}%
\pgfpathlineto{\pgfqpoint{1.425411in}{0.374834in}}%
\pgfpathlineto{\pgfqpoint{1.416818in}{0.383999in}}%
\pgfpathlineto{\pgfqpoint{1.413687in}{0.396269in}}%
\pgfpathlineto{\pgfqpoint{1.405556in}{0.407420in}}%
\pgfpathlineto{\pgfqpoint{1.394734in}{0.416096in}}%
\pgfpathlineto{\pgfqpoint{1.384986in}{0.440984in}}%
\pgfpathlineto{\pgfqpoint{1.377115in}{0.468203in}}%
\pgfpathlineto{\pgfqpoint{1.370657in}{0.478966in}}%
\pgfpathlineto{\pgfqpoint{1.359460in}{0.488026in}}%
\pgfpathlineto{\pgfqpoint{1.356662in}{0.494603in}}%
\pgfpathlineto{\pgfqpoint{1.346179in}{0.498686in}}%
\pgfpathlineto{\pgfqpoint{1.335831in}{0.516112in}}%
\pgfpathlineto{\pgfqpoint{1.326382in}{0.514774in}}%
\pgfpathlineto{\pgfqpoint{1.316787in}{0.519126in}}%
\pgfpathlineto{\pgfqpoint{1.303186in}{0.518285in}}%
\pgfpathlineto{\pgfqpoint{1.289369in}{0.525467in}}%
\pgfpathlineto{\pgfqpoint{1.285472in}{0.518637in}}%
\pgfpathlineto{\pgfqpoint{1.269325in}{0.518466in}}%
\pgfpathlineto{\pgfqpoint{1.261104in}{0.505517in}}%
\pgfpathlineto{\pgfqpoint{1.253954in}{0.489484in}}%
\pgfpathlineto{\pgfqpoint{1.238750in}{0.472274in}}%
\pgfpathlineto{\pgfqpoint{1.221524in}{0.479827in}}%
\pgfpathlineto{\pgfqpoint{1.219080in}{0.484818in}}%
\pgfpathlineto{\pgfqpoint{1.208611in}{0.488625in}}%
\pgfpathlineto{\pgfqpoint{1.205396in}{0.493833in}}%
\pgfpathlineto{\pgfqpoint{1.191499in}{0.499090in}}%
\pgfpathlineto{\pgfqpoint{1.183725in}{0.509836in}}%
\pgfpathlineto{\pgfqpoint{1.174646in}{0.515021in}}%
\pgfpathlineto{\pgfqpoint{1.166834in}{0.524068in}}%
\pgfpathlineto{\pgfqpoint{1.160749in}{0.539391in}}%
\pgfpathlineto{\pgfqpoint{1.161430in}{0.560329in}}%
\pgfpathlineto{\pgfqpoint{1.154291in}{0.570909in}}%
\pgfpathlineto{\pgfqpoint{1.153469in}{0.582369in}}%
\pgfpathlineto{\pgfqpoint{1.137630in}{0.599527in}}%
\pgfpathlineto{\pgfqpoint{1.128414in}{0.603189in}}%
\pgfpathlineto{\pgfqpoint{1.118634in}{0.619199in}}%
\pgfpathlineto{\pgfqpoint{1.110300in}{0.625560in}}%
\pgfpathlineto{\pgfqpoint{1.099711in}{0.641049in}}%
\pgfpathlineto{\pgfqpoint{1.088876in}{0.647753in}}%
\pgfpathlineto{\pgfqpoint{1.081780in}{0.664929in}}%
\pgfpathlineto{\pgfqpoint{1.073883in}{0.669271in}}%
\pgfpathclose%
\pgfusepath{fill}%
\end{pgfscope}%
\begin{pgfscope}%
\pgfpathrectangle{\pgfqpoint{0.100000in}{0.100000in}}{\pgfqpoint{2.862985in}{1.832778in}}%
\pgfusepath{clip}%
\pgfsetbuttcap%
\pgfsetmiterjoin%
\definecolor{currentfill}{rgb}{0.717186,0.879508,0.520185}%
\pgfsetfillcolor{currentfill}%
\pgfsetlinewidth{0.000000pt}%
\definecolor{currentstroke}{rgb}{0.000000,0.000000,0.000000}%
\pgfsetstrokecolor{currentstroke}%
\pgfsetstrokeopacity{0.000000}%
\pgfsetdash{}{0pt}%
\pgfpathmoveto{\pgfqpoint{0.990915in}{1.012202in}}%
\pgfpathlineto{\pgfqpoint{1.069080in}{1.000850in}}%
\pgfpathlineto{\pgfqpoint{1.099362in}{0.996184in}}%
\pgfpathlineto{\pgfqpoint{1.184722in}{0.986032in}}%
\pgfpathlineto{\pgfqpoint{1.242175in}{0.980039in}}%
\pgfpathlineto{\pgfqpoint{1.291895in}{0.975633in}}%
\pgfpathlineto{\pgfqpoint{1.289118in}{0.943987in}}%
\pgfpathlineto{\pgfqpoint{1.287612in}{0.944071in}}%
\pgfpathlineto{\pgfqpoint{1.283057in}{0.889756in}}%
\pgfpathlineto{\pgfqpoint{1.278172in}{0.834767in}}%
\pgfpathlineto{\pgfqpoint{1.272520in}{0.777201in}}%
\pgfpathlineto{\pgfqpoint{1.263807in}{0.695029in}}%
\pgfpathlineto{\pgfqpoint{1.261021in}{0.662722in}}%
\pgfpathlineto{\pgfqpoint{1.209801in}{0.667891in}}%
\pgfpathlineto{\pgfqpoint{1.165452in}{0.672249in}}%
\pgfpathlineto{\pgfqpoint{1.104118in}{0.679215in}}%
\pgfpathlineto{\pgfqpoint{1.071132in}{0.683307in}}%
\pgfpathlineto{\pgfqpoint{1.069190in}{0.676811in}}%
\pgfpathlineto{\pgfqpoint{1.073883in}{0.669271in}}%
\pgfpathlineto{\pgfqpoint{1.034240in}{0.674419in}}%
\pgfpathlineto{\pgfqpoint{0.985330in}{0.681439in}}%
\pgfpathlineto{\pgfqpoint{0.980878in}{0.653775in}}%
\pgfpathlineto{\pgfqpoint{0.936264in}{0.660490in}}%
\pgfpathlineto{\pgfqpoint{0.946747in}{0.729118in}}%
\pgfpathlineto{\pgfqpoint{0.959723in}{0.812676in}}%
\pgfpathlineto{\pgfqpoint{0.971262in}{0.885469in}}%
\pgfpathlineto{\pgfqpoint{0.981281in}{0.950172in}}%
\pgfpathlineto{\pgfqpoint{0.990915in}{1.012202in}}%
\pgfpathclose%
\pgfusepath{fill}%
\end{pgfscope}%
\begin{pgfscope}%
\pgfpathrectangle{\pgfqpoint{0.100000in}{0.100000in}}{\pgfqpoint{2.862985in}{1.832778in}}%
\pgfusepath{clip}%
\pgfsetbuttcap%
\pgfsetmiterjoin%
\definecolor{currentfill}{rgb}{0.428835,0.679123,0.210611}%
\pgfsetfillcolor{currentfill}%
\pgfsetlinewidth{0.000000pt}%
\definecolor{currentstroke}{rgb}{0.000000,0.000000,0.000000}%
\pgfsetstrokecolor{currentstroke}%
\pgfsetstrokeopacity{0.000000}%
\pgfsetdash{}{0pt}%
\pgfpathmoveto{\pgfqpoint{2.233749in}{0.610119in}}%
\pgfpathlineto{\pgfqpoint{2.159178in}{0.601837in}}%
\pgfpathlineto{\pgfqpoint{2.093446in}{0.596734in}}%
\pgfpathlineto{\pgfqpoint{2.092624in}{0.588678in}}%
\pgfpathlineto{\pgfqpoint{2.098661in}{0.581011in}}%
\pgfpathlineto{\pgfqpoint{2.106031in}{0.576507in}}%
\pgfpathlineto{\pgfqpoint{2.105919in}{0.564647in}}%
\pgfpathlineto{\pgfqpoint{2.103966in}{0.556723in}}%
\pgfpathlineto{\pgfqpoint{2.097462in}{0.551013in}}%
\pgfpathlineto{\pgfqpoint{2.088396in}{0.551618in}}%
\pgfpathlineto{\pgfqpoint{2.079825in}{0.558693in}}%
\pgfpathlineto{\pgfqpoint{2.078297in}{0.571283in}}%
\pgfpathlineto{\pgfqpoint{2.071913in}{0.578606in}}%
\pgfpathlineto{\pgfqpoint{2.067568in}{0.552398in}}%
\pgfpathlineto{\pgfqpoint{2.052803in}{0.554889in}}%
\pgfpathlineto{\pgfqpoint{2.047669in}{0.600717in}}%
\pgfpathlineto{\pgfqpoint{2.042108in}{0.649015in}}%
\pgfpathlineto{\pgfqpoint{2.044173in}{0.718697in}}%
\pgfpathlineto{\pgfqpoint{2.045455in}{0.766513in}}%
\pgfpathlineto{\pgfqpoint{2.048181in}{0.839470in}}%
\pgfpathlineto{\pgfqpoint{2.042346in}{0.845890in}}%
\pgfpathlineto{\pgfqpoint{2.052965in}{0.847429in}}%
\pgfpathlineto{\pgfqpoint{2.114830in}{0.851430in}}%
\pgfpathlineto{\pgfqpoint{2.175022in}{0.856843in}}%
\pgfpathlineto{\pgfqpoint{2.190858in}{0.801373in}}%
\pgfpathlineto{\pgfqpoint{2.201658in}{0.760657in}}%
\pgfpathlineto{\pgfqpoint{2.211339in}{0.726571in}}%
\pgfpathlineto{\pgfqpoint{2.216938in}{0.712845in}}%
\pgfpathlineto{\pgfqpoint{2.225753in}{0.700073in}}%
\pgfpathlineto{\pgfqpoint{2.224330in}{0.693573in}}%
\pgfpathlineto{\pgfqpoint{2.230627in}{0.690406in}}%
\pgfpathlineto{\pgfqpoint{2.223185in}{0.680553in}}%
\pgfpathlineto{\pgfqpoint{2.220672in}{0.663006in}}%
\pgfpathlineto{\pgfqpoint{2.227945in}{0.642483in}}%
\pgfpathlineto{\pgfqpoint{2.226935in}{0.621831in}}%
\pgfpathlineto{\pgfqpoint{2.233749in}{0.610119in}}%
\pgfpathclose%
\pgfusepath{fill}%
\end{pgfscope}%
\begin{pgfscope}%
\pgfpathrectangle{\pgfqpoint{0.100000in}{0.100000in}}{\pgfqpoint{2.862985in}{1.832778in}}%
\pgfusepath{clip}%
\pgfsetbuttcap%
\pgfsetmiterjoin%
\definecolor{currentfill}{rgb}{0.451903,0.698501,0.225375}%
\pgfsetfillcolor{currentfill}%
\pgfsetlinewidth{0.000000pt}%
\definecolor{currentstroke}{rgb}{0.000000,0.000000,0.000000}%
\pgfsetstrokecolor{currentstroke}%
\pgfsetstrokeopacity{0.000000}%
\pgfsetdash{}{0pt}%
\pgfpathmoveto{\pgfqpoint{2.052803in}{0.554889in}}%
\pgfpathlineto{\pgfqpoint{2.037652in}{0.550553in}}%
\pgfpathlineto{\pgfqpoint{2.026358in}{0.553337in}}%
\pgfpathlineto{\pgfqpoint{2.005384in}{0.546666in}}%
\pgfpathlineto{\pgfqpoint{1.989565in}{0.538074in}}%
\pgfpathlineto{\pgfqpoint{1.983075in}{0.553600in}}%
\pgfpathlineto{\pgfqpoint{1.976384in}{0.560107in}}%
\pgfpathlineto{\pgfqpoint{1.973157in}{0.567894in}}%
\pgfpathlineto{\pgfqpoint{1.977888in}{0.588972in}}%
\pgfpathlineto{\pgfqpoint{1.933060in}{0.586211in}}%
\pgfpathlineto{\pgfqpoint{1.874932in}{0.583691in}}%
\pgfpathlineto{\pgfqpoint{1.878390in}{0.588938in}}%
\pgfpathlineto{\pgfqpoint{1.874116in}{0.598847in}}%
\pgfpathlineto{\pgfqpoint{1.878375in}{0.600866in}}%
\pgfpathlineto{\pgfqpoint{1.877430in}{0.610386in}}%
\pgfpathlineto{\pgfqpoint{1.884854in}{0.619619in}}%
\pgfpathlineto{\pgfqpoint{1.888916in}{0.637542in}}%
\pgfpathlineto{\pgfqpoint{1.902503in}{0.649350in}}%
\pgfpathlineto{\pgfqpoint{1.897467in}{0.660818in}}%
\pgfpathlineto{\pgfqpoint{1.907102in}{0.667419in}}%
\pgfpathlineto{\pgfqpoint{1.898791in}{0.679370in}}%
\pgfpathlineto{\pgfqpoint{1.892776in}{0.706727in}}%
\pgfpathlineto{\pgfqpoint{1.895057in}{0.711695in}}%
\pgfpathlineto{\pgfqpoint{1.898659in}{0.721223in}}%
\pgfpathlineto{\pgfqpoint{1.895307in}{0.731223in}}%
\pgfpathlineto{\pgfqpoint{1.896946in}{0.735772in}}%
\pgfpathlineto{\pgfqpoint{1.891048in}{0.752960in}}%
\pgfpathlineto{\pgfqpoint{1.908459in}{0.783821in}}%
\pgfpathlineto{\pgfqpoint{1.909269in}{0.792035in}}%
\pgfpathlineto{\pgfqpoint{1.917656in}{0.798007in}}%
\pgfpathlineto{\pgfqpoint{1.926605in}{0.817725in}}%
\pgfpathlineto{\pgfqpoint{1.926180in}{0.825739in}}%
\pgfpathlineto{\pgfqpoint{1.937326in}{0.833932in}}%
\pgfpathlineto{\pgfqpoint{1.933848in}{0.838785in}}%
\pgfpathlineto{\pgfqpoint{2.000643in}{0.842683in}}%
\pgfpathlineto{\pgfqpoint{2.042346in}{0.845890in}}%
\pgfpathlineto{\pgfqpoint{2.048181in}{0.839470in}}%
\pgfpathlineto{\pgfqpoint{2.045455in}{0.766513in}}%
\pgfpathlineto{\pgfqpoint{2.044173in}{0.718697in}}%
\pgfpathlineto{\pgfqpoint{2.042108in}{0.649015in}}%
\pgfpathlineto{\pgfqpoint{2.047669in}{0.600717in}}%
\pgfpathlineto{\pgfqpoint{2.052803in}{0.554889in}}%
\pgfpathclose%
\pgfusepath{fill}%
\end{pgfscope}%
\begin{pgfscope}%
\pgfpathrectangle{\pgfqpoint{0.100000in}{0.100000in}}{\pgfqpoint{2.862985in}{1.832778in}}%
\pgfusepath{clip}%
\pgfsetbuttcap%
\pgfsetmiterjoin%
\definecolor{currentfill}{rgb}{0.541869,0.765705,0.307958}%
\pgfsetfillcolor{currentfill}%
\pgfsetlinewidth{0.000000pt}%
\definecolor{currentstroke}{rgb}{0.000000,0.000000,0.000000}%
\pgfsetstrokecolor{currentstroke}%
\pgfsetstrokeopacity{0.000000}%
\pgfsetdash{}{0pt}%
\pgfpathmoveto{\pgfqpoint{2.175022in}{0.856843in}}%
\pgfpathlineto{\pgfqpoint{2.240820in}{0.864070in}}%
\pgfpathlineto{\pgfqpoint{2.280410in}{0.868948in}}%
\pgfpathlineto{\pgfqpoint{2.302958in}{0.872540in}}%
\pgfpathlineto{\pgfqpoint{2.293618in}{0.857809in}}%
\pgfpathlineto{\pgfqpoint{2.293623in}{0.850863in}}%
\pgfpathlineto{\pgfqpoint{2.303238in}{0.847117in}}%
\pgfpathlineto{\pgfqpoint{2.309776in}{0.841070in}}%
\pgfpathlineto{\pgfqpoint{2.319657in}{0.840323in}}%
\pgfpathlineto{\pgfqpoint{2.328892in}{0.823299in}}%
\pgfpathlineto{\pgfqpoint{2.338910in}{0.811294in}}%
\pgfpathlineto{\pgfqpoint{2.356702in}{0.801205in}}%
\pgfpathlineto{\pgfqpoint{2.361758in}{0.793157in}}%
\pgfpathlineto{\pgfqpoint{2.377624in}{0.784454in}}%
\pgfpathlineto{\pgfqpoint{2.378817in}{0.778143in}}%
\pgfpathlineto{\pgfqpoint{2.391688in}{0.765593in}}%
\pgfpathlineto{\pgfqpoint{2.401546in}{0.762019in}}%
\pgfpathlineto{\pgfqpoint{2.408528in}{0.750169in}}%
\pgfpathlineto{\pgfqpoint{2.411618in}{0.736849in}}%
\pgfpathlineto{\pgfqpoint{2.421858in}{0.731700in}}%
\pgfpathlineto{\pgfqpoint{2.430102in}{0.717383in}}%
\pgfpathlineto{\pgfqpoint{2.432761in}{0.706853in}}%
\pgfpathlineto{\pgfqpoint{2.444841in}{0.702371in}}%
\pgfpathlineto{\pgfqpoint{2.434704in}{0.682961in}}%
\pgfpathlineto{\pgfqpoint{2.430889in}{0.671491in}}%
\pgfpathlineto{\pgfqpoint{2.433378in}{0.666114in}}%
\pgfpathlineto{\pgfqpoint{2.429032in}{0.657178in}}%
\pgfpathlineto{\pgfqpoint{2.427185in}{0.644824in}}%
\pgfpathlineto{\pgfqpoint{2.419456in}{0.642523in}}%
\pgfpathlineto{\pgfqpoint{2.423537in}{0.631212in}}%
\pgfpathlineto{\pgfqpoint{2.423277in}{0.616862in}}%
\pgfpathlineto{\pgfqpoint{2.409539in}{0.617400in}}%
\pgfpathlineto{\pgfqpoint{2.399969in}{0.620018in}}%
\pgfpathlineto{\pgfqpoint{2.395828in}{0.615151in}}%
\pgfpathlineto{\pgfqpoint{2.398918in}{0.603653in}}%
\pgfpathlineto{\pgfqpoint{2.398246in}{0.590285in}}%
\pgfpathlineto{\pgfqpoint{2.392209in}{0.589262in}}%
\pgfpathlineto{\pgfqpoint{2.387325in}{0.601747in}}%
\pgfpathlineto{\pgfqpoint{2.337593in}{0.598433in}}%
\pgfpathlineto{\pgfqpoint{2.274885in}{0.595011in}}%
\pgfpathlineto{\pgfqpoint{2.243260in}{0.592785in}}%
\pgfpathlineto{\pgfqpoint{2.233749in}{0.610119in}}%
\pgfpathlineto{\pgfqpoint{2.226935in}{0.621831in}}%
\pgfpathlineto{\pgfqpoint{2.227945in}{0.642483in}}%
\pgfpathlineto{\pgfqpoint{2.220672in}{0.663006in}}%
\pgfpathlineto{\pgfqpoint{2.223185in}{0.680553in}}%
\pgfpathlineto{\pgfqpoint{2.230627in}{0.690406in}}%
\pgfpathlineto{\pgfqpoint{2.224330in}{0.693573in}}%
\pgfpathlineto{\pgfqpoint{2.225753in}{0.700073in}}%
\pgfpathlineto{\pgfqpoint{2.216938in}{0.712845in}}%
\pgfpathlineto{\pgfqpoint{2.211339in}{0.726571in}}%
\pgfpathlineto{\pgfqpoint{2.201658in}{0.760657in}}%
\pgfpathlineto{\pgfqpoint{2.190858in}{0.801373in}}%
\pgfpathlineto{\pgfqpoint{2.175022in}{0.856843in}}%
\pgfpathclose%
\pgfusepath{fill}%
\end{pgfscope}%
\begin{pgfscope}%
\pgfpathrectangle{\pgfqpoint{0.100000in}{0.100000in}}{\pgfqpoint{2.862985in}{1.832778in}}%
\pgfusepath{clip}%
\pgfsetbuttcap%
\pgfsetmiterjoin%
\definecolor{currentfill}{rgb}{0.524337,0.754325,0.286736}%
\pgfsetfillcolor{currentfill}%
\pgfsetlinewidth{0.000000pt}%
\definecolor{currentstroke}{rgb}{0.000000,0.000000,0.000000}%
\pgfsetstrokecolor{currentstroke}%
\pgfsetstrokeopacity{0.000000}%
\pgfsetdash{}{0pt}%
\pgfpathmoveto{\pgfqpoint{2.302958in}{0.872540in}}%
\pgfpathlineto{\pgfqpoint{2.329049in}{0.885519in}}%
\pgfpathlineto{\pgfqpoint{2.343499in}{0.890420in}}%
\pgfpathlineto{\pgfqpoint{2.406366in}{0.896952in}}%
\pgfpathlineto{\pgfqpoint{2.413024in}{0.894789in}}%
\pgfpathlineto{\pgfqpoint{2.421819in}{0.885949in}}%
\pgfpathlineto{\pgfqpoint{2.422304in}{0.878098in}}%
\pgfpathlineto{\pgfqpoint{2.480165in}{0.886600in}}%
\pgfpathlineto{\pgfqpoint{2.545869in}{0.839501in}}%
\pgfpathlineto{\pgfqpoint{2.533547in}{0.826658in}}%
\pgfpathlineto{\pgfqpoint{2.516640in}{0.796833in}}%
\pgfpathlineto{\pgfqpoint{2.521453in}{0.790423in}}%
\pgfpathlineto{\pgfqpoint{2.512434in}{0.777981in}}%
\pgfpathlineto{\pgfqpoint{2.503462in}{0.776573in}}%
\pgfpathlineto{\pgfqpoint{2.503429in}{0.769097in}}%
\pgfpathlineto{\pgfqpoint{2.496939in}{0.761274in}}%
\pgfpathlineto{\pgfqpoint{2.488862in}{0.759687in}}%
\pgfpathlineto{\pgfqpoint{2.490625in}{0.752736in}}%
\pgfpathlineto{\pgfqpoint{2.486118in}{0.747400in}}%
\pgfpathlineto{\pgfqpoint{2.475337in}{0.742790in}}%
\pgfpathlineto{\pgfqpoint{2.462196in}{0.732304in}}%
\pgfpathlineto{\pgfqpoint{2.464677in}{0.725520in}}%
\pgfpathlineto{\pgfqpoint{2.456407in}{0.721260in}}%
\pgfpathlineto{\pgfqpoint{2.449432in}{0.725739in}}%
\pgfpathlineto{\pgfqpoint{2.444331in}{0.706267in}}%
\pgfpathlineto{\pgfqpoint{2.432761in}{0.706853in}}%
\pgfpathlineto{\pgfqpoint{2.430102in}{0.717383in}}%
\pgfpathlineto{\pgfqpoint{2.421858in}{0.731700in}}%
\pgfpathlineto{\pgfqpoint{2.411618in}{0.736849in}}%
\pgfpathlineto{\pgfqpoint{2.408528in}{0.750169in}}%
\pgfpathlineto{\pgfqpoint{2.401546in}{0.762019in}}%
\pgfpathlineto{\pgfqpoint{2.391688in}{0.765593in}}%
\pgfpathlineto{\pgfqpoint{2.378817in}{0.778143in}}%
\pgfpathlineto{\pgfqpoint{2.377624in}{0.784454in}}%
\pgfpathlineto{\pgfqpoint{2.361758in}{0.793157in}}%
\pgfpathlineto{\pgfqpoint{2.356702in}{0.801205in}}%
\pgfpathlineto{\pgfqpoint{2.338910in}{0.811294in}}%
\pgfpathlineto{\pgfqpoint{2.328892in}{0.823299in}}%
\pgfpathlineto{\pgfqpoint{2.319657in}{0.840323in}}%
\pgfpathlineto{\pgfqpoint{2.309776in}{0.841070in}}%
\pgfpathlineto{\pgfqpoint{2.303238in}{0.847117in}}%
\pgfpathlineto{\pgfqpoint{2.293623in}{0.850863in}}%
\pgfpathlineto{\pgfqpoint{2.293618in}{0.857809in}}%
\pgfpathlineto{\pgfqpoint{2.302958in}{0.872540in}}%
\pgfpathclose%
\pgfusepath{fill}%
\end{pgfscope}%
\begin{pgfscope}%
\pgfpathrectangle{\pgfqpoint{0.100000in}{0.100000in}}{\pgfqpoint{2.862985in}{1.832778in}}%
\pgfusepath{clip}%
\pgfsetbuttcap%
\pgfsetmiterjoin%
\definecolor{currentfill}{rgb}{0.568166,0.782776,0.339792}%
\pgfsetfillcolor{currentfill}%
\pgfsetlinewidth{0.000000pt}%
\definecolor{currentstroke}{rgb}{0.000000,0.000000,0.000000}%
\pgfsetstrokecolor{currentstroke}%
\pgfsetstrokeopacity{0.000000}%
\pgfsetdash{}{0pt}%
\pgfpathmoveto{\pgfqpoint{1.711515in}{0.926094in}}%
\pgfpathlineto{\pgfqpoint{1.749856in}{0.926440in}}%
\pgfpathlineto{\pgfqpoint{1.800580in}{0.927364in}}%
\pgfpathlineto{\pgfqpoint{1.887257in}{0.930116in}}%
\pgfpathlineto{\pgfqpoint{1.936833in}{0.932737in}}%
\pgfpathlineto{\pgfqpoint{1.942150in}{0.926155in}}%
\pgfpathlineto{\pgfqpoint{1.941818in}{0.919204in}}%
\pgfpathlineto{\pgfqpoint{1.929808in}{0.907211in}}%
\pgfpathlineto{\pgfqpoint{1.926904in}{0.900644in}}%
\pgfpathlineto{\pgfqpoint{1.960213in}{0.903107in}}%
\pgfpathlineto{\pgfqpoint{1.960195in}{0.890974in}}%
\pgfpathlineto{\pgfqpoint{1.949509in}{0.885789in}}%
\pgfpathlineto{\pgfqpoint{1.949660in}{0.877545in}}%
\pgfpathlineto{\pgfqpoint{1.945647in}{0.871619in}}%
\pgfpathlineto{\pgfqpoint{1.943020in}{0.859033in}}%
\pgfpathlineto{\pgfqpoint{1.945732in}{0.848649in}}%
\pgfpathlineto{\pgfqpoint{1.933848in}{0.838785in}}%
\pgfpathlineto{\pgfqpoint{1.937326in}{0.833932in}}%
\pgfpathlineto{\pgfqpoint{1.926180in}{0.825739in}}%
\pgfpathlineto{\pgfqpoint{1.926605in}{0.817725in}}%
\pgfpathlineto{\pgfqpoint{1.917656in}{0.798007in}}%
\pgfpathlineto{\pgfqpoint{1.909269in}{0.792035in}}%
\pgfpathlineto{\pgfqpoint{1.908459in}{0.783821in}}%
\pgfpathlineto{\pgfqpoint{1.891048in}{0.752960in}}%
\pgfpathlineto{\pgfqpoint{1.896946in}{0.735772in}}%
\pgfpathlineto{\pgfqpoint{1.895307in}{0.731223in}}%
\pgfpathlineto{\pgfqpoint{1.898659in}{0.721223in}}%
\pgfpathlineto{\pgfqpoint{1.895057in}{0.711695in}}%
\pgfpathlineto{\pgfqpoint{1.847446in}{0.709729in}}%
\pgfpathlineto{\pgfqpoint{1.785638in}{0.708711in}}%
\pgfpathlineto{\pgfqpoint{1.743006in}{0.708326in}}%
\pgfpathlineto{\pgfqpoint{1.742799in}{0.741894in}}%
\pgfpathlineto{\pgfqpoint{1.732212in}{0.744093in}}%
\pgfpathlineto{\pgfqpoint{1.725225in}{0.741221in}}%
\pgfpathlineto{\pgfqpoint{1.719643in}{0.746495in}}%
\pgfpathlineto{\pgfqpoint{1.720219in}{0.782052in}}%
\pgfpathlineto{\pgfqpoint{1.721468in}{0.857760in}}%
\pgfpathlineto{\pgfqpoint{1.715413in}{0.902092in}}%
\pgfpathlineto{\pgfqpoint{1.711515in}{0.926094in}}%
\pgfpathclose%
\pgfusepath{fill}%
\end{pgfscope}%
\begin{pgfscope}%
\pgfpathrectangle{\pgfqpoint{0.100000in}{0.100000in}}{\pgfqpoint{2.862985in}{1.832778in}}%
\pgfusepath{clip}%
\pgfsetbuttcap%
\pgfsetmiterjoin%
\definecolor{currentfill}{rgb}{0.673356,0.851057,0.467128}%
\pgfsetfillcolor{currentfill}%
\pgfsetlinewidth{0.000000pt}%
\definecolor{currentstroke}{rgb}{0.000000,0.000000,0.000000}%
\pgfsetstrokecolor{currentstroke}%
\pgfsetstrokeopacity{0.000000}%
\pgfsetdash{}{0pt}%
\pgfpathmoveto{\pgfqpoint{1.743006in}{0.708326in}}%
\pgfpathlineto{\pgfqpoint{1.785638in}{0.708711in}}%
\pgfpathlineto{\pgfqpoint{1.847446in}{0.709729in}}%
\pgfpathlineto{\pgfqpoint{1.895057in}{0.711695in}}%
\pgfpathlineto{\pgfqpoint{1.892776in}{0.706727in}}%
\pgfpathlineto{\pgfqpoint{1.898791in}{0.679370in}}%
\pgfpathlineto{\pgfqpoint{1.907102in}{0.667419in}}%
\pgfpathlineto{\pgfqpoint{1.897467in}{0.660818in}}%
\pgfpathlineto{\pgfqpoint{1.902503in}{0.649350in}}%
\pgfpathlineto{\pgfqpoint{1.888916in}{0.637542in}}%
\pgfpathlineto{\pgfqpoint{1.884854in}{0.619619in}}%
\pgfpathlineto{\pgfqpoint{1.877430in}{0.610386in}}%
\pgfpathlineto{\pgfqpoint{1.878375in}{0.600866in}}%
\pgfpathlineto{\pgfqpoint{1.874116in}{0.598847in}}%
\pgfpathlineto{\pgfqpoint{1.878390in}{0.588938in}}%
\pgfpathlineto{\pgfqpoint{1.874932in}{0.583691in}}%
\pgfpathlineto{\pgfqpoint{1.933060in}{0.586211in}}%
\pgfpathlineto{\pgfqpoint{1.977888in}{0.588972in}}%
\pgfpathlineto{\pgfqpoint{1.973157in}{0.567894in}}%
\pgfpathlineto{\pgfqpoint{1.976384in}{0.560107in}}%
\pgfpathlineto{\pgfqpoint{1.983075in}{0.553600in}}%
\pgfpathlineto{\pgfqpoint{1.989565in}{0.538074in}}%
\pgfpathlineto{\pgfqpoint{1.969065in}{0.541651in}}%
\pgfpathlineto{\pgfqpoint{1.961502in}{0.547531in}}%
\pgfpathlineto{\pgfqpoint{1.952492in}{0.547811in}}%
\pgfpathlineto{\pgfqpoint{1.943019in}{0.534934in}}%
\pgfpathlineto{\pgfqpoint{1.944907in}{0.529071in}}%
\pgfpathlineto{\pgfqpoint{1.975169in}{0.525523in}}%
\pgfpathlineto{\pgfqpoint{1.983091in}{0.518783in}}%
\pgfpathlineto{\pgfqpoint{1.995984in}{0.515315in}}%
\pgfpathlineto{\pgfqpoint{1.990353in}{0.507336in}}%
\pgfpathlineto{\pgfqpoint{1.980867in}{0.500376in}}%
\pgfpathlineto{\pgfqpoint{2.005906in}{0.484726in}}%
\pgfpathlineto{\pgfqpoint{2.017569in}{0.482281in}}%
\pgfpathlineto{\pgfqpoint{2.022419in}{0.469488in}}%
\pgfpathlineto{\pgfqpoint{2.011438in}{0.467135in}}%
\pgfpathlineto{\pgfqpoint{1.990799in}{0.483224in}}%
\pgfpathlineto{\pgfqpoint{1.982729in}{0.485461in}}%
\pgfpathlineto{\pgfqpoint{1.978824in}{0.491817in}}%
\pgfpathlineto{\pgfqpoint{1.966854in}{0.489208in}}%
\pgfpathlineto{\pgfqpoint{1.963107in}{0.481018in}}%
\pgfpathlineto{\pgfqpoint{1.965480in}{0.471892in}}%
\pgfpathlineto{\pgfqpoint{1.957422in}{0.466506in}}%
\pgfpathlineto{\pgfqpoint{1.950068in}{0.479774in}}%
\pgfpathlineto{\pgfqpoint{1.937187in}{0.475818in}}%
\pgfpathlineto{\pgfqpoint{1.932360in}{0.467901in}}%
\pgfpathlineto{\pgfqpoint{1.923213in}{0.470153in}}%
\pgfpathlineto{\pgfqpoint{1.917364in}{0.480150in}}%
\pgfpathlineto{\pgfqpoint{1.901847in}{0.483451in}}%
\pgfpathlineto{\pgfqpoint{1.898917in}{0.488653in}}%
\pgfpathlineto{\pgfqpoint{1.882751in}{0.497715in}}%
\pgfpathlineto{\pgfqpoint{1.878738in}{0.505638in}}%
\pgfpathlineto{\pgfqpoint{1.865197in}{0.502400in}}%
\pgfpathlineto{\pgfqpoint{1.866917in}{0.509654in}}%
\pgfpathlineto{\pgfqpoint{1.850086in}{0.502213in}}%
\pgfpathlineto{\pgfqpoint{1.854607in}{0.494492in}}%
\pgfpathlineto{\pgfqpoint{1.841610in}{0.489928in}}%
\pgfpathlineto{\pgfqpoint{1.824395in}{0.492437in}}%
\pgfpathlineto{\pgfqpoint{1.789551in}{0.504372in}}%
\pgfpathlineto{\pgfqpoint{1.762666in}{0.502003in}}%
\pgfpathlineto{\pgfqpoint{1.760335in}{0.517680in}}%
\pgfpathlineto{\pgfqpoint{1.763712in}{0.524834in}}%
\pgfpathlineto{\pgfqpoint{1.763423in}{0.536215in}}%
\pgfpathlineto{\pgfqpoint{1.760101in}{0.538927in}}%
\pgfpathlineto{\pgfqpoint{1.761227in}{0.551919in}}%
\pgfpathlineto{\pgfqpoint{1.771019in}{0.569988in}}%
\pgfpathlineto{\pgfqpoint{1.772272in}{0.576816in}}%
\pgfpathlineto{\pgfqpoint{1.770658in}{0.592942in}}%
\pgfpathlineto{\pgfqpoint{1.763454in}{0.600538in}}%
\pgfpathlineto{\pgfqpoint{1.759693in}{0.614721in}}%
\pgfpathlineto{\pgfqpoint{1.755065in}{0.617977in}}%
\pgfpathlineto{\pgfqpoint{1.757137in}{0.625668in}}%
\pgfpathlineto{\pgfqpoint{1.751243in}{0.637164in}}%
\pgfpathlineto{\pgfqpoint{1.743890in}{0.643403in}}%
\pgfpathlineto{\pgfqpoint{1.743006in}{0.708326in}}%
\pgfpathclose%
\pgfusepath{fill}%
\end{pgfscope}%
\begin{pgfscope}%
\pgfpathrectangle{\pgfqpoint{0.100000in}{0.100000in}}{\pgfqpoint{2.862985in}{1.832778in}}%
\pgfusepath{clip}%
\pgfsetbuttcap%
\pgfsetmiterjoin%
\definecolor{currentfill}{rgb}{0.673356,0.851057,0.467128}%
\pgfsetfillcolor{currentfill}%
\pgfsetlinewidth{0.000000pt}%
\definecolor{currentstroke}{rgb}{0.000000,0.000000,0.000000}%
\pgfsetstrokecolor{currentstroke}%
\pgfsetstrokeopacity{0.000000}%
\pgfsetdash{}{0pt}%
\pgfpathmoveto{\pgfqpoint{1.857025in}{0.493901in}}%
\pgfpathlineto{\pgfqpoint{1.863193in}{0.497570in}}%
\pgfpathlineto{\pgfqpoint{1.870679in}{0.493233in}}%
\pgfpathlineto{\pgfqpoint{1.866504in}{0.487265in}}%
\pgfpathlineto{\pgfqpoint{1.857025in}{0.493901in}}%
\pgfpathclose%
\pgfusepath{fill}%
\end{pgfscope}%
\begin{pgfscope}%
\pgfpathrectangle{\pgfqpoint{0.100000in}{0.100000in}}{\pgfqpoint{2.862985in}{1.832778in}}%
\pgfusepath{clip}%
\pgfsetbuttcap%
\pgfsetmiterjoin%
\definecolor{currentfill}{rgb}{0.524337,0.754325,0.286736}%
\pgfsetfillcolor{currentfill}%
\pgfsetlinewidth{0.000000pt}%
\definecolor{currentstroke}{rgb}{0.000000,0.000000,0.000000}%
\pgfsetstrokecolor{currentstroke}%
\pgfsetstrokeopacity{0.000000}%
\pgfsetdash{}{0pt}%
\pgfpathmoveto{\pgfqpoint{2.105919in}{0.564647in}}%
\pgfpathlineto{\pgfqpoint{2.106031in}{0.576507in}}%
\pgfpathlineto{\pgfqpoint{2.098661in}{0.581011in}}%
\pgfpathlineto{\pgfqpoint{2.092624in}{0.588678in}}%
\pgfpathlineto{\pgfqpoint{2.093446in}{0.596734in}}%
\pgfpathlineto{\pgfqpoint{2.159178in}{0.601837in}}%
\pgfpathlineto{\pgfqpoint{2.233749in}{0.610119in}}%
\pgfpathlineto{\pgfqpoint{2.243260in}{0.592785in}}%
\pgfpathlineto{\pgfqpoint{2.274885in}{0.595011in}}%
\pgfpathlineto{\pgfqpoint{2.337593in}{0.598433in}}%
\pgfpathlineto{\pgfqpoint{2.387325in}{0.601747in}}%
\pgfpathlineto{\pgfqpoint{2.392209in}{0.589262in}}%
\pgfpathlineto{\pgfqpoint{2.398246in}{0.590285in}}%
\pgfpathlineto{\pgfqpoint{2.398918in}{0.603653in}}%
\pgfpathlineto{\pgfqpoint{2.395828in}{0.615151in}}%
\pgfpathlineto{\pgfqpoint{2.399969in}{0.620018in}}%
\pgfpathlineto{\pgfqpoint{2.409539in}{0.617400in}}%
\pgfpathlineto{\pgfqpoint{2.423277in}{0.616862in}}%
\pgfpathlineto{\pgfqpoint{2.425389in}{0.606589in}}%
\pgfpathlineto{\pgfqpoint{2.429615in}{0.600719in}}%
\pgfpathlineto{\pgfqpoint{2.432915in}{0.587877in}}%
\pgfpathlineto{\pgfqpoint{2.443139in}{0.568003in}}%
\pgfpathlineto{\pgfqpoint{2.443190in}{0.562628in}}%
\pgfpathlineto{\pgfqpoint{2.458302in}{0.539291in}}%
\pgfpathlineto{\pgfqpoint{2.459764in}{0.534471in}}%
\pgfpathlineto{\pgfqpoint{2.482434in}{0.501512in}}%
\pgfpathlineto{\pgfqpoint{2.478842in}{0.500979in}}%
\pgfpathlineto{\pgfqpoint{2.488399in}{0.477537in}}%
\pgfpathlineto{\pgfqpoint{2.514653in}{0.436782in}}%
\pgfpathlineto{\pgfqpoint{2.529896in}{0.406818in}}%
\pgfpathlineto{\pgfqpoint{2.534991in}{0.401571in}}%
\pgfpathlineto{\pgfqpoint{2.543722in}{0.382763in}}%
\pgfpathlineto{\pgfqpoint{2.546750in}{0.352646in}}%
\pgfpathlineto{\pgfqpoint{2.547955in}{0.330127in}}%
\pgfpathlineto{\pgfqpoint{2.546500in}{0.315709in}}%
\pgfpathlineto{\pgfqpoint{2.541819in}{0.305401in}}%
\pgfpathlineto{\pgfqpoint{2.543994in}{0.291885in}}%
\pgfpathlineto{\pgfqpoint{2.538954in}{0.281186in}}%
\pgfpathlineto{\pgfqpoint{2.523991in}{0.272413in}}%
\pgfpathlineto{\pgfqpoint{2.514259in}{0.273053in}}%
\pgfpathlineto{\pgfqpoint{2.507940in}{0.268467in}}%
\pgfpathlineto{\pgfqpoint{2.506129in}{0.280712in}}%
\pgfpathlineto{\pgfqpoint{2.495668in}{0.283944in}}%
\pgfpathlineto{\pgfqpoint{2.486290in}{0.301030in}}%
\pgfpathlineto{\pgfqpoint{2.485231in}{0.308810in}}%
\pgfpathlineto{\pgfqpoint{2.468448in}{0.313535in}}%
\pgfpathlineto{\pgfqpoint{2.457633in}{0.312512in}}%
\pgfpathlineto{\pgfqpoint{2.451514in}{0.323801in}}%
\pgfpathlineto{\pgfqpoint{2.444508in}{0.344093in}}%
\pgfpathlineto{\pgfqpoint{2.434786in}{0.348206in}}%
\pgfpathlineto{\pgfqpoint{2.429515in}{0.359834in}}%
\pgfpathlineto{\pgfqpoint{2.416625in}{0.366687in}}%
\pgfpathlineto{\pgfqpoint{2.409174in}{0.375259in}}%
\pgfpathlineto{\pgfqpoint{2.396320in}{0.398564in}}%
\pgfpathlineto{\pgfqpoint{2.394809in}{0.414966in}}%
\pgfpathlineto{\pgfqpoint{2.401813in}{0.425461in}}%
\pgfpathlineto{\pgfqpoint{2.401113in}{0.432754in}}%
\pgfpathlineto{\pgfqpoint{2.393019in}{0.433533in}}%
\pgfpathlineto{\pgfqpoint{2.386256in}{0.438606in}}%
\pgfpathlineto{\pgfqpoint{2.382553in}{0.432412in}}%
\pgfpathlineto{\pgfqpoint{2.389043in}{0.427365in}}%
\pgfpathlineto{\pgfqpoint{2.384319in}{0.418289in}}%
\pgfpathlineto{\pgfqpoint{2.376677in}{0.425822in}}%
\pgfpathlineto{\pgfqpoint{2.377564in}{0.446710in}}%
\pgfpathlineto{\pgfqpoint{2.381250in}{0.463657in}}%
\pgfpathlineto{\pgfqpoint{2.379355in}{0.492736in}}%
\pgfpathlineto{\pgfqpoint{2.371738in}{0.499662in}}%
\pgfpathlineto{\pgfqpoint{2.367907in}{0.508549in}}%
\pgfpathlineto{\pgfqpoint{2.354785in}{0.508352in}}%
\pgfpathlineto{\pgfqpoint{2.341814in}{0.522976in}}%
\pgfpathlineto{\pgfqpoint{2.333117in}{0.527359in}}%
\pgfpathlineto{\pgfqpoint{2.330538in}{0.536620in}}%
\pgfpathlineto{\pgfqpoint{2.322031in}{0.539872in}}%
\pgfpathlineto{\pgfqpoint{2.314981in}{0.550098in}}%
\pgfpathlineto{\pgfqpoint{2.296358in}{0.558444in}}%
\pgfpathlineto{\pgfqpoint{2.281879in}{0.558658in}}%
\pgfpathlineto{\pgfqpoint{2.275576in}{0.555452in}}%
\pgfpathlineto{\pgfqpoint{2.277143in}{0.545445in}}%
\pgfpathlineto{\pgfqpoint{2.270570in}{0.545919in}}%
\pgfpathlineto{\pgfqpoint{2.250353in}{0.531927in}}%
\pgfpathlineto{\pgfqpoint{2.226068in}{0.526381in}}%
\pgfpathlineto{\pgfqpoint{2.225669in}{0.533243in}}%
\pgfpathlineto{\pgfqpoint{2.220287in}{0.539949in}}%
\pgfpathlineto{\pgfqpoint{2.205848in}{0.549208in}}%
\pgfpathlineto{\pgfqpoint{2.185121in}{0.558700in}}%
\pgfpathlineto{\pgfqpoint{2.162608in}{0.563725in}}%
\pgfpathlineto{\pgfqpoint{2.158382in}{0.570558in}}%
\pgfpathlineto{\pgfqpoint{2.150265in}{0.564859in}}%
\pgfpathlineto{\pgfqpoint{2.140491in}{0.563602in}}%
\pgfpathlineto{\pgfqpoint{2.106467in}{0.554637in}}%
\pgfpathlineto{\pgfqpoint{2.105919in}{0.564647in}}%
\pgfpathclose%
\pgfusepath{fill}%
\end{pgfscope}%
\begin{pgfscope}%
\pgfpathrectangle{\pgfqpoint{0.100000in}{0.100000in}}{\pgfqpoint{2.862985in}{1.832778in}}%
\pgfusepath{clip}%
\pgfsetbuttcap%
\pgfsetmiterjoin%
\definecolor{currentfill}{rgb}{0.524337,0.754325,0.286736}%
\pgfsetfillcolor{currentfill}%
\pgfsetlinewidth{0.000000pt}%
\definecolor{currentstroke}{rgb}{0.000000,0.000000,0.000000}%
\pgfsetstrokecolor{currentstroke}%
\pgfsetstrokeopacity{0.000000}%
\pgfsetdash{}{0pt}%
\pgfpathmoveto{\pgfqpoint{2.485737in}{0.501980in}}%
\pgfpathlineto{\pgfqpoint{2.494540in}{0.491740in}}%
\pgfpathlineto{\pgfqpoint{2.484372in}{0.491096in}}%
\pgfpathlineto{\pgfqpoint{2.485737in}{0.501980in}}%
\pgfpathclose%
\pgfusepath{fill}%
\end{pgfscope}%
\begin{pgfscope}%
\pgfpathrectangle{\pgfqpoint{0.100000in}{0.100000in}}{\pgfqpoint{2.862985in}{1.832778in}}%
\pgfusepath{clip}%
\pgfsetbuttcap%
\pgfsetmiterjoin%
\definecolor{currentfill}{rgb}{0.690888,0.862438,0.488351}%
\pgfsetfillcolor{currentfill}%
\pgfsetlinewidth{0.000000pt}%
\definecolor{currentstroke}{rgb}{0.000000,0.000000,0.000000}%
\pgfsetstrokecolor{currentstroke}%
\pgfsetstrokeopacity{0.000000}%
\pgfsetdash{}{0pt}%
\pgfpathmoveto{\pgfqpoint{1.967522in}{1.671190in}}%
\pgfpathlineto{\pgfqpoint{1.962746in}{1.661894in}}%
\pgfpathlineto{\pgfqpoint{1.951351in}{1.656468in}}%
\pgfpathlineto{\pgfqpoint{1.946416in}{1.649166in}}%
\pgfpathlineto{\pgfqpoint{1.940635in}{1.654425in}}%
\pgfpathlineto{\pgfqpoint{1.967522in}{1.671190in}}%
\pgfpathclose%
\pgfusepath{fill}%
\end{pgfscope}%
\begin{pgfscope}%
\pgfpathrectangle{\pgfqpoint{0.100000in}{0.100000in}}{\pgfqpoint{2.862985in}{1.832778in}}%
\pgfusepath{clip}%
\pgfsetbuttcap%
\pgfsetmiterjoin%
\definecolor{currentfill}{rgb}{0.690888,0.862438,0.488351}%
\pgfsetfillcolor{currentfill}%
\pgfsetlinewidth{0.000000pt}%
\definecolor{currentstroke}{rgb}{0.000000,0.000000,0.000000}%
\pgfsetstrokecolor{currentstroke}%
\pgfsetstrokeopacity{0.000000}%
\pgfsetdash{}{0pt}%
\pgfpathmoveto{\pgfqpoint{1.971396in}{1.615244in}}%
\pgfpathlineto{\pgfqpoint{1.983054in}{1.626120in}}%
\pgfpathlineto{\pgfqpoint{2.001037in}{1.628969in}}%
\pgfpathlineto{\pgfqpoint{1.996084in}{1.621405in}}%
\pgfpathlineto{\pgfqpoint{1.983756in}{1.610466in}}%
\pgfpathlineto{\pgfqpoint{1.976549in}{1.596426in}}%
\pgfpathlineto{\pgfqpoint{1.973496in}{1.604041in}}%
\pgfpathlineto{\pgfqpoint{1.968076in}{1.605134in}}%
\pgfpathlineto{\pgfqpoint{1.967532in}{1.612009in}}%
\pgfpathlineto{\pgfqpoint{1.971396in}{1.615244in}}%
\pgfpathclose%
\pgfusepath{fill}%
\end{pgfscope}%
\begin{pgfscope}%
\pgfpathrectangle{\pgfqpoint{0.100000in}{0.100000in}}{\pgfqpoint{2.862985in}{1.832778in}}%
\pgfusepath{clip}%
\pgfsetbuttcap%
\pgfsetmiterjoin%
\definecolor{currentfill}{rgb}{0.690888,0.862438,0.488351}%
\pgfsetfillcolor{currentfill}%
\pgfsetlinewidth{0.000000pt}%
\definecolor{currentstroke}{rgb}{0.000000,0.000000,0.000000}%
\pgfsetstrokecolor{currentstroke}%
\pgfsetstrokeopacity{0.000000}%
\pgfsetdash{}{0pt}%
\pgfpathmoveto{\pgfqpoint{2.017832in}{1.482495in}}%
\pgfpathlineto{\pgfqpoint{2.014727in}{1.485941in}}%
\pgfpathlineto{\pgfqpoint{2.017991in}{1.495641in}}%
\pgfpathlineto{\pgfqpoint{2.008301in}{1.496259in}}%
\pgfpathlineto{\pgfqpoint{2.010870in}{1.504623in}}%
\pgfpathlineto{\pgfqpoint{2.009272in}{1.517943in}}%
\pgfpathlineto{\pgfqpoint{2.000578in}{1.522561in}}%
\pgfpathlineto{\pgfqpoint{1.991480in}{1.531907in}}%
\pgfpathlineto{\pgfqpoint{1.977460in}{1.534643in}}%
\pgfpathlineto{\pgfqpoint{1.963816in}{1.534565in}}%
\pgfpathlineto{\pgfqpoint{1.950397in}{1.541198in}}%
\pgfpathlineto{\pgfqpoint{1.905502in}{1.550866in}}%
\pgfpathlineto{\pgfqpoint{1.900596in}{1.561110in}}%
\pgfpathlineto{\pgfqpoint{1.891845in}{1.564606in}}%
\pgfpathlineto{\pgfqpoint{1.908375in}{1.572449in}}%
\pgfpathlineto{\pgfqpoint{1.917713in}{1.582234in}}%
\pgfpathlineto{\pgfqpoint{1.935106in}{1.584888in}}%
\pgfpathlineto{\pgfqpoint{1.945809in}{1.594855in}}%
\pgfpathlineto{\pgfqpoint{1.951425in}{1.595254in}}%
\pgfpathlineto{\pgfqpoint{1.955710in}{1.602362in}}%
\pgfpathlineto{\pgfqpoint{1.966157in}{1.610783in}}%
\pgfpathlineto{\pgfqpoint{1.967067in}{1.604831in}}%
\pgfpathlineto{\pgfqpoint{1.971774in}{1.603602in}}%
\pgfpathlineto{\pgfqpoint{1.975316in}{1.596504in}}%
\pgfpathlineto{\pgfqpoint{1.975954in}{1.584403in}}%
\pgfpathlineto{\pgfqpoint{1.986599in}{1.591634in}}%
\pgfpathlineto{\pgfqpoint{1.999002in}{1.593145in}}%
\pgfpathlineto{\pgfqpoint{2.009590in}{1.589356in}}%
\pgfpathlineto{\pgfqpoint{2.023949in}{1.569658in}}%
\pgfpathlineto{\pgfqpoint{2.039631in}{1.572806in}}%
\pgfpathlineto{\pgfqpoint{2.046003in}{1.567535in}}%
\pgfpathlineto{\pgfqpoint{2.056252in}{1.567064in}}%
\pgfpathlineto{\pgfqpoint{2.063061in}{1.576518in}}%
\pgfpathlineto{\pgfqpoint{2.075982in}{1.584876in}}%
\pgfpathlineto{\pgfqpoint{2.088400in}{1.587493in}}%
\pgfpathlineto{\pgfqpoint{2.103808in}{1.587766in}}%
\pgfpathlineto{\pgfqpoint{2.115068in}{1.594222in}}%
\pgfpathlineto{\pgfqpoint{2.124258in}{1.591248in}}%
\pgfpathlineto{\pgfqpoint{2.126215in}{1.577594in}}%
\pgfpathlineto{\pgfqpoint{2.145930in}{1.575456in}}%
\pgfpathlineto{\pgfqpoint{2.156638in}{1.581847in}}%
\pgfpathlineto{\pgfqpoint{2.164061in}{1.567449in}}%
\pgfpathlineto{\pgfqpoint{2.178225in}{1.550799in}}%
\pgfpathlineto{\pgfqpoint{2.158401in}{1.550895in}}%
\pgfpathlineto{\pgfqpoint{2.152137in}{1.548837in}}%
\pgfpathlineto{\pgfqpoint{2.143551in}{1.551511in}}%
\pgfpathlineto{\pgfqpoint{2.142974in}{1.539985in}}%
\pgfpathlineto{\pgfqpoint{2.127384in}{1.548993in}}%
\pgfpathlineto{\pgfqpoint{2.107335in}{1.551746in}}%
\pgfpathlineto{\pgfqpoint{2.101816in}{1.542973in}}%
\pgfpathlineto{\pgfqpoint{2.075584in}{1.538707in}}%
\pgfpathlineto{\pgfqpoint{2.072566in}{1.531276in}}%
\pgfpathlineto{\pgfqpoint{2.054312in}{1.529046in}}%
\pgfpathlineto{\pgfqpoint{2.048794in}{1.521492in}}%
\pgfpathlineto{\pgfqpoint{2.039139in}{1.519466in}}%
\pgfpathlineto{\pgfqpoint{2.031427in}{1.501552in}}%
\pgfpathlineto{\pgfqpoint{2.021663in}{1.484203in}}%
\pgfpathlineto{\pgfqpoint{2.017832in}{1.482495in}}%
\pgfpathclose%
\pgfusepath{fill}%
\end{pgfscope}%
\begin{pgfscope}%
\pgfpathrectangle{\pgfqpoint{0.100000in}{0.100000in}}{\pgfqpoint{2.862985in}{1.832778in}}%
\pgfusepath{clip}%
\pgfsetbuttcap%
\pgfsetmiterjoin%
\definecolor{currentfill}{rgb}{0.690888,0.862438,0.488351}%
\pgfsetfillcolor{currentfill}%
\pgfsetlinewidth{0.000000pt}%
\definecolor{currentstroke}{rgb}{0.000000,0.000000,0.000000}%
\pgfsetstrokecolor{currentstroke}%
\pgfsetstrokeopacity{0.000000}%
\pgfsetdash{}{0pt}%
\pgfpathmoveto{\pgfqpoint{2.074029in}{1.274203in}}%
\pgfpathlineto{\pgfqpoint{2.083339in}{1.284002in}}%
\pgfpathlineto{\pgfqpoint{2.087580in}{1.298209in}}%
\pgfpathlineto{\pgfqpoint{2.092623in}{1.306443in}}%
\pgfpathlineto{\pgfqpoint{2.095716in}{1.317649in}}%
\pgfpathlineto{\pgfqpoint{2.096678in}{1.339993in}}%
\pgfpathlineto{\pgfqpoint{2.092021in}{1.361407in}}%
\pgfpathlineto{\pgfqpoint{2.076633in}{1.394197in}}%
\pgfpathlineto{\pgfqpoint{2.080800in}{1.404507in}}%
\pgfpathlineto{\pgfqpoint{2.075372in}{1.418758in}}%
\pgfpathlineto{\pgfqpoint{2.084684in}{1.438466in}}%
\pgfpathlineto{\pgfqpoint{2.083166in}{1.460442in}}%
\pgfpathlineto{\pgfqpoint{2.089653in}{1.463207in}}%
\pgfpathlineto{\pgfqpoint{2.090468in}{1.473664in}}%
\pgfpathlineto{\pgfqpoint{2.101982in}{1.480375in}}%
\pgfpathlineto{\pgfqpoint{2.108494in}{1.479291in}}%
\pgfpathlineto{\pgfqpoint{2.110324in}{1.468077in}}%
\pgfpathlineto{\pgfqpoint{2.115407in}{1.467638in}}%
\pgfpathlineto{\pgfqpoint{2.120044in}{1.483813in}}%
\pgfpathlineto{\pgfqpoint{2.118468in}{1.496393in}}%
\pgfpathlineto{\pgfqpoint{2.121483in}{1.503617in}}%
\pgfpathlineto{\pgfqpoint{2.137785in}{1.511080in}}%
\pgfpathlineto{\pgfqpoint{2.130356in}{1.513735in}}%
\pgfpathlineto{\pgfqpoint{2.127936in}{1.520146in}}%
\pgfpathlineto{\pgfqpoint{2.133288in}{1.531404in}}%
\pgfpathlineto{\pgfqpoint{2.143841in}{1.535292in}}%
\pgfpathlineto{\pgfqpoint{2.156095in}{1.528650in}}%
\pgfpathlineto{\pgfqpoint{2.167648in}{1.528567in}}%
\pgfpathlineto{\pgfqpoint{2.173015in}{1.520808in}}%
\pgfpathlineto{\pgfqpoint{2.181106in}{1.521347in}}%
\pgfpathlineto{\pgfqpoint{2.206083in}{1.510561in}}%
\pgfpathlineto{\pgfqpoint{2.211074in}{1.500120in}}%
\pgfpathlineto{\pgfqpoint{2.205624in}{1.496504in}}%
\pgfpathlineto{\pgfqpoint{2.207303in}{1.488679in}}%
\pgfpathlineto{\pgfqpoint{2.212687in}{1.485198in}}%
\pgfpathlineto{\pgfqpoint{2.215688in}{1.475579in}}%
\pgfpathlineto{\pgfqpoint{2.215268in}{1.452144in}}%
\pgfpathlineto{\pgfqpoint{2.208159in}{1.446546in}}%
\pgfpathlineto{\pgfqpoint{2.206564in}{1.434236in}}%
\pgfpathlineto{\pgfqpoint{2.193377in}{1.422859in}}%
\pgfpathlineto{\pgfqpoint{2.194159in}{1.409089in}}%
\pgfpathlineto{\pgfqpoint{2.205717in}{1.404265in}}%
\pgfpathlineto{\pgfqpoint{2.218756in}{1.421446in}}%
\pgfpathlineto{\pgfqpoint{2.219833in}{1.427678in}}%
\pgfpathlineto{\pgfqpoint{2.236094in}{1.438035in}}%
\pgfpathlineto{\pgfqpoint{2.246435in}{1.433238in}}%
\pgfpathlineto{\pgfqpoint{2.252924in}{1.422395in}}%
\pgfpathlineto{\pgfqpoint{2.263393in}{1.384749in}}%
\pgfpathlineto{\pgfqpoint{2.268940in}{1.372839in}}%
\pgfpathlineto{\pgfqpoint{2.267203in}{1.367915in}}%
\pgfpathlineto{\pgfqpoint{2.267372in}{1.351150in}}%
\pgfpathlineto{\pgfqpoint{2.257281in}{1.352759in}}%
\pgfpathlineto{\pgfqpoint{2.251581in}{1.340231in}}%
\pgfpathlineto{\pgfqpoint{2.250791in}{1.331714in}}%
\pgfpathlineto{\pgfqpoint{2.243169in}{1.326246in}}%
\pgfpathlineto{\pgfqpoint{2.241475in}{1.309636in}}%
\pgfpathlineto{\pgfqpoint{2.230390in}{1.288653in}}%
\pgfpathlineto{\pgfqpoint{2.169752in}{1.279623in}}%
\pgfpathlineto{\pgfqpoint{2.169393in}{1.283592in}}%
\pgfpathlineto{\pgfqpoint{2.128831in}{1.279317in}}%
\pgfpathlineto{\pgfqpoint{2.074029in}{1.274203in}}%
\pgfpathclose%
\pgfusepath{fill}%
\end{pgfscope}%
\begin{pgfscope}%
\pgfpathrectangle{\pgfqpoint{0.100000in}{0.100000in}}{\pgfqpoint{2.862985in}{1.832778in}}%
\pgfusepath{clip}%
\pgfsetbuttcap%
\pgfsetmiterjoin%
\definecolor{currentfill}{rgb}{0.859516,0.942330,0.747405}%
\pgfsetfillcolor{currentfill}%
\pgfsetlinewidth{0.000000pt}%
\definecolor{currentstroke}{rgb}{0.000000,0.000000,0.000000}%
\pgfsetstrokecolor{currentstroke}%
\pgfsetstrokeopacity{0.000000}%
\pgfsetdash{}{0pt}%
\pgfpathmoveto{\pgfqpoint{0.542548in}{0.431818in}}%
\pgfpathlineto{\pgfqpoint{0.541337in}{0.430752in}}%
\pgfpathlineto{\pgfqpoint{0.535781in}{0.432173in}}%
\pgfpathlineto{\pgfqpoint{0.536196in}{0.434969in}}%
\pgfpathlineto{\pgfqpoint{0.539341in}{0.436249in}}%
\pgfpathlineto{\pgfqpoint{0.543824in}{0.442694in}}%
\pgfpathlineto{\pgfqpoint{0.544569in}{0.447051in}}%
\pgfpathlineto{\pgfqpoint{0.547090in}{0.448192in}}%
\pgfpathlineto{\pgfqpoint{0.551190in}{0.448197in}}%
\pgfpathlineto{\pgfqpoint{0.555477in}{0.462333in}}%
\pgfpathlineto{\pgfqpoint{0.557157in}{0.463496in}}%
\pgfpathlineto{\pgfqpoint{0.555651in}{0.466260in}}%
\pgfpathlineto{\pgfqpoint{0.552125in}{0.463464in}}%
\pgfpathlineto{\pgfqpoint{0.541766in}{0.467233in}}%
\pgfpathlineto{\pgfqpoint{0.535583in}{0.473582in}}%
\pgfpathlineto{\pgfqpoint{0.538280in}{0.476479in}}%
\pgfpathlineto{\pgfqpoint{0.537053in}{0.480565in}}%
\pgfpathlineto{\pgfqpoint{0.537809in}{0.485392in}}%
\pgfpathlineto{\pgfqpoint{0.536479in}{0.491035in}}%
\pgfpathlineto{\pgfqpoint{0.536266in}{0.496684in}}%
\pgfpathlineto{\pgfqpoint{0.542217in}{0.495739in}}%
\pgfpathlineto{\pgfqpoint{0.543395in}{0.497404in}}%
\pgfpathlineto{\pgfqpoint{0.546355in}{0.497243in}}%
\pgfpathlineto{\pgfqpoint{0.549764in}{0.491217in}}%
\pgfpathlineto{\pgfqpoint{0.553051in}{0.486806in}}%
\pgfpathlineto{\pgfqpoint{0.550418in}{0.483935in}}%
\pgfpathlineto{\pgfqpoint{0.554970in}{0.482833in}}%
\pgfpathlineto{\pgfqpoint{0.556007in}{0.484241in}}%
\pgfpathlineto{\pgfqpoint{0.553573in}{0.489956in}}%
\pgfpathlineto{\pgfqpoint{0.548447in}{0.494692in}}%
\pgfpathlineto{\pgfqpoint{0.549074in}{0.496676in}}%
\pgfpathlineto{\pgfqpoint{0.544984in}{0.501691in}}%
\pgfpathlineto{\pgfqpoint{0.549161in}{0.502710in}}%
\pgfpathlineto{\pgfqpoint{0.545542in}{0.513448in}}%
\pgfpathlineto{\pgfqpoint{0.548484in}{0.514936in}}%
\pgfpathlineto{\pgfqpoint{0.549202in}{0.517511in}}%
\pgfpathlineto{\pgfqpoint{0.547288in}{0.520569in}}%
\pgfpathlineto{\pgfqpoint{0.551239in}{0.521335in}}%
\pgfpathlineto{\pgfqpoint{0.551948in}{0.523648in}}%
\pgfpathlineto{\pgfqpoint{0.556660in}{0.520176in}}%
\pgfpathlineto{\pgfqpoint{0.557707in}{0.523941in}}%
\pgfpathlineto{\pgfqpoint{0.560183in}{0.525597in}}%
\pgfpathlineto{\pgfqpoint{0.572233in}{0.527901in}}%
\pgfpathlineto{\pgfqpoint{0.575932in}{0.527086in}}%
\pgfpathlineto{\pgfqpoint{0.579031in}{0.531659in}}%
\pgfpathlineto{\pgfqpoint{0.586384in}{0.534658in}}%
\pgfpathlineto{\pgfqpoint{0.592047in}{0.533521in}}%
\pgfpathlineto{\pgfqpoint{0.594631in}{0.529538in}}%
\pgfpathlineto{\pgfqpoint{0.594495in}{0.523212in}}%
\pgfpathlineto{\pgfqpoint{0.597394in}{0.521778in}}%
\pgfpathlineto{\pgfqpoint{0.603306in}{0.521978in}}%
\pgfpathlineto{\pgfqpoint{0.610878in}{0.523624in}}%
\pgfpathlineto{\pgfqpoint{0.610785in}{0.521601in}}%
\pgfpathlineto{\pgfqpoint{0.620357in}{0.515832in}}%
\pgfpathlineto{\pgfqpoint{0.627052in}{0.517379in}}%
\pgfpathlineto{\pgfqpoint{0.629022in}{0.519381in}}%
\pgfpathlineto{\pgfqpoint{0.631133in}{0.524508in}}%
\pgfpathlineto{\pgfqpoint{0.633750in}{0.527792in}}%
\pgfpathlineto{\pgfqpoint{0.633984in}{0.532712in}}%
\pgfpathlineto{\pgfqpoint{0.631587in}{0.534647in}}%
\pgfpathlineto{\pgfqpoint{0.635136in}{0.536396in}}%
\pgfpathlineto{\pgfqpoint{0.641098in}{0.533768in}}%
\pgfpathlineto{\pgfqpoint{0.643040in}{0.535390in}}%
\pgfpathlineto{\pgfqpoint{0.643382in}{0.542352in}}%
\pgfpathlineto{\pgfqpoint{0.639528in}{0.540897in}}%
\pgfpathlineto{\pgfqpoint{0.636855in}{0.543309in}}%
\pgfpathlineto{\pgfqpoint{0.632516in}{0.544812in}}%
\pgfpathlineto{\pgfqpoint{0.625629in}{0.544561in}}%
\pgfpathlineto{\pgfqpoint{0.619030in}{0.546848in}}%
\pgfpathlineto{\pgfqpoint{0.618699in}{0.552545in}}%
\pgfpathlineto{\pgfqpoint{0.612797in}{0.558597in}}%
\pgfpathlineto{\pgfqpoint{0.605313in}{0.561138in}}%
\pgfpathlineto{\pgfqpoint{0.598780in}{0.572853in}}%
\pgfpathlineto{\pgfqpoint{0.599316in}{0.577189in}}%
\pgfpathlineto{\pgfqpoint{0.602833in}{0.579165in}}%
\pgfpathlineto{\pgfqpoint{0.602315in}{0.583269in}}%
\pgfpathlineto{\pgfqpoint{0.603132in}{0.587352in}}%
\pgfpathlineto{\pgfqpoint{0.605915in}{0.584557in}}%
\pgfpathlineto{\pgfqpoint{0.609590in}{0.585723in}}%
\pgfpathlineto{\pgfqpoint{0.609407in}{0.589270in}}%
\pgfpathlineto{\pgfqpoint{0.604311in}{0.596390in}}%
\pgfpathlineto{\pgfqpoint{0.602820in}{0.603734in}}%
\pgfpathlineto{\pgfqpoint{0.607085in}{0.604294in}}%
\pgfpathlineto{\pgfqpoint{0.613468in}{0.603843in}}%
\pgfpathlineto{\pgfqpoint{0.615478in}{0.601387in}}%
\pgfpathlineto{\pgfqpoint{0.621201in}{0.602266in}}%
\pgfpathlineto{\pgfqpoint{0.626520in}{0.601112in}}%
\pgfpathlineto{\pgfqpoint{0.630115in}{0.599083in}}%
\pgfpathlineto{\pgfqpoint{0.632604in}{0.600010in}}%
\pgfpathlineto{\pgfqpoint{0.639212in}{0.597983in}}%
\pgfpathlineto{\pgfqpoint{0.639273in}{0.596187in}}%
\pgfpathlineto{\pgfqpoint{0.644305in}{0.596794in}}%
\pgfpathlineto{\pgfqpoint{0.651053in}{0.592414in}}%
\pgfpathlineto{\pgfqpoint{0.650791in}{0.589506in}}%
\pgfpathlineto{\pgfqpoint{0.647580in}{0.588164in}}%
\pgfpathlineto{\pgfqpoint{0.644167in}{0.584972in}}%
\pgfpathlineto{\pgfqpoint{0.643755in}{0.580588in}}%
\pgfpathlineto{\pgfqpoint{0.651106in}{0.574997in}}%
\pgfpathlineto{\pgfqpoint{0.649930in}{0.572318in}}%
\pgfpathlineto{\pgfqpoint{0.655161in}{0.570260in}}%
\pgfpathlineto{\pgfqpoint{0.656445in}{0.566900in}}%
\pgfpathlineto{\pgfqpoint{0.663083in}{0.569570in}}%
\pgfpathlineto{\pgfqpoint{0.666007in}{0.566071in}}%
\pgfpathlineto{\pgfqpoint{0.667356in}{0.568189in}}%
\pgfpathlineto{\pgfqpoint{0.663383in}{0.575083in}}%
\pgfpathlineto{\pgfqpoint{0.664795in}{0.577148in}}%
\pgfpathlineto{\pgfqpoint{0.665421in}{0.582594in}}%
\pgfpathlineto{\pgfqpoint{0.663838in}{0.584991in}}%
\pgfpathlineto{\pgfqpoint{0.665025in}{0.588211in}}%
\pgfpathlineto{\pgfqpoint{0.668915in}{0.587908in}}%
\pgfpathlineto{\pgfqpoint{0.668110in}{0.582429in}}%
\pgfpathlineto{\pgfqpoint{0.665627in}{0.580602in}}%
\pgfpathlineto{\pgfqpoint{0.665803in}{0.573231in}}%
\pgfpathlineto{\pgfqpoint{0.670258in}{0.572350in}}%
\pgfpathlineto{\pgfqpoint{0.670312in}{0.566870in}}%
\pgfpathlineto{\pgfqpoint{0.675178in}{0.563315in}}%
\pgfpathlineto{\pgfqpoint{0.677415in}{0.565557in}}%
\pgfpathlineto{\pgfqpoint{0.674970in}{0.572380in}}%
\pgfpathlineto{\pgfqpoint{0.672442in}{0.574106in}}%
\pgfpathlineto{\pgfqpoint{0.669955in}{0.572970in}}%
\pgfpathlineto{\pgfqpoint{0.667767in}{0.574420in}}%
\pgfpathlineto{\pgfqpoint{0.668149in}{0.580725in}}%
\pgfpathlineto{\pgfqpoint{0.671800in}{0.583304in}}%
\pgfpathlineto{\pgfqpoint{0.675395in}{0.583066in}}%
\pgfpathlineto{\pgfqpoint{0.673905in}{0.586634in}}%
\pgfpathlineto{\pgfqpoint{0.668518in}{0.589366in}}%
\pgfpathlineto{\pgfqpoint{0.661383in}{0.600253in}}%
\pgfpathlineto{\pgfqpoint{0.664944in}{0.605475in}}%
\pgfpathlineto{\pgfqpoint{0.666926in}{0.610393in}}%
\pgfpathlineto{\pgfqpoint{0.667082in}{0.620269in}}%
\pgfpathlineto{\pgfqpoint{0.666319in}{0.628601in}}%
\pgfpathlineto{\pgfqpoint{0.664142in}{0.634072in}}%
\pgfpathlineto{\pgfqpoint{0.664576in}{0.639541in}}%
\pgfpathlineto{\pgfqpoint{0.670144in}{0.644285in}}%
\pgfpathlineto{\pgfqpoint{0.675681in}{0.649922in}}%
\pgfpathlineto{\pgfqpoint{0.689564in}{0.638405in}}%
\pgfpathlineto{\pgfqpoint{0.698640in}{0.637556in}}%
\pgfpathlineto{\pgfqpoint{0.706132in}{0.640424in}}%
\pgfpathlineto{\pgfqpoint{0.712708in}{0.644201in}}%
\pgfpathlineto{\pgfqpoint{0.714105in}{0.646000in}}%
\pgfpathlineto{\pgfqpoint{0.730340in}{0.650103in}}%
\pgfpathlineto{\pgfqpoint{0.731749in}{0.647098in}}%
\pgfpathlineto{\pgfqpoint{0.737266in}{0.644638in}}%
\pgfpathlineto{\pgfqpoint{0.747793in}{0.645392in}}%
\pgfpathlineto{\pgfqpoint{0.749716in}{0.644720in}}%
\pgfpathlineto{\pgfqpoint{0.754448in}{0.646515in}}%
\pgfpathlineto{\pgfqpoint{0.756350in}{0.643411in}}%
\pgfpathlineto{\pgfqpoint{0.760894in}{0.640575in}}%
\pgfpathlineto{\pgfqpoint{0.763167in}{0.640896in}}%
\pgfpathlineto{\pgfqpoint{0.767566in}{0.637488in}}%
\pgfpathlineto{\pgfqpoint{0.775000in}{0.638115in}}%
\pgfpathlineto{\pgfqpoint{0.782425in}{0.640874in}}%
\pgfpathlineto{\pgfqpoint{0.788489in}{0.632051in}}%
\pgfpathlineto{\pgfqpoint{0.788019in}{0.630071in}}%
\pgfpathlineto{\pgfqpoint{0.781328in}{0.627796in}}%
\pgfpathlineto{\pgfqpoint{0.783115in}{0.624355in}}%
\pgfpathlineto{\pgfqpoint{0.789275in}{0.627976in}}%
\pgfpathlineto{\pgfqpoint{0.791790in}{0.627100in}}%
\pgfpathlineto{\pgfqpoint{0.792860in}{0.622768in}}%
\pgfpathlineto{\pgfqpoint{0.789183in}{0.620936in}}%
\pgfpathlineto{\pgfqpoint{0.791809in}{0.615461in}}%
\pgfpathlineto{\pgfqpoint{0.795623in}{0.616386in}}%
\pgfpathlineto{\pgfqpoint{0.801157in}{0.613560in}}%
\pgfpathlineto{\pgfqpoint{0.806450in}{0.605718in}}%
\pgfpathlineto{\pgfqpoint{0.801793in}{0.603540in}}%
\pgfpathlineto{\pgfqpoint{0.805907in}{0.597761in}}%
\pgfpathlineto{\pgfqpoint{0.802441in}{0.596410in}}%
\pgfpathlineto{\pgfqpoint{0.806753in}{0.590922in}}%
\pgfpathlineto{\pgfqpoint{0.812012in}{0.591506in}}%
\pgfpathlineto{\pgfqpoint{0.814781in}{0.589037in}}%
\pgfpathlineto{\pgfqpoint{0.821785in}{0.585396in}}%
\pgfpathlineto{\pgfqpoint{0.830820in}{0.572470in}}%
\pgfpathlineto{\pgfqpoint{0.830616in}{0.569379in}}%
\pgfpathlineto{\pgfqpoint{0.844166in}{0.559477in}}%
\pgfpathlineto{\pgfqpoint{0.844458in}{0.555975in}}%
\pgfpathlineto{\pgfqpoint{0.848069in}{0.550735in}}%
\pgfpathlineto{\pgfqpoint{0.859163in}{0.548115in}}%
\pgfpathlineto{\pgfqpoint{0.863277in}{0.546252in}}%
\pgfpathlineto{\pgfqpoint{0.867148in}{0.540476in}}%
\pgfpathlineto{\pgfqpoint{0.867646in}{0.535415in}}%
\pgfpathlineto{\pgfqpoint{0.870973in}{0.531291in}}%
\pgfpathlineto{\pgfqpoint{0.871972in}{0.526640in}}%
\pgfpathlineto{\pgfqpoint{0.875090in}{0.524941in}}%
\pgfpathlineto{\pgfqpoint{0.859646in}{0.497067in}}%
\pgfpathlineto{\pgfqpoint{0.827590in}{0.439200in}}%
\pgfpathlineto{\pgfqpoint{0.780557in}{0.354280in}}%
\pgfpathlineto{\pgfqpoint{0.762143in}{0.321025in}}%
\pgfpathlineto{\pgfqpoint{0.766065in}{0.316559in}}%
\pgfpathlineto{\pgfqpoint{0.767793in}{0.317984in}}%
\pgfpathlineto{\pgfqpoint{0.771296in}{0.312785in}}%
\pgfpathlineto{\pgfqpoint{0.776159in}{0.314293in}}%
\pgfpathlineto{\pgfqpoint{0.782614in}{0.311178in}}%
\pgfpathlineto{\pgfqpoint{0.778433in}{0.306310in}}%
\pgfpathlineto{\pgfqpoint{0.781662in}{0.299741in}}%
\pgfpathlineto{\pgfqpoint{0.780995in}{0.296543in}}%
\pgfpathlineto{\pgfqpoint{0.786145in}{0.279677in}}%
\pgfpathlineto{\pgfqpoint{0.783966in}{0.272007in}}%
\pgfpathlineto{\pgfqpoint{0.793805in}{0.273906in}}%
\pgfpathlineto{\pgfqpoint{0.796221in}{0.272662in}}%
\pgfpathlineto{\pgfqpoint{0.798844in}{0.274685in}}%
\pgfpathlineto{\pgfqpoint{0.800709in}{0.278509in}}%
\pgfpathlineto{\pgfqpoint{0.803381in}{0.280701in}}%
\pgfpathlineto{\pgfqpoint{0.814784in}{0.280424in}}%
\pgfpathlineto{\pgfqpoint{0.817361in}{0.273064in}}%
\pgfpathlineto{\pgfqpoint{0.815149in}{0.266603in}}%
\pgfpathlineto{\pgfqpoint{0.817662in}{0.264471in}}%
\pgfpathlineto{\pgfqpoint{0.818983in}{0.260825in}}%
\pgfpathlineto{\pgfqpoint{0.818695in}{0.253967in}}%
\pgfpathlineto{\pgfqpoint{0.822048in}{0.249051in}}%
\pgfpathlineto{\pgfqpoint{0.823886in}{0.241346in}}%
\pgfpathlineto{\pgfqpoint{0.822942in}{0.237184in}}%
\pgfpathlineto{\pgfqpoint{0.824080in}{0.232966in}}%
\pgfpathlineto{\pgfqpoint{0.825722in}{0.208539in}}%
\pgfpathlineto{\pgfqpoint{0.823370in}{0.206615in}}%
\pgfpathlineto{\pgfqpoint{0.826472in}{0.203957in}}%
\pgfpathlineto{\pgfqpoint{0.824034in}{0.200723in}}%
\pgfpathlineto{\pgfqpoint{0.826233in}{0.198087in}}%
\pgfpathlineto{\pgfqpoint{0.824817in}{0.193586in}}%
\pgfpathlineto{\pgfqpoint{0.827984in}{0.192366in}}%
\pgfpathlineto{\pgfqpoint{0.832058in}{0.185522in}}%
\pgfpathlineto{\pgfqpoint{0.835085in}{0.183460in}}%
\pgfpathlineto{\pgfqpoint{0.835877in}{0.180445in}}%
\pgfpathlineto{\pgfqpoint{0.837631in}{0.179139in}}%
\pgfpathlineto{\pgfqpoint{0.837265in}{0.176636in}}%
\pgfpathlineto{\pgfqpoint{0.841046in}{0.174822in}}%
\pgfpathlineto{\pgfqpoint{0.840112in}{0.170006in}}%
\pgfpathlineto{\pgfqpoint{0.836858in}{0.167483in}}%
\pgfpathlineto{\pgfqpoint{0.835380in}{0.162795in}}%
\pgfpathlineto{\pgfqpoint{0.834988in}{0.156518in}}%
\pgfpathlineto{\pgfqpoint{0.833180in}{0.155917in}}%
\pgfpathlineto{\pgfqpoint{0.827271in}{0.150415in}}%
\pgfpathlineto{\pgfqpoint{0.822019in}{0.148962in}}%
\pgfpathlineto{\pgfqpoint{0.819513in}{0.151181in}}%
\pgfpathlineto{\pgfqpoint{0.821935in}{0.157071in}}%
\pgfpathlineto{\pgfqpoint{0.821195in}{0.159962in}}%
\pgfpathlineto{\pgfqpoint{0.824665in}{0.161426in}}%
\pgfpathlineto{\pgfqpoint{0.827909in}{0.169947in}}%
\pgfpathlineto{\pgfqpoint{0.826539in}{0.177225in}}%
\pgfpathlineto{\pgfqpoint{0.824346in}{0.178943in}}%
\pgfpathlineto{\pgfqpoint{0.817012in}{0.178387in}}%
\pgfpathlineto{\pgfqpoint{0.818490in}{0.176486in}}%
\pgfpathlineto{\pgfqpoint{0.813181in}{0.171028in}}%
\pgfpathlineto{\pgfqpoint{0.811471in}{0.174063in}}%
\pgfpathlineto{\pgfqpoint{0.811916in}{0.178027in}}%
\pgfpathlineto{\pgfqpoint{0.814963in}{0.177954in}}%
\pgfpathlineto{\pgfqpoint{0.817677in}{0.180866in}}%
\pgfpathlineto{\pgfqpoint{0.819282in}{0.185118in}}%
\pgfpathlineto{\pgfqpoint{0.825331in}{0.183918in}}%
\pgfpathlineto{\pgfqpoint{0.820488in}{0.186190in}}%
\pgfpathlineto{\pgfqpoint{0.820915in}{0.189233in}}%
\pgfpathlineto{\pgfqpoint{0.818899in}{0.190607in}}%
\pgfpathlineto{\pgfqpoint{0.818074in}{0.194752in}}%
\pgfpathlineto{\pgfqpoint{0.819505in}{0.196911in}}%
\pgfpathlineto{\pgfqpoint{0.816093in}{0.203617in}}%
\pgfpathlineto{\pgfqpoint{0.816381in}{0.207883in}}%
\pgfpathlineto{\pgfqpoint{0.812383in}{0.211682in}}%
\pgfpathlineto{\pgfqpoint{0.810220in}{0.214756in}}%
\pgfpathlineto{\pgfqpoint{0.813207in}{0.217685in}}%
\pgfpathlineto{\pgfqpoint{0.814313in}{0.222063in}}%
\pgfpathlineto{\pgfqpoint{0.813326in}{0.224937in}}%
\pgfpathlineto{\pgfqpoint{0.814481in}{0.228366in}}%
\pgfpathlineto{\pgfqpoint{0.813009in}{0.239543in}}%
\pgfpathlineto{\pgfqpoint{0.810807in}{0.244614in}}%
\pgfpathlineto{\pgfqpoint{0.808288in}{0.246519in}}%
\pgfpathlineto{\pgfqpoint{0.809360in}{0.253116in}}%
\pgfpathlineto{\pgfqpoint{0.808604in}{0.258268in}}%
\pgfpathlineto{\pgfqpoint{0.809974in}{0.261572in}}%
\pgfpathlineto{\pgfqpoint{0.807140in}{0.262035in}}%
\pgfpathlineto{\pgfqpoint{0.806380in}{0.254009in}}%
\pgfpathlineto{\pgfqpoint{0.803196in}{0.245055in}}%
\pgfpathlineto{\pgfqpoint{0.800639in}{0.247094in}}%
\pgfpathlineto{\pgfqpoint{0.800297in}{0.250037in}}%
\pgfpathlineto{\pgfqpoint{0.797168in}{0.254057in}}%
\pgfpathlineto{\pgfqpoint{0.798758in}{0.256688in}}%
\pgfpathlineto{\pgfqpoint{0.798264in}{0.262613in}}%
\pgfpathlineto{\pgfqpoint{0.794215in}{0.260530in}}%
\pgfpathlineto{\pgfqpoint{0.795009in}{0.257316in}}%
\pgfpathlineto{\pgfqpoint{0.794247in}{0.253236in}}%
\pgfpathlineto{\pgfqpoint{0.789702in}{0.253244in}}%
\pgfpathlineto{\pgfqpoint{0.787665in}{0.255773in}}%
\pgfpathlineto{\pgfqpoint{0.784590in}{0.256255in}}%
\pgfpathlineto{\pgfqpoint{0.782521in}{0.258978in}}%
\pgfpathlineto{\pgfqpoint{0.778961in}{0.266454in}}%
\pgfpathlineto{\pgfqpoint{0.777807in}{0.272259in}}%
\pgfpathlineto{\pgfqpoint{0.778546in}{0.274726in}}%
\pgfpathlineto{\pgfqpoint{0.776822in}{0.280160in}}%
\pgfpathlineto{\pgfqpoint{0.773967in}{0.283305in}}%
\pgfpathlineto{\pgfqpoint{0.769029in}{0.291344in}}%
\pgfpathlineto{\pgfqpoint{0.765447in}{0.298326in}}%
\pgfpathlineto{\pgfqpoint{0.769102in}{0.298431in}}%
\pgfpathlineto{\pgfqpoint{0.771252in}{0.299897in}}%
\pgfpathlineto{\pgfqpoint{0.771675in}{0.304435in}}%
\pgfpathlineto{\pgfqpoint{0.761334in}{0.304983in}}%
\pgfpathlineto{\pgfqpoint{0.752539in}{0.314807in}}%
\pgfpathlineto{\pgfqpoint{0.754852in}{0.317392in}}%
\pgfpathlineto{\pgfqpoint{0.750599in}{0.317833in}}%
\pgfpathlineto{\pgfqpoint{0.742152in}{0.326967in}}%
\pgfpathlineto{\pgfqpoint{0.728770in}{0.331547in}}%
\pgfpathlineto{\pgfqpoint{0.726968in}{0.336915in}}%
\pgfpathlineto{\pgfqpoint{0.724191in}{0.339711in}}%
\pgfpathlineto{\pgfqpoint{0.724065in}{0.342336in}}%
\pgfpathlineto{\pgfqpoint{0.721976in}{0.344286in}}%
\pgfpathlineto{\pgfqpoint{0.724717in}{0.347544in}}%
\pgfpathlineto{\pgfqpoint{0.718465in}{0.347858in}}%
\pgfpathlineto{\pgfqpoint{0.716876in}{0.352075in}}%
\pgfpathlineto{\pgfqpoint{0.714030in}{0.353870in}}%
\pgfpathlineto{\pgfqpoint{0.719729in}{0.356037in}}%
\pgfpathlineto{\pgfqpoint{0.716765in}{0.359902in}}%
\pgfpathlineto{\pgfqpoint{0.711311in}{0.361524in}}%
\pgfpathlineto{\pgfqpoint{0.712160in}{0.368663in}}%
\pgfpathlineto{\pgfqpoint{0.710596in}{0.371190in}}%
\pgfpathlineto{\pgfqpoint{0.707780in}{0.372970in}}%
\pgfpathlineto{\pgfqpoint{0.705372in}{0.371185in}}%
\pgfpathlineto{\pgfqpoint{0.703950in}{0.372864in}}%
\pgfpathlineto{\pgfqpoint{0.699800in}{0.373345in}}%
\pgfpathlineto{\pgfqpoint{0.700134in}{0.379823in}}%
\pgfpathlineto{\pgfqpoint{0.694119in}{0.375448in}}%
\pgfpathlineto{\pgfqpoint{0.694246in}{0.366430in}}%
\pgfpathlineto{\pgfqpoint{0.687788in}{0.364702in}}%
\pgfpathlineto{\pgfqpoint{0.683572in}{0.361101in}}%
\pgfpathlineto{\pgfqpoint{0.680182in}{0.360342in}}%
\pgfpathlineto{\pgfqpoint{0.676514in}{0.363971in}}%
\pgfpathlineto{\pgfqpoint{0.673248in}{0.363490in}}%
\pgfpathlineto{\pgfqpoint{0.673854in}{0.366805in}}%
\pgfpathlineto{\pgfqpoint{0.669211in}{0.365020in}}%
\pgfpathlineto{\pgfqpoint{0.664909in}{0.364820in}}%
\pgfpathlineto{\pgfqpoint{0.659874in}{0.363010in}}%
\pgfpathlineto{\pgfqpoint{0.649656in}{0.364028in}}%
\pgfpathlineto{\pgfqpoint{0.646376in}{0.363358in}}%
\pgfpathlineto{\pgfqpoint{0.644009in}{0.366339in}}%
\pgfpathlineto{\pgfqpoint{0.639640in}{0.366637in}}%
\pgfpathlineto{\pgfqpoint{0.638731in}{0.370500in}}%
\pgfpathlineto{\pgfqpoint{0.641335in}{0.372604in}}%
\pgfpathlineto{\pgfqpoint{0.644390in}{0.372406in}}%
\pgfpathlineto{\pgfqpoint{0.646940in}{0.369537in}}%
\pgfpathlineto{\pgfqpoint{0.648619in}{0.373944in}}%
\pgfpathlineto{\pgfqpoint{0.646276in}{0.378827in}}%
\pgfpathlineto{\pgfqpoint{0.651579in}{0.383125in}}%
\pgfpathlineto{\pgfqpoint{0.656893in}{0.384709in}}%
\pgfpathlineto{\pgfqpoint{0.660593in}{0.387292in}}%
\pgfpathlineto{\pgfqpoint{0.662990in}{0.390076in}}%
\pgfpathlineto{\pgfqpoint{0.664327in}{0.394638in}}%
\pgfpathlineto{\pgfqpoint{0.668552in}{0.393396in}}%
\pgfpathlineto{\pgfqpoint{0.677714in}{0.394034in}}%
\pgfpathlineto{\pgfqpoint{0.679680in}{0.389032in}}%
\pgfpathlineto{\pgfqpoint{0.682894in}{0.388814in}}%
\pgfpathlineto{\pgfqpoint{0.688742in}{0.381156in}}%
\pgfpathlineto{\pgfqpoint{0.688753in}{0.383947in}}%
\pgfpathlineto{\pgfqpoint{0.686710in}{0.384855in}}%
\pgfpathlineto{\pgfqpoint{0.682971in}{0.393996in}}%
\pgfpathlineto{\pgfqpoint{0.684790in}{0.395238in}}%
\pgfpathlineto{\pgfqpoint{0.679797in}{0.400007in}}%
\pgfpathlineto{\pgfqpoint{0.675334in}{0.400936in}}%
\pgfpathlineto{\pgfqpoint{0.671081in}{0.399287in}}%
\pgfpathlineto{\pgfqpoint{0.663798in}{0.400571in}}%
\pgfpathlineto{\pgfqpoint{0.660344in}{0.398261in}}%
\pgfpathlineto{\pgfqpoint{0.652478in}{0.396476in}}%
\pgfpathlineto{\pgfqpoint{0.651759in}{0.393907in}}%
\pgfpathlineto{\pgfqpoint{0.645716in}{0.393207in}}%
\pgfpathlineto{\pgfqpoint{0.644165in}{0.390270in}}%
\pgfpathlineto{\pgfqpoint{0.640709in}{0.388119in}}%
\pgfpathlineto{\pgfqpoint{0.636533in}{0.388489in}}%
\pgfpathlineto{\pgfqpoint{0.634448in}{0.386408in}}%
\pgfpathlineto{\pgfqpoint{0.631818in}{0.386506in}}%
\pgfpathlineto{\pgfqpoint{0.629304in}{0.388378in}}%
\pgfpathlineto{\pgfqpoint{0.624909in}{0.388044in}}%
\pgfpathlineto{\pgfqpoint{0.623874in}{0.386318in}}%
\pgfpathlineto{\pgfqpoint{0.619258in}{0.388090in}}%
\pgfpathlineto{\pgfqpoint{0.615339in}{0.383518in}}%
\pgfpathlineto{\pgfqpoint{0.619097in}{0.379095in}}%
\pgfpathlineto{\pgfqpoint{0.620543in}{0.375457in}}%
\pgfpathlineto{\pgfqpoint{0.619210in}{0.372376in}}%
\pgfpathlineto{\pgfqpoint{0.613870in}{0.370250in}}%
\pgfpathlineto{\pgfqpoint{0.610776in}{0.371980in}}%
\pgfpathlineto{\pgfqpoint{0.606882in}{0.371043in}}%
\pgfpathlineto{\pgfqpoint{0.606413in}{0.368343in}}%
\pgfpathlineto{\pgfqpoint{0.601673in}{0.369298in}}%
\pgfpathlineto{\pgfqpoint{0.602722in}{0.366600in}}%
\pgfpathlineto{\pgfqpoint{0.595731in}{0.366031in}}%
\pgfpathlineto{\pgfqpoint{0.591205in}{0.369456in}}%
\pgfpathlineto{\pgfqpoint{0.588787in}{0.367278in}}%
\pgfpathlineto{\pgfqpoint{0.582354in}{0.369597in}}%
\pgfpathlineto{\pgfqpoint{0.580752in}{0.367248in}}%
\pgfpathlineto{\pgfqpoint{0.577563in}{0.366203in}}%
\pgfpathlineto{\pgfqpoint{0.575005in}{0.368957in}}%
\pgfpathlineto{\pgfqpoint{0.566257in}{0.368272in}}%
\pgfpathlineto{\pgfqpoint{0.566266in}{0.364211in}}%
\pgfpathlineto{\pgfqpoint{0.562752in}{0.363136in}}%
\pgfpathlineto{\pgfqpoint{0.560018in}{0.365224in}}%
\pgfpathlineto{\pgfqpoint{0.556404in}{0.364064in}}%
\pgfpathlineto{\pgfqpoint{0.551727in}{0.363787in}}%
\pgfpathlineto{\pgfqpoint{0.544759in}{0.366724in}}%
\pgfpathlineto{\pgfqpoint{0.540769in}{0.363696in}}%
\pgfpathlineto{\pgfqpoint{0.535063in}{0.367743in}}%
\pgfpathlineto{\pgfqpoint{0.533197in}{0.366209in}}%
\pgfpathlineto{\pgfqpoint{0.531887in}{0.362246in}}%
\pgfpathlineto{\pgfqpoint{0.528218in}{0.359922in}}%
\pgfpathlineto{\pgfqpoint{0.522236in}{0.362582in}}%
\pgfpathlineto{\pgfqpoint{0.512326in}{0.369153in}}%
\pgfpathlineto{\pgfqpoint{0.509438in}{0.369737in}}%
\pgfpathlineto{\pgfqpoint{0.507808in}{0.368368in}}%
\pgfpathlineto{\pgfqpoint{0.503449in}{0.369299in}}%
\pgfpathlineto{\pgfqpoint{0.501322in}{0.371344in}}%
\pgfpathlineto{\pgfqpoint{0.496969in}{0.372446in}}%
\pgfpathlineto{\pgfqpoint{0.491023in}{0.372510in}}%
\pgfpathlineto{\pgfqpoint{0.488662in}{0.374950in}}%
\pgfpathlineto{\pgfqpoint{0.493187in}{0.377474in}}%
\pgfpathlineto{\pgfqpoint{0.492131in}{0.379966in}}%
\pgfpathlineto{\pgfqpoint{0.489022in}{0.379406in}}%
\pgfpathlineto{\pgfqpoint{0.487470in}{0.377447in}}%
\pgfpathlineto{\pgfqpoint{0.480192in}{0.374639in}}%
\pgfpathlineto{\pgfqpoint{0.476512in}{0.375594in}}%
\pgfpathlineto{\pgfqpoint{0.474964in}{0.381297in}}%
\pgfpathlineto{\pgfqpoint{0.470373in}{0.378845in}}%
\pgfpathlineto{\pgfqpoint{0.469039in}{0.386073in}}%
\pgfpathlineto{\pgfqpoint{0.466864in}{0.384210in}}%
\pgfpathlineto{\pgfqpoint{0.467228in}{0.381429in}}%
\pgfpathlineto{\pgfqpoint{0.462226in}{0.382163in}}%
\pgfpathlineto{\pgfqpoint{0.467429in}{0.387010in}}%
\pgfpathlineto{\pgfqpoint{0.472702in}{0.384130in}}%
\pgfpathlineto{\pgfqpoint{0.473833in}{0.385576in}}%
\pgfpathlineto{\pgfqpoint{0.479380in}{0.383843in}}%
\pgfpathlineto{\pgfqpoint{0.479033in}{0.385816in}}%
\pgfpathlineto{\pgfqpoint{0.494168in}{0.386469in}}%
\pgfpathlineto{\pgfqpoint{0.501551in}{0.383391in}}%
\pgfpathlineto{\pgfqpoint{0.504408in}{0.380355in}}%
\pgfpathlineto{\pgfqpoint{0.503561in}{0.377705in}}%
\pgfpathlineto{\pgfqpoint{0.506429in}{0.376702in}}%
\pgfpathlineto{\pgfqpoint{0.513558in}{0.380885in}}%
\pgfpathlineto{\pgfqpoint{0.522779in}{0.381153in}}%
\pgfpathlineto{\pgfqpoint{0.531155in}{0.378631in}}%
\pgfpathlineto{\pgfqpoint{0.535899in}{0.379054in}}%
\pgfpathlineto{\pgfqpoint{0.537665in}{0.375142in}}%
\pgfpathlineto{\pgfqpoint{0.540318in}{0.379455in}}%
\pgfpathlineto{\pgfqpoint{0.541950in}{0.380389in}}%
\pgfpathlineto{\pgfqpoint{0.549498in}{0.382264in}}%
\pgfpathlineto{\pgfqpoint{0.552304in}{0.381118in}}%
\pgfpathlineto{\pgfqpoint{0.555338in}{0.382479in}}%
\pgfpathlineto{\pgfqpoint{0.558958in}{0.381767in}}%
\pgfpathlineto{\pgfqpoint{0.559812in}{0.383352in}}%
\pgfpathlineto{\pgfqpoint{0.567363in}{0.390625in}}%
\pgfpathlineto{\pgfqpoint{0.570561in}{0.391703in}}%
\pgfpathlineto{\pgfqpoint{0.572451in}{0.395629in}}%
\pgfpathlineto{\pgfqpoint{0.582481in}{0.398668in}}%
\pgfpathlineto{\pgfqpoint{0.583737in}{0.400987in}}%
\pgfpathlineto{\pgfqpoint{0.569522in}{0.404566in}}%
\pgfpathlineto{\pgfqpoint{0.570081in}{0.408415in}}%
\pgfpathlineto{\pgfqpoint{0.568619in}{0.413499in}}%
\pgfpathlineto{\pgfqpoint{0.565309in}{0.413077in}}%
\pgfpathlineto{\pgfqpoint{0.563057in}{0.406479in}}%
\pgfpathlineto{\pgfqpoint{0.560443in}{0.405936in}}%
\pgfpathlineto{\pgfqpoint{0.558742in}{0.408071in}}%
\pgfpathlineto{\pgfqpoint{0.560726in}{0.417413in}}%
\pgfpathlineto{\pgfqpoint{0.558750in}{0.421328in}}%
\pgfpathlineto{\pgfqpoint{0.554922in}{0.425943in}}%
\pgfpathlineto{\pgfqpoint{0.556762in}{0.429465in}}%
\pgfpathlineto{\pgfqpoint{0.548683in}{0.429591in}}%
\pgfpathlineto{\pgfqpoint{0.548867in}{0.430958in}}%
\pgfpathlineto{\pgfqpoint{0.544046in}{0.432871in}}%
\pgfpathlineto{\pgfqpoint{0.542548in}{0.431818in}}%
\pgfpathclose%
\pgfusepath{fill}%
\end{pgfscope}%
\begin{pgfscope}%
\pgfpathrectangle{\pgfqpoint{0.100000in}{0.100000in}}{\pgfqpoint{2.862985in}{1.832778in}}%
\pgfusepath{clip}%
\pgfsetbuttcap%
\pgfsetmiterjoin%
\definecolor{currentfill}{rgb}{0.859516,0.942330,0.747405}%
\pgfsetfillcolor{currentfill}%
\pgfsetlinewidth{0.000000pt}%
\definecolor{currentstroke}{rgb}{0.000000,0.000000,0.000000}%
\pgfsetstrokecolor{currentstroke}%
\pgfsetstrokeopacity{0.000000}%
\pgfsetdash{}{0pt}%
\pgfpathmoveto{\pgfqpoint{0.528131in}{0.499862in}}%
\pgfpathlineto{\pgfqpoint{0.527132in}{0.498106in}}%
\pgfpathlineto{\pgfqpoint{0.529775in}{0.494368in}}%
\pgfpathlineto{\pgfqpoint{0.525597in}{0.490871in}}%
\pgfpathlineto{\pgfqpoint{0.525713in}{0.487863in}}%
\pgfpathlineto{\pgfqpoint{0.523824in}{0.487153in}}%
\pgfpathlineto{\pgfqpoint{0.517331in}{0.491639in}}%
\pgfpathlineto{\pgfqpoint{0.512703in}{0.501458in}}%
\pgfpathlineto{\pgfqpoint{0.512393in}{0.504295in}}%
\pgfpathlineto{\pgfqpoint{0.513583in}{0.507691in}}%
\pgfpathlineto{\pgfqpoint{0.518598in}{0.502496in}}%
\pgfpathlineto{\pgfqpoint{0.520132in}{0.503498in}}%
\pgfpathlineto{\pgfqpoint{0.524386in}{0.502545in}}%
\pgfpathlineto{\pgfqpoint{0.528131in}{0.499862in}}%
\pgfpathclose%
\pgfusepath{fill}%
\end{pgfscope}%
\begin{pgfscope}%
\pgfpathrectangle{\pgfqpoint{0.100000in}{0.100000in}}{\pgfqpoint{2.862985in}{1.832778in}}%
\pgfusepath{clip}%
\pgfsetbuttcap%
\pgfsetmiterjoin%
\definecolor{currentfill}{rgb}{0.859516,0.942330,0.747405}%
\pgfsetfillcolor{currentfill}%
\pgfsetlinewidth{0.000000pt}%
\definecolor{currentstroke}{rgb}{0.000000,0.000000,0.000000}%
\pgfsetstrokecolor{currentstroke}%
\pgfsetstrokeopacity{0.000000}%
\pgfsetdash{}{0pt}%
\pgfpathmoveto{\pgfqpoint{0.451126in}{0.386109in}}%
\pgfpathlineto{\pgfqpoint{0.446562in}{0.385368in}}%
\pgfpathlineto{\pgfqpoint{0.443260in}{0.386895in}}%
\pgfpathlineto{\pgfqpoint{0.441816in}{0.389380in}}%
\pgfpathlineto{\pgfqpoint{0.443453in}{0.392945in}}%
\pgfpathlineto{\pgfqpoint{0.447096in}{0.392002in}}%
\pgfpathlineto{\pgfqpoint{0.453086in}{0.394117in}}%
\pgfpathlineto{\pgfqpoint{0.455184in}{0.391196in}}%
\pgfpathlineto{\pgfqpoint{0.457217in}{0.391735in}}%
\pgfpathlineto{\pgfqpoint{0.462141in}{0.389917in}}%
\pgfpathlineto{\pgfqpoint{0.464268in}{0.387608in}}%
\pgfpathlineto{\pgfqpoint{0.461672in}{0.381567in}}%
\pgfpathlineto{\pgfqpoint{0.459090in}{0.380286in}}%
\pgfpathlineto{\pgfqpoint{0.456769in}{0.380976in}}%
\pgfpathlineto{\pgfqpoint{0.451126in}{0.386109in}}%
\pgfpathclose%
\pgfusepath{fill}%
\end{pgfscope}%
\begin{pgfscope}%
\pgfpathrectangle{\pgfqpoint{0.100000in}{0.100000in}}{\pgfqpoint{2.862985in}{1.832778in}}%
\pgfusepath{clip}%
\pgfsetbuttcap%
\pgfsetmiterjoin%
\definecolor{currentfill}{rgb}{0.859516,0.942330,0.747405}%
\pgfsetfillcolor{currentfill}%
\pgfsetlinewidth{0.000000pt}%
\definecolor{currentstroke}{rgb}{0.000000,0.000000,0.000000}%
\pgfsetstrokecolor{currentstroke}%
\pgfsetstrokeopacity{0.000000}%
\pgfsetdash{}{0pt}%
\pgfpathmoveto{\pgfqpoint{0.410323in}{0.391598in}}%
\pgfpathlineto{\pgfqpoint{0.403857in}{0.394682in}}%
\pgfpathlineto{\pgfqpoint{0.397873in}{0.395600in}}%
\pgfpathlineto{\pgfqpoint{0.395987in}{0.396913in}}%
\pgfpathlineto{\pgfqpoint{0.398135in}{0.399499in}}%
\pgfpathlineto{\pgfqpoint{0.401946in}{0.396269in}}%
\pgfpathlineto{\pgfqpoint{0.405142in}{0.396714in}}%
\pgfpathlineto{\pgfqpoint{0.410341in}{0.398951in}}%
\pgfpathlineto{\pgfqpoint{0.410761in}{0.401832in}}%
\pgfpathlineto{\pgfqpoint{0.417146in}{0.400406in}}%
\pgfpathlineto{\pgfqpoint{0.415651in}{0.396069in}}%
\pgfpathlineto{\pgfqpoint{0.419719in}{0.395751in}}%
\pgfpathlineto{\pgfqpoint{0.416306in}{0.390592in}}%
\pgfpathlineto{\pgfqpoint{0.413001in}{0.392264in}}%
\pgfpathlineto{\pgfqpoint{0.410323in}{0.391598in}}%
\pgfpathclose%
\pgfusepath{fill}%
\end{pgfscope}%
\begin{pgfscope}%
\pgfpathrectangle{\pgfqpoint{0.100000in}{0.100000in}}{\pgfqpoint{2.862985in}{1.832778in}}%
\pgfusepath{clip}%
\pgfsetbuttcap%
\pgfsetmiterjoin%
\definecolor{currentfill}{rgb}{0.859516,0.942330,0.747405}%
\pgfsetfillcolor{currentfill}%
\pgfsetlinewidth{0.000000pt}%
\definecolor{currentstroke}{rgb}{0.000000,0.000000,0.000000}%
\pgfsetstrokecolor{currentstroke}%
\pgfsetstrokeopacity{0.000000}%
\pgfsetdash{}{0pt}%
\pgfpathmoveto{\pgfqpoint{0.398962in}{0.402261in}}%
\pgfpathlineto{\pgfqpoint{0.396586in}{0.400969in}}%
\pgfpathlineto{\pgfqpoint{0.390110in}{0.402955in}}%
\pgfpathlineto{\pgfqpoint{0.385271in}{0.401805in}}%
\pgfpathlineto{\pgfqpoint{0.382813in}{0.405073in}}%
\pgfpathlineto{\pgfqpoint{0.383885in}{0.406569in}}%
\pgfpathlineto{\pgfqpoint{0.387552in}{0.406819in}}%
\pgfpathlineto{\pgfqpoint{0.390357in}{0.405880in}}%
\pgfpathlineto{\pgfqpoint{0.393414in}{0.407367in}}%
\pgfpathlineto{\pgfqpoint{0.398100in}{0.405378in}}%
\pgfpathlineto{\pgfqpoint{0.398962in}{0.402261in}}%
\pgfpathclose%
\pgfusepath{fill}%
\end{pgfscope}%
\begin{pgfscope}%
\pgfpathrectangle{\pgfqpoint{0.100000in}{0.100000in}}{\pgfqpoint{2.862985in}{1.832778in}}%
\pgfusepath{clip}%
\pgfsetbuttcap%
\pgfsetmiterjoin%
\definecolor{currentfill}{rgb}{0.859516,0.942330,0.747405}%
\pgfsetfillcolor{currentfill}%
\pgfsetlinewidth{0.000000pt}%
\definecolor{currentstroke}{rgb}{0.000000,0.000000,0.000000}%
\pgfsetstrokecolor{currentstroke}%
\pgfsetstrokeopacity{0.000000}%
\pgfsetdash{}{0pt}%
\pgfpathmoveto{\pgfqpoint{0.322628in}{0.450357in}}%
\pgfpathlineto{\pgfqpoint{0.322456in}{0.447153in}}%
\pgfpathlineto{\pgfqpoint{0.317729in}{0.445139in}}%
\pgfpathlineto{\pgfqpoint{0.313412in}{0.450180in}}%
\pgfpathlineto{\pgfqpoint{0.318253in}{0.451860in}}%
\pgfpathlineto{\pgfqpoint{0.322628in}{0.450357in}}%
\pgfpathclose%
\pgfusepath{fill}%
\end{pgfscope}%
\begin{pgfscope}%
\pgfpathrectangle{\pgfqpoint{0.100000in}{0.100000in}}{\pgfqpoint{2.862985in}{1.832778in}}%
\pgfusepath{clip}%
\pgfsetbuttcap%
\pgfsetmiterjoin%
\definecolor{currentfill}{rgb}{0.859516,0.942330,0.747405}%
\pgfsetfillcolor{currentfill}%
\pgfsetlinewidth{0.000000pt}%
\definecolor{currentstroke}{rgb}{0.000000,0.000000,0.000000}%
\pgfsetstrokecolor{currentstroke}%
\pgfsetstrokeopacity{0.000000}%
\pgfsetdash{}{0pt}%
\pgfpathmoveto{\pgfqpoint{0.283638in}{0.468371in}}%
\pgfpathlineto{\pgfqpoint{0.284642in}{0.469764in}}%
\pgfpathlineto{\pgfqpoint{0.290440in}{0.469524in}}%
\pgfpathlineto{\pgfqpoint{0.291458in}{0.471597in}}%
\pgfpathlineto{\pgfqpoint{0.294002in}{0.469788in}}%
\pgfpathlineto{\pgfqpoint{0.292533in}{0.466308in}}%
\pgfpathlineto{\pgfqpoint{0.290372in}{0.466014in}}%
\pgfpathlineto{\pgfqpoint{0.285207in}{0.467106in}}%
\pgfpathlineto{\pgfqpoint{0.283638in}{0.468371in}}%
\pgfpathclose%
\pgfusepath{fill}%
\end{pgfscope}%
\begin{pgfscope}%
\pgfpathrectangle{\pgfqpoint{0.100000in}{0.100000in}}{\pgfqpoint{2.862985in}{1.832778in}}%
\pgfusepath{clip}%
\pgfsetbuttcap%
\pgfsetmiterjoin%
\definecolor{currentfill}{rgb}{0.859516,0.942330,0.747405}%
\pgfsetfillcolor{currentfill}%
\pgfsetlinewidth{0.000000pt}%
\definecolor{currentstroke}{rgb}{0.000000,0.000000,0.000000}%
\pgfsetstrokecolor{currentstroke}%
\pgfsetstrokeopacity{0.000000}%
\pgfsetdash{}{0pt}%
\pgfpathmoveto{\pgfqpoint{0.275415in}{0.479156in}}%
\pgfpathlineto{\pgfqpoint{0.275063in}{0.482294in}}%
\pgfpathlineto{\pgfqpoint{0.278088in}{0.482062in}}%
\pgfpathlineto{\pgfqpoint{0.277878in}{0.486470in}}%
\pgfpathlineto{\pgfqpoint{0.280921in}{0.484091in}}%
\pgfpathlineto{\pgfqpoint{0.279868in}{0.480653in}}%
\pgfpathlineto{\pgfqpoint{0.275415in}{0.479156in}}%
\pgfpathclose%
\pgfusepath{fill}%
\end{pgfscope}%
\begin{pgfscope}%
\pgfpathrectangle{\pgfqpoint{0.100000in}{0.100000in}}{\pgfqpoint{2.862985in}{1.832778in}}%
\pgfusepath{clip}%
\pgfsetbuttcap%
\pgfsetmiterjoin%
\definecolor{currentfill}{rgb}{0.859516,0.942330,0.747405}%
\pgfsetfillcolor{currentfill}%
\pgfsetlinewidth{0.000000pt}%
\definecolor{currentstroke}{rgb}{0.000000,0.000000,0.000000}%
\pgfsetstrokecolor{currentstroke}%
\pgfsetstrokeopacity{0.000000}%
\pgfsetdash{}{0pt}%
\pgfpathmoveto{\pgfqpoint{0.541980in}{0.591016in}}%
\pgfpathlineto{\pgfqpoint{0.537154in}{0.592058in}}%
\pgfpathlineto{\pgfqpoint{0.536118in}{0.595283in}}%
\pgfpathlineto{\pgfqpoint{0.537776in}{0.598507in}}%
\pgfpathlineto{\pgfqpoint{0.541491in}{0.600190in}}%
\pgfpathlineto{\pgfqpoint{0.545744in}{0.591916in}}%
\pgfpathlineto{\pgfqpoint{0.549601in}{0.591548in}}%
\pgfpathlineto{\pgfqpoint{0.552479in}{0.588995in}}%
\pgfpathlineto{\pgfqpoint{0.552422in}{0.585482in}}%
\pgfpathlineto{\pgfqpoint{0.550318in}{0.583797in}}%
\pgfpathlineto{\pgfqpoint{0.553740in}{0.573408in}}%
\pgfpathlineto{\pgfqpoint{0.557484in}{0.569595in}}%
\pgfpathlineto{\pgfqpoint{0.553656in}{0.568385in}}%
\pgfpathlineto{\pgfqpoint{0.550557in}{0.572663in}}%
\pgfpathlineto{\pgfqpoint{0.543823in}{0.571338in}}%
\pgfpathlineto{\pgfqpoint{0.545918in}{0.575572in}}%
\pgfpathlineto{\pgfqpoint{0.544943in}{0.577258in}}%
\pgfpathlineto{\pgfqpoint{0.545514in}{0.581671in}}%
\pgfpathlineto{\pgfqpoint{0.544180in}{0.589648in}}%
\pgfpathlineto{\pgfqpoint{0.541980in}{0.591016in}}%
\pgfpathclose%
\pgfusepath{fill}%
\end{pgfscope}%
\begin{pgfscope}%
\pgfpathrectangle{\pgfqpoint{0.100000in}{0.100000in}}{\pgfqpoint{2.862985in}{1.832778in}}%
\pgfusepath{clip}%
\pgfsetbuttcap%
\pgfsetmiterjoin%
\definecolor{currentfill}{rgb}{0.859516,0.942330,0.747405}%
\pgfsetfillcolor{currentfill}%
\pgfsetlinewidth{0.000000pt}%
\definecolor{currentstroke}{rgb}{0.000000,0.000000,0.000000}%
\pgfsetstrokecolor{currentstroke}%
\pgfsetstrokeopacity{0.000000}%
\pgfsetdash{}{0pt}%
\pgfpathmoveto{\pgfqpoint{0.711519in}{0.353334in}}%
\pgfpathlineto{\pgfqpoint{0.704523in}{0.353412in}}%
\pgfpathlineto{\pgfqpoint{0.705062in}{0.356924in}}%
\pgfpathlineto{\pgfqpoint{0.707862in}{0.358180in}}%
\pgfpathlineto{\pgfqpoint{0.711519in}{0.353334in}}%
\pgfpathclose%
\pgfusepath{fill}%
\end{pgfscope}%
\begin{pgfscope}%
\pgfpathrectangle{\pgfqpoint{0.100000in}{0.100000in}}{\pgfqpoint{2.862985in}{1.832778in}}%
\pgfusepath{clip}%
\pgfsetbuttcap%
\pgfsetmiterjoin%
\definecolor{currentfill}{rgb}{0.859516,0.942330,0.747405}%
\pgfsetfillcolor{currentfill}%
\pgfsetlinewidth{0.000000pt}%
\definecolor{currentstroke}{rgb}{0.000000,0.000000,0.000000}%
\pgfsetstrokecolor{currentstroke}%
\pgfsetstrokeopacity{0.000000}%
\pgfsetdash{}{0pt}%
\pgfpathmoveto{\pgfqpoint{0.701812in}{0.356514in}}%
\pgfpathlineto{\pgfqpoint{0.696800in}{0.355301in}}%
\pgfpathlineto{\pgfqpoint{0.692038in}{0.353224in}}%
\pgfpathlineto{\pgfqpoint{0.690567in}{0.355671in}}%
\pgfpathlineto{\pgfqpoint{0.698757in}{0.357328in}}%
\pgfpathlineto{\pgfqpoint{0.701812in}{0.356514in}}%
\pgfpathclose%
\pgfusepath{fill}%
\end{pgfscope}%
\begin{pgfscope}%
\pgfpathrectangle{\pgfqpoint{0.100000in}{0.100000in}}{\pgfqpoint{2.862985in}{1.832778in}}%
\pgfusepath{clip}%
\pgfsetbuttcap%
\pgfsetmiterjoin%
\definecolor{currentfill}{rgb}{0.859516,0.942330,0.747405}%
\pgfsetfillcolor{currentfill}%
\pgfsetlinewidth{0.000000pt}%
\definecolor{currentstroke}{rgb}{0.000000,0.000000,0.000000}%
\pgfsetstrokecolor{currentstroke}%
\pgfsetstrokeopacity{0.000000}%
\pgfsetdash{}{0pt}%
\pgfpathmoveto{\pgfqpoint{0.624990in}{0.354028in}}%
\pgfpathlineto{\pgfqpoint{0.625182in}{0.351599in}}%
\pgfpathlineto{\pgfqpoint{0.622294in}{0.350090in}}%
\pgfpathlineto{\pgfqpoint{0.614235in}{0.353552in}}%
\pgfpathlineto{\pgfqpoint{0.612986in}{0.352159in}}%
\pgfpathlineto{\pgfqpoint{0.611036in}{0.358747in}}%
\pgfpathlineto{\pgfqpoint{0.615074in}{0.359664in}}%
\pgfpathlineto{\pgfqpoint{0.616940in}{0.356957in}}%
\pgfpathlineto{\pgfqpoint{0.620992in}{0.360345in}}%
\pgfpathlineto{\pgfqpoint{0.624990in}{0.354028in}}%
\pgfpathclose%
\pgfusepath{fill}%
\end{pgfscope}%
\begin{pgfscope}%
\pgfpathrectangle{\pgfqpoint{0.100000in}{0.100000in}}{\pgfqpoint{2.862985in}{1.832778in}}%
\pgfusepath{clip}%
\pgfsetbuttcap%
\pgfsetmiterjoin%
\definecolor{currentfill}{rgb}{0.859516,0.942330,0.747405}%
\pgfsetfillcolor{currentfill}%
\pgfsetlinewidth{0.000000pt}%
\definecolor{currentstroke}{rgb}{0.000000,0.000000,0.000000}%
\pgfsetstrokecolor{currentstroke}%
\pgfsetstrokeopacity{0.000000}%
\pgfsetdash{}{0pt}%
\pgfpathmoveto{\pgfqpoint{0.806774in}{0.219412in}}%
\pgfpathlineto{\pgfqpoint{0.800017in}{0.216896in}}%
\pgfpathlineto{\pgfqpoint{0.796437in}{0.216745in}}%
\pgfpathlineto{\pgfqpoint{0.799350in}{0.222082in}}%
\pgfpathlineto{\pgfqpoint{0.801840in}{0.224073in}}%
\pgfpathlineto{\pgfqpoint{0.801760in}{0.229600in}}%
\pgfpathlineto{\pgfqpoint{0.804465in}{0.239578in}}%
\pgfpathlineto{\pgfqpoint{0.806957in}{0.241473in}}%
\pgfpathlineto{\pgfqpoint{0.812710in}{0.238729in}}%
\pgfpathlineto{\pgfqpoint{0.811870in}{0.237158in}}%
\pgfpathlineto{\pgfqpoint{0.812306in}{0.229604in}}%
\pgfpathlineto{\pgfqpoint{0.810440in}{0.227566in}}%
\pgfpathlineto{\pgfqpoint{0.809640in}{0.222201in}}%
\pgfpathlineto{\pgfqpoint{0.806774in}{0.219412in}}%
\pgfpathclose%
\pgfusepath{fill}%
\end{pgfscope}%
\begin{pgfscope}%
\pgfpathrectangle{\pgfqpoint{0.100000in}{0.100000in}}{\pgfqpoint{2.862985in}{1.832778in}}%
\pgfusepath{clip}%
\pgfsetbuttcap%
\pgfsetmiterjoin%
\definecolor{currentfill}{rgb}{0.859516,0.942330,0.747405}%
\pgfsetfillcolor{currentfill}%
\pgfsetlinewidth{0.000000pt}%
\definecolor{currentstroke}{rgb}{0.000000,0.000000,0.000000}%
\pgfsetstrokecolor{currentstroke}%
\pgfsetstrokeopacity{0.000000}%
\pgfsetdash{}{0pt}%
\pgfpathmoveto{\pgfqpoint{0.791809in}{0.244146in}}%
\pgfpathlineto{\pgfqpoint{0.795768in}{0.237310in}}%
\pgfpathlineto{\pgfqpoint{0.795034in}{0.235680in}}%
\pgfpathlineto{\pgfqpoint{0.800329in}{0.234022in}}%
\pgfpathlineto{\pgfqpoint{0.798809in}{0.227738in}}%
\pgfpathlineto{\pgfqpoint{0.796503in}{0.228029in}}%
\pgfpathlineto{\pgfqpoint{0.790445in}{0.238982in}}%
\pgfpathlineto{\pgfqpoint{0.791307in}{0.234052in}}%
\pgfpathlineto{\pgfqpoint{0.790275in}{0.231213in}}%
\pgfpathlineto{\pgfqpoint{0.787709in}{0.229950in}}%
\pgfpathlineto{\pgfqpoint{0.786251in}{0.231311in}}%
\pgfpathlineto{\pgfqpoint{0.785437in}{0.237623in}}%
\pgfpathlineto{\pgfqpoint{0.785917in}{0.241717in}}%
\pgfpathlineto{\pgfqpoint{0.784208in}{0.243504in}}%
\pgfpathlineto{\pgfqpoint{0.788722in}{0.248675in}}%
\pgfpathlineto{\pgfqpoint{0.791684in}{0.250209in}}%
\pgfpathlineto{\pgfqpoint{0.792670in}{0.248193in}}%
\pgfpathlineto{\pgfqpoint{0.795760in}{0.249753in}}%
\pgfpathlineto{\pgfqpoint{0.797999in}{0.245516in}}%
\pgfpathlineto{\pgfqpoint{0.802840in}{0.240126in}}%
\pgfpathlineto{\pgfqpoint{0.802179in}{0.237684in}}%
\pgfpathlineto{\pgfqpoint{0.799491in}{0.235024in}}%
\pgfpathlineto{\pgfqpoint{0.797005in}{0.236278in}}%
\pgfpathlineto{\pgfqpoint{0.791809in}{0.244146in}}%
\pgfpathclose%
\pgfusepath{fill}%
\end{pgfscope}%
\begin{pgfscope}%
\pgfpathrectangle{\pgfqpoint{0.100000in}{0.100000in}}{\pgfqpoint{2.862985in}{1.832778in}}%
\pgfusepath{clip}%
\pgfsetbuttcap%
\pgfsetmiterjoin%
\definecolor{currentfill}{rgb}{0.859516,0.942330,0.747405}%
\pgfsetfillcolor{currentfill}%
\pgfsetlinewidth{0.000000pt}%
\definecolor{currentstroke}{rgb}{0.000000,0.000000,0.000000}%
\pgfsetstrokecolor{currentstroke}%
\pgfsetstrokeopacity{0.000000}%
\pgfsetdash{}{0pt}%
\pgfpathmoveto{\pgfqpoint{0.595353in}{0.341071in}}%
\pgfpathlineto{\pgfqpoint{0.590449in}{0.340433in}}%
\pgfpathlineto{\pgfqpoint{0.583542in}{0.337778in}}%
\pgfpathlineto{\pgfqpoint{0.581931in}{0.339221in}}%
\pgfpathlineto{\pgfqpoint{0.589146in}{0.340851in}}%
\pgfpathlineto{\pgfqpoint{0.587212in}{0.342273in}}%
\pgfpathlineto{\pgfqpoint{0.582739in}{0.343315in}}%
\pgfpathlineto{\pgfqpoint{0.581544in}{0.346649in}}%
\pgfpathlineto{\pgfqpoint{0.583795in}{0.349521in}}%
\pgfpathlineto{\pgfqpoint{0.585690in}{0.356112in}}%
\pgfpathlineto{\pgfqpoint{0.592614in}{0.357340in}}%
\pgfpathlineto{\pgfqpoint{0.595966in}{0.354492in}}%
\pgfpathlineto{\pgfqpoint{0.599115in}{0.357248in}}%
\pgfpathlineto{\pgfqpoint{0.603064in}{0.356572in}}%
\pgfpathlineto{\pgfqpoint{0.605902in}{0.351482in}}%
\pgfpathlineto{\pgfqpoint{0.608305in}{0.356006in}}%
\pgfpathlineto{\pgfqpoint{0.615126in}{0.345423in}}%
\pgfpathlineto{\pgfqpoint{0.612662in}{0.344577in}}%
\pgfpathlineto{\pgfqpoint{0.611150in}{0.342283in}}%
\pgfpathlineto{\pgfqpoint{0.614007in}{0.340190in}}%
\pgfpathlineto{\pgfqpoint{0.609576in}{0.338227in}}%
\pgfpathlineto{\pgfqpoint{0.606364in}{0.339308in}}%
\pgfpathlineto{\pgfqpoint{0.600296in}{0.342890in}}%
\pgfpathlineto{\pgfqpoint{0.595353in}{0.341071in}}%
\pgfpathclose%
\pgfusepath{fill}%
\end{pgfscope}%
\begin{pgfscope}%
\pgfpathrectangle{\pgfqpoint{0.100000in}{0.100000in}}{\pgfqpoint{2.862985in}{1.832778in}}%
\pgfusepath{clip}%
\pgfsetbuttcap%
\pgfsetmiterjoin%
\definecolor{currentfill}{rgb}{0.859516,0.942330,0.747405}%
\pgfsetfillcolor{currentfill}%
\pgfsetlinewidth{0.000000pt}%
\definecolor{currentstroke}{rgb}{0.000000,0.000000,0.000000}%
\pgfsetstrokecolor{currentstroke}%
\pgfsetstrokeopacity{0.000000}%
\pgfsetdash{}{0pt}%
\pgfpathmoveto{\pgfqpoint{0.787691in}{0.199220in}}%
\pgfpathlineto{\pgfqpoint{0.786719in}{0.208994in}}%
\pgfpathlineto{\pgfqpoint{0.789340in}{0.212583in}}%
\pgfpathlineto{\pgfqpoint{0.786538in}{0.212832in}}%
\pgfpathlineto{\pgfqpoint{0.787562in}{0.219812in}}%
\pgfpathlineto{\pgfqpoint{0.788770in}{0.223955in}}%
\pgfpathlineto{\pgfqpoint{0.787720in}{0.229522in}}%
\pgfpathlineto{\pgfqpoint{0.790287in}{0.230600in}}%
\pgfpathlineto{\pgfqpoint{0.791124in}{0.232039in}}%
\pgfpathlineto{\pgfqpoint{0.793653in}{0.231493in}}%
\pgfpathlineto{\pgfqpoint{0.795842in}{0.227033in}}%
\pgfpathlineto{\pgfqpoint{0.796249in}{0.222947in}}%
\pgfpathlineto{\pgfqpoint{0.793571in}{0.210709in}}%
\pgfpathlineto{\pgfqpoint{0.787691in}{0.199220in}}%
\pgfpathclose%
\pgfusepath{fill}%
\end{pgfscope}%
\begin{pgfscope}%
\pgfpathrectangle{\pgfqpoint{0.100000in}{0.100000in}}{\pgfqpoint{2.862985in}{1.832778in}}%
\pgfusepath{clip}%
\pgfsetbuttcap%
\pgfsetmiterjoin%
\definecolor{currentfill}{rgb}{0.859516,0.942330,0.747405}%
\pgfsetfillcolor{currentfill}%
\pgfsetlinewidth{0.000000pt}%
\definecolor{currentstroke}{rgb}{0.000000,0.000000,0.000000}%
\pgfsetstrokecolor{currentstroke}%
\pgfsetstrokeopacity{0.000000}%
\pgfsetdash{}{0pt}%
\pgfpathmoveto{\pgfqpoint{0.787110in}{0.229056in}}%
\pgfpathlineto{\pgfqpoint{0.786975in}{0.224290in}}%
\pgfpathlineto{\pgfqpoint{0.784778in}{0.222127in}}%
\pgfpathlineto{\pgfqpoint{0.782238in}{0.222926in}}%
\pgfpathlineto{\pgfqpoint{0.785448in}{0.229848in}}%
\pgfpathlineto{\pgfqpoint{0.787110in}{0.229056in}}%
\pgfpathclose%
\pgfusepath{fill}%
\end{pgfscope}%
\begin{pgfscope}%
\pgfpathrectangle{\pgfqpoint{0.100000in}{0.100000in}}{\pgfqpoint{2.862985in}{1.832778in}}%
\pgfusepath{clip}%
\pgfsetbuttcap%
\pgfsetmiterjoin%
\definecolor{currentfill}{rgb}{0.859516,0.942330,0.747405}%
\pgfsetfillcolor{currentfill}%
\pgfsetlinewidth{0.000000pt}%
\definecolor{currentstroke}{rgb}{0.000000,0.000000,0.000000}%
\pgfsetstrokecolor{currentstroke}%
\pgfsetstrokeopacity{0.000000}%
\pgfsetdash{}{0pt}%
\pgfpathmoveto{\pgfqpoint{0.814076in}{0.207898in}}%
\pgfpathlineto{\pgfqpoint{0.812202in}{0.199923in}}%
\pgfpathlineto{\pgfqpoint{0.807835in}{0.197402in}}%
\pgfpathlineto{\pgfqpoint{0.802107in}{0.199724in}}%
\pgfpathlineto{\pgfqpoint{0.803685in}{0.202824in}}%
\pgfpathlineto{\pgfqpoint{0.803653in}{0.204970in}}%
\pgfpathlineto{\pgfqpoint{0.805783in}{0.208616in}}%
\pgfpathlineto{\pgfqpoint{0.804286in}{0.207882in}}%
\pgfpathlineto{\pgfqpoint{0.805377in}{0.209663in}}%
\pgfpathlineto{\pgfqpoint{0.803304in}{0.213709in}}%
\pgfpathlineto{\pgfqpoint{0.805657in}{0.214455in}}%
\pgfpathlineto{\pgfqpoint{0.811104in}{0.209721in}}%
\pgfpathlineto{\pgfqpoint{0.814076in}{0.207898in}}%
\pgfpathclose%
\pgfusepath{fill}%
\end{pgfscope}%
\begin{pgfscope}%
\pgfpathrectangle{\pgfqpoint{0.100000in}{0.100000in}}{\pgfqpoint{2.862985in}{1.832778in}}%
\pgfusepath{clip}%
\pgfsetbuttcap%
\pgfsetmiterjoin%
\definecolor{currentfill}{rgb}{0.859516,0.942330,0.747405}%
\pgfsetfillcolor{currentfill}%
\pgfsetlinewidth{0.000000pt}%
\definecolor{currentstroke}{rgb}{0.000000,0.000000,0.000000}%
\pgfsetstrokecolor{currentstroke}%
\pgfsetstrokeopacity{0.000000}%
\pgfsetdash{}{0pt}%
\pgfpathmoveto{\pgfqpoint{0.794962in}{0.193796in}}%
\pgfpathlineto{\pgfqpoint{0.792492in}{0.192836in}}%
\pgfpathlineto{\pgfqpoint{0.793902in}{0.201252in}}%
\pgfpathlineto{\pgfqpoint{0.796866in}{0.203216in}}%
\pgfpathlineto{\pgfqpoint{0.796144in}{0.209521in}}%
\pgfpathlineto{\pgfqpoint{0.797348in}{0.212317in}}%
\pgfpathlineto{\pgfqpoint{0.799731in}{0.213365in}}%
\pgfpathlineto{\pgfqpoint{0.802232in}{0.210622in}}%
\pgfpathlineto{\pgfqpoint{0.804379in}{0.206901in}}%
\pgfpathlineto{\pgfqpoint{0.797478in}{0.199610in}}%
\pgfpathlineto{\pgfqpoint{0.794962in}{0.193796in}}%
\pgfpathclose%
\pgfusepath{fill}%
\end{pgfscope}%
\begin{pgfscope}%
\pgfpathrectangle{\pgfqpoint{0.100000in}{0.100000in}}{\pgfqpoint{2.862985in}{1.832778in}}%
\pgfusepath{clip}%
\pgfsetbuttcap%
\pgfsetmiterjoin%
\definecolor{currentfill}{rgb}{0.859516,0.942330,0.747405}%
\pgfsetfillcolor{currentfill}%
\pgfsetlinewidth{0.000000pt}%
\definecolor{currentstroke}{rgb}{0.000000,0.000000,0.000000}%
\pgfsetstrokecolor{currentstroke}%
\pgfsetstrokeopacity{0.000000}%
\pgfsetdash{}{0pt}%
\pgfpathmoveto{\pgfqpoint{0.815172in}{0.202568in}}%
\pgfpathlineto{\pgfqpoint{0.816382in}{0.196760in}}%
\pgfpathlineto{\pgfqpoint{0.810974in}{0.197425in}}%
\pgfpathlineto{\pgfqpoint{0.815172in}{0.202568in}}%
\pgfpathclose%
\pgfusepath{fill}%
\end{pgfscope}%
\begin{pgfscope}%
\pgfpathrectangle{\pgfqpoint{0.100000in}{0.100000in}}{\pgfqpoint{2.862985in}{1.832778in}}%
\pgfusepath{clip}%
\pgfsetbuttcap%
\pgfsetmiterjoin%
\definecolor{currentfill}{rgb}{0.859516,0.942330,0.747405}%
\pgfsetfillcolor{currentfill}%
\pgfsetlinewidth{0.000000pt}%
\definecolor{currentstroke}{rgb}{0.000000,0.000000,0.000000}%
\pgfsetstrokecolor{currentstroke}%
\pgfsetstrokeopacity{0.000000}%
\pgfsetdash{}{0pt}%
\pgfpathmoveto{\pgfqpoint{0.817752in}{0.193694in}}%
\pgfpathlineto{\pgfqpoint{0.818638in}{0.190111in}}%
\pgfpathlineto{\pgfqpoint{0.820335in}{0.188965in}}%
\pgfpathlineto{\pgfqpoint{0.820026in}{0.185593in}}%
\pgfpathlineto{\pgfqpoint{0.817656in}{0.184459in}}%
\pgfpathlineto{\pgfqpoint{0.815732in}{0.189493in}}%
\pgfpathlineto{\pgfqpoint{0.817752in}{0.193694in}}%
\pgfpathclose%
\pgfusepath{fill}%
\end{pgfscope}%
\begin{pgfscope}%
\pgfpathrectangle{\pgfqpoint{0.100000in}{0.100000in}}{\pgfqpoint{2.862985in}{1.832778in}}%
\pgfusepath{clip}%
\pgfsetbuttcap%
\pgfsetmiterjoin%
\definecolor{currentfill}{rgb}{0.859516,0.942330,0.747405}%
\pgfsetfillcolor{currentfill}%
\pgfsetlinewidth{0.000000pt}%
\definecolor{currentstroke}{rgb}{0.000000,0.000000,0.000000}%
\pgfsetstrokecolor{currentstroke}%
\pgfsetstrokeopacity{0.000000}%
\pgfsetdash{}{0pt}%
\pgfpathmoveto{\pgfqpoint{0.813239in}{0.195068in}}%
\pgfpathlineto{\pgfqpoint{0.811768in}{0.190907in}}%
\pgfpathlineto{\pgfqpoint{0.809484in}{0.191060in}}%
\pgfpathlineto{\pgfqpoint{0.808043in}{0.194492in}}%
\pgfpathlineto{\pgfqpoint{0.810390in}{0.196117in}}%
\pgfpathlineto{\pgfqpoint{0.813239in}{0.195068in}}%
\pgfpathclose%
\pgfusepath{fill}%
\end{pgfscope}%
\begin{pgfscope}%
\pgfpathrectangle{\pgfqpoint{0.100000in}{0.100000in}}{\pgfqpoint{2.862985in}{1.832778in}}%
\pgfusepath{clip}%
\pgfsetbuttcap%
\pgfsetmiterjoin%
\definecolor{currentfill}{rgb}{0.859516,0.942330,0.747405}%
\pgfsetfillcolor{currentfill}%
\pgfsetlinewidth{0.000000pt}%
\definecolor{currentstroke}{rgb}{0.000000,0.000000,0.000000}%
\pgfsetstrokecolor{currentstroke}%
\pgfsetstrokeopacity{0.000000}%
\pgfsetdash{}{0pt}%
\pgfpathmoveto{\pgfqpoint{0.807076in}{0.173613in}}%
\pgfpathlineto{\pgfqpoint{0.809114in}{0.171384in}}%
\pgfpathlineto{\pgfqpoint{0.809563in}{0.166686in}}%
\pgfpathlineto{\pgfqpoint{0.810636in}{0.165321in}}%
\pgfpathlineto{\pgfqpoint{0.808884in}{0.161070in}}%
\pgfpathlineto{\pgfqpoint{0.806578in}{0.158676in}}%
\pgfpathlineto{\pgfqpoint{0.805352in}{0.154236in}}%
\pgfpathlineto{\pgfqpoint{0.802832in}{0.161156in}}%
\pgfpathlineto{\pgfqpoint{0.802414in}{0.167439in}}%
\pgfpathlineto{\pgfqpoint{0.801004in}{0.170567in}}%
\pgfpathlineto{\pgfqpoint{0.796147in}{0.173102in}}%
\pgfpathlineto{\pgfqpoint{0.800703in}{0.178390in}}%
\pgfpathlineto{\pgfqpoint{0.797666in}{0.180650in}}%
\pgfpathlineto{\pgfqpoint{0.800341in}{0.182902in}}%
\pgfpathlineto{\pgfqpoint{0.803893in}{0.191549in}}%
\pgfpathlineto{\pgfqpoint{0.800212in}{0.194597in}}%
\pgfpathlineto{\pgfqpoint{0.801630in}{0.197524in}}%
\pgfpathlineto{\pgfqpoint{0.806388in}{0.194928in}}%
\pgfpathlineto{\pgfqpoint{0.806870in}{0.190831in}}%
\pgfpathlineto{\pgfqpoint{0.805070in}{0.188498in}}%
\pgfpathlineto{\pgfqpoint{0.808638in}{0.185341in}}%
\pgfpathlineto{\pgfqpoint{0.809761in}{0.180424in}}%
\pgfpathlineto{\pgfqpoint{0.809609in}{0.173528in}}%
\pgfpathlineto{\pgfqpoint{0.807076in}{0.173613in}}%
\pgfpathclose%
\pgfusepath{fill}%
\end{pgfscope}%
\begin{pgfscope}%
\pgfpathrectangle{\pgfqpoint{0.100000in}{0.100000in}}{\pgfqpoint{2.862985in}{1.832778in}}%
\pgfusepath{clip}%
\pgfsetbuttcap%
\pgfsetmiterjoin%
\definecolor{currentfill}{rgb}{0.859516,0.942330,0.747405}%
\pgfsetfillcolor{currentfill}%
\pgfsetlinewidth{0.000000pt}%
\definecolor{currentstroke}{rgb}{0.000000,0.000000,0.000000}%
\pgfsetstrokecolor{currentstroke}%
\pgfsetstrokeopacity{0.000000}%
\pgfsetdash{}{0pt}%
\pgfpathmoveto{\pgfqpoint{0.815789in}{0.190924in}}%
\pgfpathlineto{\pgfqpoint{0.814679in}{0.188532in}}%
\pgfpathlineto{\pgfqpoint{0.815722in}{0.184648in}}%
\pgfpathlineto{\pgfqpoint{0.813053in}{0.184321in}}%
\pgfpathlineto{\pgfqpoint{0.810255in}{0.186119in}}%
\pgfpathlineto{\pgfqpoint{0.811163in}{0.190080in}}%
\pgfpathlineto{\pgfqpoint{0.815789in}{0.190924in}}%
\pgfpathclose%
\pgfusepath{fill}%
\end{pgfscope}%
\begin{pgfscope}%
\pgfpathrectangle{\pgfqpoint{0.100000in}{0.100000in}}{\pgfqpoint{2.862985in}{1.832778in}}%
\pgfusepath{clip}%
\pgfsetbuttcap%
\pgfsetmiterjoin%
\definecolor{currentfill}{rgb}{0.859516,0.942330,0.747405}%
\pgfsetfillcolor{currentfill}%
\pgfsetlinewidth{0.000000pt}%
\definecolor{currentstroke}{rgb}{0.000000,0.000000,0.000000}%
\pgfsetstrokecolor{currentstroke}%
\pgfsetstrokeopacity{0.000000}%
\pgfsetdash{}{0pt}%
\pgfpathmoveto{\pgfqpoint{0.803351in}{0.191195in}}%
\pgfpathlineto{\pgfqpoint{0.797461in}{0.188553in}}%
\pgfpathlineto{\pgfqpoint{0.795459in}{0.189503in}}%
\pgfpathlineto{\pgfqpoint{0.799476in}{0.192989in}}%
\pgfpathlineto{\pgfqpoint{0.803351in}{0.191195in}}%
\pgfpathclose%
\pgfusepath{fill}%
\end{pgfscope}%
\begin{pgfscope}%
\pgfpathrectangle{\pgfqpoint{0.100000in}{0.100000in}}{\pgfqpoint{2.862985in}{1.832778in}}%
\pgfusepath{clip}%
\pgfsetbuttcap%
\pgfsetmiterjoin%
\definecolor{currentfill}{rgb}{0.859516,0.942330,0.747405}%
\pgfsetfillcolor{currentfill}%
\pgfsetlinewidth{0.000000pt}%
\definecolor{currentstroke}{rgb}{0.000000,0.000000,0.000000}%
\pgfsetstrokecolor{currentstroke}%
\pgfsetstrokeopacity{0.000000}%
\pgfsetdash{}{0pt}%
\pgfpathmoveto{\pgfqpoint{0.815854in}{0.165422in}}%
\pgfpathlineto{\pgfqpoint{0.814449in}{0.169400in}}%
\pgfpathlineto{\pgfqpoint{0.817419in}{0.170378in}}%
\pgfpathlineto{\pgfqpoint{0.818319in}{0.174668in}}%
\pgfpathlineto{\pgfqpoint{0.821538in}{0.174673in}}%
\pgfpathlineto{\pgfqpoint{0.821524in}{0.177660in}}%
\pgfpathlineto{\pgfqpoint{0.825997in}{0.176953in}}%
\pgfpathlineto{\pgfqpoint{0.826760in}{0.168856in}}%
\pgfpathlineto{\pgfqpoint{0.824165in}{0.163694in}}%
\pgfpathlineto{\pgfqpoint{0.820315in}{0.160469in}}%
\pgfpathlineto{\pgfqpoint{0.815854in}{0.165422in}}%
\pgfpathclose%
\pgfusepath{fill}%
\end{pgfscope}%
\begin{pgfscope}%
\pgfpathrectangle{\pgfqpoint{0.100000in}{0.100000in}}{\pgfqpoint{2.862985in}{1.832778in}}%
\pgfusepath{clip}%
\pgfsetbuttcap%
\pgfsetmiterjoin%
\definecolor{currentfill}{rgb}{0.859516,0.942330,0.747405}%
\pgfsetfillcolor{currentfill}%
\pgfsetlinewidth{0.000000pt}%
\definecolor{currentstroke}{rgb}{0.000000,0.000000,0.000000}%
\pgfsetstrokecolor{currentstroke}%
\pgfsetstrokeopacity{0.000000}%
\pgfsetdash{}{0pt}%
\pgfpathmoveto{\pgfqpoint{0.814114in}{0.168482in}}%
\pgfpathlineto{\pgfqpoint{0.815383in}{0.164715in}}%
\pgfpathlineto{\pgfqpoint{0.812255in}{0.164257in}}%
\pgfpathlineto{\pgfqpoint{0.814114in}{0.168482in}}%
\pgfpathclose%
\pgfusepath{fill}%
\end{pgfscope}%
\begin{pgfscope}%
\pgfpathrectangle{\pgfqpoint{0.100000in}{0.100000in}}{\pgfqpoint{2.862985in}{1.832778in}}%
\pgfusepath{clip}%
\pgfsetbuttcap%
\pgfsetmiterjoin%
\definecolor{currentfill}{rgb}{0.859516,0.942330,0.747405}%
\pgfsetfillcolor{currentfill}%
\pgfsetlinewidth{0.000000pt}%
\definecolor{currentstroke}{rgb}{0.000000,0.000000,0.000000}%
\pgfsetstrokecolor{currentstroke}%
\pgfsetstrokeopacity{0.000000}%
\pgfsetdash{}{0pt}%
\pgfpathmoveto{\pgfqpoint{0.816751in}{0.162942in}}%
\pgfpathlineto{\pgfqpoint{0.816315in}{0.158141in}}%
\pgfpathlineto{\pgfqpoint{0.813015in}{0.158538in}}%
\pgfpathlineto{\pgfqpoint{0.814083in}{0.160858in}}%
\pgfpathlineto{\pgfqpoint{0.816751in}{0.162942in}}%
\pgfpathclose%
\pgfusepath{fill}%
\end{pgfscope}%
\begin{pgfscope}%
\definecolor{textcolor}{rgb}{0.000000,0.000000,0.000000}%
\pgfsetstrokecolor{textcolor}%
\pgfsetfillcolor{textcolor}%
\pgftext[x=1.245194in,y=1.896122in,left,base]{\color{textcolor}\setmainfont{Lato}\rmfamily\fontsize{9.000000}{10.800000}\selectfont November 2023}%
\end{pgfscope}%
\begin{pgfscope}%
\pgfpathrectangle{\pgfqpoint{3.525000in}{0.100000in}}{\pgfqpoint{2.862985in}{1.832778in}}%
\pgfusepath{clip}%
\pgfsetbuttcap%
\pgfsetmiterjoin%
\definecolor{currentfill}{rgb}{0.760477,0.899270,0.588082}%
\pgfsetfillcolor{currentfill}%
\pgfsetlinewidth{0.000000pt}%
\definecolor{currentstroke}{rgb}{0.000000,0.000000,0.000000}%
\pgfsetstrokecolor{currentstroke}%
\pgfsetstrokeopacity{0.000000}%
\pgfsetdash{}{0pt}%
\pgfpathmoveto{\pgfqpoint{4.480609in}{0.529755in}}%
\pgfpathlineto{\pgfqpoint{4.472380in}{0.529889in}}%
\pgfpathlineto{\pgfqpoint{4.466780in}{0.536088in}}%
\pgfpathlineto{\pgfqpoint{4.462972in}{0.552550in}}%
\pgfpathlineto{\pgfqpoint{4.471822in}{0.554612in}}%
\pgfpathlineto{\pgfqpoint{4.480813in}{0.552394in}}%
\pgfpathlineto{\pgfqpoint{4.489963in}{0.545909in}}%
\pgfpathlineto{\pgfqpoint{4.489959in}{0.536940in}}%
\pgfpathlineto{\pgfqpoint{4.480609in}{0.529755in}}%
\pgfpathclose%
\pgfusepath{fill}%
\end{pgfscope}%
\begin{pgfscope}%
\pgfpathrectangle{\pgfqpoint{3.525000in}{0.100000in}}{\pgfqpoint{2.862985in}{1.832778in}}%
\pgfusepath{clip}%
\pgfsetbuttcap%
\pgfsetmiterjoin%
\definecolor{currentfill}{rgb}{0.760477,0.899270,0.588082}%
\pgfsetfillcolor{currentfill}%
\pgfsetlinewidth{0.000000pt}%
\definecolor{currentstroke}{rgb}{0.000000,0.000000,0.000000}%
\pgfsetstrokecolor{currentstroke}%
\pgfsetstrokeopacity{0.000000}%
\pgfsetdash{}{0pt}%
\pgfpathmoveto{\pgfqpoint{4.530694in}{0.432023in}}%
\pgfpathlineto{\pgfqpoint{4.521488in}{0.435528in}}%
\pgfpathlineto{\pgfqpoint{4.521446in}{0.442404in}}%
\pgfpathlineto{\pgfqpoint{4.510326in}{0.448219in}}%
\pgfpathlineto{\pgfqpoint{4.511680in}{0.462860in}}%
\pgfpathlineto{\pgfqpoint{4.514901in}{0.469769in}}%
\pgfpathlineto{\pgfqpoint{4.521701in}{0.463410in}}%
\pgfpathlineto{\pgfqpoint{4.532760in}{0.464472in}}%
\pgfpathlineto{\pgfqpoint{4.529967in}{0.444236in}}%
\pgfpathlineto{\pgfqpoint{4.530694in}{0.432023in}}%
\pgfpathclose%
\pgfusepath{fill}%
\end{pgfscope}%
\begin{pgfscope}%
\pgfpathrectangle{\pgfqpoint{3.525000in}{0.100000in}}{\pgfqpoint{2.862985in}{1.832778in}}%
\pgfusepath{clip}%
\pgfsetbuttcap%
\pgfsetmiterjoin%
\definecolor{currentfill}{rgb}{0.760477,0.899270,0.588082}%
\pgfsetfillcolor{currentfill}%
\pgfsetlinewidth{0.000000pt}%
\definecolor{currentstroke}{rgb}{0.000000,0.000000,0.000000}%
\pgfsetstrokecolor{currentstroke}%
\pgfsetstrokeopacity{0.000000}%
\pgfsetdash{}{0pt}%
\pgfpathmoveto{\pgfqpoint{4.569333in}{0.386887in}}%
\pgfpathlineto{\pgfqpoint{4.557651in}{0.388876in}}%
\pgfpathlineto{\pgfqpoint{4.551191in}{0.398877in}}%
\pgfpathlineto{\pgfqpoint{4.550279in}{0.407100in}}%
\pgfpathlineto{\pgfqpoint{4.569333in}{0.386887in}}%
\pgfpathclose%
\pgfusepath{fill}%
\end{pgfscope}%
\begin{pgfscope}%
\pgfpathrectangle{\pgfqpoint{3.525000in}{0.100000in}}{\pgfqpoint{2.862985in}{1.832778in}}%
\pgfusepath{clip}%
\pgfsetbuttcap%
\pgfsetmiterjoin%
\definecolor{currentfill}{rgb}{0.760477,0.899270,0.588082}%
\pgfsetfillcolor{currentfill}%
\pgfsetlinewidth{0.000000pt}%
\definecolor{currentstroke}{rgb}{0.000000,0.000000,0.000000}%
\pgfsetstrokecolor{currentstroke}%
\pgfsetstrokeopacity{0.000000}%
\pgfsetdash{}{0pt}%
\pgfpathmoveto{\pgfqpoint{4.545394in}{0.388631in}}%
\pgfpathlineto{\pgfqpoint{4.551519in}{0.383659in}}%
\pgfpathlineto{\pgfqpoint{4.552848in}{0.375110in}}%
\pgfpathlineto{\pgfqpoint{4.540519in}{0.375583in}}%
\pgfpathlineto{\pgfqpoint{4.545394in}{0.388631in}}%
\pgfpathclose%
\pgfusepath{fill}%
\end{pgfscope}%
\begin{pgfscope}%
\pgfpathrectangle{\pgfqpoint{3.525000in}{0.100000in}}{\pgfqpoint{2.862985in}{1.832778in}}%
\pgfusepath{clip}%
\pgfsetbuttcap%
\pgfsetmiterjoin%
\definecolor{currentfill}{rgb}{0.760477,0.899270,0.588082}%
\pgfsetfillcolor{currentfill}%
\pgfsetlinewidth{0.000000pt}%
\definecolor{currentstroke}{rgb}{0.000000,0.000000,0.000000}%
\pgfsetstrokecolor{currentstroke}%
\pgfsetstrokeopacity{0.000000}%
\pgfsetdash{}{0pt}%
\pgfpathmoveto{\pgfqpoint{4.573521in}{0.340880in}}%
\pgfpathlineto{\pgfqpoint{4.561217in}{0.345303in}}%
\pgfpathlineto{\pgfqpoint{4.564744in}{0.356026in}}%
\pgfpathlineto{\pgfqpoint{4.559586in}{0.366952in}}%
\pgfpathlineto{\pgfqpoint{4.561206in}{0.374463in}}%
\pgfpathlineto{\pgfqpoint{4.571107in}{0.377326in}}%
\pgfpathlineto{\pgfqpoint{4.570916in}{0.365292in}}%
\pgfpathlineto{\pgfqpoint{4.582627in}{0.358618in}}%
\pgfpathlineto{\pgfqpoint{4.589326in}{0.342480in}}%
\pgfpathlineto{\pgfqpoint{4.581861in}{0.336266in}}%
\pgfpathlineto{\pgfqpoint{4.573521in}{0.340880in}}%
\pgfpathclose%
\pgfusepath{fill}%
\end{pgfscope}%
\begin{pgfscope}%
\pgfpathrectangle{\pgfqpoint{3.525000in}{0.100000in}}{\pgfqpoint{2.862985in}{1.832778in}}%
\pgfusepath{clip}%
\pgfsetbuttcap%
\pgfsetmiterjoin%
\definecolor{currentfill}{rgb}{0.760477,0.899270,0.588082}%
\pgfsetfillcolor{currentfill}%
\pgfsetlinewidth{0.000000pt}%
\definecolor{currentstroke}{rgb}{0.000000,0.000000,0.000000}%
\pgfsetstrokecolor{currentstroke}%
\pgfsetstrokeopacity{0.000000}%
\pgfsetdash{}{0pt}%
\pgfpathmoveto{\pgfqpoint{4.529084in}{0.231768in}}%
\pgfpathlineto{\pgfqpoint{4.523222in}{0.243877in}}%
\pgfpathlineto{\pgfqpoint{4.525818in}{0.252022in}}%
\pgfpathlineto{\pgfqpoint{4.537315in}{0.263296in}}%
\pgfpathlineto{\pgfqpoint{4.538404in}{0.271938in}}%
\pgfpathlineto{\pgfqpoint{4.545928in}{0.291072in}}%
\pgfpathlineto{\pgfqpoint{4.565718in}{0.294341in}}%
\pgfpathlineto{\pgfqpoint{4.569506in}{0.305541in}}%
\pgfpathlineto{\pgfqpoint{4.578232in}{0.301341in}}%
\pgfpathlineto{\pgfqpoint{4.595961in}{0.269908in}}%
\pgfpathlineto{\pgfqpoint{4.596185in}{0.259645in}}%
\pgfpathlineto{\pgfqpoint{4.591173in}{0.250763in}}%
\pgfpathlineto{\pgfqpoint{4.597532in}{0.230369in}}%
\pgfpathlineto{\pgfqpoint{4.593089in}{0.227081in}}%
\pgfpathlineto{\pgfqpoint{4.577955in}{0.228037in}}%
\pgfpathlineto{\pgfqpoint{4.559756in}{0.234119in}}%
\pgfpathlineto{\pgfqpoint{4.544460in}{0.236631in}}%
\pgfpathlineto{\pgfqpoint{4.540375in}{0.233063in}}%
\pgfpathlineto{\pgfqpoint{4.529084in}{0.231768in}}%
\pgfpathclose%
\pgfusepath{fill}%
\end{pgfscope}%
\begin{pgfscope}%
\pgfpathrectangle{\pgfqpoint{3.525000in}{0.100000in}}{\pgfqpoint{2.862985in}{1.832778in}}%
\pgfusepath{clip}%
\pgfsetbuttcap%
\pgfsetmiterjoin%
\definecolor{currentfill}{rgb}{0.859516,0.942330,0.747405}%
\pgfsetfillcolor{currentfill}%
\pgfsetlinewidth{0.000000pt}%
\definecolor{currentstroke}{rgb}{0.000000,0.000000,0.000000}%
\pgfsetstrokecolor{currentstroke}%
\pgfsetstrokeopacity{0.000000}%
\pgfsetdash{}{0pt}%
\pgfpathmoveto{\pgfqpoint{4.205273in}{1.824988in}}%
\pgfpathlineto{\pgfqpoint{4.190426in}{1.766159in}}%
\pgfpathlineto{\pgfqpoint{4.179869in}{1.724007in}}%
\pgfpathlineto{\pgfqpoint{4.167317in}{1.673299in}}%
\pgfpathlineto{\pgfqpoint{4.164852in}{1.660801in}}%
\pgfpathlineto{\pgfqpoint{4.166462in}{1.649338in}}%
\pgfpathlineto{\pgfqpoint{4.164329in}{1.638779in}}%
\pgfpathlineto{\pgfqpoint{4.076709in}{1.661643in}}%
\pgfpathlineto{\pgfqpoint{4.068794in}{1.658987in}}%
\pgfpathlineto{\pgfqpoint{4.062021in}{1.661292in}}%
\pgfpathlineto{\pgfqpoint{4.030660in}{1.662045in}}%
\pgfpathlineto{\pgfqpoint{4.020053in}{1.659055in}}%
\pgfpathlineto{\pgfqpoint{4.009520in}{1.660065in}}%
\pgfpathlineto{\pgfqpoint{4.005055in}{1.664746in}}%
\pgfpathlineto{\pgfqpoint{3.976979in}{1.663725in}}%
\pgfpathlineto{\pgfqpoint{3.972545in}{1.671374in}}%
\pgfpathlineto{\pgfqpoint{3.964127in}{1.675404in}}%
\pgfpathlineto{\pgfqpoint{3.951880in}{1.677860in}}%
\pgfpathlineto{\pgfqpoint{3.930740in}{1.674264in}}%
\pgfpathlineto{\pgfqpoint{3.922937in}{1.677787in}}%
\pgfpathlineto{\pgfqpoint{3.910905in}{1.687156in}}%
\pgfpathlineto{\pgfqpoint{3.914522in}{1.705519in}}%
\pgfpathlineto{\pgfqpoint{3.913254in}{1.714897in}}%
\pgfpathlineto{\pgfqpoint{3.903718in}{1.724852in}}%
\pgfpathlineto{\pgfqpoint{3.897610in}{1.724234in}}%
\pgfpathlineto{\pgfqpoint{3.893229in}{1.734304in}}%
\pgfpathlineto{\pgfqpoint{3.882849in}{1.738374in}}%
\pgfpathlineto{\pgfqpoint{3.875305in}{1.737856in}}%
\pgfpathlineto{\pgfqpoint{3.872864in}{1.748198in}}%
\pgfpathlineto{\pgfqpoint{3.880380in}{1.747098in}}%
\pgfpathlineto{\pgfqpoint{3.879783in}{1.761472in}}%
\pgfpathlineto{\pgfqpoint{3.876507in}{1.770008in}}%
\pgfpathlineto{\pgfqpoint{3.877951in}{1.780932in}}%
\pgfpathlineto{\pgfqpoint{3.881884in}{1.789149in}}%
\pgfpathlineto{\pgfqpoint{3.881662in}{1.804741in}}%
\pgfpathlineto{\pgfqpoint{3.879563in}{1.810398in}}%
\pgfpathlineto{\pgfqpoint{3.883188in}{1.828548in}}%
\pgfpathlineto{\pgfqpoint{3.882142in}{1.840246in}}%
\pgfpathlineto{\pgfqpoint{3.878517in}{1.845854in}}%
\pgfpathlineto{\pgfqpoint{3.879070in}{1.864192in}}%
\pgfpathlineto{\pgfqpoint{3.884342in}{1.877678in}}%
\pgfpathlineto{\pgfqpoint{3.890077in}{1.874378in}}%
\pgfpathlineto{\pgfqpoint{3.909064in}{1.854764in}}%
\pgfpathlineto{\pgfqpoint{3.932071in}{1.844031in}}%
\pgfpathlineto{\pgfqpoint{3.943852in}{1.842761in}}%
\pgfpathlineto{\pgfqpoint{3.954951in}{1.838191in}}%
\pgfpathlineto{\pgfqpoint{3.958101in}{1.822814in}}%
\pgfpathlineto{\pgfqpoint{3.948375in}{1.819901in}}%
\pgfpathlineto{\pgfqpoint{3.941929in}{1.807782in}}%
\pgfpathlineto{\pgfqpoint{3.949497in}{1.808455in}}%
\pgfpathlineto{\pgfqpoint{3.952549in}{1.813895in}}%
\pgfpathlineto{\pgfqpoint{3.963221in}{1.820697in}}%
\pgfpathlineto{\pgfqpoint{3.962670in}{1.810596in}}%
\pgfpathlineto{\pgfqpoint{3.955544in}{1.808967in}}%
\pgfpathlineto{\pgfqpoint{3.956727in}{1.795983in}}%
\pgfpathlineto{\pgfqpoint{3.948050in}{1.782976in}}%
\pgfpathlineto{\pgfqpoint{3.942496in}{1.789161in}}%
\pgfpathlineto{\pgfqpoint{3.925991in}{1.785798in}}%
\pgfpathlineto{\pgfqpoint{3.925123in}{1.778232in}}%
\pgfpathlineto{\pgfqpoint{3.930813in}{1.773622in}}%
\pgfpathlineto{\pgfqpoint{3.939493in}{1.773169in}}%
\pgfpathlineto{\pgfqpoint{3.951495in}{1.783019in}}%
\pgfpathlineto{\pgfqpoint{3.960994in}{1.783837in}}%
\pgfpathlineto{\pgfqpoint{3.961358in}{1.794741in}}%
\pgfpathlineto{\pgfqpoint{3.966254in}{1.810756in}}%
\pgfpathlineto{\pgfqpoint{3.977856in}{1.822666in}}%
\pgfpathlineto{\pgfqpoint{3.973988in}{1.831891in}}%
\pgfpathlineto{\pgfqpoint{3.976630in}{1.841823in}}%
\pgfpathlineto{\pgfqpoint{3.971483in}{1.851722in}}%
\pgfpathlineto{\pgfqpoint{3.978818in}{1.864381in}}%
\pgfpathlineto{\pgfqpoint{3.979917in}{1.871996in}}%
\pgfpathlineto{\pgfqpoint{3.973548in}{1.877001in}}%
\pgfpathlineto{\pgfqpoint{3.974609in}{1.889804in}}%
\pgfpathlineto{\pgfqpoint{4.050918in}{1.866620in}}%
\pgfpathlineto{\pgfqpoint{4.131952in}{1.843896in}}%
\pgfpathlineto{\pgfqpoint{4.205273in}{1.824988in}}%
\pgfpathclose%
\pgfusepath{fill}%
\end{pgfscope}%
\begin{pgfscope}%
\pgfpathrectangle{\pgfqpoint{3.525000in}{0.100000in}}{\pgfqpoint{2.862985in}{1.832778in}}%
\pgfusepath{clip}%
\pgfsetbuttcap%
\pgfsetmiterjoin%
\definecolor{currentfill}{rgb}{0.859516,0.942330,0.747405}%
\pgfsetfillcolor{currentfill}%
\pgfsetlinewidth{0.000000pt}%
\definecolor{currentstroke}{rgb}{0.000000,0.000000,0.000000}%
\pgfsetstrokecolor{currentstroke}%
\pgfsetstrokeopacity{0.000000}%
\pgfsetdash{}{0pt}%
\pgfpathmoveto{\pgfqpoint{3.962056in}{1.845512in}}%
\pgfpathlineto{\pgfqpoint{3.965813in}{1.831026in}}%
\pgfpathlineto{\pgfqpoint{3.971706in}{1.826214in}}%
\pgfpathlineto{\pgfqpoint{3.967004in}{1.820089in}}%
\pgfpathlineto{\pgfqpoint{3.962421in}{1.829064in}}%
\pgfpathlineto{\pgfqpoint{3.962056in}{1.845512in}}%
\pgfpathclose%
\pgfusepath{fill}%
\end{pgfscope}%
\begin{pgfscope}%
\pgfpathrectangle{\pgfqpoint{3.525000in}{0.100000in}}{\pgfqpoint{2.862985in}{1.832778in}}%
\pgfusepath{clip}%
\pgfsetbuttcap%
\pgfsetmiterjoin%
\definecolor{currentfill}{rgb}{0.382699,0.640369,0.181084}%
\pgfsetfillcolor{currentfill}%
\pgfsetlinewidth{0.000000pt}%
\definecolor{currentstroke}{rgb}{0.000000,0.000000,0.000000}%
\pgfsetstrokecolor{currentstroke}%
\pgfsetstrokeopacity{0.000000}%
\pgfsetdash{}{0pt}%
\pgfpathmoveto{\pgfqpoint{4.244624in}{1.815480in}}%
\pgfpathlineto{\pgfqpoint{4.326194in}{1.797175in}}%
\pgfpathlineto{\pgfqpoint{4.403016in}{1.781647in}}%
\pgfpathlineto{\pgfqpoint{4.462121in}{1.770817in}}%
\pgfpathlineto{\pgfqpoint{4.513650in}{1.762160in}}%
\pgfpathlineto{\pgfqpoint{4.565293in}{1.754214in}}%
\pgfpathlineto{\pgfqpoint{4.609272in}{1.748018in}}%
\pgfpathlineto{\pgfqpoint{4.653320in}{1.742337in}}%
\pgfpathlineto{\pgfqpoint{4.697432in}{1.737171in}}%
\pgfpathlineto{\pgfqpoint{4.739001in}{1.732781in}}%
\pgfpathlineto{\pgfqpoint{4.733284in}{1.669428in}}%
\pgfpathlineto{\pgfqpoint{4.724775in}{1.583815in}}%
\pgfpathlineto{\pgfqpoint{4.720312in}{1.539783in}}%
\pgfpathlineto{\pgfqpoint{4.713893in}{1.480441in}}%
\pgfpathlineto{\pgfqpoint{4.668453in}{1.485403in}}%
\pgfpathlineto{\pgfqpoint{4.616429in}{1.491308in}}%
\pgfpathlineto{\pgfqpoint{4.544194in}{1.501195in}}%
\pgfpathlineto{\pgfqpoint{4.511931in}{1.505797in}}%
\pgfpathlineto{\pgfqpoint{4.432445in}{1.518274in}}%
\pgfpathlineto{\pgfqpoint{4.405083in}{1.523280in}}%
\pgfpathlineto{\pgfqpoint{4.399380in}{1.490854in}}%
\pgfpathlineto{\pgfqpoint{4.396262in}{1.493168in}}%
\pgfpathlineto{\pgfqpoint{4.390403in}{1.508767in}}%
\pgfpathlineto{\pgfqpoint{4.382871in}{1.507806in}}%
\pgfpathlineto{\pgfqpoint{4.377480in}{1.499920in}}%
\pgfpathlineto{\pgfqpoint{4.366293in}{1.498866in}}%
\pgfpathlineto{\pgfqpoint{4.364045in}{1.502638in}}%
\pgfpathlineto{\pgfqpoint{4.353566in}{1.502181in}}%
\pgfpathlineto{\pgfqpoint{4.348297in}{1.505835in}}%
\pgfpathlineto{\pgfqpoint{4.340931in}{1.500183in}}%
\pgfpathlineto{\pgfqpoint{4.323038in}{1.505269in}}%
\pgfpathlineto{\pgfqpoint{4.315230in}{1.502538in}}%
\pgfpathlineto{\pgfqpoint{4.311029in}{1.514658in}}%
\pgfpathlineto{\pgfqpoint{4.310833in}{1.526190in}}%
\pgfpathlineto{\pgfqpoint{4.298499in}{1.534484in}}%
\pgfpathlineto{\pgfqpoint{4.300510in}{1.546436in}}%
\pgfpathlineto{\pgfqpoint{4.291728in}{1.566180in}}%
\pgfpathlineto{\pgfqpoint{4.292319in}{1.582964in}}%
\pgfpathlineto{\pgfqpoint{4.284772in}{1.591178in}}%
\pgfpathlineto{\pgfqpoint{4.277690in}{1.583911in}}%
\pgfpathlineto{\pgfqpoint{4.268050in}{1.579980in}}%
\pgfpathlineto{\pgfqpoint{4.260712in}{1.588082in}}%
\pgfpathlineto{\pgfqpoint{4.261905in}{1.601019in}}%
\pgfpathlineto{\pgfqpoint{4.270705in}{1.606116in}}%
\pgfpathlineto{\pgfqpoint{4.268366in}{1.614312in}}%
\pgfpathlineto{\pgfqpoint{4.283663in}{1.653985in}}%
\pgfpathlineto{\pgfqpoint{4.271601in}{1.654979in}}%
\pgfpathlineto{\pgfqpoint{4.270361in}{1.662098in}}%
\pgfpathlineto{\pgfqpoint{4.261575in}{1.668238in}}%
\pgfpathlineto{\pgfqpoint{4.262137in}{1.675101in}}%
\pgfpathlineto{\pgfqpoint{4.257502in}{1.680433in}}%
\pgfpathlineto{\pgfqpoint{4.249957in}{1.699791in}}%
\pgfpathlineto{\pgfqpoint{4.243063in}{1.703739in}}%
\pgfpathlineto{\pgfqpoint{4.238126in}{1.720446in}}%
\pgfpathlineto{\pgfqpoint{4.239291in}{1.731551in}}%
\pgfpathlineto{\pgfqpoint{4.230080in}{1.752004in}}%
\pgfpathlineto{\pgfqpoint{4.244624in}{1.815480in}}%
\pgfpathclose%
\pgfusepath{fill}%
\end{pgfscope}%
\begin{pgfscope}%
\pgfpathrectangle{\pgfqpoint{3.525000in}{0.100000in}}{\pgfqpoint{2.862985in}{1.832778in}}%
\pgfusepath{clip}%
\pgfsetbuttcap%
\pgfsetmiterjoin%
\definecolor{currentfill}{rgb}{0.524337,0.754325,0.286736}%
\pgfsetfillcolor{currentfill}%
\pgfsetlinewidth{0.000000pt}%
\definecolor{currentstroke}{rgb}{0.000000,0.000000,0.000000}%
\pgfsetstrokecolor{currentstroke}%
\pgfsetstrokeopacity{0.000000}%
\pgfsetdash{}{0pt}%
\pgfpathmoveto{\pgfqpoint{6.213801in}{1.489012in}}%
\pgfpathlineto{\pgfqpoint{6.212168in}{1.495954in}}%
\pgfpathlineto{\pgfqpoint{6.203094in}{1.502035in}}%
\pgfpathlineto{\pgfqpoint{6.179117in}{1.580468in}}%
\pgfpathlineto{\pgfqpoint{6.166192in}{1.618318in}}%
\pgfpathlineto{\pgfqpoint{6.176583in}{1.629222in}}%
\pgfpathlineto{\pgfqpoint{6.186683in}{1.655908in}}%
\pgfpathlineto{\pgfqpoint{6.191720in}{1.664437in}}%
\pgfpathlineto{\pgfqpoint{6.188152in}{1.668018in}}%
\pgfpathlineto{\pgfqpoint{6.186957in}{1.691707in}}%
\pgfpathlineto{\pgfqpoint{6.191463in}{1.698969in}}%
\pgfpathlineto{\pgfqpoint{6.189488in}{1.715849in}}%
\pgfpathlineto{\pgfqpoint{6.207474in}{1.771228in}}%
\pgfpathlineto{\pgfqpoint{6.215578in}{1.771530in}}%
\pgfpathlineto{\pgfqpoint{6.218942in}{1.761638in}}%
\pgfpathlineto{\pgfqpoint{6.226144in}{1.758798in}}%
\pgfpathlineto{\pgfqpoint{6.239749in}{1.770433in}}%
\pgfpathlineto{\pgfqpoint{6.250417in}{1.777306in}}%
\pgfpathlineto{\pgfqpoint{6.273922in}{1.765204in}}%
\pgfpathlineto{\pgfqpoint{6.295157in}{1.698008in}}%
\pgfpathlineto{\pgfqpoint{6.299194in}{1.681476in}}%
\pgfpathlineto{\pgfqpoint{6.308495in}{1.679537in}}%
\pgfpathlineto{\pgfqpoint{6.321110in}{1.668309in}}%
\pgfpathlineto{\pgfqpoint{6.320396in}{1.661767in}}%
\pgfpathlineto{\pgfqpoint{6.328989in}{1.654013in}}%
\pgfpathlineto{\pgfqpoint{6.335970in}{1.658462in}}%
\pgfpathlineto{\pgfqpoint{6.350577in}{1.641426in}}%
\pgfpathlineto{\pgfqpoint{6.344049in}{1.627792in}}%
\pgfpathlineto{\pgfqpoint{6.335294in}{1.627522in}}%
\pgfpathlineto{\pgfqpoint{6.328271in}{1.615351in}}%
\pgfpathlineto{\pgfqpoint{6.319762in}{1.613625in}}%
\pgfpathlineto{\pgfqpoint{6.313516in}{1.607145in}}%
\pgfpathlineto{\pgfqpoint{6.294860in}{1.600195in}}%
\pgfpathlineto{\pgfqpoint{6.283568in}{1.588995in}}%
\pgfpathlineto{\pgfqpoint{6.277832in}{1.597034in}}%
\pgfpathlineto{\pgfqpoint{6.272620in}{1.591272in}}%
\pgfpathlineto{\pgfqpoint{6.274047in}{1.567979in}}%
\pgfpathlineto{\pgfqpoint{6.269920in}{1.558753in}}%
\pgfpathlineto{\pgfqpoint{6.260918in}{1.561273in}}%
\pgfpathlineto{\pgfqpoint{6.259465in}{1.551814in}}%
\pgfpathlineto{\pgfqpoint{6.255576in}{1.547924in}}%
\pgfpathlineto{\pgfqpoint{6.247791in}{1.548872in}}%
\pgfpathlineto{\pgfqpoint{6.248281in}{1.540028in}}%
\pgfpathlineto{\pgfqpoint{6.236482in}{1.542540in}}%
\pgfpathlineto{\pgfqpoint{6.230038in}{1.530287in}}%
\pgfpathlineto{\pgfqpoint{6.232472in}{1.523875in}}%
\pgfpathlineto{\pgfqpoint{6.228639in}{1.513212in}}%
\pgfpathlineto{\pgfqpoint{6.222604in}{1.505369in}}%
\pgfpathlineto{\pgfqpoint{6.221086in}{1.488995in}}%
\pgfpathlineto{\pgfqpoint{6.213801in}{1.489012in}}%
\pgfpathclose%
\pgfusepath{fill}%
\end{pgfscope}%
\begin{pgfscope}%
\pgfpathrectangle{\pgfqpoint{3.525000in}{0.100000in}}{\pgfqpoint{2.862985in}{1.832778in}}%
\pgfusepath{clip}%
\pgfsetbuttcap%
\pgfsetmiterjoin%
\definecolor{currentfill}{rgb}{0.524337,0.754325,0.286736}%
\pgfsetfillcolor{currentfill}%
\pgfsetlinewidth{0.000000pt}%
\definecolor{currentstroke}{rgb}{0.000000,0.000000,0.000000}%
\pgfsetstrokecolor{currentstroke}%
\pgfsetstrokeopacity{0.000000}%
\pgfsetdash{}{0pt}%
\pgfpathmoveto{\pgfqpoint{6.298211in}{1.595412in}}%
\pgfpathlineto{\pgfqpoint{6.303538in}{1.600991in}}%
\pgfpathlineto{\pgfqpoint{6.308609in}{1.595747in}}%
\pgfpathlineto{\pgfqpoint{6.299500in}{1.588798in}}%
\pgfpathlineto{\pgfqpoint{6.298211in}{1.595412in}}%
\pgfpathclose%
\pgfusepath{fill}%
\end{pgfscope}%
\begin{pgfscope}%
\pgfpathrectangle{\pgfqpoint{3.525000in}{0.100000in}}{\pgfqpoint{2.862985in}{1.832778in}}%
\pgfusepath{clip}%
\pgfsetbuttcap%
\pgfsetmiterjoin%
\definecolor{currentfill}{rgb}{0.263975,0.526567,0.121415}%
\pgfsetfillcolor{currentfill}%
\pgfsetlinewidth{0.000000pt}%
\definecolor{currentstroke}{rgb}{0.000000,0.000000,0.000000}%
\pgfsetstrokecolor{currentstroke}%
\pgfsetstrokeopacity{0.000000}%
\pgfsetdash{}{0pt}%
\pgfpathmoveto{\pgfqpoint{4.720312in}{1.539783in}}%
\pgfpathlineto{\pgfqpoint{4.724775in}{1.583815in}}%
\pgfpathlineto{\pgfqpoint{4.733284in}{1.669428in}}%
\pgfpathlineto{\pgfqpoint{4.739001in}{1.732781in}}%
\pgfpathlineto{\pgfqpoint{4.785821in}{1.728389in}}%
\pgfpathlineto{\pgfqpoint{4.845718in}{1.723620in}}%
\pgfpathlineto{\pgfqpoint{4.900466in}{1.720093in}}%
\pgfpathlineto{\pgfqpoint{4.950039in}{1.717585in}}%
\pgfpathlineto{\pgfqpoint{5.024019in}{1.715047in}}%
\pgfpathlineto{\pgfqpoint{5.029068in}{1.694197in}}%
\pgfpathlineto{\pgfqpoint{5.027261in}{1.684232in}}%
\pgfpathlineto{\pgfqpoint{5.026666in}{1.663751in}}%
\pgfpathlineto{\pgfqpoint{5.030101in}{1.648413in}}%
\pgfpathlineto{\pgfqpoint{5.037970in}{1.625784in}}%
\pgfpathlineto{\pgfqpoint{5.037952in}{1.597294in}}%
\pgfpathlineto{\pgfqpoint{5.039408in}{1.564197in}}%
\pgfpathlineto{\pgfqpoint{5.041408in}{1.555277in}}%
\pgfpathlineto{\pgfqpoint{5.047257in}{1.545485in}}%
\pgfpathlineto{\pgfqpoint{5.049186in}{1.530226in}}%
\pgfpathlineto{\pgfqpoint{5.048353in}{1.520039in}}%
\pgfpathlineto{\pgfqpoint{4.986333in}{1.521383in}}%
\pgfpathlineto{\pgfqpoint{4.941215in}{1.523670in}}%
\pgfpathlineto{\pgfqpoint{4.875074in}{1.527191in}}%
\pgfpathlineto{\pgfqpoint{4.809849in}{1.531824in}}%
\pgfpathlineto{\pgfqpoint{4.766407in}{1.535329in}}%
\pgfpathlineto{\pgfqpoint{4.720312in}{1.539783in}}%
\pgfpathclose%
\pgfusepath{fill}%
\end{pgfscope}%
\begin{pgfscope}%
\pgfpathrectangle{\pgfqpoint{3.525000in}{0.100000in}}{\pgfqpoint{2.862985in}{1.832778in}}%
\pgfusepath{clip}%
\pgfsetbuttcap%
\pgfsetmiterjoin%
\definecolor{currentfill}{rgb}{0.299039,0.569012,0.128797}%
\pgfsetfillcolor{currentfill}%
\pgfsetlinewidth{0.000000pt}%
\definecolor{currentstroke}{rgb}{0.000000,0.000000,0.000000}%
\pgfsetstrokecolor{currentstroke}%
\pgfsetstrokeopacity{0.000000}%
\pgfsetdash{}{0pt}%
\pgfpathmoveto{\pgfqpoint{4.701565in}{1.355506in}}%
\pgfpathlineto{\pgfqpoint{4.708725in}{1.429317in}}%
\pgfpathlineto{\pgfqpoint{4.713893in}{1.480441in}}%
\pgfpathlineto{\pgfqpoint{4.720312in}{1.539783in}}%
\pgfpathlineto{\pgfqpoint{4.766407in}{1.535329in}}%
\pgfpathlineto{\pgfqpoint{4.809849in}{1.531824in}}%
\pgfpathlineto{\pgfqpoint{4.875074in}{1.527191in}}%
\pgfpathlineto{\pgfqpoint{4.941215in}{1.523670in}}%
\pgfpathlineto{\pgfqpoint{4.986333in}{1.521383in}}%
\pgfpathlineto{\pgfqpoint{5.048353in}{1.520039in}}%
\pgfpathlineto{\pgfqpoint{5.044154in}{1.507786in}}%
\pgfpathlineto{\pgfqpoint{5.035772in}{1.498167in}}%
\pgfpathlineto{\pgfqpoint{5.042191in}{1.487092in}}%
\pgfpathlineto{\pgfqpoint{5.049268in}{1.484725in}}%
\pgfpathlineto{\pgfqpoint{5.052626in}{1.478363in}}%
\pgfpathlineto{\pgfqpoint{5.051855in}{1.431914in}}%
\pgfpathlineto{\pgfqpoint{5.050565in}{1.366544in}}%
\pgfpathlineto{\pgfqpoint{5.045791in}{1.349375in}}%
\pgfpathlineto{\pgfqpoint{5.050061in}{1.339878in}}%
\pgfpathlineto{\pgfqpoint{5.041489in}{1.319470in}}%
\pgfpathlineto{\pgfqpoint{5.050519in}{1.303018in}}%
\pgfpathlineto{\pgfqpoint{5.042847in}{1.304277in}}%
\pgfpathlineto{\pgfqpoint{5.037612in}{1.314513in}}%
\pgfpathlineto{\pgfqpoint{5.018930in}{1.321529in}}%
\pgfpathlineto{\pgfqpoint{5.007148in}{1.327701in}}%
\pgfpathlineto{\pgfqpoint{4.987377in}{1.328214in}}%
\pgfpathlineto{\pgfqpoint{4.980515in}{1.322591in}}%
\pgfpathlineto{\pgfqpoint{4.958137in}{1.333663in}}%
\pgfpathlineto{\pgfqpoint{4.956420in}{1.337164in}}%
\pgfpathlineto{\pgfqpoint{4.878322in}{1.340860in}}%
\pgfpathlineto{\pgfqpoint{4.830874in}{1.343580in}}%
\pgfpathlineto{\pgfqpoint{4.791682in}{1.346649in}}%
\pgfpathlineto{\pgfqpoint{4.726922in}{1.352768in}}%
\pgfpathlineto{\pgfqpoint{4.701565in}{1.355506in}}%
\pgfpathclose%
\pgfusepath{fill}%
\end{pgfscope}%
\begin{pgfscope}%
\pgfpathrectangle{\pgfqpoint{3.525000in}{0.100000in}}{\pgfqpoint{2.862985in}{1.832778in}}%
\pgfusepath{clip}%
\pgfsetbuttcap%
\pgfsetmiterjoin%
\definecolor{currentfill}{rgb}{0.690888,0.862438,0.488351}%
\pgfsetfillcolor{currentfill}%
\pgfsetlinewidth{0.000000pt}%
\definecolor{currentstroke}{rgb}{0.000000,0.000000,0.000000}%
\pgfsetstrokecolor{currentstroke}%
\pgfsetstrokeopacity{0.000000}%
\pgfsetdash{}{0pt}%
\pgfpathmoveto{\pgfqpoint{4.689283in}{1.230524in}}%
\pgfpathlineto{\pgfqpoint{4.647647in}{1.234326in}}%
\pgfpathlineto{\pgfqpoint{4.556892in}{1.245527in}}%
\pgfpathlineto{\pgfqpoint{4.507522in}{1.252613in}}%
\pgfpathlineto{\pgfqpoint{4.454571in}{1.260214in}}%
\pgfpathlineto{\pgfqpoint{4.409967in}{1.267352in}}%
\pgfpathlineto{\pgfqpoint{4.361009in}{1.275734in}}%
\pgfpathlineto{\pgfqpoint{4.372141in}{1.337483in}}%
\pgfpathlineto{\pgfqpoint{4.383416in}{1.400800in}}%
\pgfpathlineto{\pgfqpoint{4.399380in}{1.490854in}}%
\pgfpathlineto{\pgfqpoint{4.405083in}{1.523280in}}%
\pgfpathlineto{\pgfqpoint{4.432445in}{1.518274in}}%
\pgfpathlineto{\pgfqpoint{4.511931in}{1.505797in}}%
\pgfpathlineto{\pgfqpoint{4.544194in}{1.501195in}}%
\pgfpathlineto{\pgfqpoint{4.616429in}{1.491308in}}%
\pgfpathlineto{\pgfqpoint{4.668453in}{1.485403in}}%
\pgfpathlineto{\pgfqpoint{4.713893in}{1.480441in}}%
\pgfpathlineto{\pgfqpoint{4.708725in}{1.429317in}}%
\pgfpathlineto{\pgfqpoint{4.701565in}{1.355506in}}%
\pgfpathlineto{\pgfqpoint{4.695419in}{1.292777in}}%
\pgfpathlineto{\pgfqpoint{4.689283in}{1.230524in}}%
\pgfpathclose%
\pgfusepath{fill}%
\end{pgfscope}%
\begin{pgfscope}%
\pgfpathrectangle{\pgfqpoint{3.525000in}{0.100000in}}{\pgfqpoint{2.862985in}{1.832778in}}%
\pgfusepath{clip}%
\pgfsetbuttcap%
\pgfsetmiterjoin%
\definecolor{currentfill}{rgb}{0.428835,0.679123,0.210611}%
\pgfsetfillcolor{currentfill}%
\pgfsetlinewidth{0.000000pt}%
\definecolor{currentstroke}{rgb}{0.000000,0.000000,0.000000}%
\pgfsetstrokecolor{currentstroke}%
\pgfsetstrokeopacity{0.000000}%
\pgfsetdash{}{0pt}%
\pgfpathmoveto{\pgfqpoint{5.450520in}{1.315700in}}%
\pgfpathlineto{\pgfqpoint{5.397758in}{1.311952in}}%
\pgfpathlineto{\pgfqpoint{5.319114in}{1.308600in}}%
\pgfpathlineto{\pgfqpoint{5.316125in}{1.316544in}}%
\pgfpathlineto{\pgfqpoint{5.298686in}{1.322483in}}%
\pgfpathlineto{\pgfqpoint{5.294839in}{1.333702in}}%
\pgfpathlineto{\pgfqpoint{5.293228in}{1.347582in}}%
\pgfpathlineto{\pgfqpoint{5.297158in}{1.354693in}}%
\pgfpathlineto{\pgfqpoint{5.290951in}{1.361517in}}%
\pgfpathlineto{\pgfqpoint{5.289454in}{1.369657in}}%
\pgfpathlineto{\pgfqpoint{5.287453in}{1.387663in}}%
\pgfpathlineto{\pgfqpoint{5.281505in}{1.397444in}}%
\pgfpathlineto{\pgfqpoint{5.270940in}{1.402933in}}%
\pgfpathlineto{\pgfqpoint{5.259444in}{1.411992in}}%
\pgfpathlineto{\pgfqpoint{5.253501in}{1.422710in}}%
\pgfpathlineto{\pgfqpoint{5.242838in}{1.427026in}}%
\pgfpathlineto{\pgfqpoint{5.236570in}{1.434051in}}%
\pgfpathlineto{\pgfqpoint{5.228976in}{1.435243in}}%
\pgfpathlineto{\pgfqpoint{5.215419in}{1.445673in}}%
\pgfpathlineto{\pgfqpoint{5.217621in}{1.457675in}}%
\pgfpathlineto{\pgfqpoint{5.217197in}{1.480500in}}%
\pgfpathlineto{\pgfqpoint{5.221378in}{1.486783in}}%
\pgfpathlineto{\pgfqpoint{5.217621in}{1.496270in}}%
\pgfpathlineto{\pgfqpoint{5.211004in}{1.498105in}}%
\pgfpathlineto{\pgfqpoint{5.211550in}{1.506437in}}%
\pgfpathlineto{\pgfqpoint{5.219761in}{1.519600in}}%
\pgfpathlineto{\pgfqpoint{5.236020in}{1.530009in}}%
\pgfpathlineto{\pgfqpoint{5.235007in}{1.567025in}}%
\pgfpathlineto{\pgfqpoint{5.243143in}{1.572586in}}%
\pgfpathlineto{\pgfqpoint{5.250842in}{1.568882in}}%
\pgfpathlineto{\pgfqpoint{5.266526in}{1.574304in}}%
\pgfpathlineto{\pgfqpoint{5.296034in}{1.587932in}}%
\pgfpathlineto{\pgfqpoint{5.299878in}{1.583713in}}%
\pgfpathlineto{\pgfqpoint{5.294289in}{1.564600in}}%
\pgfpathlineto{\pgfqpoint{5.302604in}{1.568794in}}%
\pgfpathlineto{\pgfqpoint{5.316845in}{1.564606in}}%
\pgfpathlineto{\pgfqpoint{5.325596in}{1.561110in}}%
\pgfpathlineto{\pgfqpoint{5.330502in}{1.550866in}}%
\pgfpathlineto{\pgfqpoint{5.375397in}{1.541198in}}%
\pgfpathlineto{\pgfqpoint{5.388816in}{1.534565in}}%
\pgfpathlineto{\pgfqpoint{5.402460in}{1.534643in}}%
\pgfpathlineto{\pgfqpoint{5.416480in}{1.531907in}}%
\pgfpathlineto{\pgfqpoint{5.425578in}{1.522561in}}%
\pgfpathlineto{\pgfqpoint{5.434272in}{1.517943in}}%
\pgfpathlineto{\pgfqpoint{5.435870in}{1.504623in}}%
\pgfpathlineto{\pgfqpoint{5.433301in}{1.496259in}}%
\pgfpathlineto{\pgfqpoint{5.442991in}{1.495641in}}%
\pgfpathlineto{\pgfqpoint{5.439727in}{1.485941in}}%
\pgfpathlineto{\pgfqpoint{5.442832in}{1.482495in}}%
\pgfpathlineto{\pgfqpoint{5.445919in}{1.473339in}}%
\pgfpathlineto{\pgfqpoint{5.436476in}{1.468491in}}%
\pgfpathlineto{\pgfqpoint{5.430993in}{1.454984in}}%
\pgfpathlineto{\pgfqpoint{5.429273in}{1.445433in}}%
\pgfpathlineto{\pgfqpoint{5.434533in}{1.443798in}}%
\pgfpathlineto{\pgfqpoint{5.441268in}{1.450960in}}%
\pgfpathlineto{\pgfqpoint{5.446994in}{1.463372in}}%
\pgfpathlineto{\pgfqpoint{5.454741in}{1.467683in}}%
\pgfpathlineto{\pgfqpoint{5.460563in}{1.462075in}}%
\pgfpathlineto{\pgfqpoint{5.454832in}{1.445104in}}%
\pgfpathlineto{\pgfqpoint{5.453034in}{1.431928in}}%
\pgfpathlineto{\pgfqpoint{5.454728in}{1.422457in}}%
\pgfpathlineto{\pgfqpoint{5.449406in}{1.417099in}}%
\pgfpathlineto{\pgfqpoint{5.446748in}{1.404007in}}%
\pgfpathlineto{\pgfqpoint{5.448920in}{1.390247in}}%
\pgfpathlineto{\pgfqpoint{5.442640in}{1.369897in}}%
\pgfpathlineto{\pgfqpoint{5.442772in}{1.359728in}}%
\pgfpathlineto{\pgfqpoint{5.447734in}{1.337693in}}%
\pgfpathlineto{\pgfqpoint{5.450951in}{1.333913in}}%
\pgfpathlineto{\pgfqpoint{5.450520in}{1.315700in}}%
\pgfpathclose%
\pgfusepath{fill}%
\end{pgfscope}%
\begin{pgfscope}%
\pgfpathrectangle{\pgfqpoint{3.525000in}{0.100000in}}{\pgfqpoint{2.862985in}{1.832778in}}%
\pgfusepath{clip}%
\pgfsetbuttcap%
\pgfsetmiterjoin%
\definecolor{currentfill}{rgb}{0.428835,0.679123,0.210611}%
\pgfsetfillcolor{currentfill}%
\pgfsetlinewidth{0.000000pt}%
\definecolor{currentstroke}{rgb}{0.000000,0.000000,0.000000}%
\pgfsetstrokecolor{currentstroke}%
\pgfsetstrokeopacity{0.000000}%
\pgfsetdash{}{0pt}%
\pgfpathmoveto{\pgfqpoint{5.470302in}{1.494291in}}%
\pgfpathlineto{\pgfqpoint{5.471480in}{1.485502in}}%
\pgfpathlineto{\pgfqpoint{5.460690in}{1.462343in}}%
\pgfpathlineto{\pgfqpoint{5.455896in}{1.469051in}}%
\pgfpathlineto{\pgfqpoint{5.470302in}{1.494291in}}%
\pgfpathclose%
\pgfusepath{fill}%
\end{pgfscope}%
\begin{pgfscope}%
\pgfpathrectangle{\pgfqpoint{3.525000in}{0.100000in}}{\pgfqpoint{2.862985in}{1.832778in}}%
\pgfusepath{clip}%
\pgfsetbuttcap%
\pgfsetmiterjoin%
\definecolor{currentfill}{rgb}{0.428835,0.679123,0.210611}%
\pgfsetfillcolor{currentfill}%
\pgfsetlinewidth{0.000000pt}%
\definecolor{currentstroke}{rgb}{0.000000,0.000000,0.000000}%
\pgfsetstrokecolor{currentstroke}%
\pgfsetstrokeopacity{0.000000}%
\pgfsetdash{}{0pt}%
\pgfpathmoveto{\pgfqpoint{4.164329in}{1.638779in}}%
\pgfpathlineto{\pgfqpoint{4.166462in}{1.649338in}}%
\pgfpathlineto{\pgfqpoint{4.164852in}{1.660801in}}%
\pgfpathlineto{\pgfqpoint{4.167317in}{1.673299in}}%
\pgfpathlineto{\pgfqpoint{4.179869in}{1.724007in}}%
\pgfpathlineto{\pgfqpoint{4.190426in}{1.766159in}}%
\pgfpathlineto{\pgfqpoint{4.205273in}{1.824988in}}%
\pgfpathlineto{\pgfqpoint{4.244624in}{1.815480in}}%
\pgfpathlineto{\pgfqpoint{4.230080in}{1.752004in}}%
\pgfpathlineto{\pgfqpoint{4.239291in}{1.731551in}}%
\pgfpathlineto{\pgfqpoint{4.238126in}{1.720446in}}%
\pgfpathlineto{\pgfqpoint{4.243063in}{1.703739in}}%
\pgfpathlineto{\pgfqpoint{4.249957in}{1.699791in}}%
\pgfpathlineto{\pgfqpoint{4.257502in}{1.680433in}}%
\pgfpathlineto{\pgfqpoint{4.262137in}{1.675101in}}%
\pgfpathlineto{\pgfqpoint{4.261575in}{1.668238in}}%
\pgfpathlineto{\pgfqpoint{4.270361in}{1.662098in}}%
\pgfpathlineto{\pgfqpoint{4.271601in}{1.654979in}}%
\pgfpathlineto{\pgfqpoint{4.283663in}{1.653985in}}%
\pgfpathlineto{\pgfqpoint{4.268366in}{1.614312in}}%
\pgfpathlineto{\pgfqpoint{4.270705in}{1.606116in}}%
\pgfpathlineto{\pgfqpoint{4.261905in}{1.601019in}}%
\pgfpathlineto{\pgfqpoint{4.260712in}{1.588082in}}%
\pgfpathlineto{\pgfqpoint{4.268050in}{1.579980in}}%
\pgfpathlineto{\pgfqpoint{4.277690in}{1.583911in}}%
\pgfpathlineto{\pgfqpoint{4.284772in}{1.591178in}}%
\pgfpathlineto{\pgfqpoint{4.292319in}{1.582964in}}%
\pgfpathlineto{\pgfqpoint{4.291728in}{1.566180in}}%
\pgfpathlineto{\pgfqpoint{4.300510in}{1.546436in}}%
\pgfpathlineto{\pgfqpoint{4.298499in}{1.534484in}}%
\pgfpathlineto{\pgfqpoint{4.310833in}{1.526190in}}%
\pgfpathlineto{\pgfqpoint{4.311029in}{1.514658in}}%
\pgfpathlineto{\pgfqpoint{4.315230in}{1.502538in}}%
\pgfpathlineto{\pgfqpoint{4.323038in}{1.505269in}}%
\pgfpathlineto{\pgfqpoint{4.340931in}{1.500183in}}%
\pgfpathlineto{\pgfqpoint{4.348297in}{1.505835in}}%
\pgfpathlineto{\pgfqpoint{4.353566in}{1.502181in}}%
\pgfpathlineto{\pgfqpoint{4.364045in}{1.502638in}}%
\pgfpathlineto{\pgfqpoint{4.366293in}{1.498866in}}%
\pgfpathlineto{\pgfqpoint{4.377480in}{1.499920in}}%
\pgfpathlineto{\pgfqpoint{4.382871in}{1.507806in}}%
\pgfpathlineto{\pgfqpoint{4.390403in}{1.508767in}}%
\pgfpathlineto{\pgfqpoint{4.396262in}{1.493168in}}%
\pgfpathlineto{\pgfqpoint{4.399380in}{1.490854in}}%
\pgfpathlineto{\pgfqpoint{4.383416in}{1.400800in}}%
\pgfpathlineto{\pgfqpoint{4.372141in}{1.337483in}}%
\pgfpathlineto{\pgfqpoint{4.283212in}{1.354650in}}%
\pgfpathlineto{\pgfqpoint{4.235177in}{1.364205in}}%
\pgfpathlineto{\pgfqpoint{4.190246in}{1.374085in}}%
\pgfpathlineto{\pgfqpoint{4.148250in}{1.383593in}}%
\pgfpathlineto{\pgfqpoint{4.099647in}{1.395314in}}%
\pgfpathlineto{\pgfqpoint{4.124717in}{1.498158in}}%
\pgfpathlineto{\pgfqpoint{4.126069in}{1.505674in}}%
\pgfpathlineto{\pgfqpoint{4.137217in}{1.525368in}}%
\pgfpathlineto{\pgfqpoint{4.125673in}{1.537205in}}%
\pgfpathlineto{\pgfqpoint{4.128070in}{1.548828in}}%
\pgfpathlineto{\pgfqpoint{4.132464in}{1.551738in}}%
\pgfpathlineto{\pgfqpoint{4.140313in}{1.563715in}}%
\pgfpathlineto{\pgfqpoint{4.151752in}{1.572002in}}%
\pgfpathlineto{\pgfqpoint{4.152377in}{1.578133in}}%
\pgfpathlineto{\pgfqpoint{4.159282in}{1.584259in}}%
\pgfpathlineto{\pgfqpoint{4.165015in}{1.595725in}}%
\pgfpathlineto{\pgfqpoint{4.176774in}{1.607856in}}%
\pgfpathlineto{\pgfqpoint{4.175939in}{1.619838in}}%
\pgfpathlineto{\pgfqpoint{4.167899in}{1.626494in}}%
\pgfpathlineto{\pgfqpoint{4.164329in}{1.638779in}}%
\pgfpathclose%
\pgfusepath{fill}%
\end{pgfscope}%
\begin{pgfscope}%
\pgfpathrectangle{\pgfqpoint{3.525000in}{0.100000in}}{\pgfqpoint{2.862985in}{1.832778in}}%
\pgfusepath{clip}%
\pgfsetbuttcap%
\pgfsetmiterjoin%
\definecolor{currentfill}{rgb}{0.498039,0.737255,0.254902}%
\pgfsetfillcolor{currentfill}%
\pgfsetlinewidth{0.000000pt}%
\definecolor{currentstroke}{rgb}{0.000000,0.000000,0.000000}%
\pgfsetstrokecolor{currentstroke}%
\pgfsetstrokeopacity{0.000000}%
\pgfsetdash{}{0pt}%
\pgfpathmoveto{\pgfqpoint{6.111015in}{1.438080in}}%
\pgfpathlineto{\pgfqpoint{6.107320in}{1.449744in}}%
\pgfpathlineto{\pgfqpoint{6.100463in}{1.485127in}}%
\pgfpathlineto{\pgfqpoint{6.091598in}{1.496013in}}%
\pgfpathlineto{\pgfqpoint{6.083822in}{1.515590in}}%
\pgfpathlineto{\pgfqpoint{6.086375in}{1.530104in}}%
\pgfpathlineto{\pgfqpoint{6.084327in}{1.540801in}}%
\pgfpathlineto{\pgfqpoint{6.077683in}{1.551308in}}%
\pgfpathlineto{\pgfqpoint{6.077410in}{1.562882in}}%
\pgfpathlineto{\pgfqpoint{6.073552in}{1.575365in}}%
\pgfpathlineto{\pgfqpoint{6.108077in}{1.583862in}}%
\pgfpathlineto{\pgfqpoint{6.152922in}{1.595966in}}%
\pgfpathlineto{\pgfqpoint{6.154711in}{1.589025in}}%
\pgfpathlineto{\pgfqpoint{6.151856in}{1.577975in}}%
\pgfpathlineto{\pgfqpoint{6.158559in}{1.569133in}}%
\pgfpathlineto{\pgfqpoint{6.154999in}{1.557930in}}%
\pgfpathlineto{\pgfqpoint{6.140852in}{1.543846in}}%
\pgfpathlineto{\pgfqpoint{6.144741in}{1.533267in}}%
\pgfpathlineto{\pgfqpoint{6.141859in}{1.510906in}}%
\pgfpathlineto{\pgfqpoint{6.137824in}{1.498298in}}%
\pgfpathlineto{\pgfqpoint{6.142972in}{1.462444in}}%
\pgfpathlineto{\pgfqpoint{6.142278in}{1.449830in}}%
\pgfpathlineto{\pgfqpoint{6.147258in}{1.445782in}}%
\pgfpathlineto{\pgfqpoint{6.111015in}{1.438080in}}%
\pgfpathclose%
\pgfusepath{fill}%
\end{pgfscope}%
\begin{pgfscope}%
\pgfpathrectangle{\pgfqpoint{3.525000in}{0.100000in}}{\pgfqpoint{2.862985in}{1.832778in}}%
\pgfusepath{clip}%
\pgfsetbuttcap%
\pgfsetmiterjoin%
\definecolor{currentfill}{rgb}{0.428835,0.679123,0.210611}%
\pgfsetfillcolor{currentfill}%
\pgfsetlinewidth{0.000000pt}%
\definecolor{currentstroke}{rgb}{0.000000,0.000000,0.000000}%
\pgfsetstrokecolor{currentstroke}%
\pgfsetstrokeopacity{0.000000}%
\pgfsetdash{}{0pt}%
\pgfpathmoveto{\pgfqpoint{5.050565in}{1.366544in}}%
\pgfpathlineto{\pgfqpoint{5.051855in}{1.431914in}}%
\pgfpathlineto{\pgfqpoint{5.052626in}{1.478363in}}%
\pgfpathlineto{\pgfqpoint{5.049268in}{1.484725in}}%
\pgfpathlineto{\pgfqpoint{5.042191in}{1.487092in}}%
\pgfpathlineto{\pgfqpoint{5.035772in}{1.498167in}}%
\pgfpathlineto{\pgfqpoint{5.044154in}{1.507786in}}%
\pgfpathlineto{\pgfqpoint{5.048353in}{1.520039in}}%
\pgfpathlineto{\pgfqpoint{5.049186in}{1.530226in}}%
\pgfpathlineto{\pgfqpoint{5.047257in}{1.545485in}}%
\pgfpathlineto{\pgfqpoint{5.041408in}{1.555277in}}%
\pgfpathlineto{\pgfqpoint{5.039408in}{1.564197in}}%
\pgfpathlineto{\pgfqpoint{5.037952in}{1.597294in}}%
\pgfpathlineto{\pgfqpoint{5.037970in}{1.625784in}}%
\pgfpathlineto{\pgfqpoint{5.030101in}{1.648413in}}%
\pgfpathlineto{\pgfqpoint{5.026666in}{1.663751in}}%
\pgfpathlineto{\pgfqpoint{5.027261in}{1.684232in}}%
\pgfpathlineto{\pgfqpoint{5.029068in}{1.694197in}}%
\pgfpathlineto{\pgfqpoint{5.024019in}{1.715047in}}%
\pgfpathlineto{\pgfqpoint{5.058393in}{1.714359in}}%
\pgfpathlineto{\pgfqpoint{5.110604in}{1.713911in}}%
\pgfpathlineto{\pgfqpoint{5.110890in}{1.737609in}}%
\pgfpathlineto{\pgfqpoint{5.124180in}{1.735000in}}%
\pgfpathlineto{\pgfqpoint{5.130549in}{1.706104in}}%
\pgfpathlineto{\pgfqpoint{5.135245in}{1.695714in}}%
\pgfpathlineto{\pgfqpoint{5.146922in}{1.695407in}}%
\pgfpathlineto{\pgfqpoint{5.149536in}{1.691882in}}%
\pgfpathlineto{\pgfqpoint{5.165825in}{1.690320in}}%
\pgfpathlineto{\pgfqpoint{5.168562in}{1.683156in}}%
\pgfpathlineto{\pgfqpoint{5.179788in}{1.684767in}}%
\pgfpathlineto{\pgfqpoint{5.188507in}{1.691467in}}%
\pgfpathlineto{\pgfqpoint{5.203542in}{1.691216in}}%
\pgfpathlineto{\pgfqpoint{5.212845in}{1.685828in}}%
\pgfpathlineto{\pgfqpoint{5.213915in}{1.680774in}}%
\pgfpathlineto{\pgfqpoint{5.222767in}{1.679717in}}%
\pgfpathlineto{\pgfqpoint{5.228544in}{1.665932in}}%
\pgfpathlineto{\pgfqpoint{5.232272in}{1.674408in}}%
\pgfpathlineto{\pgfqpoint{5.242441in}{1.674929in}}%
\pgfpathlineto{\pgfqpoint{5.245013in}{1.668328in}}%
\pgfpathlineto{\pgfqpoint{5.256455in}{1.665315in}}%
\pgfpathlineto{\pgfqpoint{5.262765in}{1.655808in}}%
\pgfpathlineto{\pgfqpoint{5.276759in}{1.658760in}}%
\pgfpathlineto{\pgfqpoint{5.292158in}{1.670427in}}%
\pgfpathlineto{\pgfqpoint{5.297772in}{1.660142in}}%
\pgfpathlineto{\pgfqpoint{5.323039in}{1.662957in}}%
\pgfpathlineto{\pgfqpoint{5.333840in}{1.655898in}}%
\pgfpathlineto{\pgfqpoint{5.340126in}{1.658418in}}%
\pgfpathlineto{\pgfqpoint{5.345167in}{1.654444in}}%
\pgfpathlineto{\pgfqpoint{5.330213in}{1.645016in}}%
\pgfpathlineto{\pgfqpoint{5.308866in}{1.636615in}}%
\pgfpathlineto{\pgfqpoint{5.287725in}{1.619824in}}%
\pgfpathlineto{\pgfqpoint{5.269396in}{1.597742in}}%
\pgfpathlineto{\pgfqpoint{5.255521in}{1.584687in}}%
\pgfpathlineto{\pgfqpoint{5.243367in}{1.575713in}}%
\pgfpathlineto{\pgfqpoint{5.235007in}{1.567025in}}%
\pgfpathlineto{\pgfqpoint{5.236020in}{1.530009in}}%
\pgfpathlineto{\pgfqpoint{5.219761in}{1.519600in}}%
\pgfpathlineto{\pgfqpoint{5.211550in}{1.506437in}}%
\pgfpathlineto{\pgfqpoint{5.211004in}{1.498105in}}%
\pgfpathlineto{\pgfqpoint{5.217621in}{1.496270in}}%
\pgfpathlineto{\pgfqpoint{5.221378in}{1.486783in}}%
\pgfpathlineto{\pgfqpoint{5.217197in}{1.480500in}}%
\pgfpathlineto{\pgfqpoint{5.217621in}{1.457675in}}%
\pgfpathlineto{\pgfqpoint{5.215419in}{1.445673in}}%
\pgfpathlineto{\pgfqpoint{5.228976in}{1.435243in}}%
\pgfpathlineto{\pgfqpoint{5.236570in}{1.434051in}}%
\pgfpathlineto{\pgfqpoint{5.242838in}{1.427026in}}%
\pgfpathlineto{\pgfqpoint{5.253501in}{1.422710in}}%
\pgfpathlineto{\pgfqpoint{5.259444in}{1.411992in}}%
\pgfpathlineto{\pgfqpoint{5.270940in}{1.402933in}}%
\pgfpathlineto{\pgfqpoint{5.281505in}{1.397444in}}%
\pgfpathlineto{\pgfqpoint{5.287453in}{1.387663in}}%
\pgfpathlineto{\pgfqpoint{5.289454in}{1.369657in}}%
\pgfpathlineto{\pgfqpoint{5.233388in}{1.367621in}}%
\pgfpathlineto{\pgfqpoint{5.185597in}{1.366622in}}%
\pgfpathlineto{\pgfqpoint{5.142055in}{1.365982in}}%
\pgfpathlineto{\pgfqpoint{5.095992in}{1.366052in}}%
\pgfpathlineto{\pgfqpoint{5.050565in}{1.366544in}}%
\pgfpathclose%
\pgfusepath{fill}%
\end{pgfscope}%
\begin{pgfscope}%
\pgfpathrectangle{\pgfqpoint{3.525000in}{0.100000in}}{\pgfqpoint{2.862985in}{1.832778in}}%
\pgfusepath{clip}%
\pgfsetbuttcap%
\pgfsetmiterjoin%
\definecolor{currentfill}{rgb}{0.880738,0.951557,0.781546}%
\pgfsetfillcolor{currentfill}%
\pgfsetlinewidth{0.000000pt}%
\definecolor{currentstroke}{rgb}{0.000000,0.000000,0.000000}%
\pgfsetstrokecolor{currentstroke}%
\pgfsetstrokeopacity{0.000000}%
\pgfsetdash{}{0pt}%
\pgfpathmoveto{\pgfqpoint{3.777592in}{1.488751in}}%
\pgfpathlineto{\pgfqpoint{3.773152in}{1.496921in}}%
\pgfpathlineto{\pgfqpoint{3.775987in}{1.517893in}}%
\pgfpathlineto{\pgfqpoint{3.781459in}{1.528853in}}%
\pgfpathlineto{\pgfqpoint{3.778724in}{1.543679in}}%
\pgfpathlineto{\pgfqpoint{3.784432in}{1.549914in}}%
\pgfpathlineto{\pgfqpoint{3.794847in}{1.566720in}}%
\pgfpathlineto{\pgfqpoint{3.803677in}{1.576864in}}%
\pgfpathlineto{\pgfqpoint{3.816585in}{1.598959in}}%
\pgfpathlineto{\pgfqpoint{3.826485in}{1.623028in}}%
\pgfpathlineto{\pgfqpoint{3.836995in}{1.645601in}}%
\pgfpathlineto{\pgfqpoint{3.839119in}{1.655029in}}%
\pgfpathlineto{\pgfqpoint{3.853617in}{1.682000in}}%
\pgfpathlineto{\pgfqpoint{3.856409in}{1.693845in}}%
\pgfpathlineto{\pgfqpoint{3.862582in}{1.706284in}}%
\pgfpathlineto{\pgfqpoint{3.864788in}{1.721458in}}%
\pgfpathlineto{\pgfqpoint{3.876573in}{1.728857in}}%
\pgfpathlineto{\pgfqpoint{3.890546in}{1.732587in}}%
\pgfpathlineto{\pgfqpoint{3.897610in}{1.724234in}}%
\pgfpathlineto{\pgfqpoint{3.903718in}{1.724852in}}%
\pgfpathlineto{\pgfqpoint{3.913254in}{1.714897in}}%
\pgfpathlineto{\pgfqpoint{3.914522in}{1.705519in}}%
\pgfpathlineto{\pgfqpoint{3.910905in}{1.687156in}}%
\pgfpathlineto{\pgfqpoint{3.922937in}{1.677787in}}%
\pgfpathlineto{\pgfqpoint{3.930740in}{1.674264in}}%
\pgfpathlineto{\pgfqpoint{3.951880in}{1.677860in}}%
\pgfpathlineto{\pgfqpoint{3.964127in}{1.675404in}}%
\pgfpathlineto{\pgfqpoint{3.972545in}{1.671374in}}%
\pgfpathlineto{\pgfqpoint{3.976979in}{1.663725in}}%
\pgfpathlineto{\pgfqpoint{4.005055in}{1.664746in}}%
\pgfpathlineto{\pgfqpoint{4.009520in}{1.660065in}}%
\pgfpathlineto{\pgfqpoint{4.020053in}{1.659055in}}%
\pgfpathlineto{\pgfqpoint{4.030660in}{1.662045in}}%
\pgfpathlineto{\pgfqpoint{4.062021in}{1.661292in}}%
\pgfpathlineto{\pgfqpoint{4.068794in}{1.658987in}}%
\pgfpathlineto{\pgfqpoint{4.076709in}{1.661643in}}%
\pgfpathlineto{\pgfqpoint{4.164329in}{1.638779in}}%
\pgfpathlineto{\pgfqpoint{4.167899in}{1.626494in}}%
\pgfpathlineto{\pgfqpoint{4.175939in}{1.619838in}}%
\pgfpathlineto{\pgfqpoint{4.176774in}{1.607856in}}%
\pgfpathlineto{\pgfqpoint{4.165015in}{1.595725in}}%
\pgfpathlineto{\pgfqpoint{4.159282in}{1.584259in}}%
\pgfpathlineto{\pgfqpoint{4.152377in}{1.578133in}}%
\pgfpathlineto{\pgfqpoint{4.151752in}{1.572002in}}%
\pgfpathlineto{\pgfqpoint{4.140313in}{1.563715in}}%
\pgfpathlineto{\pgfqpoint{4.132464in}{1.551738in}}%
\pgfpathlineto{\pgfqpoint{4.128070in}{1.548828in}}%
\pgfpathlineto{\pgfqpoint{4.125673in}{1.537205in}}%
\pgfpathlineto{\pgfqpoint{4.137217in}{1.525368in}}%
\pgfpathlineto{\pgfqpoint{4.126069in}{1.505674in}}%
\pgfpathlineto{\pgfqpoint{4.124717in}{1.498158in}}%
\pgfpathlineto{\pgfqpoint{4.099647in}{1.395314in}}%
\pgfpathlineto{\pgfqpoint{4.046916in}{1.408826in}}%
\pgfpathlineto{\pgfqpoint{3.996046in}{1.421943in}}%
\pgfpathlineto{\pgfqpoint{3.965358in}{1.430483in}}%
\pgfpathlineto{\pgfqpoint{3.925927in}{1.441715in}}%
\pgfpathlineto{\pgfqpoint{3.863032in}{1.461511in}}%
\pgfpathlineto{\pgfqpoint{3.794661in}{1.482798in}}%
\pgfpathlineto{\pgfqpoint{3.777592in}{1.488751in}}%
\pgfpathclose%
\pgfusepath{fill}%
\end{pgfscope}%
\begin{pgfscope}%
\pgfpathrectangle{\pgfqpoint{3.525000in}{0.100000in}}{\pgfqpoint{2.862985in}{1.832778in}}%
\pgfusepath{clip}%
\pgfsetbuttcap%
\pgfsetmiterjoin%
\definecolor{currentfill}{rgb}{0.474971,0.717878,0.240138}%
\pgfsetfillcolor{currentfill}%
\pgfsetlinewidth{0.000000pt}%
\definecolor{currentstroke}{rgb}{0.000000,0.000000,0.000000}%
\pgfsetstrokecolor{currentstroke}%
\pgfsetstrokeopacity{0.000000}%
\pgfsetdash{}{0pt}%
\pgfpathmoveto{\pgfqpoint{6.147258in}{1.445782in}}%
\pgfpathlineto{\pgfqpoint{6.142278in}{1.449830in}}%
\pgfpathlineto{\pgfqpoint{6.142972in}{1.462444in}}%
\pgfpathlineto{\pgfqpoint{6.137824in}{1.498298in}}%
\pgfpathlineto{\pgfqpoint{6.141859in}{1.510906in}}%
\pgfpathlineto{\pgfqpoint{6.144741in}{1.533267in}}%
\pgfpathlineto{\pgfqpoint{6.140852in}{1.543846in}}%
\pgfpathlineto{\pgfqpoint{6.154999in}{1.557930in}}%
\pgfpathlineto{\pgfqpoint{6.158559in}{1.569133in}}%
\pgfpathlineto{\pgfqpoint{6.151856in}{1.577975in}}%
\pgfpathlineto{\pgfqpoint{6.154711in}{1.589025in}}%
\pgfpathlineto{\pgfqpoint{6.152922in}{1.595966in}}%
\pgfpathlineto{\pgfqpoint{6.154459in}{1.610832in}}%
\pgfpathlineto{\pgfqpoint{6.157332in}{1.615421in}}%
\pgfpathlineto{\pgfqpoint{6.166192in}{1.618318in}}%
\pgfpathlineto{\pgfqpoint{6.179117in}{1.580468in}}%
\pgfpathlineto{\pgfqpoint{6.203094in}{1.502035in}}%
\pgfpathlineto{\pgfqpoint{6.212168in}{1.495954in}}%
\pgfpathlineto{\pgfqpoint{6.213801in}{1.489012in}}%
\pgfpathlineto{\pgfqpoint{6.218590in}{1.486207in}}%
\pgfpathlineto{\pgfqpoint{6.218226in}{1.473617in}}%
\pgfpathlineto{\pgfqpoint{6.213151in}{1.473415in}}%
\pgfpathlineto{\pgfqpoint{6.202871in}{1.465574in}}%
\pgfpathlineto{\pgfqpoint{6.199906in}{1.457709in}}%
\pgfpathlineto{\pgfqpoint{6.172390in}{1.450901in}}%
\pgfpathlineto{\pgfqpoint{6.147258in}{1.445782in}}%
\pgfpathclose%
\pgfusepath{fill}%
\end{pgfscope}%
\begin{pgfscope}%
\pgfpathrectangle{\pgfqpoint{3.525000in}{0.100000in}}{\pgfqpoint{2.862985in}{1.832778in}}%
\pgfusepath{clip}%
\pgfsetbuttcap%
\pgfsetmiterjoin%
\definecolor{currentfill}{rgb}{0.451903,0.698501,0.225375}%
\pgfsetfillcolor{currentfill}%
\pgfsetlinewidth{0.000000pt}%
\definecolor{currentstroke}{rgb}{0.000000,0.000000,0.000000}%
\pgfsetstrokecolor{currentstroke}%
\pgfsetstrokeopacity{0.000000}%
\pgfsetdash{}{0pt}%
\pgfpathmoveto{\pgfqpoint{5.286865in}{1.172908in}}%
\pgfpathlineto{\pgfqpoint{5.272339in}{1.187293in}}%
\pgfpathlineto{\pgfqpoint{5.225906in}{1.184615in}}%
\pgfpathlineto{\pgfqpoint{5.153489in}{1.182229in}}%
\pgfpathlineto{\pgfqpoint{5.080637in}{1.183365in}}%
\pgfpathlineto{\pgfqpoint{5.075523in}{1.192282in}}%
\pgfpathlineto{\pgfqpoint{5.077611in}{1.201038in}}%
\pgfpathlineto{\pgfqpoint{5.076620in}{1.216037in}}%
\pgfpathlineto{\pgfqpoint{5.072770in}{1.230565in}}%
\pgfpathlineto{\pgfqpoint{5.073217in}{1.238237in}}%
\pgfpathlineto{\pgfqpoint{5.065033in}{1.246905in}}%
\pgfpathlineto{\pgfqpoint{5.066814in}{1.258967in}}%
\pgfpathlineto{\pgfqpoint{5.054260in}{1.282791in}}%
\pgfpathlineto{\pgfqpoint{5.050519in}{1.303018in}}%
\pgfpathlineto{\pgfqpoint{5.041489in}{1.319470in}}%
\pgfpathlineto{\pgfqpoint{5.050061in}{1.339878in}}%
\pgfpathlineto{\pgfqpoint{5.045791in}{1.349375in}}%
\pgfpathlineto{\pgfqpoint{5.050565in}{1.366544in}}%
\pgfpathlineto{\pgfqpoint{5.095992in}{1.366052in}}%
\pgfpathlineto{\pgfqpoint{5.142055in}{1.365982in}}%
\pgfpathlineto{\pgfqpoint{5.185597in}{1.366622in}}%
\pgfpathlineto{\pgfqpoint{5.233388in}{1.367621in}}%
\pgfpathlineto{\pgfqpoint{5.289454in}{1.369657in}}%
\pgfpathlineto{\pgfqpoint{5.290951in}{1.361517in}}%
\pgfpathlineto{\pgfqpoint{5.297158in}{1.354693in}}%
\pgfpathlineto{\pgfqpoint{5.293228in}{1.347582in}}%
\pgfpathlineto{\pgfqpoint{5.294839in}{1.333702in}}%
\pgfpathlineto{\pgfqpoint{5.298686in}{1.322483in}}%
\pgfpathlineto{\pgfqpoint{5.316125in}{1.316544in}}%
\pgfpathlineto{\pgfqpoint{5.319114in}{1.308600in}}%
\pgfpathlineto{\pgfqpoint{5.328684in}{1.299681in}}%
\pgfpathlineto{\pgfqpoint{5.332588in}{1.290454in}}%
\pgfpathlineto{\pgfqpoint{5.342285in}{1.284265in}}%
\pgfpathlineto{\pgfqpoint{5.343802in}{1.276813in}}%
\pgfpathlineto{\pgfqpoint{5.341914in}{1.265532in}}%
\pgfpathlineto{\pgfqpoint{5.336979in}{1.262151in}}%
\pgfpathlineto{\pgfqpoint{5.335490in}{1.251413in}}%
\pgfpathlineto{\pgfqpoint{5.321306in}{1.242884in}}%
\pgfpathlineto{\pgfqpoint{5.302818in}{1.238212in}}%
\pgfpathlineto{\pgfqpoint{5.301124in}{1.227469in}}%
\pgfpathlineto{\pgfqpoint{5.308256in}{1.219794in}}%
\pgfpathlineto{\pgfqpoint{5.308547in}{1.210143in}}%
\pgfpathlineto{\pgfqpoint{5.302793in}{1.202558in}}%
\pgfpathlineto{\pgfqpoint{5.299773in}{1.191288in}}%
\pgfpathlineto{\pgfqpoint{5.289771in}{1.187558in}}%
\pgfpathlineto{\pgfqpoint{5.290408in}{1.174999in}}%
\pgfpathlineto{\pgfqpoint{5.286865in}{1.172908in}}%
\pgfpathclose%
\pgfusepath{fill}%
\end{pgfscope}%
\begin{pgfscope}%
\pgfpathrectangle{\pgfqpoint{3.525000in}{0.100000in}}{\pgfqpoint{2.862985in}{1.832778in}}%
\pgfusepath{clip}%
\pgfsetbuttcap%
\pgfsetmiterjoin%
\definecolor{currentfill}{rgb}{0.620761,0.816917,0.403460}%
\pgfsetfillcolor{currentfill}%
\pgfsetlinewidth{0.000000pt}%
\definecolor{currentstroke}{rgb}{0.000000,0.000000,0.000000}%
\pgfsetstrokecolor{currentstroke}%
\pgfsetstrokeopacity{0.000000}%
\pgfsetdash{}{0pt}%
\pgfpathmoveto{\pgfqpoint{6.188120in}{1.409488in}}%
\pgfpathlineto{\pgfqpoint{6.174362in}{1.407318in}}%
\pgfpathlineto{\pgfqpoint{6.111166in}{1.392957in}}%
\pgfpathlineto{\pgfqpoint{6.110061in}{1.394631in}}%
\pgfpathlineto{\pgfqpoint{6.111015in}{1.438080in}}%
\pgfpathlineto{\pgfqpoint{6.147258in}{1.445782in}}%
\pgfpathlineto{\pgfqpoint{6.172390in}{1.450901in}}%
\pgfpathlineto{\pgfqpoint{6.199906in}{1.457709in}}%
\pgfpathlineto{\pgfqpoint{6.202871in}{1.465574in}}%
\pgfpathlineto{\pgfqpoint{6.213151in}{1.473415in}}%
\pgfpathlineto{\pgfqpoint{6.218226in}{1.473617in}}%
\pgfpathlineto{\pgfqpoint{6.224896in}{1.462179in}}%
\pgfpathlineto{\pgfqpoint{6.218843in}{1.445482in}}%
\pgfpathlineto{\pgfqpoint{6.217955in}{1.435691in}}%
\pgfpathlineto{\pgfqpoint{6.230207in}{1.436611in}}%
\pgfpathlineto{\pgfqpoint{6.235763in}{1.431914in}}%
\pgfpathlineto{\pgfqpoint{6.248224in}{1.412717in}}%
\pgfpathlineto{\pgfqpoint{6.254417in}{1.410382in}}%
\pgfpathlineto{\pgfqpoint{6.264790in}{1.411248in}}%
\pgfpathlineto{\pgfqpoint{6.272009in}{1.417772in}}%
\pgfpathlineto{\pgfqpoint{6.276808in}{1.411950in}}%
\pgfpathlineto{\pgfqpoint{6.246631in}{1.396025in}}%
\pgfpathlineto{\pgfqpoint{6.245688in}{1.407452in}}%
\pgfpathlineto{\pgfqpoint{6.231980in}{1.389693in}}%
\pgfpathlineto{\pgfqpoint{6.227160in}{1.386633in}}%
\pgfpathlineto{\pgfqpoint{6.220438in}{1.396881in}}%
\pgfpathlineto{\pgfqpoint{6.218598in}{1.398285in}}%
\pgfpathlineto{\pgfqpoint{6.212340in}{1.401600in}}%
\pgfpathlineto{\pgfqpoint{6.206882in}{1.415045in}}%
\pgfpathlineto{\pgfqpoint{6.188120in}{1.409488in}}%
\pgfpathclose%
\pgfusepath{fill}%
\end{pgfscope}%
\begin{pgfscope}%
\pgfpathrectangle{\pgfqpoint{3.525000in}{0.100000in}}{\pgfqpoint{2.862985in}{1.832778in}}%
\pgfusepath{clip}%
\pgfsetbuttcap%
\pgfsetmiterjoin%
\definecolor{currentfill}{rgb}{0.367320,0.627451,0.171242}%
\pgfsetfillcolor{currentfill}%
\pgfsetlinewidth{0.000000pt}%
\definecolor{currentstroke}{rgb}{0.000000,0.000000,0.000000}%
\pgfsetstrokecolor{currentstroke}%
\pgfsetstrokeopacity{0.000000}%
\pgfsetdash{}{0pt}%
\pgfpathmoveto{\pgfqpoint{4.778739in}{1.159333in}}%
\pgfpathlineto{\pgfqpoint{4.783777in}{1.221808in}}%
\pgfpathlineto{\pgfqpoint{4.755243in}{1.224124in}}%
\pgfpathlineto{\pgfqpoint{4.689283in}{1.230524in}}%
\pgfpathlineto{\pgfqpoint{4.695419in}{1.292777in}}%
\pgfpathlineto{\pgfqpoint{4.701565in}{1.355506in}}%
\pgfpathlineto{\pgfqpoint{4.726922in}{1.352768in}}%
\pgfpathlineto{\pgfqpoint{4.791682in}{1.346649in}}%
\pgfpathlineto{\pgfqpoint{4.830874in}{1.343580in}}%
\pgfpathlineto{\pgfqpoint{4.878322in}{1.340860in}}%
\pgfpathlineto{\pgfqpoint{4.956420in}{1.337164in}}%
\pgfpathlineto{\pgfqpoint{4.958137in}{1.333663in}}%
\pgfpathlineto{\pgfqpoint{4.980515in}{1.322591in}}%
\pgfpathlineto{\pgfqpoint{4.987377in}{1.328214in}}%
\pgfpathlineto{\pgfqpoint{5.007148in}{1.327701in}}%
\pgfpathlineto{\pgfqpoint{5.018930in}{1.321529in}}%
\pgfpathlineto{\pgfqpoint{5.037612in}{1.314513in}}%
\pgfpathlineto{\pgfqpoint{5.042847in}{1.304277in}}%
\pgfpathlineto{\pgfqpoint{5.050519in}{1.303018in}}%
\pgfpathlineto{\pgfqpoint{5.054260in}{1.282791in}}%
\pgfpathlineto{\pgfqpoint{5.066814in}{1.258967in}}%
\pgfpathlineto{\pgfqpoint{5.065033in}{1.246905in}}%
\pgfpathlineto{\pgfqpoint{5.073217in}{1.238237in}}%
\pgfpathlineto{\pgfqpoint{5.072770in}{1.230565in}}%
\pgfpathlineto{\pgfqpoint{5.076620in}{1.216037in}}%
\pgfpathlineto{\pgfqpoint{5.077611in}{1.201038in}}%
\pgfpathlineto{\pgfqpoint{5.075523in}{1.192282in}}%
\pgfpathlineto{\pgfqpoint{5.080637in}{1.183365in}}%
\pgfpathlineto{\pgfqpoint{5.087651in}{1.167151in}}%
\pgfpathlineto{\pgfqpoint{5.094364in}{1.160551in}}%
\pgfpathlineto{\pgfqpoint{5.102367in}{1.146251in}}%
\pgfpathlineto{\pgfqpoint{5.079686in}{1.146016in}}%
\pgfpathlineto{\pgfqpoint{5.030647in}{1.146774in}}%
\pgfpathlineto{\pgfqpoint{4.976466in}{1.148434in}}%
\pgfpathlineto{\pgfqpoint{4.921966in}{1.150526in}}%
\pgfpathlineto{\pgfqpoint{4.841835in}{1.154903in}}%
\pgfpathlineto{\pgfqpoint{4.778739in}{1.159333in}}%
\pgfpathclose%
\pgfusepath{fill}%
\end{pgfscope}%
\begin{pgfscope}%
\pgfpathrectangle{\pgfqpoint{3.525000in}{0.100000in}}{\pgfqpoint{2.862985in}{1.832778in}}%
\pgfusepath{clip}%
\pgfsetbuttcap%
\pgfsetmiterjoin%
\definecolor{currentfill}{rgb}{0.739254,0.890042,0.553941}%
\pgfsetfillcolor{currentfill}%
\pgfsetlinewidth{0.000000pt}%
\definecolor{currentstroke}{rgb}{0.000000,0.000000,0.000000}%
\pgfsetstrokecolor{currentstroke}%
\pgfsetstrokeopacity{0.000000}%
\pgfsetdash{}{0pt}%
\pgfpathmoveto{\pgfqpoint{5.822084in}{1.347781in}}%
\pgfpathlineto{\pgfqpoint{5.847175in}{1.371705in}}%
\pgfpathlineto{\pgfqpoint{5.850294in}{1.380209in}}%
\pgfpathlineto{\pgfqpoint{5.857646in}{1.387490in}}%
\pgfpathlineto{\pgfqpoint{5.852129in}{1.398090in}}%
\pgfpathlineto{\pgfqpoint{5.845212in}{1.404290in}}%
\pgfpathlineto{\pgfqpoint{5.843213in}{1.415277in}}%
\pgfpathlineto{\pgfqpoint{5.868959in}{1.426558in}}%
\pgfpathlineto{\pgfqpoint{5.890274in}{1.430126in}}%
\pgfpathlineto{\pgfqpoint{5.901728in}{1.430365in}}%
\pgfpathlineto{\pgfqpoint{5.910453in}{1.426054in}}%
\pgfpathlineto{\pgfqpoint{5.918958in}{1.429903in}}%
\pgfpathlineto{\pgfqpoint{5.939702in}{1.434217in}}%
\pgfpathlineto{\pgfqpoint{5.946870in}{1.439785in}}%
\pgfpathlineto{\pgfqpoint{5.957498in}{1.452111in}}%
\pgfpathlineto{\pgfqpoint{5.967167in}{1.457558in}}%
\pgfpathlineto{\pgfqpoint{5.967849in}{1.462782in}}%
\pgfpathlineto{\pgfqpoint{5.962774in}{1.474701in}}%
\pgfpathlineto{\pgfqpoint{5.966443in}{1.481714in}}%
\pgfpathlineto{\pgfqpoint{5.961503in}{1.489256in}}%
\pgfpathlineto{\pgfqpoint{5.953939in}{1.489782in}}%
\pgfpathlineto{\pgfqpoint{5.972870in}{1.512563in}}%
\pgfpathlineto{\pgfqpoint{5.975119in}{1.521245in}}%
\pgfpathlineto{\pgfqpoint{5.990018in}{1.543386in}}%
\pgfpathlineto{\pgfqpoint{6.003850in}{1.555368in}}%
\pgfpathlineto{\pgfqpoint{6.013347in}{1.560371in}}%
\pgfpathlineto{\pgfqpoint{6.044425in}{1.567413in}}%
\pgfpathlineto{\pgfqpoint{6.073552in}{1.575365in}}%
\pgfpathlineto{\pgfqpoint{6.077410in}{1.562882in}}%
\pgfpathlineto{\pgfqpoint{6.077683in}{1.551308in}}%
\pgfpathlineto{\pgfqpoint{6.084327in}{1.540801in}}%
\pgfpathlineto{\pgfqpoint{6.086375in}{1.530104in}}%
\pgfpathlineto{\pgfqpoint{6.083822in}{1.515590in}}%
\pgfpathlineto{\pgfqpoint{6.091598in}{1.496013in}}%
\pgfpathlineto{\pgfqpoint{6.100463in}{1.485127in}}%
\pgfpathlineto{\pgfqpoint{6.107320in}{1.449744in}}%
\pgfpathlineto{\pgfqpoint{6.111015in}{1.438080in}}%
\pgfpathlineto{\pgfqpoint{6.110061in}{1.394631in}}%
\pgfpathlineto{\pgfqpoint{6.111166in}{1.392957in}}%
\pgfpathlineto{\pgfqpoint{6.119238in}{1.346227in}}%
\pgfpathlineto{\pgfqpoint{6.123765in}{1.341969in}}%
\pgfpathlineto{\pgfqpoint{6.114028in}{1.332508in}}%
\pgfpathlineto{\pgfqpoint{6.118832in}{1.327078in}}%
\pgfpathlineto{\pgfqpoint{6.114609in}{1.318865in}}%
\pgfpathlineto{\pgfqpoint{6.114645in}{1.315368in}}%
\pgfpathlineto{\pgfqpoint{6.109366in}{1.312213in}}%
\pgfpathlineto{\pgfqpoint{6.106795in}{1.305239in}}%
\pgfpathlineto{\pgfqpoint{6.107610in}{1.324417in}}%
\pgfpathlineto{\pgfqpoint{6.091246in}{1.328639in}}%
\pgfpathlineto{\pgfqpoint{6.065655in}{1.337368in}}%
\pgfpathlineto{\pgfqpoint{6.062035in}{1.341668in}}%
\pgfpathlineto{\pgfqpoint{6.051342in}{1.342698in}}%
\pgfpathlineto{\pgfqpoint{6.045190in}{1.349098in}}%
\pgfpathlineto{\pgfqpoint{6.041871in}{1.361844in}}%
\pgfpathlineto{\pgfqpoint{6.033151in}{1.363429in}}%
\pgfpathlineto{\pgfqpoint{6.027319in}{1.370116in}}%
\pgfpathlineto{\pgfqpoint{5.953076in}{1.355148in}}%
\pgfpathlineto{\pgfqpoint{5.917594in}{1.347810in}}%
\pgfpathlineto{\pgfqpoint{5.863721in}{1.337977in}}%
\pgfpathlineto{\pgfqpoint{5.824926in}{1.331432in}}%
\pgfpathlineto{\pgfqpoint{5.822084in}{1.347781in}}%
\pgfpathclose%
\pgfusepath{fill}%
\end{pgfscope}%
\begin{pgfscope}%
\pgfpathrectangle{\pgfqpoint{3.525000in}{0.100000in}}{\pgfqpoint{2.862985in}{1.832778in}}%
\pgfusepath{clip}%
\pgfsetbuttcap%
\pgfsetmiterjoin%
\definecolor{currentfill}{rgb}{0.739254,0.890042,0.553941}%
\pgfsetfillcolor{currentfill}%
\pgfsetlinewidth{0.000000pt}%
\definecolor{currentstroke}{rgb}{0.000000,0.000000,0.000000}%
\pgfsetstrokecolor{currentstroke}%
\pgfsetstrokeopacity{0.000000}%
\pgfsetdash{}{0pt}%
\pgfpathmoveto{\pgfqpoint{6.120139in}{1.301344in}}%
\pgfpathlineto{\pgfqpoint{6.108657in}{1.297767in}}%
\pgfpathlineto{\pgfqpoint{6.106726in}{1.301055in}}%
\pgfpathlineto{\pgfqpoint{6.110411in}{1.312094in}}%
\pgfpathlineto{\pgfqpoint{6.117237in}{1.313181in}}%
\pgfpathlineto{\pgfqpoint{6.116626in}{1.316656in}}%
\pgfpathlineto{\pgfqpoint{6.122754in}{1.321872in}}%
\pgfpathlineto{\pgfqpoint{6.140472in}{1.326028in}}%
\pgfpathlineto{\pgfqpoint{6.148360in}{1.332319in}}%
\pgfpathlineto{\pgfqpoint{6.166079in}{1.337561in}}%
\pgfpathlineto{\pgfqpoint{6.174161in}{1.335643in}}%
\pgfpathlineto{\pgfqpoint{6.180959in}{1.344095in}}%
\pgfpathlineto{\pgfqpoint{6.191234in}{1.345211in}}%
\pgfpathlineto{\pgfqpoint{6.173704in}{1.328725in}}%
\pgfpathlineto{\pgfqpoint{6.120139in}{1.301344in}}%
\pgfpathclose%
\pgfusepath{fill}%
\end{pgfscope}%
\begin{pgfscope}%
\pgfpathrectangle{\pgfqpoint{3.525000in}{0.100000in}}{\pgfqpoint{2.862985in}{1.832778in}}%
\pgfusepath{clip}%
\pgfsetbuttcap%
\pgfsetmiterjoin%
\definecolor{currentfill}{rgb}{0.760477,0.899270,0.588082}%
\pgfsetfillcolor{currentfill}%
\pgfsetlinewidth{0.000000pt}%
\definecolor{currentstroke}{rgb}{0.000000,0.000000,0.000000}%
\pgfsetstrokecolor{currentstroke}%
\pgfsetstrokeopacity{0.000000}%
\pgfsetdash{}{0pt}%
\pgfpathmoveto{\pgfqpoint{5.862178in}{1.192485in}}%
\pgfpathlineto{\pgfqpoint{5.812526in}{1.184275in}}%
\pgfpathlineto{\pgfqpoint{5.803519in}{1.241018in}}%
\pgfpathlineto{\pgfqpoint{5.790146in}{1.324661in}}%
\pgfpathlineto{\pgfqpoint{5.822084in}{1.347781in}}%
\pgfpathlineto{\pgfqpoint{5.824926in}{1.331432in}}%
\pgfpathlineto{\pgfqpoint{5.863721in}{1.337977in}}%
\pgfpathlineto{\pgfqpoint{5.917594in}{1.347810in}}%
\pgfpathlineto{\pgfqpoint{5.953076in}{1.355148in}}%
\pgfpathlineto{\pgfqpoint{6.027319in}{1.370116in}}%
\pgfpathlineto{\pgfqpoint{6.033151in}{1.363429in}}%
\pgfpathlineto{\pgfqpoint{6.041871in}{1.361844in}}%
\pgfpathlineto{\pgfqpoint{6.045190in}{1.349098in}}%
\pgfpathlineto{\pgfqpoint{6.051342in}{1.342698in}}%
\pgfpathlineto{\pgfqpoint{6.062035in}{1.341668in}}%
\pgfpathlineto{\pgfqpoint{6.065655in}{1.337368in}}%
\pgfpathlineto{\pgfqpoint{6.061981in}{1.334051in}}%
\pgfpathlineto{\pgfqpoint{6.058670in}{1.322314in}}%
\pgfpathlineto{\pgfqpoint{6.050913in}{1.309125in}}%
\pgfpathlineto{\pgfqpoint{6.056120in}{1.303390in}}%
\pgfpathlineto{\pgfqpoint{6.051859in}{1.297245in}}%
\pgfpathlineto{\pgfqpoint{6.054263in}{1.283769in}}%
\pgfpathlineto{\pgfqpoint{6.061936in}{1.276695in}}%
\pgfpathlineto{\pgfqpoint{6.080144in}{1.265202in}}%
\pgfpathlineto{\pgfqpoint{6.065481in}{1.248986in}}%
\pgfpathlineto{\pgfqpoint{6.065275in}{1.242831in}}%
\pgfpathlineto{\pgfqpoint{6.053338in}{1.234894in}}%
\pgfpathlineto{\pgfqpoint{6.040138in}{1.233404in}}%
\pgfpathlineto{\pgfqpoint{6.036873in}{1.226505in}}%
\pgfpathlineto{\pgfqpoint{6.000158in}{1.218561in}}%
\pgfpathlineto{\pgfqpoint{5.957316in}{1.209933in}}%
\pgfpathlineto{\pgfqpoint{5.927868in}{1.204659in}}%
\pgfpathlineto{\pgfqpoint{5.862178in}{1.192485in}}%
\pgfpathclose%
\pgfusepath{fill}%
\end{pgfscope}%
\begin{pgfscope}%
\pgfpathrectangle{\pgfqpoint{3.525000in}{0.100000in}}{\pgfqpoint{2.862985in}{1.832778in}}%
\pgfusepath{clip}%
\pgfsetbuttcap%
\pgfsetmiterjoin%
\definecolor{currentfill}{rgb}{0.647059,0.833987,0.435294}%
\pgfsetfillcolor{currentfill}%
\pgfsetlinewidth{0.000000pt}%
\definecolor{currentstroke}{rgb}{0.000000,0.000000,0.000000}%
\pgfsetstrokecolor{currentstroke}%
\pgfsetstrokeopacity{0.000000}%
\pgfsetdash{}{0pt}%
\pgfpathmoveto{\pgfqpoint{6.111166in}{1.392957in}}%
\pgfpathlineto{\pgfqpoint{6.174362in}{1.407318in}}%
\pgfpathlineto{\pgfqpoint{6.188120in}{1.409488in}}%
\pgfpathlineto{\pgfqpoint{6.197226in}{1.373686in}}%
\pgfpathlineto{\pgfqpoint{6.195785in}{1.367275in}}%
\pgfpathlineto{\pgfqpoint{6.166533in}{1.355951in}}%
\pgfpathlineto{\pgfqpoint{6.149071in}{1.351992in}}%
\pgfpathlineto{\pgfqpoint{6.141651in}{1.343114in}}%
\pgfpathlineto{\pgfqpoint{6.118832in}{1.327078in}}%
\pgfpathlineto{\pgfqpoint{6.114028in}{1.332508in}}%
\pgfpathlineto{\pgfqpoint{6.123765in}{1.341969in}}%
\pgfpathlineto{\pgfqpoint{6.119238in}{1.346227in}}%
\pgfpathlineto{\pgfqpoint{6.111166in}{1.392957in}}%
\pgfpathclose%
\pgfusepath{fill}%
\end{pgfscope}%
\begin{pgfscope}%
\pgfpathrectangle{\pgfqpoint{3.525000in}{0.100000in}}{\pgfqpoint{2.862985in}{1.832778in}}%
\pgfusepath{clip}%
\pgfsetbuttcap%
\pgfsetmiterjoin%
\definecolor{currentfill}{rgb}{0.541869,0.765705,0.307958}%
\pgfsetfillcolor{currentfill}%
\pgfsetlinewidth{0.000000pt}%
\definecolor{currentstroke}{rgb}{0.000000,0.000000,0.000000}%
\pgfsetstrokecolor{currentstroke}%
\pgfsetstrokeopacity{0.000000}%
\pgfsetdash{}{0pt}%
\pgfpathmoveto{\pgfqpoint{6.188120in}{1.409488in}}%
\pgfpathlineto{\pgfqpoint{6.206882in}{1.415045in}}%
\pgfpathlineto{\pgfqpoint{6.212340in}{1.401600in}}%
\pgfpathlineto{\pgfqpoint{6.218598in}{1.398285in}}%
\pgfpathlineto{\pgfqpoint{6.210875in}{1.392622in}}%
\pgfpathlineto{\pgfqpoint{6.211849in}{1.375988in}}%
\pgfpathlineto{\pgfqpoint{6.195785in}{1.367275in}}%
\pgfpathlineto{\pgfqpoint{6.197226in}{1.373686in}}%
\pgfpathlineto{\pgfqpoint{6.188120in}{1.409488in}}%
\pgfpathclose%
\pgfusepath{fill}%
\end{pgfscope}%
\begin{pgfscope}%
\pgfpathrectangle{\pgfqpoint{3.525000in}{0.100000in}}{\pgfqpoint{2.862985in}{1.832778in}}%
\pgfusepath{clip}%
\pgfsetbuttcap%
\pgfsetmiterjoin%
\definecolor{currentfill}{rgb}{0.568166,0.782776,0.339792}%
\pgfsetfillcolor{currentfill}%
\pgfsetlinewidth{0.000000pt}%
\definecolor{currentstroke}{rgb}{0.000000,0.000000,0.000000}%
\pgfsetstrokecolor{currentstroke}%
\pgfsetstrokeopacity{0.000000}%
\pgfsetdash{}{0pt}%
\pgfpathmoveto{\pgfqpoint{6.051223in}{1.229031in}}%
\pgfpathlineto{\pgfqpoint{6.053338in}{1.234894in}}%
\pgfpathlineto{\pgfqpoint{6.065275in}{1.242831in}}%
\pgfpathlineto{\pgfqpoint{6.065481in}{1.248986in}}%
\pgfpathlineto{\pgfqpoint{6.080144in}{1.265202in}}%
\pgfpathlineto{\pgfqpoint{6.061936in}{1.276695in}}%
\pgfpathlineto{\pgfqpoint{6.054263in}{1.283769in}}%
\pgfpathlineto{\pgfqpoint{6.051859in}{1.297245in}}%
\pgfpathlineto{\pgfqpoint{6.056120in}{1.303390in}}%
\pgfpathlineto{\pgfqpoint{6.050913in}{1.309125in}}%
\pgfpathlineto{\pgfqpoint{6.058670in}{1.322314in}}%
\pgfpathlineto{\pgfqpoint{6.061981in}{1.334051in}}%
\pgfpathlineto{\pgfqpoint{6.065655in}{1.337368in}}%
\pgfpathlineto{\pgfqpoint{6.091246in}{1.328639in}}%
\pgfpathlineto{\pgfqpoint{6.107610in}{1.324417in}}%
\pgfpathlineto{\pgfqpoint{6.106795in}{1.305239in}}%
\pgfpathlineto{\pgfqpoint{6.096854in}{1.290696in}}%
\pgfpathlineto{\pgfqpoint{6.105045in}{1.288556in}}%
\pgfpathlineto{\pgfqpoint{6.113550in}{1.282315in}}%
\pgfpathlineto{\pgfqpoint{6.114051in}{1.265245in}}%
\pgfpathlineto{\pgfqpoint{6.111477in}{1.253161in}}%
\pgfpathlineto{\pgfqpoint{6.113190in}{1.243235in}}%
\pgfpathlineto{\pgfqpoint{6.098407in}{1.209740in}}%
\pgfpathlineto{\pgfqpoint{6.093094in}{1.194100in}}%
\pgfpathlineto{\pgfqpoint{6.085676in}{1.201703in}}%
\pgfpathlineto{\pgfqpoint{6.075838in}{1.200407in}}%
\pgfpathlineto{\pgfqpoint{6.051218in}{1.214632in}}%
\pgfpathlineto{\pgfqpoint{6.048699in}{1.222249in}}%
\pgfpathlineto{\pgfqpoint{6.051223in}{1.229031in}}%
\pgfpathclose%
\pgfusepath{fill}%
\end{pgfscope}%
\begin{pgfscope}%
\pgfpathrectangle{\pgfqpoint{3.525000in}{0.100000in}}{\pgfqpoint{2.862985in}{1.832778in}}%
\pgfusepath{clip}%
\pgfsetbuttcap%
\pgfsetmiterjoin%
\definecolor{currentfill}{rgb}{0.524337,0.754325,0.286736}%
\pgfsetfillcolor{currentfill}%
\pgfsetlinewidth{0.000000pt}%
\definecolor{currentstroke}{rgb}{0.000000,0.000000,0.000000}%
\pgfsetstrokecolor{currentstroke}%
\pgfsetstrokeopacity{0.000000}%
\pgfsetdash{}{0pt}%
\pgfpathmoveto{\pgfqpoint{5.462013in}{1.021868in}}%
\pgfpathlineto{\pgfqpoint{5.460862in}{1.037407in}}%
\pgfpathlineto{\pgfqpoint{5.465336in}{1.045116in}}%
\pgfpathlineto{\pgfqpoint{5.462387in}{1.048893in}}%
\pgfpathlineto{\pgfqpoint{5.473140in}{1.061261in}}%
\pgfpathlineto{\pgfqpoint{5.472524in}{1.064241in}}%
\pgfpathlineto{\pgfqpoint{5.482971in}{1.085818in}}%
\pgfpathlineto{\pgfqpoint{5.480859in}{1.096224in}}%
\pgfpathlineto{\pgfqpoint{5.473283in}{1.107118in}}%
\pgfpathlineto{\pgfqpoint{5.478542in}{1.120373in}}%
\pgfpathlineto{\pgfqpoint{5.471726in}{1.207597in}}%
\pgfpathlineto{\pgfqpoint{5.466818in}{1.268787in}}%
\pgfpathlineto{\pgfqpoint{5.473598in}{1.263717in}}%
\pgfpathlineto{\pgfqpoint{5.481163in}{1.263854in}}%
\pgfpathlineto{\pgfqpoint{5.499029in}{1.274203in}}%
\pgfpathlineto{\pgfqpoint{5.553831in}{1.279317in}}%
\pgfpathlineto{\pgfqpoint{5.594393in}{1.283592in}}%
\pgfpathlineto{\pgfqpoint{5.594752in}{1.279623in}}%
\pgfpathlineto{\pgfqpoint{5.604011in}{1.195770in}}%
\pgfpathlineto{\pgfqpoint{5.611981in}{1.117748in}}%
\pgfpathlineto{\pgfqpoint{5.608545in}{1.114085in}}%
\pgfpathlineto{\pgfqpoint{5.613822in}{1.104993in}}%
\pgfpathlineto{\pgfqpoint{5.613795in}{1.098441in}}%
\pgfpathlineto{\pgfqpoint{5.606265in}{1.096795in}}%
\pgfpathlineto{\pgfqpoint{5.597839in}{1.090478in}}%
\pgfpathlineto{\pgfqpoint{5.592134in}{1.092965in}}%
\pgfpathlineto{\pgfqpoint{5.583595in}{1.088922in}}%
\pgfpathlineto{\pgfqpoint{5.586237in}{1.080801in}}%
\pgfpathlineto{\pgfqpoint{5.577460in}{1.072642in}}%
\pgfpathlineto{\pgfqpoint{5.575037in}{1.063205in}}%
\pgfpathlineto{\pgfqpoint{5.569002in}{1.061656in}}%
\pgfpathlineto{\pgfqpoint{5.564498in}{1.054509in}}%
\pgfpathlineto{\pgfqpoint{5.565088in}{1.047311in}}%
\pgfpathlineto{\pgfqpoint{5.559790in}{1.042250in}}%
\pgfpathlineto{\pgfqpoint{5.551817in}{1.043031in}}%
\pgfpathlineto{\pgfqpoint{5.545757in}{1.050797in}}%
\pgfpathlineto{\pgfqpoint{5.535489in}{1.043318in}}%
\pgfpathlineto{\pgfqpoint{5.536209in}{1.036783in}}%
\pgfpathlineto{\pgfqpoint{5.524789in}{1.032965in}}%
\pgfpathlineto{\pgfqpoint{5.521919in}{1.037771in}}%
\pgfpathlineto{\pgfqpoint{5.510801in}{1.032321in}}%
\pgfpathlineto{\pgfqpoint{5.506739in}{1.024520in}}%
\pgfpathlineto{\pgfqpoint{5.493349in}{1.032520in}}%
\pgfpathlineto{\pgfqpoint{5.472127in}{1.027220in}}%
\pgfpathlineto{\pgfqpoint{5.462013in}{1.021868in}}%
\pgfpathclose%
\pgfusepath{fill}%
\end{pgfscope}%
\begin{pgfscope}%
\pgfpathrectangle{\pgfqpoint{3.525000in}{0.100000in}}{\pgfqpoint{2.862985in}{1.832778in}}%
\pgfusepath{clip}%
\pgfsetbuttcap%
\pgfsetmiterjoin%
\definecolor{currentfill}{rgb}{0.954248,0.966936,0.935640}%
\pgfsetfillcolor{currentfill}%
\pgfsetlinewidth{0.000000pt}%
\definecolor{currentstroke}{rgb}{0.000000,0.000000,0.000000}%
\pgfsetstrokecolor{currentstroke}%
\pgfsetstrokeopacity{0.000000}%
\pgfsetdash{}{0pt}%
\pgfpathmoveto{\pgfqpoint{3.965358in}{1.430483in}}%
\pgfpathlineto{\pgfqpoint{3.996046in}{1.421943in}}%
\pgfpathlineto{\pgfqpoint{4.046916in}{1.408826in}}%
\pgfpathlineto{\pgfqpoint{4.099647in}{1.395314in}}%
\pgfpathlineto{\pgfqpoint{4.148250in}{1.383593in}}%
\pgfpathlineto{\pgfqpoint{4.190246in}{1.374085in}}%
\pgfpathlineto{\pgfqpoint{4.235177in}{1.364205in}}%
\pgfpathlineto{\pgfqpoint{4.222195in}{1.302942in}}%
\pgfpathlineto{\pgfqpoint{4.210630in}{1.248537in}}%
\pgfpathlineto{\pgfqpoint{4.191658in}{1.160754in}}%
\pgfpathlineto{\pgfqpoint{4.177414in}{1.094498in}}%
\pgfpathlineto{\pgfqpoint{4.169719in}{1.057523in}}%
\pgfpathlineto{\pgfqpoint{4.159852in}{1.009534in}}%
\pgfpathlineto{\pgfqpoint{4.153025in}{0.999798in}}%
\pgfpathlineto{\pgfqpoint{4.147531in}{0.999473in}}%
\pgfpathlineto{\pgfqpoint{4.143612in}{1.007971in}}%
\pgfpathlineto{\pgfqpoint{4.134613in}{1.011057in}}%
\pgfpathlineto{\pgfqpoint{4.124916in}{1.009976in}}%
\pgfpathlineto{\pgfqpoint{4.122164in}{1.003024in}}%
\pgfpathlineto{\pgfqpoint{4.121877in}{0.978886in}}%
\pgfpathlineto{\pgfqpoint{4.118973in}{0.973374in}}%
\pgfpathlineto{\pgfqpoint{4.120633in}{0.953995in}}%
\pgfpathlineto{\pgfqpoint{4.114569in}{0.941085in}}%
\pgfpathlineto{\pgfqpoint{4.065397in}{1.016802in}}%
\pgfpathlineto{\pgfqpoint{4.016971in}{1.090379in}}%
\pgfpathlineto{\pgfqpoint{3.991763in}{1.128960in}}%
\pgfpathlineto{\pgfqpoint{3.970735in}{1.162262in}}%
\pgfpathlineto{\pgfqpoint{3.944238in}{1.203456in}}%
\pgfpathlineto{\pgfqpoint{3.914244in}{1.249665in}}%
\pgfpathlineto{\pgfqpoint{3.926576in}{1.293524in}}%
\pgfpathlineto{\pgfqpoint{3.951391in}{1.381486in}}%
\pgfpathlineto{\pgfqpoint{3.965358in}{1.430483in}}%
\pgfpathclose%
\pgfusepath{fill}%
\end{pgfscope}%
\begin{pgfscope}%
\pgfpathrectangle{\pgfqpoint{3.525000in}{0.100000in}}{\pgfqpoint{2.862985in}{1.832778in}}%
\pgfusepath{clip}%
\pgfsetbuttcap%
\pgfsetmiterjoin%
\definecolor{currentfill}{rgb}{0.367320,0.627451,0.171242}%
\pgfsetfillcolor{currentfill}%
\pgfsetlinewidth{0.000000pt}%
\definecolor{currentstroke}{rgb}{0.000000,0.000000,0.000000}%
\pgfsetstrokecolor{currentstroke}%
\pgfsetstrokeopacity{0.000000}%
\pgfsetdash{}{0pt}%
\pgfpathmoveto{\pgfqpoint{4.235177in}{1.364205in}}%
\pgfpathlineto{\pgfqpoint{4.283212in}{1.354650in}}%
\pgfpathlineto{\pgfqpoint{4.372141in}{1.337483in}}%
\pgfpathlineto{\pgfqpoint{4.361009in}{1.275734in}}%
\pgfpathlineto{\pgfqpoint{4.409967in}{1.267352in}}%
\pgfpathlineto{\pgfqpoint{4.454571in}{1.260214in}}%
\pgfpathlineto{\pgfqpoint{4.446820in}{1.211396in}}%
\pgfpathlineto{\pgfqpoint{4.438608in}{1.158753in}}%
\pgfpathlineto{\pgfqpoint{4.427613in}{1.089632in}}%
\pgfpathlineto{\pgfqpoint{4.427328in}{1.083837in}}%
\pgfpathlineto{\pgfqpoint{4.415915in}{1.012202in}}%
\pgfpathlineto{\pgfqpoint{4.368932in}{1.019484in}}%
\pgfpathlineto{\pgfqpoint{4.344991in}{1.024292in}}%
\pgfpathlineto{\pgfqpoint{4.258455in}{1.039504in}}%
\pgfpathlineto{\pgfqpoint{4.225868in}{1.045930in}}%
\pgfpathlineto{\pgfqpoint{4.169719in}{1.057523in}}%
\pgfpathlineto{\pgfqpoint{4.177414in}{1.094498in}}%
\pgfpathlineto{\pgfqpoint{4.191658in}{1.160754in}}%
\pgfpathlineto{\pgfqpoint{4.210630in}{1.248537in}}%
\pgfpathlineto{\pgfqpoint{4.222195in}{1.302942in}}%
\pgfpathlineto{\pgfqpoint{4.235177in}{1.364205in}}%
\pgfpathclose%
\pgfusepath{fill}%
\end{pgfscope}%
\begin{pgfscope}%
\pgfpathrectangle{\pgfqpoint{3.525000in}{0.100000in}}{\pgfqpoint{2.862985in}{1.832778in}}%
\pgfusepath{clip}%
\pgfsetbuttcap%
\pgfsetmiterjoin%
\definecolor{currentfill}{rgb}{0.781699,0.908497,0.622222}%
\pgfsetfillcolor{currentfill}%
\pgfsetlinewidth{0.000000pt}%
\definecolor{currentstroke}{rgb}{0.000000,0.000000,0.000000}%
\pgfsetstrokecolor{currentstroke}%
\pgfsetstrokeopacity{0.000000}%
\pgfsetdash{}{0pt}%
\pgfpathmoveto{\pgfqpoint{3.777592in}{1.488751in}}%
\pgfpathlineto{\pgfqpoint{3.794661in}{1.482798in}}%
\pgfpathlineto{\pgfqpoint{3.863032in}{1.461511in}}%
\pgfpathlineto{\pgfqpoint{3.925927in}{1.441715in}}%
\pgfpathlineto{\pgfqpoint{3.965358in}{1.430483in}}%
\pgfpathlineto{\pgfqpoint{3.951391in}{1.381486in}}%
\pgfpathlineto{\pgfqpoint{3.926576in}{1.293524in}}%
\pgfpathlineto{\pgfqpoint{3.914244in}{1.249665in}}%
\pgfpathlineto{\pgfqpoint{3.944238in}{1.203456in}}%
\pgfpathlineto{\pgfqpoint{3.970735in}{1.162262in}}%
\pgfpathlineto{\pgfqpoint{3.991763in}{1.128960in}}%
\pgfpathlineto{\pgfqpoint{4.016971in}{1.090379in}}%
\pgfpathlineto{\pgfqpoint{4.065397in}{1.016802in}}%
\pgfpathlineto{\pgfqpoint{4.114569in}{0.941085in}}%
\pgfpathlineto{\pgfqpoint{4.112595in}{0.933576in}}%
\pgfpathlineto{\pgfqpoint{4.118526in}{0.921617in}}%
\pgfpathlineto{\pgfqpoint{4.119699in}{0.905259in}}%
\pgfpathlineto{\pgfqpoint{4.128313in}{0.897324in}}%
\pgfpathlineto{\pgfqpoint{4.128654in}{0.890930in}}%
\pgfpathlineto{\pgfqpoint{4.113230in}{0.883673in}}%
\pgfpathlineto{\pgfqpoint{4.105882in}{0.876407in}}%
\pgfpathlineto{\pgfqpoint{4.103591in}{0.860364in}}%
\pgfpathlineto{\pgfqpoint{4.099871in}{0.851614in}}%
\pgfpathlineto{\pgfqpoint{4.092035in}{0.844239in}}%
\pgfpathlineto{\pgfqpoint{4.084240in}{0.825061in}}%
\pgfpathlineto{\pgfqpoint{4.095200in}{0.815071in}}%
\pgfpathlineto{\pgfqpoint{4.093785in}{0.806836in}}%
\pgfpathlineto{\pgfqpoint{4.084798in}{0.801105in}}%
\pgfpathlineto{\pgfqpoint{4.078603in}{0.802117in}}%
\pgfpathlineto{\pgfqpoint{4.005470in}{0.812266in}}%
\pgfpathlineto{\pgfqpoint{3.951454in}{0.819928in}}%
\pgfpathlineto{\pgfqpoint{3.953816in}{0.828649in}}%
\pgfpathlineto{\pgfqpoint{3.950300in}{0.843128in}}%
\pgfpathlineto{\pgfqpoint{3.949945in}{0.857733in}}%
\pgfpathlineto{\pgfqpoint{3.947647in}{0.866279in}}%
\pgfpathlineto{\pgfqpoint{3.940568in}{0.878498in}}%
\pgfpathlineto{\pgfqpoint{3.920206in}{0.906691in}}%
\pgfpathlineto{\pgfqpoint{3.913537in}{0.910145in}}%
\pgfpathlineto{\pgfqpoint{3.904944in}{0.910090in}}%
\pgfpathlineto{\pgfqpoint{3.906908in}{0.918984in}}%
\pgfpathlineto{\pgfqpoint{3.902839in}{0.930106in}}%
\pgfpathlineto{\pgfqpoint{3.882853in}{0.935633in}}%
\pgfpathlineto{\pgfqpoint{3.870671in}{0.945869in}}%
\pgfpathlineto{\pgfqpoint{3.869661in}{0.952132in}}%
\pgfpathlineto{\pgfqpoint{3.855633in}{0.967637in}}%
\pgfpathlineto{\pgfqpoint{3.842280in}{0.970604in}}%
\pgfpathlineto{\pgfqpoint{3.829908in}{0.978484in}}%
\pgfpathlineto{\pgfqpoint{3.813637in}{0.981210in}}%
\pgfpathlineto{\pgfqpoint{3.806685in}{0.991719in}}%
\pgfpathlineto{\pgfqpoint{3.813287in}{1.008354in}}%
\pgfpathlineto{\pgfqpoint{3.811268in}{1.012095in}}%
\pgfpathlineto{\pgfqpoint{3.816766in}{1.025962in}}%
\pgfpathlineto{\pgfqpoint{3.806986in}{1.033342in}}%
\pgfpathlineto{\pgfqpoint{3.810159in}{1.046721in}}%
\pgfpathlineto{\pgfqpoint{3.804935in}{1.050141in}}%
\pgfpathlineto{\pgfqpoint{3.800424in}{1.062787in}}%
\pgfpathlineto{\pgfqpoint{3.794981in}{1.066644in}}%
\pgfpathlineto{\pgfqpoint{3.794572in}{1.075778in}}%
\pgfpathlineto{\pgfqpoint{3.790307in}{1.082219in}}%
\pgfpathlineto{\pgfqpoint{3.783881in}{1.103911in}}%
\pgfpathlineto{\pgfqpoint{3.776850in}{1.114301in}}%
\pgfpathlineto{\pgfqpoint{3.778378in}{1.131939in}}%
\pgfpathlineto{\pgfqpoint{3.786647in}{1.133714in}}%
\pgfpathlineto{\pgfqpoint{3.792031in}{1.143280in}}%
\pgfpathlineto{\pgfqpoint{3.788778in}{1.153656in}}%
\pgfpathlineto{\pgfqpoint{3.779984in}{1.155388in}}%
\pgfpathlineto{\pgfqpoint{3.775596in}{1.160241in}}%
\pgfpathlineto{\pgfqpoint{3.768491in}{1.178092in}}%
\pgfpathlineto{\pgfqpoint{3.771812in}{1.184503in}}%
\pgfpathlineto{\pgfqpoint{3.769457in}{1.196440in}}%
\pgfpathlineto{\pgfqpoint{3.774672in}{1.211900in}}%
\pgfpathlineto{\pgfqpoint{3.780782in}{1.206207in}}%
\pgfpathlineto{\pgfqpoint{3.778009in}{1.199485in}}%
\pgfpathlineto{\pgfqpoint{3.788597in}{1.188830in}}%
\pgfpathlineto{\pgfqpoint{3.787921in}{1.204652in}}%
\pgfpathlineto{\pgfqpoint{3.783384in}{1.208896in}}%
\pgfpathlineto{\pgfqpoint{3.788570in}{1.222808in}}%
\pgfpathlineto{\pgfqpoint{3.808185in}{1.220593in}}%
\pgfpathlineto{\pgfqpoint{3.805545in}{1.225766in}}%
\pgfpathlineto{\pgfqpoint{3.792599in}{1.225257in}}%
\pgfpathlineto{\pgfqpoint{3.786441in}{1.233093in}}%
\pgfpathlineto{\pgfqpoint{3.778689in}{1.226135in}}%
\pgfpathlineto{\pgfqpoint{3.774576in}{1.214504in}}%
\pgfpathlineto{\pgfqpoint{3.763597in}{1.230154in}}%
\pgfpathlineto{\pgfqpoint{3.759369in}{1.233003in}}%
\pgfpathlineto{\pgfqpoint{3.760980in}{1.250036in}}%
\pgfpathlineto{\pgfqpoint{3.757634in}{1.260081in}}%
\pgfpathlineto{\pgfqpoint{3.751568in}{1.269517in}}%
\pgfpathlineto{\pgfqpoint{3.739145in}{1.298432in}}%
\pgfpathlineto{\pgfqpoint{3.743207in}{1.304826in}}%
\pgfpathlineto{\pgfqpoint{3.743159in}{1.325033in}}%
\pgfpathlineto{\pgfqpoint{3.749864in}{1.336303in}}%
\pgfpathlineto{\pgfqpoint{3.751401in}{1.353918in}}%
\pgfpathlineto{\pgfqpoint{3.745079in}{1.374099in}}%
\pgfpathlineto{\pgfqpoint{3.736692in}{1.386945in}}%
\pgfpathlineto{\pgfqpoint{3.738188in}{1.398536in}}%
\pgfpathlineto{\pgfqpoint{3.761740in}{1.426591in}}%
\pgfpathlineto{\pgfqpoint{3.762910in}{1.436161in}}%
\pgfpathlineto{\pgfqpoint{3.773523in}{1.454416in}}%
\pgfpathlineto{\pgfqpoint{3.774995in}{1.471719in}}%
\pgfpathlineto{\pgfqpoint{3.771583in}{1.476134in}}%
\pgfpathlineto{\pgfqpoint{3.777592in}{1.488751in}}%
\pgfpathclose%
\pgfusepath{fill}%
\end{pgfscope}%
\begin{pgfscope}%
\pgfpathrectangle{\pgfqpoint{3.525000in}{0.100000in}}{\pgfqpoint{2.862985in}{1.832778in}}%
\pgfusepath{clip}%
\pgfsetbuttcap%
\pgfsetmiterjoin%
\definecolor{currentfill}{rgb}{0.647059,0.833987,0.435294}%
\pgfsetfillcolor{currentfill}%
\pgfsetlinewidth{0.000000pt}%
\definecolor{currentstroke}{rgb}{0.000000,0.000000,0.000000}%
\pgfsetstrokecolor{currentstroke}%
\pgfsetstrokeopacity{0.000000}%
\pgfsetdash{}{0pt}%
\pgfpathmoveto{\pgfqpoint{5.611981in}{1.117748in}}%
\pgfpathlineto{\pgfqpoint{5.604011in}{1.195770in}}%
\pgfpathlineto{\pgfqpoint{5.594752in}{1.279623in}}%
\pgfpathlineto{\pgfqpoint{5.655390in}{1.288653in}}%
\pgfpathlineto{\pgfqpoint{5.671483in}{1.284465in}}%
\pgfpathlineto{\pgfqpoint{5.679202in}{1.279918in}}%
\pgfpathlineto{\pgfqpoint{5.688874in}{1.281179in}}%
\pgfpathlineto{\pgfqpoint{5.701627in}{1.273651in}}%
\pgfpathlineto{\pgfqpoint{5.725380in}{1.284855in}}%
\pgfpathlineto{\pgfqpoint{5.738489in}{1.285229in}}%
\pgfpathlineto{\pgfqpoint{5.753797in}{1.302307in}}%
\pgfpathlineto{\pgfqpoint{5.769372in}{1.312695in}}%
\pgfpathlineto{\pgfqpoint{5.790146in}{1.324661in}}%
\pgfpathlineto{\pgfqpoint{5.803519in}{1.241018in}}%
\pgfpathlineto{\pgfqpoint{5.797306in}{1.235646in}}%
\pgfpathlineto{\pgfqpoint{5.801317in}{1.230712in}}%
\pgfpathlineto{\pgfqpoint{5.802915in}{1.219896in}}%
\pgfpathlineto{\pgfqpoint{5.799803in}{1.209737in}}%
\pgfpathlineto{\pgfqpoint{5.799391in}{1.194887in}}%
\pgfpathlineto{\pgfqpoint{5.796468in}{1.175560in}}%
\pgfpathlineto{\pgfqpoint{5.781678in}{1.158322in}}%
\pgfpathlineto{\pgfqpoint{5.775452in}{1.154663in}}%
\pgfpathlineto{\pgfqpoint{5.770610in}{1.157813in}}%
\pgfpathlineto{\pgfqpoint{5.758627in}{1.141383in}}%
\pgfpathlineto{\pgfqpoint{5.759750in}{1.125564in}}%
\pgfpathlineto{\pgfqpoint{5.746406in}{1.129368in}}%
\pgfpathlineto{\pgfqpoint{5.740087in}{1.113569in}}%
\pgfpathlineto{\pgfqpoint{5.743281in}{1.102373in}}%
\pgfpathlineto{\pgfqpoint{5.738281in}{1.100719in}}%
\pgfpathlineto{\pgfqpoint{5.737581in}{1.091873in}}%
\pgfpathlineto{\pgfqpoint{5.725308in}{1.088305in}}%
\pgfpathlineto{\pgfqpoint{5.718930in}{1.095447in}}%
\pgfpathlineto{\pgfqpoint{5.710724in}{1.098217in}}%
\pgfpathlineto{\pgfqpoint{5.708796in}{1.105459in}}%
\pgfpathlineto{\pgfqpoint{5.701376in}{1.104185in}}%
\pgfpathlineto{\pgfqpoint{5.696515in}{1.097524in}}%
\pgfpathlineto{\pgfqpoint{5.689561in}{1.095181in}}%
\pgfpathlineto{\pgfqpoint{5.677277in}{1.099895in}}%
\pgfpathlineto{\pgfqpoint{5.670483in}{1.094286in}}%
\pgfpathlineto{\pgfqpoint{5.660827in}{1.100922in}}%
\pgfpathlineto{\pgfqpoint{5.642303in}{1.102959in}}%
\pgfpathlineto{\pgfqpoint{5.636715in}{1.115014in}}%
\pgfpathlineto{\pgfqpoint{5.629636in}{1.120355in}}%
\pgfpathlineto{\pgfqpoint{5.622773in}{1.116936in}}%
\pgfpathlineto{\pgfqpoint{5.611981in}{1.117748in}}%
\pgfpathclose%
\pgfusepath{fill}%
\end{pgfscope}%
\begin{pgfscope}%
\pgfpathrectangle{\pgfqpoint{3.525000in}{0.100000in}}{\pgfqpoint{2.862985in}{1.832778in}}%
\pgfusepath{clip}%
\pgfsetbuttcap%
\pgfsetmiterjoin%
\definecolor{currentfill}{rgb}{0.824145,0.926951,0.690504}%
\pgfsetfillcolor{currentfill}%
\pgfsetlinewidth{0.000000pt}%
\definecolor{currentstroke}{rgb}{0.000000,0.000000,0.000000}%
\pgfsetstrokecolor{currentstroke}%
\pgfsetstrokeopacity{0.000000}%
\pgfsetdash{}{0pt}%
\pgfpathmoveto{\pgfqpoint{5.411225in}{0.966856in}}%
\pgfpathlineto{\pgfqpoint{5.404121in}{0.972628in}}%
\pgfpathlineto{\pgfqpoint{5.398328in}{0.969882in}}%
\pgfpathlineto{\pgfqpoint{5.390922in}{0.983711in}}%
\pgfpathlineto{\pgfqpoint{5.394705in}{0.992402in}}%
\pgfpathlineto{\pgfqpoint{5.389252in}{1.002186in}}%
\pgfpathlineto{\pgfqpoint{5.388959in}{1.009888in}}%
\pgfpathlineto{\pgfqpoint{5.381591in}{1.012632in}}%
\pgfpathlineto{\pgfqpoint{5.378174in}{1.018436in}}%
\pgfpathlineto{\pgfqpoint{5.363768in}{1.025684in}}%
\pgfpathlineto{\pgfqpoint{5.351256in}{1.034614in}}%
\pgfpathlineto{\pgfqpoint{5.345442in}{1.041355in}}%
\pgfpathlineto{\pgfqpoint{5.345322in}{1.049578in}}%
\pgfpathlineto{\pgfqpoint{5.353102in}{1.065365in}}%
\pgfpathlineto{\pgfqpoint{5.355495in}{1.077440in}}%
\pgfpathlineto{\pgfqpoint{5.349157in}{1.084263in}}%
\pgfpathlineto{\pgfqpoint{5.340763in}{1.086834in}}%
\pgfpathlineto{\pgfqpoint{5.330579in}{1.081212in}}%
\pgfpathlineto{\pgfqpoint{5.326269in}{1.090863in}}%
\pgfpathlineto{\pgfqpoint{5.324094in}{1.103945in}}%
\pgfpathlineto{\pgfqpoint{5.309086in}{1.115607in}}%
\pgfpathlineto{\pgfqpoint{5.306087in}{1.120782in}}%
\pgfpathlineto{\pgfqpoint{5.292379in}{1.132498in}}%
\pgfpathlineto{\pgfqpoint{5.288083in}{1.141015in}}%
\pgfpathlineto{\pgfqpoint{5.284210in}{1.157906in}}%
\pgfpathlineto{\pgfqpoint{5.286865in}{1.172908in}}%
\pgfpathlineto{\pgfqpoint{5.290408in}{1.174999in}}%
\pgfpathlineto{\pgfqpoint{5.289771in}{1.187558in}}%
\pgfpathlineto{\pgfqpoint{5.299773in}{1.191288in}}%
\pgfpathlineto{\pgfqpoint{5.302793in}{1.202558in}}%
\pgfpathlineto{\pgfqpoint{5.308547in}{1.210143in}}%
\pgfpathlineto{\pgfqpoint{5.308256in}{1.219794in}}%
\pgfpathlineto{\pgfqpoint{5.301124in}{1.227469in}}%
\pgfpathlineto{\pgfqpoint{5.302818in}{1.238212in}}%
\pgfpathlineto{\pgfqpoint{5.321306in}{1.242884in}}%
\pgfpathlineto{\pgfqpoint{5.335490in}{1.251413in}}%
\pgfpathlineto{\pgfqpoint{5.336979in}{1.262151in}}%
\pgfpathlineto{\pgfqpoint{5.341914in}{1.265532in}}%
\pgfpathlineto{\pgfqpoint{5.343802in}{1.276813in}}%
\pgfpathlineto{\pgfqpoint{5.342285in}{1.284265in}}%
\pgfpathlineto{\pgfqpoint{5.332588in}{1.290454in}}%
\pgfpathlineto{\pgfqpoint{5.328684in}{1.299681in}}%
\pgfpathlineto{\pgfqpoint{5.319114in}{1.308600in}}%
\pgfpathlineto{\pgfqpoint{5.397758in}{1.311952in}}%
\pgfpathlineto{\pgfqpoint{5.450520in}{1.315700in}}%
\pgfpathlineto{\pgfqpoint{5.449554in}{1.304601in}}%
\pgfpathlineto{\pgfqpoint{5.458546in}{1.289293in}}%
\pgfpathlineto{\pgfqpoint{5.462322in}{1.276213in}}%
\pgfpathlineto{\pgfqpoint{5.466818in}{1.268787in}}%
\pgfpathlineto{\pgfqpoint{5.471726in}{1.207597in}}%
\pgfpathlineto{\pgfqpoint{5.478542in}{1.120373in}}%
\pgfpathlineto{\pgfqpoint{5.473283in}{1.107118in}}%
\pgfpathlineto{\pgfqpoint{5.480859in}{1.096224in}}%
\pgfpathlineto{\pgfqpoint{5.482971in}{1.085818in}}%
\pgfpathlineto{\pgfqpoint{5.472524in}{1.064241in}}%
\pgfpathlineto{\pgfqpoint{5.473140in}{1.061261in}}%
\pgfpathlineto{\pgfqpoint{5.462387in}{1.048893in}}%
\pgfpathlineto{\pgfqpoint{5.465336in}{1.045116in}}%
\pgfpathlineto{\pgfqpoint{5.460862in}{1.037407in}}%
\pgfpathlineto{\pgfqpoint{5.462013in}{1.021868in}}%
\pgfpathlineto{\pgfqpoint{5.456580in}{1.012333in}}%
\pgfpathlineto{\pgfqpoint{5.460999in}{1.001067in}}%
\pgfpathlineto{\pgfqpoint{5.442483in}{0.994944in}}%
\pgfpathlineto{\pgfqpoint{5.440777in}{0.988286in}}%
\pgfpathlineto{\pgfqpoint{5.445833in}{0.979847in}}%
\pgfpathlineto{\pgfqpoint{5.443506in}{0.974346in}}%
\pgfpathlineto{\pgfqpoint{5.423652in}{0.981141in}}%
\pgfpathlineto{\pgfqpoint{5.417101in}{0.981825in}}%
\pgfpathlineto{\pgfqpoint{5.408934in}{0.971495in}}%
\pgfpathlineto{\pgfqpoint{5.411225in}{0.966856in}}%
\pgfpathclose%
\pgfusepath{fill}%
\end{pgfscope}%
\begin{pgfscope}%
\pgfpathrectangle{\pgfqpoint{3.525000in}{0.100000in}}{\pgfqpoint{2.862985in}{1.832778in}}%
\pgfusepath{clip}%
\pgfsetbuttcap%
\pgfsetmiterjoin%
\definecolor{currentfill}{rgb}{0.802922,0.917724,0.656363}%
\pgfsetfillcolor{currentfill}%
\pgfsetlinewidth{0.000000pt}%
\definecolor{currentstroke}{rgb}{0.000000,0.000000,0.000000}%
\pgfsetstrokecolor{currentstroke}%
\pgfsetstrokeopacity{0.000000}%
\pgfsetdash{}{0pt}%
\pgfpathmoveto{\pgfqpoint{5.989161in}{1.156689in}}%
\pgfpathlineto{\pgfqpoint{5.983697in}{1.164803in}}%
\pgfpathlineto{\pgfqpoint{5.986780in}{1.169343in}}%
\pgfpathlineto{\pgfqpoint{5.994331in}{1.164244in}}%
\pgfpathlineto{\pgfqpoint{5.989161in}{1.156689in}}%
\pgfpathclose%
\pgfusepath{fill}%
\end{pgfscope}%
\begin{pgfscope}%
\pgfpathrectangle{\pgfqpoint{3.525000in}{0.100000in}}{\pgfqpoint{2.862985in}{1.832778in}}%
\pgfusepath{clip}%
\pgfsetbuttcap%
\pgfsetmiterjoin%
\definecolor{currentfill}{rgb}{0.824145,0.926951,0.690504}%
\pgfsetfillcolor{currentfill}%
\pgfsetlinewidth{0.000000pt}%
\definecolor{currentstroke}{rgb}{0.000000,0.000000,0.000000}%
\pgfsetstrokecolor{currentstroke}%
\pgfsetstrokeopacity{0.000000}%
\pgfsetdash{}{0pt}%
\pgfpathmoveto{\pgfqpoint{6.036873in}{1.226505in}}%
\pgfpathlineto{\pgfqpoint{6.040138in}{1.233404in}}%
\pgfpathlineto{\pgfqpoint{6.053338in}{1.234894in}}%
\pgfpathlineto{\pgfqpoint{6.051223in}{1.229031in}}%
\pgfpathlineto{\pgfqpoint{6.046868in}{1.221542in}}%
\pgfpathlineto{\pgfqpoint{6.049819in}{1.212619in}}%
\pgfpathlineto{\pgfqpoint{6.061465in}{1.201927in}}%
\pgfpathlineto{\pgfqpoint{6.064168in}{1.190668in}}%
\pgfpathlineto{\pgfqpoint{6.077585in}{1.176648in}}%
\pgfpathlineto{\pgfqpoint{6.082849in}{1.177248in}}%
\pgfpathlineto{\pgfqpoint{6.089406in}{1.156201in}}%
\pgfpathlineto{\pgfqpoint{6.088331in}{1.155991in}}%
\pgfpathlineto{\pgfqpoint{6.087135in}{1.155756in}}%
\pgfpathlineto{\pgfqpoint{6.057891in}{1.150158in}}%
\pgfpathlineto{\pgfqpoint{6.042245in}{1.205802in}}%
\pgfpathlineto{\pgfqpoint{6.036873in}{1.226505in}}%
\pgfpathclose%
\pgfusepath{fill}%
\end{pgfscope}%
\begin{pgfscope}%
\pgfpathrectangle{\pgfqpoint{3.525000in}{0.100000in}}{\pgfqpoint{2.862985in}{1.832778in}}%
\pgfusepath{clip}%
\pgfsetbuttcap%
\pgfsetmiterjoin%
\definecolor{currentfill}{rgb}{0.647059,0.833987,0.435294}%
\pgfsetfillcolor{currentfill}%
\pgfsetlinewidth{0.000000pt}%
\definecolor{currentstroke}{rgb}{0.000000,0.000000,0.000000}%
\pgfsetstrokecolor{currentstroke}%
\pgfsetstrokeopacity{0.000000}%
\pgfsetdash{}{0pt}%
\pgfpathmoveto{\pgfqpoint{5.763724in}{1.037893in}}%
\pgfpathlineto{\pgfqpoint{5.754483in}{1.038233in}}%
\pgfpathlineto{\pgfqpoint{5.745968in}{1.044098in}}%
\pgfpathlineto{\pgfqpoint{5.734465in}{1.061926in}}%
\pgfpathlineto{\pgfqpoint{5.724672in}{1.071367in}}%
\pgfpathlineto{\pgfqpoint{5.727231in}{1.078659in}}%
\pgfpathlineto{\pgfqpoint{5.725308in}{1.088305in}}%
\pgfpathlineto{\pgfqpoint{5.737581in}{1.091873in}}%
\pgfpathlineto{\pgfqpoint{5.738281in}{1.100719in}}%
\pgfpathlineto{\pgfqpoint{5.743281in}{1.102373in}}%
\pgfpathlineto{\pgfqpoint{5.740087in}{1.113569in}}%
\pgfpathlineto{\pgfqpoint{5.746406in}{1.129368in}}%
\pgfpathlineto{\pgfqpoint{5.759750in}{1.125564in}}%
\pgfpathlineto{\pgfqpoint{5.758627in}{1.141383in}}%
\pgfpathlineto{\pgfqpoint{5.770610in}{1.157813in}}%
\pgfpathlineto{\pgfqpoint{5.775452in}{1.154663in}}%
\pgfpathlineto{\pgfqpoint{5.781678in}{1.158322in}}%
\pgfpathlineto{\pgfqpoint{5.796468in}{1.175560in}}%
\pgfpathlineto{\pgfqpoint{5.799391in}{1.194887in}}%
\pgfpathlineto{\pgfqpoint{5.799803in}{1.209737in}}%
\pgfpathlineto{\pgfqpoint{5.802915in}{1.219896in}}%
\pgfpathlineto{\pgfqpoint{5.801317in}{1.230712in}}%
\pgfpathlineto{\pgfqpoint{5.797306in}{1.235646in}}%
\pgfpathlineto{\pgfqpoint{5.803519in}{1.241018in}}%
\pgfpathlineto{\pgfqpoint{5.812526in}{1.184275in}}%
\pgfpathlineto{\pgfqpoint{5.862178in}{1.192485in}}%
\pgfpathlineto{\pgfqpoint{5.867325in}{1.160097in}}%
\pgfpathlineto{\pgfqpoint{5.885300in}{1.181474in}}%
\pgfpathlineto{\pgfqpoint{5.889536in}{1.179343in}}%
\pgfpathlineto{\pgfqpoint{5.895767in}{1.191730in}}%
\pgfpathlineto{\pgfqpoint{5.904399in}{1.187846in}}%
\pgfpathlineto{\pgfqpoint{5.911925in}{1.188584in}}%
\pgfpathlineto{\pgfqpoint{5.916891in}{1.197195in}}%
\pgfpathlineto{\pgfqpoint{5.924098in}{1.201985in}}%
\pgfpathlineto{\pgfqpoint{5.934119in}{1.197766in}}%
\pgfpathlineto{\pgfqpoint{5.940712in}{1.199224in}}%
\pgfpathlineto{\pgfqpoint{5.950158in}{1.182877in}}%
\pgfpathlineto{\pgfqpoint{5.947427in}{1.170500in}}%
\pgfpathlineto{\pgfqpoint{5.922755in}{1.184435in}}%
\pgfpathlineto{\pgfqpoint{5.918134in}{1.172996in}}%
\pgfpathlineto{\pgfqpoint{5.919659in}{1.167732in}}%
\pgfpathlineto{\pgfqpoint{5.909159in}{1.149941in}}%
\pgfpathlineto{\pgfqpoint{5.904203in}{1.146402in}}%
\pgfpathlineto{\pgfqpoint{5.901965in}{1.138556in}}%
\pgfpathlineto{\pgfqpoint{5.893517in}{1.139377in}}%
\pgfpathlineto{\pgfqpoint{5.887449in}{1.117963in}}%
\pgfpathlineto{\pgfqpoint{5.884053in}{1.113049in}}%
\pgfpathlineto{\pgfqpoint{5.875326in}{1.114688in}}%
\pgfpathlineto{\pgfqpoint{5.866399in}{1.121442in}}%
\pgfpathlineto{\pgfqpoint{5.866091in}{1.111076in}}%
\pgfpathlineto{\pgfqpoint{5.857449in}{1.093654in}}%
\pgfpathlineto{\pgfqpoint{5.855300in}{1.081251in}}%
\pgfpathlineto{\pgfqpoint{5.848669in}{1.072991in}}%
\pgfpathlineto{\pgfqpoint{5.843643in}{1.059800in}}%
\pgfpathlineto{\pgfqpoint{5.843396in}{1.047589in}}%
\pgfpathlineto{\pgfqpoint{5.835633in}{1.045308in}}%
\pgfpathlineto{\pgfqpoint{5.826826in}{1.038395in}}%
\pgfpathlineto{\pgfqpoint{5.818347in}{1.037079in}}%
\pgfpathlineto{\pgfqpoint{5.816393in}{1.031228in}}%
\pgfpathlineto{\pgfqpoint{5.802730in}{1.025236in}}%
\pgfpathlineto{\pgfqpoint{5.795094in}{1.030341in}}%
\pgfpathlineto{\pgfqpoint{5.786560in}{1.020647in}}%
\pgfpathlineto{\pgfqpoint{5.771890in}{1.023505in}}%
\pgfpathlineto{\pgfqpoint{5.762893in}{1.033679in}}%
\pgfpathlineto{\pgfqpoint{5.763724in}{1.037893in}}%
\pgfpathclose%
\pgfusepath{fill}%
\end{pgfscope}%
\begin{pgfscope}%
\pgfpathrectangle{\pgfqpoint{3.525000in}{0.100000in}}{\pgfqpoint{2.862985in}{1.832778in}}%
\pgfusepath{clip}%
\pgfsetbuttcap%
\pgfsetmiterjoin%
\definecolor{currentfill}{rgb}{0.498039,0.737255,0.254902}%
\pgfsetfillcolor{currentfill}%
\pgfsetlinewidth{0.000000pt}%
\definecolor{currentstroke}{rgb}{0.000000,0.000000,0.000000}%
\pgfsetstrokecolor{currentstroke}%
\pgfsetstrokeopacity{0.000000}%
\pgfsetdash{}{0pt}%
\pgfpathmoveto{\pgfqpoint{6.087135in}{1.155756in}}%
\pgfpathlineto{\pgfqpoint{6.086760in}{1.144322in}}%
\pgfpathlineto{\pgfqpoint{6.082360in}{1.138701in}}%
\pgfpathlineto{\pgfqpoint{6.079551in}{1.126221in}}%
\pgfpathlineto{\pgfqpoint{6.066870in}{1.120469in}}%
\pgfpathlineto{\pgfqpoint{6.056229in}{1.118793in}}%
\pgfpathlineto{\pgfqpoint{6.059330in}{1.127000in}}%
\pgfpathlineto{\pgfqpoint{6.042947in}{1.133877in}}%
\pgfpathlineto{\pgfqpoint{6.029638in}{1.142487in}}%
\pgfpathlineto{\pgfqpoint{6.037878in}{1.155358in}}%
\pgfpathlineto{\pgfqpoint{6.031274in}{1.165363in}}%
\pgfpathlineto{\pgfqpoint{6.032051in}{1.173993in}}%
\pgfpathlineto{\pgfqpoint{6.027276in}{1.176362in}}%
\pgfpathlineto{\pgfqpoint{6.023402in}{1.190382in}}%
\pgfpathlineto{\pgfqpoint{6.032564in}{1.209388in}}%
\pgfpathlineto{\pgfqpoint{6.025681in}{1.212492in}}%
\pgfpathlineto{\pgfqpoint{6.023896in}{1.203115in}}%
\pgfpathlineto{\pgfqpoint{6.014070in}{1.200505in}}%
\pgfpathlineto{\pgfqpoint{6.015103in}{1.169655in}}%
\pgfpathlineto{\pgfqpoint{6.013306in}{1.159726in}}%
\pgfpathlineto{\pgfqpoint{6.018235in}{1.145575in}}%
\pgfpathlineto{\pgfqpoint{6.025818in}{1.138767in}}%
\pgfpathlineto{\pgfqpoint{6.022379in}{1.134492in}}%
\pgfpathlineto{\pgfqpoint{6.030119in}{1.128246in}}%
\pgfpathlineto{\pgfqpoint{6.032917in}{1.118107in}}%
\pgfpathlineto{\pgfqpoint{6.018730in}{1.126503in}}%
\pgfpathlineto{\pgfqpoint{6.009752in}{1.125412in}}%
\pgfpathlineto{\pgfqpoint{6.002780in}{1.134049in}}%
\pgfpathlineto{\pgfqpoint{5.995682in}{1.134890in}}%
\pgfpathlineto{\pgfqpoint{5.985594in}{1.130564in}}%
\pgfpathlineto{\pgfqpoint{5.981681in}{1.135955in}}%
\pgfpathlineto{\pgfqpoint{5.986823in}{1.147270in}}%
\pgfpathlineto{\pgfqpoint{5.989161in}{1.156689in}}%
\pgfpathlineto{\pgfqpoint{5.994331in}{1.164244in}}%
\pgfpathlineto{\pgfqpoint{5.986780in}{1.169343in}}%
\pgfpathlineto{\pgfqpoint{5.983697in}{1.164803in}}%
\pgfpathlineto{\pgfqpoint{5.976154in}{1.169411in}}%
\pgfpathlineto{\pgfqpoint{5.962796in}{1.172482in}}%
\pgfpathlineto{\pgfqpoint{5.964008in}{1.179234in}}%
\pgfpathlineto{\pgfqpoint{5.957953in}{1.183150in}}%
\pgfpathlineto{\pgfqpoint{5.950158in}{1.182877in}}%
\pgfpathlineto{\pgfqpoint{5.940712in}{1.199224in}}%
\pgfpathlineto{\pgfqpoint{5.934119in}{1.197766in}}%
\pgfpathlineto{\pgfqpoint{5.924098in}{1.201985in}}%
\pgfpathlineto{\pgfqpoint{5.916891in}{1.197195in}}%
\pgfpathlineto{\pgfqpoint{5.911925in}{1.188584in}}%
\pgfpathlineto{\pgfqpoint{5.904399in}{1.187846in}}%
\pgfpathlineto{\pgfqpoint{5.895767in}{1.191730in}}%
\pgfpathlineto{\pgfqpoint{5.889536in}{1.179343in}}%
\pgfpathlineto{\pgfqpoint{5.885300in}{1.181474in}}%
\pgfpathlineto{\pgfqpoint{5.867325in}{1.160097in}}%
\pgfpathlineto{\pgfqpoint{5.862178in}{1.192485in}}%
\pgfpathlineto{\pgfqpoint{5.927868in}{1.204659in}}%
\pgfpathlineto{\pgfqpoint{5.957316in}{1.209933in}}%
\pgfpathlineto{\pgfqpoint{6.000158in}{1.218561in}}%
\pgfpathlineto{\pgfqpoint{6.036873in}{1.226505in}}%
\pgfpathlineto{\pgfqpoint{6.042245in}{1.205802in}}%
\pgfpathlineto{\pgfqpoint{6.057891in}{1.150158in}}%
\pgfpathlineto{\pgfqpoint{6.087135in}{1.155756in}}%
\pgfpathclose%
\pgfusepath{fill}%
\end{pgfscope}%
\begin{pgfscope}%
\pgfpathrectangle{\pgfqpoint{3.525000in}{0.100000in}}{\pgfqpoint{2.862985in}{1.832778in}}%
\pgfusepath{clip}%
\pgfsetbuttcap%
\pgfsetmiterjoin%
\definecolor{currentfill}{rgb}{0.498039,0.737255,0.254902}%
\pgfsetfillcolor{currentfill}%
\pgfsetlinewidth{0.000000pt}%
\definecolor{currentstroke}{rgb}{0.000000,0.000000,0.000000}%
\pgfsetstrokecolor{currentstroke}%
\pgfsetstrokeopacity{0.000000}%
\pgfsetdash{}{0pt}%
\pgfpathmoveto{\pgfqpoint{4.764821in}{0.971059in}}%
\pgfpathlineto{\pgfqpoint{4.716895in}{0.975633in}}%
\pgfpathlineto{\pgfqpoint{4.667175in}{0.980039in}}%
\pgfpathlineto{\pgfqpoint{4.609722in}{0.986032in}}%
\pgfpathlineto{\pgfqpoint{4.524362in}{0.996184in}}%
\pgfpathlineto{\pgfqpoint{4.494080in}{1.000850in}}%
\pgfpathlineto{\pgfqpoint{4.415915in}{1.012202in}}%
\pgfpathlineto{\pgfqpoint{4.427328in}{1.083837in}}%
\pgfpathlineto{\pgfqpoint{4.427613in}{1.089632in}}%
\pgfpathlineto{\pgfqpoint{4.438608in}{1.158753in}}%
\pgfpathlineto{\pgfqpoint{4.446820in}{1.211396in}}%
\pgfpathlineto{\pgfqpoint{4.454571in}{1.260214in}}%
\pgfpathlineto{\pgfqpoint{4.507522in}{1.252613in}}%
\pgfpathlineto{\pgfqpoint{4.556892in}{1.245527in}}%
\pgfpathlineto{\pgfqpoint{4.647647in}{1.234326in}}%
\pgfpathlineto{\pgfqpoint{4.689283in}{1.230524in}}%
\pgfpathlineto{\pgfqpoint{4.755243in}{1.224124in}}%
\pgfpathlineto{\pgfqpoint{4.783777in}{1.221808in}}%
\pgfpathlineto{\pgfqpoint{4.778739in}{1.159333in}}%
\pgfpathlineto{\pgfqpoint{4.774188in}{1.099180in}}%
\pgfpathlineto{\pgfqpoint{4.770525in}{1.050201in}}%
\pgfpathlineto{\pgfqpoint{4.764821in}{0.971059in}}%
\pgfpathclose%
\pgfusepath{fill}%
\end{pgfscope}%
\begin{pgfscope}%
\pgfpathrectangle{\pgfqpoint{3.525000in}{0.100000in}}{\pgfqpoint{2.862985in}{1.832778in}}%
\pgfusepath{clip}%
\pgfsetbuttcap%
\pgfsetmiterjoin%
\definecolor{currentfill}{rgb}{0.647059,0.833987,0.435294}%
\pgfsetfillcolor{currentfill}%
\pgfsetlinewidth{0.000000pt}%
\definecolor{currentstroke}{rgb}{0.000000,0.000000,0.000000}%
\pgfsetstrokecolor{currentstroke}%
\pgfsetstrokeopacity{0.000000}%
\pgfsetdash{}{0pt}%
\pgfpathmoveto{\pgfqpoint{5.398847in}{0.935544in}}%
\pgfpathlineto{\pgfqpoint{5.400463in}{0.942786in}}%
\pgfpathlineto{\pgfqpoint{5.408837in}{0.941156in}}%
\pgfpathlineto{\pgfqpoint{5.411225in}{0.966856in}}%
\pgfpathlineto{\pgfqpoint{5.408934in}{0.971495in}}%
\pgfpathlineto{\pgfqpoint{5.417101in}{0.981825in}}%
\pgfpathlineto{\pgfqpoint{5.423652in}{0.981141in}}%
\pgfpathlineto{\pgfqpoint{5.443506in}{0.974346in}}%
\pgfpathlineto{\pgfqpoint{5.445833in}{0.979847in}}%
\pgfpathlineto{\pgfqpoint{5.440777in}{0.988286in}}%
\pgfpathlineto{\pgfqpoint{5.442483in}{0.994944in}}%
\pgfpathlineto{\pgfqpoint{5.460999in}{1.001067in}}%
\pgfpathlineto{\pgfqpoint{5.456580in}{1.012333in}}%
\pgfpathlineto{\pgfqpoint{5.462013in}{1.021868in}}%
\pgfpathlineto{\pgfqpoint{5.472127in}{1.027220in}}%
\pgfpathlineto{\pgfqpoint{5.493349in}{1.032520in}}%
\pgfpathlineto{\pgfqpoint{5.506739in}{1.024520in}}%
\pgfpathlineto{\pgfqpoint{5.510801in}{1.032321in}}%
\pgfpathlineto{\pgfqpoint{5.521919in}{1.037771in}}%
\pgfpathlineto{\pgfqpoint{5.524789in}{1.032965in}}%
\pgfpathlineto{\pgfqpoint{5.536209in}{1.036783in}}%
\pgfpathlineto{\pgfqpoint{5.535489in}{1.043318in}}%
\pgfpathlineto{\pgfqpoint{5.545757in}{1.050797in}}%
\pgfpathlineto{\pgfqpoint{5.551817in}{1.043031in}}%
\pgfpathlineto{\pgfqpoint{5.559790in}{1.042250in}}%
\pgfpathlineto{\pgfqpoint{5.565088in}{1.047311in}}%
\pgfpathlineto{\pgfqpoint{5.564498in}{1.054509in}}%
\pgfpathlineto{\pgfqpoint{5.569002in}{1.061656in}}%
\pgfpathlineto{\pgfqpoint{5.575037in}{1.063205in}}%
\pgfpathlineto{\pgfqpoint{5.577460in}{1.072642in}}%
\pgfpathlineto{\pgfqpoint{5.586237in}{1.080801in}}%
\pgfpathlineto{\pgfqpoint{5.583595in}{1.088922in}}%
\pgfpathlineto{\pgfqpoint{5.592134in}{1.092965in}}%
\pgfpathlineto{\pgfqpoint{5.597839in}{1.090478in}}%
\pgfpathlineto{\pgfqpoint{5.606265in}{1.096795in}}%
\pgfpathlineto{\pgfqpoint{5.613795in}{1.098441in}}%
\pgfpathlineto{\pgfqpoint{5.613822in}{1.104993in}}%
\pgfpathlineto{\pgfqpoint{5.608545in}{1.114085in}}%
\pgfpathlineto{\pgfqpoint{5.611981in}{1.117748in}}%
\pgfpathlineto{\pgfqpoint{5.622773in}{1.116936in}}%
\pgfpathlineto{\pgfqpoint{5.629636in}{1.120355in}}%
\pgfpathlineto{\pgfqpoint{5.636715in}{1.115014in}}%
\pgfpathlineto{\pgfqpoint{5.642303in}{1.102959in}}%
\pgfpathlineto{\pgfqpoint{5.660827in}{1.100922in}}%
\pgfpathlineto{\pgfqpoint{5.670483in}{1.094286in}}%
\pgfpathlineto{\pgfqpoint{5.677277in}{1.099895in}}%
\pgfpathlineto{\pgfqpoint{5.689561in}{1.095181in}}%
\pgfpathlineto{\pgfqpoint{5.696515in}{1.097524in}}%
\pgfpathlineto{\pgfqpoint{5.701376in}{1.104185in}}%
\pgfpathlineto{\pgfqpoint{5.708796in}{1.105459in}}%
\pgfpathlineto{\pgfqpoint{5.710724in}{1.098217in}}%
\pgfpathlineto{\pgfqpoint{5.718930in}{1.095447in}}%
\pgfpathlineto{\pgfqpoint{5.725308in}{1.088305in}}%
\pgfpathlineto{\pgfqpoint{5.727231in}{1.078659in}}%
\pgfpathlineto{\pgfqpoint{5.724672in}{1.071367in}}%
\pgfpathlineto{\pgfqpoint{5.734465in}{1.061926in}}%
\pgfpathlineto{\pgfqpoint{5.745968in}{1.044098in}}%
\pgfpathlineto{\pgfqpoint{5.754483in}{1.038233in}}%
\pgfpathlineto{\pgfqpoint{5.763724in}{1.037893in}}%
\pgfpathlineto{\pgfqpoint{5.746661in}{1.018314in}}%
\pgfpathlineto{\pgfqpoint{5.729865in}{1.006463in}}%
\pgfpathlineto{\pgfqpoint{5.723683in}{0.997053in}}%
\pgfpathlineto{\pgfqpoint{5.723788in}{0.991945in}}%
\pgfpathlineto{\pgfqpoint{5.714691in}{0.988030in}}%
\pgfpathlineto{\pgfqpoint{5.712084in}{0.980665in}}%
\pgfpathlineto{\pgfqpoint{5.693132in}{0.973246in}}%
\pgfpathlineto{\pgfqpoint{5.686414in}{0.968427in}}%
\pgfpathlineto{\pgfqpoint{5.685505in}{0.967399in}}%
\pgfpathlineto{\pgfqpoint{5.630991in}{0.962230in}}%
\pgfpathlineto{\pgfqpoint{5.598070in}{0.959461in}}%
\pgfpathlineto{\pgfqpoint{5.544049in}{0.956404in}}%
\pgfpathlineto{\pgfqpoint{5.476748in}{0.949746in}}%
\pgfpathlineto{\pgfqpoint{5.465629in}{0.951286in}}%
\pgfpathlineto{\pgfqpoint{5.467942in}{0.939936in}}%
\pgfpathlineto{\pgfqpoint{5.398847in}{0.935544in}}%
\pgfpathclose%
\pgfusepath{fill}%
\end{pgfscope}%
\begin{pgfscope}%
\pgfpathrectangle{\pgfqpoint{3.525000in}{0.100000in}}{\pgfqpoint{2.862985in}{1.832778in}}%
\pgfusepath{clip}%
\pgfsetbuttcap%
\pgfsetmiterjoin%
\definecolor{currentfill}{rgb}{0.428835,0.679123,0.210611}%
\pgfsetfillcolor{currentfill}%
\pgfsetlinewidth{0.000000pt}%
\definecolor{currentstroke}{rgb}{0.000000,0.000000,0.000000}%
\pgfsetstrokecolor{currentstroke}%
\pgfsetstrokeopacity{0.000000}%
\pgfsetdash{}{0pt}%
\pgfpathmoveto{\pgfqpoint{4.764821in}{0.971059in}}%
\pgfpathlineto{\pgfqpoint{4.770525in}{1.050201in}}%
\pgfpathlineto{\pgfqpoint{4.774188in}{1.099180in}}%
\pgfpathlineto{\pgfqpoint{4.778739in}{1.159333in}}%
\pgfpathlineto{\pgfqpoint{4.841835in}{1.154903in}}%
\pgfpathlineto{\pgfqpoint{4.921966in}{1.150526in}}%
\pgfpathlineto{\pgfqpoint{4.976466in}{1.148434in}}%
\pgfpathlineto{\pgfqpoint{5.030647in}{1.146774in}}%
\pgfpathlineto{\pgfqpoint{5.079686in}{1.146016in}}%
\pgfpathlineto{\pgfqpoint{5.102367in}{1.146251in}}%
\pgfpathlineto{\pgfqpoint{5.112350in}{1.138104in}}%
\pgfpathlineto{\pgfqpoint{5.120170in}{1.139746in}}%
\pgfpathlineto{\pgfqpoint{5.120414in}{1.132640in}}%
\pgfpathlineto{\pgfqpoint{5.114607in}{1.120348in}}%
\pgfpathlineto{\pgfqpoint{5.115237in}{1.112582in}}%
\pgfpathlineto{\pgfqpoint{5.120797in}{1.107461in}}%
\pgfpathlineto{\pgfqpoint{5.125905in}{1.096790in}}%
\pgfpathlineto{\pgfqpoint{5.136265in}{1.090665in}}%
\pgfpathlineto{\pgfqpoint{5.135919in}{1.050444in}}%
\pgfpathlineto{\pgfqpoint{5.136224in}{0.957954in}}%
\pgfpathlineto{\pgfqpoint{5.066779in}{0.958274in}}%
\pgfpathlineto{\pgfqpoint{4.993651in}{0.959554in}}%
\pgfpathlineto{\pgfqpoint{4.939818in}{0.961453in}}%
\pgfpathlineto{\pgfqpoint{4.894882in}{0.963187in}}%
\pgfpathlineto{\pgfqpoint{4.819179in}{0.967651in}}%
\pgfpathlineto{\pgfqpoint{4.764821in}{0.971059in}}%
\pgfpathclose%
\pgfusepath{fill}%
\end{pgfscope}%
\begin{pgfscope}%
\pgfpathrectangle{\pgfqpoint{3.525000in}{0.100000in}}{\pgfqpoint{2.862985in}{1.832778in}}%
\pgfusepath{clip}%
\pgfsetbuttcap%
\pgfsetmiterjoin%
\definecolor{currentfill}{rgb}{0.541869,0.765705,0.307958}%
\pgfsetfillcolor{currentfill}%
\pgfsetlinewidth{0.000000pt}%
\definecolor{currentstroke}{rgb}{0.000000,0.000000,0.000000}%
\pgfsetstrokecolor{currentstroke}%
\pgfsetstrokeopacity{0.000000}%
\pgfsetdash{}{0pt}%
\pgfpathmoveto{\pgfqpoint{5.786621in}{0.981538in}}%
\pgfpathlineto{\pgfqpoint{5.709707in}{0.970734in}}%
\pgfpathlineto{\pgfqpoint{5.686414in}{0.968427in}}%
\pgfpathlineto{\pgfqpoint{5.693132in}{0.973246in}}%
\pgfpathlineto{\pgfqpoint{5.712084in}{0.980665in}}%
\pgfpathlineto{\pgfqpoint{5.714691in}{0.988030in}}%
\pgfpathlineto{\pgfqpoint{5.723788in}{0.991945in}}%
\pgfpathlineto{\pgfqpoint{5.723683in}{0.997053in}}%
\pgfpathlineto{\pgfqpoint{5.729865in}{1.006463in}}%
\pgfpathlineto{\pgfqpoint{5.746661in}{1.018314in}}%
\pgfpathlineto{\pgfqpoint{5.763724in}{1.037893in}}%
\pgfpathlineto{\pgfqpoint{5.762893in}{1.033679in}}%
\pgfpathlineto{\pgfqpoint{5.771890in}{1.023505in}}%
\pgfpathlineto{\pgfqpoint{5.786560in}{1.020647in}}%
\pgfpathlineto{\pgfqpoint{5.795094in}{1.030341in}}%
\pgfpathlineto{\pgfqpoint{5.802730in}{1.025236in}}%
\pgfpathlineto{\pgfqpoint{5.816393in}{1.031228in}}%
\pgfpathlineto{\pgfqpoint{5.818347in}{1.037079in}}%
\pgfpathlineto{\pgfqpoint{5.826826in}{1.038395in}}%
\pgfpathlineto{\pgfqpoint{5.835633in}{1.045308in}}%
\pgfpathlineto{\pgfqpoint{5.843396in}{1.047589in}}%
\pgfpathlineto{\pgfqpoint{5.843643in}{1.059800in}}%
\pgfpathlineto{\pgfqpoint{5.848669in}{1.072991in}}%
\pgfpathlineto{\pgfqpoint{5.855300in}{1.081251in}}%
\pgfpathlineto{\pgfqpoint{5.857449in}{1.093654in}}%
\pgfpathlineto{\pgfqpoint{5.866091in}{1.111076in}}%
\pgfpathlineto{\pgfqpoint{5.866399in}{1.121442in}}%
\pgfpathlineto{\pgfqpoint{5.875326in}{1.114688in}}%
\pgfpathlineto{\pgfqpoint{5.884053in}{1.113049in}}%
\pgfpathlineto{\pgfqpoint{5.887449in}{1.117963in}}%
\pgfpathlineto{\pgfqpoint{5.893517in}{1.139377in}}%
\pgfpathlineto{\pgfqpoint{5.901965in}{1.138556in}}%
\pgfpathlineto{\pgfqpoint{5.904203in}{1.146402in}}%
\pgfpathlineto{\pgfqpoint{5.909159in}{1.149941in}}%
\pgfpathlineto{\pgfqpoint{5.919659in}{1.167732in}}%
\pgfpathlineto{\pgfqpoint{5.918134in}{1.172996in}}%
\pgfpathlineto{\pgfqpoint{5.922755in}{1.184435in}}%
\pgfpathlineto{\pgfqpoint{5.947427in}{1.170500in}}%
\pgfpathlineto{\pgfqpoint{5.950158in}{1.182877in}}%
\pgfpathlineto{\pgfqpoint{5.957953in}{1.183150in}}%
\pgfpathlineto{\pgfqpoint{5.964008in}{1.179234in}}%
\pgfpathlineto{\pgfqpoint{5.962796in}{1.172482in}}%
\pgfpathlineto{\pgfqpoint{5.976154in}{1.169411in}}%
\pgfpathlineto{\pgfqpoint{5.983697in}{1.164803in}}%
\pgfpathlineto{\pgfqpoint{5.989161in}{1.156689in}}%
\pgfpathlineto{\pgfqpoint{5.986823in}{1.147270in}}%
\pgfpathlineto{\pgfqpoint{5.982103in}{1.146498in}}%
\pgfpathlineto{\pgfqpoint{5.979365in}{1.132285in}}%
\pgfpathlineto{\pgfqpoint{5.985362in}{1.126728in}}%
\pgfpathlineto{\pgfqpoint{5.993804in}{1.131220in}}%
\pgfpathlineto{\pgfqpoint{6.001637in}{1.121734in}}%
\pgfpathlineto{\pgfqpoint{6.019137in}{1.120029in}}%
\pgfpathlineto{\pgfqpoint{6.022152in}{1.114573in}}%
\pgfpathlineto{\pgfqpoint{6.038348in}{1.109266in}}%
\pgfpathlineto{\pgfqpoint{6.036349in}{1.103004in}}%
\pgfpathlineto{\pgfqpoint{6.038456in}{1.085305in}}%
\pgfpathlineto{\pgfqpoint{6.044993in}{1.078622in}}%
\pgfpathlineto{\pgfqpoint{6.035347in}{1.075802in}}%
\pgfpathlineto{\pgfqpoint{6.036626in}{1.068246in}}%
\pgfpathlineto{\pgfqpoint{6.046916in}{1.061840in}}%
\pgfpathlineto{\pgfqpoint{6.047844in}{1.055521in}}%
\pgfpathlineto{\pgfqpoint{6.042034in}{1.050772in}}%
\pgfpathlineto{\pgfqpoint{6.030312in}{1.062022in}}%
\pgfpathlineto{\pgfqpoint{6.029153in}{1.053814in}}%
\pgfpathlineto{\pgfqpoint{6.038966in}{1.049911in}}%
\pgfpathlineto{\pgfqpoint{6.039948in}{1.045875in}}%
\pgfpathlineto{\pgfqpoint{6.053408in}{1.051209in}}%
\pgfpathlineto{\pgfqpoint{6.063722in}{1.052620in}}%
\pgfpathlineto{\pgfqpoint{6.074287in}{1.031303in}}%
\pgfpathlineto{\pgfqpoint{6.073108in}{1.031073in}}%
\pgfpathlineto{\pgfqpoint{6.068334in}{1.030086in}}%
\pgfpathlineto{\pgfqpoint{6.066927in}{1.029792in}}%
\pgfpathlineto{\pgfqpoint{6.065997in}{1.029610in}}%
\pgfpathlineto{\pgfqpoint{6.003091in}{1.016561in}}%
\pgfpathlineto{\pgfqpoint{5.946826in}{1.005045in}}%
\pgfpathlineto{\pgfqpoint{5.869033in}{0.991644in}}%
\pgfpathlineto{\pgfqpoint{5.802964in}{0.982918in}}%
\pgfpathlineto{\pgfqpoint{5.786621in}{0.981538in}}%
\pgfpathclose%
\pgfusepath{fill}%
\end{pgfscope}%
\begin{pgfscope}%
\pgfpathrectangle{\pgfqpoint{3.525000in}{0.100000in}}{\pgfqpoint{2.862985in}{1.832778in}}%
\pgfusepath{clip}%
\pgfsetbuttcap%
\pgfsetmiterjoin%
\definecolor{currentfill}{rgb}{0.541869,0.765705,0.307958}%
\pgfsetfillcolor{currentfill}%
\pgfsetlinewidth{0.000000pt}%
\definecolor{currentstroke}{rgb}{0.000000,0.000000,0.000000}%
\pgfsetstrokecolor{currentstroke}%
\pgfsetstrokeopacity{0.000000}%
\pgfsetdash{}{0pt}%
\pgfpathmoveto{\pgfqpoint{6.066870in}{1.120469in}}%
\pgfpathlineto{\pgfqpoint{6.079551in}{1.126221in}}%
\pgfpathlineto{\pgfqpoint{6.069424in}{1.096596in}}%
\pgfpathlineto{\pgfqpoint{6.062727in}{1.080908in}}%
\pgfpathlineto{\pgfqpoint{6.057902in}{1.089786in}}%
\pgfpathlineto{\pgfqpoint{6.066482in}{1.111041in}}%
\pgfpathlineto{\pgfqpoint{6.066870in}{1.120469in}}%
\pgfpathclose%
\pgfusepath{fill}%
\end{pgfscope}%
\begin{pgfscope}%
\pgfpathrectangle{\pgfqpoint{3.525000in}{0.100000in}}{\pgfqpoint{2.862985in}{1.832778in}}%
\pgfusepath{clip}%
\pgfsetbuttcap%
\pgfsetmiterjoin%
\definecolor{currentfill}{rgb}{0.382699,0.640369,0.181084}%
\pgfsetfillcolor{currentfill}%
\pgfsetlinewidth{0.000000pt}%
\definecolor{currentstroke}{rgb}{0.000000,0.000000,0.000000}%
\pgfsetstrokecolor{currentstroke}%
\pgfsetstrokeopacity{0.000000}%
\pgfsetdash{}{0pt}%
\pgfpathmoveto{\pgfqpoint{5.398847in}{0.935544in}}%
\pgfpathlineto{\pgfqpoint{5.395785in}{0.935101in}}%
\pgfpathlineto{\pgfqpoint{5.392896in}{0.934898in}}%
\pgfpathlineto{\pgfqpoint{5.394133in}{0.926031in}}%
\pgfpathlineto{\pgfqpoint{5.386682in}{0.917093in}}%
\pgfpathlineto{\pgfqpoint{5.385213in}{0.903107in}}%
\pgfpathlineto{\pgfqpoint{5.351904in}{0.900644in}}%
\pgfpathlineto{\pgfqpoint{5.354808in}{0.907211in}}%
\pgfpathlineto{\pgfqpoint{5.366818in}{0.919204in}}%
\pgfpathlineto{\pgfqpoint{5.367150in}{0.926155in}}%
\pgfpathlineto{\pgfqpoint{5.361833in}{0.932737in}}%
\pgfpathlineto{\pgfqpoint{5.312257in}{0.930116in}}%
\pgfpathlineto{\pgfqpoint{5.225580in}{0.927364in}}%
\pgfpathlineto{\pgfqpoint{5.174856in}{0.926440in}}%
\pgfpathlineto{\pgfqpoint{5.136515in}{0.926094in}}%
\pgfpathlineto{\pgfqpoint{5.136224in}{0.957954in}}%
\pgfpathlineto{\pgfqpoint{5.135919in}{1.050444in}}%
\pgfpathlineto{\pgfqpoint{5.136265in}{1.090665in}}%
\pgfpathlineto{\pgfqpoint{5.125905in}{1.096790in}}%
\pgfpathlineto{\pgfqpoint{5.120797in}{1.107461in}}%
\pgfpathlineto{\pgfqpoint{5.115237in}{1.112582in}}%
\pgfpathlineto{\pgfqpoint{5.114607in}{1.120348in}}%
\pgfpathlineto{\pgfqpoint{5.120414in}{1.132640in}}%
\pgfpathlineto{\pgfqpoint{5.120170in}{1.139746in}}%
\pgfpathlineto{\pgfqpoint{5.112350in}{1.138104in}}%
\pgfpathlineto{\pgfqpoint{5.102367in}{1.146251in}}%
\pgfpathlineto{\pgfqpoint{5.094364in}{1.160551in}}%
\pgfpathlineto{\pgfqpoint{5.087651in}{1.167151in}}%
\pgfpathlineto{\pgfqpoint{5.080637in}{1.183365in}}%
\pgfpathlineto{\pgfqpoint{5.153489in}{1.182229in}}%
\pgfpathlineto{\pgfqpoint{5.225906in}{1.184615in}}%
\pgfpathlineto{\pgfqpoint{5.272339in}{1.187293in}}%
\pgfpathlineto{\pgfqpoint{5.286865in}{1.172908in}}%
\pgfpathlineto{\pgfqpoint{5.284210in}{1.157906in}}%
\pgfpathlineto{\pgfqpoint{5.288083in}{1.141015in}}%
\pgfpathlineto{\pgfqpoint{5.292379in}{1.132498in}}%
\pgfpathlineto{\pgfqpoint{5.306087in}{1.120782in}}%
\pgfpathlineto{\pgfqpoint{5.309086in}{1.115607in}}%
\pgfpathlineto{\pgfqpoint{5.324094in}{1.103945in}}%
\pgfpathlineto{\pgfqpoint{5.326269in}{1.090863in}}%
\pgfpathlineto{\pgfqpoint{5.330579in}{1.081212in}}%
\pgfpathlineto{\pgfqpoint{5.340763in}{1.086834in}}%
\pgfpathlineto{\pgfqpoint{5.349157in}{1.084263in}}%
\pgfpathlineto{\pgfqpoint{5.355495in}{1.077440in}}%
\pgfpathlineto{\pgfqpoint{5.353102in}{1.065365in}}%
\pgfpathlineto{\pgfqpoint{5.345322in}{1.049578in}}%
\pgfpathlineto{\pgfqpoint{5.345442in}{1.041355in}}%
\pgfpathlineto{\pgfqpoint{5.351256in}{1.034614in}}%
\pgfpathlineto{\pgfqpoint{5.363768in}{1.025684in}}%
\pgfpathlineto{\pgfqpoint{5.378174in}{1.018436in}}%
\pgfpathlineto{\pgfqpoint{5.381591in}{1.012632in}}%
\pgfpathlineto{\pgfqpoint{5.388959in}{1.009888in}}%
\pgfpathlineto{\pgfqpoint{5.389252in}{1.002186in}}%
\pgfpathlineto{\pgfqpoint{5.394705in}{0.992402in}}%
\pgfpathlineto{\pgfqpoint{5.390922in}{0.983711in}}%
\pgfpathlineto{\pgfqpoint{5.398328in}{0.969882in}}%
\pgfpathlineto{\pgfqpoint{5.404121in}{0.972628in}}%
\pgfpathlineto{\pgfqpoint{5.411225in}{0.966856in}}%
\pgfpathlineto{\pgfqpoint{5.408837in}{0.941156in}}%
\pgfpathlineto{\pgfqpoint{5.400463in}{0.942786in}}%
\pgfpathlineto{\pgfqpoint{5.398847in}{0.935544in}}%
\pgfpathclose%
\pgfusepath{fill}%
\end{pgfscope}%
\begin{pgfscope}%
\pgfpathrectangle{\pgfqpoint{3.525000in}{0.100000in}}{\pgfqpoint{2.862985in}{1.832778in}}%
\pgfusepath{clip}%
\pgfsetbuttcap%
\pgfsetmiterjoin%
\definecolor{currentfill}{rgb}{0.690888,0.862438,0.488351}%
\pgfsetfillcolor{currentfill}%
\pgfsetlinewidth{0.000000pt}%
\definecolor{currentstroke}{rgb}{0.000000,0.000000,0.000000}%
\pgfsetstrokecolor{currentstroke}%
\pgfsetstrokeopacity{0.000000}%
\pgfsetdash{}{0pt}%
\pgfpathmoveto{\pgfqpoint{4.114569in}{0.941085in}}%
\pgfpathlineto{\pgfqpoint{4.120633in}{0.953995in}}%
\pgfpathlineto{\pgfqpoint{4.118973in}{0.973374in}}%
\pgfpathlineto{\pgfqpoint{4.121877in}{0.978886in}}%
\pgfpathlineto{\pgfqpoint{4.122164in}{1.003024in}}%
\pgfpathlineto{\pgfqpoint{4.124916in}{1.009976in}}%
\pgfpathlineto{\pgfqpoint{4.134613in}{1.011057in}}%
\pgfpathlineto{\pgfqpoint{4.143612in}{1.007971in}}%
\pgfpathlineto{\pgfqpoint{4.147531in}{0.999473in}}%
\pgfpathlineto{\pgfqpoint{4.153025in}{0.999798in}}%
\pgfpathlineto{\pgfqpoint{4.159852in}{1.009534in}}%
\pgfpathlineto{\pgfqpoint{4.169719in}{1.057523in}}%
\pgfpathlineto{\pgfqpoint{4.225868in}{1.045930in}}%
\pgfpathlineto{\pgfqpoint{4.258455in}{1.039504in}}%
\pgfpathlineto{\pgfqpoint{4.344991in}{1.024292in}}%
\pgfpathlineto{\pgfqpoint{4.368932in}{1.019484in}}%
\pgfpathlineto{\pgfqpoint{4.415915in}{1.012202in}}%
\pgfpathlineto{\pgfqpoint{4.406281in}{0.950172in}}%
\pgfpathlineto{\pgfqpoint{4.396262in}{0.885469in}}%
\pgfpathlineto{\pgfqpoint{4.384723in}{0.812676in}}%
\pgfpathlineto{\pgfqpoint{4.371747in}{0.729118in}}%
\pgfpathlineto{\pgfqpoint{4.361264in}{0.660490in}}%
\pgfpathlineto{\pgfqpoint{4.286176in}{0.672414in}}%
\pgfpathlineto{\pgfqpoint{4.253185in}{0.678064in}}%
\pgfpathlineto{\pgfqpoint{4.238447in}{0.686887in}}%
\pgfpathlineto{\pgfqpoint{4.142356in}{0.744864in}}%
\pgfpathlineto{\pgfqpoint{4.070244in}{0.788862in}}%
\pgfpathlineto{\pgfqpoint{4.072646in}{0.796660in}}%
\pgfpathlineto{\pgfqpoint{4.078603in}{0.802117in}}%
\pgfpathlineto{\pgfqpoint{4.084798in}{0.801105in}}%
\pgfpathlineto{\pgfqpoint{4.093785in}{0.806836in}}%
\pgfpathlineto{\pgfqpoint{4.095200in}{0.815071in}}%
\pgfpathlineto{\pgfqpoint{4.084240in}{0.825061in}}%
\pgfpathlineto{\pgfqpoint{4.092035in}{0.844239in}}%
\pgfpathlineto{\pgfqpoint{4.099871in}{0.851614in}}%
\pgfpathlineto{\pgfqpoint{4.103591in}{0.860364in}}%
\pgfpathlineto{\pgfqpoint{4.105882in}{0.876407in}}%
\pgfpathlineto{\pgfqpoint{4.113230in}{0.883673in}}%
\pgfpathlineto{\pgfqpoint{4.128654in}{0.890930in}}%
\pgfpathlineto{\pgfqpoint{4.128313in}{0.897324in}}%
\pgfpathlineto{\pgfqpoint{4.119699in}{0.905259in}}%
\pgfpathlineto{\pgfqpoint{4.118526in}{0.921617in}}%
\pgfpathlineto{\pgfqpoint{4.112595in}{0.933576in}}%
\pgfpathlineto{\pgfqpoint{4.114569in}{0.941085in}}%
\pgfpathclose%
\pgfusepath{fill}%
\end{pgfscope}%
\begin{pgfscope}%
\pgfpathrectangle{\pgfqpoint{3.525000in}{0.100000in}}{\pgfqpoint{2.862985in}{1.832778in}}%
\pgfusepath{clip}%
\pgfsetbuttcap%
\pgfsetmiterjoin%
\definecolor{currentfill}{rgb}{0.524337,0.754325,0.286736}%
\pgfsetfillcolor{currentfill}%
\pgfsetlinewidth{0.000000pt}%
\definecolor{currentstroke}{rgb}{0.000000,0.000000,0.000000}%
\pgfsetstrokecolor{currentstroke}%
\pgfsetstrokeopacity{0.000000}%
\pgfsetdash{}{0pt}%
\pgfpathmoveto{\pgfqpoint{5.144643in}{0.746495in}}%
\pgfpathlineto{\pgfqpoint{5.130724in}{0.750800in}}%
\pgfpathlineto{\pgfqpoint{5.112888in}{0.764459in}}%
\pgfpathlineto{\pgfqpoint{5.104966in}{0.767394in}}%
\pgfpathlineto{\pgfqpoint{5.099932in}{0.761495in}}%
\pgfpathlineto{\pgfqpoint{5.093575in}{0.761197in}}%
\pgfpathlineto{\pgfqpoint{5.085532in}{0.766205in}}%
\pgfpathlineto{\pgfqpoint{5.072911in}{0.759788in}}%
\pgfpathlineto{\pgfqpoint{5.068418in}{0.762948in}}%
\pgfpathlineto{\pgfqpoint{5.055999in}{0.755462in}}%
\pgfpathlineto{\pgfqpoint{5.042897in}{0.756843in}}%
\pgfpathlineto{\pgfqpoint{5.032791in}{0.761708in}}%
\pgfpathlineto{\pgfqpoint{5.025714in}{0.759877in}}%
\pgfpathlineto{\pgfqpoint{5.014386in}{0.766758in}}%
\pgfpathlineto{\pgfqpoint{5.008084in}{0.754622in}}%
\pgfpathlineto{\pgfqpoint{5.001640in}{0.765059in}}%
\pgfpathlineto{\pgfqpoint{4.988911in}{0.761010in}}%
\pgfpathlineto{\pgfqpoint{4.977778in}{0.770922in}}%
\pgfpathlineto{\pgfqpoint{4.968049in}{0.762933in}}%
\pgfpathlineto{\pgfqpoint{4.963177in}{0.770275in}}%
\pgfpathlineto{\pgfqpoint{4.956154in}{0.772659in}}%
\pgfpathlineto{\pgfqpoint{4.951877in}{0.779738in}}%
\pgfpathlineto{\pgfqpoint{4.942692in}{0.781753in}}%
\pgfpathlineto{\pgfqpoint{4.937392in}{0.776425in}}%
\pgfpathlineto{\pgfqpoint{4.928387in}{0.783330in}}%
\pgfpathlineto{\pgfqpoint{4.924189in}{0.781753in}}%
\pgfpathlineto{\pgfqpoint{4.909265in}{0.787349in}}%
\pgfpathlineto{\pgfqpoint{4.899917in}{0.787973in}}%
\pgfpathlineto{\pgfqpoint{4.895733in}{0.799856in}}%
\pgfpathlineto{\pgfqpoint{4.888240in}{0.798376in}}%
\pgfpathlineto{\pgfqpoint{4.879657in}{0.801308in}}%
\pgfpathlineto{\pgfqpoint{4.874005in}{0.799614in}}%
\pgfpathlineto{\pgfqpoint{4.861884in}{0.812959in}}%
\pgfpathlineto{\pgfqpoint{4.858507in}{0.812086in}}%
\pgfpathlineto{\pgfqpoint{4.861574in}{0.866207in}}%
\pgfpathlineto{\pgfqpoint{4.864916in}{0.933195in}}%
\pgfpathlineto{\pgfqpoint{4.810064in}{0.936263in}}%
\pgfpathlineto{\pgfqpoint{4.755936in}{0.940356in}}%
\pgfpathlineto{\pgfqpoint{4.714118in}{0.943987in}}%
\pgfpathlineto{\pgfqpoint{4.716895in}{0.975633in}}%
\pgfpathlineto{\pgfqpoint{4.764821in}{0.971059in}}%
\pgfpathlineto{\pgfqpoint{4.819179in}{0.967651in}}%
\pgfpathlineto{\pgfqpoint{4.894882in}{0.963187in}}%
\pgfpathlineto{\pgfqpoint{4.939818in}{0.961453in}}%
\pgfpathlineto{\pgfqpoint{4.993651in}{0.959554in}}%
\pgfpathlineto{\pgfqpoint{5.066779in}{0.958274in}}%
\pgfpathlineto{\pgfqpoint{5.136224in}{0.957954in}}%
\pgfpathlineto{\pgfqpoint{5.136515in}{0.926094in}}%
\pgfpathlineto{\pgfqpoint{5.140413in}{0.902092in}}%
\pgfpathlineto{\pgfqpoint{5.146468in}{0.857760in}}%
\pgfpathlineto{\pgfqpoint{5.145219in}{0.782052in}}%
\pgfpathlineto{\pgfqpoint{5.144643in}{0.746495in}}%
\pgfpathclose%
\pgfusepath{fill}%
\end{pgfscope}%
\begin{pgfscope}%
\pgfpathrectangle{\pgfqpoint{3.525000in}{0.100000in}}{\pgfqpoint{2.862985in}{1.832778in}}%
\pgfusepath{clip}%
\pgfsetbuttcap%
\pgfsetmiterjoin%
\definecolor{currentfill}{rgb}{0.690888,0.862438,0.488351}%
\pgfsetfillcolor{currentfill}%
\pgfsetlinewidth{0.000000pt}%
\definecolor{currentstroke}{rgb}{0.000000,0.000000,0.000000}%
\pgfsetstrokecolor{currentstroke}%
\pgfsetstrokeopacity{0.000000}%
\pgfsetdash{}{0pt}%
\pgfpathmoveto{\pgfqpoint{5.665820in}{0.864070in}}%
\pgfpathlineto{\pgfqpoint{5.665861in}{0.878097in}}%
\pgfpathlineto{\pgfqpoint{5.678051in}{0.883484in}}%
\pgfpathlineto{\pgfqpoint{5.678562in}{0.892095in}}%
\pgfpathlineto{\pgfqpoint{5.689476in}{0.902614in}}%
\pgfpathlineto{\pgfqpoint{5.703119in}{0.904768in}}%
\pgfpathlineto{\pgfqpoint{5.715113in}{0.916361in}}%
\pgfpathlineto{\pgfqpoint{5.727639in}{0.921543in}}%
\pgfpathlineto{\pgfqpoint{5.737040in}{0.937098in}}%
\pgfpathlineto{\pgfqpoint{5.745602in}{0.935969in}}%
\pgfpathlineto{\pgfqpoint{5.763697in}{0.950142in}}%
\pgfpathlineto{\pgfqpoint{5.773253in}{0.950409in}}%
\pgfpathlineto{\pgfqpoint{5.777288in}{0.961208in}}%
\pgfpathlineto{\pgfqpoint{5.784883in}{0.968727in}}%
\pgfpathlineto{\pgfqpoint{5.786621in}{0.981538in}}%
\pgfpathlineto{\pgfqpoint{5.802964in}{0.982918in}}%
\pgfpathlineto{\pgfqpoint{5.869033in}{0.991644in}}%
\pgfpathlineto{\pgfqpoint{5.946826in}{1.005045in}}%
\pgfpathlineto{\pgfqpoint{6.003091in}{1.016561in}}%
\pgfpathlineto{\pgfqpoint{6.065997in}{1.029610in}}%
\pgfpathlineto{\pgfqpoint{6.083973in}{1.004931in}}%
\pgfpathlineto{\pgfqpoint{6.074236in}{1.006506in}}%
\pgfpathlineto{\pgfqpoint{6.061737in}{1.003461in}}%
\pgfpathlineto{\pgfqpoint{6.049421in}{0.990884in}}%
\pgfpathlineto{\pgfqpoint{6.040570in}{0.991789in}}%
\pgfpathlineto{\pgfqpoint{6.039450in}{0.984319in}}%
\pgfpathlineto{\pgfqpoint{6.055449in}{0.990226in}}%
\pgfpathlineto{\pgfqpoint{6.071548in}{0.992662in}}%
\pgfpathlineto{\pgfqpoint{6.077526in}{0.989421in}}%
\pgfpathlineto{\pgfqpoint{6.085571in}{0.993116in}}%
\pgfpathlineto{\pgfqpoint{6.089722in}{0.990527in}}%
\pgfpathlineto{\pgfqpoint{6.093396in}{0.978199in}}%
\pgfpathlineto{\pgfqpoint{6.085750in}{0.974382in}}%
\pgfpathlineto{\pgfqpoint{6.080502in}{0.959350in}}%
\pgfpathlineto{\pgfqpoint{6.075000in}{0.953497in}}%
\pgfpathlineto{\pgfqpoint{6.058132in}{0.954777in}}%
\pgfpathlineto{\pgfqpoint{6.059039in}{0.963604in}}%
\pgfpathlineto{\pgfqpoint{6.049827in}{0.959748in}}%
\pgfpathlineto{\pgfqpoint{6.047823in}{0.952387in}}%
\pgfpathlineto{\pgfqpoint{6.054137in}{0.944761in}}%
\pgfpathlineto{\pgfqpoint{6.056282in}{0.931819in}}%
\pgfpathlineto{\pgfqpoint{6.045950in}{0.924444in}}%
\pgfpathlineto{\pgfqpoint{6.057125in}{0.921851in}}%
\pgfpathlineto{\pgfqpoint{6.062181in}{0.927273in}}%
\pgfpathlineto{\pgfqpoint{6.073362in}{0.927934in}}%
\pgfpathlineto{\pgfqpoint{6.067616in}{0.916224in}}%
\pgfpathlineto{\pgfqpoint{6.060579in}{0.910582in}}%
\pgfpathlineto{\pgfqpoint{6.039292in}{0.904933in}}%
\pgfpathlineto{\pgfqpoint{6.017436in}{0.885111in}}%
\pgfpathlineto{\pgfqpoint{6.003960in}{0.865580in}}%
\pgfpathlineto{\pgfqpoint{6.003880in}{0.857648in}}%
\pgfpathlineto{\pgfqpoint{5.998497in}{0.846702in}}%
\pgfpathlineto{\pgfqpoint{5.970869in}{0.839501in}}%
\pgfpathlineto{\pgfqpoint{5.905165in}{0.886600in}}%
\pgfpathlineto{\pgfqpoint{5.847304in}{0.878098in}}%
\pgfpathlineto{\pgfqpoint{5.846819in}{0.885949in}}%
\pgfpathlineto{\pgfqpoint{5.838024in}{0.894789in}}%
\pgfpathlineto{\pgfqpoint{5.831366in}{0.896952in}}%
\pgfpathlineto{\pgfqpoint{5.768499in}{0.890420in}}%
\pgfpathlineto{\pgfqpoint{5.754049in}{0.885519in}}%
\pgfpathlineto{\pgfqpoint{5.727958in}{0.872540in}}%
\pgfpathlineto{\pgfqpoint{5.705410in}{0.868948in}}%
\pgfpathlineto{\pgfqpoint{5.665820in}{0.864070in}}%
\pgfpathclose%
\pgfusepath{fill}%
\end{pgfscope}%
\begin{pgfscope}%
\pgfpathrectangle{\pgfqpoint{3.525000in}{0.100000in}}{\pgfqpoint{2.862985in}{1.832778in}}%
\pgfusepath{clip}%
\pgfsetbuttcap%
\pgfsetmiterjoin%
\definecolor{currentfill}{rgb}{0.594464,0.799846,0.371626}%
\pgfsetfillcolor{currentfill}%
\pgfsetlinewidth{0.000000pt}%
\definecolor{currentstroke}{rgb}{0.000000,0.000000,0.000000}%
\pgfsetstrokecolor{currentstroke}%
\pgfsetstrokeopacity{0.000000}%
\pgfsetdash{}{0pt}%
\pgfpathmoveto{\pgfqpoint{5.665820in}{0.864070in}}%
\pgfpathlineto{\pgfqpoint{5.600022in}{0.856843in}}%
\pgfpathlineto{\pgfqpoint{5.539830in}{0.851430in}}%
\pgfpathlineto{\pgfqpoint{5.477965in}{0.847429in}}%
\pgfpathlineto{\pgfqpoint{5.467346in}{0.845890in}}%
\pgfpathlineto{\pgfqpoint{5.425643in}{0.842683in}}%
\pgfpathlineto{\pgfqpoint{5.358848in}{0.838785in}}%
\pgfpathlineto{\pgfqpoint{5.370732in}{0.848649in}}%
\pgfpathlineto{\pgfqpoint{5.368020in}{0.859033in}}%
\pgfpathlineto{\pgfqpoint{5.370647in}{0.871619in}}%
\pgfpathlineto{\pgfqpoint{5.374660in}{0.877545in}}%
\pgfpathlineto{\pgfqpoint{5.374509in}{0.885789in}}%
\pgfpathlineto{\pgfqpoint{5.385195in}{0.890974in}}%
\pgfpathlineto{\pgfqpoint{5.385213in}{0.903107in}}%
\pgfpathlineto{\pgfqpoint{5.386682in}{0.917093in}}%
\pgfpathlineto{\pgfqpoint{5.394133in}{0.926031in}}%
\pgfpathlineto{\pgfqpoint{5.392896in}{0.934898in}}%
\pgfpathlineto{\pgfqpoint{5.395785in}{0.935101in}}%
\pgfpathlineto{\pgfqpoint{5.398847in}{0.935544in}}%
\pgfpathlineto{\pgfqpoint{5.467942in}{0.939936in}}%
\pgfpathlineto{\pgfqpoint{5.465629in}{0.951286in}}%
\pgfpathlineto{\pgfqpoint{5.476748in}{0.949746in}}%
\pgfpathlineto{\pgfqpoint{5.544049in}{0.956404in}}%
\pgfpathlineto{\pgfqpoint{5.598070in}{0.959461in}}%
\pgfpathlineto{\pgfqpoint{5.630991in}{0.962230in}}%
\pgfpathlineto{\pgfqpoint{5.685505in}{0.967399in}}%
\pgfpathlineto{\pgfqpoint{5.686414in}{0.968427in}}%
\pgfpathlineto{\pgfqpoint{5.709707in}{0.970734in}}%
\pgfpathlineto{\pgfqpoint{5.786621in}{0.981538in}}%
\pgfpathlineto{\pgfqpoint{5.784883in}{0.968727in}}%
\pgfpathlineto{\pgfqpoint{5.777288in}{0.961208in}}%
\pgfpathlineto{\pgfqpoint{5.773253in}{0.950409in}}%
\pgfpathlineto{\pgfqpoint{5.763697in}{0.950142in}}%
\pgfpathlineto{\pgfqpoint{5.745602in}{0.935969in}}%
\pgfpathlineto{\pgfqpoint{5.737040in}{0.937098in}}%
\pgfpathlineto{\pgfqpoint{5.727639in}{0.921543in}}%
\pgfpathlineto{\pgfqpoint{5.715113in}{0.916361in}}%
\pgfpathlineto{\pgfqpoint{5.703119in}{0.904768in}}%
\pgfpathlineto{\pgfqpoint{5.689476in}{0.902614in}}%
\pgfpathlineto{\pgfqpoint{5.678562in}{0.892095in}}%
\pgfpathlineto{\pgfqpoint{5.678051in}{0.883484in}}%
\pgfpathlineto{\pgfqpoint{5.665861in}{0.878097in}}%
\pgfpathlineto{\pgfqpoint{5.665820in}{0.864070in}}%
\pgfpathclose%
\pgfusepath{fill}%
\end{pgfscope}%
\begin{pgfscope}%
\pgfpathrectangle{\pgfqpoint{3.525000in}{0.100000in}}{\pgfqpoint{2.862985in}{1.832778in}}%
\pgfusepath{clip}%
\pgfsetbuttcap%
\pgfsetmiterjoin%
\definecolor{currentfill}{rgb}{0.717186,0.879508,0.520185}%
\pgfsetfillcolor{currentfill}%
\pgfsetlinewidth{0.000000pt}%
\definecolor{currentstroke}{rgb}{0.000000,0.000000,0.000000}%
\pgfsetstrokecolor{currentstroke}%
\pgfsetstrokeopacity{0.000000}%
\pgfsetdash{}{0pt}%
\pgfpathmoveto{\pgfqpoint{4.498883in}{0.669271in}}%
\pgfpathlineto{\pgfqpoint{4.494190in}{0.676811in}}%
\pgfpathlineto{\pgfqpoint{4.496132in}{0.683307in}}%
\pgfpathlineto{\pgfqpoint{4.529118in}{0.679215in}}%
\pgfpathlineto{\pgfqpoint{4.590452in}{0.672249in}}%
\pgfpathlineto{\pgfqpoint{4.634801in}{0.667891in}}%
\pgfpathlineto{\pgfqpoint{4.686021in}{0.662722in}}%
\pgfpathlineto{\pgfqpoint{4.688807in}{0.695029in}}%
\pgfpathlineto{\pgfqpoint{4.697520in}{0.777201in}}%
\pgfpathlineto{\pgfqpoint{4.703172in}{0.834767in}}%
\pgfpathlineto{\pgfqpoint{4.708057in}{0.889756in}}%
\pgfpathlineto{\pgfqpoint{4.712612in}{0.944071in}}%
\pgfpathlineto{\pgfqpoint{4.714118in}{0.943987in}}%
\pgfpathlineto{\pgfqpoint{4.755936in}{0.940356in}}%
\pgfpathlineto{\pgfqpoint{4.810064in}{0.936263in}}%
\pgfpathlineto{\pgfqpoint{4.864916in}{0.933195in}}%
\pgfpathlineto{\pgfqpoint{4.861574in}{0.866207in}}%
\pgfpathlineto{\pgfqpoint{4.858507in}{0.812086in}}%
\pgfpathlineto{\pgfqpoint{4.861884in}{0.812959in}}%
\pgfpathlineto{\pgfqpoint{4.874005in}{0.799614in}}%
\pgfpathlineto{\pgfqpoint{4.879657in}{0.801308in}}%
\pgfpathlineto{\pgfqpoint{4.888240in}{0.798376in}}%
\pgfpathlineto{\pgfqpoint{4.895733in}{0.799856in}}%
\pgfpathlineto{\pgfqpoint{4.899917in}{0.787973in}}%
\pgfpathlineto{\pgfqpoint{4.909265in}{0.787349in}}%
\pgfpathlineto{\pgfqpoint{4.924189in}{0.781753in}}%
\pgfpathlineto{\pgfqpoint{4.928387in}{0.783330in}}%
\pgfpathlineto{\pgfqpoint{4.937392in}{0.776425in}}%
\pgfpathlineto{\pgfqpoint{4.942692in}{0.781753in}}%
\pgfpathlineto{\pgfqpoint{4.951877in}{0.779738in}}%
\pgfpathlineto{\pgfqpoint{4.956154in}{0.772659in}}%
\pgfpathlineto{\pgfqpoint{4.963177in}{0.770275in}}%
\pgfpathlineto{\pgfqpoint{4.968049in}{0.762933in}}%
\pgfpathlineto{\pgfqpoint{4.977778in}{0.770922in}}%
\pgfpathlineto{\pgfqpoint{4.988911in}{0.761010in}}%
\pgfpathlineto{\pgfqpoint{5.001640in}{0.765059in}}%
\pgfpathlineto{\pgfqpoint{5.008084in}{0.754622in}}%
\pgfpathlineto{\pgfqpoint{5.014386in}{0.766758in}}%
\pgfpathlineto{\pgfqpoint{5.025714in}{0.759877in}}%
\pgfpathlineto{\pgfqpoint{5.032791in}{0.761708in}}%
\pgfpathlineto{\pgfqpoint{5.042897in}{0.756843in}}%
\pgfpathlineto{\pgfqpoint{5.055999in}{0.755462in}}%
\pgfpathlineto{\pgfqpoint{5.068418in}{0.762948in}}%
\pgfpathlineto{\pgfqpoint{5.072911in}{0.759788in}}%
\pgfpathlineto{\pgfqpoint{5.085532in}{0.766205in}}%
\pgfpathlineto{\pgfqpoint{5.093575in}{0.761197in}}%
\pgfpathlineto{\pgfqpoint{5.099932in}{0.761495in}}%
\pgfpathlineto{\pgfqpoint{5.104966in}{0.767394in}}%
\pgfpathlineto{\pgfqpoint{5.112888in}{0.764459in}}%
\pgfpathlineto{\pgfqpoint{5.130724in}{0.750800in}}%
\pgfpathlineto{\pgfqpoint{5.144643in}{0.746495in}}%
\pgfpathlineto{\pgfqpoint{5.150225in}{0.741221in}}%
\pgfpathlineto{\pgfqpoint{5.157212in}{0.744093in}}%
\pgfpathlineto{\pgfqpoint{5.167799in}{0.741894in}}%
\pgfpathlineto{\pgfqpoint{5.168006in}{0.708326in}}%
\pgfpathlineto{\pgfqpoint{5.168890in}{0.643403in}}%
\pgfpathlineto{\pgfqpoint{5.176243in}{0.637164in}}%
\pgfpathlineto{\pgfqpoint{5.182137in}{0.625668in}}%
\pgfpathlineto{\pgfqpoint{5.180065in}{0.617977in}}%
\pgfpathlineto{\pgfqpoint{5.184693in}{0.614721in}}%
\pgfpathlineto{\pgfqpoint{5.188454in}{0.600538in}}%
\pgfpathlineto{\pgfqpoint{5.195658in}{0.592942in}}%
\pgfpathlineto{\pgfqpoint{5.197272in}{0.576816in}}%
\pgfpathlineto{\pgfqpoint{5.196019in}{0.569988in}}%
\pgfpathlineto{\pgfqpoint{5.186227in}{0.551919in}}%
\pgfpathlineto{\pgfqpoint{5.185101in}{0.538927in}}%
\pgfpathlineto{\pgfqpoint{5.188423in}{0.536215in}}%
\pgfpathlineto{\pgfqpoint{5.188712in}{0.524834in}}%
\pgfpathlineto{\pgfqpoint{5.185335in}{0.517680in}}%
\pgfpathlineto{\pgfqpoint{5.180085in}{0.514985in}}%
\pgfpathlineto{\pgfqpoint{5.174944in}{0.505637in}}%
\pgfpathlineto{\pgfqpoint{5.181496in}{0.496583in}}%
\pgfpathlineto{\pgfqpoint{5.168777in}{0.496405in}}%
\pgfpathlineto{\pgfqpoint{5.134766in}{0.480844in}}%
\pgfpathlineto{\pgfqpoint{5.141259in}{0.490157in}}%
\pgfpathlineto{\pgfqpoint{5.133396in}{0.495178in}}%
\pgfpathlineto{\pgfqpoint{5.131753in}{0.503714in}}%
\pgfpathlineto{\pgfqpoint{5.116406in}{0.488845in}}%
\pgfpathlineto{\pgfqpoint{5.123217in}{0.478682in}}%
\pgfpathlineto{\pgfqpoint{5.113501in}{0.465742in}}%
\pgfpathlineto{\pgfqpoint{5.108273in}{0.466013in}}%
\pgfpathlineto{\pgfqpoint{5.103354in}{0.451920in}}%
\pgfpathlineto{\pgfqpoint{5.087798in}{0.440835in}}%
\pgfpathlineto{\pgfqpoint{5.073265in}{0.436840in}}%
\pgfpathlineto{\pgfqpoint{5.047917in}{0.426447in}}%
\pgfpathlineto{\pgfqpoint{5.045291in}{0.432244in}}%
\pgfpathlineto{\pgfqpoint{5.029501in}{0.426908in}}%
\pgfpathlineto{\pgfqpoint{5.039213in}{0.417822in}}%
\pgfpathlineto{\pgfqpoint{5.023894in}{0.409929in}}%
\pgfpathlineto{\pgfqpoint{5.007371in}{0.398011in}}%
\pgfpathlineto{\pgfqpoint{5.003403in}{0.403551in}}%
\pgfpathlineto{\pgfqpoint{4.989881in}{0.395248in}}%
\pgfpathlineto{\pgfqpoint{5.003135in}{0.392095in}}%
\pgfpathlineto{\pgfqpoint{4.993142in}{0.379039in}}%
\pgfpathlineto{\pgfqpoint{4.988264in}{0.382924in}}%
\pgfpathlineto{\pgfqpoint{4.981713in}{0.376684in}}%
\pgfpathlineto{\pgfqpoint{4.989878in}{0.371233in}}%
\pgfpathlineto{\pgfqpoint{4.985050in}{0.363265in}}%
\pgfpathlineto{\pgfqpoint{4.979117in}{0.344404in}}%
\pgfpathlineto{\pgfqpoint{4.970559in}{0.326238in}}%
\pgfpathlineto{\pgfqpoint{4.974401in}{0.314331in}}%
\pgfpathlineto{\pgfqpoint{4.980932in}{0.286258in}}%
\pgfpathlineto{\pgfqpoint{4.981520in}{0.274896in}}%
\pgfpathlineto{\pgfqpoint{4.991610in}{0.259994in}}%
\pgfpathlineto{\pgfqpoint{4.983829in}{0.260865in}}%
\pgfpathlineto{\pgfqpoint{4.976272in}{0.253347in}}%
\pgfpathlineto{\pgfqpoint{4.967420in}{0.259256in}}%
\pgfpathlineto{\pgfqpoint{4.964243in}{0.265150in}}%
\pgfpathlineto{\pgfqpoint{4.951645in}{0.267892in}}%
\pgfpathlineto{\pgfqpoint{4.932402in}{0.268229in}}%
\pgfpathlineto{\pgfqpoint{4.918224in}{0.279391in}}%
\pgfpathlineto{\pgfqpoint{4.905350in}{0.281256in}}%
\pgfpathlineto{\pgfqpoint{4.897540in}{0.290127in}}%
\pgfpathlineto{\pgfqpoint{4.881183in}{0.293697in}}%
\pgfpathlineto{\pgfqpoint{4.872239in}{0.322221in}}%
\pgfpathlineto{\pgfqpoint{4.863094in}{0.333648in}}%
\pgfpathlineto{\pgfqpoint{4.864648in}{0.344506in}}%
\pgfpathlineto{\pgfqpoint{4.858978in}{0.352440in}}%
\pgfpathlineto{\pgfqpoint{4.862536in}{0.363286in}}%
\pgfpathlineto{\pgfqpoint{4.859601in}{0.371234in}}%
\pgfpathlineto{\pgfqpoint{4.850411in}{0.374834in}}%
\pgfpathlineto{\pgfqpoint{4.841818in}{0.383999in}}%
\pgfpathlineto{\pgfqpoint{4.838687in}{0.396269in}}%
\pgfpathlineto{\pgfqpoint{4.830556in}{0.407420in}}%
\pgfpathlineto{\pgfqpoint{4.819734in}{0.416096in}}%
\pgfpathlineto{\pgfqpoint{4.809986in}{0.440984in}}%
\pgfpathlineto{\pgfqpoint{4.802115in}{0.468203in}}%
\pgfpathlineto{\pgfqpoint{4.795657in}{0.478966in}}%
\pgfpathlineto{\pgfqpoint{4.784460in}{0.488026in}}%
\pgfpathlineto{\pgfqpoint{4.781662in}{0.494603in}}%
\pgfpathlineto{\pgfqpoint{4.771179in}{0.498686in}}%
\pgfpathlineto{\pgfqpoint{4.760831in}{0.516112in}}%
\pgfpathlineto{\pgfqpoint{4.751382in}{0.514774in}}%
\pgfpathlineto{\pgfqpoint{4.741787in}{0.519126in}}%
\pgfpathlineto{\pgfqpoint{4.728186in}{0.518285in}}%
\pgfpathlineto{\pgfqpoint{4.714369in}{0.525467in}}%
\pgfpathlineto{\pgfqpoint{4.710472in}{0.518637in}}%
\pgfpathlineto{\pgfqpoint{4.694325in}{0.518466in}}%
\pgfpathlineto{\pgfqpoint{4.686104in}{0.505517in}}%
\pgfpathlineto{\pgfqpoint{4.678954in}{0.489484in}}%
\pgfpathlineto{\pgfqpoint{4.663750in}{0.472274in}}%
\pgfpathlineto{\pgfqpoint{4.646524in}{0.479827in}}%
\pgfpathlineto{\pgfqpoint{4.644080in}{0.484818in}}%
\pgfpathlineto{\pgfqpoint{4.633611in}{0.488625in}}%
\pgfpathlineto{\pgfqpoint{4.630396in}{0.493833in}}%
\pgfpathlineto{\pgfqpoint{4.616499in}{0.499090in}}%
\pgfpathlineto{\pgfqpoint{4.608725in}{0.509836in}}%
\pgfpathlineto{\pgfqpoint{4.599646in}{0.515021in}}%
\pgfpathlineto{\pgfqpoint{4.591834in}{0.524068in}}%
\pgfpathlineto{\pgfqpoint{4.585749in}{0.539391in}}%
\pgfpathlineto{\pgfqpoint{4.586430in}{0.560329in}}%
\pgfpathlineto{\pgfqpoint{4.579291in}{0.570909in}}%
\pgfpathlineto{\pgfqpoint{4.578469in}{0.582369in}}%
\pgfpathlineto{\pgfqpoint{4.562630in}{0.599527in}}%
\pgfpathlineto{\pgfqpoint{4.553414in}{0.603189in}}%
\pgfpathlineto{\pgfqpoint{4.543634in}{0.619199in}}%
\pgfpathlineto{\pgfqpoint{4.535300in}{0.625560in}}%
\pgfpathlineto{\pgfqpoint{4.524711in}{0.641049in}}%
\pgfpathlineto{\pgfqpoint{4.513876in}{0.647753in}}%
\pgfpathlineto{\pgfqpoint{4.506780in}{0.664929in}}%
\pgfpathlineto{\pgfqpoint{4.498883in}{0.669271in}}%
\pgfpathclose%
\pgfusepath{fill}%
\end{pgfscope}%
\begin{pgfscope}%
\pgfpathrectangle{\pgfqpoint{3.525000in}{0.100000in}}{\pgfqpoint{2.862985in}{1.832778in}}%
\pgfusepath{clip}%
\pgfsetbuttcap%
\pgfsetmiterjoin%
\definecolor{currentfill}{rgb}{0.568166,0.782776,0.339792}%
\pgfsetfillcolor{currentfill}%
\pgfsetlinewidth{0.000000pt}%
\definecolor{currentstroke}{rgb}{0.000000,0.000000,0.000000}%
\pgfsetstrokecolor{currentstroke}%
\pgfsetstrokeopacity{0.000000}%
\pgfsetdash{}{0pt}%
\pgfpathmoveto{\pgfqpoint{4.415915in}{1.012202in}}%
\pgfpathlineto{\pgfqpoint{4.494080in}{1.000850in}}%
\pgfpathlineto{\pgfqpoint{4.524362in}{0.996184in}}%
\pgfpathlineto{\pgfqpoint{4.609722in}{0.986032in}}%
\pgfpathlineto{\pgfqpoint{4.667175in}{0.980039in}}%
\pgfpathlineto{\pgfqpoint{4.716895in}{0.975633in}}%
\pgfpathlineto{\pgfqpoint{4.714118in}{0.943987in}}%
\pgfpathlineto{\pgfqpoint{4.712612in}{0.944071in}}%
\pgfpathlineto{\pgfqpoint{4.708057in}{0.889756in}}%
\pgfpathlineto{\pgfqpoint{4.703172in}{0.834767in}}%
\pgfpathlineto{\pgfqpoint{4.697520in}{0.777201in}}%
\pgfpathlineto{\pgfqpoint{4.688807in}{0.695029in}}%
\pgfpathlineto{\pgfqpoint{4.686021in}{0.662722in}}%
\pgfpathlineto{\pgfqpoint{4.634801in}{0.667891in}}%
\pgfpathlineto{\pgfqpoint{4.590452in}{0.672249in}}%
\pgfpathlineto{\pgfqpoint{4.529118in}{0.679215in}}%
\pgfpathlineto{\pgfqpoint{4.496132in}{0.683307in}}%
\pgfpathlineto{\pgfqpoint{4.494190in}{0.676811in}}%
\pgfpathlineto{\pgfqpoint{4.498883in}{0.669271in}}%
\pgfpathlineto{\pgfqpoint{4.459240in}{0.674419in}}%
\pgfpathlineto{\pgfqpoint{4.410330in}{0.681439in}}%
\pgfpathlineto{\pgfqpoint{4.405878in}{0.653775in}}%
\pgfpathlineto{\pgfqpoint{4.361264in}{0.660490in}}%
\pgfpathlineto{\pgfqpoint{4.371747in}{0.729118in}}%
\pgfpathlineto{\pgfqpoint{4.384723in}{0.812676in}}%
\pgfpathlineto{\pgfqpoint{4.396262in}{0.885469in}}%
\pgfpathlineto{\pgfqpoint{4.406281in}{0.950172in}}%
\pgfpathlineto{\pgfqpoint{4.415915in}{1.012202in}}%
\pgfpathclose%
\pgfusepath{fill}%
\end{pgfscope}%
\begin{pgfscope}%
\pgfpathrectangle{\pgfqpoint{3.525000in}{0.100000in}}{\pgfqpoint{2.862985in}{1.832778in}}%
\pgfusepath{clip}%
\pgfsetbuttcap%
\pgfsetmiterjoin%
\definecolor{currentfill}{rgb}{0.367320,0.627451,0.171242}%
\pgfsetfillcolor{currentfill}%
\pgfsetlinewidth{0.000000pt}%
\definecolor{currentstroke}{rgb}{0.000000,0.000000,0.000000}%
\pgfsetstrokecolor{currentstroke}%
\pgfsetstrokeopacity{0.000000}%
\pgfsetdash{}{0pt}%
\pgfpathmoveto{\pgfqpoint{5.658749in}{0.610119in}}%
\pgfpathlineto{\pgfqpoint{5.584178in}{0.601837in}}%
\pgfpathlineto{\pgfqpoint{5.518446in}{0.596734in}}%
\pgfpathlineto{\pgfqpoint{5.517624in}{0.588678in}}%
\pgfpathlineto{\pgfqpoint{5.523661in}{0.581011in}}%
\pgfpathlineto{\pgfqpoint{5.531031in}{0.576507in}}%
\pgfpathlineto{\pgfqpoint{5.530919in}{0.564647in}}%
\pgfpathlineto{\pgfqpoint{5.528966in}{0.556723in}}%
\pgfpathlineto{\pgfqpoint{5.522462in}{0.551013in}}%
\pgfpathlineto{\pgfqpoint{5.513396in}{0.551618in}}%
\pgfpathlineto{\pgfqpoint{5.504825in}{0.558693in}}%
\pgfpathlineto{\pgfqpoint{5.503297in}{0.571283in}}%
\pgfpathlineto{\pgfqpoint{5.496913in}{0.578606in}}%
\pgfpathlineto{\pgfqpoint{5.492568in}{0.552398in}}%
\pgfpathlineto{\pgfqpoint{5.477803in}{0.554889in}}%
\pgfpathlineto{\pgfqpoint{5.472669in}{0.600717in}}%
\pgfpathlineto{\pgfqpoint{5.467108in}{0.649015in}}%
\pgfpathlineto{\pgfqpoint{5.469173in}{0.718697in}}%
\pgfpathlineto{\pgfqpoint{5.470455in}{0.766513in}}%
\pgfpathlineto{\pgfqpoint{5.473181in}{0.839470in}}%
\pgfpathlineto{\pgfqpoint{5.467346in}{0.845890in}}%
\pgfpathlineto{\pgfqpoint{5.477965in}{0.847429in}}%
\pgfpathlineto{\pgfqpoint{5.539830in}{0.851430in}}%
\pgfpathlineto{\pgfqpoint{5.600022in}{0.856843in}}%
\pgfpathlineto{\pgfqpoint{5.615858in}{0.801373in}}%
\pgfpathlineto{\pgfqpoint{5.626658in}{0.760657in}}%
\pgfpathlineto{\pgfqpoint{5.636339in}{0.726571in}}%
\pgfpathlineto{\pgfqpoint{5.641938in}{0.712845in}}%
\pgfpathlineto{\pgfqpoint{5.650753in}{0.700073in}}%
\pgfpathlineto{\pgfqpoint{5.649330in}{0.693573in}}%
\pgfpathlineto{\pgfqpoint{5.655627in}{0.690406in}}%
\pgfpathlineto{\pgfqpoint{5.648185in}{0.680553in}}%
\pgfpathlineto{\pgfqpoint{5.645672in}{0.663006in}}%
\pgfpathlineto{\pgfqpoint{5.652945in}{0.642483in}}%
\pgfpathlineto{\pgfqpoint{5.651935in}{0.621831in}}%
\pgfpathlineto{\pgfqpoint{5.658749in}{0.610119in}}%
\pgfpathclose%
\pgfusepath{fill}%
\end{pgfscope}%
\begin{pgfscope}%
\pgfpathrectangle{\pgfqpoint{3.525000in}{0.100000in}}{\pgfqpoint{2.862985in}{1.832778in}}%
\pgfusepath{clip}%
\pgfsetbuttcap%
\pgfsetmiterjoin%
\definecolor{currentfill}{rgb}{0.673356,0.851057,0.467128}%
\pgfsetfillcolor{currentfill}%
\pgfsetlinewidth{0.000000pt}%
\definecolor{currentstroke}{rgb}{0.000000,0.000000,0.000000}%
\pgfsetstrokecolor{currentstroke}%
\pgfsetstrokeopacity{0.000000}%
\pgfsetdash{}{0pt}%
\pgfpathmoveto{\pgfqpoint{5.477803in}{0.554889in}}%
\pgfpathlineto{\pgfqpoint{5.462652in}{0.550553in}}%
\pgfpathlineto{\pgfqpoint{5.451358in}{0.553337in}}%
\pgfpathlineto{\pgfqpoint{5.430384in}{0.546666in}}%
\pgfpathlineto{\pgfqpoint{5.414565in}{0.538074in}}%
\pgfpathlineto{\pgfqpoint{5.408075in}{0.553600in}}%
\pgfpathlineto{\pgfqpoint{5.401384in}{0.560107in}}%
\pgfpathlineto{\pgfqpoint{5.398157in}{0.567894in}}%
\pgfpathlineto{\pgfqpoint{5.402888in}{0.588972in}}%
\pgfpathlineto{\pgfqpoint{5.358060in}{0.586211in}}%
\pgfpathlineto{\pgfqpoint{5.299932in}{0.583691in}}%
\pgfpathlineto{\pgfqpoint{5.303390in}{0.588938in}}%
\pgfpathlineto{\pgfqpoint{5.299116in}{0.598847in}}%
\pgfpathlineto{\pgfqpoint{5.303375in}{0.600866in}}%
\pgfpathlineto{\pgfqpoint{5.302430in}{0.610386in}}%
\pgfpathlineto{\pgfqpoint{5.309854in}{0.619619in}}%
\pgfpathlineto{\pgfqpoint{5.313916in}{0.637542in}}%
\pgfpathlineto{\pgfqpoint{5.327503in}{0.649350in}}%
\pgfpathlineto{\pgfqpoint{5.322467in}{0.660818in}}%
\pgfpathlineto{\pgfqpoint{5.332102in}{0.667419in}}%
\pgfpathlineto{\pgfqpoint{5.323791in}{0.679370in}}%
\pgfpathlineto{\pgfqpoint{5.317776in}{0.706727in}}%
\pgfpathlineto{\pgfqpoint{5.320057in}{0.711695in}}%
\pgfpathlineto{\pgfqpoint{5.323659in}{0.721223in}}%
\pgfpathlineto{\pgfqpoint{5.320307in}{0.731223in}}%
\pgfpathlineto{\pgfqpoint{5.321946in}{0.735772in}}%
\pgfpathlineto{\pgfqpoint{5.316048in}{0.752960in}}%
\pgfpathlineto{\pgfqpoint{5.333459in}{0.783821in}}%
\pgfpathlineto{\pgfqpoint{5.334269in}{0.792035in}}%
\pgfpathlineto{\pgfqpoint{5.342656in}{0.798007in}}%
\pgfpathlineto{\pgfqpoint{5.351605in}{0.817725in}}%
\pgfpathlineto{\pgfqpoint{5.351180in}{0.825739in}}%
\pgfpathlineto{\pgfqpoint{5.362326in}{0.833932in}}%
\pgfpathlineto{\pgfqpoint{5.358848in}{0.838785in}}%
\pgfpathlineto{\pgfqpoint{5.425643in}{0.842683in}}%
\pgfpathlineto{\pgfqpoint{5.467346in}{0.845890in}}%
\pgfpathlineto{\pgfqpoint{5.473181in}{0.839470in}}%
\pgfpathlineto{\pgfqpoint{5.470455in}{0.766513in}}%
\pgfpathlineto{\pgfqpoint{5.469173in}{0.718697in}}%
\pgfpathlineto{\pgfqpoint{5.467108in}{0.649015in}}%
\pgfpathlineto{\pgfqpoint{5.472669in}{0.600717in}}%
\pgfpathlineto{\pgfqpoint{5.477803in}{0.554889in}}%
\pgfpathclose%
\pgfusepath{fill}%
\end{pgfscope}%
\begin{pgfscope}%
\pgfpathrectangle{\pgfqpoint{3.525000in}{0.100000in}}{\pgfqpoint{2.862985in}{1.832778in}}%
\pgfusepath{clip}%
\pgfsetbuttcap%
\pgfsetmiterjoin%
\definecolor{currentfill}{rgb}{0.498039,0.737255,0.254902}%
\pgfsetfillcolor{currentfill}%
\pgfsetlinewidth{0.000000pt}%
\definecolor{currentstroke}{rgb}{0.000000,0.000000,0.000000}%
\pgfsetstrokecolor{currentstroke}%
\pgfsetstrokeopacity{0.000000}%
\pgfsetdash{}{0pt}%
\pgfpathmoveto{\pgfqpoint{5.600022in}{0.856843in}}%
\pgfpathlineto{\pgfqpoint{5.665820in}{0.864070in}}%
\pgfpathlineto{\pgfqpoint{5.705410in}{0.868948in}}%
\pgfpathlineto{\pgfqpoint{5.727958in}{0.872540in}}%
\pgfpathlineto{\pgfqpoint{5.718618in}{0.857809in}}%
\pgfpathlineto{\pgfqpoint{5.718623in}{0.850863in}}%
\pgfpathlineto{\pgfqpoint{5.728238in}{0.847117in}}%
\pgfpathlineto{\pgfqpoint{5.734776in}{0.841070in}}%
\pgfpathlineto{\pgfqpoint{5.744657in}{0.840323in}}%
\pgfpathlineto{\pgfqpoint{5.753892in}{0.823299in}}%
\pgfpathlineto{\pgfqpoint{5.763910in}{0.811294in}}%
\pgfpathlineto{\pgfqpoint{5.781702in}{0.801205in}}%
\pgfpathlineto{\pgfqpoint{5.786758in}{0.793157in}}%
\pgfpathlineto{\pgfqpoint{5.802624in}{0.784454in}}%
\pgfpathlineto{\pgfqpoint{5.803817in}{0.778143in}}%
\pgfpathlineto{\pgfqpoint{5.816688in}{0.765593in}}%
\pgfpathlineto{\pgfqpoint{5.826546in}{0.762019in}}%
\pgfpathlineto{\pgfqpoint{5.833528in}{0.750169in}}%
\pgfpathlineto{\pgfqpoint{5.836618in}{0.736849in}}%
\pgfpathlineto{\pgfqpoint{5.846858in}{0.731700in}}%
\pgfpathlineto{\pgfqpoint{5.855102in}{0.717383in}}%
\pgfpathlineto{\pgfqpoint{5.857761in}{0.706853in}}%
\pgfpathlineto{\pgfqpoint{5.869841in}{0.702371in}}%
\pgfpathlineto{\pgfqpoint{5.859704in}{0.682961in}}%
\pgfpathlineto{\pgfqpoint{5.855889in}{0.671491in}}%
\pgfpathlineto{\pgfqpoint{5.858378in}{0.666114in}}%
\pgfpathlineto{\pgfqpoint{5.854032in}{0.657178in}}%
\pgfpathlineto{\pgfqpoint{5.852185in}{0.644824in}}%
\pgfpathlineto{\pgfqpoint{5.844456in}{0.642523in}}%
\pgfpathlineto{\pgfqpoint{5.848537in}{0.631212in}}%
\pgfpathlineto{\pgfqpoint{5.848277in}{0.616862in}}%
\pgfpathlineto{\pgfqpoint{5.834539in}{0.617400in}}%
\pgfpathlineto{\pgfqpoint{5.824969in}{0.620018in}}%
\pgfpathlineto{\pgfqpoint{5.820828in}{0.615151in}}%
\pgfpathlineto{\pgfqpoint{5.823918in}{0.603653in}}%
\pgfpathlineto{\pgfqpoint{5.823246in}{0.590285in}}%
\pgfpathlineto{\pgfqpoint{5.817209in}{0.589262in}}%
\pgfpathlineto{\pgfqpoint{5.812325in}{0.601747in}}%
\pgfpathlineto{\pgfqpoint{5.762593in}{0.598433in}}%
\pgfpathlineto{\pgfqpoint{5.699885in}{0.595011in}}%
\pgfpathlineto{\pgfqpoint{5.668260in}{0.592785in}}%
\pgfpathlineto{\pgfqpoint{5.658749in}{0.610119in}}%
\pgfpathlineto{\pgfqpoint{5.651935in}{0.621831in}}%
\pgfpathlineto{\pgfqpoint{5.652945in}{0.642483in}}%
\pgfpathlineto{\pgfqpoint{5.645672in}{0.663006in}}%
\pgfpathlineto{\pgfqpoint{5.648185in}{0.680553in}}%
\pgfpathlineto{\pgfqpoint{5.655627in}{0.690406in}}%
\pgfpathlineto{\pgfqpoint{5.649330in}{0.693573in}}%
\pgfpathlineto{\pgfqpoint{5.650753in}{0.700073in}}%
\pgfpathlineto{\pgfqpoint{5.641938in}{0.712845in}}%
\pgfpathlineto{\pgfqpoint{5.636339in}{0.726571in}}%
\pgfpathlineto{\pgfqpoint{5.626658in}{0.760657in}}%
\pgfpathlineto{\pgfqpoint{5.615858in}{0.801373in}}%
\pgfpathlineto{\pgfqpoint{5.600022in}{0.856843in}}%
\pgfpathclose%
\pgfusepath{fill}%
\end{pgfscope}%
\begin{pgfscope}%
\pgfpathrectangle{\pgfqpoint{3.525000in}{0.100000in}}{\pgfqpoint{2.862985in}{1.832778in}}%
\pgfusepath{clip}%
\pgfsetbuttcap%
\pgfsetmiterjoin%
\definecolor{currentfill}{rgb}{0.498039,0.737255,0.254902}%
\pgfsetfillcolor{currentfill}%
\pgfsetlinewidth{0.000000pt}%
\definecolor{currentstroke}{rgb}{0.000000,0.000000,0.000000}%
\pgfsetstrokecolor{currentstroke}%
\pgfsetstrokeopacity{0.000000}%
\pgfsetdash{}{0pt}%
\pgfpathmoveto{\pgfqpoint{5.727958in}{0.872540in}}%
\pgfpathlineto{\pgfqpoint{5.754049in}{0.885519in}}%
\pgfpathlineto{\pgfqpoint{5.768499in}{0.890420in}}%
\pgfpathlineto{\pgfqpoint{5.831366in}{0.896952in}}%
\pgfpathlineto{\pgfqpoint{5.838024in}{0.894789in}}%
\pgfpathlineto{\pgfqpoint{5.846819in}{0.885949in}}%
\pgfpathlineto{\pgfqpoint{5.847304in}{0.878098in}}%
\pgfpathlineto{\pgfqpoint{5.905165in}{0.886600in}}%
\pgfpathlineto{\pgfqpoint{5.970869in}{0.839501in}}%
\pgfpathlineto{\pgfqpoint{5.958547in}{0.826658in}}%
\pgfpathlineto{\pgfqpoint{5.941640in}{0.796833in}}%
\pgfpathlineto{\pgfqpoint{5.946453in}{0.790423in}}%
\pgfpathlineto{\pgfqpoint{5.937434in}{0.777981in}}%
\pgfpathlineto{\pgfqpoint{5.928462in}{0.776573in}}%
\pgfpathlineto{\pgfqpoint{5.928429in}{0.769097in}}%
\pgfpathlineto{\pgfqpoint{5.921939in}{0.761274in}}%
\pgfpathlineto{\pgfqpoint{5.913862in}{0.759687in}}%
\pgfpathlineto{\pgfqpoint{5.915625in}{0.752736in}}%
\pgfpathlineto{\pgfqpoint{5.911118in}{0.747400in}}%
\pgfpathlineto{\pgfqpoint{5.900337in}{0.742790in}}%
\pgfpathlineto{\pgfqpoint{5.887196in}{0.732304in}}%
\pgfpathlineto{\pgfqpoint{5.889677in}{0.725520in}}%
\pgfpathlineto{\pgfqpoint{5.881407in}{0.721260in}}%
\pgfpathlineto{\pgfqpoint{5.874432in}{0.725739in}}%
\pgfpathlineto{\pgfqpoint{5.869331in}{0.706267in}}%
\pgfpathlineto{\pgfqpoint{5.857761in}{0.706853in}}%
\pgfpathlineto{\pgfqpoint{5.855102in}{0.717383in}}%
\pgfpathlineto{\pgfqpoint{5.846858in}{0.731700in}}%
\pgfpathlineto{\pgfqpoint{5.836618in}{0.736849in}}%
\pgfpathlineto{\pgfqpoint{5.833528in}{0.750169in}}%
\pgfpathlineto{\pgfqpoint{5.826546in}{0.762019in}}%
\pgfpathlineto{\pgfqpoint{5.816688in}{0.765593in}}%
\pgfpathlineto{\pgfqpoint{5.803817in}{0.778143in}}%
\pgfpathlineto{\pgfqpoint{5.802624in}{0.784454in}}%
\pgfpathlineto{\pgfqpoint{5.786758in}{0.793157in}}%
\pgfpathlineto{\pgfqpoint{5.781702in}{0.801205in}}%
\pgfpathlineto{\pgfqpoint{5.763910in}{0.811294in}}%
\pgfpathlineto{\pgfqpoint{5.753892in}{0.823299in}}%
\pgfpathlineto{\pgfqpoint{5.744657in}{0.840323in}}%
\pgfpathlineto{\pgfqpoint{5.734776in}{0.841070in}}%
\pgfpathlineto{\pgfqpoint{5.728238in}{0.847117in}}%
\pgfpathlineto{\pgfqpoint{5.718623in}{0.850863in}}%
\pgfpathlineto{\pgfqpoint{5.718618in}{0.857809in}}%
\pgfpathlineto{\pgfqpoint{5.727958in}{0.872540in}}%
\pgfpathclose%
\pgfusepath{fill}%
\end{pgfscope}%
\begin{pgfscope}%
\pgfpathrectangle{\pgfqpoint{3.525000in}{0.100000in}}{\pgfqpoint{2.862985in}{1.832778in}}%
\pgfusepath{clip}%
\pgfsetbuttcap%
\pgfsetmiterjoin%
\definecolor{currentfill}{rgb}{0.524337,0.754325,0.286736}%
\pgfsetfillcolor{currentfill}%
\pgfsetlinewidth{0.000000pt}%
\definecolor{currentstroke}{rgb}{0.000000,0.000000,0.000000}%
\pgfsetstrokecolor{currentstroke}%
\pgfsetstrokeopacity{0.000000}%
\pgfsetdash{}{0pt}%
\pgfpathmoveto{\pgfqpoint{5.136515in}{0.926094in}}%
\pgfpathlineto{\pgfqpoint{5.174856in}{0.926440in}}%
\pgfpathlineto{\pgfqpoint{5.225580in}{0.927364in}}%
\pgfpathlineto{\pgfqpoint{5.312257in}{0.930116in}}%
\pgfpathlineto{\pgfqpoint{5.361833in}{0.932737in}}%
\pgfpathlineto{\pgfqpoint{5.367150in}{0.926155in}}%
\pgfpathlineto{\pgfqpoint{5.366818in}{0.919204in}}%
\pgfpathlineto{\pgfqpoint{5.354808in}{0.907211in}}%
\pgfpathlineto{\pgfqpoint{5.351904in}{0.900644in}}%
\pgfpathlineto{\pgfqpoint{5.385213in}{0.903107in}}%
\pgfpathlineto{\pgfqpoint{5.385195in}{0.890974in}}%
\pgfpathlineto{\pgfqpoint{5.374509in}{0.885789in}}%
\pgfpathlineto{\pgfqpoint{5.374660in}{0.877545in}}%
\pgfpathlineto{\pgfqpoint{5.370647in}{0.871619in}}%
\pgfpathlineto{\pgfqpoint{5.368020in}{0.859033in}}%
\pgfpathlineto{\pgfqpoint{5.370732in}{0.848649in}}%
\pgfpathlineto{\pgfqpoint{5.358848in}{0.838785in}}%
\pgfpathlineto{\pgfqpoint{5.362326in}{0.833932in}}%
\pgfpathlineto{\pgfqpoint{5.351180in}{0.825739in}}%
\pgfpathlineto{\pgfqpoint{5.351605in}{0.817725in}}%
\pgfpathlineto{\pgfqpoint{5.342656in}{0.798007in}}%
\pgfpathlineto{\pgfqpoint{5.334269in}{0.792035in}}%
\pgfpathlineto{\pgfqpoint{5.333459in}{0.783821in}}%
\pgfpathlineto{\pgfqpoint{5.316048in}{0.752960in}}%
\pgfpathlineto{\pgfqpoint{5.321946in}{0.735772in}}%
\pgfpathlineto{\pgfqpoint{5.320307in}{0.731223in}}%
\pgfpathlineto{\pgfqpoint{5.323659in}{0.721223in}}%
\pgfpathlineto{\pgfqpoint{5.320057in}{0.711695in}}%
\pgfpathlineto{\pgfqpoint{5.272446in}{0.709729in}}%
\pgfpathlineto{\pgfqpoint{5.210638in}{0.708711in}}%
\pgfpathlineto{\pgfqpoint{5.168006in}{0.708326in}}%
\pgfpathlineto{\pgfqpoint{5.167799in}{0.741894in}}%
\pgfpathlineto{\pgfqpoint{5.157212in}{0.744093in}}%
\pgfpathlineto{\pgfqpoint{5.150225in}{0.741221in}}%
\pgfpathlineto{\pgfqpoint{5.144643in}{0.746495in}}%
\pgfpathlineto{\pgfqpoint{5.145219in}{0.782052in}}%
\pgfpathlineto{\pgfqpoint{5.146468in}{0.857760in}}%
\pgfpathlineto{\pgfqpoint{5.140413in}{0.902092in}}%
\pgfpathlineto{\pgfqpoint{5.136515in}{0.926094in}}%
\pgfpathclose%
\pgfusepath{fill}%
\end{pgfscope}%
\begin{pgfscope}%
\pgfpathrectangle{\pgfqpoint{3.525000in}{0.100000in}}{\pgfqpoint{2.862985in}{1.832778in}}%
\pgfusepath{clip}%
\pgfsetbuttcap%
\pgfsetmiterjoin%
\definecolor{currentfill}{rgb}{0.594464,0.799846,0.371626}%
\pgfsetfillcolor{currentfill}%
\pgfsetlinewidth{0.000000pt}%
\definecolor{currentstroke}{rgb}{0.000000,0.000000,0.000000}%
\pgfsetstrokecolor{currentstroke}%
\pgfsetstrokeopacity{0.000000}%
\pgfsetdash{}{0pt}%
\pgfpathmoveto{\pgfqpoint{5.168006in}{0.708326in}}%
\pgfpathlineto{\pgfqpoint{5.210638in}{0.708711in}}%
\pgfpathlineto{\pgfqpoint{5.272446in}{0.709729in}}%
\pgfpathlineto{\pgfqpoint{5.320057in}{0.711695in}}%
\pgfpathlineto{\pgfqpoint{5.317776in}{0.706727in}}%
\pgfpathlineto{\pgfqpoint{5.323791in}{0.679370in}}%
\pgfpathlineto{\pgfqpoint{5.332102in}{0.667419in}}%
\pgfpathlineto{\pgfqpoint{5.322467in}{0.660818in}}%
\pgfpathlineto{\pgfqpoint{5.327503in}{0.649350in}}%
\pgfpathlineto{\pgfqpoint{5.313916in}{0.637542in}}%
\pgfpathlineto{\pgfqpoint{5.309854in}{0.619619in}}%
\pgfpathlineto{\pgfqpoint{5.302430in}{0.610386in}}%
\pgfpathlineto{\pgfqpoint{5.303375in}{0.600866in}}%
\pgfpathlineto{\pgfqpoint{5.299116in}{0.598847in}}%
\pgfpathlineto{\pgfqpoint{5.303390in}{0.588938in}}%
\pgfpathlineto{\pgfqpoint{5.299932in}{0.583691in}}%
\pgfpathlineto{\pgfqpoint{5.358060in}{0.586211in}}%
\pgfpathlineto{\pgfqpoint{5.402888in}{0.588972in}}%
\pgfpathlineto{\pgfqpoint{5.398157in}{0.567894in}}%
\pgfpathlineto{\pgfqpoint{5.401384in}{0.560107in}}%
\pgfpathlineto{\pgfqpoint{5.408075in}{0.553600in}}%
\pgfpathlineto{\pgfqpoint{5.414565in}{0.538074in}}%
\pgfpathlineto{\pgfqpoint{5.394065in}{0.541651in}}%
\pgfpathlineto{\pgfqpoint{5.386502in}{0.547531in}}%
\pgfpathlineto{\pgfqpoint{5.377492in}{0.547811in}}%
\pgfpathlineto{\pgfqpoint{5.368019in}{0.534934in}}%
\pgfpathlineto{\pgfqpoint{5.369907in}{0.529071in}}%
\pgfpathlineto{\pgfqpoint{5.400169in}{0.525523in}}%
\pgfpathlineto{\pgfqpoint{5.408091in}{0.518783in}}%
\pgfpathlineto{\pgfqpoint{5.420984in}{0.515315in}}%
\pgfpathlineto{\pgfqpoint{5.415353in}{0.507336in}}%
\pgfpathlineto{\pgfqpoint{5.405867in}{0.500376in}}%
\pgfpathlineto{\pgfqpoint{5.430906in}{0.484726in}}%
\pgfpathlineto{\pgfqpoint{5.442569in}{0.482281in}}%
\pgfpathlineto{\pgfqpoint{5.447419in}{0.469488in}}%
\pgfpathlineto{\pgfqpoint{5.436438in}{0.467135in}}%
\pgfpathlineto{\pgfqpoint{5.415799in}{0.483224in}}%
\pgfpathlineto{\pgfqpoint{5.407729in}{0.485461in}}%
\pgfpathlineto{\pgfqpoint{5.403824in}{0.491817in}}%
\pgfpathlineto{\pgfqpoint{5.391854in}{0.489208in}}%
\pgfpathlineto{\pgfqpoint{5.388107in}{0.481018in}}%
\pgfpathlineto{\pgfqpoint{5.390480in}{0.471892in}}%
\pgfpathlineto{\pgfqpoint{5.382422in}{0.466506in}}%
\pgfpathlineto{\pgfqpoint{5.375068in}{0.479774in}}%
\pgfpathlineto{\pgfqpoint{5.362187in}{0.475818in}}%
\pgfpathlineto{\pgfqpoint{5.357360in}{0.467901in}}%
\pgfpathlineto{\pgfqpoint{5.348213in}{0.470153in}}%
\pgfpathlineto{\pgfqpoint{5.342364in}{0.480150in}}%
\pgfpathlineto{\pgfqpoint{5.326847in}{0.483451in}}%
\pgfpathlineto{\pgfqpoint{5.323917in}{0.488653in}}%
\pgfpathlineto{\pgfqpoint{5.307751in}{0.497715in}}%
\pgfpathlineto{\pgfqpoint{5.303738in}{0.505638in}}%
\pgfpathlineto{\pgfqpoint{5.290197in}{0.502400in}}%
\pgfpathlineto{\pgfqpoint{5.291917in}{0.509654in}}%
\pgfpathlineto{\pgfqpoint{5.275086in}{0.502213in}}%
\pgfpathlineto{\pgfqpoint{5.279607in}{0.494492in}}%
\pgfpathlineto{\pgfqpoint{5.266610in}{0.489928in}}%
\pgfpathlineto{\pgfqpoint{5.249395in}{0.492437in}}%
\pgfpathlineto{\pgfqpoint{5.214551in}{0.504372in}}%
\pgfpathlineto{\pgfqpoint{5.187666in}{0.502003in}}%
\pgfpathlineto{\pgfqpoint{5.185335in}{0.517680in}}%
\pgfpathlineto{\pgfqpoint{5.188712in}{0.524834in}}%
\pgfpathlineto{\pgfqpoint{5.188423in}{0.536215in}}%
\pgfpathlineto{\pgfqpoint{5.185101in}{0.538927in}}%
\pgfpathlineto{\pgfqpoint{5.186227in}{0.551919in}}%
\pgfpathlineto{\pgfqpoint{5.196019in}{0.569988in}}%
\pgfpathlineto{\pgfqpoint{5.197272in}{0.576816in}}%
\pgfpathlineto{\pgfqpoint{5.195658in}{0.592942in}}%
\pgfpathlineto{\pgfqpoint{5.188454in}{0.600538in}}%
\pgfpathlineto{\pgfqpoint{5.184693in}{0.614721in}}%
\pgfpathlineto{\pgfqpoint{5.180065in}{0.617977in}}%
\pgfpathlineto{\pgfqpoint{5.182137in}{0.625668in}}%
\pgfpathlineto{\pgfqpoint{5.176243in}{0.637164in}}%
\pgfpathlineto{\pgfqpoint{5.168890in}{0.643403in}}%
\pgfpathlineto{\pgfqpoint{5.168006in}{0.708326in}}%
\pgfpathclose%
\pgfusepath{fill}%
\end{pgfscope}%
\begin{pgfscope}%
\pgfpathrectangle{\pgfqpoint{3.525000in}{0.100000in}}{\pgfqpoint{2.862985in}{1.832778in}}%
\pgfusepath{clip}%
\pgfsetbuttcap%
\pgfsetmiterjoin%
\definecolor{currentfill}{rgb}{0.594464,0.799846,0.371626}%
\pgfsetfillcolor{currentfill}%
\pgfsetlinewidth{0.000000pt}%
\definecolor{currentstroke}{rgb}{0.000000,0.000000,0.000000}%
\pgfsetstrokecolor{currentstroke}%
\pgfsetstrokeopacity{0.000000}%
\pgfsetdash{}{0pt}%
\pgfpathmoveto{\pgfqpoint{5.282025in}{0.493901in}}%
\pgfpathlineto{\pgfqpoint{5.288193in}{0.497570in}}%
\pgfpathlineto{\pgfqpoint{5.295679in}{0.493233in}}%
\pgfpathlineto{\pgfqpoint{5.291504in}{0.487265in}}%
\pgfpathlineto{\pgfqpoint{5.282025in}{0.493901in}}%
\pgfpathclose%
\pgfusepath{fill}%
\end{pgfscope}%
\begin{pgfscope}%
\pgfpathrectangle{\pgfqpoint{3.525000in}{0.100000in}}{\pgfqpoint{2.862985in}{1.832778in}}%
\pgfusepath{clip}%
\pgfsetbuttcap%
\pgfsetmiterjoin%
\definecolor{currentfill}{rgb}{0.451903,0.698501,0.225375}%
\pgfsetfillcolor{currentfill}%
\pgfsetlinewidth{0.000000pt}%
\definecolor{currentstroke}{rgb}{0.000000,0.000000,0.000000}%
\pgfsetstrokecolor{currentstroke}%
\pgfsetstrokeopacity{0.000000}%
\pgfsetdash{}{0pt}%
\pgfpathmoveto{\pgfqpoint{5.530919in}{0.564647in}}%
\pgfpathlineto{\pgfqpoint{5.531031in}{0.576507in}}%
\pgfpathlineto{\pgfqpoint{5.523661in}{0.581011in}}%
\pgfpathlineto{\pgfqpoint{5.517624in}{0.588678in}}%
\pgfpathlineto{\pgfqpoint{5.518446in}{0.596734in}}%
\pgfpathlineto{\pgfqpoint{5.584178in}{0.601837in}}%
\pgfpathlineto{\pgfqpoint{5.658749in}{0.610119in}}%
\pgfpathlineto{\pgfqpoint{5.668260in}{0.592785in}}%
\pgfpathlineto{\pgfqpoint{5.699885in}{0.595011in}}%
\pgfpathlineto{\pgfqpoint{5.762593in}{0.598433in}}%
\pgfpathlineto{\pgfqpoint{5.812325in}{0.601747in}}%
\pgfpathlineto{\pgfqpoint{5.817209in}{0.589262in}}%
\pgfpathlineto{\pgfqpoint{5.823246in}{0.590285in}}%
\pgfpathlineto{\pgfqpoint{5.823918in}{0.603653in}}%
\pgfpathlineto{\pgfqpoint{5.820828in}{0.615151in}}%
\pgfpathlineto{\pgfqpoint{5.824969in}{0.620018in}}%
\pgfpathlineto{\pgfqpoint{5.834539in}{0.617400in}}%
\pgfpathlineto{\pgfqpoint{5.848277in}{0.616862in}}%
\pgfpathlineto{\pgfqpoint{5.850389in}{0.606589in}}%
\pgfpathlineto{\pgfqpoint{5.854615in}{0.600719in}}%
\pgfpathlineto{\pgfqpoint{5.857915in}{0.587877in}}%
\pgfpathlineto{\pgfqpoint{5.868139in}{0.568003in}}%
\pgfpathlineto{\pgfqpoint{5.868190in}{0.562628in}}%
\pgfpathlineto{\pgfqpoint{5.883302in}{0.539291in}}%
\pgfpathlineto{\pgfqpoint{5.884764in}{0.534471in}}%
\pgfpathlineto{\pgfqpoint{5.907434in}{0.501512in}}%
\pgfpathlineto{\pgfqpoint{5.903842in}{0.500979in}}%
\pgfpathlineto{\pgfqpoint{5.913399in}{0.477537in}}%
\pgfpathlineto{\pgfqpoint{5.939653in}{0.436782in}}%
\pgfpathlineto{\pgfqpoint{5.954896in}{0.406818in}}%
\pgfpathlineto{\pgfqpoint{5.959991in}{0.401571in}}%
\pgfpathlineto{\pgfqpoint{5.968722in}{0.382763in}}%
\pgfpathlineto{\pgfqpoint{5.971750in}{0.352646in}}%
\pgfpathlineto{\pgfqpoint{5.972955in}{0.330127in}}%
\pgfpathlineto{\pgfqpoint{5.971500in}{0.315709in}}%
\pgfpathlineto{\pgfqpoint{5.966819in}{0.305401in}}%
\pgfpathlineto{\pgfqpoint{5.968994in}{0.291885in}}%
\pgfpathlineto{\pgfqpoint{5.963954in}{0.281186in}}%
\pgfpathlineto{\pgfqpoint{5.948991in}{0.272413in}}%
\pgfpathlineto{\pgfqpoint{5.939259in}{0.273053in}}%
\pgfpathlineto{\pgfqpoint{5.932940in}{0.268467in}}%
\pgfpathlineto{\pgfqpoint{5.931129in}{0.280712in}}%
\pgfpathlineto{\pgfqpoint{5.920668in}{0.283944in}}%
\pgfpathlineto{\pgfqpoint{5.911290in}{0.301030in}}%
\pgfpathlineto{\pgfqpoint{5.910231in}{0.308810in}}%
\pgfpathlineto{\pgfqpoint{5.893448in}{0.313535in}}%
\pgfpathlineto{\pgfqpoint{5.882633in}{0.312512in}}%
\pgfpathlineto{\pgfqpoint{5.876514in}{0.323801in}}%
\pgfpathlineto{\pgfqpoint{5.869508in}{0.344093in}}%
\pgfpathlineto{\pgfqpoint{5.859786in}{0.348206in}}%
\pgfpathlineto{\pgfqpoint{5.854515in}{0.359834in}}%
\pgfpathlineto{\pgfqpoint{5.841625in}{0.366687in}}%
\pgfpathlineto{\pgfqpoint{5.834174in}{0.375259in}}%
\pgfpathlineto{\pgfqpoint{5.821320in}{0.398564in}}%
\pgfpathlineto{\pgfqpoint{5.819809in}{0.414966in}}%
\pgfpathlineto{\pgfqpoint{5.826813in}{0.425461in}}%
\pgfpathlineto{\pgfqpoint{5.826113in}{0.432754in}}%
\pgfpathlineto{\pgfqpoint{5.818019in}{0.433533in}}%
\pgfpathlineto{\pgfqpoint{5.811256in}{0.438606in}}%
\pgfpathlineto{\pgfqpoint{5.807553in}{0.432412in}}%
\pgfpathlineto{\pgfqpoint{5.814043in}{0.427365in}}%
\pgfpathlineto{\pgfqpoint{5.809319in}{0.418289in}}%
\pgfpathlineto{\pgfqpoint{5.801677in}{0.425822in}}%
\pgfpathlineto{\pgfqpoint{5.802564in}{0.446710in}}%
\pgfpathlineto{\pgfqpoint{5.806250in}{0.463657in}}%
\pgfpathlineto{\pgfqpoint{5.804355in}{0.492736in}}%
\pgfpathlineto{\pgfqpoint{5.796738in}{0.499662in}}%
\pgfpathlineto{\pgfqpoint{5.792907in}{0.508549in}}%
\pgfpathlineto{\pgfqpoint{5.779785in}{0.508352in}}%
\pgfpathlineto{\pgfqpoint{5.766814in}{0.522976in}}%
\pgfpathlineto{\pgfqpoint{5.758117in}{0.527359in}}%
\pgfpathlineto{\pgfqpoint{5.755538in}{0.536620in}}%
\pgfpathlineto{\pgfqpoint{5.747031in}{0.539872in}}%
\pgfpathlineto{\pgfqpoint{5.739981in}{0.550098in}}%
\pgfpathlineto{\pgfqpoint{5.721358in}{0.558444in}}%
\pgfpathlineto{\pgfqpoint{5.706879in}{0.558658in}}%
\pgfpathlineto{\pgfqpoint{5.700576in}{0.555452in}}%
\pgfpathlineto{\pgfqpoint{5.702143in}{0.545445in}}%
\pgfpathlineto{\pgfqpoint{5.695570in}{0.545919in}}%
\pgfpathlineto{\pgfqpoint{5.675353in}{0.531927in}}%
\pgfpathlineto{\pgfqpoint{5.651068in}{0.526381in}}%
\pgfpathlineto{\pgfqpoint{5.650669in}{0.533243in}}%
\pgfpathlineto{\pgfqpoint{5.645287in}{0.539949in}}%
\pgfpathlineto{\pgfqpoint{5.630848in}{0.549208in}}%
\pgfpathlineto{\pgfqpoint{5.610121in}{0.558700in}}%
\pgfpathlineto{\pgfqpoint{5.587608in}{0.563725in}}%
\pgfpathlineto{\pgfqpoint{5.583382in}{0.570558in}}%
\pgfpathlineto{\pgfqpoint{5.575265in}{0.564859in}}%
\pgfpathlineto{\pgfqpoint{5.565491in}{0.563602in}}%
\pgfpathlineto{\pgfqpoint{5.531467in}{0.554637in}}%
\pgfpathlineto{\pgfqpoint{5.530919in}{0.564647in}}%
\pgfpathclose%
\pgfusepath{fill}%
\end{pgfscope}%
\begin{pgfscope}%
\pgfpathrectangle{\pgfqpoint{3.525000in}{0.100000in}}{\pgfqpoint{2.862985in}{1.832778in}}%
\pgfusepath{clip}%
\pgfsetbuttcap%
\pgfsetmiterjoin%
\definecolor{currentfill}{rgb}{0.451903,0.698501,0.225375}%
\pgfsetfillcolor{currentfill}%
\pgfsetlinewidth{0.000000pt}%
\definecolor{currentstroke}{rgb}{0.000000,0.000000,0.000000}%
\pgfsetstrokecolor{currentstroke}%
\pgfsetstrokeopacity{0.000000}%
\pgfsetdash{}{0pt}%
\pgfpathmoveto{\pgfqpoint{5.910737in}{0.501980in}}%
\pgfpathlineto{\pgfqpoint{5.919540in}{0.491740in}}%
\pgfpathlineto{\pgfqpoint{5.909372in}{0.491096in}}%
\pgfpathlineto{\pgfqpoint{5.910737in}{0.501980in}}%
\pgfpathclose%
\pgfusepath{fill}%
\end{pgfscope}%
\begin{pgfscope}%
\pgfpathrectangle{\pgfqpoint{3.525000in}{0.100000in}}{\pgfqpoint{2.862985in}{1.832778in}}%
\pgfusepath{clip}%
\pgfsetbuttcap%
\pgfsetmiterjoin%
\definecolor{currentfill}{rgb}{0.673356,0.851057,0.467128}%
\pgfsetfillcolor{currentfill}%
\pgfsetlinewidth{0.000000pt}%
\definecolor{currentstroke}{rgb}{0.000000,0.000000,0.000000}%
\pgfsetstrokecolor{currentstroke}%
\pgfsetstrokeopacity{0.000000}%
\pgfsetdash{}{0pt}%
\pgfpathmoveto{\pgfqpoint{5.392522in}{1.671190in}}%
\pgfpathlineto{\pgfqpoint{5.387746in}{1.661894in}}%
\pgfpathlineto{\pgfqpoint{5.376351in}{1.656468in}}%
\pgfpathlineto{\pgfqpoint{5.371416in}{1.649166in}}%
\pgfpathlineto{\pgfqpoint{5.365635in}{1.654425in}}%
\pgfpathlineto{\pgfqpoint{5.392522in}{1.671190in}}%
\pgfpathclose%
\pgfusepath{fill}%
\end{pgfscope}%
\begin{pgfscope}%
\pgfpathrectangle{\pgfqpoint{3.525000in}{0.100000in}}{\pgfqpoint{2.862985in}{1.832778in}}%
\pgfusepath{clip}%
\pgfsetbuttcap%
\pgfsetmiterjoin%
\definecolor{currentfill}{rgb}{0.673356,0.851057,0.467128}%
\pgfsetfillcolor{currentfill}%
\pgfsetlinewidth{0.000000pt}%
\definecolor{currentstroke}{rgb}{0.000000,0.000000,0.000000}%
\pgfsetstrokecolor{currentstroke}%
\pgfsetstrokeopacity{0.000000}%
\pgfsetdash{}{0pt}%
\pgfpathmoveto{\pgfqpoint{5.396396in}{1.615244in}}%
\pgfpathlineto{\pgfqpoint{5.408054in}{1.626120in}}%
\pgfpathlineto{\pgfqpoint{5.426037in}{1.628969in}}%
\pgfpathlineto{\pgfqpoint{5.421084in}{1.621405in}}%
\pgfpathlineto{\pgfqpoint{5.408756in}{1.610466in}}%
\pgfpathlineto{\pgfqpoint{5.401549in}{1.596426in}}%
\pgfpathlineto{\pgfqpoint{5.398496in}{1.604041in}}%
\pgfpathlineto{\pgfqpoint{5.393076in}{1.605134in}}%
\pgfpathlineto{\pgfqpoint{5.392532in}{1.612009in}}%
\pgfpathlineto{\pgfqpoint{5.396396in}{1.615244in}}%
\pgfpathclose%
\pgfusepath{fill}%
\end{pgfscope}%
\begin{pgfscope}%
\pgfpathrectangle{\pgfqpoint{3.525000in}{0.100000in}}{\pgfqpoint{2.862985in}{1.832778in}}%
\pgfusepath{clip}%
\pgfsetbuttcap%
\pgfsetmiterjoin%
\definecolor{currentfill}{rgb}{0.673356,0.851057,0.467128}%
\pgfsetfillcolor{currentfill}%
\pgfsetlinewidth{0.000000pt}%
\definecolor{currentstroke}{rgb}{0.000000,0.000000,0.000000}%
\pgfsetstrokecolor{currentstroke}%
\pgfsetstrokeopacity{0.000000}%
\pgfsetdash{}{0pt}%
\pgfpathmoveto{\pgfqpoint{5.442832in}{1.482495in}}%
\pgfpathlineto{\pgfqpoint{5.439727in}{1.485941in}}%
\pgfpathlineto{\pgfqpoint{5.442991in}{1.495641in}}%
\pgfpathlineto{\pgfqpoint{5.433301in}{1.496259in}}%
\pgfpathlineto{\pgfqpoint{5.435870in}{1.504623in}}%
\pgfpathlineto{\pgfqpoint{5.434272in}{1.517943in}}%
\pgfpathlineto{\pgfqpoint{5.425578in}{1.522561in}}%
\pgfpathlineto{\pgfqpoint{5.416480in}{1.531907in}}%
\pgfpathlineto{\pgfqpoint{5.402460in}{1.534643in}}%
\pgfpathlineto{\pgfqpoint{5.388816in}{1.534565in}}%
\pgfpathlineto{\pgfqpoint{5.375397in}{1.541198in}}%
\pgfpathlineto{\pgfqpoint{5.330502in}{1.550866in}}%
\pgfpathlineto{\pgfqpoint{5.325596in}{1.561110in}}%
\pgfpathlineto{\pgfqpoint{5.316845in}{1.564606in}}%
\pgfpathlineto{\pgfqpoint{5.333375in}{1.572449in}}%
\pgfpathlineto{\pgfqpoint{5.342713in}{1.582234in}}%
\pgfpathlineto{\pgfqpoint{5.360106in}{1.584888in}}%
\pgfpathlineto{\pgfqpoint{5.370809in}{1.594855in}}%
\pgfpathlineto{\pgfqpoint{5.376425in}{1.595254in}}%
\pgfpathlineto{\pgfqpoint{5.380710in}{1.602362in}}%
\pgfpathlineto{\pgfqpoint{5.391157in}{1.610783in}}%
\pgfpathlineto{\pgfqpoint{5.392067in}{1.604831in}}%
\pgfpathlineto{\pgfqpoint{5.396774in}{1.603602in}}%
\pgfpathlineto{\pgfqpoint{5.400316in}{1.596504in}}%
\pgfpathlineto{\pgfqpoint{5.400954in}{1.584403in}}%
\pgfpathlineto{\pgfqpoint{5.411599in}{1.591634in}}%
\pgfpathlineto{\pgfqpoint{5.424002in}{1.593145in}}%
\pgfpathlineto{\pgfqpoint{5.434590in}{1.589356in}}%
\pgfpathlineto{\pgfqpoint{5.448949in}{1.569658in}}%
\pgfpathlineto{\pgfqpoint{5.464631in}{1.572806in}}%
\pgfpathlineto{\pgfqpoint{5.471003in}{1.567535in}}%
\pgfpathlineto{\pgfqpoint{5.481252in}{1.567064in}}%
\pgfpathlineto{\pgfqpoint{5.488061in}{1.576518in}}%
\pgfpathlineto{\pgfqpoint{5.500982in}{1.584876in}}%
\pgfpathlineto{\pgfqpoint{5.513400in}{1.587493in}}%
\pgfpathlineto{\pgfqpoint{5.528808in}{1.587766in}}%
\pgfpathlineto{\pgfqpoint{5.540068in}{1.594222in}}%
\pgfpathlineto{\pgfqpoint{5.549258in}{1.591248in}}%
\pgfpathlineto{\pgfqpoint{5.551215in}{1.577594in}}%
\pgfpathlineto{\pgfqpoint{5.570930in}{1.575456in}}%
\pgfpathlineto{\pgfqpoint{5.581638in}{1.581847in}}%
\pgfpathlineto{\pgfqpoint{5.589061in}{1.567449in}}%
\pgfpathlineto{\pgfqpoint{5.603225in}{1.550799in}}%
\pgfpathlineto{\pgfqpoint{5.583401in}{1.550895in}}%
\pgfpathlineto{\pgfqpoint{5.577137in}{1.548837in}}%
\pgfpathlineto{\pgfqpoint{5.568551in}{1.551511in}}%
\pgfpathlineto{\pgfqpoint{5.567974in}{1.539985in}}%
\pgfpathlineto{\pgfqpoint{5.552384in}{1.548993in}}%
\pgfpathlineto{\pgfqpoint{5.532335in}{1.551746in}}%
\pgfpathlineto{\pgfqpoint{5.526816in}{1.542973in}}%
\pgfpathlineto{\pgfqpoint{5.500584in}{1.538707in}}%
\pgfpathlineto{\pgfqpoint{5.497566in}{1.531276in}}%
\pgfpathlineto{\pgfqpoint{5.479312in}{1.529046in}}%
\pgfpathlineto{\pgfqpoint{5.473794in}{1.521492in}}%
\pgfpathlineto{\pgfqpoint{5.464139in}{1.519466in}}%
\pgfpathlineto{\pgfqpoint{5.456427in}{1.501552in}}%
\pgfpathlineto{\pgfqpoint{5.446663in}{1.484203in}}%
\pgfpathlineto{\pgfqpoint{5.442832in}{1.482495in}}%
\pgfpathclose%
\pgfusepath{fill}%
\end{pgfscope}%
\begin{pgfscope}%
\pgfpathrectangle{\pgfqpoint{3.525000in}{0.100000in}}{\pgfqpoint{2.862985in}{1.832778in}}%
\pgfusepath{clip}%
\pgfsetbuttcap%
\pgfsetmiterjoin%
\definecolor{currentfill}{rgb}{0.673356,0.851057,0.467128}%
\pgfsetfillcolor{currentfill}%
\pgfsetlinewidth{0.000000pt}%
\definecolor{currentstroke}{rgb}{0.000000,0.000000,0.000000}%
\pgfsetstrokecolor{currentstroke}%
\pgfsetstrokeopacity{0.000000}%
\pgfsetdash{}{0pt}%
\pgfpathmoveto{\pgfqpoint{5.499029in}{1.274203in}}%
\pgfpathlineto{\pgfqpoint{5.508339in}{1.284002in}}%
\pgfpathlineto{\pgfqpoint{5.512580in}{1.298209in}}%
\pgfpathlineto{\pgfqpoint{5.517623in}{1.306443in}}%
\pgfpathlineto{\pgfqpoint{5.520716in}{1.317649in}}%
\pgfpathlineto{\pgfqpoint{5.521678in}{1.339993in}}%
\pgfpathlineto{\pgfqpoint{5.517021in}{1.361407in}}%
\pgfpathlineto{\pgfqpoint{5.501633in}{1.394197in}}%
\pgfpathlineto{\pgfqpoint{5.505800in}{1.404507in}}%
\pgfpathlineto{\pgfqpoint{5.500372in}{1.418758in}}%
\pgfpathlineto{\pgfqpoint{5.509684in}{1.438466in}}%
\pgfpathlineto{\pgfqpoint{5.508166in}{1.460442in}}%
\pgfpathlineto{\pgfqpoint{5.514653in}{1.463207in}}%
\pgfpathlineto{\pgfqpoint{5.515468in}{1.473664in}}%
\pgfpathlineto{\pgfqpoint{5.526982in}{1.480375in}}%
\pgfpathlineto{\pgfqpoint{5.533494in}{1.479291in}}%
\pgfpathlineto{\pgfqpoint{5.535324in}{1.468077in}}%
\pgfpathlineto{\pgfqpoint{5.540407in}{1.467638in}}%
\pgfpathlineto{\pgfqpoint{5.545044in}{1.483813in}}%
\pgfpathlineto{\pgfqpoint{5.543468in}{1.496393in}}%
\pgfpathlineto{\pgfqpoint{5.546483in}{1.503617in}}%
\pgfpathlineto{\pgfqpoint{5.562785in}{1.511080in}}%
\pgfpathlineto{\pgfqpoint{5.555356in}{1.513735in}}%
\pgfpathlineto{\pgfqpoint{5.552936in}{1.520146in}}%
\pgfpathlineto{\pgfqpoint{5.558288in}{1.531404in}}%
\pgfpathlineto{\pgfqpoint{5.568841in}{1.535292in}}%
\pgfpathlineto{\pgfqpoint{5.581095in}{1.528650in}}%
\pgfpathlineto{\pgfqpoint{5.592648in}{1.528567in}}%
\pgfpathlineto{\pgfqpoint{5.598015in}{1.520808in}}%
\pgfpathlineto{\pgfqpoint{5.606106in}{1.521347in}}%
\pgfpathlineto{\pgfqpoint{5.631083in}{1.510561in}}%
\pgfpathlineto{\pgfqpoint{5.636074in}{1.500120in}}%
\pgfpathlineto{\pgfqpoint{5.630624in}{1.496504in}}%
\pgfpathlineto{\pgfqpoint{5.632303in}{1.488679in}}%
\pgfpathlineto{\pgfqpoint{5.637687in}{1.485198in}}%
\pgfpathlineto{\pgfqpoint{5.640688in}{1.475579in}}%
\pgfpathlineto{\pgfqpoint{5.640268in}{1.452144in}}%
\pgfpathlineto{\pgfqpoint{5.633159in}{1.446546in}}%
\pgfpathlineto{\pgfqpoint{5.631564in}{1.434236in}}%
\pgfpathlineto{\pgfqpoint{5.618377in}{1.422859in}}%
\pgfpathlineto{\pgfqpoint{5.619159in}{1.409089in}}%
\pgfpathlineto{\pgfqpoint{5.630717in}{1.404265in}}%
\pgfpathlineto{\pgfqpoint{5.643756in}{1.421446in}}%
\pgfpathlineto{\pgfqpoint{5.644833in}{1.427678in}}%
\pgfpathlineto{\pgfqpoint{5.661094in}{1.438035in}}%
\pgfpathlineto{\pgfqpoint{5.671435in}{1.433238in}}%
\pgfpathlineto{\pgfqpoint{5.677924in}{1.422395in}}%
\pgfpathlineto{\pgfqpoint{5.688393in}{1.384749in}}%
\pgfpathlineto{\pgfqpoint{5.693940in}{1.372839in}}%
\pgfpathlineto{\pgfqpoint{5.692203in}{1.367915in}}%
\pgfpathlineto{\pgfqpoint{5.692372in}{1.351150in}}%
\pgfpathlineto{\pgfqpoint{5.682281in}{1.352759in}}%
\pgfpathlineto{\pgfqpoint{5.676581in}{1.340231in}}%
\pgfpathlineto{\pgfqpoint{5.675791in}{1.331714in}}%
\pgfpathlineto{\pgfqpoint{5.668169in}{1.326246in}}%
\pgfpathlineto{\pgfqpoint{5.666475in}{1.309636in}}%
\pgfpathlineto{\pgfqpoint{5.655390in}{1.288653in}}%
\pgfpathlineto{\pgfqpoint{5.594752in}{1.279623in}}%
\pgfpathlineto{\pgfqpoint{5.594393in}{1.283592in}}%
\pgfpathlineto{\pgfqpoint{5.553831in}{1.279317in}}%
\pgfpathlineto{\pgfqpoint{5.499029in}{1.274203in}}%
\pgfpathclose%
\pgfusepath{fill}%
\end{pgfscope}%
\begin{pgfscope}%
\pgfpathrectangle{\pgfqpoint{3.525000in}{0.100000in}}{\pgfqpoint{2.862985in}{1.832778in}}%
\pgfusepath{clip}%
\pgfsetbuttcap%
\pgfsetmiterjoin%
\definecolor{currentfill}{rgb}{0.717186,0.879508,0.520185}%
\pgfsetfillcolor{currentfill}%
\pgfsetlinewidth{0.000000pt}%
\definecolor{currentstroke}{rgb}{0.000000,0.000000,0.000000}%
\pgfsetstrokecolor{currentstroke}%
\pgfsetstrokeopacity{0.000000}%
\pgfsetdash{}{0pt}%
\pgfpathmoveto{\pgfqpoint{3.967548in}{0.431818in}}%
\pgfpathlineto{\pgfqpoint{3.966337in}{0.430752in}}%
\pgfpathlineto{\pgfqpoint{3.960781in}{0.432173in}}%
\pgfpathlineto{\pgfqpoint{3.961196in}{0.434969in}}%
\pgfpathlineto{\pgfqpoint{3.964341in}{0.436249in}}%
\pgfpathlineto{\pgfqpoint{3.968824in}{0.442694in}}%
\pgfpathlineto{\pgfqpoint{3.969569in}{0.447051in}}%
\pgfpathlineto{\pgfqpoint{3.972090in}{0.448192in}}%
\pgfpathlineto{\pgfqpoint{3.976190in}{0.448197in}}%
\pgfpathlineto{\pgfqpoint{3.980477in}{0.462333in}}%
\pgfpathlineto{\pgfqpoint{3.982157in}{0.463496in}}%
\pgfpathlineto{\pgfqpoint{3.980651in}{0.466260in}}%
\pgfpathlineto{\pgfqpoint{3.977125in}{0.463464in}}%
\pgfpathlineto{\pgfqpoint{3.966766in}{0.467233in}}%
\pgfpathlineto{\pgfqpoint{3.960583in}{0.473582in}}%
\pgfpathlineto{\pgfqpoint{3.963280in}{0.476479in}}%
\pgfpathlineto{\pgfqpoint{3.962053in}{0.480565in}}%
\pgfpathlineto{\pgfqpoint{3.962809in}{0.485392in}}%
\pgfpathlineto{\pgfqpoint{3.961479in}{0.491035in}}%
\pgfpathlineto{\pgfqpoint{3.961266in}{0.496684in}}%
\pgfpathlineto{\pgfqpoint{3.967217in}{0.495739in}}%
\pgfpathlineto{\pgfqpoint{3.968395in}{0.497404in}}%
\pgfpathlineto{\pgfqpoint{3.971355in}{0.497243in}}%
\pgfpathlineto{\pgfqpoint{3.974764in}{0.491217in}}%
\pgfpathlineto{\pgfqpoint{3.978051in}{0.486806in}}%
\pgfpathlineto{\pgfqpoint{3.975418in}{0.483935in}}%
\pgfpathlineto{\pgfqpoint{3.979970in}{0.482833in}}%
\pgfpathlineto{\pgfqpoint{3.981007in}{0.484241in}}%
\pgfpathlineto{\pgfqpoint{3.978573in}{0.489956in}}%
\pgfpathlineto{\pgfqpoint{3.973447in}{0.494692in}}%
\pgfpathlineto{\pgfqpoint{3.974074in}{0.496676in}}%
\pgfpathlineto{\pgfqpoint{3.969984in}{0.501691in}}%
\pgfpathlineto{\pgfqpoint{3.974161in}{0.502710in}}%
\pgfpathlineto{\pgfqpoint{3.970542in}{0.513448in}}%
\pgfpathlineto{\pgfqpoint{3.973484in}{0.514936in}}%
\pgfpathlineto{\pgfqpoint{3.974202in}{0.517511in}}%
\pgfpathlineto{\pgfqpoint{3.972288in}{0.520569in}}%
\pgfpathlineto{\pgfqpoint{3.976239in}{0.521335in}}%
\pgfpathlineto{\pgfqpoint{3.976948in}{0.523648in}}%
\pgfpathlineto{\pgfqpoint{3.981660in}{0.520176in}}%
\pgfpathlineto{\pgfqpoint{3.982707in}{0.523941in}}%
\pgfpathlineto{\pgfqpoint{3.985183in}{0.525597in}}%
\pgfpathlineto{\pgfqpoint{3.997233in}{0.527901in}}%
\pgfpathlineto{\pgfqpoint{4.000932in}{0.527086in}}%
\pgfpathlineto{\pgfqpoint{4.004031in}{0.531659in}}%
\pgfpathlineto{\pgfqpoint{4.011384in}{0.534658in}}%
\pgfpathlineto{\pgfqpoint{4.017047in}{0.533521in}}%
\pgfpathlineto{\pgfqpoint{4.019631in}{0.529538in}}%
\pgfpathlineto{\pgfqpoint{4.019495in}{0.523212in}}%
\pgfpathlineto{\pgfqpoint{4.022394in}{0.521778in}}%
\pgfpathlineto{\pgfqpoint{4.028306in}{0.521978in}}%
\pgfpathlineto{\pgfqpoint{4.035878in}{0.523624in}}%
\pgfpathlineto{\pgfqpoint{4.035785in}{0.521601in}}%
\pgfpathlineto{\pgfqpoint{4.045357in}{0.515832in}}%
\pgfpathlineto{\pgfqpoint{4.052052in}{0.517379in}}%
\pgfpathlineto{\pgfqpoint{4.054022in}{0.519381in}}%
\pgfpathlineto{\pgfqpoint{4.056133in}{0.524508in}}%
\pgfpathlineto{\pgfqpoint{4.058750in}{0.527792in}}%
\pgfpathlineto{\pgfqpoint{4.058984in}{0.532712in}}%
\pgfpathlineto{\pgfqpoint{4.056587in}{0.534647in}}%
\pgfpathlineto{\pgfqpoint{4.060136in}{0.536396in}}%
\pgfpathlineto{\pgfqpoint{4.066098in}{0.533768in}}%
\pgfpathlineto{\pgfqpoint{4.068040in}{0.535390in}}%
\pgfpathlineto{\pgfqpoint{4.068382in}{0.542352in}}%
\pgfpathlineto{\pgfqpoint{4.064528in}{0.540897in}}%
\pgfpathlineto{\pgfqpoint{4.061855in}{0.543309in}}%
\pgfpathlineto{\pgfqpoint{4.057516in}{0.544812in}}%
\pgfpathlineto{\pgfqpoint{4.050629in}{0.544561in}}%
\pgfpathlineto{\pgfqpoint{4.044030in}{0.546848in}}%
\pgfpathlineto{\pgfqpoint{4.043699in}{0.552545in}}%
\pgfpathlineto{\pgfqpoint{4.037797in}{0.558597in}}%
\pgfpathlineto{\pgfqpoint{4.030313in}{0.561138in}}%
\pgfpathlineto{\pgfqpoint{4.023780in}{0.572853in}}%
\pgfpathlineto{\pgfqpoint{4.024316in}{0.577189in}}%
\pgfpathlineto{\pgfqpoint{4.027833in}{0.579165in}}%
\pgfpathlineto{\pgfqpoint{4.027315in}{0.583269in}}%
\pgfpathlineto{\pgfqpoint{4.028132in}{0.587352in}}%
\pgfpathlineto{\pgfqpoint{4.030915in}{0.584557in}}%
\pgfpathlineto{\pgfqpoint{4.034590in}{0.585723in}}%
\pgfpathlineto{\pgfqpoint{4.034407in}{0.589270in}}%
\pgfpathlineto{\pgfqpoint{4.029311in}{0.596390in}}%
\pgfpathlineto{\pgfqpoint{4.027820in}{0.603734in}}%
\pgfpathlineto{\pgfqpoint{4.032085in}{0.604294in}}%
\pgfpathlineto{\pgfqpoint{4.038468in}{0.603843in}}%
\pgfpathlineto{\pgfqpoint{4.040478in}{0.601387in}}%
\pgfpathlineto{\pgfqpoint{4.046201in}{0.602266in}}%
\pgfpathlineto{\pgfqpoint{4.051520in}{0.601112in}}%
\pgfpathlineto{\pgfqpoint{4.055115in}{0.599083in}}%
\pgfpathlineto{\pgfqpoint{4.057604in}{0.600010in}}%
\pgfpathlineto{\pgfqpoint{4.064212in}{0.597983in}}%
\pgfpathlineto{\pgfqpoint{4.064273in}{0.596187in}}%
\pgfpathlineto{\pgfqpoint{4.069305in}{0.596794in}}%
\pgfpathlineto{\pgfqpoint{4.076053in}{0.592414in}}%
\pgfpathlineto{\pgfqpoint{4.075791in}{0.589506in}}%
\pgfpathlineto{\pgfqpoint{4.072580in}{0.588164in}}%
\pgfpathlineto{\pgfqpoint{4.069167in}{0.584972in}}%
\pgfpathlineto{\pgfqpoint{4.068755in}{0.580588in}}%
\pgfpathlineto{\pgfqpoint{4.076106in}{0.574997in}}%
\pgfpathlineto{\pgfqpoint{4.074930in}{0.572318in}}%
\pgfpathlineto{\pgfqpoint{4.080161in}{0.570260in}}%
\pgfpathlineto{\pgfqpoint{4.081445in}{0.566900in}}%
\pgfpathlineto{\pgfqpoint{4.088083in}{0.569570in}}%
\pgfpathlineto{\pgfqpoint{4.091007in}{0.566071in}}%
\pgfpathlineto{\pgfqpoint{4.092356in}{0.568189in}}%
\pgfpathlineto{\pgfqpoint{4.088383in}{0.575083in}}%
\pgfpathlineto{\pgfqpoint{4.089795in}{0.577148in}}%
\pgfpathlineto{\pgfqpoint{4.090421in}{0.582594in}}%
\pgfpathlineto{\pgfqpoint{4.088838in}{0.584991in}}%
\pgfpathlineto{\pgfqpoint{4.090025in}{0.588211in}}%
\pgfpathlineto{\pgfqpoint{4.093915in}{0.587908in}}%
\pgfpathlineto{\pgfqpoint{4.093110in}{0.582429in}}%
\pgfpathlineto{\pgfqpoint{4.090627in}{0.580602in}}%
\pgfpathlineto{\pgfqpoint{4.090803in}{0.573231in}}%
\pgfpathlineto{\pgfqpoint{4.095258in}{0.572350in}}%
\pgfpathlineto{\pgfqpoint{4.095312in}{0.566870in}}%
\pgfpathlineto{\pgfqpoint{4.100178in}{0.563315in}}%
\pgfpathlineto{\pgfqpoint{4.102415in}{0.565557in}}%
\pgfpathlineto{\pgfqpoint{4.099970in}{0.572380in}}%
\pgfpathlineto{\pgfqpoint{4.097442in}{0.574106in}}%
\pgfpathlineto{\pgfqpoint{4.094955in}{0.572970in}}%
\pgfpathlineto{\pgfqpoint{4.092767in}{0.574420in}}%
\pgfpathlineto{\pgfqpoint{4.093149in}{0.580725in}}%
\pgfpathlineto{\pgfqpoint{4.096800in}{0.583304in}}%
\pgfpathlineto{\pgfqpoint{4.100395in}{0.583066in}}%
\pgfpathlineto{\pgfqpoint{4.098905in}{0.586634in}}%
\pgfpathlineto{\pgfqpoint{4.093518in}{0.589366in}}%
\pgfpathlineto{\pgfqpoint{4.086383in}{0.600253in}}%
\pgfpathlineto{\pgfqpoint{4.089944in}{0.605475in}}%
\pgfpathlineto{\pgfqpoint{4.091926in}{0.610393in}}%
\pgfpathlineto{\pgfqpoint{4.092082in}{0.620269in}}%
\pgfpathlineto{\pgfqpoint{4.091319in}{0.628601in}}%
\pgfpathlineto{\pgfqpoint{4.089142in}{0.634072in}}%
\pgfpathlineto{\pgfqpoint{4.089576in}{0.639541in}}%
\pgfpathlineto{\pgfqpoint{4.095144in}{0.644285in}}%
\pgfpathlineto{\pgfqpoint{4.100681in}{0.649922in}}%
\pgfpathlineto{\pgfqpoint{4.114564in}{0.638405in}}%
\pgfpathlineto{\pgfqpoint{4.123640in}{0.637556in}}%
\pgfpathlineto{\pgfqpoint{4.131132in}{0.640424in}}%
\pgfpathlineto{\pgfqpoint{4.137708in}{0.644201in}}%
\pgfpathlineto{\pgfqpoint{4.139105in}{0.646000in}}%
\pgfpathlineto{\pgfqpoint{4.155340in}{0.650103in}}%
\pgfpathlineto{\pgfqpoint{4.156749in}{0.647098in}}%
\pgfpathlineto{\pgfqpoint{4.162266in}{0.644638in}}%
\pgfpathlineto{\pgfqpoint{4.172793in}{0.645392in}}%
\pgfpathlineto{\pgfqpoint{4.174716in}{0.644720in}}%
\pgfpathlineto{\pgfqpoint{4.179448in}{0.646515in}}%
\pgfpathlineto{\pgfqpoint{4.181350in}{0.643411in}}%
\pgfpathlineto{\pgfqpoint{4.185894in}{0.640575in}}%
\pgfpathlineto{\pgfqpoint{4.188167in}{0.640896in}}%
\pgfpathlineto{\pgfqpoint{4.192566in}{0.637488in}}%
\pgfpathlineto{\pgfqpoint{4.200000in}{0.638115in}}%
\pgfpathlineto{\pgfqpoint{4.207425in}{0.640874in}}%
\pgfpathlineto{\pgfqpoint{4.213489in}{0.632051in}}%
\pgfpathlineto{\pgfqpoint{4.213019in}{0.630071in}}%
\pgfpathlineto{\pgfqpoint{4.206328in}{0.627796in}}%
\pgfpathlineto{\pgfqpoint{4.208115in}{0.624355in}}%
\pgfpathlineto{\pgfqpoint{4.214275in}{0.627976in}}%
\pgfpathlineto{\pgfqpoint{4.216790in}{0.627100in}}%
\pgfpathlineto{\pgfqpoint{4.217860in}{0.622768in}}%
\pgfpathlineto{\pgfqpoint{4.214183in}{0.620936in}}%
\pgfpathlineto{\pgfqpoint{4.216809in}{0.615461in}}%
\pgfpathlineto{\pgfqpoint{4.220623in}{0.616386in}}%
\pgfpathlineto{\pgfqpoint{4.226157in}{0.613560in}}%
\pgfpathlineto{\pgfqpoint{4.231450in}{0.605718in}}%
\pgfpathlineto{\pgfqpoint{4.226793in}{0.603540in}}%
\pgfpathlineto{\pgfqpoint{4.230907in}{0.597761in}}%
\pgfpathlineto{\pgfqpoint{4.227441in}{0.596410in}}%
\pgfpathlineto{\pgfqpoint{4.231753in}{0.590922in}}%
\pgfpathlineto{\pgfqpoint{4.237012in}{0.591506in}}%
\pgfpathlineto{\pgfqpoint{4.239781in}{0.589037in}}%
\pgfpathlineto{\pgfqpoint{4.246785in}{0.585396in}}%
\pgfpathlineto{\pgfqpoint{4.255820in}{0.572470in}}%
\pgfpathlineto{\pgfqpoint{4.255616in}{0.569379in}}%
\pgfpathlineto{\pgfqpoint{4.269166in}{0.559477in}}%
\pgfpathlineto{\pgfqpoint{4.269458in}{0.555975in}}%
\pgfpathlineto{\pgfqpoint{4.273069in}{0.550735in}}%
\pgfpathlineto{\pgfqpoint{4.284163in}{0.548115in}}%
\pgfpathlineto{\pgfqpoint{4.288277in}{0.546252in}}%
\pgfpathlineto{\pgfqpoint{4.292148in}{0.540476in}}%
\pgfpathlineto{\pgfqpoint{4.292646in}{0.535415in}}%
\pgfpathlineto{\pgfqpoint{4.295973in}{0.531291in}}%
\pgfpathlineto{\pgfqpoint{4.296972in}{0.526640in}}%
\pgfpathlineto{\pgfqpoint{4.300090in}{0.524941in}}%
\pgfpathlineto{\pgfqpoint{4.284646in}{0.497067in}}%
\pgfpathlineto{\pgfqpoint{4.252590in}{0.439200in}}%
\pgfpathlineto{\pgfqpoint{4.205557in}{0.354280in}}%
\pgfpathlineto{\pgfqpoint{4.187143in}{0.321025in}}%
\pgfpathlineto{\pgfqpoint{4.191065in}{0.316559in}}%
\pgfpathlineto{\pgfqpoint{4.192793in}{0.317984in}}%
\pgfpathlineto{\pgfqpoint{4.196296in}{0.312785in}}%
\pgfpathlineto{\pgfqpoint{4.201159in}{0.314293in}}%
\pgfpathlineto{\pgfqpoint{4.207614in}{0.311178in}}%
\pgfpathlineto{\pgfqpoint{4.203433in}{0.306310in}}%
\pgfpathlineto{\pgfqpoint{4.206662in}{0.299741in}}%
\pgfpathlineto{\pgfqpoint{4.205995in}{0.296543in}}%
\pgfpathlineto{\pgfqpoint{4.211145in}{0.279677in}}%
\pgfpathlineto{\pgfqpoint{4.208966in}{0.272007in}}%
\pgfpathlineto{\pgfqpoint{4.218805in}{0.273906in}}%
\pgfpathlineto{\pgfqpoint{4.221221in}{0.272662in}}%
\pgfpathlineto{\pgfqpoint{4.223844in}{0.274685in}}%
\pgfpathlineto{\pgfqpoint{4.225709in}{0.278509in}}%
\pgfpathlineto{\pgfqpoint{4.228381in}{0.280701in}}%
\pgfpathlineto{\pgfqpoint{4.239784in}{0.280424in}}%
\pgfpathlineto{\pgfqpoint{4.242361in}{0.273064in}}%
\pgfpathlineto{\pgfqpoint{4.240149in}{0.266603in}}%
\pgfpathlineto{\pgfqpoint{4.242662in}{0.264471in}}%
\pgfpathlineto{\pgfqpoint{4.243983in}{0.260825in}}%
\pgfpathlineto{\pgfqpoint{4.243695in}{0.253967in}}%
\pgfpathlineto{\pgfqpoint{4.247048in}{0.249051in}}%
\pgfpathlineto{\pgfqpoint{4.248886in}{0.241346in}}%
\pgfpathlineto{\pgfqpoint{4.247942in}{0.237184in}}%
\pgfpathlineto{\pgfqpoint{4.249080in}{0.232966in}}%
\pgfpathlineto{\pgfqpoint{4.250722in}{0.208539in}}%
\pgfpathlineto{\pgfqpoint{4.248370in}{0.206615in}}%
\pgfpathlineto{\pgfqpoint{4.251472in}{0.203957in}}%
\pgfpathlineto{\pgfqpoint{4.249034in}{0.200723in}}%
\pgfpathlineto{\pgfqpoint{4.251233in}{0.198087in}}%
\pgfpathlineto{\pgfqpoint{4.249817in}{0.193586in}}%
\pgfpathlineto{\pgfqpoint{4.252984in}{0.192366in}}%
\pgfpathlineto{\pgfqpoint{4.257058in}{0.185522in}}%
\pgfpathlineto{\pgfqpoint{4.260085in}{0.183460in}}%
\pgfpathlineto{\pgfqpoint{4.260877in}{0.180445in}}%
\pgfpathlineto{\pgfqpoint{4.262631in}{0.179139in}}%
\pgfpathlineto{\pgfqpoint{4.262265in}{0.176636in}}%
\pgfpathlineto{\pgfqpoint{4.266046in}{0.174822in}}%
\pgfpathlineto{\pgfqpoint{4.265112in}{0.170006in}}%
\pgfpathlineto{\pgfqpoint{4.261858in}{0.167483in}}%
\pgfpathlineto{\pgfqpoint{4.260380in}{0.162795in}}%
\pgfpathlineto{\pgfqpoint{4.259988in}{0.156518in}}%
\pgfpathlineto{\pgfqpoint{4.258180in}{0.155917in}}%
\pgfpathlineto{\pgfqpoint{4.252271in}{0.150415in}}%
\pgfpathlineto{\pgfqpoint{4.247019in}{0.148962in}}%
\pgfpathlineto{\pgfqpoint{4.244513in}{0.151181in}}%
\pgfpathlineto{\pgfqpoint{4.246935in}{0.157071in}}%
\pgfpathlineto{\pgfqpoint{4.246195in}{0.159962in}}%
\pgfpathlineto{\pgfqpoint{4.249665in}{0.161426in}}%
\pgfpathlineto{\pgfqpoint{4.252909in}{0.169947in}}%
\pgfpathlineto{\pgfqpoint{4.251539in}{0.177225in}}%
\pgfpathlineto{\pgfqpoint{4.249346in}{0.178943in}}%
\pgfpathlineto{\pgfqpoint{4.242012in}{0.178387in}}%
\pgfpathlineto{\pgfqpoint{4.243490in}{0.176486in}}%
\pgfpathlineto{\pgfqpoint{4.238181in}{0.171028in}}%
\pgfpathlineto{\pgfqpoint{4.236471in}{0.174063in}}%
\pgfpathlineto{\pgfqpoint{4.236916in}{0.178027in}}%
\pgfpathlineto{\pgfqpoint{4.239963in}{0.177954in}}%
\pgfpathlineto{\pgfqpoint{4.242677in}{0.180866in}}%
\pgfpathlineto{\pgfqpoint{4.244282in}{0.185118in}}%
\pgfpathlineto{\pgfqpoint{4.250331in}{0.183918in}}%
\pgfpathlineto{\pgfqpoint{4.245488in}{0.186190in}}%
\pgfpathlineto{\pgfqpoint{4.245915in}{0.189233in}}%
\pgfpathlineto{\pgfqpoint{4.243899in}{0.190607in}}%
\pgfpathlineto{\pgfqpoint{4.243074in}{0.194752in}}%
\pgfpathlineto{\pgfqpoint{4.244505in}{0.196911in}}%
\pgfpathlineto{\pgfqpoint{4.241093in}{0.203617in}}%
\pgfpathlineto{\pgfqpoint{4.241381in}{0.207883in}}%
\pgfpathlineto{\pgfqpoint{4.237383in}{0.211682in}}%
\pgfpathlineto{\pgfqpoint{4.235220in}{0.214756in}}%
\pgfpathlineto{\pgfqpoint{4.238207in}{0.217685in}}%
\pgfpathlineto{\pgfqpoint{4.239313in}{0.222063in}}%
\pgfpathlineto{\pgfqpoint{4.238326in}{0.224937in}}%
\pgfpathlineto{\pgfqpoint{4.239481in}{0.228366in}}%
\pgfpathlineto{\pgfqpoint{4.238009in}{0.239543in}}%
\pgfpathlineto{\pgfqpoint{4.235807in}{0.244614in}}%
\pgfpathlineto{\pgfqpoint{4.233288in}{0.246519in}}%
\pgfpathlineto{\pgfqpoint{4.234360in}{0.253116in}}%
\pgfpathlineto{\pgfqpoint{4.233604in}{0.258268in}}%
\pgfpathlineto{\pgfqpoint{4.234974in}{0.261572in}}%
\pgfpathlineto{\pgfqpoint{4.232140in}{0.262035in}}%
\pgfpathlineto{\pgfqpoint{4.231380in}{0.254009in}}%
\pgfpathlineto{\pgfqpoint{4.228196in}{0.245055in}}%
\pgfpathlineto{\pgfqpoint{4.225639in}{0.247094in}}%
\pgfpathlineto{\pgfqpoint{4.225297in}{0.250037in}}%
\pgfpathlineto{\pgfqpoint{4.222168in}{0.254057in}}%
\pgfpathlineto{\pgfqpoint{4.223758in}{0.256688in}}%
\pgfpathlineto{\pgfqpoint{4.223264in}{0.262613in}}%
\pgfpathlineto{\pgfqpoint{4.219215in}{0.260530in}}%
\pgfpathlineto{\pgfqpoint{4.220009in}{0.257316in}}%
\pgfpathlineto{\pgfqpoint{4.219247in}{0.253236in}}%
\pgfpathlineto{\pgfqpoint{4.214702in}{0.253244in}}%
\pgfpathlineto{\pgfqpoint{4.212665in}{0.255773in}}%
\pgfpathlineto{\pgfqpoint{4.209590in}{0.256255in}}%
\pgfpathlineto{\pgfqpoint{4.207521in}{0.258978in}}%
\pgfpathlineto{\pgfqpoint{4.203961in}{0.266454in}}%
\pgfpathlineto{\pgfqpoint{4.202807in}{0.272259in}}%
\pgfpathlineto{\pgfqpoint{4.203546in}{0.274726in}}%
\pgfpathlineto{\pgfqpoint{4.201822in}{0.280160in}}%
\pgfpathlineto{\pgfqpoint{4.198967in}{0.283305in}}%
\pgfpathlineto{\pgfqpoint{4.194029in}{0.291344in}}%
\pgfpathlineto{\pgfqpoint{4.190447in}{0.298326in}}%
\pgfpathlineto{\pgfqpoint{4.194102in}{0.298431in}}%
\pgfpathlineto{\pgfqpoint{4.196252in}{0.299897in}}%
\pgfpathlineto{\pgfqpoint{4.196675in}{0.304435in}}%
\pgfpathlineto{\pgfqpoint{4.186334in}{0.304983in}}%
\pgfpathlineto{\pgfqpoint{4.177539in}{0.314807in}}%
\pgfpathlineto{\pgfqpoint{4.179852in}{0.317392in}}%
\pgfpathlineto{\pgfqpoint{4.175599in}{0.317833in}}%
\pgfpathlineto{\pgfqpoint{4.167152in}{0.326967in}}%
\pgfpathlineto{\pgfqpoint{4.153770in}{0.331547in}}%
\pgfpathlineto{\pgfqpoint{4.151968in}{0.336915in}}%
\pgfpathlineto{\pgfqpoint{4.149191in}{0.339711in}}%
\pgfpathlineto{\pgfqpoint{4.149065in}{0.342336in}}%
\pgfpathlineto{\pgfqpoint{4.146976in}{0.344286in}}%
\pgfpathlineto{\pgfqpoint{4.149717in}{0.347544in}}%
\pgfpathlineto{\pgfqpoint{4.143465in}{0.347858in}}%
\pgfpathlineto{\pgfqpoint{4.141876in}{0.352075in}}%
\pgfpathlineto{\pgfqpoint{4.139030in}{0.353870in}}%
\pgfpathlineto{\pgfqpoint{4.144729in}{0.356037in}}%
\pgfpathlineto{\pgfqpoint{4.141765in}{0.359902in}}%
\pgfpathlineto{\pgfqpoint{4.136311in}{0.361524in}}%
\pgfpathlineto{\pgfqpoint{4.137160in}{0.368663in}}%
\pgfpathlineto{\pgfqpoint{4.135596in}{0.371190in}}%
\pgfpathlineto{\pgfqpoint{4.132780in}{0.372970in}}%
\pgfpathlineto{\pgfqpoint{4.130372in}{0.371185in}}%
\pgfpathlineto{\pgfqpoint{4.128950in}{0.372864in}}%
\pgfpathlineto{\pgfqpoint{4.124800in}{0.373345in}}%
\pgfpathlineto{\pgfqpoint{4.125134in}{0.379823in}}%
\pgfpathlineto{\pgfqpoint{4.119119in}{0.375448in}}%
\pgfpathlineto{\pgfqpoint{4.119246in}{0.366430in}}%
\pgfpathlineto{\pgfqpoint{4.112788in}{0.364702in}}%
\pgfpathlineto{\pgfqpoint{4.108572in}{0.361101in}}%
\pgfpathlineto{\pgfqpoint{4.105182in}{0.360342in}}%
\pgfpathlineto{\pgfqpoint{4.101514in}{0.363971in}}%
\pgfpathlineto{\pgfqpoint{4.098248in}{0.363490in}}%
\pgfpathlineto{\pgfqpoint{4.098854in}{0.366805in}}%
\pgfpathlineto{\pgfqpoint{4.094211in}{0.365020in}}%
\pgfpathlineto{\pgfqpoint{4.089909in}{0.364820in}}%
\pgfpathlineto{\pgfqpoint{4.084874in}{0.363010in}}%
\pgfpathlineto{\pgfqpoint{4.074656in}{0.364028in}}%
\pgfpathlineto{\pgfqpoint{4.071376in}{0.363358in}}%
\pgfpathlineto{\pgfqpoint{4.069009in}{0.366339in}}%
\pgfpathlineto{\pgfqpoint{4.064640in}{0.366637in}}%
\pgfpathlineto{\pgfqpoint{4.063731in}{0.370500in}}%
\pgfpathlineto{\pgfqpoint{4.066335in}{0.372604in}}%
\pgfpathlineto{\pgfqpoint{4.069390in}{0.372406in}}%
\pgfpathlineto{\pgfqpoint{4.071940in}{0.369537in}}%
\pgfpathlineto{\pgfqpoint{4.073619in}{0.373944in}}%
\pgfpathlineto{\pgfqpoint{4.071276in}{0.378827in}}%
\pgfpathlineto{\pgfqpoint{4.076579in}{0.383125in}}%
\pgfpathlineto{\pgfqpoint{4.081893in}{0.384709in}}%
\pgfpathlineto{\pgfqpoint{4.085593in}{0.387292in}}%
\pgfpathlineto{\pgfqpoint{4.087990in}{0.390076in}}%
\pgfpathlineto{\pgfqpoint{4.089327in}{0.394638in}}%
\pgfpathlineto{\pgfqpoint{4.093552in}{0.393396in}}%
\pgfpathlineto{\pgfqpoint{4.102714in}{0.394034in}}%
\pgfpathlineto{\pgfqpoint{4.104680in}{0.389032in}}%
\pgfpathlineto{\pgfqpoint{4.107894in}{0.388814in}}%
\pgfpathlineto{\pgfqpoint{4.113742in}{0.381156in}}%
\pgfpathlineto{\pgfqpoint{4.113753in}{0.383947in}}%
\pgfpathlineto{\pgfqpoint{4.111710in}{0.384855in}}%
\pgfpathlineto{\pgfqpoint{4.107971in}{0.393996in}}%
\pgfpathlineto{\pgfqpoint{4.109790in}{0.395238in}}%
\pgfpathlineto{\pgfqpoint{4.104797in}{0.400007in}}%
\pgfpathlineto{\pgfqpoint{4.100334in}{0.400936in}}%
\pgfpathlineto{\pgfqpoint{4.096081in}{0.399287in}}%
\pgfpathlineto{\pgfqpoint{4.088798in}{0.400571in}}%
\pgfpathlineto{\pgfqpoint{4.085344in}{0.398261in}}%
\pgfpathlineto{\pgfqpoint{4.077478in}{0.396476in}}%
\pgfpathlineto{\pgfqpoint{4.076759in}{0.393907in}}%
\pgfpathlineto{\pgfqpoint{4.070716in}{0.393207in}}%
\pgfpathlineto{\pgfqpoint{4.069165in}{0.390270in}}%
\pgfpathlineto{\pgfqpoint{4.065709in}{0.388119in}}%
\pgfpathlineto{\pgfqpoint{4.061533in}{0.388489in}}%
\pgfpathlineto{\pgfqpoint{4.059448in}{0.386408in}}%
\pgfpathlineto{\pgfqpoint{4.056818in}{0.386506in}}%
\pgfpathlineto{\pgfqpoint{4.054304in}{0.388378in}}%
\pgfpathlineto{\pgfqpoint{4.049909in}{0.388044in}}%
\pgfpathlineto{\pgfqpoint{4.048874in}{0.386318in}}%
\pgfpathlineto{\pgfqpoint{4.044258in}{0.388090in}}%
\pgfpathlineto{\pgfqpoint{4.040339in}{0.383518in}}%
\pgfpathlineto{\pgfqpoint{4.044097in}{0.379095in}}%
\pgfpathlineto{\pgfqpoint{4.045543in}{0.375457in}}%
\pgfpathlineto{\pgfqpoint{4.044210in}{0.372376in}}%
\pgfpathlineto{\pgfqpoint{4.038870in}{0.370250in}}%
\pgfpathlineto{\pgfqpoint{4.035776in}{0.371980in}}%
\pgfpathlineto{\pgfqpoint{4.031882in}{0.371043in}}%
\pgfpathlineto{\pgfqpoint{4.031413in}{0.368343in}}%
\pgfpathlineto{\pgfqpoint{4.026673in}{0.369298in}}%
\pgfpathlineto{\pgfqpoint{4.027722in}{0.366600in}}%
\pgfpathlineto{\pgfqpoint{4.020731in}{0.366031in}}%
\pgfpathlineto{\pgfqpoint{4.016205in}{0.369456in}}%
\pgfpathlineto{\pgfqpoint{4.013787in}{0.367278in}}%
\pgfpathlineto{\pgfqpoint{4.007354in}{0.369597in}}%
\pgfpathlineto{\pgfqpoint{4.005752in}{0.367248in}}%
\pgfpathlineto{\pgfqpoint{4.002563in}{0.366203in}}%
\pgfpathlineto{\pgfqpoint{4.000005in}{0.368957in}}%
\pgfpathlineto{\pgfqpoint{3.991257in}{0.368272in}}%
\pgfpathlineto{\pgfqpoint{3.991266in}{0.364211in}}%
\pgfpathlineto{\pgfqpoint{3.987752in}{0.363136in}}%
\pgfpathlineto{\pgfqpoint{3.985018in}{0.365224in}}%
\pgfpathlineto{\pgfqpoint{3.981404in}{0.364064in}}%
\pgfpathlineto{\pgfqpoint{3.976727in}{0.363787in}}%
\pgfpathlineto{\pgfqpoint{3.969759in}{0.366724in}}%
\pgfpathlineto{\pgfqpoint{3.965769in}{0.363696in}}%
\pgfpathlineto{\pgfqpoint{3.960063in}{0.367743in}}%
\pgfpathlineto{\pgfqpoint{3.958197in}{0.366209in}}%
\pgfpathlineto{\pgfqpoint{3.956887in}{0.362246in}}%
\pgfpathlineto{\pgfqpoint{3.953218in}{0.359922in}}%
\pgfpathlineto{\pgfqpoint{3.947236in}{0.362582in}}%
\pgfpathlineto{\pgfqpoint{3.937326in}{0.369153in}}%
\pgfpathlineto{\pgfqpoint{3.934438in}{0.369737in}}%
\pgfpathlineto{\pgfqpoint{3.932808in}{0.368368in}}%
\pgfpathlineto{\pgfqpoint{3.928449in}{0.369299in}}%
\pgfpathlineto{\pgfqpoint{3.926322in}{0.371344in}}%
\pgfpathlineto{\pgfqpoint{3.921969in}{0.372446in}}%
\pgfpathlineto{\pgfqpoint{3.916023in}{0.372510in}}%
\pgfpathlineto{\pgfqpoint{3.913662in}{0.374950in}}%
\pgfpathlineto{\pgfqpoint{3.918187in}{0.377474in}}%
\pgfpathlineto{\pgfqpoint{3.917131in}{0.379966in}}%
\pgfpathlineto{\pgfqpoint{3.914022in}{0.379406in}}%
\pgfpathlineto{\pgfqpoint{3.912470in}{0.377447in}}%
\pgfpathlineto{\pgfqpoint{3.905192in}{0.374639in}}%
\pgfpathlineto{\pgfqpoint{3.901512in}{0.375594in}}%
\pgfpathlineto{\pgfqpoint{3.899964in}{0.381297in}}%
\pgfpathlineto{\pgfqpoint{3.895373in}{0.378845in}}%
\pgfpathlineto{\pgfqpoint{3.894039in}{0.386073in}}%
\pgfpathlineto{\pgfqpoint{3.891864in}{0.384210in}}%
\pgfpathlineto{\pgfqpoint{3.892228in}{0.381429in}}%
\pgfpathlineto{\pgfqpoint{3.887226in}{0.382163in}}%
\pgfpathlineto{\pgfqpoint{3.892429in}{0.387010in}}%
\pgfpathlineto{\pgfqpoint{3.897702in}{0.384130in}}%
\pgfpathlineto{\pgfqpoint{3.898833in}{0.385576in}}%
\pgfpathlineto{\pgfqpoint{3.904380in}{0.383843in}}%
\pgfpathlineto{\pgfqpoint{3.904033in}{0.385816in}}%
\pgfpathlineto{\pgfqpoint{3.919168in}{0.386469in}}%
\pgfpathlineto{\pgfqpoint{3.926551in}{0.383391in}}%
\pgfpathlineto{\pgfqpoint{3.929408in}{0.380355in}}%
\pgfpathlineto{\pgfqpoint{3.928561in}{0.377705in}}%
\pgfpathlineto{\pgfqpoint{3.931429in}{0.376702in}}%
\pgfpathlineto{\pgfqpoint{3.938558in}{0.380885in}}%
\pgfpathlineto{\pgfqpoint{3.947779in}{0.381153in}}%
\pgfpathlineto{\pgfqpoint{3.956155in}{0.378631in}}%
\pgfpathlineto{\pgfqpoint{3.960899in}{0.379054in}}%
\pgfpathlineto{\pgfqpoint{3.962665in}{0.375142in}}%
\pgfpathlineto{\pgfqpoint{3.965318in}{0.379455in}}%
\pgfpathlineto{\pgfqpoint{3.966950in}{0.380389in}}%
\pgfpathlineto{\pgfqpoint{3.974498in}{0.382264in}}%
\pgfpathlineto{\pgfqpoint{3.977304in}{0.381118in}}%
\pgfpathlineto{\pgfqpoint{3.980338in}{0.382479in}}%
\pgfpathlineto{\pgfqpoint{3.983958in}{0.381767in}}%
\pgfpathlineto{\pgfqpoint{3.984812in}{0.383352in}}%
\pgfpathlineto{\pgfqpoint{3.992363in}{0.390625in}}%
\pgfpathlineto{\pgfqpoint{3.995561in}{0.391703in}}%
\pgfpathlineto{\pgfqpoint{3.997451in}{0.395629in}}%
\pgfpathlineto{\pgfqpoint{4.007481in}{0.398668in}}%
\pgfpathlineto{\pgfqpoint{4.008737in}{0.400987in}}%
\pgfpathlineto{\pgfqpoint{3.994522in}{0.404566in}}%
\pgfpathlineto{\pgfqpoint{3.995081in}{0.408415in}}%
\pgfpathlineto{\pgfqpoint{3.993619in}{0.413499in}}%
\pgfpathlineto{\pgfqpoint{3.990309in}{0.413077in}}%
\pgfpathlineto{\pgfqpoint{3.988057in}{0.406479in}}%
\pgfpathlineto{\pgfqpoint{3.985443in}{0.405936in}}%
\pgfpathlineto{\pgfqpoint{3.983742in}{0.408071in}}%
\pgfpathlineto{\pgfqpoint{3.985726in}{0.417413in}}%
\pgfpathlineto{\pgfqpoint{3.983750in}{0.421328in}}%
\pgfpathlineto{\pgfqpoint{3.979922in}{0.425943in}}%
\pgfpathlineto{\pgfqpoint{3.981762in}{0.429465in}}%
\pgfpathlineto{\pgfqpoint{3.973683in}{0.429591in}}%
\pgfpathlineto{\pgfqpoint{3.973867in}{0.430958in}}%
\pgfpathlineto{\pgfqpoint{3.969046in}{0.432871in}}%
\pgfpathlineto{\pgfqpoint{3.967548in}{0.431818in}}%
\pgfpathclose%
\pgfusepath{fill}%
\end{pgfscope}%
\begin{pgfscope}%
\pgfpathrectangle{\pgfqpoint{3.525000in}{0.100000in}}{\pgfqpoint{2.862985in}{1.832778in}}%
\pgfusepath{clip}%
\pgfsetbuttcap%
\pgfsetmiterjoin%
\definecolor{currentfill}{rgb}{0.717186,0.879508,0.520185}%
\pgfsetfillcolor{currentfill}%
\pgfsetlinewidth{0.000000pt}%
\definecolor{currentstroke}{rgb}{0.000000,0.000000,0.000000}%
\pgfsetstrokecolor{currentstroke}%
\pgfsetstrokeopacity{0.000000}%
\pgfsetdash{}{0pt}%
\pgfpathmoveto{\pgfqpoint{3.953131in}{0.499862in}}%
\pgfpathlineto{\pgfqpoint{3.952132in}{0.498106in}}%
\pgfpathlineto{\pgfqpoint{3.954775in}{0.494368in}}%
\pgfpathlineto{\pgfqpoint{3.950597in}{0.490871in}}%
\pgfpathlineto{\pgfqpoint{3.950713in}{0.487863in}}%
\pgfpathlineto{\pgfqpoint{3.948824in}{0.487153in}}%
\pgfpathlineto{\pgfqpoint{3.942331in}{0.491639in}}%
\pgfpathlineto{\pgfqpoint{3.937703in}{0.501458in}}%
\pgfpathlineto{\pgfqpoint{3.937393in}{0.504295in}}%
\pgfpathlineto{\pgfqpoint{3.938583in}{0.507691in}}%
\pgfpathlineto{\pgfqpoint{3.943598in}{0.502496in}}%
\pgfpathlineto{\pgfqpoint{3.945132in}{0.503498in}}%
\pgfpathlineto{\pgfqpoint{3.949386in}{0.502545in}}%
\pgfpathlineto{\pgfqpoint{3.953131in}{0.499862in}}%
\pgfpathclose%
\pgfusepath{fill}%
\end{pgfscope}%
\begin{pgfscope}%
\pgfpathrectangle{\pgfqpoint{3.525000in}{0.100000in}}{\pgfqpoint{2.862985in}{1.832778in}}%
\pgfusepath{clip}%
\pgfsetbuttcap%
\pgfsetmiterjoin%
\definecolor{currentfill}{rgb}{0.717186,0.879508,0.520185}%
\pgfsetfillcolor{currentfill}%
\pgfsetlinewidth{0.000000pt}%
\definecolor{currentstroke}{rgb}{0.000000,0.000000,0.000000}%
\pgfsetstrokecolor{currentstroke}%
\pgfsetstrokeopacity{0.000000}%
\pgfsetdash{}{0pt}%
\pgfpathmoveto{\pgfqpoint{3.876126in}{0.386109in}}%
\pgfpathlineto{\pgfqpoint{3.871562in}{0.385368in}}%
\pgfpathlineto{\pgfqpoint{3.868260in}{0.386895in}}%
\pgfpathlineto{\pgfqpoint{3.866816in}{0.389380in}}%
\pgfpathlineto{\pgfqpoint{3.868453in}{0.392945in}}%
\pgfpathlineto{\pgfqpoint{3.872096in}{0.392002in}}%
\pgfpathlineto{\pgfqpoint{3.878086in}{0.394117in}}%
\pgfpathlineto{\pgfqpoint{3.880184in}{0.391196in}}%
\pgfpathlineto{\pgfqpoint{3.882217in}{0.391735in}}%
\pgfpathlineto{\pgfqpoint{3.887141in}{0.389917in}}%
\pgfpathlineto{\pgfqpoint{3.889268in}{0.387608in}}%
\pgfpathlineto{\pgfqpoint{3.886672in}{0.381567in}}%
\pgfpathlineto{\pgfqpoint{3.884090in}{0.380286in}}%
\pgfpathlineto{\pgfqpoint{3.881769in}{0.380976in}}%
\pgfpathlineto{\pgfqpoint{3.876126in}{0.386109in}}%
\pgfpathclose%
\pgfusepath{fill}%
\end{pgfscope}%
\begin{pgfscope}%
\pgfpathrectangle{\pgfqpoint{3.525000in}{0.100000in}}{\pgfqpoint{2.862985in}{1.832778in}}%
\pgfusepath{clip}%
\pgfsetbuttcap%
\pgfsetmiterjoin%
\definecolor{currentfill}{rgb}{0.717186,0.879508,0.520185}%
\pgfsetfillcolor{currentfill}%
\pgfsetlinewidth{0.000000pt}%
\definecolor{currentstroke}{rgb}{0.000000,0.000000,0.000000}%
\pgfsetstrokecolor{currentstroke}%
\pgfsetstrokeopacity{0.000000}%
\pgfsetdash{}{0pt}%
\pgfpathmoveto{\pgfqpoint{3.835323in}{0.391598in}}%
\pgfpathlineto{\pgfqpoint{3.828857in}{0.394682in}}%
\pgfpathlineto{\pgfqpoint{3.822873in}{0.395600in}}%
\pgfpathlineto{\pgfqpoint{3.820987in}{0.396913in}}%
\pgfpathlineto{\pgfqpoint{3.823135in}{0.399499in}}%
\pgfpathlineto{\pgfqpoint{3.826946in}{0.396269in}}%
\pgfpathlineto{\pgfqpoint{3.830142in}{0.396714in}}%
\pgfpathlineto{\pgfqpoint{3.835341in}{0.398951in}}%
\pgfpathlineto{\pgfqpoint{3.835761in}{0.401832in}}%
\pgfpathlineto{\pgfqpoint{3.842146in}{0.400406in}}%
\pgfpathlineto{\pgfqpoint{3.840651in}{0.396069in}}%
\pgfpathlineto{\pgfqpoint{3.844719in}{0.395751in}}%
\pgfpathlineto{\pgfqpoint{3.841306in}{0.390592in}}%
\pgfpathlineto{\pgfqpoint{3.838001in}{0.392264in}}%
\pgfpathlineto{\pgfqpoint{3.835323in}{0.391598in}}%
\pgfpathclose%
\pgfusepath{fill}%
\end{pgfscope}%
\begin{pgfscope}%
\pgfpathrectangle{\pgfqpoint{3.525000in}{0.100000in}}{\pgfqpoint{2.862985in}{1.832778in}}%
\pgfusepath{clip}%
\pgfsetbuttcap%
\pgfsetmiterjoin%
\definecolor{currentfill}{rgb}{0.717186,0.879508,0.520185}%
\pgfsetfillcolor{currentfill}%
\pgfsetlinewidth{0.000000pt}%
\definecolor{currentstroke}{rgb}{0.000000,0.000000,0.000000}%
\pgfsetstrokecolor{currentstroke}%
\pgfsetstrokeopacity{0.000000}%
\pgfsetdash{}{0pt}%
\pgfpathmoveto{\pgfqpoint{3.823962in}{0.402261in}}%
\pgfpathlineto{\pgfqpoint{3.821586in}{0.400969in}}%
\pgfpathlineto{\pgfqpoint{3.815110in}{0.402955in}}%
\pgfpathlineto{\pgfqpoint{3.810271in}{0.401805in}}%
\pgfpathlineto{\pgfqpoint{3.807813in}{0.405073in}}%
\pgfpathlineto{\pgfqpoint{3.808885in}{0.406569in}}%
\pgfpathlineto{\pgfqpoint{3.812552in}{0.406819in}}%
\pgfpathlineto{\pgfqpoint{3.815357in}{0.405880in}}%
\pgfpathlineto{\pgfqpoint{3.818414in}{0.407367in}}%
\pgfpathlineto{\pgfqpoint{3.823100in}{0.405378in}}%
\pgfpathlineto{\pgfqpoint{3.823962in}{0.402261in}}%
\pgfpathclose%
\pgfusepath{fill}%
\end{pgfscope}%
\begin{pgfscope}%
\pgfpathrectangle{\pgfqpoint{3.525000in}{0.100000in}}{\pgfqpoint{2.862985in}{1.832778in}}%
\pgfusepath{clip}%
\pgfsetbuttcap%
\pgfsetmiterjoin%
\definecolor{currentfill}{rgb}{0.717186,0.879508,0.520185}%
\pgfsetfillcolor{currentfill}%
\pgfsetlinewidth{0.000000pt}%
\definecolor{currentstroke}{rgb}{0.000000,0.000000,0.000000}%
\pgfsetstrokecolor{currentstroke}%
\pgfsetstrokeopacity{0.000000}%
\pgfsetdash{}{0pt}%
\pgfpathmoveto{\pgfqpoint{3.747628in}{0.450357in}}%
\pgfpathlineto{\pgfqpoint{3.747456in}{0.447153in}}%
\pgfpathlineto{\pgfqpoint{3.742729in}{0.445139in}}%
\pgfpathlineto{\pgfqpoint{3.738412in}{0.450180in}}%
\pgfpathlineto{\pgfqpoint{3.743253in}{0.451860in}}%
\pgfpathlineto{\pgfqpoint{3.747628in}{0.450357in}}%
\pgfpathclose%
\pgfusepath{fill}%
\end{pgfscope}%
\begin{pgfscope}%
\pgfpathrectangle{\pgfqpoint{3.525000in}{0.100000in}}{\pgfqpoint{2.862985in}{1.832778in}}%
\pgfusepath{clip}%
\pgfsetbuttcap%
\pgfsetmiterjoin%
\definecolor{currentfill}{rgb}{0.717186,0.879508,0.520185}%
\pgfsetfillcolor{currentfill}%
\pgfsetlinewidth{0.000000pt}%
\definecolor{currentstroke}{rgb}{0.000000,0.000000,0.000000}%
\pgfsetstrokecolor{currentstroke}%
\pgfsetstrokeopacity{0.000000}%
\pgfsetdash{}{0pt}%
\pgfpathmoveto{\pgfqpoint{3.708638in}{0.468371in}}%
\pgfpathlineto{\pgfqpoint{3.709642in}{0.469764in}}%
\pgfpathlineto{\pgfqpoint{3.715440in}{0.469524in}}%
\pgfpathlineto{\pgfqpoint{3.716458in}{0.471597in}}%
\pgfpathlineto{\pgfqpoint{3.719002in}{0.469788in}}%
\pgfpathlineto{\pgfqpoint{3.717533in}{0.466308in}}%
\pgfpathlineto{\pgfqpoint{3.715372in}{0.466014in}}%
\pgfpathlineto{\pgfqpoint{3.710207in}{0.467106in}}%
\pgfpathlineto{\pgfqpoint{3.708638in}{0.468371in}}%
\pgfpathclose%
\pgfusepath{fill}%
\end{pgfscope}%
\begin{pgfscope}%
\pgfpathrectangle{\pgfqpoint{3.525000in}{0.100000in}}{\pgfqpoint{2.862985in}{1.832778in}}%
\pgfusepath{clip}%
\pgfsetbuttcap%
\pgfsetmiterjoin%
\definecolor{currentfill}{rgb}{0.717186,0.879508,0.520185}%
\pgfsetfillcolor{currentfill}%
\pgfsetlinewidth{0.000000pt}%
\definecolor{currentstroke}{rgb}{0.000000,0.000000,0.000000}%
\pgfsetstrokecolor{currentstroke}%
\pgfsetstrokeopacity{0.000000}%
\pgfsetdash{}{0pt}%
\pgfpathmoveto{\pgfqpoint{3.700415in}{0.479156in}}%
\pgfpathlineto{\pgfqpoint{3.700063in}{0.482294in}}%
\pgfpathlineto{\pgfqpoint{3.703088in}{0.482062in}}%
\pgfpathlineto{\pgfqpoint{3.702878in}{0.486470in}}%
\pgfpathlineto{\pgfqpoint{3.705921in}{0.484091in}}%
\pgfpathlineto{\pgfqpoint{3.704868in}{0.480653in}}%
\pgfpathlineto{\pgfqpoint{3.700415in}{0.479156in}}%
\pgfpathclose%
\pgfusepath{fill}%
\end{pgfscope}%
\begin{pgfscope}%
\pgfpathrectangle{\pgfqpoint{3.525000in}{0.100000in}}{\pgfqpoint{2.862985in}{1.832778in}}%
\pgfusepath{clip}%
\pgfsetbuttcap%
\pgfsetmiterjoin%
\definecolor{currentfill}{rgb}{0.717186,0.879508,0.520185}%
\pgfsetfillcolor{currentfill}%
\pgfsetlinewidth{0.000000pt}%
\definecolor{currentstroke}{rgb}{0.000000,0.000000,0.000000}%
\pgfsetstrokecolor{currentstroke}%
\pgfsetstrokeopacity{0.000000}%
\pgfsetdash{}{0pt}%
\pgfpathmoveto{\pgfqpoint{3.966980in}{0.591016in}}%
\pgfpathlineto{\pgfqpoint{3.962154in}{0.592058in}}%
\pgfpathlineto{\pgfqpoint{3.961118in}{0.595283in}}%
\pgfpathlineto{\pgfqpoint{3.962776in}{0.598507in}}%
\pgfpathlineto{\pgfqpoint{3.966491in}{0.600190in}}%
\pgfpathlineto{\pgfqpoint{3.970744in}{0.591916in}}%
\pgfpathlineto{\pgfqpoint{3.974601in}{0.591548in}}%
\pgfpathlineto{\pgfqpoint{3.977479in}{0.588995in}}%
\pgfpathlineto{\pgfqpoint{3.977422in}{0.585482in}}%
\pgfpathlineto{\pgfqpoint{3.975318in}{0.583797in}}%
\pgfpathlineto{\pgfqpoint{3.978740in}{0.573408in}}%
\pgfpathlineto{\pgfqpoint{3.982484in}{0.569595in}}%
\pgfpathlineto{\pgfqpoint{3.978656in}{0.568385in}}%
\pgfpathlineto{\pgfqpoint{3.975557in}{0.572663in}}%
\pgfpathlineto{\pgfqpoint{3.968823in}{0.571338in}}%
\pgfpathlineto{\pgfqpoint{3.970918in}{0.575572in}}%
\pgfpathlineto{\pgfqpoint{3.969943in}{0.577258in}}%
\pgfpathlineto{\pgfqpoint{3.970514in}{0.581671in}}%
\pgfpathlineto{\pgfqpoint{3.969180in}{0.589648in}}%
\pgfpathlineto{\pgfqpoint{3.966980in}{0.591016in}}%
\pgfpathclose%
\pgfusepath{fill}%
\end{pgfscope}%
\begin{pgfscope}%
\pgfpathrectangle{\pgfqpoint{3.525000in}{0.100000in}}{\pgfqpoint{2.862985in}{1.832778in}}%
\pgfusepath{clip}%
\pgfsetbuttcap%
\pgfsetmiterjoin%
\definecolor{currentfill}{rgb}{0.717186,0.879508,0.520185}%
\pgfsetfillcolor{currentfill}%
\pgfsetlinewidth{0.000000pt}%
\definecolor{currentstroke}{rgb}{0.000000,0.000000,0.000000}%
\pgfsetstrokecolor{currentstroke}%
\pgfsetstrokeopacity{0.000000}%
\pgfsetdash{}{0pt}%
\pgfpathmoveto{\pgfqpoint{4.136519in}{0.353334in}}%
\pgfpathlineto{\pgfqpoint{4.129523in}{0.353412in}}%
\pgfpathlineto{\pgfqpoint{4.130062in}{0.356924in}}%
\pgfpathlineto{\pgfqpoint{4.132862in}{0.358180in}}%
\pgfpathlineto{\pgfqpoint{4.136519in}{0.353334in}}%
\pgfpathclose%
\pgfusepath{fill}%
\end{pgfscope}%
\begin{pgfscope}%
\pgfpathrectangle{\pgfqpoint{3.525000in}{0.100000in}}{\pgfqpoint{2.862985in}{1.832778in}}%
\pgfusepath{clip}%
\pgfsetbuttcap%
\pgfsetmiterjoin%
\definecolor{currentfill}{rgb}{0.717186,0.879508,0.520185}%
\pgfsetfillcolor{currentfill}%
\pgfsetlinewidth{0.000000pt}%
\definecolor{currentstroke}{rgb}{0.000000,0.000000,0.000000}%
\pgfsetstrokecolor{currentstroke}%
\pgfsetstrokeopacity{0.000000}%
\pgfsetdash{}{0pt}%
\pgfpathmoveto{\pgfqpoint{4.126812in}{0.356514in}}%
\pgfpathlineto{\pgfqpoint{4.121800in}{0.355301in}}%
\pgfpathlineto{\pgfqpoint{4.117038in}{0.353224in}}%
\pgfpathlineto{\pgfqpoint{4.115567in}{0.355671in}}%
\pgfpathlineto{\pgfqpoint{4.123757in}{0.357328in}}%
\pgfpathlineto{\pgfqpoint{4.126812in}{0.356514in}}%
\pgfpathclose%
\pgfusepath{fill}%
\end{pgfscope}%
\begin{pgfscope}%
\pgfpathrectangle{\pgfqpoint{3.525000in}{0.100000in}}{\pgfqpoint{2.862985in}{1.832778in}}%
\pgfusepath{clip}%
\pgfsetbuttcap%
\pgfsetmiterjoin%
\definecolor{currentfill}{rgb}{0.717186,0.879508,0.520185}%
\pgfsetfillcolor{currentfill}%
\pgfsetlinewidth{0.000000pt}%
\definecolor{currentstroke}{rgb}{0.000000,0.000000,0.000000}%
\pgfsetstrokecolor{currentstroke}%
\pgfsetstrokeopacity{0.000000}%
\pgfsetdash{}{0pt}%
\pgfpathmoveto{\pgfqpoint{4.049990in}{0.354028in}}%
\pgfpathlineto{\pgfqpoint{4.050182in}{0.351599in}}%
\pgfpathlineto{\pgfqpoint{4.047294in}{0.350090in}}%
\pgfpathlineto{\pgfqpoint{4.039235in}{0.353552in}}%
\pgfpathlineto{\pgfqpoint{4.037986in}{0.352159in}}%
\pgfpathlineto{\pgfqpoint{4.036036in}{0.358747in}}%
\pgfpathlineto{\pgfqpoint{4.040074in}{0.359664in}}%
\pgfpathlineto{\pgfqpoint{4.041940in}{0.356957in}}%
\pgfpathlineto{\pgfqpoint{4.045992in}{0.360345in}}%
\pgfpathlineto{\pgfqpoint{4.049990in}{0.354028in}}%
\pgfpathclose%
\pgfusepath{fill}%
\end{pgfscope}%
\begin{pgfscope}%
\pgfpathrectangle{\pgfqpoint{3.525000in}{0.100000in}}{\pgfqpoint{2.862985in}{1.832778in}}%
\pgfusepath{clip}%
\pgfsetbuttcap%
\pgfsetmiterjoin%
\definecolor{currentfill}{rgb}{0.717186,0.879508,0.520185}%
\pgfsetfillcolor{currentfill}%
\pgfsetlinewidth{0.000000pt}%
\definecolor{currentstroke}{rgb}{0.000000,0.000000,0.000000}%
\pgfsetstrokecolor{currentstroke}%
\pgfsetstrokeopacity{0.000000}%
\pgfsetdash{}{0pt}%
\pgfpathmoveto{\pgfqpoint{4.231774in}{0.219412in}}%
\pgfpathlineto{\pgfqpoint{4.225017in}{0.216896in}}%
\pgfpathlineto{\pgfqpoint{4.221437in}{0.216745in}}%
\pgfpathlineto{\pgfqpoint{4.224350in}{0.222082in}}%
\pgfpathlineto{\pgfqpoint{4.226840in}{0.224073in}}%
\pgfpathlineto{\pgfqpoint{4.226760in}{0.229600in}}%
\pgfpathlineto{\pgfqpoint{4.229465in}{0.239578in}}%
\pgfpathlineto{\pgfqpoint{4.231957in}{0.241473in}}%
\pgfpathlineto{\pgfqpoint{4.237710in}{0.238729in}}%
\pgfpathlineto{\pgfqpoint{4.236870in}{0.237158in}}%
\pgfpathlineto{\pgfqpoint{4.237306in}{0.229604in}}%
\pgfpathlineto{\pgfqpoint{4.235440in}{0.227566in}}%
\pgfpathlineto{\pgfqpoint{4.234640in}{0.222201in}}%
\pgfpathlineto{\pgfqpoint{4.231774in}{0.219412in}}%
\pgfpathclose%
\pgfusepath{fill}%
\end{pgfscope}%
\begin{pgfscope}%
\pgfpathrectangle{\pgfqpoint{3.525000in}{0.100000in}}{\pgfqpoint{2.862985in}{1.832778in}}%
\pgfusepath{clip}%
\pgfsetbuttcap%
\pgfsetmiterjoin%
\definecolor{currentfill}{rgb}{0.717186,0.879508,0.520185}%
\pgfsetfillcolor{currentfill}%
\pgfsetlinewidth{0.000000pt}%
\definecolor{currentstroke}{rgb}{0.000000,0.000000,0.000000}%
\pgfsetstrokecolor{currentstroke}%
\pgfsetstrokeopacity{0.000000}%
\pgfsetdash{}{0pt}%
\pgfpathmoveto{\pgfqpoint{4.216809in}{0.244146in}}%
\pgfpathlineto{\pgfqpoint{4.220768in}{0.237310in}}%
\pgfpathlineto{\pgfqpoint{4.220034in}{0.235680in}}%
\pgfpathlineto{\pgfqpoint{4.225329in}{0.234022in}}%
\pgfpathlineto{\pgfqpoint{4.223809in}{0.227738in}}%
\pgfpathlineto{\pgfqpoint{4.221503in}{0.228029in}}%
\pgfpathlineto{\pgfqpoint{4.215445in}{0.238982in}}%
\pgfpathlineto{\pgfqpoint{4.216307in}{0.234052in}}%
\pgfpathlineto{\pgfqpoint{4.215275in}{0.231213in}}%
\pgfpathlineto{\pgfqpoint{4.212709in}{0.229950in}}%
\pgfpathlineto{\pgfqpoint{4.211251in}{0.231311in}}%
\pgfpathlineto{\pgfqpoint{4.210437in}{0.237623in}}%
\pgfpathlineto{\pgfqpoint{4.210917in}{0.241717in}}%
\pgfpathlineto{\pgfqpoint{4.209208in}{0.243504in}}%
\pgfpathlineto{\pgfqpoint{4.213722in}{0.248675in}}%
\pgfpathlineto{\pgfqpoint{4.216684in}{0.250209in}}%
\pgfpathlineto{\pgfqpoint{4.217670in}{0.248193in}}%
\pgfpathlineto{\pgfqpoint{4.220760in}{0.249753in}}%
\pgfpathlineto{\pgfqpoint{4.222999in}{0.245516in}}%
\pgfpathlineto{\pgfqpoint{4.227840in}{0.240126in}}%
\pgfpathlineto{\pgfqpoint{4.227179in}{0.237684in}}%
\pgfpathlineto{\pgfqpoint{4.224491in}{0.235024in}}%
\pgfpathlineto{\pgfqpoint{4.222005in}{0.236278in}}%
\pgfpathlineto{\pgfqpoint{4.216809in}{0.244146in}}%
\pgfpathclose%
\pgfusepath{fill}%
\end{pgfscope}%
\begin{pgfscope}%
\pgfpathrectangle{\pgfqpoint{3.525000in}{0.100000in}}{\pgfqpoint{2.862985in}{1.832778in}}%
\pgfusepath{clip}%
\pgfsetbuttcap%
\pgfsetmiterjoin%
\definecolor{currentfill}{rgb}{0.717186,0.879508,0.520185}%
\pgfsetfillcolor{currentfill}%
\pgfsetlinewidth{0.000000pt}%
\definecolor{currentstroke}{rgb}{0.000000,0.000000,0.000000}%
\pgfsetstrokecolor{currentstroke}%
\pgfsetstrokeopacity{0.000000}%
\pgfsetdash{}{0pt}%
\pgfpathmoveto{\pgfqpoint{4.020353in}{0.341071in}}%
\pgfpathlineto{\pgfqpoint{4.015449in}{0.340433in}}%
\pgfpathlineto{\pgfqpoint{4.008542in}{0.337778in}}%
\pgfpathlineto{\pgfqpoint{4.006931in}{0.339221in}}%
\pgfpathlineto{\pgfqpoint{4.014146in}{0.340851in}}%
\pgfpathlineto{\pgfqpoint{4.012212in}{0.342273in}}%
\pgfpathlineto{\pgfqpoint{4.007739in}{0.343315in}}%
\pgfpathlineto{\pgfqpoint{4.006544in}{0.346649in}}%
\pgfpathlineto{\pgfqpoint{4.008795in}{0.349521in}}%
\pgfpathlineto{\pgfqpoint{4.010690in}{0.356112in}}%
\pgfpathlineto{\pgfqpoint{4.017614in}{0.357340in}}%
\pgfpathlineto{\pgfqpoint{4.020966in}{0.354492in}}%
\pgfpathlineto{\pgfqpoint{4.024115in}{0.357248in}}%
\pgfpathlineto{\pgfqpoint{4.028064in}{0.356572in}}%
\pgfpathlineto{\pgfqpoint{4.030902in}{0.351482in}}%
\pgfpathlineto{\pgfqpoint{4.033305in}{0.356006in}}%
\pgfpathlineto{\pgfqpoint{4.040126in}{0.345423in}}%
\pgfpathlineto{\pgfqpoint{4.037662in}{0.344577in}}%
\pgfpathlineto{\pgfqpoint{4.036150in}{0.342283in}}%
\pgfpathlineto{\pgfqpoint{4.039007in}{0.340190in}}%
\pgfpathlineto{\pgfqpoint{4.034576in}{0.338227in}}%
\pgfpathlineto{\pgfqpoint{4.031364in}{0.339308in}}%
\pgfpathlineto{\pgfqpoint{4.025296in}{0.342890in}}%
\pgfpathlineto{\pgfqpoint{4.020353in}{0.341071in}}%
\pgfpathclose%
\pgfusepath{fill}%
\end{pgfscope}%
\begin{pgfscope}%
\pgfpathrectangle{\pgfqpoint{3.525000in}{0.100000in}}{\pgfqpoint{2.862985in}{1.832778in}}%
\pgfusepath{clip}%
\pgfsetbuttcap%
\pgfsetmiterjoin%
\definecolor{currentfill}{rgb}{0.717186,0.879508,0.520185}%
\pgfsetfillcolor{currentfill}%
\pgfsetlinewidth{0.000000pt}%
\definecolor{currentstroke}{rgb}{0.000000,0.000000,0.000000}%
\pgfsetstrokecolor{currentstroke}%
\pgfsetstrokeopacity{0.000000}%
\pgfsetdash{}{0pt}%
\pgfpathmoveto{\pgfqpoint{4.212691in}{0.199220in}}%
\pgfpathlineto{\pgfqpoint{4.211719in}{0.208994in}}%
\pgfpathlineto{\pgfqpoint{4.214340in}{0.212583in}}%
\pgfpathlineto{\pgfqpoint{4.211538in}{0.212832in}}%
\pgfpathlineto{\pgfqpoint{4.212562in}{0.219812in}}%
\pgfpathlineto{\pgfqpoint{4.213770in}{0.223955in}}%
\pgfpathlineto{\pgfqpoint{4.212720in}{0.229522in}}%
\pgfpathlineto{\pgfqpoint{4.215287in}{0.230600in}}%
\pgfpathlineto{\pgfqpoint{4.216124in}{0.232039in}}%
\pgfpathlineto{\pgfqpoint{4.218653in}{0.231493in}}%
\pgfpathlineto{\pgfqpoint{4.220842in}{0.227033in}}%
\pgfpathlineto{\pgfqpoint{4.221249in}{0.222947in}}%
\pgfpathlineto{\pgfqpoint{4.218571in}{0.210709in}}%
\pgfpathlineto{\pgfqpoint{4.212691in}{0.199220in}}%
\pgfpathclose%
\pgfusepath{fill}%
\end{pgfscope}%
\begin{pgfscope}%
\pgfpathrectangle{\pgfqpoint{3.525000in}{0.100000in}}{\pgfqpoint{2.862985in}{1.832778in}}%
\pgfusepath{clip}%
\pgfsetbuttcap%
\pgfsetmiterjoin%
\definecolor{currentfill}{rgb}{0.717186,0.879508,0.520185}%
\pgfsetfillcolor{currentfill}%
\pgfsetlinewidth{0.000000pt}%
\definecolor{currentstroke}{rgb}{0.000000,0.000000,0.000000}%
\pgfsetstrokecolor{currentstroke}%
\pgfsetstrokeopacity{0.000000}%
\pgfsetdash{}{0pt}%
\pgfpathmoveto{\pgfqpoint{4.212110in}{0.229056in}}%
\pgfpathlineto{\pgfqpoint{4.211975in}{0.224290in}}%
\pgfpathlineto{\pgfqpoint{4.209778in}{0.222127in}}%
\pgfpathlineto{\pgfqpoint{4.207238in}{0.222926in}}%
\pgfpathlineto{\pgfqpoint{4.210448in}{0.229848in}}%
\pgfpathlineto{\pgfqpoint{4.212110in}{0.229056in}}%
\pgfpathclose%
\pgfusepath{fill}%
\end{pgfscope}%
\begin{pgfscope}%
\pgfpathrectangle{\pgfqpoint{3.525000in}{0.100000in}}{\pgfqpoint{2.862985in}{1.832778in}}%
\pgfusepath{clip}%
\pgfsetbuttcap%
\pgfsetmiterjoin%
\definecolor{currentfill}{rgb}{0.717186,0.879508,0.520185}%
\pgfsetfillcolor{currentfill}%
\pgfsetlinewidth{0.000000pt}%
\definecolor{currentstroke}{rgb}{0.000000,0.000000,0.000000}%
\pgfsetstrokecolor{currentstroke}%
\pgfsetstrokeopacity{0.000000}%
\pgfsetdash{}{0pt}%
\pgfpathmoveto{\pgfqpoint{4.239076in}{0.207898in}}%
\pgfpathlineto{\pgfqpoint{4.237202in}{0.199923in}}%
\pgfpathlineto{\pgfqpoint{4.232835in}{0.197402in}}%
\pgfpathlineto{\pgfqpoint{4.227107in}{0.199724in}}%
\pgfpathlineto{\pgfqpoint{4.228685in}{0.202824in}}%
\pgfpathlineto{\pgfqpoint{4.228653in}{0.204970in}}%
\pgfpathlineto{\pgfqpoint{4.230783in}{0.208616in}}%
\pgfpathlineto{\pgfqpoint{4.229286in}{0.207882in}}%
\pgfpathlineto{\pgfqpoint{4.230377in}{0.209663in}}%
\pgfpathlineto{\pgfqpoint{4.228304in}{0.213709in}}%
\pgfpathlineto{\pgfqpoint{4.230657in}{0.214455in}}%
\pgfpathlineto{\pgfqpoint{4.236104in}{0.209721in}}%
\pgfpathlineto{\pgfqpoint{4.239076in}{0.207898in}}%
\pgfpathclose%
\pgfusepath{fill}%
\end{pgfscope}%
\begin{pgfscope}%
\pgfpathrectangle{\pgfqpoint{3.525000in}{0.100000in}}{\pgfqpoint{2.862985in}{1.832778in}}%
\pgfusepath{clip}%
\pgfsetbuttcap%
\pgfsetmiterjoin%
\definecolor{currentfill}{rgb}{0.717186,0.879508,0.520185}%
\pgfsetfillcolor{currentfill}%
\pgfsetlinewidth{0.000000pt}%
\definecolor{currentstroke}{rgb}{0.000000,0.000000,0.000000}%
\pgfsetstrokecolor{currentstroke}%
\pgfsetstrokeopacity{0.000000}%
\pgfsetdash{}{0pt}%
\pgfpathmoveto{\pgfqpoint{4.219962in}{0.193796in}}%
\pgfpathlineto{\pgfqpoint{4.217492in}{0.192836in}}%
\pgfpathlineto{\pgfqpoint{4.218902in}{0.201252in}}%
\pgfpathlineto{\pgfqpoint{4.221866in}{0.203216in}}%
\pgfpathlineto{\pgfqpoint{4.221144in}{0.209521in}}%
\pgfpathlineto{\pgfqpoint{4.222348in}{0.212317in}}%
\pgfpathlineto{\pgfqpoint{4.224731in}{0.213365in}}%
\pgfpathlineto{\pgfqpoint{4.227232in}{0.210622in}}%
\pgfpathlineto{\pgfqpoint{4.229379in}{0.206901in}}%
\pgfpathlineto{\pgfqpoint{4.222478in}{0.199610in}}%
\pgfpathlineto{\pgfqpoint{4.219962in}{0.193796in}}%
\pgfpathclose%
\pgfusepath{fill}%
\end{pgfscope}%
\begin{pgfscope}%
\pgfpathrectangle{\pgfqpoint{3.525000in}{0.100000in}}{\pgfqpoint{2.862985in}{1.832778in}}%
\pgfusepath{clip}%
\pgfsetbuttcap%
\pgfsetmiterjoin%
\definecolor{currentfill}{rgb}{0.717186,0.879508,0.520185}%
\pgfsetfillcolor{currentfill}%
\pgfsetlinewidth{0.000000pt}%
\definecolor{currentstroke}{rgb}{0.000000,0.000000,0.000000}%
\pgfsetstrokecolor{currentstroke}%
\pgfsetstrokeopacity{0.000000}%
\pgfsetdash{}{0pt}%
\pgfpathmoveto{\pgfqpoint{4.240172in}{0.202568in}}%
\pgfpathlineto{\pgfqpoint{4.241382in}{0.196760in}}%
\pgfpathlineto{\pgfqpoint{4.235974in}{0.197425in}}%
\pgfpathlineto{\pgfqpoint{4.240172in}{0.202568in}}%
\pgfpathclose%
\pgfusepath{fill}%
\end{pgfscope}%
\begin{pgfscope}%
\pgfpathrectangle{\pgfqpoint{3.525000in}{0.100000in}}{\pgfqpoint{2.862985in}{1.832778in}}%
\pgfusepath{clip}%
\pgfsetbuttcap%
\pgfsetmiterjoin%
\definecolor{currentfill}{rgb}{0.717186,0.879508,0.520185}%
\pgfsetfillcolor{currentfill}%
\pgfsetlinewidth{0.000000pt}%
\definecolor{currentstroke}{rgb}{0.000000,0.000000,0.000000}%
\pgfsetstrokecolor{currentstroke}%
\pgfsetstrokeopacity{0.000000}%
\pgfsetdash{}{0pt}%
\pgfpathmoveto{\pgfqpoint{4.242752in}{0.193694in}}%
\pgfpathlineto{\pgfqpoint{4.243638in}{0.190111in}}%
\pgfpathlineto{\pgfqpoint{4.245335in}{0.188965in}}%
\pgfpathlineto{\pgfqpoint{4.245026in}{0.185593in}}%
\pgfpathlineto{\pgfqpoint{4.242656in}{0.184459in}}%
\pgfpathlineto{\pgfqpoint{4.240732in}{0.189493in}}%
\pgfpathlineto{\pgfqpoint{4.242752in}{0.193694in}}%
\pgfpathclose%
\pgfusepath{fill}%
\end{pgfscope}%
\begin{pgfscope}%
\pgfpathrectangle{\pgfqpoint{3.525000in}{0.100000in}}{\pgfqpoint{2.862985in}{1.832778in}}%
\pgfusepath{clip}%
\pgfsetbuttcap%
\pgfsetmiterjoin%
\definecolor{currentfill}{rgb}{0.717186,0.879508,0.520185}%
\pgfsetfillcolor{currentfill}%
\pgfsetlinewidth{0.000000pt}%
\definecolor{currentstroke}{rgb}{0.000000,0.000000,0.000000}%
\pgfsetstrokecolor{currentstroke}%
\pgfsetstrokeopacity{0.000000}%
\pgfsetdash{}{0pt}%
\pgfpathmoveto{\pgfqpoint{4.238239in}{0.195068in}}%
\pgfpathlineto{\pgfqpoint{4.236768in}{0.190907in}}%
\pgfpathlineto{\pgfqpoint{4.234484in}{0.191060in}}%
\pgfpathlineto{\pgfqpoint{4.233043in}{0.194492in}}%
\pgfpathlineto{\pgfqpoint{4.235390in}{0.196117in}}%
\pgfpathlineto{\pgfqpoint{4.238239in}{0.195068in}}%
\pgfpathclose%
\pgfusepath{fill}%
\end{pgfscope}%
\begin{pgfscope}%
\pgfpathrectangle{\pgfqpoint{3.525000in}{0.100000in}}{\pgfqpoint{2.862985in}{1.832778in}}%
\pgfusepath{clip}%
\pgfsetbuttcap%
\pgfsetmiterjoin%
\definecolor{currentfill}{rgb}{0.717186,0.879508,0.520185}%
\pgfsetfillcolor{currentfill}%
\pgfsetlinewidth{0.000000pt}%
\definecolor{currentstroke}{rgb}{0.000000,0.000000,0.000000}%
\pgfsetstrokecolor{currentstroke}%
\pgfsetstrokeopacity{0.000000}%
\pgfsetdash{}{0pt}%
\pgfpathmoveto{\pgfqpoint{4.232076in}{0.173613in}}%
\pgfpathlineto{\pgfqpoint{4.234114in}{0.171384in}}%
\pgfpathlineto{\pgfqpoint{4.234563in}{0.166686in}}%
\pgfpathlineto{\pgfqpoint{4.235636in}{0.165321in}}%
\pgfpathlineto{\pgfqpoint{4.233884in}{0.161070in}}%
\pgfpathlineto{\pgfqpoint{4.231578in}{0.158676in}}%
\pgfpathlineto{\pgfqpoint{4.230352in}{0.154236in}}%
\pgfpathlineto{\pgfqpoint{4.227832in}{0.161156in}}%
\pgfpathlineto{\pgfqpoint{4.227414in}{0.167439in}}%
\pgfpathlineto{\pgfqpoint{4.226004in}{0.170567in}}%
\pgfpathlineto{\pgfqpoint{4.221147in}{0.173102in}}%
\pgfpathlineto{\pgfqpoint{4.225703in}{0.178390in}}%
\pgfpathlineto{\pgfqpoint{4.222666in}{0.180650in}}%
\pgfpathlineto{\pgfqpoint{4.225341in}{0.182902in}}%
\pgfpathlineto{\pgfqpoint{4.228893in}{0.191549in}}%
\pgfpathlineto{\pgfqpoint{4.225212in}{0.194597in}}%
\pgfpathlineto{\pgfqpoint{4.226630in}{0.197524in}}%
\pgfpathlineto{\pgfqpoint{4.231388in}{0.194928in}}%
\pgfpathlineto{\pgfqpoint{4.231870in}{0.190831in}}%
\pgfpathlineto{\pgfqpoint{4.230070in}{0.188498in}}%
\pgfpathlineto{\pgfqpoint{4.233638in}{0.185341in}}%
\pgfpathlineto{\pgfqpoint{4.234761in}{0.180424in}}%
\pgfpathlineto{\pgfqpoint{4.234609in}{0.173528in}}%
\pgfpathlineto{\pgfqpoint{4.232076in}{0.173613in}}%
\pgfpathclose%
\pgfusepath{fill}%
\end{pgfscope}%
\begin{pgfscope}%
\pgfpathrectangle{\pgfqpoint{3.525000in}{0.100000in}}{\pgfqpoint{2.862985in}{1.832778in}}%
\pgfusepath{clip}%
\pgfsetbuttcap%
\pgfsetmiterjoin%
\definecolor{currentfill}{rgb}{0.717186,0.879508,0.520185}%
\pgfsetfillcolor{currentfill}%
\pgfsetlinewidth{0.000000pt}%
\definecolor{currentstroke}{rgb}{0.000000,0.000000,0.000000}%
\pgfsetstrokecolor{currentstroke}%
\pgfsetstrokeopacity{0.000000}%
\pgfsetdash{}{0pt}%
\pgfpathmoveto{\pgfqpoint{4.240789in}{0.190924in}}%
\pgfpathlineto{\pgfqpoint{4.239679in}{0.188532in}}%
\pgfpathlineto{\pgfqpoint{4.240722in}{0.184648in}}%
\pgfpathlineto{\pgfqpoint{4.238053in}{0.184321in}}%
\pgfpathlineto{\pgfqpoint{4.235255in}{0.186119in}}%
\pgfpathlineto{\pgfqpoint{4.236163in}{0.190080in}}%
\pgfpathlineto{\pgfqpoint{4.240789in}{0.190924in}}%
\pgfpathclose%
\pgfusepath{fill}%
\end{pgfscope}%
\begin{pgfscope}%
\pgfpathrectangle{\pgfqpoint{3.525000in}{0.100000in}}{\pgfqpoint{2.862985in}{1.832778in}}%
\pgfusepath{clip}%
\pgfsetbuttcap%
\pgfsetmiterjoin%
\definecolor{currentfill}{rgb}{0.717186,0.879508,0.520185}%
\pgfsetfillcolor{currentfill}%
\pgfsetlinewidth{0.000000pt}%
\definecolor{currentstroke}{rgb}{0.000000,0.000000,0.000000}%
\pgfsetstrokecolor{currentstroke}%
\pgfsetstrokeopacity{0.000000}%
\pgfsetdash{}{0pt}%
\pgfpathmoveto{\pgfqpoint{4.228351in}{0.191195in}}%
\pgfpathlineto{\pgfqpoint{4.222461in}{0.188553in}}%
\pgfpathlineto{\pgfqpoint{4.220459in}{0.189503in}}%
\pgfpathlineto{\pgfqpoint{4.224476in}{0.192989in}}%
\pgfpathlineto{\pgfqpoint{4.228351in}{0.191195in}}%
\pgfpathclose%
\pgfusepath{fill}%
\end{pgfscope}%
\begin{pgfscope}%
\pgfpathrectangle{\pgfqpoint{3.525000in}{0.100000in}}{\pgfqpoint{2.862985in}{1.832778in}}%
\pgfusepath{clip}%
\pgfsetbuttcap%
\pgfsetmiterjoin%
\definecolor{currentfill}{rgb}{0.717186,0.879508,0.520185}%
\pgfsetfillcolor{currentfill}%
\pgfsetlinewidth{0.000000pt}%
\definecolor{currentstroke}{rgb}{0.000000,0.000000,0.000000}%
\pgfsetstrokecolor{currentstroke}%
\pgfsetstrokeopacity{0.000000}%
\pgfsetdash{}{0pt}%
\pgfpathmoveto{\pgfqpoint{4.240854in}{0.165422in}}%
\pgfpathlineto{\pgfqpoint{4.239449in}{0.169400in}}%
\pgfpathlineto{\pgfqpoint{4.242419in}{0.170378in}}%
\pgfpathlineto{\pgfqpoint{4.243319in}{0.174668in}}%
\pgfpathlineto{\pgfqpoint{4.246538in}{0.174673in}}%
\pgfpathlineto{\pgfqpoint{4.246524in}{0.177660in}}%
\pgfpathlineto{\pgfqpoint{4.250997in}{0.176953in}}%
\pgfpathlineto{\pgfqpoint{4.251760in}{0.168856in}}%
\pgfpathlineto{\pgfqpoint{4.249165in}{0.163694in}}%
\pgfpathlineto{\pgfqpoint{4.245315in}{0.160469in}}%
\pgfpathlineto{\pgfqpoint{4.240854in}{0.165422in}}%
\pgfpathclose%
\pgfusepath{fill}%
\end{pgfscope}%
\begin{pgfscope}%
\pgfpathrectangle{\pgfqpoint{3.525000in}{0.100000in}}{\pgfqpoint{2.862985in}{1.832778in}}%
\pgfusepath{clip}%
\pgfsetbuttcap%
\pgfsetmiterjoin%
\definecolor{currentfill}{rgb}{0.717186,0.879508,0.520185}%
\pgfsetfillcolor{currentfill}%
\pgfsetlinewidth{0.000000pt}%
\definecolor{currentstroke}{rgb}{0.000000,0.000000,0.000000}%
\pgfsetstrokecolor{currentstroke}%
\pgfsetstrokeopacity{0.000000}%
\pgfsetdash{}{0pt}%
\pgfpathmoveto{\pgfqpoint{4.239114in}{0.168482in}}%
\pgfpathlineto{\pgfqpoint{4.240383in}{0.164715in}}%
\pgfpathlineto{\pgfqpoint{4.237255in}{0.164257in}}%
\pgfpathlineto{\pgfqpoint{4.239114in}{0.168482in}}%
\pgfpathclose%
\pgfusepath{fill}%
\end{pgfscope}%
\begin{pgfscope}%
\pgfpathrectangle{\pgfqpoint{3.525000in}{0.100000in}}{\pgfqpoint{2.862985in}{1.832778in}}%
\pgfusepath{clip}%
\pgfsetbuttcap%
\pgfsetmiterjoin%
\definecolor{currentfill}{rgb}{0.717186,0.879508,0.520185}%
\pgfsetfillcolor{currentfill}%
\pgfsetlinewidth{0.000000pt}%
\definecolor{currentstroke}{rgb}{0.000000,0.000000,0.000000}%
\pgfsetstrokecolor{currentstroke}%
\pgfsetstrokeopacity{0.000000}%
\pgfsetdash{}{0pt}%
\pgfpathmoveto{\pgfqpoint{4.241751in}{0.162942in}}%
\pgfpathlineto{\pgfqpoint{4.241315in}{0.158141in}}%
\pgfpathlineto{\pgfqpoint{4.238015in}{0.158538in}}%
\pgfpathlineto{\pgfqpoint{4.239083in}{0.160858in}}%
\pgfpathlineto{\pgfqpoint{4.241751in}{0.162942in}}%
\pgfpathclose%
\pgfusepath{fill}%
\end{pgfscope}%
\begin{pgfscope}%
\definecolor{textcolor}{rgb}{0.000000,0.000000,0.000000}%
\pgfsetstrokecolor{textcolor}%
\pgfsetfillcolor{textcolor}%
\pgftext[x=4.670194in,y=1.896122in,left,base]{\color{textcolor}\setmainfont{Lato}\rmfamily\fontsize{9.000000}{10.800000}\selectfont November 2022}%
\end{pgfscope}%
\begin{pgfscope}%
\pgfpathrectangle{\pgfqpoint{2.893540in}{0.169444in}}{\pgfqpoint{0.858895in}{1.282944in}}%
\pgfusepath{clip}%
\pgfsetbuttcap%
\pgfsetmiterjoin%
\definecolor{currentfill}{rgb}{0.152941,0.392157,0.098039}%
\pgfsetfillcolor{currentfill}%
\pgfsetlinewidth{0.000000pt}%
\definecolor{currentstroke}{rgb}{0.000000,0.000000,0.000000}%
\pgfsetstrokecolor{currentstroke}%
\pgfsetstrokeopacity{0.000000}%
\pgfsetdash{}{0pt}%
\pgfpathmoveto{\pgfqpoint{3.142620in}{0.169444in}}%
\pgfpathlineto{\pgfqpoint{3.108264in}{0.169444in}}%
\pgfpathlineto{\pgfqpoint{3.108264in}{0.192354in}}%
\pgfpathlineto{\pgfqpoint{3.142620in}{0.192354in}}%
\pgfpathlineto{\pgfqpoint{3.142620in}{0.169444in}}%
\pgfpathclose%
\pgfusepath{fill}%
\end{pgfscope}%
\begin{pgfscope}%
\pgfpathrectangle{\pgfqpoint{2.893540in}{0.169444in}}{\pgfqpoint{0.858895in}{1.282944in}}%
\pgfusepath{clip}%
\pgfsetbuttcap%
\pgfsetmiterjoin%
\definecolor{currentfill}{rgb}{0.176317,0.420454,0.102960}%
\pgfsetfillcolor{currentfill}%
\pgfsetlinewidth{0.000000pt}%
\definecolor{currentstroke}{rgb}{0.000000,0.000000,0.000000}%
\pgfsetstrokecolor{currentstroke}%
\pgfsetstrokeopacity{0.000000}%
\pgfsetdash{}{0pt}%
\pgfpathmoveto{\pgfqpoint{3.142620in}{0.192771in}}%
\pgfpathlineto{\pgfqpoint{3.108264in}{0.192771in}}%
\pgfpathlineto{\pgfqpoint{3.108264in}{0.215680in}}%
\pgfpathlineto{\pgfqpoint{3.142620in}{0.215680in}}%
\pgfpathlineto{\pgfqpoint{3.142620in}{0.192771in}}%
\pgfpathclose%
\pgfusepath{fill}%
\end{pgfscope}%
\begin{pgfscope}%
\pgfpathrectangle{\pgfqpoint{2.893540in}{0.169444in}}{\pgfqpoint{0.858895in}{1.282944in}}%
\pgfusepath{clip}%
\pgfsetbuttcap%
\pgfsetmiterjoin%
\definecolor{currentfill}{rgb}{0.205536,0.455825,0.109112}%
\pgfsetfillcolor{currentfill}%
\pgfsetlinewidth{0.000000pt}%
\definecolor{currentstroke}{rgb}{0.000000,0.000000,0.000000}%
\pgfsetstrokecolor{currentstroke}%
\pgfsetstrokeopacity{0.000000}%
\pgfsetdash{}{0pt}%
\pgfpathmoveto{\pgfqpoint{3.142620in}{0.216097in}}%
\pgfpathlineto{\pgfqpoint{3.108264in}{0.216097in}}%
\pgfpathlineto{\pgfqpoint{3.108264in}{0.239007in}}%
\pgfpathlineto{\pgfqpoint{3.142620in}{0.239007in}}%
\pgfpathlineto{\pgfqpoint{3.142620in}{0.216097in}}%
\pgfpathclose%
\pgfusepath{fill}%
\end{pgfscope}%
\begin{pgfscope}%
\pgfpathrectangle{\pgfqpoint{2.893540in}{0.169444in}}{\pgfqpoint{0.858895in}{1.282944in}}%
\pgfusepath{clip}%
\pgfsetbuttcap%
\pgfsetmiterjoin%
\definecolor{currentfill}{rgb}{0.228912,0.484121,0.114033}%
\pgfsetfillcolor{currentfill}%
\pgfsetlinewidth{0.000000pt}%
\definecolor{currentstroke}{rgb}{0.000000,0.000000,0.000000}%
\pgfsetstrokecolor{currentstroke}%
\pgfsetstrokeopacity{0.000000}%
\pgfsetdash{}{0pt}%
\pgfpathmoveto{\pgfqpoint{3.142620in}{0.239423in}}%
\pgfpathlineto{\pgfqpoint{3.108264in}{0.239423in}}%
\pgfpathlineto{\pgfqpoint{3.108264in}{0.262333in}}%
\pgfpathlineto{\pgfqpoint{3.142620in}{0.262333in}}%
\pgfpathlineto{\pgfqpoint{3.142620in}{0.239423in}}%
\pgfpathclose%
\pgfusepath{fill}%
\end{pgfscope}%
\begin{pgfscope}%
\pgfpathrectangle{\pgfqpoint{2.893540in}{0.169444in}}{\pgfqpoint{0.858895in}{1.282944in}}%
\pgfusepath{clip}%
\pgfsetbuttcap%
\pgfsetmiterjoin%
\definecolor{currentfill}{rgb}{0.258131,0.519493,0.120185}%
\pgfsetfillcolor{currentfill}%
\pgfsetlinewidth{0.000000pt}%
\definecolor{currentstroke}{rgb}{0.000000,0.000000,0.000000}%
\pgfsetstrokecolor{currentstroke}%
\pgfsetstrokeopacity{0.000000}%
\pgfsetdash{}{0pt}%
\pgfpathmoveto{\pgfqpoint{3.142620in}{0.262749in}}%
\pgfpathlineto{\pgfqpoint{3.108264in}{0.262749in}}%
\pgfpathlineto{\pgfqpoint{3.108264in}{0.285659in}}%
\pgfpathlineto{\pgfqpoint{3.142620in}{0.285659in}}%
\pgfpathlineto{\pgfqpoint{3.142620in}{0.262749in}}%
\pgfpathclose%
\pgfusepath{fill}%
\end{pgfscope}%
\begin{pgfscope}%
\pgfpathrectangle{\pgfqpoint{2.893540in}{0.169444in}}{\pgfqpoint{0.858895in}{1.282944in}}%
\pgfusepath{clip}%
\pgfsetbuttcap%
\pgfsetmiterjoin%
\definecolor{currentfill}{rgb}{0.287351,0.554864,0.126336}%
\pgfsetfillcolor{currentfill}%
\pgfsetlinewidth{0.000000pt}%
\definecolor{currentstroke}{rgb}{0.000000,0.000000,0.000000}%
\pgfsetstrokecolor{currentstroke}%
\pgfsetstrokeopacity{0.000000}%
\pgfsetdash{}{0pt}%
\pgfpathmoveto{\pgfqpoint{3.142620in}{0.286076in}}%
\pgfpathlineto{\pgfqpoint{3.108264in}{0.286076in}}%
\pgfpathlineto{\pgfqpoint{3.108264in}{0.308985in}}%
\pgfpathlineto{\pgfqpoint{3.142620in}{0.308985in}}%
\pgfpathlineto{\pgfqpoint{3.142620in}{0.286076in}}%
\pgfpathclose%
\pgfusepath{fill}%
\end{pgfscope}%
\begin{pgfscope}%
\pgfpathrectangle{\pgfqpoint{2.893540in}{0.169444in}}{\pgfqpoint{0.858895in}{1.282944in}}%
\pgfusepath{clip}%
\pgfsetbuttcap%
\pgfsetmiterjoin%
\definecolor{currentfill}{rgb}{0.313495,0.582238,0.136794}%
\pgfsetfillcolor{currentfill}%
\pgfsetlinewidth{0.000000pt}%
\definecolor{currentstroke}{rgb}{0.000000,0.000000,0.000000}%
\pgfsetstrokecolor{currentstroke}%
\pgfsetstrokeopacity{0.000000}%
\pgfsetdash{}{0pt}%
\pgfpathmoveto{\pgfqpoint{3.142620in}{0.309402in}}%
\pgfpathlineto{\pgfqpoint{3.108264in}{0.309402in}}%
\pgfpathlineto{\pgfqpoint{3.108264in}{0.332312in}}%
\pgfpathlineto{\pgfqpoint{3.142620in}{0.332312in}}%
\pgfpathlineto{\pgfqpoint{3.142620in}{0.309402in}}%
\pgfpathclose%
\pgfusepath{fill}%
\end{pgfscope}%
\begin{pgfscope}%
\pgfpathrectangle{\pgfqpoint{2.893540in}{0.169444in}}{\pgfqpoint{0.858895in}{1.282944in}}%
\pgfusepath{clip}%
\pgfsetbuttcap%
\pgfsetmiterjoin%
\definecolor{currentfill}{rgb}{0.351942,0.614533,0.161399}%
\pgfsetfillcolor{currentfill}%
\pgfsetlinewidth{0.000000pt}%
\definecolor{currentstroke}{rgb}{0.000000,0.000000,0.000000}%
\pgfsetstrokecolor{currentstroke}%
\pgfsetstrokeopacity{0.000000}%
\pgfsetdash{}{0pt}%
\pgfpathmoveto{\pgfqpoint{3.142620in}{0.332728in}}%
\pgfpathlineto{\pgfqpoint{3.108264in}{0.332728in}}%
\pgfpathlineto{\pgfqpoint{3.108264in}{0.355638in}}%
\pgfpathlineto{\pgfqpoint{3.142620in}{0.355638in}}%
\pgfpathlineto{\pgfqpoint{3.142620in}{0.332728in}}%
\pgfpathclose%
\pgfusepath{fill}%
\end{pgfscope}%
\begin{pgfscope}%
\pgfpathrectangle{\pgfqpoint{2.893540in}{0.169444in}}{\pgfqpoint{0.858895in}{1.282944in}}%
\pgfusepath{clip}%
\pgfsetbuttcap%
\pgfsetmiterjoin%
\definecolor{currentfill}{rgb}{0.390388,0.646828,0.186005}%
\pgfsetfillcolor{currentfill}%
\pgfsetlinewidth{0.000000pt}%
\definecolor{currentstroke}{rgb}{0.000000,0.000000,0.000000}%
\pgfsetstrokecolor{currentstroke}%
\pgfsetstrokeopacity{0.000000}%
\pgfsetdash{}{0pt}%
\pgfpathmoveto{\pgfqpoint{3.142620in}{0.356055in}}%
\pgfpathlineto{\pgfqpoint{3.108264in}{0.356055in}}%
\pgfpathlineto{\pgfqpoint{3.108264in}{0.378964in}}%
\pgfpathlineto{\pgfqpoint{3.142620in}{0.378964in}}%
\pgfpathlineto{\pgfqpoint{3.142620in}{0.356055in}}%
\pgfpathclose%
\pgfusepath{fill}%
\end{pgfscope}%
\begin{pgfscope}%
\pgfpathrectangle{\pgfqpoint{2.893540in}{0.169444in}}{\pgfqpoint{0.858895in}{1.282944in}}%
\pgfusepath{clip}%
\pgfsetbuttcap%
\pgfsetmiterjoin%
\definecolor{currentfill}{rgb}{0.421146,0.672664,0.205690}%
\pgfsetfillcolor{currentfill}%
\pgfsetlinewidth{0.000000pt}%
\definecolor{currentstroke}{rgb}{0.000000,0.000000,0.000000}%
\pgfsetstrokecolor{currentstroke}%
\pgfsetstrokeopacity{0.000000}%
\pgfsetdash{}{0pt}%
\pgfpathmoveto{\pgfqpoint{3.142620in}{0.379381in}}%
\pgfpathlineto{\pgfqpoint{3.108264in}{0.379381in}}%
\pgfpathlineto{\pgfqpoint{3.108264in}{0.402291in}}%
\pgfpathlineto{\pgfqpoint{3.142620in}{0.402291in}}%
\pgfpathlineto{\pgfqpoint{3.142620in}{0.379381in}}%
\pgfpathclose%
\pgfusepath{fill}%
\end{pgfscope}%
\begin{pgfscope}%
\pgfpathrectangle{\pgfqpoint{2.893540in}{0.169444in}}{\pgfqpoint{0.858895in}{1.282944in}}%
\pgfusepath{clip}%
\pgfsetbuttcap%
\pgfsetmiterjoin%
\definecolor{currentfill}{rgb}{0.459592,0.704960,0.230296}%
\pgfsetfillcolor{currentfill}%
\pgfsetlinewidth{0.000000pt}%
\definecolor{currentstroke}{rgb}{0.000000,0.000000,0.000000}%
\pgfsetstrokecolor{currentstroke}%
\pgfsetstrokeopacity{0.000000}%
\pgfsetdash{}{0pt}%
\pgfpathmoveto{\pgfqpoint{3.142620in}{0.402707in}}%
\pgfpathlineto{\pgfqpoint{3.108264in}{0.402707in}}%
\pgfpathlineto{\pgfqpoint{3.108264in}{0.425617in}}%
\pgfpathlineto{\pgfqpoint{3.142620in}{0.425617in}}%
\pgfpathlineto{\pgfqpoint{3.142620in}{0.402707in}}%
\pgfpathclose%
\pgfusepath{fill}%
\end{pgfscope}%
\begin{pgfscope}%
\pgfpathrectangle{\pgfqpoint{2.893540in}{0.169444in}}{\pgfqpoint{0.858895in}{1.282944in}}%
\pgfusepath{clip}%
\pgfsetbuttcap%
\pgfsetmiterjoin%
\definecolor{currentfill}{rgb}{0.498039,0.737255,0.254902}%
\pgfsetfillcolor{currentfill}%
\pgfsetlinewidth{0.000000pt}%
\definecolor{currentstroke}{rgb}{0.000000,0.000000,0.000000}%
\pgfsetstrokecolor{currentstroke}%
\pgfsetstrokeopacity{0.000000}%
\pgfsetdash{}{0pt}%
\pgfpathmoveto{\pgfqpoint{3.142620in}{0.426033in}}%
\pgfpathlineto{\pgfqpoint{3.108264in}{0.426033in}}%
\pgfpathlineto{\pgfqpoint{3.108264in}{0.448943in}}%
\pgfpathlineto{\pgfqpoint{3.142620in}{0.448943in}}%
\pgfpathlineto{\pgfqpoint{3.142620in}{0.426033in}}%
\pgfpathclose%
\pgfusepath{fill}%
\end{pgfscope}%
\begin{pgfscope}%
\pgfpathrectangle{\pgfqpoint{2.893540in}{0.169444in}}{\pgfqpoint{0.858895in}{1.282944in}}%
\pgfusepath{clip}%
\pgfsetbuttcap%
\pgfsetmiterjoin%
\definecolor{currentfill}{rgb}{0.533103,0.760015,0.297347}%
\pgfsetfillcolor{currentfill}%
\pgfsetlinewidth{0.000000pt}%
\definecolor{currentstroke}{rgb}{0.000000,0.000000,0.000000}%
\pgfsetstrokecolor{currentstroke}%
\pgfsetstrokeopacity{0.000000}%
\pgfsetdash{}{0pt}%
\pgfpathmoveto{\pgfqpoint{3.142620in}{0.449360in}}%
\pgfpathlineto{\pgfqpoint{3.108264in}{0.449360in}}%
\pgfpathlineto{\pgfqpoint{3.108264in}{0.472269in}}%
\pgfpathlineto{\pgfqpoint{3.142620in}{0.472269in}}%
\pgfpathlineto{\pgfqpoint{3.142620in}{0.449360in}}%
\pgfpathclose%
\pgfusepath{fill}%
\end{pgfscope}%
\begin{pgfscope}%
\pgfpathrectangle{\pgfqpoint{2.893540in}{0.169444in}}{\pgfqpoint{0.858895in}{1.282944in}}%
\pgfusepath{clip}%
\pgfsetbuttcap%
\pgfsetmiterjoin%
\definecolor{currentfill}{rgb}{0.576932,0.788466,0.350404}%
\pgfsetfillcolor{currentfill}%
\pgfsetlinewidth{0.000000pt}%
\definecolor{currentstroke}{rgb}{0.000000,0.000000,0.000000}%
\pgfsetstrokecolor{currentstroke}%
\pgfsetstrokeopacity{0.000000}%
\pgfsetdash{}{0pt}%
\pgfpathmoveto{\pgfqpoint{3.142620in}{0.472686in}}%
\pgfpathlineto{\pgfqpoint{3.108264in}{0.472686in}}%
\pgfpathlineto{\pgfqpoint{3.108264in}{0.495596in}}%
\pgfpathlineto{\pgfqpoint{3.142620in}{0.495596in}}%
\pgfpathlineto{\pgfqpoint{3.142620in}{0.472686in}}%
\pgfpathclose%
\pgfusepath{fill}%
\end{pgfscope}%
\begin{pgfscope}%
\pgfpathrectangle{\pgfqpoint{2.893540in}{0.169444in}}{\pgfqpoint{0.858895in}{1.282944in}}%
\pgfusepath{clip}%
\pgfsetbuttcap%
\pgfsetmiterjoin%
\definecolor{currentfill}{rgb}{0.620761,0.816917,0.403460}%
\pgfsetfillcolor{currentfill}%
\pgfsetlinewidth{0.000000pt}%
\definecolor{currentstroke}{rgb}{0.000000,0.000000,0.000000}%
\pgfsetstrokecolor{currentstroke}%
\pgfsetstrokeopacity{0.000000}%
\pgfsetdash{}{0pt}%
\pgfpathmoveto{\pgfqpoint{3.142620in}{0.496012in}}%
\pgfpathlineto{\pgfqpoint{3.108264in}{0.496012in}}%
\pgfpathlineto{\pgfqpoint{3.108264in}{0.518922in}}%
\pgfpathlineto{\pgfqpoint{3.142620in}{0.518922in}}%
\pgfpathlineto{\pgfqpoint{3.142620in}{0.496012in}}%
\pgfpathclose%
\pgfusepath{fill}%
\end{pgfscope}%
\begin{pgfscope}%
\pgfpathrectangle{\pgfqpoint{2.893540in}{0.169444in}}{\pgfqpoint{0.858895in}{1.282944in}}%
\pgfusepath{clip}%
\pgfsetbuttcap%
\pgfsetmiterjoin%
\definecolor{currentfill}{rgb}{0.655825,0.839677,0.445905}%
\pgfsetfillcolor{currentfill}%
\pgfsetlinewidth{0.000000pt}%
\definecolor{currentstroke}{rgb}{0.000000,0.000000,0.000000}%
\pgfsetstrokecolor{currentstroke}%
\pgfsetstrokeopacity{0.000000}%
\pgfsetdash{}{0pt}%
\pgfpathmoveto{\pgfqpoint{3.142620in}{0.519338in}}%
\pgfpathlineto{\pgfqpoint{3.108264in}{0.519338in}}%
\pgfpathlineto{\pgfqpoint{3.108264in}{0.542248in}}%
\pgfpathlineto{\pgfqpoint{3.142620in}{0.542248in}}%
\pgfpathlineto{\pgfqpoint{3.142620in}{0.519338in}}%
\pgfpathclose%
\pgfusepath{fill}%
\end{pgfscope}%
\begin{pgfscope}%
\pgfpathrectangle{\pgfqpoint{2.893540in}{0.169444in}}{\pgfqpoint{0.858895in}{1.282944in}}%
\pgfusepath{clip}%
\pgfsetbuttcap%
\pgfsetmiterjoin%
\definecolor{currentfill}{rgb}{0.699654,0.868128,0.498962}%
\pgfsetfillcolor{currentfill}%
\pgfsetlinewidth{0.000000pt}%
\definecolor{currentstroke}{rgb}{0.000000,0.000000,0.000000}%
\pgfsetstrokecolor{currentstroke}%
\pgfsetstrokeopacity{0.000000}%
\pgfsetdash{}{0pt}%
\pgfpathmoveto{\pgfqpoint{3.142620in}{0.542665in}}%
\pgfpathlineto{\pgfqpoint{3.108264in}{0.542665in}}%
\pgfpathlineto{\pgfqpoint{3.108264in}{0.565574in}}%
\pgfpathlineto{\pgfqpoint{3.142620in}{0.565574in}}%
\pgfpathlineto{\pgfqpoint{3.142620in}{0.542665in}}%
\pgfpathclose%
\pgfusepath{fill}%
\end{pgfscope}%
\begin{pgfscope}%
\pgfpathrectangle{\pgfqpoint{2.893540in}{0.169444in}}{\pgfqpoint{0.858895in}{1.282944in}}%
\pgfusepath{clip}%
\pgfsetbuttcap%
\pgfsetmiterjoin%
\definecolor{currentfill}{rgb}{0.739254,0.890042,0.553941}%
\pgfsetfillcolor{currentfill}%
\pgfsetlinewidth{0.000000pt}%
\definecolor{currentstroke}{rgb}{0.000000,0.000000,0.000000}%
\pgfsetstrokecolor{currentstroke}%
\pgfsetstrokeopacity{0.000000}%
\pgfsetdash{}{0pt}%
\pgfpathmoveto{\pgfqpoint{3.142620in}{0.565991in}}%
\pgfpathlineto{\pgfqpoint{3.108264in}{0.565991in}}%
\pgfpathlineto{\pgfqpoint{3.108264in}{0.588901in}}%
\pgfpathlineto{\pgfqpoint{3.142620in}{0.588901in}}%
\pgfpathlineto{\pgfqpoint{3.142620in}{0.565991in}}%
\pgfpathclose%
\pgfusepath{fill}%
\end{pgfscope}%
\begin{pgfscope}%
\pgfpathrectangle{\pgfqpoint{2.893540in}{0.169444in}}{\pgfqpoint{0.858895in}{1.282944in}}%
\pgfusepath{clip}%
\pgfsetbuttcap%
\pgfsetmiterjoin%
\definecolor{currentfill}{rgb}{0.767551,0.902345,0.599462}%
\pgfsetfillcolor{currentfill}%
\pgfsetlinewidth{0.000000pt}%
\definecolor{currentstroke}{rgb}{0.000000,0.000000,0.000000}%
\pgfsetstrokecolor{currentstroke}%
\pgfsetstrokeopacity{0.000000}%
\pgfsetdash{}{0pt}%
\pgfpathmoveto{\pgfqpoint{3.142620in}{0.589317in}}%
\pgfpathlineto{\pgfqpoint{3.108264in}{0.589317in}}%
\pgfpathlineto{\pgfqpoint{3.108264in}{0.612227in}}%
\pgfpathlineto{\pgfqpoint{3.142620in}{0.612227in}}%
\pgfpathlineto{\pgfqpoint{3.142620in}{0.589317in}}%
\pgfpathclose%
\pgfusepath{fill}%
\end{pgfscope}%
\begin{pgfscope}%
\pgfpathrectangle{\pgfqpoint{2.893540in}{0.169444in}}{\pgfqpoint{0.858895in}{1.282944in}}%
\pgfusepath{clip}%
\pgfsetbuttcap%
\pgfsetmiterjoin%
\definecolor{currentfill}{rgb}{0.802922,0.917724,0.656363}%
\pgfsetfillcolor{currentfill}%
\pgfsetlinewidth{0.000000pt}%
\definecolor{currentstroke}{rgb}{0.000000,0.000000,0.000000}%
\pgfsetstrokecolor{currentstroke}%
\pgfsetstrokeopacity{0.000000}%
\pgfsetdash{}{0pt}%
\pgfpathmoveto{\pgfqpoint{3.142620in}{0.612643in}}%
\pgfpathlineto{\pgfqpoint{3.108264in}{0.612643in}}%
\pgfpathlineto{\pgfqpoint{3.108264in}{0.635553in}}%
\pgfpathlineto{\pgfqpoint{3.142620in}{0.635553in}}%
\pgfpathlineto{\pgfqpoint{3.142620in}{0.612643in}}%
\pgfpathclose%
\pgfusepath{fill}%
\end{pgfscope}%
\begin{pgfscope}%
\pgfpathrectangle{\pgfqpoint{2.893540in}{0.169444in}}{\pgfqpoint{0.858895in}{1.282944in}}%
\pgfusepath{clip}%
\pgfsetbuttcap%
\pgfsetmiterjoin%
\definecolor{currentfill}{rgb}{0.838293,0.933103,0.713264}%
\pgfsetfillcolor{currentfill}%
\pgfsetlinewidth{0.000000pt}%
\definecolor{currentstroke}{rgb}{0.000000,0.000000,0.000000}%
\pgfsetstrokecolor{currentstroke}%
\pgfsetstrokeopacity{0.000000}%
\pgfsetdash{}{0pt}%
\pgfpathmoveto{\pgfqpoint{3.142620in}{0.635970in}}%
\pgfpathlineto{\pgfqpoint{3.108264in}{0.635970in}}%
\pgfpathlineto{\pgfqpoint{3.108264in}{0.658879in}}%
\pgfpathlineto{\pgfqpoint{3.142620in}{0.658879in}}%
\pgfpathlineto{\pgfqpoint{3.142620in}{0.635970in}}%
\pgfpathclose%
\pgfusepath{fill}%
\end{pgfscope}%
\begin{pgfscope}%
\pgfpathrectangle{\pgfqpoint{2.893540in}{0.169444in}}{\pgfqpoint{0.858895in}{1.282944in}}%
\pgfusepath{clip}%
\pgfsetbuttcap%
\pgfsetmiterjoin%
\definecolor{currentfill}{rgb}{0.866590,0.945406,0.758785}%
\pgfsetfillcolor{currentfill}%
\pgfsetlinewidth{0.000000pt}%
\definecolor{currentstroke}{rgb}{0.000000,0.000000,0.000000}%
\pgfsetstrokecolor{currentstroke}%
\pgfsetstrokeopacity{0.000000}%
\pgfsetdash{}{0pt}%
\pgfpathmoveto{\pgfqpoint{3.142620in}{0.659296in}}%
\pgfpathlineto{\pgfqpoint{3.108264in}{0.659296in}}%
\pgfpathlineto{\pgfqpoint{3.108264in}{0.682206in}}%
\pgfpathlineto{\pgfqpoint{3.142620in}{0.682206in}}%
\pgfpathlineto{\pgfqpoint{3.142620in}{0.659296in}}%
\pgfpathclose%
\pgfusepath{fill}%
\end{pgfscope}%
\begin{pgfscope}%
\pgfpathrectangle{\pgfqpoint{2.893540in}{0.169444in}}{\pgfqpoint{0.858895in}{1.282944in}}%
\pgfusepath{clip}%
\pgfsetbuttcap%
\pgfsetmiterjoin%
\definecolor{currentfill}{rgb}{0.901961,0.960784,0.815686}%
\pgfsetfillcolor{currentfill}%
\pgfsetlinewidth{0.000000pt}%
\definecolor{currentstroke}{rgb}{0.000000,0.000000,0.000000}%
\pgfsetstrokecolor{currentstroke}%
\pgfsetstrokeopacity{0.000000}%
\pgfsetdash{}{0pt}%
\pgfpathmoveto{\pgfqpoint{3.142620in}{0.682622in}}%
\pgfpathlineto{\pgfqpoint{3.108264in}{0.682622in}}%
\pgfpathlineto{\pgfqpoint{3.108264in}{0.705532in}}%
\pgfpathlineto{\pgfqpoint{3.142620in}{0.705532in}}%
\pgfpathlineto{\pgfqpoint{3.142620in}{0.682622in}}%
\pgfpathclose%
\pgfusepath{fill}%
\end{pgfscope}%
\begin{pgfscope}%
\pgfpathrectangle{\pgfqpoint{2.893540in}{0.169444in}}{\pgfqpoint{0.858895in}{1.282944in}}%
\pgfusepath{clip}%
\pgfsetbuttcap%
\pgfsetmiterjoin%
\definecolor{currentfill}{rgb}{0.915033,0.962322,0.845675}%
\pgfsetfillcolor{currentfill}%
\pgfsetlinewidth{0.000000pt}%
\definecolor{currentstroke}{rgb}{0.000000,0.000000,0.000000}%
\pgfsetstrokecolor{currentstroke}%
\pgfsetstrokeopacity{0.000000}%
\pgfsetdash{}{0pt}%
\pgfpathmoveto{\pgfqpoint{3.142620in}{0.705948in}}%
\pgfpathlineto{\pgfqpoint{3.108264in}{0.705948in}}%
\pgfpathlineto{\pgfqpoint{3.108264in}{0.728858in}}%
\pgfpathlineto{\pgfqpoint{3.142620in}{0.728858in}}%
\pgfpathlineto{\pgfqpoint{3.142620in}{0.705948in}}%
\pgfpathclose%
\pgfusepath{fill}%
\end{pgfscope}%
\begin{pgfscope}%
\pgfpathrectangle{\pgfqpoint{2.893540in}{0.169444in}}{\pgfqpoint{0.858895in}{1.282944in}}%
\pgfusepath{clip}%
\pgfsetbuttcap%
\pgfsetmiterjoin%
\definecolor{currentfill}{rgb}{0.925490,0.963552,0.869666}%
\pgfsetfillcolor{currentfill}%
\pgfsetlinewidth{0.000000pt}%
\definecolor{currentstroke}{rgb}{0.000000,0.000000,0.000000}%
\pgfsetstrokecolor{currentstroke}%
\pgfsetstrokeopacity{0.000000}%
\pgfsetdash{}{0pt}%
\pgfpathmoveto{\pgfqpoint{3.142620in}{0.729275in}}%
\pgfpathlineto{\pgfqpoint{3.108264in}{0.729275in}}%
\pgfpathlineto{\pgfqpoint{3.108264in}{0.752184in}}%
\pgfpathlineto{\pgfqpoint{3.142620in}{0.752184in}}%
\pgfpathlineto{\pgfqpoint{3.142620in}{0.729275in}}%
\pgfpathclose%
\pgfusepath{fill}%
\end{pgfscope}%
\begin{pgfscope}%
\pgfpathrectangle{\pgfqpoint{2.893540in}{0.169444in}}{\pgfqpoint{0.858895in}{1.282944in}}%
\pgfusepath{clip}%
\pgfsetbuttcap%
\pgfsetmiterjoin%
\definecolor{currentfill}{rgb}{0.938562,0.965090,0.899654}%
\pgfsetfillcolor{currentfill}%
\pgfsetlinewidth{0.000000pt}%
\definecolor{currentstroke}{rgb}{0.000000,0.000000,0.000000}%
\pgfsetstrokecolor{currentstroke}%
\pgfsetstrokeopacity{0.000000}%
\pgfsetdash{}{0pt}%
\pgfpathmoveto{\pgfqpoint{3.142620in}{0.752601in}}%
\pgfpathlineto{\pgfqpoint{3.108264in}{0.752601in}}%
\pgfpathlineto{\pgfqpoint{3.108264in}{0.775511in}}%
\pgfpathlineto{\pgfqpoint{3.142620in}{0.775511in}}%
\pgfpathlineto{\pgfqpoint{3.142620in}{0.752601in}}%
\pgfpathclose%
\pgfusepath{fill}%
\end{pgfscope}%
\begin{pgfscope}%
\pgfpathrectangle{\pgfqpoint{2.893540in}{0.169444in}}{\pgfqpoint{0.858895in}{1.282944in}}%
\pgfusepath{clip}%
\pgfsetbuttcap%
\pgfsetmiterjoin%
\definecolor{currentfill}{rgb}{0.951634,0.966628,0.929642}%
\pgfsetfillcolor{currentfill}%
\pgfsetlinewidth{0.000000pt}%
\definecolor{currentstroke}{rgb}{0.000000,0.000000,0.000000}%
\pgfsetstrokecolor{currentstroke}%
\pgfsetstrokeopacity{0.000000}%
\pgfsetdash{}{0pt}%
\pgfpathmoveto{\pgfqpoint{3.142620in}{0.775927in}}%
\pgfpathlineto{\pgfqpoint{3.108264in}{0.775927in}}%
\pgfpathlineto{\pgfqpoint{3.108264in}{0.798837in}}%
\pgfpathlineto{\pgfqpoint{3.142620in}{0.798837in}}%
\pgfpathlineto{\pgfqpoint{3.142620in}{0.775927in}}%
\pgfpathclose%
\pgfusepath{fill}%
\end{pgfscope}%
\begin{pgfscope}%
\pgfpathrectangle{\pgfqpoint{2.893540in}{0.169444in}}{\pgfqpoint{0.858895in}{1.282944in}}%
\pgfusepath{clip}%
\pgfsetbuttcap%
\pgfsetmiterjoin%
\definecolor{currentfill}{rgb}{0.962092,0.967859,0.953633}%
\pgfsetfillcolor{currentfill}%
\pgfsetlinewidth{0.000000pt}%
\definecolor{currentstroke}{rgb}{0.000000,0.000000,0.000000}%
\pgfsetstrokecolor{currentstroke}%
\pgfsetstrokeopacity{0.000000}%
\pgfsetdash{}{0pt}%
\pgfpathmoveto{\pgfqpoint{3.142620in}{0.799254in}}%
\pgfpathlineto{\pgfqpoint{3.108264in}{0.799254in}}%
\pgfpathlineto{\pgfqpoint{3.108264in}{0.822163in}}%
\pgfpathlineto{\pgfqpoint{3.142620in}{0.822163in}}%
\pgfpathlineto{\pgfqpoint{3.142620in}{0.799254in}}%
\pgfpathclose%
\pgfusepath{fill}%
\end{pgfscope}%
\begin{pgfscope}%
\pgfpathrectangle{\pgfqpoint{2.893540in}{0.169444in}}{\pgfqpoint{0.858895in}{1.282944in}}%
\pgfusepath{clip}%
\pgfsetbuttcap%
\pgfsetmiterjoin%
\definecolor{currentfill}{rgb}{0.970934,0.959785,0.965552}%
\pgfsetfillcolor{currentfill}%
\pgfsetlinewidth{0.000000pt}%
\definecolor{currentstroke}{rgb}{0.000000,0.000000,0.000000}%
\pgfsetstrokecolor{currentstroke}%
\pgfsetstrokeopacity{0.000000}%
\pgfsetdash{}{0pt}%
\pgfpathmoveto{\pgfqpoint{3.142620in}{0.822580in}}%
\pgfpathlineto{\pgfqpoint{3.108264in}{0.822580in}}%
\pgfpathlineto{\pgfqpoint{3.108264in}{0.845490in}}%
\pgfpathlineto{\pgfqpoint{3.142620in}{0.845490in}}%
\pgfpathlineto{\pgfqpoint{3.142620in}{0.822580in}}%
\pgfpathclose%
\pgfusepath{fill}%
\end{pgfscope}%
\begin{pgfscope}%
\pgfpathrectangle{\pgfqpoint{2.893540in}{0.169444in}}{\pgfqpoint{0.858895in}{1.282944in}}%
\pgfusepath{clip}%
\pgfsetbuttcap%
\pgfsetmiterjoin%
\definecolor{currentfill}{rgb}{0.974625,0.945636,0.960631}%
\pgfsetfillcolor{currentfill}%
\pgfsetlinewidth{0.000000pt}%
\definecolor{currentstroke}{rgb}{0.000000,0.000000,0.000000}%
\pgfsetstrokecolor{currentstroke}%
\pgfsetstrokeopacity{0.000000}%
\pgfsetdash{}{0pt}%
\pgfpathmoveto{\pgfqpoint{3.142620in}{0.845906in}}%
\pgfpathlineto{\pgfqpoint{3.108264in}{0.845906in}}%
\pgfpathlineto{\pgfqpoint{3.108264in}{0.868816in}}%
\pgfpathlineto{\pgfqpoint{3.142620in}{0.868816in}}%
\pgfpathlineto{\pgfqpoint{3.142620in}{0.845906in}}%
\pgfpathclose%
\pgfusepath{fill}%
\end{pgfscope}%
\begin{pgfscope}%
\pgfpathrectangle{\pgfqpoint{2.893540in}{0.169444in}}{\pgfqpoint{0.858895in}{1.282944in}}%
\pgfusepath{clip}%
\pgfsetbuttcap%
\pgfsetmiterjoin%
\definecolor{currentfill}{rgb}{0.979239,0.927951,0.954479}%
\pgfsetfillcolor{currentfill}%
\pgfsetlinewidth{0.000000pt}%
\definecolor{currentstroke}{rgb}{0.000000,0.000000,0.000000}%
\pgfsetstrokecolor{currentstroke}%
\pgfsetstrokeopacity{0.000000}%
\pgfsetdash{}{0pt}%
\pgfpathmoveto{\pgfqpoint{3.142620in}{0.869232in}}%
\pgfpathlineto{\pgfqpoint{3.108264in}{0.869232in}}%
\pgfpathlineto{\pgfqpoint{3.108264in}{0.892142in}}%
\pgfpathlineto{\pgfqpoint{3.142620in}{0.892142in}}%
\pgfpathlineto{\pgfqpoint{3.142620in}{0.869232in}}%
\pgfpathclose%
\pgfusepath{fill}%
\end{pgfscope}%
\begin{pgfscope}%
\pgfpathrectangle{\pgfqpoint{2.893540in}{0.169444in}}{\pgfqpoint{0.858895in}{1.282944in}}%
\pgfusepath{clip}%
\pgfsetbuttcap%
\pgfsetmiterjoin%
\definecolor{currentfill}{rgb}{0.983852,0.910265,0.948328}%
\pgfsetfillcolor{currentfill}%
\pgfsetlinewidth{0.000000pt}%
\definecolor{currentstroke}{rgb}{0.000000,0.000000,0.000000}%
\pgfsetstrokecolor{currentstroke}%
\pgfsetstrokeopacity{0.000000}%
\pgfsetdash{}{0pt}%
\pgfpathmoveto{\pgfqpoint{3.142620in}{0.892559in}}%
\pgfpathlineto{\pgfqpoint{3.108264in}{0.892559in}}%
\pgfpathlineto{\pgfqpoint{3.108264in}{0.915468in}}%
\pgfpathlineto{\pgfqpoint{3.142620in}{0.915468in}}%
\pgfpathlineto{\pgfqpoint{3.142620in}{0.892559in}}%
\pgfpathclose%
\pgfusepath{fill}%
\end{pgfscope}%
\begin{pgfscope}%
\pgfpathrectangle{\pgfqpoint{2.893540in}{0.169444in}}{\pgfqpoint{0.858895in}{1.282944in}}%
\pgfusepath{clip}%
\pgfsetbuttcap%
\pgfsetmiterjoin%
\definecolor{currentfill}{rgb}{0.987543,0.896117,0.943406}%
\pgfsetfillcolor{currentfill}%
\pgfsetlinewidth{0.000000pt}%
\definecolor{currentstroke}{rgb}{0.000000,0.000000,0.000000}%
\pgfsetstrokecolor{currentstroke}%
\pgfsetstrokeopacity{0.000000}%
\pgfsetdash{}{0pt}%
\pgfpathmoveto{\pgfqpoint{3.142620in}{0.915885in}}%
\pgfpathlineto{\pgfqpoint{3.108264in}{0.915885in}}%
\pgfpathlineto{\pgfqpoint{3.108264in}{0.938795in}}%
\pgfpathlineto{\pgfqpoint{3.142620in}{0.938795in}}%
\pgfpathlineto{\pgfqpoint{3.142620in}{0.915885in}}%
\pgfpathclose%
\pgfusepath{fill}%
\end{pgfscope}%
\begin{pgfscope}%
\pgfpathrectangle{\pgfqpoint{2.893540in}{0.169444in}}{\pgfqpoint{0.858895in}{1.282944in}}%
\pgfusepath{clip}%
\pgfsetbuttcap%
\pgfsetmiterjoin%
\definecolor{currentfill}{rgb}{0.992157,0.878431,0.937255}%
\pgfsetfillcolor{currentfill}%
\pgfsetlinewidth{0.000000pt}%
\definecolor{currentstroke}{rgb}{0.000000,0.000000,0.000000}%
\pgfsetstrokecolor{currentstroke}%
\pgfsetstrokeopacity{0.000000}%
\pgfsetdash{}{0pt}%
\pgfpathmoveto{\pgfqpoint{3.142620in}{0.939211in}}%
\pgfpathlineto{\pgfqpoint{3.108264in}{0.939211in}}%
\pgfpathlineto{\pgfqpoint{3.108264in}{0.962121in}}%
\pgfpathlineto{\pgfqpoint{3.142620in}{0.962121in}}%
\pgfpathlineto{\pgfqpoint{3.142620in}{0.939211in}}%
\pgfpathclose%
\pgfusepath{fill}%
\end{pgfscope}%
\begin{pgfscope}%
\pgfpathrectangle{\pgfqpoint{2.893540in}{0.169444in}}{\pgfqpoint{0.858895in}{1.282944in}}%
\pgfusepath{clip}%
\pgfsetbuttcap%
\pgfsetmiterjoin%
\definecolor{currentfill}{rgb}{0.982930,0.846136,0.921107}%
\pgfsetfillcolor{currentfill}%
\pgfsetlinewidth{0.000000pt}%
\definecolor{currentstroke}{rgb}{0.000000,0.000000,0.000000}%
\pgfsetstrokecolor{currentstroke}%
\pgfsetstrokeopacity{0.000000}%
\pgfsetdash{}{0pt}%
\pgfpathmoveto{\pgfqpoint{3.142620in}{0.962537in}}%
\pgfpathlineto{\pgfqpoint{3.108264in}{0.962537in}}%
\pgfpathlineto{\pgfqpoint{3.108264in}{0.985447in}}%
\pgfpathlineto{\pgfqpoint{3.142620in}{0.985447in}}%
\pgfpathlineto{\pgfqpoint{3.142620in}{0.962537in}}%
\pgfpathclose%
\pgfusepath{fill}%
\end{pgfscope}%
\begin{pgfscope}%
\pgfpathrectangle{\pgfqpoint{2.893540in}{0.169444in}}{\pgfqpoint{0.858895in}{1.282944in}}%
\pgfusepath{clip}%
\pgfsetbuttcap%
\pgfsetmiterjoin%
\definecolor{currentfill}{rgb}{0.975548,0.820300,0.908189}%
\pgfsetfillcolor{currentfill}%
\pgfsetlinewidth{0.000000pt}%
\definecolor{currentstroke}{rgb}{0.000000,0.000000,0.000000}%
\pgfsetstrokecolor{currentstroke}%
\pgfsetstrokeopacity{0.000000}%
\pgfsetdash{}{0pt}%
\pgfpathmoveto{\pgfqpoint{3.142620in}{0.985864in}}%
\pgfpathlineto{\pgfqpoint{3.108264in}{0.985864in}}%
\pgfpathlineto{\pgfqpoint{3.108264in}{1.008773in}}%
\pgfpathlineto{\pgfqpoint{3.142620in}{1.008773in}}%
\pgfpathlineto{\pgfqpoint{3.142620in}{0.985864in}}%
\pgfpathclose%
\pgfusepath{fill}%
\end{pgfscope}%
\begin{pgfscope}%
\pgfpathrectangle{\pgfqpoint{2.893540in}{0.169444in}}{\pgfqpoint{0.858895in}{1.282944in}}%
\pgfusepath{clip}%
\pgfsetbuttcap%
\pgfsetmiterjoin%
\definecolor{currentfill}{rgb}{0.966321,0.788005,0.892042}%
\pgfsetfillcolor{currentfill}%
\pgfsetlinewidth{0.000000pt}%
\definecolor{currentstroke}{rgb}{0.000000,0.000000,0.000000}%
\pgfsetstrokecolor{currentstroke}%
\pgfsetstrokeopacity{0.000000}%
\pgfsetdash{}{0pt}%
\pgfpathmoveto{\pgfqpoint{3.142620in}{1.009190in}}%
\pgfpathlineto{\pgfqpoint{3.108264in}{1.009190in}}%
\pgfpathlineto{\pgfqpoint{3.108264in}{1.032100in}}%
\pgfpathlineto{\pgfqpoint{3.142620in}{1.032100in}}%
\pgfpathlineto{\pgfqpoint{3.142620in}{1.009190in}}%
\pgfpathclose%
\pgfusepath{fill}%
\end{pgfscope}%
\begin{pgfscope}%
\pgfpathrectangle{\pgfqpoint{2.893540in}{0.169444in}}{\pgfqpoint{0.858895in}{1.282944in}}%
\pgfusepath{clip}%
\pgfsetbuttcap%
\pgfsetmiterjoin%
\definecolor{currentfill}{rgb}{0.957093,0.755709,0.875894}%
\pgfsetfillcolor{currentfill}%
\pgfsetlinewidth{0.000000pt}%
\definecolor{currentstroke}{rgb}{0.000000,0.000000,0.000000}%
\pgfsetstrokecolor{currentstroke}%
\pgfsetstrokeopacity{0.000000}%
\pgfsetdash{}{0pt}%
\pgfpathmoveto{\pgfqpoint{3.142620in}{1.032516in}}%
\pgfpathlineto{\pgfqpoint{3.108264in}{1.032516in}}%
\pgfpathlineto{\pgfqpoint{3.108264in}{1.055426in}}%
\pgfpathlineto{\pgfqpoint{3.142620in}{1.055426in}}%
\pgfpathlineto{\pgfqpoint{3.142620in}{1.032516in}}%
\pgfpathclose%
\pgfusepath{fill}%
\end{pgfscope}%
\begin{pgfscope}%
\pgfpathrectangle{\pgfqpoint{2.893540in}{0.169444in}}{\pgfqpoint{0.858895in}{1.282944in}}%
\pgfusepath{clip}%
\pgfsetbuttcap%
\pgfsetmiterjoin%
\definecolor{currentfill}{rgb}{0.949712,0.729873,0.862976}%
\pgfsetfillcolor{currentfill}%
\pgfsetlinewidth{0.000000pt}%
\definecolor{currentstroke}{rgb}{0.000000,0.000000,0.000000}%
\pgfsetstrokecolor{currentstroke}%
\pgfsetstrokeopacity{0.000000}%
\pgfsetdash{}{0pt}%
\pgfpathmoveto{\pgfqpoint{3.142620in}{1.055842in}}%
\pgfpathlineto{\pgfqpoint{3.108264in}{1.055842in}}%
\pgfpathlineto{\pgfqpoint{3.108264in}{1.078752in}}%
\pgfpathlineto{\pgfqpoint{3.142620in}{1.078752in}}%
\pgfpathlineto{\pgfqpoint{3.142620in}{1.055842in}}%
\pgfpathclose%
\pgfusepath{fill}%
\end{pgfscope}%
\begin{pgfscope}%
\pgfpathrectangle{\pgfqpoint{2.893540in}{0.169444in}}{\pgfqpoint{0.858895in}{1.282944in}}%
\pgfusepath{clip}%
\pgfsetbuttcap%
\pgfsetmiterjoin%
\definecolor{currentfill}{rgb}{0.937793,0.689504,0.837985}%
\pgfsetfillcolor{currentfill}%
\pgfsetlinewidth{0.000000pt}%
\definecolor{currentstroke}{rgb}{0.000000,0.000000,0.000000}%
\pgfsetstrokecolor{currentstroke}%
\pgfsetstrokeopacity{0.000000}%
\pgfsetdash{}{0pt}%
\pgfpathmoveto{\pgfqpoint{3.142620in}{1.079169in}}%
\pgfpathlineto{\pgfqpoint{3.108264in}{1.079169in}}%
\pgfpathlineto{\pgfqpoint{3.108264in}{1.102078in}}%
\pgfpathlineto{\pgfqpoint{3.142620in}{1.102078in}}%
\pgfpathlineto{\pgfqpoint{3.142620in}{1.079169in}}%
\pgfpathclose%
\pgfusepath{fill}%
\end{pgfscope}%
\begin{pgfscope}%
\pgfpathrectangle{\pgfqpoint{2.893540in}{0.169444in}}{\pgfqpoint{0.858895in}{1.282944in}}%
\pgfusepath{clip}%
\pgfsetbuttcap%
\pgfsetmiterjoin%
\definecolor{currentfill}{rgb}{0.923183,0.641061,0.804152}%
\pgfsetfillcolor{currentfill}%
\pgfsetlinewidth{0.000000pt}%
\definecolor{currentstroke}{rgb}{0.000000,0.000000,0.000000}%
\pgfsetstrokecolor{currentstroke}%
\pgfsetstrokeopacity{0.000000}%
\pgfsetdash{}{0pt}%
\pgfpathmoveto{\pgfqpoint{3.142620in}{1.102495in}}%
\pgfpathlineto{\pgfqpoint{3.108264in}{1.102495in}}%
\pgfpathlineto{\pgfqpoint{3.108264in}{1.125405in}}%
\pgfpathlineto{\pgfqpoint{3.142620in}{1.125405in}}%
\pgfpathlineto{\pgfqpoint{3.142620in}{1.102495in}}%
\pgfpathclose%
\pgfusepath{fill}%
\end{pgfscope}%
\begin{pgfscope}%
\pgfpathrectangle{\pgfqpoint{2.893540in}{0.169444in}}{\pgfqpoint{0.858895in}{1.282944in}}%
\pgfusepath{clip}%
\pgfsetbuttcap%
\pgfsetmiterjoin%
\definecolor{currentfill}{rgb}{0.911496,0.602307,0.777086}%
\pgfsetfillcolor{currentfill}%
\pgfsetlinewidth{0.000000pt}%
\definecolor{currentstroke}{rgb}{0.000000,0.000000,0.000000}%
\pgfsetstrokecolor{currentstroke}%
\pgfsetstrokeopacity{0.000000}%
\pgfsetdash{}{0pt}%
\pgfpathmoveto{\pgfqpoint{3.142620in}{1.125821in}}%
\pgfpathlineto{\pgfqpoint{3.108264in}{1.125821in}}%
\pgfpathlineto{\pgfqpoint{3.108264in}{1.148731in}}%
\pgfpathlineto{\pgfqpoint{3.142620in}{1.148731in}}%
\pgfpathlineto{\pgfqpoint{3.142620in}{1.125821in}}%
\pgfpathclose%
\pgfusepath{fill}%
\end{pgfscope}%
\begin{pgfscope}%
\pgfpathrectangle{\pgfqpoint{2.893540in}{0.169444in}}{\pgfqpoint{0.858895in}{1.282944in}}%
\pgfusepath{clip}%
\pgfsetbuttcap%
\pgfsetmiterjoin%
\definecolor{currentfill}{rgb}{0.896886,0.553864,0.743253}%
\pgfsetfillcolor{currentfill}%
\pgfsetlinewidth{0.000000pt}%
\definecolor{currentstroke}{rgb}{0.000000,0.000000,0.000000}%
\pgfsetstrokecolor{currentstroke}%
\pgfsetstrokeopacity{0.000000}%
\pgfsetdash{}{0pt}%
\pgfpathmoveto{\pgfqpoint{3.142620in}{1.149147in}}%
\pgfpathlineto{\pgfqpoint{3.108264in}{1.149147in}}%
\pgfpathlineto{\pgfqpoint{3.108264in}{1.172057in}}%
\pgfpathlineto{\pgfqpoint{3.142620in}{1.172057in}}%
\pgfpathlineto{\pgfqpoint{3.142620in}{1.149147in}}%
\pgfpathclose%
\pgfusepath{fill}%
\end{pgfscope}%
\begin{pgfscope}%
\pgfpathrectangle{\pgfqpoint{2.893540in}{0.169444in}}{\pgfqpoint{0.858895in}{1.282944in}}%
\pgfusepath{clip}%
\pgfsetbuttcap%
\pgfsetmiterjoin%
\definecolor{currentfill}{rgb}{0.882276,0.505421,0.709419}%
\pgfsetfillcolor{currentfill}%
\pgfsetlinewidth{0.000000pt}%
\definecolor{currentstroke}{rgb}{0.000000,0.000000,0.000000}%
\pgfsetstrokecolor{currentstroke}%
\pgfsetstrokeopacity{0.000000}%
\pgfsetdash{}{0pt}%
\pgfpathmoveto{\pgfqpoint{3.142620in}{1.172474in}}%
\pgfpathlineto{\pgfqpoint{3.108264in}{1.172474in}}%
\pgfpathlineto{\pgfqpoint{3.108264in}{1.195383in}}%
\pgfpathlineto{\pgfqpoint{3.142620in}{1.195383in}}%
\pgfpathlineto{\pgfqpoint{3.142620in}{1.172474in}}%
\pgfpathclose%
\pgfusepath{fill}%
\end{pgfscope}%
\begin{pgfscope}%
\pgfpathrectangle{\pgfqpoint{2.893540in}{0.169444in}}{\pgfqpoint{0.858895in}{1.282944in}}%
\pgfusepath{clip}%
\pgfsetbuttcap%
\pgfsetmiterjoin%
\definecolor{currentfill}{rgb}{0.870588,0.466667,0.682353}%
\pgfsetfillcolor{currentfill}%
\pgfsetlinewidth{0.000000pt}%
\definecolor{currentstroke}{rgb}{0.000000,0.000000,0.000000}%
\pgfsetstrokecolor{currentstroke}%
\pgfsetstrokeopacity{0.000000}%
\pgfsetdash{}{0pt}%
\pgfpathmoveto{\pgfqpoint{3.142620in}{1.195800in}}%
\pgfpathlineto{\pgfqpoint{3.108264in}{1.195800in}}%
\pgfpathlineto{\pgfqpoint{3.108264in}{1.218710in}}%
\pgfpathlineto{\pgfqpoint{3.142620in}{1.218710in}}%
\pgfpathlineto{\pgfqpoint{3.142620in}{1.195800in}}%
\pgfpathclose%
\pgfusepath{fill}%
\end{pgfscope}%
\begin{pgfscope}%
\pgfpathrectangle{\pgfqpoint{2.893540in}{0.169444in}}{\pgfqpoint{0.858895in}{1.282944in}}%
\pgfusepath{clip}%
\pgfsetbuttcap%
\pgfsetmiterjoin%
\definecolor{currentfill}{rgb}{0.851365,0.395925,0.644675}%
\pgfsetfillcolor{currentfill}%
\pgfsetlinewidth{0.000000pt}%
\definecolor{currentstroke}{rgb}{0.000000,0.000000,0.000000}%
\pgfsetstrokecolor{currentstroke}%
\pgfsetstrokeopacity{0.000000}%
\pgfsetdash{}{0pt}%
\pgfpathmoveto{\pgfqpoint{3.142620in}{1.219126in}}%
\pgfpathlineto{\pgfqpoint{3.108264in}{1.219126in}}%
\pgfpathlineto{\pgfqpoint{3.108264in}{1.242036in}}%
\pgfpathlineto{\pgfqpoint{3.142620in}{1.242036in}}%
\pgfpathlineto{\pgfqpoint{3.142620in}{1.219126in}}%
\pgfpathclose%
\pgfusepath{fill}%
\end{pgfscope}%
\begin{pgfscope}%
\pgfpathrectangle{\pgfqpoint{2.893540in}{0.169444in}}{\pgfqpoint{0.858895in}{1.282944in}}%
\pgfusepath{clip}%
\pgfsetbuttcap%
\pgfsetmiterjoin%
\definecolor{currentfill}{rgb}{0.832141,0.325183,0.606997}%
\pgfsetfillcolor{currentfill}%
\pgfsetlinewidth{0.000000pt}%
\definecolor{currentstroke}{rgb}{0.000000,0.000000,0.000000}%
\pgfsetstrokecolor{currentstroke}%
\pgfsetstrokeopacity{0.000000}%
\pgfsetdash{}{0pt}%
\pgfpathmoveto{\pgfqpoint{3.142620in}{1.242453in}}%
\pgfpathlineto{\pgfqpoint{3.108264in}{1.242453in}}%
\pgfpathlineto{\pgfqpoint{3.108264in}{1.265362in}}%
\pgfpathlineto{\pgfqpoint{3.142620in}{1.265362in}}%
\pgfpathlineto{\pgfqpoint{3.142620in}{1.242453in}}%
\pgfpathclose%
\pgfusepath{fill}%
\end{pgfscope}%
\begin{pgfscope}%
\pgfpathrectangle{\pgfqpoint{2.893540in}{0.169444in}}{\pgfqpoint{0.858895in}{1.282944in}}%
\pgfusepath{clip}%
\pgfsetbuttcap%
\pgfsetmiterjoin%
\definecolor{currentfill}{rgb}{0.816763,0.268589,0.576855}%
\pgfsetfillcolor{currentfill}%
\pgfsetlinewidth{0.000000pt}%
\definecolor{currentstroke}{rgb}{0.000000,0.000000,0.000000}%
\pgfsetstrokecolor{currentstroke}%
\pgfsetstrokeopacity{0.000000}%
\pgfsetdash{}{0pt}%
\pgfpathmoveto{\pgfqpoint{3.142620in}{1.265779in}}%
\pgfpathlineto{\pgfqpoint{3.108264in}{1.265779in}}%
\pgfpathlineto{\pgfqpoint{3.108264in}{1.288689in}}%
\pgfpathlineto{\pgfqpoint{3.142620in}{1.288689in}}%
\pgfpathlineto{\pgfqpoint{3.142620in}{1.265779in}}%
\pgfpathclose%
\pgfusepath{fill}%
\end{pgfscope}%
\begin{pgfscope}%
\pgfpathrectangle{\pgfqpoint{2.893540in}{0.169444in}}{\pgfqpoint{0.858895in}{1.282944in}}%
\pgfusepath{clip}%
\pgfsetbuttcap%
\pgfsetmiterjoin%
\definecolor{currentfill}{rgb}{0.797539,0.197847,0.539177}%
\pgfsetfillcolor{currentfill}%
\pgfsetlinewidth{0.000000pt}%
\definecolor{currentstroke}{rgb}{0.000000,0.000000,0.000000}%
\pgfsetstrokecolor{currentstroke}%
\pgfsetstrokeopacity{0.000000}%
\pgfsetdash{}{0pt}%
\pgfpathmoveto{\pgfqpoint{3.142620in}{1.289105in}}%
\pgfpathlineto{\pgfqpoint{3.108264in}{1.289105in}}%
\pgfpathlineto{\pgfqpoint{3.108264in}{1.312015in}}%
\pgfpathlineto{\pgfqpoint{3.142620in}{1.312015in}}%
\pgfpathlineto{\pgfqpoint{3.142620in}{1.289105in}}%
\pgfpathclose%
\pgfusepath{fill}%
\end{pgfscope}%
\begin{pgfscope}%
\pgfpathrectangle{\pgfqpoint{2.893540in}{0.169444in}}{\pgfqpoint{0.858895in}{1.282944in}}%
\pgfusepath{clip}%
\pgfsetbuttcap%
\pgfsetmiterjoin%
\definecolor{currentfill}{rgb}{0.778316,0.127105,0.501499}%
\pgfsetfillcolor{currentfill}%
\pgfsetlinewidth{0.000000pt}%
\definecolor{currentstroke}{rgb}{0.000000,0.000000,0.000000}%
\pgfsetstrokecolor{currentstroke}%
\pgfsetstrokeopacity{0.000000}%
\pgfsetdash{}{0pt}%
\pgfpathmoveto{\pgfqpoint{3.142620in}{1.312431in}}%
\pgfpathlineto{\pgfqpoint{3.108264in}{1.312431in}}%
\pgfpathlineto{\pgfqpoint{3.108264in}{1.335341in}}%
\pgfpathlineto{\pgfqpoint{3.142620in}{1.335341in}}%
\pgfpathlineto{\pgfqpoint{3.142620in}{1.312431in}}%
\pgfpathclose%
\pgfusepath{fill}%
\end{pgfscope}%
\begin{pgfscope}%
\pgfpathrectangle{\pgfqpoint{2.893540in}{0.169444in}}{\pgfqpoint{0.858895in}{1.282944in}}%
\pgfusepath{clip}%
\pgfsetbuttcap%
\pgfsetmiterjoin%
\definecolor{currentfill}{rgb}{0.751403,0.095886,0.473664}%
\pgfsetfillcolor{currentfill}%
\pgfsetlinewidth{0.000000pt}%
\definecolor{currentstroke}{rgb}{0.000000,0.000000,0.000000}%
\pgfsetstrokecolor{currentstroke}%
\pgfsetstrokeopacity{0.000000}%
\pgfsetdash{}{0pt}%
\pgfpathmoveto{\pgfqpoint{3.142620in}{1.335758in}}%
\pgfpathlineto{\pgfqpoint{3.108264in}{1.335758in}}%
\pgfpathlineto{\pgfqpoint{3.108264in}{1.358667in}}%
\pgfpathlineto{\pgfqpoint{3.142620in}{1.358667in}}%
\pgfpathlineto{\pgfqpoint{3.142620in}{1.335758in}}%
\pgfpathclose%
\pgfusepath{fill}%
\end{pgfscope}%
\begin{pgfscope}%
\pgfpathrectangle{\pgfqpoint{2.893540in}{0.169444in}}{\pgfqpoint{0.858895in}{1.282944in}}%
\pgfusepath{clip}%
\pgfsetbuttcap%
\pgfsetmiterjoin%
\definecolor{currentfill}{rgb}{0.709112,0.075894,0.440600}%
\pgfsetfillcolor{currentfill}%
\pgfsetlinewidth{0.000000pt}%
\definecolor{currentstroke}{rgb}{0.000000,0.000000,0.000000}%
\pgfsetstrokecolor{currentstroke}%
\pgfsetstrokeopacity{0.000000}%
\pgfsetdash{}{0pt}%
\pgfpathmoveto{\pgfqpoint{3.142620in}{1.359084in}}%
\pgfpathlineto{\pgfqpoint{3.108264in}{1.359084in}}%
\pgfpathlineto{\pgfqpoint{3.108264in}{1.381994in}}%
\pgfpathlineto{\pgfqpoint{3.142620in}{1.381994in}}%
\pgfpathlineto{\pgfqpoint{3.142620in}{1.359084in}}%
\pgfpathclose%
\pgfusepath{fill}%
\end{pgfscope}%
\begin{pgfscope}%
\pgfpathrectangle{\pgfqpoint{2.893540in}{0.169444in}}{\pgfqpoint{0.858895in}{1.282944in}}%
\pgfusepath{clip}%
\pgfsetbuttcap%
\pgfsetmiterjoin%
\definecolor{currentfill}{rgb}{0.666820,0.055902,0.407536}%
\pgfsetfillcolor{currentfill}%
\pgfsetlinewidth{0.000000pt}%
\definecolor{currentstroke}{rgb}{0.000000,0.000000,0.000000}%
\pgfsetstrokecolor{currentstroke}%
\pgfsetstrokeopacity{0.000000}%
\pgfsetdash{}{0pt}%
\pgfpathmoveto{\pgfqpoint{3.142620in}{1.382410in}}%
\pgfpathlineto{\pgfqpoint{3.108264in}{1.382410in}}%
\pgfpathlineto{\pgfqpoint{3.108264in}{1.405320in}}%
\pgfpathlineto{\pgfqpoint{3.142620in}{1.405320in}}%
\pgfpathlineto{\pgfqpoint{3.142620in}{1.382410in}}%
\pgfpathclose%
\pgfusepath{fill}%
\end{pgfscope}%
\begin{pgfscope}%
\pgfpathrectangle{\pgfqpoint{2.893540in}{0.169444in}}{\pgfqpoint{0.858895in}{1.282944in}}%
\pgfusepath{clip}%
\pgfsetbuttcap%
\pgfsetmiterjoin%
\definecolor{currentfill}{rgb}{0.632987,0.039908,0.381084}%
\pgfsetfillcolor{currentfill}%
\pgfsetlinewidth{0.000000pt}%
\definecolor{currentstroke}{rgb}{0.000000,0.000000,0.000000}%
\pgfsetstrokecolor{currentstroke}%
\pgfsetstrokeopacity{0.000000}%
\pgfsetdash{}{0pt}%
\pgfpathmoveto{\pgfqpoint{3.142620in}{1.405736in}}%
\pgfpathlineto{\pgfqpoint{3.108264in}{1.405736in}}%
\pgfpathlineto{\pgfqpoint{3.108264in}{1.428646in}}%
\pgfpathlineto{\pgfqpoint{3.142620in}{1.428646in}}%
\pgfpathlineto{\pgfqpoint{3.142620in}{1.405736in}}%
\pgfpathclose%
\pgfusepath{fill}%
\end{pgfscope}%
\begin{pgfscope}%
\pgfpathrectangle{\pgfqpoint{2.893540in}{0.169444in}}{\pgfqpoint{0.858895in}{1.282944in}}%
\pgfusepath{clip}%
\pgfsetbuttcap%
\pgfsetmiterjoin%
\definecolor{currentfill}{rgb}{0.590696,0.019915,0.348020}%
\pgfsetfillcolor{currentfill}%
\pgfsetlinewidth{0.000000pt}%
\definecolor{currentstroke}{rgb}{0.000000,0.000000,0.000000}%
\pgfsetstrokecolor{currentstroke}%
\pgfsetstrokeopacity{0.000000}%
\pgfsetdash{}{0pt}%
\pgfpathmoveto{\pgfqpoint{3.142620in}{1.429063in}}%
\pgfpathlineto{\pgfqpoint{3.108264in}{1.429063in}}%
\pgfpathlineto{\pgfqpoint{3.108264in}{1.451972in}}%
\pgfpathlineto{\pgfqpoint{3.142620in}{1.451972in}}%
\pgfpathlineto{\pgfqpoint{3.142620in}{1.429063in}}%
\pgfpathclose%
\pgfusepath{fill}%
\end{pgfscope}%
\begin{pgfscope}%
\pgfpathrectangle{\pgfqpoint{2.893540in}{0.169444in}}{\pgfqpoint{0.858895in}{1.282944in}}%
\pgfusepath{clip}%
\pgfsetbuttcap%
\pgfsetmiterjoin%
\definecolor{currentfill}{rgb}{0.556863,0.003922,0.321569}%
\pgfsetfillcolor{currentfill}%
\pgfsetlinewidth{0.000000pt}%
\definecolor{currentstroke}{rgb}{0.000000,0.000000,0.000000}%
\pgfsetstrokecolor{currentstroke}%
\pgfsetstrokeopacity{0.000000}%
\pgfsetdash{}{0pt}%
\pgfpathmoveto{\pgfqpoint{3.142620in}{1.452389in}}%
\pgfpathlineto{\pgfqpoint{3.108264in}{1.452389in}}%
\pgfpathlineto{\pgfqpoint{3.108264in}{1.475299in}}%
\pgfpathlineto{\pgfqpoint{3.142620in}{1.475299in}}%
\pgfpathlineto{\pgfqpoint{3.142620in}{1.452389in}}%
\pgfpathclose%
\pgfusepath{fill}%
\end{pgfscope}%
\begin{pgfscope}%
\pgfpathrectangle{\pgfqpoint{2.893540in}{0.169444in}}{\pgfqpoint{0.858895in}{1.282944in}}%
\pgfusepath{clip}%
\pgfsetbuttcap%
\pgfsetmiterjoin%
\definecolor{currentfill}{rgb}{0.205536,0.455825,0.109112}%
\pgfsetfillcolor{currentfill}%
\pgfsetlinewidth{0.000000pt}%
\definecolor{currentstroke}{rgb}{0.000000,0.000000,0.000000}%
\pgfsetstrokecolor{currentstroke}%
\pgfsetstrokeopacity{0.000000}%
\pgfsetdash{}{0pt}%
\pgfpathmoveto{\pgfqpoint{3.048142in}{0.217911in}}%
\pgfpathlineto{\pgfqpoint{3.065319in}{0.217911in}}%
\pgfpathlineto{\pgfqpoint{3.065319in}{0.241669in}}%
\pgfpathlineto{\pgfqpoint{3.048142in}{0.241669in}}%
\pgfpathlineto{\pgfqpoint{3.048142in}{0.217911in}}%
\pgfpathclose%
\pgfusepath{fill}%
\end{pgfscope}%
\begin{pgfscope}%
\pgfpathrectangle{\pgfqpoint{2.893540in}{0.169444in}}{\pgfqpoint{0.858895in}{1.282944in}}%
\pgfusepath{clip}%
\pgfsetbuttcap%
\pgfsetmiterjoin%
\definecolor{currentfill}{rgb}{0.263975,0.526567,0.121415}%
\pgfsetfillcolor{currentfill}%
\pgfsetlinewidth{0.000000pt}%
\definecolor{currentstroke}{rgb}{0.000000,0.000000,0.000000}%
\pgfsetstrokecolor{currentstroke}%
\pgfsetstrokeopacity{0.000000}%
\pgfsetdash{}{0pt}%
\pgfpathmoveto{\pgfqpoint{3.030964in}{0.266235in}}%
\pgfpathlineto{\pgfqpoint{3.065319in}{0.266235in}}%
\pgfpathlineto{\pgfqpoint{3.065319in}{0.289994in}}%
\pgfpathlineto{\pgfqpoint{3.030964in}{0.289994in}}%
\pgfpathlineto{\pgfqpoint{3.030964in}{0.266235in}}%
\pgfpathclose%
\pgfusepath{fill}%
\end{pgfscope}%
\begin{pgfscope}%
\pgfpathrectangle{\pgfqpoint{2.893540in}{0.169444in}}{\pgfqpoint{0.858895in}{1.282944in}}%
\pgfusepath{clip}%
\pgfsetbuttcap%
\pgfsetmiterjoin%
\definecolor{currentfill}{rgb}{0.293195,0.561938,0.127566}%
\pgfsetfillcolor{currentfill}%
\pgfsetlinewidth{0.000000pt}%
\definecolor{currentstroke}{rgb}{0.000000,0.000000,0.000000}%
\pgfsetstrokecolor{currentstroke}%
\pgfsetstrokeopacity{0.000000}%
\pgfsetdash{}{0pt}%
\pgfpathmoveto{\pgfqpoint{3.013786in}{0.290469in}}%
\pgfpathlineto{\pgfqpoint{3.065319in}{0.290469in}}%
\pgfpathlineto{\pgfqpoint{3.065319in}{0.314227in}}%
\pgfpathlineto{\pgfqpoint{3.013786in}{0.314227in}}%
\pgfpathlineto{\pgfqpoint{3.013786in}{0.290469in}}%
\pgfpathclose%
\pgfusepath{fill}%
\end{pgfscope}%
\begin{pgfscope}%
\pgfpathrectangle{\pgfqpoint{2.893540in}{0.169444in}}{\pgfqpoint{0.858895in}{1.282944in}}%
\pgfusepath{clip}%
\pgfsetbuttcap%
\pgfsetmiterjoin%
\definecolor{currentfill}{rgb}{0.321184,0.588697,0.141715}%
\pgfsetfillcolor{currentfill}%
\pgfsetlinewidth{0.000000pt}%
\definecolor{currentstroke}{rgb}{0.000000,0.000000,0.000000}%
\pgfsetstrokecolor{currentstroke}%
\pgfsetstrokeopacity{0.000000}%
\pgfsetdash{}{0pt}%
\pgfpathmoveto{\pgfqpoint{3.048142in}{0.314702in}}%
\pgfpathlineto{\pgfqpoint{3.065319in}{0.314702in}}%
\pgfpathlineto{\pgfqpoint{3.065319in}{0.338460in}}%
\pgfpathlineto{\pgfqpoint{3.048142in}{0.338460in}}%
\pgfpathlineto{\pgfqpoint{3.048142in}{0.314702in}}%
\pgfpathclose%
\pgfusepath{fill}%
\end{pgfscope}%
\begin{pgfscope}%
\pgfpathrectangle{\pgfqpoint{2.893540in}{0.169444in}}{\pgfqpoint{0.858895in}{1.282944in}}%
\pgfusepath{clip}%
\pgfsetbuttcap%
\pgfsetmiterjoin%
\definecolor{currentfill}{rgb}{0.398078,0.653287,0.190927}%
\pgfsetfillcolor{currentfill}%
\pgfsetlinewidth{0.000000pt}%
\definecolor{currentstroke}{rgb}{0.000000,0.000000,0.000000}%
\pgfsetstrokecolor{currentstroke}%
\pgfsetstrokeopacity{0.000000}%
\pgfsetdash{}{0pt}%
\pgfpathmoveto{\pgfqpoint{3.030964in}{0.363027in}}%
\pgfpathlineto{\pgfqpoint{3.065319in}{0.363027in}}%
\pgfpathlineto{\pgfqpoint{3.065319in}{0.386785in}}%
\pgfpathlineto{\pgfqpoint{3.030964in}{0.386785in}}%
\pgfpathlineto{\pgfqpoint{3.030964in}{0.363027in}}%
\pgfpathclose%
\pgfusepath{fill}%
\end{pgfscope}%
\begin{pgfscope}%
\pgfpathrectangle{\pgfqpoint{2.893540in}{0.169444in}}{\pgfqpoint{0.858895in}{1.282944in}}%
\pgfusepath{clip}%
\pgfsetbuttcap%
\pgfsetmiterjoin%
\definecolor{currentfill}{rgb}{0.436524,0.685582,0.215532}%
\pgfsetfillcolor{currentfill}%
\pgfsetlinewidth{0.000000pt}%
\definecolor{currentstroke}{rgb}{0.000000,0.000000,0.000000}%
\pgfsetstrokecolor{currentstroke}%
\pgfsetstrokeopacity{0.000000}%
\pgfsetdash{}{0pt}%
\pgfpathmoveto{\pgfqpoint{3.030964in}{0.387260in}}%
\pgfpathlineto{\pgfqpoint{3.065319in}{0.387260in}}%
\pgfpathlineto{\pgfqpoint{3.065319in}{0.411018in}}%
\pgfpathlineto{\pgfqpoint{3.030964in}{0.411018in}}%
\pgfpathlineto{\pgfqpoint{3.030964in}{0.387260in}}%
\pgfpathclose%
\pgfusepath{fill}%
\end{pgfscope}%
\begin{pgfscope}%
\pgfpathrectangle{\pgfqpoint{2.893540in}{0.169444in}}{\pgfqpoint{0.858895in}{1.282944in}}%
\pgfusepath{clip}%
\pgfsetbuttcap%
\pgfsetmiterjoin%
\definecolor{currentfill}{rgb}{0.474971,0.717878,0.240138}%
\pgfsetfillcolor{currentfill}%
\pgfsetlinewidth{0.000000pt}%
\definecolor{currentstroke}{rgb}{0.000000,0.000000,0.000000}%
\pgfsetstrokecolor{currentstroke}%
\pgfsetstrokeopacity{0.000000}%
\pgfsetdash{}{0pt}%
\pgfpathmoveto{\pgfqpoint{3.030964in}{0.411493in}}%
\pgfpathlineto{\pgfqpoint{3.065319in}{0.411493in}}%
\pgfpathlineto{\pgfqpoint{3.065319in}{0.435252in}}%
\pgfpathlineto{\pgfqpoint{3.030964in}{0.435252in}}%
\pgfpathlineto{\pgfqpoint{3.030964in}{0.411493in}}%
\pgfpathclose%
\pgfusepath{fill}%
\end{pgfscope}%
\begin{pgfscope}%
\pgfpathrectangle{\pgfqpoint{2.893540in}{0.169444in}}{\pgfqpoint{0.858895in}{1.282944in}}%
\pgfusepath{clip}%
\pgfsetbuttcap%
\pgfsetmiterjoin%
\definecolor{currentfill}{rgb}{0.515571,0.748635,0.276125}%
\pgfsetfillcolor{currentfill}%
\pgfsetlinewidth{0.000000pt}%
\definecolor{currentstroke}{rgb}{0.000000,0.000000,0.000000}%
\pgfsetstrokecolor{currentstroke}%
\pgfsetstrokeopacity{0.000000}%
\pgfsetdash{}{0pt}%
\pgfpathmoveto{\pgfqpoint{2.876362in}{0.435727in}}%
\pgfpathlineto{\pgfqpoint{3.065319in}{0.435727in}}%
\pgfpathlineto{\pgfqpoint{3.065319in}{0.459485in}}%
\pgfpathlineto{\pgfqpoint{2.876362in}{0.459485in}}%
\pgfpathlineto{\pgfqpoint{2.876362in}{0.435727in}}%
\pgfpathclose%
\pgfusepath{fill}%
\end{pgfscope}%
\begin{pgfscope}%
\pgfpathrectangle{\pgfqpoint{2.893540in}{0.169444in}}{\pgfqpoint{0.858895in}{1.282944in}}%
\pgfusepath{clip}%
\pgfsetbuttcap%
\pgfsetmiterjoin%
\definecolor{currentfill}{rgb}{0.550634,0.771396,0.318570}%
\pgfsetfillcolor{currentfill}%
\pgfsetlinewidth{0.000000pt}%
\definecolor{currentstroke}{rgb}{0.000000,0.000000,0.000000}%
\pgfsetstrokecolor{currentstroke}%
\pgfsetstrokeopacity{0.000000}%
\pgfsetdash{}{0pt}%
\pgfpathmoveto{\pgfqpoint{2.962252in}{0.459960in}}%
\pgfpathlineto{\pgfqpoint{3.065319in}{0.459960in}}%
\pgfpathlineto{\pgfqpoint{3.065319in}{0.483718in}}%
\pgfpathlineto{\pgfqpoint{2.962252in}{0.483718in}}%
\pgfpathlineto{\pgfqpoint{2.962252in}{0.459960in}}%
\pgfpathclose%
\pgfusepath{fill}%
\end{pgfscope}%
\begin{pgfscope}%
\pgfpathrectangle{\pgfqpoint{2.893540in}{0.169444in}}{\pgfqpoint{0.858895in}{1.282944in}}%
\pgfusepath{clip}%
\pgfsetbuttcap%
\pgfsetmiterjoin%
\definecolor{currentfill}{rgb}{0.638293,0.828297,0.424683}%
\pgfsetfillcolor{currentfill}%
\pgfsetlinewidth{0.000000pt}%
\definecolor{currentstroke}{rgb}{0.000000,0.000000,0.000000}%
\pgfsetstrokecolor{currentstroke}%
\pgfsetstrokeopacity{0.000000}%
\pgfsetdash{}{0pt}%
\pgfpathmoveto{\pgfqpoint{2.945074in}{0.508284in}}%
\pgfpathlineto{\pgfqpoint{3.065319in}{0.508284in}}%
\pgfpathlineto{\pgfqpoint{3.065319in}{0.532043in}}%
\pgfpathlineto{\pgfqpoint{2.945074in}{0.532043in}}%
\pgfpathlineto{\pgfqpoint{2.945074in}{0.508284in}}%
\pgfpathclose%
\pgfusepath{fill}%
\end{pgfscope}%
\begin{pgfscope}%
\pgfpathrectangle{\pgfqpoint{2.893540in}{0.169444in}}{\pgfqpoint{0.858895in}{1.282944in}}%
\pgfusepath{clip}%
\pgfsetbuttcap%
\pgfsetmiterjoin%
\definecolor{currentfill}{rgb}{0.682122,0.856747,0.477739}%
\pgfsetfillcolor{currentfill}%
\pgfsetlinewidth{0.000000pt}%
\definecolor{currentstroke}{rgb}{0.000000,0.000000,0.000000}%
\pgfsetstrokecolor{currentstroke}%
\pgfsetstrokeopacity{0.000000}%
\pgfsetdash{}{0pt}%
\pgfpathmoveto{\pgfqpoint{3.013786in}{0.532518in}}%
\pgfpathlineto{\pgfqpoint{3.065319in}{0.532518in}}%
\pgfpathlineto{\pgfqpoint{3.065319in}{0.556276in}}%
\pgfpathlineto{\pgfqpoint{3.013786in}{0.556276in}}%
\pgfpathlineto{\pgfqpoint{3.013786in}{0.532518in}}%
\pgfpathclose%
\pgfusepath{fill}%
\end{pgfscope}%
\begin{pgfscope}%
\pgfpathrectangle{\pgfqpoint{2.893540in}{0.169444in}}{\pgfqpoint{0.858895in}{1.282944in}}%
\pgfusepath{clip}%
\pgfsetbuttcap%
\pgfsetmiterjoin%
\definecolor{currentfill}{rgb}{0.725106,0.883891,0.531180}%
\pgfsetfillcolor{currentfill}%
\pgfsetlinewidth{0.000000pt}%
\definecolor{currentstroke}{rgb}{0.000000,0.000000,0.000000}%
\pgfsetstrokecolor{currentstroke}%
\pgfsetstrokeopacity{0.000000}%
\pgfsetdash{}{0pt}%
\pgfpathmoveto{\pgfqpoint{3.048142in}{0.556751in}}%
\pgfpathlineto{\pgfqpoint{3.065319in}{0.556751in}}%
\pgfpathlineto{\pgfqpoint{3.065319in}{0.580509in}}%
\pgfpathlineto{\pgfqpoint{3.048142in}{0.580509in}}%
\pgfpathlineto{\pgfqpoint{3.048142in}{0.556751in}}%
\pgfpathclose%
\pgfusepath{fill}%
\end{pgfscope}%
\begin{pgfscope}%
\pgfpathrectangle{\pgfqpoint{2.893540in}{0.169444in}}{\pgfqpoint{0.858895in}{1.282944in}}%
\pgfusepath{clip}%
\pgfsetbuttcap%
\pgfsetmiterjoin%
\definecolor{currentfill}{rgb}{0.760477,0.899270,0.588082}%
\pgfsetfillcolor{currentfill}%
\pgfsetlinewidth{0.000000pt}%
\definecolor{currentstroke}{rgb}{0.000000,0.000000,0.000000}%
\pgfsetstrokecolor{currentstroke}%
\pgfsetstrokeopacity{0.000000}%
\pgfsetdash{}{0pt}%
\pgfpathmoveto{\pgfqpoint{3.013786in}{0.580985in}}%
\pgfpathlineto{\pgfqpoint{3.065319in}{0.580985in}}%
\pgfpathlineto{\pgfqpoint{3.065319in}{0.604743in}}%
\pgfpathlineto{\pgfqpoint{3.013786in}{0.604743in}}%
\pgfpathlineto{\pgfqpoint{3.013786in}{0.580985in}}%
\pgfpathclose%
\pgfusepath{fill}%
\end{pgfscope}%
\begin{pgfscope}%
\pgfpathrectangle{\pgfqpoint{2.893540in}{0.169444in}}{\pgfqpoint{0.858895in}{1.282944in}}%
\pgfusepath{clip}%
\pgfsetbuttcap%
\pgfsetmiterjoin%
\definecolor{currentfill}{rgb}{0.788774,0.911572,0.633602}%
\pgfsetfillcolor{currentfill}%
\pgfsetlinewidth{0.000000pt}%
\definecolor{currentstroke}{rgb}{0.000000,0.000000,0.000000}%
\pgfsetstrokecolor{currentstroke}%
\pgfsetstrokeopacity{0.000000}%
\pgfsetdash{}{0pt}%
\pgfpathmoveto{\pgfqpoint{3.030964in}{0.605218in}}%
\pgfpathlineto{\pgfqpoint{3.065319in}{0.605218in}}%
\pgfpathlineto{\pgfqpoint{3.065319in}{0.628976in}}%
\pgfpathlineto{\pgfqpoint{3.030964in}{0.628976in}}%
\pgfpathlineto{\pgfqpoint{3.030964in}{0.605218in}}%
\pgfpathclose%
\pgfusepath{fill}%
\end{pgfscope}%
\begin{pgfscope}%
\pgfpathrectangle{\pgfqpoint{2.893540in}{0.169444in}}{\pgfqpoint{0.858895in}{1.282944in}}%
\pgfusepath{clip}%
\pgfsetbuttcap%
\pgfsetmiterjoin%
\definecolor{currentfill}{rgb}{0.824145,0.926951,0.690504}%
\pgfsetfillcolor{currentfill}%
\pgfsetlinewidth{0.000000pt}%
\definecolor{currentstroke}{rgb}{0.000000,0.000000,0.000000}%
\pgfsetstrokecolor{currentstroke}%
\pgfsetstrokeopacity{0.000000}%
\pgfsetdash{}{0pt}%
\pgfpathmoveto{\pgfqpoint{3.048142in}{0.629309in}}%
\pgfpathlineto{\pgfqpoint{3.065319in}{0.629309in}}%
\pgfpathlineto{\pgfqpoint{3.065319in}{0.653067in}}%
\pgfpathlineto{\pgfqpoint{3.048142in}{0.653067in}}%
\pgfpathlineto{\pgfqpoint{3.048142in}{0.629309in}}%
\pgfpathclose%
\pgfusepath{fill}%
\end{pgfscope}%
\begin{pgfscope}%
\pgfpathrectangle{\pgfqpoint{2.893540in}{0.169444in}}{\pgfqpoint{0.858895in}{1.282944in}}%
\pgfusepath{clip}%
\pgfsetbuttcap%
\pgfsetmiterjoin%
\definecolor{currentfill}{rgb}{0.859516,0.942330,0.747405}%
\pgfsetfillcolor{currentfill}%
\pgfsetlinewidth{0.000000pt}%
\definecolor{currentstroke}{rgb}{0.000000,0.000000,0.000000}%
\pgfsetstrokecolor{currentstroke}%
\pgfsetstrokeopacity{0.000000}%
\pgfsetdash{}{0pt}%
\pgfpathmoveto{\pgfqpoint{3.048142in}{0.653542in}}%
\pgfpathlineto{\pgfqpoint{3.065319in}{0.653542in}}%
\pgfpathlineto{\pgfqpoint{3.065319in}{0.677300in}}%
\pgfpathlineto{\pgfqpoint{3.048142in}{0.677300in}}%
\pgfpathlineto{\pgfqpoint{3.048142in}{0.653542in}}%
\pgfpathclose%
\pgfusepath{fill}%
\end{pgfscope}%
\begin{pgfscope}%
\pgfpathrectangle{\pgfqpoint{2.893540in}{0.169444in}}{\pgfqpoint{0.858895in}{1.282944in}}%
\pgfusepath{clip}%
\pgfsetbuttcap%
\pgfsetmiterjoin%
\definecolor{currentfill}{rgb}{0.912418,0.962015,0.839677}%
\pgfsetfillcolor{currentfill}%
\pgfsetlinewidth{0.000000pt}%
\definecolor{currentstroke}{rgb}{0.000000,0.000000,0.000000}%
\pgfsetstrokecolor{currentstroke}%
\pgfsetstrokeopacity{0.000000}%
\pgfsetdash{}{0pt}%
\pgfpathmoveto{\pgfqpoint{3.013786in}{0.702009in}}%
\pgfpathlineto{\pgfqpoint{3.065319in}{0.702009in}}%
\pgfpathlineto{\pgfqpoint{3.065319in}{0.725767in}}%
\pgfpathlineto{\pgfqpoint{3.013786in}{0.725767in}}%
\pgfpathlineto{\pgfqpoint{3.013786in}{0.702009in}}%
\pgfpathclose%
\pgfusepath{fill}%
\end{pgfscope}%
\begin{pgfscope}%
\pgfpathrectangle{\pgfqpoint{2.893540in}{0.169444in}}{\pgfqpoint{0.858895in}{1.282944in}}%
\pgfusepath{clip}%
\pgfsetbuttcap%
\pgfsetmiterjoin%
\definecolor{currentfill}{rgb}{0.152941,0.392157,0.098039}%
\pgfsetfillcolor{currentfill}%
\pgfsetlinewidth{0.000000pt}%
\definecolor{currentstroke}{rgb}{0.000000,0.000000,0.000000}%
\pgfsetstrokecolor{currentstroke}%
\pgfsetstrokeopacity{0.000000}%
\pgfsetdash{}{0pt}%
\pgfpathmoveto{\pgfqpoint{3.374522in}{0.169444in}}%
\pgfpathlineto{\pgfqpoint{3.340166in}{0.169444in}}%
\pgfpathlineto{\pgfqpoint{3.340166in}{0.192354in}}%
\pgfpathlineto{\pgfqpoint{3.374522in}{0.192354in}}%
\pgfpathlineto{\pgfqpoint{3.374522in}{0.169444in}}%
\pgfpathclose%
\pgfusepath{fill}%
\end{pgfscope}%
\begin{pgfscope}%
\pgfpathrectangle{\pgfqpoint{2.893540in}{0.169444in}}{\pgfqpoint{0.858895in}{1.282944in}}%
\pgfusepath{clip}%
\pgfsetbuttcap%
\pgfsetmiterjoin%
\definecolor{currentfill}{rgb}{0.176317,0.420454,0.102960}%
\pgfsetfillcolor{currentfill}%
\pgfsetlinewidth{0.000000pt}%
\definecolor{currentstroke}{rgb}{0.000000,0.000000,0.000000}%
\pgfsetstrokecolor{currentstroke}%
\pgfsetstrokeopacity{0.000000}%
\pgfsetdash{}{0pt}%
\pgfpathmoveto{\pgfqpoint{3.374522in}{0.192771in}}%
\pgfpathlineto{\pgfqpoint{3.340166in}{0.192771in}}%
\pgfpathlineto{\pgfqpoint{3.340166in}{0.215680in}}%
\pgfpathlineto{\pgfqpoint{3.374522in}{0.215680in}}%
\pgfpathlineto{\pgfqpoint{3.374522in}{0.192771in}}%
\pgfpathclose%
\pgfusepath{fill}%
\end{pgfscope}%
\begin{pgfscope}%
\pgfpathrectangle{\pgfqpoint{2.893540in}{0.169444in}}{\pgfqpoint{0.858895in}{1.282944in}}%
\pgfusepath{clip}%
\pgfsetbuttcap%
\pgfsetmiterjoin%
\definecolor{currentfill}{rgb}{0.205536,0.455825,0.109112}%
\pgfsetfillcolor{currentfill}%
\pgfsetlinewidth{0.000000pt}%
\definecolor{currentstroke}{rgb}{0.000000,0.000000,0.000000}%
\pgfsetstrokecolor{currentstroke}%
\pgfsetstrokeopacity{0.000000}%
\pgfsetdash{}{0pt}%
\pgfpathmoveto{\pgfqpoint{3.374522in}{0.216097in}}%
\pgfpathlineto{\pgfqpoint{3.340166in}{0.216097in}}%
\pgfpathlineto{\pgfqpoint{3.340166in}{0.239007in}}%
\pgfpathlineto{\pgfqpoint{3.374522in}{0.239007in}}%
\pgfpathlineto{\pgfqpoint{3.374522in}{0.216097in}}%
\pgfpathclose%
\pgfusepath{fill}%
\end{pgfscope}%
\begin{pgfscope}%
\pgfpathrectangle{\pgfqpoint{2.893540in}{0.169444in}}{\pgfqpoint{0.858895in}{1.282944in}}%
\pgfusepath{clip}%
\pgfsetbuttcap%
\pgfsetmiterjoin%
\definecolor{currentfill}{rgb}{0.228912,0.484121,0.114033}%
\pgfsetfillcolor{currentfill}%
\pgfsetlinewidth{0.000000pt}%
\definecolor{currentstroke}{rgb}{0.000000,0.000000,0.000000}%
\pgfsetstrokecolor{currentstroke}%
\pgfsetstrokeopacity{0.000000}%
\pgfsetdash{}{0pt}%
\pgfpathmoveto{\pgfqpoint{3.374522in}{0.239423in}}%
\pgfpathlineto{\pgfqpoint{3.340166in}{0.239423in}}%
\pgfpathlineto{\pgfqpoint{3.340166in}{0.262333in}}%
\pgfpathlineto{\pgfqpoint{3.374522in}{0.262333in}}%
\pgfpathlineto{\pgfqpoint{3.374522in}{0.239423in}}%
\pgfpathclose%
\pgfusepath{fill}%
\end{pgfscope}%
\begin{pgfscope}%
\pgfpathrectangle{\pgfqpoint{2.893540in}{0.169444in}}{\pgfqpoint{0.858895in}{1.282944in}}%
\pgfusepath{clip}%
\pgfsetbuttcap%
\pgfsetmiterjoin%
\definecolor{currentfill}{rgb}{0.258131,0.519493,0.120185}%
\pgfsetfillcolor{currentfill}%
\pgfsetlinewidth{0.000000pt}%
\definecolor{currentstroke}{rgb}{0.000000,0.000000,0.000000}%
\pgfsetstrokecolor{currentstroke}%
\pgfsetstrokeopacity{0.000000}%
\pgfsetdash{}{0pt}%
\pgfpathmoveto{\pgfqpoint{3.374522in}{0.262749in}}%
\pgfpathlineto{\pgfqpoint{3.340166in}{0.262749in}}%
\pgfpathlineto{\pgfqpoint{3.340166in}{0.285659in}}%
\pgfpathlineto{\pgfqpoint{3.374522in}{0.285659in}}%
\pgfpathlineto{\pgfqpoint{3.374522in}{0.262749in}}%
\pgfpathclose%
\pgfusepath{fill}%
\end{pgfscope}%
\begin{pgfscope}%
\pgfpathrectangle{\pgfqpoint{2.893540in}{0.169444in}}{\pgfqpoint{0.858895in}{1.282944in}}%
\pgfusepath{clip}%
\pgfsetbuttcap%
\pgfsetmiterjoin%
\definecolor{currentfill}{rgb}{0.287351,0.554864,0.126336}%
\pgfsetfillcolor{currentfill}%
\pgfsetlinewidth{0.000000pt}%
\definecolor{currentstroke}{rgb}{0.000000,0.000000,0.000000}%
\pgfsetstrokecolor{currentstroke}%
\pgfsetstrokeopacity{0.000000}%
\pgfsetdash{}{0pt}%
\pgfpathmoveto{\pgfqpoint{3.374522in}{0.286076in}}%
\pgfpathlineto{\pgfqpoint{3.340166in}{0.286076in}}%
\pgfpathlineto{\pgfqpoint{3.340166in}{0.308985in}}%
\pgfpathlineto{\pgfqpoint{3.374522in}{0.308985in}}%
\pgfpathlineto{\pgfqpoint{3.374522in}{0.286076in}}%
\pgfpathclose%
\pgfusepath{fill}%
\end{pgfscope}%
\begin{pgfscope}%
\pgfpathrectangle{\pgfqpoint{2.893540in}{0.169444in}}{\pgfqpoint{0.858895in}{1.282944in}}%
\pgfusepath{clip}%
\pgfsetbuttcap%
\pgfsetmiterjoin%
\definecolor{currentfill}{rgb}{0.313495,0.582238,0.136794}%
\pgfsetfillcolor{currentfill}%
\pgfsetlinewidth{0.000000pt}%
\definecolor{currentstroke}{rgb}{0.000000,0.000000,0.000000}%
\pgfsetstrokecolor{currentstroke}%
\pgfsetstrokeopacity{0.000000}%
\pgfsetdash{}{0pt}%
\pgfpathmoveto{\pgfqpoint{3.374522in}{0.309402in}}%
\pgfpathlineto{\pgfqpoint{3.340166in}{0.309402in}}%
\pgfpathlineto{\pgfqpoint{3.340166in}{0.332312in}}%
\pgfpathlineto{\pgfqpoint{3.374522in}{0.332312in}}%
\pgfpathlineto{\pgfqpoint{3.374522in}{0.309402in}}%
\pgfpathclose%
\pgfusepath{fill}%
\end{pgfscope}%
\begin{pgfscope}%
\pgfpathrectangle{\pgfqpoint{2.893540in}{0.169444in}}{\pgfqpoint{0.858895in}{1.282944in}}%
\pgfusepath{clip}%
\pgfsetbuttcap%
\pgfsetmiterjoin%
\definecolor{currentfill}{rgb}{0.351942,0.614533,0.161399}%
\pgfsetfillcolor{currentfill}%
\pgfsetlinewidth{0.000000pt}%
\definecolor{currentstroke}{rgb}{0.000000,0.000000,0.000000}%
\pgfsetstrokecolor{currentstroke}%
\pgfsetstrokeopacity{0.000000}%
\pgfsetdash{}{0pt}%
\pgfpathmoveto{\pgfqpoint{3.374522in}{0.332728in}}%
\pgfpathlineto{\pgfqpoint{3.340166in}{0.332728in}}%
\pgfpathlineto{\pgfqpoint{3.340166in}{0.355638in}}%
\pgfpathlineto{\pgfqpoint{3.374522in}{0.355638in}}%
\pgfpathlineto{\pgfqpoint{3.374522in}{0.332728in}}%
\pgfpathclose%
\pgfusepath{fill}%
\end{pgfscope}%
\begin{pgfscope}%
\pgfpathrectangle{\pgfqpoint{2.893540in}{0.169444in}}{\pgfqpoint{0.858895in}{1.282944in}}%
\pgfusepath{clip}%
\pgfsetbuttcap%
\pgfsetmiterjoin%
\definecolor{currentfill}{rgb}{0.390388,0.646828,0.186005}%
\pgfsetfillcolor{currentfill}%
\pgfsetlinewidth{0.000000pt}%
\definecolor{currentstroke}{rgb}{0.000000,0.000000,0.000000}%
\pgfsetstrokecolor{currentstroke}%
\pgfsetstrokeopacity{0.000000}%
\pgfsetdash{}{0pt}%
\pgfpathmoveto{\pgfqpoint{3.374522in}{0.356055in}}%
\pgfpathlineto{\pgfqpoint{3.340166in}{0.356055in}}%
\pgfpathlineto{\pgfqpoint{3.340166in}{0.378964in}}%
\pgfpathlineto{\pgfqpoint{3.374522in}{0.378964in}}%
\pgfpathlineto{\pgfqpoint{3.374522in}{0.356055in}}%
\pgfpathclose%
\pgfusepath{fill}%
\end{pgfscope}%
\begin{pgfscope}%
\pgfpathrectangle{\pgfqpoint{2.893540in}{0.169444in}}{\pgfqpoint{0.858895in}{1.282944in}}%
\pgfusepath{clip}%
\pgfsetbuttcap%
\pgfsetmiterjoin%
\definecolor{currentfill}{rgb}{0.421146,0.672664,0.205690}%
\pgfsetfillcolor{currentfill}%
\pgfsetlinewidth{0.000000pt}%
\definecolor{currentstroke}{rgb}{0.000000,0.000000,0.000000}%
\pgfsetstrokecolor{currentstroke}%
\pgfsetstrokeopacity{0.000000}%
\pgfsetdash{}{0pt}%
\pgfpathmoveto{\pgfqpoint{3.374522in}{0.379381in}}%
\pgfpathlineto{\pgfqpoint{3.340166in}{0.379381in}}%
\pgfpathlineto{\pgfqpoint{3.340166in}{0.402291in}}%
\pgfpathlineto{\pgfqpoint{3.374522in}{0.402291in}}%
\pgfpathlineto{\pgfqpoint{3.374522in}{0.379381in}}%
\pgfpathclose%
\pgfusepath{fill}%
\end{pgfscope}%
\begin{pgfscope}%
\pgfpathrectangle{\pgfqpoint{2.893540in}{0.169444in}}{\pgfqpoint{0.858895in}{1.282944in}}%
\pgfusepath{clip}%
\pgfsetbuttcap%
\pgfsetmiterjoin%
\definecolor{currentfill}{rgb}{0.459592,0.704960,0.230296}%
\pgfsetfillcolor{currentfill}%
\pgfsetlinewidth{0.000000pt}%
\definecolor{currentstroke}{rgb}{0.000000,0.000000,0.000000}%
\pgfsetstrokecolor{currentstroke}%
\pgfsetstrokeopacity{0.000000}%
\pgfsetdash{}{0pt}%
\pgfpathmoveto{\pgfqpoint{3.374522in}{0.402707in}}%
\pgfpathlineto{\pgfqpoint{3.340166in}{0.402707in}}%
\pgfpathlineto{\pgfqpoint{3.340166in}{0.425617in}}%
\pgfpathlineto{\pgfqpoint{3.374522in}{0.425617in}}%
\pgfpathlineto{\pgfqpoint{3.374522in}{0.402707in}}%
\pgfpathclose%
\pgfusepath{fill}%
\end{pgfscope}%
\begin{pgfscope}%
\pgfpathrectangle{\pgfqpoint{2.893540in}{0.169444in}}{\pgfqpoint{0.858895in}{1.282944in}}%
\pgfusepath{clip}%
\pgfsetbuttcap%
\pgfsetmiterjoin%
\definecolor{currentfill}{rgb}{0.498039,0.737255,0.254902}%
\pgfsetfillcolor{currentfill}%
\pgfsetlinewidth{0.000000pt}%
\definecolor{currentstroke}{rgb}{0.000000,0.000000,0.000000}%
\pgfsetstrokecolor{currentstroke}%
\pgfsetstrokeopacity{0.000000}%
\pgfsetdash{}{0pt}%
\pgfpathmoveto{\pgfqpoint{3.374522in}{0.426033in}}%
\pgfpathlineto{\pgfqpoint{3.340166in}{0.426033in}}%
\pgfpathlineto{\pgfqpoint{3.340166in}{0.448943in}}%
\pgfpathlineto{\pgfqpoint{3.374522in}{0.448943in}}%
\pgfpathlineto{\pgfqpoint{3.374522in}{0.426033in}}%
\pgfpathclose%
\pgfusepath{fill}%
\end{pgfscope}%
\begin{pgfscope}%
\pgfpathrectangle{\pgfqpoint{2.893540in}{0.169444in}}{\pgfqpoint{0.858895in}{1.282944in}}%
\pgfusepath{clip}%
\pgfsetbuttcap%
\pgfsetmiterjoin%
\definecolor{currentfill}{rgb}{0.533103,0.760015,0.297347}%
\pgfsetfillcolor{currentfill}%
\pgfsetlinewidth{0.000000pt}%
\definecolor{currentstroke}{rgb}{0.000000,0.000000,0.000000}%
\pgfsetstrokecolor{currentstroke}%
\pgfsetstrokeopacity{0.000000}%
\pgfsetdash{}{0pt}%
\pgfpathmoveto{\pgfqpoint{3.374522in}{0.449360in}}%
\pgfpathlineto{\pgfqpoint{3.340166in}{0.449360in}}%
\pgfpathlineto{\pgfqpoint{3.340166in}{0.472269in}}%
\pgfpathlineto{\pgfqpoint{3.374522in}{0.472269in}}%
\pgfpathlineto{\pgfqpoint{3.374522in}{0.449360in}}%
\pgfpathclose%
\pgfusepath{fill}%
\end{pgfscope}%
\begin{pgfscope}%
\pgfpathrectangle{\pgfqpoint{2.893540in}{0.169444in}}{\pgfqpoint{0.858895in}{1.282944in}}%
\pgfusepath{clip}%
\pgfsetbuttcap%
\pgfsetmiterjoin%
\definecolor{currentfill}{rgb}{0.576932,0.788466,0.350404}%
\pgfsetfillcolor{currentfill}%
\pgfsetlinewidth{0.000000pt}%
\definecolor{currentstroke}{rgb}{0.000000,0.000000,0.000000}%
\pgfsetstrokecolor{currentstroke}%
\pgfsetstrokeopacity{0.000000}%
\pgfsetdash{}{0pt}%
\pgfpathmoveto{\pgfqpoint{3.374522in}{0.472686in}}%
\pgfpathlineto{\pgfqpoint{3.340166in}{0.472686in}}%
\pgfpathlineto{\pgfqpoint{3.340166in}{0.495596in}}%
\pgfpathlineto{\pgfqpoint{3.374522in}{0.495596in}}%
\pgfpathlineto{\pgfqpoint{3.374522in}{0.472686in}}%
\pgfpathclose%
\pgfusepath{fill}%
\end{pgfscope}%
\begin{pgfscope}%
\pgfpathrectangle{\pgfqpoint{2.893540in}{0.169444in}}{\pgfqpoint{0.858895in}{1.282944in}}%
\pgfusepath{clip}%
\pgfsetbuttcap%
\pgfsetmiterjoin%
\definecolor{currentfill}{rgb}{0.620761,0.816917,0.403460}%
\pgfsetfillcolor{currentfill}%
\pgfsetlinewidth{0.000000pt}%
\definecolor{currentstroke}{rgb}{0.000000,0.000000,0.000000}%
\pgfsetstrokecolor{currentstroke}%
\pgfsetstrokeopacity{0.000000}%
\pgfsetdash{}{0pt}%
\pgfpathmoveto{\pgfqpoint{3.374522in}{0.496012in}}%
\pgfpathlineto{\pgfqpoint{3.340166in}{0.496012in}}%
\pgfpathlineto{\pgfqpoint{3.340166in}{0.518922in}}%
\pgfpathlineto{\pgfqpoint{3.374522in}{0.518922in}}%
\pgfpathlineto{\pgfqpoint{3.374522in}{0.496012in}}%
\pgfpathclose%
\pgfusepath{fill}%
\end{pgfscope}%
\begin{pgfscope}%
\pgfpathrectangle{\pgfqpoint{2.893540in}{0.169444in}}{\pgfqpoint{0.858895in}{1.282944in}}%
\pgfusepath{clip}%
\pgfsetbuttcap%
\pgfsetmiterjoin%
\definecolor{currentfill}{rgb}{0.655825,0.839677,0.445905}%
\pgfsetfillcolor{currentfill}%
\pgfsetlinewidth{0.000000pt}%
\definecolor{currentstroke}{rgb}{0.000000,0.000000,0.000000}%
\pgfsetstrokecolor{currentstroke}%
\pgfsetstrokeopacity{0.000000}%
\pgfsetdash{}{0pt}%
\pgfpathmoveto{\pgfqpoint{3.374522in}{0.519338in}}%
\pgfpathlineto{\pgfqpoint{3.340166in}{0.519338in}}%
\pgfpathlineto{\pgfqpoint{3.340166in}{0.542248in}}%
\pgfpathlineto{\pgfqpoint{3.374522in}{0.542248in}}%
\pgfpathlineto{\pgfqpoint{3.374522in}{0.519338in}}%
\pgfpathclose%
\pgfusepath{fill}%
\end{pgfscope}%
\begin{pgfscope}%
\pgfpathrectangle{\pgfqpoint{2.893540in}{0.169444in}}{\pgfqpoint{0.858895in}{1.282944in}}%
\pgfusepath{clip}%
\pgfsetbuttcap%
\pgfsetmiterjoin%
\definecolor{currentfill}{rgb}{0.699654,0.868128,0.498962}%
\pgfsetfillcolor{currentfill}%
\pgfsetlinewidth{0.000000pt}%
\definecolor{currentstroke}{rgb}{0.000000,0.000000,0.000000}%
\pgfsetstrokecolor{currentstroke}%
\pgfsetstrokeopacity{0.000000}%
\pgfsetdash{}{0pt}%
\pgfpathmoveto{\pgfqpoint{3.374522in}{0.542665in}}%
\pgfpathlineto{\pgfqpoint{3.340166in}{0.542665in}}%
\pgfpathlineto{\pgfqpoint{3.340166in}{0.565574in}}%
\pgfpathlineto{\pgfqpoint{3.374522in}{0.565574in}}%
\pgfpathlineto{\pgfqpoint{3.374522in}{0.542665in}}%
\pgfpathclose%
\pgfusepath{fill}%
\end{pgfscope}%
\begin{pgfscope}%
\pgfpathrectangle{\pgfqpoint{2.893540in}{0.169444in}}{\pgfqpoint{0.858895in}{1.282944in}}%
\pgfusepath{clip}%
\pgfsetbuttcap%
\pgfsetmiterjoin%
\definecolor{currentfill}{rgb}{0.739254,0.890042,0.553941}%
\pgfsetfillcolor{currentfill}%
\pgfsetlinewidth{0.000000pt}%
\definecolor{currentstroke}{rgb}{0.000000,0.000000,0.000000}%
\pgfsetstrokecolor{currentstroke}%
\pgfsetstrokeopacity{0.000000}%
\pgfsetdash{}{0pt}%
\pgfpathmoveto{\pgfqpoint{3.374522in}{0.565991in}}%
\pgfpathlineto{\pgfqpoint{3.340166in}{0.565991in}}%
\pgfpathlineto{\pgfqpoint{3.340166in}{0.588901in}}%
\pgfpathlineto{\pgfqpoint{3.374522in}{0.588901in}}%
\pgfpathlineto{\pgfqpoint{3.374522in}{0.565991in}}%
\pgfpathclose%
\pgfusepath{fill}%
\end{pgfscope}%
\begin{pgfscope}%
\pgfpathrectangle{\pgfqpoint{2.893540in}{0.169444in}}{\pgfqpoint{0.858895in}{1.282944in}}%
\pgfusepath{clip}%
\pgfsetbuttcap%
\pgfsetmiterjoin%
\definecolor{currentfill}{rgb}{0.767551,0.902345,0.599462}%
\pgfsetfillcolor{currentfill}%
\pgfsetlinewidth{0.000000pt}%
\definecolor{currentstroke}{rgb}{0.000000,0.000000,0.000000}%
\pgfsetstrokecolor{currentstroke}%
\pgfsetstrokeopacity{0.000000}%
\pgfsetdash{}{0pt}%
\pgfpathmoveto{\pgfqpoint{3.374522in}{0.589317in}}%
\pgfpathlineto{\pgfqpoint{3.340166in}{0.589317in}}%
\pgfpathlineto{\pgfqpoint{3.340166in}{0.612227in}}%
\pgfpathlineto{\pgfqpoint{3.374522in}{0.612227in}}%
\pgfpathlineto{\pgfqpoint{3.374522in}{0.589317in}}%
\pgfpathclose%
\pgfusepath{fill}%
\end{pgfscope}%
\begin{pgfscope}%
\pgfpathrectangle{\pgfqpoint{2.893540in}{0.169444in}}{\pgfqpoint{0.858895in}{1.282944in}}%
\pgfusepath{clip}%
\pgfsetbuttcap%
\pgfsetmiterjoin%
\definecolor{currentfill}{rgb}{0.802922,0.917724,0.656363}%
\pgfsetfillcolor{currentfill}%
\pgfsetlinewidth{0.000000pt}%
\definecolor{currentstroke}{rgb}{0.000000,0.000000,0.000000}%
\pgfsetstrokecolor{currentstroke}%
\pgfsetstrokeopacity{0.000000}%
\pgfsetdash{}{0pt}%
\pgfpathmoveto{\pgfqpoint{3.374522in}{0.612643in}}%
\pgfpathlineto{\pgfqpoint{3.340166in}{0.612643in}}%
\pgfpathlineto{\pgfqpoint{3.340166in}{0.635553in}}%
\pgfpathlineto{\pgfqpoint{3.374522in}{0.635553in}}%
\pgfpathlineto{\pgfqpoint{3.374522in}{0.612643in}}%
\pgfpathclose%
\pgfusepath{fill}%
\end{pgfscope}%
\begin{pgfscope}%
\pgfpathrectangle{\pgfqpoint{2.893540in}{0.169444in}}{\pgfqpoint{0.858895in}{1.282944in}}%
\pgfusepath{clip}%
\pgfsetbuttcap%
\pgfsetmiterjoin%
\definecolor{currentfill}{rgb}{0.838293,0.933103,0.713264}%
\pgfsetfillcolor{currentfill}%
\pgfsetlinewidth{0.000000pt}%
\definecolor{currentstroke}{rgb}{0.000000,0.000000,0.000000}%
\pgfsetstrokecolor{currentstroke}%
\pgfsetstrokeopacity{0.000000}%
\pgfsetdash{}{0pt}%
\pgfpathmoveto{\pgfqpoint{3.374522in}{0.635970in}}%
\pgfpathlineto{\pgfqpoint{3.340166in}{0.635970in}}%
\pgfpathlineto{\pgfqpoint{3.340166in}{0.658879in}}%
\pgfpathlineto{\pgfqpoint{3.374522in}{0.658879in}}%
\pgfpathlineto{\pgfqpoint{3.374522in}{0.635970in}}%
\pgfpathclose%
\pgfusepath{fill}%
\end{pgfscope}%
\begin{pgfscope}%
\pgfpathrectangle{\pgfqpoint{2.893540in}{0.169444in}}{\pgfqpoint{0.858895in}{1.282944in}}%
\pgfusepath{clip}%
\pgfsetbuttcap%
\pgfsetmiterjoin%
\definecolor{currentfill}{rgb}{0.866590,0.945406,0.758785}%
\pgfsetfillcolor{currentfill}%
\pgfsetlinewidth{0.000000pt}%
\definecolor{currentstroke}{rgb}{0.000000,0.000000,0.000000}%
\pgfsetstrokecolor{currentstroke}%
\pgfsetstrokeopacity{0.000000}%
\pgfsetdash{}{0pt}%
\pgfpathmoveto{\pgfqpoint{3.374522in}{0.659296in}}%
\pgfpathlineto{\pgfqpoint{3.340166in}{0.659296in}}%
\pgfpathlineto{\pgfqpoint{3.340166in}{0.682206in}}%
\pgfpathlineto{\pgfqpoint{3.374522in}{0.682206in}}%
\pgfpathlineto{\pgfqpoint{3.374522in}{0.659296in}}%
\pgfpathclose%
\pgfusepath{fill}%
\end{pgfscope}%
\begin{pgfscope}%
\pgfpathrectangle{\pgfqpoint{2.893540in}{0.169444in}}{\pgfqpoint{0.858895in}{1.282944in}}%
\pgfusepath{clip}%
\pgfsetbuttcap%
\pgfsetmiterjoin%
\definecolor{currentfill}{rgb}{0.901961,0.960784,0.815686}%
\pgfsetfillcolor{currentfill}%
\pgfsetlinewidth{0.000000pt}%
\definecolor{currentstroke}{rgb}{0.000000,0.000000,0.000000}%
\pgfsetstrokecolor{currentstroke}%
\pgfsetstrokeopacity{0.000000}%
\pgfsetdash{}{0pt}%
\pgfpathmoveto{\pgfqpoint{3.374522in}{0.682622in}}%
\pgfpathlineto{\pgfqpoint{3.340166in}{0.682622in}}%
\pgfpathlineto{\pgfqpoint{3.340166in}{0.705532in}}%
\pgfpathlineto{\pgfqpoint{3.374522in}{0.705532in}}%
\pgfpathlineto{\pgfqpoint{3.374522in}{0.682622in}}%
\pgfpathclose%
\pgfusepath{fill}%
\end{pgfscope}%
\begin{pgfscope}%
\pgfpathrectangle{\pgfqpoint{2.893540in}{0.169444in}}{\pgfqpoint{0.858895in}{1.282944in}}%
\pgfusepath{clip}%
\pgfsetbuttcap%
\pgfsetmiterjoin%
\definecolor{currentfill}{rgb}{0.915033,0.962322,0.845675}%
\pgfsetfillcolor{currentfill}%
\pgfsetlinewidth{0.000000pt}%
\definecolor{currentstroke}{rgb}{0.000000,0.000000,0.000000}%
\pgfsetstrokecolor{currentstroke}%
\pgfsetstrokeopacity{0.000000}%
\pgfsetdash{}{0pt}%
\pgfpathmoveto{\pgfqpoint{3.374522in}{0.705948in}}%
\pgfpathlineto{\pgfqpoint{3.340166in}{0.705948in}}%
\pgfpathlineto{\pgfqpoint{3.340166in}{0.728858in}}%
\pgfpathlineto{\pgfqpoint{3.374522in}{0.728858in}}%
\pgfpathlineto{\pgfqpoint{3.374522in}{0.705948in}}%
\pgfpathclose%
\pgfusepath{fill}%
\end{pgfscope}%
\begin{pgfscope}%
\pgfpathrectangle{\pgfqpoint{2.893540in}{0.169444in}}{\pgfqpoint{0.858895in}{1.282944in}}%
\pgfusepath{clip}%
\pgfsetbuttcap%
\pgfsetmiterjoin%
\definecolor{currentfill}{rgb}{0.925490,0.963552,0.869666}%
\pgfsetfillcolor{currentfill}%
\pgfsetlinewidth{0.000000pt}%
\definecolor{currentstroke}{rgb}{0.000000,0.000000,0.000000}%
\pgfsetstrokecolor{currentstroke}%
\pgfsetstrokeopacity{0.000000}%
\pgfsetdash{}{0pt}%
\pgfpathmoveto{\pgfqpoint{3.374522in}{0.729275in}}%
\pgfpathlineto{\pgfqpoint{3.340166in}{0.729275in}}%
\pgfpathlineto{\pgfqpoint{3.340166in}{0.752184in}}%
\pgfpathlineto{\pgfqpoint{3.374522in}{0.752184in}}%
\pgfpathlineto{\pgfqpoint{3.374522in}{0.729275in}}%
\pgfpathclose%
\pgfusepath{fill}%
\end{pgfscope}%
\begin{pgfscope}%
\pgfpathrectangle{\pgfqpoint{2.893540in}{0.169444in}}{\pgfqpoint{0.858895in}{1.282944in}}%
\pgfusepath{clip}%
\pgfsetbuttcap%
\pgfsetmiterjoin%
\definecolor{currentfill}{rgb}{0.938562,0.965090,0.899654}%
\pgfsetfillcolor{currentfill}%
\pgfsetlinewidth{0.000000pt}%
\definecolor{currentstroke}{rgb}{0.000000,0.000000,0.000000}%
\pgfsetstrokecolor{currentstroke}%
\pgfsetstrokeopacity{0.000000}%
\pgfsetdash{}{0pt}%
\pgfpathmoveto{\pgfqpoint{3.374522in}{0.752601in}}%
\pgfpathlineto{\pgfqpoint{3.340166in}{0.752601in}}%
\pgfpathlineto{\pgfqpoint{3.340166in}{0.775511in}}%
\pgfpathlineto{\pgfqpoint{3.374522in}{0.775511in}}%
\pgfpathlineto{\pgfqpoint{3.374522in}{0.752601in}}%
\pgfpathclose%
\pgfusepath{fill}%
\end{pgfscope}%
\begin{pgfscope}%
\pgfpathrectangle{\pgfqpoint{2.893540in}{0.169444in}}{\pgfqpoint{0.858895in}{1.282944in}}%
\pgfusepath{clip}%
\pgfsetbuttcap%
\pgfsetmiterjoin%
\definecolor{currentfill}{rgb}{0.951634,0.966628,0.929642}%
\pgfsetfillcolor{currentfill}%
\pgfsetlinewidth{0.000000pt}%
\definecolor{currentstroke}{rgb}{0.000000,0.000000,0.000000}%
\pgfsetstrokecolor{currentstroke}%
\pgfsetstrokeopacity{0.000000}%
\pgfsetdash{}{0pt}%
\pgfpathmoveto{\pgfqpoint{3.374522in}{0.775927in}}%
\pgfpathlineto{\pgfqpoint{3.340166in}{0.775927in}}%
\pgfpathlineto{\pgfqpoint{3.340166in}{0.798837in}}%
\pgfpathlineto{\pgfqpoint{3.374522in}{0.798837in}}%
\pgfpathlineto{\pgfqpoint{3.374522in}{0.775927in}}%
\pgfpathclose%
\pgfusepath{fill}%
\end{pgfscope}%
\begin{pgfscope}%
\pgfpathrectangle{\pgfqpoint{2.893540in}{0.169444in}}{\pgfqpoint{0.858895in}{1.282944in}}%
\pgfusepath{clip}%
\pgfsetbuttcap%
\pgfsetmiterjoin%
\definecolor{currentfill}{rgb}{0.962092,0.967859,0.953633}%
\pgfsetfillcolor{currentfill}%
\pgfsetlinewidth{0.000000pt}%
\definecolor{currentstroke}{rgb}{0.000000,0.000000,0.000000}%
\pgfsetstrokecolor{currentstroke}%
\pgfsetstrokeopacity{0.000000}%
\pgfsetdash{}{0pt}%
\pgfpathmoveto{\pgfqpoint{3.374522in}{0.799254in}}%
\pgfpathlineto{\pgfqpoint{3.340166in}{0.799254in}}%
\pgfpathlineto{\pgfqpoint{3.340166in}{0.822163in}}%
\pgfpathlineto{\pgfqpoint{3.374522in}{0.822163in}}%
\pgfpathlineto{\pgfqpoint{3.374522in}{0.799254in}}%
\pgfpathclose%
\pgfusepath{fill}%
\end{pgfscope}%
\begin{pgfscope}%
\pgfpathrectangle{\pgfqpoint{2.893540in}{0.169444in}}{\pgfqpoint{0.858895in}{1.282944in}}%
\pgfusepath{clip}%
\pgfsetbuttcap%
\pgfsetmiterjoin%
\definecolor{currentfill}{rgb}{0.970934,0.959785,0.965552}%
\pgfsetfillcolor{currentfill}%
\pgfsetlinewidth{0.000000pt}%
\definecolor{currentstroke}{rgb}{0.000000,0.000000,0.000000}%
\pgfsetstrokecolor{currentstroke}%
\pgfsetstrokeopacity{0.000000}%
\pgfsetdash{}{0pt}%
\pgfpathmoveto{\pgfqpoint{3.374522in}{0.822580in}}%
\pgfpathlineto{\pgfqpoint{3.340166in}{0.822580in}}%
\pgfpathlineto{\pgfqpoint{3.340166in}{0.845490in}}%
\pgfpathlineto{\pgfqpoint{3.374522in}{0.845490in}}%
\pgfpathlineto{\pgfqpoint{3.374522in}{0.822580in}}%
\pgfpathclose%
\pgfusepath{fill}%
\end{pgfscope}%
\begin{pgfscope}%
\pgfpathrectangle{\pgfqpoint{2.893540in}{0.169444in}}{\pgfqpoint{0.858895in}{1.282944in}}%
\pgfusepath{clip}%
\pgfsetbuttcap%
\pgfsetmiterjoin%
\definecolor{currentfill}{rgb}{0.974625,0.945636,0.960631}%
\pgfsetfillcolor{currentfill}%
\pgfsetlinewidth{0.000000pt}%
\definecolor{currentstroke}{rgb}{0.000000,0.000000,0.000000}%
\pgfsetstrokecolor{currentstroke}%
\pgfsetstrokeopacity{0.000000}%
\pgfsetdash{}{0pt}%
\pgfpathmoveto{\pgfqpoint{3.374522in}{0.845906in}}%
\pgfpathlineto{\pgfqpoint{3.340166in}{0.845906in}}%
\pgfpathlineto{\pgfqpoint{3.340166in}{0.868816in}}%
\pgfpathlineto{\pgfqpoint{3.374522in}{0.868816in}}%
\pgfpathlineto{\pgfqpoint{3.374522in}{0.845906in}}%
\pgfpathclose%
\pgfusepath{fill}%
\end{pgfscope}%
\begin{pgfscope}%
\pgfpathrectangle{\pgfqpoint{2.893540in}{0.169444in}}{\pgfqpoint{0.858895in}{1.282944in}}%
\pgfusepath{clip}%
\pgfsetbuttcap%
\pgfsetmiterjoin%
\definecolor{currentfill}{rgb}{0.979239,0.927951,0.954479}%
\pgfsetfillcolor{currentfill}%
\pgfsetlinewidth{0.000000pt}%
\definecolor{currentstroke}{rgb}{0.000000,0.000000,0.000000}%
\pgfsetstrokecolor{currentstroke}%
\pgfsetstrokeopacity{0.000000}%
\pgfsetdash{}{0pt}%
\pgfpathmoveto{\pgfqpoint{3.374522in}{0.869232in}}%
\pgfpathlineto{\pgfqpoint{3.340166in}{0.869232in}}%
\pgfpathlineto{\pgfqpoint{3.340166in}{0.892142in}}%
\pgfpathlineto{\pgfqpoint{3.374522in}{0.892142in}}%
\pgfpathlineto{\pgfqpoint{3.374522in}{0.869232in}}%
\pgfpathclose%
\pgfusepath{fill}%
\end{pgfscope}%
\begin{pgfscope}%
\pgfpathrectangle{\pgfqpoint{2.893540in}{0.169444in}}{\pgfqpoint{0.858895in}{1.282944in}}%
\pgfusepath{clip}%
\pgfsetbuttcap%
\pgfsetmiterjoin%
\definecolor{currentfill}{rgb}{0.983852,0.910265,0.948328}%
\pgfsetfillcolor{currentfill}%
\pgfsetlinewidth{0.000000pt}%
\definecolor{currentstroke}{rgb}{0.000000,0.000000,0.000000}%
\pgfsetstrokecolor{currentstroke}%
\pgfsetstrokeopacity{0.000000}%
\pgfsetdash{}{0pt}%
\pgfpathmoveto{\pgfqpoint{3.374522in}{0.892559in}}%
\pgfpathlineto{\pgfqpoint{3.340166in}{0.892559in}}%
\pgfpathlineto{\pgfqpoint{3.340166in}{0.915468in}}%
\pgfpathlineto{\pgfqpoint{3.374522in}{0.915468in}}%
\pgfpathlineto{\pgfqpoint{3.374522in}{0.892559in}}%
\pgfpathclose%
\pgfusepath{fill}%
\end{pgfscope}%
\begin{pgfscope}%
\pgfpathrectangle{\pgfqpoint{2.893540in}{0.169444in}}{\pgfqpoint{0.858895in}{1.282944in}}%
\pgfusepath{clip}%
\pgfsetbuttcap%
\pgfsetmiterjoin%
\definecolor{currentfill}{rgb}{0.987543,0.896117,0.943406}%
\pgfsetfillcolor{currentfill}%
\pgfsetlinewidth{0.000000pt}%
\definecolor{currentstroke}{rgb}{0.000000,0.000000,0.000000}%
\pgfsetstrokecolor{currentstroke}%
\pgfsetstrokeopacity{0.000000}%
\pgfsetdash{}{0pt}%
\pgfpathmoveto{\pgfqpoint{3.374522in}{0.915885in}}%
\pgfpathlineto{\pgfqpoint{3.340166in}{0.915885in}}%
\pgfpathlineto{\pgfqpoint{3.340166in}{0.938795in}}%
\pgfpathlineto{\pgfqpoint{3.374522in}{0.938795in}}%
\pgfpathlineto{\pgfqpoint{3.374522in}{0.915885in}}%
\pgfpathclose%
\pgfusepath{fill}%
\end{pgfscope}%
\begin{pgfscope}%
\pgfpathrectangle{\pgfqpoint{2.893540in}{0.169444in}}{\pgfqpoint{0.858895in}{1.282944in}}%
\pgfusepath{clip}%
\pgfsetbuttcap%
\pgfsetmiterjoin%
\definecolor{currentfill}{rgb}{0.992157,0.878431,0.937255}%
\pgfsetfillcolor{currentfill}%
\pgfsetlinewidth{0.000000pt}%
\definecolor{currentstroke}{rgb}{0.000000,0.000000,0.000000}%
\pgfsetstrokecolor{currentstroke}%
\pgfsetstrokeopacity{0.000000}%
\pgfsetdash{}{0pt}%
\pgfpathmoveto{\pgfqpoint{3.374522in}{0.939211in}}%
\pgfpathlineto{\pgfqpoint{3.340166in}{0.939211in}}%
\pgfpathlineto{\pgfqpoint{3.340166in}{0.962121in}}%
\pgfpathlineto{\pgfqpoint{3.374522in}{0.962121in}}%
\pgfpathlineto{\pgfqpoint{3.374522in}{0.939211in}}%
\pgfpathclose%
\pgfusepath{fill}%
\end{pgfscope}%
\begin{pgfscope}%
\pgfpathrectangle{\pgfqpoint{2.893540in}{0.169444in}}{\pgfqpoint{0.858895in}{1.282944in}}%
\pgfusepath{clip}%
\pgfsetbuttcap%
\pgfsetmiterjoin%
\definecolor{currentfill}{rgb}{0.982930,0.846136,0.921107}%
\pgfsetfillcolor{currentfill}%
\pgfsetlinewidth{0.000000pt}%
\definecolor{currentstroke}{rgb}{0.000000,0.000000,0.000000}%
\pgfsetstrokecolor{currentstroke}%
\pgfsetstrokeopacity{0.000000}%
\pgfsetdash{}{0pt}%
\pgfpathmoveto{\pgfqpoint{3.374522in}{0.962537in}}%
\pgfpathlineto{\pgfqpoint{3.340166in}{0.962537in}}%
\pgfpathlineto{\pgfqpoint{3.340166in}{0.985447in}}%
\pgfpathlineto{\pgfqpoint{3.374522in}{0.985447in}}%
\pgfpathlineto{\pgfqpoint{3.374522in}{0.962537in}}%
\pgfpathclose%
\pgfusepath{fill}%
\end{pgfscope}%
\begin{pgfscope}%
\pgfpathrectangle{\pgfqpoint{2.893540in}{0.169444in}}{\pgfqpoint{0.858895in}{1.282944in}}%
\pgfusepath{clip}%
\pgfsetbuttcap%
\pgfsetmiterjoin%
\definecolor{currentfill}{rgb}{0.975548,0.820300,0.908189}%
\pgfsetfillcolor{currentfill}%
\pgfsetlinewidth{0.000000pt}%
\definecolor{currentstroke}{rgb}{0.000000,0.000000,0.000000}%
\pgfsetstrokecolor{currentstroke}%
\pgfsetstrokeopacity{0.000000}%
\pgfsetdash{}{0pt}%
\pgfpathmoveto{\pgfqpoint{3.374522in}{0.985864in}}%
\pgfpathlineto{\pgfqpoint{3.340166in}{0.985864in}}%
\pgfpathlineto{\pgfqpoint{3.340166in}{1.008773in}}%
\pgfpathlineto{\pgfqpoint{3.374522in}{1.008773in}}%
\pgfpathlineto{\pgfqpoint{3.374522in}{0.985864in}}%
\pgfpathclose%
\pgfusepath{fill}%
\end{pgfscope}%
\begin{pgfscope}%
\pgfpathrectangle{\pgfqpoint{2.893540in}{0.169444in}}{\pgfqpoint{0.858895in}{1.282944in}}%
\pgfusepath{clip}%
\pgfsetbuttcap%
\pgfsetmiterjoin%
\definecolor{currentfill}{rgb}{0.966321,0.788005,0.892042}%
\pgfsetfillcolor{currentfill}%
\pgfsetlinewidth{0.000000pt}%
\definecolor{currentstroke}{rgb}{0.000000,0.000000,0.000000}%
\pgfsetstrokecolor{currentstroke}%
\pgfsetstrokeopacity{0.000000}%
\pgfsetdash{}{0pt}%
\pgfpathmoveto{\pgfqpoint{3.374522in}{1.009190in}}%
\pgfpathlineto{\pgfqpoint{3.340166in}{1.009190in}}%
\pgfpathlineto{\pgfqpoint{3.340166in}{1.032100in}}%
\pgfpathlineto{\pgfqpoint{3.374522in}{1.032100in}}%
\pgfpathlineto{\pgfqpoint{3.374522in}{1.009190in}}%
\pgfpathclose%
\pgfusepath{fill}%
\end{pgfscope}%
\begin{pgfscope}%
\pgfpathrectangle{\pgfqpoint{2.893540in}{0.169444in}}{\pgfqpoint{0.858895in}{1.282944in}}%
\pgfusepath{clip}%
\pgfsetbuttcap%
\pgfsetmiterjoin%
\definecolor{currentfill}{rgb}{0.957093,0.755709,0.875894}%
\pgfsetfillcolor{currentfill}%
\pgfsetlinewidth{0.000000pt}%
\definecolor{currentstroke}{rgb}{0.000000,0.000000,0.000000}%
\pgfsetstrokecolor{currentstroke}%
\pgfsetstrokeopacity{0.000000}%
\pgfsetdash{}{0pt}%
\pgfpathmoveto{\pgfqpoint{3.374522in}{1.032516in}}%
\pgfpathlineto{\pgfqpoint{3.340166in}{1.032516in}}%
\pgfpathlineto{\pgfqpoint{3.340166in}{1.055426in}}%
\pgfpathlineto{\pgfqpoint{3.374522in}{1.055426in}}%
\pgfpathlineto{\pgfqpoint{3.374522in}{1.032516in}}%
\pgfpathclose%
\pgfusepath{fill}%
\end{pgfscope}%
\begin{pgfscope}%
\pgfpathrectangle{\pgfqpoint{2.893540in}{0.169444in}}{\pgfqpoint{0.858895in}{1.282944in}}%
\pgfusepath{clip}%
\pgfsetbuttcap%
\pgfsetmiterjoin%
\definecolor{currentfill}{rgb}{0.949712,0.729873,0.862976}%
\pgfsetfillcolor{currentfill}%
\pgfsetlinewidth{0.000000pt}%
\definecolor{currentstroke}{rgb}{0.000000,0.000000,0.000000}%
\pgfsetstrokecolor{currentstroke}%
\pgfsetstrokeopacity{0.000000}%
\pgfsetdash{}{0pt}%
\pgfpathmoveto{\pgfqpoint{3.374522in}{1.055842in}}%
\pgfpathlineto{\pgfqpoint{3.340166in}{1.055842in}}%
\pgfpathlineto{\pgfqpoint{3.340166in}{1.078752in}}%
\pgfpathlineto{\pgfqpoint{3.374522in}{1.078752in}}%
\pgfpathlineto{\pgfqpoint{3.374522in}{1.055842in}}%
\pgfpathclose%
\pgfusepath{fill}%
\end{pgfscope}%
\begin{pgfscope}%
\pgfpathrectangle{\pgfqpoint{2.893540in}{0.169444in}}{\pgfqpoint{0.858895in}{1.282944in}}%
\pgfusepath{clip}%
\pgfsetbuttcap%
\pgfsetmiterjoin%
\definecolor{currentfill}{rgb}{0.937793,0.689504,0.837985}%
\pgfsetfillcolor{currentfill}%
\pgfsetlinewidth{0.000000pt}%
\definecolor{currentstroke}{rgb}{0.000000,0.000000,0.000000}%
\pgfsetstrokecolor{currentstroke}%
\pgfsetstrokeopacity{0.000000}%
\pgfsetdash{}{0pt}%
\pgfpathmoveto{\pgfqpoint{3.374522in}{1.079169in}}%
\pgfpathlineto{\pgfqpoint{3.340166in}{1.079169in}}%
\pgfpathlineto{\pgfqpoint{3.340166in}{1.102078in}}%
\pgfpathlineto{\pgfqpoint{3.374522in}{1.102078in}}%
\pgfpathlineto{\pgfqpoint{3.374522in}{1.079169in}}%
\pgfpathclose%
\pgfusepath{fill}%
\end{pgfscope}%
\begin{pgfscope}%
\pgfpathrectangle{\pgfqpoint{2.893540in}{0.169444in}}{\pgfqpoint{0.858895in}{1.282944in}}%
\pgfusepath{clip}%
\pgfsetbuttcap%
\pgfsetmiterjoin%
\definecolor{currentfill}{rgb}{0.923183,0.641061,0.804152}%
\pgfsetfillcolor{currentfill}%
\pgfsetlinewidth{0.000000pt}%
\definecolor{currentstroke}{rgb}{0.000000,0.000000,0.000000}%
\pgfsetstrokecolor{currentstroke}%
\pgfsetstrokeopacity{0.000000}%
\pgfsetdash{}{0pt}%
\pgfpathmoveto{\pgfqpoint{3.374522in}{1.102495in}}%
\pgfpathlineto{\pgfqpoint{3.340166in}{1.102495in}}%
\pgfpathlineto{\pgfqpoint{3.340166in}{1.125405in}}%
\pgfpathlineto{\pgfqpoint{3.374522in}{1.125405in}}%
\pgfpathlineto{\pgfqpoint{3.374522in}{1.102495in}}%
\pgfpathclose%
\pgfusepath{fill}%
\end{pgfscope}%
\begin{pgfscope}%
\pgfpathrectangle{\pgfqpoint{2.893540in}{0.169444in}}{\pgfqpoint{0.858895in}{1.282944in}}%
\pgfusepath{clip}%
\pgfsetbuttcap%
\pgfsetmiterjoin%
\definecolor{currentfill}{rgb}{0.911496,0.602307,0.777086}%
\pgfsetfillcolor{currentfill}%
\pgfsetlinewidth{0.000000pt}%
\definecolor{currentstroke}{rgb}{0.000000,0.000000,0.000000}%
\pgfsetstrokecolor{currentstroke}%
\pgfsetstrokeopacity{0.000000}%
\pgfsetdash{}{0pt}%
\pgfpathmoveto{\pgfqpoint{3.374522in}{1.125821in}}%
\pgfpathlineto{\pgfqpoint{3.340166in}{1.125821in}}%
\pgfpathlineto{\pgfqpoint{3.340166in}{1.148731in}}%
\pgfpathlineto{\pgfqpoint{3.374522in}{1.148731in}}%
\pgfpathlineto{\pgfqpoint{3.374522in}{1.125821in}}%
\pgfpathclose%
\pgfusepath{fill}%
\end{pgfscope}%
\begin{pgfscope}%
\pgfpathrectangle{\pgfqpoint{2.893540in}{0.169444in}}{\pgfqpoint{0.858895in}{1.282944in}}%
\pgfusepath{clip}%
\pgfsetbuttcap%
\pgfsetmiterjoin%
\definecolor{currentfill}{rgb}{0.896886,0.553864,0.743253}%
\pgfsetfillcolor{currentfill}%
\pgfsetlinewidth{0.000000pt}%
\definecolor{currentstroke}{rgb}{0.000000,0.000000,0.000000}%
\pgfsetstrokecolor{currentstroke}%
\pgfsetstrokeopacity{0.000000}%
\pgfsetdash{}{0pt}%
\pgfpathmoveto{\pgfqpoint{3.374522in}{1.149147in}}%
\pgfpathlineto{\pgfqpoint{3.340166in}{1.149147in}}%
\pgfpathlineto{\pgfqpoint{3.340166in}{1.172057in}}%
\pgfpathlineto{\pgfqpoint{3.374522in}{1.172057in}}%
\pgfpathlineto{\pgfqpoint{3.374522in}{1.149147in}}%
\pgfpathclose%
\pgfusepath{fill}%
\end{pgfscope}%
\begin{pgfscope}%
\pgfpathrectangle{\pgfqpoint{2.893540in}{0.169444in}}{\pgfqpoint{0.858895in}{1.282944in}}%
\pgfusepath{clip}%
\pgfsetbuttcap%
\pgfsetmiterjoin%
\definecolor{currentfill}{rgb}{0.882276,0.505421,0.709419}%
\pgfsetfillcolor{currentfill}%
\pgfsetlinewidth{0.000000pt}%
\definecolor{currentstroke}{rgb}{0.000000,0.000000,0.000000}%
\pgfsetstrokecolor{currentstroke}%
\pgfsetstrokeopacity{0.000000}%
\pgfsetdash{}{0pt}%
\pgfpathmoveto{\pgfqpoint{3.374522in}{1.172474in}}%
\pgfpathlineto{\pgfqpoint{3.340166in}{1.172474in}}%
\pgfpathlineto{\pgfqpoint{3.340166in}{1.195383in}}%
\pgfpathlineto{\pgfqpoint{3.374522in}{1.195383in}}%
\pgfpathlineto{\pgfqpoint{3.374522in}{1.172474in}}%
\pgfpathclose%
\pgfusepath{fill}%
\end{pgfscope}%
\begin{pgfscope}%
\pgfpathrectangle{\pgfqpoint{2.893540in}{0.169444in}}{\pgfqpoint{0.858895in}{1.282944in}}%
\pgfusepath{clip}%
\pgfsetbuttcap%
\pgfsetmiterjoin%
\definecolor{currentfill}{rgb}{0.870588,0.466667,0.682353}%
\pgfsetfillcolor{currentfill}%
\pgfsetlinewidth{0.000000pt}%
\definecolor{currentstroke}{rgb}{0.000000,0.000000,0.000000}%
\pgfsetstrokecolor{currentstroke}%
\pgfsetstrokeopacity{0.000000}%
\pgfsetdash{}{0pt}%
\pgfpathmoveto{\pgfqpoint{3.374522in}{1.195800in}}%
\pgfpathlineto{\pgfqpoint{3.340166in}{1.195800in}}%
\pgfpathlineto{\pgfqpoint{3.340166in}{1.218710in}}%
\pgfpathlineto{\pgfqpoint{3.374522in}{1.218710in}}%
\pgfpathlineto{\pgfqpoint{3.374522in}{1.195800in}}%
\pgfpathclose%
\pgfusepath{fill}%
\end{pgfscope}%
\begin{pgfscope}%
\pgfpathrectangle{\pgfqpoint{2.893540in}{0.169444in}}{\pgfqpoint{0.858895in}{1.282944in}}%
\pgfusepath{clip}%
\pgfsetbuttcap%
\pgfsetmiterjoin%
\definecolor{currentfill}{rgb}{0.851365,0.395925,0.644675}%
\pgfsetfillcolor{currentfill}%
\pgfsetlinewidth{0.000000pt}%
\definecolor{currentstroke}{rgb}{0.000000,0.000000,0.000000}%
\pgfsetstrokecolor{currentstroke}%
\pgfsetstrokeopacity{0.000000}%
\pgfsetdash{}{0pt}%
\pgfpathmoveto{\pgfqpoint{3.374522in}{1.219126in}}%
\pgfpathlineto{\pgfqpoint{3.340166in}{1.219126in}}%
\pgfpathlineto{\pgfqpoint{3.340166in}{1.242036in}}%
\pgfpathlineto{\pgfqpoint{3.374522in}{1.242036in}}%
\pgfpathlineto{\pgfqpoint{3.374522in}{1.219126in}}%
\pgfpathclose%
\pgfusepath{fill}%
\end{pgfscope}%
\begin{pgfscope}%
\pgfpathrectangle{\pgfqpoint{2.893540in}{0.169444in}}{\pgfqpoint{0.858895in}{1.282944in}}%
\pgfusepath{clip}%
\pgfsetbuttcap%
\pgfsetmiterjoin%
\definecolor{currentfill}{rgb}{0.832141,0.325183,0.606997}%
\pgfsetfillcolor{currentfill}%
\pgfsetlinewidth{0.000000pt}%
\definecolor{currentstroke}{rgb}{0.000000,0.000000,0.000000}%
\pgfsetstrokecolor{currentstroke}%
\pgfsetstrokeopacity{0.000000}%
\pgfsetdash{}{0pt}%
\pgfpathmoveto{\pgfqpoint{3.374522in}{1.242453in}}%
\pgfpathlineto{\pgfqpoint{3.340166in}{1.242453in}}%
\pgfpathlineto{\pgfqpoint{3.340166in}{1.265362in}}%
\pgfpathlineto{\pgfqpoint{3.374522in}{1.265362in}}%
\pgfpathlineto{\pgfqpoint{3.374522in}{1.242453in}}%
\pgfpathclose%
\pgfusepath{fill}%
\end{pgfscope}%
\begin{pgfscope}%
\pgfpathrectangle{\pgfqpoint{2.893540in}{0.169444in}}{\pgfqpoint{0.858895in}{1.282944in}}%
\pgfusepath{clip}%
\pgfsetbuttcap%
\pgfsetmiterjoin%
\definecolor{currentfill}{rgb}{0.816763,0.268589,0.576855}%
\pgfsetfillcolor{currentfill}%
\pgfsetlinewidth{0.000000pt}%
\definecolor{currentstroke}{rgb}{0.000000,0.000000,0.000000}%
\pgfsetstrokecolor{currentstroke}%
\pgfsetstrokeopacity{0.000000}%
\pgfsetdash{}{0pt}%
\pgfpathmoveto{\pgfqpoint{3.374522in}{1.265779in}}%
\pgfpathlineto{\pgfqpoint{3.340166in}{1.265779in}}%
\pgfpathlineto{\pgfqpoint{3.340166in}{1.288689in}}%
\pgfpathlineto{\pgfqpoint{3.374522in}{1.288689in}}%
\pgfpathlineto{\pgfqpoint{3.374522in}{1.265779in}}%
\pgfpathclose%
\pgfusepath{fill}%
\end{pgfscope}%
\begin{pgfscope}%
\pgfpathrectangle{\pgfqpoint{2.893540in}{0.169444in}}{\pgfqpoint{0.858895in}{1.282944in}}%
\pgfusepath{clip}%
\pgfsetbuttcap%
\pgfsetmiterjoin%
\definecolor{currentfill}{rgb}{0.797539,0.197847,0.539177}%
\pgfsetfillcolor{currentfill}%
\pgfsetlinewidth{0.000000pt}%
\definecolor{currentstroke}{rgb}{0.000000,0.000000,0.000000}%
\pgfsetstrokecolor{currentstroke}%
\pgfsetstrokeopacity{0.000000}%
\pgfsetdash{}{0pt}%
\pgfpathmoveto{\pgfqpoint{3.374522in}{1.289105in}}%
\pgfpathlineto{\pgfqpoint{3.340166in}{1.289105in}}%
\pgfpathlineto{\pgfqpoint{3.340166in}{1.312015in}}%
\pgfpathlineto{\pgfqpoint{3.374522in}{1.312015in}}%
\pgfpathlineto{\pgfqpoint{3.374522in}{1.289105in}}%
\pgfpathclose%
\pgfusepath{fill}%
\end{pgfscope}%
\begin{pgfscope}%
\pgfpathrectangle{\pgfqpoint{2.893540in}{0.169444in}}{\pgfqpoint{0.858895in}{1.282944in}}%
\pgfusepath{clip}%
\pgfsetbuttcap%
\pgfsetmiterjoin%
\definecolor{currentfill}{rgb}{0.778316,0.127105,0.501499}%
\pgfsetfillcolor{currentfill}%
\pgfsetlinewidth{0.000000pt}%
\definecolor{currentstroke}{rgb}{0.000000,0.000000,0.000000}%
\pgfsetstrokecolor{currentstroke}%
\pgfsetstrokeopacity{0.000000}%
\pgfsetdash{}{0pt}%
\pgfpathmoveto{\pgfqpoint{3.374522in}{1.312431in}}%
\pgfpathlineto{\pgfqpoint{3.340166in}{1.312431in}}%
\pgfpathlineto{\pgfqpoint{3.340166in}{1.335341in}}%
\pgfpathlineto{\pgfqpoint{3.374522in}{1.335341in}}%
\pgfpathlineto{\pgfqpoint{3.374522in}{1.312431in}}%
\pgfpathclose%
\pgfusepath{fill}%
\end{pgfscope}%
\begin{pgfscope}%
\pgfpathrectangle{\pgfqpoint{2.893540in}{0.169444in}}{\pgfqpoint{0.858895in}{1.282944in}}%
\pgfusepath{clip}%
\pgfsetbuttcap%
\pgfsetmiterjoin%
\definecolor{currentfill}{rgb}{0.751403,0.095886,0.473664}%
\pgfsetfillcolor{currentfill}%
\pgfsetlinewidth{0.000000pt}%
\definecolor{currentstroke}{rgb}{0.000000,0.000000,0.000000}%
\pgfsetstrokecolor{currentstroke}%
\pgfsetstrokeopacity{0.000000}%
\pgfsetdash{}{0pt}%
\pgfpathmoveto{\pgfqpoint{3.374522in}{1.335758in}}%
\pgfpathlineto{\pgfqpoint{3.340166in}{1.335758in}}%
\pgfpathlineto{\pgfqpoint{3.340166in}{1.358667in}}%
\pgfpathlineto{\pgfqpoint{3.374522in}{1.358667in}}%
\pgfpathlineto{\pgfqpoint{3.374522in}{1.335758in}}%
\pgfpathclose%
\pgfusepath{fill}%
\end{pgfscope}%
\begin{pgfscope}%
\pgfpathrectangle{\pgfqpoint{2.893540in}{0.169444in}}{\pgfqpoint{0.858895in}{1.282944in}}%
\pgfusepath{clip}%
\pgfsetbuttcap%
\pgfsetmiterjoin%
\definecolor{currentfill}{rgb}{0.709112,0.075894,0.440600}%
\pgfsetfillcolor{currentfill}%
\pgfsetlinewidth{0.000000pt}%
\definecolor{currentstroke}{rgb}{0.000000,0.000000,0.000000}%
\pgfsetstrokecolor{currentstroke}%
\pgfsetstrokeopacity{0.000000}%
\pgfsetdash{}{0pt}%
\pgfpathmoveto{\pgfqpoint{3.374522in}{1.359084in}}%
\pgfpathlineto{\pgfqpoint{3.340166in}{1.359084in}}%
\pgfpathlineto{\pgfqpoint{3.340166in}{1.381994in}}%
\pgfpathlineto{\pgfqpoint{3.374522in}{1.381994in}}%
\pgfpathlineto{\pgfqpoint{3.374522in}{1.359084in}}%
\pgfpathclose%
\pgfusepath{fill}%
\end{pgfscope}%
\begin{pgfscope}%
\pgfpathrectangle{\pgfqpoint{2.893540in}{0.169444in}}{\pgfqpoint{0.858895in}{1.282944in}}%
\pgfusepath{clip}%
\pgfsetbuttcap%
\pgfsetmiterjoin%
\definecolor{currentfill}{rgb}{0.666820,0.055902,0.407536}%
\pgfsetfillcolor{currentfill}%
\pgfsetlinewidth{0.000000pt}%
\definecolor{currentstroke}{rgb}{0.000000,0.000000,0.000000}%
\pgfsetstrokecolor{currentstroke}%
\pgfsetstrokeopacity{0.000000}%
\pgfsetdash{}{0pt}%
\pgfpathmoveto{\pgfqpoint{3.374522in}{1.382410in}}%
\pgfpathlineto{\pgfqpoint{3.340166in}{1.382410in}}%
\pgfpathlineto{\pgfqpoint{3.340166in}{1.405320in}}%
\pgfpathlineto{\pgfqpoint{3.374522in}{1.405320in}}%
\pgfpathlineto{\pgfqpoint{3.374522in}{1.382410in}}%
\pgfpathclose%
\pgfusepath{fill}%
\end{pgfscope}%
\begin{pgfscope}%
\pgfpathrectangle{\pgfqpoint{2.893540in}{0.169444in}}{\pgfqpoint{0.858895in}{1.282944in}}%
\pgfusepath{clip}%
\pgfsetbuttcap%
\pgfsetmiterjoin%
\definecolor{currentfill}{rgb}{0.632987,0.039908,0.381084}%
\pgfsetfillcolor{currentfill}%
\pgfsetlinewidth{0.000000pt}%
\definecolor{currentstroke}{rgb}{0.000000,0.000000,0.000000}%
\pgfsetstrokecolor{currentstroke}%
\pgfsetstrokeopacity{0.000000}%
\pgfsetdash{}{0pt}%
\pgfpathmoveto{\pgfqpoint{3.374522in}{1.405736in}}%
\pgfpathlineto{\pgfqpoint{3.340166in}{1.405736in}}%
\pgfpathlineto{\pgfqpoint{3.340166in}{1.428646in}}%
\pgfpathlineto{\pgfqpoint{3.374522in}{1.428646in}}%
\pgfpathlineto{\pgfqpoint{3.374522in}{1.405736in}}%
\pgfpathclose%
\pgfusepath{fill}%
\end{pgfscope}%
\begin{pgfscope}%
\pgfpathrectangle{\pgfqpoint{2.893540in}{0.169444in}}{\pgfqpoint{0.858895in}{1.282944in}}%
\pgfusepath{clip}%
\pgfsetbuttcap%
\pgfsetmiterjoin%
\definecolor{currentfill}{rgb}{0.590696,0.019915,0.348020}%
\pgfsetfillcolor{currentfill}%
\pgfsetlinewidth{0.000000pt}%
\definecolor{currentstroke}{rgb}{0.000000,0.000000,0.000000}%
\pgfsetstrokecolor{currentstroke}%
\pgfsetstrokeopacity{0.000000}%
\pgfsetdash{}{0pt}%
\pgfpathmoveto{\pgfqpoint{3.374522in}{1.429063in}}%
\pgfpathlineto{\pgfqpoint{3.340166in}{1.429063in}}%
\pgfpathlineto{\pgfqpoint{3.340166in}{1.451972in}}%
\pgfpathlineto{\pgfqpoint{3.374522in}{1.451972in}}%
\pgfpathlineto{\pgfqpoint{3.374522in}{1.429063in}}%
\pgfpathclose%
\pgfusepath{fill}%
\end{pgfscope}%
\begin{pgfscope}%
\pgfpathrectangle{\pgfqpoint{2.893540in}{0.169444in}}{\pgfqpoint{0.858895in}{1.282944in}}%
\pgfusepath{clip}%
\pgfsetbuttcap%
\pgfsetmiterjoin%
\definecolor{currentfill}{rgb}{0.556863,0.003922,0.321569}%
\pgfsetfillcolor{currentfill}%
\pgfsetlinewidth{0.000000pt}%
\definecolor{currentstroke}{rgb}{0.000000,0.000000,0.000000}%
\pgfsetstrokecolor{currentstroke}%
\pgfsetstrokeopacity{0.000000}%
\pgfsetdash{}{0pt}%
\pgfpathmoveto{\pgfqpoint{3.374522in}{1.452389in}}%
\pgfpathlineto{\pgfqpoint{3.340166in}{1.452389in}}%
\pgfpathlineto{\pgfqpoint{3.340166in}{1.475299in}}%
\pgfpathlineto{\pgfqpoint{3.374522in}{1.475299in}}%
\pgfpathlineto{\pgfqpoint{3.374522in}{1.452389in}}%
\pgfpathclose%
\pgfusepath{fill}%
\end{pgfscope}%
\begin{pgfscope}%
\pgfpathrectangle{\pgfqpoint{2.893540in}{0.169444in}}{\pgfqpoint{0.858895in}{1.282944in}}%
\pgfusepath{clip}%
\pgfsetbuttcap%
\pgfsetmiterjoin%
\definecolor{currentfill}{rgb}{0.263975,0.526567,0.121415}%
\pgfsetfillcolor{currentfill}%
\pgfsetlinewidth{0.000000pt}%
\definecolor{currentstroke}{rgb}{0.000000,0.000000,0.000000}%
\pgfsetstrokecolor{currentstroke}%
\pgfsetstrokeopacity{0.000000}%
\pgfsetdash{}{0pt}%
\pgfpathmoveto{\pgfqpoint{3.408878in}{0.266235in}}%
\pgfpathlineto{\pgfqpoint{3.426056in}{0.266235in}}%
\pgfpathlineto{\pgfqpoint{3.426056in}{0.289994in}}%
\pgfpathlineto{\pgfqpoint{3.408878in}{0.289994in}}%
\pgfpathlineto{\pgfqpoint{3.408878in}{0.266235in}}%
\pgfpathclose%
\pgfusepath{fill}%
\end{pgfscope}%
\begin{pgfscope}%
\pgfpathrectangle{\pgfqpoint{2.893540in}{0.169444in}}{\pgfqpoint{0.858895in}{1.282944in}}%
\pgfusepath{clip}%
\pgfsetbuttcap%
\pgfsetmiterjoin%
\definecolor{currentfill}{rgb}{0.293195,0.561938,0.127566}%
\pgfsetfillcolor{currentfill}%
\pgfsetlinewidth{0.000000pt}%
\definecolor{currentstroke}{rgb}{0.000000,0.000000,0.000000}%
\pgfsetstrokecolor{currentstroke}%
\pgfsetstrokeopacity{0.000000}%
\pgfsetdash{}{0pt}%
\pgfpathmoveto{\pgfqpoint{3.408878in}{0.290469in}}%
\pgfpathlineto{\pgfqpoint{3.426056in}{0.290469in}}%
\pgfpathlineto{\pgfqpoint{3.426056in}{0.314227in}}%
\pgfpathlineto{\pgfqpoint{3.408878in}{0.314227in}}%
\pgfpathlineto{\pgfqpoint{3.408878in}{0.290469in}}%
\pgfpathclose%
\pgfusepath{fill}%
\end{pgfscope}%
\begin{pgfscope}%
\pgfpathrectangle{\pgfqpoint{2.893540in}{0.169444in}}{\pgfqpoint{0.858895in}{1.282944in}}%
\pgfusepath{clip}%
\pgfsetbuttcap%
\pgfsetmiterjoin%
\definecolor{currentfill}{rgb}{0.359631,0.620992,0.166321}%
\pgfsetfillcolor{currentfill}%
\pgfsetlinewidth{0.000000pt}%
\definecolor{currentstroke}{rgb}{0.000000,0.000000,0.000000}%
\pgfsetstrokecolor{currentstroke}%
\pgfsetstrokeopacity{0.000000}%
\pgfsetdash{}{0pt}%
\pgfpathmoveto{\pgfqpoint{3.408878in}{0.338936in}}%
\pgfpathlineto{\pgfqpoint{3.494767in}{0.338936in}}%
\pgfpathlineto{\pgfqpoint{3.494767in}{0.362694in}}%
\pgfpathlineto{\pgfqpoint{3.408878in}{0.362694in}}%
\pgfpathlineto{\pgfqpoint{3.408878in}{0.338936in}}%
\pgfpathclose%
\pgfusepath{fill}%
\end{pgfscope}%
\begin{pgfscope}%
\pgfpathrectangle{\pgfqpoint{2.893540in}{0.169444in}}{\pgfqpoint{0.858895in}{1.282944in}}%
\pgfusepath{clip}%
\pgfsetbuttcap%
\pgfsetmiterjoin%
\definecolor{currentfill}{rgb}{0.398078,0.653287,0.190927}%
\pgfsetfillcolor{currentfill}%
\pgfsetlinewidth{0.000000pt}%
\definecolor{currentstroke}{rgb}{0.000000,0.000000,0.000000}%
\pgfsetstrokecolor{currentstroke}%
\pgfsetstrokeopacity{0.000000}%
\pgfsetdash{}{0pt}%
\pgfpathmoveto{\pgfqpoint{3.408878in}{0.363027in}}%
\pgfpathlineto{\pgfqpoint{3.477589in}{0.363027in}}%
\pgfpathlineto{\pgfqpoint{3.477589in}{0.386785in}}%
\pgfpathlineto{\pgfqpoint{3.408878in}{0.386785in}}%
\pgfpathlineto{\pgfqpoint{3.408878in}{0.363027in}}%
\pgfpathclose%
\pgfusepath{fill}%
\end{pgfscope}%
\begin{pgfscope}%
\pgfpathrectangle{\pgfqpoint{2.893540in}{0.169444in}}{\pgfqpoint{0.858895in}{1.282944in}}%
\pgfusepath{clip}%
\pgfsetbuttcap%
\pgfsetmiterjoin%
\definecolor{currentfill}{rgb}{0.436524,0.685582,0.215532}%
\pgfsetfillcolor{currentfill}%
\pgfsetlinewidth{0.000000pt}%
\definecolor{currentstroke}{rgb}{0.000000,0.000000,0.000000}%
\pgfsetstrokecolor{currentstroke}%
\pgfsetstrokeopacity{0.000000}%
\pgfsetdash{}{0pt}%
\pgfpathmoveto{\pgfqpoint{3.408878in}{0.387260in}}%
\pgfpathlineto{\pgfqpoint{3.443233in}{0.387260in}}%
\pgfpathlineto{\pgfqpoint{3.443233in}{0.411018in}}%
\pgfpathlineto{\pgfqpoint{3.408878in}{0.411018in}}%
\pgfpathlineto{\pgfqpoint{3.408878in}{0.387260in}}%
\pgfpathclose%
\pgfusepath{fill}%
\end{pgfscope}%
\begin{pgfscope}%
\pgfpathrectangle{\pgfqpoint{2.893540in}{0.169444in}}{\pgfqpoint{0.858895in}{1.282944in}}%
\pgfusepath{clip}%
\pgfsetbuttcap%
\pgfsetmiterjoin%
\definecolor{currentfill}{rgb}{0.474971,0.717878,0.240138}%
\pgfsetfillcolor{currentfill}%
\pgfsetlinewidth{0.000000pt}%
\definecolor{currentstroke}{rgb}{0.000000,0.000000,0.000000}%
\pgfsetstrokecolor{currentstroke}%
\pgfsetstrokeopacity{0.000000}%
\pgfsetdash{}{0pt}%
\pgfpathmoveto{\pgfqpoint{3.408878in}{0.411493in}}%
\pgfpathlineto{\pgfqpoint{3.511945in}{0.411493in}}%
\pgfpathlineto{\pgfqpoint{3.511945in}{0.435252in}}%
\pgfpathlineto{\pgfqpoint{3.408878in}{0.435252in}}%
\pgfpathlineto{\pgfqpoint{3.408878in}{0.411493in}}%
\pgfpathclose%
\pgfusepath{fill}%
\end{pgfscope}%
\begin{pgfscope}%
\pgfpathrectangle{\pgfqpoint{2.893540in}{0.169444in}}{\pgfqpoint{0.858895in}{1.282944in}}%
\pgfusepath{clip}%
\pgfsetbuttcap%
\pgfsetmiterjoin%
\definecolor{currentfill}{rgb}{0.515571,0.748635,0.276125}%
\pgfsetfillcolor{currentfill}%
\pgfsetlinewidth{0.000000pt}%
\definecolor{currentstroke}{rgb}{0.000000,0.000000,0.000000}%
\pgfsetstrokecolor{currentstroke}%
\pgfsetstrokeopacity{0.000000}%
\pgfsetdash{}{0pt}%
\pgfpathmoveto{\pgfqpoint{3.408878in}{0.435727in}}%
\pgfpathlineto{\pgfqpoint{3.511945in}{0.435727in}}%
\pgfpathlineto{\pgfqpoint{3.511945in}{0.459485in}}%
\pgfpathlineto{\pgfqpoint{3.408878in}{0.459485in}}%
\pgfpathlineto{\pgfqpoint{3.408878in}{0.435727in}}%
\pgfpathclose%
\pgfusepath{fill}%
\end{pgfscope}%
\begin{pgfscope}%
\pgfpathrectangle{\pgfqpoint{2.893540in}{0.169444in}}{\pgfqpoint{0.858895in}{1.282944in}}%
\pgfusepath{clip}%
\pgfsetbuttcap%
\pgfsetmiterjoin%
\definecolor{currentfill}{rgb}{0.550634,0.771396,0.318570}%
\pgfsetfillcolor{currentfill}%
\pgfsetlinewidth{0.000000pt}%
\definecolor{currentstroke}{rgb}{0.000000,0.000000,0.000000}%
\pgfsetstrokecolor{currentstroke}%
\pgfsetstrokeopacity{0.000000}%
\pgfsetdash{}{0pt}%
\pgfpathmoveto{\pgfqpoint{3.408878in}{0.459960in}}%
\pgfpathlineto{\pgfqpoint{3.477589in}{0.459960in}}%
\pgfpathlineto{\pgfqpoint{3.477589in}{0.483718in}}%
\pgfpathlineto{\pgfqpoint{3.408878in}{0.483718in}}%
\pgfpathlineto{\pgfqpoint{3.408878in}{0.459960in}}%
\pgfpathclose%
\pgfusepath{fill}%
\end{pgfscope}%
\begin{pgfscope}%
\pgfpathrectangle{\pgfqpoint{2.893540in}{0.169444in}}{\pgfqpoint{0.858895in}{1.282944in}}%
\pgfusepath{clip}%
\pgfsetbuttcap%
\pgfsetmiterjoin%
\definecolor{currentfill}{rgb}{0.594464,0.799846,0.371626}%
\pgfsetfillcolor{currentfill}%
\pgfsetlinewidth{0.000000pt}%
\definecolor{currentstroke}{rgb}{0.000000,0.000000,0.000000}%
\pgfsetstrokecolor{currentstroke}%
\pgfsetstrokeopacity{0.000000}%
\pgfsetdash{}{0pt}%
\pgfpathmoveto{\pgfqpoint{3.408878in}{0.484193in}}%
\pgfpathlineto{\pgfqpoint{3.426056in}{0.484193in}}%
\pgfpathlineto{\pgfqpoint{3.426056in}{0.507952in}}%
\pgfpathlineto{\pgfqpoint{3.408878in}{0.507952in}}%
\pgfpathlineto{\pgfqpoint{3.408878in}{0.484193in}}%
\pgfpathclose%
\pgfusepath{fill}%
\end{pgfscope}%
\begin{pgfscope}%
\pgfpathrectangle{\pgfqpoint{2.893540in}{0.169444in}}{\pgfqpoint{0.858895in}{1.282944in}}%
\pgfusepath{clip}%
\pgfsetbuttcap%
\pgfsetmiterjoin%
\definecolor{currentfill}{rgb}{0.638293,0.828297,0.424683}%
\pgfsetfillcolor{currentfill}%
\pgfsetlinewidth{0.000000pt}%
\definecolor{currentstroke}{rgb}{0.000000,0.000000,0.000000}%
\pgfsetstrokecolor{currentstroke}%
\pgfsetstrokeopacity{0.000000}%
\pgfsetdash{}{0pt}%
\pgfpathmoveto{\pgfqpoint{3.408878in}{0.508284in}}%
\pgfpathlineto{\pgfqpoint{3.511945in}{0.508284in}}%
\pgfpathlineto{\pgfqpoint{3.511945in}{0.532043in}}%
\pgfpathlineto{\pgfqpoint{3.408878in}{0.532043in}}%
\pgfpathlineto{\pgfqpoint{3.408878in}{0.508284in}}%
\pgfpathclose%
\pgfusepath{fill}%
\end{pgfscope}%
\begin{pgfscope}%
\pgfpathrectangle{\pgfqpoint{2.893540in}{0.169444in}}{\pgfqpoint{0.858895in}{1.282944in}}%
\pgfusepath{clip}%
\pgfsetbuttcap%
\pgfsetmiterjoin%
\definecolor{currentfill}{rgb}{0.682122,0.856747,0.477739}%
\pgfsetfillcolor{currentfill}%
\pgfsetlinewidth{0.000000pt}%
\definecolor{currentstroke}{rgb}{0.000000,0.000000,0.000000}%
\pgfsetstrokecolor{currentstroke}%
\pgfsetstrokeopacity{0.000000}%
\pgfsetdash{}{0pt}%
\pgfpathmoveto{\pgfqpoint{3.408878in}{0.532518in}}%
\pgfpathlineto{\pgfqpoint{3.494767in}{0.532518in}}%
\pgfpathlineto{\pgfqpoint{3.494767in}{0.556276in}}%
\pgfpathlineto{\pgfqpoint{3.408878in}{0.556276in}}%
\pgfpathlineto{\pgfqpoint{3.408878in}{0.532518in}}%
\pgfpathclose%
\pgfusepath{fill}%
\end{pgfscope}%
\begin{pgfscope}%
\pgfpathrectangle{\pgfqpoint{2.893540in}{0.169444in}}{\pgfqpoint{0.858895in}{1.282944in}}%
\pgfusepath{clip}%
\pgfsetbuttcap%
\pgfsetmiterjoin%
\definecolor{currentfill}{rgb}{0.725106,0.883891,0.531180}%
\pgfsetfillcolor{currentfill}%
\pgfsetlinewidth{0.000000pt}%
\definecolor{currentstroke}{rgb}{0.000000,0.000000,0.000000}%
\pgfsetstrokecolor{currentstroke}%
\pgfsetstrokeopacity{0.000000}%
\pgfsetdash{}{0pt}%
\pgfpathmoveto{\pgfqpoint{3.408878in}{0.556751in}}%
\pgfpathlineto{\pgfqpoint{3.426056in}{0.556751in}}%
\pgfpathlineto{\pgfqpoint{3.426056in}{0.580509in}}%
\pgfpathlineto{\pgfqpoint{3.408878in}{0.580509in}}%
\pgfpathlineto{\pgfqpoint{3.408878in}{0.556751in}}%
\pgfpathclose%
\pgfusepath{fill}%
\end{pgfscope}%
\begin{pgfscope}%
\pgfpathrectangle{\pgfqpoint{2.893540in}{0.169444in}}{\pgfqpoint{0.858895in}{1.282944in}}%
\pgfusepath{clip}%
\pgfsetbuttcap%
\pgfsetmiterjoin%
\definecolor{currentfill}{rgb}{0.760477,0.899270,0.588082}%
\pgfsetfillcolor{currentfill}%
\pgfsetlinewidth{0.000000pt}%
\definecolor{currentstroke}{rgb}{0.000000,0.000000,0.000000}%
\pgfsetstrokecolor{currentstroke}%
\pgfsetstrokeopacity{0.000000}%
\pgfsetdash{}{0pt}%
\pgfpathmoveto{\pgfqpoint{3.408878in}{0.580985in}}%
\pgfpathlineto{\pgfqpoint{3.460411in}{0.580985in}}%
\pgfpathlineto{\pgfqpoint{3.460411in}{0.604743in}}%
\pgfpathlineto{\pgfqpoint{3.408878in}{0.604743in}}%
\pgfpathlineto{\pgfqpoint{3.408878in}{0.580985in}}%
\pgfpathclose%
\pgfusepath{fill}%
\end{pgfscope}%
\begin{pgfscope}%
\pgfpathrectangle{\pgfqpoint{2.893540in}{0.169444in}}{\pgfqpoint{0.858895in}{1.282944in}}%
\pgfusepath{clip}%
\pgfsetbuttcap%
\pgfsetmiterjoin%
\definecolor{currentfill}{rgb}{0.788774,0.911572,0.633602}%
\pgfsetfillcolor{currentfill}%
\pgfsetlinewidth{0.000000pt}%
\definecolor{currentstroke}{rgb}{0.000000,0.000000,0.000000}%
\pgfsetstrokecolor{currentstroke}%
\pgfsetstrokeopacity{0.000000}%
\pgfsetdash{}{0pt}%
\pgfpathmoveto{\pgfqpoint{3.408878in}{0.605218in}}%
\pgfpathlineto{\pgfqpoint{3.460411in}{0.605218in}}%
\pgfpathlineto{\pgfqpoint{3.460411in}{0.628976in}}%
\pgfpathlineto{\pgfqpoint{3.408878in}{0.628976in}}%
\pgfpathlineto{\pgfqpoint{3.408878in}{0.605218in}}%
\pgfpathclose%
\pgfusepath{fill}%
\end{pgfscope}%
\begin{pgfscope}%
\pgfpathrectangle{\pgfqpoint{2.893540in}{0.169444in}}{\pgfqpoint{0.858895in}{1.282944in}}%
\pgfusepath{clip}%
\pgfsetbuttcap%
\pgfsetmiterjoin%
\definecolor{currentfill}{rgb}{0.859516,0.942330,0.747405}%
\pgfsetfillcolor{currentfill}%
\pgfsetlinewidth{0.000000pt}%
\definecolor{currentstroke}{rgb}{0.000000,0.000000,0.000000}%
\pgfsetstrokecolor{currentstroke}%
\pgfsetstrokeopacity{0.000000}%
\pgfsetdash{}{0pt}%
\pgfpathmoveto{\pgfqpoint{3.408878in}{0.653542in}}%
\pgfpathlineto{\pgfqpoint{3.443233in}{0.653542in}}%
\pgfpathlineto{\pgfqpoint{3.443233in}{0.677300in}}%
\pgfpathlineto{\pgfqpoint{3.408878in}{0.677300in}}%
\pgfpathlineto{\pgfqpoint{3.408878in}{0.653542in}}%
\pgfpathclose%
\pgfusepath{fill}%
\end{pgfscope}%
\begin{pgfscope}%
\pgfpathrectangle{\pgfqpoint{2.893540in}{0.169444in}}{\pgfqpoint{0.858895in}{1.282944in}}%
\pgfusepath{clip}%
\pgfsetbuttcap%
\pgfsetmiterjoin%
\definecolor{currentfill}{rgb}{0.949020,0.966321,0.923645}%
\pgfsetfillcolor{currentfill}%
\pgfsetlinewidth{0.000000pt}%
\definecolor{currentstroke}{rgb}{0.000000,0.000000,0.000000}%
\pgfsetstrokecolor{currentstroke}%
\pgfsetstrokeopacity{0.000000}%
\pgfsetdash{}{0pt}%
\pgfpathmoveto{\pgfqpoint{3.408878in}{0.774567in}}%
\pgfpathlineto{\pgfqpoint{3.426056in}{0.774567in}}%
\pgfpathlineto{\pgfqpoint{3.426056in}{0.798325in}}%
\pgfpathlineto{\pgfqpoint{3.408878in}{0.798325in}}%
\pgfpathlineto{\pgfqpoint{3.408878in}{0.774567in}}%
\pgfpathclose%
\pgfusepath{fill}%
\end{pgfscope}%
\begin{pgfscope}%
\definecolor{textcolor}{rgb}{0.000000,0.000000,0.000000}%
\pgfsetstrokecolor{textcolor}%
\pgfsetfillcolor{textcolor}%
\pgftext[x=3.237099in,y=0.169444in,,]{\color{textcolor}\setmainfont{Lato}\rmfamily\fontsize{8.000000}{9.600000}\selectfont 1}%
\end{pgfscope}%
\begin{pgfscope}%
\definecolor{textcolor}{rgb}{0.000000,0.000000,0.000000}%
\pgfsetstrokecolor{textcolor}%
\pgfsetfillcolor{textcolor}%
\pgftext[x=3.237099in,y=0.311994in,,]{\color{textcolor}\setmainfont{Lato}\rmfamily\fontsize{8.000000}{9.600000}\selectfont 2}%
\end{pgfscope}%
\begin{pgfscope}%
\definecolor{textcolor}{rgb}{0.000000,0.000000,0.000000}%
\pgfsetstrokecolor{textcolor}%
\pgfsetfillcolor{textcolor}%
\pgftext[x=3.237099in,y=0.739642in,,]{\color{textcolor}\setmainfont{Lato}\rmfamily\fontsize{8.000000}{9.600000}\selectfont 5}%
\end{pgfscope}%
\begin{pgfscope}%
\definecolor{textcolor}{rgb}{0.000000,0.000000,0.000000}%
\pgfsetstrokecolor{textcolor}%
\pgfsetfillcolor{textcolor}%
\pgftext[x=3.237099in,y=1.167290in,,]{\color{textcolor}\setmainfont{Lato}\rmfamily\fontsize{8.000000}{9.600000}\selectfont 8}%
\end{pgfscope}%
\begin{pgfscope}%
\definecolor{textcolor}{rgb}{0.000000,0.000000,0.000000}%
\pgfsetstrokecolor{textcolor}%
\pgfsetfillcolor{textcolor}%
\pgftext[x=3.237099in,y=1.452389in,,]{\color{textcolor}\setmainfont{Lato}\rmfamily\fontsize{8.000000}{9.600000}\selectfont 10}%
\end{pgfscope}%
\end{pgfpicture}%
\makeatother%
\endgroup%
 \\
\footnotesize{Source: Bureau of Labor Statistics}
\end{minipage}
\newpage
\begin{minipage}{0.76\textwidth}
\subsubsection*{Jobless Claims}
\small \input{text/icsa.txt}

\input{text/ccsa.txt}
\end{minipage}
\vspace{1mm}

\begin{minipage}{0.345\textwidth}
\normalsize \textbf{New UI Claims}\\
\footnotesize{\textit{initial claims per week, in millions,}}\\
\footnotesize{\textit{not seasonally adjusted}}\\
\hspace*{-2mm} \begin{tikzpicture}
	\begin{axis}[\bbar{y}{0}, \ltdateaxisticks ytick={0, 1, 2, 3, 4, 5, 6}, 
		enlarge y limits={0.05}, ymin=0.25, xmin={2018-01-01},
		xticklabel={`\short{\year}}, height=7.2cm, width=6.5cm]
	\rbars
	\thinline{blue!50!purple!80!black}{date}{pua_ic}{data/fed_uic.csv}
	\stdline{cyan!80!blue}{date}{VALUE}{data/icsa.csv}
	\stdnode{4.2cm}{0.35cm}{\footnotesize \color{blue!50!purple!80!black}PUA}
	\stdnode{3.2cm}{0.7cm}{\footnotesize \color{cyan!80!blue}State}
	\end{axis}
\end{tikzpicture}\\
\footnotesize{Source: Department of Labor} \hfill \tbllink{icsa.csv}
\end{minipage} \hspace{9mm}
\begin{minipage}{0.35\textwidth}
\normalsize \textbf{Continued UI Claims}\\
\footnotesize{\textit{insured unemployed, in millions,}}\\
\footnotesize{\textit{not seasonally adjusted}}\\
\hspace*{-2mm} \begin{tikzpicture}
	\begin{axis}[\bbar{y}{0}, \ltdateaxisticks ytick={0, 5, 10, 15, 20}, 
		enlarge y limits={0.05}, ymin=1, xmin={2018-01-01},
		yticklabel style={text width=1.0em},
		xticklabel={`\short{\year}}, height=7.2cm, width=6.5cm]
	\rbars
	\thinline{green!50!blue}{date}{fed_cc}{data/fed_uic.csv}
	\stdline{green!90!blue}{date}{VALUE}{data/ccsa.csv}
	\stdnode{4.9cm}{4.2cm}{\footnotesize \color{green!50!blue}PUA+\\ \footnotesize \color{green!50!blue}PEUC}
	\stdnode{3.0cm}{0.8cm}{\footnotesize \color{green!90!blue}State}
	\end{axis}
\end{tikzpicture}\\
\footnotesize{Source: Department of Labor} \hfill \tbllink{ccsa.csv}
\end{minipage}
\vspace{1mm}

\begin{minipage}{0.76\textwidth}
\small In response to the COVID-19 pandemic, traditional state-run unemployment insurance was temporarily boosted by federal programs that expanded eligibility for benefits and increased the amount of benefit payments. These programs were ended on September 6, 2021, reducing the income of millions of families. 

%\input{text/fed_uic.txt}

\end{minipage}
\newpage
\hypertarget{labp}{}
\begin{minipage}{0.76\textwidth}
\subsubsection*{Labor Force Participation}
\small People who currently have a job, are looking for a job, or are on temporary layoff are all considered labor force participants. The share of the US population that participates in the labor force at a given point in time, or the \textbf{labor force participation rate}, is affected by many factors, including demographics and short- and long-term economic conditions.

Over the past two decades, the age composition of the population has lowered labor force participation rates, all else equal. \input{text/lfpr_adj.txt}

\input{text/lfpr_text.txt}
\vspace{2mm}

\normalsize \textbf{Labor Force Participation Rate}\\
\footnotesize{\textit{labor force as share of age 16 and older population, percent\\demographically adjusted rates are reweighted to the January 2000 age and sex distribution}}\\
\hspace*{-2mm} \begin{tikzpicture}
	\begin{axis}[\bbar{y}{0}, \dateaxisticks ytick={50, 55, 60, 65, 70, 75, 80}, 
		yticklabel style={text width=1.0em}, 
		width=12.25cm, height=9.2cm,
		xticklabel={`\short{\year}}, clip=false]
	\rbars
	\stdline{green!70!lime}{date}{ADJWGT_Total}{data/lfpr_adj.csv}
	\stdline{green!70!blue}{date}{TotLFPR}{data/lfpr.csv}
	\stdline{orange!40!yellow}{date}{ADJWGT_Women}{data/lfpr_adj.csv}
	\stdline{orange!90!red}{date}{WomenLFPR}{data/lfpr.csv}
	\stdline{blue!40!cyan}{date}{ADJWGT_Men}{data/lfpr_adj.csv}
	\stdline{blue!90!cyan}{date}{MenLFPR}{data/lfpr.csv}
	\input{text/lfpr_nodes.txt}
	\input{text/lfpr_nodes_adj.txt}
	\absnode{{1994-01-01}}{68.2}{\footnotesize \color{green!70!blue}\textbf{Total}}
	\absnode{{2012-01-01}}{67}{\footnotesize \color{green!70!lime}\textbf{Total (Adjusted)}}
	\absnode{{1998-01-01}}{74.3}{\footnotesize \color{blue!90!cyan}\textbf{Men}}
	\absnode{{2012-01-01}}{74.0}{\footnotesize \color{blue!40!cyan}\textbf{Men (Adjusted)}}
	\absnode{{1997-01-01}}{59.4}{\footnotesize \color{orange!90!red}\textbf{Women}}
	\absnode{{2011-01-01}}{61.2}{\footnotesize \color{orange!40!yellow}\textbf{Women (Adjusted)}}
	\end{axis}
\end{tikzpicture}\\
\footnotesize{Sources: BLS (Latest: \input{text/lfpr_blsdate.txt}\unskip); CPS (Latest: \input{text/lfpr_cpsdate.txt}\unskip)} \hfill \tbllink{lfpr.csv}
\end{minipage}
\newpage
\begin{minipage}{0.76\textwidth}
\subsubsection*{Reasons for Labor Force Non-participation}
\vspace{-0.5mm}

\small The Current Population Survey (CPS) asks those who are not employed or looking for work about their major activities and \textbf{reasons for not participating in the labor market}. \href{https://www.frbatlanta.org/chcs/human-capital-currents/2015/0612-measuring-labor-market-status-using-basic-data.aspx}{Answers} vary by age in intuitive ways, and are influenced by labor market conditions. 
\end{minipage}

\begin{minipage}{0.28\textwidth}
\normalsize \textbf{Nonparticipants}\\
\footnotesize{\textit{share of age 16+ population, percent}}\\
\hspace*{-3mm} \begin{tikzpicture}
  	\begin{axis}[\barplotnogrid axis y line=left, \barylab{1.8cm}{1.5ex}
    	width=3.6cm, bar width=1.8ex, height=5.6cm,
    	enlarge y limits={abs=0.3cm}, enlarge x limits=0.05, \bbar{x}{0},
		yticklabels from table={\nilf}{name}, xmin=0, 
		yticklabel style={font=\footnotesize}, nodes near coords, 
		nodes near coords style={/pgf/number format/.cd, fixed zerofill,
			precision=1, assume math mode},
		legend style={text=black!70, at={(0.26, 1.14)}, anchor=north, legend columns=-1, 
				fill=none, draw=none,
		        /tikz/every even column/.append style={column sep=0.4cm}}]
  	\addplot[fill=green!80!blue!60!lime, draw=none] 
  		table [y expr=-\coordindex, x index=2] {\nilf};
  	\addplot[fill=green!80!blue!72!black, draw=none] 
  		table [y expr=-\coordindex, x index=1] {\nilf};
  	\legend{ \input{text/nilf_mo2.txt}, \input{text/nilf_mo1.txt}}
  	\end{axis}
\end{tikzpicture}\\
\footnotesize{Source: Author's CPS calculations}
\end{minipage} \hspace{4mm}
\begin{minipage}{0.44\textwidth}
\small \input{text/nilfbasic.txt}
\end{minipage}
\vspace{1mm}

\begin{minipage}{0.76\textwidth}
\small While the recession of 2001 appears mild in measures of expenditure, it was followed by a substantial reduction in the share of the population receiving labor income. The economy was losing jobs at an alarming rate long after the 2001 recession had officially ended, though labor market weakness was \href{https://www.aeaweb.org/articles?id=10.1257/jep.30.2.179}{partially masked} by a major housing bubble. Seven years after the recession of 2001, the housing bubble collapsed, causing the great recession, which pushed many more people out of the labor force.
\end{minipage}

\begin{minipage}{0.29\textwidth}
\small \input{text/nilf_01.txt} 
\end{minipage} \hspace{5mm}
\begin{minipage}{0.42\textwidth}
\normalsize \textbf{Contribution to Labor Force Participation}\\
\footnotesize{\textit{cumulative percentage point contribution to age 18-64}}\\
\footnotesize{\textit{labor force participation rate since March 2001}}\\
\hspace*{-3mm} \begin{tikzpicture}
	\begin{axis}[\bbar{y}{0}, \dateaxisticks ytick={}, width=7.2cm, height=8.2cm,
		xticklabel={`\short{\year}}, clip=false, enlarge y limits={0.02}, 
		yticklabel style={text width=1.9em}]
	\rebars
	\stdline{violet!70!purple}{date}{Demographics}{data/nilf.csv}
	\thickline{blue}{date}{Discouraged}{data/nilf.csv}
	\thickline{green!80!blue}{date}{Disabled/Ill}{data/nilf.csv}
	\thickline{red}{date}{Family}{data/nilf.csv}
	\thickline{orange}{date}{Retired}{data/nilf.csv}
	\thickline{cyan}{date}{School}{data/nilf.csv}
	\input{text/nilf_nodes.txt}
	\absnode{{2013-01-1}}{-0.75}{\footnotesize \color{green!80!blue}\textbf{Disabled/Ill}}
	\absnode{{2009-10-01}}{-0.2}{\footnotesize \color{blue}\textbf{Discouraged}}
	\absnode{{2013-06-01}}{0.32}{\footnotesize \color{red}\textbf{Family/Care}}
	\absnode{{2015-01-01}}{0.78}{\footnotesize \color{orange}\textbf{Retired}}
	\absnode{{2016-01-01}}{-1.58}{\footnotesize \color{cyan}\textbf{School}}	
	\absnode{{2017-08-01}}{-1.18}{\footnotesize \color{violet!80!purple}\textbf{Age/Sex}}	
	\end{axis}
\end{tikzpicture}\\
\footnotesize{Source: Author's Calculations from CPS} \hfill \tbllink{nilf.csv}

\footnotesize Series in the chart are adjusted so that the distribution of the age 18--64 population by age and sex is constant and equal to its March 2001 value. The total effect of this adjustment on labor force participation is included separately in the chart, as \color{violet!80!purple}\textbf{Age/Sex}.
\end{minipage}
\newpage
\begin{minipage}{0.76\textwidth}
\small Young people's participation in the labor market, by working or looking for work, is affected by trends in educational attainment and by economic conditions. & Dec  2023 & Nov  2023 & Oct  2023 & Dec  2022 &2019&2015&2010&1994\\  \hspace{1mm}Total,  19  to  21 &26.5&26.6&26.6&27.0&27.1&27.0&24.5&15.4\\  \hspace{3mm}Men   &26.5&26.5&26.5&26.0&26.6&27.0&24.2&15.8\\  \hspace{3mm}Women   &26.5&26.6&26.7&27.9&27.5&27.0&24.7&15.1\\  \hspace{1mm}Total,  22  to  24 &11.2&11.2&11.2&11.5&11.5&11.7&10.6&6.5\\  \hspace{3mm}Men &11.5&11.4&11.3&11.2&11.7&12.0&10.8&6.7\\  \hspace{3mm}Women &10.9&11.0&11.2&11.8&11.4&11.3&10.4&6.4\\  \hspace{1mm}Total,  25  to  27 &4.4&4.4&4.3&4.5&4.6&5.2&4.6&2.9\\  \hspace{3mm}Men     &4.1&4.1&4.1&4.4&4.7&5.0&4.4&2.5\\  \hspace{3mm}Women     &4.7&4.6&4.6&4.6&4.4&5.4&4.7&3.2\\ 
\vspace{2mm}

\normalsize \textbf{Reason for Labor Force Non-participation, by Age}\\
\footnotesize{\textit{share of age group population, percent, 12-month moving average}}

\normalsize \hspace{18mm} Age 19--21 \hspace{44mm} Age 22--24 \footnotesize
\vspace{-1mm}

\hspace*{-2mm} \begin{tikzpicture}
	\begin{axis}[\bbar{y}{0}, \shdateaxisticks 
		clip=false, width=6.0cm, height=7.3cm, 
		xticklabel={`\short{\year}}, enlarge y limits={0.01}, 
        enlarge x limits={0.02}, ymax=43, ymin=0, legend columns=1, 
        legend cell align={left},
		legend style={fill=white, at={(0.43, 1.0)}}]
	\rebars
	\thickline{orange!90!yellow}{date}{Total_19-21}{data/young_school.csv}
	\stdline{blue!50!cyan}{date}{School_19-21}{data/young_school.csv}
	\stdline{violet!85!black}{date}{Other_19-21}{data/young_school.csv}
	\input{text/young_nilf19-21_node.txt}
	\legend{Total, School, Other};
	\end{axis}
\end{tikzpicture}
\hspace{5mm}
\begin{tikzpicture}
	\begin{axis}[\bbar{y}{0}, \shdateaxisticks 
		clip=false, width=6.0cm, height=7.3cm, 
        xticklabel={`\short{\year}}, enlarge y limits={0.01}, 
        enlarge x limits={0.02}, ymax=43, ymin=0]
	\rebars
	\thickline{orange!90!yellow}{date}{Total_22-24}{data/young_school.csv}
	\stdline{blue!50!cyan}{date}{School_22-24}{data/young_school.csv}
	\stdline{violet!85!black}{date}{Other_22-24}{data/young_school.csv}
	\input{text/young_nilf22-24_node.txt}
	\end{axis}
\end{tikzpicture}\\
\footnotesize{Source: Author's Calculations from CPS} \hfill \tbllink{young_school.csv}
\end{minipage}

\newpage
\hypertarget{labf}{}
\begin{minipage}{0.76\textwidth}
\subsection*{Labor Force Flows}
\small The current population survey interviews households up to eight times over 16 months, allowing insight into the labor force status of the same individual over time, and in particular, into \textbf{flows between employment, unemployment, and other categories}. The Bureau of Labor Statistics \href{https://www.bls.gov/cps/cps_flows.htm}{publish} many monthly indicators based on labor force flows, and others can be calculated directly from the public use data.

Among newly-employed workers, some were looking for work or otherwise considered unemployed the prior month, while others were not in the labor force. \input{text/lf_flow.txt} 
\vspace{2mm}

\normalsize \textbf{Newly Employed, Not Previously Looking For Work}\\
\footnotesize{\textit{share of newly employed that were not looking for work in the prior month}}\\
\hspace*{-2mm} \begin{tikzpicture}
	\begin{axis}[\bbar{y}{0}, \dateaxisticks ytick={40, 50, 60, 70}, 
		yticklabel style={text width=1.0em}, xticklabel={`\short{\year}}, 
		clip=false, height=5.4cm, width=12.2cm,
		legend columns=1, legend cell align={left}, reverse legend,
		legend style={at={(0.43, 1.03)}}]
	\rbars
	\stdline{lime!80!green}{date}{total}{data/lf_flow.csv}
	\thickline{green!60!teal!80!black}{date}{quarterly}{data/lf_flow_q.csv}
	\input{text/lf_flow_node.txt}
	\legend{Monthly, 3-Month Moving Average};
	\end{axis}
\end{tikzpicture}\\
\footnotesize{Source: Bureau of Labor Statistics} \hfill \tbllink{lf_flow.csv} \ \tbllink{lf_flow_q.csv}
\vspace{4mm}

\small The great recession worsened job-finding prospects for those not in the labor force (NILF) due to disability or illness. As a result, the flow into employment for people age 25 to 54 who are out of the labor force due to a disability or illness slowed considerably. These prospects first recovered to pre-recession levels around 2017. 
\end{minipage}

\begin{minipage}{0.28\textwidth}
\small \input{text/disflow.txt}
\end{minipage} \hspace{5mm} \begin{minipage}{0.44\textwidth}
\normalsize \textbf{Flow, Disability to Work}\\
\footnotesize{\textit{NILF disability/illness, age 25 to 54, }}\\
\footnotesize{\textit{share employed one year later}}\\
\hspace*{-2mm} \begin{tikzpicture}
	\begin{axis}[\bbar{y}{0}, \shdateaxisticks ytick={6, 8}, clip=false,
		xticklabel={`\short{\year}}, height=4.2cm, width=7.3cm]
	\rebars
	\thickline{blue}{date}{Share}{data/disflow.csv}
	\input{text/disflow_node.txt}
	\end{axis}
\end{tikzpicture}\\
\footnotesize{Source: Author's Calculations from CPS} \hfill \tbllink{disflow.csv}
\end{minipage}
\newpage
\begin{minipage}{0.76\textwidth}
\subsubsection*{Job Openings and Labor Turnover Survey}
\small Some types of \textbf{turnover} in the labor market are healthy and mean people are better able find a new job if they do not like the one they have. Additionally, the job prospects outside of a firm affect the bargaining power of the workers inside of the firm. The Bureau of Labor Statistics \href{https://www.bls.gov/news.release/pdf/jolts.pdf}{report} the number of job openings, hires, and separations in several industry groups on a monthly basis. Separations include layoffs, voluntarily leaving a job (\textit{quits}), and other separations such as retirements, transfers to other locations, or separations due to death or disability. 

\small \input{text/jolts2.txt}
\vspace{1mm}

\normalsize \textbf{Job Turnover}\\
\footnotesize{\textit{job openings, hires, quits, and layoffs per month, in millions}}
\vspace{4.0cm}

\hspace{4mm} \begin{tikzpicture}[overlay]
	\begin{axis}[\bbar{y}{0}, \dateaxisticks ytick={0, 2, 4, 6, 8, 10, 12, 14}, 
		yticklabel style={text width=1.2em}, xmin={2001-01-01},
		enlarge y limits={0.06}, ymin=0.6, ymax=11.4, clip=false,
		xticklabel={`\short{\year}}, height=6.0cm, width=12.4cm]
	\rbars
	\stdline{orange!75!yellow}{date}{Layoffs}{data/jolts.csv}
	\stdline{cyan}{date}{Hires}{data/jolts.csv}
	\stdline{green!65!black}{date}{Quits}{data/jolts.csv}
	\stdline{blue!90!black}{date}{Openings}{data/jolts.csv}
	\absnode{{2005-01-01}}{6.6}{\footnotesize \color{cyan}\textbf{Hires}}
	\absnode{{2015-01-01}}{7.6}{\footnotesize \color{blue!90!black}\textbf{Openings}}
	\absnode{{2015-01-01}}{4.2}{\footnotesize \color{green!65!black}\textbf{Quits}}
	\absnode{{2004-08-01}}{1.8}{\footnotesize \color{orange!75!yellow}\textbf{Layoffs}}
	\input{text/jolts_nodes.txt}
	\end{axis}
\end{tikzpicture}

\footnotesize{Source: Bureau of Labor Statistics} \hfill \tbllink{jolts.csv} 
\vspace{3mm}

\small The number of people who voluntarily separate from a job in a given month, divided by the total number employed is the \textbf{quits rate}. The rate typically increases when workers are confident enough to leave one job for another one, and a high quits rate, particularly in low-paying industries, can be a sign of a tight labor market.
\end{minipage}
\vspace{1mm}

\begin{minipage}{0.4\textwidth}
\normalsize \textbf{Quits Rate}\\
\footnotesize{\textit{voluntary separations as share of workforce}}\\
\hspace*{-2mm} \begin{tikzpicture}
	\begin{axis}[\bbar{y}{0}, \shdateaxisticks ytick={0, 2, 4, 6}, clip=false,
		xticklabel={`\short{\year}}, ymin=0.3, height=5.6cm, width=6.7cm]
	\rebars
	\thinline{red!50!purple}{date}{AFS_QU}{data/quits.csv}
	\absnode{{2010-10-01}}{6}{\scriptsize \color{red!50!purple}\textbf{Accommodations}\\
	\scriptsize \color{red!50!purple}\textbf{\& Food Services}}
	\input{text/quits_afs_node.txt}
	\stdline{violet!90!black}{date}{TOT_QU}{data/quits.csv}
	\absnode{{2013-01-01}}{1.6}{\scriptsize \color{violet}\textbf{Total nonfarm}}
	\input{text/quits_tot_node.txt}
	\end{axis}
\end{tikzpicture}\\
\footnotesize{Source: Bureau of Labor Statistics} \hfill \tbllink{quits.csv}
\end{minipage}
\hspace{6mm}
\begin{minipage}{0.31\textwidth}
\small The quits rate is cyclical within the accommodations and food services industries (which includes restaurants), and tends to rise when a tight labor market pulls people out of restaurant jobs and into higher paying jobs in other industries. 

\input{text/quits_afs.txt}
\end{minipage}
\newpage
\begin{minipage}{0.76\textwidth}
\small A high ratio of job openings to unemployed indicates a tight labor market, for example from low levels of unemployment, or if completing a new hire is taking more time. \input{text/unjo.txt}
\vspace{1mm}

\normalsize \textbf{Job Openings Per Unemployed Person}\\
\footnotesize{\textit{job openings divided by total unemployment}}\\
\hspace*{-2mm} \begin{tikzpicture}
	\begin{axis}[\bbar{y}{0}, \dateaxisticks ytick={0, 0.5, 1.0, 1.5, 2.0}, 
		yticklabel style={text width=1.3em}, 
		ymin=0.08, ymax=1.9,
		xticklabel={`\short{\year}}, height=4.4cm, width=12.4cm, clip=false]
	\rebars
	\thickline{violet}{date}{unjo}{data/unjo.csv}
	\input{text/unjo_node.txt}
	\end{axis}
\end{tikzpicture}\\
\footnotesize{Source: Bureau of Labor Statistics} \hfill \tbllink{unjo.csv}
\end{minipage}
\vspace{5mm}

\begin{minipage}{0.42\textwidth}
\noindent \normalsize \textbf{Monthly Quits Rate by Industry}\\
\footnotesize{\textit{share of employment, percent}}\\
\vspace{11.4cm}

\hspace{20mm} \begin{tikzpicture}[overlay]
	\begin{axis}[clip=false, stack negative=separate, ymajorgrids=false,
            yticklabels from table={\quits}{name},
			y axis line style={opacity=0}, x axis line style={opacity=0}, 
			yticklabel style={black, font=\scriptsize, 
			align=left, anchor=east}, nodes near coords align={above},    
			nodes near coords style={/pgf/number format/.cd,fixed zerofill, 
				precision=1, assume math mode},  
            ytick=data, tickwidth=0pt, xmin=0, bar width=1.8ex, 
            enlarge y limits={0.02}, enlarge x limits={0.02}, legend cell align={left}, 
            width=6.0cm, height=12.8cm, legend style={legend columns=2, at={(0.94, 1.09)}},
            \barylab{2.4cm}{1.2ex}]
	\addplot[xbar stacked, bar shift=0pt, fill=black!3!white, draw=none] 
		table [y expr=-\coordindex, x index=1] {\quits};
	\addplot[xbar stacked, bar shift=0pt, fill=cyan!65!blue!30!white, draw=cyan!65!blue!30!white, area legend] 
		table [y expr=-\coordindex, x index=2] {\quits};	
	\addplot[xbar stacked, bar shift=0pt, fill=blue!90!purple!85!white, draw=cyan!65!blue!30!white, area legend] 
		table [y expr=-\coordindex, x index=3] {\quits};
	\addplot[xbar stacked, bar shift=0pt, fill=cyan!65!blue!30!white, draw=cyan!65!blue!30!white, area legend] 
		table [y expr=-\coordindex, x index=4] {\quits};
	\addplot[xbar stacked, bar shift=0pt, fill=black!3!white, draw=none] 
		table [y expr=-\coordindex, x index=5] {\quits};	
	\addplot[only marks, mark=square*, mark options={fill=white, 
		draw=blue!90!purple!85!white, scale=1.0}] 
		table [y expr=-\coordindex, x index=7] {\quits};        
	\addplot[only marks, mark=diamond*, mark options={fill=white, 
		draw=red, scale=1.8}] 
		table [y expr=-\coordindex, x index=6] {\quits};
	\addplot [nodes near coords={\pgfmathprintnumber{\pgfkeysvalueof{/data point/x}}}, only marks, mark=none] table[y=y, x index=6] {\quits};
	\legend{, \ 2000-- range, \ March 2020-- range,, , \ 3-month avg., \ \input{text/jolts_ltdate.txt}}
	\end{axis}
\end{tikzpicture}
\vspace{2mm}

\footnotesize{Source: Bureau of Labor Statistics} \hfill \tbllink{quits_ind.csv}
\end{minipage}\hspace{12mm}
\begin{minipage}{0.42\textwidth}
\noindent \normalsize \textbf{Monthly Job Openings Rate by Industry}\\
\footnotesize{\textit{share of employment, percent}}\\
\vspace{11.4cm}

\hspace{20mm} \begin{tikzpicture}[overlay]
	\begin{axis}[clip=false, stack negative=separate, ymajorgrids=false,
            yticklabels from table={\opens}{name},
			y axis line style={opacity=0}, x axis line style={opacity=0}, 
			yticklabel style={black, font=\scriptsize, align=left, anchor=east},  
			nodes near coords align={above},    
			nodes near coords style={/pgf/number format/.cd,fixed zerofill, 
				precision=1, assume math mode},  
            ytick=data, tickwidth=0pt, xmin=0,bar width=1.8ex, 
            enlarge y limits={0.02}, enlarge x limits={0.02}, legend cell align={left}, 
            width=6.0cm, height=12.8cm, legend style={legend columns=2, at={(0.94, 1.09)}},
            \barylab{2.4cm}{1.2ex}]
	\addplot[xbar stacked, bar shift=0pt, fill=black!3!white, draw=none] 
		table [y expr=-\coordindex, x index=1] {\opens};
	\addplot[xbar stacked, bar shift=0pt, fill=brown!28!white, draw=brown!28!white, area legend] 
		table [y expr=-\coordindex, x index=2] {\opens};	
	\addplot[xbar stacked, bar shift=0pt, fill=orange!80!yellow, draw=brown!28!white, area legend] 
		table [y expr=-\coordindex, x index=3] {\opens};
	\addplot[xbar stacked, bar shift=0pt, fill=brown!28!white, draw=brown!28!white, area legend] 
		table [y expr=-\coordindex, x index=4] {\opens};
	\addplot[xbar stacked, bar shift=0pt, fill=black!3!white, draw=none] 
		table [y expr=-\coordindex, x index=5] {\opens};	
	\addplot[only marks, mark=square*, mark options={fill=white, 
		draw=orange!80!yellow, scale=1.0}] 
		table [y expr=-\coordindex, x index=7] {\opens};         
	\addplot[only marks, mark=diamond*, mark options={fill=white, 
		draw=violet!70!black, scale=1.8}] 
		table [y expr=-\coordindex, x index=6] {\opens};
	\addplot [nodes near coords={\pgfmathprintnumber{\pgfkeysvalueof{/data point/x}}}, only marks, mark=none] table[y=y, x index=6] {\opens};
	\legend{, \ 2000-- range, \ March 2020-- range,, , \ 3-month avg., \ \input{text/jolts_ltdate.txt}}
	\end{axis}
\end{tikzpicture}
\vspace{2mm}

\footnotesize{Source: Bureau of Labor Statistics} \hfill \tbllink{openings_ind.csv}
\end{minipage}
\newpage
\begin{minipage}{0.76\textwidth}
\subsubsection*{Job Switching}
\small Job switching is important for getting people into the jobs where they are most productive. Individuals boost labor productivity by switching to a more productive industry or moving from a less-productive firm to a more-productive firm. 

The current population survey asks whether individuals have the same employer as they did the previous month. The rate at which people say they have \textbf{changed employers} had fallen to below two percent after the great recession, from an average of around three percent during the late 1990s. 

\input{text/jobswitch.txt} 
\vspace{1mm}

\normalsize \textbf{New Employer}\\
\footnotesize{\textit{percent of workforce, seasonally adjusted}}\\
\hspace*{-2mm} \begin{tikzpicture}
	\begin{axis}[\dateaxisticks ytick={1.5, 2, 2.5, 3, 3.5, 4}, 
		enlarge y limits={0.05}, clip=false, legend cell align={left},
		legend style={at={(0.38, 0.32)}}, legend columns=1, 
		yticklabel style={text width=1.5em},
		xticklabel={`\short{\year}}, height=5.4cm, width=12.2cm]
	\rebars
	\stdline{green!90!blue}{date}{Monthly}{data/jobswitch.csv}
	\thickline{blue!80!black}{date}{MA3}{data/jobswitch.csv}
	\input{text/jobswitch_node.txt}
	\legend{Monthly, 3-month moving average};
	\end{axis}
\end{tikzpicture}\\
\footnotesize{Source: Author's calculations from CPS} \hfill \tbllink{jobswitch.csv} 
\vspace{5mm}

\small The monthly Survey of Consumer Expectations from the Federal Reserve Bank of New York asks people how likely they are to either \href{https://www.newyorkfed.org/microeconomics/sce#/jobsep-1}{lose or leave a job}. \textbf{Expected separations} are specifically the average perceived probability, or likelihood, of separating from a job in the next 12 months. 

\begin{minipage}{0.42\textwidth}
\small \input{text/sce_job_sep.txt}
\end{minipage}\hfill
\begin{minipage}{0.52\textwidth}
\normalsize \textbf{Expected Separations}\\
\footnotesize{\textit{mean probability of separation, percent}}\\
\hspace*{-2mm} \begin{tikzpicture}
	\begin{axis}[\ltdateaxisticks  \bbar{y}{0}, 
		enlarge y limits={0.05}, clip=false, max space between ticks=20,	
		yticklabel style={text width=1.2em}, ymin=1, ymax=28,
		xticklabel={`\short{\year}}, height=6.2cm, width=6.5cm]
	\rbar
	\stdline{cyan}{date}{Losing}{data/sce_job_separation.csv}
	\absnode{{2015-10-01}}{13}{\footnotesize \color{cyan}\textbf{Lose Job}}
	\stdline{violet!50!blue}{date}{Leaving}{data/sce_job_separation.csv}
	\absnode{{2013-10-01}}{25.3}{\footnotesize \color{violet!50!blue}\textbf{Voluntary}}
	\input{text/sce_job_sep_nodes.txt}
	\end{axis}
\end{tikzpicture}\\
\footnotesize{Source: Federal Reserve Bank of New York} \hfill \tbllink{sce_job_separation.csv} 
\end{minipage}
\end{minipage}
\newpage
\hypertarget{labh}{}
\subsection*{Hours Worked}
\begin{minipage}{0.33\textwidth}
\small The Bureau of Labor Statistics (BLS) report \textbf{hours worked per week} in both Current Employment Statistics (CES) and Labor Force Statistics (LFS).

\input{text/hours_tot.txt} 

\input{text/hours_lfs2.txt} 

\input{text/hours_ces.txt} 
\end{minipage}\hspace{5mm} \begin{minipage}{0.39\textwidth}
\normalsize \textbf{Hours Worked}\\
\footnotesize{\textit{average hours per week}}\\
\hspace*{-2mm} \begin{tikzpicture}
	\begin{axis}[\shdateaxisticks ytick={20, 25, 30, 35, 40}, 
		yticklabel style={text width=0.7em}, 
		enlarge y limits={0.06}, xmin={1989-12-30},
		xticklabel={`\short{\year}}, height=11.8cm, width=7.3cm, clip=false]
	\rbars
	\thinline{blue}{date}{TOTLFS}{data/hours.csv}
	\thinline{lime!90!green}{date}{SERVNSA}{data/hours.csv}
	\stdline{green!90!blue!70!black}{date}{SERVSA}{data/hours.csv}
	\stdline{orange}{date}{PNS}{data/hours.csv}
	\thinline{red!90!black}{date}{PTECONSA}{data/hours.csv}
	\stdnode{1.35cm}{9.5cm}{\scriptsize \color{blue}\textbf{Total, Actual (LFS)}}
	\stdnode{1.45cm}{6.2cm}{{\scriptsize \color{orange}\textbf{Production \&}}\\{\scriptsize \color{orange}\textbf{Non-Supervisory}}\\{\scriptsize \color{orange}\textbf{(CES)}}}
	\stdnode{1.5cm}{2.15cm}{{\scriptsize \color{red!90!black}\textbf{Part Time}}\\{\scriptsize \color{red!90!black}\textbf{Economic Reasons}}\\{\scriptsize \color{red!90!black}\textbf{(LFS)}}}
	\stdnode{4.45cm}{7.8cm}{{\scriptsize \color{green!90!blue!70!black}\textbf{Services}}\\{\scriptsize \color{green!90!blue!70!black}\textbf{Occupations}}\\{\scriptsize \color{green!90!blue!70!black}\textbf{(LFS)}}}
	\end{axis}
\end{tikzpicture}\\
\footnotesize{Source: BLS} \hfill \tbllink{hours.csv}
\end{minipage}
\vspace{3mm}

\begin{minipage}{0.76\textwidth}
\normalsize \textbf{Hours Worked, Various Measures}\\
\footnotesize{\textit{average hours per week, seasonally adjusted}}\\
\rowcolors{1}{}{black!5} \setlength{\tabcolsep}{3.1pt} \color{black!90}
	{\renewcommand{\arraystretch}{1.55}
		\begin{tabular}{p{52mm} R{6.7mm} R{6.7mm} R{6.7mm} R{6.7mm} R{6.7mm} R{6.7mm} 
		   R{6.7mm}}
			 \input data/hoursworked_table.tex \hline
		\end{tabular}}\vspace{-2mm}
		
\footnotesize{Source: Bureau of Labor Statistics; Author}
\end{minipage}
\newpage 
\begin{minipage}{0.76\textwidth}
\small BLS also \href{https://www.bls.gov/lpc/}{report} a quarterly index tracking aggregate hours worked in nonfarm businesses (see\cbox{cyan!80!blue}). \input{text/tot_hours.txt}
\vspace{2mm}

\normalsize \textbf{Aggregate Hours Worked}\\
\footnotesize{\textit{nonfarm businesses, quarterly growth at seasonally adjusted annual rate, percent}}\\
\hspace*{-2mm} \begin{tikzpicture}
	\begin{axis}[\bbar{y}{0}, \dateaxisticks ytick={-50, -40, -30, -20, -10, 0, 10, 20, 30, 40}, 
		yticklabel style={text width=1.5em}, width=12.8cm, height=7.2cm,
		xticklabel={`\short{\year}}, clip=false]
	\rbars
	\sbar{cyan!80!blue}{date}{hours}{data/lprod.csv}
	\end{axis}
\end{tikzpicture}\\
\footnotesize{Source: Bureau of Labor Statistics} \hfill \tbllink{lprod.csv}
\end{minipage}
\newpage
\begin{minipage}{0.76\textwidth}
\subsection*{Nonstandard Work Arrangements}
\small Many workers do not have standard work arrangements, either by choice or as the result of not being able to find standard work arrangements. Many workers work part-time, part-year, or both. Some workers have more than one job. Additionally, a portion of the workforce is considered self-employed. 
\subsubsection*{Part-time Work}
\small Many people work fewer than 35 hours per week, and the reasons for doing so vary. The Bureau of Labor Statistics classify part-time workers who would prefer full-time work as involuntary or \textbf{part time for economic reasons}. This group is comprised of people who don't have enough hours because of slack business conditions or who are unable to find full-time work. 

Voluntary part-time workers or those \textbf{part-time for non-economic reasons} do not want full-time jobs. The category includes those who work fewer hours for health, childcare, personal or family reasons, those who are retired or have a limit on earnings, and those with jobs where full-time is less than 35 hours per week. 
\vspace{0.5mm}

\begin{minipage}{0.56\textwidth}
\normalsize \textbf{Part Time, by Reason}\\
\footnotesize{\textit{percent of labor force}}\\
\hspace*{-2mm} \begin{tikzpicture}
	\begin{axis}[\bbar{y}{0}, \shdateaxisticks 
		enlarge y limits={0.1}, ymin=1.2, yticklabel style={text width=1.0em},
		xticklabel={`\short{\year}}, height=6.2cm, width=6.8cm, clip=false]
	\rebars
	\stdline{red}{date}{PTECON}{data/parttime.csv}
	\stdline{orange!80!yellow}{date}{PTNONECON}{data/parttime.csv}
	\absnode{{1994-01-01}}{12.6}{\footnotesize \color{orange}\textbf{Non-economic}}
	\absnode{{1994-01-01}}{5}{\footnotesize \color{red}\textbf{Economic}}
	\input{text/pt_nodes.txt}
	\end{axis}
\end{tikzpicture}\\
\footnotesize{Source: Bureau of Labor Statistics} \hfill \tbllink{parttime.csv} 
\end{minipage}\hfill
\begin{minipage}{0.38\textwidth}
\small \input{text/parttime.txt}
\end{minipage}
\vspace{1mm}

\subsubsection*{More Than One Job}
\small Multiple jobholding rates were fairly stable from 2000 to 2019. The household survey used to identify people with more than one job asks about a specific reference week. As a result, the multiple jobholding rate is not intending to capture people who rely on a patchwork of multiple jobs over time, but work at one job in the survey reference week. 
\vspace{0.5mm}

\begin{minipage}{0.56\textwidth}
\normalsize \textbf{Multiple Jobholders}\\
\footnotesize{\textit{percent of workers, 3-month moving average}}\\
\hspace*{-2mm} \begin{tikzpicture}
	\begin{axis}[\shdateaxisticks clip=false,
		enlarge y limits={0.16}, 
		xticklabel={`\short{\year}}, height=4.6cm, width=7.0cm]
	\rebars
	\thickline{cyan!50!blue}{date}{MJH}{data/mjh.csv}
	\input{text/mjh_node.txt}
	\end{axis}
\end{tikzpicture}\\
\footnotesize{Source: Bureau of Labor Statistics} \hfill \tbllink{mjh.csv} 
\end{minipage} \hfill
\begin{minipage}{0.38\textwidth}
\small \input{text/mjh.txt}
\end{minipage}
\end{minipage}
\newpage
\begin{minipage}{0.76\textwidth}
\subsubsection*{Self-employment}
\small Workers are considered \textbf{self-employed} if they work for profit or fees in their own business, profession, trade, or farm. Some self-employed have incorporated their business, and are similar to wage and salary workers in that they are paid by their business. Self-employment can offer more flexibility than traditional jobs, in some cases, but can also be less stable. The category includes people who work for profit but do not make any profits, for example. 

\input{text/selfemp1.txt}
\vspace{0.5mm}

\begin{minipage}{0.56\textwidth}
\normalsize \textbf{Self Employed}\\
\footnotesize{\textit{percent of labor force}}\\
\hspace*{-2mm} \begin{tikzpicture}
	\begin{axis}[\bbar{y}{0}, \shdateaxisticks clip=false,
		enlarge y limits={0.18}, ymin=1.2,
		xticklabel={`\short{\year}}, height=4.4cm, width=7.0cm]
	\rbars
	\stdline{purple}{date}{uninc}{data/selfemp.csv}
	\stdline{green!80!black}{date}{inc}{data/selfemp.csv}
	\stdline{green!80!black}{date}{INC}{data/selfemp.csv}
	\absnode{{2009-06-01}}{7.85}{\footnotesize \color{purple}\textbf{Unincorporated}}
	\absnode{{2010-01-01}}{3.4}{\footnotesize \color{green!80!black}\textbf{Incorporated}}
	\input{text/selfemp_nodes.txt}
	\end{axis}
\end{tikzpicture}\\
\footnotesize{Source: Bureau of Labor Statistics, Author} \hfill \tbllink{selfemp.csv} 
\end{minipage} \hfill
\begin{minipage}{0.38\textwidth}
\small \input{text/selfemp2.txt}
\end{minipage}
\end{minipage}
\newpage
\begin{minipage}{0.76\textwidth}
\subsection*{Wages}
\hypertarget{labw}{}
\small The \textbf{usual wages of full-time workers} can be measured at various points in the income distribution using the Current Population Survey. BLS report these data by decile and quartile, with the most commonly used measure being the median usual weekly earnings.

\input{text/uwe_basic.txt}
\vspace{3mm}

\normalsize \textbf{Weekly Earnings, First Decile}\\
\footnotesize{\textit{full-time, wage and salary earners, age 16+}} \hspace{10mm} \small 	
	\colorline{lime!65!green}{ \ bd CPS} \hspace{3mm} \colorline{blue!65!black}{ \ BLS}\\
\hspace*{-2mm} \begin{tikzpicture}
	\begin{axis}[\bbar{y}{0}, \dateaxisticks ytick={200, 300, 400, 500, 600}, 
		enlarge y limits={0.05},
		clip=false, width=6.35cm, height=5.2cm, xmin=1988-12-15, 
		yticklabel style={text width=1.5em}, 
		xtick={{1989-01-01}, {2000-01-01}, {2010-01-01}, {2020-01-01}},
        minor xtick={}, xticklabels={`89, `00, `10, `20}]
	\rbars
	\stdline{lime!65!green}{date}{All}{data/uwe_bd.csv}
	\stdline{blue!65!black}{date}{p10uwe}{data/uwe_bls.csv}
	\node[right] at (axis cs:1989-06-01,515) {\footnotesize Levels, USD};
	\end{axis}
\end{tikzpicture} \hfill
\begin{tikzpicture}
	\begin{axis}[\bbar{y}{0}, \dateaxisticks ytick={-5, 0, 5, 10}, enlarge y limits={0.05},
		clip=false, width=6.5cm, height=5.2cm, xmin=1989-12-15,
		xtick={{1990-01-01}, {2000-01-01}, {2010-01-01}, {2020-01-01}},
        minor xtick={}, xticklabels={`90, `00, `10, `20}]
	\rbars
	\stdline{lime!65!green}{date}{All}{data/uwe_bd_gr.csv}
	\stdline{blue!65!black}{date}{p10uwe}{data/uwe_bls_gr.csv}
	\node[right] at (axis cs:1990-01-01,10.8) {\footnotesize One-year growth, percent};
	\end{axis}
\end{tikzpicture}\\
\footnotesize{Source: Bureau of Labor Statistics, Author's Calculations} \hfill \tbllink{uwe_bd.csv}*
\vspace{5mm}

\normalsize \textbf{Usual Weekly Earnings}\\
\footnotesize{\textit{full-time, wage and salary earners, age 16+, nominal USD}}\\
\rowcolors{1}{}{black!5} \setlength{\tabcolsep}{3.1pt} \color{black!90}
		{\renewcommand{\arraystretch}{1.55}
		 \begin{tabular}{p{19mm} R{8.8mm} R{8.8mm} R{8.8mm} R{8.8mm} R{8.8mm} R{8.8mm} 
		   R{8.8mm} R{8.8mm} R{8.8mm} R{8.8mm}}
			 \input data/wage_dist_bls2.tex \hline
		\end{tabular}}
\vspace{-2mm}
		
\footnotesize{Source: Bureau of Labor Statistics}
\vspace{4mm}

\normalsize \textbf{Weekly Earnings Growth}\\
\footnotesize{\textit{full-time, wage and salary earners, age 16+, one-year growth, percent}}\\
\rowcolors{1}{}{black!5} \setlength{\tabcolsep}{3.1pt} \color{black!90}
	{\renewcommand{\arraystretch}{1.55}
		\begin{tabular}{p{19mm} R{8.8mm} R{8.8mm} R{8.8mm} R{8.8mm} R{8.8mm} R{8.8mm} 
		   R{8.8mm} R{8.8mm} R{8.8mm} R{8.8mm}}
			 \input data/wage_dist_bls.tex \hline
		\end{tabular}}\vspace{-2mm}
				
\footnotesize{Source: Bureau of Labor Statistics}
\end{minipage}
\newpage
\begin{minipage}{0.76\textwidth}
\subsubsection*{Nominal Hourly Wages}
\small \input{text/ahe_summary.txt}
\vspace{2mm}

\normalsize \textbf{Average Hourly Earnings}\\
\footnotesize{\textit{nominal one-year growth, percent}}\\
\hspace*{-2mm} \begin{tikzpicture}
	\begin{axis}[\bbar{y}{0}, \dateaxisticks ytick={0, 2, 4, 6, 8}, 
		xticklabel={`\short{\year}}, clip=false, ymin=0, enlarge y limits={0.01}, 
		legend columns=1, legend cell align={left}, height=6.0cm, 
		legend style={at={(0.46, 1.0)}, fill=white}]
	\rbars
	\thickline{blue!80!black}{date}{PNS}{data/ahe.csv}
	\thinline{magenta}{date}{ALL}{data/ahe.csv}
	\legend{Production and Non-supervisory, All Employees (2007-onward)};
	\end{axis}
\end{tikzpicture}\\
\footnotesize{Source: Bureau of Labor Statistics} \hfill \tbllink{ahe.csv}
\vspace{5mm}

\small \input{text/ahe_comp.txt} 
\vspace{1mm}

\normalsize \textbf{Average Hourly Earnings Growth by Industry}\\
\footnotesize{\textit{production and non-supervisory workers, as of \input{text/ahe_bar_date.txt}}}\\
\footnotesize{\textit{one-year growth, percent}}\\
\hspace*{-2mm} \begin{tikzpicture}
	\begin{axis}[\barplotnogrid axis y line=left, 
		\barylab{4.1cm}{1.5ex}
		width=6.4cm, bar width=2.8ex, 
		height=9.0cm, xtick={0}, xmajorgrids,
		enlarge y limits={abs=3mm}, enlarge x limits=0.02, 
		xmin=0, clip=false,
		yticklabels from table={\ahe}{name},
		yticklabel style={font=\footnotesize},
		nodes near coords style={style={black!80, fill=white,
		 xshift=1.0mm, inner sep=0.5}, /pgf/number format/.cd, fixed zerofill,precision=1, 
		assume math mode}]
	\addplot[fill=blue!60!cyan, draw=none, nodes near coords] table [y expr=-\coordindex, x index=1] {\ahe};
	\input{text/cpi_lt_line.txt} 
	\end{axis}
\end{tikzpicture}\\ 
\footnotesize{Source: Bureau of Labor Statistics} \hspace{48mm} \tbllink{ahe_ind.csv}
\end{minipage}
\newpage
\begin{minipage}{0.76\textwidth}
\small The US has increasingly invested in education, boosting productivity and earnings. To examine the return on investment, this section discusses the relationship between \textbf{wages and education}, over the long-term and in recent data. 

\input{text/educ_wage1.txt}
\vspace{1mm}

\normalsize \textbf{Real Earnings by Level of Education}\\
\footnotesize{\textit{median usual weekly earnings, full-time wage and salary workers age 25--54}}\\
\footnotesize{\textit{\input{text/rw_cpsdate.txt}dollars, three-month moving average}}
\vspace{4.4cm}

\hspace{7mm} \begin{tikzpicture}[overlay]
	\begin{axis}[\dateaxisticks ytick={600, 800, 1000, 1200, 1400, 1600}, 
		yticklabel style={text width=2.0em}, ymin=630,
		xticklabel={`\short{\year}}, clip=false, height=6.4cm, width=11.5cm]
	\rbars
	\thickline{blue!90!cyan}{date}{LTHS/HS}{data/educ_wage.csv}
	\thickline{green!80!blue}{date}{COLL/ADV}{data/educ_wage.csv}
	\thickline{cyan}{date}{SC}{data/educ_wage.csv}
	\thickline{orange!70!yellow}{date}{Total}{data/educ_wage.csv}
	\absnode{{1991-01-01}}{720}{\footnotesize \color{blue!90!cyan}High School Degree or Less}
	\absnode{{1991-06-01}}{1670}{\footnotesize \color{green!80!blue}Bachelor or \\\footnotesize \color{green!80!blue}Advanced\\\footnotesize \color{green!80!blue}Degree}
	\absnode{{2001-10-01}}{950}{\footnotesize \color{cyan}Some College}
	\absnode{{2012-02-01}}{1220}{\footnotesize \color{orange!70!yellow}Total,\\\footnotesize \color{orange!70!yellow}Any Education}
	\input{text/educ_wage_nodes.txt}
	\end{axis}
\end{tikzpicture}

\footnotesize{Source: Author's Calculations from Current Population Survey} \hfill \tbllink{educ_wage.csv}
\vspace{2mm}

\small \input{text/educ_wage2.txt}
\vspace{1mm}

\begin{minipage}{0.53\textwidth}
\normalsize \textbf{Real Wage Growth Since 2019}\\
\footnotesize{\textit{median usual weekly earnings, percent}}
\vspace{-7mm}

\begin{tikzpicture}
  	\begin{axis}[width=7.7cm, bar width=22pt, height=5.4cm, ytick=false,
  		ymajorticks=false,
    	enlarge x limits={abs=0.8cm}, enlarge y limits={lower, 0.16}, 
    	symbolic x coords={High school degree or less, 
    		Some college or associate degree, Bachelor's degree or more, 
    		Any education level}, clip=false, 
    	axis line style={draw=none}, xtick style={draw=none},
    	xtick=data, yticklabel style={text width=0.7em}, 
		xticklabels from table={\web}{group}, 
		xticklabel style={black!70, font=\footnotesize, text width=1.5cm, align=center}, 
		nodes near coords style={/pgf/number format/.cd, fixed zerofill,
			precision=1, assume math mode}]
  	\addplot[ybar, bar shift=0pt, fill=none, draw=none, nodes near coords] 
  		table [y index=1] {\web};
    \addplot[ybar stacked, bar shift=0pt, fill=green!70!blue, draw=none] 
  		table [y index=3] {\web};
   	\addplot[ybar stacked, bar shift=0pt, fill=green!50!blue, draw=none] 
  		table [y index=2] {\web};
  	\node[rotate=90, anchor=west, shift={(0cm,0.54cm)}] at (axis cs:Any education level,2.21){\scriptsize \color{green!50!blue}Composition};
  	\end{axis}
\end{tikzpicture}
\vspace{-2mm}

\footnotesize{\color{black!60}Total percent change in median usual weekly earnings from 2019 average to \input{text/rw_cpsdate.txt}\unskip, by highest level of education. Full-time wage and salary workers age 25--54. Adjusted for inflation by regional CPI-U.\\{\color{green!50!blue}Composition} is the real wage growth explained by the increase in the education of the population.}

\footnotesize{Source: Author's Calculations from CPS} \hfill \tbllink{educ_wage_bar.csv}
\end{minipage}\hfill \begin{minipage}{0.4\textwidth}
\small \input{text/educ_wage3.txt}
\end{minipage}
\end{minipage}
\newpage
\begin{minipage}{0.76\textwidth}
\subsubsection*{Employment Cost Index}
\small The Bureau of Labor Statistics \href{https://www.bls.gov/news.release/eci.nr0.htm}{report} the overall hourly labor costs faced by employers, using an index that is not influenced by short-term changes in the industry and occupation composition of the US workforce. Employer costs are reported separately for total compensation, wages and salaries, and benefits.

Benefits include health insurance, retirement, vacation, sick leave, and transportation benefits. Benefits access and participation vary, even within the same firm. The benefits costs in the index are averages \href{https://www.bls.gov/opub/mlr/cwc/benefit-cost-concepts-and-the-limitations-of-ecec-measurement.pdf}{computed} across all workers, including the workers who do not have benefits. 
\end{minipage}

\begin{minipage}{0.41\textwidth}
\normalsize \textbf{Employment Cost Growth}\\
\footnotesize{\textit{private industry wage and salary workers,}}\\
\footnotesize{\textit{12-month percent change}}\\
\hspace*{-2mm} \begin{tikzpicture}
	\begin{axis}[\bbar{y}{0}, \dateaxisticks ytick={0, 2, 4, 6, 8}, 
		enlarge y limits={0.01}, ymin=0, ymax=8,
		xticklabel={`\short{\year}}, width=6.9cm, height=5.8cm, clip=false]
	\rebars
	\stdline{cyan!40!blue}{date}{Be}{data/eci.csv}
	\stdline{green!80!blue}{date}{WS}{data/eci.csv}
	\input{text/eci_nodes.txt}
	\absnode{{2016-01-01}}{5.6}{\scriptsize \color{green!80!blue}\textbf{Wages and}\\\scriptsize \color{green!80!blue}\textbf{Salaries}}
	\absnode{{2004-08-01}}{7.6}{\scriptsize \color{cyan!40!blue}\textbf{Benefits}}
	\end{axis}
\end{tikzpicture}\\
\footnotesize{Source: Bureau of Labor Statistics} \hfill \tbllink{eci.csv}
\end{minipage} \hspace{6mm} \begin{minipage}{0.3\textwidth}
\small \input{text/eci.txt}
\end{minipage}
\newpage
\begin{minipage}{0.76\textwidth}
\subsubsection*{Wage Growth Tracker}
\small The Federal Reserve Bank of Atlanta \href{https://www.frbatlanta.org/chcs/wage-growth-tracker}{publish} a wage growth tracker that captures the distribution of changes in the wages of the same people over one year. This approach avoids some of the compositional changes that affect aggregate wage growth measures, though the sample used to calculate the data is affected by changes to respondents' employment status, and by survey response rates. 

\input{text/atl_wgt.txt}
\end{minipage}
\vspace{0.5mm}

\begin{minipage}{0.3\textwidth}
\normalsize \textbf{Wage Growth Tracker}\\
\footnotesize{\textit{median one-year wage growth, percent}}
\end{minipage} \hspace{7mm}
\begin{minipage}{0.44\textwidth}
\footnotesize \colorline{blue!75!cyan}{ \ 3-month moving average} \\ \colorline{orange!60!white}{ \ Unsmoothed}
\end{minipage}
\vspace{-1mm}

\begin{minipage}{0.76\textwidth}
\normalsize \ \ \ \ Hourly Wage \hspace{4.1cm} \normalsize Weekly Wage
\vspace{-2mm}

\hspace*{-3mm} \begin{tikzpicture}
	\begin{axis}[\bbar{y}{0}, \shdateaxisticks ytick={0, 1, 2, 3, 4, 5, 6, 7}, 
		enlarge y limits={0.01}, ymin=0, ymax=7.0, yticklabel style={text width=0.7em}, 
		xticklabel={`\short{\year}}, height=5.4cm, width=6.7cm, clip=false]
	\rebars
	\thinline{orange!60!white}{date}{bd_cps}{data/atl_wgt.csv}
	\stdline{blue!75!cyan}{date}{3ma}{data/atl_wgt.csv}
	\end{axis}
\end{tikzpicture} \hfill
\begin{tikzpicture}
	\begin{axis}[\bbar{y}{0}, \shdateaxisticks ytick={0, 1, 2, 3, 4, 5, 6, 7}, 
		enlarge y limits={0.01}, ymin=0, ymax=7.0, yticklabel style={text width=0.7em}, 
		xticklabel={`\short{\year}}, height=5.4cm, width=6.7cm, clip=false]
	\rebars
	\thinline{orange!60!white}{date}{wk}{data/atl_wgt.csv}
	\stdline{blue!75!cyan}{date}{wk3ma}{data/atl_wgt.csv}
	\end{axis}
\end{tikzpicture}\\
\footnotesize{Source: Author's Calculations} \hfill \tbllink{atl_wgt.csv}
\vspace{4mm}

\small By observing the same person's wage at two points, one year apart, we see how many people do not receive a wage increase. The Atlanta Fed measures this as the share of individuals who have one-year hourly wage growth of between -0.5 and 0.5 percent. The Atlanta Fed approach is replicated using the bd CPS, and smoothed with a 3-month moving average. 

\input{text/atl_zwc.txt}
\vspace{1mm}

\normalsize \textbf{Zero Wage Change}\\
\footnotesize{\textit{share of individuals with one-year wage growth between -0.5 and 0.5 percent}}\\
\hspace*{-2mm} \begin{tikzpicture}
	\begin{axis}[\bbar{y}{0}, \dateaxisticks ytick={8, 10, 12, 14, 16, 18}, 
		yticklabel style={text width=1.0em}, 
		enlarge y limits={0.1}, 
		xticklabel={`\short{\year}}, height=4.8cm, width=12.2cm, clip=false]
	\rebars
	\stdline{red}{date}{zwc}{data/atl_wgt.csv}
	\input{text/zwc_node.txt}
	\end{axis}
\end{tikzpicture}\\
\footnotesize{Source: Author's Calculations} \hfill \tbllink{atl_wgt.csv}
\end{minipage}
\newpage
\begin{minipage}{0.76\textwidth}
\subsection*{Labor Productivity}
\hypertarget{labprod}{}
\small Labor productivity is \href{https://www.bls.gov/news.release/prod2.nr0.htm}{reported} by the Bureau of Labor Statistics and measured as \textbf{real output per hour of work} in the nonfarm business sector. Economic theory suggests that labor productivity is particularly important for long-term real economic growth. The measure captures the rate at which people, with all of the resources and equipment and infrastructure available to them, are able to produce goods and services with their work. An increase in labor productivity means real wages can increase without putting upward pressure on inflation. Alternatively, an increase in productivity means a society can meet its material needs with less work.

\input{text/lprod.txt}
\vspace{3mm}

\normalsize \textbf{Labor Productivity Growth}\\
\footnotesize{\textit{quarterly growth at seasonally adjusted annual rate, percent}}\\
\hspace*{-2mm} \begin{tikzpicture}
	\begin{axis}[\bbar{y}{0}, \dateaxisticks ytick={-5, 0, 5, 10}, 
		xticklabel={`\short{\year}}, clip=false, height=5.2cm]
	\rbars
	\sbar{teal}{date}{value}{data/lprod.csv}
	\end{axis}
\end{tikzpicture}\\
\footnotesize{Source: Bureau of Labor Statistics} \hfill \tbllink{lprod.csv}
\vspace{3mm}

\small In the short-term, productivity growth is heavily affected by changes in the composition of the workforce, and by and volatility in both the number of hours worked and in production. In the longer-term, the level of business net investment in equipment and other capital goods, particularly relative to the size of the workforce, affects productivity growth. Such investment allows more goods and services to be produced by the same hours of work. Yet efforts to stimulate business investment directly through reducing corporate income taxes do not seem to have worked. 

One theory of what drives long-term trends, sometimes called the \textit{Kaldor-Verdoorn Law}, states that demand, and the capacity to meet that demand, determine productivity growth. An economy facing real resource constraints, where demand for goods and services exceeds the capacity to provide these services, is more likely to find ways to produce goods and services more efficiently. As one example, businesses invest more in labor-saving technologies when faced with a tight labor market. 
\end{minipage}
\newpage
\begin{minipage}{0.76\textwidth}
\subsection*{Union Membership}
\hypertarget{labun}{}
\small Membership in \textbf{unions and employee associations} has diminished in the United States over the past fifty years. Unionized jobs typically offer higher wages and better benefits and union membership tends to increase wages and benefits even in nonunion jobs. Many researchers argue that lower union membership increases income inequality. 

\input{text/union.txt}
\vspace{3mm}

\normalsize \textbf{Union Membership and Coverage}\\
\footnotesize{\textit{union or employee association share of jobs, percent}}\\
\hspace*{-2mm} \begin{tikzpicture}
	\begin{axis}[\bbar{y}{0}, \dateaxisticks ytick={0, 5, 10, 15, 20}, ymin=1,
		height=6.2cm, width=12.6cm, enlarge y limits={0.11},
		yticklabel style={text width=1.0em}, 
		xticklabel={`\short{\year}}, clip=false, legend style={at={(0.95, 1.0)}}]
	\rbars
	\ctsbar{violet}{date}{Membership}{data/union.csv}{1.9pt}
	\ctsbar{magenta!70!purple}{date}{Diff}{data/union.csv}{1.9pt}	
	\legend{Member, Covered Non-member};
	\end{axis}
\end{tikzpicture}\\
\footnotesize{Source: Author's Calculations from Current Population Survey} \hfill \tbllink{union.csv}
\end{minipage}
\vspace{2mm}

\begin{minipage}{0.43\textwidth}
\normalsize \textbf{Union Membership Rate by Industry}\\
\footnotesize{\textit{union or employee association member, percent}}\\
\hspace*{-2mm} \begin{tikzpicture}
	\begin{axis}[clip=false, stack negative=separate, ymajorgrids=false,
            yticklabels from table={\unmem}{name},
			y axis line style={opacity=0},   
		    x axis line style={opacity=0}, 
			yticklabel style={black, font=\footnotesize, 
			align=left, anchor=east},  
			nodes near coords align={above},    
			nodes near coords style={/pgf/number format/.cd,fixed zerofill, 
				precision=1, assume math mode},  
            ytick=data, tickwidth=0pt, xmin=0,
            bar width=2.0ex, 
            extra x ticks={0}, extra x tick style={
                grid style={black}, xticklabel=\empty,},
            enlarge y limits={0.08}, enlarge x limits={0.02},
            width=5.2cm, height=7.2cm, legend style={at={(0.7, 1.07)}},
            \barylab{2.7cm}{1.5ex}]
	\addplot[xbar stacked, bar shift=0pt, fill=none, draw=none] 
		table [y expr=-\coordindex, x index=1] {\unmem};
	\addplot[xbar stacked, bar shift=0pt, fill=violet!60!blue, draw=none] 
		table [y expr=-\coordindex, x index=2] {\unmem};		        
	\addplot[only marks, mark=diamond*, mark options={fill=orange, 
		draw=orange, scale=2}] 
		table [y expr=-\coordindex, x index=3] {\unmem};
	\addplot [nodes near coords={\pgfmathprintnumber{\pgfkeysvalueof{/data point/x}}}, only marks, mark=none] table[y=y, x index=3] {\unmem};
	\legend{, 30-year range, \input{text/cps_ltdate.txt}}
	\end{axis}
\end{tikzpicture}\\
\footnotesize{Source: Author's Calculations from CPS} \hspace{12mm} \tbllink{union_ind.csv}
\end{minipage}\hspace{4mm}
\begin{minipage}{0.3\textwidth}
\small Union membership rates vary substantially by industry. \input{text/union_ind.txt}
\end{minipage}
\newpage
\hypertarget{cap}{}
\begin{minipage}{0.76\textwidth}
\section*{Financial Markets}
\vspace{-1mm}
\small The US equity and capital markets provide funding for borrowers' activities and provide a source of income and capital gains to lenders. This section discusses equity markets, interest rates and bond markets, and money and monetary policy. 
\hypertarget{capeq}{}
\subsection*{Equity Markets}
\small Equity markets, or \textbf{stock markets}, provide a method for businesses to raise capital by selling shares, which represent ownership claims on the business. Equity markets also provide a place for people to buy and sell existing shares. Investors purchase shares in hopes that the price will go up, allowing them to sell the shares at a higher price and receive capital gains, or to gain access to a stream of dividends, which are payments from businesses to shareholders. 

In the US, there are several stock market indices that track the share price of a basket of companies. These measures are weighted by the market capitalization of the companies in the basket, which is the share price multiplied by the number of shares. Market capitalization measures the market value of the company. Note that stock market indices do not provide information about dividends, and the dividend yield varies substantially between companies and between the baskets or groups of companies used in stock market indices.
\vspace{2mm}

\normalsize \textbf{Stock Market Indices}\\
\footnotesize{\textit{adjusted close, index value, in thousands \hspace{30mm} 2018--onward, in detail}}\\
\hspace*{-3mm} \begin{tikzpicture}
	\begin{groupplot}[group style={group name=equities, group size=2 by 4, 
    x descriptions at=edge bottom, horizontal sep=40pt, vertical sep=8pt},]
	\nextgroupplot[\shdateaxisticks ytick={0, 1, 2, 3, 4, 5}, 
		enlarge y limits={0.1}, yticklabel style={text width=1.2em}, 
		xticklabel={`\short{\year}}, height=4.0cm, width=7.8cm, clip=false, xmin={1989-01-01},
		xticklabels=\empty, xtick style={draw=none},]
	\rbars
	\coordinate (c1) at ({2017-12-31},2);
  	\coordinate (c2) at (rel axis cs:1,1);
	\draw [black!25!white, dashed] (c1) rectangle (c2);
	\stdline{green!80!blue!90!black}{date}{GSPC}{data/equity_indices.csv}
	\node[below, align=left, shift=({rel axis cs:0.13,0.98}), fill=white, inner sep=0.9]
		{\normalsize \color{green!80!blue!90!black}\textbf{S\&P 500}};
	\nextgroupplot[\ltdateaxisticks ytick={2, 3, 4, 5}, 
		enlarge y limits={0.1}, %name=ax2, 
		yticklabel style={text width=0.6em}, yticklabel style={text width=1.2em},
		xticklabel={`\short{\year}}, height=4.0cm, width=5.2cm, clip=false, 
		xticklabels=\empty, xtick style={draw=none},]
	\rbar
	\thickline{black!25}{date}{GSPC_MA}{data/equity_indices_lt.csv}
	\thinline{green!80!blue!90!black}{date}{GSPC_close}{data/equity_indices_lt.csv}
	\absnodeml{{2021-01-01}}{3.35}{\scriptsize \color{black!35!white}\textbf{moving average}}{1.3}
	\nextgroupplot[\shdateaxisticks 
		enlarge y limits={0.1}, yticklabel style={text width=1.2em}, 
		xticklabel={`\short{\year}}, height=4.0cm, width=7.8cm, clip=false, xmin={1989-01-01},
		xticklabels=\empty, xtick style={draw=none},]
	\rbars
	\coordinate (d1) at ({2017-12-31},6);
  	\coordinate (d2) at (rel axis cs:1,1);
	\draw [black!25!white, dashed] (d1) rectangle (d2);
	\stdline{blue}{date}{IXIC}{data/equity_indices.csv}
	\node[below, align=left, shift=({rel axis cs:0.12,0.98}), fill=white, inner sep=0.9]
		{\normalsize \color{blue}\textbf{Nasdaq}};
	\nextgroupplot[\ltdateaxisticks 
		enlarge y limits={0.1}, %name=ax4, 
		yticklabel style={text width=0.6em}, yticklabel style={text width=1.2em}, 
		xticklabel={`\short{\year}}, height=4.0cm, width=5.2cm, clip=false, 
		xticklabels=\empty, xtick style={draw=none},]
	\rbar
	\thickline{black!25}{date}{IXIC_MA}{data/equity_indices_lt.csv}
	\thinline{blue}{date}{IXIC_close}{data/equity_indices_lt.csv}
	\absnodeml{{2021-01-01}}{10.8}{\scriptsize \color{black!35!white}\textbf{moving average}}{1.3}
	\nextgroupplot[\shdateaxisticks 
		enlarge y limits={0.1}, yticklabel style={text width=1.2em}, xmin={1989-01-01},
		xticklabel={`\short{\year}}, height=4.0cm, width=7.8cm, clip=false, 
		xticklabels=\empty, xtick style={draw=none},]
	\rbars
	\coordinate (e1) at ({2017-12-31},18);
  	\coordinate (e2) at (rel axis cs:1,1);
	\draw [black!25!white, dashed] (e1) rectangle (e2);
	\stdline{red}{date}{DJI}{data/equity_indices.csv}
	\node[below, align=left, shift=({rel axis cs:0.12,0.98}), fill=white, inner sep=0.9]
		{\normalsize \color{red}\textbf{Dow 30}};
	\nextgroupplot[\ltdateaxisticks 
		enlarge y limits={0.1}, %name=ax4, 
		yticklabel style={text width=0.6em}, yticklabel style={text width=1.2em}, 
		xticklabel={`\short{\year}}, height=4.0cm, width=5.2cm, clip=false, 
		xticklabels=\empty, xtick style={draw=none},]
	\rbar
	\thickline{black!25}{date}{DJI_MA}{data/equity_indices_lt.csv}
	\thinline{red}{date}{DJI_close}{data/equity_indices_lt.csv}
	\absnodeml{{2021-01-01}}{27}{\scriptsize \color{black!35!white}\textbf{moving average}}{1.3}	
	\nextgroupplot[\shdateaxisticks 
		enlarge y limits={0.1}, yticklabel style={text width=1.2em}, xmin={1989-01-01},
		xticklabel={`\short{\year}}, height=4.0cm, width=7.8cm, clip=false]
	\rbars
	\coordinate (f1) at ({2017-12-31},1);
  	\coordinate (f2) at (rel axis cs:1,1);
	\draw [black!25!white, dashed] (f1) rectangle (f2);
	\stdline{violet}{date}{RUT}{data/equity_indices.csv}
	\node[below, align=left, shift=({rel axis cs:0.18,0.98}), fill=white, inner sep=0.9]
		{\normalsize \color{violet}\textbf{Russell 3000}};
	\nextgroupplot[\ltdateaxisticks 
		enlarge y limits={0.1}, %name=ax4, 
		yticklabel style={text width=0.6em}, yticklabel style={text width=1.2em}, 
		xticklabel={`\short{\year}}, height=4.0cm, width=5.2cm, clip=false]
	\rbar
	\thickline{black!25}{date}{RUT_MA}{data/equity_indices_lt.csv}
	\thinline{violet}{date}{RUT_close}{data/equity_indices_lt.csv}
	\absnodeml{{2021-01-01}}{1.6}{\scriptsize \color{black!35!white}\textbf{moving average}}{1.3}	
	\end{groupplot}
	\draw [black!25!white, dashed] (c1) -- (equities c2r1.south west);
	\draw [black!25!white, dashed] (c2) -- (equities c2r1.north west);
	\draw [black!25!white, dashed] (d1) -- (equities c2r2.south west);
	\draw [black!25!white, dashed] (d2) -- (equities c2r2.north west);
	\draw [black!25!white, dashed] (e1) -- (equities c2r3.south west);
	\draw [black!25!white, dashed] (e2) -- (equities c2r3.north west);
	\draw [black!25!white, dashed] (f1) -- (equities c2r4.south west);
	\draw [black!25!white, dashed] (f2) -- (equities c2r4.north west);
\end{tikzpicture}\\
\footnotesize{Source: Yahoo! Finance} \hfill \tbllink{equity_indices.csv} 
\end{minipage}
\newpage
\begin{minipage}{0.76\textwidth}
\small The \href{https://us.spindices.com/indices/equity/sp-500}{S\&P 500} (see {\color{green!80!blue!90!black}\textbf{---}}) is a market-cap-weighted stock market index based on 500 large companies listed on US exchanges. \input{text/eq_sp500.txt}

The \href{https://www.nasdaq.com/market-activity/index/comp}{Nasdaq} composite index (see {\color{blue}\textbf{---}}) includes nearly all companies listed on the Nasdaq stock exchange, and is heavily-weighted towards large tech companies. \input{text/eq_nasdaq.txt}

The \href{https://finance.yahoo.com/quote/\%5EDJI/}{Dow Jones} (see {\color{red}\textbf{---}})  industrial average is an index based on 30 large and prominent companies listed on US exchanges. The measure is used as a proxy for the performance of the largest companies, and \input{text/eq_dow.txt}

Lastly, the \href{https://www.ftserussell.com/products/indices/russell-us}{Russell 3000} (see {\color{violet}\textbf{---}}) is a broad measure of the US stock market that seeks to be a benchmark of the performance of the overall market. \input{text/eq_russell.txt}
\vspace{1mm}

\normalsize \textbf{Stock Market Indices}\\
\footnotesize{\ \textit{\hspace{18mm} adjusted close \hspace{16mm} annual index returns, not including dividends}
\vspace{-2mm}

\rowcolors{1}{}{black!5} \setlength{\tabcolsep}{2.0pt} \color{black!90}
{\renewcommand{\arraystretch}{1.54}
	\begin{tabular}{p{20mm} R{12.0mm} R{12.0mm} R{10.0mm} R{10.0mm} R{10.0mm} R{10.0mm} R{10.0mm} R{10.0mm}}
		\input data/equities.tex \hline
	\end{tabular}
}}
		
\vspace{-2mm}
\footnotesize{Source: Yahoo! Finance}
\vspace{3mm}

\subsubsection*{Real Return}
\small \input{text/sp500rr.txt}  
\vspace{2mm}

\normalsize \textbf{S\&P 500 Real Return}\\
\footnotesize{\textit{trailing period annualized real return}}
\vspace{3.5cm}

\hspace*{4mm} \begin{tikzpicture}[overlay]
	\begin{axis}[\bbar{y}{0}, \dateaxisticks ytick={-5, 0, 5, 10, 15}, 
		yticklabel style={text width=1.0em}, 
		xticklabel={`\short{\year}}, clip=false, height=5.4cm, width=12.4cm,
		enlarge y limits={0.11}, legend style={at={(0.98, 1.12)}, fill=white}]
	\rbars
	\thickline{blue}{date}{r10}{data/sp500rr.csv}
	\thickline{orange}{date}{r15}{data/sp500rr.csv}
	\thickline{green!80!blue}{date}{r20}{data/sp500rr.csv}
	\input{text/sp500rr_nodes.txt}
	\legend{10-year, 15-year, 20-year};
	\end{axis}
\end{tikzpicture}

\footnotesize{Source: Shiller, Baker, Author's Calculations} \hfill \tbllink{sp500rr.csv}
\end{minipage}
\newpage
\begin{minipage}{0.76\textwidth}
\small The total return on corporate equities includes both dividends and stock buybacks. US nonfinancial corporations have used stock buybacks to return money to investors, often at the expense of domestic fixed investment. To the extent that markets respond rationally, however, buybacks are no different from dividends. 
\vspace{1mm}

\normalsize \textbf{Corporate Equity Payout}\\
\footnotesize{\textit{nonfinancial corporation net dividends paid and net issuance of equities, percent of GDP}}\\
\hspace*{-2mm} \begin{tikzpicture}
	\begin{axis}[\bbar{y}{0}, \dateaxisticks ytick={0, 5, 10}, 
		xticklabel={`\short{\year}}, clip=false, legend style={at={(0.44, 0.97)}, fill=white}]
	\rbars
	\sbar{blue!90!purple}{date}{Dividends}{data/eq_payout.csv}	
	\sbar{magenta}{date}{Buybacks}{data/eq_payout.csv}	
	\legend{Net Dividends, Net Issuance};
	\end{axis}
\end{tikzpicture}\\
\footnotesize{Source: Federal Reserve, Bureau of Economic Analysis} \hfill \tbllink{eq_payout.csv}

\subsubsection*{Valuation}
\small The \href{http://www.econ.yale.edu/~shiller/data.htm}{cyclically-adjusted price to earnings ratio} (CAPE) compares the current price of the S\&P 500 to the previous ten-years of total S\&P 500 returns, including dividends and buybacks (treated as dividends). Valuations often use recent or forecasted earnings. Robert Shiller's CAPE covers a normal business cycle so that valuations are less-affected by the idiosyncrasies of current economic conditions. 
\vspace{1mm}

\normalsize \textbf{Price to Earnings Ratio}\\
\footnotesize{\textit{10-year total return PE ratio}}\\
\hspace*{-2mm} \begin{tikzpicture}
	\begin{axis}[\bbar{y}{0}, \dateaxisticks ytick={10, 20, 30, 40, 50, 60}, 
		xticklabel={`\short{\year}}, clip=false, yticklabel style={text width=1.4em}, 
		height=4.4cm, width=12.3cm]
	\rbars
	\stdline{blue!80!black}{date}{TR CAPE}{data/catrpe.csv}
	\input{text/cape_node.txt}
	\end{axis}
\end{tikzpicture}\\
\footnotesize{Source: Shiller} \hfill \tbllink{catrpe.csv}

\subsubsection*{Volatility}
\small The Chicago Board Options Exchange uses S\&P 500 options data to \href{http://www.cboe.com/vix}{identify} expectations of future volatility. \input{text/vixbasic.txt}
\vspace{1mm}

\normalsize \textbf{S\&P 500 Volatility Index}\\
\footnotesize{\textit{index, monthly average shown}}\\
\hspace*{-2mm} \begin{tikzpicture}
	\begin{axis}[\bbar{y}{0}, \dateaxisticks ytick={0, 20, 40, 60}, 
		xticklabel={`\short{\year}}, clip=false, ymin=4,
		yticklabel style={text width=1.0em}, width=12.3cm,
		height=4.4cm, enlarge y limits={0.08}]
	\rbars
	\stdline{magenta}{date}{value}{data/vix.csv}
	\input{text/vix_node.txt}
	\end{axis}
\end{tikzpicture}\\
\footnotesize{Source: Chicago Board Options Exchange} \hfill \tbllink{vix.csv}
\end{minipage}
\newpage
\begin{minipage}{0.76\textwidth}
\subsection*{Interest Rates}
\hypertarget{capint}{}
\small The US Federal Reserve System (Fed) has a congressional \href{https://www.federalreserve.gov/faqs/money_12848.htm}{mandate} to promote price stability and maximum employment. In practice, a Fed committee (the FOMC) determines the federal funds rate, which aims to influence interest rates in the broader economy. There are several channels through which interest rates affect broader economic activity, for example, lower interest rates stimulate investment in capital goods and the production of these capital goods employs people, who in turn spend their wages on other goods and services. Through its influence on interest rates, the Fed's monetary policy can aim to be neutral or to stimulate or slow the economy.

The FOMC cut \href{https://www.federalreserve.gov/releases/h15/}{interest rates} three times in 2019, for a total reduction of 75 basis points. Responding to the economic shock of the coronavirus, the FOMC cut rates twice in March 2020, by 150 basis points, bringing the lower bound of the federal funds rate range to zero. With at near zero, the Fed adopted several additional measures to increase liquidity in the global financial system.

In March 2022, the FOMC raised the base interest rate by 25 basis points. \input{text/rates_ff.txt}
\vspace{2mm}

\normalsize \textbf{Effective Federal Funds Rate}\\
\footnotesize{\textit{percent, monthly average except for latest value }}\\
\hspace*{-2mm} \begin{tikzpicture}
	\begin{axis}[\bbar{y}{0}, \dateaxisticks ytick={-10, -8, -6, -4, -2, 0, 2, 4, 6, 8, 10}, 
		yticklabel style={text width=1.2em}, height=7.6cm, width=12.6cm, xmax=2025-01-02,
		xticklabel={`\short{\year}}, clip=false]
	\rbars
	\input{text/ff_proj_bar.txt}
	\thickline{cyan!80!white}{date}{Value}{data/taylor.csv}
	\stdline{blue!60!black}{date}{Fed Funds}{data/rates.csv}
	\addplot[only marks, mark=square*, mark options={fill=purple, 
		draw=white, scale=1.4}] 
		table [x=date, y=VALUE, col sep=comma] {data/sep.csv};
	\node[text width=0.8cm, anchor=west, fill=white] at 
		(axis description cs: 0.925, 0.9){\scriptsize \color{purple}\textbf{SEP}};
	\input{text/rates_ff_proj_node.txt}
	\input{text/rates_ff_proj_date.txt}
	\absnode{{1992-01-01}}{7.5}{\footnotesize \color{blue!60!black}\textbf{Fed Funds Rate}}
	\absnode{{2015-06-01}}{-1.8}{\scriptsize \color{cyan!80!white}\textbf{Modified}\\\scriptsize \color{cyan!80!white}\textbf{Taylor Rule}}
	\end{axis}
\end{tikzpicture}\\
\footnotesize{Source: Federal Reserve, Author's Calculations} \hfill \tbllink{rates.csv}
\vspace{2mm}

\small Economist John Taylor described a rule for setting the federal funds rate based on inflation and output. Versions of this \textit{Taylor rule} track the actual federal funds rate fairly closely during the 1990s and 2000s. Former Fed Chair Ben Bernanke \href{https://www.brookings.edu/blog/ben-bernanke/2015/04/28/the-taylor-rule-a-benchmark-for-monetary-policy/}{described} a modified Taylor rule based on core PCE inflation and a stronger response to the output gap (see {\color{cyan!80!white}\textbf{---}}). \input{text/rates_ff_proj.txt}
\end{minipage}
\newpage
\begin{minipage}{0.76\textwidth}
\subsubsection*{Real Interest Rates}
\small Lenders charger higher interest rates to compensate for expected inflation. Real interest rates, which are adjusted for expected inflation, can offer insight into economic and financial conditions. Low real interest rates encourage borrowing and consumption and increased economic activity while high real interest rates discourage borrowing and encourage saving. 

\input{text/real_rates_basic.txt} 
\vspace{2mm} 

\normalsize \textbf{Real Treasury Yields}\\
\footnotesize{\textit{derived from inflation-indexed treasuries, percent}}\\
\hspace*{-2mm} \begin{tikzpicture}
	\begin{axis}[\bbar{y}{0}, \dateaxisticks ytick={-2, -1, 0, 1, 2, 3, 4}, 
		xticklabel={\year}, height=5.2cm, width=12.1cm, clip=false]
	\recbars
	\thickline{green!80!blue}{date}{Ten-year}{data/real_rates.csv}
	\thickline{violet!50!blue!60!white}{date}{Five-year}{data/real_rates.csv}
	\absnode{{2004-05-01}}{0.8}{\small \color{violet!50!blue!60!white}\textbf{Five-year}}
	\absnode{{2009-10-01}}{2.2}{\small \color{green!80!blue}\textbf{Ten-year}}
	\input{text/real_yield_nodes.txt}
	\end{axis}
\end{tikzpicture}\\
\footnotesize{Source: Federal Reserve} \hfill \tbllink{real_rates.csv}
\vspace{3mm}

\small The Cleveland Fed calculates real yields using a model based on treasury yields, inflation, and financial-market- and survey-based information. The previous approach has limitations, as the market for treasury inflation-indexed securities is relatively small and can be influenced by monetary policy. The Cleveland Fed model estimates real yields across the term structure and avoids some limitations of the previous approach.

\input{text/real_yield_model.txt}
\vspace{2mm}

\normalsize \textbf{Real Treasury Yields}\\
\footnotesize{\textit{model-based, percent}}
\vspace{3.2cm}

\hspace{4mm} \begin{tikzpicture}[overlay]
	\begin{axis}[\bbar{y}{0}, \dateaxisticks  ytick={-6, -4, -2, 0, 2, 4, 6, 8},
		yticklabel style={text width=1.4em}, enlarge y limits={0.08}, ymax=4.4,
		xticklabel={`\short{\year}}, clip=false, height=5.2cm, width=12.0cm]
	\rbars
	\stdline{green!68!black}{date}{Real Rate 1-year}{data/real_rates2.csv}
	\stdline{blue}{date}{Real Rate 10-year}{data/real_rates2.csv}
	\absnode{{1995-01-01}}{5.0}{\small \color{blue}\textbf{Ten-year}}
	\absnode{{2012-06-01}}{-2.5}{\small \color{green!68!black}\textbf{One-year}}
	\input{text/real_yield_nodes2.txt}
	\end{axis}
\end{tikzpicture}
\vspace{1mm}

\footnotesize{Source: Federal Reserve of Bank of Cleveland}  \hfill \tbllink{real_rates2.csv}
\end{minipage}
\newpage
\begin{minipage}{0.76\textwidth} 
\small Changes in nominal treasury yields can be \textbf{decomposed into changes in expected inflation and changes in real yields}. Changes in real yields reflect changes in expected path of the federal funds rate and the economic outlook. Federal Reserve Bank of Cleveland models \href{https://www.clevelandfed.org/en/our-research/indicators-and-data/inflation-expectations.aspx}{identify} inflation expectations across the term structure, which can be used to identify changes in real yields. 

\input{text/real_yield_model_ch.txt}
\vspace{2mm}

\normalsize \textbf{Recent Changes}\\
\footnotesize{\textit{percentage points}}
\vspace{4.1cm}

\hspace{3mm} \begin{tikzpicture}[overlay]
	\begin{groupplot}[group style={group size=2 by 1, horizontal sep=62pt}]
	\nextgroupplot[\bbar{y}{0}, width=5.6cm, height=5.6cm, 
		axis line style={draw=none}, xtick=data, clip=false, ytick={0}, 
		xticklabel style={black!90, font=\small, text width=1.5cm, align=center}, 
	 	enlarge x limits={0.1}, transpose legend,
	 	reverse legend,
	 	xticklabels from table={\rych}{name}, enlarge y limits={lower, 0.3},
	 	nodes near coords style={style={black!75, anchor=south,
		 yshift=1.5mm, inner sep=0.5}, /pgf/number format/.cd, fixed zerofill,
			precision=2, assume math mode}]
  	\addplot[scatter, mark=*, only marks, mark options={fill=blue, 
		draw=blue, scale=1.5}, nodes near coords] 
  		table [x expr=\coordindex, y=Total] {\rych};
	\addplot[ybar stacked, bar width=20pt, draw opacity=0, fill=black!12!white] table [x expr=\coordindex, y=IE]{\rych};
	\addplot[ybar stacked, bar width=20pt, draw opacity=0, fill=lime!50!green!80!black] table [x expr=\coordindex, y=RY]{\rych};
	\stdnode{1.6cm}{3.6cm}{\normalsize \color{black!90} \textbf{Three-month change}}
	\input{text/inf_exp_ch_nodes.txt}
	\nextgroupplot[\bbar{y}{0}, width=5.6cm, height=5.6cm, 
		axis line style={draw=none}, xtick=data, clip=false, ytick={0}, 
		xticklabel style={black!90, font=\small, text width=1.5cm, align=center}, 
	 	enlarge x limits={0.1}, transpose legend, reverse legend,
	 	xticklabels from table={\rytch}{name}, enlarge y limits={lower, 0.12},
	 	nodes near coords style={style={black!75, anchor=south,
		 yshift=1.5mm, inner sep=0.5}, /pgf/number format/.cd, fixed zerofill,
			precision=2, assume math mode}, legend cell align={left}, 
		legend style={legend columns=2, at={(0.48, 1.34)}}]
  	\addplot[scatter, mark=*, only marks, mark options={fill=blue, 
		draw=blue, scale=1.5}, nodes near coords] 
  		table [x expr=\coordindex, y=Total] {\rytch};
	\addplot[ybar stacked, bar width=20pt, draw opacity=0, fill=black!12!white] table [x expr=\coordindex, y=IE]{\rytch};
	\addplot[ybar stacked, bar width=20pt, draw opacity=0, fill=lime!50!green!80!black] table [x expr=\coordindex, y=RY]{\rytch};
	\stdnode{1.6cm}{4.0cm}{\normalsize \color{black!90} \textbf{Three-year change}}
	\input{text/inf_exp_ch_nodes2.txt}
	\legend{\ Total, \ Inflation Expectations,\ Real Yields};
	\end{groupplot}
\end{tikzpicture}
\vspace{3mm}

\footnotesize{Latest data from model as of: {\input{text/frbcle_ry_date.txt}}}\\
\footnotesize{Source: Federal Reserve, Federal Reserve of Bank of Cleveland} \hfill \tbllink{inf_exp_ch.csv}
\end{minipage}
\newpage
\begin{minipage}{0.76\textwidth}
\subsubsection*{Government Bonds}
\small \input{text/rates_basic.txt} 
\vspace{1mm}

\normalsize \textbf{Treasury Constant Maturity Yields}\\
\footnotesize{\textit{percent, monthly average shown except for latest value }}\\
\hspace*{-2mm} \begin{tikzpicture}
	\begin{axis}[\bbar{y}{0}, \dateaxisticks ytick={0, 2, 4, 6, 8, 10}, ymin=0.2,
		yticklabel style={text width=1.0em}, 
		xticklabel={`\short{\year}}, clip=false, height=5.2cm]
	\rbars
	\stdline{blue!70!black}{date}{Ten-year}{data/rates.csv}
	\stdline{cyan!80!white}{date}{Two-year}{data/rates.csv}
	\absnode{{2004-03-01}}{1.8}{\scriptsize \color{cyan!90!white}\textbf{Two-year}}
	\absnode{{2014-05-01}}{3.6}{\scriptsize \color{blue!70!black}\textbf{Ten-year}}
	\end{axis}
\end{tikzpicture}\\
\footnotesize{Source: Federal Reserve} \hfill \tbllink{rates.csv}
\vspace{3mm}

\normalsize \textbf{Selected US Treasury Rates}\\
\footnotesize{\textit{constant maturity yield, percent \hspace{18mm}averages}\\
\hspace*{-3mm} \rowcolors{1}{}{black!5} \setlength{\tabcolsep}{2.0pt} \color{black!90}
		{\renewcommand{\arraystretch}{1.6}
\begin{tabular}{p{17mm} R{10.2mm} R{10.2mm} R{10.2mm} R{9.5mm}
		 		 R{9.5mm} R{9.5mm} R{9.5mm} R{9.5mm} R{9.5mm}}
			 \input data/treasury_rates.tex \hline
		\end{tabular}}}

\vspace{-2mm}		
\footnotesize{Source: Federal Reserve}
\end{minipage}
\newpage
\begin{minipage}{0.76\textwidth}
\small The \textbf{Treasury yield curve} \href{https://www.treasury.gov/resource-center/data-chart-center/interest-rates/Pages/TextView.aspx?data=yield}{shows} the interest rates on different maturities of US Treasury bonds and bills, at a given point in time. The yield curve summarizes the term structure of interest rates, how much it costs to borrow for different periods of time, and has traditionally been considered an indicator of how markets view short-term economic conditions relative to longer-term conditions. 

The yield curve is normally upward sloping as investors expect to be compensated for lending for a longer period of time. The shape of the yield curve changes over time and is affected by several factors, including the term premium, the monetary policy of the Federal Reserve, and expectations about future inflation. The curve can become steeper, for example, if interest rates or inflation is expected to be higher in the future. 
\end{minipage}
\vspace{1mm}

\begin{minipage}{0.39\textwidth}
\normalsize \textbf{Treasury Yield Curve}\\
\footnotesize{\textit{constant maturity yield, percent}}\\
\hspace*{-2mm} \begin{tikzpicture}[trim axis right]
	\begin{axis}[height=7.0cm, width=7.0cm,	xmajorgrids, ymax=3.75,
		axis line style={draw=none}, nodes near coords style={style={black!80, fill=white,
		 yshift=1.5mm, inner sep=0.5}, /pgf/number format/.cd, fixed zerofill, 
		 precision=2, assume math mode}, legend style={at={(0.5, 0.92)}, legend columns=1},
		xtick=data, xticklabels={1M, 3M, 1Y, 2Y, 5Y, 10Y, 20Y, 30Y}, legend cell align={left},
		reverse legend, xticklabel style={black!70, font=\small\bfseries}, 
		enlarge y limits={0.08}, enlarge x limits={0.02}, \bbar{y}{0}]
	\addplot[thick, mark=*, color=blue!65!black!30] 
		table [x=number, y=fiveyear, col sep=comma] {data/yc.csv};
	\addplot[thick, mark=*, color=blue!65!black!60] 
		table [x=number, y=oneyear, col sep=comma] {data/yc.csv};			
	\addplot[ultra thick, mark=*, nodes near coords, 
    		color=blue!60!black] table [x=number, y=value, col sep=comma] {data/yc.csv};
    \input{text/yc_date.txt}
    \legend{Five Years Ago, One Year Ago, Most Recent}	
	\end{axis}
\end{tikzpicture}\\
\footnotesize{Source: Federal Reserve} 
\end{minipage} \hspace{4mm}
\begin{minipage}{0.33\textwidth}
\small The yield curve can also become \textit{inverted} when yields on shorter-term debt are higher than yields on longer-term debt. An inverted yield curve can be a sign of worsening economic conditions. For example, short term rates may exceed longer-term rates if the Federal Reserve is expected to lower interest rates in the future, or if inflation is expected to fall due to weakened economic conditions. 

\input{text/yc_inversion.txt}
\end{minipage}
\vspace{3mm}

\begin{minipage}{0.76\textwidth}
\small Another measure of the term structure of interest rates is the \textit{spread} between treasuries with different maturities. \textbf{Treasury yield spreads} can be used to track changes in the term structure over time.

\input{text/spread_basic.txt}
\vspace{2mm}

\normalsize \textbf{Treasury Yield Spreads}\\
\footnotesize{\textit{percentage points}}\\
\hspace*{-3mm} \begin{tikzpicture}
    \begin{groupplot}[group style={group size=2 by 1, horizontal sep=52pt,}]
    \nextgroupplot[\bbar{y}{0}, \ltdateaxisticks ytick={0, 1, 2, 3}, ymin=-0.65, ymax=2.3,
    	yticklabel style={text width=0.8em},
     xticklabel={\year}, height=5.6cm, width=5.9cm, clip=false]
    \rbar
	\thinline{blue!70!cyan!80!white}{date}{Ten-3M}{data/spread.csv}
	\node[text width=3.8cm, anchor=west] at (axis description cs: 0, 0.95) 
		{\small \color{blue!70!cyan!80!white}\textbf{10-Year - 3-Month}};
	\input{text/spread_node.txt}
    \nextgroupplot[\bbar{y}{0}, \ltdateaxisticks ytick={0, 1, 2, 3}, ymin=-0.65, ymax=2.3,
    	yticklabel style={text width=0.8em},
     xticklabel={\year}, height=5.6cm, width=5.9cm, clip=false]
	\rbar
	\thinline{red!60!violet!90!white}{date}{Ten-2Y}{data/spread.csv}
	\node[text width=3.8cm, anchor=west] at (axis description cs: 0, 0.95) 
		{\small \color{red!60!violet!90!white}\textbf{10-Year - 2-Year}};
	\input{text/spread_node2.txt}
	\end{groupplot}
	\end{tikzpicture}\\
\footnotesize{Source: Federal Reserve} \hfill \tbllink{spread.csv} 
\end{minipage}
\newpage
\begin{minipage}{0.76\textwidth}
\subsubsection*{Corporate Bonds}

\small The US Treasury \href{https://www.treasury.gov/resource-center/economic-policy/corp-bond-yield/Pages/Corp-Yield-Bond-Curve-Papers.aspx}{publish} a yield curve for \textbf{corporate bonds} based on the market-weighted average of bonds rated AAA, AA, and A. \input{text/hqcb.txt}
\vspace{1mm} 

\normalsize \textbf{High Quality Corporate Bonds, 10-Year}\\
\footnotesize{\textit{spot rate, percent, monthly average}}\\
\hspace*{-2mm} \begin{tikzpicture}
	\begin{axis}[\bbar{y}{0}, \dateaxisticks ytick={0, 2, 4, 6, 8, 10}, 
		yticklabel style={text width=1.0em}, height=5.0cm, 
		xticklabel={`\short{\year}}, clip=false, ymin=0.5, width=12.0cm]
	\rbars
	\stdline{blue!60!violet}{date}{VALUE}{data/hqcb.csv}
	\input{text/hqcb_node.txt}
	\end{axis}
\end{tikzpicture}\\
\footnotesize{Source: Treasury} \hfill \tbllink{hqcb.csv}
\vspace{6mm} 

\small The yield of treasury bonds is sometimes considered the risk-free interest rate, as treasury bonds are backed by the full faith and credit of the US and its taxbase. The credit risk is higher for other borrowers, who pay a higher interest rate, as a result. 

Bonds issued by US corporations that are rated below investment grade (a rating below BBB) are \href{https://fred.stlouisfed.org/series/BAMLH0A0HYM2EY}{tracked} by the ICE BofA high yield index. \input{text/highyield.txt}
\vspace{1mm}

\normalsize \textbf{ICE BofA High Yield Index Effective Yield}\\
\footnotesize{\textit{percent}}\\
\hspace*{-2mm} \begin{tikzpicture}
	\begin{axis}[\bbar{y}{0}, \dateaxisticks ytick={0, 5, 10, 15, 20, 25},
		enlarge y limits={0.1}, 
		yticklabel style={text width=1.2em}, width=12.0cm, height=4.6cm,
		xticklabel={`\short{\year}}, ymin=1.8, clip=false]
	\rebars
	\stdline{red!50!purple}{date}{VALUE}{data/highyield.csv}
	\input{text/highyield_node.txt}
	\end{axis}
\end{tikzpicture}\\
\footnotesize{Source: Ice Data Indices} \hfill \tbllink{highyield.csv}
\end{minipage}
\newpage
\hypertarget{capmm}{}
\begin{minipage}{0.76\textwidth}
\subsection*{Money and Monetary Policy}
\small The Federal Reserve \href{https://www.federalreserve.gov/releases/h6/current/default.htm}{publish} data on the \textbf{money supply}. A broad measure of the amount of money, called M2, includes cash and deposits such as savings accounts and checking accounts, as well as time deposits smaller than \$100,000, and retail accounts in money market funds. \input{text/m2lvl.txt}

A large increase in the amount of money held by individuals and institutions can be the result of a higher rate of saving, a larger government sector financial deficit, an increase in the money supply, a change in preferences for liquidity, or something else. \input{text/m2chg.txt}
\vspace{2mm}

\normalsize \textbf{M2 Money Stock Growth}\\
\footnotesize{\textit{not seasonally adjusted, one-year percent change}}
\vspace{3.2cm}

\hspace*{4mm} \begin{tikzpicture}[overlay]
	\begin{axis}[\bbar{y}{0}, \dateaxisticks ytick={-5, 0, 5, 10, 15, 20, 25, 30}, 
		yticklabel style={text width=1.0em}, ymax=23.8,
		xticklabel={`\short{\year}}, clip=false, width=12.1cm, height=5.2cm]
	\rbars
	\thickline{green!80!blue}{date}{value}{data/m2.csv}
	\input{text/m2_node.txt}
	\end{axis}
\end{tikzpicture}

\footnotesize{Source: Federal Reserve} \hfill \tbllink{m2.csv}
\end{minipage}
\newpage
\begin{minipage}{0.76\textwidth}
\subsubsection*{Fed Asset Purchases}
\small During periods where the Fed funds rate is at or near zero the Fed has engaged in \textbf{large scale asset purchases} in an effort to further improve financial market conditions. These asset purchases show up on the Fed balance sheet, which is \href{https://www.federalreserve.gov/releases/h41/current/default.htm}{reported} weekly. 

\input{text/fed_assets1.txt}

\input{text/fed_assets2.txt}
\end{minipage}

\begin{minipage}{0.44\textwidth}
\normalsize \textbf{Federal Reserve Balance Sheet}\\
\footnotesize{\textit{total assets, trillion of USD}}\\
\hspace*{-2mm} \begin{tikzpicture}
	\begin{axis}[\bbar{y}{0}, \dateaxisticks ytick={0, 1, 2, 3, 4, 5, 6, 7, 8, 9}, 
		enlarge y limits={0.1}, ymin=0.7, ymax=8.4,
		xticklabel={`\short{\year}}, height=6.2cm, width=8.0cm, clip=false]
	\recbars
	\stdline{blue!80!black}{date}{RESPPMA_N.WW}{data/fed_assets.csv}
	\end{axis}
\end{tikzpicture}\\
\footnotesize{Source: Federal Reserve} \hfill \tbllink{fed_assets.csv}
\end{minipage}\hspace{7mm}
\begin{minipage}{0.27\textwidth}
\small \input{text/fed_assets3.txt}\input{text/fed_assets4.txt}
\end{minipage}
\vspace{8mm}

\begin{minipage}{0.76\textwidth}
\normalsize \textbf{Federal Reserve Assets}\\
\footnotesize{\textit{billions of US dollars}\\
\rowcolors{1}{}{black!5} \setlength{\tabcolsep}{3.1pt} \color{black!90}
{\renewcommand{\arraystretch}{1.54}
	\begin{tabular}{p{40mm} R{13.0mm} R{13.0mm} R{13.0mm} R{13.0mm} R{13.0mm}}
		\input data/fed_bal_sheet.tex \hline
	\end{tabular}
}}
		
\vspace{-2mm}
\footnotesize{Source: Federal Reserve}
\end{minipage}
\newpage
\begin{minipage}{0.76\textwidth}
\section*{Prices}
\hypertarget{pr}{}
\vspace{-1mm}
\small Changes in prices affect the amount of goods and services that can be purchased by a fixed income. When measuring changes in prices, researchers consider both the quantity that can be purchased by a unit of currency, and also changes in item quality. To understand the overall change in prices paid or charged by a certain group, such as consumers or manufacturers, researchers create a representative ``basket'' of the goods and services relevant to the group, and track the changes in the basket, and the price of the basket, over time. The end result of these methods is a price index.

Researchers can then use the price index to calculate the rate of inflation. Perhaps the most common way of measuring inflation is to calculate the 12-month percent change in the index values. This measures how prices in a given month compare to prices during the same month, one year prior. 
\vspace{1mm}

\normalsize \textbf{Price Growth, Various Measures}\\
\footnotesize{\textit{12-month percent change}}\\
\hspace*{-3mm} \rowcolors{1}{}{black!5} \setlength{\tabcolsep}{3.1pt} \color{black!90}
	{\renewcommand{\arraystretch}{1.6}
		\begin{tabular}{p{31mm} R{8.0mm} R{8.0mm} R{8.0mm} R{8.0mm} R{8.0mm} R{8.0mm} 
		   R{9.5mm} R{9.5mm}}
			 \input data/prices_12m.tex \hline
		\end{tabular}}\vspace{-1mm}
		
\footnotesize{Source: BLS, BEA, Federal Reserve Bank of Dallas}
\vspace{2mm}

\small The 12-month percent change in prices is affected not only by recent price changes, but also by price changes that happened over the past year. Results using the 12-month measure are smoothed by including a longer time interval.

While the chartbook uses less-volatile 12-month inflation rates in most cases, the \textbf{one-month rate} can be more useful for examining short-term trends, for example by eliminating the base effects from changes in prices a year ago. \input{text/cpi_monthly.txt}
\vspace{1mm}

\normalsize \textbf{CPI One-Month Change}\\
\footnotesize{\textit{monthly percent change, annualized}}\\
\hspace*{-3mm} \begin{tikzpicture}
	\begin{axis}[\bbar{y}{0}, date coordinates in=x, axis line style={draw=none},
		enlarge y limits={0.06}, enlarge x limits={0.02}, width=13.0cm, height=5.0cm,
		yticklabel style={text width=1.0em}, ymax=14, ymin=-4, 
		nodes near coords,
		nodes near coords style={/pgf/number format/.cd, fixed zerofill,
			precision=1, assume math mode}, xtick=data,
		xticklabels from table={\cpimo}{label}]
	\addplot[ybar, fill=blue, draw opacity=0] table [x=date, y=ALL_S, col sep=comma]{\cpimo};
	\end{axis}
\end{tikzpicture}\\
\footnotesize{Source: Bureau of Labor Statistics} \hfill \tbllink{cpi_monthly.csv}
\end{minipage}
\newpage
\hypertarget{prin}{}
\begin{minipage}{0.76\textwidth}
\subsection*{Consumer Price Index}
\small \input{text/cpi_main.txt}
\vspace{1mm}

\normalsize \textbf{Consumer Price Index}\\
\footnotesize{\textit{annual growth, percent, not seasonally adjusted}}\\
\vspace{28mm}

\noindent \hspace{2mm} \begin{tikzpicture}[overlay]
	\begin{axis}[\bbar{y}{0}, \dateaxisticks ytick={-2, 0, 2, 4, 6, 8}, 
		xticklabel={`\short{\year}}, clip=false, 
		legend style={at={(0.98, 1.14)}},
		height=5.2cm, width=12.4cm, enlarge y limits={0.12}]
	\rbars
	\stdline{blue!60!cyan}{date}{ALL}{data/cpi.csv}
	\stdline{gray}{date}{CORE}{data/cpi.csv}
	\input{text/cpi_node.txt}
	\legend{All-items, Core};
	\end{axis}
\end{tikzpicture}

\footnotesize{Source: Bureau of Labor Statistics} \hfill \tbllink{cpi.csv}
\vspace{2mm}

\small Recent changes in prices can be broad-based, that is derived from many prices changes at roughly the same rate, or narrow-based and driven by large changes in a subset of prices. Identifying each major spending category's contribution to overall inflation gives insight into whether inflation is broad-based and also into which groups of people face higher or lower rates of inflation.
\end{minipage}

\begin{minipage}{0.36\textwidth}
\small \input{text/cpicomp.txt}
\end{minipage} \hspace{6mm}
\begin{minipage}{0.35\textwidth}
\normalsize \textbf{Contribution of Major Categories}\\
\footnotesize{\textit{percentage points, not seasonally adjusted}}\\
\hspace*{-2mm} \begin{tikzpicture}
  	\begin{axis}[\barplotnogrid axis y line=left, \barylab{2.0cm}{1.5ex}
    	width=4.6cm, bar width=1.6ex, height=8.8cm,
    	enlarge y limits={abs=0.35cm}, enlarge x limits=0.16, \bbar{x}{0},
		yticklabels from table={\cpi}{name}, nodes near coords, 
		yticklabel style={font=\footnotesize, xshift=-4pt},
		nodes near coords style={/pgf/number format/.cd, fixed zerofill,
			precision=1, assume math mode},
		legend style={text=black!70, at={(0,1.07)}, anchor=north, legend columns=-1, 
				fill=none, draw=none,
		        /tikz/every even column/.append style={column sep=0.4cm}}]
  	\addplot[fill=cyan!85!white, draw=none] 
  		table [y expr=-\coordindex, x index=2] {\cpi};
  	\addplot[fill=blue!80!black, draw=none] 
  		table [y expr=-\coordindex, x index=1] {\cpi};
 	\legend{\input{text/cpi_mo2.txt}, \input{text/cpi_mo1.txt}}
  	\end{axis}
\end{tikzpicture}\\
\footnotesize{Source: Bureau of Labor Statistics} \hfill \tbllink{cpi_comp.csv}
\end{minipage}
\newpage
\begin{minipage}{0.76\textwidth}
\small The prices of some items are more volatile than others. Food and energy prices, for example, are sometimes separated from the rest of the CPI basket, which is referred to as the \textit{core}, because swings in food and energy prices are larger and more frequent. 

Core inflation includes core goods, core services other than shelter, and shelter. Core goods inflation was barely existent from 2013 through the start of the pandemic. Core goods prices are disproportionately affected by import prices and by changes in the quality of goods, for example from technological improvement. In contrast, domestic wage growth affects the prices for core services more than the other categories. Shelter prices are affected by housing supply and construction. 

\input{text/cpi_decomp.txt}
\vspace{1mm}

\normalsize \textbf{CPI Decomposition}\\
\footnotesize{\textit{contribution to one-year change in CPI-U, percentage points}}\\
\hspace*{-2mm} \begin{tikzpicture}
	\begin{axis}[\bbar{y}{0}, \tydateaxisticks ytick={-4, -2, 0, 2, 4, 6, 8}, 
		height=6.2cm, width=12.9cm, enlarge x limits={0.01}, 
		yticklabel style={text width=1.2em},
        xticklabel={`\short{\year}}, legend cell align={left},
		legend style={legend columns=2, at={(0.53, 0.995)}}]
	\rbars
	\ctsbar{blue!85!black}{date}{core_goods}{data/cpi_decomp.csv}{2.4pt}
	\ctsbar{green!60!black}{date}{core_services}{data/cpi_decomp.csv}{2.4pt}
	\ctsbar{cyan!50!white}{date}{shelter}{data/cpi_decomp.csv}{2.4pt}
	\ctsbar{orange!80!red}{date}{food_energy}{data/cpi_decomp.csv}{2.4pt}
	\legend{Core Goods, Core Services ex. Shelter, Shelter, Food \& Energy};
	\end{axis}
\end{tikzpicture}\\
\footnotesize{Source: Bureau of Labor Statistics, Author's Calculations} \hfill \tbllink{cpi_decomp.csv}
\vspace{4mm}

\subsubsection*{Relative Prices}
\small Some categories of spending are a small part of the overall price basket used to calculate the CPI, but have seen large cumulative increases in prices. Additionally, the price basket is based on averages across individuals, with some individuals dedicating a much larger share of spending to certain categories and other individuals have no expenses in certain categories. For example, day care costs are paid generally only for a few years of a child's life and only some households contain day-care-age children. But within those households, day care is a large share of overall spending. 

The following section and table cover one-year inflation rates for different categories of goods and services, including some smaller categories. The table also shows the cumulative change in prices for each category since February 2020, the last month of data before the COVID-19 pandemic shutdown in the US. Additionally, the weight that a category has in the overall index--the category's share of the basket of goods and services used to calculate the CPI--is included as the last column in the table. This weight comes from each category's share of overall consumer spending during the most recent reference period, and is updated by changes in prices since the reference period. 
\end{minipage}
\newpage
\begin{minipage}{0.76\textwidth}
\small \input{text/cpicomp2.txt}
\vspace{1.5mm}

\normalsize \textbf{Selected CPI Categories}\\
\footnotesize{\textit{one-year percent change}\\
\hspace*{-3mm} \rowcolors{1}{}{black!5} \setlength{\tabcolsep}{2.6pt} \color{black!90}
		{\renewcommand{\arraystretch}{1.5}
\begin{tabular}{p{42mm} R{8.0mm} R{8.0mm} R{8.0mm} R{8.0mm}
		 		 R{8.0mm} R{10.0mm} R{12.0mm}} % 
			 \input data/cpi_comp.tex \hline
		\end{tabular}}}
\vspace{-2mm}		
		
\footnotesize{Source: Bureau of Labor Statistics}
\end{minipage}
\newpage
\hypertarget{prie}{}
\begin{minipage}{0.76\textwidth}
\subsection*{Inflation Expectations}
\vspace{-1mm}

\small Researchers gain insight on expected changes in prices through regular \textbf{surveys of consumers} and through \textbf{market data}. One market-based measure is known as the \href{https://fred.stlouisfed.org/series/T5YIE}{inflation breakeven} and is calculated as the difference between the yield on a nominal treasury bond and the yield on a treasury inflation-protected bond of the same maturity. This difference represents the amount of inflation markets have priced in, on average, for the maturity of the bond.
\end{minipage}

\begin{minipage}{0.38\textwidth}
\normalsize \textbf{5-year Expected Average Inflation}\\
\footnotesize{\textit{expected average annual rate, percent}}\\
\hspace*{-2mm} \begin{tikzpicture}
	\begin{axis}[\ltdateaxisticks ytick={0, 1, 2, 3}, 
		enlarge y limits={0.1}, ymin=0.28, \bbar{y}{0}, clip=false,
		xticklabel={`\short{\year}}, height=6.6cm, width=6.6cm]
	\rbar
	\thinline{blue!70!black}{date}{VALUE}{data/infbreak.csv}
	\stdline{violet!60!magenta}{date}{PX5_MD}{data/infumich.csv}
	\input{text/infbreak_node.txt}
	\input{text/infumich_node.txt}
	\absnode{{2017-01-01}}{1.5}{\footnotesize \color{blue!70!black}\textbf{Market}}
	\absnode{{2016-10-01}}{2.9}{\footnotesize \color{violet!60!magenta}\textbf{Survey}}
	\end{axis}
\end{tikzpicture}\\
\footnotesize{Source: FRED; U. of Michigan} \hfill \tbllink{infbreak.csv}
\end{minipage} \hspace{6mm} 
\begin{minipage}{0.33\textwidth}
\small \input{text/inf_exp_cons.txt}

\input{text/inf_exp_mkts.txt}
\end{minipage}
\vspace{1mm}

\begin{minipage}{0.76\textwidth}

\small Both survey- and market-based estimates of expected inflation offer techniques for distinguishing between near-term inflation and expected medium-term inflation. The survey-based measure asks about inflation over the next year. \input{text/inf_exp_surv_st.txt}
\vspace{1mm}

\normalsize \textbf{Survey of Expected Inflation}\\
\footnotesize{\textit{annual rate, percent}}\\
\vspace{22mm}

\hspace{4mm} \begin{tikzpicture}[overlay]
	\begin{axis}[\bbar{y}{0}, \dateaxisticks ytick={0, 1, 2, 3, 4, 5, 6}, 
		xticklabel={`\short{\year}}, clip=false, 
		legend style={at={(0.98, 1.16)}}, ymin=0.24,
		yticklabel style={text width=1.0em},
		width=12.9cm, height=4.5cm]
	\rbars
	\stdline{violet!60!magenta}{date}{PX5_MD}{data/infumichlt.csv}
	\stdline{green}{date}{PX_MD}{data/infumichlt.csv}
	\legend{Next 5 years, Next year};
	\end{axis}
\end{tikzpicture}
\vspace{1mm}

\footnotesize{Source: University of Michigan} \hfill \tbllink{infumichlt.csv}
\end{minipage}
\vspace{1mm}

\begin{minipage}{0.47\textwidth}
\normalsize \textbf{Market Expected Inflation}\\
\footnotesize{\textit{annual rate, percent}}\\
\hspace*{-2mm} \begin{tikzpicture}
	\begin{axis}[\ltdateaxisticks ytick={0, 1, 2, 3, 4}, 
		enlarge y limits={0.1}, ymin=0.24, \bbar{y}{0}, clip=false,
		xticklabel={`\short{\year}}, height=4.5cm, width=8.4cm]
		\rbar
		\thinline{blue!70!black}{date}{5_year_breakeven}{data/infbreak_comb.csv}
		\thinline{cyan}{date}{5_year_5_year_forward}{data/infbreak_comb.csv}
		\absnode{{2015-05-01}}{0.8}{\footnotesize \color{blue!70!black}\textbf{Next 5 years}}
		\absnode{{2016-06-01}}{2.85}{\footnotesize \color{cyan}\textbf{Following 5 years}}
	\end{axis}
\end{tikzpicture}\\
\footnotesize{Source: FRED} \hfill \tbllink{infbreak_comb.csv}
\end{minipage} \hspace{6mm} 
\begin{minipage}{0.24\textwidth}
\small The market-based measure can be used to calculate expected inflation over the five years starting five years from now. \input{text/inf_exp_mkts_55.txt}
\end{minipage}
\newpage
\begin{minipage}{0.76\textwidth}
\subsection*{PCE Price Index}
\small The Personal Consumption Expenditure (PCE) price index from the Bureau of Economic Analysis \href{https://www.bea.gov/data/personal-consumption-expenditures-price-index}{captures} both changes in the price of goods and services as well as monthly changes in consumer behavior. The index is additionally updated over time to the latest methodology. 

\input{text/pce_inf_basic.txt}
\vspace{1mm}

\normalsize \textbf{Personal Consumption Expenditure Price Index}\\
\footnotesize{\textit{annual growth, percent}}\\
\vspace{34mm}

\hspace{3mm} \begin{tikzpicture}[overlay]
	\begin{axis}[\bbar{y}{0}, \dateaxisticks ytick={-2, 0, 2, 4, 6}, 
		xticklabel={`\short{\year}}, clip=false, ymin=-2, ymax=6.5, 
		enlarge y limits={0.01}, width=12.4cm, 
		legend style={at={(1.0, 1.12)}, fill=white},
		height=5.6cm]
	\rbars
	\stdline{blue!60!black}{date}{CORE}{data/pce_pi.csv}
	\stdline{orange!80!yellow}{date}{PCE}{data/pce_pi.csv}
	\input{text/pce_pi_node.txt}
	\legend{Core, All-items};
	\end{axis}
\end{tikzpicture}

\footnotesize{Source: Bureau of Economic Analysis} \hfill \tbllink{pce_pi.csv}
\vspace{5mm}

\small The Federal Reserve Bank of Dallas \href{https://www.dallasfed.org/research/pce}{publish} a variation of the PCE price index called the trimmed-mean index. The most volatile categories in the current month's index are removed, or \textit{trimmed}, to smooth the data. As a result, the most extreme categories, which vary from month-to-month, do not affect inflation rates calculated using the trimmed-mean index. 

\input{text/pce_tm_basic.txt}
\vspace{1mm}

\normalsize \textbf{Personal Consumption Expenditure Price Index, Trimmed Mean}\\
\footnotesize{\textit{highest and lowest categories excluded, annual growth, percent}}\\
\hspace{-2mm} \begin{tikzpicture}
	\begin{axis}[\bbar{y}{0}, \dateaxisticks ytick={0, 1, 2, 3, 4}, 
		xticklabel={`\short{\year}}, clip=false, ymin=0.15, width=12.2cm,
		height=5.0cm]
	\rbars
	\stdline{violet!60!magenta}{date}{12-month}{data/pce_tm12.csv}
	\input{text/pce_tm_node.txt}
	\end{axis}
\end{tikzpicture}\\
\footnotesize{Source: Federal Reserve Bank of Dallas} \hfill \tbllink{pce_tm12.csv}
\end{minipage} 
\newpage
\begin{minipage}{0.76\textwidth}
\subsection*{Producer Prices}

\small \input{text/ppi_main.txt}
\vspace{1mm}

\normalsize \textbf{Producer Price Index}\\
\footnotesize{\textit{all commodities, annual growth, percent}}\\
\hspace*{-2mm} \begin{tikzpicture}
	\begin{axis}[\dateaxisticks ytick={-20, -10,  0, 10, 20}, 
		enlarge y limits={0.05}, \bbar{y}{0}, clip=false,
		yticklabel style={text width=1.6em},
		xticklabel={`\short{\year}}, height=6.0cm, width=12.1cm]
    	\rbars
		\stdline{green!80!blue}{date}{PPIACO}{data/ppi.csv}
		\input{text/ppi_node.txt}
	\end{axis}
\end{tikzpicture}\\
\footnotesize{Source: Bureau of Labor Statistics} \hfill \tbllink{ppi.csv}
\vspace{4mm}

\subsection*{Import and Export Prices}
\small \input{text/mxpi.txt}
\vspace{1mm}

\normalsize \textbf{Import/Export Price Index}\\
\footnotesize{\textit{annual growth, percent}}\\
\vspace{30mm}

\hspace{4mm} \begin{tikzpicture}[overlay]
	\begin{axis}[\bbar{y}{0}, \dateaxisticks ytick={-20, -10,  0, 10, 20}, 
		xticklabel={`\short{\year}}, clip=false, 
		legend style={at={(0.98, 1.14)}},
		yticklabel style={text width=1.5em},
		height=5.4cm, width=12.1cm, enlarge y limits={0.08}]
	\rbars
	\stdline{cyan!85!yellow}{date}{Imports}{data/mxpi.csv}
	\stdline{red!25!orange}{date}{Exports}{data/mxpi.csv}
	\input{text/mxpi_nodes.txt}
	\legend{Imports, Exports};
	\end{axis}
\end{tikzpicture}
\vspace{1mm}

\footnotesize{Source: Bureau of Labor Statistics} \hfill \tbllink{mxpi.csv}
\end{minipage}
\newpage
\begin{minipage}{0.76\textwidth}
\subsection*{Commodity Prices}
\hypertarget{prco}{}
\small \input{text/wti.txt}
\vspace{1mm}

\normalsize \textbf{Oil Price}\\
\footnotesize{\textit{USD, west Texas intermediate crude, monthly average and latest value}}\\
\hspace*{-2mm} \begin{tikzpicture}
	\begin{axis}[\bbar{y}{0}, \dateaxisticks ytick={0, 25, 50, 75, 100, 125, 150}, 
		yticklabel style={text width=1.5em}, width=11.7cm, height=5.6cm,
		xticklabel={`\short{\year}}, ymin=7, clip=false]
	\rbars
	\stdline{red!80!purple}{date}{VALUE}{data/wti.csv}
	\input{text/oil_node.txt}
	\end{axis}
\end{tikzpicture}\\
\footnotesize{Source: Energy Information Administration} \hfill \tbllink{wti.csv}
\vspace{10mm}

\small \input{text/gold.txt}
\vspace{1mm}

\normalsize \textbf{Gold Price}\\
\footnotesize{\textit{USD, monthly average and latest value}}\\
\hspace*{-2mm} \begin{tikzpicture}
	\begin{axis}[\bbar{y}{0}, \dateaxisticks ytick={0, 500, 1000, 1500, 2000}, 
		yticklabel style={text width=2.0em}, clip=false, yticklabel style={text width=2.4em}, 
		xticklabel={`\short{\year}}, ymin=95, width=11.6cm, height=5.4cm]
	\rbars
	\stdline{orange!50!yellow}{date}{AM}{data/gold.csv}
	\input{text/gold_node.txt}
	\end{axis}
\end{tikzpicture}\\
\footnotesize{Source: London Bullion Market Association} \hfill \tbllink{gold.csv}
\end{minipage}
\newpage
\begin{minipage}{0.76\textwidth}
\subsection*{Acknowledgments}
\small Gabriel Mathy, Iordan Koulov, Lara Merling, Kevin Cashman, Rebecca Watts, Dean Baker, Matt Bruenig, Eileen Appelbaum, John Schmitt, Mark Weisbrot, Ben Zipperer, Alberto Rodelgo, Yevgeniya Korniyenko, Magali Pinat, Robert Blecker, Teasri Thiruvadanthai, Rainer K\"ohler, Gersenda Varisco, Venkat Josyula, Tom Augspurger, Claudia Sahm, Mike Sieferling, Andrew Paciorek, David Dorn, Skanda Amarnath, Ernie Tedeschi, Brad Setser, and Vikas Sharma.
\end{minipage}
\end{document}
